
\clearpage
\setcounter{section}{4}
\refstepcounter{section}
\section*{Chapter~\thesection\quad Bounds on connected correlations on the torus: clauses (iii-a) and (iii-c)\addcontentsline{toc}{section}{\protect\numberline{\thesection}Bounds on connected correlations on the torus: clauses (iii-a) and (iii-c)}\label{ch:5}}
%% ---- begin inlined file: chapters/c05 ----
\begin{definition}[Notation for the correlation theorem]\label{5.1:not-correlation-notation}
The following notation is fixed for the correlation theorem and its proof. The letters $(T)$,
$(P)$, $(L)$ name the three sites of one renormalisation step treated in the statement on the
three sites: the output site of Ba{\l}aban's large-field operation $\mathbf{T}$, the site of
the preparatory stages, and the localisation site.

(a) \emph{Sources, shifts, slots.} The $z_i$ are the source parameters and $w=\sum_i z_i f_i$ is the
complex source; $\zeta$ is the constant bare shift and $\zeta_k$ the running shift; the $\xi_j$ are
independent slot variables, and $s:=1-\xi g^2$; $\mathfrak{q}$ is the complex coupling slot;
$x_k:=g_k^{-2}$. Ba{\l}aban's decoupling parameters are written $\mathsf{s}(\Delta)$ (his letters
$\sigma(\Delta)$, $\tau(Y)$, $t(Y)$ occur only inside the quotation of the one formula of his in
which they are printed). Ba{\l}aban's point
$z$ is written $y$, and his functions $\zeta_{\square},\zeta_0,\zeta_1$ of
\cite[(3.4)--(3.7)]{B-CMP109} are written $\hat\chi_{\square},\hat\chi_0,\hat\chi_1$.

(b) \emph{The numbers $\beta$.} Ba{\l}aban's single small positive constant $\beta$ of
\cite{B-CMP119} (for which $1/4$ is given there as an admissible value; it is the coefficient
letter of the spaces \cite[(2.34)--(2.39)]{B-CMP119}, \cite[(1.24), (1.55)--(1.70)]{B-CMP122a},
\cite[(1.66)--(1.69)]{B-CMP122b}, and a rate letter of \cite{B-CMP119}) is instantiated at
\begin{equation*}
\beta_{\mathrm{pr}}:=1/5
\end{equation*}
in every object of Ba{\l}aban's taken by statement or as a source-free auxiliary object. At
$\beta=1/4$ the coefficient of \cite[(1.64)]{B-CMP122a} has infimum $1-4\beta=0$ (the complex radius
of the chain space then shrinks as $2^{-\lfloor N/2\rfloor}$ with the number $N$ of preparatory
stages), whereas at $\beta_{\mathrm{pr}}$ it is at least $1/5$ for all $N$, by the lemma on
Ba{\l}aban's constant $\beta$. The number
\begin{equation*}
\beta_{s}:=1/4
\end{equation*}
is the correlation theorem's own bookkeeping number for the radii and rates of the families it
proves anew (the radii $\hat\alpha_k$, the ratios $w_{k,j}$, the constant $C_w$, the fresh rate
$(1+4\beta_s)\kappa$, the collar proposition's output rate $(1+2\beta_s)\kappa$ at $(L)$); it is
not Ba{\l}aban's $\beta$ and is fixed before $L$. The three places where $\beta_{\mathrm{pr}}$ and
$\beta_s$ meet are treated in the interface lemma for the two numbers $\beta$. The
unindexed $\hat\alpha\in\,]0,1[$ is Ba{\l}aban's $1-\beta$ of \cite[(3.67)]{B-CMP119}, while the
indexed $\hat\alpha_k$ are radii. $\beta[\cdot]$ is the extraction functional, with constant
$C_\beta$; $t:=L^{j-n}$; $\mathfrak{r}_j,\mathfrak{r}_{\circ},\mathfrak{t}_j$ are the radii and the
ratio of the radii lemma; $r_{\mathrm{pr}}$ is the exponent in Ba{\l}aban's lower bound
$R_0\ge(\log g_0^{-2})^{r}$ of \cite{B-CMP119}; $\mathfrak{e}_\mu(x)$ is the comb-gauge remainder;
$C_{(\sigma)},\Gamma_{(\sigma)}$ are the operators $C^{(k)}(Z_0,\sigma)$, $\Gamma_k(Z_0,\sigma)$;
\begin{equation*}
\mathsf{n}_j:=\sup_\zeta\|Q^{(j)}_\zeta\|^{\flat},\qquad
\mathsf{q}_{k+1}:=\beta[Q^{(k+1)}_\zeta];
\end{equation*}
$q$ is the exponent of \cite[(2.28)]{B-CMP119} and $\bar q$ the power of $M$ in
\cite[(1.36)]{B-CMP116} (printed there as $q$); $A(c,U)$ is the action and $A_0,A_1$ are
Ba{\l}aban's constants; $\upsilon$ is the shrink factor of the Cauchy circles of the last lattice.

(c) \emph{The decay fraction $\delta$.} In \cite{B-CMP116} the number $\delta$ of \cite[(1.32)]{B-CMP116},
$0<\delta<1$, is carried free through the estimates up to \cite[(2.41)]{B-CMP116} under lower
bounds only (``$\delta\kappa$ sufficiently large'', a lower bound of the form
$(\cdot)\kappa_1\ge(1-2\delta)\kappa$ on a fixed multiple of $\kappa_1$,
$\tfrac18(\kappa_1-1)\ge(1-3\delta)\kappa$), and is fixed only after \cite[(2.41)]{B-CMP116} by the
assumption $(1-10\delta)\tfrac12L=1$. The correlation theorem does not make that assumption:
$\delta$ is any fixed number with
\begin{equation}\label{5.1:eq-correlation-notation-c}
0<\delta\le\delta_J(L):=\frac{1}{10}\Bigl(1-\frac{2(1+4\beta_s)}{L}\Bigr)
=\frac{1}{10}\Bigl(1-\frac{4}{L}\Bigr),
\qquad
\mathsf{r}(\delta):=(1-10\delta)\tfrac12L,
\end{equation}
chosen once, immediately after $L$, as a function of $L$ only. For $m=1,\dots,10$,
\begin{equation}\label{5.1:eq-correlation-notation-1}
\bar\kappa_m:=(1-m\delta)\tfrac12L\kappa
\end{equation}
are the rates of \cite[(2.35)--(2.41)]{B-CMP116} at level $k+1$. Since $\mathsf{r}$ is decreasing
in $\delta$ and $1-10\delta_J(L)=4/L$, every admissible $\delta$ gives
\begin{equation*}
\mathsf{r}(\delta)\ge\mathsf{r}(\delta_J(L))=\frac4L\cdot\frac L2=2=1+4\beta_s;
\end{equation*}
as $\bar\kappa_{10}=\mathsf{r}(\delta)\kappa$, as $\bar\kappa_m$ is decreasing in $m$ (because
$\delta>0$), and as $\bar\kappa_7\le\tfrac12L\kappa$ (because $1-7\delta\le1$), this gives
\begin{equation}\label{5.1:eq-correlation-notation-2}
2\kappa\le\bar\kappa_{10}\le\bar\kappa_9\le\bar\kappa_7\le\tfrac12L\kappa.
\end{equation}
By the lemma on the window for $\delta$, \cite[(2.41)]{B-CMP116} then reads, with $C_3$,
$\varepsilon_1$ the constants of \cite{B-CMP116},
\begin{equation}\label{5.1:eq-correlation-notation-3}
|E^{(k+1)}(X)|\le O(1)C_3\varepsilon_1e^{-\mathsf{r}(\delta)\kappa d_{k+1}(X)},
\qquad \mathsf{r}(\delta)\kappa\ge(1+4\beta_s)\kappa,
\end{equation}
the output rate at $(T)$ being stated as the inequality
$\mathsf{r}(\delta)\kappa\ge(1+4\beta_s)\kappa$, never as an equality. The window in
\eqref{5.1:eq-correlation-notation-c} is non-empty if and only if $L>4$, a one-sided lower
condition on $L$. The number $\delta_\star$ used along the route of the uniqueness theorem
satisfies $\delta_\star\le\delta_J(L)$ and is therefore an admissible value, so that a single
$\delta$ serves both the correlation theorem and the uniqueness theorem.

(d) \emph{Logarithms and functions of the running coupling.} The logarithms are written
\begin{equation*}
\mathsf{T}:=\log g^{-2},\quad \mathsf{T}_k:=\log g_k^{-2},\quad \mathsf{T}_\gamma:=\log\gamma^{-2},
\quad \check{\mathsf{T}}_k:=\log\min_{i\le k}x_i,
\end{equation*}
and $\mathsf{T}_*$ denotes a threshold. $R(x)$ is the least power of $L$ not smaller than
$(\log x)^{r_{\mathrm{pr}}}$ and $\check R_k:=R(\min_{i\le k}x_i)\ge\check{\mathsf{T}}_k^{\,r_{\mathrm{pr}}}$;
$R_{\min}:=\check R_K=\min_k\check R_k$. For $k\le K$ one has $x_i>g^{-2}$ for $i\le k$, by the
standing condition on the run of couplings in the setting of the correlation theorem, so that
$\check{\mathsf T}_k\ge\mathsf T$, and
$\mathsf T\ge\mathsf T_\gamma$. By \eqref{5.1:eq-correlation-notation-d} below, with $\lambda$ the
rate constant of the reference recursion named in (g),
$x_i-x_k\ge\lambda(k-i)-2c_2\ge-2c_2$ for $i<k$, so $\min_{i\le k}x_i\ge x_k-2c_2$; together with
$x_k\ge4c_2$, a condition of the list defining $\gamma_J$, this gives
$\min_{i\le k}x_i\ge\tfrac12x_k$, hence
\begin{equation*}
\mathsf T_k\le\check{\mathsf T}_k+\log2\le\check{\mathsf T}_k+1.
\end{equation*}
$N(g)$ is the number of preparatory stages at coupling $g$;
with
\begin{equation*}
C_N:=2^{r_{\mathrm{pr}}}L,
\end{equation*}
fixed with $(L,r_{\mathrm{pr}})$, one has $N(g_k)\le C_N\mathsf{T}_k^{r_{\mathrm{pr}}}$, since
$R(x)<L\max\{1,(\log x)^{r_{\mathrm{pr}}}\}$ and Ba{\l}aban's $N$ is bounded by $R$ at the scale of
the chain \cite{B-CMP122b} (it is assumed there that $N\le R_k$). $N_0(g)$ is the number of zones
(as in \cite{B-CMP122a}); by the zone-count lemma,
$N_0(g)\le4+2r_{\mathrm{pr}}\log_L\log g^{-2}$, and
\begin{equation}\label{5.1:eq-correlation-notation-4}
\mathsf{N}_0(\mathsf{T}):=4+2r_{\mathrm{pr}}\log_L(\mathsf{T}+1).
\end{equation}
For $n\le k'\le K$, \eqref{5.1:eq-correlation-notation-d} gives $x_n\ge x_{k'}-2c_2$, and with
$x_{k'}\ge4c_2$ this gives $x_{k'}\le2x_n\le e\,x_n$, that is $\log x_{k'}\le\mathsf T_n+1$; hence
$\mathsf{N}_0(\mathsf{T}_n)\ge4+2r_{\mathrm{pr}}\log_L\log g_{k'}^{-2}\ge N_0(g_{k'})$ for every
$k'\in[n,K]$.
$\mathsf{K}_{\mathrm{val}}(k)$ and $\mathsf{K}_{\mathrm{cell}}(k,n)$ are the value constants of the
collar; by the collar proposition and the modified volume bound,
\begin{equation}\label{5.1:eq-correlation-notation-5}
\mathsf{K}_{\mathrm{val}}(k)\le\mathsf{K}^0+N(g_k)C_0^{\sharp},\qquad
\mathsf{K}_{\mathrm{cell}}(k,n)\le C_{\sharp}\bigl[\mathsf{K}_{\mathrm{cell}}^0
+3(\mathsf{N}_0(\mathsf{T}_n)+3)C_0^{\sharp}\bigr]
\end{equation}
at a cell of zone $n\le k$, where $C_0^{\sharp}:=3c_G+1$ is the constant of the stage norms and
$\mathsf{K}^0,\mathsf{K}_{\mathrm{cell}}^0$
are $g$-free numbers depending on $(d,L,M,\kappa,\delta)$, fixed after $M$; the second bound counts
the chain of the present step only, each term charged once at a cell of its own stage region.
$\alpha^{\sharp}(k)$ is the smallness of the strips, and
$m_{\mathrm{lo}}(k)$ and $m_{\sharp}(k)=m_{\mathrm{hi}}(k):=\lfloor\tfrac12(\check R_k-1)\rfloor$ are
the lower and upper ends of the window of collar widths. The functions $N(g_k)$, $N_0$, $\check R_k$,
$\mathsf{K}_{\mathrm{val}}$, $\mathsf{K}_{\mathrm{cell}}$, $\varepsilon_{\mathrm{br}}$ of (f), $\alpha^\sharp$, $m_{\mathrm{lo}}$,
$m_\sharp$ and $\gamma_0^{(k)}$ of (p) are functions of the running
coupling (and, for $\mathsf{K}_{\mathrm{cell}}$, of the zone level $n\le k$), evaluated after $\gamma$
is fixed; none occurs in the order of choice before $\gamma$ or in the list of conditions defining
the threshold $\gamma_J$, and none is
asserted to be $g$-free. They are controlled only through conditions of the form
$\sup_{\mathsf{T}\ge\mathsf{T}_\gamma}(\cdot)\le c$, $c$ a number fixed before $\gamma$ --- the
strip, bracket and volume conditions of the list defining $\gamma_J$ --- and these conditions
themselves contain no such function.

(e) \emph{The frozen pile.} For a live arrest $\mathfrak{a}$ at step $k_a$, with component
$X_{\mathfrak{a}}$, live piece $W_{\mathfrak{a}}:=X_{\mathfrak{a}}\cap\Omega^c_{k_a}$, $n_a$ performed
stages and $n^+_a$ the stage attached to the arrest in its definition ($n^+_a=n_a+1$ on the
change-of-variables
branch), $\mathsf{P}^{\mathrm{fr}}(\mathfrak{a})$
is the sum of the supremum sizes of its frozen terms $\mathfrak{F}(\mathfrak{a})$ --- the chain terms
$\mathbb{C}^{(i)}_{k_a}(X,\mathcal{S})$, $1\le i\le n_a$, each meeting its own band, and the
change-of-variables terms of stage $n^+_a$ --- each charged once at an anchor cube of the scale of
its own stage; by the modified volume bound it satisfies
\begin{equation}\label{5.1:eq-correlation-notation-6}
\mathsf{P}^{\mathrm{fr}}(\mathfrak{a})\le c_{\mathsf{P}}\sum_{n\le n^+_a}\check R^{\,d+1}_{j_n+1}
\sum_X\bigl(d'_{j_n+1}(X)+1\bigr),
\end{equation}
$j_n$ the level integrated at stage $n$, $X$ ranging over the components of
$Z_{j_n+1}=\Lambda^c_{j_n+1}\cap Z$ \cite{B-CMP122a} contained in $X_{\mathfrak{a}}$, $d'$ Ba{\l}aban's
linear size in $M\check R$-cubes \cite{B-CMP122b}, and
\begin{equation*}
c_{\mathsf{P}}:=2^{d+3+r_{\mathrm{pr}}}3^dC_NC_0^{\sharp}+27\mathsf{L}_0^6L^d\mathsf{c}_d\mathsf{K}^{\flat},
\qquad \mathsf{c}_d:=2\cdot3^d,
\end{equation*}
a number fixed after $M$; here $\mathsf{L}_0$ is Ba{\l}aban's auxiliary number of
\cite[(1.24)]{B-CMP122a}, chosen after $L$ inside Ba{\l}aban's window
$2\le\mathsf{L}_0<\tfrac12L$ \cite[p.~182]{B-CMP122a} with $\mathsf{L}_0^2\le\tfrac13L$, and with
in addition $\mathsf{L}_0\ge3$; and $\mathsf{K}^{\flat}$ is a
constant of the collar proposition fixed after $M$. $\mathsf{P}^{\mathrm{fr}}(\mathfrak{a})$ depends
on $g_{k_a}$ through $N(g_{k_a})$ only.
Like the functions of (d) it is evaluated on the run and occurs in no member of the list defining
$\gamma_J$; unlike them it is dominated by no condition on $\gamma$: it enters directly in the third
term $O(1)M^dR_j^{d+1}d'_j(Z_j)$ of Ba{\l}aban's volume bound, in the form of the modified volume
bound.

(f) \emph{The brackets' activity at $(L)$.} By the collar proposition, $\varepsilon_{\mathrm{br}}(k)$
is bounded by
\begin{equation}\label{5.1:eq-correlation-notation-7}
\begin{split}
\varepsilon_{\mathrm{br}}(k)\le{}&\bigl[C_\sharp O(1)+8C_A'C_NC_0^\sharp(\mathsf{T}_k+1)^{r_{\mathrm{pr}}}
+8C_A'\mathsf{K}^{\mathrm{arr}}\bigr]\theta'(g_k)\\
&+\varepsilon^{\mathrm{rel}}(k)+\varepsilon^{\mathrm{int}}(k)+\varepsilon^{\mathrm{ext}}(k),
\qquad
\varepsilon^{\mathrm{ext}}(k):=\varepsilon^{\mathrm{ext},\sigma}(k)+\varepsilon^{\mathrm{ext},\mathsf{t}}(k).
\end{split}
\end{equation}
Its six members are: Ba{\l}aban's own brackets, the present chain's, and the exterior arrested
left-overs (the three members carrying $\theta'(g_k)$; they consist of the pieces $v_p$ of the
amplitude lemma, which carry $\mathsf t$, of the plain terms of two classes of the standing domain
convention, called here the first and the second amplitude class, read with the second, third and
fourth links of their four-link expansion empty); the relative pieces
$\varepsilon^{\mathrm{rel}}$ of the wrapped terms; the pieces $\varepsilon^{\mathrm{int}}$ of the
plain interior first-part terms whose filing domain leaves $Z$; and $\varepsilon^{\mathrm{ext}}$,
where $\varepsilon^{\mathrm{ext},\sigma}$ collects the $\mathsf{t}$-free non-empty continuations of
the plain first-part terms with $K_v\setminus Z\neq\emptyset$, and
$\varepsilon^{\mathrm{ext},\mathsf{t}}$ collects the non-empty continuations of the
$\mathsf{t}$-carrying terms $v_p$ of the plain terms of the first and the second amplitude class;
in terms of the three members carrying $\theta'(g_k)$ in \eqref{5.1:eq-correlation-notation-7},
\begin{equation}\label{5.1:eq-correlation-notation-8}
\begin{split}
\varepsilon^{\mathrm{ext},\mathsf{t}}(k):={}&(1+C_\Sigma C_e)^3eS_d\times\bigl(\text{the three
members carrying }\theta'(g_k),\\
&\text{read at the weight index }a+7,\text{ the first with the constant }
\mathsf{C}_1^{\mathrm{ex}}\bigr)\\
&-\bigl(\text{the three members carrying }\theta'(g_k)\bigr),
\end{split}
\end{equation}
where $a$ is the weight index of the plain terms, $C_\Sigma$, $C_e$, $S_d$ are the summation
constants of the exterior lemma, and
$\mathsf{C}_1^{\mathrm{ex}}:=\mathsf{C}^{\Sigma,4}+\mathsf{C}^{\Sigma}_{a+7}$, the sum of the constant
of the summation of the terms of the fourth kind and of the constant of the pointed summation at
weight index $a+7$. The six members are exhaustive: every piece of every term at $(L)$ lies in
exactly one of them or among the filed values $\mathfrak{v}^{\mathrm{val}}$, and the source-free
localised members of $\mathbf{C}_k$ lie among $\mathfrak{v}^{\mathbf{C}}$. By
\eqref{5.1:eq-correlation-notation-8}, the sum of the three members carrying $\theta'(g_k)$ and of
$\varepsilon^{\mathrm{ext},\mathsf{t}}(k)$ is $(1+C_\Sigma C_e)^3eS_d$ times these three members
read at the weight index $a+7$, the first with the constant $\mathsf{C}_1^{\mathrm{ex}}$, which
is at most $\tfrac12$ under the bracket condition of the list defining $\gamma_J$; by the interior
lemma, $\varepsilon^{\mathrm{int}}(k)+\varepsilon^{\mathrm{ext},\sigma}(k)\le\tfrac12$ under the
interior condition of that list; and $\varepsilon^{\mathrm{rel}}(k)\le1$ under its relative
condition at $(L)$. Adding the three groups,
\begin{equation}\label{5.1:eq-correlation-notation-9}
\varepsilon_{\mathrm{br}}(k)\le\tfrac12+1+\tfrac12=2.
\end{equation}

(g) \emph{The flow constants.} The reference recursion, with its rate constant $\lambda$, gives the
two-sided bound in the first line
of the display below, for $i<j$; the flow constants of the envelope estimates at $(L)$ are read in
this chapter as in its second line:
\begin{equation}\label{5.1:eq-correlation-notation-d}
\begin{gathered}
\lambda(j-i)-2c_2\le x_i-x_j\le3\lambda(j-i)+2c_2\qquad(i<j),\\
c_5:=2c_2,\qquad\lambda^{\sharp}:=\lambda,\qquad B^{\sharp}:=b_+:=3\lambda+2c_2.
\end{gathered}
\end{equation}
With these values, the first inequality at $i=k'$, $j=k$ is the lower flow bound
$x_{k'}\ge x_k+\lambda^\sharp(k-k')-c_5$ for $k'\le k$ (the case $k'=k$ being trivial), and the
second at $i=m$, $j=k'$ gives the upper flow bound $x_m\le x_{k'}+B^\sharp(k'-m)$ for
$m\le k'\le k$, because $3\lambda(k'-m)+2c_2\le B^\sharp(k'-m)$ when $k'-m\ge1$, the case $m=k'$
being trivial. Inside the induction, at step $k$, these bounds are used for pairs of indices not
exceeding $k$ only, that is, as the first clause of the induction hypothesis at step $k$; inside
conditions on $\gamma$ the constant $B^\sharp$ is read at the value $b_+$, a number fixed before
$\gamma$, in the same way as the bound $B^\sharp\le\bar B$ of the setting of the correlation
theorem is read at its right member $\bar B$. The exponent $r$ of the envelope estimates is
$r_{\mathrm{pr}}\ge2$; their $N_k$ is
$N(g_k)$, majorised there by $L(\log x_k)^{r_{\mathrm{pr}}}$ and in this chapter by
$C_N\mathsf{T}_k^{r_{\mathrm{pr}}}$, each condition being read with its own majorant. Since
$\min_{i\le k}x_i\ge x_k-2c_2=x_k-c_5$ (see (d)),
\begin{equation*}
(\log(x_k-c_5))^{r_{\mathrm{pr}}}\le\check{\mathsf{T}}_k^{\,r_{\mathrm{pr}}}\le\check R_k.
\end{equation*}
With
\begin{equation*}
E(x):=\exp\bigl(-[(\log(x-c_5))^{r_{\mathrm{pr}}}-2]\bigr),\qquad
\mu(x):=\varepsilon_\star(x-c_5)^{-2},\qquad \mu_k:=\mu(x_k),
\end{equation*}
this gives, for every $u>0$,
\begin{equation}\label{5.1:eq-correlation-notation-10}
e^{-u\check R_k}\le e^{-u(\log(x_k-c_5))^{r_{\mathrm{pr}}}}=E(x_k)^{u}e^{-2u}.
\end{equation}
The envelope letters $\alpha_*(x)$, $\varepsilon_\star$, $\varepsilon_k(s)$, $K_u$, $B_H$, $c_4$,
$\varpi_{\mathrm{fr}}$, $C_{\mathrm{fr}}$ are
those of the envelope estimates at the localisation site, taken by name; $\mathsf{e}(s)$ is the
margin at the law site $s$, one of $\varepsilon_\star$ and
$\varepsilon^{\mathrm{fr}}(W):=\varpi_{\mathrm{fr}}\alpha_{*,k_W}/(4C_{\mathrm{fr}})$.

(h) \emph{Rates and envelope functions.} $\delta_1^{\mathrm{sc}}$ is the layer rate per unit of the
scaled distance $d_{\mathrm{sc}}$, the distance defined in terms of $M_1$-cubes on the corresponding
scales \cite{B-CMP122a}; every $\delta_1$ multiplying $R_k$, $R_{\min}$, $R(\mathsf T)$ or a
$d_{\mathrm{sc}}$ in this chapter is $\delta_1^{\mathrm{sc}}$. Further,
\begin{equation*}
\delta_L:=\tfrac12\min\{\delta_1^{\mathrm{sc}},\delta_G\},\qquad
\delta_{\natural}:=\min\{\delta_1^{\mathrm{sc}},\delta_L,1\},
\qquad\text{so}\qquad \delta_\natural\le\delta_L\le\tfrac12\delta_1^{\mathrm{sc}},
\end{equation*}
$\delta_L$ and $C_L$ being the rate and the constant of the localisation of the current at the
localisation site, and $\delta_G$ the rate of the kernel localised there. $C_\natural$, $K_{\mathfrak Q}$
and $\alpha^+_k:=\bar\alpha(x_k)$ are the constants and the envelope of the same estimates:
$\varepsilon_{\mathfrak Q}(k)\le K_{\mathfrak Q}\alpha^+_k$, $\bar\alpha$ an explicit envelope
function with $\tilde\alpha_0+\tilde\alpha_1\le\bar\alpha$, and the radius of the tube at $(L)$
satisfies $\tilde\alpha_1(k)\le\alpha^+_k$. The explicit functions
\begin{equation*}
\mathsf{a}_J(x):=2.1\,C_{\theta'}^{-1}(\log x+2)^{q-p_1}\cdot\tfrac12e^{-2}C_\natural\bar\alpha(x),
\qquad \mathsf{a}_T(x):=B_H\alpha_*(x)
\end{equation*}
majorise $2.1K_u\mathsf{w}_L\mathsf{U}_L$ and the corresponding quantity at $(T)$; here $p_1$ is an
exponent of the envelope estimates, taken by name, and $C_{\theta'}$ the free constant of the ratio
$\theta'$. The numbers
$K_D$ (the constant of linearity of the layer), $c_\chi\ge1$, $\varphi_{\flat}$ and
\begin{equation*}
C^J_{\mathrm{age}}:=2\sqrt2(1+b_+)^{1/2}\sup_{u\ge0}(1+u)^{1/2}e^{-\frac12\delta_\natural u}
\end{equation*}
are the numbers of the off-location Cauchy-disc statement; they are fixed before $\gamma$, $K_D$
before $C_{\theta'}$.

(i) \emph{The positional factor.} At site $(L)$, for every term that carries the amplitude
$\mathsf t$ of the amplitude lemma, and for each localisation datum $p=(\square,Y_0,\text{walk})$ of
such a term $v$ with read set $S_v$, put
\begin{equation*}
\operatorname{Sten}_2(S_v):=\bigl\{y:\ |y-y_0|\le2\ell_i\ \text{for some }y_0\in S_v\text{ and
some closed cell }\Delta\ni y_0,\ i:=\operatorname{lev}\Delta\bigr\},
\end{equation*}
the cell $\Delta$ being arbitrary (in $\operatorname{St}_v$ or not) and $\ell_i$ the length unit of
level $i$, and
\begin{equation}\label{5.1:eq-correlation-notation-a}
\eta_p:=\begin{cases}
1,&\operatorname{Sten}_2(S_v)\cap\square^{\zeta}\neq\emptyset,\\
e^{-\delta_Pd_{\mathfrak{B}_s}(\square^{\zeta},S_v)},&\text{otherwise},
\end{cases}
\qquad \square^{\zeta}:=\operatorname{supp}\hat\chi_{\square},\quad\delta_P:=\tfrac18\delta_0,
\end{equation}
where $d_{\mathfrak{B}_s}$ is the scaled contour distance of the bracket's determining set, with
$d_{\mathfrak{B}_s}\ge d^{\wedge}$. At the sites $(T)$, $(P)$, $\eta_p$ is instead the factor
$e^{-\delta_1M\operatorname{dist}(\square,S_v)}$, $\delta_1M\ge\kappa_1$, of the resummation
geometry at those sites, taken by statement. One has $\square^\zeta\subset\square\subset\tilde\square$,
$\tilde\square$ Ba{\l}aban's fattened block \cite[p.~271]{B-CMP109}, and $\square^\zeta$ lies at
distance at least
$\tfrac13M$ inside $\square$. A block is \emph{touching} if $\operatorname{Sten}_2(S_v)$ meets
$\square^\zeta$ ($\eta_p=1$) and \emph{receding} otherwise ($\eta_p<1$); at a receding block
$\hat\chi_\square$ vanishes, and is in particular constant, on $\operatorname{Sten}_2(S_v)$. Thus
$\eta_p\le1$ is a function of the pair $(\square,S_v)$, equal to $1$ at every touching block; it
bounds the size on $S_v$ of the unit shift's piece $\delta\mathbb{H}_p$ in units of the unit shift
(property (k1)) and is not that size itself. For $y\in S_v$ one has
$d_{\mathfrak B_s}(\square^\zeta,S_v)\le d_{\mathfrak B_s}(y,\square^\zeta)$, so
$e^{-\delta_Pd_{\mathfrak B_s}(y,\square^\zeta)}\le\eta_p$ at a
receding block, and $e^{-\delta_Pd_{\mathfrak B_s}(y,\square^\zeta)}\le1=\eta_p$ at a touching one;
in particular
$\eta_p\ge e^{-2\delta_Pd_{\mathfrak B_s}(\square^\zeta,S_v)}$.

The block sums at $(L)$ are bounded as follows. For an anchor cell $\Delta$ write
$\eta(\square,\Delta)$ for the factor \eqref{5.1:eq-correlation-notation-a} of the pair
$(\square,\Delta)$, that is, with $S_v$ replaced by $\Delta$, and $y_*$ for the base point of
$\Delta$. If $\eta(\square,\Delta)=1$, there are $y\in\square^\zeta$ and $y_0\in\Delta$, with $y_0$
held by a closed cell of some level $i$ and
$|y-y_0|\le2\ell_i$; by the stencil lemma, $i\le m+1$ and $|y-y_0|\le2L\ell_m<\tfrac13M\ell_m$, $m$
the zone of the block, so $y_0$ is an interior point of $\square$, hence $\Delta$ meets $\square$ and
$\Delta\in\Sigma_\square\cup\partial^+\Sigma_\square$, where $\Sigma_\square$ is the set of cells
meeting $\square$ and $\partial^+\Sigma_\square$ the set of their neighbours. A block of zone $n$
meets cells of the levels $n-1,n,n+1$ only (the level property of the cell geometry), whence
$\#\Sigma_\square\le(2L)^d$; a cell touches at
most $\hat b_t$ others (the adjacency bound of the cell geometry), whence
$\#\partial^+\Sigma_\square\le\hat b_t\#\Sigma_\square$; and the cell-sum
lemma, applicable at the rate $\tfrac12\delta_P\ge\tfrac1{16}\delta_P$ under its weight floor,
gives, since $d_{\mathfrak B_s}\ge d^\wedge$ and
$\square^\zeta\subset\Sigma_\square$,
\begin{equation*}
\sum_{\Delta\notin\Sigma_\square\cup\partial^+\Sigma_\square}
e^{-\frac12\delta_Pd_{\mathfrak B_s}(\square^\zeta,y_*)}
\le C_{\mathrm{cell}}\,\#\partial^+\Sigma_\square .
\end{equation*}
Bounding $\eta^{1/2}\le1$ on $\Sigma_\square\cup\partial^+\Sigma_\square$ and using the last display
off it, one obtains
\begin{equation}\label{5.1:eq-correlation-notation-b}
\sup_\square\sum_\Delta\eta(\square,\Delta)^{1/2}
\le(1+\hat b_t)(2L)^d+C_{\mathrm{cell}}\hat b_t(2L)^d
=C_H^{\wedge\prime}:=(2L)^d\bigl(1+\hat b_t+\hat b_tC_{\mathrm{cell}}\bigr).
\end{equation}
Per anchor cell, the cell-sum lemma at $\Sigma:=\Delta$ gives
\begin{equation*}
\sup_\Delta\sum_\square\max\bigl\{\mathbf{1}[\Delta\in\Sigma_\square\cup\partial^+\Sigma_\square],
e^{-\frac12\delta_Pd^{\wedge}(\square^\zeta,\Delta)}\bigr\}
\le(1+\hat b_t)C_{\mathrm{blk}}+C_{\mathrm{blk}}(1+\hat b_tC_{\mathrm{cell}}),
\end{equation*}
$C_{\mathrm{blk}}:=2^{d+1}+(L+1)^d$, the number $3\cdot2^d$ of blocks covering a point being at
most $C_{\mathrm{blk}}$. Both are $(d,L)$-numbers: they contain no $M$, $M_1$, $\kappa_1$, $N_0$,
$\check R$, and no level, age, zone or band count. The coarse blocks $\square_m$ over a fine cell
are receding, since $\square_m^\zeta$ keeps $\tfrac13M$ cubes of level $n'+m-1$ away from zone
$n'+m-2$, $n'$ the level of the fine cell, and
\begin{equation*}
e^{-\frac12\delta_Pd_{\mathfrak B_s}(\square_m^\zeta,S_v)}
\le\exp\bigl(-\tfrac12\delta_P\cdot(\text{zone width})\cdot(m-2)_+\bigr)\le e^{-8(m-2)_+},
\end{equation*}
since, by the zone-width floor of the response theorem,
$d_{\mathfrak B_s}\ge(\text{zone width})\times(\text{number of zones crossed})$, and
$\delta_P\times(\text{zone width})\ge16$. For the
pointed $p$ of the first and the second amplitude class,
$\eta_p\le e^{-\delta_Pd_{\mathfrak B_s}(\tilde\square,S_v)}$ at every
block: at a receding block because $\square^\zeta\subset\tilde\square$ gives
$d_{\mathfrak B_s}(\square^\zeta,S_v)\ge d_{\mathfrak B_s}(\tilde\square,S_v)$; at a touching block
because, by the stencil lemma,
$\square$ meets a member of $\operatorname{St}_v$, so that $\tilde\square\supset\square$ meets $S_v$
and $d_{\mathfrak B_s}(\tilde\square,S_v)=0$. Hence the block lemma of the pointed summation, whose
shells are taken in $d_{\mathfrak B_s}(\tilde\square,S_v)$, bounds $\sum_\square\eta_p^{1/2}$ for
these $p$.

(j) \emph{The response piece.} $\hat\chi_\square$ is Ba{\l}aban's $\zeta_\square$
\cite[(3.4)--(3.7)]{B-CMP109}: $\operatorname{supp}\hat\chi_\square\subset\square$,
$|\hat\chi_\square|\le1$, $\sum_\square\hat\chi_\square=1$ on the support of the datum, and
$\hat\chi_\square(x)=\prod_\mu\phi(M^{-1}(x_\mu-y_\mu))$, with $\phi$ a fixed function on
$\mathbb{R}$ equal to $1$ on $|t|\le\tfrac13$ and to $0$ on $|t|\ge\tfrac23$, $\square$ a union of
$2^d$ neighbouring cubes of $\pi_k$ around $y$ \cite[p.~270]{B-CMP109}. It multiplies the response:
in Ba{\l}aban's
notation ($t$ his interpolation variable, not the scale ratio of (b), $\eta$ his scale parameter
of \cite{B-CMP109})
\begin{equation*}
\begin{split}
&\sum_\square\int_0^1dt\,\Bigl\langle\frac{\delta}{\delta\mathbf{H}}\mathbf{E}^{(j)}
\bigl(U_j(\exp iQ_j(\eta t\mathbf{H}_k(B'))\bar U^j_{k+1})\bigr),\hat\chi_\square\mathbf{H}_k(B')
\Bigr\rangle\\
&\quad=\sum_\square\int_0^1dt\,\frac{d}{dt_\square}\mathbf{E}^{(j)}\bigl(U_j(\exp iQ_j(\eta(t+
t_\square\hat\chi_\square)\mathbf{H}_k(B'))\bar U^j_{k+1})\bigr)\Big|_{t_\square=0}
\end{split}
\end{equation*}
\cite[(3.4)]{B-CMP109}; $B'$ is never cut. At $(L)$ the layer is
$(t+\mathsf{t}\hat\chi_\square)\mathbb{H}_s$, with $s$ the index of the layer (unrelated to the
slot combination $s$ of (a)), and
\begin{equation*}
\delta\mathbb{H}_p:=\hat\chi_\square\mathbb{H}_s
\end{equation*}
($s=4$ for the terms whose layer is the fourth, a class of the domain convention unrelated to the
two amplitude classes of (f):
$\delta\mathbb{H}_p=\hat\chi_\square\mathbb{H}_4$,
$\mathbb{H}_4=H'_k(g_kB)$ with $H'_k$ the response map of the induction format (cf.\
\cite[(1.20)]{B-CMP122b} for Ba{\l}aban's $\mathbf{H}''_k$ and its first-order operator
$H''_{1,k,Z}$)), linear
in $\mathsf t$, supported in
$\tilde\square$, with
$\sum_\square\hat\chi_\square\mathbb{H}_s=\mathbb{H}_s$, so that the pieces of the amplitude lemma
sum to the bracket. Where a term is read in Ba{\l}aban's box representation, the cut response
enters as the datum $Q_j(\eta(t+t_\square\hat\chi_\square)\mathbf{H}_k(B'))=B_\square$,
$Q_j(\eta t\mathbf{H}_k(B'))=B(t)$ of the level-$j$ box minimiser $\mathbf{H}_j$
\cite[(3.6)]{B-CMP109}, and the piece is its linear answer
\begin{equation*}
\delta\mathbf{H}_j=\Bigl\langle\frac{\delta\mathbf{H}_j}{\delta B}(B(t)),\Bigl\langle
\frac{\delta Q_j}{\delta\mathbf{A}}(\eta t\mathbf{H}_k(B')),L^j\eta\hat\chi_\square\mathbf{H}_k(B')
\Bigr\rangle\Bigr\rangle
\end{equation*}
\cite[(3.8)]{B-CMP109}, which satisfies on $X$
\begin{equation*}
|\delta\mathbf{H}_j|\le B_3\exp\bigl(-\tfrac12\delta_0\operatorname{dist}^{(\xi)}(X,\tilde\square)
-\tfrac12\delta_0M(L^j\eta)^{-1}\bigr)\cdot2L^j\eta|\hat\chi_\square\mathbf{H}_k(B')|
\end{equation*}
\cite[(3.9)]{B-CMP109}. Moreover, with the constant $c_\chi\ge1$ of the response theorem,
$N_y(\hat\chi_\square\mathbb{H}_s)\le c_\chi N_y(\mathbb{H}_s)$; for $\mathbb H_4$ the decay estimate
used is \cite[(1.45)]{B-CMP122a}, on the ball where $\sup|B|\le4\delta'_k$.

(k) \emph{Profiles.} $e_p(y)$ is the size of the piece $\delta\mathbb{H}_p$ at the point $y$ in units
of the site's unit shift $\mathsf{U}_{\mathrm{site}}$:
$\xi_p(y):=\Theta_0^{\mathrm{loc}}(\delta\mathbb{H}_p)(y)=\mathsf{U}_{\mathrm{site}}\,e_p(y)$. Then
(k1) $e_p\le\eta_p$ on $S_v$; and (k2) $e_p$ decays from the sources of the shift at the site's own
rate. At $(T)$,
$e_p(y)\le e^{-\delta_1^{\mathrm{sc}}(d_{\mathrm{sc}}(y,\Omega)-1)}$, $\Omega$ the fluctuation domain
of the step, in the unit $\mathsf{U}_T:=B_HK_Rg_k\sup|\mathsf{B}|$, with $\mathsf B$ the complex
fluctuation field and $K_R$ the constant, proportional to $1/(1-\theta_S)$, of the complex radius
lemma at $(T)$, which enters the radius condition
$(1+\mathsf w)K_Rg_kT^{\mathbb C}\le\alpha_{*,k}$ of (l): the piece decays from $\square$
by \cite[(3.9)]{B-CMP109}, its datum $\hat\chi_\square\mathbf{H}(B')$ carries
$e^{-\delta_1^{\mathrm{sc}}d_{\mathrm{sc}}(\square,\Omega)}$ since $B'$ lives on $\Omega$, and the
triangle inequality, with the loss $-1$ for the passage to the top block, gives the bound; this
is a statement at $(T)$ only. At $(P)$, (k2) is given by the stage's band. At $(L)$, for a live law
site $s\in S_v$ of the term (a point, unrelated to the slot combination of (a) and to the layer
index of (j)) --- of class $(\alpha)$, $s\notin Z$, or of class $(\beta)$,
$s\in Z_h\subset D_2$, in which case the term is a first-part term presented at one of the three
chain backgrounds and $D_2$ is the second domain of that presentation --- put
$\mathfrak{F}_{\mathfrak s}(s):=\Lambda^{\sim}$ in class $(\alpha)$
and $\mathfrak{F}_{\mathfrak s}(s):=\mathsf{S}$ in class $(\beta)$, $\mathsf S$ the neighbourhood of
$\Omega''^{\sim2}_{h+1}$ in $\Lambda$ of width $2M_1+3$ blocks; then
\begin{equation*}
e_p(s)\le e^{-\delta_\natural(d_{\mathrm{sc}}(s,\mathfrak{F}_{\mathfrak s}(s))-1)}
\tag*{$(\mathrm{k2})^{L}$}
\end{equation*}
in the unit $\mathsf{U}_L:=C_LB_H\sum_{s'}\sup|\arg_{s'}|$, the sum running over the law sites
$s'$ of the suprema of the presented arguments $\arg_{s'}$ there; $\mathsf U_L$ is the majorant of
all the presented
layers, with $K_u\mathsf{U}_L\le C_LK_{\mathfrak Q}\alpha^+_k$. It is kept apart from the unit shift
of the correlation theorem itself, $\mathsf{U}_J:=6B_H\delta'_k\le\mathsf{U}_L$, the unit of the
$B$-layer, for which alone $\mathsf{w}_L\mathsf{U}_J\le\tilde\alpha_1(k)$ holds. $(\mathrm{k2})^{L}$ is
proved later in this chapter, in the treatment of the localisation site; there
$d_{\mathrm{sc}}(s,\mathfrak{F}_{\mathfrak s}(s))-1\ge\check R_k-2$ in both classes, increased by
$5R_{\min}\mathsf a_k(s)$, $\mathsf a_k(s)$ the age exponent, at the frame sites $s\in\Sigma_W$.

(l) \emph{Radii.} $\rho_A(v)$ is the Cauchy radius of the amplitude lemma in the rescaled amplitude
$\mathsf t$: $\rho_A:=\mathsf{w}/\eta_p$ for the pointed terms at the sites $(T)$, $(P)$;
\begin{equation*}
\rho_A:=\mathsf{w}_L/(\eta_p\mathsf{f}_v)^{1/2},\qquad
\mathsf{f}_v:=\sup_{y\in S_v}\min\{1,c_\chi\mathsf{f}(y)\},
\end{equation*}
for the pointed terms at $(L)$; and $\rho_A:=\mathsf{w}/\eta_p^{1/2}$ for every term of the fourth
kind. At $(L)$, $\mathsf{w}=\mathsf{w}_L:=1/\theta'(g_k)$. For the pointed terms at $(L)$, $S_v$ is
their read set: $\mathfrak{r}_v$ for the compact ones, the cells meeting $X_v$ for long singles and
tails; the latter are read at $\mathsf t=0$ only. For the pieces read on their own domain,
$\rho_A(v):=\mathsf{w}_0r^{\mathrm{own}}(\square)$, $r^{\mathrm{own}}(\square)$ the own-read radius
of the block, with the factor
$\eta^{\mathrm{own}}_p:=e^{-\frac14\delta_{\mathrm{own}}d_j(X_v,\square^\zeta)}$,
$\delta_{\mathrm{own}}$ the decay rate of the clause for pieces read on their own domain, taken by
name. At $(T)$,
\begin{equation*}
\mathsf{w}_T:=\mathsf{w}^{\mathbb C}:=1/\theta^{\mathbb C}_k-1,\qquad
\theta^{\mathbb C}_k:=K_Rg_kT^{\mathbb C}/\alpha_{*,k},
\end{equation*}
the largest $\mathsf w$ with $(1+\mathsf w)K_Rg_kT^{\mathbb C}\le\alpha_{*,k}$ (indeed
$(1+\mathsf w)\theta^{\mathbb C}_k\le1$ if and only if $\mathsf w\le1/\theta^{\mathbb C}_k-1$), while
Ba{\l}aban's real-circle radius is written
$\mathsf{w}_T^{\mathrm{real}}:=\alpha_2/(4B_0C_1e^{16\kappa_1}\delta_k)$.

(m) \emph{Prefactors of the increments.} $C_P$ is the common prefactor of the increments of the
$\mathbf{R}$-layers: $P,P',P''\le C_Pe^{-3\delta_1^{\mathrm{sc}}R_k}$, where
\begin{equation*}
\begin{gathered}
P:=K_WB_HK_\Lambda(M)(1+\varrho_k)(1+\beta_0)^2C_{1/2}e^{-5\delta_1^{\mathrm{sc}}R_k},\\
C_{1/2}:=\sup_{x\ge0}(2+x)^{1/2}e^{-0.6x},\qquad
C_{1/2}':=\sup_{x\ge0}(2+x)^{1/2}e^{-0.3x},\qquad C_P':=C_P\,C_{1/2}'/C_{1/2}.
\end{gathered}
\end{equation*}
Here $K_W$, $K_\Lambda(M)$ and $\varrho_k$ are the constants of the increment proposition, taken by
name. The coupling ratio coarse over fine inside $P$ is
$\alpha_{\cdot,k}/\alpha_{\cdot,m}\le(1+\beta_0)^2(k-m)^{1/2}$ by \cite[(2.6), (2.7)]{B-CMP119}. The
shell factor $2^{-(k-1-m)}$ comes from the floor on $\delta_1^{\mathrm{sc}}R_{\min}$ of the order
of choice, which gives $10\delta_1^{\mathrm{sc}}R_{\min}\ge2$:
\begin{equation*}
e^{-10\delta_1^{\mathrm{sc}}R_{\min}(k-1-m)}\le e^{-2(k-1-m)}\le e^{-0.6(k-1-m)}2^{-(k-1-m)},
\end{equation*}
since $e^{-1.4}\le\tfrac12$. At half rate,
\begin{equation*}
e^{-5\delta_1^{\mathrm{sc}}R_{\min}(k-1-m)}\le e^{-(k-1-m)}=2^{-(k-1-m)}e^{-(1-\ln2)(k-1-m)},
\end{equation*}
with $1-\ln2\ge0.3$ and $(k-m)^{1/2}e^{-0.3(k-1-m)}\le C_{1/2}'$: the shell factor survives the
halving, and $C_P'$ is the corresponding prefactor.

(n) \emph{Labels.} Every boundary-type term $F$ of level $j$ (the terms $\mathbf{B}^{(j)}$ born at
$\mathbf T$, the terms $\mathbf{B}'^{(j)}$ born at $\mathbf{R}$, the chain terms
$\mathbb{C}^{(n)}_k$, the frozen chain terms, the change-of-variables terms) carries source-free
data: its label $A_F\in\mathbf{D}_j$, a datum of the term and not a function of its true domain,
and a finite set $\mathcal{P}_F$ of pairs $(Y,m)$, $m\le j$, live at the current time, each with its
strip $Y^{\sharp}$ touching $A_F$. $F$ reads $((\mathbf U,\mathbf J),\mathbf A,w)$ on its true
domain
\begin{equation*}
X_F:=A_F\cup\bigcup_{(Y,m)\in\mathcal{P}_F}Y^{\sharp}
\end{equation*}
only. $F$ is \emph{wrapped} if $\mathcal{P}_F\neq\emptyset$ and \emph{plain} if
$\mathcal{P}_F=\emptyset$ (then $A_F=X_F$ is Ba{\l}aban's $X$). Every length, count, closure and
weight of a boundary-type term is taken on the label: $d_j(A_F)$, and $K_v:=[A_v]_k$, the cubes of
$\pi_k$ meeting $A_v$, with $d_k(K_v)\le L^{-(k-j)}d_j(A_v)$; nothing is asserted about $d_j(X_F)$.
At birth the label is $A:=\operatorname{sk}(X)$, the skeleton of $X$ off the strips born at
$\mathbf R_k$, with $\#\{X:\operatorname{sk}X=A\}\le e^{\lambda_\theta(d_k(A)+1)}$; for an output of
$\mathbf T$ it is $X$ itself. For a union $Y$ of cubes of $\pi_k$,
\begin{equation*}
Y^{+}:=Y\cup\bigcup\bigl\{Y_p^{\sharp}:(Y_p,m)\text{ live},\ Y_p^{\sharp}\text{ touches }Y\bigr\}.
\end{equation*}
With $W_F:=\bigcup_{\mathcal P_F}Y^{\sharp}$ one has $d_{j,W_F}(X_F)\le d_j(A_F)$, the length on the
left being that of \cite[(1.67)]{B-CMP122b}; hence a bound
$|\mathbf{B}^{(j)}|\le B_0e^{-\kappa d_j(A_F)}$ on the label implies Ba{\l}aban's form
$|\mathbf{B}^{(j)}(X)|\le B_0\exp(-\kappa d_{j,W_F}(X))$ at $X=X_F$.

(o) \emph{Further constants.} $c_1^{\flat}(g):=C_{99}\,c_1$, where $c_1$ is Ba{\l}aban's constant:
$(\alpha^\sharp)^{1/3}$ where the domain $X'$ of his bound meets the union $\bigcup Y_i^{\sharp}$
of the strips about the large-field components, and $e^{-p_0(g)}$ otherwise \cite{B-CMP122b}; and
$C_{99}$ is
the constant $O(1)$ of
\begin{equation*}
|\mathbf{R}'^{(k)}(X,(\mathbf U,\mathbf J))|\le O(1)c_1\exp\bigl(-(1+\tfrac12\beta)\kappa
d_{k,\cup Y_i}(X)\bigr)
\end{equation*}
\cite[(1.99)]{B-CMP122b}, the prefactor of the convergent series \cite[(7.13)]{B-CMP122a}, a pure
constant fixed before $\gamma$; every display
$e^{\lambda_\theta}c_1^\flat(g_j)e^{-(1+\frac14\beta_s)\kappa d_j(A)}$ carries this constant once.
$\lambda_\theta:=(1+3^d)\mathsf{c}_d\log2$ is a $d$-only number, and the condition
$e^{\lambda_\theta}\gamma^{\kappa_0}\le\min\{1,B_0\}$, with $\kappa_0$ the exponent of the bound
$c_1^\flat\le g_k^{\kappa_0}$ required of $c_1^\flat$, belongs to the list defining $\gamma_J$.
$\mathsf{K}^{\mathrm{arr}}$ is the sum of the convergent series of the collar proposition.
$C_0^{(2.28)}$ is Ba{\l}aban's free constant $C_0$ of \cite[(2.28)]{B-CMP119}, bounded below by
$2.5A_0$ divided by the constant of the spaces on which the live arrested pieces are read.
The room number of the newborn terms of level $k$ is
$\mathsf{c}_e:=\theta_S$, where $\theta_S\in[\tfrac12,\tfrac89]$ is the layer-schedule fraction on
the old layers $m<k$ of the space after $\mathbf T$ (factor $1+\beta$ on the top layer,
$s^{(k)}_m+\theta_S\beta2^{-(k-m)}$ on the layer $m<k$), an absolute number fixed first in the order
of choice, and
\begin{equation*}
\mathsf{w}_0:=\min\{\tfrac12,\theta_S\}/12=1/24,
\end{equation*}
since $\theta_S\ge\tfrac12$. The absolute constant of the absolute conditions of the amplitude
lemma is written $c_{\mathrm{abs}}$.

(p) \emph{The regularity ceiling.} At density $k$,
\begin{equation*}
\gamma_0^{(k)}:=5\max_{j\le k}w_*(j)\bar\alpha_j,\qquad w_*(j)=\mathsf{L}_0^{2N_0(g_j)},
\end{equation*}
a function of the running coupling, not one number for all densities; here $\bar\alpha_j$ is the
radius at level $j$ entering the ceiling (the indexed letter; the unindexed $\bar\alpha$ of (h) is
the envelope function), and $\alpha(\mathsf T)$ denotes Ba{\l}aban's radius as a
function of the logarithm $\mathsf T$ of (d). The weights satisfy
\begin{equation*}
\mathsf{L}_0^{2N_0}\le\bigl(\tfrac13\bigr)^{N_0}L(L+1)(N_0+1)^{\beta_0}R_k<L(L+1)N_0^{-1}R_k
\end{equation*}
\cite{B-CMP122a}, $N_0$ the number of zones of (d). Put
\begin{equation*}
\Psi(\mathsf T):=L^2(L+1)C_\Psi\mathsf T^{\,r_{\mathrm{pr}}+q}e^{-\mathsf T/2},
\end{equation*}
$C_\Psi$ a fixed constant; $\Psi$ is decreasing beyond a threshold depending on
$(r_{\mathrm{pr}},q)$, and $\Psi(\mathsf T-c_\beta)\le(1+\beta_0)\Psi(\mathsf T)$, where
$c_\beta:=2\ln(1+\beta_0)$. Since
$\bar\alpha_j\le(1+\beta_0)\alpha(\mathsf T_j)$ and $\mathsf T_j\ge\mathsf T_k-c_\beta$ for $j\le k$,
by \cite[(2.6)]{B-CMP119}, these bounds give the majorant
\begin{equation*}
\gamma_0^{(k)}\le\bar\gamma_0(\mathsf T_k):=5(1+\beta_0)^2\Psi(\mathsf T_k).
\end{equation*}
Every condition on $\gamma_0^{(k)}$ is read
at $\bar\gamma_0(\mathsf T)$ inside one supremum over $\mathsf T\ge\mathsf T_\gamma$.

(q) \emph{Two standing readings.} First: the profile quantities of (k), defined through
$\Theta_0^{\mathrm{loc}}$ at law sites, majorise the four-entry cell norm $\mathsf{N}_y$ of
\cite[(3.14)]{B-CMP109} (with entries $N^\Delta_y(A')$, $N_y(\mathbb H)$,
$\ell^3|P_1(\sigma)A'|$, $\ell^3|\mathfrak{m}_H(\sigma)X|$) of the layer $\delta\mathbb{H}_p$ and of
the undilated layers, at every cell $y$ of $S_v$ and at every current block holding a site of
$S_v$. At $(T)$ this holds by the Lipschitz clause of the response theorem read at the exponential
weights in the distance $d_{\mathfrak B_s}$ to $\square$ and in the distance $d^\wedge$ to $\Omega$,
at the rate of the weight class of that theorem (not the decay fraction $\delta$ of (c)), with the
constant
$B_\sharp$ enlarged to $B_H$, the weight of the kernel bounded by $1$, and $\mathsf{w}_T$ of (l);
at $(L)$ the entries in $\mathbf U$ hold by the same clause on the supported arguments
\cite[(3.9)]{B-CMP109}, with the resummation of the localisation of the current for the first
argument, and the entries in $\mathbf J$, at touching and at receding blocks, by the
pointwise product of the corollary of the current-recursion lemma,
\begin{equation}\label{5.2:eq-three-sites-c}
N_y(\hat\chi_\square\mathbb{H}_s)\le20B_\sharp q_\square(y),\qquad
q_\square(y)\le p_s(y)e^{-\delta_Pd_{\mathfrak B_s}(y,\square^\zeta)}\le p_s(y)\eta_p,
\end{equation}
at every block and every $y\in S_v$, the member of the third entry carrying $\mathsf t$ being at
most $20K_\chi|\mathsf t|q_\square(y)$, $K_\chi$ the constant of the $\mathsf t$-entry of the
current-recursion lemma, the case of a receding block being supplied by the first clause of the
current-recursion lemma; here
$p_s(y):=\mathsf{a}_se^{-\delta_Pd^\wedge(y,A_s)}$ and
$q_\square(y):=\sup_{y'\subset\square^\zeta}p_s(y')e^{-2\delta_Pd_{\mathfrak B_s}(y,y')}$. Second: in
the notation
of the off-location Cauchy-disc statement --- $\mathfrak Q(\mathfrak a)$ the family of terms
attached there to a live arrest $\mathfrak a$ and filed in the slot of the amplitude lemma, $e$
the datum of a presentation with the two sets $S^{+}(e)$ and $\tilde S(e)$ so denoted there,
$p(z)$ the point presented at the parameter $z$, $g(z)$ (not the coupling) the presented quantity
of that statement, $\mathfrak K_G(\cdot)$ the space of global pairs over a given set of the
global-pair statement to which it refers, and $\operatorname{pol}$ the size functional applied to
$g(z)$ there --- at every
presentation at $(L)$ that is not deep in the sense of that statement, of a member of
$\mathfrak{Q}(\mathfrak a)$ or of a stored term, the
presented point $p(z)$ is the restriction of a global pair lying in a space of the same kind over
some $Y\supseteq S^{+}(e)$, hence
\begin{equation}\label{5.2:eq-half-disc-containment-g}
p(z)\in\mathfrak{K}_G(\tilde S(e)),\qquad \operatorname{pol}g(z)<\varepsilon_\star .
\end{equation}
\end{definition}

\begin{remark}
The following letters are distinct and are never identified. $T(g)=A_1(\log g^{-2})^{p_0}$ and
$T_k=\mathfrak{b}_k$ are the thresholds of the setting of the correlation theorem, while
$\mathsf T$, $\mathsf T_k$, $\mathsf T_\gamma$, $\mathsf T_*$ are logarithms. $\delta_1^{\mathrm{sc}}$
of (h) is not the rate per lattice unit in the bound $\eta_p=e^{-\delta_1M\operatorname{dist}(\square,S_v)}$,
$\delta_1M\ge\kappa_1$, which is used at the sites $(T)$, $(P)$ only and nowhere at $(L)$; nor the
rate of the kernel lemma of the fluctuation integral; nor Ba{\l}aban's $\delta_1$ of
\cite{B-CMP116}, written $\delta_1^{116}$. The $N$ of the torus side is read only by
$\|\cdot\|_\sharp$, $\mathsf Q$ and the torus floor condition, while $N(g)$ and ``$n\le N$'' in the
collar
proposition count Ba{\l}aban's preparatory stages. $\mathsf{L}_0$ is Ba{\l}aban's auxiliary number
and not the constant $L_0$ of Theorem~0. $\gamma_0^{(k)}$ of (p) is neither the lower bound
$\gamma_0$ of $C^*\Delta_kC$ \cite{B-CMP99b} nor the minimum of the two rates called $\gamma_0$ in
the modified volume bound. $C_0^{(2.28)}$ of (o) is neither the $C_0$ of the order of choice nor
$C_0^\sharp=3c_G+1$. The collar constant $c_J$ is not the constant
$c_J^{\mathrm{rec}}:=\max\{K_Wc_\chi,(c_aK_Wc_\chi)^{1/2}\}$ of the current-recursion lemma, where
$c_a:=20K_PC_\rho B_\sharp/B_H$ and $c_{\mathrm{rec}}:=\max\{c_\chi,c_J^{\mathrm{rec}},c_a\}$.
$C_\sigma$ is a constant, $C_{(\sigma)}$ an
operator. The block size and radius of Ba{\l}aban's earlier papers are written $M_1$ and
$R^{\mathrm{prev}}$ ($M_1$ being the size of the large cubes for which the theorems of the previous
papers are valid \cite{B-CMP119}; cf.\ \cite{B-CMP116}); their lower thresholds depending on $(d,L)$
are written $M_E$, $R_E$, and
the associated $K_1$ is written $K_1^X$. The constant $C_{\mathfrak d}^\flat M^d$ of the charge
corollary is written $\kappa_{\mathfrak d}$, since Ba{\l}aban's $\kappa_1$ of \cite{B-CMP116} is fixed
before $M$. $\eta_p$ of \eqref{5.1:eq-correlation-notation-a} is the factor of the pointed terms
carrying $\mathsf t$; the factor of the pieces read on their own domain is $\eta^{\mathrm{own}}_p$.
\end{remark}

\begin{definition}[Lattices, sources and the torus floor]\label{5.1:def-torus-floor}
Throughout, $L$ is the blocking factor and $K$ the number of renormalisation steps.

\emph{Lattices.} Let $\mathbb{T}^4$ be the torus of side $1$. The bare lattice on $\mathbb{T}^4$ has
spacing $\varepsilon = L^{-K}N_T^{-1}$, where $N_T\ge 1$ is an integer. After $k$ steps the unit
lattice is $T^{(k)}$, and the bare lattice has spacing $\eta = L^{-k}$ in the units of $T^{(k)}$.
Fix an integer $m\ge 0$ independent of $K$. Put $K' := K-m\ge 0$, the last level of the run of
the correlation theorem, and $N := N_TL^m$, the number of sites per side of $T^{(K')}$.

\emph{Radii and block sizes.} For an inverse coupling $x\ge 1$ and for the inverse couplings
$x_i$ of the levels $i$ put
\begin{equation}\label{5.1:eq-torus-floor-1}
R(x) := \min\bigl\{L^j : j\in\mathbb{Z}_{\ge 0},\ L^j\ge (\log x)^{r_{\mathrm{pr}}}\bigr\},
\qquad
\check{R}_k := R\bigl(\min_{i\le k}x_i\bigr).
\end{equation}
Thus $R(x)$ is Ba{\l}aban's radius read at the inverse coupling $x$: in \cite{B-CMP119},
$R_0$ is the smallest power of $L$ with $R_0\ge(\log g_0^{-2})^{r_{\mathrm{pr}}}$. Accordingly,
the radius $R_k$ of \cite{B-CMP119} at level $k$ is read as $\check{R}_k$ in what follows. The block
sizes of \cite{B-CMP119} are powers of $L$ with $M_1<M_2<M$, where $M=L^{m'}$; the
exponent $m'$ is not the exponent $m$ fixed above.

\emph{The torus constant and the torus floor.} With $X=g^{-2}$ and the constants $\lambda$,
$c_2$ and $b_+=3\lambda+2c_2$ of the reference recursion for the $x_k$, put
\begin{equation}\label{5.1:eq-torus-floor-2}
C_{\mathbb{T}} := L^{11}M,
\end{equation}
\begin{equation}\label{5.1:eq-torus-floor-3}
N_{\mathbb{T}}(g,m) := C_{\mathbb{T}}\,\sup_{u\in\mathbb{Z}_{\ge 0}}
L^{-u}\,R\bigl(X+b_++4c_2+3\lambda(m+u)\bigr).
\end{equation}
The \emph{torus floor} is the condition
\begin{equation}\label{5.1:eq-torus-floor-a}
L^m\ \ge\ N_{\mathbb{T}}(g,m),
\end{equation}
and $m_{\mathbb{T}}(g)$ denotes the least integer $m\ge 0$ satisfying
\eqref{5.1:eq-torus-floor-a}. The integer $N_T\ge 1$ is free. Condition
\eqref{5.1:eq-torus-floor-a} is the only condition the correlation theorem places on the torus.

\emph{Action.} The gauge group is $\mathrm{SU}(N)$, with the trace normalised so that the trace
of the identity matrix is $1$, and $\mathrm{SL}(N,\mathbb{C})$ is its complexification. The
integer inside $\mathrm{SU}(N)$ and $\mathrm{SL}(N,\mathbb{C})$, here, in the fifth item below and
in its proof, is the size of the matrices, and is not the site number $N=N_TL^m$ fixed above. For a
plaquette $p$ put
\begin{equation*}
a_p(U) := 1-\tfrac12\operatorname{tr}\bigl(U+U^{-1}\bigr)(\partial p),
\qquad
A(c,U) := \sum_p c(p)\,a_p(U),
\end{equation*}
for a weight $c$ on plaquettes.

\emph{Sources.} Let $f_1,\dots,f_n\colon\mathbb{T}^4\to\mathbb{R}$ be Lipschitz, let
$w=\sum_i z_if_i$ with $z\in\mathbb{C}^n$, let $x_0$ be the bare inverse coupling, let $E$ be
a constant, independent of $U$ and of $z$, entering the exponent of the bare density additively,
let $r>0$, and let $\lambda_* := (40M)^{-1}$. For a function $w$ on $\mathbb{T}^4$ we write
$A(x_0-w,U)$ for $A(c,U)$ with the weight $c(p):=x_0-w(x(p))$, where $x(p)$ is the barycentre
of $p$. Put
\begin{equation}\label{5.1:eq-torus-floor-b}
\begin{gathered}
\rho_0^w := \exp\bigl[-A(x_0-w,U)-E\bigr],\qquad Z(z) := \int\rho_0^w\,dU,\\
\|f\|_\sharp := \max\Bigl\{\|f\|_\infty,\ \frac{\operatorname{Lip}f}{\lambda_*N}\Bigr\},
\qquad
\mathbb{D} := \Bigl\{z\in\mathbb{C}^n : \sum_i|z_i|\,\|f_i\|_\sharp<r\Bigr\}.
\end{gathered}
\end{equation}
Since $T^{(K')}$ has spacing $1/N$ on $\mathbb{T}^4$, the number $\operatorname{Lip}f/N$ is the
Lipschitz constant of $f$ measured in the lattice units of $T^{(K')}$; as $1/\lambda_*=40M$,
$\|f\|_\sharp$ is the norm $\max\{\|f\|_\infty,40M\cdot\operatorname{Lip}f\}$ of the general
conventions, the Lipschitz constant being taken in the lattice units of $T^{(K')}$.

These objects have the following properties.
\begin{enumerate}
\item $\check{R}_k$ is nonincreasing in $k$. The constant $C_{\mathbb{T}}$ is a power of $L$
and $C_{\mathbb{T}}\ge 100LM$. Every partition cube used in \cite{B-CMP119} at level $k$ has a
side that is a power of $L$ not exceeding $C_{\mathbb{T}}\check{R}_k$. These cubes are those
of sides $M\check{R}_k$, $LM_2\check{R}_k$, $LM\check{R}_{k+1}$, $L^2M_2\check{R}_{k+1}$,
$M$, $M_1$ and $M_2$. The cube of side $100M\check{R}_k$ of \cite{B-CMP122a} is not a
partition cube: it is a size condition on a component of a large-field region, and
$100M\check{R}_k$ is not a power of $L$. This cube has side at most
$C_{\mathbb{T}}\check{R}_k/L$.
\item $N_{\mathbb{T}}(g,m)$ is finite and is a power of $L$. It depends only on $g$, $m$,
$L$, $M$, $\lambda$, $c_2$ and $r_{\mathrm{pr}}$, and it satisfies
$N_{\mathbb{T}}(g,m+1)\le L\,N_{\mathbb{T}}(g,m)$.
\item $m_{\mathbb{T}}(g)$ exists, and \eqref{5.1:eq-torus-floor-a} holds if and only if
$m\ge m_{\mathbb{T}}(g)$.
\item (Divisibility.) Assume \eqref{5.1:eq-torus-floor-a}, and assume that the inverse
couplings $x_k$, $k\le K'$, satisfy \eqref{5.1:eq-correlation-theorem-a} of the correlation
theorem. At every level
$k=K'-u\le K'$, the torus $T^{(k)}$ has $N_TL^{m+u}$ sites per side. Every partition-cube side
$s$ listed in the first item divides $N_TL^{m+u}$. Hence $T^{(k)}$ is partitioned into cubes of
each of these sides, the partitions being nested $L$-adically. For each side $s$ listed in the
first item we write $\pi_k(s)$ for the partition of $T^{(k)}$ into cubes of side $s$, all these
partitions being taken from a common origin of $T^{(k)}$. The cube of side $100M\check{R}_k$
fits into $T^{(k)}$.
\item $a_p\ge 0$ on $\mathrm{SU}(N)$. The action $A(c,U)$ is linear in $c$, gauge invariant
plaquette by plaquette, and entire on $\mathrm{SL}(N,\mathbb{C})$.
\end{enumerate}

Condition \eqref{5.1:eq-torus-floor-a} is a lower bound on the block exponent $m$; it places
no upper bound on $K$. No inequality on the size of $N$ alone could replace it. Indeed, for
$m=0$ and $N_T\ge N_{\mathbb{T}}(g,0)$ prime to $L$, the torus carries no partition into
cubes of side $M$.
\end{definition}

\begin{proof}
\emph{First item.} The set over which the minimum defining $\check{R}_{k+1}$ is taken
contains that for $\check{R}_k$. Hence $\min_{i\le k+1}x_i\le\min_{i\le k}x_i$. By
\eqref{5.1:eq-torus-floor-1}, $R$ is nondecreasing on $[1,\infty)$, so
$\check{R}_{k+1}\le\check{R}_k$.

Since $M$ is a power of $L$, so is $C_{\mathbb{T}}=L^{11}M$. Since $L$ is an integer greater
than one, $L^{10}\ge 2^{10}\ge 100$, and therefore $C_{\mathbb{T}}=L^{10}\cdot LM\ge 100LM$.

Now let $s$ be one of the listed sides. It is a product of powers of $L$, hence a power of
$L$. Two facts give the bound $s\le C_{\mathbb{T}}\check{R}_k$. First,
$\check{R}_{k+1}\le\check{R}_k$ and $\check{R}_k\ge 1$. Second, $LM_2\le M$, because $M_2$
and $M$ are powers of $L$ with $M_2<M$. We check each side in turn:
\begin{itemize}
\item $M\check{R}_k\le C_{\mathbb{T}}\check{R}_k$;
\item $LM_2\check{R}_k\le M\check{R}_k$;
\item $LM\check{R}_{k+1}\le LM\check{R}_k$;
\item $L^2M_2\check{R}_{k+1}\le LM\check{R}_k$;
\item $M_1<M$ and $M_2<M$, while $M\le M\check{R}_k$ since $\check{R}_k\ge 1$.
\end{itemize}
Each right-hand side is at most $L^{11}M\check{R}_k=C_{\mathbb{T}}\check{R}_k$. For the cube of
\cite{B-CMP122a}: since $L$ is odd, $100=2^2\cdot 5^2$ is not a power of $L$, and as
$M\check{R}_k$ is a power of $L$, $100M\check{R}_k$ is not a power of $L$. Moreover
$100M\check{R}_k\le L^{10}M\check{R}_k=C_{\mathbb{T}}\check{R}_k/L$.

\emph{Second item.} If $t\ge 1$ and $L^j$ is the least power of $L$ with $L^j\ge t$, then
$L^j<Lt$. Applying this with $t=\max\{1,(\log y)^{r_{\mathrm{pr}}}\}$, for which the least power
of $L$ not below $t$ is $R(y)$, gives
$R(y)\le L\max\{1,(\log y)^{r_{\mathrm{pr}}}\}$. Now take
$y=X+b_++4c_2+3\lambda(m+u)$. Since $y$ is affine in $u$, the quantity
$(\log y)^{r_{\mathrm{pr}}}$ grows like a power of $\log u$, whereas $L^u$ grows exponentially
in $u$; hence
\begin{equation*}
L^{-u}R(y)\le L^{1-u}\max\bigl\{1,(\log y)^{r_{\mathrm{pr}}}\bigr\}\to 0
\qquad(u\to\infty),
\end{equation*}
that is, the terms $L^{-u}R(y)$ in \eqref{5.1:eq-torus-floor-3} tend to $0$ as $u\to\infty$.

Each term is an integer power of $L$, and the term at $u=0$ is positive. Only finitely many
terms exceed the term at $u=0$. Hence the supremum is a maximum, it is finite, and it is a
power of $L$. Its product with $C_{\mathbb{T}}$ is therefore a power of $L$.

The right side of \eqref{5.1:eq-torus-floor-3} involves only $X=g^{-2}$, $m$, $L$, $M$,
$\lambda$, $c_2$, $r_{\mathrm{pr}}$ and $b_+$, the last being determined by $\lambda$ and
$c_2$.

For the last assertion, substitute $u=v-1$:
\begin{align*}
N_{\mathbb{T}}(g,m+1)
&= C_{\mathbb{T}}\sup_{v\ge 1}L^{1-v}R\bigl(X+b_++4c_2+3\lambda(m+v)\bigr)\\
&\le L\,C_{\mathbb{T}}\sup_{v\ge 0}L^{-v}R\bigl(X+b_++4c_2+3\lambda(m+v)\bigr)
= L\,N_{\mathbb{T}}(g,m).
\end{align*}

\emph{Third item.} Put $X' := X+b_++4c_2$. Setting $v=m+u$ in
\eqref{5.1:eq-torus-floor-3} gives
\begin{equation*}
L^{-m}N_{\mathbb{T}}(g,m)=C_{\mathbb{T}}\sup_{v\ge m}L^{-v}R\bigl(X'+3\lambda v\bigr).
\end{equation*}
So $L^{-m}N_{\mathbb{T}}(g,m)$ is $C_{\mathbb{T}}$ times the tail supremum of a sequence that
tends to $0$, since, as shown in the proof of the second item,
\begin{equation*}
R\bigl(X'+3\lambda v\bigr)\le L\max\bigl\{1,\bigl(\log(X'+3\lambda v)\bigr)^{r_{\mathrm{pr}}}\bigr\}.
\end{equation*}
The tail supremum tends to $0$ as $m\to\infty$; hence $L^{-m}N_{\mathbb{T}}(g,m)\to 0$, and
\eqref{5.1:eq-torus-floor-a} holds for some $m$. Hence $m_{\mathbb{T}}(g)$ exists.

Now suppose \eqref{5.1:eq-torus-floor-a} holds at $m$. By the second item,
\begin{equation*}
L^{m+1}\ge L\,N_{\mathbb{T}}(g,m)\ge N_{\mathbb{T}}(g,m+1),
\end{equation*}
so the condition holds at $m+1$. The set of $m$ satisfying \eqref{5.1:eq-torus-floor-a} is
therefore closed upwards. Its least element is $m_{\mathbb{T}}(g)$, so the set is exactly
$\{m\ge m_{\mathbb{T}}(g)\}$.

\emph{Fourth item.} On $\mathbb{T}^4$ of side $1$, the bare lattice has $L^KN_T$ sites per
side. The unit lattice $T^{(k)}$ has $L^{K-k}N_T$ sites per side, and
$K-k=K-K'+u=m+u$. So $T^{(k)}$ has $N_TL^{m+u}$ sites per side.

Let $s$ be one of the partition-cube sides of the first item. We claim
\begin{equation}\label{5.1:eq-torus-floor-4}
\begin{aligned}
s\ \le\ C_{\mathbb{T}}\check{R}_k\ &\le\ C_{\mathbb{T}}R(x_k)
\ \le\ C_{\mathbb{T}}R\bigl(X+b_++4c_2+3\lambda(m+u)\bigr)\\
&\le\ N_{\mathbb{T}}(g,m)\,L^u\ \le\ L^{m+u}.
\end{aligned}
\end{equation}
The inequalities are justified as follows.
\begin{itemize}
\item The first is the first item.
\item The second holds because $\min_{i\le k}x_i\le x_k$ and $R$ is nondecreasing.
\item The third is \eqref{5.1:eq-correlation-theorem-a} of the correlation theorem, read at
level $k=K'-u$, together with the monotonicity of $R$.
\item The fourth is the term of index $u$ in \eqref{5.1:eq-torus-floor-3}, multiplied by
$L^u$.
\item The fifth is \eqref{5.1:eq-torus-floor-a}, multiplied by $L^u$.
\end{itemize}

Both $s$ and $L^{m+u}$ are powers of $L$, and $s\le L^{m+u}$. Hence
$s\mid L^{m+u}\mid N_TL^{m+u}$, and $T^{(k)}$ is partitioned into cubes of side $s$. This
gives the partition into cubes of side $LM_2\check{R}_k$ compatible with all other
partitions of the lattice, as used in \cite{B-CMP119}. It likewise gives the partitions of
sides $M\check{R}_k$, $LM\check{R}_{k+1}$ and $L^2M_2\check{R}_{k+1}$.

The partitions are nested $L$-adically. All of them are taken from a common origin of
$T^{(k)}$; since $s\mid s'$ for any two of the sides $s\le s'$, both powers of $L$, each cube of
side $s'$ is then a union of cubes of side $s$.

Finally, by the first item and \eqref{5.1:eq-torus-floor-4},
\begin{equation*}
100M\check{R}_k\le C_{\mathbb{T}}\check{R}_k/L\le L^{m+u-1}<L^{m+u},
\end{equation*}
so the cube of side $100M\check{R}_k$ fits into $T^{(k)}$.

\emph{Fifth item.} Elements of $\mathrm{SU}(N)$ are unitary, so for $V\in\mathrm{SU}(N)$ we have
$\operatorname{tr}V^{-1}=\overline{\operatorname{tr}V}$. Hence
$\tfrac12\operatorname{tr}(V+V^{-1})=\operatorname{Re}\operatorname{tr}V$. The trace is
normalised so that the trace of the identity matrix is $1$, so $\operatorname{tr}V$ is the mean
of the eigenvalues of $V$, which lie on the unit circle. Therefore $|\operatorname{tr}V|\le 1$ and
\begin{equation*}
a_p(U)=1-\operatorname{Re}\operatorname{tr}U(\partial p)\ge 0.
\end{equation*}

Linearity of $A(c,U)$ in $c$ is immediate from its definition.

Under a gauge transformation, $U(\partial p)$ is replaced by $hU(\partial p)h^{-1}$, where
$h$ is the value of the transformation at the initial corner of $\partial p$. The trace is
invariant under conjugation, so each $a_p$ is gauge invariant, and hence so is each term of
$A(c,U)$.

Finally, $\mathrm{SL}(N,\mathbb{C})$ consists of complex matrices of determinant one. On such
matrices $U^{-1}$ is the adjugate of $U$, a polynomial in the entries of $U$. Hence $a_p$, and
with it $A(c,U)$, is a polynomial in the entries of the bond variables, and so is entire on
$\mathrm{SL}(N,\mathbb{C})$.

\emph{The closing example.} Take $m=0$ and $N_T$ prime to $L$. Then $N=N_T$ has no factor
$L$. The side $M$ is a positive power of $L$, so $M\nmid N$, and no partition of the torus
into cubes of side $M$ exists, however large $N_T$ is taken.
\end{proof}

\begin{definition}[The source-free scaffold and its strips]\label{5.1:def-source-free-scaffold}
Let $T$ be as in \cite[(0.1)]{B-CMP119}, and let $T^{(k)}$, $M$, $R(x)$ and
$\check{R}_k=R(\min_{i\le k}x_i)$ be as in Definition~\ref{5.1:def-torus-floor}. In this definition
$w=\sum_i z_i f_i$ is the complex source, $\zeta_k$ the running shift, and $g_k$ the real running
coupling; $N$ denotes Ba{\l}aban's number of preparatory stages, not the number of sites per side of
Definition~\ref{5.1:def-torus-floor}.

\smallskip
\noindent\textup{(a)} \emph{The scaffold.} All identities inserted and all choices made in the
expansion are those of $z=0$, taken at the real coupling $g_k$. These are: the decompositions of
unity and the thresholds and radii
\begin{equation}\label{5.1:eq-source-free-scaffold-a}
\begin{gathered}
\varepsilon_k,\quad \delta_k,\quad \delta_k',\quad \check{R}_k,\quad N,\quad p_i(g_k),\quad
\text{the radii of \cite[(2.28)]{B-CMP119}};\\
H,\ G,\ C,\ Q,\ \Delta^{(k)},\ C^{(k)};\qquad
B' = g_k C B - h\tilde{D}(g_k C B);
\end{gathered}
\end{equation}
the localisation domains, their labels and the two classes of components of
\cite{B-CMP122a}; the backgrounds, which are the
minimisers of the unweighted action $A$; the operators and the scaling $B'$, in
Ba{\l}aban's notation, displayed in the second
line of \eqref{5.1:eq-source-free-scaffold-a}, where $B$, $h$ and $\tilde{D}$ are as in the
scaling of \cite{B-CMP119} and $h$ is not a history; the trees and the graph averaging; the surrogate of
\cite[(1.100)]{B-CMP122a}, with its factor and its denominator; the relabelling of
\cite[(1.62)]{B-CMP122a}. Equally taken at $z=0$ and at the real $g_k$ are: the refinement
of the creating step of Definition~\ref{5.2:def-rim-refinement}, with its threshold
$\mathfrak{b}^{\flat}_k$, its collar
width $\varpi$, its cascade and its labels $I_j$ (not to be confused with the index ranges
$I(\mathfrak{a})$ of part (c)); the index $h$ and the parameter $r^{\flat}$ fixed
together with that refinement, and $\varpi_\gamma(g_k)$; the count
$\mathsf{N}(\cdot)$; the scaffold rule $\mathcal{S}(\cdot)$; the upper boxes $\mathfrak{b}^{\sharp}_j$
and $\mathfrak{X}^{\sharp}$ of Definition~\ref{5.2:def-rim-datum}. Only partitions already present
are used.

\smallskip
\noindent\textup{(b)} \emph{Components and strips (a reading convention).} The components of
$\Lambda_k^c$, hence of $Z$, are the unions of the equivalence classes of the closed cubes of the
partition composing $\Lambda_k^c$ under the relation generated by ``intersect or
touch'', contact at a corner included. For each $k$ the strip width, in $M$-cubes, is
\begin{equation}\label{5.1:eq-source-free-scaffold-1}
m_\sharp = m_\sharp(k) := \Big\lfloor \tfrac12\big(\check{R}_k - 1\big)\Big\rfloor ,
\end{equation}
the upper end of the window of widths allowed by the strip lemma. The set $Z^\sharp$ is $Z$ together
with $m_\sharp$ layers of $M$-cubes of the present lattice, and each component $C$ of $Z$ (here and
below the letter $C$ denotes a component of $Z$, not Ba{\l}aban's operator $C$ of
\eqref{5.1:eq-source-free-scaffold-a}) gives a strip
$C^\sharp$, namely $C$ together with $m_\sharp$ layers. We put $d_\sharp := d_{k,Z^\sharp}$. Ba{\l}aban's
$Z^{\sim}$, $X^{\sim}$, $Y_i^{\sim}$ in
\cite[(1.68)--(1.69), (1.71)--(1.72), (1.90)--(1.99)]{B-CMP122b}, and in the passage of
\cite{B-CMP122b} preceding them, are read as $Z^\sharp$, $X^\sharp$, $Y_i^\sharp$, with $d_{k,Z}$
replaced by $d_\sharp$ (the dictionary of the strip lemma). At
\cite[(2.30), (2.41)]{B-CMP119}, and at the corresponding later place of \cite{B-CMP122b}, the mark
$\sim$ keeps Ba{\l}aban's meaning, one layer of $M\check{R}_k$-cubes, and is written $\boxplus$.

\smallskip
\noindent\textup{(c)} \emph{Live arrested pieces and the arrest re-definition.}
The arrests, the
re-definition below, and the symbols it uses without fixing them here --- the standard factor
schedule $s^{(k')}_{\ell}$ (not the slot combination $s$), the attempt-space factors
$\varphi^{(n),(k)}_{\ell}$, the attempt spaces $\mathfrak{S}^{(k)}_n$, the spaces $\mathfrak{L}_l$,
the local quantities $Q_a(x)$, $a\le 7$, and their radii $r_a(\ell)$, the index $h_a$ and the stage count
$n^{+}$ of an arrest, and the
region $Z_h$ --- are those of the definition of arrests on which the arrest theorem rests.
An arrest $\mathfrak{a}$
is the attempt of the operation $\mathbf{R}$ of an earlier step $k_a$ on a component $X_{\mathfrak{a}}$,
interrupted after $n_a<N$ performed stages, or excluded after $N$ stages; its piece is
\begin{equation}\label{5.1:eq-source-free-scaffold-2}
W_{\mathfrak{a}} = X_{\mathfrak{a}} \cap \Omega^c_{k_a},
\end{equation}
and it is \emph{live} from the arrest to its dissolution. For such an arrest put
\begin{equation}\label{5.1:eq-source-free-scaffold-3}
\sigma_{\mathfrak{a}}(\ell) := \sum_{i\in I(\mathfrak{a})} 2^{-|\ell - i|},
\qquad I(\mathfrak{a}) := \{h_a+1,\dots,k_a+1\}.
\end{equation}
On every live piece $W_{\mathfrak{a}}$, from the arrest to its dissolution, the following re-definition,
called the arrest re-definition, is in force; the spaces and domains it prescribes are called
arrest spaces, and the space it puts in place of a given one is the arrest reading of
that space. Here $(\ell,\mathfrak{w},\mathcal{Q})(x)$ denotes the level, the weight and the family of
cube conditions at $x$. Below, $k'$ denotes the level at which a space is written while an arrest is
live, and $k$ the step to which an attempt space $\mathfrak{S}^{(k)}_n$ belongs.
\begin{enumerate}
\item[\textup{(0)}] \emph{Structure freeze.} Every space written after the arrest --- the standard
spaces $\tilde{U}_{k'}(X,\tilde\alpha)$, the spaces written after $\mathbf{T}$, the source spaces
$\tilde{U}^c_k(Y_4,\tilde\alpha)$ of $\mathbf{R}$, the attempt spaces of sibling arrests and of
enclosing attempts, the spaces $\mathfrak{L}_l$, the tubes, and the analyticity domains of newly born
terms --- reads at $x\in W_{\mathfrak{a}}$ the frozen $(\ell,\mathfrak{w},\mathcal{Q})(x)$ of the
innermost live
arrest holding $x$; by the third clause of the definition of arrests, the frozen level $\ell(x)$ is the
actual level of $x$. The constant third factor of the threshold product
$F\cdot \mathfrak{w}(x)\cdot c$, written $c$ in the definition of arrests, is set equal to $1$.
\item[\textup{(1)}] \emph{Factor.} With $\beta$ Ba{\l}aban's enlargement fraction of
\cite{B-CMP119}, evaluated below at $\beta=\beta_{\mathrm{pr}}=1/5$,
\begin{equation*}
F^{(k')}(x) := s^{(k')}_{\ell(x)} - \beta \sum_{\mathfrak{a}\ \mathrm{live},\ x\in W_{\mathfrak{a}}}
\sigma_{\mathfrak{a}}(\ell(x)).
\end{equation*}
The attempt spaces $\mathfrak{S}^{(k),\mathrm{arrest}}_n$ are $\mathfrak{S}^{(k)}_n$ at points of no live
piece; at a point $x$ of a live piece they carry the frozen structure and the factor
\begin{equation*}
\varphi^{(n),(k)}_{\ell(x)} - \beta \sum_{\mathfrak{a}\ \mathrm{live},\ x\in W_{\mathfrak{a}},\ k_a<k}
\sigma_{\mathfrak{a}}(\ell(x)),
\end{equation*}
pieces of the same class contributing no further term inside the run.
\item[\textup{(2)}] \emph{Spaces.} The space $\tilde{U}^{\mathrm{arrest}}_{k'}(X)$ is the set of
configurations satisfying
\begin{equation}\label{5.1:eq-source-free-scaffold-4}
Q_a(x) < F^{(k')}(x)\, \mathfrak{w}(x)\, r_a(\ell(x)) \quad (a\le 7),
\end{equation}
together with the cube conditions of $\mathcal{Q}(x)$ at weight $\mathfrak{w}(x)$; off live pieces it is
\cite[(2.34)--(2.39)]{B-CMP119} at $k'$, verbatim. The two standing readings of these spaces are
unchanged. These spaces are open and invariant under gauge transformations.
\item[\textup{(3)}] \emph{Analyticity domains.} A term born at an operation $\mathfrak{O}$ has as its
domain the arrest reading of its domain in Ba{\l}aban's papers:
\cite[(2.41)(ii)]{B-CMP119} with its
factor enlarged as prescribed by the enlargement statement of the induction in this chapter, the
enlarged factor on a live piece being
$F^{(k'+1)}(x)+\theta_S\beta 2^{-(k'+1-\ell)}$; \cite[(1.68)]{B-CMP122b}; and
$\mathfrak{D}^{\flat}_n$ with $\mathfrak{L}_n := \mathfrak{S}^{(k),\mathrm{arrest}}_n$. Frozen terms keep
their domains.
\item[\textup{(4)}] \emph{Dissolution.} At the completed $\mathbf{R}$ at step $k$ on the component of
$X_{\mathfrak{a}}$, where $W_{\mathfrak{a}}\subset Z_h\subset\Lambda$ by the definition of arrests, drop
$\mathfrak{a}$ and every
piece nested with it; on $\Lambda$ the standard structure of the new $\mathbb{B}_k$ applies.
\end{enumerate}
The arrest spaces read on live pieces are $\tilde{U}^{\mathrm{arrest}}_{k'}$ and
$\mathfrak{S}^{(k),\mathrm{arrest}}_n$.

\smallskip
\noindent\textup{(d)} \emph{The ceiling.} With $w_*(j) := \mathsf{L}_0^{2N_0(g_j)}$, and with
$\mathsf{T}_k$ the value at level $k$ of the logarithmic coupling variable in which the envelope
bounds for $w_*(j)$ and $\bar{\alpha}_j$ are stated,
\begin{equation}\label{5.1:eq-source-free-scaffold-c}
\gamma_0^{(k)} := 5\max_{j\le k} w_*(j)\,\bar{\alpha}_j ,
\end{equation}
and its majorant is
\begin{equation*}
\bar{\gamma}_0(\mathsf{T}) := 5(1+\beta_0)^2\,\Psi(\mathsf{T}),
\end{equation*}
where $\beta_0$ is the constant and $\Psi$ the function of the envelope bounds for $w_*(j)$ and
$\bar{\alpha}_j$.

\smallskip
\noindent\textup{(e)} \emph{Where the source enters.} The complex source $w$ enters at exactly four
places:
\begin{equation}\label{5.1:eq-source-free-scaffold-b}
\begin{aligned}
&\mathrm{(E1)}\quad \text{the Wilson weight } x_k(\cdot)-\zeta_k;\\
&\mathrm{(E2)}\quad \text{the exponential gauge-fixing coefficient of the block } y,\ \text{namely}\\
&\qquad\qquad a_k(y) := x_k - \zeta_k(p_y),\ \text{with constant }
\log\mathfrak{z}(a_k(y));\\
&\mathrm{(E3)}\quad \text{the terms};\\
&\mathrm{(E4)}\quad \text{the constants}.
\end{aligned}
\end{equation}
Here $\zeta_k(p_y)$ is the running shift at the plaquette $p_y$ attached to the block $y$.
\end{definition}

\begin{remark}[Components, strips and arrest spaces]
\cite{B-CMP122a} states:
\begin{quotation}
if in a component of $Z$ we finish the integrations for some $n$, because some new large field has
been introduced, then we leave the new terms $\mathbb{B}^{(n)}_k$ as a part of all boundary terms, and
it is easy to generalize the definition of the spaces in such a case for the next steps
\end{quotation}
The arrest re-definition is the generalisation of the definition of the spaces used in this chapter in
that case; it is the first hypothesis of the arrest theorem.

\smallskip
\noindent\emph{The components and the strips.} Two closed cubes of the partition at sup-distance zero
intersect; hence cubes in distinct classes of the relation in part (b) of
Definition~\ref{5.1:def-source-free-scaffold} are at positive sup-distance, and in
particular distinct components of $Z$ share no face. The collar $\boxplus$ of \cite{B-CMP119} is
$\check{R}_k$ cubes of side $M$ wide (Ba{\l}aban's $R_k$, read here as $\check{R}_k$), and the
$\boxplus$-collars of two components one $M\check{R}_k$-cube apart overlap. The window
$m_\sharp\le\lfloor\frac12(\check{R}_k-1)\rfloor$ makes the strips $C^\sharp$ pairwise non-touching
boxes, one about each component, lying inside $\boxplus$; this is the first clause of the strip lemma.
In \cite[(1.68)--(1.69), (1.71)--(1.72), (1.90)--(1.99)]{B-CMP122b} the symbols $Z^{\sim}$,
$X^{\sim}$, $Y_i^{\sim}$ are used here in the sense of (b), and every inequality of
\cite{B-CMP122b} is used in the form stated there. At the width
\eqref{5.1:eq-source-free-scaffold-1} the strip smallness condition ensures
$\alpha^{\sharp}\le\alpha_*$, by the strip proposition. Both ends of the window are scaffold functions
of the coupling $g_k$ of the run, like $N$ and $\check{R}_k$: the lower end $m_{\mathrm{lo}}(k)$ is
defined in the strip proposition from $\mathsf{K}^0+N(g_k)C_0^{\sharp}$, is dominated inside the
supremum of the strip smallness condition, and enters none of the conditions defining $\gamma_J$. The
uses of the convention (b) in the proof of the correlation theorem are listed in the strip proposition,
in the per-cell volume bound, and in a remark following the boundary lemma.

\smallskip
\noindent\emph{The arrest spaces.} By the lower bound for the arrest factors one has
\begin{equation}\label{5.1:eq-source-free-scaffold-5}
F^{(k')} \ge \varphi_{\mathrm{arrest}} := 1 - 4.51\,\beta = 0.098
\qquad\text{at } \beta=\beta_{\mathrm{pr}},
\end{equation}
and every attempt-space factor satisfies the same bound. The spaces $\mathcal{U}^{\flat}_j$ of this
chapter and the spaces \cite[(2.34)--(2.39)]{B-CMP119} are among the standard spaces and the spaces
written after $\mathbf{T}$ named in clause (0) of part (c) of
Definition~\ref{5.1:def-source-free-scaffold}. The five clauses (0)--(4) depend only on the real
scaffold at $z=0$: levels, weights, cube families, factors, and the status live or dissolved. They
therefore belong
to the scaffold of part (a) of that definition, independent of $z$ and of the slot variables $\xi_j$
and the complex coupling
slot $\mathfrak{q}$, as does the convention (b). By the lemma on the real
integration domain, under its two standing conditions, the real integration domain lies inside every
arrest space with margin $\frac12$. The sealing theorem carries the same proviso: the
arrest re-definition, jointly with one further condition stated there, is one of its
conjuncts.

Accordingly, in this chapter the spaces $\tilde{U}^c_k(Y_4,\tilde\alpha_0,\tilde\alpha_1)$ at the site
(L), the spaces $\mathcal{U}^{\flat}_j$, the tube of the response maps, and the spaces
$\tilde{U}^{(n)c}_k$ at the sites (T) and (P), in the sense of Definition~\ref{5.2:def-sites-units},
mean, at the points of a live piece, their arrest readings
$\tilde{U}^{\mathrm{arrest}}$; this applies
to the source spaces of live pieces and to the spaces written after the arrest, as in clause (0). The
analyticity domain $\mathfrak{D}_v$ of a term $v$ at a cell of a live piece foreign to $v$ is
Ba{\l}aban's
$U^c_j(X,\alpha_{\cdot,j})$ of \cite[(1.11)--(1.16)]{B-CMP109}, whatever the entry time.

\smallskip
\noindent\emph{How bounds reach the arrest spaces.} The arrest spaces are not
subfamilies of the standard
ones: their weights $\mathfrak{w}(x)=\mathsf{L}_0^{2(\cdot)}$ and their factors differ (the sealing
theorem checks no more than $F\cdot \mathfrak{w}(x)\cdot c\le 3w_*$).
No bound stated on a standard space is
therefore transported to an arrest space by inclusion. The bounds reach them in the
following four ways.
\begin{enumerate}
\item[\textup{(i)}] Every term born after an arrest is stated on arrest domains, by clause (3).
Its bounds, namely \cite[(1.68), (1.70)]{B-CMP122b}, \cite[(2.41)(ii)]{B-CMP119} and those of
$\mathfrak{D}^{\flat}_n$, are therefore bounds on the arrest reading by construction.
These domains are legitimate by clause (C) of the arrest theorem.
\item[\textup{(ii)}] Consider an older term $\tau$: a frozen term of $\mathfrak{F}(\mathfrak{a})$, a
frozen term of $\mathfrak{F}(\mathfrak{a}_1)$ for a sibling arrest $\mathfrak{a}_1$, a term of
$\mathbb{B}''_{k_a}$, or a term of the expansion of \cite{B-CMP119} born at a step $\le k_a$. By the
extended containment
lemma for older terms, applied to terms whose condition at $x$ contains the class domain (compare
the monotone reading below), the domain of $\tau$ contains on the live piece the present arrest
source space. The domain of $\tau$ exceeds
$\tilde{U}^{\mathrm{arrest}}_{k'}$ at $x$ by the margin $\rho_n(x;k')$ of that lemma, where
\begin{equation}\label{5.1:eq-source-free-scaffold-6}
\rho_n(x;k') \ge \beta\, 2^{-(k_a+1-\ell)} > 0 ,
\end{equation}
and by clause (D) of the arrest theorem
\begin{equation}\label{5.1:eq-source-free-scaffold-7}
\tilde{U}^{\mathrm{arrest}}_{k'} \subset \mathfrak{S}^{(k_a),\mathrm{arrest}}_{n^+}
\subset \mathfrak{S}^{(k_a),\mathrm{arrest}}_0 .
\end{equation}
\item[\textup{(iii)}] Let $\tau$ be an arrested left-over at the site (L), (T) or (P), and let $v$ be
the unit carrying $\tau$. For such $\tau$ the undilated half of the hypothesis
$\mathsf{R}_p\subset\mathfrak{D}_v$ of the amplitude lemma, at $\mathsf{t}=0$, is clause (A) of the
arrest theorem: the image, under the map of $\mathfrak{O}$, of the arrest
source space of
$\mathfrak{O}$ satisfies the condition of $\tau$ at $x$, with the room, the portion of it consumed, and
the ratio listed in the table of rooms. The same holds, at the case of the table of rooms that
clause (A) designates for them, for every term with a domain
as in the extended containment lemma, that is for the terms of the expansion of \cite{B-CMP119} born at
a step $\le k_a$, terms born later being covered by clause (C). Clause (A) is a statement on the map
of $\mathfrak{O}$ itself, and it is read with the
following monotone reading.

\smallskip
\textit{Monotone reading.} Let $x$ be a cell of a foreign live piece $W_{\mathfrak{a}'}$, where
$\mathfrak{a}'$ is the innermost arrest at $x$, and let $x$ be met by the localisation domain $X_v$ of a
unit $v$ of scale $j_v\le k_{a'}$. Let
$\mu(x)\ge j_v$ denote, in this reading only, the pre-attempt level, read off the regions $\Lambda_j$
(the unprimed ones) before
the attempt; there is an integer $\iota\ge 0$, an offset of levels, such that
\begin{equation*}
\ell(x) = \mu(x) + \iota, \qquad \mathfrak{w}(x) \le \mathsf{L}_0^{2\iota}.
\end{equation*}
Then the condition of $v$ at $x$ contains the stage-zero class domain of $\mathfrak{a}'$ at
$(\mu(x),1)$, by two clauses of the lemma on stage-zero class domains; so the
extended containment lemma applies at $x$. Clause (A) of the
arrest theorem at that case of the table of rooms --- by the extended containment
lemma, its part
(a) for the offset $\iota=0$ and its part (b) for $\iota\ge1$, and by part (a) of the arrest weight
lemma ---
places the image of the arrest source space, under the operation of the table of rooms to which
clause (A) refers, in
that class domain, within the room that the table assigns to that case at that operation.
Hence that image lies in $\mathfrak{D}_v$. No constant in this
reading depends on $k$, $k_{a'}$, $\iota$, $N_0$ or an age.

\smallskip
The $\mathsf{t}$-disc half of the same hypothesis at the site (L) is supplied by two cases of the
statement on the amplitude disc at the site (L): its case of arrested pieces and a second case of
that statement stated with it. They give the entries of the table of rooms for the operation just
named at halved far factors, under the halved-far-factor condition, the
pieces dissolved at this site (L) being taken from the sealing theorem.
\item[\textup{(iv)}] The scaffold maps, independent of $z$, acting on spaces that are not subfamilies of
Ba{\l}aban's --- the operators $H$ and $G$, the backgrounds, and $\mathsf{R}_p$ on the arrest
tube --- take their bounds from clause (D) of the arrest theorem: every arrest
space lies in the set
$\mathfrak{N}\subset\mathfrak{K}_P$ of that theorem (here $\mathfrak{N}$ is not a unit size). This
clause is read with $w_*$ pointwise: its comparison
$F\cdot \mathfrak{w}(x)\le w_*$, the check inside clause (D), is taken at
$\max\{w_*,\mathfrak{w}(x)\}$, and the ceiling is
$\gamma_0^{(k)}$ of
\eqref{5.1:eq-source-free-scaffold-c}; by the envelope bounds for $w_*(j)$ and $\bar{\alpha}_j$,
$\gamma_0^{(k)}\le\bar{\gamma}_0(\mathsf{T}_k)$. The frozen weight
$\mathfrak{w}_{\mathfrak{a}}=w_*(k_a)=\mathsf{L}_0^{2N_0(g_{k_a})}$ is formed with the coupling at the
arrest, while the present weight is $w_*=\mathsf{L}_0^{2N_0(g_k)}$. Compared at the present
coupling alone, the condition would therefore carry the age; and a ceiling given by one number over all
densities would be a floor on the coupling. The maximum over densities $j\le k$ in
\eqref{5.1:eq-source-free-scaffold-c} covers every frozen weight
$\mathfrak{w}_{\mathfrak{a}}=w_*(k_a)$ with
$k_a\le k$, the two cases of the weight comparison holding by definition at the
factor $5$. The majorant $\bar{\gamma}_0(\mathsf{T}_k)$ depends on $\mathsf{T}_k$ alone. Consequently
each of the
smallness conditions in which the ceiling enters, at the places where they are imposed, among them the
condition that a fixed constant multiple of the ceiling $\gamma_0^{(k)}$ be at most $\frac12$,
becomes a single inequality of degree $-\infty$, in the sense in which degrees are
fixed with the order of choice of the constants,
\begin{equation*}
\sup_{\mathsf{T}\ge\mathsf{T}_\gamma}\big[\cdots\bar{\gamma}_0(\mathsf{T})\cdots\big]\le \mathrm{const},
\end{equation*}
whose bracket is the left member of the condition with $\gamma_0^{(k)}$ replaced by
$\bar{\gamma}_0(\mathsf{T})$ and $\mathsf{T}_k$ by $\mathsf{T}$. Here
$\mathsf{T}_\gamma := \log\gamma^{-2}$, the value of the variable $\mathsf{T}$ at the coupling
$g=\gamma$.
\end{enumerate}
By (i)--(iv), every bound used on an arrest space is obtained on that space itself.
\end{remark}

\begin{definition}[Regions, labels and the expansion rule at the three sites]
\label{5.1:def-expansion-rule}
The objects and rules fixed in this definition are source-free scaffold declarations of the
same kind as the strip convention of the source-free scaffold
(Definition~\ref{5.1:def-source-free-scaffold}). They are made at $z=0$ and at the real
coupling $g_k$, and one and the same choice is used in every expansion that reads them. The
three sites of a step at which a bracket is expanded are written (T), (L) and (P): (T) is
the site of the operation $\mathbf{T}$, (P) that of the preparatory stages, and (L) the
localisation site of the operation $\mathbf{R}_k$. Sites and units are as in
Definition~\ref{5.2:def-sites-units}.

\emph{Returned pairs and strips.} A \emph{returned pair} is a pair $(Y,m)$ in which $Y$ is a
component that the operation $\mathbf{R}_m$, once completed, throws back as a domain of the
second class (Definition~\ref{5.1:def-source-free-scaffold}). A returned pair is \emph{live}
until the completed operation $\mathbf{R}$ on the component that holds it. The \emph{strip}
$Y^{\sharp}$ of a returned pair is formed by the strip convention of
Definition~\ref{5.1:def-source-free-scaffold}.

\emph{Cores.} For each component $C$ of $Z$ at $\mathbf{R}_k$ the \emph{core} is
$\mathsf{C}_C:=Z_{h'}^{\sim}\cap C$. Here the levels $h$, $h'$ and the set $\Lambda_C$ are
those attached to $C$ in the geometry statement of the relative small-field machine. By that
statement the core has the following properties:
\begin{itemize}
\item it is a box whose sides are at most $104L\check{R}_{h+1}$ level-$h$ units;
\item it is a union of cells of levels at most $h$;
\item the cells of $\mathsf{C}_C$ that touch $\mathsf{C}_C^c$ are of level $h$;
\item it satisfies the chain of inclusions
\begin{equation}\label{5.1:eq-expansion-rule-1}
\mathsf{C}_C\subset Z''_{h+1}\subset Z''_k\cap C\subset\Lambda_C .
\end{equation}
\end{itemize}

\emph{Labels and the remaining scaffold.} The labels and returned-pair sets
$(A_F,\mathcal{P}_F)$, the skeleton $\operatorname{sk}(\cdot)$, the corridors, the probes and
the old blocks are declared at the same coupling and by the same rule. The \emph{grains}
$\pi^{\flat}$ of the preparatory stages form a source-free partition. They are read only by
the resummation lemma of the preparatory stages
(Lemma~\ref{5.2:lem-animal-resummation}) and by three summation lemmas of the relative
small-field machine; they declare nothing and are no clause of the expansion.

\emph{Wrapped and plain units; the two parts.} For a unit $v$ write $\mathcal{P}_v$ for its
returned-pair set. The unit is \emph{wrapped} if $\mathcal{P}_v\neq\emptyset$, so that it
carries the strips of its returned pairs, and \emph{plain} if $\mathcal{P}_v=\emptyset$.
Following the division of boundary terms in \cite{B-CMP122b}, $v$ is of the \emph{first part}
if its localisation domain $X_v$ meets $Z$, and of the \emph{second part} if
$X_v\cap Z=\emptyset$.

\emph{Clause} (S). Clause (S) is the expansion clause of the relative small-field machine; we
repeat its four parts.
\begin{itemize}
\item[(T)] At site (T), for a wrapped term $F$, the bracket is expanded by the localisation
expansion of the relative small-field machine over the family $\sigma_0:=\pi^{T}$ of all
cells of the multiscale cover, the exempted family being empty. This replaces
the expansion of \cite[(3.6)--(3.7)]{B-CMP109} and the core $\bar S_v$. It is a choice that
departs from Ba{\l}aban's construction, where the family $\sigma_0$ at this place consists of
the cubes off a set $X_0\supseteq X$ \cite{B-CMP116,B-CMP119}.
\item[(L$_1$)] At site (L) of $\mathbf{R}_k$ the following is done for every wrapped unit,
whether or not its strips are live. In Ba{\l}aban's first operation (each of the at most four
brackets of a second-part unit, the one bracket of a first-part unit), the bracket is
expanded by the localisation expansion of the relative small-field machine in the decoupling
parameters $\mathsf{s}$ alone, over the family $\sigma_0$ of the cells disjoint with
$\Lambda^{\sim}$; the cells meeting $\Lambda^{\sim}$ are exempted.

This is a second choice departing from Ba{\l}aban's construction, and it has two halves. In
Ba{\l}aban's first operation one applies either the equality \cite[(1.59)]{B-CMP122b} or
\cite[(1.60)]{B-CMP122b}, and the
localisation operation is determined by the cubes $\Delta$ disjoint with $\square_0$, or with
$X$, and with $\Lambda^{\sim}$. Here $\square_0$, or $X$, is removed from the exempted set;
this is the first half. Ba{\l}aban's amplitude $t_{\square}$ is not taken; this is the second
half. That amplitude belongs to the bracket telescoped in the partition of unity
$\hat{\chi}_{\square}$ and estimated by Cauchy's formula in $t_{\square}$, as in
\cite[(3.4)--(3.7)]{B-CMP109} and \cite[(1.30)]{B-CMP116}.

The three later localisations run over Ba{\l}aban's own families, $\sigma_0$ being the family
of cells disjoint with $\Lambda$. No cell meeting $\Lambda$ is ever cut in a chain that a unit
reads.
\item[(L$_2$)] For each link index $t$ the set
\begin{equation}\label{5.1:eq-expansion-rule-2}
\begin{split}
Y^{\mathrm{tr}}_t:=\ &\operatorname{lab}_t\cup\{\text{components of } Z \text{ meeting }
\operatorname{lab}_t\}\\
&\cup\bigcup\bigl\{Y^{\sharp}:(Y,m)\in\mathcal{P}_F\ \text{live after}\ \mathbf{R}_k\bigr\}
\end{split}
\end{equation}
enters in two places, and nowhere else. It enters ($\alpha$) the domains $Y_4(\mathfrak{p})$
and $Y^{+}$ of the pieces, and ($\beta$) the reading sets $\mathcal{R}_{t-1}$. None of
Ba{\l}aban's determining sets $\mathbf{B}''_k(Y^{\mathrm{tr}}_t)$, $\mathbf{B}_1(\cdot)$,
$\mathbf{B}_k(\cdot)$ is formed. The representation \cite[(1.61)--(1.63)]{B-CMP122b} is not
substituted. Links two to four are the global layers of the carriers, the carrier pairs
$\mathrm{chain}(\mathfrak{x})$ of (b) below; Ba{\l}aban's reading of these links is not used.
\item[(P)] At site (P) the expansion is not changed: the preparatory expansion is kept as
it stands, and only the domain function is thinned. The resummation lemma of the preparatory
stages (Lemma~\ref{5.2:lem-animal-resummation}) is a resummation inside a proof and declares
nothing.
\end{itemize}

\emph{The rule at} (L). At site (L) every interior first-part unit ($X_v\subset Z$), plain
or wrapped, is expanded by parts (L$_1$) and (L$_2$) of clause (S). For such a unit no
partition $\hat{\chi}_{\square}$ is inserted and no amplitude $\mathsf{t}$ is taken. The
amplitude $\mathsf{t}$ of the first-stage bound by two contour integrals
(Lemma~\ref{5.2:lem-first-stage-bound}), which we call the amplitude lemma below, in its
fourth round only, the round on the pair formed by the third and fourth backgrounds of the
chain, is used on the plain first-part units with $K_v\setminus Z\neq\emptyset$. Accordingly
the expansion of a unit $v$ at (L) is fixed in exactly one way, by the case to which $v$
belongs:
\begin{equation}\label{5.1:eq-expansion-rule-a}
\begin{aligned}
&\text{(a)} &&v\ \text{wrapped, of either part:}\\
& &&\qquad\text{parts (L}_1\text{), (L}_2\text{) of clause (S);}\\
&\text{(b)} &&v\ \text{plain, of the first part, } X_v\subset Z:\\
& &&\qquad\text{parts (L}_1\text{), (L}_2\text{) of clause (S), at }
\mathcal{P}_v=\emptyset;\\
&\text{(c)} &&v\ \text{plain, of the first part, } K_v\setminus Z\neq\emptyset:\\
& &&\qquad\text{one } \mathsf{t}\text{-bracket, fourth round, on the pair formed}\\
& &&\qquad\text{by the third and fourth backgrounds of the chain;}\\
&\text{(d)} &&v\ \text{plain, of the second part, } X_v\cap Z=\emptyset:\\
& &&\qquad\text{the amplitude lemma on the four brackets.}
\end{aligned}
\end{equation}
The cases in \eqref{5.1:eq-expansion-rule-a} are specified as follows. The rows, and the
words \emph{long single}, \emph{tail} and \emph{Wilson piece}, refer to the table of plain
units at (L); the words \emph{exposed}, \emph{pointed} and \emph{compact} are those of
Definition~\ref{5.2:def-exposed-compact-units};
$p$ indexes the pieces $v_p$ that the amplitude lemma produces from a unit $v$.

(a) \emph{Wrapped units, of either part.} The pieces are
$\mathfrak{p}=(v;s;\mathsf{y}_1,\dots,\mathsf{y}_4)$, as in clause (c) of the proposition on
wrapped pieces, with $s$ as in that clause. Their sum is
$\varepsilon^{\mathrm{rel}}(k)$.

(b) \emph{Plain interior first-part units} ($X_v\subset Z$; all rows of the table of plain
units). Parts (L$_1$) and (L$_2$) are read at $\mathcal{P}_v=\emptyset$. At the first
operation this is a choice departing from Ba{\l}aban's construction. Link one is
Ba{\l}aban's one bracket, expanded by the fundamental theorem of calculus in
$\mathsf{s}^{(1)}$ alone over the cells disjoint with $\Lambda^{\sim}$. Links two to four run
over the cells disjoint with $\Lambda$, which are Ba{\l}aban's own families. The pieces are
$\mathfrak{p}=(v;\mathsf{y}_1,\dots,\mathsf{y}_4)$. Write $\mathfrak{W}$ for the parameter
polydisc, whose points we write $\mathfrak{x}=(\mathfrak{x}_1,\mathfrak{x}_2,\vec{\mathsf{s}})$,
$\mathfrak{x}_1$ being the variable of the holomorphic site map, $\mathfrak{x}_2$ the second
carrier parameter and $\vec{\mathsf{s}}$ the decoupling parameters; write $f_v$ for the
function of the unit
$v$ on $\mathfrak{W}$, $\mathrm{chain}(\mathfrak{x})$ for the carrier pair at
$\mathfrak{x}\in\mathfrak{W}$, and $\mathcal{U}(\mathfrak{x})$ for its configuration slot.
For a vector $\vec{\mathsf{s}}$ of decoupling parameters alone, $\mathrm{chain}(\vec{\mathsf{s}})$
denotes the carrier pair at the point of $\mathfrak{W}$ with these decoupling parameters, the
other two components being held fixed; $\vec{0}$ and $\vec{1}$ are the vectors all of whose
entries are $0$, respectively $1$. The all-empty tuple is the \emph{value}
$\mathrm{val}_v=f_v(\mathrm{chain}(\vec{0}))$. The sum is taken in the summation lemma for
plain interior units, whose setting is Definition~\ref{5.2:def-plain-interior-units}.

\emph{Reading of smooth far cores.} For the smooth interior far cores of the first three rows
of the table of plain units, $f_v(\mathfrak{x})$ is read at the box minimiser of the level-$p_v$
averages of the configuration slot $\mathcal{U}(\mathfrak{x})$ of $\mathrm{chain}(\mathfrak{x})$
on $Q_v$; we write $\mathrm{ext}_v(\mathfrak{x})$ for this smoothed slot.
Rough cores, and every other unit, are read on $\mathrm{chain}(\mathfrak{x})$ itself. There is
one weight $\bar{\mathsf{w}}_v$, taken at the fine weight $\mathsf{L}_0^{2\mathfrak{f}}$, where
$\mathfrak{f}$ is the fine exponent fixed in the clause of the summation lemma for plain
interior units that states this reading. The reading is made this way
because in Ba{\l}aban's construction the configuration $\mathbf{U}$ enters these terms through
$M^{\cdot}(\mathbf{U})$, the block averages of $\mathbf{U}$ \cite[p.~375]{B-CMP122b}. The same
reading applies to the smooth far cores of the plain second-part units of (d) and of their
$\mathsf{t}$-pieces, on the joint polydisc, with $p_v:=n-2-s^{(9/8)}(\bar{\mathsf{w}}_v)$,
where $n$ is the zone of the unit and $s^{(9/8)}(\cdot)$ is the function of
Definition~\ref{5.2:def-smoothing-depth}; their subtracted copies are read at the same
$\mathrm{ext}_v$. In the display of the clause of the summation lemma for plain interior
units that states this reading, the background datum $\mathsf{V}$ is evaluated at the
smoothed slot $\mathrm{ext}_v(\mathfrak{x})$. The real values are unchanged.

\emph{Subtracted copies.} At (L) Ba{\l}aban adds certain terms to complete the effective
action, subtracts them from the remaining expression, and subtracts them from the
corresponding terms $V(Y)$ \cite{B-CMP122b}. These copies are members of the value family
$\mathcal{V}$, and they are treated in the copy lemma. When
they are smooth interior far cores they are read at the same $\mathrm{ext}_v$ with
$(n,\mathsf{W})=(k,3)$, that is, at zone $k$ and weight $\bar{\mathsf{w}}=3$; they lie in
the level-$k$ collar of $C$. Rough copies, and every other copy, are read at the
configuration slot $\mathcal{U}(\mathfrak{x})$ of $\mathrm{chain}(\mathfrak{x})$ itself, the
configuration $\mathbf{U}$ of the first component $\mathfrak{x}_1$. The
declaration is one and the same for the copy and for its original. It is made for the same
reason as the reading above: in \cite{B-CMP122b} the configuration $\mathbf{U}$ enters these
terms through $M^{\cdot}(\mathbf{U})$.

\emph{Kept members of the inner returned pairs.} The domain list of this definition, and that
of the disc-width rule, carries the kept members of the inner returned pairs, written
$\mathfrak{q}_C$ and $V^{\sharp}(C^{\sharp};\cdot)$, as terms on which the membership clause of
the inherited space $\mathfrak{I}_n$ is imposed. They are not terms of the inherited family.

(c) \emph{Plain first-part units with} $K_v\setminus Z\neq\emptyset$. The amplitude lemma is
applied with one $\mathsf{t}$-bracket, in the fourth round, on the pair formed by the third
and the fourth backgrounds of the chain \cite{B-CMP122b}.

\emph{Fourth-row and long units.} This paragraph concerns the fourth-row units of (c) and of
(d), and the long singles and tails of (c) and of (d).
\begin{itemize}
\item The following are kept: the pieces $v_p$ of the amplitude lemma, their classes under
this definition, and the link-one domain $Y(K_p)$ of each piece, equal to
$Y(\operatorname{lab}(\mathfrak{b}))$ for the bundle $\mathfrak{b}$ containing $p$ introduced
in \eqref{5.1:eq-expansion-rule-4} below.
\item Their filing at (L) is by $\square$-bundles
(Definition~\ref{5.2:def-boundary-bundles}): they are read in bundles and summed by
the summation of boundary terms at finer zones. For the fourth-row units and the long units
of (c) the bundle convention of that summation is used, with the traced corridor hull
\eqref{5.1:eq-expansion-rule-4} of (c${}'$) below and $\mathcal{R}_1:=Y_0(\mathfrak{b})$ as
in \eqref{5.1:eq-expansion-rule-5}. For the long units of (d) the long-bundle convention of
that summation is used, with the corridor hull \eqref{5.1:eq-expansion-rule-3} of (c${}'$)
and the reading set $\mathcal{R}_1(p):=Y_0(p)$, where $Y_0(p)$ is as in
Definition~\ref{5.2:def-boundary-bundles}.
\item On these units the hypothesis of the amplitude lemma is read at $\mathsf{t}=0$, and
its gain $\mathsf{w}_L^{-1}\eta_p^{1/2}$ is not used. Their smallness is the walk factor
$e^{-(9\check{R}_k-2)}$, and $e^{-\frac12(9\check{R}_k-2)}$ for (c); it is not the
amplitude $\mathsf{t}$.
\item Only the filing of the non-empty continuations of these rows is changed, source-free,
by the traced corridor hull and $\mathcal{R}_1$ as in (c${}'$) below. The exponent is not
changed.
\end{itemize}

(c${}'$) \emph{The exterior class $\mathfrak{L}^{\mathrm{ex}}$.} After the bracket of (c),
two terms are localised at links two to four over Ba{\l}aban's own families: the first member
of the term, that is $f_v$ evaluated at the third background of the chain, and the
$\mathsf{t}$-carrying $Y_0$-term $v_p$ that the bracket creates. The family used is the family
$\sigma_0$ of cells disjoint with $\Lambda$. In
\cite{B-CMP122b}, after the bracket, all terms depend on the background field $U^0_k$, each
has a localisation domain $X$ intersecting $Z$, and the localisation operation is applied to
Ba{\l}aban's function $H'_k$ with the cubes disjoint with $\Lambda$. These terms are read by
part (L$_2$) of clause (S) at $\mathcal{P}_v=\emptyset$. Their sets $Y^{\mathrm{tr}}_t$ enter
the domains of the pieces and the reading sets, and nothing else. No determining set is
formed, and the representation \cite[(1.61)--(1.63)]{B-CMP122b} is not substituted.

The pieces of $v$ are $(v;\mathsf{y}_2,\mathsf{y}_3,\mathsf{y}_4)$. Their link-one label is
\begin{itemize}
\item $\operatorname{lab}_1:=K_v\cup\mathfrak{h}_v$ for an \emph{exposed} unit (short singles
of the second, third and boundary rows; Wilson pieces);
\item $\operatorname{lab}_1:=K_v$ for the other units.
\end{itemize}
The pieces of a pointed $p$ are $(p;\mathsf{y}_2,\mathsf{y}_3,\mathsf{y}_4)$, with the
following link-one labels.
\begin{itemize}
\item For a \emph{compact} pointed $p$,
\begin{equation}\label{5.1:eq-expansion-rule-3}
\operatorname{lab}_1:=K_p\cup\mathfrak{h}_v,\qquad
K_p:=[\operatorname{St}_v\cup\mathsf{y}]_k\cup\operatorname{Hub}(\mathsf{y}).
\end{equation}
Here $K_p$ is the $\pi_k$-hull of the cells of $Y_0(p)$, with corridors for the never-cut
blocks $\Lambda^{\sim}_C$ in place of their interiors. This is the \emph{corridor hull}.
\item For the $p$ of the long singles and tails of (d), $\operatorname{lab}_1:=K_p$.
\item For the $p$ of a fourth-row unit of (c) or (d), and of the long singles and tails of
(c), the label is read per bundle $\mathfrak{b}=(v;s;\tilde{\mathsf{y}})$. The bundle
collects the pointed pieces $v_p$, $p=(\square,\mathsf{y};s)$, that have the common datum $s$
and whose $\mathsf{y}$ has the trace $\tilde{\mathsf{y}}$ described below (notation as in
\eqref{5.1:eq-expansion-rule-7} below):
\begin{equation}\label{5.1:eq-expansion-rule-4}
\operatorname{lab}_1:=\operatorname{lab}(\mathfrak{b}):=K_v\cup[\tilde{\mathsf{y}}]_k
\cup\operatorname{Hub}(\tilde{\mathsf{y}})\cup\operatorname{Att}(v,\tilde{\mathsf{y}}).
\end{equation}
This is the forest of the corridor hull on the trace $\tilde{\mathsf{y}}$ of $\mathsf{y}$ in
the unit's own grain, whose level is $j_v$; we call it the \emph{traced corridor hull}. Here
$\operatorname{Att}(v,\tilde{\mathsf{y}})$ is the attachment set of $\tilde{\mathsf{y}}$ to
$v$ fixed with the corridor-hull convention in Definition~\ref{5.2:def-boundary-bundles}. The
label comes with the reading set
\begin{equation}\label{5.1:eq-expansion-rule-5}
\begin{split}
\mathcal{R}_1(p):=Y_0(\mathfrak{b}(p)):=\ &\bigcup\operatorname{St}^{j_v}_v
\cup\bigcup\tilde{\mathsf{y}}\\
&\cup\{\text{the hubs of } \mathfrak{H}_v \text{ and those sharing a wall with }
\tilde{\mathsf{y}}\},
\end{split}
\end{equation}
where $\mathfrak{H}_v$ is the family of hubs of $v$ fixed in the same definition.
\end{itemize}

\emph{Properties of the bundle filing.} The bundles are a resummation inside a proof, in the
bundle filing lemma and its corollary.
\begin{itemize}
\item The classes fixed above and the pieces $v_p$ are unchanged by the bundling.
\item There is one $Y(\operatorname{lab}_1)$ per bundle.
\item The non-empty continuations $(\operatorname{lab}_t,\mathcal{R}_t,Y_4)$, $t=2,3,4$, are
assigned per bundle, in the bundle corollary: they are functions of
$(v,\tilde{\mathsf{y}},\mathsf{y}_2,\mathsf{y}_3,\mathsf{y}_4)$ only. Hence each filed sum
$\sum_{Y_4(\mathfrak{p})=Y_4}\mathfrak{p}$ is a sum of whole bundles.
\item This is a declared source-free re-filing in which every exponent is unchanged. The sum
$\sum_{\mathfrak{p}}\mathfrak{p}$ is the same function under any filing rule, and
$\|\cdot\|_{a,1}$ is bounded for this filing by the bundle corollary.
\item $\mathcal{R}_1(p)\supseteq Y_0(p)$ is an over-cover. It is admissible because the strip
theorem accepts any family of numbers (or of functions, in the sup norm), together with the
source-localisation identity, by which the extra tuples vanish. The boundary layer of the
hubs is read inside the bundle corollary and is not added to $\mathcal{R}_1$.
\end{itemize}

\emph{Tuples of} (c${}'$). The all-empty tuple $(v;\emptyset,\emptyset,\emptyset)$ is
$\mathrm{val}_v$, in the value family $\mathcal{V}$; the all-empty tuple
$(p;\emptyset,\emptyset,\emptyset)$ is the piece $v_p$ of the amplitude lemma,
as filed in the three members of $\varepsilon_{\mathrm{br}}(k)$ in which the pieces $v_p$
themselves are estimated; we call these the three $\theta'$-members, the name referring to the
ratio $\theta'(g_k)=1/\mathsf{w}_L$ of the amplitude lemma. The member
$\varepsilon^{\mathrm{ext},\mathsf{t}}(k)$ below is not among them. The
non-empty tuples form the class
$\mathfrak{L}^{\mathrm{ex}}=\mathfrak{L}^{c}\sqcup\mathfrak{L}^{Y_0}$. They are summed by the
summation lemma for exterior pieces (Lemma~\ref{5.2:lem-exterior-pieces}) into
\begin{equation}\label{5.1:eq-expansion-rule-6}
\varepsilon^{\mathrm{ext}}(k)=\varepsilon^{\mathrm{ext},\sigma}(k)
+\varepsilon^{\mathrm{ext},\mathsf{t}}(k),
\end{equation}
the sixth member of $\varepsilon_{\mathrm{br}}(k)$. In \cite{B-CMP122b} these are the terms
whose domains $Y'_i$ connect $Y_1\setminus Z$ with $\Lambda$. The condition
$Y_1\setminus Z\neq\emptyset$ is exactly what makes the unit exterior.

\emph{The $Y_0$-terms.} The brackets of (c) and of (d) create $Y_0$-terms. Ba{\l}aban regards
them as boundary terms connected with the corresponding components of $Z$, to be treated
later together with the other terms of the first part \cite{B-CMP122b}. These terms enter at
links two to four only. They take no second bracket at the next first operation: there is
one amplitude $\mathsf{t}$ per term presented to the first operation, and after it all terms
depend on the background field $U^0_k$. Their non-empty tuples at links two to four make up
$\mathfrak{L}^{Y_0}$.

(d) \emph{Plain second-part units} ($X_v\cap Z=\emptyset$). The amplitude lemma is applied on
the four brackets \cite{B-CMP122b}. The $Y_0$-terms so created are those of (c${}'$), taken
at links two to four.

\emph{Long singles and tails of} (d). The pieces $v_p$ of the amplitude lemma and their
classes are kept. At (L) they are read in $\square$-bundles
\begin{equation}\label{5.1:eq-expansion-rule-7}
\mathfrak{b}=(v;s;\mathsf{y})=\sum_{\square}v_{(\square,\mathsf{y};s)},\qquad
\mathsf{y}\subset T_v ,
\end{equation}
at the grain $\pi^{\mathfrak{s}}$, where $T_v$ is the set fixed for $v$ in
Definition~\ref{5.2:def-boundary-bundles}; this is the case of the one bundle convention of
that definition in which the
trace is the identity. The hypothesis of the amplitude lemma is read at $\mathsf{t}=0$ only,
and its gain is not used. The corridor hull of (c${}'$) and $\mathcal{R}_1:=Y_0(p)$ are
taken. These bundles are summed by the summation of boundary terms at finer zones,
together with the long units of (c) and the fourth-row units.

The quantities $\mathsf{W}$, $C_A'$, the fattened block $\tilde{\square}$, $\eta_p$ and
$\rho_A$ are read, among the units of (c) and (d), only by the compact pointed $p$.

\emph{The compact read set.} For the compact pointed $p$ of (c) and (d) at (L), the never-cut
read set is
\begin{equation}\label{5.1:eq-expansion-rule-8}
\mathfrak{r}_v:=\bigcup\operatorname{St}(\Delta_{y_v}).
\end{equation}
This is the star of the anchor cell, taken in own-scale cells of levels $n-1,n,n+1$, where
$n$ is the zone of $v$. It satisfies
\begin{equation*}
\#\operatorname{St}\le c_S:=(5L)^d,\qquad
\#\partial\operatorname{St}\le 2dL^{d-1}c_S .
\end{equation*}
This is a choice departing from Ba{\l}aban's construction. There the exempt set is the cube
$\square_0=\square^{\sim5}$, the operation $\sim$ being defined in terms of $M$-cubes of the
$L^{-n}$-lattice, and the expansion is taken with $\sigma_0$ the cubes disjoint with
$\square_0$ and $\Lambda^{\sim}$ \cite[p.~372, after (1.59)]{B-CMP122b}.

The exempt family of Ba{\l}aban's expansion is at our disposal: by the lemma on the choice of
$\sigma_0$ any subfamily may be taken, with the same constants. The identity
\eqref{5.1:eq-expansion-rule-9} below holds for every choice. Everything the unit reads lies
inside $\mathfrak{r}_v$, by the inclusion clause of the exempt-star lemma
(Lemma~\ref{5.2:lem-exempt-star}): the box of the cube
condition, $Q_v\subset\Omega_{n-1}$, a Wilson block, and the constituents of a rough core. The
only price of the choice is that the counting constant of the read set becomes $c_S$. It also
makes true by definition that the read set consists of own-scale cells of levels $n-1,n,n+1$
only.

\emph{Consequences.}
\begin{enumerate}
\item[(1)] No piece carries two auxiliary discs. The $\mathsf{t}$-disc of the amplitude
lemma lives on (c) and (d). The disc in $\hat\lambda$ of the bound for deep units lives on
the units of (a) and (b) of the deep class $\mathfrak{L}^{\mathrm{dp}}$.
\item[(2)] Clause (S) changes the localisation of terms and never the exponent. For every
unit $v$ whose bracket is expanded by the localisation expansion, whatever the choice of the
exempted cells,
\begin{equation}\label{5.1:eq-expansion-rule-9}
\sum_{\mathfrak{p}}\mathfrak{p}+\mathrm{val}_v
=f_v\bigl(\mathrm{chain}(\vec{1})\bigr)
\end{equation}
holds exactly. The exponent estimate with which the summation of the pieces opens, and
Lemma~\ref{5.2:lem-exempt-star}, are not affected.
\item[(3)] For (a) and (b), the pieces satisfy the membership condition of the source class
at (L). This condition is implied by two covering conditions, stated on the product polydisc
$|\mathsf{s}^{(t)}(\Delta)|\le e^{\kappa_1}$ of the decoupling parameters of all four links.
The covering lemma that accompanies the bound for deep units is read here with link one
included: its clause (a), at the fourth background of the chain, covers the pieces in which
only $\mathsf{y}_1$ is non-empty, which part (L$_1$) creates.
\item[(4)] What (b) gives up is the gain $\mathsf{w}_L^{-1}$ of the amplitude lemma on
interior units. Their exterior pieces are small by $e^{-(\check{R}_k-2)}$, by the
exterior-piece sum for plain units (Lemma~\ref{5.2:lem-plain-exterior-sum}). Their interior
pieces are filed with the values at the price $e^4$, by the volume sum of one unit's pieces
(Corollary~\ref{5.2:cor-unit-volume-sum}) and the volume bound.
\item[(5)] \emph{Exhaustiveness.} At (L) every piece of every unit lies in exactly one of the
following:
\begin{equation}\label{5.1:eq-expansion-rule-10}
\begin{gathered}
\mathfrak{v}^{\mathrm{val}},\qquad \text{the three }\theta'\text{-members},\qquad
\varepsilon^{\mathrm{rel}},\\
\varepsilon^{\mathrm{int}},\qquad \text{the volume bound},\qquad
\varepsilon^{\mathrm{ext}}.
\end{gathered}
\end{equation}
These receive, respectively:
\begin{itemize}
\item $\mathfrak{v}^{\mathrm{val}}$: the all-empty tuples of (b) and (c), and of the
first-part wrapped units of (a);
\item the three $\theta'$-members: the $\mathsf{t}$-pieces $v_p$ of (c), (d) at the
all-empty tuple of links two to four;
\item $\varepsilon^{\mathrm{rel}}$: every non-empty tuple of (a);
\item $\varepsilon^{\mathrm{int}}$: the non-empty tuples of (b) whose filing domain
leaves $Z$;
\item the volume bound: the non-empty tuples of (b) that stay inside $C$;
\item $\varepsilon^{\mathrm{ext}}$: the tuples of
$\mathfrak{L}^{\mathrm{ex}}=\mathfrak{L}^{c}\sqcup\mathfrak{L}^{Y_0}$ of (c${}'$).
\end{itemize}
This is the table of members of the bracket bound.
\end{enumerate}
\end{definition}

\begin{proof}
\emph{The cases (b) and (c) are complementary.} Let $v$ be a plain first-part unit. Its
domain $X_v$ is connected, and $K_v=[X_v]_k$ is the union of the $\pi_k$-cubes that meet
$X_v$. The partitions nest, so every cell of $X_v$ lies in a $\pi_k$-cube, and
$X_v\subset K_v$.

Suppose first that $X_v\subset Z$. Being connected, $X_v$ lies in one component $C$ of $Z$.
Every $\pi_k$-cube of $K_v$ meets $X_v\subset Z$. Since $Z$ is a union of $\pi_k$-cubes,
such a cube is contained in $Z$. It meets $X_v\subset C$, and the components of $Z$ are the
classes of intersecting or touching (Definition~\ref{5.1:def-source-free-scaffold}); so the
cube lies in $C$. Hence $K_v\subset C$, and in particular $K_v\setminus Z=\emptyset$.

Conversely, suppose $K_v\setminus Z=\emptyset$, that is $K_v\subset Z$. Then
$X_v\subset K_v\subset Z$.

Thus $X_v\subset Z$ holds if and only if $K_v\setminus Z=\emptyset$, and then $K_v$ lies in
the component $C$ of $Z$ that contains $X_v$. Every plain first-part unit therefore lies in
exactly one of (b) and (c).

\emph{The cases (a)--(d) are mutually exclusive and cover all units.} A unit is either
wrapped, and then it is in (a), or plain. A plain unit is of the first part or of the second
part ($X_v\cap Z=\emptyset$), and these alternatives exclude each other. The second-part plain
units form (d). The first-part plain units split into (b) and (c) by the complementarity just
proved. So the rule \eqref{5.1:eq-expansion-rule-a} assigns to each unit at (L) exactly one
expansion.

\emph{Consequence} (1). By the rule at (L), an interior first-part unit, plain or wrapped,
is expanded by parts (L$_1$) and (L$_2$) of clause (S) with no $\hat{\chi}_{\square}$ and no
$\mathsf{t}$. By part (L$_1$), no wrapped unit takes Ba{\l}aban's amplitude at the first
operation. Hence no piece of (a) or (b) carries the $\mathsf{t}$-disc, which occurs only in
(c) and (d). The disc in $\hat\lambda$ of the bound for deep units is used only for the
units of (a) and (b) of the deep class $\mathfrak{L}^{\mathrm{dp}}$. Since (a)--(d) are
mutually exclusive, no piece carries both discs.

\emph{Consequence} (2). Let $v$ be a unit whose bracket is expanded by the localisation
expansion. By the source-localisation identity, the sum of all its tuples equals the value of
$f_v$ at $\mathrm{chain}(\vec{1})$, for any choice of the exempted cells; this applies
to the choices of clause (S) and of the compact read set alike. The all-empty tuple is
$\mathrm{val}_v$, so \eqref{5.1:eq-expansion-rule-9} holds. The choices of clause (S) and of
the compact read set only change which cells are exempted. The left side of
\eqref{5.1:eq-expansion-rule-9} is therefore a re-arrangement of the same function, so the
localisation of terms changes and the exponent does not. For the same reason the exponent
estimate with which the summation of the pieces opens and Lemma~\ref{5.2:lem-exempt-star} are
not affected: neither depends on the choice of exempted cells made by clause (S).

\emph{Consequence} (3). For (a) and (b), the membership condition of the source class at (L)
is implied by the two covering conditions. These are stated on the product polydisc
$|\mathsf{s}^{(t)}(\Delta)|\le e^{\kappa_1}$, which is the domain of the decoupling
parameters of all four links under clause (S). Part (L$_1$) creates pieces in which only
$\mathsf{y}_1$ is non-empty. They are covered by clause (a) of the covering lemma that
accompanies the bound for deep units, at the fourth background of the chain, that lemma being
read with link one included.

\emph{Consequence} (4). On an interior unit of (b) no $\mathsf{t}$ is taken, by the rule at
(L), so the gain $\mathsf{w}_L^{-1}$ of the amplitude lemma is not available there. Its
exterior pieces are bounded by $e^{-(\check{R}_k-2)}$ in
Lemma~\ref{5.2:lem-plain-exterior-sum}. Its interior pieces are filed with the values, at the
price $e^4$, by Corollary~\ref{5.2:cor-unit-volume-sum} and the volume bound.

\emph{Consequence} (5). By the mutual exclusiveness of (a)--(d), it is enough to check that
within each case every piece is filed in exactly one entry of
\eqref{5.1:eq-expansion-rule-10}.

In case (a), consider first a first-part wrapped unit. Its all-empty tuple is a value by
clause (b) of the proposition on wrapped pieces, and so is a member of $\mathcal{V}$, which
contains the boundary terms born in $\mathbf{R}$ together with their wrapped descendants
born in $\mathbf{T}$; it is filed in $\mathfrak{v}^{\mathrm{val}}$. A second-part wrapped
unit has no all-empty tuple, since $\mathsf{y}_1\neq\emptyset$ by clause (c) of the same
proposition. Every non-empty tuple of (a) is filed in $\varepsilon^{\mathrm{rel}}$.

In case (b), the all-empty tuple is $\mathrm{val}_v$, filed in
$\mathfrak{v}^{\mathrm{val}}$. A non-empty tuple whose filing domain leaves $Z$ is filed in
$\varepsilon^{\mathrm{int}}$. One that stays inside $C$ is filed with the values, under the
volume bound together with Corollary~\ref{5.2:cor-unit-volume-sum}; there it is absorbed in
the supremum that defines $\|\cdot\|_a$ and is assigned no filing domain $Y_4$.

In case (c), the all-empty tuple is $\mathrm{val}_v$, filed in $\mathfrak{v}^{\mathrm{val}}$.
The $\mathsf{t}$-pieces $v_p$ at the all-empty tuple of links two to four are filed in the
three $\theta'$-members. The non-empty tuples of links two to four lie in
$\mathfrak{L}^{\mathrm{ex}}$ and are filed in $\varepsilon^{\mathrm{ext}}$, by (c${}'$).

In case (d), the pointed pieces $v_p$ are the $Y_0$-terms created by the four brackets. At
the all-empty tuple of links two to four they are filed in the same three $\theta'$-members.
They enter at links two to four only, and their non-empty tuples
$(p;\mathsf{y}_2,\mathsf{y}_3,\mathsf{y}_4)$ form
$\mathfrak{L}^{Y_0}\subset\mathfrak{L}^{\mathrm{ex}}$ by (c${}'$); they are filed in
$\varepsilon^{\mathrm{ext}}$.

In each case the entries are distinguished by conditions that exclude one another: the tuple
is all-empty or not; a non-empty tuple of (b) leaves $Z$ or stays inside $C$. This is the
table of members of the bracket bound, and the filing above is exhaustive and without
repetition.
\end{proof}

\begin{definition}[Analyticity spaces, enlarged radii and totals]\label{5.1:def-analyticity-spaces}
Let $j\ge 0$ be a scale, $X\in\mathfrak{D}_j$ a localisation domain of scale $j$
\cite[p.~254]{B-CMP109}, and write $x_j=g_j^{-2}$.

\emph{Spaces.} With Ba{\l}aban's complex configuration spaces $U^{c}_{j}(X,\cdot,\cdot)$ of
\cite[(1.11)--(1.16)]{B-CMP109}, put
\begin{equation}\label{5.1:eq-analyticity-spaces-1}
\mathcal{U}_j(X):=U^{c}_{j}(X,\alpha_0,\alpha_1),\qquad
\mathcal{U}^\flat_j(X):=U^{c}_{j}(X,\alpha_{0,j},\alpha_{1,j}).
\end{equation}
Here $\alpha_0,\alpha_1$ are Ba{\l}aban's absolute radii, which do not depend on the coupling. As in
\cite[(2.28)]{B-CMP119}, $\alpha_{i,j}:=g_jC^{(2.28)}_i(\log x_j)^{q_i}$ for $i=0,1$; Ba{\l}aban
requires only that $q_0,q_1$ be integers greater than $1$ and that $C^{(2.28)}_0,C^{(2.28)}_1$ be
sufficiently large positive numbers \cite{B-CMP119}. We write $C^{(2.28)}_0,C^{(2.28)}_1$ for these
two constants of \cite[(2.28)]{B-CMP119}; they are distinct from the constants $C_0,C_1$ among
Ba{\l}aban's auxiliary parameters $\mathfrak{p}$. The second condition in
\eqref{5.1:eq-analyticity-spaces-a} below constrains $C^{(2.28)}_1$ in terms of $C^{(2.28)}_0$ and
$M$. We take a common exponent $q:=q_0=q_1$, so that
$\alpha_{i,j}=\alpha_i(g_j)=g_jC^{(2.28)}_i(\log x_j)^{q}$ with
$\alpha_i(g):=gC^{(2.28)}_i(\log g^{-2})^{q}$ for
$0<g\le\gamma_J$, and we write $\alpha_{\cdot}(g)$ for the pair $(\alpha_0(g),\alpha_1(g))$. We
impose the lower bounds
\begin{equation}\label{5.1:eq-analyticity-spaces-a}
q:=q_0=q_1\ \ge\ p_0+\bigl(1+d+d\beta_0^{\mathrm{pr}}\bigr)r_{\mathrm{pr}}+1,
\qquad C^{(2.28)}_1\ \ge\ 2B_3^2\,O(1)\,M\,C^{(2.28)}_0 .
\end{equation}
In \eqref{5.1:eq-analyticity-spaces-a}, $p_0$ is the fixed exponent of the degree function
$p_0(\cdot)$, and $\beta_0^{\mathrm{pr}}$, with $\beta_0^{\mathrm{pr}}\le\frac12$, is the exponent
of Ba{\l}aban's recursion \cite[(2.8)--(2.9)]{B-CMP119}. At $d=4$ the first condition in
\eqref{5.1:eq-analyticity-spaces-a} is implied by $q\ge p_0+7r_{\mathrm{pr}}+1$. It gives $q>p_0$
and, since the exponent $p_1$ of \cite{B-CMP122a} satisfies $p_1<p_0$, also
\begin{equation*}
q\ \ge\ p_1+\bigl(1+d+d\beta_0^{\mathrm{pr}}\bigr)r_{\mathrm{pr}}+1 .
\end{equation*}
The constants $B_3$ and $O(1)$ are those of \cite[(3.32)]{B-CMP109} and of the proviso below
\cite[(3.37)]{B-CMP109}. The bound on $C^{(2.28)}_1$ is taken from the proviso
\begin{equation*}
B_3^2\,O(1)\,M\alpha_0<\tfrac12\alpha_1,
\end{equation*}
which is imposed in \cite[p.~277, below (3.37)]{B-CMP109}, read for the radii $\alpha_{i,j}$.
In \cite[(3.32)]{B-CMP109} the companion term $2(\alpha_2+B_3O(1)M\alpha_0)$, with $\alpha_2$ the
radius so denoted there, carries $B_3$ to the first power; the squared constant is that of the
proviso. Later in this chapter this floor is used in a ratio form carrying the additional factor
$\mathsf{a}/(1-\beta)$,
\begin{equation*}
\frac{C^{(2.28)}_1}{C^{(2.28)}_0}\ \ge\ \frac{2\mathsf{a}B_3^2\,O(1)\,M}{1-\beta},
\end{equation*}
where $\mathsf{a}:=A_0/A_1$ and $\beta$ is Ba{\l}aban's enlargement fraction, the number $\beta$
of the enlargement after the large-field operation $\mathbf{T}$ described below.
Both conditions in \eqref{5.1:eq-analyticity-spaces-a} are lower bounds only, and no upper bound
is imposed on $q$: each quantity through which a later condition depends on $q$ --- the ratio
$\theta(g)$ of shift to radius, the site ratio $\theta'$, the amplitude-disc radii
$\mathsf{w}_T,\mathsf{w}_P$, the polydisc radius $\mathfrak{r}_j$, the factors of the preparatory
stages, and a factor with exponent $2q$ in the comparison of the radii at consecutive levels ---
either decreases as $q$ increases or enters through an upper bound on $\gamma$ fixed after $q$.
We also set $C_*:=\min\{C^{(2.28)}_0,C^{(2.28)}_1\}$ and
\begin{equation*}
\alpha_{*,j}:=\min\{\alpha_{0,j},\alpha_{1,j}\}=g_jC_*(\log x_j)^{q}.
\end{equation*}

\emph{Enlarged radii.} For each level $k$ put
\begin{equation}\label{5.1:eq-analyticity-spaces-b}
\hat\alpha_k:=(1+2\beta_s)\,\alpha_{\cdot}(g_k),
\end{equation}
componentwise. The pair $\hat\alpha_k$ is determined by $g_k$ alone, so it is available before
$g_{k+1}$ is defined.

Ba{\l}aban \cite{B-CMP119} states that ``the newly created expressions $E^{(k)}, R^{(k)}, B^{(k)}$
are defined on slightly larger spaces, with the coefficients in their definition bigger by $\beta$
multiplied by a corresponding number''. Following \cite[(1.34)]{B-CMP116}, we take this
number to be $1$. Ba{\l}aban's enlargement after the large-field operation $\mathbf{T}$ at level
$k+1$ is then $(1+\beta)\alpha_{\cdot}(g_{k+1})$, read at $\beta=\beta_{\mathrm{pr}}\le\beta_s$.

The factor $1+2\beta_s$ in \eqref{5.1:eq-analyticity-spaces-b}, rather than $1+\beta_s$, is
needed because $(1+\beta_s)\alpha_{\cdot}(g_k)$ falls short of
$(1+\beta_s)\alpha_{\cdot}(g_{k+1})$ by a relative amount of order $g_k^2$; this comparison of the
radii at consecutive levels is established later, where the ratios
$\alpha_{\cdot}(g_n)/\alpha_{\cdot}(g_j)$ of the radii are estimated.

If the corresponding number were some $c_{\mathrm{pr}}\ne1$, one would instead put
$\hat\alpha_k:=(1+(c_{\mathrm{pr}}+1)\beta_s)\alpha_{\cdot}(g_k)$. The squared ratio
$(\hat\alpha_k/\alpha_{\cdot}(g_j))^2$, denoted $w_{k,j}$ below, would then satisfy
$w_{k,j}\le C_wx_j/x_k$ with
\begin{equation*}
C_w:=\tfrac{10}{9}\bigl(1+(c_{\mathrm{pr}}+1)\beta_s\bigr)^2,
\end{equation*}
and the constant $C_{\mathfrak{m}}$ of the coupling flow would become
\begin{equation*}
C_w\Bigl(1+\frac{2}{\lambda(\kappa_0-4)}\Bigr),
\end{equation*}
where $\lambda=c_0/2$ and $\kappa_0$ are the constants of the coupling flow entering the
definition of $C_{\mathfrak{m}}$. Every statement below holds with these constants in place of the
factor $1+2\beta_s$ and of the values $C_w=\frac{10}{9}(1+2\beta_s)^2=\frac52$ and
$C_{\mathfrak{m}}$ that they take at $c_{\mathrm{pr}}=1$, uniformly in $K$, for every
$c_{\mathrm{pr}}$, and these constants do not depend on the complex source $w$. We keep
$c_{\mathrm{pr}}=1$.

\emph{Families.} An \emph{$E$-family at scale $j$} is a family
$F=\{F(X,\cdot)\}_{X\in\mathfrak{D}_j}$,
or, in the \emph{pointed} case, a family $\{F(X,\cdot,y)\}_{y\in X}$, satisfying the following
three conditions.
\begin{itemize}
\item[(f1)] $F(X,\cdot)$ is analytic on $\mathcal{U}_j(X)$ (respectively on
$\mathcal{U}^\flat_j(X)$, and we then speak of a family on $\mathcal{U}^\flat_j$). It depends
only on the restriction of $(\mathbf{U},\mathbf{J})$ to $X$. The terms whose domain $X$ does not
wrap around the torus do not depend on the torus.
\item[(f2)] $F(X,\cdot)$ is $G^{c}$-gauge invariant in the sense of \cite[(1.19)]{B-CMP109}.
\item[(f3)] The \emph{total} $F(U):=\sum_X F(X,U,J(U))$ is invariant under the Euclidean
symmetries of $T^{(j)}$. In the pointed case the requirement is $F(rU,ry)=F(U,y)$ for every such
symmetry $r$, where $F(U,y):=\sum_{X\ni y}F(X,U,J(U),y)$ \cite[(2.29)]{B-CMP119}.
\end{itemize}

\emph{Norms, marginal families.} Put
$\|F\|_j:=\sup_X e^{\kappa d_j(X)}|F(X,\cdot)|$, where $|F(X,\cdot)|$ is the supremum of the
modulus over $\mathcal{U}_j(X)$; $\|F\|^\flat_j$ is defined by the same formula with the supremum
over $\mathcal{U}^\flat_j(X)$.

A \emph{marginal family} is a family of the form \cite[(2.30)--(2.32)]{B-CMP119} with analyticity
spaces $\mathcal{U}^\flat_j$.

\emph{Equivalence, coefficients, histories.} Let $F(1)$ denote the value at the trivial
configuration. For $j\ge1$ let $L^{(j)}$ be the representation
\cite[(1.7)]{B-CMP109} of $\log Z^{(j-1)}(U_j)-\log Z^{(j-1)}(1)$. The constant $E_0$ is fixed
by setting $\frac12E_0$ equal to the constant of \cite{B-CMP116} that bounds $\|L^{(j)}\|_j$
uniformly in $j$.
Two families $F,F'$ are called \emph{equivalent}, written $F\simeq F'$, when the relation in the
first three lines of \eqref{5.1:eq-analyticity-spaces-c} holds; the same display fixes $\beta[F]$,
$F^{\mathrm{ren}}$, $F^{\mathrm{vac}}$, $\beta^0_j$ and $\mathbb{B}_k$, and restates the bound on
$L^{(j)}$:
\begin{equation}\label{5.1:eq-analyticity-spaces-c}
\begin{gathered}
F\simeq F'\ :\Longleftrightarrow\ \text{the totals of $F$ and $F'$ differ by a constant}\\
\text{independent of the configuration,}\\
\text{on the real regular configurations of every torus},\\
\beta[F]:=\text{the coefficient of \cite[(1.20)--(1.22)]{B-CMP109}},\\
F^{\mathrm{ren}}:=F-F(1)-\beta[F]\,A,\qquad F^{\mathrm{vac}}:=F-F(1),\\
\|L^{(j)}\|_j\le\tfrac12E_0,\qquad \beta^0_j:=\beta[L^{(j)}],\\
\mathbb{B}_k:=\bigl\{h=(H^{(j)})_{j\le k}:\\
H^{(j)}\text{ an $E$-family at scale $j$ on }\mathcal{U}_j,\ \|H^{(j)}\|_j\le E_0\bigr\}.
\end{gathered}
\end{equation}
In the definition of $\mathbb{B}_k$ the families $H^{(j)}$ may take complex values.

The integrand of each renormalisation step depends on the terms of earlier scales only through
their totals. Accordingly, it is the $\simeq$-classes of families that are carried from one step
to the next, and a localisation representing a class is chosen at each use, as justified in the
lemma on the choice of localisations stated later.
\end{definition}

\begin{proof}
We prove the elementary assertions contained in the definition.

\emph{The floor at $d=4$.} Since $\beta_0^{\mathrm{pr}}\le\frac12$ and $r_{\mathrm{pr}}>0$, at
$d=4$ we have
\begin{equation*}
(1+d+d\beta_0^{\mathrm{pr}})\,r_{\mathrm{pr}}\le(1+4+2)\,r_{\mathrm{pr}}=7r_{\mathrm{pr}}.
\end{equation*}
Hence $q\ge p_0+7r_{\mathrm{pr}}+1$ suffices for the first condition in
\eqref{5.1:eq-analyticity-spaces-a}.

\emph{Consequences of the floor.} The exponent $\beta_0^{\mathrm{pr}}$ in
\cite[(2.8)--(2.9)]{B-CMP119} is positive; hence, since also $r_{\mathrm{pr}}>0$, the quantity
$(1+d+d\beta_0^{\mathrm{pr}})r_{\mathrm{pr}}+1$ is at least $1$. The first condition in
\eqref{5.1:eq-analyticity-spaces-a} therefore gives $q\ge p_0+1>p_0$.

Ba{\l}aban's exponents satisfy $p_1<p_0$ \cite{B-CMP122a}. Consequently
\eqref{5.1:eq-analyticity-spaces-a} also gives
\begin{equation*}
q\ \ge\ p_1+\bigl(1+d+d\beta_0^{\mathrm{pr}}\bigr)r_{\mathrm{pr}}+1 .
\end{equation*}

\emph{The smaller radius.} Since $x_j=g_j^{-2}>1$, we have $g_j(\log x_j)^{q}>0$, so
\begin{align*}
\alpha_{*,j}
&=\min\bigl\{g_jC^{(2.28)}_0(\log x_j)^{q},\ g_jC^{(2.28)}_1(\log x_j)^{q}\bigr\}\\
&=g_j\min\{C^{(2.28)}_0,C^{(2.28)}_1\}(\log x_j)^{q}=g_jC_*(\log x_j)^{q}.
\end{align*}

\emph{Ba{\l}aban's enlargement.} Since $\beta_{\mathrm{pr}}\le\beta_s$ and both components of
$\alpha_{\cdot}(g_{k+1})$ are positive, we have
\begin{equation*}
(1+\beta_{\mathrm{pr}})\,\alpha_{\cdot}(g_{k+1})\le(1+\beta_s)\,\alpha_{\cdot}(g_{k+1})
\end{equation*}
componentwise. Since $U^c_j(X,\cdot,\cdot)$ increases with the radii, every inclusion of the space
with radii $(1+\beta_s)\alpha_{\cdot}(g_{k+1})$ into a larger space therefore holds a fortiori for
the space with Ba{\l}aban's radii $(1+\beta_{\mathrm{pr}})\alpha_{\cdot}(g_{k+1})$.
This a fortiori containment is the one used in part (ii) of
Lemma~\ref{5.2:lem-enlargement-fractions}.
\end{proof}

\begin{definition}[Functions of the coupling and radius ratios]\label{5.1:not-coupling-functions}
For $0<g\le\gamma_J$ and for levels $n$, $j$, $k$ put
\begin{equation}\label{5.1:eq-coupling-functions-a}
\begin{aligned}
T(g) &:= A_1(\log g^{-2})^{p_0}, &\qquad r_T &:= \frac{\varepsilon_1}{2T(g)},
&\qquad \tau &:= \frac{g}{r_T},\\
\vartheta &:= E^{*}\exp\Bigl(-\tfrac14\gamma_2T(g)^2\Bigr),
& \mathrm{e} &:= 2E^{*}\tau+\vartheta, & &\\
\rho_{n,j} &:= \frac{\alpha_{\cdot}(g_n)}{\alpha_{\cdot}(g_j)},
& w_{k,j} &:= \Bigl(\frac{\hat\alpha_k}{\alpha_j}\Bigr)^{2}, & &
\end{aligned}
\end{equation}
where $\alpha_{\cdot}(g)$ stands for either of the coupling-dependent radii $gC_i(\log g^{-2})^{q}$,
$i=0,1$, of Definition~\ref{5.1:def-analyticity-spaces}, whose values at $g=g_j$ are
$\alpha_{0,j}$ and $\alpha_{1,j}$, and $\alpha_j:=\alpha_{\cdot}(g_j)$, together with
\begin{equation}\label{5.1:eq-coupling-functions-1}
\psi := \mathrm{e}\,g^{2},\qquad \theta(g) := C_\theta(\log g^{-2})^{p_0-q}.
\end{equation}
Here $C_\theta$ is the constant for which $\theta(g)$ is the ratio of the shift to the radius in
\cite[(2.28)]{B-CMP119}. At $g=g_k$ one has $T(g_k)=\delta_k/g_k$, where $\delta_k$ is
the quantity so denoted in \cite{B-CMP119}. A subscript $k$ denotes
evaluation at $g_k$: $\mathrm{e}_k:=\mathrm{e}(g_k)$, $T_k:=T(g_k)$, and likewise for the other
functions above. Provided $\log\gamma_J^{-2}\ge 2p_0$, the functions $r_T$, $\tau$,
$\vartheta$, $\mathrm{e}$, $\psi$ and $\theta$
are nondecreasing in $g$ on $\mathopen{]}0,\gamma_J\mathclose{]}$, while $T$ is nonincreasing there.
Further, $k_1(k)$ denotes the largest $j$ with $L^{-(k-j)}M\check{R}_j\le L$ \cite{B-CMP119}.

The ratios satisfy
\begin{equation}\label{5.1:eq-coupling-functions-3}
\rho_{n,j}^2=\frac{x_j}{x_n}\Bigl(\frac{\log x_n}{\log x_j}\Bigr)^{2q}
\end{equation}
and
\begin{equation}\label{5.1:eq-coupling-functions-2}
w_{k,j} = (1+2\beta_s)^2\rho_{k,j}^2 \le \frac52\,\frac{x_j}{x_k},
\end{equation}
where the inequality is asserted for levels $j$, $k$ with $(\log x_k/\log x_j)^{2q}\le 10/9$,
in particular whenever $1<x_k\le x_j$.
\end{definition}

\begin{proof}
The constants $A_1$, $\varepsilon_1$, $E^{*}$, $\gamma_2$ and $C_\theta$ are positive. Assume
$\log\gamma_J^{-2}\ge 2p_0$. For
$0<g\le\gamma_J$ we have $\log g^{-2}\ge\log\gamma_J^{-2}\ge 2p_0>0$. Since $\log g^{-2}$
decreases as $g$ increases and $p_0>0$, the function $T$ is nonincreasing and positive on
$\mathopen{]}0,\gamma_J\mathclose{]}$. Hence $r_T=\varepsilon_1/(2T)$ and
$\vartheta=E^{*}\exp(-\tfrac14\gamma_2T^2)$ are nondecreasing there. Next,
$\tau=(2A_1/\varepsilon_1)\,g(\log g^{-2})^{p_0}$, and since the derivative of $\log g^{-2}$ in $g$
is $-2/g$,
\begin{equation*}
\frac{d}{dg}\Bigl(g(\log g^{-2})^{p_0}\Bigr)
=(\log g^{-2})^{p_0-1}\bigl(\log g^{-2}-2p_0\bigr).
\end{equation*}
This is nonnegative wherever $\log g^{-2}\ge 2p_0$, hence on all of
$\mathopen{]}0,\gamma_J\mathclose{]}$, so $\tau$ is nondecreasing there. Then
$\mathrm{e}=2E^{*}\tau+\vartheta$ is nondecreasing as a sum of nondecreasing functions with
$E^{*}>0$, and it is nonnegative; so $\psi=\mathrm{e}\,g^2$ is nondecreasing as a product of two
nonnegative nondecreasing functions. Finally $q>p_0$ by
Definition~\ref{5.1:def-analyticity-spaces}, so the exponent $p_0-q$ in $\theta$ is negative; since
$\log g^{-2}$ is positive and decreasing in $g$, $\theta$ is nondecreasing.

By Definition~\ref{5.1:def-analyticity-spaces} the radius at level $j$ is $g_jC_i(\log x_j)^q$, with
$x_j=g_j^{-2}$. For either $i$ the constant $C_i$ cancels in the ratio, so
$\rho_{n,j}=(g_n/g_j)(\log x_n/\log x_j)^q$ does not depend on $i$. Squaring and using
$g_n^2/g_j^2=x_j/x_n$ gives \eqref{5.1:eq-coupling-functions-3}.

The enlarged radius at level $k$ is $\hat\alpha_k=(1+2\beta_s)\alpha_{\cdot}(g_k)$, by
\eqref{5.1:eq-analyticity-spaces-b}. Divide this identity by $\alpha_j$ and square. By the
definition of $\rho_{k,j}$ in \eqref{5.1:eq-coupling-functions-a}, this gives
\begin{equation*}
w_{k,j}=(1+2\beta_s)^2\rho_{k,j}^2 ,
\end{equation*}
which is the equality in \eqref{5.1:eq-coupling-functions-2}.

Since $\beta_s=1/4$, we have $(1+2\beta_s)^2=9/4$. Insert \eqref{5.1:eq-coupling-functions-3}
with $n=k$. This gives
\begin{equation*}
w_{k,j}=\frac94\,\frac{x_j}{x_k}\Bigl(\frac{\log x_k}{\log x_j}\Bigr)^{2q}.
\end{equation*}
On the range of levels for which the inequality in \eqref{5.1:eq-coupling-functions-2} is
asserted, the logarithmic factor does not exceed $10/9$ by hypothesis. When $1<x_k\le x_j$ this
holds because then $0<\log x_k\le\log x_j$ and $q>0$, so that the factor is at most $1$. Since
$\tfrac94\cdot\tfrac{10}{9}=\tfrac52$, this gives the inequality in
\eqref{5.1:eq-coupling-functions-2}.
\end{proof}

\begin{definition}[The reference recursion, first passage and flow constants]
\label{5.1:def-reference-recursion}
Let $g$ be the reference coupling and put $X := g^{-2}$. Given $x_0 > X$, the
\emph{reference recursion} is
\begin{equation}\label{5.1:eq-reference-recursion-1}
x_{k+1} := x_k - b_{k+1}(0),
\end{equation}
where $b_{k+1}(\zeta)$ is the coefficient appearing in the recursion
\begin{equation*}
\zeta_{k+1} = \zeta_k + b_{k+1}(\zeta) - b_{k+1}(0)
\end{equation*}
of the running-shift statement of the correlation theorem, a function of the shift
$\zeta$, and $b_{k+1}(0)$ is its value at $\zeta = 0$. The \emph{first passage} $K$ is
the largest $k$ such that $x_j > X$ for all $j \le k$.

Let $c_0$ and $c_2$ be the two constants of the bound on the flow of the running coupling
from which the lower bound $\underline{\lambda} = c_0/4$ and the constant $c_5 = 2c_2$ are
formed, and put $\lambda := c_0/2$. For each $k$ put
\begin{equation}\label{5.1:eq-reference-recursion-2}
y_k := \min_{i \le k} x_i - 2c_2 .
\end{equation}
The sequence $(y_k)$ is nonincreasing. Indeed, the minimum in
\eqref{5.1:eq-reference-recursion-2} is taken over the indices $i \le k$, and
this set of indices grows with $k$, so the minimum cannot increase.

The \emph{flow constants} are, for $\varrho > 0$,
\begin{equation}\label{5.1:eq-reference-recursion-a}
\begin{aligned}
\mathfrak{m}_k &:= C_{\mathfrak{m}}\, y_k^{-(\kappa_0-2)/2},
 &\qquad C_{\mathfrak{m}} &:= \frac{5}{2} + \frac{5}{\lambda(\kappa_0-4)},\\
v_k &:= 2C_\beta C_t C_{\mathfrak{m}}\, y_k^{-(\kappa_0-4)/2},
 &\qquad \varphi'_k &:= \frac{16C_\beta}{9\sigma_{\mathrm{rad}}}\,\psi\bigl(y_k^{-1/2}\bigr),\\
\varphi''_k &:= 2v_k,
 &\qquad \mathcal{D}_k(\varrho) &:= \varphi'_k + \frac{\varphi''_k}{\varrho}.
\end{aligned}
\end{equation}
The symbols in \eqref{5.1:eq-reference-recursion-a} are as follows.
\begin{itemize}
\item $\kappa_0$ is the exponent in the running-shift statement of the correlation
theorem that enters there through the power $g^{\kappa_0-6}$ of the reference coupling,
in the bound on the sum over $k<K$ of the quantities $\mathcal{D}_k(\varrho)$.
\item $C_\beta$ is the constant of the extraction functional in
\eqref{5.2:eq-extraction-cost-a}.
\item $C_t$ is the constant that accompanies the radius $\mathfrak{r}_j$ of the polydiscs
contained in the analyticity spaces $\mathcal{U}^{\flat}_j$ at coupling-dependent radii,
and the ratio $\mathfrak{t}_j$ of that radius; it is fixed together with them.
\item $\sigma_{\mathrm{rad}}$ is the constant fixing the radii $\sigma_{\mathrm{rad}} x_j$
of the polydiscs of the running-shift statement.
\item $\psi$ is the function of the coupling fixed in the notation on functions of the
coupling and radius ratios (Notation~\ref{5.1:not-coupling-functions}). It is evaluated
here at the coupling value $y_k^{-1/2}$.
\end{itemize}
\end{definition}

\begin{definition}[The admitted inputs and the order of choice of the constants]
\label{5.1:def-admitted-inputs}
The proof of the correlation theorem rests on the scaffold of
Definition~\ref{5.1:def-source-free-scaffold}, on the regions, labels and expansion rule of
Definition~\ref{5.1:def-expansion-rule}, on the spaces of
Definition~\ref{5.1:def-analyticity-spaces}, and on the inputs and the order of choice fixed in
(a)--(g) below. In (a)--(g) a \emph{true pair} is a pair $(U,J(U))$ in which $U$ is a scaffold
minimiser, in the sense of Definition~\ref{5.1:def-source-free-scaffold}, of a determining set equal
to or coarser than that of the term's own sequence $(\Omega_j)_{j\le k+1}$, taken on the support of
the term's characteristic functions. Thus $U_{k+1}$, every $U_m$ with $m\ge k+1$, the preparatory
configurations $U_k^{(n)}$ and the inner backgrounds of \cite[(1.71)]{B-CMP122b} give true pairs.

\begin{enumerate}
\item[(a)] \emph{Ba{\l}aban's source-free run.} The run at $z=0$ of
\cite[Theorem~1]{B-CMP122b}, together with the base theorem of the source-free run, taken in proof
form under the three declarations (D1)--(D3) stated with that base theorem, is read on the family
declared in
Definition~\ref{5.1:def-source-free-scaffold}. Of this run only two kinds of statements are used:
identities, and $z$-free statements about configuration maps, domains, operators and numbers.
Every estimate on terms of the expansion is proved again in this book. Among the $z$-free
statements used are the following; one further such statement, the regularity of the background
at the scale of its data \cite[(2.19)--(2.22)]{B-CMP119}, is used in (g).
\begin{enumerate}
\item[(a1)] The comb-gauge statements (c1)--(c3) of the lemma on the comb gauge, taken at the own
box of $X$, with the rider stated with them.
\item[(a2)] The extensions of (c1)--(c3), at the level of data, to the chains of the sites (P) and
(L) of Definition~\ref{5.2:def-sites-units}. Of the operation \cite[(1.65)]{B-CMP122b} what is used
here are the configuration maps: the
presented carrier pair $\bigl((\hat{\mathcal{U}}^{\mathfrak{b}})^h,
\operatorname{Ad}_h\hat{\mathcal{J}}^{\mathfrak{b}}\bigr)$ of the chain background indexed by
$\mathfrak{b}$, as a holomorphic function of
$(\mathbf{U},\mathbf{J})$. The values on that pair of the units of the fourth class of the budget
table are presented with their live slots rotated by the site map
$\mathfrak{Q}^{\mathfrak{b}}(\zeta)$. These values are obtained through the statement on orbit data
at complex rotations described in (d) below, which is an input of the correlation theorem at the
three sites (L), (T), (P), with the Cauchy-disc rider at site (L).
\item[(a3)] The geometry of Ba{\l}aban's regions used by the boundary lemma, part (iii): $Z$ is a
union of components of $\Lambda_k^c$, and every region is a union of cubes of the nested
$L$-adic partitions \cite{B-CMP122a}, \cite{B-CMP119}.
\item[(a4)] The $z$-free resummation geometry at site (P), with its constant $K_{\mathrm{loc}}$.
It is admitted by statement, as a further $z$-free hypothesis beside the statements of the run.
\item[(a5)] The two $z$-free statements of (b) below, the random-walk bounds of (b3) and the
far-field bound of (b4), which are statements of
Ba{\l}aban about his own operators, supports and thresholds at the real scaffold. Neither is an
estimate on a term of the expansion.
\end{enumerate}

\item[(b)] \emph{The four inputs at scaffold objects.} Let
$\mathcal{A}:=C^*\Delta^{(k)}C$ be the form of \cite[(3.15), (2.21)]{B-CMP119} in the exponential
block gauge of \cite{B-CMP109}, and let $\mathfrak{X}$ be the set of free fluctuation bonds of
$\Omega_{k+1}$. All inputs (b1)--(b4) are evaluated only at scaffold objects: the conditional mean
of the rim datum, the covariance $C_{\Lambda^*}$ of (b3), the form $\mathcal{A}$ and the supports
are those of the frozen creating step, the operation $\mathbf{T}$ at which the region was created.
The complex source $w=\sum_i z_i f_i$ enters only the potentials, and there only under majorants
that do not depend on $w$.

(b1) At every true pair and for every $k$ and every sequence $\{\Omega_j\}$, the form
$\mathcal{A}$ is real symmetric on $\ell^2(\mathfrak{X};\mathfrak{g})$ and satisfies the first
inequality of \eqref{5.1:eq-admitted-inputs-a}, with
$\gamma_0:=\gamma_0^X(d,L)/c_\Phi(d,L)^2>0$. This is the operator of \cite[Sect.~E]{B-CMP99b} at
Ba{\l}aban's Dirichlet domain $\Omega^{(k)}_{k+1}$, and not at the domain $\Lambda^*$ (both in the
notation of the coercivity proposition below): at
$\Lambda^*$ one obtains only the compression $\mathcal{A}_{\Lambda^*\Lambda^*}\ge\gamma_0$, which
bounds $C_{\Lambda^*}$ but gives no information about the Schur complement on the rim. From (b1) one
obtains $\mathcal{A}_{\bar\Lambda\bar\Lambda}\ge\gamma_0$ for every set of bonds
$\bar\Lambda\subset\mathfrak{X}$, in particular for $\bar\Lambda=\bar\Lambda^*$, the set of bonds
touching $\Lambda^*$, and for $\bar\Lambda=I:=\bar\Lambda^*\cap Z_0$, with $Z_0$ as in (b3),
since these are compressions. One also obtains $S,S_{cc}\ge\gamma_0$ for the Schur complements $S$,
$S_{cc}$ on the rim of the rim-transfer lemmas, because by those lemmas
\begin{equation*}
\langle x,Sx\rangle=\min_y\bigl\langle (y,x),\mathcal{A}(y,x)\bigr\rangle .
\end{equation*}
The lower bound (b1) is a consequence of the coercivity proposition for the constrained
fluctuation form. Its hypotheses hold at every true pair. The
hypothesis on the domain holds by \cite{B-CMP119}: $\Omega_{k+1}$ is a union of
$LM\check R_{k+1}$-cubes of the lattice $T_\eta$ of spacing $\eta$, that is, the image under
Ba{\l}aban's rescaling $B^k$ of a union of big blocks of $T^{(k+1)}$, lying outside the enlargement
${Z_k'}^{\sim8}$ of Ba{\l}aban's set $Z_k'$ by eight layers of $M\check R_k$-cubes (four, one, one
and two layers), with $Z_k=\Lambda_k^c\supseteq\Omega_k^c$; and
$\{\Omega_j\}_{j\le k+1}$
is Ba{\l}aban's own sequence of domains, the one that determines $U_{k+1}$; the distance clause of
that hypothesis is the one stated with the coercivity proposition. The regularity of the
term's own minimiser holds by \cite[Theorem~1]{B-CMP102b}. The smallness hypothesis holds by the
lemma on the unipotent shift $\Phi$, under $M_1\alpha_{0,k}\le a_X(d,L)$, which is a ceiling on
$\gamma$ because the radius $\alpha_{0,k}$ of Definition~\ref{5.1:def-analyticity-spaces} depends on
the coupling. For $U=1$ a lower bound of this kind, with a constant depending on $d$ and $L$ only,
is \cite[Lemma~2.4, p.~246]{B-CMP96}; at every true pair the bound (b1) used here is the one given
by the coercivity proposition, and Corollary~\ref{5.2:cor-true-pair-positivity} records (b1) and
(b2) at every true pair.

(b2) At every true pair the kernel of $\mathcal{A}$, which is the member $\mathcal{M}=\mathcal{A}$
of the family of (b3) below at $\mathsf{s}=\mathbb{1}$, obeys the second
inequality of \eqref{5.1:eq-admitted-inputs-a}, where $\rho(b,b')$ is the distance between the bonds
$b$ and $b'$. The constants $B_{\mathcal{A}},\delta_{\mathcal{A}}$
are $(d,L)$-numbers, and the bound is undilated. It is a kernel bound on the operator
$\Delta^{(k)}$ of \cite[(3.156)]{B-CMP99b} itself. The bound (b2) is
Corollary~\ref{5.2:cor-kernel-decay}, the corollary to the coercivity proposition.

(b3) \emph{The random-walk bounds.} Let $Z_0\in\mathbf{D}_k$, with $\tilde Z_0$, $Z_0'$,
$\sigma_0$ as in \cite[p.~13]{B-CMP116}.
Let $\mathcal{M}$ be one of the following, where $I$ is as in (b1):
\begin{equation*}
\mathcal{A},\qquad C_{\Lambda^*}:=\mathcal{A}_{\Lambda^*\Lambda^*}^{-1},\qquad
C_{\Lambda^*}^{1/2},\qquad
C_I:=\mathcal{A}_{II}^{-1} .
\end{equation*}
Then the following hold.
\begin{itemize}
\item For every union $Z\supseteq Z_0'$ of cubes, the operator
$\mathcal{M}(\mathsf{s};\mathbf{U},\mathbf{J})$ is holomorphic on
$\{|\mathsf{s}(\Delta)|\le e^{\kappa_1}\}$ times Ba{\l}aban's analyticity space over $Z$, under
$\delta_1M\ge\kappa_1$
\cite[(1.7)--(1.11)]{B-CMP116}, and $\mathcal{M}(\mathbb{1})=\mathcal{M}$.
\item If $\mathsf{s}=0$ off $Z$, the kernel of
$\mathcal{M}(\mathsf{s};\mathbf{U},\mathbf{J})$
vanishes unless both bonds lie in one component of $Z$, and the kernel depends on
$(\mathbf{U},\mathbf{J})$ restricted to $Z$ only.
\item The reference operator is $\mathcal{M}_0:=\mathcal{M}(\mathsf{s}\equiv 0;U|_{Z_0'},0)$, where
$U|_{Z_0'}$ is the restriction of $U$ to $Z_0'$. It is Ba{\l}aban's own
reference, obtained by setting $\mathsf{s}=0$, $\mathbf{U}=U$ and $J=0$ \cite{B-CMP116}, and it is
real.
\end{itemize}
Put
\begin{equation*}
\mathsf{d}(b,b'):=\inf_{b_1\subset\cup\sigma_0}\bigl[\rho(b,b_1)+\rho(b_1,b')\bigr],
\qquad B_W:=Be^{16\kappa_1},
\end{equation*}
(this $\mathsf{d}$ is not the separation $\mathsf{d}$ of the supports of the test functions)
with $B,\delta_W$ numbers depending on $(d,L)$ \cite[(1.7), (1.11)]{B-CMP116}, and with
$\alpha_0,\alpha_1$ the radii of Ba{\l}aban's analyticity space of the first item, as in
\cite[(2.16)]{B-CMP116}. Then clauses (I)
and (II) of \eqref{5.1:eq-admitted-inputs-a} hold. Clause (III) holds in two forms.
\begin{itemize}
\item Its verbatim form holds at every pair with $\mathsf{d}(b,b')\ge\tfrac34M$. In particular it
holds whenever $b$ or $b'$ is a bond of $Z_0$, since then $\rho(b,\cup\sigma_0)\ge M$ and hence
$\mathsf{d}(b,b')\ge M$.
\item Its shifted form holds at all pairs and imposes no condition on $M$. The shift $M/2$ is
the width allowance of one localization domain of the walk expansion touching $\sigma_0$. It
corresponds to the additional factor $\exp(-\delta_0 d(y,y',\Omega))$ that
\cite[Theorem~3.14]{B-CMP99b} gives for walks with a domain touching the cubes carrying the
decoupling parameters, after
the adjustment of $\delta_0$ made there; compare \cite[Theorem~3.15]{B-CMP99b} and
\cite[(2.16)]{B-CMP116}. This is the form in which the random-walk bounds at the box of the
operator theorem at free current, derived from that theorem, state the comparison; the verbatim
form within $M/4$ of $\cup\sigma_0$, for the inverse members
$C_{\Lambda^*}$, $C_{\Lambda^*}^{1/2}$, $C_I$, would impose an upper bound on $M$ and is not used.
\end{itemize}
The covariances $C_{\Lambda^*}=C^{(k)}(\Lambda^*)$ and $C_I$ of \cite[(3.158)]{B-CMP99b},
\cite[(2.7)]{B-CMP116}, written in the exponential gauge, are equal, by part (b) of the coercivity
proposition, to $\tilde\Phi\tilde C^{\Phi}\tilde\Phi^*$ at Ba{\l}aban's domain, in the notation of
that proposition, where $\tilde C^{\Phi}$ and $\pi_f$ are defined. Here $\tilde\Phi$ is the unipotent
operator $1+\pi_f\bar D E\mathfrak{a}$. Its range is at most two $L$-blocks, its entries are at most
$4dL$, and it is analytic in $\mathbf{U}$. Accordingly the family of (b3) is taken as
\begin{equation*}
\mathcal{M}(\mathsf{s};\mathbf{U},\mathbf{J})
:=\tilde\Phi(\mathbf{U})\,\tilde C^{\Phi}(\mathsf{s};\mathbf{U},\mathbf{J})\,\tilde\Phi(\mathbf{U})^*,
\end{equation*}
with (I)--(III) holding at constants multiplied by a number $C(d,L)$, and with
$\mathcal{M}(\mathbb{1})=\mathcal{M}$ at true pairs.

(b4) \emph{The far-field bound.} On the support of Ba{\l}aban's characteristic functions at
step $k+1$, the exterior
fluctuation field obeys the inequality on $\|A_k\|$ of \eqref{5.1:eq-admitted-inputs-a}
\cite{B-CMP119}.
Here
\begin{equation*}
\mathfrak{b}_k=g_k^{-1}\delta_k=A_1(\log g_k^{-2})^{p_0}=T(g_k),\qquad
C_{\mathrm{far}}=C(d,L,B_3)\,\mathsf{a},\qquad \mathsf{a}:=A_0/A_1,
\end{equation*}
where $T(g)=A_1(\log g^{-2})^{p_0}$ is the threshold function, $\varepsilon_k=g_kp_0(g_k)$, with the
degree function $p_0(\cdot)$, and
$\delta_k=(A_1/A_0)\varepsilon_k$. The letter $\mathfrak{b}_k$ here is the small-field threshold; it
is unrelated to the chain background $\mathfrak{b}$ of (a2). This is Ba{\l}aban's regime
in which $A_0$ is much larger than $A_1$.

In display form, (b1)--(b4), together with the rate of (c) below, read:
\begin{equation}\label{5.1:eq-admitted-inputs-a}
\begin{aligned}
&\mathcal{A}\ \ge\ \gamma_0, \qquad
|\mathcal{A}(b,b')|\ \le\ B_{\mathcal{A}}\,e^{-\delta_{\mathcal{A}}\rho(b,b')},\\
&\text{(I)}\quad |\mathcal{M}_0(b,b')|\ \le\ B\,e^{-\delta_W\rho(b,b')},\qquad
\text{(II)}\quad |\mathcal{M}(\mathsf{s})(b,b')|\ \le\ B_W\,e^{-\delta_W\rho(b,b')},\\
&\text{(III)}\quad \bigl|[\mathcal{M}(\mathsf{s};\mathbf{U},\mathbf{J})-\mathcal{M}_0](b,b')\bigr|\\
&\qquad\le\ B_W\bigl[e^{-\delta_W\mathsf{d}(b,b')}+(\alpha_0+\alpha_1)
e^{-\delta_W\rho(b,b')}\bigr]
\quad\text{if }\mathsf{d}(b,b')\ge\tfrac34 M,\\
&\phantom{\text{(III)}}\quad \bigl|[\mathcal{M}(\mathsf{s};\mathbf{U},\mathbf{J})-\mathcal{M}_0](b,b')\bigr|\\
&\qquad\le\ B_W\bigl[e^{-\delta_W[\mathsf{d}(b,b')-M/2]_+}+(\alpha_0+\alpha_1)\bigr]
e^{-\delta_W\rho(b,b')}\quad\text{at all pairs},\\
&\|A_k\|\ \le\ C_{\mathrm{far}}\,\mathfrak{b}_k ,\\
&\kappa_{\mathrm{out}}\ \ge\ \mathsf{r}(\delta)\kappa\ \ge\ (1+4\beta_s)\kappa,
\qquad \mathsf{r}(\delta)=(1-10\delta)\tfrac12L .
\end{aligned}
\end{equation}
The quantities $\gamma_0^X,c_\Phi,B_{\mathcal{A}},\delta_{\mathcal{A}},B,\delta_W$ are
$(d,L)$-numbers. They are placed immediately after $(L,\delta)$ in the order
\eqref{5.1:eq-admitted-inputs-b}.

\item[(c)] \emph{The identities and $V$-free estimates of \cite{B-CMP116}.} These are the identities
and $V$-free estimates of \cite[(2.1)--(2.13), (2.27)--(2.41)]{B-CMP116}, in the form in which they
are run in the lemma on the decoupled Gaussian integral and in the lemma on the covariance
expansion, together with the tolerance lemma for the tilted Gaussian. Their part (ii) is the
corollary stated with the rim transfer. Their part (iv) consists of the displays
\cite[(2.37)--(2.41)]{B-CMP116}, taken at $\delta\le\delta_J(L)$ without Ba{\l}aban's normalisation
$(1-10\delta)\tfrac12L=1$, with the thresholds at $e^{5\rho_7}$ and $e^{5\rho_9}$, where
$\rho_m=(1-m\delta)\tfrac12L\kappa$ are the rates of \cite[(2.35)--(2.41)]{B-CMP116}. The fresh rate
$\kappa_{\mathrm{out}}$ they give is bounded in the last row of \eqref{5.1:eq-admitted-inputs-a}; its
second inequality holds because $\delta\le\delta_J(L)=(1-4/L)/10$ is equivalent to
$\mathsf{r}(\delta)\ge2=1+4\beta_s$ (Lemma~\ref{5.2:lem-decay-fraction}).
Two further bounds used beside these inputs are not admitted by statement but proved here: the
bound on the polydisc $\{|\mathsf{B}|<T^{\mathbb{C}}\}$, which is the complex-field lemma and is
obtained by running two lemmas proved before it, and the lemma on homogeneous containment of the
response map (Lemma~\ref{5.2:lem-response-containment}), derived from \cite[(1.17)--(1.22)]{B-CMP116},
\cite[(3.14)--(3.17)]{B-CMP109}, clause (c2) of the lemma on the comb gauge and the layer-schedule
lemma for the preparatory spaces.

\item[(d)] \emph{Statements of this book used as inputs.} The following statements of this book are
used as inputs, each in the form in which it is stated in this book. The floors on
$L,\kappa,\kappa_1,M_1,M$ and the ceilings on $\gamma$ of supremum form that these statements carry
are imposed as follows: the floors at the corresponding places of \eqref{5.1:eq-admitted-inputs-b},
the ceilings as members of $\gamma_J$, by name. In particular the condition set of the inheritance
theorem is imported whole: its floors at their places in \eqref{5.1:eq-admitted-inputs-b}, its
ceilings into $\gamma_J$, one of its constants being multiplied by $24=1/\mathsf{w}_0$.
\begin{itemize}
\item The homogeneity lemma for the polydisc radii $\mathfrak{r}_j$ of the analyticity spaces. It
follows from the homogeneity of
the clauses (i)--(iii) stated with it and \cite{B-CMP109}, and it is $z$-free.
\item The two-slot representation theorem for the preparatory stages and its corollary. They
concern Ba{\l}aban's expansion \cite[p.~188, (1.60)]{B-CMP122a}. They hold with rates
$a_k=(1+3\beta_{\mathrm{pr}})\kappa$ and $a_m=(1-\beta_{\mathrm{pr}})\kappa$, and with
$\kappa_1\ge16\kappa+1$ taken as a floor and never as an equality. Their sum over families at
fixed data holds beside a preparation's window layers only above a threshold of the coupling, so
it is not used as it stands. It is used only through the resummation lemma at site (P)
(Lemma~\ref{5.2:lem-animal-resummation}), with the floor on $\kappa_1$ stated there.
\item The sum identity for the coefficients $\beta^0_l$: there are numbers $s_i$ with
$|s_i|\le c_2$ such that
\begin{equation*}
\sum_{l=i+1}^{j}\beta^0_l=(j-i)c_0+s_i-s_j .
\end{equation*}
\item The lemma on the covariance expansion used in (c), with its three companion lemmas.
\item The product lemma.
\item The lemma on the hierarchical cluster expansion, with its corollary.
\item The insertion proposition and its companion lemma.
\item The rim-transfer lemmas, which contain the Combes--Thomas bound of \cite{CT73}.
\item The mean-field lemmas: the shift lemma, the one-step transfer lemma, the volume lemma, and
the rim inequality and the margin conditions of those lemmas, the latter used in (f).
\item The one-sided rim bound, whose Gaussian part is exact.
\item The coercivity proposition and its corollary, which give (b1) and (b2).
\item The geometry lemma for the family of configurations declared in
Definition~\ref{5.1:def-source-free-scaffold}, which is used for
\cite[(1.80)--(1.89)]{B-CMP122b}.
\item The strip lemmas.
\item The localised containment corollary at site (P).
\item The inheritance theorem. Every chain term $\mathbb{C}^{(n)}_k(X,\mathcal{S})$ of the two-slot
representation is a stored term on $\mathfrak{D}^\flat_n(X)\times\mathfrak{D}_{\mathcal{S}}(X)$. It
is holomorphic in the pair $(\mathbb{U},\mathbb{J})$ and in the weak slots, real-bounded on the
strong boxes, jointly gauge invariant and local in $X$, and it
has stage norm at most $C_0^\sharp$. The inheritance theorem comes with its corollaries on room,
transfer and sources. It is used for one thing: the units $\mathbb{C}^{(n)}_k(X,\mathcal{S};w)$ of
$\mathbf{B}''_k$ are stored terms on $\mathfrak{D}^\flat_n(X)\times\mathfrak{D}_{\mathcal{S}}(X)$,
with the factor $\mathcal{W}_k(X)$ added for sources, and this domain contains the configurations of
the chain of site (L) at $\mathfrak{T}$-locations. This supplies the content of the inductive
assumption that \cite{B-CMP122b} invokes, with reference to \cite[(1.70)]{B-CMP122a}, for the terms
of $\mathbf{B}''_k$.
\item The arrest theorem. It is stated under its own hypotheses, which comprise floors on $L$, on
$M$, on the ratio $M/M_1$ (among them $\delta_1M/M_1\ge2$), on further products of rates and
scales, and on Ba{\l}aban's free constant $C_0^{(2.28)}$, the condition $q>p_0$, ceilings on
$\gamma$ of supremum form, and the inductive inputs it lists; and with the arrest spaces in force
from each arrest on. Under these, for every live arrest $\mathfrak{a}$, every frozen term of
$\mathfrak{F}(\mathfrak{a})$, every later operation and every point of the term's domain, the image
of that operation's arrest source satisfies the frozen term's clause, with the room, spend and
ratio of its table.
Further, the real integration domain lies in every arrest space, the analyticity domains assigned
to newborn terms are
legitimate, off live pieces every space is the one of the source-free run, and on a live piece the
arrest spaces are nested. No constant depends on $k$,
$k'-k_a$, $N$, $N_0$, the nest depth, $K$ or a slot, and $L$ enters through floors only. The
containment of the arrested left-overs $\mathcal{C}_{\mathrm{arr}}$ is this theorem's and not the
inheritance theorem's.
\item The statement on orbit data at complex rotations, proved in this book as an implication from
the deep-unit lemma with the budget of radius spends that accompanies it, the budget of the
Cauchy-disc rider of the correlation theorem at site (L), the expansion rule of
Definition~\ref{5.1:def-expansion-rule}, the containment statements at the three sites, the
resummation lemma at site (P) and the filing proposition, and, for its second
conjunct, also from the relative theorem and the source-tube statement. The boundary terms and the
chain terms have orbit data, of radius $\varepsilon_\star$ at new slots and of the radius that
statement assigns to older law sites, with the same bounds in the rotated norm.
\item The statement on the change-of-variables terms and the never-integrated slots of an
interrupted preparation. These terms are not chain terms. Their analyticity domain, including the
one at site (L), and
their analyticity are that statement's; their Gaussian factor is bounded by the Gaussian-factor
lemma. Their size is at most $1$, since in \cite{B-CMP119} they are bounded by
any positive power of $g_k$.
\item The relative theorem for $\mathbf{R}$-born boundary terms, proved in this book as an
implication; among its antecedents is the statement on window data at the preparatory site, whose
proof is outstanding; its route runs through the window lemma, the own-level membership theorem
with its corollary and the remainder theorem, under a closure condition on the paired charts. Its
bound, for every label $A$ of a boundary-type term, is
\begin{equation*}
\mathfrak{N}^{A,\mathbb{C}}(\mathbf{B}'(A))\le e^{\lambda_\theta}c_1^\flat
e^{-(1+\frac14\beta_s)\kappa d_k(A)} ,
\end{equation*}
and it propagates from level $k$ to level $k+1$ together with the inductive hypothesis of the
correlation theorem.
\item The deep-unit lemma. There is a number $C_{3F}$, fixed before $\gamma$, such that for every
deep unit $v$ and all link families $(\mathsf{y}_t)$ that are not all empty,
\begin{equation*}
|\mathfrak{p}|\le C_{3F}\,e^{-\frac12\delta_P d^\wedge(X_v,\mathsf{C}_C^c)}\,\mathfrak{N}(v)\,
e^{-(\kappa_1-1)\sum_t|\mathsf{y}_t|}.
\end{equation*}
This lemma proves the additional decay factor that \cite{B-CMP122b} states for the pieces of a
unit in which one of the later link families is not empty.
\item The source-tube statement. It is the conjunction of four clauses: the form clause, the
oscillation clause, the deep-unit lemma in the direction of the surrounding field, and the clause
on all admissible splittings. In the oscillation clause the two-point oscillation of the old terms
$V^\sharp$ over the source tube is at most $\Sigma^V(k)$, and $\Sigma^V(k)$ plus twice the bound
of the form clause on the kept form is at most $C_\Sigma\check R_k^{p_\Sigma}$, where
$p_\Sigma=2d+6+r_{\mathrm{pr}}+\lceil 2p_1/r_{\mathrm{pr}}\rceil$. No decay to zero of that
oscillation is asserted. The form clause is the theorem on the kept form, which rests on the
solvability theorem for \cite[(1.40)]{B-CMP122b} and, through it, on the proof in this book of
Ba{\l}aban's \cite[Proposition~1]{B-CMP122b} at the base configuration; the oscillation clause is
the base theorem for the old terms, which uses on the core the statement on the Wilson brackets.
\end{itemize}
Also used by statement, $z$-free, are the charge and budget statements that close clause (i) of
Theorem~0, transferred to the family declared in Definition~\ref{5.1:def-source-free-scaffold} by
means of its geometry lemma; the low-scale
lemma together with Ba{\l}aban's \cite[Corollary~3]{B-CMP119} in its integrated form; and
the three gain displays of \cite[pp.~380--381, 383]{B-CMP122b}.

\item[(e)] \emph{The order of choice.} The constants are chosen in the order
\eqref{5.1:eq-admitted-inputs-b}:
\begin{equation}\label{5.1:eq-admitted-inputs-b}
\begin{aligned}
&\theta_S\in[\tfrac12,\tfrac89]\ \to\ L\ \to\ \delta\in\,]0,\delta_J(L)]\ \to\
(\gamma_0^X,c_\Phi,\gamma_0,B_{\mathcal{A}},\delta_{\mathcal{A}},B,\delta_W,K_0,\delta_1,\mu_0;\
M_E,R_E,a_X)\\
&\to(\Theta_1,\varpi,\Theta_\flat)\ \to\ (\lambda,\ c_2\to\gamma^\circ)\ \to\
(r_{\mathrm{pr}}\to p_0\to p_1\to q;\ \kappa_0)\ \to\
(\bar S_I,\bar S_\omega,\bar S_\partial)\\
&\to\kappa\ \to\ \kappa_1\ \to\ B_H\ \to\ M_1\ \to\ M\ \to\ [\,\text{the constants of (e1)--(e9)}\,]
\ \to\ \mathsf{a}:=A_0/A_1\ \to\ C_{\mathrm{far}}\\
&\to\bigl(\gamma_2,\ c_\eta\to\varepsilon_\star,\ \varpi_{\mathrm{fr}},C_{\mathrm{fr}};\
\text{the }\alpha\text{'s};\ E_0\to C_V\to C_{\mathfrak{K}}\to c_M\to\alpha_4,\alpha_6\to
\varepsilon_1\to E^*\to B_0\\
&\qquad\to\varepsilon_V^*\to\bar C_{\mathsf{W}},\bar C_\sharp\to K_R\to c_V\to c_A\to
\mathsf{K}^0,\mathsf{K}^{0\prime},\mathsf{K}^0_{\mathrm{cell}},\alpha_*\bigr)\ \to\
\gamma\le\gamma^\circ .
\end{aligned}
\end{equation}
Here $\theta_S$ is absolute, the $(d,L)$-numbers and $(d,L)$-thresholds form the fourth group,
$\gamma^\circ=(8c_2+4)^{-1/2}$, and $B_H$ is taken on $|\mathsf{s}|\le e^{\kappa_1}$. Every arrow is
one-sided: a constant depends only on constants to its left, and nothing returns to $M$,
$\kappa_1$ or $L$. The constants $C_V$, $C_{\mathfrak{K}}$, $c_M$ are the constants depending on
$\kappa$ that precede $\varepsilon_1$ in the order of choice of the uniqueness theorem.

Together with the exponents the following constants are chosen. The constants
\begin{equation*}
C_N=2^{r_{\mathrm{pr}}}L,\quad c_d,\quad \mathsf{k}_d,\quad C_{\mathrm{fib}}=2e^{2\mathsf{k}_d},
\quad c_{\mathrm{tr}},\quad C_0,\quad C_0^\sharp=3c_G+1
\end{equation*}
depend on $d$, on $(d,L)$ or on $(L,r_{\mathrm{pr}})$ only; the constant $C_0$ of this list is not
Ba{\l}aban's free constant $C_0^{(2.28)}$ of (d); here $C_0^\sharp$ is the bound on the
stage norm of the chain terms in the inheritance theorem of (d).
The first summand
$2^{d+3+r_{\mathrm{pr}}}3^dC_NC_0^\sharp$ of $c_{\mathsf{P}}$ is chosen with them. The
constants $C_A$, $K_{CE}$ and $C_U$ (the last containing the factor $e^{2^d\kappa}(3M)^4$) are
chosen at their own places after $(\kappa,M)$.

The constants $C_\theta$ and $C_{\theta'}$ are chosen at their own places after
$(\kappa_1,M,\mathsf{a},B_0)$, on which they depend through the requirement that Ba{\l}aban's
real-circle radius of \cite[(1.22)]{B-CMP116}, with $\alpha_2$ and $C_1$ his letters there and
$\delta_k=(A_1/A_0)\varepsilon_k$ as in (b4), satisfy
\begin{equation*}
\mathsf{w}_T^{(1.22)}:=\frac{\alpha_2}{4B_0C_1e^{16\kappa_1}\delta_k}\ \ge\ \frac{1}{\theta(g_k)},
\end{equation*}
where $\theta(g)$ is the ratio of shift to radius at site (T), whose constant is $C_\theta$; in
this book the radius $\mathsf{w}_T:=\mathsf{w}^{\mathbb{C}}$ is the
radius of the lemma on homogeneous containment of the response map
(Lemma~\ref{5.2:lem-response-containment}). Five $\gamma$-ceilings of the filing
argument and of the modified volume bound, a further proposition of the filing argument and the
supply clause of the filing lemma use $C_\theta$ and $C_{\theta'}$. The constant $C_{\theta'}$ is
free, subject to its floor only; that floor depends on
$K_D,c_\chi,\varphi_\flat,B_H,A_1,C_*,b_+$, with $C_*$ a constant of the Cauchy-disc bound at site
(L) entering that floor, and
$C_{\theta'}$ comes after the constants
$K_D,c_\chi,\varphi_\flat,C^J_{\mathrm{age}}$. The ratio $\mathsf{w}_L:=1/\theta'(g_k)$ depends on
$C_{\theta'}$ and is fixed after it.

The bracket after $M$ consists of the following constants. Each is chosen after the constants it
depends on and before $\gamma$, and none depends on $\gamma$.
\begin{enumerate}
\item[(e1)] The constants of the resummation at the three sites. The first is
\begin{equation*}
C_A':=C_A^{(1/2)}+\tfrac12e\,C_H^{\wedge\prime},\qquad
C_H^{\wedge\prime}:=(2L)^d(1+\hat b_t+\hat b_tC_{\mathrm{cell}}),
\end{equation*}
where $C_A^{(1/2)}$ is $C_A$ taken at the half rate $\tfrac12\delta_1M$. The first summand of $C_A'$
serves sites (T), (P) and depends on $(d,L,M,\kappa_1)$; the second serves site (L) and depends on
$(d,L)$. With them come the
prefactors $C_P,C_P'$ with $C_{1/2},C_{1/2}'$.
\item[(e2)] The $(d,L,M,\kappa,\kappa_1)$-numbers of the statement on orbit data,
\begin{equation*}
\delta_L,\ C_L,\ \delta_\natural,\ B_H^S,\ C_\natural,\ K_{\mathfrak{Q}},\ c_4,\ C_A^{(1/2)},\ C_b .
\end{equation*}
Nothing before $\varpi_{\mathrm{fr}},C_{\mathrm{fr}},\delta_\natural,C_\natural$ depends on them.
The numbers $K_D,c_\chi,\varphi_\flat,C^J_{\mathrm{age}}$ are chosen after $M$ and $(\lambda,c_2)$
and before $C_{\theta'}$. Beside $c_\chi$ stand $c_a$, $c_J^{\mathrm{rec}}$ and
$c_{\mathrm{rec}}:=\max\{c_\chi,c_J^{\mathrm{rec}},c_a\}$.
\item[(e3)] The constants of the deep-unit lemma,
\begin{equation*}
r^0,\ C_{3F}=2/r^0,\ K_1,\ K_2,\ K_3,\ B_O,\ K_\natural^O,\ c_{\mathsf{U}},\ \mathsf{w}_{\boxplus},\
c_B,\ C_{\mathsf{a}},\ C_{\mathrm{age}},
\end{equation*}
and, beside $C_{\mathsf{a}}$, the single value
$C_{\mathsf{a}}^0:=C_{\mathsf{a}}\cdot\max\{1,c_{\mathcal{J}}/(2K^{\mathrm{reg}})\}$, where
$c_{\mathcal{J}}$ and $K^{\mathrm{reg}}$ are the constants of the zone-zero chain clause.
\item[(e4)] The constants of the relative theorem,
\begin{equation*}
c_\rho,\ h_\sharp,\ c_l=1+234d,\ \lambda_*(d,L),\ C_\Sigma,\ C_e,\ C_{\mathrm{cell}},\
C_{\mathsf{a}}(1),\ C^P_{\mathsf{a}},\ C_Q,\ \hat b_t,\ b_\delta,
\end{equation*}
where $\lambda_*(d,L)$ is the animal constant of the sum lemma (not the number
$\lambda_*=(40M)^{-1}$ of the test-function norm),
and, at site (P), $N^\flat,\mathsf{x},K_2^\flat,c_{\mathsf{p}},c_\star,K_\flat''$.
\item[(e5)] The constants of the pointed-unit lemma. They are $C_\uparrow$, after
$(L,r_{\mathrm{pr}},\lambda,c_2,C_N)$; then $C^<$ and $C_\circledast$, after $(\kappa,\delta)$;
then
\begin{equation*}
c_{\mathrm{reg}},\ c_\vee,\ c_\epsilon,\ s(1),\ C^{\mathrm{sm}},\ C^{\mathrm{mar}},\ c_5^{W};
\end{equation*}
then $\mathsf{K}^\flat$, after $M$; then $C^{\mathrm{pt}}$, after $r^0$; then
$C^+:=\max\{C_{3F},C^{\mathrm{pt}}\}$; and last $C^{\mathrm{pl}}_{\mathsf{a}}$. Here $c_5^W$ is the
number of the fifth class of the size table, distinct from $c_5=2c_2$. The exterior lemma adds the
number $eS_d$ and $\mathsf{K}^{0\prime}$ beside $\mathsf{K}^0$.
\item[(e6)] The constants of the exterior carrier lemma and of the sum of the fourth class.
The numbers $S_d'$, $c_{\mathfrak{h}}$, $h_\sharp$, $c_l^\circ$ are absolute or $d$-only. The
constants $\hat b$, $\hat b_t$, $c_\boxplus$, $c_{\mathfrak{r}}=c_S$, $N_\partial$ and
$C_\square^0$ come after $L$, and $\lambda_\Sigma$, $\lambda_\Sigma'$, $\epsilon_\partial$ come
after $\kappa_1$. Then come
\begin{equation*}
\mathsf{K}^{\flat,\mathrm{ex}}:=e^{(a+9)c_lc_{\mathfrak{h}}}\,27S_d'\,\mathsf{K}^\flat ,
\end{equation*}
after $(\kappa,\mathsf{K}^\flat)$, where $a$ is the rate of the exterior carrier lemma;
$\mathsf{C}_{\mathrm{fw}}$, after
$(L,\mathsf{L}_0,C_\uparrow,\delta_\natural,c_\chi)$; and the family $\mathsf{C}^\Sigma_r$ of
constants of the exterior carrier lemma, indexed by a rate $r$.
\item[(e7)] The constants of the long units,
\begin{equation*}
\mathsf{K}^{\mathrm{LG}}:=32\,\bar{\mathfrak{n}}\,M^d\,\mathsf{c}_d\,K(d)\,L^{2d},\qquad
\mathsf{K}^{\mathrm{lg}}:=\mathsf{K}^{\mathrm{LG}} .
\end{equation*}
Here $K(d)$ is the $d$-only constant of the count of tree domains, $\mathsf{c}_d=2\cdot3^d$ is the
cube-count constant, and
$\bar{\mathfrak{n}}$ is the ceiling of the sizes of the long
units, $\mathfrak{n}_v\le\bar{\mathfrak{n}}e^{-\kappa^{\mathrm{lg}}_jd_j(X_v)}$, where
$\kappa^{\mathrm{lg}}_j$ is the
decay rate at level $j$ of the sizes of the long units. The number $\bar{\mathfrak{n}}$
comes after $(E_0,C_B)$, and $\mathsf{K}^{\mathrm{LG}}$ comes after $(L,M,E_0,C_B)$ and before
$\mathsf{C}^{\mathrm{lg},c},\mathsf{C}^{\mathrm{lg},d}$. Both constants are $g$-free, and neither
depends on $K$, $k$, a level, an age, $N$, $N_0$, $\check R$, a window, an arrest, a hub, a
position, a class or a weight.
\item[(e8)] The constants of the sum of the fourth class. They are $S_d^{(1/2)}$, $c_l^\circ$,
$\hat b$, $\epsilon_\partial$, $\lambda_\Sigma$, $\lambda_\Sigma'$ (after $\kappa_1$) and the
following, in which $p_{\mathsf{b}}$ and $C_{\mathsf{b}}$ are the exponent and the constant of the
per-cell anchor bound of that sum and $\varsigma_0$ is the auxiliary decay rate of that sum:
\begin{equation*}
\begin{aligned}
c^<_4&:=\sum_{u\ge1}(1+u)^{p_{\mathsf{b}}}L^{(d+1)u}e^{-(\varsigma_0-1)(L^u-1)}
\quad\text{(after }(L,\kappa,p_{\mathsf{b}})),\\
S_{\mathsf{b}}&:=\sum_{s\ge0}(1+s)^{p_{\mathsf{b}}}e^{-s}
\quad\text{(after }(L,\kappa,p_{\mathsf{b}})),\\
C^{\mathrm{win}}_4&:=\sup_{\varrho\ge1}(LC_\uparrow\varrho)^{d+1}e^{-(2\varrho-1)}
\quad\text{(after }(L,C_\uparrow)),\qquad
\mathsf{m}_4:=\max\{1,c^<_4,C^{\mathrm{win}}_4\},\\
\mathsf{C}^{4,d}&:=32e\,dL^{d-1}S_d\,\mathsf{m}_4C_{\mathsf{b}}S_{\mathsf{b}}
\quad\text{(after }(M,C_{\mathsf{b}})),\\
\mathsf{C}^{4,c}&:=4e\,S_d^{(1/2)}(1+C_{\mathsf{b}}c^<_4)\quad\text{(after }(M,C_{\mathsf{b}})),\\
\mathsf{C}^{\mathrm{lg},c}&:=2e\,S_d^{(1/2)}\mathsf{K}^{\mathrm{LG}},\qquad
\mathsf{C}^{\mathrm{lg},d}:=16e\,dL^{d-1}S_d\,\mathsf{K}^{\mathrm{LG}}
\quad\text{(after }(L,\mathsf{K}^{\mathrm{LG}})),\\
\mathsf{b}^\natural&:=2K(d)B_0+b_\delta+C_NC_0^\sharp2^{r_{\mathrm{pr}}}
\bigl(1+(1+c_0)^{r_{\mathrm{pr}}}\bigr)\quad\text{(at the place of }B_0),
\end{aligned}
\end{equation*}
and, at the place of $C_{\theta'}$ (after $\kappa_1,M,\mathsf{a},B_0$),
\begin{equation*}
\begin{aligned}
\mathsf{c}_{\theta'}&:=e^2C_{\theta'}^{-1}\sup_{\mathsf{T}\ge1}(\mathsf{T}+1)^{q-p_1}
e^{-\mathsf{T}^{r_{\mathrm{pr}}}},\\
\mathsf{C}^{\Sigma,4}&:=\mathsf{c}_{\theta'}\bigl[(\mathsf{C}^{4,d}+\mathsf{C}^{4,c})
\mathsf{b}^\natural+\mathsf{C}^{\mathrm{lg},c}+\mathsf{C}^{\mathrm{lg},d}\bigr],\qquad
\mathsf{C}_1^{\mathrm{ex}}:=\mathsf{C}^{\Sigma,4}+\mathsf{C}^\Sigma_{a+7}.
\end{aligned}
\end{equation*}
Only the two $\gamma$-ceilings on the bracket activity at site (L) that contain its member
$\varepsilon^{\mathrm{ext},\mathsf{t}}$ depend on $\mathsf{C}_1^{\mathrm{ex}}$, through
$\varepsilon^{\mathrm{ext},\mathsf{t}}$; no floor and no earlier constant does. The product
$\theta'\mathsf{C}^{\Sigma,4}$ and the floor on $C_{\theta'}$ do not depend on $C_{\theta'}$, so
there is no circle.
\item[(e9)] The constants $K_{\mathsf{J}}$, $C^\wedge_{\mathrm{age}}$, $\mathsf{a}^\flat$ and
$C_H^\wedge$: $K_{\mathsf{J}}$ is placed after $\kappa_1$ and before $C_{\theta'}$,
$C^\wedge_{\mathrm{age}}$ after $(\lambda,c_2)$, and $\mathsf{a}^\flat$ and $C_H^\wedge$ after $L$.
With them come:
\begin{itemize}
\item $\mathsf{c}_{\mathbf{C}}=1+c_IM^5$, after $M$;
\item the increment $27\mathsf{L}_0^6\mathsf{K}^\flat$ of $\mathsf{K}^0_{\mathrm{cell}}$, after
$\mathsf{K}^\flat$, where $\mathsf{L}_0<\tfrac12L$;
\item $c_V^S:=4e^43^d$, which is absolute;
\item $C_{99}$, which is absolute, and $\lambda_\theta$, which depends on $d$ only;
\item the second summand $27\mathsf{L}_0^6L^d\mathsf{c}_d\mathsf{K}^\flat$ of $c_{\mathsf{P}}$,
placed with $\mathsf{K}^\flat$ after $M$. Hence $c_{\mathsf{P}}$ is determined at the place of
$\mathsf{K}^\flat$, and only the third member of the modified volume bound depends on it.
\end{itemize}
\end{enumerate}
The numbers $\varepsilon_\star$ and $\varpi_{\mathrm{fr}}$ depend on $(\Theta_\flat,c_\eta)$ and
on $(\theta_S,\beta_{\mathrm{pr}})$ respectively.

The following functions of the run's coupling are not entries of
\eqref{5.1:eq-admitted-inputs-b}:
\begin{equation*}
N,\ N_0,\ \check R_k,\ \check{\mathsf{T}}_k,\ \mathsf{K}_{\mathrm{val}}(k),\
\mathsf{K}_{\mathrm{cell}}(k,n),\ \varepsilon_{\mathrm{br}}(k),\ \alpha^\sharp(k),\
m_{\mathrm{lo}}(k),\ m_\sharp(k),
\end{equation*}
together with the frozen piles $\mathsf{P}(\mathfrak{a})$. They are evaluated at the run's own
$g_k\le\gamma$, after $\gamma$ is fixed, and they occur in no member of $\gamma_J$. They are
dominated inside four $\gamma$-ceilings of the filing
proposition and of the modified volume bound. These ceilings are suprema over
$\mathsf{T}\ge\mathsf{T}_\gamma$ of expressions in constants fixed before $\gamma$ only. The piles
$\mathsf{P}(\mathfrak{a})$ are the exception: they are placed in the third member of the modified
volume bound, with the constant $c_{\mathsf{P}}$ and with no ceiling.

\item[(f)] \emph{One-sided margins.} Ba{\l}aban's finitely many parameter inequalities are imposed in
the following modified form.
\begin{itemize}
\item They are imposed with activity constants multiplied by $C_\sharp$, the budget constant, which
is independent of $g$ and of $K$.
\item Every unnamed $O(1)$ multiplying $|\Gamma_j|$ or $|Z_j|$ is replaced by $C_\sharp O(1)+1$.
This is the number for the brackets and for the first three, the fifth and the boundary entries of
the last column of the budget table.
\item The per-cell constant of the volume bound \cite[(1.69)]{B-CMP122b}, with the values of this
step's chain in the fourth class piled in, is not this number. It is the function
$\mathsf{K}_{\mathrm{cell}}(k,n)$ of the modified volume bound, which a $\gamma$-ceiling of the
modified volume bound places
inside Ba{\l}aban's per-domain allowance $O(1)\log g_j^{-2}|Z_j|$ \cite{B-CMP122b}.
\item The frozen piles $\mathsf{P}(\mathfrak{a})$ of the live arrested pieces are placed in the
third member $O(1)M^dR_j^{d+1}d_j'(Z_j)$ of the same allowance, with the constant $c_{\mathsf{P}}$.
\item In the parameter inequality of \cite[\S3]{B-CMP122b} that carries the summand $+1$, this
summand is replaced by $+2$.
\item The condition $\gamma_0A_1^2p_0(\gamma)\ge8$ of the first line of
\eqref{5.1:eq-admitted-inputs-c}, in which $\gamma_0$ is the constant so written in Ba{\l}aban's
inequality and $p_0(\gamma)$
is the degree function $p_0(\cdot)$ at $\gamma$, is imposed, and each
preparatory factor of \cite[pp.~381--383]{B-CMP122b}, raised to the power $\tfrac34$, is still at
most $e^{-p_0(g_j)}$.
\end{itemize}
These are ceilings on $\gamma$ placed after $(L,M,\text{exponents})$, and each implies Ba{\l}aban's
inequality. The modified volume bounds of the filing proposition display both allowances in this
form, and so does the volume bound of the last-lattice step. For the creating step run on the
family declared in Definition~\ref{5.1:def-source-free-scaffold}, every
occurrence of these conditions is taken in the form of
the second and third lines of \eqref{5.1:eq-admitted-inputs-c}. In that
form the logarithm carries
the per-cell pile $\mathsf{K}_{\mathrm{cell}}$ of the values, by the $\gamma$-ceiling of the
modified volume bound named in the third item above, inside Ba{\l}aban's
second
member $O(1)\log g_j^{-2}|Z_j|\le O(1)M^dR_j^{d+1}d_j'(Z_j)$ \cite{B-CMP122b}. The geometry lemma
for that family has $c_j=C_{\mathfrak{d}}^\flat M^d\check R_j^{d+1}$ with the power
$\check R_j\ge\log g_j^{-2}$ needed for this. These are still ceilings on $\gamma$, of degree
$2p_0>(d+2)r_{\mathrm{pr}}$.
\begin{equation}\label{5.1:eq-admitted-inputs-c}
\begin{aligned}
&O(1)\ \mapsto\ C_\sharp O(1)+1,\qquad \gamma_0A_1^2p_0(\gamma)\ \ge\ 8;\\
&\gamma_0A_\flat^2p_0(\gamma)\ \ge\ 8,\qquad A_\flat^2:=A_1^2/(256\,\Theta_1^2);\\
&O(1)\ \text{on }|Z_j|,\ d_j'(Z_j)\ \mapsto\
199^d3^{d+1}\bigl(C_\sharp O(1)+1\bigr)\bigl(1+2^{-d}\log g_j^{-2}\bigr).
\end{aligned}
\end{equation}

\item[(g)] \emph{Unmatched counterterm pieces at complex sources.} Consider an unmatched
counterterm piece of the counterterm decomposition, that is, one with $y\notin\Lambda_j^{\circledR}$,
the region on which the counterterm is matched, lying in the collar
$\operatorname{supp}\varphi_j\subset\Omega_j$ of the cut-off $\varphi_j$. Ba{\l}aban's construction
uses the sign $e^{-\beta_jA(h_y\varphi_j,U_k)}\le1$ for such a piece. At complex
$\hat w\in D(0,2r)$, where $r>0$ is the source radius of the correlation theorem, the following
bound is used instead:
\begin{equation}\label{5.1:eq-admitted-inputs-d}
\bigl|e^{-b_j(\hat w)A(h_y\varphi_j,U_k)}\bigr|
\ \le\ e^{-b_j(0)A(h_y\varphi_j,U_k)}\cdot e^{2r\,d_{j-1}(2r)\,A(h_y\varphi_j,U_k)} ,
\end{equation}
in which the excess factor is at most $e$ per unit $j$-cell of the collar. It is one more unnamed
$O(1)$ per collar cell, contained in $C_\sharp O(1)+1$ of \eqref{5.1:eq-admitted-inputs-c}. Without
using the sign, one also has, per unit $j$-cell of the collar, with $\varepsilon(g):=g\,p_0(g)$ and
$C$ the constant of the per-cell bound on the action of the background used in the proof below,
\begin{equation*}
\bigl|e^{-b_j(\hat w)A(h_y\varphi_j,U_k)}\bigr|\ \le\ e^{(b_++\frac23r)C\varepsilon(\gamma)^2}
\ \le\ e .
\end{equation*}
\end{enumerate}
\end{definition}

\begin{proof}
Items (a)--(f) are declarations and attributions; we prove the bound
\eqref{5.1:eq-admitted-inputs-d} and the two per-cell bounds of (g).

\emph{The real part of $b_j$.} Clause (b) of the inductive hypothesis of the correlation theorem at
level $k$ gives, for $\hat w\in D(0,2r)$,
\begin{equation*}
\operatorname{Re}b_j(\hat w)\ \ge\ b_j(0)-2r\,d_{j-1}(2r).
\end{equation*}

\emph{The bound \eqref{5.1:eq-admitted-inputs-d}.} At a piece where Ba{\l}aban uses
$e^{-\beta_jA(h_y\varphi_j,U_k)}\le1$, the quantity $A(h_y\varphi_j,U_k)$ is real and non-negative;
this is what that sign expresses. Hence
\begin{equation*}
\bigl|e^{-b_j(\hat w)A(h_y\varphi_j,U_k)}\bigr|
=e^{-\operatorname{Re}b_j(\hat w)\,A(h_y\varphi_j,U_k)} .
\end{equation*}
Inserting the lower bound on $\operatorname{Re}b_j(\hat w)$ gives
\eqref{5.1:eq-admitted-inputs-d}. The first factor on its right-hand side is the factor of the run
at $z=0$ itself.

\emph{The excess factor.} On $\operatorname{supp}h_y$ the background is regular at the scale of its
data. This is a $z$-free statement of (a), namely \cite[(2.19)--(2.22)]{B-CMP119}. Hence, by the
per-cell bound on the action of the background, with $C$ the constant of that bound and
$\varepsilon(g)=g\,p_0(g)$, $A(h_y\varphi_j,U_k)\le C\varepsilon(\gamma)^2$
per unit $j$-cell. The excess factor is therefore at most
\begin{equation*}
\exp\bigl[2r\,d_{j-1}(2r)\,C\varepsilon(\gamma)^2\bigr]\ \le\ e
\end{equation*}
per cell, since $\sum_jd_{j-1}\le\tfrac13$. This is one more unnamed $O(1)$ per collar cell, and it
is contained in $C_\sharp O(1)+1$.

\emph{The bound without the sign.} The bound $|b_j|\le b_++\tfrac23r$, established in the inductive
step of the correlation theorem, together with the per-cell bound
$A(h_y\varphi_j,U_k)\le C\varepsilon(\gamma)^2$, yields, per unit $j$-cell of the collar, the last
bound displayed in (g).
\end{proof}

\begin{definition}[Classes of complex plaquette weights]\label{5.1:def-weight-classes}
Let $r>0$ be the radius in the definition of $\mathbb{D}$, and let $\lambda_*$, $K'$, $N$,
$\varepsilon$, $T^{(j)}$ and $\pi_j$ be as in Definition~\ref{5.1:def-torus-floor}.
For $0\le j\le K'$ put
\begin{equation}\label{5.1:eq-weight-classes-1}
\mathfrak{l}_j := r\lambda_* L^{-(K'-j)/2},
\qquad
r_j := r\Bigl(1+\tfrac14\sum_{i=j}^{K'} L^{-(K'-i)/2}\Bigr).
\end{equation}
For a region $X\subset\mathbb{T}^4$ let $X^{[j]}$ be $X$ together with one layer of
$\pi_j$-cubes around it, and let $\operatorname{plaq}(X^{[j]})$ be the set of plaquettes of
$X^{[j]}$. For $w\colon \operatorname{plaq}(X^{[j]})\to\mathbb{C}$ let $\operatorname{Lip}_j w$ be
the Lipschitz constant of $w$ with distances measured in units of the spacing
$L^j\varepsilon$ of $T^{(j)}$:
\begin{equation*}
\operatorname{Lip}_j w := \sup_{x(p)\neq x(p')}
\frac{L^j\varepsilon\,|w(p)-w(p')|}{|x(p)-x(p')|},
\end{equation*}
the supremum over $p,p'\in\operatorname{plaq}(X^{[j]})$, with $|\cdot|$ the distance on
$\mathbb{T}^4$. Define
\begin{equation}\label{5.1:eq-weight-classes-a}
\begin{gathered}
\mathcal{W}_j(X) := \bigl\{ w\colon \operatorname{plaq}(X^{[j]})\to\mathbb{C} \,:\,
|w(p)|<r_j \text{ for all } p,\ \operatorname{Lip}_j w<\mathfrak{l}_j \bigr\},\\
\mathcal{W}_j := \mathcal{W}_j(\mathbb{T}^4),
\qquad
\omega_{k,j} := L^{-(k-j)/2}.
\end{gathered}
\end{equation}
Then $r_j\in[r,2r[$; $\mathcal{W}_k\subset\mathcal{W}_j$ for $j\le k$;
$r_j-r_k\ge 10M\mathfrak{l}_j$ for $j<k$; and $\sum_i z_i f_i$, read as the plaquette
function $p\mapsto\sum_i z_i f_i(x(p))$, lies in $\mathcal{W}_{K'}$ for $z\in\mathbb{D}$.
\end{definition}

\begin{proof}
Since $L$ is an odd blocking factor, $L\ge 3$, so $L^{-1/2}<3/4$. The sum in
\eqref{5.1:eq-weight-classes-1} is $\sum_{n=0}^{K'-j}L^{-n/2}$, which is at least $1$ and
less than $(1-L^{-1/2})^{-1}<4$; hence $r\le r_j<2r$.

Let $j\le k$. The sum defining $r_j$ contains that defining $r_k$, so $r_j\ge r_k$; and
$\mathfrak{l}_k=L^{(k-j)/2}\mathfrak{l}_j$. The definition of $\operatorname{Lip}_j$ gives
$\operatorname{Lip}_j w=L^{j-k}\operatorname{Lip}_k w$. For $w\in\mathcal{W}_k$ (both
classes live on $\operatorname{plaq}(\mathbb{T}^4)$) we get $|w(p)|<r_k\le r_j$ and
\begin{equation*}
\operatorname{Lip}_j w < L^{j-k}\mathfrak{l}_k = L^{-(k-j)/2}\mathfrak{l}_j\le\mathfrak{l}_j ,
\end{equation*}
so $w\in\mathcal{W}_j$.

Let $j<k$. Keeping only the term $i=j$ and using $\lambda_*=(40M)^{-1}$,
\begin{equation*}
r_j-r_k=\frac r4\sum_{i=j}^{k-1}L^{-(K'-i)/2}\ \ge\ \frac r4\,L^{-(K'-j)/2}
=10M\, r\lambda_* L^{-(K'-j)/2}=10M\mathfrak{l}_j .
\end{equation*}

Let $z\in\mathbb{D}$ and $w(p)=\sum_i z_i f_i(x(p))$. Since
$\|f_i\|_\infty\le\|f_i\|_\sharp$,
\begin{equation*}
|w(p)|\le\sum_i|z_i|\,\|f_i\|_\infty\le\sum_i|z_i|\,\|f_i\|_\sharp<r\le r_{K'} .
\end{equation*}
Further $|w(p)-w(p')|\le\sum_i|z_i|\operatorname{Lip}f_i\,|x(p)-x(p')|$, and
$L^{K'}\varepsilon=L^{-m}N_T^{-1}=N^{-1}$. Since
$\operatorname{Lip}f_i\le\lambda_*N\|f_i\|_\sharp$,
\begin{equation*}
\operatorname{Lip}_{K'}w\le\frac1N\sum_i|z_i|\operatorname{Lip}f_i
\le\lambda_*\sum_i|z_i|\,\|f_i\|_\sharp<\lambda_* r=\mathfrak{l}_{K'} .
\end{equation*}
Hence $w\in\mathcal{W}_{K'}$.
\end{proof}

\begin{theorem}[The correlation theorem]\label{5.1:thm-correlation-theorem}
Assume the admitted inputs of Definition~\ref{5.1:def-admitted-inputs}, namely:
\begin{enumerate}
\item[(a)] Ba{\l}aban's renormalisation-group construction at zero source,
\cite[Theorem~1]{B-CMP122b}, of which Definition~\ref{5.1:def-admitted-inputs} admits the
identities and the source-free statements, every estimate on terms being re-proved;
\item[(b)] the source-free statements of Ba{\l}aban about configuration maps, domains and
operators that are admitted there by statement: the configuration statements on a
localisation domain's own box, together with the supplementary clause admitted with them,
and the random-walk statement and the
far-field statement of \eqref{5.1:eq-admitted-inputs-a};
\item[(c)] the statements listed there as inputs: the lower bound and the kernel bound for
the quadratic form $\mathcal{A}$ and the identities and source-free estimates taken from
\cite[(2.1)--(2.13), (2.27)--(2.41)]{B-CMP116}, all displayed in
\eqref{5.1:eq-admitted-inputs-a}, together with the
further statements listed in Definition~\ref{5.1:def-admitted-inputs}.
\end{enumerate}
Assume further the statement on the summation of boundary terms at finer zones.
Assume also the source-free statement on the resummation geometry at the preparatory site,
which Definition~\ref{5.1:def-admitted-inputs} admits as a further source-free hypothesis,
and the statement on the window data of the blocks at the preparatory site, on which the
relative bound for the boundary-type terms depends and whose proof is outstanding.

Fix $L$. Then fix Ba{\l}aban's auxiliary number $\mathsf{L}_0$ of \cite[(1.24)]{B-CMP122a}
(a number distinct from the constant $L_0$ of Theorem~0) subject to
\begin{equation*}
3\le\mathsf{L}_0<\tfrac12 L,\qquad \mathsf{L}_0^2\le L/3 ;
\end{equation*}
then the margin parameter $\delta$ of \cite[(1.32)]{B-CMP116} subject to
\begin{equation*}
\delta\in\,]0,\delta_J(L)],\qquad \delta_J(L):=\tfrac{1}{10}\bigl(1-4/L\bigr);
\end{equation*}
then $\kappa$ subject to $\kappa\ge\max\{10\log L,\ \kappa_d,\ C_\delta(d)/\delta\}$, the
member $10\log L$ being Ba{\l}aban's own floor of \cite{B-CMP119}, where
$\kappa_d$ and $C_\delta(d)$ are numbers depending on $d$ only; then $\kappa_1$ subject to
\begin{equation*}
\begin{split}
\kappa_1\ge\kappa_1^J(L,\kappa):=\max\bigl\{&16\kappa+1,\ 4L\kappa+1,\
1+c_l(2\kappa+4)+\lambda_\ast(d,L),\\
&2+a_k+3^dc_\flat+\log\bigl(2e^2\hat b_t^{\,2}N^\flat\bigr),\
4+c_\ast+\mathsf{c}_{\mathrm{an}}+2c_\flat\bigr\}
\end{split}
\end{equation*}
and to Ba{\l}aban's own floor $\tfrac12(\kappa_1-1)\ge 2L\kappa$ of \cite{B-CMP116}. Here
$c_l:=1+234d$; $\lambda_\ast(d,L)$ is an animal-counting constant depending on $d$ and $L$
only, always written with its arguments and distinct from $\lambda_*=(40M)^{-1}$ of
Definition~\ref{5.1:def-torus-floor}; $a_k:=(1+3\beta_{\mathrm{pr}})\kappa$ with
$\beta_{\mathrm{pr}}:=1/5$;
$\hat b_t:=(L+2)^d+3^d$; $N^\flat:=(10L+2)^d$; and $c_\flat$, $c_\ast$,
$\mathsf{c}_{\mathrm{an}}$ are numbers depending on $d$, $L$ and $\kappa$ only, not on $M$.
Then fix the block sizes, which are powers of $L$: $M_1$ with $M_1\ge M_E(d,L)$ and
$L^3M_1\ge R_E(d,L)$, and $M$ with $M>M_2:=L^2M_1$ (so that $M\ge L^3M_1$), subject to
\begin{equation*}
M\ge\max\{3L,2L^2\},\qquad
\delta_1M/M_1\ge 2,\qquad \delta_1M\ge\kappa_1,
\end{equation*}
where $\delta_1:=\min\{\mu_0/2,\delta_W/2\}$ with $\mu_0$, $\delta_W$ the $(d,L)$-numbers of
the order of choice; subject to the floor that $\tfrac14\delta_P$ times the weight per
stratum of the scaled distance $d^{\wedge}$ is at least $110d(a_k+1)$, where $\delta_P$ is
the rate of the exponential weights $e^{-\delta_P d^{\wedge}}$ built on $d^{\wedge}$;
and subject to the further one-sided lower bounds on $M_1$, on $M$ and on
$M/M_1$ (the last through the same weight per stratum), each imposed after $\kappa_1$ has
been fixed, that are
displayed in the subsections of this chapter where the rim and the localised remainder
corollary at the preparatory site are constructed. Fix $\kappa_0\ge 7$ and the exponents
subject to
\eqref{5.1:eq-analyticity-spaces-a}, the remaining parameters being as in hypothesis~(a).
The analysis constants and the $(d,L)$-numbers and functions of the rim construction
are chosen in the order of choice \eqref{5.1:eq-admitted-inputs-b}, which begins with the
absolute number $\theta_S\in[\tfrac12,\tfrac89]$, in which $L$, $\delta$, $\kappa$,
$\kappa_1$, $M_1$, $M$ are followed by $\mathsf{a}:=A_0/A_1$, $C_{\mathrm{far}}$ and the
remaining constants, and in which $\gamma$, the ceiling on the coupling, comes last.
The number $r$ is fixed before $\gamma$: fix $r>0$.

Then there is $\gamma_J=\gamma_J(L,\delta,M,M_1,\kappa,\kappa_1,\kappa_0,\dots,r)>0$ with
\begin{equation*}
\gamma_J\le\gamma^\circ:=(8c_2+4)^{-1/2}
\end{equation*}
such that for every $g\le\gamma_J$, every $x_0>g^{-2}$, with $K=K(x_0,g)$ the first-passage
index of Definition~\ref{5.1:def-reference-recursion}, every integer $m$ with
$m_{\mathbb{T}}(g)\le m\le K$ (the lower bound being the torus floor
$L^m\ge N_{\mathbb{T}}(g,m)$ of \eqref{5.1:eq-torus-floor-a}) and every integer
$N_T\ge 1$, with $N:=N_TL^m$ and $K':=K-m$, the following hold.
\begin{enumerate}
\item[(i)] The reference recursion of Definition~\ref{5.1:def-reference-recursion} is
defined; the first line of the display below holds for $0\le i<j\le K$, the chain of
inequalities at the beginning of its second line holds for every $k$ with $0\le k\le K$,
and $x_K$ lies in the interval stated there:
\begin{equation}\label{5.1:eq-correlation-theorem-a}
\begin{gathered}
\lambda(j-i)-2c_2\le x_i-x_j\le 3\lambda(j-i)+2c_2,\\
|b_{k+1}(0)-\beta^0_{k+1}|\le C_\beta\mathrm{e}_k+v_k\le\lambda,\qquad
x_K\in\,]X,X+b_+],\quad b_+:=3\lambda+2c_2 .
\end{gathered}
\end{equation}
\item[(ii)] (The running-shift statement.) For $k<K$ the coefficient $b_{k+1}$ and the
running shift $\zeta_{k+1}$ are analytic functions of the bare shift $\zeta$ on
$D(0,2r)$, real for real $\zeta$; the function $f_{k+1}$ of the variables
$(\xi_0,\dots,\xi_k)$ is holomorphic on the polydisc
$\mathfrak{P}_k:=\prod_{j\le k}D(0,\sigma_{\mathrm{rad}}x_j)$, with $\sigma_{\mathrm{rad}}$
the radius constant of Definition~\ref{5.1:def-reference-recursion}, and the term
$\beta\bigl[\mathsf{P}^{(k+1)}-L^{(k+1)}\bigr]$ below is its value at $\xi_j=\zeta_j$,
$j\le k$; here $\mathsf{P}^{(k+1)}$ and $Q^{(k+1)}$ are formed from the three families at
scale $k+1$ of Definition~\ref{5.2:def-three-families}, and $\beta[\cdot]$,
$L^{(k+1)}$, $\beta^0_{k+1}$ and $\|\cdot\|^\flat$ are as in
\eqref{5.1:eq-analyticity-spaces-c}; and
\begin{equation}\label{5.1:eq-correlation-theorem-b}
\begin{aligned}
&\zeta_{k+1}=\zeta_k+b_{k+1}(\zeta)-b_{k+1}(0),\\
&b_{k+1}=\beta^0_{k+1}+\beta\bigl[\mathsf{P}^{(k+1)}-L^{(k+1)}\bigr]
+\beta\bigl[Q^{(k+1)}\bigr],\\
&\beta\bigl[\mathsf{P}^{(k+1)}-L^{(k+1)}\bigr]=f_{k+1}(\zeta_0,\dots,\zeta_k),
\qquad |f_{k+1}|\le C_\beta\mathrm{e}_k,\\
&\sup_{\frac12\mathfrak{P}_k}\ \sum_{j\le k}\sigma_{\mathrm{rad}}x_j
\Bigl|\frac{\partial f_{k+1}}{\partial\xi_j}\Bigr|\le\frac43\,C_\beta\mathrm{e}_k,\\
&\sup_{\zeta\in D(0,2r)}\bigl|\beta[Q^{(k+1)}_\zeta]\bigr|\le v_k,\qquad
\|Q^{(k+1)}_\zeta\|^\flat\le\mathfrak{m}_k,\\
&|b_{k+1}(\zeta)-b_{k+1}(0)|\le d_k(\varrho)|\zeta|\ \text{ on }D(0,\varrho)
\qquad\text{for every }\varrho\in[r,2r],\\
&\sum_{k<K}d_k(\varrho)\le C_{\mathfrak{D}}\Bigl[\lambda^{-1}g(\log g^{-2})^{p_0}
+\frac{g^{\kappa_0-6}}{\lambda\varrho}\Bigr]\le\frac13
\qquad\text{for every }\varrho\in[r,2r],\\
&|\zeta_k-\zeta|\le\tfrac13|\zeta|,\qquad
\sum_k\operatorname{Lip}_{\bar D(0,r)}b_{k+1}\le\tfrac23,\\
&N_K:=\frac{d\zeta_K}{d\zeta}(0)\in\bigl[\tfrac23,\tfrac43\bigr].
\end{aligned}
\end{equation}
Here $C_{\mathfrak{D}}$ is a constant.
All dependence on $\zeta$ --- through every earlier slot factor, every older term and
every remnant --- is contained in these statements.
\item[(iii)] For $k\le K'$ and $w\in\mathcal{W}_k$, the density $\rho^w_k$ defined by the
first member of the display below, in which $T^w$ and $R^\flat$ are the two operations
whose composite $R^\flat T^w$ is one step of the induction on the level, has the induction
format with sources of Definition~\ref{5.2:def-induction-format} at level $k$;
its terms are analytic in $w$, a term localised in a domain $X$ depends on $w$ only
through its restriction to $X^+$
(Definition~\ref{5.1:def-weight-classes}), and the terms satisfy the bounds of
Ba{\l}aban's construction with the enlarged activity constants of
\eqref{5.1:eq-admitted-inputs-c}. Moreover the
second and the third relation of the display hold, the third for every plaquette $p$,
where $\zeta_k(p)$ is the running shift at level $k$ at the plaquette $p$ for the
source $w$:
\begin{equation}\label{5.1:eq-correlation-theorem-c}
\rho^w_k:=(R^\flat T^w)^k\rho_0^w,\qquad
\int\rho^w_k\,dV_k=\int\rho_0^w\,dU,\qquad |\zeta_k(p)|\le\tfrac43\sup|w|.
\end{equation}
\item[(iv)] (The torus bound.) For $n\ge 2$ and test functions $f_1,\dots,f_n$ whose
supports are pairwise at distance at least $\mathsf{d}>0$, with $Z(z)$ the integral of
\eqref{5.1:eq-torus-floor-b} at the source $w=\sum_i z_if_i$,
\begin{equation}\label{5.1:eq-correlation-theorem-d}
\begin{gathered}
\Bigl|\partial_{z_1}\cdots\partial_{z_n}\log Z(z)\big|_{z=0}\Bigr|
\le\bigl[C_{UV}(\mathsf{d})+\log 2\cdot(2\mathsf{Q}+2)^n\bigr]
\Bigl(\frac nr\Bigr)^n\prod_{i=1}^n\|f_i\|_\sharp,\\
\mathsf{Q}=e^{(E_+^\sharp+C_{\mathrm{low}})N^4},\qquad
C_{UV}(\mathsf{d})\le C\mathsf{d}^{-4}\bigl(1+\log^+\mathsf{d}^{-1}\bigr),
\end{gathered}
\end{equation}
with a constant $C$, where $E_+^\sharp$ and $C_{\mathrm{low}}$ are constants fixed in the
treatment of the last lattice. Hence, the bound not depending on $K$, the derivatives
$\partial_{z_1}\cdots\partial_{z_n}\log Z(z)|_{z=0}$ have subsequential limits as
$K\to\infty$ at fixed $m$ and $N_T$, and these limits are connected Schwinger functions
of the smeared action density $\operatorname{tr}F\cdot F$ off the diagonals; they are
symmetric, lattice-covariant and reflection positive
(Proposition~\ref{5.12:prop-subsequential-limits} and
Lemma~\ref{5.12:lem-limit-positivity}).
\end{enumerate}
None of the constants depends on $K$, $k$, $m$, the torus or the position; $N$ enters only
through $\|\cdot\|_\sharp$, through $\mathsf{Q}$ and, through $m$, through the torus floor
\eqref{5.1:eq-torus-floor-a}. The frozen whole-lattice objects --- $b_j$, $\zeta_j$, and
the families $\mathsf{P}_\xi$, $E_\zeta$, $P^\flat_\zeta$, $Q_\zeta=E_\zeta-P^\flat_\zeta$
and $R_\zeta$ constructed in the induction on the level, the reference sequence
$x_j$ and $K$ --- are numbers and families that do not depend on the torus; for the levels
$K'<k\le K+1$, at which the physical torus is too small, they are constructed on the
auxiliary tori of Definition~\ref{5.2:def-auxiliary-torus}, whose side is a power of $L$.
\end{theorem}

Clauses (i)--(iii) are proved by the induction on the level $k$ built on
Definitions~\ref{5.2:def-auxiliary-torus} and~\ref{5.2:def-induction-format}; clause (iv)
is proved in the treatment of the last lattice, and the properties of the limits stated
there are Proposition~\ref{5.12:prop-subsequential-limits} and
Lemma~\ref{5.12:lem-limit-positivity}.

\begin{lemma}[Bounds for a bounded analytic function on a disc]\label{5.2:lem-disc-bounds}
Let $c\in\mathbb{C}$, $R>0$, and let $D(c,R)=\{z\in\mathbb{C}:|z-c|<R\}$ be the open disc of centre $c$
and radius $R$. Let $f$ be analytic on $D(c,R)$, with values either in $\mathbb{C}$ or in a family
$f=(f_\iota)_{\iota\in I}$ of complex functions each analytic on $D(c,R)$, normed by
$|f(z)|:=\sup_{\iota\in I}|f_\iota(z)|$, with derivatives $f^{(n)}=(f^{(n)}_\iota)_{\iota\in I}$ taken termwise.
Let $M\ge 0$ satisfy $|f(z)|\le M$ for all $z\in D(c,R)$. Then for every $z\in D(c,R)$
\begin{equation}\label{5.2:eq-disc-bounds-a}
\begin{gathered}
\text{(a)}\quad |f(z)-f(c)|\le \frac{2M|z-c|}{R};\\
\text{(b)}\quad |f(z)-f(c)|\le \frac{2M|z-c|^2}{R^2}\quad\text{if } f'(c)=0;\\
\text{(c)}\quad \Bigl|f(z)-\sum_{n=0}^{2}\frac{f^{(n)}(c)}{n!}(z-c)^n\Bigr|
\le \frac{M\mathsf{u}^3}{1-\mathsf{u}},\qquad \mathsf{u}:=\frac{|z-c|}{R}.
\end{gathered}
\end{equation}
\end{lemma}

\begin{proof}
The quantities $|f(c)|$, $|f(z)-f(c)|$ and
$|f(z)-\sum_{n\le 2}f^{(n)}(c)(z-c)^n/n!|$ are, in the family case, the suprema over $\iota\in I$ of the
same quantities formed with $f_\iota$, and $|f_\iota|\le M$ on $D(c,R)$ for every $\iota$; the
hypothesis $f'(c)=0$ in (b) means $f_\iota'(c)=0$ for every $\iota$. It therefore suffices to prove
each bound for a complex-valued $f$ and then take the supremum over $\iota$. Let $f$ be complex-valued.

(a) and (b). Let $n=1$ in case (a) and $n=2$ in case (b), and put $g(z):=(f(z)-f(c))/(z-c)^n$ for
$z\ne c$. The function $f(z)-f(c)$ vanishes at $c$, and in case (b) its derivative $f'(c)$ also
vanishes, so it has a zero of order at least $n$ at $c$; hence $g$ has a removable singularity at
$c$ and extends to an analytic function on $D(c,R)$. Fix $z\in D(c,R)$ and choose $r$ with
$|z-c|<r<R$. On the circle $|\zeta-c|=r$ we have
\begin{equation*}
|g(\zeta)|\le\frac{|f(\zeta)|+|f(c)|}{r^n}\le\frac{2M}{r^n},
\end{equation*}
so by the maximum principle on the closed disc $|\zeta-c|\le r$ also $|g(z)|\le 2M/r^n$. Letting
$r\uparrow R$ gives $|g(z)|\le 2M/R^n$, that is $|f(z)-f(c)|\le 2M|z-c|^n/R^n$, which is (a) for
$n=1$ and (b) for $n=2$.

(c). On $D(c,R)$ the function $f$ equals its Taylor series
$f(z)=\sum_{n\ge 0}a_n(z-c)^n$ with $a_n=f^{(n)}(c)/n!$. For $0<r<R$, Cauchy's formula
\begin{equation*}
a_n=\frac{1}{2\pi i}\oint_{|\zeta-c|=r}\frac{f(\zeta)}{(\zeta-c)^{n+1}}\,d\zeta
\end{equation*}
and $|f|\le M$ give $|a_n|\le Mr^{-n}$; letting $r\uparrow R$ gives $|a_n|\le MR^{-n}$. Since
$|z-c|<R$, we have $0\le\mathsf{u}<1$, and therefore
\begin{equation*}
\Bigl|f(z)-\sum_{n=0}^{2}a_n(z-c)^n\Bigr|=\Bigl|\sum_{n\ge 3}a_n(z-c)^n\Bigr|
\le\sum_{n\ge 3}M\mathsf{u}^n=\frac{M\mathsf{u}^3}{1-\mathsf{u}},
\end{equation*}
which is (c).
\end{proof}

\begin{lemma}[The $\ell^1$ gradient bound on a polydisc]\label{5.2:lem-polydisc-gradient}
For $r>0$ write $D(0,r):=\{t\in\mathbb{C}:|t|<r\}$. Let $n\ge 1$, let
$\mathfrak{d}_1,\dots,\mathfrak{d}_n>0$ and $\omega\ge 0$, and put
\begin{equation*}
\mathfrak{P}:=\prod_{j=1}^{n}D(0,\mathfrak{d}_j).
\end{equation*}
Let $f$ be holomorphic on $\mathfrak{P}$ with $|f|\le\omega$ on $\mathfrak{P}$. Then, for every
$\xi=(\xi_1,\dots,\xi_n)$ with $|\xi_j|\le\tfrac12\mathfrak{d}_j$ for $j=1,\dots,n$,
\begin{equation}\label{5.2:eq-polydisc-gradient-1}
\sum_{j=1}^{n}\mathfrak{d}_j\,|\partial_j f(\xi)|\le\frac{4}{3}\,\omega .
\end{equation}
\end{lemma}

\begin{proof}
If $\omega=0$, then $f$ vanishes identically on $\mathfrak{P}$, all its partial derivatives vanish, and
\eqref{5.2:eq-polydisc-gradient-1} holds. Let $\omega>0$. The point $\xi$ lies in $\mathfrak{P}$,
since $|\xi_j|\le\frac12\mathfrak{d}_j<\mathfrak{d}_j$.

Rescaling. For $u=(u_1,\dots,u_n)\in D(0,1)^n$ put
$g(u):=f(\mathfrak{d}_1u_1,\dots,\mathfrak{d}_nu_n)/\omega$. Then $g$ is holomorphic on $D(0,1)^n$ with
values in the closed unit disc $\overline{D(0,1)}$, and by the chain rule
\begin{equation}\label{5.2:eq-polydisc-gradient-2}
\partial_j g(u)=\frac{\mathfrak{d}_j}{\omega}\,(\partial_j f)(\mathfrak{d}_1u_1,\dots,\mathfrak{d}_nu_n),
\qquad j=1,\dots,n .
\end{equation}
Put $a_j:=\xi_j/\mathfrak{d}_j$ and $a:=(a_1,\dots,a_n)$; then $|a_j|\le\frac12$, and
\eqref{5.2:eq-polydisc-gradient-2} at $u=a$ reads $\partial_jg(a)=\mathfrak{d}_j\,\partial_jf(\xi)/\omega$.

Recentring. For $\alpha\in D(0,1)$ let $\varphi_\alpha(t):=(t+\alpha)/(1+\bar\alpha t)$. This is a
M\"obius transformation mapping $D(0,1)$ into itself, with $\varphi_\alpha(0)=\alpha$ and
\begin{equation*}
\varphi_\alpha'(t)=\frac{(1+\bar\alpha t)-(t+\alpha)\bar\alpha}{(1+\bar\alpha t)^2}
=\frac{1-|\alpha|^2}{(1+\bar\alpha t)^2},
\qquad\text{so}\qquad \varphi_\alpha'(0)=1-|\alpha|^2 .
\end{equation*}
Define, for $t=(t_1,\dots,t_n)\in D(0,1)^n$,
\begin{equation*}
G(t):=g\bigl(\varphi_{a_1}(t_1),\dots,\varphi_{a_n}(t_n)\bigr).
\end{equation*}
Then $G$ is holomorphic on $D(0,1)^n$ with values in $\overline{D(0,1)}$, and, since the $j$-th
argument of $g$ depends on $t_j$ alone, the chain rule gives
\begin{equation}\label{5.2:eq-polydisc-gradient-3}
\partial_jG(0)=\varphi_{a_j}'(0)\,\partial_jg(a)=(1-|a_j|^2)\,\partial_jg(a).
\end{equation}

Rotation. Choose $\lambda_j\in\mathbb{C}$ with $|\lambda_j|=1$ and $\lambda_j\,\partial_jG(0)\ge 0$
(take $\lambda_j=1$ if $\partial_jG(0)=0$). For $t\in D(0,1)$ put
$h(t):=G(\lambda_1t,\dots,\lambda_nt)$; this is defined because $|\lambda_jt|=|t|<1$, and $h$ is
holomorphic on $D(0,1)$ with values in $\overline{D(0,1)}$. By the chain rule and the choice of the
$\lambda_j$,
\begin{equation*}
h'(0)=\sum_{j=1}^{n}\lambda_j\,\partial_jG(0)=\sum_{j=1}^{n}|\partial_jG(0)| .
\end{equation*}

One-variable Schwarz lemma. For $0<r<1$ Cauchy's formula gives
$h'(0)=\frac{1}{2\pi i}\oint_{|t|=r}h(t)\,t^{-2}\,dt$, hence $|h'(0)|\le r^{-1}\sup_{|t|=r}|h(t)|\le r^{-1}$;
letting $r\to1$ yields $|h'(0)|\le 1$. Combining this with \eqref{5.2:eq-polydisc-gradient-3},
with $1-|a_j|^2\ge\frac34$ (because $|a_j|\le\frac12$), and with
$\partial_jg(a)=\mathfrak{d}_j\,\partial_jf(\xi)/\omega$, we obtain
\begin{equation*}
1\ge|h'(0)|=\sum_{j=1}^{n}(1-|a_j|^2)\,|\partial_jg(a)|
\ge\frac34\sum_{j=1}^{n}\frac{\mathfrak{d}_j\,|\partial_jf(\xi)|}{\omega}.
\end{equation*}
Multiplying by $\frac43\omega$ gives \eqref{5.2:eq-polydisc-gradient-1}.
\end{proof}

\begin{lemma}[A two-constants bound on the half disc]\label{5.2:lem-two-constants}
For $c\in\mathbb{C}$ and $r>0$ write $D(c,r):=\{z\in\mathbb{C}: |z-c|<r\}$. Let $\sigma>0$, let $0\le m\le M$
be real numbers (in this lemma only, $m$ and $M$ denote these two bounds), and let $\varphi$ be analytic on
$D(1,\sigma)$ with
\begin{equation*}
|\varphi(z)|\le M \quad (z\in D(1,\sigma)), \qquad
|\varphi(x)|\le m \quad (x\in\,]1-\sigma,1+\sigma[\,).
\end{equation*}
Then $|\varphi(z)|\le m^{1/3}M^{2/3}$ for every $z\in D(1,\sigma/2)$.
\end{lemma}

\begin{proof}
If $m=0$, then $\varphi$ vanishes on the segment $]1-\sigma,1+\sigma[$, hence on the connected open set
$D(1,\sigma)$ by the identity theorem, and the claim holds. Let $m>0$.

Let $D_+:=\{z\in D(1,\sigma):\operatorname{Im}z>0\}$ be the upper half-disc. Its boundary consists of
the diameter $]1-\sigma,1+\sigma[$, the open upper half of the circle $|z-1|=\sigma$, and the two points
$1\pm\sigma$. Let $\varpi$ be the harmonic measure of the diameter with respect to $D_+$. We compute it
through a conformal map. Put $u:=(z-1)/\sigma$; this maps $D_+$ onto the upper half of the unit disc
and the diameter onto $]-1,1[$. For $u$ in the upper half of the unit disc,
\begin{equation*}
\operatorname{Re}\frac{1+u}{1-u}=\frac{1-|u|^2}{|1-u|^2}>0,\qquad
\operatorname{Im}\frac{1+u}{1-u}=\frac{2\operatorname{Im}u}{|1-u|^2}>0,
\end{equation*}
so $u\mapsto(1+u)/(1-u)$ maps it onto the open first quadrant, sending $]-1,1[$ onto $]0,\infty[$
and the open upper unit semicircle onto the positive imaginary axis. Squaring maps the first quadrant
onto the upper half-plane, $]0,\infty[$ onto $]0,\infty[$ and the positive imaginary axis onto
$]-\infty,0[$. Hence
\begin{equation*}
\psi(u):=\Bigl(\frac{1+u}{1-u}\Bigr)^{2}
\end{equation*}
maps $D_+$ (in the variable $u$) conformally onto the upper half-plane, the diameter onto $]0,\infty[$
and the arc onto $]-\infty,0[$, and it extends holomorphically across the open diameter and the open arc.
On the upper half-plane $\omega(\zeta):=1-\tfrac{1}{\pi}\arg\zeta$, with $\arg\zeta\in\,]0,\pi[$, is
harmonic (as $\arg\zeta=\operatorname{Im}\log\zeta$), continuous up to $\mathbb{R}\setminus\{0\}$,
equal to $1$ on $]0,\infty[$ and to $0$ on $]-\infty,0[$. By conformal invariance,
\begin{equation}\label{5.2:eq-two-constants-1}
\varpi=\omega\circ\psi=1-\frac{2}{\pi}\arg\frac{1+u}{1-u},
\end{equation}
a harmonic function on $D_+$ with values in $[0,1]$, tending to $1$ at each point of the diameter and
to $0$ at each point of the open arc.

Set $h:=\varpi\log m+(1-\varpi)\log M$, harmonic on $D_+$, with $\log m\le h\le\log M$. The function
$\log|\varphi|$ is subharmonic on $D_+$ (with value $-\infty$ at zeros of $\varphi$) and
$\log|\varphi|\le\log M$ there; so $\log|\varphi|-h$ is subharmonic and bounded above by
$\log M-\log m$. At each point of the open diameter its upper limit is at most $\log m-\log m=0$,
because $\varphi$ is continuous there with $|\varphi|\le m$; at each point of the open arc its upper
limit is at most $\log M-\log M=0$. The exceptional boundary set $\{1\pm\sigma\}$ is finite, so the
extended maximum principle for subharmonic functions bounded above gives
\begin{equation}\label{5.2:eq-two-constants-2}
\log|\varphi|\le\varpi\log m+(1-\varpi)\log M\qquad\text{on }D_+ .
\end{equation}

Now let $z\in D_+\cap D(1,\sigma/2)$, that is $|u|<1/2$ and $\operatorname{Im}u>0$. The points $1+u$
and $1-u$ lie in the disc of centre $1$ and radius $|u|<1$, every point of which has argument of
absolute value at most $\arcsin|u|$. Both arguments thus lie in $]-\pi/2,\pi/2[$, the argument of the
quotient is their difference, and
\begin{equation*}
\arg\frac{1+u}{1-u}\le|\arg(1+u)|+|\arg(1-u)|\le 2\arcsin\tfrac12=\frac{\pi}{3}.
\end{equation*}
By \eqref{5.2:eq-two-constants-1}, $\varpi\ge 1-\frac{2}{\pi}\cdot\frac{\pi}{3}=\frac13$ there. Since
$\log M-\log m\ge0$, \eqref{5.2:eq-two-constants-2} yields
\begin{equation*}
\log|\varphi|\le\log M-\varpi(\log M-\log m)\le\tfrac13\log m+\tfrac23\log M,
\end{equation*}
that is $|\varphi(z)|\le m^{1/3}M^{2/3}$.

The same argument applies on the lower half-disc, since $z\mapsto\overline{\varphi(\bar z)}$ is
analytic on $D(1,\sigma)$ and satisfies the same two bounds. At points of the diameter,
$|\varphi|\le m=m^{1/3}m^{2/3}\le m^{1/3}M^{2/3}$ because $m\le M$. This covers all of $D(1,\sigma/2)$.
\end{proof}

\begin{lemma}[The fluctuation potential on the ball of histories]\label{5.2:lem-fluctuation-potential}
Let $h=(H^{(j)})_{j\le k}\in\mathbb{B}_k$, the ball of complex histories of \eqref{5.1:eq-analyticity-spaces-c}. Then there are functions
\begin{equation*}
\mathcal{V}'(Y,\mathbf{U},\mathbf{J},\mathsf{A})[h],\qquad Y\in D_k ,
\end{equation*}
with the following properties.
\begin{enumerate}
\item[(a)] Each $\mathcal{V}'(Y,\cdot)[h]$ is linear in $h$.
\item[(b)] Each $\mathcal{V}'(Y,\cdot)[h]$ is analytic on the product of the space
\begin{equation*}
U^c_{k+1}\bigl(Y,(1+\beta_{\mathrm{pr}})\alpha_0,(1+\beta_{\mathrm{pr}})\alpha_1,\alpha_0\bigr)
\end{equation*}
of \cite[Lemma~1]{B-CMP116}, Ba{\l}aban's enlargement fraction being read at
$\beta_{\mathrm{pr}}=1/5$, with the set of $\mathfrak{g}^{\mathbb{C}}$-valued fields $\mathsf{A}$
satisfying $|\mathsf{A}|<\varepsilon_1$ on $Y$.
\item[(c)] $\mathcal{V}'(Y,\cdot)[h]$ is localised in $Y$, is invariant under the transformations \cite[(3.29)]{B-CMP109}, and vanishes at $\mathsf{A}=0$.
\item[(d)] The functions sum to the difference of renormalised totals:
\begin{equation}\label{5.2:eq-fluctuation-potential-1}
\sum_{Y}\mathcal{V}'(Y,U_{k+1},J_{k+1},\mathsf{A})[h]
= h^{\mathrm{ren}}\bigl(U_k(e^{iB'}V^{(k)})\bigr)-h^{\mathrm{ren}}(U_{k+1}),
\end{equation}
where $h^{\mathrm{ren}}=\bigl((H^{(j)})^{\mathrm{ren}}\bigr)_{j\le k}$ is the history of the
renormalised families of \eqref{5.1:eq-analyticity-spaces-c}, evaluated through its totals in
the way the families $E^{(j)}$ enter Ba{\l}aban's effective action, and
\begin{equation*}
B'=C\mathsf{A}-\mathcal{M}(C\mathsf{A}),
\end{equation*}
where $C$ is the operator of \cite[Lemma~1]{B-CMP116} and $\mathcal{M}$ denotes the linear
operator written there with the letters $h\tilde{D}$; the operator $\mathcal{M}$ does not
depend on the history $h$.
\item[(e)] They obey the bound
\begin{equation}\label{5.2:eq-fluctuation-potential-2}
|\mathcal{V}'(Y)|\le \bar\varepsilon\, e^{-(1-2\delta)\kappa d_k(Y)},
\qquad
\bar\varepsilon:=E_0\,\varepsilon_1\,C_1\,M^{\bar q}\,e^{C_2\kappa_1},
\end{equation}
where $C_1$, $C_2$, $\bar q$ are the absolute constants of \cite[(1.36)]{B-CMP116} (these are
not the radius constants $C_0$, $C_1$ of Definition~\ref{5.1:def-analyticity-spaces}).
\item[(f)] If $h$ depends holomorphically on parameters, inside $\mathbb{B}_k$, then so does $\mathcal{V}'(Y,\cdot)[h]$.
\end{enumerate}
Likewise, let $B$ denote the fluctuation field of \cite[Lemma~1]{B-CMP116}, related to the
scaled field by $\mathsf{A}=g_kB$. The $\Psi$-pieces of the step, defined in the lemma of this
book on the $\Psi$- and $\Phi$-pieces of the step, have, multiplied by $g^{-2}$, the form
\begin{equation}\label{5.2:eq-fluctuation-potential-3}
\tfrac12\bigl\langle B,\mathcal{Q}^{\Psi}(Y,\mathsf{A})B\bigr\rangle ,
\end{equation}
with $\mathcal{Q}^{\Psi}(Y,\mathsf{A})$ obeying \cite[(1.43)]{B-CMP109} with a right member linear in $|\mathsf{A}|$, as in \cite[(1.40)]{B-CMP109}; and the $\Phi$-pieces $\mathcal{V}^{\Phi}(Y,\mathsf{A})$ obey the bound \cite[(1.36)]{B-CMP116} and vanish at $\mathsf{A}=0$.
\end{lemma}

\begin{proof}
The assertion is \cite[Lemma~1]{B-CMP116}, read with the given history $h$ in place of the families $E^{(j)}$ of Ba{\l}aban's induction and in the variable $\mathsf{A}$ introduced in step (iii) below. The proof has three steps: (i) the logical form of Ba{\l}aban's argument, which identifies what that lemma assumes of the $E^{(j)}$; (ii) the nine stages of the argument in \cite{B-CMP109} and \cite{B-CMP116} that lead to it, with what each stage reads of the input families $H^{(j)}$; (iii) the passage to the variable $\mathsf{A}$.

\smallskip
\noindent\textit{(i) Logical form.}
Ba{\l}aban's Theorem~3 \cite[Theorem~3]{B-CMP109} is proved as the implication ``the inductive assumptions hold at all levels $\le k$ $\Rightarrow$ they hold at level $k+1$''. This is how the paper is organised: the section describing the effective actions ends with ``This completes the description of the inductive assumptions''; the inductive step opens with ``We assume that after $k$ steps we have obtained the action $A_k$ described in the previous section''; the representation reached in it is described by ``all terms in this representation satisfy the required properties, except possibly the last term (2.13)''; and ``the proof of Theorem~3 is reduced to \dots\ the representation (1.7) with terms having the analytic extensions satisfying the bound (1.18)'' \cite{B-CMP109}. The second paper of the series closes with the sentence ``completes the proof of the inductive assumptions for the action $A_{k+1}$'' \cite{B-CMP116}. In \cite{B-CMP119} the same reading is stated outright: ``the space of all densities satisfying the conditions of the inductive assumption \dots\ RT transforms the space with the index $k$ into the space with the index $k+1$. This generalization does not seem to be useful, or interesting now.''

The inductive assumptions placed on the families $E^{(j)}$, $j\le k$, are: condition (f1) of Definition~\ref{5.1:def-analyticity-spaces}; the bound \cite[(1.18)]{B-CMP109}; conditions (f2) and (f3) of the same definition; and the identity $\beta_j=\beta[E^{(j)}]$ for the coefficient of the counterterm. Apart from the reality of the values, these are the conditions defining $\mathbb{B}_k$ in \eqref{5.1:eq-analyticity-spaces-c}; the last of them is taken over, for $h\in\mathbb{B}_k$, by the definition of $h^{\mathrm{ren}}$, as stage (B9) below shows. The statement \cite[Lemma~1]{B-CMP116} is an intermediate statement of this implication, its hypotheses being the inductive assumptions at levels $\le k$, and it is the statement claimed here. It therefore suffices to show that the argument leading to it reads of the $E^{(j)}$ only (f1), (f2), (f3), the norm bound and the functional $\beta[\cdot]$; that it reads them linearly; and that every estimate in it remains valid for complex values.

\smallskip
\noindent\textit{(ii) What each stage reads of the input.}
We go through the stages of the argument in their order in \cite{B-CMP109} and \cite{B-CMP116}.

\smallskip
\noindent (B1) In \cite[(3.3)--(3.4)]{B-CMP109} the input families are composed with the shift of the configuration and telescoped in Ba{\l}aban's interpolation parameters $t_{\square}$. This is an identity and uses nothing of the input but the gauge invariance (f2), invoked in the words ``Making use of the gauge invariance''.

\smallskip
\noindent (B2) For localisation domains $X\not\subset\tilde\square^2$ the argument runs through \cite[(3.7)]{B-CMP109}; the extension \cite[(3.13)]{B-CMP109} is analytic on the spaces \cite[(3.14)]{B-CMP109}, and ``on these spaces \dots\ the function (3.13) satisfies the inequality (1.18)''; the Cauchy estimate \cite[(3.15)]{B-CMP109} follows, and the bound \cite[(3.17)]{B-CMP109} has right member $E_0e^{-\kappa d_j(X)}$ times the factor $6\alpha_2^{-1}L^j\eta B_0B_3$ and the further constant factors printed there. What is read of the input is the analyticity (f1) and the norm bound. Ba{\l}aban states this himself: the bound ``follows only from the fact that $H_j(B(t))$ satisfies (3.14) with $\frac13\alpha_2$, and $\delta H_j$ satisfies (3.9)''.

\smallskip
\noindent (B3) For $X\subset\tilde\square^2$ the argument is \cite[(3.18)--(3.28)]{B-CMP109}, which consists of gauge changes and reads (f2). The lemma invoked there, \cite[Lemma~4]{B-CMP109}, is a statement about maps of configurations and reads nothing of the input.

\smallskip
\noindent (B4) The Taylor expansion \cite[(3.34)]{B-CMP109} is followed by the estimate of the remainder \cite[(3.54)]{B-CMP109} by Cauchy's formula; this reads the norm bound. The derivative in $t_{\square}$ ``yields a factor $O(1)\alpha_2^{-1}\varepsilon_1$''.

\smallskip
\noindent (B5) In \cite[(4.1)--(4.5)]{B-CMP109} the terms are split into multilinear parts, whose kernels are estimated by Cauchy's formula, with bounds of the form ``$E_0\exp(-\kappa d_j(X))\exp(-\delta_0\mathrm{dist}\dots)$''. This reads the norm bound.

\smallskip
\noindent (B6) The Ward--Takahashi identities used next hold for ``a function $E(V)$ defined and analytic on a domain of small gauge field configurations \dots\ gauge invariant'', and \cite[(4.9)]{B-CMP109} is stated as ``holding for all $\mathfrak{g}^{c}$-valued functions $\lambda$, and small, $\mathfrak{g}^{c}$-valued configurations $B$''; the consequence \cite[(4.14)]{B-CMP109} is drawn from the semisimplicity of the Lie algebra. This reads (f1) and (f2).

\smallskip
\noindent (B7) The relations \cite[(4.20)--(4.34)]{B-CMP109} are algebra starting from \cite[(4.15)]{B-CMP109}; the bounds \cite[(4.22)]{B-CMP109} are obtained from \cite[(4.17)--(4.18)]{B-CMP109}, and \cite[(4.33)]{B-CMP109} from invariance under constant gauge transformations. This reads (f2) and the norm bound.

\smallskip
\noindent (B8) The resummation \cite[(4.35)--(4.37)]{B-CMP109} is followed by the analysis of
Ba{\l}aban's quantity $\Pi$ in \cite[Section~5]{B-CMP109}, carried out ``using the above
properties only''. There \cite[(5.12)--(5.15)]{B-CMP109} are derived from
\cite[(5.2)]{B-CMP109}, which uses the Euclidean invariance (f3), from the symmetry of a second
derivative, from \cite[(5.9)]{B-CMP109}, which uses (f2), and from \cite[(5.10)]{B-CMP109},
which uses the norm bound. The relations \cite[(5.30)--(5.42)]{B-CMP109} are linear equations between Taylor coefficients. The reality of the values is remarked in that section and is never used.

\smallskip
\noindent (B9) The counterterm is $\beta_j$ times the explicit expansions \cite[(4.42), (4.44), (4.45)]{B-CMP109}, and ``both groups of terms cancel'' because $\beta_j$ is ``equal to this coefficient''. In the present setting this equality holds by the definition of $h^{\mathrm{ren}}$ in \eqref{5.1:eq-analyticity-spaces-c}, which subtracts from each family $F$ the term $\beta[F]A$.

\smallskip
After these stages the argument of \cite[(1.8)--(1.32)]{B-CMP116} performs the decoupling in
the parameters $\mathsf{s}(\Delta)$, the Cauchy estimates on circles in these parameters, the sums over localisation domains $X$ in \cite[(1.26)]{B-CMP116}, the sums over the cubes $\square'$ (``the factor $(L^j\eta)^5$ \dots\ $(6L)^4L^j\eta$, and the sum over $j$ is bounded by $2(6L)^4$''), and the sums over $Y_0$ and $\square_0$. None of these reads anything of the input beyond the bounds already obtained.

\smallskip
Every stage is thus a substitution, a differentiation, an integration against kernels independent of the $H^{(j)}$, or an application of the linear functional $\beta[\cdot]$. Each of these operations is linear in the input, and so is their composition; hence $\mathcal{V}'(Y,\cdot)[h]$ is linear in $h$, which is (a). Every estimate bounds a modulus by Cauchy's formula in terms of $\sup|H^{(j)}(X,\cdot)|$, that is, in terms of the norm $\|H^{(j)}\|_j\le E_0$ of \eqref{5.1:eq-analyticity-spaces-c}; the argument therefore applies verbatim to complex-valued families, and it delivers the analyticity (b), the properties (c), the identity \eqref{5.2:eq-fluctuation-potential-1} and the bound \eqref{5.2:eq-fluctuation-potential-2} exactly as \cite[Lemma~1]{B-CMP116} and \cite[(1.36)]{B-CMP116} deliver them for the real families $E^{(j)}$. Finally, let $h$ depend holomorphically on parameters inside $\mathbb{B}_k$. By linearity, integration of $h$ over a closed contour in a parameter commutes with each stage, and the bounds above, which hold uniformly on $\mathbb{B}_k$, justify the interchange; so the contour integrals of $\mathcal{V}'(Y,\cdot)[h]$ vanish, and Morera's theorem gives (f).

\smallskip
\noindent\textit{(iii) The variable $\mathsf{A}$.}
The expression localised in \eqref{5.2:eq-fluctuation-potential-1} depends on the pair $(g_k,B)$ only through the product $g_kB$ \cite[(3.2)]{B-CMP109}. Moreover, \cite[Lemma~1]{B-CMP116} treats the fluctuation as a free variable: ``in the preceding considerations $B'$ was a variable field'' \cite{B-CMP116}. The domain of that lemma in the fluctuation variable is $\{|B|<\varepsilon_1g_k^{-1}\}$, a complex polydisc, which the substitution $B=\mathsf{A}/g_k$ maps onto $\{|\mathsf{A}|<\varepsilon_1\}$. Writing $V'_k$ for the functions of \cite[Lemma~1]{B-CMP116}, we set
\begin{equation*}
\mathcal{V}'(Y,\cdot,\mathsf{A}):=V'_k(Y,\cdot,\mathsf{A}/g_k).
\end{equation*}
With this definition the properties obtained in step (ii) become (a)--(f) in the variable $\mathsf{A}$.

\smallskip
\noindent\textit{The $\Psi$- and $\Phi$-pieces.}
The second assertion is the third clause of the lemma of this book on the $\Psi$- and
$\Phi$-pieces of the step, read in the variable $\mathsf{A}=g_kB$ exactly as in step (iii).
\end{proof}

\begin{definition}[Constants of the step and thresholds at the output rate]
\label{5.2:def-step-constants}
Let $\delta$ be the margin parameter of \cite[(1.32)]{B-CMP116}, fixed with
$0<\delta\le\delta_J(L)$, where $\delta_J(L):=(1-4/L)/10$, and put
\begin{equation*}
\rho_m:=(1-m\delta)\tfrac12 L\kappa\quad(m=1,\dots,10),\qquad
\mathsf{r}(\delta):=(1-10\delta)\tfrac12 L ,
\end{equation*}
so that $\rho_{10}=\mathsf{r}(\delta)\kappa$ and $\rho_m\le\frac12L\kappa$ for every $m$. The
constants $\varepsilon_1,\alpha_4,\alpha_6,\gamma_2,\kappa_1$ are those of the list in
\cite[Lemma~3]{B-CMP116}, the small-field cut of the step being $\delta_k=g_kT(g_k)$ with
$T(g)=A_1(\log g^{-2})^{p_0}$. The numbers $C_1,C_2$ are the absolute constants of
\cite[(1.36)]{B-CMP116}, and $\bar q$ is the exponent of $M$ there, printed $q$ in
\cite[(1.36)]{B-CMP116}; $E_0,\alpha_0,\alpha_1,\alpha_5,\delta_0$ are the constants so
named in \cite{B-CMP116}. Every $O(1)$ below is an absolute constant; the one in the
definition of $C_\sigma$ is taken $\ge 1$, and the $O(1)$ in $C_3$, in (d2) and in the
merged factor $a$ of the proof is one and the same constant, that of the bracket of
\cite[(2.37)]{B-CMP116}.

\begin{enumerate}
\item[(a)] The quantity $\varepsilon_2$ of \cite[(2.28), p.~19]{B-CMP116} and the constant
$C_3$ of \cite{B-CMP116} are written
\begin{equation}\label{5.2:eq-step-constants-1}
\begin{aligned}
\varepsilon_2^{\mathrm{pr}}&:=2E_0\varepsilon_1C_1\alpha_4^{-1}\alpha_6^{-1}M^{\bar q}
e^{C_2\kappa_1},\\
C_3&:=2(L+2)^4O(1)\,2E_0C_1\alpha_4^{-1}\alpha_6^{-1}M^{\bar q}e^{C_2\kappa_1},
\end{aligned}
\end{equation}
so that $C_3\varepsilon_1=2(L+2)^4O(1)\,\varepsilon_2^{\mathrm{pr}}$; in particular $C_3$ is
proportional to $E_0$.
\item[(b)] For a kernel $\mathsf{K}(b,b')$ on bonds, with $b_-$ the initial point of the
oriented bond $b$,
\begin{equation*}
\|\mathsf{K}\|_\delta:=\sup_{b,b'}|\mathsf{K}(b,b')|\,e^{\frac12\delta_0|b_--b'_-|}.
\end{equation*}
The symbols $C^{(k)}(Z_0,\sigma)$ and $\Gamma_k(Z_0,\sigma)$ denote the operators of
\cite[(2.14)]{B-CMP116}, $Z_0$ a localisation domain and $\sigma$ the decoupling
parameters; the argument $0$ means vanishing decoupling parameters.
\item[(c)] The constants of the step are
\begin{equation}\label{5.2:eq-step-constants-a}
\begin{aligned}
E^*&:=O(1)\,\alpha_6^{-1}C_3\varepsilon_1,\qquad
\varepsilon_{\mathrm{pr}}:=O(1)e^{-\frac13\delta_0M}+O(\alpha_0+\alpha_1),\\
C_\sigma&:=4\Bigl(1+O(1)\sup_{k,Z_0}\bigl[\|C^{(k)}(Z_0,0)\|_\delta
+\|C^{(k)}(Z_0,0)^{-1}\|_\delta+\|\Gamma_k(Z_0,0)\|_\delta\bigr]\Bigr),\\
\sigma_0&:=\min\bigl\{\tfrac12,\ (C_\sigma^3(LM)^4)^{-1}\bigr\},\qquad
\varepsilon^s:=C_\sigma\bigl(2\varepsilon_{\mathrm{pr}}+C_\sigma\sigma_0/4\bigr),\\
\alpha_5^s&:=\alpha_5+O(1)\,\varepsilon^s .
\end{aligned}
\end{equation}
Here $\varepsilon_{\mathrm{pr}}$ is the constant of \cite[(2.16)]{B-CMP116}; $C_\sigma$ is
finite and independent of $k$. The letter $s$ denotes the complex Gaussian parameter,
$|s-1|<\sigma_0$, of the cluster expansion of the fluctuation integral, and $\varepsilon^s$,
$\alpha_5^s$ are the error constants at such $s$.
\item[(d)] The following thresholds are imposed.
\begin{enumerate}
\item[(d1)] $E^*\le E_0/16$.
\item[(d2)] In place of the assumption $2(L+2)^4O(1)\varepsilon_2e^{5\kappa}\le1$ of
\cite[(2.37)]{B-CMP116}:
\begin{equation}\label{5.2:eq-step-constants-2}
2(L+2)^4O(1)\,\alpha_6^{-1}\varepsilon_2^{\mathrm{pr}}\,e^{5\rho_7}\le 1 .
\end{equation}
\item[(d3)] In place of the smallness of $C_3\varepsilon_1e^{5\kappa}$ in
\cite[(2.39)--(2.40)]{B-CMP116}: the activity $\alpha_6^{-1}C_3\varepsilon_1e^{5\rho_9}$ is
small enough for the Koteck\'y--Preiss criterion \cite{KP86} in the cluster expansion of
\cite[(2.39)--(2.40)]{B-CMP116}.
\end{enumerate}
We impose (d2) and (d3) with $e^{\frac52L\kappa}$ in place of $e^{5\rho_7}$ and
$e^{5\rho_9}$; since $\rho_m\le\frac12L\kappa$, this implies them for every $\delta$ in
$]0,\delta_J(L)]$. The assumption $\varepsilon_2^{\mathrm{pr}}e^{5\kappa}\le1$ of
\cite[(2.28)]{B-CMP116} is kept as printed.
\item[(e)] Order of choice. (d1)--(d3) are ceilings on $\varepsilon_2^{\mathrm{pr}}$, which is
proportional to $\varepsilon_1$, hence ceilings on $\varepsilon_1$. They are imposed after
$E_0$, $\alpha_4$, $\alpha_6$, at the place of $\varepsilon_1$, last in the list of
\cite[Lemma~3]{B-CMP116}, and through the small-field cut $\delta_k=g_kT(g_k)$ they are
ceilings on $\gamma$.
\end{enumerate}
\end{definition}

\begin{proof}[Verification of the claims of Definition~\ref{5.2:def-step-constants} and
derivation of the thresholds]
We establish the properties asserted in Definition~\ref{5.2:def-step-constants} and show
why the thresholds have the stated form.

\emph{The merging count.} Three smallness assumptions of \cite{B-CMP116} use the inequality
\cite[(2.27)]{B-CMP116}, $\sum_Y(d(Y)+5)\ge d(Y_0)+5$, to merge $n$ exponentials in the sizes
of localisation domains into a single one: the assumption of \cite[(2.28)]{B-CMP116}
(rate $(1-4\delta)\kappa$, level $k$), that of \cite[(2.37)]{B-CMP116} (rate $\rho_7$,
level $k+1$), and that on $C_3\varepsilon_1e^{5\kappa}$ in \cite[(2.39)--(2.40)]{B-CMP116}
(rate $\rho_9$). Let $Y_1,\dots,Y_n$ be merged into $Y_0$ as in \cite[(2.27)]{B-CMP116}, and
let $a>0$, $\rho\ge0$. Write $d_i:=d(Y_i)$ and $d(\cup):=d(Y_0)$. By
\cite[(2.27)]{B-CMP116},
\begin{equation*}
\sum_{i\le n}d_i\ \ge\ d(\cup)+5-5n\ =\ d(\cup)-5(n-1).
\end{equation*}
Hence
\begin{equation}\label{5.2:eq-step-constants-3}
\prod_{i\le n}a\,e^{-\rho d_i}=a^n e^{-\rho\sum_i d_i}
\le a^ne^{-\rho d(\cup)+5\rho(n-1)}
=\bigl(a e^{5\rho}\bigr)^{n-1}\,a\,e^{-\rho d(\cup)} .
\end{equation}
Every merged domain beyond the first therefore costs $e^{5\rho}$, with $\rho$ the rate in
play; the term $n=1$ pays nothing. If $ae^{5\rho}\le1$, the product is at most
$a\,e^{-\rho d(\cup)}$.

\emph{The merge cost at level $k+1$.} At level $k$ the rate is
$(1-4\delta)\kappa\le\kappa$, so the cost in \eqref{5.2:eq-step-constants-3} is at most the
factor $e^{5\kappa}$ of the assumption of \cite[(2.28)]{B-CMP116}. The paper
\cite[(2.37)]{B-CMP116} states its assumption with the factor $e^{5\kappa}$ and with
$\varepsilon_2$; the rate in play there is $\rho_7$. At the value of $\delta$ fixed
in \cite{B-CMP116}, for which $(1-10\delta)\frac12L=1$ and the rate of
\cite[(2.41)]{B-CMP116} is $\kappa$, that is $\delta=(1-2/L)/10$, we have
$1-7\delta=(3+14/L)/10$. Therefore
\begin{equation*}
\rho_7=\frac{3L+14}{20}\,\kappa=\Bigl(1+\frac{3(L-2)}{20}\Bigr)\kappa>\kappa .
\end{equation*}
Since $\rho_7$ decreases in $\delta$ and $\delta_J(L)<(1-2/L)/10$, the excess is larger for
$\delta\le\delta_J(L)$. The rate in play at \cite[(2.37)]{B-CMP116} is $\rho_7$, and the one
at \cite[(2.39)--(2.40)]{B-CMP116} is $\rho_9$; the forms imposed here are therefore (d2)
with $e^{5\rho_7}$ and (d3) with $e^{5\rho_9}$.

\emph{The second factor $\alpha_6^{-1}$.} By \cite[(2.29)]{B-CMP116},
\begin{equation*}
\sum_D\prod_{Y\in D}\alpha_6\exp(-\delta\kappa d_k(Y))\le1,
\end{equation*}
the sum running over families $D$ of distinct localisation domains of level $k$,
for $\kappa$ sufficiently large and $\alpha_6$ sufficiently small. The factor $\alpha_6^{-1}$
in $\varepsilon_2^{\mathrm{pr}}$ is the one extracted for this estimate at level $k$. In the
passage leading to \cite[(2.37)]{B-CMP116}, the sum over all families of distinct
localisation domains $Z_i'$ is again estimated by \cite[(2.29)]{B-CMP116}, now at level
$k+1$. For this, each merged factor must lend $\alpha_6e^{-\delta\frac12L\kappa d}$, where
$d=d_{k+1}(Z')$ denotes the size of the localisation domain $Z'$ concerned. Since
$(1-6\delta)\frac12L=\delta\frac12L+(1-7\delta)\frac12L$,
\begin{equation*}
\begin{aligned}
&2(L+2)^4O(1)\varepsilon_2^{\mathrm{pr}}e^{-(1-6\delta)\frac12L\kappa d}\\
&\qquad=\bigl[\alpha_6e^{-\delta\frac12L\kappa d}\bigr]\cdot
\bigl[2(L+2)^4O(1)\varepsilon_2^{\mathrm{pr}}\alpha_6^{-1}e^{-\rho_7 d}\bigr].
\end{aligned}
\end{equation*}
The first bracket is spent in \cite[(2.29)]{B-CMP116}. The second is what
\eqref{5.2:eq-step-constants-3} merges, with $a=2(L+2)^4O(1)\alpha_6^{-1}
\varepsilon_2^{\mathrm{pr}}$ and $\rho=\rho_7$. Under \eqref{5.2:eq-step-constants-2} we
have $ae^{5\rho_7}\le1$, and the factor kept after the merge is
\begin{equation*}
a=\frac{2(L+2)^4O(1)\,\varepsilon_2^{\mathrm{pr}}}{\alpha_6}=\frac{C_3\varepsilon_1}{\alpha_6},
\end{equation*}
by \eqref{5.2:eq-step-constants-1}. The paper \cite[(2.37), p.~20]{B-CMP116} states the
bracket and the constant $C_3$ without this factor $\alpha_6^{-1}$ and with the factor
$e^{5\kappa}$; here \eqref{5.2:eq-step-constants-2} is used.

\emph{The cluster expansion at \cite[(2.39)--(2.40)]{B-CMP116}.} The activity there is the
kept factor times the merge cost at rate $\rho_9$, namely
$\alpha_6^{-1}C_3\varepsilon_1e^{5\rho_9}$. The smallness required of it is that of the
Koteck\'y--Preiss criterion \cite{KP86}, which is (d3). The volume factor
$e^{O(1)(LM)^{-4}|X|}$ is not
charged to this activity. It is paid separately by $e^{-\frac12\delta L\kappa d_{k+1}(X)}$
together with \cite[(2.30)]{B-CMP116}. This uses a lower bound on $\kappa$ that is part of
the requirement $\kappa\ge C_\delta(4)/\delta$, the numerator $C_\delta(4)$ depending only on
the dimension. Since
$\rho_7,\rho_9\le\frac12L\kappa$, the versions of (d2) and
(d3) with $e^{\frac52L\kappa}$ imply them for every admissible $\delta$. Both conditions
bound $\varepsilon_2^{\mathrm{pr}}$, which by \eqref{5.2:eq-step-constants-1} is a multiple of
$\varepsilon_1$ by a factor fixed once $E_0$, $\alpha_4$, $\alpha_6$, $\kappa_1$, $M$ are
fixed. This gives item~(e).

\emph{The output constant.} With (d2) and (d3), the bounds of \cite[Lemma~3]{B-CMP116} read
\begin{equation*}
\begin{aligned}
|H(Z)|&\le\alpha_6^{-1}C_3\varepsilon_1\,e^{-\rho_8 d_{k+1}(Z)},\\
|E^{(k+1)}(X)|&\le O(1)\,\alpha_6^{-1}C_3\varepsilon_1\,e^{-\mathsf{r}(\delta)\kappa\,
d_{k+1}(X)},
\end{aligned}
\end{equation*}
where $H(Z)$ denotes the polymer activity of \cite[(2.11), Lemma~3]{B-CMP116},
the second bound being \cite[(2.41)]{B-CMP116}. With $E^*$ as in
\eqref{5.2:eq-step-constants-a}, therefore,
$|E^{(k+1)}(X)|\le E^*e^{-\mathsf{r}(\delta)\kappa d_{k+1}(X)}$. Compared with the printed
constant $O(1)C_3\varepsilon_1$ of \cite[(2.41)]{B-CMP116}, $E^*$ carries one more absolute
factor, $\alpha_6^{-1}$. Because
$E^*$ is fixed after $\varepsilon_1$, it enters later only through ceilings on $\gamma$.

If instead every factor, including the first, is charged the merge cost $e^{5\rho}$, the
kept factor also carries $e^{5\rho_9}$. Under \eqref{5.2:eq-step-constants-3} this
is not needed, since the term $n=1$ pays no merge.

The same count may be run with a constant in place of $\varepsilon_2^{\mathrm{pr}}$ that
carries no factor $\alpha_6^{-1}$. Both extractions of
$\alpha_6$ are then explicit, and the member of \cite[(2.37)]{B-CMP116} and the output
constant carry $\alpha_6^{-2}$.

\emph{The remaining clauses.} The thresholds of the lemma bounding the output
kernels at the rate $\mathsf{r}(\delta)\kappa$, and of the final step of the proof of the lemma
of this section on the cluster expansion of the fluctuation integral at complex Gaussian
parameter $s$ (which applies \cite[(2.27)--(2.41)]{B-CMP116}), are imposed with
\eqref{5.2:eq-step-constants-2} and (d3) in place of the corresponding assumptions of
\cite{B-CMP116}. Their displays are written in terms of $\delta$ and
hold for every $\delta$ with $0<\delta\le\delta_J(L)$; the particular value
$\delta=(1-2/L)/10$, at which $\mathsf{r}(\delta)=1$, plays no role.

\emph{Condition (d1).} In \cite[p.~21]{B-CMP116} it is assumed that
$O(1)C_3\varepsilon_1\le\frac12E_0$. The paper remarks that this assumption is inessential,
because $C_3\varepsilon_1$ is small anyway and $E_0$ can be taken so that it holds. The
constant $E^*$ is compared with $E_0$ in the same way, by (d1). Condition (d1) is stricter
than that assumption and is a choice made here, not the printed condition. Since $C_3$ is
proportional to $E_0$, (d1) is a
ceiling on $\varepsilon_1$ once $E_0$ and $\alpha_6$ are fixed. It gives in particular
\begin{equation*}
\tfrac12E_0+E^*\le\tfrac12E_0+\tfrac1{16}E_0=\tfrac9{16}E_0 .
\end{equation*}

\emph{The constant $C_\sigma$.} The operators $C^{(k)}(Z_0,0)$, $C^{(k)}(Z_0,0)^{-1}$ and
$\Gamma_k(Z_0,0)$ are the kernels about which \cite[(2.16)]{B-CMP116} perturbs. Their
$\|\cdot\|_\delta$-bounds are those of \cite{B-CMP99b}, uniform in $k$ and $Z_0$. Hence
$C_\sigma$ is finite and independent of $k$. Since the absolute constant in its definition
is at least $1$,
\begin{equation*}
\|C^{(k)}(Z_0,0)\|_\delta,\ \|C^{(k)}(Z_0,0)^{-1}\|_\delta,\ \|\Gamma_k(Z_0,0)\|_\delta
\ \le\ C_\sigma/4 .
\end{equation*}
The norm $\|C^{(k)}(Z_0,0)\|_\delta$ is included in $C_\sigma$ because the lemma on the
cluster expansion of the fluctuation integral at complex Gaussian parameter $s$ uses it
twice. It enters the determinant $\det(1+C^{(k)}(Z_0,0)\mathsf{E})$,
where $\mathsf{E}$ is the perturbation of $C^{(k)}(Z_0,0)^{-1}$ considered there. It also
enters the condition $\alpha_5^s\|C^{(k)}(Z_0,0)\|<\frac12$, where the norm is the operator
norm; by the Schur test, $\|\mathsf{K}\|\le c\,\|\mathsf{K}\|_\delta$ for every kernel
$\mathsf{K}$, with $c$ depending only on $\delta_0$ and the dimension, and the condition is
estimated through $\|\cdot\|_\delta$ below. Its inclusion is of the same
nature as the condition $\gamma_2\le(LM)^{-4}$ of \cite{B-CMP116}.

\emph{The cube in $\sigma_0$.} The bound $\|C^{(k)}(Z_0,0)\|_\delta\le C_\sigma/4$ multiplies
the errors due to $s$ once in the determinant step of the proof of the lemma on the cluster
expansion of the fluctuation integral at complex Gaussian parameter $s$; this is why
$\varepsilon^s$ contains the term $C_\sigma^2\sigma_0/4$. It multiplies them once more in the
printed condition $\alpha_5\|C^{(k)}(Z_0,0)\|<\frac12$ of the real Gaussian computation of
that lemma, with $\alpha_5^s$ in place of $\alpha_5$. The part of
$\alpha_5^s\|C^{(k)}(Z_0,0)\|_\delta$ due to $s$ is therefore $O(1)C_\sigma^3\sigma_0/4$.
Since $\sigma_0\le(C_\sigma^3(LM)^4)^{-1}$, it satisfies
\begin{equation*}
O(1)\,C_\sigma^3\sigma_0/4\ \le\ O(1)/\bigl(4(LM)^4\bigr).
\end{equation*}
The exponent $3$ is used nowhere else.
\end{proof}

\begin{lemma}[The fluctuation integral at a complex Gaussian parameter]
\label{5.2:lem-complex-fluctuation}
Let $T\ge P_0$, where $P_0$ is the lower bound on the threshold fixed in the preparatory
proposition of the step. Let $\sigma_0$, $E^{*}$ and $\gamma_2$ be the constants of
Definition~\ref{5.2:def-step-constants} (see \eqref{5.2:eq-step-constants-a}). Let
$D(1,\sigma_0)\subset\mathbb{C}$ denote the open disc of radius $\sigma_0$ about $1$, and let
$s\in D(1,\sigma_0)$. Let $(a_0,a_1)$ denote either $(\alpha_0,\alpha_1)$ or the pair of
enlarged radii $\hat\alpha_k$ of \eqref{5.1:eq-analyticity-spaces-b}, and let
$\beta_{\mathrm{pr}}=1/5$. For every localisation domain $Y\in\mathfrak{D}_k$ let
$V(Y,\mathbf{U},\mathbf{J},B)$ be analytic on
\begin{equation*}
U^c_{k+1}\bigl(Y,(1+\beta_{\mathrm{pr}})a_0,(1+\beta_{\mathrm{pr}})a_1,a_0\bigr)
\times\{\,|B|<T\ \text{on}\ Y\,\},
\end{equation*}
which is the domain \cite[(1.34)]{B-CMP116} at the threshold $T$, and let $V$ have the form
\cite[(1.42)]{B-CMP116} and satisfy the bounds \cite[(1.43)]{B-CMP116} and
\cite[(1.36)]{B-CMP116} with Ba{\l}aban's constants. Suppressing the dependence on
$(\mathbf{U},\mathbf{J})$, put
\begin{equation*}
\mathsf{I}(s):=\int d\mu_{C^{(k)}/s}(B)\,\prod_b\chi\bigl(|B(b)|<T\bigr)\,
\exp\sum_{Y}V(Y,B),
\end{equation*}
where $d\mu_{C^{(k)}/s}$ is the Gaussian measure with covariance $C^{(k)}/s$ (complex for
non-real $s$), and let $\Phi(X)$, $X\in\mathfrak{D}_{k+1}$, be the functions obtained from
$\mathsf{I}(s)$ by the expansions and the resummation \cite[(2.1)--(2.13)]{B-CMP116}. Then
$\Phi(X)$ is analytic in $(\mathbf{U},\mathbf{J})\in U^c_{k+1}(X,a_0,a_1)$, in
$s\in D(1,\sigma_0)$, and in every parameter on which the $V$ depend analytically with the
bounds above, and
\begin{equation}\label{5.2:eq-complex-fluctuation-a}
\begin{gathered}
|\Phi(X)|\le E^{*}e^{-\kappa d_{k+1}(X)},\qquad
\exp\sum_X\Phi(X)=\mathsf{I}(s)\ \ \text{on a finite torus},\\
|\Phi(X)|\le e^{-\frac14\gamma_2T^2}E^{*}e^{-\kappa d_{k+1}(X)}\ \ \text{if all}\ V\equiv0.
\end{gathered}
\end{equation}
\end{lemma}

\begin{proof}
Throughout, $s^{1/2}$ denotes the principal branch on $D(1,\sigma_0)$; since
$\sigma_0\le\frac12$, $\operatorname{Re}s^{1/2}>0$ there.

\emph{Expansions acting on the integrand.} The operations \cite[(2.1)--(2.4)]{B-CMP116} are
the Mayer expansion of $\exp\sum_YV(Y,B)$, the inclusion--exclusion expansion of the product
of characteristic functions $\chi$, and the resummation of the result into functions
$F(Z_0,B)$, in which the bonds of the set $P$ are required to lie in the interior of $Z_0$
\cite{B-CMP116}. These are identities between integrands. They do not involve the measure, and
they therefore hold unchanged for every value of $s$.

\emph{Conditioning at real positive $s$.} Let $\mathcal{A}:=C^*\Delta_kC$ be the quadratic
form of Ba{\l}aban's fluctuation measure, so that for $s>0$ the measure $d\mu_{C^{(k)}/s}$ is
the Gaussian measure of the positive form $s\mathcal{A}$. Let $I$ be the set of bonds of $Z_0$,
let $E$ be the set of the remaining bonds, and let
\begin{equation*}
C_I:=(\mathcal{A}_{II})^{-1}=C^{(k)}(Z_0).
\end{equation*}
Conditioning on $B_E$ for the positive form $s\mathcal{A}$, that is, the identity
\cite[(2.5)--(2.6)]{B-CMP116} written for $s\mathcal{A}$ in place of $\mathcal{A}$, gives for
every function $F$ of $B_I$
\begin{equation*}
\begin{split}
\int d\mu_{C^{(k)}/s}(B)\,F(B_I)
&=\int d\mu_{C^{(k)}/s}(B')\,
e^{-\frac{s}{2}\langle\mathcal{A}_{IE}B',\,C_I\mathcal{A}_{IE}B'\rangle}\\
&\qquad\cdot\int d\mu_{C_I/s}(B)\,e^{-s\langle B,\mathcal{A}_{IE}B'\rangle}F(B).
\end{split}
\end{equation*}
Substitute $B'=s^{-1/2}(C^{(k)})^{1/2}X$. With Ba{\l}aban's operator $\Gamma_k$, for which
$\mathcal{A}_{IE}(C^{(k)})^{1/2}X=\Gamma_kX$, we get
\begin{equation*}
s\,\mathcal{A}_{IE}B'=s^{1/2}\,\Gamma_kX.
\end{equation*}
The first exponent becomes
\begin{equation*}
-\tfrac{s}{2}\,s^{-1}\langle\Gamma_kX,C_I\Gamma_kX\rangle
=-\tfrac12\langle\Gamma_kX,C_I\Gamma_kX\rangle,
\end{equation*}
which does not depend on $s$.

\emph{The terms of the expansion.} The further steps \cite[(2.8)]{B-CMP116}, leading to
\cite[(2.14), p.~15]{B-CMP116}, are three. They are the random-walk expansions
\cite{B-CMP99b}, the decoupling by the parameters $\mathsf{s}(\Delta)$ attached to the cells
$\Delta$, and the Cauchy representations of the derivatives in these parameters. All three act
on the operators $\Delta_k$, $C^{(k)}(Z_0)$ and $(C^{(k)})^{1/2}$, and a scalar factor is
unaffected by them. Hence the term of the expansion is Ba{\l}aban's term
\cite[(2.14)]{B-CMP116} with $s$ inserted, namely
\begin{equation}\label{5.2:eq-complex-fluctuation-1}
\begin{split}
&\prod_{\Delta}\int d\mathsf{s}(\Delta)\,\frac{1}{2\pi i}
\oint\frac{d\sigma(\Delta)}{(\sigma(\Delta)-\mathsf{s}(\Delta))^2}
\prod_{Y\in D}\int dt(Y)\,\frac{1}{2\pi i}\oint\frac{d\tau(Y)}{(\tau(Y)-t(Y))^2}\\
&\quad\cdot\int d\mu_0(X)\big|_{Z}\,
\exp\Bigl(-\tfrac12\bigl\langle\Gamma_k(Z_0,\sigma)X,\,
C^{(k)}(Z_0,\sigma)\,\Gamma_k(Z_0,\sigma)X\bigr\rangle\Bigr)\\
&\quad\cdot\int d\mu_{C^{(k)}(Z_0,\sigma)/s}(B)\,
\exp\bigl(-s^{1/2}\langle B,\Gamma_k(Z_0,\sigma)X\rangle\bigr)\\
&\qquad\qquad\cdot(-1)^{|P|}\chi_{Y_0}\chi^c_P
\exp\Bigl[\sum_{Y\in D}\tau(Y)V(Y,B)\Bigr].
\end{split}
\end{equation}
Here $D$, $Z$, $Z_0$, $Y_0$, $P$, $\chi_{Y_0}$, $\chi^c_P$, the decoupling parameters
$\sigma(\Delta)$ and the variables $t(Y)$, $\tau(Y)$ are those of \cite[(2.14)]{B-CMP116}, and
the measure $d\mu_0(X)$ is restricted to $Z$ as there. The parameter $s$ occurs in
\eqref{5.2:eq-complex-fluctuation-1} at exactly two places: in the covariance
$C^{(k)}(Z_0,\sigma)/s$ and in the factor $s^{1/2}$ of the last exponent.

\emph{Moduli.} Write $C_{(\sigma)}:=C^{(k)}(Z_0,\sigma)$ and
$\Gamma_{(\sigma)}:=\Gamma_k(Z_0,\sigma)$ for the operators, and $C_0:=C_{(0)}$,
$\Gamma_0:=\Gamma_{(0)}$. We bound the term \eqref{5.2:eq-complex-fluctuation-1} by taking
moduli under the integrals, with $B$ and $X$ real. The result is Ba{\l}aban's bound
\cite[(2.15)]{B-CMP116} with three changes.
\begin{itemize}
\item The inverse covariance $C_{(\sigma)}^{-1}$ is replaced by $sC_{(\sigma)}^{-1}$, both in
the determinant quotient $|\det(\cdot)/\det(\operatorname{Re}\cdot)|^{1/2}$ and in the
measure $d\mu_{(\operatorname{Re}\cdot)^{-1}}$.
\item $\operatorname{Re}\Gamma_{(\sigma)}$ is replaced by
$\operatorname{Re}[s^{1/2}\Gamma_{(\sigma)}]$: for real $B$ and $X$,
\begin{equation*}
\bigl|\exp\bigl(-s^{1/2}\langle B,\Gamma_{(\sigma)}X\rangle\bigr)\bigr|
=\exp\bigl(-\langle B,\operatorname{Re}[s^{1/2}\Gamma_{(\sigma)}]X\rangle\bigr).
\end{equation*}
\item The last factor is $\exp\bigl[\sum_{Y\in D}|\tau(Y)|\,|V(Y,B)|\bigr]$.
\end{itemize}
The phase of $V$ never enters. The variables $\tau(Y)$ are complex already in
\cite[(2.15)]{B-CMP116}, so the last factor has the same form there.

\emph{Errors and determinants.} As in \cite{B-CMP116}, the operators are replaced by the
corresponding operators at $\sigma(Z)=0$, $\mathbf{U}=U$, $\mathbf{J}=0$, and the error is
estimated; in addition we replace $s$ by $1$. The operators $C_0$ and $\Gamma_0$ are taken at
$\sigma=0$, $\mathbf{J}=0$ and at $\mathbf{U}=U$, the real configuration about which the complex
domain $U^c_{k+1}(X,a_0,a_1)$ is taken; they are the operators of Ba{\l}aban's real Gaussian
computation, so $C_0$ and $\Gamma_0$ are real and $C_0$ is positive. Define
\begin{equation*}
\mathsf{E}:=s\bigl(C_{(\sigma)}^{-1}-C_0^{-1}\bigr)+(s-1)\,C_0^{-1},
\qquad\text{so that}\qquad
sC_{(\sigma)}^{-1}=C_0^{-1}+\mathsf{E}.
\end{equation*}
We use the norm $\|\cdot\|_\delta$ and the constants $\varepsilon_{\mathrm{pr}}$, $C_\sigma$
of Definition~\ref{5.2:def-step-constants}. Three facts enter the bound on $\mathsf{E}$:
$|s|\le1+\sigma_0\le\frac32$; by \cite[(2.16)]{B-CMP116},
$\|C_{(\sigma)}^{-1}-C_0^{-1}\|_\delta\le\varepsilon_{\mathrm{pr}}$; and
$\|C_0^{-1}\|_\delta\le C_\sigma/4$ by the definition of $C_\sigma$. Hence
\begin{equation*}
\|\mathsf{E}\|_\delta\le\tfrac32\varepsilon_{\mathrm{pr}}+\sigma_0\|C_0^{-1}\|_\delta
\le\tfrac32\varepsilon_{\mathrm{pr}}+C_\sigma\sigma_0/4.
\end{equation*}
The error operator $R_1$ is Ba{\l}aban's and does not involve $s$. The other two are
\begin{equation*}
-R_2:=\operatorname{Re}\mathsf{E},\qquad
R_3:=\operatorname{Re}\bigl[s^{1/2}(\Gamma_{(\sigma)}-\Gamma_0)\bigr]
+\bigl(\operatorname{Re}s^{1/2}-1\bigr)\Gamma_0 .
\end{equation*}
Since $\operatorname{Re}s^{1/2}>0$, we have $|s^{1/2}+1|\ge1$ and so
$|s^{1/2}-1|=|s-1|/|s^{1/2}+1|\le\sigma_0$. Moreover $\|\Gamma_0\|_\delta\le C_\sigma/4$ by
the definition of $C_\sigma$, and
$\|\Gamma_{(\sigma)}-\Gamma_0\|_\delta\le\varepsilon_{\mathrm{pr}}$ by
\cite[(2.16)]{B-CMP116}. Therefore the bounds \cite[(2.16)]{B-CMP116} hold for $R_1$, $R_2$,
$R_3$ with $\varepsilon_{\mathrm{pr}}$ replaced by
$2\varepsilon_{\mathrm{pr}}+C_\sigma\sigma_0/4$.

For the determinants we use the following elementary bound. Let $\mathsf{N}$ be an $n\times n$
complex matrix all of whose rows have $\ell^1$-norm at most $\mathsf{a}\le\frac12$. Then
\begin{equation*}
e^{-2n\mathsf{a}}\le|\det(1+\mathsf{N})|\le e^{n\mathsf{a}}.
\end{equation*}
For the upper bound, Hadamard's inequality bounds $|\det(1+\mathsf{N})|$ by the product of the
Euclidean norms of the rows of $1+\mathsf{N}$. Each of these is at most the $\ell^1$-norm of
the row, which is at most $1+\mathsf{a}\le e^{\mathsf{a}}$. For the lower bound, $1+\mathsf{N}$
is invertible with $(1+\mathsf{N})^{-1}=1+\mathsf{N}'$, where
$\mathsf{N}'=\sum_{m\ge1}(-\mathsf{N})^m$. The largest row $\ell^1$-norm is submultiplicative,
so the rows of $\mathsf{N}'$ have $\ell^1$-norm at most
$\mathsf{a}/(1-\mathsf{a})\le2\mathsf{a}$. The upper bound applied to $1+\mathsf{N}'$ then
gives $|\det(1+\mathsf{N})|^{-1}\le e^{2n\mathsf{a}}$.

Since $C_0$ is real, $\operatorname{Re}(sC_{(\sigma)}^{-1})=C_0^{-1}+\operatorname{Re}\mathsf{E}$.
The determinant quotient of the moduli bound is therefore
\begin{equation*}
\frac{\det\bigl(sC_{(\sigma)}^{-1}\bigr)}{\det\bigl(\operatorname{Re}(sC_{(\sigma)}^{-1})\bigr)}
=\frac{\det(1+C_0\mathsf{E})}{\det(1+C_0\operatorname{Re}\mathsf{E})},
\end{equation*}
and the quotient entering the measure is treated in the same way. Here $n\le12|Z_0|$. Using
$\|C_0\|_\delta\le C_\sigma/4$, which holds by the definition of $C_\sigma$, the row sums
satisfy
\begin{equation*}
\mathsf{a}\le O(1)\,\|C_0\|_\delta\,\|\mathsf{E}\|_\delta\le O(1)\,\varepsilon^s,
\end{equation*}
where $\varepsilon^s=C_\sigma\bigl(2\varepsilon_{\mathrm{pr}}+C_\sigma\sigma_0/4\bigr)$ is the
constant of \eqref{5.2:eq-step-constants-a}. Consequently the bound
\cite[(2.17)]{B-CMP116} holds in the form $\exp\bigl(O(1)\varepsilon^s|Z_0|\bigr)$.

Next come \cite[(2.18)--(2.20), p.~16]{B-CMP116}. There the radii of the circles in the
variables $\tau(Y)$ are fixed by
\begin{equation*}
\frac{1}{|\tau(Y)|}=E_0\varepsilon_1C_1\alpha_4^{-1}M^{\bar q}e^{C_2\kappa_1}
e^{-(1-3\delta)\kappa d_k(Y)},
\end{equation*}
and it follows that
\begin{equation*}
\sum_{Y\in D}|\tau(Y)|\,|V(Y,B)|
\le\tfrac12O(1)\,\alpha_4\sum_{b\subset Y_0}|B(b)|^2+O(1)\,\alpha_4M^{-4}|Y_0|.
\end{equation*}
These steps use only the form \cite[(1.42)]{B-CMP116} and the bounds
\cite[(1.43), (1.36)]{B-CMP116}, which are the hypotheses of the present lemma on the $V$. The
bound \cite[(2.21)]{B-CMP116} holds in the same way, with the constant $\varepsilon^s$. The
bound \cite[(2.22)]{B-CMP116}, taken at the threshold $T$, is the bound on characteristic
functions that the threshold lemma of the step supplies at $T$; it does not involve $s$.

\emph{The real Gaussian integral.} After these replacements, \cite[(2.23)--(2.26),
p.~17]{B-CMP116} is Ba{\l}aban's real Gaussian computation at
$(\sigma,\mathbf{U},\mathbf{J},s)=(0,U,0,1)$, with the constant $\alpha_5$ of that
computation, of which $\gamma_2$ is one contribution, replaced by
\begin{equation*}
\alpha_5^s:=\alpha_5+O(1)\varepsilon^s
=\alpha_5+O(1)\,C_\sigma\varepsilon_{\mathrm{pr}}+O(1)\,C_\sigma^2\sigma_0 .
\end{equation*}
The computation runs as follows. The $B$-integral is \cite[(2.24)]{B-CMP116}. Its term of order
zero cancels the first quadratic form in the first exponential of \cite[(2.23)]{B-CMP116};
both of these are independent of $s$. The remainder is $\frac12O(\alpha_5^s)\|ZX\|^2$ in the
notation of \cite[(2.23)]{B-CMP116}. Finally, \cite[(2.25)]{B-CMP116} is bounded by
$\exp\bigl(O(\alpha_5^s)|Z|\bigr)$.

Ba{\l}aban's computation requires two conditions \cite{B-CMP116}: that
$\alpha_5\|C^{(k)}(Z_0,0)\|<\frac12$, and that $(LM)^4\alpha_0$, $(LM)^4\alpha_1$,
$(LM)^4\alpha_4$, $(LM)^4\gamma_2,\dots$ are bounded by a constant. Both hold with $\alpha_5$
replaced by $\alpha_5^s$, for two reasons. The first is the choice
$\sigma_0\le(C_\sigma^3(LM)^4)^{-1}$. Indeed, the part of
$\alpha_5^s\|C^{(k)}(Z_0,0)\|_\delta$ that comes from $s$ is
\begin{equation*}
O(1)\,C_\sigma^2\sigma_0\cdot C_\sigma/4\le O(1)/\bigl(4(LM)^4\bigr).
\end{equation*}
The second is Ba{\l}aban's own smallness of $\varepsilon_{\mathrm{pr}}$ ($M$ large,
$\alpha$ small).

Hence \cite[(2.26)]{B-CMP116} holds with the factor $e^{-\frac12\gamma_2T^2|P|}$, uniformly in
$s\in D(1,\sigma_0)$. Here $\gamma_2$ is one of the contributions to $\alpha_5$. The factor
$\exp\bigl(\gamma_2\sum_{b\in P}|B(b)|^2\bigr)$ absorbed into the Gaussian integral is
controlled by the real part of the quadratic form. Since $C_0$ is real and positive and
$\operatorname{Re}s\ge1-\sigma_0$, that real part satisfies
\begin{equation*}
\operatorname{Re}\bigl(sC_0^{-1}\bigr)\ge(1-\sigma_0)\,C_0^{-1}.
\end{equation*}
The condition needed for the absorption is therefore contained in
$\alpha_5^s\|C_0\|<\frac12$.

\emph{Continuation in $s$.} Each term \eqref{5.2:eq-complex-fluctuation-1} is a
finite-dimensional integral of an integrand analytic in $s$, and the $s$-free majorant just
integrated bounds it. By Morera's theorem each term is therefore analytic on $D(1,\sigma_0)$.
The same holds for the activities $H^s(Z)$ of \cite[(2.11)]{B-CMP116}, which are finite sums of
such terms. Because the majorant involves $V$ only through $|\tau(Y)|\,|V(Y,B)|$ and the bounds
\cite[(1.43)]{B-CMP116} and \cite[(1.36)]{B-CMP116}, the same argument applies to every
parameter on which the $V$ depend analytically with these bounds.

The integral $\mathsf{I}(s)$ is analytic in $s$, since its domain of integration is bounded. The
polymer identity \cite[(2.11)]{B-CMP116},
\begin{equation*}
\mathsf{I}(s)=\sum\prod_iH^s(Z_i),
\end{equation*}
holds for $s\in\,]1-\sigma_0,1+\sigma_0[$. There $s>0$, and it consists of Ba{\l}aban's
identities for the positive form $s\mathcal{A}$. On a finite torus both sides are analytic on
$D(1,\sigma_0)$, so the identity holds on the whole disc. No contour of integration is
deformed across the sharp cut-offs $\chi$.

\emph{Resummation.} The remaining steps \cite[(2.27)--(2.41)]{B-CMP116} use the terms only
through the bound \cite[(2.26)]{B-CMP116}, which has been established above uniformly in
$s\in D(1,\sigma_0)$. The conclusions of the resummation lemma for these steps therefore hold
verbatim, with the constant $E^{*}$ of \eqref{5.2:eq-step-constants-a} and the rate $\kappa$.
They give the analyticity of $\Phi(X)$ on $U^c_{k+1}(X,a_0,a_1)$, the bound
$|\Phi(X)|\le E^{*}e^{-\kappa d_{k+1}(X)}$, and, on a finite torus, the identity
$\exp\sum_X\Phi(X)=\mathsf{I}(s)$. Together with the continuation in $s$ and in the parameters,
this proves all clauses of the lemma except the last.

For the last clause, suppose that all $V$ vanish identically. Then only the empty family $D$
contributes, and the contributing terms have $P\neq\emptyset$: as Ba{\l}aban remarks, when
$Y_0$ is empty the term carries the exponential factor \cite{B-CMP116}. This factor,
$e^{-\frac12\gamma_2T^2|P|}$ in \cite[(2.26)]{B-CMP116} with $P\neq\emptyset$, yields the
factor $e^{-\frac14\gamma_2T^2}$ in the bound, which is the last clause of
\eqref{5.2:eq-complex-fluctuation-a}.
\end{proof}

\begin{proposition}[The step with a complex factor near one]\label{5.2:prop-complex-factor-step}
Let $k\ge 0$ and put $T=T_k=T(g_k)$, where $T(\cdot)$, $\mathsf{r}$, $\tau$, $\vartheta$ and $\mathrm{e}$ are the
functions of the coupling of \eqref{5.1:eq-coupling-functions-a}, with $\mathsf{r}_k=\mathsf{r}(g_k)$,
$\vartheta_k=\vartheta(g_k)$ and $\mathrm{e}_k=\mathrm{e}(g_k)$. The constants $E^*$ and $\sigma_0$ are those of
\eqref{5.2:eq-step-constants-a}, and $D(c,R)$ denotes the open disc of radius $R$ about $c\in\mathbb{C}$.

For $(\mathfrak{q},s)\in D(0,\mathsf{r}_k)\times D(1,\sigma_0)$ and $h\in\mathbb{B}_k$ (the ball of complex histories of
\eqref{5.1:eq-analyticity-spaces-c}), let $\Phi_{k+1}(\mathfrak{q},s;h)$ be the family $\Phi(X)$,
$X$ in the class $\mathfrak{D}_{k+1}$ of localisation domains at level $k+1$, which
Lemma~\ref{5.2:lem-complex-fluctuation} attaches, with cutoff $T=T_k$, to the potentials
\begin{equation}\label{5.2:eq-complex-factor-step-1}
V(Y):=s\cdot\tfrac12\langle B,\mathcal{Q}^{\Psi}(Y,\mathfrak{q}B)B\rangle+\mathcal{V}^{\Phi}(Y,\mathfrak{q}B)
+\mathcal{V}'(Y,\mathfrak{q}B)[h],\qquad Y\in\mathfrak{D}_k .
\end{equation}
Here $\tfrac12\langle B,\mathcal{Q}^{\Psi}(Y,\mathsf{A})B\rangle$ and $\mathcal{V}^{\Phi}(Y,\mathsf{A})$ are the
$g_k^{-2}$-weighted $\Psi$-pieces and the $\Phi$-pieces of the step, written as functions of the scaled
fluctuation field $\mathsf{A}$, and
$\mathcal{V}'$ is the fluctuation potential of Lemma~\ref{5.2:lem-fluctuation-potential}; the dependence on the
configuration $(\mathbf{U},\mathbf{J})$ is suppressed. Then the following hold.
\begin{enumerate}
\item[(a)] $\Phi_{k+1}$ is a family on $\mathcal{U}$ obeying the first bound of
\eqref{5.2:eq-complex-factor-step-a}; it is analytic in $(\mathfrak{q},s)$ and along holomorphic families in
$\mathbb{B}_k$.
\item[(b)] At $\mathfrak{q}=0$ the second bound of \eqref{5.2:eq-complex-factor-step-a} holds.
\item[(c)] Consequently the third bound of \eqref{5.2:eq-complex-factor-step-a} holds for all
$s\in D(1,\sigma_0)$ and all $h\in\mathbb{B}_k$.
\item[(d)] At $\mathfrak{q}=g_k$, $\Phi_{k+1}$ is an E-family, and its total
$\sum_X\Phi_{k+1}(g_k,s;h)(X)$ is given by the identity that closes \eqref{5.2:eq-complex-factor-step-a}.
\end{enumerate}
\begin{equation}\label{5.2:eq-complex-factor-step-a}
\begin{split}
&\|\Phi_{k+1}(\mathfrak{q},s;h)\|_{k+1}\le E^*,\qquad
\|\Phi_{k+1}(0,s;h)\|_{k+1}\le\vartheta_k,\\
&\|\Phi_{k+1}(g_k,s;h)\|_{k+1}\le\mathrm{e}_k ,\\
&\sum_X\Phi_{k+1}(g_k,s;h)(X)=\log\int d\mu_{C^{(k)}/s}(B)\prod_{b}\chi\bigl(|B(b)|<T_k\bigr)\\
&\qquad\times\exp\Bigl[\,s\,g_k^{-2}\sum_i\Psi_i+\sum_i\Phi_i\\
&\qquad\qquad+\bigl\{h^{\mathrm{ren}}\bigl(U_k(e^{iB'}V^{(k)})\bigr)-h^{\mathrm{ren}}(U_{k+1})\bigr\}\Bigr].
\end{split}
\end{equation}
Here $\Psi_i$, $\Phi_i$ and $V^{(k)}$ are the objects of the step in Lemma~\ref{5.3:lem-coupling-shift}
(the terms $\Phi_i$ of Ba{\l}aban's step are not to be confused with the family $\Phi_{k+1}$), and $B'$ is
the field of the summation identity of Lemma~\ref{5.2:lem-fluctuation-potential} at $\mathsf{A}=g_kB$.
\end{proposition}

\begin{proof}
We first check the hypotheses of Lemma~\ref{5.2:lem-complex-fluctuation}. Fix
$(\mathfrak{q},s)\in D(0,\mathsf{r}_k)\times D(1,\sigma_0)$ and $h\in\mathbb{B}_k$. On the support of
$\prod_b\chi(|B(b)|<T_k)$ we have $|B|<T_k$. Since $\mathsf{r}=\varepsilon_1/(2T)$ by
\eqref{5.1:eq-coupling-functions-a}, it follows that
\begin{equation*}
|\mathfrak{q}|\,|B|<\mathsf{r}_kT_k=\varepsilon_1/2 .
\end{equation*}
Thus $\mathsf{A}:=\mathfrak{q}B$ is a field with values in the complexification $\mathfrak{g}^{c}$ of the Lie
algebra of $\mathrm{SU}(N)$, with $|\mathsf{A}|<\varepsilon_1$ on $Y$.
This is the domain on which $\mathcal{Q}^{\Psi}$ and $\mathcal{V}^{\Phi}$ are analytic, and on which
$\mathcal{V}'(Y,\mathbf{U},\mathbf{J},\mathsf{A})[h]$ is analytic by Lemma~\ref{5.2:lem-fluctuation-potential}.

The bound \cite[(1.43)]{B-CMP116} on the quadratic kernel is linear in $|\mathsf{A}|$, so multiplying the
kernel by $s$ multiplies the bound by $|s|$. On $D(1,\sigma_0)$ we have $|s|\le 3/2$, and
$|\mathsf{A}|<\varepsilon_1/2$ on the support; hence
\begin{equation*}
|s|\,|\mathsf{A}|<\tfrac32\cdot\tfrac{\varepsilon_1}{2}=\tfrac34\,\varepsilon_1<\varepsilon_1 ,
\end{equation*}
and the kernel $s\,\mathcal{Q}^{\Psi}(Y,\mathfrak{q}B)$ obeys that bound with Ba{\l}aban's constant. The
$\Phi$-pieces $\mathcal{V}^{\Phi}(Y,\mathsf{A})$ lie within the bound
\cite[(1.36)]{B-CMP116}. Lemma~\ref{5.2:lem-fluctuation-potential} supplies $\mathcal{V}'(Y,\mathfrak{q}B)[h]$
with its analyticity, its bound, its linearity in $h$, and the transfer to $\mathcal{V}'$ of holomorphic
dependence of $h$ on parameters. The potentials \eqref{5.2:eq-complex-factor-step-1} therefore satisfy the
hypotheses of Lemma~\ref{5.2:lem-complex-fluctuation} with $T=T_k$.

(a) This is the conclusion of Lemma~\ref{5.2:lem-complex-fluctuation} (displayed in
\eqref{5.2:eq-complex-fluctuation-a}). The output $\Phi(X)$ is analytic on the configuration space and in $s$. It
is also analytic in every parameter on which the potentials depend analytically with the required bounds. Here
these parameters are $\mathfrak{q}$, and the parameters of a holomorphic family $h$ in $\mathbb{B}_k$. Finally,
$|\Phi(X)|\le E^*e^{-\kappa d_{k+1}(X)}$, which is the bound $\|\Phi_{k+1}\|_{k+1}\le E^*$.

(b) At $\mathfrak{q}=0$ we have $\mathsf{A}=0$, and all potentials vanish for every $(s,h)$. Indeed,
$\mathcal{V}'$ vanishes at $\mathsf{A}=0$ by Lemma~\ref{5.2:lem-fluctuation-potential}, and $\mathcal{V}^{\Phi}$
vanishes at $\mathsf{A}=0$. The kernel $\mathcal{Q}^{\Psi}$ vanishes there as well, its bound being linear in
$|\mathsf{A}|$. The last conclusion of Lemma~\ref{5.2:lem-complex-fluctuation} (all $V\equiv 0$) gives
$|\Phi(X)|\le e^{-\frac14\gamma_2T_k^2}E^*e^{-\kappa d_{k+1}(X)}$. Since
$\vartheta_k=E^*e^{-\frac14\gamma_2T_k^2}$ by \eqref{5.1:eq-coupling-functions-a}, this is the second bound of
\eqref{5.2:eq-complex-factor-step-a}.

(c) Fix $s\in D(1,\sigma_0)$ and $h\in\mathbb{B}_k$. By (a), $\mathfrak{q}\mapsto\Phi_{k+1}(\mathfrak{q},s;h)$ is
analytic on $D(0,\mathsf{r}_k)$ with values in the sup-normed family, and it is bounded there by $E^*$. We apply
Lemma~\ref{5.2:lem-disc-bounds}, clause (a) of \eqref{5.2:eq-disc-bounds-a}, termwise, with $c=0$,
$R=\mathsf{r}_k$, $M=E^*$ and $z=g_k$. This gives
\begin{equation*}
\|\Phi_{k+1}(g_k,s;h)-\Phi_{k+1}(0,s;h)\|_{k+1}\le 2E^*g_k/\mathsf{r}_k .
\end{equation*}
Adding the bound of (b) and using $\tau_k=g_k/\mathsf{r}_k$ and $\mathrm{e}=2E^*\tau+\vartheta$ from
\eqref{5.1:eq-coupling-functions-a}, we obtain
\begin{equation*}
\|\Phi_{k+1}(g_k,s;h)\|_{k+1}\le 2E^*g_k/\mathsf{r}_k+\vartheta_k=2E^*\tau_k+\vartheta_k=\mathrm{e}_k .
\end{equation*}

(d) Take $\mathfrak{q}=g_k$. An E-family is a family subject to the conditions (f1)--(f3) that enter the
definition of $\mathbb{B}_k$ in \eqref{5.1:eq-analyticity-spaces-c}; beyond the analyticity and the bound of
(a), it remains to verify (f2) and (f3). Condition (f2) is carried through the construction
\cite[(2.1)--(2.13)]{B-CMP116} in the passage of \cite{B-CMP116} that immediately precedes the completion
of the proof of the inductive assumptions for the effective action after $k+1$ steps. For condition (f3),
Ba{\l}aban remarks, in the
discussion of \cite[(2.12)]{B-CMP109}, that all the expressions there, together with the measure, are
invariant. Here the Gaussian measure carries the scalar factor $s$, and the inputs are invariant
(Lemma~\ref{5.2:lem-fluctuation-potential} for $\mathcal{V}'$). Neither circumstance affects the arguments
establishing (f2) and (f3), so both conditions hold for $\Phi_{k+1}(g_k,s;h)$.

For the total, Lemma~\ref{5.2:lem-complex-fluctuation} gives $\exp\sum_X\Phi_{k+1}(g_k,s;h)(X)=\mathsf{I}(s)$
on a finite torus, with $\mathsf{I}(s)$ the integral of $\exp\sum_YV(Y,B)$ against
$d\mu_{C^{(k)}/s}(B)\prod_b\chi(|B(b)|<T_k)$. By the definition of $\mathcal{Q}^{\Psi}$ and
$\mathcal{V}^{\Phi}$, the $g_k^{-2}$-weighted $\Psi$-pieces of the step are
$\tfrac12\langle B,\mathcal{Q}^{\Psi}(Y,\mathsf{A})B\rangle$ and its $\Phi$-pieces are
$\mathcal{V}^{\Phi}(Y,\mathsf{A})$, with $\mathsf{A}=g_kB$; hence
\begin{equation*}
\sum_Y s\cdot\tfrac12\langle B,\mathcal{Q}^{\Psi}(Y,g_kB)B\rangle=s\,g_k^{-2}\sum_i\Psi_i,\qquad
\sum_Y\mathcal{V}^{\Phi}(Y,g_kB)=\sum_i\Phi_i .
\end{equation*}
At $\mathfrak{q}=g_k$ the sum over $Y$ is therefore as follows:
\begin{itemize}
\item the first two terms of \eqref{5.2:eq-complex-factor-step-1} sum to $s\,g_k^{-2}\sum_i\Psi_i+\sum_i\Phi_i$;
\item the third sums to $h^{\mathrm{ren}}(U_k(e^{iB'}V^{(k)}))-h^{\mathrm{ren}}(U_{k+1})$, by the summation
identity of Lemma~\ref{5.2:lem-fluctuation-potential}.
\end{itemize}
Taking the logarithm gives the identity that closes \eqref{5.2:eq-complex-factor-step-a}.
\end{proof}

\begin{proposition}[Box representation, trivial-point containment and the comb gauge;
{\cite[(3.18), (3.25), (3.26), (3.32), Lemma~4, pp.~276--277, (4.18)]{B-CMP109}};
{\cite[(2.28), (3.68)]{B-CMP119}}]
\label{5.2:prop-box-comb-gauge}
Let $j\le n$ and put $t:=L^{j-n}$. The \emph{admissible backgrounds at scale $n$} are the following
configurations:
\begin{itemize}
\item the background $U_k$;
\item the background $U_k(e^{iB'}V^{(k)})$, together with its deformations in Ba{\l}aban's
interpolation parameters (one global parameter and one for each cube of his cover) and in the
decoupling parameters $\mathsf{s}(\Delta)$, taken on the circles of his expansion;
\item the backgrounds $U^{(n')}_k$ of the preparatory integrals of
\cite[(1.51)--(1.63)]{B-CMP122a};
\item the configurations of the chain \cite[(1.51)--(1.58)]{B-CMP122b};
\item the extensions of all of these to pairs $(\mathbf{U},\mathbf{J})$ on Ba{\l}aban's complex tubes.
\end{itemize}
At the localisation site the carrier pairs built on these backgrounds enter the first three, the
fifth and the boundary groups of estimates of that site through the containment statement for the
carriers on the treated components, and through its exterior form on the exterior hulls of the
pointed units of the exterior class.

In addition, let $Q$ be any box; let $\mathbf{U}$ lie in the ball over $Q$ of the lemma on the
polydisc radius inside the analyticity space, enlarged by the fixed factor $c''$; let
$M^n$ be Ba{\l}aban's averaging map to level $n$ (the superscript is a level, not a power of the
parameter $M$), and let $Q^{\sim-2}$ be the box $Q$ with the two
layers of $\pi_n$-cubes nearest to its boundary removed. Then $U_n(Q,M^n(\mathbf{U}))$ is an
admissible background at scale $n$ on $Q^{\sim-2}$, that is, (c1)--(c3) below apply to it for
domains $X$ and boxes $\square_0$ contained in $Q^{\sim-2}$; this is the content of the lemma on
admissibility of box minimisers of averaged data.

For a cube $\square'\in\pi_j$ and an integer $a\ge0$, $\square'^{\sim a}$ denotes $\square'$ enlarged
by $a$ layers of cubes of $\pi_j$. Let $a_j$ stand for either pair of radii of the analyticity
spaces of Definition~\ref{5.1:def-analyticity-spaces}, the absolute pair $(\alpha_0,\alpha_1)$ or
the coupling-dependent pair $(\alpha_{0,j},\alpha_{1,j})$, and let $\alpha_{\cdot}(g_n)$ be as in
Notation~\ref{5.1:not-coupling-functions}. Then the following three statements hold; none of them
involves the complex source $w=\sum_i z_if_i$. Their principal relations are
\begin{equation}\label{5.2:eq-box-comb-gauge-a}
\begin{gathered}
\text{(c1)}\quad \mathfrak{b}\big|_X = U_j(\square_0,\mathsf{V})\big|_X\ \text{modulo gauge},\qquad
|\partial\mathsf{V}-1|\le\mathfrak{a}_n t^2,\\
\text{(c2)}\quad \bigl(U_j(\square_0,e^{iB}),J_j(\square_0,e^{iB})\bigr)\big|_X\in U^c_j(X,a_j)
\quad\text{for } B \text{ complex},\ |B|<\alpha_{3,j},\\
\text{(c3)}\quad B^{\mathrm{ax}}_\mu(x)=\sum_{\kappa<\mu}(x_\kappa-y_\kappa)F_{\kappa\mu}(y)
+\mathfrak{e}_\mu(x),
\end{gathered}
\end{equation}
with $\square_0$, $\mathsf{V}$, $B$, $\alpha_{3,j}$, $B^{\mathrm{ax}}$ and $y$ as specified in
(c1)--(c3) below, in the following precise sense.
\begin{enumerate}
\item[(c1)] (Local form.) Let $X$ be a domain and let $\square_0\supset X$ be a box of one of the
following two kinds:
\begin{itemize}
\item a box about $X$ with one layer of $\pi_j$-cubes of margin, namely $X\subset\square'^{\sim a}$
and $\square_0=\square'^{\sim a+1}$ for some $\square'\in\pi_j$, where either $a=n-j$
(Ba{\l}aban's printed representation) or $a=1+n-j$ (the same statement one cube size up, taken
as a hypothesis, see the remarks after the proposition);
\item any box containing $X$ with two layers of $\pi_j$-cubes of margin.
\end{itemize}
Let $\mathfrak{b}$ be an admissible background at scale $n$, and let $\mathsf{V}$ denote its
level-$j$ data on the determining set of $\square_0$, in the sense of
\cite[(3.18), (3.25)]{B-CMP109}. Then, restricted to $X$, $\mathfrak{b}$
equals, modulo gauge, Ba{\l}aban's box function $U_j(\square_0,\mathsf{V})$ (whose companion in the
pair of \cite[(1.11)--(1.16)]{B-CMP109} is $J_j(\square_0,\mathsf{V})$), and
$|\partial\mathsf{V}-1|\le\mathfrak{a}_nt^2$, where
\begin{equation*}
\mathfrak{a}_n:=C_{\mathfrak{a}}\,\alpha_{\cdot}(g_n)
\end{equation*}
with the constant $C_{\mathfrak{a}}$ fixed as follows. Let $\alpha_0$ be the first absolute radius of
Definition~\ref{5.1:def-analyticity-spaces}, $\alpha_2$ Ba{\l}aban's further radius of
\cite{B-CMP109}, entering the bound \cite[(3.32)]{B-CMP109}, and $B_3$ the constant of
\cite[Proposition~9]{B-CMP102b}. The bound of
\cite[(3.32)]{B-CMP109} is $2(\alpha_2+B_3O(1)M\alpha_0)=C_{\mathfrak{a}}\alpha_0$, since $\alpha_2$
is proportional to $\alpha_0$; $\mathfrak{a}_n$ is this bound with $\alpha_0$ replaced by
$\alpha_{\cdot}(g_n)$.
\item[(c2)] (Containment at the trivial point.) Put $\alpha_{3,j}:=c_3a_j$, where $c_3$ is the
constant of \cite[Lemma~4]{B-CMP109} and the product is taken with the radius of the pair $a_j$ to
which Ba{\l}aban's $\alpha_3$ is proportional. Let $B$ be a complex
configuration on the determining set of $\square_0$ with $|B|<\alpha_{3,j}$. Then the pair
$\bigl(U_j(\square_0,e^{iB}),J_j(\square_0,e^{iB})\bigr)$, restricted to $X$, lies in
$U^c_j(X,a_j)$, and it depends analytically on $B$. Let $\alpha_1$ be the second absolute radius
of Definition~\ref{5.1:def-analyticity-spaces}, and let $\alpha_3$ (an absolute radius, not to be
confused with $\alpha_{3,j}$) and $\beta_{\mathrm{pr}}$ be Ba{\l}aban's further radius and parameter
of \cite[p.~277]{B-CMP109}, where the parameter is written $\beta$. The restrictions under which
Ba{\l}aban proves this are
\begin{align*}
B_3^2\,O(1)M\alpha_0&<\tfrac12\alpha_1, & O(1)M\alpha_1&\le\beta_{\mathrm{pr}}, &
\bigl(4B_3^2\,O(1)M\bigr)^2\alpha_0&\le\beta_{\mathrm{pr}},\\
B_3\alpha_3&<\tfrac12\alpha_1, & 2B_3\alpha_3&\le\beta_{\mathrm{pr}} L^{-2}\alpha_0 . & &
\end{align*}
Each of them bounds a ratio of radii or bounds a radius from above. Being of this kind, all of them
are satisfied by the radii of \cite[(2.28)]{B-CMP119}; Ba{\l}aban assumes them at that place.
\item[(c3)] (Comb gauge.) Let $y$ be a point. Consider Ba{\l}aban's axial gauge centred at $y$ on
the box $\square_0$ of (c1), assumed of side $D\le 5M/t$.
In this gauge $\mathsf{V}=e^{iB^{\mathrm{ax}}}$, and for every bond $b$
\begin{equation}\label{5.2:eq-box-comb-gauge-1}
|B^{\mathrm{ax}}(b)|\le 2d\,(1+|b-y|)\,\mathfrak{a}_n t^2 .
\end{equation}
Moreover $B^{\mathrm{ax}}$ has the expansion in the last line of \eqref{5.2:eq-box-comb-gauge-a}.
Its coefficients satisfy $|F_{\kappa\mu}(y)|\le\mathfrak{a}_nt^2$, and the remainder satisfies
\begin{equation}\label{5.2:eq-box-comb-gauge-2}
|\mathfrak{e}_\mu(x)|\le C\,(1+|x-y|)^2\,\mathfrak{a}_n\,t^{2+\hat\alpha},
\end{equation}
where $\hat\alpha$ is the H\"older exponent of \cite[(4.18)]{B-CMP109}.
\end{enumerate}
\end{proposition}

This is \cite[(3.18), (3.25), (3.26), (3.32), Lemma~4, pp.~276--277, (4.18)]{B-CMP109} together
with \cite[(2.28), (3.68)]{B-CMP119}, with one exception in (c1), the one-layer box at
$a=1+n-j$, and one added family of admissible backgrounds, $U_n(Q,M^n(\mathbf{U}))$ on
$Q^{\sim-2}$; both are described below. Statements (c1)--(c3) are cited from these places, their
proofs being Ba{\l}aban's, and they are used under clause~(i) of the admitted inputs of
Definition~\ref{5.1:def-admitted-inputs}.

Item by item, the references are as follows.
\begin{itemize}
\item The last family of admissible backgrounds, $U_n(Q,M^n(\mathbf{U}))$ on $Q^{\sim-2}$, is not
taken from Ba{\l}aban's papers; it is the content of the lemma on admissibility of box minimisers
of averaged data, which derives it from (c1) applied on nested boxes,
from \cite[Proposition~9]{B-CMP102b} at the trivial configuration $V_0=1$, and
from \cite[(3.27), (3.32)]{B-CMP109}.
\item The representation in (c1) is \cite[(3.18), (3.25)]{B-CMP109}.
\item The one-layer alternative of (c1) is the representation of
\cite[Sect.~3, before (3.68)]{B-CMP119}, written there
for $a=n-j$, that is, for the cube $\square'^{\sim n-j+1}$ about a domain contained in
$\square'^{\sim n-j}$, $\square'\in\pi_j$. At $a=1+n-j$ the one-layer representation is not
printed; it is the same configuration statement one cube size up, and it is taken as such under
that clause of the admitted inputs. It can be dispensed with by using instead the box
$\square'^{\sim a+2}$ of the second alternative, see the next item. With $\nu:=1+n-j$, the price
of this replacement is that the index $\nu+2$ of the enlarged boxes in clauses (iii) and (iv) of
the lemma on the groups cut at the localisation site becomes $\nu+3$, with constants independent
of $K$; no constant of the correlation theorem changes.
\item The two-layer alternative of (c1) is the construction of clause~(iv) in the definition of
the spaces $U^c_j$, \cite[(1.11)--(1.16)]{B-CMP109}, performed there for an arbitrary domain. In
particular, whenever $X\subset\square'^{\sim a}$, the box $\square'^{\sim a+2}$ is admissible in
(c1) under the second alternative, whatever $a$ is.
\item Statement (c2) is \cite[Lemma~4, pp.~276--277]{B-CMP109}, and its restrictions are those
printed on \cite[p.~277]{B-CMP109}.
\item The axial gauge of (c3) is that of \cite[Sect.~F]{B-CMP102b}, \cite[(3.26)]{B-CMP109} and
\cite[(3.68)]{B-CMP119}. In \cite[(3.68)]{B-CMP119} the main term is written as
$\sum_{\kappa<\mu}(x_\kappa-z_\kappa)F_{\kappa\mu}(z)$.
\item The bound \eqref{5.2:eq-box-comb-gauge-2} on the remainder rests on the H\"older bound
\cite[(4.18)]{B-CMP109}.
\end{itemize}

The factor $2d(1+|b-y|)$ in \eqref{5.2:eq-box-comb-gauge-1} has the following origin. In the axial
gauge centred at $y$, a bond variable is a product of at most $d(1+|b-y|)$ conjugated plaquette
variables, and the conjugations distort the bound by a factor at most $e^{O(M\alpha_1)}\le2$.

\begin{remark}[The box representation on a domain's own box]\label{5.2:rem-own-box}
The local form \eqref{5.2:eq-box-comb-gauge-a} of Proposition~\ref{5.2:prop-box-comb-gauge} is
used on three kinds of box $\square_0$ containing a domain $X$ at scale $j$. Ba{\l}aban's
construction uses it on its own boxes: on the fattened cube $\tilde{\square}^5$ of
\cite[(3.18)]{B-CMP109}, at the scale of the field, and in \cite{B-CMP119} on
$\square^{\sim n-j+1}$ about
$X \subset \square^{\sim n-j}$, that is, with one $\pi_j$-layer of margin. The marginal bound uses
it on the own box of $X$, of side $D_X$, with two $\pi_j$-layers of margin. The irrelevance lemma
for a renormalised pointed group uses it on $\square^{\sim\nu+1}$ about its core
$\square^{\sim\nu}$, again with one $\pi_j$-layer of margin, as in \cite{B-CMP119}. In each of
these cases the representation holds: if $\mathfrak{b}$ is an admissible background at scale
$n \ge j$ and $\square_0$ is either $\square'^{\sim a+1}$ for a cube $\square' \in \pi_j$ with
$X \subset \square'^{\sim a}$, where $a = n-j$ or $a = \nu = 1+n-j$, or any box containing $X$
with two $\pi_j$-layers of margin, then, modulo gauge,
\begin{equation}\label{5.2:eq-own-box-a}
  \mathfrak{b} = U_j(\square_0,\mathsf{V}) \quad\text{on } X,
\end{equation}
where $\mathsf{V}$ denotes the level-$j$ data of $\mathfrak{b}$ on the determining set of
$\square_0$. The statement concerns configurations only; the complex source $w$ plays no role in
it.

The choice of the own box of $X$ in the marginal bound is what removes the logarithm from the
marginal weight. It is also what makes the condition $\kappa_0 \ge 7$ suffice in the proposition
that sums the marginal units over the scales $j$; here $\kappa_0$ is the exponent of the factor
$g_j^{\kappa_0}$ in Ba{\l}aban's bound for the marginal terms in \cite{B-CMP119}. Consider
instead a box of side comparable to
$M\nu$ centred at a point $y$, where $\nu = 1+n-j$. On such a box $\sup|B^{\mathrm{ax}}|$, the
supremum of the field of the comb gauge, carries the factor $\nu$ of \cite[(3.68)]{B-CMP119}.
The weight then becomes $\nu^2\rho_{n,j}^2$, and the sum over $j$ of the marginal units would
need $\kappa_0 > 10$ after extraction. In the irrelevance lemma for a renormalised pointed group
the factor $\nu$ is harmless, because its units carry the factor $t^{4+\hat\alpha}$, where
$t = L^{j-n}$.
\end{remark}

\begin{proof}
The argument has two parts.

First, the admissible backgrounds are, by their construction in
\cite[(3.3)--(3.4)]{B-CMP109} and \cite[(97)]{B-CMP98}, of the form $U_j$ evaluated
on level-$j$ data. Thus $\mathfrak{b}$ is a minimiser determined by its level-$j$ data.

Second, the restriction of such a minimiser to a box with margin is the minimiser on that box of
its own determining data, that is, $U_j(\square_0,\mathsf{V})$ with $\mathsf{V}$ as in
\eqref{5.2:eq-own-box-a}. We treat the two kinds of margin separately.

With one $\pi_j$-layer of margin and $a = n-j$ this is Ba{\l}aban's own representation in
\cite{B-CMP119}. It is taken there on the box $\square^{\sim n-j+1}$ about
$X \subset \square^{\sim n-j}$. With $a = \nu$, the configuration used by the irrelevance lemma
about its core, the representation is part of the statement of
Proposition~\ref{5.2:prop-box-comb-gauge} as cited; it may also be replaced by the two-layer case
on the box $\square^{\sim\nu+2}$, which contains the core $\square^{\sim\nu}$ with two
$\pi_j$-layers of margin and is covered below.

With two $\pi_j$-layers of margin it is the construction that Ba{\l}aban performs for an
arbitrary domain $X$ in clause (iv) of the definition of the spaces $U^c_j$ in \cite{B-CMP109}.
We quote that clause:
\begin{quotation}
(iv) We consider $X$ as $\Omega_0$, and we construct the sequence $\{\Omega_n\}$,
$n = 1,\dots,j$, of maximal possible domains satisfying the conditions (1.3)--(1.6) [14]
\dots{} we construct the functions $U_n(V)$ in the axial gauges, for regular $G^c$-valued
configurations $V$. We consider the pair
$(U_n(M^{\cdot}(\mathbf{U})),J_n(M^{\cdot}(\mathbf{U})))$, $n = 1,\dots,j$ \dots{} This pair
satisfies the bounds $|\partial U_n(M^{\cdot}(\mathbf{U}))-1| < \alpha_0\xi^2$,
$|J_n(M^{\cdot}(\mathbf{U}))| < \alpha_0(L^n\xi)^2$ on $X^{\sim -2}$, and the same bounds hold
for $U$ instead of $\mathbf{U}$. The domain $X^{\sim -2}$ is obtained from $X$ by taking away two
layers of cubes from $\pi_j$ which are closest to the boundary of $X$
\end{quotation}
Here [14] is \cite{B-CMP99a}, and the conditions are \cite[(1.3)--(1.6)]{B-CMP99a}.

We apply this construction with the box $\square_0$ in the role of the domain. It produces the
box minimiser of the data of the configuration on $\square_0$, together with the quoted bounds on
$\square_0^{\sim -2}$. The set $\square_0^{\sim -2}$ is obtained from $\square_0$ by removing the
two layers of $\pi_j$-cubes closest to its boundary. It therefore contains $X$ whenever
$\square_0$ contains $X$ with two $\pi_j$-layers of margin.

The quoted bounds come from \cite[Proposition~9]{B-CMP102b} through the argument of
\cite{B-CMP109} that follows the clause. The constant $B_3$ of \cite[Proposition~9]{B-CMP102b}
that enters them does not depend on the box. Clause (iv) is a defining condition of the spaces
$U^c_j$, and the spaces in use here satisfy it. Hence the representation
\eqref{5.2:eq-own-box-a}, with its bounds, holds on any box containing $X$ with two
$\pi_j$-layers of margin, and in particular on the own box of $X$ used in the marginal bound.
\end{proof}

\begin{remark}[The configuration statements along the preparatory and large-field chains]
\label{5.2:rem-chain-configurations}
Everything in this remark is read in the source-free scaffold of
Definition~\ref{5.1:def-source-free-scaffold}: all inserted identities, choices, domains and
backgrounds are those at $z=0$ and at the real couplings $g_k$. Here $N(g_k)$ is Ba{\l}aban's
number of preparatory stages, $Z$ is the large-field region at level $k$, $\Omega_k$ is the
domain of that name at level $k$ in \cite{B-CMP122b},
and $\mathfrak{T}:=\Omega_k\cup Z^c$. The configuration statements (c1)--(c3) are those of
Proposition~\ref{5.2:prop-box-comb-gauge}, displayed in \eqref{5.2:eq-box-comb-gauge-a}.
Assume the following.
\begin{enumerate}
\item[(a)] Ba{\l}aban's configuration statement in \cite{B-CMP122b}, taken by statement: the
configurations of \cite[(1.51)--(1.58)]{B-CMP122b} are box functions of level data, and the
values of these configurations lie
``in the space $\tilde U''^{c}_k(Y_1,\tilde\alpha_0,\tilde\alpha_1)$ \dots contained in the
analyticity domain of the term $E(X)$'', and hence in the analyticity domain of each of the
terms other than the chain terms $\mathbb{C}^{(n)}_k(X,\mathcal{S};w)$ of $\mathbf{B}''_k$ and
the arrested left-overs (the terms $\mathcal{C}_{\mathrm{arr}}$ left frozen when an attempt at
the operation $\mathbf{R}$ at an earlier step on a live piece of $Z$ was finished early,
Definition~\ref{5.1:def-source-free-scaffold}), namely ``all new terms created by the previous
localization operations \dots all the old terms''.
\item[(b)] The inherited-domain theorem, with its corollary on containment, the claim
accompanying that corollary, and the companion corollary at complex sources $w$. For the chain
terms $\mathbb{C}^{(n)}_k(X,\mathcal{S};w)$ of $\mathbf{B}''_k$, $n\le N(g_k)$, the theorem gives
the inherited domain $\mathfrak{D}^{\flat}_n(X)\times\mathfrak{D}_{\mathcal{S}}(X)$; the corollary
and the claim give the containment in it of the values of the configurations of hypothesis~(a).
\item[(c)] The theorem on arrested pieces. Here the arrested-piece re-definition is the
re-definition of the spaces put in force on every live arrested piece in
\eqref{5.1:eq-source-free-scaffold-a}, and the arrested-piece reading of a space is that space as
re-defined there. The theorem states: for an arrested left-over term and a point $x$, the
image, under the map of the operation at which the term is read, of the source of that
operation in its arrested-piece reading lies in the frozen domain of the term at $x$, with
positive room.
\item[(d)] The live-slot reading at the three sites of the renormalisation step --- the
localisation site (L), the site (T) of the operation $\mathbf{T}$ and the preparatory site (P) ---
and in particular the third clause of that statement: at these sites the live slots of a
unit are read at the presented rotation $\mathbf{A}^{[\mathfrak{Q}^{\mathfrak{b}}(\zeta)]}$ on the
orbit datum.
\item[(e)] The corollary in which the configurations that we denote $\Phi(b)$ are constructed,
taken by statement: each configuration $\Phi(b)$ constructed in that corollary is a function
$U_j(\square_0,\mathsf{V})$ of its own
level-$j$ data $\mathsf{V}$ on the determining set of a box $\square_0$ with margin.
\end{enumerate}
Then the following hold.
\begin{enumerate}
\item[(i)] The configuration statements (c1)--(c3) hold at the level of the scaffold data for the
backgrounds $U^{(n')}_k$, $n'\le N(g_k)$, of the preparatory integrals
\cite[(1.51)--(1.63)]{B-CMP122a}. They hold also for the configurations $\Phi(b)$ of
hypothesis~(e).
\item[(ii)] The configuration statements (c1)--(c3) hold at the level of the scaffold data for the
configurations of \cite[(1.51)--(1.58)]{B-CMP122b}: the background $U_k$, the new background
with its new determining set, and their interpolations. Their values lie in the analyticity
domain of every term of the expansion: for the terms named in hypothesis~(a) because the space
quoted there is contained in their analyticity domains, for the chain terms of $\mathbf{B}''_k$
in the inherited domain of
hypothesis~(b), and for the arrested left-overs in the frozen domain of hypothesis~(c). For these
configurations a box meeting $Z$ is admissible.
\end{enumerate}
\end{remark}

\begin{proof}
\emph{The preparatory chain.} The backgrounds $U^{(n')}_k$, $n'\le N(g_k)$, of
\cite[(1.51)--(1.63)]{B-CMP122a} are Ba{\l}aban's minimisers for the successively enlarged
determining sets. Each of them is therefore again a function $U_j(\square_0,\mathsf{V})$ of its own
level data $\mathsf{V}$ on the determining set of a box $\square_0$ with margin, as in
Proposition~\ref{5.2:prop-box-comb-gauge}. Hence, by clause~(iv) of
\cite[(1.11)--(1.16)]{B-CMP109}, as used in Remark~\ref{5.2:rem-own-box} and displayed in
\eqref{5.2:eq-own-box-a}, (c1)--(c3) hold for the $U^{(n')}_k$. The configurations
$\Phi(b)$ of hypothesis~(e) are of the same form, so the same clause (iv) applies to them and
gives (c1)--(c3) for them. All data entering
these minimisers are those of the scaffold of Definition~\ref{5.1:def-source-free-scaffold} at
$z=0$, so the statements hold at the level of the scaffold data.

\emph{The large-field chain.} By hypothesis~(a) the configurations of
\cite[(1.51)--(1.58)]{B-CMP122b} are box
functions of level data. This covers $U_k$, the new background with its new determining set
of \cite{B-CMP122b}, and their interpolations. It remains to place their values in the
analyticity domain of every term. There are three cases.

\emph{Case~1: the chain terms $\mathbb{C}^{(n)}_k(X,\mathcal{S};w)$ of $\mathbf{B}''_k$,
$n\le N(g_k)$.} These are among the units of the summation of boundary terms at finer zones
(Theorem~\ref{5.4:thm-finer-zone-sum}). A unit of index $k$ may
reach $X\cap\mathfrak{T}$ outside $Z$. For these units the analyticity domain is the inherited
domain $\mathfrak{D}^{\flat}_n(X)\times\mathfrak{D}_{\mathcal{S}}(X)$ of hypothesis~(b). The
containment of the chain's values in it is the corollary on containment together with its claim.

The corollary and the claim supply the following; there $r_a$ denotes the threshold of the entry
$a$ of the conditions that define the source $\tilde U^c_k(Y_4,\tilde\alpha_0,\tilde\alpha_1)$,
$r_{\min}(y)$ the least threshold at $y$, and $m(y)$ the level attached to the location $y$,
all as in the corollary on containment.
For the entries $a\notin\{5,8\}$ set
$\mathrm{room}^{\flat}(y):=\min\{\tfrac12 r_{\min}(y),\beta_{\mathrm{pr}}\}$. Over the
source $\tilde U^c_k(Y_4)$ at locations in $\mathfrak{T}$ this satisfies
\begin{equation*}
\mathrm{room}^{\flat}(y)\ \ge\ \tfrac12\,\beta_{\mathrm{pr}}\,\theta_S\,
2^{-(k-m(y))} .
\end{equation*}
On $\Omega_k$ it is at least $\tfrac12\beta_{\mathrm{pr}}$, under the top-layer clause of the lemma
on the real integration domain. The increments of the $\mathbf{R}$-layers spend at most one
quarter of this room, under the condition
\begin{equation*}
16\max\{P,P',P''\}\ \le\ \beta_{\mathrm{pr}}\,\theta_S
\end{equation*}
on the prefactors $P$, $P'$, $P''$ of these increments;
it is one of the conditions under which the threshold $\gamma_J$ is chosen. The entries $5$
and $8$ keep, at locations in $\mathfrak{T}$, the
thresholds of the source and are never incremented. The box clause holds at every triple of the
second of the two readings named in the clause on spaces of the
arrested-piece re-definition \eqref{5.1:eq-source-free-scaffold-a}. The
scale-drop clauses hold with
$1.1\,r_a\le 1.5\,r_a$. The claim's
clauses $(\alpha)$--$(\gamma)$ hold for the pair presented by $h$ at real slot data, where $h$
is the presenting element of the claim, acting on the current slot by
$\operatorname{Ad}_h$; in this remark the letter $h$ denotes only this element. The
companion corollary gives them at $w$. The live slots are then read at the presented rotation
$\mathbf{A}^{[\mathfrak{Q}^{\mathfrak{b}}(\zeta)]}$ on the orbit datum, by hypothesis~(d).

The map under which the source is carried into this domain is that of
\cite[(1.65)]{B-CMP122b}. It is written so that the carrier pair presented by
$h$ is a holomorphic function of $(\mathbf{U},\mathbf{J})$, by the holomorphy theorems for the
presented carrier pair. The live slots are presented rotated by $\mathfrak{Q}^{\mathfrak{b}}(\zeta)$.
The map is read through the live-slot reading of hypothesis~(d) at the sites (L), (T) and (P); at
the localisation site (L) it is read together with the supplementary condition attached to the
live-slot reading at that site.

At this point \cite{B-CMP122b} refers to an inductive inequality: ``It is clear for terms of
$\mathbf{B}''_k$, by the inductive assumptions \dots (see the inequality (1.70) [IV] \dots)'',
the bracketed [IV] being a reference of the bibliography of \cite{B-CMP122b}. For these units the
containment is taken here instead from hypothesis~(b); this includes the units
$\mathbb{C}^{(N(g_k))}_k(X)$ with $X\cap(\Omega_k\setminus Z)\neq\emptyset$.

\emph{Case~2: the arrested left-overs $\mathcal{C}_{\mathrm{arr}}$.} Let $\mathfrak{a}$ be a live
arrest: the attempt at the operation $\mathbf{R}$ at an earlier step $k_a<k$ on a live piece
$W_{\mathfrak{a}}\subset Z$ was finished early. Its left-overs are of two kinds. The
first are the frozen chain terms $\mathbb{C}^{(i)}_{k_a}(X,\mathcal{S})$, $i\le n^+_a$,
where $n^+_a$ is the bound on the stage index attached to the arrest $\mathfrak{a}$ in the
definition of a live arrest.
The second are the change-of-variables terms that the definition of a live arrest freezes together
with them; their frozen form is that of the freezing statement for change-of-variables terms.
Both are left ``as a part of all boundary
terms'' \cite{B-CMP122a} when the attempt is finished early in this way. For them Ba{\l}aban
writes ``it is easy to
generalize the definition of the spaces'' \cite{B-CMP122a}. Both kinds occur also among the units
of the summation of boundary terms at finer zones (Theorem~\ref{5.4:thm-finer-zone-sum}) and
among the values at the localisation site,
by the second step of the lemma on the values of the chain units, by clause (i$'$) of the
proposition on the strip widths and by clause ($\beta'$) of the per-zone-cell volume count.

For these terms the analyticity domain is the domain frozen at the arrest, taken in the
arrested-piece reading of Definition~\ref{5.1:def-source-free-scaffold}. The source is the
arrested source $\tilde U^{\mathrm{arrest}}_k$, that is
$\tilde U^c_k(Y_4,\tilde\alpha_0,\tilde\alpha_1)$ in its arrested-piece reading on the live piece,
since the five clauses
of the arrested-piece re-definition are in force by \eqref{5.1:eq-source-free-scaffold-a}. The
containment of the
chain's values in that domain is hypothesis~(c), the theorem on arrested pieces. Its proof treats
the live piece $W_{\mathfrak{a}}$ by setting the lemma on the room of arrested pieces against the
increments of the increment lemma. These increments arise at every later step of the operation
$\mathbf{R}$, in the first or the second part of the map \cite[(1.65)]{B-CMP122b} of that step.
Off $W_{\mathfrak{a}}$ it uses the
inherited-domain theorem with $(Z,k,n)$ replaced by $(W_{\mathfrak{a}},k_a,i)$, and it treats two
further cases. The scope clause of the corollary on containment sends the rooms of these terms to
the arrested-piece reading.

\emph{Case~3: every other term.} These are ``all new terms created by the previous localization
operations \dots all the old terms'' \cite{B-CMP122b}. For them the containment is
hypothesis~(a), taken by statement.

In all three cases the configurations are box functions of level data whose values lie in the
analyticity domain of the term. Hence (c1)--(c3) hold along the chain at the level of the scaffold
data. Finally, these configurations differ from $U_k$ only through changes ``made inside the
domain $\Lambda$'' \cite{B-CMP122b}, where $\Lambda$ is Ba{\l}aban's domain of the operation
$\mathbf{R}$ in \cite{B-CMP122b}. A box meeting $Z$ is therefore admissible. Ba{\l}aban's own
box $\tilde\square^{\sim 5}$ at \cite[(1.59)]{B-CMP122b} is of this kind.
\end{proof}

\begin{lemma}[A polydisc inside the analyticity space]\label{5.2:lem-polydisc-containment}
Let $X$ be a localisation domain at scale $j$, and let $\mathcal{U}^\flat_j(X)$ and the radii
$\alpha_{0,j},\alpha_{1,j}$ be as in Definition~\ref{5.1:def-analyticity-spaces}. There is a
constant $c_T>0$, independent of $j$, of $X$ and of the coupling, such that, with
$\mathfrak{r}_j$ defined in \eqref{5.2:eq-polydisc-containment-a} below,
\begin{equation*}
\bigl\{(e^{i\xi\mathcal{A}},\mathbf{J}) :
\max\bigl(|\mathcal{A}|,|\nabla^{\xi}\mathcal{A}|,|\mathbf{J}|\bigr)\le\mathfrak{r}_j
\text{ on } X\bigr\}\subset\mathcal{U}^\flat_j(X).
\end{equation*}
Here a pair on $X$ is written in the form $(e^{i\xi\mathcal{A}},\mathbf{J})$ of
\cite[(1.11)--(1.16)]{B-CMP109}: $\mathcal{A}$ is the bond field in the exponent, $\xi$ the
factor multiplying it there, and $\nabla^{\xi}\mathcal{A}$ the lattice derivative of
$\mathcal{A}$ that enters the conditions of \cite[(1.11)--(1.16)]{B-CMP109}; in this lemma the
letters $\mathcal{A}$ and $\xi$ have this meaning only. Moreover, with $\alpha_0,\alpha_1$ the
absolute radii of $\mathcal{U}_j(X)$ and with a constant $C_t$ depending only on
$\alpha_0,\alpha_1$ and on the constants $C_0,C_1$ of the radii $\alpha_{i,j}$ of
Definition~\ref{5.1:def-analyticity-spaces}, set
\begin{equation}\label{5.2:eq-polydisc-containment-a}
\begin{gathered}
\mathfrak{r}_j := c_T\min\{\alpha_{0,j},\alpha_{1,j}\},\qquad
\mathfrak{r}_{\circ} := c_T\min\{\alpha_0,\alpha_1\},\qquad
\mathfrak{t}_j := \mathfrak{r}_j/\mathfrak{r}_{\circ};\\
\text{then}\quad \mathfrak{t}_j^{-2}\le C_t\,x_j
\quad\text{for every } j \text{ with } \log x_j\ge 1 .
\end{gathered}
\end{equation}
\end{lemma}

\begin{proof}
By Definition~\ref{5.1:def-analyticity-spaces},
$\mathcal{U}^\flat_j(X) = U^c_j(X,\alpha_{0,j},\alpha_{1,j})$. This is Ba{\l}aban's space of
\cite[(1.11)--(1.16)]{B-CMP109} with the radii $(\alpha_{0,j},\alpha_{1,j})$. Membership is
given by the conditions (i)--(iv) there.

\emph{Conditions} (i)--(iii). Each of them bounds $\mathcal{A}$, $\nabla^{\xi}\mathcal{A}$ or
$\mathbf{J}$ on $X$ by a quantity that is homogeneous of degree one in the pair of radii. Hence
there is a constant $c_T>0$ with the following property. For every pair of radii
$(\alpha_0',\alpha_1')$, the bound
\begin{equation*}
\max\bigl(|\mathcal{A}|,|\nabla^{\xi}\mathcal{A}|,|\mathbf{J}|\bigr)
\le c_T\min\{\alpha_0',\alpha_1'\}\quad\text{on } X
\end{equation*}
implies (i)--(iii) at $(\alpha_0',\alpha_1')$. Homogeneity is exactly what makes $c_T$
independent of the radii, and hence of $j$, $X$ and the coupling. We apply this at
$(\alpha_0',\alpha_1') = (\alpha_{0,j},\alpha_{1,j})$. Every pair in the set on the left
therefore satisfies (i)--(iii) with the radii of $\mathcal{U}^\flat_j(X)$.

\emph{Condition} (iv). For small radii it is not an independent requirement. Ba{\l}aban
observes in \cite{B-CMP109}, following the conditions (i)--(iv), that for $\alpha_0'$
sufficiently small the estimates (i)--(iii) imply condition (iv). We apply this at the radii
$(\alpha_{0,j},\alpha_{1,j})$. Together with the previous step, this proves the containment.

\emph{The ratio.} By Definition~\ref{5.1:def-analyticity-spaces},
$\alpha_{i,j} = g_jC_i(\log x_j)^q$ for $i=0,1$. Hence
\begin{equation*}
\mathfrak{r}_j = c_T\,g_j\min\{C_0,C_1\}(\log x_j)^q .
\end{equation*}
The absolute radii do not depend on $g_j$. Dividing by $\mathfrak{r}_{\circ}$ and using
$g_j^{-2}=x_j$ gives
\begin{equation*}
\mathfrak{t}_j^{-2}
=\Bigl(\frac{\min\{\alpha_0,\alpha_1\}}{\min\{C_0,C_1\}}\Bigr)^{2}x_j\,(\log x_j)^{-2q}.
\end{equation*}
For $j$ with $\log x_j\ge 1$, as in the last clause of the statement, the factor
$(\log x_j)^{-2q}$ is at most one. This gives the last line of
\eqref{5.2:eq-polydisc-containment-a} with
$C_t := \bigl(\min\{\alpha_0,\alpha_1\}/\min\{C_0,C_1\}\bigr)^{2}$.
\end{proof}

\begin{lemma}[The extracted coefficient depends on the total only]\label{5.2:lem-coefficient-total}
Let $F$ and $F'$ be $E$-families at scale $j$ in the sense of
Definition~\ref{5.1:def-analyticity-spaces}. Each of them may be plain or pointed, their
localisations may differ, and the radii of the spaces on which their terms are analytic may
differ. If $F\simeq F'$, then
\begin{equation*}
\beta[F]=\beta[F'].
\end{equation*}
In particular, $\beta[F]$ depends only on the terms $F(X,\cdot)$ whose domains $X$ do not wrap
around the torus $T^{(j)}$, and it does not depend on the torus.
\end{lemma}

\begin{proof}
Recall the objects involved, as fixed in \eqref{5.1:eq-analyticity-spaces-c}. The coefficient
$\beta[F]$ is defined by \cite[(1.20)--(1.22)]{B-CMP109}. The relation $F\simeq F'$ means that
the totals of $F$ and $F'$ agree up to a constant on the real regular configurations of every
torus. Thus, for each torus, whatever its side, the difference of the totals takes one and the
same value on all real regular configurations of that torus:
\begin{equation}\label{5.2:eq-coefficient-total-1}
F(U)-F'(U)=F(V)-F'(V)
\end{equation}
for all real regular configurations $U$, $V$ on that torus; this value may depend on the torus.
Here $F(U)$ and $F'(U)$ are the totals
$\sum_X F(X,U,J(U))$ and $\sum_X F'(X,U,J(U))$.

We first treat plain families. The quantity in \cite[(1.20)]{B-CMP109} is a second derivative of
the total at $B=0$, where $B$ is the field variable in which \cite[(1.20)]{B-CMP109}
differentiates the total. This derivative is taken along real configurations in a neighbourhood
of $B=0$; these are regular configurations, so \eqref{5.2:eq-coefficient-total-1} holds on them.
By \eqref{5.2:eq-coefficient-total-1}, the difference of the two totals is constant on these
configurations, so its derivatives vanish. Hence, on every finite torus $T^{(j)}$, the
second-derivative tensor built from the total of $F$ coincides with the tensor built from the
total of $F'$.

Let $N$ denote here the number of sites per side of $T^{(j)}$ (in this proof only; $N$ is not the
rank of the gauge group). Some terms $F(X,\cdot)$ have domains
$X$ that wrap around the torus. Their contribution to these tensors is $O(e^{-cN})$, with a
constant $c>0$, and the same holds for $F'$.

The defining expression \cite[(1.22)]{B-CMP109} is formed after the limit $T^{(j)}\nearrow
\mathbb{Z}^4$. This limit exists by the localised representation \cite[(1.7)]{B-CMP109}, as is
stated in \cite[at (1.22)]{B-CMP109}. Since the finite-torus tensors of $F$ and $F'$
coincide for every torus, their limits coincide, and $\beta[F]=\beta[F']$.

Consider now a pointed family $\{F(X,\cdot,y)\}_{y\in X}$, with pointed sums
\begin{equation*}
F(U,y)=\sum_{X\ni y}F(X,U,y).
\end{equation*}
For such a family the number \cite[(3.61)]{B-CMP119} is formed from these sums. By
\cite{B-CMP119}, this number coincides with the number defined by \cite[(1.22)]{B-CMP109}, so
that $\beta[F]$ is unambiguous. The argument there uses
three ingredients: the translation invariance of the family, the identity
\cite[(4.15)]{B-CMP109}, and linear algebra over $\mathbb{C}$. Hence, whenever $F$ or $F'$ is
pointed, its coefficient is again the number of \cite[(1.22)]{B-CMP109} formed from its total. The
argument of the preceding paragraphs applies without change.

In that argument, $F$ and $F'$ enter only through their totals on real regular configurations of
each torus and through the existence of the limit. Nothing depends on the choice of localisation,
on whether the family is pointed, or on the radii of the spaces on which the terms are analytic.
This proves the first assertion.

For the second assertion, note that the contribution of the wrapping terms is $O(e^{-cN})$. This
quantity tends to zero as $T^{(j)}\nearrow\mathbb{Z}^4$, that is, as $N\to\infty$. Hence it does
not enter the limit in \cite[(1.22)]{B-CMP109}, and $\beta[F]$ is determined by the terms
$F(X,\cdot)$ whose domains do not wrap. By the requirement in
Definition~\ref{5.1:def-analyticity-spaces} that the non-wrapping terms of an $E$-family be
independent of the torus, these terms do not depend on the torus.
Therefore $\beta[F]$ is independent of the torus.
\end{proof}

\begin{remark}[A second route to the driven coefficient with a real Gaussian]
\label{5.2:rem-real-gaussian-route}
Let $\Phi_{k+1}=\Phi_{k+1}(\mathfrak{q},s;h)$ be the family of
Proposition~\ref{5.2:prop-complex-factor-step}, with $h\in\mathbb{B}_k$ fixed and suppressed from
the notation. The driven flow lemma (Lemma~\ref{5.3:lem-driven-flow}) reads the step only through
the scalar $f:=\beta[\Phi_{k+1}]$ and through the membership in $\mathbb{B}_k$ of the family
$\mathsf{P}_\xi$, which carries the memory of the earlier steps and enters the step as the
history $h$. The bound
on this scalar that the running-shift statement in \eqref{5.1:eq-correlation-theorem-b} consumes
can be obtained a second time, with the Gaussian measure kept real. At $\mathfrak{q}=g_k$,
\begin{equation}\label{5.2:eq-real-gaussian-route-1}
  |f|\le C_\beta\bigl(2E^*\tau_k+E^*e^{-\gamma_2T_k^2/24}\bigr)
\end{equation}
holds on the polydisc of radii $\tfrac12\sigma_0x_j$, $0\le j\le k$, in the slot variables
$\xi_j$. On this
polydisc each slot combination $1-\xi_jg_j^2$ lies in $D(1,\sigma_0/2)$, because
$|\xi_jg_j^2|<\tfrac12\sigma_0x_jg_j^2=\tfrac12\sigma_0$. Consequently the correlation theorem
(Theorem~\ref{5.1:thm-correlation-theorem})
follows with $\sigma_0$ replaced by $\sigma_0/2$, and with $\vartheta_k$ replaced by
$E^*e^{-\gamma_2T_k^2/24}$.
\end{remark}

\begin{proof}
Throughout, $T=T_k$ and $(\mathfrak{q},s)\in D(0,\mathsf{r}_k)\times D(1,\sigma_0)$, as in
Proposition~\ref{5.2:prop-complex-factor-step}.

\emph{The real measure and the extra form.} We do not pass to the complex Gaussian measure
$d\mu_{C^{(k)}/s}$. We keep the real measure $d\mu_{C^{(k)}}$ and carry the quadratic form
\begin{equation*}
  \tfrac12(1-s)\langle B,\Delta^{(k)}B\rangle
\end{equation*}
in the potential. This form is localised by \cite[(1.10)]{B-CMP116} on the random walks of
$\Delta^{(k)}$. Let $C_\Delta$ denote the constant of the resulting localised expansion. Since
$|1-s|<\sigma_0$, the localised form is one more form of the type \cite[(1.42)]{B-CMP116}, of size
$\sigma_0C_\Delta$. We require $\sigma_0C_\Delta$ to be at most the constant of
\cite[(1.43)]{B-CMP116}; this is a condition on $\sigma_0$.

\emph{The family $\Phi^{\mathrm{out}}$.} Under this condition \cite[Lemma~3]{B-CMP116} applies as
stated to the augmented potential, namely the potential of
Proposition~\ref{5.2:prop-complex-factor-step} together with the above form; this is the
statement of this chapter which carries \cite[Lemma~3]{B-CMP116} over to potentials depending
holomorphically on parameters and which gives the analytic dependence of the resulting family on
those parameters. It
yields a family, which we denote $\Phi^{\mathrm{out}}=\Phi^{\mathrm{out}}(\mathfrak{q},s)$,
with the following properties:
\begin{equation*}
  \Phi^{\mathrm{out}}\simeq\Phi_{k+1}
  \quad\text{(in the sense of \eqref{5.1:eq-analyticity-spaces-c})},
  \qquad
  \|\Phi^{\mathrm{out}}\|\le E^*.
\end{equation*}
Moreover $\Phi^{\mathrm{out}}$ is analytic in $(\mathfrak{q},s)$.

By Lemma~\ref{5.2:lem-coefficient-total}, $\beta[\Phi^{\mathrm{out}}]=\beta[\Phi_{k+1}]$.
The functional $\beta$ is linear and satisfies $|\beta[F]|\le C_\beta\|F\|$ by
Lemma~\ref{5.2:lem-extraction-cost}. Hence the scalar
$(\mathfrak{q},s)\mapsto\beta[\Phi^{\mathrm{out}}(\mathfrak{q},s)]$ is analytic, and
\begin{equation}\label{5.2:eq-real-gaussian-route-2}
  |\beta[\Phi^{\mathrm{out}}(\mathfrak{q},s)]|\le C_\beta E^*
  \qquad\text{on } D(0,\mathsf{r}_k)\times D(1,\sigma_0).
\end{equation}

\emph{The value at $\mathfrak{q}=0$ on the real diameter.} At $\mathfrak{q}=0$ the potentials of
Proposition~\ref{5.2:prop-complex-factor-step} vanish, and only the quadratic form above remains.
Take $s$ real, $s\in\,]1-\sigma_0,1+\sigma_0[$. The substitution $B\mapsto s^{-1/2}B$ gives
\begin{equation*}
  \text{total of }\Phi^{\mathrm{out}}(0,s)
  =\mathrm{const}+\log\int d\mu_{C^{(k)}}\prod\chi\bigl(|B|<T\sqrt{s}\,\bigr).
\end{equation*}
By the statement supplying the family of a cut-off real Gaussian integral together with its norm
bound, applied at the threshold $T\sqrt{s}$ with $s>1-\sigma_0$, the family of this integral has
norm at most $E^*e^{-\gamma_2(1-\sigma_0)T^2/4}$.

By Lemma~\ref{5.2:lem-coefficient-total}, $\beta[\Phi^{\mathrm{out}}(0,s)]$ equals the coefficient
of that family. Lemma~\ref{5.2:lem-extraction-cost} therefore gives
\begin{equation}\label{5.2:eq-real-gaussian-route-3}
  |\beta[\Phi^{\mathrm{out}}(0,s)]|\le C_\beta E^*e^{-\frac14\gamma_2(1-\sigma_0)T^2},
  \qquad s\in\,]1-\sigma_0,1+\sigma_0[.
\end{equation}
On the whole disc $D(1,\sigma_0)$ the bound \eqref{5.2:eq-real-gaussian-route-2} at $\mathfrak{q}=0$
gives $|\beta[\Phi^{\mathrm{out}}(0,s)]|\le C_\beta E^*$.

\emph{The half disc.} Apply Lemma~\ref{5.2:lem-two-constants} to
$\varphi(s):=\beta[\Phi^{\mathrm{out}}(0,s)]$ on $D(1,\sigma_0)$, with $M:=C_\beta E^*$ and $m$ the
right side of \eqref{5.2:eq-real-gaussian-route-3}, so that $m\le M$. We require in addition
$\sigma_0\le\frac12$, a second condition on $\sigma_0$; then $(1-\sigma_0)/12\ge1/24$. For
$s\in D(1,\sigma_0/2)$ it follows that
\begin{equation}\label{5.2:eq-real-gaussian-route-4}
\begin{aligned}
  |\beta[\Phi^{\mathrm{out}}(0,s)]|
  &\le m^{1/3}M^{2/3}
   =C_\beta E^*e^{-\frac1{12}\gamma_2(1-\sigma_0)T^2}\\
  &\le C_\beta E^*e^{-\gamma_2T^2/24}.
\end{aligned}
\end{equation}

\emph{The coupling slot.} Fix $s\in D(1,\sigma_0/2)$. The function
$\mathfrak{q}\mapsto\beta[\Phi^{\mathrm{out}}(\mathfrak{q},s)]$ is analytic on $D(0,\mathsf{r}_k)$ and
bounded there by $C_\beta E^*$, by \eqref{5.2:eq-real-gaussian-route-2}. Part~(a) of
Lemma~\ref{5.2:lem-disc-bounds}, \eqref{5.2:eq-disc-bounds-a}, with centre $0$, radius
$\mathsf{r}_k$ and $z=g_k$, gives
\begin{equation*}
  \bigl|\beta[\Phi^{\mathrm{out}}(g_k,s)]-\beta[\Phi^{\mathrm{out}}(0,s)]\bigr|
  \le\frac{2C_\beta E^*g_k}{\mathsf{r}_k}=2C_\beta E^*\tau_k,
\end{equation*}
since $\tau_k=g_k/\mathsf{r}_k$. We have $f=\beta[\Phi_{k+1}(g_k,s)]=\beta[\Phi^{\mathrm{out}}(g_k,s)]$.
Combining the last display with \eqref{5.2:eq-real-gaussian-route-4} at $T=T_k$ therefore gives
the bound \eqref{5.2:eq-real-gaussian-route-1} for every $h\in\mathbb{B}_k$ and every
$s\in D(1,\sigma_0/2)$. On the polydisc of radii $\frac12\sigma_0x_j$, $0\le j\le k$, the last slot
gives $s=1-\xi_kg_k^2\in D(1,\sigma_0/2)$, as noted in Remark~\ref{5.2:rem-real-gaussian-route},
and the earlier slots enter the
step only through the membership $\mathsf{P}_\xi\in\mathbb{B}_k$; hence
\eqref{5.2:eq-real-gaussian-route-1} holds there.

\emph{Conclusion.} The bound \eqref{5.2:eq-real-gaussian-route-1} has the form
$C_\beta(2E^*\tau_k+\vartheta_k)$ of Proposition~\ref{5.2:prop-complex-factor-step}(c), with
$\vartheta_k$ replaced by $E^*e^{-\gamma_2T_k^2/24}$. This bound, on the
polydisc of radii $\frac12\sigma_0x_j$, together with the membership $\mathsf{P}_\xi\in\mathbb{B}_k$,
is all that the driven flow lemma reads. The correlation theorem
(Theorem~\ref{5.1:thm-correlation-theorem}) therefore follows with $\sigma_0$
replaced by $\sigma_0/2$ and with $\vartheta_k$ so replaced.

It suffices that the argument sees the step only through its total. The reason is that the memory
of earlier steps is carried in the history slot $h=\mathsf{P}_\xi$ of the potential, and not by
re-localising differences of families.
\end{proof}

\begin{lemma}[Extraction costs the inverse square of the radius]
\label{5.2:lem-extraction-cost}
Let $F=\{F(X,\cdot)\}_{X}$ be an $E$-family at scale $j$ in the sense of
Definition~\ref{5.1:def-analyticity-spaces}, and let $\beta[F]$ be its marginal coefficient,
fixed in \eqref{5.1:eq-analyticity-spaces-c} by \cite[(1.20)--(1.22)]{B-CMP109}. Let
$\mathfrak{t}_j=\mathfrak{r}_j/\mathfrak{r}_\circ$ be the ratio of radii of
\eqref{5.2:eq-polydisc-containment-a} (Lemma~\ref{5.2:lem-polydisc-containment}).
There is a constant $C_\beta$, given by \eqref{5.2:eq-extraction-cost-a} below in terms of
the constant $B_3$ of \cite[Prop.~9]{B-CMP102b}, the constants $C$, $\delta_1$ of
\cite[(5.10)]{B-CMP109} and the absolute radius
$\mathfrak{r}_\circ$ of Lemma~\ref{5.2:lem-polydisc-containment}, such that the following hold.
\begin{enumerate}
\item[(1)] If $F$ is a family on the spaces $\mathcal{U}_j(X)$, then
$|\beta[F]|\le C_\beta\|F\|_j$.
\item[(2)] If $F$ is a family on the spaces $\mathcal{U}^\flat_j(X)$, then
\begin{equation}\label{5.2:eq-extraction-cost-1}
|\beta[F]|\le C_\beta\,\mathfrak{t}_j^{-2}\,\|F\|^\flat_j .
\end{equation}
\item[(3)] The map $F\mapsto\beta[F]$ is linear.
\end{enumerate}
\end{lemma}

\begin{proof}
Clause~(1) is the first clause of the extraction bound for families at absolute radii, which
is used here by statement, with the constant $C_\beta$
of \eqref{5.2:eq-extraction-cost-a}; its proof is not repeated here. Clause~(3) holds because
\cite[(1.20)--(1.22)]{B-CMP109} obtains $\beta[F]$ from second derivatives of $F$ at the
identity configuration, and these depend linearly on $F$. It remains to prove~(2). Let
$F$ be a family on the spaces $\mathcal{U}^\flat_j(X)$, and write
$\mathfrak{n}:=\|F\|^\flat_j=\sup_X e^{\kappa d_j(X)}|F(X,\cdot)|$.

\emph{Reduction to a second derivative of $F$.}
For a Lie-algebra-valued field $B$, \cite[(4.2)]{B-CMP109} gives
\begin{equation*}
F\bigl(X,U_j(e^{iB})\bigr)=F\bigl(X,e^{i\xi H_j(B)}\bigr),
\end{equation*}
where $U_j(\cdot)$ and $H_j$ are the maps of that display and $\xi$ is the parameter of that
display, the same one that enters the set of the polydisc lemma. We differentiate twice at
$B=0$. By the chain rule, the second derivative of the
composite is $D^2F$ applied to the first derivatives of $H_j$, plus terms carrying the first
derivative $DF(X;1,0)$ of $F(X;\cdot,\cdot)$ at $(\mathbf{U},\mathbf{J})=(1,0)$. By
\cite[(4.14)]{B-CMP109} we have $DF(X;1,0)=0$, so those terms vanish. This is
\cite[(4.35)]{B-CMP109}:
\begin{equation}\label{5.2:eq-extraction-cost-2}
\frac{\partial^2}{\partial B(x)\,\partial B(x')}F\bigl(X,U_j(e^{iB})\bigr)\Big|_{B=0}
=D^2F(X;1,0)\bigl[\mathsf{h}_x,\mathsf{h}_{x'}\bigr].
\end{equation}
Here $\mathsf{h}_x:=\delta H_j/\delta B(x)$, taken together with its $\mathbf{J}$-entry.

\emph{Decay of the kernel.}
Configurations on $X$ are written $(e^{i\xi\mathcal{A}},\mathbf{J})$, with $\mathcal{A}$ a
Lie-algebra-valued bond field on $X$ and $\mathbf{J}$ a current on $X$ (in this sense the
letter $\mathcal{A}$ is used only here and in Lemma~\ref{5.2:lem-polydisc-containment}). For
such a pair put
\begin{equation*}
\|(\mathcal{A},\mathbf{J})\|:=\max_X\max\bigl(|\mathcal{A}|,|\nabla^\xi\mathcal{A}|,
|\mathbf{J}|\bigr),
\end{equation*}
the norm in which Lemma~\ref{5.2:lem-polydisc-containment} describes its set; we write
$\|\cdot\|$ for it. This is Ba{\l}aban's
regularity norm with its $\xi$-weights included, not a bondwise supremum. The set in the
polydisc lemma is its closed ball of radius $\mathfrak{r}_j$. By
\cite[Prop.~9, (190)]{B-CMP102b}, which \cite{B-CMP109} invokes in the discussion following
\cite[(4.35)]{B-CMP109}, together with the
statement that the same bounds hold for the derivatives and for the second-order operators,
there are constants $B_3$ and $\delta_0>0$ such that
\begin{equation}\label{5.2:eq-extraction-cost-3}
\|\mathsf{h}_x\|\le B_3\,e^{-\delta_0\operatorname{dist}(x,X)} ,
\end{equation}
that is, the maximum over $X$ of $|\mathsf{h}_x|$, $|\nabla^\xi\mathsf{h}_x|$ and the
modulus of the $\mathbf{J}$-entry of $\mathsf{h}_x$ is at most
$B_3e^{-\delta_0\operatorname{dist}(x,X)}$.

\emph{Cauchy estimate.}
By the polydisc lemma (Lemma~\ref{5.2:lem-polydisc-containment}), the ball of radius
$\mathfrak{r}_j$ in $\|\cdot\|$ is contained in $\mathcal{U}^\flat_j(X)$. On it $F(X;\cdot)$
is therefore analytic and bounded by $\mathfrak{n}e^{-\kappa d_j(X)}$.

Fix nonzero fields $\mathsf{h},\mathsf{h}'$ and consider the function of the complex
parameters $u,u'$ with $|u|,|u'|\le\mathfrak{r}_j/2$ obtained by evaluating $F(X;\cdot)$ at
the point with coordinate
$u\,\mathsf{h}/\|\mathsf{h}\|+u'\,\mathsf{h}'/\|\mathsf{h}'\|$. That coordinate has norm at
most $|u|+|u'|\le\mathfrak{r}_j$, so the function is analytic in the closed bidisc and
bounded there by $\mathfrak{n}e^{-\kappa d_j(X)}$. Its mixed derivative at the origin is
$D^2F(X;1,0)[\mathsf{h},\mathsf{h}']/(\|\mathsf{h}\|\,\|\mathsf{h}'\|)$. The Cauchy
estimate in two variables, with both radii equal to $\mathfrak{r}_j/2$, therefore gives
\begin{equation}\label{5.2:eq-extraction-cost-4}
\bigl|D^2F(X;1,0)[\mathsf{h},\mathsf{h}']\bigr|
\le 4\,\mathfrak{n}\,e^{-\kappa d_j(X)}\,\|\mathsf{h}\|\,\|\mathsf{h}'\|\,\mathfrak{r}_j^{-2}.
\end{equation}

\emph{Summation over localisation domains and sites.}
Insert \eqref{5.2:eq-extraction-cost-3} into \eqref{5.2:eq-extraction-cost-4} with
$\mathsf{h}=\mathsf{h}_x$ and $\mathsf{h}'=\mathsf{h}_{x'}$, and then use
\eqref{5.2:eq-extraction-cost-2}. Summing over the localisation domains $X$ gives
\begin{equation*}
\begin{split}
&\sum_X\biggl|\frac{\partial^2 F(X,U_j(e^{iB}))}{\partial B(x)\,\partial B(x')}\Big|_{B=0}\biggr|\\
&\qquad\le \frac{4\,\mathfrak{n}\,B_3^2}{\mathfrak{r}_j^{2}}
\sum_X e^{-\kappa d_j(X)-\delta_0(\operatorname{dist}(x,X)+\operatorname{dist}(x',X))}.
\end{split}
\end{equation*}
As in \cite[(5.10)]{B-CMP109}, there are constants $C$ and $\delta_1>0$ such that the last
sum over $X$ is at most $Ce^{-\delta_1|x-x'|}$.

The coefficient $\beta[F]$ of \cite[(1.20)--(1.22)]{B-CMP109} is bounded by the second
moment, in $x'-x$, of the second derivatives of the total at the identity, for $x$ fixed.
Write $S(\delta_1):=\sum_{x\in\mathbb{Z}^4}|x|^2e^{-\delta_1|x|}$. The sum over the sites
$x'$ of the torus in the next display is majorised by the corresponding sum over
$\mathbb{Z}^4$, so that
\begin{equation*}
\begin{split}
|\beta[F]|
&\le\sum_{x'}|x-x'|^2\sum_X
\biggl|\frac{\partial^2 F(X,U_j(e^{iB}))}{\partial B(x)\,\partial B(x')}\Big|_{B=0}\biggr|\\
&\le 4\,\mathfrak{n}\,B_3^2\,C\,\mathfrak{r}_j^{-2}\sum_{x'}|x-x'|^2e^{-\delta_1|x-x'|}
\le4B_3^2\,C\,S(\delta_1)\,\mathfrak{r}_j^{-2}\,\mathfrak{n}.
\end{split}
\end{equation*}
Now set
\begin{equation}\label{5.2:eq-extraction-cost-a}
C_\beta:=4B_3^2\,C\,S(\delta_1)\,\mathfrak{r}_\circ^{-2}.
\end{equation}
Since $\mathfrak{t}_j=\mathfrak{r}_j/\mathfrak{r}_\circ$, we have
$\mathfrak{r}_j^{-2}=\mathfrak{r}_\circ^{-2}\mathfrak{t}_j^{-2}$. The last display therefore
reads $|\beta[F]|\le C_\beta\mathfrak{t}_j^{-2}\mathfrak{n}$, which is
\eqref{5.2:eq-extraction-cost-1}.
\end{proof}

\begin{lemma}[The marginal bound on a domain's own box]
\label{5.2:lem-marginal-bound}
Let $F(X,\cdot)$ have the properties (f1), (f2) of the families on the analyticity spaces fixed
with \eqref{5.1:eq-analyticity-spaces-b}, on $U^c_j(X,a_j)$, with $|F|\le\mathfrak{n}$
there. Let $\mathfrak{b}$ be an admissible background at scale $n\ge j$, let $t=L^{j-n}$, and put
\begin{equation}\label{5.2:eq-marginal-bound-1}
D_X:=\operatorname{diam}_j X+5M\le 4M\bigl(d_j(X)+5\bigr),
\end{equation}
where $\operatorname{diam}_j X$ is the diameter of $X$ measured on the level-$j$ lattice.
\begin{enumerate}
\item[(i)] If $X\subset\tilde\square^2$, then
\begin{equation}\label{5.2:eq-marginal-bound-a}
\bigl|F(X,\mathfrak{b})-F(X,1)\bigr|\le 2\mathfrak{n}\bigl(C_M D_X t^2\rho_{n,j}\bigr)^2,
\qquad C_M:=\frac{2dC_{\mathfrak{a}}}{c_3}.
\end{equation}
\item[(ii)] If $X\not\subset\tilde\square^2$, then $d_j(X)\ge t^{-1}$.
\item[(iii)] If $X\subset\tilde\square^2$ and $w\mapsto\mathfrak{b}_w$ is a holomorphic family of
admissible backgrounds on $|w|<\mathsf{w}$, then $|F(X,\mathfrak{b}_w)-F(X,\mathfrak{b}_0)|$ is at
most $2|w|/\mathsf{w}$ times the right side of \eqref{5.2:eq-marginal-bound-a}.
\end{enumerate}
Here $\rho_{n,j}=\alpha_\cdot(g_n)/\alpha_\cdot(g_j)$ is the radius ratio of
\eqref{5.1:eq-coupling-functions-a}, and the radius $a_j$ is the level-$j$ radius
$\alpha_\cdot(g_j)$ of the same family.
\end{lemma}

\begin{proof}
(i) Let $\square_0$ be a box of side $D_X$ containing $X$ with two $\pi_j$-layers of margin, and
fix the centre $y\in X$ of the comb gauge. This is the box representation on the domain's own box
of Remark~\ref{5.2:rem-own-box}, that is, the case of item (c1) of
\eqref{5.2:eq-box-comb-gauge-a} in Proposition~\ref{5.2:prop-box-comb-gauge} in which $\square_0$
is any box containing $X$ with two $\pi_j$-layers of margin. On $\square_0$ we use the comb gauge,
item (c3) of \eqref{5.2:eq-box-comb-gauge-a}, centred at $y$. By (c1), modulo gauge
$\mathfrak{b}$ is $U_j(\square_0,\mathsf{V})$, where $\mathsf{V}$ is its level-$j$ data on the
determining set of $\square_0$. In the comb gauge, $\mathsf{V}=e^{iB^{\mathrm{ax}}}$. Together with
property (f2) of $F$, this gives $F(X,\mathfrak{b})=\varphi(1)$, where
\begin{equation}\label{5.2:eq-marginal-bound-2}
\varphi(\sigma):=F\Bigl(X,\,U_j\bigl(\square_0,e^{i\sigma B^{\mathrm{ax}}}\bigr),\,
J_j\bigl(\square_0,e^{i\sigma B^{\mathrm{ax}}}\bigr)\Bigr).
\end{equation}
Put
\begin{equation}\label{5.2:eq-marginal-bound-3}
S:=\frac{\alpha_{3,j}}{\sup|B^{\mathrm{ax}}|},
\end{equation}
the supremum being over the bonds of $\square_0$. For $|\sigma|<S$ the complex field
$B'=\sigma B^{\mathrm{ax}}$ satisfies $|B'|<\alpha_{3,j}$. Hence, by the trivial-point
containment, item (c2) of \eqref{5.2:eq-box-comb-gauge-a}, the pair in
\eqref{5.2:eq-marginal-bound-2}, restricted to $X$, lies in
$U^c_j(X,a_j)$, analytically in $\sigma$. Therefore $\varphi$ is analytic on the disc
$|\sigma|<S$ and satisfies $|\varphi|\le\mathfrak{n}$ there.

At $\sigma=0$ the data are trivial, $\mathsf{V}=1$, so $\varphi(0)=F(X,1)$.

Next, $\varphi'(0)=0$. Put
\begin{equation*}
E(\mathsf{V}):=F\bigl(X,U_j(\square_0,\mathsf{V}),J_j(\square_0,\mathsf{V})\bigr).
\end{equation*}
This is a gauge-invariant analytic function of small unit-lattice configurations, that is, a
function in the setting of \cite[(4.7)]{B-CMP109}. Moreover
$\varphi(\sigma)=E(e^{i\sigma B^{\mathrm{ax}}})$. By \cite[(4.14)]{B-CMP109}, applied to $E$,
the derivative of $\varphi$ at $\sigma=0$ vanishes.

We now bound $|\varphi(1)-\varphi(0)|$. Suppose first $S>1$. Lemma~\ref{5.2:lem-disc-bounds},
part (b) of \eqref{5.2:eq-disc-bounds-a}, applies to $\varphi$ on the disc $D(0,S)$ with bound
$\mathfrak{n}$ and $\varphi'(0)=0$; at $z=1$ it gives
\begin{equation}\label{5.2:eq-marginal-bound-4}
\bigl|F(X,\mathfrak{b})-F(X,1)\bigr|=|\varphi(1)-\varphi(0)|\le 2\mathfrak{n}S^{-2}.
\end{equation}
Suppose instead $S\le1$. Then \eqref{5.2:eq-marginal-bound-4} is trivial, because
$|F(X,\mathfrak{b})|+|F(X,1)|\le 2\mathfrak{n}\le 2\mathfrak{n}S^{-2}$.

It remains to bound $S^{-1}$. By (c2), $\alpha_{3,j}=c_3a_j$. By (c3), a bond $b$ of $\square_0$
satisfies
\begin{equation*}
|B^{\mathrm{ax}}(b)|\le 2d\,(1+|b-y|)\,\mathfrak{a}_nt^2.
\end{equation*}
On $\square_0$, whose side is $D_X$ and whose comb gauge is centred at $y\in X$, the factor
$1+|b-y|$ is bounded by $D_X$. Hence
\begin{equation}\label{5.2:eq-marginal-bound-5}
S^{-1}\le\frac{2dD_X\mathfrak{a}_nt^2}{c_3a_j}
=\frac{2dC_{\mathfrak{a}}}{c_3}\,D_Xt^2\,\frac{\alpha_\cdot(g_n)}{a_j}
=C_MD_Xt^2\rho_{n,j}.
\end{equation}
The middle equality uses $\mathfrak{a}_n=C_{\mathfrak{a}}\alpha_\cdot(g_n)$ from (c1), where
$\alpha_\cdot$ is the radius function through which $\rho_{n,j}$ is defined in
\eqref{5.1:eq-coupling-functions-a}. The last equality uses $a_j=\alpha_\cdot(g_j)$, as fixed in
the statement.
Inserting \eqref{5.2:eq-marginal-bound-5} into
\eqref{5.2:eq-marginal-bound-4} gives \eqref{5.2:eq-marginal-bound-a}.

(ii) This is a geometric statement about a domain $X$ and the box $\tilde\square^2$. It is the
geometry set up in \cite{B-CMP119}: a domain not contained in $\tilde\square^2$ has
$d_j(X)\ge t^{-1}$.

(iii) Let $B_{\mathrm{(i)}}$ denote the right side of \eqref{5.2:eq-marginal-bound-a}. Each
$\mathfrak{b}_w$ with $|w|<\mathsf{w}$ is an admissible background at scale $n$, so part (i)
applies to it and gives $|f(w)|\le B_{\mathrm{(i)}}$, where
\begin{equation*}
f(w):=F(X,\mathfrak{b}_w)-F(X,1).
\end{equation*}
The function $f$ is analytic on $|w|<\mathsf{w}$, since the family is holomorphic and $F(X,\cdot)$
has (f1), (f2). Lemma~\ref{5.2:lem-disc-bounds}, part (a) of \eqref{5.2:eq-disc-bounds-a}, on
the disc $D(0,\mathsf{w})$ with bound $B_{\mathrm{(i)}}$ gives
\begin{equation*}
\bigl|F(X,\mathfrak{b}_w)-F(X,\mathfrak{b}_0)\bigr|=|f(w)-f(0)|
\le\frac{2|w|}{\mathsf{w}}\,B_{\mathrm{(i)}}.
\qedhere
\end{equation*}
\end{proof}

\begin{remark}
The argument uses no covariance and no domain other than $X$, and it produces no logarithm,
because the box $\square_0$ is the box of $X$ itself. The only cross-scale number in
\eqref{5.2:eq-marginal-bound-a} is $\rho_{n,j}^2$. By \eqref{5.1:eq-coupling-functions-a} it
satisfies $\rho^2_{n,j}\le(1+o(1))x_j/x_n$. In \cite{B-CMP119} the corresponding factor is written
$O(1)(L^jL^{-n})^4g_j^{\kappa_0}$, with $\kappa_0$ the exponent of the coupling $g_j$ fixed in
that paper; here it is expressed through $\rho_{n,j}$. Suppose
the background lives on the enlarged radius $\hat\alpha_k$ of \eqref{5.1:eq-analyticity-spaces-b},
as after the frozen step of the large-field operation $\mathbf{T}$ at level $k+1$, and suppose
$a_j=\alpha_j$. Then the
number is $(\hat\alpha_k/\alpha_j)^2=w_{k,j}$, with $w_{k,j}$ as in
\eqref{5.1:eq-coupling-functions-a}.
\end{remark}

\begin{lemma}[Irrelevance of a renormalised pointed group]\label{5.2:lem-pointed-irrelevance}
Let $F$ be a pointed $E$-family at scale $j$ on radii $a_j$ with $\|F\|\le\mathfrak{n}$, let
$\mathfrak{b}$ be a background admissible at scale $n\ge j$ near the point $y$, and put
$\nu:=1+n-j$ and $t:=L^{j-n}$. Then the quantity $\mathfrak{i}_y(\mathfrak{b})$ defined in the first
line of the following display obeys the bound in its second line:
\begin{equation}\label{5.2:eq-pointed-irrelevance-a}
\begin{aligned}
\mathfrak{i}_y(\mathfrak{b})&:=\sum_{X\ni y}\bigl[F(X,\mathfrak{b},y)-F(X,1,y)\bigr]
-\beta[F]\,A(h_y,\mathfrak{b}),\\
|\mathfrak{i}_y(\mathfrak{b})|&\le C_I\,\mathfrak{n}\,\nu^3\rho_{n,j}^2(1+\rho_{n,j})\,t^{4+\hat\alpha},
\end{aligned}
\end{equation}
where $\rho_{n,j}$ is the radius ratio of \eqref{5.1:eq-coupling-functions-a} and $\hat\alpha$ is the
exponent of the remainder of the comb gauge in \eqref{5.2:eq-box-comb-gauge-a}.
\end{lemma}

\begin{proof}
We follow the second proof of the corresponding estimate in \cite[Sect.~3]{B-CMP119}, carried out here at
general radii $a_j$ and for complex-valued kernels. We write $\rho:=\rho_{n,j}$. Since
$\|F\|\le\mathfrak{n}$, each $F(X,\cdot,y)$ is analytic on its domain and bounded there by
$\mathfrak{n}e^{-\kappa d_j(X)}$, $\kappa$ being the decay rate of the family; we put
$C(\kappa):=\sum_{X\ni y}e^{-\kappa d_j(X)}$. Let $\square\in\pi_j$ be the cube containing $y$, let
$\square_0:=\square^{\sim\nu+1}$, and let $D:=M(2\nu+3)$ be the side of $\square_0$. In the tensor
computations of the core below, the Greek indices $\kappa,\lambda,\mu,\nu'$ run over the four lattice
directions; the indices $\kappa$ and $\lambda$ there are not the decay rate $\kappa$ nor the flow
constant $\lambda$ of \eqref{5.1:eq-coupling-functions-a}.

\emph{Near levels.} Suppose that $C_MDt^2\rho>\tfrac12$, where $C_M$ is the constant of the marginal
bound on a domain's own box (Lemma~\ref{5.2:lem-marginal-bound}). Since
\begin{equation*}
\rho^2\le 1+(3\lambda\nu+2c_2)g^2+o(1),
\end{equation*}
with the flow constants $\lambda$ and $c_2$, this can happen only for $\nu\le\nu_*(L,M)$, a number
independent of $K$. In that case we use the trivial bound: every summand in
\eqref{5.2:eq-pointed-irrelevance-a} satisfies
$|F(X,\mathfrak{b},y)-F(X,1,y)|\le2\mathfrak{n}e^{-\kappa d_j(X)}$, so that
\begin{equation*}
|\mathfrak{i}_y(\mathfrak{b})|\le 2C(\kappa)\mathfrak{n}+|\beta[F]|\,|A(h_y,\mathfrak{b})|
\le C\mathfrak{n}(1+\rho^2).
\end{equation*}
Since $\rho\ge\tfrac12$ and $\nu\le\nu_*$, this lies inside the claimed bound as soon as
$C_I\ge CL^{5\nu_*}$. From now on we assume $C_MDt^2\rho\le\tfrac12$.

\emph{Tail.} Let $X\ni y$ with $X\not\subset\square^{\sim\nu}$. Then $d_j(X)\ge\nu-1=n-j$. Under the
floor $\kappa\ge10\log L$ on the decay rate, which is the floor displayed in the correlation theorem
(Theorem~\ref{5.1:thm-correlation-theorem}),
\begin{equation*}
e^{-\frac12\kappa d_j(X)}\le e^{-\frac12\kappa(n-j)}\le L^{-5(n-j)}=t^5 ,
\end{equation*}
which is the inequality $\exp(-\tfrac12\kappa(n-j))<(L^jL^{-n})^5$ used in
\cite[Sect.~3]{B-CMP119}. Hence
\begin{equation*}
|F(X,\mathfrak{b},y)-F(X,1,y)|\le2\mathfrak{n}e^{-\kappa d_j(X)}
\le2\mathfrak{n}e^{-\frac12\kappa d_j(X)}t^5,
\end{equation*}
and the sum of these terms over all $X\ni y$ with $X\not\subset\square^{\sim\nu}$ is at most
$2C(\tfrac12\kappa)\mathfrak{n}t^5$.

\emph{Core: the comb gauge.} Let $X\ni y$ with $X\subset\square^{\sim\nu}$. We apply the box
representation, trivial-point containment and the comb gauge
(Proposition~\ref{5.2:prop-box-comb-gauge}, display \eqref{5.2:eq-box-comb-gauge-a}): the box
representation in its form with one $\pi_j$-layer of margin, with $X\subset\square^{\sim\nu}$ and
$\square_0=\square^{\sim\nu+1}$; the containment at the trivial point; and the comb gauge centred at
$y$ on $\square_0$. The latter applies because the side of $\square_0$ satisfies
\begin{equation*}
D=M(2\nu+3)\le5ML^{\nu-1}=5M/t,
\end{equation*}
since the blocking factor $L$ is an odd integer greater than $1$. Thus, modulo gauge, $\mathfrak{b}$ is
$U_j(\square_0,\mathsf{V})$ with $\mathsf{V}=e^{iB^{\mathrm{ax}}}$. As in the proof of
Lemma~\ref{5.2:lem-marginal-bound} we set
\begin{equation*}
\varphi_X(\sigma):=F\bigl(X,U_j(\square_0,e^{i\sigma B^{\mathrm{ax}}}),
J_j(\square_0,e^{i\sigma B^{\mathrm{ax}}}),y\bigr),
\end{equation*}
so that $\varphi_X(1)=F(X,\mathfrak{b},y)$ and $\varphi_X(0)=F(X,1,y)$. By the containment at the
trivial point, $\varphi_X$ is analytic
and bounded by $\mathfrak{n}e^{-\kappa d_j(X)}$ on $|\sigma|<S:=\alpha_{3,j}/\sup|B^{\mathrm{ax}}|$; as
in that proof, the comb gauge gives $S^{-1}\le C_MDt^2\rho\le\tfrac12$, and $\varphi_X'(0)=0$. We apply
Lemma~\ref{5.2:lem-disc-bounds}(c) (display \eqref{5.2:eq-disc-bounds-a}) to $\varphi_X$ on the disc
of centre $0$ and radius $S$, at $z=1$, with $\mathsf{u}=S^{-1}\le\tfrac12$; the first-order term
vanishes, and $\mathsf{u}^3/(1-\mathsf{u})\le2\mathsf{u}^3$, so
\begin{equation}\label{5.2:eq-pointed-irrelevance-1}
\bigl|\varphi_X(1)-\varphi_X(0)-\tfrac12\varphi_X''(0)\bigr|
\le2\mathfrak{n}e^{-\kappa d_j(X)}(C_MD\rho)^3t^6 .
\end{equation}

\emph{Core: the second derivative.} For bond configurations $B_1,B_2$ on $\square_0$ let
$\langle F^{(2)}(X,y);B_1,B_2\rangle$ be the derivative $\partial^2/\partial\sigma_1\partial\sigma_2$
at $\sigma_1=\sigma_2=0$ of
\begin{equation*}
(\sigma_1,\sigma_2)\mapsto F\bigl(X,U_j(\square_0,e^{i(\sigma_1B_1+\sigma_2B_2)}),
J_j(\square_0,e^{i(\sigma_1B_1+\sigma_2B_2)}),y\bigr).
\end{equation*}
This is a symmetric bilinear form and $\varphi_X''(0)=\langle F^{(2)}(X,y);B^{\mathrm{ax}},
B^{\mathrm{ax}}\rangle$. On the polydisc $|\sigma_i|<\alpha_{3,j}/(2|B_i|_\infty)$, $i=1,2$, one has
$|\sigma_1B_1+\sigma_2B_2|<\alpha_{3,j}$ bond by bond, so by the containment at the trivial point the
function above is analytic
there and bounded by $\mathfrak{n}e^{-\kappa d_j(X)}$; the Cauchy estimate in the two variables gives
\begin{equation}\label{5.2:eq-pointed-irrelevance-2}
|\langle F^{(2)}(X,y);B_1,B_2\rangle|
\le4\mathfrak{n}e^{-\kappa d_j(X)}|B_1|_\infty|B_2|_\infty\,\alpha_{3,j}^{-2}.
\end{equation}
We now insert the expansion of the comb gauge: $B^{\mathrm{ax}}=B^{\mathrm{m}}+\mathfrak{e}$, with main
term
$B^{\mathrm{m}}_\mu(x):=\sum_{\kappa<\mu}(x_\kappa-y_\kappa)F_{\kappa\mu}(y)$. By bilinearity and
symmetry,
\begin{equation*}
\langle F^{(2)};B^{\mathrm{ax}},B^{\mathrm{ax}}\rangle-\langle F^{(2)};B^{\mathrm{m}},B^{\mathrm{m}}\rangle
=\langle F^{(2)};\mathfrak{e},B^{\mathrm{ax}}\rangle+\langle F^{(2)};B^{\mathrm{m}},\mathfrak{e}\rangle .
\end{equation*}
On $\square_0$ one has $|x-y|\le D$, hence by the comb gauge
$|B^{\mathrm{ax}}|_\infty,|B^{\mathrm{m}}|_\infty\le CD\mathfrak{a}_nt^2$ and
$|\mathfrak{e}|_\infty\le CD^2\mathfrak{a}_nt^{2+\hat\alpha}$. With
\eqref{5.2:eq-pointed-irrelevance-2} and the comparison $\mathfrak{a}_n/\alpha_{3,j}\le C\rho$ which
underlies $S^{-1}\le C_MDt^2\rho$, replacing $B^{\mathrm{ax}}$ by its main term costs at most
$C\mathfrak{n}e^{-\kappa d_j(X)}D^3\rho^2t^{4+\hat\alpha}$.

\emph{Core: replacement and extension.} Next we replace $U_j(\square_0,\cdot)$ by the map $U_j$
without the restriction to $\square_0$, and extend the sum over the core to all $X\ni y$ in
$L^{-j}\mathbb{Z}^4$. By \cite[(4.35)--(4.36)]{B-CMP109} and \cite[Sect.~3]{B-CMP119} the errors carry
the factor $e^{-\delta_0M\nu}$, $\delta_0$ being the decay constant of the kernel bounds in
\cite[(4.35)--(4.36)]{B-CMP109}, and the factor $e^{-\frac12\kappa\nu}$ respectively, each at most
$t^5$; for $X$ near
$\partial\square^{\sim\nu}$ the smallness is supplied by $e^{-\kappa d_j(X)}$, the joint rate being
$\min\{\delta_0M,\tfrac12\kappa\}\nu$.

\emph{Core: the tensor.} By \cite[(4.33)]{B-CMP109} the resulting kernel
$\sum_{X\ni y}F^{(2)}(X,y)$ at the trivial configuration is diagonal in colour: at bonds $(x,\mu)$,
$(x',\nu')$ and colour components $a,b$ it equals $\Pi_{\mu\nu'}(x,x',y)\,\delta_{ab}$. Hence the
quadratic part is
\begin{equation}\label{5.2:eq-pointed-irrelevance-3}
\sum_{X\ni y}\tfrac12\langle F^{(2)}(X,y);B^{\mathrm{m}},B^{\mathrm{m}}\rangle
=\tfrac12\sum_{\kappa<\mu,\ \lambda<\nu'}\Pi_{\mu\nu',\kappa\lambda}
\operatorname{tr}F_{\kappa\mu}(y)F_{\lambda\nu'}(y),
\end{equation}
where
\begin{equation}\label{5.2:eq-pointed-irrelevance-4}
\Pi_{\mu\nu',\kappa\lambda}:=\sum_{x,x'}\Pi_{\mu\nu'}(x,x',y)(x_\kappa-y_\kappa)(x'_\lambda-y_\lambda),
\end{equation}
the series being absolutely convergent, with $|\Pi_{\mu\nu',\kappa\lambda}|\le
C\mathfrak{n}\alpha_{3,j}^{-2}$. The numbers \eqref{5.2:eq-pointed-irrelevance-4} are independent of
$y$ by property (f3) of the family (translations). They are antisymmetric in the pair $(\mu,\kappa)$
and in the pair $(\nu',\lambda)$ by the first Ward--Takahashi identity
\cite[(3.51)--(3.55)]{B-CMP119}. By the lattice reflections they vanish unless $\mu=\nu'$ and
$\kappa=\lambda$, and by the permutations of the axes the non-vanishing ones are all equal
\cite[(3.58)--(3.61)]{B-CMP119}; these are statements of the representation theory of the hypercubic
group acting on the complex tensor \eqref{5.2:eq-pointed-irrelevance-4}. Denoting the common value by
$\beta_F$, the quadratic part \eqref{5.2:eq-pointed-irrelevance-3} equals
$\beta_F\cdot\tfrac12\sum_{\kappa<\mu}\operatorname{tr}F_{\kappa\mu}(y)^2$, and $\beta_F=\beta[F]$ by
Lemma~\ref{5.2:lem-coefficient-total}.

\emph{Core: the subtracted term.} The same analysis applied to the entire pointed covariant family
$A(h_y,\cdot)$, which has radius $O(1)$ and extracted coefficient $1$, is
\cite[(3.66)]{B-CMP119}:
\begin{equation}\label{5.2:eq-pointed-irrelevance-5}
A(h_y,\mathfrak{b})=\tfrac12\sum_{\mu<\nu'}\operatorname{tr}F_{\mu\nu'}(y)^2
+O\bigl(\nu^3\mathfrak{a}_n^2t^{4+\hat\alpha}\bigr).
\end{equation}
This is multiplied by $\beta[F]$, and $|\beta[F]|=|\beta_F|\le C\mathfrak{n}\alpha_{3,j}^{-2}$ by the
bound on \eqref{5.2:eq-pointed-irrelevance-4}; with $\mathfrak{a}_n/\alpha_{3,j}\le C\rho$ the remainder
of \eqref{5.2:eq-pointed-irrelevance-5} contributes at most $C\mathfrak{n}\nu^3\rho^2t^{4+\hat\alpha}$.

\emph{Conclusion.} The marginal parts, $\beta[F]\cdot\tfrac12\sum_{\kappa<\mu}\operatorname{tr}
F_{\kappa\mu}(y)^2$ from the core and from \eqref{5.2:eq-pointed-irrelevance-5}, cancel at the point
$y$. What remains of $\mathfrak{i}_y(\mathfrak{b})$ consists of five contributions: the tail
contribution; the remainders \eqref{5.2:eq-pointed-irrelevance-1} summed over the core; the cost of
inserting the expansion of the comb gauge; the errors of the replacement and the extension, which
carry the factor $t^5$; and the remainder of \eqref{5.2:eq-pointed-irrelevance-5} multiplied by
$\beta[F]$. Using $D\le5M\nu$, the first, second, third and fifth of these are bounded, in this order,
by
\begin{align*}
&2C(\tfrac12\kappa)\mathfrak{n}t^5,\\
&2C(\kappa)\mathfrak{n}(C_MD\rho)^3t^6\le C\mathfrak{n}\nu^3\rho^3t^6,\\
&C(\kappa)C\mathfrak{n}D^3\rho^2t^{4+\hat\alpha}\le C\mathfrak{n}\nu^3\rho^2t^{4+\hat\alpha},\\
&C\mathfrak{n}\nu^3\rho^2t^{4+\hat\alpha}.
\end{align*}
Together these give the bound $C_I\mathfrak{n}\nu^3\rho^2(1+\rho)t^{4+\hat\alpha}$.
\end{proof}

\begin{remark}
Nothing attached to a point other than $y$ enters the argument: for a source frozen cell by cell at
the value $\hat w=w(p_y)$ of the source at the plaquette $p_y$ assigned to the point $y$ by the
cell-wise freezing, the lemma applies to the family frozen at that number. The cancellation
uses only the core $\{X\ni y:\ X\subset\square^{\sim\nu}\}$ and the term $\beta[F]A(h_y,\cdot)$, on one
and the same background. The tails are bounded one by one in the proof, on whatever admissible
backgrounds they sit and whether or not they are present, and the extension of the sum over $X$ to
all $X\ni y$ concerns the kernels $F^{(2)}(X,y)$ of the whole frozen family at the trivial
configuration. The groups of terms that Ba{\l}aban's construction cuts at the localisation step are
treated in the lemma on the groups cut at the localisation step, later in this chapter.
\end{remark}

\begin{lemma}[Complex plaquette weights in the expansions]\label{5.2:lem-complex-weights}
Let $c$ be a complex weight, that is, a function from the plaquettes to $\mathbb{C}$, and let
$a_p(\mathbf{U}) = 1-\tfrac12\operatorname{tr}(\mathbf{U}+\mathbf{U}^{-1})(\partial p)$ and
$A(c,\mathbf{U}) = \sum_p c(p)\,a_p(\mathbf{U})$ be as in Definition~\ref{5.1:def-torus-floor},
for complex configurations $\mathbf{U}$ as well as $SU(N)$-valued $U$.
\begin{enumerate}
\item[(a)] Every expansion of Ba{\l}aban's construction that is applied to a Wilson term produces a
sum $\sum_p c(p)\,\tau_p$ in which each $\tau_p$ is analytic and independent of $c$
\cite[(1.17)]{B-CMP122b}.
\item[(b)] Let $c\mathbb{1}_{\mathrm{cell}}$ denote the weight equal to $c$ on the plaquettes of a
unit cell of scale $n$ and to $0$ elsewhere. Suppose that $\alpha(L^{k-n}\eta)^2\le\tfrac14$, that
$|\mathbf{U}(\partial p)-1|$ is at most $\alpha(L^{k-n}\eta)^2$ for every plaquette $p$ of the cell,
and that $H$, a field on the bonds of the cell with Hermitian traceless values (so that
$e^{i\eta H}$ is $SU(N)$-valued), satisfies $|H|,|\nabla H|\le a\le 1$ in the norms of
scale $n$. Then
\begin{gather*}
|A(c\mathbb{1}_{\mathrm{cell}},\mathbf{U})| \le 4\sup|c|\,\alpha^2,\\
\bigl|A(c\mathbb{1}_{\mathrm{cell}},e^{i\eta H}\mathbf{U})
 - A(c\mathbb{1}_{\mathrm{cell}},\mathbf{U})\bigr|
 \le C_S\sup|c|\,(\alpha+a)\,a .
\end{gather*}
\item[(c)] Let $\bar c\ge 0$ be a weight and $\vartheta$ a number with $|c(p)|\le\vartheta\,\bar c(p)$
for every plaquette $p$ (in this lemma $\vartheta\ge0$ and $\mathfrak{b}\ge0$ are local letters:
a ratio of weights and a numerical bound). Then for every $SU(N)$-valued configuration $U$
\begin{equation*}
\bigl|e^{A(c,U)}\bigr| \le e^{\vartheta A(\bar c,U)} ,
\end{equation*}
and a bound $\sum_p \bar c(p)\,|\tau_p|\le\mathfrak{b}$ of Ba{\l}aban's construction for the terms
of (a) gives $\bigl|\sum_p c(p)\,\tau_p\bigr|\le\vartheta\,\mathfrak{b}$.
\end{enumerate}
\end{lemma}

\begin{proof}
Throughout, $|\cdot|$ is the operator norm on $N\times N$ matrices; since $\operatorname{tr}\mathbb{1}=1$,
one has $|\operatorname{tr}Y|\le|Y|$ for every matrix $Y$.

(a) The weight enters $A(c,\mathbf{U})=\sum_p c(p)a_p(\mathbf{U})$ only through the constants $c(p)$,
each multiplying the single function $a_p$ of the configuration; that is, $a_p$ enters linearly.
An expansion applied to the Wilson term therefore acts plaquette by plaquette, and its result is
$\sum_p c(p)\tau_p$, where $\tau_p$ is what the expansion produces from $a_p$ alone. Thus $\tau_p$
does not depend on $c$, and it is analytic by \cite[(1.17)]{B-CMP122b}.

(b) The lattice in the cell has spacing $L^{k-n}\eta=L^{-n}$, since $\eta=L^{-k}$ by
Definition~\ref{5.1:def-torus-floor}; as the cell is a cube of side one containing this lattice,
$L^{k-n}\eta\le1$. Fix a plaquette $p$ of the cell and write $W=\mathbf{U}(\partial p)$. Since
$(W-1)(W^{-1}-1)=2-W-W^{-1}$ and $W^{-1}-1=-W^{-1}(W-1)$, cyclicity of the trace gives
\begin{equation}\label{5.2:eq-complex-weights-1}
a_p(W)=\tfrac12\operatorname{tr}\bigl[(W-1)(W^{-1}-1)\bigr]
 =-\tfrac12\operatorname{tr}\bigl[W^{-1}(W-1)^2\bigr].
\end{equation}
By hypothesis, $|W-1|\le\alpha(L^{k-n}\eta)^2\le 1/4$. The Neumann series then
gives $|W^{-1}|\le(1-|W-1|)^{-1}\le 4/3$, and also $|W|\le 5/4$. By \eqref{5.2:eq-complex-weights-1},
\begin{equation*}
|a_p(W)| \le \tfrac12\cdot\tfrac43\,\alpha^2(L^{k-n}\eta)^4=\tfrac23\,\alpha^2(L^{k-n}\eta)^4 .
\end{equation*}
The plaquettes of the cell are those of a lattice of spacing $L^{k-n}\eta$ in a cube of side one.
There are $(L^{k-n}\eta)^{-4}$ sites, each carrying six plaquettes (one for each pair of coordinate
directions), so the cell contains $6(L^{k-n}\eta)^{-4}$ plaquettes. Summing,
\begin{equation*}
|A(c\mathbb{1}_{\mathrm{cell}},\mathbf{U})|
 \le \sup|c|\cdot 6(L^{k-n}\eta)^{-4}\cdot\tfrac23\,\alpha^2(L^{k-n}\eta)^4
 = 4\sup|c|\,\alpha^2 .
\end{equation*}

For the second bound, the hypotheses on $H$ and $\nabla H$ in the norms of scale $n$ allow the
holonomy of the deformed configuration around $p$ to be written as
\begin{equation*}
(e^{i\eta H}\mathbf{U})(\partial p)=e^{i\mathsf{X}}W,
\end{equation*}
where $\mathsf{X}$ is a matrix with $\operatorname{tr}\mathsf{X}=0$ and
$|\mathsf{X}|\le 4a(L^{k-n}\eta)^2$. Since $(e^{i\mathsf{X}}W)^{-1}=W^{-1}e^{-i\mathsf{X}}$,
the definition of $a_p$ gives
\begin{equation*}
a_p(e^{i\mathsf{X}}W)-a_p(W)
 =-\tfrac12\operatorname{tr}\bigl[(e^{i\mathsf{X}}-1)W+W^{-1}(e^{-i\mathsf{X}}-1)\bigr].
\end{equation*}
Write $e^{\pm i\mathsf{X}}-1=\pm i\mathsf{X}+R_\pm$, where
$|R_\pm|\le e^{|\mathsf{X}|}-1-|\mathsf{X}|\le\tfrac12|\mathsf{X}|^2e^{|\mathsf{X}|}$.

The term of first order in $\mathsf{X}$ is $-\tfrac{i}{2}\operatorname{tr}[\mathsf{X}(W-W^{-1})]$.
Replacing $W$ and $W^{-1}$ by $1$ in its two summands would give the contributions
$-\tfrac{i}{2}\operatorname{tr}\mathsf{X}$ and $+\tfrac{i}{2}\operatorname{tr}\mathsf{X}$, which cancel
(and vanish separately, as $\operatorname{tr}\mathsf{X}=0$). Since
$W-W^{-1}=(W-1)-(W^{-1}-1)$, the first-order term equals
$-\tfrac{i}{2}\operatorname{tr}\bigl[\mathsf{X}\bigl((W-1)-(W^{-1}-1)\bigr)\bigr]$, and therefore
\begin{equation*}
\Bigl|\tfrac{i}{2}\operatorname{tr}\bigl[\mathsf{X}(W-W^{-1})\bigr]\Bigr|
 \le\tfrac12|\mathsf{X}|\,\bigl(1+|W^{-1}|\bigr)|W-1|\le\tfrac76|\mathsf{X}|\,|W-1| .
\end{equation*}
This first-order term is not required to vanish; it is kept and estimated by $|\mathsf{X}||W-1|$.

The remaining terms are bounded by
\begin{equation*}
\tfrac12\cdot\tfrac12|\mathsf{X}|^2e^{|\mathsf{X}|}\bigl(|W|+|W^{-1}|\bigr)
 \le\tfrac{31}{48}\,e^{|\mathsf{X}|}|\mathsf{X}|^2 .
\end{equation*}
Insert $|\mathsf{X}|\le 4a(L^{k-n}\eta)^2\le 4$ (recall $a\le1$ and $L^{k-n}\eta\le 1$) and
$|W-1|\le\alpha(L^{k-n}\eta)^2$. This gives
\begin{equation*}
\bigl|a_p(e^{i\mathsf{X}}W)-a_p(W)\bigr|
 \le\Bigl(\tfrac{14}{3}\,\alpha+\tfrac{31}{3}\,e^{4}a\Bigr)a\,(L^{k-n}\eta)^4 .
\end{equation*}
Multiplying by $\sup|c|$ and by the number $6(L^{k-n}\eta)^{-4}$ of plaquettes of the cell, we obtain
the bound $\sup|c|\,(28\alpha+62e^4a)a$, hence the second bound of (b) with an absolute constant
$C_S$; for instance $C_S=62e^4$ will do, since $28\alpha a+62e^4a^2\le 62e^4(\alpha+a)a$.

(c) Write $W=U(\partial p)$ for an $SU(N)$-valued configuration $U$. Then $W^{-1}=W^{*}$, so
$a_p(U)=\tfrac12\operatorname{tr}[(W-1)(W-1)^{*}]$ by
\eqref{5.2:eq-complex-weights-1}; hence $a_p(U)\ge0$. Therefore
\begin{equation*}
\bigl|e^{A(c,U)}\bigr|=e^{\operatorname{Re}A(c,U)}
 =e^{\sum_p\operatorname{Re}c(p)\,a_p(U)}
 \le e^{\sum_p|c(p)|\,a_p(U)}\le e^{\vartheta\sum_p\bar c(p)\,a_p(U)}
 =e^{\vartheta A(\bar c,U)} .
\end{equation*}
Likewise,
\begin{equation*}
\Bigl|\sum_p c(p)\tau_p\Bigr|\le\sum_p|c(p)|\,|\tau_p|
 \le\vartheta\sum_p\bar c(p)\,|\tau_p|\le\vartheta\,\mathfrak{b} .
\end{equation*}
\end{proof}

In the expansions, part (c) is applied with $c=\zeta_j$, the running shift at level $j$, and
$\bar c=x_j(\cdot)$, the inverse coupling at level $j$ read plaquette by plaquette: the plaquette
function $x_j(\cdot)$ is assembled from the inverse couplings $x_i$ by cut-off functions
$\varphi_i$, each of which passes from one member of the sequence $(x_i)$ to the next across a
transition layer. By the running-shift statement of the correlation theorem one has
$\sup|\zeta_j|\le\tfrac83 r$, where $r$ is the radius of the ball $\mathbb{D}$ in
\eqref{5.1:eq-torus-floor-b}. Moreover, $x_j(p)>g^{-2}$ for every plaquette $p$: $x_j(p)$ lies
between two consecutive members of the sequence $(x_i)$, since the transition layers of the functions
$\varphi_i$ are disjoint, and each member $x_i$ entering at level $j$ exceeds $g^{-2}$, the levels
of the run preceding the first passage $K(g_0,g)$. No information on the signs of the $\beta_i$ is
used. With $\theta_W:=\tfrac83\,r\,g^2$ this gives
\begin{equation*}
|\zeta_j(p)|\le\tfrac83\,r=\theta_W\,g^{-2}<\theta_W\,x_j(p),
\end{equation*}
so (c) holds with $\vartheta=\theta_W$ and $\bar c=x_j(\cdot)$.

\begin{lemma}[Freezing the weight at one plaquette]\label{5.2:lem-weight-freezing}
Let $j<k$, and let $F(X;\cdot)$ be analytic on $\mathcal{W}_j(X)$ with $|F(X;w')|\le\mathfrak{n}$ for
every $w'\in\mathcal{W}_j(X)$. Let $w\in\mathcal{W}_k(X)$ and $c\in\operatorname{plaq}(X^+)$, and write
$w(c)$ also for the constant weight on $\operatorname{plaq}(X^+)$ whose value is $w(c)$. Then
\begin{equation}\label{5.2:eq-weight-freezing-1}
\bigl|F(X;w)-F(X;w(c))\bigr|
\le 2\mathfrak{n}\,\omega_{k,j}\Bigl(1+\frac{\operatorname{diam}_j(X^+)}{10M}\Bigr).
\end{equation}
Here $\operatorname{diam}_j(X^+)$ denotes the diameter of $\operatorname{plaq}(X^+)$ in units of the
level-$j$ lattice, the distance with respect to which $\operatorname{Lip}_j$ is taken.
\end{lemma}

\begin{proof}
We use the classes of complex plaquette weights of Definition~\ref{5.1:def-weight-classes}, in
the notation of \eqref{5.1:eq-weight-classes-a}. For $\sigma\in\mathbb{C}$ put
\begin{equation*}
w_\sigma := w(c)+\sigma\bigl(w-w(c)\bigr),
\qquad
S := \omega_{k,j}^{-1}\min\Bigl\{1,\frac{10M}{\operatorname{diam}_j(X^+)}\Bigr\},
\end{equation*}
so that $w_0=w(c)$ and $w_1=w$; moreover
$S\,\omega_{k,j}\le\min\{1,\,10M/\operatorname{diam}_j(X^+)\}$.

We first bound the level-$j$ Lipschitz constant of $w$. A Lipschitz constant for distances
measured in units of the level-$j$ lattice is $L^{-(k-j)}$ times the one for the level-$k$
lattice, so $\operatorname{Lip}_j w=L^{-(k-j)}\operatorname{Lip}_k w$. Since
$w\in\mathcal{W}_k(X)$ we have $\operatorname{Lip}_k w<\mathfrak{l}_k$, and by the definition of
$\mathfrak{l}_j$ one has $\mathfrak{l}_k=L^{(k-j)/2}\mathfrak{l}_j$. Hence
\begin{equation}\label{5.2:eq-weight-freezing-2}
\operatorname{Lip}_j w = L^{-(k-j)}\operatorname{Lip}_k w
< L^{-(k-j)}L^{(k-j)/2}\,\mathfrak{l}_j = \omega_{k,j}\,\mathfrak{l}_j .
\end{equation}

We show that $w_\sigma\in\mathcal{W}_j(X)$ whenever $|\sigma|<S$. Since $w(c)$ is constant,
\begin{equation*}
\operatorname{Lip}_j(w_\sigma)=|\sigma|\operatorname{Lip}_j w
=|\sigma|L^{-(k-j)}\operatorname{Lip}_k w
< S\,\omega_{k,j}\,\mathfrak{l}_j\le\mathfrak{l}_j ,
\end{equation*}
by \eqref{5.2:eq-weight-freezing-2} and $S\omega_{k,j}\le1$. For the size, let
$p\in\operatorname{plaq}(X^+)$. Then $|w(c)|<r_k$ because $w\in\mathcal{W}_k(X)$, and
$|w(p)-w(c)|\le\operatorname{Lip}_j w\cdot\operatorname{diam}_j(X^+)$, so
\begin{align*}
|w_\sigma(p)|
&\le |w(c)|+|\sigma|\,|w(p)-w(c)|
< r_k + S\,\omega_{k,j}\,\mathfrak{l}_j\operatorname{diam}_j(X^+)\\
&\le r_k+10M\mathfrak{l}_j\le r_j ,
\end{align*}
where we used \eqref{5.2:eq-weight-freezing-2}, then
$S\omega_{k,j}\le 10M/\operatorname{diam}_j(X^+)$, and finally the inequality
$r_j-r_k\ge 10M\mathfrak{l}_j$, valid for $j<k$, of Definition~\ref{5.1:def-weight-classes}.
Thus $w_\sigma\in\mathcal{W}_j(X)$ for $|\sigma|<S$.

Consequently $f(\sigma):=F(X;w_\sigma)$, the composition of $F(X;\cdot)$ with an affine map of
$\sigma$, is analytic on the disc $D(0,S)$ and satisfies $|f|\le\mathfrak{n}$ there. Note also
that
\begin{equation*}
\frac1S=\omega_{k,j}\max\Bigl\{1,\frac{\operatorname{diam}_j(X^+)}{10M}\Bigr\}
\le\omega_{k,j}\Bigl(1+\frac{\operatorname{diam}_j(X^+)}{10M}\Bigr).
\end{equation*}

If $S>1$, then $1\in D(0,S)$, and bound (a) of Lemma~\ref{5.2:lem-disc-bounds}
(see \eqref{5.2:eq-disc-bounds-a}), applied to $f$ with centre $0$, radius $S$, bound
$\mathfrak{n}$ and $z=1$, gives
\begin{equation*}
|F(X;w)-F(X;w(c))|=|f(1)-f(0)|\le\frac{2\mathfrak{n}}{S},
\end{equation*}
and \eqref{5.2:eq-weight-freezing-1} follows from the bound on $1/S$.

If $S\le1$, both $w(c)=w_0$ and $w$ lie in $\mathcal{W}_j(X)$: the first because $0\in D(0,S)$,
the second because $|w|<r_k\le r_j$ and, by \eqref{5.2:eq-weight-freezing-2},
$\operatorname{Lip}_j w<\omega_{k,j}\mathfrak{l}_j\le\mathfrak{l}_j$. Hence
\begin{equation*}
|F(X;w)-F(X;w(c))|\le 2\mathfrak{n}\le\frac{2\mathfrak{n}}{S},
\end{equation*}
and \eqref{5.2:eq-weight-freezing-1} follows again from the bound on $1/S$.
\end{proof}

The constant that accompanies this lemma in later estimates is
\begin{equation}\label{5.2:eq-weight-freezing-a}
C_B := \sup_X\, 2\Bigl(1+\frac{\operatorname{diam}_j(X^+)}{10M}\Bigr)
e^{-\frac14\delta\kappa d_j(X)} \le 5 .
\end{equation}

\begin{definition}[Sites, units and response maps]\label{5.2:def-sites-units}
A \emph{site} is a place at which Ba{\l}aban's construction expands terms of the effective action about a
new background. There are three sites.
\begin{itemize}
\item[(T)] The expansion \cite[(3.16)]{B-CMP119}, in the form of the creating step used in the proof
of the correlation theorem (Theorem~\ref{5.1:thm-correlation-theorem}), together with
\cite[(3.20)--(3.47)]{B-CMP119}. Both are read in
the splitting of these terms given by the splitting lemma of this chapter.
\item[(P)] The $N(g)$ preparatory integrals \cite[(1.60)--(1.63)]{B-CMP122a}, where $N(g)$ denotes
the number of preparatory stages. Both slots of the bracket in
\cite[(1.60)]{B-CMP122a} are included: the response $\Phi(B'_j)$, and also the integration
variable $A_j$.
\item[(L)] The localisation brackets \cite[(1.59)--(1.66)]{B-CMP122b}. At this site there is no
integral. The exterior fluctuation fields enter as parameters, as in \cite[(1.49)]{B-CMP122b} and in
the lemma on the exterior fields.
\end{itemize}

A \emph{unit} at a site is a triple
\begin{equation}\label{5.2:eq-sites-units-1}
v = (j, S_v, f_v).
\end{equation}
Here $j$ is a level and $S_v$ is a subset of the level-$j$ lattice, the read set of the unit. The
function $f_v$ is analytic on a
domain $\mathfrak{D}_v$ of level-$j$ configurations on $S_v$, and it is gauge invariant in the
following sense.
\begin{itemize}
\item[(a)] Let $v$ belong to the first three classes, to the fifth class or to the class $\partial$ of
the table of units of this chapter. Then $f_v$ is gauge invariant as a function of its own
configuration. This holds in the plain sense of \cite[(2.27)(i)]{B-CMP119}, according to which the
corresponding term depends on $U_k$ restricted to its localisation domain $X$.
\item[(b)] Let $v$ belong to the fourth class of that table. Then $f_v$ is gauge invariant jointly
with its live exterior slots, under $SU(N)$-valued gauge transformations, as in
\cite[(2.41)(iii)]{B-CMP119}. For such a unit the invariance asserted is this joint one, and not the
invariance of $f_v$ as a function of its own configuration alone.
\end{itemize}

Each site comes with its \emph{response maps} $\mathsf{R}_p(\mathsf{t},\sigma)$. These maps are
independent of the source variables $z$ of the complex source $w=\sum_i z_i f_i$. Here
$p=(\square,Y_0)$ denotes Ba{\l}aban's localisation datum, consisting of a block $\square$ and a set
$Y_0$.
\begin{itemize}
\item $\mathsf{t}$ is a complex amplitude.
\item $\sigma=(\sigma(\Delta))_\Delta$ are the decoupling parameters, one for each cell $\Delta$.
\item $\mathsf{w}_{\mathrm{site}}$ denotes the radius of the amplitude disc attached to the site, that
is, $\mathsf{w}_T$, $\mathsf{w}_P$ or $\mathsf{w}_L$ at (T), (P) or (L) respectively.
\end{itemize}
A response map is a holomorphic map
\begin{equation}\label{5.2:eq-sites-units-2}
\begin{split}
\mathsf{R}_p \colon\ & \{\mathsf{t} : |\mathsf{t}| \le \mathsf{w}_{\mathrm{site}}\}
\times \{\sigma : |\sigma(\Delta)| \le e^{\kappa_1} \text{ for every cell } \Delta\}\\
&\times (\text{the output tube of the site}) \times \operatorname{supp}\chi
\ \longrightarrow\ \{\text{admissible backgrounds}\}.
\end{split}
\end{equation}
At the site (T) the domain in \eqref{5.2:eq-sites-units-2} carries one further factor:
\begin{equation}\label{5.2:eq-sites-units-3}
\{\mathsf{B} : \mathsf{B} \text{ is } \mathfrak{g}^{\mathbb{C}}\text{-valued on the free bonds }
\mathfrak{X},\ |\mathsf{B}| < T^{\mathbb{C}}\}.
\end{equation}
Here $\mathfrak{g}^{\mathbb{C}}$ is the complexified Lie algebra of $SU(N)$, and $T^{\mathbb{C}}$ is the
bound on the fluctuation field $\mathsf{B}$; in $T^{\mathbb{C}}$ the superscript is
part of the name and does not denote a complexification. At (T) the maps are taken
under the homogeneous containment given by the homogeneous-containment lemma of this chapter.

In \eqref{5.2:eq-sites-units-2} the amplitude $\mathsf{t}$ multiplies the shift piece, which is the
product of $\hat{\chi}_{\square}$ with the response. Here $\hat{\chi}_{\square}$ is the function
denoted $\zeta_{\square}$ in Ba{\l}aban's construction. Finally, $\eta_p \le 1$ denotes the decay
factor between the block $\square$ and the set $S_v$.
\end{definition}

\begin{proposition}[The three expansion sites and their disc radii]\label{5.2:prop-three-sites}
Let the sites $(T)$, $(P)$, $(L)$, the units $v=(j,S_v,f_v)$ with their analyticity domains
$\mathfrak{D}_v$, the localisation data $p=(\square,Y_0)$, the decay factors $\eta_p\le 1$ and the
response maps $\mathsf{R}_p(\mathsf{t},\sigma)$ be as in Definition~\ref{5.2:def-sites-units}. A unit
listed in the fourth entry of the table \eqref{5.2:eq-unit-table-a} is called a \emph{type-four unit};
the units of the other entries are those of types one to three and five and those of the entry marked
$\partial$. A \emph{live law site} of a type-four unit $v$ is a point $s$ at which a live exterior
slot of $v$ is read; it is of the \emph{outer class} if $s\notin Z$ and of the \emph{inner class} if
$s\in Z_{k-N_k}$, $Z$ being the current large-field region. A \emph{frame} $W$ is a region of the margin
theorem with a creation step $k_W$, so that $W\subset Z_{k_W}$, with a set $\Sigma_W\subset W$ of law
sites and its own chart size (hypothesis (vi)).
Assume:
\begin{enumerate}
\item[(i)] the orbit-datum bound for type-four units, clauses (a), (b), (c), (g), together with the
definition of the orbit datum $c^\sharp_v$ on which it rests (its holomorphy property and its
real-slice property), the complex induction hypothesis at level $k$, and the riding-rotation lemma;
\item[(ii)] the inheritance theorem for the chain terms $\mathbb{C}^{(n)}_k(X,\mathcal{S};w)$, its
corollary on the inherited domain $\mathfrak{D}^{\flat}_n(X)$, and the chain-containment claim;
\item[(iii)] the presented-chain containment theorem: for every unit $\mathfrak{E}$ of the consumer
class $\mathcal{C}$, every $\zeta\in\mathfrak{S}_{\mathrm{src}}$, every real configuration
$\varpi$ in the support of $\Pi$, every $|\sigma|\le e^{\kappa_1}$, every
$\mathfrak{b}\in\{\mathfrak{b}_2,\mathfrak{b}_3,\mathfrak{b}_4\}$ and the pairs
$(\mathfrak{b}_2,\mathfrak{b}_3)$, $(\mathfrak{b}_3,\mathfrak{b}_4)$, the restriction to $X$ of the
presented carrier pair $((\hat{\mathcal{U}}^{\mathfrak{b}})^h,\operatorname{Ad}_h
\hat{\mathcal{J}}^{\mathfrak{b}})$ lies in $\mathfrak{D}_{\mathfrak{E}}$ with each entry at most
$1/16$ of its threshold on $X\cap(\Lambda^{\natural}\cup\mathfrak{N}^{\mathrm{rec}})$ (the
containment clause), and $(\zeta,\sigma,t)\mapsto F^{\mathfrak{b}}_{\mathfrak{E}}$ is holomorphic
(the holomorphy clause), together with the consequences that theorem draws from these two clauses
and its remaining clause with the three conditions under which that clause holds;
here $\Pi$ and the region $\mathfrak{N}^{\mathrm{rec}}$ are those of that theorem, and
$F^{\mathfrak{b}}_{\mathfrak{E}}$ is the function of $(\zeta,\sigma,t)$, $t\in[0,1]$ being
Ba{\l}aban's interpolation parameter, which that theorem attaches
to $\mathfrak{E}$ and $\mathfrak{b}$; its inheritance conjunct, under the fourth condition of that
theorem, is hypothesis (ii) for the terms of $\mathbf{B}''_k$, its arrest conjunct, under the
condition of the theorem on arrested left-overs, is hypothesis (iv) for the arrested left-overs
$\mathcal{C}_{\mathrm{arr}}\subset\mathcal{C}$, and the live slots of every type-four unit are read
through hypothesis (i);
\item[(iv)] the theorem on arrested left-overs, together with its condition;
\item[(v)] the layer-increment condition on the
three $\mathbf{R}$-layer increments $P$, $P'$, $P''$ (each at most
$C_Pe^{-3\delta_1^{\mathrm{sc}}R_k}$, where $R_k:=R(x_k)$) of the layer-increment proposition, in the
form in which it is read on the disc of site $(L)$,
\begin{equation}\label{5.2:eq-three-sites-1}
\sup_{\mathsf{T}\ge\mathsf{T}_\gamma-1}32\,C_{\theta'}^{-1}(\mathsf{T}+2)^{q-p_1}\,C_P'\,
e^{-\frac32\delta_1^{\mathrm{sc}}R(e^{\mathsf{T}})}\le\beta_{\mathrm{pr}}\theta_S ,
\end{equation}
where $R(e^{\mathsf{T}})$ is $R$ at $x=e^{\mathsf{T}}$, that is, at the coupling $g$ with
$\log g^{-2}=\mathsf{T}$; applied at $\mathsf{T}=\mathsf{T}_k-1\ge\mathsf{T}_\gamma-1$, with
$\mathsf{w}_L\ge1$ and $C_P'\ge C_P$ ($\mathsf{w}_L$, $C_P'$ as below), it contains the undilated
condition $16\max\{P,P',P''\}\le\beta_{\mathrm{pr}}\theta_S$;
\item[(vi)] the ceiling: with
\begin{equation*}
\mathsf{a}_J(x):=2.1\,C_{\theta'}^{-1}(\log x+2)^{q-p_1}\cdot\tfrac12 e^{-2}C_{\natural}\,\bar\alpha(x),
\qquad \mathsf{a}_T(x):=B_H\,\alpha_*(x),
\end{equation*}
for all $x\ge\gamma^{-2}$
\begin{equation}\label{5.2:eq-three-sites-b}
\begin{gathered}
(x-c_5)^2E(x)^{\delta_\natural/2}\max\{\mathsf{a}_J(x),2.1K_u\mathsf{a}_T(x)\}
\max\Bigl\{\frac{1}{\varepsilon_\star},\frac{8e\,C_{\mathrm{fr}}(1+B^\sharp)^{1/2}}
{\varpi_{\mathrm{fr}}\,\alpha_*(x)}\Bigr\}\le 1,\\
\mathsf{a}_J(x)E(x)^{\delta_\natural/2}\le\tfrac13K_u\min\{10^{-2}/K_u,\,c_4\},
\end{gathered}
\end{equation}
where $B^\sharp$ is read as $b_+$, $C_{\mathrm{fr}}$ and $\varpi_{\mathrm{fr}}$ are the constants of
the chart size $\mathsf{e}(s)=\varpi_{\mathrm{fr}}\alpha_{*,k_W}/4C_{\mathrm{fr}}$ at a law
site $s\in\Sigma_W$ of a frame $W$, $\mathsf{a}_T$ is the amplitude of site
$(T)$, and $C_{\mathrm{fr}}$, $\varpi_{\mathrm{fr}}$, $\varepsilon_\star$,
$\delta_\natural$, $C_\natural$, $K_u$, $B_H$, $c_4$, $C_{\theta'}$ are fixed before $\gamma$;
\item[(vii)] the standing conditions: the bracket condition $4C_A'\theta'(g_k)\le1$; the
live-region conditions $5\delta_\natural R_{\min}\ge2$ and $N_k\le x_k$, with $R_{\min}$ the
minimal-radius constant of the distance lemma, the frame-shell condition of that lemma, and the upper
half of the reference flow bound; the relation $(\log(x_k-c_5))^{r_{\mathrm{pr}}}\le\check R_k$,
which with the definition of the decay function $E$ gives
$e^{-\frac12\delta_\natural(\check R_k-2)}\le E(x_k)^{\delta_\natural/2}$; the summed-margin bound
$\sum_{k'}\eta_{k'}(s)\le\frac14\mathsf{e}(s)$, $\eta_{k'}(s)$ the rotation allowance of step $k'$ at
$s$; the rate condition $10\delta_1^{\mathrm{sc}}R_{\min}\ge2$ of the layer-increment proposition;
Ba{\l}aban's condition $\delta_1M/M_1\ge2$ \cite[\S1]{B-CMP122a}; the condition
\eqref{5.2:eq-half-disc-containment-d} on $C_{\theta'}$ and the recession condition of the
$\mathbf{J}$-recession lemma;
\item[(viii)] the containment on the half disc outside the small-field region: at the clause
locations of $\mathfrak{D}_v$ other than the $\mathfrak{T}$-locations, the response maps of
type-four units lie in $\mathfrak{D}_v$ on the disc $|\mathsf{t}|<\rho_A$ defined below, under the
conditions \eqref{5.2:eq-half-disc-containment-c}, \eqref{5.2:eq-half-disc-containment-d},
\eqref{5.2:eq-half-disc-containment-e}.
\end{enumerate}
The supremand of \eqref{5.2:eq-three-sites-1} tends to zero as $\mathsf{T}\to\infty$, and each left
member of \eqref{5.2:eq-three-sites-b} tends to zero as $x\to\infty$: the factor
$E(x)^{\delta_\natural/2}$ obeys
\begin{equation*}
E(x)^{\delta_\natural/2}\le e^{\delta_\natural}(x-c_5)^{-\frac12\delta_\natural\log(x-c_5)},
\end{equation*}
which lies below every power of $x$ for $r_{\mathrm{pr}}\ge2$ and $\log(x-c_5)\ge1$ (the latter for
all $x\ge\gamma^{-2}$ as soon as $\gamma^{-2}\ge c_5+e$), while the remaining factors are powers of $x$
times polylogarithms ($\bar\alpha/\alpha_*$ and $(\log x+2)^{q-p_1}$ being polylogarithmic); so, for
$r_{\mathrm{pr}}\ge2$, conditions (v) and (vi) hold for all $\gamma$ in a neighbourhood of
$0$, and neither involves $K$, $k$, an age or the volume.
Then the following hold.

\smallskip\noindent\emph{Site $(T)$.} The expansion is that of \cite[(3.4), (3.7), (3.25)]{B-CMP109}
and \cite[(1.9)--(1.22)]{B-CMP116} on the multiscale cover of \cite[\S3]{B-CMP119}, with containment
by \cite[Lemma~4, (3.16), (3.43)]{B-CMP119}. On the polydisc $|\mathsf{B}|<T^{\mathbb{C}}$,
$T^{\mathbb{C}}:=\mathfrak{b}_k+(\Theta_\flat+1)\varrho$, the amplitude disc has radius
\begin{equation}\label{5.2:eq-three-sites-2}
\mathsf{w}_T:=\mathsf{w}^{\mathbb{C}}:=\frac{1}{\theta^{\mathbb{C}}_k}-1,\qquad
\theta^{\mathbb{C}}_k:=\frac{K_Rg_kT^{\mathbb{C}}}{\alpha_{*,k}},
\end{equation}
(the quantities of \eqref{5.2:eq-activity-threshold-a}), the largest $\mathsf{w}$ with
$(1+\mathsf{w})K_Rg_kT^{\mathbb{C}}\le\alpha_{*,k}$; $\mathsf{w}^{\mathbb{C}}$ stays
bounded as $\gamma\to0$ and is at least $1$ under the condition of the complex-radius lemma,
and this radius is the one used for every amplitude disc at $(T)$. Ba{\l}aban's radius on the real
circle is $\mathsf{w}_T^{(1.22)}=\alpha_2/(4B_0C_1e^{16\kappa_1}\delta_k)\ge1/\theta_k$, with the
constants $B_0$, $C_1$, $\alpha_2$, $\theta_k$, $\delta_k$ of
\cite[(1.22)]{B-CMP116}; the disc is run at $\mathsf{w}^{\mathbb{C}}$ because the hypothesis
$(1+\eta_p|\mathsf{t}|)K_Rg_k\sup|\mathsf{B}|\le\alpha_{*,k}$ of
\eqref{5.2:eq-response-containment-a} must hold at $|\mathsf{t}|\le\mathsf{w}/\eta_p$, and it fails on
the disc of radius $\mathsf{w}_T^{(1.22)}/\eta_p$: with $K_R\ge16B_0C_1e^{16\kappa_1}/c_2$ and
$\alpha_2\ge c_2\alpha_*$ one has $\mathsf{w}_T^{(1.22)}K_R\delta_k\ge4\alpha_*$.
The type-four units met at $(T)$ (older $\mathbf{B}^{(j)}$, $\mathbf{B}'^{(j)}$, and chain terms left
on live arrested pieces) carry live exterior slots which, at real data, are read rotated by the
$G$-valued map $u=u(b,\mathbb{U})$ from the Landau to the axial gauge of the response, in
Ba{\l}aban's form ``$(M_j(\exp i\xi H)M_j(U^{(n+1)}))^{u_j}$'' \cite[\S1]{B-CMP122a}, joint
invariance transporting $u$ onto the gauge transformation of the background. On the amplitude discs,
where $(\mathsf{t},\sigma,\mathsf{B})$ are complex, the local gauge factor $\mathfrak{u}$ is not
$G$-valued; there the unit is read as its orbit datum
$c^\sharp_v(S_v;\,\cdot\,;\mathbf{A},\mathbf{\Phi},\,\cdot\,)$, evaluated at the carrier pair in
the Landau gauge and at the rotation argument multiplied by $\mathfrak{u}$ ($\mathbf{\Phi}$ being
the further holomorphic coordinate which the orbit datum carries), and bounded by clause
(c) of hypothesis (i) through the complex-radius corollary and the riding-rotation lemma, at the
price of the factor $\eta_p^{1/2}$ in place of $\eta_p$.

\smallskip\noindent\emph{Site $(P)$.} At the $n$-th preparatory stage the response $\Phi(b)$ to an
admissible input $b$ satisfies
$\Phi(b)\in\tilde U^{(n)c}_k$, and the same holds when the layer $\eta\mathbb{H}(b)$ is replaced by
$\sum_\square t_\square\zeta_\square\eta\mathbb{H}(b)$, for dilations up to the bound
$r_*\propto1/\vartheta_0$, all blocks at once, for both slots of the bracket of
\cite[(1.60)]{B-CMP122a} (the response and the integration variable $A_j$), each local in the
parameters $t_\square$ (the localised-response
corollary, whose constant is $\vartheta_0$, whose dilation bound is $r_*$, and whose factor is
$\eta$); by the proposition on the expansion at the preparatory site, $\mathsf{w}_P\ge1/\theta'$ with
$\theta'\propto(\log g_k^{-2})^{p_1-q}$, and the
integration slot $A_j$ of \cite[(1.60)]{B-CMP122a} is expanded by the second M\"obius family at weak
radius $\varrho_j=c_A\alpha_{*,j}/g_j$, with $c_A$ as in \eqref{5.2:eq-activity-threshold-a} and
$r_A:=\varrho/\mathfrak{b}_k$, which grows like
$(\log g^{-2})^{q-p_0}$ and tends to infinity, the strong cubes held real in the upper box
$\mathfrak{b}^\sharp$;
\cite[(1.9)]{B-CMP122a} is read as \eqref{5.2:eq-second-stage-activities-a}; at every stage the
earlier stages' chain terms and the frozen terms of the live arrested pieces are read as orbit data,
with the rotation by $u$, on the polydiscs in the parameters $t_\square$ of the preparatory
expansion, through clauses (b), (c)
of hypothesis (i) and the polydisc corollary; no constant of the preparatory expansion changes.

\smallskip\noindent\emph{Site $(L)$.} The four expansions
\cite[(1.51), (1.54), (1.57), (1.58)]{B-CMP122b}, with \cite[(1.61)--(1.65)]{B-CMP122b}, are each
``an analytic function of the variables $(\mathbf{U},\mathbf{J})$ in the space
$\tilde U^c_k(Y_4,\tilde\alpha_0,\tilde\alpha_1)$, with values in the space
$\tilde U''^{c}_k(Y_1,\tilde\alpha_0,\tilde\alpha_1)$ \dots\ contained in the analyticity domain of
the term $E(X)$ \dots\ hence \dots\ of all terms'' \cite[\S1]{B-CMP122b}. For the terms of
$\mathbf{B}''_k$, ``the sum of the new boundary terms created by the $N$ preliminary integrations''
\cite[\S1]{B-CMP122b}, that is, for the chain terms $\mathbb{C}^{(n)}_k(X,\mathcal{S};w)$ with
$n\le N$, Ba{\l}aban rests this containment on an inequality of the preliminary integrations:
``It is clear for terms of $\mathbf{B}''_k$, by the inductive assumptions \dots\ (see the
inequality (1.70) [IV] and the assumptions connected with it)'' \cite[\S1]{B-CMP122b}; here, by
hypothesis (ii)
--- $\mathbb{C}^{(n)}_k(X,\mathcal{S})$ being a stored term with slot datum on
$\mathfrak{D}^{\flat}_n(X)\times\mathfrak{D}_{\mathcal{S}}(X)$ ---, the configurations of the chains
and their response dilations lie in the inherited domain $\mathfrak{D}^{\flat}_n(X)$:
\begin{itemize}
\item[($\alpha$)] on $X\cap Z^{\mathrm{on}}$, the part $Z^{\mathrm{on}}$ of the large-field region
on which the inherited domain imposes its own list of clauses $\mathfrak{L}_n$, they obey
$\mathfrak{L}_n$;
\item[($\beta$)] at the $\mathfrak{T}$-locations $y\in X\cap\mathfrak{T}$ the fifth and eighth
entries of the list of entries defining the configuration space are those
of the input space $\tilde U^c_k(Y_4)$, the scale-drop clauses see tube entries
$\le1.1r_a\le1.5r_a$ ($r_a$ the tube radius read by those clauses), and the other entries (the
composite entries attached to every region $X'\subset X$ of the triples over which the composite
clauses run included) exceed the input space's thresholds only by the $\mathbf{R}$-layer increments,
which are at most $\frac14\mathrm{room}^\flat(y)$ under hypothesis (v); here, with $r_{\min}(y)$
the minimal radius at $y$ and $m(y)$ the shell index of $y$,
$\mathrm{room}^\flat(y):=\min\{\frac12r_{\min}(y),\beta_{\mathrm{pr}}\}$, which is
$\ge\frac12\beta_{\mathrm{pr}}\theta_S2^{-(k-m(y))}$ on the untouched shells $m(y)<k$ and
$\ge\frac12\beta_{\mathrm{pr}}$ on $\Omega_k\cap\mathfrak{T}$ under the top layer of the layer
schedule, $\Omega_k$ being the small-field region at level $k$; the composite entries are controlled
for all triples at once, the fat regularity bound
$1.1<1+2\beta_{\mathrm{pr}}-2\mathsf{w}_0\beta_{\mathrm{pr}}$ holds, and the fat clauses of the
inherited domain are met;
\item[($\gamma$)] on $X\cap\Lambda^{\natural}$ they keep the margin one half.
\end{itemize}

\smallskip\noindent\emph{Type-four units at $(L)$.} Let $v$ be $\mathbb{C}^{(n)}_k(X,\mathcal{S};w)$,
$\mathbf{B}^{(j)}$, $\mathbf{B}'^{(j)}$ or a frozen member of $\mathfrak{F}(\mathfrak{a})$, with live
exterior slots $\mathbf{A}$, carrying a localisation datum $p$. For a positive function
$\varepsilon$ on $S_v$,
$\mathfrak{G}_\varepsilon(S_v)$ denotes the polydisc of rotation arguments $\mathfrak{o}$ on $S_v$
with polar size $\operatorname{pol}(\mathfrak{o}(s'))<\varepsilon(s')$ at every $s'\in S_v$. The
orbit datum of $v$ at the presented rotation $\mathfrak{Q}^{\mathfrak{b}}(\zeta,\mathsf{t},\sigma)$,
\begin{equation}\label{5.2:eq-three-sites-3}
f_v\circ\mathsf{R}_p(\mathsf{t},\sigma):=c^\sharp_v\bigl(S_v;\mathsf{R}_p(\mathsf{t},\sigma;
\mathbf{U},\mathbf{J});\mathbf{A},\mathbf{\Phi},\mathfrak{Q}^{\mathfrak{b}}(\zeta,\mathsf{t},\sigma);
w\bigr),
\end{equation}
coincides with Ba{\l}aban's term on the real slice, is holomorphic in
$(\mathsf{t},\sigma,(\mathbf{U},\mathbf{J}),w)$ jointly with the weak slot coordinates, the strong
and silent parts of the slots being held real and fixed, and obeys
$|f_v\circ\mathsf{R}_p-f_v(1)|\le\mathfrak{N}(v):=2\mathfrak{N}^{\mathbb{C}}_{\varepsilon_k}(v)$,
where $f_v(1)$ is the value at the trivial configuration and
$\mathfrak{N}^{\mathbb{C}}_{\varepsilon_k}(v)$ is the supremum of the orbit datum over the chart
$\mathfrak{G}_{\varepsilon_k}(S_v)$, $\varepsilon_k=\varepsilon_k(s)$ being the radius function of
the chart of rotation arguments at level $k$. Put
\begin{equation*}
\mathsf{w}_L:=\frac{1}{\theta'(g_k)}=C_{\theta'}^{-1}\mathsf{T}_k^{\,q-p_1},\qquad
\rho_A:=\frac{\mathsf{w}_L}{\eta_p^{1/2}} ,
\end{equation*}
and
\begin{equation*}
\varepsilon'(s):=\varepsilon_k(s)-\eta_0(s)-\mu(s),\qquad
\mu(s):=\frac{\mu_k\mathsf{e}(s)}{\varepsilon_\star},\qquad
\mu_k:=\varepsilon_\star(x_k-c_5)^{-2},
\end{equation*}
where $\eta_0(s)$ is the supremum, over the parameters
$(\sigma,\mathbf{U},\mathbf{J},\mathbf{\Phi},\zeta,w)$ and over $t\in[0,1]$, of the polar size of the
undilated presented rotation
$\mathfrak{Q}^{\mathfrak{b}}(\zeta,0,\sigma)(s)$, and $\mathsf{e}(s)$ is as in (vi) at the law sites
of frames and $\mathsf{e}(s)=\varepsilon_\star$ at the other law sites.
Then on $|\mathsf{t}|<\rho_A$, $|\sigma(\Delta)|\le e^{\kappa_1}$, the response map lies in
$\mathfrak{D}_v$, the orbit datum read through the $\mathsf{t}$-dilated rotation is holomorphic
over the chart $\mathfrak{G}_{\varepsilon'}(S_v)$ of the presentation and bounded there by
$\mathfrak{N}(v)$, and Ba{\l}aban's piece
\begin{equation*}
v_p=\prod_{\Delta}\int_0^1d\mathsf{s}(\Delta)\,\partial_{\mathsf{s}(\Delta)}\,
\partial_{\mathsf{t}}\big|_0\,f_v(\mathsf{R}_p(\mathsf{t},\mathsf{s}))
\end{equation*}
(\cite[(3.7)]{B-CMP109}, \cite[(1.10), (1.23)]{B-CMP116}; the product over the cells $\Delta$ of
$Y_0$ off $\bar S_v$, the union of the own-scale exempt family of cells of $v$, which contains
$S_v$, $m(p)$ their number) obeys
\begin{equation}\label{5.2:eq-three-sites-a}
|v_p|\le2\,\mathsf{w}_L^{-1}\,\eta_p^{1/2}\,e^{-(\kappa_1-1)m(p)}\,\mathfrak{N}(v).
\end{equation}

\smallskip\noindent\emph{Complex sources.} The orbit datum of $\mathbb{C}^{(n)}_k(X,\mathcal{S};w)$
is the orbit datum of the family of the complex-source corollary, holomorphic in
$w\in\mathcal{W}_k(X)$ jointly, with the same $\mathfrak{N}^{\mathbb{C}}$.

\smallskip\noindent\emph{Scope.} At $(L)$ the backgrounds read are
$\mathfrak{b}\in\{\mathfrak{b}_2,\mathfrak{b}_3,\mathfrak{b}_4\}$ and the pairs
$(\mathfrak{b}_2,\mathfrak{b}_3)$, $(\mathfrak{b}_3,\mathfrak{b}_4)$; the amplitude $\mathsf{t}$
rides on the pair $(\mathfrak{b}_3,\mathfrak{b}_4)$ alone for a plain unit of the first part with
$K_v\not\subset Z$, on each of the four brackets for a plain unit of the second part, and on no
bracket of an interior or wrapped unit; consumers all of whose live law sites are of the outer class
are read also at $\mathfrak{b}_0$, $\mathfrak{b}_1$.
\end{proposition}

\begin{proof}
\emph{Site $(T)$.} By \eqref{5.2:eq-three-sites-2},
$(1+\mathsf{w}^{\mathbb{C}})K_Rg_kT^{\mathbb{C}}=\alpha_{*,k}$, so $\mathsf{w}^{\mathbb{C}}$ is the
largest $\mathsf{w}$ for which the hypothesis of \eqref{5.2:eq-response-containment-a} holds on
$\sup|\mathsf{B}|<T^{\mathbb{C}}$ at every $|\mathsf{t}|\le\mathsf{w}/\eta_p$. For Ba{\l}aban's
radius, $K_R\ge16B_0C_1e^{16\kappa_1}/c_2$ and $\alpha_2\ge c_2\alpha_*$ give
\begin{equation*}
\mathsf{w}_T^{(1.22)}K_R\delta_k=\frac{\alpha_2K_R}{4B_0C_1e^{16\kappa_1}}
\ge\frac{4\alpha_2}{c_2}\ge4\alpha_* .
\end{equation*}
Since $T^{\mathbb{C}}\ge\mathfrak{b}_k=g_k^{-1}\delta_k$, the left member
$(1+\eta_p|\mathsf{t}|)K_Rg_k\sup|\mathsf{B}|$ of that hypothesis, at
$|\mathsf{t}|=\mathsf{w}_T^{(1.22)}/\eta_p$ and $\sup|\mathsf{B}|$ close to $T^{\mathbb{C}}$, comes
arbitrarily close to a number at least $\mathsf{w}_T^{(1.22)}K_R\delta_k\ge4\alpha_*$, with $\alpha_*$
read at level $k$; this exceeds $\alpha_{*,k}$, so the hypothesis fails on the disc of radius
$\mathsf{w}_T^{(1.22)}/\eta_p$.
The type-four units at $(T)$ are read as orbit data exactly as at $(L)$ below and bounded by clause
(c) of hypothesis (i) through the complex-radius corollary and the riding-rotation lemma, with the
disc radius divided by $\eta_p^{1/2}$; the price of this is stated at the end of the discussion of
the disc below. At $(P)$ the earlier chain terms and the frozen terms are read in the same way
through clauses (b), (c) of hypothesis (i) and the polydisc corollary, on the polydiscs of the
preparatory expansion; no condition of that expansion is added or changed.

\smallskip\noindent\emph{Site $(L)$: the undilated chains.} By hypothesis (ii),
$\mathbb{C}^{(n)}_k(X,\mathcal{S})$ is a stored term with slot datum on
$\mathfrak{D}^{\flat}_n(X)\times\mathfrak{D}_{\mathcal{S}}(X)$ and norm at most $C_0^\sharp$, and by
the corollary on the inherited domain $\mathfrak{D}^{\flat}_n(X)$ contains every pair which extends
to $X^+$ under the fat clauses, obeys $\mathfrak{L}_n$ on $X\cap Z^{\mathrm{on}}$, whose fifth and
eighth entries at $\mathfrak{T}$-locations obey the clauses of $\tilde U^c_k(Y_4)$, and whose other
entries there exceed the thresholds of $\tilde U^c_k(Y_4)$ by less than $\mathrm{room}^\flat(y)$.
The chain-containment claim checks these conditions on the chain configurations, which gives
($\alpha$)--($\gamma$) for the undilated chain, the parameters $\sigma$ and the segments
$t\in[0,1]$ included; the containment clause of hypothesis (iii) covers $|\sigma|\le e^{\kappa_1}$
and the segments in $t$, and no amplitude disc. On the amplitude disc the containment is the
content of the containment step below. The pair placed in $\mathfrak{D}^{\flat}_n(X)$ is the carrier
pair in its presentation $h$, restricted to $X$.

\smallskip\noindent\emph{The orbit datum.} Ba{\l}aban's term sits at the carrier regauged by the
product $Q_{\mathfrak{b}}$ of the gauge transformations of the step, which is not the identity:
``Notice that we have to keep the gauge transformation $u_k$ in the above configuration. We cannot
remove it'' \cite[\S1]{B-CMP122b}. A type-four unit is invariant under $G$-valued maps only jointly
with its slots \cite[(2.41)(iii)]{B-CMP119}; hence, writing $\mathbf{A}^{[q]}$ for the slots
transformed by $q$, $U''_k$ for the configuration of \cite[(1.49)]{B-CMP122b} and $Q_4$ for the
regauging $Q_{\mathfrak{b}}$ at $\mathfrak{b}=\mathfrak{b}_4$: at a product datum with the slot held
real and unrotated, the last link is
$f_v((U''_k)^{Q_4};\mathbf{A})=f_v(U''_k;\mathbf{A}^{[Q_4^{-1}]})$, which is not
$\mathbf{B}(U''_k,\mathbf{A})$ of \cite[(1.49)]{B-CMP122b}, and the slot which restores Ba{\l}aban's
identity, $\mathbf{A}^{[\mathfrak{Q}^{\mathfrak{b}}(\zeta)]}$, is complex off the real slice. The
unit is therefore read as \eqref{5.2:eq-three-sites-3}: the orbit datum of the definition in
hypothesis (i), with clauses (b), (c) for the chain terms, clauses (a), (g) for $\mathbf{B}^{(j)}$,
$\mathbf{B}'^{(j)}$ of either part, and, for the members of $\mathfrak{F}(\mathfrak{a})$, hypothesis
(iv) together with clauses (a), (g) of hypothesis (i), which cover every boundary-type unit of
$\mathcal{C}$, and the complex
induction hypothesis at level $k$ (which also covers the older chain terms), taken at the presented
carrier pair $\mathsf{R}_p$, holomorphic in $(\mathsf{t},\sigma,\mathbf{U},\mathbf{J})$ by
Definition~\ref{5.2:def-sites-units} and the containment clause of hypothesis (iii), and at the
rotation $\mathfrak{Q}^{\mathfrak{b}}(\zeta,\mathsf{t},\sigma)$, the holomorphic
site map of the interpolated configuration. By the real-slice property it equals Ba{\l}aban's piece
$\mathfrak{E}(X;\mathfrak{b}^{\mathrm{pr}},\mathcal{J}(\mathfrak{b}^{\mathrm{pr}});\mathbf{A})$ on
the real slice; by the holomorphy property it is holomorphic in
$(\mathsf{t},\sigma,(\mathbf{U},\mathbf{J}),w)$ jointly with the weak slot coordinates at fixed real
strong and silent parts; by clause (c) it obeys
$|f_v\circ\mathsf{R}_p-f_v(1)|\le2\mathfrak{N}^{\mathbb{C}}_{\varepsilon_k}(v)=\mathfrak{N}(v)$ on the
image of $\mathsf{R}_p$ times $\mathfrak{G}_{\varepsilon_k}(X)$. This is the size of type-four units in
\eqref{5.2:eq-unit-table-a} with $\mathfrak{N}^{\mathbb{C}}\ge\mathfrak{N}$ in place of the real-slice
supremum; the constants $C_0^\sharp$, $B_0$, $e^{\lambda_\theta}c_1^\flat$ are those of clauses
(a)--(c). A unit carrying no datum and no amplitude (a wrapped unit or a plain interior unit) is
read at $\mathsf{t}=0$ at its decoupled chain. The units of types one to three and five and those of
the entry marked $\partial$ carry no slot: such a unit depends only on the restriction of $U_k$ to
$X$ (``it depends on $U_k$ restricted to $X$'', \cite[(2.27)(i)]{B-CMP119}).

\smallskip\noindent\emph{The dilated piece.} Since
$\mathfrak{Q}^{\mathfrak{b}}$ depends on the configuration, it depends on $\mathsf{t}$; we bound the
polar size of its $\mathsf{t}$-variation at each live law site $s$ of $v$ (a point at which a live
slot of $v$ is read). Let $\delta\mathbb{H}_p:=\hat\chi_\square\mathbb{H}_i$, $\mathbb{H}_i$ the
layer of the $i$-th bracket, be the piece which
$\mathsf{R}_p(\mathsf{t},\cdot)$ multiplies by $\mathsf{t}$: it is linear in $\mathsf{t}$, the
interpolation being Ba{\l}aban's substitution ``we substitute
$A=H_j(B(t))+t_\square\langle\cdots\rangle$'' \cite[\S3]{B-CMP109}, the block
function $\hat\chi_\square$ multiplying the response as in \cite[(3.4), (3.6)]{B-CMP109}, and
$|\hat\chi_\square|\le1$. Put
\begin{equation*}
\xi_p(s):=\Theta_0^{\mathrm{loc}}(\delta\mathbb{H}_p)(s)=\mathsf{U}_L\,e_p(s),\qquad
\mathsf{U}_L:=C_LB_H\sum_{\mathrm{arg}}\sup|\mathrm{arg}|,
\end{equation*}
$\mathrm{arg}$ running over the arguments of the presented layers and
$\mathsf{U}_L$ being the majorant of all presented layers: their total datum size, with the
constant $B_H$ of the layer-datum bound (into which the composition constant $C_b\ge1$ of
\cite[(3.9)]{B-CMP109} with that bound is absorbed) and the constant $C_L\ge1$ of the second-part
localisation lemma, which majorises the field $f_1$ there,
$|f_1|\le C\sup|\hat{\mathcal{J}}|e^{-\delta_Gd_{\mathrm{sc}}(\cdot,\mathsf{S})}$, $C\le C_L$. Since
$K_uB_H\sum\sup|\mathrm{arg}|\le\varepsilon_{\mathfrak{Q}}(k)\le K_{\mathfrak{Q}}\alpha^+_k$ by the
rotation-size bound at level $k$, in which $\varepsilon_{\mathfrak{Q}}(k)$ is the bound on the size
of the presented rotations at level $k$, one has $K_u\mathsf{U}_L\le C_LK_{\mathfrak{Q}}
\alpha^+_k$.

\smallskip\noindent\emph{Two bounds on the profile.} (a) $\xi_p(s)\le\eta_p\mathsf{U}_L$ for $s\in
S_v$: this is the definition of $\eta_p$ (the piece seen on $S_v$ has size at most $\eta_p$ times the
unit shift $\mathsf{U}_J$), and $\mathsf{U}_J\le\mathsf{U}_L$; the live law sites of the slots of $v$
lie in $S_v=X$, and the top block $B^{\mathrm{top}}(s)$ holding $s$ lies in $S_v$.
(b) $\xi_p(s)\le\mathsf{U}_Le^{-\delta_\natural(d_{\mathrm{sc}}(s,\mathfrak{F}_{\mathfrak{s}}(s))-1)}$,
where $\mathfrak{F}_{\mathfrak{s}}(s)=\Lambda^{\sim}$ for $s$ of the outer class ($s\notin Z$) and
$\mathfrak{F}_{\mathfrak{s}}(s)=\mathsf{S}$ for $s$ of the inner class
($s\in Z_{k-N_k}\subset\Lambda$, $\Lambda$ being the region of the operation $\mathbf{R}$ of
\cite[\S1]{B-CMP122b}).
This follows from the two whole-layer lines of the layer-reading lemma, applied to one
$\square$-summand. \emph{Outer class.} The supported arguments of the four brackets have their
supports in $\Lambda^{\sim}$. For such an argument $\hat\chi_\square\mathbb{H}(\mathrm{arg})$ decays
from $\square$ by \cite[(3.9)]{B-CMP109}, and its datum at $\square$ carries
$e^{-\delta_1^{\mathrm{sc}}d_{\mathrm{sc}}(\square,\operatorname{supp}\mathrm{arg})}$ by the
layer-datum bound; with $d_{\mathrm{sc}}(s,\square)+d_{\mathrm{sc}}(\square,\Lambda^{\sim})\ge
d_{\mathrm{sc}}(s,\Lambda^{\sim})$ and $-1$ for the passage to the top block,
\begin{equation*}
\xi_p(s)\le B_H\sum\sup|\mathrm{arg}|\,e^{-\delta_1^{\mathrm{sc}}(d_{\mathrm{sc}}(s,\Lambda^{\sim})-1)}
\le\mathsf{U}_Le^{-\delta_\natural(d_{\mathrm{sc}}(s,\Lambda^{\sim})-1)},
\end{equation*}
as $\delta_\natural\le\delta_1^{\mathrm{sc}}$ and $C_L\ge1$. The one argument which is not supported,
the field $G''_k\theta J_{1,2}$ of a second-part unit (``The field $J_{1,2}$ is not localized'',
\cite[\S1]{B-CMP122b}), is treated by the resummation of the inner class below, carried out at the
outer site with $d_{\mathrm{sc}}(s,\mathsf{S})\ge d_{\mathrm{sc}}(s,\Lambda^{\sim})$ ($\mathsf{S}\subset
\Lambda^{\sim}$), rate $\delta_L\ge\delta_\natural$ and constant $C_L$ inside $\mathsf{U}_L$.
\emph{Inner class.} The only unit carrying an amplitude and an inner site is a plain first-part unit
with $K_v\not\subset Z$; its one amplitude bracket is the fourth, the pair
$(\mathfrak{b}_3,\mathfrak{b}_4)$, so that $\delta\mathbb{H}_p$ is a $\square$-piece of
$\mathbb{H}_4=H'_k(g_kB)$, of amplitude $\mathsf{U}_J$, and the third bracket's layer rides undilated
inside $a_0$ below; the line is given for the third and fourth brackets because $a_0$ needs both.
The first two brackets have their weights identically $1$ on $\Lambda^{\sim}\ni s$, whatever the
configuration, so their pieces do not rotate $s$. At the third and fourth brackets the second
argument of $\mathbb{H}_3$ and the argument $g_kB$ of $\mathbb{H}_4$ are supported in $\mathsf{S}$,
and the first argument of $\mathbb{H}_3$ is the field $f_1$ above (``The field $J_{1,2}$ is not
localized, but in the composition $G''_kJ_{1,2}$ we may localize it'', \cite[\S1]{B-CMP122b}).
Cutting $f_1$ into $M$-cubes $y$, the layer-datum bound and \cite[(3.9)]{B-CMP109} give
\begin{align*}
&\sum_y e^{-\delta_1^{\mathrm{sc}}[d_{\mathrm{sc}}(s,\square)+d_{\mathrm{sc}}(\square,y)]
-\delta_Gd_{\mathrm{sc}}(y,\mathsf{S})}\\
&\qquad\le\sum_ye^{-\delta_1^{\mathrm{sc}}d_{\mathrm{sc}}(s,y)-\delta_Gd_{\mathrm{sc}}(y,\mathsf{S})}
\le C_\rho(\delta_L)e^{-\delta_Ld_{\mathrm{sc}}(s,\mathsf{S})},
\end{align*}
the last step being the resummation of the second-part localisation lemma, which uses
$\min\{\delta_1^{\mathrm{sc}},\delta_G\}=2\delta_L$ and $d_{\mathrm{sc}}(s,y)+d_{\mathrm{sc}}(y,
\mathsf{S})\ge\max\{d_{\mathrm{sc}}(s,\mathsf{S}),d_{\mathrm{sc}}(s,y)\}$, with
$C_L:=C_QC_\rho(\delta_L)C$, where $C_\rho(\delta_L)$ is the constant of that resummation and $C_Q$
a constant of that lemma. Hence, as $\delta_\natural\le\delta_L$,
\begin{equation*}
\xi_p(s)\le C_LB_H\sum\sup|\mathrm{arg}|\,e^{-\delta_L(d_{\mathrm{sc}}(s,\mathsf{S})-1)}
\le\mathsf{U}_Le^{-\delta_\natural(d_{\mathrm{sc}}(s,\mathsf{S})-1)}.
\end{equation*}
In both classes the majorant of the $\square$-piece is one summand of the sum over $\square$ which
gives the whole layer's majorant at $s$; only the lines of the layer-reading lemma and
\cite[(3.9)]{B-CMP109} are used. Both bounds hold in the four-entry cell norm $\mathsf{N}_y$ of
\cite[(3.14)]{B-CMP109}, at every cell $y$ of $S_v$ and every current block holding a site of $S_v$:
for the $U$-entries by the layer-response theorem read at the weight
$e^{-\delta d(\cdot,\square)}$, respectively $e^{-\delta d^{\wedge}(\cdot,\mathfrak{F}_{\mathfrak{s}})}$
(at $(T)$, $\mathfrak{F}_{\mathfrak{s}}=\Omega_k$, the small-field region), $\delta$ being the decay
rate of the weights of that theorem and its constant $B_\sharp$ absorbed into $B_H$; for the
$\mathbf{J}$-entries at $(T)$ by the same line, the layer's datum being supported in $\square$ and
read at the weight $e^{-\delta d(\cdot,\square)}$, the response itself being non-local at $(T)$,
\cite[(3.9)]{B-CMP109}; at $(L)$ by the pointwise product of the
$\mathbf{J}$-recession lemma at touching and receding blocks. Writing $N^\Delta_y$, $N_y$ for the two
$U$-entries of $\mathsf{N}_y$, $K_\chi$ for the constant of the $\mathbf{J}$-recession lemma attached
to the block functions $\hat\chi_\square$, and $p_i(y)$, $\mathsf{a}_i$ for the weights of its
pointwise product for the layer $\mathbb{H}_i$,
\begin{equation*}
\begin{gathered}
N_y(\hat\chi_\square\mathbb{H}_i)\le20B_\sharp\,q_\square(y),\qquad
\bigl(\text{$\mathsf{t}$-member of the third entry}\bigr)\le20K_\chi|\mathsf{t}|\,q_\square(y),\\
q_\square(y)\le p_i(y)\,e^{-\delta_Pd(y,\square^{\zeta})}\le p_i(y)\,\eta_p\quad\text{on }S_v,
\end{gathered}
\end{equation*}
in the unit $20K_\chi\mathsf{a}_i$ of the third entry: the far factor at $y$ times the block decay
to $y$, with no junction of rates; the last inequality holds for $y\in S_v$ because at site $(L)$
$\eta_p=e^{-\delta_Pd(\square^{\zeta},S_v)}$ and $d(y,\square^{\zeta})\ge d(\square^{\zeta},S_v)$.

\smallskip\noindent\emph{The disc and the Schwarz bound.} On $|\mathsf{t}|\le\rho_A$, by (a), (b)
and $\min\{a,b\}\le(ab)^{1/2}$,
\begin{align*}
|\mathsf{t}|\xi_p(s)&\le\mathsf{w}_L\eta_p^{-1/2}\mathsf{U}_L\min\bigl\{\eta_p,
e^{-\delta_\natural(d_{\mathrm{sc}}(s,\mathfrak{F}_{\mathfrak{s}}(s))-1)}\bigr\}\\
&\le\mathsf{w}_L\mathsf{U}_Le^{-\frac12\delta_\natural(d_{\mathrm{sc}}(s,\mathfrak{F}_{\mathfrak{s}}(s))-1)}:
\end{align*}
the dilated piece decays at half rate at every live law site, in either class. Let
$Z_s(\mathsf{t}):=\log[\mathfrak{Q}^{\mathfrak{b}}(\zeta,0,\sigma)^{-1}
\mathfrak{Q}^{\mathfrak{b}}(\zeta,\mathsf{t},\sigma)](s)$, holomorphic in $\mathsf{t}$ with
$Z_s(0)=0$. By the local gauge-factor lemma, $\mathfrak{u}(s)$ reads the layer on
$B^{\mathrm{top}}(s)$ alone, holomorphically on the ball of radius
$\varrho_0:=\min\{10^{-2}/K_u,c_4\}$, with $|\log\mathfrak{u}(s)|\le K_u
\Theta_0^{\mathrm{loc}}(\text{layer})(s)$. Let
\begin{equation*}
a_0(s):=\Theta_0^{\mathrm{loc}}(\text{undilated layers})(s)
\le\mathsf{U}_Le^{-\delta_\natural(d_{\mathrm{sc}}(s,\mathfrak{F}_{\mathfrak{s}}(s))-1)}
\end{equation*}
(the whole-layer line of which (b) is the form uniform in $\square$; in the inner class $a_0$ is the
third and fourth brackets' layers only) and $\xi:=\xi_p(s)$. By the block-local Schwarz bound of
the composition lemma for layers (the lemma which also gives the layer-datum bound), $Z_s$ is
holomorphic on $|\mathsf{t}|\xi<\varrho_0-a_0$, with
$|Z_s|\le1.05K_u(\varrho_0+a_0)$ there, so by the Schwarz lemma
\begin{equation*}
|Z_s(\mathsf{t})|\le1.05K_u\frac{\varrho_0+a_0}{\varrho_0-a_0}|\mathsf{t}|\xi
\le2.1K_u|\mathsf{t}|\xi_p(s)
\end{equation*}
provided $a_0\le\varrho_0/3$ (then the quotient is at most $2$) and $a_0+|\mathsf{t}|\xi\le\varrho_0$.
Both follow from the second line of \eqref{5.2:eq-three-sites-b}: the bracket condition
$4C_A'\theta'(g_k)\le1$ of hypothesis (vii) gives $\mathsf{w}_L\ge1$, hence, with the distances below,
\begin{equation*}
K_ua_0,\ K_u|\mathsf{t}|\xi\le K_u\mathsf{w}_L\mathsf{U}_Le^{-\frac12\delta_\natural(\check R_k-2)}
\le\frac{\mathsf{a}_J(x_k)}{2.1}E(x_k)^{\delta_\natural/2}\le\frac{K_u\varrho_0}{6.3}.
\end{equation*}

\smallskip\noindent\emph{The two unit shifts.} $\mathsf{w}_L=1/\theta'(g_k)$ is the room
$\tilde\alpha_1(k)$ of the tube $\tilde U^c_k(Y_4,\tilde\alpha_0,\tilde\alpha_1)$ over the unit shift
$\mathsf{U}_J:=6B_H\delta'_k$ of the fluctuation layer alone (whose argument is at most
$6\delta'_k$): $\mathsf{w}_L\mathsf{U}_J\le\tilde\alpha_1(k)\le\alpha^+_k$ is the defining ratio of
$\theta'$, and it concerns that layer only. The three $\Lambda$-functionals have data of the order of
the tube radius, so $\mathsf{U}_L$, which majorises them, is proportional to $\alpha$, and
$\mathsf{w}_L\mathsf{U}_L$ is a polylogarithm times $\alpha$: from
$K_u\mathsf{U}_L\le C_LK_{\mathfrak{Q}}\alpha^+_k\le\frac12e^{-2}C_\natural\bar\alpha(x_k)$, where
$C_\natural:=2e^2\max\{C_LK_{\mathfrak{Q}},K_uB_H,K_u\}$, and
$\mathsf{w}_L\le C_{\theta'}^{-1}(\log x_k+2)^{q-p_1}$ ($\log x_k+2\ge\mathsf{T}_k+1$),
\begin{equation*}
2.1K_u\mathsf{w}_L\mathsf{U}_L\le\mathsf{a}_J(x_k),\qquad
|Z_s(\mathsf{t})|\le\vartheta(s):=\mathsf{a}_J(x_k)\,
e^{-\frac12\delta_\natural(d_{\mathrm{sc}}(s,\mathfrak{F}_{\mathfrak{s}}(s))-1)} .
\end{equation*}

\smallskip\noindent\emph{Distances, and $\vartheta\le\mu$.} By the distance lemma for live law
sites: in the outer class $d_{\mathrm{sc}}(s,\Lambda^{\sim})\ge5\check R_k$, increased by
$5R_{\min}(k-1-k_W)_+$ at a site of a frame; in the inner class $d_{\mathrm{sc}}(s,\mathsf{S})\ge
\check R_k-1$, increased by $5R_{\min}(k-N_k-k_W)_+$ at a site of a frame. As $k-1\ge k-N_k$,
both increments are at least $5R_{\min}\mathsf{a}_k(s)$, $\mathsf{a}_k(s):=(k-N_k-k_W)_+$. The law
sites of frames and of arrested variables inside $Z$ are of the inner class
($\Sigma_W\subset W\subset Z_{k_W}\subset Z_{k-N_k}$, $W_{\mathfrak{a}}\subset Z_{k-N_k}$); those of a
frame outside $Z$ read by a presented term are of the outer class, with the numbers just given
($5\check R_k-1\ge\check R_k-2$). So $d_{\mathrm{sc}}(s,\mathfrak{F}_{\mathfrak{s}}(s))-1\ge
\check R_k-2$ at every live law site, plus $5R_{\min}\mathsf{a}_k(s)$ at a frame's. Since
$(\log(x_k-c_5))^{r_{\mathrm{pr}}}\le\check R_k$ by hypothesis (vii), the relation between
$\check R_k$ and the decay function gives
$e^{-\frac12\delta_\natural(\check R_k-2)}\le E(x_k)^{\delta_\natural/2}$, and
$5\delta_\natural R_{\min}\ge2$ (hypothesis (vii)) gives
$e^{-\frac12\delta_\natural5R_{\min}\mathsf{a}}\le e^{-\mathsf{a}}$; hence
\begin{equation*}
\vartheta(s)\le\mathsf{a}_J(x_k)E(x_k)^{\delta_\natural/2}\ \text{(pointwise site)},\qquad
\vartheta(s)\le\mathsf{a}_J(x_k)E(x_k)^{\delta_\natural/2}e^{-\mathsf{a}_k(s)}\
\text{($s\in\Sigma_W$)}.
\end{equation*}
Let $\mu(s):=\mu_k\mathsf{e}(s)/\varepsilon_\star$, the margin reserved for the Cauchy estimate in the
rotation argument $\mathfrak{o}$ on the chart which the presented-bracket corollary subtracts for
this presentation, with
$\mu_k=\varepsilon_\star(x_k-c_5)^{-2}$, $\mathsf{e}(s)=\varepsilon_\star$ at a pointwise site and
$\mathsf{e}(s)=\varpi_{\mathrm{fr}}\alpha_{*,k_W}/4C_{\mathrm{fr}}$ at a frame's. At a pointwise
site
\begin{equation*}
\frac{\vartheta(s)}{\mu(s)}\le\frac{\mathsf{a}_J(x_k)E(x_k)^{\delta_\natural/2}(x_k-c_5)^2}
{\varepsilon_\star}\le1
\end{equation*}
by the first line of \eqref{5.2:eq-three-sites-b}. At $s\in\Sigma_W$,
\begin{equation*}
\frac{\vartheta(s)}{\mu(s)}\le\mathsf{a}_J(x_k)E(x_k)^{\delta_\natural/2}e^{-\mathsf{a}_k(s)}
(x_k-c_5)^2\,\frac{4C_{\mathrm{fr}}}{\varpi_{\mathrm{fr}}\alpha_{*,k_W}} .
\end{equation*}
By the coupling-ratio bound of the margin theorem with its level parameter at $k_W$, with the upper
half of the reference
flow bound and $N_k\le x_k$ (hypothesis (vii)), $q_*$ being the exponent of the logarithm in
$\alpha_*$,
\begin{align*}
\frac{\alpha_{*,k}}{\alpha_{*,k_W}}
&=\Bigl(\frac{x_{k_W}}{x_k}\Bigr)^{1/2}\Bigl(\frac{\log x_k}{\log x_{k_W}}\Bigr)^{q_*}
\le2\Bigl(1+\frac{B^\sharp(k-k_W)}{x_k}\Bigr)^{1/2}\\
&\le2(1+B^\sharp)^{1/2}(1+\mathsf{a}_k(s))^{1/2},
\end{align*}
using $k-k_W\le N_k+\mathsf{a}_k(s)$; and $(1+\mathsf{a})^{1/2}e^{-\mathsf{a}}\le1$ for
$\mathsf{a}\ge0$. So
\begin{equation*}
\frac{\vartheta(s)}{\mu(s)}\le\mathsf{a}_J(x_k)E(x_k)^{\delta_\natural/2}(x_k-c_5)^2\,
\frac{8C_{\mathrm{fr}}(1+B^\sharp)^{1/2}}{\varpi_{\mathrm{fr}}\alpha_*(x_k)}\le1
\end{equation*}
by the first line of
\eqref{5.2:eq-three-sites-b}: the current coupling against the frame's chart size is compensated by
the shell factor contained in $d_{\mathrm{sc}}$. The second line of \eqref{5.2:eq-three-sites-b} also
gives $\vartheta(s)\le10^{-2}/3$, the size hypothesis $\vartheta\le10^{-2}$ of the riding-rotation
lemma.

\smallskip\noindent\emph{The riding-rotation lemma at $\rho_A$.} Let $\mathfrak{M}$ be the polydisc
of the holomorphic coordinates $(\sigma,\mathbf{U},\mathbf{J},\mathbf{\Phi},\zeta,w)$, $\iota$ a
point of it, $\mathfrak{o}$ the rotation argument, and
$f(\mathsf{t},\iota,\mathfrak{o}):=c^\sharp_v(S_v;\mathsf{R}_p(\mathsf{t},\sigma;\mathbf{U},
\mathbf{J});\mathbf{A},\mathbf{\Phi},\mathfrak{o};w)-f_v(1)$, holomorphic on
$\{|\mathsf{t}|<\rho_A\}\times\mathfrak{M}\times\mathfrak{G}_{\varepsilon_k(s)}(S_v)$ with
$|f|\le\mathfrak{N}(v)$ (the orbit datum above; $\mathsf{R}_p\subset\mathfrak{D}_v$ on the disc by the
containment step below). Take $u(\mathsf{t},\iota):=\mathfrak{Q}^{\mathfrak{b}}(\zeta,\mathsf{t},
\sigma)$; $\eta_0(s)$ the polar size of $\mathfrak{Q}^{\mathfrak{b}}(\zeta,0,\sigma)(s)$, read as its
supremum over $\mathfrak{M}$ and over $t\in[0,1]$ (the undilated rotation, part of the rotation
allowance $\eta^L_k(s)$); $\vartheta(s)$ as above, $|\log[u(0)(s)^{-1}u(\mathsf{t})(s)]|=
|Z_s(\mathsf{t})|\le\vartheta(s)\le10^{-2}$; and $\mu(s)$ as above. Then
$\varepsilon'(s):=\varepsilon_k(s)-\eta_0(s)-\mu(s)>0$ by the summed-margin bound of hypothesis (vii)
($\sum_{k'}\eta_{k'}(s)\le\frac14\mathsf{e}(s)$, hence $\frac34\mathsf{e}\le\varepsilon_k\le
\mathsf{e}$), and $\mathfrak{G}_{\varepsilon'}(S_v)$ is the chart on which the presented-bracket
corollary reads this presentation ($\eta_0\le\eta^L$). The lemma gives: $F(\mathsf{t},\iota,
\mathfrak{o}'):=f(\mathsf{t},\iota,u(\mathsf{t},\iota)\mathfrak{o}')$ is holomorphic on
$\{|\mathsf{t}|<\rho'\}\times\mathfrak{M}\times\mathfrak{G}_{\varepsilon'}(S_v)$ with
$|F|\le\mathfrak{N}$, where $\rho':=\rho_A\min\{1,\min_s\mu(s)/\vartheta(s)\}$, and
$|\partial_{\mathsf{t}}F(0,\iota,\mathfrak{o}')|\le\mathfrak{N}/\rho_A$ times the loss factor
$\max\{1,\max_s\vartheta(s)/\mu(s)\}$ of the lemma. Since $\vartheta\le\mu$ at every live law site,
the loss factor equals $1$ and $\rho'=\rho_A$: the orbit datum read through the dilated rotation is
holomorphic on the full disc over $\mathfrak{G}_{\varepsilon'}$ and bounded by $\mathfrak{N}(v)$, and
$|\partial_{\mathsf{t}}F(0)|\le\mathfrak{N}\eta_p^{1/2}/\mathsf{w}_L$. The chart radius is reduced
by $\eta_0+\mu$, both already subtracted for this presentation, and the dilated rotation enters only
the bound, through the loss factor $1$; no further margin is subtracted, and the margin theorem
applies unchanged: $\eta_{k'}(s)\le4\mu_{k'}\mathsf{e}(s)/\varepsilon_\star$ per step and
\begin{equation*}
\sum_{k'}\frac{\eta_{k'}(s)}{\mathsf{e}(s)}\le4\bigl[y^{-2}+(\lambda^\sharp y)^{-1}\bigr]+\frac1{16}
\le\frac4{32}+\frac1{16}=\frac3{16}<\frac14
\end{equation*}
for every $K$, $y$ here denoting the reference variable of the margin theorem (not a cell), at least
its lower bound $y_\flat$. The margin $\mu$ of
a presentation is the gap between the charts $\varepsilon_k-\eta_0$ and $\varepsilon'$; the disc here
uses at most $\mu$.

\smallskip\noindent\emph{The piece.} By the Cauchy estimate at radius $\rho_A$, with the constant
$2$ for closed discs, $|\partial_{\mathsf{t}}F(0,\iota,\mathfrak{o}')|\le2\mathfrak{N}(v)
\eta_p^{1/2}/\mathsf{w}_L$ uniformly for $|\sigma(\Delta)|\le e^{\kappa_1}$; the lemma bounding the
pieces $v_p$ of a unit,
whose hypotheses $\mathsf{R}_p\subset\mathfrak{D}_v$ and $|f_v\circ\mathsf{R}_p-f_v(1)|\le
\mathfrak{N}(v)$ on the disc are established here, turns this into
\eqref{5.2:eq-three-sites-a}. The price is $\eta_p\mapsto\eta_p^{1/2}$, paid where $\eta_p$ is
summed: in the resummation over $\square$ ($C_A\mapsto C_A'$; at $(T)$, $(P)$ the decay from $\square$
to $S_v$ at rate $\frac12\delta_1^{116}M\ge\frac12\kappa_1\ge8\kappa$ per $M$-cube instead of
$\delta_1^{116}M$;
at $(L)$, where $\eta_p:=e^{-\delta_Pd(\square^{\zeta},S_v)}$ equals $1$ at the blocks whose
$\square^{\zeta}$ meets the $2$-stencil of $S_v$, by
$\sum_\square\eta_p^{1/2}\le C_H^{\wedge\prime}:=(2L)^d(1+\hat b_t+\hat b_tC_{\mathrm{cell}})$, which
is the block-sum clause of the $\mathbf{J}$-recession lemma at rate $\frac12\delta_P$ on
$\square^{\zeta}\subset\square$, the touching blocks being counted inside
$\#\Sigma_\square\le(2L)^d$), a $(d,L,M,\kappa_1)$-number fixed before $\gamma$, the output rate and the
walk factor $e^{-(\kappa_1-1)m(p)}$ untouched; and in the second step of the complex-radius lemma,
where $e^{-\delta_0cMR_k}$, $\delta_0$ and $c$ being the rate and the constant of that step, becomes
$e^{-\frac12\delta_0cMR_k}$, still tending to zero. At $(L)$ the
type-four units among the plain units carrying a localisation datum in the decoupling convention are
resummed exactly over the blocks at $\mathsf{t}=0$ by the summation theorem for type-four units,
which uses the containment $\mathsf{R}_p\subset\mathfrak{D}_v$ and the bound
$|f_v\circ\mathsf{R}_p-f_v(1)|\le\mathfrak{N}(v)$ at $\mathsf{t}=0$ only, and not the gain
$\mathsf{w}_L^{-1}\eta_p^{1/2}$ of \eqref{5.2:eq-three-sites-a}; there the
disc acts only through the compact pointed data of those plain units (the pointed units of
Definition~\ref{5.4:def-pointed-units}) and through the pointwise product of the
$\mathbf{J}$-recession lemma.

\smallskip\noindent\emph{Containment on the disc.} It is checked at each clause location of
$\mathfrak{D}_v$: at the $\mathfrak{T}$-locations here; at the remaining clause locations by
hypothesis (viii), under
\eqref{5.2:eq-half-disc-containment-c}--\eqref{5.2:eq-half-disc-containment-e}. The line used at all
of them: on
$|\mathsf{t}|\le\rho_A$, at $y\in S_v$, with $\mathsf{U}\in\{\mathsf{U}_J,\mathsf{U}_L\}$ the
piece's amplitude and $\mathsf{f}(y)\le1$ the far factor of the undilated layer at $y$,
\begin{equation*}
|\mathsf{t}|N_y(\delta\mathbb{H}_p)\le\mathsf{w}_L\eta_p^{-1/2}\mathsf{U}\min\{\eta_p,\mathsf{f}(y)\}
\le\mathsf{w}_L\mathsf{U}\min\{1,\mathsf{f}(y)^{1/2}\},
\end{equation*}
since $\mathsf{f}\eta_p^{-1/2}=\mathsf{f}^{1/2}(\mathsf{f}/\eta_p)^{1/2}\le\mathsf{f}^{1/2}$ if
$\mathsf{f}\le\eta_p$, $\eta_p^{1/2}\le\mathsf{f}^{1/2}$ if $\mathsf{f}\ge\eta_p$, and $\eta_p\le1$.
Under the convention $\eta_p=e^{-\delta_Pd(\square^{\zeta},S_v)}$ at site $(L)$, equal to $1$
exactly at touching blocks, this line is used at touching blocks only; at a receding
block the support $\square^{\zeta}$ of $\hat\chi_\square$ does not meet $S_v$, so
$\hat\chi_\square\mathbb{H}_i$ vanishes on $S_v$, and what moves there is the
$\mathbf{J}$-tail $-\mathsf{t}T_p$, with $K_P$ the constant of the $\mathbf{J}$-recession lemma,
$\ell$ the linear size of the cell $y$, and
\begin{equation*}
\ell^3|T_p|\le20K_PB_\sharp p_i(y)e^{-\delta_Pd(y,\square^{\zeta})}\le20K_PB_\sharp p_i(y)\eta_p,
\end{equation*}
so $|\mathsf{t}|\ell^3|T_p|\le20K_PB_\sharp\mathsf{w}_L\eta_p^{1/2}p_i(y)$ with the full far factor,
and the third entry moves by at most $\theta_3/16$, $\theta_3$ its threshold, by the
$\mathbf{J}$-recession lemma under the conditions of hypothesis (vii). For the pointed units of
Definition~\ref{5.4:def-pointed-units}, at
radius $\mathsf{w}_L/(\eta_p\mathsf{f}_v)^{1/2}$, the $U$-entries move by at most
$c_\chi\mathsf{w}_L\mathsf{U}\min\{1,\mathsf{f}(y)^{1/2}\}$ at touching blocks, the extra
$\mathsf{f}_v^{-1/2}$ paid by the far factor, and the $\mathbf{J}$-entry at every block with half the
far factor kept.

At $y\in X\cap\mathfrak{T}$ the increments induced by the dilated piece are those of the
layer-increment proposition with the $\square$-piece in place of the layer. Its majorant has supports
in a set $\mathcal{Z}$ with $d_{\mathrm{sc}}(y,\mathcal{Z})\ge5R_k+10R_{\min}(k-1-m)$ on another
component's shell $m$ (the shell index, not the number $m(p)$), $\ge5R_k$ on $\Omega_k\setminus Z$,
$\ge3R_k$ on the collars; with (a) and
$\min\{a,b\}\le(ab)^{1/2}$ the dilated increment is the undilated one with the prefactor
$\mathsf{w}_L$ and the exponent halved, on the shells
\begin{equation*}
\begin{split}
P^{(1/2)}:={}&\mathsf{w}_LK_WB_HK_\Lambda(M)(1+\varrho_k)(1+\beta_0)^2(k-m)^{1/2}\\
&\times e^{-\frac52\delta_1^{\mathrm{sc}}R_k}e^{-5\delta_1^{\mathrm{sc}}R_{\min}(k-1-m)},
\end{split}
\end{equation*}
$K_W$, $K_\Lambda(M)$, $\varrho_k$, $\beta_0$ being the constants of the layer-increment proposition
and the ratio $\alpha_{\cdot,k}/\alpha_{\cdot,m}\le(1+\beta_0)^2(k-m)^{1/2}$, $m<k$, being
\cite[(2.6), (2.7)]{B-CMP119}, and $P'^{(1/2)}\propto e^{-\frac52\delta_1^{\mathrm{sc}}R_k}$,
$P''^{(1/2)}\propto e^{-\frac32\delta_1^{\mathrm{sc}}R_k}$; the undilated remainder adds
$P\cdot2^{-(k-1-m)}\le P^{(1/2)}$ (as $\mathsf{w}_L\ge1$), so the total is at most $2P^{(1/2)}$, and
likewise for $P'$, $P''$. The shell factor survives the halving: the rate condition
$10\delta_1^{\mathrm{sc}}R_{\min}\ge2$ of hypothesis (vii) gives
\begin{equation*}
e^{-5\delta_1^{\mathrm{sc}}R_{\min}(k-1-m)}\le e^{-(k-1-m)}
=2^{-(k-1-m)}e^{-(1-\ln2)(k-1-m)},\qquad 1-\ln2\ge0.3,
\end{equation*}
and
\begin{equation*}
(k-m)^{1/2}e^{-0.3(k-1-m)}\le C_{1/2}':=\sup_{x\ge0}(2+x)^{1/2}e^{-0.3x};
\end{equation*}
Ba{\l}aban uses
the same device: ``the product of $1+(j-i)^{1/2}$ and the exponential factor in (1.47) can be
estimated by $\exp(-(j-i))$'' \cite[\S1]{B-CMP122a}. So on shell $m$ the increment over the threshold
is at most $2\mathsf{w}_LC_P'e^{-\frac32\delta_1^{\mathrm{sc}}R_k}2^{-(k-1-m)}$, $C_P'=C_PC_{1/2}'/
C_{1/2}$. On strata left thin by an earlier live arrest (``separated by one layer of $M$-cubes \dots\
one layer of $L^{-1}M$-cubes, and so on'', \cite[\S1]{B-CMP122a}) the same line is read per crossed
stratum: the condition $\delta_1M/M_1\ge2$ of hypothesis (vii) (``We choose $M$ large enough, so that
$\delta(M/M_1)\ge2$'', \cite[\S1]{B-CMP122a}) gives $e^{-2}$ per stratum undilated and
$e^{-1}=2^{-1}e^{-(1-\ln2)}$ at half rate, the same two factors. Against
$\frac14\mathrm{room}^\flat(y)\ge\frac1{16}\beta_{\mathrm{pr}}\theta_S2^{-(k-1-m)}$ (respectively
$\frac18\beta_{\mathrm{pr}}$ off the shells), the increment is at most
$\frac14\mathrm{room}^\flat$ as soon as $32\mathsf{w}_LC_P'e^{-\frac32\delta_1^{\mathrm{sc}}R_k}\le
\beta_{\mathrm{pr}}\theta_S$, which follows from \eqref{5.2:eq-three-sites-1} applied at
$\mathsf{T}=\mathsf{T}_k-1\ge\mathsf{T}_\gamma-1$, since
$\mathsf{w}_L\le C_{\theta'}^{-1}(\mathsf{T}_k+1)^{q-p_1}$ and $R_k\ge R(e^{\mathsf{T}_k-1})$; the shell
index enters only through $2^{-(k-1-m)}$ on both sides. The price of the rotation is hypothesis (vi),
read with the conditions $5\delta_\natural R_{\min}\ge2$, $N_k\le x_k$, the upper half of the
reference flow bound and the frame-shell condition of the distance lemma, all in hypothesis (vii); the
live-margin condition and the condition controlling the margin of the complex-radius corollary at
depth $8\check R_k$ stay as they are, the latter not being used at $(L)$.

\smallskip\noindent\emph{Complex sources.} The orbit-datum bound of hypothesis (i) is stated without
complex sources; $w$ enters $f_v$ through
the complex-source corollary under its two conditions, as one more holomorphic polydisc parameter
under $w$-free majorants. The holomorphy property of the orbit datum is a statement at fixed values
of every parameter it does not name, and the riding-rotation lemma passes parameters
($w\in\mathfrak{M}$); so the orbit datum of $\mathbb{C}^{(n)}_k(X,\mathcal{S};w)$ is that of the
family of the complex-source corollary, holomorphic in $w\in\mathcal{W}_k(X)$ jointly, with the
same $\mathfrak{N}^{\mathbb{C}}$.

\smallskip\noindent\emph{Scope.} The backgrounds listed are those of hypothesis (iii); the
brackets carrying $\mathsf{t}$ are fixed by the decoupling convention; for consumers all of whose
live law sites are of the outer class the presentation gauge $h^{\mathfrak{b}}$ equals $1$ within
distance $\frac12$ of those sites,
which allows the reading at $\mathfrak{b}_0$, $\mathfrak{b}_1$ with the clauses of the site map. For
the families of the decoupling convention's source clause the containment and holomorphy clauses of
hypothesis (iii) hold by the source-anchored containment lemma: the decoupling set $\sigma_0$ is a
free parameter of the base configuration on which the containment theorem of hypothesis (iii) is
stated, the decoupled operators reading it through one number, and the
position conditions ask $\sigma_0\cap\Lambda=\emptyset$ of the presented chain's families, which they
obey. Hypothesis (i) is thus used at all three sites.
\end{proof}

\begin{lemma}[Containment on the half disc outside the small-field region]
\label{5.2:lem-half-disc-containment}
Fix $K$ and a level $k\le K$, and work at the site $(L)$ of
Proposition~\ref{5.2:prop-three-sites}, with
$\mathfrak{T}:=\Omega_k\cup Z^c$. By \emph{units with live exterior slots} we mean the terms
$\mathbb{C}^{(n)}_k(X,\mathcal{S};w)$, $\mathbf{B}^{(j)}$, $\mathbf{B}'^{(j)}$ and
the frozen members of $\mathfrak{F}(\mathfrak{a})$. Let $K_D\ge\max\{K_W,3,2K_S^b\}$ be the constant
of the linearised comparison \eqref{5.2:eq-half-disc-containment-b} below, and put
\begin{equation*}
C^J_{\mathrm{age}}:=2\sqrt{2}\,(1+b_+)^{1/2}\sup_{u\ge 0}(1+u)^{1/2}e^{-\frac12\delta_\natural u}.
\end{equation*}
Here $K_W$ is the constant of the layer increments at the $\mathfrak{T}$-locations of
Proposition~\ref{5.2:prop-three-sites}, $K_S^b$ the constant of the linearised stability statement
for composite entries used in the proof, $C^b$ and $a_\flat$ the constants of the stability condition
of that statement, $R_{\min}$ a lower bound for all the radii $R_k$, $\tilde\alpha_{0,j}$ (at level
$j$) and
$\tilde\alpha_1(k)$ the radii of the analyticity tube of the source space at $(L)$, and $C_*$, $q$,
$A_1$, $p_1$ the constants and exponents in the forms $\alpha_*=C_*g\mathsf{T}^{q}$ and
$\delta'=gA_1\mathsf{T}^{p_1}$ of the radius and of the fluctuation bound.
Assume the following.
\begin{enumerate}
\item[(1)] $5\delta_\natural R_{\min}\ge 2$.
\item[(2)] $N_k\le x_k$.
\item[(3)] The comparison of the running couplings with the reference recursion, in both its
halves, holds for all pairs of levels not exceeding $k$.
\item[(4)] $\delta_P w_{\mathrm{sc}}\ge 1\ge\delta_\natural$, where $w_{\mathrm{sc}}$ is the factor in
the comparison $d^{\wedge}\ge w_{\mathrm{sc}}\,d_{\mathrm{sc}}$ of the scaled distances.
\item[(5)] The linearised stability condition at the largest increment majorant:
\begin{equation}\label{5.2:eq-half-disc-containment-f}
\sup_{x\ge\gamma^{-2}}C^b\bigl(a_\flat+c_\chi\bar\alpha(x)\bigr)\le 1.
\end{equation}
\item[(6)] $\delta_1M/M_1\ge 2$, where $\delta_1$ is Ba{\l}aban's decay rate $\delta$ in the choice of
$M$ so large that $\delta(M/M_1)\ge 2$ in \cite{B-CMP122a}.
\item[(7)] $\mathsf{L}_0\ge 3$.
\item[(8)] The lower bound on the free constant $C_{\theta'}$:
\begin{equation}\label{5.2:eq-half-disc-containment-d}
C_{\theta'}\ge 192\,(1+b_+)^{1/2}K_Dc_\chi\varphi_\flat^{-1}B_HA_1/C_*.
\end{equation}
\item[(9)] The upper bound on $\gamma$ for points of $Z_h$, where $h:=k-N_k$:
\begin{equation}\label{5.2:eq-half-disc-containment-c}
\sup_{x\ge\gamma^{-2}}16\,K_Dc_\chi\varphi_\flat^{-1}C^J_{\mathrm{age}}\,
\frac{\bar\alpha(x)E(x)^{\delta_\natural/2}}{\alpha_*(x)}\le 1.
\end{equation}
\item[(10)] The upper bound on $\gamma$ for the $\mathfrak{T}$-locations and the live arrested pieces:
\begin{equation}\label{5.2:eq-half-disc-containment-e}
\sup_{\mathsf{T}\ge\mathsf{T}_\gamma-1}32\,C_{\theta'}^{-1}(\mathsf{T}+2)^{q-p_1}C_P'
e^{-\frac32\delta_1^{\mathrm{sc}}R(\mathsf{T})}\le\beta_{\mathrm{pr}}\theta_S.
\end{equation}
\item[(11)] The containments of the undilated chain:
\begin{enumerate}
\item[(11a)] the carrier pair presented at $(L)$, restricted to $X$, lies in
$\mathfrak{D}_{\mathfrak{E}}$ for every member $\mathfrak{E}$ of the consumer class, as a field on the
whole domain of $\mathbf{B}''_k$, and at $\mathsf{t}=0$ each entry of that pair is at most
$1/16$ of the corresponding threshold of $\mathfrak{D}_v$ on
$X\cap(\Lambda^\natural\cup\mathfrak{N}^{\mathrm{rec}})$;
\item[(11b)] on $X\cap Z^{\mathrm{on}}$ the undilated chain configurations, with $\sigma$ and the
segments $t\in[0,1]$, obey the clauses of the inherited domain with the margin $\frac12$ of the
available room, and on the rings $Z^{\mathrm{on}}\setminus\Lambda^\natural$ the ratio of
threshold to entry is at least $(1-4\beta_{\mathrm{pr}})\mathsf{L}_0^2$;
\item[(11c)] on live arrested pieces the undilated chain lies within the rooms of the room table
for live arrests in its first three zones, read with the monotonicity clause at a live piece of
another component, and in the fourth zone the undilated bound holds within $1/64$ of every
threshold, by the own-bound and the bending statement for the stored families, and within $1/16$,
by the containment corollary for the stored families, for
every threshold whose level is the current level of the point or a finer birth index;
\item[(11d)] on $X\cap\mathfrak{T}$ the undilated chain configurations obey the clauses of the
inherited domain in the form used in the containment part of
Proposition~\ref{5.2:prop-three-sites}: the entries $5,8$ are those of the source space, and the
other entries exceed those of the source space by the layer increments only, which are at most a
quarter of the available room; and on $X\cap\Lambda^\natural$ they obey the margin $\frac12$.
\end{enumerate}
\end{enumerate}
Then for every plain unit $v$ with live exterior slots, of class (c) or
(d) of the classification of units at $(L)$, that carries the amplitude $\mathsf{t}$, with
$\rho_A=\mathsf{w}_L/\eta_p^{1/2}$, every datum $p$, all $|\sigma(\Delta)|\le e^{\kappa_1}$, every
$t\in[0,1]$ and every $|\mathsf{t}|\le\rho_A$,
\begin{equation*}
\mathsf{R}_p(\mathsf{t},\sigma)^h\upharpoonright S_v\in\mathfrak{D}_v .
\end{equation*}
More precisely: off $\mathfrak{T}$ the dilation adds to each entry at most $1/16$ of its threshold
(for the composite entries $e=2,4$: $Q_e^+\le\frac{17}{16}Q_e+\frac1{16}\theta_e$). Hence on
$X\cap(\Lambda^\natural\cup\mathfrak{N}^{\mathrm{rec}})$ each entry is at most
$\frac1{16}+\frac1{16}$ of its threshold, and each of the entries $2,4$ at most
$\frac1{16}\cdot\frac{17}{16}+\frac1{16}$ of it; these fractions are below
$1-\varpi_{\mathrm{fr}}$, where $\varpi_{\mathrm{fr}}\le 1/80$ is the margin fraction of the stored
domain, and below $\frac12$, hence inside
the diminished thresholds $\theta^{(i)}_T\ge\frac12\theta_T$. Elsewhere on $Z^{\mathrm{on}}$ each
entry stays inside the $\frac12$-margin of (11b). On live pieces each entry stays within the rooms
of the room table, at the ratios: at most $\theta_S/4$ in the first zone, at most $\frac14$ in the
second, within the room in part (a) of the third, and less than $0.78$ in part (b) of the third.
None of the constants and none of the upper
bounds on $\gamma$ in (1)--(11) depends on $K$, $k$, an age, $N$, $N_0$, $\check{R}_k$, the width of
a window, the number of arrests or the volume.
\end{lemma}

\begin{proof}
\emph{The units concerned.} By the classification of units at $(L)$, the units that carry the
amplitude $\mathsf{t}$ there are plain, of two kinds. (II) Second-part units, for which
$X\cap Z=\emptyset$; they have four brackets, each of amplitude at most $\mathsf{U}_L$, and all
their clause locations lie in $\mathfrak{T}$, which is treated in the containment part of
Proposition~\ref{5.2:prop-three-sites}, or on live arrested pieces of other components, treated
below. (I) First-part units with $K_v\setminus Z\ne\emptyset$; they carry one $\mathsf{t}$-bracket,
that of the pair $(\mathfrak{b}_3,\mathfrak{b}_4)$, with
$\delta\mathbb{H}_p=\hat\chi_\square\mathbb{H}_4$, $\mathbb{H}_4=H'_k(g_kB)$, whose datum is
supported in the set $\mathbb{B}_0$ of \cite{B-CMP122b}, where
\begin{equation*}
\mathbb{B}_0\subset\Lambda_0:=\Lambda\cap\Omega''^{\sim2}_{h+1}\subset\mathsf{S}
\end{equation*}
(\cite{B-CMP122b}, \cite{B-CMP122a}).

\emph{The layer bound.} For the units (I), \cite[(1.45)]{B-CMP122a} states that the field in the
argument is bounded by $4\delta'_j$ and is localised, with
$\sup L^i\eta|H|$ and $\sup(L^i\eta)^2|\nabla H|$ both bounded by
$B_3\exp(-\delta d(y,\cdot))\,4\delta'_j$; both norms are used, each in the units of its own level.
With the bound
$N_y(\hat\chi_\square\mathbb{H})\le c_\chi N_y(\mathbb{H})$ for the block cut-off in
the cell norm this gives $N_y(\delta\mathbb{H}_p)\le c_\chi\mathsf{U}_J\mathsf{f}(y)$, where
\begin{equation*}
\mathsf{f}(y):=e^{-\delta_\natural(d_{\mathrm{sc}}(y,\mathsf{S})-1)_+},
\qquad\delta_\natural\le\delta_1^{\mathrm{sc}}.
\end{equation*}
The bound of the profile by $\eta_p$ times the unit shift holds in the cell norm $\mathsf{N}_y$,
hence in its entry $N_y$, at every cell of $S_v$, by \eqref{5.2:eq-three-sites-c}. For
$|\mathsf{t}|\le\rho_A=\mathsf{w}_L\eta_p^{-1/2}$ these
two bounds give
\begin{equation*}
|\mathsf{t}|\,N_y(\delta\mathbb{H}_p)
\le\mathsf{w}_L\eta_p^{-1/2}c_\chi\mathsf{U}_J\min\{\eta_p,\mathsf{f}(y)\}
\le c_\chi\mathsf{w}_L\mathsf{U}_J\min\{1,\mathsf{f}(y)^{1/2}\}.
\end{equation*}
The second inequality is elementary. If $\mathsf{f}\le\eta_p$, then
$\mathsf{f}\eta_p^{-1/2}=\mathsf{f}^{1/2}(\mathsf{f}/\eta_p)^{1/2}\le\mathsf{f}^{1/2}$. If
$\mathsf{f}\ge\eta_p$, then $\eta_p\cdot\eta_p^{-1/2}=\eta_p^{1/2}\le\mathsf{f}^{1/2}$. In both
cases $\eta_p\le 1$ gives the bound by $1$.

This is the one-line estimate of the containment part of Proposition~\ref{5.2:prop-three-sites},
read on the half disc \eqref{5.2:eq-three-sites-a}. In the same proposition the radius
$\mathsf{w}_L=1/\theta'(g_k)$ is fixed so that
$\mathsf{w}_L\mathsf{U}_J\le\tilde\alpha_1(k)\le\bar\alpha(x_k)$. Hence
\begin{equation}\label{5.2:eq-half-disc-containment-a}
|\mathsf{t}|\,N_y(\delta\mathbb{H}_p)\le\varkappa(y)
:=c_\chi\bar\alpha(x_k)\min\{1,\mathsf{f}(y)^{1/2}\}
\quad\text{on }|\mathsf{t}|\le\rho_A,\ \text{at every }y\in S_v .
\end{equation}
The function $\varkappa$ lies in the class $\mathfrak{M}^{\wedge}_{\delta_P/2}(\mathbf{B}''_k)$ of
majorants admitted by the linearised stability statement used below. Indeed, by the triangle
inequality for $d_{\mathrm{sc}}$, for cells $y,y'$ one has
$\mathsf{f}(y)^{1/2}\le\mathsf{f}(y')^{1/2}e^{\frac12\delta_\natural d_{\mathrm{sc}}(y,y')}$, and the
same holds for $\min\{1,\mathsf{f}^{1/2}\}$. By (4) and $d^\wedge\ge w_{\mathrm{sc}}\,d_{\mathrm{sc}}$,
\begin{equation*}
\tfrac12\delta_\natural d_{\mathrm{sc}}\le\tfrac12\delta_P w_{\mathrm{sc}}\,d_{\mathrm{sc}}
\le\tfrac12\delta_P d^\wedge .
\end{equation*}
Hence $\varkappa(y)\le\varkappa(y')e^{\frac12\delta_P d^\wedge(y,y')}$ for all cells $y,y'$, which is
the property of $\varkappa$ that membership in this class requires.

\emph{The linearised comparison.} The containment of a configuration in $\mathfrak{D}_v$ reads a
layer only through the following comparison, the linearised comparison lemma for layers, which we
use by statement; it is valid at every cell $y$ and at every triple of the
composite clauses. Here $\pi(y)$ is the level of $\mathbf{B}''_k$ at $y$. The index $n(y)\le\pi(y)$
is the one that the clause of $v$ reads at $y$: for a frame it is $k_W$, for a frozen term it is
its own freezing level. We put
\begin{equation*}
b(y):=\pi(y)-n(y)\ge 0,\qquad \varrho_n(y):=L^{-b(y)} .\end{equation*}
The factor $\varrho_n(y)$ is the scale drop of a level-$\pi$ layer read in level-$n$ units: the
weighted norm $\Theta^{(n)}_\nu$ at level $n$ is $\Theta^{(\pi)}_\nu$ multiplied by the
$(1+\nu)$-th power of the ratio of the length unit of level $n$ to that of level $\pi$. The
comparison is
\begin{equation}\label{5.2:eq-half-disc-containment-b}
\begin{aligned}
&\text{entries }5,8:\ \text{unchanged};\\
&\text{entries }1,3,6,7:\ \text{change, in threshold units,}\
\le K_D\varrho_n(y)\varkappa(y)/(\varphi_\flat\alpha_{*,n(y)});\\
&\text{entries }e=2,4:\ Q_e^+\le Q_e(1+\pi_\times)+\pi_{\mathrm{lin}}\mathbf{r}_e,\\
&\qquad\pi_{\mathrm{lin}}:=K_S^b\varkappa(y)/\tilde\alpha_{0,\pi(y)},\quad
\pi_\times\le\pi_{\mathrm{lin}}L^{-2b(y)}.
\end{aligned}
\end{equation}
The bounds for the entries $2,4$ are the
linearised stability statement for composite entries, which we use by statement. That statement
holds for every $1$-form with
majorant in $\mathfrak{M}^{\wedge}_{\delta_P/2}(\mathbf{B}''_k)$ under the condition
\eqref{5.2:eq-half-disc-containment-f}, read at
$\sup_y\varkappa(y)=c_\chi\bar\alpha(x_k)$. The quadratic piece of the comparison lemma for
composite layers is below its linear piece for layers not exceeding the threshold.

The thresholds are bounded from below by the threshold clause of the containment statement for
the stored families, which we also use by statement. Every threshold at $y$ whose level is the
current level of $y$ or a finer birth index has the form $F\,w'\,c\,\mathbf{r}_e$, with $w',c\ge1$ and
$F\ge\varphi_\flat$. In particular $\theta_a(y)\ge\varphi_\flat\alpha_{*,n(y)}$ in the unit of the
entry. The number $\pi_{\mathrm{lin}}$ is bounded by the same chain of inequalities as the
increments, read at $b=0$.

\emph{The base.} At $\mathsf{t}=0$ each entry is at most $1/16$ of its threshold on
$X\cap(\Lambda^\natural\cup\mathfrak{N}^{\mathrm{rec}})$, by (11a).

\emph{Case $(\beta)$: $y\in S_v\cap Z_h$.} Here $n(y)\le\pi(y)\le h$. This is the locus of every
clause location of class $(\beta)$, and $Z_h\cap\mathfrak{T}=\emptyset$. A contour from $y$ to
$\mathsf{S}$ first crosses the levels $\pi(y)+1,\dots,h-1$ inside $Z_h$ and then the moat between
$Z_h$ and $\mathsf{S}$.

Let $t(y)$ be the number of thick crossed levels. These are zones made by the operation
$\mathbf{T}$, each at least $10R_{\min}$ own-scale $M$-cubes wide, since each renormalisation step
adds at least ten layers of $MR_k$-cubes (\cite{B-CMP122a}). Let
\begin{equation*}
u(y):=(h-\pi(y)-1)_+-t(y)
\end{equation*}
be the number of the other crossed levels. These were re-made by an earlier attempt, with one layer
of $M$-cubes, one layer of $L^{-1}M$-cubes, and so on (\cite{B-CMP122a}), and each is at least one
own-scale $M$-cube wide, by the width bound for levels re-made by an attempt. The crossings lie on
disjoint arcs of the contour, so their contributions add:
\begin{equation*}
d_{\mathrm{sc}}(y,\mathsf{S})-1\ge(\check{R}_k-2)+10R_{\min}\,t(y)+u(y).
\end{equation*}
The first member comes from the distance lemma for live law sites of class $(\beta)$, whose proof
is geometric and holds at every point of $Z_h$. No increment per level beyond this count is claimed
at a point that is not a frame site.

Since $(r)_+\ge r$ for real $r$ and $\delta_\natural>0$, we have
$\mathsf{f}(y)^{1/2}\le e^{-\frac12\delta_\natural(d_{\mathrm{sc}}(y,\mathsf{S})-1)}$. We use the
distance bound just obtained; the inequality
$e^{-\frac12\delta_\natural(\check{R}_k-2)}\le E(x_k)^{\delta_\natural/2}$, used in the proof of
Proposition~\ref{5.2:prop-three-sites} and resting there on
$(\log(x_k-c_5))^{r_{\mathrm{pr}}}\le\check{R}_k$; and (1), which gives
$e^{-5\delta_\natural R_{\min}t}\le e^{-2t}$. Together these yield
\begin{equation*}
\mathsf{f}(y)^{1/2}\le E(x_k)^{\delta_\natural/2}\,e^{-2t}\,e^{-\frac12\delta_\natural u},
\qquad t=t(y),\ u=u(y).
\end{equation*}

\emph{The age count.} By the coupling-ratio bound of the flow theorem, which follows from the
comparison (3) together with (2), one has, for $0\le n\le k$,
\begin{equation}\label{5.2:eq-half-disc-containment-1}
\frac{\alpha_{*,k}}{\alpha_{*,n}}\le 2\Bigl(1+b_+\frac{k-n}{x_k}\Bigr)^{1/2}.
\end{equation}
Since $h=k-N_k$, we have
\begin{equation*}
k-n(y)=N_k+(h-\pi(y))+b(y).
\end{equation*}
Moreover $h-\pi(y)\le 1+(h-\pi(y)-1)_+=1+t+u$, so $k-n(y)\le N_k+1+t+u+b$. By
\eqref{5.2:eq-half-disc-containment-1} with $n=n(y)$,
\begin{equation*}
\frac{\alpha_{*,k}}{\alpha_{*,n(y)}}\le 2\Bigl(1+b_+\frac{N_k+1+t+u+b}{x_k}\Bigr)^{1/2}.
\end{equation*}
By (2), and with $x_k\ge 1$, the quotient in the bracket is at most $2+t+u+b$. Since
$2+t+u+b\ge 1$, we get $1+b_+(2+t+u+b)\le(1+b_+)(2+t+u+b)$. Expanding the product gives
$2+t+u+b\le 2(1+t)(1+u)(1+b)$. Hence
\begin{equation*}
\frac{\alpha_{*,k}}{\alpha_{*,n(y)}}
\le 2(1+b_+)^{1/2}(2+t+u+b)^{1/2}
\le 2\sqrt{2}(1+b_+)^{1/2}(1+t)^{1/2}(1+u)^{1/2}(1+b)^{1/2}.
\end{equation*}
Two elementary facts close the count: $(1+t)^{1/2}e^{-2t}\le 1$, and
$(1+b)^{1/2}L^{-b}\le 1$ for integers $b\ge0$. Multiplying,
\begin{equation*}
\varrho_n(y)\,\mathsf{f}(y)^{1/2}\,\frac{\alpha_{*,k}}{\alpha_{*,n(y)}}
\le 2\sqrt{2}(1+b_+)^{1/2}(1+u)^{1/2}e^{-\frac12\delta_\natural u}E(x_k)^{\delta_\natural/2}
\le C^J_{\mathrm{age}}E(x_k)^{\delta_\natural/2}.
\end{equation*}
The constant $C^J_{\mathrm{age}}$ depends on $(b_+,\delta_\natural)$ only. Half of the decay pays
the depth and the scale drop pays the birth index. No age, $N$, $N_0$, window or arrest count
survives. The thin levels are paid one at a time, so no comparison between the numbers of thin and
thick levels is used.

By \eqref{5.2:eq-half-disc-containment-b}, with
$\varkappa(y)\le c_\chi\bar\alpha(x_k)\mathsf{f}(y)^{1/2}$ and $\bar\alpha(x_k)=(\bar\alpha/\alpha_*)(x_k)\,\alpha_{*,k}$,
every additive increment at $y$, in threshold units, satisfies
\begin{equation*}
\text{increment}\le K_Dc_\chi\varphi_\flat^{-1}C^J_{\mathrm{age}}\,
\frac{\bar\alpha}{\alpha_*}(x_k)\,E(x_k)^{\delta_\natural/2}\le\frac1{16}.
\end{equation*}
The last inequality is \eqref{5.2:eq-half-disc-containment-c} read at $x=x_k$. There $\bar\alpha$
stands for the bound $\mathsf{w}_L\mathsf{U}_J\le\tilde\alpha_1(k)\le\bar\alpha(x_k)$, and the
condition is of the same kind as \eqref{5.2:eq-three-sites-b}: a supremum over $x\ge\gamma^{-2}$ of
a function of $x$ that decays faster than every power of $x^{-1}$ as $x\to\infty$.

\emph{The near case: $y\in S_v\cap(Z^{\mathrm{on}}\cup\Lambda^\natural)\setminus Z_h$.} Here
$\pi(y)\in[h,k]$, so $k-n(y)\le N_k+b$. By \eqref{5.2:eq-half-disc-containment-1} and (2), exactly
as above,
\begin{equation*}
\varrho_n(y)\frac{\alpha_{*,k}}{\alpha_{*,n(y)}}
\le 2(1+b_+)^{1/2}(1+b)^{1/2}L^{-b}\le 2(1+b_+)^{1/2}.
\end{equation*}
Using $\min\{1,\mathsf{f}^{1/2}\}\le 1$ and the bound $\mathsf{w}_L\mathsf{U}_J$ in place of
$\bar\alpha(x_k)$ in \eqref{5.2:eq-half-disc-containment-a}, the increment is at most
\begin{equation*}
2(1+b_+)^{1/2}K_Dc_\chi\varphi_\flat^{-1}\,\frac{\mathsf{w}_L\mathsf{U}_J}{\alpha_{*,k}}.
\end{equation*}

This quotient does not depend on the coupling. We have $\mathsf{U}_J=6B_H\delta'_k$ and
$\mathsf{w}_L=1/\theta'(g_k)$, and the three functions of the coupling are
\begin{equation*}
\delta'=gA_1\mathsf{T}^{p_1},\qquad \theta'=C_{\theta'}\mathsf{T}^{p_1-q},\qquad
\alpha_*=C_*g\mathsf{T}^{q}.
\end{equation*}
Therefore
\begin{equation*}
\frac{\mathsf{w}_L\mathsf{U}_J}{\alpha_{*,k}}=\frac{6B_H\delta'_k}{\theta'(g_k)\alpha_{*,k}}
=\frac{6B_HA_1}{C_{\theta'}C_*}.
\end{equation*}
So the increment is at most $1/16$ as soon as
\begin{equation*}
C_{\theta'}\ge 16\cdot 2\cdot 6\,(1+b_+)^{1/2}K_Dc_\chi\varphi_\flat^{-1}B_HA_1/C_*,
\end{equation*}
which is
\eqref{5.2:eq-half-disc-containment-d}. This is a lower bound on the free constant $C_{\theta'}$,
chosen after $K_D$. Raising $C_{\theta'}$ shrinks $\mathsf{w}_L$ and enlarges linearly the left
member of the smallness condition on the brackets, which is an upper bound on $\gamma$ read after
$C_{\theta'}$. The condition is one-sided and puts no upper bound on $C_{\theta'}$.

\emph{The case of live arrested pieces: $y$ in a live arrested piece, or
$y\in X_\tau\cap\mathfrak{T}_{k_a}$ off it.}
Here $k_a$ is the step of the arrest $\mathfrak{a}$, $\mathfrak{T}_{k_a}$ the region of
$\mathfrak{T}$-locations at that step, and $X_\tau$ the localisation domain of the term $\tau$ whose
clause is read at $y$.
These are the increments of the second kind in the bookkeeping of live arrests, at the site $(L)$
of a chain completing on another component, for units (I) and (II). The increment of that kind in
the live-arrest lemma is linear in the layer. It reads the geometry only through a decay factor of
at most $e^{-2}$ per crossed stratum, which (6) provides, thin strata included.

In what follows $\ell$ denotes the level from which the crossed strata are counted and $m$ the
shell index of the second zone of the room table.
On the half disc every far factor is halved. Per stratum the factor becomes
\begin{equation*}
e^{-1}=2^{-1}e^{-(1-\ln 2)},
\end{equation*}
and $(k-\ell)^{1/2}e^{-0.3(k-1-\ell)_+}\le C_{1/2}'$, as at the $\mathfrak{T}$-locations of
Proposition~\ref{5.2:prop-three-sites}. Hence the spend is at most
\begin{equation*}
2\mathsf{w}_L\bar P^{(1/2)}2^{-(k-1-\ell)_+},\qquad
\bar P^{(1/2)}:=C_P'e^{-\frac32\delta_1^{\mathrm{sc}}R_k}.
\end{equation*}
We compare it with the rooms of the room table, with $\beta=\beta_{\mathrm{pr}}$, the fraction of
the room table.
\begin{itemize}
\item In the first zone the room is $\beta 2^{-(k_a+1-\ell)}$, by the extended descent lemma, by
descent in $s$ alone. Since $2^{-(k-1-\ell)_+}2^{k_a+1-\ell}\le 4$, the ratio is at most
$8\mathsf{w}_L\bar P^{(1/2)}/\beta\le\theta_S/4$.
\item In the second zone the room is $\frac12\theta_S\beta2^{-(k_a-m)}$. Since
$2^{-(k-1-m)_+}\le 2\cdot2^{-(k_a-m)}$, the ratio is at most
$8\mathsf{w}_L\bar P^{(1/2)}/(\theta_S\beta)\le\frac14$.
\item In part (a) of the third zone the share is $0.3$, against the room
$\beta/8+\beta\theta_S/16$.
\item In part (b) of the third zone, whose undilated ratio is $\frac23$ and not at most
$\frac14$, the spend is $0.49+\beta_{\mathrm{pr}}\theta_S/16<0.78$ of the room.
\end{itemize}
All four comparisons hold under $32\mathsf{w}_L\bar P^{(1/2)}\le\beta_{\mathrm{pr}}\theta_S$. With
$\mathsf{w}_L$ read at its upper function and $R$ at its lower, as in
Proposition~\ref{5.2:prop-three-sites}, this is \eqref{5.2:eq-half-disc-containment-e}. That is the
$\mathfrak{T}$-condition of that proposition, read for these increments at halved far factors.

\emph{Case of dissolved pieces: $y$ in a piece $W_{\mathfrak{c}}\subset Z_h$ dissolved at this
site.} This covers the fourth zone, and the part of the third zone at $(L)$ with
$y\in\Lambda^\natural\cup\mathfrak{N}^{\mathrm{rec}}$. The undilated bound at such $y$ is taken from
the stored families' own bound, bending statement and containment corollary, as in (11c): by the
own-bound and the bending statement it holds within $1/64$, and by the containment corollary within
$1/16$, of every threshold at the current level of $y$ or a finer birth index.
On the half disc the increment added by the dilation at $y$ is bounded verbatim as in case
$(\beta)$, with $\tau$ the term whose clause is read at $y$, $n(y)$ the index that this clause
reads, and $k_c$ the level index of the dissolved piece $W_{\mathfrak{c}}$, so that
$n(y)\le\pi(y)\le k_c\le h$, and with the same $t$, $u$, $b$. It is therefore at most $1/16$ by
\eqref{5.2:eq-half-disc-containment-c}.

\emph{Conclusion.} The clause locations of $v$ lie in $\mathfrak{T}$, in $Z_h$, in
$(Z^{\mathrm{on}}\cup\Lambda^\natural)\setminus Z_h$, on live arrested pieces or on dissolved
pieces. The $\mathfrak{T}$-locations are covered by the containment part of
Proposition~\ref{5.2:prop-three-sites}, under (11d) and (10); the others by the four cases above.

Off $\mathfrak{T}$ each additive increment is at most $1/16$ of its threshold. For the entries
$2,4$ the number $\pi_{\mathrm{lin}}$ is bounded in the same way, and
$\pi_\times\le\pi_{\mathrm{lin}}L^{-2b}\le\pi_{\mathrm{lin}}$; the composite entries then obey
$Q_e^+\le\frac{17}{16}Q_e+\frac1{16}\theta_e$.

On $X\cap(\Lambda^\natural\cup\mathfrak{N}^{\mathrm{rec}})$ the base gives at most $\frac1{16}$.
After dilation each entry is at most $\frac18$ of its threshold, and the entries $2,4$ at most
$\frac1{16}\cdot\frac{17}{16}+\frac1{16}=\frac{33}{256}$. Both numbers are less than
$1-\varpi_{\mathrm{fr}}$, since $\varpi_{\mathrm{fr}}\le 1/80$, and less than $\frac12$, so the
entries lie inside the stored domain and inside the diminished thresholds
$\theta^{(i)}_T\ge\frac12\theta_T$.

Elsewhere on $Z^{\mathrm{on}}$ each entry stays inside the $\frac12$-margin of (11b); on the rings
$Z^{\mathrm{on}}\setminus\Lambda^\natural$ this is seen as follows. There
(11b) and (7) give
\begin{equation*}
\text{threshold}/\text{entry}\ge(1-4\beta_{\mathrm{pr}})\mathsf{L}_0^2=\tfrac15\cdot 9=\tfrac95
\end{equation*}
at $\mathsf{L}_0=3$, and more for larger $\mathsf{L}_0$. So the entry is at most
$\frac59\theta$ and the room at least $\frac49\theta$. Half the room is
$\frac29\theta>\frac1{16}\theta$, which is at least the share of the dilation; so the entries stay
inside the $\frac12$-margin of (11b). This is where (7) is used.

On live pieces the entries stay within the rooms of the room table at the ratios of the case of
live arrested pieces.
Finally, no constant in (1)--(11) depends on $K$, $k$, an age, $N$, $N_0$, $\check{R}_k$, a window,
an arrest count or the volume. In the age count the quantities $t,u,b$ enter only through factors
that are bounded by constants. In the case of live arrested pieces the indices $\ell$ and $m$ enter
only through factors
$2^{-(\cdot)}$ that appear on both sides of each comparison.
\end{proof}

\begin{remark}[Scope]
The lemma is formulated for units with live exterior slots, with
$\rho_A=\mathsf{w}_L/\eta_p^{1/2}$; from this radius the factor $\min\{1,\mathsf{f}^{1/2}\}$ in
\eqref{5.2:eq-half-disc-containment-a} is derived. When the units with live exterior slots are
summed at
$(L)$, only the containment of the undilated chain at $\mathsf{t}=0$ is used. The pointed pieces of
the plain units of classes (c), (d) --- the constituents listed as items $1$--$3$, $5$ and the
boundary-type item of the classification of units at $(L)$, and the far cores of class (d) --- are
read at $(L)$ on
the radius $\rho_A:=\mathsf{w}_L/(\eta_p\mathsf{f}_v)^{1/2}$. Here
$\eta_p:=e^{-\delta_Pd(\square^\zeta,S_v)}$, which equals $1$ where the $2$-stencil of $S_v$ meets
$\square^\zeta$. The lemma enters at $(L)$ through the compact ones among these pieces, and it is
read for the pointed pieces as follows.

(i) At a touching block the bound
\begin{equation*}
|\mathsf{t}|\,N_y(\delta\mathbb{H}_p)
\le\mathsf{w}_L\mathsf{f}_v^{-1/2}c_\chi\mathsf{U}F_s(y)
\le c_\chi\mathsf{w}_L\mathsf{U}\min\{1,\mathsf{f}(y)^{1/2}\}
\end{equation*}
holds at every cell $y$ of $S_v$, because $F_s\le\mathsf{f}\le\mathsf{f}_v$ on $S_v$. The extra
factor $\mathsf{f}_v^{-1/2}$ of the radius is paid by $F_s/\mathsf{f}_v\le1$, never by the half
exponent. The cases of the proof read a layer only through such a majorant, in
\eqref{5.2:eq-half-disc-containment-b}, so they hold verbatim.

(ii) At a receding block the $\mathbf{U}$-slot is $\mathsf{t}$-free on $S_v$. The one moving
quantity is the $\mathbf{J}$-tail $-\mathsf{t}T_p$, which is bounded by the lemma on the
$\mathbf{J}$-tail at receding blocks, at every clause index $n\le\pi(y)$.

(iii) For a compact $p$ the relevant cases are the near case, the $\mathfrak{T}$-locations and the
case of live arrested pieces. Here
$p$ is a short single, a Wilson piece or a core, with $S_v=\mathfrak{r}_v$, clause index $\pi(y)$
and $b(y)=0$, read at the auxiliary consumer term $\mathfrak{E}^{\oplus}_v$. Cases $(\beta)$ and
dissolved pieces are void, since there is no clause location in $Z_h$.

For a long single or a tail of (c) or (d) neither the four cases of the proof nor the
$\mathfrak{T}$-locations of Proposition~\ref{5.2:prop-three-sites} are used at $(L)$. These are read
at $\mathsf{t}=0$ only, with $S_v$ the set of cells meeting $X_v$ and $\mathfrak{D}_v$ the space
$U^c_j(X_v,\alpha_{\cdot,j})$ of \cite{B-CMP122b}, at every cell of $X_v$, cells of foreign windows
included, and with clause index $n(y):=j$. For them the undilated base is that of the lemma itself:
(11a) at the unit's own family, together with (11b), (11c) and (11d).

When the lemma is read for the pointed pieces, its constants are replaced once by majorants:
\begin{itemize}
\item $c_\chi$ by $c_{\mathrm{rec}}:=\max\{c_\chi,c_J^{\mathrm{rec}},c_a\}$ in
\eqref{5.2:eq-half-disc-containment-d} and \eqref{5.2:eq-half-disc-containment-c};
\item $C_P'$ by $c_{\mathrm{rec}}C_P'$ in \eqref{5.2:eq-half-disc-containment-e};
\item the supremum $c_\chi\bar\alpha(x)$ by $e^{2d+2}c_\chi\bar\alpha(x)$ in
\eqref{5.2:eq-half-disc-containment-f}.
\end{itemize}
At the sites $(T)$ and $(P)$ no layer of site $(L)$ is dilated, and the pointed pieces keep the
radius $\mathsf{w}/\eta_p$ there.

Interior first-part units and wrapped units carry no amplitude $\mathsf{t}$ at $(L)$: their
membership is given by the sourced containment lemma at $(L)$ together with (11a), on
$|\mathsf{s}(\Delta)|\le e^{\kappa_1}$, and on the disc of the deep interior expansion by the
containment lemma for deep interior units. The brackets of the first three of the four expansions
at $(L)$ carry an amplitude $\mathsf{t}$ on second-part units only.
\end{remark}

\begin{remark}[The domain clause at non-deep presentations]
Consider a non-deep presentation at
$(L)$ of a stored unit or of a member of the family $\mathfrak{Q}(\mathfrak{a})$ attached to the
arrest $\mathfrak{a}$ in the gauge-orbit domain theorem, that is, a piece $e$ of that theorem whose
set $S^+(e)$ is not contained in the core $\mathsf{C}_C$ of the component. The presented point $p(z)$
is the carrier pair of the $(L)$-chain with its presented layers at $z=(\mathsf{t},\sigma)$ on the
disc of the unit. It is the restriction to $S_v\supseteq S^+(e)$ of a global pair lying in one of
the configuration spaces of Ba{\l}aban's construction over a domain $Y\supseteq S^+(e)$, in the
sense of the extension clause of the inherited domain.

For the undilated chain this holds by (11a), since
\begin{equation*}
((\hat{\mathcal{U}}^{\mathfrak{b}})^h,\operatorname{Ad}_h\hat{\mathcal{J}}^{\mathfrak{b}})
\upharpoonright X\in\mathfrak{D}_{\mathfrak{E}}
\end{equation*}
for every consumer class member $\mathfrak{E}$, a
field on the whole domain of $\mathbf{B}''_k$. Together with the extension clause of the inherited
domain and the fact that the members of the arrest spaces are restrictions of pairs of
$\mathfrak{G}(\mathbf{B}''_k)$, this gives the global extension.

For the dilated chain it holds by the cases $(\beta)$, the near case, the case of live arrested
pieces and dissolved pieces above,
together with the $\mathfrak{T}$-locations of Proposition~\ref{5.2:prop-three-sites}: each entry is
within its threshold at every $y\in S_v$ and at every block. For the $\mathbf{J}$-entry this uses
the lemma on the $\mathbf{J}$-tail at receding blocks with its pointwise-product corollary. Hence,
by the gauge-orbit domain theorem, the presented point satisfies
\eqref{5.2:eq-half-disc-containment-g}, that is, it lies in the domain
$\mathfrak{K}_G(\tilde S(e))$ of that theorem over the set $\tilde S(e)$ attached to the piece $e$:
\begin{equation*}
p(z)\in\mathfrak{K}_G(\tilde S(e)),\qquad \operatorname{pol}g(z)<\varepsilon_\star .
\end{equation*}
The second bound follows from $\vartheta\le\mu<\varepsilon_\star$ in
Proposition~\ref{5.2:prop-three-sites}, together with the summability of the live margins.
\end{remark}

\begin{lemma}[The first-stage bound by two contour integrals]\label{5.2:lem-first-stage-bound}
Let $v=(j,S_v,f_v)$ be a unit at one of the three expansion sites (T), (P), (L) of
Proposition~\ref{5.2:prop-three-sites}. Here $f_v$ is analytic on the domain $\mathfrak{D}_v$, and
$\mathsf{R}_p(\mathsf{t},\sigma)$ is the response map attached to Ba{\l}aban's localisation data
$p=(\square,Y_0)$, a block $\square$ and a localisation domain $Y_0$, at that site. Write
$\mathsf{w}_{\mathrm{site}}\in\{\mathsf{w}_T,\mathsf{w}_P,\mathsf{w}_L\}$ for
the radius of the amplitude disc at the site. Ba{\l}aban's piece of $f_v$ with data $p$ is
\begin{equation}\label{5.2:eq-first-stage-bound-1}
v_p=\Big[\prod_{\Delta\subset Y_0\setminus\bar S_v}\int_0^1 d\mathsf{s}(\Delta)\,
\partial_{\mathsf{s}(\Delta)}\Big]\,\partial_{\mathsf{t}}\big|_{\mathsf{t}=0}
f_v\big(\mathsf{R}_p(\mathsf{t},\mathsf{s})\big)
\end{equation}
(\cite[(3.7)]{B-CMP109}, \cite[(1.10), (1.23)]{B-CMP116}), and $v_p$ is linear in $f_v$. Put
$m(p):=\#\{\Delta\}$, the number of cells in the product \eqref{5.2:eq-first-stage-bound-1};
the set $\bar S_v\supset S_v$ is the union of the exempt family of $v$, made precise at (L) below.

Assume the following on the disc $|\mathsf{t}|<\rho_A(v)$ in the rescaled amplitude, and for
the decoupling parameters continued to complex values $\sigma(\Delta)$ with
$|\sigma(\Delta)|\le e^{\kappa_1}$:
\begin{equation*}
\mathsf{R}_p\subset\mathfrak{D}_v,\qquad
|f_v\circ\mathsf{R}_p-f_v(1)|\le\mathfrak{N}(v).
\end{equation*}
Then the following hold.
\begin{enumerate}
\item[(i)] For the units of rows 1--3, 5 and $\partial$ of the table \eqref{5.2:eq-unit-table-a}
(those carrying no exterior slot), at the sites (T) and (P), with
$\rho_A=\mathsf{w}_{\mathrm{site}}/\eta_p$:
\begin{equation}\label{5.2:eq-first-stage-bound-2}
|v_p|\le 2\,\mathsf{w}_{\mathrm{site}}^{-1}\,\eta_p\,e^{-(\kappa_1-1)m(p)}\,\mathfrak{N}(v).
\end{equation}
\item[(ii)] For every row-4 unit of that table (those carrying live exterior slots), with
$\rho_A=\mathsf{w}_{\mathrm{site}}/\eta_p^{1/2}$:
\begin{equation}\label{5.2:eq-first-stage-bound-3}
|v_p|\le 2\,\mathsf{w}_{\mathrm{site}}^{-1}\,\eta_p^{1/2}\,e^{-(\kappa_1-1)m(p)}\,\mathfrak{N}(v).
\end{equation}
Here $f_v\circ\mathsf{R}_p$ means the orbit datum of the unit at the presented rotation. At (L)
it is the row-4 reading of Proposition~\ref{5.2:prop-three-sites},
\eqref{5.2:eq-three-sites-a}. At (T) and (P) it is the reading with the
exterior slots rotated by the gauge map of the response, as described with the three expansion
sites. Moreover $\mathfrak{N}(v)=2\mathfrak{N}^{\mathbb{C}}(v)$, by the complex bound on the
orbit datum.
\item[(iii)] At (L), take any compact pointed datum $p$ of a plain unit of the filing classes
(c) or (d) of the decomposition of units at (L) (compact pointed data are the short singles,
Wilson pieces and cores, as opposed to the long singles and tails; for them the read set is
$S_v=\mathfrak{r}_v$). These are the constituents of rows 1--3, 5 and $\partial$ and the far cores of
class (d). Suppose the hypothesis holds with $\tfrac{81}{64}\mathfrak{N}^{\flat}(v)$ in place of
$\mathfrak{N}(v)$ on the disc of radius $\rho_A=\mathsf{w}_L/(\eta_p\mathsf{f}_v)^{1/2}$. Then
\begin{equation}\label{5.2:eq-first-stage-bound-4}
|v_p|\le 2\,\mathsf{w}_L^{-1}(\eta_p\mathsf{f}_v)^{1/2}\,e^{-(\kappa_1-1)m(p)}\cdot
\tfrac{81}{64}\,\mathfrak{N}^{\flat}(v).
\end{equation}
At (L), in clause (ii) as well as here, $\eta_p:=e^{-\delta_P d(\square^{\zeta},S_v)}$, and
$\eta_p:=1$ where the
$2$-stencil $\operatorname{Sten}_2(S_v)$ meets $\square^{\zeta}$. Finally
$\mathsf{f}_v:=\sup_{S_v}\min\{1,c_\chi\mathsf{f}\}\le 1$, where $\mathsf{f}$ is the far factor
of the undilated layer. This form of the hypothesis is the one supplied by the pointed
containment statement.
\end{enumerate}
At (L) the cells are made precise as follows. The cells $\Delta$ are cells in proper scales,
namely the members of $\sigma_0(v):=\pi^{\mathfrak{s}}\setminus\operatorname{St}_v$, where
$\operatorname{St}_v$ is an own-scale exempt family with
$S_v\subset\bigcup\operatorname{St}_v=\bar S_v$. Writing $p=(\square,\mathsf{y};s)$, with
$\square$ the block, $\mathsf{y}$ the family of cells $\Delta$ and $s$ the index of the response
layer $\mathbb{H}_s$ whose block piece $\hat{\chi}_{\square}\mathbb{H}_s$ is dilated, one has
$m(p)=|\mathsf{y}|$ (\cite{B-CMP119}, \cite{B-CMP122b}).

The piece $v_p$ is analytic, uniformly, in every parameter in which $f_v$ is analytic. This
includes the complex source $w$, under the holomorphy of the orbit data in
$w\in\mathcal{W}_k(X)$ stated with the three expansion sites, $X$ being the domain of the
term $v$.
\end{lemma}

\begin{proof}
\emph{Scaling of the amplitude.} The amplitude $\mathsf{t}$ is measured in rescaled units, scaled
by $\eta_p$. By the definition of $\eta_p$, the piece seen on $S_v$ has size at most $\eta_p$
times the unit shift of the site. Hence, on $|\mathsf{t}|<\mathsf{w}_{\mathrm{site}}/\eta_p$,
the dilated piece on $S_v$ is at most $\mathsf{w}_{\mathrm{site}}$ times the unit shift.

\emph{Contour integral in $\mathsf{t}$.} Fix $\sigma$ with $|\sigma(\Delta)|\le e^{\kappa_1}$ for
every $\Delta$, and put
\begin{equation*}
g(\mathsf{t}):=f_v\big(\mathsf{R}_p(\mathsf{t},\sigma)\big)-f_v(1).
\end{equation*}
By hypothesis $\mathsf{R}_p\subset\mathfrak{D}_v$ on the disc $|\mathsf{t}|<\rho_A$. Since $f_v$
is analytic on $\mathfrak{D}_v$ and $\mathsf{R}_p$ is holomorphic in $\mathsf{t}$, the function
$g$ is analytic on $|\mathsf{t}|<\rho_A$, with $|g|\le\mathfrak{N}(v)$ there. The constant
$f_v(1)$ does not contribute to the derivative, so
$\partial_{\mathsf{t}}|_{\mathsf{t}=0}f_v(\mathsf{R}_p(\mathsf{t},\sigma))=g'(0)$.

Cauchy's inequality on the circles $|\mathsf{t}|=r<\rho_A$, with $r\uparrow\rho_A$, gives
$|g'(0)|\le\mathfrak{N}(v)/\rho_A$. We record the bound with the constant $2$ (Cauchy's inequality
on the circle of radius $\rho_A/2$), which is all that is used below. Hence
\begin{equation}\label{5.2:eq-first-stage-bound-5}
|g'(0)|\le\frac{2\,\mathfrak{N}(v)}{\rho_A}=
\begin{cases}
2\,\mathfrak{N}(v)\,\eta_p/\mathsf{w}_{\mathrm{site}}, & \rho_A=\mathsf{w}_{\mathrm{site}}/\eta_p,\\[2pt]
2\,\mathfrak{N}(v)\,\eta_p^{1/2}/\mathsf{w}_{\mathrm{site}}, &
\rho_A=\mathsf{w}_{\mathrm{site}}/\eta_p^{1/2}.
\end{cases}
\end{equation}
The first line serves clause (i), the second clause (ii). The same line, on the disc of radius
$\mathsf{w}_L/(\eta_p\mathsf{f}_v)^{1/2}$ and with $\tfrac{81}{64}\mathfrak{N}^{\flat}(v)$ in place
of $\mathfrak{N}(v)$, gives
\begin{equation*}
|g'(0)|\le 2\,\mathsf{w}_L^{-1}(\eta_p\mathsf{f}_v)^{1/2}\cdot\tfrac{81}{64}\,\mathfrak{N}^{\flat}(v)
\end{equation*}
for clause (iii). This bound is uniform in $\sigma$ on the polydisc
$\{|\sigma(\Delta)|\le e^{\kappa_1}\}$.

\emph{Contour integrals in the decoupling parameters.} The function
$\sigma\mapsto\partial_{\mathsf{t}}|_{\mathsf{t}=0}f_v(\mathsf{R}_p(\mathsf{t},\sigma))$ is
analytic on the polydisc, since $\mathsf{R}_p$ is holomorphic in $\sigma$. It is bounded there by
the right member of \eqref{5.2:eq-first-stage-bound-5}. Take the derivatives
$\partial_{\mathsf{s}(\Delta)}$, $\Delta\subset Y_0\setminus\bar S_v$, one at a time, at a real
point $\mathsf{s}(\Delta)\in[0,1]$. For each one apply Cauchy's formula on the circle
$|\sigma(\Delta)|=e^{\kappa_1}$, keeping the other variables fixed in the polydisc. For $h$
analytic on $|\sigma|\le e^{\kappa_1}$ and $0\le a\le 1$,
\begin{equation*}
|h'(a)|\le\frac{e^{\kappa_1}}{(e^{\kappa_1}-a)^2}\sup_{|\sigma|=e^{\kappa_1}}|h|
\le e^{\kappa_1}(e^{\kappa_1}-1)^{-2}\sup_{|\sigma|=e^{\kappa_1}}|h|.
\end{equation*}
Moreover
\begin{equation*}
e^{\kappa_1}(e^{\kappa_1}-1)^{-2}=e^{-\kappa_1}(1-e^{-\kappa_1})^{-2}\le e^{-(\kappa_1-1)}
\quad\text{as soon as}\quad 1-e^{-\kappa_1}\ge e^{-1/2}.
\end{equation*}
This holds for $\kappa_1\ge 1$, and that follows from the floor $\kappa_1\ge 16\kappa+1$ used in
the resummation below. Each of the $m(p)$ derivatives therefore contributes a factor
$e^{-(\kappa_1-1)}$, uniformly in the remaining variables. Each integral
$\int_0^1 d\mathsf{s}(\Delta)$ contributes at most $1$. With \eqref{5.2:eq-first-stage-bound-5},
this proves \eqref{5.2:eq-first-stage-bound-2}, \eqref{5.2:eq-first-stage-bound-3} and
\eqref{5.2:eq-first-stage-bound-4}. Since $v_p$ is obtained from $f_v\circ\mathsf{R}_p$ by the
contour integrals above, whose contours and bounds do not depend on the remaining parameters,
$v_p$ is analytic in every parameter in which $f_v\circ\mathsf{R}_p$ is, with the same bounds;
for $w$ this is the holomorphy of the orbit data in $w\in\mathcal{W}_k(X)$.

\emph{Containment on the smaller disc.} For a unit with live exterior slots, the hypothesis
$\mathsf{R}_p\subset\mathfrak{D}_v$ on the disc $|\mathsf{t}|<\mathsf{w}_{\mathrm{site}}/\eta_p^{1/2}$
is supplied by the site, as follows.
\begin{itemize}
\item At (L), at $\mathfrak{T}$-locations, by the containment clause of the row-4 reading at (L)
in Proposition~\ref{5.2:prop-three-sites}.
\item At (L), off the $\mathfrak{T}$-locations, by Lemma~\ref{5.2:lem-half-disc-containment}.
\item At (L), in both cases, over the containment of the undilated chain and of the live
arrested pieces.
\item At (T), by the homogeneous containment lemma.
\item At (P), by the local dilation corollary at that site.
\end{itemize}
At (L) the units with live exterior slots of classes (c) and (d) use this lemma only at the
centre $\mathsf{t}=0$, as explained after the proof. The disc, the reading
\eqref{5.2:eq-three-sites-a} and Lemma~\ref{5.2:lem-half-disc-containment} supply the compact
pointed $p$ at (L), and the units with live exterior slots at (T) and (P).

\emph{The rotation riding with $\mathsf{t}$.} For a unit with live exterior slots, the presented
rotation depends on the configuration, hence on $\mathsf{t}$. For a positive function $\varepsilon$
on the live law sites, write $\mathfrak{G}_{\varepsilon}(S_v)$ for the polydisc of gauge rotations
on $S_v$ whose polar part at each live law site $s$ has size less than $\varepsilon(s)$; the
orbit datum is defined on the chart $\mathfrak{G}_{\varepsilon_k}(S_v)$, $\varepsilon_k(s)$ being
its radius at level $k$. The presented rotation stays inside $\mathfrak{G}_{\varepsilon_k}$, by
the riding-rotation clause of the row-4 reading at (L) in
Proposition~\ref{5.2:prop-three-sites}. Concretely, the riding-rotation lemma is applied as it
stands with the following data:
\begin{itemize}
\item $\eta_0(s)$, the polar size of the undilated presentation rotation at $s$;
\item the margin $\mu(s)=\mu_k\mathsf{e}(s)/\varepsilon_\star$, already reserved for this
presentation in the chart bound;
\item $\vartheta(s)$, the bound on the $\mathsf{t}$-variation of the rotation at the live law
site $s$.
\end{itemize}
Under \eqref{5.2:eq-three-sites-b}, that riding-rotation clause gives
$\vartheta(s)\le\mu(s)$ at every live law site. Hence
\begin{equation*}
\Lambda:=\max\Big\{1,\max_s\frac{\vartheta(s)}{\mu(s)}\Big\}=1,
\qquad
\rho_A\cdot\min\Big\{1,\min_s\frac{\mu(s)}{\vartheta(s)}\Big\}=\rho_A.
\end{equation*}
So the disc returned by the riding-rotation lemma is the full disc $|\mathsf{t}|<\rho_A$.

The function $g$ then lives on the chart $\mathfrak{G}_{\varepsilon'}(S_v)$, with
$\varepsilon'(s)=\varepsilon_k(s)-\eta_0(s)-\mu(s)$, which every later use of this presentation
employs.
The chart radius is reduced by $\eta_0+\mu$ and by nothing else. The bound
\eqref{5.2:eq-first-stage-bound-5} therefore holds with the rotation riding along.
\end{proof}

Clause (ii) is used at (T) and (P). At (L) it is not used for the following units:
\begin{itemize}
\item the units with live exterior slots of the filing classes (c) and (d);
\item the long singles and tails of classes (c) and (d).
\end{itemize}
These units are treated by the summation of boundary terms at finer zones. That summation uses
from the present lemma only its hypothesis at $\mathsf{t}=0$: $\mathsf{R}_p\subset\mathfrak{D}_v$
and $|f_v\circ\mathsf{R}_p-f_v(1)|\le\mathfrak{N}(v)$, the same for every block $\square$. It does
not use the gain of (ii).

Their pieces $v_p$ are resummed exactly over the blocks $\square$, and over the fine cells of
each $\pi_{j_v}$-cube. Write $\Phi_t$, $t\in[0,1]$, for the interpolated configuration of the
response. By the exact block-resummation lemma (linearity of the differential of $f_v$ and
$\sum_\square\hat{\chi}_{\square}=1$),
\begin{equation*}
\sum_{\square}\partial_{\mathsf{t}}\big|_{\mathsf{t}=0}
f_v\big(\Phi_{t+\mathsf{t}\hat{\chi}_{\square}}\big)=\frac{d}{dt}f_v(\Phi_{t}).
\end{equation*}
Each bundle $\mathfrak{b}$ is then bounded by $2\mathfrak{N}(v)e^{-(\kappa_1-1)|\tilde{\mathsf{y}}|}$
with no amplitude disc. The bundles are summed by the walk-factor summation theorem in the
unit's own grain.

Lemma~\ref{5.2:lem-first-stage-bound} is the one-scale first-stage estimate with the
oscillation $\omega(v)\le 4\mathfrak{N}(v)/\mathsf{w}_{\mathrm{site}}$. Nothing in its proof
depends on the site or on the term. The only effect of clause (ii) is the replacement
$\eta_p\mapsto\eta_p^{1/2}$ in the resummation that follows.

\emph{The resummed bound.} Let $V(Y)$ denote the localised term collecting the pieces $v_p$ with
localisation domain $Y_0=Y$. The pieces are resummed:
\begin{itemize}
\item by the one-scale resummation of the first stage;
\item at site (P), by the M\"obius resummation at that site, with its two summation conditions;
\item by the corresponding $z$-free bounds for the sums over $p$, $\square$ and $Y_0$.
\end{itemize}
This gives
\begin{equation}\label{5.2:eq-first-stage-bound-6}
|V(Y)|\le C_A'\,\mathsf{W}\,e^{-(1-2\delta)\kappa d(Y)},
\end{equation}
with
\begin{equation}\label{5.2:eq-first-stage-bound-a}
\begin{gathered}
\mathsf{W}:=\sup_{\Delta}\sum_{v:\,S_v\cap\Delta\neq\emptyset}
4\,\mathsf{w}_{\mathrm{site}}^{-1}\,\mathfrak{N}(v)\,
e^{(1-\frac12\delta)\kappa d(\bar S_v)},\\
C_A':=C_A^{(1/2)}+\tfrac12 e\,C_H^{\wedge\prime},\qquad
C_H^{\wedge\prime}:=(2L)^d\big(1+\hat b_t+\hat b_tC_{\mathrm{cell}}\big).
\end{gathered}
\end{equation}

\emph{Pieces summed elsewhere at (L).} Two families of pieces at (L) enter neither
$\mathsf{W}$ nor $C_A'$.
\begin{itemize}
\item A compact pointed $p$ of class (c) or (d). Its sum over $(v,s,\square,\mathsf{y})$ is the
summation of compact pointed pieces. There the factor $\eta_p^{1/2}$ is summed by the
block-shell lemma, the walk factor by the walk-factor summation theorem, and the factor
$\mathsf{f}_v^{1/2}$ cell by cell with the weight $\bar{\mathsf{w}}^3$, independently of
$\square$.
\item The long singles and tails of classes (c) and (d), together with the units with live
exterior slots. They are summed in $\square$-bundles by the summation of boundary terms at finer
zones, with no factor $\eta_p^{1/2}$, no use of the block-shell lemma and no shell weight
$\mathsf{b}_\square$.
\end{itemize}
Consequently, at (L), $\mathsf{W}$ and $C_A'$ act on no unit of classes (c) and (d), and
$\eta_p$, $\rho_A$ enter only through the compact pointed $p$.

\emph{The first summand of $C_A'$: sites (T) and (P).} At these sites $C_A^{(1/2)}$ is $C_A$, the
constant of the one-scale resummation, read at the half decay rate $\tfrac12\delta_1^{116}M$
between $\square$ and $S_v$. In the one-scale resummation, the sum over $\square$ at fixed $S_v$
is geometric in $\eta_p=e^{-\delta_1^{116}M\cdot\operatorname{dist}(\square,S_v)}$. With
$\eta_p^{1/2}$ in its place it is the same sum at rate $\tfrac12\delta_1^{116}M$ per $M$-cube.
Since $\delta_1^{116}M\ge\kappa_1\ge 16\kappa+1$,
\begin{equation*}
\tfrac12\delta_1^{116}M\ge\tfrac12\kappa_1\ge 8\kappa,
\end{equation*}
so the sum converges, with a constant of the same kind.

\emph{The second summand of $C_A'$: site (L).} At (L) the one-scale resummation uses $\eta$
only through $C_H:=\sup_\square\sum_\Delta\eta(\square,\Delta)$, where
$\eta(\square,\Delta):=e^{-\delta_Pd(\square^{\zeta},\Delta)}$, set equal to $1$ when the
$2$-stencil of the cell $\Delta$ meets $\square^{\zeta}$, is the factor $\eta_p$ of clause (iii)
with the read set replaced by the anchor cell $\Delta$. With $\eta^{1/2}$ in place of $\eta$ it
uses it through $C_H^{(1/2)}:=\sup_\square\sum_\Delta\eta(\square,\Delta)^{1/2}$. The passage to
$e^{-(1-\delta)\kappa d_k(Y)}$ and the shape factors come from the walk
factor, since $Y_0\supset\bar S_v\cup\square$ is connected. The second summand
$\tfrac12eC_H^{\wedge\prime}$ of $C_A'$ is the price at (L) of the bound
$C_H^{(1/2)}\le C_H^{\wedge\prime}$, which we now prove.

\begin{proof}[Proof of $C_H^{(1/2)}\le C_H^{\wedge\prime}$]
Consider an anchor cell $\Delta$ with $\eta(\square,\Delta)=1$. It holds a point $x_0$ such that
$|x-x_0|\le 2\ell_i$ for some $x\in\square^{\zeta}$ and some cell of level $i$ holding $x_0$,
$\ell_i$ denoting the side of a cell of level $i$. By the zone lemma, $i\le m+1$ and
$|x-x_0|\le 2L\ell_m<\tfrac13M\ell_m$, where $m$ is the zone of the block, the last inequality by
the floor $M>6L$, a condition on $M$ fixed after $L$. Hence $x_0$ is an interior point of
$\square$, and so $\Delta\in\Sigma_\square\cup\partial^{+}\Sigma_\square$. Here $\Sigma_\square$
is the set of cells meeting $\square$ with positive measure, and $\partial^{+}\Sigma_\square$ is
the set of their neighbours.

A block of zone $n$ meets cells of the levels $n-1$, $n$, $n+1$ only. Therefore, by the counting
clause for blocks and cells, $\#\Sigma_\square\le(2L)^d$, and by the neighbour count
$\#\partial^{+}\Sigma_\square\le\hat b_t\,\#\Sigma_\square$. A receding anchor cell satisfies
\begin{equation*}
\eta^{1/2}(\square,\Delta)\le e^{-\frac12\delta_P d^{\wedge}(\Sigma_\square,\Delta)},
\end{equation*}
because $d(\square^{\zeta},\Delta)\ge d^{\wedge}(\square^{\zeta},\Delta)\ge
d^{\wedge}(\Sigma_\square,\Delta)$, using $\square^{\zeta}\subset\Sigma_\square$. By the
cell-sum bound at $\beta=\tfrac12\delta_P\ge\delta_P/16$, under its floor condition on
$\delta_Pw_{\wedge}$, where $w_{\wedge}$ denotes the constant in the comparison
$d^{\wedge}\ge w_{\wedge}d_{\mathrm{sc}}$ of the scaled distances, these factors sum to at most
$C_{\mathrm{cell}}\,\#\partial^{+}\Sigma_\square$. Hence
\begin{equation*}
C_H^{(1/2)}=\sup_\square\sum_\Delta\eta(\square,\Delta)^{1/2}
\le(2L)^d(1+\hat b_t)+C_{\mathrm{cell}}\hat b_t(2L)^d=C_H^{\wedge\prime}.
\end{equation*}
This is the block-count bound of the anchor recursion at $\beta=\tfrac12\delta_P$, on
$\square^{\zeta}\subset\square$. For the sum over blocks at a fixed anchor cell, the cell-sum bound
at $\Sigma:=\Delta$ and the covering number $3\cdot 2^d\le C_{\mathrm{blk}}$ on $\square$ give
\begin{equation*}
\sup_\Delta\sum_\square\max\big\{\mathbf{1}[\Delta\in\Sigma_\square\cup\partial^{+}\Sigma_\square],
\,e^{-\frac12\delta_Pd^{\wedge}(\square^{\zeta},\Delta)}\big\}
\le(1+\hat b_t)C_{\mathrm{blk}}+C_{\mathrm{blk}}(1+\hat b_tC_{\mathrm{cell}}).
\end{equation*}
Cells and blocks of all levels enter. The bound depends on none of $M$, $M_1$, $\kappa_1$, and on
no level, age, zone or band count.

Consider a cell lying low in the thin band of an arrest. It lies in $\tilde{\square}$ for the
blocks of the higher band zones. Its stencil $\operatorname{Sten}_2$, however, does not meet
$\square^{\zeta}\subset\square$ for those blocks, which are therefore receding for it. The set
$\square_m^{\zeta}$ keeps $L/3$ cubes of its own sub-level off the next finer zone, so each such
block is damped by the zone separation bound, using $\delta_Pw_{\wedge}\ge 16$:
\begin{equation*}
e^{-\frac12\delta_Pd(\square_m^{\zeta},S_v)}\le e^{-8(m-2)_+}.
\end{equation*}
The number of such blocks in zone $m$ is at most $3\cdot 2^d$, so their contributions sum over
$m$ to at most
\begin{equation*}
3\cdot 2^d\big(2+(1-e^{-8})^{-1}\big),
\end{equation*}
a number independent of the band; no additional term arises from these blocks.
\end{proof}

\emph{Place of $C_A'$ in the order of choice.} The number $C_A'\ge C_A$ depends on
$(d,L,M,\kappa_1)$; its (L)-summand depends on $(d,L)$ only. It is fixed before $\gamma$, at the
place of $C_A$ in the order of choice. It replaces $C_A$ in $C_{\mathsf{W}}$, $\mathsf{K}^0$,
$\varepsilon_{\mathrm{br}}$ and in the constant $\gamma_V$ of \eqref{5.2:eq-complex-localisation-a},
where it is written $C_A$ with this meaning. The output rate $(1-2\delta)\kappa$ and the walk
factor $e^{-(\kappa_1-1)m(p)}$ are unchanged. No condition on $M_1$ arises. The constant
$C_H^{\wedge\prime}$ is used again by the block-count bound of the anchor recursion and by the
bound on the members of class (d) attached to live arrests and filed with the pieces of the
present lemma.

\emph{Units that enter no $\mathsf{W}$.} Let $v$ be a boundary-type unit with
$\mathcal{P}_v\neq\emptyset$, at (T) or at (L), with live or dead strips. For such $v$ the first
stage is the source-expansion lemma, under its source condition. It is summed by its summation
lemma together with the anchoring lemma, and $v$ enters no $\mathsf{W}$ and no
$\hat{\chi}_{\square}$-telescoping. There
\begin{equation*}
F(\mathsf{R}(\mathbb{1}))-F(\mathsf{R}(0))=\sum_{\mathsf{y}\in\mathfrak{Y}(F)}F_{\mathsf{y}},\qquad
F_{\mathsf{y}}:=\int_{[0,1]^{\mathsf{y}}}\prod_{\Delta\in\mathsf{y}}d\mathsf{s}(\Delta)\,
\partial_{\mathsf{s}(\Delta)}F(\mathsf{R}(\mathsf{s}))\Big|_{\mathsf{s}=0\text{ off }\mathsf{y}},
\end{equation*}
and
\begin{equation*}
|F_{\mathsf{y}}|\le\mathfrak{N}(F)\,e^{-(\kappa_1-1)|\mathsf{y}|}.
\end{equation*}
There is no amplitude disc and no $\eta_p$; the distance is seen through the vanishing clause of
the source-expansion lemma. The totals of these pieces are as follows.
\begin{itemize}
\item At (T), the total is $V^{\mathrm{rel}}(Y)$, with
$|V^{\mathrm{rel}}(Y)|\le\varepsilon^{\mathrm{rel}}_k e^{-\kappa_Vd_k(Y)}$ under the relative
smallness condition at (T), as in \eqref{5.2:eq-complex-localisation-a}.
\item At (L), the total is bounded by the constant $\varepsilon^{\mathrm{rel}}(k)$ of the bracket
proposition, under the relative smallness condition at (L).
\item At (P), the expansion is unchanged.
\end{itemize}
At (L), no plain interior unit of the first part enters $\mathsf{W}$ either: its pieces are those
of the interior-piece lemma.

Hence, at (L), the supremum defining $\mathsf{W}$ contains no unit of classes (c) and (d). Their
units with live exterior slots and their long units are summed by the summation of boundary terms
at finer zones, as described after the proof of Lemma~\ref{5.2:lem-first-stage-bound}; the compact
pointed $p$ are summed by the summation of compact pointed pieces.

So the one-scale resummation, $\mathsf{W}$ and $C_A'$ act at (L) on no unit of classes (c) and
(d), and at (T) and (P) on the plain units. No length of a true domain containing a strip is
ever weighed in \eqref{5.2:eq-first-stage-bound-6}.

\begin{lemma}[The second stage sees activities only]\label{5.2:lem-second-stage-activities}
Assume:
\begin{itemize}
\item the cluster-expansion lemma at site (T), with its clauses (B1)--(B4);
\item the second and third steps of the collar-expansion proposition, read through the two-slot
machine clause, the two-slot M\"obius lemma, the exterior-slot polydisc lemma and its
corollary, and the localised all-blocks dilation corollary;
\item the decay-fraction lemma for $\delta_J(L)$.
\end{itemize}
The second stage of the expansion is Ba{\l}aban's cluster expansion
\cite[(2.1)--(2.41)]{B-CMP116}, run after the first stage of
Lemma~\ref{5.2:lem-first-stage-bound}. It is used at two sites. At site (T) it is run on the
integral over $\Lambda_{k+1}$ at $A_k=0$, where the measures \cite[(3.27), (3.29)]{B-CMP119}
carry no tilt; the rim datum $x:=A_k|_{\mathfrak{r}}$ is treated in the rim version of this
lemma. At site (P) it is run for the conditional covariance on the collar, written
$\mathcal{C}$ as in the collar-expansion proposition.

Let the potentials $V(Y,\cdot,B)$ satisfy the following hypotheses:
\begin{itemize}
\item they are analytic on Ba{\l}aban's spaces and gauge invariant;
\item they are of the form \cite[(1.42)]{B-CMP116}, with \cite[(1.43), (1.36)]{B-CMP116};
\item their activity constant $\bar\varepsilon'$ satisfies $2\bar\varepsilon'/\alpha_4\le\nu_0$,
the threshold hypothesis of the cluster-expansion lemma at site (T).
\end{itemize}
With $\varepsilon_2:=2\bar\varepsilon'/\alpha_4$, which is free of $\alpha_6$, the member of that
hypothesis used below reads
\begin{equation}\label{5.2:eq-second-stage-activities-1}
2(L+2)^4c_{37}\,\alpha_6^{-2}\,\varepsilon_2\,e^{5\rho_7}\le 1 .
\end{equation}
Here $c_{37}$ is Ba{\l}aban's constant of \cite{B-CMP116}.
One factor $\alpha_6^{-1}$ accounts for the use of \cite[(2.29)]{B-CMP116} at level $k$;
it is carried by the constant denoted $\varepsilon_2$ in \cite[(2.29)]{B-CMP116}, which is to be
distinguished from the $\varepsilon_2$ defined above. The other accounts for
its use at level $k+1$.
The factor $e^{5\rho_7}$ is the reading of the factor $\exp 5\kappa$ of the level-$k$ estimate
of \cite{B-CMP116}. The constant $c_{41}$ of \cite[(2.39)--(2.40)]{B-CMP116} carries the
factor $\alpha_6^{-1}e^{5\rho_9}$. These factors are as fixed in the list of constants of the
analysis of Ba{\l}aban's step.

Then at sites (T) and (P) the following hold.

The outputs are analytic on Ba{\l}aban's output spaces and in all parameters. These outputs
include the pinned kernels \cite[(3.47)]{B-CMP119}, the kernels of the boundary terms
$\mathbf{B}$ and of $R''$, and the chain terms $\mathbb{C}^{(n+1)}_k(X)$. An output with
localisation domain $X$ is bounded by
\begin{equation}\label{5.2:eq-second-stage-activities-2}
c_{41}\,2(L+2)^4c_{37}\,\alpha_6^{-2}\,\bigl(2\bar\varepsilon'/\alpha_4\bigr)\,e^{-\kappa d(X)},
\end{equation}
where $d(X)$ is the tree length of $X$ at the scale of the class of localisation domains to
which the output belongs, and with Ba{\l}aban's additional factors for those outputs which
carry them in \cite{B-CMP119}, \cite{B-CMP122a}.

At site (T), consider the outputs at level $k+1$: $E^{(k+1)}$, the pinned pieces,
$R''^{(k+1)}$, $\mathbf{B}^{(k+1)}|_{A_k=0}$ and the filed rim terms
$\mathfrak{K}(X,\mathcal{S})$. Each is bounded by the constant of
\eqref{5.2:eq-second-stage-activities-2} times $e^{-\mathsf{r}(\delta)\kappa d_{k+1}(X)}$,
where
\begin{equation}\label{5.2:eq-second-stage-activities-3}
\mathsf{r}(\delta)\kappa=(1-10\delta)\tfrac12L\kappa\ \ge\ (1+4\beta_s)\kappa=2\kappa .
\end{equation}
Here $\delta\in\,]0,\delta_J(L)]$ is the margin parameter of \cite{B-CMP116}. These rates do
not depend on $V$ and hold uniformly in the complex source $w=\sum_iz_if_i$.

An output with domain $X$ reads only the $V(Y)$ with $Y\subset X$.

Finally, let $V=V^u+\lambda V^m$, let $\bar\varepsilon'_m$ be the activity constant of $V^m$,
and let $t_\lambda$ denote the output at the value $\lambda$. Then the marked sub-sum satisfies
\begin{equation*}
|t_1-t_0|\le K_{CE}\,\bar\varepsilon'_m\,C_A^{-1}\,e^{-\kappa d(X)} ,
\end{equation*}
with $d(X)$ as in \eqref{5.2:eq-second-stage-activities-2}.
\end{lemma}

\begin{proof}
We first locate the potentials in the second stage.

In \cite[(2.1)]{B-CMP116} the potentials $V(Y)$ enter only as Mayer factors of the
polymers $Y$. In \cite[(2.14)]{B-CMP116} these factors are realised by Cauchy integrals in
the variables $\tau(Y)$. Next, \cite[(2.15)]{B-CMP116} bounds
$\exp[\sum_Y\tau(Y)V(Y)]$ by $\exp[\sum_Y|\tau(Y)|\,|V(Y)|]$. Then \cite[(2.18)]{B-CMP116}
fixes the radii $|\tau(Y)|$ by the activity.

After \cite[(2.20)]{B-CMP116} no potential occurs. Indeed, \cite[(2.26)]{B-CMP116} is the
product $\prod_{Y\in D}[2/|\tau(Y)|]$ times factors which do not depend on $V$. The estimates
\cite[(2.27)--(2.41)]{B-CMP116} depend on the potentials only through
$\varepsilon_2=2\bar\varepsilon'/\alpha_4$, and monotonically so.

Everything else in the second stage is independent of the complex source $w$. This covers
the reference Gaussian measure, its complex perturbation by the parameters
$(\sigma,\mathbf{U},\mathbf{J})$, the characteristic functions $\chi$ and the random walks. At
site (T) this holds for the centred fluctuation $A-m(x)$, $x=A_k|_{\mathfrak{r}}$, whose mean
$m(x)=Tx$ is exact, $T$ being the transfer operator of the rim construction.

Since $V$ enters only through Mayer factors, an output whose domain is $X$ reads $V(Y)$ only
for the polymers $Y$ of the clusters filed at $X$. These polymers satisfy $Y\subset X$.

\emph{Site (T).} The hypothesis $2\bar\varepsilon'/\alpha_4\le\nu_0$ is the threshold of the
cluster-expansion lemma at site (T), including its member
\eqref{5.2:eq-second-stage-activities-1}. At this site the assertion is therefore that lemma's
clauses (B1)--(B4), taken by statement.

Two of Ba{\l}aban's remarks in \cite{B-CMP119} enter here. The first concerns the pinned
factor of \cite[(3.47)]{B-CMP119}: although it differs from the other factors, it satisfies
the same bound as they do. The second concerns the per-point kernels: by
\cite[p.~278]{B-CMP119}, they satisfy the bounds \cite[(2.42)]{B-CMP116} with the constant
$B_0$ replaced by $O(p_0^3(g_k))$, where $p_0(\cdot)$ is the degree function. The exponent $3$
is not used below. As in
\cite{B-CMP119}, this factor is absorbed by $g_k$, since multiplication by $g_k$ yields small
bounds.

The restriction to $A_k=0$ is that of the statement: for $A_k\neq0$ the measures
\cite[(3.27), (3.29)]{B-CMP119} carry a bilinear tilt in $A$ and $A_k$, which the rim version
of this lemma treats.

For the rate of the level-$(k+1)$ outputs we use clause (B3) of the cluster-expansion lemma at
site (T). Before the pinned factor is inserted, it gives
\begin{equation*}
|t_\lambda(X)|\le c_{41}\,C_3(\lambda)\,\varepsilon_1\,
e^{-(1-10\delta)\frac12L\kappa d_{k+1}(X)} ,
\end{equation*}
where $C_3(\lambda)$ and $\varepsilon_1$ are the constants of that clause; by the first of
Ba{\l}aban's remarks above, the pinned factor satisfies the same bound as the other factors.
Since $\delta\le\delta_J(L)=(1-4/L)/10$, we have $1-10\delta\ge4/L$. Hence, because
$\beta_s=1/4$,
\begin{equation*}
\mathsf{r}(\delta)=(1-10\delta)\tfrac12L\ \ge\ \tfrac4L\cdot\tfrac L2=2=1+4\beta_s .
\end{equation*}
This is
\eqref{5.2:eq-second-stage-activities-3}. By the decay-fraction lemma, this rate does not
depend on $V$ and holds uniformly in $w$.

\emph{Site (P).} Here the same displays of \cite{B-CMP116} are run by the second and third
steps of the collar-expansion proposition, for the conditional covariance $\mathcal{C}$ on
the collar. On the collar cubes
$\square'\subset\Omega^{\sim-1}_{j+1}\setminus\Lambda_{j+1}$, the characteristic function
\cite[(1.9)]{B-CMP122a} is read as
\begin{equation}\label{5.2:eq-second-stage-activities-a}
\chi\bigl(\{\,|A_j(b)|<\mathfrak{b}^{\sharp}_j(b)\,\}\bigr).
\end{equation}
The outer layer keeps the threshold $\mathfrak{b}_j$ of \cite[(1.28)]{B-CMP122a}, not
$C_{\mathrm{far}}\mathfrak{b}_j$.

The estimate \cite[(2.22)]{B-CMP116} is applied bond by bond, at thresholds
$T_b\ge\mathfrak{b}^{\flat}_j$. The floor and the parameter $\vartheta$ of the bond-threshold
statement are read at $\mathfrak{b}^{\flat}_j=\mathfrak{b}_j/(4\Theta_1)$. This is
admissible once $(\log g_j^{-2})^{p_0-p_1}\ge4\Theta_1$, one of the $\gamma$-ceilings of the
order of choice; here $p_1$ denotes the auxiliary exponent of Ba{\l}aban's parameter list
$\mathfrak{p}$ which enters that ceiling.

Ba{\l}aban formulates the expansion of \cite{B-CMP116} also ``for integrals conditioned to
subdomains of the lattice''. In \cite{B-CMP122a} he describes the collar integrals as being of
the same type as an integral already treated in his earlier papers.

Two further ingredients enter at this site. The first is the filing geometry. The second is
the accumulation in norm form over the $N(g)$ preparatory stages. By the stage-counting lemma,
at most $N_0+3$ stages meet a point. In the weighted norm $\|\cdot\|^{A}$ the corresponding
member is the one of the norm-form second-stage proposition, under the factor
$\mathsf{w}_P^{-1}$; the unweighted member $3(N_0+3)C_0$ of the stage count, $C_0$ being the
constant of that count, is unchanged. All of these
ingredients are independent of $V$.

The second and third steps of the collar-expansion proposition are read through the two-slot
machinery as follows. These steps are the norm-form second-stage proposition: its polymers are
sets of fluctuation supports only, and it reads its elements only through their data
$(\mathsf{b},\mathsf{d},\nu)$ in the sense of that proposition. It is applied here to the
elements of the two-slot machine clause and of the two-slot M\"obius lemma. Their M\"obius
inversion runs over the response cubes and over the $A_j$-cubes of the second slot of
\cite[(1.60)]{B-CMP122a}. The remaining inputs are as follows:
\begin{itemize}
\item the exterior-slot polydisc lemma and its corollary carry the exterior slots $A_{j'}$,
$j'>j$, as polydisc parameters;
\item the localised all-blocks dilation corollary supplies the dilation of all blocks;
\item the stored-term datum $A^{\flat}$ is supplied at its creation by the creation corollary
for the stored-term datum;
\item the resulting bound is the one stated in the anchored output corollary.
\end{itemize}

The second step consumes the elements $\mathbb{V}(Y)$ of the norm-form second-stage
proposition only through their holomorphy and the bound
\begin{equation*}
|\mathbb{V}(Y)|\le\mathfrak{N}(\cdot)\,
\exp\Bigl[\kappa_A^{\mathrm{pr}}(g_k)|\mathcal{S}_k|
+\kappa_A(k)\sum_{j'<k}|\mathcal{S}_{j'}|\Bigr]\,e^{-\Theta(Y)},
\end{equation*}
where $\Theta(Y)$ is the decay exponent of that proposition. The weight in this bound
satisfies
\begin{equation*}
\mathfrak{N}(\cdot)\,
\exp\Bigl[\kappa_A^{\mathrm{pr}}(g_k)|\mathcal{S}_k|
+\kappa_A(k)\sum_{j'<k}|\mathcal{S}_{j'}|\Bigr]
\ \ge\ \mathfrak{N}(\cdot)\,e^{\kappa_A(k)|\mathcal{S}|}.
\end{equation*}
This is the weight that the creation corollary for the stored-term datum delivers: the
newest slot at $\kappa_A^{\mathrm{pr}}(g_k)$, the older slots at $\kappa_A(k)$. Every later
use takes the weight $e^{\kappa_A(k)|\mathcal{S}|}$, which the delivered weight dominates.

\emph{Analyticity in $\lambda$.} The integrand is entire in $\lambda$. Let $r_\lambda>0$ be the
margin in the activity constant which clause (a) of the Schwarz-lemma estimate takes as its
input. We apply that clause to
$\lambda\mapsto t_\lambda$ on the disc $|\lambda|<r_\lambda/\bar\varepsilon'_m$. Together with
the marked-cluster corollary, this gives the bound on the marked sub-sum $t_1-t_0$.
\end{proof}

\begin{proposition}[The table of units and their sizes]\label{5.2:prop-unit-table}
Fix a level $k$. Let the source be either a complex source $w\in\mathcal{W}_k$ or the constant bare
shift $\zeta$. For a point $y$ let $\tilde\square\in\pi_n$ be the current-zone cube of $y$ and
$\square\in\pi_j$ its $j$-cube, and put $\nu:=1+n-j$ and $t:=L^{j-n}$. In what follows
$\rho:=\rho_{n,j}$. A point $y$ is called \emph{near} if $\square^{\sim\nu+1}\cap\partial Z\neq\emptyset$,
and \emph{far} otherwise. Assume the following.
\begin{itemize}
\item[(a)] The inductive hypothesis of Theorem~\ref{5.1:thm-correlation-theorem} at level $k$ for the
chosen source, $w$ or the constant $\zeta$.
\item[(b)] The frozen-group lemma, and the cut-group lemma in its parts (i)--(iv), with the constant
$C_\partial'$ of \eqref{5.2:eq-cut-groups-a}.
\item[(c)] The freezing lemma, with constant $C_B$, and the marginal-difference lemma in its parts
(i)--(iii), with constant $C_M$.
\item[(d)] Part (b) of the weight-free expansion lemma, and the auxiliary estimate entering the
counterterm piece of class~$\partial$.
\item[(e)] Clause (f1) of the definition of the families of effective-action terms.
\item[(f)] The slot-parameter lemma. It carries the exterior slots of the stored terms as holomorphic
polydisc parameters, together with the rotation presented at the site.
\item[(g)] The first-stage bound by two contour integrals, Lemma~\ref{5.2:lem-first-stage-bound}. It is
used for the units of class~4 in the form \eqref{5.2:eq-three-sites-a}, at the radius
$\rho_A=\mathsf{w}/\eta_p^{1/2}$.
\item[(h)] The creation corollary for stored boundary terms. It supplies the stored datum of
\eqref{5.2:eq-stored-terms-a} and the bound \eqref{5.2:eq-created-terms-bound-a}.
\item[(i)] The output corollary of the preparatory stages, together with the inheritance theorem for
the domains $\mathfrak{D}^{\flat}_{l+1}(X)$ of the preparatory stages $l$.
\item[(j)] The orbit-datum bound, clauses (a)--(c).
\item[(k)] The containment theorem for arrested pieces.
\item[(l)] The label bound for the boundary terms born in the operation $\mathbf{R}$.
\item[(m)] The reference-flow comparison.
\item[(n)] The chain-stage counting lemma and the nesting lemma for live arrests.
\item[(o)] Summation of boundary terms at finer zones (Theorem~\ref{5.4:thm-finer-zone-sum}), and the
stage-count bound at complex fluctuation fields.
\end{itemize}
Then at each of the three sites (T), (P), (L) of Definition~\ref{5.1:def-expansion-rule} the units $v$
fall into the six classes described below. For
each class, the second column of \eqref{5.2:eq-unit-table-a} gives the size $\mathfrak{N}(v)$, and the
third column gives the contribution per unit $n$-cell, summed over $j$, except in class~4, where it is
the accumulated sum of the sizes of the class-4 units meeting a point, per level-$m$ cell.
The statements supplying each entry are named class by class in the proof.
\begin{equation}\label{5.2:eq-unit-table-a}
\begin{array}{lll}
 & \mathfrak{N}(v) & \text{per unit $n$-cell, summed over $j$}\\[3pt]
1 & \begin{array}[t]{l}C_IE_0\nu^3\rho^2(1+\rho)\,t^{4+\hat\alpha};\\
 \text{a tail: }2E_0e^{-\kappa d_j(X)}\end{array}
  & C_IE_0S_I\\[2pt]
2 & \begin{array}[t]{l}2C_BE_0\,\omega_{k,j}(C_MD_Xt^2\rho)^2\\
 \quad\cdot e^{-(1-\frac14\delta)\kappa d_j(X)}\end{array} & CE_0S_\omega\\[2pt]
3 & 2g_j^{\kappa_0}(C_MD_Xt^2\rho)^2e^{-\kappa d_j(X)}
  & C\sum_j\rho_{n,j}^2g_j^{\kappa_0}\le C\mathfrak{m}_n\\[2pt]
4 & \begin{array}[t]{l}2\,\mathfrak{N}^{\mathbb{C}}(\cdot)\ (\mathbf{B},\mathbb{C});\\
 2\,\sup|\cdot|\ (V^{(k)})\end{array}
  & \le 6(\mathsf{N}_0(\mathsf{T}_m)+3)C_0^{\sharp}\\[2pt]
5 & \tfrac83rC_S(2M)^4\mathfrak{a}_n^2;\ \ CL^4r|B'|^2 & \le Crg_n^2(\log g_n^{-2})^{2q}\\[2pt]
\partial & \begin{array}[t]{l}
 X\subset\tilde\square^{\sim2}:\\
 \quad 32C_M^2M^2E_0\rho^2t^4(d_j(X)+5)^2e^{-\kappa d_j(X)}\\
 X\not\subset\tilde\square^{\sim2}:\ 2E_0t^4e^{-(1-\theta_\nu)\kappa d_j(X)}\\
 \text{counterterm: }(b_++\tfrac23r)\cdot4\mathfrak{a}_n^2t^4\end{array}
 & \begin{array}[t]{l}\le 27d\bigl[C_\partial'M^3E_0S_\partial^{\sharp}\\
 \qquad+21M(b_++r)\mathfrak{a}_n^2\bigr]\end{array}
\end{array}
\end{equation}
In class $\partial$ we put $\theta_1:=0$ and $\theta_\nu:=2/15$ for $\nu\ge2$. The sums in the third
column are
\begin{equation}\label{5.2:eq-unit-table-1}
\begin{gathered}
S_I:=\sum_\nu\nu^3\rho^2(1+\rho)L^{-\hat\alpha(\nu-1)},\qquad
S_\omega:=1+\sum_{a\ge1}L^{-a/2}\rho^2_{\cdot,\cdot-a},\\
S_\partial^{\sharp}:=\sum_\nu(\nu+2)\rho^2_{n,n+1-\nu}L^{-(\nu-1)}.
\end{gathered}
\end{equation}
Here $\rho^2_{\cdot,\cdot-a}$ is the squared radius ratio between two levels $a$ steps apart.

The six classes are the following.
\begin{itemize}
\item[1.] The frozen group $\mathfrak{i}^{(j)}_y$ at the frozen source value $\hat w=w(p_y)$. Its core
consists of the terms with domain $X\ni y$, where $X$ lies in $2^d$ current cubes at the sites (T) and
(P), and $X\subset\square^{\sim\nu}$ (the domain of the frozen-group lemma) at the site (L) for far
points $y$. Together with the core goes the counterterm piece $b_j(\hat w)A(h_y,\cdot)$, and all of
these are taken on one background. The tails are taken singly: each is a unit of its own, on any
admissible background.
\item[2.] The differences $D^{(j)}(X,\cdot,y;w)-D^{(j)}(X,1,y;w)$.
\item[3.] The differences $R^{(j)}(X,\cdot;w)-R^{(j)}(X,1;w)$.
\item[4.] The boundary terms $\mathbf{B}^{(j)}$ and the chain terms $\mathbb{C}^{(l)}_k$ of the
preparatory stages $l$, stored with the slot datum of \eqref{5.2:eq-stored-terms-a}. This class also
contains:
\begin{itemize}
\item the terms $\mathbf{B}'^{(j)}$ born in $\mathbf{R}$, with their wrapped descendants born in
$\mathbf{T}$ (wrapped units carry no localisation datum and no amplitude $\mathsf{t}$ at the site (L));
\item the arrested left-overs $\mathcal{C}_{\mathrm{arr}}$, namely the frozen
$\mathbb{C}^{(i)}_{k_1}$ with $k_1<k$ and the change-of-variables terms of the arrest construction
\cite{B-CMP122a};
\item the $V^{(k)}$-terms.
\end{itemize}
The live exterior slots of every unit of this class are read at the rotation that the site presents:
$\mathbf{A}^{[\mathfrak{Q}^{\mathfrak{b}}(\zeta,\mathsf{t},\sigma)]}$ at (L), and
$\mathbf{A}^{[\mathfrak{u}\cdot\mathsf{g}]}$ at (T) and (P), where $\mathsf{g}$ is the rotation
variable of the orbit datum and $\mathfrak{u}$ is the local gauge factor of the response; at real
data $\mathfrak{u}$ is the Landau-to-axial gauge map, under which the slots transform. In other
words, each such unit is its orbit datum, and
it is bounded by $\mathfrak{N}^{\mathbb{C}}$. In this class $\mathfrak{N}^{\mathbb{C}}$ is the
supremum of the orbit datum, weighted with the weights of \eqref{5.2:eq-created-terms-bound-a}. For the
terms $\mathbf{B}'^{(j)}$ born in $\mathbf{R}$ the size is taken on the label $A$, and not at the plain
length of the true domain $X_F$:
\begin{equation}\label{5.2:eq-unit-table-2}
\mathfrak{N}^{A,\mathbb{C}}\bigl(\mathbf{B}'^{(j)}(A,\cdot)\bigr)
\le e^{\lambda_\theta}c_1^{\flat}(g_j)\,e^{-(1+\frac14\beta_s)\kappa d_j(A)} .
\end{equation}
\item[5.] The terms $f_\square:=A(\zeta_k\hat\chi_\square,\cdot)$ and $\zeta_k(p_y)G_y(B')$.
\item[$\partial$.] This class occurs at the site (L) only, for near points $y$. It consists of each
constituent $E^{(j)}_{\hat w}(X,\cdot,y)-E^{(j)}_{\hat w}(X,1,y)$ with $X\ni y$, taken singly as a unit
of its own on its own background, and of the counterterm piece $b_j(\hat w)A(h_y,\cdot)$, also taken
singly.
\end{itemize}
Finally, $S_I$, $S_\omega$ and $S_\partial^{\sharp}$ are independent of $K$ and bounded uniformly in
$g\le\gamma_J$. The only sum in \eqref{5.2:eq-unit-table-a} without a geometric weight is
$\sum_j\rho^2_{n,j}g_j^{\kappa_0}$; it is governed by a separate condition among the admitted inputs
of Definition~\ref{5.1:def-admitted-inputs}.
\end{proposition}

\begin{proof}
We first establish the second column of \eqref{5.2:eq-unit-table-a}, class by class, from hypotheses
(a)--(o). We then derive the third column and the uniformity.

\emph{Class 1.} For the core together with the counterterm piece on one background, the frozen-group
lemma applied under hypothesis (a) gives the size $C_IE_0\nu^3\rho^2(1+\rho)t^{4+\hat\alpha}$. A tail,
taken singly on an admissible background, has size $2E_0e^{-\kappa d_j(X)}$. Part (i) of the cut-group
lemma shows that the factor $t^5$ is extracted after the weight.

\emph{Class 2.} The freezing lemma and the marginal-difference lemma together give the class-2 entry of
the second column of \eqref{5.2:eq-unit-table-a} as a bound on the difference of the
$D$-terms.

\emph{Class 3.} The marginal-difference lemma gives the class-3 entry of the second column of
\eqref{5.2:eq-unit-table-a} for the difference of the $R$-terms.

\emph{Class 4.} The size of a boundary or chain term is twice $\mathfrak{N}^{\mathbb{C}}$, the supremum
of its orbit datum. By clauses (a)--(c) of the orbit-datum bound, this supremum dominates the
real-slice size, with the same constants. The bound carries two weights:
\begin{itemize}
\item the strong weight $e^{\kappa_A^{\mathrm{pr}}(g_{j_0})|\mathcal{S}_{j_0}|}$ on the newest slot,
whose level we denote by $j_0$;
\item the weight $e^{\kappa_A(k)|\mathcal{S}_{j'}|}$ on the older slots.
\end{itemize}
These are the weights that the induction delivers in \eqref{5.2:eq-created-terms-bound-a}. The
individual terms are bounded as follows.
\begin{itemize}
\item A created boundary term $\mathbf{B}^{(j)}$ with domain $X$ is bounded by
$B_0e^{-\kappa_{\mathrm{out}}d_j(X)}$, by \eqref{5.2:eq-created-terms-bound-a}, with the constant
$B_0$ and the tree length $d_j(X)$ at the term's own scale $j$ as in that display. Its slot datum is
supplied at creation by the creation corollary.
\item A chain term satisfies $\|\mathbb{C}^{(l+1)}\|^{A}_{(l+1)}\le C_0^{\sharp}$ by the output
corollary of the preparatory stages. This holds at the two rates of that corollary, written there
$a_k:=(1+3\beta_{\mathrm{pr}})\kappa$ and $a_m:=(1-\beta_{\mathrm{pr}})\kappa$ (the subscripts are
labels, not levels), on the output corollary's own configuration space
$\tilde U^{(l+1)c,\Sigma}_k(X)$ over $X$, and also on the inherited domain
$\mathfrak{D}^{\flat}_{l+1}(X)$ of the inheritance theorem. The latter is the domain read at the site
(L).
\item A frozen $\mathbb{C}^{(i)}_{k_1}$ in $\mathcal{C}_{\mathrm{arr}}$ carries the same datum and the
same constant $C_0^{\sharp}$, at the rate $a_m$. This is clause (a) of the proposition on the operation
$\mathbf{R}$ at step $k_1$, read on its frozen arrest domain. Its containment in the configuration
spaces on which the live arrested pieces are read is given by the containment theorem for arrested
pieces.
\item A frozen change-of-variables term has size at most $1$: it is ``bounded by any positive power of
$g_k$'' \cite{B-CMP119}. Its domain is that of the arrest construction.
\item A term $\mathbf{B}'^{(j)}$ born in $\mathbf{R}$ is bounded on its label $A$, and not at the plain
length of its true domain $X_F$, by \eqref{5.2:eq-unit-table-2}.
This is the label bound of hypothesis (l). It rests on the decay factor
$\exp(-(1+\tfrac12\beta)\kappa d_{k,\cup Y_i}(X))$ of \cite[(1.99)]{B-CMP122b}, in Ba{\l}aban's
notation there ($\beta=\beta_{\mathrm{pr}}$, the $Y_i$ being the large-field regions of that formula,
$d_{k,Z}$ the tree length off $Z$), and on the skeleton lemma, under the condition
$\tfrac14\beta_s\kappa\ge\lambda_\theta$: the quarter of $\beta_s$ pays for the fibres of the skeleton
$\operatorname{sk}$. It further uses clause (f2) of the relative-term statement at level $k$,
clause (e) of the proposition on the operation $\mathbf{R}$, and clause (d)$'$ of the inductive scheme
of the correlation theorem.
\item A term $V^{(k)}$ is bounded trivially by twice its supremum norm.
\end{itemize}
The unit read at the presented rotation is holomorphic in the slots. This follows from clause (f1)
(hypothesis (e)) and the slot-parameter lemma, which carries the presented rotation through the
definition of the orbit datum and its rotation lemma.

At the sites (T) and (P), the first-stage bound is applied at $\rho_A=\mathsf{w}/\eta_p^{1/2}$, in the
form \eqref{5.2:eq-three-sites-a}. At the site (L), the hypothesis of the first-stage bound is read at
$\mathsf{t}=0$ only. The $\mathsf{t}$-pieces of this class are summed in bundles by
Theorem~\ref{5.4:thm-finer-zone-sum}, and the gain $\mathsf{w}_L^{-1}\eta_p^{1/2}$ is not used on
this class at (L).

For the third column of class 4 we use the geometry of \cite[(2.47)]{B-CMP119}, and we count the
accumulated terms at each unit's own level. By the chain-stage counting lemma, at most $N_0+3$ of the
stage families of one chain (in the sense of that lemma) meet a given point. The sum over $X$ of each
family at an anchor cube (again in the sense of that lemma) is its norm, which is at most
$C_0^{\sharp}$. For one chain the contribution per cell at a point is at most
$3(N_0+3)C_0^{\sharp}$, which, with $N_0\le\mathsf{N}_0(\mathsf{T}_n)$ by the reference-flow
comparison, is the second member $3(\mathsf{N}_0(\mathsf{T}_n)+3)C_0^{\sharp}$ of
$\mathsf{K}_{\mathrm{cell}}$ used below. This applies to the chain of the present step. It applies in
the same way, at step $k_a$, to the frozen chain of each live arrest $\mathfrak{a}$ containing the
point.

The chains at a point are the present step's attempt and the nested live arrests, by the nesting lemma
for live arrests. Their attempt ranges (in the sense of the nesting lemma) are pairwise disjoint except
for one shared index. Hence at most
two of them have a stage family at a given level $m$. By the reference-flow comparison,
$N_0(g_{k'})\le\mathsf{N}_0(\mathsf{T}_m)$ for the step $k'\in\{k,k_a\}$ of the chain concerned and
every $m\le k'$. The contribution per level-$m$ cell at a
point is therefore at most
\begin{equation*}
2\cdot3(N_0+3)C_0^{\sharp}\le 6(\mathsf{N}_0(\mathsf{T}_m)+3)C_0^{\sharp},
\end{equation*}
and no further accumulation and no tail enters.

For the zone-cell counts of the volume bound, the chain of the present step is charged once per term,
at a cell of its own stage region. This is the second member $3(\mathsf{N}_0(\mathsf{T}_n)+3)C_0^{\sharp}$
of $\mathsf{K}_{\mathrm{cell}}$, by clause $(\alpha)$ of the per-zone-cell bookkeeping statement. The
frozen terms of the live arrests are charged extensively as $\mathsf{P}(\mathfrak{a})$, by clause
$(\beta')$ of the same statement. At $|\mathsf{B}|<T^{\mathbb{C}}$ the count is paid by distance at half
rate, by the stage-count bound at complex fluctuation fields with $\eta_p^{1/2}$.

\emph{Class 5.} Part (b) of the weight-free expansion lemma gives the size
$\tfrac83rC_S(2M)^4\mathfrak{a}_n^2$ for $f_\square$. It gives $CL^4r|B'|^2$ for $\zeta_k(p_y)G_y(B')$.

\emph{Class $\partial$.} Here the bounds are provided by parts (i)--(iii) of the marginal-difference
lemma, part (b) of the weight-free expansion lemma together with the auxiliary estimate entering the
counterterm piece, and parts (ii)--(iv) of the cut-group lemma. For a constituent with
$X\subset\tilde\square^{\sim2}$ they give
the first bound of the class-$\partial$ entry in the second column of \eqref{5.2:eq-unit-table-a}. For
a constituent with $X\not\subset\tilde\square^{\sim2}$ they give the second bound of that entry. For the
counterterm piece they give its third bound.

\emph{The third column.} A unit $n$-cell contains $t^{-4}$ points $y$, respectively at most $t^{-4}$
cubes. This factor cancels the factor $t^4$ in classes 1--3 and 5.

Since $t=L^{j-n}=L^{-(\nu-1)}$, the class-1 size multiplied by $t^{-4}$ is
\begin{equation*}
t^{-4}\cdot C_IE_0\nu^3\rho^2(1+\rho)t^{4+\hat\alpha}=C_IE_0\nu^3\rho^2(1+\rho)L^{-\hat\alpha(\nu-1)}.
\end{equation*}
Summing over $j$, that is over $\nu$, gives
$C_IE_0S_I$ with $S_I$ as in \eqref{5.2:eq-unit-table-1}.

In class 2, $t^{-4}(C_MD_Xt^2\rho)^2=C_M^2D_X^2\rho^2$, and the sum over $j$ is the entry
$CE_0S_\omega$. In class 3 the same cancellation leaves $C\sum_j\rho^2_{n,j}g_j^{\kappa_0}\le
C\mathfrak{m}_n$. The class-5 entry is $\le Crg_n^2(\log g_n^{-2})^{2q}$. The class-$\partial$ entry is
the one in the third column of \eqref{5.2:eq-unit-table-a}, with $S_\partial^{\sharp}$
as in \eqref{5.2:eq-unit-table-1}. The class-4 entry was obtained above.

\emph{Uniformity.} By the reference-flow comparison,
\begin{equation*}
\rho^2_{n,j}\le(1+o(1))\bigl(1+(3\lambda(n-j)+2c_2)g^2\bigr),\qquad n-j=\nu-1 .
\end{equation*}
This bound grows at most linearly in $\nu$, respectively in $a$, and holds for all $g\le\gamma_J$. The
weights $L^{-\hat\alpha(\nu-1)}$, $L^{-a/2}$ and $L^{-(\nu-1)}$ in \eqref{5.2:eq-unit-table-1} are
geometric. Hence the series $S_I$, $S_\omega$ and $S_\partial^{\sharp}$ converge, with bounds
independent of $K$ and uniform in $g\le\gamma_J$. The remaining sum $\sum_j\rho^2_{n,j}g_j^{\kappa_0}$
of class 3 carries no such weight; it is governed by a separate condition among the admitted inputs of
Definition~\ref{5.1:def-admitted-inputs}.
\end{proof}

At site (L) the effective action is split term by term in \cite{B-CMP122b}:
``We divide the terms in the second exponential in (1.49) into two parts: the first part
includes the first three terms, and these terms of the rest of the effective action, for
which their localization domains intersect the region $Z$, the second part includes the
remaining terms, i.e., the terms with localization domains disjoint with $Z$''. The
brackets \cite[(1.59)]{B-CMP122b} are formed for pairs $(X,y)$ with $X\subset Z^c$:
``Consider a term $E^{(j)}(X,U''_k,z)$ with $X\subset Z^c$, and let
$z\in\Omega_n\setminus\Omega_{n+1}$ for $n\geq j$'' (Ba{\l}aban's point $z$ is our point
$y$; the letter $z$ occurs below only inside quotations, in this sense). Being cut is
therefore a property of the pair $(X,y)$, and not of the point $y$. Near
$\partial Z$ some domains $X\ni y$ run through the chain \cite[(1.51)--(1.58)]{B-CMP122b}
and belong to the second part. The others are first-part terms, treated in
\cite[(1.61)--(1.65)]{B-CMP122b}. Later, at the final background $U_k$, Ba{\l}aban adds
``all the missing terms with localization domains intersecting the domain $Z$ \dots\ and we
subtract them from the remaining expression'' \cite{B-CMP122b}. The frozen groups of the
table in Proposition~\ref{5.2:prop-unit-table} need the core and the counterterm piece on
one background; this is the closing remark of the core estimate. The lemma below says at
which points $y$ this holds on both sides of $\partial Z$, and it bounds everything else
singly on a thin layer about $\partial Z$.

Throughout, $Z$ is the region of \cite{B-CMP122b} at site (L), level $k$. It is a union
of connected components of $\Lambda_k^c$: ``Each term in the expansion (2.18) \dots\ has a
large field region $\Lambda_k^c$. It is a union of connected components \dots\ Let us
denote the union of the above class of components by $Z$'' \cite{B-CMP122a}. Since a
component shares no face with another component, $\partial Z\subset\partial\Lambda_k$.

The set $\Lambda_k$ is a union of $M\check{R}_k$-cubes of the current unit lattice. In
\cite{B-CMP119}, $\Lambda_j^{\circledR}$ denotes ``the set which is obtained by removing one
layer of the $MR_j$-cubes from $\Lambda_j$''. Further, ``Each renormalization step adds at
least ten layers of $MR_k$-cubes'' \cite{B-CMP122a}. All regions are unions of cubes of
the nested $L$-adic partitions: ``the domains $\Omega_j^c$ are unions of
$L^{-(k-j)}MR_j$ cubes \dots\ the complements of these sets form an admissible sequence
of domains based on partitions into $M$-cubes in the corresponding scales \dots\ all the
regions connected with the last $N$ steps are rectangular parallelepipeds''
\cite{B-CMP122a}.

Hence $\partial Z$ lies on coordinate hyperplanes $\{x_\mu=c\}$ of the $M$-grid of the
current scale. In the units of any zone $n\le k$ these hyperplanes are spaced at least
$ML^{k-n}\ge M$ apart, that is, for any $j\le n$, at least $M/t$ apart in $j$-units, with
$t=L^{j-n}$ as in the lemma below.
None of this depends on the complex source $w=\sum_i z_i f_i$.

\begin{lemma}[Cut groups at the boundary of the large-field region]\label{5.2:lem-cut-groups}
Work in the setting of Proposition~\ref{5.2:prop-unit-table} at site (L), level $k$, with
$Z$ as above. Fix $j\le n\le k$, and put
\[
t=L^{j-n},\qquad \nu=1+n-j,\qquad \rho=\rho_{n,j},\qquad F=E^{(j)}_{\hat w},
\]
a family of size $\mathfrak{n}=E_0$, with $E_0$ the size bound of the frozen families used
in the table of Proposition~\ref{5.2:prop-unit-table}. Let
$y\in\Lambda_j^{\circledR}\cap(\Omega_n\setminus\Omega_{n+1})$ be a point of $T^{(j)}$.
Let $\square\in\pi_j$ and $\tilde\square\in\pi_n$ be the cubes containing $y$.

Call $y$ \emph{far} if $\square^{\sim(\nu+1)}\cap\partial Z=\emptyset$, so that the cube
$\square^{\sim(\nu+1)}$ lies in $Z$ or in $Z^c$. Call $y$ \emph{near} otherwise. This
marking does not depend on the complex source $w=\sum_i z_i f_i$: it is the same for all
values of the $z_i$.

Take as the \emph{core} of the frozen group $\mathfrak{i}^{(j)}_y$ the family used in the
core estimate, $\{X\ni y:\ X\subset\square^{\sim\nu}\}$, together with the piece
$b_j(\hat w)A(h_y,\cdot)$. Every other constituent $X\ni y$ is called a \emph{tail} and is
a unit of its own. Below, $1$ denotes the trivial configuration. Then:
\begin{enumerate}
\item[(i)] (Far points.) If $y$ is far, the core, the box $\square_0=\square^{\sim(\nu+1)}$
of the core estimate and $\operatorname{supp}h_y$ lie on one side of $\partial Z$.
Ba{\l}aban's split sends the whole core and the counterterm piece to one and the same
part. There they are functions of one background along one chain. Every constituent that
the split sends to the other part is a tail. So is every constituent of a subtracted copy
at a far point $y\in Z^c$.

The core and the counterterm piece obey the core estimate
\begin{equation}\label{5.2:eq-cut-groups-1}
\Bigl|\sum_{X\ni y,\ X\subset\square^{\sim\nu}}\bigl[F(X,\mathfrak{b},y)-F(X,1,y)\bigr]
-\beta[F]A(h_y,\mathfrak{b})\Bigr|
\le C_IE_0\nu^3\rho^2(1+\rho)\,t^{4+\hat\alpha}
\end{equation}
on every admissible background $\mathfrak{b}$. Along holomorphic families they obey the
form of clause (iii) of the background-variation lemma.

Each tail obeys
\begin{equation}\label{5.2:eq-cut-groups-2}
|F(X,\mathfrak{b}',y)-F(X,1,y)|\le 2E_0e^{-\kappa d_j(X)}=:\mathfrak{N}(\mathrm{tail})
\end{equation}
on whatever admissible background $\mathfrak{b}'$, whether or not $X$ is present in the
part at hand. The factor $t^5$ is extracted after the resummation weight (see the proof).
These are the sizes of the frozen groups in the table of
Proposition~\ref{5.2:prop-unit-table}.

\item[(ii)] (Near points.) Take arbitrary admissible backgrounds at scale $n$, one for each
constituent. Then
\begin{align*}
|F(X,\mathfrak{b}_X,y)-F(X,1,y)|&\le
32C_M^2M^2E_0\rho^2t^4(d_j(X)+5)^2e^{-\kappa d_j(X)}
&&\text{if } X\subset\tilde\square^{\sim 2},\\
|F(X,\mathfrak{b}_X,y)-F(X,1,y)|&\le 2E_0t^4e^{-(1-\theta_\nu)\kappa d_j(X)}
&&\text{if } X\not\subset\tilde\square^{\sim 2},
\end{align*}
where $\theta_1:=0$ and $\theta_\nu:=2/15$ for $\nu\ge 2$. The counterterm piece obeys
\[
|b_j(\hat w)A(h_y,\mathfrak{b})|\le 4\bigl(b_++\tfrac23 r\bigr)\mathfrak{a}_n^2t^4 .
\]
Now sum over $X\ni y$ with the resummation weight of site (L). This is the weight
$e^{(1+3\beta_s)\kappa d_k(\bar X)}$ of \cite[(1.66)]{B-CMP122b}, where $\bar X=Y_0$ is the
closure of $X$ at the current scale: ``The inductive assumption yields the factor \dots\
$\exp(-\beta\kappa d_j(X))\exp(-(1+3\beta)\kappa d_k(Y_0))$ for $j<k$, where $Y_0$ is the
smallest domain from $\mathfrak{D}_k$ containing $X$, and $\exp(-(1+4\beta)\kappa d_k(X))$
for $j=k$, from the improved bounds after the $k$-th renormalization transformation
$\mathbf{T}$'' \cite{B-CMP122b}. For $j=k$ the sizes carry the improved exponent
$(1+4\beta_s)\kappa$ of the level-$k$ families. This is the improved level-$k$ bound, which
rests on the decay-fraction lemma. Then
\begin{equation}\label{5.2:eq-cut-groups-a}
\begin{gathered}
\sum_{X\ni y}e^{(1+3\beta_s)\kappa d_k(\bar X)}
\sup|F(X,\mathfrak{b}_X,y)-F(X,1,y)|\le C_\partial'M^2E_0\rho^2t^4,\\
C_\partial':=(32C_M^2+8)\,C(\kappa),\qquad
C(\kappa):=\sup_{j,y}\sum_{X\in\mathfrak{D}_j,\ X\ni y}(d_j(X)+5)^2e^{-\frac14\delta\kappa d_j(X)} .
\end{gathered}
\end{equation}
The constant $C(\kappa)$ is finite and independent of $j$ and $K$. It is a sum of the type
of \cite[(1.26)]{B-CMP116},
``$\sum_{X\in\mathfrak{D}_j,\,X\supset\square'}\exp(-\kappa d_j(X))\le O(1)$ \dots\ for
$\kappa$ sufficiently large'', taken at the rate $\frac14\delta\kappa$. This is the rate at
which the units of the second entry of the table are already summed. The other factor
$e^{-\frac14\delta\kappa d_j(X)}$ is kept.
\item[(iii)] (Count.) A unit $n$-cell contains at most $27dM(\nu+2)t^{-3}$ near points of
$T^{(j)}$.
\item[(iv)] (Sum.) Per unit $n$-cell, summed over $j\le n$, the near units contribute at
most
\[
27d\bigl[C_\partial'M^3E_0S_\partial^{\sharp}+21M(b_++r)\mathfrak{a}_n^2\bigr]
\]
to the last column of the table \eqref{5.2:eq-unit-table-a}. Here
\begin{equation}\label{5.2:eq-cut-groups-3}
S_\partial^{\sharp}:=\sum_{\nu\ge 1}(\nu+2)\rho^2_{n,n+1-\nu}L^{-(\nu-1)}
\le(1+o(1))\sum_{\nu\ge1}(\nu+2)\bigl(1+(3\lambda\nu+2c_2)g^2\bigr)L^{-(\nu-1)} ,
\end{equation}
where $o(1)$ is the error term supplied by the reference comparison of the couplings.
The sum $S_\partial^{\sharp}$ is independent of $K$ and bounded uniformly in
$g\le\gamma_J$.
\end{enumerate}
\end{lemma}

\begin{proof}
\emph{Part (i): the supports.} By construction $y\in\square$. In $j$-units,
$\operatorname{supp}h_y\subset y+[-1,1]^4$. Indeed, the functions $h_y$ are those of
\cite[(3.40)]{B-CMP119}, ``$h_z(x)=\prod_\mu h(x_\mu-z_\mu)$ \dots\ They form a
decomposition of unity''. As in \cite{B-CMP119}, they are taken ``on the lattice
$T_{L^{-j}}$, and with $h(t)=\max\{1-|t|,0\}$'' (in this quotation $t$ is a real
variable). Hence
\[
\square^{\sim(\nu+1)}\supset\square^{\sim\nu}\cup\operatorname{supp}h_y .
\]
If $y$ is far, $\square^{\sim(\nu+1)}$ lies in $Z^c$ or in $Z$. Together with the display,
this gives the first assertion of (i). We treat the two cases in turn.

\emph{Part (i): the case $\square^{\sim(\nu+1)}\subset Z^c$.} Every core constituent
$X\subset\square^{\sim\nu}$ is disjoint from $Z$, and so is $\operatorname{supp}h_y$.
So all of them belong to the second part. The data of the chain \cite[(1.59)]{B-CMP122b}
hang on $\tilde\square$ and $\tilde\square_0=\tilde\square^{\sim5}$, not on $X$: ``we take
the cube $\tilde\square\in\pi_n$ containing the point $z$, and we consider the case
$X\subset\tilde\square^{\sim2}$ \dots\ where $\tilde\square_0=\tilde\square^{\sim5}$ (the
operation $\sim$ is defined here in terms of $M$-cubes of the $L^{-n}$-lattice)''
\cite{B-CMP122b}.

The core lies in $\tilde\square^{\sim2}$. Its points are within $(\nu+1)M$ $j$-units of
$y$, and
\[
(\nu+1)M\ j\text{-units}=(\nu+1)ML^{-(\nu-1)}\ n\text{-units}\le 2M\ n\text{-units}.
\]
So all core constituents belong to the class $X\subset\tilde\square^{\sim2}$ of
\cite[(1.59)]{B-CMP122b}. The piece $A(h_y,\cdot)$ has
$\operatorname{supp}h_y\subset\tilde\square^{\sim1}\subset\tilde\square_0$. Together with
this piece, the core constituents are functions of one background along the chain. They
therefore form one unit in the sense of the first-stage bound
(Lemma~\ref{5.2:lem-first-stage-bound}).

\emph{Part (i): the case $\square^{\sim(\nu+1)}\subset Z$.} Every $X\ni y$ meets $Z$,
because it contains $y$, and $\operatorname{supp}h_y\subset Z$. So all of these belong to
the first part, and nothing lies elsewhere. The set $\square^{\sim(\nu+1)}$ is connected
and contained in $Z$. Hence the core lies inside the component $C$ of $Z$ containing $y$.

For a core constituent $X$, the only component of $Z$ meeting $X$ is $C$, and $X\subset C$.
Ba{\l}aban takes the smallest localization domain $Y_1$ ``containing $X$, and containing
these components of $Z$, which intersect $X$'' \cite{B-CMP122b}. For every core $X$ this
is therefore the smallest localization domain containing $C$. The same holds for the piece,
since $\operatorname{supp}h_y\subset C$. Call this common domain $Y_1(C)$. The
representation \cite[(1.61)]{B-CMP122b} and its successors are then common to the core and
the piece, so again they form one unit.

\emph{Part (i): the subtracted copies.} The subtracted copies are the pairs $(X,y)$ with
$X\cap Z\ne\emptyset$, at the single background $U_k$.
\begin{itemize}
\item For far $y\in Z$ they are whole groups, and the core estimate applies to them as it
stands.
\item For far $y\in Z^c$ the core constituents do not meet $Z$. So every constituent of
such a copy is a tail.
\end{itemize}

\emph{Part (i): what the proof of the core estimate uses.} We now read the proof of the
core estimate.
\begin{itemize}
\item Its clause for the levels $n$ close to $j$ is proved term by term.
\item The tail class $X\not\subset\square^{\sim\nu}$ is estimated term by term. The
estimate uses $\sup|F(X,\cdot,y)|\le E_0e^{-\kappa d_j(X)}$, $d_j(X)\ge\nu-1$ and
$\kappa\ge10\log L$. This is Ba{\l}aban's own second class: ``For $X$ in the second class
the difference of the corresponding expressions can be estimated by
$2E_0\exp(-\frac12\kappa(n-j))\exp(-\frac12\kappa d_j(X))$, and
$\exp(-\frac12\kappa(n-j))<(L^jL^{-n})^5$, hence these terms can be treated as before''
\cite{B-CMP119}. This estimate sees neither the background on which a tail sits nor
whether the tail is present.
\item The estimate of the core reads the background on $\square_0$ only. It uses the three
conditions on the background in the core estimate, the first of them with Ba{\l}aban's
one-layer margin, and the comb gauge centred at $y$.
\item The step that extends $\sum_X$ to all $X\ni y$ of $L^{-j}\mathbb{Z}^4$ concerns the
second-order kernels $F^{(2)}(X,y)$ of the whole frozen family $E^{(j)}_{\hat w}$ at the
trivial configuration. These are contracted with the main term
$(x_\kappa-y_\kappa)F_{\kappa\mu}(y)$ (here $\kappa$ is a coordinate index), and they exist
regardless of $Z$.
\end{itemize}
Therefore the core estimate \eqref{5.2:eq-cut-groups-1} holds for the core and the
counterterm piece alone, with the constant $C_I$. The tails may be distributed over the
two parts, or be absent, in any way.

\emph{Part (i): the tail floor.} Let $X$ be a tail. Then $X$ is connected; in
\cite{B-CMP109}, ``two consecutive cubes have a common wall''. It contains the cube of $y$
and a cube outside $\square^{\sim\nu}$. So its shortest tree has $j$-length at least
$\nu M$. By the definition ``A length of a shortest graph in this class, divided by $M$, is
the linear size of $X$, and is denoted by $d_j(X)$'' \cite{B-CMP109}, we get
$d_j(X)\ge\nu$.

\emph{Part (i): extraction of $t^5$.} The factor $t^5$ is taken from
$\mathfrak{N}(\mathrm{tail})=2E_0e^{-\kappa d_j(X)}$ after the resummation weight. By the
argument in part (ii) below,
\[
d_k(\bar X)\le L^{-(k-j)}d_j(X)\le t\,d_j(X).
\]
Consequently:
\begin{itemize}
\item at sites (T), (P) the weight of the first-stage bound
(Lemma~\ref{5.2:lem-first-stage-bound}) is at most
$e^{(1-\frac12\delta)\kappa t d_j(X)}\le e^{\kappa td_j(X)}$;
\item at site (L), for $j<k$, the weight of \cite[(1.66)]{B-CMP122b} is at most
$e^{\frac74\kappa td_j(X)}$, using $k-j\ge n-j$.
\end{itemize}

Let $\nu\ge2$. Then $t\le\frac13$, since $L\ge3$, and we use $\kappa\ge10\log L$. What
remains after Ba{\l}aban's factor $t^5$ is extracted is as follows, site by site.

At sites (T), (P) we treat the tails of the frozen groups. They satisfy the floor
$d_j(X)\ge\nu-1$ of the tail class of the core estimate and of Ba{\l}aban's second class
quoted above ($n-j=\nu-1$). The core used at sites (T), (P) respects this floor: by the
tail-size hypothesis of that core, its tails satisfy
\[
d_j(X)>L^{\nu-1}/2-\sqrt d-1\ge\nu-1
\]
under the lower bound on $L$ assumed for that core. Since $1-t\ge\frac23$, and since
$d_j(X)\ge\nu-1$ gives
\[
e^{-\frac12\kappa d_j(X)}\le L^{-5d_j(X)}\le L^{-5(\nu-1)}=t^5,
\]
what remains is
\[
e^{-(1-t)\kappa d_j(X)}\le e^{-\kappa d_j(X)/6}e^{-\frac12\kappa d_j(X)}
\le t^5e^{-\kappa d_j(X)/6}.
\]

At site (L) we treat the tails of this lemma, which satisfy $d_j(X)\ge\nu$. What remains
is
\[
e^{-(1-\frac74t)\kappa d_j(X)}\le t^5e^{-\kappa d_j(X)/6}.
\]
This is checked in two cases.
\begin{itemize}
\item At $\nu=2$ we have $1-\frac{7}{4L}-\frac16\ge\frac14$, and, since $d_j(X)\ge2$,
\[
e^{-\frac14\kappa d_j(X)}\le L^{-5d_j(X)/2}\le L^{-5}=t^5 .
\]
\item At $\nu\ge3$ we have $t\le\frac19$ and $1-\frac7{36}-\frac16=\frac{23}{36}$.
Since $d_j(X)\ge\nu$,
\[
e^{-\frac{23}{36}\kappa d_j(X)}\le L^{-6d_j(X)}\le L^{-5(\nu-1)}=t^5 .
\]
\end{itemize}

For $\nu=1$, so $t=1$, nothing is extracted. What remains is:
\begin{itemize}
\item $e^{-\frac12\delta\kappa d_j(X)}$ under the weight of the first-stage bound;
\item $e^{-\frac5{12}\kappa d_j(X)}$ under the weight of \cite[(1.66)]{B-CMP122b} for
$j<k$, because $L^{-(k-j)}\le\frac13$;
\item $e^{-\frac14\kappa d_k(X)}$ for $j=k$.
\end{itemize}
In the last case the level-$k$ sizes carry the improved exponent
$(1+4\beta_s)\kappa=2\kappa$ of the improved level-$k$ bound, which rests on the
decay-fraction lemma; \cite{B-CMP122b} uses the same rate in the sentence quoted in (ii).
This exponent is set against the weight $e^{\frac74\kappa d_k(X)}$. This is one level, with
no scale sum. After the sum over $X$ it is at most $CE_0$ per cell, where $C$ is the bound
of the sum in \cite[(1.26)]{B-CMP116} taken at the rate $\frac14\kappa$; this contribution
is contained in the per-cell bound $C_IE_0S_I$ of the frozen groups in the table.

\emph{Part (ii): the first class, $X\subset\tilde\square^{\sim2}$.} Apply clause (i) of the
background-variation lemma (clause (iii) along families) with
$\mathfrak{n}=E_0e^{-\kappa d_j(X)}$ and $D_X\le4M(d_j(X)+5)$. This gives
\[
|F(X,\mathfrak{b}_X,y)-F(X,1,y)|
\le2E_0e^{-\kappa d_j(X)}\bigl(4C_MM(d_j(X)+5)t^2\rho\bigr)^2
=32C_M^2M^2E_0\rho^2t^4(d_j(X)+5)^2e^{-\kappa d_j(X)} .
\]
The box of that clause may meet $Z$. At site (L) the changes of the chain are ``made inside
the domain $\Lambda$'' only \cite{B-CMP122b}, and Ba{\l}aban's own box
$\tilde\square^{\sim5}$ at \cite[(1.59)]{B-CMP122b} may meet $Z$ as well.

\emph{Part (ii): the second class, $X\not\subset\tilde\square^{\sim2}$.} We use clause (ii)
of the background-variation lemma, the lower bound $d_j(X)\ge t^{-1}=L^{\nu-1}$ of
\cite[(3.48)]{B-CMP119}, and the trivial bound $2E_0e^{-\kappa d_j(X)}$. For $\nu\ge2$ we
have $L^{\nu-1}\ge3(\nu-1)$, since $L\ge3$, and $\kappa\ge10\log L$. So
\[
e^{-\theta_\nu\kappa d_j(X)}\le e^{-\frac2{15}\cdot10\log L\cdot3(\nu-1)}
=L^{-4(\nu-1)}=t^4 .
\]
Hence
\[
2E_0e^{-\kappa d_j(X)}
=2E_0e^{-\theta_\nu\kappa d_j(X)}e^{-(1-\theta_\nu)\kappa d_j(X)}
\le 2E_0t^4e^{-(1-\theta_\nu)\kappa d_j(X)}.
\]
For $\nu=1$ we have $t^4=1$ and $\theta_1=0$, and the bound is the trivial one.

\emph{Part (ii): the counterterm piece.} By the counterterm-coefficient bound,
$|b_j(\hat w)|\le b_++\frac23r$. Next,
$A(h_y,\mathfrak{b})=\sum_ph_y(p)a_p(\mathfrak{b})$. Since the $h_y$ form a
translation-covariant decomposition of unity, $\sum_ph_y(p)$ equals the number of
plaquettes of one $j$-cell. For each fine plaquette, $a_p(\mathfrak{b})$ is bounded by
part (b) of the Wilson-term lemma. That bound uses the regularity of the
admissible background on fine plaquettes, clauses (i) and (iii) of its defining tube, and
not the level-$j$ data of the first condition of the core estimate. Hence
$|A(h_y,\mathfrak{b})|\le4\mathfrak{a}_n^2t^4$.

The same bound holds for any weight $c$ with $|c|\le h_y$ in place of $h_y$. So it does
not matter how Ba{\l}aban files the piece: whole or plaquette by plaquette, with the Wilson
term (expanded linearly in the weight, by part (a) of the Wilson-term lemma) or with the
group.

\emph{Part (ii): the weighted sum.} We have $\bar X\supseteq X$. A shortest tree of $X$ has
$j$-length $Md_j(X)$ and lies in $X\subset\bar X$. It meets every current cube of
$\bar X$, since each of these contains a cube of $X$. Its length in current units is
$ML^{j-k}d_j(X)\le Mt\,d_j(X)$. By the definition of $d_k$ in \cite{B-CMP109},
\[
d_k(\bar X)\le L^{-(k-j)}d_j(X)\le t\,d_j(X).
\]
Hence, for $j<k$, with $\beta_s=\frac14$ and $L\ge3$, the weight of
\cite[(1.66)]{B-CMP122b} satisfies
\[
e^{\frac74\kappa L^{-(k-j)}d_j(X)}\le e^{\frac7{12}\kappa d_j(X)} .
\]
This weight is set against the exponential factor of each class.
\begin{itemize}
\item Against $e^{-\kappa d_j(X)}$ (the first class, and the second class at $\nu=1$) it
leaves at most $e^{-\frac5{12}\kappa d_j(X)}$.
\item Against $e^{-\frac{13}{15}\kappa d_j(X)}$ (the second class at $\nu\ge2$) it leaves
$e^{-\frac{17}{60}\kappa d_j(X)}$.
\item For $j=k$, so $\nu=1$, the level-$k$ sizes carry the improved exponent
$(1+4\beta_s)\kappa=2\kappa$ in place of $\kappa$. This is the improved level-$k$ bound,
resting on the decay-fraction lemma; see the sentence of \cite{B-CMP122b} quoted in (ii),
``$\exp(-(1+4\beta)\kappa d_k(X))$ for $j=k$, from the improved bounds after the $k$-th
renormalization transformation $\mathbf{T}$''. The weight $e^{\frac74\kappa d_k(X)}$ then
leaves $e^{-\frac14\kappa d_k(X)}$.
\end{itemize}
In every case the remainder is at most $e^{-\frac12\delta\kappa d_j(X)}$. This uses
$\delta<\frac13$, which holds because the rate $(1-3\delta)\kappa$ of
\cite[(2.18)]{B-CMP116} is positive. Of the remainder, the factor
$e^{-\frac14\delta\kappa d_j(X)}$ is summed by $C(\kappa)$ and the other factor
$e^{-\frac14\delta\kappa d_j(X)}$ is kept.

Summing over $X\ni y$, the first class contributes at most
$32C_M^2M^2E_0\rho^2t^4C(\kappa)$. The second class contributes at most
\[
2E_0t^4C(\kappa)\le8\rho^2E_0t^4C(\kappa)M^2 .
\]
Here $\rho\ge\frac12$: the function $u/(\log u)^{2q}$ is increasing, and $x_j\ge x_n-2c_2$
by the reference comparison of the couplings, so $\rho_{n,j}^2\ge1-o(1)$. We also use
$M\ge1$. The total is at most
\[
(32C_M^2+8)C(\kappa)M^2E_0\rho^2t^4=C_\partial'M^2E_0\rho^2t^4,
\]
which is \eqref{5.2:eq-cut-groups-a}.

\emph{Part (iii).} Let $y$ be near. Then a face of $\partial Z$ meets
$\square^{\sim(\nu+1)}$. All points of $\square^{\sim(\nu+1)}$ lie within sup-distance
\[
(2\nu+3)M/2+M/2=M(\nu+2)
\]
of $y$, in $j$-units. So $y$ lies in the slab $\{|x_\mu-c|\le M(\nu+2)\}$ of a grid
hyperplane $\{x_\mu=c\}$ carrying that face.

A unit $n$-cell has side $t^{-1}$ in $j$-units and $t^{-1}$ points of $T^{(j)}$ per axis.
One slab meets it in at most
\[
t^{-3}\bigl(2M(\nu+2)+1\bigr)\le3M(\nu+2)t^{-3}
\]
lattice points. The hyperplanes of one direction whose slab meets the cell have $c$ in an
interval of length $t^{-1}+2M(\nu+2)$, and they are spaced at least $M/t$ apart. So they
number at most
\[
1+\bigl(t^{-1}+2M(\nu+2)\bigr)\frac tM=1+\frac1M+2(\nu+2)L^{-(\nu-1)}\le9 .
\]
The last inequality uses $M\ge3$ and $L\ge3$; the middle expression is largest at $\nu=1$.
Over the $d$ directions this gives at most $27dM(\nu+2)t^{-3}$ near points. If
$27dM(\nu+2)t\ge1$, this exceeds the trivial count $t^{-4}$, and there is nothing to prove.
Faces of a finer grid do not occur, since $\partial Z$ lies on the grid of the current
scale.

\emph{Part (iv).} Combine (ii) with (iii). Per level, the near units contribute at most
\begin{align*}
&27dM(\nu+2)t^{-3}\bigl[C_\partial'M^2E_0\rho^2
+4\bigl(b_++\tfrac23r\bigr)\mathfrak{a}_n^2\bigr]t^4\\
&\qquad=27dM(\nu+2)\,t\,\bigl[C_\partial'M^2E_0\rho^2
+4\bigl(b_++\tfrac23r\bigr)\mathfrak{a}_n^2\bigr],
\end{align*}
with $t=L^{-(\nu-1)}$. Summing over $j\le n$ is summing over $\nu\ge1$. For $L\ge3$,
\[
\sum_{\nu\ge1}(\nu+2)L^{-(\nu-1)}\le\sum_{m\ge0}(m+3)3^{-m}=\frac92+\frac34=\frac{21}4 .
\]
We also have $4\cdot\frac{21}4=21$ and $\frac23r\le r$. The first term therefore sums to
$27dC_\partial'M^3E_0S_\partial^{\sharp}$, and the second is at most
$27d\cdot21M(b_++r)\mathfrak{a}_n^2$.

The bound \eqref{5.2:eq-cut-groups-3} on $S_\partial^{\sharp}$ follows from the reference
comparison of the couplings, as under the table:
$\rho^2_{n,j}\le(1+o(1))\bigl(1+(3\lambda(n-j)+2c_2)g^2\bigr)$ with $n-j=\nu-1\le\nu$.
The series on the right of \eqref{5.2:eq-cut-groups-3} converges geometrically for
$L\ge3$. It is independent of $K$ and bounded uniformly in $g\le\gamma_J$.

A unit may be anchored at a point $y$ away from the cell $\Delta$ that its support touches.
This anchoring is absorbed by the kept factor $e^{-\frac14\delta\kappa d_j(X)}$, because by
(ii) $C(\kappa)$ sums only at the rate $\frac14\delta\kappa$. The argument is the same as
for the first three entries of the table: an anchor at $j$-distance $s$ from $\Delta$
forces $d_j(X)\ge s/M-1$.
\end{proof}

\begin{remark}[Remarks on the cut groups]\label{5.2:rem-cut-group-remarks}
We keep the notation of the lemma on cut groups at the boundary of the large-field region
(Lemma~\ref{5.2:lem-cut-groups}): $j\le n\le k$, $t=L^{j-n}$, $\nu=1+n-j$, $F=E^{(j)}_{\hat w}$
with $\mathfrak{n}=E_0$, and $\rho=\rho_{n,j}$. Rows refer to the table of units and their
sizes \eqref{5.2:eq-unit-table-a}. The sites (T), (P), (L) are those of the first-stage bound
by two contour integrals (Lemma~\ref{5.2:lem-first-stage-bound}).
\begin{enumerate}
\item[(1)] Among the constituents of the cut groups, those that enter
Lemma~\ref{5.2:lem-first-stage-bound} as units are the constituents of the plain units of
clauses (c) and (d) of the classification of units at site (L), the only classes of that
classification that carry an amplitude $\mathsf{t}$. Among them, the
near constituents are of the type of rows 2--3, and the near counterterm pieces are of the
type of row 5.
\item[(2)] The smaller core of Lemma~\ref{5.2:lem-cut-groups} is used at site (L) only, and
its use there is legitimate.
\item[(3)] Lemma~\ref{5.2:lem-first-stage-bound} covers brackets
$f(\mathfrak{b}_1)-f(\mathfrak{b}_2)$. The constituents sent to the first part of the split in
the proof of clause (i) of Lemma~\ref{5.2:lem-cut-groups} whose expansion is trivial, and the
subtracted copies of that proof (the pairs $(X,y)$ with $X\cap Z\ne\emptyset$, taken at the
single background $U_k$), are values $F(X,\mathfrak{b},y)-F(X,1,y)$ at one admissible
$\mathfrak{b}$. They are bounded by the same
sizes, without the factor $\mathsf{w}_L^{-1}$. They pass through no $\mathsf{W}$. Their
constant is the one of \cite[\S1]{B-CMP122b} for whole groups and tails, and at most $C_\sharp$
times it for the near singles.
\item[(4)] The sharper scale sum
$S_\partial=\sum_\nu\rho^2L^{-(\nu-1)}$ in place of $S_\partial^{\sharp}$ is not asserted;
it is not needed.
\item[(5)] The bounds of Lemma~\ref{5.2:lem-cut-groups} are uniform in $K$.
\item[(6)] Two consistency checks hold: a family whose constituents are single cubes,
and the family $c_jA_X$, where $A_X$ denotes the Wilson action restricted to the plaquettes
contained in $X$ and $c_j$ denotes here a scalar coefficient depending on the scale $j$ only
(not a geometric constant).
\item[(7)] Older boundaries $\partial\Lambda_j$, $j<k$, need no separate treatment, and no
term-wise split occurs at sites (T) and (P).
\end{enumerate}
\end{remark}

\begin{proof}
\emph{(1).} Interior and wrapped constituents are clauses (a) and (b) of the classification
of units at site (L). They are treated by the site-(L) statements for interior and wrapped
constituents, and they carry neither $\mathsf{t}$ nor $\mathsf{w}_L^{-1}$.

The near constituents are of the type of rows 2--3 for three reasons:
\begin{itemize}
\item their read set is $S_v=X$;
\item $f_v$ is gauge invariant by property (f2) of the families;
\item $f_v(1)$ is subtracted.
\end{itemize}
Here $S_v=X$ is the read set of Lemma~\ref{5.2:lem-first-stage-bound} at sites (T) and (P);
at site (L) it is the read set of those constituents that are not pointed pieces.

At site (L) the read set of a pointed piece $p$ is instead the one fixed by the read-set
statement for pointed units. For a compact $p$ it is
$\mathfrak{r}_v=\bigcup\operatorname{St}(\Delta_{y_v})$, the union of the cells of the star
$\operatorname{St}(\Delta_{y_v})$ of the anchor cell $\Delta_{y_v}$. For a
long $p$ it is the set of cells meeting $X_v$.

The bound $|f_v\circ\mathsf{R}_p-f_v(1)|\le\mathfrak{N}(v)$ holds on the whole response disc.
The reason is that
$\mathsf{R}_p$ maps into admissible backgrounds, by clause (iii) of the background-variation
lemma.

\emph{(2).} By the proof of clause (i) of Lemma~\ref{5.2:lem-cut-groups}, the set
$\square^{\sim\nu}\cup\operatorname{supp}h_y$ lies in $\square^{\sim\nu+1}$. The set
$\square^{\sim\nu}\cup\operatorname{supp}h_y$ has side, in $n$-units, at most
\begin{equation*}
(2\nu+1)ML^{-(\nu-1)}\le 3M \qquad (L\ge 3).
\end{equation*}
Hence $\bar S_v$ lies in at most $3^d$ current cubes. Read on $\bar S_v$, as at sites (T) and
(P), the unit formed at a far point $y$ by the core and the counterterm piece therefore has
weight at most $e^{3^d\kappa}$ in
Lemma~\ref{5.2:lem-first-stage-bound}. This factor lies inside the $O(1)$ of
\eqref{5.2:eq-comparison-constant-a}; it is a constant of the same kind as the factor
$e^{2^d\kappa}$ in the constant $C_U$ of the three families per scale, fixed in the proposition
of site (T) in which $C_U$ appears.

At site (L), however, a far core is a pointed piece $p$ read on $\mathfrak{r}_v$. This set is
the star, and it
consists of at most $c_S=(5L)^d$ own-scale cells, not of $3^d$ current cubes. Such a piece is
summed by the summation of boundary terms at finer zones
(Theorem~\ref{5.4:thm-finer-zone-sum}). By the corollary of the read-set statement for pointed
units, its anchor weight is the
cube $\bar{\mathsf{w}}^3$ of the threshold weight, paid by $\mathsf{f}_v^{1/2}$, and its label
weight is that of the walk. It enters neither $\mathsf{W}$ nor $C_A'$ of
\eqref{5.2:eq-first-stage-bound-a}.

An interior or wrapped far unit carries no $\mathsf{w}_L^{-1}$. Its pieces are those of
clause (b), respectively (a), of the classification. They are collected in the member
$\varepsilon^{\mathrm{int}}$, together with the volume corollary for plain units, respectively
in the member $\varepsilon^{\mathrm{rel}}$, at those members' own weights and ceilings.

Consider the constituents lying between the core of the core estimate for frozen groups and
the core of the hypercube core estimate, that is, those $X\ni y$ contained in $2^d$ current
cubes but not in $\square^{\sim\nu}$. They become tails carrying the factor $t^5$. Per cell
they contribute $t^{-4}\cdot CE_0t^5$, where $C$ denotes a constant, and this lies inside
$C_IS_I$.

\emph{(3).} Take the values listed in (3). They are whole groups for far $y\in Z$, tails for
far $y\in Z^c$, and singles at near $y$.

We first treat values inside $Z^{\sim}$. Under the standing convention of this chapter,
$Z^{\sim}$ is read as $Z^{\sharp}$. By clause (c) of the strip theorem and its strip form of
\cite[(1.69)]{B-CMP122b}, the per-cell volume bound has the same constant on strip cells as
on the cells of $Z$. Indeed, the strip adds to a component of $Z$ at most $2^d-1$ times as
many cells as that component contains.
The per-cell volume input is supplied, for any cell, by clause (iv) of
Lemma~\ref{5.2:lem-cut-groups} and by the last column of \eqref{5.2:eq-unit-table-a}.

This input includes the row-4 values of the chain of the present step. Each such value is
charged once per term, at a cell of its own stage region. By the chain-count lemma their total
charge per zone-$n$ cell is at most
\begin{equation*}
3\bigl(N_0(g_k)+3\bigr)C_0^{\sharp}\le 3\bigl(\mathsf{N}_0(\mathsf{T}_n)+3\bigr)C_0^{\sharp}.
\end{equation*}
The per-cell constant is the function $\mathsf{K}_{\mathrm{cell}}(k,n)$ of
the zone-cell volume statement, entered through its volume clause.

To this are added the frozen terms $\mathfrak{F}(\mathfrak{a})$ of the live arrested pieces,
namely the chain terms $\mathbb{C}^{(i)}$ and the further frozen terms attached to the arrest.
They are counted
extensively, as $\mathsf{P}(\mathfrak{a})$ per live arrest inside the component, by clause
$(\beta')$ of the zone-cell volume statement. They enter the third member of the right-hand
side of the volume bound of \cite[(1.69)]{B-CMP122b}, discussed next, with the constant
$2c_{\mathsf{P}}M^{-d}$, a number fixed after $M$.

Values inside $Z^{\sim}$ enter only the volume bound at complex source $w$, the form of
\cite[(1.69)]{B-CMP122b} displayed in clause (c) of the proposition on the
$\mathbf{R}$-operation at site (L). In \cite[(1.69)]{B-CMP122b} this bound reads
$O(1)\sum_j|\Gamma_j^0\cap Y|$ when $Y$ is a component of
$Z^{\sim}$. Clause (iv) of Lemma~\ref{5.2:lem-cut-groups} and the last column of the table
supply it. For the units of rows 1--3, 5 and $\partial$ the supplied bound carries the extra
factor $C_\sharp$. Row 4's piles make
up the per-cell constant $\mathsf{K}_{\mathrm{cell}}(k,n)$. The two forms of this volume
bound, at complex source and in strip form, displayed in the proposition on the
$\mathbf{R}$-operation at site (L), clause (c), carry this constant. They do not carry the
constant $C_\sharp O(1)+1$ of \eqref{5.2:eq-comparison-constant-a}.

Now let $Y\setminus Z^{\sim}\ne\emptyset$, that is, $Y\setminus Z^{\sharp}\ne\emptyset$ under
the same convention, by clause (c) of the proposition on the $\mathbf{R}$-operation at site
(L). The smallness is then obtained as in \cite[\S1]{B-CMP122b}: ``For the remaining terms we
use a part of the exponential factor to get a small constant in the bound.''

For values this part is the units' own rate. It is $e^{-\kappa d_j(X)}$ for $j<k$. At $j=k$ it
is the improved level-$k$ rate $e^{-(1+4\beta_s)\kappa d_k(X)}$, given by the decay-fraction
lemma. This is the first, inductive, factor of \cite[(1.66)]{B-CMP122b}, and it is the
size at which these terms are entered in \eqref{5.2:eq-unit-table-a}. The weight of
the strip theorem is met by these own rates.

The second factor of \cite[(1.66)]{B-CMP122b}, the $(\kappa_1-1)$-decay per localisation
cube, lives on pieces. It is supplied by the proofs of the source-piece statements. The third
factor of \cite[(1.66)]{B-CMP122b}, the decay factor in $\mathbf{H}$, is charged to no class,
and neither is \cite[(1.67)]{B-CMP122b}. Where a depth factor is needed, the depth lemma is
used in their place.

Since $\mathfrak{n}=E_0$, as in \cite[\S1]{B-CMP122b}, the constant for whole groups and tails
is Ba{\l}aban's. For the near singles it is at most $C_\sharp$ times it.

\emph{(4).} The sharper sum would require re-inserting the quadratic main terms of the
missing constituents into the core of the core estimate for frozen groups, with boxes inside
$Z^c$ chosen per constituent.

\emph{(5).} The only scale sum is
\begin{equation*}
\sum_{\nu}(\nu+2)\rho^2L^{-(\nu-1)}.
\end{equation*}
This sum contains no $K$ and, by clause (iv) of Lemma~\ref{5.2:lem-cut-groups}, it is bounded
uniformly in $g\le\gamma_J$; no other sum over scales occurs.
The radius $r$ of the ball $\mathbb{D}$ (Definition~\ref{5.1:def-weight-classes}) enters
through the bound $|b_j(\hat w)|\le b_++\tfrac23r$ of clause (ii) of
Lemma~\ref{5.2:lem-cut-groups}, as $r\mathfrak{a}_n^2\asymp rg_n^2(\log g_n^{-2})^{2q}$.
This is of the kind of row 5, and it gives an infrared
ceiling on $g$ chosen after $r$. The numbers $27d$, $3^d$, $e^{3^d\kappa}$, $C(\kappa)$ and
the powers of $M$ are constants inside $C_\sharp$. They are met by $g$, which is chosen after
$(L,M,\kappa)$.

\emph{(6).} First check: a family all of whose constituents are single cubes. No pair is cut,
because $\pi_j$ refines the partition of $Z$; only $\operatorname{supp}h_y$ straddles
$\partial Z$, and only for $y$ adjacent to $\partial Z$. The contribution is at most a
constant multiple of $t^{-3}$ points times $t^4$, consistent with (iii)--(iv) of
Lemma~\ref{5.2:lem-cut-groups}.

Second check: the family $c_jA_X$. The near layer contributes $c_j\rho^2(\nu+2)t$ per cell,
and the far groups cancel at $y$. This is consistent with the true size.

\emph{(7).} Older boundaries $\partial\Lambda_j$ are protected by the trimmed region
$\Lambda_j'$, obtained from $\Lambda_j$ by removing one layer of $M\check{R}_j$-cubes. The
domains
$X\ni y$ missing there have $d_j(X)\ge\check{R}_j$, so they are tails. At sites (T) and (P)
the whole action sits under one integral, and no term-wise split occurs. No other part of the
expansion is affected by the cut at $\partial Z$.
\end{proof}

\begin{definition}[The comparison constant and the budget at each site]
\label{5.2:def-comparison-constant}
Let $S_I$, $S_{\omega}$ and $S_{\partial}^{\sharp}$ be the per-cell scale sums of the table of units
(Proposition~\ref{5.2:prop-unit-table}, display~\eqref{5.2:eq-unit-table-a}). Denote by
$S_I^{\flat}$ the sum $S_I$ evaluated with $\rho=\rho_{n,j}$, the radius ratio of
Notation~\ref{5.1:not-coupling-functions}. Denote by $S_I^{\mathrm{pr}}$ the sum $S_I$ evaluated at
$\rho\equiv 1$, which is the corresponding sum of Ba{\l}aban's construction. For integers $\nu\ge 0$ put
$\bar\rho_\nu^{2}:=(4/3)^{2q}\,(5/4+\lambda \nu)$, and
\begin{equation}\label{5.2:eq-comparison-constant-1}
\begin{aligned}
\bar S_I&:=\sum_{\nu\ge 1}\nu^{3}\bar\rho_\nu^{2}(1+\bar\rho_\nu)\,L^{-\hat\alpha(\nu-1)},\qquad
\bar S_{\omega}:=\sum_{a\ge 0}L^{-a/2}\bar\rho_a^{2},\\
\bar S_{\partial}&:=\sum_{\nu\ge 1}(\nu+2)\,\bar\rho_\nu^{2}\,L^{-(\nu-1)} .
\end{aligned}
\end{equation}
These are numbers determined by $L,\lambda,c_2,q,\hat\alpha$. Define
\begin{equation}\label{5.2:eq-comparison-constant-a}
\begin{gathered}
\gamma^{\circ}:=(8c_2+4)^{-1/2},\qquad
C_{\sharp}:=\sup
\Bigl(1+\frac{S_I^{\flat}+S_{\omega}+S_{\partial}^{\sharp}}{S_I^{\mathrm{pr}}}\Bigr),\\
\bar C_{\sharp}:=1+\frac{\bar S_I+\bar S_{\omega}+\bar S_{\partial}}{S_I^{\mathrm{pr}}},\qquad
\bar C_{\mathsf{W}}:=O(1)\,E_0\bigl(C_I\bar S_I+C\bar S_{\omega}\bigr).
\end{gathered}
\end{equation}
In the definition of $C_{\sharp}$ the supremum is taken over all $0<g\le\gamma^{\circ}$, all
sequences $(x_i)$ obeying the reference recursion (Definition~\ref{5.1:def-reference-recursion}) with
$x_n>g^{-2}$, and all levels $n$; $C_I$, $E_0$
and $C$ are the constants of the first and second entries of the table of units. The constant $O(1)$
in $\bar C_{\mathsf{W}}$ is the one in the constant $C_{\mathsf{W}}$ of the first-stage bound
(display~\eqref{5.2:eq-first-stage-bound-a}), which is of the form
$O(1)E_0(C_IS_I^{\flat}+CS_{\omega})$. The number $C_{\sharp}$ is the
\emph{comparison constant}. The sums in its definition do not depend on $K$. The ceiling
$\gamma^{\circ}$ depends on the flow constant $c_2$ alone. In the order of choice of the constants
(Definition~\ref{5.1:def-admitted-inputs}) it
precedes every constant that depends on $C_{\sharp}$, on the constant $C_{\mathsf{W}}$ attached to the
weight supremum $\mathsf{W}$ of the first-stage bound (display~\eqref{5.2:eq-first-stage-bound-a}), or
on $c_A$, so that no constant is made to depend on $\gamma_J$ before $\gamma_J$ has been defined. The
comparison constant is fixed before
the smallness condition on the coupling. The further powers of $M$, the number $27d$, the constant
$C_{\partial}'$ of the cut-group lemma (Lemma~\ref{5.2:lem-cut-groups},
display~\eqref{5.2:eq-cut-groups-a}) and the factor $e^{3^{d}\kappa}$ are contained in the constant
$O(1)$ of~\eqref{5.2:eq-comparison-constant-3} below.

Let $0<g\le\gamma^{\circ}$, let $(x_i)$ be a sequence obeying the reference recursion with
$x_n>g^{-2}$, and let $j\le n$. Then
\begin{equation}\label{5.2:eq-comparison-constant-2}
\begin{gathered}
\rho_{n,j}^{2}\le\bar\rho_{n-j}^{2},\qquad S_I^{\flat}\le\bar S_I,\qquad
S_{\omega}\le\bar S_{\omega},\qquad S_{\partial}^{\sharp}\le\bar S_{\partial},\\
C_{\sharp}\le\bar C_{\sharp},\qquad C_{\mathsf{W}}\le\bar C_{\mathsf{W}} .
\end{gathered}
\end{equation}
Every later use of $C_{\sharp}$ and $C_{\mathsf{W}}$ is as an upper bound only. Therefore $\bar C_{\sharp}$
and $\bar C_{\mathsf{W}}$ may be used in their place.

Consider each of the three expansion sites (T), (P), (L) and the radii $\mathsf{w}_T$, $\mathsf{w}_P$,
$\mathsf{w}_L$ of their amplitude discs, display~\eqref{5.2:eq-three-sites-a}; write
$\mathsf{w}_{\mathrm{site}}$ for the radius at the site under consideration.
Write $\bar\varepsilon^{w}$ for the constant of the history potential on the ball at the complex
source $w$, and $\bar\varepsilon^{\mathrm{pr}}$ for the same constant in Ba{\l}aban's construction.
At each such site,
\begin{equation}\label{5.2:eq-comparison-constant-3}
\begin{aligned}
\bar\varepsilon^{w}
&\le C_A\,\mathsf{w}_{\mathrm{site}}^{-1}\,O(1)\Bigl[E_0\bigl(C_IS_I+CS_{\omega}
+CM^{3}S_{\partial}^{\sharp}\bigr)+\mathfrak{m}_n+B_0+(N_0+3)C_0\Bigr]\\
&\qquad+C(r+b_+)\,g^{2}(\log g^{-2})^{2q}\\
&\le C_{\sharp}\,\bar\varepsilon^{\mathrm{pr}}+\tfrac14\nu_0\alpha_4 .
\end{aligned}
\end{equation}
The letters in~\eqref{5.2:eq-comparison-constant-3} are those of
Proposition~\ref{5.2:prop-unit-table} and Lemma~\ref{5.2:lem-first-stage-bound}, and $\nu_0$,
$\alpha_4$ are those of the smallness condition on the coupling. Here $S_I$ is evaluated at
$\rho=\rho_{n,j}$, that is, $S_I=S_I^{\flat}$.
Per cell of level $m$, the unit's own level, the term $(N_0+3)C_0$ stands for
$6(\mathsf{N}_0(\mathsf{T}_m)+3)C_0^{\sharp}$.
The right member $C_{\sharp}\bar\varepsilon^{\mathrm{pr}}+\tfrac14\nu_0\alpha_4$ lies within the flat
form of the smallness condition on the coupling for $0<g\le\gamma_J(r)$. The first inequality
of~\eqref{5.2:eq-comparison-constant-3} is the
\emph{budget at the site}.
\end{definition}

\begin{proof}
We begin with the majorants~\eqref{5.2:eq-comparison-constant-2}. Let $g$, $(x_i)$, $n$ and $j$ be as
stated. The reference recursion gives two bounds:
\begin{equation*}
\frac{x_j}{x_n}\le 1+\frac14+\frac34\lambda(n-j),\qquad \log x_j\ge\log x_n-\frac13 .
\end{equation*}
Inserted into the ratio $\rho_{n,j}$ of the radii at the levels $n$ and $j$, these give
$\rho_{n,j}^{2}\le\bar\rho_{n-j}^{2}$.

Since $\lambda>0$, $\bar\rho_\nu$ increases with $\nu$. We compare the sums term by term.
\begin{itemize}
\item In $S_I^{\flat}$ the summand of index $\nu=1+n-j$ carries $\rho=\rho_{n,n+1-\nu}$. Hence
$\rho\le\bar\rho_{\nu-1}\le\bar\rho_\nu$. The factor $\nu^{3}\rho^{2}(1+\rho)$ increases with $\rho$, so
$S_I^{\flat}\le\bar S_I$.
\item In $S_{\omega}$ the summand of index $a\ge 1$ carries $\rho_{n,n-a}^{2}\le\bar\rho_a^{2}$, the
upper index being the level $n$ at which $x_n>g^{-2}$ is assumed. The
term $1$ is at most $\bar\rho_0^{2}=(5/4)(4/3)^{2q}$. Hence $S_{\omega}\le\bar S_{\omega}$.
\item In $S_{\partial}^{\sharp}$ the summand of index $\nu$ carries
$\rho_{n,n+1-\nu}^{2}\le\bar\rho_{\nu-1}^{2}\le\bar\rho_\nu^{2}$. Hence
$S_{\partial}^{\sharp}\le\bar S_{\partial}$.
\end{itemize}
The three majorants depend neither on $g$, nor on the sequence, nor on $K$. Taking the supremum over
$0<g\le\gamma^{\circ}$ in~\eqref{5.2:eq-comparison-constant-a} therefore gives
$C_{\sharp}\le\bar C_{\sharp}$. Since $C_{\mathsf{W}}$ is of the form $O(1)E_0(C_IS_I^{\flat}+CS_{\omega})$,
the same majorants $S_I^{\flat}\le\bar S_I$ and $S_{\omega}\le\bar S_{\omega}$ give
$C_{\mathsf{W}}\le\bar C_{\mathsf{W}}$.

We next describe the weights at site (L) and the scope of the constant $O(1)$. At site (L), under the
weight of \cite[(1.66)]{B-CMP122b}, a far unit acquires the factor $e^{(7/4)3^{d}\kappa}$, while the
units of the second and third entries of the table of units keep the decay rate
\begin{equation*}
(5/12-\delta/4)\kappa d_j\quad\text{when } j<k,\qquad (1/4-\delta/4)\kappa d_k\quad\text{when } j=k.
\end{equation*}
For the units that remain in $\mathsf{W}$ at site (L), both the growth factor of the far unit and the
loss incurred by the second and third entries keeping only these rates lie inside
the $O(1)$ of~\eqref{5.2:eq-comparison-constant-3} and are met by $\mathsf{w}_L^{-1}$; the pointed far
cores are treated differently, as described below. This applies to the units that are estimated at
site (L) by the first-stage bound. To say which these are, two independent distinctions among the
units at site (L) are used: a unit is exterior if $K_v\setminus Z\neq\emptyset$ and interior if
$K_v\subset Z$; and, independently of this, a unit is plain or wrapped. Pointed, short and long units
are those of Definition~\ref{5.4:def-pointed-units}. The first-stage bound is set up at site (L) for
the plain units of the exterior class and of the class of long single constituents and tails; these are
the only units carrying the amplitude $\mathsf{t}$, and the wrapped units and
the interior units carry neither $\mathsf{t}$ nor the factor $\mathsf{w}_L^{-1}$. Of the units carrying
$\mathsf{t}$, the pointed ones and the long ones are nevertheless summed by other means, described
next, and at site (L) the $O(1)$ carries only what the first-stage bound keeps in $\mathsf{W}$ and in
$C_A'$ there.

The pointed units among the units carrying $\mathsf{t}$ are handled by the summation of pointed units.
These are
the far cores read on the read set $\mathfrak{r}_v$ and the constituents of the entries $1$--$3$, $5$
and $\partial$. They contribute a constant of the summation of pointed units which does not depend on
$\mathsf{T}$ and is contained in
$\mathsf{C}_1^{\mathrm{ex}}$. Their factor $e^{(7/4)3^{d}\kappa}$ is
replaced by the cube of the threshold weight $\bar{\mathsf{w}}$, paid against $\mathsf{f}_v^{1/2}$.

The long single constituents and the tails are summed together with the long exterior plain units.
This is done by the summation of boundary terms at finer zones
(Theorem~\ref{5.4:thm-finer-zone-sum}), through the bundle frame for long units off the treated
class (Lemma~\ref{5.4:lem-bundle-frame}) and the bundle sum for long units off the treated class
(Theorem~\ref{5.4:thm-bundle-sum}), with the number $\mathsf{C}^{\mathrm{lg},d}$ contained in
$\mathsf{C}^{\Sigma,4}$. These units
undergo no resummation of the first-stage bound, and carry none of $\mathsf{W}$, $C_A'$, $\eta_p$,
$\rho_A$. The units of the
fourth entry at levels $n'<k$, and in windows other than their own, are likewise handled by the
summation of boundary terms at finer zones.

For the plain interior pointed units, the same $O(1)$ sits in the factor
$\mathsf{K}^{\flat}(\mathsf{w}^{\oplus})^{3}$ of the column bound of the summation lemma for plain
interior units. The exterior pieces of these units carry the factor $e^{-(\check{R}_k-2)}$; the others
are absorbed in $c_V^{S}\mathsf{K}_{\mathrm{cell}}$.

We now obtain the budget. Collect the last column of the table of units per unit cell, the entry
$\partial$ entering through Lemma~\ref{5.2:lem-cut-groups}. Combine this with the factor
$C_A\mathsf{w}_{\mathrm{site}}^{-1}$ supplied by Lemma~\ref{5.2:lem-first-stage-bound}. The result is
the first inequality of~\eqref{5.2:eq-comparison-constant-3} at every site. In the second inequality
the scale sums enter through the ratio that defines $C_{\sharp}$:
by~\eqref{5.2:eq-comparison-constant-a}, for $0<g\le\gamma^{\circ}$,
\begin{equation*}
S_I^{\flat}+S_{\omega}+S_{\partial}^{\sharp}\le(C_{\sharp}-1)\,S_I^{\mathrm{pr}},
\end{equation*}
where $S_I^{\mathrm{pr}}$ is the sum of Ba{\l}aban's construction, the scale sum occurring in the
corresponding quantity $\bar\varepsilon^{\mathrm{pr}}$ of that construction; the factor $M^{3}$ in front
of $S_{\partial}^{\sharp}$ is among the powers of $M$ that, in passing to the second inequality, are
absorbed into the constant $O(1)$. The right member
$C_{\sharp}\bar\varepsilon^{\mathrm{pr}}+\tfrac14\nu_0\alpha_4$ lies within the flat form of the
smallness condition on the coupling for $0<g\le\gamma_J(r)$.

Next we read the term $(N_0+3)C_0$. The frozen terms $\mathfrak{F}(\mathfrak{a})$ of the live arrested
pieces are units of the fourth entry at each of the sites (T), (P), (L). They consist of the chain terms
$\mathbb{C}^{(i)}$ and the change-of-variables terms of \cite{B-CMP122a}. Their pile is counted at the
unit's own level~$m$.

Consider the chains at a point: the attempt of the present step and the live arrests nested there.
At most two of them have a stage family at level $m$. Indeed, by the lemma on attempt ranges, their
attempt ranges are pairwise disjoint except for one index. Let $k'$ be the step of such a chain. It
contributes at most $3(N_0(g_{k'})+3)C_0^{\sharp}$ at an anchor of level $m$, by the lemma on the
stage families of one chain applied at step $k'$.

Moreover, $N_0(g_{k'})\le\mathsf{N}_0(\mathsf{T}_m)$ for $m\le k'$, because the reference recursion
gives $x_m\ge x_{k'}-2c_2$. Hence, per cell of level $m$, the term $(N_0+3)C_0$ stands for
$6(\mathsf{N}_0(\mathsf{T}_m)+3)C_0^{\sharp}$. The change-of-variables terms add at most $1$ per block
of the cover of their stage \cite{B-CMP119}. No contribution beyond these two chains arises at the
cell.

The function $\mathsf{N}_0(\mathsf{T})$ enters only as an upper bound, so that replacing it by a
larger function only enlarges the bounds. At sites (T) and (P) the whole member stays inside
the $O(1)$ that is met by $\mathsf{w}_{\mathrm{site}}^{-1}$ on the real small-field domain of
Ba{\l}aban's construction; on the polydisc
$\|\mathsf{B}\|<T^{\mathbb{C}}$ it is paid by distance instead, at half rate.

We turn to the exterior units at site (L), those with $K_v\setminus Z\neq\emptyset$. Their brackets
are counted in the third member of the bracket bound of the bracket proposition at site (L), the
member that carries $\theta'(g_k)$. This uses the factor
\begin{equation*}
4C_A'\,\mathsf{w}_L^{-1}\eta_p^{1/2}e^{-(\kappa_1-1)m(p)}\le 4C_A'\,\theta'(g_k)
\end{equation*}
of the half-power form~\eqref{5.2:eq-three-sites-a} of the first-stage bound
(Lemma~\ref{5.2:lem-first-stage-bound}), applied to a bracket,
together with the coupling-free constant $\mathsf{K}^{\mathrm{arr}}$ of the collar. It adds
$8C_A'\mathsf{K}^{\mathrm{arr}}\theta'(g_k)$ to $\varepsilon_{\mathrm{br}}(k)$.

This member is a majorant that is not needed, for the following reason. The $\mathsf{t}$-pieces of these
units are pieces of the fourth entry among the exterior plain units. They are summed in bundles by the
summation of
boundary terms at finer zones, into
$\varepsilon^{\Sigma4}(k)\le\theta'(g_k)\mathsf{C}^{\Sigma,4}$. The hypothesis of the first-stage bound
is read at $\mathsf{t}=0$, and its gain $\mathsf{w}_L^{-1}\eta_p^{1/2}$ is not used.

Two further contributions come from the summation lemma for exterior units. The non-empty
continuations of the $\mathsf{t}$-pieces along the second to fourth links multiply that member by
$(1+C_{\Sigma}C_e)^{3}$, where $C_{\Sigma}$ and $C_e$ are the constants of the summation lemma for
exterior units. The $\mathsf{t}$-free pieces of the first term $f_v(\mathfrak{b}_3)$ along the
second to fourth links enter $\varepsilon^{\mathrm{ext},\sigma}(k)$.

The interior units ($X_v\subset Z$, hence $K_v\subset Z$) carry no $\mathsf{t}$ and no factor
$\mathsf{w}_L^{-1}$. They enter
the norm $\|\cdot\|_a$ as well. Indeed, the supremum in $\|\cdot\|_a$ restricts the filing domain
$Y_4$ of a piece, not the domain of the unit, and an interior unit has pieces whose filing domain $Y_4$
leaves $Z$. This is seen from \cite[\S1, after (1.66)]{B-CMP122b}:
\begin{quotation}
$Y_4\setminus Z=(Y_0\setminus Z)\cup\bigcup_{i=2}^{4}(Y'_i\setminus Z)$; therefore the above
exponential provides a complete exponential factor for the domain $Y_4\setminus Z$. The domains $Y'_i$
connect $Y_1\setminus Z$ with $\Lambda$, and the third term in the exponential yields a connection of
$X$ with $Y'_{i_0}$.
\end{quotation}
Here the $Y'_i$ are domains that ``contain at least one component of $\Lambda$, and intersect the
boundary layer $Y_1\setminus Y_1^{0}$'' \cite[\S1]{B-CMP122b}.

The tuples of an interior unit all of whose entries are empty are values. Its pieces whose filing
domain stays inside the component of the large-field region $Z$ that carries the unit are counted with
the values. They enter the per-cell constant
$\mathsf{K}_{\mathrm{cell}}$ of the volume bound and the extensive pile $\mathsf{P}(\mathfrak{a})$, at
the price $e^{4}$ of the corollary on the pieces of plain units: all pieces of one unit $v$ sum to at
most $\mathfrak{N}^{\flat}(v)e^{4d_k(K_v)+4}$, with $\mathfrak{N}^{\flat}=\mathfrak{N}$ on the fourth
entry. Its pieces whose filing domain leaves $Z$ form the member $\varepsilon^{\mathrm{int}}(k)$ of the
bracket bound for plain units. For wrapped units they lie inside $\varepsilon^{\mathrm{rel}}(k)$. Both
follow from the piece display of the bracket proposition at site (L).

For a unit buried at depth $d^{\wedge}$ in the core, the piece bound is supplied by the depth lemma:
such a unit responds to a change made outside the core by at most a constant of the depth lemma times
$e^{-\delta_Pd^{\wedge}/2}$. This bound takes the
place of the third factor of \cite[(1.66)]{B-CMP122b}, and its
pointed version serves the entries
$1$--$3$ and $5$.

The summation is carried out as follows. For the plain units it is the summation lemma for plain
interior units, under its own ceiling. For the wrapped units it is the lemma on deep wrapped units
together with the proposition on wrapped units at site (L), under their ceiling at site (L).
Consequently no per-cell frozen pile of a cell of zone $n$ is multiplied by $\theta'(g_k)$, and no
factor indexed by the level difference $k-n$ arises. The depth is paid by the factor of the depth
lemma, split between the cell bound and the
pile at a cell, so that nothing formed at earlier levels of the core enters the bound.

Finally, on the real small-field domain of Ba{\l}aban's construction the factor
$\mathsf{w}_{\mathrm{site}}^{-1}$ tends to zero as $g_k\to 0$. It pays
every member of~\eqref{5.2:eq-comparison-constant-3}, including $(N_0+3)C_0$ and $B_0$: one has
$\theta(g_k)(N_0+3)\to 0$, since $N_0$ grows only like $\log\log g^{-2}$ \cite{B-CMP122a}.

At site (L) the budget of the brackets, with the pile of chain terms of all $N(g_k)$ preparatory
stages, is the function $\varepsilon_{\mathrm{br}}(k)$ of the bracket proposition. It is dominated
inside the ceiling of $\gamma_J$ given by a supremum over $\mathsf{T}$. For the values, which carry no
factor $\mathsf{w}_L^{-1}$, the same pile is a function of $g_k$. It is controlled by the ceilings of
the bracket proposition and of the volume count.

On the polydisc $\|\mathsf{B}\|<T^{\mathbb{C}}$, where $\mathsf{w}^{\mathbb{C}}$ is bounded, the members
of the fourth entry are paid by distance instead. This uses the second step of the lemma on the
polydisc and the distance estimate at half rate. Under the flat variant of the step on the support
$\mathfrak{X}$, the chain-term member of~\eqref{5.2:eq-comparison-constant-3} reads $C_0^{\sharp}$, the
constant of the corollary bounding the chain terms in
the A-slot norm. At site (P) the operative budget is the A-slot condition among the conditions defining
$\gamma_J$, not $\mathsf{w}_P^{-1}$ alone.
\end{proof}

\begin{corollary}[The first stage inside the integral and the sourced preparatory step]
\label{5.2:cor-sourced-preparatory}
Let $w\in\mathcal{W}_k$. Assume the following statements:
\begin{itemize}
\item the first-stage bound by two contour integrals, Lemma~\ref{5.2:lem-first-stage-bound};
\item the bound on the renormalised pointed group, valid on the whole of
  $\mathsf{R}_p(\{|\mathsf{t}|\le\mathsf{w}\})$;
\item the table of units and their sizes, Proposition~\ref{5.2:prop-unit-table}, that is,
  the table \eqref{5.2:eq-unit-table-a};
\item the statement that the second stage sees activities only,
  Lemma~\ref{5.2:lem-second-stage-activities}, applied after the two-slot M\"obius expansion
  of the dependence on the $A$-slots;
\item the preparatory-step lemma in its two-slot form, both slots of the bracket
  \cite[(1.60)]{B-CMP122a} varied;
\item the first-stage lemma for boundary-type units;
\item the orbit-datum statement.
\end{itemize}
Then the following hold.
\begin{enumerate}
\item[(i)] Since the bound on the renormalised pointed group holds on the whole of
  $\mathsf{R}_p(\{|\mathsf{t}|\le\mathsf{w}\})$, the expansion centred at the point $y$ of each
  pointed unit (the frozen groups $\mathfrak{i}^{(j)}_y$ of the first entry of
  \eqref{5.2:eq-unit-table-a}) is carried out inside the integral. Ba{\l}aban's step
  ``we differentiate these terms with respect to the parameter $t_{\square}$'' in
  \cite{B-CMP119} is the differentiation in the amplitude $\mathsf{t}$ of the first-stage
  bound. In particular, the standing hypothesis of the opening
  section of this chapter on the expansion centred at $y$ is discharged.
\item[(ii)] The sourced form of the preparatory-step lemma holds. Its bracket is
  \begin{equation}\label{5.2:eq-sourced-preparatory-1}
  \{F+W\}\bigl(\Phi(B'_j),A_j\bigr)-\{F+W\}\bigl(\Phi(0),A_j=0\bigr),
  \end{equation}
  as in \cite[(1.60), p.~188]{B-CMP122a}. Its second, third and fourth clauses are the
  following entries of the table \eqref{5.2:eq-unit-table-a}.
  \begin{itemize}
  \item The second clause, bounding the boundary and chain terms, is the fourth entry, in the
    anchored-domain norm $\|\cdot\|^{A}$. Write $A^{\flat}$ for the $A$-slot datum with which
    that entry stores the $\mathbf{B}$- and $\mathbb{C}$-terms. The $A_j$-dependence of these
    terms is the stored datum $A^{\flat}$, and at site (P) of the three expansion sites of
    \eqref{5.2:eq-three-sites-a} it is expanded by the two-slot M\"obius expansion.
  \item The third clause, bounding the pointed groups and the differences of the $D^{(j)}$-
    and $R^{(j)}$-terms, is the first three entries. The entry marked $\partial$ is idle at
    site (P): there is one integral, with no term-wise split.
  \item The fourth clause, bounding the source terms, is the fifth entry.
  \end{itemize}
  In particular, the standing hypothesis of the opening section of this chapter on the
  sourced preparatory step is discharged.
\item[(iii)] For each stage $n$ of the preparatory operation at level $k$, the chain terms of
  stage $n+1$ decompose as
  \begin{equation}\label{5.2:eq-sourced-preparatory-2}
  \mathbb{C}^{(n+1)}_k=\sum_{(X,\mathcal{S})}\mathbb{C}^{(n+1)}_k(X,\mathcal{S};w),
  \qquad
  \bigl\|\mathbb{C}^{(n+1)}_k\bigr\|^{A}_{(n+1)}\le C_0^{\sharp},
  \end{equation}
  the norm being that of the weighted second-stage output bound.
\item[(iv)] At site (L) the activities $V(Y;w)$ are obtained in two parts; here the units at
  site (L) fall into the four classes (a)--(d) of the classification of units at that site:
  the wrapped units (a), the plain interior units (b), and the plain units (c), (d). The
  first-stage bound supplies the plain units of the classes (c) and (d). The first-stage lemma
  for boundary-type units supplies the wrapped units and the plain interior units, the classes
  (a) and (b). Two readings of the slots are used.
  \begin{itemize}
  \item The exterior $A$-slots of the units of the first three entries, of the fifth entry
    and of the entry marked $\partial$ enter as holomorphic parameters, by the parameter
    lemma for the exterior slots.
  \item The live exterior $A$-slots of the units of the fourth entry, in the sense of that
    entry of \eqref{5.2:eq-unit-table-a}, are evaluated at the rotation the site presents,
    through the orbit-datum statement, as \eqref{5.2:eq-three-sites-a} says.
  \end{itemize}
\end{enumerate}
\end{corollary}

\begin{proof}
We take the four conclusions in turn.

\emph{Conclusion} (i). The first-stage bound, Lemma~\ref{5.2:lem-first-stage-bound}, writes
the piece $v_p$ of a unit $f_v$ with localisation datum $p$ through two operations: the
derivative at $\mathsf{t}=0$ of $\mathsf{t}\mapsto f_v(\mathsf{R}_p(\mathsf{t},\mathsf{s}))$,
and the derivatives in the decoupling parameters $\mathsf{s}(\Delta)$. Of $f_v$ it asks two
things only:
\begin{itemize}
\item that $\mathsf{R}_p\subset\mathfrak{D}_v$;
\item that $|f_v\circ\mathsf{R}_p-f_v(1)|\le\mathfrak{N}(v)$ on the amplitude disc.
\end{itemize}
For the frozen groups $\mathfrak{i}^{(j)}_y$ of the first entry of
\eqref{5.2:eq-unit-table-a}, the table takes the size $\mathfrak{N}(v)$ from the bound on the
renormalised pointed group. By hypothesis, that bound holds on all of
$\mathsf{R}_p(\{|\mathsf{t}|\le\mathsf{w}\})$. Hence the expansion centred at $y$ is
performed on the integrand itself, inside the integral. The differentiation with respect to
the parameter written $t_{\square}$ in \cite{B-CMP119} is exactly the differentiation in the
amplitude $\mathsf{t}$ of Lemma~\ref{5.2:lem-first-stage-bound}. This discharges the standing
hypothesis on the expansion centred at $y$.

\emph{Conclusion} (ii). We apply Lemma~\ref{5.2:lem-first-stage-bound} to the units of
\eqref{5.2:eq-unit-table-a}, with the sizes $\mathfrak{N}(v)$ listed there. This gives the
second, third and fourth clauses of the preparatory-step lemma in its sourced form, with both
slots of the bracket \eqref{5.2:eq-sourced-preparatory-1} varied. The sourced form is the
preparatory-step lemma of the hypothesis, with the inputs of its second, third and fourth
clauses replaced, at the source $w$, by the corresponding entries of
\eqref{5.2:eq-unit-table-a} as bounded by Lemma~\ref{5.2:lem-first-stage-bound}.

The second clause, on the boundary and chain terms, is read off the fourth entry, in the norm
$\|\cdot\|^{A}$. There the $\mathbf{B}$- and $\mathbb{C}$-terms are stored together with the
datum $A^{\flat}$, and their dependence on $A_j$ enters only through $A^{\flat}$. At site (P)
this dependence is expanded by the two-slot M\"obius expansion.

The third clause, on the pointed groups and the response differences, is read off the first
three entries, the frozen groups and the differences of the $D^{(j)}$- and $R^{(j)}$-terms.
The entry marked $\partial$ contributes nothing at site (P): the integral there is one integral
and is not split term by term.

The fourth clause, on the source terms, is read off the fifth entry. These are the two terms
of that entry.

This discharges the standing hypothesis on the sourced preparatory step.

\emph{Conclusion} (iii). By Lemma~\ref{5.2:lem-second-stage-activities}, the outputs of the
second stage, the chain terms $\mathbb{C}^{(n+1)}_k(X)$ among them, are analytic on the output
spaces and in all parameters. An output with domain $X$ reads only activities $V(Y)$ with
$Y\subset X$. We apply the lemma after the two-slot M\"obius expansion: its elements are then
filed by the pairs $(X,\mathcal{S})$. This gives the decomposition
\eqref{5.2:eq-sourced-preparatory-2}, and the norm bound by $C_0^{\sharp}$ there is the one
the weighted second-stage output bound delivers.

\emph{Conclusion} (iv). At site (L) the units fall into the classes (a)--(d) of the
classification of units at that site.
\begin{itemize}
\item For the plain units of the classes (c) and (d), the pieces are those of
  Lemma~\ref{5.2:lem-first-stage-bound}.
\item For the wrapped units (a) and the plain interior units (b), the first stage is the
  first-stage lemma for boundary-type units.
\end{itemize}
Together these two give $V(Y;w)$.

The slots are read as follows. The exterior $A$-slots of the units of the first three
entries, of the fifth entry and of the entry marked $\partial$ are carried as parameters,
by the parameter lemma for the exterior slots. The live exterior $A$-slots of the units of
the fourth entry, in the sense of that entry of \eqref{5.2:eq-unit-table-a}, are evaluated at
the rotation that site (L) presents: by the orbit-datum statement, such a unit is its orbit
datum, bounded by $\mathfrak{N}^{\mathbb{C}}$, which is the reading fixed in
\eqref{5.2:eq-three-sites-a}.
\end{proof}

In \cite[p.~276, ll.~31--42]{B-CMP119} it is stated that the action in \cite[(3.29)]{B-CMP119} differs from the action analysed before only by its additional dependence on the field $tA_k$, where $t$, $A$ and $A_k$ denote the variables so named in \cite[(3.29)]{B-CMP119}. It is further stated that this dependence was already taken into account in the preceding considerations. The exponent of \cite[(3.29)]{B-CMP119} (see \cite[p.~272]{B-CMP119}) is
\begin{equation*}
-t\langle A,\mathcal{A}A_k\rangle+\mathbf{P}^{(k)}\bigl(g_k,(tA_k,A)\bigr)+\cdots,
\end{equation*}
the dots standing for the further terms written there, and $\mathcal{A}$ being the constrained fluctuation form of item~(b) of the definition below.

This book uses that statement for $\mathbf{P}^{(k)}$ and for those further terms. They are read as functions of the single field $g_kC(tA_k\oplus A)$, and the preceding considerations are linear in the supremum of that field. The corresponding bound is the one-field bound for $\mathbf{P}^{(k)}$ and the further terms, which estimates these terms through the supremum of that one field.

The bilinear term $-t\langle A,\mathcal{A}A_k\rangle$, taken against the small-field characteristic function $\chi_k$ (written $\chi^{(k)}$ in \cite{B-CMP119}), is treated in this section by a separate argument. Its coefficient is of order one and its size is of order $\mathfrak{b}^2$, with $\mathfrak{b}$ as in the definition below. The Gaussian part of this term is computed exactly by the theorem on the Gaussian part of the bilinear term and its corollary. The sentences of the same content at \cite[p.~272]{B-CMP119} and at \cite[p.~279]{B-CMP119} are used in this book in the same reading. The following definition introduces the objects with which the bilinear term is treated in this section.

\begin{definition}[The rim datum and the upper box]\label{5.2:def-rim-datum}
Fix $k\ge 0$ and consider step $k+1$, the renormalisation step from level $k$ to level $k+1$. For $k=0$ this is the step of \cite[\S1]{B-CMP119}. Let $\mathfrak{g}$ be the Lie algebra of $SU(N)$.
\begin{enumerate}
\item[(a)] $\mathfrak{X}$ is the set of bonds of $\Omega_{k+1}$ other than the bonds $b_0(c)$, one for each block $c$; these are the bonds that are not free in Ba{\l}aban's construction of the step (\cite[\S1]{B-CMP119} for $k=0$). The set $\mathfrak{X}$ coincides with the set of free bonds of the domain $\Omega^{(k)}_{k+1}$, in the sense of the definition of the free bonds of a domain. This is the domain of step $k+1$ on which the lower bound $\gamma_0$ for the form $C^*\Delta^{(k)}C$ is in force at level $k$, by the level-$k$ lower-bound statement and the exponential-gauge proposition.
\item[(b)] The constrained fluctuation form is
\begin{equation*}
\mathcal{A}:=C^*\Delta^{(k)}C\quad\text{on }\ell^2(\mathfrak{X};\mathfrak{g}).
\end{equation*}
Here $\langle B',\Delta^{(k)}B'\rangle=\mathsf{G}+\mathsf{E}$ in the exponential gauge fixing of \cite{B-CMP109}. In the notation of the exponential-gauge proposition, where the form is written $\mathsf{G}(B)+\langle B,\Delta_kB\rangle$, the operator $C$ above is $C_{\flat}$.
\item[(c)] The number $\mathfrak{b}$ is
\begin{equation*}
\mathfrak{b}:=\mathfrak{b}_k:=g_k^{-1}\delta_k=A_1\bigl(\log g_k^{-2}\bigr)^{p_0},
\end{equation*}
where $\delta_k$ is Ba{\l}aban's small-field parameter of level $k$.
That is, $\mathfrak{b}$ equals the small-field threshold $T_k$, so that $\delta_k=g_kT_k$; the operator $T$ of (d) is unrelated to this threshold.
\item[(d)] Fix one output term of the expansion of step $k+1$ treated in this section. The output terms of that expansion are indexed by labels $(R,I^*)$; to the fixed term belong a region $\Lambda^*$, equal to $\Lambda_{k+1}$, and a family of cubes, its $R^*$-cubes, fixed with the expansion. Put $\Lambda:=\Lambda^*$.
\begin{itemize}
\item $\bar\Lambda$ is the set of bonds of $\mathfrak{X}$ with at least one end-point in $\Lambda$.
\item $\mathfrak{r}:=\mathfrak{X}\setminus\bar\Lambda$.
\item The \emph{rim datum} of the term is the restriction $x:=A_k\lceil\mathfrak{r}$ of $A_k$ to $\mathfrak{r}$.
\end{itemize}
With respect to the decomposition $\mathfrak{X}=\bar\Lambda\sqcup\mathfrak{r}$ write $\mathcal{A}_{\Lambda\Lambda}$, $\mathcal{A}_{\Lambda\mathfrak{r}}$, $\mathcal{A}_{\mathfrak{r}\Lambda}$, $\mathcal{A}_{\mathfrak{r}\mathfrak{r}}$ for the blocks of $\mathcal{A}$; an index $\Lambda$ refers to the bonds of $\bar\Lambda$. Put
\begin{align*}
C_\Lambda&:=\mathcal{A}_{\Lambda\Lambda}^{-1}, & T&:=-C_\Lambda\mathcal{A}_{\Lambda\mathfrak{r}}, & m(x)&:=Tx,
\end{align*}
and let $S$ be the Schur complement,
\begin{equation*}
S:=\mathcal{A}_{\mathfrak{r}\mathfrak{r}}-\mathcal{A}_{\mathfrak{r}\Lambda}C_\Lambda\mathcal{A}_{\Lambda\mathfrak{r}}.
\end{equation*}
\item[(e)] The notation for the rim integral is as follows.
\begin{itemize}
\item $\mu_{C_\Lambda}$ is the Gaussian measure with covariance $C_\Lambda$ on fields $\psi$ on $\bar\Lambda$.
\item For a condition $P$, $\chi(P)$ equals $1$ if $P$ holds and $0$ otherwise.
\item $x\oplus(\psi+m(x))$ is the field on $\mathfrak{X}$ equal to $x$ on $\mathfrak{r}$ and to $\psi+m(x)$ on $\bar\Lambda$.
\item $V(Y;\cdot)$, for $Y$ in the class $\mathbf{D}_k$ of localisation domains, are the terms of the exponent in \cite[p.~276, ll.~31--32]{B-CMP119}.
\end{itemize}
The \emph{rim integral} $\mathsf{I}(x)$, the \emph{upper box} $\mathfrak{b}^{\sharp}$ and the set $\mathfrak{X}^{\sharp}$ are defined by
\begin{equation}\label{5.2:eq-rim-datum-a}
\begin{aligned}
\mathsf{I}(x)&:=\int d\mu_{C_\Lambda}(\psi)\prod_{b\in\bar\Lambda}
\chi\bigl(|\psi_b+m_b(x)|<\mathfrak{b}\bigr)\\
&\qquad\times\exp\sum_{\substack{Y\in\mathbf{D}_k\\ Y\cap\Lambda\neq\emptyset}}
V\bigl(Y;x\oplus(\psi+m(x))\bigr),\\
\mathfrak{b}^{\sharp}(b')&:=
\begin{cases}
\mathfrak{b}^{\flat} & \text{if } \operatorname{dist}_{\rho}(b',\Lambda)\le\varpi,\\
C_{\mathrm{far}}\,\mathfrak{b} & \text{if } b' \text{ lies in an } R^*\text{-cube or off }
\Omega^{\sim-1}_{k+1},\\
\mathfrak{b} & \text{otherwise},
\end{cases}\\
\mathfrak{X}^{\sharp}&:=\bigl\{x \text{ real}:\ |x(b')|\le\mathfrak{b}^{\sharp}(b')
\ \text{for all } b'\in\mathfrak{r}\bigr\}.
\end{aligned}
\end{equation}
In the definition of $\mathsf{I}(x)$ the sum runs over $Y\in\mathbf{D}_k$ with $Y\cap\Lambda\neq\emptyset$, as in \cite[p.~276, ll.~31--32]{B-CMP119}. In the definition of $\mathfrak{b}^{\sharp}$, the lower variant $\mathfrak{b}^{\flat}$ of the small-field box, the width $\varpi$ and the distance $\operatorname{dist}_{\rho}$ are fixed in the rim construction of step $k+1$, and $\Omega^{\sim-1}_{k+1}$ is the region so denoted there; $C_{\mathrm{far}}$ is the constant of the far-bond bound. In the definition of $\mathfrak{X}^{\sharp}$, \emph{real} means that $x$ takes values in $\mathfrak{g}$ and not in its complexification.
\end{enumerate}
\end{definition}

With these definitions, by the Gaussian decomposition identity for the $\Lambda$-integral, applied with $C_\Lambda$, $T$ and $m(x)$ of (d), the $\Lambda$-integral of \cite[(3.25)]{B-CMP119} equals a factor independent of $x$ times $\mathsf{I}(x)$.

The set $\mathfrak{X}^{\sharp}$ contains $x=0$. The factors of the output term of item~(d) that depend on $x$, its rim factors as described in the rim-factor lemma, have support, as functions of $x$, inside $\mathfrak{X}^{\sharp}$. This is the content of clauses (d) and (e) of the rim-factor lemma, together with the far-bond bound, whose constant is $C_{\mathrm{far}}$.

\begin{definition}[Constants of the rim construction in order]\label{5.2:not-rim-constants}
The objects $\mathfrak{X}$, $\mathcal{A}$, $\mathfrak{b}=\mathfrak{b}_k$, $\Lambda$, $\bar\Lambda$, $\mathfrak{r}$,
the rim datum $x$ and the transfer operator $T$ are those of
Definition~\ref{5.2:def-rim-datum}. A bond of $\mathfrak{r}$ is called a \emph{rim bond}. The
constants of the rim construction are chosen in the order
\begin{equation}\label{5.2:eq-rim-constants-a}
\begin{split}
&\theta_S \to L \to \delta \to \{\text{constants depending on } (d,L)\} \to \cdots \to \kappa\\
&\quad\to \kappa_1 \to M_1 \to M \to \mathsf{a} := A_0/A_1 \to C_{\mathrm{far}}\\
&\quad\to \bigl(\gamma_2\text{-family};\ \text{the radii } \alpha;\ \ldots \to \varepsilon_1 \to \cdots\\
&\qquad\to \varepsilon_V^{*} \to \bar C_{\mathsf{W}},\, \bar C_{\sharp} \to K_R \to c_V \to c_A\\
&\qquad\to \mathsf{K}^0,\, \mathsf{K}^{0\prime},\, \mathsf{K}^0_{\mathrm{cell}},\, \alpha_{*}\bigr)
\to \gamma ,
\end{split}
\end{equation}
each letter being chosen after the letters to its left, $\theta_S$ first; the dots stand for
the intermediate letters of the order of choice of the constants. Of the letters inside the
parentheses, those not treated below are constants defined where they are first used in this
book; only their place in the order is fixed here.

\emph{Distances.} The number $\rho(b,b')$ is the bond distance in the current unit lattice, the
distance in which the kernel lemma for the transfer operators states its bounds: the least number
of steps in a chain of successively adjacent bonds leading from $b$ to $b'$,
two bonds being adjacent when they share an endpoint. For a rate
$\delta'>0$ put
\begin{equation*}
C_\rho(\delta') := \sup_b \sum_{b'} e^{-\delta'\rho(b,b')} ,
\end{equation*}
a number depending only on $d$ and $\delta'$; the argument $\delta'$ is a rate and is not the
letter $\delta$ of \eqref{5.2:eq-rim-constants-a}. With
$N_\rho(n) := \sup_b \#\{b' : \rho(b,b') = n\}$, the factorisation statement for the rim integral
writes $C_\rho(\delta')$ for the number
\begin{equation}\label{5.2:eq-rim-constants-1}
1 + \sum_{n \ge 1} N_\rho(n)\, e^{-\delta'(n-1)} ,
\end{equation}
which majorises the supremum above; either number may be read wherever $C_\rho(\delta')$
occurs.

\emph{Constants depending on $(d,L)$.} Put $\gamma_0 := \gamma_0^X(d,L)/c_\Phi(d,L)^2$, with
$\gamma_0^X$ and $c_\Phi$ the constants of the coercivity estimate for the constrained
fluctuation form; with this value the positivity bound in
\eqref{5.2:eq-true-pair-positivity-a} follows from that estimate. The pair
$(B_{\mathcal{A}},\delta_{\mathcal{A}})$ consists of the constants of the decay bound in
\eqref{5.2:eq-true-pair-positivity-a}, and $(B,\delta_W)$ of those of the zeroth-order
random-walk bound. The rate $\mu_0$ is the one given by the Combes--Thomas lemma as a function of
$(\gamma_0, B_{\mathcal{A}}, \delta_{\mathcal{A}})$. Let $(K_0,\delta_1)$ be a pair such that
\begin{equation*}
|T(b,b')|,\ |T_0(b,b')| \le K_0\, e^{-\delta_1\rho(b,b')}
\quad\text{for all } b, b' ,
\end{equation*}
where $T_0$ denotes the zeroth-order transfer operator of the kernel lemma for the transfer
operators. By that kernel lemma and the zeroth-order random-walk bound,
\begin{equation*}
K_0 := \max\Bigl\{\frac{2B_{\mathcal{A}}}{\gamma_0}\, C_\rho(\mu_0/2),\
B^2 C_\rho(\delta_W/2)\Bigr\},
\qquad
\delta_1 := \min\{\mu_0/2,\ \delta_W/2\}
\end{equation*}
is such a pair; it is fixed, and $\Theta_1$ below is defined with it. Put
\begin{equation*}
\Theta_1 := \max\{1,\ K_0 C_\rho(\delta_1)\},
\qquad
\mathfrak{b}^{\flat} := \mathfrak{b}/(4\Theta_1) .
\end{equation*}
By the corollary on the gauge mode, $\Theta_1 \ge \theta_K \ge d(L-1)/2$, with $\theta_K$ the
constant of that corollary, so that the gauge mode at slope $\mathfrak{b}^{\flat}$ has mean at
most $\mathfrak{b}/4$ for every $L$. Put $K_1 := K_0 e^{\delta_1} C_\rho(\delta_1/2)$, and let
$\varpi$ be the least integer with $K_1 e^{-\delta_1\varpi/2} \le 1/16$ (the gauge-mode lemma,
a statement distinct from that corollary, requires of $\varpi$ only that
$\varpi \ge (L-3)/2$). Put $\delta_T := \delta_W/4$,
$\delta_\flat := \min\{\delta_1, \delta_T\}$, $\Theta_\flat := (K_0+1)C_\rho(\delta_\flat)$,
\begin{equation}\label{5.2:eq-rim-constants-2}
\tau_\flat(u) := (K_0+1)\, e^{\delta_\flat}\, C_\rho(\delta_\flat/2)\, e^{-\delta_\flat u/2},
\end{equation}
and $c_0' := \sup \|C_0^{-1}\|$, with $C_0$ the operator so denoted in the decoupling-parameter
lemma (an operator, not the auxiliary parameter $C_0$), the supremum being taken over the
parameters on which $C_0$ depends in that lemma.

\emph{After $\kappa_1$.} Once $\kappa_1$ is fixed, put $B_W := B e^{16\kappa_1}$.

\emph{After $M$.} Put
\begin{equation*}
\varepsilon_T := 2B_W^2\, C_\rho(\delta_T)\bigl(e^{-\delta_W M/4} + \alpha_0 + \alpha_1\bigr),
\qquad
\gamma_2 := \mathsf{c}_2 (LM)^{-4} ;
\end{equation*}
here $\mathsf{c}_2$ is the constant of the upper bound on $\gamma_2$ in \cite{B-CMP116}; that
upper bound is here taken as the definition of $\gamma_2$, so
that there is no circularity with $M$, and it is this number that is meant wherever $\gamma_2$
occurs in this chapter, in particular in the analysis constants of the step. Put further
\begin{equation*}
c_\eta := \Bigl[\frac{\gamma_2}{16(c_0'+1)}\Bigr]^{1/2} \le \frac{1}{16},
\qquad
\mathsf{p}(g) := \frac{\gamma_2 A_1^2 (\log g^{-2})^{2p_0}}{64} ,
\end{equation*}
so that $\mathsf{p}(g_k) = \gamma_2\mathfrak{b}^2/64$. The inequality $c_\eta \le 1/16$ is a
condition on $M$: the choice of $M$ is required to satisfy $(LM)^4 \ge 16\mathsf{c}_2/(c_0'+1)$,
which is equivalent to it.

\emph{After $L$, together with $\kappa_1$.} The parameter $M_1$ is chosen with
$M_1 \ge M_E(d,L)$ and $L^3 M_1 \ge R_E(d,L)$, where $M_E$ and $R_E$ are the thresholds that the
coercivity estimate imposes on its two parameters, denoted $(M,R)$ in that estimate, read here
at $(M,R) := (M_1, L^3)$ (these are that estimate's own letters, not the $M$ of
\eqref{5.2:eq-rim-constants-a} nor a ball radius); $M_1$ also meets the further condition on
$M_1$ imposed later in this chapter.

\emph{Then.} Next in the order, $\mathsf{a}$ is chosen with
$\mathsf{a} \ge \max\{c(d,L),\ c'(d,L)M^2\}$ as in \cite[p.~381]{B-CMP122b};
$C_{\mathrm{far}} := C(d,L,B_3)\,\mathsf{a}$, where $C(d,L,B_3)$ is the constant of the
far-region bound and $B_3$ the constant so denoted in that bound; the radii
$\alpha$ with $M_1\alpha_0 \le a_X(d,L)$, the condition of the coercivity estimate (at the
running radii of \cite{B-CMP119} this is a ceiling in $\gamma$); and $c_A$, one of the letters
inside the parentheses of \eqref{5.2:eq-rim-constants-a} that are not defined here: it enters
$\varrho$ below, and only its place in \eqref{5.2:eq-rim-constants-a} is fixed here.

\emph{Functions of $g$.} Let $\alpha_{*,k}$ be the coupling-dependent radius $\alpha_*$ at
$g_k$, namely $\alpha_{*,k} = g_k C_*(\log g_k^{-2})^{q}$, where $C_*$ is its constant and $q$
the auxiliary exponent; put
\begin{equation}\label{5.2:eq-rim-constants-3}
\varrho := c_A\alpha_{*,k}/g_k,
\qquad
r_A := \varrho/\mathfrak{b} = c_A\,\frac{C_*}{A_1}\,(\log g_k^{-2})^{q-p_0},
\qquad
T^{\mathbb{C}} := \mathfrak{b} + (\Theta_\flat+1)\varrho ,
\end{equation}
where $r_A \to \infty$ as $g_k \to 0$ because $q > p_0$. Let $\varpi_\gamma$ be the least
integer $\ge \varpi$ with $\tau_\flat(\varpi_\gamma)\, r_A \le \tfrac12 c_\eta$. With
$C_\tau := 2(K_0+1)e^{\delta_\flat}C_\rho(\delta_\flat/2)$, for $g_k < 1$,
\begin{equation}\label{5.2:eq-rim-constants-4}
\varpi_\gamma \le \varpi + 1
+ \frac{2}{\delta_\flat}\max\Bigl\{0,\ \log\frac{C_\tau r_A}{c_\eta}\Bigr\} ,
\end{equation}
so that $\varpi_\gamma$ grows at most like a constant multiple of $\log\log g_k^{-2}$. Let
$\mathsf{Q}^A_k$ be the set of $\pi_{k+1}$-cubes holding a rim bond. For $P \subset \bar\Lambda$
the \emph{scaffold rule} is
\begin{equation}\label{5.2:eq-rim-constants-5}
\mathcal{S}(P) := \bigl\{ \Delta \in \mathsf{Q}^A_k :\ \Delta \text{ holds a rim bond } b'
\text{ with } \rho(b,b') \le \varpi_\gamma \text{ for some } b \in P \bigr\},
\end{equation}
a function of $(P, \Lambda, g_k)$. Put $N_A(g) := \sup_b |\mathcal{S}(\{b\})|$, so that
\begin{equation*}
N_A(g) \le \bigl(3 + 2\varpi_\gamma/(LM)\bigr)^d .
\end{equation*}
Let $\mathsf{N}(b)$ be the set of $M$-cubes meeting
$\{b'' : \rho(b,b'') \le \varpi_\gamma\} \cup \bigcup \mathcal{S}(\{b\})$ if
$\mathcal{S}(\{b\}) \ne \emptyset$, and the set of $M$-cubes holding $b$ otherwise; it is
connected. Put $N_Z := \sup_b \#\mathsf{N}(b)$; then $N_Z \le c_d(L + \varpi_\gamma/M)^d$ with a
constant $c_d$ depending only on $d$. Finally
\begin{equation*}
\kappa_A^{\mathrm{pr}}(g) := \mathsf{p}(g)/N_A(g),
\qquad
\kappa_A(k) := \min_{i \le k} \kappa_A^{\mathrm{pr}}(g_i) ;
\end{equation*}
the number $\kappa_A(k)$ is non-increasing in $k$ and does not depend on the ages of the
large-field regions present at level $k$ (the numbers of steps since the steps that created
them). In
$\kappa_A^{\mathrm{pr}}(g_i)$ the quantities $r_A$, $\varpi_\gamma$ and $N_A$ are evaluated with
$g_k$ replaced by $g_i$.
\end{definition}

\begin{proof}
We verify the properties stated above.

\emph{The majorant of $C_\rho$.} The bond distance takes values in the non-negative integers and
vanishes only for $b' = b$. Hence for every $b$
\begin{equation*}
\sum_{b'} e^{-\delta'\rho(b,b')}
= 1 + \sum_{n\ge 1}\#\{b' : \rho(b,b') = n\}\, e^{-\delta' n}
\le 1 + \sum_{n \ge 1} N_\rho(n)\, e^{-\delta'(n-1)} ,
\end{equation*}
and taking the supremum over $b$ gives that \eqref{5.2:eq-rim-constants-1} majorises
$C_\rho(\delta')$.

\emph{$\Theta_1$ and $\mathfrak{b}^{\flat}$.} By definition $\Theta_1 \ge 1$, so
$\mathfrak{b}^{\flat} \le \mathfrak{b}/4$. The inequalities
$\Theta_1 \ge \theta_K \ge d(L-1)/2$ and the bound on the mean of the gauge mode at slope
$\mathfrak{b}^{\flat}$ are the content of the corollary on the gauge mode, applied with the
pair $(K_0,\delta_1)$ fixed above.

\emph{Existence of $\varpi$ and $\varpi_\gamma$.} Since $K_0 > 0$ and $C_\rho(\cdot) \ge 1$
(the term $b' = b$), $K_1 > 0$, and $u \mapsto K_1e^{-\delta_1 u/2}$ is strictly decreasing, tends
to $+\infty$ as $u \to -\infty$ and to $0$ as $u \to +\infty$; hence the least integer $\varpi$
with $K_1 e^{-\delta_1\varpi/2} \le 1/16$ exists. Moreover $\varpi \ge 1$. Indeed, $L$ is an odd
integer greater than $1$, so $d(L-1)/2 \ge d > 1$ and $\Theta_1 > 1$ by the inequality
$\Theta_1 \ge d(L-1)/2$ above; hence $\Theta_1 = K_0C_\rho(\delta_1)$. Since $e^{\delta_1} \ge 1$
and $C_\rho(\delta_1/2) \ge C_\rho(\delta_1)$ (each term of the sum decreases as the rate
increases), $K_1 \ge K_0C_\rho(\delta_1) = \Theta_1 > 1/16$, so that
$K_1e^{-\delta_1 u/2} > 1/16$ for every integer $u \le 0$. Likewise, for $g_k < 1$ one has
$r_A > 0$, and by \eqref{5.2:eq-rim-constants-2} the function $u \mapsto \tau_\flat(u)\, r_A$ is
strictly decreasing with limit $0$ as $u \to \infty$; so integers $u \ge \varpi$ with
$\tau_\flat(u) r_A \le \tfrac12 c_\eta$ exist, and $\varpi_\gamma$, the least of them, is
well defined and satisfies $\varpi_\gamma \ge \varpi \ge 1$.

\emph{The identity for $r_A$ and its growth.} Inserting $\alpha_{*,k}$ into
\eqref{5.2:eq-rim-constants-3} gives $\varrho = c_A C_*(\log g_k^{-2})^{q}$; dividing by
$\mathfrak{b} = A_1(\log g_k^{-2})^{p_0}$ gives the stated expression for $r_A$. Since
$q > p_0$, the exponent $q - p_0$ is positive and $r_A \to \infty$ as $g_k \to 0$.

\emph{The bound \eqref{5.2:eq-rim-constants-4}.} By \eqref{5.2:eq-rim-constants-2},
$\tau_\flat(u) r_A = \tfrac12 C_\tau r_A e^{-\delta_\flat u/2}$ with $C_\tau$ as in the
definition, so $\tau_\flat(u) r_A \le \tfrac12 c_\eta$ holds if and only if
\begin{equation*}
u \ge u_0 := \frac{2}{\delta_\flat}\log\frac{C_\tau r_A}{c_\eta} .
\end{equation*}
Every integer $u \ge \max\{\varpi, u_0\}$ therefore qualifies, and the least such integer is at
most $\max\{\varpi, u_0\} + 1$. If $u_0 < 0$, every integer $u \ge \varpi$ qualifies, so
$\varpi_\gamma = \varpi$, and \eqref{5.2:eq-rim-constants-4} holds because its right-hand side is
at least $\varpi + 1$. If $u_0 \ge 0$, that is, if $C_\tau r_A \ge c_\eta$, the maximum in
\eqref{5.2:eq-rim-constants-4} equals $\log(C_\tau r_A/c_\eta)$, and $\varpi \ge 1$ by the
above, hence $\max\{\varpi,u_0\} \le \varpi + u_0$, which is
\eqref{5.2:eq-rim-constants-4}. As
$\log r_A = \log(c_AC_*/A_1) + (q-p_0)\log\log g_k^{-2}$, the right-hand side of
\eqref{5.2:eq-rim-constants-4} is at most a constant multiple of $\log\log g_k^{-2}$ for
$g_k$ small.

\emph{$c_\eta$ and $\mathsf{p}$.} By definition, $c_\eta \le 1/16$ is equivalent to
$\gamma_2 \le (c_0'+1)/16$, that is, to $(LM)^4 \ge 16\mathsf{c}_2/(c_0'+1)$. Since $c_0'$ is
fixed with the constants depending on $(d,L)$, and $M$ after them in
\eqref{5.2:eq-rim-constants-a}, this is a one-sided lower bound on $M$, which the choice of $M$
is required to meet. At $g = g_k$ one has
$A_1^2(\log g_k^{-2})^{2p_0} = \mathfrak{b}^2$, so $\mathsf{p}(g_k) = \gamma_2\mathfrak{b}^2/64$.

\emph{Dependence of the scaffold rule.} With the set $\mathfrak{X}$ of
Definition~\ref{5.2:def-rim-datum} fixed, the set of rim bonds
$\mathfrak{r} = \mathfrak{X} \setminus \bar\Lambda$ and with it $\mathsf{Q}^A_k$ are determined
by $\Lambda$, and
$\varpi_\gamma$ is determined by $g_k$ through $r_A$; hence \eqref{5.2:eq-rim-constants-5}
depends only on $(P, \Lambda, g_k)$.

\emph{The counts $N_A$ and $N_Z$.} Fix $b$ and put
$\mathcal{B}_b := \{b'' : \rho(b,b'') \le \varpi_\gamma\}$. In each coordinate direction, every
coordinate of an endpoint of a bond of $\mathcal{B}_b$ differs by at most $\varpi_\gamma$ from the
corresponding coordinate of an endpoint of $b$, and the endpoints of $b$ span an interval of
length at most $1$; so the projection of $\mathcal{B}_b$ onto each coordinate axis lies in an
interval $I$ of length at most $2\varpi_\gamma + 1$. A partition of the line into intervals of
length $a$ has at most $\ell/a + 2$ members meeting a given interval of length $\ell$. A cube of
$\mathcal{S}(\{b\})$ is a $\pi_{k+1}$-cube, of side $LM$, holding a bond of $\mathcal{B}_b$, so
its projection meets $I$; hence in each direction at most
\begin{equation*}
\frac{2\varpi_\gamma + 1}{LM} + 2 \le 3 + \frac{2\varpi_\gamma}{LM}
\end{equation*}
projections occur, using $LM \ge 1$, and $|\mathcal{S}(\{b\})| \le (3 + 2\varpi_\gamma/(LM))^d$.
If $\mathcal{S}(\{b\}) \ne \emptyset$, the projection of
$\mathcal{B}_b \cup \bigcup\mathcal{S}(\{b\})$ onto each axis lies in an interval of length at
most $2\varpi_\gamma + 1 + 2LM$, so the number of $M$-cubes meeting this set is at most the $d$-th
power of
\begin{equation*}
\frac{2\varpi_\gamma + 1 + 2LM}{M} + 2
\le \frac{2\varpi_\gamma}{M} + 2L + 3
\le 5\Bigl(L + \frac{\varpi_\gamma}{M}\Bigr),
\end{equation*}
using $M \ge 1$ and $L \ge 1$. If $\mathcal{S}(\{b\}) = \emptyset$, the $M$-cubes holding $b$
number at most
\begin{equation*}
(1/M + 2)^d \le 3^d \le 5^d L^d \le 5^d\bigl(L + \varpi_\gamma/M\bigr)^d ,
\end{equation*}
the last step because $\varpi_\gamma \ge 1$. Taking suprema over $b$ gives the bound on
$N_A$, and the bound on $N_Z$ with the number $5^d$, which depends only on $d$, as the constant.

\emph{Connectedness of $\mathsf{N}(b)$.} Along a shortest chain of adjacent bonds from $b$ to a
bond $b'' \in \mathcal{B}_b$ every bond is at distance at most $\rho(b,b'') \le \varpi_\gamma$
from $b$, so $\mathcal{B}_b$ is connected; each cube of $\mathcal{S}(\{b\})$ holds a bond of
$\mathcal{B}_b$, so $\mathcal{B}_b \cup \bigcup \mathcal{S}(\{b\})$ is connected, and the
$M$-cubes meeting a connected set form a connected family. In the other case the cubes holding
the single bond $b$ all contain $b$.

\emph{Monotonicity of $\kappa_A$.} By definition
\begin{equation*}
\kappa_A(k+1) = \min\{\kappa_A(k),\ \kappa_A^{\mathrm{pr}}(g_{k+1})\} \le \kappa_A(k) .
\end{equation*}
No age of a large-field region enters the definitions of $\mathsf{p}$, $N_A$,
$\kappa_A^{\mathrm{pr}}$ and $\kappa_A$, so $\kappa_A(k)$ does not depend on these ages.
\end{proof}

\begin{definition}[Stored terms with a strong-set datum]\label{5.2:def-stored-terms}
Let $S$ be a set of the lattice, the set over which stored terms are taken in the earlier
definition of stored terms; the letter is unrelated to the Schur complement of
Definition~\ref{5.2:def-rim-datum}.
Let $\mathfrak{U}_c(S)$ denote, for a stored term $c$ as below, the configuration space over $S$ of
the pairs $(\mathbf{U},\mathbf{J})$ of Ba{\l}aban's construction on which stored terms are defined in
the earlier definition of stored terms; the present definition refines that one by a strong-set
datum. Let $\mathcal{W}=\mathcal{W}_k$ be the ball of complex sources
$w=\sum_i z_i f_i$ at the current level $k$. For each level $j'$ we use three objects:
\begin{itemize}
\item $\mathsf{Q}^A_{j'}$, the family of cubes holding a rim bond, as fixed among the constants
of the rim construction in Definition~\ref{5.2:not-rim-constants}, read at level $j'$;
\item $\varrho_{j'}$, the radius $\varrho$ fixed there, read at level $j'$;
\item $\mathfrak{b}^{\sharp}_{j'}(b)$, the upper variant of the small-field box $\mathfrak{b}$ of
Definition~\ref{5.2:def-rim-datum}, read at level $j'$ and at the bond $b$.
\end{itemize}
A \emph{stored term with a strong-set datum} is a function
\begin{equation*}
c=c\bigl(S;(\mathbf{U},\mathbf{J}),(A_{j'})_{j'};w\bigr)
\end{equation*}
together with a declared finite family $\mathcal{S}=(\mathcal{S}_{j'})_{j'}$ of \emph{strong sets}.
Each $\mathcal{S}_{j'}$ is a set of $\mathsf{Q}^A_{j'}$-cubes meeting $S$. The family
$\mathcal{S}$ determines the domain
\begin{equation}\label{5.2:eq-stored-terms-a}
\mathfrak{D}_{\mathcal{S}}(S):=\left\{(A_{j'})_{j'}:\
\begin{aligned}
&A_{j'}(b)\ \text{real},\ |A_{j'}(b)|\le\mathfrak{b}^{\sharp}_{j'}(b),
&&\text{if $b$ lies in a cube of }\mathcal{S}_{j'};\\
&A_{j'}(b)\ \text{complex},\ |A_{j'}(b)|<\varrho_{j'},
&&\text{otherwise}
\end{aligned}
\right\}.
\end{equation}
In \eqref{5.2:eq-stored-terms-a} the bond $b$ runs over all bonds on which $A_{j'}$ is defined,
and ``otherwise'' refers to those among them that do not lie in a cube of $\mathcal{S}_{j'}$.
The following variables are real parameters on which $c$ may depend and which are not displayed
among its arguments:
\begin{itemize}
\item the variables on the bonds of the cubes of $R^{\sim 2}\cap S$, in the notation of
\cite[(1.22)]{B-CMP119}: $R$ is the region entering the smooth cut-off there, and $R^{\sim 2}$
is its two-fold enlargement;
\item all slots of those components, in the decomposition into components used in that earlier
definition of stored terms, which lie outside the two classes of components singled out there.
\end{itemize}
They take real values only and are not among the coordinates in which holomorphy is required
below.

The function $c$ is required to be continuous on
\begin{equation*}
\mathfrak{U}_c(S)\times\mathfrak{D}_{\mathcal{S}}(S)\times\mathcal{W}
\end{equation*}
and holomorphic in every complex coordinate of this product. The size of $c$ and the size of
its strong-set datum are
\begin{equation}\label{5.2:eq-stored-terms-1}
\mathfrak{N}(c):=\sup_{\mathfrak{U}_c(S)\times\mathfrak{D}_{\mathcal{S}}(S)\times\mathcal{W}}|c|,
\qquad
|\mathcal{S}|:=\sum_{j'}|\mathcal{S}_{j'}|,
\end{equation}
where $|\mathcal{S}_{j'}|$ is the number of cubes in $\mathcal{S}_{j'}$.

The configuration $(A_{j'})_{j'}=0$ belongs to $\mathfrak{D}_{\mathcal{S}}(S)$, so every stored
term has a value at $A=0$.

The size $\mathfrak{N}(c)$ in \eqref{5.2:eq-stored-terms-1} is a plain supremum. It carries no
exponential weight in the tree length of $S$. The decay rate of a stored term is carried outside
$\mathfrak{N}(c)$: it appears in the bounds that use $\mathfrak{N}(c)$ and in the anchored-domain
norm $\|\cdot\|^{A}$.
\end{definition}

\begin{definition}[The inductive bound on created boundary terms]\label{5.2:def-created-terms-bound}
Let $\delta$, with $0<\delta\le\delta_J(L)$ and $\delta_J(L)=(1-4/L)/10$, and $\kappa$ be as fixed in
Notation~\ref{5.2:not-rim-constants}. Put $\mathsf{r}(\delta):=(1-10\delta)L/2$, and
define the output rate
\begin{equation*}
\kappa_{\mathrm{out}}:=\mathsf{r}(\delta)\,\kappa=(1-10\delta)\tfrac12 L\kappa .
\end{equation*}
For each level $i$ let $\kappa_A^{\mathrm{pr}}(g_i)$ be the rate fixed in
Notation~\ref{5.2:not-rim-constants}, and set
$\kappa_A(k):=\min_{i\le k}\kappa_A^{\mathrm{pr}}(g_i)$. Let $B_0$ be a constant. We say that the
\emph{inductive bound on created boundary terms} holds if the following three requirements are met
for every step $k'$ and every localisation domain $X\in\mathbf{D}_{k'}$.
\begin{enumerate}
\item The terms of $\mathbf{B}^{(k')}$ created under the exact refinement of the rim cut-offs
\eqref{5.2:eq-rim-refinement-a} are finite sums
$\mathbf{B}^{(k')}(X)=\sum_{\mathcal{S}}\mathbf{B}^{(k')}(X,\mathcal{S})$, in which each
$\mathbf{B}^{(k')}(X,\mathcal{S})$ is a stored term with strong-set datum in the sense of
Definition~\ref{5.2:def-stored-terms}, on the domain \eqref{5.2:eq-stored-terms-a}. Each family
$\mathcal{S}$ is declared by the rule fixed in Notation~\ref{5.2:not-rim-constants}, or is a union
of families so declared.
\item With $\mathfrak{N}$, $|\mathcal{S}_{j'}|$ and $|\mathcal{S}|=\sum_{j'}|\mathcal{S}_{j'}|$
as in Definition~\ref{5.2:def-stored-terms},
\begin{equation}\label{5.2:eq-created-terms-bound-a}
\begin{split}
\sum_{\mathcal{S}}&\exp\Big[\kappa_A^{\mathrm{pr}}(g_{k'-1})\,|\mathcal{S}_{k'-1}|
 +\kappa_A(k'-1)\sum_{j'<k'-1}|\mathcal{S}_{j'}|\Big]\\
&\qquad\times\mathfrak{N}\big(\mathbf{B}^{(k')}(X,\mathcal{S})\big)
 \le B_0\,e^{-\kappa_{\mathrm{out}}\,d_{k'}(X)} .
\end{split}
\end{equation}
The newest slot $k'-1$ thus carries the weight $\kappa_A^{\mathrm{pr}}(g_{k'-1})$. All older slots
carry the single weight $\kappa_A(k'-1)$ of the level.
\item Call a created term $\mathbf{B}^{(k')}(X)$ \emph{wrapped} if the cluster that created it
contained a relative potential $V^{\mathrm{rel}}(Y)$, on some region $Y$, in the sense of the
strong-set decoupling statement. Such a term is in the
format of boundary-type terms, clauses (f0)--(f4), with the following data. Its label is $A:=X$,
where $X$ is its own localisation domain in $\mathbf{D}_{k'}$; this is the same $X$ as in
\eqref{5.2:eq-created-terms-bound-a}, in clause (c) of the theorem on the creating step under the
exact refinement of the rim cut-offs \eqref{5.2:eq-rim-refinement-a},
and in clause (iii) of the domain-counting lemma. The output is filed on its own domain because the
domain-counting lemma sees neither the potential $V^{\mathrm{rel}}$ nor where it reads the
background. Its pocket set $\mathcal{P}$ is the set of pockets of its constituents. It reads
$((\mathbf{U},\mathbf{J}),\mathbf{A},w)$ on
\begin{equation*}
X^{+}:=X\cup\{\text{the live strips touching }X\}.
\end{equation*}
Pockets and live strips are those of the format of boundary-type terms.
No created term is bounded at the plain length of a domain that contains a whole strip.
\end{enumerate}
\end{definition}

We now check the lower bound on the output rate. On the window $0<\delta\le\delta_J(L)$ we have
$10\delta\le 1-4/L$, hence $1-10\delta\ge 4/L$ and $\mathsf{r}(\delta)\ge 2$. Since
$\beta_s=1/4$, this gives
\begin{equation*}
\kappa_{\mathrm{out}}=\mathsf{r}(\delta)\kappa\ \ge\ 2\kappa=(1+4\beta_s)\kappa .
\end{equation*}
The quantity $\kappa_{\mathrm{out}}$ is the rate at which freshly created terms come out of the
cluster expansion of \cite[(2.37)--(2.41)]{B-CMP116} when $\delta\le\delta_J(L)$, those displays
being read with the thresholds taken at $e^{5\rho_7}$ and $e^{5\rho_9}$, where
$\rho_m:=(1-m\delta)\tfrac12 L\kappa$.
The window for $\delta$ is the one fixed by the lemma on the threshold $\delta_J(L)$. Because
$d_{k'}(X)\ge0$, the bound
\eqref{5.2:eq-created-terms-bound-a} implies the same bound with
$B_0e^{-(1+4\beta_s)\kappa d_{k'}(X)}$ as right member. In what follows the requirement is read
in this weaker form.

Later uses of \eqref{5.2:eq-created-terms-bound-a} at a level $k\ge k'-1$ weight a family
$\mathcal{S}$, whose slots are those appearing in \eqref{5.2:eq-created-terms-bound-a}, by
$e^{\kappa_A(k)|\mathcal{S}|}$. This weight is dominated by the weight in the display. Indeed,
$\kappa_A$ is non-increasing in the level, and by its definition as a minimum over $i\le k'-1$,
\begin{equation*}
\kappa_A(k)\le\kappa_A(k'-1)\le\kappa_A^{\mathrm{pr}}(g_{k'-1}) .
\end{equation*}
Hence
\begin{equation*}
e^{\kappa_A(k)|\mathcal{S}|}
\le\exp\Big[\kappa_A^{\mathrm{pr}}(g_{k'-1})\,|\mathcal{S}_{k'-1}|
 +\kappa_A(k'-1)\sum_{j'<k'-1}|\mathcal{S}_{j'}|\Big].
\end{equation*}

The bound \eqref{5.2:eq-created-terms-bound-a} implies the first inductive hypothesis of the
anchored-domain format, in its refined form, for the terms created by $\mathbf{T}$, read, at a
level $k\ge k'-1$, with the rate parameter of that hypothesis set equal to $\kappa_A(k)$, on its
tight datum. Wrapped terms are read on their
labels: for a wrapped term with label $A_T$ one takes $K_T:=(A_T)^{+}$ as its reading domain,
replaces $d_j(X)$ by $d_j(A_T)$, and, in the length clause (f2) of the format of boundary-type
terms, replaces $\mathsf{d}(e)$ by the relative length $\mathsf{d}^{\mathrm{rel}}(\hat e)$, with
$e$, $\hat e$ and the length $\mathsf{d}(e)$ as in that clause.
Consequently the slot of that hypothesis that is reserved for the classes $\mathbb{R}'$,
$\mathbf{B}'$ of \cite{B-CMP122b} bounds an $\mathbf{R}$-born term at the site (P) of the
expansion rule (Definition~\ref{5.1:def-expansion-rule}) at the relative length only.

That the three requirements hold at every step is established below, in the corollary giving the
inductive bound for created boundary terms.

\begin{lemma}[Random-walk bounds for the conditioned covariances {\cite[Theorem~3.15]{B-CMP99b}}]\label{5.2:lem-random-walk-bounds}
Let $\mathfrak{X}$, the form $\mathcal{A} = C^{*}\Delta_k C$ on $\ell^2(\mathfrak{X};\mathfrak{g})$, the domain $\Lambda = \Lambda^{*}$ and $C_\Lambda = \mathcal{A}_{\Lambda\Lambda}^{-1}$
be as in Definition~\ref{5.2:def-rim-datum}, where the operator $\Delta_k$ is written $\Delta^{(k)}$. The operators $C_\Lambda$ and $C_I$ are the operators of the form \cite[(3.158)]{B-CMP99b} on the domains $\Lambda^{*}$ and $\Lambda^{*}\cap Z_0$ respectively. The operator $C_I$ depends on the decoupling parameters $\sigma$, and we then write $C_I(\sigma)$. Let $C_\Lambda^{1/2}(\sigma)$ be the operator of \cite[(2.7)]{B-CMP116}, and let $\Gamma_{(\sigma)} := \Gamma_k(Z_0,\sigma)$ be the operator introduced with the decoupling parameters in this chapter. A subscript $0$ denotes the value at the reference $\sigma \equiv 0$; thus $C_{I,0} := C_I(0)$, and $\Gamma_{(0)}$ is $\Gamma_{(\sigma)}$ at $\sigma \equiv 0$.

Then the four operators $\mathcal{A}$, $C_\Lambda$, $C_I$ and $C_\Lambda^{1/2}$ have the random-walk expansions of \cite[Theorem~3.15]{B-CMP99b}. This holds also at the bonds where the boundary of $\Lambda^{*}$ meets $\tilde Z_0$, and clauses (a)--(c) of the random-walk input of this chapter hold for the kernels of these operators; in terms of the cited papers they read as follows.
\begin{enumerate}
\item[(a)] The kernels of the four operators satisfy the exponential bound at pairs of bonds, with the constants $B$ and $\delta_W$, of clause~(a) of the random-walk input of this chapter; this bound is a consequence of the walk-wise bound \cite[(1.7)]{B-CMP116}.
\item[(b)] The kernels of the four operators satisfy the bound \cite[(1.11)]{B-CMP116}.
\item[(c)] Consider the walks touching $\sigma_0$. At a pair of bonds $b,b'$, their contribution to the kernels of the four operators is bounded by a constant independent of $M$ times the shifted factor
\begin{equation}\label{5.2:eq-random-walk-bounds-1}
e^{-\delta_W[\mathsf{d}(b,b') - M/2]_+}.
\end{equation}
At pairs with $\mathsf{d}(b,b') \ge \tfrac34 M$ it is bounded by such a constant times $e^{-\delta_W \mathsf{d}(b,b')}$, where the rate is renamed within the existentially chosen $\delta_W$. The contributions of the remaining walks differ from their values at $(\mathbf{U},0)$ by the Cauchy estimate in the variables $(A',\mathbf{J})$ of \cite[(2.16)]{B-CMP116}, at ratio $O(\alpha_0+\alpha_1)$.
\end{enumerate}
\end{lemma}

This is \cite[Theorem~3.15]{B-CMP99b}, together with \cite[(3.154), (3.158), (3.185)]{B-CMP99b}, read for the operators above. The proof is Ba{\l}aban's and is not reproduced. The following facts are taken by statement:
\begin{itemize}
\item the bound \cite[(1.7)]{B-CMP116} with the region $X_0$ of \cite[(1.7)]{B-CMP116} prescribed so that $\tilde X_0 = \tilde Z_0$;
\item the bound for the $\sigma$-family of resolvents of \cite[(2.7)]{B-CMP116} on the Dirichlet domain;
\item the bounds at the reference $\sigma \equiv 0$ and for the operator $C_\Lambda^{1/2}(\sigma)$.
\end{itemize}
The exponential-gauge form of the operators is that of the exponential-gauge proposition of this chapter.

We now explain why the cited theorem applies to the four operators.

\emph{Identification of the operators.} The form $\mathcal{A}$ is expanded through all the operators that determine $\Delta_k$, as in \cite{B-CMP116}. Ba{\l}aban's covariance $C^{(k)}(\Lambda_{k+1})$ equals $\mathcal{A}_{\Lambda\Lambda}^{-1}$: this follows from the two equations of Ba{\l}aban's construction that define this covariance and the $\Lambda$-integral, once that integral is reduced as stated with Definition~\ref{5.2:def-rim-datum}. In the exponential gauge and under the constraint of \cite{B-CMP109}, $C_\Lambda$ and $C_I$ are the covariances of the exponential-gauge representation, written there $\tilde\Phi\tilde C^{\Phi}\tilde\Phi^{*}$, taken on the domains $\Lambda^{*}$ and $\Lambda^{*}\cap Z_0$. In \cite{B-CMP116} the operator $C[x + C^{*}\Delta_k C]^{-1}C^{*}$, with $x$ the parameter of \cite{B-CMP116} and not the rim datum, is represented by the integral \cite[(3.185)]{B-CMP99b} with the measure $\mu(QA)$ of \cite{B-CMP116}. At $\sigma = \mathbb{1}$ the operator $C_\Lambda^{1/2}(\sigma)$ of \cite[(2.7)]{B-CMP116} is the square root of $C_\Lambda$ supplied by the square-root lemma.

\emph{The hypothesis on the domain.} The whole hypothesis that \cite[Theorem~3.15]{B-CMP99b} places on its domain is the sentence of \cite[p.~427]{B-CMP99b}:
\begin{quote}
$\Lambda \subset \Lambda_k = \Omega_k^{(k)}$, $\Lambda$ is a union of big blocks, and a distance between $\Lambda$ and $\Lambda_k^c$ is bigger than $RM$
\end{quote}
Here Ba{\l}aban's $\Lambda_k$ of \cite{B-CMP99b} is the domain $\Omega_k$ of \cite{B-CMP119}, by the lemma of this chapter that identifies these two domains. In the quoted sentence $R$ and $M$ are the parameters of \cite{B-CMP99b}, read here as $L^{3}$ and $M_1$; in particular the block size $M$ of \cite{B-CMP99b} is the size $M_1$ of \cite{B-CMP119}. Since $M \ge L^{3}M_1$, the $M$-cubes of $T^{(k)}$ are unions of $M_1$-cubes of $T^{(k+1)}$.

The domains $\Lambda^{*}$ and $\Lambda^{*}\cap Z_0$ are unions of $M$-cubes inside $\Omega^{\sim -1}_{k+1}$. They lie at a depth of at least seven layers of $LM\check R_{k+1}$-cubes inside $\Omega_k$, by \cite[(3.2)--(3.5)]{B-CMP119}. The hypothesis is therefore met, whether or not the boundary of $\Lambda^{*}$ crosses $\tilde Z_0$.

This is where the domain condition among the conditions \eqref{5.2:eq-form-positivity-a} is met for the covariances $C_\Lambda$ and $C_I$. For the positivity of $\mathcal{A}$ on $\ell^2(\mathfrak{X})$, the condition is met when the domain of \cite[Theorem~3.15]{B-CMP99b} is taken to be $\Omega_{k+1}$, as in the statement of the inputs of this chapter and the exponential-gauge proposition.

The theorem assumes no relation between its domain and the localisation domains of the walks. In \cite[p.~13]{B-CMP116} these domains are chosen
\begin{quote}
to a large extent arbitrarily, they have to be only unions of connected families of $M_1$-cubes
\end{quote}
The whole-lattice setting of \cite{B-CMP116} is the case in which $\Lambda^{*}$ is the whole lattice. In Ba{\l}aban's own words \cite[p.~277]{B-CMP119},
\begin{quote}
the restriction of the integration to $\Lambda_{k+1}$, or the larger class of localization domains, do not matter
\end{quote}
Hence nothing beyond the stated generality of the theorem is used where the boundary of $\Lambda$ meets $\tilde Z_0$.

\emph{The clauses.} Clauses (a)--(c) follow from the walk-wise bound \cite[(1.7)]{B-CMP116}.

Clause (b) is \cite[(1.11)]{B-CMP116}.

For (c), take a walk touching $\sigma_0$. It has a localisation domain meeting a $\sigma$-cube. Its length is therefore at least $\mathsf{d}(b,b')$ minus the width of that one domain. This is the argument of \cite[(3.154)]{B-CMP99b}:
\begin{quote}
Remaining terms correspond to walks \dots\ for which at least one localization $X_i$ intersects $\Omega^c$. Then the exponential factor in (3.108) gives the factor (3.154) (after adjusting a definition of $\delta_0$)
\end{quote}
Here \cite[(3.108)]{B-CMP99b} is the exponential factor of the walk expansion and \cite[(3.154)]{B-CMP99b} the resulting decay factor.

At general pairs this argument gives the shifted factor \eqref{5.2:eq-random-walk-bounds-1}, with a constant independent of $M$. The width allowance $M/2$ is at least the diameter of one localisation domain at scale $M_1$ of \cite{B-CMP116}; for a walk of the expansion in the operator theorem at free current it is at least the diameter of one $M/4$-hull.

At pairs with $\mathsf{d} \ge \tfrac34 M$ we have $\mathsf{d} - M/2 \ge \mathsf{d}/3$, hence
\begin{equation*}
e^{-\delta_W[\mathsf{d}-M/2]_+} \le e^{-(\delta_W/3)\mathsf{d}}
\end{equation*}
there, which is the bound by $e^{-\delta_W\mathsf{d}}$ once $\delta_W/3$ is renamed $\delta_W$. Clause (c) states the unshifted factor $e^{-\delta_W\mathsf{d}(b,b')}$ only at pairs with $\mathsf{d}(b,b') \ge \tfrac34 M$.

The remaining walks differ from their values at $(\mathbf{U},0)$ by the Cauchy estimate \cite[(2.16)]{B-CMP116} in the variables $(A',\mathbf{J})$, at ratio $O(\alpha_0+\alpha_1)$.

\begin{corollary}
In the setting of Lemma~\ref{5.2:lem-random-walk-bounds}, the bounds of clause~(ii) of the input of this chapter on the dependence of $C_I$ and $\Gamma_{(\sigma)}$ on the decoupling parameters hold; with the number $\varepsilon$ of that clause, they read
\begin{equation}\label{5.2:eq-random-walk-bounds-2}
\|C_I(\sigma) - C_{I,0}\| \le \varepsilon, \qquad \|\Gamma_{(\sigma)} - \Gamma_{(0)}\| \le \varepsilon .
\end{equation}
\end{corollary}

\begin{proof}
The corollary follows from the bounds above in two steps. One first applies the two-product estimate of the conditional-mean lemma (the lemma containing \eqref{5.2:eq-conditional-mean-a}). One then applies the Schur test to the difference kernels, in which the rows are indexed by bonds in $Z_0$ and the pairs of bonds entering satisfy $\mathsf{d}(b,b') \ge M$.
\end{proof}

Throughout, $k$ is fixed and $U$ is a configuration satisfying the two standing
hypotheses on the configuration in this section. The second of them fixes the
radius $\alpha_0$ which enters the proof of (G5) below; it is the radius formed from the
regularity bounds \cite[Theorem~1, (9)--(10)]{B-CMP102b} of the minimal configurations. We write
$\mathcal{M}$ for the block-averaging operation of \cite{B-CMP109} (denoted $M$ there). We put
$V:=\mathcal{M}^k(U)$, a configuration on $T^{(k)}$, and $\mathsf{W}:=\mathcal{M}(V)$, a
configuration on $T^{(k+1)}$. The configuration $V$ is the configuration
$V^{(k)}=\mathcal{M}^k(U_{k+1})$ of \cite{B-CMP119} (written $M^k(U_{k+1})$ there), with $U$ in
the place of the minimal configuration $U_{k+1}$ determined by
the sequence of domains of \cite{B-CMP119}. As in \cite{B-CMP119}, the formulas of
\cite[Sect.~2]{B-CMP109} used below are read with the operators determined by this sequence of
domains.

A bond $b$ has end-points $b_-$, $b_+$, and gauge transformations act by
$U^v(b)=v(b_-)U(b)v(b_+)^{-1}$. For $u\in SU(N)$ and $X\in\mathfrak{su}(N)$ we put
$R(u)X:=uXu^{-1}$. Here $\mathfrak{su}(N)$ denotes the traceless Hermitian $N\times N$ matrices,
so that $e^{iX}\in SU(N)$. The norm $|X|=(\operatorname{tr}X^*X)^{1/2}$ on $\mathfrak{su}(N)$ is
preserved by every $R(u)$.

For a site $y$ of $T^{(k+1)}$, $B(y)$ denotes the block of $L^d$ sites of $T^{(k)}$ with centre
$y$. A bond $b$ lies inside $B(y)$, written $b\subset B(y)$, if both its end-points lie in
$B(y)$. For $x\in B(y)$, $G(y,x)$ is the family of contours of \cite[(0.11)]{B-CMP109}, and for any
configuration $Y$ on $T^{(k)}$ we write
\[
Y(y,x):=\mathcal{M}\big(\{Y(\Gamma)\}_{\Gamma\in G(y,x)}\big),
\]
as in \cite[(0.11)]{B-CMP109}, with the convention $Y(y,y)=1$. In particular $V(y,y)=1$.
For a bond $c$ of $T^{(k+1)}$, $b_0(c)$ denotes the bond of $T^{(k)}$ joining the blocks
$B(c_-)$ and $B(c_+)$ whose variable $B'(b_0(c))$ is eliminated by means of the averaging
constraint of \cite[(2.4)]{B-CMP109} (the operator written $Q^{\flat}$ below).

$\Lambda$ denotes a set of sites of $T^{(k)}$, and $\Lambda'$ is the set of block centres
lying in $\Lambda$. A $\Lambda$-bond is a bond with at least one end-point in $\Lambda$.
For a function $\mu$ we write $\mu|_{\Lambda'}$ for its restriction to
$\Lambda'$. For a function $F$ of $B$, analytic near $B=0$ with $F(0)=0$, its \emph{linear part}
is its differential at $B=0$.

Let $B$ be an $\mathfrak{su}(N)$-valued function on $\Lambda$-bonds, $\mu$ one on the sites of
$\Lambda$, and $\nu$ one on $\Lambda'$, each extended by $0$ elsewhere. Write $e^{iB}V$ for the
configuration $b\mapsto e^{iB(b)}V(b)$. For bonds $b$ of $T^{(k)}$ and $c$ of $T^{(k+1)}$ define
\begin{align*}
(\bar D\mu)(b)&:=R(V(b))\mu(b_+)-\mu(b_-),\\
(\bar D_1\nu)(c)&:=L^{-1}\big[R(\mathsf{W}(c))\nu(c_+)-\nu(c_-)\big].
\end{align*}

Next, $\mathfrak{g}(B)(y,x)$ is the linear part of
\[
B\mapsto \tfrac1i\log\big[(e^{iB}V)(y,x)\,V(y,x)^{-1}\big].
\]
This is the variable $\tilde B'(y,x)$ of \cite[(2.5)]{B-CMP109}. Likewise $Q^{\flat}B(c)$ is
$L^{-1}$ times the linear part of
\[
B\mapsto\tfrac1i\log\big[\mathcal{M}(e^{iB}V)(c)\,\mathsf{W}(c)^{-1}\big].
\]
Thus $LQ^{\flat}$ is the linear part of $\tilde Q$ in \cite[(2.4)]{B-CMP109}, and $Q^{\flat}$ is
the operator $Q$ there.

With $y$ ranging over the sites of $T^{(k+1)}$ and $x$ over $B(y)$, we further put
\begin{align*}
\mathfrak{a}(B)(y)&:=L^{-d}\sum_{x\in B(y)}\mathfrak{g}(B)(y,x),\\
(Q_1'\mu)(y)&:=L^{-d}\sum_{x\in B(y)}R(V(y,x))\mu(x),\\
(E\nu)(x)&:=R(V(y,x))^{-1}\nu(y),\\
(\mathfrak{d}\mu)(y,x)&:=R(V(y,x))\mu(x)-\mu(y).
\end{align*}
The operator $Q_1'$ is that of \cite[(3.19)]{B-CMP99b}. We also put
\[
Q_1:=Q^{\flat}+\bar D_1\mathfrak{a}.
\]
This is the linear part of the one-step average of \cite[(3.113)]{B-CMP99b}.

We set $\mathcal{N}:=\ker Q_1'$, and we let $\mathfrak{F}$ be the space of
$\mathfrak{su}(N)$-valued functions on the pairs $(y,x)$ with $x\in B(y)$, $x\neq y$. We put
\[
\mathsf{G}(B):=\|\mathfrak{g}(B)\|^2
=\sum_{y}\sum_{x\in B(y),\,x\neq y}|\mathfrak{g}(B)(y,x)|^2 .
\]
Thus $\mathsf{G}(B)=G^{(2)}(B)$, and the quadratic part of the gauge-fixing term in
\cite[(2.5)]{B-CMP109} is $\tfrac12\mathsf{G}(B)$.

Let $Q$, $G_1$ and the constant $a$ be the objects so named in
\cite[(3.128)--(3.129)]{B-CMP99b}; this $Q$ is not
the operator $Q^{\flat}$, and the operators $R$ and $D^*$ of \cite{B-CMP99b} used below are not
the adjoint action $R(u)$ above. Put
\[
\Delta^{\circ}_k:=(QG_1Q^*)^{-1}-a,
\qquad
\mathsf{E}(B):=\langle B,\Delta^{\circ}_kB\rangle .
\]
By \cite[(3.128)--(3.129)]{B-CMP99b}, where $H_1$ is characterised by $QH_1B=B$ and
$RD^*H_1B=0$, one has $\mathsf{E}(B)=\langle H_1B,\Delta_1H_1B\rangle$; the operators $H_1$,
$\Delta_1$ are those of the quadratic term of \cite[(2.8)]{B-CMP109}. Thus $\mathsf{E}$ is a
quadratic form
and contains no term in the current; in particular it does not contain the term of
\cite[(1.5)]{B-CMP109} which is linear in $J$. Finally, $\mathfrak{j}:=H_1^*J$, with $J$ the
current of the linear term of \cite[(2.8)]{B-CMP109}.

In the notation of \cite{B-CMP99b}, with $j$ the scale, $\eta$ the lattice spacing and $M$ the
cube size of that paper (here $M_1$), the operator $H$ and the current $J$ of that paper satisfy
the bound
\[
|(H^*J)(b)|\le O(1)\,M\alpha_0\,(L^j\eta)^{-3}
\]
of \cite{B-CMP99b}.

\begin{lemma}[Block averages and five covariance identities]\label{5.2:lem-covariance-identities}
Let $U$ satisfy the two standing hypotheses on the configuration in this
section, and let $V$, $\mathsf{W}$,
$\Lambda$, $\Lambda'$ and the operators above be as defined. For all $B$ on $\Lambda$-bonds,
all $\mu$ on the sites of $\Lambda$, all $\nu$ on $\Lambda'$, all sites $y$ of $T^{(k+1)}$ and
all $x\in B(y)$:
\begin{equation}\label{5.2:eq-covariance-identities-a}
\begin{aligned}
&\text{(G1)}\quad \mathfrak{g}(B+\bar D\mu)=\mathfrak{g}(B)+\mathfrak{d}\mu,\\
&\text{(G2)}\quad Q^{\flat}\bar D\mu=\bar D_1(\mu|_{\Lambda'}),\\
&\text{(G3)}\quad \mathfrak{a}(\bar D\mu)=Q_1'\mu-\mu|_{\Lambda'},
   \qquad Q_1\bar D\mu=\bar D_1Q_1'\mu,\\
&\text{(G4)}\quad \mathfrak{d}E=0,\qquad Q_1'E=1,\qquad (E\nu)|_{\Lambda'}=\nu,\\
&\text{(G5)}\quad |\mathfrak{g}(B)(y,x)|,\ |\mathfrak{a}(B)(y)|,\ |(E\mathfrak{a}(B))(x)|
   \le 2dL\sup_{b\subset B(y)}|B(b)|.
\end{aligned}
\end{equation}
Moreover, in (G5), $\mathfrak{g}(B)(y,x)$ and $\mathfrak{a}(B)(y)$ depend only on the values of
$B$ on bonds inside $B(y)$. None of these bonds is one of the bonds $b_0(c)$.

The second identity of (G3) is \cite[(3.114)]{B-CMP99b}.
\end{lemma}

\begin{proof}
\emph{The infinitesimal gauge transformation.} Fix $\mu$ and put $v_t:=e^{it\mu}$ for real $t$,
so that $v_t=1$ off $\Lambda$. For every $u\in SU(N)$ and $X\in\mathfrak{su}(N)$ we use
$u\,e^{-itX}u^{-1}=e^{-itR(u)X}$.
This gives
\begin{align*}
(e^{itB}V)^{v_t}(b)
&=e^{it\mu(b_-)}e^{itB(b)}V(b)e^{-it\mu(b_+)}\\
&=e^{it\mu(b_-)}e^{itB(b)}e^{-itR(V(b))\mu(b_+)}V(b).
\end{align*}
The product of the three exponentials is $\exp\{it(B(b)+\mu(b_-)-R(V(b))\mu(b_+))+O(t^2)\}$.
Hence
\[
(e^{itB}V)^{v_t}(b)=e^{it(B-\bar D\mu)(b)+O(t^2)}V(b).
\]
Equivalently, $(e^{itB}V)^{v_t}=e^{itB_t}V$ for a family $B_t$ of functions on
$\Lambda$-bonds with $B_t=B-\bar D\mu+O(t)$. On bonds which are not $\Lambda$-bonds both sides
equal $V(b)$, since there $\mu$, $B$ and $\bar D\mu$ vanish.

\emph{Covariance of the averages.} Every $\Gamma\in G(y,x)$ transforms as
$Y^{v}(\Gamma)=v(y)Y(\Gamma)v(x)^{-1}$. By \cite[(0.6)]{B-CMP109}, so does the variable of
\cite[(0.11)]{B-CMP109}:
\[
Y^{v}(y,x)=v(y)\,Y(y,x)\,v(x)^{-1}.
\]
By the gauge covariance of the averages \cite[Sect.~2]{B-CMP109}, the block average transforms
as
\[
\mathcal{M}(Y^{v})(c)=v(c_-)\,\mathcal{M}(Y)(c)\,v(c_+)^{-1}.
\]
We apply both relations to $Y=e^{itB}V$ and $v=v_t$, and use the previous step.

\emph{Proof of (G1) in \eqref{5.2:eq-covariance-identities-a}.} By the definition of
$\mathfrak{g}$,
\[
(e^{iB''}V)(y,x)V(y,x)^{-1}=1+i\,\mathfrak{g}(B'')(y,x)+O(|B''|^2)
\quad\text{as } B''\to 0.
\]
With $B''=tB_t$ the left side of
\[
(e^{itB_t}V)(y,x)\,V(y,x)^{-1}=v_t(y)\,(e^{itB}V)(y,x)\,v_t(x)^{-1}V(y,x)^{-1}
\]
is $1+it\,\mathfrak{g}(B-\bar D\mu)(y,x)+O(t^2)$. The right side equals
\[
v_t(y)\,\big[(e^{itB}V)(y,x)V(y,x)^{-1}\big]\,e^{-itR(V(y,x))\mu(x)}.
\]
Its expansion is
\[
1+it\big[\mu(y)+\mathfrak{g}(B)(y,x)-R(V(y,x))\mu(x)\big]+O(t^2).
\]
Comparing the derivatives at $t=0$ gives
\[
\mathfrak{g}(B-\bar D\mu)(y,x)=\mathfrak{g}(B)(y,x)+\mu(y)-R(V(y,x))\mu(x)
=\mathfrak{g}(B)(y,x)-(\mathfrak{d}\mu)(y,x).
\]
Since $\bar D$ and $\mathfrak{d}$ are linear in $\mu$, replacing $\mu$ by $-\mu$ gives (G1) for
$x\neq y$. For $x=y$ both sides of (G1) vanish. Indeed $Y(y,y)=1$ for every configuration gives
$\mathfrak{g}(\cdot)(y,y)=0$, and $V(y,y)=1$ gives $(\mathfrak{d}\mu)(y,y)=\mu(y)-\mu(y)=0$.

\emph{Proof of (G2).} By the definition of $Q^{\flat}$,
\[
\mathcal{M}(e^{iB''}V)(c)\mathsf{W}(c)^{-1}=1+iLQ^{\flat}B''(c)+O(|B''|^2).
\]
We multiply the covariance relation
\[
\mathcal{M}(e^{itB_t}V)(c)=v_t(c_-)\,\mathcal{M}(e^{itB}V)(c)\,v_t(c_+)^{-1}
\]
on the right by $\mathsf{W}(c)^{-1}$. We then use
\[
\mathsf{W}(c)v_t(c_+)^{-1}\mathsf{W}(c)^{-1}=e^{-itR(\mathsf{W}(c))\mu(c_+)}
\]
and compare derivatives at $t=0$ as before. This yields
\[
LQ^{\flat}(B-\bar D\mu)(c)=LQ^{\flat}B(c)+\mu(c_-)-R(\mathsf{W}(c))\mu(c_+).
\]
By linearity of $Q^{\flat}$,
\[
LQ^{\flat}\bar D\mu(c)=R(\mathsf{W}(c))\mu(c_+)-\mu(c_-).
\]
The end-points $c_\pm$ are block centres, so $\mu(c_\pm)=(\mu|_{\Lambda'})(c_\pm)$, both members
being zero when $c_\pm\notin\Lambda$. The right side is therefore
$L(\bar D_1(\mu|_{\Lambda'}))(c)$, which is (G2).

\emph{Proof of (G3).} We apply (G1) at $B=0$. Since $\mathfrak{g}(0)=0$, this gives
$\mathfrak{g}(\bar D\mu)(y,x)=(\mathfrak{d}\mu)(y,x)$ for every $x\in B(y)$, the term $x=y$
vanishing on both sides because $V(y,y)=1$. We average over $x\in B(y)$, which has $L^d$ sites:
\[
\mathfrak{a}(\bar D\mu)(y)=L^{-d}\sum_{x\in B(y)}\big[R(V(y,x))\mu(x)-\mu(y)\big]
=(Q_1'\mu)(y)-\mu(y).
\]
Since $\mu$ vanishes off $\Lambda$, we have $\mu(y)=(\mu|_{\Lambda'})(y)$ for every block centre
$y$, and this is the first identity. By
the definition of $Q_1$, then by (G2) and the first identity,
\begin{align*}
Q_1\bar D\mu&=Q^{\flat}\bar D\mu+\bar D_1\mathfrak{a}(\bar D\mu)\\
&=\bar D_1(\mu|_{\Lambda'})+\bar D_1\big(Q_1'\mu-\mu|_{\Lambda'}\big)\\
&=\bar D_1Q_1'\mu .
\end{align*}
This is the relation \cite[(3.114)]{B-CMP99b}.

\emph{Proof of (G4).} Since $V(y,y)=1$, we have $(E\nu)(y)=\nu(y)$, which is
$(E\nu)|_{\Lambda'}=\nu$. For $x\in B(y)$,
\[
(\mathfrak{d}E\nu)(y,x)=R(V(y,x))R(V(y,x))^{-1}\nu(y)-(E\nu)(y)=\nu(y)-\nu(y)=0.
\]
Likewise,
\[
(Q_1'E\nu)(y)=L^{-d}\sum_{x\in B(y)}R(V(y,x))R(V(y,x))^{-1}\nu(y)=\nu(y).
\]

\emph{Proof of (G5).} The contours of $G(y,x)$ stay inside the block $B(y)$, and each has at most
$d(L-1)$ bonds. This is the lemma on the contours $G(y,x)$. Hence
$\mathfrak{g}(B)(y,x)$, and with it $\mathfrak{a}(B)(y)$, depends only on the values of $B$ on
bonds inside $B(y)$.

A bond $b_0(c)$ joins two distinct blocks, so it is not inside any $B(y)$. Therefore
$\mathfrak{a}(B)$ does not depend on the values of $B$ on the bonds $b_0(c)$.

Now let $\Gamma$ consist of the bonds $b_1,\dots,b_n$, traversed in this order, with
$n\le d(L-1)$. Expanding the ordered product $(e^{iB}V)(\Gamma)$ to first order, the linear part
of $\frac1i\log[(e^{iB}V)(\Gamma)V(\Gamma)^{-1}]$ is a sum of $n$ terms. Each term is
$B(b_j)$ conjugated by the adjoint action of a product of the $V(b_i)$, $i<j$ (with the
orientation of $b_j$ in $\Gamma$). Since each $R(u)$ preserves $|\cdot|$, this linear part has
norm at most $d(L-1)\sup_{b\subset B(y)}|B(b)|$.

By the lemma identifying the block contour variables with the averages of
\cite[(0.11)]{B-CMP109}, the linear part of $\mathcal{M}$ at $V$ is the mean of these linear
parts over the contours $\Gamma\in G(y,x)$, up to a term $O(L^2\alpha_0)$ times a linear
function of them, with $\alpha_0$ the radius fixed by the second standing hypothesis. This gives
\[
|\mathfrak{g}(B)(y,x)|\le 2dL\sup_{b\subset B(y)}|B(b)|.
\]
The bound for $\mathfrak{a}(B)(y)$ follows, since it is an average of $L^d$ such values. Finally,
for $x\in B(y)$,
\[
|(E\mathfrak{a}(B))(x)|=|R(V(y,x))^{-1}\mathfrak{a}(B)(y)|=|\mathfrak{a}(B)(y)|.
\]
\end{proof}

The operators $Q_1$ and $Q^{\flat}$ are different. By (G2) and (G3),
\[
(Q_1-Q^{\flat})\bar D\mu=\bar D_1\big(Q_1'\mu-\mu|_{\Lambda'}\big),
\]
and this does not vanish in general. The one-step constraint $Q_1$ of
\cite[(3.113)--(3.114)]{B-CMP99b} is built from block means $Q_1'\mu$. The constraint
$Q^{\flat}$ of \cite[(2.4)]{B-CMP109} is built from the values at the block centres: by (G2),
$Q^{\flat}\bar D\mu$ is a difference of centre values of $\mu$; this is the content of the lemma
on the averaging constraint for exact forms, which expresses $Q^{\flat}$ on lattice gradients as
$L^{-1}$ times differences of centre values, and on which the form $\mathcal{A}$ below is built.
By the definition
$Q_1=Q^{\flat}+\bar D_1\mathfrak{a}$, the two agree on every $B$ with $\mathfrak{a}(B)=0$.
Let $B'=CB$ be the solution of the constraint $\delta(Q^{\flat}B')$ obtained in
\cite{B-CMP109} by eliminating the variables $B'(b_0(c))$, and let $\Delta^{(k)}$ be the operator
of \cite{B-CMP109} for which the covariance $C^{(k)}=C^{(k)}(U_{k+1})$ of the fluctuation
integral there is $C^{(k)}=(C^*\Delta^{(k)}C)^{-1}$. The quadratic form
$\mathcal{A}:=C^*\Delta^{(k)}C$ is then built on the constraint $Q^{\flat}$, whereas the
proposition on the fluctuation form under the one-step constraint is formulated on
$\{Q_1B=0\}$. On the slice on which the axial gauge condition is imposed by $\delta$-functions
in Ba{\l}aban's construction one has $\mathfrak{a}=0$, and there the two constraints agree.

\begin{lemma}[Inverting the block difference map]\label{5.2:lem-difference-inverse}
Let $V$, the blocks $B(y)$, the operators $\bar D$, $\mathfrak{g}$, $Q_1'$, $\mathfrak{d}$ and the spaces
$\mathcal{N}=\ker Q_1'$ and $\mathfrak{F}$ be as in Lemma~\ref{5.2:lem-covariance-identities}. Then the
map $\mathfrak{d}\colon\mathcal{N}\to\mathfrak{F}$ is bijective, and
\begin{equation*}
\|\mathfrak{d}^{-1}\|_{\infty\to\infty}\le 2,
\qquad
\|\mathfrak{d}\mu\|_2\ge\|\mu\|_2\quad(\mu\in\mathcal{N}).
\end{equation*}
The map $\mu_*:=-\mathfrak{d}^{-1}\mathfrak{g}$, with $\mathfrak{g}(B)$ regarded as an element of
$\mathfrak{F}$, has the following properties:
\begin{equation}\label{5.2:eq-difference-inverse-a}
\begin{aligned}
&\text{(m1)}\quad \mu_*\ \text{is linear and block-local;}\\
&\text{(m2)}\quad |\mu_*(B)|\le 4dL\,\sup|B|,\ \text{and each bond feeds at most $L^d$ sites;}\\
&\text{(m3)}\quad \mu_*\ \text{is $\mathcal{N}$-valued;}\\
&\text{(m4)}\quad \mu_*(\bar D v)=-v\quad\text{for } v\in\mathcal{N}.
\end{aligned}
\end{equation}
Moreover, for every $B$ and every $\mu\in\mathcal{N}$,
\begin{equation}\label{5.2:eq-difference-inverse-1}
\mathfrak{g}(B+\bar D\mu)=\mathfrak{d}\bigl(\mu-\mu_*(B)\bigr).
\end{equation}
\end{lemma}

\begin{proof}
Every site of $\Lambda$ lies in exactly one block $B(y)$, $y\in\Lambda'$, and $V(y,y)=1$;
$R(V(y,x))$ acts isometrically on the Lie algebra, and $|\cdot|$ denotes its norm.
Fix $f\in\mathfrak{F}$. For $\mu$ on the sites of $\Lambda$, the equation $\mathfrak{d}\mu=f$ says
that for every block $B(y)$ and every $x\in B(y)$ with $x\neq y$,
\begin{equation*}
R(V(y,x))\,\mu(x)=\mu(y)+f(y,x).
\end{equation*}
For such $\mu$, inserting this into the definition of $Q_1'$ (the term $x=y$ contributes $\mu(y)$,
and $B(y)$ has $L^d$ sites) gives
\begin{equation*}
(Q_1'\mu)(y)=L^{-d}\Bigl[\mu(y)+\sum_{x\neq y}\bigl(\mu(y)+f(y,x)\bigr)\Bigr]
=\mu(y)+L^{-d}\sum_{x\in B(y),\,x\neq y}f(y,x).
\end{equation*}
Hence $Q_1'\mu=0$ fixes $\mu(y)=-L^{-d}\sum_{x\neq y}f(y,x)$, so that
$|\mu(y)|\le\sup|f|$, and then
\begin{equation*}
\mu(x)=R(V(y,x))^{-1}\bigl(\mu(y)+f(y,x)\bigr),\qquad |\mu(x)|\le 2\sup|f|.
\end{equation*}
Conversely, the $\mu$ so defined satisfies $\mathfrak{d}\mu=f$ by construction, and
$Q_1'\mu=0$ by the formula for $(Q_1'\mu)(y)$ above. This gives existence and uniqueness of
$\mu\in\mathcal{N}$ with $\mathfrak{d}\mu=f$, and the bound $\|\mathfrak{d}^{-1}\|_{\infty\to\infty}\le 2$.

For the $\ell^2$ bound let $\mu\in\mathcal{N}$ and $f=\mathfrak{d}\mu$. On one block, using
the isometry of $R(V(y,x))$,
\begin{equation*}
\begin{split}
\sum_{x\in B(y)}|\mu(x)|^2
&=|\mu(y)|^2+\sum_{x\neq y}|\mu(y)+f(y,x)|^2\\
&=L^d|\mu(y)|^2+2\Bigl\langle\mu(y),\sum_{x\neq y}f(y,x)\Bigr\rangle+\sum_{x\neq y}|f(y,x)|^2\\
&=\sum_{x\neq y}|f(y,x)|^2-L^d|\mu(y)|^2,
\end{split}
\end{equation*}
where the last step uses $\sum_{x\neq y}f(y,x)=-L^d\mu(y)$. Summing over the blocks gives
$\|\mu\|_2^2\le\|\mathfrak{d}\mu\|_2^2$.

(m1): $\mathfrak{g}$ is linear, being a linear part, and so is $\mathfrak{d}^{-1}$; the inverse
constructed above acts block by block, and by the bound in \eqref{5.2:eq-covariance-identities-a} on
the reading set of $\mathfrak{g}$, the values $\mathfrak{g}(B)(y,\cdot)$ read only bonds inside $B(y)$.
(m2): by the same bound in \eqref{5.2:eq-covariance-identities-a},
$|\mathfrak{g}(B)|\le 2dL\sup_{B(y)}|B|$, and $\|\mathfrak{d}^{-1}\|_{\infty\to\infty}\le 2$ gives
the constant $4dL$; since $\mu_*(B)$ on $B(y)$ depends only on bonds inside $B(y)$, a bond influences
$\mu_*(B)$ only at the $L^d$ sites of the block containing it. (m3): $\mathfrak{d}^{-1}$ takes values
in $\mathcal{N}$. (m4): for $v\in\mathcal{N}$ the covariance identity
$\mathfrak{g}(B+\bar D\mu)=\mathfrak{g}(B)+\mathfrak{d}\mu$ of \eqref{5.2:eq-covariance-identities-a},
taken at $B=0$ (where $\mathfrak{g}(0)=0$ by linearity), gives $\mathfrak{g}(\bar D v)=\mathfrak{d}v$,
so $\mu_*(\bar Dv)=-\mathfrak{d}^{-1}\mathfrak{d}v=-v$.

Finally, for $\mu\in\mathcal{N}$ the same covariance identity gives
$\mathfrak{g}(B+\bar D\mu)=\mathfrak{g}(B)+\mathfrak{d}\mu$, and
$\mathfrak{g}(B)=-\mathfrak{d}\mu_*(B)$ by the definition of $\mu_*$; this is
\eqref{5.2:eq-difference-inverse-1}.
\end{proof}

The construction is the algebra of Ba{\l}aban's gauge-fixing construction with the contours
$\Gamma_{y,x}$ replaced by the pairs $(y,x)$; in the present setting $\mu_*(B)$ is the field $\mu(B)$
of \cite[(3.169)]{B-CMP99b}.

\begin{lemma}[Almost-invariance of the fluctuation form]\label{5.2:lem-almost-invariance}
Let $\mathfrak{B}$ be the determining set of $(\Omega_j)_{j\le k}$, and let $\mathfrak{B}'$ be a
determining set equal to or coarser than $\mathfrak{B}$, that is,
$M_{\mathfrak{B}'}=\mathsf{M}\circ M_{\mathfrak{B}}$ for a map $\mathsf{M}$. Let
$U=U_{\mathfrak{B}'}(\mathsf{V}')$ be a scaffold minimiser for $\mathfrak{B}'$, in the sense of
Definition~\ref{5.1:def-source-free-scaffold}: a minimiser of the action $A$ over the regular
configurations $U'$ with $M_{\mathfrak{B}'}U'=\mathsf{V}'$. This covers
$U_{k+1}$, every later $U_m$ with $m\ge k+1$, the preparatory minimisers $U_k^{(n)}$, and the
inner backgrounds of \cite[(1.71)]{B-CMP122b}. Let $\Lambda\subset\Omega_k$ lie at depth at
least $R_{\Lambda}M_1$ inside $\Omega_k$, where $R_{\Lambda}$ is the depth constant taken over
from \cite{B-CMP99b}, and let the Lie algebra $\mathfrak{g}$, the field
$V:=M^k(U)$ on $T^{(k)}$, the $\Lambda$-bonds, the operator $\bar D$, the adjoint action $R$, the
fluctuation form $\mathsf{E}$ and the current $\mathfrak{j}$ be as in
Lemma~\ref{5.2:lem-covariance-identities}. Then for every $\mathfrak{g}$-valued $B$ on the
$\Lambda$-bonds and every $\mathfrak{g}$-valued $\mu$ on the sites of $\Lambda$,
\begin{equation}\label{5.2:eq-almost-invariance-1}
\langle \bar D\mu,\mathfrak{j}\rangle=0,\qquad
\mathsf{E}(B-\bar D\mu)=\mathsf{E}(B)-2\langle \bar q(B,\mu),\mathfrak{j}\rangle ,
\end{equation}
where, for a bond $b$ with initial point $b_-$ and final point $b_+$,
\begin{equation}\label{5.2:eq-almost-invariance-2}
\bar q(B,\mu)(b):=-\tfrac{i}{2}\bigl[B(b),\,\mu(b_-)+R(V(b))\mu(b_+)\bigr]
-\tfrac{i}{2}\bigl[\mu(b_-),\,(\bar D\mu)(b)\bigr].
\end{equation}
\end{lemma}

\begin{proof}
We first identify $U$ as a minimiser for the finer determining set $\mathfrak{B}$. Let $U'$ be
regular with $M_{\mathfrak{B}}U'=M_{\mathfrak{B}}U$. Then
\begin{equation*}
M_{\mathfrak{B}'}U'=\mathsf{M}(M_{\mathfrak{B}}U')=\mathsf{M}(M_{\mathfrak{B}}U)
=M_{\mathfrak{B}'}U=\mathsf{V}'.
\end{equation*}
Hence the set $\{M_{\mathfrak{B}}=M_{\mathfrak{B}}U\}$ is contained in
$\{M_{\mathfrak{B}'}=\mathsf{V}'\}$, and it contains $U$. Since $U$ minimises the action on the
larger set, it minimises it on the smaller one. The constrained problem has a unique critical
orbit by \cite[Theorem~1]{B-CMP102b}, so $U=U_{\mathfrak{B}}(M_{\mathfrak{B}}U)$. Define
\begin{equation}\label{5.2:eq-almost-invariance-3}
f(\mathsf{V}):=A(U_{\mathfrak{B}}(\mathsf{V}))
=\min\{A(U') : U'\ \text{regular},\ M_{\mathfrak{B}}(U')=\mathsf{V}\}.
\end{equation}

Next we show that $f$ is gauge invariant. Let $v$ take values in $SU(N)$ on $\Lambda$ and equal
$1$ elsewhere. Let $u$ be any gauge transformation on the sites of the fine lattice underlying
$T^{(k)}$, with $u=v$ on $T^{(k)}\cap\Lambda$ and $u=1$ on the other sites of $\mathfrak{B}$.
For a constraint datum $\mathsf{V}$ we write $\mathsf{V}^{v}$ for the datum whose component on a
bond $b$ of $T^{(k)}$ is $v(b_-)\mathsf{V}(b)v(b_+)^{-1}$ and whose other components are those of
$\mathsf{V}$. Then three facts hold.
\begin{itemize}
\item The constraint is covariant: $M_{\mathfrak{B}}(U'^{\,u})=M_{\mathfrak{B}}(U')^{v}$.
\item The action is invariant: $A(U'^{\,u})=A(U')$.
\item The class of regular configurations is invariant under $u$ \cite{B-CMP102b}.
\end{itemize}
The lower constraints, imposed on $\Lambda_j$ for $j<k$, are untouched by $u$, since
$\Lambda\subset\Omega_k$ lies at depth $R_{\Lambda}M_1$ inside $\Omega_k$. Therefore
$U'\mapsto U'^{\,u}$ maps $\{U'\ \text{regular},\ M_{\mathfrak{B}}U'=\mathsf{V}\}$ bijectively
onto $\{U'\ \text{regular},\ M_{\mathfrak{B}}U'=\mathsf{V}^{v}\}$ and preserves $A$. By
\eqref{5.2:eq-almost-invariance-3},
\begin{equation*}
f(\mathsf{V}^{v})=f(\mathsf{V}).
\end{equation*}

We now expand $f$ to second order. The analysis of \cite[\S G]{B-CMP102b}, namely
\cite[(170)--(171), (174), (178)]{B-CMP102b}, is written for an arbitrary constraint datum
satisfying \cite[(7)]{B-CMP102b} and the corresponding constrained minimiser, not only for the
one of \cite[(2.6)--(2.8)]{B-CMP109}, which is an instance of it. By \cite[(181)]{B-CMP102b}
(see also \cite{B-CMP119}) we therefore obtain the expansion at the datum
$\mathsf{V}=M_{\mathfrak{B}}U$, whose component on the bonds of $T^{(k)}$ is $V=M^k(U)$. For
$\mathfrak{g}$-valued $X$ on the $\Lambda$-bonds, $e^{iX}V$ denotes the datum obtained from
$M_{\mathfrak{B}}U$ by replacing this component $V$ with $e^{iX}V$, the lower components being
left unchanged; then
\begin{equation}\label{5.2:eq-almost-invariance-4}
f(e^{iX}V)=f(V)+\langle X,\mathfrak{j}\rangle+\tfrac12\,\mathsf{E}(X)+O(\|X\|^3).
\end{equation}

We compute the gauge transform of $e^{itB}V$. Put $v_t:=e^{it\mu}$ and, for a bond $b$,
$\mu_-:=\mu(b_-)$ and $\mu_+':=R(V(b))\mu(b_+)$. Since $R$ is the adjoint action, we have
$V(b)e^{-it\mu(b_+)}=e^{-it\mu_+'}V(b)$. Hence
\begin{equation*}
(e^{itB}V)^{v_t}(b)=e^{it\mu_-}e^{itB(b)}e^{-it\mu_+'}\,V(b).
\end{equation*}
The Baker--Campbell--Hausdorff formula to second order gives
\begin{equation*}
\log\bigl(e^{X_1}e^{X_2}e^{X_3}\bigr)=X_1+X_2+X_3
+\tfrac12\bigl([X_1,X_2]+[X_1,X_3]+[X_2,X_3]\bigr)+\cdots,
\end{equation*}
where the dots stand for terms of total degree at least three in $X_1,X_2,X_3$. We apply
it with $X_1=it\mu_-$, $X_2=itB(b)$ and $X_3=-it\mu_+'$. Since $(it)^2/2=-t^2/2$, this yields
\begin{equation*}
\begin{split}
\log\bigl(e^{it\mu_-}e^{itB}e^{-it\mu_+'}\bigr)
&=it(B-\bar D\mu)
-\tfrac{t^2}{2}\bigl([\mu_-,B]-[\mu_-,\mu_+']-[B,\mu_+']\bigr)+O(t^3)\\
&=it(B-\bar D\mu)+\tfrac{t^2}{2}\bigl([B,\mu_-+\mu_+']+[\mu_-,\mu_+']\bigr)+O(t^3).
\end{split}
\end{equation*}
Here $\bar D\mu=\mu_+'-\mu_-$ was used for the linear term. For the quadratic term we use
\begin{equation*}
[\mu_-,(\bar D\mu)(b)]=[\mu_-,\mu_+'-\mu_-]=[\mu_-,\mu_+'].
\end{equation*}
Dividing by $i$ and comparing with \eqref{5.2:eq-almost-invariance-2}, we get
$(e^{itB}V)^{v_t}=e^{iX_t}V$ with
\begin{equation}\label{5.2:eq-almost-invariance-5}
X_t=t(B-\bar D\mu)+t^2\,\bar q(B,\mu)+O(t^3).
\end{equation}
On bonds that are not $\Lambda$-bonds, $v_t=1$ at both end-points and $B=0$, so nothing changes
there.

Finally we compare orders. By gauge invariance, $f(e^{iX_t}V)=f(e^{itB}V)$. We expand both
sides with \eqref{5.2:eq-almost-invariance-4}. Since $\mathsf{E}$ is a quadratic form, we have
$\mathsf{E}(X_t)=t^2\,\mathsf{E}(B-\bar D\mu)+O(t^3)$. Using \eqref{5.2:eq-almost-invariance-5}, this gives
\begin{equation*}
\begin{split}
f(e^{iX_t}V)&=f(V)+t\langle B-\bar D\mu,\mathfrak{j}\rangle+t^2\langle\bar q(B,\mu),\mathfrak{j}\rangle
+\tfrac{t^2}{2}\,\mathsf{E}(B-\bar D\mu)+O(t^3),\\
f(e^{itB}V)&=f(V)+t\langle B,\mathfrak{j}\rangle+\tfrac{t^2}{2}\,\mathsf{E}(B)+O(t^3).
\end{split}
\end{equation*}
Equating the coefficients of $t$ gives $\langle\bar D\mu,\mathfrak{j}\rangle=0$. Equating the
coefficients of $t^2$ gives
\begin{equation*}
\langle\bar q(B,\mu),\mathfrak{j}\rangle+\tfrac12\mathsf{E}(B-\bar D\mu)=\tfrac12\mathsf{E}(B).
\end{equation*}
This is the second identity in \eqref{5.2:eq-almost-invariance-1}.

The argument uses only the second-order expansion \eqref{5.2:eq-almost-invariance-4}, gauge
covariance and the Baker--Campbell--Hausdorff formula; no positivity of $\mathsf{E}$ enters. On
the bonds $b$ of $T^{(k)}$ the current satisfies
$\|\mathfrak{j}(b)\|\le C_H M_1\alpha_0$, by the bound on the current taken over from
\cite{B-CMP99b}, where $C_H$ and $\alpha_0$ are the constants of that bound.
\end{proof}

\begin{corollary}[Gauge invariance on the constraint space]\label{5.2:cor-constrained-invariance}
In the setting of Lemma~\ref{5.2:lem-almost-invariance} (almost-invariance of the fluctuation
form), let $B$ be a field on the $\Lambda$-bonds and $\mu$ a
field on $\Lambda$ with $\mu|_{\Lambda'} = 0$, and assume $\tilde{D}^{(2)} = hC^{(2)}$. Here
$\Lambda'$, the linear map $Q^{\flat}$, the quadratic maps $C^{(2)}$ and $\tilde{D}^{(2)}$, and the
operator $h$ with
$LQ^{\flat}h = 1$, where $L$ is the blocking factor, are those of the definition of the
fluctuation operator $\Delta_k$. Then, exactly,
\begin{equation*}
\langle B-\bar D\mu,\Delta_k(B-\bar D\mu)\rangle = \langle B,\Delta_k B\rangle .
\end{equation*}
\end{corollary}

\begin{proof}
Recall that $\mathsf{M}$ denotes the one-step block-averaging map and that, in exponential
coordinates at $\mathsf{M}(V)$,
the terms of first and of second order in $X$ of $\mathsf{M}(e^{iX}V)$ are $LQ^{\flat}X$ and
$C^{(2)}(X)$, with $L$ the blocking factor.

By \cite{B-CMP109}, the current $\mathfrak{j}$ is orthogonal to the kernel of $Q^{\flat}$. In finite
dimension the orthogonal complement of the kernel of $Q^{\flat}$ is the range of its adjoint.
Hence $\mathfrak{j} = Q^{\flat *}\omega$ for some $\omega$.

Let $v_t := e^{it\mu}$. As in the proof of Lemma~\ref{5.2:lem-almost-invariance},
\begin{equation*}
(e^{itB}V)^{v_t} = e^{iX_t}V,\qquad X_t = t(B-\bar D\mu)+t^2\bar q(B,\mu)+O(t^3).
\end{equation*}
Since $v_t$ is the identity on $\Lambda'$ by the hypothesis $\mu|_{\Lambda'} = 0$, the covariance
of the block-averaging map under gauge transformations used in the proof of
Lemma~\ref{5.2:lem-almost-invariance} gives
\begin{equation*}
\mathsf{M}\bigl((e^{itB}V)^{v_t}\bigr) = \mathsf{M}(e^{itB}V).
\end{equation*}
The left side equals $\mathsf{M}(e^{iX_t}V)$. Its term of order $t^2$ consists of two parts: the
first-order term evaluated at the $t^2$-part of $X_t$, which is $LQ^{\flat}\bar q(B,\mu)$, and
the second-order term evaluated at the $t$-part of $X_t$, which is $C^{(2)}(B-\bar D\mu)$. The
term of order $t^2$ on the right side is $C^{(2)}(B)$. Equating the two gives
\begin{equation*}
LQ^{\flat}\bar q(B,\mu) = C^{(2)}(B)-C^{(2)}(B-\bar D\mu).
\end{equation*}

Since $LQ^{\flat}h$ is the identity \cite{B-CMP109}, applying $LQ^{\flat}$ to the field
\begin{equation*}
h\bigl(C^{(2)}(B)-C^{(2)}(B-\bar D\mu)\bigr)
\end{equation*}
returns $C^{(2)}(B)-C^{(2)}(B-\bar D\mu)$; this field therefore has the same image under
$LQ^{\flat}$ as $\bar q(B,\mu)$, and hence the same image under
$Q^{\flat}$. A field enters a pairing with $\mathfrak{j} = Q^{\flat *}\omega$ only through its image
under $Q^{\flat}$, since $\langle Y,\mathfrak{j}\rangle = \langle Q^{\flat}Y,\omega\rangle$ for every
field $Y$ on the $\Lambda$-bonds.
Therefore
\begin{equation*}
-2\langle \bar q(B,\mu),\mathfrak{j}\rangle
= 2\bigl\langle hC^{(2)}(B-\bar D\mu)-hC^{(2)}(B),\mathfrak{j}\bigr\rangle .
\end{equation*}

Insert this into the second identity of Lemma~\ref{5.2:lem-almost-invariance}. This gives
\begin{equation*}
\mathsf{E}(B-\bar D\mu)
= \mathsf{E}(B)+2\bigl\langle hC^{(2)}(B-\bar D\mu)-hC^{(2)}(B),\mathfrak{j}\bigr\rangle .
\end{equation*}
Since $\tilde{D}^{(2)} = hC^{(2)}$, this rearranges to
\begin{equation*}
\mathsf{E}(B-\bar D\mu)-2\langle \tilde{D}^{(2)}(B-\bar D\mu),\mathfrak{j}\rangle
= \mathsf{E}(B)-2\langle \tilde{D}^{(2)}(B),\mathfrak{j}\rangle .
\end{equation*}
By the definition of the fluctuation operator $\Delta_k$, for every field $X$ on the $\Lambda$-bonds
\begin{equation*}
\langle X,\Delta_k X\rangle = \mathsf{E}(X)-2\langle \tilde{D}^{(2)}(X),\mathfrak{j}\rangle .
\end{equation*}
The two sides of the preceding identity are therefore
$\langle B-\bar D\mu,\Delta_k(B-\bar D\mu)\rangle$ and $\langle B,\Delta_k B\rangle$, which is the
assertion.
\end{proof}

\begin{lemma}[The unipotent map between the two constraints]\label{5.2:lem-unipotent-map}
Work in the setting of Lemma~\ref{5.2:lem-covariance-identities} (block averages and five covariance
identities), and put
\begin{equation}\label{5.2:eq-unipotent-map-1}
\Phi := 1+\bar D E\mathfrak{a},
\end{equation}
acting on $\mathfrak{g}$-valued fields on $\Lambda$-bonds. For a bond $b$ write $b_{+}$, $b_{-}$ for
its end-points.
\begin{enumerate}
\item[(a)] $\mathfrak{a}\bar D E=0$. Hence $(\bar D E\mathfrak{a})^2=0$, $\Phi^{-1}=1-\bar D E\mathfrak{a}$
and $\det\Phi=1$. The maps $\Phi^{\pm1}$ send fields supported on $\Lambda$-bonds to fields supported
on $\Lambda$-bonds, and
\begin{equation}\label{5.2:eq-unipotent-map-2}
|(\Phi^{\pm1}B)(b)|\le(1+4dL)\sup|B|;
\end{equation}
moreover $(\Phi^{\pm1}B)(b)$ depends on $B$ only through $B(b)$ and the values of $B$ on bonds inside
the at most two blocks containing $b_{+}$ and $b_{-}$.
\item[(b)] $\mathfrak{g}\Phi=\mathfrak{g}$, $\mathfrak{a}\Phi=\mathfrak{a}$ and $Q^{\flat}\Phi=Q_1$. In
particular $\Phi$ maps $\{Q_1=0\}$ onto $\{Q^{\flat}=0\}$.
\item[(c)] Let $U$ be as in Lemma~\ref{5.2:lem-almost-invariance} (almost-invariance of the
fluctuation form), and let $\tilde D^{(2)}$ satisfy, with a constant $C_D$, the hypothesis on
quadratic maps in \eqref{5.2:eq-form-positivity-a}. Then, for $B_1$ a $\mathfrak{g}$-valued field on
$\Lambda$-bonds,
\begin{equation}\label{5.2:eq-unipotent-map-3}
\langle\Phi B_1,\Delta_k\Phi B_1\rangle
=\langle B_1,\Delta^{\circ}_k B_1\rangle-2\langle\tilde D^{(2),\Phi}(B_1),\mathfrak{j}\rangle,
\end{equation}
where
\begin{equation}\label{5.2:eq-unipotent-map-4}
\tilde D^{(2),\Phi}(B_1):=\tilde D^{(2)}(\Phi B_1)+\bar q(B_1,-E\mathfrak{a}B_1).
\end{equation}
The map $\tilde D^{(2),\Phi}$ satisfies the same hypothesis with the constant
\begin{equation}\label{5.2:eq-unipotent-map-5}
C_D^{\Phi}:=(1+4dL)^2C_D+c_{\mathfrak{g}}(2dL+4d^2L^2),
\end{equation}
where $c_{\mathfrak{g}}$ denotes the norm of the bracket, $|[X,Y]|\le c_{\mathfrak{g}}|X|\,|Y|$.
\end{enumerate}
\end{lemma}

\begin{proof}
(a) Let $\nu$ be a function on $\Lambda'$. The third covariance identity of
Lemma~\ref{5.2:lem-covariance-identities}, applied to $\mu=E\nu$, gives
$\mathfrak{a}(\bar D E\nu)=Q_1'E\nu-(E\nu)|_{\Lambda'}$. By the fourth identity, $Q_1'E=1$ and
$(E\nu)|_{\Lambda'}=\nu$, so that
\begin{equation*}
\mathfrak{a}\bar D E\nu=\nu-\nu=0 .
\end{equation*}
Hence $\mathfrak{a}\bar D E=0$, and therefore
\begin{equation*}
(\bar D E\mathfrak{a})^2=\bar D E(\mathfrak{a}\bar D E)\mathfrak{a}=0 .
\end{equation*}
Consequently $(1+\bar D E\mathfrak{a})(1-\bar D E\mathfrak{a})=1-(\bar D E\mathfrak{a})^2=1$, and
likewise in the other order, so $\Phi^{-1}=1-\bar D E\mathfrak{a}$. Since $\bar D E\mathfrak{a}$ is
nilpotent, $\Phi$ is unipotent and $\det\Phi=1$.

The field $E\mathfrak{a}B$ lives on the sites of $\Lambda$. By the definition
$(\bar D\mu)(b)=R(V(b))\mu(b_{+})-\mu(b_{-})$ for $\mu$ on the sites of $\Lambda$ (and zero
elsewhere), $(\bar D\mu)(b)$ vanishes unless an end-point of $b$ lies in $\Lambda$. Thus
$\bar D E\mathfrak{a}B$ is supported on $\Lambda$-bonds, and so are $\Phi^{\pm1}B$ whenever $B$ is.

By the fifth covariance identity, $\mathfrak{g}$ and $\mathfrak{a}$ read only bonds inside the block
$B(y)$, and $|E\mathfrak{a}(B)|\le 2dL\sup_{B(y)}|B|$. As $R(V(b))$ acts isometrically on
$\mathfrak{g}$, $|(\bar D\mu)(b)|\le|\mu(b_{+})|+|\mu(b_{-})|$, hence
\begin{equation*}
|(\bar D E\mathfrak{a}B)(b)|\le 4dL\sup|B| ,
\end{equation*}
and \eqref{5.2:eq-unipotent-map-2} follows from $\Phi^{\pm1}B=B\pm\bar D E\mathfrak{a}B$. The same
identity shows that $(\bar D E\mathfrak{a}B)(b)$ reads $E\mathfrak{a}B$ at $b_{+}$ and $b_{-}$ only,
and $(E\mathfrak{a}B)(x)$, for $x\in B(y)$, reads $B$ only on bonds inside $B(y)$. This is the locality
statement.

(b) Let $\nu=\mathfrak{a}B_1$. By the first covariance identity,
$\mathfrak{g}(B_1+\bar D E\nu)=\mathfrak{g}(B_1)+\mathfrak{d}E\nu$. By the fourth identity,
$\mathfrak{d}E=0$, so $\mathfrak{g}(\Phi B_1)=\mathfrak{g}(B_1)$. The relation
$\mathfrak{a}\Phi=\mathfrak{a}$ follows from $\mathfrak{a}\Phi=\mathfrak{a}+\mathfrak{a}\bar D E\mathfrak{a}$
and (a).

By the second covariance identity,
$Q^{\flat}\bar D E\mathfrak{a}B_1=\bar D_1\bigl((E\mathfrak{a}B_1)|_{\Lambda'}\bigr)$. By the fourth,
$(E\mathfrak{a}B_1)|_{\Lambda'}=\mathfrak{a}B_1$. Hence
\begin{equation*}
Q^{\flat}\Phi B_1=Q^{\flat}B_1+\bar D_1\mathfrak{a}B_1=Q_1B_1 ,
\end{equation*}
by the definition $Q_1:=Q^{\flat}+\bar D_1\mathfrak{a}$.

For the inverse, the third covariance identity gives $Q_1\bar D\mu=\bar D_1Q_1'\mu$. With
$\mu=E\mathfrak{a}B$ and $Q_1'E=1$ this yields
\begin{equation*}
Q_1\Phi^{-1}=Q_1-\bar D_1Q_1'E\mathfrak{a}=Q_1-\bar D_1\mathfrak{a}=Q^{\flat}.
\end{equation*}
Thus if $Q_1B_1=0$ then $Q^{\flat}\Phi B_1=0$. Conversely, if $Q^{\flat}B=0$ with $B$ on
$\Lambda$-bonds, then $B_1:=\Phi^{-1}B$ lies on $\Lambda$-bonds by (a), satisfies
$Q_1B_1=Q^{\flat}B=0$, and $\Phi B_1=B$. So $\Phi$ maps $\{Q_1=0\}$ onto $\{Q^{\flat}=0\}$.

(c) We apply Lemma~\ref{5.2:lem-almost-invariance} with $B=B_1$ and $\mu=-E\mathfrak{a}B_1$. Its
hypotheses hold: $U$ is as required there, $B_1$ lives on $\Lambda$-bonds, and $\mu$ lives on the
sites of $\Lambda$ by (a). Since
\begin{equation*}
B-\bar D\mu=B_1+\bar D E\mathfrak{a}B_1=\Phi B_1,
\end{equation*}
that lemma gives
\begin{equation}\label{5.2:eq-unipotent-map-6}
\mathsf{E}(\Phi B_1)=\mathsf{E}(B_1)-2\langle\bar q(B_1,-E\mathfrak{a}B_1),\mathfrak{j}\rangle .
\end{equation}
We insert \eqref{5.2:eq-unipotent-map-6} into the identity \cite[(3.156)]{B-CMP99b}, taken at
$B=\Phi B_1$. The gauge-fixing form enters that identity through $\mathsf{G}(\Phi B_1)$, and by (b)
\begin{equation*}
\mathsf{G}(\Phi B_1)=\|\mathfrak{g}(\Phi B_1)\|^2=\|\mathfrak{g}(B_1)\|^2=\mathsf{G}(B_1).
\end{equation*}
The result is \eqref{5.2:eq-unipotent-map-3}, with $\tilde D^{(2),\Phi}$ given directly by
\eqref{5.2:eq-unipotent-map-4}.

It remains to check the hypothesis on $\tilde D^{(2),\Phi}$. Since $\Phi$ and $E\mathfrak{a}$ are
linear, and $\tilde D^{(2)}$ and $\bar q$ are homogeneous quadratic, $\tilde D^{(2),\Phi}$ is a
homogeneous quadratic map; its bilinear form is obtained by polarisation. By (a) and the locality of
$E\mathfrak{a}$, its value at $b$ reads $B_1$ at $b$ and in the blocks of $b_{\pm}$, composed with
the range of $\tilde D^{(2)}$.

For the bounds, \eqref{5.2:eq-unipotent-map-2} and the hypothesis on $\tilde D^{(2)}$ give
\begin{equation*}
|\tilde D^{(2)}(\Phi B_1,\Phi B_2)(b)|\le(1+4dL)^2C_D\sup|B_1|\,\sup|B_2|,
\end{equation*}
the suprema being taken over the bonds read. By its definition in
Lemma~\ref{5.2:lem-almost-invariance}, each of the two brackets in $\bar q(B,\mu)(b)$ carries a factor
$-i/2$ and one argument that is a sum of two terms. This gives
\begin{equation*}
|\bar q(B,\mu)(b)|\le c_{\mathfrak{g}}\bigl(|B(b)|\,\|\mu\|_\infty+\|\mu\|_\infty^2\bigr).
\end{equation*}
By the fifth covariance identity, $\|E\mathfrak{a}B_1\|_\infty\le 2dL\sup|B_1|$. Hence
\begin{equation*}
|\bar q(B_1,-E\mathfrak{a}B_1)(b)|\le c_{\mathfrak{g}}(2dL+4d^2L^2)(\sup|B_1|)^2 .
\end{equation*}
Each of these quadratic expressions is the diagonal of a bilinear expression
$\beta(B_1,B_2)$ whose terms obey the same bound with $\sup|B_1|\,\sup|B_2|$ in place of
$(\sup|B_1|)^2$. The polarised form $\tfrac12(\beta(B_1,B_2)+\beta(B_2,B_1))$ then obeys that bound
with the same constant, so no polarisation factor arises. Adding the two contributions gives the
hypothesis for $\tilde D^{(2),\Phi}$ with the constant $C_D^{\Phi}$ of \eqref{5.2:eq-unipotent-map-5}.
\end{proof}

\begin{proposition}[Positivity of the constrained fluctuation form]
\label{5.2:prop-form-positivity}
Let $k\ge 0$, and let $(\Lambda,\{\Omega_j\}_{j\le k},U)$ and a map $\tilde D^{(2)}$ be given. A
$\Lambda$-bond is a bond of the unit lattice with at least one end-point in $\Lambda$. For a bond $c$ of
$T^{(k+1)}$, $b_0(c)$ is the bond of \cite[\S2]{B-CMP109} at which the constraint is solved; it joins two
blocks. The other $\Lambda$-bonds are called free. Let $V$, $\mathsf{W}$, $\bar D$, $\bar D_1$,
$\mathfrak{g}$, $Q^{\flat}$, $\mathfrak{a}$, $Q_1'$, $E$,
$\mathfrak{d}$, $Q_1=Q^{\flat}+\bar D_1\mathfrak{a}$, $\mathcal{N}=\ker Q_1'$, $\mathsf{G}=\|\mathfrak{g}(\cdot)\|^2$,
$\Delta^{\circ}_k$ and the current $\mathfrak{j}=H_1^*J$ be the objects fixed together with the covariance
identities \eqref{5.2:eq-covariance-identities-a}. Let $\mu_*$ be the map of
Lemma~\ref{5.2:lem-difference-inverse}, $\bar q(B,\mu)$ the second-order correction of
Lemma~\ref{5.2:lem-almost-invariance}, and $\Phi=1+\bar DE\mathfrak{a}$, $\tilde D^{(2),\Phi}$, $C_D^{\Phi}$ as in
Lemma~\ref{5.2:lem-unipotent-map}.

The following statements are assumed:
\begin{itemize}
\item the exponential-gauge covariance proposition on the block-mean constraint space $\{Q_1B=0\}$, with
its estimate lemma and its walk-expansion clause, namely the existence of the covariance with its kernel
bound, its lower bound, and its walk expansion and analyticity;
\item the inversion of the block difference map, Lemma~\ref{5.2:lem-difference-inverse};
\item the almost-invariance of the fluctuation form, Lemma~\ref{5.2:lem-almost-invariance};
\item the unipotent map between the two constraints, Lemma~\ref{5.2:lem-unipotent-map}.
\end{itemize}
The data $(\Lambda,\{\Omega_j\}_{j\le k},U)$ and $\tilde D^{(2)}$ satisfy the following conditions, whose
quantitative parts are collected in \eqref{5.2:eq-form-positivity-a}.
\begin{itemize}
\item[$(\tilde{\mathrm{H}}1)$] $\Lambda=B(\Lambda')\subset\Omega_k^{(k)}$, where $\Lambda'\subset T^{(k+1)}$ is
a union of big blocks, of size $M_1$; here $B(y)$ denotes the $L$-block of a site $y$, $B(\Lambda')$ the
union of the blocks $B(y)$, $y\in\Lambda'$, $B^k$ the $k$-fold iterate of this block map, and
$\Omega_k^{(k)}$ the set $\Omega_k$ read on $T^{(k)}$, as in \cite[(2.1)--(2.2)]{B-CMP96}. The distance of
$\Lambda$ from the complement of $\Omega_k^{(k)}$ exceeds $R_{\mathrm{depth}}M_1$, with $R_{\mathrm{depth}}$
the depth constant of \cite{B-CMP99b}, as displayed in \eqref{5.2:eq-form-positivity-a}, so that the
extended sequence
$\tilde\Omega_j:=\Omega_j$ $(j\le k)$, $\tilde\Omega_{k+1}:=B^k(\Lambda)$ obeys
\cite[(2.1)--(2.2)]{B-CMP96}.
\item[$(\tilde{\mathrm{H}}2)$] $U$ is regular in the sense of \cite[(3.35)--(3.36)]{B-CMP99b}, with
constant $\alpha_0$, with respect to the cube classes of both $\{\Omega_j\}$ and $\{\tilde\Omega_j\}$.
Moreover $U$ is a scaffold minimiser as in Lemma~\ref{5.2:lem-almost-invariance}. The block size $M_1$ and
the depth $R_{\mathrm{depth}}M_1$ satisfy the one-sided thresholds $M_E(d,L)$, $R_E(d,L)$ of the
exponential-gauge covariance proposition, as displayed.
\end{itemize}
The hypothesis on $D$ is imposed on $\tilde D^{(2)}$: it is a homogeneous quadratic map from unit-lattice
bond fields on $\Lambda$ to functions on $\Lambda$-bonds, with $\tilde D^{(2)}(B_1,B_2)$ its polar bilinear
form. It is block-local at range $c_d$ $L$-blocks: with $N(b)$ the set
of bonds in the $L$-blocks within that range of $b$, it obeys the displayed bound.
Finally $M_1\alpha_0\le a_X(d,L)$, where $a_X$ is the smallness threshold $a_E$ of the exponential-gauge
covariance proposition with $C_D$ replaced by $C_D^{\Phi}$:
\begin{equation}\label{5.2:eq-form-positivity-a}
\begin{aligned}
&(\tilde{\mathrm{H}}1)\quad
  \operatorname{dist}\bigl(\Lambda,(\Omega_k^{(k)})^{c}\bigr)>R_{\mathrm{depth}}M_1,\\
&(\tilde{\mathrm{H}}2)\quad
  M_1\ge M_E(d,L),\qquad R_{\mathrm{depth}}M_1\ge R_E(d,L),\\
&\text{(hypothesis on $D$)}\quad
  |\tilde D^{(2)}(B_1,B_2)(b)|\le C_D\sup_{N(b)}|B_1|\,\sup_{N(b)}|B_2|,\\
&\phantom{\text{(hypothesis on $D$)}}\quad M_1\alpha_0\le a_X(d,L).
\end{aligned}
\end{equation}
Let $\Delta_k$ be the fluctuation form of \cite[(3.156)]{B-CMP99b} with second-order term
$\tilde D^{(2)}$, the form of Lemma~\ref{5.2:lem-unipotent-map}(c), so that
\[
\langle B,\Delta_kB\rangle=\langle B,\Delta^{\circ}_kB\rangle-2\langle\tilde D^{(2)}(B),\mathfrak{j}\rangle.
\]
Write $\Delta^{(k)}_{\exp}$ for the form $\mathsf{G}(B)+\langle B,\Delta_kB\rangle$, which carries the
exponential gauge-fixing term. Let $C_{Q_1}$ be the solution operator of $Q_1B=0$ in the variables
$B(b_0(c))$. Then the following hold.

\smallskip
(a) \emph{Constraint $Q_1$.} The exponential-gauge covariance proposition holds verbatim when its linearised
average, difference map and translation map are replaced by
$\mathfrak{g}$, $\mathfrak{d}$, $\mu_*$. The covariance $C^{(k)}_{\exp}(\Lambda)$ exists; by that
proposition, on the free $\Lambda$-bonds
$C^{(k)}_{\exp}(\Lambda)=(C_{Q_1}^*\Delta^{(k)}_{\exp}C_{Q_1})^{-1}$, in the sense in which
\cite[\S2]{B-CMP109} writes $C^{(k)}=(C^*\Delta^{(k)}C)^{-1}$. The form
$\Delta^{(k)}_{\exp}$ is positive definite on the space of fields $B$ with $B=0$ off $\Lambda$ and $Q_1B=0$.
With
$B_0$, $\delta_0$ the constants of \cite[(3.187)]{B-CMP99b} and $K_1^X$ the constant of the
exponential-gauge covariance proposition,
\begin{equation}\label{5.2:eq-form-positivity-1}
\begin{gathered}
|C^{(k)}_{\exp}(\Lambda;y,y')|\le B_0^Xe^{-\delta_0|y-y'|},\qquad B_0^X:=2B_0+c_dL^{d+3},\\
C_{Q_1}^*\Delta^{(k)}_{\exp}C_{Q_1}\ge (B_0^XK_1^X)^{-1}=:\gamma_0^X(d,L).
\end{gathered}
\end{equation}
Moreover $C^{(k)}_{\exp}(\Lambda)$ has the walk expansion and the analyticity of that proposition's
walk-expansion clause.

\smallskip
(b) \emph{The constraint of \cite{B-CMP109}.} Put
$\langle B,\Delta^{(k)}B\rangle:=\mathsf{G}(B)+\langle B,\Delta_kB\rangle$, so that
$\Delta^{(k)}=\Delta^{(k)}_{\exp}$; the notation $\Delta^{(k)}$ is that of \cite{B-CMP109}, the notation
$\Delta^{(k)}_{\exp}$ that of part (a). Let $C_{\flat}$ be the solution
operator of $Q^{\flat}B=0$ in the variables $B(b_0(c))$; this is the operator $C$ of \cite[\S2]{B-CMP109}
and \cite{B-CMP119}. Let $\mathcal{A}:=C_{\flat}^*\Delta^{(k)}C_{\flat}$ act on the free $\Lambda$-bonds, and
let $\pi_f$ denote restriction to free bonds. Define, with $C_d$ a constant depending only on $d$,
\begin{equation}\label{5.2:eq-form-positivity-2}
\tilde\Phi:=1+\pi_f\bar DE\mathfrak{a},\qquad c_{\Phi}:=\|\tilde\Phi\|_{2\to2}\le 1+C_dL^{(d+1)/2}.
\end{equation}
Let $\tilde C^{\Phi}$ be the free-bond covariance of part (a) with $\tilde D^{(2)}$ replaced by
$\tilde D^{(2),\Phi}$. Then $\mathcal{A}$ is real symmetric, and
\begin{equation}\label{5.2:eq-form-positivity-3}
\begin{gathered}
\mathcal{A}\ge\gamma_0:=\gamma_0^X/c_{\Phi}^2,\qquad
\mathcal{A}^{-1}=\tilde\Phi\,\tilde C^{\Phi}\,\tilde\Phi^*,\\
|\mathcal{A}^{-1}(b,b')|\le B_0^{\flat}e^{-\delta_0|b-b'|},\qquad B_0^{\flat}:=C(d,L)B_0^X.
\end{gathered}
\end{equation}
The right member of the middle identity carries the walk expansion and the analytic extension of part (a).

All constants depend only on $d$ and $L$. They are uniform in $k$, $U$, $\Lambda$, $\{\Omega_j\}$, the
labels and the volume.
\end{proposition}

\begin{proof}
\emph{Part (a).} The proof of the exponential-gauge covariance proposition constructs a map $\mu(\cdot)$,
the map of \cite[(3.169)]{B-CMP99b}. That construction is read in the averaging convention in which every
block variable is the average \cite[(0.11)]{B-CMP109}. In this convention the construction is
Lemma~\ref{5.2:lem-difference-inverse}, and $\mu(\cdot)=\mu_*$. The remaining parts of that proof, and its
estimate lemma, use $\mu(\cdot)$ at the following places and nowhere else.

First, they use the Gaussian identity \cite[(3.168)]{B-CMP99b}, which in the present notation reads
\begin{equation}\label{5.2:eq-form-positivity-4}
\int d(\mu\upharpoonright\Lambda)\;\delta(Q_1'\mu)\,e^{-\frac12\|\mathfrak{g}(B+\bar D\mu)\|^2}F(\mu)
=\int_{\mathcal{N}}d\mu'\,e^{-\frac12\|\mathfrak{d}\mu'\|^2}F(\mu_*(B)+\mu').
\end{equation}
Here $\upharpoonright$ denotes restriction, and $d(\mu\upharpoonright\Lambda)$ is Lebesgue measure on the
values of $\mu$ at the sites of $\Lambda$.
This follows from the last clause of Lemma~\ref{5.2:lem-difference-inverse}: the factor $\delta(Q_1'\mu)$
restricts the integration to $\mu\in\mathcal{N}$, with the measure on $\mathcal{N}$ induced by it, and on
$\mathcal{N}$ that clause gives $\mathfrak{g}(B+\bar D\mu)=\mathfrak{d}(\mu-\mu_*(B))$. Since
$\mu_*(B)\in\mathcal{N}$ by clause
(m3) of \eqref{5.2:eq-difference-inverse-a} in Lemma~\ref{5.2:lem-difference-inverse},
the translation $\mu'=\mu-\mu_*(B)$ maps $\mathcal{N}$ onto itself and preserves its measure.

Second, they use linearity and block-locality of $\mu(\cdot)$, which is clause (m1) of
\eqref{5.2:eq-difference-inverse-a}. They use
$\mu(\bar Dv)=-v$ for $v\in\mathcal{N}$, which is clause (m4) of \eqref{5.2:eq-difference-inverse-a}. They
use the commutation relation of
\cite[(3.114)]{B-CMP99b}, which here reads $Q_1\bar D\mu=\bar D_1Q_1'\mu$. This follows from
\eqref{5.2:eq-covariance-identities-a}, namely $Q^{\flat}\bar D\mu=\bar D_1(\mu\upharpoonright\Lambda')$ and
$\mathfrak{a}(\bar D\mu)=Q_1'\mu-\mu\upharpoonright\Lambda'$.

Third, they use $Q_1\bar D\mu=0$ for $\mu\in\mathcal{N}$. This is the same relation, since $Q_1'\mu=0$ on
$\mathcal{N}$.

Fourth, they require that the values of $\mu(\cdot)$ there used lie in $\mathcal{N}$; this is clause (m3) of
\eqref{5.2:eq-difference-inverse-a}.

Fifth, they use sup bounds on $\mu(\cdot)$ and on the operators built from it, $\bar D\mu(\cdot)$ among them.
These are supplied by clause (m2) of \eqref{5.2:eq-difference-inverse-a}, $|\mu_*(B)|\le 4dL\sup|B|$. Only
the numerical constants change.

Sixth, the walk-expansion clause uses the analyticity of $\mu(\cdot)$ in $V$. Since $\mu_*=-\mathfrak{d}^{-1}\mathfrak{g}$
is built from the block average, this follows from the analyticity of the average $M$ entering
\cite[(0.11)]{B-CMP109}.

Thus the proof uses only the properties (m1)--(m4) of $\mu(\cdot)$ in
\eqref{5.2:eq-difference-inverse-a}, the identity \eqref{5.2:eq-form-positivity-4}, the commutation
relation and analyticity in $V$; the map $\mu_*$ has all of them.
The estimates of the exponential-gauge covariance proposition and of its estimate lemma are therefore taken
here by statement, read in the averaging convention \cite[(0.11)]{B-CMP109}. The hypotheses of that
proposition are $(\tilde{\mathrm{H}}1)$, $(\tilde{\mathrm{H}}2)$, the hypothesis on $D$ and its smallness
condition. The first
three hold by \eqref{5.2:eq-form-positivity-a}, and the smallness condition is imposed as
$M_1\alpha_0\le a_X(d,L)$, with $a_X$ its threshold $a_E$ at $C_D$ replaced by $C_D^{\Phi}$. This gives
\eqref{5.2:eq-form-positivity-1} and the remaining assertions of (a).

\emph{Part (b): the measure.} Let $\varphi$ be a source field on the $\Lambda$-bonds. Substitute
$B=\Phi B_1$ in
\[
\int dB\;\delta(Q^{\flat}B)\,e^{-\frac12\langle B,\Delta^{(k)}B\rangle+\langle B,\varphi\rangle}.
\]
By Lemma~\ref{5.2:lem-unipotent-map}(a), $\Phi^{\pm1}$ preserve the property of being supported on
$\Lambda$-bonds. Since $\det\Phi=1$, the substitution is a linear bijection of the integration space with
$dB=dB_1$. By Lemma~\ref{5.2:lem-unipotent-map}(b), $Q^{\flat}\Phi=Q_1$, so
$\delta(Q^{\flat}\Phi B_1)=\delta(Q_1B_1)$. By the same lemma $\mathfrak{g}\Phi=\mathfrak{g}$, so
$\mathsf{G}(\Phi B_1)=\mathsf{G}(B_1)$.

The background $U$ is a scaffold minimiser as in Lemma~\ref{5.2:lem-almost-invariance} by
$(\tilde{\mathrm{H}}2)$, and $\tilde D^{(2)}$ satisfies the hypothesis on $D$. Hence
Lemma~\ref{5.2:lem-unipotent-map}(c)
gives
\[
\langle\Phi B_1,\Delta_k\Phi B_1\rangle
=\langle B_1,\Delta^{\circ}_kB_1\rangle-2\langle\tilde D^{(2),\Phi}(B_1),\mathfrak{j}\rangle.
\]
Multiplied by $-\frac12$, the last term becomes
$\langle\tilde D^{(2),\Phi}(B_1),\mathfrak{j}\rangle=\langle H_1\tilde D^{(2),\Phi}(B_1),J\rangle$, because
$\mathfrak{j}=H_1^*J$. Finally $\langle\Phi B_1,\varphi\rangle=\langle B_1,\Phi^*\varphi\rangle$. The
integral therefore equals
\begin{multline}\label{5.2:eq-form-positivity-5}
\int dB_1\;\delta(Q_1B_1)\exp\Bigl[-\tfrac12\mathsf{G}(B_1)-\tfrac12\langle B_1,\Delta^{\circ}_kB_1\rangle\\
+\langle H_1\tilde D^{(2),\Phi}(B_1),J\rangle+\langle B_1,\Phi^*\varphi\rangle\Bigr],
\end{multline}
with the sign of the term in $J$ as in \cite[(3.155), (3.161)]{B-CMP99b}.

This is the integral of part (a) with $\tilde D^{(2)}$ replaced by $\tilde D^{(2),\Phi}$ and source
$\Phi^*\varphi$. By Lemma~\ref{5.2:lem-unipotent-map}(c), $\tilde D^{(2),\Phi}$ satisfies the hypothesis on
$D$ with constant
$C_D^{\Phi}$. The smallness condition $M_1\alpha_0\le a_X(d,L)$ is the threshold of part (a) at
$C_D^{\Phi}$, so part (a) applies to $\tilde D^{(2),\Phi}$. In the exponential-gauge covariance proposition
the constant $C_D$ enters only through the threshold $a_E$ and an auxiliary constant. It does not enter
$\gamma_0^X$, $B_0^X$ or $\delta_0$, so these keep their values.

Denote by $\Delta^{(k),\Phi}_{\exp}$ the form $\Delta^{(k)}_{\exp}$ built with $\tilde D^{(2),\Phi}$. By
Lemma~\ref{5.2:lem-unipotent-map}(b),(c) and the two equalities above,
$\langle\Phi B_1,\Delta^{(k)}\Phi B_1\rangle=\langle B_1,\Delta^{(k),\Phi}_{\exp}B_1\rangle$.

\emph{Part (b): positivity.} Let $\tilde B$ be a field on the free $\Lambda$-bonds. Put
\[
B:=C_{\flat}\tilde B,\qquad B_1:=\Phi^{-1}B,\qquad \tilde B_1:=\pi_fB_1.
\]
Since $Q^{\flat}\Phi=Q_1$ by Lemma~\ref{5.2:lem-unipotent-map}(b), $Q_1\Phi^{-1}=Q^{\flat}$. Hence
$Q_1B_1=Q^{\flat}B=0$, that is, $B_1\in\{Q_1=0\}$.

By \eqref{5.2:eq-covariance-identities-a}, $\mathfrak{g}$ and $\mathfrak{a}$ read only bonds inside the
blocks $B(y)$, and none of these is a bond $b_0(c)$. Hence $\mathfrak{a}\pi_f=\mathfrak{a}$. As
$\pi_fB=\tilde B$, we get $\mathfrak{a}B=\mathfrak{a}\tilde B$. Lemma~\ref{5.2:lem-unipotent-map}(a) gives
$\Phi^{-1}=1-\bar DE\mathfrak{a}$, so
\[
\tilde B_1=\tilde B-\pi_f\bar DE\mathfrak{a}\tilde B.
\]
From $\mathfrak{a}\Phi=\mathfrak{a}$ we get $\mathfrak{a}\Phi^{-1}=\mathfrak{a}$. Therefore
\[
\mathfrak{a}\tilde B_1=\mathfrak{a}B_1=\mathfrak{a}B=\mathfrak{a}\tilde B,
\]
and consequently
\[
\tilde\Phi\tilde B_1=\tilde B-\pi_f\bar DE\mathfrak{a}\tilde B+\pi_f\bar DE\mathfrak{a}\tilde B=\tilde B.
\]

The map $\tilde\Phi$ is unipotent. Indeed, by $\mathfrak{a}\pi_f=\mathfrak{a}$ and
$\mathfrak{a}\bar DE=0$ (Lemma~\ref{5.2:lem-unipotent-map}(a)),
\[
(\pi_f\bar DE\mathfrak{a})^2=\pi_f\bar DE\,(\mathfrak{a}\bar DE)\,\mathfrak{a}=0.
\]

The operators $Q_1=Q^{\flat}+\bar D_1\mathfrak{a}$ and $Q^{\flat}$ have the same coefficient of
$B(b_0(c))$, because $\mathfrak{a}$ reads no bond $b_0(c)$. So $C_{Q_1}$ exists exactly as $C_{\flat}$ does.
A field in $\{Q_1=0\}$ is determined by its free values, which gives
$B_1=C_{Q_1}\pi_fB_1=C_{Q_1}\tilde B_1$.

Combining these facts with part (a) at $\tilde D^{(2),\Phi}$ and with \eqref{5.2:eq-form-positivity-2},
\begin{equation}\label{5.2:eq-form-positivity-6}
\begin{aligned}
\langle\tilde B,\mathcal{A}\tilde B\rangle
&=\langle\Phi B_1,\Delta^{(k)}\Phi B_1\rangle
=\langle C_{Q_1}\tilde B_1,\Delta^{(k),\Phi}_{\exp}C_{Q_1}\tilde B_1\rangle\\
&\ge\gamma_0^X\|\tilde B_1\|^2\ge\gamma_0^Xc_{\Phi}^{-2}\|\tilde B\|^2.
\end{aligned}
\end{equation}
The last inequality holds because $\|\tilde B\|=\|\tilde\Phi\tilde B_1\|\le c_{\Phi}\|\tilde B_1\|$. This
proves $\mathcal{A}\ge\gamma_0$.

\emph{Part (b): symmetry and the inverse.} $\mathcal{A}=C_{\flat}^*\Delta^{(k)}C_{\flat}$ is real symmetric,
since $\Delta^{(k)}$ is the symmetric operator of the real quadratic form $\langle B,\Delta^{(k)}B\rangle$.
Put $\mathcal{A}^{\Phi}:=C_{Q_1}^*\Delta^{(k),\Phi}_{\exp}C_{Q_1}$ on the free bonds. By part (a) its inverse
is $\tilde C^{\Phi}$. Since $\tilde B_1=\tilde\Phi^{-1}\tilde B$, the first line of
\eqref{5.2:eq-form-positivity-6} is an identity of quadratic forms between symmetric operators. Hence
\[
\mathcal{A}=(\tilde\Phi^{-1})^*\mathcal{A}^{\Phi}\tilde\Phi^{-1},
\qquad
\mathcal{A}^{-1}=\tilde\Phi\,\tilde C^{\Phi}\,\tilde\Phi^*.
\]
In this representation the right member carries the walk expansion and the analytic extension of part (a),
applied to $\tilde C^{\Phi}$.

\emph{Part (b): the kernel.} By Lemma~\ref{5.2:lem-unipotent-map}(a), $\tilde\Phi-1$ has range at most two
blocks. By \eqref{5.2:eq-covariance-identities-a}, $|E\mathfrak{a}B|\le 2dL\sup|B|$; the difference
$\bar D$ at most doubles this, so the entries of $\tilde\Phi-1$ are bounded by $4dL$. By the Schur test
these two facts give
\[
c_{\Phi}\le 1+\|\tilde\Phi-1\|_{2\to2}\le 1+c(d,L),
\]
with a number $c(d,L)$ depending only on $d$ and $L$.

For the kernel, write
\[
\mathcal{A}^{-1}(b,b')=\sum_{b_1,b_2}\tilde\Phi(b,b_1)\,\tilde C^{\Phi}(b_1,b_2)\,\tilde\Phi^*(b_2,b').
\]
Only bonds $b_1$ within two blocks of $b$, and $b_2$ within two blocks of $b'$, contribute. Their number,
and the excess of $|b-b'|$ over $|b_1-b_2|$, are bounded by numbers depending only on $d$ and $L$. Each
entry of $\tilde\Phi$ is bounded by $1+4dL$. Part (a) bounds $|\tilde C^{\Phi}(b_1,b_2)|$ by
$B_0^Xe^{-\delta_0|b_1-b_2|}$. This gives the kernel bound in \eqref{5.2:eq-form-positivity-3}, where
$C(d,L)$ collects the number of terms, the square of the entry bound and the factor $e^{\delta_0\ell}$, with
$\ell$ the bound on the distance shift.
\end{proof}

In the construction of \cite{B-CMP109} the map $\tilde D^{(2)}$ is $hC^{(2)}$. Here $C^{(2)}$ is the
second-order term in the expansion of the block average, and $h$ is the operator with $LQ^{\flat}h=1$; see
\cite[(1.5)]{B-CMP109} and \cite[(148)]{B-CMP98}.

\begin{corollary}[Kernel decay of the fluctuation form]\label{5.2:cor-kernel-decay}
Assume the following.
\begin{itemize}
\item[(i)] Let $k\ge 0$; let $(\Lambda,\{\Omega_j\}_{j\le k},U)$ satisfy the hypotheses
$(\tilde{\mathrm{H}}1)$ and $(\tilde{\mathrm{H}}2)$ of
Proposition~\ref{5.2:prop-form-positivity}, let $\tilde D^{(2)}$ satisfy its hypothesis
$(\mathrm{HD})$, and let $M_1\alpha_0\le a_X(d,L)$. Let $\mathcal{A}=C_{\flat}^*\Delta^{(k)}C_{\flat}$
be the form on the free $\Lambda$-bonds defined in its part (b), where
$\langle B,\Delta^{(k)}B\rangle=\mathsf{G}(B)+\langle B,\Delta_kB\rangle$.
\item[(ii)] The kernel bound for $(QG_1Q^*)^{-1}$, namely clause (iii) of the theorem on
propagators in a background field, which is the estimate of \cite{B-CMP99b} for the operator
$G_1$ of \cite[(3.138)]{B-CMP99b}, holds. Here, in the notation of
\cite{B-CMP99b}, $Q$ is the averaging operator and $G_1$ the operator given by
\cite[(3.138)]{B-CMP99b}; $c,c'$ are bonds of the lattice on which $Q$ takes its values,
$\ell_c,\ell_{c'}\le 1$ the
local units attached to them, $d(c,c')$ the distance of \cite{B-CMP99b} measured in these local
units, and $\delta_1>0$, $a_2(d,L)$ constants fixed there; $L^2M$ and $R$ are the block size and
the radius of \cite{B-CMP99b}, the letters $M$, $R$, $\delta_1$ and the distance $d(\cdot,\cdot)$
being used, in this corollary and its proof only, in the sense of \cite{B-CMP99b}.
The bound states: for $U$ satisfying \cite[(3.35)--(3.36)]{B-CMP99b} with
$L^2M\alpha_0\le a_2(d,L)$ and $R\ge 13L$,
\begin{equation}\label{5.2:eq-kernel-decay-1}
\|(QG_1Q^*)^{-1}(c,c')\|\le O(1)\,\ell_c^{-2}\,\ell_{c'}^{-d}\,e^{-\delta_1 d(c,c')} ,
\end{equation}
with $O(1)$ understood in the sense of \cite{B-CMP99b}.
\item[(iii)] Clause (c) of the lemma on the unipotent map between the two constraints
(Lemma~\ref{5.2:lem-unipotent-map}) holds; in particular $\tilde D^{(2),\Phi}$ satisfies
$(\mathrm{HD})$ with the constant $C_D^{\Phi}$ given there.
\item[(iv)] The covariance identity (G5) of \eqref{5.2:eq-covariance-identities-a} holds; in
particular the form $\mathsf{G}$ is block-diagonal.
\end{itemize}
Let $\mathsf{d}(b,b')$ denote the distance between the bonds $b$ and $b'$. Then there are constants
$B_{\mathcal{A}}$ and $\delta_{\mathcal{A}}>0$, depending only on $d$ and $L$, such that for all
free $\Lambda$-bonds $b,b'$
\begin{equation}\label{5.2:eq-kernel-decay-2}
\|\mathcal{A}(b,b')\|\le B_{\mathcal{A}}\,e^{-\delta_{\mathcal{A}}\mathsf{d}(b,b')} .
\end{equation}
\end{corollary}

\begin{proof}
Writing $\mathsf{G}$ also for the symmetric operator of the quadratic form $\mathsf{G}(B)$, we have
by definition $\mathcal{A}=C_{\flat}^*\mathsf{G}C_{\flat}+C_{\flat}^*\Delta_kC_{\flat}$, and we
bound the kernels of the operators that enter.

The operator $\Delta_k$ is $(QG_1Q^*)^{-1}$ minus $a$ times the identity, where $a$ is the
constant multiplying the averaging term of Ba{\l}aban's renormalization transformation, minus a
local term, namely the symmetric operator of the block-local quadratic form coming from
$\tilde D^{(2)}$.
For the first term we apply hypothesis (ii). Its three conditions are met in the present setting.
The configuration $U$ satisfies \cite[(3.35)--(3.36)]{B-CMP99b} by $(\tilde{\mathrm{H}}2)$. The
condition $L^2M\alpha_0\le a_2(d,L)$ of \cite{B-CMP99b} is met by the hypothesis
$M_1\alpha_0\le a_X(d,L)$ of (i). The radius
$R$ of \cite{B-CMP99b} is here $R_{\ast}$, the radius multiplying $M_1$ in the distance condition
of $(\tilde{\mathrm{H}}1)$, and the condition $R\ge 13L$ holds because
$R_{\ast}=L^3\ge 13L$, $L$ being large. The arguments $c,c'$ relevant to $\mathcal{A}$ lie in the
region $\Lambda\subset\Omega_k^{(k)}$ by $(\tilde{\mathrm{H}}1)$, where the local units equal $1$,
so we read \eqref{5.2:eq-kernel-decay-1} at $\ell_c=\ell_{c'}=1$. Since all local units are at
most $1$,
\begin{equation*}
d(c,c')\ge|c-c'| .
\end{equation*}
Hence the kernel of $(QG_1Q^*)^{-1}$ is bounded by $O(1)\,e^{-\delta_1|c-c'|}$.

The local term, the form coming from $\tilde D^{(2)}$, is block-local. By hypothesis (iii),
$\tilde D^{(2),\Phi}$ satisfies $(\mathrm{HD})$ with constant $C_D^{\Phi}$; since
$C_D\le C_D^{\Phi}$ by the expression for $C_D^{\Phi}$ given in
Lemma~\ref{5.2:lem-unipotent-map}(c), the map $\tilde D^{(2)}$ satisfies $(\mathrm{HD})$ with
constant $C_D^{\Phi}$ as well. The entries of the symmetric operator of the form coming from
$\tilde D^{(2)}$ are bounded by $2C_D^{\Phi}C_HM_1\alpha_0$, where $C_H$ denotes a constant
depending only on $d$ and $L$.

By hypothesis (iv), that is, by (G5), the form $\mathsf{G}$ is block-diagonal.

Finally, let $e$ be a bond of $T^{(k+1)}$ with endpoints $e_-,e_+\in\Lambda'$, let $b_0(e)$ be the
constrained bond of Proposition~\ref{5.2:prop-form-positivity}(b), and write $B(e)$ for the set
of bonds joining the blocks $B(e_-)$ and $B(e_+)$. The operator $C_{\flat}$ is the solution
operator of $Q^{\flat}B=0$ in the bonds $b_0(e)$: at free bonds $C_{\flat}\tilde B=\tilde B$, and
the component $(C_{\flat}\tilde B)(b_0(e))$ depends only on the restriction of $\tilde B$ to the
bonds in $B(e_-)\cup B(e_+)\cup B(e)$, as in \cite{B-CMP119}, with coefficients bounded by a
constant $C(d,L)$. The same holds for $C_{\flat}^*$.

It remains to assemble these bounds. The operators $C_{\flat}$ and $C_{\flat}^*$ depend only on
the bonds just described and have entries at most $C(d,L)$. The operator $a$ times the identity is
diagonal, $\mathsf{G}$ is block-diagonal, and the local term is block-local with the bound above.
The remaining operator
$(QG_1Q^*)^{-1}$ has kernel at most $O(1)\,e^{-\delta_1|c-c'|}$. The kernel of $\mathcal{A}$
between $b$ and $b'$ is therefore a sum, over a number of bonds bounded in terms of $d$ and $L$,
of products of entries of these operators. Each product contains at most one entry of
$(QG_1Q^*)^{-1}$, and its other factors connect bonds at distance at most $C(d,L)$. Hence a
product without such an entry vanishes unless $\mathsf{d}(b,b')\le C(d,L)$, and a product
containing the entry $(QG_1Q^*)^{-1}(c,c')$ has $|c-c'|\ge\mathsf{d}(b,b')-C(d,L)$. This gives
\eqref{5.2:eq-kernel-decay-2} with $\delta_{\mathcal{A}}=\delta_1$ and with $B_{\mathcal{A}}$
depending only on $d$ and $L$.
\end{proof}

\begin{lemma}[Square root of the conditioned covariance]\label{5.2:lem-covariance-root}
Assume the hypotheses of Proposition~\ref{5.2:prop-form-positivity}, with region $\Lambda$, and let
$\mathcal{A}$ be the constrained fluctuation quadratic form on the free $\Lambda$-bonds $\mathfrak{X}$, with
lower bound $\gamma_0$ and decay constants $(B_{\mathcal{A}},\delta_{\mathcal{A}})$ as in
\eqref{5.2:eq-true-pair-positivity-a}. Let $\rho(b,b')$ be the bond distance in the current unit
lattice, and let $\mu_0$ be the Combes--Thomas rate determined by $(\gamma_0,B_{\mathcal{A}},\delta_{\mathcal{A}})$,
both as fixed in \eqref{5.2:eq-rim-constants-a}. For $\bar\Lambda\subset\mathfrak{X}$ let
$\mathcal{A}_{\bar\Lambda\bar\Lambda}$ be the restriction of $\mathcal{A}$ to functions supported on
$\bar\Lambda$, and let $C_\Lambda:=(\mathcal{A}_{\bar\Lambda\bar\Lambda})^{-1}$ be the conditioned
covariance of \cite[(2.7)]{B-CMP116} at $\sigma=\mathbb{1}$. Then for every
$\bar\Lambda\subset\mathfrak{X}$
\begin{equation}\label{5.2:eq-covariance-root-1}
C_\Lambda^{1/2}=\frac{1}{\pi}\int_0^\infty dx\,x^{-1/2}\,\bigl(x+\mathcal{A}_{\bar\Lambda\bar\Lambda}\bigr)^{-1},
\end{equation}
and, for all bonds $b,b'\in\bar\Lambda$,
\begin{equation}\label{5.2:eq-covariance-root-2}
\bigl\|C_\Lambda^{1/2}(b,b')\bigr\|\le 2\gamma_0^{-1/2}\,e^{-\mu_0\rho(b,b')}.
\end{equation}
This is the case $\sigma=\mathbb{1}$ of \cite[(2.7)]{B-CMP116}, read on $\Lambda$.
\end{lemma}

\begin{proof}
By \eqref{5.2:eq-true-pair-positivity-a}, $\mathcal{A}$ is real symmetric and $\mathcal{A}\ge\gamma_0>0$.
For $v$ supported on $\bar\Lambda$ one has
$\langle v,\mathcal{A}_{\bar\Lambda\bar\Lambda}v\rangle=\langle v,\mathcal{A}v\rangle$, so that
$\mathcal{A}_{\bar\Lambda\bar\Lambda}$ is real symmetric and satisfies
$\mathcal{A}_{\bar\Lambda\bar\Lambda}\ge\gamma_0$. It is therefore invertible, and $C_\Lambda$ is
defined. By the spectral theorem, as in \cite[(2.7)]{B-CMP116}, it suffices to verify
\eqref{5.2:eq-covariance-root-1} on each
eigenvalue $a\ge\gamma_0$ of $\mathcal{A}_{\bar\Lambda\bar\Lambda}$. The substitution $x=at^2$ gives
\begin{equation}\label{5.2:eq-covariance-root-3}
\int_0^\infty\frac{x^{-1/2}\,dx}{x+a}=2a^{-1/2}\int_0^\infty\frac{dt}{1+t^2}=\pi a^{-1/2}.
\end{equation}
Thus the right member of \eqref{5.2:eq-covariance-root-1} equals
$(\mathcal{A}_{\bar\Lambda\bar\Lambda})^{-1/2}=C_\Lambda^{1/2}$. The integral converges in operator
norm, since
\begin{equation*}
\bigl\|(x+\mathcal{A}_{\bar\Lambda\bar\Lambda})^{-1}\bigr\|\le(x+\gamma_0)^{-1}.
\end{equation*}

For the kernel bound, fix $x\ge0$ and apply the Combes--Thomas lemma of this chapter
(cf.\ \cite{CT73}) to
$x+\mathcal{A}_{\bar\Lambda\bar\Lambda}$. Its hypotheses hold. The operator is bounded below by
$\gamma_0+x$. Its off-diagonal entries are entries of $\mathcal{A}$, so they obey the decay bound
with constants $(B_{\mathcal{A}},\delta_{\mathcal{A}})$ of \eqref{5.2:eq-true-pair-positivity-a}.

Fix a bond $b_0\in\bar\Lambda$ and let $D$ be multiplication by $e^{\mu_0\rho(b_0,\cdot)}$.
Conjugating by $D$ leaves the multiple $x$ of the identity unchanged:
\begin{equation*}
D(x+\mathcal{A}_{\bar\Lambda\bar\Lambda})D^{-1}-(x+\mathcal{A}_{\bar\Lambda\bar\Lambda})
=D\mathcal{A}_{\bar\Lambda\bar\Lambda}D^{-1}-\mathcal{A}_{\bar\Lambda\bar\Lambda}.
\end{equation*}
Hence the conjugation error of the Combes--Thomas lemma at the rate $\mu_0$,
\begin{equation*}
\varepsilon(\mu_0):=\sup_{b_0}\bigl\|D\mathcal{A}_{\bar\Lambda\bar\Lambda}D^{-1}
-\mathcal{A}_{\bar\Lambda\bar\Lambda}\bigr\|,
\end{equation*}
does not depend on $x$. By the
choice of $\mu_0$ in \eqref{5.2:eq-rim-constants-a} it satisfies
\begin{equation*}
\varepsilon(\mu_0)\le\tfrac12\gamma_0\le\tfrac12(\gamma_0+x).
\end{equation*}
The Combes--Thomas lemma therefore gives, for all $x\ge0$,
\begin{equation}\label{5.2:eq-covariance-root-4}
\bigl\|(x+\mathcal{A}_{\bar\Lambda\bar\Lambda})^{-1}(b,b')\bigr\|\le 2(\gamma_0+x)^{-1}e^{-\mu_0\rho(b,b')}.
\end{equation}

Take the $(b,b')$ entry of \eqref{5.2:eq-covariance-root-1} and insert \eqref{5.2:eq-covariance-root-4}
under the integral. Then apply \eqref{5.2:eq-covariance-root-3} with $a=\gamma_0$:
\begin{equation*}
\begin{split}
\bigl\|C_\Lambda^{1/2}(b,b')\bigr\|
&\le\frac{1}{\pi}\int_0^\infty dx\,x^{-1/2}\bigl\|(x+\mathcal{A}_{\bar\Lambda\bar\Lambda})^{-1}(b,b')\bigr\|\\
&\le\frac{2}{\pi}\,e^{-\mu_0\rho(b,b')}\int_0^\infty\frac{x^{-1/2}\,dx}{\gamma_0+x}
=2\gamma_0^{-1/2}e^{-\mu_0\rho(b,b')}.
\end{split}
\end{equation*}
This is \eqref{5.2:eq-covariance-root-2}.

Lemma~\ref{5.2:lem-covariance-root} concerns only the case $\sigma=\mathbb{1}$. The family of
covariances in the decoupling parameters $\sigma$, its difference bound and the reference point
$\sigma\equiv0$ are taken as stated from the random-walk expansion of this section.
\end{proof}

\begin{corollary}[Positivity and decay at every true pair]\label{5.2:cor-true-pair-positivity}
Let $k\ge 0$, and let $(\Omega_j)_{j\le k+1}$ be the sequence of regions of a term at level $k+1$.
A \emph{true pair} is a pair $(U,J(U))$ with the following properties. First, $U$ is a scaffold
minimiser, in the sense of Lemma~\ref{5.2:lem-almost-invariance}, for a determining set equal to
or coarser than the determining set of $(\Omega_j)_{j\le k+1}$. Second, the pair is taken on the
support of the characteristic functions of the term. In part (b) of
Proposition~\ref{5.2:prop-form-positivity}
take the following data: the domain is $\Omega^{(k)}_{k+1}$; the regions are $(\Omega_j)_{j\le k}$;
the background is $U$; the block size is $M_1$; the range is $R^{\mathrm{bg}}=L^3$, where
$R^{\mathrm{bg}}$ denotes the parameter $R$ of the construction of \cite{B-CMP99b}; and the
quadratic map is the map $\tilde D^{(2)}=hC^{(2)}$ of \cite{B-CMP109}, where $C^{(2)}$ is the
operator of \cite[(148)]{B-CMP98} and $h$ is the coefficient multiplying it in \cite{B-CMP109}
(here $h$ is not a history). At these data we fix the
following objects, all as in part (b) of Proposition~\ref{5.2:prop-form-positivity}:
\begin{itemize}
\item $\mathfrak{X}$, the set of free $\Omega^{(k)}_{k+1}$-bonds;
\item $C$ (denoted $C_\flat$ in part (b) of Proposition~\ref{5.2:prop-form-positivity}), the
solution operator of the constraint $Q^{\flat}B=0$ in the distinguished bonds $b_0(c)$;
\item the form $\langle B,\Delta^{(k)}B\rangle:=\mathsf{G}(B)+\langle B,\Delta_kB\rangle$.
\end{itemize}
Put
\begin{equation*}
\mathcal{A}:=C^{*}\Delta^{(k)}C\quad\text{on }\ell^2(\mathfrak{X};\mathfrak{su}(N)),
\end{equation*}
where $\mathfrak{su}(N)$ is the Lie algebra of $\mathrm{SU}(N)$. Assume:
\begin{enumerate}
\item[(i)] Proposition~\ref{5.2:prop-form-positivity} (positivity of the constrained fluctuation
form), part (b);
\item[(ii)] Corollary~\ref{5.2:cor-kernel-decay} (kernel decay of the fluctuation form);
\item[(iii)] Lemma~\ref{5.2:lem-almost-invariance} (almost-invariance of the fluctuation form);
\item[(iv)] part (a) of Proposition~\ref{5.2:prop-form-positivity}, the covariance estimate under
the linear constraint $Q_1$ on which part (b) rests, with its uniformity clause;
\item[(v)] the Gaussian identity on the kernel of the constraint, used in the proof of part (a) of
Proposition~\ref{5.2:prop-form-positivity};
\item[(vi)] Lemma~\ref{5.2:lem-unipotent-map} on the unipotent substitution $\Phi$, in particular
its part (c);
\item[(vii)] the identification of the small-field domain, together with the monotonicity of
regularity under passage to sub-cubes;
\item[(viii)] the decay bound on the inverse of the averaged background-field propagator, as
stated in \cite{B-CMP99b};
\item[(ix)] the statement on the block-averaging map.
\end{enumerate}
Let $\alpha_0(\gamma)$ be the regularity constant of $U$ supplied by
\cite[Theorem~1]{B-CMP102b} at the present data; it is of the size
\begin{equation*}
\alpha_0(\gamma)=O(1)\,B_3L^2(1+\beta_0)\sup_{i\le k}\varepsilon_i,\qquad
\varepsilon_i:=g_i\,p_0(g_i),
\end{equation*}
with $B_3$, $\beta_0$ the constants of its derivation from the data bound
\cite[(7)]{B-CMP102b} combined with the bound of \cite{B-CMP119} recalled in the proof, and with
$p_0(\cdot)$ the degree function of the coupling. Assume
the one condition
$M_1\alpha_0(\gamma)\le a_X(d,L)$, with $a_X(d,L)$ the constant of
Proposition~\ref{5.2:prop-form-positivity}. Then at every true pair $\mathcal{A}$ is real
symmetric and
\begin{equation}\label{5.2:eq-true-pair-positivity-a}
\mathcal{A}\ge\gamma_0:=\gamma_0^X/c_\Phi^2,\qquad
\|\mathcal{A}(b,b')\|\le B_{\mathcal{A}}\,e^{-\delta_{\mathcal{A}}\rho(b,b')},
\end{equation}
with $\gamma_0^X$, $c_\Phi$ as in Proposition~\ref{5.2:prop-form-positivity} and $\rho(b,b')$ as
in Corollary~\ref{5.2:cor-kernel-decay}. The numbers $\gamma_0$, $B_{\mathcal{A}}$,
$\delta_{\mathcal{A}}$ depend only on $d$ and $L$, and are uniform in $k$, in $\{\Omega_j\}$, in
the labels and in the volume. Hence, for every $\bar\Lambda\subset\mathfrak{X}$, the compression
$\mathcal{A}_{\bar\Lambda\bar\Lambda}$ of $\mathcal{A}$ to $\ell^2(\bar\Lambda;\mathfrak{su}(N))$
satisfies $\mathcal{A}_{\bar\Lambda\bar\Lambda}\ge\gamma_0$. Moreover $S\ge\gamma_0$ and
$S_{\mathrm{cc}}\ge\gamma_0$, where
$S$ is the Schur complement of the rim construction and $S_{\mathrm{cc}}$ is a compression of $S$.
\end{corollary}

\begin{proof}
Three statements of the rim construction are used below as stated, namely the
Schur-complement statement, the cube-structure statement and the conjugation bound, together with
the support lemma for the compressed domain. Of the hypotheses, (iii), (iv), (vi), (vii) and
(viii) are invoked below by name; (v) enters through part (a) of
Proposition~\ref{5.2:prop-form-positivity}, on which part (b) rests.

Both bounds in \eqref{5.2:eq-true-pair-positivity-a} follow from part (b) of
Proposition~\ref{5.2:prop-form-positivity} and from Corollary~\ref{5.2:cor-kernel-decay},
applied at the data fixed in the statement, which are data of the construction of
\cite{B-CMP99b}. We show that the
following hold at these data, with the domain $\Lambda$ of
Proposition~\ref{5.2:prop-form-positivity} read at $\Omega^{(k)}_{k+1}$: the hypotheses
$(\tilde{\mathrm{H}}1)$, $(\tilde{\mathrm{H}}2)$ and the hypothesis on $\tilde D^{(2)}$ of
Proposition~\ref{5.2:prop-form-positivity}, displayed in \eqref{5.2:eq-form-positivity-a}, and
the thresholds of that proposition.

\smallskip
The domain. By the identification of the small-field domain, the domain of the construction of
\cite{B-CMP99b} at level $k$, inside which the domain $\Lambda$ of $(\tilde{\mathrm{H}}1)$ must
lie, is the set $\Omega_k^{(k)}$ of \cite{B-CMP119}. For
$k=0$ the identification is read off \cite{B-CMP102a}.

\smallskip
Hypothesis $(\tilde{\mathrm{H}}1)$. The domain $\Omega^{(k)}_{k+1}$ is the region $\Omega_{k+1}$
read on the lattice $T^{(k)}$, and distances below are measured in units of $T^{(k)}$. The set
$\Omega_{k+1}$ is defined in \cite[(3.5)]{B-CMP119}
as the complement of the union of two sets: the enlargement $Q'^{\sim 2}_{k+1}$, and an
enlargement of $(B_{k+1}(P^1_{k+1}))^c$. Here $Q'_{k+1}$, $P^1_{k+1}$ and $B_{k+1}$ are as in
\cite{B-CMP119}. By \cite{B-CMP119}, $\Omega_{k+1}$ is a union of $LM\check{R}_{k+1}$-cubes of
the level-$k$ lattice. Equivalently, it is a union of the cubes of the corresponding partition of
the level-$(k+1)$ lattice into $M\check{R}_{k+1}$-cubes. In the notation of
Proposition~\ref{5.2:prop-form-positivity}, this says that $\Omega_{k+1}$ is the image under the
map $B^k$ of that proposition (a map on sets of the lattice, not the bond field $B$) of a union of
$M\check{R}_{k+1}$-cubes of $T^{(k+1)}$.

Each such cube is a union of big $M_1$-blocks. Indeed, $M$ and $M_2=L^2M_1$ are powers of $L$
with $M>M_2$, so $M\ge L^3M_1$, and the partitions are compatible \cite{B-CMP119}.

Further, $\Omega_{k+1}$ lies outside $Z_k'^{\sim 8}$, the enlargement by eight layers of the set
$Z_k'$ of \cite{B-CMP119}. The enlargement consists of four layers, then one more, then one more,
then two more, as enumerated in \cite[\S 3]{B-CMP119}. The large-field
region $Z_k$ of \cite{B-CMP119}, the complement of the region $\Lambda_k$ of that paper, satisfies
$Z_k=\Lambda_k^{c}\supseteq\Omega_k^{c}$. Hence, in units of $T^{(k)}$,
\begin{equation*}
\operatorname{dist}\bigl(\Omega_{k+1},(\Omega_k^{(k)})^c\bigr)
\ge 8LM\check{R}_{k+1}\ge 8LM\ge 8L^4M_1>L^3M_1=R^{\mathrm{bg}}M_1,
\end{equation*}
where $R^{\mathrm{bg}}:=L^3$.

The inequality $L^3\ge 13L$ amounts to the lower bound $L^2\ge 13$ on $L$, which is one of the
lower bounds on $L$ secured by $L\ge L_0$. So this
$R^{\mathrm{bg}}$ meets the condition $R\ge 13L$ of the decay bound (viii) on the inverse of the
averaged background-field propagator. The
requirement $R_E\le L^3M_1$ is a lower bound on $M_1$, imposed after $L$. The distance condition
is one-sided, and no upper bound enters.

Finally, the extended sequence $\{\tilde\Omega_j\}_{j\le k+1}$ of
Proposition~\ref{5.2:prop-form-positivity} is the term's own sequence $(\Omega_j)_{j\le k+1}$
\cite{B-CMP119}. This sequence satisfies \cite[(2.1)--(2.2)]{B-CMP96}, as part of Ba{\l}aban's
standing construction.

\smallskip
Hypothesis $(\tilde{\mathrm{H}}2)$. The background $U=U_{k+1}$, or any of the other scaffold
minimisers admitted in Lemma~\ref{5.2:lem-almost-invariance}, is the minimiser belonging to its
own sequence. Its data in the sense of \cite[(7)]{B-CMP102b} are small on the support of the
characteristic functions, by \cite[(3.12)]{B-CMP119}. Therefore \cite[Theorem~1]{B-CMP102b} gives
the regularity required in $(\tilde{\mathrm{H}}2)$ on the cubes of index $k+1$, which lie in
$\Omega_{k+1}$, with the constant $\alpha_0=\alpha_0(\gamma)$ of the statement. The value
of $\alpha_0$ comes from that data bound combined with the bound
$|V_{k+1}(\partial p')-1|<2L^2\varepsilon_k$ of \cite{B-CMP119}.

On the cubes of index $j$ of older levels the regularity follows from the monotonicity of
regularity under passage to sub-cubes. Suppose a gauge transformation $u$ brings $U$ on a cube to
the form $U^u=e^{i\eta A}$, with $\eta$ the lattice spacing and $A$ a Lie-algebra valued field
satisfying
\begin{equation*}
|A|<cM_1\alpha_0\,\ell^{-1},\quad |\nabla A|<cM_1\alpha_0\,\ell^{-2},\quad
|\partial^{*}\partial A|<cM_1\alpha_0\,\ell^{-3},\qquad \ell=L^{j+1}\eta,
\end{equation*}
with $c$ the constant of the regularity conditions \cite[(3.35)--(3.36)]{B-CMP99b}, and with
their block size $M$ read at $M_1$, as in the data fixed in the statement.
Then the same $u$ on any sub-cube gives the bounds of index $j$ with constant $\alpha_0/L$. The
bookkeeping of the cubes near the boundary of $\Omega_{k+1}$ is part of that monotonicity
statement.

For the aged backgrounds $U_{\mathfrak{B}'}$ of Lemma~\ref{5.2:lem-almost-invariance}, one applies
\cite[Theorem~1]{B-CMP102b} for the determining set $\mathfrak{B}'$ and then the monotonicity. The
scaffold-minimiser clause of $(\tilde{\mathrm{H}}2)$ is the definition of a true pair.

\smallskip
The hypothesis on $\tilde D^{(2)}$. In \cite{B-CMP109} the quadratic map is
$\tilde D^{(2)}=hC^{(2)}$, with $C^{(2)}$ the operator of \cite[(148)]{B-CMP98} and $h$ the
coefficient multiplying it in \cite{B-CMP109}. It has the block-locality and the bound required in
the hypothesis on $\tilde D^{(2)}$. Part (c)
of Lemma~\ref{5.2:lem-unipotent-map} on the unipotent substitution $\Phi$ then supplies the same
properties for the transformed map, with constant $C_D^{\Phi}$.

\smallskip
Thresholds. The conditions $M_1\ge M_E(d,L)$ and $L^3M_1\ge R_E(d,L)$ are lower bounds on $M_1$,
imposed after $L$. In \cite{B-CMP119}, $M_1$ is the size of the large cubes for which the
theorems of Ba{\l}aban's earlier papers are valid. Consider next the conditions
\begin{equation*}
\begin{gathered}
M_1\alpha_0=M_1\,O(1)\,B_3L^2(1+\beta_0)\sup_{i\le k}\varepsilon_i\le a_X(d,L),\\
M\le M(\epsilon_{\mathrm{data}})=R_1M_1(a_1/\epsilon_{\mathrm{data}}).
\end{gathered}
\end{equation*}
The second is the condition of \cite[Theorem~1]{B-CMP102b}, with $R_1$, $a_1$ its constants and
$\epsilon_{\mathrm{data}}$ its smallness parameter, denoted $\varepsilon_1$ in
\cite[Theorem~1]{B-CMP102b} and written $\epsilon_{\mathrm{data}}$ here to avoid collision with
$\varepsilon_i$. At the present data this parameter is read at
the smallness of the data \cite[(7)]{B-CMP102b} on the support, which enters through
$\sup_{i\le k}\varepsilon_i$, as in the value of $\alpha_0$ above. Both conditions are therefore
upper bounds on $\gamma$, imposed after $M_1$ and $M$. They can be met because
$\varepsilon_k=g_kp_0(g_k)\to 0$.

\smallskip
The two bounds. All hypotheses of part (b) of Proposition~\ref{5.2:prop-form-positivity} hold,
so $\mathcal{A}$ is real symmetric and $\mathcal{A}\ge\gamma_0^X/c_\Phi^2=\gamma_0$. The same
hypotheses are those of Corollary~\ref{5.2:cor-kernel-decay}. That corollary gives the kernel
bound in \eqref{5.2:eq-true-pair-positivity-a}, with $(d,L)$-constants $B_{\mathcal{A}}$,
$\delta_{\mathcal{A}}$. This proves \eqref{5.2:eq-true-pair-positivity-a}.

\smallskip
Compression. We use three objects of the rim construction: the compressed domain $\Lambda^{*}$;
the corresponding set $\Lambda^{*\prime}$ of $T^{(k+1)}$; and the set
$\bar\Lambda=(\Lambda^{*})^{*}$ of free bonds attached to $\Lambda^{*}$ by the convention of that
construction. Let $\tilde B$ be supported on $\bar\Lambda$. By the support lemma for the
compressed domain, $(C\tilde B)(b_0(c))=0$ unless the cube $c$ meets $\Lambda^{*\prime}$, and the
$y$-block $\mathsf{G}_y$ of the block-diagonal form $\mathsf{G}$ satisfies
$\mathsf{G}_y(\tilde B)=0$ for $y\notin\Lambda^{*\prime}$. Hence the form of \cite{B-CMP99b} on the
domain $\Lambda^{*}$ is $\mathcal{A}_{\bar\Lambda\bar\Lambda}$, as the rim construction defines it.

Moreover, $\Lambda^{*}$ satisfies its own hypothesis $(\tilde{\mathrm{H}}1)$. By the
cube-structure statement, it is a union of $LM\check{R}_{k+1}$-cubes contained in
$\Omega_{k+1}^{\sim -1}$, the set obtained from $\Omega_{k+1}$ by removing one layer of cubes
along its boundary; hence it lies deeper in $\Omega_k^{(k)}$ than $\Omega_{k+1}$. So part
(b) of Proposition~\ref{5.2:prop-form-positivity} applies on $\Lambda^{*}$. It gives, for
$C_\Lambda=(\mathcal{A}_{\bar\Lambda\bar\Lambda})^{-1}$, the covariance of
Lemma~\ref{5.2:lem-covariance-root}, the bound
\begin{equation*}
\|C_\Lambda(b,b')\|\le B_0^{\flat}\,e^{-\delta_0|b-b'|},
\end{equation*}
with $B_0^{\flat}$ the constant and $\delta_0$ the decay rate of the covariance bound in part (b)
of Proposition~\ref{5.2:prop-form-positivity}.
This bound agrees with the conjugation bound of the rim construction.

\smallskip
Uniformity. In the sequences continued beyond level $k+1$ after the compression, the later
$\Omega_j$ are built on $(\Lambda^{*})^{c}$. By the
cube-structure statement, they are unions of the same cubes at the same distances. By the
uniformity clause of the covariance estimate under the linear constraint $Q_1$, three facts hold
for its constants: $k$ enters them nowhere; $U$ enters only through $\alpha_0$; and $\Lambda$,
$\{\Omega_j\}$ do not enter at all. Hence $\gamma_0$, $B_{\mathcal{A}}$, $\delta_{\mathcal{A}}$
are uniform in $k$, $\{\Omega_j\}$, the labels and the volume.

\smallskip
Consequences. For $v$ supported on $\bar\Lambda\subset\mathfrak{X}$ one has
\begin{equation*}
\langle v,\mathcal{A}_{\bar\Lambda\bar\Lambda}v\rangle=\langle v,\mathcal{A}v\rangle
\ge\gamma_0\|v\|^2,
\end{equation*}
so $\mathcal{A}_{\bar\Lambda\bar\Lambda}\ge\gamma_0$. The Schur-complement statement of the rim
construction then gives $S\ge\gamma_0$. Finally, $S_{\mathrm{cc}}\ge\gamma_0$ because
$S_{\mathrm{cc}}$ is a compression of $S$.

\smallskip
Order of the choices. The numbers $\gamma_0^X$, $c_\Phi$, $C_D^{\Phi}$, $B_{\mathcal{A}}$,
$\delta_{\mathcal{A}}$, $a_X$, $M_E$, $R_E$ depend only on $(d,L)$. The number $\gamma_0$
decreases polynomially in $L$. It is read only in two places: by the $(d,L)$-numbers
$\Theta_1$, $\varpi$, $\mu_0$ (the rate in Lemma~\ref{5.2:lem-covariance-root}) and $K_0$ of the
rim construction, all chosen after it; and
by upper bounds on $\gamma$, such as $\gamma_0A_{\flat}^2p_0(\gamma)\ge 8$. The conditions on
$M_1$ are lower bounds imposed after $L$. Checking the hypotheses at the present data adds one
upper bound on
$\gamma$, namely $M_1\alpha_0\le a_X$, beside the condition $M\le M(\epsilon_{\mathrm{data}})$ of
\cite{B-CMP102b}. No inequality of the form $C(L)\le c$ with $c$ independent of $L$ occurs. No
upper bound is imposed on $L$ or on $M_1$, and the condition $M\le M(\epsilon_{\mathrm{data}})$ is
met by the choice of $\gamma$.
\end{proof}

\begin{remark}[A second route to positivity, and its scope]\label{5.2:rem-positivity-scope}
\textup{(a)} A positivity bound of the same kind as the lower bound in
\eqref{5.2:eq-true-pair-positivity-a}, with a different constant, can also be reached
without the kernel representation. This second route yields the positivity part only and
gives no information on the decay of the kernel of $\mathcal{A}$. It combines gauge
invariance on the constraint space (Corollary~\ref{5.2:cor-constrained-invariance}) with the
slice positivity corollary, which bounds $\Delta_k$ from below on the slice
$\{\mathfrak{g}=0,\ Q_1=0\}$. Here, as in the definition of the fluctuation forms,
$\Delta^{(k)}$ denotes the fluctuation form in which the gauge-fixing form is added to
$\Delta_k$, and $\Delta_k$ the form without it; $\pi_f$ denotes an orthogonal projection
fixed with the fluctuation forms, of which only the fact that its norm is at most one is
used below. Put
\begin{equation*}
\gamma_0^{\mathrm{sl}} := \bigl(B_0 K_1^X\bigr)^{-1},
\end{equation*}
where $B_0$ and $K_1^X$ are the two constants entering the bound of that corollary. At a
true pair in the sense of Corollary~\ref{5.2:cor-true-pair-positivity}, and with
$\tilde D^{(2)}=hC^{(2)}$, every $B$ with $Q^{\flat}B=0$ satisfies
\begin{equation}\label{5.2:eq-positivity-scope-1}
\langle B,\Delta^{(k)}B\rangle
\ \ge\ \min\Bigl\{\frac{\gamma_0^{\mathrm{sl}}}{2},\ \frac{1}{8d}\Bigr\}\,\|\pi_f B\|^2 .
\end{equation}

\textup{(b)} In the proof of the correlation theorem
(Theorem~\ref{5.1:thm-correlation-theorem}), the lower bound and the kernel decay of
\eqref{5.2:eq-true-pair-positivity-a} are used only at true pairs, in the sense of
Corollary~\ref{5.2:cor-true-pair-positivity}. The uses are the following.
\begin{enumerate}
\item The bounds of the transfer-operator lemmas on $T=-C_\Lambda\mathcal{A}_{\Lambda\mathfrak{r}}$,
where $C_\Lambda$ is the covariance of the Gaussian integral actually performed on the
region $\Lambda$, and $\mathfrak{r}$ is the set of rim bonds. The quantities that read
$\gamma_0$ there are the first member $(2B_{\mathcal{A}}/\gamma_0)\,C_\rho(\mu_0/2)$ of the
constant $K_0$ and the numbers $\delta_1$, $\mu_0$, all in the notation of those lemmas.
\item The first two steps and clause \textup{(a)} of the theorem on the creating step on the
support $\mathfrak{X}^{\sharp}$, where the Gaussian identities used are identities for
the Gaussian field of mean zero and covariance $C_\Lambda$, that is, for the actual Gaussian.
\item The lower bounds $S,\ S_{\mathrm{cc}}\ge\gamma_0$ on the Schur complement $S$ and on its
compression $S_{\mathrm{cc}}$ in the payment estimate per cube, at the actual large-field
factor $\mathbf{T}$. They are read at the backgrounds of large-field regions created at
earlier steps; that $S$ is there the Schur complement of the actual factor follows from the
second-order identity for the actual factor at such backgrounds, the identity into which the
second-order comparison is inserted in the proof of
Corollary~\ref{5.2:cor-constrained-invariance}, and the lower bound is the second clause of
the bound on the actual factor $\mathbf{T}$ at such backgrounds.
\item The numerical gains read from $\gamma_0$ in the strip lemma and in the per-zone-cell
volume constant $\mathsf{K}_{\mathrm{cell}}(k,n)$.
\item The level-$j$ form $C^*\Delta^{(j)}_{\Lambda'_{j+1}}C$ of Ba{\l}aban's large-field
expansion \cite[(1.60)]{B-CMP122a}, where $\Lambda'_{j+1}$ and the background $U_k^{(n+1)}$
below are those of that expansion; by the second-order identity for the actual factor this
form is the Schur form of the payment estimate above,
taken at the background $U_k^{(n+1)}$, which is free of $B_j$.
\end{enumerate}
Elsewhere in that proof, at complex pairs and at pairs that are not true pairs, only bounds
on kernels and on differences of kernels are used, namely in clause \textup{(i)} of the
$M$-freeness of the random-walk constants (Proposition~\ref{5.10:prop-m-freeness}), in
the conditional-mean lemma and in the tolerance lemma. In particular, the reference transfer
operator $T_0$ of the random-walk analysis is bounded there by clause \textup{(a)} of the
random-walk kernel bound, which also supplies the second member of $K_0$.

Three statements stand independently of \eqref{5.2:eq-true-pair-positivity-a} and are not
affected by it.
\begin{enumerate}
\item[(n1)] The reference operator $C_{(0)}=C^{(k)}(Z_0,0)$ of the random-walk analysis is
real symmetric and positive definite at $\sigma\equiv 0$, $\mathbf{U}=U$ on $Z_0'$,
$\mathbf{J}=0$. This is the reference point of the tolerance lemma and of the random-walk
kernel bound; it rests on the statement taking over Ba{\l}aban's cluster-expansion bounds,
which quotes the following sentence of \cite{B-CMP116}:
\begin{quote}
For the pair $(U,0)$ the operators are symmetric, and the measure is positive \dots\ we
replace the operators by the corresponding operators with $\sigma(Z)=0$, $\mathbf{U}=U$,
$\mathbf{J}=0$.
\end{quote}
\item[(n2)] The extensions to general pairs of $\mathcal{A}$ and of the covariances
$C_\Lambda$ and $C_I(\sigma)$ are
mutually inverse at minimisers only. No step of that proof uses the inverse relation
away from minimisers; in the conditional-mean lemma the operator $\mathsf{D}$ is defined
directly as the inverse $\mathsf{D}:=C_I(\sigma)^{-1}$ of the operator $C_I(\sigma)$ of that
lemma.
\item[(n3)] The quadratic form $-\tfrac12\langle B_j,C^*\Delta^{(j)}CB_j\rangle$ in the
level-$j$ variable $B_j$, occurring at the site lettered (P) of the correlation theorem
(Theorem~\ref{5.1:thm-correlation-theorem}), is
a level-$j$ form standing beside the higher levels. It enters there through the bound on
level-$j$ forms applied at that site, and does not use $\gamma_0$.
\end{enumerate}
A use of the lower bound of \eqref{5.2:eq-true-pair-positivity-a} at a pair that is not a
true pair is not covered by Corollary~\ref{5.2:cor-true-pair-positivity}.
\end{remark}

\begin{proof}
We first prove \textup{(a)}. Let $B$ satisfy $Q^{\flat}B=0$. Here $y$ runs over the block
sites and, for each $y$, $x$ over the other sites averaged to $y$; $V(y,x)$ is the parallel
transporter of the block background from $y$ to $x$, and $\operatorname{Ad}$ is the adjoint
representation of $SU(N)$ on its Lie algebra, so that $\operatorname{Ad}(V(y,x))$ is
orthogonal. In terms of the linearised average $\mathfrak{g}(B)$, define $\mu$ by
\begin{equation*}
\mu(y):=0,\qquad \mu(x):=\operatorname{Ad}\bigl(V(y,x)\bigr)^{-1}\mathfrak{g}(B)(y,x),
\end{equation*}
and set $B_a:=B-\bar D\mu$. Then $B_a$ lies in $\{\mathfrak{g}=0,\ Q^{\flat}=0\}$. Since
$\mathfrak{g}$ and $Q^{\flat}$ are linear and $Q^{\flat}B=0$, this membership is the pair of
identities $\mathfrak{g}(\bar D\mu)=\mathfrak{g}(B)$ and $Q^{\flat}\bar D\mu=0$; these two
identities are verified from the definitions of $\mathfrak{g}$, $\bar D$ and $Q^{\flat}$,
given with the linearised averages and difference operators.

On $\{\mathfrak{g}=0\}$ the linearised average $\mathfrak{a}(B)$ vanishes, and there the
constraints $Q^{\flat}=0$ and $Q_1=0$ coincide, by the definitions of $Q^{\flat}$ and $Q_1$.
Thus the two constraint sets coincide:
\begin{equation*}
\{\mathfrak{g}=0,\ Q^{\flat}=0\}=\{\mathfrak{g}=0,\ Q_1=0\}.
\end{equation*}
The right-hand set is the slice on which the slice positivity corollary is stated. That
corollary therefore gives
\begin{equation}\label{5.2:eq-positivity-scope-2}
\langle B_a,\Delta_k B_a\rangle\ \ge\ \gamma_0^{\mathrm{sl}}\,\|B_a\|^2 .
\end{equation}
The constant $\gamma_0^{\mathrm{sl}}$ is read against $\|B_a\|^2$ itself, and not against a
norm of coordinates on the slice, because $\{\mathfrak{g}=0\}$ is not a coordinate slice.

Next we apply Corollary~\ref{5.2:cor-constrained-invariance} to the pair $(B,\mu)$. Its
hypotheses are that $\mu$ vanish on the region $\Lambda'$ of that corollary and that
$\tilde D^{(2)}=hC^{(2)}$. Both hold: $\mu$ vanishes at the block sites $y$ by construction,
and the region $\Lambda'$ of that corollary is the set of these block sites; $\tilde D^{(2)}=hC^{(2)}$
is a hypothesis of \textup{(a)} (Ba{\l}aban's choice, \cite{B-CMP109}). Both corollaries are
applied at the true pair of \textup{(a)}. The corollary gives
$\langle B_a,\Delta_kB_a\rangle=\langle B,\Delta_kB\rangle$ exactly. The
form $\Delta^{(k)}$ is, by the definition of the fluctuation forms, $\Delta_k$ plus the
gauge-fixing form, whose value at $B$ is
$\sum_{y,x}|\mathfrak{g}(B)(y,x)|^2$; since each $\operatorname{Ad}(V(y,x))$ is orthogonal
and $\mu(y)=0$, this sum equals $\|\mu\|^2$. Therefore
\begin{equation*}
\langle B,\Delta^{(k)}B\rangle=\|\mu\|^2+\langle B_a,\Delta_kB_a\rangle .
\end{equation*}
Inserting \eqref{5.2:eq-positivity-scope-2} gives
$\langle B,\Delta^{(k)}B\rangle\ge\|\mu\|^2+\gamma_0^{\mathrm{sl}}\|B_a\|^2$. On the other
hand $B=B_a+\bar D\mu$ and $\pi_f$ has norm at most one. Each bond value of $\bar D\mu$ is the
difference of the values of $\mu$ at the two endpoints of the bond, one of them transported
by an orthogonal map, and each site is an endpoint of $2d$ bonds, so
$\|\bar D\mu\|^2\le 4d\|\mu\|^2$. Hence
\begin{equation*}
\|\pi_fB\|^2\ \le\ 2\|B_a\|^2+2\|\bar D\mu\|^2\ \le\ 2\|B_a\|^2+8d\|\mu\|^2 .
\end{equation*}
With $c:=\min\{\gamma_0^{\mathrm{sl}}/2,\ 1/(8d)\}$ one has
$2c\|B_a\|^2\le\gamma_0^{\mathrm{sl}}\|B_a\|^2$ and $8dc\|\mu\|^2\le\|\mu\|^2$, and we obtain
\begin{align*}
\langle B,\Delta^{(k)}B\rangle
&\ \ge\ \|\mu\|^2+\gamma_0^{\mathrm{sl}}\|B_a\|^2
\ \ge\ c\,\bigl(2\|B_a\|^2+8d\|\mu\|^2\bigr)\\
&\ \ge\ \min\Bigl\{\frac{\gamma_0^{\mathrm{sl}}}{2},\ \frac{1}{8d}\Bigr\}\,\|\pi_fB\|^2,
\end{align*}
which is \eqref{5.2:eq-positivity-scope-1}.

We turn to the positivity in \textup{(n1)}. In the representation underlying
Corollary~\ref{5.2:cor-true-pair-positivity}, every vertex of that representation that reads
$\mathbf{J}$ vanishes at $\mathbf{J}=0$. What remains is the direct sum of the inverse of the
operator $\tilde G_0$ of that representation and the form $\mathfrak{d}^*\mathfrak{d}$ of the
difference operator $\mathfrak{d}$,
\begin{equation*}
\tilde G_0^{-1}\oplus\mathfrak{d}^*\mathfrak{d}.
\end{equation*}
This operator is positive under condition \cite[(3.35)]{B-CMP99b} alone, by
\cite[Theorem~3.11]{B-CMP99b}. Hence the positivity in \textup{(n1)} does not rest on
\eqref{5.2:eq-true-pair-positivity-a}.
\end{proof}

\begin{lemma}[Homogeneous containment of the response map]\label{5.2:lem-response-containment}
Assume the following.
\begin{enumerate}
\item[(a)] Ba{\l}aban's analyticity statements and bounds for the response field and the functions built
on it, \cite[(1.17)--(1.22)]{B-CMP116} and \cite[(3.14)--(3.17)]{B-CMP109}, taken by statement.
\item[(b)] The ratio clause of the radii lemma, which gives Ba{\l}aban's radius $\alpha_2$ the lower
bound $\alpha_2\ge c_2\,\alpha_{*,k}$ with a constant $c_2>0$.
\item[(c)] The layer estimate for the shift of the response at the layers $n\le k$. It rests on
\cite[Proposition~9, (190)]{B-CMP102b}, taken by statement.
\end{enumerate}
Then there is a constant $K_R=K_R(d,L,M,\kappa_1)$ with the following property. Fix a step $k$,
a unit $v$ at the site (T) of Definition~\ref{5.2:def-sites-units}, with read set $S_v$ and with
any of its labels. Let $p=(\square,Y_0)$ be a localisation datum and $(\mathbf{U},\mathbf{J})$ a point
of the output tube over $Y_0$. Let $\sigma$ satisfy $|\sigma(\Delta)|\le e^{\kappa_1}$ for every
cube $\Delta$. Let $\mathsf{B}$ be a $\mathfrak{g}^{\mathbb{C}}$-valued field on the free bonds of
$Y_0\cap\Omega_{k+1}$; here $\mathfrak{g}^{\mathbb{C}}$ denotes the complexification of the Lie
algebra of $SU(N)$. Let $\mathsf{t}\in\mathbb{C}$, and suppose
\begin{equation}\label{5.2:eq-response-containment-a}
(1+\eta_p|\mathsf{t}|)\,K_R\,g_k\sup|\mathsf{B}|\le\alpha_{*,k}.
\end{equation}
Then $\mathsf{R}_p(\mathsf{t},\sigma;(\mathbf{U},\mathbf{J}),\mathsf{B})\restriction S_v$ is an
admissible background, in the sense of the definition of admissible backgrounds, and depends
jointly holomorphically on its arguments. At real $\mathsf{B}$ with $\sup|\mathsf{B}|<\mathfrak{b}_k$
it coincides with the response map of Definition~\ref{5.2:def-sites-units}. Here $\eta_p\le 1$ is
the decay factor between $\square$ and $S_v$ fixed in that definition; it is read from the decay
of $\delta\mathbf{H}_j$ in \cite[(3.9)]{B-CMP109}.
\end{lemma}

\begin{proof}
The proof has four parts and a conclusion: the content of Ba{\l}aban's containment, the top level,
the layers, and the rim bonds with real data.

\emph{What Ba{\l}aban's containment is.}
In \cite[p.~6]{B-CMP116} the field $B'$, a free variable in the preceding estimates there, is
specialised to
\begin{equation*}
B'=g_kCB-hD(g_kCB),
\end{equation*}
which is \cite[(1.19)]{B-CMP116}; the operators $C$, $D$, the function $h$ and the constants
$C_1$, $B_0$, $\varepsilon_1$ used below are those of \cite[(1.17)--(1.22)]{B-CMP116}.
It satisfies $|B'|\le C_1g_k|B|$.

The function $\mathbf{H}_k(\mathsf{s}(Y_0),B')$ is analytic in the decoupling parameters
$\mathsf{s}(Y_0)$ and in $B$ on the set where $|\mathsf{s}(Y_0)|\le e^{\kappa_1}$ and
$g_k|B|<\varepsilon_1$. In the notation of this book these parameters are the $\sigma(\Delta)$.
On that set $\mathbf{H}_k(\mathsf{s}(Y_0),B')$ satisfies the estimate of
\cite[pp.~6--7]{B-CMP116} by
\begin{equation*}
4B_0C_1e^{16\kappa_1}g_k|B| ,
\end{equation*}
and the same estimate holds in
all admissible norms, provided $e^{32\kappa_1}\varepsilon_1$ is smaller than an absolute constant.
All these bounds are uniform on the domain of $(\mathbf{U},\mathbf{J})$.

On \cite[p.~7]{B-CMP116} the expression of \cite[(1.10)]{B-CMP116} is extended analytically in
$t_\square$ on the set
\begin{equation*}
|t_\square|\,4B_0C_1e^{16\kappa_1}g_k|B|\le\tfrac12\alpha_2
\end{equation*}
(the same constant appears in \cite[(1.22)]{B-CMP116} as $8B_0C_1e^{16\kappa_1}\alpha_2^{-1}$).
In the response map the complex amplitude $\mathsf{t}$ occupies the place of $t_\square$ for the
one block $\square$ of $p$, and the piece it multiplies is seen on $S_v$ with at most the factor
$\eta_p$ (Definition~\ref{5.2:def-sites-units}).

By \cite[p.~273]{B-CMP109}, the bound \cite[(3.17)]{B-CMP109} uses only two facts: that
$\mathbf{H}_j(B(t))$ satisfies \cite[(3.14)]{B-CMP109} with $\tfrac13\alpha_2$, and that
$\delta\mathbf{H}_j$ satisfies \cite[(3.9)]{B-CMP109}. The functions produced by the later
localisation again satisfy these two bounds, so \cite[(3.17)]{B-CMP109} is preserved.

By \cite[Lemma~1, p.~9]{B-CMP116} the functions are defined and analytic on a domain of the
configuration $(\mathbf{U},\mathbf{J})$ and the field $B$ whose factor in the field variable is
$\{B:|B|<\varepsilon_1g_k^{-1}\text{ on }Y\}$, where $Y$ is the region of that lemma. At complex
$B$ this factor is a complex polydisc; this is part~(iii) of the lemma of this book
recording the analyticity domains of Ba{\l}aban's functions.

Ba{\l}aban's containment is therefore a bound on one number, $g_k\sup|B|$. The estimate
$4B_0C_1e^{16\kappa_1}g_k\sup|B|$ enters twice: with coefficient $|t_\square|$ against
$\tfrac12\alpha_2$, and with coefficient $1$ against the shift margin $\tfrac13\alpha_2$
(\cite[p.~7]{B-CMP116} and \cite[(3.14)]{B-CMP109}); both readings are homogeneous of degree one
in $g_k\sup|B|$. In the response map the field $\mathsf{B}$, extended to complex values, occupies
the place of Ba{\l}aban's $B$ on the free bonds of $Y_0\cap\Omega_{k+1}$, so his estimates are read
with $g_k\sup|\mathsf{B}|$ in place of $g_k\sup|B|$. Since
$1+\eta_p|\mathsf{t}|\ge1$, hypothesis \eqref{5.2:eq-response-containment-a} implies
\begin{equation}\label{5.2:eq-response-containment-1}
K_R\,g_k\sup|\mathsf{B}|\le\alpha_{*,k}.
\end{equation}
Both readings of $g_k\sup|\mathsf{B}|$ will be taken inside
\eqref{5.2:eq-response-containment-1}.

\emph{The top level.}
By the radii of \cite[(2.28)]{B-CMP119} and hypothesis~(b),
$\alpha_2\ge c_2\alpha_{*,k}$. The conditions on the radii in \cite[p.~263]{B-CMP109} are
homogeneous of degree one in the radii, so they may be read with $\alpha_2$ bounded below by
$c_2\alpha_{*,k}$. Require
\begin{equation}\label{5.2:eq-response-containment-2}
K_R\ge 16B_0C_1e^{16\kappa_1}/c_2 .
\end{equation}

First consider the undilated part. By \eqref{5.2:eq-response-containment-1} and
\eqref{5.2:eq-response-containment-2},
\begin{equation*}
4B_0C_1e^{16\kappa_1}g_k\sup|\mathsf{B}|
\le\frac{4B_0C_1e^{16\kappa_1}\alpha_{*,k}}{K_R}
\le\frac{4B_0C_1e^{16\kappa_1}\alpha_2}{c_2K_R}
\le\tfrac14\alpha_2<\tfrac13\alpha_2 .
\end{equation*}
So the bound \cite[(3.14)]{B-CMP109} with $\tfrac13\alpha_2$ holds at the base point.

Next consider the dilated part. The amplitude $\mathsf{t}$ multiplies the piece of the response
cut by the block $\square$. On $S_v$ this piece is seen with at most the fraction $\eta_p$ of its
size. By \eqref{5.2:eq-response-containment-a} and the same chain of inequalities,
\begin{equation*}
\eta_p|\mathsf{t}|\,4B_0C_1e^{16\kappa_1}g_k\sup|\mathsf{B}|\le\tfrac14\alpha_2 .
\end{equation*}
This lies inside Ba{\l}aban's circle $\tfrac12\alpha_2$. Thus the two printed readings, on
\cite[p.~7]{B-CMP116} and \cite[p.~273]{B-CMP109}, hold as displayed.

The absolute conditions \cite[(1.13)--(1.16)]{B-CMP116} require
$e^{32\kappa_1}C_1g_k\sup|\mathsf{B}|\le c_{\mathrm{abs}}$, where $c_{\mathrm{abs}}$ is the absolute
constant of those conditions. By \eqref{5.2:eq-response-containment-1},
\begin{equation*}
e^{32\kappa_1}C_1g_k\sup|\mathsf{B}|\le e^{32\kappa_1}C_1\alpha_{*,k}/K_R ,
\end{equation*}
a number free of $\mathsf{B}$. That this number is at most $c_{\mathrm{abs}}$, at every step $k$,
is one of the standing conditions of this book, in which $c_{\mathrm{abs}}$ is fixed.
Hence the requirement holds, and the absolute conditions \cite[(1.13)--(1.16)]{B-CMP116} hold.

\emph{The layers.}
Take a layer $n<k$. The shift reaches layer $n$ with the factor
\begin{equation*}
L^{-2(k-n)}\prod_{n<m\le k}e^{-2\delta_0MR_m};
\end{equation*}
here $\delta_0$ is the decay rate of \cite[(3.9)]{B-CMP109}, $M=L^{m'}$ is Ba{\l}aban's auxiliary
parameter of that name, and
$R_m=R(x_m)$ is Ba{\l}aban's radius at step $m$.
The room available for it is the schedule's
$(1-\theta_S)\,\beta_{\mathrm{pr}}\,2^{-(k+1-n)}\,\alpha_{\cdot,n}$, where $\alpha_{\cdot,n}$ stands
for each of the radii $\alpha_{i,n}$, $\alpha_{*,n}$ at level $n$. Here
$\theta_S\in[\tfrac12,\tfrac89]$ is the layer fraction, so $1/(1-\theta_S)\le9$.

In the layer estimate of hypothesis~(c) the arithmetic is linear in
\begin{equation*}
\mathsf{v}_k:=\|C\|\,g_k\sup|\mathsf{B}| ,
\end{equation*}
where $\|C\|$ is the norm of the operator $C$ of \cite[(1.19)]{B-CMP116}; under
$|\sigma(\Delta)|\le e^{\kappa_1}$ the constant $C_H$ of the layer estimate becomes
$C_He^{16\kappa_1}$. By
\eqref{5.2:eq-response-containment-1}, $\mathsf{v}_k\le\|C\|\alpha_{*,k}/K_R$. The layer estimate
therefore applies at every layer $n\le k$ provided
\begin{equation}\label{5.2:eq-response-containment-3}
K_R\ge\max\Bigl\{\frac{4(1+\beta_0)^2\|C\|C_He^{16\kappa_1}}{(1-\theta_S)\beta_{\mathrm{pr}}},\
\frac{8\|C\|C_He^{16\kappa_1}}{\beta_{\mathrm{pr}}}\Bigr\},
\end{equation}
where $\beta_0$ is Ba{\l}aban's parameter which, through the factor $(1+\beta_0)^2$, enters the
comparison of the radii at different levels.
The first member of the maximum secures the room
$(1-\theta_S)\beta_{\mathrm{pr}}2^{-(k+1-n)}\alpha_{\cdot,n}$ at the layers $n<k$, the second the
top layer $n=k$.
This mechanism is the one described in \cite[p.~276]{B-CMP119}: as Ba{\l}aban observes there,
the localisation of the fluctuation field in $\Omega_{k+1}$, together with the exponential decay of
the propagators and of the minimisers, produces the additional small factors.

\emph{Rim bonds and real data.}
The field of \cite[(3.15)]{B-CMP119} lives on all of the free bonds $\mathfrak{X}$
\cite[p.~267]{B-CMP119}. The field $x\oplus A$ is exactly that field. It is formed from the rim
datum $x=A_k\restriction\mathfrak{r}$ on the rim bonds $\mathfrak{r}$ and the fluctuation variable
$A$ of \cite[(3.15)]{B-CMP119} on the other free bonds. Hence the bounds above cover the rim bonds
as well.

At real $\mathsf{B}$ with $\sup|\mathsf{B}|<\mathfrak{b}_k$, the map
$\mathsf{R}_p(\mathsf{t},\sigma;(\mathbf{U},\mathbf{J}),\mathsf{B})$ is the response map of
Definition~\ref{5.2:def-sites-units}. The radius of Ba{\l}aban's real circle, read from
\cite[(1.22)]{B-CMP116} at $\sup|\mathsf{B}|<\mathfrak{b}_k$, is denoted
\begin{equation*}
\mathsf{w}_T^{(1.22)}:=\frac{\alpha_2}{4B_0C_1e^{16\kappa_1}\delta_k},
\qquad \delta_k=g_k\mathfrak{b}_k ;
\end{equation*}
it is recorded for comparison only, the disc furnished by \eqref{5.2:eq-response-containment-a}
being the one used.

\emph{Conclusion.}
Joint holomorphy in $(\mathsf{t},\sigma,(\mathbf{U},\mathbf{J}),\mathsf{B})$ follows from the facts
recalled in the first part: the analyticity in the decoupling parameters and in the field of
\cite[pp.~6--7]{B-CMP116}; the analyticity in the configuration $(\mathbf{U},\mathbf{J})$ and in
the field on the domain of \cite[Lemma~1, p.~9]{B-CMP116}; and the analytic extension in the
amplitude of \cite[p.~7]{B-CMP116}. The bounds being uniform on the domain of
$(\mathbf{U},\mathbf{J})$, the map is locally bounded, and a locally bounded map which is
holomorphic in each of these variables separately is jointly holomorphic.

Define $K_R$ to be the maximum of the three numbers displayed in
\eqref{5.2:eq-response-containment-2} and \eqref{5.2:eq-response-containment-3}. Then both
readings at the top level, the absolute conditions \cite[(1.13)--(1.16)]{B-CMP116} (by
\eqref{5.2:eq-response-containment-1} and the standing condition
$e^{32\kappa_1}C_1\alpha_{*,k}/K_R\le c_{\mathrm{abs}}$), and the layer estimate at every layer
hold. Hence
$\mathsf{R}_p(\mathsf{t},\sigma;(\mathbf{U},\mathbf{J}),\mathsf{B})\restriction S_v$ is an
admissible background, jointly holomorphic, and equal to the response map of
Definition~\ref{5.2:def-sites-units} at real $\mathsf{B}$ with $\sup|\mathsf{B}|<\mathfrak{b}_k$.
\end{proof}

\begin{definition}[The activity threshold and the complex radius ratio]
\label{5.2:not-activity-threshold}
The following constants are fixed before $\gamma$, in the order of choice listed with the
constants of the rim construction, display~\eqref{5.2:eq-rim-constants-a}.

\emph{The weight constant.} Put
\begin{equation*}
C_{\mathsf{W}} := O(1)\,E_0\,\bigl(C_IS_I + C\,S_\omega\bigr),
\end{equation*}
where $E_0$, $C_I$ and $C$ are constants fixed before $\gamma$, and where $S_I$ and $S_\omega$ are
the first members of the comparison constant $C_\sharp$ of
Definition~\ref{5.2:def-comparison-constant}. They are suprema, over the couplings below the first
ceiling $\gamma^\circ$, of sums that do not depend on $K$, and they do not increase when the
ceiling is lowered. One has $C_{\mathsf{W}} \le \bar C_{\mathsf{W}}$, the constant of
display~\eqref{5.2:eq-comparison-constant-a}, and $\bar C_{\mathsf{W}}$ may be used in place of
$C_{\mathsf{W}}$ throughout.

\emph{The activity threshold.} The activity threshold $\varepsilon_V^{*}$ is the threshold on the
activity $\varepsilon_2$ in the activity lemma for the cluster expansion, read at
$\varepsilon_2 := 2\varepsilon_V$. It is a positive number, fixed so that, for every
$\varepsilon_V \in (0,\varepsilon_V^{*}]$, the two inequalities in the first line of
display~\eqref{5.2:eq-activity-threshold-a} below hold and the quantity in its second line is small
in the sense of the convergence condition of \cite[(2.39)--(2.40)]{B-CMP116}, all three with
$e^{(5/2)L\kappa}$ in place of $e^{5\rho_7}$ and $e^{5\rho_9}$; it therefore depends neither on
$g$ nor on $\delta$.
The first condition is the member of \cite[(2.28)]{B-CMP116} and gives the level-$k$ bound
\cite[(2.29)]{B-CMP116}. The second is the condition at \cite[(2.37)]{B-CMP116}, where the family
sum at level $k+1$ is estimated by means of \cite[(2.29)]{B-CMP116}. The third concerns the
activity kept after that step.

The activity $\varepsilon_V$ carries no factor of the radius $\alpha_6$ of
\cite[(2.28)]{B-CMP116}, so the powers of $\alpha_6^{-1}$ appear explicitly: one in the first
condition, which is a level-$k$ bound; two in the second, where the level-$(k+1)$ families are
summed by means of the level-$k$ bound; and two in the third, which is the constant of the kept
activity in item~$(\zeta)$ of the list of constants of the one-step analysis, where this count is
carried. Since $\varepsilon_V^{*}$ is independent
of $g$ and enters $c_V$ and $c_A$ below, these factors cannot be absorbed by the choice of
$\gamma$.

The rates are $\rho_m = (1-m\delta)\tfrac12 L\kappa$, $m = 1,\dots,10$. Under the conditions in
the first two lines of \eqref{5.2:eq-activity-threshold-a} the bounds delivered by the activity
lemma are at most $\tfrac18\min\{E^{*},B_0\}$, where $B_0$ is the constant of
\cite{B-CMP116}.

\emph{The constants $c_V$, $c_A$ and the complex radius ratio.} With $K_R$ the constant of the
homogeneous containment of the response map (Lemma~\ref{5.2:lem-response-containment}, with
hypothesis~\eqref{5.2:eq-response-containment-a}), $C_A$ the constant of the first-stage bound by
two contour integrals, display~\eqref{5.2:eq-first-stage-bound-a}, and $\Theta_\flat$, $T^{\mathbb{C}}$,
$\alpha_{*,k}$ as in display~\eqref{5.2:eq-rim-constants-a}, set
\begin{equation}\label{5.2:eq-activity-threshold-a}
\begin{gathered}
2\varepsilon_V\,\alpha_6^{-1}e^{5\rho_7} \le 1,\qquad
4(L+2)^4\,O(1)\,\alpha_6^{-2}\,\varepsilon_V\,e^{5\rho_7} \le 1,\\
4(L+2)^4\,O(1)\,\alpha_6^{-2}\,\varepsilon_V\,e^{5\rho_9},\\
c_V := \min\Bigl\{\frac{\varepsilon_V^{*}}{16K_RC_AC_{\mathsf{W}}},\ \frac{1}{4K_R}\Bigr\},
\qquad c_A := \frac{c_V}{\Theta_\flat+2},\\
\theta^{\mathbb{C}}_k := \frac{K_R\,g_k\,T^{\mathbb{C}}}{\alpha_{*,k}},\qquad
\mathsf{w}^{\mathbb{C}} := \frac{1}{\theta^{\mathbb{C}}_k}-1 .
\end{gathered}
\end{equation}
Equivalently, by the definition of $\mathsf{w}^{\mathbb{C}}$,
$(1+\mathsf{w}^{\mathbb{C}})\,K_R\,g_k\,T^{\mathbb{C}} = \alpha_{*,k}$.
Then:
\begin{enumerate}
\item[(a)] One has
\begin{equation}\label{5.2:eq-activity-threshold-2}
\theta^{\mathbb{C}}_k
= K_R\Bigl[\frac{A_1}{C_*}\bigl(\log g_k^{-2}\bigr)^{p_0-q} + (\Theta_\flat+1)\,c_A\Bigr].
\end{equation}
\item[(b)] The second member of \eqref{5.2:eq-activity-threshold-2}, which comes from $\varrho$,
satisfies $K_R(\Theta_\flat+1)c_A < \tfrac14$.
\item[(c)] $\mathsf{w}^{\mathbb{C}}$ is bounded by a constant independent of $g_k$, hence stays
bounded as $\gamma \to 0$.
\item[(d)] $\mathsf{w}^{\mathbb{C}} \ge 1$ whenever $\theta^{\mathbb{C}}_k \le \tfrac12$.
\end{enumerate}

The containment required of these radii in the statement fixing the radii of the amplitude discs
is the ratio condition $\theta^{\mathbb{C}}_k \le \tfrac12$; it does not require
$g_kT^{\mathbb{C}}$ to tend to zero. Its member in~(b)
is independent of $g$ and is met by the choice of $c_A$, before $\gamma$.

Of the demands that this containment and the estimates made under it place on the radii, only the
following are conditions on $\gamma$:
\begin{itemize}
\item the smallness of the first member $K_R(A_1/C_*)(\log g_k^{-2})^{p_0-q}$ of
\eqref{5.2:eq-activity-threshold-2};
\item the smallness of $g_k(T^{\mathbb{C}})^3$;
\item the smallness of $g_kT^{\mathbb{C}}/\varepsilon_1$, where $\varepsilon_1$ is the analysis
constant, independent of $g$, at the absolute radii of the configuration space $\mathcal{U}$.
\end{itemize}
No lower bound on $c_A$ is required other than conditions absorbed into the choice of $\gamma$,
namely $r_A \ge \max\{2,C_{\mathrm{far}}\}$ and the smallness of $K_Q/r_A$. In the order of
choice, $\varepsilon_V^{*}$, $c_V$, $c_A$ and with them $\varrho$, $r_A$, $T^{\mathbb{C}}$,
$\varpi_\gamma$ are defined after $\alpha_6$.
\end{definition}

\begin{proof}
We first show that the threshold $\varepsilon_V^{*}$ may be taken independent of $\delta$. Since
$0 \le 1-m\delta \le 1$, one has $\rho_m \le \tfrac12 L\kappa$ and hence
$e^{5\rho_m} \le e^{(5/2)L\kappa}$ for $m=7,9$. The conditions in the first two lines of
\eqref{5.2:eq-activity-threshold-a} with $e^{(5/2)L\kappa}$ in place of $e^{5\rho_7}$ and
$e^{5\rho_9}$, which fix $\varepsilon_V^{*}$, therefore imply them for every admissible $\delta$.
Thus $\varepsilon_V^{*}$ does not depend on $\delta$.

For~(a), substitute the definitions in display~\eqref{5.2:eq-rim-constants-a}. There
$T^{\mathbb{C}} = \mathfrak{b} + (\Theta_\flat+1)\varrho$ and
$\varrho = c_A\alpha_{*,k}/g_k$, so that
\begin{equation*}
\theta^{\mathbb{C}}_k
= K_R\Bigl[\frac{g_k\mathfrak{b}}{\alpha_{*,k}} + (\Theta_\flat+1)\,c_A\Bigr].
\end{equation*}
Moreover $r_A = \varrho/\mathfrak{b} = c_A(C_*/A_1)(\log g_k^{-2})^{q-p_0}$, whence
\begin{equation*}
\frac{g_k\mathfrak{b}}{\alpha_{*,k}}
= \frac{g_k\varrho}{\alpha_{*,k}\,r_A}
= \frac{c_A}{r_A}
= \frac{A_1}{C_*}\bigl(\log g_k^{-2}\bigr)^{p_0-q}.
\end{equation*}
Inserting this gives \eqref{5.2:eq-activity-threshold-2}.

For~(b), the definitions in \eqref{5.2:eq-activity-threshold-a} give $c_V \le 1/(4K_R)$ and
$c_A = c_V/(\Theta_\flat+2)$. Therefore
\begin{equation*}
K_R(\Theta_\flat+1)\,c_A
\le \frac{\Theta_\flat+1}{4(\Theta_\flat+2)}
< \frac14 .
\end{equation*}

For~(c), the first member in \eqref{5.2:eq-activity-threshold-2} is non-negative, because
$A_1, C_* > 0$ and $\log g_k^{-2} > 0$. Hence
$\theta^{\mathbb{C}}_k \ge K_R(\Theta_\flat+1)c_A > 0$, so $\mathsf{w}^{\mathbb{C}}$ is well
defined and
\begin{equation*}
\mathsf{w}^{\mathbb{C}}
\le \frac{1}{K_R(\Theta_\flat+1)\,c_A} - 1 .
\end{equation*}
The right-hand side involves only constants fixed before $\gamma$ and does not depend on $g_k$.

For~(d), if $\theta^{\mathbb{C}}_k \le \tfrac12$ then $1/\theta^{\mathbb{C}}_k \ge 2$, so
$\mathsf{w}^{\mathbb{C}} \ge 1$.
\end{proof}

\begin{lemma}[Complex localisation of the effective action]\label{5.2:lem-complex-localisation}
Let $\gamma\le\gamma_V$, the ceiling fixed after the proof below (its last member is
\eqref{5.2:eq-complex-localisation-1}). At site (T) of every step $k$, for every output term of
the exact refinement of the rim cut-offs \eqref{5.2:eq-rim-refinement-a}, and for
$w\in\mathcal{W}_k$ or $w$ constant, the localised action of \cite[(3.25)]{B-CMP119} is a sum
$\sum_Y V(Y)$ over $Y\in\mathbf{D}_k$. Here $\operatorname{space}(Y)$ denotes the first factor of
the product of spaces in \cite[(3.43)]{B-CMP119}, read layer by layer as in the layered analyticity
lemma. Each term
\begin{equation*}
V(Y)=V\bigl(Y;(\mathbf{U},\mathbf{J}),\mathsf{B};w;\mathbf{A}_{<k}\bigr)
\end{equation*}
has the following four properties.
\begin{enumerate}
\item[(a)] $V(Y)$ is holomorphic on the product of $\operatorname{space}(Y)$, the set of
$\mathfrak{g}^{\mathbb{C}}$-valued fields $\mathsf{B}$ on the free bonds of $Y$ with
$|\mathsf{B}|<T^{\mathbb{C}}$, and $\mathcal{W}_k(Y)$. The one exception is the set of bonds of
$(R^{\sim2}\cap S)^{*}$, where $R$ and $S$ are regions of the site-(T) geometry (here $R$ denotes a
region, not a radius). There $\mathsf{B}$ is not a holomorphic variable: it enters as a real
parameter, with $|\mathsf{B}|\le C_{\mathrm{far}}\mathfrak{b}$.
\item[(b)] $V(Y)$ reads $\mathsf{B}$ on $Y$ only and reads $((\mathbf{U},\mathbf{J}),\mathbf{A},w)$
on $Y^{+}$. One has $Y^{+}=Y$ for a potential made of plain units. The units composing
$V^{\mathrm{rel}}$ below read the strips of their pockets, and in (a) the sets
$\operatorname{space}(Y)$ and $\mathcal{W}_k(Y)$ are then taken over $Y^{+}$. Moreover $V(Y)$ is
invariant under the gauge transformations of \cite[(3.29)]{B-CMP109}, with
$\mathsf{B}\mapsto\operatorname{Ad}_u\mathsf{B}$.
\item[(c)] A unit $F$ of the table of units (Proposition~\ref{5.2:prop-unit-table}) is called
\emph{wrapped} if it is of the fourth kind and its pocket set $\mathcal{P}_F$ is non-empty; every
other unit of the table is called \emph{plain}. The wrapped units $F$ occur among the boundary terms
$\mathbf{B}'^{(j)}$, $j\le k$, created by
$\mathbf{R}$, among their descendants, and among the $\mathbb{C}$-terms. Let $V^{\mathrm{pl}}$ be the
sum of the potentials of the plain units, obtained from the first-stage bound in steps 1--4 of the
proof below. Let the wrapped
units be filed under the labelling rule for wrapped units, on $Y=\operatorname{lab}(F,\mathsf{y})
\in\mathbf{D}_k$. Then
\begin{equation}\label{5.2:eq-complex-localisation-a}
\begin{gathered}
|V(Y)|\le\nu(Y):=\varepsilon_V e^{-\kappa_V d_k(Y)},\qquad
\kappa_V:=(1-2\delta)\kappa,\qquad \varepsilon_V\le\varepsilon_V^{*},\\
V=V^{\mathrm{pl}}+V^{\mathrm{rel}},\qquad
V^{\mathrm{rel}}(Y):=\sum_F\ \sum_{\mathsf{y}\in\mathfrak{Y}(F):\,\operatorname{lab}(F,\mathsf{y})=Y}
F_{\mathsf{y}},\\
|V^{\mathrm{rel}}(Y)|\le\varepsilon^{\mathrm{rel}}_k e^{-\kappa_V d_k(Y)},\qquad
\varepsilon^{\mathrm{rel}}_k:=C_\Sigma C_{\mathsf{a}}(1)\,\mathsf{b}_k(0)\,
e^{-\frac14\delta\kappa(7\check{R}_k-2)}.
\end{gathered}
\end{equation}
Here $C_\Sigma$ and $C_{\mathsf{a}}(1)$ are constants, and $\mathsf{b}_k(0)$ is a quantity at
level $k$, of the construction of the wrapped units.
Every label contains a source cell and a cube with a point on a strip. Hence $d_k(Y)\ge
7\check{R}_k-2$ for every label $Y$. For the wrapped units,
therefore, (c) is stated
at the plain length $d_k(Y)$ of a label $Y$. A label is never a domain containing a whole strip.

The share of the wrapped units in $\varepsilon_V$ is $\varepsilon^{\mathrm{rel}}_k\le\frac14
\varepsilon_V^{*}$, under
the ceiling on $\gamma$ that imposes this inequality, a ceiling of degree $-\infty$.
The sum $\sum_Y V^{\mathrm{rel}}(Y)$ equals
exactly the sum of the brackets of the wrapped units.

The outputs created by $\mathbf{T}$ whose cluster contains a $V^{\mathrm{rel}}$ are boundary terms
in the standard format, with $A:=X$. Each carries the factor $e^{-\kappa_{\mathrm{out}}
d_{k+1}(A)}$ and reads $X^{+}$; this is the relative form, on
labels, of the bound on boundary terms in the inductive hypothesis.
No output with $X\subset\Lambda_{k+1}$ contains a $V^{\mathrm{rel}}$.
\item[(d)] Assume the inductive bound on created boundary terms
(Definition~\ref{5.2:def-created-terms-bound}) and the first clause of the proposition on the
outputs of the $\mathbf{R}$-operation, at all steps $\le k$. Then $V(Y)$ decomposes over older
strong sets $\mathcal{S}'=(\mathcal{S}'_{j'})_{j'<k}$, and
\begin{equation*}
V(Y)=\sum_{\mathcal{S}'}V(Y,\mathcal{S}'),\qquad
\sum_{\mathcal{S}'}\exp\Bigl[\kappa_A(k)\sum_{j'<k}|\mathcal{S}'_{j'}|\Bigr]
\sup|V(Y,\mathcal{S}')|\le\nu(Y).
\end{equation*}
The supremum is taken also over $\mathfrak{D}_{\mathcal{S}'}$. When the inputs are read as orbit
data, it is taken also over the chart $\mathfrak{G}_{\varepsilon'}(Y)$ of the rotation presented
to the live exterior slots, with $\varepsilon':=\varepsilon-\eta^{T}-\mu_k\mathsf{e}/
\varepsilon_\star$, where $\varepsilon$ is the radius of the chart and $\eta^{T}$ and
$\mu_k\mathsf{e}/\varepsilon_\star$ are the losses of radius at site (T). The weight is the
age-free one,
\begin{equation*}
\kappa_A(k)=\min_{i\le k}\kappa_A^{\mathrm{pr}}(g_i)\le\kappa_A^{\mathrm{pr}}(g_{j'}).
\end{equation*}
It is the only weight used for the older strong sets in the bounds in the A-slot norm.
The weight of \eqref{5.2:eq-created-terms-bound-a} implies this one. The arrested
$\mathbb{C}$-units stored by the first clause of the proposition on the $\mathbf{R}$-operation
carry exactly this weight.
\end{enumerate}
\end{lemma}

\begin{proof}
Throughout, $o_\gamma(1)$ denotes a quantity that tends to zero as $\gamma\to0$.

\emph{Step 1: units of the first three kinds and $f_{\square}$ in the table of units.}
Let $|\mathsf{B}|<T^{\mathbb{C}}$ and $\eta_p|\mathsf{t}|\le\mathsf{w}^{\mathbb{C}}$. By
\eqref{5.2:eq-activity-threshold-a},
\begin{equation*}
\theta^{\mathbb{C}}_k=\frac{K_Rg_kT^{\mathbb{C}}}{\alpha_{*,k}},\qquad
\mathsf{w}^{\mathbb{C}}=\frac1{\theta^{\mathbb{C}}_k}-1.
\end{equation*}
Therefore
\begin{equation}\label{5.2:eq-complex-localisation-2}
(1+\eta_p|\mathsf{t}|)\,K_Rg_k\sup|\mathsf{B}|
\le(1+\mathsf{w}^{\mathbb{C}})\,K_Rg_kT^{\mathbb{C}}
=\frac{K_Rg_kT^{\mathbb{C}}}{\theta^{\mathbb{C}}_k}=\alpha_{*,k}.
\end{equation}
This is the hypothesis of the homogeneous containment of the response map,
Lemma~\ref{5.2:lem-response-containment}. By that lemma, $\mathsf{R}_p$ restricted to $S_v$ is an
admissible background.

The sizes $\mathfrak{N}(v)$ of Proposition~\ref{5.2:prop-unit-table} hold on all admissible
backgrounds. The lemmas from which
these sizes are drawn read the background only through the three defining conditions
of an admissible background,
with $\mathfrak{a}_n$ proportional to $\alpha$. Whether the shift came from a real or a complex
field is therefore invisible to them.

Consequently the first-stage bound, Lemma~\ref{5.2:lem-first-stage-bound}, applies without change
once $\mathsf{w}_T$ is replaced by $\mathsf{w}^{\mathbb{C}}$. On the polydisc it gives
\begin{equation*}
|v_p|\le2(\mathsf{w}^{\mathbb{C}})^{-1}\eta_p\,e^{-(\kappa_1-1)m(p)}\,\mathfrak{N}(v).
\end{equation*}
Moreover $v_p$ is holomorphic in $\mathsf{B}$, since it is given by Cauchy integrals of functions
holomorphic in $\mathsf{B}$ under one common bound.

The resummation over $p$, $\square$ and $Y_0$ does not involve $\mathsf{B}$. Since
$\theta^{\mathbb{C}}_k\le\frac12$,
\begin{equation*}
\frac4{\mathsf{w}^{\mathbb{C}}}=\frac{4\theta^{\mathbb{C}}_k}{1-\theta^{\mathbb{C}}_k}
\le8\theta^{\mathbb{C}}_k.
\end{equation*}
With the resummation constant $C_A$ of the first-stage bound,
Lemma~\ref{5.2:lem-first-stage-bound}, the contribution of these units to $V(Y)$ is at most
\begin{equation*}
8C_A\theta^{\mathbb{C}}_k\bigl[C_{\mathsf{W}}+C\mathfrak{m}_k
+Crg_k^2(\log g_k^{-2})^{2q}\bigr]\,e^{-\kappa_Vd_k(Y)}.
\end{equation*}

The weight-variation pieces $A(x_k(\cdot)-x_k,\,\cdot)$ have weight $x_k(\cdot)-x_k$, where
$x_k(\cdot)$ is the position-dependent inverse coupling of the density at level $k$; this weight is
supported in $Z_k$ \cite{B-CMP119}, so that $x_k(\cdot)$ equals $x_k$ outside $Z_k$. These pieces
are of the same kind as $f_\square$ in the table,
with the shift at zone $j$ bounded by $e^{-2\delta_0M\check{R}_k(k-j)}$ times the unit shift, one
factor $e^{-2\delta_0M\check{R}_k}$ for each of the $k-j$ zones crossed. This
holds because $H_k$ decays in scaled distance and each zone is at least $2MR_n$ wide. Hence the sum
\begin{equation*}
\sum_j\bigl(3\lambda(k-j)+2c_2\bigr)e^{-2\delta_0M\check{R}_k(k-j)}
\end{equation*}
is bounded independently of $K$ and is $o_\gamma(1)$. This is Ba{\l}aban's observation in
\cite{B-CMP119} that these terms are suppressed by the exponential decay of $H_k$.

\emph{Step 2: units of the fourth kind.} These are treated under the floor condition involving the
preparation count $\mathsf{N}_0$, which lets piles of depth controlled by $\mathsf{N}_0$ be paid by
distance at half rate. The units
in question are:
\begin{itemize}
\item the boundary terms $\mathbf{B}^{(j)}$, and the plain boundary terms $\mathbf{B}'^{(j)}$
created by $\mathbf{R}$ (the wrapped ones belong to $V^{\mathrm{rel}}$);
\item the $\mathbb{C}$-terms of this step's chain;
\item the $\mathbb{C}$-terms and the change-of-variables terms left on arrested components
\cite{B-CMP122a}, which together form the arrested left-overs $\mathcal{C}_{\mathrm{arr}}$, taken
on their frozen spaces, with containment
given on these frozen spaces by the containment theorem for arrested components;
\item the terms $A(x_k(\cdot)-x_k,\,\cdot)$, supported in $Z_k$ \cite{B-CMP119};
\item $V^{(k)}(S)$.
\end{itemize}

Their sizes are sup norms. When a unit is read as an orbit datum, its size is
$\mathfrak{N}^{\mathbb{C}}$ instead. The
sizes are piled at each unit's own level $m$ and at a point.

For the $\mathbb{C}$-terms the pile is at most $6(\mathsf{N}_0(\mathsf{T}_m)+3)$ deep. Of the
chains nested at the point, at most two have a stage family at level $m$, because their ranges of
attempt indices are pairwise disjoint except for one shared index. Each of these contributes at most
$3(N_0+3)$, by the count of the stage families of one chain meeting a point, made at the chain's own
step. Finally, for the step $k'\ge m$ at which each such chain runs,
$N_0(g_{k'})\le\mathsf{N}_0(\mathsf{T}_m)$, by the comparison with the reference recursion of
Definition~\ref{5.1:def-reference-recursion}.

The remaining sizes are as follows. The $\mathbf{B}^{(j)}$ have size $B_0$. The plain
$\mathbf{B}'^{(j)}$ have size $e^{\lambda_\theta}c_1^{\flat}$; these are the terms of the first
group lying off every strip. The wrapped ones belong to $V^{\mathrm{rel}}$ and are bounded in (c),
with $d_k(Y)\ge7\check{R}_k-2$ on their labels $Y$. The change-of-variables terms have size at most
$1$ per block.

The factor $(\mathsf{w}^{\mathbb{C}})^{-1}$ of the first-stage bound does not tend to zero as
$\gamma\to0$. It therefore cannot pay for these sizes, and the payment comes from distance.

Every such unit is anchored in $Z_k\cup S$. The reasons are the condition $X\cap Z_j\neq\emptyset$
of \cite[(2.41)]{B-CMP119} (there the condition is written with the enlarged set $Z_j^{\sim}$; the
difference is immaterial here), the inclusion $\Lambda_j^c\subset\Lambda_k^c=Z_k$, the support of
the weight-variation pieces in $Z_k$ \cite{B-CMP119}, and the location of the arrested components
\cite{B-CMP122a}. The field $\mathsf{B}$, on the other hand, lives on
\begin{equation*}
\Omega_{k+1}\subset\bigl((Z'_k)^{\sim8}\bigr)^{c},
\end{equation*}
where $Z'_k$ is the region so denoted in \cite[(3.2)--(3.5)]{B-CMP119}, the eight layers being
$4+1+1+2$ from those displays. Moreover $Y\cap\Lambda
\neq\emptyset$ and $\operatorname{dist}(S,\Lambda)\ge\mathfrak{d}_{\mathrm{far}}$, where $\Lambda$
is a region and $\mathfrak{d}_{\mathrm{far}}$ the far distance fixed in the site-(T) geometry.

Hence, for each such unit $v$, either $d_j(S_v)$ or the length of the response or random walk is at
least $cR_{k+1}$, for a positive constant $c$, in units of the current lattice, that is,
$L^{k-j}cR_{k+1}$ at the unit's own
scale. This is Ba{\l}aban's remark in \cite{B-CMP122a} that each renormalization step adds at least
ten layers of $MR_k$-cubes.

Under \eqref{5.2:eq-three-sites-a}, the pieces of the units of the fourth kind carry
$\eta_p^{1/2}$. The response rate is thus halved, and under $\kappa\ge10\log L$ and the standing
lower bound on $L$ the decay length is still at least $2R_{k+1}$. A quarter of that decay gives, at
each level $j$, the bound
\begin{equation*}
\Bigl(B_0+2+6\bigl(\mathsf{N}_0(\mathsf{T}_j)+3\bigr)C_0^{\sharp}+|x_j-x_k|\Bigr)
\max\bigl\{e^{-\kappa cR_{k+1}/(2L)},\,e^{-\frac12\delta_0cMR_k}\bigr\},
\end{equation*}
which tends to zero. Here the term $2$ comes from $e^{\lambda_\theta}c_1^{\flat}+1\le2$, which holds
under the conditions on the strip smallness $\alpha^{\sharp}$, on the degree function $p_0(\cdot)$,
read with the constant $C_{99}$, and on $\lambda_\theta$, among the conditions on the constants in
Definition~\ref{5.1:def-admitted-inputs}. Because of this bound the scale-dependent function
$c_1^{\flat}=C_{99}(\alpha^{\sharp})^{1/3}$ does not enter the defining member
\eqref{5.2:eq-complex-localisation-1} of $\gamma_V$.

These bounds are summed over the shell index $s=k-j$, with one more factor
$\max\{e^{-\kappa cR_{k+1}/(2L)},e^{-\frac12\delta_0cMR_k}\}$ per shell. The series converges
because $R(x)\ge(\log x)^{r_{\mathrm{pr}}}$, because $N_0$ grows only like $\log\log g^{-2}$
\cite{B-CMP122a}, and because the costs $|x_j-x_k|\le3\lambda s+2c_2$ and
$\mathsf{T}_j\le\mathsf{T}_k+c_s$, where $c_s$ bounds the increase of $\mathsf{T}$ over $s$ shells
and grows only logarithmically in $s$, sit $s$ shells
away. The other half of the decay supplies the
factor $e^{-\kappa_Vd_k(Y)}$, with no use of the ratio $(\mathsf{w}^{\mathbb{C}})^{-1}$.

The term $V^{(k)}(S)$ is bounded by any positive power of $g$ and is analytic in $H^{(k)}$
\cite{B-CMP119}. It is only $C^\infty$ in $\mathsf{B}$ on the bonds of $(R^{\sim2}\cap S)^{*}$,
through the quantity denoted $g$ in \cite[(1.22)]{B-CMP119}. This is the exception in (a).

\emph{Step 3: $\mathbf{P}^{(k)}$ and $\zeta_kG_y$.} The $g^{-2}\Psi$-pieces are
$\frac12\langle\mathsf{B},\mathcal{Q}^{\Psi}(Y,g_k\mathsf{B})\mathsf{B}\rangle$, as in
\cite[(1.43)]{B-CMP116}. By \cite[(1.40)]{B-CMP116} the kernel $\mathcal{Q}^{\Psi}(Y,g_k\mathsf{B})$
is bounded by a constant times $|g_k\mathsf{B}|$, being a bounded polynomial in $\mathsf{B}$ only
\cite{B-CMP116}. Since the pieces are quadratic in $\mathsf{B}$ besides, they are bounded by
\begin{equation*}
C(M,\kappa_1)\,g_k(T^{\mathbb{C}})^3e^{-\kappa_Vd_k(Y)}.
\end{equation*}
The $\Phi$-pieces vanish at $\mathsf{B}=0$ and are analytic for $g_k|\mathsf{B}|<\varepsilon_1$.
By the Schwarz-type bound for functions vanishing at the origin, part (a), they are bounded by
$2g_kT^{\mathbb{C}}/\varepsilon_1$ times the bound \cite[(1.36)]{B-CMP116}. Finally,
\begin{equation*}
\zeta_k(p_y)G_y(B')\le CL^4r\bigl(C_1g_kT^{\mathbb{C}}\bigr)^2,
\end{equation*}
from the size $CL^4r|B'|^2$ of this unit in the table of units and $|B'|\le C_1g_k|\mathsf{B}|$
\cite[p.~6]{B-CMP116}. Since $T^{\mathbb{C}}$ is a polynomial in $\log g_k^{-2}$, each of these
three bounds is $o_\gamma(1)$.

\emph{Step 4: the constant $\varepsilon_V$, and clauses (b) and (d).} Collecting steps 1--3 and the
share $\varepsilon^{\mathrm{rel}}_k$ of (c), which is counted among the quantities $o_\gamma(1)$
since its ceiling has degree $-\infty$, the
only contribution that does not tend to zero with $\gamma$ comes from the term
$8C_A\theta^{\mathbb{C}}_kC_{\mathsf{W}}$ of step 1, through the member $K_R(\Theta_\flat+1)c_A$ of
$\theta^{\mathbb{C}}_k$ that is free of $g$ (see \eqref{5.2:eq-activity-threshold-a}); the other
member of $\theta^{\mathbb{C}}_k$ is controlled by $\gamma$. Hence
\begin{equation*}
\varepsilon_V\le8C_AC_{\mathsf{W}}K_R(\Theta_\flat+1)c_A+o_\gamma(1)
\le\tfrac12\varepsilon_V^{*}+o_\gamma(1).
\end{equation*}
The second inequality follows from the definitions $c_A=c_V/(\Theta_\flat+2)$ and
$c_V\le\varepsilon_V^{*}/(16K_RC_AC_{\mathsf{W}})$ in \eqref{5.2:eq-activity-threshold-a}, which
give
\begin{equation*}
8C_AC_{\mathsf{W}}K_R(\Theta_\flat+1)c_A\le8C_AC_{\mathsf{W}}K_Rc_V\le\tfrac12\varepsilon_V^{*}.
\end{equation*}
Under $\gamma\le\gamma_V$ the quantities $o_\gamma(1)$ are in total at most
$\frac12\varepsilon_V^{*}$, so
that $\varepsilon_V\le\varepsilon_V^{*}$.

Clause (b). The locality of $V(Y)$ is the last sentence of \cite[Lemma~2, p.~11]{B-CMP116}. The
invariance under the gauge transformations of \cite[(3.29)]{B-CMP109}, with
$\mathsf{B}\mapsto\operatorname{Ad}_u\mathsf{B}$, rests on the $\operatorname{Ad}$-invariance of
$|\cdot|$, so that the domain $|\mathsf{B}|<T^{\mathbb{C}}$ is mapped onto itself, and on the
decoupling in the parameters $\sigma$ of \cite[(1.10)]{B-CMP116}.

Clause (d). The first-stage bound is linear in $f_v$. Each $v_p$ is therefore filed under the older
strong set $\mathcal{S}'$ of its unit, and since the quantity $\mathsf{W}$ of
\eqref{5.2:eq-first-stage-bound-a} is a sum of sizes $\mathfrak{N}$, the weighted sum over
$\mathcal{S}'$ is obtained by running step 2 on weighted sizes in place of sizes; below,
$|\mathcal{S}|$ stands for $\sum_{j'}|\mathcal{S}_{j'}|$. For the boundary terms created by
$\mathbf{T}$ the input is
\begin{equation*}
\sum_{\mathcal{S}}e^{\kappa_A(k)|\mathcal{S}|}\,\mathfrak{N}\bigl(\mathbf{B}^{(j)}(X,\mathcal{S})\bigr)
\le B_0e^{-\kappa d_j(X)}.
\end{equation*}
This is \eqref{5.2:eq-created-terms-bound-a} at the steps $k'=j\le k$: the weights there,
$\kappa_A^{\mathrm{pr}}(g_{k'-1})$ on the newest slot and $\kappa_A(k'-1)$ on the older slots, are
both at least $\kappa_A(k)=\min_{i\le k}\kappa_A^{\mathrm{pr}}(g_i)$, because $k'-1\le k$, and
$\kappa_{\mathrm{out}}\ge\kappa$.

For the boundary terms $\mathbf{B}'^{(j)}$ created by $\mathbf{R}$ in the first group of clause (e)
of the proposition on the outputs of the $\mathbf{R}$-operation, the datum is carried by clause
$(\gamma)$ of the corollary to the inductive bound on created boundary terms. Under it the
exterior $\mathbf{A}$-slots are parameters, and by the corollary on exterior slots as parameters the
strong
sets pass to unions and the weights multiply. The relative form, on labels, of the bound on boundary
terms in the inductive hypothesis then gives, on the label $A$,
\begin{equation*}
\sum_{\mathcal{S}}e^{\kappa_A(k)|\mathcal{S}|}\,
\mathfrak{N}^{\mathbb{C}}\bigl(\mathbf{B}'^{(j)}(A,\mathcal{S})\bigr)
\le e^{\lambda_\theta}c_1^{\flat}(g_j)\,e^{-(1+\frac14\beta_s)\kappa d_j(A)}
\end{equation*}
(the proposition on the boundary terms created by $\mathbf{R}$, clauses (a) and (c)). The strip and
comb lemmas are statements about numbers indexed by $\mathbf{D}_k^{Z}$ and $\mathbf{D}_k^{\sharp}$.
Fed with the sizes $\mathfrak{N}^{A}(\mathfrak{v}(Y_4,\cdot))$ of the pre-filing family, they return
the sizes $\mathfrak{N}^{A}$ of the outputs, because every operation between the pre-filing family
and \cite[(1.99)]{B-CMP122b} is a sum, a product under which strong sets pass to unions, or an
integral $\mathbf{T}'_k(X_j)$ under which the exterior slots are parameters. The weighted sum over
$\mathcal{S}$ is thus delivered, on labels. Concordantly, the weight $e^{\kappa_A|\mathcal{S}|}$ is
carried by the filing of the terms created by $\mathbf{R}$: by clause (b) of the filing proposition
for those terms, $V(Y)$, $\mathbf{R}'^{(k)}(X)$ and $\mathbf{B}'^{(k)}(X)$ are filed under the strong
sets $\mathcal{S}'$ of their units with the weight $e^{\kappa_A(k)|\mathcal{S}|}$ of the complex
induction, termwise, exactly as the outputs of the first-stage bound are filed. These units enter
$V$ through $V^{\mathrm{rel}}$ (clause (c), under the
labelling rule for wrapped units), not through step 2. For their wrapped descendants created by
$\mathbf{T}$ the bound is $B_0e^{-\kappa_{\mathrm{out}}d_j(A)}$, with $A:=X$.

For the arrested units of the fourth kind,
\begin{equation*}
\sum_{\mathcal{S}}e^{\kappa_A(k)|\mathcal{S}|}\,
\mathfrak{N}\bigl(\mathbb{C}^{(i)}_{k_1}(X,\mathcal{S})\bigr)
\le C_0^{\sharp}e^{-a_md(X)},
\end{equation*}
where $\mathbb{C}^{(i)}_{k_1}$, $k_1<k$, are the frozen chain terms, $a_m:=(1-\beta_{\mathrm{pr}})
\kappa$ is the rate at which they are stored in the A-slot norm, and $d(X)$ is the length of $X$ at
their own level. This holds
by the first clause of the proposition on the outputs of the $\mathbf{R}$-operation at the steps
$\le k$, with the norm $\|\cdot\|^{A}$ taken at the weight $\kappa_A$; the change-of-variables
terms are at most $1$ per block, by the bound on change-of-variables terms on arrested components.
On the chart every $\mathfrak{N}$ is read as $\mathfrak{N}^{\mathbb{C}}$.
\end{proof}

\emph{The ceiling $\gamma_V$.} It is defined by the following conditions, each a ceiling on
$\gamma$ chosen after $L$, $M$, $\kappa_1$, $c_A$ and $r$, and each in the supremum form used for
all ceilings on $\gamma$.
\begin{enumerate}
\item $\theta^{\mathbb{C}}_k\le\frac12$. Its member free of $g$ is met by the choice of $c_A$, its
member depending on $\mathfrak{b}$ by $\gamma$. If $\mathsf{w}^{\mathbb{C}}\ge3$ is wanted, one
imposes $\theta^{\mathbb{C}}_k\le\frac14$ instead, since
$\mathsf{w}^{\mathbb{C}}=1/\theta^{\mathbb{C}}_k-1$.
\item $e^{32\kappa_1}C_1g_kT^{\mathbb{C}}\le c_{\mathrm{abs}}$.
\item The quantities written $o_\gamma(1)$ in the proof are in total at most
$\frac12\varepsilon_V^{*}$.
\item The member of step 2:
\begin{equation}\label{5.2:eq-complex-localisation-1}
\begin{split}
\sup_{\mathsf{T}\ge\mathsf{T}_\gamma-1}\ \sum_{s\ge0}
&\Bigl(B_0+2+6\bigl(\mathsf{N}_0(\mathsf{T}+c_s+2)+3\bigr)C_0^{\sharp}+3\lambda s+2c_2\Bigr)\\
&\times\max\bigl\{e^{-\kappa cR(\mathsf{T})/(2L)},\,e^{-\frac12\delta_0cMR(\mathsf{T})}\bigr\}^{s+1}
\le\frac{\varepsilon_V^{*}}8.
\end{split}
\end{equation}
\end{enumerate}
In \eqref{5.2:eq-complex-localisation-1}, $R(\mathsf{T})$, on the side of the decay, is taken at the
lower value $\mathsf{T}$, and $\mathsf{N}_0$, on the side of the costs, at the upper value
$\mathsf{T}+c_s+2$. The series in $s$ converges for every $\mathsf{T}$, since $c_s$ grows
logarithmically. The quantity under the supremum decreases to zero as $\mathsf{T}_\gamma\to\infty$,
so that this is a
ceiling of degree $-\infty$, with no cap on ages, on $K$ or on $k$. The share of these ceilings
needed by the complex-chart corollary is carried, by name, by two of the conditions on the constants
under which that corollary is proved.

\begin{definition}[Exact refinement of the rim cut-offs]\label{5.2:def-rim-refinement}
Consider step $k+1$ of Ba{\l}aban's construction in the notation of \cite[Sect.~3]{B-CMP119}.
Let $\mathfrak{b}=\mathfrak{b}_k$ be the small-field box of Definition~\ref{5.2:def-rim-datum},
and let $\mathfrak{b}^{\flat}$ and $\varpi$ be the constants fixed in
Notation~\ref{5.2:not-rim-constants}.
Let $\rho$ be the bond distance in the current unit lattice, and $\operatorname{dist}_\rho$ the
distance it induces.

\emph{Where the refinement is inserted.} The construction below is inserted between
\cite[(3.19)]{B-CMP119} and \cite[(3.20)]{B-CMP119}. At the first step it is inserted between
\cite[(1.19)]{B-CMP119} and \cite[(1.20)]{B-CMP119}, \cite[pp.~250--251]{B-CMP119}.

\emph{The splitting.} Fix one term of the expansion \cite[(3.16)]{B-CMP119}. Let $\mathcal{P}$ be
the partition into cubes with which the refinement identity named below works, and let $R$ be the
large-field set of the fixed term, a union of cubes of $\mathcal{P}$. For a bond $b$, let $A_k(b)$ be
the value on $b$ of the field on which the characteristic functions of the term act, and let
$\chi_b:=\chi\bigl(|A_k(b)|<\mathfrak{b}\bigr)$ be the small-field characteristic function on $b$.
On the bonds described below we write
\begin{equation}\label{5.2:eq-rim-refinement-1}
\chi_b=\chi^0_b+\chi^1_b,\qquad
\chi^0_b:=\chi\bigl(|A_k(b)|<\mathfrak{b}^{\flat}\bigr),\qquad
\chi^1_b:=\chi\bigl(\mathfrak{b}^{\flat}\le|A_k(b)|<\mathfrak{b}\bigr).
\end{equation}
A bond is called \emph{refined} once its factor $\chi_b$ has been replaced by the right member
of~\eqref{5.2:eq-rim-refinement-1}. For a set $I$ of bonds, $\operatorname{cubes}(I)$ denotes the
set of $\mathcal{P}$-cubes that contain a bond of $I$; a bond lying between two $\mathcal{P}$-cubes
contributes both of them.

\emph{The recursion.} Let $R'\mapsto\Lambda(R')$ be the map that attaches to a large-field set its
domain, and let $\operatorname{conv}$ be the hull operation, both as fixed in the refinement
identity named below. Put $R^0:=R$ and $\Lambda^0:=\Lambda(R)$. Suppose $R^i$ and $\Lambda^i$ are
given. Let $\mathfrak{r}^i$ be the
set of rim bonds $b$ of $\Lambda^i$, that is, bonds of $\mathfrak{X}$ with no end-point in
$\Lambda^i$ (the rim $\mathfrak{r}$ of Definition~\ref{5.2:def-rim-datum} read with $\Lambda^i$ in
place of $\Lambda$), with the following three properties: $b$ is not yet refined;
$b$ lies in no cube of $R^i$; and $\operatorname{dist}_\rho(b,\Lambda^i)\le\varpi$.

Refine every $b\in\mathfrak{r}^i$ and expand the product of the factors
\eqref{5.2:eq-rim-refinement-1}. In each term of this expansion, the set $I^i\subset\mathfrak{r}^i$
of \emph{violators} is the set of bonds carrying $\chi^1_b$. Put
\begin{equation}\label{5.2:eq-rim-refinement-2}
R^{i+1}:=R^i\cup\operatorname{cubes}(I^i),\qquad
\Lambda^{i+1}:=\operatorname{conv}\bigl[\Lambda(R^{i+1})\cap\Lambda^i\bigr].
\end{equation}

The recursion stops at the first stage $i$ with $I^i=\emptyset$. At that stage put $R^*:=R^i$ and
$\Lambda^*:=\Lambda^i$. Let $I^*:=I^0\cup\dots\cup I^{i-1}$ be the set of all violators, and let
$\mathsf{C}^*$ be the set of refined bonds not in $I^*$. The output of the refinement is
\begin{equation*}
(R,I^*)\longmapsto(\Lambda^*,R^*,\mathsf{C}^*).
\end{equation*}

\emph{The refined step.} Put $\Lambda_{k+1}:=\Lambda^*$. The following are taken without change
at this domain:
\begin{itemize}
\item the definitions of $R_{k+1}$ and $S_{k+1}$;
\item the changes of variables \cite[(1.22)]{B-CMP119};
\item $H^{(k)}$ of \cite{B-CMP119};
\item the formulas \cite[(3.22)--(3.25)]{B-CMP119}.
\end{itemize}
The characteristic function of \cite[(3.21)]{B-CMP119} is replaced by
\begin{equation}\label{5.2:eq-rim-refinement-a}
\chi(\Omega\cap\Lambda^c,S,I)
:=\chi^{(k)}(\mathrm{rest})\cdot\chi^0(\mathsf{C}^*)\,\chi^1(I^*)\cdot
\chi^{(k)c}(R)\,\chi'_k(S),
\end{equation}
with $\Lambda=\Lambda^*$ and $I=I^*$, where $\Omega$ denotes the region so written in
\cite[(3.21)]{B-CMP119}. Here $\chi^{(k)}(\mathrm{rest})$ is the factor $\chi^{(k)}$ of
\cite[(3.21)]{B-CMP119} taken over the bonds of $\Omega\cap\Lambda^c$ that lie neither in $R$ nor
in $S$ and were not refined. Further,
$\chi^0(\mathsf{C}^*)$ and $\chi^1(I^*)$ are the products of $\chi^0_b$ over $b\in\mathsf{C}^*$
and of $\chi^1_b$ over $b\in I^*$. The factors $\chi^{(k)c}(R)$ and $\chi'_k(S)$ are those of
\cite[(3.21)]{B-CMP119}, and $S$ is an admissible set.

\emph{Hypotheses.} Assume the following two statements.
\begin{enumerate}
\item The refinement identity: the fixed term of \cite[(3.16)]{B-CMP119} equals the sum of the
refined terms produced by the recursion~\eqref{5.2:eq-rim-refinement-2}, each refined term being
obtained from the fixed term by replacing, on every refined bond $b$, the factor $\chi_b$ by
$\chi^0_b$ or by $\chi^1_b$, as the corresponding term of the expansion of the products
\eqref{5.2:eq-rim-refinement-1} prescribes.
\item The refinement lemma, which asserts the following clauses.
\begin{enumerate}
\item[(a)] The recursion is finite.
\item[(b)] The refinement is exact: for each configuration exactly one term is non-zero.
\item[(c)] The refinement is history-free: its terms are indexed by the pairs $(R,I^*)$.
\item[(d)] Every rim bond of $\Lambda^*$ that lies in no cube of $R^*$ and lies within $\varpi$ of
$\Lambda^*$ carries $\chi^0$.
\item[(d$'$)] No bond of $\bar\Lambda^*$, the set of free bonds with at least one end-point in
$\Lambda^*$ (the set $\bar\Lambda$ of Definition~\ref{5.2:def-rim-datum} read at $\Lambda^*$), is
refined.
\item[(e)] $\Lambda^*$ is a union of cubes of side $\ell(x_{k+1})$, the localisation length at
$x_{k+1}$, contained in $\Lambda(R^*)\cap\Omega^{\sim-1}_{k+1}$, where $\Omega^{\sim-1}_{k+1}$ is
the region $\Omega_{k+1}$ of \cite[Sect.~3]{B-CMP119} with one layer removed. The small components
of the complement of $\Lambda^*$ are
parallelepipeds. The $R^*$-cubes and $(\Omega^{\sim-1}_{k+1})^c$ lie at distance at least
$\mathfrak{d}_{\mathrm{far}}$ from $\Lambda^*$, where $\mathfrak{d}_{\mathrm{far}}$ is the
separation fixed in the refinement lemma, and
\begin{equation*}
\mathfrak{d}_{\mathrm{far}}\ge MR(x_{k+1}).
\end{equation*}
\end{enumerate}
\end{enumerate}

\emph{Conclusion.} Under these hypotheses the following hold.
\begin{enumerate}
\item[(i)] The refinement is covariant under the partition group. Consequently
\cite[(2.32)]{B-CMP119}, clause (f3) of the definition of E-families and the covariance statement
for the labels apply to the labels $(R,I^*)$ of the refined expansion exactly as to the labels of
the unrefined one.
\item[(ii)] The estimate \cite[(3.19)]{B-CMP119} holds in every refined term.
\item[(iii)] The changes of variables \cite[(1.22)]{B-CMP119} act only on the bonds of the set
$S^*$ of \cite[(1.22)]{B-CMP119}.
\item[(iv)] The bound of $H^{(k)}$ on $S$ holds at $\Lambda_{k+1}=\Lambda^*$.
\item[(v)] The formulas \cite[(3.23)--(3.25)]{B-CMP119} hold without change on $\Lambda^*$.
\item[(vi)] Ba{\l}aban's sum over domains $\Lambda_{k+1}$, and for a fixed domain over admissible
sets $S_{k+1}$, becomes a sum over admissible pairs $(S,I)$. This sum is exact and history-free.
\end{enumerate}
\end{definition}

\begin{proof}
We take the conclusions in order.

\emph{Covariance, conclusion~(i).} The refinement is built from five ingredients only:
\begin{itemize}
\item the bond distance $\rho$;
\item the partition $\mathcal{P}$;
\item the small-field box $\mathfrak{b}$;
\item the single threshold $\mathfrak{b}^{\flat}$ in~\eqref{5.2:eq-rim-refinement-1};
\item the rule that a bond lying between two $\mathcal{P}$-cubes contributes both cubes to
$\operatorname{cubes}(\cdot)$.
\end{itemize}
Each ingredient is covariant under the partition group, and hence so is the map
$(R,I^*)\mapsto(\Lambda^*,R^*,\mathsf{C}^*)$. Therefore \cite[(2.32)]{B-CMP119}, clause (f3) of the
definition of E-families and the covariance statement for the labels apply to the labels $(R,I^*)$
exactly as to the labels of the unrefined expansion.

\emph{Conclusion~(ii).} The estimate \cite[(3.19)]{B-CMP119} uses, on the bonds concerned, only
the condition $|A_k|<\delta_k$, where $\delta_k$ is the threshold printed in that estimate. This
condition holds under $\chi^0_b$ and under $\chi^1_b$. Indeed,
by~\eqref{5.2:eq-rim-refinement-1} these two functions are non-negative, have disjoint supports
and sum to $\chi_b$. Hence each of them is supported in the support of $\chi_b$, and there the
condition is in force. So \cite[(3.19)]{B-CMP119} remains valid term by term after the splitting.

\emph{Conclusion~(iii).} The changes of variables \cite[(1.22)]{B-CMP119} act only on the bonds of
the set $S^*$ of \cite[(1.22)]{B-CMP119}. At this point \cite{B-CMP119} states: ``The other
characteristic functions are unchanged.''

\emph{Conclusion~(iv).} The bound of $H^{(k)}$ on $S$ requires that $\operatorname{dist}(S,\Lambda)$
be at least one layer. For Ba{\l}aban's domain $\Lambda(R)$ this separation holds. Since
$\Lambda^*\subset\Lambda(R)$, we have
\begin{equation*}
\operatorname{dist}(S,\Lambda^*)\ge\operatorname{dist}\bigl(S,\Lambda(R)\bigr).
\end{equation*}
So the separation is inherited by $\Lambda_{k+1}=\Lambda^*$, and the bound of $H^{(k)}$ holds there.

\emph{Conclusion~(v).} The formulas \cite[(3.23)--(3.25)]{B-CMP119} are taken without change with
$\Lambda_{k+1}:=\Lambda^*$.

\emph{Conclusion~(vi).} In \cite{B-CMP119} the expansion after the changes of variables is written
as ``a sum over domains $\Lambda_{k+1}$, and for a fixed domain $\Lambda_{k+1}$, \dots\ over the
admissible sets $S_{k+1}$''. After the refinement, the domain is determined by the output of the
recursion, $\Lambda_{k+1}=\Lambda^*$, and the characteristic function
is~\eqref{5.2:eq-rim-refinement-a}.
Hence this sum becomes a sum over admissible pairs $(S,I)$ with $I=I^*$.

By clause~(a) of the refinement lemma the recursion is finite, so the fixed term of
\cite[(3.16)]{B-CMP119} is replaced by a finite sum of refined terms. By clause~(b) the refinement
is exact, with
exactly one non-zero term for each configuration. So the refined sum reproduces the unrefined term
of the refinement identity. By clause~(c) the terms are indexed by $(R,I^*)$ alone. Hence the sum
over admissible $(S,I)$ is history-free.
\end{proof}

\begin{lemma}[Factorisation of the refined step over components]
\label{5.2:lem-refined-factorisation}
Let $X_1,\dots,X_n$ be the connected components of $(\Lambda^{*})^c$. The pair $(R,I^{*})$ is
realisable with output $\Lambda^{*}$, in the sense of Definition~\ref{5.2:def-rim-refinement},
if and only if, for each $l=1,\dots,n$, the pair
$(R\cap X_l,\,I^{*}\cap X_l)$, run on the seeds lying in $X_l$, is realisable with output
$X_l^c$. The factors of \eqref{5.2:eq-rim-refinement-a} are products over bonds and cubes.
Hence the sum over admissible $(S,I)$ at fixed $\Lambda^{*}$ factorises over $X_1,\dots,X_n$,
and the relations \cite[(2.19)--(2.20)]{B-CMP119} hold for $\mathbf{T}^{(k+1)}$.
\end{lemma}

\begin{proof}
The refinement of Definition~\ref{5.2:def-rim-refinement} produces sets
$\Lambda^0\supset\Lambda^1\supset\cdots$, ending at $\Lambda^{*}$. We index its stages by $i$
and the components of $(\Lambda^{*})^c$ by $l$, and we call $(\Lambda^i)^c$ the stage set at
stage $i$.

First, the complements $(\Lambda^i)^c$ increase with $i$ and are contained in
$(\Lambda^{*})^c$. A connected component $K$ of some $(\Lambda^i)^c$ is therefore a connected
subset of $(\Lambda^{*})^c$. It consequently lies in exactly one component $X_l$.

Second, consider the operations performed at a stage of Definition~\ref{5.2:def-rim-refinement}:
\begin{itemize}
\item adding a layer of big cubes around the seeds;
\item replacing a component $K$ by its bounding box inside a $100$-cube; this box is
connected, contains $K$, and lies inside the next stage set;
\item forming the collar of width $\varpi$, where $\varpi$ is smaller than one big cube;
\item adjoining the cubes of the violators, $\operatorname{cubes}(I^i)\subset K$.
\end{itemize}
Each of these operations maps subsets of $X_l$ into $X_l$. Each reads nothing beyond the
one-cube neighbourhood of the component $K$ on which it acts. Moreover, two distinct
components $X_l$ and $X_{l'}$ are not adjacent. Consequently, what a stage does inside $X_l$
is determined by the data
$R\cap X_l$, $I^{*}\cap X_l$ and by the seeds in $X_l$, and by nothing in any other
$X_{l'}$.

Third, we argue by induction on $i$. The induction hypothesis is that the stage set of the
joint run at stage $i$ is the disjoint union, over $l$, of the stage sets of the separate runs
in the $X_l$. By the second step, the operations of stage $i+1$ act on each $X_l$ separately
and exactly as in the separate run. If the separate run in some $X_l$ has already stopped,
then that component has no unrefined collar bond. It is therefore idle in the joint run: the
joint run changes nothing in it, just as the stopped separate run does. Hence the stage sets
at stage $i+1$ are again disjoint unions, and the induction closes.

It follows that the joint run stops, that is $I^i=\emptyset$, exactly when every separate run
has stopped. The final stage set of the joint run is the disjoint union of the final stage
sets of the separate runs. Therefore the output of the joint
run is $\Lambda^{*}$, with $(\Lambda^{*})^c=X_1\sqcup\dots\sqcup X_n$, if and only if the
run in each $X_l$ has output $X_l^c$, that is, has final stage set $X_l$. This is the
asserted equivalence.

By its definition, the factor in \eqref{5.2:eq-rim-refinement-a} is a product of
characteristic functions each of which depends on a single bond or on a single cube. Combined
with the equivalence, this shows that the sum over admissible $(S,I)$ at fixed $\Lambda^{*}$
factorises over $X_1,\dots,X_n$. This is the form in which
\cite[(2.19)--(2.20)]{B-CMP119} hold for $\mathbf{T}^{(k+1)}$.
\end{proof}

\begin{definition}[One-sided conditions on the rim constants]\label{5.2:not-rim-conditions}
Except for $\bar\kappa$, $\mathfrak{d}_{\mathrm{far}}$, $C_{\mathrm{rem}}$, $B_0$, $C_{\mathrm{F}1}$ and
$c_{\mathrm{F}1}$, which are explained further on, the symbols below are the constants of the rim
construction of Notation~\ref{5.2:not-rim-constants}, chosen in the order
\eqref{5.2:eq-rim-constants-a}.
The functions of the coupling
among them ($\mathsf{p}$, $\varrho$, $r_A$, $\varpi_\gamma$) and the box parameter
$\mathfrak{b}=\mathfrak{b}_k$ are read at level $k$, that is, at $g_k$. The rim construction uses
the following six conditions:
\begin{equation}\label{5.2:eq-rim-conditions-a}
\begin{aligned}
&(\mathrm{F}_{\varpi}) && M \ge \varpi,\\
&(\mathrm{F}1) && \varepsilon_T \le 1,\qquad \varepsilon_T\, C_\rho(\delta_T) \le \tfrac12\, c_\eta,\\
&(\mathrm{F}_{\mathrm{far}}) && K_1\, e^{\delta_1/2}\, e^{-\delta_1 M R_{k+1}/2}\, C_{\mathrm{far}}
 \le \tfrac{1}{16},\\
&(\mathrm{F}2) && \mathfrak{b} \ge 16,\qquad
 \log\bigl(2(1+c_\eta \mathfrak{b})\bigr) \le \mathsf{p},\qquad
 r_A \ge \max\{2, C_{\mathrm{far}}\},\\
& && \varpi_\gamma + \operatorname{diam}_\rho(\pi_{k+1}) + 1 \le \mathfrak{d}_{\mathrm{far}},\\
&(\mathrm{F}3) && \mathsf{p} \ge 5(\bar\kappa + 1)\, N_Z,\\
&(\mathrm{F}_{\mathrm{rem}}) && C_{\mathrm{rem}}\, \varrho^2\, e^{-\delta_1 M \check{R}_{k+1}/2}
 \le \tfrac14\, B_0 .
\end{aligned}
\end{equation}
In $(\mathrm{F}_{\mathrm{far}})$, $R_{k+1}$ denotes $R(x_{k+1})$, where $R(x)$ is the least power
of $L$ that is at least $(\log x)^{r_{\mathrm{pr}}}$.
In $(\mathrm{F}2)$, $\operatorname{diam}_\rho(\pi_{k+1})$ denotes the $\rho$-diameter of a
$\pi_{k+1}$-cube. It equals $\sqrt{d}\,LM$ when $\rho$ is the Euclidean distance and $dLM$ when
$\rho$ is the $\ell^1$ bond distance. Also in $(\mathrm{F}2)$, $\mathfrak{d}_{\mathrm{far}}$ is the
far distance of the far bound among the admitted inputs
(Definition~\ref{5.1:def-admitted-inputs}).
Condition $(\mathrm{F}_{\mathrm{rem}})$ is imposed for the bound of the rim construction on the
remnants between distinct components; $C_{\mathrm{rem}}$ and $B_0$ are constants appearing in
that bound. In $(\mathrm{F}3)$,
$\bar\kappa$ is the largest decay rate occurring in the construction. Here $\bar\kappa \le \tfrac12 L\kappa$:
the rate in question is $\rho_7=(1-7\delta)\tfrac12 L\kappa$ with $\delta \le \delta_J(L)$.
Condition $(\mathrm{F}3)$ also comprises the conditions \cite[(2.31), (2.33)--(2.34)]{B-CMP116}
with the quantity $\tfrac12\gamma_2\varepsilon_1^2/g_k^2$ replaced by $\mathsf{p}$.

All six conditions are one-sided. The floors $(\mathrm{F}_{\varpi})$ and $(\mathrm{F}1)$, the latter
through \eqref{5.2:eq-rim-conditions-2}, are entered among the one-sided conditions on the
parameters chosen after $L$ (Definition~\ref{5.1:def-admitted-inputs}), on whose joint
satisfiability the existence of $L_0=L_0(N)$ rests; the ceilings $(\mathrm{F}_{\mathrm{far}})$,
$(\mathrm{F}2)$, $(\mathrm{F}3)$ and $(\mathrm{F}_{\mathrm{rem}})$ are entered among the conditions
fixing $\gamma_J$.

Condition $(\mathrm{F}_{\varpi})$ is a floor on $M$. $(\mathrm{F}1)$ is secured by a floor on $M$
imposed after $\kappa_1$ and a ceiling on $\alpha_0+\alpha_1$ imposed after $M$, see
\eqref{5.2:eq-rim-conditions-2}. The
constant $C_{\mathrm{far}}$ does not occur in $(\mathrm{F}1)$.
Let $C_{\mathrm{F}1}(d,L) \ge 0$ and
$c_{\mathrm{F}1}(d,L) > 0$ be numbers satisfying the following inequality; they are auxiliary
$(d,L)$-numbers used only for $(\mathrm{F}1)$. With $B$, $c_0'$ and $\mathsf{c}_2$ the constants of
Notation~\ref{5.2:not-rim-constants} entering $B_W$, $c_\eta$ and $\gamma_2$, we require
\begin{equation}\label{5.2:eq-rim-conditions-1}
2B^2 C_\rho(\delta_T)\bigl(e^{-C_{\mathrm{F}1}(d,L)} + c_{\mathrm{F}1}(d,L)\bigr)
\le \min\Bigl\{1,\ \frac{1}{2C_\rho(\delta_T)}
\Bigl[\frac{\mathsf{c}_2}{16(c_0'+1)}\Bigr]^{1/2}\Bigr\}.
\end{equation}
Any pair of $(d,L)$-numbers with this property will do; such a pair exists: the right member is a
positive $(d,L)$-number, and the left member tends to zero as $C_{\mathrm{F}1}(d,L)\to\infty$ and
$c_{\mathrm{F}1}(d,L)\to 0$.
Then $(\mathrm{F}1)$ holds provided
\begin{equation}\label{5.2:eq-rim-conditions-2}
\frac{\delta_W M}{4} \ge 32\kappa_1 + 2\log(LM) + C_{\mathrm{F}1}(d,L),
\qquad
\alpha_0 + \alpha_1 \le c_{\mathrm{F}1}(d,L)\, e^{-32\kappa_1} (LM)^{-2}.
\end{equation}

Conditions $(\mathrm{F}_{\mathrm{far}})$, $(\mathrm{F}2)$, $(\mathrm{F}3)$ and
$(\mathrm{F}_{\mathrm{rem}})$ are ceilings on $\gamma$: each is imposed at the levels $k$ with
$g_i \le \gamma$ for all $i \le k+1$, at which
\begin{equation*}
R_{k+1} \ge \check{R}_{k+1} \ge \bigl(\log\gamma^{-2}\bigr)^{r_{\mathrm{pr}}},
\end{equation*}
and the choice of $\gamma_J$ is required to make it hold there. For $(\mathrm{F}_{\mathrm{far}})$,
the first three inequalities of $(\mathrm{F}2)$ and the displayed inequality of $(\mathrm{F}3)$,
this requirement is met by every $\gamma$ below a threshold fixed by the constants chosen
before $\gamma$.
\end{definition}

\begin{proof}
We first show that \eqref{5.2:eq-rim-conditions-2} implies $(\mathrm{F}1)$. By the list of
constants of the rim construction (Notation~\ref{5.2:not-rim-constants}) we have
\begin{equation*}
\varepsilon_T = 2B_W^2\, C_\rho(\delta_T)\bigl(e^{-\delta_W M/4} + \alpha_0 + \alpha_1\bigr),
\qquad B_W = B e^{16\kappa_1}.
\end{equation*}
We also have $c_\eta = [\gamma_2/(16(c_0'+1))]^{1/2}$ with $\gamma_2 = \mathsf{c}_2 (LM)^{-4}$, so
\begin{equation*}
\tfrac12\, c_\eta = \tfrac12 \Bigl[\frac{\mathsf{c}_2}{16(c_0'+1)}\Bigr]^{1/2} (LM)^{-2}.
\end{equation*}
The first inequality of \eqref{5.2:eq-rim-conditions-2} gives
\begin{equation*}
e^{-\delta_W M/4} \le e^{-32\kappa_1}(LM)^{-2} e^{-C_{\mathrm{F}1}(d,L)} .
\end{equation*}
Adding the second inequality of \eqref{5.2:eq-rim-conditions-2}, we get
\begin{equation*}
e^{-\delta_W M/4} + \alpha_0 + \alpha_1
\le e^{-32\kappa_1}(LM)^{-2}\bigl(e^{-C_{\mathrm{F}1}(d,L)} + c_{\mathrm{F}1}(d,L)\bigr).
\end{equation*}
Since $B_W^2 = B^2 e^{32\kappa_1}$, the factor $e^{32\kappa_1}$ cancels, and we obtain
\begin{equation*}
\varepsilon_T \le 2B^2 C_\rho(\delta_T)\bigl(e^{-C_{\mathrm{F}1}(d,L)} + c_{\mathrm{F}1}(d,L)\bigr)\,
(LM)^{-2}.
\end{equation*}
By \eqref{5.2:eq-rim-conditions-1}, the factor
$2B^2 C_\rho(\delta_T)\bigl(e^{-C_{\mathrm{F}1}(d,L)} + c_{\mathrm{F}1}(d,L)\bigr)$
is at most $1$. Since $L \ge 1$ and $M \ge 1$, this gives
$\varepsilon_T \le 1$. Again by \eqref{5.2:eq-rim-conditions-1},
\begin{equation*}
\begin{aligned}
\varepsilon_T\, C_\rho(\delta_T)
&\le 2B^2 C_\rho(\delta_T)\bigl(e^{-C_{\mathrm{F}1}(d,L)} + c_{\mathrm{F}1}(d,L)\bigr)\,
C_\rho(\delta_T)\,(LM)^{-2}\\
&\le \tfrac12\Bigl[\frac{\mathsf{c}_2}{16(c_0'+1)}\Bigr]^{1/2}(LM)^{-2} = \tfrac12\, c_\eta ,
\end{aligned}
\end{equation*}
which is the second inequality of $(\mathrm{F}1)$.

The first inequality of \eqref{5.2:eq-rim-conditions-2} is a lower bound on $M$ once $L$ and
$\kappa_1$ are fixed. The second is an upper bound on $\alpha_0+\alpha_1$, imposed when the
$\alpha$'s are chosen after $M$. Neither inequality, nor \eqref{5.2:eq-rim-conditions-1},
contains $C_{\mathrm{far}}$. Condition $(\mathrm{F}_{\varpi})$ is a lower bound on $M$ in terms of
$\varpi$, which is fixed among the $(d,L)$-numbers. This establishes the assertions about
$(\mathrm{F}_{\varpi})$ and $(\mathrm{F}1)$.

Next consider $(\mathrm{F}_{\mathrm{far}})$, $(\mathrm{F}2)$, $(\mathrm{F}3)$ and
$(\mathrm{F}_{\mathrm{rem}})$. These involve the coupling through $\mathsf{T}_k$, via
$\mathfrak{b}$, $\mathsf{p}$, $r_A$, $\varpi_\gamma$, $N_Z$ and $\varrho$, and through the radii
$R_{k+1}$ and $\check{R}_{k+1}$; and $\mathsf{T}_k \ge \mathsf{T}_\gamma
= \log\gamma^{-2}$ when $g_k \le \gamma$.

Recall that $R(x)$ is the least power of $L$ that is at least $(\log x)^{r_{\mathrm{pr}}}$; thus
$x \mapsto R(x)$ is non-decreasing and $R(x) \ge (\log x)^{r_{\mathrm{pr}}}$. Recall also
that $\check{R}_{k+1} = R(\min_{i\le k+1} x_i)$. Whenever $g_i \le \gamma$ for all $i \le k+1$,
we have $\min_{i\le k+1} x_i \ge \gamma^{-2}$, and therefore
\begin{equation*}
\check{R}_{k+1} \ge \bigl(\log \min_{i\le k+1} x_i\bigr)^{r_{\mathrm{pr}}}
\ge \bigl(\log \gamma^{-2}\bigr)^{r_{\mathrm{pr}}}.
\end{equation*}
Since $x_{k+1} \ge \min_{i\le k+1} x_i$ and $R$ is non-decreasing, also
$R_{k+1} = R(x_{k+1}) \ge \check{R}_{k+1}$.

In $(\mathrm{F}_{\mathrm{far}})$ the coupling enters only through $R_{k+1}$, and the left member
decreases as $R_{k+1}$ increases. By the two bounds just proved, when $g_i \le \gamma$ for all
$i \le k+1$ the left member is at most
\begin{equation*}
K_1\, e^{\delta_1/2}\, C_{\mathrm{far}}\,
e^{-\delta_1 M (\log\gamma^{-2})^{r_{\mathrm{pr}}}/2},
\end{equation*}
and this is at most $\tfrac1{16}$ once $\gamma$ is small enough, since $r_{\mathrm{pr}}>0$ and
$K_1$, $\delta_1$, $M$ and $C_{\mathrm{far}}$ are fixed before $\gamma$.

By Notation~\ref{5.2:not-rim-constants}, at level $k$ we have
$\mathsf{p} = \gamma_2\mathfrak{b}^2/64$ with $\mathfrak{b} = A_1\mathsf{T}_k^{p_0}$, and
$r_A = c_A (C_*/A_1)\mathsf{T}_k^{q-p_0}$ with $p_0>0$ and $q > p_0$; thus $\mathfrak{b}$,
$\mathsf{p}$ and
$r_A$ increase without bound as $\mathsf{T}_k$ increases, while $c_\eta$, $\gamma_2$ and
$C_{\mathrm{far}}$ are fixed before $\gamma$. Hence $\mathfrak{b} \ge 16$ and
$r_A \ge \max\{2, C_{\mathrm{far}}\}$ hold as soon as $\mathsf{T}_k$ is large enough, and so does
$\log(2(1+c_\eta\mathfrak{b})) \le \mathsf{p}$, whose left member grows like $\log\mathfrak{b}$ and
whose right member like $\mathfrak{b}^2$. By the bounds on $\varpi_\gamma$ and on $N_Z$ stated
in Notation~\ref{5.2:not-rim-constants}, $\varpi_\gamma$ exceeds $\varpi + 1$ by at most a
constant plus a constant multiple of $\log r_A$, and $N_Z$ is at most a constant times
$(L + \varpi_\gamma/M)^d$, all constants being fixed before $\gamma$; so for $\mathsf{T}_k \ge 1$
the number $N_Z$ is at most a constant fixed before $\gamma$ times $(1+\log\mathsf{T}_k)^d$. Since
$\bar\kappa \le \tfrac12 L\kappa$ and $\mathsf{p} = \gamma_2 A_1^2\mathsf{T}_k^{2p_0}/64$, the
displayed inequality of $(\mathrm{F}3)$ holds as soon as $\mathsf{T}_k$ is large enough. In each of
these cases the inequality holds for every $\mathsf{T}_k$ above a threshold fixed before $\gamma$,
hence at every level with $g_k \le \gamma$ once $\gamma$ is small enough: it is a ceiling on
$\gamma$.

The remaining three conditions, namely the conditions \cite[(2.31), (2.33)--(2.34)]{B-CMP116}
comprised in $(\mathrm{F}3)$, condition $(\mathrm{F}_{\mathrm{rem}})$ and the last inequality of
$(\mathrm{F}2)$, are imposed at the levels with $g_i \le \gamma$ for all $i \le k+1$, and the
choice of $\gamma_J$ is required to make them hold there; no threshold on $\gamma$ for them is
asserted in the statement. The coupling enters them as follows. In the conditions of
\cite[(2.31), (2.33)--(2.34)]{B-CMP116} the quantity $\tfrac12\gamma_2\varepsilon_1^2/g_k^2$ is
replaced by $\mathsf{p}$, a function of $\mathsf{T}_k$. In $(\mathrm{F}_{\mathrm{rem}})$, by
Notation~\ref{5.2:not-rim-constants},
$\varrho^2 = (c_A C_*)^2\mathsf{T}_k^{2q}$, read at $g_k$, increases with $\mathsf{T}_k$; the
factor $e^{-\delta_1 M\check{R}_{k+1}/2}$ decreases as $\check{R}_{k+1}$ increases and, by the
bound on $\check{R}_{k+1}$ proved above, is at most
$e^{-\delta_1 M(\log\gamma^{-2})^{r_{\mathrm{pr}}}/2}$ whenever $g_i\le\gamma$ for all
$i\le k+1$. In the last inequality of $(\mathrm{F}2)$ the left member increases with
$\mathsf{T}_k$ through $\varpi_\gamma$, which exceeds $\varpi+1$ by at most a constant plus a
constant multiple of $\log r_A$, as recalled above; the right member is the far distance
$\mathfrak{d}_{\mathrm{far}}$ of the far bound.
\end{proof}

\begin{lemma}[Conditional mean of the rim datum]\label{5.2:lem-conditional-mean}
Put
\begin{equation*}
T(\sigma) := -C_\Lambda(\sigma)\,\mathcal{A}_{\Lambda\mathfrak{r}}(\sigma),
\qquad
T_0 := -C_{\Lambda,0}\,\mathcal{A}_0 ,
\end{equation*}
where $T_0$ is the reference operator, at $\sigma \equiv 0$, of the random-walk bounds for the
conditioned covariances (Lemma~\ref{5.2:lem-random-walk-bounds}). The operator $T_0$ is real, and it
exists on the local spaces of \cite{B-CMP109} and \cite{B-CMP116}.
\begin{enumerate}
\item[(i)] For $b \in Z_0\cap\bar\Lambda$ and $b' \in \mathfrak{r}$,
\begin{equation*}
\bigl|[T(\sigma;\mathbf{U},\mathbf{J}) - T_0](b,b')\bigr|
\le \varepsilon_T\, e^{-\delta_T\rho(b,b')} .
\end{equation*}
Consequently, for each $b \in Z_0\cap\bar\Lambda$, the sum of $|T(\sigma)(b,b')|$ over
$b' \in \mathfrak{r}$ is at most $\Theta_{\flat}$, and for every $w \ge 0$ the part of this sum
over those $b'$ with $\rho(b,b') \ge w$ is at most $\tau_{\flat}(w)$.
\item[(ii)] Assume conditions $(\mathrm{F}1)$ and $(\mathrm{F}2)$ of \eqref{5.2:eq-rim-conditions-a}. Let
$P \subset Z_0\cap\bar\Lambda$. Let $x$ be a rim datum such that, on the rim bonds of the cubes of
$\mathcal{S}(P)$, $x$ is real and lies in $\mathfrak{X}^{\sharp}$; on the remaining bonds $x$ is
complex with $|x| < \varrho$, except that it is real on the silent bonds of the rim datum
\eqref{5.2:eq-rim-datum-a}. Then
$\mathsf{m} := (T(\sigma)x)\restriction Z_0$ lies in
\begin{equation}\label{5.2:eq-conditional-mean-a}
\begin{split}
\mathfrak{M}_P := \Bigl\{\mathsf{m} :\ & |\mathsf{m}_b| \le \Theta_{\flat}\varrho
\ \text{for all } b;\\
& |\operatorname{Re}\mathsf{m}_b| \le \tfrac{7}{16}\,\mathfrak{b},\
|\operatorname{Im}\mathsf{m}_b| \le c_\eta\,\mathfrak{b}\ \text{for } b \in P\Bigr\}.
\end{split}
\end{equation}
\item[(iii)] Suppose $\sigma = 0$ off $Z$. Then, for $b$ in a connected component $Z'$ of $Z$,
$\mathsf{m}_b$ depends on $(x,\mathbf{U},\mathbf{J})$ only through their restrictions to $Z'$.
In all cases $T(\mathbb{1})x = m(x)$.
\end{enumerate}
\end{lemma}

\begin{proof}
\emph{(i).} We write
\begin{equation*}
T(\sigma) - T_0
= -\bigl[C_\Lambda(\sigma) - C_{\Lambda,0}\bigr]\mathcal{A}(\sigma)
- C_{\Lambda,0}\bigl[\mathcal{A}(\sigma) - \mathcal{A}_0\bigr],
\end{equation*}
and we bound the two products separately at $(b,b')$, summing over the intermediate bond $b''$.

The family $\sigma_0$ of cubes entering clause (c) of Lemma~\ref{5.2:lem-random-walk-bounds} does
not meet $\tilde Z_0$. Hence $\rho(b,\cup\sigma_0) \ge M$, and so for every
$b''$
\begin{equation*}
\mathsf{d}(b,b'') \ge \max\{M,\rho(b,b'')\} \ge \tfrac12 M + \tfrac12\rho(b,b'').
\end{equation*}

\emph{First product.} Apply clause (c) of Lemma~\ref{5.2:lem-random-walk-bounds} to
$C_\Lambda(\sigma) - C_{\Lambda,0}$ at $(b,b'')$, and clause (b) to $\mathcal{A}(\sigma)$ at
$(b'',b')$. Clause (c) is used in its unshifted form, with the factor $e^{-\delta_W\mathsf{d}}$.
That form applies to the pair $(b,b'')$ because
\begin{equation*}
\mathsf{d}(b,b'') \ge M \ge \tfrac34 M.
\end{equation*}
Combining these with $\mathsf{d}(b,b'') \ge \tfrac12 M + \tfrac12\rho(b,b'')$ gives
\begin{equation*}
\Bigl|\bigl([C_\Lambda(\sigma) - C_{\Lambda,0}]\mathcal{A}(\sigma)\bigr)(b,b')\Bigr|
\le B_W^2\bigl(e^{-\delta_W M/2} + \alpha_0 + \alpha_1\bigr)
\sum_{b''} e^{-\frac12\delta_W\rho(b,b'') - \delta_W\rho(b'',b')} .
\end{equation*}
By the triangle inequality $\rho(b,b'') + \rho(b'',b') \ge \rho(b,b')$, we have
\begin{equation*}
\tfrac12\rho(b,b'') + \rho(b'',b')
\ge \tfrac14\rho(b,b'') + \tfrac14\bigl[\rho(b,b'') + \rho(b'',b')\bigr]
\ge \tfrac14\rho(b,b'') + \tfrac14\rho(b,b').
\end{equation*}
Summing $e^{-\frac14\delta_W\rho(b,b'')}$ over $b''$ gives at most $C_\rho(\delta_W/4)$. The first
product is therefore at most
\begin{equation*}
B_W^2\, C_\rho(\delta_W/4)\bigl(e^{-\delta_W M/2} + \alpha_0 + \alpha_1\bigr)
e^{-\frac14\delta_W\rho(b,b')} .
\end{equation*}

\emph{Second product.} Apply clause (a) of Lemma~\ref{5.2:lem-random-walk-bounds} to
$C_{\Lambda,0}$ at $(b,b'')$. Apply clause (c), in its shifted form, to
$\mathcal{A}(\sigma) - \mathcal{A}_0$ at the general pair $(b'',b')$. The shifted form carries the
factor $e^{-\delta_W[\mathsf{d}-M/2]_+}$ and holds at every pair, including pairs with
$\mathsf{d}(b'',b') < \tfrac34 M$.

In the non-$\alpha$ part, the exponential factors multiply to
\begin{equation*}
e^{-\delta_W\rho(b,b'')}\, e^{-\delta_W[\mathsf{d}(b'',b') - M/2]_+}\, e^{-\delta_W\rho(b'',b')}
= e^{-\delta_W E},
\end{equation*}
where
\begin{equation*}
E := \rho(b,b'') + [\mathsf{d}(b'',b') - M/2]_+ + \rho(b'',b').
\end{equation*}
We use two triangle inequalities:
\begin{equation*}
\rho(b,b'') + \mathsf{d}(b'',b') \ge \mathsf{d}(b,b') \ge M,
\qquad
\rho(b,b'') + \rho(b'',b') \ge \rho(b,b').
\end{equation*}
The first gives $\rho(b,b'') + \mathsf{d}(b'',b') - M/2 \ge M/2$. Since
$[\,\cdot\,]_+ \ge \cdot$, it follows that
\begin{equation*}
\rho(b,b'') + [\mathsf{d}(b'',b') - M/2]_+ \ge M/2 .
\end{equation*}
Now split $E$ as
\begin{align*}
E ={}& \bigl[\tfrac14\rho(b,b'') + \tfrac14\rho(b'',b')\bigr]
+ \tfrac12\bigl[\rho(b,b'') + [\mathsf{d}(b'',b') - M/2]_+\bigr]\\
&+ \bigl[\tfrac14\rho(b,b'') + \tfrac12[\mathsf{d}(b'',b') - M/2]_+
+ \tfrac34\rho(b'',b')\bigr].
\end{align*}
The coefficients of $\rho(b,b'')$, of $[\mathsf{d}(b'',b') - M/2]_+$ and of $\rho(b'',b')$ add up
to one in each case. We bound the three brackets in turn:
\begin{itemize}
\item the first bracket is at least $\tfrac14\rho(b,b')$, by the second triangle inequality;
\item the second is at least $\tfrac14 M$, by the inequality just derived;
\item the third is at least $\tfrac14\rho(b,b'')$, since its other terms are non-negative.
\end{itemize}
Hence
\begin{equation*}
E \ge \tfrac14\rho(b,b') + \tfrac14 M + \tfrac14\rho(b,b''),
\end{equation*}
that is,
\begin{equation*}
e^{-\delta_W\rho(b,b'')}\, e^{-\delta_W[\mathsf{d}(b'',b') - M/2]_+}\, e^{-\delta_W\rho(b'',b')}
\le e^{-\frac14\delta_W\rho(b,b'')}\, e^{-\frac14\delta_W[M + \rho(b,b')]} .
\end{equation*}
In the part proportional to $\alpha_0 + \alpha_1$ the factors are $e^{-\delta_W\rho(b,b'')}$ and
$e^{-\delta_W\rho(b'',b')}$. Here we use
\begin{equation*}
\rho(b,b'') + \rho(b'',b') \ge \tfrac12\rho(b,b'') + \tfrac12\rho(b,b'),
\end{equation*}
which follows from the second triangle inequality. This part therefore carries at most
\begin{equation*}
e^{-\frac12\delta_W\rho(b,b'')}\,e^{-\frac12\delta_W\rho(b,b')}
\le e^{-\frac14\delta_W\rho(b,b'')}\,e^{-\frac14\delta_W\rho(b,b')} .
\end{equation*}
Summing over $b''$ as before, the
second product is at most
\begin{equation*}
B B_W\, C_\rho(\delta_W/4)\bigl(e^{-\delta_W M/4} + \alpha_0 + \alpha_1\bigr)
e^{-\frac14\delta_W\rho(b,b')} .
\end{equation*}

\emph{Sum of the two products.} We use three facts:
\begin{itemize}
\item $\delta_T = \delta_W/4$;
\item $B \le B_W$, since $B_W = Be^{16\kappa_1}$;
\item $e^{-\delta_W M/2} \le e^{-\delta_W M/4}$.
\end{itemize}
Adding the two products then gives
\begin{equation*}
\bigl|[T(\sigma) - T_0](b,b')\bigr|
\le 2B_W^2\, C_\rho(\delta_T)\bigl(e^{-\delta_W M/4} + \alpha_0 + \alpha_1\bigr)
e^{-\delta_T\rho(b,b')}
= \varepsilon_T\, e^{-\delta_T\rho(b,b')} .
\end{equation*}
Here $\varepsilon_T$ is exactly the constant so named in \eqref{5.2:eq-rim-constants-a}.

\emph{Row sums and tails.} By the choice of the pair $(K_0,\delta_1)$ in
\eqref{5.2:eq-rim-constants-a}, we have $|T_0(b,b')| \le K_0 e^{-\delta_1\rho(b,b')}$. Hence
\begin{equation*}
|T(\sigma)(b,b')| \le K_0 e^{-\delta_1\rho(b,b')} + \varepsilon_T e^{-\delta_T\rho(b,b')} .
\end{equation*}
Recall from \eqref{5.2:eq-rim-constants-a} that $\delta_{\flat} = \min\{\delta_1,\delta_T\}$ and
$\Theta_{\flat} = (K_0+1)C_\rho(\delta_{\flat})$. Since $C_\rho(\delta')$ does not increase as
$\delta'$ increases, $C_\rho(\delta_1)$ and $C_\rho(\delta_T)$ are at most $C_\rho(\delta_\flat)$,
and summing over $b'$ bounds the rows at $b \in Z_0\cap\bar\Lambda$ by
$(K_0+\varepsilon_T)C_\rho(\delta_\flat)$. This is at most $\Theta_{\flat}$ when
$\varepsilon_T \le 1$, which is the first condition of $(\mathrm{F}1)$ in
\eqref{5.2:eq-rim-conditions-a}.

For the tails, bound both terms by $(K_0+\varepsilon_T)e^{-\delta_{\flat}\rho(b,b')}$ and write
$e^{-\delta_{\flat}\rho} = e^{-\delta_{\flat}\rho/2}e^{-\delta_{\flat}\rho/2}$.
On $\rho(b,b') \ge w$ the first factor is at most $e^{-\delta_{\flat}w/2}$, and the sum of the
second over $b'$ is at most $C_\rho(\delta_{\flat}/2)$. Since $e^{\delta_{\flat}} \ge 1$ and
$K_0+\varepsilon_T \le K_0+1$ under the same condition, the
tails beyond $w$ are at most $\tau_{\flat}(w)$.

\emph{(ii), first clause.} By (i), the row of $T(\sigma)$ at $b$ has sum at most
$\Theta_{\flat}$. We claim that $|x| \le \varrho$ on every bond.
\begin{itemize}
\item Off the cubes of $\mathcal{S}(P)$ this is the hypothesis.
\item On those cubes $x$ lies in $\mathfrak{X}^{\sharp}$ of \eqref{5.2:eq-rim-datum-a}, with the
bound $C_{\mathrm{far}}\mathfrak{b}$. This is at most $\varrho$, because by $(\mathrm{F}2)$
\begin{equation*}
r_A = \varrho/\mathfrak{b} \ge C_{\mathrm{far}} .
\end{equation*}
\end{itemize}
The row sum times $\sup|x|$ gives $|\mathsf{m}_b| \le \Theta_{\flat}\varrho$ for all $b$.

\emph{(ii), bounds at $b \in P$.} Let $b \in P$. Call the cubes of $\mathcal{S}(P)$ strong, and
write $x = x^{s} + x^{w}$. Here $x^{s}$ is the restriction of $x$ to the rim bonds of the strong
cubes, and $x^{w}$ is the rest.

\emph{The weak part.} Every rim bond within distance $\varpi_\gamma$ of $b$ lies in a strong cube;
this is the definition of $\mathcal{S}(P)$ in \eqref{5.2:eq-rim-constants-a}. Hence $x^{w}$
vanishes within $\varpi_\gamma$ of $b$. By the tail bound of (i) and $|x^{w}| \le \varrho$,
\begin{equation*}
|T(\sigma)x^{w}|_b \le \tau_{\flat}(\varpi_\gamma)\varrho
= \tau_{\flat}(\varpi_\gamma)\, r_A\, \mathfrak{b} \le \tfrac12 c_\eta\mathfrak{b},
\end{equation*}
by the choice of $\varpi_\gamma$ in \eqref{5.2:eq-rim-constants-a}. This bound includes the far
bonds of the rim datum \eqref{5.2:eq-rim-datum-a} and the silent bonds.

\emph{The strong part lies near $\Lambda$.} The bond $b$ touches $\Lambda$. A strong cube holds a
rim bond within $\varpi_\gamma$ of a bond of $P$, and $c_dLM$ is the diameter of a
$\pi_{k+1}$-cube. So every bond of a strong cube lies within
\begin{equation*}
\varpi_\gamma + c_dLM + 1 \le \mathfrak{d}_{\mathrm{far}}
\end{equation*}
of $\Lambda$, the inequality being $(\mathrm{F}2)$; a bond within $\mathfrak{d}_{\mathrm{far}}$ of
$\Lambda$ is not far. Thus no bond of a strong cube is far. Consequently
$|x^{s}| \le \mathfrak{b}$, and $|x^{s}| \le \mathfrak{b}^{\flat}$ within $\varpi$ of $\Lambda$, by
the refinement \eqref{5.2:eq-rim-refinement-a}.

\emph{Contribution of $T_0$.} Since $T_0$ and $x^{s}$ are real, $T_0x^{s}$ is real. Split $x^{s}$
into its bonds within $\varpi$ of $\Lambda$ and the others. The row sums of $|T_0|$ are at most
$\Theta_1$, which bounds the contribution of the first part by $\Theta_1\mathfrak{b}^{\flat}$.
Since $b$ touches $\Lambda$, the bonds of the second part lie at distance at least $\varpi-1$
from $b$, and
\begin{equation*}
\begin{split}
\sum_{b':\,\rho(b,b')\ge\varpi-1} K_0 e^{-\delta_1\rho(b,b')}
&\le K_0 e^{-\delta_1(\varpi-1)/2}\, C_\rho(\delta_1/2)\\
&\le K_0 e^{\delta_1} C_\rho(\delta_1/2)\, e^{-\delta_1\varpi/2}
= K_1 e^{-\delta_1\varpi/2},
\end{split}
\end{equation*}
by the definition of $K_1$ in \eqref{5.2:eq-rim-constants-a}. This gives
\begin{equation*}
|T_0x^{s}|_b \le \Theta_1\mathfrak{b}^{\flat} + K_1 e^{-\delta_1\varpi/2}\mathfrak{b}
\le \bigl(\tfrac14 + \tfrac1{16}\bigr)\mathfrak{b},
\end{equation*}
because $\mathfrak{b}^{\flat} = \mathfrak{b}/4\Theta_1$ and $\varpi$ is chosen so that
$K_1 e^{-\delta_1\varpi/2} \le 1/16$.

\emph{Contribution of $T(\sigma) - T_0$.} By (i), $|x^{s}| \le \mathfrak{b}$ and $(\mathrm{F}1)$,
\begin{equation*}
|[T(\sigma) - T_0]x^{s}|_b \le \varepsilon_T C_\rho(\delta_T)\mathfrak{b} \le \tfrac12 c_\eta\mathfrak{b}.
\end{equation*}
No factor $C_{\mathrm{far}}$ enters here, precisely because $|x^{s}| \le \mathfrak{b}$.

\emph{Conclusion of (ii).} We have
\begin{equation*}
\mathsf{m}_b = T_0x^{s} + [T(\sigma) - T_0]x^{s} + T(\sigma)x^{w},
\end{equation*}
and the first term is real. Therefore
\begin{equation*}
|\operatorname{Re}\mathsf{m}_b| \le \bigl(\tfrac5{16} + c_\eta\bigr)\mathfrak{b}
\le \tfrac7{16}\mathfrak{b},
\qquad
|\operatorname{Im}\mathsf{m}_b| \le c_\eta\mathfrak{b},
\end{equation*}
where $c_\eta \le 1/16$ by \eqref{5.2:eq-rim-constants-a}. Hence
$\mathsf{m} \in \mathfrak{M}_P$.

\emph{(iii).} Both assertions are contained in the random-walk bounds of
Lemma~\ref{5.2:lem-random-walk-bounds}.
\end{proof}

\begin{remark}
The bound
\begin{equation*}
|m_b(x)| \le \tfrac38\,\mathfrak{b}
\qquad\text{for } x \in \mathfrak{X}^{\sharp}
\end{equation*}
is the case $\sigma = \mathbb{1}$ of the argument of part (ii) of
Lemma~\ref{5.2:lem-conditional-mean}, where the intermediate bound
$(\tfrac5{16}+c_\eta)\mathfrak{b} \le \tfrac38\,\mathfrak{b}$, sharper than the box
\eqref{5.2:eq-conditional-mean-a}, is used,
together with condition $(\mathrm{F}_{\mathrm{far}})$ of \eqref{5.2:eq-rim-conditions-a}. This
bound on the mean $m(x)$ at $\sigma = \mathbb{1}$ serves the real splitting of the rim cut-off.

Consider the gauge mode at slope $\mathfrak{b}^{\flat}$ of \eqref{5.2:eq-rim-constants-a}. Its mean
is $d(L-1)\mathfrak{b}/(8\Theta_1)$. By \eqref{5.2:eq-rim-constants-a},
\begin{equation*}
\Theta_1 \ge \tfrac12\, d(L-1),
\end{equation*}
and therefore
\begin{equation*}
\frac{d(L-1)\,\mathfrak{b}}{8\Theta_1}
\le \frac{d(L-1)\,\mathfrak{b}}{4\,d(L-1)}
= \tfrac14\,\mathfrak{b}.
\end{equation*}
This holds for every $L$.
\end{remark}

\begin{theorem}[Filing of rim-dependent boundary terms]\label{5.2:thm-rim-filing}
Let $\delta\in\,]0,\delta_J(L)]$, where $\delta_J(L)=(1-4/L)/10$. Put
$\mathsf{r}(\delta):=(1-10\delta)\tfrac12 L$, and let $\gamma\le\gamma_V$, the ceiling of
Lemma~\ref{5.2:lem-complex-localisation}. Assume:
\begin{itemize}
\item positivity and decay at every true pair, that is, Corollary~\ref{5.2:cor-true-pair-positivity}
and \eqref{5.2:eq-true-pair-positivity-a};
\item the random-walk bounds for the interpolated operators $C_\Lambda(\sigma)$, $\mathcal{A}(\sigma)$
and for the real reference operator $T_0$ of the transfer operator $T(\sigma)$, one of the admitted
inputs of Definition~\ref{5.1:def-admitted-inputs}, in the letters of
Lemma~\ref{5.2:lem-conditional-mean};
\item the far-bond bound that fixes the constant $C_{\mathrm{far}}$, one of the admitted inputs of
Definition~\ref{5.1:def-admitted-inputs};
\item the homogeneity lemma for the constants of the step;
\item the cut-off identity for a shifted sharp cut-off, and the one-term bound;
\item the tolerance lemma;
\item the resummation lemma for Ba{\l}aban's cluster bounds;
\item Ba{\l}aban's estimates \cite[(2.27)--(2.41)]{B-CMP116};
\item the complex localisation of the effective action, Lemma~\ref{5.2:lem-complex-localisation};
\item the one-sided conditions \eqref{5.2:eq-rim-conditions-a} on the rim constants.
\end{itemize}
Consider, at step $k+1$, an output term of the refined step \eqref{5.2:eq-rim-refinement-a}. Let
$\Lambda=\Lambda_{k+1}$ be the region determined by the second member of its pair of labels. Let $x$ be
its rim datum, $\mathfrak{X}^{\sharp}$ its
upper box and $\mathsf{I}(x)$ its rim integral, all as in \eqref{5.2:eq-rim-datum-a}. Then there are
functions $\mathfrak{K}(X,\mathcal{S};\cdot)$, indexed by the domains $X\in\mathbf{D}_{k+1}$ with
$X\cap\Lambda\neq\emptyset$ and by the subsets $\mathcal{S}$ of the set of $\pi_{k+1}$-cubes in $X$ that
hold a rim bond. These functions have the following properties.
\begin{enumerate}
\item[(a)] On a finite torus, at every true pair in the sense of
Corollary~\ref{5.2:cor-true-pair-positivity}, and for every $x\in\mathfrak{X}^{\sharp}$,
\begin{equation}\label{5.2:eq-rim-filing-1}
\log\mathsf{I}(x)=\sum_{X,\mathcal{S}}\mathfrak{K}(X,\mathcal{S};x).
\end{equation}
\item[(b)] Each $\mathfrak{K}(X,\mathcal{S})$ is a stored term with a strong-set datum, with strong set
$\mathcal{S}$, in the sense of Definition~\ref{5.2:def-stored-terms}. It is holomorphic on
\begin{equation*}
\operatorname{space}(X)\times\bigl\{x:\ x(b)\in\mathbb{C},\ |x(b)|<\varrho
\text{ for } b \text{ in no cube of } \mathcal{S}\bigr\}\times\mathcal{W}_k(X),
\end{equation*}
where $\operatorname{space}(X)$ is the configuration space of
Lemma~\ref{5.2:lem-complex-localisation}(a), read over $X$.
It is continuous in the values of $x$ on the cubes of $\mathcal{S}$ as these range over the real box
$|x(b)|\le\mathfrak{b}^{\sharp}(b)$. It is local in $X$: it depends on
$((\mathbf{U},\mathbf{J}),\mathbf{A},w)$ only through their restrictions to $X$ (to $X^{+}$ in the
wrapped case of (c)). It is invariant under the gauge
transformations of \cite[(2.41)(iii)]{B-CMP119}.
\item[(c)] For every such $X$, with $\mathfrak{N}$ the size of Definition~\ref{5.2:def-stored-terms}
and $B_0$ the constant of \eqref{5.2:eq-created-terms-bound-a},
\begin{equation}\label{5.2:eq-rim-filing-2}
\sum_{\mathcal{S}}e^{\kappa_A^{\mathrm{pr}}(g_k)|\mathcal{S}|}\,
\mathfrak{N}\bigl(\mathfrak{K}(X,\mathcal{S})\bigr)
\le\tfrac18 B_0\,e^{-\kappa_{\mathrm{out}}d_{k+1}(X)},
\end{equation}
where $\kappa_{\mathrm{out}}:=\mathsf{r}(\delta)\kappa\ge(1+4\beta_s)\kappa$. If the output term is
wrapped, that is, if its creating cluster contains a relative potential $V^{\mathrm{rel}}(Y)$ as in
\eqref{5.2:eq-complex-localisation-a}, then, by the proposition on the format of wrapped outputs, $X$ in
(b) and (c) is its label. In that case \emph{local in $X$} means that the term depends on
$((\mathbf{U},\mathbf{J}),\mathbf{A},w)$ only through their restrictions to the enlargement $X^{+}$ of
$X$, as in the inductive bound \eqref{5.2:eq-created-terms-bound-a}.
\item[(d)] Put $\mathbf{B}^{\mathrm{sl}}(X,\mathcal{S};x):=\mathfrak{K}(X,\mathcal{S};x)-\mathfrak{K}(X,\mathcal{S};0)$.
Then $\mathbf{B}^{\mathrm{sl}}(X,\mathcal{S};x)$ vanishes unless $X$ contains a rim bond.
\item[(e)] None of the constants in (a)--(d) depends on $k$, $K$, $m$, on the pair of labels of the
output term, on
$\Lambda$, on the sequence $(\Omega_j)_j$, on the volume, or on the ages of large-field regions. The
radius $r$ in the definition of the classes of complex plaquette weights
(Definition~\ref{5.1:def-weight-classes}) enters only through $\gamma_V$.
\end{enumerate}
\end{theorem}

\begin{proof}
\emph{Real data.} The potentials $V(Y)$, with $Y\in\mathbf{D}_k$ and $Y\cap\Lambda\neq\emptyset$, enter
$\mathsf{I}(x)$ through $\exp\sum_Y V(Y)$. This factor is expanded as
\begin{equation*}
\exp\sum_Y V(Y)=\sum_{D}\prod_{Y\in D}\bigl(e^{V(Y)}-1\bigr),
\end{equation*}
where $D$ runs over the finite families of distinct such $Y$. For a given $D$ put
$Y_0:=\bigcup_{Y\in D}Y$. Let $\bar\Lambda$ be the free bonds with at least one end-point in
$\Lambda$. Write $\chi_b(\psi_b+m_b):=\chi(|\psi_b+m_b(x)|<\mathfrak{b})$ for $b\in\bar\Lambda$. Put
$r_0:=\tfrac12\mathfrak{b}-1$, and define
\begin{equation*}
h(y):=1\ \ (y\le r_0),\qquad h(y):=r_0+1-y\ \ (r_0\le y\le\tfrac12\mathfrak{b}),
\qquad h(y):=0\ \ (y\ge\tfrac12\mathfrak{b}).
\end{equation*}
For $b\in Y_0\cap\bar\Lambda$ we write $\chi_b(\psi_b+m_b)=h(|\psi_b|)+\chi^{\mathrm{rem}}_b$. By the real
form of the cut-off identity, $\chi^{\mathrm{rem}}_b:=\chi_b(\psi_b+m_b)-h(|\psi_b|)$ takes values in
$[0,1]$ whenever $m_b$ is real with $|m_b|\le 7\mathfrak{b}/16$; this hypothesis is verified where it is
needed, at the bonds of $P$, in the next step. For
$b\in\bar\Lambda\setminus Y_0$ we write $\chi_b=1-\chi^c_b$, with
$\chi^c_b:=\chi(|\psi_b+m_b(x)|\ge\mathfrak{b})$, as in \cite[(2.3)]{B-CMP116}.

We now expand the product of these factors over $b\in\bar\Lambda$. This gives a sum over pairs
$(P',P'')$. Here $P'\subset Y_0\cap\bar\Lambda$ is the set of bonds at which $\chi^{\mathrm{rem}}_b$ is
chosen, and
$P''\subset\bar\Lambda\setminus Y_0$ is the set at which $-\chi^c_b$ is chosen. Put
$P:=P'\sqcup P''$ and $Z_0:=Y_0\cup\bigcup_{b\in P}\mathsf{N}(b)$, and group the terms by $Z_0$. The
result is the representation of \cite[(2.4)]{B-CMP116}:
\begin{equation*}
\mathsf{I}(x)=\sum_{Z_0}\int d\mu_{C_\Lambda}(\psi)\,
F\bigl(Z_0,\psi|_{I};m|_{Z_0},x\bigr).
\end{equation*}
Here $F$ is the sum of the terms with the given $Z_0$, $I$ is the set of bonds of
\cite[(2.4)]{B-CMP116} on which $F$ reads the fluctuation field, $\psi|_I$ is the restriction of
$\psi$ to $I$, and $F$ is bounded.

\emph{Conditioning and decoupling.} The displays \cite[(2.5)--(2.6)]{B-CMP116} are identities for a
Gaussian field $\psi$ of covariance $C_\Lambda$. They follow from the first two steps of the
Gaussian conditioning lemma, applied to the quadratic form $\mathcal{A}_{\Lambda\Lambda}$ in the role
of that lemma's operator. In them the rim
datum $x$ occurs in $F$ only.

In \cite[(2.8)]{B-CMP116} we interpolate, with the same decoupling parameters, both the operators and
the mean $\mathsf{m}=T(\sigma)x|_{Z_0}$. Let $G$ be the interpolated expression. It has three
properties.
\begin{itemize}
\item $G$ is holomorphic in $\sigma$. This follows from clause~(a) of the one-term bound, which
applies because Lemma~\ref{5.2:lem-conditional-mean} places $\mathsf{m}$ in $\mathfrak{M}_P$.
\item At $\sigma\equiv 1$, that is, with all decoupling parameters equal to one, $G$ is the
uninterpolated expression, since $T(1)x=m(x)$ by Lemma~\ref{5.2:lem-conditional-mean}(iii).
\item When $\sigma=0$ off $Z$, $G$ factorises over the connected components of $Z$, because each
$\mathsf{N}(b)$ is connected. Each factor depends on $(x,\mathbf{U},\mathbf{J})$ only through their
restrictions to its component, by Lemma~\ref{5.2:lem-conditional-mean}(iii).
\end{itemize}
Hence \cite[(2.9)--(2.11)]{B-CMP116} hold, with $H(Z;x)=\sum_P H(Z,P;x)$. For real values of the
decoupling parameters, all objects are real integrals over real fields, and, for $b\in P$,
$|\mathsf{m}_b|\le 7\mathfrak{b}/16$ by Lemma~\ref{5.2:lem-conditional-mean}(ii); this is the
hypothesis under which $\chi^{\mathrm{rem}}_b\in[0,1]$.

\emph{One term.} Fix $D$, $P$ and $Z\supseteq Z_0'$, with $Z_0'$ the set attached to $Z_0$ in the
decoupling expansion of \cite[(2.8)]{B-CMP116}, and consider the corresponding term of
$H(Z,P;x)$. The one-sided conditions \eqref{5.2:eq-rim-conditions-a} include the hypotheses of
Lemma~\ref{5.2:lem-conditional-mean}(ii). That clause therefore gives $\mathsf{m}\in\mathfrak{M}_P$ for
complex $\sigma$. On $\mathfrak{M}_P$, clause~(c) of the one-term bound, which rests on the tolerance
lemma, gives the factors
$\prod_{D}2\nu(Y)$, with $\nu(Y)$ as in \eqref{5.2:eq-complex-localisation-a}, $e^{-3\mathsf{p}|P|}$ and
$e^{2c_F|Z|}$. Cauchy's estimate for the derivatives in
the decoupling parameters $\sigma(\Delta)$, taken on the circles $|\sigma(\Delta)|=e^{\kappa_1}$,
contributes the factor $e^{-(\kappa_1-1)(LM)^{-4}|Z\setminus Z_0'|}$. Together,
\begin{equation}\label{5.2:eq-rim-filing-a}
|\text{term}|\le\prod_{Y\in D}2\nu(Y)\cdot e^{-3\mathsf{p}|P|}
\cdot e^{-(\kappa_1-1)(LM)^{-4}|Z\setminus Z_0'|}\cdot e^{2c_F|Z|}.
\end{equation}

\emph{Resummation.} In \cite[p.~18]{B-CMP116} the definition of $Z_0$ is used once, to bound $|P|$
from below by a multiple of $M^{-4}|Z_0\setminus Y_0|$. Here the corresponding inequality is
\begin{equation*}
|P|\ge N_Z^{-1}M^{-4}|Z_0\setminus Y_0|,
\end{equation*}
since $Z_0\setminus Y_0\subset\bigcup_{b\in P}\mathsf{N}(b)$ and each $\mathsf{N}(b)$ consists of at
most $N_Z$ cubes of side $M$.

We use one of the three factors $e^{-\mathsf{p}|P|}$ in \eqref{5.2:eq-rim-filing-a} as follows. Four
fifths of it, $e^{-\frac45\mathsf{p}|P|}$, are used in \cite[(2.31)--(2.34)]{B-CMP116}. This is
possible because the conditions \eqref{5.2:eq-rim-conditions-a} contain
$\mathsf{p}\ge5(\bar\kappa+1)N_Z$, with $\bar\kappa$ the largest rate in play, together with
Ba{\l}aban's conditions at \cite[(2.31), (2.33)--(2.34)]{B-CMP116} with $\tfrac12\gamma_2\varepsilon_1^2/g_k^2$
replaced by $\mathsf{p}$. The remaining fifth, $e^{-\frac15\mathsf{p}|P|}$, sums over $P$. Since
$P=P'\sqcup P''$, it factorises as $e^{-\frac15\mathsf{p}|P''|}\,e^{-\frac15\mathsf{p}|P'|}$: the first
factor sums over $P''$, and the second over the set $P'\subset Y_0\cap\bar\Lambda$, which the split of
the real data adds. Since $Y_0$ holds at most $d|Y_0|$ bonds, this last sum is bounded by
\begin{equation*}
\sum_{P'\subset Y_0\cap\bar\Lambda}e^{-\frac15\mathsf{p}|P'|}
\le\prod_{b\in Y_0\cap\bar\Lambda}\bigl(1+e^{-\mathsf{p}/5}\bigr)\le\exp\bigl(e^{-\mathsf{p}/5}d|Y_0|\bigr).
\end{equation*}
It is absorbed by means of \cite[(2.30)]{B-CMP116}.

The estimates \cite[(2.27)--(2.30), (2.35)--(2.41)]{B-CMP116} see the term only through
\eqref{5.2:eq-rim-filing-a}. We apply them by the resummation lemma, with the smallness parameter
$\varepsilon_2$ of that lemma set equal to $2\varepsilon_V$,
where $\varepsilon_V\le\varepsilon_V^{*}$ is the constant of Lemma~\ref{5.2:lem-complex-localisation},
read at $\delta\le\delta_J(L)$, in three particular ways:
\begin{itemize}
\item with the supplementary condition that \cite{B-CMP116} imposes on these displays dropped, the
displays being used for every $\delta\le\delta_J(L)$;
\item with the thresholds of \cite[(2.37), (2.39)--(2.40)]{B-CMP116} taken at $e^{5\rho_7}$ and
$e^{5\rho_9}$, where $\rho_m:=(1-m\delta)\tfrac12L\kappa$;
\item with these two thresholds being the two quantities from which $\varepsilon_V^{*}$ is fixed in
the order of choice of the constants of the rim construction,
Notation~\ref{5.2:not-rim-constants}.
\end{itemize}
The output rate is then that of \cite[(2.41)]{B-CMP116}, namely $\mathsf{r}(\delta)\kappa$. From
$\delta\le(1-4/L)/10$ we get $1-10\delta\ge4/L$, and therefore
\begin{equation*}
\mathsf{r}(\delta)\kappa=(1-10\delta)\tfrac12L\kappa\ge2\kappa.
\end{equation*}
This rate is $\kappa_{\mathrm{out}}$ of (c). Since $\beta_s=1/4$, we have $(1+4\beta_s)\kappa=2\kappa$,
so $\kappa_{\mathrm{out}}\ge(1+4\beta_s)\kappa$; this is the lemma on the decay fraction.

\emph{Filing.} By definition, $\mathcal{S}(P)=\bigcup_{b\in P}\mathcal{S}(\{b\})$ and
$N_A=\sup_b|\mathcal{S}(\{b\})|$, so $|\mathcal{S}(P)|\le N_A|P|$. Since
$\kappa_A^{\mathrm{pr}}=\mathsf{p}/N_A$, it follows that
$e^{\kappa_A^{\mathrm{pr}}|\mathcal{S}(P)|}\le e^{\mathsf{p}|P|}$. This is paid by the second factor
$e^{-\mathsf{p}|P|}$ of \eqref{5.2:eq-rim-filing-a}; the third is not used.

The product $\prod_iH(Z_i,P_i)$ in \cite[(2.13)]{B-CMP116} is assigned the strong set
$\bigcup_i\mathcal{S}(P_i)\subset X$. This assignment is consistent for three reasons:
\begin{itemize}
\item the domains of the factors combine as fibre products;
\item the weights are submultiplicative, since
\begin{equation*}
e^{\kappa_A^{\mathrm{pr}}(g_k)|\bigcup_i\mathcal{S}(P_i)|}
\le\prod_ie^{\kappa_A^{\mathrm{pr}}(g_k)|\mathcal{S}(P_i)|};
\end{equation*}
\item the weights of older strong sets are carried inside $\nu(Y)$, by
Lemma~\ref{5.2:lem-complex-localisation}(d).
\end{itemize}

\emph{Conclusion.}
\begin{itemize}
\item Clause~(b), apart from gauge invariance, is the holomorphy of $G$ established in the
conditioning and decoupling step, together with the filing step; clause~(c) is
\eqref{5.2:eq-rim-filing-a} after the resummation and filing steps.
\item Clause~(a) is \cite[(2.11)--(2.13)]{B-CMP116}, read with the functions $H(Z,P;x)$ obtained above.
\item Gauge invariance. The modulus $|\cdot|$, the function $h$, the box $\mathfrak{b}^{\sharp}$ and the
rule $\mathcal{S}(\cdot)$ are $\operatorname{Ad}$-invariant. Moreover
$T(\sigma;\mathbf{U}^u)\operatorname{Ad}_u=\operatorname{Ad}_uT(\sigma;\mathbf{U})$ by
\cite[p.~21]{B-CMP116}. Hence each $\mathfrak{K}(X,\mathcal{S})$ is invariant under the gauge
transformations of \cite[(2.41)(iii)]{B-CMP119}, which is the last assertion of (b).
\item Clause~(d). If $X$ contains no rim bond, then $\mathsf{m}\equiv0$ and the potentials $V$ do not
depend on $x$. Hence $\mathfrak{K}(X,\mathcal{S};x)=\mathfrak{K}(X,\mathcal{S};0)$.
\item Clause~(e). The constants entering (a)--(d) are those of Notation~\ref{5.2:not-rim-constants}
and of the conditions \eqref{5.2:eq-rim-conditions-a}. Of these, $\varpi_\gamma$, $N_A$ and $N_Z$ are
functions of $g_k$; they enter only on the side of the estimates that is controlled by
$\mathsf{p}(g_k)$, through $\mathcal{S}(\cdot)$, $\kappa_A^{\mathrm{pr}}=\mathsf{p}/N_A$, the lower bound
on $|P|$ and \eqref{5.2:eq-rim-conditions-a}. The radius $r$ of
Definition~\ref{5.1:def-weight-classes} enters only through the hypothesis $\gamma\le\gamma_V$.
\end{itemize}
\end{proof}

\begin{lemma}[The boundary-term bound with the rim tilt]\label{5.2:lem-tilted-bound}
Work at the site $(\mathrm{T})$ of Definition~\ref{5.1:def-expansion-rule}, where the large-field
operation $\mathbf{T}$ acts, with $\Lambda_{k+1} := \Lambda^{*}$. Let the complex source $w$ lie in
$\mathcal{W}_k$ or be constant. Let $x\in\mathfrak{X}^{\sharp}$ be the rim datum and $\mathsf{I}(x)$ the
rim integral of \eqref{5.2:eq-rim-datum-a}.
\begin{enumerate}
\item[(i)] $\log\mathsf{I}(0)$ is decomposed by the expansion \cite[(3.26)--(3.47)]{B-CMP119}, run at
$A_k = 0$. There the measures \cite[(3.27), (3.29)]{B-CMP119} are untilted and the first expression
of \cite[(3.28)]{B-CMP119} is absent. Lemma~\ref{5.2:lem-first-stage-bound} and
Lemma~\ref{5.2:lem-second-stage-activities} apply as they stand. They give $E^{(k+1)}$,
$R''^{(k+1)}$ and $\mathbf{B}^{(k+1)}|_{A_k=0}$.
\item[(ii)] With $\mathbf{B}^{\mathrm{sl}}$ the slope part of the filed rim terms of
Theorem~\ref{5.2:thm-rim-filing},
\begin{equation}\label{5.2:eq-tilted-bound-1}
\log\mathsf{I}(x)-\log\mathsf{I}(0)=\sum_{X,\mathcal{S}}\mathbf{B}^{\mathrm{sl}}(X,\mathcal{S};x;w).
\end{equation}
\item[(iii)] Let $\mathbf{B}^{(k+1)}$ be the sum of three families: the terms of (i), the terms of (ii),
and the $\psi$-free terms. Here $\psi$-free means that the centred fluctuation field
$\psi := A-m(A_k)$, where $A$ is the field integrated at this step and $m$ the mean of the rim
construction, does not occur.
Then $\mathbf{B}^{(k+1)}$ obeys \cite[(2.41)(i)--(iv)]{B-CMP119} and has the support property of
\cite[p.~262]{B-CMP119}: every component of a polymer contains a $Y$ or a $P$-bond meeting
$\Lambda\subset\Omega_{k+1}$.

It also obeys the bound \cite[(2.42)]{B-CMP119} in two-sided form. This holds
\begin{itemize}
\item for every $\mathbf{A}$ in the upper box $\mathfrak{X}^{\sharp}$ of \eqref{5.2:eq-rim-datum-a},
which contains the support of the refined cut-off of \eqref{5.2:eq-rim-refinement-a};
\item on the enlarged spaces of \cite{B-CMP119};
\item at the rate $(1+4\beta_s)\kappa$;
\item analytically in $w\in\mathcal{W}_k(X)$.
\end{itemize}
Explicitly:
\begin{equation}\label{5.2:eq-tilted-bound-a}
\bigl|\mathbf{B}^{(k+1)}\bigl(X;(\mathbf{U},\mathbf{J}),\mathbf{A};w\bigr)\bigr|
\le B_0\,e^{-(1+4\beta_s)\kappa\, d_{k+1}(X)} .
\end{equation}
\end{enumerate}
For a wrapped output, in the sense of Theorem~\ref{5.2:thm-rim-filing}, $X$ denotes its label, and the
decay in \eqref{5.2:eq-tilted-bound-a} is in $d_{k+1}(X)$. The locality of
\cite[(2.41)(i)]{B-CMP119} then
means that the term depends on $((\mathbf{U},\mathbf{J}),\mathbf{A},w)$ only through their
restrictions to $X^{+}$.
\end{lemma}

\begin{proof}
The argument uses the potentials $V$ only through the properties that
Lemma~\ref{5.2:lem-second-stage-activities} consumes. It does not depend on where $V$ reads the
background field, since the strips contain no integration variable.

\emph{Decomposition.} The total is that of \cite{B-CMP119}: the expressions
\cite[(3.26)]{B-CMP119} and \cite[(3.28)]{B-CMP119} sum to $\log\mathsf{I}(x)$ plus a constant.
The $E$- and $R''$-terms have domains $X\subset\Lambda_{k+1}$. They do not depend on $\Lambda_{k+1}$
\cite[p.~278, p.~279]{B-CMP119}. By \cite[(2.27)(i), (2.30)]{B-CMP119}, they carry no
$\mathbf{A}$-slot. They are therefore the terms of (i).

\emph{What is free of the rim.} Two objects are free of the rim.
\begin{itemize}
\item The splitting $\mathsf{I}(x)=\mathsf{I}(0)\cdot[\mathsf{I}(x)/\mathsf{I}(0)]$.
\item The whole-lattice integral at the site $(\mathrm{T})$, regarded as a functional of its inputs.
\end{itemize}
For the second, $b_{k+1}=\beta[E^{(k+1)}_\zeta]$ is extracted from the whole-lattice family by the
extraction formula for the marginal coefficient, which does not involve $\Lambda_{k+1}$. This is said
of the functional only. Moreover, the $E$- and $R''$-outputs of the site $(\mathrm{T})$ carry no
$\mathbf{A}$-slot.

The inputs to which the functional is applied are those of the chain of steps carrying the refinement
\eqref{5.2:eq-rim-refinement-a} itself. These inputs are the families
$(R^{(j)}_\zeta)_{j\le k_1(k)}$, produced by $\mathbf{R}$ acting on the components that carry the
label $I$, and
hence $Q^{(k+1)}$, $b_{k+1}$, $x_k$, $g_k$, $\check{R}_k$ and $K$. The boundary terms
$\mathbf{B}^{\zeta}$ and the operations $\mathbf{T}^{\zeta}$ of the inductive hypothesis at running
shift $\zeta$ carry that refinement in the same way as the local run.

The running-shift statement of the correlation theorem (Theorem~\ref{5.1:thm-correlation-theorem}),
together with its reference-recursion clause and its sourced-format clause, is
proved for that flow with the constants unchanged. These follow from the proposition on the
$\mathbf{R}$-operation read at the datum $A_{\flat}$, together with \eqref{5.2:eq-rim-payment-a}.
The inductive bound on boundary terms is read in the form fixed for the refined step
\eqref{5.2:eq-rim-refinement-a}.

\emph{Rates.} We treat the four kinds of output in turn.
\begin{itemize}
\item The outputs of Lemma~\ref{5.2:lem-second-stage-activities} that enter (i) carry the rate
$\mathsf{r}(\delta)\kappa$. By the lemma on the decay fraction $\delta_J(L)$, this satisfies
$\mathsf{r}(\delta)\kappa\ge 2\kappa$.
\item By \eqref{5.2:eq-rim-constants-a} and \eqref{5.2:eq-three-families-a},
$L^{(k+1)}\subset E^{(k+1)}$. These terms carry the rate $\delta_0 M$ of \cite[\S2]{B-CMP116}. We have
$\delta_0M\ge 2\kappa$. Indeed, \cite[\S2]{B-CMP116} gives
\[
e^{-\frac13\delta_0M}\,e^{16\kappa_1}<1 ,
\]
that is, $\delta_0M>48\kappa_1$. Also $\kappa_1\ge 4L\kappa+1>\kappa$. Together these give
$\delta_0M>48\kappa>2\kappa$.
\item The $\psi$-free potentials carry the rate $\rho_3$. Since $\rho_m=(1-m\delta)\tfrac12 L\kappa$,
we have $\rho_3\ge\rho_{10}=\mathsf{r}(\delta)\kappa$.
\item The $\psi$-free remnants carry the rate $\delta_1M\check{R}_{k+1}$, and
$\delta_1M\check{R}_{k+1}\ge\kappa_1\check{R}_{k+1}$. Here $\delta_1=\min\{\mu_0/2,\delta_W/2\}$ is the
constant of \eqref{5.2:eq-rim-constants-a}, with $\mu_0$ and $\delta_W$ as fixed there. The inequality
$\delta_1M\ge\kappa_1$ is the floor on $M$
imposed after $\kappa_1$ among the one-sided conditions fixed with the existential constant $L_0$.
With it,
\[
\delta_1M\ge\kappa_1\ge 4L\kappa+1 ,
\]
and the size of the remnants is that of \eqref{5.2:eq-rim-conditions-a}.
\end{itemize}

\emph{Part (ii).} Apply Theorem~\ref{5.2:thm-rim-filing}(a) at $x$ and at $0$. Both points lie in
$\mathfrak{X}^{\sharp}$; note that $\mathsf{I}(0)$ is a value of a function, not a term of the
expansion. This gives
\[
\log\mathsf{I}(x)-\log\mathsf{I}(0)
=\sum_{X,\mathcal{S}}\bigl[\mathfrak{K}(X,\mathcal{S};x)-\mathfrak{K}(X,\mathcal{S};0)\bigr].
\]
By Theorem~\ref{5.2:thm-rim-filing}(d), the summand is $\mathbf{B}^{\mathrm{sl}}(X,\mathcal{S};x;w)$,
and it vanishes unless $X$ holds a rim bond. This is \eqref{5.2:eq-tilted-bound-1}.

\emph{The $\psi$-free terms.} There are two kinds.
\begin{itemize}
\item The potentials with $Y\cap\bar\Lambda=\emptyset$. These are bounded by $\nu(Y)$ of
\eqref{5.2:eq-complex-localisation-a} and re-filed in $\mathbf{D}_{k+1}$ by
\cite[(2.36)]{B-CMP119}.
\item The inter-component remnants $\langle x_c,S_{cc'}x_{c'}\rangle$ of \cite[(2.21)]{B-CMP119}
(see \cite[p.~258]{B-CMP119}). Here $x_c$ is the restriction of $x$ to the component $c$ and
$S_{cc'}$ the corresponding block of the Schur complement $S$. They are expanded in random walks by
the random-walk input among the admitted inputs of Definition~\ref{5.1:def-admitted-inputs}, which
gives
\[
\bigl|\langle x_c,S_{cc'}x_{c'}\rangle\bigr|\le C\varrho^{2}e^{-\delta_1M\check{R}_{k+1}} .
\]
This bound holds holomorphically on $|x|<\varrho$ and persists at real $|x|\le C_{\mathrm{far}}\mathfrak{b}$.
By \eqref{5.2:eq-rim-conditions-a} it is at most $\tfrac14B_0$.
\end{itemize}

\emph{Sizes.} $B_0$ is the existential constant of \cite[Theorem~1]{B-CMP119}. We bound the three
families in turn, each bound carrying the factor $e^{-\kappa_{\mathrm{out}}d_{k+1}(X)}$ with
$\kappa_{\mathrm{out}}=\mathsf{r}(\delta)\kappa$.
\begin{itemize}
\item Take $B_0$ at least four times the untilted output of
Lemma~\ref{5.2:lem-second-stage-activities}. The terms of (i) are then at most $\tfrac14B_0$.
\item By the choice of $\varepsilon_V^{*}$, the terms of (ii) are at most $2\cdot\tfrac18B_0$, each
slope term being a difference of two filed terms.
\item The $\psi$-free terms are at most $\tfrac14B_0$.
\end{itemize}
The sum is at most $\tfrac34B_0e^{-\kappa_{\mathrm{out}}d_{k+1}(X)}$. Since
$\kappa_{\mathrm{out}}\ge(1+4\beta_s)\kappa$, this implies \eqref{5.2:eq-tilted-bound-a}.

\emph{The support property in (iii).} This is the property of \cite[p.~262]{B-CMP119}: every
component of a polymer contains a $Y$ or a
$P$-bond meeting $\Lambda\subset\Omega_{k+1}$.

\emph{Analyticity in $w$.} The source $w$ enters only as a parameter of the potentials $V$. It does so
under majorants that do not depend on $w$: see the table \eqref{5.2:eq-unit-table-a}, with the constant
$\bar\varepsilon^{w}$ inside \eqref{5.2:eq-rim-payment-a}. The mean $m$, the operator $T(\sigma)$
and the function $h$ of the proof of Theorem~\ref{5.2:thm-rim-filing} do not depend on the
coefficients $z_i$ of the source $w=\sum_i z_if_i$.
\end{proof}

The refined box is needed in this lemma: by the corollary on the bilinear tilt, the terms tilted by
$A_k$ do not satisfy a bound of the form \cite[(2.42)]{B-CMP119} on the small-field box of the
cut-offs \cite[(3.21)]{B-CMP119}. The two expectations in \cite[(3.26)]{B-CMP119} are also tilted;
the proof above never uses them at $x\neq0$.

\begin{corollary}[The created-term bound at every step]\label{5.2:cor-created-bound-induction}
Assume the hypotheses of Theorem~\ref{5.2:thm-rim-filing} (filing of rim-dependent boundary terms)
at every step. Assume also clause~(a) of the flow proposition for the rim-tilted chain at the
earlier steps, both within the single induction on the level by which the correlation theorem is
proved. Then the inductive bound on created boundary terms
(Definition~\ref{5.2:def-created-terms-bound}) holds at every step $k'$, with the weights
displayed in \eqref{5.2:eq-created-terms-bound-a}: the weight
$\kappa_A^{\mathrm{pr}}(g_{k'-1})$ on the newest slot and the weight $\kappa_A(k'-1)$ on the
older slots. In particular, in the notation of Definition~\ref{5.2:def-created-terms-bound}, the
terms of $\mathbf{B}^{(k')}$ created at step $k'$ are finite sums
$\sum_{\mathcal{S}}\mathbf{B}^{(k')}(X,\mathcal{S})$ of stored terms with datum $A^{\flat}$, the
strong sets $\mathcal{S}$ being declared by the rule of the rim construction and by unions, and
\begin{equation*}
\begin{split}
&\sum_{\mathcal{S}} \exp\Bigl[\kappa_A^{\mathrm{pr}}(g_{k'-1})\,|\mathcal{S}_{k'-1}|
 +\kappa_A(k'-1)\sum_{j'<k'-1}|\mathcal{S}_{j'}|\Bigr]\\
&\qquad\times\mathfrak{N}\bigl(\mathbf{B}^{(k')}(X,\mathcal{S})\bigr)
 \le B_0\,e^{-\kappa_{\mathrm{out}}d_{k'}(X)},
\end{split}
\end{equation*}
where
\begin{equation*}
\kappa_A(k'-1)=\min_{i\le k'-1}\kappa_A^{\mathrm{pr}}(g_i).
\end{equation*}
\end{corollary}

\begin{proof}
We argue by induction on the step. Suppose the assertion holds at the steps up to $k$, and
consider step $k+1$. As in clause~(iii) of Lemma~\ref{5.2:lem-tilted-bound} (the boundary-term
bound with the rim tilt), we put
\begin{equation}\label{5.2:eq-created-bound-induction-1}
\mathbf{B}^{(k+1)}
:=\mathbf{B}^{(k+1)}\big|_{A_k=0}+\sum\mathbf{B}^{\mathrm{sl}}
+\bigl(\text{terms free of the fluctuation field}\bigr).
\end{equation}
The first summand is obtained at $A_k=0$. The second is the sum of the slope parts of the filed
rim terms. The third consists of the terms free of the Gaussian fluctuation field.

(\(\alpha\)) The newest slot. The newest slot is carried by the slope parts. By clause~(d) of
Theorem~\ref{5.2:thm-rim-filing}, each slope part is
\begin{equation*}
\mathbf{B}^{\mathrm{sl}}(X,\mathcal{S};x)=\mathfrak{K}(X,\mathcal{S};x)-\mathfrak{K}(X,\mathcal{S};0),
\end{equation*}
and it vanishes unless $X$ contains a rim bond. By clause~(b), each $\mathfrak{K}(X,\mathcal{S})$
is a stored term of the kind described there. The subtraction of the value at $x=0$ at most
doubles the size, that is,
\begin{equation*}
\mathfrak{N}\bigl(\mathbf{B}^{\mathrm{sl}}(X,\mathcal{S})\bigr)
\le 2\,\mathfrak{N}\bigl(\mathfrak{K}(X,\mathcal{S})\bigr).
\end{equation*}
Clause~(c) of the same theorem then gives
\begin{equation}\label{5.2:eq-created-bound-induction-2}
\sum_{\mathcal{S}}e^{\kappa_A^{\mathrm{pr}}(g_k)|\mathcal{S}|}\,
\mathfrak{N}\bigl(\mathbf{B}^{\mathrm{sl}}(X,\mathcal{S})\bigr)
\le \tfrac14 B_0\,e^{-\kappa_{\mathrm{out}}d_{k+1}(X)} .
\end{equation}

The first summand of \eqref{5.2:eq-created-bound-induction-1} is evaluated at $A_k=0$, so it
does not depend on the rim datum of step $k+1$.

The terms free of the fluctuation field are of two kinds. The first kind are the potentials $V(Y)$ with
$Y\cap\bar\Lambda=\emptyset$. The second kind are the inter-component remnants. Here the
components are $\infty$-connected unions of big cubes, any two of them at least one cube of side
$LMR_{k+1}$ apart, by the separation statement for rim components. Both kinds are holomorphic on
$|x|<\varrho$. The remnants are bounded by $C\varrho^2e^{-\delta_1MR_{k+1}}$, where $\delta_1$ is
the decay rate fixed for these remnants in the rim construction and $R_{k+1}$ denotes the radius
function $R(\cdot)$ of the glossary at level $k+1$. By clauses~(a)
and~(b) of the theorem on the rim-datum slot of stored terms, these terms are filed with
$\mathcal{S}=\emptyset$.

(\(\beta\)) The older slots. The rim data of the earlier steps enter step $k+1$ only as parameters
of the potentials $V$ of Lemma~\ref{5.2:lem-complex-localisation} (complex localisation of the
effective action): the Gaussian measure, the cutoff function $h$ of the first step of that proof,
$\mathsf{L}$ and $\mathsf{L}^c$, the
contours and the majorants of the five steps of the proof of Theorem~\ref{5.2:thm-rim-filing} are
all free of them. The parameter lemma and its corollary therefore apply,
after centring. Under them the strong sets combine by unions. The weights are submultiplicative
and are carried inside the bound $\nu(Y)$ of clause~(d) of
Lemma~\ref{5.2:lem-complex-localisation}, at the weight $\kappa_A(k)$. Since $k'=k+1$, this is
the weight $\kappa_A(k'-1)$ on the older slots.

(\(\gamma\)) Later sites. At the large-field operations $\mathbf{T}^{(k'')}$ of later steps, at
the R-terms of the large-field expansion and at the site (L) \cite[(1.49)]{B-CMP122b}, the exterior
rim data are
parameters of stored terms only. Where the gauge map carried by a site rotates them, they are read
at the presented rotation through the orbit datum of the term. This uses the definition of the
orbit datum together with clauses~(a) and~(g) of the rotation statement. At the sites (T) and (P)
the presented rotation is read through the presented-rotation statement for those sites, and at
(L) through $\mathfrak{Q}^{\mathfrak{b}}(\zeta)$, as in \eqref{5.2:eq-three-sites-a}.

The expansions at these sites are of the type treated in the rim-tilted expansion proposition.
The parameter lemma and its corollary carry the datum of the stored terms through them: strong
sets go to unions, radii are kept and weights are multiplied, as in the treatment of later steps
that accompanies the parameter lemma. The rotation lemma passes the rotation as one
more holomorphic coordinate.

Part~(\(\alpha\)) gives the weight $\kappa_A^{\mathrm{pr}}(g_k)$ on the newest slot. Part~(\(\beta\))
gives the weight $\kappa_A(k)$ on the older slots. Part~(\(\gamma\)) carries the datum through the
later sites. This completes the induction.
\end{proof}

\begin{remark}
Corollary~\ref{5.2:cor-created-bound-induction} is the form of the inductive bound on created
boundary terms that is read by the first induction hypothesis of the A-slot bookkeeping and by
the correlation theorem (Theorem~\ref{5.1:thm-correlation-theorem}).
\end{remark}

\begin{remark}[The strong-set weight dominates powers of the logarithm]\label{5.2:rem-strong-weight}
Write $\mathsf{T} := \log g^{-2}$, $\mathsf{T}_k := \log g_k^{-2}$ and
$\check{\mathsf{T}}_k := \log \min_{i\le k} g_i^{-2}$, and let the constants of the rim construction
be fixed as in \eqref{5.2:eq-rim-constants-a}. There is a constant $c>0$, depending only on the
constants of \eqref{5.2:eq-rim-constants-a} fixed before $g$, such that
\begin{equation}\label{5.2:eq-strong-weight-1}
\kappa_A^{\mathrm{pr}}(g) \ge c\,\frac{\gamma_2 A_1^2\, \mathsf{T}^{2p_0}}{(\log \mathsf{T})^{d}},
\qquad \mathsf{T}\ge e,
\end{equation}
so that, for every real $a$, $e^{-\kappa_A(k)}\,\mathsf{T}_k^{\,a} \to 0$ as $g_k\to 0$, uniformly
over the levels $k$ at which $g_k^{-2}=\min_{i\le k}g_i^{-2}$, that is, at which
$\check{\mathsf{T}}_k=\mathsf{T}_k$: at each such level the third member of the
bound on the A-slot family in the anchored-domain norm $\|\cdot\|^{A}$, which carries the factor
$e^{-\kappa_A(k)}$ against at most a power of $\log g_k^{-2}$, decays
faster than every power of $\log g_k^{-2}$. The estimates of the strong-set expansion depend on
$\kappa_A$ only through the factor $e^{\kappa_A(k)|\mathcal{S}|}$ and through the smallness
constant of the strong sets, at a single level $k$. In the counting of the preparatory stage the
large-field condition on the bonds of the rim is imposed at the tight threshold: with $j$ the level
of that counting, $\mathfrak{b}_j$ its small-field box and
$\mathfrak{b}^{\flat}_j := \mathfrak{b}_j/(4\Theta_1)$ as in \eqref{5.2:eq-rim-constants-a}, the
weight $\tilde{\mathfrak{l}}$ of that counting is replaced by
\begin{equation*}
2d\,(LM)^4\, e^{-\gamma_2 \mathfrak{b}_j^2/(32\Theta_1^2) + c_J + c_{\flat}}
= 2d\,(LM)^4\, e^{-\frac12\gamma_2 (\mathfrak{b}^{\flat}_j)^2 + c_J + c_{\flat}} ,
\end{equation*}
where $c_J$ and $c_{\flat}$ are the constants of that counting.
\end{remark}

\begin{proof}
By \eqref{5.2:eq-rim-constants-a},
$\mathsf{p}(g) = \gamma_2 A_1^2\mathsf{T}^{2p_0}/64$,
$\kappa_A^{\mathrm{pr}}(g) = \mathsf{p}(g)/N_A(g)$ and
$N_A(g) \le (3 + 2\varpi_\gamma/(LM))^d$.
Put $C_\tau := 2(K_0+1)e^{\delta_\flat}C_\rho(\delta_\flat/2)$, so that
$\tau_\flat(n) = \tfrac12 C_\tau e^{-\delta_\flat n/2}$ for $n\ge 0$. The inequality
$\tau_\flat(n) r_A \le \tfrac12 c_\eta$ holds exactly when
$n \ge (2/\delta_\flat)\log(C_\tau r_A/c_\eta)$. Since $\varpi$ is an integer, the least integer
$\varpi_\gamma \ge \varpi$ satisfying it obeys
\begin{equation*}
\varpi_\gamma \le \varpi + 1 + \frac{2}{\delta_\flat}\log_+\frac{C_\tau r_A}{c_\eta},
\end{equation*}
where $\log_+ := \max\{\log,0\}$; the positive part covers the case $C_\tau r_A\le c_\eta$, in
which $\varpi_\gamma=\varpi$.
Substituting $r_A = c_A (C_\ast/A_1)\mathsf{T}^{q-p_0}$ gives
\begin{equation*}
\log\frac{C_\tau r_A}{c_\eta} = \log\frac{C_\tau c_A C_\ast}{A_1 c_\eta} + (q-p_0)\log\mathsf{T},
\end{equation*}
so that for $\mathsf{T}\ge e$ there is a constant $C_\varpi$, depending on the constants fixed
before $g$, with $\varpi_\gamma \le C_\varpi\log \mathsf{T}$, and hence
$N_A(g) \le \bar{C}_\varpi (\log\mathsf{T})^d$ with $\bar{C}_\varpi := (3+2C_\varpi/(LM))^d$.
Therefore
\begin{equation}\label{5.2:eq-strong-weight-2}
\kappa_A^{\mathrm{pr}}(g) \ge \frac{\gamma_2A_1^2}{64\,\bar{C}_\varpi}\, h(\mathsf{T}),
\qquad h(\mathsf{T}) := \frac{\mathsf{T}^{2p_0}}{(\log\mathsf{T})^d},
\qquad \mathsf{T}\ge e .
\end{equation}
This is \eqref{5.2:eq-strong-weight-1} with $c := 1/(64\,\bar{C}_\varpi)$.

The function $h$ satisfies $h'(\mathsf{T}) = h(\mathsf{T})\,(2p_0\log\mathsf{T} - d)/(\mathsf{T}\log\mathsf{T})$,
so it is increasing for $\log\mathsf{T} > d/(2p_0)$, and it tends to infinity because $p_0>0$.
Since $\kappa_A(k) = \min_{i\le k}\kappa_A^{\mathrm{pr}}(g_i)$ and
$\min_{i\le k}\mathsf{T}_i = \check{\mathsf{T}}_k$, inequality \eqref{5.2:eq-strong-weight-2} gives,
once $\check{\mathsf{T}}_k \ge \max\{e, e^{d/(2p_0)}\}$,
\begin{equation*}
e^{-\kappa_A(k)}\,\mathsf{T}_k^{\,a}
\le \exp\Bigl(-\frac{\gamma_2A_1^2}{64\,\bar{C}_\varpi}\,h(\check{\mathsf{T}}_k) + a\log\mathsf{T}_k\Bigr).
\end{equation*}
At a level $k$ with $\check{\mathsf{T}}_k = \mathsf{T}_k$ the right side is a function of
$\mathsf{T}_k$ alone, and it tends to zero as $\mathsf{T}_k\to\infty$, that is, as $g_k\to0$,
because $\log\mathsf{T}_k/h(\mathsf{T}_k)\to 0$; the convergence is therefore uniform over such
levels. This is the assertion that the third member has degree $-\infty$ at each such level.
The remaining assertions of the remark, on how $\kappa_A$ enters the estimates of the strong-set
expansion and on the threshold at which the large-field condition on the rim bonds is imposed in
the counting of the preparatory stage, are conventions and need no proof.
\end{proof}

\begin{proposition}[Payment of the refined rim and the modified cover]\label{5.2:prop-rim-payment}
Let $d=4$. Fix a term of the exact refinement of the rim cut-offs
(Definition~\ref{5.2:def-rim-refinement}), with output $(\Lambda^*,R^*,\mathsf{C}^*)$ and refined
violator set $I^*$. Let $\mathfrak{X}$ be the set of free bonds, $x$ the rim datum, $S$ the Schur
complement of $\mathcal{A}$ on the rim coordinates and $S_{cc}$ its compression, as in
\eqref{5.2:eq-rim-datum-a}. Let $\mathfrak{b}$ be the small-field box of the step; when the step is
read in the notation of \cite{B-CMP122b} at level $j$, write $\mathfrak{b}_j$,
$\mathsf{T}_j=\log g_j^{-2}$ and $\check{R}_j$ for the corresponding quantities. Let $n_{\mathrm{rim}}$
be the number of rim bonds. Here $P_j$, $Q_j$, $R_j$, $Z'_{j-1}$ and $Z_i^{(j)}$ are the sets of
\cite[(1.76)]{B-CMP122b}, $Z_i^{(j)\flat}$ the corresponding sets of the refined step, $I_j$ the set
$I^*$ of the step at level $j$, $\mathfrak{s}(R)$ the set of cubes attached to $R$ in the complement
lemma of (h4) below, which is monotone in $R$, $X^{\sim n}$ the enlargement of a set $X$ by $n$
layers of cubes in the sense of \cite{B-CMP122b}, and $X^{+n}$ its enlargement by $n$ layers of cubes
of the rim construction. In the factors $(MR_j)^{-d}$ and $M^dR_j^{d+1}$ quoted below from
\cite{B-CMP122b}, $R_j$ denotes instead Ba{\l}aban's number of that name, not the region $R_j$ of
\cite[(1.76)]{B-CMP122b}. Assume:
\begin{enumerate}
\item[(h1)] positivity and decay at every true pair (Corollary~\ref{5.2:cor-true-pair-positivity}):
$\mathcal{A}\ge\gamma_0$ on $\ell^2(\mathfrak{X};\mathfrak{g})$ at every true pair, aged and inner
backgrounds included \eqref{5.2:eq-true-pair-positivity-a}; this is
Proposition~\ref{5.2:prop-form-positivity}(b) under the scaffold-minimiser hypothesis at the domain
$\Omega_{k+1}$;
\item[(h2)] the facts of the rim construction: (i) for every rim datum $x$,
\begin{equation*}
\langle x,Sx\rangle=\min_{y}\,\langle (y,x),\mathcal{A}(y,x)\rangle ,
\end{equation*}
where $y$ ranges over the remaining coordinates of $\mathfrak{X}$, and $S_{cc}$ is a compression of
$S$; (ii) the rim bound: on the support of the refined characteristic functions, whenever
$S\ge\gamma_0$ and $S_{cc}\ge\gamma_0$,
\begin{equation*}
e^{-\frac12\langle x,Sx\rangle}\le\prod_{b\in I^*}e^{-\gamma_0\mathfrak{b}^2/(64\Theta_1^2)}
\cdot e^{-\frac14\langle x,Sx\rangle},
\end{equation*}
and (iii) the constant $\Theta_1$ of the rim construction satisfies $\Theta_1\ge d(L-1)/2\ge 8$;
\item[(h3)] the properties of the exact refinement of the rim cut-offs: its terms are indexed by the
pair $(R,I^*)$, the index being the set $I^*$ itself (history-free indexing); a violator bond lying
between two cubes contributes both cubes, and $R^*=R\cup\operatorname{cubes}(I^*)$;
\item[(h4)] the complement lemma for the refined domain:
\begin{equation*}
\Lambda^{*c}\subset\mathfrak{s}(R^*)^{+99},\qquad
\#(\Lambda^{*c})\le 199^d\bigl[\#\mathfrak{s}(R)+2\cdot 3^d\,\#\operatorname{cubes}(I^*)\bigr];
\end{equation*}
\item[(h5)] the degree conditions $p_0\ge 5r_{\mathrm{pr}}$ \cite[p.~246]{B-CMP119}, and, for
Ba{\l}aban's exponents $p_1$ and $r_0$ of \cite{B-CMP122b}, the conditions $r_0\ge r_{\mathrm{pr}}$,
$p_1<p_0$ and $2p_1-(d+5)r_0>p_0$;
\item[(h6)] the summation of the large-field terms, with its constants $C_{\mathfrak{d}}$ and
$c_0:=2d+24$ (numbers depending on $d$ only; this $c_0$ is not the flow constant $c_0$ of the
glossary); the per-cube count of the large-field terms, with its
per-cube rate, its positivity constant $\lambda_0$ and its summation condition; the smallness
conditions on $\gamma$; the modified large-field bound with its volume clause and the per-cell value
constant $\mathsf{K}_{\mathrm{cell}}(k,j)$; and the bound $1-\theta_W\ge\frac34$ on the ratio of
complex to real weight.
\end{enumerate}
Set, with $\beta_0$ the constant of that name in \cite{B-CMP122b} and with $A_\flat$ as in the first
line of \eqref{5.2:eq-rim-payment-a} below,
\begin{gather*}
\mathsf{Q}_j:=\frac{\gamma_0A_\flat^2}{(1+\beta_0)^2}\,\mathsf{T}_j^{2p_0},\qquad
\mathsf{D}_j:=\frac{\mathsf{Q}_j}{2\sqrt d},\qquad
\mathsf{P}_j:=p_0(g_j)=A_0\mathsf{T}_j^{p_0},\qquad
c_j:=C^\flat_{\mathfrak{d}}M^d\check{R}_j^{\,d+1},\\
C^\flat_{\mathfrak{d}}:=C_{\mathfrak{d}}+2^{d+2}(1+dL^d),\qquad
\Gamma:=\sqrt d\,2^{d+1}219^d,\qquad \sigma:=220\sqrt d,
\end{gather*}
where $\mathfrak{b}_j^2=A_1^2\mathsf{T}_j^{2p_0}$ is the convention of \cite{B-CMP122b}, and consider
\begin{equation}\label{5.2:eq-rim-payment-a}
\begin{aligned}
&A_\flat^2:=\frac{A_1^2}{256\,\Theta_1^2},\qquad
\kappa_{\mathfrak{d}}:=C^\flat_{\mathfrak{d}}M^d,\qquad
\gamma_0A_\flat^2\,\mathsf{T}_\gamma^{p_0}\ge 8,\\
&\text{each $O(1)$ multiplying $|Z_j|$ or $d'_j(Z_j)$ is replaced by}\\
&\qquad 199^d\,3^{d+1}\bigl(C_\sharp O(1)+1\bigr)\bigl(1+2^{-d}\log g_j^{-2}\bigr),\\
&Z_j\subset\Bigl[Z'^{\sim 10}_{j-1}\cup\bigcup_i\bigl(Z_i^{(j)\flat}\bigr)^{\sim 2}\Bigr]^{+99},\\
&(\mathrm{S}^\flat)\quad \mathsf{Q}_j\ge 4\mathsf{P}_j,\\
&(\mathrm{B}^\flat)\quad \mathsf{D}_j\ge\Gamma c_j\bigl(9+4c_0+4\check{R}_j\bigr),\\
&(\mathrm{M}^\flat)\quad (1+\beta_0)^{-1}\mathsf{P}_j\ge
\Gamma c_j\bigl(5\sigma+4\sqrt d\,101^d+4c_0+4\check{R}_j+48\bigr),\\
&(\mathrm{F}^\flat)\quad d\,(LM\check{R}_j)^d\,
\exp\Bigl(-\frac{\gamma_0\mathfrak{b}_j^2}{128\,\Theta_1^2(1+\beta_0)^2}\Bigr)\le 1,\\
&\gamma_P^\flat:=\sup\bigl\{\gamma\in(0,1):\ (\mathrm{S}^\flat),\ (\mathrm{B}^\flat),\
(\mathrm{M}^\flat),\ (\mathrm{F}^\flat)\\
&\qquad\qquad\text{and the summation condition of the per-cube count in (h6)}\\
&\qquad\qquad\text{hold at every level $j$ with $g_j\le\gamma$}\bigr\}.
\end{aligned}
\end{equation}
Here $(\mathrm{S}^\flat)$ compares $\mathsf{Q}_j$, which is $(1+\beta_0)^{-2}$ times the gain per cube
of part (c) below, with four times $\mathsf{P}_j$,
$(\mathrm{B}^\flat)$ is the condition at the creation of a large-field region, $(\mathrm{M}^\flat)$ the
condition at joins and re-creations, and $(\mathrm{F}^\flat)$ states that the pre-summed factor of the
violator bonds is at most one per cell.
The \emph{modified smallness conditions} are the smallness conditions on $\gamma$ of (h6) with
$\gamma_0A_1^2\mathsf{T}_\gamma^{p_0}\ge 8$ replaced by $\gamma_0A_\flat^2\mathsf{T}_\gamma^{p_0}\ge 8$
and with the constants
$O(1)$ multiplying $|Z_j|$ and $d'_j(Z_j)$ replaced as in the second and third lines of
\eqref{5.2:eq-rim-payment-a}. The per-cube rate of the per-cube count of the large-field terms is
defined as the minimum of the per-cube rates of the present construction, and its positivity constant
as $\min\{\gamma_0,\lambda_0\}$; at $d=4$ this definition rests on the inequality
$256\Theta_1^2\ge\frac92\,6^d$, which holds
since $\Theta_1\ge d(L-1)/2\ge 8$ by (h2)(iii). Ba{\l}aban's coefficient $\kappa_1$ of the term
$-\kappa_1d_k(X)$ in the bound of the large-field terms is written $\kappa_{\mathfrak{d}}$, as in the
first line of \eqref{5.2:eq-rim-payment-a}.
Then:
\begin{enumerate}
\item[(a)] at every true pair, aged and inner backgrounds included, on the support of the refined
characteristic functions, the rim bound of (h2)(ii) holds;
\item[(b)] each factor $e^{-\gamma_0\mathfrak{b}^2/(64\Theta_1^2)}$ of that bound is the square of
$a:=e^{-\gamma_0\mathfrak{b}^2/(128\Theta_1^2)}$, and one of the two factors $a$ per bond of $I^*$
is summable over the index $I^*$:
\begin{equation*}
\sum_{I^*}\prod_{b\in I^*}a\le\exp\bigl(a\,n_{\mathrm{rim}}\bigr),
\end{equation*}
and this factor lies inside the factor $\exp O(1)(MR_j)^{-d}|Z_j|$ of \cite[p.~383]{B-CMP122b};
\item[(c)] the other half pays: in the display of \cite[p.~381]{B-CMP122b} the minimum gains the
member $A_\flat^2$ of \eqref{5.2:eq-rim-payment-a} per cube of $I^*$, and $A_1$ is replaced by
$A_\flat$ on its right-hand side;
\item[(d)] the cover in the fourth line of \eqref{5.2:eq-rim-payment-a} holds, and it is a cover of
$P_j\cup Q_j\cup R_j\cup\operatorname{cubes}(I_j)$;
\item[(e)] with $p:=\mathsf{T}_j^{p_0}$, the quarter $\frac14\gamma_0A_\flat^2p^2$ of the gain per
cube of $I^*$ read by the modified large-field bound is at least $2p$ exactly when
$\gamma_0A_\flat^2\mathsf{T}_j^{p_0}\ge 8$; this holds at every level with $g_j\le\gamma$ as soon as
$\gamma_0A_\flat^2\mathsf{T}_\gamma^{p_0}\ge 8$, the first member of the modified smallness
conditions, under which the modified large-field bound is applied;
\item[(f)] the conditions $(\mathrm{S}^\flat)$, $(\mathrm{B}^\flat)$, $(\mathrm{M}^\flat)$,
$(\mathrm{F}^\flat)$ of \eqref{5.2:eq-rim-payment-a}, which with the summation condition of (h6)
define $\gamma_P^\flat$, are ceilings on $\gamma$ chosen after $(L,M)$,
and at every level with $g_j\le\gamma<1$ the first three follow from
\begin{equation}\label{5.2:eq-rim-payment-1}
\begin{aligned}
(\mathrm{S}^\flat)&\Longleftarrow\ \mathsf{T}_\gamma^{p_0}
\ge\frac{1024(1+\beta_0)^2\Theta_1^2A_0}{\gamma_0A_1^2},\\
(\mathrm{B}^\flat)&\Longleftarrow\ \mathsf{T}_\gamma^{2p_0-(d+2)r_{\mathrm{pr}}}
\ge\frac{512\sqrt d(1+\beta_0)^2\Theta_1^2\,\Gamma(13+4c_0)C^\flat_{\mathfrak{d}}M^dL^{d+2}}
{\gamma_0A_1^2},\\
(\mathrm{M}^\flat)&\Longleftarrow\ \mathsf{T}_\gamma^{p_0-(d+2)r_{\mathrm{pr}}}
\ge\frac{(1+\beta_0)\Gamma\bigl(5\sigma+4\sqrt d\,101^d+4c_0+52\bigr)C^\flat_{\mathfrak{d}}M^dL^{d+2}}
{A_0},
\end{aligned}
\end{equation}
whose left members, read at $\mathsf{T}_j$ in place of $\mathsf{T}_\gamma$, are monotone in
$\mathsf{T}_j$;
\item[(g)] a component meeting $I_{j+1}$ is in Ba{\l}aban's class (ii), ``a new large field region
created'' \cite[p.~177]{B-CMP122a}.
\end{enumerate}
\end{proposition}

\begin{proof}
\emph{Part (a).} By (h1), $\mathcal{A}\ge\gamma_0$ on the whole of $\ell^2(\mathfrak{X};\mathfrak{g})$
at every true pair, aged and inner backgrounds included. By (h2)(i), since
$\|(y,x)\|^2=\|y\|^2+\|x\|^2$,
\begin{equation*}
\langle x,Sx\rangle=\min_y\langle (y,x),\mathcal{A}(y,x)\rangle
\ge\gamma_0\min_y\bigl(\|y\|^2+\|x\|^2\bigr)=\gamma_0\|x\|^2 ,
\end{equation*}
so $S\ge\gamma_0$. Since $S_{cc}$ is a compression of $S$, its quadratic form is the restriction of
that of $S$ to a subspace, and $S_{cc}\ge\gamma_0$. (The bound on $S$ uses the positivity of
$\mathcal{A}$ on all of $\ell^2(\mathfrak{X})$, not only on a compression.)
The hypotheses of the rim bound (h2)(ii) are thus met, and (a) follows. Besides the factors indexed
by $I^*$ displayed in (h2)(ii), the rim bound draws from the same first quarter
$\frac14\gamma_0\|x\|^2=\frac14\gamma_0\sum_b|x_b|^2$ of
$e^{-\frac12\langle x,Sx\rangle}=e^{-\frac14\langle x,Sx\rangle}e^{-\frac14\langle x,Sx\rangle}$
the factors attached to the cubes of $R$, which are not displayed there; the bonds witnessing the
cubes of $R$ and the bonds of $I^*$ are disjoint sets of coordinates of $x$, so no coordinate is
spent twice.

\emph{Part (b).} By the history-free indexing of (h3), the sum over terms at fixed $R$ is a sum over
the set $I^*$, each subset of the rim bonds occurring at most once. With $a$ as in (b),
\begin{equation*}
\sum_{I^*}\prod_{b\in I^*}a\le\sum_{I\subset\{\text{rim bonds}\}}a^{\#I}
=(1+a)^{n_{\mathrm{rim}}}\le e^{a\,n_{\mathrm{rim}}}.
\end{equation*}
Read at level $j$, condition $(\mathrm{F}^\flat)$ states that this pre-summed factor is at most one
per cell of side $LM\check{R}_j$, such a cell carrying at most $d(LM\check{R}_j)^d$ bonds; the
exponent therefore lies inside $O(1)(MR_j)^{-d}|Z_j|$ of \cite[p.~383]{B-CMP122b}.

\emph{Part (c).} The remaining factor $a$ per bond of $I^*$ gives
\begin{equation*}
\prod_{b\in I^*}a=\exp\Bigl(-\#I^*\,\frac{\gamma_0\mathfrak{b}^2}{128\,\Theta_1^2}\Bigr).
\end{equation*}
By (h3) each bond of $I^*$ contributes at most two cubes, so
$\#I^*\ge\frac12\#\operatorname{cubes}(I^*)$ and
\begin{equation*}
\prod_{b\in I^*}a\le\exp\Bigl(-\#\operatorname{cubes}(I^*)\,
\frac{\gamma_0\mathfrak{b}^2}{256\,\Theta_1^2}\Bigr).
\end{equation*}
At level $j$, $\mathfrak{b}_j^2=A_1^2\mathsf{T}_j^{2p_0}$, so the exponent per cube is
$\gamma_0A_1^2\mathsf{T}_j^{2p_0}/(256\Theta_1^2)=\gamma_0A_\flat^2\mathsf{T}_j^{2p_0}$. This is the
member $A_\flat^2$ gained per cube of $I^*$ by the minimum in the display of
\cite[p.~381]{B-CMP122b}, with $A_1$ replaced by $A_\flat$ on its right-hand side.

\emph{Part (d).} By (h4), $\Lambda^{*c}\subset\mathfrak{s}(R^*)^{+99}$. At the creating step $Z_j$
is the complement $\Lambda^{*c}$ of the refined domain. Apart from the padding
$(\cdot)^{+99}$, the right-hand side of the fourth line of \eqref{5.2:eq-rim-payment-a} is the
cover of \cite[(1.76)]{B-CMP122b} read at $R^*$ in place of $R$; its bracket contains
$\mathfrak{s}(R^*)$, so that (h4) gives the inclusion of the fourth line. It remains a cover because
$\mathfrak{s}(\cdot)$ is monotone and $R\subset R^*$ by (h3). Since
$\operatorname{cubes}(I_j)\subset R^*$, again by (h3), it covers
$P_j\cup Q_j\cup R_j\cup\operatorname{cubes}(I_j)$.

\emph{Part (e).} Of the gain delivered per cube of $I^*$, the modified large-field bound uses the
following splitting. Let $p:=\mathsf{T}_j^{p_0}>0$, the power of the logarithm of the inverse
coupling entering there. Since
$1-\theta_W\ge\frac34$ by (h6), we have
\begin{align*}
&(1-\theta_W)\gamma_0A_\flat^2p^2\ge\tfrac12\gamma_0A_\flat^2p^2+\tfrac14\gamma_0A_\flat^2p^2,\\
&\tfrac14\gamma_0A_\flat^2p^2\ge 2p\iff\gamma_0A_\flat^2p\ge 8 ,
\end{align*}
the equivalence following on division by $p/4>0$. For $\gamma<1$ and $g_j\le\gamma$ one has
$\mathsf{T}_j\ge\mathsf{T}_\gamma>0$, and $\mathsf{T}\mapsto\mathsf{T}^{p_0}$ is increasing on
$\mathsf{T}>0$; so the right-hand side of the equivalence follows from the condition in the first
line of \eqref{5.2:eq-rim-payment-a}. By part (c), the gain delivered per
cube of $I^*$ is $\gamma_0A_\flat^2\mathsf{T}_j^{2p_0}$, of which the modified bound uses one
quarter.

For the constants: by the second inequality of (h4), and since
\begin{equation*}
\#\mathfrak{s}(R)+2\cdot3^d\,\#\operatorname{cubes}(I^*)
\le 3^{d+1}\bigl(\#\mathfrak{s}(R)+\#\operatorname{cubes}(I^*)\bigr),
\end{equation*}
the volume constants acquire the factor $199^d3^{d+1}$. The comparison constant
$C_\sharp$ is that of \eqref{5.2:eq-comparison-constant-a}. The logarithm is the per-cell value
constant $\mathsf{K}_{\mathrm{cell}}(k,j)$, which the volume clause of (h6) places inside
Ba{\l}aban's term $O(1)\log g_j^{-2}|Z_j|$ of the modified bound. It is absorbed by the factor
$\check{R}_j^{\,d+1}$ of $c_j$, exactly as Ba{\l}aban's bound
$\le O(1)M^dR_j^{d+1}d'_j(Z_j)$ in \cite{B-CMP122b} absorbs it.

The constant $C^\flat_{\mathfrak{d}}$ enlarges $C_{\mathfrak{d}}$ so that it dominates every per-step
charge per unit of $d'$, the constants $O(1)$ of the smallness conditions, and $199^d3^{d+1}$. The
summand $2^{d+2}(1+dL^d)$ accounts for the term $1$ in the factor $d'_j(Z_j)+1$ of the per-step
charge of the summation of the large-field terms in (h6).

\emph{Part (f).} Note first that $\check{R}_j<L\mathsf{T}_j^{\,r_{\mathrm{pr}}}$. Indeed
$\check{R}_j=R(\min_{i\le j}x_i)$ is the least power of $L$ that is at least
$(\log\min_{i\le j}x_i)^{r_{\mathrm{pr}}}$, hence less than $L$ times that number, and
$\log\min_{i\le j}x_i\le\log x_j=\mathsf{T}_j$. Also $\check{R}_j\ge 1$. Fix a level with
$g_j\le\gamma<1$, so that $\mathsf{T}_j\ge\mathsf{T}_\gamma>0$.

$(\mathrm{S}^\flat)$: since
\begin{equation*}
\mathsf{Q}_j=\frac{\gamma_0A_1^2\mathsf{T}_j^{2p_0}}{256\,\Theta_1^2(1+\beta_0)^2},\qquad
4\mathsf{P}_j=4A_0\mathsf{T}_j^{p_0},
\end{equation*}
division by $\mathsf{T}_j^{p_0}$ gives
\begin{equation*}
(\mathrm{S}^\flat)\iff\mathsf{T}_j^{p_0}\ge\frac{1024(1+\beta_0)^2\Theta_1^2A_0}{\gamma_0A_1^2},
\end{equation*}
which holds when it holds at $\mathsf{T}_\gamma\le\mathsf{T}_j$.

$(\mathrm{B}^\flat)$, the line for the creation of a region: using $\check{R}_j\ge 1$ and then
$\check{R}_j<L\mathsf{T}_j^{\,r_{\mathrm{pr}}}$,
\begin{align*}
\Gamma c_j(9+4c_0+4\check{R}_j)&\le\Gamma(13+4c_0)C^\flat_{\mathfrak{d}}M^d\check{R}_j^{\,d+2}\\
&<\Gamma(13+4c_0)C^\flat_{\mathfrak{d}}M^dL^{d+2}\mathsf{T}_j^{(d+2)r_{\mathrm{pr}}},
\end{align*}
while
\begin{equation*}
\mathsf{D}_j=\frac{\gamma_0A_1^2\mathsf{T}_j^{2p_0}}{512\sqrt d\,\Theta_1^2(1+\beta_0)^2}.
\end{equation*}
Hence $(\mathrm{B}^\flat)$ holds as soon as
\begin{equation*}
\mathsf{T}_j^{2p_0-(d+2)r_{\mathrm{pr}}}\ge\frac{512\sqrt d(1+\beta_0)^2\Theta_1^2\,
\Gamma(13+4c_0)C^\flat_{\mathfrak{d}}M^dL^{d+2}}{\gamma_0A_1^2}.
\end{equation*}
Here the degree $2p_0$ is set against $(d+2)r_{\mathrm{pr}}$. By (h5),
$2p_0\ge 10r_{\mathrm{pr}}\ge(d+2)r_{\mathrm{pr}}$, so the left-hand side is non-decreasing in
$\mathsf{T}_j$. The condition therefore follows from its value at $\mathsf{T}_\gamma$, which is the
second line of \eqref{5.2:eq-rim-payment-1}.

$(\mathrm{M}^\flat)$, the line for joins and re-creations: it is a quantitative form of Ba{\l}aban's
requirement ``for $p_0$ large and $\gamma$ small enough'' in \cite{B-CMP122b}, together with
\cite[(1.88)]{B-CMP122b}.
Using $48+4\check{R}_j\le 52\check{R}_j$, and bounding each of the remaining summands by itself times
$\check{R}_j\ge 1$,
\begin{align*}
&\Gamma c_j\bigl(5\sigma+4\sqrt d\,101^d+4c_0+4\check{R}_j+48\bigr)\\
&\qquad<\Gamma\bigl(5\sigma+4\sqrt d\,101^d+4c_0+52\bigr)C^\flat_{\mathfrak{d}}M^dL^{d+2}
\mathsf{T}_j^{(d+2)r_{\mathrm{pr}}},
\end{align*}
while $(1+\beta_0)^{-1}\mathsf{P}_j=(1+\beta_0)^{-1}A_0\mathsf{T}_j^{p_0}$. Hence $(\mathrm{M}^\flat)$
holds as soon as the third line of \eqref{5.2:eq-rim-payment-1} holds with $\mathsf{T}_j$ in place of
$\mathsf{T}_\gamma$. Here the degree $p_0$ is set against $(d+2)r_{\mathrm{pr}}$, and
$p_0\ge 5r_{\mathrm{pr}}$ does not suffice at $d=4$. The pair of (h5) does: from $p_1<p_0$ and
$2p_1-(d+5)r_0>p_0$,
\begin{equation*}
2p_0>2p_1>p_0+(d+5)r_0,
\end{equation*}
so $p_0>(d+5)r_0=9r_0$ at $d=4$, and with $r_0\ge r_{\mathrm{pr}}$ from (h5),
\begin{equation*}
p_0>9r_0\ge 9r_{\mathrm{pr}}\ge 6r_{\mathrm{pr}}=(d+2)r_{\mathrm{pr}}.
\end{equation*}
No condition on the exponents beyond (h5) is introduced. Then $p_0-(d+2)r_{\mathrm{pr}}>0$, the
left-hand side is increasing in $\mathsf{T}_j$, and the third line of \eqref{5.2:eq-rim-payment-1}
suffices.

Each of $(\mathrm{S}^\flat)$, $(\mathrm{B}^\flat)$, $(\mathrm{M}^\flat)$, $(\mathrm{F}^\flat)$ is thus a
ceiling on $\gamma$ whose constants depend on $(L,M)$ and the fixed numbers $d$, $\gamma_0$,
$\Theta_1$, $\beta_0$, $A_0$, $A_1$, $C^\flat_{\mathfrak{d}}$, and the left members of the three
closed forms of \eqref{5.2:eq-rim-payment-1}, read at $\mathsf{T}_j$, are monotone in
$\mathsf{T}_j$. Together with the summation condition of the per-cube count of the large-field terms
in (h6) they are the five conditions of the last lines of \eqref{5.2:eq-rim-payment-a}, which define
$\gamma_P^\flat$. Finally, $\Theta_1\ge d(L-1)/2\ge 8$ by (h2)(iii), which gives
\begin{equation*}
256\,\Theta_1^2\ge 256\cdot 64=16384\ge 5832=\tfrac92\,6^4,
\end{equation*}
which is the inequality used at $d=4$ for the positivity constant of the per-cube rate.

\emph{Part (g).} A component of the large-field region that meets $I_{j+1}$ contains cubes newly
added at that step by the refinement (h3). It is therefore in Ba{\l}aban's class (ii), ``a new large
field region created'' \cite[p.~177]{B-CMP122a}.
\end{proof}

\begin{proposition}[Preparatory integrations with both bracket slots]
\label{5.2:prop-two-slot-preparation}
Let $k$ be a level and $n$ a stage of the preparatory integrations
\cite[(1.60)--(1.63)]{B-CMP122a}, and let $\mathbb{C}^{(n+1)}_k$ denote the chain terms of stage
$n+1$. Assume three statements: the two-slot corollary of the $A$-slot, part (a) of the
inheritance theorem of the $A$-slot, and the source corollary of the $A$-slot. The first asserts
conclusion (i) below under the three conditions in the first sentence of hypothesis~(1) and the
created-term clause \eqref{5.2:eq-created-terms-bound-a}. The second asserts conclusion (ii)
below under that clause and the conditions of hypothesis~(2). The third asserts conclusion (iii)
below, given conclusion (ii), under the two conditions of hypothesis~(3). Assume moreover the
following.
\begin{enumerate}
\item[(1)] The inductive hypothesis of the $A$-slot, the flat kernel condition of the $A$-slot,
read with the augmented standing condition of the $A$-slot, and the standing condition of the
$A$-slot hold. The
hypotheses of Corollary~\ref{5.2:cor-created-bound-induction} hold at every step: those of the
theorem on created boundary terms, namely the analyticity lemma for the rim potentials, parts
(a)--(d), at the radius $T^{\mathbb{C}}$; positivity and decay at every true pair,
\eqref{5.2:eq-true-pair-positivity-a}; the random-walk input and the far-bond bound of
Definition~\ref{5.1:def-admitted-inputs}; the bounds on
the untilted expansion taken from \cite{B-CMP116}, among them
\cite[(2.37)--(2.41)]{B-CMP116}; and the one-sided conditions on the rim
constants, \eqref{5.2:eq-rim-conditions-a}; together with part (a) of the proposition on the
sourced step at the earlier steps of the one induction over levels.
\item[(2)] The remaining hypotheses of the inheritance theorem of the $A$-slot, part (a), hold:
\begin{itemize}
\item the third inductive clause of the preparatory-step analysis, at the fraction
$\beta_{\mathrm{pr}}$, restricted to the set $Z^{\mathrm{on}}$ of locations at which that analysis
establishes it;
\item the room condition for the terms created at level $k$, with fraction
$\mathsf{c}_e := \theta_S \in [1/2, 8/9]$;
\item the results of \cite[\S\S3--5]{B-CMP109}, taken by statement, for the near pairs of
$\mathbb{E}$- and $\mathbb{R}$-terms;
\item the ambient hypothesis of that theorem;
\item the natural kernel condition of the $A$-slot at window width $\mathsf{w}_0 = 1/24$ and
fraction $1/2$, and the augmented standing condition of the $A$-slot also in its version at
window width $\mathsf{w}_0$ and fraction $1/2$, that is, with $K'$ replaced by $24K'$, where $K'$
is the last step of the local run.
\end{itemize}
\item[(3)] The two source conditions of the source corollary of the $A$-slot hold.
\end{enumerate}
Then the following hold.
\begin{enumerate}
\item[(i)] The chain terms decompose as
\begin{equation}\label{5.2:eq-two-slot-preparation-1}
\mathbb{C}^{(n+1)}_k = \sum_{(X,\mathcal{S})} \mathbb{C}^{(n+1)}_k(X,\mathcal{S}),
\qquad
\|\mathbb{C}^{(n+1)}_k\|^{A}_{(n+1)} \le C_0^{\sharp}.
\end{equation}
Each summand is a stored term with datum $A^{\flat}$ in the sense of
Definition~\ref{5.2:def-stored-terms}. The norm $\|\cdot\|^{A}_{(n+1)}$ is taken with the rates
$a_k = (1+3\beta_{\mathrm{pr}})\kappa$ and $a_m = (1-\beta_{\mathrm{pr}})\kappa$ (the two fixed
rates of the norm $\|\cdot\|^{A}$; the subscripts are names, not the level or a layer index).
Every constant is free of the stage $n$, of the number $N(g)$ of preparatory stages and the
number $N_0(g)$ of zones, of $\check{R}_k$, $k$, $K$, the volume and $g$. All of this holds at the
cores $\mathsf{C}_C$ of the components of Ba{\l}aban's construction.
\item[(ii)] The decomposition and the bound \eqref{5.2:eq-two-slot-preparation-1}, with the same
$C_0^{\sharp}$ and the same rates, hold with each stored term
$\mathbb{C}^{(n+1)}_k(X,\mathcal{S})$ taken on the inherited domain
$\mathfrak{D}^{\flat}_{n+1}(X)\times\mathfrak{D}_{\mathcal{S}}(X)$.
\item[(iii)] The same holds with each stored term $\mathbb{C}^{(n+1)}_k(X,\mathcal{S})$ taken on
$\mathfrak{D}^{\flat}_{n+1}(X)\times\mathfrak{D}_{\mathcal{S}}(X)\times\mathcal{W}_k(X)$. Here
$\mathfrak{D}^{\flat}_{n+1}(X)$ does not depend on the source $w$.
\end{enumerate}
\end{proposition}

\begin{proof}
\emph{Step 1: the structure of the integrand at the site $(\mathrm{P})$.}

In \cite[(1.60)]{B-CMP122a} the index $j$ labels the step whose rim field $A_j$ is integrated; we
keep it and the letters $g_j$, $B_j$, $\chi^{(j)}$, $C$, $\Lambda'_{j+1}$, $h$, $\tilde{D}$,
$V^{(j)}$, $U_k^{(n)}$, $U_k^{(n+1)}$ of that display.
We describe the integrand of \cite[p.~188, (1.60)]{B-CMP122a}, that is, the integrand at the
preparatory site $(\mathrm{P})$ of \eqref{5.2:eq-three-sites-a}. The rim field $A_j$ is an
integration variable on $Z_{j+1}\cap\Omega_{j+1}$. It is weighted by the cut-off $\chi^{(j)}$ of
\cite[(1.9)]{B-CMP122a} and by the Gaussian
\begin{equation*}
\exp\Big(-\tfrac12\big\langle A_j, C^{*}\Delta^{(j)}_{\Lambda'_{j+1}} C A_j\big\rangle\Big).
\end{equation*}
The boundary terms enter the exponent of the $B_j$-integral in two ways. The first is through
both slots, the $B_j$-slot and the $A_j$-slot, of the bracket below, written in the notation of
\cite[(1.60)]{B-CMP122a}:
\begin{equation}\label{5.2:eq-two-slot-preparation-2}
\begin{split}
\Big\{&(\mathbb{E}_k+\mathbb{R}_k+\mathbb{B}_k+\mathbb{B}_k^{(n)})
\big(U_k^{(n)}(\exp i[g_jCB_j - h\tilde{D}(\cdot)]V^{(j)})\big)\\
&-(\ldots)\big(U_k^{(n+1)}, A_j = 0\big)\Big\}.
\end{split}
\end{equation}
The second is through the term $\mathbb{P}^{(j)}(g_j,A_j,B_j)$ of \cite[(1.60)]{B-CMP122a}.

No bilinear $A_j$--$B_j$ term of order one is displayed in \cite[(1.60)]{B-CMP122a}. Its absence
is structural, for the following reasons.
\begin{itemize}
\item By \cite[(2.23)]{B-CMP119}, the effective action, denoted $A_k(X,U_k)$ in
\cite{B-CMP119} (not to be confused with the fluctuation fields $A_j$), depends on the gauge
field variable $V$ through the background field $U_k(V)$. Writing $\mathbf{A}$ for the system of
fluctuation fields $\{A_j\}$, it has the general form
\begin{equation*}
A_k(X,U_k) = -A(1/g_k^2,U_k)+E_k(U_k)+R_k(U_k)+B_k(U_k,\mathbf{A})-E_k .
\end{equation*}
So the fluctuation fields enter only the $B$-terms. The terms $E_k$, $R_k$ and the Wilson term
depend on $U_k(V)$, which is a function of $V$ alone.
\item By \cite[(2.21)]{B-CMP119}, the quadratic form in the exponential couples only fields
$A_j$ in the same component of $Z_{j+1}\cap\Omega_{j+1}$. The other terms of this quadratic
form are included in the effective action as $B$-terms. This quadratic form is the one kept in
the Gaussian displayed above.
\item The variable $B_j$ of \cite[(1.60)]{B-CMP122a} comes from the group variables $V_j$ on
$\Omega^{c}_{j+1}\setminus Z''_{j+1}$ through the critical background field. By
\cite{B-CMP122a}, the effective action expressed through this background has the form
\cite[(2.23)]{B-CMP119}, with some new boundary terms only.
\end{itemize}
Hence the only joint dependence on $(A_j,B_j)$ in the exponent of
\cite[(1.60)]{B-CMP122a} is through the bracket \eqref{5.2:eq-two-slot-preparation-2} and
$\mathbb{P}^{(j)}(g_j,A_j,B_j)$. These are exactly the objects of the two-slot clause of the
$A$-slot.

Consequently no conditional mean of either variable depends on the sharp small-field box of the
other, that is, on the region cut out by its characteristic-function cut-off, and there is no
tilt at the site $(\mathrm{P})$. The cut-off \cite[(1.9)]{B-CMP122a} is
therefore read against an unshifted Gaussian. For this reading of the cut-off two facts suffice:
\begin{itemize}
\item the lower bound $\mathfrak{b}^{\flat}_j \ge P_0$ on the small-field box, fixed in the rim
construction, where $P_0$ denotes the floor under the thresholds in
\cite[pp.~188--189]{B-CMP122a};
\item the replacement of \cite[(1.9)]{B-CMP122a} by its form
\eqref{5.2:eq-second-stage-activities-a} on the collar.
\end{itemize}
That no bilinear term occurs thus follows from the structure \cite[(2.21), (2.23)]{B-CMP119}.

\emph{Step 2: the second slot must be carried.}

The bracket \eqref{5.2:eq-two-slot-preparation-2} cannot be read as the difference of the values
of one function of $B_j$ at $B_j$ and at $0$, with $A_j$ the same in both terms: in
\eqref{5.2:eq-two-slot-preparation-2} the subtracted term is evaluated at $A_j = 0$, and that
reading drops this second slot. Nor can the second slot be paid for by a supremum bound on
$\mathbf{B}$. A
boundary term $\mathbf{B}^{(\ell)}(X)$ with $\ell - j \gtrsim N_0(g)$ depends on $A_j$ on a number
of collar cubes of order $R_{j+1}^{d}$. Each of these cubes is charged a factor $e^{c_F}$ by
\cite[(2.26)]{B-CMP116}, so that the second-stage estimate acquires a factor $e^{c_FV_0}$ per
link in the $A$-slot, $V_0$ denoting the number of collar cubes on which one link reads $A_j$.

\emph{Step 3: the two-slot expansion.}

We run \cite[(1.60)--(1.63)]{B-CMP122a} under the two-slot clause of the $A$-slot. Let $c$ be a
term of $\mathbb{B}_k$ or a chain term, with datum $(c,S,\mathcal{S})$. Its bracket is expanded by
M\"obius inversion jointly over two families of cubes.
\begin{itemize}
\item The response cubes $\square$. Their parameters are $t_{\square}$, of radius
$r_{\square} = c_M/\vartheta_0(\square)$. Here $\vartheta_0(\square)$ denotes the response bound
of the block $\square$ in the all-blocks response dilation
\eqref{5.2:eq-response-dilation-a}.
\item The weak $A_j$-cubes $q$ meeting $S$. Their parameters are $u_q$, of radius $r_A$.
\end{itemize}
On far weak cubes, where $|A_j| \le C_{\mathrm{far}}\mathfrak{b}$, the radius is
$r_A/C_{\mathrm{far}}$. This is conservative beside the bound $\mathfrak{b}$ of the rim
construction on the outer layer, and it is paid for by the factor
$C_{\mathrm{far}}$ in the choice of $\gamma_J$. Put $A_j^{u} := u_qA_j$ on weak cubes $q$ and
$A_j^{u} := A_j$ on strong cubes.

By the two-slot expansion lemma, each element $e$ of the expansion of $c$, with value $F$,
depends on $(\mathbb{U},\mathbb{J})$ only on its domain $\mathsf{d}(e)$, and on the rim only on
its rim set $\mathsf{b}(e)$; these are the sets of that lemma. Let $\mathcal{F}_{\mathrm{r}}$ be a
finite set of response cubes and $\mathcal{F}_{\mathrm{w}}$ a finite set of weak cubes, and write
$\partial_{t_{\mathcal{F}_{\mathrm{r}}}}$ and $\partial_{u_{\mathcal{F}_{\mathrm{w}}}}$ for the
products of the derivatives in $t_{\square}$, $\square \in \mathcal{F}_{\mathrm{r}}$, and in
$u_q$, $q \in \mathcal{F}_{\mathrm{w}}$. Then
\begin{equation}\label{5.2:eq-two-slot-preparation-3}
\big|\partial_{t_{\mathcal{F}_{\mathrm{r}}}}\partial_{u_{\mathcal{F}_{\mathrm{w}}}}F\big|
\le \mathfrak{N}(c)\prod_{\square\in\mathcal{F}_{\mathrm{r}}}\frac{2\vartheta_0(\square)}{c_M}
\prod_{q\in\mathcal{F}_{\mathrm{w}}}\frac{2}{r_A}.
\end{equation}

Strong cubes are not separated. The second stage charges each of them a factor $e^{c_{\flat}}$.
This charge is paid for by the strong weight $e^{\kappa_A(k)|\mathcal{S}_j(c)|}$ carried in the
norm. By the strong-weight step of the two-slot theorem, what remains is
\begin{equation*}
e^{-(\kappa_A-c_{\flat})|\mathcal{S}_j(c)|}\le \varepsilon_S/2 ,
\end{equation*}
where $\varepsilon_S$ is the smallness constant of the strong-set sum.

The proposition of the second stage in norm form sums the elements. They are summed in the
groups $\hat e$ of the grouping lemma, with its constants $K_2^{\flat}$, $K_{\flat}''$. The
wrapped elements, namely those whose creating cluster
contained a relative potential, are read on $\mathsf{d}^{\mathrm{rel}}(\hat e)$.

The holomorphic-pieces lemma and its corollary treat pieces holomorphic in an auxiliary
parameter under majorants free of that parameter. They carry the exterior slots $A_{j'}$,
$j' > j$, as
polydisc or box parameters, and they file the outputs under unions of strong sets. The result is
the two-slot theorem of the $A$-slot, and, at the cores $\mathsf{C}_C$ of the components of
Ba{\l}aban's construction, the two-slot corollary.

\emph{Step 4: conclusion (i).}

The two-slot corollary has four hypotheses:
\begin{itemize}
\item the inductive hypothesis, the flat kernel condition and the standing condition of
the $A$-slot, all of which are assumed in hypothesis~(1);
\item the created-term clause, taken to be \eqref{5.2:eq-created-terms-bound-a}.
\end{itemize}
The last is supplied by Corollary~\ref{5.2:cor-created-bound-induction}. The hypotheses of that
corollary are those listed in hypothesis~(1).

The two-slot corollary therefore gives the decomposition \eqref{5.2:eq-two-slot-preparation-1}
into stored terms with datum $A^{\flat}$, with $\|\mathbb{C}^{(n+1)}_k\|^{A}_{(n+1)} \le C_0^{\sharp}$
and the rates $a_k = (1+3\beta_{\mathrm{pr}})\kappa$, $a_m = (1-\beta_{\mathrm{pr}})\kappa$. Every
constant is free of the stage $n$, of the number $N(g)$ of preparatory stages and the number
$N_0(g)$ of zones, of $\check{R}_k$, $k$, $K$, the volume and $g$. In particular
the constants $K_2^{\flat}$, $K_{\flat}''$ of the grouping in Step~3 are free of these
quantities. This is (i).

\emph{Step 5: conclusion (ii).}

We apply part (a) of the inheritance theorem of the $A$-slot. Its hypotheses are supplied as
follows.
\begin{itemize}
\item The clause \eqref{5.2:eq-created-terms-bound-a} is Corollary~\ref{5.2:cor-created-bound-induction},
as in Step~4.
\item The third inductive clause of the preparatory-step analysis, at $\beta_{\mathrm{pr}}$, is
assumed in hypothesis~(2), restricted to the locations of $Z^{\mathrm{on}}$.
\item The room condition for the terms created at level $k$, with
$\mathsf{c}_e := \theta_S \in [1/2,8/9]$, is supplied as follows.
\begin{itemize}
\item For $\mathbb{E}^{(k)}$ and $\mathbb{R}^{(k)}$ it follows from the radius comparison of the
preparatory-step analysis,
\begin{equation*}
(1+\beta_{\mathrm{pr}})\alpha_{\cdot,k+1}\le(1+\beta_s)\alpha_{\cdot,k+1}\le\hat\alpha_k ,
\end{equation*}
which gives room at least $\beta_{\mathrm{pr}}$ on the top layer.
\item For $\mathbb{B}^{(k)}$ it follows from part (b) of the theorem on created boundary
terms, whose analyticity space, read layer
by layer, is the layer schedule with radius factor $1+\beta_{\mathrm{pr}}$ on the
top layer and $s^{(k)}_m+\theta_S\beta_{\mathrm{pr}}2^{-(k-m)}$ on the layers $m<k$, where
$s^{(k)}_m$ denotes the radius factor of layer $m$ at level $k$ before the room
$\theta_S\beta_{\mathrm{pr}}2^{-(k-m)}$ is added.
\end{itemize}
\item The results of \cite[\S\S3--5]{B-CMP109} for the near pairs of $\mathbb{E}$- and
$\mathbb{R}$-terms are used by statement. They are re-run as displayed in the inheritance
theorem, with the parameter $\ell_T$ of that re-run taken equal to the layer index $m$.
\item The ambient hypothesis, and the natural kernel condition of the $A$-slot at window width
$\mathsf{w}_0$ and fraction $1/2$, are assumed in hypothesis~(2). The whole set of conditions of
the latter enters the choice of $\gamma_J$ and the order of choice of the constants
(Definition~\ref{5.1:def-admitted-inputs}).
\end{itemize}
The inheritance theorem then gives the same stored terms, with the same $C_0^{\sharp}$ and the
same rates, on the inherited domain $\mathfrak{D}^{\flat}_{n+1}(X)\times\mathfrak{D}_{\mathcal{S}}(X)$.
This is (ii).

In this application the augmented standing condition of the $A$-slot is used in the version at
window width $\mathsf{w}_0$ and fraction $1/2$ assumed in the last item of hypothesis~(2).

\emph{Step 6: conclusion (iii).}

The source corollary of the $A$-slot applies under its two source conditions. These are assumed
in hypothesis~(3); they are supplied by part (a) of the proposition on the sourced step. The
corollary extends (ii) to
$\mathfrak{D}^{\flat}_{n+1}(X)\times\mathfrak{D}_{\mathcal{S}}(X)\times\mathcal{W}_k(X)$, and the
domain $\mathfrak{D}^{\flat}_{n+1}(X)$ is free of $w$. This is (iii).
\end{proof}

\begin{remark}
\emph{Use at the other sites.}
Conclusions (ii) and (iii) are what the containment at the localisation site $(\mathrm{L})$ of
\eqref{5.2:eq-three-sites-a} uses.
\begin{itemize}
\item For the presented pair, this is through the containment corollary and its claim at
$(\mathrm{L})$.
\item For the live slots, it is through the disc radii of \eqref{5.2:eq-three-sites-a}.
\end{itemize}
At the site $(\mathrm{P})$ the same units are read together with the gauge-map datum of that site
as orbit data, by the orbit-datum corollary at $(\mathrm{P})$.

\emph{What the correlation theorem uses.}
The correlation theorem, Theorem~\ref{5.1:thm-correlation-theorem}, uses this proposition in the
following places:
\begin{itemize}
\item the half of Lemma~\ref{5.2:lem-second-stage-activities} that concerns the site
$(\mathrm{P})$;
\item the term of its estimate that is summed in Theorem~\ref{5.4:thm-finer-zone-sum};
\item the sourced bound on the small-field density;
\item part (a) of the proposition on the sourced step;
\item the containment at the site $(\mathrm{L})$ for the units of that term.
\end{itemize}
Through parts (b)--(e) of the proposition on the sourced step, the opening display of the
section containing that proposition, the modification step and the induction over levels, it
uses only sizes and analyticity in $w$. Both are implied by
$\|\cdot\|^{A} \le C_0^{\sharp}$, since $0 \in \mathfrak{D}_{\mathcal{S}}$.

In particular, no factor $e^{c_FV_0}$ per link in the $A$-slot arises.

\emph{Robustness of the later readings of the rim function.}
Of the later readings of the rim function in \cite{B-CMP122a, B-CMP122b}, all but two are treated
in the robustness lemma. The two remaining readings are treated as follows.

First, the reading of \cite[pp.~188--189, (1.60)]{B-CMP122a} uses a floor $P_0$ under the
thresholds and the absence of a bilinear term. It is served by $\mathfrak{b}^{\flat}_j \ge P_0$
and by \eqref{5.2:eq-second-stage-activities-a}, with the second slot now carried as in
Steps~1--3 of the proof above.

Second, the reading on \cite[pp.~381, 384--387]{B-CMP122b} involves
\cite[(1.76), (1.80)--(1.89)]{B-CMP122b} and the new-region factor in
\cite[(1.85)]{B-CMP122b}. It is served by the three bounds collected in
\eqref{5.2:eq-rim-payment-a}: the payment $A_{\flat}$ per cube, the variant of
\cite[(1.76)]{B-CMP122b} displayed there, and the strip bound displayed there; with
\cite[(1.83)]{B-CMP122b} at $I=\emptyset$ taken verbatim as arithmetic. On the family of
modified covers it is closed by the refined-cover lemmas, which give that
\cite[(1.80)]{B-CMP122b} holds as it stands, that the quantity those lemmas denote by $\kappa$
(not the rate $\kappa$ of the norm $\|\cdot\|^{A}$) is non-negative at every step, and the
factor $-\kappa_{\mathfrak{d}}d_k$ at the triples $(\mathsf{D},\mathsf{P},c)$ of those lemmas.
These lemmas use neither \cite[(1.76), (1.84)]{B-CMP122b}, nor the connectedness of the
graph $G$ of \cite{B-CMP122b}, nor the adjacent passage of \cite{B-CMP122b}. Finally,
\cite[(1.83)]{B-CMP122b} at $I=\emptyset$ holds by the third geometric clause of the
refined-cover lemmas.
\end{remark}

Throughout, we work at stage $n+1$ of the operation $\mathbf{R}_k$, at level $j=h+n$, in the setting of Proposition~\ref{5.2:prop-two-slot-preparation}, which is proved as an implication from the statements named there. Blocks are denoted $\square$, and $\tilde{\square}^{4}$ is the fourfold fattening of $\square$. $A$ is a finite set of blocks, and $Q$ a set of weak fluctuation cubes, as in that proposition.

The objects of the preparatory construction named below without definition are those of the setting of Proposition~\ref{5.2:prop-two-slot-preparation}. These are stored terms, plain or wrapped; window blocks, old blocks and exceptional blocks; arrested attempts, live arrests and their raised strata; near groups, far pieces, $\mathbb{P}$-pieces and Wilson pieces. Near units, read on their own or on their box, and cell units of the Wilson class are understood as in the data-weight lemma, whose statement is quoted below.

The set $\sigma_0(A)$ consists of the cells not contained in $\bigcup_{\square\in A}\tilde{\square}^{4}$. To every $\Delta\in\sigma_0(A)$ a decoupling parameter $\mathsf{s}(\Delta)$ is attached, and we put
\begin{equation*}
\mathfrak{P}_A:=\bigl\{\mathsf{s}:\ |\mathsf{s}(\Delta)|\le e^{\kappa_1}\ \text{for all}\ \Delta\in\sigma_0(A)\bigr\}.
\end{equation*}
For a function $f$ holomorphic on $\mathfrak{P}_A$, let $\varepsilon^1_\Delta$ and $\varepsilon^0_\Delta$ denote the evaluations of the variable $\mathsf{s}(\Delta)$ at $1$ and at $0$. For $\mathsf{y}\subset\sigma_0(A)$ we write
\begin{equation}\label{5.2:eq-animal-resummation-1}
f_{\mathsf{y}}:=\prod_{\Delta\in\mathsf{y}}\bigl(\varepsilon^1_\Delta-\varepsilon^0_\Delta\bigr)
\prod_{\Delta\in\sigma_0(A)\setminus\mathsf{y}}\varepsilon^0_\Delta\, f .
\end{equation}
The elements of the two-slot expansion of Proposition~\ref{5.2:prop-two-slot-preparation} are indexed by $e=(c,A,Q,\mathsf{y})$ with $\mathsf{y}\subset\sigma_0(A)$. At fixed $(c,A,Q)$, the element with index $\mathsf{y}$ is $f_{\mathsf{y}}$ for the appropriate function $f$.

The sum over $\mathsf{y}$ cannot be bounded cell by cell, for the following reason. By \cite[(1.10), p.~179]{B-CMP122a}, a preparation leaves between $Z''_k$ and $Z''_{k_0+1}$ zones that are one cube of their own scale thick. Successive sets are separated by one layer of $M$-cubes, then one layer of $L^{-1}M$-cubes, and so on.

Consider a block $\square$ of zone $k-1$ touching $Z''_{k-1}$. Its fattening $\tilde{\square}^{4}$ crosses all these zones and ends inside zone $k_0+1$, whose cells have side $L^{-(N_0-2)}$. Let $\mathcal{N}$ be the number of cells of $\sigma_0(A)$ that share a wall with the far face of $\tilde{\square}^{4}$. Then
\begin{equation*}
\mathcal{N}\ge(10\check{R}_{k_0+1}/L)^{d-1}.
\end{equation*}
Every subset of these $\mathcal{N}$ cells is admissible in the two-slot expansion. Hence, at fixed $(c,A,Q)$ with $\square\in A$, the sum over $\mathsf{y}$ of the cell-by-cell majorants is at least $(1+e^{-(\kappa_1-1)})^{\mathcal{N}}$.

A bound of this sum by a fixed factor per block of $A$ (the factor $4^d$ used in summing the elements) would therefore need
\begin{equation*}
e^{-(\kappa_1-1)}\bigl(10\check{R}_{k_0+1}/L\bigr)^{d-1}\le\mathrm{const}.
\end{equation*}
Since $\kappa_1$ is fixed before $\gamma$ and $\check{R}$ tends to infinity as $g\downarrow0$, this would bound the coupling from below, for plain units as well. No such bound is assumed anywhere in what follows. The obstruction belongs to the cell-by-cell majorants, not to the functions themselves. We therefore group the cells of $\mathsf{y}$ by the members of a coarser partition, the grain.

\smallskip

\emph{The grain of stage $n+1$.} Let $\pi^{\flat}$ be the family consisting of:
\begin{itemize}
\item[$(\flat\mathrm{a})$] the cubes of $\pi_k$ not contained in $Z$;
\item[$(\flat\mathrm{b})$] the cubes of $\pi_{j+1}$ contained in $Z\cap\Omega^c_{j+1}$;
\item[$(\flat\mathrm{c})$] the cells of $\mathbb{B}^{(n)}_k$ contained in $Z\cap\Omega_{j+1}$.
\end{itemize}
On $Z$ these are the domains described in \cite{B-CMP122a}. There the localisation domains $X$ are not only from $\mathbf{D}_{j+1}$ but of a more general type: $X\cap\Omega^c_{j+2}\in\mathbf{D}_{j+1}$, and $X\cap(\Omega_m\setminus\Omega_{m+1})\in\mathbf{D}_m$ for $m=j+2,\dots,k$.

The grain depends neither on the sources nor on the fields; it is a datum of the scaffold of the stage. For a block $\square$ and a set $A$ of blocks put
\begin{equation*}
\hat{\mathsf{K}}_{\square}:=\bigcup\{G\in\pi^{\flat}:\ G\cap\tilde{\square}^{4}\neq\emptyset\},\qquad
\hat{\mathsf{K}}_A:=\bigcup_{\square\in A}\hat{\mathsf{K}}_{\square},
\end{equation*}
\begin{equation*}
\hat\sigma_0(A):=\{G\in\pi^{\flat}:\ G\not\subset\hat{\mathsf{K}}_A\},\qquad N^{\flat}:=(10L+2)^d .
\end{equation*}

We use the following properties of the grain, which are the content of the grain lemma. We write $\operatorname{lev}(G)$ for the level of a member $G$.
\begin{enumerate}
\item[(a)] $\pi^{\flat}$ is a partition into closed cubes, each of which is a union of cells.
\item[(b)] $j+1\le\operatorname{lev}(G)\le k$ for every member $G$. Touching members differ in level by at most one, and a member touches at most $\hat b_t$ members.
\item[(c)] Every fluctuation cube $q\in\mathsf{Q}$ is a member.
\item[(d)] For every block $\square$ of any zone, $\hat{\mathsf{K}}_{\square}\supset\tilde{\square}^{4}$. Moreover $\hat{\mathsf{K}}_{\square}$ is wall-connected, has at most $N^{\flat}$ members, and lies within $6$ level-$k$ units of $\square$.
\item[(e)] Let $\mathsf{D}$ be a connected set meeting $\bigcup\mathsf{Q}$. Adding to $\mathsf{D}$ members of $\pi^{\flat}$ that meet it changes neither its filing level $m(\mathsf{D})$ nor its hulls $[\mathsf{D}]_{m'}$, $m'\ge m(\mathsf{D})$.
\end{enumerate}
No $N$, $N_0$, $\check{R}$, age of a region or number of windows enters these properties. Every window has levels at most $j$ (those of earlier attempts at most $k_a\le h$) and is contained in the part $(\flat\mathrm{b})$. The zones of level greater than $j$ are produced by the operation $\mathbf{T}$ and are at least $8\check{R}$ units thick, measured at the scale next coarser than their own.

By property~(a), every cell of $\sigma_0(A)$ lies either in $\hat{\mathsf{K}}_A$ or in exactly one member of $\hat\sigma_0(A)$. By property~(d), every cell of $\bigcup_{\square\in A}\tilde{\square}^4$ lies in $\hat{\mathsf{K}}_A$. A set $\hat{\mathsf{y}}\subset\hat\sigma_0(A)$ is called \emph{admissible} if every wall-component of $\hat{\mathsf{y}}$ shares a wall with $\hat{\mathsf{K}}_A$. We write $\operatorname{dist}_k$ for distance measured in level-$k$ units.

\begin{lemma}[Resummation of preparatory animals over grain members]\label{5.2:lem-animal-resummation}
Let $f$ be holomorphic on the polydisc $\mathfrak{P}_A$ with $|f|\le\mathfrak{n}$ there. Assume that $f$ reads $\mathsf{s}$ only through the decoupled operators of the localisation at complex source, restricted to $\bigcup_{\square\in A}\tilde{\square}^{4}$. For $\mathsf{y}\subset\sigma_0(A)$ let $\hat{\mathsf{y}}(\mathsf{y})$ be the set of members of $\hat\sigma_0(A)$ holding a cell of $\mathsf{y}$. For $\hat{\mathsf{y}}\subset\hat\sigma_0(A)$ put
\begin{equation*}
f_{\hat{\mathsf{y}}}:=\sum_{\mathsf{y}\subset\sigma_0(A):\ \hat{\mathsf{y}}(\mathsf{y})=\hat{\mathsf{y}}}f_{\mathsf{y}},
\end{equation*}
so that the members of $\hat\sigma_0(A)$ holding a cell of $\mathsf{y}$ are exactly those of $\hat{\mathsf{y}}$.
For $s=(s(G))_{G\in\hat\sigma_0(A)}$ define
\begin{equation}\label{5.2:eq-animal-resummation-2}
\mathsf{s}_s(\Delta):=
\begin{cases}
s(G), & \Delta\subset G\in\hat\sigma_0(A),\\
1, & \Delta\subset\hat{\mathsf{K}}_A .
\end{cases}
\end{equation}
Then the following hold.
\begin{enumerate}
\item[(i)] With $s(G)=0$ for $G\notin\hat{\mathsf{y}}$,
\begin{equation}\label{5.2:eq-animal-resummation-3}
f_{\hat{\mathsf{y}}}=\int_{[0,1]^{\hat{\mathsf{y}}}}\prod_{G\in\hat{\mathsf{y}}}ds(G)\,
\partial_{s(G)}\,f(\mathsf{s}_s).
\end{equation}
Moreover,
\begin{equation*}
\sum_{\hat{\mathsf{y}}}f_{\hat{\mathsf{y}}}=\sum_{\mathsf{y}}f_{\mathsf{y}}=f(\mathbb{1}).
\end{equation*}
\item[(ii)] $|f_{\hat{\mathsf{y}}}|\le\mathfrak{n}\,e^{-(\kappa_1-1)|\hat{\mathsf{y}}|}$. Here $|\hat{\mathsf{y}}|$ is the number of members of $\pi^{\flat}$ in $\hat{\mathsf{y}}$, however many cells they hold.
\item[(iii)] $f_{\hat{\mathsf{y}}}$ reads $b$, and $(\mathbb{U},\mathbb{J})$ beyond the $\mathsf{s}$-free reading of $f$, on $\hat{\mathsf{K}}_A\cup\bigcup\hat{\mathsf{y}}$ only. It vanishes unless $\hat{\mathsf{y}}$ is admissible.
\item[(iv)] Let $\mathsf{x}$ be the weight defined in \eqref{5.2:eq-animal-resummation-5} below, and assume $2e^2\hat b_t^2N^{\flat}\mathsf{x}\le1$. Then
\begin{equation*}
\sum_{\hat{\mathsf{y}}\ \text{admissible}}(e\mathsf{x})^{|\hat{\mathsf{y}}|}\le e^{|A|}.
\end{equation*}
Let $\square_0\in A$ and let $\mathsf{E}$ be a closed set. Restrict the sum to those admissible $\hat{\mathsf{y}}$ one of whose components joins $\hat{\mathsf{K}}_{\square_0}$ to $\mathsf{E}$, each weighted by $\mathsf{x}^{|\hat{\mathsf{y}}|}$. This restricted sum is at most
\begin{equation*}
e^{-(\operatorname{dist}_k(\mathsf{E},\hat{\mathsf{K}}_{\square_0})-1)_+}\,e^{|A|}.
\end{equation*}
\item[(v)] Index the elements as $e=(c,A,Q,\mathsf{y})$. Define the fibre map $e\mapsto\hat e:=(c,A,Q,\hat{\mathsf{y}})$, where $\hat{\mathsf{y}}$ is the set of members of $\hat\sigma_0(A)$ holding a cell of $\mathsf{y}$, and put $v_{\hat e}:=\sum_{e\mapsto\hat e}v_e$. Then:
\begin{itemize}
\item $\mathsf{b}(e)=\mathsf{b}(\hat e)$, and the filing levels and the hulls of $e$ and $\hat e$ agree;
\item the exponent is unchanged: $\sum_e v_e=\sum_{\hat e}v_{\hat e}$;
\item for every $n$ and every hull $X$, the finite sum of cluster terms with $n$ polymers and hull $X$ is the same number whether it is formed from the $e$ or from the $\hat e$.
\end{itemize}
Thus every cluster term $\mathbb{C}_I(X)$ is one and the same function. The summation of the elements in norm form used in the proof of Proposition~\ref{5.2:prop-two-slot-preparation} reads an element only through $(\mathsf{b},\mathsf{d},\nu)$. Applied to the $\hat e$, it bounds $\mathbb{C}_I(X)$.
\end{enumerate}
\end{lemma}

Here $\mathsf{b}(e)$ and $\mathsf{d}(e)$ are the sets on which the element $e$ reads the rim field and the pair $(\mathbb{U},\mathbb{J})$, and $\nu(e)$ is its bound, as in the two-slot expansion.

The lemma is applied to the following functions $f$, with the indicated $\mathfrak{n}$.
\begin{enumerate}
\item[(1)] The elements $c_{A,Q}(\mathsf{s})$ of the two-slot expansion of a stored term $c$, plain or wrapped, with
\begin{equation*}
\mathfrak{n}=\mathfrak{N}^{\natural}(c)\prod_{\square\in A}K_A\vartheta(\square)\prod_{q\in Q}(2/r_A).
\end{equation*}
\item[(2)] The far elements of the stored terms of the $\mathbb{E}$- and $\mathbb{R}$-classes.
\item[(3)] At $A=\{\square\}$: the near groups, the far pieces of such a stored term with a near cube, and the $\mathbb{P}$-pieces and the Wilson pieces, with
\begin{equation*}
\mathfrak{n}=c_\alpha C_R(E_0+1)e^{C_2\kappa_1}\bigl(\vartheta(\square)+\mathfrak{p}_j\bigr).
\end{equation*}
For the near groups this size carries the further factor $\mathsf{l}^2(\square)$ of \eqref{5.2:eq-animal-resummation-13}.
\item[(4)] The differences of the functions in (1)--(3) formed by two runs of the construction on one scaffold.
\end{enumerate}
Here $K_A$ is the constant of the all-blocks response dilation, \eqref{5.2:eq-response-dilation-a}, and $C_2$ is the constant in the exponent of its radius $R_0$. The size $\mathfrak{N}^{\natural}(c)$ of a stored term in the norm of the two-slot expansion and the paying function $\mathfrak{p}_j$ of level $j$ are those of Proposition~\ref{5.2:prop-two-slot-preparation}.

\begin{proof}
\emph{(i).} The operators $\varepsilon^1_\Delta$, $\varepsilon^0_\Delta$ act on distinct variables for distinct $\Delta$ and therefore commute. Fix $\hat{\mathsf{y}}$. By the remark preceding the lemma, the sets $\mathsf{y}$ contributing to $f_{\hat{\mathsf{y}}}$ are obtained by choosing independently:
\begin{itemize}
\item an arbitrary subset of the cells of $\sigma_0(A)$ lying in $\hat{\mathsf{K}}_A$;
\item for each $G\in\hat{\mathsf{y}}$, a non-empty subset of the cells of $G$;
\item for each $G\in\hat\sigma_0(A)\setminus\hat{\mathsf{y}}$, the empty subset.
\end{itemize}
Hence the sum \eqref{5.2:eq-animal-resummation-1} over these $\mathsf{y}$ factorises. For a finite set $S$ of cells, summing over all subsets $T\subset S$ gives
\begin{equation*}
\sum_{T\subset S}\ \prod_{\Delta\in T}(\varepsilon^1_\Delta-\varepsilon^0_\Delta)\prod_{\Delta\in S\setminus T}\varepsilon^0_\Delta
=\prod_{\Delta\in S}\varepsilon^1_\Delta .
\end{equation*}
Summing over the non-empty subsets instead gives $\prod_{S}\varepsilon^1_\Delta-\prod_{S}\varepsilon^0_\Delta$.

Let $\varepsilon^1_G$ and $\varepsilon^0_G$ denote the evaluation of all cells of $G$ at $1$ and at $0$. Then
\begin{equation*}
f_{\hat{\mathsf{y}}}=\prod_{G\in\hat{\mathsf{y}}}\bigl(\varepsilon^1_G-\varepsilon^0_G\bigr)
\prod_{G\in\hat\sigma_0(A)\setminus\hat{\mathsf{y}}}\varepsilon^0_G
\prod_{\Delta\in\sigma_0(A),\,\Delta\subset\hat{\mathsf{K}}_A}\varepsilon^1_\Delta\,f .
\end{equation*}
By the definition \eqref{5.2:eq-animal-resummation-2}, the right-hand side equals $\prod_{G\in\hat{\mathsf{y}}}(\varepsilon^1_G-\varepsilon^0_G)$ applied to the function $s\mapsto f(\mathsf{s}_s)$ at $s(G)=0$ for $G\notin\hat{\mathsf{y}}$. Each $\varepsilon^1_G-\varepsilon^0_G$ now acts on the single variable $s(G)$. For $s\in[0,1]^{\hat\sigma_0(A)}$ the point $\mathsf{s}_s$ lies in $\mathfrak{P}_A$, because $e^{\kappa_1}\ge1$. The function is holomorphic there, so the fundamental theorem of calculus in each $s(G)$, $G\in\hat{\mathsf{y}}$, gives \eqref{5.2:eq-animal-resummation-3}.

Every $\mathsf{y}\subset\sigma_0(A)$ belongs to exactly one $\hat{\mathsf{y}}$, namely the set of members holding one of its cells. Hence $\sum_{\hat{\mathsf{y}}}f_{\hat{\mathsf{y}}}=\sum_{\mathsf{y}}f_{\mathsf{y}}$. By the first identity above with $S=\sigma_0(A)$, this equals $\prod_{\Delta\in\sigma_0(A)}\varepsilon^1_\Delta f=f(\mathbb{1})$.

\emph{(ii).} Consider $s$ with $|s(G)|\le e^{\kappa_1}$ for all $G\in\hat\sigma_0(A)$. Each value $\mathsf{s}_s(\Delta)$ is either some $s(G)$ or $1$, and both have modulus at most $e^{\kappa_1}$. So $s\mapsto\mathsf{s}_s$ maps this polydisc into $\mathfrak{P}_A$, and $F(s):=f(\mathsf{s}_s)$ is holomorphic on it with $|F|\le\mathfrak{n}$.

The bounds furnishing $|f|\le\mathfrak{n}$ for the functions (1)--(4) are the all-blocks response dilation, the analyticity bounds of the preparatory stage, and the bounds on the near and far pieces. They are used as they stand, because $F$ is the restriction of $f$ to a diagonal of its own polydisc: all cells of a member carry one parameter. In the random-walk representation of the decoupled operators, the weight of a walk becomes a power of $s(G)$. Cauchy's estimate and the fundamental theorem of calculus are insensitive to the degree, and the multiplicity $m(\omega)$ of a walk $\omega$ is paid by its length $d(\omega)$, by part~(a) of the localisation lemma at complex source.

Now let $s\in[0,1]^{\hat{\mathsf{y}}}$, with $s(G)=0$ off $\hat{\mathsf{y}}$. For each $G\in\hat{\mathsf{y}}$ the circle $|\sigma_G-s(G)|=e^{\kappa_1}-1$ lies in the disc $|\sigma_G|\le e^{\kappa_1}$, since $|s(G)|\le1$. The radius is positive because $\kappa_1>1$ by \eqref{5.2:eq-animal-resummation-a}. Cauchy's formula in the variables $s(G)$, $G\in\hat{\mathsf{y}}$, gives
\begin{equation*}
\Bigl|\prod_{G\in\hat{\mathsf{y}}}\partial_{s(G)}F(s)\Bigr|\le\mathfrak{n}\,(e^{\kappa_1}-1)^{-|\hat{\mathsf{y}}|}.
\end{equation*}
Since $\kappa_1\ge1$, we have $e^{\kappa_1}-1\ge e^{\kappa_1-1}(e-1)\ge e^{\kappa_1-1}$. The domain of integration in \eqref{5.2:eq-animal-resummation-3} has volume one, and (ii) follows.

\emph{(iii).} Fix $s$ with $s(G)=0$ for $G\notin\hat{\mathsf{y}}$. Then $\mathsf{s}_s$ vanishes on every cell of $\sigma_0(A)$ outside $\hat{\mathsf{K}}_A\cup\bigcup\hat{\mathsf{y}}$. Apply part~(b) of the localisation lemma at complex source at $\sigma=\mathsf{s}_s$. The decoupled operators through which $f$ reads $\mathsf{s}$ then read $b$ and $(\mathbb{U},\mathbb{J})$ only on a set
\begin{equation*}
Y(\sigma)\subset\hat{\mathsf{K}}_A\cup\bigcup\hat{\mathsf{y}}.
\end{equation*}
By \eqref{5.2:eq-animal-resummation-3}, $f_{\hat{\mathsf{y}}}$ is an integral over such $s$ of derivatives of $f(\mathsf{s}_s)$, so it inherits this property.

Now let $C$ be a wall-component of $\hat{\mathsf{y}}$ that shares no wall with $\hat{\mathsf{K}}_A$. Being a component, $C$ shares no wall with the rest of $\hat{\mathsf{y}}$ either. Every member across a wall of $\bigcup C$ therefore lies outside $\hat{\mathsf{K}}_A\cup\bigcup\hat{\mathsf{y}}$, where $\mathsf{s}_s=0$. The operators are thus decoupled between $\bigcup C$ and $\hat{\mathsf{K}}_A\supset\bigcup_{\square\in A}\tilde{\square}^{4}$, to which $f$'s reading of $\mathsf{s}$ is restricted. As in the cluster expansion of \cite[\S1]{B-CMP116}, ``the derivatives with respect to $s$ restricted to these components render the term equal to $0$''. Hence $f_{\hat{\mathsf{y}}}=0$ unless $\hat{\mathsf{y}}$ is admissible.

\emph{(iv).} Call a member of $\hat\sigma_0(A)$ a root if it shares a wall with $\hat{\mathsf{K}}_A$. By property~(d), $\hat{\mathsf{K}}_A$ has at most $N^{\flat}|A|$ members. By property~(b), each of them touches at most $\hat b_t$ members. So there are at most $\hat b_tN^{\flat}|A|$ roots.

Every component of an admissible $\hat{\mathsf{y}}$ contains a root, and distinct components contain distinct roots. By property~(b), the members touching a given member number at most $\hat b_t$. Hence the wall-connected sets of $n$ members containing a given root number at most $(e\hat b_t)^n$. Assigning to each component of $\hat{\mathsf{y}}$ one of its roots, we obtain
\begin{equation*}
\sum_{\hat{\mathsf{y}}\ \text{admissible}}(e\mathsf{x})^{|\hat{\mathsf{y}}|}
\le\prod_{\text{roots}}\Bigl(1+\sum_{n\ge1}(e^2\hat b_t\mathsf{x})^n\Bigr).
\end{equation*}
Since $\hat b_t,N^{\flat}\ge1$, the hypothesis $2e^2\hat b_t^2N^{\flat}\mathsf{x}\le1$ gives $e^2\hat b_t\mathsf{x}\le(2\hat b_tN^{\flat})^{-1}\le\frac12$. Hence
\begin{equation*}
\sum_{n\ge1}(e^2\hat b_t\mathsf{x})^n\le2e^2\hat b_t\mathsf{x}\le(\hat b_tN^{\flat})^{-1}.
\end{equation*}
The product over at most $\hat b_tN^{\flat}|A|$ roots is then at most $(1+(\hat b_tN^{\flat})^{-1})^{\hat b_tN^{\flat}|A|}\le e^{|A|}$.

For the second assertion, consider a component joining $\hat{\mathsf{K}}_{\square_0}$ to $\mathsf{E}$. By property~(b), members of $\pi^{\flat}$ are cubes of level at most $k$. Such a component therefore holds at least $(\operatorname{dist}_k(\mathsf{E},\hat{\mathsf{K}}_{\square_0})-1)_+$ members. Consequently
\begin{equation*}
\mathsf{x}^{|\hat{\mathsf{y}}|}=e^{-|\hat{\mathsf{y}}|}(e\mathsf{x})^{|\hat{\mathsf{y}}|}
\le e^{-(\operatorname{dist}_k(\mathsf{E},\hat{\mathsf{K}}_{\square_0})-1)_+}(e\mathsf{x})^{|\hat{\mathsf{y}}|},
\end{equation*}
and the first assertion bounds the remaining sum by $e^{|A|}$.

\emph{(v).} By property~(e), adding to a connected set meeting $\bigcup\mathsf{Q}$ the members of $\pi^{\flat}$ that meet it changes neither its filing level nor its hulls. Hence $\mathsf{b}(e)=\mathsf{b}(\hat e)$ for every $e\mapsto\hat e$, and the filing levels and the hulls agree. The fibres of the map $e\mapsto\hat e$ partition the set of elements, so $\sum_e v_e=\sum_{\hat e}\sum_{e\mapsto\hat e}v_e=\sum_{\hat e}v_{\hat e}$. Moreover,
\begin{equation*}
e^{v_{\hat e}}-1=\prod_{e\mapsto\hat e}\bigl(1+(e^{v_e}-1)\bigr)-1
=\sum_{\emptyset\neq D\subset\{e:\,e\mapsto\hat e\}}\ \prod_{e\in D}\bigl(e^{v_e}-1\bigr).
\end{equation*}
The first step of the summation in norm form is linear in the products $\prod_{e\in D}(e^{v_e}-1)$ and reads only $\bigcup_{e\in D}\mathsf{b}(e)$. The compatibility relation and the Ursell coefficients of the cluster expansion depend only on the supports of the polymers. Hence, for every $n$ and every hull $X$, the finite sum of the cluster terms with $n$ polymers and hull $X$ is the same number on the $e$ and on the $\hat e$. That summation reads the elements through $(\mathsf{b},\mathsf{d},\nu)$ only, so when run on the $\hat e$ it bounds the same cluster terms $\mathbb{C}_I(X)$.
\end{proof}

The lemma is thus a resummation inside the proof of Proposition~\ref{5.2:prop-two-slot-preparation}. The two-slot expansion, its elements, the exponent, the hull rule and every cluster term $\mathbb{C}_I(X)$ are unchanged; the statement at site $(\mathrm{P})$ of the expansion rule (Definition~\ref{5.1:def-expansion-rule}) in the proof of Proposition~\ref{5.2:prop-two-slot-preparation} stands, and nothing is added to the data of the stage.

\smallskip

\emph{Constants.} All constants below are fixed before $\gamma$. We set
\begin{equation}\label{5.2:eq-animal-resummation-5}
\mathsf{x}:=e^{-(\kappa_1-1)+(a_k+1)+3^dc_{\flat}} .
\end{equation}
Here $a_k=(1+3\beta_{\mathrm{pr}})\kappa$ and $c_{\flat}$ is the charge per strong cube of the two-slot expansion. In the summation of the grouped elements, $\mathsf{x}$ is the weight of one member of $\hat{\mathsf{y}}$. Let $c_{\mathsf{p}}:=212^d$ and $c_{\star}:=214^d$ be the constants of the cell-count lemma, by which the block volume factor $V$ of the summations is replaced by $c_{\mathsf{p}}V'$. We impose the floor on $\kappa_1$ and define $K_2^{\flat}$ and $K_{\flat}''$ by
\begin{equation}\label{5.2:eq-animal-resummation-a}
\begin{gathered}
\kappa_1-1\ \ge\ (a_k+1)+3^dc_{\flat}+\log\bigl(2e^2\hat b_t^2N^{\flat}\bigr),\\[2pt]
K_2^{\flat}:=K_A\exp\bigl[(a_k+1+3^dc_{\flat})N^{\flat}+1\bigr],\\[2pt]
K_{\flat}'':=e^{c_J+\sigma_{*}+1}\,c_{\star}\,\bigl(K_2^{\flat}/K_2\bigr)\,K_{\flat}\,\bigl(1+C^{\mathrm{P}}_{\mathsf{a}}\bigr).
\end{gathered}
\end{equation}
Here $K_2$ and $K_{\flat}$ are the constants of the summation in norm form in the proof of Proposition~\ref{5.2:prop-two-slot-preparation}, and $\sigma_*$ and $C^{\mathrm{P}}_{\mathsf{a}}$ are numbers of that summation, fixed before $\gamma$.
The floor is a condition on $\kappa_1$, imposed after $(L,\kappa)$; it reads no coupling. By \eqref{5.2:eq-animal-resummation-5}, it is equivalent to $2e^2\hat b_t^2N^{\flat}\mathsf{x}\le1$, the hypothesis of (iv). Under the floor, the grouped elements hung on $\hat{\mathsf{K}}_A$ cost at most a factor $e$ per block of $A$, and this factor is carried in $K_2^{\flat}$. The constant $K_2^{\flat}$ depends on $(d,L,\kappa,\kappa_1)$ only.

\smallskip

\emph{Covering of a point by the block cut-offs.} The bump $\hat{\chi}_{\square}$ of a block is Ba{\l}aban's cut-off of \cite[(3.4)--(3.7), p.~270]{B-CMP109}. It is built from a profile $\zeta$ with ``$\zeta(t)=1$ for $|t|\le1/3$, $\zeta(t)=0$ for $|t|\ge2/3$''. Hence, with $c_{\square}$ the centre of $\square$,
\begin{equation*}
\square^{\zeta}:=\operatorname{supp}\hat{\chi}_{\square}
\subset\bigl\{x:\ |x-c_{\square}|_\infty\le\tfrac23 M\ell_m\bigr\}\subset\square ,
\end{equation*}
and not merely $\square^{\zeta}\subset\tilde{\square}$. The covering numbers $\nu_d$, $\mathsf{c}_{\square}$, $\mathsf{c}_{\square\square}$ and $N_{\mathrm{cov}}$ are therefore taken on the blocks $\square$ themselves. By the block-covering count of the geometric preliminaries, a point lies in at most $3\cdot2^d\le4^d$ blocks, so $N_{\mathrm{cov}}=4^d$. Consequently $K_A=4\mathsf{K}_re^{C_2\kappa_1}$, which enters $K_2^{\flat}$, depends on $(d,L,\kappa,\kappa_1)$ only. The depth $j-k_0\le N_0-1$ of a window, which an overlap number of the fattened blocks $\tilde{\square}$ would have carried, enters no constant.

\smallskip

\emph{What site $(\mathrm{P})$ reads of the coupling.} At site $(\mathrm{P})$ the coupling enters through two conditions, both one-sided. The first is the floor on $\kappa_1$ in \eqref{5.2:eq-animal-resummation-a}, which reads no coupling. The second consists of the six conditions (q1)--(q6) below, in sum form. In them each scaffold function is replaced by its upper function of $\mathsf{T}$:
\begin{gather*}
N\le C_N(\mathsf{T}+1)^{r_{\mathrm{pr}}},\qquad N_0\le\mathsf{N}_0(\mathsf{T}),\qquad
V'\le\bar V'(\mathsf{T}):=C'_V\bigl(104L\mathsf{R}^{\uparrow}(\mathsf{T})\bigr)^d,\\
\mathfrak{l}_j\le\mathfrak{l}(\mathsf{T}),\qquad \mathsf{l}_j^2\le\mathsf{l}^2(\mathsf{T}),\qquad
\mathfrak{t}^{\mathrm{win}}\le\mathfrak{t}^{\mathrm{win}}(\mathsf{T}),\\
N^{\mathrm{fr}}\le N^{\mathrm{fr}}(\mathsf{T}),\qquad
\mathfrak{F}^{\mathrm{fi}}\le\mathfrak{F}^{\mathrm{fi}}(\mathsf{T}),\qquad \check R\ge\mathsf{T}^{r_{\mathrm{pr}}} .
\end{gather*}
Here $\mathsf{l}_j^2$ is the level-$j$ factor of the overlap bound for the near groups and $\mathsf{l}^2(\mathsf{T})$ its upper function; neither is the block factor $\mathsf{l}^2(\square)$ of \eqref{5.2:eq-animal-resummation-13}.
The paying functions are $\vartheta(\mathsf{T})$, $1/r_A(\mathsf{T})$, $\varepsilon_S(\mathsf{T})$, $\tilde{\mathfrak{l}}(\mathsf{T})$ and $\mathfrak{p}_j(\mathsf{T})$. The conditions are
\begin{equation}\label{5.2:eq-animal-resummation-6}
\begin{aligned}
&\text{(q1)}\quad \mu^{\flat}+\mu_J\le\tfrac14;\\
&\text{(q2)}\quad 4\bigl(\mathsf{M}^{\flat}+\mathsf{M}_J+2^dc_{\mathsf{p}}\bar V'\tilde{\mathfrak{l}}\bigr)\le\tfrac12;\\
&\text{(q3)}\quad 6Nc_{\mathsf{p}}\bar V'\bigl(K_2^{\flat}\vartheta+K_Q/r_A\bigr)\le1;\\
&\text{(q4)}\quad \vartheta\le c_M/2,\quad r_A\ge2,\quad \kappa_A(\mathsf{T})\ge c_{\flat}+c_J+1,\quad
\check R\ge2,\quad Ne^{-\sigma_*\mathsf{T}^{r_{\mathrm{pr}}}/L}\le1;\\
&\text{(q5)}\quad \text{the level condition under which the window-weight lemma below is stated};\\
&\text{(q6)}\quad \varepsilon^N_{\flat,A}\le\tfrac14 .
\end{aligned}
\end{equation}
The totals $\mu^{\flat}$, $\mu_J$, $\mathsf{M}^{\flat}$, $\mathsf{M}_J$ and $\varepsilon^N_{\flat,A}$ are displayed in the counting rule below. In (q4), $\kappa_A(k):=\min_{i\le k}\kappa_A^{\mathrm{pr}}(g_i)$, a one-sided quantity. Here $K_Q$ is the constant attached to a weak cube and $\varepsilon_S$ the bound left by the strong weight in the summation in norm form, as in Proposition~\ref{5.2:prop-two-slot-preparation}, and $\sigma_*$ is the number of \eqref{5.2:eq-animal-resummation-a}. When these conditions are invoked through the inheritance theorem of the preparatory construction, $\vartheta_0$ (the function denoted $\vartheta$ above) is replaced by $24\vartheta_0$, as in the definition of $\gamma_J$.

Each of (q1)--(q6) is implied, for every $k\le K$ and every $K$, by a single inequality
\begin{equation*}
\sup_{\mathsf{T}\ge\mathsf{T}_\gamma}\Phi(\mathsf{T})\le c ,
\end{equation*}
with $\Phi(\mathsf{T})\to0$ as $\mathsf{T}\to\infty$. A supremum over $\{\mathsf{T}\ge\mathsf{T}_\gamma\}$ does not increase as $\gamma$ decreases. Hence the admissible $\gamma$ form an interval $]0,\gamma_P]$ for some number $\gamma_P>0$. The conditions are thus smallness conditions on the coupling, of the kind Ba{\l}aban imposes himself, whose constants ``can be made arbitrarily small by choosing $\gamma$ \dots\ sufficiently small \dots\ do not depend on any scale'' \cite{B-CMP122a}. None of them bounds the coupling from below. The constants $K_2^{\flat}$, $K_{\flat}''$, $c_{\mathsf{p}}$, $c_{\star}$ are numbers.

The degrees in $\mathsf{T}$ are those displayed by the sum form. Each $\Phi$ is a finite sum of products of a number fixed before $\gamma$ and a function of $\mathsf{T}$. Each $\Phi$ tends to $0$ under the condition $q\ge p_0+7r_{\mathrm{pr}}+1$ on Ba{\l}aban's auxiliary exponent $q$ (Definition~\ref{5.1:def-admitted-inputs}), at $d=4$. The worst products are as follows.
\begin{itemize}
\item $\bar V'\mathfrak{t}^{\mathrm{win}}\vartheta$ and $\mathfrak{l}\mathfrak{t}^{\mathrm{win}}\vartheta$ have degree $5r_{\mathrm{pr}}+0^+-7r_{\mathrm{pr}}-1<0$.
\item $c_{\mathrm{loc}}\mathfrak{F}^{\mathrm{fi}}\vartheta$ has degree $6r_{\mathrm{pr}}+0^+-7r_{\mathrm{pr}}-1<0$.
\end{itemize}
No condition contains a product of two volume-type functions, such as $\bar V'\mathsf{l}^2$, $\bar V'\mathfrak{l}$ or $\bar V'^2$. Such products would have degree $8r_{\mathrm{pr}}+0^+-7r_{\mathrm{pr}}-1\ge0$. Each condition is majorised, degree by degree and under the same condition on $q$, by the sum form of the smallness condition of the two-slot expansion, that is, by the constant of that condition times $N\mathsf{r}\bar V+(N_0+3)\mathsf{l}^2$, where $\bar V$ is the upper function of $V$.

\smallskip

\emph{Counting stored terms at site $(\mathrm{P})$.} The bounds per block that enter the summation of the elements count stored terms at a block $\square_0$ of zone $i_0$. They compare the level $\ell$ of a stored term met by $\tilde{\square}_0$ with the zone $i_0$ of the block. One bound weighs $(4L^{i_0-\ell})^d$ anchors against the depth $d_\ell\ge2L^{i_0-\ell}-4$. The other uses that a set $S$ meeting $\tilde{\square}_0$ has filing level $m(S)\ge i_0-1$.

In Ba{\l}aban's overlap window this comparison fails. There ``the basic localization regions $Z_{h+i}\setminus Z''_{h+i}$ contain the common domain $\Omega^c_{k_0+1}\setminus\Omega''_k$ for $h+i>k_0$, hence the bounds from the corresponding terms cumulate. There are at most $N_0$ such terms'' \cite{B-CMP122a}. Moreover, ``some integration regions overlap, more exactly the regions with indices greater than $k_0$'' \cite{B-CMP122a}.

A block of zone $i>k_0+2$ lies over strata of levels in $]k_0,i-2]$, each one cube of its own level thick. A $\pi_\ell$-cube $X\subset\Omega_\ell\cap Z_\ell^{\boxplus}$ under it carries a kept boundary term $\mathbf{B}^{(\ell)}(X)$ with $d_\ell(X)=0$. There are up to $(\check R_{k_0+1}/L)^d$ such cubes per block of zone $i$, all of depth zero.

This is the factor $\mathsf{l}_j^2\ge(2L^{N_0})^d$ of the overlap bound for the near groups, and it cannot be dropped. Suppose it is read as a pile multiplying $V'$. Then the degree in $\mathsf{T}$ of the corresponding condition among (q1)--(q6) becomes non-negative, and the supremum there is infinite. The condition would then have the shape of a bound on the coupling from below. Such a bound is neither true nor needed: the pile comes from counting per block what must be counted per cube of the term's own level.

The terms are therefore counted in their own cubes, independently of the fields, by means of the following four statements. Their letters not defined in the present section ($A_{\flat}$, $B_0$, $b_\delta$, $K(d)$, $C_0^{\sharp}$, $N^{\mathrm{fr}}$, $\mathfrak{w}(c)$, $\mathsf{e}(c)$, $\mathsf{f}(c)$, $\mathcal{A}_c$, $\mathsf{q}_W$, $\mathfrak{s}_J$, $\mathfrak{s}_\Sigma$, $c_{\mathrm{att}}$, $\mathfrak{F}^{\mathrm{fi}}$ and the like) are those fixed with the stratum lemma, the local anchor-sum lemma, the window-weight lemma and the old-block anchor lemma. Throughout, $c_s:=\log(1+3\lambda s+2c_2)$.

\emph{The stratum lemma.} Work at stage $n+1$ of $\mathbf{R}_k$.
\begin{itemize}
\item Let $\square_0$ be a window block of zone $i_0$, so that $\tilde{\square}_0$ meets no live strip, no cell of level $<h$ of a treated component and no untreated component. Every cell met by $\tilde{\square}_0$ has $\mathbf{T}$-level (its base) $b\ge\mathfrak{f}(i_0)$, where $\mathfrak{f}(i_0):=\max\{h,\min\{i_0-2,k_0-1\}\}$. The base levels met lie in $[\mathfrak{f}(i_0),i_0+1]$, and there are at most $\max\{4,i_0-k_0+3\}\le N_0+2$ of them.
\item Let $\square_0$ be an old block over the raised strata of an arrested attempt $a$, with $n_0\le k_a$. The strata are those of this one attempt, and $b\ge\min\{n_0-2,k_0(k_a)-1\}$. Moreover
\begin{equation*}
L^{N_0(k_a)}\le L\cdot L\bigl(\check{\mathsf{T}}_k+1+c_{k-k_0(k_a)}\bigr)^{r_{\mathrm{pr}}},\qquad
N_0(k_a)\le4+2r_{\mathrm{pr}}\log_L\bigl(\check{\mathsf{T}}_k+2+c_{k-k_a}\bigr).
\end{equation*}
\item A stored term of the $\mathbb{E}$- or $\mathbb{R}$-class through a point $y$, with label $A_F$, is short, $d_\ell(A_F)<\check R_k-3$, only at $\ell\in\{b(y),b(y)+1\}$.
\end{itemize}
No $N$, no age, no number of attempts and no number of live regions enters.

\emph{The local anchor-sum lemma.} Assume the floor on $M/M_1$ of the old-block depth estimate. For $\beta\ge\delta_P/16$ and every member $G$ of $\pi^{\flat}$,
\begin{equation*}
\sum_{\square_0}e^{-\beta d^{\wedge}(\square_0,R_n)}\le c_{\mathrm{loc}},
\end{equation*}
the sum running over the blocks $\square_0$ such that $G\subset\hat{\mathsf{K}}_{\square_0}$ or a member of $\hat{\mathsf{K}}_{\square_0}$ touches $G$; here $c_{\mathrm{loc}}$ depends on $(d,L)$ only.

\emph{The window-weight lemma.} Let $\mathfrak{F}$ be the class of stored boundary-type terms, plain or wrapped. It consists of the $\mathbf{B}^{(\ell)}$, the $\mathbf{B}'^{(\ell)}$, the frozen chain terms $\mathbb{C}^{(i)}_{k_a}$, the chain terms $\mathbb{C}^{(i)}_k$ of the running chain, and the terms $\mathsf{q}_W$. Give $c\in\mathfrak{F}$ the weight $W(c):=\mathfrak{w}(c)e^{\mathsf{e}(c)+\mathsf{f}(c)}$. Assume the floor on $\kappa$ at site $(\mathrm{P})$, the floor on $\check R$, and the level condition (q5). Then the following hold.
\begin{itemize}
\item For every window point $y$,
\begin{equation*}
\sum_{c\in\mathfrak{F}:\,A_c\ni y}W(c)\le\mathfrak{t}^{\mathrm{win}}
:=8K(d)A_{\flat}(B_0+b_\delta)+NC_0^{\sharp}+A_{\flat}C_0^{\sharp}N^{\mathrm{fr}} .
\end{equation*}
\item For every point $x$,
\begin{equation*}
\sum_{c:\ A_c\text{ has a }\pi_k\text{-cube}\ni x}W(c)\le\mathfrak{t}^{\mathrm{top}}
:=NC_0^{\sharp}+2K(d)A_{\flat}(B_0+b_\delta).
\end{equation*}
\item At one window block, the coarse incidences weigh at most $(5L)^d\mathfrak{t}^{\mathrm{win}}$; these are the incidences whose own cube is not finer than the block's zone less one. All incidences weigh at most $2\mathfrak{l}_j\mathfrak{t}^{\mathrm{win}}$, where
\begin{gather*}
\mathfrak{l}_j:=
\begin{cases}
3(5L^3)^d, & j\le k_0+1,\\
3(5L^{N_0+1})^d, & j\ge k_0+2,
\end{cases}\\
\mathfrak{l}_j\le c_{\mathfrak{l}}\mathsf{l}_j^2,\qquad
\mathfrak{l}_j\le c'_{\mathfrak{l}}\bigl(L\mathsf{R}^{\uparrow}(\check{\mathsf{T}}_k)\bigr)^d .
\end{gather*}
\end{itemize}
The last bound puts $\mathfrak{l}_j$ in the class of $V'$. The factor $\mathfrak{l}_j$, which cannot be dropped, is displayed per block and added, never multiplied.

\emph{The old-block anchor lemma.} Let $\mathfrak{O}$ be the set of old blocks. For every $\square_0\in\mathfrak{O}$,
\begin{equation*}
\mathfrak{S}^{\mathrm{all}}(\square_0):=\sum\bigl\{W(c):\ c\in\mathfrak{F},\ \square_0\in\mathcal{A}_c\bigr\}
\le C^{\mathrm{P,all}}_{\mathsf{a}}\,e^{\frac18\delta_Pd^{\wedge}(\square_0,R_n)} .
\end{equation*}
Here $C^{\mathrm{P,all}}_{\mathsf{a}}$ is a number fixed before $\gamma$. The age $k-n_0$ and the factor $L^{dN_0(k_a)}$ of the old window sit in a product $\Pi^{\mathrm{all}}(\varrho)$ of a polynomial and a polylogarithm, paid by $e^{2-\varrho}$, which measures the zones crossed. Here and below $\varrho:=d^{\wedge}(\square_0,R_n)/(M/M_1)$ denotes the depth of the block in units of $M/M_1$. No age, no $N_0(k_a)$ and no number of attempts survives.

\emph{The counting rule.} At site $(\mathrm{P})$ every sum over stored terms is taken in the terms' own cubes, never per block.
\begin{itemize}
\item[(i)] At a point or a $\pi_k$-cube, the sum is bounded by $\mathfrak{t}^{\mathrm{win}}$ or $\mathfrak{t}^{\mathrm{top}}$, by the window-weight lemma.
\item[(ii)] At a window block, the coarse incidences are bounded by $(5L)^d\mathfrak{t}^{\mathrm{win}}$ per block. When a sum over the blocks of the collar cannot be avoided, it is taken over the blocks once, against $c_{\star}''V'$. The fine incidences are taken per own cube through the local anchor-sum lemma, with the constant $c_{\mathrm{loc}}$ and no $V'$. Each fine object is counted once:
\begin{itemize}
\item chain positions, at most $(N+1)(100L\check R^{\uparrow})^d$ per component;
\item the stored terms of the $\mathbb{E}$- and $\mathbb{R}$-classes, by their anchors;
\item the small-field terms, at most $4^d(N_0+2)(104L\check R^{\uparrow})^d$.
\end{itemize}
\item[(iii)] At an old block, the sum is bounded by the depth, through the old-block anchor lemma.
\end{itemize}

Let $\mu$ and $\mathsf{M}_*$ be the two totals bounded in the summation in norm form. By the rule,
\begin{equation*}
\mu\le\mu^{\flat}+\mu_J,\qquad \mathsf{M}_*\le\mathsf{M}^{\flat}+\mathsf{M}_J .
\end{equation*}
In $\mu^{\flat}$ the constant $K_{\flat}$ of the unmodified summation is replaced by $4e^{c_J+1/2}K_2^{\flat}c_{\mathrm{loc}}$. The factor $\mathsf{l}_j^2$ per block is replaced by
\begin{equation*}
\mathfrak{l}_j\mathfrak{t}^{\mathrm{win}}+2\mathsf{l}_j^2\mathfrak{s}_J+C^{\mathrm{P},\star}_{\mathsf{a}},\qquad
C^{\mathrm{P},\star}_{\mathsf{a}}:=C^{\mathrm{P,all}}_{\mathsf{a}}+C_{\mathrm{w}}\mathfrak{s}_J ,
\end{equation*}
by the local anchor-sum lemma. Thus $\mu$ carries no $V'$.

$\mathsf{M}^{\flat}$ carries $V'$ once. It enters through the factor of the sum over the elements of one stored term and through the coarse counts per block. The anchor is the block $\square_0\in A$ carrying the contact. The decay of $\sum_{\square_0}\vartheta(\square_0)$ is kept, with $y_{\square_0}=K_2^{\flat}\vartheta(\square_0)$. The fine incidences enter as $c_{\mathrm{loc}}(\mathfrak{F}^{\mathrm{fi}}+\mathfrak{F}^{\mathrm{fi}}_J)$. The near groups enter $\mathsf{M}^{\flat}$, per component of top level at most $j$, through the additive bracket
\begin{equation}\label{5.2:eq-animal-resummation-7}
\begin{split}
&\Bigl[2L^{2d}c_{\star}''V'+C_{\mathrm{pos}}\Bigl(2\cdot4^d(100L\check R^{\uparrow})^d\\
&\qquad+\bigl(2C_{\mathrm{nr}}+(2(7L)^d)^{-1}C^{\partial}_{\mathrm{nr}}\bigr)\mathsf{l}_j^2\Bigr)\Bigr],\\
&C^{\partial}_{\mathrm{nr}}:=7700^dL^3/200 .
\end{split}
\end{equation}
Here $C_{\mathrm{pos}}$ and $C_{\mathrm{nr}}$ are the constants of the near-group lemma and $c_{\star}''$ is the constant of the summation over the blocks of the collar; like $K_2^{\flat}$ and $c_{\mathrm{loc}}$, they are free of $n$, $N$, $N_0$, $\check R$, $k$, $K$, the volume and $g$.
No $N_0$ occurs in this bracket. No product $V'\mathsf{l}^2$, $V'\mathfrak{l}$ or $V'V'$ is ever formed.

The shares of the stored terms of the $\mathbb{E}$- and $\mathbb{R}$-classes are
\begin{equation*}
\mu_J:=4e^{c_J+3/2}K_2^{\flat}\vartheta\,c_{\mathrm{loc}}\,\mathfrak{n}^{\mathrm{kept}},\qquad
\mathsf{M}_J:=2e^{c_J+3/2}K_2^{\flat}\vartheta\,c_{\star}''V'\,\mathfrak{n}^{\mathrm{kept}} .
\end{equation*}
Here $\mathfrak{n}^{\mathrm{kept}}$, a number fixed before $\gamma$, is multiplied by the far factor of the term itself and added at old anchors. No block, $N$ or age is counted in these shares. The degrees are unchanged: $V'$ occurs once in $\mathsf{M}_J$ and not at all in $\mu_J$.

Let $\mathsf{a}_\mu:=\mu^{\flat}|_1$, $\mathsf{a}_{\mathsf{M}}:=\mathsf{M}^{\flat}|_1$, $\mathsf{a}^{J}_\mu:=\mu_J|_1$ and $\mathsf{a}^{J}_{\mathsf{M}}:=\mathsf{M}_J|_1$. The operation $|_1$ replaces by $1$ the size constants of each class of terms and keeps every count, scaffold function and number. The size constants are $NC_0^{\sharp}$; $C_0^{\sharp}$ inside $c_{\mathrm{att}}A_{\flat}C_0^{\sharp}N^{\mathrm{fr}}$; and $b_1$, $B_0$, $b_\delta$, $E_0+1$, $C_\Sigma$ inside $\mathfrak{n}^{\mathrm{kept}}$, so that $\mathfrak{n}^{\mathrm{kept}}|_1=2^d\mathfrak{s}_\Sigma$. Then
\begin{equation}\label{5.2:eq-animal-resummation-8}
\varepsilon^N_{\flat,A}:=N\bigl[8\bigl(\mathsf{a}_{\mathsf{M}}+\mathsf{a}^{J}_{\mathsf{M}}\bigr)
+4\bigl(\mathsf{a}_\mu+\mathsf{a}^{J}_\mu\bigr)\bigr],
\end{equation}
which is the quantity in (q6). Its worst degree in $\mathsf{T}$ is $6r_{\mathrm{pr}}+0^+-(7r_{\mathrm{pr}}+1)$; the degree $6r_{\mathrm{pr}}+0^+$ occurs in the totals only through the product $N^2(L\check R)^d$.

\emph{Two runs.} For two runs $\mathsf{A}$, $\mathsf{B}$ of the construction on one scaffold, the lemma is applied to the differences of member (4). It is linear in $f$, and the grain $\pi^{\flat}$ and the fibre map of (v) are data of the scaffold, common to both runs. Hence the grouped element $\hat e$ of the difference $\delta c$ of a stored term $c$ has the bound
\begin{equation*}
\nu^m(\hat e):=\sup|\delta c|\,e^{\kappa_A|\mathcal{S}_{>j}|}\prod_{\square\in A}K_A\vartheta(\square)
\prod_{q\in Q}(2/r_A)\,e^{-(\kappa_1-1)|\hat{\mathsf{y}}|}.
\end{equation*}
Every lemma used in the counting is linear in the term and anchors by the scaffold. The totals of the counting rule, taken on the $\nu^m$, therefore satisfy
\begin{align*}
\mathsf{M}_*[\nu^m]&\le\bigl(\mathsf{a}_{\mathsf{M}}+\mathsf{a}^{J}_{\mathsf{M}}\bigr)
\Bigl(\sum_{i\le n}\mathsf{s}_i+\mathsf{l}_{n+1}\Bigr),\\
\mu[\nu^m]&\le\bigl(\mathsf{a}_\mu+\mathsf{a}^{J}_\mu\bigr)\Bigl(\sum_{i\le n}\mathsf{s}_i+\mathsf{l}_{n+1}\Bigr),
\end{align*}
and the recursion
\begin{equation*}
\mathsf{s}_{n+1}\le\frac{\varepsilon^N_{\flat,A}}{N}\Bigl[\sum_{i\le n}\mathsf{s}_i+\mathsf{l}_{n+1}\Bigr]+\mathsf{z}_{n+1}
\end{equation*}
holds on the grouped elements, those of the stored terms of the $\mathbb{E}$- and $\mathbb{R}$-classes included. Here $\mathsf{s}_i$ and $\mathsf{z}_{n+1}$ are the difference sizes of the two-run comparison in the proof of Proposition~\ref{5.2:prop-two-slot-preparation}. The difference bounds of the near groups enter the totals $\mathsf{a}$ with the factor $1+C_{\mathrm{w}}$ kept at unit sizes; they rest on clause (e) of the data-weight lemma below. The quantity $\mathsf{l}_{n+1}$ is the maximum of the per-cube difference constants of the stored terms, of the difference sizes of members (3) and (4), and of
\begin{equation*}
\mathsf{l}^{\mathfrak{K}}(k):=\sup_{(Y,m)\ \text{live}}
\bigl[\mathfrak{O}^{\mathcal{K}}(\delta V(Y^{\sharp}))+\mathfrak{O}^{\mathcal{K}}(\delta\mathsf{Q}_Y)\bigr]/\check R_m^{p_\Sigma}
\le2C_\Sigma ,
\end{equation*}
the bound being that of the pair lemma. For every stored term $c$ of the class $\mathfrak{S}$, the quantity $\sup|\delta c|$ and the $\mathsf{s}_i$ are taken over the chart-paired system $\mathcal{K}^{\sharp}_c$. Its seeds carry one common chart and common slot data, and it is closed under the paired reads of the two runs at one common parameter. The paired read of this stage carries a point of the system $\mathcal{K}^{\sharp}$ of the output into $\mathcal{K}^{\sharp}_c$ by construction, memberships being taken run by run. The bound on $\nu^m$ then holds as stated, and on the diagonal $\mathsf{B}=\mathsf{A}$ every $\mathsf{s}_i$ vanishes, by the pair lemma. The value of $\mathsf{l}^{\mathfrak{K}}$ is read by no condition of the present section, since (q6) is read at unit sizes.

The small-field terms of the construction are not in $\mathfrak{F}$. These are the $\mathbb{E}^{(\ell)}$, $P^{\flat}$, $Q$, $R''$, the near groups and the $\mathbb{P}$-pieces and the Wilson pieces. For them the bounds on the near and far pieces hold at any block, since far pieces are long by diameter. The data constant of the near groups carries the factor $\mathsf{l}_j^2$ of the overlap bound. The two per-block bounds described above, the anchor-versus-depth bound and the filing-level bound, are valid at the zones produced by $\mathbf{T}$, that is, at blocks of zones $\ge j+1$ and at rim cubes. They are used nowhere else.

\smallskip

\emph{Old blocks over arrested attempts.} Let $\square_0\in\mathfrak{O}$ be an old block over the raised strata of an arrested attempt $a$, of levels in $]h_a,j_a]$; these strata are kept \cite{B-CMP122a}. At such a block, the per-block bounds for the boundary terms and for the near groups carry the tile and data factor of that attempt, namely $L^{dN_0(k_a)}$, resp.\ $L^{4N_0(k_a)}$, with $N_0(k_a)\ge N_0(k)$. They do not carry the current $\mathsf{l}_j^2$.

This factor is paid by depth. At old blocks the factor $\vartheta(\square_0)=\vartheta e^{-\frac12\delta_Pd^{\wedge}(\square_0,R_n)}$ is kept inside the sum over blocks. The sum $\sum_{\square_0}\vartheta(\square_0)$ over anchors is not bounded by replacing $\vartheta(\square_0)$ with $\vartheta$ at $\square_0\in\mathfrak{O}$.

The stratum lemma and the old-block depth estimate give, at such a block, the depth relations
\begin{equation*}
k-k_a\le N+(h-n_0),\qquad
d^{\wedge}(\square_0,R_n)\ge (M/M_1)\bigl[(h-n_0-2)_++\tfrac52\check R_h-5\bigr].
\end{equation*}
Together with the counts of the stratum lemma, they give
\begin{equation*}
(\text{count per block})\times e^{-\frac14\delta_Pd^{\wedge}(\square_0,R_n)}\le C_{\mathrm{w},\mathrm{cnt}}\le C_{\mathrm{w}},
\end{equation*}
with the constants
\begin{equation}\label{5.2:eq-animal-resummation-9}
\begin{aligned}
C_{\mathrm{w},\mathrm{cnt}}&:=\sup_{\varrho\ge1}16L^{20}
\Bigl(2^{r_{\mathrm{pr}}}(\varrho+4)+3+c_{\lceil\mathfrak{n}_*(\varrho)\rceil+\lceil\varrho\rceil}\Bigr)^{8r_{\mathrm{pr}}}
e^{2-\varrho},\\
\mathfrak{n}_*(\varrho)&:=C_N2^{r_{\mathrm{pr}}}(\varrho+4)+2,\\
C_{\mathrm{w},\mathrm{pr}}&:=\sup_{\varrho\ge1}L^{12}\bigl(\varrho+7+c_{\mathsf{n}(\varrho)}\bigr)^{6r_{\mathrm{pr}}}e^{2-\varrho},\\
\mathsf{n}(\varrho)&:=C_N2^{r_{\mathrm{pr}}}(\varrho+4)+\varrho+4,\\
C_{\mathrm{w}}&:=\max\Bigl\{C_{\mathrm{w},\mathrm{cnt}},\ C_{\mathrm{w},\mathrm{pr}},\
\kappa_Rc_W^3\bigl[2(7L)^d(1+C_{\mathrm{w},\mathrm{pr}})+C^{\mathrm{tile}}\bigr]\Bigr\}.
\end{aligned}
\end{equation}
The constants $\kappa_R$, $c_W$ and $C^{\mathrm{tile}}$ are those of the data-weight lemma, fixed below.

In the depth estimate the index of $c$ in $C_{\mathrm{w},\mathrm{cnt}}$ arises as $N+2+\lceil\varrho\rceil$. The quantity $N+2$ is a quantity of the run, not a number fixed before $\gamma$. At an old block it is majorised by depth, as follows. First, since $\check{\mathsf{T}}_k\ge1$ and $\check R_k\ge\check{\mathsf{T}}_k^{r_{\mathrm{pr}}}$ by the definition of $\check R_k$,
\begin{equation*}
N\le C_N(\check{\mathsf{T}}_k+1)^{r_{\mathrm{pr}}}\le C_N2^{r_{\mathrm{pr}}}\check R_k .
\end{equation*}
Second, $\check R_k\le\check R_h$, because $R$ is non-increasing in the level: ``in some steps the number $R_k$ decreases by the factor $L^{-1}$'' \cite{B-CMP122a}. Third, the depth bound $d^{\wedge}(\square_0,R_n)\ge (M/M_1)[\frac52\check R_h-5]$ of the stratum lemma and the old-block depth estimate gives $\check R_h\le\frac25(\varrho+5)\le\varrho+4$. Hence $N+2\le\mathfrak{n}_*(\varrho)$ at an old block.

Moreover $c_s$ grows at most linearly in $s$. So the suprema in \eqref{5.2:eq-animal-resummation-9} are finite, and all three members of the maximum are numbers fixed before $\gamma$. The first two are fixed with $(d,L,\lambda,c_2,r_{\mathrm{pr}},\mathsf{L}_0)$. The third is fixed, in addition, with $(M,\kappa)$ through $\kappa_R$, $c_W$ and $K_{\mathrm{loc}}$.

$C_{\mathrm{w}}$ is read one-sidedly, in three places:
\begin{itemize}
\item through $C^{\mathrm{P},\star}_{\mathsf{a}}=C^{\mathrm{P,all}}_{\mathsf{a}}+C_{\mathrm{w}}\mathfrak{s}_J$;
\item through the factor $1+C_{\mathrm{w}}$ kept at unit sizes;
\item through the bound $\mathfrak{d}\le1+C_{\mathrm{w}}e^{\frac18\delta_Pd^{\wedge}}$ below.
\end{itemize}
No floor reads it, and no exponent of $L$ is bounded from above. The constant thus joins the family of $K_{\flat}''$: it is one-sided, bounded by depth, and imposes no upper bound on the coupling.

\smallskip

\emph{Data weights of live arrests.} The data weights of a live older arrest persist \cite{B-CMP122a}. Let $\mathcal{W}$ be the set of raised cells of live arrests. A completed operation $\mathbf{R}$ contributes no further class of cells.

What site $(\mathrm{P})$ uses of these weights is the bound
\begin{equation}\label{5.2:eq-animal-resummation-10}
\mathfrak{d}(\square)\le1+C_{\mathrm{w}}e^{\frac18\delta_Pd^{\wedge}(\square,R_n)}\qquad\text{at every block }\square ,
\end{equation}
for the members (3) and (4) of the list following Lemma~\ref{5.2:lem-animal-resummation}. Member (4) involves two runs of the construction on one scaffold. The quantity $\mathfrak{d}(\square)$ is defined through $\mathsf{l}^2(\square)$ below. The bound on the near groups reads \eqref{5.2:eq-animal-resummation-10} from the data-weight lemma, whose constants, hypotheses and statement follow.

\emph{Constants of the data-weight lemma.} The following constants are numbers fixed after $(d,L,M,\kappa,\lambda,c_2)$ and before $\gamma$. None of them depends on an age, on $N_0$, on a depth or on $K$. The letters of the data-weight lemma not defined in the present section ($\theta_2$, $\mathbb{W}$, $\ell_p$, $\beta_0$, $\mathfrak{B}_j$, $M^{\cdot}$, $S^{\star}_v$, $\mathfrak{N}_W$, $\mathfrak{O}^+$, $t(\Delta)$, $X_{\mathfrak{c}}$, $\mathbb{D}_\square$, $r_T(\square)$, $K''$, $\mathfrak{U}_T$, $\mathsf{P}^1$, $\mathfrak{M}^{\mathrm{rd}}$, $C(\kappa)$, $\alpha_i$, $\mathsf{l}(\cdot)$, $\eta$, $\mathfrak{R}(\square)$, and $\mathsf{S}$, $\mathsf{a}$ of the rim construction) are those fixed with that lemma. We write $\mathrm{w}(\Delta)$ for the data weight of a cell $\Delta$, and $w^{\vee}(\Delta)$ for its thresholded value.
\begin{itemize}
\item $K_{\mathrm{ax}}$ is the constant $O(1)$ of \cite[(3.26)]{B-CMP109} on a box of at most nine cubes a side.
\item $C_M^{\mathrm{av}}\ge2$ is the constant of \cite[Prop.~2]{B-CMP98} at complex data.
\item $K_h$ and $c_1$ are the constants of the estimate for the window restricted to a box, on which clauses (a3) and (a4) below rest. The numbers $K_h$, $K_{\mathrm{ax}}M$ and $c_1$ are enlarged, where necessary, to be at least $1$.
\item $c_\cdot:=\sup_n\alpha_{\cdot,n}/\alpha_{0,n}$.
\end{itemize}
With these,
\begin{equation*}
K_{\mathrm{loc}}:=4\bigl(c_\cdot+3K_hK_{\mathrm{ax}}MC_M^{\mathrm{av}}\bigr)\ \ (\ge12C_M^{\mathrm{av}}),\qquad
c_W:=\max\Bigl\{1,\ 2c_1K_{\mathrm{loc}}L^3\sup_n\alpha_{0,n}/\mathfrak{a}_n\Bigr\}.
\end{equation*}
Let $C'$ be the counterterm constant. Let $C_I^{\sharp}$ be the constant of the core estimate in the analytic regime, taken with the rim. Put
\begin{equation*}
C_I':=\max\Bigl\{C_I^{\sharp},\ 8C(\kappa)(10C_MM)^{5/2},\ 4C'(10C_MM)^{1/2}\Bigr\},
\end{equation*}
and let $C_I''$ be the constant $C_I'$ enlarged in the rim construction.

Let $\kappa_R^{\mathrm{act}}$ be the maximum of $1$ and of the least number with the following property. For every block $\square$ and every unit $n$-cell $\Delta'$, consider the contribution of the near units anchored in $\Delta'$ to the total of the near groups at $\square$. These contributions are pieces and remainders. They are formed by the Cauchy lemma of the amplitude disc on the filed disc of each unit and resummed from the table $\mathfrak{N}^{\mathrm{P}}(\cdot\,;1)$ of clause (b) below. For units read on their own, the disc has the single radius $\rho_A:=\mathsf{w}_0r_T(\square)$ and the prefactor $24K''\vartheta(\square)\eta_p$. The property required is that this contribution is at most $\kappa_R^{\mathrm{act}}\mathfrak{N}_0(\square)$.

Let $\kappa_R^{\times}$ be the least number with the same property for the differences of member (4). We read $\kappa_R$ independently of sizes:
\begin{equation*}
\kappa_R:=\max\{\kappa_R^{\mathrm{act}},\kappa_R^{\times}\}.
\end{equation*}
$C_R$ is the least constant of the same kind for the prefactor, whose factor $\eta_p$ is the factor $\eta_p^{\mathrm{own}}$ of units read on their own. Both $\kappa_R$ and $C_R$ are least numbers fixed before $\gamma$; they depend on $(d,L,M,\kappa,\lambda,c_2)$ only.

Further, with $c_s$ as above,
\begin{equation*}
y_*:=\sup\Bigl\{y\ge1:\ y\le L\bigl(2c_{\log_Ly}\bigr)^{r_{\mathrm{pr}}}\Bigr\}<\infty,
\end{equation*}
\begin{equation*}
C^{\mathrm{tile}}:=3\cdot7^d\mathsf{L}_0^6L^d\Bigl[y_*^d+L^d2^{dr_{\mathrm{pr}}}
\sup_{\varrho\ge1}\bigl(\varrho+5+c_{\mathsf{n}(\varrho)}\bigr)^{dr_{\mathrm{pr}}}e^{1-\varrho}\Bigr],
\end{equation*}
with $\mathsf{n}(\varrho)$ as in \eqref{5.2:eq-animal-resummation-9}. The constant $C_{\mathrm{w}}$ is the maximum of three members displayed in \eqref{5.2:eq-animal-resummation-9}; it has one value throughout.

The data-weight lemma is stated under the one-sided smallness condition
\begin{equation}\label{5.2:eq-animal-resummation-11}
K_{\mathrm{loc}}\,\mathsf{A}_w(\gamma)\le c_h,\qquad
\mathsf{A}_w(\gamma):=\sup\bigl\{\mathsf{l}(u)\alpha_i(u):\ u\ge\gamma^{-2}\bigr\}.
\end{equation}
Here $\mathsf{A}_w(\gamma)$ decreases to $0$ as $\gamma\downarrow0$. The number $c_h$ is half the least of the absolute smallness constants of \cite[Prop.~2]{B-CMP98}, of \cite[Thm.~1, Prop.~9]{B-CMP102b} and of the Baker--Campbell--Hausdorff estimates of \cite[(3.30)--(3.32)]{B-CMP109}, each taken for the quantity it bounds. Condition \eqref{5.2:eq-animal-resummation-11} is a bound on $\gamma$ from above, in supremum form, and is one of the conditions defining $\gamma_J$. In the constants of the correlation theorem, $C_I$ is replaced by $C_I'$.

The floors are the following:
\begin{itemize}
\item $L\ge8$;
\item the floors on $M/M_1$ and on $M_1$ under which the all-blocks response dilation is stated;
\item the floor on $M/M_1$ of the old-block depth estimate and the floor on $\check R$;
\item $\kappa\ge10\log L$ and $\min\{\delta_0M,\tfrac12\kappa\}\ge5\log L$.
\end{itemize}

\smallskip

\emph{The three inputs.} The data-weight lemma is formulated as an implication from the following three statements; we quote its constants, hypotheses and conclusions, and use it here only through the bound \eqref{5.2:eq-animal-resummation-10}.

\emph{Input on membership.} For every block $\square$ and every near unit $v$ of the first three kinds that is filed as read on its own, there is a map $\mathsf{R}_p(t,\mathsf{t},\sigma)$, $t\in[0,1]$, with the following properties. It is holomorphic in $(\mathbb{U},\mathbb{J},b,\sigma,\mathsf{t})$ on
\begin{equation*}
(\text{domain of the sources})\times\bigl\{|\mathsf{t}|\le\rho_A(v):=\mathsf{w}_0r_T(\square)\bigr\}
\times\bigl\{|\sigma(\Delta)|\le e^{\kappa_1}\bigr\}.
\end{equation*}
It takes values in $\mathfrak{D}_v=\mathfrak{U}_T$, and its real jet is the stored piece. Finally, the hypothesis of the Cauchy lemma of the amplitude disc holds with $\mathfrak{N}(v):=2\mathfrak{n}_v$.

\emph{Input on remainders.} Let $\square$ be a block of any kind, $v$ a near unit read on its box, and $T$ a constituent of $v$ with $X_T\subset\mathfrak{r}_v$. Consider the remainder configurations \cite[(3.21)--(3.22)]{B-CMP109} of the box construction of the unit. They are stored as the exact difference
\begin{equation*}
\int dt\,\partial_{\mathsf{t}}\big|_{0}\bigl[T(\mathsf{P}^1)-T(\mathsf{P}^0)\bigr]
\end{equation*}
of the two box fields, which is Ba{\l}aban's real jet \cite[p.~274]{B-CMP109}. These configurations lie in $\mathfrak{D}_v$, holomorphically, on the product of the filed disc, $[0,1]$ and $\{|\sigma|\le e^{\kappa_1}\}$. Hence the Cauchy lemma of the amplitude disc and the bound of the stored remainder apply.

\emph{Input on the pair system.} Let $\mathcal{K}:=\mathcal{K}^{\mathrm{rd}}$ be the closure of the seeds of the rim construction under one list $\mathfrak{M}^{\mathrm{rd}}$ of moves. This list contains the following moves:
\begin{itemize}
\item the move supplied by the all-blocks response dilation;
\item the two moves $(\mathsf{t},\sigma)\mapsto\mathsf{P}^1$ and the closure at $(B^0,\Delta)$;
\item a move $\mathsf{M}^{\perp}$ for every cell unit of the Wilson class.
\end{itemize}
Then $\mathcal{K}^{\mathrm{rd}}$ is contained in the production closure, and the four conditions required of the pair system of two runs hold on $\mathcal{K}^{\mathrm{rd}}$. Moreover, the quantities entering $\mathsf{l}^J$ of clause (e) are supplied on the closure $\operatorname{cl}^{\mathrm{pr}}(\mathfrak{M}^{\mathrm{rd}})$.

\smallskip

\emph{The data-weight lemma.} Work in the setting of the preparatory construction, under \eqref{5.2:eq-animal-resummation-11} and the floors, and assume the three inputs. Then the following hold at every stage of every $\mathbf{R}_k$, for every block $\square$, every $(t,\mathsf{t})\in[0,1]\times\mathbb{D}_\square$, and every $\sigma$ with $|\sigma|\le e^{\kappa_1}$.

(a) \emph{Local form of the background.}
\begin{itemize}
\item[(a1)] $|\partial\mathbb{W}(x)-1|<1.4\,\theta_2(x)$ for every cell $x$ meeting $\tilde{\square}^3$.
\item[(a2)] $|\partial\mathfrak{b}'_j(x)-1|\le4\theta_2(x)$ for the same cells and every $j\le k$.
\item[(a3)] Let $Q$ be a box of at most nine cubes of $\pi_p$ a side, with $Q^{(1)}\subset\tilde{\square}^3$ and every cell meeting $Q^{(1)}$ of level at least $p$. Let $\mathrm{w}_Q$ and $\alpha_Q$ be the largest values of $\mathrm{w}(\Delta)$ and of $\alpha_{0,\ell(\Delta)}$ over those cells. In one $G^{\mathbb{C}}$-gauge on $Q^{\sim-1}$ we have $\mathfrak{b}=e^{i\eta A}$ with
\begin{equation*}
\max\bigl\{\ell_p|A|,\ \ell_p^2|\nabla A|,\ \ell_p^{2+\beta}\|\nabla A\|_\beta\bigr\}\le\mathsf{a}_Q:=K_{\mathrm{loc}}\mathrm{w}_Q\alpha_Q\qquad(0\le\beta\le\beta_0).
\end{equation*}
\item[(a4)] Let $j\le p$ and let $\square_0'\subset Q^{\sim-2}$ be a box of $\pi_j$-cubes. Then, modulo gauge,
\begin{equation*}
\mathfrak{b}'_j\lceil\square_0'=(U_j,J_j)(\square_0',\mathsf{V}'),\qquad
\mathsf{V}':=M^{\cdot}(\mathfrak{b}'_j)\lceil\mathfrak{B}_j(\square_0').
\end{equation*}
Put $t_p:=L^{j-p}$ and work in the comb gauge centred at any $y\in\square_0'$. Then $|\partial\mathsf{V}'-1|\le c_1\mathsf{a}_Qt_p^2$, and on all bonds
\begin{equation*}
|B^{\mathrm{ax}}(b)|\le2d(1+|b-y|)c_1\mathsf{a}_Qt_p^2 .
\end{equation*}
Further,
\begin{equation*}
B^{\mathrm{ax}}_\mu(x)=\sum_{\kappa<\mu}(x-y)_\kappa F_{\kappa\mu}(y)+\mathfrak{e}_\mu(x),\qquad |F|\le c_1\mathsf{a}_Qt_p^2 .
\end{equation*}
The remainder satisfies
\begin{align*}
|\mathfrak{e}_\mu(x)|&\le c_1(1+|x-y|)^2\mathsf{a}_Qt_p^{2+\beta}&&\text{on the bonds of level }j,\\
|\mathfrak{e}_\mu(b)|&\le4d(1+|b-y|)c_1\mathsf{a}_Qt_p^2&&\text{on the bonds of the graded rim (levels }i<j).
\end{align*}
\end{itemize}

(b) \emph{The table of unit sizes.} Let $z:=\Theta\rho_{n,j}$. On the filed discs the near units obey the following bounds.
\begin{itemize}
\item A core read on its box satisfies
\begin{equation*}
|\mathfrak{i}_y|\le C_I'\mathfrak{n}\nu^3z^2(1+z)t^{4+\hat\alpha}.
\end{equation*}
\item A short single unit read on its box is bounded by $2\mathfrak{n}_vS^{\star-2}_v$.
\item A Wilson unit, read on its box on $\mathfrak{b}'_j$, is bounded by $\Theta^2\mathfrak{N}_W(v)$.
\item A unit read on its own is bounded by $2\mathfrak{n}_v$, by the input on membership. A core so read has $S^{\star}_v\le1$.
\end{itemize}
Every entry of this table $\mathfrak{N}^{\mathrm{P}}(v;\Theta)$ is non-decreasing and of degree at most $3$ in $\Theta$. Hence
\begin{equation*}
\mathfrak{N}^{\mathrm{P}}\bigl(v;c_Ww^{\vee}(y)\bigr)\le c_W^3w^{\vee}(y)^3\,\mathfrak{N}^{\mathrm{P}}(v;1).
\end{equation*}
The far pieces and the remainders read no datum. The remainders lie in $\mathfrak{D}_v$ at units read on their box, by the input on remainders; there are none at units read on their own. The $\mathbb{P}$-pieces and the Wilson pieces sit at blocks meeting $R_n$, and these read no cell of $\mathcal{W}$.

(c) \emph{Geometry.} Let $\Delta\in\mathcal{W}$ and $y\in\tilde{\square}^2$ with $\operatorname{dist}_n(\Delta,y)<1$. Then
\begin{equation*}
t(\Delta)\ge\ell(y),\qquad t(\Delta)\ge n_0-1,\qquad \square\in\mathfrak{O}^+ .
\end{equation*}
Suppose $\tilde{\square}^2$ holds a point of level at most $n_0-2$. Then $\tilde{\square}^2\subset X_{\mathfrak{c}}$ crosses the thin strata of exactly one preparation $\mathfrak{c}$. Its step $t_{\mathfrak{c}}\ge n_0$ is
\begin{itemize}
\item $k_{\mathfrak{c}}$ for an arrest;
\item $k$ for the running or a sibling preparation of this $\mathbf{R}_k$;
\item $m$ for the completed preparation of a live region $(Y,m)$.
\end{itemize}
Moreover $k_0(\mathfrak{c})\le n_0-2$, and every cell inside $\tilde{\square}^2$ has level in $[k_0(\mathfrak{c}),n_0+1]$. Every weighted cell (with $\mathrm{w}>1$) within one own unit of $\tilde{\square}^2$ belongs to $\mathfrak{c}$.

No window block, the exceptional ones included, reads a cell of $\mathcal{W}$ through $\mathfrak{R}(\square)$. The fivefold fattening $\tilde{\square}^5$ of an exceptional block may meet $\mathcal{W}$; this case is covered by the count-times-weight bound below.

\emph{Count times weight.} Let $\tilde{\square}^5$ meet $\mathcal{W}$. Here $\square\in\mathfrak{O}^+$ or $\square$ is exceptional, and $\varrho\ge\check R_k-4$, with $\varrho$ the depth of $\square$ as above. Put
\begin{equation}\label{5.2:eq-animal-resummation-12}
\mathfrak{K}(\square):=\sum\bigl\{w^{\vee}(\Delta')^3:\ \Delta'\ \text{a unit $n$-cell inside}\ \tilde{\square}^2\bigr\}.
\end{equation}
Here and below $s(\square)$ denotes the raise of $\square$ in the near-group lemma, and $L_\circ^{6s(\square)}$ the tile factor of the tiling lemma at that raise.
\begin{itemize}
\item[(i)] If $\tilde{\square}^2$ holds no point of level at most $n_0-2$, then
\begin{equation*}
\mathfrak{K}(\square)\le2(7L)^d\max\bigl\{L_\circ^{6s(\square)},\ C_{\mathrm{w},\mathrm{pr}}e^{\varrho-2}\bigr\}.
\end{equation*}
\item[(ii)] Otherwise let $t:=t_{\mathfrak{c}}\ge n_0$ be as in (c), with $t=k$ at an exceptional block. Then
\begin{equation*}
\mathfrak{K}(\square)\le3\cdot7^d\mathsf{L}_0^6L^{dN_0(t)}\le C^{\mathrm{tile}}e^{\varrho-1}.
\end{equation*}
\end{itemize}

(d) \emph{The data weight of the near groups.} At every block,
\begin{equation*}
(\text{total of the near groups at }\square)\le\kappa_Rc_W^3\,\mathfrak{K}(\square)\,\mathfrak{N}_0(\square),
\end{equation*}
where $\mathfrak{N}_0(\square)$ is the size of the near groups at $\square$ without the data weight, and with $\mathfrak{K}(\square)$ as in \eqref{5.2:eq-animal-resummation-12}. It is bounded by part (a) of the tiling lemma where $\tilde{\square}^5\cap\mathcal{W}=\emptyset$, and by the count-times-weight bound where $\tilde{\square}^5$ meets $\mathcal{W}$. Define, at every block,
\begin{equation}\label{5.2:eq-animal-resummation-13}
\mathsf{l}^2(\square):=\max\bigl\{L_\circ^{6s(\square)},\ \kappa_Rc_W^3\mathfrak{K}(\square)\bigr\},\qquad
\mathfrak{d}(\square):=\mathsf{l}^2(\square)/L_\circ^{6s(\square)} .
\end{equation}
Then $\mathfrak{d}(\square)\ge1$ at every block.

At the blocks of the count-times-weight bound, let $\mathfrak{K}^{\sharp}(\square)$ be the majorant it provides. Since $\kappa_R\ge1$ and $C_{\mathrm{w}}\ge C_{\mathrm{w},\mathrm{pr}}$,
\begin{equation*}
\mathfrak{d}(\square)\le\mathfrak{d}^{\sharp}(\square)
:=\max\Bigl\{1,\ \kappa_Rc_W^3\mathfrak{K}^{\sharp}(\square)/L_\circ^{6s(\square)}\Bigr\}
\le1+C_{\mathrm{w}}e^{\frac18\delta_Pd^{\wedge}(\square,R_n)} .
\end{equation*}
At the blocks with $\tilde{\square}^5\cap\mathcal{W}=\emptyset$, the total of the near groups is at most $\kappa_Rc_W^3\mathfrak{K}(\square)\mathfrak{N}_0(\square)$, with $\mathfrak{K}$ bounded by part (a) of the tiling lemma. The factor of the near groups there is the same $\mathsf{l}^2(\square)$, its member $\kappa_Rc_W^3\mathfrak{K}(\square)$ being read through the tiling lemma. Accordingly, with $\mathfrak{n}_0$ the corresponding size of a single near group without the data weight, the bound of a near group at such a block reads
\begin{equation*}
\mathfrak{n}(\square)\le\kappa_Rc_W^3\mathfrak{K}(\square)\cdot\mathfrak{n}_0\vartheta\bigl(1+\mathfrak{p}_j/\vartheta\bigr)
e^{-\frac12\delta_Pd^{\wedge}(\square,R_n)} .
\end{equation*}
No equality $\mathsf{l}^2(\square)=L_\circ^{6s(\square)}$ is asserted at these blocks. The data weight of the near groups thus enters at every block through the one quantity $\mathsf{l}^2(\square)$, of degree $0$.

(e) \emph{Two runs.} Let $\mathsf{A}$ and $\mathsf{B}$ be two runs of the construction on one scaffold. Let $\mathcal{K}$ be their pair system, satisfying the four pair-system conditions; all four are met on $\mathcal{K}:=\mathcal{K}^{\mathrm{rd}}$ by the input on the pair system. Put
\begin{equation*}
\delta f(p):=f^{\mathsf{B}}(P^{\mathsf{B}})-f^{\mathsf{A}}(P^{\mathsf{A}}),
\end{equation*}
with $\mathsf{S}$ and $\mathsf{a}$ as in the rim construction. Let $\mathsf{R}^{\times}_p$ be the pair produced by the corresponding move of $\mathfrak{M}^{\mathrm{rd}}$, for units read on their box and for units read on their own respectively. Then, for every near unit and every $p\in\mathcal{K}_{n+1}(\hat X)$ (the class of the stage), on the common disc,
\begin{equation*}
\bigl|\delta(f_v\circ\mathsf{R}^{\times}_p)-\delta f_v(\mathbb{1})\bigr|
\le w^{\vee}(y)^3\,\mathfrak{N}^{\mathrm{P}}(v;c_W)\big|_{\mathfrak{n}\mapsto\mathsf{S}+\mathfrak{n}\mathsf{a},\ C_I'\mapsto C_I''} .
\end{equation*}
Hence the differences of the functions of member (3) at $A=\{\square\}$ belong to the class of member (4), with the bound
\begin{equation*}
\mathsf{l}^2(\square)\cdot\mathfrak{n}_0\big|_{E_0+1\mapsto\mathsf{l}^J}\cdot\bigl(\vartheta(\square)+\mathfrak{p}_j\mathbb{1}\bigr),
\end{equation*}
with the same majorant $\mathsf{l}^2(\square)$. Here $\mathsf{l}^J:=\max\{\mathsf{l}^J_1,\mathsf{l}^J_2\}$ is defined once. The quantity $\mathsf{l}^J_1$ is formed from the sums over $j$ of $\mathsf{b}_j+b_+\mathsf{a}_j$ and from $\sup_v\mathsf{S}^{\perp}_v$, the supremum of the difference sizes in the $\circ$-slot; the quantity $\mathsf{l}^J_2$ is bounded by the series $\bar{\mathfrak{m}}$ of the pair system. Both are fixed, with their constants, together with the rim construction, whose constants are listed in \eqref{5.2:eq-rim-constants-a}; the value of $\mathsf{l}^J$ is read by no condition of the present section.
The quantity $\mathsf{l}^J$ is one of those over which $\mathsf{l}_{n+1}$ is the maximum.

\smallskip

\emph{The mechanism.} The window of a block is read cell by cell, at each cell's own level and actual weight. This is clause (a). On a box $Q$ of at most nine cubes of $\pi_p$, the window restricted to $Q$ is the minimiser of $Q$ itself, and \cite[(3.26)--(3.32)]{B-CMP109} apply with $K_{\mathrm{loc}}\cdot(\max_Q\mathrm{w})\cdot\alpha$.

A near unit feels only $w^{\vee}(y)^3$, the weights at less than one own unit from its anchor. By clause (c), such cells belong to live arrests of step at least $n_0-1$. The exceptional blocks and every window block read none of them.

At old and exceptional blocks the tile is paid by depth. By the count-times-weight bound,
\begin{equation*}
\sum_{\tilde{\square}^2}(w^{\vee})^3\le3\cdot7^d\mathsf{L}_0^6L^{dN_0(t)}\le C^{\mathrm{tile}}e^{\varrho-1},
\end{equation*}
and $C_{\mathrm{w}}$ of \eqref{5.2:eq-animal-resummation-9} has degree $0$.

At the blocks free of $\mathcal{W}$ the tile is paid by count, through $\mathsf{l}_j^2$. This uses part (b) of the tiling lemma and the boundary count lemma, whose number $C^{\partial}_{\mathrm{nr}}$ of \eqref{5.2:eq-animal-resummation-7} is free of $N_0$; each fine cell is counted once, and no depth is used. At the blocks of the count-times-weight bound, \eqref{5.2:eq-animal-resummation-10} is the content of (d); at the remaining blocks it is the content of the lemma on the data-weight ratio $\mathfrak{d}$, parts (i) and (ii), together with the old-block depth estimate.

The quantity $\mathsf{l}^2(\square)$ of \eqref{5.2:eq-animal-resummation-13} is the one factor of the near groups at every block. By (e), the same $\mathsf{l}^2(\square)$ multiplies the difference sizes of member (4). No age of a region enters, and no condition bounds the coupling from below.

The per-block bounds described above, which compare a stored term's level with the zone of the block, are used at site $(\mathrm{P})$ only at the zones produced by $\mathbf{T}$; every other sum over stored terms is taken by the counting rule.

Throughout the following corollary we work at stage $n+1$ of the preparatory integrations at level $k$ and scale $j$, in the frame $\mathfrak{B} := \mathbb{B}^{(n)}_k$, with the notation of the setting of the response map. We recall that notation.

\emph{The configuration variables.}
\begin{itemize}
\item The pair $(\mathbb{U},\mathbb{J})$ lies in $\mathfrak{K}_P \supset \tilde U^{(n+1)c,\Sigma\flat}_k(\hat X)$, where $\hat X$ is the domain over which the spaces $\tilde U^{(\cdot)c,\Sigma\flat}_k(\hat X)$ of the setting are taken; the pair is a field on the whole domain of $\mathfrak{B}$.
\item The decoupling parameters satisfy $|\sigma|\le e^{\kappa_1}$.
\item $R_n$ is the region of stage $n$ of the setting, and $\mathcal{C}_j\subset R_n$ is the set on which the datum $b$ lives. The datum $b$ is a complex field on $\mathcal{C}_j$ with $\sup|b|\le 1.1\,C_1\delta'_j$, where $\delta'_j$ is the constant at scale $j$ of the setting which bounds this datum.
\item $\mathbb{H}$ is the layer of the response-layer theorem at $(\mathfrak{B};\mathbb{U},\mathbb{J};\sigma;(0,b))$.
\end{itemize}

\emph{Further notation.} For a cell $y$ we write $\ell$ for the side of $y$, $\tilde\Delta(y)$ for the fattened cell of $y$ and $m(y)$ for the level of $y$; $N_y$ and $N^{\Delta}_y$ are the cell norms at $y$ of the setting of the response map, and $A'$ is the field of that name in the response-layer theorem. The classes $\mathfrak{M}_{\delta}$ and $\mathfrak{M}_{2\delta}$ are the classes of weights of rates $\delta$ and $2\delta$ of that setting, and $\mathfrak{M}^{\wedge}_{\delta/2}(\mathfrak{B})$ is its class of weights of rate $\delta/2$ on the frame $\mathfrak{B}$. The space $\mathfrak{K}_P$ is the domain of configuration pairs of that setting, $\eta$ is the factor multiplying the layer in the exponent of the first slot, as in that setting, and $C_2$ is the constant of that setting which multiplies $\kappa_1$ in the factor $e^{-C_2\kappa_1}$ below.

\emph{The rate and the distances.} In this corollary $\delta := \delta_0/8$ (this is the rate written $\delta_P$ elsewhere in this chapter; the parameter $\delta$ of \cite[(1.32)]{B-CMP116} does not occur in this corollary), and $d^{\wedge}$ is the scaled distance of the setting. We put
\[
p(y) := \sup|b|\; e^{-\delta d^{\wedge}(y,\mathcal{C}_j)} .
\]

\emph{The blocks.} The response cubes $\square$ carry bumps $\zeta_\square$ with
\[
\operatorname{supp}\zeta_\square\subset\tilde\square,\qquad \sum_\square\zeta_\square=1,
\]
and sets $\square^{\zeta}$. We write $d^-_\square := d^{\wedge}(\square,R_n)$. The map $\chi\mapsto\mathcal{J}^{\sharp}[\chi]$ is the increment of the second slot defined in the second-slot lemma, and for a function $\chi$ on the lattice we put
\[
\Phi^{\chi} := \bigl(e^{i\eta\chi\mathbb{H}}\,\mathbb{U},\ \mathbb{J}+\mathcal{J}^{\sharp}[\chi]\bigr),
\qquad \Phi^{1} := \Phi^{\chi}\big|_{\chi\equiv 1}.
\]

\emph{The local constants.} Put
\[
\mathsf{c}_{\square} := \sup_y\sum_\square e^{-\delta d(y,\tilde\square)},
\]
a number depending only on $d$ and $L$. With the constant $\mathsf{c}_{\vartheta}\ge 1$ of the setting of the response map, the constant $K_S^{b}\ge 1$ of the bump lemma of hypothesis (h) below, and with the numbers $K_\chi$, $B_\sharp$ of the response-layer theorem (which depend on $\kappa_1$), we set $\theta := 1/2$ and
\begin{equation}\label{5.2:eq-response-dilation-a}
\begin{gathered}
\mathsf{K}_r := 2^{9}\,\mathsf{c}_{\square}\mathsf{c}_{\vartheta}\max\{K_\chi,\,K_S^{b}B_\sharp\},\qquad
R_0 := \frac{e^{-C_2\kappa_1}}{2\mathsf{K}_r\vartheta},\qquad
r_\square^{\sharp} := R_0\,e^{\frac12\delta d^-_\square},\\
c_M := \frac{e^{-C_2\kappa_1}}{2\mathsf{K}_r},\qquad
K_A := \frac{2}{c_M} = 4\mathsf{K}_r e^{C_2\kappa_1}.
\end{gathered}
\end{equation}
Here $\vartheta$ is the smallness parameter of the two-slot expansion of Proposition~\ref{5.2:prop-two-slot-preparation}, and $\vartheta(\square) := \vartheta\,e^{-\frac12\delta d^-_\square}$ denotes its value at the block $\square$.

\emph{Entries and room.} The space $\tilde U^{(n)c,\Sigma\flat}_k(\hat X)$ of the setting of the response map is defined by a list of bounds on a configuration, numbered one to seven, together with a bound on cubes; we call them its entries, the last being the cube entry. The space $\tilde U^{(n+1)c,\Sigma\flat}_k(\hat X)$ is defined by the same entries with the stage-$(n+1)$ bounds. The room of an entry at a point $x$ is the allowance, in the setting of the response map, by which a configuration may exceed the entry's stage-$(n+1)$ bound at $x$ and still obey its stage-$n$ bound there; at a point of level $m$ it is the factor $\rho_m$ of that setting times the normalisation of the entry (for entry three, $\rho_m\alpha_{0,m}$). Throughout the corollary $\rho_m$ denotes this room factor; it is not the decay rate of the cluster expansion denoted by the same letter, which does not occur here.

\begin{corollary}[The all-blocks response dilation]\label{5.2:cor-response-dilation}

\emph{Floors.} Let $\varrho_0 := \lceil 8\delta^{-1}\ln L\rceil$; this number is distinct from the radius $\varrho$. Assume
\begin{equation}\label{5.2:eq-response-dilation-1}
\frac{M}{M_1}\ge 3L(\varrho_0+2L),\qquad \frac{\delta M}{M_1}\ge 16\ln L+4,\qquad
\delta_1M\ge 16\kappa_1,\qquad L\ge 8,
\end{equation}
together with the floor on $M_1$ of the setting of the response map. The first two floors in \eqref{5.2:eq-response-dilation-1} and the floor on $M_1$ are floors of the setting of the response map; the floor $\delta_1M\ge 16\kappa_1$, in which $\delta_1$ is the constant so denoted in \cite{B-CMP116}, is a floor of the response-layer theorem.

\emph{Ceilings on $\gamma$.} Assume the ceilings on $\gamma$ of the setting of the response map, together with
\begin{equation}\label{5.2:eq-response-dilation-2}
4\mathsf{K}_r e^{C_2\kappa_1}\vartheta\le 1,\qquad
\hat b_{\gamma} := 22\,B_\sharp\mathsf{c}_{\square}R_0C_1\delta'_j\le \frac1{10}.
\end{equation}

\emph{Named antecedents.} Assume moreover the following; (a) and (c)--(e) are clauses of the second-slot lemma.
\begin{enumerate}
\item[(a)] \emph{Holomorphy of the layer map.} $\Phi^{\chi}$ is holomorphic in $(\mathbb{U},\mathbb{J},b,\sigma,\chi)$.
\item[(b)] \emph{The layer bound and the product estimate.} The layer bound of the setting of the response map holds: for every cell $y$,
\begin{equation}\label{5.2:eq-response-dilation-3}
N_y(\mathbb{H})\le B_\sharp\,p(y),\qquad N^{\Delta}_y(A')\le B_\sharp\,p(y),
\end{equation}
and likewise for the further cell norms of the layer listed in that bound. Moreover the product estimate of that setting, $N_y(\chi\mathbb{H})\le 2a(y)N_y(\mathbb{H})$, holds, with $a$ as in (c). Finally, part (a) of the response-layer theorem gives that at $b=0$ the map $\Phi_t$ defined below does not depend on $t$.
\item[(c)] \emph{The second-slot bound.} Suppose $q\in\mathfrak{M}_{2\delta}$, $a\,p\le q$, $2aB_\sharp p\le 1/10$ and $\mathsf{p}\le 1$. Then, with $y$ the cell of $x$,
\begin{equation}\label{5.2:eq-response-dilation-4}
\ell^{3}\,|\mathcal{J}^{\sharp}[\chi](x)|\le K_\chi\,[a(y)p(y)+q(y)],\qquad
a(y) := \sup_{\tilde\Delta(y)}\max\{|\chi|,\ \ell|\nabla\chi|,\ \ell^{2}|\nabla\nabla\chi|\}.
\end{equation}
\item[(d)] \emph{The majorant of a combination of bumps.} Let $\chi=\sum_\square c_\square\zeta_\square$ with arbitrary amplitudes $c_\square$, and put
\[
q_\square(y) := \sup_{y'\in\square^{\zeta}}p(y')\,e^{-2\delta d(y,y')},\qquad
q := 10\sum_\square|c_\square|\,q_\square .
\]
Then $q\in\mathfrak{M}_{2\delta}$ and $a\,p\le q$. Moreover
\[
a(y)\le 5\sum_{\square:\,\square^{\zeta}\cap\tilde\Delta(y)\neq\emptyset}|c_\square|.
\]
\item[(e)] \emph{Growth of the amplitudes.} If $|c_\square|\le R_0e^{\theta\delta d^-_\square}$, then at every $y$
\begin{equation}\label{5.2:eq-response-dilation-5}
10|c_\square|\,q_\square(y)\le 10R_0\sup|b|\;
e^{-(1-\theta)\delta d^{\wedge}(y,\mathcal{C}_j)}\,e^{-(1+\theta)\delta d(y,\tilde\square)} .
\end{equation}
\item[(f)] \emph{The absorption inequality.} With $\alpha_{\cdot,m} := \min\{\alpha_{0,m},\alpha_{1,m}\}$, at every $y$,
\begin{equation}\label{5.2:eq-response-dilation-6}
C_1\delta'_j\,e^{-\frac14\delta d^{\wedge}(y,\mathcal{C}_j)}\le \vartheta\,\rho_{m(y)}\,\alpha_{\cdot,m(y)} .
\end{equation}
\item[(g)] \emph{Weights and first-order lines.} The function $\hat b_\gamma e^{-\frac12\delta d^{\wedge}(\cdot,\mathcal{C}_j)}$ lies in $\mathfrak{M}^{\wedge}_{\delta/2}(\mathfrak{B})$; this is the first clause of the weight-class lemma, read through part (a) of the general form of the bump lemma with weight function identically $1$. Moreover the first-order lines of the entry estimate hold: if the exponent $\chi\mathbb{H}$ of the first slot satisfies $N_y(\chi\mathbb{H})\le\varphi(y)$ at every $y$ for a weight $\varphi\in\mathfrak{M}^{\wedge}_{\delta/2}(\mathfrak{B})$, and the two domain conditions of the entry estimate hold (the membership condition, that the pair $(\mathbb{U},\mathbb{J})$ lies in $\mathfrak{K}_P$; and the size condition, which for $\varphi=\hat b_\gamma e^{-\frac12\delta d^{\wedge}(\cdot,\mathcal{C}_j)}$ consists of the ceilings on $\gamma$ of the setting of the response map together with $\hat b_\gamma\le\alpha_{\cdot,j}$), then entries one, six and seven increase at a point of level $m$ by at most $3\varphi/\alpha_{\cdot,m}$, in the normalisation in which their room is $\rho_m$.
\item[(h)] \emph{The bump lemma and the weighted space.} The bump lemma in the frame $\mathfrak{B}$ holds, together with its general form and its closing estimate, in the following form. Let $\mathcal{L}_n$ be its index set of pairs $(n',p)$ and $\hat X^{\flat}$ its set of points; for $(n',p)\in\mathcal{L}_n$ and $x\in\hat X^{\flat}$ let $m$ be the level at which the bump lemma takes its radius $\tilde\alpha_{0,m}$, and $l_x$ the level of the point $x$. If, for a weight $\varphi$ as in (g), the quantity $\pi(x):=K_S^{b}\varphi(x)/\tilde\alpha_{0,m}$ satisfies $\pi(x)\le Q\rho_m$ with $Q:=11/2^{9}$ for every $(n',p)\in\mathcal{L}_n$ and $x\in\hat X^{\flat}$, then, with $\pi_\times:=\pi L^{-2(m-l_x)}$, the closing estimate gives that entries two and four of the configuration exceed their stage-$(n+1)$ bounds at $x$ by at most $\frac18 e^{-\frac14\delta d^{\wedge}(x,\mathcal{C}_j)}$ of their room, the square in $\pi_\times$ paying the factor $2^{\Delta}$ of the $\Sigma\flat$-weighting. The two lemmas on the weighted space of the setting of the response map, on which the spaces with the $\Sigma\flat$-weighting rest, hold, and so does the identification of the spaces written with $\Sigma$ and with $\Sigma\flat$.
\end{enumerate}

\emph{The map.} Let $t=(t_\square)$ satisfy $|t_\square|\le r_\square^{\sharp}$ simultaneously for all blocks $\square$. Put
\begin{equation}\label{5.2:eq-response-dilation-7}
\chi_t := \sum_\square t_\square\zeta_\square,\qquad
\Phi_t := \bigl(e^{i\eta\chi_t\mathbb{H}}\,\mathbb{U},\ \mathbb{J}+\chi_t\,\mathcal{J}^{\sharp}[1]\bigr).
\end{equation}

\emph{Conclusions.} Then:
\begin{enumerate}
\item[(i)] $\Phi_t\in\tilde U^{(n)c,\Sigma\flat}_k(\hat X)$, holomorphically in $(\mathbb{U},\mathbb{J},b,\sigma,t)$.
\item[(ii)] Every entry of $\Phi_t$ at $x$ exceeds its stage-$(n+1)$ bound by at most $\frac18 e^{-\frac14\delta d^{\wedge}(x,\mathcal{C}_j)}$ of its room.
\item[(iii)] The restriction of $\Phi_t$ to a set $S$ depends on $t_\square$ only for those $\square$ with $\operatorname{supp}\zeta_\square$ meeting the stencil of $S$. These blocks form the index set $\mathcal{A}_c$ of the M\"obius expansions of Proposition~\ref{5.2:prop-two-slot-preparation}.
\item[(iv)] $\Phi_t=\Phi^{1}$ when $t_\square=1$ for all $\square$, $\Phi_0$ is the identity, and at $b=0$ the map $\Phi_t$ does not depend on $t$.
\item[(v)] $r_\square^{\sharp}\ge 2$.
\end{enumerate}

In other words, the response-dilation hypothesis of Proposition~\ref{5.2:prop-two-slot-preparation} holds on the map $\Phi_t$ itself, with radius $r_\square := c_M/\vartheta(\square)$. Here the $\vartheta(\square)$ are the values of $\vartheta$ at the blocks, and
\[
r_\square\le r_\square^{\sharp}=R_0e^{\frac12\delta d^-_\square},\qquad R_0=\frac{c_M}{\vartheta},\qquad r_\square\ge 2 .
\]
\end{corollary}

\begin{proof}
\emph{Step one: one weight for all blocks.} We apply hypothesis (d) with $c_\square=t_\square$. It gives that
\[
q := 10\sum_\square|t_\square|q_\square
\]
lies in $\mathfrak{M}_{2\delta}$ and that $a\,p\le q$, where $a$ is formed from $\chi_t$:
\[
a(y) = \sup_{\tilde\Delta(y)}\max\{|\chi_t|,\ell|\nabla\chi_t|,\ell^{2}|\nabla\nabla\chi_t|\}.
\]
By the last bound of (d), $a(y)\le 5\sum_{\square:\,\square^{\zeta}\cap\tilde\Delta(y)\neq\emptyset}|t_\square|$; the $\zeta_\square$ are the bumps of the setting, with their profile.

Since $|t_\square|\le r_\square^{\sharp}=R_0e^{\frac12\delta d^-_\square}$, hypothesis (e) applies termwise with $\theta=\frac12$:
\[
10|t_\square|\,q_\square(y)\le 10R_0\sup|b|\;e^{-\frac12\delta d^{\wedge}(y,\mathcal{C}_j)}
\,e^{-\frac32\delta d(y,\tilde\square)} .
\]
Summing over $\square$, we use that $10\sup|b|\le 11\,C_1\delta'_j$ and that
\[
\sum_\square e^{-\frac32\delta d(y,\tilde\square)}\le\sum_\square e^{-\delta d(y,\tilde\square)}\le\mathsf{c}_{\square}.
\]
Together they give, at every $y$,
\begin{equation}\label{5.2:eq-response-dilation-8}
q(y)\le\bar q(y) := 11\,\mathsf{c}_{\square}R_0C_1\delta'_j\,e^{-\frac12\delta d^{\wedge}(y,\mathcal{C}_j)} .
\end{equation}
The amplitudes grow at half the rate at which the layer decays. The tails of the far blocks are summed by the constant $\mathsf{c}_{\square}$, which $\mathsf{K}_r$ already contains.

\emph{Step two: the first slot.} By the product estimate of (b), the layer bound (b), and step one,
\[
N_y(\chi_t\mathbb{H})\le 2a(y)N_y(\mathbb{H})\le 2B_\sharp\,a(y)p(y)\le 2B_\sharp\,q(y).
\]
Since $2B_\sharp\bar q=\hat b_\gamma e^{-\frac12\delta d^{\wedge}(\cdot,\mathcal{C}_j)}$, this yields
\[
N_y(\chi_t\mathbb{H})\le \varphi_\gamma(y) := \hat b_\gamma\,e^{-\frac12\delta d^{\wedge}(y,\mathcal{C}_j)} .
\]
By the first clause of the weight-class lemma, read through part (a) of the general form of the bump lemma with weight function identically $1$, we have $\varphi_\gamma\in\mathfrak{M}^{\wedge}_{\delta/2}(\mathfrak{B})$.

The domain conditions are met as follows.
\begin{itemize}
\item The membership condition, that $(\mathbb{U},\mathbb{J})$ lies in $\mathfrak{K}_P$, holds by the setting.
\item For the size condition, together with the ceilings of the setting of the response map, apply the absorption inequality \eqref{5.2:eq-response-dilation-6} on $\mathcal{C}_j$. Using $R_0\vartheta\le 1/(2\mathsf{K}_r)$ and $\mathsf{K}_r\ge 2^{9}\mathsf{c}_{\square}B_\sharp$,
\[
\hat b_\gamma\le 22B_\sharp\mathsf{c}_{\square}R_0\vartheta\,\alpha_{\cdot,j}
\le\frac{11B_\sharp\mathsf{c}_{\square}}{\mathsf{K}_r}\,\alpha_{\cdot,j}\le\alpha_{\cdot,j}.
\]
Since $\hat b_\gamma$ is of the size of $\gamma$ times a polylogarithmic factor, as in the first-order estimates, this is a ceiling on $\gamma$.
\end{itemize}

Finally, write $C_1\delta'_je^{-\frac12\delta d^{\wedge}}=C_1\delta'_je^{-\frac14\delta d^{\wedge}}\cdot e^{-\frac14\delta d^{\wedge}}$. Applying \eqref{5.2:eq-response-dilation-6} at $y$ to the first factor, together with $R_0\vartheta\le 1/(2\mathsf{K}_r)$, gives
\begin{equation}\label{5.2:eq-response-dilation-9}
\varphi_\gamma(y)\le\frac{11B_\sharp\mathsf{c}_{\square}}{\mathsf{K}_r}\,\rho_{m(y)}\alpha_{\cdot,m(y)}\,
e^{-\frac14\delta d^{\wedge}(y,\mathcal{C}_j)} .
\end{equation}

\emph{Step three: the entries.}

Entry five and the cube entry. These are unmoved, since the field $\bar U$ of the setting of the response map is kept by $\Phi_t$.

Entries one, six and seven. We apply the first-order lines of (g) with $\varphi=\varphi_\gamma$. They bound the increase of each of these entries by $3\varphi_\gamma/\alpha_{\cdot,m}$. By \eqref{5.2:eq-response-dilation-9} and $\mathsf{K}_r\ge 2^{9}\mathsf{c}_{\square}B_\sharp$, this is at most
\[
\frac{33B_\sharp\mathsf{c}_{\square}}{\mathsf{K}_r}\le\frac{33}{2^{9}}<\frac18
\]
of $\rho_m$, with the factor $e^{-\frac14\delta d^{\wedge}}$ kept.

Entries two and four. We apply the bump lemma of (h) in the frame $\mathfrak{B}$, for every $(n',p)\in\mathcal{L}_n$ and $x\in\hat X^{\flat}$, with
\[
\pi(x) := \frac{K_S^{b}\varphi_\gamma(x)}{\tilde\alpha_{0,m}}
\le\frac{11K_S^{b}B_\sharp\mathsf{c}_{\square}}{\mathsf{K}_r}\,\rho_m\le\frac{11}{2^{9}}\,\rho_m,
\qquad
\pi_\times := \pi\,L^{-2(m-l_x)} .
\]
The first inequality is \eqref{5.2:eq-response-dilation-9} with the factor $e^{-\frac14\delta d^{\wedge}}\le 1$ dropped and with the radius $\alpha_{\cdot,m}$ there read as the radius $\tilde\alpha_{0,m}$ of the bump lemma; the second inequality uses $\mathsf{K}_r\ge 2^{9}\mathsf{c}_{\square}K_S^{b}B_\sharp$. These entries are closed by the closing estimate of the bump lemma at $Q := 11/2^{9}$; the square in $\pi_\times$ pays the factor $2^{\Delta}$ of the $\Sigma\flat$-weighting.

Entry three, the second slot. This is the point of the corollary. We apply hypothesis (c) at $\chi\equiv 1$, and check its hypotheses: in this application the function $a$ of (c) is identically $1$; we take $q := p\in\mathfrak{M}_{\delta}\subset\mathfrak{M}_{2\delta}$; by the ceilings on $\gamma$,
\[
2B_\sharp p\le 2.2\,B_\sharp C_1\delta'_j\le \frac1{10};
\]
and $\mathsf{p}\le 1$.

So (c) gives $\ell^{3}|\mathcal{J}^{\sharp}[1](x)|\le 2K_\chi p(y)$. With $a$ now the function of step one, formed from $\chi_t$, we have $|\chi_t(x)|\le a(y)$ by its definition and, by (d) and step one,
\[
a\,p\le q\le\bar q .
\]

Hence, using \eqref{5.2:eq-response-dilation-6} and $R_0\vartheta\le 1/(2\mathsf{K}_r)$,
\begin{align*}
\ell^{3}|\chi_t\,\mathcal{J}^{\sharp}[1](x)|
&\le 2K_\chi\,\bar q(y)
= 22K_\chi\mathsf{c}_{\square}R_0C_1\delta'_j\,e^{-\frac12\delta d^{\wedge}}\\
&\le 22K_\chi\mathsf{c}_{\square}R_0\vartheta\,\rho_m\alpha_{\cdot,m}\,e^{-\frac14\delta d^{\wedge}}
\le\frac{11K_\chi\mathsf{c}_{\square}}{\mathsf{K}_r}\,\rho_m\alpha_{\cdot,m}\,e^{-\frac14\delta d^{\wedge}}\\
&\le\frac{11}{2^{9}}\,\rho_m\alpha_{0,m}\,e^{-\frac14\delta d^{\wedge}(x,\mathcal{C}_j)} .
\end{align*}
The last inequality uses $\mathsf{K}_r\ge 2^{9}\mathsf{c}_{\square}K_\chi$ and $\alpha_{\cdot,m}\le\alpha_{0,m}$. This is less than $\frac18$ of the room, with the factor $e^{-\frac14\delta d^{\wedge}(x,\mathcal{C}_j)}$ kept. Steps two and three, with membership in the $\Sigma\flat$-weighted space read through the two lemmas on the weighted space of (h), give (i) apart from holomorphy, and (ii).

\emph{Step four: holomorphy, locality, special values, radius.}

Holomorphy. By (a), the first slot $e^{i\eta\chi\mathbb{H}}\mathbb{U}$ is holomorphic in $(\mathbb{U},\mathbb{J},b,\sigma,\chi)$. Also by (a), at fixed $\chi\equiv 1$ the increment $\mathcal{J}^{\sharp}[1]$ is holomorphic in $(\mathbb{U},\mathbb{J},b,\sigma)$. Since $\chi_t$ is linear in $t$ and the space is open, $\Phi_t$ is holomorphic in $(\mathbb{U},\mathbb{J},b,\sigma,t)$.

Locality. Both slots read $\chi_t$ at the point: the first through $e^{i\eta\chi_t\mathbb{H}}$, the second as the product $\chi_t\cdot\mathcal{J}^{\sharp}[1]$. Hence the restriction of $\Phi_t$ to $S$ depends on $t_\square$ only if $\operatorname{supp}\zeta_\square$ meets the stencil of $S$.

Special values. Since $\sum_\square\zeta_\square=1$, we have $\chi_t\equiv 1$ when $t_\square=1$ for all $\square$, so $\Phi_t=\Phi^{1}$ there by the definitions. Since $\chi_0=0$, $\Phi_0$ is the identity. At $b=0$, the map $\Phi_t$ does not depend on $t$, by part (a) of the response-layer theorem as recorded in (b).

Radius. The first inequality of \eqref{5.2:eq-response-dilation-2} is equivalent to
\[
R_0=\frac{e^{-C_2\kappa_1}}{2\mathsf{K}_r\vartheta}\ge 2 .
\]
Since $d^-_\square\ge 0$, it follows that $r_\square^{\sharp}\ge R_0\ge 2$. Since $\vartheta(\square)=\vartheta e^{-\frac12\delta d^-_\square}$, we have
\[
r_\square=\frac{c_M}{\vartheta(\square)}=R_0e^{\frac12\delta d^-_\square}=r_\square^{\sharp},
\]
hence $r_\square\le r_\square^{\sharp}$ and $r_\square\ge 2$.
\end{proof}

\emph{Identification with the response-dilation hypothesis.} The corollary places $\Phi_t$ in $\tilde U^{(n)c,\Sigma\flat}_k(\hat X)$. The space written with the superscript $\Sigma$ in the response-dilation hypothesis of Proposition~\ref{5.2:prop-two-slot-preparation} is this same space, by the identification in hypothesis (h).

The map of that hypothesis has second slot
\[
\mathbb{J}+\sum_\square t_\square\zeta_\square\,\mathcal{J}(\mathsf{s};b),
\]
where $\mathcal{J}(\mathsf{s};b)$ is the increment of $\mathbb{J}$ under the substitution \cite[(3.10)--(3.12)]{B-CMP109}. By the recursive properties of the second-slot increment, this increment is $\mathcal{J}^{\sharp}[1](\sigma;b)$ at $\sigma=\mathsf{s}$, as in \cite[(3.13)]{B-CMP109}. Hence that map coincides with $\Phi_t$, and the response-dilation hypothesis holds on the map of the two-slot expansion itself.

Consequently, the M\"obius-inversion lemma of the preparatory expansion and its form with both bracket slots, on which Proposition~\ref{5.2:prop-two-slot-preparation} rests, apply to $\Phi_t$ without change. Since $\Phi_t$ is local in $t$, the index set $\mathcal{A}_c$ suffices. Entry three is obtained at first order from the second-slot bound (c) at $\chi\equiv 1$. No far block enters through a non-local second slot, because the second slot of $\Phi_t$ is a pointwise product. No dilation statement for $\Phi^{\chi}$ at general $\chi$ is used.

The pieces $\zeta_\square\mathcal{J}^{\sharp}[1](\sigma)$ are localised in $\sigma$ by the localisation statement for the decoupled operators. The field $\mathbb{J}$ is a free variable \cite[\S3]{B-CMP109}, so intermediate configurations need only membership in the space.

\emph{Constants and conditions.}
\begin{itemize}
\item \emph{Where $c_M$ and $K_A$ enter.} They enter only the constants $K_2^{\flat}$ and $K_{\flat}''$ of \eqref{5.2:eq-animal-resummation-a}, hence only the smallness condition on the coupling displayed there.
\item \emph{Where $\kappa_1$ enters.} It occurs only in $e^{-C_2\kappa_1}$ (in $R_0$ and $c_M$) and inside $B_\sharp$ and $K_\chi$, hence inside $\mathsf{K}_r$. Raising $\kappa_1$ shrinks a radius fixed before $\gamma$, and nothing else.
\item \emph{Order of choice.} The floors \eqref{5.2:eq-response-dilation-1} are imposed as follows: on $M$, after $(L,M_1)$, for the first two; on $M_1$, after $(d,L)$, for the floor of the setting of the response map; on $M$, after $\kappa_1$, for $\delta_1M\ge 16\kappa_1$; and $L\ge 8$. The ceilings of the setting of the response map and \eqref{5.2:eq-response-dilation-2} are imposed on $\gamma$, last, with $\vartheta\to 0$.
\item \emph{$\delta_0$ and $L$.} No condition compares $\delta_0$ with $L$. $L$ enters $K_S^{b}$ through $c_B := e^{\frac12\delta dL(\varrho_0+2L)}$, a number depending on $L$ which multiplies $\vartheta$ only.
\item \emph{No upper bound on $L$ and no lower bound on the coupling} is imposed by these conditions.
\end{itemize}

\begin{definition}[Three families per scale and the coefficient split]
\label{5.2:def-three-families}
Throughout this definition we work on the whole lattice $T_\eta$, with no small-field restriction: the
regions $\Omega_i$ and $\Lambda_i$ of Ba{\l}aban's construction coincide with $T_\eta$ for every $i$.
Let $\xi=(\xi_0,\xi_1,\dots)\in\mathfrak{P}$ and $\zeta\in D(0,2r)$, and write $\zeta_i=\zeta_i(\zeta)$
for the running shifts determined by $\zeta$. For each scale $j$ we define three families at scale
$j+1$, and a fourth as their difference.

\begin{enumerate}
\item[(a)] The family $\mathsf{P}^{(j+1)}_\xi$ is
\begin{equation}\label{5.2:eq-three-families-1}
\mathsf{P}^{(j+1)}_\xi:=L^{(j+1)}+\Phi_{j+1}\bigl(g_j,s_j;\mathsf{P}^{(\le j)}_\xi\bigr),
\qquad s_j:=1-\xi_j g_j^2 ,
\end{equation}
where $\mathsf{P}^{(\le j)}_\xi$ stands for $(\mathsf{P}^{(i)}_\xi)_{i\le j}$, and $\Phi_{j+1}$ is the family of
the proposition defining the families $\Phi_{k+1}(\mathfrak{q},s;h)$, read at $\mathfrak{q}=g_j$: its radii
are free of $g$, and the parameter $s$ is carried in the Gaussian measure.
\item[(b)] The family $E^{(j+1)}_\zeta(X,\cdot,y)$ is the pointed output, on the radii $\hat\alpha_j$,
of the frozen step $\mathbf{T}$, that is, of ``the last logarithm'' of \cite[(3.28)]{B-CMP119}, which is
``formally given by the same formula as the term $E^{(k+1)}$ in [I] [see (2.12), (2.13)]''
\cite{B-CMP119}, where [I] is \cite{B-CMP109}. It is pointed as in
\cite[(3.30)--(3.42)]{B-CMP119} and expanded as in \cite[(3.44)--(3.47)]{B-CMP119}; the cut is
that of \cite[(3.16)]{B-CMP119},
\begin{equation}\label{5.2:eq-three-families-2}
\prod_{b}\chi\bigl(|B(b)|<T_j\bigr),
\end{equation}
the product running over the bonds $b$, and the Gaussian measure is real. The source enters the
integrand through the vertex
\begin{equation}\label{5.2:eq-three-families-3}
\mathfrak{S}_j:=A\bigl(\zeta_j,U_j(e^{iB'}V^{(j)})\bigr)-A\bigl(\zeta_j,U_{j+1}\bigr)
+\zeta_j\sum_y G_y(B'),
\end{equation}
and the inputs are
\begin{equation}\label{5.2:eq-three-families-4}
\sum_{i\le j}\bigl(E^{(i)}_\zeta\bigr)^{\mathrm{ren}}
+\sum_{i\le k_1(j)}\bigl(R^{(i)}_\zeta\bigr)^{\mathrm{vac}} .
\end{equation}
The inputs take this form under the convention of \cite{B-CMP119}: ``It has the same
covariance property as the expression determined by $E_k$, and we include it into $E_k$''.
\item[(c)] The family $P^{\flat(j+1)}_\zeta$ is given by the same construction as in (b), with inputs
$\sum_{i\le j}(P^{\flat(i)}_\zeta)^{\mathrm{ren}}$ only and with the same numbers $\zeta_i(\zeta)$.
\item[(d)] $Q^{(j+1)}_\zeta:=E^{(j+1)}_\zeta-P^{\flat(j+1)}_\zeta$.
\end{enumerate}

With these definitions, $P^{\flat(j)}_\zeta\simeq\mathsf{P}^{(j)}_\xi$ at $\xi_i=\zeta_i(\zeta)$ for every $j$,
and for every $k$
\begin{equation}\label{5.2:eq-three-families-a}
\beta\bigl[E^{(k+1)}_\zeta\bigr]=\beta^0_{k+1}+f_{k+1}(\zeta_0,\dots,\zeta_k)
+\beta\bigl[Q^{(k+1)}_\zeta\bigr],
\end{equation}
where
\begin{equation}\label{5.2:eq-three-families-5}
f_{k+1}(\xi):=\beta\bigl[\Phi_{k+1}\bigl(g_k,s_k;\mathsf{P}^{(\le k)}_\xi\bigr)\bigr],
\end{equation}
and $f_{k+1}(\zeta_0,\dots,\zeta_k)$ denotes $f_{k+1}$ at $\xi_i=\zeta_i(\zeta)$, $i\le k$.
\end{definition}

\begin{proof}
We prove the equivalence $P^{\flat(j)}_\zeta\simeq\mathsf{P}^{(j)}_\xi$ at $\xi_i=\zeta_i(\zeta)$ by induction
on $j$. Assume that it holds for every $i\le j$, and set $\xi_i=\zeta_i(\zeta)$ for $i\le j$. We show
that the integral defining $P^{\flat(j+1)}_\zeta$ in (c) and the integral defining
$\mathsf{P}^{(j+1)}_\xi$ in \eqref{5.2:eq-three-families-1}, through
$\Phi_{j+1}(g_j,s_j;\mathsf{P}^{(\le j)}_\xi)$, have one integrand.

By the Gaussian lemma at complex parameter $s$, the integral defining
$\Phi_{j+1}(g_j,s_j;\cdot)$ carries the Gaussian weight $e^{-(s/2)\langle B,\Delta B\rangle}$ with
$s=s_j$, and this weight factorises as
\begin{equation*}
e^{-(s/2)\langle B,\Delta B\rangle}
=e^{-\frac12\langle B,\Delta B\rangle}\,e^{-\frac12(s-1)\langle B,\Delta B\rangle}.
\end{equation*}
The first factor is the real Gaussian weight of the step $\mathbf{T}$ used in (b) and (c). On the
other side, for constant $\zeta_j$ the vertex \eqref{5.2:eq-three-families-3} is
\begin{equation*}
\mathfrak{S}_j=(s-1)\Bigl[-\tfrac12\langle B,\Delta B\rangle+g^{-2}\sum_i\Psi_i\Bigr],
\qquad s-1=-\zeta_j g_j^2 ,
\end{equation*}
the term linear in the fluctuation field vanishing. Since $\xi_j=\zeta_j(\zeta)$, the parameter
$s_j=1-\xi_j g_j^2$ of \eqref{5.2:eq-three-families-1} satisfies $s_j-1=-\zeta_j g_j^2$, so the
number $s-1$ multiplying the quadratic form in the vertex is the number $s-1$ in the second
factor of the Gaussian weight. It is through the Gaussian lemma at complex parameter $s$, the
factorisation above and this form of the vertex that the Gaussian and vertex parts of the two
integrands are identified.

It remains to compare the inputs. In (c) the input is $\sum_{i\le j}(P^{\flat(i)}_\zeta)^{\mathrm{ren}}$;
in \eqref{5.2:eq-three-families-1} it is $\mathsf{P}^{(\le j)}_\xi$. The inputs are read as totals, and
by the induction hypothesis $P^{\flat(i)}_\zeta\simeq\mathsf{P}^{(i)}_\xi$ at $\xi_i=\zeta_i(\zeta)$ for each
$i\le j$; the field-free constants produced along the way drop. Hence the two integrals have one
integrand, and $P^{\flat(j+1)}_\zeta\simeq\mathsf{P}^{(j+1)}_\xi$ at $\xi_i=\zeta_i(\zeta)$, which completes
the induction.

For \eqref{5.2:eq-three-families-a}, fix $k$. By (d), $E^{(k+1)}_\zeta=P^{\flat(k+1)}_\zeta+Q^{(k+1)}_\zeta$.
By the equivalence just proved at scale $k+1$, $P^{\flat(k+1)}_\zeta\simeq\mathsf{P}^{(k+1)}_\xi$ at
$\xi_i=\zeta_i(\zeta)$, and by \eqref{5.2:eq-three-families-1} at $j=k$,
\begin{equation*}
\mathsf{P}^{(k+1)}_\xi=L^{(k+1)}+\Phi_{k+1}\bigl(g_k,s_k;\mathsf{P}^{(\le k)}_\xi\bigr).
\end{equation*}
The extraction lemma for $\beta[\cdot]$, applied to this decomposition of $E^{(k+1)}_\zeta$, gives
the coefficient $\beta^0_{k+1}$ from $L^{(k+1)}$, the coefficient
$\beta[\Phi_{k+1}(g_k,s_k;\mathsf{P}^{(\le k)}_\xi)]=f_{k+1}(\xi)$ at $\xi_i=\zeta_i(\zeta)$, that is
$f_{k+1}(\zeta_0,\dots,\zeta_k)$ by \eqref{5.2:eq-three-families-5}, and the coefficient
$\beta[Q^{(k+1)}_\zeta]$ from the remainder. This is \eqref{5.2:eq-three-families-a}.
\end{proof}

\begin{proposition}[Smallness of the difference family]\label{5.2:prop-difference-smallness}
Let $E^{(k+1)}_\zeta$, $P^{\flat(k+1)}_\zeta$ and $Q^{(k+1)}_\zeta := E^{(k+1)}_\zeta - P^{\flat(k+1)}_\zeta$ be
the families of Definition~\ref{5.2:def-three-families}. Assume the comparison with the reference
recursion for all pairs of levels not exceeding $k$. Assume moreover that on the disc $D(0,2r)$
of the variable $\zeta$
\begin{equation*}
|\zeta_j| \le \tfrac{8}{3}\,r,\qquad
\|Q^{(j)}_\zeta\|^{\flat} \le \mathfrak{m}_{j-1}\quad (j\le k),\qquad
\|R^{(j)}_\zeta\|^{\flat} \le g_j^{\kappa_0}\quad (j\le k_1(k)),
\end{equation*}
all these quantities being analytic in $\zeta$. Then $E^{(k+1)}_\zeta$ and $P^{\flat(k+1)}_\zeta$ are
pointed E-families on the radii $\hat\alpha_k$, of norm at most $E_0$, analytic in $\zeta$, and
\begin{equation}\label{5.2:eq-difference-smallness-1}
\|Q^{(k+1)}_\zeta\|^{\flat} \le \mathfrak{m}_k,
\qquad
\bigl|\beta[Q^{(k+1)}_\zeta]\bigr| \le v_k .
\end{equation}
\end{proposition}

\begin{proof}
We work at the frozen large-field step $\mathbf{T}$ whose output defines $E^{(k+1)}_\zeta$ in
Definition~\ref{5.2:def-three-families}, on the whole lattice. There is no large-field region $Z$ and no
term-wise splitting of the families; in particular the near-point units of the table of units
\eqref{5.2:eq-unit-table-a}, which occur at the localisation site only, are absent.

We divide the units entering the step into unmarked and marked ones. The unmarked units are:
the $z$-free effective potentials (Ba{\l}aban's $P^{(k)}$), controlled by
\cite[Lemma~2]{B-CMP116} at the parameter $\delta_k$; the source units
$f_{\square} := A(\zeta_k\hat{\chi}_{\square},\cdot)$ and
$\zeta_k(p_y)G_y(B')$ of the table \eqref{5.2:eq-unit-table-a}; and the frozen groups built from
the families $P^{\flat(j)}$, which enter the table as frozen groups of size $\mathfrak{n} = E_0$. By
the first-stage bound (Lemma~\ref{5.2:lem-first-stage-bound}) these unmarked units are of the
format \cite[(1.42)]{B-CMP116} with activity constant at most $\tfrac14\nu_0\alpha_4$, where
$\nu_0$ and $\alpha_4$ are the fixed constants that set the admissible activity in that format.

The marked units are of two kinds. The first consists of the frozen groups
$\mathfrak{i}_y[Q^{(j)}]$, each carrying its own marginal coefficient $\beta[Q^{(j)}]$. This
division of the frozen groups is legitimate: the extraction functional $\beta[\cdot]$ is linear,
so renormalisation is additive and, since $E^{(j)} = P^{\flat(j)} + Q^{(j)}$,
\begin{equation*}
\bigl(E^{(j)}\bigr)^{\mathrm{ren}} = \bigl(P^{\flat(j)}\bigr)^{\mathrm{ren}} + \bigl(Q^{(j)}\bigr)^{\mathrm{ren}} .
\end{equation*}
These groups are bounded by the bound for frozen renormalised groups, applied with
$\mathfrak{n} = \mathsf{n}_j := \sup_\zeta \|Q^{(j)}_\zeta\|^{\flat}$. The second kind consists of
the units formed from the families $R^{(j)}$; they are bounded by the bound for differences of
$R$-terms, applied at the pair of levels $(n,j) = (k+1,j)$ on the radii $\hat\alpha_k$, with the
weight $w_{k,j}$.

By the first-stage bound (Lemma~\ref{5.2:lem-first-stage-bound}), counting at most $t^{-4}$
anchors per cell, the marked part $\mathsf{W}^m$ of the weight supremum $\mathsf{W}$ of
\eqref{5.2:eq-first-stage-bound-a} satisfies
\begin{equation}\label{5.2:eq-difference-smallness-2}
\mathsf{W}^m \le C_U\,\theta_k\Bigl[S_I \max_{j\le k}\mathsf{n}_j
+ \sum_{j\le k_1(k)} w_{k,j}\,g_j^{\kappa_0}\Bigr],
\end{equation}
where
\begin{equation*}
C_U := 4(3M)^4 e^{2^d\kappa}\max\{C_I,\, C\,C_M^2\}.
\end{equation*}
Since $x_j = g_j^{-2}$, we have $g_j^{\kappa_0} = x_j^{-\kappa_0/2}$, and the sum in
\eqref{5.2:eq-difference-smallness-2} is bounded by
\begin{equation*}
\mu_k := \sum_{j\le k} w_{k,j}\,x_j^{-\kappa_0/2} .
\end{equation*}

We now apply the second-stage lemma (Lemma~\ref{5.2:lem-second-stage-activities}) together with
the corollary comparing two runs of the step, which bounds the difference of their outputs.
Applied to the two runs of the step --- with all inputs, whose output is $E^{(k+1)}_\zeta$, and
with the unmarked inputs only, whose output is $P^{\flat(k+1)}_\zeta$
(Definition~\ref{5.2:def-three-families}) --- this gives
\begin{equation*}
\|Q^{(k+1)}_\zeta\|^{\flat} \le K_{CE}\,\mathsf{W}^m .
\end{equation*}
By \eqref{5.2:eq-difference-smallness-2} and the bound on its sum by $\mu_k$,
\begin{equation*}
K_{CE}\,\mathsf{W}^m \le K_{CE}C_U\theta_k\,S_I\max_{j\le k}\mathsf{n}_j + K_{CE}C_U\theta_k\,\mu_k .
\end{equation*}
Suppose that
\begin{equation}\label{5.2:eq-difference-smallness-3}
K_{CE}\,C_U\,\theta_k\,(1+S_I) \le \tfrac12 .
\end{equation}
As $S_I\ge 0$, condition \eqref{5.2:eq-difference-smallness-3} gives both
$K_{CE}C_U\theta_kS_I \le \frac12$ and $K_{CE}C_U\theta_k \le \frac12$, and hence
\begin{equation*}
\|Q^{(k+1)}\|^{\flat} \le \tfrac12\max_{j\le k}\mathsf{n}_j + \tfrac12\,\mu_k .
\end{equation*}
Condition \eqref{5.2:eq-difference-smallness-3} is met once the coupling lies below a ceiling: by
the smallness statement for $\theta$ among the conditions on the reference recursion, $\theta_k$
tends to zero as the coupling tends to zero, so \eqref{5.2:eq-difference-smallness-3} amounts to
an upper bound on $g$. When $\kappa \ge 10\log L$, the factor $e^{2^d\kappa}$ contained in $C_U$
is a power of $L$, and the condition is then met in the same way, by an upper bound on $g$. By the
hypothesis $\mathsf{n}_j \le \mathfrak{m}_{j-1}$ for $j\le k$,
\begin{equation*}
\max_{j\le k}\mathsf{n}_j \le \mathfrak{m}_{k-1} \le \mathfrak{m}_k ,
\end{equation*}
and $\mu_k \le \mathfrak{m}_k$ by the condition on the constants of the reference recursion that
bounds the sum $\mu_k$ by $\mathfrak{m}_k$.
Together these give $\|Q^{(k+1)}_\zeta\|^{\flat} \le \frac12\mathfrak{m}_k + \frac12\mathfrak{m}_k
= \mathfrak{m}_k$, the first bound in \eqref{5.2:eq-difference-smallness-1}.

It remains to verify the covariance property (f3) of the definition of E-families for the two
sub-sums, the marked and the unmarked, into which the inputs of the step were divided. The marking
depends only on the levels of the units, and the restriction $j \le k_1(k)$ is Ba{\l}aban's
merging criterion. Hence both sub-sums are covariant, exactly as in \cite[(3.39)]{B-CMP119}.
There the covariance follows from the transformation laws \cite[(2.29), (2.32)]{B-CMP119} and
from the definitions.

Finally we extract the marginal coefficient. We apply the extraction lemma for the marginal
coefficient at the radii $\hat\alpha_k$. Here $\mathfrak{t}^{-2} \le C_t x_k$, and the factor
$(1+2\beta_s)^{-2}$ is discarded, since it is at most one. This gives
\begin{equation*}
\bigl|\beta[Q^{(k+1)}]\bigr| \le C_\beta\,C_t\,x_k\,\mathfrak{m}_k
\le 2C_\beta\,C_t\,y_k\,\mathfrak{m}_k \le v_k ,
\end{equation*}
where the second inequality uses $x_k \le 2y_k$, and the last is the condition on the constants
of the reference recursion that bounds $2C_\beta C_t y_k\mathfrak{m}_k$ by $v_k$. This is the
second bound in \eqref{5.2:eq-difference-smallness-1}.
\end{proof}

\begin{remark}[Renormalising the remainder by its own coefficient]
\label{5.2:rem-remainder-renormalisation}
This remark explains why, in the induction of Proposition~\ref{5.2:prop-difference-smallness}, the
bounds assumed at the levels $j\le k$ reproduce themselves at level $k+1$ without circularity: the
estimate of $Q^{(k+1)}_\zeta$ uses only those bounds. The two kinds of memory carried by the
induction are treated differently.

\emph{The memory of the remainder.} The remainder families $Q^{(j)}_\zeta$ of the earlier levels
enter the bound on $\|Q^{(k+1)}_\zeta\|^{\flat}$ in the supremum norm, through the sizes
$\mathsf{n}_j := \sup_\zeta \|Q^{(j)}_\zeta\|^{\flat}$, and they carry the geometric weights of the
scale sum $S_I$. This is possible because each $Q^{(j)}_\zeta$ is renormalised by its own
coefficient $\beta[Q^{(j)}_\zeta]$. Suppose instead that the remainder is left unrenormalised, so
that its memory is summed in $\ell^1$. Then the bound picks up, step after step, the product
\begin{equation*}
  \prod_i \bigl(1 + c\,\theta_i\bigr),
\end{equation*}
with $\theta_i := \theta(g_i)$ and $c$ a constant. The logarithm of this product is of order $K$
divided by a quantity that is only polynomial in a logarithm, so that the product grows
exponentially in the number $K$ of steps, up to that divisor in the exponent.
This is the content of the proposition on the unrenormalised treatment of the remainder.
The coefficient extracted at step $k+1$ is therefore forced to be
\begin{equation*}
  b_{k+1} = \beta\bigl[P^{\flat(k+1)}_\zeta\bigr] + \beta\bigl[Q^{(k+1)}_\zeta\bigr].
\end{equation*}

\emph{The memory of the families $R^{(j)}_\zeta$.} These families enter for $j \le k_1(k)$, and the
weight with which $R^{(j)}_\zeta$ enters at level $k+1$ has no decay in $k-j$; none is needed. That
weight is $w_{k,j}$, and
\begin{equation*}
  w_{k,j} \asymp \frac{x_j}{x_{k+1}} .
\end{equation*}
This weight is paid by the factor $g_j^{\kappa_0}$ in the hypothesis
$\|R^{(j)}_\zeta\|^{\flat} \le g_j^{\kappa_0}$ of Proposition~\ref{5.2:prop-difference-smallness}.
\end{remark}

\begin{lemma}[Modulus bound for the sourced large-field operation]\label{5.2:lem-sourced-modulus}
Throughout, a quantity is called $z$-free if it does not depend on the complex sources $z_i$
of $w=\sum_i z_i f_i$. Let $k\le K$. Let $X$ be a large-field component at level $k$, of any
history. Let $w\in\mathcal{W}_k(X)$, or let $w$ be constant. Let $(\mathbf{U},\mathbf{J})$ lie
in the tube on which $\mathbf{T}'^{\,w}_k(X,\cdot)$ is defined, and let $F$ be bounded.

Let $\mathfrak{M}_k(X)$ denote the right member of \cite[(1.79)]{B-CMP122b}, under the
conditions of hypothesis~(a) below, read at $A_{\flat}$ and with the modified cover, in the
following sense. In the minimum of \cite[p.~381]{B-CMP122b} the member $A_{\flat}^2$ of
\eqref{5.2:eq-sourced-modulus-1} joins, and $A_1$ is replaced by $A_{\flat}$ on its right. The
cover \cite[(1.76)]{B-CMP122b} is replaced by
\begin{equation}\label{5.2:eq-sourced-modulus-3}
Z_j\subset\Bigl[(Z'_{j-1})^{\sim 10}\cup\bigcup_i\bigl(Z_i^{(j)\flat}\bigr)^{\sim 2}\Bigr]^{+99},
\end{equation}
the union now extending over $P_j\cup Q_j\cup R_j$ and the cubes of $I_j$. Here $P_j$, $Q_j$,
$R_j$ and $Z'_{j-1}$ are Ba{\l}aban's sets of \cite[(1.76)]{B-CMP122b}, the $Z_i^{(j)\flat}$ are
the corresponding sets of the rim construction, and $I_j$ is the set of bonds carrying the form
gain $\exp(-\gamma_0\mathfrak{b}_j^2/(64\,\Theta_1^2))$ of the rim construction, with
$\mathfrak{b}_j=g_j^{-1}\delta_j$ (see \eqref{5.2:eq-sourced-modulus-7}).
Assume the following.
\begin{enumerate}
\item[(a)] \emph{Smallness.} Ba{\l}aban's smallness conditions on the auxiliary parameters,
under which the right member of \cite[(1.79)]{B-CMP122b} is printed, hold in the form that the
rim construction requires. This form contains the following:
\begin{itemize}
\item the inequality below, for every $j\le k$; here and throughout, $p_0^2(g)$ and $p_0^2$
stand for $(\log g^{-2})^{2p_0}$, Ba{\l}aban's convention in \cite{B-CMP122b}, that is, the
square of the degree function $p_0(\cdot)$ without its factor $A_0$:
\begin{equation}\label{5.2:eq-sourced-modulus-1}
\tfrac14\,\gamma_0A_{\flat}^2\,p_0^2(g_j)\ \ge\ 2p_0(g_j),
\qquad A_{\flat}^2:=\frac{A_1^2}{256\,\Theta_1^2};
\end{equation}
\item the strict degree inequality $2p_1-(d+5)r_{\mathrm{deg}}>p_0$ of \cite{B-CMP122b}, where
$p_1$ and $r_{\mathrm{deg}}$ denote Ba{\l}aban's exponents in that inequality;
\item a last clause controlling the unmatched counterterm pieces at complex $\hat w$;
\item constants $O(1)$, all of which are contained in the constant $C_{\mathfrak{d}}^{\flat}$
of (c) below.
\end{itemize}
\item[(b)] \emph{Containment of cascades.} On the family of large-field data of the rim
construction, every cascade $\mathcal{C}$ of a component $Z$, with its witness
graph, satisfies $\mathcal{C}\subset Z\subset\mathcal{C}^{\sim 109}$; that is, $Z$ contains
the cascade and lies in its enlargement by $109$ layers of cubes.
\item[(c)] \emph{Charge invariance and its corollary.} On the same family the charge-invariance
theorem for large-field components holds. It holds for every component of the complement
$(\Lambda^{*})^{c}$ of the set $\Lambda^{*}$ of that theorem, taken before the operation
$\mathbf{R}$, that is, for the components about to be removed. Its constants are the numbers
\begin{equation*}
\Gamma=\sqrt d\,2^{d+1}219^d,\qquad \sigma=220\sqrt d,\qquad c_0=2d+24 ,
\end{equation*}
which depend on $d$ only; this $\sigma$ is unrelated to the factor $\sigma(g_kB)$ of
\cite[(1.71)]{B-CMP122b}, and this $c_0$ is not the flow constant $c_0$ of the glossary (the
one with $\underline{\lambda}=c_0/4$). Its
deposit is
\begin{equation*}
\mathsf{D}=\frac{\gamma_0A_{\flat}^2p_0^2}{2\sqrt d\,(1+\beta_0)^2},
\end{equation*}
with $\beta_0$ the parameter of \cite[(1.89)]{B-CMP122b}, and with $\gamma_0$ there the smaller
of the lower bound $\gamma_0$ of the constrained fluctuation forms and a constant $\lambda_0$
of that theorem; and the per-cube rate $\mathsf{Q}_n$ is defined as the minimum of the per-cube
rates of the construction. Its charge is
\begin{equation*}
c_n=C_{\mathfrak{d}}^{\flat}M^d\check{R}_n^{d+1},\qquad
C_{\mathfrak{d}}^{\flat}:=C_{\mathfrak{d}}+2^{d+2}(1+dL^d).
\end{equation*}
The theorem holds under its standing conditions, among them one of the ceilings on $\gamma$
collected in $\gamma_P^{\flat}$, and under $L\ge 5$. Its corollary
gives:
\begin{itemize}
\item[(c1)] \cite[(1.80)]{B-CMP122b} as printed;
\item[(c2)] $\kappa_n(Z)\ge\Gamma c_n\ge 0$ at every time $n$ of the clock of that theorem, the
couplings $g_n$ being continued as constants beyond the index $k$; here $\kappa_n(Z)$ is the
quantity which that theorem attaches to the component $Z$ at time $n$. Consequently, for a
component $X$ of class (i) or (ii),
\begin{equation}\label{5.2:eq-sourced-modulus-2}
\mathfrak{M}_k(X)\ \le\ \exp\bigl(-2\mathsf{P}_{j(X)}\bigr)
\ \le\ \exp\bigl(-2(1+\beta_0)^{-1}p_0(g_k)\bigr),
\end{equation}
where $\mathsf{P}_{j(X)}$ is the payment which that theorem attaches to $X$; this is
\cite[(1.89)]{B-CMP122b} on this family;
\item[(c3)] $\kappa_n(Z)\ge\sqrt d\,c_n|Z|\ge\kappa_{\mathfrak{d}}d_n(Z)$, with
$\kappa_{\mathfrak{d}}:=C_{\mathfrak{d}}^{\flat}M^d$. This is Ba{\l}aban's additional term
$-\kappa_1d_k(X)$ of \cite{B-CMP122b}, with coefficient $\kappa_{\mathfrak{d}}$.
\end{itemize}
\item[(d)] \emph{Delay statements.} The delay which the rim construction imposes is paid
for within the bounds, and the case $I_j\neq\emptyset$ is admitted as an event of the
construction.
\item[(e)] \emph{Internal sequences.} The sums over internal sequences, those over the sets
$I_j$ included, are bounded by $c_A+1$ per cell, a factor contained in $c_n$.
\end{enumerate}
Then
\begin{equation}\label{5.2:eq-sourced-modulus-4}
\bigl|\mathbf{T}'^{\,w}_k(X,(\mathbf{U},\mathbf{J}))F\bigr|\ \le\ e\,\mathfrak{M}_k(X)\,\sup|F| .
\end{equation}
If moreover the history of $X$ belongs to the family of hypothesis~(b), then, for $X$ of
class (i) or (ii),
\begin{equation}\label{5.2:eq-sourced-modulus-a}
\bigl|\mathbf{T}'^{\,w}_k(X,(\mathbf{U},\mathbf{J}))F\bigr|
\ \le\ e\,\exp\bigl(-2(1+\beta_0)^{-1}p_0(g_k)\bigr)\,\sup|F| .
\end{equation}
\end{lemma}

\begin{proof}
We start from the representation \cite[(1.71)]{B-CMP122b} of the large-field operation. Its
factors are the following:
\begin{itemize}
\item the measures, the factor $\sigma(g_kB)$, the tree $\delta$-functions and the
characteristic functions $\chi$, all of which are $z$-free;
\item the inner operations $\mathbf{T}_h$;
\item the quadratic form $-\tfrac12\langle DH''B,\hat{\chi}DH''B\rangle$;
\item the raw Wilson term $-A\bigl((x''_k-\zeta_k)\hat{\chi}_1,U''_{h+2}\bigr)$;
\item the term $V(X^{\sim};w)$;
\item constants.
\end{itemize}
As observed in \cite{B-CMP122b}, only the quadratic form in $B$ and the term $V(X^{\sim})$
depend on $(\mathbf{U},\mathbf{J})$. We bound the modulus of each factor at the point
$((\mathbf{U},\mathbf{J}),w)$.

(1) \emph{The term $V$.} We have $|e^{V}|\le e^{\sup|V|}$. The supremum is bounded by
\eqref{5.2:eq-resummation-input-a}, which is a statement on the tube, at every
$w\in\mathcal{W}_k$. Ba{\l}aban's passage through $\sigma$, in which the relevant quantity is
real and positive, and the bounds \cite[(1.74)--(1.75)]{B-CMP122b} serve only the optional
bounds \cite[(1.103)--(1.104)]{B-CMP122b}. None of them is used here.

(2) \emph{The quadratic form.} Its complex part is $z$-free and small, as shown in
\cite{B-CMP122b}.

(3) \emph{Raw Wilson factors.} Recall $A(c,U)=\sum_p c(p)\,a_p(U)$ with
$a_p(U)=1-\tfrac12\operatorname{tr}(U+U^{-1})(\partial p)$. The configuration $U''_{h+2}$ is
$\mathrm{SU}(N)$-valued, so the action there is non-negative \cite{B-CMP122b}. Each
$a_p(U''_{h+2})$ is
real and non-negative. Hence the modulus of $e^{-A(c,U''_{h+2})}$ is
$e^{-A(\operatorname{Re}c,\,U''_{h+2})}$.

For $c=(x''_k-\zeta_k)\hat{\chi}_1$, clause (c) of the complex-weight lemma therefore gives
\begin{equation}\label{5.2:eq-sourced-modulus-5}
\Bigl|e^{-A((x''_k-\zeta_k)\hat{\chi}_1,\,U''_{h+2})}\Bigr|
\ \le\ e^{-(1-\theta_W)A(x''_k\hat{\chi}_1,\,U''_{h+2})} .
\end{equation}
The loss $1-\theta_W$ is taken on these raw Wilson factors only. The same bound is applied
to the raw Wilson factors inside the inner operations $\mathbf{T}_h$.

(4) \emph{The $A_j$-integrals.} The quadratic forms of \cite[(2.21)]{B-CMP122b} are
$z$-free. The exponent $\mathfrak{S}_j$ of the weight factor $e^{\mathfrak{S}_j}$ of the
$A_j$-integrals is at most one per cell, by clause (b) of the
complex-weight lemma; the upper restriction of that clause persists on the set $S_{j+1}$ of
the next level of the construction. The rim factors $\chi^0$,
$\chi^1$ of the rim construction take values in $[0,1]$, so
their moduli are at most one.

(5) \emph{Constants.} The constants contribute $O(1)|Z_j|$. This lies inside Ba{\l}aban's
allowance \cite{B-CMP122b}, by which every large-field domain $Z_j$ contributes a constant
bounded as
\begin{equation}\label{5.2:eq-sourced-modulus-6}
\begin{split}
O(1)\log g_j^{-2}\,|Z_j|
&\le O(1)\log g_j^{-2}\,(MR_j)^d\,d'_j(Z_j)\\
&\le O(1)\,M^dR_j^{d+1}\,d'_j(Z_j).
\end{split}
\end{equation}
Under the standing reading convention of this chapter, the collar $Z^{\sim}$ is replaced by
the strip $Z^{\sharp}$ and $d_{k,Z}$ by $d_{k,Z^{\sharp}}$. Under it, the strip $X^{\sharp}$
has at most $2^dN(X)$ cells, where $N(X)$ is the number of cells of $X$; this is the first
clause of the strip lemma. So the per-cell constant of the strip form of
\eqref{5.2:eq-resummation-input-a} carries one further factor $2^d$, as the corollary of the
strip construction for the present bound states. That per-cell constant is at most a multiple
of $\log g_j^{-2}$, and so lies inside the first member of \eqref{5.2:eq-sourced-modulus-6};
the frozen terms left by chains of large-field operations stopped at earlier steps are bounded
extensively, at the scale of each frozen term, inside its third member. Both bounds are proved
in this section, in the lemma on the per-cell value constant and in the lemma on the frozen
piles of live arrests.

(6) \emph{Unmatched counterterm pieces.} The unmatched counterterm pieces at complex $\hat w$
are controlled by the last clause of the smallness conditions in hypothesis~(a).

(7) \emph{Gains.} The gains are those of \cite[pp.~380--381]{B-CMP122b}. They are the factor
\begin{equation*}
\exp\Bigl(-\gamma_0(2g_j^2)^{-1}\bigl(B_3^{-1}\varepsilon_j\bigr)^2
(LM_2R_j)^{-d}|P_j|\Bigr),
\end{equation*}
the factor attached to $Q_j$, and the factor
\begin{equation*}
\exp\Bigl(-\gamma_0g_j^{-2}\delta_j^2(LM_2R_j)^{-d}|R_j|\Bigr).
\end{equation*}
Here $|P_j|$ and $|R_j|$ are the volumes of the regions $P_j$, $R_j$ entering the cover
\eqref{5.2:eq-sourced-modulus-3}, while $R_j$ inside $(LM_2R_j)$, and in
\eqref{5.2:eq-sourced-modulus-6}, is Ba{\l}aban's radius.
With the rim construction there is in addition the form gain
\begin{equation}\label{5.2:eq-sourced-modulus-7}
\exp\bigl(-\gamma_0\mathfrak{b}_j^2/(64\,\Theta_1^2)\bigr)\qquad\text{per bond of }I_j ,
\qquad \mathfrak{b}_j=g_j^{-1}\delta_j .
\end{equation}
This form gain rests on the lower bound $\gamma_0$ for the constrained fluctuation forms at
the true pairs of \cite[(2.21)]{B-CMP122b}, whose backgrounds there are the $z$-free
minimisers of the variational problem.

Each gain is a coefficient times a $z$-free lower bound of a non-negative, coupling-free
functional. For the Wilson gains the coefficient is
$x_j(p)-\operatorname{Re}\zeta_j(p)\ge(1-\theta_W)\,x_j(p)$. For the form gain it is unchanged.
In \cite[p.~381]{B-CMP122b} the gains are then combined into
\begin{equation*}
\begin{split}
&\exp\Bigl(-\gamma_0\min\bigl\{\tfrac12B_3^{-2}A_0^2,\,2A_1^2,\,A_1^2\bigr\}\,
p_0^2(g_j)\,(d'_j+1)\Bigr)\\
&\qquad\le\exp\Bigl(-\tfrac12\gamma_0A_1^2\,p_0^2(g_j)\,(d'_j+1)-2p_0(g_j)\Bigr).
\end{split}
\end{equation*}
Here the minimum contains the further member $A_{\flat}^2$, the smaller one, and the Wilson
gains carry the factor $1-\theta_W\ge\tfrac34$. The same right member, with $A_1$ replaced by
$A_{\flat}$, then follows from the inequality \eqref{5.2:eq-sourced-modulus-1} of
hypothesis~(a).

(8) \emph{Preparatory factors and sums.} The preparatory factors are controlled by the strict
degree inequality $2p_1-(d+5)r_{\mathrm{deg}}>p_0$ of hypothesis~(a). The group integrals are
bounded by
one, and the sums over sequences are $z$-free \cite{B-CMP122b}. On the family of
hypothesis~(b), the sums over internal sequences, those over the sets $I_j$ included, are
bounded by
hypothesis~(e).

(9) \emph{Assembly.} The charge-invariance theorem of hypothesis~(c) is arithmetic on the
right member so obtained, at the deposit $\mathsf{D}$, the payments $\mathsf{P}_j$ and the
charges $c_n$. The constant $C_{\mathfrak{d}}^{\flat}$ contains the $O(1)$ constants of the
smallness conditions. That theorem is a statement about a finite history and the couplings
$(g_n)_{n\le k}$, continued as constants in the time parameter; it reads no coupling of index
larger than $k$. There is therefore no circularity with the reference recursion of
Definition~\ref{5.1:def-reference-recursion}.

The following are not used anywhere in this argument:
\begin{itemize}
\item the cover \cite[(1.76)]{B-CMP122b};
\item the bound \cite[(1.84)]{B-CMP122b};
\item the connectedness of the graph $G$ derived from it in \cite{B-CMP122b};
\item the argument of \cite[p.~387]{B-CMP122b}.
\end{itemize}
Collecting the moduli bounds (1)--(6) and the gains (7)--(8) in the representation
\cite[(1.71)]{B-CMP122b} gives the right member $\mathfrak{M}_k(X)$ of
\cite[(1.79)]{B-CMP122b}, read at $A_{\flat}$ with the cover
\eqref{5.2:eq-sourced-modulus-3}. This gives \eqref{5.2:eq-sourced-modulus-4}.

The analyticity of $\mathbf{T}'^{\,w}_k(X,(\mathbf{U},\mathbf{J}))F$ in $(\mathbf{U},\mathbf{J})$
and in $w$ follows from Morera's theorem under the majorant so obtained.

Finally, let the history of $X$ belong to the family of hypothesis~(b), and let $X$ be of
class (i) or (ii). Hypothesis~(c2) gives \eqref{5.2:eq-sourced-modulus-2}, which is
\cite[(1.89)]{B-CMP122b} on this family. Inserting it into \eqref{5.2:eq-sourced-modulus-4}
yields \eqref{5.2:eq-sourced-modulus-a}.
\end{proof}

\emph{The per-cell value constant.}\label{5.2:prf-cell-constant}
Let $k\le K$ and $0\le n\le k$. Denote by $\mathsf{K}_{\mathrm{cell}}(k,n)$ the per-cell constant
of the volume clause of \eqref{5.2:eq-resummation-input-a}, in its strip reading, at a cell of
zone $n$. The cells are counted in level-$n$ units, as the sum $\sum_j|\Gamma^0_j\cap Y|$ over a
component $Y$ of the enlarged large-field region $Z^{\sim}$ counts them
in \cite[(1.69)]{B-CMP122b}. Thus $\mathsf{K}_{\mathrm{cell}}(k,n)$ is the value constant of
\eqref{5.2:eq-resummation-input-e}. It is formed by summing over the rows the last column of the
table of per-cell sizes, with the factor $C_{\sharp}$ on the first three rows, the fifth row and
the boundary row $\partial$ of that table, and by adding the chain-row pile of the chain of the
present step. We show that
\begin{equation}\label{5.2:eq-cell-constant-a}
\mathsf{K}_{\mathrm{cell}}(k,n)\le C_{\sharp}\bigl[\mathsf{K}_{\mathrm{cell}}^0
+3\bigl(\mathsf{N}_0(\mathsf{T}_n)+3\bigr)C_0^{\sharp}\bigr],
\qquad
\mathsf{N}_0(\mathsf{T}):=4+2r_{\mathrm{pr}}\log_L(\mathsf{T}+1),
\end{equation}
where $\mathsf{T}_n=\log g_n^{-2}$. Here $\mathsf{K}_{\mathrm{cell}}^0$ is a number that depends on
$d,L,M,\kappa,\delta$ only and is free of $g$. The frozen terms of the live arrested pieces are
not part of $\mathsf{K}_{\mathrm{cell}}(k,n)$. They form the separate functional
$\mathsf{P}(\mathfrak{a})$ of \eqref{5.2:eq-resummation-input-e}, which is counted extensively.

The right member of \eqref{5.2:eq-cell-constant-a} has two members. The first comes from the rows
other than the row of chain terms. The second comes from the chain of the present step. We treat
them in this order and then combine them.

\emph{The first member.} The rows other than the chain row contribute the number
$\mathsf{K}_{\mathrm{cell}}^0$, and it enters in four ways.
\begin{itemize}
\item It enters through the per-cell scale sums of those rows and through a constant depending on
$\kappa$ only.
\item It enters through the constants in the bound on the stored boundary terms $\mathbf{B}$.
\item It enters through the $\mathbf{R}$-born boundary terms $\mathbf{B}'$, which are taken on their
labels. For these the prefactor obeys $e^{\lambda_\theta}c_1^{\flat}\le 1$, in two steps.
First, $c_1^{\flat}=C_{99}c_1\le g^{\kappa_0}$. This follows from the third-root ceiling on the
strip smallness in \eqref{5.2:eq-resummation-window-a} and from the ceiling on $\gamma$ that bounds
$C_{99}e^{-p_0(g_k)}$, both read with the constant $C_{99}$. Second,
$e^{\lambda_\theta}g^{\kappa_0}\le 1$ by the ceiling on $\gamma$ that bounds this product.
Together,
\begin{equation*}
e^{\lambda_\theta}c_1^{\flat}\le e^{\lambda_\theta}g^{\kappa_0}\le 1.
\end{equation*}
With the prefactor bounded by $1$, the contribution of the $\mathbf{B}'$ terms to
$\mathsf{K}_{\mathrm{cell}}^0$ is free of $g$.
\item It enters through the charge of the boundary-type terms. Every such term $F$ of level $j$ is
charged on its label $A_F$, at a level-$j$ cube of $A_F$ lying in $(Z_j^{\boxplus})^{(1)}$. Here
$Z_j$ and $Z_j^{\boxplus}$ are the sets of the current label. By
\eqref{5.2:eq-current-label-anchors-a} such a cube exists for every stored boundary-type term
alive at the current time.
\end{itemize}
Write $N_j(\cdot)$ for the number of level-$j$ cubes of a set. Counting the anchors in
$(Z_j^{\boxplus})^{(1)}$ instead of $Z_j^{\boxplus}$ costs the factor $3^d$ on the member of the
chain row, since, by clause (a) of the volume proposition for wrapped units,
\begin{equation*}
N_j\bigl((Z_j^{\boxplus})^{(1)}\bigr)\le 3^d N_j(Z_j^{\boxplus}).
\end{equation*}
This factor does not appear in \eqref{5.2:eq-cell-constant-a}. It is contained in
$c_V^S=4e^{4}\,3^{d}$, the number by which $\mathsf{K}_{\mathrm{cell}}(k,n)$ and
$\mathsf{P}(\mathfrak{a})$ are multiplied when they are read in \eqref{5.2:eq-resummation-input-e}.

\emph{The second member: the chain of the present step, per cell.} By the lemma bounding the
number of stages whose boundary families meet a point, at most $N_0(g_k)+3$ boundary families of
the stages of the present step's chain are present at any point. This bound rests on the
observation in \cite{B-CMP122a} that ``the basic localization regions $Z_{h+i}\setminus Z'_{h+i}$
contain the common domain \dots\ hence the bounds from the corresponding terms cumulate. There are
at most $N_0$ such terms''. The count used below is the lemma's, $N_0(g_k)+3$. Moreover, by
\cite{B-CMP122a},
\begin{equation*}
N_0(g)\le 4+2r_{\mathrm{pr}}\log_L\log g^{-2}.
\end{equation*}

Each term $\mathbb{C}^{(i)}_k(X,\mathcal{S};w)$ of localization level $m$ is charged exactly once.
It is charged at an anchor cube $\square(X)\in\pi_m$ of the set
$X\cap\Omega_m\cap\Omega^c_{m+1}$.

\emph{The anchor exists.} The set $X\cap\Omega_m\cap\Omega^c_{m+1}$ is non-empty by
\cite[(1.63)]{B-CMP122a}. That display gives $X\subset\Omega^c_{m+1}$ for some $m\ge j+1$,
$X\cap\Omega_m\neq\emptyset$ and $X\in\mathbb{D}_m$, where $j$ is the integrated level of the
stage that produced the term.

\emph{The anchor is a cell of the old determining set.} The anchor cube is a cell of zone $m$ of
$\{\Gamma^0_j\}$, the set obtained after the last renormalization transformation. It is with this
set that \cite[(1.69)]{B-CMP122b} counts. For the running chain, the bound
\cite[(1.69)]{B-CMP122b} is exactly this charge.

\emph{Each stage contributes at most $C_0^{\sharp}$ at an anchor.} Fix an anchor. The sum over $X$
of the terms of one stage charged there is the supremand of the anchored-domain norm
$\|\cdot\|^{A}$ at that anchor, by the identification of the sum over domains at an anchor with
that norm in the payment lemma for chain terms. This norm dominates the supremum of the terms,
since $0\in\mathfrak{D}_{\mathcal{S}}(X)$. By clause (a) of the proposition on the
$\mathbf{R}$-step at complex source, the stored chain terms of each stage have norm
$\|\mathbb{C}^{(i)}_k\|^{A}\le C_0^{\sharp}$. Hence each stage contributes at most $C_0^{\sharp}$ at
the anchor.

\emph{Conclusion per cell.} A cell of zone $m$ is charged by at most $N_0(g_k)+3$ stages. In the
form in which the lemma bounding the number of stages meeting a point states its bound, this gives
at most
\begin{equation*}
3\bigl(N_0(g_k)+3\bigr)C_0^{\sharp}
\end{equation*}
per cell.

\emph{The count at the coupling of the cell's zone.} Let $n\le k$ be the zone of the cell. We show
that $N_0(g_k)\le\mathsf{N}_0(\mathsf{T}_n)$. Write $x_n=g_n^{-2}$, so that
$\mathsf{T}_n=\log x_n$ and $\mathsf{T}_k=\log x_k$. The running couplings satisfy
$x_n\ge x_k-2c_2$ for $n\le k$, by their comparison with the reference recursion. The list
$\gamma_J$ of conditions on $\gamma$ contains $x_k\ge 4c_2$, so that $2c_2/x_k\le\tfrac12$. Hence
\begin{equation*}
\mathsf{T}_n\ge\log(x_k-2c_2)=\mathsf{T}_k+\log\Bigl(1-\frac{2c_2}{x_k}\Bigr)
\ge\mathsf{T}_k-\log 2\ge\mathsf{T}_k-1 .
\end{equation*}
This is the same argument that gives $\mathsf{T}_k\le\check{\mathsf{T}}_k+1$. We therefore have
$\mathsf{T}_k\le\mathsf{T}_n+1$. Since $\log_L$ is increasing, it follows that
\begin{equation*}
N_0(g_k)\le 4+2r_{\mathrm{pr}}\log_L\mathsf{T}_k
\le 4+2r_{\mathrm{pr}}\log_L(\mathsf{T}_n+1)=\mathsf{N}_0(\mathsf{T}_n).
\end{equation*}
This holds for every zone $n\le k$. No relation between the zone $n$ and the localization level $m$
of the charged terms is used. The second member of $\mathsf{K}_{\mathrm{cell}}(k,n)$ is therefore at
most
\begin{equation*}
3\bigl(\mathsf{N}_0(\mathsf{T}_n)+3\bigr)C_0^{\sharp}.
\end{equation*}

\emph{Combination.} The first member is at most $C_{\sharp}\mathsf{K}_{\mathrm{cell}}^0$. We add the
second member. As $C_{\sharp}\ge1$, the second member may be multiplied by $C_{\sharp}$ as well.
This gives \eqref{5.2:eq-cell-constant-a}. The frozen terms of live arrested pieces have not been
charged to any cell. They are carried by $\mathsf{P}(\mathfrak{a})$ in
\eqref{5.2:eq-resummation-input-e}.

\medskip
\noindent\emph{The frozen contribution of live arrested regions.}\label{5.2:prf-arrested-regions}
The frozen terms of the live arrested pieces are not charged through the per-cell value constant.
They form the separate functional $\mathsf{P}(\mathfrak{a})$ of the boundary-term input at the
resummation site, \eqref{5.2:eq-resummation-input-e}, and they are bounded extensively. Each
frozen term is charged once, at a cube of the band of its own stage, and the charges are counted
stage by stage, as a volume at the term's own scale.

Let $C$ be the component of the large-field region $Z$ over which the sum in the preceding part of
this proof, on the per-cell value constant, runs. Let $\mathfrak{a}$ be a live arrest
with $W_{\mathfrak{a}}\subset C$. Here $k_a$ is the step of the arrest, and $X_{\mathfrak{a}}$ is
the component of the large-field region at step $k_a$ to which the arrest belongs. Further,
$W_{\mathfrak{a}}=X_{\mathfrak{a}}\cap\Omega^c_{k_a}$ is its live piece, $n^+_a$ is the number
of its stages, and $h_a$ is the level after which its stages begin. By the definition of live
arrests, $W_{\mathfrak{a}}$ is a component of
$\Omega^c_{k_a}\cap Z$, so it lies in a single component of $Z$, and that component is $C$.

The stages of $\mathfrak{a}$ integrate the fields $A_{h_a+1},A_{h_a+2},\dots$ in succession. The
stage $n\le n^+_a$ therefore integrates the level
\begin{equation*}
  j_n=h_a+n-1,
\end{equation*}
and distinct stages integrate distinct levels. The frozen terms of stage $n$ are of two kinds.
\begin{itemize}
\item The first kind are the chain terms $\mathbb{C}^{(n)}_{k_a}(X,\mathcal{S})$ of localisation
  levels $m\ge j_n+1$. For these $X\in\mathbb{D}_m$, and $X$ meets
  $Z_{j_n+1}\setminus Z'_{j_n+1}$ \cite[(1.63)]{B-CMP122a}, $Z'_{j_n+1}$ being the region so
  denoted there. They obey
  $\|\mathbb{C}^{(n)}_{k_a}\|^{A}\le C_0^{\sharp}$ by the bound on the chain terms at step
  $k_a$. That is, for each level $m$, the supremum over the anchor cubes $\square\in\pi_m$ of the
  weighted sum, over $X\ni\square$, of $\sup|\mathbb{C}^{(n)}_{k_a}(X,\mathcal{S})|$ is at most
  $C_0^{\sharp}$.
\item The second kind are the change-of-variables terms of stage $n$. There is one almost-local
  term for each block of the stage's cover on its band, and each has modulus at most $1$, since
  such terms are bounded by any positive power of the running coupling, here $g_{k_a}$
  \cite{B-CMP119}, \cite{B-CMP122a}.
\end{itemize}

\emph{The bands and the charge.} The band of stage $n$ is the stage region
$R^{(n)}_{\mathfrak{a}}$ of the lemma on the bands of arrested chains,
\begin{equation*}
  B_n:=R^{(n)}_{\mathfrak{a}}=\bigl(Z_{j_n+1}\setminus Z''_{j_n+1}\bigr)\cap X_{\mathfrak{a}} .
\end{equation*}
Here $Z_{j+1}=\Lambda^c_{j+1}\cap Z$ in the notation of \cite{B-CMP122a}, so $Z_{j_n+1}$ is the
large-field domain of level $j_n+1$ of the label. By that lemma, every chain
term of stage $n$ meets its own band $B_n$. We charge each $\mathbb{C}^{(n)}_{k_a}(X,\mathcal{S})$
of level $m$ once, at an anchor cube $\square(X)\in\pi_m$ of $X$ that meets $B_n$. At a fixed pair
$(m,\square)$, the sum of $\sup|\mathbb{C}^{(n)}_{k_a}(X,\mathcal{S})|$ over the terms charged
there is bounded by the supremand of the $\|\cdot\|^{A}$-norm at $\square$, hence by
$C_0^{\sharp}$. Therefore
\begin{equation*}
  \sum_{X}\sup\bigl|\mathbb{C}^{(n)}_{k_a}(X,\mathcal{S})\bigr|
  \le C_0^{\sharp}\cdot\#\bigl\{(m,\square):\ j_n+1\le m\le k_a,\
  \square\in\pi_m,\ \square\cap B_n\neq\emptyset\bigr\}.
\end{equation*}

\emph{The cube count.} The level index is related to the spacing as follows: a level $j<k$ has
spacing $L^{-(k-j)}$ in units of level $k$. Hence, for $m\ge j_n+1$, a cube of $\pi_m$ is coarser
than, or equal to, the cubes of $\pi_{j_n+1}$, and it is a union of them. The band $B_n$ is a union
of cubes of level $j_n+1$, since the regions involved are unions of $M$-cubes in the
corresponding scales \cite[pp.~178--179]{B-CMP122a}. Consequently each cube of $\pi_m$ that meets
$B_n$ contains a cube of $\pi_{j_n+1}$ lying in $B_n$. Choose one such cube for each $\square$.
Distinct cubes of $\pi_m$ are disjoint, so the map from $\square$ to the chosen cube is injective.

For a set $Y$, let $\#_{j+1}(Y)$ be the number of cubes of $\pi_{j+1}$ meeting $Y$, and let
$\#^{M\check{R}}_{j+1}(Y)$ be the number of $M\check{R}_{j+1}$-cubes of level $j+1$ meeting $Y$.
The injection gives, for every $m$,
\begin{equation*}
  \#\{\square\in\pi_m:\ \square\cap B_n\neq\emptyset\}\le\#_{j_n+1}(B_n).
\end{equation*}
The levels $m$ with $j_n+1\le m\le k_a$ number at most $k_a-j_n\le N(g_{k_a})$.

Put $j:=j_n$. The set $Z_{j+1}$ is a union of $M\check{R}_{j+1}$-cubes \cite[p.~256]{B-CMP119},
and each of these holds $\check{R}_{j+1}^d$ cubes of $\pi_{j+1}$. By the cube count for connected
unions, a connected union $X$ of $M\check{R}_{j+1}$-cubes consists of at most
$\mathsf{c}_d\,(d'_{j+1}(X)+1)$ of them, where
\begin{equation*}
  \mathsf{c}_d:=2\cdot 3^d .
\end{equation*}
Let $X$ run over the components of $Z_{j+1}$ contained in $X_{\mathfrak{a}}$. This covers all of
$Z_{j+1}\cap X_{\mathfrak{a}}$: a component of $Z_{j+1}\subset Z_{k_a}$ that meets the component
$X_{\mathfrak{a}}$ of $Z_{k_a}$ lies in it. Since $B_n\subset Z_{j+1}\cap X_{\mathfrak{a}}$, we
obtain
\begin{align*}
  \#_{j+1}(B_n)
  &\le\#_{j+1}\bigl(Z_{j+1}\cap X_{\mathfrak{a}}\bigr)
   \le\check{R}_{j+1}^d\,\#^{M\check{R}}_{j+1}\bigl(Z_{j+1}\cap X_{\mathfrak{a}}\bigr)\\
  &\le\check{R}_{j+1}^d\,\mathsf{c}_d\sum_{X}\bigl(d'_{j+1}(X)+1\bigr).
\end{align*}
By the description of the second kind of frozen terms above, the change-of-variables terms of
stage $n$ contribute at most $\#_{j_n+1}(B_n)$, which falls within the same count.

\emph{The bound on $\mathsf{P}(\mathfrak{a})$.} Set
\begin{equation}\label{5.2:eq-arrested-regions-1}
\begin{split}
  \mathsf{P}(\mathfrak{a}):=\sum_{n\le n^+_a}\Bigl[&\sum_{m,X}
  \sup\bigl|\mathbb{C}^{(n)}_{k_a}(X,\mathcal{S})\bigr|\\
  &+\sum\bigl|\text{change-of-variables terms of stage }n\bigr|\Bigr].
\end{split}
\end{equation}
Combining the three preceding estimates gives
\begin{equation}\label{5.2:eq-arrested-regions-2}
  \mathsf{P}(\mathfrak{a})\le\mathsf{c}_d\bigl(C_0^{\sharp}N(g_{k_a})+1\bigr)
  \sum_{n\le n^+_a}\check{R}^d_{j_n+1}\sum_{X}\bigl(d'_{j_n+1}(X)+1\bigr),
\end{equation}
where the inner sum is over the components $X$ of $Z_{j_n+1}$ contained in $X_{\mathfrak{a}}$.

The factor $N(g_{k_a})$, which bounds the number of localisation levels, is the only
coupling-dependent count in \eqref{5.2:eq-arrested-regions-2}. It is absorbed into one power of
$\check{R}_{j_n+1}$ by the following line:
\begin{equation}\label{5.2:eq-arrested-regions-3}
\begin{split}
  N(g_{k_a})&\le C_N\mathsf{T}_{k_a}^{r_{\mathrm{pr}}}
  \le C_N\bigl(\check{\mathsf{T}}_{k_a}+1\bigr)^{r_{\mathrm{pr}}}
  \le 2^{r_{\mathrm{pr}}}C_N\check{\mathsf{T}}_{k_a}^{r_{\mathrm{pr}}}\\
  &\le 2^{r_{\mathrm{pr}}}C_N\check{R}_{k_a}
  \le 2^{r_{\mathrm{pr}}}C_N\check{R}_{j_n+1}.
\end{split}
\end{equation}
The inequalities $N\le C_N\mathsf{T}^{r_{\mathrm{pr}}}$, $\mathsf{T}_k\le\check{\mathsf{T}}_k+1$,
$\check{\mathsf{T}}\ge 1$ and $\check{R}_k\ge\check{\mathsf{T}}_k^{r_{\mathrm{pr}}}$ used in
\eqref{5.2:eq-arrested-regions-3} are
among the standing inequalities on the scaffold quantities $N$, $\mathsf{T}_k$,
$\check{\mathsf{T}}_k$, $\check{R}_k$ fixed in this chapter.
The inequalities are justified as follows.
\begin{itemize}
\item The first is the bound $N\le C_N\mathsf{T}^{r_{\mathrm{pr}}}$ on the number of preparatory
  stages.
\item The second is $\mathsf{T}_k\le\check{\mathsf{T}}_k+1$.
\item The third uses $\check{\mathsf{T}}\ge 1$, so that
  $\check{\mathsf{T}}_{k_a}+1\le 2\check{\mathsf{T}}_{k_a}$.
\item The fourth is $\check{R}_k\ge\check{\mathsf{T}}_k^{r_{\mathrm{pr}}}$.
\item The fifth holds because $\check{R}_j=R(\min_{i\le j}x_i)$ is non-increasing in the index
  $j$: the minimum over a larger index set is smaller, and $R(\cdot)$ is non-decreasing. Since
  $j_n+1=h_a+n\le k_a$, this gives $\check{R}_{j_n+1}\ge\check{R}_{k_a}$.
\end{itemize}

We insert \eqref{5.2:eq-arrested-regions-3} into \eqref{5.2:eq-arrested-regions-2}. Because
$C_0^{\sharp}\ge 1$, $C_N\ge 1$ and $\check{R}_{j_n+1}\ge 1$,
\begin{equation*}
  \mathsf{c}_d\bigl(C_0^{\sharp}N(g_{k_a})+1\bigr)
  \le\mathsf{c}_d\cdot 2\cdot 2^{r_{\mathrm{pr}}}C_NC_0^{\sharp}\,\check{R}_{j_n+1}.
\end{equation*}
Even after the further factor $2^d$ of the strip, with which the input
\eqref{5.2:eq-resummation-input-e} is read, this constant does not exceed the first summand of
$c_{\mathsf{P}}$ below:
\begin{equation*}
  \mathsf{c}_d\cdot 2^d\cdot 2\cdot 2^{r_{\mathrm{pr}}}C_NC_0^{\sharp}
  =2^{d+2+r_{\mathrm{pr}}}3^dC_NC_0^{\sharp}
  \le 2^{d+3+r_{\mathrm{pr}}}3^dC_NC_0^{\sharp}.
\end{equation*}
This yields
\begin{equation}\label{5.2:eq-arrested-regions-a}
\begin{split}
  \mathsf{P}(\mathfrak{a})&\le c_{\mathsf{P}}\sum_{n\le n^+_a}\check{R}^{d+1}_{j_n+1}
  \sum_{X}\bigl(d'_{j_n+1}(X)+1\bigr),\\
  c_{\mathsf{P}}&=2^{d+3+r_{\mathrm{pr}}}3^dC_NC_0^{\sharp}
  +27\,\mathsf{L}_0^6L^d\mathsf{c}_d\mathsf{K}^{\flat},
\end{split}
\end{equation}
with $X$ running over the components of $Z_{j_n+1}$ contained in $X_{\mathfrak{a}}$. For the terms
of \eqref{5.2:eq-arrested-regions-1} the first summand of $c_{\mathsf{P}}$ suffices. The second
summand is the constant of the weighted mass of frozen cells,
\eqref{5.2:eq-frozen-cell-mass-a}; that mass is added to $\mathsf{P}(\mathfrak{a})$ when the
boundary-term input \eqref{5.2:eq-resummation-input-e} is read, and, since
$\check{R}_{j_n+1}\ge 1$, it is bounded by the second summand times the same double sum.

\begin{lemma}[The weighted mass of frozen cells]\label{5.2:lem-frozen-cell-mass}
Let $C$ be a component of the large-field region $Z$, and let $\mathfrak{a}$ be a live
arrest whose component $X_{\mathfrak{a}}$ satisfies $X_{\mathfrak{a}}\subset C$. Here $X_{\mathfrak{a}}$
is the component of the large-field region at the arrest's step $k_a$ that carries $\mathfrak{a}$, and
$h_a:=k_a-N(g_{k_a})$ and $k_{0,a}:=k_a-N_0(g_{k_a})\ge h_a+1$ are the levels $h$ and $k_0$ of that
step. The stages $n\le n^+_a$ of $\mathfrak{a}$ integrate the distinct levels $j_n:=h_a+n-1$. The
window of $\mathfrak{a}$ is the set of cells $\Delta'\subset X_{\mathfrak{a}}$ of level in
$]k_{0,a},k_a]$. For a cell $\Delta'$ put
$\bar{\mathsf{w}}(\Delta'):=3w^{\vee}(\Delta')$, the threshold weight of the pointed rows at $\Delta'$,
with the weights $w$ and $w^{\vee}$ of \eqref{5.2:eq-plain-interior-units-a}. Let
$\mathsf{K}^{\flat}\bar{\mathsf{w}}(\Delta')^3$ be their per-cell pile, with $\mathsf{K}^{\flat}$ as in
\eqref{5.2:eq-column-sum-pointed-a}. Recall that $Z_{j+1}=\Lambda^c_{j+1}\cap Z$ is the
level-$(j+1)$ large-field domain of the label, and put $\mathsf{c}_d:=2\cdot 3^d$. Then
\begin{equation}\label{5.2:eq-frozen-cell-mass-a}
\sum_{\Delta'}\mathsf{K}^{\flat}\bar{\mathsf{w}}(\Delta')^3
\le 27\mathsf{L}_0^6L^d\mathsf{c}_d\mathsf{K}^{\flat}
\sum_{n\le n^+_a}\check{R}^d_{j_n+1}\sum_{X}\bigl(d'_{j_n+1}(X)+1\bigr).
\end{equation}
In \eqref{5.2:eq-frozen-cell-mass-a} the cell $\Delta'$ runs over the frozen window cells
$\Delta'\subset X_{\mathfrak{a}}$ with $w(\Delta')>1$. For each $n$, the set $X$ runs over the
components of $Z_{j_n+1}$ contained in $X_{\mathfrak{a}}$.
\end{lemma}

\begin{proof}
The arrest $\mathfrak{a}$ carries on $X_{\mathfrak{a}}$ a level assignment $\ell_{\mathfrak{a}}$ and a
weight assignment $w_{\mathfrak{a}}$, fixed at step $k_a$; by the freezing statement for live arrests,
they are not changed by later steps. Hence a window cell $\Delta'$ has level
$n':=\ell_{\mathfrak{a}}(\Delta')\in\,]k_{0,a},k_a]$. Its weight is $w=\mathsf{L}_0^{2e}$, in one of the two
branches of the weight \eqref{5.2:eq-plain-interior-units-a} at the arrest's step. These branches are
those of \cite[(1.64)--(1.67), p.~190]{B-CMP122a}, with Ba{\l}aban's step parameter there (his $n$)
equal to $N(g_{k_a})$:
\begin{itemize}
\item[(i)] $e=n'-k_{0,a}$, on a stratum of level $n'<k_a$ or on the top shell off $\Omega_{k_{0,a}+1}$.
Under these cells the pre-arrest label has levels $\le k_{0,a}$. The pre-arrest label is the label
that the preparation of $\mathbf{R}_{k_a}$ coarsened.
\item[(ii)] $e=k_a-m'$, on the top zone's old ring of pre-arrest level $m'\in\,]k_{0,a},k_a]$.
\end{itemize}
An arrest may stop at a stage of level $j_a\le k_a$. In that case branch (ii) reads $e=j_a-m'$ with
$m'\in\,]k_{0,a},j_a]$. Then $e$ only decreases, and the count below is unchanged.

Passing from $w$ to $w^{\vee}$ raises $e$ by at most one. We also have $\mathsf{L}_0^2\le L/3$ and
$d=4>3$. Therefore
\begin{equation}\label{5.2:eq-frozen-cell-mass-1}
\bar{\mathsf{w}}(\Delta')^3\le 27\mathsf{L}_0^{6(e+1)}
=27\mathsf{L}_0^6\bigl(\mathsf{L}_0^2\bigr)^{3e}
\le 27\mathsf{L}_0^6(L/3)^{3e}\le 27\mathsf{L}_0^6L^{de}.
\end{equation}

A cell of level $n'$ contains $L^{d(n'-j)}$ cells of any finer level $j$. In branch (i) the pre-arrest
cells under $\Delta'$ have level $j\le k_{0,a}$. Hence $\Delta'$ contains at least
$L^{d(n'-k_{0,a})}=L^{de}$ of them. In branch (ii) the cell $\Delta'$ lies in the top zone, whose level
$n'$ is $k_a$, or $j_a$ if the arrest stopped early. It therefore contains exactly
$L^{d(n'-m')}=L^{de}$ pre-arrest cells of level $m'$. Combined with \eqref{5.2:eq-frozen-cell-mass-1},
this gives
\begin{equation*}
\mathsf{K}^{\flat}\bar{\mathsf{w}}(\Delta')^3
\le 27\mathsf{L}_0^6\mathsf{K}^{\flat}\cdot\#\{\text{pre-arrest cells in }\Delta'\}.
\end{equation*}
The window cells are disjoint, so we may sum over them. The sum is at most $27\mathsf{L}_0^6\mathsf{K}^{\flat}$
times the number of pre-arrest cells under the window of $X_{\mathfrak{a}}$, of the levels $j$
re-made there, that is, of the levels $j$ whose cells a stage of $\mathfrak{a}$ replaced by coarser
cells.

Take a pre-arrest cell of level $j$. It is a cell of zone $j$ of the pre-arrest label, so it lies off
$\Omega_{j+1}$, that is, in $\Omega^c_{j+1}\subset\Lambda^c_{j+1}$. Since it lies inside
$X_{\mathfrak{a}}\subset Z$, it lies in $Z_{j+1}=\Lambda^c_{j+1}\cap Z$.

Let $\#_{j+1}(S)$ denote the number of $\pi_{j+1}$-cubes meeting a set $S$. A level-$(j+1)$ cube
contains $L^d$ level-$j$ cubes. Hence the number of these cells is at most
$L^d\cdot\#_{j+1}(Z_{j+1}\cap X_{\mathfrak{a}})$. We now apply the cube count established in the proof
of the bound \eqref{5.2:eq-arrested-regions-a} for the frozen contribution of live arrested regions:
\begin{equation}\label{5.2:eq-frozen-cell-mass-2}
\#_{j+1}(Z_{j+1}\cap X_{\mathfrak{a}})\le\check{R}^d_{j+1}\,\mathsf{c}_d
\sum_{X}\bigl(d'_{j+1}(X)+1\bigr).
\end{equation}
Here $X$ runs over the components of $Z_{j+1}$ contained in $X_{\mathfrak{a}}$.

A level $j$ is re-made under the window only by the stage $n$ of $\mathfrak{a}$ with $j_n+1=j+1$,
because each arrest has at most one stage per level. A level that no stage re-made keeps $w=1$ and
contributes nothing to the left member of \eqref{5.2:eq-frozen-cell-mass-a}. Multiply
\eqref{5.2:eq-frozen-cell-mass-2} by $27\mathsf{L}_0^6\mathsf{K}^{\flat}L^d$, put $j=j_n$, and sum
over the stages $n\le n^+_a$. This gives \eqref{5.2:eq-frozen-cell-mass-a}.
\end{proof}

The bound \eqref{5.2:eq-frozen-cell-mass-a} enters the bound \eqref{5.2:eq-arrested-regions-a} for the
frozen pile as a further member; that is, $\mathsf{P}(\mathfrak{a})$ is replaced by
\begin{equation*}
\mathsf{P}(\mathfrak{a})+\sum_{\Delta'}\mathsf{K}^{\flat}\bar{\mathsf{w}}(\Delta')^3 .
\end{equation*}
We have $\check{R}_{j_n+1}\ge 1$, so $\check{R}^d_{j_n+1}\le\check{R}^{d+1}_{j_n+1}$. Hence this member
lies in the same third member of Ba{\l}aban's volume allowance, in which every large-field domain $Z_j$
contributes at most $O(1)M^dR_j^{d+1}d'_j(Z_j)$ \cite{B-CMP122b}. This is the origin of the second
summand $27\mathsf{L}_0^6L^d\mathsf{c}_d\mathsf{K}^{\flat}$ of $c_{\mathsf{P}}$ in
\eqref{5.2:eq-arrested-regions-a}. That summand is a number fixed after $M$, since $\mathsf{K}^{\flat}$
is. None of an age, the $N_0$ of the arrest, the number of arrests or the window size enters this
constant. The exponent $e$ is spent against the cell count, exactly as for the window of the current
step's own chain, where it is spent against the cells of the old determining set $\{\Gamma^0_j\}$. The
ceiling on $\gamma$ governing the volume sums is taken with this value of $c_{\mathsf{P}}$.

\begin{proof}[Value constants within Ba{\l}aban's per-cell allowance]\label{5.2:prf-per-cell-allowance}
In this part of the proof the frozen piles $\mathsf{P}(\mathfrak{a})$ of the live arrests and the
per-cell value constant $\mathsf{K}_{\mathrm{cell}}(k,n)$ are placed inside the allowance that
Ba{\l}aban makes for the constants of the large-field operation. That allowance is the sentence of
\cite[pp.~380--381]{B-CMP122b}
\begin{quotation}
every large field domain $Z_j$ contributes the constant
\begin{equation*}
O(1)\log g_j^{-2}|Z_j| \le O(1)\log g_j^{-2}(MR_j)^d d'_j(Z_j) \le O(1)M^dR_j^{d+1}d'_j(Z_j)
\end{equation*}
\end{quotation}
and we speak of the first, second and third member of this chain of inequalities. The constant
$O(1)$ in it is read as $C_\sharp O(1)+1$, as fixed in the conditions that accompany the
large-field bound.

\emph{The frozen piles.} Fix a component $C$ of $Z$ and consider the live arrests $\mathfrak{a}$
whose component $X_{\mathfrak{a}}$ lies in $C$. By \eqref{5.2:eq-arrested-regions-a}, the pile
$\mathsf{P}(\mathfrak{a})$ is bounded by a sum over the stages $n \le n^+_{\mathfrak{a}}$ of the
arrest, where stage $n$ integrates the level $j_n = h_{\mathfrak{a}}+n-1$, and over the components
$X \subset X_{\mathfrak{a}}$ of the large-field domain $Z_{j_n+1}$; the pair $(\mathfrak{a},n)$
charges such a component $X$ by at most
$c_{\mathsf{P}}\check{R}_{j_n+1}^{d+1}(d'_{j_n+1}(X)+1)$, the constant $c_{\mathsf{P}}$ being read,
like $\mathsf{K}_{\mathrm{cell}}(k,n)$ below, with the factor $c_V^S$ of the volume sums
\eqref{5.2:eq-resummation-input-e}.
Now fix a level $j$ and a component $X$ of the level-$j$ large-field domain
$Z_j = \Lambda_j^c \cap Z$. The pairs $(\mathfrak{a},n)$ that charge $X$ are those with
$j_n+1 = j$ and $X \subset X_{\mathfrak{a}}$, and three facts restrict them.
\begin{itemize}
\item The live arrests whose components contain $X$ are nested,
$X_{\mathfrak{a}_1} \subset X_{\mathfrak{a}_2} \subset \cdots$: by the first clause of the structure
lemma for live arrests and the last part of that clause, components that are not nested are
disjoint, whereas all of these contain $X$.
\item Their band windows are pairwise disjoint. The band window of $\mathfrak{a}$ is the set of
levels $j_n+1$, $n \le n^+_{\mathfrak{a}}$, that is
$\{h_{\mathfrak{a}}+1,\dots,h_{\mathfrak{a}}+n^+_{\mathfrak{a}}\}$, and it is contained in
$[h_{\mathfrak{a}}+1,k_{\mathfrak{a}}]$; if $X_{\mathfrak{a}} \subset X_{\mathfrak{a}'}$ are the
components of two distinct arrests $\mathfrak{a}$, $\mathfrak{a}'$ of the nested family, then the
second clause of the same lemma gives $h_{\mathfrak{a}'} \ge k_{\mathfrak{a}}$, so the window of
$\mathfrak{a}'$ begins above the window of $\mathfrak{a}$.
\item Each arrest has at most one stage per level, since the levels $j_n = h_{\mathfrak{a}}+n-1$
of its stages are distinct.
\end{itemize}
By the second fact the level $j$ lies in the band window of at most one of these arrests, and by
the third at most one stage of that arrest has $j_n+1 = j$. Hence at most one pair charges $X$.
We nevertheless keep a factor $2$ in the bound below as a margin; by the preceding count it is not needed. Thus $X$
is charged by the frozen piles by at most
$2c_{\mathsf{P}}\check{R}_j^{d+1}(d'_j(X)+1)$ in all, and summing over the components $X$ of
$Z_j$,
\begin{equation*}
\sum_{X}2c_{\mathsf{P}}\check{R}_j^{d+1}\bigl(d'_j(X)+1\bigr)
= 2c_{\mathsf{P}}M^{-d}\cdot M^d\check{R}_j^{d+1}
\Bigl(\sum_X d'_j(X)+\#\{\text{components of } Z_j\}\Bigr).
\end{equation*}
Here $d'_j(Z_j) = \sum_X d'_j(X)$, the quoted line being read component by component and summed
over the components $X$ of $Z_j$.
So every large-field domain $Z_j$ of the label is charged, beyond the share of Ba{\l}aban's own
terms and the share of the chain of the present step, at most $2c_{\mathsf{P}}M^{-d}$ times
\begin{equation*}
M^d\check{R}_j^{d+1}\bigl(d'_j(Z_j)+\#\{\text{components of } Z_j\}\bigr).
\end{equation*}
This lies inside the third member
$O(1)M^dR_j^{d+1}d'_j(Z_j)$ of the quoted line. Indeed, its $O(1)$ is $C_\sharp O(1)+1$, which
contains $2c_{\mathsf{P}}M^{-d}$ because the constant $C_{\mathfrak{d}}$, the part of
$C_{\mathfrak{d}}^{\flat} = C_{\mathfrak{d}}+2^{d+2}(1+dL^d)$ in the per-step charge
$c_n = C_{\mathfrak{d}}^{\flat}M^d\check{R}_n^{d+1}$ of the large-field contour theorem, is enlarged
to dominate every per-step per-unit-$d'$ charge. The summand $1$ per component is Ba{\l}aban's own
$(d'_j+1)$, which appears in the exponents of his gain factors at the same place of
\cite[pp.~380--381]{B-CMP122b}. No condition on $\gamma$ is used here, and none of $k$, $K$, an
age, the number of arrests, $N_0$ or the volume enters. The pile $\mathsf{P}(\mathfrak{a})$ is
extensive at each frozen term's own scale, exactly as the old boundary terms are bounded by
the level-wise volumes $O(1)\sum_j|\Gamma_j|$ in \cite[(2.47)--(2.49)]{B-CMP119} and never by a
constant per current cell.

\emph{Why the frozen piles are not charged per cell.} Inside a live piece $W_{\mathfrak{a}}$ the
processed bands of the arrest were re-levelled by the stages of $\mathfrak{a}$ itself: the basic
localisation regions contain the common domain \cite{B-CMP122a}, and the frozen structure of the
arrest assigns the current level. A cell of the current determining set there may therefore be
coarser than the cubes of a frozen term by a factor up to
$L^{N_0(g_{k_{\mathfrak{a}}})-1} = \check{R}_{k_{\mathfrak{a}}}/M$ per side
\cite{B-CMP122a}. A constant per current cell would then carry a factor
$\check{R}_{k_{\mathfrak{a}}}^d$, which is not of the form $O(1)\log g_j^{-2}$ per cell. The
extensive charge uses nothing of this.

\emph{What the other places read of the frozen piles.} At the other places of the step where they
enter, the frozen piles are read per point, at each unit's own level $m$: at the site of the
large-field operation $\mathbf{T}$ and at the preparatory site, through the chain-term entry of the
size table of all units, with the constant $6(\mathsf{N}_0(\mathsf{T}_m)+3)C_0^{\sharp}$ (at most two
chains per level at a point, the cumulation count for chain stages being applied to each), and at
the localisation site through the weight supremum $\mathsf{W}$ of Stage~A. None of these reads
$\mathsf{P}(\mathfrak{a})$.

\emph{The per-cell value constant.} By \eqref{5.2:eq-cell-constant-a}, $\mathsf{K}_{\mathrm{cell}}(k,n)$
is a function of the couplings of the run and of the zone $n$ of the cell, in the same way as the
number $N(g)$ of preparatory stages; it depends on the coupling, and it is displayed so in
\eqref{5.2:eq-resummation-input-a} and \eqref{5.2:eq-resummation-input-e}, and again in the bound on
the last lattice. We place it inside
Ba{\l}aban's per-cell allowance $O(1)\log g_j^{-2}$ of the first member of the quoted line, the
$O(1)$ there being $C_\sharp O(1)+1$. The charge $c_n = C_{\mathfrak{d}}^{\flat}M^d\check{R}_n^{d+1}$
uses one power of $\check{R}_n$ for the logarithm, exactly as Ba{\l}aban's second inequality does:
for $r_{\mathrm{pr}} \ge 2$ and $\check{\mathsf{T}}_n \ge 2$ (which holds, since
$\check{\mathsf{T}}_n \ge \mathsf{T}$ for $n \le k \le K$ by the standing convention on the run,
$\mathsf{T} \ge \mathsf{T}_\gamma$, and $\mathsf{T}_\gamma \ge 2$ is among the conditions on $\gamma$),
\begin{equation*}
\check{R}_n \ge \check{\mathsf{T}}_n^{\,r_{\mathrm{pr}}} \ge 2\check{\mathsf{T}}_n
\ge \check{\mathsf{T}}_n+1 \ge \mathsf{T}_n = \log g_n^{-2},
\end{equation*}
by the conventions fixing $\check{R}_n$, $\check{\mathsf{T}}_n$ and $r_{\mathrm{pr}}$ (which give
$\check{R}_n \ge \check{\mathsf{T}}_n^{\,r_{\mathrm{pr}}}$ and
$\mathsf{T}_n \le \check{\mathsf{T}}_n+1$), the condition $r_{\mathrm{pr}} \ge 2$ being
Ba{\l}aban's condition $r \ge 2$ \cite{B-CMP119}, where his $r$ is the exponent $r_{\mathrm{pr}}$
of $R$; the second inequality is $\check{\mathsf{T}}_n^{\,r_{\mathrm{pr}}} \ge
\check{\mathsf{T}}_n^{\,2} \ge 2\check{\mathsf{T}}_n$ for $\check{\mathsf{T}}_n \ge 2$, and the third
is $\check{\mathsf{T}}_n \ge 1$. The constant $C_{\mathfrak{d}}$ is enlarged to dominate every
per-step per-unit-$d'$ charge, read with the constants $O(1)$ of the same conditions; sharing the
allowance with Ba{\l}aban's own normalisation constants costs a factor $2$ inside that constant, of
which only the existence is used. The placement is effected by the ceiling on $\gamma$
\begin{equation}\label{5.2:eq-per-cell-allowance-a}
\begin{split}
\sup_{\mathsf{T}\ge\mathsf{T}_\gamma}\frac{1}{\mathsf{T}}\,2^dc_V^SC_\sharp\Bigl[
&\mathsf{K}^0_{\mathrm{cell}}+27\mathsf{L}_0^6\mathsf{K}^{\flat}+1\\
&+3\bigl(7+2r_{\mathrm{pr}}\log_L(\mathsf{T}+1)\bigr)C_0^{\sharp}\Bigr]
\le C_\sharp O(1)+1 .
\end{split}
\end{equation}
Only one ceiling of this kind is imposed on $\gamma$: the display restates the volume ceiling in
the list of conditions that fix the threshold $\gamma_J$, and it is that condition which is imposed. In it, $c_V^S := 4e^{4}\cdot 3^{d}$ is the factor with
which the bounds of the volume sums \eqref{5.2:eq-resummation-input-e} are read; the summand
$27\mathsf{L}_0^6\mathsf{K}^{\flat}$ is the window weight of the pointed units of the present step,
absorbed per old cell of the determining set; the weighted mass of the frozen window cells,
\eqref{5.2:eq-frozen-cell-mass-a}, does not enter here, being the second summand of $c_{\mathsf{P}}$
and carried extensively by $\mathsf{P}(\mathfrak{a})$; the summand $1$ accounts for the
strip-cell clause of \eqref{5.2:eq-resummation-input-e}, whose additional term $\varepsilon^{\mathrm{ext}}(k)$
is at most $1$, its two parts $\varepsilon^{\mathrm{ext},\sigma}(k)$ and
$\varepsilon^{\mathrm{ext},\mathsf{t}}(k)$ being at most $\tfrac12$ each by the ceilings on $\gamma$
for the interior activity and for the bracket activity, so that $\varepsilon^{\mathrm{ext}}(k)$ is
covered by the summand $1$ inside the bracket, since $C_\sharp \ge 1$; and
$3(7+2r_{\mathrm{pr}}\log_L(\mathsf{T}+1))C_0^{\sharp} = 3(\mathsf{N}_0(\mathsf{T})+3)C_0^{\sharp}$ is
the share of the chain of the present step in \eqref{5.2:eq-cell-constant-a}, since
$\mathsf{N}_0(\mathsf{T})+3 = 7+2r_{\mathrm{pr}}\log_L(\mathsf{T}+1)$. The factor $2^d$ is the one
carried by the per-cell constant under the strip reading, as shown in the proof of the per-cell
value constant. The supremand involves only quantities fixed before
$\gamma$; it is finite, and it tends to $0$ as $\mathsf{T}_\gamma \to \infty$, the bracket growing
only like $\log_L\mathsf{T}$. Hence \eqref{5.2:eq-per-cell-allowance-a} is met for all sufficiently
small $\gamma$.

\emph{Conclusion.} Let $n \le k$ be a zone whose cells are counted. By the standing convention on
the run, $x_n > g^{-2}$ for $n \le k \le K$, so $\mathsf{T}_n = \log x_n > \log g^{-2} = \mathsf{T}$,
and $\mathsf{T} \ge \mathsf{T}_\gamma$ since $g \le \gamma$. Reading
$\mathsf{K}_{\mathrm{cell}}(k,n)$ with the factor $c_V^S$, with the running window weight
$27\mathsf{L}_0^6\mathsf{K}^{\flat}$ and with the summand $1$ for the strip cells, as explained
above, the bound \eqref{5.2:eq-cell-constant-a} and \eqref{5.2:eq-per-cell-allowance-a} evaluated at
$\mathsf{T} = \mathsf{T}_n$ give
\begin{align*}
2^d\mathsf{K}_{\mathrm{cell}}(k,n)
&\le 2^dc_V^SC_\sharp\Bigl[\mathsf{K}^0_{\mathrm{cell}}+27\mathsf{L}_0^6\mathsf{K}^{\flat}+1
+3\bigl(\mathsf{N}_0(\mathsf{T}_n)+3\bigr)C_0^{\sharp}\Bigr]\\
&\le \bigl(C_\sharp O(1)+1\bigr)\mathsf{T}_n = \bigl(C_\sharp O(1)+1\bigr)\log g_n^{-2}.
\end{align*}
Hence Ba{\l}aban's line applies verbatim, zone by zone, to the chain of the present step and to the
terms counted in $\mathsf{K}^0_{\mathrm{cell}}$: each contributes one more logarithmic term inside
$O(1)\log g_j^{-2}|Z_j|$. Its third member carries the frozen piles $\mathsf{P}(\mathfrak{a})$ of
\eqref{5.2:eq-arrested-regions-a}, of every live arrest in the component, with the constant
$2c_{\mathsf{P}}M^{-d}$. Both contributions lie inside $C_{\mathfrak{d}}^{\flat}$, as before. The
counterterm pieces at complex $\hat{w}$ that are not matched are bounded by the last clause of the
conditions that accompany the large-field bound.
\end{proof}

\begin{proposition}[The large-field step at complex sources]
\label{5.2:prop-sourced-large-field-step}
Let $k\le K$ be a level, let $r$ be the radius fixed in the correlation theorem, and write
$D(0,2r)$ for the open disc of radius $2r$ about $0$ in the complex plane. Assume the following.
\begin{itemize}
\item The inductive hypothesis at complex source of the correlation theorem holds at level $k$,
after the large-field operation $\mathbf{T}$ of the $k$-th step.
\item $\sup|\zeta_j|\le \tfrac{8}{3}\,r$ for the running shifts $\zeta_j$, the supremum being the
sup norm of each shift.
\item Part (a) of the inheritance theorem for the chain terms holds, together with its corollary at
complex source. The two hypotheses of that corollary (holomorphy of the constituents in $w$ under
$w$-free majorants, and $w$-freeness of the scaffold) are established in the proof of clause (a)
below.
\item The stage-norm corollary for the chain terms holds at Ba{\l}aban's cores.
\item Lemma~\ref{5.2:lem-sourced-modulus} holds.
\item The strip theorem, its volume corollary and the combinatorial lemma for strips hold.
\item Part (b) of the charge corollary of the cascade bound holds on the modified family of
large-field regions.
\end{itemize}
Then the following hold.
\begin{itemize}
\item[(a)] For $n<N$ the chain terms of stage $n+1$ decompose as
\begin{equation}\label{5.2:eq-sourced-large-field-step-1}
\mathbb{C}^{(n+1)}_k=\sum_{(X,\mathcal{S})}\mathbb{C}^{(n+1)}_k(X,\mathcal{S};w).
\end{equation}
Here $X$ is a localisation domain as in \cite[(1.63)]{B-CMP122a}, taken among the domains of the
arrest device, and $\mathcal{S}=(\mathcal{S}_{j'})_{j<j'<k}$ is a family of declared
strong sets of rim cubes meeting $X$, where $j$ is the level in \cite[(1.63)]{B-CMP122a} for which
$X$ meets $Z_{j+1}\setminus Z'_{j+1}$. The summands are stored terms with the A-slot datum
$A^{\flat}$ of the A-slot family of the preparatory stages, and they satisfy
\begin{equation}\label{5.2:eq-sourced-large-field-step-2}
\|\mathbb{C}^{(n+1)}_k\|^{A}_{(n+1)}\le C_0^{\sharp},\qquad C_0^{\sharp}=3c_G+1 .
\end{equation}
The sum and the weight of this norm are taken on labels. The rates are
$a_k=(1+3\beta_{\mathrm{pr}})\kappa$ at level $k$ and $a_\ell=(1-\beta_{\mathrm{pr}})\kappa$ at
levels $\ell<k$. The summands have the following regularity.
\begin{itemize}
\item They are holomorphic on
\begin{equation*}
\tilde U^{(n+1)c,\Sigma}_k(X)\times\mathcal{W}_k(X)
\times\{\text{exterior weak $A$-coordinates at radii }\varrho_{j'}\}.
\end{equation*}
\item They are real-bounded and continuous on the strong upper boxes $\mathfrak{b}^{\sharp}$.
\item They are stored terms, with the same constant and rates, on
$\mathfrak{D}^{\flat}_{n+1}(X)\times\mathfrak{D}_{\mathcal{S}}(X)$.
\item They are stored terms, with the same constant and rates, on
$\mathfrak{D}^{\flat}_{n+1}(X)\times\mathfrak{D}_{\mathcal{S}}(X)\times\mathcal{W}_k(X)$,
holomorphically in $w$.
\end{itemize}
For a chain term built on wrapped constituents, $X$ is its label, and the statement on
$\mathfrak{D}^{\flat}_{n+1}(X)$ is read over $X^{+}$. Every constant in this clause is
independent of $n$, $N$, $N_0$, $\check{R}_k$, $k$, $K$, the volume and $g$. In particular the
summands obey \cite[(1.70)]{B-CMP122b}.
\item[(b)] The source part of the Wilson action at the resummation site decomposes as in
\eqref{5.3:eq-wilson-source-a}.
\item[(c)] Let $Z$ be the class of components of $\Lambda_k^c$ treated at the $k$-th step, and
let $Z^{\sim}$ be its enlargement by one layer of $M\check{R}_k$-cubes. The activities
$V(Y,(\mathbf{U},\mathbf{J});w)$ of \cite[(1.68)]{B-CMP122b} are analytic on the complex tube
of configurations $(\mathbf{U},\mathbf{J})$ over $Y$, as in Lemma~\ref{5.2:lem-sourced-modulus},
times $\mathcal{W}_k(Y)$. At complex source they satisfy the following forms of
\cite[(1.68)--(1.69)]{B-CMP122b}, with $\alpha$ the smallness parameter of
\cite[(1.68)]{B-CMP122b}. For $Y\setminus Z^{\sim}\neq\emptyset$,
\begin{equation}\label{5.2:eq-sourced-large-field-step-3}
|V(Y,(\mathbf{U},\mathbf{J});w)|\le C_{\sharp}\,\alpha\,e^{-(1+2\beta_s)\kappa\, d_{k,Z}(Y)} .
\end{equation}
For a component $Y$ of $Z^{\sim}$,
\begin{equation}\label{5.2:eq-sourced-large-field-step-4}
|V(Y,(\mathbf{U},\mathbf{J});w)|\le \sum_j\mathsf{K}_{\mathrm{cell}}(k,j)\,|\Gamma^0_j\cap Y|
+\sum_{\mathfrak{a}\ \mathrm{live},\ W_{\mathfrak{a}}\subset Y}\mathsf{P}(\mathfrak{a}) .
\end{equation}
Now replace $Z^{\sim}$ by the strip hull $Z^{\sharp}$ and $d_{k,Z}$ by
$d_{\sharp}:=d_{k,Z^{\sharp}}$, and let $\alpha^{\sharp}(k)$ be the strip smallness of the
boundary-term input at the resummation site (\ref{5.2:prf-resummation-input}). The resummed
activities $V^{\sharp}$ then satisfy, for $Y\setminus Z^{\sharp}\neq\emptyset$,
\begin{equation*}
|V^{\sharp}(Y)|\le\alpha^{\sharp}(k)\,e^{-(1+2\beta_s)\kappa\, d_{\sharp}(Y)} .
\end{equation*}
For the strip $C^{\sharp}\supset C$ of a component $C$ of $Z$ they satisfy
\begin{equation*}
|V^{\sharp}(C^{\sharp})|\le\sum_j\mathsf{K}_{\mathrm{cell}}(k,j)\,|\Gamma^0_j\cap C^{\sharp}|
+\sum_{\mathfrak{a}\ \mathrm{live},\ W_{\mathfrak{a}}\subset C}\mathsf{P}(\mathfrak{a}) .
\end{equation*}
\item[(d)] The large-field operation at complex source obeys the modulus bound
\eqref{5.2:eq-sourced-modulus-a}.
\item[(e)] After the resummation the terms fall into two groups: those whose localisation
domains meet $Z_k^{\sim}\cup\bigcup_iY_i^{\sim}$, where $Z_k=\Lambda_k^c$ is the large-field
region at level $k$, $Z_k^{\sim}$ is its enlargement by one layer of $M\check{R}_k$-cubes, and the
$Y_i$ are the new large-field regions created by $\mathbf{R}_k$ (the first group), and the rest
(the second group). The terms
$\mathbf{R}'^{(k)}(X,(\mathbf{U},\mathbf{J});w)$ of the second group are analytic on the complex
tube of configurations $(\mathbf{U},\mathbf{J})$ over $X$ times $\mathcal{W}_k(X)$ and local in
$w$, and they satisfy, with $\kappa_0$ the exponent of $g_k$ fixed in the correlation theorem,
\begin{equation*}
|\mathbf{R}'^{(k)}|\le e^{-p_0(g_k)}e^{-\kappa d_k(X)}\le \tfrac12\, g_k^{\kappa_0}e^{-\kappa d_k(X)} .
\end{equation*}
The terms of the first group, collected on labels $A$ (with $A:=\operatorname{sk}(X)$ for a
domain $X$ meeting some new strip $Y_i^{\sharp}$, and $A:=X$ for $X$ meeting none), form new
boundary terms $\mathbf{B}'^{(k)}(A,\mathcal{S})$, analytic over $X_A\times\mathcal{W}_k$, where
$X_A=A\cup\bigcup_{\mathcal{P}_A}Y^{\sharp}$ is the true domain and $\mathcal{P}_A$ is the set
consisting of the pairs $(Y_i,k)$ whose strip $Y_i^{\sharp}$ touches $A$ and of the older
large-field regions that still carry their large-field factor after $\mathbf{R}_k$ and whose
strips touch $A$, with
\begin{equation*}
\mathfrak{N}^{A,\mathbb{C}}\bigl(\mathbf{B}'^{(k)}(A,\cdot\,;w)\bigr)
\le e^{\lambda_\theta}c_1^{\flat}(g_k)\,e^{-(1+\frac14\beta_s)\kappa d_k(A)} .
\end{equation*}
At constant $w$ these terms are Euclidean covariant.
\item[(f)] The terms $R''^{(k)}$ form the second expectation of \cite[(3.26)]{B-CMP119}, with
domains $X\subset\Lambda_k$. With $\theta_k:=\theta(g_k)$ the ratio of shift to radius and
$\beta_0$ the constant fixed in the correlation theorem, they are bounded by
\begin{equation}\label{5.2:eq-sourced-large-field-step-5}
K_{CE}\,C_U\,\theta_k\cdot\bigl(k-k_1(k)\bigr)\cdot 3\,(1+\beta_0)^{\kappa_0}g_k^{\kappa_0}
\le\tfrac12\,g_k^{\kappa_0}.
\end{equation}
\end{itemize}
Consequently the family $R^{(k)}_{\zeta}$ of $R$-terms at constant shift $\zeta$ is analytic in
$\zeta\in D(0,2r)$, and in the norm $\|\cdot\|^{\flat}$ of the correlation theorem
\begin{equation}\label{5.2:eq-sourced-large-field-step-a}
\|R^{(k)}_{\zeta}\|^{\flat}\le g_k^{\kappa_0}.
\end{equation}
\end{proposition}

\begin{proof}
Clause (c) is proved below as the boundary-term input at the resummation site
(\ref{5.2:prf-resummation-input}). Clause (e) is proved below as the remnants after resummation
and the new boundary terms (\ref{5.2:prf-resummation-remnants}). Here we prove (a), (b), (d),
(f) and \eqref{5.2:eq-sourced-large-field-step-a}.

\emph{Clause (a): the stage norm.} The stage-norm corollary for the chain terms is applied at
Ba{\l}aban's cores, with the constant $C_0^{\sharp}=3c_G+1$ fixed first. It gives the following.
\begin{itemize}
\item The decomposition \eqref{5.2:eq-sourced-large-field-step-1}.
\item The summands are stored terms with the A-slot datum $A^{\flat}$.
\item The bound \eqref{5.2:eq-sourced-large-field-step-2}, at the rates
$a_k=(1+3\beta_{\mathrm{pr}})\kappa$ and $a_\ell=(1-\beta_{\mathrm{pr}})\kappa$, $\ell<k$.
\item Holomorphy on $\tilde U^{(n+1)c,\Sigma}_k(X)\times\mathcal{W}_k(X)$ times the exterior
weak $A$-coordinates at radii $\varrho_{j'}$.
\item Real-boundedness and continuity on the strong upper boxes $\mathfrak{b}^{\sharp}$.
\end{itemize}
In the stored format the sum and the weight of the norm $\|\cdot\|^{A}$ are taken on labels. A
chain term built on wrapped constituents is in that format with $X$ as its label. For such a term
the statement on $\mathfrak{D}^{\flat}_{n+1}(X)$ is therefore read over $X^{+}$.

\emph{Clause (a): the inherited domain.} The extension to
$\mathfrak{D}^{\flat}_{n+1}(X)\times\mathfrak{D}_{\mathcal{S}}(X)$ is part (a) of the inheritance
theorem for the chain terms. Here $\mathfrak{D}^{\flat}_{n+1}(X)$ is the inherited domain. It is
free of $w$, because it is cut out by three sets of conditions: the constituents' own conditions
on the spaces $\mathfrak{U}_T$ (the space of a terminal term or piece $T$, cut out by its
thresholds), the stage increments $\iota$, and each unit's own analyticity domain $\mathfrak{D}_v$.
On this domain the same terms are stored terms, with the same $C_0^{\sharp}$ and the same rates.
That theorem rests on six inputs: the zeroth-order bound for the A-slot family at the preparatory
site; the third inductive statement of the stage-norm corollary, at $\beta_{\mathrm{pr}}$,
restricted to the region $Z^{\mathrm{on}}$ to which the inheritance theorem restricts that
statement; the layer condition with fraction $\theta_S$, which is supplied for the classes
$\mathbb{E}^{(k)}$ and $\mathbb{R}^{(k)}$ by the enlargement of the third step of the stage
construction and for the class $\mathbb{B}^{(k)}$ by the layer-wise space of part (b) of the
theorem on the creating step; \cite[\S\S3--5]{B-CMP109}, by statement, for the near pairs
of $\mathbb{E}$- and $\mathbb{R}$-terms; the ambient condition; and the cardinality condition of
the A-slot family at width $\mathsf{w}_0$ and fraction $\tfrac12$, a condition of $\gamma_J$. The
statement of the inheritance theorem itself is free of $w$.

\emph{Clause (a): the source variable.} The statement holomorphic in $w$, on
$\mathfrak{D}^{\flat}_{n+1}(X)\times\mathfrak{D}_{\mathcal{S}}(X)\times\mathcal{W}_k(X)$ and with
the same constants, is the corollary of the inheritance theorem at complex source. That corollary
is conditional on (w1) and (w2), which we now establish.

(w1) The condition concerns every terminal term of the first inductive expansion of the A-slot
family, and every piece of the third and fourth steps of the stage construction. Each of them, say
$T$ with domain $X_T$, must depend on $w\in\mathcal{W}_k(X_T)$ holomorphically, jointly with its
other complex variables. This must hold
on the same $w$-free spaces $\mathfrak{U}_T\times\mathfrak{D}_{\mathcal{S}}$, and under the same
majorants: the sizes $\mathfrak{N}$, with the supremum over $w$ taken inside, and $\|T\|_\infty$
for a terminal term $T$. The thresholds $\theta_{T,a}$ that cut out these spaces are free of $w$,
being the thresholds of the source-free scaffold of the localisation site. The required holomorphy
comes from three sources.
\begin{itemize}
\item The clauses of the inductive hypothesis at complex source on the families
$E^{(j)}(\cdot\,;w)$, $D^{(j)}$, $R^{(j)}$ and $\mathbf{B}^{(j)}$ of that hypothesis, which hold
uniformly on $\mathcal{W}_j(X)$ and analytically.
\item The fact that the sizes in the table of unit sizes at the localisation site are free of $w$.
\item The decoupling lemma at complex source.
\end{itemize}

(w2) The following objects are free of $w$:
\begin{itemize}
\item the Gaussian measures of \cite[(1.60)]{B-CMP122b};
\item the characteristic functions $\chi^{(j)\sharp}\chi'^{(j)}$;
\item the response data $\mathbb{H}$, $\mathcal{J}^{\sharp}[1]$ and $\zeta_{\square}$;
\item the decoupling and interpolation parameters $\sigma$, $\mathsf{t}$, $\mathsf{u}$ of the
expansion.
\end{itemize}
They belong to the source-free scaffold of the localisation site and to the fixed data of the
response layers. The first A-slot lemma and its companion hypothesis are statements about these
source-free maps. The second A-slot lemma uses the constituents only under $w$-free majorants.

It is on the domain
$\mathfrak{D}^{\flat}_{n+1}(X)\times\mathfrak{D}_{\mathcal{S}}(X)\times\mathcal{W}_k(X)$ that the
containment at the localisation site is established for these terms. For the presented pair it
follows from the containment corollary and claim of the A-slot family. For the live slots it
follows from the orbit bound of the rotation estimates, applied to the orbit datum. The latter
requires the following further condition.

(w3) The orbit datum $c^{\sharp}$ of $\mathbb{C}^{(n+1)}_k(X,\mathcal{S};w)$ is that of the
family of the corollary at complex source. It is therefore holomorphic in $w$, jointly with the
other variables, under the same $w$-free majorants. The orbit statement of the rotation estimates
is a statement at fixed values of every parameter not named in it, and the parameter lemma of the
rotation estimates carries parameters through unchanged. Hence (w3) follows from (w1) and (w2).

\emph{Clause (a): how the corollary at complex source uses (w1) and (w2).} Its proof lets $w$
enter as one more polydisc parameter in the polydisc lemma for the A-slot family. That lemma
assumes that the pieces of the third and fourth steps are free of the parameter $\lambda$. Its
proof uses this only in the weaker form ``holomorphic in $\lambda$ under $\lambda$-free
majorants''. For $\lambda=w$ the $w$-free majorants are those of the table of unit sizes and of
the decoupling lemma at complex source. The vertex is that of the last two lines of
\cite[(1.60)]{B-CMP122b}, with the bounded weight $-\zeta_k$.

\emph{Clause (a): constants.} Every constant obtained is free of $n$, $N$, $N_0$,
$\check{R}_k$, $k$, $K$, the volume and $g$. This holds at the price of the two conditions of
$\gamma_J$ attached to the A-slot norm at the preparatory site, the second with its step count
evaluated at $24K'$, where $K'=K-m$ is the last step of the local run, together with the full
list of cardinality conditions of the A-slot family, all in $\gamma_J$.
Finally, $0\in\mathfrak{D}_{\mathcal{S}}(X)$, so $\|\cdot\|^{A}$ dominates the supremum on the
real box. Hence \eqref{5.2:eq-sourced-large-field-step-2} implies \cite[(1.70)]{B-CMP122b}. This
bound, together with the analyticity in $w$, is all that clauses (b)--(e),
Lemma~\ref{5.3:lem-wilson-source}, Lemma~\ref{5.2:lem-sourced-modulus} and the inductive step use
of clause (a).

\emph{Clause (b).} This is Lemma~\ref{5.3:lem-wilson-source}, the identity
\eqref{5.3:eq-wilson-source-a}.

\emph{Clause (d).} This is Lemma~\ref{5.2:lem-sourced-modulus}, the bound
\eqref{5.2:eq-sourced-modulus-a}. There Ba{\l}aban's right member \cite[(1.79)]{B-CMP122b} is read
at $A_{\flat}$, with the modified cover of \cite[(1.76)]{B-CMP122b}.

\emph{Clause (f): the bound.} The terms $R''^{(k)}$, namely the second expectation of
\cite[(3.26)]{B-CMP119} with $X\subset\Lambda_k$, form the marked sub-sum of the terms $R^{(j)}$
with $k_1<j<k$. Their bound is the left member of \eqref{5.2:eq-sourced-large-field-step-5}. The
factor $3$ there dominates the weight $w_{k,j}$, which satisfies
$w_{k,j}\le\tfrac52(1+o(1))$ for $k_1<j<k$.

\emph{Clause (f): control of $k-k_1(k)$.} The difference $k-k_1(k)$ is controlled by a function
of $\mathsf{T}_k$ alone. Put $c_u:=\log(1+3\lambda u+2c_2)$, and let
\begin{equation*}
\mathfrak{k}(\mathsf{T}):=\min\bigl\{u\ge1:\ L^{u-1}\ge M(\mathsf{T}+c_u)^{r_{\mathrm{pr}}}\bigr\}.
\end{equation*}
Take $u:=\mathfrak{k}(\mathsf{T}_k)$ and $j:=k-u$. We use three facts.
\begin{itemize}
\item The definition of $R(x)$, with $\check{R}_j=R(\min_{i\le j}x_i)$.
\item The reference recursion, which gives $\log x_j\le\mathsf{T}_k+c_u$.
\item The choice of $u$.
\end{itemize}
Together they give
\begin{equation*}
\begin{split}
L^{-u}M\check{R}_j &< L^{-u}M\cdot L(\log x_j)^{r_{\mathrm{pr}}}\\
&\le L^{1-u}M(\mathsf{T}_k+c_u)^{r_{\mathrm{pr}}}\\
&\le L ,
\end{split}
\end{equation*}
and so $k_1(k)\ge k-\mathfrak{k}(\mathsf{T}_k)$. The function $\mathfrak{k}$ is nondecreasing,
and, for $\mathsf{T}$ above a threshold $c_{\mathfrak{k}}$ depending only on $L$, $M$, $\lambda$,
$c_2$ and $r_{\mathrm{pr}}$,
\begin{equation*}
\mathfrak{k}(\mathsf{T})\le 2+\log_LM+r_{\mathrm{pr}}\log_L(2\mathsf{T}) .
\end{equation*}
Hence, with $C_\theta$ the constant in $\theta_k\le C_\theta\,\mathsf{T}_k^{\,p_0-q}$,
\begin{equation*}
\theta_k\cdot\bigl(k-k_1(k)\bigr)\le C_\theta\,\mathsf{T}_k^{\,p_0-q}\,\mathfrak{k}(\mathsf{T}_k).
\end{equation*}
The right member, inserted into the left member of
\eqref{5.2:eq-sourced-large-field-step-5}, is dominated by the condition of $\gamma_J$ for the
second expectation. That condition is a supremum over $\mathsf{T}\ge\mathsf{T}_\gamma$ and uses no
quantity of the history of the run. This gives the bound $\tfrac12 g_k^{\kappa_0}$ in
\eqref{5.2:eq-sourced-large-field-step-5}.

\emph{The radii.} The step at level $k$ uses Ba{\l}aban's enlarged spaces after the operation
$\mathbf{T}$ at level $k$, with radii $(1+\beta_{\mathrm{pr}})\alpha_{\cdot,k}$. We have
$\beta_{\mathrm{pr}}=1/5\le\beta_s=1/4$. We also have the enlargement condition of the third step
of the stage construction, and part (ii) of the interpolation lemma for the enlargement fractions.
These give
\begin{equation*}
(1+\beta_{\mathrm{pr}})\alpha_{\cdot,k}\le(1+\beta_s)\alpha_{\cdot,k}\le\hat\alpha_{k-1}.
\end{equation*}

\emph{The bound \eqref{5.2:eq-sourced-large-field-step-a}.} By clause (e), the second-group terms
$\mathbf{R}'^{(k)}$ are bounded by $\tfrac12 g_k^{\kappa_0}e^{-\kappa d_k(X)}$. By clause (f), the
terms $R''^{(k)}$ are bounded by $\tfrac12 g_k^{\kappa_0}$. Hence
$\|R^{(k)}_{\zeta}\|^{\flat}\le g_k^{\kappa_0}$, and $R^{(k)}_{\zeta}$ is analytic in
$\zeta\in D(0,2r)$, which is \eqref{5.2:eq-sourced-large-field-step-a}.
\end{proof}

\begin{proposition}[The boundary-term input at the resummation site]
\label{5.2:prf-resummation-input}
Let $k\le K$. Assume the inductive hypothesis of the step at complex sources at level $k$, after
the large-field operation $\mathbf{T}$, with $\sup_j|\zeta_j|\le\frac83 r$, $r$ being the radius
fixed in that hypothesis. Let $Z$ be the class of
components of $\Lambda_k^c$ treated by the operation $\mathbf{R}_k$, and let $Z^\sharp$,
$C^\sharp\supset C$ be the strips of width $m_\sharp(k)$ about $Z$ and about its components $C$.
Put
\begin{equation}\label{5.2:eq-resummation-input-f}
\begin{gathered}
\beta_\sharp:=\tfrac1{20},\qquad b:=(1+2\beta_s)\kappa=\tfrac32\kappa,\qquad
a:=b+\beta_\sharp\kappa=\tfrac{31}{20}\kappa,\\
\kappa\ \ge\ \kappa_d:=\max\Bigl\{20\bigl(c_d+2+2\sqrt d\,(2+4\beta_{\mathrm{pr}})\bigr),\
40(\mathsf{k}_d+1),\ \tfrac19\bigl(c_F+2\cdot3^d+8\bigr)\Bigr\},
\end{gathered}
\end{equation}
where $c_d:=2^d\log(2de)+1$, $\mathsf{k}_d:=1+\mathsf{c}_d(2d\log3+(1+3^d)\log2)$,
$\mathsf{c}_d:=2\cdot3^d$, and $c_F\le2\mathsf{b}_{\mathtt{f}}$ is an absolute constant; that is,
$\kappa/20\ge\max\{c_d+2+2\sqrt d(2+4\beta_{\mathrm{pr}}),\,2(\mathsf{k}_d+1)\}$ and
$9\kappa\ge c_F+2\cdot3^d+8$. The units at the localisation site are sorted into the classes
$1$--$5$ and $\partial$ of the size table of that site. Put
\begin{equation*}
\varepsilon^{\theta'}(k):=\bigl[C_\sharp O(1)+8C_A'C_NC_0^\sharp(\mathsf{T}_k+1)^{r_{\mathrm{pr}}}
+8C_A'\mathsf{K}^{\mathrm{arr}}\bigr]\theta'(g_k),
\end{equation*}
the sum of the three $\theta'$-members, and let $\varepsilon^{\theta'}_{a+7}(k)$ be the same three
members read at weight $a+7$, with
$\mathsf{C}_1^{\mathrm{ex}}:=\mathsf{C}^{\Sigma,4}+\mathsf{C}^{\Sigma}_{a+7}$ for the first member
and $\mathsf{K}^{\mathrm{arr}\prime}$ in place of $\mathsf{K}^{\mathrm{arr}}$, as in clause (e) of
the exterior-class lemma. Assume moreover the following.
\begin{itemize}
\item[(P1)] The stage-norm bound for the chain terms: for every stage $n\le N(g_k)$ of this step's
chain, $\|\mathbb{C}^{(n)}_k\|^A_{(n)}\le C_0^\sharp$, the weight being taken at the rate
$a_k:=(1+3\beta_{\mathrm{pr}})\kappa=\frac85\kappa$ for terms filed at level $k$ and at
$a_m:=(1-\beta_{\mathrm{pr}})\kappa=\frac45\kappa$ in $d_m$ for terms filed at level $m<k$, each
$\mathbb{C}^{(n)}_k(X,\mathcal{S};w)$ being holomorphic on its inherited domain times
$\mathcal{W}_k(X)$ with the same constant and rates; and the same bound and holomorphy for the
frozen chain terms of every live arrest at its own step.
\item[(P2)] The summation of boundary terms at finer zones, with its corollary: the contribution
of the class-$4$ units and of the long singles and tails of (c) below to the three
$\theta'$-members of \eqref{5.2:eq-resummation-input-b} below
is bounded by $\theta'(g_k)\mathsf{C}^{\Sigma,4}$, the pieces being summed in bundles, in every
zone, rooted at the source, without the gain $\mathsf{w}_L^{-1}$.
\item[(P3)] The statements for wrapped units (units stored in the relative format over live
strips): their piece decomposition, in which the value $\mathrm{val}_v$ is the piece with all
families empty and a second-part wrapped unit has $\mathsf{y}_1\neq\emptyset$; the relative bound,
which bounds the pieces of all wrapped units, interior ones included, by the member
$\varepsilon^{\mathrm{rel}}(k)$ below, with $\varepsilon^{\mathrm{rel}}(k)\le1$ under its two
ceilings on $\gamma$; the volume bound, by which the pieces and the value of one wrapped unit
$v$ sum to at most $4\mathfrak{N}^{A}(v)e^{4(d_k(K_v)+1)}$, $\mathfrak{N}^{A}(v)$ being the size of
$v$ in the relative format; the Cauchy bound in the decoupling parameters at the source; the
field-free geometric, cell and summation lemmas; the summation lemma for deep units; and the length
estimate of \cite[(1.67)]{B-CMP122b} for wrapped units.
\item[(P4)] The exterior-piece sum for plain units (Lemma~\ref{5.2:lem-plain-exterior-sum},
clauses (c), (d)), the volume sum of one unit's pieces
(Corollary~\ref{5.2:cor-unit-volume-sum}), and the exterior-class lemma, clauses (c), (e), for the
class $\mathfrak{L}^{\mathrm{ex}}=\mathfrak{L}^c\sqcup\mathfrak{L}^{Y_0}$.
\item[(P5)] The three-factor bound for deep units: for a deep class-$4$ unit $v$,
$X_v\subset\mathsf{C}_C$, and every piece $\mathfrak{p}=(v;\mathsf{y}_1,\dots,\mathsf{y}_4)$,
pieces with $\mathsf{y}_1$ alone non-empty included,
\begin{equation*}
|\mathfrak{p}|\le C_{3F}\,e^{-\frac12\delta_P d^{\wedge}(X_v,\mathsf{C}_C^c)}\,
\mathfrak{N}(v)\,e^{-(\kappa_1-1)\sum_t|\mathsf{y}_t|},
\end{equation*}
with $d^{\wedge}:=d^{\wedge}_{\mathbf{B}''_k}$ the multi-scale contour distance, which counts at
least the weight of each stratum it crosses.
\item[(P6)] The strip theorem, clauses (a)--(c), and its corollary on strips, for an arbitrary
family indexed by $\mathbf{D}_k^Z$.
\item[(P7)] The bounds \eqref{5.2:eq-wilson-remainders-a} of
Lemma~\ref{5.2:lem-wilson-remainders} for the family $\mathfrak{v}^{\mathbf{C}}$.
\item[(P8)] $\gamma\le\gamma_J$; in particular the bracket ceiling
$(1+C_\Sigma C_e)^3e\,S_d\,\varepsilon^{\theta'}_{a+7}(k)\le\frac12$, the ceiling of
Lemma~\ref{5.2:lem-plain-exterior-sum}(d), and the allowance
\eqref{5.2:eq-per-cell-allowance-a}, each as a supremum over $\mathsf{T}\ge\mathsf{T}_\gamma$.
\end{itemize}
For $Y_4\in\mathbf{D}_k^Z$ let $\mathfrak{v}(Y_4)$ be the sum of the pieces and values filed at
$Y_4$ at the localisation site, each domain having been enlarged by the components of $Z$ which it
meets, and split
$\mathfrak{v}=\mathfrak{v}^{\mathrm{br}}+\mathfrak{v}^{\mathrm{val}}+\mathfrak{v}^{\mathbf{C}}$
into the pieces of the brackets, the values, and the filed members of $\mathbf{C}_k$. Put
\begin{equation*}
\|\mathfrak{v}\|_a:=\sup_{Y_4\setminus Z\neq\emptyset}e^{a\,d_{k,Z}(Y_4)}|\mathfrak{v}(Y_4)|.
\end{equation*}
Here and below $\mathsf{K}_{\mathrm{cell}}(k,n)$ denotes $c_V^S$ times the per-cell constant of
\eqref{5.2:eq-cell-constant-a}, and $\mathsf{P}(\mathfrak{a})$ denotes $c_V^S$ times the frozen
contribution of \eqref{5.2:eq-arrested-regions-a}, where $c_V^S:=4e^43^d$; $n(\Delta)$ is the zone
of a cell $\Delta$. Then:
\begin{enumerate}
\item[(1)] Each $\mathfrak{v}(Y_4)$ is analytic on the tube over $Y_4$ (over $Y_4^{+}$ for the parts
coming from wrapped units) times $\mathcal{W}_k(Y_4)$, and
\begin{equation}\label{5.2:eq-resummation-input-b}
\begin{gathered}
\|\mathfrak{v}^{\mathrm{br}}\|_a\le\varepsilon_{\mathrm{br}}(k),\qquad
\varepsilon_{\mathrm{br}}(k):=\varepsilon^{\theta'}(k)+\varepsilon^{\mathrm{rel}}(k)
+\varepsilon^{\mathrm{int}}(k)+\varepsilon^{\mathrm{ext}}(k),\\
\varepsilon^{\mathrm{ext}}(k):=\varepsilon^{\mathrm{ext},\sigma}(k)
+\varepsilon^{\mathrm{ext},\mathsf{t}}(k),\qquad \varepsilon_{\mathrm{br}}(k)\le2,
\end{gathered}
\end{equation}
with the six members $\varepsilon^{\theta'}$ (three members), $\varepsilon^{\mathrm{rel}}$,
$\varepsilon^{\mathrm{int}}$, $\varepsilon^{\mathrm{ext}}$, where
\begin{align*}
\varepsilon^{\mathrm{rel}}(k)&:=e^{-(\check R_k-2)}(1+C_\Sigma C_e)^3
\bigl[(1+4h_\sharp^dC_\Sigma)\mathsf{Q}^{\star}_{a+9}+4C_\Sigma C_{\mathsf{a}}(1)\mathsf{b}_k(0)\bigr],\\
\varepsilon^{\mathrm{int}}(k)&:=e^{-(\check R_k-2)}(1+C_\Sigma C_e)^4\,\mathsf{Q}^{\star,\mathrm{pl}}_{a+9},\\
\varepsilon^{\mathrm{ext},\sigma}(k)&:=e^{-(\check R_k-2)}(1+C_\Sigma C_e)^3e\,S_d
\bigl[\mathsf{K}^{0\prime}+\mathsf{K}^{\flat,\mathrm{ex}}+N(g_k)C_0^\sharp\bigr],\\
\varepsilon^{\mathrm{ext},\mathsf{t}}(k)&:=(1+C_\Sigma C_e)^3e\,S_d\,\varepsilon^{\theta'}_{a+7}(k)
-\varepsilon^{\theta'}(k),
\end{align*}
with the constants $C_\Sigma$, $C_e$, $h_\sharp$, $\mathsf{Q}^{\star}_{a+9}$,
$C_{\mathsf{a}}(1)$, $\mathsf{b}_k(0)$ of the statements for wrapped units (P3), the constant
$\mathsf{Q}^{\star,\mathrm{pl}}_{a+9}$ of Lemma~\ref{5.2:lem-plain-exterior-sum}, and
$\mathsf{K}^{\flat,\mathrm{ex}}:=e^{(a+9)c_lc_{\mathfrak{h}}}\,27S_d'\mathsf{K}^{\flat}$; every unit
and every tuple at the localisation site is held by exactly one of these members or by
$\mathfrak{v}^{\mathrm{val}}$, $\mathfrak{v}^{\mathbf{C}}$ or the volume bound, according to the
following allocation:
\begin{itemize}
\item[(a)] Wrapped units: every non-empty tuple goes to $\varepsilon^{\mathrm{rel}}$; the all-empty
tuple of a first-part wrapped unit goes to $\mathfrak{v}^{\mathrm{val}}$. A second-part wrapped
unit has no all-empty tuple in its first link, since $\mathsf{y}_1\neq\emptyset$ for such units by
the piece decomposition (P3); its remainder is part of the new action.
\item[(b)] Plain interior units ($X_v\subset Z$): the all-empty tuple goes to the values of
\eqref{5.2:eq-resummation-input-e} (Corollary~\ref{5.2:cor-unit-volume-sum}); a non-empty tuple
whose filing domain leaves $Z$ goes to $\varepsilon^{\mathrm{int}}$; a non-empty tuple whose filing
domain stays in the component $C$ goes to \eqref{5.2:eq-resummation-input-e}, and is under no
$Y_4$ of the supremum $\|\cdot\|_a$.
\item[(c)] Plain first-part units with $K_v\setminus Z\neq\emptyset$: for the first member
$f_v(\mathfrak{b}_3)$, the all-empty tuple goes to $\mathfrak{v}^{\mathrm{val}}$ and a non-empty
tuple to $\varepsilon^{\mathrm{ext},\sigma}$; for its $Y_0$-term $v_p$, the tuple with empty
links $2$--$4$ goes to the three $\theta'$-members and a non-empty continuation in links $2$--$4$
to $\varepsilon^{\mathrm{ext},\mathsf{t}}$; for class-$4$ units the pieces $v_p$ are summed in
bundles.
\item[(d)] Plain second-part units: their four $Y_0$-terms are treated as the $v_p$ of (c).
\item[(e)] The members of $\mathbf{C}_k$ free of the sources: every tuple goes to
$\mathfrak{v}^{\mathbf{C}}$, in no member of $\varepsilon_{\mathrm{br}}$.
\end{itemize}
Further,
\begin{equation}\label{5.2:eq-resummation-input-d}
\|\mathfrak{v}^{\mathrm{val}}\|_a\le\mathsf{K}_{\mathrm{val}}(k);
\end{equation}
$\mathfrak{v}^{\mathbf{C}}$ obeys \eqref{5.2:eq-wilson-remainders-a}; and, for every component $C$
of $Z$ and every strip cell $\Delta\subset C^\sharp\setminus C$,
\begin{equation}\label{5.2:eq-resummation-input-e}
\begin{aligned}
\sum_{Y_4\subset C}|\mathfrak{v}(Y_4)|&\le\sum_n\mathsf{K}_{\mathrm{cell}}(k,n)|\Gamma^0_n\cap C|
+\sum_{\mathfrak{a}\ \mathrm{live},\,W_{\mathfrak{a}}\subset C}\mathsf{P}(\mathfrak{a}),\\
\sum_{Y_4\cap\Delta\neq\emptyset,\ Y_4\subset C^\sharp}|\mathfrak{v}(Y_4)|
&\le\mathsf{K}_{\mathrm{cell}}(k,n(\Delta))+\varepsilon^{\mathrm{ext}}(k).
\end{aligned}
\end{equation}
\item[(2)] On the tube times $\mathcal{W}_k(Y)$, the function $V(Y,(\mathbf{U},\mathbf{J});w)$ of
\cite[(1.68)--(1.69)]{B-CMP122b} obeys, with Ba{\l}aban's constant $\alpha$ there,
\begin{equation}\label{5.2:eq-resummation-input-a}
\begin{aligned}
|V(Y,(\mathbf{U},\mathbf{J});w)|&\le C_\sharp\alpha\,e^{-(1+2\beta_s)\kappa d_{k,Z}(Y)}
\qquad\text{if }Y\setminus Z^{\sim}\neq\emptyset,\\
|V(Y,(\mathbf{U},\mathbf{J});w)|&\le\sum_j\mathsf{K}_{\mathrm{cell}}(k,j)|\Gamma^0_j\cap Y|
+\sum_{\mathfrak{a}\ \mathrm{live},\,W_{\mathfrak{a}}\subset Y}\mathsf{P}(\mathfrak{a})\\
&\qquad\text{if $Y$ is a component of }Z^{\sim},
\end{aligned}
\end{equation}
which, with $Z^{\sim}$ read as $Z^\sharp$ and $d_{k,Z}$ as
$d_\sharp:=d_{k,Z^\sharp}$, reads
\begin{equation}\label{5.2:eq-resummation-input-1}
\begin{aligned}
|V^\sharp(Y)|&\le\alpha^\sharp e^{-(1+2\beta_s)\kappa d_\sharp(Y)}
&&\text{if }Y\setminus Z^\sharp\neq\emptyset,\\
|V^\sharp(C^\sharp)|&\le\sum_j\mathsf{K}_{\mathrm{cell}}(k,j)|\Gamma^0_j\cap C^\sharp|
+\sum_{\mathfrak{a}\ \mathrm{live},\,W_{\mathfrak{a}}\subset C}\mathsf{P}(\mathfrak{a})
&&\text{for a strip }C^\sharp\supset C.
\end{aligned}
\end{equation}
\end{enumerate}
\end{proposition}

\begin{proof}
\emph{The pre-filing family.} At the localisation site every term of the first part of the
effective action other than its first two, the new terms constructed at this step included, is
localised by the four
operations of \cite[(1.59)--(1.68)]{B-CMP122b}; each bracket \cite[(1.59), (1.60)]{B-CMP122b} is a
difference $f(\mathfrak{b}_1)-f(\mathfrak{b}_2)$ of one function at two admissible backgrounds.
The expansion in the decoupling parameters of the four links writes a unit $v$ as its value
$\mathrm{val}_v$, the piece with all four families empty, plus its pieces
$\mathfrak{p}=(v;\mathsf{y}_1,\dots,\mathsf{y}_4)$ with $(\mathsf{y}_t)$ not all empty. The family
$\sum_{Y_4\in\mathbf{D}_k^Z}\mathfrak{v}(Y_4)$ is Ba{\l}aban's state before the last resummation
\cite{B-CMP122b}, after the pre-filing $Y_4\mapsto Y_4\cup\{\text{components of $Z$ meeting
$Y_4$}\}$. This enlargement adds nothing off $Z$; hence it leaves $Y_4\setminus Z$ unchanged and
does not raise $d_{k,Z}(Y_4)$, although Ba{\l}aban's domains $Y'_i$ may meet $Z$ partially. The
number of domains regrouped into one enlarged domain is paid by the $\kappa_1$-volume factor of the
expansion, as in the strip theorem. The analyticity of each $\mathfrak{v}(Y_4)$ on the tube over
$Y_4$ times $\mathcal{W}_k(Y_4)$ is that of its summands; for the parts coming from wrapped units
the tube is read over $Y_4^{+}$, as the statements for wrapped units (P3) provide.

\emph{The allocation.} The allocation (a)--(e) of the statement is exhaustive: nothing else occurs
at the localisation site \cite{B-CMP122b}; the remainder of the second part
outside the brackets is, in Ba{\l}aban's words, ``the desired terms we need to reconstruct the new
action'' \cite{B-CMP122b}, and is no member of the family. In particular the three
$\theta'$-members are read only for the plain units of (c) and (d), those with
$K_v\setminus Z\neq\emptyset$ or $X\cap Z=\emptyset$, and there at the tuple with empty links
$2$--$4$, so that no unit is counted in two members.

\emph{The three $\theta'$-members.} They come from the Cauchy bound at the localisation site,
applied with the size table of that site, and $\mathsf{w}_L^{-1}\le\theta'(g_k)
=C_{\theta'}\mathsf{T}_k^{\,p_1-q}$, with Ba{\l}aban's exponents $p_1$ and $q$. The first,
$C_\sharp O(1)\theta'(g_k)$, holds the brackets of
the classes $1$--$3$, $5$, $\partial$ and the long singles and tails of (c). The second holds the
brackets of all stages $n\le N(g_k)\le C_N(\mathsf{T}_k+1)^{r_{\mathrm{pr}}}$ of this step's chain:
at a cell the weighted sum over $X$ of one stage is at most its norm, which is at most $C_0^\sharp$
by (P1); it is multiplied by the factor $2$ of class $4$ and by the factor
$4C_A'\mathsf{w}_L^{-1}\eta_p^{1/2}$ of the half-power form of the Cauchy bound, and
$\eta_p\le1$, $\mathsf{w}_L^{-1}\le\theta'(g_k)$ give
\begin{equation*}
N(g_k)\,C_0^\sharp\cdot2\cdot4C_A'\theta'(g_k)
\le8C_A'C_NC_0^\sharp(\mathsf{T}_k+1)^{r_{\mathrm{pr}}}\theta'(g_k).
\end{equation*}
This bound is crude: it ignores the finer count
$N_0+3$ of the stage-count lemma. The third holds the brackets of the exterior arrested left-overs
($K_v\setminus Z\neq\emptyset$; the frozen chain terms and the change-of-variables terms of the
frozen terms $\mathfrak{F}(\mathfrak{a})$): their sum is the collar series $\mathsf{K}^{\mathrm{arr}}$
of \eqref{5.2:eq-exterior-collar-bound-a}, by Lemma~\ref{5.2:lem-exterior-collar-bound}, a series
free of $g$, of $k$ and of the number of steps elapsed since the arrest, multiplied by
the same $2\cdot4C_A'\mathsf{w}_L^{-1}$. The class-$4$ part of the three members, that is, the
pieces $v_p$ of the class-$4$ units of (c), (d) and the long singles and tails of (c) inside
$C_\sharp O(1)$, is bounded by $\theta'(g_k)\mathsf{C}^{\Sigma,4}$ by (P2); for it the three
displayed constants are majorants only.

\emph{The wrapped member.} By (P3), $\varepsilon^{\mathrm{rel}}(k)$ bounds the pieces of all
wrapped units, of the classes $\mathfrak{L}^{\mathrm{II}}$, $\mathfrak{L}^{\circ}$,
$\mathfrak{L}^{\mathrm{sh}}$, $\mathfrak{L}^{\mathrm{dp}}$ of those statements, interior ones
included, the tube of the
terms $\mathfrak{v}(Y_4)$ being read on $Y_4^{+}$; for these units neither the Cauchy bound at the
localisation site, nor its telescoping in $\hat\chi_{\square}$, nor the weight supremum
$\mathsf{W}$ is used. Under the two ceilings of (P3), $\varepsilon^{\mathrm{rel}}(k)\le1$.

\emph{The interior member.} Let $\mathfrak{v}^{\mathrm{br,int}}(Y_4)$, for
$Y_4\setminus Z\neq\emptyset$, be the part of $\mathfrak{v}^{\mathrm{br}}(Y_4)$ made of the pieces
$\mathfrak{p}=(v;\mathsf{y}_1,\dots,\mathsf{y}_4)$, $(\mathsf{y}_t)$ not all empty, of interior
first-part units $v$, that is, $X_v\subset Z$ and hence $X_v$ inside a component $C$ of $Z$: this
chain's $\mathbb{C}^{(n)}_k$, the frozen terms $\mathfrak{F}(\mathfrak{a})$ of the live arrests
nested in $C$, the stored $\mathbf{B}^{(j)}$, $\mathbf{B}'^{(j)}$, and the pieces of the classes
$1$--$3$, $5$, $\partial$ anchored in $C$. Such pieces exist and reach $Z^c$ through the domains
$Y'_i$, which ``contain at least one component of $\Lambda$, and intersect the boundary layer
$Y_1\setminus Y_1^0$''; in Ba{\l}aban's words \cite{B-CMP122b}:
\begin{quotation}
$Y_4\setminus Z=(Y_0\setminus Z)\cup\bigcup_{i=2}^4(Y'_i\setminus Z)$ \dots\ The domains $Y'_i$
connect $Y_1\setminus Z$ with $\Lambda$, and the third term in the exponential yields a connection
of $X$ with $Y'_{i_0}$.
\end{quotation}
Split $\mathfrak{v}^{\mathrm{br,int}}=\mathfrak{v}^{\mathrm{int,pl}}+(\text{the
part of the wrapped interior units})$; the second summand lies inside the wrapped member, in the
classes $\mathfrak{L}^{\mathrm{sh}}\cup\mathfrak{L}^{\mathrm{dp}}$. For the first,
Lemma~\ref{5.2:lem-plain-exterior-sum}, clauses (c) and (d), gives, over all plain interior
first-part units (classes $1$--$3$, $5$, $\partial$; plain $\mathbf{B}$, $\mathbf{B}'$; this chain's
chain terms; frozen chain terms and change-of-variables terms), with the pieces fixed by the
allocation,
\begin{equation}\label{5.2:eq-resummation-input-c}
\|\mathfrak{v}^{\mathrm{int,pl}}\|_a\le\|\mathfrak{v}^{\mathrm{int,pl}}\|_{a,1}
\le\varepsilon^{\mathrm{int}}(k)
=e^{-(\check R_k-2)}(1+C_\Sigma C_e)^4\,\mathsf{Q}^{\star,\mathrm{pl}}_{a+9}\le\tfrac12 .
\end{equation}
Each piece enters that lemma through one of two inputs. For a deep unit it is the three-factor
bound: for a deep class-$4$ unit, $X_v\subset\mathsf{C}_C$, the bound (P5), for all four links,
which replaces Ba{\l}aban's third
factor $\exp(-\delta d((Y_1^0)^c\cap Y_1\cap Y'_{i_0},X))$ \cite{B-CMP122b}; for the deep pointed
units of the classes $1$--$3$, $5$ its pointed analogue established in the proof of
Lemma~\ref{5.2:lem-plain-exterior-sum}, at rate $\frac12\delta_P$ and with the sizes
$\mathfrak{N}^{\flat}$, the size $\mathfrak{N}(v)$ in (P5) being the supremum on the unit's own
domain and lacking the factor $t^4$ of the size table there. For a shallow unit it is the Cauchy
estimate in the decoupling parameters on the product polydisc, through the source Cauchy bound of
(P3). The
counting uses the field-free geometric, cell and summation lemmas of (P3), with the pile lemma of
this chapter in place of the anchor lemma of the statements for wrapped units; the weight
$e^{a\,d_{k,Z}(Y_4)}$ is paid by the decay
in $\mathsf{y}$ under the single condition $\kappa_1-1\ge c_l(a+9)+\lambda_\ast$, $c_l=1+234d$. The
member $\varepsilon^{\mathrm{int}}(k)$ carries no factor $\theta'(g_k)$: its smallness is
$e^{-(\check R_k-2)}$, not $\mathsf{w}_L^{-1}$, and its ceiling on $\gamma$ is that of
Lemma~\ref{5.2:lem-plain-exterior-sum}(d).

\emph{The exterior member.} By clauses (c) and (e) of the exterior-class lemma (P4), the pieces of
the class $\mathfrak{L}^{\mathrm{ex}}$ are bounded as follows:
$\varepsilon^{\mathrm{ext},\sigma}(k)$ bounds the pieces of the units of (c) free of $\mathsf{t}$
with a non-empty family in links $2$--$4$ (the class $\mathfrak{L}^c$), and
$\varepsilon^{\mathrm{ext},\mathsf{t}}(k)$ bounds the non-empty continuations in links $2$--$4$ of
the $\mathsf{t}$-pieces of (c) and (d) (the class $\mathfrak{L}^{Y_0}$).

\emph{The bound $\varepsilon_{\mathrm{br}}(k)\le2$.} By the definition of
$\varepsilon^{\mathrm{ext},\mathsf{t}}$,
\begin{equation*}
\varepsilon^{\theta'}(k)+\varepsilon^{\mathrm{ext},\mathsf{t}}(k)
=(1+C_\Sigma C_e)^3e\,S_d\,\varepsilon^{\theta'}_{a+7}(k)\le\tfrac12
\end{equation*}
by the bracket ceiling of (P8); $\varepsilon^{\mathrm{rel}}(k)\le1$ by (P3); and
$\varepsilon^{\mathrm{int}}(k)+\varepsilon^{\mathrm{ext},\sigma}(k)\le\frac12$ by the ceiling of
Lemma~\ref{5.2:lem-plain-exterior-sum}(d). Grouping the six members in these three sums gives
\begin{equation*}
\varepsilon_{\mathrm{br}}(k)
=\bigl(\varepsilon^{\theta'}+\varepsilon^{\mathrm{ext},\mathsf{t}}\bigr)(k)
+\varepsilon^{\mathrm{rel}}(k)
+\bigl(\varepsilon^{\mathrm{int}}+\varepsilon^{\mathrm{ext},\sigma}\bigr)(k)
\le\tfrac12+1+\tfrac12=2 .
\end{equation*}
Here $\varepsilon_{\mathrm{br}}(k)$ is a function of $g_k$, the ceilings are suprema over
$\mathsf{T}\ge\mathsf{T}_\gamma$, and $\mathsf{T}_k\ge\mathsf{T}\ge\mathsf{T}_\gamma$ for $k\le K$,
since $x_k>g^{-2}$ for $k\le K$ by the definition of the first passage $K$, and $g\le\gamma$;
so the bound holds at every $k\le K$. Since $\varepsilon^{\theta'}(k)\ge0$, also
$\varepsilon^{\mathrm{ext},\mathsf{t}}(k)\le\frac12$, and
$\varepsilon^{\mathrm{ext},\sigma}(k)\le\frac12$; hence $\varepsilon^{\mathrm{ext}}(k)\le1$. Each
piece of $\mathfrak{v}^{\mathrm{br}}(Y_4)$, $Y_4\setminus Z\neq\emptyset$, is held by exactly one of
the six members, by (a)--(d); the triangle inequality gives the first inequality of
\eqref{5.2:eq-resummation-input-b}.

\emph{The values.} By Lemma~\ref{5.2:lem-value-unit-filing},
$\|\mathfrak{v}^{\mathrm{val}}\|_a\le\mathsf{K}_{\mathrm{val}}(k)$, which is
\eqref{5.2:eq-resummation-input-d}, and by Lemma~\ref{5.2:lem-exterior-collar-bound},
$\mathsf{K}_{\mathrm{val}}(k)\le\mathsf{K}^0+N(g_k)C_0^\sharp$.
This is Ba{\l}aban's treatment of ``the exceptional terms, for which all the domains $Y'$ are
empty'', which are ``bounded by a constant times the same exponential factor'' \cite{B-CMP122b}.

\emph{The members of $\mathbf{C}_k$.} They obey \eqref{5.2:eq-wilson-remainders-a} by
Lemma~\ref{5.2:lem-wilson-remainders}. If $Y_4\setminus Z\neq\emptyset$, choose a cube
$\Delta\subset Y_4\setminus Z$; then $e^{a\,d_{k,Z}(Y_4)}|\mathfrak{v}^{\mathbf{C}}(Y_4)|$ is one
summand of the sum over $Y_4'\ni\Delta$ in $\|\mathfrak{v}^{\mathbf{C}}\|_{a,1}$, so
$\|\mathfrak{v}^{\mathbf{C}}\|_a\le\|\mathfrak{v}^{\mathbf{C}}\|_{a,1}\le\mathsf{c}_{\mathbf{C}}$.
With the two preceding paragraphs, the triangle inequality gives
$\|\mathfrak{v}\|_a\le\varepsilon_{\mathrm{br}}(k)+\mathsf{K}_{\mathrm{val}}(k)+\mathsf{c}_{\mathbf{C}}$.

\emph{The volume bound.} The first line of \eqref{5.2:eq-resummation-input-e} is obtained as
follows. The last column of the size table, per cell of zone $n$, is summed over the classes and
multiplied by $C_\sharp$ for the classes $1$--$3$, $5$, $\partial$. The class-$4$ values of this
step's chain are charged once per term at a cell of the term's own stage region, by the stage-count
lemma; this gives the per-cell constant of \eqref{5.2:eq-cell-constant-a}, at most
$C_\sharp[\mathsf{K}^0_{\mathrm{cell}}+3(\mathsf{N}_0(\mathsf{T}_n)+3)C_0^\sharp]$, so that, with
the factor $c_V^S$ of the statement, justified below,
\begin{equation*}
\mathsf{K}_{\mathrm{cell}}(k,n)\le c_V^S\,C_\sharp\bigl[\mathsf{K}^0_{\mathrm{cell}}
+3(\mathsf{N}_0(\mathsf{T}_n)+3)C_0^\sharp\bigr].
\end{equation*}
The frozen terms of the live arrested pieces, their
values and those pieces whose filing domain stays inside $C$, are counted extensively as
$\mathsf{P}(\mathfrak{a})$, \eqref{5.2:eq-arrested-regions-a}: each is charged once, at an anchor
of its own stage's scale. The pieces whose filing domain leaves $Z$ belong to
\eqref{5.2:eq-resummation-input-c}. The strip cells $\Delta\not\subset Z$ hold no live piece, and a
frozen term reaching such a $\Delta$ from $Z$ is exterior and is counted in the collar series of
Lemma~\ref{5.2:lem-exterior-collar-bound}.

The pieces of an interior unit carry no factor $\mathsf{w}_L^{-1}$; their volume sum is paid as
follows.
All pieces of one plain interior unit $v$ sum to at most $\mathfrak{N}^{\flat}(v)e^{4d_k(K_v)+4}$,
by Corollary~\ref{5.2:cor-unit-volume-sum}, with $\mathfrak{N}^{\flat}$ the size of the pointed
classes and $\mathfrak{N}^{\flat}=\mathfrak{N}$ on class $4$; the pieces and the value of a wrapped
unit sum to at most $4\mathfrak{N}^{A}(v)e^{4(d_k(K_v)+1)}$, by the volume bound of (P3). The factor
$e^{4d_k(K_v)}\le e^{4d_j(A_v)}$ is absorbed by four units of the unit's own decay rate, which are
available in every class, and charging on labels costs the factor $3^d$. Hence
\eqref{5.2:eq-resummation-input-e}, \eqref{5.2:eq-resummation-input-a} and the volume count of
\eqref{5.2:eq-cell-constant-a} hold with the per-cell constant and the frozen contribution
multiplied by the pure number $c_V^S:=4e^4 3^d$, that is, with $\mathsf{K}_{\mathrm{cell}}(k,n)$ and
$\mathsf{P}(\mathfrak{a})$ as in the statement; this factor is part of
the allowance \eqref{5.2:eq-per-cell-allowance-a}, and $c_{\mathsf{P}}$ is read as
$c_V^Sc_{\mathsf{P}}$. No depth factor and no three-factor bound is needed for interior sums.

For the second line of \eqref{5.2:eq-resummation-input-e}, the contribution at a strip cell
$\Delta\subset C^\sharp\setminus C$ of the last column of the size table is
$\mathsf{K}_{\mathrm{cell}}(k,n(\Delta))$. The pieces of $\mathfrak{L}^{\mathrm{ex}}$ filed at a
$Y_4\subset C^\sharp$ meeting $\Delta\not\subset Z$ lie under the supremum of
$\|\cdot\|_{a,1}$ at the $\pi_k$-cube of $\Delta$ with weight at least $1$, so their sum there is at
most $\varepsilon^{\mathrm{ext}}(k)\le1\le C_\sharp\cdot1$. This is the summand $+1$ added to
$\mathsf{K}^0_{\mathrm{cell}}$ in \eqref{5.2:eq-per-cell-allowance-a}; no further ceiling is
needed.

All these bounds are supplied at Ba{\l}aban's cores by the filing described below: $K_v=\bar X$;
for the classes $1$--$3$, $5$, $\partial$ the weight $e^{(1+3\beta_s)\kappa d_k(\bar X)}$ is
summed and the factor $e^{-\frac14\delta\kappa d_j(X)}$ is kept; the chain terms carry the norm
$\|\cdot\|^A$ at $a_k=\frac85\kappa$, of which $a=\frac{31}{20}\kappa$ is extracted. The rates are
those of \eqref{5.2:eq-resummation-input-f}. The fibre sum of the strip theorem is taken at
$\lambda_0:=\frac12\beta_\sharp\kappa$, under the member $2(\mathsf{k}_d+1)\le\kappa/20$ of
$\kappa_d$.

\emph{Filing at the localisation site.} Per unit the first estimate is the one the allocation
names: the Cauchy bound at the localisation site with the size table on the units of (c) and (d),
the source Cauchy bound of (P3) on those of (a) and (b). The three exponential factors of
\cite[(1.66)]{B-CMP122b} are obtained as follows. The length of \cite[(1.67)]{B-CMP122b} is given by
the length estimate for wrapped units of (P3), and by the tree argument of
Lemma~\ref{5.2:lem-value-unit-filing} for plain units. The $\kappa_1$-volume comes from the Cauchy
estimate in $\mathsf{s}$. The third factor is the three-factor bound --- (P5) for class-$4$ units,
its pointed analogue of Lemma~\ref{5.2:lem-plain-exterior-sum} for the classes $1$--$3$, $5$ --- at
rate $\frac12\delta_P$ in $d^{\wedge}_{\mathbf{B}''_k}$, in all four links.
Its sum is the deep-unit summation lemma of (P3) for wrapped units and
Lemma~\ref{5.2:lem-plain-exterior-sum} for plain units.

The chain terms of class $4$ reach the site with their own rate
\begin{equation*}
a_k=(1+3\beta_{\mathrm{pr}})\kappa=\tfrac85\kappa<\tfrac74\kappa=(1+3\beta_s)\kappa.
\end{equation*}
They are
therefore not summed under the weight of \cite[(1.66)]{B-CMP122b}. Their sum over $X$ at a cube is
the quantity whose supremum over cubes defines the norm $\|\cdot\|^A$, and so is at most
$C_0^\sharp$ per stage by (P1). From
$e^{-\frac85\kappa d_k(X)}$ one extracts the factor
\begin{equation*}
e^{-\frac{31}{20}\kappa d_{k,Z}(Y)}=e^{-(1+2\beta_s+\beta_\sharp)\kappa d_{k,Z}(Y)},
\end{equation*}
as Ba{\l}aban uses \cite[(1.66)--(1.67)]{B-CMP122b} (``This we use now to control the
resummations'', \cite[p.~375]{B-CMP122b}). Here $d_{k,Z}(Y)\le d_k(X)$ for the domain $Y=Y_1(X)$
under which the unit is filed, by the tree argument in the proof of
Lemma~\ref{5.2:lem-value-unit-filing}. Of the extracted rate, $(1+2\beta_s)\kappa=\frac32\kappa$ is
the output rate of \cite[(1.68)]{B-CMP122b}, and $\beta_\sharp\kappa=\kappa/20$ is the input margin
spent by the last resummation at the strip. The remaining factor $e^{-(\kappa/20)d_k(X)}$ sums the
domains, since $\#\{X\in\mathbf{D}_k:X\ni\Delta,\ d_k(X)\le n\}\le e^{c_d(n+1)}$
\cite[(1.26)]{B-CMP116} and $\kappa/20\ge c_d+2$ by \eqref{5.2:eq-resummation-input-f}. For these
terms the sum over $X$ is moreover the norm itself, so this count is a redundancy there. The
arithmetic is
\begin{equation*}
\tfrac85=\tfrac32+\tfrac1{20}+\tfrac1{20}:
\end{equation*}
one twentieth for the domain count and one for the strip's fibre sum. The floor $\kappa_d$ collects,
at the fraction $\frac1{20}$, the $d$-only conditions on $\kappa$ needed here. These are
$\kappa/20\ge c_d+2+2\sqrt d(2+4\beta_{\mathrm{pr}})$, for the stage norms and the domain count,
which also gives $\kappa\ge8(c_d+2)\ge4(c_d+2)$; $9\kappa\ge c_F+2\cdot3^d+8$; and
$\beta_\sharp\kappa\ge2(\mathsf{k}_d+1)$, for the strip's fibre sum. This is
\eqref{5.2:eq-resummation-input-f}.

For chain terms filed at a level $m<k$, with rate $a_m=(1-\beta_{\mathrm{pr}})\kappa=\frac45\kappa$
in $d_m(X)$ by (P1), the tree argument of the boundary-group lemma gives
$d_k(\bar X)\le L^{-(k-m)}d_m(X)+c_{\mathrm{tr}}$, with $c_{\mathrm{tr}}\le2d$ the constant of the
closure. Hence, as $L\ge3$,
\begin{equation*}
\tfrac{31}{20}\kappa\,d_{k,Z}(Y)\le\tfrac{31}{20}\kappa\bigl(\tfrac13d_m(X)+c_{\mathrm{tr}}\bigr),
\qquad
\bigl(\tfrac45-\tfrac{31}{60}\bigr)\kappa\,d_m(X)=\tfrac{17}{60}\kappa\,d_m(X)
\ge\tfrac{\kappa}{20}d_m(X),
\end{equation*}
and the last quantity pays the domain count at scale $m$ under the same $\kappa_d$. The factor
$e^{\frac{31}{20}c_{\mathrm{tr}}\kappa}$ joins the $O(1)$ inside $C_\sharp$, as $e^{3^d\kappa}$
does, and is met by $\mathsf{w}_L^{-1}$. The prefactor stays $C_\sharp\alpha$, met by
$\mathsf{w}_L^{-1}$, that is, by $\gamma$; it cannot be absorbed into the exponential decay, since
$C_\sharp$ contains $e^{3^d\kappa}$ and powers of $M$. The values of the value-unit remark,
chain-term values among
them, do not enter the weight supremum $\mathsf{W}$; they are filed by
Lemma~\ref{5.2:lem-value-unit-filing}, not here.

Ba{\l}aban's sentence ``For E-terms and R-terms this includes renormalization in the form described
in the proof of Theorem 2'' \cite{B-CMP122b} is carried out here as follows. For the groups at far
points it is the classes $1$--$3$, by the boundary-group lemma, clause (i): the core, its box and
the counterterm piece lie on one side of $\partial Z$, on one chain of backgrounds, and what
Ba{\l}aban sends to the other part are tails, units of their own, bounded singly. For the near
points it is the class $\partial$. The first-part values and Ba{\l}aban's subtracted copies are
treated by the boundary-group lemma and the value-unit remark. The volume bound is
\eqref{5.2:eq-resummation-input-e}, as shown above.

\emph{Conclusion.} The family $\sum_{Y_4\in\mathbf{D}_k^Z}\mathfrak{v}(Y_4)$ is indexed by
$\mathbf{D}_k^Z$, is analytic as stated, and satisfies \eqref{5.2:eq-resummation-input-b},
\eqref{5.2:eq-resummation-input-d}, \eqref{5.2:eq-wilson-remainders-a} and
\eqref{5.2:eq-resummation-input-e}. The strip theorem (P6), clauses (a)--(c), applied to it, with
$Z^{\sim}$ read as $Z^\sharp$ and $d_{k,Z}$ as $d_\sharp$, gives the bound
$|V^\sharp(Y)|\le\alpha^\sharp e^{-(1+2\beta_s)\kappa
d_\sharp(Y)}$ for $Y\setminus Z^\sharp\neq\emptyset$, with $\alpha^\sharp$ as in
\eqref{5.2:eq-resummation-window-a}, on the tube times $\mathcal{W}_k(Y)$. The corollary of the
strip theorem is stated with a single per-cell constant, in its hypothesis and in its conclusion.
Its proof is the summation of the volume bound over the cells of $C$ together with the strip's own
units. Carried out with the right member of \eqref{5.2:eq-resummation-input-e} in place of the
single constant, that summation gives, zone by zone, the second line of
\eqref{5.2:eq-resummation-input-1}. That line is thus the summed form of the corollary's last line,
not the corollary
itself. Per cell of zone $j$, $\mathsf{K}_{\mathrm{cell}}(k,j)\le2^{-d}(C_\sharp O(1)+1)\log
g_j^{-2}$ by \eqref{5.2:eq-per-cell-allowance-a}. This is one more logarithmic term inside
Ba{\l}aban's constant $O(1)\log g_j^{-2}|Z_j|$ \cite{B-CMP122b}, where Ba{\l}aban's bound for the
running chain reads $O(1)\sum_j|\Gamma^0_j\cap Y|$ \cite[(1.69)]{B-CMP122b}. The frozen
contributions obey, by \eqref{5.2:eq-arrested-regions-a} and the factor $c_V^S$,
\begin{equation*}
\mathsf{P}(\mathfrak{a})\le c_V^Sc_{\mathsf{P}}\sum_n\check R^{d+1}_{j_n+1}\sum_X
\bigl(d'_{j_n+1}(X)+1\bigr),
\end{equation*}
the sum running over the stages $n$ of the arrest, $j_n$ being the level integrated at stage $n$
and $X$ the components of $Z_{j_n+1}$ contained in the arrest's component; they are extensive at
each frozen term's own
scale. They lie inside the third member $O(1)M^dR_j^{d+1}d'_j(Z_j)$ of the same bound, with the
constant $2c_V^Sc_{\mathsf{P}}M^{-d}$. Read in Ba{\l}aban's notation, with $Z^{\sim}$ for $Z^\sharp$,
$d_{k,Z}$ for $d_\sharp$ and $C_\sharp\alpha$ for $\alpha^\sharp$, these are the two lines of
\eqref{5.2:eq-resummation-input-a}.
\end{proof}

\begin{lemma}[Localised remainders of the expanded Wilson actions
{\cite[(1.11), (1.23), (1.27), (1.31), (1.47)--(1.49), pp.~374--375]{B-CMP122b}}]
\label{5.2:lem-wilson-remainders}
In the setting of Definition~\ref{5.1:def-source-free-scaffold}, with the constants chosen in the
order of Definition~\ref{5.1:def-admitted-inputs}, let $k\le K$.
Let $Z$ be the class of components of $\Lambda_k^c$ treated by the large-field
operation $\mathbf{R}$ at step $k$, and let $\mathbf{C}_k$ be the sum of the new terms of the
effective action that arise from expanding Wilson actions, including the term $I(U_0)$ from the
denominator, in the sense of \cite[p.~374]{B-CMP122b}. The members of $\mathbf{C}_k$ are:
\begin{itemize}
\item[(Ca)] the remainders of the expansions \cite[(1.17), (1.20), (1.24), (1.29), (1.36),
(1.46)--(1.48)]{B-CMP122b} on $Z$;
\item[(Cb)] their counterparts on $\Omega_k$ off $Z$, cut by the decomposition of unity
connected with the partition $\pi_k$;
\item[(Cc)] the term $I(U_0)$.
\end{itemize}
These members are localised by the same four operations that are applied to every term of the
first part of the effective action, namely \cite[(1.59) or (1.60), and (1.61)--(1.68)]{B-CMP122b}.
Let $\mathfrak{P}$ be the parameter polydisc of the localisation, with points
$m=(\zeta,\varpi,\vec{\mathsf{s}})$: the shift $\zeta$, the variable $\varpi$ of the rim
construction, and the decoupling parameters $\vec{\mathsf{s}}$.
For $Y_4$ in the class $\mathbf{D}_k^Z$ of pre-filed localisation domains of level $k$ (a domain
of $\mathbf{D}_k$ enlarged by the components of $Z$ that it meets) let
$\mathfrak{v}^{\mathbf{C}}(Y_4)$ be the sum of the pieces and values filed at $Y_4$ of the part of
(Ca)--(Cc) that remains after the term $\mathfrak{c}_w$ of the source part of the Wilson action
(Lemma~\ref{5.3:lem-wilson-source}), which carries the dependence on the complex source
$w=\sum_i z_i f_i$, is removed. The filing is as follows:
\begin{itemize}
\item a member on cells of a component $C$ of $Z$ is filed as a plain interior term of the
first part. It carries no block cutoff $\hat{\chi}_{\square}$ and no amplitude $\mathsf{t}$, and
it is anchored at its own cell $\Delta'\subset C$ of the multi-scale partition $\mathbf{B}''_k$
produced by the preparatory operations, its
anchor-and-weight domain $K_v$ being the union of the $\pi_k$-cubes that meet $\Delta'$;
\item a member that depends on the shift $\zeta$ itself, respectively on $\mathsf{V}$,
the $k$-fold block average of $\mathbf{U}$ restricted to $C\setminus\Lambda$, at level-$k$ sets is
filed as a value at $Y_4 = C$, without link families; here $\Lambda$ is the region whose volume
$|\Lambda|$ enters the bound \cite[(1.11)]{B-CMP122b} for $I(U_0)$, which Ba{\l}aban identifies
with one of its components; so identified, $\Lambda$ is a rectangular parallelepiped contained in
a cube of size $100M$. These are the part
of the new Wilson action on $Z$ and the members containing $U_0$, among them $I(U_0)$;
\item a member of (Cb) is filed at its own $\pi_k$-cube off $Z$.
\end{itemize}
Put
\begin{equation*}
a := (1+2\beta_s+\beta_\sharp)\kappa = \tfrac{31}{20}\,\kappa ,
\qquad \beta_s = \tfrac14,\quad \beta_\sharp = \tfrac1{20}.
\end{equation*}
For $Y\in\mathbf{D}_k^Z$ let $d_{k,Z}(Y)$ be the relative linear size of $Y$ off $Z$ of
\cite[(1.67)]{B-CMP122b}: the length, in units of $M$, of a shortest tree graph contained in
$Y$ and intersecting all $M$-cubes of the components of $Y\setminus Z$. Let $C^\sharp\supset C$
be the strip of width
$m_\sharp(k)=\lfloor\frac12(\check{R}_k-1)\rfloor$ about a component $C$ of $Z$. Set
\begin{equation*}
\|\mathfrak{v}^{\mathbf{C}}\|_{a,1}
:= \sup_{\Delta\not\subset Z}\ \sum_{Y_4\ni\Delta} e^{a\,d_{k,Z}(Y_4)}
\sup\bigl|\mathfrak{v}^{\mathbf{C}}(Y_4)\bigr| ,
\end{equation*}
where $\Delta$ runs over the $\pi_k$-cubes; here and in \eqref{5.2:eq-wilson-remainders-a} the
supremum of $|\mathfrak{v}^{\mathbf{C}}(Y_4)|$ is taken over the tube over $Y_4$, that is, over
the space of complex configurations over $Y_4$ on which the terms filed at $Y_4$ are analytic
(the analyticity spaces of Definition~\ref{5.1:def-analyticity-spaces}).

Then each $\mathfrak{v}^{\mathbf{C}}(Y_4)$ is independent of the source $w$ and analytic on the
tube over $Y_4$, and the family $(\mathfrak{v}^{\mathbf{C}}(Y_4))_{Y_4}$ is jointly invariant
under $\mathrm{SU}(N)$-valued gauge transformations acting simultaneously on all the background
fields, in the gauge-transformation law of the construction that produces these background fields.
Let $\Pi$ be the law of $\varpi$. Moreover, uniformly in real
$\varpi\in\operatorname{supp}\Pi$,
\begin{equation}\label{5.2:eq-wilson-remainders-a}
\begin{gathered}
\|\mathfrak{v}^{\mathbf{C}}\|_{a,1}\le\mathsf{c}_{\mathbf{C}},\\
\sum_{Y_4\subset C^\sharp}\sup\bigl|\mathfrak{v}^{\mathbf{C}}(Y_4)\bigr|\le\mathsf{c}_{\mathbf{C}}
\quad\text{for every component $C$ of $Z$,}
\end{gathered}
\end{equation}
where $\mathsf{c}_{\mathbf{C}} := 1+c_I M^5$. Here $c_I$ is the constant written $O(1)$ in
\cite[(1.11)]{B-CMP122b}; it is a number fixed after $M$.
\end{lemma}

The bounds \eqref{5.2:eq-wilson-remainders-a} are an input of the correlation theorem
(Theorem~\ref{5.1:thm-correlation-theorem}) taken by statement, among the admitted inputs of
Definition~\ref{5.1:def-admitted-inputs}; they are cited from the loci of \cite{B-CMP122b} given
in the heading, in the reading stated next; the proof is Ba{\l}aban's and is not reproduced.
We read Ba{\l}aban's phrase ``their sum is also small'' of the absolute sum over localisation
domains: both members of \eqref{5.2:eq-wilson-remainders-a} are absolute sums over localisation
domains, the first carrying, on the domains leaving $Z$, the weight $e^{a\,d_{k,Z}(Y_4)}$ of
\cite[(1.66)--(1.68)]{B-CMP122b}.
The individual loci are as follows.

\emph{The object.} In \cite[(1.49)]{B-CMP122b} the first part of the effective action includes
the first three terms of that formula, of which $\mathbf{C}_k$ is the third. Ba{\l}aban describes
the terms collected in $\mathbf{C}_k$ as follows \cite[p.~374]{B-CMP122b}:
``many new terms in the effective action. They have been obtained by expanding some Wilson
actions \dots\ easily localizable, but they depend on several different background fields \dots\
they are all small, in fact their sum is also small, although we need boundedness only, and they
depend on background fields restricted to a neighborhood of the region $Z$, therefore they are
boundary terms \dots\ We denote the sum of all these terms, including the term $I(U_0)$ from the
denominator, by $\mathbf{C}_k$''. The members of (Cb) are the terms obtained with
``a decomposition of unity connected with the partition $\pi_k$, and terms of the obtained sum
are small, and exponentially small in the distance to $\Omega_k^c$'' \cite{B-CMP122b}.

\emph{The sizes.} The sizes of the members of (Ca) are
\cite[(1.23), (1.27), (1.31), (1.47)]{B-CMP122b}, and after \cite[(1.47)]{B-CMP122b}
Ba{\l}aban writes: ``The last bound above is small under the usual conditions on $N$.
Bounds for the other terms in the expansion are similar. Thus we have'' the bounds
\cite[(1.48)]{B-CMP122b}, ``where the terms in
$O(1)$ are small, and depend analytically on the field $\mathbf{U}$ restricted to the domain
$Z$''. The size of the member (Cc) is \cite[(1.11)]{B-CMP122b}: ``the analytically extended
$I(U_0)$ satisfies the bound $|I(U_0)| < O(1)\log M|\Lambda| < O(1)M^5$''.

There is one term $I(U_0)$ per component of $Z$. The reason is that the construction of
\cite{B-CMP122b} proceeds by ``identifying $\Lambda$ with one of its components'', and ``$\Lambda$
is a rectangular parallelepiped contained in a cube of the size $100M$''. In addition, the
penultimate member of \cite[(1.47)]{B-CMP122b} is a sum over a whole component. This is why the
constant $\mathsf{c}_{\mathbf{C}}$ in \eqref{5.2:eq-wilson-remainders-a} is a constant per
component.

The analyticity on the tube rests on the sentence ``the bounds are preserved if we replace this
field by the corresponding $G^c$-valued field. This replacement will be described more precisely
later, when we will discuss localization'' \cite{B-CMP122b}; here $G^c$ is the complexification
of the gauge group $G=\mathrm{SU}(N)$.

\emph{The localisation and the weight.} The localisation pages include $\mathbf{C}_k$
explicitly \cite[p.~375]{B-CMP122b}: ``Now we consider terms of the first part of the effective
action in (1.49), except the first two terms, and all the above constructed new terms with the
localization domains $Y_0$. The first operation is exactly the same as in the previous case, i.e.,
we apply either the equality (1.59), or (1.60), and the localization operation is \dots''. Thus
the members of $\mathbf{C}_k$ pass through \cite[(1.59) or (1.60), and (1.61)--(1.68)]{B-CMP122b}.
For the domains that leave $Z$, the weight $e^{a\,d_{k,Z}(Y_4)}$ in
$\|\mathfrak{v}^{\mathbf{C}}\|_{a,1}$ is the weight of \cite[(1.66)--(1.68)]{B-CMP122b}; for the
members of (Cb) the decay off $Z$ is the exponential smallness in the distance to $\Omega_k^c$
quoted above. Only boundedness of these terms is used below, in accordance with the words
``although we need boundedness only'' quoted above.

\emph{Members on the cells of a component.} For the members of (Ca) on the cells of a component
$C$ of $Z$ the weighted bound below follows from the lemma on the Wilson densities read along the
chain of backgrounds and the summation lemma for link families; it supplements, and does not
replace, \eqref{5.2:eq-wilson-remainders-a}, which is used as stated.
Let $\mathcal{U}(m)$ be the $U$-slot of the carrier pair at the point $m$ of $\mathfrak{P}$.
Let $\Pi_0$ (unrelated to the law $\Pi$) be the set of plaquettes all four of whose
bonds are zone-$0$ bonds of the multi-scale partition $\mathbf{B}''_k$ produced by the
preparatory operations. Since the level $h$ of the field $U''_{h+2}$ in
\cite[(1.49)]{B-CMP122b}, which is also the index in the radius $\check{R}_{h+1}$ below,
satisfies $h\ge1$ by \cite[(ii)]{B-CMP122a}, these
plaquettes lie in the region $\mathsf{F}$ fixed in the lemma on the Wilson densities read along
the chain of backgrounds
and, at $C$, in the core $\mathsf{C}=\mathsf{C}_C$ of $C$; their summands cancel identically in
the difference \cite[(1.59) or (1.60)]{B-CMP122b} of the density evaluated on two backgrounds,
because the zone-$0$ datum on which they depend is the same in both.
The exclusion of the plaquettes of $\Pi_0$ from the presentation of these densities is a reading
adopted in that lemma.
On cells $\Delta'\subset C$ of the multi-scale partition the Wilson densities read along the chain
are
\begin{equation*}
f^W_{\Delta'}(m) := \sum_{p:\ x(p)\in\Delta',\ p\notin\Pi_0} c(p)\,
a_p\bigl(\mathcal{U}(m)\bigr),
\qquad |c(p)|\le x''(p) := x_{\pi(x(p))},
\end{equation*}
with coefficients $c(p)$, $a_p$ the plaquette density, and $\pi(y)$ the index of the zone
containing the point $y$. They are holomorphic on $\mathfrak{P}$, and by the same lemma
they satisfy:
\begin{enumerate}
\item[(i)] for $\Delta'\subset C$, $\Delta'\not\subset\mathsf{C}$, with
$\mathfrak{N}^W(\Delta')$ the majorant of the cell density supplied by that lemma,
\begin{equation*}
\sup_{\mathfrak{P}}|f^W_{\Delta'}|\le\mathfrak{N}^W(\Delta');
\end{equation*}
\item[(ii)] for $\Delta'\subset\mathsf{C}$ and $m,m'\in\mathfrak{P}$ with the same component
$\varpi$, with $\mathfrak{O}^W(\Delta')$ the majorant of the oscillation supplied by that lemma,
\begin{equation*}
|f^W_{\Delta'}(m)-f^W_{\Delta'}(m')|\le\mathfrak{O}^W(\Delta'),
\end{equation*}
with $\mathfrak{O}^W(\Delta')=0$ at every level-$0$ cell $\Delta'$ of $\mathsf{F}$ that touches
no level-$1$ cell;
\item[(iii)] with $N(g_k)$ the number of preparatory stages at step $k$, $\mathsf{n}_\partial(k)$
and $\mathsf{n}_{\mathrm{sh}}(k)$ the two cell counts of that lemma, and $C_W$ a number fixed
before $\gamma$ (a constant of that lemma, unrelated to the constant $C_W(s)$ of the two-point
witness theorem), independent of the volume of the core, of the ages of the large-field regions
involved, and of the further datum $D$ of that lemma,
\begin{equation*}
\begin{gathered}
\sum_{\Delta'\subset\mathsf{C}}\mathfrak{O}^W(\Delta')
\le C_W(1+N(g_k))\check{\mathsf{T}}_k^{2q}\,\mathsf{n}_\partial(k),\\
\sum_{\Delta'\subset C,\ \Delta'\not\subset\mathsf{C}}2\,\mathfrak{N}^W(\Delta')
\le C_W\,\mathsf{n}_{\mathrm{sh}}(k)\,\check{R}_{h+1}^2\,\check{\mathsf{T}}_k^{2q}.
\end{gathered}
\end{equation*}
\end{enumerate}
Read these members as plain interior units of the first part with $d_k(K_v)=0$. The pieces of
these members whose filing domains leave $Z$ have the $\|\cdot\|_{a,1}$-norm of their sum at most
\begin{equation*}
e^{-(\check{R}_k-2)}(1+C_\Sigma C_e)^4
\Bigl(C_W(1+N(g_k))\check{\mathsf{T}}_k^{2q}\,\mathsf{n}_\partial(k)
+C_W\,\mathsf{n}_{\mathrm{sh}}(k)\,\check{R}_{h+1}^2\,\check{\mathsf{T}}_k^{2q}\Bigr),
\end{equation*}
where $C_\Sigma$, $C_e$ are the constants of the summation lemma for the link families of a
piece.
This follows from the bound for pieces leaving $Z$ in the localisation of plain units, applied
verbatim, with the reference value taken at the origin of the parameters: a non-empty link
family contains a source cell at distance at least $9\check{R}_k-1$ from $Z^c$, so that the
label reached after the four link families has level-$k$ length at least $\check{R}_k-2$; and the
summation lemma for link families is applied four times under the floor
\begin{equation*}
\kappa_1-1\ge c_l(a+9)+\lambda_*,\qquad c_l=1+234d ,
\end{equation*}
with $\lambda_*$ as in the summation lemma for link families.
The right member is of degree $-\infty$: a function of polynomial type in
$\check{\mathsf{T}}_k$ and $\check{R}_{h+1}$ multiplied by $e^{-(\check{R}_k-2)}$; no condition
on $\gamma$ is taken from
it. What is taken from \cite{B-CMP122b} as printed, with no such supplement, is: the members
(Cb), the member (Cc) together with the members containing $U_0$, the joint invariance, and the
identification of Ba{\l}aban's intermediate localisations beyond their linearity in the weights.

\emph{Place in the filing.} The part of $\mathbf{C}_k$ depending on $w$ is the term
$\mathfrak{c}_w$ of the source part of the Wilson action (Lemma~\ref{5.3:lem-wilson-source},
display~\eqref{5.3:eq-wilson-source-a}). The present lemma concerns the part independent of $w$.
Every piece of this part, for every tuple of link families, is filed in $\mathfrak{v}^{\mathbf{C}}$
and in none of the members of $\varepsilon_{\mathrm{br}}(k)$. Accordingly, the lemma enters the
subsequent estimates in four places:
\begin{itemize}
\item in the strip resummation theorem for the pre-filing family $\mathfrak{v}$, which holds
$\mathfrak{v}^{\mathbf{C}}$ beside the bracket and value terms, the weighted supremum norm
\begin{equation*}
\|\mathfrak{v}\|_a := \sup_{Y_4\setminus Z\neq\emptyset}e^{a\,d_{k,Z}(Y_4)}|\mathfrak{v}(Y_4)|
\end{equation*}
of the pre-filing family is at most
$\varepsilon_{\mathrm{br}}(k)+\mathsf{K}_{\mathrm{val}}(k)+\mathsf{c}_{\mathbf{C}}$, the first
two summands being the bracket bound and the value bound of that theorem,
so that $\mathsf{K}^0$ is read as $\mathsf{K}^0+\mathsf{c}_{\mathbf{C}}$ in the width function
$m_{\mathrm{lo}}(k)$ and in the two ceilings on $\gamma$ that carry the number
$2+\mathsf{K}^0$. The constant $\mathsf{c}_{\mathbf{C}}$ is a number fixed after $M$, and the six
members of $\varepsilon_{\mathrm{br}}(k)$ and their ceilings do not change;
\item in the volume bounds for the filed terms inside a component of $Z$ and its strip,
$\mathsf{c}_{\mathbf{C}}$ is added once per component of $Z$ and
charged at one level-$k$ cell of the collar of that component, so that
$\mathsf{K}_{\mathrm{cell}}^0$ is read as $\mathsf{K}_{\mathrm{cell}}^0+\mathsf{c}_{\mathbf{C}}$;
\item in the theorem on the base values of the strip sums, the oscillation of
$\sum_{Y_4\subset C^\sharp}\mathfrak{v}^{\mathbf{C}}(Y_4)$ enters its sum, and the second
bound in \eqref{5.2:eq-wilson-remainders-a} gives the majorant $2\mathsf{c}_{\mathbf{C}}$;
\item in the bound of the first and third brackets of $\mathfrak{c}_w$ in
\eqref{5.3:eq-wilson-source-a}, which is $\theta_W$ times Ba{\l}aban's bounds for the
corresponding members of $\mathbf{C}_k$, by the comparison lemma for complex against real
plaquette weights, in its clause on the expanded Wilson actions.
\end{itemize}

\begin{lemma}[Filing of the value units]\label{5.2:lem-value-unit-filing}
Work at the localisation site of step $k$, with $Z$ the class of components of $\Lambda_k^c$
treated there, and let
\begin{equation*}
a := (1+2\beta_s+\beta_{\sharp})\kappa = \tfrac{31}{20}\,\kappa
\end{equation*}
be the weight of the norm $\|\cdot\|_a$. Let $\mathcal{V}$ be the finite set of value units at
this site, in the sense of the remark on the values at the localisation site. These are:
\begin{itemize}
\item the constituents of the first part of the effective action whose localisation expansion is
trivial, that is, all of whose domains $Y'_i$ are empty; these are Ba{\l}aban's ``exceptional
terms, for which all the domains $Y'$ are empty'', which he bounds ``by a constant times the
same exponential factor'' \cite[\S1]{B-CMP122b};
\item the copies of terms that Ba{\l}aban's construction subtracts at the localisation site
\cite[\S1]{B-CMP122b};
\item the values of boundary terms and of chain terms, namely
\begin{itemize}
\item the boundary terms $\mathbf{B}^{(j)}$ born at an operation $\mathbf{T}$;
\item the boundary terms $\mathbf{B}'^{(j)}$ born at an operation $\mathbf{R}$, together with
their wrapped descendants born at a later $\mathbf{T}$;
\item the chain terms $\mathbb{C}^{(n)}_k$ of the chains of this step;
\item the frozen chain terms $\mathbb{C}^{(i)}_{k_1}$, $k_1<k$, and the frozen
change-of-variables terms of the live pieces of chains arrested at earlier steps.
\end{itemize}
\end{itemize}
Each $v\in\mathcal{V}$ carries the following data.

\emph{Label and weight domain.} The unit $v$ has a label $A_v\in\mathbf{D}_m$, $m\le k$, and an
\emph{anchor-and-weight domain}
\begin{equation*}
\mathfrak{K}_v := K_v := [A_v]_k \in \mathbf{D}_k ,
\end{equation*}
the union of the $\pi_k$-cubes meeting $A_v$; only units whose anchor-and-weight domain meets
$Z$ belong to $\mathcal{V}$. For a plain unit $A_v$ is its localisation domain $X$, and
$\mathfrak{K}_v=\bar X$ is the closure of $X$ at level $k$. The bound for the closure of a label
at level $k$ and the tree
argument for the closure of a localisation domain, run on $A_v\in\mathbf{D}_m$, give
\begin{equation*}
d_k(K_v)\le L^{-(k-m)}d_m(A_v)+c_{\mathrm{tr}} .
\end{equation*}
With $\mathfrak{K}_v$ as above, and as modified below for exposed pointed units, the
\emph{filing domain} of $v$ is
\begin{equation*}
Y_1(v) := \mathfrak{K}_v\cup\{\text{components of } Z \text{ meeting } \mathfrak{K}_v\}
\in\mathbf{D}_k^Z .
\end{equation*}

\emph{Exposed pointed units.} An exposed pointed unit of $\mathcal{V}$ is a short single from the
second, third or boundary family of the table of unit sizes, or a Wilson piece astride
$\partial Z$. For such a unit one sets
\begin{equation*}
\mathfrak{K}_v := K_v\cup\mathfrak{h}_v = \operatorname{lab}_1 ,
\end{equation*}
the label $\operatorname{lab}_1$ of exposed units fixed in the localisation conventions, so that
\begin{equation*}
Y_1(v)=Y(\operatorname{lab}_1)=X^{\oplus}_v ,
\end{equation*}
where $Y(\cdot)$ adjoins to a union of cubes the components of $Z$ it meets, as in the
definition of $Y_1(v)$. Here $f_v$ denotes the unit as a function of the carrier,
$\operatorname{chain}(\vec{\mathsf{s}})$ the carrier at the decoupling parameters
$\vec{\mathsf{s}}$, and $f_v(1)$ the value of the unit on the undecoupled carrier, which is the
value entered for $v$ in the table of unit sizes. The value
$\mathrm{val}_v=f_v(\operatorname{chain}(\vec 0))$ is then filed where the bound
$|\mathrm{val}_v-f_v(1)|\le\mathfrak{N}^{\flat}(v)$ is proved, and the carrier estimate from which
that bound follows is applied to the auxiliary consumer term $\mathfrak{E}^{\oplus}_v$ inside its
hypothesis that the unit's domain lie in the filing domain (written there $X\subset Y_4$), for the
tuple with all families empty as well. In this case
$\mathfrak{K}_v$ is connected and
\begin{equation*}
d_k(\mathfrak{K}_v)\le d_k(K_v)+c_l c_{\mathfrak{h}} ,
\end{equation*}
and the factor $e^{a c_l c_{\mathfrak{h}}}$ that this costs the weight is contained in
$\mathsf{K}^{\flat,\mathrm{ex}} = e^{(a+9)c_lc_{\mathfrak{h}}}\,27S_d'\mathsf{K}^{\flat}$.

\emph{Value.} The unit $v$ has a value $\mathrm{val}_v$, its piece with all families empty. The
value is analytic on the tube over $Y_1(v)^{+}\supset X_v$, and
$\|v\| := \sup|\mathrm{val}_v|$ on that tube.

\emph{Filing.} Let $\mathcal{P}_v$ denote the set of large-field regions, each with its strip
$Y^{\sharp}$, that the format of boundary-type terms attaches to $v$. The strips of members of
$\mathcal{P}_v$ that are no longer live lie in components of $Z$ meeting $K_v$, and the live ones
touch $A_v\subset K_v$; hence $X_v\subset Y_1(v)^{+}$. Define
\begin{equation*}
\mathfrak{v}^{\mathrm{val}}(Y_1) := \sum_{v:\,Y_1(v)=Y_1}\mathrm{val}_v ,
\end{equation*}
which is analytic over $Y_1^{+}$. The strip resummation, its corollary, the combinatorial
lemma of the last resummation, and the combinatorics of \cite[(1.93)--(1.99)]{B-CMP122b} read
$\mathfrak{v}^{\mathrm{val}}$ on $Y_1$, the strip resummation being stated for any family indexed
by $\mathbf{D}_k^Z$. By part~(a) of the proposition on boundary terms over live strips, the older
live strips enter no sum, no length and no compatibility relation there. By part~(a) of the
proposition on the localisation of wrapped units and part~(b) of the proposition on the volume
sums of wrapped units, a unit for which a member of $\mathcal{P}_v$ ceases to be live at the
present $\mathbf{R}_k$ is a first-part unit of that component and ceases here. A deep unit, one
whose true domain $X_v$ lies in the core $\mathsf{C}_C$ of a component $C$ of $Z$, has
$K_v\subset Z$.

\emph{Exterior units and the value constant.} Set
$\mathcal{V}^{\mathrm{ext}} := \{v\in\mathcal{V}:\ \mathfrak{K}_v\setminus Z\neq\emptyset\}$;
in particular no deep unit lies in $\mathcal{V}^{\mathrm{ext}}$. Set
\begin{equation}\label{5.2:eq-value-unit-filing-a}
\mathsf{K}_{\mathrm{val}}(k) := \sup_{\Delta\not\subset Z}
\sum_{v\in\mathcal{V}^{\mathrm{ext}},\ \mathfrak{K}_v\supset\Delta} e^{a d_k(\mathfrak{K}_v)}\,\|v\| .
\end{equation}

Then the value input \eqref{5.2:eq-resummation-input-d} of the boundary-term input at the
resummation site (Proposition~\ref{5.2:prf-resummation-input}) holds in the form
\begin{equation}\label{5.2:eq-value-unit-filing-1}
\|\mathfrak{v}^{\mathrm{val}}\|_a := \sup_{Y_1\setminus Z\neq\emptyset}
e^{a d_{k,Z}(Y_1)}\,|\mathfrak{v}^{\mathrm{val}}(Y_1)| \le \mathsf{K}_{\mathrm{val}}(k).
\end{equation}
\end{lemma}

\begin{proof}
Here two sets \emph{meet} when they share a $\pi_k$-cube, and no connector cubes are added in
forming $Y_1(v)$. By definition $Y_1(v)\in\mathbf{D}_k^Z$. The components of $Z$ adjoined to
$\mathfrak{K}_v$ lie in $Z$, so
\begin{equation}\label{5.2:eq-value-unit-filing-2}
Y_1(v)\setminus Z = \mathfrak{K}_v\setminus Z .
\end{equation}

Fix $Y_1$ with $Y_1\setminus Z\neq\emptyset$. Since $Y_1\setminus Z$ is a non-empty union of
$\pi_k$-cubes, we may choose a fixed enumeration of the current cubes and let $\Delta_*$ be the
first cube of $Y_1\setminus Z$. Then $\Delta_*\not\subset Z$.

Let $v$ belong to the fibre $\{v:\ Y_1(v)=Y_1\}$. By \eqref{5.2:eq-value-unit-filing-2},
$\mathfrak{K}_v\setminus Z = Y_1\setminus Z\neq\emptyset$, so $v\in\mathcal{V}^{\mathrm{ext}}$.
Moreover
\begin{equation*}
\mathfrak{K}_v\supset\mathfrak{K}_v\setminus Z = Y_1\setminus Z\ni\Delta_* .
\end{equation*}

Next we compare the two lengths. Recall from \cite[(1.66)--(1.67)]{B-CMP122b} that
$d_{k,Z}(Y)$ is the length of a shortest tree graph contained in $Y$ and intersecting all
$M$-cubes of the components of $Y\setminus Z$. It thus measures only the part of a tree needed
off $Z$. Take an optimal tree for $d_k(\mathfrak{K}_v)$. It lies in
$\mathfrak{K}_v\subset Y_1(v)=Y_1$, and it meets every cube of $\mathfrak{K}_v$. Since
$\mathfrak{K}_v\supseteq Y_1\setminus Z$, it meets in particular every cube of the components of
$Y_1\setminus Z$. It is therefore admissible in the definition of $d_{k,Z}(Y_1)$, and
\begin{equation}\label{5.2:eq-value-unit-filing-3}
d_{k,Z}(Y_1)\le d_k(\mathfrak{K}_v)\qquad\text{for every } v \text{ with } Y_1(v)=Y_1 .
\end{equation}

At every point of the tube over $Y_1^{+}$ (which is $Y_1(v)^{+}$ for every $v$ of the fibre) we
have $|\mathfrak{v}^{\mathrm{val}}(Y_1)|\le\sum_{v:\,Y_1(v)=Y_1}|\mathrm{val}_v|$, and each term is
at most $\|v\|$; hence the same bound holds for the supremum over the tube. Combining this with
\eqref{5.2:eq-value-unit-filing-3}, and with the fact that every member of the fibre lies in
$\mathcal{V}^{\mathrm{ext}}$ and has $\mathfrak{K}_v\supset\Delta_*$ (all terms being
non-negative), we obtain
\begin{align*}
e^{a d_{k,Z}(Y_1)}\,|\mathfrak{v}^{\mathrm{val}}(Y_1)|
&\le \sum_{v:\,Y_1(v)=Y_1} e^{a d_k(\mathfrak{K}_v)}\,\|v\| \\
&\le \sum_{v\in\mathcal{V}^{\mathrm{ext}},\ \mathfrak{K}_v\supset\Delta_*}
e^{a d_k(\mathfrak{K}_v)}\,\|v\| \\
&\le \mathsf{K}_{\mathrm{val}}(k).
\end{align*}
The last inequality is \eqref{5.2:eq-value-unit-filing-a} at $\Delta=\Delta_*\not\subset Z$.
Taking the supremum over $Y_1$ with $Y_1\setminus Z\neq\emptyset$ gives
\eqref{5.2:eq-value-unit-filing-1}.

The weight in \eqref{5.2:eq-value-unit-filing-a} is taken on the label, in which the unit's own
bound decays. Thus no live strip is ever weighed at plain length, whether it lies inside $Z$ or
belongs to an older member of a set $\mathcal{P}_v$ beside it.

No sum over anchors and no count of domains at the reduced rate $\kappa/20$ occurs in this
argument. For a plain unit, $A_v=X$ and $\mathfrak{K}_v=\bar X$, as in the statement.

Both the notion of an exterior unit (membership in $\mathcal{V}^{\mathrm{ext}}$) and the collar
bound \eqref{5.2:eq-exterior-collar-bound-b} are read on $\mathfrak{K}_v$: a unit is exterior if
and only if its label leaves $Z$. The collar bound is a statement about a connected union of
cubes having a cube in $\Omega_k^c$, and the label $A_v\in\mathbf{D}_m$ is such a union. Indeed,
$A_v$ meets $Z_m^{\boxplus}$ or has a point on a strip of level at most $m$, and these sets lie in
$\Omega^c_{m+1}\subset\Omega_k^c$ for $m<k$. So the label crosses the collar like any domain, as
part~(b) of the proposition on the localisation of wrapped units records.
\end{proof}

\begin{lemma}[The collar bound for exterior value units]\label{5.2:lem-exterior-collar-bound}
Let $k\le K$, and assume the standing smallness conditions on $\gamma$ of this chapter (in particular
$\mathsf{T}_\gamma\ge2$). Let $Z$, the finite set $\mathcal{V}$ of value units at the localisation site
of step $k$, their labels $A_v$, their anchor-and-weight domains $\mathfrak{K}_v$, their
values $\mathrm{val}_v$ with $\|v\|:=\sup|\mathrm{val}_v|$, the exterior class
$\mathcal{V}^{\mathrm{ext}}=\{v\in\mathcal{V}:\mathfrak{K}_v\setminus Z\neq\emptyset\}$ and the
constant $\mathsf{K}_{\mathrm{val}}(k)$ be as in the filing of the value units
(Lemma~\ref{5.2:lem-value-unit-filing} and \eqref{5.2:eq-value-unit-filing-a}). Thus
$\mathfrak{K}_v=K_v=[A_v]_k$, except for the exposed pointed units (the short singles of lines $2$,
$3$, $\partial$ of the size table and the Wilson pieces astride $\partial Z$), for which
$\mathfrak{K}_v=K_v\cup\mathfrak{h}_v$, with $\mathfrak{h}_v$ the $\pi_k$-neighbourhood of the unit's
read set, and $d_k(\mathfrak{K}_v)\le d_k(K_v)+c_lc_{\mathfrak{h}}$. For a plain unit, that is, a unit
not carried in the relative format over live strips, the label is its domain,
$A_v=X$, and $K_v=\bar X:=[X]_k$; thus $\mathfrak{K}_v=\bar X$ unless the unit is an exposed pointed
unit, in which case $\mathfrak{K}_v=\bar X\cup\mathfrak{h}_v$. Put
$a:=\tfrac{31}{20}\kappa=(1+2\beta_s+\beta_\sharp)\kappa$, for a unit of level $j\le k$ write
$s:=k-j$ for the age of its level, put $c_s:=\log(1+3\lambda s+2c_2)$, and define
\begin{equation}\label{5.2:eq-exterior-collar-bound-1}
\mathsf{K}'_{\mathrm{val}}(k):=\sup_{\Delta\not\subset Z}
\sum_{v\in\mathcal{V}^{\mathrm{ext}},\ \mathfrak{K}_v\supset\Delta}
e^{(a+8)d_k(\mathfrak{K}_v)}\,\mathfrak{N}(v).
\end{equation}
Assume the following.
\begin{itemize}
\item \emph{Stage norms} (the stage-norm corollary for the preparatory chains). For every step
$k_1\le k$ and every stage $n$ of a chain of step $k_1$, the
chain terms satisfy $\|\mathbb{C}^{(n)}_{k_1}\|^{A}_{(n)}\le C_0^{\sharp}$, with rate
$a_{k_1}=(1+3\beta_{\mathrm{pr}})\kappa=\tfrac85\kappa$ at the top level $k_1$, in level-$k_1$
units, and rate $a_m=(1-\beta_{\mathrm{pr}})\kappa=\tfrac45\kappa$ at levels $m<k_1$, in level-$m$
units. At every cube $\square$ of a given level the norm dominates the sum over the domains
$X\ni\square$ of that level of the exponential weight at the rate of that level times the supremum of
$|\mathbb{C}^{(n)}_{k_1}(X)|$ on the tube; this supremum is bounded by $\|\cdot\|^{A}$ because
$0\in\mathfrak{D}_{\mathcal{S}}$. The number of stages of a chain of step $k_1$ is at most $N(g_{k_1})$.
\item \emph{Frozen terms of live arrests} (the definition of an arrest and the arrest lemmas). Every
live arrest $\mathfrak{a}$ of top level
$k_a=k_1<k$ has a live piece $W_{\mathfrak{a}}\subset\Omega^c_{k_1}$; its frozen terms are the chain
terms $\mathbb{C}^{(i)}_{k_1}(X,\mathcal{S})$, $i\le n^+_{\mathfrak{a}}$, stored at step $k_1$ with
the stage norms above, together with the change-of-variables terms of its stages. Live arrests
whose components share a cube are nested, $\mathfrak{a}_1,\mathfrak{a}_2,\dots$, with
$k_{\mathfrak{a}_{t+1}}\ge k_{\mathfrak{a}_t}+N$, with $N\ge1$ the number of preparatory stages; the
components of
non-nested live arrests are disjoint, two live pieces of the same top level $k_1$ being distinct
components of $\Omega^c_{k_1}$, and two such components are at least one $LM\check{R}_{k_1}$-block
apart. For an arrest whose last stage is the
$N(g_{k_1})$-th, the pieces of levels $m<k_1$ may have domains reaching the level-$k_1$ bonds
adjacent to $W_{\mathfrak{a}}$.
\item \emph{Change-of-variables terms} (the decay statement for the change-of-variables terms). Each
stage of a live arrest carries one almost-local family,
with one term per block of the stage's cover on its band $B_n\subset W_{\mathfrak{a}}$; each term has
size at most $1$ and decays at rate at least $a_m=\tfrac45\kappa$ in the length $d_\ell$ of domains at
its own level $\ell\le k_1$.
\item \emph{Stored format of the $\mathbf{T}$-born boundary terms} (the stored boundary-term format
at level $k$). For $j\le k$ and $X\in\mathbf{D}_j$, the $\mathcal{S}$-weighted size of
$\mathbf{B}^{(j)}(X,\cdot)$ in the sense of the stored format is at most
$B_0e^{-\kappa_{\mathrm{out}}d_j(X)}$, pointwise in $X$, in the manner of
\cite[(2.42)]{B-CMP119}; here $B_0$ is the size constant of that format, and
$\kappa_{\mathrm{out}}\ge2\kappa$ for every $j\le k$ by the decay-fraction lemma.
\item \emph{Relative format of the $\mathbf{R}$-born boundary terms} (the induction hypothesis at
level $k$ in the proof of the correlation theorem, Theorem~\ref{5.1:thm-correlation-theorem}, in its
clause for these terms, and the relative-format statements
for boundary terms stored on labels). The $\mathbf{R}$-born boundary terms
$\mathbf{B}'^{(j)}$, $j<k$, and their wrapped $\mathbf{T}$-born descendants (wrapped meaning: carried
in the relative format over live strips) are stored on labels $A\in\mathbf{D}_j$, with one term for
each kind of term, each label $A$ and each $\mathcal{S}$. The fibres of the skeleton map were summed
at birth into the
factor $e^{\lambda_\theta}$ under $\tfrac14\beta_s\kappa\ge\lambda_\theta$. These terms obey
\begin{equation}\label{5.2:eq-exterior-collar-bound-2}
\mathfrak{N}^{A,\mathbb{C}}\le e^{\lambda_\theta}c_1^{\flat}(g_j)\,
e^{-(1+\frac14\beta_s)\kappa d_j(A)}\le B_0e^{-\kappa d_j(A)},
\end{equation}
respectively $\mathfrak{N}^{A,\mathbb{C}}\le B_0e^{-\kappa d_j(A)}$ when stored. Further,
$d_k([A]_k)\le L^{-(k-j)}d_j(A)$, and every label $A\in\mathbf{D}_j$ is connected and has a cube
touching a strip of level $\le j$ or meeting $Z_j^{\boxplus}$. A term meeting the class treated by an
$\mathbf{R}$-operation, or one of whose strips is treated by it, is first-part there and ceases; in
particular a unit one of whose strips is treated by the present step is a first-part unit of that
component, with $K_v=[A_v]_k$.
\end{itemize}
Then the following hold.
\begin{enumerate}
\item[(a)] If $v\in\mathcal{V}^{\mathrm{ext}}$ and its domain $X\in\mathbf{D}_j$ (for a wrapped unit,
its label $A_v\in\mathbf{D}_j$), $j\le k$, is a connected union of level-$j$ cubes one of which lies
in $\Omega_k^c$, then
\begin{equation}\label{5.2:eq-exterior-collar-bound-b}
d_j(X)\ge(2\check{R}_k-1)L^{k-j},
\end{equation}
with $X$ replaced by $A_v$ for a wrapped unit.
\item[(b)] Define the constants $\mathsf{K}^{\mathrm{arr}}$ and $\mathsf{K}^{\flat,\mathrm{ex}}$ by the
last two lines of the following display. Then the two bounds in its first line hold:
\begin{equation}\label{5.2:eq-exterior-collar-bound-a}
\begin{gathered}
\mathsf{K}_{\mathrm{val}}(k)\le\mathsf{K}^0+N(g_k)C_0^{\sharp},\qquad
\mathsf{K}'_{\mathrm{val}}(k)\le\mathsf{K}^{0\prime}+N(g_k)C_0^{\sharp},\\
\mathsf{K}^{\mathrm{arr}}:=\sum_{s\ge1}2C_NC_0^{\sharp}e^{ac_{\mathrm{tr}}}
(2+c_s)^{r_{\mathrm{pr}}}L^{(d-1)s}\,\frac{2}{e\kappa}\,e^{-\frac12\kappa L^s},\\
\mathsf{K}^{\flat,\mathrm{ex}}:=e^{(a+9)c_lc_{\mathfrak{h}}}\,27S_d'\,\mathsf{K}^{\flat}.
\end{gathered}
\end{equation}
Here $N(g_k)\le C_N\mathsf{T}_k^{r_{\mathrm{pr}}}$. The constant $\mathsf{K}^0$ is independent of $g$
and of $k$ and depends on $(d,L,M,\kappa,\delta)$; it is fixed after $M$ in the order of choice of
the auxiliary parameters. It is built from the constants of the size table, of the boundary-term
format and of the stage norms, together with $C_N$, $C_0^{\sharp}$, $c_{\mathrm{tr}}$ and the floors
on $\kappa$ and $\delta$, and it contains $\mathsf{K}^{\mathrm{arr}}$ and $\mathsf{K}^{\flat,\mathrm{ex}}$
as summands. The constant $\mathsf{K}^{0\prime}$ is the $(d,L,M,\kappa,\delta)$-number produced by the
same construction at the weight $a+8$ with the sizes $\mathfrak{N}(v)$ in place of the values.
Moreover, for every cube $\Delta\not\subset Z$, the sum of $e^{ad_k(\mathfrak{K}_v)}\|v\|$ over the
units $v\in\mathcal{V}^{\mathrm{ext}}$ with $\mathfrak{K}_v\supset\Delta$ that are frozen terms of
live arrests (frozen chain terms and change-of-variables terms) is at most $\mathsf{K}^{\mathrm{arr}}$;
the series $\mathsf{K}^{\mathrm{arr}}$ converges and depends on
$(d,L,\kappa,\lambda,c_2,r_{\mathrm{pr}},C_N,C_0^{\sharp},c_{\mathrm{tr}})$ only, and in particular it
is independent of $g$, of $k$ and of the ages $s$.
\end{enumerate}
\end{lemma}

\begin{proof}
\emph{The collar.} The set $Z$ is a class of components of $\Lambda_k^c$. (In \cite{B-CMP122a} the
treated class is written $Z_k$ and identified with $Z$; we write $Z$ throughout.)
Moreover $\Omega_k^c\subset\Lambda_k^c$.
The domains $\Omega_j,\Lambda_j$, which are determined by the $j$-th renormalization transformation
and not by the $\mathbf{R}$-operation, are unions of $MR_j$-cubes, and the distance between their
boundaries is at least $2MR_j$ \cite[pp.~256--257]{B-CMP119}. Between $\partial\Omega_k$ and
$\partial\Lambda_k$ there is therefore a collar of width at least $2\check{R}_k$ current $M$-cubes.

The following sets all lie in $\Omega_k^c$: every zone $n<k$; every older large-field region
$Z_j=\Lambda_j^c$, $j<k$; every live arrested piece $W_{\mathfrak{a}}\subset\Omega^c_{k_a}$; and every
strip $\bigcup Y_i^{\sharp}$ of an older level. Indeed
$\Lambda_j^c\subset\Omega_{j+1}^c\subset\Omega_k^c$ \cite{B-CMP119}.

Let $C$ be the component of $\Lambda_k^c$ containing a cube of $X$ lying in $\Omega_k^c$. If
$C\subset Z$,
a unit of $\mathcal{V}^{\mathrm{ext}}$ must reach a cube $\Delta\not\subset Z$. If $C\not\subset Z$, it
must meet $Z$, since $K_v$ meets $Z$ by the definition of $\mathcal{V}$. In either case the connected
domain leaves $C$ and so crosses the collar. A tree crossing $2\check{R}_k$ level-$k$ cubes has length
at least $2\check{R}_k-1$ in level-$k$ units, that is, at least $(2\check{R}_k-1)L^{k-j}$ in level-$j$
units. This is \eqref{5.2:eq-exterior-collar-bound-b}.

Exteriority and \eqref{5.2:eq-exterior-collar-bound-b} are read on $\mathfrak{K}_v$: a unit is
exterior if and only if its label leaves $Z$. A label $A_v\in\mathbf{D}_m$, $m<k$, is a connected union
of cubes with a cube in $\Omega_k^c$. The reason is that it meets $Z_m^{\boxplus}$ or has a point on a
strip of level $\le m$, and these lie in $\Omega^c_{m+1}\subset\Omega_k^c$. Indeed, a level-$m$ strip
reaches at most $m_\sharp(m)\le\frac12(\check{R}_m-1)$ level-$m$ cubes beyond $\Lambda_m^c$. Moreover
$\operatorname{dist}(\Lambda_m^c,\Omega_{m+1})\ge 8L\check{R}_{m+1}$ \cite{B-CMP119}, and
$8L\check{R}_{m+1}\ge 8\check{R}_m$ since $\check{R}_m\le L\check{R}_{m+1}$.
Hence a label crosses the
collar like any other domain, which proves~(a).

For a plain unit born at level $j<k$ we have $K_v=\bar X$, and part (ii) of the boundary-line lemma
gives $d_k(\bar X)\le L^{-s}d_j(X)+c_{\mathrm{tr}}$ with
$c_{\mathrm{tr}}\le 2d$. Hence, using $L\ge3$,
\begin{equation}\label{5.2:eq-exterior-collar-bound-3}
e^{ad_k(K_v)}\le e^{ac_{\mathrm{tr}}}e^{\frac{31}{20}\kappa L^{-s}d_j(X)}
\le e^{ac_{\mathrm{tr}}}e^{\frac{31}{60}\kappa d_j(X)} .
\end{equation}
This estimate is used in each of the classes below with $s\ge1$.
The sum anchored at a level-$k$ cube $\Delta$ runs, at level $j$, over the $L^{ds}$ level-$j$
$M$-cubes $\square'\subset\Delta$.

\emph{Chain values of the present step.} A chain term of a stage integrating level $j'$ is localised
at some level $m\ge j'+1$: $X\in\mathbb{D}_m$, $X\subset\Omega^c_{m+1}$, $X\cap\Omega_m\neq\emptyset$
and $X\cap(Z_{j'+1}\setminus Z'_{j'+1})\neq\emptyset$ \cite[(1.63)]{B-CMP122a}.

Consider first the pieces filed at $m<k$. They lie in
$\Omega^c_{m+1}\subset\Omega_k^c\subset\Lambda_k^c$, are connected, and meet the region of $Z$. So
they lie inside one component of $\Lambda_k^c$. Either this component belongs to $Z$, and then
$K_v\subset Z$, or it does not, and then $K_v\cap Z=\emptyset$. In both cases the unit is not in
$\mathcal{V}^{\mathrm{ext}}$.

The chain units of the present step in $\mathcal{V}^{\mathrm{ext}}$ are therefore those with $m=k$.
The stage-norm hypothesis delivers them at $a_k=(1+3\beta_{\mathrm{pr}})\kappa=\tfrac85\kappa$, and
$\tfrac85\kappa\ge a+\kappa/20$. For each stage, by that hypothesis,
\begin{equation}\label{5.2:eq-exterior-collar-bound-4}
\sum_{X\ni\Delta}e^{ad_k(X)}\sup|\mathbb{C}^{(n)}_k(X)|\le\|\mathbb{C}^{(n)}_k\|^{A}_{(n)}
\le C_0^{\sharp},
\end{equation}
since the $X$-sum is the norm's supremand and the supremum on the tube is bounded by $\|\cdot\|^{A}$.
There are at most $N(g_k)$ stages. The contribution of this class to $\mathsf{K}_{\mathrm{val}}(k)$ is
therefore at most $N(g_k)C_0^{\sharp}$.

\emph{Chain values left by chains arrested at earlier steps.} Let $k_1=k-s<k$ with $s\ge1$. The chains
arrested at step $k_1$ leave their new terms as part of all boundary terms \cite{B-CMP122a}. By the
hypothesis on frozen terms, these are the frozen chain terms $\mathbb{C}^{(i)}_{k_1}(X,\mathcal{S})$,
$i\le n^+_{\mathfrak{a}}$, of a live piece $W_{\mathfrak{a}}$ with $k_a=k_1$. They are stored with
$\|\cdot\|^{A}\le C_0^{\sharp}$ per stage, at rate $\tfrac85\kappa$ in level-$k_1$ units at the top
level and at rate $\tfrac45\kappa$ in level-$m$ units at levels $m<k_1$.

Pieces filed at $m<k_1$ lie in $\Omega^c_{m+1}\subset\Omega^c_{k_1}\subset\Lambda_k^c$ and are
connected. They therefore lie inside one component of $\Lambda_k^c$, so $K_v\subset Z$ or
$K_v\cap Z=\emptyset$, and they are not in $\mathcal{V}^{\mathrm{ext}}$.

Pieces at the top level $k_1$ meet $Z_{j'+1}\setminus Z'_{j'+1}\subset\Lambda^c_{k_1}\subset\Omega_k^c$,
with $j'$ as above.
If such a piece is in $\mathcal{V}^{\mathrm{ext}}$, it leaves its component of $\Lambda_k^c$, and (a)
gives $d_{k_1}(X)\ge(2\check{R}_k-1)L^s$. By \eqref{5.2:eq-exterior-collar-bound-3}, and since
$\tfrac85-\tfrac{31}{60}=\tfrac{13}{12}$,
\begin{align*}
e^{ad_k(\bar X)}|\mathbb{C}|
&\le e^{ac_{\mathrm{tr}}}e^{\frac{31}{60}\kappa d_{k_1}(X)}|\mathbb{C}|
 = e^{ac_{\mathrm{tr}}}e^{-\frac{13}{12}\kappa d_{k_1}(X)}\,
   e^{\frac85\kappa d_{k_1}(X)}|\mathbb{C}|\\
&\le e^{ac_{\mathrm{tr}}}e^{-\frac{13}{12}\kappa(2\check{R}_k-1)L^s}\,
   e^{\frac85\kappa d_{k_1}(X)}|\mathbb{C}|
 \le e^{ac_{\mathrm{tr}}}e^{-\kappa\check{R}_kL^s}\,e^{\frac85\kappa d_{k_1}(X)}|\mathbb{C}| .
\end{align*}
The last step uses $\tfrac{13}{12}(2\check{R}-1)\ge\check{R}$, which is $14\check{R}\ge13$.

At each of the $L^{ds}$ level-$k_1$ cubes $\square'\subset\Delta$, the $X$-sum of
$e^{\frac85\kappa d_{k_1}(X)}|\mathbb{C}|$ is the norm, so it is at most $C_0^{\sharp}$ per stage.

The stages number at most $N(g_{k_1})\le C_N\mathsf{T}_{k_1}^{r_{\mathrm{pr}}}$. By the reference
recursion (Definition~\ref{5.1:def-reference-recursion}), $x_{k_1}\le x_k+3\lambda s+2c_2$. Since
$x_k\ge1$, this gives
\[
\mathsf{T}_{k_1}=\log x_{k_1}\le\log(x_k+3\lambda s+2c_2)\le\mathsf{T}_k+c_s .
\]

Among the pieces meeting the component of $\Delta$, a nested family of live arrests contains at most
one of top level $k_1$, since $k_{\mathfrak{a}_{t+1}}\ge k_{\mathfrak{a}_t}+N$ with $N\ge1$
by the hypothesis on frozen terms. By the same hypothesis, two non-nested live pieces of the same top
level inside one component of $\Lambda_k^c$ are distinct components of $\Omega^c_{k_1}$; they are at
least one $LM\check{R}_{k_1}$-block apart. A second such piece therefore costs a further factor
$e^{-\frac45\kappa\check{R}_{k_1}}$ and is covered by the factor $2$ below.

Hence the class contributes at most
\begin{equation}\label{5.2:eq-exterior-collar-bound-5}
\sum_{s\ge1}2C_NC_0^{\sharp}e^{ac_{\mathrm{tr}}}(\mathsf{T}_k+c_s)^{r_{\mathrm{pr}}}
L^{ds}e^{-\kappa\check{R}_kL^s}.
\end{equation}

\emph{Absorption.} Suppose $\mathsf{T}_k\ge2$. Then
\[
(\mathsf{T}_k+c_s)^{r_{\mathrm{pr}}}\le(2+c_s)^{r_{\mathrm{pr}}}(\mathsf{T}_k-1)^{r_{\mathrm{pr}}}
\le(2+c_s)^{r_{\mathrm{pr}}}\check{\mathsf{T}}_k^{r_{\mathrm{pr}}}
\le(2+c_s)^{r_{\mathrm{pr}}}\check{R}_k .
\]
The first inequality holds because
$(2+c_s)(\mathsf{T}-1)-(\mathsf{T}+c_s)=(\mathsf{T}-2)(1+c_s)\ge0$. The other two use
$\mathsf{T}_k-1\le\check{\mathsf{T}}_k$ and $\check{R}_k\ge\check{\mathsf{T}}_k^{r_{\mathrm{pr}}}$,
from the standing bounds on $\check{R}_k$ and on $\check{\mathsf{T}}_k$ of this chapter.

Next, since $\check{R}_k\ge1$,
\[
\check{R}_ke^{-\kappa\check{R}_kL^s}
\le e^{-\frac12\kappa L^s}\sup_{R\ge0}Re^{-\frac12\kappa RL^s}
=\frac{2}{e\kappa L^s}\,e^{-\frac12\kappa L^s}.
\]
Inserting both bounds into \eqref{5.2:eq-exterior-collar-bound-5} shows that the class is at most
$\mathsf{K}^{\mathrm{arr}}$ as defined in \eqref{5.2:eq-exterior-collar-bound-a}.

The series $\mathsf{K}^{\mathrm{arr}}$ converges and depends only on $(d,L,\kappa,\lambda,c_2,
r_{\mathrm{pr}},C_N,C_0^{\sharp},c_{\mathrm{tr}})$. It is independent of $g$, of $k$ and of the age:
no bound on $s$ is used, and the factor $L^s$ supplied by the collar does the work. The condition
$\mathsf{T}_k\ge2$ holds because $k\le K$ gives $x_k>g^{-2}$ by first passage
(Definition~\ref{5.1:def-reference-recursion}), hence
$\mathsf{T}_k\ge\mathsf{T}\ge\mathsf{T}_\gamma$, and $\mathsf{T}_\gamma\ge2$ is contained in the
condition $\log g^{-2}>4p_0$ among the smallness conditions on $\gamma$. The constant
$\mathsf{K}^{\mathrm{arr}}$ is a summand of $\mathsf{K}^0$.

Consider the pieces at levels $m<k_1$ of an arrest whose last stage is the $N(g_{k_1})$-th, whose
domains reach the level-$k_1$
bonds adjacent to $W_{\mathfrak{a}}$. They fall within the same count, because they too must cross the
collar to be exterior.

\emph{Change-of-variables terms of the frozen left-overs.} By hypothesis, each stage of each live arrest
carries one almost-local family, with one term per block of the stage's cover on its band
$B_n\subset W_{\mathfrak{a}}\subset\Omega^c_{k_1}$. Ba{\l}aban states that these terms depend almost
locally on the fields and can be bounded by any positive power of the coupling at their own level
\cite{B-CMP119}; we use only
the bound $1$. The pieces $\mathsf{q}_{\mathsf{p}}$ of the Gaussian kept by a live arrest are carried
by the rotation at the localisation site, with size at most
$e^{4\varepsilon_\star}\bar{\mathsf{q}}_{\mathsf{p}}$, $\bar{\mathsf{q}}_{\mathsf{p}}$ being their
stored size; they are not entries of the size table and no bound for them is used in the present
count.

Such a term is exterior only if its domain $X$ leaves the component of $\Lambda_k^c$ containing
$W_{\mathfrak{a}}$. It then crosses the same collar, and at its own level $\ell\le k_1$, by (a),
$d_\ell(X)\ge(2\check{R}_k-1)L^{k-\ell}\ge(2\check{R}_k-1)L^s$. The weight is
$e^{ad_k(\mathfrak{K}_v)}\le e^{ac_{\mathrm{tr}}}e^{\frac{31}{60}\kappa d_\ell(X)}$. Against the decay
at rate at least $\tfrac45\kappa$ this leaves $\tfrac{17}{60}\kappa\,d_\ell(X)$, as for the non-top
chain terms treated below. Of the decay rate only its excess over $\tfrac{31}{60}\kappa$ is used.

We show that these terms enter the series of $\mathsf{K}^{\mathrm{arr}}$ with $C_0^{\sharp}$ replaced
by $C_0^{\sharp}+1$. Fix an arrest of top level $k_1=k-s$ and one of its stages, let $\ell\le k_1$ be
the level of that stage's terms, and put $t:=k-\ell\ge s$. Of the spare
$\tfrac{17}{60}\kappa\,d_\ell(X)$, the part $\tfrac{3}{60}\kappa=\kappa/20$ sums over the domains
through a level-$\ell$ cube $\square'\subset\Delta$ by the count of \cite[(1.26)]{B-CMP116}, at a cost
at most $\sum_{n\ge0}e^{c_d(n+1)}e^{-(c_d+2)n}\le2e^{c_d}$, since $\kappa/20\ge c_d+2$ by the first
member of $\kappa_d$. The remaining part is
$\tfrac{14}{60}\kappa\,d_\ell(X)\ge c_\ast L^t$ with $c_\ast:=\tfrac{14}{60}\kappa(2\check{R}_k-1)$,
by (a).
There are $L^{dt}$ level-$\ell$ cubes $\square'\subset\Delta$. So the stage contributes at most
\[
2e^{c_d}e^{ac_{\mathrm{tr}}}L^{dt}e^{-c_\ast L^t}\le2e^{c_d}e^{ac_{\mathrm{tr}}}L^{ds}e^{-c_\ast L^s}.
\]
For the second inequality put $u:=t-s\ge0$. Since $L\ge3$ we have $L^u-1\ge2u$; since $s\ge1$ gives
$L^s\ge3$, and $c_\ast\ge\tfrac{14}{60}\kappa$, it follows that
\[
c_\ast(L^t-L^s)=c_\ast L^s(L^u-1)\ge\tfrac{14}{60}\kappa\cdot3\cdot2u
=\tfrac75\kappa u\ge14u\log L\ge du\log L
\]
at $d=4$, by the floor $\kappa\ge10\log L$. The stages number at most
$N(g_{k_1})\le C_N(\mathsf{T}_k+c_s)^{r_{\mathrm{pr}}}\le C_N(2+c_s)^{r_{\mathrm{pr}}}\check{R}_k$, by the
absorption above, and a second arrest of the same top level is covered by the factor $2$ as there.
Since the summand of $\mathsf{K}^{\mathrm{arr}}$ is
$2C_NC_0^{\sharp}e^{ac_{\mathrm{tr}}}(2+c_s)^{r_{\mathrm{pr}}}$ times
$L^{(d-1)s}\,\frac{2}{e\kappa}\,e^{-\frac12\kappa L^s}$,
the change-of-variables terms of age $s$ add at most that summand with $C_0^{\sharp}$ replaced by $1$,
provided
\begin{equation}\label{5.2:eq-exterior-collar-bound-6}
2e^{c_d}\check{R}_kL^{ds}e^{-\frac{14}{60}\kappa(2\check{R}_k-1)L^s}
\le L^{(d-1)s}\,\frac{2}{e\kappa}\,e^{-\frac12\kappa L^s}.
\end{equation}
To prove \eqref{5.2:eq-exterior-collar-bound-6}, note first that
\[
\check{R}_k\ge\check{\mathsf{T}}_k^{r_{\mathrm{pr}}}\ge2^{r_{\mathrm{pr}}}\ge4,
\]
because $\check{\mathsf{T}}_k\ge\mathsf{T}_\gamma\ge2$ for $k\le K$ by first passage and
$r_{\mathrm{pr}}\ge2$.
Taking logarithms, \eqref{5.2:eq-exterior-collar-bound-6} is equivalent to
\[
c_d+1+\log\kappa+s\log L\le\tfrac{14}{60}\kappa(2\check{R}_k-1)L^s-\tfrac12\kappa L^s
-\log\check{R}_k .
\]
Since $\kappa\ge1$ by the floor $\kappa_d$, the right member, as a function of $\check{R}_k\ge1$, has
derivative $\tfrac{28}{60}\kappa L^s-1/\check{R}_k>0$; it is therefore at least its value at
$\check{R}_k=4$, namely $\tfrac{68}{60}\kappa L^s-\log4$. On the other side $c_d\le\kappa/20-2$,
$1+\log\kappa\le\kappa$ and $s\log L\le L^s\log L\le\kappa L^s/10$. It therefore suffices that
$\tfrac{21}{20}\kappa-2+\log4\le\tfrac{62}{60}\kappa L^s$, which holds because $L^s\ge3$ makes the
right member at least $\tfrac{31}{10}\kappa$ and $\log4<2$.
This replacement is absorbed by the factor $2$ in $2C_NC_0^{\sharp}$, since
$C_0^{\sharp}\ge1$.

\emph{Boundary terms, non-top arrested chain terms and the remaining lines.} We treat in turn the
$\mathbf{T}$-born terms, the
$\mathbf{R}$-born terms, non-top chain terms of arrests, and the remaining lines of the size table.

\smallskip
\emph{$\mathbf{T}$-born boundary terms $\mathbf{B}^{(j)}$.} By the hypothesis on their stored format,
their bound holds pointwise in $X$: the $\mathcal{S}$-weighted size of $\mathbf{B}^{(j)}(X,\cdot)$ in
the sense of the stored format is at most $B_0e^{-\kappa_{\mathrm{out}}d_j(X)}$.
The sum over $X\ni\square$ is not part of this bound. It is paid out of the spare rate by the count
$\#\{X\in\mathbf{D}_j\ni\square:d_j(X)\le n\}\le e^{c_d(n+1)}$ of \cite[(1.26)]{B-CMP116}.

Case $j=k$. Since $\kappa_{\mathrm{out}}\ge2\kappa$,
\[
\sum_{X\ni\Delta}e^{ad_k(X)}|\mathbf{B}^{(k)}(X)|
\le B_0\sum_{n\ge0}e^{c_d(n+1)}e^{-(2-\frac{31}{20})\kappa n}\le 2e^{c_d}B_0 .
\]
This holds as soon as $\tfrac{9}{20}\kappa\ge c_d+1$, which follows from the first member
$\kappa/20\ge c_d+2$ of the floor $\kappa_d$.

Case $j<k$, $s\ge1$. The domain meets $Z_j=\Lambda_j^c\subset\Omega_k^c$ by \cite[(2.41)]{B-CMP119}.
Since the unit is exterior, (a) gives $d_j(X)\ge(2\check{R}_k-1)L^s$. Per unit of $d_j(X)$ the spare
rate is $(2-\tfrac{31}{60})\kappa=\tfrac{89}{60}\kappa$, by $\kappa_{\mathrm{out}}\ge2\kappa$ and
\eqref{5.2:eq-exterior-collar-bound-3}.
Of this, $\kappa/20$ pays the scale-$j$ count of domains at each $\square'$ by
\cite[(1.26)]{B-CMP116} and $\kappa_d$. The remaining rate gives at least
\[
\tfrac{86}{60}\kappa d_j(X)\ge\tfrac{86}{60}\kappa(2\check{R}_k-1)L^s ,
\]
which pays the $L^{ds}$ cubes $\square'\subset\Delta$ and leaves $e^{-\kappa\check{R}_kL^s}$, because
of the floor $\kappa\ge10\log L$ on $\kappa$. Summing over $s$ gives at most
$2e^{c_d}B_0e^{ac_{\mathrm{tr}}}\sum_{s\ge1}e^{-\kappa L^s}$, which is independent of $g$.

\smallskip
\emph{$\mathbf{R}$-born boundary terms $\mathbf{B}'^{(j)}$, $j<k$, and their wrapped $\mathbf{T}$-born
descendants.} By the relative-format hypothesis, their bound \eqref{5.2:eq-exterior-collar-bound-2} is
taken on labels. It derives from Ba{\l}aban's relative bound \cite[(1.99)]{B-CMP122b}, in which the
decay is measured by $d_{k,\cup Y_i}$.

For such a unit in $\mathcal{V}^{\mathrm{ext}}$, (a) holds of the label:
$d_j(A_v)\ge(2\check{R}_k-1)L^s$. As shown in the proof of (a), the label $A_v\in\mathbf{D}_j$ is
connected and rooted in $\Omega_k^c$: it has a cube touching a strip of level $\le j$ or meeting
$Z_j^{\boxplus}$, and both lie inside $\Omega^c_{j+1}\subset\Omega_k^c$.
The label must leave its component of $\Lambda_k^c$ and so crosses the collar.

The weight is at most
\[
e^{ad_k(\mathfrak{K}_v)}\le e^{\frac{31}{20}\kappa L^{-s}d_j(A_v)}
\le e^{\frac{31}{60}\kappa d_j(A_v)},
\]
using $s\ge1$, $L\ge3$ and $d_k([A]_k)\le L^{-s}d_j(A)$.

We take the spare rate at the stored rate $\kappa$, the weakest one available. The fresh rate
$(1+\tfrac14\beta_s)\kappa=\tfrac{17}{16}\kappa$ would leave $\tfrac{131}{240}\kappa$, which is more.
At rate $\kappa$ the spare is $(1-\tfrac{31}{60})\kappa=\tfrac{29}{60}\kappa$ per unit of $d_j(A_v)$.
It is spent as follows.
\begin{itemize}
\item $\tfrac{3}{60}\kappa=\kappa/20$ sums the labels through a scale-$j$ cube
$\square'\subset\Delta$, by \cite[(1.26)]{B-CMP116} alone:
$\#\{A\in\mathbf{D}_j\ni\square':d_j(A)\le n\}\le e^{c_d(n+1)}$, with $c_d+2\le\kappa/20$ by the first
member of $\kappa_d$. There is one term for each kind of term, each label and each $\mathcal{S}$, and
no fibre count remains, since the fibres of the skeleton map were summed at birth into
$e^{\lambda_\theta}$.
\item The rest gives
\[
\tfrac{26}{60}\kappa d_j(A_v)\ge\tfrac{26}{60}\kappa(2\check{R}_k-1)L^s
\ge2\cdot\tfrac{13}{60}\kappa\check{R}_kL^s .
\]
One half pays the $L^{ds}$ cubes $\square'\subset\Delta$:
\[
\tfrac{13}{60}\cdot10\log L\cdot L^s=\tfrac{13}{6}L^s\log L\ge ds\log L,
\]
which at $d=4$ follows from $L^s\ge\tfrac{24}{13}s$, true since $L\ge3$, with the floor
$\kappa\ge10\log L$. The other half,
$e^{-\frac{13}{60}\kappa\check{R}_kL^s}$, is spare.
\end{itemize}

The prefactor $e^{\lambda_\theta}c_1^{\flat}(g_j)$ is at most $1$. This follows, at every level
$j\le k$, from the smallness conditions on $\gamma$, read with the constant $C_{99}$, which give
$c_1^{\flat}(g_j)\le g_j^{\kappa_0}$, and from the condition $e^{\lambda_\theta}g_j^{\kappa_0}\le1$;
here $\kappa_0$ is the exponent of the coupling in the bounds $g^{\kappa_0}$ of the induction
hypothesis. The sum over $s$ is independent
of $g$ and enters $\mathsf{K}^0$ with the constant
$2e^{c_d}\max\{1,B_0\}\sum_{s\ge1}e^{-\frac{13}{60}\kappa L^s}$.

There is no $\mathbf{B}'^{(k)}$ at the localisation site of step $k$: it is created after that site,
at \cite[(1.90)--(1.99)]{B-CMP122b}. A unit one of whose strips is treated by the present step is a
first-part unit of that component with $K_v=[A_v]_k$; if it is exterior, it is counted among the lines
treated here.

\smallskip
\emph{Arrested chain terms at non-top levels.} By the first part of the treatment of arrested chains,
these are never exterior: $X\subset\Omega^c_{m+1}\subset\Omega^c_{k_1}$ confines them to one component
of $\Lambda_k^c$. Were one exterior, its rate would be $a_m=\tfrac45\kappa$ and the spare
$(\tfrac{48}{60}-\tfrac{31}{60})\kappa=\tfrac{17}{60}\kappa$. Of this, $\tfrac{3}{60}\kappa$ sums over
$X$. The rest gives
\[
\tfrac{14}{60}\kappa(2\check{R}_k-1)L^s\ge2\cdot\tfrac{7}{60}\kappa\check{R}_kL^s .
\]
One half pays $L^{ds}$, since $\tfrac76L^s\log L\ge ds\log L$ when $L^s\ge4s$, which holds for every
$s\ge1$ once $L\ge4$; this is implied by the conditions $\mathsf{L}_0\ge3$ and
$\mathsf{L}_0^2\le L/3$ on the auxiliary number $\mathsf{L}_0$ of \cite[(1.24)]{B-CMP122a}, as fixed
with the auxiliary parameters. The
other half is spare. These terms thus lie inside the series of $\mathsf{K}^{\mathrm{arr}}$. This
completes the proof of the last assertion of (b).

\smallskip
\emph{The remaining lines of the size table.} These are the lines $1$--$3$, $5$ and $\partial$:
terminal terms of the expansion, renormalised groups, tails and boundary pieces, with Ba{\l}aban's
subtracted copies. Part (ii) of the boundary-line lemma sums them at weight
$(1+3\beta_s)\kappa=\tfrac74\kappa\ge a$, and at weight $2\kappa$ when $j=k$ by the decay-fraction
lemma, with the per-cell sizes of the size table.

The kept factor $e^{-\frac14\delta\kappa d_j(X)}$ sums over the anchors $y$ of the units with
$X\ni\Delta$. In zone $k$ there are at most $M^d(2s'+3)^d$ cells at distance $s'$ current cubes from
$\Delta$, against $e^{-\frac14\delta\kappa(s'-1)}$; this sum is independent of $g$.

Anchors in zones $n<k$ lie in $\Omega_k^c$, at least $2\check{R}_k$ current cubes inside their
component of $\Lambda_k^c$. A unit of $\mathcal{V}^{\mathrm{ext}}$ must leave that component, so by (a)
it has $d_j(X)\ge(2\check{R}_k-1)L^s$ in level-$j$ units. The kept factor then splits as
\[
e^{-\frac14\delta\kappa d_j(X)}\le e^{-\frac14\delta\kappa(2\check{R}_k-1)L^s}\,
e^{-\frac14\delta\kappa(d_j(X)-(2\check{R}_k-1)L^s)} .
\]
The number of scale-$j$ anchors within $s'$ further current cubes is bounded by a fixed multiple of
$(ML^s)^d(2\check{R}_k+3+s')^d$. Setting this against the displayed factor and summing over $s'$ and
over $j\le k$, with the scale factors of part (ii) of the boundary-line lemma, leaves a
$(d,L,M,\kappa,\delta)$-number.
This number is uniform in $\check{R}_k\ge1$, because $\sup_{R\ge0}e^{-bR}R^d<\infty$ for every $b>0$,
uniform in the numbers
of stages $N(g_\cdot)$ and of zones $N_0(g_\cdot)$, and uniform in the ages.

\emph{Exposed pointed units.} The exposed pointed units of $\mathcal{V}^{\mathrm{ext}}$ are the short
singles of lines $2$, $3$ and $\partial$ of the size table, and the Wilson pieces astride $\partial Z$.
With $f_v$ the unit's function of the chain parameter, their
values $\mathrm{val}_v=f_v(\mathrm{chain}(\vec0))$ differ from the value $f_v(1)$ entered in the size
table by at most $\mathfrak{N}^{\flat}(v)$ on the cut carrier, by the value-difference bound for
exposed pointed units. Summed per cell by the column bound for
exposed pointed units at $\bar{\mathsf{w}}=3$, and anchored as in the exterior-piece bound, they
contribute the summand $\mathsf{K}^{\flat,\mathrm{ex}}$ of \eqref{5.2:eq-exterior-collar-bound-a} to
$\mathsf{K}^0$. Their weight is taken on $\mathfrak{K}_v=K_v\cup\mathfrak{h}_v$; by
$d_k(\mathfrak{K}_v)\le d_k(K_v)+c_lc_{\mathfrak{h}}$ this enlargement costs the weight the flat factor
$e^{ac_lc_{\mathfrak{h}}}$, which is contained in the factor $e^{(a+9)c_lc_{\mathfrak{h}}}$ of
$\mathsf{K}^{\flat,\mathrm{ex}}$. This is a number fixed after $M$.

\emph{Total.} The contributions of the arrested chains, the change-of-variables terms, the non-chain
lines, the boundary terms and the exposed pointed units add up to at most $\mathsf{K}^0$. This constant
is independent of $g$ and fixed after $M$; it depends on $\kappa$ only through $\kappa\ge\kappa_d$,
$C_\delta(d)/\delta$ and the floors. Together with the bound $N(g_k)C_0^{\sharp}$ for the chain values
of the present step, this proves the first inequality of \eqref{5.2:eq-exterior-collar-bound-a}, with
$N(g_k)\le C_N\mathsf{T}_k^{r_{\mathrm{pr}}}$.

\emph{The weight $a+8$.} The arguments above read only the sizes and the rates of the units. We run
them again at the activity weight $a+8\le a_k=\tfrac85\kappa$ in place of $a$, with the sizes
$\mathfrak{N}(v)$ in place of the values. The rate budgets close at $a+8$:
\begin{itemize}
\item For the chain values of the present step, the $X$-sum is the stage norm at $a_k$.
\item For the arrested chains and for the boundary terms, the weight is at most
$e^{\frac{8}{15}\kappa d_j}$. The spare is $\tfrac{28}{60}=\tfrac{3}{60}+\tfrac{25}{60}$, with
$\tfrac{25}{12}L^s\ge ds$ since $L\ge3$.
\item For the terms read at $a_m=\tfrac45\kappa$, the spare is
$\tfrac{16}{60}=\tfrac{3}{60}+\tfrac{13}{60}$, with $\tfrac{13}{12}L^s\ge4s$ for every $s\ge1$ once
$L\ge4$; this is implied by the conditions $\mathsf{L}_0\ge3$ and $\mathsf{L}_0^2\le L/3$, as fixed
with the auxiliary parameters.
\item The remaining lines are summed at $\tfrac74\kappa$ by part (ii) of the boundary-line lemma for the
sizes of the size table, and the sizes $\mathfrak{N}^{\flat}$ of the exposed pointed lines by the
column bound for exposed pointed units.
\end{itemize}
This gives $\mathsf{K}'_{\mathrm{val}}(k)\le\mathsf{K}^{0\prime}+N(g_k)C_0^{\sharp}$, the second
inequality of \eqref{5.2:eq-exterior-collar-bound-a}.
\end{proof}

\begin{lemma}[The resummation window is non-empty]\label{5.2:lem-resummation-window}
Let $\beta_{\sharp} := 1/20$ be the input margin and $\mathsf{k}_d$ the counting constant of the floor
$\kappa_d$ in \eqref{5.2:eq-resummation-input-f}, and put $C_{\mathrm{fib}} := 2e^{2\mathsf{k}_d}$. Let
$\alpha_{\ast}$ be the least of the finitely many ceilings on $\alpha$, independent of $g$, that are
imposed in this chapter when Ba{\l}aban's resummations \cite[(1.93)--(1.99)]{B-CMP122b} are
carried out; it is a number fixed before $\gamma$. Let $\mathsf{K}^0$ be the constant of the collar bound
\eqref{5.2:eq-exterior-collar-bound-a}, which contains $\mathsf{K}^{\flat,\mathrm{ex}}$, and let
$\mathsf{c}_{\mathbf{C}} \ge 0$ be the constant that bounds the filed $z$-free members of $\mathbf{C}_k$;
let $\varepsilon_{\mathrm{br}}(k)$ be the six-member bound of
\eqref{5.2:eq-resummation-input-b}, and $\mathsf{K}_{\mathrm{val}}(k)$ the value constant of
\eqref{5.2:eq-value-unit-filing-a}. Let $\kappa_0$ be the exponent of the bound $g_k^{\kappa_0}$ in
Proposition~\ref{5.2:prf-resummation-input} (the boundary-term input at the resummation site), and
$\mathsf{T}_\gamma := \log\gamma^{-2}$. For $0 \le k \le K$ and $u \ge 0$ we define
$m_{\mathrm{lo}}(k)$, $m_{\sharp}(k)$, $\alpha^{\sharp}(k)$ and $F_{\sharp}(u)$ by the first four rows
of the following display, and we call its last two rows, which are conditions on $\gamma$, the strip
ceiling and the cube-root ceiling:
\begin{equation}\label{5.2:eq-resummation-window-a}
\begin{aligned}
m_{\mathrm{lo}}(k) &:= 1 + \Bigl\lceil \frac{2}{\beta_{\sharp}\kappa}
  \log\frac{C_{\mathrm{fib}}\bigl(2+\mathsf{K}^0+\mathsf{c}_{\mathbf{C}}+N(g_k)C_0^{\sharp}\bigr)}
  {\alpha_{\ast}} \Bigr\rceil,\\
m_{\sharp}(k) &:= \bigl\lfloor \tfrac12(\check{R}_k-1) \bigr\rfloor,\\
\alpha^{\sharp}(k) &:= C_{\mathrm{fib}}\bigl(\varepsilon_{\mathrm{br}}(k)+\mathsf{K}_{\mathrm{val}}(k)\bigr)
  e^{-\frac12\beta_{\sharp}\kappa(m_{\sharp}(k)-1)},\\
F_{\sharp}(u) &:= C_{\mathrm{fib}}\bigl(2+\mathsf{K}^0+\mathsf{c}_{\mathbf{C}}
  +C_NC_0^{\sharp}(u+1)^{r_{\mathrm{pr}}}\bigr)
  e^{-\frac14\beta_{\sharp}\kappa(u^{r_{\mathrm{pr}}}-4)},\\
&\sup_{u\ge\mathsf{T}_\gamma} F_{\sharp}(u) \le \alpha_{\ast}
  \qquad\text{(the strip ceiling)},\\
&F_{\sharp}(u)^{1/3} \le e^{-\frac{\kappa_0}{2}(u+1)}\quad\text{for all } u\ge\mathsf{T}_\gamma
  \qquad\text{(the cube-root ceiling)}.
\end{aligned}
\end{equation}
Like $N(g_k)$, $N_0(g_k)$ and $\check{R}_k$, the number $m_{\mathrm{lo}}(k)$ is evaluated along the
sequence of effective couplings after $\gamma$ has been fixed, and it enters none of the conditions
fixed before $\gamma$.

Assume the hypotheses of the collar bound for exterior value units
(Lemma~\ref{5.2:lem-exterior-collar-bound}), so that \eqref{5.2:eq-exterior-collar-bound-a} holds at
every level $k \le K$; assume the ceiling $\varepsilon^{\mathrm{rel}}(k) \le 1$ on the relative member
of \eqref{5.2:eq-resummation-input-b}, provided by the bound for the wrapped units; and assume that
$\gamma$ satisfies the conditions defining $\gamma_J$, among them
the strip ceiling and the cube-root ceiling. Then for every $0 \le k \le K$:
\begin{enumerate}
\item[(i)] the strip smallness obeys
\begin{equation*}
\alpha^{\sharp}(k) \le C_{\mathrm{fib}}\bigl(2+\mathsf{K}^0+\mathsf{c}_{\mathbf{C}}+N(g_k)C_0^{\sharp}\bigr)
  e^{-\frac14\beta_{\sharp}\kappa(\check{R}_k-4)}
\le F_{\sharp}(\check{\mathsf{T}}_k) \le \alpha_{\ast};
\end{equation*}
\item[(ii)] $m_{\mathrm{lo}}(k) \le m_{\sharp}(k)$, that is, the window of admissible strip widths is
non-empty, and $2m_{\mathrm{lo}}(k)+1 \le \check{R}_k$;
\item[(iii)] $\alpha^{\sharp}(k) \le \sup_{u\ge\mathsf{T}_\gamma}F_{\sharp}(u)$, and
this supremum tends to zero as $\gamma \to 0$; thus $\alpha^{\sharp}(k) \to 0$ as $\gamma\to 0$
uniformly in $k$;
\item[(iv)] $\alpha^{\sharp}(k)^{1/3} \le g_k^{\kappa_0}$;
\item[(v)] with the width $m_{\sharp}(k)$, the strips $C^{\sharp}$ about the components $C$ of $Z$
are pairwise non-touching boxes contained in $Z^{\boxplus}$, the enlargement of $Z$ by one layer of
$M\check{R}_k$-cubes, each contained in one component of $\Lambda_k^c$, and
\begin{equation*}
\#\{\text{level-$k$ cubes in } C^{\sharp}\}
\le \min\bigl\{2^d\,\#\{\text{level-$k$ cubes in } C\},\,(101\check{R}_k-1)^d\bigr\};
\end{equation*}
the filing onto the strips is linear and moves localisation domains but never anchors, so the
volume bound at the strips holds cell by cell as before the filing, and neither the pairings of terms
nor the coupling dials are affected. Ba{\l}aban's localisation domains $Y_1$ are kept as they are; no
hull of them is taken.
\end{enumerate}
\end{lemma}

\begin{proof}
Fix $0 \le k \le K$. We first collect the inputs.

By definition $\check{R}_k = R(\min_{i\le k}x_i)$ is the least power of $L$ that is at least
$(\log\min_{i\le k}x_i)^{r_{\mathrm{pr}}}$, and $\check{\mathsf{T}}_k = \log\min_{i\le k}x_i$; hence
$\check{R}_k \ge \check{\mathsf{T}}_k^{\,r_{\mathrm{pr}}}$, and $\check{R}_k$ is an integer. By the
first-passage bound, $\check{\mathsf{T}}_k \ge \mathsf{T}$ for $k \le K$, where
$\mathsf{T} = \log g^{-2}$, and $\mathsf{T} \ge \mathsf{T}_\gamma$; hence
$\check{\mathsf{T}}_k \ge \mathsf{T}_\gamma$. By the
comparison with the reference recursion, $x_i \ge x_k - 2c_2$ for $i \le k$, so
$\min_{i\le k}x_i \ge x_k - 2c_2$; since $x_k \ge 4c_2$ (one of the conditions defining $\gamma_J$),
$x_k - 2c_2 \ge x_k/2$ and $\log(x_k-2c_2) \ge \log x_k - \log 2 \ge \log x_k - 1$. As
$\mathsf{T}_k = \log x_k$, this gives
\begin{equation}\label{5.2:eq-resummation-window-1}
\mathsf{T}_k \le \check{\mathsf{T}}_k + 1 .
\end{equation}
The number of preparatory stages satisfies $N(g_k) \le C_N\mathsf{T}_k^{r_{\mathrm{pr}}}$, hence, by
\eqref{5.2:eq-resummation-window-1},
$N(g_k) \le C_N(\check{\mathsf{T}}_k+1)^{r_{\mathrm{pr}}}$.

Next, $\varepsilon_{\mathrm{br}}(k) \le 2$. Indeed, of the six members in
\eqref{5.2:eq-resummation-input-b}, the three members carrying $\theta'(g_k)$ together with
$\varepsilon^{\mathrm{ext},\mathsf{t}}(k)$ are at most $\tfrac12$ by the ceiling on the
$\theta'$-carrying members and on $\varepsilon^{\mathrm{ext},\mathsf{t}}$ in the list of conditions
defining $\gamma_J$, $\varepsilon^{\mathrm{rel}}(k) \le 1$ by the ceiling on the relative member
assumed in the statement, and
$\varepsilon^{\mathrm{int}}(k)+\varepsilon^{\mathrm{ext},\sigma}(k) \le \tfrac12$ by the ceiling in
\eqref{5.2:eq-plain-exterior-sum-a}, which is also among the conditions defining $\gamma_J$; the
ceilings defining $\gamma_J$ are
suprema over $u \ge \mathsf{T}_\gamma$ and apply at $k \le K$ because there $\mathsf{T}_k \ge \mathsf{T}$
and $\mathsf{T} \ge \mathsf{T}_\gamma$, by the first-passage bound. By the collar bound for exterior
value units (Lemma~\ref{5.2:lem-exterior-collar-bound}), in the form
\eqref{5.2:eq-exterior-collar-bound-a}, we have
$\mathsf{K}_{\mathrm{val}}(k) \le \mathsf{K}^0 + N(g_k)C_0^{\sharp}$. Together,
\begin{equation}\label{5.2:eq-resummation-window-2}
\varepsilon_{\mathrm{br}}(k)+\mathsf{K}_{\mathrm{val}}(k) \le 2+\mathsf{K}^0+N(g_k)C_0^{\sharp}.
\end{equation}

Finally, since $\check{R}_k$ is an integer, $\tfrac12(\check{R}_k-1)$ is an integer or a half-integer,
so its integer part is at least $\tfrac12(\check{R}_k-1)-\tfrac12$; hence
\begin{equation}\label{5.2:eq-resummation-window-3}
m_{\sharp}(k)-1 = \bigl\lfloor\tfrac12(\check{R}_k-1)\bigr\rfloor - 1 \ge \tfrac12(\check{R}_k-4).
\end{equation}

(i) From the definition of $\alpha^{\sharp}(k)$, \eqref{5.2:eq-resummation-window-2},
\eqref{5.2:eq-resummation-window-3} and $\mathsf{c}_{\mathbf{C}} \ge 0$ we obtain
\begin{equation*}
\alpha^{\sharp}(k) \le C_{\mathrm{fib}}\bigl(2+\mathsf{K}^0+\mathsf{c}_{\mathbf{C}}+N(g_k)C_0^{\sharp}\bigr)
  e^{-\frac14\beta_{\sharp}\kappa(\check{R}_k-4)} .
\end{equation*}
In the prefactor we use $N(g_k) \le C_N(\check{\mathsf{T}}_k+1)^{r_{\mathrm{pr}}}$, and in the
exponent $\check{R}_k \ge \check{\mathsf{T}}_k^{\,r_{\mathrm{pr}}}$; the right member is therefore at
most $F_{\sharp}(\check{\mathsf{T}}_k)$. Since $\check{\mathsf{T}}_k \ge \mathsf{T}_\gamma$,
the strip ceiling gives $F_{\sharp}(\check{\mathsf{T}}_k) \le \alpha_{\ast}$. This is (i).

(ii) Read backwards, (i) states that
\begin{equation*}
C_{\mathrm{fib}}\bigl(2+\mathsf{K}^0+\mathsf{c}_{\mathbf{C}}+N(g_k)C_0^{\sharp}\bigr)
e^{-\frac14\beta_{\sharp}\kappa(\check{R}_k-4)}
\le\alpha_{\ast};
\end{equation*}
taking logarithms and using \eqref{5.2:eq-resummation-window-3},
\begin{equation*}
\tfrac12\beta_{\sharp}\kappa\bigl(m_{\sharp}(k)-1\bigr)
\ge \tfrac14\beta_{\sharp}\kappa(\check{R}_k-4)
\ge \log\frac{C_{\mathrm{fib}}\bigl(2+\mathsf{K}^0+\mathsf{c}_{\mathbf{C}}+N(g_k)C_0^{\sharp}\bigr)}
{\alpha_{\ast}} .
\end{equation*}
Thus $m_{\sharp}(k)-1$ is at least $2(\beta_{\sharp}\kappa)^{-1}$ times the logarithm in the
definition of $m_{\mathrm{lo}}(k)$; being an integer, it is at least the least integer not smaller
than that number, namely $m_{\mathrm{lo}}(k)-1$. So
$m_{\mathrm{lo}}(k) \le m_{\sharp}(k)$, and the window is non-empty at every
$k \le K$. Moreover $m_{\sharp}(k) \le \tfrac12(\check{R}_k-1)$, whence
$2m_{\mathrm{lo}}(k)+1 \le 2m_{\sharp}(k)+1 \le \check{R}_k$.

(iii) By (i) and $\check{\mathsf{T}}_k \ge \mathsf{T}_\gamma$,
$\alpha^{\sharp}(k) \le F_{\sharp}(\check{\mathsf{T}}_k) \le \sup_{u\ge\mathsf{T}_\gamma}F_{\sharp}(u)$,
a bound independent of $k$. The function $F_{\sharp}$ is $e^{\beta_{\sharp}\kappa}$ times a polynomial
in $u$ times $e^{-\frac14\beta_{\sharp}\kappa u^{r_{\mathrm{pr}}}}$; it is bounded on
$[0,\infty)$ and tends to zero as $u\to\infty$, so the supremum over
$u \ge \mathsf{T}_\gamma$ is finite and tends to zero as $\mathsf{T}_\gamma \to \infty$, that
is, as $\gamma \to 0$. In particular the strip ceiling holds for all $\gamma$ below a threshold
depending only on $d, L, M, \kappa, \delta, \lambda, c_2, r_{\mathrm{pr}}, C_N, C_0^{\sharp},
c_{\mathrm{tr}}, \alpha_{\ast}$ and on the constants entering $\mathsf{K}^0$ and
$\mathsf{c}_{\mathbf{C}}$, all of which are fixed before
$\gamma$. This agrees with the statement in \cite[p.~375]{B-CMP122b} that $\alpha$ can be made
arbitrarily small for $\gamma$ small enough. Since $r_{\mathrm{pr}} \ge 2$, the quantity
$F_{\sharp}(u)^{1/3}e^{\frac{\kappa_0}{2}(u+1)}$ tends to zero as $u \to \infty$, so the cube-root
ceiling, like the strip ceiling, holds for all $\gamma$ below a threshold.

(iv) Applying (i), then the cube-root ceiling at $u := \check{\mathsf{T}}_k \ge \mathsf{T}_\gamma$, and
finally \eqref{5.2:eq-resummation-window-1}, we find
\begin{equation*}
\alpha^{\sharp}(k)^{1/3} \le F_{\sharp}(\check{\mathsf{T}}_k)^{1/3}
\le e^{-\frac{\kappa_0}{2}(\check{\mathsf{T}}_k+1)}
\le e^{-\frac{\kappa_0}{2}\mathsf{T}_k} = g_k^{\kappa_0},
\end{equation*}
the last equality because $\mathsf{T}_k = \log g_k^{-2}$.

(v) The geometry of the strips and the bound on the number of their cubes are the conclusions of
the strip lemma read at the width $m_{\sharp}(k)$; the statements that the volume bound at the strips
holds cell by cell and that the pairings of terms and the coupling dials are unaffected are the
conclusions of the strip corollary, the filing onto the strips being linear and moving localisation
domains but never anchors.
\end{proof}

\begin{lemma}[Anchors lie in the current label]\label{5.2:lem-current-label-anchors}
Let $k$ be the current step. Throughout, $Z_j:=\Lambda_j^c$, $Z_j^{\boxplus}$ and $\Omega_j$ denote the
unprimed sets of the \emph{current label}, that is, the regions $\Lambda_j$, $\Omega_j$ as they stand at
the current time. For a union $Y$ of cubes of $\pi_j$, $Y^{\boxplus}$ is $Y$ enlarged by one layer of
$M\check{R}_j$-cubes and $Y^{(1)}$ is $Y$ enlarged by the cubes of $\pi_j$ adjacent to it, as in the
format of stored boundary-type terms. The reading of $Z_j$, $Z_j^{\boxplus}$, $\Omega_j$ in the current
label applies wherever these sets occur:
\begin{itemize}
\item in the two anchoring conditions of the format of stored boundary-type terms;
\item in the anchoring of stored terms, in the collar geometry of exterior value units
\eqref{5.2:eq-exterior-collar-bound-b}, in the bound on the collar value constant
$\mathsf{K}_{\mathrm{val}}(k)$ derived from it, and in the geometry of plain units;
\item in the volume charge of a boundary-type term $F$ at a cube of $\pi_j$ of its domain label $A_F$
lying in $(Z_j^{\boxplus})^{(1)}$.
\end{itemize}
A stored term is
\emph{plain} if it is not carried over a live strip.
The following hold, all sets being those of the current label:
\begin{equation}\label{5.2:eq-current-label-anchors-a}
\begin{gathered}
\text{(i)}\quad (D^{\boxplus})^{(1)}\subset Z_{j+1},\\
\text{(ii)}\quad \text{there is a cube $Q_F$ of $\pi_j$ with }
Q_F\subset A_F\cap(Z_j^{\boxplus})^{(1)}.
\end{gathered}
\end{equation}
Here (i) holds for every $j<k$ and every component $D$ of $Z_j$, and (ii) holds for every stored
boundary-type term $F$ of level $j\le k$ alive at the current time; if such an $F$ is plain, moreover
$A_F$ shares a cube with $Z_j^{\boxplus}$.

The inclusion $(Z_j^{\boxplus})^{(1)}\subset Z_{j+1}$, with $Z_j$ read in a term's label at birth and
$Z_{j+1}$ in the current label, is false across a component healed in between. It is never read in
that way here.
\end{lemma}

\begin{proof}
\emph{Part (i).} Let $D$ be a component of the current $Z_j$, $j<k$. Then $D$ is large-field at
level $j$ in the current label.

The step $\mathbf{T}_{j+1}$ made $Z_{j+1}\supseteq D$ with at least ten layers of
$M\check{R}_{j+1}$-cubes about $D$: ``Each renormalization step adds at least ten layers of
$MR_k$-cubes'' \cite[p.~177]{B-CMP122a}. An $M\check{R}_{j+1}$-cube has side $L\check{R}_{j+1}$
cubes of $\pi_j$. By the comparison $\check{R}_j\le L\check{R}_{j+1}$ of consecutive radii, these ten
layers amount to at least
\begin{equation*}
10L\check{R}_{j+1}\ \ge\ 10\check{R}_j\ \ge\ \check{R}_j+1
\end{equation*}
layers of cubes of $\pi_j$. Now $D^{\boxplus}$ adds to $D$ one layer of $M\check{R}_j$-cubes, that is,
$\check{R}_j$ layers of cubes of $\pi_j$, and the operation $(\cdot)^{(1)}$ adds one further layer.
Hence $(D^{\boxplus})^{(1)}$ lies inside the layers that $\mathbf{T}_{j+1}$ added, so it lies in
$Z_{j+1}$.

It remains to see that no later operation undoes this while $D$ is current. Suppose a later completed
$\mathbf{R}$-operation took that part of $Z_{j+1}$ away. It then re-set the whole component to one
level: ``We define $\mathbf{B}_k$ as equal to $\mathbf{B}''_k$ on $Z^c$, and to $\{Z^{(k)}\}$ on
$Z$'' \cite[(1.33)]{B-CMP122b}. That re-setting takes $D$ with it, and then $D$ is not current,
contrary to the hypothesis. A preparation re-makes $\Omega''_j$, $Z''_j$ inside its own component
\cite[pp.~178--179]{B-CMP122a}, and the unprimed $\Lambda_j$ stay as they are. This proves (i).

\emph{Part (ii).} We argue by induction over the completed $\mathbf{R}$-operations since the birth
of $F$.

\emph{Birth.} We distinguish according to the kind of term.
\begin{itemize}
\item A boundary term $\mathbf{B}^{(j)}(X)$ born under $\mathbf{T}$ has $X\not\subset\Lambda_j$. Thus
$X$ shares a cube with $Z_j$ itself; the $E$- and $R''$-terms created together with it are those
whose domain lies in $\Lambda_j$, as in the description of the $E$- and $R''$-terms of the small-field
step in this chapter (written $\Lambda_{k+1}$ in \cite[(3.26), (3.28)]{B-CMP119}, where
the created level is $k+1$).
\item A term $\mathbf{B}'^{(j)}$ born under $\mathbf{R}$ meets ``the large field region
$Z_k^{\sim}\cup\bigcup Y_i^{\sim}$'' \cite[p.~390]{B-CMP122b} of the new label; there $k$ denotes
the step of birth, here $j$, the $Y_i$ are the regions that this operation adds to the large field
region, and the enlargement $\sim$ is read as $\boxplus$.
\item A chain term, or a frozen chain term, of level $m$ meets the basic localization region of its
stage, a subset of $Z_{j'+1}$ written in \cite[(1.63)]{B-CMP122a} as the difference of $Z_{j'+1}$ and
a primed set; here $j'$ is the level integrated at that stage and $m\ge j'+1$. A change-of-variables
term of a frozen family meets the band of its stage.
\item A term carried over live strips touches a live strip, as the format of such terms states
directly after its anchoring condition.
\end{itemize}

\emph{Survival of a completed $\mathbf{R}_{k'}$, $k'\ge j$, with treated class $Z'$.} The class $Z'$
is the class of components of $\Lambda_{k'}^c$ on which $\mathbf{R}_{k'}$ operates.

A term that survives $\mathbf{R}_{k'}$ belongs there to the second part, so $A_F\cap Z'=\emptyset$:
``the second part includes the remaining terms, i.e., the terms with localization domains disjoint
with $Z$'' \cite[p.~369]{B-CMP122b}. A term that meets $Z'$ belongs to the first part and ceases.
So does a term one of whose large-field regions in $\mathcal{P}_F$ is healed by $\mathbf{R}_{k'}$;
here $\mathcal{P}_F$ is the set of live large-field regions over whose strips $F$ is carried. Both
are instances of the ceasing clause of the format of terms carried over live strips.

By the induction hypothesis, in the label before $\mathbf{R}_{k'}$ there is a cube $Q_F$ of $\pi_j$
with $Q_F\subset A_F\cap(D^{\boxplus})^{(1)}$ for some component $D$ of the then-current $Z_j$; call
$Q_F$ the anchor of $F$. There are three cases.

\emph{Case $D\not\subset Z'$.} Then $D$ stays current: ``The new large field region is equal to
$Z_k\cup\bigcup Y_i$'' \cite[p.~390]{B-CMP122b}. The anchor $Q_F$ remains an anchor in the new label.

\emph{Case $D\subset Z'$ and $j<k'$.} By (i), applied in the label before $\mathbf{R}_{k'}$, and
since $j+1\le k'$, we have
\begin{equation*}
(D^{\boxplus})^{(1)}\subset Z_{j+1}=\Lambda_{j+1}^c\subset\Lambda_{k'}^c .
\end{equation*}
The set $(D^{\boxplus})^{(1)}$ is connected and contains $D$. It therefore lies in the component of
$\Lambda_{k'}^c$ that holds $D$, and this component belongs to $Z'$ because $D\subset Z'$. Since
$Q_F\subset A_F\cap(D^{\boxplus})^{(1)}$, this gives
\begin{equation*}
Q_F\subset A_F\cap Z'=\emptyset,
\end{equation*}
which is impossible.

\emph{Case $D\subset Z'$ and $j=k'$.} Then $F$ was born under $\mathbf{T}_{k'}$, the other kinds being
excluded: the outputs of $\mathbf{R}_{k'}$ are born after it; the chain terms of step $k'$ are
first-part and cease; the component $D$ of a frozen term is that of its arrest, which is not in $Z'$.
Hence, by the birth case for boundary terms born under $\mathbf{T}$, $A_F$ shares a cube with $Z_{k'}$
itself.
That cube lies in a component $D'$ of $Z_{k'}$ with $D'\not\subset Z'$, because
$A_F\cap Z'=\emptyset$. The anchor is taken there, and $D'$ stays current by the first case.

This completes the induction. In particular, no stored term is anchored at a component that has
since been healed.
\end{proof}

\begin{proof}[Remnants after resummation and new boundary terms]
\label{5.2:prf-resummation-remnants}
This part of the proof of the $\mathbf{R}$-step at complex source establishes its clause (e). The
clauses proved before it are taken as established: clause (c), the bound \cite[(1.68)]{B-CMP122b} and
the volume bound that follows it there, both at complex source, including the pre-filing family of the
boundary-term input at the resummation site, established in the part of this proof headed
\emph{The boundary-term input at the resummation site} (p.~\pageref{5.2:prf-resummation-input}); and
clause (d), the
modulus bound for the sourced large-field operation (Lemma~\ref{5.2:lem-sourced-modulus}). The argument
also takes the following statements as hypotheses:
\begin{itemize}
\item clause (b) of the corollary on the positivity of the cascade charges, which gives
\cite[(1.89)]{B-CMP122b} on the family of large-field data on which the second assertion of
Lemma~\ref{5.2:lem-sourced-modulus} is stated;
\item clauses (u1)--(u7) of the combinatorial lemma on strips;
\item clauses (b) and (c) of the skeleton lemma;
\item clause (c) of the relative-format proposition for $\mathbf{R}$-born boundary terms;
\item the bound on the pre-filing family, applied to $\mathfrak{v}+\mathfrak{v}^{\mathrm{rel}}$;
\item the inequality $\tfrac14\beta_s\kappa\ge\lambda_\theta$.
\end{itemize}

\emph{The bound in place of} \cite[(1.92)]{B-CMP122b}. Ba{\l}aban introduces that estimate with the
words ``Using (1.68), (1.73), (1.89), we obtain''. Here the three inputs are replaced as follows. The
first is replaced by clause (c). The second is replaced by clause (d), that is, by
Lemma~\ref{5.2:lem-sourced-modulus}. The third is replaced by clause (b) of the corollary on the
positivity of the cascade charges; inserted into Lemma~\ref{5.2:lem-sourced-modulus}, it gives
\eqref{5.2:eq-sourced-modulus-a}. The factor $e$ in front of \eqref{5.2:eq-sourced-modulus-a}, absent
in \cite{B-CMP122b} because $|\sigma|$ is bounded there by $1$, is paid here by a summand $2$ in the
constants of the smallness conditions under which the steps \cite[(1.93)--(1.99)]{B-CMP122b} are
carried out.

\emph{The combinatorics of} \cite[(1.93)--(1.99)]{B-CMP122b}. These steps are combinatorics in $\alpha$
and $c_1$, and they do not involve the source variables $z$. In clause (c) the strips $Z^\sharp$
replace the enlargements of Ba{\l}aban's construction. Accordingly, these steps are run with the
dictionary
\begin{equation}\label{5.2:eq-resummation-remnants-1}
\begin{gathered}
Z^{\sim},\ X^{\sim},\ Y_i^{\sim}\ \longmapsto\ Z^\sharp,\ X^\sharp,\ Y_i^\sharp,\qquad
d_{k,Z}\ \longmapsto\ d_\sharp := d_{k,Z^\sharp},\\
\beta/(3\cdot 2^d)\ \longmapsto\ \beta_v := \beta/(2\mathsf{c}_d).
\end{gathered}
\end{equation}
The term $2\kappa(100\check{R}_k)^d$ displayed in \cite{B-CMP122b} still absorbs the strip, that is, it
still dominates the contribution of the strip $Z^\sharp$ that replaces the enlargement $Z^{\sim}$. The
one-layer enlargement $\boxplus$ is kept in three places: in \cite[(2.30), (2.41)]{B-CMP119}, and in
the definition of the first group in \cite{B-CMP122b}. Read through \eqref{5.2:eq-resummation-remnants-1},
the conclusions of \cite[(1.93)--(1.99)]{B-CMP122b} are clauses (u1)--(u7) of the combinatorial lemma on
strips. Those clauses hold under finitely many smallness conditions on $\alpha$, free of $g$. The least
of these is the number $\alpha_\ast$ which bounds $\alpha^\sharp(k)$ in
\eqref{5.2:eq-resummation-window-a}.

\emph{The second group.} It follows that $\mathbf{R}'^{(k)}(X,(\mathbf{U},\mathbf{J});w)$ is analytic on
the product of the tube over $X$ with $\mathcal{W}_k(X)$, and local in $w$. For the terms of the second
group it satisfies
\begin{equation}\label{5.2:eq-resummation-remnants-2}
\bigl|\mathbf{R}'^{(k)}(X,(\mathbf{U},\mathbf{J});w)\bigr|
\le e^{-p_0(g_k)}e^{-\kappa d_k(X)}
\le \tfrac12\, g_k^{\kappa_0}e^{-\kappa d_k(X)} .
\end{equation}
The second inequality rests on the smallness condition on $e^{-p_0(g_k)}$ recalled below, in the
paragraph on the prefactor.
Here the degree function $p_0(\cdot)$ enters through the sentence of \cite{B-CMP122b}: ``By their
definition the linear size of the domain $X$ in (1.99) can be replaced by $d_k(X)$, and $c_1$ is equal
to $\exp(-p_0(g_k))$''. That sentence is said of the second group only.

\emph{The first group.} In \cite{B-CMP122b} the first group consists of the terms whose localisation
domains $X$ meet $Z_k^{\sim}\cup\bigcup_i Y_i^{\sim}$: ``To the first group we assign all the terms
with the localization domains $X$ intersecting the large field region \dots{} The terms of the first
group are new boundary terms, and they are denoted by $\mathbf{B}'^{(k)}(X)$''. Here $\sim$ denotes the
one-layer enlargement $\boxplus$ and is not replaced by $\sharp$. Denote these terms by
$\mathbf{B}'^{(k)}_0(X,\mathcal{S})$. For them Ba{\l}aban proves the relative bound
\cite[(1.99), p.~390]{B-CMP122b}:
\begin{equation}\label{5.2:eq-resummation-remnants-3}
\bigl|\mathbf{R}'^{(k)}(X,(\mathbf{U},\mathbf{J}))\bigr|
\le O(1)\,c_1\exp\bigl(-(1+\tfrac12\beta)\kappa\, d_{k,\cup Y_i}(X)\bigr).
\end{equation}
We restate this bound in the notation of this book. Let $\mathsf{I} := \{Y_i^\sharp\}$ be the set of
strips born at $\mathbf{R}_k$. For $X\in\mathbf{D}_k^{\mathsf{I}}$ put $A(X) := \operatorname{sk}(X)$, and
put $A(X) := X$ if $X$ meets no $Y_i^\sharp$. For a fixed label $A$ set
\begin{equation}\label{5.2:eq-resummation-remnants-4}
\mathbf{B}'^{(k)}(A,\mathcal{S}) := \sum_{X:\,A(X)=A}\mathbf{B}'^{(k)}_0(X,\mathcal{S}).
\end{equation}
Let $\mathcal{P}_A$ be the set of pairs $(Y_i,k)$ with $Y_i^\sharp$ touching $A$, together with the older
pairs $(Y,m)$ still live after $\mathbf{R}_k$ whose strips $Y^\sharp$ touch $A$. Then
\begin{equation}\label{5.2:eq-resummation-remnants-a}
\mathfrak{N}^{A,\mathbb{C}}\bigl(\mathbf{B}'^{(k)}(A,\cdot\,;w)\bigr)
\le e^{\lambda_\theta}\,c_1^{\flat}(g_k)\,e^{-(1+\frac14\beta_s)\kappa d_k(A)} .
\end{equation}
Moreover, $\mathbf{B}'^{(k)}(A,\cdot\,;w)$ is analytic on the product of the tube over
$X_A := A\cup\bigcup_{\mathcal{P}_A}Y^\sharp$ with $\mathcal{W}_k$, and it is in the relative format
(f0)--(f4) of the boundary terms. The definition \eqref{5.2:eq-resummation-remnants-4}, the set
$\mathcal{P}_A$ and the format clauses are those of clause (c) of the relative-format proposition; the
bound \eqref{5.2:eq-resummation-remnants-a} is obtained here, as in that clause, from clause (u7) of the
combinatorial lemma on strips, clauses (b) and (c) of the skeleton lemma and the inequality
$\frac14\beta_s\kappa\ge\lambda_\theta$, in three steps.

First, clause (u7) of the combinatorial lemma on strips assigns the first group to $\mathbf{B}'$, with
its size relative to the collar in every reading. Clause (u7) rests on
\eqref{5.2:eq-resummation-remnants-3}, and the lemma is applied given the bound on the pre-filing family
for $\mathfrak{v}+\mathfrak{v}^{\mathrm{rel}}$.

Second, clauses (u3)--(u7) carry a symbolic $\beta\le\frac14$. We run them at $\beta := \beta_s$, since
their input is the rate $(1+2\beta_s)\kappa$ of $V^\sharp$ in clause (c). They then deliver the rate
$(1+\frac12\beta_s)\kappa$ in the length $d_{\mathrm{out}} := d_{k,\cup_i Y_i^\sharp}$. Combined with
clauses (b) and (c) of the skeleton lemma, this gives
\begin{equation}\label{5.2:eq-resummation-remnants-5}
\mathfrak{N}^{A,\mathbb{C}}\bigl(\mathbf{B}'^{(k)}(A,\cdot\,;w)\bigr)
\le e^{\lambda_\theta(d_k(A)+1)}\,c_1^{\flat}(g_k)\,e^{-(1+\frac12\beta_s)\kappa d_k(A)} .
\end{equation}

Third, the exponents in \eqref{5.2:eq-resummation-remnants-5} can be rewritten exactly:
\begin{equation*}
\lambda_\theta(d_k(A)+1)-(1+\tfrac12\beta_s)\kappa d_k(A)
=\lambda_\theta-(1+\tfrac14\beta_s)\kappa d_k(A)
-\bigl(\tfrac14\beta_s\kappa-\lambda_\theta\bigr)d_k(A).
\end{equation*}
Since $\frac14\beta_s\kappa\ge\lambda_\theta$, the last term is not positive, and
\eqref{5.2:eq-resummation-remnants-5} implies \eqref{5.2:eq-resummation-remnants-a}.

\emph{Clause (f2) of the format.} Among the smallness conditions whose least member is $\alpha_\ast$
there is one of the form $B_0/O(1)$, with $O(1)$ the constant of
\eqref{5.2:eq-resummation-remnants-3}; it is imposed here in the form $B_0e^{-\lambda_\theta}/C_{99}$.
With this reading the right member of \eqref{5.2:eq-resummation-remnants-a} is at most
$B_0e^{-\kappa d_k(A)}$, which is clause (f2).

\emph{The prefactor.} The prefactor is $c_1^\flat(g_k) = C_{99}c_1$. Where $X$ meets
$\bigcup_iY_i^\sharp$, we have $c_1^\flat = C_{99}(\alpha^\sharp)^{1/3}$. This is at most
$g_k^{\kappa_0}$ by the condition on $(\alpha^\sharp)^{1/3}$ in \eqref{5.2:eq-resummation-window-a},
imposed in the form $C_{99}(\alpha^\sharp)^{1/3}\le g_k^{\kappa_0}$. Otherwise
$c_1^\flat = C_{99}e^{-p_0(g_k)}$, and this is at most
$\frac12 g_k^{\kappa_0}$ by the ceiling $C_{99}e^{-p_0(g_k)}\le\frac12 g_k^{\kappa_0}$. In both cases
the bound \eqref{5.2:eq-resummation-remnants-a} carries the constant of
\eqref{5.2:eq-resummation-remnants-3}.

\emph{What is and is not asserted.} The bound \eqref{5.2:eq-resummation-remnants-a} is the form in which
the inductive hypothesis at complex source carries the terms $\mathbf{B}'$, and the inductive step
closes with it. Consider a $\mathbf{B}'$ whose true domain contains a strip. For such a term the
plain-length form $B_0e^{-\kappa d_k(X_A)}$ is not asserted anywhere in this book. That the inductive
assumptions may be carried in this relative form is stated in \cite[p.~262]{B-CMP119}: ``There are
other possible forms of the
inductive assumptions for the boundary terms \dots{} (2.42) with $d_j(X)$ replaced by
$d_j(X\setminus Z_j)$. All constructions and proofs of the procedure can be carried on with these
inductive assumptions''.

\emph{Covariance.} Finally, at constant $w$ the graph averaging and the covariant inputs give the
Euclidean covariance required of the boundary terms in the inductive hypotheses, for the new terms, as
stated in \cite{B-CMP122b}: ``they are Euclidean covariant''.
\end{proof}

\begin{definition}[Plain interior units at the resummation site]\label{5.2:def-plain-interior-units}
The definition has three parts: \textup{(a)} the setting and the weights, \textup{(b)} the
convention by which the units are read, and \textup{(c)} the class of plain interior units.

\smallskip
\noindent\textup{(a)} \emph{Setting.} We work at the resummation site of the operation
$\mathbf{R}_k$ and put $N := N(g_k)$, $h := k-N$ and $k_0 := k-N_0(g_k)$. Let $Z$ be the
treated class at this site, a large-field region, and $C$ one of its components; by the box
property of the components fixed with the geometry of the treated class, $C$ is a box. The sets
$\Lambda_C$, $\Lambda_C^{\sim}$, $D_2$ and the core
\begin{equation*}
\mathsf{C} := \mathsf{C}_C := {Z'_h}^{\sim}\cap C
\end{equation*}
are those of the geometry of the treated class. They are defined there by means of its
nesting property, and they are the sets of the same names in the construction of the three
families per scale at this site. We write $\mathfrak{B} := \mathbf{B}''_k$ and $\pi(y)$ for
its level at a point $y$.

The following objects are those fixed with that geometry: cells, their level
$\operatorname{lev}$, zones, the partitions $\pi_\ell$, the hull $[\cdot]_k$, the lengths
$d_\ell$, the distances $\operatorname{dist}_\ell$, and the relations ``meet'' and ``touch''.
The sets $Z_j$ and $\Omega_j$ are always those of the current label. The scaled distance
$d^{\wedge}$, its length, which we denote $\mathfrak{w}$, and the rate $\delta_P$, with their
properties, are those fixed with the scaled distance. Put
\begin{equation*}
\ell_n := L^n\eta,\qquad a := \tfrac{31}{20}\kappa,\qquad
\varsigma_0 := \tfrac15\kappa,\qquad \varsigma_1 := \tfrac14\delta\kappa,
\end{equation*}
where $\eta$ is the length constant of the apparatus at this site, and recall $d=4$ and
$L\ge 8$.

An \emph{attempt} $\mathsf{a}$ on $C$ is one of two things:
\begin{itemize}
\item the running one, for which $k_{\mathsf a} := k$ and $X_{\mathsf a} := C$;
\item a live arrest $\mathfrak{a}$ with $X_{\mathfrak a}\subset C$, in which case we also write
$\mathsf a$ for $\mathfrak a$.
\end{itemize}
At its own step an attempt carries the indices $h_{\mathsf a}$ and
$k_{0,\mathsf a}\ge h_{\mathsf a}+1$; the latter is the index $k_0$ of that step, formed as at
the running step, so that $k_{0,\mathsf a}=k_{\mathsf a}-N_0(\mathsf a)$. The number
$N_0(\mathsf a)$ is given by
$L^{N_0(\mathsf a)} = L\check{R}_{k_{0,\mathsf a}+1}$, as in \cite{B-CMP122a}.

For a live arrest $\mathfrak a$ one has $k_{\mathfrak a}\le h$ and
\begin{equation*}
X_{\mathfrak a}\subset Z_{k_{\mathfrak a}}\cap C\subset Z_h\subset \mathsf{C}.
\end{equation*}
The ranges $]h_{\mathfrak a},k_{\mathfrak a}]$ of distinct live arrests on $C$ are pairwise
disjoint; both facts are those of the arrest construction. A \emph{window cell} of an attempt
$\mathsf a$ is a cell $\Delta\subset X_{\mathsf a}$ whose level lies in
$]k_{0,\mathsf a},k_{\mathsf a}]$.

\emph{Weights.} Let $\Delta\subset C$ be a cell of level $n$.

The \emph{window weight} $\mathsf{w}^{\mathrm{win}}(\Delta)$ equals $L^{N_0(\mathsf a)}$ if
$\Delta$ is a window cell of an attempt $\mathsf a$; such an attempt is unique, as shown in the
proof below. Otherwise $\mathsf{w}^{\mathrm{win}}(\Delta) := 1$. The weight
$\mathsf{w}^{\mathrm{win},\vee}(\Delta)$ is the maximum of $\mathsf{w}^{\mathrm{win}}$ over the
cells at $\operatorname{dist}_n\le 1$ from $\Delta$. Further,
$\varrho(\Delta) := 1+d^{\wedge}(\Delta,\mathsf{C}^c)/\mathfrak{w}$, and $\varrho(\Delta):=1$ off
$\mathsf{C}$; here $\mathfrak{w}$ is the length of the scaled distance just recalled.

The \emph{fine weight} $w(\Delta)$ is the power of Ba{\l}aban's auxiliary constant
$\mathsf{L}_0$ (recall $\mathsf{L}_0^2\le L/3$) that the clause of the domain
$\mathfrak{D}^{\flat}_N(\cdot)$ carries at $\Delta$. Here $\mathfrak{D}^{\flat}_N(\cdot)$ is the
domain of \cite[(1.64)--(1.67), p.~190]{B-CMP122a} read at $n=N$. It has two branches; in the
notation of \cite{B-CMP122a}, with Ba{\l}aban's indices $m,j,k_0$ and region $X$, they read as
follows.
\begin{itemize}
\item[(i)] $w(\Delta)=\mathsf{L}_0^{2e}$, where $e=\operatorname{lev}\Delta-k_{0,\mathsf a}$.
This branch applies on a window cell of an attempt $\mathsf a$ under which the label preceding
the attempt has levels $\le k_{0,\mathsf a}$. These cells are:
\begin{itemize}
\item the strata of levels in $]k_{0,\mathsf a},k_{\mathsf a}[$;
\item the top shell off $\Omega_{k_{0,\mathsf a}+1}$.
\end{itemize}
They correspond to Ba{\l}aban's clauses on $(\Omega'_m\setminus\Omega'_{m+1})\cap X$ for
$m=1,\dots,j-1$ and on $(\Omega'_j\setminus\Omega_{k_0+1})\cap X$ for $m=j$. Branch (i)
stands on $\Omega^c_{k_{0,\mathsf a}+1}$ only.
\item[(ii)] $w(\Delta)=\mathsf{L}_0^{2e}$, where $e=k_{\mathsf a}-m'$. This branch applies on the
old ring of the top zone whose level before the attempt is
$m'\in\,]k_{0,\mathsf a},k_{\mathsf a}]$. It corresponds to Ba{\l}aban's clauses on
$(\Omega_m\setminus\Omega_{m+1})\cap X$ for $m=k_0+1,\dots,j-1,j$.
\end{itemize}
On a live piece, $w$ is the weight $w_{\mathfrak a}$ recorded in the frozen structure of the
arrest $\mathfrak a$. Wherever no power of $\mathsf{L}_0$ stands, $w(\Delta)=1$.

In particular $w=1$ on all of $\mathfrak{T} := \Omega_k\cup Z^c$ off the live pieces. Two
facts give this. First, the collar of $C$ has old level $k$, so branch (ii) applies there with
$m'=k$. Second, the clauses inherited on the untouched shells carry no power of
$\mathsf{L}_0$; this is the inheritance statement for the domains. In one display:
\begin{equation}\label{5.2:eq-plain-interior-units-a}
\begin{aligned}
\mathsf{w}^{\mathrm{win}}(\Delta) &:=
\begin{cases}
L^{N_0(\mathsf a)} & \text{if $\Delta$ is a window cell of the attempt $\mathsf a$},\\
1 & \text{otherwise},
\end{cases}\\
w(\Delta) &:=
\begin{cases}
\mathsf{L}_0^{2e}, & e=\operatorname{lev}\Delta-k_{0,\mathsf a}\ \text{in branch (i)},\quad
e=k_{\mathsf a}-m'\ \text{in branch (ii)},\\
w_{\mathfrak a} & \text{on a live piece of $\mathfrak a$},\\
1 & \text{wherever no power of $\mathsf{L}_0$ stands},
\end{cases}\\
w^{\vee}(\Delta) &:= \max\{w(\Delta') : \operatorname{dist}_n(\Delta',\Delta)\le 1\}.
\end{aligned}
\end{equation}
On $C$ these weights satisfy
\begin{equation}\label{5.2:eq-plain-interior-units-1}
w(\Delta)\le \mathsf{w}^{\mathrm{win}}(\Delta),\qquad
w^{\vee}(\Delta)\le \mathsf{w}^{\mathrm{win},\vee}(\Delta).
\end{equation}

Now let $\Delta\not\subset C$ be a cell met by the hull of an exterior unit, namely a unit
of the exterior counterpart of the present class. There $w(\Delta)$ is the weight which the
same domain clause carries. It equals $1$ off the live pieces, and it equals the frozen weight
$w_{\mathfrak a'}$ on the window of a live arrest $\mathfrak a'$ foreign to $C$. No window
weight $\mathsf{w}^{\mathrm{win}}$ is defined off $C$.

The window weight $\mathsf{w}^{\mathrm{win}} = L^{N_0(\mathsf a)} = L\check{R}_{k_{0,\mathsf a}+1}$
is used in two ways only:
\begin{itemize}
\item as a count of the cells of a window: a window cell of level $n$ holds at most
$L^{d(n-k_{0,\mathsf a})}$ old cells;
\item as a majorant of $w$.
\end{itemize}
Every size, threshold and gauge reading carries the fine weight $w$, through
$\bar{\mathsf{w}}_v := 3w^{\vee}$ of part (c).

The parameter polydisc is $\mathfrak{P}\ni m=(\zeta,\varpi,\vec{\mathsf s})$; the decoupled
carrier $P_{\hat\lambda}$ and the radius $r^0$ are those of the construction of the three
families per scale. The pair $\operatorname{chain}(m)$ is the carrier pair of the resummation
site at the background $\mathfrak{b}_4$, and $\mathcal{U}(m)$ is its $U$-slot. For
$S\subset\mathsf{C}$ put
\begin{equation*}
\rho_S := r^0 e^{\frac12\delta_P d^{\wedge}(S,\mathsf{C}^c)}.
\end{equation*}
The units of entries (1)--(3), (5) and $(\partial)$ of the table of units at the resummation
site are invariant under the gauge transformations with values in the complexification of
$\mathrm{SU}(N)$, \cite[(1.10)]{B-CMP109}. They are therefore functions of the orbit, and neither
the presented rotation nor the hatted carrier plays any role for them.

\smallskip
\noindent\textup{(b)} \emph{The reading convention.} By the expansion convention at the
resummation site, every interior first-part unit $v$ (that is, $X_v\subset Z$), plain or
wrapped, is expanded by the two link rules of that convention. Neither Ba{\l}aban's partition
of unity $\zeta_{\square}$ nor the amplitude $\mathsf t$ enters for these units; the amplitude
$\mathsf t$ of the amplitude lemma enters only at the fourth link, and only on the plain
first-part units with $K_v\setminus Z\neq\emptyset$.

Hence, for a plain $v$ with $X_v\subset Z$, the four links are as follows.
\begin{itemize}
\item The first link is Ba{\l}aban's single interpolation bracket \cite[p.~375]{B-CMP122b}. It
is expressed by the fundamental theorem of calculus \cite[(1.9)]{B-CMP116} in the decoupling
parameters $\mathsf{s}^{(1)}$ alone, over the cells disjoint from $\Lambda^{\sim}$.
\item The second to fourth links are taken over the cells disjoint from $\Lambda$
\cite[p.~375]{B-CMP122b}.
\end{itemize}
The carriers are read as the carrier-reading convention prescribes, and the sets
$Y^{\mathrm{tr}}_t$ are Ba{\l}aban's $Y_t$. No member of the expansion meets the interior of
$\Lambda$.

The units are read as follows; the core $\mathfrak{i}^{(j)}_y$, the index $p_v$ and the box
$Q_v$ are those of part (c). Let $v$ be a smooth far core in the sense of part (c). It may
be an interior one, or the far core of a plain second-part unit of clause (d) of the expansion
convention, which belongs to the exterior counterpart. For such $v$, $f_v(m)$ means
$\mathfrak{i}^{(j)}_y$ evaluated at
\begin{equation}\label{5.2:eq-plain-interior-units-2}
\operatorname{ext}_v(m) := (U_p,J_p)\bigl(Q_v,M^p(\mathcal{U}(m)\restriction Q_v)\bigr),\qquad
p:=p_v.
\end{equation}
This is the box minimiser of the slot's own level-$p$ averages on $Q_v$, with its own current
\cite[(3.24), (3.27)]{B-CMP109}, \cite{B-CMP119}. Here $M^p(\cdot)$ denotes the level-$p$
averaging operation of Ba{\l}aban's construction; it is not a power of the constant $M$.

The same $\operatorname{ext}_v$ is read in two further places: for the $\mathsf t$-piece built
on such a unit, and for the subtracted copy of clause (d) at the undilated, uncut carrier. The
reading is on the joint $(\mathsf t,\vec{\mathsf s})$-polydisc, so that one declaration of
analyticity covers the difference. The subtracted copies' smooth far cores
\cite{B-CMP122b} are read at $\operatorname{ext}_v$ at $(n,\mathsf{W})=(k,3)$; rough copies are
read at $\zeta$.

Equivalently, by clause (c1) of the core comparison together with \cite{B-CMP102b}, $f_v(m)$
is $f_v$ at
$(U_j,J_j)(\square_0,\mathsf{V}(m))$, where $\mathsf{V}(m)$ denotes the level-$j$ data of
$\operatorname{ext}_v(m)$. Rough cores, and every other unit of $\mathfrak{L}^{\mathrm{pt}}$ and
of $\mathfrak{L}^{\mathrm{ex}}$, are read on $\operatorname{chain}(m)$ itself. Where the
amplitude lemma acts, this carrier is dilated by $(t+\mathsf t\hat\chi_{\square})\mathbb{H}_s$;
the letter $t$ in this factor is that of the amplitude lemma and is not the ratio
$t:=L^{j-n}$ of part (c).
There is no $\zeta_{\square}$, no $t_{\square}$ and no exemption, and there is no $\mathsf t$ on
an interior unit. There is one weight,
$\bar{\mathsf{w}}_v := 3w^{\vee}(\Delta_{y_v})$, with $w=\mathsf{L}_0^{2e}$ the fine weight of
part (a).

The reading \eqref{5.2:eq-plain-interior-units-2} is meaningful because of the following
statement of the admissibility lemma for smoothed box minimisers. For $\mathbf{U}$ in the ball
of the radius lemma enlarged by the factor $c''$, over any box $Q$,
\begin{equation*}
U_p\bigl(Q,M^p(\mathbf{U})\bigr)\ \text{is admissible at scale $p$ on } Q^{\sim-2},
\end{equation*}
where $Q^{\sim-2}$ denotes $Q$ with two boundary layers of $\pi_p$-cubes removed (so that
$Q_v^{\sim-2}=[\square_0]_{p_v}$).
That lemma rests on the nested-box statement of the treatment of the pointed entries in the
proof of the correlation theorem, on
\cite[Proposition~9]{B-CMP102b} at the trivial background $V_0=\mathbb{1}$, and on
\cite[(3.27), (3.32)]{B-CMP109}.

Real values are unchanged. At real $\zeta$ and $\vec{\mathsf s}=\vec{\mathbb{1}}$ the carrier
$\operatorname{chain}(m)$ is Ba{\l}aban's minimiser; this is the real-point identification of
the carrier. Apply clause (c1) of the core comparison with $(j,X,\square_0)$ replaced by
$(p,[\square_0]_p,Q_v)$; $Q_v$ carries the two layers of $\pi_p$-cubes that the clause
requires. At such real points it gives $\operatorname{ext}_v(m)=\operatorname{chain}(m)$
modulo a gauge transformation on $[\square_0]_p$. The current slot is not read, and the
support condition of the expansion holds with $\mathcal{R}_0 := Q_v\subset C$.

\smallskip
\noindent\textup{(c)} \emph{The class.} The class
$\mathfrak{L}^{\mathrm{pl}} = \mathfrak{L}^{B}\sqcup\mathfrak{L}^{\mathrm{pt}}$ consists of the
units $v$ of the table of units at the resummation site with $\mathcal{P}_v=\emptyset$ and
$X_v\subset Z$; this is the list in the filing proposition. Each $X_v$ lies in a single
component $C$, and $K_v := [X_v]_k\subset C$.

\emph{The class $\mathfrak{L}^{B}$} (entry (4)) consists of the plain stored boundary-type
terms in the format (f0)--(f4) of boundary-type terms at $\mathcal{P}_F=\emptyset$. These are:
\begin{itemize}
\item the terms $\mathbf{B}^{(j)}$ and $\mathbf{B}'^{(j)}$;
\item the chain terms $\mathbb{C}^{(n)}_k$ of this step;
\item the frozen chain terms $\mathbb{C}^{(i)}_{k_{\mathfrak a}}$;
\item the change-of-variables terms of the preparatory stages, filed with the per-cube
constant $b_\delta$.
\end{itemize}
For these, $X_v\in\mathbf{D}_{j_v}$, $\mathfrak{N}(v):=\mathfrak{N}^{A,\mathbb{C}}(v)$ and
$\hat X_v := X_v$.

\emph{The class $\mathfrak{L}^{\mathrm{pt}}$} (entries (1)--(3), (5), $(\partial)$) consists
of units of level $j\le n$, where $n=n_v$ is the zone of the unit's point $y_v$ introduced
below. The level is restricted as follows.
\begin{itemize}
\item On entries (1)--(3) and $(\partial)$, for $X$ meeting $Z$, one has $j<k$
\cite[(2.41)]{B-CMP119}.
\item In clause (d) of the expansion convention, where $X\cap Z=\emptyset$, one has $j\le k$.
At $j=k$:
\begin{itemize}
\item entries (1), (3) and $(\partial)$ decay at the rate $2\kappa$ of the lemma on the decay
fraction $\delta_J$, as in Ba{\l}aban's bound $\exp(-(1+4\beta)\kappa d_k(X))$ for $j=k$
\cite{B-CMP122b};
\item entry (2) is void;
\item by the lemma classifying the singles at the current level, every single is long or a
tail.
\end{itemize}
\end{itemize}
Each such unit carries a point $y := y_v\in X_v$ of zone $n:=n_v$, with cell $\Delta_y$. Put
$t:=L^{j-n}$ and $\nu:=1+n-j$, let $\square\in\pi_j$ be the cube of $y$, and put
\begin{equation*}
\bar{\mathsf{w}}_v := 3w^{\vee}(\Delta_y)\ \le\ 3\mathsf{w}^{\mathrm{win},\vee}(\Delta_y).
\end{equation*}
Wherever $\bar{\mathsf{w}}$ or $\mathsf{w}^{\mathrm{win},\vee}$ occurs in what follows as a
threshold factor, it is read only through $\bar{\mathsf{w}}_v$. The statement then holds with
$w^{\vee}$ in place of $\mathsf{w}^{\mathrm{win},\vee}$, and a fortiori as stated.

A point $y$ is called far or near in the sense of the classification of pointed locations at
the resummation site.
\emph{Singles} are the differences $F(X,\cdot)-F(X,1)$ with $y\in X\in\mathbf{D}_j$ and
$X\subset\Lambda_j$. They comprise three kinds:
\begin{itemize}
\item the tails of entry (1), which have far $y$ and $X\not\subset\square^{\sim\nu}$;
\item entries (2) and (3);
\item the constituents of entry $(\partial)$, which have near $y$.
\end{itemize}
Their size $\mathfrak{n}_v$ is defined entry by entry, as the table of units at the
resummation site prescribes, as one of
\begin{equation*}
E_0e^{-\kappa d_j(X)},\qquad C_BE_0\,\omega_{k,j}\,e^{-(1-\frac14\delta)\kappa d_j(X)},\qquad
g_j^{\kappa_0}e^{-\kappa d_j(X)},
\end{equation*}
with $\omega_{k,j}=L^{-(k-j)/2}$; here $E_0$ is the constant factor and $\kappa_0$ the
exponent of $g_j$ in these expressions, both fixed with the table of units. It is one letter
for all entries. At a near unit its meaning is
$\sup_{\mathfrak{D}_v}|f_v|\le\mathfrak{n}_v$, an inequality and not an identity. At $j=k$,
which occurs in clause (d) only, the sizes carry the exponent $2\kappa$ of clause (ii) of the
lemma on $\delta_J$ in place of $\kappa$.

Let $\hat X$ be the smallest box of $\pi_j$-cubes containing $X$; it lies in $C$. The single is
\emph{short} if it is not a tail and every side of $\hat X$ is at most $L^{\nu-2}$
$\pi_j$-cubes, that is, at most one cube of $\pi_{n-1}$. Otherwise it is \emph{long}. For long
singles the following hold:
\begin{equation}\label{5.2:eq-plain-interior-units-b}
\begin{aligned}
&d_j(X)\ge L^{\nu-2}-1 &&\text{for a long single that is not a tail},\\
&d_j(X)\ge\nu &&\text{for a tail},\\
&X\not\subset\square_y^{\sim\nu}\ \text{and}\ d_j(X)\ge\nu
&&\text{for every long single, when $\nu\ge 3$.}
\end{aligned}
\end{equation}

\emph{Wilson pieces} are of two kinds: the units of entry (5), for which $X_v$ is a block, and
the terms $b_jA(h_y,\cdot)$ at near $y$.

\emph{Cores} are defined for far $y$ by
\begin{equation*}
\mathfrak{i}^{(j)}_y := \sum_{X\ni y,\ X\subset\square^{\sim\nu}}
\bigl[F(X,\cdot,y)-F(X,1,y)\bigr]-b_jA(h_y,\cdot),
\end{equation*}
with $X_v := \square_0 := \square^{\sim\nu+1}\subset Z$. For $p\ge j$ let
$Q^{(p)} := [\square_0]_p$ with two layers of $\pi_p$-cubes. With the smoothing-depth function
$s(\cdot)$ of the smoothing-depth rule (this $s$ is neither the depth below the cutoff nor the
slot combination) put
\begin{equation*}
p_v := \max\bigl\{p:\ j\le p\le n-2-s(\bar{\mathsf{w}}_v),\ Q^{(p)}\subset Z\bigr\}.
\end{equation*}
If $p_v$ exists, the core is \emph{smooth}, and we put $a_v := n-p_v\ge 3$,
$b_v := p_v-j$ and $Q_v := Q^{(p_v)}$. Otherwise the core is \emph{rough}. The letter $a_v$
should not be confused with $a=\frac{31}{20}\kappa$, which occurs only in the norm
$\|\cdot\|_a$ and in $\mathsf{Q}_{a+9}$; the subscript of $Q_a$, $\theta_a$ and
$\mathbf{r}_a$ is an entry index.

\emph{Hulls.} The hull $\hat X_v$ is defined as follows:
\begin{itemize}
\item $\hat X_v := \hat X$ for a short single;
\item $\hat X_v := X$ for a long single;
\item $\hat X_v := X_v$ for a Wilson piece and for a rough core;
\item $\hat X_v := Q_v$ for a smooth core.
\end{itemize}
Put $\operatorname{rad}(v) := \max_{x\in\hat X_v}\operatorname{dist}_n(x,y)$.

A unit is \emph{deep} if $X_v\subset\mathsf{C}$, and \emph{shallow} otherwise. Put
$\hat D_v := d^{\wedge}(\hat X_v,\mathsf{C}^c)$ if $\hat X_v\subset\mathsf{C}$, and
$\hat D_v := 0$ otherwise. The chain terms $\mathbb{C}$ of this step and the constituents of
entry $(\partial)$ are never deep, by \cite[(1.63)]{B-CMP122a} and the geometry of the treated
class.
\end{definition}

\begin{proof}[Proof of the assertions made in the definition]
We prove the uniqueness of the attempt in the definition of $\mathsf{w}^{\mathrm{win}}$,
the inequalities \eqref{5.2:eq-plain-interior-units-1}, the bound on $\bar{\mathsf{w}}_v$, the
equivalent form of shortness, and \eqref{5.2:eq-plain-interior-units-b} except its clause
$X\not\subset\square_y^{\sim\nu}$ for a long single that is not a tail, which is taken from the
containment remark for long units.

\emph{Uniqueness of the attempt.} Every attempt satisfies $k_{0,\mathsf a}\ge h_{\mathsf a}+1$.
Hence
\begin{equation*}
]k_{0,\mathsf a},k_{\mathsf a}]\subset\,]h_{\mathsf a},k_{\mathsf a}].
\end{equation*}
The ranges $]h_{\mathfrak a},k_{\mathfrak a}]$ of distinct live arrests are pairwise disjoint.
Each of them is also disjoint from the running range $]h,k]$, because $k_{\mathfrak a}\le h$.
So the level $n$ of a cell lies in $]k_{0,\mathsf a},k_{\mathsf a}]$ for at most one attempt
$\mathsf a$, and $\mathsf{w}^{\mathrm{win}}(\Delta)$ is well defined.

\emph{The first inequality of \eqref{5.2:eq-plain-interior-units-1}.} Let $\Delta\subset C$.
In both branches of the fine weight the exponent satisfies $e\le k_{\mathsf a}-k_{0,\mathsf a}$:
in branch (i) because $\operatorname{lev}\Delta\le k_{\mathsf a}$, in branch (ii) because
$m'>k_{0,\mathsf a}$. Moreover $k_{\mathsf a}-k_{0,\mathsf a}=N_0(\mathsf a)$ by part (a).
Hence, since $\mathsf{L}_0^2\le L/3$, every weight $\mathsf{L}_0^{2e}$ of either
branch of an attempt $\mathsf a$ obeys
\begin{equation*}
\mathsf{L}_0^{2e}\le (L/3)^e\le L^{N_0(\mathsf a)}.
\end{equation*}
There are four cases.
\begin{itemize}
\item If $w(\Delta)=1$, the inequality holds because $\mathsf{w}^{\mathrm{win}}\ge 1$.
\item In branch (i), $\Delta$ is a window cell of $\mathsf a$, so
$\mathsf{w}^{\mathrm{win}}(\Delta)=L^{N_0(\mathsf a)}$, and the displayed inequality gives
$w(\Delta)\le\mathsf{w}^{\mathrm{win}}(\Delta)$.
\item In branch (ii) of the running attempt, take a ring cell with $e=k-m'>0$. It lies in the
running window: it is a cell of the top zone of $C=X_{\mathsf a}$, hence of current level
$k\in\,]k_0,k]$. Hence
$\mathsf{w}^{\mathrm{win}}(\Delta)=L^{N_0(\mathsf a)}$ for the running attempt, and the same
chain of inequalities applies. A ring cell with $e=0$ has $w(\Delta)=1$ and falls under the
first case.
\item On a live piece of a live arrest $\mathfrak a$, the frozen structure of $\mathfrak a$
assigns to a window cell of $\mathfrak a$, that is, to a cell of $X_{\mathfrak a}$ of level in
$]k_{0,\mathfrak a},k_{\mathfrak a}]$, a weight $\mathsf{L}_0^{2e}$ in one of the two branches
at the arrest's step, and weight $1$ elsewhere. In the first event
$\mathsf{w}^{\mathrm{win}}(\Delta)=L^{N_0(\mathfrak a)}$ and the displayed inequality with
$\mathsf a=\mathfrak a$ gives the claim; in the second the first case applies.
\end{itemize}

\emph{The second inequality of \eqref{5.2:eq-plain-interior-units-1}, and the bound on
$\bar{\mathsf w}_v$.} Take the maximum of the first inequality over the cells at
$\operatorname{dist}_n\le 1$ from $\Delta$. Multiplying by $3$ at $\Delta=\Delta_y$ gives
$\bar{\mathsf{w}}_v\le 3\mathsf{w}^{\mathrm{win},\vee}(\Delta_y)$.

\emph{Shortness.} Since $\nu-2=n-1-j$, the number $L^{\nu-2}$ is the number of $\pi_j$-cubes
along one side of a $\pi_{n-1}$-cube. So a side of at most $L^{\nu-2}$ $\pi_j$-cubes is a side
of at most one cube of $\pi_{n-1}$.

\emph{\eqref{5.2:eq-plain-interior-units-b}.}

First, let the single be long but not a tail. Then some side of $\hat X$ exceeds $L^{\nu-2}$
$\pi_j$-cubes. Because $\hat X$ is the smallest box of $\pi_j$-cubes holding $X$, the set $X$
contains two $\pi_j$-cubes at distance at least $L^{\nu-2}-1$ in that direction. The lower
bound of the tree length by the distance of two cubes of the set, which is part of the
geometry of the treated class, then gives $d_j(X)\ge L^{\nu-2}-1$.

Second, a tail has $d_j(X)\ge\nu$ by clause (i) of the tail lemma.

Third, let $\nu\ge 3$. For a tail, $X\not\subset\square_y^{\sim\nu}$ holds by definition and
$d_j(X)\ge\nu$ was just shown. For a long single that is not a tail, $L\ge 8$ gives
\begin{equation*}
d_j(X)\ge L^{\nu-2}-1\ge 8^{\nu-2}-1\ge\nu .
\end{equation*}
The last inequality holds at $\nu=3$, where it reads $7\ge 3$, and it persists for larger
$\nu$, since the left member is multiplied by at least $8$ when $\nu$ increases by one while
the right member increases by one. The clause $X\not\subset\square_y^{\sim\nu}$ for such a
single at $\nu\ge 3$ is taken by statement from the containment remark for long units, where it
is stated together with $d_j(X)\ge\nu$.
\end{proof}

\begin{lemma}[The exterior-piece sum for plain units]\label{5.2:lem-plain-exterior-sum}
In the setting and notation of Definition~\ref{5.2:def-plain-interior-units} (in particular
$N:=N(g_k)$, the number of preparatory stages, $h:=k-N$, the core $\mathsf{C}$, the classes
$\mathfrak{L}^{\mathrm{pl}}=\mathfrak{L}^{B}\sqcup\mathfrak{L}^{\mathrm{pt}}$, the labels $K_v$,
the hulls $\hat X_v$, the distinction between shallow and deep units, the depths $\hat D_v$, the
carrier $\operatorname{chain}(m)$ and the polydisc $\mathfrak{P}$), let the parameters of the
correlation theorem obey the floors of its parameter sentence, including
the floor that accompanies clause (b) of the cell lemma of the resummation of wrapped units and
the floor
\begin{equation}\label{5.2:eq-plain-exterior-sum-1}
\kappa_1-1\ge c_l(a+9)+\lambda_*,\qquad c_l=1+234d,\qquad a:=\tfrac{31}{20}\kappa .
\end{equation}
Let $\gamma\le\gamma_J$, where the ceilings defining $\gamma_J$ include those of the third-factor
lemma, the regularity ceiling \eqref{5.2:eq-carrier-regularity-a} and the ceiling
\eqref{5.2:eq-plain-exterior-sum-a} below. Consider the resummation site of the operation
$\mathbf{R}_k$, read under the convention of hypothesis (i), and assume:
\begin{enumerate}
\item[(i)] the reading convention at the resummation site: every interior first-part unit
($X_v\subset Z$), plain or wrapped, is expanded with no block bump $\hat{\chi}_{\square}$ and no
amplitude $\mathsf{t}$; a smooth far core is read at the box minimiser of the carrier's own
level-$p_v$ averages on $Q_v$, rough cores and all other units on $\operatorname{chain}(m)$
itself, with the one weight $\bar{\mathsf{w}}_v:=3w^{\vee}(\Delta_{y_v})$; the smooth far cores of
Ba{\l}aban's subtracted copies are read at the same extension; and the leaves are declared as
clause-bearers in the inherited domains;
\item[(ii)] the third-factor lemma, clauses (a) and (b), its increment bound and the polydisc
lemma, with the choice of the radius $r^0$ made in their proofs;
\item[(iii)] the geometry lemma, clauses (f) and (g), the cell lemma, clause (b), and the
summation lemma, clauses (a)--(e), of the resummation of wrapped units;
\item[(iv)] the decoupling-radius lemma, the source-anchored analyticity statement and the
source-anchored Cauchy bound at the resummation site, for every consumer, and the key statement
on the domains of consumers with its inheritance and arrest clauses;
\item[(v)] the boundary-term bound lemma of the resummation of wrapped units;
\item[(vi)] the core estimates for the real carrier at the resummation site, the near-point
lemma, and the admissibility, at every scale, of the box minimiser of the averages at that scale;
\item[(vii)] the radius lemma, the disc estimates and the complex-weight lemma, with the floors
$\kappa\ge\kappa_d\ge 20(c_d+2)$ and $\kappa\ge 10\log L$;
\item[(viii)] the own-field statement for the second background on the core $\mathsf{C}$;
\item[(ix)] the filing of the change-of-variables terms of the live arrests, for
$\mathfrak{L}^{B}$.
\end{enumerate}
Then for every $K$, every $k\le K$ and every large-field label the following hold.

(a) For $v\in\mathfrak{L}^{\mathrm{pl}}$,
$f_v(\operatorname{chain}(\vec{\mathbb{1}}))=\mathrm{val}_v+\sum\mathfrak{p}$ with
$\mathrm{val}_v:=f_v(\operatorname{chain}(\vec{0}))$ and
\begin{equation*}
\mathfrak{p}=(v;\mathsf{y}_1,\dots,\mathsf{y}_4)
:=\prod_{t=1}^{4}\prod_{\Delta\in\mathsf{y}_t}\int_0^1 d\mathsf{s}^{(t)}(\Delta)\,
\partial_{\mathsf{s}^{(t)}(\Delta)}\,
f_v\bigl(\operatorname{chain}(\vec{\mathsf{s}})\bigr)
\Big|_{\mathsf{s}^{(t)}=0\ \text{off}\ \mathsf{y}_t},
\end{equation*}
the sum running over the tuples $(\mathsf{y}_t)$, not all empty, with
$\mathsf{y}_t\in\mathfrak{Y}((\operatorname{lab}_{t-1},\mathcal{R}_{t-1}))\cup\{\emptyset\}$,
$\operatorname{lab}_0:=K_v$, and $\operatorname{lab}_t$, $\mathcal{R}_{t-1}$ those of the
four-link expansion on labels with $\mathcal{R}_0:=\hat X_v\cup X_v$. The set
\begin{equation*}
Y_4(\mathfrak{p}):=\operatorname{lab}_4\cup\{\text{components of }Z\text{ meeting }
\operatorname{lab}_4\}
\end{equation*}
lies in $\mathbf{D}_k^Z$ and $d_{k,Z}(Y_4(\mathfrak{p}))\le d_k(\operatorname{lab}_4)$;
$\mathfrak{p}$ is linear in $f_v$ and analytic over $Y_4(\mathfrak{p})$ (jointly with the source in
$\mathcal{W}_k$), on the chart fixed for boundary-type units by the chart proposition, clause (a);
no amplitude $\mathsf{t}$ occurs, and no disc moves a chart.

(b) $|\mathfrak{p}|\le\|v\|_\star\,e^{-(\kappa_1-1)\sum_t|\mathsf{y}_t|}$, where
$\|v\|_\star:=\mathfrak{N}^{\flat}(v)$ if $v$ is shallow and
$\|v\|_\star:=C^{+}e^{-\frac14\delta_P\hat D_v}\mathfrak{N}^{\flat}(v)$ if $v$ is deep; here
$\mathfrak{N}^{\flat}:=\mathfrak{N}$ on $\mathfrak{L}^{B}$, $\mathfrak{N}^{\flat}$ is given by the
table \eqref{5.2:eq-smoothing-depth-a} on $\mathfrak{L}^{\mathrm{pt}}$, and
$C^{+}:=\max\{C_{3F},C^{\mathrm{pt}}\}$, where $C_{3F}$ is the constant of the third-factor lemma
(hypothesis (ii)) and $C^{\mathrm{pt}}$ that of the depth gain for deep pointed units
(Proposition~\ref{5.2:prop-deep-unit-gain}).

(c) Put $\mathfrak{v}^{\mathrm{int,pl}}(Y_4):=\sum_{Y_4(\mathfrak{p})=Y_4}\mathfrak{p}$ for
$Y_4\setminus Z\neq\emptyset$. Then
\begin{align*}
&\|\mathfrak{v}^{\mathrm{int,pl}}\|_a\le\|\mathfrak{v}^{\mathrm{int,pl}}\|_{a,1}
:=\sup_{\Delta\not\subset Z}\sum_{Y_4\ni\Delta}e^{a\,d_{k,Z}(Y_4)}
|\mathfrak{v}^{\mathrm{int,pl}}(Y_4)|\le\varepsilon^{\mathrm{int}}(k),\\
&\varepsilon^{\mathrm{int}}(k):=e^{-(\check{R}_k-2)}(1+C_\Sigma C_e)^4\,
\mathsf{Q}^{\star,\mathrm{pl}}_{a+9},
\end{align*}
where $C_\Sigma$ and $C_e$ are the constants of the resummation of wrapped units entering its
summation and cell lemmas (hypothesis (iii)),
\begin{equation*}
\mathsf{Q}^{\star,\mathrm{pl}}_r:=\sup_{Q\in\pi_k}\sum_{v\in\mathfrak{L}^{\mathrm{pl}}:\,K_v\ni Q}
\|v\|_\star e^{r\,d_k(K_v)},
\end{equation*}
and, for $r\le\frac85\kappa$,
\begin{equation*}
\mathsf{Q}^{\star,\mathrm{pl}}_r\le\sup_C\bigl[C^{+}S^{\mathrm{dp}}(C)+S^{\mathrm{sh}}(C)\bigr].
\end{equation*}

(d) The ceiling
\begin{equation}\label{5.2:eq-plain-exterior-sum-a}
\sup_{\mathsf{T}\ge\mathsf{T}_\gamma}\mathsf{P}_{\mathrm{int}}(\mathsf{T})\,
e^{-(\mathsf{T}^{r_{\mathrm{pr}}}-2)}\le\tfrac12 ,
\end{equation}
in which $\mathsf{P}_{\mathrm{int}}$ is a polynomial-type majorant with coefficients fixed before
$\gamma$, including the summand of clause (d) of the exterior-piece sum for the exterior class
$\mathfrak{L}^{\mathrm{ex}}$, gives
$\varepsilon^{\mathrm{int}}(k)+\varepsilon^{\mathrm{ext},\sigma}(k)\le\frac12$ at every $k\le K$.
Together with the ceiling of $\gamma_J$ requiring that $(1+C_\Sigma C_e)^3\,e\,S_d$ (with $e$
Euler's number) times the three members of $\varepsilon_{\mathrm{br}}(k)$ carrying $\theta'(g_k)$,
taken at the weight $a+7$ in place of $a$ and with the constant of the first of them replaced by
$\mathsf{C}_1^{\mathrm{ex}}$, be at most $\frac12$, and with the ceiling of the resummation of
wrapped units which gives $\varepsilon^{\mathrm{rel}}(k)\le1$, the six members of
\eqref{5.2:eq-resummation-input-b} give $\varepsilon_{\mathrm{br}}(k)\le\frac12+1+\frac12=2$. The
bound \eqref{5.2:eq-resummation-input-b} follows from the exhaustiveness of the table of units and
tuples accompanying it, and not from this lemma alone.

(e) No constant or ceiling depends on $K$, $k$, an age, $N_0$, the number, nesting, size or past
of the arrests and of the old large-field regions carried by stored terms, $|\mathsf{C}|$ or
the volume; $N$ and $\check{R}_{h+1}$ enter $\mathsf{P}_{\mathrm{int}}$ polynomially, through
\begin{equation*}
N\le C_N(\check{\mathsf{T}}_k+1)^{r_{\mathrm{pr}}},\qquad
\check{R}_{h+1}\le\mathsf{R}^{\uparrow}(\check{\mathsf{T}}_k).
\end{equation*}

(f) For every $v\in\mathfrak{L}^{\mathrm{pt}}$ the map $m\mapsto f_v(m)$ is holomorphic on
$\mathfrak{P}$ and $|f_v(m)-f_v(1)|\le\mathfrak{N}^{\flat}(v)$.

(g) For every pointed member of $\mathfrak{L}^{\mathrm{ex}}$ --- a single, Wilson piece or
$\partial$-constituent $v$ of $\mathfrak{L}^{c}$; a single, Wilson piece, $\partial$-constituent or
far core underlying a $\mathsf{t}$-piece $p\in\mathfrak{L}^{Y_0}$ --- the same holds on
$\mathfrak{P}$ with the sizes $\mathfrak{N}^{\flat}$ of \eqref{5.2:eq-smoothing-depth-a} taken at
$\bar{\mathsf{w}}_v:=3w^{\vee}(\Delta_{y_v})$ and at the zone $n\le k$ of $y_v$; and, for a compact
$p$ (compact in the sense fixed in the exterior value bound on the carrier), on the joint polydisc
$\mathfrak{P}\times\{|\mathsf{t}|\le\rho_A\}$,
$\rho_A:=\mathsf{w}_L/(\eta_p\mathsf{f}_v)^{1/2}$, with $\frac{81}{64}\mathfrak{N}^{\flat}(v)$ in
place of $\mathfrak{N}^{\flat}(v)$. Here $\eta_p:=e^{-\delta_P d(\square^{\zeta},S_v)}$, and
$\eta_p:=1$ where the $2$-stencil of $S_v$ meets $\square^{\zeta}$; $S_v:=\mathfrak{r}_v$ for a
compact $p$, the containment of its $\mathbf{J}$-entry being that of the containment lemma for the
$\mathbf{J}$-entry.
\end{lemma}

\begin{proof}
\emph{Clause (a).} The expansion of $f_v(\operatorname{chain}(\vec{\mathbb{1}}))$ into
$\mathrm{val}_v$ and the pieces $\mathfrak{p}$, the membership
$Y_4(\mathfrak{p})\in\mathbf{D}_k^Z$, the linearity of $\mathfrak{p}$ in $f_v$, its analyticity
over $Y_4(\mathfrak{p})$ jointly in the source on the chart fixed for boundary-type units by the
chart proposition, clause (a), and the absence of an amplitude $\mathsf{t}$ and of a disc moving a
chart, are the four-link expansion of plain interior units, taken here by statement.

For the length, note that the adjoined components lie in $Z$, so
$Y_4(\mathfrak{p})\setminus Z=\operatorname{lab}_4\setminus Z$. An optimal tree graph of
$\operatorname{lab}_4$, whose length gives $d_k(\operatorname{lab}_4)$, is contained in
$\operatorname{lab}_4\subset Y_4(\mathfrak{p})$ and meets every cube of
$\operatorname{lab}_4\supset Y_4(\mathfrak{p})\setminus Z$. It is therefore admissible in the
definition of $d_{k,Z}(Y_4(\mathfrak{p}))$ as the length of a shortest tree graph contained in
$Y_4(\mathfrak{p})$ and meeting all $M$-cubes of the components of
$Y_4(\mathfrak{p})\setminus Z$ \cite[(1.67)]{B-CMP122b}. Hence
$d_{k,Z}(Y_4(\mathfrak{p}))\le d_k(\operatorname{lab}_4)$.

\emph{Clause (b).} For a shallow unit, the Cauchy estimate in the decoupling parameters on the
product polydisc is provided by the source-anchored analyticity statement, clause (ii), and by the
source-anchored Cauchy bound at the resummation site (hypothesis (iv)); it gives a factor
$e^{-(\kappa_1-1)}$ for each derivative $\partial_{\mathsf{s}^{(t)}(\Delta)}$,
$\Delta\in\mathsf{y}_t$, $t=1,\dots,4$, hence $e^{-(\kappa_1-1)\sum_t|\mathsf{y}_t|}$, times the
supremum on the polydisc of the function to which it is applied. For $v\in\mathfrak{L}^{B}$ it is
applied to $f_v$ itself, whose
supremum is at most $\mathfrak{N}(v)=\mathfrak{N}^{\flat}(v)$. For
$v\in\mathfrak{L}^{\mathrm{pt}}$ it is applied to $f_v-f_v(1)$: a piece with a non-empty tuple
carries at least one derivative $\partial_{\mathsf{s}^{(t)}(\Delta)}$, so it is unchanged when
$f_v$ is replaced by $f_v-f_v(1)$, and the supremum of $|f_v-f_v(1)|$ is at most
$\mathfrak{N}^{\flat}(v)$ by Proposition~\ref{5.2:prop-carrier-value-bound}. For a deep
$v\in\mathfrak{L}^{B}$, clause (b) of the third-factor lemma (hypothesis (ii)) gives
\begin{equation*}
|\mathfrak{p}|\le C_{3F}\,e^{-\frac12\delta_P d^{\wedge}(X_v,\mathsf{C}^c)}\,
\mathfrak{N}(v)\,e^{-(\kappa_1-1)\sum_t|\mathsf{y}_t|}.
\end{equation*}
For a deep $v\in\mathfrak{L}^{B}$ one has $\hat X_v=X_v\subset\mathsf{C}$, so
$d^{\wedge}(X_v,\mathsf{C}^c)=\hat D_v\ge0$, and
$e^{-\frac12\delta_P\hat D_v}\le e^{-\frac14\delta_P\hat D_v}$. For a deep
$v\in\mathfrak{L}^{\mathrm{pt}}$ the depth gain for deep pointed units,
Proposition~\ref{5.2:prop-deep-unit-gain}, gives
\begin{equation*}
|\mathfrak{p}|\le C^{\mathrm{pt}}\,e^{-\frac12\delta_P\hat D_v}\,
\mathfrak{N}^{\flat}(v)\,e^{-(\kappa_1-1)\sum_t|\mathsf{y}_t|},
\end{equation*}
with the depth $\hat D_v$ taken on the hull $\hat X_v$, and again
$e^{-\frac12\delta_P\hat D_v}\le e^{-\frac14\delta_P\hat D_v}$. As
$C_{3F},C^{\mathrm{pt}}\le C^{+}$, clause (b) follows.

\emph{Clause (c).} Let $Y_4\setminus Z\neq\emptyset$ and choose a cube $\Delta\subset Y_4$ with
$\Delta\not\subset Z$. Then $e^{a d_{k,Z}(Y_4)}|\mathfrak{v}^{\mathrm{int,pl}}(Y_4)|$ is one
summand of the sum at $\Delta$ defining $\|\mathfrak{v}^{\mathrm{int,pl}}\|_{a,1}$; taking the
supremum over such $Y_4$ gives the first inequality. The second inequality is the summation of the
pieces of plain interior units, taken here by statement, which combines clause (b) with the
summation lemma, clauses (a)--(e), the geometry lemma, clauses (f), (g), and the cell lemma,
clause (b) (hypothesis (iii)), under the floor \eqref{5.2:eq-plain-exterior-sum-1}. The bound of
$\mathsf{Q}^{\star,\mathrm{pl}}_r$, $r\le\frac85\kappa$, by
$\sup_C[C^{+}S^{\mathrm{dp}}(C)+S^{\mathrm{sh}}(C)]$ is the deep pile bound and the shallow pile
bound, taken here by statement, the deep units carrying the factor $C^{+}$ of $\|v\|_\star$. It
applies at $r=a+9$, since $a+9=\frac{31}{20}\kappa+9\le\frac85\kappa$, because
$\kappa/20\ge c_d+2\ge 9$ by the floor of hypothesis (vii) and $c_d=2^d\log(2de)+1$.

\emph{Clause (d).} By definition $\check{R}_k=R(\min_{i\le k}x_i)$ is a power of $L$ not smaller
than $\check{\mathsf{T}}_k^{r_{\mathrm{pr}}}$, and $\check{\mathsf{T}}_k\ge\mathsf{T}\ge
\mathsf{T}_\gamma$ for $k\le K$ by the correlation theorem's bound on the reference recursion. By
its construction in the statement on the interior ceiling majorant, taken here by statement, the
majorant $\mathsf{P}_{\mathrm{int}}$ dominates at $\mathsf{T}=\check{\mathsf{T}}_k$ the prefactor
of $e^{-(\check{R}_k-2)}$ in $\varepsilon^{\mathrm{int}}(k)$ and, through its summand from clause
(d) of the exterior-piece sum for the exterior class $\mathfrak{L}^{\mathrm{ex}}$, that in
$\varepsilon^{\mathrm{ext},\sigma}(k)$. Hence
\begin{equation*}
\varepsilon^{\mathrm{int}}(k)+\varepsilon^{\mathrm{ext},\sigma}(k)
\le\mathsf{P}_{\mathrm{int}}(\check{\mathsf{T}}_k)\,
e^{-(\check{\mathsf{T}}_k^{r_{\mathrm{pr}}}-2)}
\le\sup_{\mathsf{T}\ge\mathsf{T}_\gamma}\mathsf{P}_{\mathrm{int}}(\mathsf{T})\,
e^{-(\mathsf{T}^{r_{\mathrm{pr}}}-2)}\le\tfrac12
\end{equation*}
by \eqref{5.2:eq-plain-exterior-sum-a}. The six members of
\eqref{5.2:eq-resummation-input-b} are the three members carrying $\theta'(g_k)$,
$\varepsilon^{\mathrm{rel}}(k)$, $\varepsilon^{\mathrm{int}}(k)$ and
$\varepsilon^{\mathrm{ext}}(k)=\varepsilon^{\mathrm{ext},\sigma}(k)
+\varepsilon^{\mathrm{ext},\mathsf{t}}(k)$. Since $\varepsilon^{\mathrm{ext},\mathsf{t}}(k)$ is
$(1+C_\Sigma C_e)^3\,e\,S_d$ times the three $\theta'$-members taken at the weight $a+7$ in place
of $a$ and with the constant of the first of them replaced by $\mathsf{C}_1^{\mathrm{ex}}$, minus
the three $\theta'$-members, the three $\theta'$-members and
$\varepsilon^{\mathrm{ext},\mathsf{t}}(k)$ together are at most $\frac12$ by the corresponding
ceiling of $\gamma_J$; $\varepsilon^{\mathrm{rel}}(k)\le1$; and the remaining two are at most
$\frac12$. Summing, $\varepsilon_{\mathrm{br}}(k)\le2$.

\emph{Clause (e)} is part of the statement on the interior ceiling majorant, taken here by
statement: the constants of (a)--(d) are those of hypotheses (ii)--(vii), of the piles and of the
majorant $\mathsf{P}_{\mathrm{int}}$, whose coefficients are fixed before $\gamma$; the quantities
that vary along the sequence of renormalization steps enter only through $\check{R}_k$, $N$ and
$\check{R}_{h+1}$, and the last two through
the displayed bounds of the resummation of wrapped units on the number of preparatory stages and
on the radii.

\emph{Clause (f)} is Proposition~\ref{5.2:prop-carrier-value-bound}. \emph{Clause (g)} is the
exterior value bound on the carrier, in which the containment of the $\mathbf{J}$-entry of
$S_v=\mathfrak{r}_v$ for a compact $p$ is the containment lemma for the $\mathbf{J}$-entry.
\end{proof}

\begin{corollary}[The volume sum of one unit's pieces]\label{5.2:cor-unit-volume-sum}
Assume the hypotheses of the lemma on plain interior units at the resummation site
(Lemma~\ref{5.2:def-plain-interior-units}), and assume in addition:
\begin{itemize}
\item[(a)] the summation estimate over one link family of a piece, under the condition
$\kappa_1-1-3c_l\ge\lambda_*$; this condition is implied by the floor
$\kappa_1-1\ge c_l(a+9)+\lambda_*$ on $\kappa_1$ among the hypotheses of
Lemma~\ref{5.2:def-plain-interior-units}, since $a+9\ge 3$ and $c_l>0$;
\item[(b)] the comparison of sizes: away from the window cells defined below, the table size
$\mathfrak{N}^{\flat}$ exceeds the size $\mathfrak{N}$ by at most a factor that is a number fixed
before $\gamma$;
\item[(c)] the absorption of the window weights by the count of old cells: a window cell $\Delta$
(defined below) carrying the weight $w(\Delta)=\mathsf{L}_0^{2e}$ contains at least $L^{de}$ cells
of the old determining set $\{\Gamma^0_j\}$ if it lies in the window of the running attempt, and at
least $L^{de}$ pre-arrest cells if it lies in the frozen window of a live arrest.
\end{itemize}
Let $v$ be a plain interior first-part unit at the resummation site, with $K_v=[A_v]_k$ the family
of $\pi_k$-cubes meeting its label $A_v$, and with true domain $X_v\subset C$ for a component $C$ of
$Z$, and let $\mathfrak{p}=(v;\mathsf{y}_1,\dots,\mathsf{y}_4)$
run over its pieces. Call \emph{window cells} the cells of level $>k_0$ inside the component of the
running attempt (these lie outside the core $\mathsf{C}$), and the cells of level $>k_{0,\mathfrak{a}}$
inside $X_{\mathfrak{a}}$ for a live arrest $\mathfrak{a}$ (its frozen window, inside $\mathsf{C}$).
Here $k_0$, resp.\ $k_{0,\mathfrak{a}}$, is the base level of the window in the weights
\eqref{5.2:eq-plain-interior-units-a}, so that a weighted cell $\Delta$ carries
$w(\Delta)=\mathsf{L}_0^{2e}$ with its exponent $e$ fixed there ($e$ is here an integer exponent,
not to be confused with the base of the exponential in $e^{4d_k(K_v)+4}$ and $e^4$). Then:
\begin{enumerate}
\item[(i)] the pieces of $v$ obey
\begin{equation}\label{5.2:eq-unit-volume-sum-1}
\sum_{\mathfrak{p}}|\mathfrak{p}|\le\mathfrak{N}^{\flat}(v)\,e^{4d_k(K_v)+4};
\end{equation}
\item[(ii)] the pieces that stay in $C$ are contained in $e^4$ times the volume bound
\eqref{5.2:eq-resummation-input-e}, that is, under the per-cell constants
$\mathsf{K}_{\mathrm{cell}}(k,n)$ and the frozen piles $\mathsf{P}(\mathfrak{a})$, where the
constant $\mathsf{K}_{\mathrm{cell}}^0$ contains the summand $27\mathsf{L}_0^6\mathsf{K}^{\flat}$;
\item[(iii)] for every live arrest $\mathfrak{a}$ with $X_{\mathfrak{a}}\subset C$ the weighted
pointed mass of its frozen window satisfies
\begin{equation}\label{5.2:eq-unit-volume-sum-2}
\sum_{\Delta'}\mathsf{K}^{\flat}\bar{\mathsf{w}}(\Delta')^3
\le 27\mathsf{L}_0^6L^d\mathsf{c}_d\mathsf{K}^{\flat}
\sum_{n\le n^+_{\mathfrak{a}}}\check{R}^d_{j_n+1}
\sum_{X}\bigl(d'_{j_n+1}(X)+1\bigr),
\end{equation}
the sum on the left over the frozen window cells $\Delta'\subset X_{\mathfrak{a}}$ with
$w(\Delta')>1$, and the inner sum on the right over the components $X$ of $Z_{j_n+1}$ contained in
$X_{\mathfrak{a}}$, $j_n$ the
level integrated at the stage $n$ of $\mathfrak{a}$; this mass is a member of $\mathsf{P}(\mathfrak{a})$,
and the constant $c_{\mathsf{P}}$ of \eqref{5.2:eq-arrested-regions-a} contains the summand
$27\mathsf{L}_0^6L^d\mathsf{c}_d\mathsf{K}^{\flat}$.
\end{enumerate}
In the ceiling \eqref{5.2:eq-per-cell-allowance-a} the bracket, and in
\eqref{5.2:eq-arrested-regions-a} the constant $c_{\mathsf{P}}$, are multiplied by $1+e^4$, the
factor $e^4$ being that of \eqref{5.2:eq-unit-volume-sum-1}.
\end{corollary}

\begin{proof}
\emph{The sum.} Apply hypothesis (a) four times, once to each of the link families
$\mathsf{y}_4,\mathsf{y}_3,\mathsf{y}_2,\mathsf{y}_1$, under $\kappa_1-1-3c_l\ge\lambda_*$, to the
pieces of $v$ as bounded in the lemma on plain interior units at the resummation site. The four
summations together bound $\sum_{\mathfrak{p}}|\mathfrak{p}|$ by $\mathfrak{N}^{\flat}(v)$ times the
factor $e^{4d_k(K_v)+4}$, which is
\eqref{5.2:eq-unit-volume-sum-1}.

\emph{The weight.} The factor $e^{4d_k(K_v)}$ is contained in the weights carried by the column sum
of pointed sizes \eqref{5.2:eq-column-sum-pointed-a}, by the bound on stored boundary-type terms, and
by the anchored-domain norm $\|\cdot\|^{A}$. Hence the pieces of $v$ that stay in $C$ lie under
$e^4$ times the right member of \eqref{5.2:eq-resummation-input-e}, with its
$\mathsf{K}_{\mathrm{cell}}$ and $\mathsf{P}(\mathfrak{a})$. By hypothesis (b),
$\mathfrak{N}^{\flat}$ is at most a number fixed before $\gamma$ times the size $\mathfrak{N}$ except
on the window cells; there, by hypothesis (c), the weight is absorbed by the count of old cells. We
recall that count with its proof in both windows, since its branches fix the exponent $e$ used
below.

\emph{The running window.} In a window the preparation coarsens the label:
``the set $\Gamma_h$ is replaced by the two sets
$\{\Gamma_h\setminus Z''_{h+1},\ \Gamma_h\cap Z''_{h+1}\}$. The first set is combined with
$\Gamma_{h+1}$'' \cite[p.~179]{B-CMP122a}. Consequently a current cell $\Delta$ of the running
window carrying the weight $w(\Delta)=\mathsf{L}_0^{2e}$ contains at least $L^{de}$ cells of the old
determining set $\{\Gamma^0_j\}$, which \cite[(1.69)]{B-CMP122b} and the charge of the running chain
in the proof of the frozen contribution of live arrested regions (\ref{5.2:prf-arrested-regions})
count. This holds in both branches of the weights
\eqref{5.2:eq-plain-interior-units-a}.
\begin{itemize}
\item[(1)] $\Delta$ is a stratum cell of level $m\in\,]k_0,k[$, with $e=m-k_0$, or a top-zone cell on
$\Omega''_k\setminus\Omega_{k_0+1}$, with $e=N_0$. Under such a cell the old label has levels at
most $k_0$, since $Z''_k\subset\Omega^c_{k_0+1}$ and these weights are defined on
$\Omega^c_{k_0+1}$ only; a cell of level $m$ contains $L^{d(m-j)}$ cells of any level $j\le m$, so
$\Delta$ contains at least $L^{d(m-k_0)}=L^{de}$ old cells.
\item[(2)] $\Delta$ is a level-$k$ cell on the old ring of $\Omega''_k$ of level
$m'\in\,]k_0,k]$, with $e=k-m'$. Then $\Delta$ contains exactly $L^{d(k-m')}=L^{de}$ old cells,
one when $m'=k$.
\end{itemize}
On the other hand, the threshold weight of the pointed rows satisfies $\bar{\mathsf{w}}\le 3w^{\vee}$,
and $w^{\vee}$ raises the exponent $e$ by at most one, so that, using
$\mathsf{L}_0^2\le L/3$ and $d=4>3$,
\begin{equation}\label{5.2:eq-unit-volume-sum-3}
\bar{\mathsf{w}}^3\le(3w^{\vee})^3\le 27\mathsf{L}_0^{6(e+1)}
\le 27\mathsf{L}_0^6\Bigl(\frac{L}{3}\Bigr)^{3e}\le 27\mathsf{L}_0^6L^{de}.
\end{equation}
Since the per-cell pile of the pointed rows is $\mathsf{K}^{\flat}\bar{\mathsf{w}}^3$ by
\eqref{5.2:eq-column-sum-pointed-a}, the two counts give, per old cell, a cost of at most
$27\mathsf{L}_0^6\mathsf{K}^{\flat}$. As $\mathsf{L}_0<\tfrac12L$ is a number fixed before $\gamma$,
this cost is the summand $27\mathsf{L}_0^6\mathsf{K}^{\flat}$ added to
$\mathsf{K}_{\mathrm{cell}}^0$ in the bracket of the ceiling \eqref{5.2:eq-per-cell-allowance-a}.
This proves (ii).

\emph{A frozen window.} On the frozen window of a live arrest $\mathfrak{a}$ the current label
coincides with the frozen label, so the old determining set of the present step does not refine it;
the same inequality is therefore run on the pre-arrest cells, and the mass is placed in
$\mathsf{P}(\mathfrak{a})$ by the lemma on the weighted mass of frozen cells
(Lemma~\ref{5.2:lem-frozen-cell-mass}, display \eqref{5.2:eq-frozen-cell-mass-a}). A frozen weighted
cell $\Delta'$ of level $n'$ and exponent $e$ contains at least $L^{de}$ pre-arrest cells: in
branch (1) these have levels at most $k_{0,\mathfrak{a}}$ under the strata and the top shell; in
branch (2) they have level $m'$ under the ring $m'$. By \eqref{5.2:eq-unit-volume-sum-3},
\begin{equation*}
\mathsf{K}^{\flat}\bar{\mathsf{w}}(\Delta')^3
\le 27\mathsf{L}_0^6\mathsf{K}^{\flat}\cdot\#\{\text{pre-arrest cells in }\Delta'\}.
\end{equation*}
A pre-arrest cell of level $j$ lies off $\Omega_{j+1}$, hence in
$\Omega^c_{j+1}\subset\Lambda^c_{j+1}$; those under the window of $X_{\mathfrak{a}}$ lie in
$Z_{j+1}\cap X_{\mathfrak{a}}$. Writing $\#_{j+1}(S)$ for the number of level-$(j+1)$ cubes meeting
a set $S$, and using that a level-$(j+1)$ cube contains $L^d$ level-$j$ cubes, their number is
\begin{equation*}
\le L^d\,\#_{j+1}(Z_{j+1}\cap X_{\mathfrak{a}})
\le L^d\check{R}^d_{j+1}\mathsf{c}_d\sum_{X}\bigl(d'_{j+1}(X)+1\bigr),
\end{equation*}
the sum over the components $X$ of $Z_{j+1}$ in $X_{\mathfrak{a}}$; the second inequality is the
count in the proof of the frozen contribution of live arrested regions
(\ref{5.2:prf-arrested-regions}) at the stage $n$ with $j_n+1=j+1$. That stage was run: a
level $j$ not re-made under the window carries no weight. Summing over $n\le n^+_{\mathfrak{a}}$
gives \eqref{5.2:eq-unit-volume-sum-2}. This mass is added to $\mathsf{P}(\mathfrak{a})$ on the
right of \eqref{5.2:eq-resummation-input-e}, which is why $c_{\mathsf{P}}$ in
\eqref{5.2:eq-arrested-regions-a} contains $27\mathsf{L}_0^6L^d\mathsf{c}_d\mathsf{K}^{\flat}$; this
is a number fixed after $M$, and it involves no age, no number of zones of the arrest and no count
of arrests. This proves (iii). One could instead bound $\bar{\mathsf{w}}\le 3L^{N_0}$ and use
$27N_0L^{3N_0}\le L^{4N_0}$, which requires a lower bound on $L$; that route is not used.

\emph{The factor $e^4$.} The factor $e^4$ in \eqref{5.2:eq-unit-volume-sum-1} is accounted for by
multiplying the bracket of the ceiling \eqref{5.2:eq-per-cell-allowance-a}, and likewise
$c_{\mathsf{P}}$, by $1+e^4$.
\end{proof}

\begin{lemma}[The carrier stays within its prescribed domain]\label{5.2:lem-carrier-domain}
We work at the site (L) of $\mathbf{R}_k$, in the setting and notation of
Definition~\ref{5.2:def-plain-interior-units}. Thus $N := N(g_k)$ and $k_0 := k - N_0(g_k)$.
The set $C$ is a component of the treated class $Z$; it is a box, and
$\mathsf{C} = \mathsf{C}_C$ is its core. The set $\mathfrak{P}$ is the parameter polydisc.
For $m \in \mathfrak{P}$, $\operatorname{chain}(m)$ is the carrier pair at the background
$\mathfrak{b}_4$ and $\mathcal{U}(m)$ is its $U$-slot. We write $\bar U(m)$ for the $G$-valued
part of $\mathcal{U}(m)$, $\pi(y)$ for the level of $\mathbf{B}''_k$ at $y$, $\Delta_y$ for the
cell of $y$, and $\ell_n := L^n\eta$, with $\eta$ as in
Definition~\ref{5.2:def-plain-interior-units}. Here $G$ denotes the structure group, and
$\ell_n$ is the level-$n$ unit of length. Each clause of the
analyticity domains below is a list of
numbered bounds, its entries $1$ to $8$, numbered as in the linearisation table of the
three-field bound; for an entry $a$, $Q_a(\mathbf{U},\mathbf{J})(y)$ denotes the left member of
bound $a$, evaluated at the configuration $(\mathbf{U},\mathbf{J})$ and at the point $y$. We write
$U_2$ for the $G$-valued configuration on $\mathsf{C}$ fixed in the own-part bound for $U_2$
(hypothesis (d) below).

A cube $\square$ is a \emph{gauge cube of index} $\mathsf{m}$ if it has the following three
properties: it is a union of $2^d$ cubes of $\pi_{\mathsf{m}}$; $\square \subset
\Omega_{\mathsf{m}}$; and $\square$ meets zone $\mathsf{m}$. At
$\mathsf{m} = k$, the condition that $\square$ meets zone $k$ is read through the printed clause
that supplies the gauge. Here and below, \emph{old} refers to the sets of the label in force
before the attempt; the sets $\Omega_j$ without this word are those of the current label.
There are two cases.
\begin{itemize}
\item If $\square$ meets the old set $\Omega^c_{k_0+1}$, the gauge is supplied by
\cite[(1.65)]{B-CMP122a}, through its family $m' = j$, where $j = k$.
\item Otherwise, let $m'$ be the largest index with $\square$ contained in the old set
$\Omega_{m'}$. Then $\square$ meets the old set $\Omega^c_{m'+1}$, and the gauge is supplied by
\cite[(1.67)]{B-CMP122a} at that $m'$.
\end{itemize}

Put $\mathfrak{E}^C := \mathbb{C}^{(N)}_k(C)$. We assume the following four hypotheses.
\begin{enumerate}
\item[(a)] The membership part of the consumer-domain statement holds, taken with its
inheritance clauses and its arrest reading. This means three things. First, for every consumer
$\mathfrak{E} \in \mathcal{C}$ and every $m \in \mathfrak{P}$, the carrier $\operatorname{chain}(m)$
at $\mathfrak{b}_4$ lies in the domain $\mathfrak{D}_{\mathfrak{E}}$, and at the points of
$\Lambda^{\natural}$ its entries obey the thresholds of $\mathfrak{D}_{\mathfrak{E}}$ reduced by
the factor $1/16$. Second, for
$\mathfrak{E} = \mathfrak{E}^C$,
\begin{equation}\label{5.2:eq-carrier-domain-1}
\mathfrak{D}_{\mathfrak{E}^C} = \mathfrak{D}^{\flat}_N(C)
= \mathfrak{N}^{(N)} \cap \mathfrak{L}_N(C\cap\Omega_k^c) \cap \mathfrak{I}_N(C\cap\Omega_k).
\end{equation}
Here $\mathfrak{L}_N(Y)$ is the set cut out on $Y$ by the lettered clauses
\cite[(1.64)--(1.69)]{B-CMP122a} at $n = N$; $\mathfrak{I}_N(Y)$ is the set cut out on $Y$ by the
clauses \cite[(2.34)--(2.39)]{B-CMP119}, at $n = j = k$, of the constituents
$\mathbf{B}^{(k)}(X_T)$ of $\mathfrak{E}^C$, indexed by their localisation domains
$X_T \subset Y$; and $\mathfrak{N}^{(N)}$ is the stage-$N$ factor fixed in the inheritance
statement. Third, on a live piece of an arrest $\mathfrak{a}$ the frozen unit length
$\ell_{\mathfrak{a}}$, weight $w_{\mathfrak{a}}$ and domain $\mathcal{Q}_{\mathfrak{a}}$ are read,
the frozen level being the current one.
\item[(b)] The sourced-family extension holds: the membership in (a) persists for the families
carrying the decoupling parameters $\mathsf{s}$, with $\mathsf{s}$ free on the cells of $X$ as
well, where $X$ denotes the localisation domain of the consumer $\mathfrak{E}$.
\item[(c)] The cube clause of the inheritance statement holds. Let $X$ be the localisation
domain of $\mathfrak{E}$. Off $\Lambda^{\natural}$, on every cube of $2^d$ cubes of $\pi_k$
contained in $X\setminus\Lambda^{\natural}$, the configurations of $\mathfrak{D}_{\mathfrak{E}}$
are those of the source tube itself, the source tube being the analyticity domain in which the
sourced configurations are taken; hence there the cube clause \cite[(1.65), (1.67)]{B-CMP122a}
of $\mathfrak{D}_{\mathfrak{E}}$, with its $U$-part and its cube gauges, is that of the source
tube.
\item[(d)] The own-part bound for $U_2$ holds: on $\mathsf{C}$ the $G$-part $\bar U(m)$ of
$\mathcal{U}(m)$ equals $U_2$, and $U_2$ satisfies there the bounds of entries $5$ and $8$ and
the clause \cite[(2.38)]{B-CMP119} as printed.
\end{enumerate}
Then the following two conclusions hold.
\begin{enumerate}
\item[(i)] $\mathfrak{E}^C \in \mathcal{C}$. At $y \in C$ the clause index of
$\mathfrak{D}_{\mathfrak{E}^C}$ is $\pi(y)$. The threshold of entry $a$ at $y$ is
\begin{equation*}
\begin{gathered}
\theta_a(y) = F\,w(\Delta_y)\,\alpha_{i,\pi(y)}\times(\text{the level-}\pi(y)\text{ unit of entry }a),
\\
i \in \{0,1\}, \qquad \varphi_{\flat} \le F \le 1 .
\end{gathered}
\end{equation*}
Here $w(\Delta_y) = \mathsf{L}_0^{2e}$, with $e = \max\{0, m'-k_0\}$ on a cell of clause (i) of
\cite[(1.64)--(1.67)]{B-CMP122a} indexed by the level $m'$, and $e = k-m'$ on an old ring of
$\Omega''_k$ of old level $m'$ carrying clause (ii); and $e \le N_0(\mathsf{a})$, $\mathsf{a}$
being the attempt whose window contains $\Delta_y$, the running attempt if there is none.
Moreover $w(\Delta_y) = 1$ unless $\mathsf{w}^{\mathrm{win}}(\Delta_y) > 1$. Put
\begin{equation*}
\begin{gathered}
\mathbf{r}_a(n) := \alpha_{*,n}\times(\text{the level-}n\text{ unit of entry }a),
\\
c_\alpha := (1+\varrho)C_0/C_*,
\end{gathered}
\end{equation*}
where $C_*$ is the constant fixed with the thresholds of the
three-field bound. Then
\begin{equation}\label{5.2:eq-carrier-domain-a}
\begin{split}
\varphi_{\flat}\,\mathbf{r}_a(\pi(y)) \le \theta_a(y)
&\le 3c_\alpha\, w(\Delta_y)\,\mathbf{r}_a(\pi(y))\\
&\le 3c_\alpha\, \mathsf{w}^{\mathrm{win}}(\Delta_y)\,\mathbf{r}_a(\pi(y)).
\end{split}
\end{equation}
\item[(ii)] For every $m \in \mathfrak{P}$, every $y \in C$ and every $a \in \{1,3,5,6,7\}$,
\begin{equation}\label{5.2:eq-carrier-domain-b}
\begin{gathered}
Q_a(\operatorname{chain}(m))(y) < \theta_a(y),
\\
Q_a(\operatorname{chain}(m))(y) \le \tfrac{1}{16}\,\theta_a(y)
\ \text{ on } \Lambda^{\natural} \supset \mathsf{C}.
\end{gathered}
\end{equation}
Moreover, on every gauge cube $\square$ of index $\mathsf{m}$, $\bar U(m)$ has a $G$-valued
gauge $u$ with $\bar U(m)^u = e^{i\eta A}$, where both $\ell_{\mathsf{m}}|A|$ and
$\ell_{\mathsf{m}}^2|\nabla A|$ are less than
$3\mathsf{w}^{\mathrm{win}}\cdot 2BM\alpha_{0,\mathsf{m}}$; and, by restriction, the same gauge
$u$, the same identity and the same two bounds hold on every box contained in such a cube.
\end{enumerate}
\end{lemma}

\begin{proof}
\emph{The auxiliary consumer.} By the definition of the consumer class $\mathcal{C}$, this class
contains the chain terms
$\mathbb{C}^{(n)}_k$ of $\mathbf{B}''_k$ with $n \le N$. At the last stage $n = N$, condition
\cite[(1.63)]{B-CMP122a} asks three things of a domain $X$:
\begin{equation*}
X \in \mathbf{D}_k, \qquad X\cap\Omega_k \neq \emptyset, \qquad
X\cap(Z_k\setminus Z''_k) \neq \emptyset .
\end{equation*}
The box $C$ has all three properties. Hence $\mathfrak{E}^C \in \mathcal{C}$, and hypothesis (a)
applies at $\mathfrak{E} := \mathfrak{E}^C$.

\emph{The domain.} By the inheritance clauses in (a), membership in
$\mathfrak{D}_{\mathfrak{E}^C}$ is membership in the set \eqref{5.2:eq-carrier-domain-1}. Its
meaning differs on the two parts of $C$.
\begin{itemize}
\item On $C\cap\Omega_k^c$ it is membership in the factor $\mathfrak{L}_N(C\cap\Omega_k^c)$, that
is, the lettered clauses \cite[(1.64)--(1.69)]{B-CMP122a}
at $n = N$ (so $j = k$). On live pieces these are read through the arrest reading of (a).
\item On the ring $C\cap\Omega_k$ it is membership in the factor
$\mathfrak{I}_N(C\cap\Omega_k)$, that is, the clauses of the constituents
$\mathbf{B}^{(k)}(X_T)$ of $\mathfrak{E}^C$, indexed by their localisation domains
$X_T \subset C$. These are \cite[(2.34)--(2.39)]{B-CMP119} at $n = j = k$.
\end{itemize}

\emph{The clause index.} Read \cite[(1.64)--(1.67), p.~190]{B-CMP122a} at $n = N$.
\begin{itemize}
\item Clause (i) is indexed by $(\Omega''_{m'}\setminus\Omega''_{m'+1})\cap X$ and stands in
level-$m'$ units.
\item Clause (ii) stands in level-$k$ units on the old rings of $\Omega''_k$.
\item Clause (iii) is empty at $j = k$.
\end{itemize}
Therefore the clause of $\mathfrak{D}_{\mathfrak{E}^C}$ acting at $y$ has index $\pi(y)$. Its
threshold for entry $a$ has the form stated in (i): the factor $F$ satisfies
$\varphi_{\flat} \le F \le 1$, and the weight is $w(\Delta_y) = \mathsf{L}_0^{2e}$. Here
$e = \max\{0,m'-k_0\}$ on a cell of clause (i) indexed by the level $m'$, and $e = k-m'$ on an
old ring of old level $m'$ carrying clause (ii), with $e \le N_0(\mathsf{a})$ for the attempt
$\mathsf{a}$ whose window contains $\Delta_y$ (the running one if there is none),
and $w(\Delta_y) = 1$ unless $\mathsf{w}^{\mathrm{win}}(\Delta_y) > 1$. On a live piece of an
arrest $\mathfrak{a}$, the unit length, the weight and the domain are the frozen $\ell_{\mathfrak{a}}$,
$w_{\mathfrak{a}}$, $\mathcal{Q}_{\mathfrak{a}}$, and by the arrest reading in (a) the frozen
level equals the current one.

\emph{The threshold bounds.} Four facts enter:
\begin{itemize}
\item the inequality $\mathsf{L}_0^2 \le L/3$ of \cite{B-CMP122a};
\item the bound $F\,w\,c \le 3w_*$ on the thresholds of the consumer domains, established with
the consumer-domain statement, for the product of the factor, the weight and the threshold
constant $c$;
\item the definition of $\mathbf{r}_a(n)$;
\item the definition of $c_\alpha$.
\end{itemize}
The lower inequality of \eqref{5.2:eq-carrier-domain-a} rests on $\varphi_{\flat} \le F$, on
$w(\Delta_y) = \mathsf{L}_0^{2e} \ge 1$ (as $\mathsf{L}_0 \ge 3$ and $e \ge 0$) and on the
definition of $\mathbf{r}_a$; the upper one rests on the bound $F\,w\,c \le 3w_*$ and on the
definition of $c_\alpha$. So the threshold satisfies the first two inequalities of
\eqref{5.2:eq-carrier-domain-a}; the middle member carries the domain's own weight
$w(\Delta_y) = \mathsf{L}_0^{2e}$.

The last member of \eqref{5.2:eq-carrier-domain-a} majorises the middle one, because
$w(\Delta_y) \le \mathsf{w}^{\mathrm{win}}(\Delta_y)$ on $C$, by
Definition~\ref{5.2:def-plain-interior-units} and \eqref{5.2:eq-plain-interior-units-a}; this
inequality rests on $\mathsf{L}_0^2 \le L/3$, which with $0 \le e \le N_0(\mathsf{a})$ gives
\begin{equation*}
\mathsf{L}_0^{2e} \le (L/3)^e \le L^{N_0(\mathsf{a})}.
\end{equation*}
In particular, if $\mathsf{w}^{\mathrm{win}}(\Delta_y) = 1$, then $w(\Delta_y) \le 1$, while
$w(\Delta_y) = \mathsf{L}_0^{2e} \ge 1$; so $w(\Delta_y) = 1$ unless
$\mathsf{w}^{\mathrm{win}}(\Delta_y) > 1$. This proves (i).

\emph{The entry bounds.} By (a) at the consumer $\mathfrak{E}^C$ and the background
$\mathfrak{b}_4$, $\operatorname{chain}(m)$ lies in $\mathfrak{D}_{\mathfrak{E}^C}$ for every
$m \in \mathfrak{P}$. Its entries $a \in \{1,3,5,6,7\}$ therefore obey at each $y \in C$ the
thresholds $\theta_a(y)$ identified above. On $\Lambda^{\natural} \supset \mathsf{C}$ they obey
$\theta_a(y)/16$, again by (a). For the families carrying the decoupling parameters $\mathsf{s}$,
the same holds by (b), with $\mathsf{s}$ free on the cells of $X$ as well. This gives
\eqref{5.2:eq-carrier-domain-b}. Hypothesis (a), taken at $\mathfrak{E} := \mathfrak{E}^C$ and at
the background $\mathfrak{b}_4$, is the same statement as the claim, made in the proof of the
correlation theorem, that on $X\cap Z^{\mathrm{on}}$, where $X$ is the localisation domain of the
consumer, the carriers obey the lettered clauses $\mathfrak{L}_n$.

\emph{The cube clause.} On $C$ the gauge statement on gauge cubes is the cube clause of the
domain \eqref{5.2:eq-carrier-domain-1}, in which $\operatorname{chain}(m)$ lies by (a): the
clauses \cite[(1.65), (1.67)]{B-CMP122a} on $C\cap\Omega_k^c$, and the clause
\cite[(2.38)]{B-CMP119} on the ring $C\cap\Omega_k$. Two refinements complete it.
\begin{itemize}
\item Off $\Lambda^{\natural}$, hypothesis (c) guarantees that on every cube of $2^d$ cubes of
$\pi_k$ the configurations of the domain are those of the source tube itself, so that the cube
clause there is the source tube's own, with its $U$-part and its cube gauges.
\item On $\mathsf{C}$, hypothesis (d) gives $\bar U(m) = U_2$, and $U_2$ obeys
\cite[(2.38)]{B-CMP119} as printed, again by (d).
\end{itemize}
This gives the gauge $u$, the identity $\bar U(m)^u = e^{i\eta A}$ and the bounds on
$\ell_{\mathsf{m}}|A|$ and $\ell_{\mathsf{m}}^2|\nabla A|$. The gauge $u$ obtained on a gauge
cube $\square$ serves, by restriction, every box contained in $\square$: the identity
$\bar U(m)^u = e^{i\eta A}$ and the two bounds on $A$ hold pointwise on $\square$, so they hold on
every sub-box. This proves (ii).
\end{proof}

\begin{remark}
In the definition of a gauge cube, the number $2^d$ comes from the cube clauses
\cite[(1.65), (1.67)]{B-CMP122a} and \cite[(2.38)]{B-CMP119}: the constant multiplying the cube
size in these clauses is at most $2$ as printed, and we take it equal to $2$.
In Ba{\l}aban's notation, conclusion (ii) of Lemma~\ref{5.2:lem-carrier-domain} is the statement
\cite[(1.65)]{B-CMP122b}, not to be confused with the cube clause \cite[(1.65)]{B-CMP122a}. That
statement says that the extended function $U''_k$ is analytic in $(\mathbf{U},\mathbf{J})$ on
$\tilde U^c_k(Y_4,\tilde\alpha_0,\tilde\alpha_1)$. Its values lie in
$\tilde U''^{\,c}_k(Y_1,\tilde\alpha_0,\tilde\alpha_1)$, the space defined by
\cite[(1.64)--(1.69)]{B-CMP122a} for $n = N$; see \cite[p.~374]{B-CMP122b}. That statement is
not used as an input in the proof above.
\end{remark}

Throughout what follows $\pi_j$ is Ba{\l}aban's $L$-adic partition into cubes of level $j$, and a
\emph{level-$j$ unit} is the side of a cube of $\pi_j$, so that one level-$(j-1)$ unit is $1/L$
level-$j$ units. We write $\operatorname{dist}_j$ for distances measured in level-$j$ units.
$\Omega_j$ are the small-field regions. The \emph{zone} $j$ is $\Omega_j\setminus\Omega_{j+1}$.
$C$ is the box of Lemma~\ref{5.2:lem-carrier-domain}. The \emph{cells of level $j$} are the cubes
of $\pi_j$ contained in zone $j$; for a point $y$, $\Delta_y$ is the cell containing $y$.

We use the gauge cubes of index $\mathsf{m}$ defined with Lemma~\ref{5.2:lem-carrier-domain}. Such
a cube is a union of $2^d$ cubes of $\pi_{\mathsf{m}}$, is contained in $\Omega_{\mathsf{m}}$, and
meets zone $\mathsf{m}$, the last condition being read at $\mathsf{m}=k$ as it is fixed there. Every
box contained in a gauge cube is given the restriction of its gauge.

\begin{lemma}[Regular cubes around a short box]\label{5.2:lem-short-box-cubes}
Let $S\subset C$ be a box whose sides do not exceed the side of a cube of $\pi_{n-1}$ (one
level-$(n-1)$ unit), and let $S$ contain a point $y$ of zone $n$. Then:
\begin{itemize}
\item $S$ lies in a gauge cube of index $n-1$ or $n$, whose cells of that index are at
$\operatorname{dist}_n<1$ from $\Delta_y$;
\item $S$ meets cells of the levels $n-1$, $n$, $n+1$ only, and all of them are at
$\operatorname{dist}_n<1$ from $\Delta_y$.
\end{itemize}
\end{lemma}

\begin{proof}
Along each axis the projection of $S$ is an interval of length at most one level-$(n-1)$ unit.
Such an interval lies in the union of two consecutive cubes of $\pi_{n-1}$ along that axis.
Hence $S$ lies in a cube $\square'$ which is a union of $2^d$ cubes of $\pi_{n-1}$. The side of
$\square'$ is two level-$(n-1)$ units, that is $2/L$ level-$n$ units. Since $y\in S\subset\square'$,
every point of $\square'$ lies within $2/L<1$ level-$n$ unit of $y$, and
$y\in\Omega_n\setminus\Omega_{n+1}$. By the separation property of the nested small-field
regions, in the form used here (a point at $\operatorname{dist}_n<1$ from a point of zone $n$ lies
in $\Omega_{n-1}$ and not in $\Omega_{n+2}$), applied to the points of $\square'$, we obtain
\begin{equation*}
\square'\subset\Omega_{n-1},\qquad \square'\cap\Omega_{n+2}=\emptyset .
\end{equation*}

First, suppose $\square'$ meets zone $n-1$. Then $\square'$ is a union of $2^d$ cubes of
$\pi_{n-1}$, is contained in $\Omega_{n-1}$ and meets zone $n-1$. So it is a gauge cube of
index $n-1$ containing $S$.

Otherwise $\square'\subset\Omega_n$. Consider the cubes of $\pi_n$ that meet $\square'$. They form
a box $B\subset\Omega_n$, a union of cubes of $\pi_n$. Along each axis the side of $\square'$ is
$2/L<1$ level-$n$ unit, so $B$ has sides $1$ or $2$ in level-$n$ units.

Let $e$ be an axis along which $B$ has side $1$. We show that one of the two neighbouring slabs
$B+e$, $B-e$ (the translates of $B$ by one level-$n$ unit along $e$) lies inside $\Omega_n$, so
that the side of $B$ along $e$ can be doubled inside $\Omega_n$. Suppose instead that both slabs
meet $\Omega_n^c$. Two cases arise, according to the form of $\Omega_n^c$ near $C$: either
$\Omega_n^c$ is a union of cubes of its own partition, of side $\check{R}_n$ level-$n$ units, or
$\Omega_n^c\cap C=Z''_n$ is a single box, as happens for the regions connected with the last $N$
steps, $N$ being here Ba{\l}aban's number of preparatory steps of the large-field operation
\cite[pp.~177--179]{B-CMP122a} (not the rank of the gauge group).
\begin{itemize}
\item In the first case, choose cubes $Q_+$ and $Q_-$ of the partition of $\Omega_n^c$ meeting
$B+e$ and $B-e$ respectively; neither meets $B$, since $B\subset\Omega_n$. Measure the $e$-axis in
level-$n$ units, write $[b,b+1]$ for the interval of $B$ along $e$, and $I_\pm$ for the intervals
of $Q_\pm$ along $e$. The transverse intervals of $B\pm e$ are those of $B$, so the transverse
intervals of $Q_\pm$ meet those of $B$; since a box is a product of intervals, $Q_\pm$ would meet
$B$ if $I_\pm$ met $[b,b+1]$. Hence $I_+$ meets $[b+1,b+2]$ and lies to the right of $b+1$, while
$I_-$ meets $[b-1,b]$ and lies to the left of $b$. So $I_+$ and $I_-$ are distinct intervals of
the partition of the $e$-axis into intervals of length $\check{R}_n$, at a distance greater than
$1$ and at most $3$. The cubes of that partition have side $\check{R}_n\ge 38$ level-$n$ units
\cite[pp.~256--257]{B-CMP119}, \cite[pp.~177--179]{B-CMP122a}. Distinct intervals of the
partition of the $e$-axis into intervals of length $\check{R}_n$ are at distance $0$ or at least
$\check{R}_n$, so a distance greater than $1$ and at most $3$ is excluded.
\item In the second case, $\Omega_n^c\cap C=Z''_n$ is a box, since ``all the regions connected
with the last $N$ steps are rectangular parallelepipeds'' \cite[pp.~177--179]{B-CMP122a}. This
box would meet both $B+e$ and $B-e$ and miss $B$. That is impossible. A box is a product of
intervals. Meeting $B+e$ and $B-e$ forces each transverse interval to meet the transverse
interval of $B$. It also forces the interval along $e$ to contain points on both sides of the
interval of $B$ along $e$, and hence to contain that interval. So the product meets $B$.
\end{itemize}
Thus some neighbouring slab lies in $\Omega_n$. Adjoining it replaces $B$ by a box in $\Omega_n$
that is a union of cubes of $\pi_n$ and has side $2$ along $e$ and unchanged sides along the
other axes.

Proceeding axis by axis, we obtain a cube of side $2$ level-$n$ units in $\Omega_n$. It is a union
of $2^d$ cubes of $\pi_n$ and contains $B\supset\square'\supset S\ni y$. Since $y$ lies in zone
$n$, this cube meets zone $n$, so it is a gauge cube of index $n$.

It remains to check the distances in the two cases.
\begin{itemize}
\item Index $n-1$: every cell of $\square'$ lies within $2/L<1$ level-$n$ unit of $y$.
\item Index $n$: $\Delta_y$ is the cube of $\pi_n$ containing $y$; it meets $\square'$, so it is
one of the cubes forming $B$, and hence one of the $2^d$ cubes of the gauge cube obtained. Every
cube of $\pi_n$ in that cube of side $2$ is equal or adjacent to $\Delta_y$, hence at
$\operatorname{dist}_n=0$ from $\Delta_y$.
\end{itemize}
Since $y\in\Delta_y$, in both cases the cells of the index in question are at
$\operatorname{dist}_n<1$ from $\Delta_y$.

Finally, $S\subset\square'$, $\square'\subset\Omega_{n-1}$ and $\square'\cap\Omega_{n+2}=\emptyset$.
Hence $S$ lies in the union of the zones $n-1$, $n$, $n+1$. Every cell that $S$ meets contains a
point of $S$, which lies within $2/L<1$ level-$n$ unit of $y$. So every such cell is at
$\operatorname{dist}_n<1$ from $\Delta_y$.
\end{proof}

\begin{corollary}[The decoupled carrier stays admissible on the core]
\label{5.2:cor-decoupled-carrier}
Let $\mathsf{C}$ be the core of the component $C$, let $S\subset\mathsf{C}$, let $m,m'\in\mathfrak{P}$
have the same pair $(\zeta,\varpi)$, so that they differ only in the decoupling parameters
$\vec{\mathsf{s}}$, and let $P_{\hat\lambda}$, $|\hat\lambda|\le\rho_S$, be the decoupled carrier
formed from $\operatorname{chain}(m)$ and $\operatorname{chain}(m')$, with disc variable
$\hat\lambda$. Assume the following.
\begin{enumerate}
\item[(a)] Lemma~\ref{5.2:lem-carrier-domain}, the carrier stays within its prescribed domain, in
particular the bound \eqref{5.2:eq-carrier-domain-b} and the two-sided bound
\eqref{5.2:eq-carrier-domain-a} for the thresholds $\theta_a(y)$.
\item[(b)] The statement that on the core $\mathsf{C}$ every carrier has the same $G$-valued part,
the fixed configuration $U_2$ of the core ($G=SU(N)$; the $G$-part as in
Lemma~\ref{5.2:lem-carrier-domain}), and that $U_2$ obeys there the bare bounds of the entries $5$
and $8$.
\item[(c)] The cell-inclusion statement for the core: for $y\in S$ one has
$\tilde\Delta(y)\subset D_2$ and $\Delta(y)\subset\mathsf{C}$, where $\Delta(y)$ is the cell of $y$,
$\tilde\Delta(y)$ its enlargement, and $D_2$ the region fixed with the core $\mathsf{C}$ in the
setting of this section.
\item[(d)] The disc lemma for the decoupled carrier, the estimate accompanying it, and the first
case of the half-rate bound.
\item[(e)] The elementary half of the linearisation table (its entries $1$, $3$ and $5$ to $8$).
\item[(f)] The age inequality at $b=0$.
\item[(g)] The amplitude bound, together with the choice of the base radius $r^0$ (so that
$\rho_S=r^0e^{\frac12\delta_P d^{\wedge}(S,\mathsf{C}^c)}$).
\end{enumerate}
Then $P_{\hat\lambda}|_S$ has the same $G$-part $U_2$, and for $a\in\{1,3,6,7\}$ and $y\in S$
\begin{equation}\label{5.2:eq-decoupled-carrier-1}
Q_a(P_{\hat\lambda})(y)\;\le\;\tfrac{1}{16}\,\theta_a(y)+\tfrac12\,\varphi_{\flat}\,
\mathbf{r}_a(\pi(y))\;<\;\theta_a(y).
\end{equation}
\end{corollary}

\begin{proof}
The argument is the proof of the main carrier bound, taken at the level $n:=\pi(y)$ on $S$. That
proof uses the auxiliary consumer $\mathfrak{E}^C$ of Lemma~\ref{5.2:lem-carrier-domain} only
through its thresholds $\theta_a$, its base bound \eqref{5.2:eq-carrier-domain-b} and the inequality
$n(y)\le\pi(y)$, where $n(y)$ is the level at which that proof evaluates the entries at $y$; and it
depends on the consumer's localisation region (there a general region, here the box $C$) only
through the domain of the disc lemma of hypothesis (d).
We go through its steps in order.

\emph{Base.} The base of the comparison is Lemma~\ref{5.2:lem-carrier-domain} at $m'$: by
\eqref{5.2:eq-carrier-domain-b}, for every $a\in\{1,3,5,6,7\}$ the quantity $Q_a(\operatorname{chain}(m'))(y)$
is at most $\theta_a(y)/16$ on $\Lambda^{\natural}$. Since $S\subset\mathsf{C}\subset\Lambda^{\natural}$,
this holds at every $y\in S$, in particular for $a\in\{1,3,6,7\}$.

\emph{The increment majorant.} Let $y\in S$. By hypothesis (c), $\tilde\Delta(y)\subset D_2$ and
$\Delta(y)\subset\mathsf{C}$, so $y$ lies in the domain of the disc lemma. Hypothesis (d) --- the disc
lemma and its estimate, read through the first case of the half-rate bound --- gives
\begin{equation}\label{5.2:eq-decoupled-carrier-2}
|\hat\lambda|\,\mathsf{N}_y(\mathcal{K}),\ \ |\hat\lambda|\,\ell^3|\jmath|\;\le\;\varkappa(y),
\qquad
\varkappa(y):=r^0\,\mathsf{U}\,e^{\delta_P\mathsf{w}_{\boxplus}}\,
e^{-\frac12\delta_P d^{\wedge}(y,\mathsf{C}^c)} .
\end{equation}
Here $\mathsf{N}_y(\mathcal{K})$ is the cell norm at $y$ of the quantity $\mathcal{K}$ estimated in
the disc lemma, $\jmath$ is the quantity estimated there with the weight $\ell^3$, $\ell$ being
the length unit of that lemma, $\mathsf{U}$ is the layer majorant of the disc lemma,
$\mathsf{w}_{\boxplus}$ is the weight of the disc lemma, and the
notation $\varkappa(y)$ is local to this proof.
Only the first case of the half-rate bound is used: no ball $B(x)$ and no fattening of the cell enter,
because the entries $2$ and $4$ do not occur in this corollary.

\emph{Linearisation.} Apply the elementary half of the linearisation table, hypothesis (e), at
$n(y):=\pi(y)$. The entries $5$ and $8$ are unmoved; the entries $1$, $6$ and $7$ move by at most
$3\varkappa(y)$ times the entry's unit; the entry $3$ moves by at most $|\hat\lambda\,\jmath|$.

\emph{Age.} The age inequality at $b=0$, hypothesis (f), gives
\begin{equation}\label{5.2:eq-decoupled-carrier-3}
e^{-\frac12\delta_P d^{\wedge}(y,\mathsf{C}^c)}\;\frac{\alpha_{*,k}}{\alpha_{*,\pi(y)}}
\;\le\;C_{\mathrm{age}} ,
\end{equation}
where $C_{\mathrm{age}}$ is the constant of the age inequality.

\emph{Size of the increments.} By the amplitude bound and the choice of $r^0$, hypothesis (g), each
of the increments just listed, bounded through \eqref{5.2:eq-decoupled-carrier-2} and
\eqref{5.2:eq-decoupled-carrier-3}, is at most $\tfrac12\varphi_{\flat}\,\mathbf{r}_a(\pi(y))$ for
the entry $a$ at $y$, where $\mathbf{r}_a(n)$ is $\alpha_{*,n}$ times the level-$n$ unit of the
entry $a$.

\emph{Conclusion.} The linearisation step leaves the entries $5$ and $8$ unmoved. The assertion that
$P_{\hat\lambda}|_S$ has the same $G$-part $U_2$ is the one supplied by hypothesis (b), the statement
that the $G$-part of the carriers on $\mathsf{C}\supset S$ is $U_2$.
For $a\in\{1,3,6,7\}$ and $y\in S$, the base bound and the size of the increments give the first
inequality in \eqref{5.2:eq-decoupled-carrier-1}. By the lower half of
\eqref{5.2:eq-carrier-domain-a}, $\varphi_{\flat}\mathbf{r}_a(\pi(y))\le\theta_a(y)$, so the right
member of the first inequality is at most $\tfrac{9}{16}\theta_a(y)$, which is less than $\theta_a(y)$
because $\theta_a(y)\ge\varphi_{\flat}\mathbf{r}_a(\pi(y))>0$.
\end{proof}

\begin{lemma}[Regularity of the carrier on a short box]\label{5.2:lem-carrier-regularity}
Put $c_{\mathrm{reg}} := 2L^2 c_\Psi c_\alpha(1+2BM)$, where $c_\alpha$ is the constant of the two-sided
bound \eqref{5.2:eq-carrier-domain-a}, $B$ is the constant of the cube clause in
\eqref{5.2:eq-carrier-domain-b}, and $c_\Psi$ is a constant of the construction, independent of the
couplings, which enters this lemma only through $c_{\mathrm{reg}}$. Let $S\subset C$ be a box whose sides
are at most one cube of
$\pi_{n-1}$ and which contains a point $y$ of zone $n$, as in Lemma~\ref{5.2:lem-short-box-cubes}, and put
$\ell := \ell_{n-1}$. Assume the condition in the first line of \eqref{5.2:eq-carrier-regularity-a}, and
define $\mathfrak{a}'$ by its second line:
\begin{equation}\label{5.2:eq-carrier-regularity-a}
\begin{gathered}
\sup_{x\ge\gamma^{-2}}\,360\,d\,L^2M\,c_{\mathrm{reg}}\,\mathsf{R}^{\uparrow}(\log x)\,\alpha_*(x)\le 1,\\
\mathfrak{a}' := 3c_{\mathrm{reg}}\,\mathsf{w}^{\mathrm{win},\vee}(\Delta_y)\,\alpha_{*,n},
\end{gathered}
\end{equation}
with $\mathsf{w}^{\mathrm{win},\vee}$ the window weight of \eqref{5.2:eq-plain-interior-units-a}. Then
$60\,d\,L\,M\,\mathfrak{a}'\le 1$, and:
\begin{enumerate}
\item[(i)] for every $m\in\mathfrak{P}$ the restriction $\operatorname{chain}(m)\restriction S$ reads, in
one $G^{\mathbb{C}}$-gauge ($G^{\mathbb{C}}$ denoting the complexification of $G=SU(N)$),
$(e^{i\eta\mathcal{A}^0},\mathbf{J}^0)$, where
\begin{equation*}
\ell\,|\mathcal{A}^0|,\quad \ell^2\,|\nabla\mathcal{A}^0|,\quad \ell^3\,|\mathbf{J}^0|\ \le\ \mathfrak{a}',
\qquad
|\partial\mathcal{U}(m)-1|\le \mathfrak{a}'\,(\eta/\ell)^2 ;
\end{equation*}
\item[(ii)] if $S\subset\mathsf{C}$, the same holds for $P_{\hat\lambda}\restriction S$ whenever
$|\hat\lambda|\le\rho_S$.
\end{enumerate}
\end{lemma}

The first line of \eqref{5.2:eq-carrier-regularity-a} is a condition on $\gamma$ through the range
$x\ge\gamma^{-2}$ of the supremum; it has the form of a supremum, its supremand has negative degree in
$\log x$, and it is one of the conditions that define the threshold $\gamma_J$.

\begin{proof}
\emph{Part} (i). By Lemma~\ref{5.2:lem-short-box-cubes}, $S$ lies in a regular cube $\square$ of index
$\mathsf{m}\in\{n-1,n\}$, whose cells of index $\mathsf{m}$ are at $\operatorname{dist}_n<1$ from
$\Delta_y$. Moreover $S$ meets cells of levels $n-1$, $n$, $n+1$ only, all at $\operatorname{dist}_n<1$
from $\Delta_y$, and $S$, being a sub-box of $\square$, inherits the gauge of $\square$. By the cube clause
of Lemma~\ref{5.2:lem-carrier-domain}, in \eqref{5.2:eq-carrier-domain-b}, the $G$-part $\bar U(m)$ has on
$\square$ a $G$-valued gauge $u$ with $\bar U(m)^u=e^{i\eta A}$ and
\begin{equation*}
\ell_{\mathsf{m}}|A|,\ \ell_{\mathsf{m}}^2|\nabla A| < 3\,\mathsf{w}^{\mathrm{win}}\cdot 2BM\,\alpha_{0,\mathsf{m}} .
\end{equation*}
Write $\mathcal{U}(m)=e^{i\eta A'}\,\bar U(m)$ bond by bond, where $\mathcal{U}(m)$ is the $U$-slot of the
carrier, and write $\nabla_{\bar U}$ for the covariant difference in the background $\bar U(m)$. In the gauge
$u$ the $U$-slot reads
\begin{equation}\label{5.2:eq-carrier-regularity-1}
\mathcal{U}(m)^u=e^{i\eta\operatorname{Ad}_uA'}\,e^{i\eta A},
\qquad
\nabla(\operatorname{Ad}_uA')=\operatorname{Ad}_u\nabla_{\bar U}A' .
\end{equation}
The field $\operatorname{Ad}_uA'$ and its differences are read in entries $6$, $7$, $8$ of
the clauses of Lemma~\ref{5.2:lem-carrier-domain}. For $a\in\{6,7\}$, the bound
\eqref{5.2:eq-carrier-domain-b} of Lemma~\ref{5.2:lem-carrier-domain} gives
$Q_a(\operatorname{chain}(m))(y')<\theta_a(y')$ at every
$y'\in S$. By \eqref{5.2:eq-carrier-domain-a},
\begin{equation*}
\theta_a(y')\le 3c_\alpha\,\mathsf{w}^{\mathrm{win}}(\Delta_{y'})\,\mathbf{r}_a(\pi(y')),
\end{equation*}
where $\pi(y')$ is the clause index at $y'$ in Lemma~\ref{5.2:lem-carrier-domain}, and
$\mathbf{r}_a(\pi(y'))$ is $\alpha_{*,\pi(y')}$ times the level-$\pi(y')$ unit of entry $a$.

Two facts convert these bounds to the level of $S$:
\begin{itemize}
\item The points of $S$ lie in cells of levels $n-1$, $n$, $n+1$. Reading entries $6$, $7$, $8$ in
level-$(n-1)$ units costs at most a factor $L^2$.
\item By the comparison of the radii of adjacent levels, the $\alpha$'s of the levels $n-1$, $n$, $n+1$
are at most $2\alpha_{\cdot,n}$.
\end{itemize}
The Baker--Campbell--Hausdorff bound, in the case where the background variable of that bound is the
identity $1$, then merges the two exponentials in
\eqref{5.2:eq-carrier-regularity-1} into a single $e^{i\eta\mathcal{A}^0}$. This bound is applied to the
two factors of \eqref{5.2:eq-carrier-regularity-1}, each bounded as just described. With the constant
$c_{\mathrm{reg}}$ this yields $\ell|\mathcal{A}^0|,\ \ell^2|\nabla\mathcal{A}^0|\le\mathfrak{a}'$.

Entries $3$ and $1$ are used as they stand. Entry $3$ gives the bound $\ell^3|\mathbf{J}^0|\le\mathfrak{a}'$
on the current. Entry $1$ gives the plaquette bound
$|\partial\mathcal{U}(m)-1|\le\mathfrak{a}'(\eta/\ell)^2$.

\emph{Part} (ii). Let $S\subset\mathsf{C}$ and $|\hat\lambda|\le\rho_S$. By
Corollary~\ref{5.2:cor-decoupled-carrier}, $P_{\hat\lambda}\restriction S$ has the same $G$-part $U_2$ as the
carrier, $U_2$ being the $G$-part of the carrier on $\mathsf{C}$.
It also satisfies $Q_a(P_{\hat\lambda})(y')<\theta_a(y')$ for $a\in\{1,3,6,7\}$ and $y'\in S$.
The gauge $u$ depends only on the $G$-part, so it does not depend on $\hat\lambda$. The argument of
part (i) therefore applies to $P_{\hat\lambda}\restriction S$ with the same $u$.

\emph{The ceiling} $60\,d\,L\,M\,\mathfrak{a}'\le 1$. Suppose a weighted cell of a live arrest
$\mathfrak{a}$ lies at $\operatorname{dist}_n\le 1$ from $\Delta_y$. We first show that this forces
$n\le k_a$. If $n>k_a$, then $X_{\mathfrak{a}}$ together with its collar lies in
$Z_{n-1}$, by the collar property of live arrests. On the other hand
$\operatorname{dist}_n(Z_{n-1},\Omega_n)\ge 1$, by the cases of the first clause of the geometry of the
plain units. As $y$ is a point of zone $n$, we have $\Delta_y\subset\Omega_n\subset Z_{n-1}^{\,c}$, while
the weighted cell lies in $X_{\mathfrak{a}}$, inside its collar, inside $Z_{n-1}$; separated in this way
from $Z_{n-1}^{\,c}$, the cell is at $\operatorname{dist}_n>1$ from $\Delta_y$, contrary to the
assumption; hence $n\le k_a$.

Next, let $k_{0,a}$ be the index in the definition of the window weight of the arrest $\mathfrak{a}$ in
\eqref{5.2:eq-plain-interior-units-a}. By the bound on the minimal radius at step $k_a$,
\begin{equation*}
\mathsf{w}^{\mathrm{win}}=L\check R_{k_{0,a}+1}\le L\,\mathsf{R}^{\uparrow}(\check{\mathsf{T}}_{k_a}),
\end{equation*}
and $\check{\mathsf{T}}_{k_a}\le\log x_{k_a}$ by the definition of
$\check{\mathsf{T}}_k$; since $\mathsf{R}^{\uparrow}$ is nondecreasing,
$\mathsf{R}^{\uparrow}(\check{\mathsf{T}}_{k_a})\le\mathsf{R}^{\uparrow}(\log x_{k_a})$. Moreover,
$\alpha_{*,n}\le 2\alpha_*(x_{k_a})$. These bounds hold for every live
arrest $\mathfrak{a}$ with a weighted cell at $\operatorname{dist}_n\le 1$ from $\Delta_y$, and they are
applied to the window weight entering $\mathsf{w}^{\mathrm{win},\vee}(\Delta_y)$ in
\eqref{5.2:eq-plain-interior-units-a}. Hence
\begin{align*}
60\,d\,L\,M\,\mathfrak{a}'
&=180\,d\,L\,M\,c_{\mathrm{reg}}\,\mathsf{w}^{\mathrm{win},\vee}(\Delta_y)\,\alpha_{*,n}\\
&\le 180\,d\,L\,M\,c_{\mathrm{reg}}\cdot L\,\mathsf{R}^{\uparrow}(\log x_{k_a})
\cdot 2\alpha_*(x_{k_a})\\
&=360\,d\,L^2M\,c_{\mathrm{reg}}\,\mathsf{R}^{\uparrow}(\log x_{k_a})\,
\alpha_*(x_{k_a})\le 1 .
\end{align*}
The last inequality is the first line of \eqref{5.2:eq-carrier-regularity-a}, evaluated at
$x=x_{k_a}$.
\end{proof}

\begin{remark}[Unit sizes free of scale powers]\label{5.2:rem-scale-free-sizes}
In this remark, and in the count of $M^4t^{-4}$ points of $T^{(j)}$ per cell made in the column
bound for pointed units of this section, a \emph{cell} is a cube of the partition $\pi_n$, that
is, a cube of side $M$ in level-$n$ lattice units. Cells are not to be confused with the
$MR_j$-cubes that serve the regions $\Omega_j$ in \cite[p.~178]{B-CMP122a}. For a unit of level
$j$ read at a cell of level $n \ge j$ we write $t = L^{j-n} \le 1$.

In the first three rows of the bound on the three families of units, the size of a unit
$\mathfrak{E}$ is
\begin{equation}\label{5.2:eq-scale-free-sizes-1}
\mathfrak{N} = \sup_{\mathfrak{D}_{\mathfrak{E}}} |\mathfrak{E}| ,
\end{equation}
and it carries no power of $t$. If each unit in these rows is bounded only by its size
\eqref{5.2:eq-scale-free-sizes-1}, the units anchored in one cell of the core $\mathsf{C}$
contribute in total
\begin{equation}\label{5.2:eq-scale-free-sizes-2}
M^4 E_0 \sum_{j \le n} L^{4(n-j)} ,
\end{equation}
a quantity that depends on the number of steps; here $E_0$ is the constant of the exponential
bounds on single units. A bound uniform in the depth therefore needs
the product of the powers of $t$ with the depth. This product is stated in
\cite[p.~374]{B-CMP122b}: ``We discuss here only exponential factors, the other factors,
like the powers of $L^j\eta$ \dots\ were discussed thoroughly before'' and ``For E-terms and
R-terms this includes renormalization''.
It is not an input of the present chapter: the product bound for rows one to
three is supplied by the increment bound for deep pointed units of this section. For the fourth
row no such question arises, since there $\mathfrak{N} = \mathfrak{N}^{A,\mathbb{C}}$ is itself
the size that enters the budget estimate.
\end{remark}

\begin{proof}
For \eqref{5.2:eq-scale-free-sizes-2}, fix a level $j \le n$. One level-$n$ lattice unit is
$L^{n-j}$ level-$j$ lattice units. A cell, of side $M$ in level-$n$ units, therefore has side
$ML^{n-j}$ in level-$j$ units. Hence it contains
\begin{equation*}
(ML^{n-j})^4 = M^4 L^{4(n-j)} = M^4 t^{-4}
\end{equation*}
points of $T^{(j)}$. When the units anchored at these points, summed over their localization
domains with the exponential weight, are bounded only by their sizes
\eqref{5.2:eq-scale-free-sizes-1}, no power of $t$ offsets the number of points, and level $j$
contributes $M^4E_0L^{4(n-j)}$ to the cell, that is, $E_0$ per point. Summing over the levels
$j \le n$ gives \eqref{5.2:eq-scale-free-sizes-2}. The sum has one term for each level and every
term $L^{4(n-j)}$ is at least $1$, so it is at least the number of levels $j \le n$. The quantity
\eqref{5.2:eq-scale-free-sizes-2} is therefore not bounded uniformly in the number of steps.
\end{proof}

\begin{lemma}[Second-order bound in a local gauge]\label{5.2:lem-local-gauge-bound}
Let $j\le n-1$ be levels, let $t:=L^{j-n}$ and $t':=Lt$. Write $\xi$ for the spacing of the
level-$j$ lattice, so that its
bonds are $b=\langle x,x+\xi e_\mu\rangle$. Write $G^{\mathbb{C}}=\operatorname{SL}(N,\mathbb{C})$ for the
complexification of $\operatorname{SU}(N)$, and $\mathfrak{g}$ for the Lie algebra of $\operatorname{SU}(N)$.
Let $F(X,\cdot)$ be analytic on $U^c_j(X,a_j)$, Ba{\l}aban's space of complex configurations
over $X$ at the radii $a_j$, invariant there under the gauge
transformations of \cite[(1.10)]{B-CMP122b}, and bounded by $|F|\le\mathfrak{n}$. Let $y\in X\subset S$,
where $S$ is a box whose sides, in level-$j$ lattice units, are at most $D$, with $D\le M/t'$.
Suppose that in one $G^{\mathbb{C}}$-gauge the restriction of $(\mathbf{U},\mathbf{J})$ to $S$ reads
$(e^{i\xi\mathcal{A}^0},\mathbf{J}^0)$ with, in level-$j$ units,
\begin{equation*}
|\mathcal{A}^0|\le t'\mathfrak{a}',\qquad |\nabla\mathcal{A}^0|\le t'^2\mathfrak{a}',\qquad
|\mathbf{J}^0|\le t'^3\mathfrak{a}',
\end{equation*}
and that $60dLM\mathfrak{a}'\le 1$. Then there is a gauge in which the pair reads
$(e^{i\xi\mathcal{A}},\mathbf{J}')$ with
\begin{equation}\label{5.2:eq-local-gauge-bound-a}
|\mathcal{A}|\le 4dDt'^2\mathfrak{a}',\qquad |\nabla\mathcal{A}|\le 16t'^2\mathfrak{a}',\qquad
|\mathbf{J}'|\le 2t'^3\mathfrak{a}',
\end{equation}
and
\begin{equation}\label{5.2:eq-local-gauge-bound-1}
\bigl|F(X,(\mathbf{U},\mathbf{J}))-F(X,1)\bigr|
\le 2\mathfrak{n}\Bigl(\frac{16dDt'^2\mathfrak{a}'}{\mathfrak{r}_j}\Bigr)^{2}.
\end{equation}
\end{lemma}

\begin{proof}
We work in the gauge of the hypothesis, in which the restriction of the pair to $S$ is
$(e^{i\xi\mathcal{A}^0},\mathbf{J}^0)$.

\emph{The gauge transformation.} We measure coordinates, $\mathcal{A}^0$ and $\mathbf{J}^0$ in
level-$j$ units, in which $\xi=1$. Put $c_\mu:=\mathcal{A}^0_\mu(y)$, let
$s(x):=\sum_\mu(x_\mu-y_\mu)c_\mu$ for $x\in S$, and let $g:=e^{is}$. For $x\in S$ each
$|x_\mu-y_\mu|$ is at most $D$ and each $|c_\mu|$ is at most $t'\mathfrak{a}'$. Since
$Dt'\le M$ and $60dLM\mathfrak{a}'\le 1$,
\begin{equation*}
|s|\le dDt'\mathfrak{a}'\le dM\mathfrak{a}'\le\tfrac16 .
\end{equation*}
Let $b=\langle x,x+\xi e_\mu\rangle$ be a bond in $S$. Then $s(x+\xi e_\mu)=s(x)+\xi c_\mu$.
Define $\tilde c_\mu$ (depending on $x$) by
\begin{equation*}
e^{-i\xi\tilde c_\mu}:=e^{-i(s+\xi c_\mu)}e^{is},
\end{equation*}
with $s=s(x)$. By this definition,
\begin{equation*}
g(x+\xi e_\mu)^{-1}=e^{-i(s+\xi c_\mu)}=e^{-i\xi\tilde c_\mu}e^{-is},
\end{equation*}
and since $g(x)=e^{is}$ we obtain the identity
\begin{equation}\label{5.2:eq-local-gauge-bound-2}
g(x)e^{i\xi\mathcal{A}^0(b)}g(x+\xi e_\mu)^{-1}
=\operatorname{Ad}_{e^{is}}\bigl[e^{i\xi\mathcal{A}^0(b)}e^{-i\xi\tilde c_\mu}\bigr].
\end{equation}
Duhamel's formula for the derivative of the exponential map, with $|s|\le 1/6$, gives
$|\tilde c_\mu-c_\mu|\le 4|s||c_\mu|$.

\emph{The first bound in \eqref{5.2:eq-local-gauge-bound-a}.} Write the left member of
\eqref{5.2:eq-local-gauge-bound-2} as $e^{i\xi\mathcal{A}(b)}$. Merge the two exponentials inside
the bracket by the Baker--Campbell--Hausdorff formula and use $\mathcal{A}^0_\mu(y)=c_\mu$:
\begin{equation*}
\mathcal{A}(b)=\operatorname{Ad}_{e^{is}}\Bigl[\mathcal{A}^0_\mu(x)-\mathcal{A}^0_\mu(y)
+(c_\mu-\tilde c_\mu)+O\bigl(\xi|\mathcal{A}^0(b)||\tilde c_\mu|\bigr)\Bigr].
\end{equation*}
The first difference is estimated along a lattice path from $y$ to $x$ inside the box $S$. Such a
path has length at most $dD$ level-$j$ units, and along it the difference quotient
$\nabla\mathcal{A}^0$ is at most $t'^2\mathfrak{a}'$, so
\begin{equation*}
|\mathcal{A}^0_\mu(x)-\mathcal{A}^0_\mu(y)|\le dDt'^2\mathfrak{a}'.
\end{equation*}
The second term obeys
\begin{equation*}
|c_\mu-\tilde c_\mu|\le 4|s||c_\mu|\le 4dDt'^2\mathfrak{a}'^{\,2}.
\end{equation*}
The third term is at most a constant multiple of $\xi|\mathcal{A}^0||\tilde c_\mu|$, where
$|\mathcal{A}^0|\le t'\mathfrak{a}'$ and $|\tilde c_\mu|\le(1+4|s|)t'\mathfrak{a}'$. Conjugation by
$e^{is}$ changes norms by at most the factor $e^{2|s|}\le e^{1/3}<2$. Since
$\mathfrak{a}'\le 1/(60dLM)\le 1/60$, the second term is at most $dDt'^2\mathfrak{a}'/15$. The
third term is at most a fixed multiple of $t'^2\mathfrak{a}'^{\,2}(1+4|s|)$, the multiple being the
absolute constant of the quadratic Baker--Campbell--Hausdorff remainder for arguments of norm at
most $1/6$; since $\mathfrak{a}'\le 1/(60dLM)$ and $D\ge 1$ (the box $S$ contains the bond $b$),
it is at most $dDt'^2\mathfrak{a}'/2$. With the factor $e^{2|s|}<2$ the three contributions add
up to at most
\begin{equation*}
2\Bigl(1+\tfrac1{15}+\tfrac12\Bigr)dDt'^2\mathfrak{a}'\le 4dDt'^2\mathfrak{a}'.
\end{equation*}
This is the first inequality in \eqref{5.2:eq-local-gauge-bound-a}.

\emph{The second bound in \eqref{5.2:eq-local-gauge-bound-a}.} Take the lattice difference of the
expression for $\mathcal{A}(b)$ displayed above in a second direction. Only three quantities enter:
$|\nabla\mathcal{A}^0|$, $|c|^2$ and $|c||\mathcal{A}|$. Indeed
\begin{equation*}
|\nabla\operatorname{Ad}_{e^{is}}|\le 4|c|,\qquad |\nabla\tilde c|\le 8|c|^2 .
\end{equation*}
The product $|c||\mathcal{A}|$ is controlled by the first bound together with $Dt'\le M$. This
gives the second inequality in \eqref{5.2:eq-local-gauge-bound-a}.

\emph{The third bound in \eqref{5.2:eq-local-gauge-bound-a}.} The current $\mathbf{J}'$ is
$\mathbf{J}^0$ carried by the same gauge transformation $g=e^{is}$, which conjugates it by the
values of $g$; its norm therefore changes by at most the factor $e^{2|s|}\le e^{1/3}<2$ found
above. With $|\mathbf{J}^0|\le t'^3\mathfrak{a}'$ this gives $|\mathbf{J}'|\le 2t'^3\mathfrak{a}'$.

\emph{The second-order bound.} Put
\begin{equation*}
S_0:=\frac{\mathfrak{r}_j}{16dDt'^2\mathfrak{a}'},\qquad
\varphi(\sigma):=F\bigl(X,(e^{i\sigma\xi\mathcal{A}},\sigma\mathbf{J}')\bigr).
\end{equation*}
Let $|\sigma|<S_0$. By \eqref{5.2:eq-local-gauge-bound-a}, the quantities $|\sigma\mathcal{A}|$,
$|\sigma\nabla\mathcal{A}|$ and $|\sigma\mathbf{J}'|$ are bounded by
\begin{equation*}
\frac{\mathfrak{r}_j}{4},\qquad \frac{\mathfrak{r}_j}{dD},\qquad
\frac{2t'}{16dD}\,\mathfrak{r}_j,
\end{equation*}
respectively, and the last is at most $\mathfrak{r}_j$, because $t'\le 1$ (as $j\le n-1$) and
$D\ge 1$ give $2t'\le 2\le 16dD$. Hence the pair
$(e^{i\sigma\xi\mathcal{A}},\sigma\mathbf{J}')$ satisfies the hypotheses of the Taylor-ball lemma at
radius $\mathfrak{r}_j$, and the Taylor-ball lemma's ball is contained in the domain
$U^c_j(X,a_j)$, on which $F(X,\cdot)$ is analytic and bounded by $\mathfrak{n}$. So
$\varphi$ is analytic with $|\varphi|\le\mathfrak{n}$ on $|\sigma|<S_0$.

The endpoint values are as follows. By the gauge invariance of $F(X,\cdot)$,
$\varphi(1)=F(X,(\mathbf{U},\mathbf{J}))$. At $\sigma=0$ the argument is the unit configuration with
zero current, so $\varphi(0)=F(X,1)$.

Next, $\varphi'(0)=0$. We have $\varphi'(0)=DF(X;1,0)[(i\xi\mathcal{A},\mathbf{J}')]$, where
$DF(X;1,0)$ is the derivative of $F(X,\cdot)$ at the unit configuration with zero current.
A constant gauge transformation $h$ fixes that configuration. The invariance of $F(X,\cdot)$ under
constant gauge transformations therefore makes the linear functional $DF(X;1,0)$
invariant under the constant adjoint action $\operatorname{Ad}_h$ applied to both slots.
Differentiating in $h$, it annihilates $[\lambda,\cdot\,]$ for every $\lambda\in\mathfrak{g}$, bond by
bond. Since $[\mathfrak{g},\mathfrak{g}]=\mathfrak{g}$, every element supported on a single bond is a sum of
such commutators, so $DF(X;1,0)$ vanishes. This is Ba{\l}aban's argument for
\cite[(4.14)]{B-CMP109}.

If $S_0>1$, apply part (b) of the Schwarz-type lemma to $\varphi$. It concerns a function analytic
and bounded
by $\mathfrak{n}$ on the disc of radius $S_0$ whose derivative vanishes at $0$, and gives
\begin{equation*}
|\varphi(1)-\varphi(0)|\le 2\mathfrak{n}S_0^{-2}.
\end{equation*}
This is \eqref{5.2:eq-local-gauge-bound-1}.

If $S_0\le 1$, the right member of \eqref{5.2:eq-local-gauge-bound-1} is at least $2\mathfrak{n}$.
The bound is then trivial provided $(\mathbf{U},\mathbf{J})$ lies in $U^c_j(X,a_j)$, where
$|F(X,\cdot)|\le\mathfrak{n}$: this membership is presupposed by the left member
$F(X,(\mathbf{U},\mathbf{J}))$ of \eqref{5.2:eq-local-gauge-bound-1}, $F(X,\cdot)$ being given on
$U^c_j(X,a_j)$, and the unit configuration lies in that domain.
\end{proof}

Here $t'=L^{j-(n-1)}$ is the ratio of level $j$ to level $n-1$, and the hypothesis records only the
value $\mathcal{A}^0(y)$ and the gradient $\nabla\mathcal{A}^0$ of a field of size $\mathfrak{a}'$ at
level $n-1$, read in level-$j$ units; no minimiser of a variational problem is used, and no
further level of data is read.

\begin{remark}[No H\"older bound off the regular region]\label{5.2:rem-no-holder-bound}
Let $y$ be a point of zone $n$ and $j\le n$ a level, and put $t:=L^{j-n}$ and $\nu:=1+n-j$.
Let $C$ be the component of the treated class that contains $y$, and let
$m=(\zeta,\varpi,\vec{\mathsf{s}})\in\mathfrak{P}$ be a point of the parameter polydisc.
\begin{enumerate}
\item[(a)] The gain $t^{\hat\alpha}$ in the second part of the lemma that bounds the renormalised
pointed group $\mathfrak{i}^{(j)}_y$ comes from the comb-gauge remainder bound in the definition of
an admissible background: for all $x$ and $\mu$ as in that definition,
\begin{equation}\label{5.2:eq-no-holder-bound-1}
|\mathfrak{e}_\mu(x)|\le C_{\mathfrak{e}}\,(1+|x-y|)^2\,\mathfrak{a}_n\,t^{2+\hat\alpha},
\end{equation}
where $C_{\mathfrak{e}}$ denotes the constant printed in that definition; this is
\cite[(4.18)]{B-CMP109}, drawn there from the H\"older norm in \cite[(3.32)]{B-CMP109}, which
that paper obtains from \cite[(3.27)]{B-CMP109} for the field
$H_{k+1}(\square_0,\tfrac1i\log V)$, in the notation of that paper, $V$ being the averages of
$\mathbf{U}$ \cite[(3.24)]{B-CMP109}.
\item[(b)] None of the conditions \cite[(1.11)--(1.16)]{B-CMP109}, \cite[(2.34)--(2.39)]{B-CMP119},
\cite[(1.64)--(1.69)]{B-CMP122a} contains a H\"older clause.
\item[(c)] Off the regular region $\Lambda^{\natural}$, and in the collar $\mathfrak{c}$ of the
component $C$, the carrier pair $\operatorname{chain}(m)$ is the product of the layers with the raw
configuration $\mathbb{U}$, by the lemma representing
the carrier there. At complex $\zeta$ the level-$j$ data of $\operatorname{chain}(m)$ obey only the
first of the comb-gauge bounds in the definition of an admissible background, and a core
$\mathfrak{i}^{(j)}_y$ evaluated on these data has the marginal bound $\nu^2t^4$ and no better bound.
\item[(d)] Bounds stated through the quantity $U_2$ belonging to the region $D_2$ fixed with the
geometry of $C$, and through the cell norms $\mathsf{N}_y(\mathbb{H}_s)$ of the layers, hold on $D_2$.
They say nothing on $C\setminus\Lambda^{\natural}$, which contains most of the shallow units in the
sense of Definition~\ref{5.2:def-plain-interior-units}.
\item[(e)] The convention of Definition~\ref{5.2:def-plain-interior-units} evaluates a smooth far core
$v$, with level $p_v$ and box $Q_v$ as in that definition, at the box minimiser of the level-$p_v$
averages on $Q_v$ of the $\mathbf{U}$-slot $\mathcal{U}(m)$ of $\operatorname{chain}(m)$. This is
Ba{\l}aban's own device \cite[(3.24), (3.27)]{B-CMP109}, applied to the whole slot.
\end{enumerate}
\end{remark}

\begin{proof}
We take the assertions in order.

For (a), the factor $t^{\hat\alpha}$ enters the bound of $\mathfrak{i}^{(j)}_y$ through the remainder
$\mathfrak{e}_\mu$ of the comb gauge, that is, through \eqref{5.2:eq-no-holder-bound-1}; the citations
given in (a) trace it to \cite[(3.27)]{B-CMP109}.

For (b), we go through the conditions \cite[(1.11)--(1.16)]{B-CMP109},
\cite[(2.34)--(2.39)]{B-CMP119} and \cite[(1.64)--(1.69)]{B-CMP122a} in turn. None of them contains a
H\"older clause; so the H\"older norm of \cite[(3.32)]{B-CMP109} is not supplied by these conditions.

For (c), off $\Lambda^{\natural}$ and in the collar $\mathfrak{c}$ of the component $C$ the carrier
pair $\operatorname{chain}(m)$ is the product of the layers with the raw configuration $\mathbb{U}$,
by the lemma representing the carrier there. Hence at complex $\zeta$ the level-$j$ data of
$\operatorname{chain}(m)$ obey only the first of the comb-gauge bounds in the definition of an
admissible background, and by (b) no further condition supplies a H\"older bound.

The improvement of a core $\mathfrak{i}^{(j)}_y$ over the marginal size comes from the cancellation at
$y$, which requires the field at $x$ to be close to the field at $y$. A merely bounded field gives no
such closeness. Hence the bound $\nu^2t^4$ is the best the data yield.

For (d), bounds stated through $U_2$ and $\mathsf{N}_y(\mathbb{H}_s)$ are statements about $D_2$; such
bounds say nothing on $C\setminus\Lambda^{\natural}$, and on it the carrier is as described in (c).

For (e), with $p:=p_v$, the convention of Definition~\ref{5.2:def-plain-interior-units} evaluates a
smooth far core $v$ at the box minimiser, with its current, of the level-$p$ block averages of the
restriction of $\mathcal{U}(m)$ to $Q_v$. Together with its current, this minimiser is the object of
\cite[(3.24), (3.27)]{B-CMP109}, here formed from the averages of the slot $\mathcal{U}(m)$ itself.
The convention thus applies Ba{\l}aban's device to the whole slot.
\end{proof}

\begin{lemma}[The smoothed core bound at intermediate scale]\label{5.2:lem-smoothed-core-bound}
Let $y$ be a point of zone $n$, let $j\le n$, $\nu := 1+n-j$, and let the core
$\mathfrak{i}^{(j)}_y$, the box $\square_0$ and, for $p\ge j$, the box
$Q^{(p)}$, that is, the box $[\square_0]_p$ enlarged by two layers of $\pi_p$-cubes,
be as in Definition~\ref{5.2:def-plain-interior-units}.
Assume the following four statements.
\begin{itemize}
\item[(a)] The admissibility lemma for box minimisers, with its number $c''\le 1$: for every
box $Q$ and every $\mathbf{U}$ in $c''$ times the ball of the polydisc lemma over $Q$, the
configuration $U_p(Q,M^p(\mathbf{U}))$ is an admissible background at scale $p$ on
$Q^{\sim -2}$.
\item[(b)] Clause~(i) of the pointed-group lemma (the bound on the constituents of a core over
an admissible background).
\item[(c)] The polydisc lemma, with its radii $\mathfrak{r}_p$, at level $p$.
\item[(d)] Clause~(b) of the Schwarz bound.
\end{itemize}
Let $j\le p$, $\nu_p := 1+p-j\le \nu$, $t_p := L^{j-p}$ and $Q := Q^{(p)}$. Let $\mathcal{U}$ be a
configuration whose restriction to $Q$ reads, in one complex gauge (a gauge transformation with
values in $G^{\mathbb{C}}=\mathrm{SL}(N,\mathbb{C})$, the complexification of $\mathrm{SU}(N)$),
$e^{i\eta\mathcal{A}}$ with
\begin{equation}\label{5.2:eq-smoothed-core-bound-1}
\max\{\ell_p|\mathcal{A}|,\ \ell_p^2|\nabla\mathcal{A}|\}\le \epsilon\, c''\mathfrak{r}_p,
\qquad \epsilon\le 1 ,
\end{equation}
and put $\operatorname{ext}_p(\mathcal{U}) := (U_p,J_p)(Q,M^p(\mathcal{U}))$, the box minimiser on
$Q$ with data the level-$p$ averages $M^p(\mathcal{U})$ of $\mathcal{U}$, together with its current.
Then
\begin{equation}\label{5.2:eq-smoothed-core-bound-2}
\bigl|\mathfrak{i}^{(j)}_y(\operatorname{ext}_p(\mathcal{U}))\bigr|\le 2\epsilon^2\,\mathfrak{N}_p ,
\end{equation}
where
\begin{equation}\label{5.2:eq-smoothed-core-bound-a}
\mathfrak{N}_p := (C_I+2K(d))\,E_0\,\nu_p^3\,\bar\rho_{p-j}^{\,2}\,(1+\bar\rho_{p-j})\,
t_p^{4+\hat\alpha};
\end{equation}
here $E_0$ is the size constant of the pointed rows (the prefactor of the row-one sizes
$E_0e^{-\kappa d_j(X)}$), and $K(d)$ is the constant, depending on $d$ only, that bounds the sum
over domains in clause~(i) of the pointed-group lemma.
\end{lemma}

\begin{proof}
\emph{Reduction to the given gauge.} For a gauge transformation $g$ of the configuration on $Q$
there is an induced gauge transformation $\tilde g$ with
\begin{equation*}
\operatorname{ext}_p(\mathcal{U}^g)=\operatorname{ext}_p(\mathcal{U})^{\tilde g};
\end{equation*}
this is the covariance of the averaging operation, \cite[(11)]{B-CMP98}, in its extension to
complex gauge transformations, combined with \cite[(3.29)]{B-CMP109}. The core
$\mathfrak{i}^{(j)}_y$ is invariant under gauge transformations of its argument. Hence the left
member of \eqref{5.2:eq-smoothed-core-bound-2} does not change if $\mathcal{U}$ is replaced by a
gauge transform of it, and we may take $\mathcal{U}=e^{i\eta\mathcal{A}}$ on $Q$.

\emph{The auxiliary function.} For complex $\sigma$ with $|\sigma|<1/\epsilon$ set
\begin{equation*}
\Phi(\sigma) := \mathfrak{i}^{(j)}_y\bigl(\operatorname{ext}_p(e^{i\sigma\eta\mathcal{A}})\bigr).
\end{equation*}
For such $\sigma$ the field $\sigma\mathcal{A}$ satisfies
\eqref{5.2:eq-smoothed-core-bound-1} with $\epsilon|\sigma|<1$ in place of $\epsilon$. By the
polydisc lemma (c) at level $p$, the configuration $e^{i\sigma\eta\mathcal{A}}$ therefore lies in
$c''$ times the ball of that lemma over the box $Q$. By the admissibility lemma (a), applied to
this box, the minimiser $U_p(Q,M^p(e^{i\sigma\eta\mathcal{A}}))$ is an admissible background at
scale $p$ on $Q^{\sim-2}$, and so $\operatorname{ext}_p(e^{i\sigma\eta\mathcal{A}})$, which is this
minimiser together with its current, is an admissible background at scale $p$ on $Q^{\sim-2}$.
Moreover $Q^{\sim-2}\supseteq\square^{\sim\nu}$, the region on which the
constituents of the core are read. The minimiser and its current depend analytically on $\sigma$
by \cite[(3.24), (3.27)]{B-CMP109}. Hence $\Phi$ is analytic on the disc $|\sigma|<1/\epsilon$.

\emph{The bound on the disc.} We apply clause~(i) of the pointed-group lemma (b) with
$(n,\nu,t)$ replaced by $(p,\nu_p,t_p)$, at the admissible background
$\operatorname{ext}_p(e^{i\sigma\eta\mathcal{A}})$. It splits the constituents into two groups.

The first group consists of the constituents with $X\subset\square^{\sim\nu_p}$. Together with the
counterterm, they are bounded by
\begin{equation*}
C_I\,E_0\,\nu_p^3\,\rho_{p,j}^{\,2}\,(1+\rho_{p,j})\,t_p^{4+\hat\alpha},
\end{equation*}
where $\rho_{p,j}$ is the radii ratio of that clause, read at the levels $(p,j)$.

The second group consists of all other constituents. They have $d_j(X)\ge\nu_p$, and each of them
is bounded singly by $2E_0e^{-\kappa d_j(X)}$. Summed over $X$, using $\kappa-c_d-1\ge 9\log L$,
they contribute at most $2K(d)E_0t_p^5$.

Adding the two contributions gives $|\Phi(\sigma)|\le\mathfrak{N}_p$ for $|\sigma|<1/\epsilon$, with
$\mathfrak{N}_p$ as in \eqref{5.2:eq-smoothed-core-bound-a}: the first contribution is at most the
term of $\mathfrak{N}_p$ carrying $C_I$, since $\rho_{p,j}\le\bar\rho_{p-j}$; the second is at most
the term carrying $2K(d)$, since $t_p\le 1$ and $\hat\alpha\le 1$ give
$t_p^5\le t_p^{4+\hat\alpha}$, and since $\nu_p^3\bar\rho_{p-j}^{\,2}(1+\bar\rho_{p-j})\ge 1$.

\emph{Vanishing to second order.} At $\sigma=0$ the argument is the trivial configuration, and
$\operatorname{ext}_p$ of it is the trivial background; every constituent $F(X,\cdot)-F(X,1)$
vanishes there, and so does the counterterm, since the plaquette density vanishes at
$\mathbb{1}$. Hence $\Phi(0)=0$. Both parts also vanish to first order. For the constituents, the
first derivative $DF(X;1,0)$ at the identity vanishes, by the argument given in
the proof of the second-order bound in a local gauge (Lemma~\ref{5.2:lem-local-gauge-bound}):
$DF(X;1,0)$ is invariant under a constant $\operatorname{Ad}_g$ acting on both slots, so it
annihilates, bond by bond, the commutator with every constant element of $\mathfrak{g}$, the Lie
algebra of $\mathrm{SU}(N)$, and $[\mathfrak{g},\mathfrak{g}]=\mathfrak{g}$. For the counterterm,
the plaquette density
\begin{equation*}
a_{p'}(U)=1-\tfrac12\operatorname{tr}(U+U^{-1})(\partial p')
\end{equation*}
of a plaquette $p'$ vanishes to second order at $U(\partial p')=\mathbb{1}$, so the counterterm
contributes nothing to first order. Hence $\Phi'(0)=0$.

\emph{Conclusion.} We apply clause~(b) of the Schwarz bound (d) to $\Phi$ on the disc of radius
$1/\epsilon$, on which $|\Phi|\le\mathfrak{N}_p$ and $\Phi'(0)=0$. Since $\Phi(0)=0$, this gives
\begin{equation*}
|\Phi(1)| = |\Phi(1)-\Phi(0)|\le 2\,\mathfrak{N}_p\,\epsilon^2 .
\end{equation*}
Since $\Phi(1)=\mathfrak{i}^{(j)}_y(\operatorname{ext}_p(\mathcal{U}))$, this is
\eqref{5.2:eq-smoothed-core-bound-2}.
\end{proof}

\begin{definition}[The smoothing depth and pointed-unit sizes]\label{5.2:def-smoothing-depth}
Let $v$ be a pointed unit of the class described in Definition~\ref{5.2:def-plain-interior-units}.
Let $y=y_v$ be its point, of zone $n$, let $j$ be its level, and put $t=L^{j-n}$ and
$\nu=1+n-j$. Let $\bar{\mathsf{w}}_v$ be its weight, and write $\ell_p:=L^p\eta$. Put
\begin{equation*}
\mathfrak{a}'_v:=c_{\mathrm{reg}}\,\bar{\mathsf{w}}_v\,\alpha_{*,n}.
\end{equation*}
Let $c''\le 1$ be the constant of the admissibility statement for box minimisers. That
statement says that $(U_p,J_p)(Q,M^p(\mathbf{U}))$ is an admissible background at scale $p$ on
$Q^{\sim-2}$, the box $Q$ with two boundary layers of cubes of $\pi_p$ removed, whenever
$\mathbf{U}$ lies, over the box $Q$, in $c''$ times the ball of radius $\mathfrak{r}_p$ of the
lemma fixing the polydisc radii $\mathfrak{r}_j$, read at level $p$.
Define
\begin{equation*}
c_\epsilon:=4dM\,c_{\mathrm{reg}}\,\sup_p\frac{\alpha_{*,p}}{c''\,\mathfrak{r}_p}.
\end{equation*}

\emph{Cores.} Let $v$ be a core. Let $\square$ be the cube of $\pi_j$ containing $y$, put
$\square_0=\square^{\sim\nu+1}$, and for $p\ge j$ let $Q^{(p)}$ be $[\square_0]_p$ with two
layers of cubes of $\pi_p$, as in Definition~\ref{5.2:def-plain-interior-units}. For an integer
$p$ with $j\le p\le n-3$ put $a_v:=n-p$. Then:
\begin{itemize}
\item[(i)] $Q^{(p)}$ has at most
\begin{equation*}
(2\nu+3)L^{-(p-j)}+6\;\le\;2a_v+11\;\le\;L^{a_v-1}
\end{equation*}
cubes of $\pi_p$ on a side.
\item[(ii)] Let $m\in\mathfrak{P}$, and assume $Q^{(p)}\subset C$. Let $\mathcal{A}$ be the
field in which the $U$-slot
$\mathcal{U}(m)$ of the carrier pair of $m$ reads, in the gauge given by the second-order
bound in a local gauge (Lemma~\ref{5.2:lem-local-gauge-bound}), $e^{i\eta\mathcal{A}}$ on
$Q^{(p)}$. Put
\begin{equation*}
\epsilon:=\frac{\max\{\ell_p|\mathcal{A}|,\ \ell_p^2|\nabla\mathcal{A}|\}}{c''\,\mathfrak{r}_p}.
\end{equation*}
Then
\begin{equation}\label{5.2:eq-smoothing-depth-1}
\epsilon\;\le\;c_\epsilon\,\bar{\mathsf{w}}_v\,(2a_v+11)\,\bar\rho_{a_v}\,L^{-2(a_v-1)}.
\end{equation}
\end{itemize}

\emph{Sizes.} Put
\begin{equation*}
c_{\vee}:=\max\Bigl\{1,\ \sup_{j\le n}\frac{8L^2c_3a_jc_{\mathrm{reg}}\,\alpha_{*,n}}
{\mathfrak{a}_n\,\mathfrak{r}_j}\Bigr\}.
\end{equation*}
This is a quotient of the constants fixed in the statement that fixes $C_M$, $c_3$ and
$\mathfrak{a}_n$, read in its case $q_0=q_1$, in which
$C_M\rho_{n,j}=2d\mathfrak{a}_n/(c_3a_j)$; here $a_j$ is the radius
of the analyticity space $U^c_j(X,a_j)$ on which the singles are analytic. Put further
\begin{equation*}
c_5^{W}:=\max\Bigl\{1,\ \frac{6L^4}{C_S}\Bigl(\frac{c_{\mathrm{reg}}\,\alpha_{*,n}}
{\mathfrak{a}_n}\Bigr)^2\Bigr\}.
\end{equation*}
The number $c_5^{W}$ is a constant of the Wilson pieces, depending on $d,L,M$; it is not the
flow constant $c_5=2c_2$.

For a single with set $X$, $D_X$ denotes the longest side of its hull $\hat X$, measured in
level-$j$ lattice units; $\hat X$ is made of cubes of $\pi_j$, of side $M$ in these units, and
$D_X\le 4M(d_j(X)+5)$.

The \emph{smoothing depth} $s(\bar{\mathsf{w}})$ and the sizes $\mathfrak{N}^{\flat}(v)$ are:
\begin{equation}\label{5.2:eq-smoothing-depth-a}
\begin{aligned}
s(\bar{\mathsf{w}})&:=\min\bigl\{s\in\mathbb{Z},\ s\ge 1:\ c_\epsilon\bar{\mathsf{w}}(2a+11)
\bar\rho_a L^{-2(a-1)}\le 1\\
&\text{for every integer } a\ge 2+s\bigr\},\\
s^{(9/8)}(\bar{\mathsf{w}})&:=s\bigl(\tfrac98\bar{\mathsf{w}}\bigr);\\
\mathfrak{N}^{\flat}(v)&:=2\mathfrak{n}_v\min\bigl\{1,\ (\bar{\mathsf{w}}_vc_{\vee}C_MD_Xt^2
\rho_{n,j})^2\bigr\}
&&\text{($v$ a short single)},\\
\mathfrak{N}^{\flat}(v)&:=2\mathfrak{n}_v &&\text{($v$ a long single or a tail)},\\
\mathfrak{N}^{\flat}(v)&:=c_5^{W}\,\bar{\mathsf{w}}_v^2\,\mathfrak{N}(v)
&&\text{($v$ a Wilson piece, $X_v$ a block)},\\
\mathfrak{N}^{\flat}(v)&:=6L^4(b_++r)\,\mathfrak{a}_v'^{\,2}\,t^4
&&\text{($v$ the Wilson piece $b_jA(h_y,\cdot)$)},\\
\mathfrak{N}^{\flat}(v)&:=C^{\mathrm{sm}}E_0\,\bar{\mathsf{w}}_v^2(2a_v+11)^2\bar\rho_{a_v}^2\\
&\qquad\times(b_v+1)^3\bar\rho_{b_v}^2(1+\bar\rho_{b_v})L^{-\hat\alpha b_v}\,t^4
&&\text{($v$ a smooth core)},\\
\mathfrak{N}^{\flat}(v)&:=\mathfrak{N}^{\mathrm{mar}}(v) &&\text{($v$ a rough core)}.
\end{aligned}
\end{equation}
Here $C^{\mathrm{sm}}:=2c_\epsilon^2L^4(C_I+2K(d))$, with $C_I$ and $K(d)$ the constants of
\eqref{5.2:eq-smoothed-core-bound-a}. For a smooth core, $a_v=n-p_v$ and
$b_v=p_v-j$. In the entries carrying $b_++r$, $b_+=3\lambda+2c_2$ is the flow constant and $r$
is the constant accompanying it in the bound for the coefficients $b_j$ of the pieces
$b_jA(h_y,\cdot)$. The marginal size is
\begin{equation*}
\mathfrak{N}^{\mathrm{mar}}(v):=\sum_{X\ni y,\ X\subset\square^{\sim\nu}}\mathfrak{N}^{\flat}(X)
+6L^4(b_++r)\,\mathfrak{a}_v'^{\,2}\,t^4,
\end{equation*}
where $\mathfrak{N}^{\flat}(X)$ is the entry of \eqref{5.2:eq-smoothing-depth-a} for the single
$F(X,\cdot)-F(X,1)$ with set $X$. There is a
constant $C^{\mathrm{mar}}$, depending only on $d$, $L$, $M$, $\kappa$, $c_{\vee}C_M$ and
$c_{\mathrm{reg}}\sup_n\alpha_{*,n}$, such that
\begin{equation}\label{5.2:eq-smoothing-depth-2}
\mathfrak{N}^{\mathrm{mar}}(v)\;\le\;C^{\mathrm{mar}}\,\bar{\mathsf{w}}_v^2\,
\bigl[E_0\nu^2\bar\rho_\nu^2+b_++r\bigr]\,t^4 .
\end{equation}

The depth $s(\bar{\mathsf{w}})$ depends on $\bar{\mathsf{w}}$ only through the product
$c_\epsilon\bar{\mathsf{w}}$. The bounds on the depth stated next hold under the growth
condition
\begin{equation}\label{5.2:eq-smoothing-depth-3}
(2a+13)\,\bar\rho_{a+1}\;\le\;L\,(2a+11)\,\bar\rho_a\qquad\text{for every integer } a\ge3
\end{equation}
on the majorants $\bar\rho$ of the radius ratios (Notation~\ref{5.1:not-coupling-functions});
this condition is a hypothesis of these bounds wherever they are used. Under
\eqref{5.2:eq-smoothing-depth-3}, for every $\bar{\mathsf{w}}\ge1$ the set over which the
minimum defining $s(\bar{\mathsf{w}})$ is taken is non-empty, and
\begin{equation*}
s(\bar{\mathsf{w}})\;\le\;s(1)+\lceil\log_L\bar{\mathsf{w}}\rceil\;<\;
s(1)+\log_L\bar{\mathsf{w}}+1 .
\end{equation*}

On the joint polydisc of the lemma fixing the amplitude radius $\rho_A$, the exterior
counterpart of the regularity of the carrier on a short box
(Lemma~\ref{5.2:lem-carrier-regularity}) holds on the exterior hulls with $\mathfrak{a}'_v$
replaced by $\tfrac98\mathfrak{a}'_v$.
There $s$ is evaluated at the dilated argument: for a smooth
far core of a plain second-part unit on the exterior hulls, for its $\mathsf{t}$-piece and for
its subtracted copy, one and the same level $p_v:=n-2-s^{(9/8)}(\bar{\mathsf{w}}_v)$ is taken.
Moreover, under \eqref{5.2:eq-smoothing-depth-3}, for $\bar{\mathsf{w}}\ge1$,
\begin{equation*}
s^{(9/8)}(\bar{\mathsf{w}})\;\le\;s(\bar{\mathsf{w}})+1\;\le\;
s(1)+\lceil\log_L\bar{\mathsf{w}}\rceil+1 .
\end{equation*}
\end{definition}

\begin{proof}
\emph{Assertion (i).} Write $b:=p-j\ge 0$. Then $\nu=1+a_v+b$, so
\begin{equation*}
(2\nu+3)L^{-b}=(2a_v+2b+5)L^{-b}.
\end{equation*}
For $b=0$ this equals $2a_v+5$. For $b\ge1$ it is at most $2a_v+5$, because
$2b\le(2a_v+5)(L^b-1)$. Adding $6$ gives the first inequality of~(i).

For the second inequality, note that $a_v\ge3$, since $p\le n-3$. At $a_v=3$ it reads
$17\le L^2$, which holds since $L\ge 8$ in the setting of
Definition~\ref{5.2:def-plain-interior-units}. When $a_v$ increases by one, the left member
increases by $2$ and the right member is multiplied by $L\ge 8$.

Since $L^{a_v-1}$ cubes of $\pi_p$ in a row span one side of a cube of
$\pi_{p+a_v-1}=\pi_{n-1}$, the sides of
$Q^{(p)}$ are at most one cube of $\pi_{n-1}$. Moreover $Q^{(p)}$ contains $y$. When
$Q^{(p)}\subset C$, as assumed in~(ii), $Q^{(p)}$ is therefore a box of the kind considered in
Lemma~\ref{5.2:lem-short-box-cubes}.

\emph{Assertion (ii).} Take $S=Q^{(p)}$ in the regularity of the carrier on a short box
(Lemma~\ref{5.2:lem-carrier-regularity}), applied with the fine weight $w^{\vee}(\Delta_y)$ in
place of $\mathsf{w}^{\mathrm{win},\vee}(\Delta_y)$, which the lemma permits since the weight
enters it only through $\bar{\mathsf{w}}_v=3w^{\vee}(\Delta_y)$; its radius
\eqref{5.2:eq-carrier-regularity-a} is then $\mathfrak{a}'_v$. It provides
$60dLM\mathfrak{a}'_v\le1$ and one
$G^{\mathbb{C}}$-gauge in which the carrier pair of $m$ on $S$ reads
$(e^{i\eta\mathcal{A}^0},\mathbf{J}^0)$, with
\begin{equation*}
\ell_{n-1}|\mathcal{A}^0|,\quad \ell_{n-1}^2|\nabla\mathcal{A}^0|,\quad
\ell_{n-1}^3|\mathbf{J}^0|\;\le\;\mathfrak{a}'_v .
\end{equation*}

Apply Lemma~\ref{5.2:lem-local-gauge-bound} in level-$p$ units, with
$t'=\ell_p/\ell_{n-1}=L^{-(a_v-1)}$ and box side $D=M(2a_v+11)$ (the length
$M(2a_v+11)\ell_p$). Its hypotheses are the three bounds just displayed, rewritten in
level-$p$ units, together with $D\le M/t'$. The latter is the second inequality of~(i).
The bounds \eqref{5.2:eq-local-gauge-bound-a} then give
\begin{equation*}
\ell_p|\mathcal{A}|\le 4dDt'^2\mathfrak{a}'_v=4dM(2a_v+11)L^{-2(a_v-1)}\mathfrak{a}'_v
\end{equation*}
and
\begin{equation*}
\ell_p^2|\nabla\mathcal{A}|\le16L^{-2(a_v-1)}\mathfrak{a}'_v
\le 4dM(2a_v+11)L^{-2(a_v-1)}\mathfrak{a}'_v .
\end{equation*}

Divide by $c''\mathfrak{r}_p$ and insert $\mathfrak{a}'_v=c_{\mathrm{reg}}\bar{\mathsf{w}}_v
\alpha_{*,n}$. Write $\alpha_{*,n}=\alpha_{*,p}\cdot(\alpha_{*,n}/\alpha_{*,p})$, and use
$\alpha_{*,n}/\alpha_{*,p}\le\bar\rho_{a_v}$, the defining property of the majorant $\bar\rho$
of the radius ratios (Notation~\ref{5.1:not-coupling-functions}). This yields
\begin{equation*}
\epsilon\le 4dMc_{\mathrm{reg}}\frac{\alpha_{*,p}}{c''\mathfrak{r}_p}\,\bar{\mathsf{w}}_v
(2a_v+11)\bar\rho_{a_v}L^{-2(a_v-1)},
\end{equation*}
which is at most the right member of \eqref{5.2:eq-smoothing-depth-1}.

\emph{Role of the depth.} For a smooth core, $p=p_v\le n-2-s(\bar{\mathsf{w}}_v)$, so
$a_v\ge 2+s(\bar{\mathsf{w}}_v)$. Moreover $Q^{(p_v)}\subset Z$ is a box containing $y$, and
$y$ lies in $X_v$, hence in the component $C$ of $Z$ that contains $X_v$; so
$Q^{(p_v)}\subset C$ and (ii) applies. By \eqref{5.2:eq-smoothing-depth-1} and the definition
of $s$ in \eqref{5.2:eq-smoothing-depth-a}, $\epsilon\le1$. With the factor $c''$ contained in
$c_\epsilon$, this is the hypothesis of the smoothed core bound
(Lemma~\ref{5.2:lem-smoothed-core-bound}) on $c''$ times the ball of radius $\mathfrak{r}_p$.

\emph{Existence and growth of the depth.} For integers $a\ge3$ write
$\phi_\epsilon(a):=c_\epsilon(2a+11)\bar\rho_aL^{-2(a-1)}$, so that $s(\bar{\mathsf{w}})$ is
the least integer $s\ge1$ with $\bar{\mathsf{w}}\,\phi_\epsilon(a)\le1$ for every integer
$a\ge2+s$. Since
\begin{equation*}
\frac{\phi_\epsilon(a+1)}{\phi_\epsilon(a)}=\frac{(2a+13)\,\bar\rho_{a+1}}{(2a+11)\,\bar\rho_a}\,
L^{-2},
\end{equation*}
condition \eqref{5.2:eq-smoothing-depth-3} says exactly that
$\phi_\epsilon(a+1)\le L^{-1}\phi_\epsilon(a)$ for every integer $a\ge3$. Assume it. Then
$\phi_\epsilon(a)\le L^{3-a}\phi_\epsilon(3)$ for every $a\ge3$, and the right member tends to
$0$ as $a\to\infty$; hence for every $\bar{\mathsf{w}}\ge1$ there is an integer $s\ge1$ with
$\bar{\mathsf{w}}\,\phi_\epsilon(a)\le1$ for all integers $a\ge2+s$, and the minimum defining
$s(\bar{\mathsf{w}})$ exists. Now let $\bar{\mathsf{w}}\ge1$ and put
$k:=\lceil\log_L\bar{\mathsf{w}}\rceil\ge0$, so that $L^k\ge\bar{\mathsf{w}}$. Let $a$ be an
integer with $a\ge2+s(1)+k$. The integers $a-k,\dots,a-1$ are all at least $2+s(1)\ge3$, so
$k$ applications of $\phi_\epsilon(b+1)\le L^{-1}\phi_\epsilon(b)$ give
$\phi_\epsilon(a)\le L^{-k}\phi_\epsilon(a-k)$. Since $a-k\ge2+s(1)$, the definition of $s(1)$
gives $\phi_\epsilon(a-k)\le1$. Hence
\begin{equation*}
\bar{\mathsf{w}}\,\phi_\epsilon(a)\le\bar{\mathsf{w}}L^{-k}\le1 .
\end{equation*}
As $s(1)+k$ is an integer at least $1$, it belongs to the set over which the minimum is taken,
and $s(\bar{\mathsf{w}})\le s(1)+k$. Finally $\lceil\log_L\bar{\mathsf{w}}\rceil
<\log_L\bar{\mathsf{w}}+1$.

\emph{The dilated depth.} On the joint polydisc, the bounds of the regularity statement hold
with $\mathfrak{a}'_v$ replaced by $\tfrac98\mathfrak{a}'_v$. The argument of (ii) then gives
the right member of
\eqref{5.2:eq-smoothing-depth-1} multiplied by $\tfrac98$, that is, the same expression with
$\bar{\mathsf{w}}_v$ replaced by $\tfrac98\bar{\mathsf{w}}_v$. With
$p_v=n-2-s^{(9/8)}(\bar{\mathsf{w}}_v)$ one has $a_v\ge2+s(\tfrac98\bar{\mathsf{w}}_v)$.
By the definition of $s$ at the argument $\tfrac98\bar{\mathsf{w}}_v$, the dilated $\epsilon$ is
still at most $1$, and the disc of the smoothed core bound does not leave $c''$ times the ball
of radius $\mathfrak{r}_p$.

Assume \eqref{5.2:eq-smoothing-depth-3} and let $\bar{\mathsf{w}}\ge1$. For an integer
$a\ge2+s(\bar{\mathsf{w}})+1$ one has $a-1\ge2+s(\bar{\mathsf{w}})\ge3$, so
$\phi_\epsilon(a)\le L^{-1}\phi_\epsilon(a-1)\le L^{-1}\bar{\mathsf{w}}^{-1}$ by the
definition of $s(\bar{\mathsf{w}})$. Since $L\ge8>\tfrac98$, this gives
$\tfrac98\bar{\mathsf{w}}\,\phi_\epsilon(a)\le1$. Hence
$s^{(9/8)}(\bar{\mathsf{w}})=s(\tfrac98\bar{\mathsf{w}})\le s(\bar{\mathsf{w}})+1$, and the
bound $s(\bar{\mathsf{w}})\le s(1)+\lceil\log_L\bar{\mathsf{w}}\rceil$ established above gives
\begin{equation*}
s^{(9/8)}(\bar{\mathsf{w}})\le s(\bar{\mathsf{w}})+1\le
s(1)+\lceil\log_L\bar{\mathsf{w}}\rceil+1 .
\end{equation*}

\emph{The smooth core entry.} Let $v$ be smooth, $p=p_v$. The smoothed core bound gives
$|\mathfrak{i}^{(j)}_y|\le2\epsilon^2\mathfrak{N}_p$, with $\mathfrak{N}_p$ as in
\eqref{5.2:eq-smoothed-core-bound-a}. Its parameters are
\begin{equation*}
\nu_p=1+p-j=b_v+1\le\nu,\qquad t_p=L^{-b_v},\qquad \bar\rho_{p-j}=\bar\rho_{b_v}.
\end{equation*}
Square \eqref{5.2:eq-smoothing-depth-1} and insert it:
\begin{equation*}
\begin{split}
2\epsilon^2\mathfrak{N}_p\le{}&2c_\epsilon^2(C_I+2K(d))E_0\,\bar{\mathsf{w}}_v^2(2a_v+11)^2
\bar\rho_{a_v}^2(b_v+1)^3\bar\rho_{b_v}^2(1+\bar\rho_{b_v})\\
&\times L^{-4(a_v-1)}t_p^{4+\hat\alpha}.
\end{split}
\end{equation*}
Since $t=L^{-(a_v+b_v)}$, we have $t_p=L^{a_v}t$. Therefore
\begin{equation*}
L^{-4(a_v-1)}t_p^{4+\hat\alpha}=L^{-4(a_v-1)}L^{(4+\hat\alpha)a_v}t^{4+\hat\alpha}
=L^4t^4L^{-\hat\alpha b_v},
\end{equation*}
and $2\epsilon^2\mathfrak{N}_p$ is at most the smooth-core entry of
\eqref{5.2:eq-smoothing-depth-a}, with $C^{\mathrm{sm}}=2c_\epsilon^2L^4(C_I+2K(d))$.

\emph{The rough core bound.} Let $v$ be rough, and consider the sets $X\ni y$ with
$X\subset\square^{\sim\nu}$. None of them is a tail, since $X\subset\square^{\sim\nu}$. The
singles $F(X,\cdot)-F(X,1)$ entering the core have prefactor
$\mathfrak{n}_v=E_0e^{-\kappa d_j(X)}$ (Definition~\ref{5.2:def-plain-interior-units}); write
$\mathfrak{n}(X)$ for it. By the
summation of localisation domains containing a given point, \cite[(1.26)]{B-CMP116}, there is a
number $C(\kappa)$, depending only on $d$ and $\kappa$, with
\begin{equation*}
\sum_{X\ni y}e^{-\kappa d_j(X)}\bigl(d_j(X)+5\bigr)^2\le C(\kappa) .
\end{equation*}

If $\nu\ge4$, then $2\nu+1\le L^{\nu-2}$ because $L\ge8$. The hull of each such $X$ lies in
$\square^{\sim\nu}$, a box of $2\nu+1\le L^{\nu-2}$ cubes of $\pi_j$ a side, so each such $X$
is short, and its entry in \eqref{5.2:eq-smoothing-depth-a} is the short-single entry. Using
$D_X\le4M(d_j(X)+5)$ and $\rho_{n,j}\le\bar\rho_{\nu-1}\le\bar\rho_\nu$ (the defining property
of the majorant $\bar\rho$ of Notation~\ref{5.1:not-coupling-functions}, which is
non-decreasing),
\begin{equation*}
\begin{split}
\sum_{X\ni y,\ X\subset\square^{\sim\nu}}\mathfrak{N}^{\flat}(X)
&\le2(4Mc_{\vee}C_M)^2\bar{\mathsf{w}}_v^2\bar\rho_\nu^2t^4E_0
\sum_{X\ni y}e^{-\kappa d_j(X)}\bigl(d_j(X)+5\bigr)^2\\
&\le2C(\kappa)(4Mc_{\vee}C_M)^2\,\bar{\mathsf{w}}_v^2E_0\nu^2\bar\rho_\nu^2t^4,
\end{split}
\end{equation*}
since $\nu\ge1$.

If $\nu\le3$, the short-single and long-single entries of
\eqref{5.2:eq-smoothing-depth-a}, each at most $2\mathfrak{n}(X)$, are used. Since
$t=L^{1-\nu}\ge L^{-2}$, we have $1\le L^8t^4$, and since every cell weight $w(\Delta)$ of
\eqref{5.2:eq-plain-interior-units-a} is at least $1$, we have
$\bar{\mathsf{w}}_v=3w^{\vee}(\Delta_y)\ge3$; hence
\begin{equation*}
\sum_{X\ni y,\ X\subset\square^{\sim\nu}}\mathfrak{N}^{\flat}(X)\le2E_0
\sum_{X\ni y}e^{-\kappa d_j(X)}\le2C(\kappa)L^8\,\bar{\mathsf{w}}_v^2E_0t^4,
\end{equation*}
and, since $\nu^2\bar\rho_\nu^2\ge1$ (as $\bar\rho\ge1$,
Notation~\ref{5.1:not-coupling-functions}), the factor $2C(\kappa)L^8$ is
absorbed in $C^{\mathrm{mar}}$.

Finally, by the definition of $\mathfrak{a}'_v$, the Wilson term of
$\mathfrak{N}^{\mathrm{mar}}(v)$ satisfies
\begin{equation*}
\begin{split}
6L^4(b_++r)\,\mathfrak{a}_v'^{\,2}\,t^4
&=6L^4c_{\mathrm{reg}}^2\alpha_{*,n}^2\,\bar{\mathsf{w}}_v^2(b_++r)t^4\\
&\le6L^4(c_{\mathrm{reg}}\sup_n\alpha_{*,n})^2\,\bar{\mathsf{w}}_v^2(b_++r)t^4 .
\end{split}
\end{equation*}
This gives \eqref{5.2:eq-smoothing-depth-2}.
\end{proof}

\begin{proposition}[Holomorphy and the value bound on the carrier]\label{5.2:prop-carrier-value-bound}
Let the parameters obey the floors of the correlation theorem,
Theorem~\ref{5.1:thm-correlation-theorem}, together with the two further floors attached to
$\kappa_1$ and to $w$ in the order of choice of the constants fixed in
Definition~\ref{5.1:def-admitted-inputs}. Let $\gamma\le\gamma_J$, and let
$\gamma$ obey the ceilings of the three-family bound, the regularity ceiling of
Lemma~\ref{5.2:lem-carrier-regularity} and the interior ceiling in
\eqref{5.2:eq-plain-exterior-sum-a}. Let the resummation site be run under the reading convention of
Definition~\ref{5.2:def-plain-interior-units}, by which the values $f_v$ of the pointed units are
defined, and assume the statements listed as hypotheses of Lemma~\ref{5.2:lem-carrier-domain}.
Write $\mathfrak{P}\ni m=(\zeta,\varpi,\vec{\mathsf{s}})$ for the parameter polydisc,
$\operatorname{chain}(m)$ for the carrier pair and $\mathcal{U}(m)$ for its $U$-slot.

Then, for every $K$, every $k\le K$ and every label, the following holds for every pointed unit
$v\in\mathfrak{L}^{\mathrm{pt}}$: the map $m\mapsto f_v(m)$ is holomorphic on $\mathfrak{P}$, and
\begin{equation}\label{5.2:eq-carrier-value-bound-1}
|f_v(m)-f_v(1)|\le\mathfrak{N}^{\flat}(v)\qquad(m\in\mathfrak{P}),
\end{equation}
where $\mathfrak{N}^{\flat}(v)$ is the entry of the size table \eqref{5.2:eq-smoothing-depth-a}
for the class of $v$.
\end{proposition}

\begin{proof}
\emph{Holomorphy.} The singles and the tails are read directly on the carrier pair: their value is
the family itself evaluated at $\operatorname{chain}(m)$. For these units holomorphy is the content
of the two analyticity statements for families read on the carrier, which are among the
hypotheses of Lemma~\ref{5.2:lem-carrier-domain}, applied on the unit's own
analyticity domain. These statements apply because
of the hypothesis of Lemma~\ref{5.2:lem-carrier-domain} on the placement of families. That hypothesis
says
that the $\mathbf{E}$- and $\mathbf{R}$-families with $X\cap Z\neq\emptyset$ belong to the consumer
class $\mathcal{C}$. It also says that the frozen families and the $\mathbf{D}$-families share one
domain, $U^c_j(X,a_j)$.

For a Wilson piece, $A(c,\cdot)$ is an entire function of the plaquette variables, so there is
nothing to prove. For a smooth core, let $p=p_v$ be the smoothing level of
Definition~\ref{5.2:def-smoothing-depth} and $\epsilon$ the smallness parameter of
Lemma~\ref{5.2:lem-smoothed-core-bound}. After a gauge transformation, the restriction of
$\mathcal{U}(m)$ to $Q_v$ lies in the ball on which the hypothesis of
Lemma~\ref{5.2:lem-smoothed-core-bound} holds, namely within $c''$ times the polydisc of radius
$\mathfrak{r}_p$ over the box. The reason is that $\epsilon\le 1$, by the choice of the smoothing
depth in
Definition~\ref{5.2:def-smoothing-depth}. On that ball the extension $\operatorname{ext}_p$ of
Lemma~\ref{5.2:lem-smoothed-core-bound} is an admissible background, analytic in its argument.
Hence $f_v$ is holomorphic in $m$. A rough
core is a sum of terms of the preceding kinds, and its holomorphy holds termwise.

\emph{The bound \eqref{5.2:eq-carrier-value-bound-1}.} We go through the classes of the size table
\eqref{5.2:eq-smoothing-depth-a}.

\emph{Short single.} Here $f_v(m)=F(X,\operatorname{chain}(m))$ and $f_v(1)=F(X,1)$, with
$|F|\le\mathfrak{n}_v$; write $t=L^{j-n}$, $t':=Lt$, and take for the regularity constant
$\mathfrak{a}'$ of Lemma~\ref{5.2:lem-local-gauge-bound} the value $\mathfrak{a}'_v$.
Take for the box $S$ of Lemma~\ref{5.2:lem-local-gauge-bound} the hull
$\hat X$ of the unit. Read in level-$j$ units, clause (i) of Lemma~\ref{5.2:lem-carrier-regularity}
is precisely the hypothesis of Lemma~\ref{5.2:lem-local-gauge-bound} on the pair
$(\mathbf{U},\mathbf{J})$ restricted to $S$. The side of $S$ obeys $D\le D_X$.
Lemma~\ref{5.2:lem-local-gauge-bound}
therefore gives
\begin{equation*}
|F(X,(\mathbf{U},\mathbf{J}))-F(X,1)|
\le 2\mathfrak{n}_v\bigl(16\,d\,D_X\,t'^2\mathfrak{a}'_v/\mathfrak{r}_j\bigr)^2 .
\end{equation*}
The constant $c_\vee$ of the size table converts the right member into the table's entry for short
singles. The trivial bound $2\mathfrak{n}_v$ accounts for the minimum with $1$ in that entry.

\emph{Long single and tail.} On the unit's analyticity domain $\mathfrak{D}_{\mathfrak{E}}$ we have
$|F|\le\mathfrak{n}_v$. Both values in \eqref{5.2:eq-carrier-value-bound-1} are therefore bounded
by $\mathfrak{n}_v$, and their difference is at most $2\mathfrak{n}_v$. This is the table's entry.

\emph{Wilson piece.} Let $\eta$ be the spacing of the lattice carrying $\mathcal{U}$, and let
$\ell_j$ be the side of a level-$j$ unit. The bound rests on the following facts.
\begin{itemize}
\item Near $\mathbb{1}$ the plaquette density $a$ satisfies $|a(W)|\le|W-1|^2$.
\item The support of the piece lies within one level-$j$ unit of $y$, or, for a block piece, in a
block. Such a block consists of cells of levels at least $n-1$ at distance $\operatorname{dist}_n$
at most $1$ from $\Delta_y$.
\item The first entry of \eqref{5.2:eq-carrier-domain-b}, read pointwise, gives for every
plaquette of the support, with $W$ its plaquette variable,
\begin{equation*}
|W-1|\le\mathfrak{a}'_v\,(\eta/\ell_{n-1})^2 .
\end{equation*}
\item A level-$j$ cell contains $6(\ell_j/\eta)^4$ plaquettes.
\end{itemize}
Since $\ell_{n-1}=L^{\,n-1-j}\ell_j$, that is $\ell_j/\ell_{n-1}=Lt$, these facts give for one
level-$j$ cell
\begin{equation*}
\sum|a(W)|\le 6(\ell_j/\eta)^4\,\mathfrak{a}'^2_v\,(\eta/\ell_{n-1})^4=6L^4t^4\mathfrak{a}'^2_v ,
\end{equation*}
the sum running over the plaquettes of the cell. This is the product $6L^4\mathfrak{a}'^2_vt^4$ of
the table's entry for $b_jA(h_y,\cdot)$, where it is multiplied by $(b_++r)$. For the row-five
(Wilson) pieces the same count, with $\mathfrak{a}'_v=c_{\mathrm{reg}}\bar{\mathsf{w}}_v\alpha_{*,n}$,
enters the table through the constant $c_5^{W}$, whose definition contains the factor
$6L^4(c_{\mathrm{reg}}\alpha_{*,n}/\mathfrak{a}_n)^2/C_S$.

\emph{Smooth core.} Here $p=p_v$. The hypothesis $\epsilon\le 1$ of
Lemma~\ref{5.2:lem-smoothed-core-bound} holds by the choice of the smoothing depth in
Definition~\ref{5.2:def-smoothing-depth}. The core vanishes at the trivial carrier, so
$f_v(1)=0$, and that lemma bounds $|f_v(m)|$ by $2\epsilon^2\mathfrak{N}_p$. Insert the bound on
$\epsilon$ of Definition~\ref{5.2:def-smoothing-depth}. Here $p=n-a_v$ and $t=L^{j-n}$, so that, with
$b_v=p-j$,
\begin{equation*}
L^{-4(a_v-1)}L^{(4+\hat\alpha)a_v}t^{4+\hat\alpha}=L^4t^4L^{-\hat\alpha b_v},
\end{equation*}
and this gives the table's entry for smooth cores, with
$C^{\mathrm{sm}}=2c_\epsilon^2L^4(C_I+2K(d))$.

\emph{Rough core.} The table's entry is the sum of the entries of its constituents
$X\ni y$, $X\subset\square^{\sim\nu}$, and of its summand $6L^4(b_++r)\mathfrak{a}'^2_vt^4$, the
entry for $b_jA(h_y,\cdot)$. The bound holds termwise, by the
cases above.
\end{proof}

The argument uses no bound that is established only at the real point, no identity of the
$\mathbf{J}$-slot at any configuration other than the trivial one $\mathbb{1}$, and no H\"older
norm of a layer.

\begin{proposition}[The depth gain for deep pointed units]\label{5.2:prop-deep-unit-gain}
Assume the following statements.
\begin{itemize}
\item[(a)] Clause (a) of the decay lemma for decoupled families, with its constant $C_{\mathrm{dec}}$:
for a unit whose size $\mathfrak{N}$ is the supremum of its modulus over its analyticity domain
$\mathfrak{D}_{\mathfrak{E}}$ and whose hull has depth $D$, and for $m,m'\in\mathfrak{P}$ with the
same $(\zeta,\varpi)$, the values of the unit at $m$ and at $m'$ differ by at most
$C_{\mathrm{dec}}e^{-\frac12\delta_PD}\,\mathfrak{N}$.
\item[(b)] The disc lemma for the decoupled carrier: for a unit $F(X,\cdot)$ the map
$\hat\lambda\mapsto F(X,P_{\hat\lambda})$ is holomorphic on the disc $|\hat\lambda|<\rho_X$. The
radius $\rho_S$ is non-increasing in the set $S$, and for the hull $\hat X_v$ of a deep
$v\in\mathfrak{L}^{\mathrm{pt}}$ it satisfies
\begin{equation}\label{5.2:eq-deep-unit-gain-3}
\rho_{\hat X_v}\ge r^0e^{\frac12\delta_P\hat D_v}.
\end{equation}
\item[(c)] The endpoint identification for the decoupled carrier: for $m,m'\in\mathfrak{P}$ with the
same $(\zeta,\varpi)$, the decoupled carrier $P_{\hat\lambda}$ formed for the pair $m,m'$ reduces at
$\hat\lambda=0$ to the carrier at $m'$ and at $\hat\lambda=1$ to the carrier at $m$.
\item[(d)] Clause (ii) of the source-polydisc statement, which bounds the pieces of a function
holomorphic on a product polydisc by a bound on its deviation from a centre value.
\item[(e)] Clause (a) of the Schwarz-type bound: if $G$ is holomorphic on $|\hat\lambda|<\rho$ with
$\rho>1$ and $|G|\le\mathfrak{N}$ there, then $|G(1)-G(0)|\le 2\mathfrak{N}/\rho$.
\end{itemize}
Put
\begin{equation}\label{5.2:eq-deep-unit-gain-1}
C^{\mathrm{pt}}:=\max\{C_{\mathrm{dec}},\,2/r^0,\,2\}.
\end{equation}
Then the following holds for every deep $v\in\mathfrak{L}^{\mathrm{pt}}$ and all
$m,m'\in\mathfrak{P}$ that have the same $(\zeta,\varpi)$:
\begin{equation}\label{5.2:eq-deep-unit-gain-2}
|f_v(m)-f_v(m')|\le C^{\mathrm{pt}}e^{-\frac12\delta_P\hat D_v}\,\mathfrak{N}^{\flat}(v).
\end{equation}
Here $\mathfrak{N}^{\flat}(v)$ is the size given by the table of
Definition~\ref{5.2:def-smoothing-depth}, display \eqref{5.2:eq-smoothing-depth-a}.
Consequently the piece bound for deep units in \eqref{5.2:eq-plain-exterior-sum-a} holds. The
constant $C^{\mathrm{pt}}$ and the bound \eqref{5.2:eq-deep-unit-gain-2} do not depend on the age of
any large-field region, on the number $N(g)$ of preparatory stages, the number $N_0(g)$ of zones,
or the radius $\check{R}_k$, or on any slot variable.
\end{proposition}

\begin{proof}
Write $\rho:=\rho_{\hat X_v}$, let $j$ be the level of $v$, and put $y:=y_v$. By hypothesis (b),
$\rho$ obeys the lower bound \eqref{5.2:eq-deep-unit-gain-3}, and $\rho\le\rho_S$ for every set
$S\subset\hat X_v$.

\emph{The case $\hat D_v=0$ or $\rho\le 1$.} Apply Proposition~\ref{5.2:prop-carrier-value-bound}
at both points. It gives $|f_v(m)-f_v(1)|\le\mathfrak{N}^{\flat}(v)$ and
$|f_v(m')-f_v(1)|\le\mathfrak{N}^{\flat}(v)$. The triangle inequality then gives
$|f_v(m)-f_v(m')|\le 2\mathfrak{N}^{\flat}(v)$.

If $\hat D_v=0$, then $2\le C^{\mathrm{pt}}=C^{\mathrm{pt}}e^{0}$, and
\eqref{5.2:eq-deep-unit-gain-2} follows. If $\rho\le1$, then
$r^0e^{\frac12\delta_P\hat D_v}\le 1$ by \eqref{5.2:eq-deep-unit-gain-3}, hence
\begin{equation*}
2\le (2/r^0)\,e^{-\frac12\delta_P\hat D_v}\le C^{\mathrm{pt}}e^{-\frac12\delta_P\hat D_v},
\end{equation*}
where the second inequality holds because $2/r^0\le C^{\mathrm{pt}}$ by
\eqref{5.2:eq-deep-unit-gain-1}. Again \eqref{5.2:eq-deep-unit-gain-2} follows.

There remains the case $\hat D_v>0$ and $\rho>1$, with the hull $\hat X_v$ contained in the core
$\mathsf{C}$. We treat the entries of the table in turn.

\emph{Long single units and tails.} For these entries the size is the supremum of $|F|$ over
$\mathfrak{D}_{\mathfrak{E}}$, which is at most $\mathfrak{n}_v$. Apply hypothesis (a) with
$\mathfrak{N}=\mathfrak{n}_v$ and depth $D=\hat D_v$. It bounds the difference by
$C_{\mathrm{dec}}e^{-\frac12\delta_P\hat D_v}\mathfrak{n}_v$. This is at most the right member of
\eqref{5.2:eq-deep-unit-gain-2}, because $C_{\mathrm{dec}}\le C^{\mathrm{pt}}$ and because the table
gives $\mathfrak{N}^{\flat}(v)=2\mathfrak{n}_v\ge\mathfrak{n}_v$ for these entries.

\emph{Short single units.} Let $X$ be the support of the unit and put
\begin{equation*}
G(\hat\lambda):=F(X,P_{\hat\lambda})-F(X,1).
\end{equation*}
The properties of $G$ used below are these.
\begin{itemize}
\item By hypothesis (b), $G$ is holomorphic on $|\hat\lambda|<\rho_X$. Since the hull $\hat X_v$
contains $X$, hypothesis (b) gives $\rho\le\rho_X$, so $G$ is holomorphic on $|\hat\lambda|<\rho$.
\item On this disc $|G|\le\mathfrak{N}^{\flat}(v)$. This follows from clause (ii) of
Lemma~\ref{5.2:lem-carrier-regularity} together with Lemma~\ref{5.2:lem-local-gauge-bound}; the
short-single entry of the table is the smaller of $2\mathfrak{n}_v$ and the bound of that lemma with
$D\le D_X$, the constants converted by the constant $c_\vee$ of
Definition~\ref{5.2:def-smoothing-depth}.
\item By hypothesis (c), $G(0)$ and $G(1)$ are the values at $m'$ and at $m$, each diminished by the
same number $F(X,1)$. Hence $G(1)-G(0)=f_v(m)-f_v(m')$.
\end{itemize}
Hypothesis (e) now gives $|f_v(m)-f_v(m')|\le 2\mathfrak{N}^{\flat}(v)/\rho$. By
\eqref{5.2:eq-deep-unit-gain-3} and \eqref{5.2:eq-deep-unit-gain-1},
\begin{equation*}
\frac{2\mathfrak{N}^{\flat}(v)}{\rho}\le\frac{2}{r^0}e^{-\frac12\delta_P\hat D_v}\mathfrak{N}^{\flat}(v)
\le C^{\mathrm{pt}}e^{-\frac12\delta_P\hat D_v}\mathfrak{N}^{\flat}(v),
\end{equation*}
which is \eqref{5.2:eq-deep-unit-gain-2}.

\emph{Wilson pieces.} The argument is the one just given for short single units. The only change is
the bound on the disc: together with clause (ii) of Lemma~\ref{5.2:lem-carrier-regularity}, the first
entry of the pointwise increment bound for the plaquette variables of the carrier, which enters the
Wilson case of Proposition~\ref{5.2:prop-carrier-value-bound}, holds for $P_{\hat\lambda}$ on the
whole disc $|\hat\lambda|<\rho$ and replaces Lemma~\ref{5.2:lem-local-gauge-bound} in the bound
$|G|\le\mathfrak{N}^{\flat}(v)$.

\emph{Smooth cores.} Let $p:=p_v$ and let $\operatorname{ext}_p$ be as in
Lemma~\ref{5.2:lem-smoothed-core-bound}. Let $\mathcal{U}_{\hat\lambda}$ denote the $U$-slot of
$P_{\hat\lambda}$ restricted to $Q_v$, and put
\begin{equation*}
G(\hat\lambda):=\mathfrak{i}^{(j)}_y\bigl(\operatorname{ext}_p(\mathcal{U}_{\hat\lambda})\bigr).
\end{equation*}
By clause (ii) of Lemma~\ref{5.2:lem-carrier-regularity} and the bound
\eqref{5.2:eq-local-gauge-bound-a}, for $|\hat\lambda|<\rho$ the field $\mathcal{U}_{\hat\lambda}$
satisfies, modulo gauge, the hypothesis of Lemma~\ref{5.2:lem-smoothed-core-bound} with $\epsilon\le1$;
it therefore stays in the open ball on which $\operatorname{ext}_p$ is analytic, as in the holomorphy
part of Proposition~\ref{5.2:prop-carrier-value-bound}. Hence $G$ is holomorphic on
$|\hat\lambda|<\rho$, and Lemma~\ref{5.2:lem-smoothed-core-bound} gives
$|G|\le2\epsilon^2\mathfrak{N}_p$ on the whole disc. By the bound on $\epsilon$ in
Definition~\ref{5.2:def-smoothing-depth} and the identity
\begin{equation*}
L^{-4(a_v-1)}L^{(4+\hat\alpha)a_v}t^{4+\hat\alpha}=L^4t^4L^{-\hat\alpha b_v}
\end{equation*}
recorded with the table there, $2\epsilon^2\mathfrak{N}_p$ is at most the smooth-core entry
$\mathfrak{N}^{\flat}(v)$ of the table; hence $|G|\le\mathfrak{N}^{\flat}(v)$ on the whole disc. By
hypothesis (c) the carrier at $\hat\lambda=0$ and at $\hat\lambda=1$ is the carrier at $m'$ and at
$m$, so $G(0)=f_v(m')$ and $G(1)=f_v(m)$ by the definition of the pointed value of a smooth core
through $\operatorname{ext}_p$ applied to the whole slot. Hypothesis (e) then bounds
$|f_v(m)-f_v(m')|$ by $2\mathfrak{N}^{\flat}(v)/\rho$, and the last display of the short case
turns this into \eqref{5.2:eq-deep-unit-gain-2}.

\emph{Rough cores.} For a rough core the table sets $\mathfrak{N}^{\flat}(v)=\mathfrak{N}^{\mathrm{mar}}(v)$,
the sum of the sizes $\mathfrak{N}^{\flat}(X)$ of the constituents $X\ni y$, $X\subset\square^{\sim\nu}$,
and of the Wilson term $6L^4(b_++r)\mathfrak{a}'^2_vt^4$ of the table of
Definition~\ref{5.2:def-smoothing-depth}. Apply the preceding arguments to each term
separately, on the disc $|\hat\lambda|<\rho$, and add. This is legitimate: the hull $\hat X$ of
each constituent satisfies $\hat X\subset\hat X_v$, so $\rho\le\rho_{\hat X}$ by hypothesis (b);
hence the function $G$ of each constituent is holomorphic on $|\hat\lambda|<\rho$, and the short or
the Wilson argument applies to it.

\emph{Pieces.} Apply hypothesis (d) on the product polydisc, with centre value
$c:=f_v(\operatorname{chain}(\vec 0))$; the bound fed into it is \eqref{5.2:eq-deep-unit-gain-2} read at the
centre, that is, with $c$ in place of $f_v(m')$. It gives the piece bound for deep units in
\eqref{5.2:eq-plain-exterior-sum-a}.

For each piece the Cauchy disc used is the disc in $\hat\lambda$. Nested inside it is the disc in
$\sigma$ of Lemma~\ref{5.2:lem-local-gauge-bound} or of Lemma~\ref{5.2:lem-smoothed-core-bound}. No
slot variable enters.
\end{proof}

\begin{lemma}[The column sum of pointed sizes per cell]\label{5.2:lem-column-sum-pointed}
Let the component $C$, the class $\mathfrak{L}^{\mathrm{pt}}$ of pointed units with their points
$y_v$, the level-$k$ closures $K_v$ of their domains and their radii $\operatorname{rad}(v)$, the
rate $\varsigma_1$ and the
window weights $\mathsf{w}^{\mathrm{win},\vee}$ be those of the plain interior units at the
resummation site (Definition~\ref{5.2:def-plain-interior-units} and
\eqref{5.2:eq-plain-interior-units-a}). Let $\mathfrak{N}^{\flat}(v)$ be the sizes of the table
\eqref{5.2:eq-smoothing-depth-a}. For every cell $\Delta'\subset C$ and every
$r\le\frac{8}{5}\kappa$,
\begin{equation}\label{5.2:eq-column-sum-pointed-a}
\sum_{v\in\mathfrak{L}^{\mathrm{pt}}:\ y_v\in\Delta'}\mathfrak{N}^{\flat}(v)\,
e^{r\,d_k(K_v)}\,e^{\varsigma_1\operatorname{rad}(v)}
\le\mathsf{K}^{\flat}\,\mathsf{w}^{\mathrm{win},\vee}(\Delta')^{3}.
\end{equation}
Here $\mathsf{K}^{\flat}$ is a number that does not depend on $g$ and is fixed after $M$.
\end{lemma}

\begin{proof}
Throughout, $\bar{\mathsf{w}}$ stands for the weight $\bar{\mathsf{w}}_v$ of
Definition~\ref{5.2:def-plain-interior-units} of the units pointed in $\Delta'$.
For a unit $v$ of level $j$ and zone $n$ we write $t=L^{j-n}$ and $\nu=1+n-j$.

\emph{Exponential weights of the singles.} Let $v$ be a single with domain $X\ni y=y_v$. Here
$j<k$ and $L\ge 8$. The comparison of tree lengths across levels therefore gives
\[
d_k(K_v)\le\tfrac18\,d_j(X).
\]
As $r\le\frac85\kappa$, it follows that $e^{r d_k(K_v)}\le e^{\frac15\kappa d_j(X)}$. Moreover
$\operatorname{rad}(v)\le d_j(X)+1$, so that
\[
e^{\varsigma_1\operatorname{rad}(v)}\le e^{\varsigma_1}\,e^{\varsigma_1 d_j(X)},
\qquad \varsigma_1=\tfrac14\delta\kappa .
\]
Here and below, the rows are those of the classification of the pointed units in
Definition~\ref{5.2:def-plain-interior-units}: rows~$1$, $2$, $3$ and $\partial$ consist of
singles (row~$1$ of the tails, row~$2$ of the singles whose sizes carry the factor
$\omega_{k,j}$, row~$3$ of those whose sizes carry a power of $g_j$, row~$\partial$ of the
constituents at near points), and row~$5$ consists of the Wilson pieces.
The size of every row decays at least like $e^{-(1-\frac14\delta)\kappa d_j(X)}$. After these two
factors are absorbed, every row keeps the decay
$\exp\bigl(-(\frac45-\frac12\delta)\kappa\,d_j(X)\bigr)$. Since
$\delta\le\delta_J(L)\le\frac1{10}$, we have $(\frac45-\frac12\delta)\kappa\ge\frac35\kappa$.
Of this decay, the portion $\frac1{20}\kappa\,d_j(X)$, with $\frac1{20}\kappa\ge c_d+2$, pays for
the sums over domains by \cite[(1.26)]{B-CMP116}. The portion $\frac12\kappa\,d_j(X)$, with
$\frac12\kappa\ge 5\log L$, is spare.

\emph{Cores and Wilson pieces.} For these units $d_k(K_v)\le 6^d$ and $\operatorname{rad}(v)\le 5$.
The two exponentials therefore contribute at most the factor
$e^{\frac85\kappa 6^d+5\varsigma_1}$, which does not depend on $g$.

\emph{Counting points.} A cell is a cube of $\pi_n$, that is, a cube of side $M$ at level $n$.
Hence $\Delta'$ contains $(ML^{n-j})^4=M^4t^{-4}$ points of $T^{(j)}$.

\emph{Singles at a fixed point.} Let $\mathfrak{n}_j$ denote the prefactor of the row in the
definition of $\mathfrak{n}_v$, namely the factor multiplying the exponential in $d_j(X)$
(Definition~\ref{5.2:def-plain-interior-units}).

Short singles are treated first. Insert $D_X\le 4M(d_j(X)+5)$ into the entry for short singles in
\eqref{5.2:eq-smoothing-depth-a}. The two exponential weights are absorbed into the decay as
above. Then sum over the short singles $v$ of a row with point $y_v=y$, whose domains $X$
contain $y$, against the retained decay, using the constant $C(\kappa)$ of the domain sums of
this chapter. This gives
\begin{align*}
&\sum_{\substack{v\ \mathrm{short}\\ X\ni y}}\mathfrak{N}^{\flat}(v)\,e^{r\,d_k(K_v)}\,
e^{\varsigma_1\operatorname{rad}(v)}\\
&\qquad\le 2C(\kappa)\,(4M\bar{\mathsf{w}}c_{\vee}C_M)^2\,\bar\rho_\nu^{\,2}\,t^4\,
\mathfrak{n}_j .
\end{align*}

Next come the long singles. A long non-tail has $d_j(X)\ge L^{\nu-2}-1$
by \eqref{5.2:eq-plain-interior-units-b}. After the two exponential weights are absorbed, the
spare portion of the decay and the sum over domains, with the constant $K(d)$ of the domain sums
of this chapter, then give, for the sum over the long non-tail singles $v$ with point $y$ and
domain $X\ni y$,
\begin{align*}
\sum_{\substack{v\ \mathrm{long}\\ X\ni y}}\mathfrak{N}^{\flat}(v)\,e^{r\,d_k(K_v)}\,
e^{\varsigma_1\operatorname{rad}(v)}
&\le 2K(d)\,\mathfrak{n}_j\,e^{-\frac12\kappa(L^{\nu-2}-1)_+}\\
&\le C_L\,t^5\,\mathfrak{n}_j ,
\end{align*}
with a constant $C_L$ depending only on $L$ (the last bound uses $\frac12\kappa\ge 5\log L$).

Finally, the tails. For a tail one has $d_j(X)\ge\nu$, by
\eqref{5.2:eq-plain-interior-units-b}. With the row's own weight, the sum over the tails
$v$ with point $y$ and domain $X\ni y$ satisfies
\[
\sum_{\substack{v\ \mathrm{tail}\\ X\ni y}}\mathfrak{N}^{\flat}(v)\,e^{r\,d_k(K_v)}\,
e^{\varsigma_1\operatorname{rad}(v)}
\le 2K(d)\,E_0\,t^5 .
\]

\emph{Sums over the levels.} After multiplication by the $M^4t^{-4}$ points of $\Delta'$, the
factor $t^5$ of the long singles and of the tails becomes $M^4t=M^4L^{-(n-j)}$. The sums over $j$
are performed row by row, as follows.
\begin{itemize}
\item Row~$2$ is summed by $\omega_{k,j}\le L^{-(n-j)/2}$ and the scale-sum bound for $\omega$.
\item Row~$3$ is summed by the summability statement for row~$3$ over the levels.
\item Row~$\partial$ is summed by part~(iii) of the lemma on the constituents of row~$\partial$,
by which at most
$27dM(\nu+2)t^{-3}$ points of $\Delta'$ are near points, together with the scale-sum bound for
row~$\partial$.
\item The contribution of row~$5$ is bounded independently of $K$.
\end{itemize}

\emph{Smooth cores with $a_v=a_0:=2+s(\bar{\mathsf{w}})$.} Recall that
$p_v\le n-2-s(\bar{\mathsf{w}})$, so that $a_v\ge a_0$ for every smooth core. With $b=p_v-j$, the
entry for smooth cores in \eqref{5.2:eq-smoothing-depth-a}, times the $M^4t^{-4}$ points, is
summed over $b$. Apart from the factor $e^{\frac85\kappa 6^d+5\varsigma_1}$ of the exponentials,
the result is at most
\[
C^{\mathrm{sm}}E_0M^4\,\bar{\mathsf{w}}^2(2a_0+11)^2\bar\rho_{a_0}^{\,2}
\sum_{b\ge 0}(b+1)^3\bar\rho_b^{\,2}(1+\bar\rho_b)L^{-\hat\alpha b},
\]
and the series converges.

\emph{Smooth cores with $a_v>a_0$: the rim.} Here
\[
j\le p_v+1\le n-2-s(\bar{\mathsf{w}}).
\]
By the maximality in the definition of $p_v$, the box $Q^{(p_v+1)}=Q^{(n-a_v+1)}$ is therefore not
contained in $Z$. This puts $y$ within
\[
(\nu+2)L^{-(\nu-1)}+3L^{-(a_v-1)}
\]
level-$n$ units of a face of $\partial Z$, where $n=k$. At most $2dM^4t^{-4}$ times this quantity
of points of $\Delta'$ lie within that distance of $\partial Z$. Since $\nu=1+a_v+b_v$, the sum
over $a=a_v$ and $b=b_v$ is
\[
\begin{split}
\sum_{a,b}\bigl[(a+b+3)L^{-(a+b)}+3L^{-(a-1)}\bigr]&(2a+11)^2\bar\rho_a^{\,2}\\
&\times(b+1)^3\bar\rho_b^{\,2}(1+\bar\rho_b)L^{-\hat\alpha b}<\infty .
\end{split}
\]

\emph{Rough cores.} A core is rough when no $p$ with $j\le p\le n-2-s(\bar{\mathsf{w}})$ has
$Q^{(p)}\subset Z$. Hence either $j>n-2-s(\bar{\mathsf{w}})$, that is
$\nu\le 2+s(\bar{\mathsf{w}})$, or $Q^{(j)}=\square^{\sim\nu+3}\not\subset Z$.

In the first case the rough cores contribute at most
$C^{\mathrm{mar}}\bar{\mathsf{w}}^2(3+s)^4$ in all, with $s=s(\bar{\mathsf{w}})$. In the second
case at most $2d(\nu+4)M^4t^{-3}$ points $y$ of $\Delta'$ satisfy
$\square^{\sim\nu+3}\not\subset Z$, and
\[
\sum_{\nu}(\nu+4)\,\nu^2\bar\rho_\nu^{\,2}L^{-(\nu-1)}<\infty .
\]

\emph{Conclusion.} By Definition~\ref{5.2:def-smoothing-depth},
$s(\bar{\mathsf{w}})\le s(1)+\log_L\bar{\mathsf{w}}$. Moreover $\bar{\mathsf{w}}\ge 3$, since
$\bar{\mathsf{w}}_v=3w^{\vee}(\Delta_{y_v})$ with $w^{\vee}\ge 1$. Hence the polynomials in $a_0$
and $s$ occurring above are at most $c\,\bar{\mathsf{w}}$, with $c$ independent of $g$, and every
contribution above is at most a number independent of $g$, fixed after $M$, times
$\bar{\mathsf{w}}^3$. Collecting the contributions, the left member of
\eqref{5.2:eq-column-sum-pointed-a} is at most such a number times $\bar{\mathsf{w}}^3$. Since
$\bar{\mathsf{w}}\le 3\mathsf{w}^{\mathrm{win},\vee}(\Delta')$, this gives
\eqref{5.2:eq-column-sum-pointed-a}.

Every sum over scales above is bounded independently of $K$. The rim and the layer of rough cores
grazing $\partial Z$ are paid for by their thinness, as row~$\partial$ is.
\end{proof}

\begin{definition}[Exposed and compact pointed units of the exterior classes]
\label{5.2:def-exposed-compact-units}
Fix $K$, a level $k\le K$ and the current global label, and consider the localisation site of the
step $\mathbf{R}_k$. Throughout, $Z$ denotes, as in \cite[\S1]{B-CMP122a}, the union of a class of
components of $\Lambda_k^c$ (the treated ones), and not a single component. With
$Z_k:=\Lambda_k^c$, the large-field operation factorises as
$\mathbf{T}_k(Z_k)=\mathbf{T}_k(Z_k\cap Z^c)\,\mathbf{T}_k(Z)$ \cite[\S1]{B-CMP122a}. $C$ is the
component under treatment, and $\mathfrak{T}:=\Omega_k\cup Z^c$.

We recall the exterior classes of the lemma on the sum of the exterior pieces. Consider the
decomposition of the plain units into classes~(c) and~(d). The class $\mathfrak{L}^{c}$ is formed of
plain units $v$ of class~(c) whose domain lies astride $\partial Z$:
\begin{equation*}
X_v\cap C\neq\emptyset,\qquad K_v=[X_v]_k\not\subset Z .
\end{equation*}
The class $\mathfrak{L}^{Y_0}$ is formed of the $\mathsf{t}$-pieces $v_p$ of the plain units of
classes~(c) and~(d). A unit $v$ of these classes has a level $j\le k$, with $X_v\in\mathbf{D}_j$, and a
point $y_v\in X_v$, as for the plain pointed units on $C$.

\begin{enumerate}
\item The table of sizes assigns to each pointed unit a size (the powers $t^4$, $t^2$,
$t^{4+\hat\alpha}$). This size is a stronger bound than the unit's own supremum $\mathfrak{n}_v$ on its
own analyticity domain, and that supremum is all that containment at the unit's own family yields. On
$C$ the additional smallness is supplied by five results:
\begin{itemize}
\item the current bound;
\item the gauges of the gauge cubes of the gauge-cube lemma on $C$;
\item the regularity lemma;
\item the marginal bound for singles;
\item the integral bound for smooth cores.
\end{itemize}
A unit is called \emph{exposed} if this additional smallness must be supplied off $C$ as well.

In $\mathfrak{L}^c$ the exposed units are:
\begin{itemize}
\item the short singles of types~$2$ and~$3$ of the table of sizes;
\item the short singles among the near constituents of type~$\partial$;
\item the Wilson pieces (type~$5$, the terms $b_jA(h_y,\cdot)$).
\end{itemize}
The following are not exposed:
\begin{itemize}
\item the units of type~$4$: their size is their norm, and containment at their own analyticity
domain suffices;
\item the tails and the long singles: their size is $\mathfrak{N}^{\flat}=2\mathfrak{n}_v$, the
trivial bound.
\end{itemize}
No core lies in $\mathfrak{L}^c$.

\item In $\mathfrak{L}^{Y_0}$ every pointed $p$ is exposed, namely:
\begin{itemize}
\item the $\mathsf{t}$-piece of a single, a Wilson piece or a $\partial$-constituent of a plain unit
of class~(c);
\item the $\mathsf{t}$-piece of a single, a Wilson piece, a $\partial$-constituent or a far core of a
plain unit of class~(d).
\end{itemize}
Each such $p$ is taken at every tuple of the later links, the all-empty tuple included. At the
all-empty tuple the piece is $v_p$ itself, evaluated at $\operatorname{chain}(\vec 0)$.

For class~(d) one has $X\cap Z=\emptyset$, and the point $y$ may lie in any zone $n\le k$ of the
current global label. These are the terms $\mathbf{E}^{(j)}(X,U''_k,z)$ with $X\subset Z^c$ and point
$z\in\Omega_n\setminus\Omega_{n+1}$, $n\ge j$, of \cite[p.~372]{B-CMP122b}; the point $z$ there is the
point $y$ here. Such a unit sits in one of three places:
\begin{itemize}
\item on $\Omega_k$;
\item on the shells of levels $n<k$ of an untreated component off $Z$;
\item in or beside the frozen window of a \emph{foreign live arrest}.
\end{itemize}
A foreign live arrest is an arrested attempt $\mathfrak{a}'$ on a component of $\Lambda_k^c$ that is
now untreated, so that $X_{\mathfrak{a}'}\subset Z^c$.

\item A pointed $p$ of $\mathfrak{L}^{Y_0}$, and its unit $v$, is:
\begin{itemize}
\item \emph{compact} if $v$ is a short single, a Wilson piece or a core (smooth or rough);
\item \emph{long} if $v$ is a long single or a tail.
\end{itemize}
For long units the hull is $\hat X_v:=X$. They correspond to the second case of
\cite[p.~372]{B-CMP122b}: the terms with localisation domains not contained in $\square^{\sim2}$,
which are treated there through the identity \cite[(1.60)]{B-CMP122b}. For these terms the
localisation family consists of the cubes disjoint with $X$, which takes the place of $\square_0$
\cite[p.~372]{B-CMP122b}.

The following objects are introduced below: $\bar{\mathsf{w}}_v$, $\mathfrak{r}_v$,
$\mathfrak{h}_v$, $\operatorname{lab}_1$, $X^{\oplus}_v$ and $\mathfrak{E}^{\oplus}_v$. They, and the
geometric and analytic statements of the carrier bound that use them, are defined only for the
exposed units of $\mathfrak{L}^c$ and for the compact $p$.

A long $p$ is different. It is controlled on its unit's own analyticity domain on all of $X_v$; it
carries the trivial size $\mathfrak{N}^{\flat}=2\mathfrak{n}_v$ and no weight. It is treated
separately, in the clause on long singles and tails of the carrier bound.

\item Let $v$ be such a unit; this includes the unit underlying a compact $p$. Put $y:=y_v$, and let
\begin{equation*}
n:=\pi(y)
\end{equation*}
be the zone index of $y$ in the current global label $\{\Omega''_i\}_i$, which is the label of the
whole lattice at this site. The index $n$ may be any integer in $[j,k]$, and:
\begin{itemize}
\item For units of types $1$--$3$ and $\partial$ whose domain meets $Z$ one has $j<k$
\cite[\S2]{B-CMP119}.
\item In class~(d) one has $j\le k$. At $j=k$ the units of types $1$, $3$, $\partial$ decay at the
rate $2\kappa$ of the lemma on the decay fraction $\delta_J$, type~$2$ is void, and every single is
long or a tail.
\item For the units astride $\partial C$ one has $n=k$, by clause~(i) of the geometry of the
enlarged consumer.
\end{itemize}
The following objects are defined exactly as in the definition of the plain pointed units on $C$ and
in the core lemma on $C$:
\begin{itemize}
\item the cell $\Delta_y$ of $y$, of level $n$;
\item $t$, $\nu$, $\square$, $\hat X$, $Q^{(p)}$, $a_v$, $b_v$;
\item the distinctions short/long and smooth/rough;
\item the hull $\hat X_v$ and $\operatorname{rad}(v)$.
\end{itemize}
The only change is that $C$'s label is replaced by the global one; cells, $\pi_n$ and
$\mathrm{dist}_n$ are measured in $\pi$-cubes of the level named. Moreover
\begin{equation*}
\bar{\mathsf{w}}_v:=3\,w^{\vee}(\Delta_y),
\end{equation*}
where $w^{\vee}(\Delta_y)$ is the maximum of the threshold weight $w(\Delta)$ over the cells $\Delta$
at $\mathrm{dist}_n\le1$ from $\Delta_y$. Here $w(\Delta)$ is the cell weight entering the thresholds
$\theta_a$ of the analyticity clauses of the current bound, not the complex source
$w=\sum_i z_if_i$.

There is one exception. For a far core of class~(d) one sets
\begin{equation}\label{5.2:eq-exposed-compact-units-1}
p_v:=n-2-s^{(9/8)}(\bar{\mathsf{w}}_v),\qquad
s^{(9/8)}(\bar{\mathsf{w}}):=s\bigl(\tfrac98\,\bar{\mathsf{w}}\bigr).
\end{equation}
Here $s(\cdot)$ is the smoothing-depth function of the core lemma on $C$. It is a function of the
weight, unrelated to the bracket index. The core is smooth if $p_v\ge j$ and rough otherwise.

One value of $p_v$ serves $v$, its $\mathsf{t}$-piece and its subtracted copy. As a result, the disc of
the integral bound holds on the $\tfrac98$-dilated carrier as well.

No clause $Q^{(p)}\subset Z$ is imposed. The box $Q_v$ lies in the one-layer enlargement of
$[\square_0]_k$. It is kept inside the domain of the auxiliary consumer defined below, through
$\mathfrak{h}_v\subset\operatorname{lab}_1$, and not inside $Z$.

\item The weight is as follows.
\begin{itemize}
\item One has $w(\Delta)=1$ on every cell of $\mathfrak{T}$ off the live pieces. Consequently
$\bar{\mathsf{w}}_v=3$ in three cases:
\begin{itemize}
\item every exposed unit of $\mathfrak{L}^c$;
\item every pointed $p$ of class~(c);
\item every pointed $p$ of class~(d) whose $\Delta_y$ has no neighbour on a live piece.
\end{itemize}
\item On a foreign live arrest $\mathfrak{a}'$ the frozen threshold data
$(\ell_{\mathfrak{a}'},w_{\mathfrak{a}'},\mathcal{Q}_{\mathfrak{a}'})$ stand, by the freezing clauses
of the statements on live arrests. There $w_{\mathfrak{a}'}=\mathsf{L}_0^{2e}$ with $e\le N_0(a')$.
We write $k_{a'}$ for the step of the arrest, and $k_{0,a'}$ for the step $k_0$ attached to the
attempt $\mathfrak{a}'$ in the statements on live arrests. When $w^{\vee}(\Delta_y)$ is attained on
the window of $X_{\mathfrak{a}'}$,
\begin{equation}\label{5.2:eq-exposed-compact-units-2}
\begin{aligned}
\bar{\mathsf{w}}_v
&\le 3\,\mathsf{L}_0^{2(e+1)}
\le 3\,\mathsf{L}_0^{2}\,L^{N_0(a')}
=3\,\mathsf{L}_0^{2}\,L\,\check R_{k_{0,a'}+1}\\
&\le 3\,\mathsf{L}_0^{2}\,L\,C_{\uparrow}\,\check R_{k_{a'}}
=:\bar{\mathsf{w}}^{\mathrm{fw}}(k_{a'}) .
\end{aligned}
\end{equation}
\end{itemize}
The right member of \eqref{5.2:eq-exposed-compact-units-2} depends on the step $k_{a'}$ of the arrest,
and no constant majorises it. It is paid unit by unit by the factor $\mathsf{f}_v^{1/2}$ of the
$\mathsf{t}$-piece's own bound in the carrier bound.

The units of class~(d) carry $\mathsf{t}$-pieces only, by the definition of class~(d). Their first
member $\mathbf{E}^{(j)}(X,U^0_k,z)$ in \cite[(1.60)]{B-CMP122b} is kept whole.

\item The \emph{read set} of $v$ is
\begin{equation}\label{5.2:eq-exposed-compact-units-3}
\mathfrak{r}_v:=\bigcup\operatorname{St}(\Delta_y).
\end{equation}
Here, for a cell $\Delta'$ of level $n':=\operatorname{lev}\Delta'$, $\operatorname{St}(\Delta')$ is the
wall-component of $\Delta'$ in the family
\begin{equation*}
\{\text{cells }\Delta:\ \operatorname{lev}\Delta\ge n'-1,\ \Delta\text{ meets }Q_5(\Delta')\},
\qquad
Q_5(\Delta'):=\{x:\ \mathrm{dist}_{n'}(x,\Delta')\le2\}.
\end{equation*}
The read set has the following properties.
\begin{itemize}
\item It consists of at most $(5L)^d=:c_S$ cells, of levels $n-1$, $n$, $n+1$.
\item It depends on $v$ through $\Delta_y$ alone.
\item It contains everything that is read for $v$ by the current bound, the regularity lemma, the
marginal and integral bounds and the Wilson line.
\item Every box $V\subset Q_5(\Delta_y)\cap\Omega_{n-1}$ meeting $\Delta_y$ lies in
$\mathfrak{r}_v$; in particular $\{\mathrm{dist}_n(\cdot,\Delta_y)\le1\}\subset\mathfrak{r}_v$.
\end{itemize}
The boxes read for $v$ are the following. All of them are boxes meeting $\Delta_y$ inside
$Q_5(\Delta_y)\cap\Omega_{n-1}$.
\begin{itemize}
\item The hull $\hat X_v$. It is a short single's box, a core's $Q_v$, or a Wilson piece's block of
side at most $2$; the block of type~$5$ has $(2M)^4$ sites.
\item The box $S\supset\hat X_v$ of the extended gauge-cube statement. Its sides are at most one
cube of $\pi_{n-1}$, and it holds $y$.
\item The gauge cube of the extended gauge-cube statement. It has index $n-1$ or $n$, and
its cells are at $\mathrm{dist}_n<1$ from $\Delta_y$.
\end{itemize}
A rough core has no hull of its own: it is read termwise.

\item Define
\begin{equation*}
\mathfrak{h}_v:=\{\Delta\in\pi_k:\ \mathrm{dist}_k(\Delta,\hat X_v)\le2\}.
\end{equation*}
This is a set of at most $7^d$ cells of $\pi_k$, which may or may not lie in $Z$. It satisfies
$[\mathfrak{r}_v]_k\subset\mathfrak{h}_v$. The label of $v$ is
\begin{align*}
\operatorname{lab}_1&:=K_v\cup\mathfrak{h}_v &&(v\in\mathfrak{L}^c,\ K_v=[X_v]_k),\\
\operatorname{lab}_1&:=K_p\cup\mathfrak{h}_v &&(p\in\mathfrak{L}^{Y_0}\text{ compact}),
\end{align*}
where
\begin{equation}\label{5.2:eq-exposed-compact-units-4}
K_p:=[\operatorname{St}_v\cup\mathsf{y}]_k\cup\operatorname{Hub}(\mathsf{y})\supseteq[X_v]_k .
\end{equation}
In \eqref{5.2:eq-exposed-compact-units-4} the letters are those of the summation of the pointed
pieces:
\begin{itemize}
\item $\operatorname{St}_v=\operatorname{St}(\Delta_y)$ is the exempt family of $v$;
\item $\mathsf{y}$ is the link-one walk of $p$;
\item $\operatorname{Hub}(\mathsf{y})$ is the set of corridor cubes bridging the never-cut blocks.
\end{itemize}

Since $\mathfrak{h}_v\subset\operatorname{lab}_1$, its cells inside $Z$ included, every component of
$Z$ holding a cell of $\mathfrak{h}_v$ meets $\operatorname{lab}_1$. Here ``meets'' means that it has a
common interior point with $\operatorname{lab}_1$. Such a component lies in
$Y(\operatorname{lab}_1)$, defined below. Near $\partial C$ those cells are $C$'s own. The lengths
$d_k$ and $d_{k,Z}$ are charged for all of $\mathfrak{h}_v$.

The label $\operatorname{lab}_1$ is connected, and
\begin{equation}\label{5.2:eq-exposed-compact-units-5}
\begin{aligned}
d_k(\operatorname{lab}_1)&\le d_k(K_v)+c_l\,c_{\mathfrak{h}} &&(v\in\mathfrak{L}^c),\\
d_k(\operatorname{lab}_1)&\le d_k(K_p)+c_l\,c_{\mathfrak{h}} &&(p\in\mathfrak{L}^{Y_0}\text{ compact}),
\end{aligned}
\qquad c_{\mathfrak{h}}:=7^d.
\end{equation}

\item The \emph{enlarged domain} of $v$ (resp.\ of $p$) is
\begin{equation*}
X^{\oplus}_v:=Y(\operatorname{lab}_1)
:=\operatorname{lab}_1\cup\{\text{components of }Z\text{ meeting }\operatorname{lab}_1\}
\in\mathbf{D}_k^Z .
\end{equation*}
It is connected. It contains a whole component of $Z$ in the following cases:
\begin{itemize}
\item for $\mathfrak{L}^c$ and for the compact $p$ of class~(c), it contains $C$;
\item for a compact $p$ of class~(d) whose piece does not vanish, it contains every component $C_p$
whose hub $\Lambda^{\sim}_{C_p}$ has a source cell held by the walk $\mathsf{y}$; equivalently, every
component $C_p$ whose $\Lambda^{\sim}_{C_p}$ is held by the domain $Y_0(p)$.
\end{itemize}
In each of these cases $X^{\oplus}_v$ satisfies the three clauses of \cite[(1.63)]{B-CMP122a}. For a
vanishing piece nothing is required. The \emph{auxiliary consumer} of $v$ is
\begin{equation*}
\mathfrak{E}^{\oplus}_v:=\mathbb{C}^{(N)}_k(X^{\oplus}_v)\in\mathcal{C}.
\end{equation*}
For every tuple $\mathfrak{p}$ of $v$ (resp.\ of $p$),
\begin{equation}\label{5.2:eq-exposed-compact-units-6}
X^{\oplus}_v\subset Y_4(\mathfrak{p}).
\end{equation}
Hence two hypotheses are met at $\mathfrak{E}:=\mathfrak{E}^{\oplus}_v$:
\begin{itemize}
\item the hypothesis ``$\mathfrak{E}\in\mathcal{C}$ with domain $X\subset Y_4$'' of the containment
statement for the consumer class;
\item the hypothesis ``stage $\le N$, $X\subset Y_4$'' of the claim on inherited clauses.
\end{itemize}

\item None of the exposed units of $\mathfrak{L}^c$ and of the pointed $p$ of $\mathfrak{L}^{Y_0}$ is
deep: $\hat X_v\not\subset\mathsf{C}_C$. Their size is the exterior size $\mathfrak{N}^{\mathrm{ex}}(v)$
of the carrier bound. Neither the disc in $\hat\lambda$ nor the proposition on deep pointed units is
used for them.
\end{enumerate}
\end{definition}

\begin{proof}
We verify the assertions made in the definition, in the order in which they occur.

\emph{No core lies in $\mathfrak{L}^c$.} A core is formed only at a far point. Indeed, by clause~(i)
of the lemma on the boundary type, the points near $\partial Z$ give constituents of type $\partial$.
Hence the hull $\square^{\sim\nu+1}$ of a core lies on one side of $\partial Z$, and no core has its
domain astride $\partial Z$.

\emph{The weight on $\mathfrak{T}$.} We show that $w(\Delta)=1$ on every cell of
$\mathfrak{T}=\Omega_k\cup Z^c$ off the live pieces. There are three parts.
\begin{itemize}
\item On $\Omega_k$: the threshold weight at level $k$ is $\mathsf{L}_0^0=1$, by clause~(ii) of
\cite[p.~190]{B-CMP122a} read at the level $k$.
\item On the collar $C\cap\Omega_k$: this collar has old level $k$.
\item On the untouched shells of the untreated components: by the inheritance theorem for stored
domains, the inherited clauses of $\mathfrak{I}_N$ are those of the source space
$\tilde U^c_k(Y_4,\tilde\alpha_0,\tilde\alpha_1)$ of the containment statement for the consumer
class, where $\tilde\alpha_0,\tilde\alpha_1$ are its radii. These clauses carry no power of
$\mathsf{L}_0$, cf.\ \cite[(2.34)--(2.39)]{B-CMP119}.
\end{itemize}
By clause~(i) of the geometry of the enlarged consumer, two kinds of units lie at $\partial C$, in
$C$'s collar or in $\Omega_k$ beside it: the exposed units of $\mathfrak{L}^c$, and the pointed $p$ of
class~(c). There every cell at $\mathrm{dist}_n\le1$ from $\Delta_y$ has weight $1$, so
$\bar{\mathsf{w}}_v=3$. The same holds for a pointed $p$ of class~(d) whose $\Delta_y$ has no
neighbour on a live piece.

\emph{The bound \eqref{5.2:eq-exposed-compact-units-2}.} Let $w^{\vee}(\Delta_y)$ be attained on the
window of a foreign live arrest $\mathfrak{a}'$.
\begin{enumerate}
\item By the freezing clauses of the statements on live arrests, the frozen weight there is
$w_{\mathfrak{a}'}=\mathsf{L}_0^{2e}$ with $e\le N_0(a')$. Hence
$\bar{\mathsf{w}}_v\le3\,\mathsf{L}_0^{2(e+1)}$.
\item Since $\mathsf{L}_0^2\le L/3$ and $e\le N_0(a')$,
\begin{equation*}
\mathsf{L}_0^{2(e+1)}=\mathsf{L}_0^{2}\,(\mathsf{L}_0^{2})^{e}
\le\mathsf{L}_0^{2}\,(L/3)^{N_0(a')}
\le\mathsf{L}_0^{2}\,L^{N_0(a')} .
\end{equation*}
\item The equality $L^{N_0(a')}=L\,\check R_{k_{0,a'}+1}$ is the definition of the number $N_0(a')$
of zones of the arrest.
\item The proof of the window-weight bound \eqref{5.2:eq-plain-fine-units-a} establishes, at the
arrest's step, the inequality
\begin{equation*}
\check R_{k_{0,a'}+1}\le\mathsf{R}^{\uparrow}(\check{\mathsf{T}}_{k_{a'}})
\le C_{\uparrow}\check R_{k_{a'}} .
\end{equation*}
This gives the last inequality of \eqref{5.2:eq-exposed-compact-units-2}.
\end{enumerate}
This yields \eqref{5.2:eq-exposed-compact-units-2}.

\emph{The read set.} The following properties of $\mathfrak{r}_v=\bigcup\operatorname{St}(\Delta_y)$
are those of the star lemma of the summation of the pointed pieces.
\begin{itemize}
\item The members of $\operatorname{St}(\Delta_y)$ have levels $n-1$, $n$, $n+1$. The lower bound
holds by definition. The upper bound holds by the admissibility of the label,
$\mathrm{dist}_{i+1}(\Omega_i^c,\Omega_{i+1})\ge1$ for every $i$.
\item Their number is at most $(5L)^d=c_S$.
\item Every box $V\subset Q_5(\Delta_y)\cap\Omega_{n-1}$ meeting $\Delta_y$ lies in
$\bigcup\operatorname{St}(\Delta_y)$; in particular $\{\mathrm{dist}_n(\cdot,\Delta_y)\le1\}$ does.
\item If $y_v\in\Delta'$ then $[\operatorname{St}(\Delta')]_k\subset\mathfrak{h}_v$. This is the
inclusion $[\mathfrak{r}_v]_k\subset\mathfrak{h}_v$.
\end{itemize}
The set $\operatorname{St}(\Delta_y)$ is defined from $\Delta_y$ only, so $\mathfrak{r}_v$ depends on
$v$ through $\Delta_y$ alone.

Three boxes are read for $v$: the hull $\hat X_v$, the box $S\supset\hat X_v$ of the extended
gauge-cube statement, and its gauge cube. Each of them meets $\Delta_y$: the hull and $S$
hold $y$, and the gauge cube contains $S\ni y$. The hull and $S$ lie within one level-$n$ unit of
$\Delta_y$, for the following reasons.
\begin{itemize}
\item $S$ and a short single's box have sides at most one cube of $\pi_{n-1}$.
\item A Wilson piece's block holds $y$ and lies within one level-$n$ unit of $\Delta_y$.
\item A core's $Q_v$ lies within one level-$n$ unit of $y$.
\end{itemize}
So these boxes lie in $Q_5(\Delta_y)$. By admissibility at level $n-1$, the cube
$\{\mathrm{dist}_n(\cdot,\Delta_y)\le1\}$ lies in $\Omega_{n-1}$, hence so do they. The cells of the
gauge cube have level $n-1$ or $n$ and are at $\mathrm{dist}_n<1$ from $\Delta_y$, so the gauge
cube lies in $\Omega_{n-1}\cap Q_5(\Delta_y)$. By the third property above, the three boxes therefore
lie in $\mathfrak{r}_v$.

\emph{Components of $Z$ at $\mathfrak{h}_v$.} The cells of $\mathfrak{h}_v$ lying in $Z$ belong to
$\operatorname{lab}_1$. A component of $Z$ holding such a cell therefore has a common interior point
with $\operatorname{lab}_1$, that is, it meets $\operatorname{lab}_1$. It is thus adjoined in
$Y(\operatorname{lab}_1)$.

\emph{Connectedness of $\operatorname{lab}_1$ and \eqref{5.2:eq-exposed-compact-units-5}.}
\begin{itemize}
\item $[X_v]_k\subset K_v$ by definition, and $[X_v]_k\subset K_p$ by
\eqref{5.2:eq-exposed-compact-units-4}.
\item $[X_v]_k\subset\mathfrak{h}_v$, since $X_v\subset\hat X_v$.
\item Hence $[X_v]_k\subset K_v\cap\mathfrak{h}_v$, resp.\ $[X_v]_k\subset K_p\cap\mathfrak{h}_v$.
Since $K_v$ (resp.\ $K_p$) and $\mathfrak{h}_v$ are connected and meet, their union
$\operatorname{lab}_1$ is connected.
\item Clause~(a) of the multiscale resummation lemma states that the tree length grows by at most
$c_l$ per cell added adjacent to the set. The set $\mathfrak{h}_v$ is connected and meets $K_v$
(resp.\ $K_p$). So its cells outside $K_v$ (resp.\ $K_p$), at most $c_{\mathfrak{h}}=7^d$ of them, can
be added one at a time, each adjacent to the set already built. This gives
\eqref{5.2:eq-exposed-compact-units-5}.
\end{itemize}

\emph{The enlarged domain.} $X^{\oplus}_v$ is connected, because $\operatorname{lab}_1$ is connected
and each adjoined component of $Z$ meets it. It contains a whole component of $Z$ in the cases
stated in the definition:
\begin{itemize}
\item For $\mathfrak{L}^c$ and for the compact $p$ of class~(c): $K_v\cap C\neq\emptyset$ and
$K_v\subset\operatorname{lab}_1$. So $C$ meets $\operatorname{lab}_1$, and $C\subset X^{\oplus}_v$.
\item For the compact $p$ of class~(d): by the lemma on the vanishing of pieces outside the
source-anchored walks, the walk $\mathsf{y}$ of a non-vanishing piece holds a source cell. We place
this cell in $C_p$ and in $\operatorname{lab}_1$ as follows.
\begin{itemize}
\item It shares a wall with the hub $\Lambda^{\sim}_{C_p}\subset C_p$.
\item It lies in $Z$, since $\mathrm{dist}_k(\Lambda^{\sim},Z^c)\ge9\check R_k$.
\item Being wall-adjacent within $Z$ to $\Lambda^{\sim}_{C_p}$, it lies in $C_p$.
\item It is a cube of $[\mathsf{y}]_k\subset K_p\subset\operatorname{lab}_1$.
\end{itemize}
Hence $C_p$ meets $\operatorname{lab}_1$, and $C_p\subset X^{\oplus}_v$. Equivalently, $Y_0(p)$ holds
$\Lambda^{\sim}_{C_p}$. Indeed, the summation in \cite[p.~372]{B-CMP122b} runs over domains
$Y_0\in\mathbf{D}_k$ containing $\square_0$ and at least one of the components of $\Lambda^{\sim}$,
and $\Lambda^{\sim}\subset Z$.
\end{itemize}
A component of $Z$ satisfies the three clauses of \cite[(1.63)]{B-CMP122a}, and a superset keeps
them, exactly as in the consumer lemma on $C$. Hence $X^{\oplus}_v$ satisfies them.

\emph{The auxiliary consumer.} The consumer class $\mathcal{C}$ contains the chain terms
$\mathbb{C}^{(i)}_k$ with $i\le N$. By the inheritance theorem for stored domains, every index
$X\in\mathbf{D}_k^Z$ satisfying the clauses of \cite[(1.63)]{B-CMP122a} carries the stored domain
$\mathfrak{D}^{\flat}_N(X)$. The class is indexed by $X$, so an index astride $\partial Z$ is as
admissible as $C$. Hence $\mathfrak{E}^{\oplus}_v=\mathbb{C}^{(N)}_k(X^{\oplus}_v)\in\mathcal{C}$.

\emph{The inclusion \eqref{5.2:eq-exposed-compact-units-6}.} Along the four links the labels grow,
so $\operatorname{lab}_1\subset\operatorname{lab}_4$. The map $Y(\cdot)$ is monotone, and
$Y_4(\mathfrak{p})=Y(\operatorname{lab}_4)$ by clause~(a) of the multiscale resummation lemma. Hence
\begin{equation*}
X^{\oplus}_v=Y(\operatorname{lab}_1)\subset Y(\operatorname{lab}_4)=Y_4(\mathfrak{p}).
\end{equation*}
This is the domain hypothesis of the containment statement for the consumer class and of the claim
on inherited clauses, at $\mathfrak{E}:=\mathfrak{E}^{\oplus}_v$.

\emph{No unit is deep.} In every case $X_v\not\subset Z$:
\begin{itemize}
\item for $\mathfrak{L}^c$ and class~(c), $K_v\not\subset Z$;
\item for class~(d), $X\subset Z^c$.
\end{itemize}
The core $\mathsf{C}_C$ lies in $C\subset Z$. The hull $\hat X_v$ contains $X_v$ by the definition of
the hull; for a core the hull is $\square^{\sim\nu+1}$. Hence $\hat X_v$ is not contained in
$\mathsf{C}_C$.
\end{proof}

\begin{lemma}[Geometry of exposed units: collar, zones and admissible cube]\label{5.2:lem-exposed-geometry}
Let $v$ be an exposed unit of $\mathfrak{L}^{c}$, or the unit underlying a compact pointed
piece $p$ of $\mathfrak{L}^{Y_0}$, as in Definition~\ref{5.2:def-exposed-compact-units}.
Write $y:=y_v$, let $n:=\pi(y)$ be its zone index in the current global label, $\Delta_y$ its
cell (of level $n$), $\hat X_v$ its hull, and let $\mathfrak{r}_v$, $\mathfrak{h}_v$,
$\operatorname{lab}_1$ and $X^{\oplus}_v=Y(\operatorname{lab}_1)$ be as defined there.
\begin{enumerate}
\item[(i)] \emph{Units meeting $C$.} Let $v$ be an exposed unit of $\mathfrak{L}^{c}$ or the unit
of a compact pointed piece of class (c). If $X_v$ meets $C$, then $n=\pi(y)=k$,
$\mathfrak{r}_v\subset\Omega_k\cup(C\cap\Omega_k)$, and
\begin{equation*}
\operatorname{dist}_k(y,\Omega_k^c)\ \ge\ 2\check{R}_k-5\ \ge\ 71 .
\end{equation*}
\item[(ii)] \emph{Units of class (d).} Let $v$ be of class (d), so that $y\in X_v\subset Z^c$.
Then $y$ lies in a zone $n\le k$ of the global label: on $\Omega_k$, or on a shell of an
untreated component, or in or beside the window of a foreign live arrest inside such a
component. No further property of $n$ is asserted. The cells of $\mathfrak{r}_v$ have levels
$n-1$, $n$, $n+1$ only; and $\mathfrak{h}_v$ is contained in the union of $Z^c$ with the
treated components lying within two $\pi_k$-cubes of $\hat X_v$; each such component lies
inside $X^{\oplus}_v$.
\item[(iii)] \emph{Admissible cube.} Let $S\subset\mathfrak{r}_v$ be a box holding $y$ whose
sides are at most one cube of $\pi_{n-1}$. Then $S$ lies in an admissible cube (in the sense of
the admissible-cube statement established on the component $C$, read at the zone index $n$ of
$y$) of index $n-1$ or $n$, whose cells of that index are at $\operatorname{dist}_n<1$ from
$\Delta_y$; this admissible cube lies inside
$\mathfrak{r}_v\subset\mathfrak{h}_v\subset X^{\oplus}_v$. For a far core of class (d), the box
$Q_v=Q^{(p_v)}\subset\mathfrak{r}_v$ is itself such a box $S$.
\end{enumerate}
\end{lemma}

\begin{proof}
\emph{(i).} The zones of $C$ of level $<k$ lie in
$Z''_k\subset\Lambda_C:=(\Omega^{\sim 4}_{k_0+1})^c\cap C$, where $\Omega^{\sim m}$ denotes
$\Omega$ enlarged by $m$ layers of $M$-cubes: in Ba{\l}aban's construction
$Z''_k=(\Omega^{\sim 5}_{k_0+1})^c\cap Z$, and $\Lambda=(\Omega^{\sim 4}_{k_0+1})^c\cap Z$ is
obtained by adding one layer of $M$-cubes to $Z''_k$ \cite{B-CMP122a}. Hence $C\setminus\Lambda_C$
contains $C\cap\Omega_k$ less four layers of $M$-cubes of level $k_0+1$. On the other hand, the
boundaries of $\Omega_j$ and $\Lambda_j$ are at distance at least $2MR_j$
\cite[p.~256]{B-CMP119}. That statement concerns regions $\Omega_k$, $\Lambda_k$ made by the
operation $\mathbf{T}$, and at $C$ these regions are of that kind by condition (ii) of the
selection of the treated components \cite{B-CMP122a}. Therefore
\begin{equation*}
\operatorname{dist}_k(\partial C,\Omega_k^c)\ \ge\ 2\check{R}_k
\end{equation*}
on both sides of $\partial C$. So the collar $C\cap\Omega_k$ of $C$ has thickness at least
$2\check{R}_k$ and its old level is $k$.

Next, the hull $\hat X_v$ has diameter at most two level-$n$ units, according to the kind of
unit:
\begin{itemize}
\item for a short single it is at most one cube of $\pi_{n-1}$;
\item for a Wilson piece it is one level-$j$ unit, or a block of side at most $2$;
\item for a smooth core it is at most $L^{a_v-1}$ cubes of $\pi_{p_v}$, by the definition of
cores and their smoothing boxes.
\end{itemize}
For a rough core the box $\square^{\sim\nu+1}$ would measure $(2\nu+3)L^{-(\nu-1)}$ level-$n$
units a side, that is, five units at $\nu=1$, which is $Q_5(\Delta_y)$. A rough core is
therefore not assigned a hull but is read term by term, each constituent being a short single
(compact) or a long one (treated separately). This does not affect the present lemma. No core
occurs among the units of case (i), and for a single term or boundary constituent whose domain
meets $Z$ one has $j<k$ \cite{B-CMP119}; at $n=k$ this gives $\nu=1+n-j\ge 2$, and the box
$\square^{\sim\nu+1}$ is then less than one level-$n$ unit. In class (d) one has $j\le k$, and
at $j=k$:
\begin{itemize}
\item the terms of type $\mathbf{E}$ and $\mathbf{R}$ and the boundary constituents decay at
the rate $2\kappa$;
\item the terms of type $\mathbf{D}$ are absent;
\item every single is long or a tail.
\end{itemize}

Finally, distinct components of $Z_k=\Lambda_k^c$ are unions of non-touching
$M\check{R}_k$-cubes of one partition \cite[p.~256]{B-CMP119}, \cite{B-CMP122a}. They are at
least $\check{R}_k$ units apart \cite{B-CMP119}, and $\check{R}_k\ge 38$ by the standing lower
bound on $\check{R}$.

Now let $X_v$ meet $C$. For an exposed unit of $\mathfrak{L}^{c}$, and for the unit of a compact
pointed piece of class (c), one also has $K_v=[X_v]_k\not\subset Z$. These two facts, together
with the diameter bound on the hull and the separation of the components, force $y$ to lie
within two units of $\partial C$. By the display above, and since $2\check{R}_k>2$, every point
within two units of $\partial C$, on either side of it, lies in $\Omega_k$. Hence $n=k$, and
$\operatorname{dist}_k(y,\Omega_k^c)\ge 2\check{R}_k-5\ge 71$. At $n=k$ no member of
$\operatorname{St}(\Delta_y)$ has level $k+1$, since no cell of the label has level above $k$;
so $\mathfrak{r}_v\subset\{\operatorname{dist}_k(\cdot,\Delta_y)\le 2\}$ by clause (e5) of the
lemma on the star of a cell. Every point of $\mathfrak{r}_v$ is therefore within $3<71$
level-$k$ units of $y$, hence lies in $\Omega_k$, which gives
$\mathfrak{r}_v\subset\Omega_k\cup(C\cap\Omega_k)$.

The restriction of (i) to compact pieces is needed: the domain of a long single or tail may join
a point $y\in Z''_k$ to $\Lambda_k$, and for such a point the conclusion $n=k$ fails.

\emph{(ii).} Here $X_v\subset Z^c$, and $Z$ is the union of a class of components of
$\Lambda_k^c$, namely the treated ones \cite{B-CMP122a}. So $y$ lies in a zone $n\le k$ of the
global label, in one of three places:
\begin{itemize}
\item on $\Omega_k$;
\item on a shell of an untreated component, made by the operation $\mathbf{T}$, each zone of
which consists of at least ten layers of its own $M\check{R}$-cubes \cite{B-CMP122a};
\item in or beside the window of a foreign live arrest inside such a component, whose zones
were re-made by the stages of that arrest and are frozen since, by the freezing property of
live arrests.
\end{itemize}

The levels of the cells of $\mathfrak{r}_v$ are bounded as follows. The lower end $n-1$ holds
by definition of $\operatorname{St}(\Delta_y)$. The upper end $n+1$ is clause (e1) of the lemma
on the star of a cell, which rests on the zone-separation property of every label,
$\operatorname{dist}_{j+1}(\Omega_j^c,\Omega_{j+1})\ge 1$, taken at $j=n+1$.

The claim on $\mathfrak{h}_v$ is unwound from the definitions. First,
\begin{equation*}
\mathfrak{h}_v=\{\Delta\in\pi_k:\operatorname{dist}_k(\Delta,\hat X_v)\le 2\},
\end{equation*}
so a cell of $\mathfrak{h}_v$ contained in $Z$ lies in a treated component within two
$\pi_k$-cubes of $\hat X_v$. Second, such a component meets
$\operatorname{lab}_1\supseteq\mathfrak{h}_v$. Since
\begin{equation*}
X^{\oplus}_v=Y(\operatorname{lab}_1)
=\operatorname{lab}_1\cup\{\text{components of }Z\text{ meeting }\operatorname{lab}_1\},
\end{equation*}
the component lies inside $X^{\oplus}_v$.

\emph{(iii).} The admissible-cube statement on $C$ is applied verbatim at the zone index $n$ of
$y$. Its proof uses three facts, each of which holds here.
\begin{itemize}
\item The zone-separation property, which holds in every label and hence globally.
\item The $\check{R}_n$-grid of an $\Omega_n^c$ made by the operation $\mathbf{T}$, which yields
the first alternative of that statement (two cubes of its own partition at gap $1$); this is
true of the nested zones of $C$ and of those of an untreated component alike.
\item The box property of a zone re-made inside the component $C'$ of an attempt, where
$\Omega_n^c\cap C'=Z''_n$ is a box. In Ba{\l}aban's construction all regions connected with the
last preparatory steps are rectangular parallelepipeds, and so are the sets $Z''_j$
\cite{B-CMP122a}. This gives the second alternative, and it holds for the running attempt and
for a live arrest, on $C$ or foreign, alike. For a live arrest the structure is frozen
by the freezing property of live arrests.
\end{itemize}
Nothing specific to $C$ is used. At the boundary of $C$, in case (i), the first alternative
cannot even occur within distance $71>1$ of $y$, so no obstruction arises there.

The admissible cube is a box meeting $\Delta_y$ inside $Q_5(\Delta_y)\cap\Omega_{n-1}$. By
Definition~\ref{5.2:def-exposed-compact-units}, with clause (e5) of the lemma on the star of a
cell, every such box lies in $\mathfrak{r}_v$. Moreover $[\mathfrak{r}_v]_k\subset\mathfrak{h}_v$
and $\mathfrak{h}_v\subset\operatorname{lab}_1\subset X^{\oplus}_v$. This gives the chain
$\mathfrak{r}_v\subset\mathfrak{h}_v\subset X^{\oplus}_v$.

For a far core of class (d), consider $Q_v=Q^{(p_v)}$. By the definition of cores and their
smoothing boxes, its sides are at most $L^{a_v-1}$ cubes of $\pi_{p_v}$, which is at most one
cube of $\pi_{n-1}$. It lies within one level-$n$ unit of $y$, hence in
$\{\operatorname{dist}_n(\cdot,\Delta_y)\le 1\}$. By the zone-separation property at $j=n-1$,
this cube lies in $\Omega_{n-1}$, so $Q_v$ lies in $\Omega_{n-1}$; and by clause (e5) of the
lemma on the star of a cell the same cube lies in $\mathfrak{r}_v$, so
$Q_v\subset\mathfrak{r}_v$. Hence $Q_v$ is a box $S$ as above.
\end{proof}

\begin{lemma}[Bounds for exposed units at every zone]\label{5.2:lem-exposed-zone-bound}
Fix $K$, a level $k\le K$ and a label.
Let $v$ be an exposed unit, or the unit underlying a compact pointed piece $p$, in the sense of
Definition~\ref{5.2:def-exposed-compact-units}. Let $y:=y_v$, let $n:=\pi(y)$ be the zone index
of $y$ in the current global label, and let $\Delta_y$ be the cell of $y$. Let
$\mathfrak{r}_v$, $\mathfrak{h}_v$, $\operatorname{lab}_1$, $X^{\oplus}_v=Y(\operatorname{lab}_1)$,
$\mathfrak{E}^{\oplus}_v=\mathbb{C}^{(N)}_k(X^{\oplus}_v)$ and $\bar{\mathsf{w}}_v=3w^{\vee}(\Delta_y)$
be as in that definition. For a tuple $\mathfrak{p}$ of $v$, or of $p$, let $\operatorname{lab}_4$
and $Y_4(\mathfrak{p})=Y(\operatorname{lab}_4)$ be its label and its domain, as there; let
$\mathfrak{T}:=\Omega_k\cup Z^c$. Fix such a tuple $\mathfrak{p}$ and write
$Y_4:=Y_4(\mathfrak{p})$; the hypotheses below are assumed for every such tuple. The input space is
$\tilde U^c_k(Y_4,\tilde\alpha_0,\tilde\alpha_1)$, with radii $\tilde\alpha_0,\tilde\alpha_1$; for a
point $y'$ and $a\in\{1,3,5,6,7\}$, $r_a(y')$ denotes the $a$-th threshold of the clause of the
input space at $y'$. Assume the following.
\begin{enumerate}
\item[(A)] \emph{The carrier containment statement.} For every $\mathfrak{E}\in\mathcal{C}$ whose
domain $X$ lies in $Y_4$, every $\zeta\in\mathfrak{S}_{\mathrm{src}}=\tilde U^c_k(Y_4,\tilde\alpha_0,
\tilde\alpha_1)$, every real argument $\Phi$ of the site map, every $|\sigma|\le e^{\kappa_1}$ and
every $\mathfrak{b}\in\{\mathfrak{b}_2,\mathfrak{b}_3,\mathfrak{b}_4\}$, the presented pair, rotated
by the gauge transformation $h$ of the presentation, satisfies
\begin{equation*}
\bigl((\hat{\mathcal{U}}^{\mathfrak{b}})^h,\operatorname{Ad}_h\hat{\mathcal{J}}^{\mathfrak{b}}\bigr)
\lceil X\in\mathfrak{D}_{\mathfrak{E}}.
\end{equation*}
\item[(A$'$)] \emph{The containment statement at the localisation site.} The containment of (A)
holds with $\mathsf{s}$ free on the cells of $X$ as well; and, for the brackets
$(\mathfrak{b}_0,\mathfrak{b}_1)$ and $(\mathfrak{b}_1,\mathfrak{b}_2)$ of the units whose
localisation domain is disjoint from $Z$, the same containment holds.
\item[(B)] \emph{The inheritance theorem.} For every chain term $\mathfrak{E}$ of stage at most $N$
in $\mathcal{C}$ whose domain $X\in\mathbf{D}^Z_k$ lies in $Y_4$ and satisfies the three
admissibility clauses of a chain-term index, the analyticity domain $\mathfrak{D}_{\mathfrak{E}}$ of
$\mathfrak{E}$ is the stored domain $\mathfrak{D}^{\flat}_N(X)$ carried by the index $X$, the
intersection of $\mathfrak{N}^{(N)}$,
$\mathfrak{L}_N(X\cap Z^{\mathrm{on}})$ and $\mathfrak{I}_N(X\cap\mathfrak{T})$; here
$\mathfrak{N}^{(N)}$, $\mathfrak{L}_N(\cdot)$ and $\mathfrak{I}_N(\cdot)$ are the sets fixed with
this theorem, and $Z^{\mathrm{on}}$ is the region, fixed with this theorem, on which
$\mathfrak{L}_N$ is supported. Off the
live pieces, the entries of the inherited clause exceed those of the input space by the stage
increments only.
\item[(C)] \emph{The increment corollary.} For $X$ as in (B), at every point $y'$ of
$X\cap\mathfrak{T}$ off the live pieces, the stage increments are at most
$\frac14\mathrm{room}^{\flat}(y')<r_a(y')$. The $U$-part and the cube gauges are kept, and
$\mathbf{R}$ changes nothing on $Z^c$.
\item[(D)] \emph{The inherited-clause claim.} For $X$ as in (B), on $X\cap\mathfrak{T}$ off every
live piece, the clause of $\mathfrak{I}_N$ is taken at the shell index of $y'$, which equals
$\pi(y')\le k$. Its
fifth and eighth entries and its cube gauges are those of the input variable $\zeta$.
\item[(E)] \emph{The freezing properties of live arrests.} On a live piece $W_{\mathfrak{a}'}$ of a
foreign live arrest $\mathfrak{a}'$, and on the set $\mathfrak{T}_{k_{a'}}$ of its component beside
it (the region of $\mathfrak{T}$-locations of that component at the step $k_{a'}$ of the arrest),
the input space is read with the frozen structure of the arrest, consisting of a frozen clause
index, the frozen weight $w_{\mathfrak{a}'}$ and the frozen family of cubes
$\mathcal{Q}_{\mathfrak{a}'}$, the factor $c$ of the frozen structure being identically $1$. Its
members are restrictions of pairs of
$\mathfrak{G}(\mathfrak{B})$, the configuration pairs over the determining set $\mathfrak{B}$. At a
cell $y'$ of these sets the frozen clause, as the theorem on arrested attempts gives it on live
pieces, has frozen clause index equal to $\pi(y')$, thresholds
$3c_\alpha w_{\mathfrak{a}'}\mathbf{r}_a$ and the cube gauges of $\mathcal{Q}_{\mathfrak{a}'}$.
\item[(F)] \emph{The consumer domain on treated components.} This is the domain used in the carrier
bound on $C$. On a cell of $C$, or of a treated component, it has clause index $\pi(y')$ and
thresholds
\begin{equation*}
\varphi_{\flat}\mathbf{r}_a\le\theta_a\le 3c_\alpha w(\Delta_{y'})\mathbf{r}_a(\pi(y')).
\end{equation*}
It is read through the frozen structure on live pieces, and it carries the cube clause.
\item[(G)] \emph{The regularity ceiling}
\begin{equation*}
\sup_{x\ge\gamma^{-2}}360\,dL^2M\,c_{\mathrm{reg}}\,\mathsf{R}^{\uparrow}(\log x)\,\alpha_*(x)\le1 .
\end{equation*}
\end{enumerate}
Then the following hold.
\begin{enumerate}
\item[(i)] For every $m\in\mathfrak{P}$, every $y'\in\mathfrak{r}_v$ and every
$a\in\{1,3,5,6,7\}$,
\begin{equation}\label{5.2:eq-exposed-zone-bound-2}
Q_a(\mathrm{chain}(m))(y')<\theta_a(y')\le 3c_\alpha\,w(\Delta_{y'})\,\mathbf{r}_a(\pi(y')).
\end{equation}
Here $w(\Delta_{y'})=1$ unless $\Delta_{y'}$ is a window cell of $C$ or of a foreign live arrest.
Moreover, let $\square\subset\mathfrak{r}_v$ be an admissible cube of
Lemma~\ref{5.2:lem-exposed-geometry}, of index $\mathsf{m}\in\{n-1,n\}$, and let $\ell_{\mathsf{m}}$
be the length unit of level $\mathsf{m}$. Then the gauge-field part of the clause provides on
$\square$ an $SU(N)$-valued gauge transformation $u$ whose field $A$ satisfies
\begin{equation*}
\ell_{\mathsf{m}}|A|,\ \ell_{\mathsf{m}}^2|\nabla A|<3w\cdot2BM\alpha_{0,\mathsf{m}},
\end{equation*}
where $w$ is the weight entering the clause on $\square$ and $B$ is the constant of the cube clause.
\item[(ii)] Let $S\subset\mathfrak{r}_v$ be a box holding $y$, with sides at most one cube of
$\pi_{n-1}$, as in Lemma~\ref{5.2:lem-exposed-geometry}. Put $\ell:=\ell_{n-1}$, the length unit of
level $n-1$, and
\begin{equation}\label{5.2:eq-exposed-zone-bound-3}
\mathfrak{a}'_v:=c_{\mathrm{reg}}\,\bar{\mathsf{w}}_v\,\alpha_{*,n}.
\end{equation}
Then for every $m\in\mathfrak{P}$, in the gauge $u$ of (i) restricted to $S$,
$\mathrm{chain}(m)\lceil S=(e^{i\eta\mathcal{A}^0},\mathbf{J}^0)$, with $\eta$ the lattice spacing
of Ba{\l}aban's construction, and
\begin{equation*}
\ell|\mathcal{A}^0|\le\mathfrak{a}'_v,\qquad \ell^2|\nabla\mathcal{A}^0|\le\mathfrak{a}'_v,
\qquad \ell^3|\mathbf{J}^0|\le\mathfrak{a}'_v .
\end{equation*}
Moreover $60dLM\mathfrak{a}'_v\le1$.
\item[(iii)] With the bounds of (ii) as input, the conclusions of the carrier bound on $C$ for
singles (boxes of side $D\le D_X$), for Wilson pieces (pointwise in the first entry) and, for a
smooth far core of a unit whose localisation domain is disjoint from $Z$, at the extension
$\mathrm{ext}_{p_v}$ on $Q_v\subset\mathfrak{r}_v$, hold verbatim, read at the zone index $n$ and
the cell $\Delta_y$ in place of the level $k$ and its cell.
\end{enumerate}
\end{lemma}

\begin{proof}
\emph{The auxiliary term satisfies the hypothesis of \textup{(A)}.} By
Definition~\ref{5.2:def-exposed-compact-units}, $\mathfrak{E}^{\oplus}_v\in\mathcal{C}$, and its
domain is $X^{\oplus}_v=Y(\operatorname{lab}_1)$. Let $\mathfrak{p}$ be any tuple of $v$, or of the
piece $p$. Then $\operatorname{lab}_1\subset\operatorname{lab}_4$, the map $Y(\cdot)$ is monotone,
and $Y_4(\mathfrak{p})=Y(\operatorname{lab}_4)$. Hence $X^{\oplus}_v\subset Y_4(\mathfrak{p})$ for
every tuple $\mathfrak{p}$. Moreover $X^{\oplus}_v$ contains a whole component of $Z$, and so
satisfies the three admissibility clauses of a chain-term index
(Definition~\ref{5.2:def-exposed-compact-units}). So the hypothesis ``$\mathfrak{E}\in\mathcal{C}$
with domain $X\subset Y_4$'' of (A) holds at $\mathfrak{E}:=\mathfrak{E}^{\oplus}_v$, with
$\zeta\in\mathfrak{S}_{\mathrm{src}}=\tilde U^c_k(Y_4,\tilde\alpha_0,\tilde\alpha_1)$. The term
$\mathfrak{E}^{\oplus}_v=\mathbb{C}^{(N)}_k(X^{\oplus}_v)$ has stage $N$ and domain contained in
$Y_4$, so it lies in the scope of (B), (C) and (D) as well.

\emph{Applying \textup{(A)}.} For $\mathfrak{b}\in\{\mathfrak{b}_2,\mathfrak{b}_3,\mathfrak{b}_4\}$,
(A), with $\mathsf{s}$ free on the cells of $X^{\oplus}_v$ by (A$'$), gives that the presented pair
restricted to $X^{\oplus}_v$ lies in $\mathfrak{D}_{\mathfrak{E}^{\oplus}_v}$. For the brackets
$(\mathfrak{b}_0,\mathfrak{b}_1)$ and $(\mathfrak{b}_1,\mathfrak{b}_2)$ of the units whose
localisation domain is disjoint from $Z$, the same
conclusion follows from (A$'$). By (B),
\begin{equation*}
\mathfrak{D}_{\mathfrak{E}^{\oplus}_v}=\mathfrak{D}^{\flat}_N(X^{\oplus}_v)
=\mathfrak{N}^{(N)}\cap\mathfrak{L}_N(X^{\oplus}_v\cap Z^{\mathrm{on}})
\cap\mathfrak{I}_N(X^{\oplus}_v\cap\mathfrak{T}).
\end{equation*}

\emph{Where the points $y'$ lie.} By Definition~\ref{5.2:def-exposed-compact-units} and
Lemma~\ref{5.2:lem-exposed-geometry}, $\mathfrak{r}_v\subset\mathfrak{h}_v\subset X^{\oplus}_v$.
So every $y'\in\mathfrak{r}_v$ is a point of the domain of $\mathfrak{E}^{\oplus}_v$. Since
$\mathrm{chain}(m)$ is the carrier pair of (A) at the parameter $m$, presented in the gauge $h$, the
entries $Q_a(\mathrm{chain}(m))(y')$ are below the thresholds $\theta_a(y')$ of
$\mathfrak{D}^{\flat}_N(X^{\oplus}_v)$ at $y'$. This is the first inequality
of~\eqref{5.2:eq-exposed-zone-bound-2}. It remains to identify the clause of this domain at $y'$,
in three cases.
By Lemma~\ref{5.2:lem-exposed-geometry}, every cell of $\mathfrak{r}_v$ lies either in a treated
component contained in $X^{\oplus}_v$, or in $Z^c\subset\mathfrak{T}$, off or on a foreign live
piece; these are the three cases below.

($\alpha$) \emph{$y'$ lies on a cell of $C$, or of a treated component contained in
$X^{\oplus}_v$.} There the domain is the consumer domain of (F). Its clause index is $\pi(y')$ and
its thresholds obey the display of (F). On live
pieces it is read through the frozen structure, and it carries the cube clause.

($\beta$) \emph{$y'$ lies in $\mathfrak{r}_v\cap\mathfrak{T}$ off every live piece.} This covers
$\Omega_k$ and the shells of the untreated components, at any shell index at most $k$. By (D), the
domain there is the inherited clause of $\mathfrak{I}_N$ at the shell index of $y'$, which is
$\pi(y')$.

The fifth and eighth entries, and the cube gauges, are those of $\zeta$. Now
$\zeta\in\tilde U^c_k(Y_4,\tilde\alpha_0,\tilde\alpha_1)$ satisfies Ba{\l}aban's level-$\pi(y')$
clauses, of the type of \cite[(2.34)--(2.39)]{B-CMP119}, at every cell of zone $\pi(y')$ of
$Y_4\supset\mathfrak{r}_v$. These clauses include a gauge transformation $u$ defined on
$\square\cap X$ for every cube $\square$ of the cover entering these clauses
\cite[(2.38)]{B-CMP119}. This applies in
particular to the admissible cube, which lies in $\mathfrak{r}_v\subset Y_4$.

The clause of $\zeta$ at $y'$ has threshold $r_a(y')$. The other entries exceed those of $\zeta$
only by the stage increments of (B). By (C) these are at most
$\frac14\mathrm{room}^{\flat}(y')<r_a(y')$; the $U$-part and the cube gauges are kept, and
$\mathbf{R}$ changes nothing on $Z^c$. Adding the increment, which is less than $r_a(y')$, and
using the threshold convention $r_a\le c_\alpha w\mathbf{r}_a$, here with $w=1$, since every cell
of $\mathfrak{T}$ off the live pieces has weight $1$
(Definition~\ref{5.2:def-exposed-compact-units}), every threshold of
the inherited clause is therefore at most
\begin{equation*}
2r_a(y')\le 2c_\alpha\mathbf{r}_a(\pi(y'))<3c_\alpha\mathbf{r}_a(\pi(y')).
\end{equation*}

($\gamma$) \emph{$y'$ lies on a cell of a foreign live piece
$W_{\mathfrak{a}'}\subset Z^c\cap X^{\oplus}_v$, or of the set $\mathfrak{T}_{k_{a'}}$ of its
component beside it.} By (E), the input space there is itself read with the frozen structure of the
arrest, with frozen weight $w_{\mathfrak{a}'}$, frozen family $\mathcal{Q}_{\mathfrak{a}'}$ and
the factor $c$ identically $1$, and its members are restrictions of pairs of
$\mathfrak{G}(\mathfrak{B})$. The frozen
clause index at $y'$ is $\pi(y')$, the threshold is $3c_\alpha w_{\mathfrak{a}'}\mathbf{r}_a$, and
the cube gauges are those of the frozen family. $\mathfrak{I}_N$ inherits exactly this clause.

\emph{The weight.} In each case the threshold is at most
$3c_\alpha w(\Delta_{y'})\mathbf{r}_a(\pi(y'))$. By Definition~\ref{5.2:def-exposed-compact-units},
$w(\Delta_{y'})=1$ unless $\Delta_{y'}$ is a window cell of $C$ or of a foreign live arrest. This
proves~\eqref{5.2:eq-exposed-zone-bound-2}.

\emph{The gauge clause.} The three cases also supply, on every admissible cube
$\square\subset\mathfrak{r}_v$ of index $\mathsf{m}\in\{n-1,n\}$, the cube gauge of the clause
there: that of (F) in case ($\alpha$), that of $\zeta$ in case ($\beta$), and that of the frozen
family in case ($\gamma$). This gives the $SU(N)$-valued gauge transformation $u$ with the stated
bounds on $A$ and $\nabla A$.

\emph{The gauge $h$ plays no role.} The exposed units belong to the types $1$--$3$, $5$ and
$\partial$ of the table of unit sizes, and units of these types are functions of the gauge orbit of
the pair. Hence the gauge
transformation $h$ in (A) plays no role, and (i) is proved.

\emph{Part (ii): the box.} By Lemma~\ref{5.2:lem-exposed-geometry}, $S$ lies in an admissible cube of
index $n-1$ or $n$ whose cells of that index are at $\mathrm{dist}_n<1$ from $\Delta_y$. The cells
of $\mathfrak{r}_v$ have levels $n-1$, $n$, $n+1$.

\emph{Part (ii): the regularity argument.} With (i) in place of the corresponding bounds on $C$,
the proof of the first clause of the regularity lemma runs on $S$ at the zone index $n$ without
change. The steps are the following.
\begin{itemize}
\item Entries $6$, $7$ and $8$ are expressed in level-$(n-1)$ units, at the cost of a factor $L^2$.
\item The radii $\alpha$ of the levels $n-1$, $n$, $n+1$ lie within a factor $2$ of one another, by
the comparison of radii at neighbouring levels.
\item The Baker--Campbell--Hausdorff bound is used at the group element $V=\mathbb{1}$.
\item The cube gauge $u$ of (i) on the admissible cube is restricted to $S$.
\end{itemize}
The conclusion of that proof is that, in this gauge, $\mathrm{chain}(m)\lceil S$ equals
$(e^{i\eta\mathcal{A}^0},\mathbf{J}^0)$, with $\ell|\mathcal{A}^0|$, $\ell^2|\nabla\mathcal{A}^0|$ and
$\ell^3|\mathbf{J}^0|$ bounded by $\mathfrak{a}'_v$ of~\eqref{5.2:eq-exposed-zone-bound-3},
$\ell=\ell_{n-1}$.

\emph{Part (ii): the ceiling.} Suppose $w^{\vee}(\Delta_y)>1$, that is, a cell of weight larger
than $1$ (a weighted cell) lies at $\mathrm{dist}_n\le1$ from
$\Delta_y$. It lies in the window of an attempt $\mathfrak{a}$. The attempt may be the running
one, or a live arrest, belonging to the present component or to a foreign one. Let $k_a$, $k_{0a}$
and $N_0(a)$ be the step, the preparatory step index and the number of zones of $\mathfrak{a}$; in
these three symbols the letter $a$ refers to $\mathfrak{a}$. Then $n\le k_a$. Since
$\mathsf{L}_0^2\le L/3$ and $L^{N_0(a)}=L\check{R}_{k_{0a}+1}$,
\begin{equation*}
\bar{\mathsf{w}}_v=3w^{\vee}(\Delta_y)\le3\mathsf{L}_0^{2N_0(a)}\le3L^{N_0(a)}
=3L\check{R}_{k_{0a}+1}\le3L\,\mathsf{R}^{\uparrow}(\check{\mathsf{T}}_{k_a}).
\end{equation*}
The last step is the radius bound $\check{R}_{k_{0a}+1}\le\mathsf{R}^{\uparrow}(\check{\mathsf{T}}_{k_a})$
at the attempt's step. Moreover $\check{\mathsf{T}}_{k_a}\le\log x_{k_a}$ and
$\alpha_{*,n}\le2\alpha_*(x_{k_a})$. Hence
\begin{align*}
60dLM\mathfrak{a}'_v&\le 60dLM\,c_{\mathrm{reg}}\cdot3L\,\mathsf{R}^{\uparrow}(\log x_{k_a})\cdot
2\alpha_*(x_{k_a})\\
&=360\,dL^2M\,c_{\mathrm{reg}}\,\mathsf{R}^{\uparrow}(\log x_{k_a})\,\alpha_*(x_{k_a}).
\end{align*}
The range $x\ge\gamma^{-2}$ of the supremum in the ceiling (G) contains
$x=x_{k_a}$, whatever the age of the attempt. So $60dLM\mathfrak{a}'_v\le1$. Where
$w^{\vee}(\Delta_y)=1$, that is, no weighted cell
lies at $\mathrm{dist}_n\le1$ from $\Delta_y$, one has $\bar{\mathsf{w}}_v=3$, and
$\mathfrak{a}'_v=3c_{\mathrm{reg}}\alpha_{*,n}$ is the regularity radius of the carrier bound on $C$
at the zone index $n$, whose ceiling line is the ceiling (G) verbatim; so
$60dLM\mathfrak{a}'_v\le1$ in this case as well.

\emph{Part (iii).} With (ii) on the box $S$ at the zone index $n$, the remaining bounds are those of
the carrier bound on $C$, read at the zone index $n$ and the cell $\Delta_y$ in place of the level
$k$ and its cell.
\begin{itemize}
\item The bound for singles on boxes with $D\le D_X$.
\item The Wilson bound, pointwise in the first entry.
\item For a smooth far core of a unit whose localisation domain is disjoint from $Z$, the core
bound at the extension $\mathrm{ext}_{p_v}$ of
the pointed convention. Here one uses the bound for the renormalised pointed group valid over an
arbitrary box, applied to $Q_v$, with $\epsilon\le1$ on $c''$ times the polydisc of the
polydisc-radius lemma. By
Lemma~\ref{5.2:lem-exposed-geometry}, $Q_v=Q^{(p_v)}\subset\mathfrak{r}_v$ is itself a box of the
kind treated in (ii): its sides are at most one $\pi_{n-1}$-cube, it lies within one level-$n$
unit of $y$, and hence it lies in $\Omega_{n-1}$.
\end{itemize}
\end{proof}

\begin{lemma}[The carrier bound for exterior pointed units]\label{5.2:lem-carrier-exterior}
Let the exposed units of the class $\mathfrak{L}^{c}$, the pointed pieces $p$ of the class
$\mathfrak{L}^{Y_0}$ together with their units $v$, and the division of these pieces into
compact and long ones be as in Definition~\ref{5.2:def-exposed-compact-units}. The read set
$\mathfrak{r}_v$, the neighbourhood $\mathfrak{h}_v$, the label $\operatorname{lab}_1$, the
auxiliary consumer $\mathfrak{E}^{\oplus}_v$ with domain $X^{\oplus}_v$, the weight
$\bar{\mathsf{w}}_v=3w^{\vee}(\Delta_y)$ and the regularity radius
$\mathfrak{a}'_v=c_{\mathrm{reg}}\bar{\mathsf{w}}_v\alpha_{*,n}$ are those fixed there and in
Lemmas~\ref{5.2:lem-exposed-geometry} and~\ref{5.2:lem-exposed-zone-bound}. A unit is called
\emph{first-part} if it is plain, its domain $X_v$ meets $Z$ and $K_v=[X_v]_k\not\subset Z$, and
\emph{second-part} if $X_v\cap Z=\emptyset$. Assume the following.
\begin{itemize}
\item The third, the fourth, the sixth (in its primed form) and the seventh clause of the
inductive hypothesis of the correlation theorem. The fourth clause is used together with the
inheritance theorem and its corollary, in the following form: for every consumer
$\mathfrak{E}\in\mathcal{C}$ whose domain lies in $Y_4$, every $\zeta\in\mathfrak{S}_{\mathrm{src}}$,
every real $\Phi$ admitted by that clause, every $|\sigma|\le e^{\kappa_1}$ and every chain background
$\mathfrak{b}\in\{\mathfrak{b}_2,\mathfrak{b}_3,\mathfrak{b}_4\}$, the presented pair
$((\hat{\mathcal{U}}^{\mathfrak{b}})^h,\operatorname{Ad}_h\hat{\mathcal{J}}^{\mathfrak{b}})$,
restricted to that domain, lies in $\mathfrak{D}_{\mathfrak{E}}$ and depends jointly
holomorphically on $(\zeta,\sigma)$ and on the amplitude $\mathsf{t}$ of (c0), along the
segments $t\in[0,1]$. The sixth clause is used with
the smoothing estimate for the extension $\operatorname{ext}_v$ over an arbitrary box. The seventh
clause contains $\kappa\ge 2(d+1)\log L$.
\item The regularity ceiling
\begin{equation}\label{5.2:eq-carrier-exterior-c}
\sup_{x\ge\gamma^{-2}}\,405\,d\,L^{2}M\,c_{\mathrm{reg}}\,
\mathsf{R}^{\uparrow}(\log x)\,\alpha_*(x)\le 1,
\end{equation}
together with the same display with $360$ in place of $405$.
\item The reading convention for cores: a smooth far core, and its subtracted copy, are read
at one and the same extension $\operatorname{ext}_v$.
\item On live pieces: the freezing of the structure $(\operatorname{lev}_{\mathfrak{a}},w_{\mathfrak{a}},
\mathcal{Q}_{\mathfrak{a}})$ of a live arrest $\mathfrak{a}$ since its step, and the fact that the
members of the arrest-read source space are restrictions of configuration pairs of the
admissible family over the determining set $\mathfrak{B}$.
\item The containment supplement at the localisation site, with its collar clause (locations
in the collar and ring of $C$), its $\mathfrak{T}$-location clause ($\mathfrak{T}:=\Omega_k\cup
Z^{c}$) and its foreign-piece clause, under their ceilings: the floor
\begin{equation}\label{5.2:eq-carrier-exterior-1}
C_{\theta'}\ge 192\,(1+b_+)^{1/2}K_D\,c_{\mathrm{rec}}\,\varphi_{\flat}^{-1}B_HA_1/C_*
\end{equation}
(the supplement's floor with $c_\chi$ enlarged to
$c_{\mathrm{rec}}:=\max\{c_\chi,c_J^{\mathrm{rec}},c_a\}$; here $\alpha_*=C_*g\mathsf{T}^{q}$ and
$\delta'=gA_1\mathsf{T}^{p_1}$), and the ceiling on the member of the supplement's
$\mathfrak{T}$-location clause, which serves the foreign-piece clause as well, with $C_P'$
replaced by $c_{\mathrm{rec}}C_P'$:
\begin{equation}\label{5.2:eq-carrier-exterior-2}
\sup_{\mathsf{T}\ge\mathsf{T}_\gamma-1}32\,C_{\theta'}^{-1}(\mathsf{T}+2)^{q-p_1}\,
c_{\mathrm{rec}}C_P'\,e^{-(3/2)\delta_1^{\mathrm{sc}}R(\mathsf{T})}\le\beta_{\mathrm{pr}}\theta_S;
\end{equation}
and, for the touching-block containment of (c5)(i) only, all five clauses of the supplement
(these three and the two clauses for clause locations inside the large-field region $Z_{h}$ at
the supplement's base level $h$ of the component $C$, the levels of the collar and ring of $C$
lying in $[h,k]$), together with the supplement's half-exponent membership ceiling, its constant
$c_\chi$ replaced by $c_{\mathrm{rec}}$.
\item For (c2): the receding-block tail lemma with its first and second corollaries, under the
floors on $\kappa_1$, on the levels and on $\gamma$ that accompany the fourth clause, the
amplitude bounds $\mathsf{a}_4\le6\delta'_k$ and $\mathsf{a}_s\le\mathsf{a}^{\flat}$, a number,
where $\mathsf{a}_s$ is the supremum over $\mathfrak{P}$ of the source datum of the bracket $s$
(for $s=3$ together with its current term), the floor
\begin{equation}\label{5.2:eq-carrier-exterior-3}
\tfrac{1}{16}\,\delta_P\,\mathfrak{w}\ge c_{\boxplus}\log\hat b_t+1,\qquad \mathfrak{w}:=M/M_1,
\end{equation}
the nesting of the site geometry with its hub clauses (in particular the hubs
$\Lambda^{\sim}_C$ are boxes of sides at most $h_\sharp:=112$, at mutual distance at least
$19\check R_k$, and $\operatorname{dist}_k(\Lambda^{\sim},Z^{c})\ge 9\check R_k$), the reference
comparison of the couplings for all pairs of levels $\le k$ (with $N_k\le x_k$), and
\begin{gather}
C_{\theta'}\ge 96\,K_{\mathsf{J}}C^{\wedge}_{\mathrm{age}}A_1/(\beta_{\mathrm{pr}}\theta_S
\varphi_{\flat}C_*),\label{5.2:eq-carrier-exterior-4}\\
\sup_{x\ge\gamma^{-2}}16\,K_{\mathsf{J}}C^{\wedge}_{\mathrm{age}}\,\mathsf{a}^{\flat}
(\beta_{\mathrm{pr}}\theta_S\varphi_{\flat}C_{\theta'})^{-1}(\log x+2)^{q-p_1}
\alpha_*(x)^{-1}E(x)\le 1,\label{5.2:eq-carrier-exterior-5}
\end{gather}
where $C^{\wedge}_{\mathrm{age}}:=2\sqrt{2}\,(1+b_+)^{1/2}$.
\end{itemize}
All of these are read with the positional factor $\eta_p$ of (c0) below.

Then for every $K$, every $k\le K$, every label and every exposed unit $v$, whatever the zone
$n$ of its point $y$, the following hold.
\begin{itemize}
\item[(a)] The map $m\mapsto f_v(m)$ is holomorphic on $\mathfrak{P}$.
\item[(b)] For every $m\in\mathfrak{P}$,
\begin{equation*}
|f_v(m)-f_v(1)|\le\mathfrak{N}^{\mathrm{ex}}(v):=\mathfrak{N}^{\flat}(v),
\end{equation*}
the entry of the size table at this $\bar{\mathsf{w}}_v$ and this $n$, with $\mathfrak{a}'_v$ as
above, a rough core read termwise, the Wilson piece with the constant $c_5^{W}$, no depth factor,
and $\bar{\mathsf{w}}_v=3$ unless $\Delta_y$ neighbours the window of a foreign live arrest.
Hence, by the Cauchy bounds in the cut decoupling parameters, the bound of the exterior activity
lemma holds on these units with $\mathfrak{N}$ replaced by $\mathfrak{N}^{\mathrm{ex}}$; and the
value $\mathrm{val}_v=f_v(\mathrm{chain}(\vec 0))$ of an exposed unit of $\mathfrak{L}^{c}$ obeys
$|\mathrm{val}_v-f_v(1)|\le\mathfrak{N}^{\flat}(v)$.
\item[(c)] Every compact pointed piece $p$ of a first-part or second-part unit is read on the
amplitude disc $|\mathsf{t}|\le\rho_A$ of (c0), and the following hold.
\begin{itemize}
\item[(c0)] The radius, the positional factor, the read set and the far factors are
\begin{equation}\label{5.2:eq-carrier-exterior-a}
\begin{gathered}
\rho_A:=\frac{\mathsf{w}_L}{(\eta_p\mathsf{f}_v)^{1/2}},\qquad
\eta_p:=\begin{cases}1,&\operatorname{Sten}_2(S_v)\cap\square^{\zeta}\neq\emptyset,\\
e^{-\delta_P d(\square^{\zeta},S_v)},&\text{otherwise},\end{cases}\\
S_v:=\mathfrak{r}_v\ \ (p\ \text{compact}),\qquad
S_v:=\{\text{cells meeting }X_v\}\ \ (p\ \text{long}),\\
\mathsf{f}(y'):=e^{-\delta_{\natural}(d_{\mathrm{sc}}(y',\mathfrak{F}_{\mathfrak{s}})-1)_+},\qquad
F_s(y'):=\Theta_s(y')/\mathsf{U}\le\mathsf{f}(y'),\\
\mathsf{f}_v:=\sup_{y'\in S_v}\min\{1,c_\chi\mathsf{f}(y')\}\in\,]0,1].
\end{gathered}
\end{equation}
Here $\square^{\zeta}:=\operatorname{supp}\hat{\chi}_{\square}$; the block $\square$
\emph{touches} when $\eta_p=1$ and \emph{recedes} otherwise; a cell $y'$ of $S_v$ means a unit
block meeting $S_v$; $\mathfrak{F}:=\mathfrak{F}_{\mathfrak{s}}$ is $\mathsf{S}$ for the
first-part piece of the fourth bracket and $\Lambda^{\sim}$ for the second-part pieces of the
four brackets, the third bracket's source term read along the pieces of its localisation;
$\Theta_s(y')$ is the majorant of the cell norms of the undilated layer $\mathbb{H}_s$ at $y'$; and
$\mathsf{U}=\mathsf{U}_J$ for the first-part piece, $\mathsf{U}=\mathsf{U}_L$ for the four brackets
of a second-part piece. Consequently $F_s\le\mathsf{f}\le\mathsf{f}_v$ on $S_v$ (as
$c_\chi\ge1$ and $\mathsf{f}\le1$), $\rho_A\ge\mathsf{w}_L/\eta_p^{1/2}$, and
$\mathsf{f}^{\wedge}_v:=e^{-\delta_Pd^{\wedge}(S_v,\mathfrak{F})}\le\mathsf{f}_v$, because
$d^{\wedge}\ge \mathfrak{w}\,d_{\mathrm{sc}}$ and $\delta_P\mathfrak{w}\ge1\ge\delta_{\natural}$.
\item[(c1)] The block piece $\delta\mathbb{H}_p=\hat{\chi}_{\square}\mathbb{H}_s$ is supported
in $\square^{\zeta}\subset\square$ (see \cite[p.~270]{B-CMP109}), and the local entries at a cell
read the layer on that cell's stencil, which lies in $\operatorname{Sten}_2(S_v)$. Hence at a
cell $y'$ of $S_v$ the local members of the entries $Q_a$, $a\in\{1,3,6,7\}$, of the dilated
carrier move only at touching blocks; at a receding block they do not move at all, the
$\mathbf{J}$-member of $Q_3$ being the only mover there, treated in (c2)$(\alpha)$. At a touching block,
on $|\mathsf{t}|\le\rho_A$,
\begin{equation}\label{5.2:eq-carrier-exterior-b}
\begin{split}
\mathsf{N}_{y'}(\mathsf{t}\hat{\chi}_{\square}\mathbb{H}_s)
&\le|\mathsf{t}|\,c_\chi\mathsf{U}F_s(y')
\le\mathsf{w}_L\mathsf{f}_v^{-1/2}c_\chi\mathsf{U}F_s(y')\\
&\le c_\chi\mathsf{w}_L\mathsf{U}\,F_s(y')^{1/2}
\le c_\chi\mathsf{w}_L\mathsf{U}\,\min\{1,\mathsf{f}(y')^{1/2}\},
\end{split}
\end{equation}
the far factor $\mathsf{f}$ kept for $\mathsf{U}=\mathsf{U}_L$; this is the supplement's
half-exponent line with its half exponent intact, the extra factor $\mathsf{f}_v^{-1/2}$ of the
radius being paid by $F_s\le\mathsf{f}_v$ and never by the supplement's own far factor. Whence
three clauses of the supplement hold at $b(y')=0$, the clause index of the consumer being
$\pi(y')$, exactly as for boundary terms. Its collar clause covers first-part pieces at the collar
and ring of $C$, with $\mathsf{U}=\mathsf{U}_J$, the ratio $\mathsf{w}_L\mathsf{U}_J/\alpha_{*,k}$
being free of the coupling, under \eqref{5.2:eq-carrier-exterior-1}, the factor
$\mathsf{f}^{1/2}\le1$ dropped. Its $\mathfrak{T}$-location clause covers first- and second-part
pieces at $\mathfrak{T}$-locations, $\mathsf{U}\in\{\mathsf{U}_J,\mathsf{U}_L\}$: the block piece
stands in place of the layer, the prefactor is multiplied by $\mathsf{w}_L$ and the exponent is
halved, which is what $c_\chi\mathsf{w}_L\mathsf{U}F_s(y')^{1/2}$ is; the shell factor
$2^{-(k-1-m(y'))}$ of an untouched shell of index $m(y')$ survives the halving, the constant being
$c_{\mathrm{rec}}C_P'$, and the increment is at most $\tfrac14\mathrm{room}^{\flat}$ under
\eqref{5.2:eq-carrier-exterior-2}. Its foreign-piece clause covers second-part pieces on or
beside a foreign live piece: the far factors are halved to $e^{-1}$ per stratum, with the rooms of
the foreign arrest, under its ceiling. The constants of the three clauses are enlarged once by
$c_{\mathrm{rec}}$: \eqref{5.2:eq-carrier-exterior-1} and the membership ceiling carry
$c_{\mathrm{rec}}$ in place of $c_\chi$, \eqref{5.2:eq-carrier-exterior-2} carries
$c_{\mathrm{rec}}C_P'$ in place of $C_P'$, and the block constant of the supplement's
half-exponent step bound is multiplied by $e^{2d+2}$. Then every additive increment is at most
$\theta_a/16$, and for $a\in\{2,4\}$ one has $Q_a^{+}\le\tfrac{17}{16}Q_a+\theta_a/16$. The two
clauses for clause locations in $Z_{h}$ are void for these pieces. At receding blocks the
entries of the $\mathbf{U}$-slot do not move on $S_v$ at all,
since $\operatorname{Sten}_2(S_v)$ misses $\square^{\zeta}$ and so $\hat{\chi}_{\square}$ is
constant on every stencil of $S_v$.
\item[(c2)] The dilated carrier is the family $\mathsf{R}_p(\mathsf{t})=\Phi_s^{\chi}$ at
$\chi:=t+\mathsf{t}\hat{\chi}_{\square}$, whose current
\begin{equation*}
\begin{split}
\mathcal{J}^{\sharp}_s[\chi]:=&\,(\Delta_{W_s}-P_1(\sigma_s)+\mathsf{K})(\chi A'_s)
-\chi\,\mathfrak{m}_H(\sigma_s)X_s\\
&-[D^{*}D,\chi]\,H(\sigma_s)X_s+F(W_s,\chi\mathbb{H}_s),
\end{split}
\end{equation*}
with $\Delta_{W_s}$, $P_1(\sigma_s)$, $\mathsf{K}$, $\mathfrak{m}_H(\sigma_s)$, $H(\sigma_s)$ and
$F$ the operators and the function of the layer construction in the proof of the fourth clause,
at the bracket $s$ (the construction of the response layers $\mathbb{H}_s$),
has exactly one member non-local in $\chi$, namely $-P_1(\sigma_s)(\chi A'_s)$, and it is linear
in $\chi$.

$(\alpha)$ At a receding block, $\mathsf{R}_p(\mathsf{t})\lceil S_v$ (here and below the sign
$\lceil S$ denotes restriction to $S$) is affine in $\mathsf{t}$: its
$\mathbf{U}$-slot is free of $\mathsf{t}$ and its $\mathbf{J}$-slot is its value at $\mathsf{t}=0$ minus
$\mathsf{t}\,T_p\lceil S_v$, $T_p:=P_1(\sigma_s)(\hat{\chi}_{\square}A'_s)$, which is the whole
$\mathsf{t}$-dependence there. A reader of the first slot only (a Wilson piece; a smooth core
read at $\operatorname{ext}_v$) has $v_p=0$; the pointed pieces that read $\mathbf{J}$ are the
singles, through the single-term bound $F(X,(\mathbf{U},\mathbf{J}))$ and the regularity bound
$\ell^{3}|\mathbf{J}^0|\le\mathfrak{a}'_v$, a restriction belonging to the present classification
and not to the layer estimate of that construction, whose inequality transfers, since its proof
reads only the class of the kernel and $\|\mathsf{d}_s\|_{p_s}\le1$. The entry
$Q_3$ is the only mover: $Q_4=|J_{p,X}(M'(\mathbf{U}))|$ is a function of the
$\mathbf{U}$-slot (see \cite{B-CMP122a} and, for the model, \cite[(1.8)]{B-CMP109}). The tail obeys, at every
cell $y'$ and $x\in\Delta(y')$,
\begin{equation}\label{5.2:eq-carrier-exterior-6}
\ell(y')^{3}|T_p(x)|\le 20K_PB_\sharp\,p_s(y')\,e^{-\delta_Pd(y',\square^{\zeta})}
\le 20K_PB_\sharp\,\mathsf{a}_s\,e^{-\delta_Pd^{\wedge}(y',\mathfrak{F})}
e^{-\delta_Pd(y',\square^{\zeta})},
\end{equation}
the far factor at $y'$ times the block decay to $y'$, both at the rate $\delta_P$, and not a
minimum; whence, on $|\mathsf{t}|\le\rho_A$, for every $y'\in S_v$ and every clause index
$n\le\pi(y')$, in the presented pair,
\begin{equation}\label{5.2:eq-carrier-exterior-7}
|\mathsf{t}|\,\ell_n^{3}|T_p(y')|
\le K_{\mathsf{J}}\mathsf{w}_L\mathsf{a}_s\,L^{-3(\pi(y')-n)}\,\eta_p^{1/2}
e^{-\frac12\delta_Pd^{\wedge}(y',\mathfrak{F})}
\le\varepsilon_{\mathsf{J}}\,2^{-u(y')}\varphi_{\flat}\alpha_{*,n},
\end{equation}
with $K_P$, $B_\sharp$, $K_\chi$ the constants of the layer bounds of that construction,
$K_{\mathsf{J}}:=42\max\{K_PB_\sharp,K_\chi\}$, $\varepsilon_{\mathsf{J}}:=
\beta_{\mathrm{pr}}\theta_S/16$, $p_s(y'):=\mathsf{a}_se^{-\delta_Pd^{\wedge}(y',A_s)}$ and
$u(y'):=d^{\wedge}(y',\mathfrak{F})/\mathfrak{w}$, in the notation of the receding-block tail lemma. So $Q_3$
moves by at most $\beta_{\mathrm{pr}}\theta_S/16\le1/16$ of every threshold $\theta_a$, at
$y'\notin Z$ by at most $\tfrac14\mathrm{room}^{\flat}(y')$, and nothing else moves.

$(\beta)$ At a touching block ($\eta_p=1$, $|\mathsf{t}|\le\mathsf{w}_L\mathsf{f}_v^{-1/2}$) the
$\mathsf{t}$-member of $Q_3$ is at most $20K_\chi|\mathsf{t}|q_{\square}(y')\le
20K_\chi|\mathsf{t}|p_s(y')$, at every block and every $y'$, and the two-sided display
\eqref{5.2:eq-carrier-exterior-7} holds with the same $K_{\mathsf{J}}$, since
\begin{equation*}
p_s(y')\le\mathsf{a}_se^{-\frac12\delta_Pd^{\wedge}(y',\mathfrak{F})}(\mathsf{f}^{\wedge}_v)^{1/2}
\le\mathsf{a}_se^{-\frac12\delta_Pd^{\wedge}(y',\mathfrak{F})}\mathsf{f}_v^{1/2}
\end{equation*}
on $S_v$.

So, by (c1) and (c2), at every cell of $S_v$ and every block each moving entry gains at most
$1/16$ of its threshold, and the clauses of the supplement deliver
$\mathsf{R}_p(\mathsf{t},\sigma)\lceil S_v\in\mathfrak{D}_v$ at every block.
\item[(c3)] By (c1) and (c2), every entry of the dilated carrier at $S_v$ is less than
$\theta_a+\theta_a/16+\theta_a/16\le\tfrac98\theta_a$ on $\mathfrak{P}\times\{|\mathsf{t}|\le
\rho_A\}$, the undilated entries being less than $\theta_a$. So the regularity bound of
Lemma~\ref{5.2:lem-exposed-zone-bound} runs at $\tfrac98\mathfrak{a}'_v$ under
\eqref{5.2:eq-carrier-exterior-c}; the single-term, smoothing and Wilson-line estimates are
quadratic in $\mathfrak{a}'$, and the depth function of a second-part core is read at
$\tfrac98\bar{\mathsf{w}}_v$; hence
\begin{equation*}
|f_v\circ(\text{dilated carrier})-f_v(1)|\le\tfrac{81}{64}\mathfrak{N}^{\flat}(v)
\end{equation*}
there. By
Cauchy's inequality in $\mathsf{t}$ at radius $\rho_A$ (the derivative in $\mathsf{t}$ kills
$f_v(1)$), with the walk factor of the amplitude lemma and the Cauchy bounds in the cut
decoupling parameters,
\begin{equation}\label{5.2:eq-carrier-exterior-d}
\begin{gathered}
|(p;\mathsf{y})|\le\mathfrak{N}^{\mathrm{ex}}_A(p)\,e^{-(\kappa_1-1)\sum_t|\mathsf{y}_t|},\\
\mathfrak{N}^{\mathrm{ex}}_A(p):=2\mathsf{w}_L^{-1}(\eta_p\mathsf{f}_v)^{1/2}\cdot
\tfrac{81}{64}\,\mathfrak{N}^{\flat}(v)\cdot e^{-(\kappa_1-1)m(p)},
\end{gathered}
\end{equation}
one disc for all tuples of $p$: the amplitude lemma's bound for boundary terms with
$\eta_p^{1/2}$ replaced by $(\eta_p\mathsf{f}_v)^{1/2}\le\eta_p^{1/2}$ and the size
$\tfrac{81}{64}\mathfrak{N}^{\flat}(v)$. The sum of these bounds over $(v,s,\square,\mathsf{y})$
is performed by the summation of compact pointed pieces, not by the one-scale resummation:
$\eta_p^{1/2}\le e^{-\frac12\delta_Pd(\tilde{\square},S_v)}$ (at a receding block
$\square^{\zeta}\subset\tilde{\square}$; at a touching one $S_v$ meets $\tilde{\square}$) sums the
blocks in shells, the walk factor sums the cells in their own scales, and the factor
$\mathsf{f}_v^{1/2}\le1$, independent of the block, is kept.
\item[(c4)] Under the geometric clause
\begin{equation}\label{5.2:eq-carrier-exterior-e}
\begin{gathered}
d_{\mathrm{sc}}(y',\Lambda^{\sim})\ge 2\check R_{k_{a'}}-6\ge\check R_{k_{a'}}-5
\quad\text{for every cell $y'$ of $S_v$,}\\
\bar{\mathsf{w}}_v\le\bar{\mathsf{w}}^{\mathrm{fw}}(k_{a'}):=
3\mathsf{L}_0^{2}LC_{\uparrow}\check R_{k_{a'}},
\end{gathered}
\end{equation}
which holds whenever the weighted cell of $p$ (the cell at $\operatorname{dist}_n\le1$ from
$\Delta_y$ where $w^{\vee}(\Delta_y)>1$ is attained) lies in the window of a foreign live
arrest $\mathfrak{a}'$ of step $k_{a'}$, one has
$\mathsf{f}_v\le c_\chi e^{-\delta_{\natural}(\check R_{k_{a'}}-6)_+}$ and
\begin{equation}\label{5.2:eq-carrier-exterior-f}
\mathsf{f}_v^{1/2}\,\bar{\mathsf{w}}_v^{3}\le\mathsf{C}_{\mathrm{fw}}:=
c_\chi^{1/2}(3\mathsf{L}_0^{2}LC_{\uparrow})^{3}\,
\sup_{\iota\ge0}(\iota+6)^{3}e^{-\frac12\delta_{\natural}(\iota-1)_+},
\end{equation}
a number depending on $(d,L,\mathsf{L}_0,C_{\uparrow},\delta_{\natural},c_\chi)$ only, in which no
$k_{a'}$, $k-k_{a'}$, $N_0(a')$, window size or arrest count survives. Where
$\bar{\mathsf{w}}_v=3$, $\mathsf{f}_v^{1/2}\bar{\mathsf{w}}_v^{3}\le27\le\mathsf{C}_{\mathrm{fw}}$,
since $\mathsf{f}_v\le1$ and $\mathsf{C}_{\mathrm{fw}}\ge27\cdot6^{3}c_\chi^{1/2}$.
Second-part units carry only $\mathsf{t}$-pieces (their undilated first member is kept whole), so
every exterior piece that meets a foreign weight carries the factor $(\eta_p\mathsf{f}_v)^{1/2}$;
first-part pieces lie at $\partial C$, where $\bar{\mathsf{w}}=3$. Containment on the foreign piece
is that of the supplement's foreign-piece clause, with the bound
$\varepsilon_{\mathsf{J}}2^{-u(y')}\le\tfrac14\mathrm{room}^{\flat}$ of (c2) for the
$\mathbf{J}$-entry.
\item[(c5)] Let $p$ be long: $v=F(X,\cdot)-F(X,1)$, $y\in X\in\mathbf{D}_j$, a long single of
rows $2$, $3$ or the boundary row of the size table, or a row-$1$ tail, of a plain first-part or
second-part unit. Put $S_v:=$ the cells meeting $X_v$, $\operatorname{lab}_1:=K_p$,
$\mathfrak{D}_v:=U^{c}_j(X_v,\alpha_{\cdot,j})$ of \cite[(1.11)--(1.16)]{B-CMP109} at every cell of
$X_v$, foreign-window cells included, whatever the entry time (at $j=k$ with the radii
$\hat{\alpha}_{k-1}$), and $\mathfrak{N}^{\flat}(v)=2\mathfrak{n}_v$; no
$\bar{\mathsf{w}}_v$, $\mathfrak{r}_v$, $\mathfrak{h}_v$, $\mathfrak{E}^{\oplus}_v$ is used, and
$\eta_p$, $\mathsf{f}_v$, $\rho_A$ of (c0) are taken on this $S_v$. Then:
(i) $\mathsf{R}_p(\mathsf{t},\sigma)\lceil X\in\mathfrak{D}_v$ on
$\mathfrak{P}\times\{|\mathsf{t}|\le\rho_A\}$ at every block: at a touching block by the five
clauses of the supplement read at the unit's own clause index $n(y')=j$, with
$b(y')=\pi(y')-j\ge0$ and the scale ratio of the supplement equal to $L^{-b(y')}$, the half-exponent
line
\begin{equation*}
|\mathsf{t}|\mathsf{N}_{y'}(\delta\mathbb{H}_p)\le c_\chi\mathsf{w}_L\mathsf{U}\min\{1,
\mathsf{f}(y')^{1/2}\}
\end{equation*}
holding at every $y'$ of $S_v\supseteq X$ because $F_s\le\mathsf{f}\le
\mathsf{f}_v$ there, and the two clauses for $Z_{h}$ being in force for first-part $X$ reaching
$Z_{h}$; at a receding block by (c2)$(\alpha)$, which holds at every index $n\le\pi(y')$ with the
scale drop $L^{-3(\pi(y')-n)}$; over the undilated base, supplied for first-part units by the
fourth clause at the unit's own family, for the brackets of a second-part unit by the chain
clauses at the source site (with the chain theorem beside them), by the three cases of the
inheritance claim, and on a foreign live piece by the first-zone bracket of the live-arrest
containment statement, through its lemmas for unraised and for raised cells; that base is read at
a foreign live piece with the following
monotonicity: at a cell $x$ of a foreign live piece $W_{\mathfrak{a}'}$ (the innermost arrest
at $x$) met by $X_v$, with $j_v\le k_{a'}$, $\operatorname{lev}^{0}(x)\ge j_v$ the level at $x$
before the attempt and
frozen structure datum $(\operatorname{lev}_{\mathfrak{a}'}(x),w_{\mathfrak{a}'}(x))$ at $x$, where
$\operatorname{lev}_{\mathfrak{a}'}(x)=\operatorname{lev}^{0}(x)+g_x$ and
$w_{\mathfrak{a}'}(x)\le\mathsf{L}_0^{2g_x}$
for an integer $g_x\ge0$, the unit's
clause at $x$ contains the stage-zero class of $\mathfrak{a}'$ at $(\operatorname{lev}^{0}(x),1)$,
in which the
live-arrest containment statement places the image of the arrest-read source within its room;
hence the image lies
in $\mathfrak{D}_v$, and no constant depends on $k$, $k_{a'}$, $g_x$, $N_0$ or an age.
(ii) $|f_v\circ\mathsf{R}_p-f_v(1)|\le2\mathfrak{n}_v$ there, whence
\begin{equation*}
|(p;\mathsf{y})|\le2\mathsf{w}_L^{-1}(\eta_p\mathsf{f}_v)^{1/2}\cdot2\mathfrak{n}_v\cdot
e^{-(\kappa_1-1)m(p)}e^{-(\kappa_1-1)\sum_t|\mathsf{y}_t|},
\end{equation*}
the display \eqref{5.2:eq-carrier-exterior-d} at $\mathfrak{N}^{\flat}=2\mathfrak{n}_v$ without the
factor $\tfrac{81}{64}$. (iii) $\#\operatorname{St}_v\le\mathsf{c}_d(d_j(X)+1)$. The long pieces,
first-part and second-part, are summed in $\square$-bundles at $\mathsf{t}=0$ by the summation of
boundary terms at finer zones, the blocks resummed exactly; there only (i) at $\mathsf{t}=0$, the
undilated size $2\mathfrak{n}_v$ of (ii), and (iii) enter.
\end{itemize}
\item[(d)] A smooth far core of a second-part unit is read at $\operatorname{ext}_v$, with
$p_v:=n-2-s^{(9/8)}(\bar{\mathsf{w}}_v)$, $s^{(9/8)}(\bar{\mathsf{w}}):=s(\tfrac98\bar{\mathsf{w}})$,
$s(\cdot)$ the depth function of the smooth-core entry of the size table; the smoothing estimate on
$\mathfrak{P}\times\{|\mathsf{t}|\le\rho_A\}$ gives $2\epsilon^{2}\mathfrak{N}_{p_v}$ with the
dilated $\epsilon\le1$, that is, the smooth-core entry of the size table read at
$\tfrac98\bar{\mathsf{w}}_v$, at most $\tfrac{81}{64}$ times its $\bar{\mathsf{w}}^{2}$-factor at
$a_v\ge2+s^{(9/8)}$. Rough cores are read termwise, their constituents being short singles or long
ones, all of them at $\nu\le2$ in the latter case (c5). The per-cell polynomial in $s$ is read at
$s^{(9/8)}\le s(1)+\log_L\bar{\mathsf{w}}_v+1$, whence the third power:
$(\tfrac98)^{3}=\tfrac{729}{512}\le\tfrac32$ wherever the dilated sizes are summed.
\end{itemize}
\end{lemma}

\begin{proof}[Proof of Lemma~\ref{5.2:lem-carrier-exterior}: clauses (a), (b), (c1) and (c2)]
\label{5.2:prf-carrier-proof-local}
Throughout, $v$ is an exposed unit, $y_v$ its point, and, where clauses (c1) and (c2) are
concerned, $p$ is a compact pointed
piece of $v$ with read set $S_v=\mathfrak{r}_v$. The letters $\eta_p$, $\mathsf{f}_v$,
$\mathsf{f}_v^{\wedge}$, $F_s$ and $\rho_A=\mathsf{w}_L/(\eta_p\mathsf{f}_v)^{1/2}$ are those of
\eqref{5.2:eq-carrier-exterior-a}. A block $\square$ of the cover of $\mathfrak{B}_s$
\emph{touches} if $\operatorname{Sten}_2(S_v)$ meets $\square^{\zeta}=\operatorname{supp}
\hat{\chi}_{\square}$, so that $\eta_p=1$, and \emph{recedes} otherwise, so that
$\eta_p=e^{-\delta_P d(\square^{\zeta},S_v)}<1$. For a cell $y$ we write $\ell(y)=\ell_{\pi(y)}$,
where $\ell_n$ denotes the length unit of level $n$, so that $\ell_n=L^{-(\pi(y)-n)}\ell(y)$ for
$n\le\pi(y)$; the far set is $\mathfrak{F}=\mathsf{S}$ for class (c) and $\mathfrak{F}=\Lambda^{\sim}$
for class (d); and $u(y):=M_1\,d^{\wedge}(y,\mathfrak{F})/M$, where $M/M_1$ is the scale ratio of
the determining sets.

\emph{Clause (a).} The carrier $m\mapsto\mathrm{chain}(m)$, restricted to the domain
$X^{\oplus}_v$ of the auxiliary consumer $\mathfrak{E}^{\oplus}_v$, is holomorphic on
$\mathfrak{P}$ with values in $\mathfrak{D}_{\mathfrak{E}^{\oplus}_v}$. This is the consumer
containment theorem, in both of its clauses (containment and holomorphy), applied at
$\mathfrak{E}:=\mathfrak{E}^{\oplus}_v$. Its hypothesis asks that $\mathfrak{E}$ belong to the
consumer class $\mathcal{C}$ and that its domain lie in $Y_4(\mathfrak{p})$. Both were
established with the definitions preceding Lemma~\ref{5.2:lem-carrier-exterior}:
$\mathfrak{E}^{\oplus}_v\in\mathcal{C}$, and
$X^{\oplus}_v\subset Y_4(\mathfrak{p})$ for every tuple $\mathfrak{p}$ of $v$ (resp.\ of $p$).
The unit $f_v$ is analytic on its own domain, by the hypothesis of the lemma on the analyticity
of the $\mathbf{E}$- and $\mathbf{R}$-families; the families of class (d) share the domain
$U^c_j(X,\alpha_{\cdot,j})$. Composition gives the holomorphy of $m\mapsto f_v(m)$ on
$\mathfrak{P}$. For a smooth core the unit is read through the extension $\mathrm{ext}_p$, which
is analytic on the polydisc of the radii lemma dilated by the factor $c''$, by the
hypothesis of Lemma~\ref{5.2:lem-carrier-exterior} on the smoothing extension. Composition again
gives (a).

\emph{Clause (b).} By the bounds for exposed units at every zone
(Lemma~\ref{5.2:lem-exposed-zone-bound}) we have three statements at the zone index
$n:=\pi(y_v)$ of $y_v$:
the current-entry bound, the cube statement, and the per-zone regularity statement. Together they
supply on $S_v$, at zone $n$, every input that the size estimates of the interior carrier bound
read at zone $k$. Those estimates therefore apply word for word at zone $n$. They cover:
\begin{itemize}
\item the short single;
\item the long single and the tail;
\item the Wilson piece;
\item the smooth core, together with clause (d);
\item the rough core, read termwise.
\end{itemize}
Here $c_\vee$ denotes the constant, fixed with the interior carrier bound, by which lengths
measured in level-$j$ units are converted into level-$n$ units. This
gives $|f_v(m)-f_v(1)|\le\mathfrak{N}^{\flat}(v)$ for every $m\in\mathfrak{P}$. The value
$\mathrm{val}_v=f_v(\mathrm{chain}(\vec 0))$ is $f_v$ at the point of $\mathfrak{P}$ that
corresponds to the all-zero tuple, so the bound on the value is a special case.

\emph{Clause (c1).} Fix a cell $y$ of $S_v$ and a block $\square$; from here on $y$ denotes such a
cell, not the point $y_v$. We measure the increment of an
entry under the dilation in units of the threshold $\theta_a(y)$ read at the consumer
$\mathfrak{E}^{\oplus}_v$. By clauses $(\alpha)$--$(\gamma)$ of the current-entry bound, the
clause index of that consumer at $y$ is $\pi(y)$ itself. Hence, in the linearisation estimate of
the containment rider at the localisation site, the index drop is $b(y)=0$ and the scale factor is
$\varrho_n=1$. That estimate is linear in the cell majorant of the layer, and gives for this
increment the bound
\[
K_D\,\varphi_{\flat}^{-1}\,
\mathsf{N}_y(\mathsf{t}\hat{\chi}_{\square}\mathbb{H}_s)/\alpha_{*,\pi(y)} .
\]
Next, $\Theta_s$, a weight of the class $\mathfrak{M}^{\wedge}_{\delta_P/2}$ of majorants of rate
$\delta_P/2$ in $d^{\wedge}$, dominates the cell norms of $\mathbb{H}_s$
everywhere; this domination is the hypothesis of the layer-pricing proposition, supplied by the
majorant property of the localised current. The first inequality below is the multiplier bound of
the block cutoff $\hat{\chi}_{\square}$, with constant $c_\chi$. Hence
\[
\mathsf{N}_y(\hat{\chi}_{\square}\mathbb{H}_s)\le c_\chi\,\mathsf{N}_y(\mathbb{H}_s)
\le c_\chi\,\Theta_s(y)=c_\chi\,\mathsf{U}\,F_s(y).
\]
Moreover $\hat{\chi}_{\square}\mathbb{H}_s$ vanishes on the stencil of $y$ unless
$\operatorname{Sten}_2(S_v)$ meets $\square^{\zeta}=\operatorname{supp}\hat{\chi}_{\square}$. In
that case the block touches, $\eta_p=1$ by \eqref{5.2:eq-carrier-exterior-a}, and
$\rho_A=\mathsf{w}_L\mathsf{f}_v^{-1/2}$. On $|\mathsf{t}|\le\rho_A$ we therefore obtain
\begin{align*}
\mathsf{N}_y(\mathsf{t}\hat{\chi}_{\square}\mathbb{H}_s)
&\le |\mathsf{t}|\,c_\chi\,\mathsf{U}\,F_s(y)
 \le \mathsf{w}_L\,\mathsf{f}_v^{-1/2}\,c_\chi\,\mathsf{U}\,F_s(y)\\
&= c_\chi\,\mathsf{w}_L\,\mathsf{U}\,F_s(y)^{1/2}\,
   \bigl[F_s(y)/\mathsf{f}_v\bigr]^{1/2}
 \le c_\chi\,\mathsf{w}_L\,\mathsf{U}\,F_s(y)^{1/2}\\
&\le c_\chi\,\mathsf{w}_L\,\mathsf{U}\,\min\{1,\mathsf{f}(y)^{1/2}\}.
\end{align*}
The penultimate inequality uses $F_s\le\mathsf{f}\le\mathsf{f}_v$ on $S_v$, the last uses
$F_s\le\mathsf{f}\le 1$.
This is the displayed line \eqref{5.2:eq-carrier-exterior-b}.

For a piece of class (c) we have $\mathsf{U}=\mathsf{U}_J$ and
$\mathsf{w}_L\mathsf{U}_J\le\tilde{\alpha}_1\le\bar{\alpha}$, where $\tilde{\alpha}_1$ is the
second radius of the source space $\tilde U^c_k(Y_4,\tilde\alpha_0,\tilde\alpha_1)$. The line just
proved is then the half-exponent line of the containment rider, which bounds
$|\mathsf{t}|\mathsf{N}_y$ by $c_\chi\bar{\alpha}\min\{1,\mathsf{f}(y)^{1/2}\}$. At $y$ in the
collar or ring of $C$ the collar case applies. Those cells have levels in an interval $[h,k]$, and
here the level is $k$ by part (i) of the geometry of exposed units. The collar case bounds the
increment by
\[
2(1+b_+)^{1/2}K_D\,c_\chi\,\varphi_{\flat}^{-1}\,
\mathsf{w}_L\mathsf{U}_J/\alpha_{*,k}.
\]
With $\delta'_k=g_kA_1\mathsf{T}_k^{p_1}$, $\theta'(g_k)=C_{\theta'}\mathsf{T}_k^{p_1-q}$ and
$\alpha_{*,k}=C_*g_k\mathsf{T}_k^{q}$, where $p_1$ and $q$ denote the exponents of $\mathsf{T}_k$ in
$\delta'_k$ and in $\alpha_{*,k}$, the ratio is free of the coupling:
\[
\frac{\mathsf{w}_L\mathsf{U}_J}{\alpha_{*,k}}
=\frac{6B_H\,g_kA_1\mathsf{T}_k^{p_1}}
{C_{\theta'}\mathsf{T}_k^{p_1-q}\,C_*g_k\mathsf{T}_k^{q}}
=\frac{6B_HA_1}{C_{\theta'}C_*}.
\]
The increment is therefore at most $1/16$ under the floor
\[
C_{\theta'}\ge 192\,(1+b_+)^{1/2}K_D\,c_\chi\,\varphi_{\flat}^{-1}B_HA_1/C_* ,
\]
since $2\cdot 6\cdot 16=192$. At $y\in\Omega_k\cap\mathfrak{T}$, with shell index $m=k$, the case
of locations in $\mathfrak{T}$ applies in the same way.

For a piece of class (d) we have $\mathsf{U}=\mathsf{U}_L$, the total majorant of the four
brackets. The quotient $\mathsf{w}_L\mathsf{U}_L/\alpha_*$ is a power of the logarithm, and it is
not what is spent. Instead, the line $c_\chi\mathsf{w}_L\mathsf{U}_LF_s(y)^{1/2}$ is the one line
of the case for locations in $\mathfrak{T}$, which bounds the profile $|\mathsf{t}|\,\xi_p(y)$ of
the dilated piece by $\mathsf{w}_L$ times the majorant of the layer times its far factor at half
exponent, $e^{-\frac12\delta_1^{\mathrm{sc}}d_{\mathrm{sc}}(y,\mathfrak{F}_{\mathfrak{s}})}$,
with the $\square$-piece in place of the layer. Indeed, at a location in $\mathfrak{T}$ the ratio
$F_s$ is the majorant of the layer-pricing proposition divided by $\mathsf{U}$. So the conclusion
of that case holds word for word at every location $y\in S_v$ in $\mathfrak{T}$, whatever the
shell index $m(y)\le k$:
\[
\frac{\text{increment}}{\text{threshold}}
\le 2\,\mathsf{w}_L\,c_{\mathrm{rec}}C_P'\,
e^{-\frac32\delta_1^{\mathrm{sc}}R_k}\,2^{-(k-1-m(y))}
\le\tfrac14\,\mathrm{room}^{\flat}(y).
\]
This is taken under the supremum ceiling of that case with $C_P'$ replaced by
$c_{\mathrm{rec}}C_P'$, and per crossed stratum on thin strata, where $\delta_1M/M_1\ge 2$. On and
beside a
foreign live piece the foreign-arrest case holds word for word. Its halved far factors $e^{-1}$ per
stratum are exactly the per-stratum reading of $F_s^{1/2}$, and its rooms are those tabulated for
the live arrest, under that case's ceiling.

Entries $2$ and $4$ obey $Q_e^+\le\frac{17}{16}Q_e+\theta_e/16$, where $Q_e^+$ is the dilated
value of entry $e$. This follows from the linearisation estimate and the summation estimate of the
rider at $b=0$, under its half-exponent ceiling with its constant multiplied by $e^{2d+2}$, at
touching blocks only.

At a receding block $\operatorname{Sten}_2(S_v)\cap\square^{\zeta}=\emptyset$, so
$\hat{\chi}_{\square}\equiv 0$ on every stencil of $S_v$. The entries of the carrier's
$U$-slot on $S_v$ are then those of the undilated chain, which are bounded by the current-entry
bound. Entry $3$ moves, and its motion is controlled by (c2)$(\alpha)$ below.

\emph{Clause (c2)$(\alpha)$.} This is the receding-block current lemma, clauses (a) and (b),
together with its corollary for an arbitrary positional weight, read at the weight $\eta_p$ of
\eqref{5.2:eq-carrier-exterior-a}. That corollary asks of the weight $\eta_p\in\,]0,1]$ only that
$\eta_p\ge e^{-2\delta_P d(\square^{\zeta},S_v)}$. This holds: at a receding block
$\eta_p=e^{-\delta_P d(\square^{\zeta},S_v)}$, and at a touching block $\eta_p=1$.

By clause (a) of the receding-block current lemma, at a receding block the $\mathbf{J}$-slot of the
dilated carrier on $S_v$ is its value at $\mathsf{t}=0$ minus $\mathsf{t}T_p$, where
$T_p=P_1(\sigma_s)(\hat{\chi}_{\square}A'_s)$. At every cell $y$ and $x\in\Delta(y)$ one has
\begin{align*}
\ell(y)^3|T_p(x)|
&\le 20K_PB_\sharp\,p_s(y)\,e^{-\delta_Pd(y,\square^{\zeta})}\\
&\le 20K_PB_\sharp\,\mathsf{a}_s\,
e^{-\delta_Pd^{\wedge}(y,\mathfrak{F})}\,e^{-\delta_Pd(y,\square^{\zeta})} ,
\end{align*}
where $K_P$, $B_\sharp$, and $K_\chi$ below, are the constants of the receding-block current
lemma, and $p_s(y)=\mathsf{a}_se^{-\delta_Pd^{\wedge}(y,A_s)}$ with $A_s$ the set of cells carrying
the datum of bracket $s$. The second inequality is the source-location bound
$p_s\le\mathsf{a}_se^{-\delta_Pd^{\wedge}(\cdot,\mathfrak{F})}$, which holds because $A_s$ lies in
$\mathfrak{F}$.

Let the block recede and $y\in S_v$. Then $d(\square^{\zeta},S_v)\le d(y,\square^{\zeta})$, and
$D_{\mathfrak{F}}:=d^{\wedge}(S_v,\mathfrak{F})\le d^{\wedge}(y,\mathfrak{F})$. Since
$\mathsf{f}_v^{\wedge}=e^{-\delta_PD_{\mathfrak{F}}}\le\mathsf{f}_v$,
\[
\rho_A=\mathsf{w}_L\,e^{\frac12\delta_Pd(\square^{\zeta},S_v)}\,\mathsf{f}_v^{-1/2}
\le\mathsf{w}_L\,e^{\frac12\delta_Pd(\square^{\zeta},S_v)}\,e^{\frac12\delta_PD_{\mathfrak{F}}} .
\]
Multiplying the bound just displayed by $\rho_A$ gives
\begin{align*}
\rho_A\,\ell(y)^3|T_p(x)|
&\le 20K_PB_\sharp\,\mathsf{a}_s\mathsf{w}_L\,
 e^{\delta_P[\frac12d(\square^{\zeta},S_v)-d(y,\square^{\zeta})]}\,
 e^{\delta_P[\frac12D_{\mathfrak{F}}-d^{\wedge}(y,\mathfrak{F})]}\\
&\le 20K_PB_\sharp\,\mathsf{a}_s\mathsf{w}_L\,\eta_p^{1/2}\,
 e^{-\frac12\delta_Pd^{\wedge}(y,\mathfrak{F})}.
\end{align*}
In the second line we used
\[
\tfrac12d(\square^{\zeta},S_v)-d(y,\square^{\zeta})\le-\tfrac12d(y,\square^{\zeta})
\le-\tfrac12d(\square^{\zeta},S_v),
\]
whose exponential is $\eta_p^{1/2}$, and
$\frac12D_{\mathfrak{F}}-d^{\wedge}(y,\mathfrak{F})\le-\frac12d^{\wedge}(y,\mathfrak{F})$.

At a touching block $\rho_A=\mathsf{w}_L\mathsf{f}_v^{-1/2}$. The same display holds there, with
$e^{-\delta_Pd(y,\square^{\zeta})}\le 1=\eta_p^{1/2}$. To pass to a clause index $n\le\pi(y)$ we
use $\ell_n^3=L^{-3(\pi(y)-n)}\ell(y)^3$. Reading in the presented pair costs a factor at most
$1.05$, and taking for a cell $y$ of $S_v$ a unit block meeting $S_v$ costs a factor at most $2$.
These factors are absorbed in $K_{\mathsf{J}}=42\max\{K_PB_\sharp,K_\chi\}$, and this gives the
first inequality of the two-sided display of clause (c2).

The second inequality is clause (b) of the receding-block current lemma. Its level clause and its
age bound give, with $b:=\pi(y)-n$, $u:=u(y)$ and $C^{\wedge}_{\mathrm{age}}=2\sqrt2(1+b_+)^{1/2}$,
\[
\alpha_{*,k}/\alpha_{*,n}\le C^{\wedge}_{\mathrm{age}}(1+u)^{1/2}(1+b)^{1/2},
\qquad (1+b)^{1/2}L^{-3b}\le 1 .
\]
The factor $e^{-\frac12\delta_Pd^{\wedge}(y,\mathfrak{F})}=e^{-\frac12\delta_P(M/M_1)u}$ is split as
\[
e^{-\frac12\delta_P(M/M_1)u}
=e^{-\frac18\delta_P(M/M_1)u}\,e^{-\frac18\delta_P(M/M_1)u}\,e^{-\frac14\delta_P(M/M_1)u}.
\]
Lemma~\ref{5.2:lem-carrier-exterior} assumes the floor
$\frac1{16}\delta_PM/M_1\ge c_{\boxplus}\log\hat{b}_t+1\ge 1$, so
$\delta_PM/M_1\ge 16$. The first factor is then at most $e^{-2u}$, which absorbs $(1+u)^{1/2}$ since
$(1+u)^{1/2}e^{-2u}\le 1$. The second is at most $e^{-2u}\le 2^{-u}$.

For $s=4$ one has
$\mathsf{w}_L\mathsf{a}_4/\alpha_{*,k}\le 6A_1/(C_{\theta'}C_*)$, free of the coupling, and the floor
\[
C_{\theta'}\ge 96\,K_{\mathsf{J}}C^{\wedge}_{\mathrm{age}}A_1/
(\beta_{\mathrm{pr}}\theta_S\varphi_{\flat}C_*)
\]
assumed in Lemma~\ref{5.2:lem-carrier-exterior} closes the inequality. For $s\le 3$ (class (d))
the third factor gives
$e^{-\frac14\delta_Pd^{\wedge}(y,\Lambda^{\sim})}\le e^{-\check{R}_k}\le E(x_k)$. For $y\notin Z$
this is the level clause. For $y\in S_v\cap Z$ it is the separation
$\operatorname{dist}_k(\Lambda^{\sim},Z^c)\ge 9\check{R}_k$ of the hubs from $Z^c$: by the star
lemma $S_v$ lies within four level-$k$ units of the cell $\Delta_{y_v}$, hence within five
level-$k$ units of the point $y_v\in Z^c$, so every cell of $S_v\cap Z$ is at
distance at least $9\check{R}_k-5$ from $\Lambda^{\sim}$. With $\mathsf{a}_s\le\mathsf{a}^{\flat}$,
the ceiling
\[
\sup_{x\ge\gamma^{-2}}16K_{\mathsf{J}}C^{\wedge}_{\mathrm{age}}\mathsf{a}^{\flat}
(\beta_{\mathrm{pr}}\theta_S\varphi_{\flat}C_{\theta'})^{-1}(\log x+2)^{q-p_1}
\alpha_*(x)^{-1}E(x)\le 1
\]
assumed in Lemma~\ref{5.2:lem-carrier-exterior} closes it. Thus entry $3$ moves by at most
$\varepsilon_{\mathsf{J}}=\beta_{\mathrm{pr}}\theta_S/16\le 1/16$ of every threshold, and at
$y\notin Z$ by at most $\frac14\mathrm{room}^{\flat}(y)$. By clause (a) of the same lemma no other
entry moves.

\emph{Clause (c2)$(\beta)$.} Let the block touch, so that $\eta_p=1$ and
$|\mathsf{t}|\le\mathsf{w}_L\mathsf{f}_v^{-1/2}$. We use the pointwise corollary of the
receding-block current lemma, valid at every block and every $y$. Wherever
$10|\mathsf{t}|B_\sharp p_s+2B_\sharp p_s\le 1/10$ on the enlarged cell $\tilde{\Delta}(y)$, it
bounds the $\mathsf{t}$-member of entry $3$ by $20K_\chi|\mathsf{t}|q_{\square}(y)$, with
$q_{\square}\le p_s$. As in clause (c2)$(\alpha)$,
$p_s(y)\le\mathsf{a}_se^{-\delta_Pd^{\wedge}(y,\mathfrak{F})}$ since $A_s\subset\mathfrak{F}$.

For $y\in S_v$ one has $d^{\wedge}(y,\mathfrak{F})\ge d^{\wedge}(S_v,\mathfrak{F})$, hence
$e^{-\frac12\delta_Pd^{\wedge}(y,\mathfrak{F})}\le(\mathsf{f}_v^{\wedge})^{1/2}\le
\mathsf{f}_v^{1/2}$. Therefore
\begin{align*}
|\mathsf{t}|\,p_s(y)
&\le\mathsf{w}_L\,\mathsf{f}_v^{-1/2}\,\mathsf{a}_s\,e^{-\delta_Pd^{\wedge}(y,\mathfrak{F})}\\
&=\mathsf{w}_L\,\mathsf{a}_s\,e^{-\frac12\delta_Pd^{\wedge}(y,\mathfrak{F})}\cdot
 \bigl[e^{-\frac12\delta_Pd^{\wedge}(y,\mathfrak{F})}\mathsf{f}_v^{-1/2}\bigr]
\le\mathsf{w}_L\,\mathsf{a}_s\,e^{-\frac12\delta_Pd^{\wedge}(y,\mathfrak{F})}.
\end{align*}
This is the first inequality of the display of clause (c2) at $\eta_p=1$, with $20K_\chi$ in place
of $20K_PB_\sharp$; it is covered by $K_{\mathsf{J}}\ge 42K_\chi$. The second inequality of that
display, clause (b) of the receding-block current lemma, reads only $(y,n,s)$ and not the block.
So it gives the bound $\varepsilon_{\mathsf{J}}2^{-u(y)}\varphi_{\flat}\alpha_{*,n}$ here as well.

The hypothesis of the pointwise corollary holds by the same display:
$10B_\sharp|\mathsf{t}|p_s\le 10B_\sharp\varepsilon_{\mathsf{J}}\varphi_{\flat}\alpha_*/K_{\mathsf{J}}
\le\alpha_*$, since $K_{\mathsf{J}}\ge 42B_\sharp$. The weight $p_s$ varies by a factor at most
$e^{2\delta_P}\le 2$ on $\tilde{\Delta}(y)$. The bound $\alpha_*\le 1/400$ follows a fortiori from
the regularity ceiling assumed in Lemma~\ref{5.2:lem-carrier-exterior}. The remaining floor on
$\gamma$ among the hypotheses of the receding-block current lemma is assumed as well.

No factor is spent twice. Of $e^{-\delta_Pd^{\wedge}(y,\mathfrak{F})}$, one half is taken by the
disc $|\mathsf{t}|\le\rho_A$. The other half is spent on containment: $\frac18$ on the age
polynomial, $\frac18$ on $2^{-u}$ and $\frac14$ on $E$. The Cauchy gain is the disc's own half.
Clause (c4) and the summation lemma for the pointed pieces spend only that gain.
\end{proof}

\begin{proof}[Proof of Lemma~\ref{5.2:lem-carrier-exterior}: long units and far cores]
\label{5.2:prf-carrier-proof-far}
This part of the proof establishes clauses (c3), (c4), (c5) and (d) of
Lemma~\ref{5.2:lem-carrier-exterior}. Clauses (a), (b), (c1) and (c2) are proved in the
first part of this proof, and they are used below in the form proved there. Throughout, $p$
is a pointed piece of the exterior classes, $v$ its unit, $y_v$ the point of $v$,
$\Delta_{y_v}$ the cell of $y_v$, $n=\pi(y_v)$ in the current global label, and $\rho_A$,
$\eta_p$, $\mathsf f_v$, $S_v$ are as defined in \eqref{5.2:eq-carrier-exterior-a}; the letters
$y$, $y'$ denote running cells of $S_v$.

\smallskip\noindent
\emph{Clause (c3): size and Cauchy estimate.} Let $p$ be compact, so that $S_v=\mathfrak r_v$.
By (c1) and (c2), at every cell of $S_v$ and at every block $\square$, each entry of the
dilated carrier that moves at all moves by at most $\theta_a/16$, one sixteenth of its
threshold $\theta_a$. For the
entries of the U-slot this is (c1): they move through the layer datum, which by
\eqref{5.2:eq-carrier-exterior-b} is at most
$c_\chi\mathsf w_L\mathsf U\min\{1,\mathsf f(y)^{1/2}\}$, and through the containment
clauses of the off-diagonal containment statement for the collar and ring, for
$\mathfrak T$-locations and for live arrests, applied to
that datum. For entry~3 it is (c2); for the present clause a gain of at most $\theta_3/16$ is
all that is needed of (c2). The exterior current estimate gives, for the undilated carrier,
every entry below $\theta_a$ on $S_v$, with
$\theta_a(y')\le 3c_\alpha w(\Delta_{y'})\mathbf r_a(\pi(y'))$.
Adding the increments allowed by (c1) and by (c2), every entry $Q_a$ of the dilated carrier at
a cell $y'$ of $S_v$ satisfies, on the whole of $\mathfrak P\times\{|\mathsf t|\le\rho_A\}$,
\begin{equation*}
Q_a(y')<\theta_a(y')+\tfrac{1}{16}\theta_a(y')+\tfrac{1}{16}\theta_a(y')
\le\tfrac98\,\theta_a(y')\le\tfrac98\cdot 3c_\alpha w(\Delta_{y'})\,\mathbf r_a(\pi(y')) .
\end{equation*}
Hence the exterior regularity estimate, which on the box $S\subset\mathfrak r_v$ at the zone
index $n$ requires the entries to lie below $\theta_a$ and returns the radius
$\mathfrak a'_v=c_{\mathrm{reg}}\bar{\mathsf w}_v\alpha_{*,n}$, runs at the radius
$\tfrac98\mathfrak a'_v$, provided its hypothesis
\begin{equation*}
60\,dLM\cdot\tfrac98\,\mathfrak a'_v\le 1
\end{equation*}
holds. This is supplied by \eqref{5.2:eq-carrier-exterior-c}. Indeed, the regularity
condition on the auxiliary parameters, of which \eqref{5.2:eq-carrier-exterior-c} is the
dilated form,
yields $60dLM\mathfrak a'_v\le 1$ through the ceiling line of the regularity estimate. That
line runs as follows. Either $\bar{\mathsf w}_v=3$, or the weighted cell lies in the window of
an attempt with $n\le k_a$. In the second case
$\bar{\mathsf w}_v\le 3L\mathsf R^{\uparrow}(\check{\mathsf T}_{k_a})$ with
$\check{\mathsf T}_{k_a}\le\log x_{k_a}$, and $\alpha_{*,n}\le 2\alpha_*(x_{k_a})$, so that the
supremum over $x\ge\gamma^{-2}$ covers $x=x_{k_a}$ whatever the age of the attempt. The
constants in that line multiply to $60\cdot 3\cdot 2=360$ in front of
$dL^2Mc_{\mathrm{reg}}\mathsf R^{\uparrow}(\log x)\alpha_*(x)$. With the extra factor $9/8$
the constant becomes $\tfrac98\cdot 360=405$, which is the constant of
\eqref{5.2:eq-carrier-exterior-c}.

On the resulting box the size estimate for singles, the Wilson-line estimate (entry~1 read
pointwise) and, for the far cores of class (d), the core estimate read as in clause~(d) below,
apply at $\tfrac98\mathfrak a'_v$. The size table of the pointed units is quadratic in
$\mathfrak a'$, and the core's depth parameter $s$ is read at $s^{(9/8)}$. Therefore, on the
joint polydisc,
\begin{equation}\label{5.2:eq-carrier-proof-far-1}
\bigl|f_v\circ\mathsf R_p(\mathsf t,\cdot)-f_v(1)\bigr|
\le\bigl(\tfrac98\bigr)^2\mathfrak N^{\flat}(v)=\tfrac{81}{64}\,\mathfrak N^{\flat}(v).
\end{equation}

The carrier is jointly holomorphic in $(\zeta,\sigma,\mathsf t)$, by the analyticity of the
source map at the localisation site, read in the product chart of the decoupling parameters.
Hence $\mathsf t\mapsto f_v\circ\mathsf R_p(\mathsf t,\cdot)-f_v(1)$ is holomorphic on
$|\mathsf t|\le\rho_A$ and bounded there by $\tfrac{81}{64}\mathfrak N^{\flat}(v)$, by
\eqref{5.2:eq-carrier-proof-far-1}. The derivative in $\mathsf t$ does not see the constant
$f_v(1)$, and Cauchy's inequality at radius $\rho_A=\mathsf w_L/(\eta_p\mathsf f_v)^{1/2}$ gives
\begin{equation*}
\Bigl|\partial_{\mathsf t}\bigl(f_v\circ\mathsf R_p(\mathsf t,\cdot)-f_v(1)\bigr)
\big|_{\mathsf t=0}\Bigr|
\le\frac{\tfrac{81}{64}\mathfrak N^{\flat}(v)}{\rho_A}
=\mathsf w_L^{-1}(\eta_p\mathsf f_v)^{1/2}\cdot\tfrac{81}{64}\,\mathfrak N^{\flat}(v).
\end{equation*}
From here the proof of the amplitude-disc estimate for localisation pieces runs as it does
for the plain boundary-type units: after the rescaling of $\mathsf t$ made there, the walk
factor $e^{-(\kappa_1-1)m(p)}$ and the factor $2$ are supplied by that proof itself. The
Cauchy estimate in the link parameters is then applied in each decoupling parameter
$\mathsf s(\Delta)$ of the cut families $\mathsf y_t$, on the circles
$|\sigma-\mathsf s(\Delta)|=e^{\kappa_1}-1\ge e^{\kappa_1-1}$. It multiplies the bound by
$e^{-(\kappa_1-1)\sum_t|\mathsf y_t|}$, and this is the display
\eqref{5.2:eq-carrier-exterior-d}.

The sums over $\square$, over the walks $\mathsf y$ and over the units $v$ are not taken here.
They are those of the sum of pointed pieces behind one coarse anchor,
Corollary~\ref{5.2:cor-anchor-sum}, carried out with the block-shell lemma, the source-rooted
sum and the rim-mass lemma, Lemma~\ref{5.2:lem-rim-mass}. The factor $\mathsf f_v^{1/2}$ does
not depend on $\square$, and it is carried unchanged into the per-cell bound
\eqref{5.2:eq-rim-mass-a}.

\smallskip\noindent
\emph{Clause (c4): proof of \eqref{5.2:eq-carrier-exterior-e}, and the constant
$\mathsf C_{\mathrm{fw}}$.} Let $\Delta'$ be the weighted cell of $p$, that is, the cell at
$\operatorname{dist}_n\le 1$ from $\Delta_{y_v}$ on which $w^{\vee}(\Delta_{y_v})>1$ is
attained. Suppose that $\Delta'\subset W_{\mathfrak a'}\subset X_{\mathfrak a'}$, the window of
a live arrest $\mathfrak a'$ on a component
$C'$ of $\Lambda_k^c$ that is not in $Z$. Here $k_{a'}$ denotes the step of $\mathfrak a'$,
$k_{0,a'}$ the initial index of its attempt, and
$N_0(a')=\log_L\bigl(L\check R_{k_{0,a'}+1}\bigr)$ the number of zones of that attempt, so that
$L^{N_0(a')}=L\check R_{k_{0,a'}+1}$. The arrest is foreign: were $C'$ in $Z$, then
$X_{\mathfrak a'}\subset Z_{k_{a'}}\cap C'$ would lie at distance at least $\check R_k$ inside
$C'\subset Z$. But $\Delta_{y_v}\subset Z^c\cup(\text{the collar of }C)$, so no cell of a unit
of class (c) or (d) could neighbour it.

By the definition of a live arrest, $X_{\mathfrak a'}$ is the arrest's component of the
large-field region $\Lambda^c_{k_{a'}}$ at its step. The small-field regions decrease,
$\Lambda_k\subset\Lambda_{k_{a'}}$, so $\Lambda^c_k\supseteq\Lambda^c_{k_{a'}}$. Moreover
$\Lambda^{\sim}\subset C$, where $C\neq C'$ are components of $\Lambda^c_k$. Hence
$\Lambda^{\sim}$ lies outside the component $X_{\mathfrak a'}$ of $\Lambda^c_{k_{a'}}$.

Let $y$ be a cell of $S_v$. By the first lines of the proof of
Lemma~\ref{5.2:lem-rim-mass}, $y$ lies within $4$ level-$n$ units of $\Delta'$, or, on a member
of level $n+1\le k_{a'}$, within $2$ level-$k_{a'}$ units. Here $n\le k_{a'}$: by the ceiling
line recalled in (c3), a weighted cell at $\operatorname{dist}_n\le1$ from $\Delta_{y_v}$ lies in
the window of an attempt of step $k_a$ with $n\le k_a$; applied to $\mathfrak a'$ this gives
$n\le k_{a'}$, which is also the first statement of the proof of Lemma~\ref{5.2:lem-rim-mass}.
Hence $y\in X_{\mathfrak a'}$, within $5$ level-$k_{a'}$ cubes of $W_{\mathfrak a'}$.

Every contour from $y$ to $\Lambda^{\sim}$ therefore crosses the level-$k_{a'}$ collar
$\Omega_{k_{a'}}\setminus\Lambda_{k_{a'}}$ about that component. For this collar Ba{\l}aban's
statement that ``the distance between their boundaries is at least equal to $2MR_j$''
\cite[p.~256]{B-CMP119} applies at $j=k_{a'}$: the regions were made by the operation
$\mathbf T$ at step $k_{a'}$ and have stayed unchanged since, by the freezing property of live
arrests \cite{B-CMP122a}. The collar is thus at least $2\check R_{k_{a'}}$ own-scale $M$-cubes
thick, and at most $6$ of them can already lie between $y$ and the inner side of the collar.
Therefore
\begin{equation*}
d_{\mathrm{sc}}(y,\Lambda^{\sim})\ge 2\check R_{k_{a'}}-6\ge\check R_{k_{a'}}-5
\qquad(\check R\ge 38).
\end{equation*}

The same bound holds in the currency of $d^{\wedge}$. By the definition of a live arrest, its
collar is $\mathfrak c_{\mathfrak a}:=X_{\mathfrak a}\cap\Omega_{k_a}$, of level $k_a$, made of
two layers of $MR_{k_a}$-cubes. By the contour lemma for live arrests, a contour from
$X_{\mathfrak a}^c$ to the piece moreover crosses $\mathfrak c_{\mathfrak a}$, and its
$d^{\wedge}$-length there is at least $2R_{k_a}\,M/M_1$. We have
$\Lambda^{\sim}\subset Z\subset X_{\mathfrak a'}^c$, since the arrest's component is untreated,
and $S_v$ lies within $5$ level-$k_{a'}$ units of the weighted cell. The passage from a point of
$W_{\mathfrak a'}$ to the infimum over $S_v$ uses $n\le k_{a'}$: the cells of $S_v$ are of level
at most $k_{a'}$, hence no coarser than the collar's cubes \cite{B-CMP122a}, so that at least
$2\check R_{k_{a'}}-5$ collar cubes remain to be crossed. This gives
\begin{equation*}
d^{\wedge}(S_v,\Lambda^{\sim})\ge\frac{M}{M_1}\bigl(2\check R_{k_{a'}}-5\bigr).
\end{equation*}

For class (d) the source set is $\mathfrak F_{\mathfrak s}=\Lambda^{\sim}$. The pieces of class
(c) never meet a foreign window, by part~(i) of the exterior geometry statement. Since
$\mathsf f(y)=e^{-\delta_{\natural}(d_{\mathrm{sc}}(y,\Lambda^{\sim})-1)_+}$ and
$d_{\mathrm{sc}}(y,\Lambda^{\sim})-1\ge\check R_{k_{a'}}-6$, we obtain
\begin{equation*}
\mathsf f(y)\le e^{-\delta_{\natural}(\check R_{k_{a'}}-6)_+}\quad\text{on }S_v,
\qquad
\mathsf f_v=\sup_{y\in S_v}\min\{1,c_\chi\mathsf f(y)\}
\le c_\chi e^{-\delta_{\natural}(\check R_{k_{a'}}-6)_+}.
\end{equation*}

For the weight, let $w_{\mathfrak a'}=\mathsf L_0^{2g}$ with $g\le N_0(a')$ be the frozen weight
of the arrest, and use $\mathsf L_0^2\le L/3$. Then
\begin{align*}
\bar{\mathsf w}_v=3w^{\vee}(\Delta_{y_v})=3w_{\mathfrak a'}(\Delta')
&\le 3\mathsf L_0^{2(g+1)}\le 3\mathsf L_0^2(L/3)^{N_0(a')}\le 3\mathsf L_0^2L^{N_0(a')}\\
&=3\mathsf L_0^2L\check R_{k_{0,a'}+1}\le 3\mathsf L_0^2LC_{\uparrow}\check R_{k_{a'}}.
\end{align*}
The last step is the window line of \eqref{5.2:eq-plain-fine-units-a} at the arrest's step:
\begin{equation*}
\check R_{k_{0,a'}+1}\le\mathsf R^{\uparrow}(\check{\mathsf T}_{k_{a'}})
\le C_{\uparrow}\check{\mathsf T}^{\,r_{\mathrm{pr}}}_{k_{a'}}\le C_{\uparrow}\check R_{k_{a'}}.
\end{equation*}
This proves \eqref{5.2:eq-carrier-exterior-e}.

Combining the two bounds,
\begin{equation*}
\mathsf f_v^{1/2}\bar{\mathsf w}_v^{\,3}
\le c_\chi^{1/2}\bigl(3\mathsf L_0^2LC_{\uparrow}\bigr)^3\,\check R_{k_{a'}}^{\,3}\,
e^{-\frac12\delta_{\natural}(\check R_{k_{a'}}-6)_+}.
\end{equation*}
Put $\varrho:=\check R_{k_{a'}}-5\ge 0$. Then $\check R_{k_{a'}}\le\varrho+6$ and
$(\check R_{k_{a'}}-6)_+=(\varrho-1)_+$. The right member is therefore at most
\begin{equation*}
c_\chi^{1/2}\bigl(3\mathsf L_0^2LC_{\uparrow}\bigr)^3
\sup_{\varrho\ge 0}(\varrho+6)^3e^{-\frac12\delta_{\natural}(\varrho-1)_+}
=\mathsf C_{\mathrm{fw}},
\end{equation*}
the constant of \eqref{5.2:eq-carrier-exterior-f}. This proves (c4).

\smallskip\noindent
\emph{Clause (c5): long singles and tails.} Let $p$ be long, $j$ the level of its unit $v$ and
$X:=X_v\in\mathbf D_j$, with
$S_v=\{\text{cells meeting }X_v\}$ and with $\eta_p$, $\mathsf f_v$, $\rho_A$ taken on this
$S_v$.

(i) At a touching block, the five containment clauses of the off-diagonal containment
statement are read within the scope that statement permits: at the unit's clause index
$n(y)=j$, with
$b(y)=\pi(y)-j\ge 0$. These clauses read the layer only through a cell majorant of decay rate
$\delta_P/2$, and otherwise only through their own geometry, namely the linearity of the
threshold increments and the row-sum estimates at index $b$. Neither depends on the choice of
read set. Since $\mathsf f_v$ is the supremum of $\min\{1,c_\chi\mathsf f\}$ over $S_v$ and
$c_\chi\ge1$, one has $F_s(y)\le\mathsf f(y)\le\mathsf f_v$ at every $y\in S_v$, so the
computation of (c1),
\begin{equation*}
|\mathsf t|\,\mathsf N_y(\delta\mathbb H_p)
\le\mathsf w_L\mathsf f_v^{-1/2}c_\chi\mathsf U F_s(y)
\le c_\chi\mathsf w_L\mathsf U\min\{1,\mathsf f(y)^{1/2}\},
\end{equation*}
holds at every $y\in S_v$, and the value $c_\chi\mathsf w_L\mathsf U\min\{1,\mathsf f(y)^{1/2}\}$
of that majorant is available on the larger read set $S_v$. At a receding block, the
containment is clause (c2), in its part
for receding blocks, which is stated at every index $n\le\pi(y)$ with the scale drop
$L^{-3(\pi(y)-n)}$.

(ii) The Cauchy estimate in $\mathsf t$ at radius $\rho_A$ runs as in (c3). Here $f_v$ is
holomorphic on $\mathfrak D_v$ and bounded there by $\mathfrak n_v$, by the definition of the
families of localised terms; here $\mathfrak D_v=U^c_j(X_v,\alpha_{\cdot,j})$ is the space of
\cite[(1.11)--(1.16)]{B-CMP109}. Hence
$|f_v\circ\mathsf R_p-f_v(1)|\le 2\mathfrak n_v$ on the disc, and the display of (c5)(ii)
follows as \eqref{5.2:eq-carrier-exterior-d} did, with $\mathfrak N^{\flat}(v)=2\mathfrak n_v$
and without the factor $81/64$.

(iii) For class (c) the count is the one given in the statement. The cells meeting
$X\subset\Lambda_j$ have level $\ge j$, each of them holds a cube of $X$, and
$X\in\mathbf D_j$ has at most $\mathsf c_d(d_j(X)+1)$ cubes, by the cube-count inequality
$N_j(X)\le\mathsf c_d(d_j(X)+1)$; hence
$\#\operatorname{St}_v\le\mathsf c_d(d_j(X)+1)$. For class (d) the long singles and tails are
counted and summed with the long units of the sum of boundary-type units in block bundles. That
summation uses the lemma bounding the mass of long units that meet one cell, its corollary for
class (d), and the counts of the exempt family of a long unit. The long pieces of class (c) are
likewise summed in block bundles, by the class-(c) clause of that summation of boundary-type
units.

\smallskip\noindent
\emph{Clause (d): far cores.} A smooth far core of class (d) is read through
$\operatorname{ext}_v$, the extension map of the pointed reading of smooth cores, with
$p_v=n-2-s^{(9/8)}(\bar{\mathsf w}_v)$. We follow the smooth-core line of the core estimate,
without its disc variable $\hat\lambda$. Let $\mathcal U(\mathsf t,m)$ denote the U-slot of the
dilated cut carrier at $(\mathsf t,m)$, and put
\begin{equation*}
G(\mathsf t,m):=\mathfrak i^{(j)}_{y_v}\Bigl(\operatorname{ext}_v\bigl(\mathcal U(\mathsf t,m)
\restriction_{Q_v}\bigr)\Bigr).
\end{equation*}
Three inputs are combined. The first is part (i) of the exterior regularity estimate on
$Q_v\subset\mathfrak r_v$ at the radius $\tfrac98\mathfrak a'_v$, whose hypothesis was verified
in (c3). The second is the regularity input of the smooth-core line. The third is the core
estimate with $\epsilon\le 1$; this bound holds because $p_v$ is defined through $s^{(9/8)}$,
that is, at the dilated weight $\tfrac98\bar{\mathsf w}_v$. Together they give
\begin{equation*}
|G|\le 2\epsilon^2\mathfrak N_{p_v}\quad\text{on the whole joint polydisc }
\mathfrak P\times\{|\mathsf t|\le\rho_A\}.
\end{equation*}
The real values are unchanged. At the real chain, $\operatorname{ext}_v$ coincides with the
chain modulo gauge on $[\square_0]_{p_v}$ with two layers of $\pi_{p_v}$-cubes, by the
real-chain clause of the core extension. Finally, the inequality
$(\tfrac98)^3=\tfrac{729}{512}\le\tfrac32$, used wherever the dilated sizes are summed, is the
arithmetic
\begin{equation*}
729\cdot 2=1458\le 1536=512\cdot 3. \qedhere
\end{equation*}
\end{proof}

\begin{lemma}[The column bound per cell for exposed units]\label{5.2:lem-exposed-column}
Let $\Delta'$ be a cell of the current global partition which satisfies one of the following:
\begin{itemize}
\item it lies off $C$ and off the other treated components, and its level is then the level of its zone;
\item it is a cell of the collar of $C$.
\end{itemize}
In both cases let $n'$ denote the level of $\Delta'$ (for a cell of the collar of $C$, $n'=k$).
Assume that $\Delta'$ contains the point $y_v$ of at least one exposed unit $v$ in the sense of Definition~\ref{5.2:def-exposed-compact-units}. Put $\bar{\mathsf{w}}(\Delta') := 3w^{\vee}(\Delta')$. Then for every $r \le \tfrac85\kappa$
\begin{equation}\label{5.2:eq-exposed-column-a}
\sum_{v\ \text{exposed}:\ y_v\in\Delta'} \mathfrak{N}^{\flat}(v)\,
e^{r\,d_k(K_v)}\,e^{\varsigma_1\operatorname{rad}(v)}
\;\le\; \mathsf{K}^{\flat}\,\bar{\mathsf{w}}(\Delta')^{3}.
\end{equation}
Here $\bar{\mathsf{w}}(\Delta') = 3$ unless $\Delta'$ neighbours the window of a foreign live arrest $\mathfrak{a}'$. In that case
\begin{equation*}
\bar{\mathsf{w}}(\Delta') \le \bar{\mathsf{w}}^{\mathrm{fw}}(k_{a'})
= 3\mathsf{L}_0^2 L C_{\uparrow}\check{R}_{k_{a'}},
\end{equation*}
as stated in \eqref{5.2:eq-carrier-exterior-e}. Moreover, in either case, the same sum with each term weighted by $\mathsf{f}_v^{1/2}$ (the factor carried by the bounds of the $\mathsf{t}$-pieces) satisfies
\begin{equation}\label{5.2:eq-exposed-column-1}
\sum_{v\ \text{exposed}:\ y_v\in\Delta'}
\mathsf{f}_v^{1/2}\,\mathfrak{N}^{\flat}(v)\,
e^{r\,d_k(K_v)}\,e^{\varsigma_1\operatorname{rad}(v)}
\;\le\; \mathsf{C}_{\mathrm{fw}}\,\mathsf{K}^{\flat}.
\end{equation}
\end{lemma}

The bound \eqref{5.2:eq-exposed-column-a} is a bound per cell, each cell taken in its own scale. It is not a bound per cube of $\pi_k$. The passage from cells to a single $\pi_k$-anchor is made by the source-rooted summation lemma below, never by counting level-$k$ cells.

\begin{proof}
We run the proof of the column bound per cell on the component $C$ at the zone of the cell $\Delta'$. That proof is local at the point $y_v$ and at the zone of $y_v$, so it applies word for word at the zone $n'$ of $\Delta'$.

Let $j \le n'$ be the level of a unit with point in $\Delta'$, and put $t := L^{j-n'}$. The cell $\Delta'$ holds $M^4t^{-4}$ points of the lattice $T^{(j)}$. The row sums around $y$ are bounded row by row, as follows.
\begin{itemize}
\item Row~$2$. Its factor is $\omega_{k,j} = L^{-(k-j)/2}$, and since $n' \le k$ we have $\omega_{k,j} \le L^{-(n'-j)/2}$. Hence the row-two sum is dominated, a fortiori, by the row-two scale series $\bar{S}_\omega$ of the column bound on $C$, taken at zone $n'$.
\item Row~$3$. It is bounded by the row-three estimate that the column bound on $C$ uses.
\item The boundary row. It is bounded by the near-point count $27dM(\nu+2)t^{-3}$ of part~(iii) of the boundary-row lemma.
\item Row~$5$. Its bound does not depend on $K$.
\end{itemize}
For cores and Wilson pieces one has $d_k(K_v) \le 6^d$ and $\operatorname{rad}(v) \le 5$. As in the proof on $C$, every row keeps a decay rate of at least $(\tfrac45 - \tfrac12\delta)\kappa$.

Consider next the far cores of the units whose localisation domain is disjoint from $Z$. For such a unit the zone $\pi(y_v)$ of its point is the level $n'$ of $\Delta'$, and, with $p_v$ and the function $s(\cdot)$ as fixed with the letters of the exposed units (Definition~\ref{5.2:def-exposed-compact-units}) and $s^{(9/8)}(\bar{\mathsf{w}}) := s(\tfrac98\bar{\mathsf{w}})$, one always has
\begin{equation*}
p_v = n' - 2 - s^{(9/8)}(\bar{\mathsf{w}}_v),
\end{equation*}
and there is no clause requiring the box $Q_v$ of the core to lie in $Z$. Hence the rim case of the proof on $C$ does not occur. The smooth-core series
\begin{equation*}
\sum_{b}(b+1)^3\,\bar{\rho}_b^{\,2}\,(1+\bar{\rho}_b)\,L^{-\hat{\alpha} b}
\end{equation*}
and the bound for rough cores are taken as there. The polynomial in $s^{(9/8)}$ that enters these bounds is controlled by
\begin{equation*}
s^{(9/8)} \le s(1) + \log_L \bar{\mathsf{w}} + 1,
\end{equation*}
so it is at most a constant multiple of $\bar{\mathsf{w}}$. The smooth-core entry of the size table carries the factor $\bar{\mathsf{w}}_v^{2}$; multiplied by this polynomial, bounded by $c\,\bar{\mathsf{w}}_v$, it gives the third power of $\bar{\mathsf{w}}(\Delta')$ in \eqref{5.2:eq-exposed-column-a}.

Every sum over scales is a sum over the levels $j \le n'$ at fixed $n'$. Hence $\mathsf{K}^{\flat}$ depends neither on $K$ nor on $k$. This proves \eqref{5.2:eq-exposed-column-a}.

It remains to prove the weighted clause \eqref{5.2:eq-exposed-column-1}. For every exposed $v$ with $y_v \in \Delta'$ the cell $\Delta_{y_v}$ of the point $y_v$ is $\Delta'$, so $\bar{\mathsf{w}}_v = 3w^{\vee}(\Delta_{y_v}) = \bar{\mathsf{w}}(\Delta')$. All terms are non-negative. Hence the left member of \eqref{5.2:eq-exposed-column-1} is at most the supremum of $\mathsf{f}_v^{1/2}$ over the units it contains, times the left member of \eqref{5.2:eq-exposed-column-a}. By \eqref{5.2:eq-exposed-column-a} it is therefore bounded by
\begin{equation*}
\sup_{v}\,\bigl(\mathsf{f}_v^{1/2}\,\bar{\mathsf{w}}(\Delta')^3\bigr)\,\mathsf{K}^{\flat}
= \sup_{v}\,\bigl(\mathsf{f}_v^{1/2}\,\bar{\mathsf{w}}_v^{3}\bigr)\,\mathsf{K}^{\flat}.
\end{equation*}
We now bound the supremum in the two cases of the statement.
\begin{itemize}
\item Suppose $\Delta'$ neighbours the window of a foreign live arrest, where $w^{\vee}(\Delta')$ is attained. Then the carrier bound for exterior pointed units (Lemma~\ref{5.2:lem-carrier-exterior}), in the form \eqref{5.2:eq-carrier-exterior-f}, gives $\mathsf{f}_v^{1/2}\bar{\mathsf{w}}_v^{3} \le \mathsf{C}_{\mathrm{fw}}$ for each such $v$.
\item Off the foreign windows, $\bar{\mathsf{w}}_v = 3$ and $\mathsf{f}_v \le 1$. Since $c_\chi \ge 1$ this gives
\begin{equation*}
\mathsf{f}_v^{1/2}\,\bar{\mathsf{w}}_v^{3} = 27\,\mathsf{f}_v^{1/2} \le 27 \le 27\,c_\chi^{1/2} \le \mathsf{C}_{\mathrm{fw}}.
\end{equation*}
For the last inequality, evaluating the supremum in the definition \eqref{5.2:eq-carrier-exterior-f} of $\mathsf{C}_{\mathrm{fw}}$ at $\varrho = 0$ gives
\begin{equation*}
\mathsf{C}_{\mathrm{fw}} \ge c_\chi^{1/2}\,(3\mathsf{L}_0^2 L C_{\uparrow})^3\,6^3 ,
\end{equation*}
and $3\mathsf{L}_0^2 L C_{\uparrow} \ge 3$, so that $\mathsf{C}_{\mathrm{fw}} \ge 27\cdot 6^3\,c_\chi^{1/2} \ge 27\,c_\chi^{1/2}$.
\end{itemize}
In both cases the supremum is at most $\mathsf{C}_{\mathrm{fw}}$, which proves \eqref{5.2:eq-exposed-column-1}.
\end{proof}

\begin{definition}[Notation for the source-rooted sum over cut cells]
\label{5.2:not-rooted-sum-notation}
The notation below serves to sum the compact pointed pieces of
Definition~\ref{5.2:def-exposed-compact-units} whose underlying plain unit $v$ either has its
domain meeting $Z$ with $K_v\not\subset Z$ (we call these the units of class (c)) or
satisfies $X_v\cap Z=\emptyset$ (the units of class (d)). The
pieces of every zone $n'\le k$ are summed behind one
$\pi_k$-anchor. The sum is organised as follows. The fine cells lying under a coarse cube are
never counted. A piece that does not vanish is tied to a hub $\Lambda^{\sim}_C$ (defined
below), a never-cut block, by
a connected family of cells. Each cell of this family is taken in its own scale and costs a
factor $e^{-(\kappa_1-1)}$. The sum is rooted at the source:
\begin{itemize}
\item the anchor is traded for a source cell, at the price $S_d$;
\item the connected family of own-scale cells through that source cell is counted, $(e\hat b)^n$
families of $n$ members against the weight $e^{-\lambda_\Sigma n}$;
\item a unit is then one whose exempt family has a given cell on its rim; since the exempt
family of a compact unit depends only on its anchor cell $\Delta'$, at most $N_\partial$ anchor
cells occur for a given rim cell;
\item the mass of the units anchored at one cell is bounded by the column bound per cell
(Lemma~\ref{5.2:lem-exposed-column}) together with the geometric clause
\eqref{5.2:eq-carrier-exterior-e};
\item the blocks $\square$ are summed, in shells, through the factor $\eta_p^{1/2}$.
\end{itemize}

\medskip\noindent\emph{Site geometry.}
We work at the localisation site of $\mathbf{R}_k$, in the current global zone labelling. The
following notions and facts are those of the geometry of the localisation site, whose
statements we take over:
\begin{itemize}
\item cells and their levels $\operatorname{lev}\Delta$, and zones;
\item the partitions $\pi_\ell$;
\item the blocks $\square$ of zone $n$ and their fattenings $\tilde\square$;
\item the distances $\operatorname{dist}_\ell$ and the counting functions $N_\ell$;
\item the $\pi_k$-hull $[\cdot]_k$, the class $\mathbf{D}_k$ of localisation domains and its
tree length $d_k$;
\item the relations ``meet'' (a common interior point) and ``touch''.
\end{itemize}
The admissibility conditions hold in every zone labelling that occurs, the frozen windows of
live arrests included. They say that wall- or corner-neighbouring cells differ in level by at
most one, and that
\begin{equation}\label{5.2:eq-rooted-sum-notation-1}
\operatorname{dist}_{j+1}(\Omega_j^c,\Omega_{j+1})\ge 1 .
\end{equation}
We write $\pi^{\mathfrak{s}}$ for the family of the cells disjoint with $\Lambda^{\sim}$. It is
one family for the four brackets $s=1,\dots,4$, because the determining sets
$\mathfrak{B}_s$ differ only inside $\Lambda$. Indeed, the determining set restricted to
$\Lambda^c$, and the field there, coincide with those of the new configuration, so the
expansions are connected with changes made inside $\Lambda$ \cite{B-CMP122b}.

The \emph{hubs} are the sets $\Lambda^{\sim}_C$, the part of $\Lambda^{\sim}$ lying in a
component $C$ of $Z$. They are boxes of $\pi_k$-cubes with
sides at most $h_\sharp:=112$, at mutual distances at least $19\check R_k$. Moreover
$\operatorname{dist}_k(\Lambda^{\sim},Z^c)\ge 9\check R_k$, and every cell disjoint with
$\Lambda$ and lying in $Z$ has level $k$; these facts of the hub geometry are taken over as
stated there.
The \emph{source cells} are the members of $\pi^{\mathfrak{s}}$ sharing a wall with a hub.
A source cell lies in $Z$, has level $k$ and is a $\pi_k$-cube.
Two members are \emph{adjacent} if they share a wall, or if both share a wall with one hub;
``connected'' and ``component'' refer to this adjacency. A member has at most
\[
\hat b:=2dL^{d-1}+2d(h_\sharp+2)^{d-1}
\]
adjacent members. For a member $\Delta$, $\rho(\Delta)$ denotes the least cardinality of a
connected subfamily holding $\Delta$ and a source cell. The site geometry gives
\begin{equation}\label{5.2:eq-rooted-sum-notation-2}
\rho(\Delta)\ge k-\operatorname{lev}\Delta+1,\qquad
\rho(\Delta)\ge \operatorname{dist}_k\Bigl(\Delta,\ \bigcup\{\text{source cells}\}\Bigr).
\end{equation}
It also gives that the connected subfamilies with $n$ members through a given member number at
most $(e\hat b)^n$. Three further facts on tree lengths are used:
\begin{itemize}
\item if $X\subset Y$ in $\mathbf{D}_k$, then $d_k(Y)\le d_k(X)+N_k(Y\setminus X)$;
\item a tree graph intersecting two cubes $Q,Q'$ has length at least
$\operatorname{dist}(Q,Q')$;
\item cubes at distance $x$ are joined by a wall-chain with fewer than $d(x+1)$ intermediate
cubes, lying inside the smallest box holding both.
\end{itemize}

Recall that $\delta_P=\tfrac18\delta_0$. We write $d$ and $d^{\wedge}$ for the
contour distances of $\mathfrak{B}_s$; a contour distance is always written with its two
arguments, $d(\cdot,\cdot)$ or $d^{\wedge}(\cdot,\cdot)$, so that no confusion with the
dimension $d$ arises. The contours $\Gamma$ and their lengths
$\operatorname{len}_{\mathfrak{B}}(\Gamma)$ are those of the geometry of $\mathfrak{B}_s$ from
which these distances are formed. The distances satisfy
$d^{\wedge}(\cdot,\cdot)\le d(\cdot,\cdot)$, and
$d^{\wedge}(y,y')\ge \tfrac{M}{M_1}(|i-i'|-1)_+$ for points $y,y'$ in cells of levels $i,i'$.
A cell touches
at most $\hat b_t:=(L+2)^d+3^d$ others. With $c_{\boxplus}:=3^d(L+5)$, a contour $\Gamma$ with
$\operatorname{len}_{\mathfrak{B}}(\Gamma)\le x\,M/M_1$ meets at most $c_{\boxplus}(x+1)$ cells,
and these cells are connected by touching. The further constants are:
\begin{itemize}
\item $S_d$, the constant of the site geometry for which, for every $Q_*\in\pi_k$,
\[
\sum_{Q\in\pi_k}e^{-\operatorname{dist}_k(Q_*,Q)}\le S_d\le\sum_{x\ge0}(2x+3)^de^{-x};
\]
\item $c_l:=1+234d$ and $c_l^{\circ}:=1+d(h_\sharp+1)\le c_l$;
\item $\lambda_*(d,L)$, the constant of the site geometry with
$L^de^{-\frac12(\lambda_*(d,L)-\log(e\hat b))}\le\tfrac12$; its argument $(d,L)$ is always
written, to distinguish it from the constant $\lambda_*$ of the test-function norm;
\item $a:=\tfrac{31}{20}\kappa$ and $c_{\mathfrak{h}}=7^d$.
\end{itemize}

\medskip\noindent\emph{Conditions on the parameters and derived constants.}
The source-rooted sum uses no conditions on the parameters beyond the following three lower
bounds, which are among the conditions fixed with the order of choice of the constants
(Definition~\ref{5.1:def-admitted-inputs}); the radius condition enters only through the hub
geometry.
In the same display we fix the derived constants, for $0\le r\le a+9$:
\begin{equation}\label{5.2:eq-rooted-sum-notation-a}
\begin{gathered}
\kappa_1-1\ge c_l(2\kappa+4)+\lambda_*(d,L),\qquad
\tfrac1{16}\delta_P\tfrac{M}{M_1}\ge c_{\boxplus}\log\hat b_t+1,\qquad \kappa\ge180,\\
\mathsf{q}_{\square}:=\hat b_t^{\,c_{\boxplus}}e^{-\frac12\delta_PM/M_1}\le e^{-8},\\
\lambda_\Sigma:=\kappa_1-2-c_l^{\circ}(r+1),\qquad
\lambda_\Sigma':=\lambda_\Sigma-\log(e\hat b),\qquad
\epsilon_\partial:=2e^{-\lambda_\Sigma'},\\
c_S:=(5L)^d,\qquad N_\partial:=5(L+2)^{2d}.
\end{gathered}
\end{equation}
Then, for $0\le r\le a+9$,
\begin{equation}\label{5.2:eq-rooted-sum-notation-3}
\lambda_\Sigma'\ge(\lambda_*(d,L)-\log(e\hat b))+75c_l-1,
\end{equation}
so that $\lambda_\Sigma'\ge3$ and $e^{-\lambda_\Sigma'}\le\tfrac14L^{-2d}e^{-(75c_l-1)}$.

\medskip\noindent\emph{Pieces; the cut cells are cells in proper scales.}
Let $v$ be a plain unit of class (c) or of class (d):
\begin{itemize}
\item class (c), the units whose domain meets $Z$ with $K_v\not\subset Z$, is read in one
bracket, $s=4$;
\item class (d) has $X_v\cap Z=\emptyset$ and $y_v$ in any zone $n'\le k$, and is read in the
brackets $s=4,\dots,1$.
\end{itemize}
Put $\Delta':=\Delta_{y_v}$ and $n':=\operatorname{lev}\Delta'$. The \emph{exempt family}
$\operatorname{St}_v$ of $v$ is a finite wall-connected family of members of
$\pi^{\mathfrak{s}}$ with $S_v\subset\bigcup\operatorname{St}_v=:\bar S_v$, where $S_v$ is the
read set of \eqref{5.2:eq-carrier-exterior-a}. In the amplitude-disc lemma at the localisation
site, the piece is
\begin{equation}\label{5.2:eq-rooted-sum-notation-4}
v_p=\prod_{\Delta\subset Y_0(p)\setminus\bar S_v}\int_0^1 d\mathsf{s}(\Delta)\,
\partial_{\mathsf{s}(\Delta)}\,\partial_{\mathsf{t}}\big|_0\,
f_v\bigl(\mathsf{R}_p(\mathsf{t},\mathsf{s})\bigr),
\qquad m(p):=\#\{\Delta\}.
\end{equation}
Here $Y_0(p)$ is defined below. The cells $\Delta$ occurring here are the members of
$\sigma_0(v):=\pi^{\mathfrak{s}}\setminus\operatorname{St}_v$, that is, cells in proper scales.
A piece is $p=(\square,\mathsf{y};s)$, where:
\begin{itemize}
\item $\mathsf{y}\subset\sigma_0(v)$ is finite; it is the \emph{walk}, the link-one family of
cut cells;
\item $m(p)=|\mathsf{y}|$;
\item $\square$ is a block of the cover of $\mathfrak{B}_s$.
\end{itemize}
The families $(\mathsf{y}_2,\mathsf{y}_3,\mathsf{y}_4)$ of links two to four do not occur in
the sum. We put
\[
Y_0(p):=\bar S_v\cup\bigcup\mathsf{y}\cup\{\text{the hubs sharing a wall with }\mathsf{y}\}.
\]
This agrees with Ba{\l}aban's construction. The expansions are built for families of cubes in
proper scales, a cube contained in $\Omega_j\setminus\Omega_{j+1}$ having side $ML^j$ lattice
spacings, and are then
resummed in localisation domains $Y_0\in\mathbf{D}_k$ \cite{B-CMP119}. At the localisation
site the expansion is taken with $\sigma_0$ the set of cubes disjoint with the cube
$\square^{\sim5}$ and $\Lambda^{\sim}$, and the sum runs over domains $Y_0\in\mathbf{D}_k$
containing $\square^{\sim5}$ and
at least one component of $\Lambda^{\sim}$; the other terms vanish by the localisation of the
field $B$ \cite[p.~372]{B-CMP122b}. Thus each exempt cell is taken at its own scale. This
choice costs
nothing analytically: the decoupling lemma for the site holds for an arbitrary subfamily
$\sigma_0$ with the same constants, and the polydisc of the decoupling parameters runs over
all cells. If the cut cells were instead taken to be $\pi_k$-cubes, with
$\bar S_v:=[S_v]_k$, the walk could not resolve the cells inside one $\pi_k$-cube.

\medskip\noindent\emph{The star.}
Recall from Definition~\ref{5.2:def-exposed-compact-units} that
$\mathfrak{r}_v:=\bigcup\operatorname{St}(\Delta_{y_v})$. Here
$\operatorname{St}(\Delta')$ is the wall-component of $\Delta'$ in the family of cells
$\Delta$ with $\operatorname{lev}\Delta\ge n'-1$ that meet
$Q_5(\Delta'):=\{x:\operatorname{dist}_{n'}(x,\Delta')\le2\}$, and
$n'=\operatorname{lev}\Delta'$. For a compact pointed unit we put
\[
S_v:=\bar S_v:=\mathfrak{r}_v=\bigcup\operatorname{St}(\Delta_{y_v}),\qquad
\operatorname{St}_v:=\operatorname{St}(\Delta_{y_v}).
\]
Hence $\operatorname{St}_v$, $S_v$, $\eta_p$, $\mathsf{f}_v$ and $\bar{\mathsf{w}}_v$ depend on
$v$ through $\Delta'$ alone. We write accordingly $\eta(\square,\Delta')$,
$\mathsf{f}(\Delta')$ and $\bar{\mathsf{w}}(\Delta')=3w^{\vee}(\Delta')$. A rough core is read
term by term: each constituent is a short single, which is compact, or a long one; at
$\nu\le2$ every constituent is long.

\smallskip\noindent\emph{What a core reads.}
The hull $X_v:=\square_0:=\square^{\sim\nu+1}$ of a core determines $K_v=[X_v]_k$ and hence
enters the label through $K_v\subset\mathfrak{h}_v$; it is not the region on which the unit
depends.
\begin{itemize}
\item A far core of class (d) depends on its constituents $X\ni y$ with
$X\subset\square^{\sim\nu}$, that is, on points within $\nu L^{-(\nu-1)}\le1$ level-$n'$ unit
of $\Delta_{y_v}$.
\item A Wilson piece depends on a block of $2^d$ cubes of $\pi_j$, $j\le n'$, holding $y$,
within one such unit.
\end{itemize}
Add the cube supplied by the cube clause of the carrier bound for exterior pointed units
(Lemma~\ref{5.2:lem-carrier-exterior}), of index $n'-1$ or $n'$, whose cells of that index lie
at $\operatorname{dist}_{n'}$-distance less than $1$ from $\Delta_{y_v}$, and use that
$Q_v\subset\Omega_{n'-1}$. Then everything read lies in
\[
\{\operatorname{dist}_{n'}(\cdot,\Delta_{y_v})\le1\}\subset\bigcup\operatorname{St}(\Delta_{y_v})
=\mathfrak{r}_v ,
\]
by the star lemma of this section. Consequently the gain of the bound
\eqref{5.2:eq-rooted-sum-notation-6} below is legitimate at a block that touches
$\mathfrak{r}_v$ away from the points actually read, since nothing is read there. Enlarging
$S_v$ beyond the set actually read only shrinks $\rho_A$. Moreover, the receding-block estimate
for the current holds for every positional factor $\eta_p\in\,]0,1]$ with
$\eta_p\ge e^{-2\delta_Pd(\square^{\zeta},S_v)}$. The hull $\square_0$ enters $K_v\subset\mathfrak{h}_v$
only. The choice $S_v:=\mathfrak{r}_v$ is a stated departure from the exempt cube
$\square^{\sim5}$ used in \cite[p.~372]{B-CMP122b}.

For any exempt family $\operatorname{St}$, its rim $\partial\operatorname{St}$ is the family of
members of $\pi^{\mathfrak{s}}\setminus\operatorname{St}$ sharing a wall with a member of
$\operatorname{St}$. The constants $c_S$ and $N_\partial$ are those of
\eqref{5.2:eq-rooted-sum-notation-a}.

\medskip\noindent\emph{Labels: the walk's hull with corridors.}
Let $A\in\mathbf{D}_k$ with $[\operatorname{St}_v]_k\subset A$, and let $\mathsf{y}$ be such
that $\operatorname{St}_v\cup\mathsf{y}$ is connected. Fix a spanning forest of $\mathsf{y}$.
For every edge of the forest joining two members that share no wall, take the intermediate
cubes of a wall-chain between the two, chosen as in the third tree-length fact above;
$\operatorname{Hub}(\mathsf{y})$ is the union of these intermediate cubes. We put
\begin{equation}\label{5.2:eq-rooted-sum-notation-5}
\operatorname{lab}(A,\mathsf{y}):=A\cup[\mathsf{y}]_k\cup\operatorname{Hub}(\mathsf{y}),\qquad
K_p:=[\operatorname{St}_v\cup\mathsf{y}]_k\cup\operatorname{Hub}(\mathsf{y}).
\end{equation}
Thus $K_p$ is the $\pi_k$-hull of the walk, with corridors across the never-cut blocks but
without their interiors. The label of a compact $p$ is
\[
\operatorname{lab}_1(p):=K_p\cup\mathfrak{h}_v=\operatorname{lab}(\mathfrak{h}_v,\mathsf{y}).
\]
For a long single or tail the exempt family $\operatorname{St}_v$ is the family of cells of
$\pi^{\mathfrak{s}}$ meeting $X_v$, and its label is $\operatorname{lab}_1(p):=K_p$. If $X_v$
is disjoint with $\Lambda^{\sim}$ (in particular for class (d), since
$\Lambda^{\sim}\subset Z$), then $[\operatorname{St}_v]_k=K_v$ by the counting
proposition for long units; hence in that case
$\operatorname{lab}_1(p)=K_p=\operatorname{lab}(K_v,\mathsf{y})$.

The hubs and $\operatorname{Hub}(\mathsf{y})$ lie in components of $Z$ which the source cells
of $\mathsf{y}$ already meet, so $K_p$ and $[Y_0(p)]_k$ agree off $Z$. Hence leaving the
interiors of the hubs out of $K_p$ does not change the domain $Y(\operatorname{lab}_1)$, which
adds to $\operatorname{lab}_1$ every component of $Z$ meeting it. The site geometry gives
$d_{k,Z}(Y(\operatorname{lab}_1))\le d_k(\operatorname{lab}_1)$. The reading condition for
links two to four holds at a hub: a cube of the label $\operatorname{lab}_{t-1}$ built
at link $t-1$ ($t=2,3,4$) shares a wall with that hub, so a member of $\mathsf{y}_t$ sharing a
wall with that hub lies within distance $h_\sharp$ of $\operatorname{lab}_{t-1}$, and
$h_\sharp$ does not exceed the radius of the reading condition.

\medskip\noindent\emph{Inputs.}
Three inputs are used. The first is on the pieces. It is the amplitude-disc bound for pointed
pieces in the form \eqref{5.2:eq-carrier-exterior-d}, with $\eta_p$ as in
\eqref{5.2:eq-carrier-exterior-a}:
\begin{equation}\label{5.2:eq-rooted-sum-notation-6}
|v_p|\le\mathfrak{N}^{\mathrm{ex}}_A(p)
=2\mathsf{w}_L^{-1}(\eta_p\mathsf{f}_v)^{1/2}\,\tfrac{81}{64}\,\mathfrak{N}^{\flat}(v)\,
e^{-(\kappa_1-1)m(p)},\qquad \mathsf{w}_L=1/\theta'(g_k).
\end{equation}
Here $\eta_p=1$ if the $2$-stencil $\operatorname{Sten}_2(S_v)$ meets $\square^{\zeta}$, and
$\eta_p=e^{-\delta_Pd(\square^{\zeta},S_v)}$ otherwise. Hence
$\eta_p\le e^{-\delta_Pd(\tilde\square,S_v)}$ at every block; this comparison is the first
clause of the corollary to the touching-block lemma of this section. The block-sum lemma of
this section is applied through this majorant, in shells of $d(\tilde\square,S_v)$.

The second input is on the cells. It is the column bound per cell
\eqref{5.2:eq-exposed-column-a} of Lemma~\ref{5.2:lem-exposed-column}. In addition, the
dilated sizes sum to at most $\tfrac32$ times its right member, since
$(\tfrac98)^3=\tfrac{729}{512}\le\tfrac32$.

The third input is on the weights. It is the geometric clause
\eqref{5.2:eq-carrier-exterior-e} of Lemma~\ref{5.2:lem-carrier-exterior}. Suppose
$w^{\vee}(\Delta')>1$ is attained on a
foreign live arrest $\mathfrak{a}'$. Then $n'\le k_{a'}$, and
$d_{\mathrm{sc}}(y,\Lambda^{\sim})\ge2\check R_{k_{a'}}-6$ for every $y$ of $S_v$. Moreover
\begin{equation}\label{5.2:eq-rooted-sum-notation-7}
\bar{\mathsf{w}}(\Delta')\le3\mathsf{L}_0^2LC_{\uparrow}\check R_{k_{a'}},\qquad\text{hence}\qquad
\mathsf{f}(\Delta')^{1/2}\,\bar{\mathsf{w}}(\Delta')^3\le\mathsf{C}_{\mathrm{fw}},
\end{equation}
with $\mathsf{C}_{\mathrm{fw}}$ as in Lemma~\ref{5.2:lem-carrier-exterior}. Elsewhere
$\bar{\mathsf{w}}=3$ and $27c_\chi^{1/2}\le\mathsf{C}_{\mathrm{fw}}$.
\end{definition}

\begin{proof}
We verify the assertions made in the definition.

\emph{Source cells.} A source cell is a member of $\pi^{\mathfrak{s}}$ sharing a wall with a hub
$\Lambda^{\sim}_C$. Since $\operatorname{dist}_k(\Lambda^{\sim},Z^c)\ge9\check R_k$, it lies in
$Z$. It is disjoint with $\Lambda^{\sim}\supset\Lambda$ and lies in $Z$, so it has level $k$ by
the hub geometry. Being a cell of level $k$, it is a $\pi_k$-cube.

\emph{The constant $c_l^{\circ}$.} Since $h_\sharp=112$, we have
\[
c_l^{\circ}=1+d(h_\sharp+1)=1+113d\le1+234d=c_l .
\]

\emph{The bound on $\mathsf{q}_{\square}$.} The second floor in
\eqref{5.2:eq-rooted-sum-notation-a} gives
$\tfrac12\delta_P\tfrac{M}{M_1}\ge8c_{\boxplus}\log\hat b_t+8$. Hence
\[
\mathsf{q}_{\square}=\hat b_t^{\,c_{\boxplus}}e^{-\frac12\delta_PM/M_1}
\le\hat b_t^{\,c_{\boxplus}}\hat b_t^{-8c_{\boxplus}}e^{-8}
=\hat b_t^{-7c_{\boxplus}}e^{-8}\le e^{-8},
\]
because $\hat b_t\ge1$.

\emph{Proof of \eqref{5.2:eq-rooted-sum-notation-3}.} Let $0\le r\le a+9$. By definition,
$\lambda_\Sigma'=\kappa_1-2-c_l^{\circ}(r+1)-\log(e\hat b)$. The first floor in
\eqref{5.2:eq-rooted-sum-notation-a} gives $\kappa_1-2\ge c_l(2\kappa+4)+\lambda_*(d,L)-1$.
Since $c_l^{\circ}\le c_l$ and $r+1\le a+10$,
\[
\lambda_\Sigma'\ge c_l\bigl(2\kappa+4-(a+10)\bigr)+\lambda_*(d,L)-\log(e\hat b)-1 .
\]
Now $2\kappa+4-(a+10)=\tfrac9{20}\kappa-6$, and $\kappa\ge180$ gives
$\tfrac9{20}\kappa-6\ge81-6=75$. This proves \eqref{5.2:eq-rooted-sum-notation-3}.

\emph{Consequences.} The defining property of $\lambda_*(d,L)$ gives
$e^{-\frac12(\lambda_*(d,L)-\log(e\hat b))}\le\tfrac12L^{-d}<1$, so
$\lambda_*(d,L)-\log(e\hat b)>0$. Since $75c_l-1\ge3$,
\eqref{5.2:eq-rooted-sum-notation-3} yields $\lambda_\Sigma'\ge3$. Squaring the same property
gives $e^{-(\lambda_*(d,L)-\log(e\hat b))}\le\tfrac14L^{-2d}$. Inserting this into
\eqref{5.2:eq-rooted-sum-notation-3},
\[
e^{-\lambda_\Sigma'}\le e^{-(\lambda_*(d,L)-\log(e\hat b))}e^{-(75c_l-1)}
\le\tfrac14L^{-2d}e^{-(75c_l-1)}.
\]

\emph{The label of a compact piece.} For a compact $p$ the star lemma of this section gives
$[\operatorname{St}_v]_k\subset\mathfrak{h}_v$. Hence, by
\eqref{5.2:eq-rooted-sum-notation-5},
\begin{align*}
K_p\cup\mathfrak{h}_v
&=[\operatorname{St}_v]_k\cup[\mathsf{y}]_k\cup\operatorname{Hub}(\mathsf{y})\cup\mathfrak{h}_v\\
&=\mathfrak{h}_v\cup[\mathsf{y}]_k\cup\operatorname{Hub}(\mathsf{y})\\
&=\operatorname{lab}(\mathfrak{h}_v,\mathsf{y}).
\end{align*}
The hypothesis $[\operatorname{St}_v]_k\subset A$ of the definition of $\operatorname{lab}(A,\mathsf{y})$
is met at $A=\mathfrak{h}_v$ by the same inclusion.
\end{proof}

For a finite wall-connected family $\operatorname{St}\subset\pi^{\mathfrak{s}}$ that contains no source cell, put
\begin{equation}\label{5.2:eq-piece-source-cell-1}
\mathfrak{Y}(\operatorname{St}):=\bigl\{\mathsf{y}\subset\pi^{\mathfrak{s}}\setminus\operatorname{St}\ \text{finite}:\
\operatorname{St}\cup\mathsf{y}\ \text{is connected and}\ \mathsf{y}\ \text{contains a source cell}\bigr\}.
\end{equation}
Cells, levels, hubs (the boxes $\Lambda^{\sim}_C$), source cells and adjacency are those of
Notation~\ref{5.2:not-rooted-sum-notation}. Two members are adjacent if they share a wall, or if both
share a wall with one hub, and ``connected'' refers to this adjacency. By this geometry, the
components of the sets $Y(\sigma_s)$ introduced in hypothesis (a) below are the unions of their
wall-components joined through hubs. For a family
$\operatorname{St}$, the rim $\partial\operatorname{St}$ is the set of members of
$\pi^{\mathfrak{s}}\setminus\operatorname{St}$ that share a wall with a member of
$\operatorname{St}$.

\begin{lemma}[A non-vanishing piece is tied to a source cell]\label{5.2:lem-piece-source-cell}
Assume the following four statements, at the localisation site of the step and for every
bracket $s\in\{1,\dots,4\}$. Here $\sigma_s$ denotes the configuration of the decoupling
parameters $\mathsf{s}(\Delta)$ of the bracket $s$ on the cells of a cut family
$\sigma_0\subset\pi^{\mathfrak{s}}$ (arbitrary in (a), and $\sigma_0=\sigma_0(v)$, defined below,
in the conclusion), and $P_1(\sigma_s)$,
$H(\sigma_s)$, $\mathfrak{m}_H(\sigma_s)$ are the operators of the bracket at that configuration.
\begin{enumerate}
\item[(a)] The localisation lemma for the decoupling parameters, valid uniformly for every
subfamily $\sigma_0\subset\pi^{\mathfrak{s}}$. Let $\mathsf{s}=0$ off a finite
$\mathsf{y}\subset\sigma_0$. Let $Y(\sigma_s)$ be the union of the never-cut set
$\bigl(\bigcup(\pi^{\mathfrak{s}}\setminus\sigma_0)\bigr)\cup\Lambda^{\sim}$ with
$\bigcup\mathsf{y}$. Then on each component of $Y(\sigma_s)$ the objects
$A'_s,X_s,\mathbb{H}_s=A'_s-H(\sigma_s)X_s$ depend on $\mathsf{s}$, on the datum $\mathsf{d}_s$
of the bracket (the field datum of the bracket $s=3$ included) and on the input configuration
$(\mathbf{U},\mathbf{J})$ only through their restrictions to that component. They vanish
identically on every component that does not meet the support of $\mathsf{d}_s$. The same
locality holds for the background pair $(W_s,J_s)$ of the bracket: on each component of
$Y(\sigma_s)$ it depends on $\mathsf{s}$ only through the restriction of $\mathsf{s}$ to that
component.
\item[(b)] The kernel-locality statement. The kernels of $P_1(\sigma_s)$, $H(\sigma_s)$ and
$\mathfrak{m}_H(\sigma_s)$ vanish unless both arguments lie in one component of $Y(\sigma_s)$.
Ba{\l}aban proves the same property for the operators $H$, $G$, $H_0$ of
\cite[\S1, after~(1.9)]{B-CMP116} as functions of the decoupling parameters.
\item[(c)] The support statement for the bracket data. The datum $\mathsf{d}_s$ is supported in
the region $\Lambda\subset\Lambda^{\sim}$ in which the changes of the step are made.
\item[(d)] Locality of the carrier in $\chi$. For every multiplier $\chi$ on the cells (below,
$\chi=t+\mathsf{t}\hat{\chi}_{\square}$), the first slot $e^{i\eta\chi\mathbb{H}_s}W_s$ and
every member of the current \eqref{5.2:eq-piece-source-cell-3} below except
$-P_1(\sigma_s)(\chi A'_s)$ depend on $\chi$ locally; in particular, on a component of
$Y(\sigma_s)$ on which $A'_s,X_s,\mathbb{H}_s$ vanish identically they are independent of $\chi$.
\end{enumerate}
Let $v$ be a plain unit ($\mathcal{P}_v=\emptyset$) of one of the two classes of units at the
localisation site:
\begin{itemize}
\item a unit with $K_v\setminus Z\neq\emptyset$, with the single bracket $s=4$; or
\item a unit with $X_v\cap Z=\emptyset$, with the brackets $s=4,\dots,1$.
\end{itemize}
Let $\operatorname{St}_v\subset\pi^{\mathfrak{s}}$ be its exempt family: finite,
wall-connected, and with $S_v\subset\bar S_v:=\bigcup\operatorname{St}_v$, where $S_v$ is the
read set of $v$. Suppose that $\operatorname{St}_v\cup\partial\operatorname{St}_v$ contains no
source cell.

The expansion runs over families of cells in proper scales: the cut cells are the members of
$\sigma_0(v):=\pi^{\mathfrak{s}}\setminus\operatorname{St}_v$. A piece is
$p=(\square,\mathsf{y};s)$, with $\mathsf{y}\subset\sigma_0(v)$ finite and $\square$ a block of
the cover of $\mathfrak{B}_s$, and
\begin{equation}\label{5.2:eq-piece-source-cell-2}
v_p=\prod_{\Delta\in\mathsf{y}}\int_0^1 d\mathsf{s}(\Delta)\,\partial_{\mathsf{s}(\Delta)}\,
\partial_{\mathsf{t}}\big|_{\mathsf{t}=0}\,
f_v\bigl(\mathsf{R}_p(\mathsf{t},\mathsf{s})\bigr)\Big|_{\mathsf{s}=0\ \text{off}\ \mathsf{y}},
\qquad m(p)=|\mathsf{y}| .
\end{equation}
Here $f_v$ is the function of the carrier that defines the unit $v$, and
$\mathsf{R}_p(\mathsf{t},\mathsf{s})=\Phi_s^{\chi}(\mathsf{s})$ at
$\chi:=t+\mathsf{t}\hat{\chi}_{\square}$ for a
fixed $t\in[0,1]$. Explicitly,
$\Phi_s^{\chi}=(e^{i\eta\chi\mathbb{H}_s}W_s,\ J_s+\mathcal{J}^{\sharp}_s[\chi])$ with
\begin{equation}\label{5.2:eq-piece-source-cell-3}
\begin{split}
\mathcal{J}^{\sharp}_s[\chi]
&=(\Delta_{W_s}-P_1(\sigma_s)+\mathsf{K})(\chi A'_s)
-\chi\,\mathfrak{m}_H(\sigma_s)X_s\\
&\quad-[D^*D,\chi]\,H(\sigma_s)X_s+F(W_s,\chi\mathbb{H}_s).
\end{split}
\end{equation}
Here $\Delta_{W_s}$, $\mathsf{K}$, $D^*D$ and $\mathfrak{m}_H(\sigma_s)$ are the operators, and
$F(W_s,\cdot)$ the remainder, of the current at the localisation site.
Then $v_p=0$ unless $\mathsf{y}\in\mathfrak{Y}(\operatorname{St}_v)$.
\end{lemma}

\begin{proof}
The set $\mathfrak{Y}(\operatorname{St}_v)$ is defined, because $\operatorname{St}_v$ is finite,
wall-connected and contains no source cell.

Apply hypothesis (a) to the subfamily $\sigma_0(v)$; this is admissible because (a) holds for
every subfamily. The never-cut set is $\bar S_v\cup\Lambda^{\sim}$. At $\mathsf{s}=0$ off
$\mathsf{y}$ we therefore have
$Y(\sigma_s)=\bar S_v\cup\Lambda^{\sim}\cup\bigcup\mathsf{y}$.

Two consequences follow. By (a), on each component of $Y(\sigma_s)$ the objects
$A'_s,X_s,\mathbb{H}_s$ depend on $\mathsf{s}$, on the datum and on the input configuration only
through their restrictions to that component, and they vanish identically on a component off the
support of the datum; the background pair $(W_s,J_s)$ depends on $\mathsf{s}$ there only through
its restriction to that component. By (b), the
kernels of $P_1(\sigma_s)$, $H(\sigma_s)$, $\mathfrak{m}_H(\sigma_s)$ vanish unless both
arguments lie in one component of $Y(\sigma_s)$.

The function $f_v\circ\mathsf{R}_p$ depends on the carrier only through its restriction to
$S_v\subset\bar S_v$. Since
$\operatorname{St}_v$ is wall-connected, $\bar S_v$ lies in one component of $Y(\sigma_s)$; call
it $Y^v$.

($\alpha$) Let $\Delta\in\mathsf{y}$ with $\Delta\not\subset Y^v$. By the two consequences above,
$f_v\circ\mathsf{R}_p$ does not depend on $\mathsf{s}(\Delta)$, so $\partial_{\mathsf{s}(\Delta)}$
annihilates it and \eqref{5.2:eq-piece-source-cell-2} gives $v_p=0$. This is Ba{\l}aban's
observation in \cite[\S1, after~(1.9)]{B-CMP116} that the derivatives with respect to the
parameters of components
other than the one read by the term make the term vanish. Hence, if $v_p\neq0$, then
$\bigcup\mathsf{y}\subset Y^v$. By the geometry recalled before the lemma, the components of
$Y(\sigma_s)$ are unions of wall-components joined through hubs. Since $\bar S_v$ and
$\bigcup\mathsf{y}$ both lie in $Y^v$, the family $\operatorname{St}_v\cup\mathsf{y}$ is
connected.

($\beta$) Suppose $Y^v$ contains no hub. Then $Y^v$ does not meet $\Lambda^{\sim}$. By (c) the
datum, which is supported in $\Lambda\subset\Lambda^{\sim}$, vanishes on $Y^v$, and by (a)
$A'_s\equiv X_s\equiv\mathbb{H}_s\equiv0$ on $Y^v$. Consequently, by hypothesis (d), the first
slot $e^{i\eta\chi\mathbb{H}_s}W_s$ is independent of $\chi$ there, and so is every member of
\eqref{5.2:eq-piece-source-cell-3} other than the excepted member
$-P_1(\sigma_s)(\chi A'_s)$, which is handled next.

The excepted member splits as
\begin{equation*}
-P_1(\sigma_s)(\chi A'_s)=-t\,P_1(\sigma_s)A'_s-\mathsf{t}\,P_1(\sigma_s)(\hat{\chi}_{\square}A'_s),
\end{equation*}
the last summand being the tail. By (b), both summands, restricted to $S_v$, depend only on
$A'_s$ on $Y^v$, and therefore vanish. Thus $\mathsf{R}_p(\mathsf{t},\mathsf{s})$ restricted to
$S_v$ does not depend on $\mathsf{t}$. Hence
$\partial_{\mathsf{t}}|_{\mathsf{t}=0}f_v(\mathsf{R}_p(\mathsf{t},\mathsf{s}))=0$ and $v_p=0$.
This is Ba{\l}aban's remark that the other terms of the expansion vanish because of the
localisation of the field $B$ \cite[p.~372]{B-CMP122b}.

It follows that if $v_p\neq0$, then $Y^v$ contains a hub $H$. Because $Y^v$ is a union of
wall-components of $Y(\sigma_s)$ joined through hubs, $H$ shares a wall with a member of
$\operatorname{St}_v\cup\mathsf{y}$. It cannot share a wall with a member of
$\operatorname{St}_v$: that member would then be a source cell in $\operatorname{St}_v$, contrary
to the hypothesis. So $H$ shares a wall with a member of $\mathsf{y}$, and that member is a
source cell.

Combining these facts: $\mathsf{y}\subset\pi^{\mathfrak{s}}\setminus\operatorname{St}_v$ is
finite, $\operatorname{St}_v\cup\mathsf{y}$ is connected, and $\mathsf{y}$ contains a source
cell. That is, $\mathsf{y}\in\mathfrak{Y}(\operatorname{St}_v)$.
\end{proof}

\begin{lemma}[The exempt star of a cell]\label{5.2:lem-exempt-star}
In the setting and notation fixed in \ref{5.2:not-rooted-sum-notation} (notation for the
source-rooted sum over cut cells), let $\Delta'$ be a cell of the current global label and
$n' := \operatorname{lev}\Delta'$. Put
\begin{equation}\label{5.2:eq-exempt-star-1}
Q_5(\Delta') := \{x : \operatorname{dist}_{n'}(x,\Delta') \le 2\},
\end{equation}
and let $\operatorname{St}(\Delta')$ be the wall-component of $\Delta'$ in the family of cells $\Delta$
with $\operatorname{lev}\Delta \ge n'-1$ that meet $Q_5(\Delta')$. Here a cell meets a set if the two
have a common interior point. Let $\partial\operatorname{St}(\Delta')$ be the family of members of
$\pi^{\mathfrak{s}}\setminus\operatorname{St}(\Delta')$ that share a wall with a member of
$\operatorname{St}(\Delta')$, and let $c_S$, $N_\partial$, $\epsilon_\partial$ be as in display
\eqref{5.2:eq-rooted-sum-notation-a} of that notation. Then the following hold.
\begin{enumerate}
\item[(e1)] The members of $\operatorname{St}(\Delta')$ have levels $n'-1$, $n'$, $n'+1$.
\item[(e2)] $\#\operatorname{St}(\Delta') \le c_S$.
\item[(e3)] $\#\partial\operatorname{St}(\Delta') \le 2dL^{d-1}c_S$, and
$\epsilon_\partial\,\#\partial\operatorname{St}(\Delta') \le 1$.
\item[(e4)] Every cell $\Delta_1$ lies in $\partial\operatorname{St}(\Delta')$ for at most $N_\partial$
cells $\Delta'$.
\item[(e5)] Every box $V \subset Q_5(\Delta')\cap\Omega_{n'-1}$ meeting $\Delta'$ lies in
$\bigcup\operatorname{St}(\Delta')$; in particular
$\{\operatorname{dist}_{n'}(\cdot,\Delta') \le 1\} \subset \bigcup\operatorname{St}(\Delta')$.
Moreover
\begin{equation*}
\bigcup\operatorname{St}(\Delta') \subset \{\operatorname{dist}_{n'}(\cdot,\Delta') \le L+2\},
\end{equation*}
with $2$ in place of $L+2$ if no member has level $n'+1$. Finally, let $v$ be a compact pointed unit,
with point $y_v$ lying in its hull $\hat X_v$, and let
$\mathfrak{h}_v = \{\Delta\in\pi_k : \operatorname{dist}_k(\Delta,\hat X_v)\le 2\}$ be the
$\pi_k$-neighbourhood of that hull; then $y_v \in \Delta'$ implies
$[\operatorname{St}(\Delta')]_k \subset \mathfrak{h}_v$.
\item[(e6)] If $\Delta' \not\subset Z$, or $\Delta'$ lies within $2$ level-$k$ units of $\partial Z$,
then $\operatorname{St}(\Delta')\cup\partial\operatorname{St}(\Delta') \subset \pi^{\mathfrak{s}}$
and this family holds no source cell.
\end{enumerate}
\end{lemma}

\begin{proof}
Throughout we use two properties of the current global label recorded in
\ref{5.2:not-rooted-sum-notation}: wall- or corner-neighbouring cells differ in level by at most $1$
(the gradation property), and $\operatorname{dist}_{j+1}(\Omega_j^c,\Omega_{j+1}) \ge 1$ for every $j$
(the separation property).

(e1) A member has level $\ge n'-1$ by definition. For the upper end, note that
$\Delta' \subset \Omega^c_{n'+1}$, since $\Delta'$ has level $n'$. By the separation property at
$j = n'+1$, the set
$\Omega_{n'+2}$ lies at distance at least one level-$(n'+2)$ unit from $\Omega^c_{n'+1}$, that is at
distance at least $L^2$ level-$n'$ units, which is more than $2$ level-$n'$ units. A cell of level
$\ge n'+2$ lies in $\Omega_{n'+2}$, whereas every point of $Q_5(\Delta')$ is within $2$ level-$n'$
units of $\Delta'$. Hence no cell of level $\ge n'+2$ meets $Q_5(\Delta')$, and every member has
level at most $n'+1$.

(e2) The partitions are nested, and $Q_5(\Delta')$ is a union of $\pi_{n'}$-cubes, hence of
$\pi_{n'-1}$-cubes. A member has level $\ge n'-1$ and is therefore a union of $\pi_{n'-1}$-cubes.
It has a common interior point with $Q_5(\Delta')$, so it contains a cube of $\pi_{n'-1}$ lying in
$Q_5(\Delta')$. Distinct members have disjoint interiors and therefore contain distinct such cubes.
The set $Q_5(\Delta')$ is a cube of side $5$ level-$n'$ units, that is $5L$ level-$(n'-1)$ units, and
so holds $(5L)^d$ cubes of $\pi_{n'-1}$. Hence $\#\operatorname{St}(\Delta') \le (5L)^d = c_S$.

(e3) By the gradation property, a cell has at most $2dL^{d-1}$ wall-neighbours. Therefore
$\#\partial\operatorname{St}(\Delta') \le 2dL^{d-1}\#\operatorname{St}(\Delta') \le 2dL^{d-1}c_S$
by (e2). The notation for the source-rooted sum gives $\epsilon_\partial = 2e^{-\lambda_\Sigma'}$ with
$e^{-\lambda_\Sigma'} \le \tfrac14 L^{-2d}e^{-(75c_l-1)}$. Hence
\begin{equation*}
\epsilon_\partial\cdot 2dL^{d-1}(5L)^d \le d\,5^dL^{-1}e^{-(75c_l-1)} \le 1 ,
\end{equation*}
where the last inequality holds because $c_l = 1+234d \ge d\log 5+\log d+1$.

(e4) Let $\Delta_1$ share a wall with some $\Delta_2 \in \operatorname{St}(\Delta')$, and put
$n_1 := \operatorname{lev}\Delta_1$. By (e1), $\operatorname{lev}\Delta_2 \in [n'-1,n'+1]$. By the
gradation property,
$n_1$ differs from $\operatorname{lev}\Delta_2$ by at most $1$. Hence $n' \in [n_1-2,n_1+2]$, which
allows five levels. Measured in level-$n'$ units, $\Delta_2$ has side at most $L$ and meets
$Q_5(\Delta')$; since $\Delta_1$ touches $\Delta_2$, this gives
$\operatorname{dist}_{n'}(\Delta_1,\Delta') \le L+2$. In the same units $\Delta_1$ has side at most
$L^2$. The cell $\Delta'$, of level $n'$, is a $\pi_{n'}$-cube. The $\pi_{n'}$-cubes within $L+2$ of
$\Delta_1$ number at most $(L^2+2L+6)^d \le (L+2)^{2d}$. Summing over the five admissible values of
$n'$, at most $5(L+2)^{2d} = N_\partial$ cells $\Delta'$ have $\Delta_1 \in \partial\operatorname{St}(\Delta')$.

(e5) Since $V \subset \Omega_{n'-1}$, the cells meeting $V$ have level $\ge n'-1$. Since
$V \subset Q_5(\Delta')$, they meet $Q_5(\Delta')$. The box $V$ is convex. Join a generic interior
point of $V$ lying in such a cell to an interior point of $V\cap\Delta'$ by a segment. This segment
stays in $V$ and avoids all $(d-2)$-faces, so consecutive cells along it share a wall. Hence every
cell meeting $V$ is wall-connected to $\Delta'$ within the family, that is, it belongs to
$\operatorname{St}(\Delta')$. By the separation property at $j = n'-1$, the cube
$\{\operatorname{dist}_{n'}(\cdot,\Delta') \le 1\}$ lies in $\Omega_{n'-1}$. It is a box in
$Q_5(\Delta')$ meeting $\Delta'$, which gives the particular case.

For the outer bound, consider first a member of level $\le n'$. It lies in a single
$\pi_{n'}$-cube, and $Q_5(\Delta')$ is a union of such cubes, so it lies in $Q_5(\Delta')$. A member
of level $n'+1$ has side $L$ and meets $Q_5(\Delta')$, so it lies within $L+2$ of $\Delta'$. If no
member has level $n'+1$, everything lies in $Q_5(\Delta')$, i.e.\ within $2$.

Now let $y_v \in \Delta'$ and $x \in \bigcup\operatorname{St}(\Delta')$. If $n' = k$, no member has
level $k+1$, since no cell has level above $k$. Then $x$ lies in a $\pi_k$-cube of $Q_5(\Delta')$,
which is at $\operatorname{dist}_k \le 1$
from $\Delta' \ni y_v$. Hence this cube is at $\operatorname{dist}_k \le 2$ from the hull
$\hat X_v \ni y_v$. If $n' < k$, then by the outer bound $x$ is within $L+3$ level-$n'$ units of
$y_v$. Therefore
\begin{equation*}
|x-y_v|_\infty \le (L+3)/L \le 2 \quad\text{level-}k\text{ units}.
\end{equation*}
In both cases the $\pi_k$-cube of $x$ belongs to $\mathfrak{h}_v$.

(e6) A member of $\partial\operatorname{St}(\Delta')$ has level $\le \min\{n'+2,k\}$, by (e1) and
the gradation property and since levels do not exceed $k$; hence it has side
at most $1$ level-$k$ unit. Together with (e5), this shows that every member of
$\operatorname{St}(\Delta')\cup\partial\operatorname{St}(\Delta')$ lies within $4$ level-$k$ units of
$\Delta'$. Since $\Delta'\not\subset Z$ or $\Delta'$ lies within $2$ level-$k$ units of $\partial Z$,
and $\Delta'$ has side at most $1$ level-$k$ unit, every such member lies within $7$ level-$k$ units
of $Z^c$. By the hub geometry recorded in \ref{5.2:not-rooted-sum-notation},
$\operatorname{dist}_k(\Lambda^{\sim},Z^c) \ge 9\check{R}_k$, and $\check{R}_k \ge 38$ by the lower
bound on $\check{R}_k$, so $9\check{R}_k \ge 342$. Consequently no member meets
$\Lambda^{\sim}$, so the family lies in $\pi^{\mathfrak{s}}$. Moreover, a source cell is a
$\pi_k$-cube sharing a wall with a hub, hence lies within $1$ level-$k$ unit of $\Lambda^{\sim}$ and
so at $\operatorname{dist}_k \ge 9\check{R}_k-1 \ge 341$ from $Z^c$; hence none of these members is a
source cell.
\end{proof}

\begin{lemma}[The block sum in shells around a family of cells]\label{5.2:lem-shell-block-sum}
In the notation for the source-rooted sum over cut cells fixed in~\ref{5.2:not-rooted-sum-notation}, and under the parameter conditions \eqref{5.2:eq-rooted-sum-notation-a}, let $s\le 4$ and let $\mathrm{St}$ be a finite family of cells of $\mathfrak{B}_s$. Put $S:=\bigcup\mathrm{St}$, write $n_-:=\min_{\Delta\in\mathrm{St}}\operatorname{lev}\Delta$ for the least level of a member of $\mathrm{St}$, and set
\begin{equation}\label{5.2:eq-shell-block-sum-a}
C^0_{\square}:=2(L+9)^d\,\hat b_t^{\,2c_{\boxplus}},\qquad
\mathsf{b}_{\square}(\mathrm{St}):=C^0_{\square}\,\#\mathrm{St}\,(\#\mathrm{St}+3).
\end{equation}
Let the sums below run over all blocks $\square$ of the cover of $\mathfrak{B}_s$, including those inside the hubs. Then
\begin{equation}\label{5.2:eq-shell-block-sum-1}
\sum_{\square}e^{-\frac12\delta_P d(\tilde\square,S)}\le C^0_{\square}\,\#\mathrm{St}\,(k-n_-+3).
\end{equation}
Suppose moreover that $\mathrm{St}\subset\pi^{\mathfrak{s}}$ is wall-connected and holds no source cell. Let $\mathfrak{Y}(\mathrm{St})$ be the set of finite families $\mathsf{y}\subset\pi^{\mathfrak{s}}\setminus\mathrm{St}$ such that $\mathrm{St}\cup\mathsf{y}$ is connected and $\mathsf{y}$ holds a source cell. Then for every $\mathsf{y}\in\mathfrak{Y}(\mathrm{St})$
\begin{equation}\label{5.2:eq-shell-block-sum-2}
e^{-|\mathsf{y}|}\sum_{\square}e^{-\frac12\delta_P d(\tilde\square,S)}\le\mathsf{b}_{\square}(\mathrm{St}).
\end{equation}
\end{lemma}

\begin{proof}
Throughout, $d$ and $d^{\wedge}$ are the two contour distances attached to $\mathfrak{B}_s$. We use the following properties of the cell partition and of the label, recorded with the notation fixed in~\ref{5.2:not-rooted-sum-notation}.
\begin{itemize}
\item Comparison of distances: $d\ge d^{\wedge}$, and $d^{\wedge}(y,y')\ge w(|m-m'|-1)_+$ for points $y,y'$ in cells of levels $m,m'$.
\item Nesting: $\Omega_{i+1}\subset\Omega_i$ for every $i$.
\item Gradation: wall- or corner-neighbouring cells differ in level by at most one.
\item Separation: $\operatorname{dist}_{j+1}(\Omega_j^c,\Omega_{j+1})\ge 1$.
\item Contour count: a contour $\Gamma$ with $\operatorname{len}_{\mathfrak{B}}(\Gamma)\le xw$ meets at most $c_{\boxplus}(x+1)$ cells, and these cells form a family connected by touching.
\end{itemize}
We also use that a cell touches at most $\hat b_t$ other cells, which is recorded there as well. Since $d\ge d^{\wedge}$, every block satisfies
\begin{equation*}
e^{-\frac12\delta_P d(\tilde\square,S)}\le e^{-\frac12\delta_P d^{\wedge}(\tilde\square,S)}.
\end{equation*}

\emph{Shells.} For a block $\square$ put $\mathsf{i}=\mathsf{i}(\square):=\lfloor d^{\wedge}(\tilde\square,S)/w\rfloor$, the shell index of $\square$. Then $d^{\wedge}(\tilde\square,S)\ge w\mathsf{i}$, so
\begin{equation*}
e^{-\frac12\delta_P d(\tilde\square,S)}\le e^{-\frac12\delta_P w\mathsf{i}}.
\end{equation*}
Since $d^{\wedge}(\tilde\square,S)<(\mathsf{i}+1)w$, there is a contour $\Gamma$ from a point of $S$ to a point $z\in\tilde\square$ with $\operatorname{len}_{\mathfrak{B}}(\Gamma)<(\mathsf{i}+1)w$. By the contour count at $x=\mathsf{i}+1$, $\Gamma$ meets at most $c_{\boxplus}(\mathsf{i}+2)$ cells. These cells are connected by touching, and the first of them is a member of $\mathrm{St}$, because $\Gamma$ starts in $S$. Let $\Delta''\ni z$ be the last of them and put $n'':=\operatorname{lev}\Delta''$.

\emph{(i) Cells per shell.} Each cell touches at most $\hat b_t$ others. Hence the cells reached from $\mathrm{St}$ by touching-chains of at most $c_{\boxplus}(\mathsf{i}+2)$ cells number at most
\begin{equation*}
\#\mathrm{St}\sum_{m<c_{\boxplus}(\mathsf{i}+2)}\hat b_t^{\,m}\le\#\mathrm{St}\cdot\hat b_t^{\,c_{\boxplus}(\mathsf{i}+2)},
\end{equation*}
where the last step uses $\hat b_t=(L+2)^d+3^d\ge 2$. The cell $\Delta''$ attached to $\square$ is one of these cells.

\emph{(ii) Levels.} The starting point of $\Gamma$ lies in a member of $\mathrm{St}$, whose level is at least $n_-$, and $z\in\Delta''$. The length $\operatorname{len}_{\mathfrak{B}}(\Gamma)$ is at least the $d^{\wedge}$-distance of its endpoints, so by the comparison of distances,
\begin{equation*}
(\mathsf{i}+1)w>\operatorname{len}_{\mathfrak{B}}(\Gamma)\ge w(n_--n''-1).
\end{equation*}
Since levels are integers, $n''\ge n_--\mathsf{i}-1$.

\emph{(iii) Zones over one cell.} Let $\square$ have zone $n$, with $z\in\tilde\square\cap\Delta''$. By the definition of the fattened block, $\tilde\square\subset\{\operatorname{dist}_n(\cdot,\Delta_n)\le 2\}$ for a cell $\Delta_n$ of level $n$. The blocks are unions of $2^d$ neighbouring cubes of the partition of their zone, as in the cover of \cite{B-CMP119}. Suppose $n\le n''-2$. Then $\Delta''\subset\Omega_{n''}\subset\Omega_{n+2}$, the second inclusion by nesting, and $\Delta_n\subset\Omega^c_{n+1}$. Separation at $j=n+1$ puts these sets at least one level-$(n+2)$ unit apart, that is, at least $L^2>2$ level-$n$ units. This contradicts $z\in\Delta''\cap\tilde\square$. Hence $n\in[n''-1,k]$, and by (ii) the admitted zones number at most
\begin{equation*}
k-n''+2\le k-n_-+\mathsf{i}+3 .
\end{equation*}
No upper bound on $n$ better than $k$ follows from the separation property.

\emph{(iv) Blocks per zone, counted per reached cell $\Delta''$.} What is needed is $\#\{\square\text{ of zone }n:\tilde\square\cap\Delta''\ne\emptyset\}$ for every admitted $n$. The $\tilde\square$ of zone $n$ are cubes of side $4$ at grid positions of $\pi_n$.
\begin{itemize}
\item For $n\ge n''$, the partitions nest, so $\Delta''$ lies in one cube $P\in\pi_n$. Each $\tilde\square$ meeting $\Delta''$ contains or touches $P$, so at most $6^d$ of them meet $\Delta''$.
\item For the one finer admitted zone $n=n''-1$, the cell $\Delta''$ has side at most $L$ in level-$(n''-1)$ units and $\tilde\square$ has side $4$. So at most $(L+9)^d$ grid positions give a $\tilde\square$ meeting $\Delta''$. This case occurs generically at every interface of zone $n''$ with zone $n''-1$, which the gradation property allows.
\end{itemize}
Put $a_{\square}:=k-n_-+3\ge 3$. Since $6^d\le(L+9)^d$, and by (iii),
\begin{equation*}
\sum_{n\text{ admitted}}\#\{\square\text{ of zone }n:\tilde\square\cap\Delta''\ne\emptyset\}
\le(L+9)^d+6^d(k-n''+1)\le(L+9)^d(a_{\square}+\mathsf{i}).
\end{equation*}

\emph{Sum.} By (i) and (iv), the blocks with shell index $\mathsf{i}$ number at most
\begin{equation*}
\#\mathrm{St}\cdot\hat b_t^{\,c_{\boxplus}(\mathsf{i}+2)}\cdot(L+9)^d(a_{\square}+\mathsf{i}).
\end{equation*}
Each of them contributes at most $e^{-\frac12\delta_P w\mathsf{i}}$. Recall $\mathsf{q}_{\square}=\hat b_t^{\,c_{\boxplus}}e^{-\frac12\delta_P w}\le e^{-8}$ from \eqref{5.2:eq-rooted-sum-notation-a}. Therefore
\begin{align*}
\sum_{\square}e^{-\frac12\delta_P d(\tilde\square,S)}
&\le\sum_{\mathsf{i}\ge 0}e^{-\frac12\delta_P w\mathsf{i}}\cdot\#\mathrm{St}\cdot\hat b_t^{\,c_{\boxplus}(\mathsf{i}+2)}
\cdot(L+9)^d(a_{\square}+\mathsf{i})\\
&=(L+9)^d\hat b_t^{\,2c_{\boxplus}}\#\mathrm{St}\sum_{\mathsf{i}\ge0}\mathsf{q}_{\square}^{\mathsf{i}}(a_{\square}+\mathsf{i})\\
&\le(L+9)^d\hat b_t^{\,2c_{\boxplus}}\#\mathrm{St}\cdot 2a_{\square}.
\end{align*}
The last step uses that, for $\mathsf{q}\le e^{-8}$ and $a\ge3$,
\begin{equation*}
\sum_{\mathsf{i}\ge0}\mathsf{q}^{\mathsf{i}}(a+\mathsf{i})=\frac{a}{1-\mathsf{q}}+\frac{\mathsf{q}}{(1-\mathsf{q})^2}\le 2a .
\end{equation*}
Since $2(L+9)^d\hat b_t^{\,2c_{\boxplus}}=C^0_{\square}$ by \eqref{5.2:eq-shell-block-sum-a}, this is \eqref{5.2:eq-shell-block-sum-1}.

\emph{Second display.} Let $\mathsf{y}\in\mathfrak{Y}(\mathrm{St})$. Then $\mathrm{St}\cup\mathsf{y}$ is connected, and it holds a member of level $n_-$ and a source cell, which has level $k$. Along a path of adjacent members the level changes by at most one at each step:
\begin{itemize}
\item for wall-neighbours this is the gradation property;
\item two members adjacent through a hub both share a wall with that hub, so both are source cells of level $k$.
\end{itemize}
Hence $\#\mathrm{St}+|\mathsf{y}|\ge k-n_-+1$. Since $e^{-m}\le 1$ and $me^{-m}\le 1$ for $m\ge0$, we get
\begin{equation*}
(k-n_-+3)e^{-|\mathsf{y}|}\le(\#\mathrm{St}+2+|\mathsf{y}|)e^{-|\mathsf{y}|}
\le\#\mathrm{St}+2+|\mathsf{y}|e^{-|\mathsf{y}|}\le\#\mathrm{St}+3.
\end{equation*}
Multiplying \eqref{5.2:eq-shell-block-sum-1} by $e^{-|\mathsf{y}|}$ gives
\begin{equation*}
e^{-|\mathsf{y}|}\sum_{\square}e^{-\frac12\delta_P d(\tilde\square,S)}\le C^0_{\square}\,\#\mathrm{St}\,(\#\mathrm{St}+3)=\mathsf{b}_{\square}(\mathrm{St}),
\end{equation*}
which is \eqref{5.2:eq-shell-block-sum-2}.
\end{proof}

Let $s\le 4$ index one of the four pairs of consecutive chain backgrounds at the localisation site. Let $\mathfrak{B}_s$ be its determining set, with its cells and its cover by blocks (Notation~\ref{5.2:not-rooted-sum-notation}). Let $\mathrm{St}$ be a finite family of cells of $\mathfrak{B}_s$, and put $S:=\bigcup\mathrm{St}$. We fix the following conventions:
\begin{itemize}
\item $\ell_j$ is the spacing of the level-$j$ lattice, so that $\ell_{j+1}=L\ell_j$ and a cube of $\pi_j$ has side $M\ell_j$.
\item $\operatorname{dist}_j$ is the distance in the norm $|\cdot|$ measured in units $M\ell_j$.
\item $|\cdot|$ is the maximum norm.
\item $\operatorname{dist}$ without subscript is the distance in the norm $|\cdot|$.
\item A cell of level $i$ is a cube of $\pi_i$ lying in the zone $\Omega_i\setminus\Omega_{i+1}$.
\item A block of zone $m$ is a union of $2^d$ neighbouring cubes of $\pi_m$. It is thus a cube of side $2M\ell_m$, with centre $c_{\square}$. Every block of zone $m$ of the cover contains a cell of level $m$.
\item The support $\square^{\zeta}:=\operatorname{supp}\hat\chi_{\square}$ keeps a margin of $\tfrac13 M\ell_m$, i.e.\ $\tfrac13 M$ lattice spacings of level $m$, inside $\square$ \cite[p.~270]{B-CMP109}, that is,
\begin{equation}\label{5.2:eq-touching-block-4}
|x-c_{\square}|\le \tfrac23 M\ell_m\qquad (x\in\square^{\zeta}).
\end{equation}
\item Two cubes meet if they have a common interior point.
\end{itemize}
We write:
\begin{itemize}
\item $\tilde\square$ for the fattened block of a block $\square$ of the cover of $\mathfrak{B}_s$ (Notation~\ref{5.2:not-rooted-sum-notation}), so that $\square^{\zeta}\subset\square\subset\tilde\square$;
\item $\operatorname{Sten}_2(S)$ for the $2$-stencil of $S$ (Notation~\ref{5.2:not-rooted-sum-notation}): the points within $2\ell_i$ of a point of $S$ lying in a cell of level $i$;
\item $d(\cdot,\cdot)$ for the distance of $\mathfrak{B}_s$; the letter $d$ standing alone is the dimension, $d=4$.
\end{itemize}

In the notation for the source-rooted sum over cut cells (Notation~\ref{5.2:not-rooted-sum-notation}) we set
\begin{equation}\label{5.2:eq-touching-block-1}
\eta^{\zeta}(\square,S):=
\begin{cases}
1, & \operatorname{Sten}_2(S)\cap\square^{\zeta}\neq\emptyset,\\[2pt]
e^{-\delta_P\, d(\square^{\zeta},S)}, & \text{otherwise},
\end{cases}
\qquad
\eta^{\sim}(\square,S):=e^{-\delta_P\, d(\tilde\square,S)} .
\end{equation}
At $S=S_v$ the first of these is the positional factor $\eta_p$. In the first case we call the block $\square$ touching for $S$.

We further put
\begin{equation}\label{5.2:eq-touching-block-2}
C_{\mathrm{blk}}:=2^{d+1}+(L+1)^d .
\end{equation}
This is the bound on the number of blocks met by one cell, provided by the site geometry on which the source-rooted sum rests. The same geometry guarantees that a block of zone $n$ meets cells of the levels $n-1$, $n$, $n+1$ only.

Besides these facts we use two admissibility conditions of that site geometry. They hold at every pair of consecutive levels of the current label, the frozen windows of live arrests included:
\begin{equation}\label{5.2:eq-touching-block-3}
\Omega_{i+1}\subset\Omega_i ,
\end{equation}
\begin{equation}\label{5.2:eq-touching-block-6}
\operatorname{dist}_{j+1}\bigl(\Omega_j^{c},\Omega_{j+1}\bigr)\ge 1 .
\end{equation}

The lemma below is used for touching blocks: there a common point of $\operatorname{Sten}_2(S)$ and $\square^{\zeta}$ provides $x\in\square^{\zeta}$ and a point $x_0\in S$ in a cell of level $i$ with $|x-x_0|\le 2\ell_i$, and the lemma then says that $\square$ meets $S$.

\begin{lemma}[A touching block contains the base point in its interior]\label{5.2:lem-touching-block}
Let $\mathrm{St}$ be a finite family of cells of $\mathfrak{B}_s$ and $S=\bigcup\mathrm{St}$. Assume \eqref{5.2:eq-touching-block-3}, \eqref{5.2:eq-touching-block-6}, $L\ge 2$ and $M>6L$. Let $\square$ be a block of zone $m$ and $x\in\square^{\zeta}$. Let $x_0\in S$ (the base point), and let $\Delta\ni x_0$ be a closed cell of level $i$ with $|x-x_0|\le 2\ell_i$. Then:
\begin{enumerate}
\item $i\le m+1$;
\item $x_0$ is an interior point of $\square$;
\item $\square$ meets every member of $\mathrm{St}$ that contains $x_0$.
\end{enumerate}
\end{lemma}

\begin{proof}
The block $\square$ contains a cell $\Delta_m$ of level $m$ (see the conventions above), so $\Delta_m\subset\Omega_m\setminus\Omega_{m+1}\subset\Omega_{m+1}^{c}$. Both $x$ and $\Delta_m$ lie in $\square$, whose side is $2M\ell_m$. With $|x-x_0|\le 2\ell_i$ this gives
\begin{equation}\label{5.2:eq-touching-block-5}
\operatorname{dist}(x_0,\Delta_m)\le 2M\ell_m+2\ell_i .
\end{equation}

Suppose that $i\ge m+2$. The cell $\Delta$ has level $i$, so $\Delta\subset\Omega_i$. Since $i-1\ge m+1$, condition \eqref{5.2:eq-touching-block-3} gives $\Omega_{i-1}\subset\Omega_{m+1}$, hence $\Delta_m\subset\Omega_{m+1}^{c}\subset\Omega_{i-1}^{c}$. Condition \eqref{5.2:eq-touching-block-6} at $j=i-1$ reads $\operatorname{dist}_i(\Omega_{i-1}^{c},\Omega_i)\ge 1$, that is, distance at least $M\ell_i$. Therefore $\operatorname{dist}(x_0,\Delta_m)\ge M\ell_i$.

Combining this with \eqref{5.2:eq-touching-block-5} gives $(M-2)\ell_i\le 2M\ell_m$. Now $\ell_i\ge L^2\ell_m\ge 4\ell_m$, because $i\ge m+2$ and $L\ge 2$; moreover $M-2>0$. Hence
\begin{equation*}
4(M-2)\ell_m\le (M-2)\ell_i\le 2M\ell_m ,
\end{equation*}
that is, $M\le 4$. This contradicts $M>6L\ge 12$. Hence $i\le m+1$, which is the first claim.

Since $i\le m+1$ we have $\ell_i\le L\ell_m$. Using $M>6L$,
\begin{equation*}
|x-x_0|\le 2\ell_i\le 2L\ell_m<\tfrac13 M\ell_m .
\end{equation*}
By \eqref{5.2:eq-touching-block-4},
\begin{equation*}
|x_0-c_{\square}|\le |x_0-x|+|x-c_{\square}|
<\tfrac13 M\ell_m+\tfrac23 M\ell_m=M\ell_m .
\end{equation*}
Now $M\ell_m$ is half the side of $\square$, so $x_0$ is an interior point of $\square$. This is the second claim.

Finally, let $\Delta_2\in\mathrm{St}$ contain $x_0$. The cell $\Delta_2$ is a closed cube, hence the closure of its interior. Choose a neighbourhood of $x_0$ contained in the interior of $\square$. It contains interior points of $\Delta_2$, and these are interior points of $\square$ as well. So $\square$ and $\Delta_2$ have a common interior point, i.e.\ they meet. This is the third claim.
\end{proof}

The hypothesis $M>6L$ imposes no new condition: it follows from lower bounds on $M$ that are already imposed, each chosen after $L$.

First, the choice of the auxiliary parameters in the correlation theorem requires $M\ge L^3M_1$, where $M_1\ge 1$. Since $L$ is odd and $L\ge 2$, we have $L\ge 3$, so $L^3\ge 9L>6L$ and $M>6L$.

Second, put $c_{\boxplus}:=3^d(L+5)$ and $\hat b_t:=(L+2)^d+3^d$; a cell touches (is in closed contact with) at most $\hat b_t$ others. The lower bound on $M/M_1$ imposed in the same choice of parameters reads
\begin{equation*}
\tfrac1{16}\,\delta_P\, M/M_1\ge c_{\boxplus}\log\hat b_t+1 .
\end{equation*}
Here $\delta_P\le 1$: under the normalisation of the decay rates fixed with the bounds on the background layers one has $e^{d\delta_P}\le 2$, hence $d\delta_P\le\log 2<1$; and $\log\hat b_t\ge 1$. Hence
\begin{equation*}
M/M_1\ge 16\,\delta_P^{-1}\bigl(c_{\boxplus}\log\hat b_t+1\bigr)>16\cdot 3^d(L+5),
\end{equation*}
and, since $M_1\ge 1$, $M\ge M/M_1$, so $M>6L$.

Third, the lower bound $M/M_1\ge 3L(\varrho_0+2L)$, imposed in the bounds on the background layers, where $\varrho_0\ge 0$ denotes the constant fixed with those bounds, gives the same floor: it yields $M/M_1\ge 6L^2>6L$, and again $M\ge M/M_1>6L$.

All three are one-sided lower bounds.

\begin{corollary}[Comparison of the two block-proximity distances]\label{5.2:cor-proximity-comparison}
Assume the hypotheses of Lemmas~\ref{5.2:lem-touching-block} and~\ref{5.2:lem-shell-block-sum}. Let $s\le 4$, let $\mathrm{St}$ be a finite family of cells of $\mathfrak{B}_s$, and put
\begin{equation*}
S:=\bigcup\mathrm{St},\qquad n_-:=\min_{\Delta\in\mathrm{St}}\operatorname{lev}\Delta,\qquad n_+:=\max_{\Delta\in\mathrm{St}}\operatorname{lev}\Delta,
\end{equation*}
where $\operatorname{lev}\Delta$ denotes the level of the cell $\Delta$.
A block $\square$ of the cover of $\mathfrak{B}_s$ is said to \emph{touch} $S$ if $\operatorname{Sten}_2(S)\cap\square^{\zeta}\neq\emptyset$. Put $\eta^{\zeta}(\square,S):=1$ if $\square$ touches $S$ and $\eta^{\zeta}(\square,S):=e^{-\delta_P d(\square^{\zeta},S)}$ otherwise, and $\eta^{\sim}(\square,S):=e^{-\delta_P d(\tilde{\square},S)}$. Then:
\begin{enumerate}
\item[(t1)] at every block, $\eta^{\zeta}(\square,S)\le\eta^{\sim}(\square,S)$ and $\eta^{\zeta}(\square,S)\le e^{-\delta_P d^{\wedge}(\square,S)}$;
\item[(t2)] with $C_{\mathrm{blk}}:=2^{d+1}+(L+1)^d$, the blocks touching $S$ number at most $C_{\mathrm{blk}}\,\#\mathrm{St}$, and their zones lie in $[n_--1,n_++1]$;
\item[(t3)] both displays \eqref{5.2:eq-shell-block-sum-a} of Lemma~\ref{5.2:lem-shell-block-sum}, with the constant $C_{\square}^0$, hold with $\sum_\square e^{-\frac12\delta_P d(\tilde{\square},S)}$ replaced by $\sum_\square\eta^{\zeta}(\square,S)^{1/2}$:
\begin{equation}\label{5.2:eq-proximity-comparison-1}
\begin{gathered}
\sum_{\square}\eta^{\zeta}(\square,S)^{1/2}\le C_{\square}^0\,\#\mathrm{St}\,(k-n_-+3),\\
e^{-|\mathsf{y}|}\sum_{\square}\eta^{\zeta}(\square,S)^{1/2}\le\mathsf{b}_{\square}(\mathrm{St})
\quad\text{for }\mathsf{y}\in\mathfrak{Y}(\mathrm{St}).
\end{gathered}
\end{equation}
\end{enumerate}
\end{corollary}

\begin{proof}
(t1) Let $\square$ touch $S$. Then, by the definition of the $2$-stencil of $S$, there are a point $x\in\square^{\zeta}$ and a point $x_0\in S$, lying in a closed cell $\Delta$ of some level $i$, with $|x-x_0|\le 2\ell_i$. This is the situation of Lemma~\ref{5.2:lem-touching-block}. By that lemma $\square$ meets every member of $\mathrm{St}$ that holds $x_0$, so $\square$ meets $S$. Since $\tilde{\square}\supset\square$, also $\tilde{\square}$ meets $S$. Hence $d(\tilde{\square},S)=d^{\wedge}(\square,S)=0$, and both inequalities read $1\le 1$.

Let $\square$ not touch $S$. Then $\square^{\zeta}\subset\square\subset\tilde{\square}$. A contour starting in $\square^{\zeta}$ starts in $\square$ and in $\tilde{\square}$, so $d(\square^{\zeta},S)\ge d(\tilde{\square},S)$. This gives $\eta^{\zeta}(\square,S)\le\eta^{\sim}(\square,S)$. Further, the restricted contour distance never exceeds the contour distance, $d^{\wedge}\le d$. The distance $d^{\wedge}(\cdot,S)$ does not increase when its first argument is enlarged. Therefore
\begin{equation*}
d(\square^{\zeta},S)\ge d^{\wedge}(\square^{\zeta},S)\ge d^{\wedge}(\square,S),
\end{equation*}
which is the second inequality.

(t2) By the first paragraph, a touching block meets a member of $\mathrm{St}$. We use the block--cell incidence bound of the site geometry: a block of zone $n$ meets cells of levels $n-1,n,n+1$ only, and a cell meets at most $C_{\mathrm{blk}}$ blocks. Assign to each touching block a member of $\mathrm{St}$ that it meets. Each member is assigned at most $C_{\mathrm{blk}}$ times, so the touching blocks number at most $C_{\mathrm{blk}}\,\#\mathrm{St}$. If the assigned member has level $i\in[n_-,n_+]$, the zone $n$ of the block satisfies $|n-i|\le1$. Hence $n\in[n_--1,n_++1]$.

(t3) By (t1), $\eta^{\zeta}(\square,S)^{1/2}\le\eta^{\sim}(\square,S)^{1/2}=e^{-\frac12\delta_P d(\tilde{\square},S)}$ at every block. Summing over all blocks and applying the two displays \eqref{5.2:eq-shell-block-sum-a} of Lemma~\ref{5.2:lem-shell-block-sum} gives \eqref{5.2:eq-proximity-comparison-1}.

Finally, by the first paragraph of the proof of (t1), every block that touches $S$ satisfies $\tilde{\square}\cap S\neq\emptyset$; the blocks touching $S$ thus form a subclass of the blocks whose fattened block $\tilde{\square}$ meets $S$.
\end{proof}

\begin{lemma}[The block sum at the support scale]\label{5.2:lem-support-block-sum}
Let $s\le 4$ and let $\operatorname{St}$ be a finite family of cells of $\mathfrak{B}_s$, with $S:=\bigcup\operatorname{St}$. Write $w:=M/M_1$. Assume the following hypotheses:
\begin{itemize}
\item in every label, wall- or corner-neighbouring cells differ in level by at most one;
\item $\operatorname{dist}_{j+1}(\Omega_j^c,\Omega_{j+1})\ge 1$ for every $j$;
\item the zones are nested: $\Omega_{j+1}\subset\Omega_j$ for every $j$;
\item a block of zone $n$ meets only cells of the levels $n-1,n,n+1$, and a cell meets at most $C_{\mathrm{blk}}:=2^{d+1}+(L+1)^d$ blocks;
\item for every integer $\mathsf{i}\ge0$, a contour $\Gamma$ of $\mathfrak{B}_s$-length $\operatorname{len}_{\mathfrak{B}_s}(\Gamma)\le\mathsf{i}w$ meets at most $c_{\boxplus}(\mathsf{i}+1)$ cells, and these cells form a family connected by touching;
\item $\tfrac1{16}\delta_P w\ge c_{\boxplus}\log\hat b_t+1$;
\item $M>6L$.
\end{itemize}
For a block $\square$ of the cover of $\mathfrak{B}_s$ recall the following.
\begin{itemize}
\item $\square^{\zeta}:=\operatorname{supp}\hat{\chi}_{\square}$.
\item $\square$ \emph{touches} $S$ if the $2$-stencil $\operatorname{Sten}_2(S)$ meets $\square^{\zeta}$.
\item $\eta^{\zeta}(\square,S):=1$ if $\square$ touches $S$, and $\eta^{\zeta}(\square,S):=e^{-\delta_P d(\square^{\zeta},S)}$ otherwise.
\end{itemize}
Then, summing over all blocks of the cover of $\mathfrak{B}_s$, those inside the hubs included,
\begin{equation}\label{5.2:eq-support-block-sum-1}
\sum_{\square}\eta^{\zeta}(\square,S)^{1/2}\le C_{\square}^{\zeta}\,\#\operatorname{St},
\qquad
C_{\square}^{\zeta}:=2C_{\mathrm{blk}}\,\hat b_t^{\,2c_{\boxplus}}\le C_{\square}^0 .
\end{equation}
Here $C_{\square}^0$ is the constant of \eqref{5.2:eq-shell-block-sum-a}.
\end{lemma}

\begin{proof}
Let $d(\cdot,\cdot)\ge d^{\wedge}(\cdot,\cdot)$ be the two contour distances of $\mathfrak{B}_s$, and let $w=M/M_1$ be the ratio fixed in the statement. To a block $\square$ assign the integer
\[
\mathsf{i}(\square):=
\begin{cases}
0 & \text{if $\square$ touches $S$,}\\
\lfloor d^{\wedge}(\square^{\zeta},S)/w\rfloor & \text{otherwise.}
\end{cases}
\]
Then $\eta^{\zeta}(\square,S)^{1/2}\le e^{-\frac12\delta_P w\,\mathsf{i}(\square)}$. At a touching block both sides of this inequality equal $1$. At any other block we have
\[
\eta^{\zeta}(\square,S)=e^{-\delta_P d(\square^{\zeta},S)}\le e^{-\delta_P d^{\wedge}(\square^{\zeta},S)},
\]
and $d^{\wedge}(\square^{\zeta},S)\ge w\,\mathsf{i}(\square)$.

\emph{Claim.} The block $\square$ meets a cell $\Delta''$ which is joined to a member of $\operatorname{St}$ by a chain of at most $c_{\boxplus}(\mathsf{i}(\square)+2)$ successively touching cells.

We argue in two cases.
\begin{itemize}
\item Let $\square$ touch $S$. Since $\operatorname{Sten}_2(S)$ meets $\square^{\zeta}$, there are a point $x\in\square^{\zeta}$ and a point $x_0\in S$, lying in a cell of some level $i$, with $|x-x_0|\le2\ell_i$, where $\ell_i$ is the lattice spacing at level $i$. By Lemma~\ref{5.2:lem-touching-block}, the block $\square$ meets every member of $\operatorname{St}$ containing $x_0$. Take $\Delta''$ to be such a member, so that $\Delta''\in\operatorname{St}$.
\item Otherwise put $\mathsf{i}:=\mathsf{i}(\square)$. Since $d^{\wedge}(\square^{\zeta},S)<(\mathsf{i}+1)w$, there is a contour $\Gamma$ from a point of $S$ to a point $z\in\square^{\zeta}$ of $\mathfrak{B}_s$-length $\operatorname{len}_{\mathfrak{B}_s}(\Gamma)<(\mathsf{i}+1)w$. By the contour--cell count assumed above, $\Gamma$ meets at most $c_{\boxplus}(\mathsf{i}+2)$ cells, and these form a family connected by touching, because a connected set is not covered by two disjoint closed sets. One of these cells is a member of $\operatorname{St}$, since $\Gamma$ starts in $S$. One of them, $\Delta''$, holds $z$. As $\square^{\zeta}\subset\operatorname{int}\square$, the point $z$ is interior to $\square$, so $\Delta''$ meets $\square$.
\end{itemize}
This proves the claim.

\emph{Count.} Fix an integer $\mathsf{i}\ge0$. A cell touches at most $\hat b_t$ other cells; this follows from the first hypothesis, the level condition on neighbouring cells. Hence the cells joined to $\operatorname{St}$ by chains of at most $c_{\boxplus}(\mathsf{i}+2)$ successively touching cells number at most
\[
\#\operatorname{St}\sum_{l<c_{\boxplus}(\mathsf{i}+2)}\hat b_t^{\,l}\le\#\operatorname{St}\,\hat b_t^{\,c_{\boxplus}(\mathsf{i}+2)} .
\]
By the block--cell incidence assumed above, a cell meets at most $C_{\mathrm{blk}}$ blocks. By the claim, therefore,
\[
\#\{\square:\ \mathsf{i}(\square)=\mathsf{i}\}\le C_{\mathrm{blk}}\,\#\operatorname{St}\,\hat b_t^{\,c_{\boxplus}(\mathsf{i}+2)} .
\]
Recall $\mathsf{q}_{\square}=\hat b_t^{\,c_{\boxplus}}e^{-\frac12\delta_P w}$ from \eqref{5.2:eq-rooted-sum-notation-a}. The hypothesis $\tfrac1{16}\delta_P w\ge c_{\boxplus}\log\hat b_t+1$ gives $\tfrac12\delta_P w\ge 8c_{\boxplus}\log\hat b_t+8$, hence $\mathsf{q}_{\square}\le e^{-8}$. Summing over the values of $\mathsf{i}(\square)$,
\begin{align*}
\sum_{\square}\eta^{\zeta}(\square,S)^{1/2}
&\le\sum_{\mathsf{i}\ge0}C_{\mathrm{blk}}\,\#\operatorname{St}\,\hat b_t^{\,c_{\boxplus}(\mathsf{i}+2)}\,e^{-\frac12\delta_P w\mathsf{i}}\\
&=C_{\mathrm{blk}}\,\hat b_t^{\,2c_{\boxplus}}\,\#\operatorname{St}\sum_{\mathsf{i}\ge0}\mathsf{q}_{\square}^{\,\mathsf{i}}\\
&\le 2C_{\mathrm{blk}}\,\hat b_t^{\,2c_{\boxplus}}\,\#\operatorname{St}.
\end{align*}

It remains to compare $C_{\square}^{\zeta}$ with $C_{\square}^0$. By the binomial expansion,
\[
(L+9)^d\ge(L+1)^d+8d(L+1)^{d-1}\ge(L+1)^d+2^{d+1}=C_{\mathrm{blk}}.
\]
Hence $C_{\square}^{\zeta}\le 2(L+9)^d\hat b_t^{\,2c_{\boxplus}}=C_{\square}^0$.
\end{proof}

\begin{remark}
\begin{enumerate}
\item[(1)] The proof above is the proof of the block sum in shells (Lemma~\ref{5.2:lem-shell-block-sum}) with $\tilde{\square}$ replaced by $\square^{\zeta}$. In that proof one step counts the zones over a single cell, of which there are at most $k-n''+2$, $n''$ being the level of the cell reached. Here this step is replaced by the three zones $n-1,n,n+1$ of the block-cell incidence, which are absorbed into $C_{\mathrm{blk}}$. Consequently the factor $k-n_-+3$, with $n_-:=\min_{\operatorname{St}}\operatorname{lev}$, of \eqref{5.2:eq-shell-block-sum-a} does not occur at $\square^{\zeta}$.
\item[(2)] The family $\operatorname{St}$ in Lemma~\ref{5.2:lem-support-block-sum} is arbitrary. For a compact pointed piece with exempt family $\operatorname{St}_v$ and read set $S_v=\bigcup\operatorname{St}_v$ one has $\eta_p=\eta^{\zeta}(\square,S_v)$, so the lemma gives the block sum at each bracket:
\[
\sum_{\square}\eta_p^{1/2}\le C_{\square}^{\zeta}\#\operatorname{St}_v\le C_{\square}^0\#\operatorname{St}_v\le\mathsf{b}_{\square}(\operatorname{St}_v).
\]
That sum nevertheless keeps Lemma~\ref{5.2:lem-shell-block-sum} as its majorant, through Corollary~\ref{5.2:cor-proximity-comparison}, with every constant unchanged. The count of touching blocks for long singles and tails, at most $C_{\mathrm{blk}}\#\operatorname{St}_v$, rests on the same block-cell incidence.
\end{enumerate}
\end{remark}

\begin{theorem}[The source-rooted sum over families of cut cells]\label{5.2:thm-rooted-sum}
Let the family $\pi^{\mathfrak{s}}$ of cells in proper scales disjoint with $\Lambda^{\sim}$, the
hubs with their side bound $h_\sharp$, the source cells, the adjacency relation with its bound
$\hat b$, the function $\rho$, the $\pi_k$-hull $[\cdot]_k$, the number $N_k$ of $\pi_k$-cubes of a
set, the distance $\operatorname{dist}_k$, the class $\mathbf{D}_k$ with its tree length $d_k$, the
constants $S_d$, $c_l^{\circ}$, $a=\tfrac{31}{20}\kappa$ and, for $r\le a+9$, the numbers
$\lambda_\Sigma$, $\lambda_\Sigma'$, $\epsilon_\partial$ be as in
Notation~\ref{5.2:not-rooted-sum-notation}, under the lower bounds on $\kappa_1$ and $\kappa$ stated
in \eqref{5.2:eq-rooted-sum-notation-a}. For a finite family $\operatorname{St}\subset\pi^{\mathfrak{s}}$
let $\partial\operatorname{St}$ be the set of members of $\pi^{\mathfrak{s}}\setminus\operatorname{St}$
sharing a wall with a member of $\operatorname{St}$; if $\operatorname{St}$ is wall-connected and holds
no source cell, put
\begin{equation*}
\mathfrak{Y}(\operatorname{St}):=\bigl\{\mathsf{y}\subset\pi^{\mathfrak{s}}\setminus\operatorname{St}
\ \text{finite}:\ \operatorname{St}\cup\mathsf{y}\ \text{connected},\ \mathsf{y}\ \text{holds a
source cell}\bigr\},
\end{equation*}
and, for $A\in\mathbf{D}_k$ with $[\operatorname{St}]_k\subset A$ and $\mathsf{y}$ with
$\operatorname{St}\cup\mathsf{y}$ connected,
$\operatorname{lab}(A,\mathsf{y}):=A\cup[\mathsf{y}]_k\cup\operatorname{Hub}(\mathsf{y})$, with
$\operatorname{Hub}(\mathsf{y})$ as in Notation~\ref{5.2:not-rooted-sum-notation}.

A \emph{unit datum} is a triple $\mathsf{u}=(\operatorname{St}_{\mathsf{u}},A_{\mathsf{u}},\mu_{\mathsf{u}})$
with $\operatorname{St}_{\mathsf{u}}\subset\pi^{\mathfrak{s}}$ finite and wall-connected,
$\operatorname{St}_{\mathsf{u}}\cup\partial\operatorname{St}_{\mathsf{u}}$ without source cell,
$A_{\mathsf{u}}\in\mathbf{D}_k$ with $[\operatorname{St}_{\mathsf{u}}]_k\subset A_{\mathsf{u}}$, and
$\mu_{\mathsf{u}}\ge0$. For a family $\mathbb{D}$ of unit data and $r\ge0$ let
$\mathsf{M}^{\partial}_r[\mathbb{D}]$ be defined by the first line of \eqref{5.2:eq-rooted-sum-a},
that is, the total weight of the unit data that have one given cell $\Delta_1$ on their rim. Then
for $r\le a+9$ and every $Q_*\in\pi_k$ the inequality in the second line holds:
\begin{equation}\label{5.2:eq-rooted-sum-a}
\begin{gathered}
\mathsf{M}^{\partial}_r[\mathbb{D}]:=\sup_{\Delta_1\in\pi^{\mathfrak{s}}}e^{-\rho(\Delta_1)}
\sum_{\mathsf{u}\in\mathbb{D}:\ \Delta_1\in\partial\operatorname{St}_{\mathsf{u}}}
\mu_{\mathsf{u}}\,e^{r\,d_k(A_{\mathsf{u}})}\,e^{\epsilon_\partial\#\partial\operatorname{St}_{\mathsf{u}}},
\\
\sum_{\mathsf{u}\in\mathbb{D}}\ \sum_{\substack{\mathsf{y}\in\mathfrak{Y}(\operatorname{St}_{\mathsf{u}}):\\
\operatorname{lab}(A_{\mathsf{u}},\mathsf{y})\ni Q_*}}
\mu_{\mathsf{u}}\,e^{-(\kappa_1-2)|\mathsf{y}|}\,e^{r\,d_k(\operatorname{lab}(A_{\mathsf{u}},\mathsf{y}))}
\le S_d\,\mathsf{M}^{\partial}_{r+1}[\mathbb{D}].
\end{gathered}
\end{equation}
None of the constants entering \eqref{5.2:eq-rooted-sum-a} ($S_d$, $c_l^{\circ}$, $\hat b$,
$\epsilon_\partial$) depends on $K$, $k$, the levels of the cells, the ages of large-field regions,
the number $N(g_k)$ of preparatory stages, the number $N_0$ of zones, $\check R_k$, the number or
thickness of the zones, the windows of arrested large-field operations or the number of such
operations, the hubs, or the volume.
\end{theorem}

\begin{proof}
All terms are non-negative, so enlarging a range of summation only increases a sum. Fix once and
for all an order on $\pi^{\mathfrak{s}}$. Throughout, $\mathsf{u}\in\mathbb{D}$ and
$\mathsf{y}\in\mathfrak{Y}(\operatorname{St}_{\mathsf{u}})$; we write
$\operatorname{lab}:=\operatorname{lab}(A_{\mathsf{u}},\mathsf{y})$.

\emph{(a) The label.} Every component of $\mathsf{y}$ has a member in
$\partial\operatorname{St}_{\mathsf{u}}$. Indeed, $\operatorname{St}_{\mathsf{u}}\cup\mathsf{y}$ is
connected, so from a member of a component $\mathsf{y}'$ of $\mathsf{y}$ there is a path of
successively adjacent members of $\operatorname{St}_{\mathsf{u}}\cup\mathsf{y}$ ending in
$\operatorname{St}_{\mathsf{u}}$. Members of distinct components of $\mathsf{y}$ are not adjacent,
so the path leaves $\mathsf{y}'$ only into $\operatorname{St}_{\mathsf{u}}$; that step is across a
wall, since adjacency through a hub joins two members sharing a wall with one hub, i.e.\ two source
cells, and $\operatorname{St}_{\mathsf{u}}$ holds no source cell. The last member of $\mathsf{y}'$
on the path therefore shares a wall with a member of $\operatorname{St}_{\mathsf{u}}$: it lies in
$\partial\operatorname{St}_{\mathsf{u}}$.

Next, $\operatorname{lab}\in\mathbf{D}_k$. The $\pi_k$-cubes holding two wall-neighbouring members
coincide or share a wall; the edges of the fixed spanning forest of $\mathsf{y}$ which join two
members sharing no wall are bridged by the chains forming $\operatorname{Hub}(\mathsf{y})$; each
component of $\mathsf{y}$ has, by the preceding paragraph, a member sharing a wall with a member of
$\operatorname{St}_{\mathsf{u}}$, whose $\pi_k$-cube lies in
$[\operatorname{St}_{\mathsf{u}}]_k\subset A_{\mathsf{u}}$; and $A_{\mathsf{u}}\in\mathbf{D}_k$.
The forest has fewer than $|\mathsf{y}|$ edges; the two ends of a bridged edge are source cells at
one hub, hence at most $h_\sharp$ apart, so by the wall-chain bound of
Notation~\ref{5.2:not-rooted-sum-notation} (two $\pi_k$-cubes at distance $x$ are joined by a
wall-chain with fewer than $d(x+1)$ intermediate cubes) each such edge contributes fewer than
$d(h_\sharp+1)$ cubes to $\operatorname{Hub}(\mathsf{y})$. Since $[\mathsf{y}]_k$ has at most
$|\mathsf{y}|$ cubes,
$N_k(\operatorname{lab}\setminus A_{\mathsf{u}})\le(1+d(h_\sharp+1))|\mathsf{y}|
=c_l^{\circ}|\mathsf{y}|$, and the growth bound of the tree length in
Notation~\ref{5.2:not-rooted-sum-notation}, $d_k(Y)\le d_k(X)+N_k(Y\setminus X)$ for $X\subset Y$
in $\mathbf{D}_k$, applied to $A_{\mathsf{u}}\subset\operatorname{lab}$, gives
\begin{equation}\label{5.2:eq-rooted-sum-2}
d_k(\operatorname{lab}(A_{\mathsf{u}},\mathsf{y}))\le d_k(A_{\mathsf{u}})+c_l^{\circ}|\mathsf{y}|.
\end{equation}

\emph{(b) The root.} Let $\Delta_{\mathfrak{F}}$ be the first source cell of $\mathsf{y}$; it is a
$\pi_k$-cube of $\operatorname{lab}$. If $Q_*\in\operatorname{lab}$, a tree graph realising
$d_k(\operatorname{lab})$ intersects both $Q_*$ and $\Delta_{\mathfrak{F}}$, so, since a tree graph
meeting two cubes is at least as long as their distance
(Notation~\ref{5.2:not-rooted-sum-notation}),
$d_k(\operatorname{lab})\ge\operatorname{dist}_k(Q_*,\Delta_{\mathfrak{F}})$. Writing
$r\,d_k=(r+1)d_k-d_k$ and using \eqref{5.2:eq-rooted-sum-2},
\begin{equation*}
e^{r\,d_k(\operatorname{lab})}\le e^{-\operatorname{dist}_k(Q_*,\Delta_{\mathfrak{F}})}\,
e^{(r+1)d_k(A_{\mathsf{u}})}\,e^{c_l^{\circ}(r+1)|\mathsf{y}|}.
\end{equation*}
Since $\kappa_1-2-c_l^{\circ}(r+1)=\lambda_\Sigma$, and since the source cells are distinct
$\pi_k$-cubes with $\sum_{Q\in\pi_k}e^{-\operatorname{dist}_k(Q_*,Q)}\le S_d$, the left member of
the inequality in \eqref{5.2:eq-rooted-sum-a} is at most
$S_d\sup_{\Delta_{\mathfrak{F}}}T(\Delta_{\mathfrak{F}})$,
where, dropping the conditions $Q_*\in\operatorname{lab}$ and ``first'',
\begin{equation*}
T(\Delta_{\mathfrak{F}}):=\sum_{\mathsf{u}\in\mathbb{D}}\mu_{\mathsf{u}}\,
e^{(r+1)d_k(A_{\mathsf{u}})}\sum_{\substack{\mathsf{y}\in\mathfrak{Y}(\operatorname{St}_{\mathsf{u}}):\\
\mathsf{y}\ni\Delta_{\mathfrak{F}}}}e^{-\lambda_\Sigma|\mathsf{y}|}.
\end{equation*}

\emph{(c) Trunk and rim.} Let $\mathsf{y}_1$ be the component of $\mathsf{y}$ containing
$\Delta_{\mathfrak{F}}$. The other components form a set of pairwise disjoint connected families,
each with a member in $\partial\operatorname{St}_{\mathsf{u}}$ by (a); they need not contain a
source cell, and nothing below uses one. For $\Delta\in\pi^{\mathfrak{s}}$ let $\varphi(\Delta)$ be
the sum of $e^{-\lambda_\Sigma|\mathsf{y}_2|}$ over the connected families
$\mathsf{y}_2\ni\Delta$. By the bound of Notation~\ref{5.2:not-rooted-sum-notation} on connected
families (at most $(e\hat b)^n$ of them have $n$ members and pass through a given member) and
$\lambda_\Sigma'=\lambda_\Sigma-\log(e\hat b)$,
\begin{equation*}
\varphi(\Delta)\le\sum_{n\ge1}(e\hat b)^n e^{-\lambda_\Sigma n}
=\sum_{n\ge1}e^{-\lambda_\Sigma' n}\le 2e^{-\lambda_\Sigma'}=\epsilon_\partial,
\end{equation*}
because $\lambda_\Sigma'\ge3$ for $r\le a+9$ by \eqref{5.2:eq-rooted-sum-notation-a}, so that
$e^{-\lambda_\Sigma'}\le\tfrac12$. File each of the other components at its first member in
$\partial\operatorname{St}_{\mathsf{u}}$; disjoint families have distinct such members. Hence the
sum over the other components of the product of their weights is at most
\begin{equation*}
\prod_{\Delta\in\partial\operatorname{St}_{\mathsf{u}}}(1+\varphi(\Delta))
\le e^{\epsilon_\partial\#\partial\operatorname{St}_{\mathsf{u}}}.
\end{equation*}

\emph{(d) Interchange.} By (a), $\mathsf{y}_1$ has a member
$\Delta_1\in\partial\operatorname{St}_{\mathsf{u}}$, and $\rho(\Delta_1)\le|\mathsf{y}_1|$ because
$\mathsf{y}_1$ is a connected family holding $\Delta_1$ and a source cell. By (c), inserting
$1\le e^{|\mathsf{y}_1|}e^{-\rho(\Delta_1)}$ and summing over $\Delta_1\in\mathsf{y}_1$ instead of
choosing one,
\begin{equation*}
\begin{split}
T(\Delta_{\mathfrak{F}})&\le\sum_{\substack{\mathsf{y}_1\ni\Delta_{\mathfrak{F}}\\ \text{connected}}}
e^{-\lambda_\Sigma|\mathsf{y}_1|}\sum_{\Delta_1\in\mathsf{y}_1}
\sum_{\mathsf{u}:\ \Delta_1\in\partial\operatorname{St}_{\mathsf{u}}}\mu_{\mathsf{u}}\,
e^{(r+1)d_k(A_{\mathsf{u}})}\,e^{\epsilon_\partial\#\partial\operatorname{St}_{\mathsf{u}}}\\
&\le\mathsf{M}^{\partial}_{r+1}[\mathbb{D}]\sum_{n\ge1}(e\hat b)^n e^{-\lambda_\Sigma n}\,n\,e^{n}
=\mathsf{M}^{\partial}_{r+1}[\mathbb{D}]\sum_{n\ge1}n\,e^{-(\lambda_\Sigma'-1)n}
\le\mathsf{M}^{\partial}_{r+1}[\mathbb{D}],
\end{split}
\end{equation*}
where the same bound on connected families counts the connected $\mathsf{y}_1$ with $n$ members
through $\Delta_{\mathfrak{F}}$, and
the last step uses $\lambda_\Sigma'-1\ge2$, so that
$\sum_{n\ge1}n e^{-2n}=e^{-2}(1-e^{-2})^{-2}<1$. Together with (b) this is the inequality in
\eqref{5.2:eq-rooted-sum-a}.

The fine cells are thus reached from the source: a connected family of $n$ members of
$\pi^{\mathfrak{s}}$ through a given level-$k$ cell is one of at most $(e\hat b)^n$, whatever levels
it descends to, and the unit datum is then one of those with a given cell on their rim. No count of
the cells of a given level inside a $\pi_k$-cube enters anywhere.
\end{proof}

\begin{lemma}[Rim mass of the compact pointed units]\label{5.2:lem-rim-mass}
Let $v$ run over the compact pointed units of the classes (c) and (d) of the carrier bound for
exterior pointed units (Lemma~\ref{5.2:lem-carrier-exterior}): short singles, Wilson pieces and
cores. A rough core is read termwise, so that its constituents enter as short singles or as long
units. For such $v$ let $\mathfrak{N}_{\mathrm{dil}}(v)$ denote its dilated size, that is, the size
of $v$ read on the $(9/8)$-dilated carrier of part (c3) of Lemma~\ref{5.2:lem-carrier-exterior}, so
that $\mathfrak{N}_{\mathrm{dil}}(v)\le\tfrac{81}{64}\mathfrak{N}^{\flat}(v)$. Let
$\mathbb{D}^{\mathrm{sh}}$ be the family of triples
$(\operatorname{St}(\Delta_{y_v}),\mathfrak{h}_v,\mu_v)$ with
\begin{equation*}
\mu_v:=\mathsf{b}_{\square}(\operatorname{St}_v)\,\mathsf{f}_v^{1/2}\,\mathfrak{N}_{\mathrm{dil}}(v),
\end{equation*}
$\mathsf{b}_{\square}$ as in \eqref{5.2:eq-shell-block-sum-a}. The rim mass
\eqref{5.2:eq-rooted-sum-a} is formed for $\mathbb{D}^{\mathrm{sh}}$ as for a family of unit data;
only its definition is used below. Here, for a cell $\Delta'$, $\mathsf{f}(\Delta')$ and
$\bar{\mathsf{w}}(\Delta')=3w^{\vee}(\Delta')$ denote the common values of $\mathsf{f}_v$ and
$\bar{\mathsf{w}}_v$ over the units $v$ with $y_v\in\Delta'$ (by hypothesis (i) below these depend
on $v$ through $\Delta_{y_v}$ alone). Assume:
\begin{enumerate}
\item[(i)] the read set of a compact pointed unit is its star: $S_v=\mathfrak{r}_v$ and
$\operatorname{St}_v=\operatorname{St}(\Delta_{y_v})$, where
$\mathfrak{r}_v=\bigcup\operatorname{St}(\Delta_{y_v})$;
\item[(ii)] the column bound per cell for exposed units (Lemma~\ref{5.2:lem-exposed-column},
display \eqref{5.2:eq-exposed-column-a}), together with the fact that the dilated sizes of the
units with $y_v$ in a cell sum to at most $\tfrac32$ times its right member, because
$(9/8)^3\le 3/2$;
\item[(iii)] the bound \eqref{5.2:eq-carrier-exterior-e} of
Lemma~\ref{5.2:lem-carrier-exterior}: if, for a cell $\Delta'$ with $n':=\operatorname{lev}\Delta'$,
the value $w^{\vee}(\Delta')>1$ is attained on a foreign live arrest $\mathfrak{a}'$, then
\begin{gather*}
n'\le k_{a'},\qquad
\bar{\mathsf{w}}(\Delta')\le 3\mathsf{L}_0^2LC_{\uparrow}\check{R}_{k_{a'}},\\
d_{\mathrm{sc}}(y,\Lambda^{\sim})\ge 2\check{R}_{k_{a'}}-6
\quad\text{for every }y\in\textstyle\bigcup\operatorname{St}(\Delta'),
\end{gather*}
where $\bigcup\operatorname{St}(\Delta')$ is the common read set, by (i), of the units $v$ with
$y_v\in\Delta'$; whence $\mathsf{f}(\Delta')^{1/2}\bar{\mathsf{w}}(\Delta')^3\le\mathsf{C}_{\mathrm{fw}}$
with $\mathsf{C}_{\mathrm{fw}}$ as in \eqref{5.2:eq-carrier-exterior-f}; otherwise
$\bar{\mathsf{w}}(\Delta')=3$, and $27c_\chi^{1/2}\le\mathsf{C}_{\mathrm{fw}}$.
\end{enumerate}
Then for every cell $\Delta'$
\begin{equation}\label{5.2:eq-rim-mass-a}
\sum_{v\in\mathbb{D}^{\mathrm{sh}}:\,y_v\in\Delta'}\mathsf{f}_v^{1/2}\,\mathfrak{N}_{\mathrm{dil}}(v)
\;\le\;\tfrac32\,\mathsf{f}(\Delta')^{1/2}\,\mathsf{K}^{\flat}\,\bar{\mathsf{w}}(\Delta')^3
\;\le\;\tfrac32\,\mathsf{C}_{\mathrm{fw}}\,\mathsf{K}^{\flat},
\end{equation}
and, for every $r\ge 0$, the rim mass \eqref{5.2:eq-rooted-sum-a} of $\mathbb{D}^{\mathrm{sh}}$
obeys
\begin{equation}\label{5.2:eq-rim-mass-1}
\mathsf{M}^{\partial}_r[\mathbb{D}^{\mathrm{sh}}]\;\le\;\mathsf{M}^{\mathrm{sh}}_r
:=N_\partial\,e^{1+rc_{\mathfrak{h}}}\,C_{\square}^0\,c_S(c_S+3)\,
\tfrac32\,\mathsf{C}_{\mathrm{fw}}\,\mathsf{K}^{\flat}.
\end{equation}
\end{lemma}

\begin{proof}
Fix a cell $\Delta'$ and put $n':=\operatorname{lev}\Delta'$. By hypothesis (i), every compact
pointed unit $v$ with $y_v\in\Delta'$ has $S_v=\bigcup\operatorname{St}(\Delta')$ and
$\operatorname{St}_v=\operatorname{St}(\Delta')$.

\emph{The distance clause of \eqref{5.2:eq-carrier-exterior-e} at the read set of (i).}
Suppose that $w^{\vee}(\Delta')>1$ is attained on a foreign live arrest $\mathfrak{a}'$, that is,
a live arrest on a component of $\Lambda_k^c$ which does not belong to $Z$, with component
$X_{\mathfrak{a}'}$, live piece $W_{\mathfrak{a}'}$ and collar $\mathfrak{c}_{\mathfrak{a}'}$.
Let $\Delta^{*}\subset W_{\mathfrak{a}'}$ be the weighted cell, at $\operatorname{dist}_{n'}\le 1$
from $\Delta'$. Then $n'\le k_{a'}$, by the window statement for weighted cells: a weighted cell at
$\operatorname{dist}_{n'}\le 1$ from $\Delta'$ lies in an attempt's window, and its zone index $n'$
is at most that attempt's step. This is also the remark used for the regularity bound in the proof
of Lemma~\ref{5.2:lem-carrier-exterior}. No member of $\operatorname{St}(\Delta')$ has level
$k_{a'}+1$: the collar $\mathfrak{c}_{\mathfrak{a}'}$, which by the collar statement for live
arrests has level $k_{a'}$ and consists of two layers of $M\check{R}_{k_{a'}}$-cubes, is
$2\check{R}_{k_{a'}}\ge 76$ level-$k_{a'}$ units thick about $W_{\mathfrak{a}'}$.

By clause (e5) of Lemma~\ref{5.2:lem-exempt-star}, $\bigcup\operatorname{St}(\Delta')$ lies in
$\{\operatorname{dist}_{n'}(\cdot,\Delta')\le L+2\}$, and in
$\{\operatorname{dist}_{n'}(\cdot,\Delta')\le 2\}$ if no member has level $n'+1$. In the latter
case a point of $\bigcup\operatorname{St}(\Delta')$ lies within $2+1+1=4$ level-$n'$ units of
$\Delta^{*}$: two from $\Delta'$, one across $\Delta'$, and one from $\Delta'$ to $\Delta^{*}$. If
a member has level $n'+1\le k_{a'}$, a point of $\bigcup\operatorname{St}(\Delta')$ lies within
$(L+4)/L\le 2$ level-$k_{a'}$ units of $\Delta^{*}$. In either case, since $n'\le k_{a'}$, every
point of $\bigcup\operatorname{St}(\Delta')$ lies within $5$ level-$k_{a'}$ units of $\Delta^{*}$.

Next, $\Lambda^{\sim}\cap X_{\mathfrak{a}'}=\emptyset$: $\Lambda^{\sim}$ lies in the components of
$Z$, while $X_{\mathfrak{a}'}$ lies in a component of $\Lambda_k^c$ that does not belong to $Z$. By
the contour statement for live arrests, a contour joining a point of $X_{\mathfrak{a}'}^c$ to the
live piece $W_{\mathfrak{a}'}$ crosses the whole collar $\mathfrak{c}_{\mathfrak{a}'}$, that is,
all its $2\check{R}_{k_{a'}}$ level-$k_{a'}$ layers. Since every point of
$\bigcup\operatorname{St}(\Delta')$ lies within five level-$k_{a'}$ units of
$\Delta^{*}\subset W_{\mathfrak{a}'}$, such a point lies in $X_{\mathfrak{a}'}$ and at most five
layers into the collar, while $\Lambda^{\sim}\subset X_{\mathfrak{a}'}^c$. Hence a contour from a
point of $\bigcup\operatorname{St}(\Delta')$ to $\Lambda^{\sim}$ still has at least
$2\check{R}_{k_{a'}}-5\ge 2\check{R}_{k_{a'}}-6$ own-scale cubes of the collar to cross, that is,
$d_{\mathrm{sc}}(y,\Lambda^{\sim})\ge 2\check{R}_{k_{a'}}-6$ for
$y\in\bigcup\operatorname{St}(\Delta')$. This is the distance clause of
\eqref{5.2:eq-carrier-exterior-e} at $S_v=\bigcup\operatorname{St}(\Delta')$.
The majorant of the weight in \eqref{5.2:eq-carrier-exterior-e} does not involve $S_v$.

The units concerned belong to (d), whose source set is $\Lambda^{\sim}$; the units of (c) meet no
foreign window. From the distance clause, $d_{\mathrm{sc}}(y,\Lambda^{\sim})-1\ge\check{R}_{k_{a'}}-6$
for $y\in\bigcup\operatorname{St}(\Delta')$ (as $\check{R}_{k_{a'}}\ge 1$), so
$\mathsf{f}(y)\le e^{-\delta_{\natural}(\check{R}_{k_{a'}}-6)_+}$ there; $\mathsf{f}(\Delta')$ being
the supremum of $\min\{1,c_\chi\mathsf{f}(y)\}$ over $\bigcup\operatorname{St}(\Delta')$, this gives
$\mathsf{f}(\Delta')\le c_\chi e^{-\delta_{\natural}(\check{R}_{k_{a'}}-6)_+}$. From the weight
clause, $\bar{\mathsf{w}}(\Delta')^3\le(3\mathsf{L}_0^2LC_{\uparrow})^3\check{R}_{k_{a'}}^3$.
Put $\varrho:=\check{R}_{k_{a'}}-5\ge 0$. Since $\check{R}_{k_{a'}}=\varrho+5\le\varrho+6$,
$(\check{R}_{k_{a'}}-6)_+=(\varrho-1)_+$ and $2\check{R}_{k_{a'}}-6=2\varrho+4\ge(\varrho-1)_+$,
\begin{align*}
\check{R}_{k_{a'}}^3\,e^{-\frac12\delta_{\natural}(\check{R}_{k_{a'}}-6)_+}
&\;\le\;(\varrho+6)^3\,e^{-\frac12\delta_{\natural}(\varrho-1)_+},\\
\check{R}_{k_{a'}}^3\,e^{-\frac12\delta_{\natural}(2\check{R}_{k_{a'}}-6)}
&\;\le\;(\varrho+6)^3\,e^{-\frac12\delta_{\natural}(\varrho-1)_+},
\end{align*}
and the right member is the quantity whose supremum over $\varrho\ge 0$ enters
$\mathsf{C}_{\mathrm{fw}}$ in \eqref{5.2:eq-carrier-exterior-f}. Hence
\begin{equation*}
\mathsf{f}(\Delta')^{1/2}\bar{\mathsf{w}}(\Delta')^3
\le c_\chi^{1/2}(3\mathsf{L}_0^2LC_{\uparrow})^3\,\check{R}_{k_{a'}}^3\,
e^{-\frac12\delta_{\natural}(\check{R}_{k_{a'}}-6)_+}
\le\mathsf{C}_{\mathrm{fw}},
\end{equation*}
which is the last clause of (iii) at this read set. The same clause covers a cell
$\Delta'$ of the collar itself ($n'=k_{a'}$, $w(\Delta')=1$, $\bar{\mathsf{w}}(\Delta')>3$): a
contour from it leaves $X_{\mathfrak{a}'}$ through the rest of the collar.

\emph{Proof of \eqref{5.2:eq-rim-mass-a}.} By (i), $\mathsf{f}_v=\mathsf{f}(\Delta')$ and
$\bar{\mathsf{w}}_v=\bar{\mathsf{w}}(\Delta')$ for all $v$ with $y_v\in\Delta'$, so
$\mathsf{f}(\Delta')^{1/2}$ may be taken out of the sum. The remaining sum of dilated sizes runs
over a subfamily of the units summed in \eqref{5.2:eq-exposed-column-a}. Take that display at
$r=0$: the weight $e^{rd_k(K_v)}$ is then $1$, and $e^{\varsigma_1\operatorname{rad}(v)}\ge 1$ may be
dropped. By hypothesis (ii) the sum of dilated sizes is at most
$\tfrac32\mathsf{K}^{\flat}\bar{\mathsf{w}}(\Delta')^3$, which gives the first inequality.

For the second inequality, consider first the case where $w^{\vee}(\Delta')>1$ is attained on a
foreign live arrest. Then hypothesis (iii), whose distance clause was established above at the
read set of (i), gives $\mathsf{f}(\Delta')^{1/2}\bar{\mathsf{w}}(\Delta')^3\le\mathsf{C}_{\mathrm{fw}}$.
Otherwise $\bar{\mathsf{w}}(\Delta')=3$ and $\mathsf{f}(\Delta')\le 1$, so, since $c_\chi\ge 1$,
\begin{equation*}
\mathsf{f}(\Delta')^{1/2}\bar{\mathsf{w}}(\Delta')^3\le 27\le 27c_\chi^{1/2}\le\mathsf{C}_{\mathrm{fw}}.
\end{equation*}

\emph{Proof of \eqref{5.2:eq-rim-mass-1}.} By the definition \eqref{5.2:eq-rooted-sum-a}, with
$\operatorname{St}_{\mathsf{u}}=\operatorname{St}(\Delta_{y_v})$ and $A_{\mathsf{u}}=\mathfrak{h}_v$,
\begin{equation*}
\mathsf{M}^{\partial}_r[\mathbb{D}^{\mathrm{sh}}]
=\sup_{\Delta_1\in\pi^{\mathfrak{s}}}e^{-\rho(\Delta_1)}
\sum_{v:\,\Delta_1\in\partial\operatorname{St}(\Delta_{y_v})}
\mu_v\,e^{rd_k(\mathfrak{h}_v)}\,e^{\epsilon_\partial\#\partial\operatorname{St}(\Delta_{y_v})}.
\end{equation*}
We bound the factors one by one.
\begin{itemize}
\item Since $\mathfrak{h}_v$ is a connected union of at most $c_{\mathfrak{h}}=7^d$ cubes of
$\pi_k$, the growth bound $d_k(Y)\le d_k(X)+N_k(Y\setminus X)$ for $X\subset Y$ in
$\mathbf{D}_k$, applied with $X$ a single cube of $\mathfrak{h}_v$ ($d_k(X)=0$), gives
$d_k(\mathfrak{h}_v)\le c_{\mathfrak{h}}-1$, hence $e^{rd_k(\mathfrak{h}_v)}\le e^{rc_{\mathfrak{h}}}$.
\item By \eqref{5.2:eq-shell-block-sum-a},
$\mathsf{b}_{\square}(\operatorname{St}_v)
=C_{\square}^0\#\operatorname{St}_v(\#\operatorname{St}_v+3)$. Clause (e2) of
Lemma~\ref{5.2:lem-exempt-star} gives $\#\operatorname{St}(\Delta_{y_v})\le c_S$, so
$\mathsf{b}_{\square}(\operatorname{St}_v)\le C_{\square}^0c_S(c_S+3)$.
\item Clause (e3) of that lemma gives $\epsilon_\partial\#\partial\operatorname{St}(\Delta_{y_v})\le 1$,
so the rim factor is at most $e$.
\item Since $\rho\ge 0$, $e^{-\rho(\Delta_1)}\le 1$.
\end{itemize}

Now group the units $v$ according to the cell $\Delta'=\Delta_{y_v}$. By clause (e4) of
Lemma~\ref{5.2:lem-exempt-star}, a given cell $\Delta_1$ lies in $\partial\operatorname{St}(\Delta')$
for at most $N_\partial$ cells $\Delta'$. For each such $\Delta'$, the remaining sum
$\sum_{y_v\in\Delta'}\mathsf{f}_v^{1/2}\mathfrak{N}_{\mathrm{dil}}(v)$ is at most
$\tfrac32\mathsf{C}_{\mathrm{fw}}\mathsf{K}^{\flat}$ by \eqref{5.2:eq-rim-mass-a}. Collecting the
factors,
\begin{equation*}
\mathsf{M}^{\partial}_r[\mathbb{D}^{\mathrm{sh}}]\le N_\partial\cdot C_{\square}^0c_S(c_S+3)\cdot
e^{rc_{\mathfrak{h}}}\cdot e\cdot\tfrac32\mathsf{C}_{\mathrm{fw}}\mathsf{K}^{\flat}
=\mathsf{M}^{\mathrm{sh}}_r. \qedhere
\end{equation*}
\end{proof}

\begin{corollary}[The sum of pointed pieces behind one coarse anchor]\label{5.2:cor-anchor-sum}
Let $v$ run over the compact pointed units of the exterior class $\mathfrak{L}^{Y_0}$. These are the
short singles and the Wilson pieces of plain units of class (c) (one bracket, $s=4$), and the short
singles, the Wilson pieces and the far cores of plain units of class (d) ($X_v\cap Z=\emptyset$,
brackets $s=4,\dots,1$), a rough core being treated term by term through its short constituents. For
such a unit, $p=(\square,\mathsf{y};s)$ runs over its pieces, $\Delta':=\Delta_{y_v}$ is its anchor
cell and $n':=\operatorname{lev}\Delta'$. Assume
the following.
\begin{itemize}
\item The amplitude-disc bound of Lemma~\ref{5.2:lem-carrier-exterior} (the carrier bound for exterior
pointed units) for the pieces at the localisation site, valid on the disc of radius $\rho_A$ of
\eqref{5.2:eq-carrier-exterior-a}. It states
$|v_p|\le\mathfrak{N}^{\mathrm{ex}}_A(p)$, with $\mathfrak{N}^{\mathrm{ex}}_A(p)$ as in
\eqref{5.2:eq-carrier-exterior-d}. Here $\eta_p=1$ if the $2$-stencil of $S_v$ meets
$\square^{\zeta}$, and $\eta_p=e^{-\delta_P d(\square^{\zeta},S_v)}$ otherwise. Finally
$\mathsf{w}_L=1/\theta'(g_k)$.
\item The column bound per cell for exposed units (Lemma~\ref{5.2:lem-exposed-column}). The dilated
sizes $\tfrac{81}{64}\mathfrak{N}^{\flat}(v)$ (for a far core, with its depth parameter taken at
$\tfrac98\bar{\mathsf{w}}_v$) sum to at most $\tfrac32$ times its right member, since the dilation
enters at most through its third power and $(9/8)^3\le\tfrac32$.
\item The bound for a weight attained on a foreign window, proved in
Lemma~\ref{5.2:lem-carrier-exterior}. If $w^{\vee}(\Delta')>1$ is attained on the window of a live
arrest $\mathfrak{a}'$ on an untreated component, then three things hold:
\begin{equation*}
n'\le k_{a'},\qquad
d_{\mathrm{sc}}(y,\Lambda^{\sim})\ge 2\check{R}_{k_{a'}}-6 \text{ on } S_v,\qquad
\bar{\mathsf{w}}(\Delta')\le 3\mathsf{L}_0^2LC_{\uparrow}\check{R}_{k_{a'}}.
\end{equation*}
Hence $\mathsf{f}(\Delta')^{1/2}\bar{\mathsf{w}}(\Delta')^3\le\mathsf{C}_{\mathrm{fw}}$, with
$\mathsf{C}_{\mathrm{fw}}$ as in \eqref{5.2:eq-carrier-exterior-f}, where
$\mathsf{f}(\Delta'):=\mathsf{f}_v$, which by the second convention below depends on $v$ only through
$\Delta'$. Elsewhere $\bar{\mathsf{w}}=3$
and $27c_\chi^{1/2}\le\mathsf{C}_{\mathrm{fw}}$.
\item The site geometry at the localisation site and its floors, as fixed in the notation for the
source-rooted sum over cut cells (Notation~\ref{5.2:not-rooted-sum-notation}; see
\eqref{5.2:eq-rooted-sum-notation-a}).
\item The decoupling lemma for an arbitrary exempt subfamily of cells: with the decoupling parameters
set to zero off the cut cells $\mathsf{y}$, on a connected component of the decoupled domain of a
piece (the union of its exempt star, the hubs and the cells of $\mathsf{y}$) the fields
of the expansion depend only on the decoupling parameters and the data on that component, and they
vanish identically on a component that carries no datum.
\item The locality of the decoupled propagators: the kernels of the operators of the decoupled
expansion at parameters $\mathsf{s}$ vanish unless both arguments lie in one connected component of
the decoupled domain.
\item Three conventions on the expansion.
  \begin{itemize}
  \item The cells available for cutting are the members of
  $\pi^{\mathfrak{s}}\setminus\operatorname{St}_v$, cells in proper scales; the cut cells of $p$ form
  the finite family $\mathsf{y}\subset\pi^{\mathfrak{s}}\setminus\operatorname{St}_v$, and
  $m(p)=|\mathsf{y}|$.
  \item For a compact unit, $\operatorname{St}_v:=\operatorname{St}(\Delta_{y_v})$ and
  $S_v:=\bigcup\operatorname{St}_v=\mathfrak{r}_v$.
  \item $K_p:=[\operatorname{St}_v\cup\mathsf{y}]_k\cup\operatorname{Hub}(\mathsf{y})$ and
  $\operatorname{lab}_1(p):=K_p\cup\mathfrak{h}_v$; this equals
  $\operatorname{lab}(\mathfrak{h}_v,\mathsf{y})$, since
  $[\operatorname{St}_v]_k\subset\mathfrak{h}_v$ by the clause on boxes of
  Lemma~\ref{5.2:lem-exempt-star}.
  \end{itemize}
\end{itemize}
The block sums are supplied by the block sum in shells around a family of cells
(Lemma~\ref{5.2:lem-shell-block-sum}).

Then for every $K$, every $k\le K$, every label, every $r\le a+9$, with $a=\tfrac{31}{20}\kappa$ as in
Notation~\ref{5.2:not-rooted-sum-notation}, and every $Q_*\in\pi_k$:
\begin{equation}\label{5.2:eq-anchor-sum-a}
\begin{gathered}
\sum_{p\ \text{compact pointed of (c), (d)}:\ \operatorname{lab}_1(p)\ni Q_*}
\mathfrak{N}^{\mathrm{ex}}_A(p)\,e^{r\,d_k(\operatorname{lab}_1(p))}
\le\theta'(g_k)\,\mathsf{C}^{\Sigma}_r,\\
\mathsf{C}^{\Sigma}_r:=8S_d\,\mathsf{M}^{\mathrm{sh}}_{r+1}
=8S_dN_{\partial}\,e^{1+(r+1)c_{\mathfrak{h}}}\,C_{\square}^0\,c_S(c_S+3)\,
\tfrac32\,\mathsf{C}_{\mathrm{fw}}\mathsf{K}^{\flat}.
\end{gathered}
\end{equation}
The sum runs over the units $v$ of every level $j$ and every zone $n'\le k$. These units lie on
$\Omega_k$, on the shells of untreated components, or in or beside the frozen windows of foreign live
arrests. The sum also runs over $s$, $\square$ and $\mathsf{y}$.

The constant is a product of numbers depending on $(d,L,\kappa)$ with $\mathsf{C}_{\mathrm{fw}}$ and
$\mathsf{K}^{\flat}$. It depends on none of $K$, $k$, an age, $N$, $N_0$, $\check{R}$, a window, a
number of arrests, or the volume.

Moreover, every term has
\begin{equation*}
|\mathsf{y}|\ge\max\{k-n'-1,\ 9\check{R}_k-8\},
\end{equation*}
and the left member is in fact at most $e^{-(9\check{R}_k-8)}\theta'(g_k)\mathsf{C}^{\Sigma}_r$.
This sharper bound is not used in what follows.
\end{corollary}

\begin{proof}
\emph{Which walks occur.} Let $v$ be a compact pointed unit and $\Delta':=\Delta_{y_v}$.
\begin{itemize}
\item If $v$ is of class (d), then $y_v\in X_v\subset Z^c$, so $\Delta'\not\subset Z$.
\item If $v$ is of class (c), then $y_v$ lies within $2$ level-$k$ units of $\partial C$, where $C$ is
the treated component of $Z$ (as in Lemma~\ref{5.2:lem-carrier-exterior}); since
$\partial C\subset\partial Z$, it lies within $2$ level-$k$ units of $\partial Z$.
This is the geometry of units meeting $C$ in Lemma~\ref{5.2:lem-carrier-exterior}.
\end{itemize}
In both cases the exempt star of a cell (Lemma~\ref{5.2:lem-exempt-star}), in its last clause, shows
that $\operatorname{St}_v\cup\partial\operatorname{St}_v$ holds no source cell. Hence
Lemma~\ref{5.2:lem-piece-source-cell} applies. It is proved from the decoupling lemma and the locality
of the kernels listed among the hypotheses, and it gives $v_p=0$ unless
$\mathsf{y}\in\mathfrak{Y}(\operatorname{St}_v)$. Only such walks therefore occur.

The bound $\mathfrak{N}^{\mathrm{ex}}_A(p)$ is nonnegative. Summing it over all blocks $\square$ of
the cover of $\mathfrak{B}_s$, rather than over those of the pieces present, only adds terms.

\emph{Splitting the walk factor.} By the first convention, $m(p)=|\mathsf{y}|$. We write
\begin{equation*}
e^{-(\kappa_1-1)|\mathsf{y}|}=e^{-|\mathsf{y}|}\,e^{-(\kappa_1-2)|\mathsf{y}|}.
\end{equation*}

\emph{The block sum.} Neither $\operatorname{lab}_1(p)=\operatorname{lab}(\mathfrak{h}_v,\mathsf{y})$
nor $\mathsf{f}_v$ depends on $\square$.

Next, $\eta_p^{1/2}\le e^{-\frac12\delta_P d(\tilde{\square},S_v)}$ at every block:
\begin{itemize}
\item At a receding block, $\square^{\zeta}\subset\tilde{\square}$, so
$d(\square^{\zeta},S_v)\ge d(\tilde{\square},S_v)$.
\item At a touching block, $\eta_p=1$ and $d(\tilde{\square},S_v)=0$, by the comparison of the two
block-proximity distances (Corollary~\ref{5.2:cor-proximity-comparison}).
\end{itemize}

By the second convention, $S_v=\bigcup\operatorname{St}_v$. The second display of
Lemma~\ref{5.2:lem-shell-block-sum}, applied at each $s$ to $\operatorname{St}=\operatorname{St}_v$ and
$\mathsf{y}\in\mathfrak{Y}(\operatorname{St}_v)$, gives
\begin{equation*}
e^{-|\mathsf{y}|}\sum_{\square}\eta_p^{1/2}\le\mathsf{b}_{\square}(\operatorname{St}_v).
\end{equation*}

Now count the remaining factors:
\begin{itemize}
\item A unit of class (c) has one bracket and a unit of class (d) has four, so the sum over $s$
contributes at most a factor $4$.
\item The display \eqref{5.2:eq-carrier-exterior-d} carries the factor $2$.
\item $(\eta_p\mathsf{f}_v)^{1/2}=\eta_p^{1/2}\mathsf{f}_v^{1/2}$.
\end{itemize}
Put $\mu_v:=\mathsf{b}_{\square}(\operatorname{St}_v)\,\mathsf{f}_v^{1/2}\,
\tfrac{81}{64}\mathfrak{N}^{\flat}(v)$. The left member is then at most
\begin{equation*}
8\,\mathsf{w}_L^{-1}\sum_{v}\
\sum_{\mathsf{y}\in\mathfrak{Y}(\operatorname{St}_v):\
\operatorname{lab}(\mathfrak{h}_v,\mathsf{y})\ni Q_*}
\mu_v\,e^{-(\kappa_1-2)|\mathsf{y}|}\,e^{r\,d_k(\operatorname{lab}(\mathfrak{h}_v,\mathsf{y}))}.
\end{equation*}

\emph{The source-rooted sum.} By Lemma~\ref{5.2:lem-exempt-star}, in its clause on boxes,
$[\operatorname{St}_v]_k\subset\mathfrak{h}_v$. The family
$\mathbb{D}^{\mathrm{sh}}:=\{(\operatorname{St}(\Delta_{y_v}),\mathfrak{h}_v,\mu_v)\}$ is therefore a
family of unit data: each star is wall-connected, and neither the star nor its rim holds a source
cell, as shown above.

The source-rooted sum over families of cut cells (Theorem~\ref{5.2:thm-rooted-sum},
\eqref{5.2:eq-rooted-sum-a}) applies for $r\le a+9$. It bounds the double sum by
$S_d\,\mathsf{M}^{\partial}_{r+1}[\mathbb{D}^{\mathrm{sh}}]$.

The rim mass of the compact pointed units (Lemma~\ref{5.2:lem-rim-mass}, with
\eqref{5.2:eq-rim-mass-a}) is proved from the column bound and the foreign-window bound assumed
above. It gives $\mathsf{M}^{\partial}_{r+1}[\mathbb{D}^{\mathrm{sh}}]\le\mathsf{M}^{\mathrm{sh}}_{r+1}$.

Since $\mathsf{w}_L^{-1}=\theta'(g_k)$, the left member is at most
$8S_d\,\mathsf{M}^{\mathrm{sh}}_{r+1}\,\theta'(g_k)=\theta'(g_k)\mathsf{C}^{\Sigma}_r$, as in
\eqref{5.2:eq-anchor-sum-a}.

\emph{The lower bounds on $|\mathsf{y}|$.} By Lemma~\ref{5.2:lem-exempt-star}, the members of
$\operatorname{St}_v$ have levels $n'-1$, $n'$, $n'+1$. The walk $\mathsf{y}$ holds a source cell,
which has level $k$, and $\operatorname{St}_v\cup\mathsf{y}$ is connected. Along an adjacency the level
changes by at most $1$: wall-neighbours differ by at most $1$, and the members joined through a hub are
source cells of level $k$. So each of the levels $n'+2,\dots,k$ occurs in $\mathsf{y}$, and
$|\mathsf{y}|\ge k-n'-1$.

In the proof of Theorem~\ref{5.2:thm-rooted-sum}, the component of $\mathsf{y}$ holding its first
source cell (the rooted component) has a member $\Delta_1\in\partial\operatorname{St}_v$. Being a
connected family holding $\Delta_1$ and a source cell, it has at least $\rho(\Delta_1)$ members.
Combining:
\begin{itemize}
\item the lower bound of $\rho$ by the level-$k$ distance to the union of the source cells, as stated
in Notation~\ref{5.2:not-rooted-sum-notation};
\item the hub geometry $\operatorname{dist}_k(\Lambda^{\sim},Z^c)\ge 9\check{R}_k$;
\item the fact that $\Delta_1$ reaches at most $3$ level-$k$ units beyond $\Delta'$, and $\Delta'$ at
most $3$ units into $Z$;
\end{itemize}
one obtains that the level-$k$ distance from $\Delta_1$ to the union of the source cells is at least
$9\check{R}_k-8$, and hence
\begin{equation*}
|\mathsf{y}|\ge\rho(\Delta_1)\ge 9\check{R}_k-8=:n_0 .
\end{equation*}

Hence, in the interchange step of that proof, the sum over the size $n$ of the rooted component may be
restricted to $n\ge n_0$. The notation gives $\lambda_{\Sigma}'\ge(\lambda_*-\log(e\hat b))+75c_l-1$,
and $\lambda_*-\log(e\hat b)>0$ by the condition on $\lambda_*$. So $\lambda_{\Sigma}'-1\ge 4$, and
\begin{equation*}
\sum_{n\ge n_0}n\,e^{-(\lambda_{\Sigma}'-1)n}
\le e^{-n_0}\sum_{n\ge 1}n\,e^{-3n}\le e^{-n_0}.
\end{equation*}
This factor carries through the rest of the argument unchanged and gives the sharper bound.
\end{proof}

\begin{remark}[The decay factors in the pointed sum]\label{5.2:rem-decay-factors}
We note where each decay factor is spent in the summation of the compact pointed pieces behind one coarse anchor (Corollary~\ref{5.2:cor-anchor-sum}). Each factor is spent exactly once.

A compact pointed piece $p=(\square,\mathsf{y};s)$ of a unit $v$ enters that corollary with the bound
\begin{equation*}
 \mathfrak{N}^{\mathrm{ex}}_A(p)
 =2\mathsf{w}_L^{-1}(\eta_p\mathsf{f}_v)^{1/2}\,\tfrac{81}{64}\,\mathfrak{N}^{\flat}(v)\,
 e^{-(\kappa_1-1)|\mathsf{y}|}.
\end{equation*}
Here $\mathsf{y}$ is the walk, i.e.\ the family of cut cells in proper scales. Its factors are used as follows.
\begin{enumerate}
\item The factor $\mathsf{f}_v^{1/2}$ pays the cube $\bar{\mathsf{w}}_v^3$ of the threshold weight, and nothing else. It does so cell by cell, through the bound $\mathsf{f}(\Delta')^{1/2}\bar{\mathsf{w}}(\Delta')^3\le\mathsf{C}_{\mathrm{fw}}$ at every cell $\Delta'$, which holds trivially where $\bar{\mathsf{w}}(\Delta')=3$ and comes from the collar geometry of the foreign live arrest otherwise. It is never spent against a count of cells and never against an age.
\item The factor $\eta_p^{1/2}$ is used through its majorant $e^{-\frac12\delta_P d(\tilde{\square},S_v)}$. It pays the sum over the blocks $\square$ at fixed read set $S_v$, organised in shells of $d(\tilde{\square},S_v)$; only the block lemma reads it.
\item For each member of the walk, the rate $\kappa_1-1$ is shared out by \eqref{5.2:eq-decay-factors-1} below.
 \begin{itemize}
 \item The first summand $1$ pays the tower of coarse fattened blocks $\tilde{\square}$ that stand over a fine cell in the shell organisation of the block lemma. This tower has height up to $k-n'$, where $n'$ is the level of the cell $\Delta_{y_v}$ of the unit. It is genuinely present for the shells of $d(\tilde{\square},S_v)$, although, measured by the positional decay factor $\eta_p$ itself, these blocks recede.
 \item The term $c_l^{\circ}(r+1)$ pays the growth of the label and the anchor.
 \item The term $\log(e\hat b)$ pays the entropy of the connected families; the refinement count $2dL^{d-1}$ sits inside $\hat b$.
 \item The second summand $1$ pays $e^{\rho(\Delta_1)}$; here $\Delta_1$ is the cell through which the walk reaches the rim of the star of the unit, and $\rho(\Delta_1)$ is the least cardinality of a connected family of cells holding $\Delta_1$ and a source cell.
 \item The remainder $\lambda_\Sigma'-1$ is the net rate per cell of the walk, against which the sum over the lengths of the walk converges.
 \end{itemize}
 The later links $\mathsf{y}_2,\mathsf{y}_3,\mathsf{y}_4$ have their own families and their own Cauchy factors.
\end{enumerate}
The other halves of $\eta_p$ and of $\mathsf{f}_v$ are spent by the containment of the dilated carrier on the disc $|\mathsf{t}|\le\rho_A=\mathsf{w}_L(\eta_p\mathsf{f}_v)^{-1/2}$. The one-scale summation of localisation domains is not used on these pieces, and neither are the Stage-A weight supremum $\mathsf{W}$ and the constant $C_A'$.

The source-rooted summation theorem is blind to the nature of the units. The boundary-type units are not summed here but by the summation of boundary terms at finer zones. That summation runs the grain version of the source-rooted theorem, with the rate dial $\beta=3$, in the unit's own grain $\pi^{\mathfrak{s}}_{j_v}$ on the data $(\operatorname{St}^{j_v}_v,K_v,2\mathfrak{N}(v))$, with the $\pi_{j_v}$-cubes of $X_v$ exempt. It does so only after the fine rim and the blocks have been summed exactly. Its rim mass comes from the rim-mass bound for the units of a given level $j_v$ with $X_v\cap Z=\emptyset$, under the bound for the boundary-type terms, the anchoring statement and the label geometry, without the block constant $\mathsf{b}_{\square}$ and without any level count of the kind described next.

The factor $e^{-\rho(\Delta_1)}$ in the rim mass $\mathsf{M}^{\partial}$ is what absorbs the factor $k-\operatorname{lev}\Delta+1$, the number of levels from that of a counted cell $\Delta$ up to $k$, that a direct count of the cells stacked at one anchor would produce. That count is not formed on the long singles and tails, whose mass is counted per cell directly; these too are carried by the summation of boundary terms at finer zones, not by Corollary~\ref{5.2:cor-anchor-sum}. The compact units do not need $e^{-\rho(\Delta_1)}$, because a cell lies on the rim of the star $\operatorname{St}(\Delta')$ for at most $N_\partial=5(L+2)^{2d}$ cells $\Delta'$. The summation of boundary terms at finer zones uses the stronger factor $e^{-3\rho(\Delta_1)}$, with $\rho$ computed in the unit's own grain.

\emph{Witnesses.} Three situations show that no hidden count survives.
\begin{enumerate}
\item Let $Q\in\pi_k$ lie over zone $k-s$ of an untreated component. It presents $3^dL^{ds}$ cells, each of mass of the order of $\mathsf{K}^{\flat}$. Each piece anchored there needs a trunk, i.e.\ the connected component of its walk holding the first source cell, of $n\ge\max\{s-1,9\check{R}_k-8\}$ own-scale cells issued from a source cell. The pairs (trunk, cell) of length $n$ at one source number at most $(e\hat b)^n\,n\,N_\partial$, each at weight $e^{-\lambda_\Sigma n}$. Together the $L^{ds}$ cells weigh less than one.
\item The age factor $L^{d(k-k_{a'})}$, $k_{a'}$ the step of the foreign arrest $\mathfrak{a}'$ behind the anchor, which the collar of that arrest cannot pay, is never formed.
\item Zones of a level $i$ that are only one cube of $\pi_{i+1}$ thick (thin strata) occur here wherever they arise: beside foreign live arrests, inside the sets $\mathcal{P}_F$ carried by boundary-type terms $F$, and where a completed operation $\mathbf{R}$ has left them. About them only three geometric facts are read:
 \begin{itemize}
 \item wall- or corner-neighbouring cells differ in level by at most one;
 \item $\operatorname{dist}_{j+1}(\Omega_j^c,\Omega_{j+1})\ge1$;
 \item the hubs $\Lambda^{\sim}_C$ are boxes of side at most $h_\sharp=112$, at least $19\check{R}_k$ apart, with $\operatorname{dist}_k(\Lambda^{\sim},Z^c)\ge9\check{R}_k$, and every cell disjoint from $\Lambda$ and lying in $Z$ has level $k$.
 \end{itemize}
 Nothing else about the strata needs to be named.
\end{enumerate}
\end{remark}

\begin{proof}
\emph{The share of $\kappa_1-1$.} By definition
\begin{equation*}
 \lambda_\Sigma=\kappa_1-2-c_l^{\circ}(r+1),\qquad \lambda_\Sigma'=\lambda_\Sigma-\log(e\hat b),
\end{equation*}
so that
\begin{equation}\label{5.2:eq-decay-factors-1}
 \kappa_1-1=1+c_l^{\circ}(r+1)+\log(e\hat b)+1+(\lambda_\Sigma'-1).
\end{equation}
We bound the remainder from below for $r\le a+9$, where $a=\frac{31}{20}\kappa$. Let $\lambda_*(d,L)$ be the entropy constant of the site geometry, a number satisfying $L^de^{-\frac12(\lambda_*(d,L)-\log(e\hat b))}\le\frac12$. The floor $\kappa_1-1\ge c_l(2\kappa+4)+\lambda_*(d,L)$ and $c_l^{\circ}\le c_l$ give
\begin{align*}
 \lambda_\Sigma'
 &\ge c_l(2\kappa+4)+\lambda_*(d,L)-1-c_l(a+10)-\log(e\hat b)\\
 &= c_l\bigl(\tfrac{9}{20}\kappa-6\bigr)+\lambda_*(d,L)-\log(e\hat b)-1 .
\end{align*}
Since $\kappa\ge180$, we have $\frac{9}{20}\kappa-6\ge75$. Hence $\lambda_\Sigma'-1\ge(\lambda_*(d,L)-\log(e\hat b))+75c_l-2$, which is what remains as the net rate per cell of the walk.

\emph{The first witness.} A cell of the cube $Q$ has level $n'=k-s$. Put
\begin{equation*}
 n_0:=\max\{s-1,\,9\check{R}_k-8\}.
\end{equation*}
Every trunk has at least $n_0$ members. Its member $\Delta_1$ on the rim of the star $\operatorname{St}(\Delta')$ has level at most $n'+2$, since the members of the star have levels $n'-1$, $n'$, $n'+1$ and wall-neighbouring cells differ in level by at most one; its source cell has level $k$; and adjacent members of a connected family differ in level by at most one, the edges through a hub joining cells of level $k$. So the trunk holds at least $k-(n'+2)+1=s-1$ members. Moreover the trunk is a connected family holding $\Delta_1$ and a source cell, so its number of members is at least $\rho(\Delta_1)$, and
\begin{equation*}
 \rho(\Delta_1)\ \ge\ \operatorname{dist}_k\bigl(\Delta_1,\textstyle\bigcup\text{source cells}\bigr)\ \ge\ 9\check{R}_k-8 .
\end{equation*}
The first inequality is a general property of $\rho$ in the site geometry; the second follows from $\operatorname{dist}_k(\Lambda^{\sim},Z^c)\ge9\check{R}_k$, since $\Delta_1$ reaches at most $3$ level-$k$ units beyond $\Delta'$ and $\Delta'$ lies at most $3$ units inside $Z$.

The pieces are counted from the source as follows.
\begin{itemize}
\item The trunk is the connected component of $\mathsf{y}$ holding its first source cell. The connected families of $n$ members through a given member number at most $(e\hat b)^n$.
\item The trunk has a member $\Delta_1$ on the rim of the star of the unit, and $\Delta_1$ has at most $n$ positions on a trunk of $n$ members.
\item A given cell $\Delta_1$ lies on the rim $\partial\operatorname{St}(\Delta')$ for at most $N_\partial$ cells $\Delta'$, and the star of a compact unit depends on the unit only through its cell $\Delta'$.
\end{itemize}
At one source cell the total weight, in units of the per-cell mass, is therefore at most
\begin{equation*}
 \sum_{n\ge n_0}n\,N_\partial\,(e\hat b)^n e^{-\lambda_\Sigma n}
 = N_\partial\sum_{n\ge n_0}n\,e^{-\lambda_\Sigma' n}.
\end{equation*}
Put $x:=e^{-\lambda_\Sigma'}$. By the bound just proved and the defining property of $\lambda_*(d,L)$,
\begin{equation*}
 x\le\tfrac14L^{-2d}e^{-(75c_l-1)}\le\tfrac14 .
\end{equation*}
Hence $\sum_{n\ge1}nx^n=x(1-x)^{-2}\le\frac{16}{9}x$. With $N_\partial=5(L+2)^{2d}$ the total at one source cell is at most
\begin{equation*}
 \tfrac{20}{9}\bigl(1+2/L\bigr)^{2d}e^{-(75c_l-1)}\le\tfrac{20}{9}\,3^{2d}e^{-(75c_l-1)}<1,
\end{equation*}
because $c_l=1+234d$. The passage from the anchor to the source cells multiplies this only by the further factor $S_d$ of the source-rooted summation theorem, a number depending on $d$ alone.
\end{proof}

\begin{definition}[Bundles and trace labels for boundary terms]\label{5.2:def-boundary-bundles}
We work at the localisation site of the operation $\mathbf{R}_k$. Two closed sets are said to \emph{meet} if they have a common interior point. Let $v$ be one of the following; in both cases $K_v:=[X_v]_k$, where $[A]_k$ denotes the union of the cubes of $\pi_k$ meeting $A$.
\begin{itemize}
\item a plain boundary-type unit, of level $j_v\le k$ with $X_v\in\mathbf{D}_{j_v}$ and in the format of boundary-type terms, either of the first part, with $K_v\setminus Z\neq\emptyset$ and the single bracket index $s=4$, or of the second part, with $X_v\cap Z=\emptyset$ and $s=4,\dots,1$; here $\mathfrak{N}(v):=2\mathfrak{N}^{\mathbb{C}}_{\varepsilon_k}(v)$, twice the orbit-datum size of $v$;
\item a long single or a tail of the first part, with $X_v\subset\Lambda_{j_v}$; here $\mathfrak{N}(v):=\mathfrak{N}^{\flat}(v)=2\mathfrak{n}_v$.
\end{itemize}
We write $\pi^{\mathfrak{s}}$ for the family of the cells, in proper scales, not meeting $\Lambda^{\sim}$. Put $j:=j_v$. The hubs are the components $\Lambda^{\sim}_C$ of $\Lambda^{\sim}$. A source cell is a member of $\pi^{\mathfrak{s}}$ that shares a wall with a hub. We write $\operatorname{St}_v$ for the cells of $\pi^{\mathfrak{s}}$ meeting $X_v$, and we set $T_v:=\pi^{\mathfrak{s}}\setminus\operatorname{St}_v$. The pieces $v_{(\square,\mathsf{y};s)}$, for $\square$ a block of the cover of the determining set $\mathfrak{B}_s$ of the bracket $s$ and $\mathsf{y}\subset T_v$ finite, are
\begin{equation*}
v_{(\square,\mathsf{y};s)}:=\prod_{\Delta\in\mathsf{y}}\int_0^1 d\mathsf{s}(\Delta)\,\partial_{\mathsf{s}(\Delta)}
\int_0^1 dt\,\partial_{\mathsf{t}}\big|_{0}\,
f_v\bigl(\Phi_s^{t+\mathsf{t}\hat{\chi}_{\square}}(\mathsf{s})\bigr)\Big|_{\mathsf{s}=0\text{ off }\mathsf{y}},
\end{equation*}
\begin{equation*}
\Phi_s^{\chi}:=\bigl(e^{i\eta\chi\mathbb{H}_s}W_s,\ J_s+\mathcal{J}^{\sharp}_s[\chi]\bigr),
\end{equation*}
where $f_v$ is the function carried by the unit $v$ (for the boundary-type units of the first kind, $f_v\circ\Phi$ means the orbit datum of $v$ evaluated at the presented rotation $\mathfrak{Q}^{\mathfrak{b}}[\Phi]$), $\eta$ is the lattice spacing, $\mathbb{H}_s$ is the response layer of the bracket $s$, $(W_s,J_s)$ is its configuration, $\mathcal{J}^{\sharp}_s[\chi]$ is the current of the dilated layer $\chi\mathbb{H}_s$, $\hat{\chi}_{\square}$ is the bump of the block $\square$, $\mathsf{t}$ is the complex amplitude multiplying the block piece $\hat{\chi}_{\square}\mathbb{H}_s$ of the layer, and $\mathsf{s}(\Delta)$ are the decoupling parameters.

(i) \emph{Grain and trace.} The grain at level $j$ is
\begin{multline*}
\pi^{\mathfrak{s}}_j:=\{\Delta\in\pi^{\mathfrak{s}}:\ \operatorname{lev}\Delta\ge j\}\\
\cup\{Q\in\pi_j:\ Q\subset\operatorname{cl}\bigl((\Omega^{\mathrm{cu}}_j)^{c}\bigr),\ Q\text{ not meeting }\Lambda^{\sim}\},
\end{multline*}
where $\Omega^{\mathrm{cu}}_j$ is the union of the cells of level at least $j$ and $\operatorname{lev}\Delta$ is the level of the cell $\Delta$.
For $\Delta\in\pi^{\mathfrak{s}}$, $\operatorname{tr}_j\Delta$ is the member of $\pi^{\mathfrak{s}}_j$ holding $\Delta$. For a family $\mathsf{y}$ of cells, $\operatorname{tr}_j\mathsf{y}:=\{\operatorname{tr}_j\Delta:\Delta\in\mathsf{y}\}$.

We write $\operatorname{St}^j_v$ for the members of $\pi^{\mathfrak{s}}_j$ meeting $X_v$, and we set $T^j_v:=\pi^{\mathfrak{s}}_j\setminus\operatorname{St}^j_v$. We write $\mathfrak{H}_v$ for the hubs that are met by $X_v$ or share a wall with a member of $\operatorname{St}^j_v$.

Two members of $\pi^{\mathfrak{s}}_j$ are adjacent if they share a wall or both share a wall with one and the same hub; connectedness, components and spanning forests of a family of members are taken for this adjacency. We write $\partial\operatorname{St}^j_v$ for the members of $T^j_v$ sharing a wall with a member of $\operatorname{St}^j_v$, and $\partial^{+}\operatorname{St}^j_v$ for $\partial\operatorname{St}^j_v$ together with the members of $T^j_v$ sharing a wall with a hub of $\mathfrak{H}_v$. We write $\mathfrak{Y}_v$ for the family of the finite $\tilde{\mathsf{y}}\subset T^j_v$ such that every component of $\tilde{\mathsf{y}}$ has a member in $\partial^{+}\operatorname{St}^j_v$, and such that $\mathfrak{H}_v\neq\emptyset$ or $\tilde{\mathsf{y}}$ holds a source cell.

(ii) \emph{Label.} For $\tilde{\mathsf{y}}\subset T^j_v$ finite put
\begin{equation*}
\operatorname{lab}(v,\tilde{\mathsf{y}}):=K_v\cup[\tilde{\mathsf{y}}]_k\cup\operatorname{Hub}(\tilde{\mathsf{y}})
\cup\operatorname{Att}(v,\tilde{\mathsf{y}}).
\end{equation*}
The two added sets are defined as follows.
\begin{itemize}
\item $\operatorname{Hub}(\tilde{\mathsf{y}})$: fix a spanning forest of $\tilde{\mathsf{y}}$. For every edge of it that joins two members sharing no wall, take the intermediate $\pi_k$-cubes of a wall-chain of $\pi_k$-cubes, with fewer than $d(x+1)$ intermediate cubes, joining the two inside the smallest box holding both, $x$ being their $\operatorname{dist}_k$-distance; such a chain exists by the site geometry of the localisation site. $\operatorname{Hub}(\tilde{\mathsf{y}})$ is the union of these cubes.
\item $\operatorname{Att}(v,\tilde{\mathsf{y}})$: consider every component of $\tilde{\mathsf{y}}$ none of whose members shares a wall with $\operatorname{St}^j_v$. For each such component, take such a wall-chain of $\pi_k$-cubes from its first source cell, in a fixed enumeration of $\pi^{\mathfrak{s}}$, at a hub $H\in\mathfrak{H}_v$ to a nearest cube of $K_v$ that lies in $H$ or shares a wall with $H$. $\operatorname{Att}(v,\tilde{\mathsf{y}})$ is the union of these chains. (No chain is taken for a component without such a source cell; by the fine-rim lemma, clause (c), the bundle below vanishes unless $\tilde{\mathsf{y}}\in\mathfrak{Y}_v$, and then every such component has one.)
\end{itemize}

(iii) \emph{Bundles.} The bundle $\mathfrak{b}=(v;s;\tilde{\mathsf{y}})$ is
\begin{equation*}
\mathfrak{b}:=\sum_{\square}\ \sum_{\mathsf{y}\subset T_v:\ \operatorname{tr}_j\mathsf{y}=\tilde{\mathsf{y}}}
v_{(\square,\mathsf{y};s)},
\qquad
\|\mathfrak{b}\|:=2\mathfrak{N}(v)\,e^{-(\kappa_1-1)|\tilde{\mathsf{y}}|},
\end{equation*}
$\square$ running over the blocks of the cover of $\mathfrak{B}_s$. Its label is $\operatorname{lab}(\mathfrak{b}):=\operatorname{lab}(v,\tilde{\mathsf{y}})$. Its reading set is
\begin{multline*}
Y_0(\mathfrak{b}):=\bigcup\operatorname{St}^j_v\cup\bigcup\tilde{\mathsf{y}}\\
\cup\{\text{the hubs of }\mathfrak{H}_v\text{ and those sharing a wall with }\tilde{\mathsf{y}}\}.
\end{multline*}
The exempt set, at the level of cells, is $\bigcup\operatorname{St}_v\cup\Lambda^{\sim}=\bigcup\operatorname{St}^j_v\cup\Lambda^{\sim}$, the two unions $\bigcup\operatorname{St}_v$ and $\bigcup\operatorname{St}^j_v$ being equal by the grain lemma. Hence the bundles resum the pieces $v_{(\square,\mathsf{y};s)}$ themselves.

On the units of the two kinds listed above, three filing conventions (iv)--(vi) are in force; (vii) extends them to the long singles and tails of the second part.

(iv) \emph{Reading per bundle.} The per-piece bound, and the bound on the continuations of the amplitude pieces, in the lemma summing the pieces of links two to four of the plain units whose label leaves $Z$, are read per bundle $\mathfrak{b}=(v;s;\tilde{\mathsf{y}})$. This is a resummation inside a proof. The classes of units, the choice of exempt families in proper scales, and the pieces $v_{(\square,\mathsf{y};s)}$ are left as they are.

(v) \emph{Trace label and reading set.} At the first link the label and the reading set of a piece are those of its bundle. The spanning forest of the label is taken on the trace:
\begin{gather*}
\operatorname{lab}_1(p):=K_p:=\operatorname{lab}\mathfrak{b}(p)
=K_v\cup[\tilde{\mathsf{y}}]_k\cup\operatorname{Hub}(\tilde{\mathsf{y}})\cup\operatorname{Att}(v,\tilde{\mathsf{y}}),\\
\mathcal{R}_1(p):=Y_0(\mathfrak{b}(p)),
\end{gather*}
where $p=(\square,\mathsf{y};s)$ and $\tilde{\mathsf{y}}=\operatorname{tr}_j\mathsf{y}$. The later labels are $\operatorname{lab}_t:=\operatorname{lab}((\operatorname{lab}_{t-1},\mathcal{R}_{t-1}),\mathsf{y}_t)$ for $t=2,3,4$, with $\operatorname{lab}(\cdot,\cdot)$ the label rule at links two to four of the lemma summing the pieces of links two to four of the plain units whose label leaves $Z$, and $Y_4:=Y(\operatorname{lab}_4)$, where $Y(A):=A\cup\{\text{components of }Z\text{ meeting }A\}$. This is a re-filing of the non-empty continuations. By the declaration lemma for the bundle labels and by the fine-rim lemma, clause (e), the triple $(\operatorname{lab}_t,\mathcal{R}_t,Y_4)$ is a function of $(v,\tilde{\mathsf{y}},\mathsf{y}_2,\mathsf{y}_3,\mathsf{y}_4)$, there is one $Y(\operatorname{lab}_1)$ per bundle, and each filed sum $\sum_{Y_4(\mathfrak{p})=Y_4}\mathfrak{p}$ is a sum of whole bundles. The re-filing is independent of the complex source, and every exponent is left untouched.

(vi) \emph{No amplitude gain.} On these units the gain $\mathsf{w}_L^{-1}\eta_p^{1/2}$ of the Cauchy bound in the amplitude variable $\mathsf{t}$ is not used, and the blocks $\square$ are resummed as well. The reason is that there are at least $2^{-d}L^{d(j-n)}$ blocks of zone $n$ with $\square^{\zeta}\subset X_v$, $\square^{\zeta}$ being the support of $\hat{\chi}_{\square}$, per cube of $\pi_j$ in $X_v$ lying over zone $n$, and for the boundary-type units of the first kind the amplitude bound carries no far factor to pay for them.

The smallness comes instead from the length of the label. For $\tilde{\mathsf{y}}\in\mathfrak{Y}_v$ the fine-rim lemma gives $d_k(\operatorname{lab}\mathfrak{b})\ge 9\check{R}_k-2$, which gives the factor $e^{-(9\check{R}_k-2)}$ for second-part units and $e^{-\frac12(9\check{R}_k-2)}$ for first-part units. This factor is converted into $\theta'(g_k)\mathsf{C}^{\Sigma,4}$, by the corollary to the summation of boundary terms at finer zones (Theorem~\ref{5.4:thm-finer-zone-sum}) that bounds the sum of the right members of that theorem by $\theta'(g_k)\mathsf{C}^{\Sigma,4}$, only where the bound is read against $\theta'(g_k)$. These conventions change only the organisation of the sums inside proofs; no hypothesis and no conclusion of the correlation theorem is changed.

Accordingly, for the plain first-part units with $K_v\setminus Z\neq\emptyset$ the amplitude expansion of the bracket $s=4$ is still performed, but its pieces are bounded and filed only through their bundles.

(vii) \emph{Second-part long singles and tails.} For the long singles and tails of the second part, the same conventions are used at the grain $\pi^{\mathfrak{s}}$ itself. The bundle is $\mathfrak{b}(v;s;\mathsf{y})=\sum_{\square}v_{(\square,\mathsf{y};s)}$ with $\mathsf{y}\subset T_v$ finite, and $\|\mathfrak{b}\|=2\mathfrak{N}^{\flat}(v)e^{-(\kappa_1-1)|\mathsf{y}|}$. The label and the reading set at the first link are
\begin{equation*}
\operatorname{lab}_1=K_p=\operatorname{lab}(K_v,\mathsf{y}):=K_v\cup[\mathsf{y}]_k\cup\operatorname{Hub}(\mathsf{y}),
\end{equation*}
\begin{equation*}
\mathcal{R}_1=Y_0(p):=\bigcup\operatorname{St}_v\cup\bigcup\mathsf{y}\cup
\{\text{the hubs sharing a wall with }\mathsf{y}\},
\end{equation*}
with $\operatorname{Hub}(\mathsf{y})$ formed as in (ii), adjacency in $\pi^{\mathfrak{s}}$ being defined in the same way and the spanning forest being taken on $\mathsf{y}$ itself. The hypothesis of the Cauchy bound in the amplitude variable $\mathsf{t}$ (the containment of the dilated configuration in the domain of $f_v$ and the bound on the deviation of $f_v$ from its value at the trivial pair) is read at the centre $\mathsf{t}=0$ of its disc only.

Conventions (vi) and (vii) are declared departures from \cite[p.~372]{B-CMP122b}, which sends both the terms with localisation domains not contained in $\square^{\sim 2}$ and the boundary terms of the second group through the identity quoted below and re-expands the bracket:
\begin{quote}
for the boundary terms in the second group, we write the following simpler identity
\begin{equation*}
\begin{split}
\mathbf{E}^{(j)}(X,U''_k,z)&=\mathbf{E}^{(j)}(X,U^0_k,z)\\
&\quad+\bigl[\mathbf{E}^{(j)}(X,\exp i\eta\mathbf{H}''_k(g_kB)U^0_k,z)
-\mathbf{E}^{(j)}(X,U^0_k,z)\bigr]
\end{split}
\tag*{\cite[(1.60)]{B-CMP122b}}
\end{equation*}
and we repeat all the above described operations for the expression in the square bracket
\end{quote}
In \cite{B-CMP122b}, as in \cite{B-CMP119}, the bracket is expanded again in the blocks and the amplitude is taken. Here it is not. By the fine-rim lemma, clause (b), the sum over the blocks telescopes to the bracket of \cite[(1.60)]{B-CMP122b}; the bundle is bounded there, the amplitude factor $\mathsf{w}_L^{-1}$ being left unused.
\end{definition}

\begin{lemma}[Plain units finer than a cell and the window weights]
\label{5.2:lem-plain-fine-units}
We work at the localisation site of the operation $\mathbf{R}_k$, in the current global label,
with $C$ the component of the large-field region under treatment. For an attempt $\mathfrak{a}$
we write $k_a$ for its step, $X_{\mathfrak{a}}$ for its component, and $h_a$, $k_{0a}$ for the two
levels fixed in its construction \cite[(1.10)--(1.12)]{B-CMP122a}; the corresponding indices of
the running attempt are written $h$ and $k_0$, and $\mathsf{C}:=(Z'_h)^{\sim}\cap C$ is the core
of $C$. Fix $r\le\frac85\kappa$ and, for a unit $v$, put $W(v):=\mathfrak{N}^{\flat}(v)\,e^{r d_k(K_v)}$;
for the boundary-type units of hypothesis (a) below the stored size $\mathfrak{N}^{\flat}(v)$ in
$W(v)$ is $\mathfrak{N}^{A,\mathbb{C}}(v)$.
For a cell $\Delta\subset C$ of level $n:=\operatorname{lev}\Delta$, let $\mathsf{S}^{\mathrm{dp}}(\Delta)$,
resp.\ $\mathsf{S}^{\mathrm{sh}}(\Delta)$, be the sum of $W(v)$ over the units
$v\in\mathfrak{L}^{\mathrm{pl}}$ of the class $\mathfrak{L}^{\mathrm{dp}}$ of deep units, resp.\ of
the class $\mathfrak{L}^{\mathrm{sh}}$ of shallow units, whose hull $\hat X_v$ meets $\Delta$. Put
\begin{equation}\label{5.2:eq-plain-fine-units-1}
\begin{gathered}
C^{<}:=C_{\mathsf{b}}\sum_{u\ge1}(1+u)^{p_{\mathsf{b}}}L^{du}e^{-\varsigma_0(L^u-1)},\\
C_{\circledast}:=\sum_{i\ge0}\bigl[(2i+5)^d+(L+2i+4)^d+i(3i+4)^d\bigr](2+i)^3e^{-\varsigma_1 i},
\end{gathered}
\end{equation}
with $C_{\mathsf{b}}$, $p_{\mathsf{b}}$ as in hypothesis (a) below. Assume the following.
\begin{enumerate}
\item[(a)] \emph{The per-cube bound for plain boundary-type units.} For $j\le k$ and
$\square'\in\pi_j$,
\begin{equation}\label{5.2:eq-plain-fine-units-b}
\begin{gathered}
\sum_{\substack{F\in\mathfrak{L}^{\mathrm{B}}:\ j_F=j,\\ X_F\ni\square'}}
W(F)\,e^{\varsigma_0 d_j(X_F)\mathbf{1}[j<k]}\le\mathsf{b}(k-j),\\
\mathsf{b}(s):=\mathsf{b}_k(s)+(s+1)b_\delta,\qquad
\mathsf{b}_k(s):=b_0+(s+1)C_NC_0^{\sharp}(\mathsf{T}_k+c_s)^{r_{\mathrm{pr}}},\\
b_0:=2K(d)B_0+N(g_k)C_0^{\sharp},\\
\mathsf{b}(s+u)\le C_{\mathsf{b}}(1+u)^{p_{\mathsf{b}}}\,\mathsf{b}(s),\qquad
p_{\mathsf{b}}:=1+r_{\mathrm{pr}},\qquad C_{\mathsf{b}}:=(1+3\lambda)^{r_{\mathrm{pr}}},
\end{gathered}
\end{equation}
the function $\mathsf{b}$ being non-decreasing. Here $B_0$ is the constant of the bound
\cite[(2.41)]{B-CMP119} for boundary terms, $K(d)$ the constant, depending on $d$ only, of the
companion count to \cite[(2.42)]{B-CMP119}, and $b_\delta$, $c_s$ are the constants of the
per-cube bound for stored boundary-type terms. The first line of
\eqref{5.2:eq-plain-fine-units-b} is taken for every plain stored boundary-type unit alive at the
localisation site, whatever its position: contained in $Z$, disjoint from $Z$ at any zone,
astride $\partial Z$, or in a foreign frozen window.
\item[(b)] \emph{The anchoring statement.} First, for $j<k$ one has
$Z_j^{\boxplus}\subset Z_{j+1}$, both sets being those of the current label. Second, every
$F\in\mathfrak{L}^{\mathrm{B}}$ of level $j$ has a $\pi_j$-cube meeting $Z_j^{\boxplus}$ (for the
boundary-type terms listed last in the storage format, on their band).
\item[(c)] \emph{The preparation-count bound.} For every attempt $\mathfrak{a}$, at its step
$k_a$, one has $\check{R}_{k_{0a}+1}\le\mathsf{R}^{\uparrow}(\check{\mathsf{T}}_{k_a})$ and
$\check{\mathsf{T}}_{k_a}\ge1$; the function $\mathsf{R}^{\uparrow}$, fixed in the statement
bounding the number $N(g)$ of preparatory stages, obeys, for $\mathsf{T}\ge1$,
\begin{equation*}
\mathsf{R}^{\uparrow}(\mathsf{T})\le L\Bigl(\mathsf{T}+1
+\log\bigl(1+3\lambda C_N(\mathsf{T}+1)^{r_{\mathrm{pr}}}+2c_2\bigr)\Bigr)^{r_{\mathrm{pr}}};
\end{equation*}
and the
layer by which $(Z'_{k_{0a}})^{\sim}$ exceeds $Z'_{k_{0a}}$ is one level-$k_a$ unit thick.
\item[(d)] \emph{The geometry of the core.} One has $h+1\le k_0$; all cells of $\mathsf{C}$, and
all cells within one own unit of $\mathsf{C}$, have level at most $h$; $\mathsf{C}\subset Z''_{h+1}$;
and $Z_h\cap C$ lies at least $L\check{R}_{h+1}$ level-$h$ units from $\mathsf{C}^c$.
\item[(e)] \emph{The window weights of the running attempt.} The cells carrying window weight of
the running attempt have levels $>k_0$ and lie off $Z''_{h+1}$, and such a cell has
$\mathsf{w}^{\mathrm{win}}\le L\check{R}_{h+1}$.
\end{enumerate}
Then the following hold.
\begin{enumerate}
\item[(i)] Let $j\le k$ and let $\Delta\subset C$ be a cell of level $n>j$. Either
$\operatorname{dist}_n(Z_j^{\boxplus},\Delta)\ge1$, or $\Delta$ lies in a window: there is an
attempt $\mathfrak{a}$ with $\Delta\subset X_{\mathfrak{a}}$ and $k_{0a}\le j<n\le k_a$.
\item[(ii)] Let $\Delta\subset C$ be a cell with $\mathsf{w}^{\mathrm{win}}(\Delta)>1$. If the
weight is that of the running attempt,
then $\Delta\not\subset\mathsf{C}$ and $\mathsf{w}^{\mathrm{win}}(\Delta)\le L\check{R}_{h+1}$. If it
is that of a live arrest $\mathfrak{a}$, then $\Delta\subset\mathsf{C}$, and the following collar
claim holds: $Z_{k_a}\cap C$ has a collar contained in $\mathsf{C}$, at least $\check{R}_{k_a}$
level-$k_a$ units wide and made of cells of level at most $k_a$; consequently, for $x$ in this
collar and $E$ its outer edge, the first line of \eqref{5.2:eq-plain-fine-units-a} holds, and the
weight obeys its second line:
\begin{equation}\label{5.2:eq-plain-fine-units-a}
\begin{gathered}
d^{\wedge}(x,\mathsf{C}^c)\ge \frac{M}{M_1}\operatorname{dist}_{k_a}(x,E),\\
\mathsf{w}^{\mathrm{win}}(\Delta)\le L\,C_{\uparrow}\,\check{R}_{k_a}\le L\,C_{\uparrow}\,\varrho(\Delta),\\
C_{\uparrow}:=L\bigl(2+r_{\mathrm{pr}}+\log(1+3\lambda C_N+2c_2)\bigr)^{r_{\mathrm{pr}}}.
\end{gathered}
\end{equation}
\item[(iii)] Let $\Delta',\Delta\subset C$ be cells, $n':=\operatorname{lev}\Delta'$ and
$i:=\operatorname{dist}_{n'}(\Delta',\Delta)$. If $\mathsf{w}^{\mathrm{win},\vee}(\Delta')$ is
attained on a live arrest, then
\begin{equation}\label{5.2:eq-plain-fine-units-3}
\mathsf{w}^{\mathrm{win},\vee}(\Delta')\le 2L\,C_{\uparrow}\,\varrho(\Delta)\,(2+i);
\end{equation}
otherwise $\mathsf{w}^{\mathrm{win},\vee}(\Delta')\le L\check{R}_{h+1}$, with
$\mathsf{w}^{\mathrm{win},\vee}(\Delta')=1$ if moreover $\Delta'\subset\mathsf{C}$.
\end{enumerate}
\end{lemma}

\begin{proof}
\emph{(i).} By the first part of hypothesis (b), $Z_j^{\boxplus}\subset Z_{j+1}$, both sets being
those of the current label. Behind this inclusion stand the facts that each renormalization step
adds at least ten layers of $MR_k$-cubes \cite{B-CMP122a}, and that the two enlargements made by
$\mathbf{T}_{k+1}$ add four and two layers of $LMR_{k+1}$-cubes \cite{B-CMP119}; of these only the
consequence that at least one level-$n$ unit is added is used, together with
$\check{R}_j\le L\check{R}_{j+1}$. Applying the inclusion successively at the levels
$j,j+1,\dots,n-2$, and using $Z_i\subset Z_i^{\boxplus}$, we obtain
\begin{equation*}
Z_j^{\boxplus}\subset Z_{j+1}\subset Z_{j+1}^{\boxplus}\subset\dots\subset Z_{n-1}\subset
Z_{n-1}^{\boxplus}.
\end{equation*}
Since $\Delta$ has level $n$, $\Delta\subset\Omega_n$.

Suppose first that $\Omega_n$ is, at $\Delta$, as the operation $\mathbf{T}_n$ makes it. Then
$\operatorname{dist}_{n-1}(Z_{n-1},\Omega_n)\ge 8L\check{R}_n$, and the layer added by $\boxplus$, one
layer of $M\check{R}_{n-1}$-cubes, is $\check{R}_{n-1}\le L\check{R}_n$ level-$(n-1)$ units thick.
Hence $Z_{n-1}^{\boxplus}$, and with it
$Z_j^{\boxplus}$, lies at least $7L\check{R}_n$ level-$(n-1)$ units, that is at least $7\check{R}_n$
level-$n$ units, away from $\Omega_n\supset\Delta$; since $7\check{R}_n\ge1$, the first alternative
of (i) holds.

Otherwise $\Omega_n$ is re-made at $\Delta$. This happens only through an attempt $\mathfrak{a}$
with $h_a<n\le k_a$, inside $X_{\mathfrak{a}}$ \cite{B-CMP122a}, and through at most one such
attempt; there $\Omega''_n=X_{\mathfrak{a}}\setminus Z''_n$ \cite[(1.11)--(1.12)]{B-CMP122a}, where
$Z'_i$ and $Z''_i$ denote the sets constructed by the attempt
\cite[(1.10)--(1.12)]{B-CMP122a}, with $Z'_i\supset Z_i$. We
distinguish three cases.

If $n\le k_{0a}$, then $Z''_n=(Z'_{n-1})^{\sim3}$: three layers of $LM\check{R}_n$-cubes, that is
$3\check{R}_n$ level-$n$ units, about $Z'_{n-1}\supset Z_{n-1}$. The layer added by $\boxplus$ to
$Z_{n-1}$ takes at most $L\check{R}_n$ level-$(n-1)$ units, that is $\check{R}_n$ level-$n$ units;
this leaves at least $2\check{R}_n\ge1$
level-$n$ units between $Z_{n-1}^{\boxplus}\supset Z_j^{\boxplus}$ and $\Delta$: the first
alternative holds.

If $n>k_{0a}>j$, the chain of inclusions above, run up to the level $k_{0a}$, gives
$Z_j^{\boxplus}\subset Z_{k_{0a}}\subset Z'_{k_{0a}}$, while
\begin{equation*}
Z''_n\supseteq Z''_{k_{0a}+1}=(Z'_{k_{0a}})^{\sim}
\end{equation*}
by \cite[(1.10)]{B-CMP122a}. By the second part of hypothesis (c) the layer added by $\sim$ here is
one level-$k_a$ unit thick, and since $n\le k_a$ a level-$k_a$ unit is at least one level-$n$ unit;
so $\operatorname{dist}_n(Z_j^{\boxplus},\Delta)\ge1$, and again the first alternative holds.

There remains the case $n>k_{0a}$ and $j\ge k_{0a}$. Together with $n\le k_a$ and
$\Delta\subset X_{\mathfrak{a}}$ this is the window alternative $k_{0a}\le j<n\le k_a$. The strips
play no role here: no plain boundary-type unit has its localisation domain in a strip, by the
description of the domain class in the strip construction.

\emph{(ii), running attempt.} By hypothesis (e) the cells carrying window weight of the running
attempt have levels $>k_0$, and $k_0\ge h+1$ by hypothesis (d); by (e) they lie off $Z''_{h+1}$,
which contains $\mathsf{C}$ by (d). Hence
$\Delta\not\subset\mathsf{C}$, and $\mathsf{w}^{\mathrm{win}}(\Delta)\le L\check{R}_{h+1}$ by
hypothesis (e).

\emph{(ii), live arrest: the collar claim.} If $k_a=h$: by hypothesis (d) the set $Z_h\cap C$ lies
at least
$L\check{R}_{h+1}$ level-$h$ units from $\mathsf{C}^c$, and $L\check{R}_{h+1}\ge\check{R}_h$; by
hypothesis (d) all cells of $\mathsf{C}$ have level at most $h=k_a$. If $k_a<h$: one has
\begin{equation*}
\operatorname{dist}_{k_a}(Z_{k_a},\Omega_{k_a+1})\ge 3L\check{R}_{k_a+1}\ge 3\check{R}_{k_a},
\end{equation*}
where the first inequality holds with $8$ in place of $3$ if $\Omega_{k_a+1}$ is made by the
operation $\mathbf{T}$, and with $3$ if it is re-made by an attempt, which we denote
$\mathfrak{a}_1$, with $h_{\mathfrak{a}_1}=k_a$; this collar lies in
$\Omega^c_{k_a+1}\cap C\subset Z_h\cap C\subset\mathsf{C}$ (the last inclusion by hypothesis (d)),
and its cells have level at most $k_a$. In
both cases the first line of \eqref{5.2:eq-plain-fine-units-a} follows from the comparison of the
contour distance $d^{\wedge}$ with the level distances: a contour that runs through cells of level
at most $k_a$ across $D$ level-$k_a$ units has $d^{\wedge}$-length at least $\frac{M}{M_1}D$; in the
collar all cells have level at most $k_a$. This proves the collar claim.

A cell $\Delta$ carrying the window weight of the live arrest $\mathfrak{a}$ lies in
$X_{\mathfrak{a}}$, the arrest's component of the large-field region at its step $k_a$, so
$\Delta\subset Z_{k_a}\cap C$. By the first part of hypothesis (b), iterated as in the proof of (i)
(trivially if $k_a=h$), $Z_{k_a}\subset Z_h$, and $Z_h\cap C\subset\mathsf{C}$ by hypothesis (d);
hence $\Delta\subset\mathsf{C}$. A contour from a point of $\Delta$ to $\mathsf{C}^c$ crosses the
whole collar, which is at least $\check{R}_{k_a}$ level-$k_a$ units wide; at its entry point $x$
into the collar $\operatorname{dist}_{k_a}(x,E)\ge\check{R}_{k_a}$, so by the first line of
\eqref{5.2:eq-plain-fine-units-a} the part of the contour beyond $x$ has $d^{\wedge}$-length at least
$\frac{M}{M_1}\check{R}_{k_a}$. This gives $\check{R}_{k_a}\le\varrho(\Delta)$, $\varrho(\Delta)$
being the depth of $\Delta$ in $\mathsf{C}$ in the contour distance $d^{\wedge}$.

\emph{(ii), live arrest: the bound.} By the first part of hypothesis (c), at the step $k_a$,
\begin{equation}\label{5.2:eq-plain-fine-units-4}
\check{R}_{k_{0a}+1}\le\mathsf{R}^{\uparrow}(\check{\mathsf{T}}_{k_a})
\le C_{\uparrow}\,\check{\mathsf{T}}_{k_a}^{\,r_{\mathrm{pr}}}\le C_{\uparrow}\,\check{R}_{k_a}.
\end{equation}
The last inequality holds because $\check{R}_{k_a}$ is the least power of $L$ not below
$\check{\mathsf{T}}_{k_a}^{\,r_{\mathrm{pr}}}$. The second inequality follows from the bound on
$\mathsf{R}^{\uparrow}$ in hypothesis (c), at $\mathsf{T}=\check{\mathsf{T}}_{k_a}\ge1$, and the
elementary bound, valid for $\mathsf{T}\ge1$,
\begin{equation}\label{5.2:eq-plain-fine-units-5}
\mathsf{T}+1+\log\bigl(1+3\lambda C_N(\mathsf{T}+1)^{r_{\mathrm{pr}}}+2c_2\bigr)
\le\bigl(2+r_{\mathrm{pr}}+\log(1+3\lambda C_N+2c_2)\bigr)\mathsf{T}:
\end{equation}
raising \eqref{5.2:eq-plain-fine-units-5} to the power $r_{\mathrm{pr}}$ and multiplying by $L$
turns the bound of hypothesis (c) into
$\mathsf{R}^{\uparrow}(\mathsf{T})\le C_{\uparrow}\mathsf{T}^{r_{\mathrm{pr}}}$.
To verify \eqref{5.2:eq-plain-fine-units-5}: since $(\mathsf{T}+1)^{r_{\mathrm{pr}}}\ge1$, the
argument of the logarithm is at most $(1+3\lambda C_N+2c_2)(\mathsf{T}+1)^{r_{\mathrm{pr}}}$, so the
logarithm is at most
\begin{equation*}
\log(1+3\lambda C_N+2c_2)+r_{\mathrm{pr}}\log(\mathsf{T}+1)
\le\log(1+3\lambda C_N+2c_2)\,\mathsf{T}+r_{\mathrm{pr}}\mathsf{T},
\end{equation*}
using $\log(\mathsf{T}+1)\le\mathsf{T}$, $\mathsf{T}\ge1$ and
$\log(1+3\lambda C_N+2c_2)\ge0$; and $\mathsf{T}+1\le2\mathsf{T}$. With
$\check{R}_{k_a}\le\varrho(\Delta)$ from the collar claim, \eqref{5.2:eq-plain-fine-units-4} gives
the second line of \eqref{5.2:eq-plain-fine-units-a}.

\emph{(iii).} The cell at which $\mathsf{w}^{\mathrm{win},\vee}(\Delta')$ is attained lies at
$\operatorname{dist}_{n'}\le1$ from $\Delta'$. Suppose it lies in the window of a live arrest
$\mathfrak{a}$. Then $n'\le k_a$ by the regularity lemma for pointed units, and
$\operatorname{dist}_{k_a}(\Delta',X_{\mathfrak{a}})\le1$; by (ii),
$\mathsf{w}^{\mathrm{win},\vee}(\Delta')\le LC_{\uparrow}\check{R}_{k_a}$. If
$\operatorname{dist}_{k_a}(\Delta,X_{\mathfrak{a}})\le\frac12\check{R}_{k_a}$, then $\Delta$ lies in
the collar of the collar claim and $\varrho(\Delta)\ge\frac12\check{R}_{k_a}$, so that
\begin{equation*}
\mathsf{w}^{\mathrm{win},\vee}(\Delta')\le LC_{\uparrow}\check{R}_{k_a}
\le 2LC_{\uparrow}\varrho(\Delta)\le 2LC_{\uparrow}\varrho(\Delta)(2+i).
\end{equation*}
Otherwise $i\ge\frac12\check{R}_{k_a}-2$, that is $\check{R}_{k_a}\le2(2+i)$, and
$\mathsf{w}^{\mathrm{win},\vee}(\Delta')\le LC_{\uparrow}\check{R}_{k_a}\le2LC_{\uparrow}(2+i)$, which
gives \eqref{5.2:eq-plain-fine-units-3}. If the weight is not attained on a live arrest, it is that
of the running attempt, and $\mathsf{w}^{\mathrm{win},\vee}(\Delta')\le L\check{R}_{h+1}$ by (ii).
For the last clause: by hypothesis (d) the cells within one own unit of $\mathsf{C}$ have level at
most $h\le k_0$, while by hypothesis (e) the cells
carrying window weight of the running attempt have levels $>k_0$; hence, in this case,
$\mathsf{w}^{\mathrm{win},\vee}(\Delta')=1$ for $\Delta'\subset\mathsf{C}$.
\end{proof}

Hypothesis (a) is the per-cube bound for stored boundary-type terms together with the bound for
the boundary-type terms listed last in the storage format. The former is stated for all stored
terms of their format and every $\square'\in\pi_j$, and its proof uses, besides
\cite[(1.26)]{B-CMP116}, only statements independent of the position of the term; this is why the
first line of \eqref{5.2:eq-plain-fine-units-b} is taken for the whole class. The summation of
boundary terms at finer zones uses this bound for the whole class, with
$\mathfrak{N}(v)=2\mathfrak{N}^{\mathbb{C}}_{\varepsilon_k}(v)=2\mathfrak{N}^{A,\mathbb{C}}(v)$.
The second part of hypothesis (b) rests on \cite[(2.41)]{B-CMP119}, \cite[(1.63)]{B-CMP122b} and
\cite{B-CMP122a}; a term surviving a completed exit from a region $X'$ belongs to the class of terms
whose localisation domain is disjoint from $X'$, $X_F\cap X'=\emptyset$, and its anchor touches an
unhealed part, which lies inside the current $Z_j$.

\begin{lemma}[Plain mass meeting one cell]\label{5.2:lem-plain-mass-cell}
Use the notation of Lemma~\ref{5.2:lem-plain-fine-units}: the level $k$, the function
$\mathsf{b}(s)$ with the growth bound $\mathsf{b}(s+u)\le C_{\mathsf{b}}(1+u)^{p_{\mathsf{b}}}\mathsf{b}(s)$
of \eqref{5.2:eq-plain-fine-units-b}, the weights
$W(v)=\mathfrak{N}^{\flat}(v)e^{r d_k(K_v)}$, the component $C$ with its core $\mathsf{C}$ and the
level $h$ attached to it, the radii $\check{R}_j$, the depth $\varrho(\Delta)$ of a cell, the deep
and shallow plain units, the window weights
$\mathsf{w}^{\mathrm{win}}$, and the constants $\mathsf{K}^{\flat}$, $C_{\uparrow}$, $C^{<}$,
$C_{\circledast}$, $C_{\mathsf{b}}$, $p_{\mathsf{b}}$. Let $\Delta\subset C$ be a cell of level $n$,
and let
$\mathsf{S}^{\mathrm{dp}}(\Delta)$, $\mathsf{S}^{\mathrm{sh}}(\Delta)$ be the sums of $W(v)$ over the
deep, respectively shallow, units $v\in\mathfrak{L}^{\mathrm{pl}}$ whose hull $\hat X_v$ meets $\Delta$.
Put
\begin{equation*}
\mathsf{w}^{\oplus}:=
\begin{cases}
2LC_{\uparrow}\varrho(\Delta) & \text{for } \mathsf{S}^{\mathrm{dp}}, \ \Delta\subset\mathsf{C},\\
2LC_{\uparrow}\check{R}_{h+1} & \text{for } \mathsf{S}^{\mathrm{sh}}, \ \Delta\not\subset\mathsf{C}.
\end{cases}
\end{equation*}
Then, with $\mathsf{S}^{\cdot}$ standing for either sum in its case,
\begin{equation}\label{5.2:eq-plain-mass-cell-a}
\begin{split}
\mathsf{S}^{\cdot}(\Delta)\le\Big[&(k-n+1)\big(1+C_{\circledast}\mathsf{K}^{\flat}(\mathsf{w}^{\oplus})^3\big)
+C^{<}\\
&+2C_{\mathsf{b}}\big(1+\log_L\mathsf{w}^{\mathrm{win}}(\Delta)\big)^{p_{\mathsf{b}}}
\mathsf{w}^{\mathrm{win}}(\Delta)^d\Big]\,\mathsf{b}(k-n).
\end{split}
\end{equation}
\end{lemma}

\begin{proof}
The class $\mathfrak{L}^{\mathrm{pl}}$ is the disjoint union of the boundary-type units
$\mathfrak{L}^{B}$ and the pointed units $\mathfrak{L}^{\mathrm{pt}}$. We treat the units
$F\in\mathfrak{L}^{B}$ according to their level $j=j_F$, and then the pointed units.

\emph{Boundary-type units of level $j\ge n$.} The domain $X_F$ is a union of
$\pi_j$-cubes and meets $\Delta$; since $j\ge n$, the $\pi_j$-cube $\square'$ containing $\Delta$
lies in $X_F$. The bound \eqref{5.2:eq-plain-fine-units-b}, applied at this $\square'$ and with the
factor $e^{\varsigma_0 d_j(X_F)\mathbb{1}[j<k]}\ge 1$ dropped, bounds the sum of $W(F)$ over these
$F$ by $\mathsf{b}(k-j)$. As $\mathsf{b}$ is non-decreasing and $k-j\le k-n$, each level contributes
at most $\mathsf{b}(k-n)$, and there are $k-n+1$ levels $j\in\{n,\dots,k\}$: in total at most
$(k-n+1)\mathsf{b}(k-n)$.

\emph{Boundary-type units of level $j=n-u<n$, no window.} By Lemma~\ref{5.2:lem-plain-fine-units},
every $F\in\mathfrak{L}^{B}$ of level $j$ has a $\pi_j$-cube meeting $Z_j^{\boxplus}$; since $X_F$
meets $\Delta$ and $j<n$, it also has a $\pi_j$-cube contained in $\Delta$. By clause~(i) of
Lemma~\ref{5.2:lem-plain-fine-units}, when no window occurs, $\operatorname{dist}_n(Z_j^{\boxplus},\Delta)\ge 1$,
so these two cubes are at least $L^u-1$ level-$j$ units apart. By the lower bound of the size $d_j$
of a localisation domain by the distance of two of its cubes, $d_j(X_F)\ge L^u-1$. The cell
$\Delta$ contains $L^{du}$ cubes of $\pi_j$; for each of them, \eqref{5.2:eq-plain-fine-units-b}
(with $j<k$, so that the factor $e^{\varsigma_0 d_j(X_F)}$ is present) gives
\begin{equation*}
\sum_{F} W(F)\le e^{-\varsigma_0(L^u-1)}\,\mathsf{b}(k-n+u),
\end{equation*}
the sum running over the $F$ of level $j$ whose domain contains that cube. By the growth bound in
\eqref{5.2:eq-plain-fine-units-b},
$\mathsf{b}(k-n+u)\le C_{\mathsf{b}}(1+u)^{p_{\mathsf{b}}}\mathsf{b}(k-n)$. Summing over the $L^{du}$
cubes and over $u\ge 1$ gives at most $C^{<}\mathsf{b}(k-n)$, by the definition of $C^{<}$.

\emph{Boundary-type units in a window.} Otherwise clause~(i) of Lemma~\ref{5.2:lem-plain-fine-units}
gives $\Delta\subset X_{\mathfrak{a}}$ and $k_{0\mathfrak{a}}\le j<n\le k_{\mathfrak{a}}$ for an attempt
$\mathfrak{a}$. For each such level $j$ the cell $\Delta$ contains $L^{d(n-j)}$ cubes of $\pi_j$,
each $X_F$ contains one of them, and \eqref{5.2:eq-plain-fine-units-b} bounds the contribution of
each cube by $\mathsf{b}(k-j)$. The window has at most $\log_L\mathsf{w}^{\mathrm{win}}(\Delta)$
levels below $n$, that is, $n-k_{0\mathfrak{a}}\le\log_L\mathsf{w}^{\mathrm{win}}(\Delta)$. Hence,
for $u=n-j$, the growth bound in \eqref{5.2:eq-plain-fine-units-b} gives
\begin{equation*}
\mathsf{b}(k-j)\le C_{\mathsf{b}}(1+u)^{p_{\mathsf{b}}}\mathsf{b}(k-n)
\le C_{\mathsf{b}}\big(1+\log_L\mathsf{w}^{\mathrm{win}}(\Delta)\big)^{p_{\mathsf{b}}}\mathsf{b}(k-n),
\end{equation*}
and, since $L^d\ge 2$, the geometric series satisfies
\begin{equation*}
\sum_{u=1}^{n-k_{0\mathfrak{a}}}L^{du}\le 2L^{d(n-k_{0\mathfrak{a}})}
\le 2\,\mathsf{w}^{\mathrm{win}}(\Delta)^d.
\end{equation*}
Together,
\begin{equation*}
\sum_{j=k_{0\mathfrak{a}}}^{n-1}L^{d(n-j)}\mathsf{b}(k-j)
\le 2\,\mathsf{w}^{\mathrm{win}}(\Delta)^d\,C_{\mathsf{b}}
\big(1+\log_L\mathsf{w}^{\mathrm{win}}(\Delta)\big)^{p_{\mathsf{b}}}\,\mathsf{b}(k-n).
\end{equation*}

\emph{Pointed units.} Let $v\in\mathfrak{L}^{\mathrm{pt}}$ have its anchor $y_v$ in a cell
$\Delta'$ of level $n'$, and put $i:=\operatorname{dist}_{n'}(\Delta',\Delta)$. Since
$y_v\in\Delta'$ and $\hat X_v$ meets $\Delta$, $\operatorname{rad}(v)\ge i$. By the flat collar bound
for pointed units, together with clause~(iii) of Lemma~\ref{5.2:lem-plain-fine-units} (which applies
in the deep case because the anchors of deep units lie in $\mathsf{C}$), the contribution of these
units is at most
\begin{equation*}
\mathsf{K}^{\flat}(\mathsf{w}^{\oplus})^3\sum_{\Delta'}(2+i)^3e^{-\varsigma_1 i}.
\end{equation*}
For an integer $i\ge 0$ we count the cells $\Delta'$ with
$\operatorname{dist}_{n'}(\Delta',\Delta)\in[i,i+1)$.
For $n'\ge n$ there are at most $(2i+5)^d$ cubes of $\pi_{n'}$, on each of $k-n+1$ levels. For
$n'=n-1$ there are at most $(L+2i+4)^d$. For $n'=n-u$ with $u\ge 2$, by the separation of cells of
distant levels such a $\Delta'$ is at distance at least $L^u$ of its own units from $\Delta$, so
$L^u\le i$; there are at most $(3i+4)^d$ such cells for each $u$, and the number of $u$ with
$L^u\le i$ is at most $i$. Since $k-n+1\ge 1$, the sum over $\Delta'$ is at most
\begin{equation*}
(k-n+1)\sum_{i\ge 0}\big[(2i+5)^d+(L+2i+4)^d+i(3i+4)^d\big](2+i)^3e^{-\varsigma_1 i}
=(k-n+1)\,C_{\circledast},
\end{equation*}
and as $\mathsf{b}\ge 1$ the pointed units contribute at most
$(k-n+1)C_{\circledast}\mathsf{K}^{\flat}(\mathsf{w}^{\oplus})^3\mathsf{b}(k-n)$.

Adding the four contributions gives \eqref{5.2:eq-plain-mass-cell-a}.
\end{proof}

\begin{remark}
The factor $k-n+1$ in \eqref{5.2:eq-plain-mass-cell-a} is genuinely present: under a window every
coarser level stands within two of its own units. It is paid for by the depth, as in clause~(i) of
the anchoring lemma.
\end{remark}

\begin{remark}
A bound on the boundary-type units by a constant per cell, uniform in the cell, fails inside a
window, where a factor of order $\check{R}^d$ appears. The factor
$\mathsf{w}^{\mathrm{win}}(\Delta)^d$ in \eqref{5.2:eq-plain-mass-cell-a} is that factor, with
$\check{R}$ an old radius. It is paid by the window bound \eqref{5.2:eq-plain-fine-units-a}, which,
through clause~(iii) of Lemma~\ref{5.2:lem-plain-fine-units}, also pays for the threshold weights
$\bar{\mathsf{w}}^3$ on the pointed units.
\end{remark}

\begin{lemma}[The deep sum over the core]\label{5.2:lem-deep-sum}
Work in the setting and notation of Lemma~\ref{5.2:lem-plain-fine-units} and
Lemma~\ref{5.2:lem-plain-mass-cell}. In particular:
\begin{itemize}
\item $C$ is a component and $\mathsf{C}=\mathsf{C}_C$ is its core;
\item $h$ is the level of the core, $\mathsf{C}=Z_h'^{\sim}\cap C$, and $k$ is the current level;
\item $W(v)$ is the weight of a unit $v$, $\hat D_v$ its depth and $\hat X_v$ its hull;
\item a unit $v$ is deep in $C$ in the sense of Lemma~\ref{5.2:lem-plain-mass-cell};
\item $\mathsf{b}(\cdot)$ is the non-decreasing per-cube function, and $C_{\mathsf{b}}$ and
$p_{\mathsf{b}}$ are the constants of the growth bound \eqref{5.2:eq-plain-fine-units-b};
\item $C_{\uparrow}$ is the constant of \eqref{5.2:eq-plain-fine-units-a};
\item $C^{<}$, $C_{\circledast}$ and $\mathsf{K}^{\flat}$ are the constants of the pile bound
\eqref{5.2:eq-plain-mass-cell-a}, and $N=N(g_k)$ is the number of preparatory stages;
\item $C_{\mathrm{cell}}$ is the constant of part (b) of the cell-sum lemma;
\item $\delta_P=\tfrac18\delta_0$, and $d^{\wedge}$ is the contour distance.
\end{itemize}
Put
\begin{equation}\label{5.2:eq-deep-sum-1}
\begin{split}
C^{\mathrm{pl}}_{\mathsf{a}} := C_{\mathsf{b}}\sup_{\varrho\ge 1}
\Big[&\varrho\big(1+C_{\circledast}\mathsf{K}^{\flat}(2LC_{\uparrow}\varrho)^3\big)+C^{<}\\
&+2C_{\mathsf{b}}\big(2+\log_L(C_{\uparrow}\varrho)\big)^{p_{\mathsf{b}}}(LC_{\uparrow}\varrho)^d\Big]
(1+\varrho)^{p_{\mathsf{b}}}e^{-\varrho}.
\end{split}
\end{equation}
Then
\begin{equation}\label{5.2:eq-deep-sum-2}
\begin{split}
S^{\mathrm{dp}}(C)&:=\sum_{v\ \text{deep in}\ C}W(v)\,e^{-\frac14\delta_P\hat D_v}\\
&\le e\,C^{\mathrm{pl}}_{\mathsf{a}}\,C_{\mathrm{cell}}\cdot 2d\,(104L\check{R}_{h+1}+2)^{d-1}
(2+N)\,\mathsf{b}(N+1).
\end{split}
\end{equation}
\end{lemma}

\begin{proof}
The argument is that of the deep-sum lemma for anchored units, with the pile
$\mathsf{S}^{\mathrm{dp}}$ of Lemma~\ref{5.2:lem-plain-mass-cell} in place of the cell sum used
there; we carry it out.

\emph{Reduction to cells.} Let $v$ be deep in $C$. Choose a point of $\hat X_v\cap\mathsf{C}$ that
is nearest to $\mathsf{C}^c$ in the distance $d^{\wedge}$. This point lies in a cell
$\Delta_v\subset\mathsf{C}$ which is met by $\hat X_v$, and
\begin{equation*}
d^{\wedge}(\Delta_v,\mathsf{C}^c)\le\hat D_v .
\end{equation*}
Hence
\begin{equation*}
e^{-\frac14\delta_P\hat D_v}\le e^{-\frac14\delta_P d^{\wedge}(\Delta_v,\mathsf{C}^c)} .
\end{equation*}
Because $\hat X_v$ meets $\Delta_v$, the term $W(v)$ is one of the terms of the pile
$\mathsf{S}^{\mathrm{dp}}(\Delta_v)$ of Lemma~\ref{5.2:lem-plain-mass-cell}. Grouping the deep
units according to the cell $\Delta_v$ therefore gives
\begin{equation}\label{5.2:eq-deep-sum-3}
S^{\mathrm{dp}}(C)\le\sum_{\Delta\subset\mathsf{C}}
e^{-\frac14\delta_P d^{\wedge}(\Delta,\mathsf{C}^c)}\,\mathsf{S}^{\mathrm{dp}}(\Delta).
\end{equation}

\emph{The pile of one cell.} Fix a cell $\Delta\subset\mathsf{C}$ of level $n$, and let
$\varrho=\varrho(\Delta)\ge 1$ be its depth parameter as in Lemma~\ref{5.2:lem-plain-mass-cell}.
In this proof only, write $w$ for the common width, in the distance $d^{\wedge}$, of the zones
(not the complex source $w$). The cells of $\mathsf{C}$
that touch $\mathsf{C}^c$ have level $h$. A contour from $\mathsf{C}^c$ to $\Delta$ therefore
crosses the zones of levels $h-1,\dots,n+1$, each of width $w$. As in the deep-sum lemma for
anchored units, this yields
\begin{equation}\label{5.2:eq-deep-sum-4}
k-n+1\le N+1+\varrho\le\varrho(2+N).
\end{equation}
The second inequality in \eqref{5.2:eq-deep-sum-4} is equivalent to $(N+1)(\varrho-1)\ge 0$, that
is, to $\varrho\ge 1$. Since $\mathsf{b}$ is non-decreasing, the growth bound
\eqref{5.2:eq-plain-fine-units-b} then gives
\begin{equation}\label{5.2:eq-deep-sum-5}
\mathsf{b}(k-n)\le C_{\mathsf{b}}(1+\varrho)^{p_{\mathsf{b}}}\,\mathsf{b}(N+1).
\end{equation}
We now insert three ingredients into the pile bound \eqref{5.2:eq-plain-mass-cell-a} for
$\mathsf{S}^{\mathrm{dp}}(\Delta)$: the bounds \eqref{5.2:eq-deep-sum-4} and
\eqref{5.2:eq-deep-sum-5}, and the window-weight bound \eqref{5.2:eq-plain-fine-units-a}, which
controls the window weight of $\Delta$ through $C_{\uparrow}\varrho$. For this proof only, denote
by $\Phi(\varrho)$ the square bracket in \eqref{5.2:eq-deep-sum-1}. The insertions turn the pile
bound into
\begin{equation*}
\mathsf{S}^{\mathrm{dp}}(\Delta)\le C_{\mathsf{b}}\,\Phi(\varrho)\,(1+\varrho)^{p_{\mathsf{b}}}
(2+N)\,\mathsf{b}(N+1).
\end{equation*}
For every $\varrho\ge1$, the definition of
$C^{\mathrm{pl}}_{\mathsf{a}}$ as a supremum gives
\begin{equation*}
C_{\mathsf{b}}\Phi(\varrho)(1+\varrho)^{p_{\mathsf{b}}}\le e^{\varrho}C^{\mathrm{pl}}_{\mathsf{a}} .
\end{equation*}
With this the pile bound becomes
\begin{equation}\label{5.2:eq-deep-sum-6}
\mathsf{S}^{\mathrm{dp}}(\Delta)\le e^{\varrho}\,C^{\mathrm{pl}}_{\mathsf{a}}\,(2+N)\,\mathsf{b}(N+1).
\end{equation}

\emph{Spending the decay.} Split the rate in \eqref{5.2:eq-deep-sum-3} as
$\tfrac14=\tfrac18+\tfrac18$. By the zone-width condition, which fixes the zone width $w$, we have
$\tfrac18\delta_P w\ge 2$. By the
crossing of zones counted in the preceding step, as in the deep-sum lemma for anchored units,
this gives
\begin{equation}\label{5.2:eq-deep-sum-7}
e^{-\frac18\delta_P d^{\wedge}(\Delta,\mathsf{C}^c)}\le e^{1-\varrho}.
\end{equation}
This factor cancels the factor $e^{\varrho}$ in \eqref{5.2:eq-deep-sum-6} up to the factor $e$.
Combining \eqref{5.2:eq-deep-sum-3}, \eqref{5.2:eq-deep-sum-6} and \eqref{5.2:eq-deep-sum-7}, we
obtain
\begin{equation}\label{5.2:eq-deep-sum-8}
S^{\mathrm{dp}}(C)\le e\,C^{\mathrm{pl}}_{\mathsf{a}}(2+N)\,\mathsf{b}(N+1)
\sum_{\Delta\subset\mathsf{C}}e^{-\frac18\delta_P d^{\wedge}(\Delta,\mathsf{C}^c)} .
\end{equation}
The remaining sum is bounded by part (b) of the cell-sum lemma. That part applies at every rate
not smaller than $\delta_P/16$, and here the rate is $\tfrac18\delta_P\ge\delta_P/16$. We apply it
to the set $\Sigma:=\operatorname{cl}(\mathsf{C}^c)$. Since $\mathsf{C}^c\subset\Sigma$, we have
$d^{\wedge}(\Delta,\Sigma)\le d^{\wedge}(\Delta,\mathsf{C}^c)$, so each summand is at most
$e^{-\frac18\delta_P d^{\wedge}(\Delta,\Sigma)}$. The cell-sum lemma bounds the sum of these
majorants by $C_{\mathrm{cell}}\,\#\partial^{+}\Sigma$, where $\partial^{+}\Sigma$ denotes the set
of cells not in $\Sigma$ that are adjacent to $\Sigma$; these are cells of $\mathsf{C}$ touching
$\mathsf{C}^c$. The count of these cells in part (g) of the geometry lemma for the zones then
gives
\begin{equation}\label{5.2:eq-deep-sum-9}
\begin{split}
\sum_{\Delta\subset\mathsf{C}}e^{-\frac18\delta_P d^{\wedge}(\Delta,\mathsf{C}^c)}
&\le\sum_{\Delta\subset\mathsf{C}}e^{-\frac18\delta_P d^{\wedge}(\Delta,\Sigma)}
\le C_{\mathrm{cell}}\,\#\partial^{+}\Sigma\\
&\le C_{\mathrm{cell}}\cdot 2d\,(104L\check{R}_{h+1}+2)^{d-1}.
\end{split}
\end{equation}
Inserting \eqref{5.2:eq-deep-sum-9} into \eqref{5.2:eq-deep-sum-8} gives
\eqref{5.2:eq-deep-sum-2}.
\end{proof}

The right-hand side of \eqref{5.2:eq-deep-sum-2} contains nothing about the core's past. Neither
its volume, nor the number, nesting or ages of the arrests it contains, nor any of the finer
levels $j$ enters the bound.

\begin{lemma}[The shallow sum off the core]\label{5.2:lem-shallow-sum}
Let $C$ be a component with core $\mathsf{C}=\mathsf{C}_C$, let $h$, $k$, $N$, the piles
$\mathsf{S}^{\mathrm{sh}}(\Delta)$ and the class of units shallow in $C$ be as in
Lemma~\ref{5.2:lem-plain-mass-cell}, and let $W(v)$
and $\mathsf{b}(\cdot)$ be the weight of a plain unit and the per-cube function fixed in
Lemma~\ref{5.2:lem-plain-fine-units}, \eqref{5.2:eq-plain-fine-units-b}. Let the regions
$Z''_j$ and $Z'_j$, the level $k_0$ with $h\le k_0\le k$, the exponent $N_0$, which satisfies
$L^{N_0}=L\check{R}_{k_0+1}$, and the box $\Lambda_C$ be those of the zone geometry of $C$ fixed
earlier in this chapter. Then
\begin{equation}\label{5.2:eq-shallow-sum-1}
\begin{split}
S^{\mathrm{sh}}(C) &:= \sum_{v\ \text{shallow in}\ C} W(v)\\
&\le (N+1)\,(110L\check{R}_{h+1})^d\Big[(N+1)\big(1+C_{\circledast}\mathsf{K}^{\flat}
(2LC_{\uparrow}\check{R}_{h+1})^3\big)\\
&\qquad\qquad + C^{<} + 2C_{\mathsf{b}}(1+N)^{p_{\mathsf{b}}}(L\check{R}_{h+1})^d\Big]\,
\mathsf{b}(N).
\end{split}
\end{equation}
\end{lemma}

\begin{proof}
Let $v$ be shallow in $C$, so that $X_v\not\subset\mathsf{C}$. Since $X_v\subset C$ and
$X_v\not\subset\mathsf{C}$, the set $X_v$, and hence the hull $\hat X_v\supseteq X_v$, meets a
cell $\Delta\subset\operatorname{cl}(C\setminus\mathsf{C})$. The pile
$\mathsf{S}^{\mathrm{sh}}(\Delta)$ of Lemma~\ref{5.2:lem-plain-mass-cell} is a sum of $W(v)$
over the units $v$ whose hull $\hat X_v$ meets $\Delta$. Every shallow unit therefore occurs
in at least one of these piles, and all weights $W(v)$ are non-negative; hence
\begin{equation}\label{5.2:eq-shallow-sum-2}
S^{\mathrm{sh}}(C)\le\sum_{\Delta\subset\operatorname{cl}(C\setminus\mathsf{C})}
\mathsf{S}^{\mathrm{sh}}(\Delta).
\end{equation}

We next locate the levels of these cells. The set $Z''_{h+1}\setminus\mathsf{C}$ lies in zone
$h$, and the cells of $C\setminus Z''_{h+1}$ have levels $n$ with $h<n\le k$. Consequently every
cell $\Delta\subset\operatorname{cl}(C\setminus\mathsf{C})$ has a level
$n=\operatorname{lev}\Delta$ with $h\le n\le k$, and $k-n\le N$. No such cell lies in the window
of a live arrest, because
$X_{\mathfrak{a}}\subset\mathsf{C}$ for every live arrest $\mathfrak{a}$. Therefore the window
weight of each such cell satisfies
\begin{equation*}
\mathsf{w}^{\mathrm{win}}(\Delta)\le L\check{R}_{h+1},\qquad
\log_L\mathsf{w}^{\mathrm{win}}(\Delta)\le N .
\end{equation*}
Reading the pile bound \eqref{5.2:eq-plain-mass-cell-a} of
Lemma~\ref{5.2:lem-plain-mass-cell} with $k-n\le N$,
$\mathsf{w}^{\mathrm{win}}(\Delta)\le L\check{R}_{h+1}$ and
$\log_L\mathsf{w}^{\mathrm{win}}(\Delta)\le N$, and using that $\mathsf{b}$ of
\eqref{5.2:eq-plain-fine-units-b} is non-decreasing, we obtain for each such cell
\begin{equation*}
\begin{split}
\mathsf{S}^{\mathrm{sh}}(\Delta)\le{}&\Big[(N+1)\big(1+C_{\circledast}\mathsf{K}^{\flat}
(2LC_{\uparrow}\check{R}_{h+1})^3\big)\\
&\quad + C^{<} + 2C_{\mathsf{b}}(1+N)^{p_{\mathsf{b}}}(L\check{R}_{h+1})^d\Big]\,
\mathsf{b}(N).
\end{split}
\end{equation*}

It remains to count the cells of $\operatorname{cl}(C\setminus\mathsf{C})$, zone by zone, each
zone measured in its own units.
\begin{itemize}
\item Zone $n=k$: its side is at most $100\check{R}_k$, by the geometric bound on the finest
zone taken from \cite{B-CMP122a}.
\item Zones $k_0\le n<k$: they lie inside $Z''_k\subset\Lambda_C$, whose side is at most
$110L^{N_0}=110L\check{R}_{k_0+1}$.
\item Zones $h\le n<k_0$: zone $n$ lies inside $Z''_{n+1}=Z'^{\sim 3}_n$. By condition (ii) of
the arrest conditions, which in the words of \cite{B-CMP122a} requires that ``the previous
regions contained in it satisfy the condition (i) on the corresponding scales'', the side of
$Z''_{n+1}$ is bounded by
\begin{equation*}
100\check{R}_n+8L\check{R}_{n+1}\le 108L\check{R}_{n+1},
\end{equation*}
where the comparison $\check{R}_n\le L\check{R}_{n+1}$ of consecutive radii has been used.
\end{itemize}
Since $\check{R}_j=R(\min_{i\le j}x_i)$ and $R$ is non-decreasing, $\check{R}_j$ is
non-increasing in $j$. With $h\le k_0$ and $h\le n$, the second and third sides above are
therefore at most $110L\check{R}_{h+1}$; so is the first, $100\check{R}_k$, since the
comparison $\check{R}_j\le L\check{R}_{j+1}$ used in the third item, taken at $j=k$, and the
same monotonicity give, for every $h\le k$,
\begin{equation*}
100\check{R}_k\le 100L\check{R}_{k+1}\le 110L\check{R}_{h+1}.
\end{equation*}
Thus a zone contributes at most $(110L\check{R}_{h+1})^d$
cells. The levels
$n$ satisfy $k-N\le n\le k$, so there are at most $N+1$ zones. Hence
$\operatorname{cl}(C\setminus\mathsf{C})$ contains at most $(N+1)(110L\check{R}_{h+1})^d$ cells.
Inserting this count and the bound on each pile into \eqref{5.2:eq-shallow-sum-2} gives
\eqref{5.2:eq-shallow-sum-1}.
\end{proof}

\begin{proof}[Proof of clauses (a)--(e) of the lemma on the pieces of plain units whose label
stays in the large-field region]\label{5.2:prf-deep-shallow-sums}
Besides the hypotheses of the lemma, the argument uses in particular the deep sum over the core
(Lemma~\ref{5.2:lem-deep-sum}), the shallow sum off the core (Lemma~\ref{5.2:lem-shallow-sum}),
the plain mass meeting one cell (Lemma~\ref{5.2:lem-plain-mass-cell}), and the lemma on the pieces of
units whose label leaves the large-field region (Lemma~\ref{5.2:lem-exterior-pieces}). It also uses
the following statements of this chapter: the source-anchoring statement, clauses
(i) and (ii); the source-carrier statement at the site $(L)$; the geometry lemma, clause (f); the
summation lemma for the four-link expansion, clauses (a), (c) and (e), and the statement
accompanying it that lists the cases of the reading property and fixes the rule for carrying
majorants; the three-factor bound,
clauses (b) and (c); the pointed three-factor proposition; and the floor condition on
$\kappa_1$. Throughout, $\mathfrak{p}=(v;\mathsf{y}_1,\dots,\mathsf{y}_4)$ is a piece of a plain
unit $v$ whose anchor $K_v$ lies in the component $C$, and $Y_4=Y_4(\mathfrak{p})$.

\emph{Clause (a).} Holomorphy of $v$ on its domain is, for the units of the class
$\mathfrak{L}^{B}$, the source-carrier statement at $(L)$, and, for the pointed units, clause (f)
of the lemma, the carrier value bound on $C$. Clause (i) of the source-anchoring statement, applied
at each of the four
links in turn, shows that a component of $\mathsf{y}_t$ which contains no source cell, or which
does not meet $\mathcal{R}_{t-1}$, contributes zero. Part (iii) of clause (a) is the locality of
the pieces. Clause (a) of the summation lemma applies to the pieces. The label
$\operatorname{lab}_4$ of the piece
contains $Y_4\setminus Z$ and is contained in $Y_4$; hence
\begin{equation}\label{5.2:eq-deep-shallow-sums-1}
d_{k,Z}(Y_4)\le d_k(\operatorname{lab}_4).
\end{equation}

We turn to the reading property. At the first link, a member meeting $\mathcal{R}_0$, which is
contained in $C$, has level $k$ by clause (f) of the geometry lemma, and its parent lies within
distance $1$ of $K_v$. Indeed $\hat X_v\neq X_v$ only for short singles and for smooth cores, and in
these two cases $\hat X_v$ differs from $X_v$ only by boxes whose sides are at most one cube of
$\pi_{n-1}$ and which meet
$X_v$. For a source cell at a hub met by $\mathcal{R}_0$ the parent lies within distance
$h_{\sharp}+1$ of $K_v$, and $h_{\sharp}+1\le c_\rho$, where $c_\rho$ is the constant of the
reading property in the summation
lemma and $h_{\sharp}$ the hub constant of the source-anchored families. For the links $t\ge 2$ the
reading property is the list of cases of the statement accompanying the
summation lemma, whose clauses concerning strips are void here. Finally
$\mathcal{R}_0\subset C\subset Y_4$. The chart clause is clause (c) of the three-factor bound (its
sixth and seventh items), which holds for every class of units alike; each piece carries one disc.

\emph{Clause (b).} For a shallow unit, the bound is clause (ii) of the source-anchoring statement,
applied with the holomorphy of $v$ established in (a) for units of the class $\mathfrak{L}^B$,
respectively with clause (f) of the lemma at $c:=f_v(1)$ for pointed units, $f_v$ being the
function that enters that clause. For a deep
unit of the class $\mathfrak{L}^B$, the bound is clause (b) of the three-factor bound, together
with $e^{-\frac12\delta_P D}\le e^{-\frac14\delta_P D}$. For a deep pointed unit, the bound is the
pointed three-factor proposition.

\emph{Clause (c).} Let $\Delta\subset Y_4$ be a cell with $\Delta\not\subset Z$. By (a),
$\Delta$ lies in $\operatorname{lab}_4$, while $K_v\subset Z$. Hence some $\mathsf{y}_t\neq\emptyset$
contains a source cell, that is, a cube of level $k$ sharing a wall with $\Lambda^{\sim}$ or with
$\Lambda$; by clause (f) of the geometry lemma such a cube lies at distance at least
$9\check{R}_k-1$ from $Z^c$. By the inequality that bounds the tree length $d_k$ of a set from
below in terms of its distance from $Z^c$, it follows that
\begin{equation}\label{5.2:eq-deep-shallow-sums-2}
d_k(\operatorname{lab}_4)\ge D:=\check{R}_k-2 .
\end{equation}
Define the reading families
\begin{equation*}
\mathfrak{L}^0:=\{(K_v,\mathcal{R}_0,\|v\|_{\star})\},\qquad
\mathfrak{L}^t:=\mathfrak{L}^{t-1}\cup(\mathfrak{L}^{t-1})',\quad t=1,\dots,4,
\end{equation*}
where $(\mathfrak{L}^{t-1})'$ is the family produced from $\mathfrak{L}^{t-1}$ at link $t$ by the
summation lemma; there is a single bracket at the first link. Here $\mathsf{Q}_r[\cdot]$ is the
majorant functional at rate $r$ of the summation lemma, and $C_\Sigma$, $C_e$ are the constants of
its clauses (c) and (e). The rate $a$, the size $\|v\|_{\star}$ of a unit and the plain majorant
functional $\mathsf{Q}^{\star,\mathrm{pl}}_r$ are those fixed in the statement of the lemma. By
\eqref{5.2:eq-deep-shallow-sums-1} and \eqref{5.2:eq-deep-shallow-sums-2}, the left member of (c)
is at most $e^{-D}\mathsf{Q}_{a+1}[\mathfrak{L}^4]$. Clauses (c) and (e) of the summation lemma give,
for $t=1,\dots,4$ at the rate $r=a+9-2t$,
\begin{equation}\label{5.2:eq-deep-shallow-sums-3}
\mathsf{Q}_r[\mathfrak{L}^t]\le(1+C_\Sigma C_e)\,\mathsf{Q}_{r+2}[\mathfrak{L}^{t-1}] .
\end{equation}
Applying \eqref{5.2:eq-deep-shallow-sums-3} four times, at the rates $a+1$, $a+3$, $a+5$, $a+7$,
we obtain
\begin{equation}\label{5.2:eq-deep-shallow-sums-4}
\mathsf{Q}_{a+1}[\mathfrak{L}^4]\le(1+C_\Sigma C_e)^4\,\mathsf{Q}_{a+9}[\mathfrak{L}^0].
\end{equation}
Each application is admissible because the floor condition on $\kappa_1$ gives
$\kappa_1-1\ge c_l(a+9)+\lambda_*$, and $a+9\le\frac85\kappa$ since $\kappa\ge180$, which is among
the conditions collected in $\kappa_d$. The majorants are carried by the rule of the statement
accompanying the summation lemma: the
intermediate reading families of a deep unit $v$ carry $\|v\|_{\star}$; only the tuples whose
label leaves $Z$ are summed, and for those tuples the bound of (b) is the bound entering the sum.
Finally $\mathsf{Q}_r[\mathfrak{L}^0]=\mathsf{Q}^{\star,\mathrm{pl}}_r$. A cube $Q\subset K_v$
determines the component $C$, so that for fixed $Q$ the sum runs over the class of units of one
component, and the whole class of that component is bounded by
$\mathsf{Q}^{\star,\mathrm{pl}}_{a+9}$. Together with \eqref{5.2:eq-deep-shallow-sums-4} this proves
(c).

\emph{Clause (d).} For $\mathsf{T}\ge\mathsf{T}_\gamma$ put
$\mathsf{N}:=C_N(\mathsf{T}+1)^{r_{\mathrm{pr}}}$,
$\mathsf{R}:=\mathsf{R}^{\uparrow}(\mathsf{T})$, and
\begin{equation}\label{5.2:eq-deep-shallow-sums-5}
\mathsf{b}^{\mathsf{T}}(s):=2K(d)B_0+\mathsf{N}C_0^{\sharp}
+(s+1)\bigl[C_NC_0^{\sharp}(\mathsf{T}+1+c_s)^{r_{\mathrm{pr}}}+b_\delta\bigr],
\end{equation}
and define
\begin{equation}\label{5.2:eq-deep-shallow-sums-a}
\begin{split}
\mathsf{P}_{\mathrm{int}}(\mathsf{T})
:={}&(1+C_\Sigma C_e)^4\Bigl\{C^{+}e\,C^{\mathrm{pl}}_{\mathsf{a}}C_{\mathrm{cell}}\,2d\,
(104L\mathsf{R}+2)^{d-1}(2+\mathsf{N})\,\mathsf{b}^{\mathsf{T}}(\mathsf{N}+1)\\
&+(\mathsf{N}+1)(110L\mathsf{R})^d\Bigl[(\mathsf{N}+1)\bigl(1+C_{\circledast}\mathsf{K}^{\flat}
(2LC_{\uparrow}\mathsf{R})^3\bigr)+C^{<}\\
&\qquad+2C_{\mathsf{b}}(1+\mathsf{N})^{p_{\mathsf{b}}}(L\mathsf{R})^d\Bigr]
\mathsf{b}^{\mathsf{T}}(\mathsf{N})\Bigr\}\\
&+(1+C_\Sigma C_e)^3\,e\,S_d\bigl[\mathsf{K}^{0\prime}+\mathsf{K}^{\flat,\mathrm{ex}}
+\mathsf{N}C_0^{\sharp}\bigr].
\end{split}
\end{equation}
Here $C^{+}$ is a constant of the order of choice; like every constant in
\eqref{5.2:eq-deep-shallow-sums-a}, it is fixed before $\gamma$.
The last summand is the one of clause (d) of Lemma~\ref{5.2:lem-exterior-pieces}; its term
$\mathsf{K}^{\flat,\mathrm{ex}}$, which equals
$e^{(a+9)c_lc_{\mathfrak{h}}}\,27S_d'\,\mathsf{K}^{\flat}$ and is a number fixed after $M$, is the
contribution of the exposed pointed units (clauses (c) and (d) of that lemma, with the constants
defined in \eqref{5.2:eq-exterior-pieces-a}).

Take $\mathsf{T}:=\check{\mathsf{T}}_k$. By the standing definitions of $\check R$ and
$\mathsf{R}^{\uparrow}$ and the comparison of $N$ with $\check R$, one has $N\le\mathsf{N}$,
$\mathsf{T}_k\le\mathsf{T}+1$ and $\check{R}_{h+1}\le\mathsf{R}$. Comparing the definition of
$\mathsf{b}(s)$ in \eqref{5.2:eq-plain-fine-units-b} with \eqref{5.2:eq-deep-shallow-sums-5}, term by
term, these three inequalities give $\mathsf{b}(s)\le\mathsf{b}^{\mathsf{T}}(s)$ for every $s$.
Every member of the bounds of Lemma~\ref{5.2:lem-deep-sum} and Lemma~\ref{5.2:lem-shallow-sum} is
non-decreasing in $N$, $\check{R}_{h+1}$ and $\mathsf{b}$; replacing them by $\mathsf{N}$,
$\mathsf{R}$ and $\mathsf{b}^{\mathsf{T}}$ therefore bounds
$(1+C_\Sigma C_e)^4\,\mathsf{Q}^{\star,\mathrm{pl}}_{a+9}$ by the first summand of
\eqref{5.2:eq-deep-shallow-sums-a} at $\mathsf{T}=\check{\mathsf{T}}_k$.
Since $\check{R}_k\ge\check{\mathsf{T}}_k^{\,r_{\mathrm{pr}}}$, the number $D$ of
\eqref{5.2:eq-deep-shallow-sums-2} satisfies $D\ge\check{\mathsf{T}}_k^{\,r_{\mathrm{pr}}}-2$.
Combining (c) with clause (d) of Lemma~\ref{5.2:lem-exterior-pieces} for $\varepsilon^{\mathrm{ext},
\sigma}(k)$, and using $\check{\mathsf{T}}_k\ge\mathsf{T}_\gamma$, we obtain
\begin{equation}\label{5.2:eq-deep-shallow-sums-6}
\varepsilon^{\mathrm{int}}(k)+\varepsilon^{\mathrm{ext},\sigma}(k)
\le\mathsf{P}_{\mathrm{int}}(\check{\mathsf{T}}_k)\,
e^{-(\check{\mathsf{T}}_k^{\,r_{\mathrm{pr}}}-2)}\le\frac12 .
\end{equation}
The last inequality is the requirement that
\begin{equation*}
\sup_{\mathsf{T}\ge\mathsf{T}_\gamma}\mathsf{P}_{\mathrm{int}}(\mathsf{T})\,
e^{-(\mathsf{T}^{r_{\mathrm{pr}}}-2)}\le\tfrac12 .
\end{equation*}
The function of $\mathsf{T}$ under the supremum
involves only constants fixed before $\gamma$, is finite, and tends to $0$ as
$\mathsf{T}\to\infty$ because $r_{\mathrm{pr}}\ge2$. The requirement is therefore a one-sided
ceiling on $\gamma$, of degree $-\infty$ in the order of choice.

\emph{Clause (e).} Clause (e) is read off from the definitions of the quantities entering it, as
displayed above.

For the volume corollary we refer to the earlier part of the lemma, where it is established.

\emph{The carrier statement on the core.} Off $\Lambda^{\natural}$ the carrier pair is the pair
\begin{equation*}
\Bigl(\prod_s e^{i\eta\mathbb{H}_s}\,\mathbb{U},\ \mathbf{J}+\sum\mathfrak{s}^{\sharp}\Bigr)
\end{equation*}
of the carrier theorem at a site, where $\eta$, the response layers $\mathbb{H}_s$, $\mathbb{U}$,
$\mathbf{J}$ and the summands $\mathfrak{s}^{\sharp}$ are the objects of that theorem. It is taken
on the direct region
$\mathfrak{Dir}_X:=X\setminus(\Lambda^{\natural})^{+1/2}$, on which the parameter $T$ of that
theorem equals $1$ and on which, by the first clause of the hybrid carrier theorem, the pair uses
at most one quarter of every room. The clauses of the output tube at the
index $\pi(y)\oplus\mathbb{H}(b)$ of the current label, in the notation of the carrier theorem, its
amplitude clause, and its clause for the
layers then give the carrier bound with the weight $w(\Delta_y)$ carried by the definition of the
auxiliary consumer $\mathfrak{E}^C$. This weight equals $1$ off the window cells, so that
$\bar{\mathsf{w}}_v=3$ off them. On the window cells, namely the cells of $X_{\mathfrak{a}}$ above
the window level fixed in the earlier part of the lemma,
one has $\bar{\mathsf{w}}_v\le3\mathsf{L}_0^{2(e+1)}$, where $e$ is the exponent of the window
weights fixed there (not Euler's number); this factor is
absorbed per old cell as displayed in the last clause of the volume corollary: the factor
$27\mathsf{L}_0^6\mathsf{K}^{\flat}$ lies inside $\mathsf{K}_{\mathrm{cell}}^0$ in both branches,
and the frozen windows of each cell present before the arrest lie inside
$\mathsf{P}(\mathfrak{a})$, by the window clause of the fifth modification of the format. This
statement is used in the last clause of the volume corollary; its
counterpart on the exterior hulls is the exterior carrier statement.

\emph{Place in the induction.} The lemma uses the induction hypothesis at level $k$, the
format at $\mathcal{P}=\emptyset$, the estimates of the order
\eqref{5.2:eq-induction-order-a} up to and including the row of clause (c), and the carrier
statement at the auxiliary consumer
$\mathfrak{E}^C$. It returns its member of the pre-filing family at the site $(L)$ of step $k$. It
uses none of the lengths, none of the terms $V$, neither the operation $\mathbf{R}'$ nor
$\alpha^{\sharp}(k)$ of step $k$. Hence its place in the order \eqref{5.2:eq-induction-order-a}
involves no circularity.
\end{proof}

\begin{lemma}[Pieces of units whose label leaves the large-field region]
\label{5.2:lem-exterior-pieces}
\emph{Classes and reading data.} Let $\mathfrak{L}^{c}$ be the class of plain first-part units $v$ of
clause (c) of the classification of plain units at the localisation site of $\mathbf{R}_k$, that
is, of the units with $K_v\setminus Z\neq\emptyset$, of every row of the size table of plain
units. The exterior arrested left-overs (the frozen terms $\mathfrak{F}(\mathfrak{a})$ of an arrest
with $K_v\setminus Z\neq\emptyset$) and the chain terms $\mathbb{C}^{(n)}_k$ of the present chain
with $K_v\setminus Z\neq\emptyset$ are included.

The reading data of $v\in\mathfrak{L}^{c}$ are the triple $(K_v,X_v,\mathfrak{N}(v))$ (its label,
its reading set and its size), where $\mathfrak{N}(v)$ is the size entered for the unit. For a
unit of the size table this is its entry in the size column; for the arrested units it is the
norm of the frozen terms, and for the chain units the stage norm. No depth factor is used, and no
depth-dependent bound for pointed cores is invoked. The exterior size of $v$ is
\begin{equation}\label{5.2:eq-exterior-pieces-1}
\mathfrak{N}^{\mathrm{ex}}(v):=
\begin{cases}
\mathfrak{N}^{\flat}(v), & \text{exposed pointed rows},\\
\mathfrak{N}^{\flat}(v)=2\mathfrak{n}_v, & \text{tails and long singles},\\
\mathfrak{N}(v), & \text{row }4,\ \text{arrested and chain units}.
\end{cases}
\end{equation}
In the first case $\mathfrak{N}^{\flat}(v)$ is the entry of the size table of pointed rows at
$\bar{\mathsf{w}}_v=3w^{\vee}(\Delta_{y_v})$.
The exposed pointed rows are the short singles of rows $2$, $3$ and $\partial$ and the Wilson
pieces.

Let $\mathfrak{L}^{Y_0}$ be the class of the $\mathsf{t}$-pieces $v_p$ (Ba{\l}aban's
$Y_0$-terms) of the plain units of clauses (c) and (d). The reading data of such a piece are
$(K_p,Y_0(p),\mathfrak{N}_A(p))$, with
\begin{equation*}
K_p:=[\operatorname{St}_v\cup\mathsf{y}]_k\cup\operatorname{Hub}(\mathsf{y}),
\end{equation*}
the $\pi_k$-hull of the exempt family and of the walk, with corridors for the never-cut blocks and
not their interiors. For a row-four piece, $\operatorname{St}_v$ is the set of cells meeting
$X_v$. Here $\mathfrak{N}_A(p)$ is the bound of the amplitude lemma on $v_p$ as it is carried by
the three $\theta'$-members of the bracket bound, with the factor
$4C_A'\mathsf{w}_L^{-1}\eta_p^{1/2}$ for row-four units. We put
$\mathfrak{L}^{\mathrm{ex}}:=\mathfrak{L}^{c}\sqcup\mathfrak{L}^{Y_0}$.

\emph{Bundles.} Three families of $\mathsf{t}$-pieces are read in bundles: those of row four (of
(c) and of (d)), those of the long singles and tails of (c), and those of the long singles and
tails of (d). The first two are read in the bundles
$\mathfrak{b}=(v;s;\tilde{\mathsf{y}})$, $\tilde{\mathsf{y}}\subset T^{j_v}_v$, in the unit's own
grain $\pi^{\mathfrak{s}}_{j_v}$, with $T^{j_v}_v$, $\operatorname{Hub}$, $\operatorname{Att}$ and
the grain as in Definition~\ref{5.2:def-boundary-bundles}; the third is read in the grain
$\pi^{\mathfrak{s}}$ of cells itself, the trace map being the identity, as described below. In
each such bundle
the sum over the blocks $\square$ and over the cells inside each member of the grain is taken
exactly first. For the first two families the reading data of a bundle are
\begin{equation*}
\bigl(\operatorname{lab}\mathfrak{b},\ Y_0(\mathfrak{b}),\ \|\mathfrak{b}\|\bigr),\qquad
\|\mathfrak{b}\|=2\mathfrak{N}(v)e^{-(\kappa_1-1)|\tilde{\mathsf{y}}|},
\end{equation*}
where $\mathfrak{N}(v)=\mathfrak{N}^{\flat}(v)=2\mathfrak{n}_v$ for the long units of (c), and with
\begin{equation*}
\operatorname{lab}_1(p):=\operatorname{lab}\mathfrak{b}
:=K_v\cup[\tilde{\mathsf{y}}]_k\cup\operatorname{Hub}(\tilde{\mathsf{y}})
\cup\operatorname{Att}(v,\tilde{\mathsf{y}}),
\qquad
\mathcal{R}_1(p):=Y_0(\mathfrak{b}).
\end{equation*}
For the third family $\operatorname{lab}_1=\operatorname{lab}(K_v,\mathsf{y})$ and
$\mathcal{R}_1=Y_0(p)$, as below.
Consequently $(\operatorname{lab}_t,\mathcal{R}_t,Y_4)$ (the link-$t$ labels and reading sets and
the filing domain of (a) below) are the bundle's for $t=2,3,4$, and each
filed sum is a sum of whole bundles. This is a re-filing of the non-empty continuations. It is
independent of the source coefficients $z_i$ and leaves every exponent untouched. For a unit of
(c) meeting $\Lambda^{\sim}$ the
inclusion $\operatorname{lab}_1\supseteq K_v$ holds by definition, since
$K_v\subset\operatorname{lab}\mathfrak{b}$.

The summands $4C_A'\mathsf{w}_L^{-1}\eta_p^{1/2}\,\mathfrak{N}$ that the three
$\theta'$-members carry for these pieces are idle majorants. The gain of the amplitude lemma is
not used on these rows; in this respect the treatment departs from Ba{\l}aban's construction.
Instead, the summation of boundary terms at finer zones bounds them by
$\varepsilon^{\Sigma4}(k)\le\theta'(g_k)\mathsf{C}^{\Sigma,4}$; see (e) below. The pieces $v_p$
and the classes of the classification are not changed: only the filing of their continuations is
the bundle's; the reading in bundles, the bundle's link-one label and reading set, and the non-use
of the amplitude gain are conventions of filing only.

\emph{Pointed pieces.} Let $p$ be a pointed piece, i.e.\ a piece of rows $1$--$3$, $5$, $\partial$
or of a far core of (d). Every such piece is read through
Lemma~\ref{5.2:lem-carrier-exterior}, at the size $\mathfrak{N}^{\mathrm{ex}}_A(p)$ of
\eqref{5.2:eq-carrier-exterior-d}:
\begin{equation*}
\mathfrak{N}^{\mathrm{ex}}_A(p)
=2\mathsf{w}_L^{-1}(\eta_p\mathsf{f}_v)^{1/2}\cdot\tfrac{81}{64}\,\mathfrak{N}^{\flat}(v)
\cdot e^{-(\kappa_1-1)m(p)}.
\end{equation*}
This is the amplitude lemma on the radius $\rho_A:=\mathsf{w}_L/(\eta_p\mathsf{f}_v)^{1/2}$ of
\eqref{5.2:eq-carrier-exterior-a}. The positional factor is
$\eta_p:=e^{-\delta_P d(\square^{\zeta},S_v)}$, where $d(\cdot,\cdot)$ is the contour distance of
the determining set $\mathfrak{B}_s$, with $\eta_p:=1$ where the $2$-stencil of $S_v$
meets $\square^{\zeta}$. The read set is $S_v=\mathfrak{r}_v$ for a compact $p$, and $S_v$ is the
set of cells meeting $X_v$ for a long one.

Each of these pieces is filed exactly once: the compact pieces are summed over
$(v,s,\square,\mathsf{y})$ by Corollary~\ref{5.2:cor-anchor-sum}, and the long ones of (c) and of
(d), together with row four, are summed in $\square$-bundles by the summation of boundary terms
at finer zones.

\emph{Enlarged labels.} The reading data of the exposed units and of the compact pieces carry the
enlarged label:
\begin{align*}
&(K_v\cup\mathfrak{h}_v,\ X_v,\ \mathfrak{N}^{\flat}(v))
&&\text{for the exposed units of }\mathfrak{L}^{c},\\
&(K_p\cup\mathfrak{h}_v,\ Y_0(p),\ \mathfrak{N}^{\mathrm{ex}}_A(p))
&&\text{for the compact pointed }p.
\end{align*}
In both cases $\operatorname{lab}_1:=K_v\cup\mathfrak{h}_v$, resp.\ $K_p\cup\mathfrak{h}_v$, with
$d_k(\operatorname{lab}_1)\le d_k(K_v)+c_lc_{\mathfrak{h}}$, resp.\
$\le d_k(K_p)+c_lc_{\mathfrak{h}}$, so that the domain
$Y(\operatorname{lab}_1)$ of the auxiliary consumer lies in $Y_4(\mathfrak{p})$.

The long singles and tails of (c) and of (d) are read in $\square$-bundles, like the row-four
units:
\begin{itemize}
\item row four and the long units of (c) are read on
$(\operatorname{lab}\mathfrak{b},Y_0(\mathfrak{b}),\|\mathfrak{b}\|)$;
\item the long units of (d) are read on
$(\operatorname{lab}(K_v,\mathsf{y}),Y_0(p),\|\mathfrak{b}\|)$, with
$\|\mathfrak{b}\|=2\mathfrak{N}^{\flat}(v)e^{-(\kappa_1-1)|\mathsf{y}|}$.
\end{itemize}

\emph{A cube off $Z$.} In every class $\operatorname{lab}_1$ holds a cube off $Z$:
\begin{itemize}
\item for the exposed units because $K_v\setminus Z\neq\emptyset$;
\item for the compact $p$ of (d) because $X_v\cap Z=\emptyset$ and $K_p\supseteq[X_v]_k$;
\item for the compact pointed $p$ of (c) because
$K_p\cup\mathfrak{h}_v\supseteq[X_v]_k=K_v$ via $\mathfrak{h}_v$ (indeed
$[X_v]_k\subset K_p\cap\mathfrak{h}_v$), and $K_v\not\subset Z$;
\item for the bundles because $\operatorname{lab}\mathfrak{b}\supseteq K_v$, resp.\
$\operatorname{lab}(K_v,\mathsf{y})\supseteq K_v$, by definition.
\end{itemize}

\emph{Hypotheses.} Assume:
\begin{enumerate}
\item[(i)] the classification of plain units at the localisation site, with its exterior
clause (c)$'$, and the reading convention for pointed units;
\item[(ii)] the hypotheses on geometry, summation, sources and membership of the lemma on the link
pieces of interior plain units, namely:
\begin{itemize}
\item the hub geometry, clauses (f) and (g), which follow from the nesting of the strips;
\item clauses (a)--(e) of the link-summation lemma;
\item clauses (i) and (ii) of the source expansion lemma;
\item the source-class statement at the localisation site;
\item the sealed membership statement, at Ba{\l}aban's family $\sigma_0$;
\end{itemize}
\item[(iii)] the hypothesis of that lemma in the form that invokes neither its interior carrier
bound nor its column bound, together with the analyticity of the smooth-core extension on every
box;
\item[(iv)] Lemma~\ref{5.2:lem-carrier-exterior}; its input on the current entry is the
receding-block current bound, read with the positional factor $\eta_p$ of
\eqref{5.2:eq-carrier-exterior-a}, under its floor on $C_{\theta'}$ and its ceiling; together
with its clause on long singles and tails and with Corollary~\ref{5.2:cor-anchor-sum}, which
imposes no condition of its own beyond the floor on $\kappa_1$, the floor on the ratio $M/M_1$
and $\kappa\ge180$, read with cells in proper scales, the star $\mathfrak{r}_v$ as read set and
the corridor label $K_p$;
\item[(v)] the summation of boundary terms at finer zones, with
Definition~\ref{5.2:def-boundary-bundles}, taken as an implication from:
\begin{itemize}
\item the containment of the undilated base for the row-four and long units;
\item the partition of unity $\sum_{\square}\hat{\chi}_{\square}=1$ over finitely many blocks;
\item the localisation of the decoupled layers;
\item the anchoring of stored boundary terms;
\item the description of raised and un-raised cells;
\item the hub geometry (f);
\end{itemize}
\item[(vi)] the off-site containment rows at clause offset $b=0$, for the compact pointed pieces.
The long pieces are read at $\mathsf{t}=0$ only, so the rows at $b>0$ are not used;
\item[(vii)] the floor on $C_{\theta'}$, the membership ceiling and the arrest ceiling at half
exponent, the floor $\kappa_1-1\ge c_l(2\kappa+4)+\lambda_{\mathrm{ent}}$, where
$\lambda_{\mathrm{ent}}=\lambda_{\mathrm{ent}}(d,L)$ is a number with
$L^de^{-\frac12(\lambda_{\mathrm{ent}}-\log e\hat b)}\le\frac12$, $\hat b$ being the bound on the
number of cells adjacent to a cell in the site geometry,
and $\kappa\ge180$ (part of the floor $\kappa\ge\kappa_d$).
\end{enumerate}
Then, for every $K$ and every $k\le K$, the statements (a)--(e) below hold.

\begin{enumerate}
\item[(a)] \emph{Pieces.} Let $v\in\mathfrak{L}^{c}$, and write $f_v(\mathfrak{b}_3)$ for $f_v$
evaluated at the third chain background $\mathfrak{b}_3$ (unrelated to the bundles $\mathfrak{b}$
above). Then
$f_v(\mathfrak{b}_3)\circ\operatorname{chain}(\vec{\mathbb{1}})=\mathrm{val}_v+\sum\mathfrak{p}$,
the sum running over the tuples
\begin{equation}\label{5.2:eq-exterior-pieces-2}
\begin{split}
\mathfrak{p}&=(v;\mathsf{y}_2,\mathsf{y}_3,\mathsf{y}_4)\\
&:=\prod_{t=2}^{4}\prod_{\Delta\in\mathsf{y}_t}\int_0^1 d\mathsf{s}^{(t)}(\Delta)\,
\partial_{\mathsf{s}^{(t)}(\Delta)}\,
f_v(\mathfrak{b}_3)\circ\operatorname{chain}(\vec{\mathsf{s}})
\Big|_{\mathsf{s}^{(t)}=0\text{ off }\mathsf{y}_t},
\end{split}
\end{equation}
with $(\mathsf{y}_t)$ not all empty and
$\mathsf{y}_t\in\mathfrak{Y}((\operatorname{lab}_{t-1},\mathcal{R}_{t-1}))\cup\{\emptyset\}$. Here
$\operatorname{lab}_1:=K_v\cup\mathfrak{h}_v$ for an exposed unit,
$\operatorname{lab}_1:=K_v$ for the others, and
$\mathcal{R}_1:=\hat{X}_v\cup X_v$. For $t=2,3,4$, $\operatorname{lab}_t$ and $\mathcal{R}_t$ are
the link-$t$ label and reading set, built from $(\operatorname{lab}_{t-1},\mathcal{R}_{t-1})$ and
$\mathsf{y}_t$ as in the lemma on the link pieces of interior plain units, and
$\mathfrak{Y}((\operatorname{lab}_{t-1},\mathcal{R}_{t-1}))$ is the class of source-anchored
families seen by that pair.

Let $p\in\mathfrak{L}^{Y_0}$. Then
$v_p\circ\operatorname{chain}(\vec{\mathbb{1}})=v_p^{\emptyset}+\sum\mathfrak{p}$, the sum running
over $\mathfrak{p}=(p;\mathsf{y}_2,\mathsf{y}_3,\mathsf{y}_4)$ defined likewise. The link-one
reading data are:
\begin{itemize}
\item $\operatorname{lab}_1:=K_p\cup\mathfrak{h}_v$, $\mathcal{R}_1:=Y_0(p)$ for a compact pointed
$p$;
\item the bundle's, for the long singles, tails and row-four units (which are summed in
$\square$-bundles by the summation of boundary terms at finer zones): for row four and the long
units of (c), $\operatorname{lab}_1:=\operatorname{lab}\mathfrak{b}$ and
$\mathcal{R}_1:=Y_0(\mathfrak{b})$; for the long units of (d),
$\operatorname{lab}_1:=\operatorname{lab}(K_v,\mathsf{y})=K_p$ and $\mathcal{R}_1:=Y_0(p)$.
\end{itemize}
Here $v_p^{\emptyset}$ is the all-empty tuple, which is the piece $v_p$ carried by the three
$\theta'$-members.

The enlarged label grows by at most $c_{\mathfrak{h}}=7^d$ cells. The families
$\mathfrak{Y}((\operatorname{lab}_{t-1},\mathcal{R}_{t-1}))$ see the larger label, so there are
more families, but each is still source-anchored and is summed by clauses (c) and (e) of the
link-summation lemma at the same rates. The tree lengths $d_k$ and $d_{k,Z}$ pay the flat factor
$e^{(a+9)c_lc_{\mathfrak{h}}}$, which is absorbed in $\mathsf{K}^{\flat,\mathrm{ex}}$ and
$\mathsf{C}_1^{\mathrm{ex}}$.

The families are Ba{\l}aban's own, the cells disjoint with $\Lambda$; in the words of
\cite{B-CMP122b}, ``we apply the localization operation to the function $H'_k$ with the set
$\sigma_0$ of cubes $\Delta$ disjoint with $\Lambda$''. The families are read at
$\mathcal{P}_v=\emptyset$, as clause (c)$'$ of the classification prescribes. Put
\begin{equation*}
Y_4(\mathfrak{p}):=\operatorname{lab}_4\cup\{\text{components of }Z\text{ meeting }
\operatorname{lab}_4\}\in\mathbf{D}_k^Z .
\end{equation*}
Then $d_{k,Z}(Y_4(\mathfrak{p}))\le d_k(\operatorname{lab}_4)$. Each $\mathfrak{p}$ is linear in
$f_v$ (resp.\ $v_p$) and analytic where $f_v$ is.

\item[(b)] \emph{Per piece.} Write $(v;\mathsf{y})$ for $(v;\mathsf{y}_2,\mathsf{y}_3,\mathsf{y}_4)$.
Then
\begin{equation}\label{5.2:eq-exterior-pieces-3}
|(v;\mathsf{y})|\le\mathfrak{N}^{\mathrm{ex}}(v)\,e^{-(\kappa_1-1)\sum_t|\mathsf{y}_t|},
\qquad
|(p;\mathsf{y})|\le\mathfrak{N}^{\mathrm{ex}}_A(p)\,e^{-(\kappa_1-1)\sum_t|\mathsf{y}_t|}.
\end{equation}
The first bound needs no auxiliary disc. For row four it holds on the set
$\mathfrak{G}_{\varepsilon'}(X_v)$, the domain of $X_v$ at the parameter $\varepsilon'$ over which
the size of a row-four unit is taken, the size being the orbit-datum supremum
$\mathfrak{N}^{\mathbb{C}}$; here
$\mathfrak{N}^{\mathrm{ex}}=\mathfrak{N}^{\flat}$ on every pointed row and
$\mathfrak{N}^{\mathrm{ex}}=\mathfrak{N}$ on row four and on the arrested and chain units.
For a bundle $\mathfrak{b}$ of row four or of the long units of (c), let
$(\mathfrak{b};\mathsf{y}_2,\mathsf{y}_3,\mathsf{y}_4)$ be the sum of the pieces
$(p;\mathsf{y}_2,\mathsf{y}_3,\mathsf{y}_4)$ over the pieces $p$ of the bundle; then
\begin{equation*}
|(\mathfrak{b};\mathsf{y}_2,\mathsf{y}_3,\mathsf{y}_4)|
\le\|\mathfrak{b}\|\,e^{-(\kappa_1-1)\sum_t|\mathsf{y}_t|},
\end{equation*}
with no disc, by the Cauchy estimates in the parameters $\tilde{s}(G)$ attached to the members $G$
of the grain, on the diagonal of the sealed polydisc, and in the link variables $\mathsf{s}$ of
links $2$--$4$ on the product $\mathfrak{P}$. For a pointed $p$ the bound holds on one disc, the
$\mathsf{t}$-disc of the amplitude lemma at radius
$\rho_A=\mathsf{w}_L/(\eta_p\mathsf{f}_v)^{1/2}$, the same for every tuple of $p$; for a
row-four piece the radius of the amplitude lemma is $\rho_A=\mathsf{w}_L/\eta_p^{1/2}$.

\item[(c)] \emph{The $\mathsf{t}$-free sum.} Put
\begin{equation*}
\mathfrak{v}^{\mathrm{ext},\sigma}(Y_4):=\sum_{\mathfrak{p}:\,Y_4(\mathfrak{p})=Y_4}\mathfrak{p},
\end{equation*}
the sum running over the pieces $\mathfrak{p}$ of the units of $\mathfrak{L}^{c}$. Then
\begin{equation}\label{5.2:eq-exterior-pieces-a}
\begin{split}
\|\mathfrak{v}^{\mathrm{ext},\sigma}\|_a
&\le\|\mathfrak{v}^{\mathrm{ext},\sigma}\|_{a,1}
:=\sup_{\Delta\not\subset Z}\sum_{Y_4\ni\Delta}e^{a\,d_{k,Z}(Y_4)}
|\mathfrak{v}^{\mathrm{ext},\sigma}(Y_4)|\\
&\le\varepsilon^{\mathrm{ext},\sigma}(k)
:=e^{-(\check{R}_k-2)}(1+C_\Sigma C_e)^3\,eS_d\,
\bigl[\mathsf{K}^{0\prime}+\mathsf{K}^{\flat,\mathrm{ex}}+N(g_k)C_0^{\sharp}\bigr],\\
\mathsf{K}^{\flat,\mathrm{ex}}&:=e^{(a+9)c_lc_{\mathfrak{h}}}\cdot 27S_d'\,\mathsf{K}^{\flat}.
\end{split}
\end{equation}
Here $C_\Sigma$, $C_e$ are the constants of clauses (c) and (e) of the link-summation lemma, each
link's cell sum costing a factor $1+C_\Sigma C_e$; $\mathsf{K}^{0\prime}$ is the supply constant
of the collar filing, recomputed at weight
$a+8$ with the sizes entered for the units; $\mathsf{K}^{\flat,\mathrm{ex}}$ is the constant of
the exposed pointed rows, given by the column bound of Lemma~\ref{5.2:lem-exposed-column} at
$\bar{\mathsf{w}}=3$ (the units of $\mathfrak{L}^{c}$ lie at the boundary of the treated
component, where the window weights equal one), multiplied by the flat factor of the enlarged
label; and
\begin{equation*}
S_d:=\sup_{Q\in\pi_k}\sum_{\Delta\in\pi_k}e^{-\operatorname{dist}_k(Q,\Delta)},\qquad
S_d':=\sum_{i\ge0}(2i+3)^de^{-(i-1)_+},\qquad
c_{\mathfrak{h}}:=7^d
\end{equation*}
are numbers depending on $d$ only. The values $\mathrm{val}_v$ of the exposed units enter the
value constant $\mathsf{K}_{\mathrm{val}}$, the lower collar width $m_{\mathrm{lo}}$ and the two
ceilings on the collar values through the bound $|\mathrm{val}_v-f_v(1)|\le\mathfrak{N}^{\flat}(v)$
of Lemma~\ref{5.2:lem-carrier-exterior}, where $f_v(1)$ is the value entered in the size table at
the trivial pair; they are summed per cell by Lemma~\ref{5.2:lem-exposed-column} at
$\bar{\mathsf{w}}=3$. Therefore
$\mathsf{K}^0\mapsto\mathsf{K}^0+\mathsf{K}^{\flat,\mathrm{ex}}$ in the supply of the values, a
number chosen after $M$.

\item[(d)] \emph{Ceiling.} The quantity $\mathsf{P}_{\mathrm{int}}$ under the supremum of the
interior ceiling \eqref{5.2:eq-deep-shallow-sums-a} receives the additional summand
\begin{equation*}
(1+C_\Sigma C_e)^3\,eS_d\bigl[\mathsf{K}^{0\prime}+\mathsf{K}^{\flat,\mathrm{ex}}
+C_N(\mathsf{T}+1)^{r_{\mathrm{pr}}}C_0^{\sharp}\bigr],
\end{equation*}
and under that ceiling
$\varepsilon^{\mathrm{int}}(k)+\varepsilon^{\mathrm{ext},\sigma}(k)\le\frac12$ at every $k\le K$.
In the interior ceiling this summand, a polynomial-type function of $\mathsf{T}$, is multiplied by
$e^{-(\check{R}_k-2)}$ with $\check{R}_k\ge\check{\mathsf{T}}_k^{r_{\mathrm{pr}}}$, so that its
contribution has degree $-\infty$.

\item[(e)] \emph{The continuations of the $\mathsf{t}$-pieces.} For every
$p\in\mathfrak{L}^{Y_0}$, the sum in $\|\cdot\|_a$ of the pieces $(p;\mathsf{y})$ over all tuples
$(\mathsf{y}_2,\mathsf{y}_3,\mathsf{y}_4)$, the empty tuple included, taken in the family form of
clause (c) of the link-summation lemma at an anchor cube not contained in $Z$, is at most
$(1+C_\Sigma C_e)^3eS_d$ times the summand of $p$ in the three $\theta'$-members, its unit read at
weight $a+7$ in place of $a$. The weight is taken on $\operatorname{lab}_1$, that is, on
$K_p\cup\mathfrak{h}_v$ for a compact pointed $p$, and on the label of the $\square$-bundle,
$\operatorname{lab}\mathfrak{b}$ resp.\ $\operatorname{lab}(K_v,\mathsf{y})$, for the long singles,
tails and row-four units, as in (a). For a pointed $p$ the summand is taken at the size and radius
of $\mathfrak{N}^{\mathrm{ex}}_A(p)$; for a bundled one at $\|\mathfrak{b}\|$.

Hence the three $\theta'$-members, read at weight $a+7$ and at the sizes and radius of the
pointed rows, are multiplied by $(1+C_\Sigma C_e)^3eS_d$, with the following changes of
constants.
\begin{itemize}
\item The numbers $8C_A'C_NC_0^{\sharp}(\mathsf{T}_k+1)^{r_{\mathrm{pr}}}$ and $4C_A'$ are
unchanged. The row-four content of the second and third members (the exposed chain terms
$\mathbb{C}^{(n)}_k$ of the present chain and the exterior arrested left-overs, which are
$\mathsf{t}$-pieces of row-four units) belongs to the summation of boundary terms at finer zones
and is contained in $\mathsf{C}^{\Sigma,4}$ below. The summands
$8C_A'C_NC_0^{\sharp}(\mathsf{T}+1)^{r_{\mathrm{pr}}}$ and $8C_A'\mathsf{K}^{\mathrm{arr}\prime}$
remain in the bracket ceiling as idle majorants; the summand $4C_A'$ remains in use.
\item The constant $C_\sharp O(1)$ of the first member is replaced by
$\mathsf{C}_1^{\mathrm{ex}}$ of \eqref{5.2:eq-exterior-pieces-b}. There
$\mathsf{C}^{\Sigma}_r$ is the constant of \eqref{5.2:eq-anchor-sum-a},
\begin{equation*}
\mathsf{C}^{\Sigma}_r=8S_dN_\partial e^{1+(r+1)c_{\mathfrak{h}}}C_{\square}^0c_S(c_S+3)
\tfrac32\,\mathsf{C}_{\mathrm{fw}}\mathsf{K}^{\flat},
\end{equation*}
for the compact pointed pieces of every zone $n'\le k$ (Corollary~\ref{5.2:cor-anchor-sum}), and
\begin{equation*}
\mathsf{C}^{\Sigma,4}:=\mathsf{c}_{\theta'}\bigl[(\mathsf{C}^{4,d}+\mathsf{C}^{4,c})
\mathsf{b}^{\natural}+\mathsf{C}^{\mathrm{lg},c}+\mathsf{C}^{\mathrm{lg},d}\bigr],
\qquad
\mathsf{c}_{\theta'}:=e^2C_{\theta'}^{-1}\sup_{\mathsf{T}\ge1}(\mathsf{T}+1)^{q-p_1}
e^{-\mathsf{T}^{r_{\mathrm{pr}}}},
\end{equation*}
for the row-four units of (c) and (d) and the long singles and tails of (c) and of (d). The long
singles and tails of (c) are summed through the per-cell bound for the long rows of class (c),
those of (d) through the exact resummation of the blocks together with the long-row theorem for
class (d), at
$\mathsf{C}^{\mathrm{lg},d}:=16e\cdot dL^{d-1}S_d\mathsf{K}^{\mathrm{LG}}$. All are read in
bundles at weight $a+7$ on $\operatorname{lab}(\mathfrak{b})$ (on
$\operatorname{lab}(K_v,\mathsf{y})$ on the long rows of (d)), by the summation of boundary terms
at finer zones with its corollary
$\varepsilon^{\Sigma4}(k)\le\theta'(g_k)\mathsf{C}^{\Sigma,4}$ and its corollary for links
$2$--$4$; $\mathsf{c}_{\theta'}$ is a number chosen after $C_{\theta'}$. The factor
$\frac32\ge(\frac98)^3$ majorises the factors of the dilated sizes, namely the $(\frac98)^2$ of
$\epsilon^2$ and the cube of the polynomial in the core's scale index $s(\cdot)$, read at
$s(\frac98\bar{\mathsf{w}})$; and
$\mathsf{C}_{\mathrm{fw}}$ is the constant of \eqref{5.2:eq-carrier-exterior-f}, a
$\mathsf{T}$-free number chosen after $M$ and before $\gamma$.
\item $\mathsf{K}^{\mathrm{arr}}$ is replaced by $\mathsf{K}^{\mathrm{arr}\prime}$, the same collar
series of the supply at weight $a+7$, free of $g$, of $k$ and of ages alike.
\end{itemize}
With these constants,
\begin{equation}\label{5.2:eq-exterior-pieces-b}
\begin{split}
\mathsf{C}_1^{\mathrm{ex}}&:=\mathsf{C}^{\Sigma,4}+\mathsf{C}^{\Sigma}_{a+7},\\
\varepsilon^{\mathrm{ext},\mathsf{t}}(k)&:=(1+C_\Sigma C_e)^3eS_d\times
\bigl[\text{the three $\theta'$-members at weight $a+7$}\bigr]\\
&\qquad-\bigl[\text{the three $\theta'$-members}\bigr],\\
\frac12&\ge\sup_{\mathsf{T}\ge\mathsf{T}_\gamma}(1+C_\Sigma C_e)^3eS_d
\bigl[\mathsf{C}_1^{\mathrm{ex}}+8C_A'C_NC_0^{\sharp}(\mathsf{T}+1)^{r_{\mathrm{pr}}}\\
&\qquad+8C_A'\mathsf{K}^{\mathrm{arr}\prime}+4C_A'\bigr]\,C_{\theta'}\mathsf{T}^{p_1-q},
\end{split}
\end{equation}
where in the second line the first member is taken at the constant
$\mathsf{C}_1^{\mathrm{ex}}$, and $\varepsilon^{\mathrm{ext},\mathsf{t}}(k)$ is read as this
difference. The third line is the bracket ceiling with its new constants. Its degree is at most
$r_{\mathrm{pr}}-(q-p_1)\le-6r_{\mathrm{pr}}-1<0$ under the condition
$q-p_1\ge7r_{\mathrm{pr}}+1$; it is a ceiling on $\gamma$ of supremum form, as before. Of the
amplitude lemma, its half-amplitude form and the off-site containment rows, what moves is the
radius of the pointed rows at the localisation site, from $\mathsf{w}/\eta_p$ to
$\mathsf{w}_L/(\eta_p\mathsf{f}_v)^{1/2}$, and nothing else. No constant and no ceiling reads $K$,
$k$, an age, $N_0$, $\check{R}$, a window, an arrest count, the cardinality of any of the sets
$\mathcal{P}_v$ or
the volume; $N$ enters only through $\mathsf{N}$, and the only floor read is the floor on
$\kappa_1$.
\end{enumerate}
\end{lemma}

Clauses (a), (b) and (c) of the lemma on pieces of units whose label leaves the large-field region (Lemma~\ref{5.2:lem-exterior-pieces}) are proved below. The proof also establishes the reach property of the link families, which clause (c) uses. Throughout, the hypotheses of that lemma are in force.

The classes $\mathfrak{L}^{c}$ and $\mathfrak{L}^{Y_0}$, their readers, the labels $\operatorname{lab}_t$ and the reading sets $\mathcal{R}_t$ are those of that lemma. So are the domains $Y_4(\mathfrak{p})$, the family $\mathfrak{v}^{\mathrm{ext},\sigma}$ and the constants of \eqref{5.2:eq-exterior-pieces-a}, namely $\varepsilon^{\mathrm{ext},\sigma}(k)$, $\mathsf{K}^{\flat,\mathrm{ex}}$, $S_d$, $S_d'$ and $c_{\mathfrak{h}}=7^d$.

The three clauses are derived from the following statements. Each is used through its statement, except where its proof is said to be used.
\begin{itemize}
\item the link expansion lemma, parts (i) and (ii);
\item the holomorphy supply at the localisation site, clauses (c1) and (c6), and, through its proof,
clauses (c2) and (c3);
\item the two consumer holomorphy statements, applied at the background $\mathfrak{b}_3$;
\item the carrier bound for exterior pointed units, Lemma~\ref{5.2:lem-carrier-exterior}, clauses
(a), (b) and (c), with the Cauchy radius $\rho_A$ of \eqref{5.2:eq-carrier-exterior-a};
\item the amplitude lemma, through its Cauchy estimate in the amplitude $\mathsf{t}$;
\item the exact block-summation lemma of the type-four sum, part (b);
\item the column bound per cell for exposed units, Lemma~\ref{5.2:lem-exposed-column}, in the form
\eqref{5.2:eq-exposed-column-a};
\item the family-summation lemma, parts (a)--(e);
\item clause (f) of the collar geometry, which follows from the nesting condition on the
large-field regions;
\item clause (c) of the rate proposition, which fixes $a=(31/20)\kappa$ and $a_k=(8/5)\kappa$;
\item the supply bound of the value filing, clauses (i$'$) and (ii), together with clause (ii) of the
boundary-type lemma;
\item the connectedness property of tree lengths: if $X$ is a connected union of level-$k$ cubes and
$\Delta,\Delta'$ are cubes of $X$, then $d_k(X)\ge \operatorname{dist}_k(\Delta,\Delta')-1$.
\end{itemize}

For a cell $\Delta$ we write $\operatorname{pa}(\Delta)$ for the cube of $\pi_k$ containing it.

By the units of types one to five and $\partial$ we mean the units carrying the corresponding types in the table of sizes of the plain units. The exposed pointed units are the short singles of types two, three and $\partial$, and the Wilson pieces.

\begin{proof}[Proof of clauses (a)--(c) of Lemma~\ref{5.2:lem-exterior-pieces}]
\label{5.2:prf-exterior-source-free}

\emph{Clause (a).}
Fix $v\in\mathfrak{L}^{c}$.

For every link $t\in\{2,3,4\}$ and every cell $\Delta$ carrying a decoupling parameter $\mathsf{s}^{(t)}(\Delta)$, apply the fundamental theorem of calculus in that parameter. It writes the value at $\mathsf{s}^{(t)}(\Delta)=1$ as the value at $\mathsf{s}^{(t)}(\Delta)=0$ plus $\int_0^1 d\mathsf{s}^{(t)}(\Delta)\,\partial_{\mathsf{s}^{(t)}(\Delta)}$. We do this link by link, as in part (i) of the link expansion lemma.

Expanding the resulting products, $f_v(\mathfrak{b}_3)\circ\operatorname{chain}(\vec{1})$ becomes the sum, over all tuples $(\mathsf{y}_2,\mathsf{y}_3,\mathsf{y}_4)$, of the pieces $\mathfrak{p}=(v;\mathsf{y}_2,\mathsf{y}_3,\mathsf{y}_4)$ defined in clause (a). In other words, the sum over all tuples is the function at $\vec{1}$.

By part (i) of the link expansion lemma, the family $\mathsf{y}_t$ at link $t$ ranges over
\[
\mathfrak{Y}((\operatorname{lab}_{t-1},\mathcal{R}_{t-1}))\cup\{\emptyset\}.
\]
The all-empty tuple contributes the function with all decoupling parameters of links $2$ to $4$ set to zero, that is, $\operatorname{val}_v$. The same expansion applied to $v_p$, $p\in\mathfrak{L}^{Y_0}$, gives $v_p\circ\operatorname{chain}(\vec{1})=v_p^{\emptyset}+\sum\mathfrak{p}$, with the link-one readers fixed in clause (a).

The families are Ba{\l}aban's own: the decoupling parameters are attached only to the cubes disjoint with $\Lambda$, as recorded in clause (a) of Lemma~\ref{5.2:lem-exterior-pieces}. Consequently no cell meeting $\Lambda$ is cut.

Consider next the holomorphy of the integrands in the decoupling parameters. Take any consumer $\mathfrak{E}\in\mathcal{C}$ whose domain $X$ satisfies $X\cap Z\neq\emptyset$, in any position. For such a consumer, holomorphy is the content of the two consumer holomorphy statements. They are applied by statement at $\mathfrak{b}=\mathfrak{b}_3$, on the unit's own analyticity domain $\mathfrak{D}_v$. That they apply there is the content of clauses (c1) and (c6) of the holomorphy supply at the localisation site.

For the exposed pointed units, two facts are supplied by clauses (a) and (b) of the carrier bound for exterior pointed units (Lemma~\ref{5.2:lem-carrier-exterior}): the holomorphy, and the $\mathsf{t}$-powered size on the cut carrier.

The domain $Y_4(\mathfrak{p})$ adds to $\operatorname{lab}_4$ only whole components of $Z$, which $d_{k,Z}$ does not count; hence $d_{k,Z}(Y_4(\mathfrak{p}))\le d_k(\operatorname{lab}_4)$. The piece $\mathfrak{p}$ is an iterated integral of derivatives of $f_v$ (resp. $v_p$) in the decoupling parameters; hence it is linear in $f_v$ (resp. $v_p$) and analytic where $f_v$ is. This proves (a).

\emph{Clause (b).}
We apply the estimates below to the unit minus its value at $\operatorname{chain}(\vec{0})$; subtracting this constant does not change any derivative in the decoupling parameters.

Let $\Delta\in\mathsf{y}_t$ be a cut cell and $\mathsf{s}(\Delta)\in[0,1]$. The circle
\[
|\sigma-\mathsf{s}(\Delta)|=e^{\kappa_1}-1
\]
lies inside the sealed polydisc of Lemma~\ref{5.2:lem-exterior-pieces}, on which $|\sigma(\Delta)|\le e^{\kappa_1}$. Its radius satisfies $e^{\kappa_1}-1\ge e^{\kappa_1-1}$.

Part (ii) of the link expansion lemma is Cauchy's estimate on these circles, in every cut parameter of links $2$ to $4$. It is applied to the unit minus its value at $\operatorname{chain}(\vec{0})$. The supremum used is the one supplied by clause (b) of Lemma~\ref{5.2:lem-carrier-exterior} for the exposed pointed units, and the assigned norm for the others. It gives
\begin{equation*}
|(v;\mathsf{y}_2,\mathsf{y}_3,\mathsf{y}_4)|
\le\mathfrak{N}^{\mathrm{ex}}(v)\prod_{t=2}^{4}\bigl(e^{\kappa_1}-1\bigr)^{-|\mathsf{y}_t|}
\le\mathfrak{N}^{\mathrm{ex}}(v)\,e^{-(\kappa_1-1)\sum_t|\mathsf{y}_t|}.
\end{equation*}

For $p\in\mathfrak{L}^{Y_0}$, the estimate in the decoupling parameters is taken jointly with the Cauchy estimate of the amplitude lemma in $\mathsf{t}$, on the disc of radius $\rho_A$. The piece is jointly holomorphic in $(\zeta,\sigma,\mathsf{t})$ by the consumer holomorphy statements, and a $\sigma$-localisation piece carries no amplitude of its own. The two estimates are therefore read on the product of the two discs.

For a pointed $p$, the size and the radius on this joint polydisc are those of clause (c) of Lemma~\ref{5.2:lem-carrier-exterior}:
\begin{itemize}
\item the radius $\rho_A=\mathsf{w}_L/(\eta_p\mathsf{f}_v)^{1/2}$ of \eqref{5.2:eq-carrier-exterior-a};
\item the local entries at every block;
\item the $\mathbf{J}$-entry, by the recursive current bound, at every block;
\item the size $(81/64)\mathfrak{N}^{\flat}(v)$ and the Cauchy estimate;
\item the clause treating the long units.
\end{itemize}
This gives the bound $|(p;\mathsf{y})|\le\mathfrak{N}^{\mathrm{ex}}_A(p)e^{-(\kappa_1-1)\sum_t|\mathsf{y}_t|}$.

Let $\mathfrak{b}$ be a bundle of type-four units or of the long units of the first class of the decomposition of Lemma~\ref{5.2:lem-exterior-pieces}. The bound
\begin{equation*}
|(\mathfrak{b};\mathsf{y}_2,\mathsf{y}_3,\mathsf{y}_4)|
\le\|\mathfrak{b}\|\,e^{-(\kappa_1-1)\sum_t|\mathsf{y}_t|},
\end{equation*}
with no auxiliary disc, is part (b) of the exact block-summation lemma of the type-four sum: it is Cauchy's estimate in the grain parameter on the diagonal of the sealed polydisc, and in the decoupling parameters of links $2$ to $4$ on the product polydisc $\mathfrak{P}$.

It remains to treat the plain units of the second class of the decomposition of Lemma~\ref{5.2:lem-exterior-pieces}. For them, the first two brackets are taken on the pairs of backgrounds $(\mathfrak{b}_0,\mathfrak{b}_1)$ and $(\mathfrak{b}_1,\mathfrak{b}_2)$. These lie below the range $\mathfrak{b}\in\{\mathfrak{b}_2,\mathfrak{b}_3,\mathfrak{b}_4\}$ of the sealed statement.

There the holomorphy used above is supplied by the proof of clauses (c2) and (c3) of the holomorphy supply at the localisation site, together with one further step of that proof. In that proof, $\sigma_0$ is left free and the argument is an induction on $s=1,\dots,4$. Its consumer class includes the $Y_0$-terms, and for a unit of the second part the induction runs over $s=4,\dots,1$. This proves (b).

\emph{The reach of the link families.}
Let $t\in\{2,3,4\}$ and let $\Delta\in\mathsf{y}_t$ be a cell that sees $\mathcal{R}_{t-1}$. Here $\rho(\Delta)$ is the reach radius of the cell $\Delta$ and $c_\rho$ the reach constant of the family-summation lemma, and $h_\sharp$ is the bound on the side of a hub in level-$k$ cubes. The reach clause of the family-summation lemma requires
\begin{equation}\label{5.2:eq-exterior-source-free-1}
\operatorname{dist}_k\bigl(\operatorname{pa}(\Delta),\operatorname{lab}_{t-1}\bigr)\le c_\rho\,\rho(\Delta),
\end{equation}
which is the condition under which that lemma operates. We verify its cases, as that clause lists them, for the present readers. There are three cases.

\begin{enumerate}
\item[(1)] The cell $\Delta$ meets $X_v$, or is a cell of $\operatorname{St}_v$, or belongs to an earlier family $\mathsf{y}_s$ with $s<t$. Then $\operatorname{pa}(\Delta)$ either lies in $\operatorname{lab}_{t-1}\supseteq\operatorname{lab}_1$ or shares a wall with a cube of it, so the distance in \eqref{5.2:eq-exterior-source-free-1} is at most $1$. Here one uses $[\operatorname{St}_v]_k\subset\mathfrak{h}_v\subset\operatorname{lab}_1$ for a compact $p$, and $[X_v]_k\subset K_p$ for a long one.
\item[(2)] The cell $\Delta$ is a source cell at a hub touched or met before. These are the only members of the families lying in $\Lambda^{\sim}\setminus\Lambda$, the hubs themselves never being cut. Since that hub was touched or met at an earlier link, the reach clause gives a cube of $\operatorname{lab}_{t-1}$ sharing a wall with it or lying in it, so the distance is at most $h_\sharp\le c_\rho$. At $t=2$ the hub is one met by $\operatorname{St}_v\cup\mathsf{y}$, with $\mathsf{y}$ the link-one family of $p$. The set $\operatorname{Hub}(\mathsf{y})\subset K_p$ joins the hubs through which the family $\mathsf{y}$ passes, and the constant $c_l$ per added cell in part (a) of the family-summation lemma pays for the corridor attached to each hub there.
\item[(3)] The cell $\Delta$ touches one of three rims: the rim of $\operatorname{lab}_{t-1}$ (distance at most $1$), the rim of a live strip (one more), or the rim of a component of $Z$ that $\operatorname{lab}_{t-1}$ meets. By clause (f) of the collar geometry, the nearest hub is at distance at least $9\check{R}_k$ from any rim of $Z$. Hence the distance is bounded as follows:
\begin{equation*}
\operatorname{dist}_k\bigl(\operatorname{pa}(\Delta),\operatorname{lab}_{t-1}\bigr)
\le 100\check{R}_k+1\le c_\rho\,(9\check{R}_k-2)\le c_\rho\,\rho(\Delta).
\end{equation*}
\end{enumerate}

Consequently the family-summation lemma, parts (a)--(e), applies at links $2$ to $4$. In it, $K_p=[\operatorname{St}_v\cup\mathsf{y}]_k\cup\operatorname{Hub}(\mathsf{y})$, and the domain $Y(\operatorname{lab}_1)$ of the auxiliary consumer is unchanged.

\emph{Clause (c): geometry.}
Only source-anchored families contribute: if a non-empty family $\mathsf{y}_t$ contains no source cell, the piece vanishes, since the derivatives with respect to the decoupling parameters, restricted to components containing no source cell, make the term vanish \cite{B-CMP116}. So we may assume that every non-empty family $\mathsf{y}_t$ contains a source cell.

By clause (f) of the collar geometry, a source cell lies in the collar of $\Lambda^{\sim}$, at $d_k$-distance at least $9\check{R}_k-1$ from $Z^{c}$. Moreover, $\operatorname{lab}_4\supseteq\operatorname{lab}_1$, and by the classes paragraph of Lemma~\ref{5.2:lem-exterior-pieces}, $\operatorname{lab}_1$ holds a cube off $Z$. Since $\operatorname{lab}_4$ is connected, the connectedness property of tree lengths (which costs one) gives
\begin{equation}\label{5.2:eq-exterior-source-free-2}
d_k(\operatorname{lab}_4)\ge 9\check{R}_k-2\ge\check{R}_k-2=:D.
\end{equation}
Only the bound by $D$ is used below. This is the same lower bound as in clause (c) of the interior plain-unit lemma and in clause (d) of the label-length proposition, and it holds for shallow and deep units alike.

\emph{Clause (c): rates.}
Put $d:=d_k(\operatorname{lab}_4)$. By clause (a), $d_{k,Z}(Y_4)\le d$, and by \eqref{5.2:eq-exterior-source-free-2}, $d\ge D$. Hence
\[
e^{a\,d_{k,Z}(Y_4)}\le e^{a d}\le e^{-D}e^{(a+1)d}.
\]

Apply parts (c) and (e) of the family-summation lemma three times, at links $4$, $3$ and $2$. At each link, the sum over the cells costs a factor $(1+C_\Sigma C_e)$ and two units of weight, so the weight rises from $a+1$ to $a+7$. This gives
\begin{equation}\label{5.2:eq-exterior-source-free-3}
\begin{split}
\|\mathfrak{v}^{\mathrm{ext},\sigma}\|_{a,1}
&=\sup_{\Delta\not\subset Z}\sum_{Y_4\ni\Delta}e^{a\,d_{k,Z}(Y_4)}
\bigl|\mathfrak{v}^{\mathrm{ext},\sigma}(Y_4)\bigr|\\
&\le e^{-D}(1+C_\Sigma C_e)^3\,\mathsf{Q}_{a+7}[\mathfrak{L}^{c}],
\end{split}
\end{equation}
where we define
\begin{equation*}
\mathsf{Q}_r[\mathfrak{L}^{c}]:=\sup_{Q\in\pi_k}\ \sum_{v\in\mathfrak{L}^{c}:\,K_v\ni Q}
\mathfrak{N}(v)\,e^{r\,d_k(K_v)}.
\end{equation*}
Since $\mathsf{Q}_{a+7}[\mathfrak{L}^{c}]$ is a supremum over all cubes $Q\in\pi_k$, the same chain bounds the left member $\|\mathfrak{v}^{\mathrm{ext},\sigma}\|_a$ of clause (c) by the same right-hand side.

This application requires $\kappa_1-1\ge c_l(a+7)+\lambda_*$. The floor condition on $\kappa_1$ in the hypotheses implies it, because
\[
a+7\le a_k=(8/5)\kappa\le 2\kappa+4 .
\]
The first of these inequalities reads $(31/20)\kappa+7\le(32/20)\kappa$, which holds for $\kappa\ge140$, and $\kappa\ge180$ is a standing hypothesis.

\emph{Clause (c): transfer to a cube off $Z$.}
Let $Q\in\pi_k$ and let $v\in\mathfrak{L}^{c}$ with $K_v\ni Q$. Choose a cube $\Delta_v\in K_v$ with $\Delta_v\not\subset Z$. Since $K_v$ is connected, the connectedness property of tree lengths gives $d_k(K_v)\ge\operatorname{dist}_k(Q,\Delta_v)-1$. Therefore
\begin{equation*}
e^{(a+7)d_k(K_v)}\le e\cdot e^{-\operatorname{dist}_k(Q,\Delta_v)}\,e^{(a+8)d_k(K_v)}.
\end{equation*}

Sum first over the position $\Delta$ of $\Delta_v$, and use $\sum_{\Delta\in\pi_k}e^{-\operatorname{dist}_k(Q,\Delta)}\le S_d$. This gives
\begin{equation}\label{5.2:eq-exterior-source-free-4}
\mathsf{Q}_{a+7}[\mathfrak{L}^{c}]\le e\,S_d\,\sup_{\Delta\not\subset Z}
\sum_{v\in\mathfrak{L}^{c}:\,K_v\ni\Delta}\mathfrak{N}(v)\,e^{(a+8)d_k(K_v)}.
\end{equation}

The supremum on the right is the display of the supply bound of the value filing at weight $a+8$. This weight is at most $a_k=(8/5)\kappa$ as soon as $\kappa\ge160$. We re-run that display with the sizes of the units, following clause (c) of the rate proposition and the corresponding clause of the supply bound:
\begin{itemize}
\item the sum over the domains $X$ is taken at $a_k$;
\item clauses (i$'$) and (ii) are taken at the weight $e^{(8/15)\kappa d_j}$, with the split $28/60=3/60+25/60$ and the inequality $(25/12)L^s\ge ds$, valid for $L\ge3$;
\item the units of the table of sizes are taken at $a_m=(4/5)\kappa$, with the split $16/60=3/60+13/60$ and the inequality $(13/12)L^s\ge4s$, valid for $L\ge4$;
\item the units of types one to three, five and $\partial$ are summed at weight $(7/4)\kappa$ by the paragraph of clause (ii) of the supply bound, that is, by clause (ii) of the boundary-type lemma at $(7/4)\kappa$, at their assigned sizes.
\end{itemize}
For the units at their assigned sizes, this bounds the supremum by $\mathsf{K}^{0\prime}+N(g_k)C_0^{\sharp}$.

\emph{Clause (c): the exposed pointed units.}
These units carry the size $\mathfrak{N}^{\mathrm{ex}}=\mathfrak{N}^{\flat}$. They are summed instead by the column bound per cell for exposed units (Lemma~\ref{5.2:lem-exposed-column}), on their enlarged label $\operatorname{lab}_1=K_v\cup\mathfrak{h}_v$.

Let such a unit have $\operatorname{lab}_1\ni\Delta$. By the location clause for exposed units, its point $y_v$ lies in a level-$k$ cell $\Delta'$ of zone $k$: either in the collar of $C$, or in $\Omega_k$ beside $\partial C$. There $\bar{\mathsf{w}}(\Delta')=3$.

Put $i:=\operatorname{dist}_k(\Delta',\Delta)$. Since $\operatorname{lab}_1$ is connected, the connectedness property of tree lengths and the estimate $d_k(\operatorname{lab}_1)\le d_k(K_v)+c_lc_{\mathfrak{h}}$ of Lemma~\ref{5.2:lem-carrier-exterior} give
\[
i\le d_k(\operatorname{lab}_1)+1\le d_k(K_v)+c_lc_{\mathfrak{h}}+1 .
\]
Hence $d_k(K_v)\ge(i-1-c_lc_{\mathfrak{h}})_+$. Together with $d_k(\operatorname{lab}_1)\le d_k(K_v)+c_lc_{\mathfrak{h}}$ this yields
\begin{equation*}
e^{(a+8)d_k(\operatorname{lab}_1)}\le e^{(a+8)c_lc_{\mathfrak{h}}}\,e^{(a+9)d_k(K_v)}\,
e^{-(i-1-c_lc_{\mathfrak{h}})_+}.
\end{equation*}

The cells $\Delta'$ at distance $i$ from $\Delta$ number at most $(2i+3)^d$. For each of them, \eqref{5.2:eq-exposed-column-a} at $r=a+9$ gives
\begin{equation*}
\sum_{y_v\in\Delta'}\mathfrak{N}^{\flat}(v)\,e^{(a+9)d_k(K_v)}\le 27\,\mathsf{K}^{\flat}.
\end{equation*}
This is admissible because $a+9\le(8/5)\kappa$, which holds for $\kappa\ge180$.

Write $c:=c_lc_{\mathfrak{h}}\ge7^d$. Summing the previous two displays over the cells $\Delta'$ and over $i$, the exposed pointed units with $\operatorname{lab}_1\ni\Delta$ contribute at most
\begin{align*}
&e^{(a+8)c}\cdot27\,\mathsf{K}^{\flat}\sum_{i\ge0}(2i+3)^d e^{-(i-1-c)_+}\\
&\qquad\le e^{(a+8)c}\,(c+2)(2c+5)^d\,S_d'\cdot27\,\mathsf{K}^{\flat}
\le e^{(a+9)c}\cdot27\,S_d'\,\mathsf{K}^{\flat}=\mathsf{K}^{\flat,\mathrm{ex}}.
\end{align*}

The first inequality is obtained by splitting the sum at $i=c+1$.
\begin{itemize}
\item The terms with $i\le c+1$ number at most $c+2$, and each is at most $(2c+5)^d$.
\item In the tail, put $i=c+1+j$ with $j\ge1$. There $(2(j+c+1)+3)\le(2c+5)(2j+1)$ and $e^{-j}\le e^{-(j-2)_+}$, so that
\[
\sum_{j\ge1}(2(j+c+1)+3)^de^{-j}\le(2c+5)^dS_d'.
\]
\item Finally, $(c+2)+S_d'\le(c+2)S_d'$, since both numbers are at least $2$.
\end{itemize}
The second inequality uses $(c+2)(2c+5)^d\le e^{c}$ for $c\ge7^d$.

The factor $e\,S_d$ of the transfer, from $Q$ to $\Delta$, and the factor $S_d'$ here, from $\Delta$ to $\Delta'$, come from two different sums. Both depend on $d$ only.

\emph{Conclusion.}
Insert the last two bounds into \eqref{5.2:eq-exterior-source-free-4}, and that into \eqref{5.2:eq-exterior-source-free-3} and into the same bound for $\|\mathfrak{v}^{\mathrm{ext},\sigma}\|_a$, with $D=\check{R}_k-2$. This gives
\begin{equation*}
\begin{split}
\max\bigl\{\|\mathfrak{v}^{\mathrm{ext},\sigma}\|_{a},\|\mathfrak{v}^{\mathrm{ext},\sigma}\|_{a,1}\bigr\}
&\le e^{-(\check{R}_k-2)}(1+C_\Sigma C_e)^3\,e\,S_d\,
\bigl[\mathsf{K}^{0\prime}+\mathsf{K}^{\flat,\mathrm{ex}}+N(g_k)C_0^{\sharp}\bigr]\\
&=\varepsilon^{\mathrm{ext},\sigma}(k),
\end{split}
\end{equation*}
which is clause (c).
\end{proof}

\begin{proof}[Proof of Lemma~\ref{5.2:lem-exterior-pieces}, clauses (d) and (e)]
\label{5.2:prf-exterior-sourced}
Clauses (a)--(c) of Lemma~\ref{5.2:lem-exterior-pieces} were proved above. The argument below
proves (d) and (e) from the following statements, each used by statement:
\begin{itemize}
\item Corollary~\ref{5.2:cor-anchor-sum} (the sum of pointed pieces behind one coarse anchor),
with its constant $\mathsf{C}^{\Sigma}_r$ of \eqref{5.2:eq-anchor-sum-a};
\item the theorem on summation of boundary terms at finer zones, its corollary in the
coupling ratio $\theta'(g_k)$, its corollary for links two to four, the lemma
on exact summation over blocks and fine cells, and the summation theorem in the unit's own
grain;
\item clauses (a), (b) and (c) of the summation lemma for source-anchored families;
\item clause (i$'$) of the supply bound for the values, and the column bounds for the pointed
rows, interior and exterior;
\item the amplitude lemma with the size table, and the inductive bound \cite[(1.66)]{B-CMP122b};
\item the carrier bound for exterior pointed units, through the size
\eqref{5.2:eq-carrier-exterior-d};
\item the deep and shallow sums inside the large-field region (proved above), through
$\mathsf{P}_{\mathrm{int}}$ of \eqref{5.2:eq-deep-shallow-sums-a}.
\end{itemize}

\emph{Clause (d).} By clause (c),
\begin{equation*}
\varepsilon^{\mathrm{ext},\sigma}(k)
= e^{-(\check{R}_k-2)}\,(1+C_\Sigma C_e)^3\, e\,S_d\,
\bigl[\mathsf{K}^{0\prime}+\mathsf{K}^{\flat,\mathrm{ex}}+N(g_k)\,C_0^{\sharp}\bigr].
\end{equation*}
The constants $\mathsf{K}^{0\prime}$ and $\mathsf{K}^{\flat,\mathrm{ex}}$ do not depend on
$\mathsf{T}$, on $k$ or on $K$. By
$N\le C_N(\mathsf{T}+1)^{r_{\mathrm{pr}}}$ the bracket is a function of polynomial type in
$\mathsf{T}$. The prefactor $e^{-(\check{R}_k-2)}$ is set against it. Since
$\check{R}_k=R(\min_{i\le k}x_i)$ is a power of $L$ not smaller than
$(\log\min_{i\le k}x_i)^{r_{\mathrm{pr}}}$, we have
$\check{R}_k\ge \check{\mathsf{T}}_k^{\,r_{\mathrm{pr}}}$, and the product has
degree $-\infty$. It therefore enters the one supremum of the ceiling condition on the interior
activity, beside $\mathsf{P}_{\mathrm{int}}$ of \eqref{5.2:eq-deep-shallow-sums-a}, as the
additional summand
\begin{equation*}
(1+C_\Sigma C_e)^3\, e\,S_d\,
\bigl[\mathsf{K}^{0\prime}+\mathsf{K}^{\flat,\mathrm{ex}}
+C_N(\mathsf{T}+1)^{r_{\mathrm{pr}}}C_0^{\sharp}\bigr].
\end{equation*}
The ceiling condition so augmented states $\varepsilon^{\mathrm{int}}(k)
+\varepsilon^{\mathrm{ext},\sigma}(k)\le \tfrac12$ at every $k\le K$, which is (d).

\emph{Clause (e), the bound at fixed $p$.} Fix $p\in\mathfrak{L}^{Y_0}$. The all-empty tuple is
$v_p$ itself, that is, the summand by which $v_p$ enters the three $\theta'$-members. It is
carried at the weight $e^{a d_k(K_p)}$, and for a pointed $p$ its size is
$\mathfrak{N}^{\mathrm{ex}}_A(p)$ of \eqref{5.2:eq-carrier-exterior-d}.

For the non-empty tuples we use two clauses of the summation lemma for source-anchored
families, whose cell-sum at each link costs the factor $1+C_\Sigma C_e$. Clause (a) gives
\begin{equation*}
d_k(\operatorname{lab}_4)\le d_k(K_p)+c_l\sum_{t}|\mathsf{y}_t| .
\end{equation*}
Clause (b) is applied three times, once for each of the links four, three and two, and bounds
the exponential factor that it attaches to link $t$ by $e^{d_k(\operatorname{lab}_{t-1})+1}$.

Both clauses are applied to the triple (label, domain, size)
$=(K_p,Y_0(p),\mathfrak{N}^{\mathrm{ex}}_A(p))$. For a
compact pointed $p$ the label $K_p$ in the displays of this paragraph is to be read as
$\operatorname{lab}_1=K_p\cup\mathfrak{h}_v$, because
Corollary~\ref{5.2:cor-anchor-sum}, used below, is stated on $\operatorname{lab}_1$. For a
bundled $p$ the triple is that of its bundle, as fixed in clause (a).

The factor $e^{-D}=e^{-(\check{R}_k-2)}$ extracted in the proof of clause (c) is not extracted
here. The smallness of the pointed pieces is that of $\mathsf{w}_L^{-1}$, which is already
contained in $\mathfrak{N}^{\mathrm{ex}}_A(p)$ through its Cauchy prefactor
$2\mathsf{w}_L^{-1}(\eta_p\mathsf{f}_v)^{1/2}$; for the bundles the corresponding factor
$\theta'(g_k)$ is supplied by \eqref{5.2:eq-exterior-sourced-4} and not by the size
$\|\mathfrak{b}\|$.

The floor condition on $\kappa_1$ contains $\kappa_1-1\ge (a+2)c_l+\lambda_*$, which is what the
three links require. Then, using $d_{k,Z}(Y_4(p;\mathsf{y}))\le d_k(\operatorname{lab}_4)$ from
clause (a) of Lemma~\ref{5.2:lem-exterior-pieces},
\begin{equation}\label{5.2:eq-exterior-sourced-1}
\begin{aligned}
\sum_{(\mathsf{y}_2,\mathsf{y}_3,\mathsf{y}_4)}|(p;\mathsf{y})|\,
e^{a\,d_{k,Z}(Y_4(p;\mathsf{y}))}
&\le e^3\,\mathfrak{N}^{\mathrm{ex}}_A(p)\,e^{(a+3)d_k(K_p)}\\
&\le (1+C_\Sigma C_e)^3\,\mathfrak{N}^{\mathrm{ex}}_A(p)\,e^{(a+7)d_k(K_p)} .
\end{aligned}
\end{equation}
Here the three applications of clause (b) produce the factor $e^3$ and raise the weight from
$a$ to $a+3$. The empty tuple is the summand $1$ in the expansion of the product over the three
links. The last inequality uses $e^3\le(1+C_\Sigma C_e)^3$ and $a+3\le a+7$.

\emph{Clause (e), the anchor.} A continuation $(p;\mathsf{y})$ is filed at $Y_4(p;\mathsf{y})$.
In the norm $\|\cdot\|_a$ it is weighted at an anchor cube $\Delta_*\in Y_4(p;\mathsf{y})$ with
$\Delta_*\not\subset Z$. This cube may lie in the part of $\operatorname{lab}_4$ contributed by
the families, off $K_p$.

Pick $\Delta_p\in K_p$ with $\Delta_p\not\subset Z$. Since $\operatorname{lab}_4$ is connected,
$d_k(\operatorname{lab}_4)\ge \operatorname{dist}_k(\Delta_*,\Delta_p)-1$, and therefore
\begin{equation*}
e^{a\,d_k(\operatorname{lab}_4)}
\le e\cdot e^{-\operatorname{dist}_k(\Delta_*,\Delta_p)}\,e^{(a+1)d_k(\operatorname{lab}_4)} .
\end{equation*}
We insert this inequality into the left member of \eqref{5.2:eq-exterior-sourced-1} and apply the
fixed-$p$ argument at the weight $a+1$. The floor condition covers this raise, since
$(a+3)c_l\le (a+7)c_l$ and the proof of clause (c) used the floor condition in the form
$\kappa_1-1\ge c_l(a+7)+\lambda_*$. We then sum first over the position of $\Delta_p$ relative
to $\Delta_*$, which costs $S_d$. The result is
\begin{equation}\label{5.2:eq-exterior-sourced-2}
\sum_{p,\,\mathsf{y}:\ Y_4(p;\mathsf{y})\ni\Delta_*}|(p;\mathsf{y})|\,
e^{a\,d_{k,Z}(Y_4(p;\mathsf{y}))}
\le e\,S_d\,\sup_{\Delta\not\subset Z}\ \sum_{p:\ K_p\ni\Delta}
e^3\,\mathfrak{N}^{\mathrm{ex}}_A(p)\,e^{(a+4)d_k(K_p)},
\end{equation}
with $a+4\le a+7$. This is the family form of clause (c) of the summation lemma, in its first
case, read at a fixed anchor. The factor $e\,S_d$ it costs is the one displayed in
$\varepsilon^{\mathrm{ext},\mathsf{t}}(k)$ and in the ceiling condition with the constants of
\eqref{5.2:eq-exterior-pieces-b}.

\emph{Clause (e), every $Y_0$-term.} The three links are treated in the same way for every
$Y_0$-term. Later brackets hold fewer lower layers, and an absent link is an empty family, so
the bound only decreases.

\emph{Clause (e), the rows after the raise.} After the raise of the weight to $a+7$ the rows'
$X$-sums keep margins. Since $a=(31/20)\kappa$, the rows summed at $(7/4)\kappa$ and at
$2\kappa$ keep the margins
\begin{equation*}
\tfrac74\kappa-a-7=\tfrac15\kappa-7
\qquad\text{and}\qquad 2\kappa-a-7,
\end{equation*}
both $\ge c_d+2$ under the $d$-only floor $\kappa\ge\kappa_d$. Both follow from the column bounds
for the pointed rows, interior and exterior,
according to which every row retains a margin of at least $(\tfrac45-\tfrac12\delta)\kappa$ in
its rate, of which the part $\kappa/20\ge c_d+2$ pays for the sum over the localisation domains.
Hence the rows' constants are those already entered at
the weight $a$.

The chain member needs no margin. Since $a+7\le a_k=(8/5)\kappa$ for $\kappa\ge140$,
monotonicity of the norm in the weight gives, for every cube $Q\in\pi_k$ and every stage $n$,
\begin{equation*}
\sum_{X\ni Q}|\mathbb{C}^{(n)}_k(X)|\,e^{(a+7)d_k(X)}\le \|\mathbb{C}^{(n)}_k\|^{A}_{a_k}
\le C_0^{\sharp}.
\end{equation*}
The collar series of clause (i$'$) of the supply bound is re-run exactly as in the proof of
clause (c).

\emph{Clause (e), the three $\theta'$-members at weight $a+7$.} The derivations of the three
$\theta'$-members close at $a+7$ in place of $a$ without change. We take them in turn.

The first member consists of Ba{\l}aban's own brackets, bounded by the amplitude lemma together
with the size table. Its ingredients are the sizes times the units' rates.
\begin{itemize}
\item Rows one to three, five and $\partial$ are carried at the rate $(7/4)\kappa$, the
inductive weight of \cite[(1.66)]{B-CMP122b}.
\item The pointed rows are carried at the size $(81/64)\mathfrak{N}^{\flat}\le
(3/2)\mathfrak{N}^{\flat}$ on the radius $\mathsf{w}_L/(\eta_p\mathsf{f}_v)^{1/2}$, with Cauchy
prefactor $2\mathsf{w}_L^{-1}(\eta_p\mathsf{f}_v)^{1/2}$. Corollary~\ref{5.2:cor-anchor-sum},
applied at $Q_*:=\Delta$ and $r=a+7\le a+8$ for the compact $p$ on
$\operatorname{lab}_1=K_p\cup\mathfrak{h}_v$, gives \eqref{5.2:eq-exterior-sourced-3} below.
\item The row-four units of both domain classes of the plain units, together with the long
singles and tails of both classes, are grouped into bundles $\mathfrak{b}$ with the triple
(label, domain, size) $=(\operatorname{lab}\mathfrak{b},Y_0(\mathfrak{b}),\|\mathfrak{b}\|)$.
On the long rows of the class whose domains are disjoint from $Z$
the triples are $(\operatorname{lab}(K_v,\mathsf{y}),Y_0(p),\|\mathfrak{b}\|)$, with the label and
$\mathcal{R}_1$ as fixed in clause (a). For them \eqref{5.2:eq-exterior-sourced-4} below holds.
\end{itemize}
For the pointed rows,
\begin{equation}\label{5.2:eq-exterior-sourced-3}
\sup_{\Delta\not\subset Z}\ \sum_{p\ \mathrm{pointed}:\ \operatorname{lab}_1(p)\ni\Delta}
\mathfrak{N}^{\mathrm{ex}}_A(p)\,e^{(a+7)d_k(\operatorname{lab}_1(p))}
\le \theta'(g_k)\,\mathsf{C}^{\Sigma}_{a+7}.
\end{equation}
Corollary~\ref{5.2:cor-anchor-sum} performs the sums over $v$ (over every zone $n'\le k$),
over $s$, over $\square$ and over $\mathsf{y}$ at once, rooted at the source cell. No cell
beneath $\Delta$ is counted separately, and no per-cell reduction to clause (c) enters.

For the bundles,
\begin{equation}\label{5.2:eq-exterior-sourced-4}
\sup_{Q_*}\ \sum_{\mathfrak{b}:\ \operatorname{lab}(\mathfrak{b})\ni Q_*}
\|\mathfrak{b}\|\,e^{(a+7)d_k(\operatorname{lab}\mathfrak{b})}
\le \varepsilon^{\Sigma4}(k)\le \theta'(g_k)\,\mathsf{C}^{\Sigma,4}.
\end{equation}
This is the theorem on summation of boundary terms at finer zones, in its clauses for the
row-four units and for the long units of both classes,
at $r=a+7$, together with its corollary in the coupling ratio. By its corollary
for links two to four, the per-bundle bound and the anchor step above hold for the bundles as
written. Inside each $\pi_{j_v}$-cube the blocks and the fine cells are summed exactly first, by
the lemma on exact summation over blocks and fine cells. Then $v$, $s$ and
$\tilde{\mathsf{y}}$ are summed by the summation theorem in the unit's own grain, rooted at the
source. No count at a single scale and no Cauchy gain $\mathsf{w}_L^{-1}$ per unit is used.

Thus the first member carries the constant $\mathsf{C}_1^{\mathrm{ex}}$ of
\eqref{5.2:eq-exterior-pieces-b} in place of $C_{\sharp}O(1)$. Against the activity weight
stand the rate $2\kappa$ of the row-four units, the rate $a_k=(8/5)\kappa$ of the chain terms
$\mathbb{C}^{(n)}_k$, and the rate $(7/4)\kappa$ of the other rows. The last suffices because
\begin{equation*}
\tfrac74\kappa\ \ge\ a+7=\tfrac{31}{20}\kappa+7\qquad\text{for }\kappa\ge 35 .
\end{equation*}

The second member consists of the brackets of the chain terms $\mathbb{C}^{(n)}_k$. Each stage's
$X$-sum at a cell is its stage norm at $a_k=(8/5)\kappa$, and $a_k\ge a+7$ for $\kappa\ge140$.

The third member consists of the exterior arrested left-overs. The collar series of clause
(i$'$) of the supply bound is re-run at $a+7\le a+8$, exactly as in the proof of clause (c), with
the margin $(25/12)L^s\ge ds$ of $\mathsf{K}^{\mathrm{arr}}$.

The row-four $\mathsf{t}$-pieces entered in the second and third members are the bundles of the
theorem on summation of boundary terms at finer zones and are counted in
$\mathsf{C}^{\Sigma,4}$. The two re-runs therefore stand in
\eqref{5.2:eq-exterior-pieces-b} only as majorants on which no estimate draws.

\emph{Clause (e), conclusion.} In each of the three members the summand for $p$ is multiplied by
$(1+C_\Sigma C_e)^3$, by \eqref{5.2:eq-exterior-sourced-1}, and nothing else changes. The anchor
step \eqref{5.2:eq-exterior-sourced-2} contributes the overall factor $e\,S_d$. Summing over $p$
gives (e).
\end{proof}

\begin{corollary}[Exterior pieces at a strip cell sum to at most one]\label{5.2:cor-strip-cell-sum}
Assume the hypotheses of the lemma on pieces of units whose label leaves the large-field region
(Lemma~\ref{5.2:lem-exterior-pieces}), and in addition:
\begin{enumerate}
\item[(i)] the interior ceiling, by which the constant $\varepsilon^{\mathrm{ext},\sigma}(k)$ of
\eqref{5.2:eq-exterior-pieces-a} satisfies $\varepsilon^{\mathrm{ext},\sigma}(k)\le 1/2$ at every
$k\le K$;
\item[(ii)] the bracket ceiling with the constants of \eqref{5.2:eq-exterior-pieces-b}, by which
the constant $\varepsilon^{\mathrm{ext},\mathsf{t}}(k)$ defined there satisfies
$\varepsilon^{\mathrm{ext},\mathsf{t}}(k)\le 1/2$.
\end{enumerate}
Put
\begin{equation}\label{5.2:eq-strip-cell-sum-1}
\varepsilon^{\mathrm{ext}}(k):=\varepsilon^{\mathrm{ext},\sigma}(k)
+\varepsilon^{\mathrm{ext},\mathsf{t}}(k),
\end{equation}
the sixth member of $\varepsilon_{\mathrm{br}}(k)$. Let $\Delta\subset C^{\sharp}\setminus C$ be a strip
cell. Then the pieces of $\mathfrak{L}^{\mathrm{ex}}=\mathfrak{L}^{c}\sqcup\mathfrak{L}^{Y_0}$ filed at
a $Y_4\subset C^{\sharp}$ meeting $\Delta$ sum there to at most $\varepsilon^{\mathrm{ext}}(k)$, and
$\varepsilon^{\mathrm{ext}}(k)\le 1$. Consequently the strip-cell clause of the volume inclusion holds
with
\begin{equation*}
\mathsf{K}_{\mathrm{cell}}(k,n(\Delta))+\varepsilon^{\mathrm{ext}}(k)
\end{equation*}
in the place of $\mathsf{K}_{\mathrm{cell}}(k,n(\Delta))$, where $n(\Delta)$ is the zone of $\Delta$, and
this addition is carried in the volume ceiling as $\mathsf{K}_{\mathrm{cell}}^0+1$.
\end{corollary}

\begin{proof}
By hypotheses (i) and (ii) and the definition \eqref{5.2:eq-strip-cell-sum-1},
\begin{equation*}
\varepsilon^{\mathrm{ext}}(k)\le \tfrac12+\tfrac12=1 .
\end{equation*}
The pieces of $\mathfrak{L}^{\mathrm{ex}}$ filed at a $Y_4\subset C^{\sharp}$ meeting $\Delta$ lie under
the supremum defining $\|\cdot\|_{a,1}$ in Lemma~\ref{5.2:lem-exterior-pieces}, taken at the
$\pi_k$-cube of $\Delta$, and there each of them carries the weight $e^{a\,d_{k,Z}(Y_4)}$, which is
at least $1$ because $a=(31/20)\kappa>0$ and $d_{k,Z}(Y_4)\ge 0$. Hence their unweighted sum at
$\Delta$ is bounded by the corresponding weighted sum.

For the class $\mathfrak{L}^{c}$ these pieces are the summands of
$\mathfrak{v}^{\mathrm{ext},\sigma}(Y_4)$, and clause (c) of
Lemma~\ref{5.2:lem-exterior-pieces} bounds their weighted sum at the cube of $\Delta$ by
$\|\mathfrak{v}^{\mathrm{ext},\sigma}\|_{a,1}\le\varepsilon^{\mathrm{ext},\sigma}(k)$. For the class
$\mathfrak{L}^{Y_0}$ the pieces are the continuations $(p;\mathsf{y}_2,\mathsf{y}_3,\mathsf{y}_4)$
with not all $\mathsf{y}_t$ empty; the all-empty tuple $v_p^{\emptyset}$ is, by clause (a) of
Lemma~\ref{5.2:lem-exterior-pieces}, the summand $v_p$ already accounted for in the three
$\theta'$-members. These continuations are summed in clause (e) of that lemma, in the family form at
an anchor off $Z$, at weight $a+7$ in place of $a$ and with the first member's constant replaced by
$\mathsf{C}_1^{\mathrm{ex}}$; the part of $\varepsilon^{\mathrm{ext}}(k)$ which that clause adds for
them is $\varepsilon^{\mathrm{ext},\mathsf{t}}(k)$ of \eqref{5.2:eq-exterior-pieces-b}, and their
weighted sum at the cube of $\Delta$ is bounded by $\varepsilon^{\mathrm{ext},\mathsf{t}}(k)$.
Adding the two classes, the pieces of $\mathfrak{L}^{\mathrm{ex}}$ at $\Delta$ sum to at
most $\varepsilon^{\mathrm{ext},\sigma}(k)+\varepsilon^{\mathrm{ext},\mathsf{t}}(k)
=\varepsilon^{\mathrm{ext}}(k)\le 1$.

Therefore the strip-cell clause of the volume inclusion reads
$\mathsf{K}_{\mathrm{cell}}(k,n(\Delta))+\varepsilon^{\mathrm{ext}}(k)$, and, since
$\varepsilon^{\mathrm{ext}}(k)\le 1$, the volume ceiling, with the weights fixed in clause (c) of the
weight proposition, carries it as $\mathsf{K}_{\mathrm{cell}}^0+1$.
\end{proof}

\begin{definition}[The physical torus and the auxiliary torus]\label{5.2:def-auxiliary-torus}
Put $K' := K - m$.
\begin{enumerate}
\item[(a)] The local run, by which we mean the renormalisation steps at the levels $k \le K'$, is
carried out on the physical torus
$\mathbb{T}_m$. On it the lower torus-size condition, with the torus-compatibility constant
$C_{\mathbb{T}}$, ensures that the cube partitions $\pi_k$ of Ba{\l}aban's construction exist
at every level, that is, that their cubes divide the torus.
\item[(b)] Let $N_{\mathrm{aux}}$ be the least power of $L$ such that
\begin{equation}\label{5.2:eq-auxiliary-torus-1}
N_{\mathrm{aux}} \ge C_{\mathbb{T}}\, R(x_0 + 2c_2).
\end{equation}
For an integer $n \ge 1$, the auxiliary torus of depth $n$ is the periodic lattice of spacing
$L^{-n} N_{\mathrm{aux}}^{-1}$ with $N_{\mathrm{aux}} L^{n}$ sites per side, so that its side has
length $1$. The auxiliary torus $\mathbb{T}_{\mathrm{aux}}$ is the auxiliary
torus of depth $K+1$, that is, the periodic lattice of spacing
\begin{equation}\label{5.2:eq-auxiliary-torus-2}
\varepsilon_{\mathrm{aux}} := L^{-(K+1)} N_{\mathrm{aux}}^{-1}
\end{equation}
with $N_{\mathrm{aux}} L^{K+1}$ sites per side.
The frozen whole-lattice run, namely clauses (b) and (c) of the level-$k$ statement of the frozen
run and the inductive hypothesis on the running shift, for all levels $k \le K+1$, is performed on
$\mathbb{T}_{\mathrm{aux}}$; the frozen families are the families of terms which this run
constructs on the whole lattice.
\item[(c)] For the levels $k \le K'$ the physical torus $\mathbb{T}_m$ carries its own copy of the
frozen families; these copies are the inputs of clause (b) of the inductive step statement in the
proof of the correlation theorem (Theorem~\ref{5.1:thm-correlation-theorem}).
\item[(d)] On $\mathbb{T}_{\mathrm{aux}}$ the partitions $\pi_k$ have room and divide the torus at
every level $k \le K+1$; the numbers $b_j(\zeta)$, the running shifts $\zeta_j(\zeta)$, the
reference values $x_j$ and the integer $K$ do not depend on the torus on which they are computed;
and the depth $K+1$ of $\mathbb{T}_{\mathrm{aux}}$ is thereby well defined.
\end{enumerate}
\end{definition}

\begin{proof}[Verification of clause~(d)]
We show that the auxiliary torus leaves room for, and is divided by, Ba{\l}aban's partitions at
every level $k \le K+1$, and that performing the frozen run on $\mathbb{T}_{\mathrm{aux}}$ rather
than on $\mathbb{T}_m$ changes none of the numbers that the local run reads.

\emph{Room and divisibility.} By \eqref{5.2:eq-auxiliary-torus-2}, the lattice of
$\mathbb{T}_{\mathrm{aux}}$ reached after $k$ blocking steps, $k \le K+1$, has spacing
$L^{k}\varepsilon_{\mathrm{aux}}$ and, by the definition of $\mathbb{T}_{\mathrm{aux}}$,
$N_{\mathrm{aux}} L^{K+1-k}$ sites per
side. By the construction of the partitions $\pi_k$, the sides of the partition cubes at level $k$
are powers of $L$ not exceeding
$C_{\mathbb{T}}\check{R}_k$, where $\check{R}_k = R(\min_{i \le k} x_i)$. The function $R$,
the least power of $L$ not smaller than $(\log x)^{r_{\mathrm{pr}}}$, is nondecreasing in $x$;
since $\min_{i \le k} x_i \le x_k$, this gives $\check{R}_k \le R(x_k)$. The bound on the
reference sequence gives $R(x_k) \le R(x_0 + 2c_2)$, and \eqref{5.2:eq-auxiliary-torus-1}
gives $C_{\mathbb{T}} R(x_0+2c_2) \le N_{\mathrm{aux}}$. Hence every partition-cube side $\ell$
at level $k$ satisfies
\begin{equation}\label{5.2:eq-auxiliary-torus-3}
\ell \le C_{\mathbb{T}}\check{R}_k \le C_{\mathbb{T}} R(x_k)
\le C_{\mathbb{T}} R(x_0 + 2c_2) \le N_{\mathrm{aux}} .
\end{equation}
This is the room. Since $N_{\mathrm{aux}}$ is itself a power of $L$, every power of $L$ not
exceeding $N_{\mathrm{aux}}$ divides $N_{\mathrm{aux}}$. Moreover $N_{\mathrm{aux}}$ divides
$N_{\mathrm{aux}}L^{K+1-k}$ for
$k \le K+1$. Therefore every partition-cube side at level $k$ divides the number of sites per side
of the level-$k$ lattice, and the partitions exist at every level $k \le K+1$. Finally, the class
cube at level $k$, which is a cube of side $100M\check{R}_k$, satisfies, since
$C_{\mathbb{T}} = L^{11}M$ and $L^{10} \ge 100$ (as $L \ge 2$), and by
\eqref{5.2:eq-auxiliary-torus-3},
\begin{equation*}
100M\check{R}_k \le L^{10}M\check{R}_k = C_{\mathbb{T}}\check{R}_k / L \le N_{\mathrm{aux}}/L ,
\end{equation*}
which is the condition under which the class cube fits, as the construction of the class cubes
requires.

\emph{Legitimacy.} By the first defining property of the families, the terms of every family whose
localisation domains do not wrap around the torus are independent of the torus on which they are
computed. By the lemma on the torus-independence of the flow numbers, it follows that each of the
following is torus-independent: every number $b_j(\zeta)$, every running shift
$\zeta_j(\zeta)$, every reference value $x_j$, and the integer $K$. In particular, at the levels
$k \le K'$ the copy of the frozen families carried by $\mathbb{T}_m$ and the copy on
$\mathbb{T}_{\mathrm{aux}}$ give the same numbers $b_j$ as inputs to clause (a) of the inductive
step statement in the proof of the correlation theorem (Theorem~\ref{5.1:thm-correlation-theorem}).

\emph{Determination of $K$.} The integer $K$, which fixes the depth of $\mathbb{T}_{\mathrm{aux}}$,
is not known before the run. For a given depth $n$, the auxiliary torus of depth $n$ produces the
numbers $b_{k+1}(0)$, $k < n$. By the torus-independence just established, these numbers do not
depend on $n$. The integer $K$ is the first-passage index determined by them, and one takes
$n = K+1$.

\emph{Absence of a dependence on the auxiliary size.} The integer $N_{\mathrm{aux}}$ enters no
constant of the estimates. The left side of the torus bound does not depend on $m$.
\end{proof}

\begin{definition}[The induction format with sources]\label{5.2:def-induction-format}
Throughout, $0\le k\le K$, $w=\sum_i z_i f_i$ is a complex source, $T^{(j)}$ is the unit lattice
after $j$ steps, $\mathcal{W}_j(X)$ is the class of complex Lipschitz weights at level $j$ over
$X$, and $A(c,U)$ is the weighted Wilson action with plaquette weight $c$.

\smallskip
\noindent\emph{Counterterm numbers.} For a level $j$ and a complex number $\zeta$, let
$E^{(j)}_{\zeta}$ be the frozen whole-lattice family at level $j$ and constant source value
$\zeta$, as given by the lemma on frozen whole-lattice families. Put
\begin{equation}\label{5.2:eq-induction-format-1}
  b_j(\zeta) := \beta\bigl[E^{(j)}_{\zeta}\bigr],
\end{equation}
the number defined in \cite[(3.61)]{B-CMP119}, read for this family. The family is a single
object: it is formed with the cut $\delta_{j-1}$ and the radii of \cite[(2.28)]{B-CMP119}, the
operation $\tilde{\mathbf{R}}$ is folded into it, and the gauge-fixing coefficient is carried
along in it. The number $b_j(\zeta)$ is the one subtracted at the counterterm steps of the
induction in the proof of the correlation theorem.

\smallskip
\noindent\emph{The format.} The sourced effective density $\rho^w_k$ at level $k$ is in the
\emph{induction format with sources} if
\begin{equation}\label{5.2:eq-induction-format-2}
  \rho^w_k=\sum_{\{\Omega_j,\Lambda_j,I_j\}}\chi_k\,\mathbf{T}^w_k\,\exp A^w_k ,
\end{equation}
as in \cite[(2.18)]{B-CMP119}, where the summation also runs over the label $I_j$ carried by
the large-field step on the refined support, and where $\chi_k$ contains the rim
functions of the rim construction, as given by clause (f) of the component factorisation lemma;
and if
\begin{equation}\label{5.2:eq-induction-format-3}
  A^w_k=-A\bigl(x_k(\cdot)-\zeta_k,U_k\bigr)+\mathbf{E}^w_k+\mathbf{D}^w_k
        +\mathbf{R}^w_k+\mathbf{B}^w_k-E^w_k ,
\end{equation}
with the following clauses (a)--(f), clause (d) being refined by (d$'$).

\begin{itemize}
\item[(a)] \emph{The running shift.} For every plaquette $p$,
\begin{equation}\label{5.2:eq-induction-format-a}
  \zeta_k(p):=w(p)+\sum_{j\le k}\varphi_j(p)\sum_{y\in T^{(j)}}h^{(j)}_y(p)
  \bigl[b_j\bigl(w(p_y)\bigr)-b_j(0)\bigr],
\end{equation}
where $b_j$ is the number \eqref{5.2:eq-induction-format-1}, the functions $h^{(j)}_y$ are those
of \cite[(3.40), (3.65)]{B-CMP119}, so that $h^{(j)}_y\ge 0$ and $\sum_y h^{(j)}_y=1$,
$\varphi_j$ is the function of \cite[(2.24)]{B-CMP119}, and $p_y$ is the plaquette attached to
the site $y$.

\item[(b)] \emph{The localised counterterm groups.} Let $\Lambda_j^{\circledR}\subset T^{(j)}$ be
the set of sites at level $j$ at which the renormalised pointed groups $\mathfrak{i}^{(j)}_y$ are
formed; it is a datum fixed in the construction of these groups. At the sites
$y\in T^{(j)}\setminus\Lambda_j^{\circledR}$ the counterterm piece stays unmatched. The term
$\mathbf{E}^w_k$ is
\begin{equation*}
  \mathbf{E}^w_k:=\sum_{j\le k}\Bigl[\sum_{y\in\Lambda_j^{\circledR}}\mathfrak{i}^{(j)}_y(U_k)
  -\sum_{y\in T^{(j)}\setminus\Lambda_j^{\circledR}}
  b_j(\hat w)\,A\bigl(h^{(j)}_y\varphi_j,U_k\bigr)\Bigr],
\end{equation*}
where in the summand indexed by $y$ the constant source value is $\hat w=w(p_y)$. The groups
$\mathfrak{i}^{(j)}_y$ are built on the frozen families $E^{(j)}_{\hat w}$,
$\hat w\in D(0,2r)$, the disc of centre $0$ and radius $2r$ in $\mathbb{C}$, where $r>0$ is the
radius of constant source values fixed with these families in the lemma on frozen whole-lattice
families; the second sum
consists of the unmatched counterterm pieces of Ba{\l}aban's construction, whose moduli are
controlled by the last clause of the smallness conditions on the sources.

\item[(c)] \emph{The local terms and their source differences.} The local terms
$E^{(j)}(X,\cdot,y;w)$ satisfy \cite[(2.27)(i)--(iv)]{B-CMP119} uniformly on
$\mathcal{W}_j(X)$, on which class each such term is defined by the formula that produces it,
with inputs $\zeta_{j-1}$ restricted to $X^+$ and old terms on sets $X'^{+}\subset X^+$. Each is
analytic in $w$, depends on $w$ only through the restriction of $w$ to $X^+$, and equals
$E^{(j)}_{\hat w}$ when $w\equiv\hat w$ on $X^+$. Put
\begin{equation*}
  D^{(j)}:=E^{(j)}(\cdot\,;w)-E^{(j)}_{w(p_y)}
\end{equation*}
and
\begin{equation*}
  \mathbf{D}^w_k:=\sum\Bigl[D^{(j)}(X,U_k,y;w)-D^{(j)}(X,1,y;w)\Bigr],
\end{equation*}
the sum running over the indices $j$, $X$, $y$ of the local terms.

\item[(d)] \emph{The terms $\mathbf{R}^w$ and $\mathbf{B}^w$.} The terms $\mathbf{R}^w$ satisfy
\cite[(2.30)--(2.31)]{B-CMP119}. The terms $\mathbf{B}^w$ satisfy \cite[(2.40)]{B-CMP119} and
the following relative complex condition at level $k$: every term has a label
$A\in\mathfrak{D}_j$ and a set of live regions $(Y,m)$ whose strips $Y^\sharp$ touch $A$; it
depends only on the fields on $A\cup\bigcup Y^\sharp$; it carries an orbit datum of radius
$\varepsilon_k$ on the domain and bound with which it is stored; and
\cite[(2.42)]{B-CMP119} holds for it in $d_j(A)$, in the size
$\mathfrak{N}^{A,\mathbb{C}}$, that is,
\begin{equation}\label{5.2:eq-induction-format-4}
  \bigl|\mathbf{B}^{(j)}(X)\bigr|\le B_0\exp\bigl(-\kappa\, d_{j,\cup Y^\sharp}(X)\bigr),
\end{equation}
in the length of \cite[(1.67)]{B-CMP122b}. This is the format of boundary terms with live
regions, given by the defining conditions of that format and by its corollary on boundary terms;
when the set of live regions is empty it is \cite[(2.41)--(2.42)]{B-CMP119} as printed there,
together with the orbit datum of the definition of orbit data. The bound holds uniformly on
$\mathcal{W}_j(X)$; the terms are analytic and local, depending only on the restriction of
$((\mathbf{U},\mathbf{J}),\mathbf{A},w)$ to $X_F$ and of $w$ to $X_F^{+}$; and
\cite[(2.32)]{B-CMP119} holds at constant sources $w\equiv\hat w$.

Where the correlation theorem asserts the format of this definition with Ba{\l}aban's bounds
under the smallness conditions on the sources, the bounds meant for $\mathbf{B}^w$ are
\cite[(2.42)]{B-CMP119} in the length of \cite[(1.67)]{B-CMP122b}; this is Ba{\l}aban's own form
of the bound for these terms (\cite[(1.99)]{B-CMP122b}; \cite{B-CMP119}).

Moreover, in the two-sided sense fixed by the corollary on \cite[(2.42)]{B-CMP119} for the rim
construction, \cite[(2.42)]{B-CMP119} holds for every rim datum $\mathbf{A}$ in the set
$\mathfrak{X}^\sharp_j$, which contains the support of the rim functions; the terms $\mathbf{B}$
are stored with the prepared rim datum $A^\flat$ of the corollary on the prepared rim datum.
In this two-sided statement the rim datum ranges over $\mathfrak{X}^\sharp_j$; the corollary on
the box of the rim construction shows that this restriction of the range is needed for the
two-sided statement.

\item[(d$'$)] \emph{The classes within $\mathbf{B}^w$.} Within $\mathbf{B}^w$ the following three
classes of terms are distinguished, each with the properties stated for it.

$(\alpha)$ The terms $\mathbf{B}'^{(j)}(A,\mathcal{S})$ of the first group produced by the
operation $\mathbf{R}$, $\mathcal{S}$ being the strong set of the term, as in clause (e) of the
proposition on the $\mathbf{R}$-operation with sources, and their descendants produced by the
operation $\mathbf{T}$, in wrapped form, are
terms in the format $(A_F,\mathcal{P}_F)$ of boundary terms with live regions: they satisfy
\cite[(2.40)--(2.41)]{B-CMP119} with Ba{\l}aban's side conditions imposed on the label, namely
$A_F\cap\Omega_j\neq\emptyset$; $A_F\cap Z_j^{\boxplus}\neq\emptyset$ or
$\mathcal{P}_F\neq\emptyset$; and joint $G$-valued invariance, $G=\mathrm{SU}(N)$.
In place of the plain length in \cite[(2.42)]{B-CMP119} they satisfy the relative bound
\begin{equation}\label{5.2:eq-induction-format-b}
\begin{aligned}
  \mathfrak{N}^{A,\mathbb{C}}\bigl(\mathbf{B}'^{(j)}(A,\cdot\,;w)\bigr)
  &\le e^{\lambda_\theta}\,c_1^{\flat}(g_j)\,
     e^{-(1+\frac14\beta_s)\kappa d_j(A)}\\
  &\le B_0\,e^{-\kappa d_j(A)},
\end{aligned}
\end{equation}
respectively, for a descendant produced by the operation $\mathbf{T}$ at step $j$, with $A:=X$,
the bound
$B_0e^{-\kappa_{\mathrm{out}}d_j(A)}$ when it is fresh and $B_0e^{-\kappa d_j(A)}$ when it is
stored. The first bound is clause (c) of the proposition on boundary terms with live regions,
which rests on \cite[(1.99)]{B-CMP122b}; the bound for descendants is clause (e) of the
proposition on the $\mathbf{T}$-operation with sources. These bounds hold uniformly on
$\mathcal{W}_j(X_F)$; the terms are analytic and local, and are stored with the prepared rim
datum $A^\flat$, by clause $(\gamma)$ of the corollary on the prepared rim datum. This is the
alternative form of
the inductive assumption that \cite{B-CMP119} allows, in the words of \cite{B-CMP119}:
``(2.42) with $d_j(X)$ replaced by $d_j(X\setminus Z_j)$ \ldots\ All constructions and proofs
\ldots\ can be carried on''; here the strips of the live regions are excised through the label.

A term does not outlive its live regions: if $(Y,m)\in\mathcal{P}_F$ ceases to be live at
$\mathbf{R}_k$, then $F$ belongs at that step to the first part of the expansion and is
resummed into $V^\sharp$, by clause (b) of the proposition on the resummation into $V^\sharp$.

Not in this class are the two members kept under $\mathbf{T}'_{k_W}(Y_i)$ whose domain is
$Y_i^\sharp$ itself, namely $V^\sharp(Y_i^\sharp)$ and the form of \cite[(1.71)]{B-CMP122b},
whose exponent, in the notation of \cite[p.~378]{B-CMP122b}, is
\begin{equation*}
  -\tfrac12\bigl\langle DH''_{1,k,X}B,\ \zeta\, DH''_{1,k,X}B\bigr\rangle .
\end{equation*}
These two members are governed by the tube statement, which is
used here by statement only and is not proved here. Its content, as used here, is the
following. The form of \cite[(1.71)]{B-CMP122b} is controlled through the coupling slot, by the
form lemma of the tube statement, which rests on the slot theorem: for a large-field component $C$
with strip $C^\sharp$ and its complex coupling slot $\mathfrak{q}_C$, evaluated at a
configuration pair, written $\zeta$ in the next display only (it is not the constant source
value of \eqref{5.2:eq-induction-format-1}), and at the pair $(U,0)$,
\begin{equation*}
  \bigl|\mathfrak{q}_C(\zeta)-\mathfrak{q}_C((U,0))\bigr|\le\Sigma^{\mathfrak{q}}(x_{k_W}),
\end{equation*}
where the majorant satisfies $\Sigma^{\mathfrak{q}}(x_{k_W})\downarrow 0$ in the sense of the tube
statement, at every admissible splitting. For
$V:=V^\sharp(C^\sharp;\cdot)$ the two-point oscillation is at most $C_\Sigma\check{R}^{p_\Sigma}$,
with a constant $C_\Sigma$ and
\begin{equation*}
  p_\Sigma=2d+6+r_{\mathrm{pr}}+\lceil 2p_1/r_{\mathrm{pr}}\rceil ,
\end{equation*}
by the oscillation theorem for $V^\sharp$, the localised members $\mathfrak{v}^{\mathbf{C}}$ of
$\mathbf{C}_k$ being listed in $V^\sharp$ under the filing of the $\mathbf{C}$-terms and the
Wilson-density count. This oscillation is paid at every later site, in own-scale cells and in
$d^{\wedge}$, through the collar of the region $Y_i$ itself, by the collar corollary. No
further property of $V$ is asserted.

$(\beta)$ The arrested left-overs $\mathfrak{F}(\mathfrak{a})$ of every live arrest
$\mathfrak{a}$, that is, the frozen terms $\mathbb{C}^{(i)}_{k_a}(X,\mathcal{S};w)$ and the
further frozen terms listed with them in the definition of arrests, are carried on the bounds
and domains fixed at their arrest, by the clause of the arrest construction on frozen terms
(frozen terms keep their bounds and domains), with the datum
of clause (a) of the proposition on the $\mathbf{R}$-operation with sources and with
$\|\cdot\|^{A}\le C_0^{\sharp}$ per stage at step $k_a$, the further frozen terms being bounded
by $1$ per block. They remain unchanged from their arrest to their dissolution, by the clause of
the arrest construction on dissolution. Their
dependence on $w$ is given by the weight corollary for arrested terms at
$w\in\mathcal{W}_{k_a}(X)\supset\mathcal{W}_k(X)$. Clause (c) of the proposition on the
$\mathbf{R}$-operation with sources takes them out of $(\beta)$, and the filing step of the
induction in the proof of the correlation theorem returns them into it.

$(\gamma)$ The format carries the orbit data used later in the large-field analysis. Every
$\mathbf{B}$-type term of $\mathbf{B}^w$ with live exterior slots, that is, the terms
$\mathbb{C}^{(n)}_j(X,\mathcal{S};w)$, the terms $\mathbf{B}'^{(j)}$ produced by $\mathbf{R}$,
the stored terms $\mathbf{B}^{(j)}$ with slots, and the frozen terms $\mathfrak{F}(\mathfrak{a})$,
is stored together with its orbit datum $c^\sharp(X;P,\mathbf{A},\boldsymbol{\Phi},g;w)$ in the
sense of the definition of orbit data. The orbit datum is holomorphic jointly on the product of
the carrier chart, the chart $\mathfrak{G}_{\varepsilon_k(s)}(X)$ of radius $\varepsilon_k(s)$ of
the definition of orbit data, and $\mathcal{W}_j(X)$; it equals the
term on the real slice; and it is bounded by $\mathfrak{N}^{\mathbb{C}}_{\varepsilon_k}(\cdot)$,
which dominates the size on the real slice and contains the supremum over $w$. Its chart radius
at each live law site $s$ is the running radius
\begin{equation*}
  \varepsilon_k(s)=\mathsf{e}(s)-\sum_{k'<k}\eta_{k'}(s)\ \ge\ \tfrac34\,\mathsf{e}(s),
\end{equation*}
where $\eta_{k'}(s)$ is the amount by which the radius is decreased at step $k'$.
The orbit datum is returned into the format at every site by the outputs of the orbit-data
analysis: at the large-field step $\mathbf{T}$ by clause (d) of the complex-orbit corollary at
the large-field step, over the chart $\mathfrak{G}_{\varepsilon'}(Y)$, with $w$ among the
variables; at the preparatory stage by
the corollary on the preparatory stage; at the localisation step by clause (b) of the proposition
of that analysis on the localisation step and by the corollary on the brackets. It is returned as
well by the filing steps
of the induction in the proof of the correlation theorem, which file the new terms together with
their orbit data as they file their bounds on the real slice. Nothing in $(\alpha)$, $(\beta)$ or
(a)--(f) is weakened: $(\gamma)$ adds a stored co-ordinate, not a bound.

\item[(e)] \emph{The large-field factor.} The factor $\mathbf{T}^w_k$ satisfies
\cite[(2.19)--(2.22)]{B-CMP119} with the two complex weights of the sourced step and with the
rim functions, is factorised over components as in the component factorisation lemma, and obeys
the modification rules of the large-field operations.

\item[(f)] \emph{The constant.}
\begin{align*}
  E^w_k=E&-\sum_{j}\ \sum_{X\subset\Lambda_j}e^{(j)}(X;w)\\
  &-\bigl(\text{the $z$-free logarithms, including }\log\mathfrak{z}(a_j(y))\bigr),
\end{align*}
where $\mathfrak{z}$ is the one-bond Haar integral of the gauge-fixing move.
\end{itemize}
No covariance is asserted for non-constant $w$.
\end{definition}

\begin{definition}[The hypothesis of the joint induction]\label{5.2:def-joint-hypothesis}
Let $1\le k\le K$. All objects below live on the tori of
Definition~\ref{5.2:def-auxiliary-torus}, the physical torus and the auxiliary torus
$\mathbb{T}_{\mathrm{aux}}$. Here $r>0$ is the radius fixed earlier in this chapter. Write
$D(0,2r)$ for the disc of radius $2r$ about the origin in
the plane of the constant shift $\zeta$. We say that \emph{the hypothesis of the joint
induction holds at level $k$} if the following four clauses hold.
\begin{enumerate}
\item[(a)] The reference statement holds for every pair to which it applies whose two
entries are both at most $k$.
\item[(b)] For every $j\le k$, the running shift $\zeta_j$ and the marginal coefficient
$b_j(\zeta)=\beta[E^{(j)}_{\zeta}]$ of Definition~\ref{5.2:def-induction-format} are
analytic functions of $\zeta$ on $D(0,2r)$ and are real for real $\zeta$. For every
$\zeta\in D(0,2r)$ they satisfy
\begin{equation*}
|b_j(\zeta)-b_j(0)|\le d_{j-1}(2r)\,|\zeta|,
\qquad
|\zeta_j-\zeta|\le \tfrac13\,|\zeta| .
\end{equation*}
Here $d_{j-1}(2r)$ is the flow function $d_i(\varrho)$ of the reference recursion taken at the
index $i=j-1$ and at $\varrho=2r$; it is not the size $d_j(X)$ of a localisation domain.
\item[(c)] For every $j\le k$, the family $\mathsf{P}^{(j)}_{\xi}$ belongs to the ball fixed
for these families and is holomorphic on $\mathfrak{P}_{j-1}$. Moreover, with
$Q^{(j)}_{\zeta}=E^{(j)}_{\zeta}-P^{\flat(j)}_{\zeta}$,
\begin{equation*}
\|E^{(j)}_{\zeta}\|^{\flat}\le E_0,
\qquad
\|P^{\flat(j)}_{\zeta}\|^{\flat}\le E_0,
\qquad
\|Q^{(j)}_{\zeta}\|^{\flat}\le \mathfrak{m}_{j-1},
\end{equation*}
where $E_0$ is the bound fixed earlier in this chapter for these families.
\item[(d)] For all $\zeta$ with $|\zeta|<2r$, the induction hypothesis of the frozen
whole-lattice run at constant shift $\zeta$ holds at level $k$. In particular, for every
$j\le k$,
\begin{equation*}
\|R^{(j)}_{\zeta}\|^{\flat}\le g_j^{\kappa_0},
\end{equation*}
where $\kappa_0$ is the exponent fixed in that induction hypothesis.
If in addition $k\le K'$, then for every $w\in\mathcal{W}_k$ the induction hypothesis with
sources of Definition~\ref{5.2:def-induction-format} holds at level $k$ on the physical
torus, with
\begin{equation*}
|\zeta_j(p)|\le \tfrac43\,\sup|w|
\end{equation*}
for all $j\le k$ and all plaquettes $p$.
\end{enumerate}
At level $k=0$ the hypothesis of the joint induction is that the density is the sourced
density $\rho_0^w$ and that the running shift is $\zeta_0=w$; nothing else is required.
\end{definition}

\begin{proof}[Joint induction: the pure step and the frozen integral]
The passage from $(\mathrm{J}_k)$ to $(\mathrm{J}_{k+1})$, the hypothesis of the joint induction
of Definition~\ref{5.2:def-joint-hypothesis}, is carried out in several steps; the first two are
the pure step and the frozen whole-lattice integral. Throughout, $k<K$,
so that $g_j<g$ for every $j\le k$. For a complex number $c$ and a radius $a>0$ we write $D(c,a)$
for the open disc of radius $a$ about $c$.

\emph{The pure step.} Let $\xi=(\xi_j)\in\mathfrak{P}_k$. Then every slot combination $s_j$
lies in the disc $D(1,\sigma_0)$. For each $j\le k$ the complex-slot proposition bounds the family
$\mathsf{P}^{(j)}_{\xi}$ a priori,
\begin{equation*}
\|\mathsf{P}^{(j)}_{\xi}\|\le \tfrac12 E_0+E^{*}\le E_0 ,
\end{equation*}
and this bound does not depend on the inputs: the bound furnished by that proposition involves
no property of the families of the earlier levels. Consequently the tuple
$\mathsf{P}^{(\le k)}_{\xi}=(\mathsf{P}^{(j)}_{\xi})_{j\le k}$ lies in the ball $\mathbb{B}_k$ of
complex histories. Since $\mathsf{P}^{(\le k)}_{\xi}\in\mathbb{B}_k$, the family
$\mathsf{P}^{(k+1)}_{\xi}$ is defined. Each of its terms is holomorphic in each variable $\xi_j$
separately, the others being held fixed, and is locally bounded on $\mathfrak{P}_k$. By Osgood's
lemma (a locally bounded function of finitely many complex variables on an open set that is
holomorphic in each variable separately is holomorphic), each term of $\mathsf{P}^{(k+1)}_{\xi}$
is jointly holomorphic in $\xi$ on $\mathfrak{P}_k$.

The marginal coefficient $f_{k+1}=\beta[\Phi_{k+1}]$ of the one-step family $\Phi_{k+1}$ is
therefore holomorphic on $\mathfrak{P}_k$, and clause~(c) of the complex-slot proposition gives
\begin{equation}\label{5.2:eq-pure-step-a}
|f_{k+1}|\le C_\beta\,\mathrm{e}_k .
\end{equation}
Finally, by the slot lemma, the family $\mathsf{P}^{(k+1)}_{\xi}$ so obtained belongs to the ball;
this is the membership statement required in clause~(c) of $(\mathrm{J}_{k+1})$.

\emph{The frozen whole-lattice integral (the bridge).} The construction of the frozen integral at
level $k+1$ uses only data of levels at most $k$, namely: the cut $T_k$; the families
$(E^{(j)}_{\zeta})_{j\le k}$; the remainders $(R^{(j)}_{\zeta})_{j\le k_1(k)}$, where $k_1(k)<k$;
and the bound $|\zeta_k|\le\frac83 r$ on the running shift. All of these are contained in
$(\mathrm{J}_k)$; for the last three the relevant clauses are as follows: the families
$E^{(j)}_{\zeta}$, $j\le k$, and their bounds are those of clause~(c) of
Definition~\ref{5.2:def-joint-hypothesis}; the remainders $R^{(j)}_{\zeta}$ with $j\le k_1(k)$ are
among those of clause~(d), since $k_1(k)<k$; and clause~(b) gives $|\zeta_k-\zeta|\le\frac13|\zeta|$
for $\zeta\in D(0,2r)$, whence
\begin{equation*}
|\zeta_k|\le|\zeta|+|\zeta_k-\zeta|\le\tfrac43|\zeta|\le\tfrac83\,r .
\end{equation*}
With these data the frozen-integral proposition produces the family $E^{(k+1)}_{\zeta}$ on the
radii $\hat\alpha_k$. The valuation statement is then applied with
\begin{equation*}
\mathsf{q}_{k+1}:=\beta\bigl[Q^{(k+1)}_{\zeta}\bigr],
\end{equation*}
which satisfies $|\mathsf{q}_{k+1}|\le v_k$ on $D(0,2r)$.

The remaining clauses of $(\mathrm{J}_{k+1})$ are established in the steps that follow.
\end{proof}

\medskip
\noindent\textit{Joint induction: the frozen averaging step.}
We continue the passage from level $k$ to level $k+1$, $k<K$, of the joint induction
(see the part on the pure step and the frozen integral above), and now
treat the operation $\mathbf{T}$ of this step together
with its domains, and the renormalisation that follows it.

The operation $\mathbf{T}$ is applied in Ba{\l}aban's form, through the decompositions
\cite[(3.2), (3.3), (3.16)]{B-CMP119}, none of which involves the sources $z$. These
decompositions are followed by the exact refinement of the averaging step, and the step is then
taken with its parameter $a_k$.

At the sites of the operation $\mathbf{T}$ four lemmas, used together with the budget inequality
for the constant $C_{\sharp}$ of this chapter, give the three output families at level $k+1$. The
four lemmas are:
\begin{itemize}
\item the amplitude lemma;
\item the complex-parameter lemma;
\item the summation lemma;
\item the flat summation lemma.
\end{itemize}
The three families are $E^{(k+1)}$, $R''^{(k+1)}$ and $\mathbf{B}^{(k+1)}$. The last is kept
as a family of stored terms
\begin{equation}\label{5.2:eq-frozen-averaging-1}
  \sum_{\mathcal{S}} \mathbf{B}^{(k+1)}(X,\mathcal{S}),
\end{equation}
which carry the datum $A^{\flat}$ and satisfy the weighted bound of the flat boundary-term
corollary. All three families are defined on Ba{\l}aban's enlarged
spaces $\tilde{U}$, and they decay at the rate $\mathsf{r}(\delta)\kappa$, where
\begin{equation}\label{5.2:eq-frozen-averaging-2}
  \mathsf{r}(\delta)\kappa \;\ge\; (1+4\beta_s)\kappa \;=\; 2\kappa .
\end{equation}

The rate in \eqref{5.2:eq-frozen-averaging-2} is supplied by the threshold lemma for the decay
fraction $\delta$. That lemma reads \cite[(2.41)]{B-CMP116} at $\delta\le\delta_J(L)$. The
inequality in \eqref{5.2:eq-frozen-averaging-2} holds for every
$\delta\in\,]0,\delta_J(L)]$. Indeed,
$\delta\le\delta_J(L)=(1-4/L)/10$ gives $1-10\delta\ge 4/L$. Hence
\begin{equation*}
  \mathsf{r}(\delta)=\frac{(1-10\delta)L}{2}\ge 2 .
\end{equation*}
Moreover $\beta_s=1/4$, so that $(1+4\beta_s)\kappa=2\kappa$. The members of the one-loop
family $L^{(k+1)}$ enter with the decay rate $\delta_0 M$, and $\delta_0 M\ge 2\kappa$.

Let $\Lambda_{k+1}$ be the region so denoted at level $k+1$ in the decompositions
\cite[(3.2), (3.3), (3.16)]{B-CMP119}. The terms localised in domains $X\subset\Lambda_{k+1}$
``do not depend on
$\Lambda_{k+1}$, and coincide with the corresponding terms arising from the expressions
defined on the whole lattice'' \cite{B-CMP119}. These terms are therefore the whole-lattice
terms already constructed in the frozen-integral part of the induction step
above.

The vacuum values of the output families are collected into the constant $e^{(k+1)}$. The
marginal part is then removed as in \cite{B-CMP119}, ``subtracting
$\beta_{k+1}(g_k)A(\varphi_{k+1},U_{k+1})$'', with $\beta_{k+1}(g_k)$ replaced by the complex
coupling increment $b_{k+1}(\zeta)$ of the flow step
(the part on the coupling increment and the reference coupling, below).
Here $A(\varphi_{k+1},U_{k+1})$ is the action of
the level-$(k+1)$ configuration in Ba{\l}aban's notation.
After the subtraction, the weight multiplying the action is
\begin{equation}\label{5.2:eq-frozen-averaging-3}
  x_{k+1}(\cdot)-\zeta_{k+1},
\end{equation}
with $x_{k+1}$ and the running shift $\zeta_{k+1}$ as in the flow step.

\begin{remark}[Joint induction: the large-field step]\label{5.2:prf-large-field-step}
We apply the $\mathbf{R}$-operation at level $k+1$, in both its frozen and its local readings,
to the density obtained from the averaging step. Specifically, we apply at level $k+1$:

\begin{itemize}
\item the proposition describing the $\mathbf{R}$-operation at complex sources, of which clauses
(a) and (e) are used below;
\item the lemma that accompanies that proposition.
\end{itemize}

The averaging step in question is either the frozen averaging step or the local averaging step under sources
(\ref{5.3:prf-local-averaging}).

Clause (e) of that proposition divides the terms produced by the $\mathbf{R}$-operation into two
groups: the first group consists of the $\mathbf{R}$-born boundary terms $\mathbf{B}'^{(k+1)}$,
and the second group of the remaining terms produced by the $\mathbf{R}$-operation.

\emph{The first group.}
Clause (e) returns the first group, the terms $\mathbf{B}'^{(k+1)}$, into $\mathbf{B}^w$. Their
bound is the relative form of clause $(\mathrm{d})'(\alpha)$ of the induction format
(Definition~\ref{5.2:def-induction-format}), that is, \eqref{5.2:eq-induction-format-b}.
This is exactly the bound that clause (e) displays. This relative bound is supplied
by the relative complex bound for the large-field operations at level $k+1$; the propagation
lemma of the present induction derives it, through the five propositions on which that lemma
rests, from that bound at level $k$ and the joint induction
statement at level $k$. Hence the format is reproduced at level $k+1$ exactly as written, in
its relative form, that is, with the length in \cite[(2.42)]{B-CMP119} read as the relative size
of \cite[(1.67)]{B-CMP122b}, and not with the plain length of \cite[(2.42)]{B-CMP119}. Wherever
Ba{\l}aban's bounds are invoked in this step, they are his printed bounds, with the single change
that in \cite[(2.42)]{B-CMP119} the length is measured by the relative size of
\cite[(1.67)]{B-CMP122b}.

\emph{The second group.}
Clause (e) returns the second group into $\mathbf{R}^w$, bounded in the form of
\cite[(2.30)--(2.32)]{B-CMP119}.

\emph{Orbit data.}
Every new term of $\mathbf{B}$-type is filed together with its orbit datum, as clause
$(\mathrm{d})'(\gamma)$ of the format requires. In the two averaging steps this datum was supplied by
clause (d) of the corollary on boundary terms at complex sources. In the present step it is
supplied as follows:
\begin{itemize}
\item at the preparatory stages, by the corresponding corollary for the preparatory boundary
terms at complex sources;
\item at the localisation site, by clause (b) of the proposition on the bracket terms and its
corollary at complex sources.
\end{itemize}

\emph{Arrests.}
Clause (a) of the proposition on the $\mathbf{R}$-operation, together with clause
$(\mathrm{d})'(\beta)$ of the format, carries the arrested left-overs of every live arrest unchanged
from level $k$ to level $k+1$. An arrest $\mathfrak{a}$ may also be made at the present step,
because a chain is finished early or is excluded. Its frozen terms $\mathfrak{F}(\mathfrak{a})$
are then filed into clause $(\mathrm{d})'(\beta)$ of the format. This filing is made by the filing rule
for arrests, read under the standing domain convention.

\emph{Conclusion.}
All terms produced by the $\mathbf{R}$-operation at level $k+1$ have now been placed:
\begin{itemize}
\item the second group in $\mathbf{R}^w$, with the bound of \cite[(2.30)--(2.32)]{B-CMP119};
\item the first group in $\mathbf{B}^w$, with the relative bound
\eqref{5.2:eq-induction-format-b} and with its orbit datum;
\item the arrested left-overs, old and new, in clause $(\mathrm{d})'(\beta)$.
\end{itemize}
The format of Definition~\ref{5.2:def-induction-format} therefore holds at level $k+1$. With it
the joint induction statement holds at level $k+1$, which completes the induction step.
\end{remark}

\begin{proposition}[Order of the joint induction and its conclusion]
\label{5.2:prf-induction-order}
Let $(J_k)$, $0\le k\le K$, be the joint hypothesis of
Definition~\ref{5.2:def-joint-hypothesis}. Consider the following statements:
\begin{enumerate}
\item the correlation theorem, whose inductive step is the passage from $(J_k)$ to $(J_{k+1})$
carried out in the pure step and the frozen integral, the coupling
increment and the reference coupling (\ref{5.3:prf-flow-reference}), the frozen averaging step, the local averaging step under sources
(\ref{5.3:prf-local-averaging}) and the large-field step (\ref{5.2:prf-large-field-step}) of the
joint induction;
\item the collar-payment statement of the large-field step;
\item the relative bound for the chain terms $\mathbb{C}^{(n)}_k$;
\item the closing lemma at the localisation site, together with its version in the shift
variable $\zeta$;
\item the tube statement, in its form part (the tube-form proposition) and in its value part, the
two-point bound supplied by its base theorem;
\item the group bound at the preparatory site;
\item the lemma on plain units inside the large-field region;
\item Lemma~\ref{5.2:lem-exterior-pieces}, on the pieces of units whose label leaves the large-field
region.
\end{enumerate}
Assume, as hypotheses, the inductive steps of these statements: for each of them, at each step $k$
and each site of that step, the statement at that site is derived from the inputs named
where that statement is proved; these inputs are recalled, position by position, in the proof below.
We write $(g)_{k-1}$ and $(g)_k$ for the positions at which step $k$ opens and closes,
$(\mathrm{T})$ for the site of the averaging steps, $\mathcal{I}^{\mathbb{C}}_k$ for the position
of the complex inductive bound at scale $k$, that is, of the inductive hypothesis of the relative
bound for the chain terms at step $k$, $(\mathrm{P})$ and
$(\mathrm{L})$ for the preparatory and the localisation site of the large-field step, and $(c)$ for
the position, after $(\mathrm{P})$, that closes the part of step $k$ read by the closing lemma.
Then these statements hold together, by one induction on pairs $(\text{step},\text{site})$, in which
step $k$ runs through the positions in the order
\begin{equation}\label{5.2:eq-induction-order-a}
(g)_{k-1}\;\to\;(\mathrm{T})\;\to\;\mathcal{I}^{\mathbb{C}}_k\;\to\;(\mathrm{P})\;\to\;(c)
\;\to\;(\mathrm{L})\;\to\;(g)_k .
\end{equation}
Each item uses only items preceding it in \eqref{5.2:eq-induction-order-a} at the same step and
items of earlier steps. In particular the induction contains no circle. The induction ends with
$(J_K)$, and $(J_K)$ implies the reference-coupling clause, the running-shift statement and the
clause on $N_K$ of the correlation theorem.
\end{proposition}

\begin{proof}
The positions in \eqref{5.2:eq-induction-order-a} are those named in the statement. The site
$(\mathrm{T})$ is the site of the averaging steps, frozen (the frozen averaging step) and local (the local averaging step under sources,
\ref{5.3:prf-local-averaging}). The sites $(\mathrm{P})$ and $(\mathrm{L})$ are the two sites of the
large-field step (\ref{5.2:prf-large-field-step}). We list, position by position, what each item
uses, and check that nothing is used before it is established.

\emph{The flow and the remainders.} The running-shift statement at levels $\le j$ is used at the
large-field step of level $j$ to produce $R^{(j)}_\zeta$. The remainder $R^{(j)}_\zeta$ enters the
running-shift statement at level $k+1$ only for $k_1(k)\ge j$, through the frozen integral of the
pure step. Thus the flow at level $k+1$ uses remainders of levels already
completed, and there is no circle.

\emph{Internal orders.} The order \eqref{5.2:eq-induction-order-a} is the order in which the
collar-payment statement lists its items, and the ordering lemma of that statement shows that this
list contains no forward pointer. The items of the relative bound for the chain terms are used in
that bound's own ordered list; the first of them includes the tube statement at the site where the
large-field component is created.

\emph{The site $(\mathrm{P})$.} At $(\mathrm{P})$ the part of the correlation theorem attached to
this site uses the following inputs:
\begin{itemize}
\item the field-free scaffold and counting lemmas of the relative bound for the chain terms, namely
the flat lemma, the resummation lemma of the preparatory stage with its five clauses, the localised
star lemma, a further counting lemma, the window lemma and the anchoring lemma for all preparatory
terms;
\item the preparatory-site corollary of the collar-payment statement, with its accompanying bound;
\item the group bound at the preparatory site, for the groups it concerns.
\end{itemize}
The scaffold and counting lemmas are used in one direction only: no statement of the correlation
theorem is an input of theirs. In the other direction, the relative bound for the chain terms uses,
among the parts of the correlation theorem, only those lettered (C) and the small-field parts. It
does not use the conclusion reached at $(\mathrm{P})$.

\emph{The site $(\mathrm{L})$.} The closing lemma uses the items established up to position $(c)$ of
\eqref{5.2:eq-induction-order-a}. At $(\mathrm{L})$ of step $k$ its inputs are:
\begin{itemize}
\item clause (a) of $(J_k)$, i.e.\ the reference-coupling clause for pairs $\le k$;
\item the sourced inductive hypothesis for the chain terms at level $k$;
\item the containment statement;
\item the constant $\sigma_0$ of the relative bound for the chain terms;
\item clause (g) of the geometry, which follows from the statement on nested regions;
\item the two statements on the sources $w$;
\item two further lemmas established before position $(c)$ of \eqref{5.2:eq-induction-order-a}.
\end{itemize}
It uses no length, label or sum of step $k$, no value bound of the tube statement, no
$\mathbf{R}'^{(k)}$ and no $\alpha^{\sharp}(k)$. The off-diagonal rider and the half-power statement
use the same inputs, together with the supplementary statement on clause (a).

\emph{The plain-unit lemma.} The lemma on plain units inside the large-field region uses the closing
lemma and field-free counting, and returns its member of the pre-filing family.

\emph{The exterior lemma: inputs.} Lemma~\ref{5.2:lem-exterior-pieces} uses:
\begin{itemize}
\item the two statements on the sources;
\item clause (f) of the geometry;
\item the sum, which carries no field;
\item the sizes, and only the sizes, supplied by the filing step;
\item the stored bound of the amplitude lemma;
\item through one of the displayed bounds of the lemma on plain units, the containment statement
read at the auxiliary consumer term $\mathfrak{E}^{\oplus}_v=\mathbb{C}^{(N)}_k(Y(\mathrm{lab}_1))$
inside $Y_4$;
\item the inheritance statement, its corollary and its claim, in clause ($\beta$);
\item the arrest statements (0)--(2) and (5) on foreign live pieces;
\item the bounds of the off-diagonal rider in the case $b=0$ of its index $b\ge 0$, for the compact
pointed pieces only;
\item clause (ii)($\alpha$) of the collar-payment statement at the site $(\mathrm{L})$;
\item the bound \eqref{5.2:eq-plain-fine-units-a} on plain units finer than a cell;
\item the third of the auxiliary bounds of this chapter;
\item the recursion lemma for the correlation theorem.
\end{itemize}
The long pointed pieces, in the two domain classes of the correlation theorem where they occur, use
no item of the rider. For them the long-piece bound at $\mathsf{t}=0$ and
the two inequalities for $\theta'$ are used instead.

\emph{The exterior lemma: the recursion lemma and the output.} The recursion lemma uses:
\begin{itemize}
\item the layer objects $\mathbb{H}_s$ and $\mathcal{J}^{\sharp}_s[\chi]$, with the input bounds of
each layer $s$ and their localisation under sources, all of which lie upstream of every
part of the correlation theorem;
\item the supplementary statement on clause (a), read at clause (a) of $(J_k)$;
\item the arrest geometry;
\item the geometry and the cell structure, which carry no field.
\end{itemize}
It does not use the base theorem of the two-point bound. That theorem rests, through its part (D),
on the exterior carrier bound, so its use here would be a circle. Neither
Lemma~\ref{5.2:lem-exterior-pieces} nor the recursion lemma uses any length, value bound,
$\mathbf{R}'^{(k)}$ or $\alpha^{\sharp}(k)$ of step $k$. Lemma~\ref{5.2:lem-exterior-pieces} returns
$\varepsilon^{\mathrm{ext}}$.

\emph{The operations $\mathbf{R}$ and $\mathbf{B}$.} Clauses (c)--(e) of the proposition on the
operation $\mathbf{R}$ use the pre-filing family with its six members. The proposition on the terms
$\mathbf{B}$ returns the newly born terms of the relative bound for the chain terms at step $k+1$.

\emph{The tube statement: its two parts.} Its form part is the tube-form proposition, whose input is
a theorem that is not among the statements listed in the proposition. Its value part is the
two-point bound supplied by its base theorem. The
base theorem has its own position within the site $(\mathrm{L})$ of step $k$, after the
lists compiled at that site. It uses:
\begin{itemize}
\item the splitting, rim and totalisation statements and the carrier-pair statement;
\item the pointed version in $\zeta$ of the closing lemma;
\item the Wilson-density rider;
\item the nesting statement, which is obtained from its ingredients (A)--(H), with the statement on
the terms $\mathbf{C}_k$ entering at its line (i).
\end{itemize}

\emph{The tube statement: where it is used and paid.} Both parts are used at the localisation site
where the component is created, and at every later site $(\mathrm{T})$, $(\mathrm{P})$,
$(\mathrm{L})$ of the component's life. At every later site the value is paid in cells, measured in
$d^{\wedge}$, through the collar belonging to the region $(Y,m)$ that carries it, by the
collar-payment statement.

\emph{Conclusion of the first part.} In every case above, an item of step $k$ uses only items
preceding it in \eqref{5.2:eq-induction-order-a} and items of steps $<k$. The pairs
$(\text{step},\text{site})$ ordered lexicographically, with the sites of a step ordered by
\eqref{5.2:eq-induction-order-a}, are well ordered. The inductive step of the correlation theorem,
the passage from $(J_k)$ to $(J_{k+1})$ carried out in the statements listed in the first item of
the proposition, is placed
in the same order as follows: the averaging steps occupy $(\mathrm{T})$ and the large-field step
occupies $(\mathrm{P})$ and $(\mathrm{L})$, while the pure step, the frozen integral and the coupling
increment use only remainders of levels already completed, as shown in the paragraph on the flow.
Therefore the induction
closes, every statement of the list holds at every step $k\le K$, and the induction ends with
$(J_K)$. From $(J_K)$ we now deduce the three clauses of the correlation theorem.

\emph{The reference-coupling clause.} Clause (a) of $(J_k)$, valid at every step $k$ reached by the
induction, makes the first passage $K$ finite. The inequality
$x_{K+1}\le X$ gives $x_K\le X+b_+$. For $\varrho\in[r,2r]$ the sum $\sum_{k<K}d_k(\varrho)$ is
estimated by the summability condition on the flow constants $d_k(\varrho)$, applied with the lower
bound
\begin{equation}\label{5.2:eq-induction-order-1}
x_k\;\ge\;x_K+\lambda(K-k)-2c_2,\qquad 0\le k<K .
\end{equation}
These facts, together with clause (a) of $(J_K)$, give the reference-coupling clause.

\emph{The running-shift statement and the clause on $N_K$: the Cauchy estimate.} For $0\le k<K$ put
$h_k:=b_{k+1}-b_{k+1}(0)$. By clause (b) of $(J_K)$, applied with $j=k+1\le K$, the function $h_k$ is
analytic on $D(0,2r)$, real on reals, and $h_k(0)=0$. On $D(0,2r)$ it satisfies
\begin{equation*}
|h_k(\zeta)|\;\le\;d_k(2r)\,|\zeta|\;\le\;2r\,d_k(2r).
\end{equation*}
Let $\zeta\in\bar D(0,r)$ and $0<r'<r$. The circle of radius $r'$ about $\zeta$ lies in $D(0,2r)$, so
Cauchy's estimate gives $|h_k'(\zeta)|\le 2r\,d_k(2r)/r'$. Letting $r'\to r$ yields
\begin{equation}\label{5.2:eq-induction-order-2}
|h_k'(\zeta)|\;\le\;2\,d_k(2r),\qquad \zeta\in\bar D(0,r).
\end{equation}
Clause (b) of $(J_K)$ together with \eqref{5.2:eq-induction-order-2} gives the running-shift
statement.

\emph{The running-shift statement and the clause on $N_K$: the Schwarz lemma.} Since $h_k(0)=0$ and
$|h_k|\le 2r\,d_k(2r)$ on $D(0,2r)$, the Schwarz lemma gives $|h_k'(0)|\le d_k(2r)$. The quantity
$N_K$ of the correlation theorem satisfies
\begin{equation}\label{5.2:eq-induction-order-3}
N_K-1\;=\;\sum_{0\le k<K}h_k'(0).
\end{equation}
Each $h_k$ is real on reals, so each $h_k'(0)$ is real. Hence $N_K$ is real, and by the triangle
inequality
\begin{equation*}
|N_K-1|\;\le\;\sum_{0\le k<K}|h_k'(0)|\;\le\;\sum_{0\le k<K}d_k(2r).
\end{equation*}
The last sum is the sum estimated above at $\varrho=2r$. The reality of $N_K$ and the last display
give the clause on $N_K$.
\end{proof}

\begin{lemma}[Admissible decay fractions and the doubled output rate]\label{5.2:lem-decay-fraction}
Let $\delta\in\left]0,\delta_J(L)\right]$, where
\begin{equation}\label{5.2:eq-decay-fraction-a}
\delta_J(L):=\frac{1}{10}\Bigl(1-\frac{4}{L}\Bigr),\qquad
\mathsf{r}(\delta):=(1-10\delta)\,\frac{L}{2},\qquad
\rho_m:=(1-m\delta)\,\frac{L}{2}\,\kappa\quad(m=1,\dots,10),
\end{equation}
so that $\mathsf{r}$ is decreasing in $\delta$ and $\mathsf{r}(\delta)\ge\mathsf{r}(\delta_J(L))=2$
throughout this window.
Here $\delta_0$, $\alpha_6$, $C_3$, $\varepsilon_1$, $\varepsilon_2$ denote the constants of the
cluster expansion of \cite{B-CMP116}, $\beta_s=\tfrac14$ is the
enlargement fraction of the families of this chapter, and $d_{k+1}(X)$ is the linear size of a
localization domain $X$ at level $k+1$. Then the following hold.
\begin{enumerate}
\item[(i)] Every statement of this chapter asserting, for an output of the large-field
operation $\mathbf{T}$, decay with the factor $e^{-\kappa d_{k+1}(X)}$ holds.
\item[(ii)] Every output kernel of the large-field operation $\mathbf{T}$ at level $k+1$,
namely $E^{(k+1)}$ without its members in $L^{(k+1)}$, the pinned pieces, $R''^{(k+1)}$,
the term $\mathbf{B}^{(k+1)}$ evaluated at $A_k=0$, where $A_k$ is the field whose restriction
to the rim bonds is the rim datum, and the filed rim terms
$\mathfrak{K}(X,\mathcal{S})$, obeys its bound with the factor
\begin{equation*}
e^{-\mathsf{r}(\delta)\kappa d_{k+1}(X)}\;\le\;e^{-(1+4\beta_s)\kappa d_{k+1}(X)} .
\end{equation*}
This bound involves the interaction $V$ of the step in \cite{B-CMP116} only through
$\varepsilon_2$, and it is
uniform in the complex source $w$. The members of $L^{(k+1)}$ carry the rate $\delta_0M$, and
$\delta_0M\ge 2\kappa$. Consequently the following requirements of this chapter are met: the
output rate $(1+4\beta_s)\kappa$ stated for the large-field operation in the inductive
statement of the step; the lower bound $\kappa_{\mathrm{out}}\ge(1+4\beta_s)\kappa$ on the
output rate, both in the general step and in the repair of the creating step on its flat
support; the clause of the resummation corollary for the chain terms at its index $\ell=k$,
which asks that the outputs of the large-field operation $\mathbf{T}$ of the current step,
that is, the kernels at level $k+1$ listed above, decay at a rate, denoted $\kappa_T$ in that
corollary, at least $(1+4\beta)\kappa$, where $\beta=\beta_{\mathrm{pr}}=\tfrac15$ is the
enlargement fraction of that corollary, so that $(1+4\beta)\kappa=\tfrac95\kappa\le2\kappa$,
together with the rate condition at $\kappa_R=\kappa$ on the chain terms that this
clause feeds;
the case $\nu=1$ of the first clause and the case $j=k$ of the second clause
of the boundary-term lemma; the sentence of the per-cell budget of the constant
$C_{\sharp}$ on the margin $(\tfrac14-\tfrac14\delta)\kappa d_k$ at $j=k$;
the case $j=k$ of the third remark on the boundary terms; and the margin $\tfrac14$ at $j=k$
in the count of near points.
\item[(iii)] The only conditions of this chapter that change when $\delta$ is taken in the
window \eqref{5.2:eq-decay-fraction-a} rather than at the value fixed in \cite{B-CMP116} are
the following. First, the thresholds, denoted $\nu_0$, of the holomorphic cluster expansion.
Second, the
members of the smallness conditions of the E-term lemma and of $\varepsilon_V^{*}$ that carry
the factors $e^{5\rho_7}$ and $e^{5\rho_9}$. Through these members the constants $c_{41}$,
$b^{\mathrm{adm}}$, $K_{CE}$ and
\begin{equation*}
E^{*}:=O(1)\,\alpha_6^{-1}C_3\varepsilon_1
\end{equation*}
depend on $\delta$. These constants are paid for by $\varepsilon_1$ and by the member
$K_{CE}C_U\theta(1+S_I)\le\tfrac12$ of the list $\gamma_J$, where $\theta=\theta(g)$ is the
ratio of shift to radius. The powers of $\alpha_6$ are
counted as $\alpha_6^{-2}$ on the members of the $\varepsilon_2$-lemma that come from
\cite[(2.37)]{B-CMP116} and on $\varepsilon_V^{*}$. The quantities $\varepsilon_V^{*}$, $c_V$,
$c_A$, $\varrho$, $r_A$, $T^{\mathbb{C}}$ and $\varpi_\gamma$ enter as definitions. Third, the
lower bound $\delta_0M\ge2\kappa$. Fourth, the lower bound $\kappa\ge C_\delta(d)/\delta$, with
$C_\delta(d)$ depending on $d$ only. No new upper bound on $\gamma$ arises, and no upper bound
on $L$, $\kappa$, $\kappa_1$, $M_1$ or $M$ arises.
\end{enumerate}
\end{lemma}

\begin{proof}
We prove (ii) first, then (i), then (iii).

\emph{Proof of} (ii). The cluster expansion of \cite{B-CMP116}, from
\cite[(1.8)]{B-CMP116} to \cite[(2.41)]{B-CMP116}, is carried out for an arbitrary
$\delta\in\left]0,1\right[$. It is subject only to lower bounds on $\kappa$ and $\kappa_1$ and
to the upper bounds on $\varepsilon_2$ collected in the list of thresholds of the analytic step
of this chapter. We go through the places where $\delta$ enters, in the order in which they
occur there.

First, at \cite[(1.32)]{B-CMP116} the parameter $\delta$ is introduced with the sole requirement
$0<\delta<1$.

Second, several estimates are asserted ``for $\delta\kappa$ sufficiently large''
\cite{B-CMP116}. For fixed $\delta$ these are lower bounds on $\kappa$. They are collected in
the lower bound $\kappa\ge C_\delta(d)/\delta$ in the proof of (iii) below.

Third, \cite{B-CMP116} uses two conditions on $\kappa_1$. One requires that the multiple of
$\kappa_1$ printed there dominate $(1-2\delta)\kappa$. The other is
\begin{equation*}
\tfrac18(\kappa_1-1)\ge(1-3\delta)\kappa ,
\end{equation*}
imposed together with a condition $\bar q\ge8$ on an integer of \cite{B-CMP116}, written $q$
there and denoted $\bar q$ here. For every $\delta>0$, both inequalities on $\kappa_1$ are
implied by the
lower bound $\kappa_1\ge16\kappa+1$, which belongs to the choice of parameters of the correlation
theorem; the condition $\bar q\ge8$ is not a condition on $\kappa_1$ and does not involve
$\delta$. For the second inequality this is the chain
\begin{equation*}
\tfrac18(\kappa_1-1)\ge2\kappa\ge(1-3\delta)\kappa .
\end{equation*}

Fourth, \cite{B-CMP116} states that ``the last exponential multiplied by
$\exp(-\frac12\delta L\kappa d_{k+1}(X))$ is bounded by $1$''. For the $\delta$ fixed in the
statement, it suffices that
\begin{equation}\label{5.2:eq-decay-fraction-1}
e^{O(1)(LM)^{-4}|X|}\le e^{\frac12\delta L\kappa\,d_{k+1}(X)} .
\end{equation}
It follows from \cite[(2.30)]{B-CMP116}, which reads $(3\cdot2^3)^{-1}M^{-4}|Y|\le d_k(Y)$,
applied at level $k+1$. By that inequality we have $(LM)^{-4}|X|\le 24\,d_{k+1}(X)$, so
\eqref{5.2:eq-decay-fraction-1} holds as soon as $\tfrac12\delta L\kappa\ge24\cdot O(1)$. This
is the lower bound
\begin{equation*}
\kappa\ge\frac{48\cdot O(1)}{L\delta}
\end{equation*}
on $\kappa$, and it is contained in the lower bound $\kappa\ge C_\delta(d)/\delta$.

The inequality \cite[(2.30)]{B-CMP116} concerns domains with $d_{k+1}(X)\ge1$; for a non-empty
domain with $d_{k+1}(X)=0$ its right member vanishes while its left member is positive. A domain
consisting of at most $16+8d$ cubes sharing a vertex has size $0$ (see \cite{B-CMP109}), and
\cite{B-CMP116} states: ``we sum over $X$ with $d_j(X)\neq0$''. We therefore use the inequality
\begin{equation*}
(LM)^{-4}|X|\le 24\,d_{k+1}(X)+16+8d ,
\end{equation*}
which holds for every $X$. Compared with \eqref{5.2:eq-decay-fraction-1}, it leaves the bounded
excess factor $e^{(16+8d)O(1)}$. This factor is absorbed into the unnamed constant $O(1)$ of
\cite[(2.41)]{B-CMP116}, and the lower bound on $\kappa$ is unchanged.

Fifth, the remaining lower bound on $\kappa_1$ imposed in this part of \cite{B-CMP116} does not
involve $\delta$.

Sixth, $\delta$ is fixed to a particular value only after \cite[(2.41)]{B-CMP116}, which reads
\begin{equation}\label{5.2:eq-decay-fraction-2}
|E^{(k+1)}(X)|\le O(1)\,C_3\varepsilon_1
\exp\Bigl(-(1-10\delta)\tfrac12 L\kappa\, d_{k+1}(X)\Bigr).
\end{equation}

Next we pass to complex parameters. In the holomorphic cluster expansion, which carries
Ba{\l}aban's bounds over to a complex parameter $\varsigma$, $\delta$ is kept symbolic
in its second and third clauses. The particular value of $\delta$ is used there exactly once,
in the last equality of
\begin{equation}\label{5.2:eq-decay-fraction-3}
|t_\varsigma(X)|\le c_{41}C_3(\varsigma)\varepsilon_1
e^{-(1-10\delta)\frac12L\kappa\, d_{k+1}(X)}
= b(\varepsilon_2(\varsigma))\,e^{-\kappa d_{k+1}(X)} .
\end{equation}
Here $t_\varsigma(X)$ denotes the output kernel at $X$ and at the parameter $\varsigma$,
$C_3(\varsigma)$ and $\varepsilon_2(\varsigma)$ denote the constants $C_3$ and $\varepsilon_2$ of
\cite{B-CMP116} formed with the data at $\varsigma$, and $b(\cdot)$ is the function of
$\varepsilon_2$ written there. The third clause of that expansion covers all the
output kernels listed in (ii). Moreover, by the $\varepsilon_2$-lemma it depends on the
interaction $V$ only through $\varepsilon_2$. The same holds for the three bounds of the
companion statement on the expansion of the step, in which the value of $\delta$ is used at two
places only.

We now drop the last equality in \eqref{5.2:eq-decay-fraction-3} and keep $\delta$ in the
window of the statement. The smallness conditions of the form $a\,e^{5\rho}\le1$, where $a$ is
the common prefactor of the factors $a\,e^{-\rho d}$ that are merged by
\cite[(2.27)]{B-CMP116}, are then imposed at the rate $\rho$ actually in play, as prescribed in
the list of thresholds of the analytic step. That is, at the two places where \cite{B-CMP116}
writes the factor $\exp 5\kappa$ in these conditions, this chapter reads them with the factor
$e^{5\rho_7}$ at the first place and $e^{5\rho_9}$ at the second, the rates in play there being
$\rho_7$ and $\rho_9$. Since $m\delta\ge0$ we have $\rho_m\le\tfrac12L\kappa$
for every $m$, hence
\begin{equation*}
e^{5\rho_m}\le e^{\frac52L\kappa},
\end{equation*}
uniformly in $\delta$.

With these readings, \eqref{5.2:eq-decay-fraction-2} and the corresponding bounds for the other
output kernels deliver the rate
\begin{equation*}
(1-10\delta)\tfrac12L\kappa=\rho_{10}=\mathsf{r}(\delta)\kappa .
\end{equation*}
The function $\mathsf{r}$ is decreasing in $\delta$, and from
\eqref{5.2:eq-decay-fraction-a} we get $10\delta_J(L)=1-4/L$. Hence, for
$\delta\le\delta_J(L)$,
\begin{equation*}
\mathsf{r}(\delta)\kappa\ge\mathsf{r}(\delta_J(L))\kappa
=\frac{4}{L}\cdot\frac{L}{2}\,\kappa=2\kappa=(1+4\beta_s)\kappa .
\end{equation*}
This is the inequality displayed in (ii). It depends on $V$ only through $\varepsilon_2$.

The repair of the creating step on its flat support runs, in its fourth step, through
\cite[(2.27)--(2.30)]{B-CMP116} and \cite[(2.35)--(2.41)]{B-CMP116}. The only change there is
that the flat analogue of \cite[(2.26)]{B-CMP116} is used. The decay condition of that repair
is read at a rate, denoted $\bar\kappa$ there, with $\bar\kappa\le\tfrac12L\kappa$. The same
argument therefore gives the same rate for the rim
terms $\mathfrak{K}(X,\mathcal{S})$ produced there.

For the members of $L^{(k+1)}$, \cite{B-CMP116} states that $\kappa$ is ``replaced by
$\delta_0M$'' with $\delta_0M\ge\kappa$. Here we need the stronger bound $\delta_0M\ge2\kappa$.
It follows from the condition
\begin{equation*}
e^{-\frac13\delta_0M}e^{16\kappa_1}<1
\end{equation*}
of \cite{B-CMP116}, which says $\delta_0M>48\kappa_1$. Combined with the lower bound
$\kappa_1\ge4L\kappa+1$ in the choice of parameters of the correlation theorem, it gives
\begin{equation*}
\delta_0M>48\kappa_1\ge48(4L\kappa+1)>2\kappa .
\end{equation*}

The sentence of \cite{B-CMP119} in which $\kappa$ is ``replaced, for example, by
$(1+4\beta)\kappa$'' concerns the same rate: with $\delta$ in the window of the statement, the
argument above gives the rate $2\kappa\ge(1+4\beta)\kappa$, where $\beta=\beta_{\mathrm{pr}}$.
That sentence is not used as an input anywhere in this chapter.

Each of the requirements listed at the end of (ii) asks that the outputs of $\mathbf{T}$ at
level $k+1$ decay at rate at least $(1+4\beta_s)\kappa=2\kappa$ (for the clause of the
resummation corollary, at least $(1+4\beta)\kappa=\tfrac95\kappa$, which does not exceed
$(1+4\beta_s)\kappa=2\kappa$), and that the members of
$L^{(k+1)}$ decay at rate at least $2\kappa$. Both have been established, so these requirements
are met. This proves (ii).

\emph{Proof of} (i). By (ii), the outputs of $\mathbf{T}$ decay at rate $\mathsf{r}(\delta)\kappa$
with $\mathsf{r}(\delta)\ge2\ge1$, and the members of $L^{(k+1)}$ decay at rate
$\delta_0M\ge2\kappa\ge\kappa$. Since $d_{k+1}(X)\ge0$, this gives
\begin{equation*}
e^{-\mathsf{r}(\delta)\kappa d_{k+1}(X)}\le e^{-\kappa d_{k+1}(X)},\qquad
e^{-\delta_0M d_{k+1}(X)}\le e^{-\kappa d_{k+1}(X)} .
\end{equation*}
Every statement made at rate $e^{-\kappa d_{k+1}(X)}$ for such an output therefore holds.

\emph{Proof of} (iii). We first list the conditions of this chapter that read $\delta$ directly
but do not change when $\delta$ is taken in the window:
\begin{itemize}
\item the conditions in which the rates $(1-2\delta)\kappa$ and $(1-\tfrac12\delta)\kappa$, and
the rate $\kappa_V:=(1-2\delta)\kappa$, of the lemma on the Cauchy radius $\rho_A$ enter;
\item the requirement $\delta<\tfrac13$ of the boundary-term lemma;
\item the requirement that the margins in the count of near points be at least $\tfrac12\delta$,
and the inequality $\tfrac14\ge\tfrac12\delta$;
\item the requirement $\delta\le\tfrac12$.
\end{itemize}
Each of these conditions is monotone in $\delta$: it is an upper bound on $\delta$. Each holds
for all $\delta<1/10$. Since $\delta\le\delta_J(L)<1/10$, they all hold for every admissible
$\delta$, and none of them changes.

The remaining conditions that read $\delta$ are the following: the requirements ``for
$\delta\kappa$ sufficiently large''; the constants $C(\kappa)$ and $C_B$, taken at the rate
$\tfrac14\delta\kappa$, of the kind appearing at \cite[(1.26)]{B-CMP116}; the inequality
\eqref{5.2:eq-decay-fraction-1}; and the estimate in the filing of the units in which the
retained factor $e^{-\frac14\delta\kappa d_j}$ pays for the anchor counts. Each of these becomes
a lower bound of the form $\kappa\ge C_\delta(d)/\delta$, which depends only on $d$ once $\delta$
is fixed. At $\delta=\delta_J(L)$ we have
\begin{equation*}
\frac{1}{\delta_J(L)}=\frac{10}{1-4/L}\le20\qquad\text{for } L\ge8,
\end{equation*}
so the lower bound is then at most $20\,C_\delta(d)$.

The window is non-empty exactly when $\delta_J(L)>0$, that is, when $L>4$. This is contained in
the lower bound $L\ge8$ in force in this chapter. The lower bound $\delta_0M\ge2\kappa$ was
obtained in the proof of (ii). The thresholds named in (iii) are those of the list of thresholds
of the analytic step, where they are read at $e^{5\rho_7}$ and $e^{5\rho_9}$ as explained in the
proof of (ii).

Altogether, the conditions that read $\delta$ are either upper bounds on $\delta$, which hold
unchanged for every admissible $\delta$, or belong to the items of (iii): the thresholds and
definitions just listed, the lower bound $\delta_0M\ge2\kappa$, and the lower bound on $\kappa$.
None of them is a new upper bound on $\gamma$, and none is an upper bound on $L$, $\kappa$,
$\kappa_1$, $M_1$ or $M$. This proves (iii).
\end{proof}

\begin{lemma}[Enlargement by one quarter against Ba{\l}aban's one fifth]
\label{5.2:lem-enlargement-fractions}
Let $\beta_{\mathrm{pr}}:=1/5$ be the enlargement fraction carried by the objects of
Ba{\l}aban's construction \cite[p.~190, (1.64)]{B-CMP122a}, and let $\beta_s:=1/4$ be the
enlargement fraction of the families of the correlation theorem. Let
$\delta\in\,]0,\delta_J(L)]$ and $\mathsf{r}(\delta)$ be as in \eqref{5.2:eq-decay-fraction-a},
and let $\beta_\sharp=1/20$ and $a=(\frac32+\beta_\sharp)\kappa$ be as in clause (c) of the
proposition bounding the resummed terms at the localisation site (the resummation proposition,
in what follows). Then the following hold.

\smallskip
\noindent\textup{(i)} Consider the chain terms $\mathbb{C}^{(n)}_k(X,\mathcal{S};w)$ of the
$\mathbf{R}$-operation of localisation level $m=k$. They are delivered by the weighted bound on
the chain terms with the weight $e^{a_kd_k(X)}$, where
\begin{equation*}
a_k:=(1+3\beta_{\mathrm{pr}})\kappa=\tfrac85\kappa .
\end{equation*}
These terms meet the output rate $(1+2\beta_s)\kappa=\frac32\kappa$ of the resummation
proposition, clause (c), with $\kappa/10$ to spare. They are
resummed at the localisation site under their own weight: the sum over $X$ is the norm
$\|\cdot\|^{A}$, and no hull is taken. Let $Y$ be the domain in which a resummed term is filed.
Then the factor $e^{-(3/2+1/20)\kappa d_{k,Z}(Y)}$ is extracted; the fraction $1/20=\beta_\sharp$
is the input margin of the strips, spent in clause (c) of that proposition under the standing
reading convention fixed at the beginning of the construction. The spare factor
$e^{-(\kappa/20)d_k(X)}$ pays the count of domains, whose constant
depends on $d$ only, under the floor $\kappa\ge\kappa_d$. In numbers,
\begin{equation}\label{5.2:eq-enlargement-fractions-a}
\tfrac85=\tfrac32+\tfrac1{20}+\tfrac1{20}.
\end{equation}
Consider next the chain terms of localisation level $m<k$. They are delivered with the weight
$e^{a_md_m(X)}$ in $d_m(X)$, where $a_m:=(1-\beta_{\mathrm{pr}})\kappa=\frac45\kappa$. These
terms are filed in a level-$k$ domain $\bar X$ with
\begin{equation}\label{5.2:eq-enlargement-fractions-1}
d_k(\bar X)\le \frac{d_m(X)}{L}+c_{\mathrm{tr}}\le \frac{d_m(X)}{3}+c_{\mathrm{tr}} .
\end{equation}
Consequently
\begin{equation*}
\bigl(\tfrac45-\tfrac{31}{60}\bigr)\kappa d_m(X)=\tfrac{17}{60}\kappa d_m(X)
\ge\tfrac{\kappa}{20}d_m(X),
\end{equation*}
and this pays the count at scale $m$. The factor $e^{(31/20)c_{\mathrm{tr}}\kappa}$ joins the
factor $O(1)$ of the constant $C_\sharp$. No part of that $O(1)$ is required to be absorbed by
$\kappa$.

\smallskip
\noindent\textup{(ii)} Let $c_{\mathrm{pr}}\ge0$ be the enlargement constant of \cite{B-CMP119}:
the enlarged spaces there have radii $(1+c_{\mathrm{pr}}\beta_{\mathrm{pr}})\alpha$ in place of
radii $\alpha$. After the large-field operation $\mathbf{T}$, the families of the
correlation theorem have radii $(1+2\beta_s)\alpha=\frac32\alpha$, the radii fixed for them at the
beginning of the construction. If Ba{\l}aban's corresponding number $c_{\mathrm{pr}}$ is carried
along, these families have instead the radii $(1+(c_{\mathrm{pr}}+1)\beta_s)\alpha$, of which
$(1+2\beta_s)\alpha$ is the value at $c_{\mathrm{pr}}=1$. The radii
$(1+(c_{\mathrm{pr}}+1)\beta_s)\alpha$ dominate $(1+c_{\mathrm{pr}}\beta_{\mathrm{pr}})\alpha$ for
every $c_{\mathrm{pr}}\ge0$. Hence every statement of Ba{\l}aban quantified over the
$(1+c_{\mathrm{pr}}\beta_{\mathrm{pr}})$-enlarged space applies to the families of the correlation
theorem, these being defined on a space containing that enlarged space; this is how such
statements are used in the last sentence of the resummation proposition.

The room fraction $\mathsf{c}_e$ does not depend on Ba{\l}aban's number $c_{\mathrm{pr}}$. It is
defined as $\mathsf{c}_e:=\theta_S$, the fraction $\theta_S\in[\frac12,\frac89]$ of the layer
schedule that the present construction fixes after the operation $\mathbf{T}$, in entry (c) of
the table of layer fractions. Hence $\mathsf{w}_0=1/24$ whatever $c_{\mathrm{pr}}$ is.

The two single-step analyticity lemmas have inputs $h\in\mathfrak{B}_k$ and a configuration $V$.
They are stated on Ba{\l}aban's $(1+\beta_{\mathrm{pr}})$-enlarged sets. Their inputs live on the
space $\mathcal{U}$ at absolute radii, which contains both enlarged sets. Moreover, the radii
$\hat\alpha_k=(1+2\beta_s)\alpha_{\cdot}(g_k)$ of the correlation theorem's families contain every
$(1+\beta_{\mathrm{pr}})\alpha_{\cdot,k+1}$-set on which those lemmas are applied. Here the dot
stands for the index of the radius. Nothing is lost in either direction.

\smallskip
\noindent\textup{(iii)} The output rate of the kernels produced by the operation $\mathbf{T}$
and the weight at the localisation site compare with Ba{\l}aban's inductive rates
\cite{B-CMP119} as follows:
\begin{align*}
\mathsf{r}(\delta)\kappa&\ge(1+4\beta_s)\kappa=2\kappa\ge(1+4\beta_{\mathrm{pr}})\kappa
=\tfrac95\kappa,\\
(1+3\beta_s)\kappa&=\tfrac74\kappa\ge(1+3\beta_{\mathrm{pr}})\kappa=\tfrac85\kappa .
\end{align*}
Thus the correlation theorem's own families satisfy Ba{\l}aban's inductive hypothesis on the
rates a fortiori.

The heavier weight $\frac74\kappa$ at the localisation site is a choice of the present
construction. It is paid by the construction's own margins. For a term of the correlation
theorem's own families whose domain has length $d_j$ at level $j\le k$, these margins are
$(\frac5{12}-\frac14\delta)\kappa d_j$ for $j<k$ and $(\frac14-\frac14\delta)\kappa d_k$ for
$j=k$, as given by the margin lemma for the boundary members, clause (ii), and by the clause at
$j=k$ of the budget for the constant $C_\sharp$. Of these margins the strips' $\kappa/20$ is spent
under the standing reading convention. Indeed
\begin{equation*}
\tfrac74\kappa\ge\tfrac85\kappa=a+\tfrac1{20}\kappa ,
\end{equation*}
so there remains $\frac3{20}\kappa\ge\frac1{20}\kappa$ for the first, second, third and fifth
rows and the row $\partial$ of that budget. This is besides the factor
$e^{-\frac14\delta\kappa d_j}$ kept for the anchor counts in the resummation proposition,
clause (c).

No family that enters by statement, rather than as one of the correlation theorem's own
families, is summed under the weight $\frac74\kappa$ of the correlation theorem, except through
\textup{(i)}.
\end{lemma}

\begin{proof}
All assertions are comparisons of the two fractions $\beta_{\mathrm{pr}}=1/5$ and $\beta_s=1/4$.
We take the clauses in order.

\smallskip
\noindent\textup{(i)} The setting is that of the resummation proposition, clauses (a) to (c).
There the chain terms of localisation level $m=k$ are
delivered with the weight $a_k=(1+3\beta_{\mathrm{pr}})\kappa$, which is the weight
$(1+3\beta)\kappa$ of Ba{\l}aban's inductive hypothesis \cite{B-CMP119} read at
$\beta=\beta_{\mathrm{pr}}$, and those of level $m<k$ with the weight
$a_m=(1-\beta_{\mathrm{pr}})\kappa$ in $d_m(X)$.

Let first $m=k$. Since $3\beta_{\mathrm{pr}}=3/5$, we have $a_k=\frac85\kappa$, and
\begin{equation*}
a_k-(1+2\beta_s)\kappa=\bigl(\tfrac{16}{10}-\tfrac{15}{10}\bigr)\kappa=\tfrac{\kappa}{10}.
\end{equation*}
This is the spare of $\kappa/10$. Writing $\frac1{10}=\frac1{20}+\frac1{20}$ gives
\eqref{5.2:eq-enlargement-fractions-a}. Since $a=(\frac32+\beta_\sharp)\kappa=\frac{31}{20}\kappa$,
the decay factor $e^{-a_kd_k(X)}$ that accompanies the weight $e^{a_kd_k(X)}$ splits
correspondingly as
\begin{equation*}
e^{-\frac85\kappa d_k(X)}=e^{-a\,d_k(X)}\,e^{-\frac{\kappa}{20}d_k(X)} .
\end{equation*}

The chain terms are summed under the weight with which they are delivered. Fix a cube. The sum,
over the domains $X$ containing it, of $|\mathbb{C}^{(n)}_k(X,\mathcal{S};w)|$ weighted by
$e^{a_kd_k(X)}$ is the norm $\|\cdot\|^{A}$ of the family. Hence no hull of $X$ enters the
resummation.

By clause (c) of the resummation proposition, the $Z$-exterior tree length of the filing domain
$Y$ is at most the length of the domain of the filed term: $d_{k,Z}(Y)\le d_k(X)$ when $m=k$, and
$d_{k,Z}(Y)\le d_k(\bar X)$ when $m<k$. Hence the first factor of the split satisfies
$e^{-a\,d_k(X)}\le e^{-a\,d_{k,Z}(Y)}$, and this provides the factor
$e^{-(3/2+1/20)\kappa d_{k,Z}(Y)}$ that clause
(c) of that proposition extracts at the filing domain $Y$. The summand $\beta_\sharp\kappa=\kappa/20$
of $a$ is the input margin of the strips, spent there under the standing reading convention.

The second factor pays the count of domains. For an integer $u\ge0$, the number of domains $X$
containing a fixed cube with $u\le d_k(X)<u+1$ is at most $e^{c_d(u+1)}$, with $c_d$ depending on
$d$ only. The spare factor therefore turns the count into
\begin{equation*}
\sum_{u\ge0}e^{c_d(u+1)}e^{-\frac{\kappa}{20}u}
\le e^{c_d}\sum_{u\ge0}e^{-(\frac{\kappa}{20}-c_d)u}\le e^{c_d}\bigl(1-e^{-2}\bigr)^{-1},
\end{equation*}
a number depending on $d$ only. The last bound holds because the floor $\kappa\ge\kappa_d$ contains
the condition $\kappa/20\ge c_d+2$, so that $\kappa/20-c_d\ge2$.

Let now $m<k$. Here $a_m=(1-\frac15)\kappa=\frac45\kappa$. The blocking factor $L$ is an odd
integer greater than one, so $1/L\le1/3$; this is the second inequality in
\eqref{5.2:eq-enlargement-fractions-1}. The first inequality there is the filing of the term in
the level-$k$ domain $\bar X$, as the resummation proposition provides it. By
\eqref{5.2:eq-enlargement-fractions-1}, the weight $e^{a\,d_k(\bar X)}$ demanded at the filing
domain satisfies
\begin{equation*}
e^{a\,d_k(\bar X)}\le e^{a c_{\mathrm{tr}}}\,e^{\frac a3 d_m(X)}
=e^{\frac{31}{20}c_{\mathrm{tr}}\kappa}\,e^{\frac{31}{60}\kappa d_m(X)} .
\end{equation*}
Against the delivered factor $e^{-\frac45\kappa d_m(X)}$ this leaves
\begin{equation*}
e^{-(\frac{48}{60}-\frac{31}{60})\kappa d_m(X)}=e^{-\frac{17}{60}\kappa d_m(X)}
\le e^{-\frac{\kappa}{20}d_m(X)},
\end{equation*}
since $\frac{17}{60}\ge\frac3{60}=\frac1{20}$. The last factor pays the count at scale $m$ by
the geometric series above, with $d_m$ in place of $d_k$. Its convergence rests on the same
condition $\kappa/20\ge c_d+2$.

What remains is the factor $e^{(31/20)c_{\mathrm{tr}}\kappa}$. It depends on $\kappa$ and
$c_{\mathrm{tr}}$ only, not on the domains. It joins the factor $O(1)$ of the constant
$C_\sharp$: it is multiplied into that factor, not compared with any exponential decay in the
domains. Therefore no lower bound on $\kappa$ is imposed to absorb that $O(1)$, and no part of it
is asked of $\kappa$.

\smallskip
\noindent\textup{(ii)} For $c_{\mathrm{pr}}\ge0$ we have
\begin{equation*}
1+(c_{\mathrm{pr}}+1)\beta_s\ge1+c_{\mathrm{pr}}\beta_{\mathrm{pr}}
\iff\frac{c_{\mathrm{pr}}+1}{4}\ge\frac{c_{\mathrm{pr}}}{5}
\iff5(c_{\mathrm{pr}}+1)\ge4c_{\mathrm{pr}} ,
\end{equation*}
and the last inequality holds for every $c_{\mathrm{pr}}\ge0$. At $c_{\mathrm{pr}}=1$ the left
member is $1+2\beta_s=\frac32$ and the right member is $\frac65$.

The enlarged spaces are defined through their radii. Consequently the space enlarged by the
factor $1+c_{\mathrm{pr}}\beta_{\mathrm{pr}}$ is contained in the space enlarged by the larger
factor. A statement of Ba{\l}aban quantified over the former therefore applies to the
correlation theorem's families, which live on a space containing it.

The room fraction is, by definition, $\mathsf{c}_e:=\theta_S$. Here $\theta_S$ is the fraction of
the layer schedule fixed after the operation $\mathbf{T}$ of the present construction, in entry
(c) of the table of layer fractions, and it is defined without reference to
$\beta_{\mathrm{pr}}$ or $c_{\mathrm{pr}}$. The number
$\mathsf{w}_0=1/24$ is fixed. It coincides with $\min\{\frac12,\mathsf{c}_e\}/12$ for every
$\theta_S\in[\frac12,\frac89]$, because then $\min\{\frac12,\theta_S\}=\frac12$. Neither quantity
involves $c_{\mathrm{pr}}$.

The two single-step analyticity lemmas are stated with hypotheses on Ba{\l}aban's
$(1+\beta_{\mathrm{pr}})$-enlarged sets. Their inputs $h\in\mathfrak{B}_k$ and $V$ are taken in
the space $\mathcal{U}$ at absolute radii. That space contains both the
$(1+\beta_{\mathrm{pr}})$-enlarged and the $(1+2\beta_s)$-enlarged sets at coupling-dependent
radii. Hence inputs drawn from either enlarged set are admissible inputs of the lemmas; the
lemmas' conclusions are used on the $(1+\beta_{\mathrm{pr}})\alpha_{\cdot,k+1}$-sets only.

It remains to see that every $(1+\beta_{\mathrm{pr}})\alpha_{\cdot,k+1}$-set on which the lemmas
are applied lies inside the set of radii $\hat\alpha_k=(1+2\beta_s)\alpha_{\cdot}(g_k)$. The
enlargement factors compare as
\begin{equation*}
1+\beta_{\mathrm{pr}}=\tfrac65<\tfrac32=1+2\beta_s .
\end{equation*}
The inclusion of the sets, in which the radii $\alpha_{\cdot,k+1}$ at level $k+1$ are compared
with the radii $\alpha_{\cdot}(g_k)$ as well, is the inclusion statement for the radii
$\hat\alpha_k$ made together with the two single-step analyticity lemmas, which we use by
statement. Thus nothing is lost in either direction.

\smallskip
\noindent\textup{(iii)} By Lemma~\ref{5.2:lem-decay-fraction}, the output rate satisfies
$\mathsf{r}(\delta)\kappa\ge(1+4\beta_s)\kappa=2\kappa$ for every
$\delta\in\,]0,\delta_J(L)]$. Further,
\begin{equation*}
2\ge\tfrac95=1+4\beta_{\mathrm{pr}},\qquad
1+3\beta_s=\tfrac74\ge\tfrac85=1+3\beta_{\mathrm{pr}},
\end{equation*}
the latter since $\frac{35}{20}\ge\frac{32}{20}$.
These two comparisons give the displayed inequalities. Ba{\l}aban's inductive hypothesis
\cite{B-CMP119} requires the rates $(1+4\beta)\kappa$ for the fresh kernels and $(1+3\beta)\kappa$
for the weight, at $\beta=\beta_{\mathrm{pr}}$. The families of the correlation theorem carry
rates at least as large, so they satisfy that hypothesis a fortiori.

The weight $\frac74\kappa$ is paid as follows. Consider a term of the correlation theorem's own
families whose domain has length $d_j$ at level $j\le k$. For $j<k$, the margin lemma for the
boundary members, clause (ii), supplies the margin $(\frac5{12}-\frac14\delta)\kappa d_j$. For
$j=k$, it supplies $(\frac14-\frac14\delta)\kappa d_k$, together with the clause at $j=k$ of the
budget for $C_\sharp$. At $j=k$ the margin $\frac14=2-\frac74$ rests on the output rate
$\mathsf{r}(\delta)\kappa\ge2\kappa$ of Lemma~\ref{5.2:lem-decay-fraction}.

From these margins the strips' $\kappa/20$ is spent under the standing reading convention. Since
\begin{equation*}
a+\tfrac1{20}\kappa=\tfrac{31}{20}\kappa+\tfrac1{20}\kappa=\tfrac85\kappa ,
\end{equation*}
we have
\begin{equation*}
\tfrac74\kappa-\bigl(a+\tfrac1{20}\kappa\bigr)=\bigl(\tfrac{35}{20}-\tfrac{32}{20}\bigr)\kappa
=\tfrac3{20}\kappa\ge\tfrac1{20}\kappa .
\end{equation*}
This remainder is available for the first, second, third and fifth rows and the row $\partial$
of the budget. The factor $e^{-\frac14\delta\kappa d_j}$ is a separate factor, at the
rate $\frac14\delta\kappa$. It is kept for the anchor counts and is supplied in clause (c) of the
resummation proposition.

Finally, the sum at the localisation site is organised in that proposition as follows. The
families that enter it by statement, rather than as the correlation theorem's own families,
enter it only as the chain terms of \textup{(i)}, and these are summed under their own weights
$a_k$ or $a_m$. None of them is summed under the weight $\frac74\kappa$. This is the last
assertion of \textup{(iii)}.
\end{proof}

\begin{lemma}[The driven flow of the running shift]\label{5.3:lem-driven-flow}
For $c\in\mathbb{C}$ and $R>0$ let $D(c,R)$ denote the open disc of radius $R$ about $c$. Let $K\ge 1$
and $\mathfrak{d}_0,\dots,\mathfrak{d}_{K-1}>0$. For $0\le k\le K-1$ put
\begin{equation*}
\mathfrak{P}_k:=\prod_{j\le k}D(0,\mathfrak{d}_j)\subset\mathbb{C}^{k+1},
\qquad
\mathfrak{d}_*(k):=\min_{j\le k}\mathfrak{d}_j .
\end{equation*}
For each $k<K$ let $f_{k+1}$ be holomorphic on $\mathfrak{P}_k$ with $|f_{k+1}|\le\omega_k$. Let
$\varrho>0$ satisfy $\tfrac83\varrho\le\mathfrak{d}_*(K-1)$. Suppose that, for each $k<K$, whenever
$\zeta_0,\zeta_1,\dots,\zeta_k$ are holomorphic functions on $D(0,\varrho)$, with $\zeta_0(\zeta):=\zeta$ and
$|\zeta_j(\zeta)-\zeta|\le\tfrac13|\zeta|$ for $j\le k$ and $\zeta\in D(0,\varrho)$, a function
$\mathsf{q}_{k+1}$ holomorphic on $D(0,\varrho)$ with $|\mathsf{q}_{k+1}|\le v_k$ is given; define then
$\zeta_{k+1}$ on $D(0,\varrho)$ by
\begin{equation}\label{5.3:eq-driven-flow-1}
\zeta_{k+1}(\zeta):=\zeta_k(\zeta)
+\bigl[f_{k+1}(\zeta_0(\zeta),\dots,\zeta_k(\zeta))-f_{k+1}(0)\bigr]
+\bigl[\mathsf{q}_{k+1}(\zeta)-\mathsf{q}_{k+1}(0)\bigr].
\end{equation}
Put
\begin{equation}\label{5.3:eq-driven-flow-2}
d_k:=\frac{16}{9}\,\frac{\omega_k}{\mathfrak{d}_*(k)}+\frac{2v_k}{\varrho}\qquad(0\le k<K).
\end{equation}
If $\sum_{k<K}d_k\le\tfrac13$, then for all $k<K$ the flow \eqref{5.3:eq-driven-flow-1} is defined,
and on $D(0,\varrho)$
\begin{equation}\label{5.3:eq-driven-flow-3}
|\zeta_{k+1}-\zeta_k|\le d_k|\zeta|,
\qquad
|\zeta_{k+1}-\zeta|\le\tfrac13|\zeta| .
\end{equation}
\end{lemma}

\begin{proof}
Write $\tfrac12\mathfrak{P}_k$ for the closed polydisc
$\{\xi=(\xi_j)_{j\le k}\in\mathbb{C}^{k+1}:|\xi_j|\le\tfrac12\mathfrak{d}_j\ \text{for all}\ j\le k\}$,
which is contained in $\mathfrak{P}_k$. For every $k<K$ the numbers $\omega_k$ and $v_k$ are
non-negative. Indeed, $|f_{k+1}|\le\omega_k$ holds on the non-empty set $\mathfrak{P}_k$. Further, the
family in which each of $\zeta_0,\dots,\zeta_k$ is the identity map $\zeta\mapsto\zeta$ consists of
holomorphic functions on $D(0,\varrho)$ with $|\zeta_j(\zeta)-\zeta|=0\le\tfrac13|\zeta|$, so the
hypothesis supplies for this family a function $\mathsf{q}_{k+1}$ with $|\mathsf{q}_{k+1}|\le v_k$ on the
non-empty disc $D(0,\varrho)$. Hence, by \eqref{5.3:eq-driven-flow-2}, $d_k\ge0$ for every $k<K$.

We prove by induction on $k\in\{0,\dots,K\}$ the following assertion $P(k)$: the functions
$\zeta_1,\dots,\zeta_k$ are defined by \eqref{5.3:eq-driven-flow-1} (with $\zeta_0(\zeta)=\zeta$), the
functions $\zeta_0,\dots,\zeta_k$ are holomorphic on $D(0,\varrho)$, and
for $\zeta\in D(0,\varrho)$ one has $|\zeta_{i+1}(\zeta)-\zeta_i(\zeta)|\le d_i|\zeta|$ for all $i<k$ and
$|\zeta_j(\zeta)-\zeta|\le\tfrac13|\zeta|$ for all $j\le k$. The assertion $P(0)$ holds, since
$\zeta_0(\zeta)=\zeta$ is holomorphic, there is no increment to bound, and $|\zeta_0-\zeta|=0$.
Assertion $P(K)$ contains all claims of the lemma.

Let $k<K$ and assume $P(k)$. By $P(k)$ the functions $\zeta_0,\dots,\zeta_k$ satisfy the hypothesis
under which $\mathsf{q}_{k+1}$ is given, so $\mathsf{q}_{k+1}$ is holomorphic on $D(0,\varrho)$ with
$|\mathsf{q}_{k+1}|\le v_k$.

\emph{The point $(\zeta_j(\zeta))_{j\le k}$ lies in the half polydisc.} Fix $\zeta\in D(0,\varrho)$ and set
$\xi:=(\zeta_j(\zeta))_{j\le k}$. For $j\le k$, $P(k)$ gives
\begin{equation}\label{5.3:eq-driven-flow-4}
|\xi_j|\le|\zeta|+|\zeta_j(\zeta)-\zeta|\le\tfrac43|\zeta|<\tfrac43\varrho
\le\tfrac12\mathfrak{d}_*(K-1)\le\tfrac12\mathfrak{d}_j ,
\end{equation}
where the penultimate inequality is the hypothesis $\tfrac83\varrho\le\mathfrak{d}_*(K-1)$, and the last
holds because $j\le k\le K-1$, so that $\mathfrak{d}_j$ is one of the numbers whose minimum is
$\mathfrak{d}_*(K-1)$. Hence $\xi\in\tfrac12\mathfrak{P}_k\subset\mathfrak{P}_k$. The map
$\zeta\mapsto(\zeta_j(\zeta))_{j\le k}$ is holomorphic from $D(0,\varrho)$ into $\mathfrak{P}_k$, and
$f_{k+1}$ is holomorphic on $\mathfrak{P}_k$. So the composite
$\zeta\mapsto f_{k+1}(\zeta_0(\zeta),\dots,\zeta_k(\zeta))$ is holomorphic on $D(0,\varrho)$, and
$\zeta_{k+1}$ of \eqref{5.3:eq-driven-flow-1} is defined and holomorphic on $D(0,\varrho)$.

\emph{The $f$-increment.} The set $\tfrac12\mathfrak{P}_k$ is convex, being a product of closed discs, and
it contains $0$ and $\xi$. Hence the segment $\{t\xi:0\le t\le1\}$ lies in
$\tfrac12\mathfrak{P}_k\subset\mathfrak{P}_k$, where $f_{k+1}$ is holomorphic. By the chain rule,
\begin{equation*}
f_{k+1}(\xi)-f_{k+1}(0)=\int_0^1\sum_{j\le k}\xi_j\,\partial_jf_{k+1}(t\xi)\,dt .
\end{equation*}
Therefore
\begin{equation}\label{5.3:eq-driven-flow-5}
|f_{k+1}(\xi)-f_{k+1}(0)|\le\max_{j\le k}|\xi_j|\cdot
\sup_{\eta\in\frac12\mathfrak{P}_k}\sum_{j\le k}|\partial_jf_{k+1}(\eta)| .
\end{equation}
Let $\eta\in\tfrac12\mathfrak{P}_k$. Since $\mathfrak{d}_j\ge\mathfrak{d}_*(k)$ for $j\le k$, the
$\ell^1$ gradient bound on a polydisc (Lemma~\ref{5.2:lem-polydisc-gradient}), applied to $f_{k+1}$ on
$\mathfrak{P}_k$ with sup bound $\omega_k$, gives
\begin{equation*}
\sum_{j\le k}|\partial_jf_{k+1}(\eta)|
\le\frac{1}{\mathfrak{d}_*(k)}\sum_{j\le k}\mathfrak{d}_j|\partial_jf_{k+1}(\eta)|
\le\frac43\,\frac{\omega_k}{\mathfrak{d}_*(k)} .
\end{equation*}
By \eqref{5.3:eq-driven-flow-4}, $\max_{j\le k}|\xi_j|\le\tfrac43|\zeta|$. Inserting both bounds into
\eqref{5.3:eq-driven-flow-5} yields
\begin{equation}\label{5.3:eq-driven-flow-6}
|f_{k+1}(\zeta_0(\zeta),\dots,\zeta_k(\zeta))-f_{k+1}(0)|
\le\frac43|\zeta|\cdot\frac43\,\frac{\omega_k}{\mathfrak{d}_*(k)}
=\frac{16}{9}\,\frac{\omega_k}{\mathfrak{d}_*(k)}\,|\zeta| .
\end{equation}

\emph{The driving increment.} The function $\mathsf{q}_{k+1}$ is holomorphic on $D(0,\varrho)$ and
bounded there by $v_k$. Part (a) of the bounds for a bounded analytic function on a disc
(Lemma~\ref{5.2:lem-disc-bounds}, \eqref{5.2:eq-disc-bounds-a}), applied with centre $0$ and radius
$\varrho$, gives
\begin{equation}\label{5.3:eq-driven-flow-7}
|\mathsf{q}_{k+1}(\zeta)-\mathsf{q}_{k+1}(0)|\le\frac{2v_k}{\varrho}\,|\zeta|
\qquad(\zeta\in D(0,\varrho)).
\end{equation}
The disc is the same $D(0,\varrho)$ at every step $k$, so the radius entering
\eqref{5.3:eq-driven-flow-7} is $\varrho$ for all $k<K$.

\emph{Summation.} By \eqref{5.3:eq-driven-flow-1}, \eqref{5.3:eq-driven-flow-6},
\eqref{5.3:eq-driven-flow-7} and the definition \eqref{5.3:eq-driven-flow-2},
\begin{equation*}
|\zeta_{k+1}(\zeta)-\zeta_k(\zeta)|
\le\Bigl(\frac{16}{9}\,\frac{\omega_k}{\mathfrak{d}_*(k)}+\frac{2v_k}{\varrho}\Bigr)|\zeta|
=d_k|\zeta| .
\end{equation*}
Combining this with the increment bounds for $i<k$ contained in $P(k)$, and using $d_i\ge0$,
\begin{equation*}
|\zeta_{k+1}(\zeta)-\zeta|
\le\sum_{i\le k}|\zeta_{i+1}(\zeta)-\zeta_i(\zeta)|
\le\Bigl(\sum_{i\le k}d_i\Bigr)|\zeta|
\le\Bigl(\sum_{i<K}d_i\Bigr)|\zeta|
\le\tfrac13|\zeta| .
\end{equation*}
Together with $P(k)$ this is $P(k+1)$. The induction is complete, and $P(K)$ gives
\eqref{5.3:eq-driven-flow-3} for every $k<K$.
\end{proof}

\begin{lemma}[Position of the coupling shift in one step]\label{5.3:lem-coupling-shift}
Let $\xi \in \mathbb{C}$, let $x_k = g_k^{-2}$, and put $a := x_k - \xi$. In the density
\cite[(2.1)]{B-CMP109}, which is \cite[(3.15)]{B-CMP119} specialised to the whole lattice,
place $a$ in front of $G(V) + A(U_k(V))$ in place of $g_k^{-2}$, so that this part of the
exponent reads $-a\,[G(V) + A(U_k(V))]$. Take every inserted identity and every choice to be
those of the source-free scaffold of Definition~\ref{5.1:def-source-free-scaffold}. Then, after
the operations of \cite[(2.2)--(2.12)]{B-CMP109}, the exponent is
\begin{equation}\label{5.3:eq-coupling-shift-1}
\begin{split}
-a\,A(U_{k+1})
&+ s\Bigl[-\tfrac12\langle B,\Delta^{(k)}B\rangle
+ g_k^{-2}\sum_i \Psi_i(g_kCB)\Bigr]\\
&+ \sum_i \Phi_i(g_kCB) + \{\dots\},
\end{split}
\end{equation}
where
\begin{equation}\label{5.3:eq-coupling-shift-2}
s := a\,g_k^{2} = 1 - \xi g_k^{2},
\end{equation}
the complex Gaussian parameter of the step (the letter is not used here for the depth below
the cutoff), and $\{\dots\}$ denotes the term in curly brackets of the exponent of
\cite[(2.12)]{B-CMP109}, which is not written out here.
\end{lemma}

\begin{proof}
The second equality in \eqref{5.3:eq-coupling-shift-2} is the identity
$a g_k^2 = (g_k^{-2} - \xi) g_k^2 = 1 - \xi g_k^2$.

We sort the terms of the exponent produced by \cite[(2.2)--(2.12)]{B-CMP109} by their origin.
The terms $\Phi_i$ (the $\operatorname{Tr}\log$ of
the Jacobians and the logarithms of the densities of the Haar measures) do not come from the
action.
They arise from the changes of variables and from the Haar measure, and are therefore
unaffected by the coefficient placed in front of $G + A$. This gives the term
$\sum_i \Phi_i(g_kCB)$ in \eqref{5.3:eq-coupling-shift-1}, without a factor $s$.

The term $A(U_{k+1})$, the quadratic form of \cite[(2.11)]{B-CMP109}, and every $\Psi_i$ come
from $-g_k^{-2}[G + A]$. They are produced by four operations:
\begin{itemize}
\item the translation to the background $V^{(k)}$;
\item the change of variables $B' = B - h\tilde{D}(B)$;
\item the scaling by the real coupling $g_k$;
\item the elimination by $\delta(QB')$.
\end{itemize}
None of these depends on the coefficient in front of $G + A$. By
Definition~\ref{5.1:def-source-free-scaffold} and its display
\eqref{5.1:eq-source-free-scaffold-a}, the background and the operators $h$, $\tilde{D}$, $C$,
$Q$, $\Delta^{(k)}$, $C^{(k)}$ are the choices at $z = 0$ and at the real $g_k$; in particular
the background is a minimiser of the unweighted action. Hence, when $g_k^{-2}$ is replaced by
$a$, each of these terms is multiplied by $a/g_k^{-2} = a g_k^2 = s$.

For the term $A(U_{k+1})$ this gives the coefficient $-a$. For the quadratic form, the
unshifted exponent contains $-\tfrac12\langle B,\Delta^{(k)}B\rangle$, which becomes
$-\tfrac{s}{2}\langle B,\Delta^{(k)}B\rangle$. For the higher-order terms, the unshifted exponent
contains $g_k^{-2}\sum_i\Psi_i(g_kCB)$, which becomes $s\, g_k^{-2}\sum_i\Psi_i(g_kCB)$. These
are the first two groups of terms in \eqref{5.3:eq-coupling-shift-1}.

There remains the term linear in the fluctuation field, the third, linear term of
\cite[(2.10)]{B-CMP109}. Of this term Ba{\l}aban states that it
\begin{quote}
vanishes now \dots\ because $\langle \delta A', J\rangle = 0$ for all $\delta A'$ satisfying
the condition
\end{quote}
where $\delta A'$ and $J$ are the letters of \cite[(2.10)]{B-CMP109} and the condition is the
one imposed there on $\delta A'$. This is a property of $U_{k+1}$ alone, and it holds
whatever coefficient multiplies the action. So the linear term vanishes here as well.

Finally, that $G$ is given the same coefficient $a$ as $A$, as in the hypothesis, is what the
clause of \eqref{5.1:eq-source-free-scaffold-b} on the exponential gauge-fixing coefficient
dictates: the coefficient of block $y$ follows the Wilson weight, namely it is $x_k$ minus the
shift, which here is $x_k - \xi = a$. The operations listed above therefore act on $G$ exactly
as on $A$. This proves \eqref{5.3:eq-coupling-shift-1}.
\end{proof}

The Gaussian part of \eqref{5.3:eq-coupling-shift-1} is
$\exp\bigl(-\tfrac{s}{2}\langle B,\Delta^{(k)}B\rangle\bigr)$. Let $dB$ be Lebesgue measure on
the $N_k$ real fluctuation variables, and put
\begin{equation*}
Z^{(k)}(U_{k+1}) := \int \exp\bigl(-\tfrac12\langle B,\Delta^{(k)}B\rangle\bigr)\,dB .
\end{equation*}
For $s$ with positive real part, $s^{-N_k/2}$ being the principal branch, the unnormalised
Gaussian $\exp\bigl(-\tfrac{s}{2}\langle B,\Delta^{(k)}B\rangle\bigr)\,dB$ is
\begin{equation}\label{5.3:eq-coupling-shift-3}
s^{-N_k/2}\, Z^{(k)}(U_{k+1})\, d\mu_{C^{(k)}/s}.
\end{equation}
Here $d\mu_{C^{(k)}/s}$ is the normalised Gaussian measure with covariance $C^{(k)}/s$ (a complex
measure when $s$ is not real). In particular the one-loop term $Z^{(k)}(U_{k+1})$ does not depend
on $s$.

\begin{lemma}[Source part of the Wilson action]\label{5.3:lem-wilson-source}
Consider a site of Ba{\l}aban's operation $\mathbf{R}$ at level $k$, with large-field region $Z$,
configurations $U_k$, $U''_k$, $U''_{h_{\mathbf{R}}+2}$ (the index written $h$ in
\cite{B-CMP122b} is denoted here $h_{\mathbf{R}}$, to keep it apart from the other uses of the
letter $h$; this index and these configurations are as in
\cite{B-CMP122b}) and cutoff functions $\hat{\chi}_0$, $\hat{\chi}$, $\hat{\chi}_1$ as in
\cite{B-CMP122b}, where they are written with the letter $\zeta$ ($\zeta_0$, $\zeta$, $\zeta_1$);
here the letter $\zeta$ is reserved for the shifts produced by the complex source $w$ ($\zeta$ the
constant bare shift, $\zeta_k$ the running shift). These satisfy
\begin{equation*}
1=(1-\hat{\chi}_0)+\hat{\chi}+\hat{\chi}_1 .
\end{equation*}
For a weight $c$ on plaquettes let $A(c,U)$ be the weighted Wilson action
\begin{equation*}
A(c,U)=\sum_p c(p)\,a_p(U),
\qquad a_p(U)=1-\tfrac12\operatorname{tr}\bigl(U+U^{-1}\bigr)(\partial p).
\end{equation*}
Let $\zeta_k$ be the running shift at a complex source $w$, and assume that
$\sup_p|\zeta_k(p)|\le\tfrac83 r$, where $r$ is the number in the standing smallness hypothesis
on the shifts under which the lemma on complex plaquette weights in the expansions
(Lemma~\ref{5.2:lem-complex-weights}) is proved. Let $\zeta_k^{\mathrm{new}}$ be the weight
obtained from clause (a) of the definition of the running shift $\zeta_k$ when the functions
$\varphi_j$ occurring in that clause are replaced by $1$ on $Z$. Define $\mathfrak{c}_w$ by the
second line of the following display. Then the identity in its first line holds exactly:
\begin{equation}\label{5.3:eq-wilson-source-a}
\begin{gathered}
A(\zeta_k,U''_k)=A(\zeta_k^{\mathrm{new}},U_k)+A(\zeta_k\hat{\chi}_1,U''_{h_{\mathbf{R}}+2})
+\mathfrak{c}_w ,\\
\begin{aligned}
\mathfrak{c}_w
:={}&\bigl[A(\zeta_k(1-\hat{\chi}_0),U''_k)-A(\zeta_k(1-\hat{\chi}_0),U_k)\bigr]
+A(\zeta_k\hat{\chi},U''_k)\\
&+\bigl[A(\zeta_k\hat{\chi}_1,U''_k)-A(\zeta_k\hat{\chi}_1,U''_{h_{\mathbf{R}}+2})\bigr]
-A\bigl(\zeta_k^{\mathrm{new}}-\zeta_k(1-\hat{\chi}_0),U_k\bigr).
\end{aligned}
\end{gathered}
\end{equation}
The four terms in the definition of $\mathfrak{c}_w$ in \eqref{5.3:eq-wilson-source-a} obey the
following bounds.
\begin{itemize}
\item[(a)] The first and the third brackets are each bounded by $\theta_W$ times Ba{\l}aban's
bounds for the corresponding expansions \cite[(1.16)--(1.17), (1.36), (1.46)--(1.47)]{B-CMP122b},
where $\theta_W$ is the constant of clauses (a) and (c) of Lemma~\ref{5.2:lem-complex-weights}.
\item[(b)] The term $A(\zeta_k\hat{\chi},U''_k)$ and the term $A(\zeta_k^{\mathrm{new}}-\zeta_k(1-\hat{\chi}_0),U_k)$ are each bounded by
\begin{equation*}
\tfrac83\, r\,C_{\mathrm{cw}}\bigl(O(\varepsilon_k)+\delta'_k\bigr)^2
\end{equation*}
per cube $\square$ of Ba{\l}aban's cover by doubled cubes that carries a function
$\hat{\chi}_{\square}$ of the partition of unity on that cover (written $\zeta_{\square}$ in
\cite[p.~270]{B-CMP109}), where $C_{\mathrm{cw}}$ denotes the constant of
clause (b) of Lemma~\ref{5.2:lem-complex-weights},
and $\varepsilon_k$, $\delta'_k$ are the parameters so denoted in \cite{B-CMP122b}.
\end{itemize}
\end{lemma}

\begin{proof}
\emph{The identity.} By its definition, $A(c,U)$ is linear in the weight $c$ for fixed $U$. In particular,
\begin{equation*}
A\bigl(\zeta_k^{\mathrm{new}}-\zeta_k(1-\hat{\chi}_0),U_k\bigr)
=A(\zeta_k^{\mathrm{new}},U_k)-A(\zeta_k(1-\hat{\chi}_0),U_k).
\end{equation*}
Substitute this into the definition of $\mathfrak{c}_w$ in \eqref{5.3:eq-wilson-source-a} and add
$A(\zeta_k^{\mathrm{new}},U_k)+A(\zeta_k\hat{\chi}_1,U''_{h_{\mathbf{R}}+2})$. Several terms then
cancel in pairs:
\begin{itemize}
\item $A(\zeta_k^{\mathrm{new}},U_k)$ cancels against its negative;
\item $A(\zeta_k\hat{\chi}_1,U''_{h_{\mathbf{R}}+2})$ cancels against the negative term of the
third bracket;
\item $-A(\zeta_k(1-\hat{\chi}_0),U_k)$ from the first bracket cancels against $+A(\zeta_k(1-\hat{\chi}_0),U_k)$ from the last term.
\end{itemize}
What remains is
\begin{equation*}
\begin{split}
&A(\zeta_k^{\mathrm{new}},U_k)+A(\zeta_k\hat{\chi}_1,U''_{h_{\mathbf{R}}+2})+\mathfrak{c}_w\\
&\qquad=A(\zeta_k(1-\hat{\chi}_0),U''_k)+A(\zeta_k\hat{\chi},U''_k)
+A(\zeta_k\hat{\chi}_1,U''_k)\\
&\qquad=A\bigl(\zeta_k[(1-\hat{\chi}_0)+\hat{\chi}+\hat{\chi}_1],U''_k\bigr)
=A(\zeta_k,U''_k).
\end{split}
\end{equation*}
The second equality is linearity in the weight once more. The last equality uses the decomposition of unity $1=(1-\hat{\chi}_0)+\hat{\chi}+\hat{\chi}_1$. This proves \eqref{5.3:eq-wilson-source-a}. No estimate enters, so the identity is exact.

\emph{The first and third brackets.} In each of these two brackets the same weight multiplies two
Wilson actions taken at two different configurations. The weight is $\zeta_k(1-\hat{\chi}_0)$ in
the first bracket and $\zeta_k\hat{\chi}_1$ in the third. Ba{\l}aban expands such differences in
\cite[(1.16)--(1.17), (1.36), (1.46)--(1.47)]{B-CMP122b}. Each of these expansions is linear in
the weight, and each is bounded by the sum over plaquettes $p$ of the modulus of the weight at $p$
times the absolute value of the corresponding term. Replace Ba{\l}aban's real weight by the
complex weight carried by the bracket. Clauses (a) and (c) of
Lemma~\ref{5.2:lem-complex-weights} then bound the bracket by $\theta_W$ times the bound in
\cite{B-CMP122b}. This gives (a).

\emph{The remaining two terms.} The terms $A(\zeta_k\hat{\chi},U''_k)$ and
$A(\zeta_k^{\mathrm{new}}-\zeta_k(1-\hat{\chi}_0),U_k)$ are not differences of this kind. Clause
(b) of Lemma~\ref{5.2:lem-complex-weights} applies to each of them and bounds it by
$\tfrac83 rC_{\mathrm{cw}}(O(\varepsilon_k)+\delta'_k)^2$ per cube $\square$ of that cover carrying a
$\hat{\chi}_{\square}$. This gives (b).

\emph{Placement of $\mathfrak{c}_w$ in $\mathbf{C}_k$.} By (a) and (b), $\mathfrak{c}_w$ joins Ba{\l}aban's sum $\mathbf{C}_k$ of small new
terms produced by the expanded Wilson actions \cite{B-CMP122b}; it is the $w$-dependent part of
that sum. Of the terms of this sum Ba{\l}aban remarks in \cite{B-CMP122b} that they
are all small, although only their boundedness is needed; the same holds for $\mathfrak{c}_w$.

Two further places in the $\mathbf{R}$-operation involve the Wilson action.
\begin{itemize}
\item Between numerator and denominator there is the cancellation of $g_k^{-2}A(\hat{\chi}_0,U_0)$
in \cite[(1.10)]{B-CMP122b}, where $U_0$ is the configuration so denoted there. This cancellation
does not depend on $z$. The source term has no counterpart of it, and, since it carries no factor
$g_k^{-2}$, needs none.
\item At \cite[(1.62)]{B-CMP122a}, a further place at which the parameter written $\beta$ in
\cite{B-CMP122a} enters, the operation consists in adding and subtracting weights independent
of $z$, so it contributes no source part.
\end{itemize}
Hence neither of these two places produces a $z$-dependent Wilson term.
\end{proof}

\begin{proof}[Joint induction: the coupling increment and the reference coupling]
\label{5.3:prf-flow-reference}
We continue the passage from the hypothesis of the joint induction at level $k$
(Definition~\ref{5.2:def-joint-hypothesis}) to the same hypothesis at level $k+1$. This passage
begins with the pure step and the frozen integral, treated above. As there, $k<K$, so that
$g_j<g$ for $j\le k$, that is, $x_j>X=g^{-2}$.

\emph{The increment as a function of the shift.} The coefficient $b_{k+1}$ of the induction
format (Definition~\ref{5.2:def-induction-format}) is decomposed as
\begin{equation}\label{5.3:eq-flow-reference-1}
b_{k+1}=\beta^0_{k+1}+f_{k+1}(\zeta_0,\dots,\zeta_k)+\mathsf{q}_{k+1}.
\end{equation}
We apply the flow lemma, which bounds the dependence on the shift of an increment of the form
\eqref{5.3:eq-flow-reference-1}, to $b_{k+1}$, with the
following data:
\begin{itemize}
\item the polydisc radii $\mathfrak{d}_j=\sigma_0 x_j$; the quantity $\mathfrak{d}_*(k)$ that the
flow lemma forms from them satisfies $\mathfrak{d}_*(k)\ge\sigma_0 y_k$, where
$\sigma_0$ is the radius of the discs $D(1,\sigma_0)$ of the pure step and $y_k$ is the constant of
Definition~\ref{5.1:def-reference-recursion};
\item the sup bound furnished by the pure-step bound \eqref{5.2:eq-pure-step-a},
\begin{equation*}
\omega_k=C_\beta\mathrm{e}_k\le C_\beta\,\mathrm{e}\bigl(y_k^{-1/2}\bigr),
\end{equation*}
in which $\mathrm{e}(\cdot)$ is the function of the coupling of
Definition~\ref{5.1:not-coupling-functions};
\item the driving term $\mathsf{q}_{k+1}=\beta[Q^{(k+1)}_\zeta]$, with $Q^{(k+1)}_\zeta$ the family of
Definition~\ref{5.2:def-three-families}, whose modulus at $\zeta=0$ is bounded by the remainder
constant $v_k$ of Definition~\ref{5.1:def-reference-recursion};
\item an arbitrary $\varrho\in[r,2r]$.
\end{itemize}
The flow lemma requires $\tfrac{8}{3}\cdot 2r\le\sigma_0 g^{-2}$. This inequality is one of the
conditions on $g$ collected in the definition of $\gamma_J$; it holds for $g\le\gamma_J$. With
these data the flow lemma gives
\begin{equation}\label{5.3:eq-flow-reference-2}
|b_{k+1}(\zeta)-b_{k+1}(0)|\le d_k(\varrho)\,|\zeta|,\qquad \zeta\in D(0,\varrho).
\end{equation}

\emph{The running shift.} Clause (a) of the hypothesis of the joint induction controls the
reference couplings for pairs of levels at most $k$. From this clause and the summability
condition on the Lipschitz constants $d_j(\varrho)$ (the constants of
Definition~\ref{5.1:def-reference-recursion}), applied with the inequality $x_k>X$, we obtain
$\sum_{j\le k}d_j(\varrho)\le\tfrac13$.

In the recursion of the flow lemma the shift advances by the increment,
$\zeta_l-\zeta_{l-1}=b_l(\zeta)-b_l(0)$ for $1\le l\le k+1$, and $\zeta_0=\zeta$ by the
hypothesis of the joint induction at level zero. Summing these relations,
\begin{equation*}
\zeta_{k+1}-\zeta=\sum_{l=1}^{k+1}\bigl(b_l(\zeta)-b_l(0)\bigr).
\end{equation*}
Each summand obeys $|b_l(\zeta)-b_l(0)|\le d_{l-1}(\varrho)|\zeta|$: for $l=k+1$ this is
\eqref{5.3:eq-flow-reference-2}, and for $l\le k$ it is the same bound obtained at the earlier
steps. Hence
\begin{equation}\label{5.3:eq-flow-reference-3}
|\zeta_{k+1}-\zeta|\le\sum_{j\le k}d_j(\varrho)\,|\zeta|\le\tfrac13|\zeta|,
\qquad \zeta\in D(0,\varrho).
\end{equation}

\emph{The increment at zero shift and the reference recursion.} Set $\zeta=0$ in
\eqref{5.3:eq-flow-reference-1}. The term $f_{k+1}$ is bounded by $C_\beta\mathrm{e}_k$, by the
pure-step bound. The remaining term $\mathsf{q}_{k+1}$ is bounded in modulus by $v_k$, as
recorded among the data of the flow lemma above. Therefore
\begin{equation}\label{5.3:eq-flow-reference-4}
|b_{k+1}(0)-\beta^0_{k+1}|\le C_\beta\mathrm{e}_k+v_k\le\lambda,
\end{equation}
where $\lambda=c_0/2$ is the step constant of Definition~\ref{5.1:def-reference-recursion} (not to
be confused with $\underline{\lambda}=c_0/4$), and the last inequality is again one of the
conditions on $g$ collected in the definition of $\gamma_J$, and holds for $g\le\gamma_J$. We define
\begin{equation*}
x_{k+1}:=x_k-b_{k+1}(0).
\end{equation*}
Unwinding this recursion gives, for $i\le k$,
\begin{equation*}
x_i-x_{k+1}=\sum_{l=i+1}^{k+1}b_l(0).
\end{equation*}
We now use the statement on the sums of the one-loop coefficients, in the form
\begin{equation*}
\sum_{l=i+1}^{k+1}\beta^0_l=(k+1-i)c_0+s_i-s_{k+1},\qquad |s_i-s_{k+1}|\le 2c_2,
\end{equation*}
where $s_i$ denotes the remainder of that statement at level $i$. The $k+1-i$ differences
$b_l(0)-\beta^0_l$ are each at most $\lambda$ in modulus: for $l=k+1$ this is
\eqref{5.3:eq-flow-reference-4}, and for $l\le k$ it is the same bound obtained at the earlier
steps. Since $c_0=2\lambda$, the sum of $(k+1-i)c_0$ and of these differences lies in
$[\lambda(k+1-i),3\lambda(k+1-i)]$, and adding $s_i-s_{k+1}$ widens this interval by at most
$2c_2$ on each side. The result is
\begin{equation}\label{5.3:eq-flow-reference-a}
\begin{split}
x_i-x_{k+1}&=(k+1-i)c_0+s_i-s_{k+1}+\sum_{l=i+1}^{k+1}\bigl(b_l(0)-\beta^0_l\bigr)\\
&\in\bigl[\lambda(k+1-i)-2c_2,\ 3\lambda(k+1-i)+2c_2\bigr].
\end{split}
\end{equation}
Together with clause (a) at level $k$, \eqref{5.3:eq-flow-reference-a} gives the reference
comparisons for all pairs of levels at most $k+1$, which is clause (a) of the hypothesis of the
joint induction at level $k+1$.

\emph{The radii at level $k+1$.} With $x_{k+1}$ in hand, the quantities $g_{k+1}$,
$\check{R}_{k+1}$, $\varepsilon_{k+1}$ and $\alpha_{\cdot,k+1}$ are defined; here
$\check{R}_{k+1}=R(\min_{i\le k+1}x_i)$, where $R(x)$ is the least power of $L$ not below
$(\log x)^{r_{\mathrm{pr}}}$, $\varepsilon_{k+1}$ is the radius of the orbit data carried by the
induction format at level $k+1$, and $\alpha_{\cdot,k+1}=\alpha_\cdot(g_{k+1})$ are the
coupling-dependent radii at level $k+1$ of Definition~\ref{5.1:def-analyticity-spaces}. The
case $i=k$ of
\eqref{5.3:eq-flow-reference-a} gives $x_k-x_{k+1}\le 3\lambda+2c_2=b_+$, that is,
$x_{k+1}\ge x_k-b_+$. This gives
\begin{equation*}
\alpha_\cdot(g_{k+1})\le\bigl(1+b_+g^2\bigr)^{1/2}\bigl(1+o(1)\bigr)\,\alpha_\cdot(g_k),
\end{equation*}
where $o(1)$ denotes a quantity tending to zero as $g\to0$, uniformly in $k$,
and hence, for $g\le\gamma_J$,
\begin{equation}\label{5.3:eq-flow-reference-5}
(1+\beta_{\mathrm{pr}})\,\alpha_{\cdot,k+1}\le(1+\beta_s)\,\alpha_{\cdot,k+1}
\le(1+2\beta_s)\,\alpha_\cdot(g_k)=\hat\alpha_k .
\end{equation}
In \eqref{5.3:eq-flow-reference-5}:
\begin{itemize}
\item the first inequality holds because Ba{\l}aban's enlargement of the spaces is by the fraction
$\beta_{\mathrm{pr}}=1/5$, and $\beta_{\mathrm{pr}}\le\beta_s=1/4$, as in part~(ii) of
Lemma~\ref{5.2:lem-enlargement-fractions} (enlargement by one quarter against Ba{\l}aban's one
fifth);
\item the second inequality holds because
\begin{equation*}
(1+\beta_s)\bigl(1+b_+g^2\bigr)^{1/2}\bigl(1+o(1)\bigr)\le 1+2\beta_s
\end{equation*}
once $g$ is below the ceiling $\gamma_J$.
\end{itemize}
Consequently the families bounded on the spaces of radii $\hat\alpha_k$ are bounded, by
restriction, on Ba{\l}aban's enlarged spaces after the operation $\mathbf{T}$ at level $k+1$
\cite[\S2]{B-CMP119}; the later steps of the induction read them there.

The definition of $g_{k+1}$ involves no circularity: $b_{k+1}(0)$ is the value of the functional
$\beta[\cdot]$ on the family $E^{(k+1)}_0$ of the induction format
(Definition~\ref{5.2:def-induction-format}), and $\beta[\cdot]$ is a derivative at the trivial
configuration; so $b_{k+1}(0)$ is determined independently of the radii $\alpha_{\cdot,k+1}$,
which are defined from $x_{k+1}$ only afterwards.

\emph{The background hypothesis at level $k+1$.} The items of the background hypothesis at level
$k+1$ that concern the admissible backgrounds and their comb-gauge data involve the couplings
through two kinds of input only. The
first is $g_j\le\gamma$. The second is the comparisons \cite[(2.6)--(2.9)]{B-CMP119}, whose
constants, printed there $\beta_0$ and $\beta'$, we write $\beta_{\mathrm{cp}}$ and
$\beta'_{\mathrm{cp}}$ so as not to confuse them with the one-loop coefficients $\beta^0_l$:
for $m<n$,
\begin{equation*}
g_m\le(1+\beta_{\mathrm{cp}})\,g_n,\qquad
g_n\le\bigl(1+g_n^2\beta'_{\mathrm{cp}}(n-m)\bigr)^{1/2}g_m .
\end{equation*}
These are exactly the reference comparisons of clause (a) of the hypothesis of the joint
induction, for pairs of levels at most $k+1$. Thus these items use no bound on the couplings of
the asymptotic-freedom kind derived in \cite{B-CMP109} beyond the two inputs just listed.
\end{proof}

\begin{proof}[Joint induction: the local averaging step under sources]
\label{5.3:prf-local-averaging}
We treat the local averaging step from level $k$ to level $k+1$ with $k+1\le K'$, where
$K'=K-m$ is the last step of the local run (Definition~\ref{5.2:def-auxiliary-torus}), on the
physical torus, for a complex source $w\in\mathcal{W}_{k+1}\subset\mathcal{W}_k$.

\emph{The operations.} The step performs the same operations as the frozen averaging step
of the joint induction. This includes the exact refinement of the rim cut-offs
(Definition~\ref{5.2:def-rim-refinement}): its rule, its thresholds and its labels do not
depend on the source coefficients $z_i$ of $w=\sum_i z_i f_i$, so they are the same at every
source. The summation lemma for the refined rim cut-offs, which the frozen averaging step
uses, holds for $w\in\mathcal{W}_k$; there $w$ enters only as a
parameter of the potentials $V$, and the majorants are independent of $w$. At level $k$ the
running shift obeys $|\zeta_k|\le\tfrac83 r$. For $k=0$ the sum in
\eqref{5.2:eq-induction-format-a} is empty, so $\zeta_0=w$ and $|w|<2r$. For $k\ge1$ this is
the bound \eqref{5.3:eq-local-averaging-2} below at level $k$, read with $\sup|w|<2r$ on
$\mathcal{W}_k$, since $\tfrac43\cdot 2r=\tfrac83 r$.
The vertex $\mathfrak{S}_k$ of the step (Definition~\ref{5.2:def-three-families}) is a
function of $g_kC\mathsf{A}$, where $\mathsf{A}$ is the scaled fluctuation field and $C$ the
operator for which the constrained fluctuation form is
$\langle\mathsf{A},C^{*}\Delta^{(k)}C\mathsf{A}\rangle$; it vanishes at $0$ and contains no
factor $g_k^{-2}$. Hence the passage of \cite{B-CMP119} that treats vertex functions of this
kind applies to it.

\emph{Properties of the outputs.} Here $X'$ denotes a localisation domain, and $X'^{+}$ its
enlargement that enters the class $\mathcal{W}_{k+1}(X')$ of
Definition~\ref{5.1:def-weight-classes}. The outputs of the step have three properties.
\begin{itemize}
\item An output term with localisation domain $X'$ is analytic in $w$ on
$\mathcal{W}_{k+1}(X')$.
\item They are local: a potential in $Y\subset X'$ reads $\zeta_k|_Y$ and the old terms in $Y$.
It therefore reads only $w|_{X'^{+}}$.
\item They are frozen at constants: if $w\equiv\hat w$ on $X'^{+}$, every potential in $X'$ is
the frozen one.
\end{itemize}

\emph{The counterterm and the running shift.} Since $|w(p_y)|<2r$, the number
$b_{k+1}(w(p_y))$ is defined. We subtract the sum
\begin{equation*}
\sum_{y} b_{k+1}(w(p_y))\,A\bigl(h^{(k+1)}_y\varphi_{k+1},U_{k+1}\bigr)
\end{equation*}
from the effective action produced by the step, and move it, together with
$+A(\zeta_k,U_{k+1})$, into the weight. The total is unchanged: what is subtracted from the
effective action is added back through the weight. Because $A(c,U)$ is linear in the weight $c$ and
$\sum_y h^{(k+1)}_y=1$, the sum splits into two parts. The first is
$b_{k+1}(0)A(\varphi_{k+1},U_{k+1})$, which does not depend on $w$; it is the part taken up
by the reference recursion $x_{k+1}=x_k-b_{k+1}(0)$ of the flow step. The second is $A$
applied to the weight
\begin{equation*}
\varphi_{k+1}\sum_y h^{(k+1)}_y\bigl[b_{k+1}(w(p_y))-b_{k+1}(0)\bigr],
\end{equation*}
and it joins $A(\zeta_k,U_{k+1})$.
The result is the definition \eqref{5.2:eq-induction-format-a} of the induction format
(Definition~\ref{5.2:def-induction-format}) at level $k+1$. Comparing that definition at
levels $k+1$ and $k$, this reads
\begin{equation}\label{5.3:eq-local-averaging-1}
\zeta_{k+1}(p)=\zeta_k(p)
+\varphi_{k+1}(p)\sum_{y\in T^{(k+1)}}h^{(k+1)}_y(p)\bigl[b_{k+1}(w(p_y))-b_{k+1}(0)\bigr].
\end{equation}

\emph{The bound on the running shift.} This is the one place in the local run where the
increment bound of the flow step is used, at the levels $j\le k+1$. In the flow step of the
joint induction, on the coupling increment and the reference coupling, the increments satisfy
$|b_j(\zeta)-b_j(0)|\le d_{j-1}(\varrho)|\zeta|$ on $D(0,\varrho)$ for every
$\varrho\in[r,2r]$. We take $\varrho=2r$; this is admissible since $|w(p_y)|<2r$.
Inserting these bounds into the definition \eqref{5.2:eq-induction-format-a} at level $k+1$
gives the first inequality in
\begin{equation}\label{5.3:eq-local-averaging-2}
\begin{split}
|\zeta_{k+1}(p)|
&\le |w(p)|+\sum_{j\le k+1}\varphi_j(p)\sum_{y}h^{(j)}_y(p)\,d_{j-1}(2r)\,|w(p_y)|\\
&\le \tfrac43\sup|w| .
\end{split}
\end{equation}
The first inequality uses $\varphi_j\ge0$ and $h^{(j)}_y\ge0$. The second uses
$0\le\varphi_j\le1$ for the cut-off functions $\varphi_j$ of the induction format,
$\sum_y h^{(j)}_y=1$, and
\begin{equation*}
\sum_{j\le k+1}d_{j-1}(2r)\le\tfrac13 ,
\end{equation*}
which is the bound on $\sum_{j\le k}d_j(\varrho)$ proved in that same flow step, read at
$\varrho=2r$.

\emph{The bound on the coupling increments.} The bound $|b_j|\le b_++\tfrac23 r$ is taken
from the flow step of the joint induction, once that step has been carried
out. In particular no Cauchy estimate at the radii of
\cite[(2.28)]{B-CMP119} is involved.

\emph{The difference term.} Set $D^{(k+1)}$ equal to the local potential minus the frozen
one, as in the induction format (Definition~\ref{5.2:def-induction-format}). At each later
level, the freezing lemma, which bounds a local potential minus the frozen one, gives the
bound on $D^{(k+1)}$ directly from the a priori bound $E_0$ that the induction format
(Definition~\ref{5.2:def-induction-format}) carries for the potentials. This bound is not
derived from the bound at the preceding level, so nothing compounds
from level to level.
\end{proof}

\begin{definition}[Notation for the coupling flow and its reference run]\label{5.14:not-flow-notation}
The following notation is used in the remaining statements of this section. Throughout, $g$ is the
reference coupling and $g_0 \in\, ]0,g_{-}(g)]$ the bare coupling, with
$g_{-}(g) = (g^{-2}+B^{\sharp})^{-1/2}$, as in the prefix of Theorem~0, and we write $X := g^{-2}$.

\emph{(a) The flow.} The running inverse square couplings, the running couplings and the first
passage of clause~(0) of Theorem~0 are
\begin{equation}\label{5.14:eq-flow-notation-a}
\begin{gathered}
x_0 := g_0^{-2}, \qquad x_{k+1} := x_k - \beta_{k+1}, \qquad g_k := x_k^{-1/2},\\
K = K(g_0,g) := \min\{k \ge 0 : x_k \le g^{-2} + B^{\sharp}\},
\end{gathered}
\end{equation}
where $\beta_{k+1}$ is the coupling increment of step $k+1$, the coefficient
\cite[(1.22)]{B-CMP109} of the complete regular part of the effective action produced by that step,
large-field remainders included, computed on the whole lattice. For such $g_0$ one has
$x_k > g^{-2} + B^{\sharp}$ for $k < K$, and $x_K$ lies in the interval $]g^{-2}, g^{-2}+B^{\sharp}]$;
hence $x_k > g^{-2}$ and $g_k < g$ for all $k \le K$.

\emph{(b) The one-loop constants.} For $l \ge 1$, $\beta^{0}_{l}$ is the coefficient, in the sense of
\cite[(1.22)]{B-CMP109}, of the Gaussian normaliser of step $l$ of Ba{\l}aban's construction. Let
$C_{\mathfrak{g}} = 2N^2$ be the Casimir constant of $\mathrm{SU}(N)$ in the realisation with
$\operatorname{tr}\mathbf{1} = 1$, and put
\begin{equation*}
c_0 := \frac{11\,C_{\mathfrak{g}}}{24\pi^2}\,\log L = \frac{11N^2}{12\pi^2}\,\log L > 0 .
\end{equation*}
The one-loop lemma provides a constant $c_2 = c_2(L,N) \ge 0$ such that
\begin{equation}\label{5.14:eq-flow-notation-b}
\Bigl|\,\sum_{l=1}^{k}\beta^{0}_{l} - k\,c_0\Bigr| \le c_2 \qquad\text{for every } k \ge 0 ;
\end{equation}
equivalently, for $0 \le i \le j$ the sum $\sum_{l=i+1}^{j}\beta^{0}_{l}$ equals $(j-i)c_0$ plus the
difference of two numbers of modulus at most $c_2$. No sign of an individual coefficient
$\beta^{0}_{l}$ is asserted. Here the coupling is normalised so that its inverse square multiplies the
Wilson action with the trace normalised to one, $g^{-2} = 2N/g_{\mathrm{std}}^{2}$; in standard units
the inverse square of the coupling grows by $11N/(24\pi^2) = \beta_0/(8\pi^2)$ per e-fold of the
scale, $\beta_0 = 11N/3$, and the factor $2N$ between the two coefficients is this normalisation.
Multiplied by $\log L$, the number of e-folds in one step, the coefficient $11N^2/(12\pi^2)$ gives
$c_0$.

\emph{(c) The functions of the correlation theorem.} Let $E^{*}$, $\varepsilon_1$, $\gamma_2$, $A_1$,
the exponent $p_0$, the extraction constant $C_t$, the per-step constant $C_{\beta}$ of the
running-shift statement, the exponent $\kappa_0 \ge 7$ and the coupling ceiling $\gamma_J$ be the
constants of the correlation theorem, Theorem~\ref{5.1:thm-correlation-theorem}; in particular
$\gamma_J \le \gamma^{\circ} := (8c_2+4)^{-1/2}$, hence $X \ge 8c_2 + 4$ for every $g \le \gamma_J$.
For $0 < g' \le \gamma_J$ put
$T(g') := A_1(\log g'^{-2})^{p_0}$, $\mathsf{r} := \varepsilon_1/(2T)$, $\tau := g'/\mathsf{r}$ and
$\vartheta_J := E^{*}e^{-\frac14\gamma_2 T^2}$; the number $\vartheta_J$ is not the shape parameter
$\vartheta$ of the class $\mathfrak{T}(\mathsf{d};\vartheta)$. With $\lambda := c_0/2$,
$y_k := \min_{i\le k}x_i - 2c_2$ and
$C_{\mathfrak{m}} := \tfrac52 + 5/(\lambda(\kappa_0-4))$ put
\begin{equation}\label{5.14:eq-flow-notation-c}
\begin{aligned}
e(g') &:= 2E^{*}\tau + \vartheta_J
= \frac{4E^{*}A_1}{\varepsilon_1}\,g'(\log g'^{-2})^{p_0}\\
&\qquad + E^{*}\exp\Bigl(-\tfrac14\gamma_2A_1^2(\log g'^{-2})^{2p_0}\Bigr),\\
e_k &:= e(g_k), \qquad v_k := 2C_{\beta}C_tC_{\mathfrak{m}}\,y_k^{-(\kappa_0-4)/2}.
\end{aligned}
\end{equation}
By Theorem~\ref{5.1:thm-correlation-theorem}, the functions $\mathsf{r}$, $\tau$, $\vartheta_J$ and $e$
are non-decreasing on $]0,\gamma_J]$ (the function $T$ is not); the sequence $y_k$ is
non-increasing in $k$. The reference run of Definition~\ref{5.1:def-reference-recursion} is: given
$x_0 > X$, $x_{k+1} := x_k - b_{k+1}(0)$, where $b_{k+1}(0)$ is the increment function at zero shift,
the coefficient \cite[(3.61)]{B-CMP119} of the complete regular part of step $k+1$; and $K_J$, the
stopping index of that definition, is the largest $k$ with $x_j > X$ for all $j \le k$. In what
follows the letter $K$ always denotes the first passage \eqref{5.14:eq-flow-notation-a}.

\emph{(d) The flow bounds of the correlation theorem.} For every $g \le \gamma_J$ and $x_0 > g^{-2}$,
Theorem~\ref{5.1:thm-correlation-theorem} gives
\begin{equation}\label{5.14:eq-flow-notation-d}
\begin{aligned}
\lambda(j-i) - 2c_2 \le x_i - x_j &\le 3\lambda(j-i) + 2c_2
&&(0 \le i < j \le K_J),\\
|b_{k+1}(0) - \beta^{0}_{k+1}| &\le C_{\beta}e_k + v_k \le \lambda
&&(0 \le k < K_J),\\
X < x_{K_J} &\le X + b_{+}, \qquad b_{+} := 3\lambda + 2c_2 .&&
\end{aligned}
\end{equation}
None of $K$, $k$, $m$, the torus and the position occurs in any of these constants.
Started at $x_0 = g_0^{-2}$, the reference run is the flow \eqref{5.14:eq-flow-notation-a}, with
$b_{k+1}(0) = \beta_{k+1}$, and $K \le K_J$; in particular the first two lines of
\eqref{5.14:eq-flow-notation-d} hold for all $0 \le i < j \le K$ and all $0 \le k < K$.

\emph{(e) The symbols of the two-point law of clause~(iv) of Theorem~0.} Put
$\bar{x}_s := g^{-2} + B^{\sharp}(s+1)$ and
$\underline{x}_s := g^{-2} + \lambda^{\sharp}s - c_5$, the window of clause~(0) for $x_{K-s}$. Let
$p_0(\cdot)$ be the degree function, $p_0(g') = A_0(\log g'^{-2})^{p_0}$, $A_0$ and the exponent
$p_0$ being the auxiliary parameters of those names in $\mathfrak{p}$, and
$C_W(s) := C_0\,p_0(\bar{x}_s^{-1/2})^2 = C_0A_0^2(\log\bar{x}_s)^{2p_0}$, which satisfies
$C_W(s) \ge 1$ and is non-decreasing in $s$; this function $C_W$, the room function of the two-point
witness theorem, is taken with these two properties from that theorem. Let
$\zeta'_k := \partial_{\zeta}\zeta_k(0)$, which is the number $N_k$ of the correlation theorem,
Theorem~\ref{5.1:thm-correlation-theorem}, let
$b_0 := (N^2-1)/2$, and let $\mathcal{K}_s(f_1,f_2)$ be the free lattice kernel, which satisfies,
for $s = s(\mathsf{d})$ and $(f_1,f_2) \in \mathfrak{T}(\mathsf{d};\vartheta)$, with $\mathsf{d}$
subject to the conditions listed after \eqref{5.14:eq-flow-notation-e},
\begin{equation*}
c_{\mathcal{K}}\|f_1\|_{\infty}\|f_2\|_{\infty} \le \mathcal{K}_s(f_1,f_2)
\le C_{\mathcal{K}}\|f_1\|_{\infty}\|f_2\|_{\infty} .
\end{equation*}
The display of the two-point witness theorem reads
\begin{equation}\label{5.14:eq-flow-notation-e}
\begin{gathered}
s(\mathsf{d}) := \min\{s \ge 0 : C_W(s)L^{-s} \le \mathsf{d}\},\\
S^{T}_{2;K,m}(f_1,f_2) = \zeta'^{\,2}_{K-s}\, b_0\, g_{K-s}^{4}\,\mathcal{K}_s(f_1,f_2)\,(1+R_K),\\
|R_K| \le \tfrac12,\qquad \zeta'_{K-s} \in [\tfrac23,\tfrac43],\\
\tfrac29\, b_0 g_{K-s}^4\mathcal{K}_s(f_1,f_2) \le S^{T}_{2;K,m}(f_1,f_2)
\le \tfrac83\, b_0 g_{K-s}^4\mathcal{K}_s(f_1,f_2),
\end{gathered}
\end{equation}
for $g \le \gamma_W^{\mathbb{T}}$, where $\gamma_W^{\mathbb{T}}$ $(\le\gamma_J)$ is the coupling
ceiling of the two-point witness theorem on the torus, $g_0 \in\, ]0,g_{-}(g)]$, $K = K(g_0,g)$,
every $m$, every $\mathsf{d}$ satisfying the depth floor of the two-point witness theorem and its
ultraviolet floor $K - s(\mathsf{d}) \ge k_W$,
$s := s(\mathsf{d})$, and $(f_1,f_2) \in \mathfrak{T}(\mathsf{d};\vartheta)$, the class of
test-function pairs, whose two-sided separation window is a condition on the test functions;
uniformly in $K$, $g_0$
and the position, and uniformly also in $m$ on
$\{m_{\mathbb{T}}(g_{-}(g)) \le m \le m_V(g)\}$, where $m_V(g)$ is the upper end of the range of
torus exponents of that theorem, beyond which the volume enters only through the number
$s_{\mathrm{vol}}(g,m)$ fixed with that theorem. The number $s(\mathsf{d})$ does not depend on $K$
and tends to infinity as
$\mathsf{d} \downarrow 0$.
\end{definition}

\begin{proof}
\emph{The flow.} By the definition of $K$ as a minimum, $x_k > g^{-2} + B^{\sharp}$ for $k < K$. By the
first-passage statement in clause~(0) of Theorem~0, $g^{-2} < x_K \le g^{-2} + B^{\sharp}$ and
$x_k > g^{-2}$ for every $k \le K$. Hence $x_k > 0$, $g_k = x_k^{-1/2}$ is defined, and $g_k < g$
for all $k \le K$.

\emph{The one-loop constants.} Since $C_{\mathfrak{g}} = 2N^2$, we have
$11C_{\mathfrak{g}}/(24\pi^2) = 11N^2/(12\pi^2)$. For $0 \le i \le j$,
\begin{equation*}
\sum_{l=i+1}^{j}\beta^{0}_{l} - (j-i)c_0
= \Bigl(\sum_{l=1}^{j}\beta^{0}_{l} - jc_0\Bigr) - \Bigl(\sum_{l=1}^{i}\beta^{0}_{l} - ic_0\Bigr),
\end{equation*}
and by \eqref{5.14:eq-flow-notation-b}, applied with $k = j$ and with $k = i$, each bracket has
modulus at most $c_2$; conversely the case $i = 0$ is \eqref{5.14:eq-flow-notation-b}. For the
normalisation: with $\operatorname{tr}\mathbf{1} = 1$ the weight $e^{-g_0^{-2}A(U)}$ gives
$g_0^{-2} = \beta_W = 2N/g_{\mathrm{std}}^2$. In the standard one-loop law the increment of
$g_{\mathrm{std}}^{-2}$ per e-fold of the scale is $11N/(24\pi^2)$, which
equals $2\beta_0/(16\pi^2) = \beta_0/(8\pi^2)$ with $\beta_0 = 11N/3$; multiplied by $2N$ it
becomes $11N^2/(12\pi^2)$, and multiplied further by $\log L$ it becomes $c_0$.

\emph{The function $e$.} Since $\tau = g'/\mathsf{r} = 2Tg'/\varepsilon_1$, one has
$2E^{*}\tau = (4E^{*}A_1/\varepsilon_1)\,g'(\log g'^{-2})^{p_0}$, and
$\vartheta_J = E^{*}\exp(-\frac14\gamma_2A_1^2(\log g'^{-2})^{2p_0})$ because
$T^2 = A_1^2(\log g'^{-2})^{2p_0}$; this is the second expression in
\eqref{5.14:eq-flow-notation-c}. The sequence $y_k$ is non-increasing because the minimum
$\min_{i\le k}x_i$ is taken over a set that grows with $k$. For every $g \le \gamma_J$,
\begin{equation*}
X = g^{-2} \ge \gamma_J^{-2} \ge (\gamma^{\circ})^{-2} = 8c_2 + 4 .
\end{equation*}

\emph{The flow bounds.} The first and third lines of \eqref{5.14:eq-flow-notation-d} are stated in
Theorem~\ref{5.1:thm-correlation-theorem}. For the second line, the running-shift statement writes
$b_{k+1} = \beta^{0}_{k+1} + f_{k+1}(\zeta_0,\dots,\zeta_k) + \mathsf{q}_{k+1}$, where $f_{k+1}$ is
the part of the increment function that is a function of the earlier shifts
$\zeta_0,\dots,\zeta_k$ and $\mathsf{q}_{k+1}$ is its remnant part, with
$|f_{k+1}| \le C_{\beta}e_k$; and $|\mathsf{q}_{k+1}| \le v_k$ on the disc
$D(0,2r) := \{\zeta \in \mathbb{C} : |\zeta| < 2r\}$ of the shift variable by the remnant-value bound
of Theorem~\ref{5.1:thm-correlation-theorem}, where $r > 0$ is the radius of the shift variable in
that theorem; this disc contains $0$. Evaluating at $\zeta = 0$ and applying the triangle inequality
gives $|b_{k+1}(0) - \beta^{0}_{k+1}| \le C_{\beta}e_k + v_k$, and the correlation theorem bounds the
right side by $\lambda$.

\emph{Identification.} The increment $b_{k+1}(0)$ is, by definition, the coupling increment
$\beta_{k+1}$ of \eqref{5.14:eq-flow-notation-a}, both being the coefficient of the complete regular
part of step $k+1$ (\cite[(1.22)]{B-CMP109} and \cite[(3.61)]{B-CMP119} are the same number), so the
reference run started at $x_0 = g_0^{-2}$ is the
sequence $(x_k)$ of \eqref{5.14:eq-flow-notation-a}. Its hypothesis $x_0 > X$ holds: $g_0 \le g_{-}(g)$
with $g_{-}(g) = (g^{-2}+B^{\sharp})^{-1/2}$ means $x_0 \ge g^{-2} + B^{\sharp} > g^{-2}$. Since
$x_j > X$ for every $j \le K$, the definition of $K_J$ gives $K \le K_J$; hence the first two lines of
\eqref{5.14:eq-flow-notation-d} hold for all $0 \le i < j \le K$ and all $0 \le k < K$.

\emph{The symbols of the two-point law of clause~(iv).} Since
$p_0(\bar{x}_s^{-1/2}) = A_0(\log\bar{x}_s)^{p_0}$, the two
expressions for $C_W(s)$ agree. As $\bar{x}_s$ grows linearly in $s$, $C_W(s)$ grows as a power of
$\log s$, so $C_W(s)L^{-s} \to 0$ and the set in the definition of $s(\mathsf{d})$ is non-empty for
every $\mathsf{d} > 0$; $s(\mathsf{d})$ involves only $g$, $B^{\sharp}$, $C_0$, $A_0$, $p_0$ and $L$,
and not $K$. Because $C_W(s) \ge 1$, $C_W(s)L^{-s} \le \mathsf{d}$ forces $L^{-s} \le \mathsf{d}$, so
$s(\mathsf{d}) \ge \log_L(1/\mathsf{d}) \to \infty$ as $\mathsf{d} \downarrow 0$. Finally,
$\zeta'_{K-s} \in [\frac23,\frac43]$ gives $\zeta'^{\,2}_{K-s} \in [\frac49,\frac{16}9]$ and
$|R_K| \le \frac12$ gives $1 + R_K \in [\frac12,\frac32]$; the product lies in $[\frac29,\frac83]$, and
the remaining factors $b_0$, $g_{K-s}^4$, $\mathcal{K}_s(f_1,f_2)$ are non-negative, which yields the
two-sided bound in \eqref{5.14:eq-flow-notation-e}.
\end{proof}

\begin{definition}[Cutoff-free majorants of the per-step defects]\label{5.14:def-defect-majorants}
Let $c_0$, $c_2$ be the constants of the one-loop lemma, and let $\lambda=c_0/2$, $C_{\beta}$,
$C_t$, $C_{\mathfrak{m}}$, $\kappa_0$ and the function $e(g')$ be as in
Notation~\ref{5.14:not-flow-notation}, display~\eqref{5.14:eq-flow-notation-c}. We put
\begin{equation}\label{5.14:eq-defect-majorants-a}
\begin{aligned}
\bar{h}_u &:= (6c_2+4+\lambda u)^{-1/2} && (u\ge 1),\\
\mathfrak{e}_u &:= \sup\{\,e(g') : 0<g'\le \bar{h}_u\,\} && (u\ge 1),\\
\mathfrak{v}_u &:= 2C_{\beta}C_tC_{\mathfrak{m}}\,(4+\lambda u)^{-(\kappa_0-4)/2}
  && (u\ge 1),\\
C_v &:= \sum_{u=1}^{\infty}\mathfrak{v}_u, &&\\
\delta_e(s) &:= \frac{1}{s}\sum_{u=1}^{s}\mathfrak{e}_u && (s\ge 1),\\
\Delta(s) &:= 2c_2 + C_v + C_{\beta}\, s\,\delta_e(s) && (s\ge 1),\\
\Delta(0) &:= 2c_2 + C_v. &&
\end{aligned}
\end{equation}
Thus $\Delta(s)=2c_2+C_v+C_{\beta}\sum_{u=1}^{s}\mathfrak{e}_u$ for every $s\ge 0$, the sum
being empty for $s=0$.
Each of these quantities depends on $(L,N,\mathfrak{p})$ only, through $c_2$, $\lambda$,
$C_{\beta}$, $C_t$, $C_{\mathfrak{m}}$, $\kappa_0$ and the parameters $A_1$,
$\varepsilon_1$, $E^{*}$, $\gamma_2$, $p_0$ that enter $e(\cdot)$. None of them depends on
$K$, $g$, $g_0$, $m$, $s(\mathsf{d})$ or a position on the lattice; they are defined
quantities, not hypotheses, and no numerical value is assigned to any of them. The sequence
$\bar{h}_u$ is strictly decreasing with $\bar{h}_u\le\tfrac12$ and $\bar{h}_u\to 0$; each
$\mathfrak{e}_u$ is finite and $u\mapsto\mathfrak{e}_u$ is non-increasing; and $C_v<\infty$.
The quantity $\delta_e(s)$ is the Ces\`aro mean of $(\mathfrak{e}_u)_{u\ge1}$. That
$\mathfrak{e}_u\to 0$ as $u\to\infty$ and $\delta_e(s)\to 0$ as $s\to\infty$, and hence
$\Delta(s)/s\to 0$, is proved in the proof of the statement on exact one-loop running below.
\end{definition}

\begin{proof}
Since $c_0>0$ we have $\lambda>0$, so $u\mapsto 6c_2+4+\lambda u$ is strictly increasing and
tends to $+\infty$; therefore $\bar{h}_u$ is strictly decreasing and $\bar{h}_u\to 0$. As
$c_2\ge 0$ and $\lambda u>0$, we have $6c_2+4+\lambda u\ge 4$, whence
$\bar{h}_u\le 4^{-1/2}=\tfrac12<1$.

By display~\eqref{5.14:eq-flow-notation-c},
\begin{equation*}
e(g') = \frac{4E^{*}A_1}{\varepsilon_1}\,g'\,(\log g'^{-2})^{p_0}
 + E^{*}\exp\!\Bigl(-\tfrac14\gamma_2A_1^{2}(\log g'^{-2})^{2p_0}\Bigr).
\end{equation*}
For $0<g'<1$ one has $\log g'^{-2}>0$, so both terms are continuous on $]0,1[$. As
$g'\downarrow 0$, $\log g'^{-2}\to+\infty$; the first term tends to $0$ because $g'$ times a
fixed power of $\log g'^{-2}$ tends to $0$, and the second tends to $0$ because its exponent
tends to $-\infty$. Hence $e$ extends to a continuous function on $[0,\bar{h}_u]$, which is a
compact subset of $[0,1[$ because $\bar{h}_u\le\tfrac12$; it is therefore bounded there, and
$\mathfrak{e}_u<\infty$. Since $\bar{h}_{u+1}<\bar{h}_u$, the interval
$]0,\bar{h}_{u+1}]$ is contained in $]0,\bar{h}_u]$, and a supremum over a smaller set is not
larger; thus $\mathfrak{e}_{u+1}\le\mathfrak{e}_u$.

For $C_v$: since $\kappa_0\ge 7$, the exponent satisfies $(\kappa_0-4)/2\ge 3/2>1$. Because
$4+\lambda u\ge 1$, raising it to the larger negative power gives a smaller number, so
\begin{equation*}
\mathfrak{v}_u\le 2C_{\beta}C_tC_{\mathfrak{m}}\,(4+\lambda u)^{-3/2}\qquad(u\ge 1).
\end{equation*}
The function $t\mapsto(4+\lambda t)^{-3/2}$ is decreasing on $[0,\infty[$, so
$(4+\lambda u)^{-3/2}\le\int_{u-1}^{u}(4+\lambda t)^{-3/2}\,dt$ for every $u\ge1$. Summing
over $u\ge1$,
\begin{equation*}
C_v \le 2C_{\beta}C_tC_{\mathfrak{m}}\sum_{u\ge 1}(4+\lambda u)^{-3/2}
 \le 2C_{\beta}C_tC_{\mathfrak{m}}\int_{0}^{\infty}(4+\lambda t)^{-3/2}\,dt
 = \frac{2C_{\beta}C_tC_{\mathfrak{m}}}{\lambda},
\end{equation*}
where the integral equals
$\bigl[-\tfrac{2}{\lambda}(4+\lambda t)^{-1/2}\bigr]_{0}^{\infty}=\tfrac{2}{\lambda}\cdot
\tfrac12=\tfrac1\lambda$. Thus $C_v<\infty$.

Finally, dividing the definition of $\Delta(s)$ by $s$ gives
\begin{equation*}
\frac{\Delta(s)}{s} = \frac{2c_2+C_v}{s} + C_{\beta}\,\delta_e(s),
\end{equation*}
so $\Delta(s)/s\to 0$ as soon as $\delta_e(s)\to 0$, the first term tending to $0$ because
$2c_2+C_v$ is finite.
\end{proof}

\begin{theorem}[Exact one-loop law of the construction's running coupling]\label{5.14:thm-one-loop-law}
Let $N\ge 2$, let the blocking factor $L$, Ba{\l}aban's auxiliary parameters $\mathfrak{p}$ and the
remaining constants be chosen as in the prefix of Theorem~0, and let hypothesis (H7) of Theorem~0 be
in force, so that the reference coupling satisfies $g\le\gamma_J$. The letters are those of
Notation~\ref{5.14:not-flow-notation}, and $\mathfrak{e}_u$, $C_v$, $\Delta$ are the majorants of
Definition~\ref{5.14:def-defect-majorants}. Assume:
\begin{enumerate}
\item[(a)] the flow statement of the correlation theorem, that is, the bounds
\eqref{5.14:eq-flow-notation-d}: for $N=2$ as part of Theorem~\ref{5.1:thm-correlation-theorem}, and
for $N\ge 3$ as part of Proposition~\ref{5.13:prop-sun-correlation};
\item[(b)] the one-loop lemma, in the form \eqref{5.14:eq-flow-notation-b} recorded in
Notation~\ref{5.14:not-flow-notation}.
\end{enumerate}
Then for every $g\le\gamma_J$, every $g_0\in\left]0,g_{-}(g)\right]$ with $K=K(g_0,g)$, every $m$ of
the family of tori, and all $0\le i<j\le K$,
\begin{equation}\label{5.14:eq-one-loop-law-a}
\bigl|x_i-x_j-(j-i)\,c_0\bigr|\le 2c_2+\sum_{k=i}^{j-1}\bigl[C_\beta\,e(g_k)+v_k\bigr].
\end{equation}
In particular, for every $0\le s\le K$,
\begin{align}
\bigl|x_{K-s}-x_K-s\,c_0\bigr|
&\le 2c_2+\sum_{u=1}^{s}\bigl[C_\beta\,e(g_{K-u})+v_{K-u}\bigr]\notag\\
&\le \Delta(s)=2c_2+C_v+C_\beta\sum_{u=1}^{s}\mathfrak{e}_u,
\label{5.14:eq-one-loop-law-b}
\end{align}
and $\Delta(s)/s\to 0$ as $s\to\infty$. Here
\begin{equation}\label{5.14:eq-one-loop-law-1}
c_0=\frac{11N^2}{12\pi^2}\,\log L .
\end{equation}
The majorant $\Delta$ depends on $(L,N,\mathfrak{p})$ only; the convergence $\Delta(s)/s\to 0$ is
therefore uniform in $K\ge s$, in $g\le\gamma_J$, in $g_0$, in $m$ and in the position of the
sources, on which no constant of the correlation theorem depends. Consequently, along every
subsequence of the corollary to clause~(0) of
Theorem~0, that is, for $g_0^{(i)}\downarrow 0$ with $K_i:=K(g_0^{(i)},g)\ge s$ and
$g_{K_i-s}\to g_{\infty,s}$, one has
\begin{equation}\label{5.14:eq-one-loop-law-c}
g_{\infty,s}^{-2}\in\bigl[\,g^{-2}+s\,c_0-\Delta(s),\;g^{-2}+B^{\sharp}+s\,c_0+\Delta(s)\,\bigr],
\end{equation}
so that $g_{\infty,s}^{-2}/(s\,c_0)\to 1$ as $s\to\infty$, and $(x_{K-s}-x_K)/s\to c_0$ as
$s\to\infty$, uniformly in $K\ge s$.
\end{theorem}

In this statement $x_k=g_k^{-2}$ is the inverse square of the running coupling of clause~(0) of
Theorem~0, as
fixed in \eqref{5.14:eq-flow-notation-a}; $c_0$ and $c_2=c_2(L,N)$ are the constants of the one-loop
lemma; $e(\cdot)$ is the function of the correlation theorem that tends to zero with the coupling, and
$v_k$ is its summable remnant bound; $C_\beta$ is the per-step constant of the running-shift
statement; $\mathfrak{e}_u$, $C_v$ and $\Delta(s)$ are the majorants of
Definition~\ref{5.14:def-defect-majorants}, \eqref{5.14:eq-defect-majorants-a}, none of which depends
on $K$, $g$, $g_0$, $m$ or the position.

Here the coupling is normalised so that its inverse square multiplies the Wilson action with the
trace normalised to one, $g^{-2}=2N/g_{\mathrm{std}}^{2}$, where $g_{\mathrm{std}}$ denotes the
coupling in the standard normalisation; in standard units the inverse square of
the coupling grows by $11N/(24\pi^2)=\beta_0/(8\pi^2)$ per e-fold of the scale, $\beta_0=11N/3$, and
the factor $2N$ between the two coefficients is this normalisation.

In words: the inverse square of the running coupling grows by $\frac{11N^2}{12\pi^2}\log L$ per
renormalisation step, up to an error whose Gaussian part is bounded by $2c_2$ cumulatively over any
number of steps and whose remaining part $C_\beta e(g_k)+v_k$ vanishes with the coupling; cumulatively
over $s$ steps the error is $o(s)$, uniformly in the cutoff and in the volume. This is the one-loop
law of $\mathrm{SU}(N)$ Yang--Mills theory with its
universal first coefficient ($11N/3$ in the standard normalisation), obtained here as a
consequence of the correlation theorem's control of the flow, not assumed. The bounds
\eqref{5.14:eq-one-loop-law-a} and \eqref{5.14:eq-one-loop-law-b} carry no range of $m$ and no volume
floor beyond the prefix of Theorem~0: they form a volume-free flow statement about the sequence of
couplings of the construction, and as such they hold verbatim for the coupling sequences of the
infinite-volume family of clause~(ii) of Theorem~0.
The proof rests on hypotheses~(a) and~(b) and on the identification of the correlation theorem's
coupling increment with that of the flow \eqref{5.14:eq-flow-notation-a}, both being
Ba{\l}aban's coefficient of the complete regular part (Notation~\ref{5.14:not-flow-notation}).

The theorem does not assert: a sign, or the one-loop value, of an individual increment $\beta_{k+1}$
or of an individual coefficient $\beta^{0}_{l}$; convergence of $g_{K-s}$ along the first passage;
any rate for the $o(s)$; anything about the Schwinger functions; nor, for $N\ge 3$, the flow
statement of Theorem~\ref{5.1:thm-correlation-theorem}, which for those $N$ enters through
hypothesis~(a) in the form of Proposition~\ref{5.13:prop-sun-correlation}.

\begin{proof}[Proof of Theorem~\ref{5.14:thm-one-loop-law}]
\emph{Step 0: one sequence, and the range of the flow statement.} Under the identification, recorded
in Notation~\ref{5.14:not-flow-notation}, of the correlation theorem's coupling increment with that of
the flow \eqref{5.14:eq-flow-notation-a}, both being Ba{\l}aban's coefficient of the complete regular
part, for every $k\le K$ the number $x_k$ of
\eqref{5.14:eq-flow-notation-a} is the reference value of the correlation theorem; the increment
function of the correlation theorem at zero shift satisfies $b_{k+1}(0)=\beta_{k+1}$; the coefficients
$\beta^{0}_{l}$ of the correlation theorem are the coefficients $\beta^{0}_{l}$ of the Gaussian
normalisers of the steps to which the one-loop lemma applies, and the one-loop bound used in the
correlation theorem is the one-loop lemma \eqref{5.14:eq-flow-notation-b}. Moreover $K\le K_J$, where
$K_J$ is the stopping index of the correlation theorem, that is, the largest $k$ with $x_{k'}>g^{-2}$
for all $k'\le k$: indeed $x_k>g^{-2}$ for every $k\le K$ by the first-passage lemma of clause~(0) of
Theorem~0. The hypotheses under which the correlation theorem
asserts the bounds \eqref{5.14:eq-flow-notation-d} are met: $g\le\gamma_J$ by (H7); and
$x_0=g_0^{-2}\ge g^{-2}+B^{\sharp}>g^{-2}$ by the choice $g_0\le g_{-}(g)$ (hypothesis (H6) of
Theorem~0), so the initial value exceeds the reference value $g^{-2}$, as
\eqref{5.14:eq-flow-notation-d} requires. Moreover no constant
in \eqref{5.14:eq-flow-notation-d} depends on $K$, $k$, $m$, the torus or the position; hence
hypothesis (H5) of Theorem~0 on the torus plays no role here, and $m$ enters nothing below. Since
$K\le K_J$, the
two-sided linear flow window in \eqref{5.14:eq-flow-notation-d} holds for all pairs $0\le i<j\le K$,
and the closeness bound in \eqref{5.14:eq-flow-notation-d},
\begin{equation}\label{5.14:eq-one-loop-law-2}
\bigl|b_{k+1}(0)-\beta^{0}_{k+1}\bigr|\le C_\beta\,e(g_k)+v_k,
\end{equation}
holds for all $0\le k\le K-1$.

\emph{Step 1: telescoping.} Fix $0\le i<j\le K$. By \eqref{5.14:eq-flow-notation-a} and Step~0,
$x_{l-1}-x_l=\beta_l=b_l(0)$ for $1\le l\le K$. Summing over $l=i+1,\dots,j$ and adding and
subtracting $\beta^{0}_{l}$ in each term,
\begin{align}
x_i-x_j&=\sum_{l=i+1}^{j}\bigl(x_{l-1}-x_l\bigr)=\sum_{l=i+1}^{j}b_l(0)\notag\\
&=\sum_{l=i+1}^{j}\beta^{0}_{l}+\sum_{l=i+1}^{j}\bigl(b_l(0)-\beta^{0}_{l}\bigr).
\label{5.14:eq-one-loop-law-3}
\end{align}

\emph{Step 2: the one-loop lemma, used at both ends.} The one-loop lemma
\eqref{5.14:eq-flow-notation-b}, applied at $k=j$ and at $k=i$, bounds each of
$\bigl|\sum_{l=1}^{j}\beta^{0}_{l}-j\,c_0\bigr|$ and $\bigl|\sum_{l=1}^{i}\beta^{0}_{l}-i\,c_0\bigr|$
by $c_2$. Writing the partial sum from $i+1$ to $j$ as the difference of the partial sums from $1$ to
$j$ and from $1$ to $i$, and using the triangle inequality,
\begin{align}
\Bigl|\sum_{l=i+1}^{j}\beta^{0}_{l}-(j-i)\,c_0\Bigr|
&=\Bigl|\Bigl(\sum_{l=1}^{j}\beta^{0}_{l}-j\,c_0\Bigr)
-\Bigl(\sum_{l=1}^{i}\beta^{0}_{l}-i\,c_0\Bigr)\Bigr|\notag\\
&\le c_2+c_2=2c_2.
\label{5.14:eq-one-loop-law-4}
\end{align}

\emph{Step 3: the closeness bound, summed.} For each $l$ with $i+1\le l\le j$ put $k:=l-1$; then
$k\in\{i,\dots,j-1\}\subset\{0,\dots,K-1\}$, so \eqref{5.14:eq-one-loop-law-2} applies at this $k$ by
Step~0 and gives
\begin{equation*}
\bigl|b_l(0)-\beta^{0}_{l}\bigr|\le C_\beta\,e(g_{l-1})+v_{l-1}.
\end{equation*}
Subtract $(j-i)c_0$ from both sides of \eqref{5.14:eq-one-loop-law-3}, apply the triangle inequality,
bound the first sum by \eqref{5.14:eq-one-loop-law-4} and each term of the second sum by the last
display, and re-index the second sum by $k=l-1$:
\begin{align}
\bigl|x_i-x_j-(j-i)\,c_0\bigr|
&\le 2c_2+\sum_{l=i+1}^{j}\bigl|b_l(0)-\beta^{0}_{l}\bigr|\notag\\
&\le 2c_2+\sum_{k=i}^{j-1}\bigl[C_\beta\,e(g_k)+v_k\bigr].
\label{5.14:eq-one-loop-law-5}
\end{align}
This is \eqref{5.14:eq-one-loop-law-a}.

Now let $1\le s\le K$ and take $i=K-s$, $j=K$ in \eqref{5.14:eq-one-loop-law-a}, and substitute
$u:=K-k$, so that $k=K-u$ and $k$ runs over $K-s,\dots,K-1$ exactly when $u$ runs over $1,\dots,s$.
The defects summed are those of the $s$ steps from level $K-s$ to level $K$. This gives the first
inequality of \eqref{5.14:eq-one-loop-law-b}. For $s=0$ the left side of
\eqref{5.14:eq-one-loop-law-b} vanishes, the sum on the right is empty, and the first inequality reads
$0\le 2c_2$, which holds because $c_2$ bounds an absolute value in \eqref{5.14:eq-flow-notation-b}.

It remains to prove the second inequality in \eqref{5.14:eq-one-loop-law-b}, the limit
$\Delta(s)/s\to 0$ with its uniformity, and the consequence \eqref{5.14:eq-one-loop-law-c}. For $s=0$
the three members of \eqref{5.14:eq-one-loop-law-b} read $0\le 2c_2\le 2c_2+C_v$, since the sums are
empty and $C_v\ge 0$ as a sum of majorants of the non-negative numbers $v_k$. Throughout what follows,
$g\le\gamma_J$, $g_0\in\left]0,g_{-}(g)\right]$, $K=K(g_0,g)$, and $1\le s\le K$.

We use the lower member of the two-sided linear flow window in \eqref{5.14:eq-flow-notation-d}. It reads
\begin{equation}\label{5.14:eq-one-loop-law-6}
x_i-x_j\ge\lambda\,(j-i)-2c_2,\qquad \lambda=\frac{c_0}{2},
\end{equation}
and by Step~0 it holds for all $0\le i<j\le K$. Since $L>1$, \eqref{5.14:eq-one-loop-law-1} gives
$c_0>0$, hence $\lambda>0$. We also use two facts from clause~(0) of Theorem~0 and one from the
correlation theorem. First, $K$ is the first index $k\ge 0$ with $x_k\le g^{-2}+B^{\sharp}$, so
$x_K\le g^{-2}+B^{\sharp}$. Second, $x_K>g^{-2}$ by the first-passage lemma of clause~(0), as in Step~0.
Third, the coupling ceiling of the correlation theorem satisfies $\gamma_J\le(8c_2+4)^{-1/2}$, so
$g\le\gamma_J$ gives
\begin{equation}\label{5.14:eq-one-loop-law-7}
g^{-2}\ge 8c_2+4.
\end{equation}

\emph{Step 4: the couplings along the last $s$ steps.} Let $1\le u\le s$ and $0\le i\le K-u$. Then
$i<K$, and \eqref{5.14:eq-one-loop-law-6} applies to the pair $(i,K)$. Since $K-i\ge u$ and
$\lambda>0$, we obtain $x_i-x_K\ge\lambda(K-i)-2c_2\ge\lambda u-2c_2$. Together with $x_K>g^{-2}$ and
\eqref{5.14:eq-one-loop-law-7},
\begin{equation}\label{5.14:eq-one-loop-law-8}
x_i>g^{-2}+\lambda u-2c_2\ge 6c_2+4+\lambda u\qquad(0\le i\le K-u).
\end{equation}
For $i=K-u$, \eqref{5.14:eq-one-loop-law-8} says
$g_{K-u}^{-2}>6c_2+4+\lambda u=\bar{h}_u^{-2}$. Hence $0<g_{K-u}<\bar{h}_u$, where $\bar{h}_u$ is as in
Definition~\ref{5.14:def-defect-majorants}. Since $\mathfrak{e}_u$ is the supremum of $e(g')$ over
$0<g'\le\bar{h}_u$ (Definition~\ref{5.14:def-defect-majorants}, \eqref{5.14:eq-defect-majorants-a}),
\begin{equation}\label{5.14:eq-one-loop-law-9}
e(g_{K-u})\le\mathfrak{e}_u\qquad(1\le u\le s).
\end{equation}

\emph{Step 5: the remnant bounds.} The remnant bound of the correlation theorem is
$v_k=2C_\beta C_tC_{\mathfrak{m}}\,y_k^{-(\kappa_0-4)/2}$, where
$y_k=\min_{i\le k}x_i-2c_2$ and $C_{\mathfrak{m}}=5/2+5/(\lambda(\kappa_0-4))$. Take $k=K-u$ with
$1\le u\le s$. Taking the minimum over $0\le i\le K-u$ in \eqref{5.14:eq-one-loop-law-8} gives
\begin{equation*}
y_{K-u}>6c_2+4+\lambda u-2c_2=4c_2+4+\lambda u>0 .
\end{equation*}
The exponent satisfies $(\kappa_0-4)/2\ge 3/2>0$ because $\kappa_0\ge 7$, so
\begin{equation}\label{5.14:eq-one-loop-law-10}
v_{K-u}\le 2C_\beta C_tC_{\mathfrak{m}}\,\bigl(4c_2+4+\lambda u\bigr)^{-(\kappa_0-4)/2}
\qquad(1\le u\le s).
\end{equation}
The right side of \eqref{5.14:eq-one-loop-law-10} is built from $(L,N,\mathfrak{p})$ only. Because the
exponent is at least $3/2>1$, its sum over $u\ge 1$ converges. By
Definition~\ref{5.14:def-defect-majorants}, $\mathfrak{v}_u$ is the majorant of $v_{K-u}$ obtained
through this lower bound on $y_{K-u}$; hence $v_{K-u}\le\mathfrak{v}_u$ for $1\le u\le s$, and
$C_v=\sum_{u\ge 1}\mathfrak{v}_u$. Therefore
\begin{equation}\label{5.14:eq-one-loop-law-11}
\sum_{u=1}^{s}v_{K-u}\le\sum_{u\ge 1}\mathfrak{v}_u=C_v .
\end{equation}

\emph{Step 6: the second inequality in \eqref{5.14:eq-one-loop-law-b}.} Multiply
\eqref{5.14:eq-one-loop-law-9} by $C_\beta\ge 0$ and sum over $u=1,\dots,s$. Add
\eqref{5.14:eq-one-loop-law-11}. This gives
\begin{equation*}
2c_2+\sum_{u=1}^{s}\bigl[C_\beta\,e(g_{K-u})+v_{K-u}\bigr]
\le 2c_2+C_v+C_\beta\sum_{u=1}^{s}\mathfrak{e}_u=\Delta(s),
\end{equation*}
which is the second inequality in \eqref{5.14:eq-one-loop-law-b}.

\emph{Step 7: $\Delta(s)/s\to 0$, uniformly.} By Definition~\ref{5.14:def-defect-majorants},
$\bar{h}_u\downarrow 0$ as $u\to\infty$; and $e(g')\to 0$ as $g'\downarrow 0$, a property of the
correlation theorem's function $e$ (hypothesis~(a)). Let $\varepsilon>0$.
Choose $\eta>0$ with $e(g')\le\varepsilon$ for $0<g'\le\eta$, and then $u_0$ with
$\bar{h}_u\le\eta$ for $u\ge u_0$. For $u\ge u_0$ the supremum defining $\mathfrak{e}_u$ is taken over
a subset of $\left]0,\eta\right]$, so $0\le\mathfrak{e}_u\le\varepsilon$. Thus $\mathfrak{e}_u\to 0$.
For $s\ge u_0$,
\begin{equation*}
\frac1s\sum_{u=1}^{s}\mathfrak{e}_u\le\frac1s\sum_{u=1}^{u_0-1}\mathfrak{e}_u+\varepsilon ,
\end{equation*}
and the first term on the right tends to $0$ as $s\to\infty$, so
\begin{equation*}
\limsup_{s\to\infty}\frac1s\sum_{u=1}^{s}\mathfrak{e}_u\le\varepsilon .
\end{equation*}
Since $\varepsilon>0$ was arbitrary, the Ces\`aro means of $(\mathfrak{e}_u)$ tend to $0$. Therefore
\begin{equation*}
\frac{\Delta(s)}{s}=\frac{2c_2+C_v}{s}+C_\beta\cdot\frac1s\sum_{u=1}^{s}\mathfrak{e}_u\longrightarrow 0
\qquad(s\to\infty).
\end{equation*}
Every quantity in $\Delta(s)$, namely $c_2$, $C_v$, $C_\beta$ and the numbers $\mathfrak{e}_u$, depends
on $(L,N,\mathfrak{p})$ only. So $\Delta(s)$, and with it the convergence $\Delta(s)/s\to 0$, is the
same for every $K\ge s$, every $g\le\gamma_J$, every $g_0$, every $m$ and every position.

\emph{Step 8: the consequences.} Let $K\ge s$. Dividing \eqref{5.14:eq-one-loop-law-b} by $s\ge 1$ gives
\begin{equation*}
\Bigl|\frac{x_{K-s}-x_K}{s}-c_0\Bigr|\le\frac{\Delta(s)}{s}.
\end{equation*}
By Step~7 the right side tends to $0$ independently of $K$. Hence $(x_{K-s}-x_K)/s\to c_0$ as
$s\to\infty$, uniformly in $K\ge s$.

Now let $g_0^{(i)}\downarrow 0$ with $K_i:=K(g_0^{(i)},g)\ge s$ and $g_{K_i-s}\to g_{\infty,s}$ along a
subsequence of the corollary to clause~(0). Apply \eqref{5.14:eq-one-loop-law-b} at $g_0=g_0^{(i)}$
and combine it with $g^{-2}<x_{K_i}\le g^{-2}+B^{\sharp}$. This gives
\begin{equation*}
g^{-2}+s\,c_0-\Delta(s)\le x_{K_i-s}=g_{K_i-s}^{-2}\le g^{-2}+B^{\sharp}+s\,c_0+\Delta(s).
\end{equation*}
The upper bound keeps $g_{K_i-s}$ above a positive number independent of $i$, so $g_{\infty,s}>0$ and
$g_{K_i-s}^{-2}\to g_{\infty,s}^{-2}$. The closed interval above does not depend on $i$, so it contains
the limit. This is \eqref{5.14:eq-one-loop-law-c}. Dividing \eqref{5.14:eq-one-loop-law-c} by
$s\,c_0>0$, the endpoints are
\begin{equation*}
1+\frac{g^{-2}}{s\,c_0}-\frac{\Delta(s)}{s\,c_0}
\qquad\text{and}\qquad
1+\frac{g^{-2}+B^{\sharp}}{s\,c_0}+\frac{\Delta(s)}{s\,c_0};
\end{equation*}
since $g$ and $B^{\sharp}$ do not depend on $s$, both tend to $1$ by Step~7. Hence
$g_{\infty,s}^{-2}/(s\,c_0)\to 1$ as $s\to\infty$.
\end{proof}

\begin{proof}[Small late couplings and vanishing mean defect]
\label{5.14:prf-mean-defect}
In this proof the steps for Theorem~\ref{5.14:thm-one-loop-law} are carried out in full detail.
The first inequality of \eqref{5.14:eq-one-loop-law-b}, proved above, is taken as given. We bound its right-hand side by $\Delta(s)$ and show that $\Delta(s)/s\to 0$.
This takes four steps:
\begin{enumerate}
\item an a priori bound on the couplings met at depth $u$ below the final level $K$;
\item summability of the remnant defects $v_{K-u}$;
\item a Ces\`aro argument for the small-field defects $C_{\beta}\,e(g_{K-u})$;
\item assembly of the bound.
\end{enumerate}
Throughout, $g\le\gamma_J$, $g_0\in\left]0,g_{-}(g)\right]$, $K=K(g_0,g)$, and the torus is any member of the family.
The letters $x_k$, $g_k=x_k^{-1/2}$, $\lambda=c_0/2$, $y_k$, $v_k$, $C_{\mathfrak{m}}$ and the function $e(\cdot)$ are those of Notation~\ref{5.14:not-flow-notation}.
The majorants $\bar{h}_u$, $\mathfrak{e}_u$, $\mathfrak{v}_u$, $C_v$, $\delta_e(s)$ and $\Delta(s)$ are those of Definition~\ref{5.14:def-defect-majorants}, display \eqref{5.14:eq-defect-majorants-a}.
As in Definition~\ref{5.1:def-reference-recursion}, $X=g^{-2}$.

\medskip
\emph{A priori size of the couplings met at depth $u$ below the final level $K$.}
Let $1\le u\le K$.

We first collect three facts.
\begin{itemize}
\item The correlation theorem (Theorem~\ref{5.1:thm-correlation-theorem}) provides the two-sided linear flow window displayed in \eqref{5.14:eq-flow-notation-d}. Its lower member, applied with $(i,j)=(K-u,K)$ (admissible since $K\le K_J$, where $K_J$ is the correlation theorem's own stopping index, Notation~\ref{5.14:not-flow-notation}), gives
\begin{equation*}
x_{K-u}\ \ge\ x_K+\lambda u-2c_2 .
\end{equation*}
\item By the definition of the first passage in \eqref{5.14:eq-flow-notation-a}, $x_K>X$.
\item Since $g\le\gamma_J$ and $\gamma_J\le\gamma^{\circ}$, where $\gamma^{\circ}=(8c_2+4)^{-1/2}$, display \eqref{5.14:eq-flow-notation-c} gives
\begin{equation*}
X=g^{-2}\ \ge\ \gamma_J^{-2}\ \ge\ (\gamma^{\circ})^{-2}=8c_2+4 .
\end{equation*}
\end{itemize}
Combining the three, and using $X-2c_2\ge 6c_2+4$, we obtain
\begin{equation}\label{5.14:eq-mean-defect-a}
\begin{gathered}
x_{K-u}\ >\ X+\lambda u-2c_2\ \ge\ 6c_2+4+\lambda u,\\
\text{hence}\quad g_{K-u}=x_{K-u}^{-1/2}\ <\ \bar{h}_u=(6c_2+4+\lambda u)^{-1/2};
\quad\text{also}\quad g_{K-u}<g\le\gamma_J .
\end{gathered}
\end{equation}
The last inequality holds by the definition of the first passage in \eqref{5.14:eq-flow-notation-a}.

Next let $0\le i\le K-u$. Then $i<K$, and the same lower member, applied with $(i,K)$, gives
\begin{equation*}
x_i\ \ge\ x_K+\lambda(K-i)-2c_2 .
\end{equation*}
Since $K-i\ge u$ and $\lambda>0$, and since $x_K>X$, this yields
\begin{equation*}
x_i\ \ge\ x_K+\lambda u-2c_2\ >\ X+\lambda u-2c_2 .
\end{equation*}
This holds for every $i\le K-u$. Taking the minimum over $i$ and subtracting $2c_2$ gives the second bound:
\begin{equation}\label{5.14:eq-mean-defect-b}
y_{K-u}=\min_{i\le K-u}x_i-2c_2\ >\ X+\lambda u-4c_2\ \ge\ 4c_2+4+\lambda u\ \ge\ 4+\lambda u\ >\ 0 .
\end{equation}
Here the second inequality uses $X\ge 8c_2+4$, and the third uses $c_2\ge 0$.

\medskip
\emph{The remnant defects are summable.}
By \eqref{5.14:eq-flow-notation-c}, $v_k=2C_{\beta}C_tC_{\mathfrak{m}}\,y_k^{-(\kappa_0-4)/2}$.
Since $\kappa_0\ge 7$, the exponent satisfies $(\kappa_0-4)/2\ge 3/2$.
Hence $v_k$ is a decreasing function of $y_k$ on $\left]0,\infty\right[$.
By \eqref{5.14:eq-mean-defect-b}, $y_{K-u}\ge 4+\lambda u$, and $4+\lambda u\ge 1$.
Therefore, for every $1\le u\le K$,
\begin{equation}\label{5.14:eq-mean-defect-1}
v_{K-u}=2C_{\beta}C_tC_{\mathfrak{m}}\,y_{K-u}^{-(\kappa_0-4)/2}
\ \le\ 2C_{\beta}C_tC_{\mathfrak{m}}\,(4+\lambda u)^{-(\kappa_0-4)/2}=\mathfrak{v}_u .
\end{equation}
The terms $\mathfrak{v}_u$ are non-negative. Since $(\kappa_0-4)/2>1$, their series converges to the finite constant $C_v$ of \eqref{5.14:eq-defect-majorants-a}.
Summing \eqref{5.14:eq-mean-defect-1} over $u$ gives, for every $s\le K$,
\begin{equation}\label{5.14:eq-mean-defect-2}
\sum_{u=1}^{s}v_{K-u}\ \le\ \sum_{u=1}^{\infty}\mathfrak{v}_u\ =\ C_v\ <\ \infty .
\end{equation}
By Definition~\ref{5.14:def-defect-majorants}, $C_v$ depends on $(L,N,\mathfrak{p})$ only.

\medskip
\emph{The small-field defects grow more slowly than $s$.}

\emph{(a) Pointwise bound.}
By \eqref{5.14:eq-mean-defect-a}, $g_{K-u}\in\left]0,\bar{h}_u\right]$.
Since $\mathfrak{e}_u$ is the supremum of $e$ over this interval,
\begin{equation}\label{5.14:eq-mean-defect-3}
e(g_{K-u})\ \le\ \mathfrak{e}_u\qquad\text{for every }1\le u\le K .
\end{equation}

\emph{(b) The sequence $(\mathfrak{e}_u)_{u\ge 1}$ is non-increasing.}
The sequence $\bar{h}_u$ decreases in $u$, so the intervals $\left]0,\bar{h}_u\right]$ shrink as $u$ grows. A supremum over a smaller set is no larger.

\emph{(c) The sequence $(\mathfrak{e}_u)_{u\ge 1}$ tends to $0$.}
By \eqref{5.14:eq-flow-notation-c},
\begin{equation*}
e(g')=\frac{4E^{*}A_1}{\varepsilon_1}\,g'\,(\log g'^{-2})^{p_0}
+E^{*}\exp\Bigl(-\tfrac14\gamma_2A_1^2(\log g'^{-2})^{2p_0}\Bigr).
\end{equation*}
As $g'\downarrow 0$ we have $\log g'^{-2}\to+\infty$, and both members tend to $0$:
\begin{itemize}
\item the first, because a power of $\log g'^{-2}$ grows more slowly than $g'^{-1}$;
\item the second, because its exponent tends to $-\infty$.
\end{itemize}
Hence, for every $\eta>0$, there is $g_\eta>0$ with $e\le\eta$ on $\left]0,g_\eta\right]$.
Since $\bar{h}_u\downarrow 0$, there is $U$ with $\bar{h}_U\le g_\eta$.
For every $u\ge U$ we then have $\left]0,\bar{h}_u\right]\subset\left]0,g_\eta\right]$, and so $\mathfrak{e}_u\le\eta$.
By (b) and what has just been shown, $(\mathfrak{e}_u)_{u\ge 1}$ is non-increasing with limit $0$; hence $\mathfrak{e}_u\ge 0$ for every $u\ge 1$.

\emph{(d) The Ces\`aro mean tends to $0$.}
Let $\eta>0$. Since the sequence is null, choose $U=U(\eta)\ge 1$ with $\mathfrak{e}_u\le\eta/2$ for all $u\ge U$.
By (b), $\mathfrak{e}_1=\max_u\mathfrak{e}_u$, and for $u\ge U+1$ we have $\mathfrak{e}_u\le\mathfrak{e}_U$.
Hence, for every $s\ge U$,
\begin{equation}\label{5.14:eq-mean-defect-4}
\begin{aligned}
\delta_e(s)&=\frac1s\sum_{u=1}^{s}\mathfrak{e}_u
=\frac1s\sum_{u=1}^{U}\mathfrak{e}_u+\frac1s\sum_{u=U+1}^{s}\mathfrak{e}_u\\
&\le\frac{U\,\mathfrak{e}_1}{s}+\mathfrak{e}_U\,\frac{s-U}{s}
\ \le\ \frac{U\,\mathfrak{e}_1}{s}+\frac{\eta}{2}.
\end{aligned}
\end{equation}
In the last step we used $\mathfrak{e}_U\le\eta/2$ and $0\le(s-U)/s\le 1$.
Put
\begin{equation*}
S(\eta):=\max\{U(\eta),\,2U(\eta)\mathfrak{e}_1/\eta\}.
\end{equation*}
For $s\ge S(\eta)$ we have $U\mathfrak{e}_1/s\le\eta/2$, and \eqref{5.14:eq-mean-defect-4} gives $0\le\delta_e(s)\le\eta$.
Since $\eta>0$ was arbitrary,
\begin{equation*}
\delta_e(s)\to 0\qquad\text{as } s\to\infty .
\end{equation*}
The choice of $U(\eta)$, and hence the threshold $S(\eta)$, depends on $(L,N,\mathfrak{p})$ and $\eta$ only, because $e$ and $\bar{h}_u$ do.

\emph{(e) Summed bound.}
Summing \eqref{5.14:eq-mean-defect-3} over $u$ and multiplying by $C_{\beta}\ge 0$ gives, for every $1\le s\le K$,
\begin{equation}\label{5.14:eq-mean-defect-5}
\sum_{u=1}^{s}C_{\beta}\,e(g_{K-u})\ \le\ C_{\beta}\sum_{u=1}^{s}\mathfrak{e}_u\ =\ C_{\beta}\,s\,\delta_e(s).
\end{equation}

\medskip
\emph{Assembly.}
Let $1\le s\le K$. We combine three inequalities:
\begin{itemize}
\item the first inequality of \eqref{5.14:eq-one-loop-law-b}, proved above;
\item the remnant bound \eqref{5.14:eq-mean-defect-2};
\item the small-field bound \eqref{5.14:eq-mean-defect-5}.
\end{itemize}
They give
\begin{equation}\label{5.14:eq-mean-defect-6}
\begin{aligned}
|x_{K-s}-x_K-s\,c_0|
&\le 2c_2+\sum_{u=1}^{s}\bigl[C_{\beta}\,e(g_{K-u})+v_{K-u}\bigr]\\
&\le 2c_2+C_v+C_{\beta}\,s\,\delta_e(s)=\Delta(s).
\end{aligned}
\end{equation}
This is the same as $\Delta(s)=2c_2+C_v+C_{\beta}\sum_{u=1}^{s}\mathfrak{e}_u$.

For $s=0$ the sums over $u$ are empty, and the second inequality reduces to $2c_2\le 2c_2+C_v$, which holds because $C_v\ge 0$.

Moreover,
\begin{equation*}
\frac{\Delta(s)}{s}=\frac{2c_2+C_v}{s}+C_{\beta}\,\delta_e(s)\ \longrightarrow\ 0
\qquad\text{as } s\to\infty ,
\end{equation*}
by the Ces\`aro step (d).

\emph{Uniformity.}
By Definition~\ref{5.14:def-defect-majorants}, $\Delta$ depends on neither $K$, nor $g$, nor $g_0$, nor the torus, nor the position. Consequently the following holds for every $g\le\gamma_J$, every $g_0\in\left]0,g_{-}(g)\right]$ with $K=K(g_0,g)$, and every torus of the family.
For every $\eta>0$ put
\begin{equation*}
S'(\eta):=\max\Bigl\{S\bigl(\eta/(2C_{\beta})\bigr),\ 2(2c_2+C_v)/\eta\Bigr\},
\end{equation*}
which depends on $(L,N,\mathfrak{p})$ and $\eta$ only. Then
\begin{equation}\label{5.14:eq-mean-defect-7}
|x_{K-s}-x_K-s\,c_0|\ \le\ \eta\,s\qquad\text{for all } S'(\eta)\le s\le K .
\end{equation}
Indeed, let $s\ge S'(\eta)$.
\begin{itemize}
\item Since $s\ge S(\eta/(2C_{\beta}))$, step (d) gives $\delta_e(s)\le\eta/(2C_{\beta})$, so $C_{\beta}\,\delta_e(s)\le\eta/2$.
\item Since $s\ge 2(2c_2+C_v)/\eta$, we have $(2c_2+C_v)/s\le\eta/2$.
\end{itemize}
Hence $\Delta(s)\le\eta\,s$, and \eqref{5.14:eq-mean-defect-6} gives \eqref{5.14:eq-mean-defect-7}.
This is the asserted uniformity.
This gives the second inequality of \eqref{5.14:eq-one-loop-law-b} and the limit $\Delta(s)/s\to 0$, with the asserted uniformity, in full detail.
\end{proof}

\begin{proof}[Proof of Theorem~\ref{5.14:thm-one-loop-law}, continued]
\label{5.14:prf-coefficient-limits}
We first identify the constant $c_0$ of the one-loop lemma, whose content is displayed in
\eqref{5.14:eq-flow-notation-b}, and fix the normalisation in which it is read.

The Lie algebra $\mathfrak{g}$ of $\mathrm{SU}(N)$ is realised as the traceless Hermitian
$N\times N$ matrices, with inner product $\langle X,Y\rangle=\operatorname{tr}(XY)$ and
$\operatorname{tr}\mathbf{1}=1$. For an orthonormal basis $(t_a)$ of $\mathfrak{g}$ for this
inner product, the Casimir constant $C_{\mathfrak{g}}$ is defined by
\begin{equation*}
\sum_a\,[t_a,[t_a,X]]=C_{\mathfrak{g}}\,X\qquad(X\in\mathfrak{g}).
\end{equation*}
Let $\operatorname{Tr}=N\operatorname{tr}$ be the usual matrix trace and let $(T_a)$ be traceless
Hermitian matrices with $\operatorname{Tr}(T_aT_b)=\tfrac12\delta_{ab}$. Then
$t_a:=\sqrt{2N}\,T_a$ satisfies
$\operatorname{tr}(t_at_b)=2\operatorname{Tr}(T_aT_b)=\delta_{ab}$, so $(t_a)$ is an orthonormal
basis. In the normalisation $\operatorname{Tr}(T_aT_b)=\tfrac12\delta_{ab}$ one has
$\sum_a[T_a,[T_a,X]]=N\,X$, which is the quadratic Casimir of the adjoint representation. Hence
\begin{equation*}
C_{\mathfrak{g}}=2N\cdot N=2N^2
\end{equation*}
for $\mathrm{SU}(N)$, and
\begin{equation}\label{5.14:eq-coefficient-limits-1}
c_0=\frac{11\,C_{\mathfrak{g}}}{24\pi^2}\,\log L=\frac{11N^2}{12\pi^2}\,\log L .
\end{equation}

The coefficient $11C_{\mathfrak{g}}/(24\pi^2)$ is the one produced by the lattice one-loop
computation in the proof of the one-loop lemma. At scale $t=L^{-k}$, and per unit of
$\log(1/t)$, the gluon fluctuation determinant contributes $5C_{\mathfrak{g}}/(12\pi^2)$ and the
ghost determinant contributes $C_{\mathfrak{g}}/(24\pi^2)$. Their total is
\begin{equation*}
\frac{5C_{\mathfrak{g}}}{12\pi^2}+\frac{C_{\mathfrak{g}}}{24\pi^2}
=\frac{11\,C_{\mathfrak{g}}}{24\pi^2}.
\end{equation*}
Since $\log(1/t)=k\log L$, this computation gives
\begin{equation*}
\sum_{l\le k}\beta^{0}_{l}=k\,c_0+O(1),
\end{equation*}
where $\beta^{0}_{l}$ is the coefficient of the Gaussian normaliser of step $l$ of
Ba{\l}aban's construction. This is the bound of the one-loop lemma with its constant $c_2$, as
displayed in \eqref{5.14:eq-flow-notation-b}.

We now relate this normalisation to the standard one. With $\operatorname{tr}\mathbf{1}=1$,
the weight $e^{-g_0^{-2}A(U)}$ is the Wilson weight with $\beta_W=g_0^{-2}$. Here $\beta_W$ is
the Wilson parameter multiplying
$\sum_p\bigl[1-\tfrac1N\operatorname{Re}\operatorname{Tr}U(\partial p)\bigr]$, and
$\beta_W=2N/g_{\mathrm{std}}^2$, where $g_{\mathrm{std}}$ is the coupling in the standard
normalisation. Thus $g_0^{-2}=2N/g_{\mathrm{std}}^2$.

Write $\mu$ for the scale and $\beta_0:=11N/3$ for the standard first coefficient of pure
$\mathrm{SU}(N)$ gauge theory. The standard continuum one-loop law reads
\begin{equation*}
\frac{d\,g_{\mathrm{std}}^{-2}}{d\ln\mu}=\frac{2\beta_0}{16\pi^2}=\frac{11N}{24\pi^2}.
\end{equation*}
Multiplying by $2N$ gives $11N^2/(12\pi^2)$ per e-fold of the scale in the normalisation
$g^{-2}=2N/g_{\mathrm{std}}^2$. One renormalisation step is $\log L$ e-folds, and
$\tfrac{11N^2}{12\pi^2}\log L=c_0$ by \eqref{5.14:eq-coefficient-limits-1}.

Here the coupling is normalised so that its inverse square multiplies the Wilson action with the
trace normalised to one, $g^{-2}=2N/g_{\mathrm{std}}^2$; in standard units the inverse square of
the coupling grows by $11N/(24\pi^2)=\beta_0/(8\pi^2)$ per e-fold of the scale,
$\beta_0=11N/3$, and the factor $2N$ between the two coefficients is this normalisation.

We next prove consequence \eqref{5.14:eq-one-loop-law-c}. Fix $g\le\gamma_J$ and an integer
$s\ge1$. Consider a subsequence as in the corollary to clause~(0) of Theorem~0 on subsequential
limits of the running couplings: bare couplings $g_0^{(i)}\in\left]0,g_{-}(g)\right]$ with
$g_0^{(i)}\downarrow0$, such that $K_i:=K(g_0^{(i)},g)\ge s$ and
$g_{K_i-s}\to g_{\infty,s}$.

By the first-passage definition \eqref{5.14:eq-flow-notation-a},
\begin{equation*}
x_{K_i}\in\left]g^{-2},\,g^{-2}+B^{\sharp}\right].
\end{equation*}
By \eqref{5.14:eq-one-loop-law-b}, applied with $K=K_i$ and this $s\le K_i$,
\begin{equation*}
x_{K_i-s}\in\bigl[x_{K_i}+s\,c_0-\Delta(s),\;x_{K_i}+s\,c_0+\Delta(s)\bigr].
\end{equation*}
Inserting the two bounds on $x_{K_i}$ yields
\begin{equation}\label{5.14:eq-coefficient-limits-2}
x_{K_i-s}\in\bigl[g^{-2}+s\,c_0-\Delta(s),\;g^{-2}+B^{\sharp}+s\,c_0+\Delta(s)\bigr].
\end{equation}
This interval does not depend on $K_i$, since $\Delta$ depends on $(L,N,\mathfrak{p})$ only.

The upper endpoint of \eqref{5.14:eq-coefficient-limits-2} gives
\begin{equation*}
g_{K_i-s}=x_{K_i-s}^{-1/2}\ge\bigl(g^{-2}+B^{\sharp}+s\,c_0+\Delta(s)\bigr)^{-1/2}>0
\end{equation*}
for all $i$. Hence $g_{\infty,s}>0$, and by continuity of $y\mapsto y^{-2}$ on $]0,\infty[$ we
have $g_{\infty,s}^{-2}=\lim_i x_{K_i-s}$. This limit lies in the closed interval
\eqref{5.14:eq-coefficient-limits-2}.

Since $c_0>0$ and $s\ge1$, dividing by $s\,c_0$ gives
\begin{equation*}
\frac{g_{\infty,s}^{-2}}{s\,c_0}\in
\Bigl[\,1+\frac{g^{-2}-\Delta(s)}{s\,c_0},\;1+\frac{g^{-2}+B^{\sharp}+\Delta(s)}{s\,c_0}\,\Bigr].
\end{equation*}
Now let $s\to\infty$ with $g$ fixed. Then $g^{-2}/s\to0$ and $B^{\sharp}/s\to0$, and
$\Delta(s)/s\to0$ by \eqref{5.14:eq-one-loop-law-b}. Both endpoints therefore tend to $1$, and
so $g_{\infty,s}^{-2}/(s\,c_0)\to1$.

Dividing \eqref{5.14:eq-one-loop-law-b} by $s$ gives, for every $K\ge s$,
\begin{equation*}
\Bigl|\frac{x_{K-s}-x_K}{s}-c_0\Bigr|\le\frac{\Delta(s)}{s}.
\end{equation*}
The right side does not depend on $K$ and tends to zero. Hence $(x_{K-s}-x_K)/s\to c_0$ as
$s\to\infty$, uniformly in $K\ge s$. This proves \eqref{5.14:eq-one-loop-law-c}.
\end{proof}

\begin{lemma}[Logarithmic tie between separation and witness depth]\label{5.14:lem-depth-separation}
Fix the reference coupling $g$, and use the letters of Notation~\ref{5.14:not-flow-notation}, in
particular the function $C_W$, the witness depth $s(\mathsf{d})$ and the window endpoint $\bar{x}_s$ of
\eqref{5.14:eq-flow-notation-e}. Limits as $\mathsf{d}\downarrow 0$ are taken at fixed $g$. Then
\begin{equation}\label{5.14:eq-depth-separation-1}
s(\mathsf{d})\,\log L=\ln(1/\mathsf{d})\,\bigl(1+o(1)\bigr)\qquad\text{as }\mathsf{d}\downarrow 0,
\end{equation}
where the term $o(1)$ depends on $g$ (through $C_0$, $B^{\sharp}$, $p_0$, $\vartheta$) but not on $K$,
$g_0$, $m$ or the position of the supports. In particular, where $s(\mathsf{d})$ enters the two-point
witness theorem, \eqref{5.14:eq-depth-separation-1} holds
uniformly in $K\ge s(\mathsf{d})+k_W$, in $g_0$, in the position, in the pair
$(f_1,f_2)\in\mathfrak{T}(\mathsf{d};\vartheta)$, and in $m$ on the range $m\le m_V(g)$ (for each torus
beyond that range, under the condition $s(\mathsf{d})\ge s_{\mathrm{vol}}(g,m)$ of the two-point witness
theorem). Here $\log$ and $\ln$ both denote the natural logarithm.
\end{lemma}

\begin{proof}
By \eqref{5.14:eq-flow-notation-e},
\begin{equation*}
s(\mathsf{d})=\min\bigl\{s\ge 0:\ C_W(s)L^{-s}\le\mathsf{d}\bigr\},
\end{equation*}
where $C_W\ge 1$ is non-decreasing in $s$.

\emph{Lower bound.} Since $s(\mathsf{d})$ belongs to the set of which it is the minimum,
$C_W(s(\mathsf{d}))L^{-s(\mathsf{d})}\le\mathsf{d}$. Taking logarithms and rearranging, and then using
$C_W(s(\mathsf{d}))\ge 1$, we obtain
\begin{equation}\label{5.14:eq-depth-separation-2}
s(\mathsf{d})\log L\ \ge\ \ln(1/\mathsf{d})+\ln C_W(s(\mathsf{d}))\ \ge\ \ln(1/\mathsf{d}).
\end{equation}
Since $L>1$, \eqref{5.14:eq-depth-separation-2} gives $s(\mathsf{d})\ge\ln(1/\mathsf{d})/\log L$; in
particular $s(\mathsf{d})\uparrow\infty$ as $\mathsf{d}\downarrow 0$. Moreover, for $\mathsf{d}<1$ the
right-hand side of \eqref{5.14:eq-depth-separation-2} is strictly positive, so the integer $s(\mathsf{d})$
satisfies $s(\mathsf{d})\ge 1$.

\emph{Upper bound.} Let $\mathsf{d}<1$, so that $s(\mathsf{d})-1\ge 0$. By minimality of $s(\mathsf{d})$ the
integer $s(\mathsf{d})-1$ does not belong to the set defining $s(\mathsf{d})$, that is,
\begin{equation*}
C_W(s(\mathsf{d})-1)\,L^{-(s(\mathsf{d})-1)}>\mathsf{d}.
\end{equation*}
Since $C_W$ is non-decreasing, $C_W(s(\mathsf{d}))\ge C_W(s(\mathsf{d})-1)$, hence
$C_W(s(\mathsf{d}))\,L^{-(s(\mathsf{d})-1)}>\mathsf{d}$. Taking logarithms,
\begin{equation*}
\bigl(s(\mathsf{d})-1\bigr)\log L<\ln(1/\mathsf{d})+\ln C_W(s(\mathsf{d})).
\end{equation*}
Combining this with \eqref{5.14:eq-depth-separation-2}, we obtain for every $\mathsf{d}\in\left]0,1\right[$
\begin{equation}\label{5.14:eq-depth-separation-a}
\ln(1/\mathsf{d})\ \le\ s(\mathsf{d})\log L\ <\ \ln(1/\mathsf{d})+\ln C_W(s(\mathsf{d}))+\log L.
\end{equation}

\emph{Growth of $C_W$.} By \eqref{5.14:eq-flow-notation-e}, $C_W(s)=C_0A_0^2(\log\bar{x}_s)^{2p_0}$ with
$\bar{x}_s=g^{-2}+B^{\sharp}(s+1)$, so
\begin{equation*}
\ln C_W(s)=\ln\bigl(C_0A_0^2\bigr)+2p_0\,\ln\log\bar{x}_s .
\end{equation*}
At fixed $g$ we have $\bar{x}_s\le(g^{-2}+B^{\sharp})(s+1)$, hence
$\log\bar{x}_s\le\log(g^{-2}+B^{\sharp})+\log(s+1)$, and $\bar{x}_s\to\infty$ as $s\to\infty$.
Therefore $\ln C_W(s)=O(\log\log\bar{x}_s)$, and, $\log\log\bar{x}_s$ growing at most like
$\log\log(s+1)$ up to an additive constant depending on $g$, also $\ln C_W(s)=o(s)$ as $s\to\infty$ at
fixed $g$.

\emph{Conclusion.} Put
\begin{equation*}
\varepsilon(\mathsf{d}):=\frac{\ln C_W(s(\mathsf{d}))+\log L}{s(\mathsf{d})\log L}.
\end{equation*}
Since $s(\mathsf{d})\to\infty$ as $\mathsf{d}\downarrow 0$, the previous paragraph gives
$\ln C_W(s(\mathsf{d}))/(s(\mathsf{d})\log L)\to 0$, and clearly $\log L/(s(\mathsf{d})\log L)\to 0$;
thus $\varepsilon(\mathsf{d})\to 0$. Dividing \eqref{5.14:eq-depth-separation-a} by $s(\mathsf{d})\log L$
gives
\begin{equation*}
1-\varepsilon(\mathsf{d})<\frac{\ln(1/\mathsf{d})}{s(\mathsf{d})\log L}\le 1 .
\end{equation*}
Hence, for $\mathsf{d}$ so small that $\varepsilon(\mathsf{d})<1$, the ratio
$s(\mathsf{d})\log L/\ln(1/\mathsf{d})$ lies in $\left[1,(1-\varepsilon(\mathsf{d}))^{-1}\right[$, and
$(1-\varepsilon(\mathsf{d}))^{-1}\to 1$. This is \eqref{5.14:eq-depth-separation-1}.

\emph{Uniformity.} The quantity $\varepsilon(\mathsf{d})$ is a function of $s(\mathsf{d})$ and of $C_W$
alone, and $s(\mathsf{d})$ is determined by $\mathsf{d}$, $L$ and $C_W$; by
\eqref{5.14:eq-flow-notation-e}, $C_W$ is determined by $g$ and by the constants $C_0$, $A_0$, $p_0$,
$B^{\sharp}$. Neither $s(\mathsf{d})$ nor $C_W$ involves $K$, $g_0$, $m$, the position of the supports or the
pair $(f_1,f_2)$. The lower bound $s(\mathsf{d})\ge\ln(1/\mathsf{d})/\log L$ from
\eqref{5.14:eq-depth-separation-2} likewise involves only $\mathsf{d}$ and $L$. Hence the term $o(1)$
in \eqref{5.14:eq-depth-separation-1} does not depend on $K$, $g_0$, $m$ or the position, and the
convergence is uniform in the ranges stated in the lemma.
\end{proof}

\begin{corollary}[One-loop law of the coupling at separation $\mathsf{d}$]
\label{5.14:cor-separation-one-loop}
Let $0<g\le\gamma_J$ be fixed. Let $g_0\in\left]0,g_{-}(g)\right]$ with $K=K(g_0,g)$, let $\mathbb{T}_m$ be
any torus of the family, and let $\mathsf{d}>0$ be a separation with $K-s(\mathsf{d})\ge k_W$, where
the witness depth $s(\mathsf{d})$ and the floor $k_W$ are those of
\eqref{5.14:eq-flow-notation-e}. Put
\begin{equation*}
g(\mathsf{d}) := g_{K-s(\mathsf{d})},
\end{equation*}
the coupling at level $K-s(\mathsf{d})$; for $g\le\gamma_W^{\mathbb{T}}$ it is the coupling that
appears in the display of part (a) of clause (iv) of Theorem~0 at cutoff $K$. Then
\begin{equation}\label{5.14:eq-separation-one-loop-a}
\frac{1}{g(\mathsf{d})^{2}}
=\frac{11N^{2}}{12\pi^{2}}\,\ln\frac{1}{\mathsf{d}}\,\bigl(1+o(1)\bigr)
\qquad\text{as } \mathsf{d}\downarrow 0,
\end{equation}
uniformly in $K\ge s(\mathsf{d})+k_W$, in $g_0$, in $m$ and in the position of the supports in
part (a); here $o(1)$ denotes a quantity whose modulus is bounded by a function of $\mathsf{d}$
alone (for the fixed reference coupling $g$ and the fixed constants of the construction) that
tends to zero as $\mathsf{d}\downarrow 0$.
Equivalently, $g(\mathsf{d})^{-2}/\ln(1/\mathsf{d})\to 11N^{2}/12\pi^{2}$. The leading term
contains neither the blocking factor $L$ nor any of the auxiliary parameters $\mathfrak{p}$;
its coefficient is the universal one-loop coefficient of $\mathrm{SU}(N)$ Yang--Mills theory.
Here the coupling is normalised so that its inverse square multiplies the Wilson action with
the trace normalised to one, $g'^{-2}=2N/g_{\mathrm{std}}^{2}$ for a coupling $g'$ of the
construction and the standard coupling $g_{\mathrm{std}}$; in standard units the inverse square
of the coupling grows by $11N/(24\pi^{2})=\beta_0/(8\pi^{2})$ per e-fold of the scale,
$\beta_0=11N/3$, and the factor $2N$ between the two coefficients is this normalisation.

The same holds for the subsequential limit couplings: if $g_{\infty,s}$ denotes a limit of
$g_{K_i-s}$ along a sequence $g_0^{(i)}\downarrow 0$ with $K_i:=K(g_0^{(i)},g)\ge s$, as in
\eqref{5.14:eq-one-loop-law-c}, then
\begin{equation}\label{5.14:eq-separation-one-loop-1}
g_{\infty,s(\mathsf{d})}^{-2}
=\frac{11N^{2}}{12\pi^{2}}\,\ln\frac{1}{\mathsf{d}}\,\bigl(1+o(1)\bigr)
\qquad\text{as } \mathsf{d}\downarrow 0 .
\end{equation}
\end{corollary}

This is a statement about the construction's sequence of couplings; nothing is asserted here
about Schwinger functions.

\begin{proof}
The floor $K-s(\mathsf{d})\ge k_W\ge 1$ gives $0\le s(\mathsf{d})\le K$, so the bound
\eqref{5.14:eq-one-loop-law-b} of the exact one-loop law (Theorem~\ref{5.14:thm-one-loop-law})
applies with $s=s(\mathsf{d})$: $|x_{K-s(\mathsf{d})}-x_K-s(\mathsf{d})c_0|\le\Delta(s(\mathsf{d}))$,
where the majorant $\Delta$ depends on $(L,N,\mathfrak{p})$ only and $\Delta(s)/s\to 0$ as
$s\to\infty$. Hence there is a number $\theta$ with $|\theta|\le 1$ such that, using
$x_k=g_k^{-2}$,
\begin{equation}\label{5.14:eq-separation-one-loop-2}
\frac{1}{g(\mathsf{d})^{2}}=x_{K-s(\mathsf{d})}
=x_K+s(\mathsf{d})\,c_0+\theta\,\Delta(s(\mathsf{d})).
\end{equation}
By the logarithmic tie between separation and witness depth
(Lemma~\ref{5.14:lem-depth-separation}), display \eqref{5.14:eq-depth-separation-a}, we have
$s(\mathsf{d})\log L=\ln(1/\mathsf{d})(1+o(1))$, so $s(\mathsf{d})\to\infty$ as
$\mathsf{d}\downarrow 0$; in particular $s(\mathsf{d})\ge 1$ for $\mathsf{d}$ small, and
$c_0=(11N^{2}/12\pi^{2})\log L>0$ by Theorem~\ref{5.14:thm-one-loop-law}. For such $\mathsf{d}$,
\eqref{5.14:eq-separation-one-loop-2} reads
\begin{equation*}
\frac{1}{g(\mathsf{d})^{2}}=s(\mathsf{d})\,c_0\,\bigl(1+\eta_1(\mathsf{d})\bigr),
\qquad
\eta_1(\mathsf{d}):=\frac{x_K+\theta\,\Delta(s(\mathsf{d}))}{s(\mathsf{d})\,c_0},
\end{equation*}
and, recalling $x_K\in\left]g^{-2},g^{-2}+B^{\sharp}\right]$ from
\eqref{5.14:eq-flow-notation-a}, so that $0<x_K\le g^{-2}+B^{\sharp}$, and $|\theta|\le 1$,
\begin{equation}\label{5.14:eq-separation-one-loop-3}
|\eta_1(\mathsf{d})|
\le\frac{g^{-2}+B^{\sharp}}{s(\mathsf{d})\,c_0}+\frac{\Delta(s(\mathsf{d}))}{s(\mathsf{d})\,c_0}.
\end{equation}
The first term tends to zero because $g$ is fixed and $s(\mathsf{d})\to\infty$; the second
because $\Delta(s)/s\to 0$. The right side of \eqref{5.14:eq-separation-one-loop-3} depends on
$\mathsf{d}$ only through $s(\mathsf{d})$, which does not depend on $K$, $g_0$, $m$ or the
position, and $\Delta$ depends on $(L,N,\mathfrak{p})$ only; the convergence is therefore uniform.

Next, by the value of $c_0$ in Theorem~\ref{5.14:thm-one-loop-law} and by
\eqref{5.14:eq-depth-separation-a},
\begin{equation*}
s(\mathsf{d})\,c_0=\frac{c_0}{\log L}\,s(\mathsf{d})\log L
=\frac{11N^{2}}{12\pi^{2}}\,s(\mathsf{d})\log L
=\frac{11N^{2}}{12\pi^{2}}\,\ln\frac{1}{\mathsf{d}}\,\bigl(1+\eta_2(\mathsf{d})\bigr),
\end{equation*}
where $\eta_2(\mathsf{d})\to 0$ as $\mathsf{d}\downarrow 0$, and by
Lemma~\ref{5.14:lem-depth-separation} $\eta_2$ depends on $g$ but not on $K$, $g_0$, $m$ or
the position. Multiplying the two factors,
\begin{equation*}
\frac{1}{g(\mathsf{d})^{2}}
=\frac{11N^{2}}{12\pi^{2}}\,\ln\frac{1}{\mathsf{d}}\,
\bigl(1+\eta_1(\mathsf{d})\bigr)\bigl(1+\eta_2(\mathsf{d})\bigr),
\end{equation*}
and $|(1+\eta_1)(1+\eta_2)-1|\le|\eta_1|+|\eta_2|+|\eta_1||\eta_2|$
tends to zero uniformly in the stated sense. This is \eqref{5.14:eq-separation-one-loop-a};
dividing by $\ln(1/\mathsf{d})>0$ for $\mathsf{d}<1$ gives the equivalent limit form.

For the subsequential limits, \eqref{5.14:eq-one-loop-law-c} gives, for every $s\ge 0$,
\begin{equation*}
g_{\infty,s}^{-2}\in
\bigl[g^{-2}+s\,c_0-\Delta(s),\;g^{-2}+B^{\sharp}+s\,c_0+\Delta(s)\bigr],
\end{equation*}
an interval free of $K$; dividing by $s\,c_0$ for $s\ge 1$,
\begin{equation*}
\frac{g_{\infty,s}^{-2}}{s\,c_0}\in
\Bigl[1+\frac{g^{-2}-\Delta(s)}{s\,c_0},\;1+\frac{g^{-2}+B^{\sharp}+\Delta(s)}{s\,c_0}\Bigr].
\end{equation*}
Take $s=s(\mathsf{d})$ and write the ratio as $1+\eta_3(\mathsf{d})$. Since
$|g^{-2}-\Delta(s)|\le g^{-2}+\Delta(s)$, both ends give
$|\eta_3(\mathsf{d})|\le(g^{-2}+B^{\sharp}+\Delta(s(\mathsf{d})))/(s(\mathsf{d})\,c_0)$,
which tends to zero by the same two facts as above. Multiplying
$g_{\infty,s(\mathsf{d})}^{-2}=s(\mathsf{d})c_0(1+\eta_3(\mathsf{d}))$ by the expression for
$s(\mathsf{d})c_0$ obtained above yields \eqref{5.14:eq-separation-one-loop-1}.
\end{proof}

\begin{corollary}[Fixed-factor one-loop bounds for the physical charge]\label{5.14:cor-physical-charge}
Assume the following two statements.
\begin{itemize}
\item[(a)] The two-point law of clause (iv) of Theorem~0 on the torus, that is, the two-sided
display of the two-point witness theorem, with the letters fixed in
Notation~\ref{5.14:not-flow-notation} and \eqref{5.14:eq-flow-notation-e}. Let
$g\le\gamma_W^{\mathbb{T}}$, $g_0\in\,]0,g_{-}(g)]$ with $K=K(g_0,g)$, and let $m$ be arbitrary.
Let $\mathsf{d}$ satisfy the depth floor of the two-point witness theorem on the torus and its
ultraviolet floor $K-s(\mathsf{d})\ge k_W$, and, when $m>m_V(g)$, also
$s(\mathsf{d})\ge s_{\mathrm{vol}}(g,m)$. Put $s:=s(\mathsf{d})$, and let
$(f_1,f_2)\in\mathfrak{T}(\mathsf{d};\vartheta)$. Then
\begin{gather*}
S^{T}_{2;K}(f_1,f_2)=\zeta'^{\,2}_{K-s}\, b_0\, g_{K-s}^4\,\mathcal{K}_s(f_1,f_2)\,(1+R_K),\\
|R_K|\le\tfrac12,\qquad \zeta'_{K-s}\in\left[\tfrac23,\tfrac43\right],\\
c_{\mathcal{K}}\|f_1\|_\infty\|f_2\|_\infty\le\mathcal{K}_s(f_1,f_2)
\le C_{\mathcal{K}}\|f_1\|_\infty\|f_2\|_\infty .
\end{gather*}
This holds uniformly in $K$, $g_0$ and the position, and uniformly in $m$ on the range
$m_{\mathbb{T}}(g_{-}(g))\le m\le m_V(g)$ of that clause.
\item[(b)] Clause (iii)(d) of Theorem~0, the exact one-loop law of the construction's running
coupling \eqref{5.14:eq-one-loop-law-b}, for $g\le\gamma_J$.
\end{itemize}
Work in the setting of (a), for the group $\mathrm{SU}(2)$, so that $N=2$, and put
$g(\mathsf{d}):=g_{K-s(\mathsf{d})}$. For $(f_1,f_2)\in\mathfrak{T}(\mathsf{d};\vartheta)$ define the
charge $g_{\mathrm{phys}}(\mathsf{d})>0$ read off the two-point function, normalised by the kernel, by
\begin{equation}\label{5.14:eq-physical-charge-1}
g_{\mathrm{phys}}(\mathsf{d})^4:=\frac{S^{T}_{2;K}(f_1,f_2)}{b_0\,\mathcal{K}_{s(\mathsf{d})}(f_1,f_2)} .
\end{equation}
Then, as $\mathsf{d}\downarrow0$ at fixed $g$,
\begin{equation}\label{5.14:eq-physical-charge-2}
\left(\tfrac38\right)^{1/2}(1+o(1))
\le\frac{g_{\mathrm{phys}}(\mathsf{d})^{-2}}{\frac{11N^2}{12\pi^2}\,\ln(1/\mathsf{d})}
\le\left(\tfrac92\right)^{1/2}(1+o(1)).
\end{equation}
The bound is uniform in $K\ge s(\mathsf{d})+k_W$, in $g_0$, in the position, in the pair
$(f_1,f_2)\in\mathfrak{T}(\mathsf{d};\vartheta)$, and in $m$ on the range
$m_{\mathbb{T}}(g_{-}(g))\le m\le m_V(g)$; for each
torus beyond that range it holds under $s(\mathsf{d})\ge s_{\mathrm{vol}}(g,m)$.
\end{corollary}

\begin{proof}
By (a), the quotient in \eqref{5.14:eq-physical-charge-1} is positive, so that
$g_{\mathrm{phys}}(\mathsf{d})$ is well defined as its positive fourth root. Indeed, the functions of
a pair in $\mathfrak{T}(\mathsf{d};\vartheta)$ are not identically zero, so
$\mathcal{K}_s(f_1,f_2)\ge c_{\mathcal{K}}\|f_1\|_\infty\|f_2\|_\infty>0$. Moreover
$\zeta'^{\,2}_{K-s}\ge\frac49$, $1+R_K\ge\frac12$, $b_0=(N^2-1)/2>0$, and $g_{K-s}>0$ by the strict
positivity of the running coupling at every scale (clause (0) of Theorem~0), so the display
of (a) gives $S^{T}_{2;K}(f_1,f_2)>0$.

Divide the display of (a) by $b_0\,g_{K-s}^4\,\mathcal{K}_s(f_1,f_2)$, with $s=s(\mathsf{d})$ and
$g_{K-s}=g(\mathsf{d})$. By \eqref{5.14:eq-physical-charge-1} this gives
\begin{equation*}
\frac{g_{\mathrm{phys}}(\mathsf{d})^4}{g(\mathsf{d})^4}=\zeta'^{\,2}_{K-s}\,(1+R_K).
\end{equation*}
Since $\zeta'_{K-s}\in[\frac23,\frac43]$ is positive, we have
$\zeta'^{\,2}_{K-s}\in[\frac49,\frac{16}9]$. Since $|R_K|\le\frac12$, we have
$1+R_K\in[\frac12,\frac32]$. Both factors are positive, so their product lies between the products of
the endpoints, $\frac49\cdot\frac12=\frac29$ and $\frac{16}9\cdot\frac32=\frac83$:
\begin{equation}\label{5.14:eq-physical-charge-3}
\frac29\le\frac{g_{\mathrm{phys}}(\mathsf{d})^4}{g(\mathsf{d})^4}\le\frac83 .
\end{equation}
The map $t\mapsto t^{-1/2}$ is decreasing on $]0,\infty[$. Raising \eqref{5.14:eq-physical-charge-3}
to the power $-\frac12$ therefore reverses the order:
\begin{equation}\label{5.14:eq-physical-charge-4}
\left(\tfrac38\right)^{1/2}=\left(\tfrac83\right)^{-1/2}
\le\frac{g_{\mathrm{phys}}(\mathsf{d})^{-2}}{g(\mathsf{d})^{-2}}
\le\left(\tfrac29\right)^{-1/2}=\left(\tfrac92\right)^{1/2}.
\end{equation}
The window \eqref{5.14:eq-physical-charge-4} holds with the constants of (a) alone. It is therefore
uniform in everything in which (a) is uniform, and in the pair $(f_1,f_2)$.

Since the torus coupling ceiling of (a) satisfies $\gamma_W^{\mathbb{T}}\le\gamma_J$, the coupling
ceiling of clause (iii), statement (b) applies at every $g\le\gamma_W^{\mathbb{T}}$.
Through Corollary~\ref{5.14:cor-separation-one-loop} it yields the one-loop law of the coupling at
separation $\mathsf{d}$, \eqref{5.14:eq-separation-one-loop-a}:
\begin{equation*}
\frac{g(\mathsf{d})^{-2}}{\frac{11N^2}{12\pi^2}\,\ln(1/\mathsf{d})}=1+o(1)
\qquad(\mathsf{d}\downarrow0),
\end{equation*}
uniformly in $K\ge s(\mathsf{d})+k_W$, $g_0$, $m$ and the position. Corollary~\ref{5.14:cor-separation-one-loop}
obtains this from (b) and from Lemma~\ref{5.14:lem-depth-separation}. For $\mathsf{d}<1$ the left-hand
side is positive. We now write
\begin{equation*}
\frac{g_{\mathrm{phys}}(\mathsf{d})^{-2}}{\frac{11N^2}{12\pi^2}\,\ln(1/\mathsf{d})}
=\frac{g_{\mathrm{phys}}(\mathsf{d})^{-2}}{g(\mathsf{d})^{-2}}\cdot
\frac{g(\mathsf{d})^{-2}}{\frac{11N^2}{12\pi^2}\,\ln(1/\mathsf{d})} .
\end{equation*}
Multiply the two inequalities of \eqref{5.14:eq-physical-charge-4} by the positive second factor,
which equals $1+o(1)$. This gives \eqref{5.14:eq-physical-charge-2}, with the same $o(1)$ as in
\eqref{5.14:eq-separation-one-loop-a}. The uniformity of \eqref{5.14:eq-physical-charge-2} is the
conjunction of the uniformity of \eqref{5.14:eq-physical-charge-4} and that of
\eqref{5.14:eq-separation-one-loop-a}. This includes the range in $m$ and the condition
$s(\mathsf{d})\ge s_{\mathrm{vol}}(g,m)$ for each torus beyond it, which enter only through (a).
\end{proof}

The two ends of the window in \eqref{5.14:eq-physical-charge-2} differ by the factor
\begin{equation*}
\frac{(9/2)^{1/2}}{(3/8)^{1/2}}=\left(\tfrac92\cdot\tfrac83\right)^{1/2}=12^{1/2}.
\end{equation*}

The same window holds for the limits on a fixed torus $\mathbb{T}_m$. Let $(K_j)_{j\ge1}$ be the
subsequence of cutoffs of clause (ii-$\mathbb{T}$-b) of Theorem~0, along which
$g_{K_j-s}\to g_{\infty,s}$, the subsequential limit coupling of clause (0). Fix $\mathsf{d}$ and a
pair $(f_1,f_2)\in\mathfrak{T}(\mathsf{d};\vartheta)$, let $S(\mathsf{d})$ be a limit point of
$S^{T}_{2;K_j,m}(f_1,f_2)$ along this subsequence, and let $\mathcal{K}^{-}$ and $\mathcal{K}^{+}$ be
the lower and the upper limit of $\mathcal{K}_s(f_1,f_2)$ along it. At every cutoff $K_j$ with
$K_j\ge s+k_W$ the display of (a) holds, and, as in the derivation of
\eqref{5.14:eq-physical-charge-3}, the factor $\zeta'^{\,2}_{K_j-s}(1+R_{K_j})$ lies in
$[\frac29,\frac83]$; hence
\begin{equation*}
\tfrac29\, b_0\, g_{K_j-s}^4\,\mathcal{K}_s(f_1,f_2)\le S^{T}_{2;K_j,m}(f_1,f_2)
\le\tfrac83\, b_0\, g_{K_j-s}^4\,\mathcal{K}_s(f_1,f_2).
\end{equation*}
Pass to the limit along a further subsequence on which $S^{T}_{2;K_j,m}(f_1,f_2)\to S(\mathsf{d})$.
Since $g_{K_j-s}^4\to g_{\infty,s}^4$, and the lower and upper limits of $\mathcal{K}_s(f_1,f_2)$ along
the further subsequence lie in $[\mathcal{K}^{-},\mathcal{K}^{+}]$, this gives
\begin{equation*}
\tfrac29\, b_0\, g_{\infty,s}^4\,\mathcal{K}^{-}\le S(\mathsf{d})\le\tfrac83\, b_0\,
g_{\infty,s}^4\,\mathcal{K}^{+},
\end{equation*}
and the kernel bounds of (a) pass to the limit as
\begin{equation*}
c_{\mathcal{K}}\|f_1\|_\infty\|f_2\|_\infty\le\mathcal{K}^{-}\le\mathcal{K}^{+}
\le C_{\mathcal{K}}\|f_1\|_\infty\|f_2\|_\infty .
\end{equation*}
Now take $s=s(\mathsf{d})$. Corollary~\ref{5.14:cor-separation-one-loop} asserts the one-loop law
also for the subsequential limits of clause (0):
$g_{\infty,s(\mathsf{d})}^{-2}=\frac{11N^2}{12\pi^2}\ln(1/\mathsf{d})\,(1+o(1))$, that is,
$g_{\infty,s(\mathsf{d})}^{4}\bigl(\frac{11N^2}{12\pi^2}\ln(1/\mathsf{d})\bigr)^2=1+o(1)$. Multiplying
the two bounds on $S(\mathsf{d})$ by $\bigl(\frac{11N^2}{12\pi^2}\ln(1/\mathsf{d})\bigr)^2$ and
dividing by $b_0\mathcal{K}^{-}$, respectively by $b_0\mathcal{K}^{+}$, we obtain
\begin{align*}
\tfrac29\,(1+o(1))&\le\frac{S(\mathsf{d})\bigl(\frac{11N^2}{12\pi^2}\ln(1/\mathsf{d})\bigr)^2}
{b_0\,\mathcal{K}^{-}},\\
\frac{S(\mathsf{d})\bigl(\frac{11N^2}{12\pi^2}\ln(1/\mathsf{d})\bigr)^2}{b_0\,\mathcal{K}^{+}}
&\le\tfrac83\,(1+o(1)).
\end{align*}
Let now $\mathcal{K}\in[\mathcal{K}^{-},\mathcal{K}^{+}]$ be any number determined along the same
subsequence, and define, for the limit, $g_{\mathrm{phys}}(\mathsf{d})>0$ by
$g_{\mathrm{phys}}(\mathsf{d})^4:=S(\mathsf{d})/(b_0\mathcal{K})$; this applies in particular to the
dilates to scale $\mathsf{d}$ of a fixed reference pair. Since
$\mathcal{K}^{-}/\mathcal{K}\ge\mathcal{K}^{-}/\mathcal{K}^{+}$ and
$\mathcal{K}^{+}/\mathcal{K}\le\mathcal{K}^{+}/\mathcal{K}^{-}$, the last two bounds give
\begin{equation*}
\tfrac29\,\frac{\mathcal{K}^{-}}{\mathcal{K}^{+}}\,(1+o(1))
\le g_{\mathrm{phys}}(\mathsf{d})^4\Bigl(\tfrac{11N^2}{12\pi^2}\ln(1/\mathsf{d})\Bigr)^2
\le\tfrac83\,\frac{\mathcal{K}^{+}}{\mathcal{K}^{-}}\,(1+o(1)),
\end{equation*}
and raising to the power $-\frac12$, which reverses the order, and using
$(1+o(1))^{-1/2}=1+o(1)$,
\begin{equation*}
\left(\tfrac38\right)^{1/2}\Bigl(\frac{\mathcal{K}^{-}}{\mathcal{K}^{+}}\Bigr)^{1/2}(1+o(1))
\le\frac{g_{\mathrm{phys}}(\mathsf{d})^{-2}}{\frac{11N^2}{12\pi^2}\,\ln(1/\mathsf{d})}
\le\left(\tfrac92\right)^{1/2}\Bigl(\frac{\mathcal{K}^{+}}{\mathcal{K}^{-}}\Bigr)^{1/2}(1+o(1)).
\end{equation*}
Thus the window of \eqref{5.14:eq-physical-charge-2} holds up to the ratio
$(\mathcal{K}^{+}/\mathcal{K}^{-})^{1/2}$, which lies in $[1,(C_{\mathcal{K}}/c_{\mathcal{K}})^{1/2}]$
by the kernel bounds above. If instead the charge is normalised by the sup norms,
$g_{\mathrm{phys}}(\mathsf{d})^4:=S(\mathsf{d})/(b_0\|f_1\|_\infty\|f_2\|_\infty)$, the bounds on
$S(\mathsf{d})$ together with $c_{\mathcal{K}}\|f_1\|_\infty\|f_2\|_\infty\le\mathcal{K}^{-}$ and
$\mathcal{K}^{+}\le C_{\mathcal{K}}\|f_1\|_\infty\|f_2\|_\infty$ give
$g_{\mathrm{phys}}(\mathsf{d})^4\bigl(\frac{11N^2}{12\pi^2}\ln(1/\mathsf{d})\bigr)^2
\in[\frac29 c_{\mathcal{K}},\frac83 C_{\mathcal{K}}]\,(1+o(1))$, and hence
\begin{equation*}
\left(\tfrac38\right)^{1/2}C_{\mathcal{K}}^{-1/2}(1+o(1))
\le\frac{g_{\mathrm{phys}}(\mathsf{d})^{-2}}{\frac{11N^2}{12\pi^2}\,\ln(1/\mathsf{d})}
\le\left(\tfrac92\right)^{1/2}c_{\mathcal{K}}^{-1/2}(1+o(1)).
\end{equation*}
In every one of these readings the inverse square of the charge read off the two-point function is
within a fixed factor of the one-loop Yang--Mills value.

Here the coupling is normalised so that its inverse square multiplies the Wilson action with the
trace normalised to one, $g^{-2}=2N/g_{\mathrm{std}}^2$. In standard units the inverse square of the
coupling grows by $11N/(24\pi^2)=\beta_0/(8\pi^2)$ per e-fold of the scale, with $\beta_0=11N/3$, and
the factor $2N$ between the two coefficients is this normalisation.

Consequently the coupling $g(\mathsf{d})$ of part (a) obeys
\begin{equation*}
1/g(\mathsf{d})^2=\frac{11N^2}{12\pi^2}\,\ln(1/\mathsf{d})\,(1+o(1))
\end{equation*}
as $\mathsf{d}\to0$, uniformly in the cutoff. The charge read off the two-point function has a fourth
power that differs from $g(\mathsf{d})^4$ by a factor in $[\frac29,\frac83]$. It therefore runs with an
inverse-square slope between $(3/8)^{1/2}$ and $(9/2)^{1/2}$ times the one-loop Yang--Mills value,
that is, within a fixed factor of it.

\begin{corollary}[The one-loop law for the coupling of the two-point display]
\label{5.14:cor-the-one-loop-law-for-the-cou}
We work in the setting of the two-point display of clause (iv)(a) of Theorem~0, display
\eqref{5.14:eq-flow-notation-e}. Precisely, let $N=2$, and let
$g\le\gamma_W^{\mathbb{T}}$, the coupling ceiling of the two-point witness theorem for the
torus family; recall that $\gamma_W^{\mathbb{T}}\le\gamma_J$. Let
$g_0\in\left]0,g_{-}(g)\right]$ and $K=K(g_0,g)$. Let $m$ be any member of the torus
family. Let $\mathsf{d}$ satisfy the depth floor and the ultraviolet floor of the two-point
witness theorem for the torus family; the latter gives
\begin{equation}\label{5.14:eq-the-one-loop-law-for-the-cou-1}
K\ \ge\ s(\mathsf{d})+k_W .
\end{equation}
Put $s:=s(\mathsf{d})$, and let $(f_1,f_2)\in\mathfrak{T}(\mathsf{d};\vartheta)$.
Define the coupling of the display at cutoff $K$ by
\begin{equation*}
g(\mathsf{d}) := g_{K-s(\mathsf{d})} .
\end{equation*}
Limits $\mathsf{d}\downarrow 0$ are taken at fixed $g$, which is fixed in the quantifier prefix
of Theorem~0 before $\mathsf{d}$. A quantity $o(1)$ is one whose absolute value is bounded by a
function of $\mathsf{d}$ alone, which tends to zero as $\mathsf{d}\downarrow 0$ once $g$, $N$,
$L$ and the fixed parameters of the construction are chosen. This function does not depend on
$K\ge s(\mathsf{d})+k_W$, on
$g_0$, on the position of the supports, on the pair
$(f_1,f_2)\in\mathfrak{T}(\mathsf{d};\vartheta)$, or on $m$ in the range $m\le m_V(g)$. For
$m>m_V(g)$ the same holds for each torus separately, under the volume condition
$s(\mathsf{d})\ge s_{\mathrm{vol}}(g,m)$ exactly as in clause (iv)(a) of Theorem~0.
Then
\begin{equation}\label{5.14:eq-the-one-loop-law-for-the-cou-a}
\begin{split}
\frac{1}{g(\mathsf{d})^{2}}
&=\frac{11N^2}{12\pi^2}\,\ln(1/\mathsf{d})\,\bigl(1+o(1)\bigr)
\qquad\text{as } \mathsf{d}\downarrow 0,\\
&\qquad\text{uniformly in } K\ge s(\mathsf{d})+k_W,\ g_0,\ m \text{ and the position.}
\end{split}
\end{equation}
Equivalently, $g(\mathsf{d})^{-2}/\ln(1/\mathsf{d})\to 11N^2/(12\pi^2)$ as
$\mathsf{d}\downarrow0$, uniformly in the same quantities.

The same holds for the limit couplings. Let $g_0^{(i)}\downarrow 0$ and
$K_i:=K(g_0^{(i)},g)$ be a sequence as in the corollary to clause (0) of Theorem~0, along
which, for every $s\ge0$, one has $K_i\ge s$ for all sufficiently large $i$ and
$g_{K_i-s}\to g_{\infty,s}$ as $i\to\infty$. Then
\begin{equation}\label{5.14:eq-the-one-loop-law-for-the-cou-2}
g_{\infty,s(\mathsf{d})}^{-2}
=\frac{11N^2}{12\pi^2}\,\ln(1/\mathsf{d})\,\bigl(1+o(1)\bigr)
\qquad\text{as } \mathsf{d}\downarrow 0 .
\end{equation}
\end{corollary}

In the leading term of \eqref{5.14:eq-the-one-loop-law-for-the-cou-a} the base $L$ and every
auxiliary parameter in $\mathfrak{p}$ have dropped out. Only the universal one-loop coefficient
$11N^2/(12\pi^2)$ of $\mathrm{SU}(N)$ Yang--Mills theory remains, in Ba{\l}aban's normalisation
$\operatorname{tr}\mathbf{1}=1$, $g_0^{-2}=2N/g_{\mathrm{std}}^{2}$, where $g_{\mathrm{std}}$ is
the coupling in the standard normalisation; it corresponds to $11N/(24\pi^2)$ per e-fold of
scale for $g_{\mathrm{std}}^{-2}$.

\begin{proof}
\emph{Applicability of the flow bound.} By
\eqref{5.14:eq-the-one-loop-law-for-the-cou-1}, $K-s(\mathsf{d})\ge k_W\ge 1$, and hence
$0\le s(\mathsf{d})\le K$. Moreover $g\le\gamma_W^{\mathbb{T}}\le\gamma_J$,
$g_0\in\left]0,g_{-}(g)\right]$ and $K=K(g_0,g)$. So the hypotheses of
Theorem~\ref{5.14:thm-one-loop-law} hold, and its bound \eqref{5.14:eq-one-loop-law-b} applies
with $s=s(\mathsf{d})$.

\emph{The expansion of $x_{K-s}$.} By definition $x_k=g_k^{-2}$, so
$g(\mathsf{d})^{-2}=x_{K-s(\mathsf{d})}$. By \eqref{5.14:eq-one-loop-law-b},
\begin{equation*}
\bigl|x_{K-s}-x_K-s\,c_0\bigr|\le\Delta(s).
\end{equation*}
Hence there is a number $\theta$ with $|\theta|\le 1$ (depending on all the data) such that
\begin{equation}\label{5.14:eq-the-one-loop-law-for-the-cou-3}
\frac{1}{g(\mathsf{d})^{2}}=x_{K-s(\mathsf{d})}
=x_K+s(\mathsf{d})\,c_0+\theta\,\Delta(s(\mathsf{d})).
\end{equation}
By \eqref{5.14:eq-flow-notation-a}, the final coupling satisfies
$x_K\in\left]g^{-2},\,g^{-2}+B^{\sharp}\right]$. In particular
$0<x_K\le g^{-2}+B^{\sharp}$.

\emph{The depth tends to infinity.} By \eqref{5.14:eq-depth-separation-1} of
Lemma~\ref{5.14:lem-depth-separation},
$s(\mathsf{d})\log L=\ln(1/\mathsf{d})\,(1+o(1))$ as $\mathsf{d}\downarrow0$. The $o(1)$ there
depends on $g$ but not on $K$, $g_0$, $m$ or the position. Since $\ln(1/\mathsf{d})\to\infty$,
it follows that $s(\mathsf{d})\to\infty$ as $\mathsf{d}\downarrow0$, uniformly in the same
quantities. In particular $s(\mathsf{d})\ge 1$ for all sufficiently small $\mathsf{d}$, and we
consider only such $\mathsf{d}$. Since $c_0>0$, we may factor
\eqref{5.14:eq-the-one-loop-law-for-the-cou-3} as
\begin{equation}\label{5.14:eq-the-one-loop-law-for-the-cou-4}
\frac{1}{g(\mathsf{d})^{2}}
=s(\mathsf{d})\,c_0\Bigl[1+\frac{x_K+\theta\,\Delta(s(\mathsf{d}))}{s(\mathsf{d})\,c_0}\Bigr].
\end{equation}

\emph{The first factor of $(1+o(1))$.} Using $0<x_K\le g^{-2}+B^{\sharp}$ and $|\theta|\le1$,
\begin{equation}\label{5.14:eq-the-one-loop-law-for-the-cou-5}
\Bigl|\frac{x_K+\theta\,\Delta(s(\mathsf{d}))}{s(\mathsf{d})\,c_0}\Bigr|
\le\frac{g^{-2}+B^{\sharp}}{s(\mathsf{d})\,c_0}
+\frac{\Delta(s(\mathsf{d}))}{s(\mathsf{d})\,c_0}.
\end{equation}
Both terms tend to zero as $\mathsf{d}\downarrow0$, for three reasons. First, $g$ is fixed.
Second, $s(\mathsf{d})\to\infty$, as just shown. Third, $\Delta(s)/s\to0$ as $s\to\infty$ by
\eqref{5.14:eq-one-loop-law-b}.

The right-hand side of \eqref{5.14:eq-the-one-loop-law-for-the-cou-5} is uniform in the stated
sense. It depends on the data only through $g$, $B^{\sharp}$, $c_0$, $s(\mathsf{d})$ and the
majorant $\Delta$. By Theorem~\ref{5.14:thm-one-loop-law}, $\Delta$ depends on $(L,N,\mathfrak{p})$
only. The convergence $s(\mathsf{d})\to\infty$ is uniform in $K$, $g_0$, $m$ and the position.
Finally, $g(\mathsf{d})$ does not involve the test functions, so the pair
$(f_1,f_2)\in\mathfrak{T}(\mathsf{d};\vartheta)$ enters nowhere. Thus the bracket in
\eqref{5.14:eq-the-one-loop-law-for-the-cou-4} is $1+o(1)$.

\emph{The second factor of $(1+o(1))$.} By \eqref{5.14:eq-flow-notation-b},
$c_0=(11N^2/12\pi^2)\log L$. Hence, by \eqref{5.14:eq-flow-notation-b} and then by
\eqref{5.14:eq-depth-separation-1},
\begin{align*}
s(\mathsf{d})\,c_0
&=\frac{c_0}{\log L}\,s(\mathsf{d})\log L
=\frac{11N^2}{12\pi^2}\,s(\mathsf{d})\log L\\
&=\frac{11N^2}{12\pi^2}\,\ln(1/\mathsf{d})\,\bigl(1+o(1)\bigr),
\end{align*}
with the $o(1)$ of Lemma~\ref{5.14:lem-depth-separation}.

\emph{Conclusion.} Inserting both factors into \eqref{5.14:eq-the-one-loop-law-for-the-cou-4}
gives $g(\mathsf{d})^{-2}=(11N^2/(12\pi^2))\ln(1/\mathsf{d})(1+a)(1+b)$, where $a$ and $b$ are
$o(1)$ in the stated uniform sense. Since
\begin{equation*}
(1+a)(1+b)=1+a+b+ab
\quad\text{and}\quad
|a+b+ab|\le|a|+|b|+|a|\,|b|,
\end{equation*}
the product is again $1+o(1)$. This proves \eqref{5.14:eq-the-one-loop-law-for-the-cou-a}. For
$\mathsf{d}<1$ we have $\ln(1/\mathsf{d})>0$, and dividing
\eqref{5.14:eq-the-one-loop-law-for-the-cou-a} by $\ln(1/\mathsf{d})$ gives the equivalent form
$g(\mathsf{d})^{-2}/\ln(1/\mathsf{d})\to 11N^2/(12\pi^2)$.

\emph{The limit couplings.} Apply \eqref{5.14:eq-one-loop-law-c} with $s=s(\mathsf{d})$ to the
sequence $g_0^{(i)}\downarrow0$, along which $K_i\ge s(\mathsf{d})$ for all sufficiently
large $i$. This gives
\begin{equation*}
g_{\infty,s}^{-2}\in\bigl[\,g^{-2}+s\,c_0-\Delta(s),\;
g^{-2}+B^{\sharp}+s\,c_0+\Delta(s)\,\bigr],\qquad s=s(\mathsf{d}).
\end{equation*}
The interval $\left]g^{-2},g^{-2}+B^{\sharp}\right]$ of \eqref{5.14:eq-flow-notation-a}
contains $x_{K_i}$ for every $i$, hence $B^{\sharp}>0$ and $g^{-2}+B^{\sharp}>0$. Since also
$g^{-2}>0$, we obtain
$|g_{\infty,s}^{-2}-s\,c_0|\le g^{-2}+B^{\sharp}+\Delta(s)$. Hence, for $s(\mathsf{d})\ge1$,
\begin{equation*}
g_{\infty,s(\mathsf{d})}^{-2}
=s(\mathsf{d})\,c_0\Bigl[1+\frac{g_{\infty,s(\mathsf{d})}^{-2}-s(\mathsf{d})\,c_0}
{s(\mathsf{d})\,c_0}\Bigr],
\end{equation*}
and the quotient in the bracket is bounded by the right-hand side of
\eqref{5.14:eq-the-one-loop-law-for-the-cou-5}. From here the two preceding steps apply word for
word, and give \eqref{5.14:eq-the-one-loop-law-for-the-cou-2}.
\end{proof}

\begin{corollary}[The charge read off the two-point function runs at the one-loop rate within a fixed factor]
\label{5.14:cor-the-charge-read-off-the-two}
Assume the following three statements.
\begin{enumerate}
\item[(a)] Clause (iv)(a) of Theorem~0, in the form of the display of the two-point witness theorem
restated below. The setting is the following. We have $N=2$,
$g\le\gamma_W^{\mathbb{T}}$, $g_0\in\left]0,g_-(g)\right]$, $K=K(g_0,g)$, every $m$, and every
$\mathsf{d}$ satisfying the depth floor and the ultraviolet floor of that clause. We put
$s:=s(\mathsf{d})$ and take any $(f_1,f_2)\in\mathfrak{T}(\mathsf{d};\vartheta)$. In this setting
\begin{equation*}
S^T_{2;K,m}(f_1,f_2)=\zeta'^{\,2}_{K-s}\,b_0\,g_{K-s}^4\,\mathcal{K}_s(f_1,f_2)\,(1+R_K),
\qquad |R_K|\le\tfrac12,\quad \zeta'_{K-s}\in\left[\tfrac23,\tfrac43\right].
\end{equation*}
This holds uniformly in $K$, $g_0$ and the position of the supports. It holds uniformly also in $m$ on
$m_{\mathbb{T}}(g_-(g))\le m\le m_V(g)$, and on each torus beyond that range under
$s(\mathsf{d})\ge s_{\mathrm{vol}}(g,m)$. Here $\mathcal{K}_s$ is the free lattice kernel, with
\begin{equation*}
c_{\mathcal{K}}\|f_1\|_\infty\|f_2\|_\infty\le\mathcal{K}_s(f_1,f_2)
\le C_{\mathcal{K}}\|f_1\|_\infty\|f_2\|_\infty ,
\end{equation*}
and $\mathcal{K}_s(f_1,f_2)>0$.
\item[(b)] Corollary~\ref{5.14:cor-the-one-loop-law-for-the-cou} (the one-loop law for the coupling of
the two-point display). In the same setting, with $g(\mathsf{d}):=g_{K-s(\mathsf{d})}$, it states
\eqref{5.14:eq-the-one-loop-law-for-the-cou-a}:
\begin{equation*}
\frac{1}{g(\mathsf{d})^2}=\frac{11N^2}{12\pi^2}\,\ln\frac{1}{\mathsf{d}}\,\bigl(1+o(1)\bigr)
\qquad(\mathsf{d}\downarrow 0\ \text{at fixed}\ g).
\end{equation*}
This holds uniformly in $K\ge s(\mathsf{d})+k_W$, $g_0$, $m$ and the position. The same law holds for
the subsequential limits $g_{\infty,s(\mathsf{d})}$ of clause (0) of Theorem~0.
\item[(c)] Clause (ii-$\mathbb{T}$-b) of Theorem~0: for each $m$ in the range of (a) there is a
subsequence of cutoffs $K_j\to\infty$ on $\mathbb{T}_m$ along which $g_{K_j-s}\to g_{\infty,s}$ for
every $s$.
\end{enumerate}
In this setting define the \emph{charge read off the two-point function} at cutoff $K$ by
\begin{equation}\label{5.14:eq-the-charge-read-off-the-two-1}
g_{\mathrm{phys}}(\mathsf{d})^4:=\frac{S^T_{2;K,m}(f_1,f_2)}{b_0\,\mathcal{K}_{s(\mathsf{d})}(f_1,f_2)} .
\end{equation}
Then the following hold.
\begin{enumerate}
\item[(1)] The two fourth powers are comparable within fixed factors:
\begin{equation}\label{5.14:eq-the-charge-read-off-the-two-2}
\frac29\le\frac{g_{\mathrm{phys}}(\mathsf{d})^4}{g(\mathsf{d})^4}\le\frac83 .
\end{equation}
\item[(2)] As $\mathsf{d}\downarrow 0$ at fixed $g$,
\begin{equation}\label{5.14:eq-the-charge-read-off-the-two-3}
\Bigl(\frac38\Bigr)^{1/2}\bigl(1+o(1)\bigr)
\le\frac{g_{\mathrm{phys}}(\mathsf{d})^{-2}}{\frac{11N^2}{12\pi^2}\ln\frac{1}{\mathsf{d}}}
\le\Bigl(\frac92\Bigr)^{1/2}\bigl(1+o(1)\bigr).
\end{equation}
This holds uniformly in $K\ge s(\mathsf{d})+k_W$, $g_0$, the position, the pair
$(f_1,f_2)\in\mathfrak{T}(\mathsf{d};\vartheta)$, and $m$ in the range of (a).
\item[(3)] (Torus limits.) Fix $m$ in the range of (a), and fix the subsequence of cutoffs
$K_j\to\infty$ on $\mathbb{T}_m$ of hypothesis (c), along which
$g_{K_j-s}\to g_{\infty,s}$ for every $s$. For $\mathsf{d}$ as in (a) put $s=s(\mathsf{d})$, take
$(f_1,f_2)\in\mathfrak{T}(\mathsf{d};\vartheta)$, and let $\mathcal{K}^-$ and $\mathcal{K}^+$ be the
lower and upper limits of $\mathcal{K}_s(f_1,f_2)$ along this subsequence. Then
\begin{equation*}
c_{\mathcal{K}}\|f_1\|_\infty\|f_2\|_\infty\le\mathcal{K}^-\le\mathcal{K}^+
\le C_{\mathcal{K}}\|f_1\|_\infty\|f_2\|_\infty ,
\end{equation*}
and every limit point $S(\mathsf{d})$ of $S^T_{2;K_j,m}(f_1,f_2)$ satisfies
\begin{equation}\label{5.14:eq-the-charge-read-off-the-two-4}
\tfrac29\,b_0\,g_{\infty,s}^4\,\mathcal{K}^-\le S(\mathsf{d})\le\tfrac83\,b_0\,g_{\infty,s}^4\,\mathcal{K}^+ .
\end{equation}
Now fix a reference pair $(\phi_1,\phi_2)$ of test functions, and for each $\mathsf{d}$ let
$f_i^{(\mathsf{d})}:=\phi_i(\,\cdot\,/\mathsf{d})$, composed with a translation placing the two
supports at separation $\mathsf{d}$, be the dilate of $\phi_i$ by the factor $\mathsf{d}$; thus
$\|f_i^{(\mathsf{d})}\|_\infty=\|\phi_i\|_\infty$. Assume that $(\phi_1,\phi_2)$ is such that
$(f_1^{(\mathsf{d})},f_2^{(\mathsf{d})})\in\mathfrak{T}(\mathsf{d};\vartheta)$ for all sufficiently
small $\mathsf{d}$, and take $(f_1,f_2)=(f_1^{(\mathsf{d})},f_2^{(\mathsf{d})})$ above. In this item
we write $g_{\mathrm{phys}}(\mathsf{d})^4$ for $S(\mathsf{d})/(b_0\mathcal{K})$, where
$\mathcal{K}\in[\mathcal{K}^-,\mathcal{K}^+]$ is arbitrary. Then, as $\mathsf{d}\downarrow0$,
\begin{align*}
\Bigl(\frac38\Bigr)^{1/2}\Bigl(\frac{\mathcal{K}^-}{\mathcal{K}^+}\Bigr)^{1/2}\bigl(1+o(1)\bigr)
&\le\frac{g_{\mathrm{phys}}(\mathsf{d})^{-2}}{\frac{11N^2}{12\pi^2}\ln\frac{1}{\mathsf{d}}}\\
&\le\Bigl(\frac92\Bigr)^{1/2}\Bigl(\frac{\mathcal{K}^+}{\mathcal{K}^-}\Bigr)^{1/2}\bigl(1+o(1)\bigr),
\end{align*}
with $(\mathcal{K}^+/\mathcal{K}^-)^{1/2}\in[1,(C_{\mathcal{K}}/c_{\mathcal{K}})^{1/2}]$. If
$S(\mathsf{d})/(b_0\|\phi_1\|_\infty\|\phi_2\|_\infty)$ is taken in place of
$g_{\mathrm{phys}}(\mathsf{d})^4$, then, as $\mathsf{d}\downarrow0$,
\begin{equation*}
\Bigl(\frac38\Bigr)^{1/2}C_{\mathcal{K}}^{-1/2}\bigl(1+o(1)\bigr)
\le\frac{g_{\mathrm{phys}}(\mathsf{d})^{-2}}{\frac{11N^2}{12\pi^2}\ln\frac{1}{\mathsf{d}}}
\le\Bigl(\frac92\Bigr)^{1/2}c_{\mathcal{K}}^{-1/2}\bigl(1+o(1)\bigr).
\end{equation*}
\end{enumerate}
In words: by (b) the coupling $g(\mathsf{d})$ of clause (iv)(a) obeys
$1/g(\mathsf{d})^2=(11N^2/12\pi^2)\ln(1/\mathsf{d})\,(1+o(1))$ as $\mathsf{d}\to0$,
uniformly in the cutoff; by (1) the charge read off the two-point function has fourth power differing
from $g(\mathsf{d})^4$ by a factor in $[\frac29,\frac83]$; and by (2) it therefore runs with an
inverse-square slope between $(3/8)^{1/2}$ and $(9/2)^{1/2}$ times the one-loop Yang--Mills value
$11N^2/12\pi^2$: within a fixed factor of it. It is not asserted that the charge read off the
two-point function has exactly the one-loop slope; that would require $R_K\to0$ and
$\zeta'_{K-s}\to1$ in the display of hypothesis (a), together with the convergence of the kernel
$\mathcal{K}_{s(\mathsf{d})}$.
\end{corollary}

\begin{proof}
\emph{Positivity.} By hypothesis (a), $\mathcal{K}_{s(\mathsf{d})}(f_1,f_2)>0$, so
\eqref{5.14:eq-the-charge-read-off-the-two-1} is well defined. Since $\zeta'^{\,2}_{K-s}\ge\frac49$
and $1+R_K\ge\frac12$, the display in (a) gives $S^T_{2;K,m}(f_1,f_2)>0$. Hence
$g_{\mathrm{phys}}(\mathsf{d})^4>0$.

\emph{Proof of \textup{(1)}.} Divide the display in (a) by $b_0\,\mathcal{K}_s(f_1,f_2)\,g_{K-s}^4$,
with $s=s(\mathsf{d})$ and $g_{K-s}=g(\mathsf{d})$. This gives
\begin{equation*}
\frac{g_{\mathrm{phys}}(\mathsf{d})^4}{g(\mathsf{d})^4}=\zeta'^{\,2}_{K-s}\,(1+R_K).
\end{equation*}
From $\zeta'_{K-s}\in[\frac23,\frac43]$ we get $\zeta'^{\,2}_{K-s}\in[\frac49,\frac{16}9]$. From
$|R_K|\le\frac12$ we get $1+R_K\in[\frac12,\frac32]$. Both factors are positive, so their product
lies between the products of the endpoints:
\begin{equation*}
\frac49\cdot\frac12=\frac29
\qquad\text{and}\qquad
\frac{16}9\cdot\frac32=\frac83 .
\end{equation*}
This is \eqref{5.14:eq-the-charge-read-off-the-two-2}.

\emph{Proof of \textup{(2)}.} The map $t\mapsto t^{-1/2}$ is decreasing on $\left]0,\infty\right[$.
Applying it to \eqref{5.14:eq-the-charge-read-off-the-two-2} reverses the order:
\begin{equation*}
\Bigl(\frac83\Bigr)^{-1/2}\le\frac{g_{\mathrm{phys}}(\mathsf{d})^{-2}}{g(\mathsf{d})^{-2}}
\le\Bigl(\frac29\Bigr)^{-1/2}.
\end{equation*}
Moreover $(8/3)^{-1/2}=(3/8)^{1/2}$ and $(2/9)^{-1/2}=(9/2)^{1/2}$.

Write the conclusion of hypothesis (b) as
\begin{equation*}
g(\mathsf{d})^{-2}=\frac{11N^2}{12\pi^2}\ln\frac1{\mathsf{d}}\,\bigl(1+\varepsilon(\mathsf{d})\bigr).
\end{equation*}
Here $\varepsilon(\mathsf{d})\to0$ as $\mathsf{d}\downarrow0$, uniformly in $K\ge s(\mathsf{d})+k_W$,
$g_0$, $m$ and the position. For $\mathsf{d}<1$ both $g(\mathsf{d})^{-2}$ and $\ln(1/\mathsf{d})$
are positive, so $1+\varepsilon(\mathsf{d})>0$. Multiplying the last chain by the positive number
\begin{equation*}
\frac{g(\mathsf{d})^{-2}}{\frac{11N^2}{12\pi^2}\ln\frac1{\mathsf{d}}}=1+\varepsilon(\mathsf{d})
\end{equation*}
gives \eqref{5.14:eq-the-charge-read-off-the-two-3} with $o(1)=\varepsilon(\mathsf{d})$.

The uniformity follows from two facts. The window $[(3/8)^{1/2},(9/2)^{1/2}]$ is a pair of
absolute numbers, valid for every admissible $K$, $g_0$, $m$, position and pair by (a). The function
$\varepsilon$ is uniform in the stated parameters by (b), and does not depend on the pair.

The two ends of the window in \eqref{5.14:eq-the-charge-read-off-the-two-3} differ by the factor
\begin{equation*}
\frac{(9/2)^{1/2}}{(3/8)^{1/2}}=\Bigl(\frac92\cdot\frac83\Bigr)^{1/2}=12^{1/2}.
\end{equation*}

\emph{Proof of \textup{(3)}.} Along the sequence $K_j$ the display in (a) and
\eqref{5.14:eq-the-charge-read-off-the-two-2} give, for every $j$ with $K_j\ge s(\mathsf{d})+k_W$,
hence for all large $j$,
\begin{equation*}
\tfrac29\,b_0\,g_{K_j-s}^4\,\mathcal{K}_s(f_1,f_2)\le S^T_{2;K_j,m}(f_1,f_2)
\le\tfrac83\,b_0\,g_{K_j-s}^4\,\mathcal{K}_s(f_1,f_2).
\end{equation*}
Here $\mathcal{K}_s(f_1,f_2)$ is evaluated at cutoff $K_j$. Since $g_{K_j-s}\to g_{\infty,s}$ by
hypothesis (c), the
lower and upper limits of the outer members are $\frac29 b_0 g_{\infty,s}^4\mathcal{K}^-$ and
$\frac83 b_0 g_{\infty,s}^4\mathcal{K}^+$. Every limit point of the middle member therefore lies
between them, which is \eqref{5.14:eq-the-charge-read-off-the-two-4}. Passing to the lower and upper
limits in the kernel bounds of (a) gives
\begin{equation*}
c_{\mathcal{K}}\|f_1\|_\infty\|f_2\|_\infty\le\mathcal{K}^-\le\mathcal{K}^+
\le C_{\mathcal{K}}\|f_1\|_\infty\|f_2\|_\infty ,
\end{equation*}
and $\mathcal{K}^->0$ because $c_{\mathcal{K}}\|f_1\|_\infty\|f_2\|_\infty>0$.

Take $s=s(\mathsf{d})$ and the pair $(f_1,f_2)=(f_1^{(\mathsf{d})},f_2^{(\mathsf{d})})$. By the part
of hypothesis (b) concerning the limits of clause (0), write
\begin{equation*}
g_{\infty,s(\mathsf{d})}^{-2}
=\frac{11N^2}{12\pi^2}\ln\frac1{\mathsf{d}}\,\bigl(1+\varepsilon_\infty(\mathsf{d})\bigr)
\end{equation*}
with $\varepsilon_\infty(\mathsf{d})\to0$. For $\mathsf{d}$ small enough the right side is finite
and positive, so $g_{\infty,s(\mathsf{d})}>0$ and
\begin{equation*}
g_{\infty,s(\mathsf{d})}^4=\Bigl(\frac{11N^2}{12\pi^2}\ln\frac1{\mathsf{d}}\Bigr)^{-2}
\bigl(1+\varepsilon_\infty(\mathsf{d})\bigr)^{-2},
\end{equation*}
where $(1+\varepsilon_\infty(\mathsf{d}))^{-2}=1+o(1)$. In particular $S(\mathsf{d})>0$ by
\eqref{5.14:eq-the-charge-read-off-the-two-4}. Dividing \eqref{5.14:eq-the-charge-read-off-the-two-4}
by $b_0\mathcal{K}$ with $\mathcal{K}\in[\mathcal{K}^-,\mathcal{K}^+]$, and using
$\mathcal{K}^-/\mathcal{K}\ge\mathcal{K}^-/\mathcal{K}^+$ and
$\mathcal{K}^+/\mathcal{K}\le\mathcal{K}^+/\mathcal{K}^-$, we obtain
\begin{equation*}
\frac29\,\frac{\mathcal{K}^-}{\mathcal{K}^+}\,g_{\infty,s}^4\le g_{\mathrm{phys}}(\mathsf{d})^4
\le\frac83\,\frac{\mathcal{K}^+}{\mathcal{K}^-}\,g_{\infty,s}^4 .
\end{equation*}
Taking the $(-\frac12)$-power reverses the order. Inserting the expression for
$g_{\infty,s(\mathsf{d})}^{-2}$ gives the window of (3), with
$(\mathcal{K}^+/\mathcal{K}^-)^{1/2}\in[1,(C_{\mathcal{K}}/c_{\mathcal{K}})^{1/2}]$ by the bounds on
$\mathcal{K}^\mp$ just derived.

For the sup-norm normalisation, use $\mathcal{K}^-\ge c_{\mathcal{K}}\|f_1\|_\infty\|f_2\|_\infty$ in
the lower bound of \eqref{5.14:eq-the-charge-read-off-the-two-4} and
$\mathcal{K}^+\le C_{\mathcal{K}}\|f_1\|_\infty\|f_2\|_\infty$ in the upper bound, together with
$\|f_i\|_\infty=\|f_i^{(\mathsf{d})}\|_\infty=\|\phi_i\|_\infty$. This gives
\begin{equation*}
\frac29\,c_{\mathcal{K}}\,g_{\infty,s}^4\le\frac{S(\mathsf{d})}{b_0\|\phi_1\|_\infty\|\phi_2\|_\infty}
\le\frac83\,C_{\mathcal{K}}\,g_{\infty,s}^4 .
\end{equation*}
The $(-\frac12)$-power and the expression for $g_{\infty,s(\mathsf{d})}^{-2}$ give, with
$S(\mathsf{d})/(b_0\|\phi_1\|_\infty\|\phi_2\|_\infty)$ in place of $g_{\mathrm{phys}}(\mathsf{d})^4$,
\begin{equation*}
\Bigl(\frac38\Bigr)^{1/2}C_{\mathcal{K}}^{-1/2}\bigl(1+o(1)\bigr)
\le\frac{g_{\mathrm{phys}}(\mathsf{d})^{-2}}{\frac{11N^2}{12\pi^2}\ln\frac{1}{\mathsf{d}}}
\le\Bigl(\frac92\Bigr)^{1/2}c_{\mathcal{K}}^{-1/2}\bigl(1+o(1)\bigr).
\end{equation*}

In each of these normalisations the inverse square of the charge is therefore within a fixed factor
of $\frac{11N^2}{12\pi^2}\ln\frac1{\mathsf{d}}$. The closing sentences of the statement restate
hypothesis (b), \eqref{5.14:eq-the-charge-read-off-the-two-2} and
\eqref{5.14:eq-the-charge-read-off-the-two-3}.
\end{proof}

\begin{definition}[Levels, zones, cells and the standing geometric axioms]
\label{5.4:not-zones-cells}
Throughout, $d \ge 1$ and $L \ge 3$ are integers, and $|\cdot| := |\cdot|_{\infty}$. The
statements of this section hold for every such $d$ and $L$; in the application $d = 4$ and
$L$ is the blocking factor. In this section the letters $m$, $m'$, $i$, $j$ denote
levels, and $m_0 \le k$ is the lowest level. The \emph{current label} is the label at the
localisation site (L) of the operation $\mathbf{R}_k$ (Definition~\ref{5.2:def-sites-units});
the current cells are defined from it as in
(S3) below. The \emph{carrier} is the domain of the fields on which $\mathbf{R}_k$ acts. The
following conventions and axioms are in force in all later statements.

\smallskip\noindent
\emph{Levels and grids.} The levels are the integers $m$ with $m_0 \le m \le k$. For such
$m$, $\pi_m$ is the family of closed cubes of the level-$m$ grid; each has side $1$ in
level-$m$ units, that is, side $L^{m-k}$ in level-$k$ units. We write
$\operatorname{dist}_m$ and $\operatorname{diam}_m$ for the set distance and the diameter
with respect to $|\cdot|_{\infty}$, measured in level-$m$ units; thus
$\operatorname{dist}_m = L^{k-m}\operatorname{dist}_k$. For a cell $\Delta$ (see (S3)) we
write $\operatorname{lev}\Delta$ for its level. The grids and the current cells satisfy
\begin{equation}\label{5.4:eq-zones-cells-a}
\begin{aligned}
&\text{(S1)}\ \ \text{for } m \le m' \text{, every } q' \in \pi_{m'} \text{ is the union}\\
&\phantom{\text{(S1)}\ \ } \text{of the } L^{d(m'-m)} \text{ cubes of } \pi_m
\text{ it contains;}\\
&\phantom{\text{(S1)}\ \ } q \in \pi_m,\ q' \in \pi_{m'},\
\operatorname{int} q \cap \operatorname{int} q' \neq \emptyset
\ \Longrightarrow\ q \subset q';\\
&\text{(S2)}\ \ q \neq q' \text{ in } \pi_m \ \Longrightarrow\
\operatorname{int} q \cap \operatorname{int} q' = \emptyset;\\
&\phantom{\text{(S2)}\ \ } q, q_1, \dots, q_n \in \pi_m,\ q \subset q_1 \cup \dots \cup q_n
\ \Longrightarrow\ q \in \{q_1, \dots, q_n\};\\
&\text{(S3)}\ \ \text{each current cell is a closed cube of exactly one level;}\\
&\phantom{\text{(S3)}\ \ } \text{a cell of level } m \text{ is a cube of } \pi_m
\text{ contained in } \operatorname{cl}(\text{zone } m)\text{;}\\
&\phantom{\text{(S3)}\ \ } \text{distinct cells have disjoint interiors; the cells cover
the carrier.}
\end{aligned}
\end{equation}
Here zone $m$ is the set defined below, before \eqref{5.4:eq-zones-cells-b}. The cells form
the partition into cubes in proper scales of \cite{B-CMP119}; the hubs are unions of cells.
The family of cells is locally finite, and we set
\begin{equation*}
\Omega_m^{\mathrm{cur}} := \bigcup\{\Delta : \Delta \text{ a cell}, \ \operatorname{lev}\Delta \ge m\},
\end{equation*}
a closed set.

\smallskip\noindent
\emph{The current label.} The label consists of sets $\Omega_m$, $\Lambda_m$. With the
conventions $\Omega_{k+1} := \emptyset$ and $\Omega_{m_0-1} :=$ the whole carrier, we put
\begin{equation*}
\text{zone } m := \Omega_m \setminus \Omega_{m+1}, \qquad U_m := \operatorname{cl}(\Omega_m^{c}),
\end{equation*}
and the following hold for every level $m$, every $i < k$ and all cells $\Delta, \Delta'$;
moreover, (A0) and (A3) hold not only in the current label but in every label occurring in
the construction:
\begin{equation}\label{5.4:eq-zones-cells-b}
\begin{aligned}
&\text{(A0)}\ \ \Omega_{m+1} \subset \Lambda_m \subset \Omega_m;
\qquad \text{(A1)}\ \ \Omega_{m+1} \subset \Omega_m;\\
&\text{(A2)}\ \ \Delta \cap \Delta' \neq \emptyset \ \Longrightarrow\
|\operatorname{lev}\Delta - \operatorname{lev}\Delta'| \le 1;\\
&\text{(A3)}\ \ \operatorname{dist}_{i+1}(U_i, \Omega_{i+1}) \ge 1 .
\end{aligned}
\end{equation}
In the form given by
Ba{\l}aban, (A3) says that $U_i$ and $\Omega_{i+1}$ are separated by at least one cube of
$\pi_{i+1}$ \cite{B-CMP122a}. The sets $\Omega_m$ are closed, so that
$\operatorname{cl}(\Omega_m) = \Omega_m$; this is part of the standing assumptions. The
following consequences hold:
\begin{itemize}
\item $\Omega_m^{\mathrm{cur}} \subset \Omega_m$;
\item a cell $\Delta$ of level $m$ satisfies
$\Delta \subset \operatorname{cl}(\text{zone } m) \subset \Omega_m$ and $\Delta \subset U_{m+1}$;
\item for $i < k$, every subset of $U_i$ and every subset of $\Omega_{i+1}$ are at
$\operatorname{dist}_{i+1}$-distance at least $1$.
\end{itemize}

\smallskip\noindent
\emph{Incidence.} Two closed sets \emph{meet} if their interiors intersect, \emph{touch} if
they intersect, and \emph{share a wall} if their intersection is $(d-1)$-dimensional. A cell
has at most $2dL^{d-1}$ cells sharing a wall with it. The constant $\hat{b}$ satisfies
$\hat{b} \ge 2dL^{d-1}$, and $\hat{b}_t := 3^d + (L+2)^d$ bounds the number of cells touching
a given cell.

\smallskip\noindent
\emph{Polymers.} $\mathbf{D}_j$ is the set of wall-connected finite families of
$\pi_j$-cubes. For $X \in \mathbf{D}_j$, $N_j(X)$ is its number of cubes and $d_j(X)$ its
tree length; for $m \ge j$, $[X]_m$ is the family of $\pi_m$-cubes containing cubes of $X$.
With a constant $\mathsf{c}_d \ge 1$ depending only on $d$,
\begin{equation}\label{5.4:eq-zones-cells-c}
\begin{aligned}
&\text{(g1)}\ \ N_j(X) \le \mathsf{c}_d\,(d_j(X) + 1);\\
&\text{(g4)}\ \ d_j(X) \ge \operatorname{dist}_j(Q, Q') \ \text{ for all cubes } Q, Q'
\text{ of } X.
\end{aligned}
\end{equation}

\smallskip\noindent
\emph{The treated class, hubs and source cells.} $Z$ is the treated class, a union of
$\pi_k$-cubes, with boundary $\partial Z$. A $\pi_k$-cube $P \not\subset Z$ satisfies
$\operatorname{int} P \cap Z = \emptyset$, hence $P \subset \operatorname{cl}(Z^c)$. The hubs
$\Lambda_C^{\sim}$, one for each treated component $C$ of $Z$, lie in $Z$ and are never cut;
$\Lambda^{\sim}$ is their union. In (Z)
below, $\Lambda$ denotes the set of the current label from which the hubs $\Lambda_C^{\sim}$
are formed. For each level $m$, $\check{R}_m$ denotes the radius of level $m$; it enters the
separations (G0), ($\Gamma$3) and ($\Gamma$5) below and is subject to ($\check{R}$1) and
($\check{R}$2) in \eqref{5.4:eq-zones-cells-g}. We assume
\begin{equation}\label{5.4:eq-zones-cells-e}
\begin{aligned}
&\text{(Z)}\ \ \operatorname{dist}_k(\Lambda^{\sim}, Z^c) \ge 9\check{R}_k;\\
&\phantom{\text{(Z)}\ \ } \text{every cell disjoint with } \Lambda \text{ lying in } Z
\text{ has level } k.
\end{aligned}
\end{equation}
$\pi^{\mathfrak{s}}$ is the family of cells disjoint with $\Lambda^{\sim}$. A \emph{source
cell} is a member of $\pi^{\mathfrak{s}}$ sharing a wall with a hub; a source cell lies within
$\operatorname{dist}_k$-distance $1$ of $\Lambda^{\sim}$, hence at $\operatorname{dist}_k$-distance
at least $9\check{R}_k - 1$ from $Z^c$; that every source cell has level $k$ is part of the
standing assumptions. Two members of
$\pi^{\mathfrak{s}}$ are \emph{adjacent} if they share a wall, or if both share a wall with
one hub. For $\Delta \in \pi^{\mathfrak{s}}$, $\rho(\Delta)$ is the least cardinality of a
connected finite family $\subset \pi^{\mathfrak{s}}$ containing $\Delta$ and a source cell.
Then:
\begin{equation}\label{5.4:eq-zones-cells-d}
\begin{aligned}
&\text{(}\rho\text{0)}\ \ \rho(\Delta) \ge n \ \text{ whenever every such family has at least }
n \text{ members;}\\
&\text{(}\rho\text{1)}\ \ \rho(\Delta) \ge k - \operatorname{lev}\Delta + 1;\\
&\text{(}\rho\text{2)}\ \ \rho(\Delta) \ge \operatorname{dist}_k\bigl(\Delta,
\textstyle\bigcup\{\text{source cells}\}\bigr);\\
&\text{(}\rho\text{3)}\ \ \text{there are at most } (e\hat{b})^n \text{ connected families of }
n \text{ members}\\
&\phantom{\text{(}\rho\text{3)}\ \ } \text{through a given member.}
\end{aligned}
\end{equation}
($\rho$0), ($\rho$1) and ($\rho$2) are proved below; ($\rho$3) is not derived here from the
definition of adjacency: it is part of the standing assumptions, with the constant $\hat{b}$
of the incidence bounds above.

\smallskip\noindent
\emph{Separation near $\partial Z$, raised cells and arrests.} The attempts (running,
arrested, excluded or completed), their live arrests, the cells raised by attempts, the
foreign and sibling components and the frozen levels referred to below are those of the
arrest construction of the large-field expansion. We assume
\begin{equation}\label{5.4:eq-zones-cells-f}
\begin{aligned}
&\text{(}\Gamma\text{5)}\ \ \operatorname{dist}_k(\partial Z, \Omega_k^c) \ge 2\check{R}_k
\ \text{ on both sides of } \partial Z,\\
&\phantom{\text{(}\Gamma\text{5)}\ \ } \text{and no cell within } 2\check{R}_k \text{ of }
\partial Z \text{ is raised}
\end{aligned}
\end{equation}
(compare \cite{B-CMP119}). Consequently every cell of level $< k$ and every raised cell has
$\operatorname{dist}_k(\cdot, \partial Z) \ge 2\check{R}_k$, and every $\pi_k$-cube at
$\operatorname{dist}_k$-distance $< 2\check{R}_k$ from $\partial Z$ is a current un-raised
cell of level $k$.

For a live arrest $\mathfrak{a}$ (of any class), with level $k_a$, let $X_{\mathfrak{a}}$ be
its component of $\Lambda_{k_a}^c$, and put, with $\Omega^{\mathbf{T}}_{k_a}$ the
$\mathbf{T}$-label set of level $k_a$,
\begin{equation*}
W_{\mathfrak{a}} := X_{\mathfrak{a}} \cap (\Omega^{\mathbf{T}}_{k_a})^c \ \ (\text{window}),
\qquad
\mathfrak{c}_{\mathfrak{a}} := X_{\mathfrak{a}} \cap \Omega^{\mathbf{T}}_{k_a} \ \ (\text{collar});
\end{equation*}
the collar consists of two layers of $M\check{R}_{k_a}$-cubes, of level $k_a$. We assume,
for every level $m$, every current cell $c$, every live arrest $\mathfrak{a}$ and every live
arrest $\mathfrak{a}'$ on a component not in $Z$ (a foreign one, or a sibling one with
$k_{a'} = k$):
\begin{equation}\label{5.4:eq-zones-cells-g}
\begin{aligned}
&\text{(G0)}\ \ \text{the components of } \Lambda_m^c \text{ are at distance at least }
\check{R}_m \text{ from one another;}\\
&\text{(}\Gamma\text{1)}\ \ c \text{ is un-raised, or raised by exactly one live attempt }
\mathfrak{t} \text{ (running or}\\
&\phantom{\text{(}\Gamma\text{1)}\ \ } \text{live-arrested), and then } c \subset X_{\mathfrak{t}},\
\operatorname{lev} c \le k_{\mathfrak{t}};\\
&\text{(}\Gamma\text{3)}\ \ \operatorname{dist}_{k_a}(W_{\mathfrak{a}}, X_{\mathfrak{a}}^c)
\ge 2\check{R}_{k_a};\\
&\text{(}\Gamma\text{3}'\text{)}\ \ \text{in the current label every cell meeting }
X_{\mathfrak{a}'} \text{ lies in } X_{\mathfrak{a}'}\\
&\phantom{\text{(}\Gamma\text{3}'\text{)}\ \ } \text{and has level} \le k_{a'}, \text{ and }
\operatorname{dist}_{k_{a'}+1}\bigl(X_{\mathfrak{a}'}, \Omega^{\mathrm{cur}}_{k_{a'}+1}\bigr) > 0;\\
&\text{(}\Gamma\text{4)}\ \ L^{N_0(\mathfrak{a}')} \le L\,C_{\uparrow}\,\check{R}_{k_{a'}};\\
&\text{(}\check{R}\text{1)}\ \ \check{R}_{m-1} \le L\,\check{R}_m \ \ (m > m_0);
\qquad
\text{(}\check{R}\text{2)}\ \ \check{R}_m \ge 38.
\end{aligned}
\end{equation}
Here $X_{\mathfrak{t}}$ and $k_{\mathfrak{t}}$ are the component and the level of the attempt
$\mathfrak{t}$, and
$N_0(\mathfrak{a}')$ is the integer attached to the arrest $\mathfrak{a}'$ in its construction.

\smallskip\noindent
\emph{Changes of the label.} The following list of the events that change the label,
referred to below as (E), is complete:
\begin{equation}\label{5.4:eq-zones-cells-h}
\begin{aligned}
&\text{(E)}\ \ \text{between the creation of a term and the site (L), the label changes}\\
&\phantom{\text{(E)}\ \ } \text{only through the events } (\alpha)\text{--}(\delta)
\text{ below.}
\end{aligned}
\end{equation}
\begin{itemize}
\item[($\alpha$)] $\mathbf{T}$-steps $m \to m+1$, which define the new top pair
$(\Omega_{m+1}, \Lambda_{m+1})$ and leave $(\Omega_i, \Lambda_i)_{i \le m}$ unchanged
\cite{B-CMP119}, \cite{B-CMP122a}.
\item[($\beta$)] Preparations of attempts $b$ (running ones, ones later arrested, excluded
ones, completed ones) with class $Z_b$ and range $]h, k_b]$. By
\cite[(1.10)--(1.12), p.~179]{B-CMP122a}, in the notation of that paper, the label
$(\Omega''_j)_j$ after the preparation is
\begin{equation*}
\begin{aligned}
\Omega''_j &= \Omega_j && (j \le h),\\
\Omega''_j &= \bigl((Z''_j)^c \cap Z_b\bigr) \cup \bigl(\Omega_j \cap Z_b^c\bigr)
&& (h < j \le k_b),\\
Z''_j &= (\Omega_j^{\sim 5})^c \cap Z_b && (h < j \le k_0),\\
Z''_j &\subset Z''_{k_b} = (\Omega_{k_0+1}^{\sim 5})^c \cap Z_b && (k_0 < j \le k_b),
\end{aligned}
\end{equation*}
where $k_0$ is the level of that paper at which the two ranges of $j$ in the last two lines
divide, and, writing $\ell$ for the level denoted $k$ in that paper, $(Z''_{\ell})^c$ and
$Z''_{\ell-1}$ are separated by one layer, and so on for the lower levels.
\item[($\gamma$)] Completions, at which, by \cite[(1.33), p.~365]{B-CMP122b}, in the
notation of that paper, $\mathbf{B} := \mathbf{B}''$ on $Z^c$ and $\{Z^{(k')}\}$ on $Z$, and
$\Gamma_j = \Gamma''_j \cap Z^c$ for $j < k'$.
\item[($\delta$)] Arrests, exclusions and freezing, which do not re-label the frozen levels
\cite{B-CMP122a}; moreover, a later attempt over a component containing $X_{\mathfrak{a}'}$
has $h \ge k_{a'}$.
\end{itemize}

\smallskip\noindent
\emph{Constants used in the sequel.} The following letters are used in the sequel. We write
\begin{equation*}
S_d := \sum_{x \ge 0} (2x+3)^d e^{-x}
\end{equation*}
and
\begin{equation}\label{5.4:eq-zones-cells-j}
S_d^{(1/2)} := \sum_{x \ge 0} (2x+3)^d e^{-x/2}
\end{equation}
for the grid sums. We put
\begin{equation*}
c_l^{\circ} = 1 + d(h_{\sharp} + 1), \qquad c_l = 1 + 234d, \qquad a := \tfrac{31}{20}\,\kappa,
\end{equation*}
and
\begin{equation*}
\lambda_{\Sigma} := \kappa_1 - 2 - c_l^{\circ}(r+1), \qquad
\lambda'_{\Sigma} := \lambda_{\Sigma} - \log(e\hat{b}), \qquad
\epsilon_{\partial} := 2e^{-\lambda'_{\Sigma}}.
\end{equation*}
Under the condition on $\kappa_1$ fixed in the order of choice of the constants
(Definition~\ref{5.1:def-admitted-inputs}), $\lambda'_{\Sigma} \ge 75c_l - 1$ for
$r \le a + 9$. We assume the condition
\begin{equation}\label{5.4:eq-zones-cells-i}
\text{(}\epsilon\text{)}\ \ 2\,\epsilon_{\partial}\,\hat{b}\,\mathsf{c}_d \le 1 .
\end{equation}
The small factor is
\begin{equation*}
\theta'(g_k) = C_{\theta'}\,\mathsf{T}_k^{\,p_1 - q},
\end{equation*}
where $p_1$ denotes the exponent entering the small factor together with the auxiliary
parameter $q$, and we put $\mathsf{w}_L := \theta'(g_k)^{-1}$. The constants
$C^{0}_{\square}$, $\mathsf{b}^{\natural}$, $\mathsf{c}_{\theta'}$,
$\mathsf{C}^{4,d}$, $\mathsf{C}^{4,c}$, $\mathsf{M}^{\mathrm{sh}}$, $C_{\uparrow}$,
$C_{\mathsf{b}}$, $p_{\mathsf{b}}$ have values fixed in the statements that introduce
them. The window weights $w$ satisfy
$w \le (\mathsf{L}_0^{w})^{2} L\, C_{\uparrow}\, \check{R}_{k_a}$ on $W_{\mathfrak{a}}$ and
$w = 1$ off the windows, and $\bar{\mathsf{w}} = 3w$; here $\mathsf{L}_0^{w}$ is the constant
of the window weights and is unrelated to the constant $L_0 = L_0(N)$. The cell weights are
$\mathsf{f}(\Delta')$, with bounding constant $c_{\chi}$, and for a unit $v$
(Definition~\ref{5.2:def-sites-units}), with its finite family of cells read by $v$, we put
\begin{equation*}
\mathsf{f}_v := c_{\chi} \max\{\mathsf{f}(\Delta') : \Delta' \text{ in that family}\}.
\end{equation*}
\end{definition}

\begin{proof}
We prove the assertions above that are derived from the axioms and conventions, in the
order in which they appear.

\emph{(S2), second assertion.} Let $q, q_1, \dots, q_n \in \pi_m$ with
$q \subset q_1 \cup \dots \cup q_n$, and suppose $q \neq q_i$ for every $i$. By the first
assertion of (S2), $\operatorname{int} q \cap \operatorname{int} q_i = \emptyset$ for every
$i$, so $\operatorname{int} q \subset \partial q_1 \cup \dots \cup \partial q_n$. A finite
union of boundaries of cubes has empty interior, while $\operatorname{int} q$ is a non-empty
open set; this is a contradiction, so $q$ is one of the $q_i$.

\emph{Local finiteness; $\Omega_m^{\mathrm{cur}}$ is closed.} By (S3) every cell is a cube of
one of the finitely many levels $m_0, \dots, k$, hence has side between $L^{m_0-k}$ and $1$ in
level-$k$ units, and distinct cells have disjoint interiors. The cells meeting a bounded set
$B$ therefore lie in the $1$-neighbourhood of $B$, a bounded set, and have disjoint interiors
of volume at least $L^{d(m_0-k)}$; hence they are finitely many. Thus the family of cells is
locally finite, and
$\Omega_m^{\mathrm{cur}}$, a locally finite union of closed sets, is closed.

\emph{Inclusions.} Let $\Delta$ be a cell of level $m$. By (S3),
$\Delta \subset \operatorname{cl}(\text{zone } m)$. Since
$\text{zone } m = \Omega_m \setminus \Omega_{m+1} \subset \Omega_m$ and $\Omega_m$ is closed,
$\operatorname{cl}(\text{zone } m) \subset \Omega_m$. Since also
$\text{zone } m \subset \Omega_{m+1}^c$, we get
$\operatorname{cl}(\text{zone } m) \subset \operatorname{cl}(\Omega_{m+1}^c) = U_{m+1}$. If
$\Delta$ has level $m' \ge m$, then $\Delta \subset \Omega_{m'} \subset \Omega_m$ by what was
just shown and by (A1) applied $m' - m$ times; taking the union over all cells of level
$\ge m$ gives $\Omega_m^{\mathrm{cur}} \subset \Omega_m$. Finally, for $i < k$, if
$A \subset U_i$ and $B \subset \Omega_{i+1}$, then
\begin{equation*}
\operatorname{dist}_{i+1}(A, B) \ge \operatorname{dist}_{i+1}(U_i, \Omega_{i+1}) \ge 1
\end{equation*}
by (A3).

\emph{Wall-neighbours.} A cell $\Delta$ of level $m$ has $2d$ faces. By (A2), a cell sharing
a wall with $\Delta$ has level $m-1$, $m$ or $m+1$. Across a given face there is either one
cell of equal or coarser level, or finer cells, of level $m-1$; in the latter case there are
at most $L^{d-1}$ of them, since a face of $\Delta$ has side $L$ in level-$(m-1)$ units and
so contains $L^{d-1}$ faces of $\pi_{m-1}$-cubes. Summing over the $2d$ faces, $\Delta$ has at
most $2dL^{d-1}$ wall-neighbours.

\emph{Touching cells.} Let $\Delta$ be a cell of level $m$. By (A2), a cell touching $\Delta$
has level $m-1$, $m$ or $m+1$. Let $\Delta'$ be a cell of level $m' \ge m$ touching $\Delta$.
By (S1), $\Delta'$ is the union of the cubes of $\pi_m$ it contains, so a common point of
$\Delta$ and $\Delta'$ lies in a cube $q \in \pi_m$ with $q \subset \Delta'$, and $q$ touches
$\Delta$. The cubes of $\pi_m$ touching $\Delta$ are the $3^d$ cubes of the cube of side $3$,
in level-$m$ units, centred at $\Delta$. Distinct cells $\Delta'$ give distinct cubes $q$,
since $q$ has non-empty interior and distinct cells have disjoint interiors. Hence at most
$3^d$ cells of level $\ge m$ touch $\Delta$. A cell of level $m-1$ touching $\Delta$ is a cube
of $\pi_{m-1}$. By (S1), $\Delta$ is the union of the $L^d$ cubes of $\pi_{m-1}$ it contains,
so a cube of $\pi_{m-1}$ touching $\Delta$ is one of the $(L+2)^d$ cubes of $\pi_{m-1}$
contained in the cube of side $L+2$, in level-$(m-1)$ units, concentric with $\Delta$.
Distinct cells being distinct cubes, at most $(L+2)^d$ cells of level $m-1$ touch $\Delta$.
Altogether at most $3^d + (L+2)^d = \hat{b}_t$ cells touch $\Delta$.

\emph{Cubes not contained in $Z$.} Let $P$ be a $\pi_k$-cube with $P \not\subset Z$, and
suppose $\operatorname{int} P \cap Z \neq \emptyset$. As $Z$ is a union of $\pi_k$-cubes,
$\operatorname{int} P$ meets some $\pi_k$-cube $Q \subset Z$; since $\operatorname{int} P$ is
open and $Q$ is the closure of its interior, $\operatorname{int} P \cap \operatorname{int} Q
\neq \emptyset$, and (S1) with $m = m' = k$ gives $P \subset Q \subset Z$, a contradiction.
Hence $\operatorname{int} P \cap Z = \emptyset$, that is, $\operatorname{int} P \subset Z^c$,
and $P = \operatorname{cl}(\operatorname{int} P) \subset \operatorname{cl}(Z^c)$.

\emph{Source cells.} Every cell has level at most $k$, hence side at most $1$ in level-$k$
units. A source cell shares a wall with a hub, so it touches $\Lambda^{\sim}$, and each of
its points lies within $\operatorname{dist}_k$-distance $1$ of $\Lambda^{\sim}$. By (Z),
$\operatorname{dist}_k(\Lambda^{\sim}, Z^c) \ge 9\check{R}_k$, and the triangle inequality gives
$\operatorname{dist}_k(\text{source cell}, Z^c) \ge 9\check{R}_k - 1$. This is positive by
($\check{R}$2), so the source cell lies in $Z$.

\emph{($\rho$0), ($\rho$1) and ($\rho$2).} ($\rho$0) is the definition of $\rho(\Delta)$ as a
least cardinality. For ($\rho$2), let $\mathcal{F}$ be a connected finite family in
$\pi^{\mathfrak{s}}$ containing $\Delta$ and a source cell. As $\mathcal{F}$ is connected, it
contains a path of distinct members from $\Delta$ to a source cell, consecutive members being
adjacent; we stop it at its first source cell and write it $\Delta = \Delta_1, \Delta_2, \dots,
\Delta_n$, so that $\Delta_n$ is a source cell and $\Delta_1, \dots, \Delta_{n-1}$ are not. An
adjacency through a hub joins two cells each sharing a wall with a hub, that is, two source
cells; since of any two consecutive members of the path at least one is not a source cell,
consecutive members share a wall. Each
member has level at most $k$, hence $\operatorname{diam}_k$ at most $1$. Since
$\Delta_{l+1}$ touches $\Delta_l$, we have
\begin{equation*}
\operatorname{dist}_k(\Delta, \Delta_{l+1}) \le \operatorname{dist}_k(\Delta, \Delta_l)
+ \operatorname{diam}_k \Delta_l \le \operatorname{dist}_k(\Delta, \Delta_l) + 1
\qquad (2 \le l < n),
\end{equation*}
and $\operatorname{dist}_k(\Delta, \Delta_2) = 0$; hence
$\operatorname{dist}_k(\Delta, \Delta_n) \le \max\{n-2, 0\} \le n$. Therefore
\begin{equation*}
\operatorname{dist}_k\bigl(\Delta, \textstyle\bigcup\{\text{source cells}\}\bigr)
\le \operatorname{dist}_k(\Delta, \Delta_n) \le n,
\end{equation*}
so $\mathcal{F}$ has at least
$\operatorname{dist}_k(\Delta, \bigcup\{\text{source cells}\})$ members, and ($\rho$0)
gives ($\rho$2). For ($\rho$1), consider the same chain. Consecutive members share a wall,
hence have a common point, so by (A2) their levels differ by at most $1$. The chain starts at
$\Delta$, of level $\operatorname{lev}\Delta$, and ends at a source cell, which has level $k$
by the standing assumption on source cells; it therefore has at least
$k - \operatorname{lev}\Delta + 1$
members, all distinct and in $\mathcal{F}$. Hence every such family $\mathcal{F}$ has at least
$k - \operatorname{lev}\Delta + 1$ members, and ($\rho$0) gives ($\rho$1).

\emph{Consequences of ($\Gamma$5).} Let $\Delta$ be a cell of level $m < k$. Then
$m + 1 \le k$, so $\Omega_k \subset \Omega_{m+1}$ by (A1), and
\begin{equation*}
\Delta \subset \operatorname{cl}(\text{zone } m) \subset \operatorname{cl}(\Omega_{m+1}^c)
\subset \operatorname{cl}(\Omega_k^c).
\end{equation*}
Since the distance to a closure equals the distance to the set,
\begin{equation*}
\operatorname{dist}_k(\Delta, \partial Z) \ge \operatorname{dist}_k(\Omega_k^c, \partial Z)
\ge 2\check{R}_k
\end{equation*}
by ($\Gamma$5). A raised cell has $\operatorname{dist}_k(\cdot, \partial Z) \ge 2\check{R}_k$
because, by ($\Gamma$5), no cell within $2\check{R}_k$ of $\partial Z$ is raised.

Now let $Q$ be a $\pi_k$-cube with $\operatorname{dist}_k(Q, \partial Z) < 2\check{R}_k$. Every
point of $\operatorname{int} Q$ lies in a cell (the cells cover the carrier, by (S3)), and a
cell $\Delta$ containing a point of $\operatorname{int} Q$ satisfies
$\operatorname{int}\Delta \cap \operatorname{int} Q \neq \emptyset$, because $\Delta$ is the
closure of its interior and $\operatorname{int} Q$ is open. If some such $\Delta$ has level
$k$, then (S1) with $m = m' = k$, applied in both directions, gives $\Delta = Q$. Otherwise
every such $\Delta$ has level $< k$ and lies in $\operatorname{cl}(\Omega_k^c)$ by the
preceding paragraph; then $\operatorname{int} Q \subset \operatorname{cl}(\Omega_k^c)$, hence
$Q \subset \operatorname{cl}(\Omega_k^c)$, and
\begin{equation*}
\operatorname{dist}_k(Q, \partial Z) \ge \operatorname{dist}_k(\Omega_k^c, \partial Z)
\ge 2\check{R}_k,
\end{equation*}
a contradiction. Hence $Q$ is a current cell of level $k$, and it is un-raised by ($\Gamma$5),
since it lies within $2\check{R}_k$ of $\partial Z$.
\end{proof}

\begin{definition}[Pointed units, short and long units, and the met-cell sums]
\label{5.4:def-pointed-units}
Throughout, $k$ is the level under consideration, $Z$ is the large-field region at level $k$, and
cells, current cells, source cells, the level $\operatorname{lev}\Delta$ of a cell $\Delta$, the level-$n$
distances $\operatorname{dist}_n$, the grids $\pi_j$, $\pi^{\mathfrak{s}}$, the coarse hulls
$[\,\cdot\,]_k$, the regions $\Lambda_j$ and the zones $\Omega_m$ are those fixed in
Notation~\ref{5.4:not-zones-cells}.
The units considered below belong to the two classes of units lettered (c) and (d) in the
representation and bound of the large-field terms; the letters (c) and (d) are names of these
classes and are distinct from the part numbers (i)--(v) of this definition.

\smallskip
\noindent\emph{(i) Pointed units.} A \emph{pointed unit} $v$ of class (c) or (d) carries:
\begin{itemize}
\item a level $j_v\le k$;
\item a region $X_v\in\mathbf{D}_{j_v}$, which satisfies $X_v\subset\Lambda_{j_v}$ in the label
  of the step at which $v$ is created (this inclusion is a property of the terms of the
  representation and bound of the large-field terms from which the units arise, and is taken
  from that statement);
\item the coarse hull $K_v:=[X_v]_k\in\mathbf{D}_k$;
\item a distinguished point $y_v\in\bigcup X_v$;
\item the size $\mathfrak{N}^{\flat}(v):=2\mathfrak{n}_v$, where $\mathfrak{n}_v$ is the size
  of $v$ fixed in Definition~\ref{5.2:def-smoothing-depth};
\item the tree length $d_v:=d_{j_v}(X_v)$ and, for $r'\ge 0$, the weight
  $w_{r'}(v):=\mathfrak{N}^{\flat}(v)\,e^{r'd_k(K_v)}$;
\item bracket indices, that is, the values taken by the index $s$ of the brackets of the
  representation and bound of the large-field terms: the single index $s=4$ if $v$ is of
  class (c), and the indices $s=4,3,2,1$ if $v$ is of class (d).
\end{itemize}
The pointed units are of two kinds, \emph{singles} and \emph{tails}; a tail arises together with
a \emph{head} from one unit, its \emph{parent}, and the union of head and tail is the
\emph{carrier} of the parent.
Let $q_v$ be the first cube of $X_v$, in a fixed enumeration of its cubes, that contains $y_v$,
and let $\Delta'_v$ be the current cell containing $q_v$; that $\Delta'_v$ is well defined is
the third clause of the lemma on current cells below, and since $y_v\in q_v\subset\Delta'_v$,
one has $y_v\in\Delta'_v$. For each $v$ let $\Delta_0(v)$ be the current cell containing $y_v$
with respect to which the short/long split of $v$ is made in the representation and bound of
the large-field terms; when that cell is the one holding $q_v$, $\Delta_0(v)=\Delta'_v$. Put
$n_0(v):=\operatorname{lev}\Delta_0(v)$ and, for a current cell $\Delta_0$, define the open set
\begin{equation}\label{5.4:eq-pointed-units-1}
\mathcal{N}(\Delta_0):=\{x:\ \operatorname{dist}_{\operatorname{lev}\Delta_0}(x,\Delta_0)<1\}.
\end{equation}

\smallskip
\noindent\emph{(ii) Short and long units.} For a current cell $\Delta_0$ of level $n_0$, the
\emph{star} of $\Delta_0$ is the family of current cells
\begin{equation*}
\mathrm{St}(\Delta_0):=\{\Delta\ \text{a current cell}:\ \operatorname{dist}_{n_0}(\Delta,\Delta_0)<1\}.
\end{equation*}
The star lemma of this section shows that its union contains
\begin{equation*}
\{x:\ \operatorname{dist}_{n_0}(x,\Delta_0)\le 1\}\cap\Omega_{n_0-1}.
\end{equation*}
A pointed unit $v$ is \emph{short} if and only if
\begin{equation}\label{5.4:eq-pointed-units-a}
\textstyle\bigcup X_v\subset\bigcup\mathrm{St}\bigl(\Delta_0(v)\bigr).
\end{equation}
The \emph{long singles} and the \emph{long tails} are the pointed units that are not short. The
lemma on long units below shows that
$\bigcup X_v\not\subset\mathcal{N}(\Delta_0(v))$ for every long single or long tail $v$.

If a tail $v$ is given with its head detached, write $X^{\mathrm{hd}}_v$ for the head and
$X^{\mathrm{tl}}_v$ for the tail; then, in \eqref{5.4:eq-pointed-units-a}, in the condition
``$X_v$ meets $\Delta_2$'' of part (iii), and wherever $X_v$ enters the bounds of this section
through $d_v=d_{j_v}(X_v)$ or through $K_v=[X_v]_k$ (in particular in the weights $w_{r'}(v)$
and in the factors $e^{\beta d_v}$ and $e^{(2+c)d_v}$), the region $X_v$ is taken to be the
carrier of the parent,
\begin{equation}\label{5.4:eq-pointed-units-b}
X_v:=X^{\mathrm{hd}}_v\cup X^{\mathrm{tl}}_v .
\end{equation}
A tree spanning $X^{\mathrm{hd}}_v\cup X^{\mathrm{tl}}_v$ spans $X^{\mathrm{tl}}_v$, so the tree
lengths computed from \eqref{5.4:eq-pointed-units-b}, and with them the weights, are not smaller
than those computed from the tail alone; and if $X^{\mathrm{tl}}_v$ meets $\Delta_2$, then so does
$X_v$. Hence every bound stated below with $X_v$ as in \eqref{5.4:eq-pointed-units-b} majorises,
term by term, the corresponding sum indexed and weighted by the tail alone.

\smallskip
\noindent\emph{(iii) The long families and the met-cell sums.} Let $\mathfrak{L}^{\mathrm{lg}}_{c}$ be
the set of long singles and long tails $v$ of class (c) with $X_v\cap Z\neq\emptyset$ (intersection
of closed sets) and $K_v\not\subset Z$; let $\mathfrak{L}^{\mathrm{lg}}_{d}$ be the set of those of
class (d) with $X_v\cap Z=\emptyset$; and put
$\mathfrak{L}^{\mathrm{lg}}:=\mathfrak{L}^{\mathrm{lg}}_{c}\sqcup\mathfrak{L}^{\mathrm{lg}}_{d}$
(the subscripts $c$ and $d$ here name the classes and are not parameters).
For a current cell $\Delta_2$, a family $\mathfrak{L}\subset\mathfrak{L}^{\mathrm{lg}}$, $r'\ge 0$
and $\beta\ge 0$, and for a current cell $\Delta'$ and a parameter $c$, define
\begin{equation}\label{5.4:eq-pointed-units-c}
\begin{aligned}
S^{(\beta)}_{r'}(\Delta_2;\mathfrak{L})&:=\sum_{v\in\mathfrak{L}:\ X_v\text{ meets }\Delta_2}
  w_{r'}(v)\,e^{\beta d_v},\qquad
  S_{r'}(\Delta_2;\mathfrak{L}):=S^{(2)}_{r'}(\Delta_2;\mathfrak{L}),\\
\mathrm{Col}^{\mathrm{file}}_{r',c}(\Delta')&:=\sum_{v\in\mathfrak{L}^{\mathrm{lg}}:\
  \Delta'_v=\Delta'} w_{r'}(v)\,e^{(2+c)d_v}\\
  &\;\le\;\mathrm{Col}_{r',c}(\Delta'):=\sum_{v\in\mathfrak{L}^{\mathrm{lg}}:\ y_v\in\Delta'}
  w_{r'}(v)\,e^{(2+c)d_v}.
\end{aligned}
\end{equation}
The inequality in \eqref{5.4:eq-pointed-units-c} holds because $y_v\in\Delta'_v$ by part (i), so
every $v$ with $\Delta'_v=\Delta'$ has $y_v\in\Delta'$, and every term of both sums is non-negative.

\smallskip
\noindent\emph{(iv) Pieces of the units of class (d).} For $v\in\mathfrak{L}^{\mathrm{lg}}_{d}$ let
$\mathrm{St}_v$, the \emph{star of the unit} $v$, be the set of current cells meeting $X_v$ (a
notion distinct from the star $\mathrm{St}(\Delta_0)$ of a current cell in part (ii)),
$S_v:=\bigcup\mathrm{St}_v$,
$T_v:=\pi^{\mathfrak{s}}\setminus\mathrm{St}_v$, and let $\partial\mathrm{St}_v$ be the set of members
of $T_v$ sharing a wall with a member of $\mathrm{St}_v$. With $\delta_P$, the distance
$\operatorname{dist}$ between sets, the set $\tilde{\square}_p$ attached to the block of a piece $p$,
$\mathsf{w}_L$,
$\kappa_1$, the weight $\mathsf{f}_v$, the function $f_v$, the blocks $\square$ with their functions
$\hat{\chi}_{\square}$ and the interpolated chain $\Phi_s^{\,\cdot}(\mathsf{s})$ as in the
representation and bound of the large-field terms, a piece is a quadruple
$p=(v;s;\square;\mathsf{y})$ with $s$ a bracket index of $v$ and $\mathsf{y}\subset T_v$ finite, and
one puts $m(p):=|\mathsf{y}|$,
\begin{equation*}
\eta_p:=e^{-\delta_P \operatorname{dist}(\tilde{\square}_p,S_v)},\qquad
\rho_A(v,p):=\mathsf{w}_L/(\eta_p\mathsf{f}_v)^{1/2},
\end{equation*}
\begin{equation}\label{5.4:eq-pointed-units-2}
v_p:=\Bigl[\prod_{\Delta\in\mathsf{y}}\int_0^1 d\mathsf{s}(\Delta)\,\partial_{\mathsf{s}(\Delta)}\Bigr]
\int_0^1 dt\;\partial_{\mathsf{t}}\big|_{\mathsf{t}=0}\,
f_v\bigl(\Phi_s^{t+\mathsf{t}\hat{\chi}_{\square}}(\mathsf{s})\bigr)
\Big|_{\mathsf{s}(\Delta)=0\ \text{for}\ \Delta\notin\mathsf{y}},
\end{equation}
\begin{equation}\label{5.4:eq-pointed-units-3}
\mathfrak{N}^{\mathrm{ex}}_{A}(p):=2\mathsf{w}_L^{-1}(\eta_p\mathsf{f}_v)^{1/2}\,
\mathfrak{N}^{\flat}(v)\,e^{-(\kappa_1-1)m(p)},
\end{equation}
and
\begin{equation*}
\mathrm{lab}^{\circ}_1(p):=\mathrm{lab}(K_v,\mathsf{y}):=K_v\cup[\mathsf{y}]_k\cup
\mathrm{Hub}(\mathsf{y}).
\end{equation*}
For a finite wall-connected $\mathrm{St}\subset\pi^{\mathfrak{s}}$
containing no source cell,
\begin{equation}\label{5.4:eq-pointed-units-d}
\mathfrak{Y}(\mathrm{St}):=\bigl\{\mathsf{y}\subset\pi^{\mathfrak{s}}\setminus\mathrm{St}\ \text{finite}:\
\mathrm{St}\cup\mathsf{y}\ \text{connected},\ \mathsf{y}\ \text{contains a source cell}\bigr\}.
\end{equation}

\smallskip
\noindent\emph{(v) Units of class (c).} For $v\in\mathfrak{L}^{\mathrm{lg}}_{c}$ one uses the grain
$\pi^{\mathfrak{s}}_j$, the sets $\mathrm{St}^j_v$, $\partial^{+}\mathrm{St}^j_v$ and
$\mathfrak{H}_v$, the bundles $\mathfrak{b}=(v;4;\tilde{\mathsf{y}})$ with norm
$\|\mathfrak{b}\|:=2\mathfrak{N}^{\flat}(v)e^{-(\kappa_1-1)|\tilde{\mathsf{y}}|}$ and label
$\mathrm{lab}(\mathfrak{b})$, and the admissible families $\mathfrak{Y}_v$, all with the values
fixed in Definition~\ref{5.2:def-boundary-bundles}, $\mathfrak{Y}_v$ being cut out in
Definition~\ref{5.2:def-boundary-bundles} by its two admissibility conditions.
\end{definition}

\begin{definition}[Age functions and the two base bounds for long units]\label{5.4:def-age-base-bounds}
Throughout, $k$ is the current level; cells, their levels $\mathrm{lev}\,\Delta'$, the regions $Z$
and $\Lambda$, the hubs $\Lambda^{\sim}$, the cores and the windows are as in
Notation~\ref{5.4:not-zones-cells}, and long units are as in
Definition~\ref{5.4:def-pointed-units}. We write $\mathfrak{L}^{\mathrm{lg}}$ for the set of long units,
that is, of all long singles and long tails. For $v\in\mathfrak{L}^{\mathrm{lg}}$ we write $j_v$ for its
level, $X_v$ for its set of cubes, $w_{r'}(v)$ for its weight at rate $r'$ and $f_v$ for the
function it carries, with analyticity domain $\mathfrak{D}_v$, and we set
$d_v:=d_{j_v}(X_v)$. We write $y_v$ for the datum of $v$ whose cell's star separates long units
from short ones, $\Delta_0(v)$ for the cell attached to $v$, and $K_v$ for the set at
level $k$ attached to $v$, all as in Definition~\ref{5.4:def-pointed-units}.

(a) \emph{Age functions.} An \emph{age function} is a non-decreasing map
$\mathsf{b}^{\mathrm{lg}}\colon\mathbb{Z}_{\ge 0}\to[1,\infty)$ for which there are real numbers
$C_{\mathsf{b}}\ge 1$ and $p_{\mathsf{b}}\ge 0$ such that, for all $x,u\in\mathbb{Z}_{\ge 0}$,
\begin{equation}\label{5.4:eq-age-base-bounds-a}
\mathsf{b}^{\mathrm{lg}}(x+u)\le C_{\mathsf{b}}\,(1+u)^{p_{\mathsf{b}}}\,\mathsf{b}^{\mathrm{lg}}(x).
\end{equation}
Three age functions used below are instances of this family:
\begin{itemize}
\item the $k$-free function $\mathsf{b}^{\mathrm{lg}}\equiv\mathsf{b}^{\mathrm{lg}}(0)$, with
$(C_{\mathsf{b}},p_{\mathsf{b}})=(1,0)$;
\item the linear function $\mathsf{b}^{\mathrm{lg}}(x)=(x+1)\,\mathsf{b}^{\mathrm{lg}}(0)$, with
$(C_{\mathsf{b}},p_{\mathsf{b}})=(1,1)$;
\item the polynomial function $\mathsf{b}(x)=\mathsf{b}_k(x)+(x+1)b_{\delta}$, where the function
$\mathsf{b}_k(\cdot)$, the number $b_{\delta}$ and the exponent $r_{\mathrm{pr}}$ are those fixed in the
representation and bound of the large-field terms; for this function
\eqref{5.4:eq-age-base-bounds-a} is the growth bound stated there, with
$C_{\mathsf{b}}=(1+3\lambda)^{r_{\mathrm{pr}}}$ and $p_{\mathsf{b}}=1+r_{\mathrm{pr}}$.
\end{itemize}
Moreover, the product of two age functions with constants $(C_1,p_1)$ and $(C_2,p_2)$ is an age
function with constants $(C_1C_2,\,p_1+p_2)$.

(b) \emph{The long base bound per own-scale cube.} We say that this bound holds if there are a
number $c_{\mathrm{sp}}>0$ and an age function $\mathsf{b}^{\mathrm{lg}}$ such that for every
$j\le k$, every $\square'\in\pi_j$ and every $r'\le\frac{8}{5}\kappa$,
\begin{equation}\label{5.4:eq-age-base-bounds-b}
\sum_{\substack{v\in\mathfrak{L}^{\mathrm{lg}}:\ j_v=j,\\ \square'\in X_v}}
w_{r'}(v)\,e^{(2+c_{\mathrm{sp}})d_v}\le\mathsf{b}^{\mathrm{lg}}(k-j).
\end{equation}
The sum runs over all long singles and long tails of level $j$ one of whose cubes is $\square'$.
No condition is imposed on $y_v$ or on $\Delta_0(v)$, nor on the position of $\square'$ or
of $K_v$ relative to $Z$, $\Lambda$, the cores, the hubs or the windows. We say that the bound
holds \emph{off $Z$} if \eqref{5.4:eq-age-base-bounds-b} is required only for cubes
$\square'$ off $Z$; where only this weaker form is assumed, the conclusions drawn below from
\eqref{5.4:eq-age-base-bounds-b} are asserted only at cells $\Delta_2\not\subset Z$, and this
restricted form suffices for every later use of \eqref{5.4:eq-age-base-bounds-b} below. The number
$c_{\mathrm{sp}}$ is subject only to the one-sided requirement $2+c_{\mathrm{sp}}\le\frac12\kappa$,
the rate $\frac12\kappa$ being the decay rate available to the long units at
$r'\le\frac{8}{5}\kappa$; thus the choice $c_{\mathrm{sp}}=1$ is admissible once
$\frac12\kappa\ge 3$, in particular under $\kappa\ge 180$.

(c) \emph{The filed long column bound.} Let $\mathsf{b}^{\mathrm{lg}}$ be an age function. We say
that this bound holds for $\mathsf{b}^{\mathrm{lg}}$ if there are a number $c\ge 2$ and a number
$\mathsf{K}^{\flat}_{\mathrm{lg}}$ such that for every current cell $\Delta'$ and every
$r'\le\frac{8}{5}\kappa$,
\begin{equation}\label{5.4:eq-age-base-bounds-c}
\mathrm{Col}^{\mathrm{file}}_{r',c}(\Delta')\le
\mathsf{b}^{\mathrm{lg}}(k-\mathrm{lev}\,\Delta')\,\mathsf{K}^{\flat}_{\mathrm{lg}},
\end{equation}
where $\mathrm{Col}^{\mathrm{file}}_{r',c}(\Delta')$ is the filed column of the long units at the
cell $\Delta'$, with parameters $r'$ and $c$, in which long units of all levels $j_v$ are summed.
The cell $\Delta'$ is arbitrary: off $Z$, in $Z\setminus\Lambda^{\sim}$, in a hub,
in a core or in a window alike. The number $\mathsf{K}^{\flat}_{\mathrm{lg}}$ depends on none of
$k$, $\mathrm{lev}\,\Delta'$, the ages, or the windows. The bound
\eqref{5.4:eq-age-base-bounds-c} is implied by the same bound for the closed-membership column
$\mathrm{Col}_{r',c}(\Delta')$, and also by the same bound for the met-cell sum
$S^{(2+c)}(\Delta';\mathfrak{L}^{\mathrm{lg}})$ of Definition~\ref{5.4:def-pointed-units}, the
latter by the lemma comparing the met-indexed and the filed forms of the column. The stronger
requirement $c\ge 3$ is imposed only in those estimates below which carry window weights.

(d) \emph{The cell-weight bound.} This is the conjunction of the two lines of the following
display: the first line, required for every cell $\Delta'$, bounds the cell weights; the second
line, required for every long unit $v$, bounds $\mathsf{f}_v$, with $\mathsf{c}_{\mathsf{f}}$
defined there:
\begin{equation}\label{5.4:eq-age-base-bounds-d}
\begin{gathered}
\mathsf{f}(\Delta')^{1/2}\,\bar{\mathsf{w}}(\Delta')^{3}\le\mathsf{C}_{fw},\qquad
\bar{\mathsf{w}}(\Delta')\ge 3,\\
\mathsf{f}_v^{1/2}\le\mathsf{c}_{\mathsf{f}}
:=c_{\chi}^{1/2}\max\{1,\ \mathsf{C}_{fw}/27\}.
\end{gathered}
\end{equation}
The second line is assumed together with the first and is not derived from the first line in
what follows; the proof below derives from the first line alone the bound on
$\mathsf{f}(\Delta')^{1/2}$ from which the constant $\mathsf{c}_{\mathsf{f}}$ is built.

(e) \emph{The analyticity hypothesis for long units.} Let
$E:=\{\sigma\in\mathbb{C}:\ |\sigma|<e^{\kappa_1}\}$. For $v\in\mathfrak{L}^{\mathrm{lg}}$, a piece
index $p$, an index $s$ of the interpolated chain, a block $\square$ and a finite family
$\mathsf{y}\subset T_v$, and for $|\mathsf{t}|<\rho_A(v,p)$, the analyticity radius of $v$ at the
piece $p$, for $t\in[0,1]$ and for $\mathsf{s}\in\bar{E}^{\mathsf{y}}$, extended by $0$ off
$\mathsf{y}$, consider
\begin{equation}\label{5.4:eq-age-base-bounds-e}
\begin{gathered}
\mathsf{g}_v(\mathsf{t},t,\mathsf{s}):=
f_v\bigl(\Phi_s^{t+\mathsf{t}\hat{\chi}_{\square}}(\mathsf{s})\bigr),\\
\bigl|\mathsf{g}_v(\mathsf{t},t,\mathsf{s})-f_v(1)\bigr|\le 2\mathfrak{n}_v=\mathfrak{N}^{\flat}(v).
\end{gathered}
\end{equation}
This hypothesis consists of the following two clauses, required for every
$v\in\mathfrak{L}^{\mathrm{lg}}$, every piece index $p$, every $s$, every block $\square$ and every
finite family $\mathsf{y}\subset T_v$.
\begin{itemize}
\item[(i)] For all $(\mathsf{t},t,\mathsf{s})$ as above, the dilated chain
$\Phi_s^{t+\mathsf{t}\hat{\chi}_{\square}}(\mathsf{s})$ lies in the analyticity domain
$\mathfrak{D}_v$ of $f_v$, and the function $\mathsf{g}_v$ of the first line of
\eqref{5.4:eq-age-base-bounds-e} is continuous, holomorphic in $\mathsf{t}$, and holomorphic in each
$\mathsf{s}(\Delta)\in E$, $\Delta\in\mathsf{y}$, with the other arguments frozen.
\item[(ii)] For all $(\mathsf{t},t,\mathsf{s})$ as above, the second line of
\eqref{5.4:eq-age-base-bounds-e} holds.
\end{itemize}
\end{definition}

\begin{proof}
We verify the assertions made in (a), and record the consequence of the first line of
\eqref{5.4:eq-age-base-bounds-d} from which the constant $\mathsf{c}_{\mathsf{f}}$ is built.

\emph{The $k$-free function.} A constant map with value $\mathsf{b}^{\mathrm{lg}}(0)\ge 1$ is
non-decreasing with values in $[1,\infty)$, and both sides of \eqref{5.4:eq-age-base-bounds-a}
equal $\mathsf{b}^{\mathrm{lg}}(0)$ when $C_{\mathsf{b}}=1$ and $p_{\mathsf{b}}=0$.

\emph{The linear function.} Since $\mathsf{b}^{\mathrm{lg}}(0)\ge 1$, the map
$x\mapsto(x+1)\mathsf{b}^{\mathrm{lg}}(0)$ is non-decreasing with values in $[1,\infty)$. For
$x,u\ge 0$ we have $ux\ge 0$, hence
\begin{equation*}
x+u+1\le x+1+ux+u=(1+u)(x+1).
\end{equation*}
Multiplying by $\mathsf{b}^{\mathrm{lg}}(0)$ gives
$\mathsf{b}^{\mathrm{lg}}(x+u)\le(1+u)\,\mathsf{b}^{\mathrm{lg}}(x)$, which is
\eqref{5.4:eq-age-base-bounds-a} with $(C_{\mathsf{b}},p_{\mathsf{b}})=(1,1)$.

\emph{The polynomial function.} Here \eqref{5.4:eq-age-base-bounds-a}, with the constants
$C_{\mathsf{b}}=(1+3\lambda)^{r_{\mathrm{pr}}}$ and $p_{\mathsf{b}}=1+r_{\mathrm{pr}}$, is the growth
bound stated in the representation and bound of the large-field terms, and is taken from there.

\emph{Products.} Let $\mathsf{b}'$, $\mathsf{b}''$ be age functions with constants $(C',p')$ and
$(C'',p'')$. Their product takes values in $[1,\infty)$, and it is non-decreasing as a product of
two non-decreasing functions with positive values. Applying \eqref{5.4:eq-age-base-bounds-a} to
each factor, and using that all four quantities involved are non-negative, $\mathsf{b}'$ and
$\mathsf{b}''$ taking values in $[1,\infty)$,
\begin{equation*}
\begin{split}
\mathsf{b}'(x+u)\,\mathsf{b}''(x+u)
&\le C'(1+u)^{p'}\mathsf{b}'(x)\cdot C''(1+u)^{p''}\mathsf{b}''(x)\\
&=C'C''\,(1+u)^{p'+p''}\,\mathsf{b}'(x)\,\mathsf{b}''(x).
\end{split}
\end{equation*}
Since $C'C''\ge 1$ and $p'+p''\ge 0$, the product is an age function with constants
$(C'C'',\,p'+p'')$.

\emph{The cell-weight constant.} Since $\bar{\mathsf{w}}(\Delta')\ge 3$, we have
$\bar{\mathsf{w}}(\Delta')^{3}\ge 27$, and the first line of \eqref{5.4:eq-age-base-bounds-d} gives
\begin{equation*}
\mathsf{f}(\Delta')^{1/2}\le\mathsf{C}_{fw}/27\le\max\{1,\ \mathsf{C}_{fw}/27\}
\end{equation*}
for every cell $\Delta'$; the constant $\mathsf{c}_{\mathsf{f}}$ of the second line is this bound
multiplied by $c_{\chi}^{1/2}$.
\end{proof}

\begin{definition}[The constants of the two transports]\label{5.4:not-transport-constants}
Let $C_{\mathsf{b}}$ and $p_{\mathsf{b}}$ be the growth constants of the age function, let
$\mathsf{b}^{\mathrm{lg}}$ be the age function of the base bound for long units
(Definition~\ref{5.4:def-age-base-bounds}), and put
\begin{equation*}
q_{\mathsf{b}} := d + p_{\mathsf{b}} .
\end{equation*}
We refer to the two transports as the first transport and the second transport.

\emph{The first transport.} Set $\delta := 2$ in general; set $\delta := 1$ if $\Delta_0(v) = \Delta'_v$ for
every $v$; and replace $\delta$ by $\delta + 1$ if axiom (g4) of
Definition~\ref{5.4:not-zones-cells} holds only with a margin of $-1$. With this $\delta$, and for $c > 0$,
\begin{equation}\label{5.4:eq-transport-constants-a}
\begin{aligned}
\delta &\in \{1, 2, 3\} \quad \text{as fixed above},\\
\mathfrak{m}_0 &:= 0, \qquad \mathfrak{m}_u := \bigl(L^{u-1} - \delta\bigr)_+ \quad (u \ge 1),\\
\mathfrak{c}_0 &:= 1, \qquad \mathfrak{c}_u := C_{\mathsf{b}}(1+u)^{p_{\mathsf{b}}} \quad (u \ge 1),\\
\mathfrak{A}_{\delta}(c) &:= \Bigl(\frac{2q_{\mathsf{b}}}{ec}\Bigr)^{q_{\mathsf{b}}} e^{(\delta - 3/2)c}
  \bigl(1 - e^{-c}\bigr)^{-1},\\
C^{LS} &:= 1 + C_{\mathsf{b}}\, 2^{p_{\mathsf{b}}} L^{d}
  \bigl[1 + \mathfrak{A}_{\delta}(c_{\mathrm{sp}})\bigr],\\
C^{LS}_{\sharp}(c) &:= e L^{d} \max\Bigl\{1,\ \Bigl(\frac{d+1}{ec}\Bigr)^{d+1} e^{\delta c}\Bigr\}.
\end{aligned}
\end{equation}
Here $c_{\mathrm{sp}}$ is a constant which is not fixed in this definition; $C^{LS}$ is defined
when $c_{\mathrm{sp}} > 0$, since $\mathfrak{A}_{\delta}(c)$ was defined for $c > 0$.

\emph{The second transport.}
\begin{equation}\label{5.4:eq-transport-constants-b}
\begin{aligned}
\theta &:= \theta_d := \frac{2}{\mathsf{c}_d},\\
\mathfrak{a}_d &:= e^{5\theta/4}\Bigl(\frac{2d\,\mathsf{c}_d}{e}\Bigr)^{d}
  \bigl(1 - e^{-\theta/4}\bigr)^{-1},\\
\mathfrak{e}_d &:= e^{-3\theta}\bigl(1 - e^{-3\theta/2}\bigr)^{-1},\\
C^{\mathrm{gen}}_{\mathsf{b}} &:= e^{2}\mathfrak{a}_d\Bigl[\,2 + C_{\mathsf{b}}\, 2^{p_{\mathsf{b}}} L^{d}
  + \Bigl(e^{\theta} C_{\mathsf{b}}\Bigl(\frac{q_{\mathsf{b}}\,\mathsf{c}_d}{e}\Bigr)^{q_{\mathsf{b}}} + 2^{d}\Bigr)
  \mathfrak{e}_d\Bigr].
\end{aligned}
\end{equation}

\emph{Common to both transports.} Set
\begin{equation}\label{5.4:eq-transport-constants-c}
\begin{aligned}
\mathfrak{K} &:= C^{LS} \ \text{(first transport)} \quad\text{or}\quad
  \mathfrak{K} := C^{\mathrm{gen}}_{\mathsf{b}}\,\mathsf{K}^{\flat}_{\mathrm{lg}} \ \text{(second transport)},\\
\mathsf{K}^{\mathrm{lg}} &:= \mathfrak{K}\cdot \mathsf{b}^{\mathrm{lg}}(0),\\
\mathfrak{s}_{\mathsf{b}} &:= \sup_{x \ge 0}\,(1+x)^{p_{\mathsf{b}}} e^{-x}
  \ \le\ \max\bigl\{1,\ p_{\mathsf{b}}^{\,p_{\mathsf{b}}} e^{1 - p_{\mathsf{b}}}\bigr\},\\
S_{\mathsf{b}} &:= \sum_{s \ge 0} (1+s)^{p_{\mathsf{b}}} e^{-s}.
\end{aligned}
\end{equation}
For $\varrho \ge 1$ let $B^{\mathrm{win}}(\varrho)$ be as in the first line of
\eqref{5.4:eq-transport-constants-d} (here $\varrho$ is the argument of $B^{\mathrm{win}}$ and
$W^{\mathrm{win}}$, and not the parameter $\varrho$ of the glossary). When the base bound for long
units carries a weight at window cells or cubes, that weight is either $B^{\mathrm{win}}(\varrho)$ or
a power of $3(\mathsf{L}_0^{w})^{2} L C_{\uparrow}\varrho$; in the latter case $p_{W}$ denotes the
exponent of that power, and in either case $W^{\mathrm{win}}(\varrho)$ denotes the weight carried,
as in the second line of \eqref{5.4:eq-transport-constants-d}. Call \emph{first case} the case in which
$\mathsf{b}^{\mathrm{lg}}(0)$ is of a radius-dependent form, depending on $\check{R}_k$ in the
manner of the age function $\mathsf{b}$, for which $\mathsf{b}(0) \le \mathsf{b}^{\natural}\check{R}_k$.
Then
\begin{equation}\label{5.4:eq-transport-constants-d}
\begin{aligned}
B^{\mathrm{win}}(\varrho) &:= 2C_{\mathsf{b}}\bigl(2 + \log_L(C_{\uparrow}\varrho)\bigr)^{p_{\mathsf{b}}}
  (L C_{\uparrow}\varrho)^{d},\\
W^{\mathrm{win}}(\varrho) &\in \Bigl\{ B^{\mathrm{win}}(\varrho),\
  \bigl(3(\mathsf{L}_0^{w})^{2} L C_{\uparrow}\varrho\bigr)^{p_{W}} \Bigr\},\\
C^{\mathrm{win}}_{\mathrm{lg}} &:= \sup_{\varrho \ge 1} W^{\mathrm{win}}(\varrho)\,
  e^{1 - (2\varrho - 2)/\mathsf{c}_d},\\
\sigma^{\mathrm{lg}} &:= \sup_{\varrho \ge 1} W^{\mathrm{win}}(\varrho)\, e^{-(7/2)\varrho - 1},\\
\mathsf{b}^{\mathrm{lg},\natural} &:=
\begin{cases}
\text{a fixed } \check{R}_k\text{-free constant with }
  \mathsf{b}^{\mathrm{lg}}(0) \le \mathsf{b}^{\mathrm{lg},\natural}\check{R}_k
  & \text{in the first case},\\
\mathsf{b}^{\mathrm{lg}}(0) & \text{otherwise},
\end{cases}
\end{aligned}
\end{equation}
where $C^{\mathrm{win}}_{\mathrm{lg}} := 0$ if the base bound for long units carries no weight at
window cells or cubes.
\end{definition}

\begin{proof}
The only assertion in Definition~\ref{5.4:not-transport-constants} is the upper bound for
$\mathfrak{s}_{\mathsf{b}}$ in \eqref{5.4:eq-transport-constants-c}. We assume
$p_{\mathsf{b}} \ge 0$ and read $0^{0} = 1$. Put
$f(x) := (1+x)^{p_{\mathsf{b}}} e^{-x}$ for $x \ge 0$. Then $f > 0$ and
\begin{equation*}
(\log f)'(x) = \frac{p_{\mathsf{b}}}{1+x} - 1 .
\end{equation*}
If $0 \le p_{\mathsf{b}} \le 1$, this derivative is at most $p_{\mathsf{b}} - 1 \le 0$ for
$x \ge 0$; so $f$ is non-increasing on
$[0,\infty)$ and $\mathfrak{s}_{\mathsf{b}} = f(0) = 1$. If $p_{\mathsf{b}} > 1$, the derivative is
positive on $[0, p_{\mathsf{b}} - 1)$ and negative on $(p_{\mathsf{b}} - 1, \infty)$, so the supremum
is attained at $x = p_{\mathsf{b}} - 1$, and
\begin{equation*}
\mathfrak{s}_{\mathsf{b}} = f(p_{\mathsf{b}} - 1) = p_{\mathsf{b}}^{\,p_{\mathsf{b}}} e^{1 - p_{\mathsf{b}}} .
\end{equation*}
In either case $\mathfrak{s}_{\mathsf{b}} \le \max\{1, p_{\mathsf{b}}^{\,p_{\mathsf{b}}} e^{1 - p_{\mathsf{b}}}\}$.
\end{proof}

\begin{lemma}[A unit's carrier stays inside its own zone set]\label{5.4:lem-carrier-in-zone}
Assume the following.
\begin{enumerate}
\item[(a)] Every pointed unit $v$ of Definition~\ref{5.4:def-pointed-units} satisfies
$\bigcup X_v\subset\Lambda_{j_v}$ in the label of $v$'s birth.
\item[(b)] Between the birth of a unit and the label in which it is read, the label changes only
through the events of the event census \eqref{5.4:eq-zones-cells-h}. For preparations the
relabelling is that of \cite[(1.10)--(1.12)]{B-CMP122a}, and for completions it is that of
\cite[(1.33)]{B-CMP122b}.
\end{enumerate}
Let $v$ be a pointed unit and put $j:=j_v$. Then the following hold.
\begin{enumerate}
\item[(i)] $\bigcup X_v\subset\Omega_j$, where $\Omega_j$ is taken in the current label.
\item[(ii)] $\bigcup X_v\subset\Omega_j^{\mathrm{cur}}$, where $\Omega_j^{\mathrm{cur}}$ denotes the
union of the current cells of level $\ge j$, and the interior of every cube of $X_v$
is disjoint from every cell of level $<j$.
\end{enumerate}
\end{lemma}

\begin{proof}
\emph{Part} (i). In the label of $v$'s birth, hypothesis (a) and the axiom (A0) of
\eqref{5.4:eq-zones-cells-b}, read in that label, give
\begin{equation*}
\bigcup X_v\subset\Lambda_j\subset\Omega_j .
\end{equation*}
It therefore suffices to show that the set $\Omega_j$ does not shrink along any event of the census
that occurs after the birth of $v$. By hypothesis (b), such an event belongs to one of the four
kinds of event $(\alpha)$--$(\delta)$ of \eqref{5.4:eq-zones-cells-h}. We treat them in turn.

$(\alpha)$ Consider a $\mathbf{T}$-step $m\to m+1$ that occurs after the birth of $v$. Then $m$
is at least the level of the birth step, which is at least $j$. By the census, such a step defines
only the new top pair $(\Omega_{m+1},\Lambda_{m+1})$ and leaves the pairs
$(\Omega_i,\Lambda_i)_{i\le m}$ unchanged. Since $j\le m$, the set $\Omega_j$ is untouched.

$(\beta)$ Consider a preparation of an attempt with class $Z_{\mathfrak{b}}$ and range
$]h,k_b]$. Let $\Omega''_j$ denote the relabelled set, with $k_0$, $\Omega_j^{\sim5}$ and $Z''_j$
as in \eqref{5.4:eq-zones-cells-h}; recall that $\Omega_j\subset\Omega_j^{\sim5}$. Three cases
arise.
\begin{itemize}
\item Let $j\le h$. Then $\Omega''_j=\Omega_j$, and there is nothing to prove.
\item Let $j\in\,]h,k_0]$. Then, since $\Omega_j\subset\Omega_j^{\sim5}$,
\begin{equation*}
Z''_j=(\Omega_j^{\sim5})^c\cap Z_{\mathfrak{b}}\subset\Omega_j^c .
\end{equation*}
\item Let $j\in\,]k_0,k_b]$. Then
\begin{equation*}
Z''_j\subset Z''_{k_b}=(\Omega^{\sim5}_{k_0+1})^c\cap Z_{\mathfrak{b}}
\subset\Omega^c_{k_0+1}\subset\Omega_j^c .
\end{equation*}
The first inclusion holds because the statement in \cite[(1.10)]{B-CMP122a} that the successive
sets are ``separated by one layer \dots\ and so on'' makes $Z''_j$ non-decreasing in $j$ on
$]k_0,k_b]$. The second inclusion holds because $\Omega_{k_0+1}\subset\Omega^{\sim5}_{k_0+1}$.
The last inclusion holds because $j\ge k_0+1$, so that axiom (A1) of
\eqref{5.4:eq-zones-cells-b} gives $\Omega_j\subset\Omega_{k_0+1}$.
\end{itemize}
In both of the last two cases we have $\Omega_j\subset(Z''_j)^c$. Hence
\begin{equation*}
\Omega''_j=\bigl((Z''_j)^c\cap Z_{\mathfrak{b}}\bigr)\cup\bigl(\Omega_j\cap Z_{\mathfrak{b}}^c\bigr)
\supseteq\bigl(\Omega_j\cap Z_{\mathfrak{b}}\bigr)\cup\bigl(\Omega_j\cap Z_{\mathfrak{b}}^c\bigr)
=\Omega_j .
\end{equation*}
This is the inclusion $(\Omega''_j)^c\subset\Omega_j^c$ of \cite{B-CMP122a}.

A later attempt may be prepared over a component that holds the region $X_{\mathfrak{a}'}$ of a
live arrest $\mathfrak{a}'$. By the arrest clause of the census, such an attempt has
$h\ge k_{a'}$. It is therefore again an event of kind $(\beta)$, and the argument just given
applies to it.

$(\gamma)$ Consider a completion at a level $k'\ge j$. It makes the class $Z$ one zone of the top
level $k'$. Indeed, let $\Gamma''_i$ denote the zone of level $i$ in the label immediately before
the completion and $\Gamma_i$ the zone of level $i$ in the label after it; by
\cite[(1.33)]{B-CMP122b}, $\Gamma_i=\Gamma''_i\cap Z^c$ for $i<k'$ and
$\Gamma_{k'}=\Gamma''_{k'}\cup Z$. Since $\Omega_i$ is the union of the zones of levels at
least $i$, it follows that for every $i\le k'$, in particular for $i=j$, the completion replaces
the set $\Omega''_i$ of the label before it by
\begin{equation*}
(\Omega''_i\cap Z^c)\cup Z\supseteq\Omega''_i\supseteq\Omega_i .
\end{equation*}
The first inclusion is the required monotonicity for this event. The label before the completion
is the one produced by the preparation of the completed attempt, $\Omega_i$ is the set before
that preparation, and the second inclusion is the conclusion of case $(\beta)$ for that
preparation.

$(\delta)$ By the census, arrests, exclusions and freezing do not relabel the frozen levels. They
leave $\Omega_j$ unchanged.

Between the birth of $v$ and the current label there are finitely many events, and each of them
replaces $\Omega_j$ by a set containing it. By induction along these events, the inclusion
$\bigcup X_v\subset\Omega_j$ that holds at birth persists, and it holds in the current label. This
proves (i).

\emph{Part} (ii). Let $Q$ be a cube of $X_v$, and let $\Delta$ be a cell of level $m<j$. Recall
that $U_j$ is the closure of $\Omega_j^c$. The cell $\Delta$ lies in the closure of the zone of
level $m$, and that zone lies in $\Omega_{m+1}^c$. Since $m+1\le j$, axiom (A1) gives
$\Omega_j\subset\Omega_{m+1}$, hence $\Omega_{m+1}^c\subset\Omega_j^c$. Therefore
\begin{equation*}
\Delta\subset\operatorname{cl}(\text{zone }m)\subset\operatorname{cl}(\Omega_{m+1}^c)
\subset\operatorname{cl}(\Omega_j^c)=U_j .
\end{equation*}
On the other hand, part (i) gives $Q\subset\Omega_j$, hence
$\operatorname{int}Q\subset\operatorname{int}\Omega_j$. The closure of the complement of a set is
disjoint from the interior of the set, so $U_j\cap\operatorname{int}\Omega_j=\emptyset$. It follows
that $\Delta\cap\operatorname{int}Q=\emptyset$. This is the second assertion of (ii).

The cube $Q$ is bounded. By local finiteness, only finitely many cells meet $Q$, and these cells
cover $Q$. By what has just been shown, those of level $<j$ are disjoint from
$\operatorname{int}Q$. Hence every point of $\operatorname{int}Q$ lies in a cell of level $\ge j$,
that is, $\operatorname{int}Q\subset\Omega_j^{\mathrm{cur}}$. The set $\Omega_j^{\mathrm{cur}}$
is closed, so
\begin{equation*}
Q=\operatorname{cl}\operatorname{int}Q\subset\Omega_j^{\mathrm{cur}} .
\end{equation*}
Since $Q$ was an arbitrary cube of $X_v$, we conclude $\bigcup X_v\subset\Omega_j^{\mathrm{cur}}$.
\end{proof}

\begin{remark}
The lemma asserts only an inclusion. This inclusion is stable along the events of the census
\eqref{5.4:eq-zones-cells-h}. The lemma does not assert that $X_v$ fails to be contained in
$\Omega^{\mathrm{cur}}_{j_v+1}$. Its role for long units is structural: every current cell
containing a cube of $X_v$ has level at least $j_v$, because a cell of level $<j_v$ is disjoint
from the interior of every cube of $X_v$. This is the fact on which the count of the family
$\mathrm{St}_v$ of cells meeting $X_v$ rests for long units, and it is the form in which the lemma
is used in the statement on the cells containing the cubes of a unit.
\end{remark}

\begin{lemma}[Cells against a carrier]\label{5.4:lem-cells-carrier}
Let $v$ be a pointed unit (Definition~\ref{5.4:def-pointed-units}) and put $j:=j_v$; the notation
(levels, zones, cells and the standing geometric axioms) is that fixed in
Definition~\ref{5.4:not-zones-cells}.
In this lemma a cell $\Delta$ is said to \emph{meet} $X_v$ if
$\operatorname{int}\Delta\cap\operatorname{int}Q\neq\emptyset$ for some cube $Q$ of $X_v$. The
family $\mathrm{St}_v$ of cells meeting $X_v$ and its rim $\partial\mathrm{St}_v$ are those of
Definition~\ref{5.4:def-pointed-units}.
\begin{enumerate}
\item[(a)] A cell having a common point with $\bigcup X_v$ has level at least $j-1$.
\item[(b)] A cell meeting $X_v$ has level at least $j$ and contains a cube of $X_v$.
\item[(c)] Every cube $Q$ of $X_v$ lies in exactly one cell, denoted $\Delta(Q)$; one has
$\operatorname{lev}\Delta(Q)\ge j$, and $\Delta(Q)$ meets $X_v$. In particular
$\Delta'_v:=\Delta(q_v)$ is well defined, and $y_v\in\Delta'_v$.
\item[(d)] With $\mathfrak{d}_v:=\mathsf{c}_d(d_v+1)$,
\begin{equation*}
\operatorname{diam}_{\infty,j}\Bigl(\bigcup X_v\Bigr)\le N_j(X_v)\le\mathfrak{d}_v ,
\end{equation*}
and, for every level $m\le j$,
\begin{equation*}
\operatorname{diam}_{\infty,m}\Bigl(\bigcup X_v\Bigr)\le L^{j-m}\,\mathfrak{d}_v .
\end{equation*}
\item[(e)] If a cell $\Delta_2$ of level $\ell$ meets $X_v$, then $j\le\ell$, and the cube of
$X_v$ given by (b) is one of the $L^{d(\ell-j)}$ $\pi_j$-cubes of $\Delta_2$.
\end{enumerate}
If moreover $v\in\mathfrak{L}^{\mathrm{lg}}_{d}$, then in addition:
\begin{enumerate}
\item[(f)] every member of $\mathrm{St}_v$ lies in a $\pi_k$-cube $P$ with $P\not\subset Z$, hence
lies in $\operatorname{cl}Z^{c}$, and is not contained in $Z$; the cells of
$\mathrm{St}_v\cup\partial\mathrm{St}_v$ miss $\Lambda^{\sim}$, and none of them is a source cell;
$\rho(\Delta_1)\ge 9\check{R}_k-2$ for $\Delta_1\in\partial\mathrm{St}_v$, and
$\rho(\Delta_1)\ge 9\check{R}_k-1$ for $\Delta_1\in\mathrm{St}_v$; finally, $\mathrm{St}_v$ is
wall-connected and $[\mathrm{St}_v]_k\subset K_v$;
\item[(g)] the following bounds hold:
\begin{equation*}
\#\mathrm{St}_v\le N_j(X_v)\le\mathfrak{d}_v,\qquad
\#\partial\mathrm{St}_v\le 2dL^{d-1}\,\#\mathrm{St}_v,\qquad
\epsilon_{\partial}\,\#\partial\mathrm{St}_v\le\tfrac12(d_v+1);
\end{equation*}
\item[(h)] no grouped cube of the grain $\pi^{\mathfrak{s}}_j$ used in the summation of boundary
terms at finer zones, and no cell of level $<j$, meets $X_v$; thus $\mathrm{St}^j_v=\mathrm{St}_v$.
Moreover $\mathfrak{H}_v=\emptyset$.
\end{enumerate}
The first assertion of (h) holds also for $v\in\mathfrak{L}^{\mathrm{lg}}_{c}$.
\end{lemma}

\begin{proof}
We use Lemma~\ref{5.4:lem-carrier-in-zone} (a unit's carrier stays inside its own zone set);
its clause (i) gives $\bigcup X_v\subset\Omega_j$, and its clause (ii) gives
$\bigcup X_v\subset\Omega_j^{\mathrm{cur}}$.

(a) Let $\Delta$ be a cell of level $m\le j-2$, and suppose that there is a point
$p\in\Delta\cap\bigcup X_v$. By axiom (A1) of~\eqref{5.4:eq-zones-cells-b}, applicable since
$m+1\le j-1$,
\begin{equation*}
\Delta\subset\operatorname{cl}\bigl(\Omega^{c}_{m+1}\bigr)\subset U_{j-1}.
\end{equation*}
On the other hand $\bigcup X_v\subset\Omega_j$ by Lemma~\ref{5.4:lem-carrier-in-zone}(i). Axiom
(A3) of~\eqref{5.4:eq-zones-cells-b}, applied at $i=j-1$, which is admissible because $j-1<k$,
gives $U_{j-1}\cap\Omega_j=\emptyset$. But $p$ lies in $\Delta\subset U_{j-1}$ and in
$\bigcup X_v\subset\Omega_j$, a contradiction. Hence every cell with a common point with
$\bigcup X_v$ has level at least $j-1$. If $j\le m_0+1$ there is no cell of level at most $j-2$
at all, and the assertion is vacuous.

(b) Let $\Delta$ be a cell and $Q$ a cube of $X_v$ with
$\operatorname{int}\Delta\cap\operatorname{int}Q\neq\emptyset$, and put
$m:=\operatorname{lev}\Delta$. Suppose first that $m<j$. By axiom (S1)
of~\eqref{5.4:eq-zones-cells-a}, $\Delta$ lies in one $\pi_j$-cube $P$ with
$\operatorname{int}\Delta\subset\operatorname{int}P$. Then
$\operatorname{int}P\cap\operatorname{int}Q\neq\emptyset$, and axiom (S2)
of~\eqref{5.4:eq-zones-cells-a} gives $P=Q$. Therefore, by
Lemma~\ref{5.4:lem-carrier-in-zone}(ii),
\begin{equation*}
\Delta\subset Q\subset\Omega_j^{\mathrm{cur}} .
\end{equation*}
The cells of level at least $j$ that intersect $\Delta$ are finitely many; call them
$\Delta''_1,\dots,\Delta''_N$. Each of them is distinct from $\Delta$, whose level is $m<j$, so
$\operatorname{int}\Delta\cap\operatorname{int}\Delta''_i=\emptyset$ for every $i$, distinct cells
having disjoint interiors. Since $\Delta\subset\Omega_j^{\mathrm{cur}}$ and
$\Omega_j^{\mathrm{cur}}$ is covered by the cells of level at least $j$, every point of
$\operatorname{int}\Delta$ lies in some $\Delta''_i$ but not in $\operatorname{int}\Delta''_i$,
hence in $\partial\Delta''_i$. Thus the non-empty open set $\operatorname{int}\Delta$ is
contained in $\bigcup_i\partial\Delta''_i$, a finite union of closed sets with empty interior;
such a union has empty interior, which is impossible. Hence $m\ge j$.

Now $m\ge j$, and by (S1) the cell $\Delta$ is the union of its $L^{d(m-j)}$ $\pi_j$-cubes $P'$.
The non-empty open set $\operatorname{int}\Delta\cap\operatorname{int}Q$ is contained in the
union of these finitely many cubes; it is not contained in the union of their boundaries, which
has empty interior, so it meets $\operatorname{int}P'$ for some $\pi_j$-cube $P'$ of $\Delta$.
Then $\operatorname{int}P'\cap\operatorname{int}Q\neq\emptyset$, and (S2) gives $P'=Q$. Hence
$Q\subset\Delta$.

(c) Let $Q$ be a cube of $X_v$. By Lemma~\ref{5.4:lem-carrier-in-zone}(ii),
$Q\subset\Omega_j^{\mathrm{cur}}$, so $Q$ is covered by the cells of level at least $j$ that
intersect $Q$; these are finitely many, say $\Delta''_1,\dots,\Delta''_N$. The union
$\bigcup_i\partial\Delta''_i$ has empty interior, so
$\operatorname{int}Q\not\subset\bigcup_i\partial\Delta''_i$, and a point of $\operatorname{int}Q$
outside this union lies in $\operatorname{int}\Delta''_i$ for some $i$. For this cell
$\Delta:=\Delta''_i$ we have $\operatorname{int}Q\cap\operatorname{int}\Delta\neq\emptyset$, that
is, $\Delta$ meets $X_v$; by (b) and its proof, $Q\subset\Delta$ and
$\operatorname{lev}\Delta\ge j$. If a second cell $\tilde\Delta\neq\Delta$ contained $Q$, then
$\operatorname{int}Q\subset\operatorname{int}\Delta\cap\operatorname{int}\tilde\Delta$, and two
distinct cells would share interior points, which is impossible. Hence $\Delta(Q):=\Delta$ is the
unique cell containing $Q$. Applied to the cube $q_v$ of $X_v$, this shows that
$\Delta'_v=\Delta(q_v)$ is well defined, and $y_v\in q_v\subset\Delta'_v$.

(d) The set $\bigcup X_v$ is connected. Measured in the units of level $j$, each cube of $X_v$
projects onto each coordinate axis as an interval of length one; the projection of
$\bigcup X_v$ on an axis is therefore a connected union of $N_j(X_v)$ such intervals, that is, an
interval of length at most $N_j(X_v)$. Taking the maximum over the axes gives
$\operatorname{diam}_{\infty,j}(\bigcup X_v)\le N_j(X_v)$. The bound
$N_j(X_v)\le\mathsf{c}_d(d_v+1)=\mathfrak{d}_v$ is (g1) of~\eqref{5.4:eq-zones-cells-c}. For
$m\le j$, a length measured in the units of level $m$ is $L^{j-m}$ times the same length measured
in the units of level $j$; this change of units gives
$\operatorname{diam}_{\infty,m}(\bigcup X_v)\le L^{j-m}\mathfrak{d}_v$.

(e) Let $\Delta_2$ be a cell of level $\ell$ meeting $X_v$. By (b), $\Delta_2$ contains a cube
$Q$ of $X_v$ and $j\le\ell$. By (S1), $\Delta_2$ is the union of its $L^{d(\ell-j)}$
$\pi_j$-cubes, and the proof of (b) showed that $Q$ coincides with one of them.

(f) Let $v\in\mathfrak{L}^{\mathrm{lg}}_{d}$ and $\Delta\in\mathrm{St}_v$. By (b), $\Delta$
contains a cube $Q$ of $X_v$; by nesting, $\Delta$ lies in one $\pi_k$-cube $P$, and then
$\operatorname{int}Q\subset\operatorname{int}P$. Were $P$ a cube of $Z$, we would have
$Q\subset Z$, which is excluded since $X_v\cap Z=\emptyset$. So $P$ is not among the cubes of $Z$;
hence $\operatorname{int}P\cap Z=\emptyset$, so $P\subset\operatorname{cl}Z^{c}$, and
$\Delta\not\subset Z$, because $\Delta\supset Q$ and $\operatorname{int}Q\subset\operatorname{int}P$
misses $Z$.

Consequently the points of a cell of $\mathrm{St}_v$ satisfy $\operatorname{dist}_k(\cdot,Z^c)=0$.
A cell of $\partial\mathrm{St}_v$ shares a wall with such a cell and has diameter at most $1$ in
the units of level $k$, so its points satisfy $\operatorname{dist}_k(\cdot,Z^c)\le 1$. The points
of a source cell satisfy $\operatorname{dist}_k(\cdot,Z^c)\ge 9\check{R}_k-1>1$, and the points of
$\Lambda^{\sim}$ satisfy $\operatorname{dist}_k(\cdot,Z^c)\ge 9\check{R}_k$, as fixed
in~\eqref{5.4:eq-zones-cells-e}. Therefore the cells of $\mathrm{St}_v\cup\partial\mathrm{St}_v$
miss $\Lambda^{\sim}$, and none of them is a source cell. Moreover, if $x\in\Delta_1$ with
$\Delta_1\in\partial\mathrm{St}_v$ and $y$ lies in a source cell, then, by the triangle
inequality for $\operatorname{dist}_k(\cdot,Z^c)$,
\begin{equation*}
|x-y|\ge\operatorname{dist}_k(y,Z^c)-\operatorname{dist}_k(x,Z^c)\ge 9\check{R}_k-2 ,
\end{equation*}
and $(\rho 2)$ of~\eqref{5.4:eq-zones-cells-d} gives $\rho(\Delta_1)\ge 9\check{R}_k-2$. For
$\Delta_1\in\mathrm{St}_v$ the same argument, with $\operatorname{dist}_k(x,Z^c)=0$, gives
$\rho(\Delta_1)\ge 9\check{R}_k-1$.

Wall-connectedness: by (c), wall-neighbouring cubes of $X_v$ lie in cells that coincide or meet
in a $(d-1)$-dimensional set, and this gives the wall-connectedness of $\mathrm{St}_v$. Finally,
a member of $\mathrm{St}_v$ lies in one $\pi_k$-cube, as shown above, and this $\pi_k$-cube
contains the member's cube of $X_v$; hence $[\mathrm{St}_v]_k\subset K_v$.

(g) By (b) each member of $\mathrm{St}_v$ contains a cube of $X_v$, and by (c) a cube of $X_v$
lies in only one cell; hence distinct members contain distinct cubes, and, with (d),
$\#\mathrm{St}_v\le N_j(X_v)\le\mathfrak{d}_v$. By axiom (A2) of~\eqref{5.4:eq-zones-cells-b},
each cell has at most $2dL^{d-1}$ wall-neighbours; every member of $\partial\mathrm{St}_v$ is a
wall-neighbour of a member of $\mathrm{St}_v$, so
\begin{equation*}
\#\partial\mathrm{St}_v\le 2dL^{d-1}\,\#\mathrm{St}_v\le 2dL^{d-1}\,\mathsf{c}_d(d_v+1)
\le\hat{b}\,\mathsf{c}_d(d_v+1),
\end{equation*}
and the condition $(\epsilon)$ of~\eqref{5.4:eq-zones-cells-i} then gives
$\epsilon_{\partial}\,\#\partial\mathrm{St}_v\le\tfrac12(d_v+1)$.

(h) A cell meeting $X_v$ has level at least $j$, by (b). A grouped cube $G$ of the grain
$\pi^{\mathfrak{s}}_j$ is a $\pi_j$-cube with
$G\subset\operatorname{cl}\bigl((\Omega_j^{\mathrm{cur}})^{c}\bigr)$. Suppose that $G$ meets
$X_v$, that is, $\operatorname{int}G\cap\operatorname{int}Q\neq\emptyset$ for a cube $Q$ of $X_v$;
then, as in the proof of (b), (S2) gives $G=Q$. But $Q\subset\Omega_j^{\mathrm{cur}}$ by
Lemma~\ref{5.4:lem-carrier-in-zone}(ii), so
$\operatorname{int}Q\subset\operatorname{int}\Omega_j^{\mathrm{cur}}$, while
$\operatorname{int}Q=\operatorname{int}G$ lies in
$\operatorname{cl}\bigl((\Omega_j^{\mathrm{cur}})^{c}\bigr)$, the complement of
$\operatorname{int}\Omega_j^{\mathrm{cur}}$; since $\operatorname{int}Q$ is non-empty, this is a
contradiction. Hence no
grouped cube of $\pi^{\mathfrak{s}}_j$ and no cell of level $<j$ meets $X_v$, which is the
equality $\mathrm{St}^j_v=\mathrm{St}_v$. This argument uses only (b) and
Lemma~\ref{5.4:lem-carrier-in-zone}, and so applies equally to $v\in\mathfrak{L}^{\mathrm{lg}}_{c}$.

It remains to show $\mathfrak{H}_v=\emptyset$. The hubs lie at $\operatorname{dist}_k$-distance
at least $9\check{R}_k$ from $Z^{c}$, in particular inside $Z$. Since $X_v\cap Z=\emptyset$, $X_v$
meets no hub. A member of $\mathrm{St}_v$ lies in $\operatorname{cl}Z^{c}$ by (f), so its points
are at $\operatorname{dist}_k$-distance $0$ from $Z^{c}$; a wall shared with a hub would consist
of points at distance $0$ and at distance at least $9\check{R}_k$ from $Z^{c}$, which is
impossible. Hence no member of $\mathrm{St}_v$ shares a wall with a hub, and
$\mathfrak{H}_v=\emptyset$.
\end{proof}

\begin{lemma}[The radius-one neighbourhood of a cell lies in either star]
\label{5.4:lem-neighbourhood-star}
Let $\Delta_0$ be a cell of level $n_0$, and let
\begin{equation*}
\mathcal{N}(\Delta_0)=\{x:\ \operatorname{dist}_{n_0}(x,\Delta_0)<1\}
\end{equation*}
be its open radius-one neighbourhood, measured in the units of level $n_0$. Consider the two forms
of the star of $\Delta_0$ that enter the split \eqref{5.4:eq-pointed-units-a} of
Definition~\ref{5.4:def-pointed-units}: the union of the cells $\Delta$ with
$\operatorname{dist}_{n_0}(\Delta,\Delta_0)<1$, and the union $\bigcup\mathrm{St}(\Delta_0)$.
\begin{enumerate}
\item[(i)] $\mathcal{N}(\Delta_0)$ is contained in the union of the cells $\Delta$ with
$\operatorname{dist}_{n_0}(\Delta,\Delta_0)<1$.
\item[(ii)] $\mathcal{N}(\Delta_0)\subset\Omega_{n_0-1}$ and
$\mathcal{N}(\Delta_0)\subset\Omega^{\mathrm{cur}}_{n_0-1}$; hence
\begin{equation}
\label{5.4:eq-neighbourhood-star-1}
\begin{split}
\mathcal{N}(\Delta_0)&\subset\{x:\ \operatorname{dist}_{n_0}(x,\Delta_0)\le 1\}\cap\Omega_{n_0-1}\\
&\subset\bigcup\mathrm{St}(\Delta_0),
\end{split}
\end{equation}
the last inclusion being the inclusion of the star stated with the split
\eqref{5.4:eq-pointed-units-a} in Definition~\ref{5.4:def-pointed-units}.
\item[(iii)] Consequently, if $v$ is a pointed unit with
$\bigcup X_v\subset\mathcal{N}(\Delta_0(v))$, then $v$ is short under either form of the star.
Contrapositively, every long single and every long tail $v$ satisfies
$\bigcup X_v\not\subset\mathcal{N}(\Delta_0(v))$.
\end{enumerate}
\end{lemma}

\begin{proof}
(i) Let $x\in\mathcal{N}(\Delta_0)$. Since the cells cover the underlying set, $x$ lies in some
cell $\Delta$. The distance between two sets is the infimum of the distances between their points,
and $x\in\Delta$; therefore
\begin{equation*}
\operatorname{dist}_{n_0}(\Delta,\Delta_0)\le\operatorname{dist}_{n_0}(x,\Delta_0)<1 .
\end{equation*}
Thus $x$ lies in a cell at $\operatorname{dist}_{n_0}$-distance less than $1$ from $\Delta_0$.

(ii) The cell $\Delta_0$, being of level $n_0$, lies in $\Omega_{n_0}$. Let $x$ satisfy
$\operatorname{dist}_{n_0}(x,\Delta_0)<1$. We first show that $x\notin U_{n_0-1}$, where
$U_{n_0-1}=\mathrm{cl}(\Omega^{c}_{n_0-1})$.

If $n_0=m_0$, then $\Omega_{n_0-1}$ is the whole underlying set, so $\Omega^{c}_{n_0-1}$ and its
closure $U_{n_0-1}$ are empty, and $x\notin U_{n_0-1}$.

Otherwise we apply the axiom (A3) among \eqref{5.4:eq-zones-cells-b} of
Notation~\ref{5.4:not-zones-cells} with index $i=n_0-1$, which is admissible since $n_0-1<k$. It
gives
\begin{equation*}
\operatorname{dist}_{n_0}(U_{n_0-1},\Omega_{n_0})\ge 1 .
\end{equation*}
Since $\Delta_0\subset\Omega_{n_0}$, the point $x$ satisfies
\begin{equation*}
\operatorname{dist}_{n_0}(x,\Omega_{n_0})\le\operatorname{dist}_{n_0}(x,\Delta_0)<1 ,
\end{equation*}
so $x\notin U_{n_0-1}$.

In both cases, $U_{n_0-1}\supseteq\Omega^{c}_{n_0-1}$ gives $x\notin\Omega^{c}_{n_0-1}$, that is,
$x\in\Omega_{n_0-1}$. Moreover, any cell containing $x$ and of level at most $n_0-2$ would lie in
$\mathrm{cl}(\Omega^{c}_{n_0-1})=U_{n_0-1}$, which does not contain $x$. Hence every cell through
$x$ has level at least $n_0-1$. Since $x$ lies in at least one cell, and $\Omega^{\mathrm{cur}}_{n_0-1}$
is the union of the cells of level at least $n_0-1$, we get $x\in\Omega^{\mathrm{cur}}_{n_0-1}$.
As $x\in\mathcal{N}(\Delta_0)$ was arbitrary, both inclusions of (ii) hold.

For the first inclusion of \eqref{5.4:eq-neighbourhood-star-1}: every $x\in\mathcal{N}(\Delta_0)$ satisfies
$\operatorname{dist}_{n_0}(x,\Delta_0)<1\le 1$ and, by what was just shown, $x\in\Omega_{n_0-1}$.
The second inclusion of \eqref{5.4:eq-neighbourhood-star-1} is the inclusion
$\bigcup\mathrm{St}(\Delta_0)\supseteq\{x:\ \operatorname{dist}_{n_0}(x,\Delta_0)\le 1\}\cap\Omega_{n_0-1}$
stated with the split \eqref{5.4:eq-pointed-units-a} in Definition~\ref{5.4:def-pointed-units},
used here as stated there.

(iii) Suppose $\bigcup X_v\subset\mathcal{N}(\Delta_0(v))$. Part (i), applied to the cell
$\Delta_0=\Delta_0(v)$ of level $n_0$, places $\bigcup X_v$ inside the union of the cells at
$\operatorname{dist}_{n_0}$-distance less than $1$ from $\Delta_0(v)$. Part (ii), in the form
\eqref{5.4:eq-neighbourhood-star-1} for the same cell, places $\bigcup X_v$ inside
$\bigcup\mathrm{St}(\Delta_0(v))$. Thus $\bigcup X_v$ lies inside the star of $\Delta_0(v)$ in each of
its two forms, and by the split \eqref{5.4:eq-pointed-units-a}, $v$ is short under either form of the star.

The long singles and the long tails are, by \eqref{5.4:eq-pointed-units-a}, the pointed units that are
not short. Hence, by the contrapositive of the preceding paragraph, every long single and every long tail
$v$ satisfies $\bigcup X_v\not\subset\mathcal{N}(\Delta_0(v))$.
\end{proof}

\begin{lemma}[Length from exit of the star]\label{5.4:lem-star-exit-length}
Let $v\in\mathfrak{L}^{\mathrm{lg}}$ be a long unit, in the sense of
Definition~\ref{5.4:def-pointed-units}, and write $j:=j_v$ and $\Delta_0:=\Delta_0(v)$. Let $n_0$
be the level of the cell $\Delta_0$, and put $R:=L^{n_0-j}$; in this lemma and its proof $R$
denotes this number only. Then there are points
\begin{equation*}
y_v\in\bigl(\textstyle\bigcup X_v\bigr)\cap\Delta_0
\qquad\text{and}\qquad
x\in\textstyle\bigcup X_v \ \text{ with } \ \operatorname{dist}_{n_0}(x,\Delta_0)\ge 1 .
\end{equation*}
If $Q_x$ is a cube of $X_v$ containing $x$, then
\begin{equation}\label{5.4:eq-star-exit-length-1}
\operatorname{dist}_j(Q_x,q_v)\ge R-2,
\end{equation}
and $\operatorname{dist}_j(Q_x,q_v)\ge R-1$ if $q_v\subset\Delta_0$, as happens for instance when
$\Delta_0=\Delta'_v$. Hence, by \eqref{5.4:eq-zones-cells-c},
\begin{equation}\label{5.4:eq-star-exit-length-2}
d_v\ge\bigl(L^{n_0-j}-\delta\bigr)_+ ,
\end{equation}
with $\delta=2$ in general and $\delta=1$ if $q_v\subset\Delta_0$.
\end{lemma}

\begin{proof}
Since $v$ is long, part~(iii) of Lemma~\ref{5.4:lem-neighbourhood-star} gives
$\bigcup X_v\not\subset\mathcal{N}(\Delta_0)$, and the display \eqref{5.4:eq-pointed-units-a},
read through that part of the lemma, provides the points $y_v\in(\bigcup X_v)\cap\Delta_0$ and
$x\in\bigcup X_v$ outside $\mathcal{N}(\Delta_0)$. As $\mathcal{N}(\Delta_0)$ is the open
neighbourhood of radius one of $\Delta_0$ measured in level-$n_0$ units, the point $x$ satisfies
$\operatorname{dist}_{n_0}(x,\Delta_0)\ge 1$. One level-$n_0$ unit is $L^{n_0-j}=R$ level-$j$
units, so
\begin{equation*}
\operatorname{dist}_j(x,\Delta_0)\ge R .
\end{equation*}

Write $|\cdot|$ for distance measured in level-$j$ units, and let $z\in Q_x$ and $w\in q_v$.
By the triangle inequality,
\begin{equation*}
|z-w|\ \ge\ |x-y_v|-|x-z|-|y_v-w| .
\end{equation*}
Since $y_v\in\Delta_0$ and $\operatorname{dist}_j(x,\Delta_0)\ge R$, we have $|x-y_v|\ge R$. The
cubes $Q_x$ and $q_v$ have diameter $1$ in level-$j$ units, so $|x-z|\le 1$ and
$|y_v-w|\le 1$. Hence $|z-w|\ge R-2$ for all $z\in Q_x$ and $w\in q_v$, and taking the infimum
over such pairs gives \eqref{5.4:eq-star-exit-length-1}.

If $q_v\subset\Delta_0$, then every $w\in q_v$ lies in $\Delta_0$, so $|x-w|\ge R$. For
$z\in Q_x$ this gives
\begin{equation*}
|z-w|\ \ge\ |x-w|-|x-z|\ \ge\ R-1 ,
\end{equation*}
and taking the infimum gives $\operatorname{dist}_j(Q_x,q_v)\ge R-1$.

Finally, apply \eqref{5.4:eq-zones-cells-c} to the pair $(Q_x,q_v)$. With the lower bounds
$\operatorname{dist}_j(Q_x,q_v)\ge R-\delta$ just obtained, where $\delta=2$ in general and
$\delta=1$ when $q_v\subset\Delta_0$, this gives \eqref{5.4:eq-star-exit-length-2}.
\end{proof}

\begin{lemma}[Length of a long unit at a met cell]\label{5.4:lem-met-cell-length}
Let $v\in\mathfrak{L}^{\mathrm{lg}}$, and let $\Delta_2$ be a cell of level $\ell$ which meets
$X_v$ and satisfies $u:=\ell-j_v\ge 1$. Then
\begin{equation}\label{5.4:eq-met-cell-length-1}
  d_v\;\ge\;\mathfrak{m}_u=\bigl(L^{u-1}-\delta\bigr)_+ ,
\end{equation}
with $\mathfrak{m}_u$ and $\delta$ as in \eqref{5.4:eq-transport-constants-a}. More precisely:
\begin{itemize}
\item[$(\alpha)$] if $n_0(v)\ge \ell-1$, then the length from exit of the star
  (Lemma~\ref{5.4:lem-star-exit-length}) gives $d_v\ge L^{u-1}-\delta$;
\item[$(\beta)$] if $n_0(v)\le \ell-2$, then $\operatorname{dist}_\ell(\Delta_0(v),\Delta_2)\ge 1$ by
  the axiom (A3) of \eqref{5.4:eq-zones-cells-b} at $i=\ell-1<k$, and
  $d_v\ge \operatorname{dist}_j(q_v,Q)\ge L^u-1$, where $j:=j_v$ and $Q\subset\Delta_2$ is the cube
  of $X_v$ provided by Lemma~\ref{5.4:lem-cells-carrier}\,(e).
\end{itemize}
\end{lemma}

\begin{proof}
Write $j:=j_v$, $\Delta_0:=\Delta_0(v)$ and $n_0:=n_0(v)$. Since $\Delta_2$ is a cell of level
$\ell$ meeting $X_v$, part (e) of the lemma on cells against a carrier
(Lemma~\ref{5.4:lem-cells-carrier}) gives $\ell\ge j$ and a cube $Q$ of $X_v$ with $Q\subset\Delta_2$.

Case $(\alpha)$: $n_0\ge\ell-1$. Then
\begin{equation*}
  n_0-j\;\ge\;\ell-1-j\;=\;u-1 .
\end{equation*}
Lemma~\ref{5.4:lem-star-exit-length} gives $d_v\ge(L^{n_0-j}-\delta)_+$. Since $L\ge 1$ and
$n_0-j\ge u-1$, we have $L^{n_0-j}\ge L^{u-1}$, and the positive part is non-decreasing, so
\begin{equation*}
  d_v\;\ge\;\bigl(L^{n_0-j}-\delta\bigr)_+\;\ge\;\bigl(L^{u-1}-\delta\bigr)_+\;=\;\mathfrak{m}_u .
\end{equation*}

Case $(\beta)$: $n_0\le\ell-2$. The cell $\Delta_0$ has level $n_0$, so it lies in the closure of
the zone of level $n_0$, and hence in the closure of $\Omega^{c}_{n_0+1}$, that is, in
$U_{n_0+1}$. Since $n_0+1\le\ell-1$, the axiom (A1) of \eqref{5.4:eq-zones-cells-b} gives
$U_{n_0+1}\subset U_{\ell-1}$. Altogether
\begin{equation*}
  \Delta_0\;\subset\;\overline{\text{zone } n_0}\;\subset\;\overline{\Omega^{c}_{n_0+1}}
  \;\subset\;U_{\ell-1} .
\end{equation*}
On the other hand, the cell $\Delta_2$ has level $\ell$, so $\Delta_2\subset\Omega_\ell$. The axiom
(A3) of \eqref{5.4:eq-zones-cells-b} at $i=\ell-1<k$ therefore gives
$\operatorname{dist}_\ell(\Delta_0,\Delta_2)\ge 1$. Passing from level-$\ell$ units to level-$j$
units multiplies distances by $L^{\ell-j}=L^u$, so this reads
\begin{equation}\label{5.4:eq-met-cell-length-2}
  \operatorname{dist}_j(\Delta_0,\Delta_2)\;\ge\;L^u .
\end{equation}
The point $y_v$ lies in $\Delta_0\cap q_v$. Let $z\in q_v$ and $w\in Q$. Since $w\in
Q\subset\Delta_2$ and $y_v\in\Delta_0$, \eqref{5.4:eq-met-cell-length-2} gives $|y_v-w|\ge L^u$ in
level-$j$ units, and since $z$ and $y_v$ lie in the same level-$j$ cube $q_v$, $|y_v-z|\le 1$.
Hence
\begin{equation*}
  |z-w|\;\ge\;|y_v-w|-|y_v-z|\;\ge\;L^u-1 ,
\end{equation*}
and taking the infimum over $z\in q_v$, $w\in Q$ yields $\operatorname{dist}_j(q_v,Q)\ge L^u-1$.
Both $q_v$ and $Q$ are cubes of $X_v$, so the axiom (g4) of \eqref{5.4:eq-zones-cells-c} applied
to the pair $(q_v,Q)$ gives
\begin{equation*}
  d_v\;\ge\;\operatorname{dist}_j(q_v,Q)\;\ge\;L^u-1 .
\end{equation*}
Finally $L^u-1\ge 0$, and $L^u-1\ge L^{u-1}-\delta$ because
$L^u-L^{u-1}=L^{u-1}(L-1)\ge 1$ for $u\ge 1$ and $L\ge 2$, while $\delta\ge 0$; hence
$d_v\ge(L^{u-1}-\delta)_+=\mathfrak{m}_u$.

The two cases exhaust all possibilities. Indeed, the cell $\Delta_0$ contains
$y_v\in\bigcup X_v$, so it has a common point with $\bigcup X_v$, and part (a) of
Lemma~\ref{5.4:lem-cells-carrier} gives $n_0\ge j-1$. If $n_0=j-1$, then case $(\alpha)$ would
require $j-1\ge\ell-1$, that is $u\le 0$; so for $u\ge 1$ the value $n_0=j-1$ satisfies
$n_0=j-1\le\ell-2$ and falls under $(\beta)$. Every other value satisfies either
$n_0\ge\ell-1\ge j$, which is case $(\alpha)$, or $n_0\le\ell-2$, which is case $(\beta)$. Thus
every $u\ge 1$ is covered.

In case $(\beta)$ only the inequality $\operatorname{dist}_\ell\ge 1$ of the axiom (A3) is used.
Consequently thin prepared strata, at which (A3) holds with equality inside a foreign frozen
window, change nothing in the argument.
\end{proof}

\begin{lemma}[The shell series]\label{5.4:lem-shell-series}
Let $\mathfrak{c}_u$ $(u\ge 0)$, $\mathfrak{A}_{\delta}(c)$ and $C^{LS}_{\sharp}(c)$ be the constants of the
two transports, defined in \eqref{5.4:eq-transport-constants-a}. For every $c>0$ and every
$\delta\ge 0$,
\begin{equation}\label{5.4:eq-shell-series-1}
\sum_{u\ge 0}\mathfrak{c}_u L^{du}e^{-c(L^{u-1}-\delta)_+}
\le 1+C_{\mathsf{b}}2^{p_{\mathsf{b}}}L^d\bigl[1+\mathfrak{A}_{\delta}(c)\bigr],
\end{equation}
and
\begin{equation}\label{5.4:eq-shell-series-2}
\sup_{u\ge 0}e^{u}L^{du}e^{-c(L^{u-1}-\delta)_+}\le C^{LS}_{\sharp}(c).
\end{equation}
\end{lemma}

\begin{proof}
Recall from \eqref{5.4:eq-transport-constants-a} that $\mathfrak{c}_0=1$ and
$\mathfrak{c}_u=C_{\mathsf{b}}(1+u)^{p_{\mathsf{b}}}$ for $u\ge 1$. We bound the terms of the series
in \eqref{5.4:eq-shell-series-1} separately for $u=0$, $u=1$ and $u\ge 2$.

For $u=0$ the term is $e^{-c(L^{-1}-\delta)_+}$; its exponent is non-positive, since $c>0$ and
$(\,\cdot\,)_+\ge 0$, so the term is at most $1$.

For $u=1$ the exponent is $-c(1-\delta)_+\le 0$, so the term is at most
$\mathfrak{c}_1L^d=C_{\mathsf{b}}2^{p_{\mathsf{b}}}L^d$. No decay is gained from this term.

Let $u\ge 2$ and put $y:=L^{u-1}$; since $L\ge 3$, we have $y\ge 3$. Then $L^{du}=L^dy^d$.
Moreover $1+u\le 2\cdot 3^{u-1}$ for every $u\ge 1$ (equality at $u=1$; passing from $u$ to $u+1$
adds $1$ on the left and $4\cdot 3^{u-1}\ge 1$ on the right), and $3^{u-1}\le L^{u-1}=y$.
Hence $1+u\le 2y$, so that $(1+u)^{p_{\mathsf{b}}}\le(2y)^{p_{\mathsf{b}}}$. Finally
$(y-\delta)_+\ge y-\delta$ and $c>0$ give
$e^{-c(y-\delta)_+}\le e^{\delta c}e^{-cy}$. Writing $q_{\mathsf{b}}:=d+p_{\mathsf{b}}$, so that
$y^d(2y)^{p_{\mathsf{b}}}=2^{p_{\mathsf{b}}}y^{q_{\mathsf{b}}}$, the $u$-th term is therefore bounded by
\begin{equation}\label{5.4:eq-shell-series-3}
\mathfrak{c}_uL^{du}e^{-c(y-\delta)_+}
\le C_{\mathsf{b}}2^{p_{\mathsf{b}}}L^de^{\delta c}\,y^{q_{\mathsf{b}}}e^{-cy}
\le C_{\mathsf{b}}2^{p_{\mathsf{b}}}L^de^{\delta c}
\Bigl(\frac{2q_{\mathsf{b}}}{ec}\Bigr)^{q_{\mathsf{b}}}e^{-cy/2}.
\end{equation}
For the second inequality we split $e^{-cy}=e^{-cy/2}e^{-cy/2}$ and use
\begin{equation*}
y^{q_{\mathsf{b}}}e^{-cy/2}\le\sup_{t>0}t^{q_{\mathsf{b}}}e^{-ct/2}
=\Bigl(\frac{2q_{\mathsf{b}}}{ec}\Bigr)^{q_{\mathsf{b}}};
\end{equation*}
the supremum is attained at $t=2q_{\mathsf{b}}/c$, where the derivative of
$q_{\mathsf{b}}\log t-ct/2$ vanishes, and its
value is $(2q_{\mathsf{b}}/c)^{q_{\mathsf{b}}}e^{-q_{\mathsf{b}}}$.

It remains to sum the factor $e^{-cy/2}=e^{-cL^{u-1}/2}$ over $u\ge 2$. With $m:=u-1$ and the
inequality $L^m\ge 3^m\ge 1+2m$ for $m\ge 1$ (Bernoulli's inequality for $(1+2)^m$),
\begin{equation}\label{5.4:eq-shell-series-4}
\sum_{u\ge 2}e^{-cL^{u-1}/2}=\sum_{m\ge 1}e^{-cL^{m}/2}
\le\sum_{m\ge 1}e^{-c(1+2m)/2}
=e^{-c/2}\,\frac{e^{-c}}{1-e^{-c}}
=e^{-3c/2}(1-e^{-c})^{-1}.
\end{equation}
Multiplying \eqref{5.4:eq-shell-series-4} by the $u$-independent prefactor in
\eqref{5.4:eq-shell-series-3} gives
\begin{equation*}
\sum_{u\ge 2}\mathfrak{c}_uL^{du}e^{-c(L^{u-1}-\delta)_+}
\le C_{\mathsf{b}}2^{p_{\mathsf{b}}}L^d\Bigl(\frac{2q_{\mathsf{b}}}{ec}\Bigr)^{q_{\mathsf{b}}}
e^{(\delta-3/2)c}(1-e^{-c})^{-1}
=C_{\mathsf{b}}2^{p_{\mathsf{b}}}L^d\,\mathfrak{A}_{\delta}(c),
\end{equation*}
the last equality being the definition of $\mathfrak{A}_{\delta}(c)$ in
\eqref{5.4:eq-transport-constants-a}. Adding the bounds $1$ for $u=0$ and
$C_{\mathsf{b}}2^{p_{\mathsf{b}}}L^d$ for $u=1$ yields \eqref{5.4:eq-shell-series-1}.

For the supremum \eqref{5.4:eq-shell-series-2}, let first $u\ge 2$ and $y=L^{u-1}$ as above.
Since $L\ge 3\ge e$, we have $e^{u}=e\cdot e^{u-1}\le e\cdot L^{u-1}=ey$. Together with
$L^{du}=L^dy^d$ and $e^{-c(y-\delta)_+}\le e^{-c(y-\delta)}$ this gives
\begin{equation*}
e^{u}L^{du}e^{-c(y-\delta)_+}
\le eL^de^{\delta c}\,y^{d+1}e^{-cy}
\le eL^de^{\delta c}\Bigl(\frac{d+1}{ec}\Bigr)^{d+1},
\end{equation*}
where we used $\sup_{t>0}t^{d+1}e^{-ct}=((d+1)/(ec))^{d+1}$, attained at $t=(d+1)/c$.
For $u\le 1$ the exponential factor $e^{-c(L^{u-1}-\delta)_+}$ is at most $1$, so the quantity
is at most $e^{u}L^{du}\le eL^d$. Both bounds are at most
$eL^d\max\{1,((d+1)/(ec))^{d+1}e^{\delta c}\}=C^{LS}_{\sharp}(c)$ by
\eqref{5.4:eq-transport-constants-a}, which proves \eqref{5.4:eq-shell-series-2}.
\end{proof}

\begin{proposition}[The per-cube transport]\label{5.4:prop-per-cube-transport}
Assume the following.
\begin{enumerate}
\item The axioms (S1)--(S3) of \eqref{5.4:eq-zones-cells-a} and (A0)--(A3) of
\eqref{5.4:eq-zones-cells-b} hold, together with (g4) of \eqref{5.4:eq-zones-cells-c}.
\item The condition on the classes and ranges of the preparation attempts holds.
\item Every pointed unit $v$ satisfies $X_v\subset\Lambda_{j_v}$ at the step $j_v$ at which it is
created.
\item The split \eqref{5.4:eq-pointed-units-a} of the pointed units into short and long units holds;
$\mathfrak{L}^{\mathrm{lg}}$ is its class of long units.
\item The per-cube base bound for long units \eqref{5.4:eq-age-base-bounds-b} holds, with its
constant $c_{\mathrm{sp}}>0$. In the form in which it is used below, it states that for every level
$j$, every cube $\square'\in\pi_j$ and the $r'$ of the statement below,
\begin{equation*}
\sum_{\substack{v\in\mathfrak{L}^{\mathrm{lg}}:\ j_v=j,\ \square'\in X_v}}
w_{r'}(v)\,e^{(2+c_{\mathrm{sp}})d_v}\le\mathsf{b}^{\mathrm{lg}}(k-j).
\end{equation*}
\end{enumerate}
Fix $K$, $k\le K$ and a label. Let $\Delta_2$ be any cell; in the off-$Z$ variant of the met-cell
sums of \eqref{5.4:eq-pointed-units-c}, let it be any cell with $\Delta_2\not\subset Z$. Write
$\ell:=\operatorname{lev}\Delta_2$ and let
$r'\le\frac{8}{5}\kappa$. Then the met-cell sum of \eqref{5.4:eq-pointed-units-c} over the long units
satisfies
\begin{equation}\label{5.4:eq-per-cube-transport-1}
\begin{aligned}
S_{r'}(\Delta_2;\mathfrak{L}^{\mathrm{lg}})
&\le\Bigl[\sum_{u\ge 0}\mathfrak{c}_u\,L^{du}\,e^{-c_{\mathrm{sp}}\mathfrak{m}_u}\Bigr]\,
\mathsf{b}^{\mathrm{lg}}(k-\ell)\\
&\le C^{LS}\,\mathsf{b}^{\mathrm{lg}}(k-\ell).
\end{aligned}
\end{equation}
Here $\mathsf{b}^{\mathrm{lg}}$ is the age function of Definition~\ref{5.4:def-age-base-bounds}, with
its growth constants $C_{\mathsf{b}}$, $p_{\mathsf{b}}$, and $\mathfrak{m}_u$, $\mathfrak{c}_u$ and
$C^{LS}$ are the constants of
\eqref{5.4:eq-transport-constants-a}. The constant $C^{LS}$ depends only on
$d,L,c_{\mathrm{sp}},C_{\mathsf{b}},p_{\mathsf{b}}$. The difference $k-\ell$ enters only as the
argument of $\mathsf{b}^{\mathrm{lg}}$. No constant in \eqref{5.4:eq-per-cube-transport-1} depends on
the position of $y_v$, on $Z$ or $\Lambda$, on a hub $\Lambda^{\sim}$, on a window, or on the radius
$\check{R}_k$.
\end{proposition}

\begin{proof}
Let $v\in\mathfrak{L}^{\mathrm{lg}}$ be a long unit such that $X_v$ meets $\Delta_2$. Lemma~\ref{5.4:lem-cells-carrier},
part (e), on cells against a carrier, applies to the cell $\Delta_2$ of level $\ell$. It gives
\begin{equation*}
u(v):=\ell-j_v\ge 0 .
\end{equation*}
It also gives a non-empty set of cubes $\square'\in\pi_{j_v}$ with $\square'\subset\Delta_2$ and
$\square'\in X_v$. For each such cube we have $\square'\in\pi_{\ell-u(v)}$.

All terms $w_{r'}(v)e^{2d_v}$ of the met-cell sum are non-negative. Each $v$ that contributes has at least one such cube $\square'$. Hence
counting each $v$ once for every such $\square'$ can only enlarge the sum. Grouping the units according
to the value $u$ of $u(v)$, we obtain
\begin{equation}\label{5.4:eq-per-cube-transport-2}
S_{r'}(\Delta_2;\mathfrak{L}^{\mathrm{lg}})
\le\sum_{u\ge 0}\ \sum_{\substack{\square'\in\pi_{\ell-u}\\ \square'\subset\Delta_2}}\
\sum_{\substack{v:\ j_v=\ell-u,\ \square'\in X_v\\ X_v\text{ meets }\Delta_2}}
w_{r'}(v)\,e^{2d_v}.
\end{equation}

Fix $u\ge 0$, a cube $\square'$ and a unit $v$ in the innermost sum of
\eqref{5.4:eq-per-cube-transport-2}; we claim $d_v\ge\mathfrak{m}_u$. If $u\ge 1$, then $\Delta_2$
is a cell of level $\ell$ meeting $X_v$ with $\ell-j_v=u\ge 1$. The lemma on the length of a long unit
at a met cell (Lemma~\ref{5.4:lem-met-cell-length}) then gives $d_v\ge\mathfrak{m}_u$. If $u=0$, the claim
is the inequality $d_v\ge 0$. Since $c_{\mathrm{sp}}>0$, the claim gives
\begin{equation*}
e^{2d_v}=e^{(2+c_{\mathrm{sp}})d_v}\,e^{-c_{\mathrm{sp}}d_v}
\le e^{-c_{\mathrm{sp}}\mathfrak{m}_u}\,e^{(2+c_{\mathrm{sp}})d_v}.
\end{equation*}

We insert this bound term by term into the innermost sum and then drop the condition that $X_v$
meets $\Delta_2$; the terms are non-negative, so this again only enlarges. The innermost sum is
thereby bounded by $e^{-c_{\mathrm{sp}}\mathfrak{m}_u}$ times the sum of
$w_{r'}(v)e^{(2+c_{\mathrm{sp}})d_v}$ over all units $v$ with $j_v=\ell-u$ and $\square'\in X_v$.
To this sum we apply
the per-cube base bound \eqref{5.4:eq-age-base-bounds-b} at the level $\ell-u$ and the cube
$\square'$. We then use \eqref{5.4:eq-age-base-bounds-a} at $s=k-\ell$; for $u=0$ nothing needs to be
applied. This bounds the innermost sum of \eqref{5.4:eq-per-cube-transport-2} as follows:
\begin{equation}\label{5.4:eq-per-cube-transport-3}
\begin{aligned}
\sum_{\substack{v:\ j_v=\ell-u,\ \square'\in X_v\\ X_v\text{ meets }\Delta_2}}
w_{r'}(v)\,e^{2d_v}
&\le e^{-c_{\mathrm{sp}}\mathfrak{m}_u}\,\mathsf{b}^{\mathrm{lg}}(k-\ell+u)\\
&\le e^{-c_{\mathrm{sp}}\mathfrak{m}_u}\,\mathfrak{c}_u\,\mathsf{b}^{\mathrm{lg}}(k-\ell).
\end{aligned}
\end{equation}

The right-hand side of \eqref{5.4:eq-per-cube-transport-3} does not depend on $\square'$. By (S1) of
\eqref{5.4:eq-zones-cells-a}, the
cubes $\square'\in\pi_{\ell-u}$ contained in the cell $\Delta_2$ of level $\ell$ are $L^{du}$ in
number. Inserting \eqref{5.4:eq-per-cube-transport-3} into \eqref{5.4:eq-per-cube-transport-2} and
summing over $\square'$ and then over $u\ge 0$ gives the first inequality of
\eqref{5.4:eq-per-cube-transport-1}.

The second inequality follows from the shell series lemma (Lemma~\ref{5.4:lem-shell-series}) applied with
$c=c_{\mathrm{sp}}>0$ and the $\delta$ entering $\mathfrak{m}_u$. That lemma bounds the bracket in
\eqref{5.4:eq-per-cube-transport-1} by
\begin{equation*}
1+C_{\mathsf{b}}2^{p_{\mathsf{b}}}L^{d}\bigl[1+\mathfrak{A}_{\delta}(c_{\mathrm{sp}})\bigr],
\end{equation*}
which is the constant $C^{LS}$ of \eqref{5.4:eq-transport-constants-a}.
\end{proof}

\begin{corollary}[The long-unit bound at a cell of the current level]
\label{5.4:cor-current-level-bound}
At every cell $\Delta_2$ of level $k$,
\begin{equation}\label{5.4:eq-current-level-bound-1}
S_{r'}(\Delta_2;\mathfrak{L}^{\mathrm{lg}}_{c})
\le S_{r'}(\Delta_2;\mathfrak{L}^{\mathrm{lg}})
\le \mathsf{K}^{\mathrm{lg}} := C^{LS}\,\mathsf{b}^{\mathrm{lg}}(0).
\end{equation}
\end{corollary}

\begin{proof}
Throughout, $\operatorname{lev}\Delta$ denotes the level of a cell $\Delta$, and $r'$ is the
exponent of the per-cube transport (Proposition~\ref{5.4:prop-per-cube-transport}), so
$r'\le (8/5)\kappa$.

The sums $S_{r'}(\Delta_2;\cdot)$ are those of \eqref{5.4:eq-pointed-units-c}. The first
inequality in \eqref{5.4:eq-current-level-bound-1} compares the sum over the long units of
class (c), $\mathfrak{L}^{\mathrm{lg}}_{c}$, with the sum over all long units,
$\mathfrak{L}^{\mathrm{lg}}$, of which they form a part.

For the second inequality, apply the per-cube transport,
Proposition~\ref{5.4:prop-per-cube-transport}, to the cell $\Delta_2$. Its level is
$\ell:=\operatorname{lev}\Delta_2=k$, so $k-\ell=0$. Since $k-\ell$ occurs in that
proposition only as the argument of the age function $\mathsf{b}^{\mathrm{lg}}$, its second
inequality gives
\begin{equation*}
S_{r'}(\Delta_2;\mathfrak{L}^{\mathrm{lg}})
\le C^{LS}\,\mathsf{b}^{\mathrm{lg}}(0).
\end{equation*}
Here $C^{LS}$ is the constant of \eqref{5.4:eq-transport-constants-a}. The right-hand side
is $\mathsf{K}^{\mathrm{lg}}$, the level-$k$ long-unit bound of
\eqref{5.4:eq-transport-constants-c} with $\mathfrak{K}=C^{LS}$.

Two features of this route matter for later use.

No filing cell inside $Z$ is consulted. The only base input is the base bound for long
units \eqref{5.4:eq-age-base-bounds-b}. It is used at pairs $(j,\square')$ with
$\square'\subset\Delta_2$, and it sums over all $v$ with $\square'\in X_v$, whatever the
cell of $y_v$. The point $y_v$ enters only through $\operatorname{lev}\Delta_0(v)$ in the
proof of the per-cube transport, and there only as a geometric fact.

Suppose $\Delta_2\not\subset Z$. A cell of level $k$ is a $\pi_k$-cube. Such a cube is
either a cube of $Z$ or interior-disjoint from $Z$. Hence every sub-cube $\square'$ of
$\Delta_2$ lies off $Z$. It therefore suffices that \eqref{5.4:eq-age-base-bounds-b} hold at
cubes off $Z$. This is the off-$Z$ delivery under which
Proposition~\ref{5.4:prop-per-cube-transport} is stated for cells not contained in $Z$.
\end{proof}

\begin{corollary}[The long-unit bound at a general cell, level by level]
\label{5.4:cor-general-cell-bound}
In the setting of the per-cube transport (Proposition~\ref{5.4:prop-per-cube-transport}), the
following holds at every cell $\Delta_2$ of level $\ell$:
\begin{equation}\label{5.4:eq-general-cell-bound-1}
S_{r'}(\Delta_2;\mathfrak{L}^{\mathrm{lg}}) \le C^{LS}\,\mathsf{b}^{\mathrm{lg}}(k-\ell).
\end{equation}

Read level by level, with the grain $\pi^{\mathfrak{s}}_j$ at level $j$, the star
$\mathrm{St}^j_v$ of the cells of this grain meeting $X_v$, its rim $\partial\mathrm{St}^j_v$,
and the units of level $j$, the following holds for every level $j\le k$ and every cell
$\Delta_1$ of the grain:
\begin{equation}\label{5.4:eq-general-cell-bound-2}
e^{-\rho_j(\Delta_1)}
\sum_{\substack{v\in\mathfrak{L}^{\mathrm{lg}}:\ j_v=j,\\ \Delta_1\in\partial\mathrm{St}^j_v}}
w_{r'}(v)\,e^{2d_v}
\le 2dL^{d-1}\,C^{LS}_{\sharp}(c_{\mathrm{sp}})\,e^{-(k-j)}\,\mathsf{b}^{\mathrm{lg}}(k-j).
\end{equation}
The exponent $\rho_j(\Delta_1)$ is the cell exponent $\rho$ of the conditions
\eqref{5.4:eq-zones-cells-d}, evaluated at the cell $\Delta_1$ with the cells and their levels
taken in the grain $\pi^{\mathfrak{s}}_j$.
Here $c_{\mathrm{sp}}$ is the decay rate in the factors $e^{-c_{\mathrm{sp}}\mathfrak{m}_u}$ of the
bound of Proposition~\ref{5.4:prop-per-cube-transport}.
The factors on the right of \eqref{5.4:eq-general-cell-bound-2} are summed over the levels by
\begin{equation}\label{5.4:eq-general-cell-bound-3}
\sum_{j\le k} e^{-(k-j)}\,\mathsf{b}^{\mathrm{lg}}(k-j)
\le C_{\mathsf{b}}\,S_{\mathsf{b}}\,\mathsf{b}^{\mathrm{lg}}(0).
\end{equation}
\end{corollary}

\begin{proof}
The bound \eqref{5.4:eq-general-cell-bound-1} is the conclusion of the per-cube transport,
Proposition~\ref{5.4:prop-per-cube-transport}. The constant $C^{LS}$ there is the one fixed in
\eqref{5.4:eq-transport-constants-a}. The factor $\mathsf{b}^{\mathrm{lg}}(k-\ell)$ in
\eqref{5.4:eq-general-cell-bound-1} is not summed over the ages $k-\ell$ here; that summation
is carried out later in this chapter.

We prove \eqref{5.4:eq-general-cell-bound-2}. Fix a level $j\le k$ and a cell $\Delta_1$ of
the grain $\pi^{\mathfrak{s}}_j$. Let $v\in\mathfrak{L}^{\mathrm{lg}}$ with $j_v=j$ and
$\Delta_1\in\partial\mathrm{St}^j_v$. Since $\Delta_1$ lies in the rim of the star
$\mathrm{St}^j_v$, it shares a wall with a member $\Delta_2$ of that star, that is, with a
cell $\Delta_2$ meeting $X_v$. The cell $\Delta_1$ shares a wall with at most $2dL^{d-1}$
cells. The sum on the left of \eqref{5.4:eq-general-cell-bound-2} is therefore bounded by the
sum, over at most $2dL^{d-1}$ cells $\Delta_2$ walled to $\Delta_1$, of
\begin{equation*}
e^{-\rho_j(\Delta_1)}\sum_{v}\,w_{r'}(v)\,e^{2d_v},
\end{equation*}
where the inner sum runs over the units $v\in\mathfrak{L}^{\mathrm{lg}}$ with $j_v=j$ and
$X_v$ meeting $\Delta_2$. We bound each such term according to the level
$\operatorname{lev}_j\Delta_2$ of $\Delta_2$ in the grain, distinguishing the cases
$\operatorname{lev}_j\Delta_2=j$ and $\operatorname{lev}_j\Delta_2=j+u$ with $u\ge 1$.

First, let $\operatorname{lev}_j\Delta_2=j$. By the second of the standing geometric axioms
\eqref{5.4:eq-zones-cells-a}, taken together with Lemma~\ref{5.4:lem-cells-carrier}\,(b), for
every unit $v$ of the inner sum the cell $\Delta_2$, which meets $X_v$, is a cube of $X_v$. The
base bound for long units, \eqref{5.4:eq-age-base-bounds-b}, then bounds the inner sum by
$\mathsf{b}^{\mathrm{lg}}(k-j)$.

Second, let $\operatorname{lev}_j\Delta_2=j+u$ with $u\ge 1$. The proof of
Proposition~\ref{5.4:prop-per-cube-transport}, restricted to the units of level $j$, bounds the
inner sum by
\begin{equation*}
L^{du}\,e^{-c_{\mathrm{sp}}\mathfrak{m}_u}\,\mathsf{b}^{\mathrm{lg}}(k-j),
\end{equation*}
with the shell length $\mathfrak{m}_u$ of \eqref{5.4:eq-transport-constants-a}.

In both cases we write $\operatorname{lev}_j\Delta_2=j+u$, with $u=0$ in the first case. The
second of the conditions \eqref{5.4:eq-zones-cells-d}, read in the grain $\pi^{\mathfrak{s}}_j$,
gives
\begin{equation*}
e^{-\rho_j(\Delta_1)} \le e^{-(k-\operatorname{lev}_j\Delta_2)} = e^{-(k-j)}\,e^{u}.
\end{equation*}
In the first case, where $u=0$, the term is therefore at most
$e^{-(k-j)}\,\mathsf{b}^{\mathrm{lg}}(k-j)$. In the second case it is at most
\begin{equation*}
e^{-(k-j)}\,e^{u}L^{du}e^{-c_{\mathrm{sp}}\mathfrak{m}_u}\,\mathsf{b}^{\mathrm{lg}}(k-j).
\end{equation*}
In both cases the factor multiplying $e^{-(k-j)}\,\mathsf{b}^{\mathrm{lg}}(k-j)$, which is $1$
when $u=0$ and $e^{u}L^{du}e^{-c_{\mathrm{sp}}\mathfrak{m}_u}$ when $u\ge 1$, is bounded by
$C^{LS}_{\sharp}(c_{\mathrm{sp}})$ through the supremum bound of the shell series lemma
(Lemma~\ref{5.4:lem-shell-series}), applied with $c=c_{\mathrm{sp}}$: the shell length
$\mathfrak{m}_u$ of \eqref{5.4:eq-transport-constants-a} is the exponent
$(L^{u-1}-\delta)_+$ of that lemma, so that $e^{u}L^{du}e^{-c_{\mathrm{sp}}\mathfrak{m}_u}$ is a
term of the supremum which the lemma bounds by $C^{LS}_{\sharp}(c_{\mathrm{sp}})$. Each term is
therefore at most
\begin{equation*}
e^{-(k-j)}\,C^{LS}_{\sharp}(c_{\mathrm{sp}})\,\mathsf{b}^{\mathrm{lg}}(k-j).
\end{equation*}
Multiplying by the number $2dL^{d-1}$ of cells $\Delta_2$ gives
\eqref{5.4:eq-general-cell-bound-2}.

Finally, put $n=k-j$ in the sum of \eqref{5.4:eq-general-cell-bound-3}. The resulting sum is at
most $\sum_{n\ge 0}e^{-n}\mathsf{b}^{\mathrm{lg}}(n)$, which the growth bound for the age
function, \eqref{5.4:eq-age-base-bounds-a}, together with the constant $S_{\mathsf{b}}$ of
\eqref{5.4:eq-transport-constants-c}, bounds by
$C_{\mathsf{b}}S_{\mathsf{b}}\,\mathsf{b}^{\mathrm{lg}}(0)$. This is
\eqref{5.4:eq-general-cell-bound-3}.
\end{proof}

\begin{remark}[Remarks on the per-cube transport]\label{5.4:rem-per-cube-remarks}
Throughout, the notation and the hypotheses are those of the per-cube transport
(Proposition~\ref{5.4:prop-per-cube-transport}).
Here $v$ denotes a pointed unit with carrier $X_v$, $j=j_v$, cube $q_v$ and point $y_v$, and
$\Delta_0(v)$ is the cell containing $y_v$; further, $\Delta_2$ denotes a cell of level $\ell$, and
$r'\le \tfrac85\kappa$. The constants $\delta$, $\mathfrak{m}_u$, $\mathfrak{c}_u$,
$\mathfrak{A}_\delta(c)$ and $C^{LS}$ are those of \eqref{5.4:eq-transport-constants-a}.
\begin{enumerate}
\item[(a)] If $y_v$ is taken in the cell containing $q_v$, that is, $\Delta_0(v)$ is the cell
$\Delta'_v=\Delta(q_v)$ of Lemma~\ref{5.4:lem-cells-carrier}\textup{(c)}, then
$\delta=1$. If the inequality (g4) of \eqref{5.4:eq-zones-cells-c} holds only with its lower bound
weakened by one, then $\delta$ is replaced by $\delta+1$. In either case the shell series lemma
(Lemma~\ref{5.4:lem-shell-series}) still applies, and only the constant $\mathfrak{A}_\delta(c)$ changes.
\item[(b)] Suppose that, of the axioms in \eqref{5.4:eq-zones-cells-c}, only (g1) is assumed and (g4) is
not. Then the proof of Proposition~\ref{5.4:prop-per-cube-transport} runs with the lower bound, for every
unit $v$ counted in the shell of index $u$,
\begin{equation}\label{5.4:eq-per-cube-remarks-1}
d_v\ \ge\ \bigl(L^{u-1}/\mathsf{c}_d-2\bigr)_+ ,
\end{equation}
and the shell series lemma is applied with $c$ replaced by $c/\mathsf{c}_d$. Accordingly,
$\mathfrak{A}_\delta(c)$ is replaced by
\begin{equation}\label{5.4:eq-per-cube-remarks-2}
\Bigl(\frac{2q\,\mathsf{c}_d}{e\,c}\Bigr)^{q}\,e^{2c}\,\bigl(1-e^{-c/\mathsf{c}_d}\bigr)^{-1},
\end{equation}
with $q$ as in the definition of $\mathfrak{A}_\delta(c)$ in \eqref{5.4:eq-transport-constants-a}.
\item[(c)] Write $S^{\mathrm{touch}}_{r'}(\Delta_2;\mathfrak{L}^{\mathrm{lg}})$ for the sum
$S_{r'}(\Delta_2;\mathfrak{L}^{\mathrm{lg}})$ in which a unit $v$ is counted whenever $X_v$ touches
$\Delta_2$, that is, whenever $\bigcup X_v$ has a common point with $\Delta_2$ (the touch reading
of ``meet''). Then
\begin{equation}\label{5.4:eq-per-cube-remarks-3}
\begin{aligned}
S^{\mathrm{touch}}_{r'}(\Delta_2;\mathfrak{L}^{\mathrm{lg}})
&\le \hat{b}_t\,C^{LS}\,\mathsf{b}^{\mathrm{lg}}(k-\ell+1)\\
&\le \hat{b}_t\,C_{\mathsf{b}}\,2^{p_{\mathsf{b}}}\,C^{LS}\,\mathsf{b}^{\mathrm{lg}}(k-\ell).
\end{aligned}
\end{equation}
\item[(d)] Let $\mathfrak{a}'$ be a live arrest of level $k_{a'}$ with window $W_{\mathfrak{a}'}$. For
each fine cube of the shell of index $u$, the proof of Proposition~\ref{5.4:prop-per-cube-transport}
forms no factor other than
$L^{du}$, and this factor is set against $e^{-c_{\mathrm{sp}}\mathfrak{m}_u}$. In particular, the volume
factor $(L\check{R}_{k_{a'}+1})^d$ is never formed. Suppose the base bound
\eqref{5.4:eq-age-base-bounds-b} holds only with the cube weight, denoted $\Theta$ in this remark only,
\begin{equation}\label{5.4:eq-per-cube-remarks-4}
\Theta(\square')=
\begin{cases}
W^{\mathrm{win}}(\check{R}_{k_{a'}}) & \text{for } \square'\subset W_{\mathfrak{a}'},\\
1 & \text{otherwise}.
\end{cases}
\end{equation}
Then
\begin{equation}\label{5.4:eq-per-cube-remarks-5}
S_{r'}(\Delta_2;\mathfrak{L}^{\mathrm{lg}})
\le C^{LS}\,W^{\mathrm{win}}(\check{R}_{k_{a'}})\,\mathsf{b}^{\mathrm{lg}}(k-\ell)
\quad\text{if }\Delta_2\subset W_{\mathfrak{a}'} ,
\end{equation}
and the unweighted bound holds otherwise. The collar lemma gives, for the exponent
$\rho(\Delta_1)$ of the decay factor $e^{-\rho(\Delta_1)}$ of that lemma,
$\rho(\Delta_1)\ge2\check{R}_{k_{a'}}-1$ at every cell $\Delta_1$ walled to a cell
$\Delta_2\subset W_{\mathfrak{a}'}$. One separate factor $e^{-\rho(\Delta_1)}$ of the decay then absorbs
the window weight, since
\begin{equation*}
W^{\mathrm{win}}(\check{R}_{k_{a'}})\,e^{-(2\check{R}_{k_{a'}}-1)}\ \le\ C^{\mathrm{win}}_{\mathrm{lg}}
\end{equation*}
by \eqref{5.4:eq-transport-constants-d}. Let $\Delta^{\partial}$ be the cell of level $k$ denoted so in
Corollary~\ref{5.4:cor-current-level-bound}; it lies within distance $2\check{R}_k$ of $\partial Z$.
At $\Delta^{\partial}$
one has $\Theta\equiv1$. Hence the constant
$\mathsf{K}^{\mathrm{lg}}$ never carries a window weight.
\item[(e)] The base bound \eqref{5.4:eq-age-base-bounds-b} alone does not imply the conclusion of
Proposition~\ref{5.4:prop-per-cube-transport}; it is hypothesis \eqref{5.4:eq-pointed-units-a} that
excludes the configuration given in the proof below. Under it, the series
$\sum_{u\ge0}\mathfrak{c}_uL^{du}e^{-c_{\mathrm{sp}}\mathfrak{m}_u}$ converges for every
$c_{\mathrm{sp}}>0$. Suppose instead that short and long units were separated by a length threshold
$d_v\ge d_*$, for a number $d_*$. Then the units above the threshold that are short in the geometric
sense of \eqref{5.4:eq-pointed-units-a} are treated with the short units, by the short-unit sum
lemma and its first corollary. The per-cube transport serves those that are long in that geometric
sense.
\item[(f)] In the direction opposite to Proposition~\ref{5.4:prop-per-cube-transport}, suppose that at
every cell $\Delta'$, in $Z$ or not, the long units $v$ whose point $y_v$ lies in $\Delta'$ obey the
per-cell long column, that is, the bound on their sum whose constant is
$\mathsf{K}^{\flat}_{\mathrm{lg}}$ (in what follows, a per-cell bound). Then these bounds give back
\eqref{5.4:eq-age-base-bounds-b} through the same distance estimates. The availability of such per-cell
bounds at cells in $Z$ is a condition on the way \eqref{5.4:eq-age-base-bounds-b} is obtained, if it is
obtained by summing per-cell bounds. It is not a condition on the per-cube transport.
\end{enumerate}
\end{remark}

\begin{proof}
(a) The shell series lemma is stated for every $c>0$ and every $\delta\ge0$. Replacing $\delta$ by
$\delta+1$, or fixing $\delta=1$, therefore leaves its hypotheses intact. The only change is in the
constant $\mathfrak{A}_\delta(c)$ on its right side.

(b) In the step of the proof of Proposition~\ref{5.4:prop-per-cube-transport} that treats a unit $v$
counted in the shell of index $u$, a cube $Q_x$ of $X_v$ and the cube $q_v$ are at set distance at least
$R-2$; here $R$ is the distance parameter of that step for the shell of index $u$, not the ball radius
$R$ of the observables. Let $D$ be
the set distance of two cubes of the grid. For their lower corners, $D$ forces some coordinate to change
by $D+1$. A wall-chain changes one corner coordinate by one at each step. Hence any wall-chain
joining $Q_x$ to $q_v$ has at least $D+1\ge R-1$ steps, and so more than $R-1$ cubes. Since $X_v$ is
wall-connected, such a chain exists inside $X_v$; its cubes are cubes of $X_v$, so the number
$N_j(X_v)$ of $\pi_j$-cubes of $X_v$ satisfies $N_j(X_v)\ge R-1$, and the cube-count axiom (g1) of
\eqref{5.4:eq-zones-cells-c} gives
\begin{equation*}
\mathsf{c}_d(d_v+1)\ \ge\ R-1 .
\end{equation*}
With $R$ as fixed in that step for the shell of index $u$, this yields the lower bound
\eqref{5.4:eq-per-cube-remarks-1}, which is used in the proof in
place of the one coming from axiom (g4) of \eqref{5.4:eq-zones-cells-c}. The weight $e^{-c\,d_v}$ is then
controlled with exponent
$c/\mathsf{c}_d$ in place of $c$, so the shell series lemma runs with $c\mapsto c/\mathsf{c}_d$. The
constant $\mathfrak{A}_\delta(c)$ is replaced by \eqref{5.4:eq-per-cube-remarks-2}.

(c) Let $v$ be counted in $S^{\mathrm{touch}}_{r'}(\Delta_2;\mathfrak{L}^{\mathrm{lg}})$, and let $Q$ be a
cube of $X_v$ touching $\Delta_2$. By part (c) of the lemma on cells against a carrier
(Lemma~\ref{5.4:lem-cells-carrier}), $Q$ lies in exactly one cell $\Delta(Q)$, and $X_v$ meets
$\Delta(Q)$. Since $Q\subset\Delta(Q)$ and $Q$ touches $\Delta_2$, the cell $\Delta(Q)$ either equals
$\Delta_2$ or touches it. Cells which touch differ in level by at most one, by hypothesis (A2) of
Proposition~\ref{5.4:prop-per-cube-transport}, so $\Delta(Q)$ has level $\ell-1$, $\ell$ or $\ell+1$.
Thus $v$ is counted, in the sense of
meeting, in $S_{r'}(\Delta;\mathfrak{L}^{\mathrm{lg}})$ for some cell $\Delta$ equal to or touching
$\Delta_2$. There are at most $\hat{b}_t$ such cells.

We apply Proposition~\ref{5.4:prop-per-cube-transport} at each such cell. This bounds its sum by
$C^{LS}\mathsf{b}^{\mathrm{lg}}(k-\operatorname{lev}\Delta)$. Since $\operatorname{lev}\Delta\ge\ell-1$,
the largest argument that occurs is $k-\ell+1$; since the age function $\mathsf{b}^{\mathrm{lg}}$ is
non-decreasing in its argument, as fixed in \eqref{5.4:eq-age-base-bounds-b}, each term is bounded by
$C^{LS}\mathsf{b}^{\mathrm{lg}}(k-\ell+1)$. Summing over the at most $\hat{b}_t$ cells gives the first
inequality of \eqref{5.4:eq-per-cube-remarks-3}. Increasing the argument of the age function by one
costs at most the factor $C_{\mathsf{b}}2^{p_{\mathsf{b}}}$, by the growth constants of the age function
fixed in \eqref{5.4:eq-age-base-bounds-b}. This gives the
second inequality.

(d) The proof of Proposition~\ref{5.4:prop-per-cube-transport} reads the base bound
\eqref{5.4:eq-age-base-bounds-b} only at pairs $(j,\square')$ with $\square'\subset\Delta_2$. By
the collar-zone clause of \eqref{5.4:eq-zones-cells-g}, the window $W_{\mathfrak{a}'}$ is a union of
cells. Hence a cube inside
$\Delta_2$ lies in $W_{\mathfrak{a}'}$ only when $\Delta_2\subset W_{\mathfrak{a}'}$. In that case each
use of the base bound carries at most the factor $W^{\mathrm{win}}(\check{R}_{k_{a'}})$, which gives
\eqref{5.4:eq-per-cube-remarks-5}. Otherwise $\Theta=1$ at every cube read, and the proposition applies
unchanged.

The cell $\Delta^{\partial}$ lies within distance $2\check{R}_k$ of $\partial Z$. It is not raised in
level, by \eqref{5.4:eq-zones-cells-f}, and it lies in no window. Hence $\Theta\equiv1$ on its cubes, and
the constant $\mathsf{K}^{\mathrm{lg}}=C^{LS}\mathsf{b}^{\mathrm{lg}}(0)$ of
Corollary~\ref{5.4:cor-current-level-bound} carries no window weight.

(e) For every $m<\ell$, take a unit $v_m$ with $d_{v_m}=0$ whose carrier $X_{v_m}$ consists of a single
$\pi_m$-cube contained in $\Delta_2$. This configuration obeys \eqref{5.4:eq-age-base-bounds-b} with
$\mathsf{b}^{\mathrm{lg}}$ constant. The sum
$S_{r'}(\Delta_2;\{v_m\}_{m<\ell})$ taken over this family alone is at least $\ell-1$. So a per-cube bound
of the form
\eqref{5.4:eq-age-base-bounds-b} by itself gives no bound on this sum uniform in $\ell$. The units
$v_m$ are short in the sense of \eqref{5.4:eq-pointed-units-a}, their carriers lying in $\Delta_2$ or
inside $\mathcal{N}(\Delta_0(v_m))$, the open neighbourhood of radius one of $\Delta_0(v_m)$ in the
units of $\Delta_0(v_m)$.
Hypothesis \eqref{5.4:eq-pointed-units-a} therefore excludes them from $\mathfrak{L}^{\mathrm{lg}}$.

Under \eqref{5.4:eq-pointed-units-a}, every unit counted in the shell of index $u\ge2$ has
$d_v\ge L^{u-1}-\delta$. By the shell series lemma, applied with $c=c_{\mathrm{sp}}$, the series of
Proposition~\ref{5.4:prop-per-cube-transport} then converges for every $c_{\mathrm{sp}}>0$.

Now consider a length threshold. The units with $d_v\ge d_*$ whose carrier satisfies the geometric
criterion of shortness in \eqref{5.4:eq-pointed-units-a} are short units, whatever their length. They are
therefore
handled with the short units, by the short-unit sum lemma and its first corollary. The per-cube transport
is applied only to those whose carrier does not satisfy that criterion.

(f) The per-cube transport passes from the base bound at the cubes $\square'\subset\Delta_2$ to the
sum over the units meeting $\Delta_2$, using the distances between the cell $\Delta_0(v)$ and these
cubes. The same distances relate a per-cell bound at $\Delta_0(v)$ to the base bound at any cube
$\square'$ of $X_v$, in the opposite direction. A per-cell bound at every cell $\Delta'$, in $Z$ or not,
therefore gives back \eqref{5.4:eq-age-base-bounds-b}. Corollary~\ref{5.4:cor-current-level-bound},
like the per-cube transport, reads \eqref{5.4:eq-age-base-bounds-b} only at cubes inside $\Delta_2$ and
uses no per-cell bound inside $Z$. The availability of such per-cell bounds at cells inside $Z$
therefore concerns only the way \eqref{5.4:eq-age-base-bounds-b} is obtained, and not the transport
itself.
\end{proof}

\begin{proposition}[The per-cell transport]\label{5.4:prop-per-cell-transport}
Let $c\ge 2$, and assume the following.
\begin{itemize}
\item The per-cell filed column bound for long units and tails, \eqref{5.4:eq-age-base-bounds-c},
with this constant $c$ and with the constant $\mathsf{K}^{\flat}_{\mathrm{lg}}$.
\item The carrier $X_v$ of every pointed unit $v$ is the set fixed in \eqref{5.4:eq-pointed-units-b}
of the definition of pointed units, short and long units, and the met-cell sums.
\item The conclusion of Lemma~\ref{5.4:lem-carrier-in-zone}: a unit's carrier stays inside its
own zone set.
\item The standing axioms (S1)--(S3) of \eqref{5.4:eq-zones-cells-a}, the axioms (A1)--(A3) of
\eqref{5.4:eq-zones-cells-b} for the current zone sets, and condition (g1) of
\eqref{5.4:eq-zones-cells-c}.
\end{itemize}
Then for every level $\ell\le k$, every cell $\Delta_2$ of level $\ell$ and every
$r'\le \tfrac{8}{5}\kappa$,
\begin{equation}\label{5.4:eq-per-cell-transport-1}
S_{r'}(\Delta_2;\mathfrak{L}^{\mathrm{lg}})
\le C^{\mathrm{gen}}_{\mathsf{b}}\,\mathsf{b}^{\mathrm{lg}}(k-\ell)\,\mathsf{K}^{\flat}_{\mathrm{lg}} ,
\end{equation}
where $S_{r'}(\Delta_2;\mathfrak{L}^{\mathrm{lg}})$ is the met-cell sum of
\eqref{5.4:eq-pointed-units-c} over the family $\mathfrak{L}^{\mathrm{lg}}$ of long units and
tails, $\mathsf{b}^{\mathrm{lg}}$ is the age function of Definition~\ref{5.4:def-age-base-bounds},
and $C^{\mathrm{gen}}_{\mathsf{b}}$ is the constant fixed in \eqref{5.4:eq-transport-constants-b}.
The constant $C^{\mathrm{gen}}_{\mathsf{b}}$ depends only on $d$, $L$, $\mathsf{c}_d$,
$C_{\mathsf{b}}$ and $p_{\mathsf{b}}$; in particular it depends neither on the large-field region
$Z$, nor on the radius $\check{R}_k$, nor on any foot or window of a unit. The level $k$ enters
\eqref{5.4:eq-per-cell-transport-1}
only through the range of levels at which axiom (A3) is applied and through the argument of
$\mathsf{b}^{\mathrm{lg}}$.
\end{proposition}

\begin{proof}
Fix $\ell\le k$, a cell $\Delta_2$ of level $\ell$, and $r'\le\frac85\kappa$. Throughout we use
coordinates of level $\ell$. In these coordinates $\Delta_2=[0,1]^d$, and the cubes of the grid
$\pi_\ell$ are the unit cubes $\prod_{i=1}^d[m_i,m_i+1]$ with integers $m_i$.
For a level $\ell'$ we write $\operatorname{dist}_{\ell'}$ and $\operatorname{diam}_{\infty,\ell'}$
for the
sup-norm distance and diameter measured in units of level $\ell'$, so that
$\operatorname{dist}_\ell=L^{t}\operatorname{dist}_{\ell+t}$ for every integer $t$, and we write
$\operatorname{lev}\Delta$ for the level of a cell $\Delta$. Put
\begin{equation*}
\mathcal{V}:=\{v\in\mathfrak{L}^{\mathrm{lg}}:\ X_v \text{ meets } \Delta_2\};
\end{equation*}
these are the units over which the met-cell sum $S_{r'}(\Delta_2;\mathfrak{L}^{\mathrm{lg}})$
runs.

Let $v\in\mathcal{V}$ and let $j_v$ denote its level. Since $X_v$ meets the cell $\Delta_2$ of level
$\ell$, part~(e) of the lemma on cells against a carrier (Lemma~\ref{5.4:lem-cells-carrier}) gives
$j_v\le\ell$ and a cube $q_2(v)$ of $X_v$ with $q_2(v)\subset\Delta_2$. By part~(c) of the same
lemma the cube $q_v$ of $X_v$ lies in the filing cell $\Delta'_v=\Delta(q_v)$, and
$\ell'_v:=\operatorname{lev}\Delta'_v$ satisfies $j_v\le\ell'_v\le k$. Set
\begin{equation*}
\delta_v:=\operatorname{dist}_\ell(\Delta'_v,\Delta_2),\qquad
s_v:=\lfloor\delta_v\rfloor\in\mathbb{Z}_{\ge 0}.
\end{equation*}

Step~1 (the exit length dominates the separation). Since $q_v\subset\Delta'_v$ and
$q_2(v)\subset\Delta_2$, we have $\operatorname{dist}_\ell(q_v,q_2(v))\ge\delta_v$. Both $q_v$
and $q_2(v)$ are cubes of $X_v$, so they lie in $\bigcup X_v$, and therefore
$\operatorname{diam}_{\infty,\ell}(\bigcup X_v)\ge\delta_v$. Part~(d) of
Lemma~\ref{5.4:lem-cells-carrier}, applied at level $\ell$, bounds this diameter by
$L^{j_v-\ell}\mathfrak{d}_v$, where $\mathfrak{d}_v=\mathsf{c}_d(d_v+1)$. Hence, using $j_v\le\ell$
and $\delta_v\ge\lfloor\delta_v\rfloor$,
\begin{equation}\label{5.4:eq-per-cell-transport-2}
\mathfrak{d}_v\ \ge\ L^{\ell-j_v}\delta_v\ \ge\ s_v .
\end{equation}

Step~2a (finer filing cells are far in exit length). Suppose $\ell'_v=\ell-u$ with $u\ge2$. The
cell $\Delta'_v$, being a cell of level $\ell-u$, lies in the closure of the zone of level
$\ell-u$. That closure lies in $U_{\ell-u+1}$, and $U_{\ell-u+1}\subset U_{\ell-1}$ because
$\ell-u+1\le\ell-1$. Further, $\Delta_2\subset\Omega_\ell$, since $\Delta_2$ is a cell of level
$\ell$. Axiom (A3), applied at level $\ell-1<k$ to $U_{\ell-1}\supset\Delta'_v$ and
$\Omega_\ell\supset\Delta_2$, gives $\delta_v\ge1$. Together with $\delta_v\ge s_v$ this yields
$\delta_v\ge\max\{1,s_v\}$. Moreover $j_v\le\ell'_v=\ell-u$, so $L^{\ell-j_v}\ge L^u$, and the
first inequality of \eqref{5.4:eq-per-cell-transport-2} gives
\begin{equation}\label{5.4:eq-per-cell-transport-3}
\mathfrak{d}_v\ \ge\ L^{u}\max\{1,s_v\}\qquad(\ell'_v=\ell-u,\ u\ge2).
\end{equation}

Step~2b (coarser filing cells are far). Suppose $\ell'_v=\ell+t$ with $t\ge2$; recall
$\ell+t\le k$. The cell $\Delta_2$ lies in $U_{\ell+1}$, and $U_{\ell+1}\subset U_{\ell+t-1}$,
while $\Delta'_v$ lies in the closure of $\Omega_{\ell+t}$. Axiom (A3), applied at level
$\ell+t-1<k$, gives $\operatorname{dist}_{\ell+t}(\Delta'_v,\Delta_2)\ge1$, that is,
$\delta_v\ge L^t$. Since $L^t$ is a positive integer,
\begin{equation}\label{5.4:eq-per-cell-transport-4}
s_v=\lfloor\delta_v\rfloor\ \ge\ L^{t}\qquad(\ell'_v=\ell+t,\ t\ge2).
\end{equation}

Step~3 (counting cells). For a level $\ell'$ and an integer $s\ge0$ let $n_{\ell'}(s)$ be the
number of cells $\Delta'$ of level $\ell'$ with
$\lfloor\operatorname{dist}_\ell(\Delta',\Delta_2)\rfloor=s$.

Level $\ell$. A cell of level $\ell$ is $\Delta'=\prod_i[m_i,m_i+1]$, and its sup-norm distance
to $[0,1]^d$ is $\max_i(|m_i|-1)_+$, an integer. If this equals $s$, then $|m_i|\le s+1$ for every
$i$, which leaves $2s+3$ values for each $m_i$. Hence $n_\ell(s)\le(2s+3)^d$. The same argument
shows that at most $(2s+3)^d$ cubes of $\pi_\ell$ lie at distance at most $s$ from $\Delta_2$.

Level $\ell-u$, $u\ge1$. By (S1), a cell $\Delta'$ of level $\ell-u$ has a unique ancestor in
$\pi_\ell$, a cube containing $\Delta'$. The ancestor's distance to $\Delta_2$ is at most
$\operatorname{dist}_\ell(\Delta',\Delta_2)<s+1$, and it is an integer, hence at most $s$. So the
ancestor is one of at most $(2s+3)^d$ cubes. Each such cube holds $L^{du}$ cubes of
$\pi_{\ell-u}$, and the map sending $\Delta'$ to its ancestor together with its position inside
the ancestor is injective. Hence $n_{\ell-u}(s)\le(2s+3)^dL^{du}$.

Level $\ell+t$, $t\ge1$. Let $\Delta'$ be such a cell, pick $x\in\Delta'$ nearest to $\Delta_2$,
and, by (S1), a cube $q$ of $\pi_\ell$ with $x\in q\subset\Delta'$. Then
\begin{equation*}
\operatorname{dist}_\ell(q,\Delta_2)\le\operatorname{dist}_\ell(x,\Delta_2)
=\operatorname{dist}_\ell(\Delta',\Delta_2)<s+1 .
\end{equation*}
The left side is an integer, so it is at most $s$, and $q$ is one of at most $(2s+3)^d$ cubes.
Distinct cells have disjoint interiors and the interior of $q$ lies in the interior of $\Delta'$,
so $\Delta'\mapsto q$ is injective. Hence $n_{\ell+t}(s)\le(2s+3)^d$.

Below, for the levels other than $\ell$, we also use the weaker bounds obtained from
$(2s+3)^d\le(2s+5)^d$.

Step~4 (the column enters). For a cell $\Delta'$ of level at most $k$ put
$s:=\lfloor\operatorname{dist}_\ell(\Delta',\Delta_2)\rfloor$ and
\begin{equation*}
m(\Delta'):=
\begin{cases}
L^u\max\{1,s\}, & \operatorname{lev}\Delta'=\ell-u \text{ with } u\ge2,\\
s, & \text{otherwise}.
\end{cases}
\end{equation*}
By Step~1 and Step~2a, $\mathfrak{d}_v\ge m(\Delta'_v)$ for every $v\in\mathcal{V}$. The constant
$\theta$ of \eqref{5.4:eq-transport-constants-b} satisfies $\theta\mathfrak{d}_v=2d_v+2$ exactly,
that is, $\theta=2/\mathsf{c}_d$. Hence
\begin{equation*}
e^{-2d_v}=e^{2}e^{-\theta\mathfrak{d}_v}\le e^{2}e^{-\theta m(\Delta'_v)} .
\end{equation*}
Since $(2+c)-2-(c-2)=2$, we have $e^{2d_v}=e^{(2+c)d_v}e^{-2d_v}e^{-(c-2)d_v}$. Therefore
\begin{equation*}
e^{2d_v}\le e^{2}e^{-\theta m(\Delta'_v)}\,e^{(2+c)d_v}\,\mathfrak{r}_v,
\qquad \mathfrak{r}_v:=e^{-(c-2)d_v}\le1 ,
\end{equation*}
where $\mathfrak{r}_v\le1$ because $c\ge2$.

We use this bound for each $v\in\mathcal{V}$ in $S_{r'}(\Delta_2;\mathfrak{L}^{\mathrm{lg}})$ and
drop $\mathfrak{r}_v\le1$. Grouping the units of $\mathcal{V}$ by their filing cell $\Delta'_v$ and
applying the column bound \eqref{5.4:eq-age-base-bounds-c} (with $c\ge2$) to each group, the group
filed at a cell $\Delta'$ contributes at most
$e^{2}e^{-\theta m(\Delta')}\,\mathsf{K}^{\flat}_{\mathrm{lg}}\,
\mathsf{b}^{\mathrm{lg}}(k-\operatorname{lev}\Delta')$. This
gives
\begin{equation}\label{5.4:eq-per-cell-transport-5}
S_{r'}(\Delta_2;\mathfrak{L}^{\mathrm{lg}})
\le e^{2}\mathsf{K}^{\flat}_{\mathrm{lg}}\sum_{\Delta'}
\mathsf{b}^{\mathrm{lg}}(k-\operatorname{lev}\Delta')\,e^{-\theta m(\Delta')} ,
\end{equation}
the sum running over the filing cells. By Step~2a, the filing cells of level $\ell-u$ with $u\ge2$
have $s\ge1$; by Step~2b, those of level $\ell+t$ with $t\ge2$ have $s\ge L^t$. We extend the sum
in \eqref{5.4:eq-per-cell-transport-5} to all cells of level at most $k$ obeying these two
constraints; this is licit since all terms are non-negative.

Step~5 (age, cell by cell). If $\operatorname{lev}\Delta'\ge\ell$, then
$k-\operatorname{lev}\Delta'\le k-\ell$, and since the age function $\mathsf{b}^{\mathrm{lg}}$ of
Definition~\ref{5.4:def-age-base-bounds} is non-decreasing,
$\mathsf{b}^{\mathrm{lg}}(k-\operatorname{lev}\Delta')\le\mathsf{b}^{\mathrm{lg}}(k-\ell)$. If
$\operatorname{lev}\Delta'=\ell-u$ with $u\ge1$, the growth bound
\eqref{5.4:eq-age-base-bounds-a} gives
\begin{equation*}
\mathsf{b}^{\mathrm{lg}}(k-\operatorname{lev}\Delta')\le C_{\mathsf{b}}(1+u)^{p_{\mathsf{b}}}
\mathsf{b}^{\mathrm{lg}}(k-\ell) ;
\end{equation*}
for $u=1$ this reads
$C_{\mathsf{b}}2^{p_{\mathsf{b}}}\mathsf{b}^{\mathrm{lg}}(k-\ell)$. Sorting the cells by level, and
using $m(\Delta')=s$ at the levels $\ell$, $\ell\pm1$ and $\ell+t$, we obtain
\begin{equation}\label{5.4:eq-per-cell-transport-6}
\sum_{\Delta'}\mathsf{b}^{\mathrm{lg}}(k-\operatorname{lev}\Delta')e^{-\theta m(\Delta')}
\le\mathsf{b}^{\mathrm{lg}}(k-\ell)\bigl[\Sigma_\ell+\Sigma_{\ell+1}+\Sigma_{\ell+}
+C_{\mathsf{b}}2^{p_{\mathsf{b}}}\Sigma_{\ell-1}+C_{\mathsf{b}}\Sigma_{\ell-}\bigr],
\end{equation}
where
\begin{align*}
\Sigma_\ell&:=\sum_{s\ge0}n_\ell(s)e^{-\theta s},\qquad
\Sigma_{\ell+1}:=\sum_{s\ge0}n_{\ell+1}(s)e^{-\theta s},\qquad
\Sigma_{\ell-1}:=\sum_{s\ge0}n_{\ell-1}(s)e^{-\theta s},\\
\Sigma_{\ell+}&:=\sum_{t\ge2}\ \sum_{s\ge L^t}n_{\ell+t}(s)e^{-\theta s},\qquad
\Sigma_{\ell-}:=\sum_{u\ge2}(1+u)^{p_{\mathsf{b}}}\sum_{s\ge0}n_{\ell-u}(s)
e^{-\theta L^u\max\{1,s\}} .
\end{align*}

Step~6 (the series). We use two elementary inequalities. For $y\ge0$, $\alpha>0$ and
$q'\ge0$, the maximum of $y^{q'}e^{-\alpha y}$ over $y\ge0$ is attained at $y=q'/\alpha$, whence
\begin{equation}\label{5.4:eq-per-cell-transport-7}
y^{q'}e^{-\alpha y}\le\Bigl(\frac{q'}{e\alpha}\Bigr)^{q'} .
\end{equation}
Since the blocking factor $L$ is an odd integer larger than one, $L\ge3$. Hence, for every integer
$u\ge1$,
\begin{equation}\label{5.4:eq-per-cell-transport-8}
L^u\ge1+u,\qquad L^u\ge3u .
\end{equation}
Indeed, $L^u\ge3^u\ge1+u$ and $3^u\ge3u$.

The basic series. Put $\Xi_1:=\sum_{s\ge0}(2s+5)^de^{-\theta s}$. With $y:=2s+5$ we have
$e^{-\theta s}=e^{5\theta/2}e^{-\theta y/2}$. By \eqref{5.4:eq-per-cell-transport-7} with
$q'=d$, $\alpha=\theta/4$, and by $4/\theta=2\mathsf{c}_d$,
\begin{align*}
(2s+5)^de^{-\theta s}
&=e^{5\theta/2}\bigl[y^de^{-\theta y/4}\bigr]e^{-\theta y/4}
\le e^{5\theta/2}\Bigl(\frac{4d}{e\theta}\Bigr)^de^{-\theta(2s+5)/4}\\
&=e^{5\theta/4}\Bigl(\frac{2d\mathsf{c}_d}{e}\Bigr)^de^{-\theta s/2} .
\end{align*}
Further,
\begin{equation*}
\sum_{s\ge0}e^{-\theta s/2}=(1-e^{-\theta/2})^{-1}\le(1-e^{-\theta/4})^{-1} .
\end{equation*}
Hence
\begin{equation*}
\Xi_1\le e^{5\theta/4}\Bigl(\frac{2d\mathsf{c}_d}{e}\Bigr)^d(1-e^{-\theta/4})^{-1}\le\mathfrak{a}_d ,
\end{equation*}
the constant $\mathfrak{a}_d$ of \eqref{5.4:eq-transport-constants-b}. Also
\begin{equation*}
\Xi_0:=\sum_{s\ge0}(2s+3)^de^{-\theta s}\le \Xi_1\le\mathfrak{a}_d .
\end{equation*}

Next put $\Xi_1':=\sum_{s\ge0}(2s+5)^de^{-\theta s/2}$. With $y:=2s+5$ we have
$e^{-\theta s/2}=e^{5\theta/4}e^{-\theta y/4}$. By \eqref{5.4:eq-per-cell-transport-7} with
$q'=d$, $\alpha=\theta/8$,
\begin{align*}
(2s+5)^de^{-\theta s/2}
&=e^{5\theta/4}\bigl[y^de^{-\theta y/8}\bigr]e^{-\theta y/8}
\le e^{5\theta/4}\Bigl(\frac{8d}{e\theta}\Bigr)^de^{-\theta(2s+5)/8}\\
&=e^{5\theta/8}\Bigl(\frac{8d}{e\theta}\Bigr)^de^{-\theta s/4} .
\end{align*}
Summing, with $\sum_{s\ge0}e^{-\theta s/4}=(1-e^{-\theta/4})^{-1}$, and using
$8d/(e\theta)=2\cdot 2d\mathsf{c}_d/e$ and $e^{5\theta/8}\le e^{5\theta/4}$, we get
\begin{equation*}
\Xi_1'\le e^{5\theta/8}\Bigl(\frac{8d}{e\theta}\Bigr)^d(1-e^{-\theta/4})^{-1}\le2^d\mathfrak{a}_d .
\end{equation*}

The levels $\ell$, $\ell\pm1$. By Step~3, $n_\ell(s)\le(2s+3)^d$ and
$n_{\ell+1}(s)\le(2s+3)^d$. Hence
\begin{equation*}
\Sigma_\ell\le \Xi_0\le\mathfrak{a}_d,\qquad\Sigma_{\ell+1}\le \Xi_0\le\mathfrak{a}_d .
\end{equation*}
By Step~3, $n_{\ell-1}(s)\le(2s+3)^dL^d$. Hence $\Sigma_{\ell-1}\le L^d\Xi_0\le L^d\mathfrak{a}_d$.

The levels $\ell+t$, $t\ge2$. For $s\ge L^t$ we have
$e^{-\theta s}=e^{-\theta s/2}e^{-\theta s/2}\le e^{-\theta L^t/2}e^{-\theta s/2}$. With
$n_{\ell+t}(s)\le(2s+3)^d\le(2s+5)^d$, the slice of $\Sigma_{\ell+}$ with fixed $t$ is at most
$e^{-\theta L^t/2}\Xi_1'\le e^{-\theta L^t/2}2^d\mathfrak{a}_d$. By the second inequality of
\eqref{5.4:eq-per-cell-transport-8},
\begin{equation*}
\sum_{t\ge2}e^{-\theta L^t/2}\le\sum_{t\ge2}e^{-3\theta t/2}=\mathfrak{e}_d ,
\end{equation*}
the constant $\mathfrak{e}_d$ of \eqref{5.4:eq-transport-constants-b}. Hence
$\Sigma_{\ell+}\le2^d\mathfrak{e}_d\mathfrak{a}_d$.

The levels $\ell-u$, $u\ge2$. First, $L^u\max\{1,s\}\ge L^u+s-1$. For $s\ge1$ this is
$(L^u-1)(s-1)\ge0$; for $s=0$ it reads $L^u\ge L^u-1$. Hence
$e^{-\theta L^u\max\{1,s\}}\le e^{\theta}e^{-\theta L^u}e^{-\theta s}$. With
$n_{\ell-u}(s)\le(2s+3)^dL^{du}$ from Step~3,
\begin{equation*}
\sum_{s\ge0}n_{\ell-u}(s)e^{-\theta L^u\max\{1,s\}}
\le e^{\theta}L^{du}e^{-\theta L^u}\Xi_0\le e^{\theta}L^{du}e^{-\theta L^u}\Xi_1 .
\end{equation*}
By the first inequality of
\eqref{5.4:eq-per-cell-transport-8}, $(1+u)^{p_{\mathsf{b}}}\le L^{p_{\mathsf{b}}u}$. Hence, by
\eqref{5.4:eq-per-cell-transport-7} with $q'=d+p_{\mathsf{b}}$, $\alpha=\theta/2$, and by
$2/\theta=\mathsf{c}_d$,
\begin{align*}
(1+u)^{p_{\mathsf{b}}}L^{du}e^{-\theta L^u}
&\le L^{(d+p_{\mathsf{b}})u}e^{-\theta L^u}
=\bigl[y^{d+p_{\mathsf{b}}}e^{-\theta y/2}\bigr]e^{-\theta y/2}\Big|_{y=L^u}\\
&\le\Bigl(\frac{(d+p_{\mathsf{b}})\mathsf{c}_d}{e}\Bigr)^{d+p_{\mathsf{b}}}e^{-\theta L^u/2} .
\end{align*}
By the second inequality of \eqref{5.4:eq-per-cell-transport-8},
\begin{equation*}
\sum_{u\ge2}e^{-\theta L^u/2}\le\sum_{u\ge2}e^{-3\theta u/2}=\mathfrak{e}_d .
\end{equation*}
Collecting, and
using $\Xi_1\le\mathfrak{a}_d$,
\begin{equation*}
\Sigma_{\ell-}\le e^{\theta}\Bigl(\frac{(d+p_{\mathsf{b}})\mathsf{c}_d}{e}\Bigr)^{d+p_{\mathsf{b}}}
\mathfrak{e}_d\mathfrak{a}_d .
\end{equation*}

Step~7 (assembly). Insert the five bounds of Step~6 into \eqref{5.4:eq-per-cell-transport-6}, and
the result into \eqref{5.4:eq-per-cell-transport-5}:
\begin{align*}
S_{r'}(\Delta_2;\mathfrak{L}^{\mathrm{lg}})
&\le e^{2}\mathsf{K}^{\flat}_{\mathrm{lg}}\,\mathsf{b}^{\mathrm{lg}}(k-\ell)\,\mathfrak{a}_d
\Bigl[2+2^d\mathfrak{e}_d+C_{\mathsf{b}}2^{p_{\mathsf{b}}}L^d\\
&\qquad\qquad
+C_{\mathsf{b}}e^{\theta}\Bigl(\frac{(d+p_{\mathsf{b}})\mathsf{c}_d}{e}\Bigr)^{d+p_{\mathsf{b}}}
\mathfrak{e}_d\Bigr]\\
&=C^{\mathrm{gen}}_{\mathsf{b}}\,\mathsf{b}^{\mathrm{lg}}(k-\ell)\,\mathsf{K}^{\flat}_{\mathrm{lg}} .
\end{align*}
In the last step we used that $e^{2}\mathfrak{a}_d$ times the bracket is the constant
$C^{\mathrm{gen}}_{\mathsf{b}}$ of \eqref{5.4:eq-transport-constants-b}. This is
\eqref{5.4:eq-per-cell-transport-1}.

The bracket and the factor $e^2\mathfrak{a}_d$ are built from $d$, $L$, $C_{\mathsf{b}}$,
$p_{\mathsf{b}}$, $\theta=2/\mathsf{c}_d$ and $\mathsf{c}_d$ alone. So
$C^{\mathrm{gen}}_{\mathsf{b}}$ depends only on $d$, $L$, $\mathsf{c}_d$, $C_{\mathsf{b}}$ and
$p_{\mathsf{b}}$. The level $k$ was used only in applying axiom (A3) at the levels $\ell-1<k$
(Step~2a) and $\ell+t-1<k$ (Step~2b), and through the argument of $\mathsf{b}^{\mathrm{lg}}$.
\end{proof}

\begin{remark}[Variant under (g4)]
The argument above has a variant under condition (g4) of
\eqref{5.4:eq-zones-cells-c}. In that variant, Step~1 is run through the inequality
\begin{equation*}
d_v\ \ge\ \operatorname{dist}_{j_v}(q_v,q_2(v))\ \ge\ L^{\ell-j_v}\delta_v .
\end{equation*}
Together with the bounds on $\delta_v$ of Step~2a, this gives $d_v\ge m(\Delta'_v)$ directly.
Step~4 then runs at $\theta:=2$, with $e^{-2d_v}\le e^{-2m(\Delta'_v)}$; the prefactor $e^{2}$ is
absent, and in Step~6 the number $2/\theta$ equals $1$ where it equalled $\mathsf{c}_d$ before.
With
\begin{equation*}
\mathfrak{a}^{\natural}_d:=e^{5/2}\Bigl(\frac{2d}{e}\Bigr)^d(1-e^{-1/2})^{-1},\qquad
\mathfrak{e}^{\natural}:=e^{-6}(1-e^{-3})^{-1},
\end{equation*}
one obtains $S_{r'}(\Delta_2;\mathfrak{L}^{\mathrm{lg}})\le C^{\mathrm{gen},\natural}_{\mathsf{b}}\,
\mathsf{b}^{\mathrm{lg}}(k-\ell)\,\mathsf{K}^{\flat}_{\mathrm{lg}}$, where
\begin{equation*}
C^{\mathrm{gen},\natural}_{\mathsf{b}}:=\mathfrak{a}^{\natural}_d\Bigl[2
+C_{\mathsf{b}}2^{p_{\mathsf{b}}}L^d
+\Bigl(e^{2}C_{\mathsf{b}}\Bigl(\frac{d+p_{\mathsf{b}}}{e}\Bigr)^{d+p_{\mathsf{b}}}
+2^d\Bigr)\mathfrak{e}^{\natural}\Bigr] .
\end{equation*}
\end{remark}

\begin{lemma}[Weighted filing columns]\label{5.4:lem-weighted-columns}
Work in the setting of the per-cell transport, Proposition~\ref{5.4:prop-per-cell-transport}, with
$S_{r'}(\Delta_2)$ and $\mathrm{Col}^{\mathrm{file}}$ as in \eqref{5.4:eq-pointed-units-c} and $\theta$ as in
\eqref{5.4:eq-transport-constants-b}, and with the level $\mathrm{lev}\,\Delta'$ and the quantity
$m(\Delta')$ of a cell $\Delta'$, the region $X_v$ and the number $d_v$ of a unit $v$, and the
exponent $c$ as fixed in that setting. Suppose that at some cells $\Delta'$ the filing column obeys the
weighted bound
\begin{equation}\label{5.4:eq-weighted-columns-1}
\mathrm{Col}^{\mathrm{file}}(\Delta')
\le G(\Delta')\,\mathsf{b}^{\mathrm{lg}}(k-\mathrm{lev}\,\Delta')\,\mathsf{K}^{\flat}_{\mathrm{lg}},
\qquad G\ge 1 .
\end{equation}
Then
\begin{equation}\label{5.4:eq-weighted-columns-2}
S_{r'}(\Delta_2)
\le e^{2}\,\mathsf{K}^{\flat}_{\mathrm{lg}}\sum_{\Delta'}
\mathsf{b}^{\mathrm{lg}}(k-\mathrm{lev}\,\Delta')\,e^{-\theta m(\Delta')}\,G(\Delta')\,
\mathfrak{r}(\Delta';\Delta_2),
\end{equation}
where the reserve factor is
\begin{equation}\label{5.4:eq-weighted-columns-3}
\mathfrak{r}(\Delta';\Delta_2)
:=\sup\bigl\{e^{-(c-2)d_v}:\ v \text{ filed at } \Delta',\ X_v\cap\Delta_2\neq\emptyset\bigr\}.
\end{equation}
When $c\ge 3$, one has
\begin{equation*}
\mathfrak{r}(\Delta';\Delta_2)
\le\sup\bigl\{e^{-d_v}:\ v \text{ filed at } \Delta',\ X_v\cap\Delta_2\neq\emptyset\bigr\}.
\end{equation*}
\end{lemma}

\begin{proof}
The bound \eqref{5.4:eq-weighted-columns-2} is the summation over filing cells that concludes the
proof of Proposition~\ref{5.4:prop-per-cell-transport}. Here it is carried out with
\eqref{5.4:eq-weighted-columns-1} in place of the unweighted column bound, so that each summand
retains the factor $G(\Delta')$. The reserve factor
$e^{-(c-2)d_v}$ of each unit $v$ filed at $\Delta'$ whose region $X_v$ meets $\Delta_2$ is bounded by
its supremum \eqref{5.4:eq-weighted-columns-3} over such units.

For the last assertion, let $c\ge 3$. Then $c-2\ge 1$, and since $d_v\ge 0$ this gives
$(c-2)d_v\ge d_v$. Hence $e^{-(c-2)d_v}\le e^{-d_v}$ for every unit $v$ entering
\eqref{5.4:eq-weighted-columns-3}. Taking suprema over these units gives the stated bound.
\end{proof}

\begin{lemma}[Met-indexed and filed per-cell bounds are one statement]\label{5.4:lem-met-filed-equivalence}
Let $\mathfrak{L}^{\mathrm{lg}}$ be the family of long units of
Definition~\ref{5.4:def-pointed-units}. Let the met-cell sums $S^{(\beta)}_{r'}(\Delta';\mathfrak{L}^{\mathrm{lg}})$ and
the filing columns be as in \eqref{5.4:eq-pointed-units-c}. We write
$\mathrm{Col}^{\mathrm{file},(\beta)}_{r'}(\Delta')$ for the first member of
\eqref{5.4:eq-met-filed-equivalence-1} below, the filing column at exponent $\beta$. The column
$\mathrm{Col}^{\mathrm{file}}_{r',c}(\Delta')$ of clause (iii) is this column with the exponent determined by
$c$ as in \eqref{5.4:eq-pointed-units-c}.
\begin{enumerate}
\item[(i)] For every cell $\Delta'$, every $r'$ and every exponent $\beta$,
\begin{equation}\label{5.4:eq-met-filed-equivalence-1}
\begin{aligned}
\sum_{v\in\mathfrak{L}^{\mathrm{lg}}:\ \Delta'_v=\Delta'} w_{r'}(v)\,e^{\beta d_v}
&\le \sum_{\substack{v\in\mathfrak{L}^{\mathrm{lg}}:\ y_v\in\Delta',\\ X_v \text{ meets } \Delta'}}
w_{r'}(v)\,e^{\beta d_v}\\
&\le S^{(\beta)}_{r'}(\Delta';\mathfrak{L}^{\mathrm{lg}}).
\end{aligned}
\end{equation}
\item[(ii)] Let $c\ge 2$, and assume the hypotheses of the per-cell transport
(Proposition~\ref{5.4:prop-per-cell-transport}) other than the per-cell long column
\eqref{5.4:eq-age-base-bounds-c}. Then for every cell $\Delta_2$ of level
$\ell\le k$ and every $r'\le\tfrac{8}{5}\kappa$,
\begin{equation}\label{5.4:eq-met-filed-equivalence-2}
S^{(2)}_{r'}(\Delta_2;\mathfrak{L}^{\mathrm{lg}})
\le C^{\mathrm{gen}}_{\mathsf{b}}\,\mathsf{b}^{\mathrm{lg}}(k-\ell)\cdot
\sup_{\Delta'}\,
\frac{\mathrm{Col}^{\mathrm{file},(2+c)}_{r'}(\Delta')}{\mathsf{b}^{\mathrm{lg}}(k-\operatorname{lev}\Delta')},
\end{equation}
where both sides may be infinite. Moreover, under the same hypotheses, the following hold.
If, for some constant $\mathsf{K}$ and every cell $\Delta'$, either
$S^{(2+c)}_{r'}(\Delta';\mathfrak{L}^{\mathrm{lg}})$ or the middle member of
\eqref{5.4:eq-met-filed-equivalence-1} at $\beta=2+c$ (the sum indexed by the position of $y_v$ in the
closed cell $\Delta'$, with $X_v$ meeting $\Delta'$) is at most
$\mathsf{K}\,\mathsf{b}^{\mathrm{lg}}(k-\operatorname{lev}\Delta')$, then the per-cell long column
\eqref{5.4:eq-age-base-bounds-c} holds with constant $\mathsf{K}$.

Conversely, if \eqref{5.4:eq-age-base-bounds-c} holds with constant $\mathsf{K}^{\flat}_{\mathrm{lg}}$, then
\begin{equation*}
S^{(2)}_{r'}(\Delta_2;\mathfrak{L}^{\mathrm{lg}})
\le C^{\mathrm{gen}}_{\mathsf{b}}\,\mathsf{K}^{\flat}_{\mathrm{lg}}\,\mathsf{b}^{\mathrm{lg}}(k-\ell)
\end{equation*}
for every cell $\Delta_2$ of level $\ell\le k$ and every $r'\le\tfrac{8}{5}\kappa$. Thus the met-indexed
and the filed per-cell bounds are one bound, stated at exponents that differ by $c$.
\item[(iii)] Assume the hypotheses of the per-cube transport
(Proposition~\ref{5.4:prop-per-cube-transport}), and let $c_{\mathrm{sp}}$ be the decay constant of that
transport, on which the constant $C^{LS}$ of \eqref{5.4:eq-transport-constants-a} depends. Let $\Delta'$ be
a cell of level $n'$, let $0\le c<c_{\mathrm{sp}}$ and let $r'\le\tfrac{8}{5}\kappa$. Then
\begin{equation}\label{5.4:eq-met-filed-equivalence-3}
\mathrm{Col}^{\mathrm{file}}_{r',c}(\Delta')
\le \bigl[1+C_{\mathsf{b}}\,2^{p_{\mathsf{b}}}L^{d}\bigl(1+\mathfrak{A}_{\delta}(c_{\mathrm{sp}}-c)\bigr)\bigr]\,
\mathsf{b}^{\mathrm{lg}}(k-n').
\end{equation}
\end{enumerate}
\end{lemma}

\begin{proof}
\emph{(i).} All three members of \eqref{5.4:eq-met-filed-equivalence-1} are sums of the terms
$w_{r'}(v)e^{\beta d_v}$, each non-negative, over subsets of $\mathfrak{L}^{\mathrm{lg}}$. It therefore
suffices to show that each index set is contained in the next.

Let $v$ be filed at $\Delta'$, that is, $\Delta'_v=\Delta'$. By clause (c) of the lemma on cells
against a carrier (Lemma~\ref{5.4:lem-cells-carrier}), $\Delta'_v=\Delta(q_v)$ is well defined and
$y_v\in\Delta'_v$. By the same clause, $X_v$ meets the cell $\Delta(Q)$ for every cube $Q$ of $X_v$, in
particular for $Q=q_v$. Hence $y_v\in\Delta'$ and $X_v$ meets $\Delta'$, which is the first inclusion.

The sum $S^{(\beta)}_{r'}(\Delta';\mathfrak{L}^{\mathrm{lg}})$ of \eqref{5.4:eq-pointed-units-c} runs over
all long units whose carrier meets $\Delta'$. This gives the second inclusion.

\emph{(ii).} If the supremum on the right of \eqref{5.4:eq-met-filed-equivalence-2} is infinite, there is
nothing to prove. Otherwise let $\mathsf{K}$ be this supremum. By the
definition of the supremum,
\begin{equation*}
\mathrm{Col}^{\mathrm{file},(2+c)}_{r'}(\Delta')
\le \mathsf{K}\,\mathsf{b}^{\mathrm{lg}}(k-\operatorname{lev}\Delta')
\quad\text{for every cell }\Delta'.
\end{equation*}
This is the per-cell long column \eqref{5.4:eq-age-base-bounds-c} with constant
$\mathsf{K}$. Together with the remaining hypotheses, which are assumed, the per-cell transport,
Proposition~\ref{5.4:prop-per-cell-transport}, then bounds the met-cell sum at $\Delta_2$ by
$C^{\mathrm{gen}}_{\mathsf{b}}\,\mathsf{b}^{\mathrm{lg}}(k-\ell)\,\mathsf{K}$, which is
\eqref{5.4:eq-met-filed-equivalence-2}.

Now suppose that, for every cell $\Delta'$, either
$S^{(2+c)}_{r'}(\Delta';\mathfrak{L}^{\mathrm{lg}})$ or the middle member of
\eqref{5.4:eq-met-filed-equivalence-1} at $\beta=2+c$ is at most
$\mathsf{K}\,\mathsf{b}^{\mathrm{lg}}(k-\operatorname{lev}\Delta')$. The first
member of \eqref{5.4:eq-met-filed-equivalence-1} at $\beta=2+c$ is the filing column at exponent $2+c$.
By (i), this column is bounded by each of the other two members. Hence
\eqref{5.4:eq-age-base-bounds-c} holds with the same constant. The met-indexed form at exponent
$2+c$ is the indexing in which the representation and bound of the large-field terms and the label-sum
bound state their per-cell long bound.

Conversely, if \eqref{5.4:eq-age-base-bounds-c} holds with constant $\mathsf{K}^{\flat}_{\mathrm{lg}}$, the
supremum in \eqref{5.4:eq-met-filed-equivalence-2} is at most $\mathsf{K}^{\flat}_{\mathrm{lg}}$. Then
\eqref{5.4:eq-met-filed-equivalence-2} gives
\begin{equation*}
S^{(2)}_{r'}(\Delta_2;\mathfrak{L}^{\mathrm{lg}})
\le C^{\mathrm{gen}}_{\mathsf{b}}\,\mathsf{K}^{\flat}_{\mathrm{lg}}\,\mathsf{b}^{\mathrm{lg}}(k-\ell)
\end{equation*}
for every cell $\Delta_2$ of level $\ell\le k$. This is the met-indexed bound at exponent $2$.

\emph{(iii).} Let $v$ be a long unit filed at $\Delta'$, and put $u:=n'-j_v$. By clause (c) of
Lemma~\ref{5.4:lem-cells-carrier}, $n'=\operatorname{lev}\Delta(q_v)\ge j_v$, so $u\ge 0$. The cube $q_v$ is
a $\pi_j$-cube, $j=j_v$, contained in $\Delta'_v=\Delta'$. A cell of level $n'$ is the union of
$L^{d(n'-j)}=L^{du}$ cubes of the grid $\pi_j$. Hence $q_v$ is one of these $L^{du}$ cubes.

The cell $\Delta_0(v)$ contains $y_v$, and $y_v\in\Delta'$ by (i). Therefore $\Delta_0(v)$ equals
$\Delta'$ or touches it. With this position of $\Delta_0(v)$ relative to $\Delta'$, the geometry of the
length bound from exit of the star (Lemma~\ref{5.4:lem-star-exit-length}) or of the length bound at a met
cell (Lemma~\ref{5.4:lem-met-cell-length}) gives
\begin{equation*}
d_v\ge \mathfrak{m}_u,
\end{equation*}
with $\mathfrak{m}_u$ as in \eqref{5.4:eq-transport-constants-a}.

Group the units filed at $\Delta'$ according to $u$. The two facts just established are the same as in
the proof of the per-cube transport, Proposition~\ref{5.4:prop-per-cube-transport}: at most $L^{du}$
positions for $q_v$, and the length lower bound $d_v\ge\mathfrak{m}_u$. For each position of $q_v$, the
units with that position are counted, as in that proof, by the base bound for long units
\eqref{5.4:eq-age-base-bounds-b}, using $d_v\ge\mathfrak{m}_u$ and $c<c_{\mathrm{sp}}$. The age bound
\eqref{5.4:eq-age-base-bounds-a} is then used as in that proof to compare the age function at the level
$j_v$ of the unit with $\mathsf{b}^{\mathrm{lg}}(k-n')$, and the sum over $u\ge 0$ is carried out as there.
This yields \eqref{5.4:eq-met-filed-equivalence-3}.
\end{proof}

\begin{lemma}[The belt at a cell beside the treated class]\label{5.4:lem-boundary-belt}
In the notation of the levels, zones and cells (Definition~\ref{5.4:not-zones-cells}) and of the
constants of the two transports (Definition~\ref{5.4:not-transport-constants}), and for $r'$
as in Proposition~\ref{5.4:prop-per-cell-transport}, let $\Delta_2$ be a cell of level $k$ with
$\Delta_2\not\subset Z$ and $\operatorname{dist}_k(\Delta_2,\partial Z)\le\check{R}_k$, put
$R:=\check{R}_k-1$, and suppose $R\ge 37$.
\begin{enumerate}
\item[(i)] Every cell $\Delta$ of level less than $k$ and every raised cell $\Delta$ satisfies
$\operatorname{dist}_k(\Delta,\Delta_2)\ge R$.
\item[(ii)] If $X_v$ meets $\Delta_2$ and the filing cell $\Delta'_v$ of $v$ has level $k-u<k$, then
$\mathfrak{d}_v\ge L^u\max\{R,s_v\}$, where $s_v$ is defined for each unit $v$ in the first step
of the proof of the per-cell transport (Proposition~\ref{5.4:prop-per-cell-transport}).
\item[(iii)] The filing shells of levels below $k$ are exponentially small in $\check{R}_k$ at
$\Delta_2$: with $\theta$ and $\mathfrak{a}_d$ of \eqref{5.4:eq-transport-constants-b},
\begin{equation}\label{5.4:eq-boundary-belt-1}
S_{r'}(\Delta_2)\le e^{2}\,\mathfrak{a}_d\Bigl[1+C_{\mathsf{b}}\,\mathfrak{w}_d\,
e^{-2\theta(\check{R}_k-2)}\Bigr]\,\mathsf{b}^{\mathrm{lg}}(0)\,\mathsf{K}^{\flat}_{\mathrm{lg}},
\end{equation}
where
\begin{equation}\label{5.4:eq-boundary-belt-2}
\mathfrak{w}_d:=e^{-\theta}\Bigl(\frac{3q\mathsf{c}_d}{2e}\Bigr)^{q}\bigl(1-e^{-\theta}\bigr)^{-1};
\end{equation}
near $\Delta_2$ the window weights and the weights of raised cells are idle.
\item[(iv)] Let $c\ge 3$, let $A\ge 0$ and $p>0$ be constants, and write $\mathfrak{r}$ for
$\mathfrak{r}(\Delta';\Delta_2)$ of Lemma~\ref{5.4:lem-weighted-columns}. If at the cells
$\Delta'$ of level $k-u<k$ the column carries a factor $G\le A(\check{R}_kL^u)^p$, then
$G\mathfrak{r}\le A\,\mathfrak{G}_p$; if at the raised cells of level $k$ it carries
$G\le A\check{R}_k^{\,p}$, then $G\mathfrak{r}\le A\,\mathfrak{G}'_p$. Here
\begin{equation}\label{5.4:eq-boundary-belt-3}
\mathfrak{G}_p:=e\Bigl(\frac{38}{37}\Bigr)^{p}\Bigl(\frac{p\mathsf{c}_d}{e}\Bigr)^{p},\qquad
\mathfrak{G}'_p:=e^{1+2/\mathsf{c}_d}\Bigl(\frac{p\mathsf{c}_d}{e}\Bigr)^{p}.
\end{equation}
\end{enumerate}
The bound for $S_{r'}(\Delta_2)$ used in what follows is not \eqref{5.4:eq-boundary-belt-1}; it
is the per-cell transport, Proposition~\ref{5.4:prop-per-cell-transport}, at $\ell=k$, with the
constant $C^{\mathrm{gen}}_{\mathsf{b}}$. The present lemma shows the room that this bound
leaves at the cells beside $Z$.
\end{lemma}

\begin{proof}
(i) By \eqref{5.4:eq-zones-cells-f} (Definition~\ref{5.4:not-zones-cells}), every cell $\Delta$
of level less than $k$ and every raised cell $\Delta$ satisfies
$\operatorname{dist}_k(\Delta,\partial Z)\ge 2\check{R}_k$. On the other hand
$\Delta_2$ lies within $\check{R}_k$ of $\partial Z$ and is a cell of level $k$, so every point
of $\Delta_2$ lies within $\check{R}_k+1$ of $\partial Z$. By the triangle inequality, such a
cell $\Delta$ satisfies
\begin{equation*}
\operatorname{dist}_k(\Delta,\Delta_2)\ge 2\check{R}_k-(\check{R}_k+1)=\check{R}_k-1=R .
\end{equation*}

(ii) Let $X_v$ meet $\Delta_2$, with $\Delta'_v$ of level $k-u<k$, and let $\delta_v$ denote the
distance from the filing cell $\Delta'_v$ to $\Delta_2$ that enters the first step of the proof
of the per-cell transport (Proposition~\ref{5.4:prop-per-cell-transport}). By (i), $\Delta'_v$
lies at distance at least $R$ from $\Delta_2$, so that $\delta_v\ge R$. That first step bounds
$\mathfrak{d}_v$ from below by $L^u\max\{\delta_v,s_v\}$; with $\delta_v\ge R$ this gives
$\mathfrak{d}_v\ge L^u\max\{R,s_v\}$.

(iii) Let $v$ be as in (ii) and write $s=s_v$. Since $\max\{R,s\}\ge R$ and $\max\{R,s\}\ge s$,
\begin{equation*}
L^u\max\{R,s\}=(L^u-1)\max\{R,s\}+\max\{R,s\}\ge (L^u-1)R+s .
\end{equation*}
Since $u\ge1$ and the blocking factor $L$ is odd and greater than one, $L^u\ge L\ge 3$; and
$R\ge 37>2/3$; hence
\begin{align*}
(L^u-1)R-2(R-1)-\tfrac23 L^u
&=L^u\bigl(R-\tfrac23\bigr)-3R+2\\
&\ge 3\bigl(R-\tfrac23\bigr)-3R+2=0 .
\end{align*}
Combining these two inequalities with (ii), and using $R-1=\check{R}_k-2$, we obtain
\begin{equation}\label{5.4:eq-boundary-belt-4}
\mathfrak{d}_v\ge 2(\check{R}_k-2)+\tfrac23 L^u+s_v .
\end{equation}
We insert \eqref{5.4:eq-boundary-belt-4} into the estimate of the filing shells in the proof of
the per-cell transport (Proposition~\ref{5.4:prop-per-cell-transport}), with the constants
$\theta$ and $\mathfrak{a}_d$ of \eqref{5.4:eq-transport-constants-b} and the growth constant
$C_{\mathsf{b}}$ of the age function; for the shells filed at cells of level below $k$ the lower
bound \eqref{5.4:eq-boundary-belt-4} produces the factor $e^{-2\theta(\check{R}_k-2)}$, and the
result is \eqref{5.4:eq-boundary-belt-1}, with $\mathfrak{w}_d$ as in
\eqref{5.4:eq-boundary-belt-2}. By (i), no raised cell and no cell of level below $k$ lies within
distance $R$ of $\Delta_2$; hence near $\Delta_2$ the window weights and the raised-cell weights
are idle.

(iv) Since $c\ge3$, Lemma~\ref{5.4:lem-weighted-columns} gives
$\mathfrak{r}\le\sup e^{-d_v}$, the supremum over $v$ filed at $\Delta'$ with $X_v$ meeting
$\Delta_2$. Take first $\Delta'$ of level $k-u<k$. By (ii) and $\mathfrak{d}_v=\mathsf{c}_d(d_v+1)$
we have $\mathsf{c}_d(d_v+1)\ge L^uR$, that is, $d_v\ge RL^u/\mathsf{c}_d-1$. As $R\ge37$,
\begin{equation*}
\check{R}_k=R+1\le R+\tfrac{R}{37}=\tfrac{38}{37}R .
\end{equation*}
For $x>0$, and since $p>0$, the function $x^pe^{-x/\mathsf{c}_d}$ attains its maximum at
$x=p\mathsf{c}_d$, so $x^pe^{-x/\mathsf{c}_d}\le(p\mathsf{c}_d/e)^p$. With $x=RL^u$ these three
facts give
\begin{align*}
G\mathfrak{r}
&\le A(\check{R}_kL^u)^p\,e^{1-RL^u/\mathsf{c}_d}
\le A\,e\Bigl(\frac{38}{37}\Bigr)^{p}(RL^u)^p e^{-RL^u/\mathsf{c}_d}\\
&\le A\,e\Bigl(\frac{38}{37}\Bigr)^{p}\Bigl(\frac{p\mathsf{c}_d}{e}\Bigr)^{p}=A\,\mathfrak{G}_p .
\end{align*}
Now let $\Delta'$ be a raised cell of level $k$. By (i) it lies at distance at least $R$ from
$\Delta_2$; for a unit $v$ filed at $\Delta'$ with $X_v$ meeting $\Delta_2$ this gives, by the
definition of $s_v$ used in the first step of the proof of the per-cell transport
(Proposition~\ref{5.4:prop-per-cell-transport}), the bound $s_v\ge R-1$. The same definition gives
$\mathfrak{d}_v=\mathsf{c}_d(d_v+1)\ge s_v$ for such $v$, that is, $d_v\ge s_v/\mathsf{c}_d-1$;
we therefore have
\begin{equation*}
d_v\ge\frac{R-1}{\mathsf{c}_d}-1=\frac{\check{R}_k}{\mathsf{c}_d}-\frac{2}{\mathsf{c}_d}-1 .
\end{equation*}
Hence, with $x=\check{R}_k$ in the same maximum bound,
\begin{equation*}
G\mathfrak{r}\le A\,e^{1+2/\mathsf{c}_d}\,\check{R}_k^{\,p}e^{-\check{R}_k/\mathsf{c}_d}
\le A\,e^{1+2/\mathsf{c}_d}\Bigl(\frac{p\mathsf{c}_d}{e}\Bigr)^{p}=A\,\mathfrak{G}'_p .
\qedhere
\end{equation*}
\end{proof}

\begin{lemma}[Polynomial against exponential]\label{5.4:lem-poly-exp-sup}
For real $p\ge 0$ put
\begin{equation*}
\mathfrak{s}(p):=\sup_{x\ge 0}\,(1+x)^{p}e^{-x}.
\end{equation*}
Then $\mathfrak{s}(p)=1$ if $p\le 1$, and $\mathfrak{s}(p)=p^{p}e^{1-p}$ if $p\ge 1$. In all cases
$\mathfrak{s}(p)\le\max\{1,p^{p}e^{1-p}\}$. Hence, with $C_{\mathsf{b}}$, $p_{\mathsf{b}}$ and
$\mathsf{b}^{\mathrm{lg}}$ as in Definition~\ref{5.4:def-age-base-bounds} and
$\mathfrak{s}_{\mathsf{b}}$ as in \eqref{5.4:eq-transport-constants-c}, for every integer $x\ge 0$
\begin{equation}\label{5.4:eq-poly-exp-sup-1}
e^{-x}\mathsf{b}^{\mathrm{lg}}(x)
\le C_{\mathsf{b}}(1+x)^{p_{\mathsf{b}}}e^{-x}\mathsf{b}^{\mathrm{lg}}(0)
\le C_{\mathsf{b}}\,\mathfrak{s}_{\mathsf{b}}\,\mathsf{b}^{\mathrm{lg}}(0).
\end{equation}
\end{lemma}

\begin{proof}
For $x\ge 0$ we have $(1+x)^{p}e^{-x}=e^{g(x)}$ with $g(x):=p\log(1+x)-x$, so that
\begin{equation*}
g'(x)=\frac{p}{1+x}-1 .
\end{equation*}
Suppose first $p\le 1$. For $x\ge 0$ we have $p/(1+x)\le p\le 1$, hence $g'(x)\le 0$. Thus $g$ is
non-increasing on $[0,\infty)$, its supremum is attained at $x=0$, and $\mathfrak{s}(p)=e^{g(0)}=1$.

Suppose next $p\ge 1$. Then $g'(x)\ge 0$ for $0\le x\le p-1$ and $g'(x)\le 0$ for $x\ge p-1$. Hence
$g$ attains its maximum on $[0,\infty)$ at $x=p-1$, where
\begin{equation*}
g(p-1)=p\log p-(p-1),
\end{equation*}
so that $\mathfrak{s}(p)=p^{p}e^{-(p-1)}=p^{p}e^{1-p}$. At $p=1$ both formulas give the value $1$.
In either case $\mathfrak{s}(p)$ equals one of the two numbers $1$ and $p^{p}e^{1-p}$, whence
$\mathfrak{s}(p)\le\max\{1,p^{p}e^{1-p}\}$.

For \eqref{5.4:eq-poly-exp-sup-1}, let $x\ge 0$ be an integer. The growth bound
\eqref{5.4:eq-age-base-bounds-a}, read at $(s,u)=(0,x)$, gives
$\mathsf{b}^{\mathrm{lg}}(x)\le C_{\mathsf{b}}(1+x)^{p_{\mathsf{b}}}\mathsf{b}^{\mathrm{lg}}(0)$.
Multiplying by $e^{-x}>0$ gives the first inequality of \eqref{5.4:eq-poly-exp-sup-1}.
For the second, the first part of the lemma at $p=p_{\mathsf{b}}$ gives
$(1+x)^{p_{\mathsf{b}}}e^{-x}\le\mathfrak{s}(p_{\mathsf{b}})$. The constant $\mathfrak{s}_{\mathsf{b}}$
of \eqref{5.4:eq-transport-constants-c} bounds $\mathfrak{s}(p_{\mathsf{b}})$ from above. Multiplying
by the non-negative factor $C_{\mathsf{b}}\mathsf{b}^{\mathrm{lg}}(0)$ gives the second inequality.
\end{proof}

\begin{theorem}[Summation of boundary terms at finer zones]\label{5.4:thm-finer-zone-sum}
Let $k$ be the current level. Assume one base bound for long units, in one of the following two
shapes.
\begin{itemize}
\item The per-cube case: the per-cube long-unit bound \eqref{5.4:eq-age-base-bounds-b} together
with the star-exit definition of long units \eqref{5.4:eq-pointed-units-a}, and the condition
(g4) among the hypotheses of the per-cube transport
(Proposition~\ref{5.4:prop-per-cube-transport}).
\item The per-cell case: the per-cell long-unit column \eqref{5.4:eq-age-base-bounds-c}, with its
exponent $c \ge 2$, and the condition (g1) and the carrier convention for pointed units among the
hypotheses of the per-cell transport (Proposition~\ref{5.4:prop-per-cell-transport}).
\end{itemize}
Assume moreover that every pointed unit $v$ satisfies $X_v \subset \Lambda_{j_v}$ in the label of its
birth, that the event census holds, that the conditions (S1)--(S3), which are hypotheses of both
transports, hold, and that the axioms (A0)--(A3), among
them the zone-separation axiom (A3), hold: in the per-cube case in the form in which the per-cube
transport (Proposition~\ref{5.4:prop-per-cube-transport}) assumes them, in the per-cell case at the
current level. Let $\mathfrak{K}$ and
$\mathsf{K}^{\mathrm{lg}} = \mathfrak{K}\,\mathsf{b}^{\mathrm{lg}}(0)$ be the constants of
\eqref{5.4:eq-transport-constants-c} matching the case, that is,
$\mathfrak{K} = C^{LS}$ in the per-cube case and
$\mathfrak{K} = C^{\mathrm{gen}}_{\mathsf{b}}\mathsf{K}^{\flat}_{\mathrm{lg}}$ in the per-cell
case. Then for every current cell $\Delta_2$ of level
$\ell \le k$ the following hold.
\begin{enumerate}
\item For every $r' \le (8/5)\kappa$ the first line below holds, and the second line holds for the
exponent $r$ of the sum $S_r$ provided $r \le (8/5)\kappa$:
\begin{equation}\label{5.4:eq-finer-zone-sum-a}
\begin{gathered}
S_{r'}(\Delta_2;\mathfrak{L}^{\mathrm{lg}}_{c}) \le S_{r'}(\Delta_2;\mathfrak{L}^{\mathrm{lg}})
\le \mathfrak{K}\,\mathsf{b}^{\mathrm{lg}}(k-\ell),\\
S_{r}(\Delta_2;\mathfrak{L}^{\mathrm{lg}}_{d}) \le \mathfrak{K}\,\mathsf{b}^{\mathrm{lg}}(k-\ell).
\end{gathered}
\end{equation}
In particular, at every cell $\Delta_2$ of level $k$ the two sums in the first line are at most
$\mathsf{K}^{\mathrm{lg}}$.
\item Let $\Delta_1$ be a rim cell of a unit $v \in \mathfrak{L}^{\mathrm{lg}}_{d}$, that is, a
cell of $\partial\mathrm{St}_v$, sharing a wall with $\Delta_2$, and let $\rho(\Delta_1)$ be the
walk exponent at $\Delta_1$. Assume the lower bound on the walk exponent at rim cells,
$\rho(\Delta_1) \ge k-\ell$, at such a rim cell. Then the age is paid by the single factor
$e^{-\rho(\Delta_1)}$:
\begin{equation}\label{5.4:eq-finer-zone-sum-1}
e^{-\rho(\Delta_1)}\,\mathsf{b}^{\mathrm{lg}}(k-\ell)
\le C_{\mathsf{b}}\,\mathfrak{s}_{\mathsf{b}}\,\mathsf{b}^{\mathrm{lg}}(0).
\end{equation}
\item At every cell of level $k$ that is not contained in $Z$ and lies within distance
$2\check{R}_k$ of $\partial Z$ (the cells at which the boundary terms sit, among them their
anchor cell $\Delta^{\partial}$),
the bound in the first line of
\eqref{5.4:eq-finer-zone-sum-a} is the number $\mathsf{K}^{\mathrm{lg}}$. In the per-cube case its
only dependence on $\check{R}_k$ is a factor $\check{R}_k$ that appears when $\mathsf{b}^{\mathrm{lg}}$
is expressed through its $\check{R}_k$-free companion $\mathsf{b}^{\mathrm{lg},\natural}$ as in
Notation~\ref{5.4:not-transport-constants}, and that
factor is paid later in this section; in the per-cell case the same holds for the factor
$C^{\mathrm{gen}}_{\mathsf{b}}\,\mathsf{b}^{\mathrm{lg}}(0)$ of $\mathsf{K}^{\mathrm{lg}}$, the
remaining factor $\mathsf{K}^{\flat}_{\mathrm{lg}}$ being the constant of the assumed column
\eqref{5.4:eq-age-base-bounds-c}. In the per-cube case the bound reads neither the
distinguished point $y_v$ of a unit, nor $Z$, nor a window, nor a filing cell (the cell at which
the per-cell column files a unit). In the per-cell case the bound reads the filing cells of the
column \eqref{5.4:eq-age-base-bounds-c}; those contained in $Z$ are controlled by the remark on
the in-class column together with the proposition on own-window weights of this section.
\end{enumerate}
\end{theorem}

\begin{proof}
The carrier lemma (Lemma~\ref{5.4:lem-carrier-in-zone}) is proved as an implication from two
statements: the containment $X_v \subset \Lambda_{j_v}$ in the label of the birth of $v$, and
the event census. Both are hypotheses here. Hence, for every pointed unit $v$ with
$j := j_v$, we have $\bigcup X_v \subset \Omega_j^{\mathrm{cur}}$, and the interior of every cube
of $X_v$ misses every cell of level $< j$. This is the form in which the zone structure enters
both transports below.

\emph{The sum bounds.} Fix a current cell $\Delta_2$ of level $\ell \le k$ and
$r' \le (8/5)\kappa$.

In the per-cube case, the per-cube transport (Proposition~\ref{5.4:prop-per-cube-transport})
gives
\begin{equation*}
S_{r'}(\Delta_2;\mathfrak{L}^{\mathrm{lg}}) \le C^{LS}\,\mathsf{b}^{\mathrm{lg}}(k-\ell).
\end{equation*}
Its hypotheses are the base bound \eqref{5.4:eq-age-base-bounds-b} with
\eqref{5.4:eq-pointed-units-a}, the birth containment, the event census, (S1)--(S3), (A0)--(A3)
in the form assumed there and (g4), all of which are assumed here.
At a cell of level $k$ this is the current-level bound,
Corollary~\ref{5.4:cor-current-level-bound}. At a cell of general level it is
Corollary~\ref{5.4:cor-general-cell-bound}.

In the per-cell case, the per-cell transport (Proposition~\ref{5.4:prop-per-cell-transport})
gives
\begin{equation*}
S_{r'}(\Delta_2;\mathfrak{L}^{\mathrm{lg}})
\le C^{\mathrm{gen}}_{\mathsf{b}}\,\mathsf{b}^{\mathrm{lg}}(k-\ell)\,
\mathsf{K}^{\flat}_{\mathrm{lg}}.
\end{equation*}
Its hypotheses are the column \eqref{5.4:eq-age-base-bounds-c} with its exponent $c \ge 2$,
the carrier convention for pointed units, the carrier lemma, (S1)--(S3), (A1)--(A3) at the
current level and (g1); the carrier lemma was obtained above, and the others are assumed here.

In either case the right-hand side is $\mathfrak{K}\,\mathsf{b}^{\mathrm{lg}}(k-\ell)$, by the
choice of $\mathfrak{K}$ among the constants of the two transports,
see \eqref{5.4:eq-transport-constants-c}. This is the second
inequality in the first line of \eqref{5.4:eq-finer-zone-sum-a}.

The sum $S_{r'}(\Delta_2;\cdot)$ of \eqref{5.4:eq-pointed-units-c} is a sum of non-negative
terms over the units of the family named in its second argument. The family
$\mathfrak{L}^{\mathrm{lg}}_{c}$ of long units of class (c) is contained in
$\mathfrak{L}^{\mathrm{lg}}$. Hence
$S_{r'}(\Delta_2;\mathfrak{L}^{\mathrm{lg}}_{c}) \le S_{r'}(\Delta_2;\mathfrak{L}^{\mathrm{lg}})$,
which is the first inequality.

The family $\mathfrak{L}^{\mathrm{lg}}_{d}$ of long units of class (d) is likewise contained in
$\mathfrak{L}^{\mathrm{lg}}$. When $r \le (8/5)\kappa$, the same transport, with $r$ in place of
$r'$, gives
\begin{equation*}
S_{r}(\Delta_2;\mathfrak{L}^{\mathrm{lg}}) \le \mathfrak{K}\,\mathsf{b}^{\mathrm{lg}}(k-\ell);
\end{equation*}
since $\mathfrak{L}^{\mathrm{lg}}_{d} \subset \mathfrak{L}^{\mathrm{lg}}$ and the terms of the sum
are non-negative, the second line of \eqref{5.4:eq-finer-zone-sum-a} follows.

If $\Delta_2$ has level $k$, then $k-\ell = 0$. The right-hand side then becomes
$\mathfrak{K}\,\mathsf{b}^{\mathrm{lg}}(0) = \mathsf{K}^{\mathrm{lg}}$, which is the particular
case stated after \eqref{5.4:eq-finer-zone-sum-a}.

\emph{The age inequality.} Put $x := k-\ell$, a non-negative integer since $\ell \le k$. By the
lower bound on the walk exponent assumed in the second assertion, $\rho(\Delta_1) \ge x$. Since
$\mathsf{b}^{\mathrm{lg}}$ takes non-negative values, this gives
$e^{-\rho(\Delta_1)}\mathsf{b}^{\mathrm{lg}}(x) \le e^{-x}\mathsf{b}^{\mathrm{lg}}(x)$.

By the lemma on polynomial against exponential (Lemma~\ref{5.4:lem-poly-exp-sup}), at the
integer $x \ge 0$,
\begin{equation*}
e^{-x}\mathsf{b}^{\mathrm{lg}}(x)
\le C_{\mathsf{b}}(1+x)^{p_{\mathsf{b}}}e^{-x}\,\mathsf{b}^{\mathrm{lg}}(0)
\le C_{\mathsf{b}}\,\mathfrak{s}_{\mathsf{b}}\,\mathsf{b}^{\mathrm{lg}}(0).
\end{equation*}
Combining the two inequalities gives \eqref{5.4:eq-finer-zone-sum-1}.

\emph{The boundary cells.} By the particular case noted after
\eqref{5.4:eq-finer-zone-sum-a}, at every cell of level $k$ the bound is
$\mathsf{K}^{\mathrm{lg}}$; in particular this holds
at the cells of level $k$ not contained in $Z$ and within $2\check{R}_k$ of $\partial Z$.

The constant $C^{LS}$ depends only on $d$, $L$, the decay rate $c_{\mathrm{sp}}$ of the per-cube
bound \eqref{5.4:eq-age-base-bounds-b}, $C_{\mathsf{b}}$ and
$p_{\mathsf{b}}$. Moreover, $k-\ell$ enters the per-cube transport only as the argument of
$\mathsf{b}^{\mathrm{lg}}$, and that transport reads neither the position of $y_v$, nor $Z$,
nor $\Lambda$, nor a hub, nor a window, nor $\check{R}_k$
(Proposition~\ref{5.4:prop-per-cube-transport}). The constant
$C^{\mathrm{gen}}_{\mathsf{b}}$ reads neither $Z$, nor $\check{R}_k$, nor a window
(Proposition~\ref{5.4:prop-per-cell-transport}).

Consequently, in the per-cube case $\mathfrak{K} = C^{LS}$ reads neither $Z$ nor $\check{R}_k$,
and in the per-cell case the factor $C^{\mathrm{gen}}_{\mathsf{b}}$ of $\mathfrak{K}$ reads
neither. In the per-cube case, therefore, the dependence of $\mathsf{K}^{\mathrm{lg}}$ on
$\check{R}_k$, and in the per-cell case that of its factor
$C^{\mathrm{gen}}_{\mathsf{b}}\,\mathsf{b}^{\mathrm{lg}}(0)$, is carried by the value
$\mathsf{b}^{\mathrm{lg}}(0)$ of the age function
$\mathsf{b}^{\mathrm{lg}}$ of the base bound, that is, by the factor $\check{R}_k$ that appears
when $\mathsf{b}^{\mathrm{lg}}$ is expressed through $\mathsf{b}^{\mathrm{lg},\natural}$ as in
Notation~\ref{5.4:not-transport-constants}. In the
per-cell case the remaining factor $\mathsf{K}^{\flat}_{\mathrm{lg}}$ of $\mathfrak{K}$ is the
constant of the column \eqref{5.4:eq-age-base-bounds-c}, and is taken from that hypothesis as
stated.

In the per-cube case the only one of the assumed bounds that is consulted is the per-cube bound
\eqref{5.4:eq-age-base-bounds-b}, at
pairs $(j,\square')$ with $\square' \subset \Delta_2$. It is summed over all units $v$ with
$\square' \in X_v$, whatever the cell of $y_v$
(Corollary~\ref{5.4:cor-current-level-bound}). Hence no filing cell and no window is read.

In the per-cell case the input is the per-cell column \eqref{5.4:eq-age-base-bounds-c}, which is
filed at cells. Its filing cells contained in $Z$ are treated in the remark on the in-class
column and the proposition on own-window weights, both later in this section.
\end{proof}

\begin{lemma}[Rim and trunk of a long unit off the treated class]
\label{5.4:lem-rim-trunk}
Let $v\in\mathfrak{L}^{\mathrm{lg}}_{d}$ be a long unit of class (d), and let $\mathrm{St}_v$ and
$\partial\mathrm{St}_v$ be the star of $v$ and its rim, as in Definition~\ref{5.4:def-pointed-units};
the levels $\operatorname{lev}$ of cells and the level $k$ are as in
Notation~\ref{5.4:not-zones-cells}.
\begin{enumerate}
\item[(i)--(ii)] Assertions (i) and (ii) are the two assertions of part (f) of the lemma on cells
against a carrier (Lemma~\ref{5.4:lem-cells-carrier}), applied to $v$.
\item[(iii)] If a cell $\Delta_1\in\partial\mathrm{St}_v$ shares a wall with a cell
$\Delta_2\in\mathrm{St}_v$, then the function $\rho$ of \eqref{5.4:eq-zones-cells-d} satisfies
$\rho(\Delta_1)\ge k-\operatorname{lev}\Delta_2$.
\item[(iv)] Let $\mathsf{y}$ belong to the class $\mathfrak{Y}(\mathrm{St}_v)$ of admissible
families fixed in \eqref{5.4:eq-pointed-units-d}. The component $\mathsf{y}_1$ of $\mathsf{y}$ that
contains the first source cell of $\mathsf{y}$ contains some cell $\Delta_1\in\partial\mathrm{St}_v$,
and
\begin{equation*}
|\mathsf{y}|\ \ge\ |\mathsf{y}_1|\ \ge\ \rho(\Delta_1)\ \ge\ 9\check{R}_k-2 .
\end{equation*}
\item[(v)] Suppose that the label-sum theorem, in its form with $\beta=1$, or the label-sum
theorem at an arbitrary grain, with its parameter $\beta$, is applied to unit data of units $v$ as
above, and suppose that $\beta\le\lambda'_{\Sigma}-3$. Then the quantity bounded in the conclusion
of the theorem so applied is at most
\begin{equation}
\label{5.4:eq-rim-trunk-1}
e^{-(9\check{R}_k-2)}\,S_d\,\mathsf{M}^{\partial}_{r+1,\beta},
\end{equation}
where $r$ and the rim supremum $\mathsf{M}^{\partial}_{r+1,\beta}$ are those of that conclusion,
and no constant of the theorem is changed.
\end{enumerate}
\end{lemma}

\begin{proof}
Assertions (i) and (ii) are part (f) of Lemma~\ref{5.4:lem-cells-carrier}; nothing is to be
proved for them here.

(iii) Let $\Delta_1\in\partial\mathrm{St}_v$ share a wall with $\Delta_2\in\mathrm{St}_v$. By the
inequality ($\rho1$) of \eqref{5.4:eq-zones-cells-d}, applied to $\Delta_1$,
\begin{equation*}
\rho(\Delta_1)\ \ge\ k-\operatorname{lev}\Delta_1+1 .
\end{equation*}
Two cells that share a wall touch, so the axiom (A2) of \eqref{5.4:eq-zones-cells-b} applies to
the pair $\Delta_1,\Delta_2$ and gives $\operatorname{lev}\Delta_1\le\operatorname{lev}\Delta_2+1$.
Inserting this into the previous inequality,
\begin{equation*}
\rho(\Delta_1)\ \ge\ k-(\operatorname{lev}\Delta_2+1)+1\ =\ k-\operatorname{lev}\Delta_2 .
\end{equation*}

(iv) Let $\mathsf{y}\in\mathfrak{Y}(\mathrm{St}_v)$. The union $\mathrm{St}_v\cup\mathsf{y}$ is
connected. The family
$\mathsf{y}_1$ is a connected component of $\mathsf{y}$, so no member of $\mathsf{y}_1$ is adjacent
to a member of another component of $\mathsf{y}$; since $\mathrm{St}_v\cup\mathsf{y}$ is
connected, some member of $\mathsf{y}_1$ must therefore be adjacent to some member of
$\mathrm{St}_v$. This adjacency is through a wall: an adjacency through a hub would make the
member of $\mathrm{St}_v$ concerned a source cell, and this is excluded by (ii). Hence the member
of $\mathsf{y}_1$ in question shares a wall with a cell of $\mathrm{St}_v$, and it is a cell
\begin{equation*}
\Delta_1\in\mathsf{y}_1\cap\partial\mathrm{St}_v .
\end{equation*}
Now $\mathsf{y}_1$ is a connected family which contains $\Delta_1$ and contains a source cell,
namely the first source cell of $\mathsf{y}$. By (ii), such a family has at least $\rho(\Delta_1)$
members, and $\rho(\Delta_1)\ge 9\check{R}_k-2$. Since $\mathsf{y}_1\subset\mathsf{y}$, also
$|\mathsf{y}|\ge|\mathsf{y}_1|$, and the chain of inequalities in (iv) follows.

(v) In the case (d) of the label-sum theorem with $\beta=1$, and in the case (d) of the label-sum
theorem at an arbitrary grain with parameter $\beta$, the trunk series reads
\begin{equation*}
\mathsf{M}^{\partial}_{r+1,\beta}\sum_{n} n\,e^{-(\lambda'_{\Sigma}-\beta)n},
\end{equation*}
with $\mathsf{M}^{\partial}_{r+1,\beta}$ the rim supremum of that conclusion,
the sum running over the values $n=|\mathsf{y}_1|$ of the size of the trunk. By (iv), for unit
data of the units $v$ considered here only the values $n\ge n_0$ occur, where
$n_0:=9\check{R}_k-2$. The proviso $\beta\le\lambda'_{\Sigma}-3$ gives
$\lambda'_{\Sigma}-\beta-2\ge 1$. Hence, for $n\ge n_0$,
\begin{equation*}
e^{-(\lambda'_{\Sigma}-\beta)n}
=e^{-(\lambda'_{\Sigma}-\beta-2)n}\,e^{-2n}
\le e^{-(\lambda'_{\Sigma}-\beta-2)n_0}\,e^{-2n},
\end{equation*}
and we obtain
\begin{align*}
\sum_{n\ge n_0} n\,e^{-(\lambda'_{\Sigma}-\beta)n}
&\le e^{-(\lambda'_{\Sigma}-\beta-2)n_0}\sum_{n\ge 1} n\,e^{-2n}
\le e^{-n_0}.
\end{align*}
In the last step we used $\lambda'_{\Sigma}-\beta-2\ge 1$ and the evaluation
\begin{equation*}
\sum_{n\ge 1} n\,e^{-2n}\ =\ \frac{e^{-2}}{(1-e^{-2})^{2}}\ \approx\ 0.181\ \le\ 1 .
\end{equation*}
The trunk series is thus bounded by $\mathsf{M}^{\partial}_{r+1,\beta}e^{-n_0}$, and this factor
$e^{-n_0}=e^{-(9\check{R}_k-2)}$ passes to the conclusion, which is
thereby bounded by \eqref{5.4:eq-rim-trunk-1}. The other steps of the two theorems are not
affected; in particular the exchange, in their proofs, of a factor $e^{\beta\rho}$, with $\rho$ a
value of the function $\rho$ of \eqref{5.4:eq-zones-cells-d}, for $e^{\beta n}$
is untouched, and no constant of the theorems is altered. Finally, since the decay rate
$\lambda'_{\Sigma}$ of the two theorems satisfies
$\lambda'_{\Sigma}\ge 75c_l-1$, the proviso $\beta\le\lambda'_{\Sigma}-3$ holds for every
$\beta\le 3$ with a wide margin.
\end{proof}

\begin{lemma}[The collar lemma for foreign and sibling windows]\label{5.4:lem-collar}
Let $\mathfrak{a}''$ be a live arrest on a component $X'' := X_{\mathfrak{a}''}$ that is not in $Z$
(that is, $\mathfrak{a}''$ is foreign, or it is a sibling with $k_{a''} = k$). Let $\Delta_2$ be a
current cell with $\Delta_2 \subset W_{\mathfrak{a}''}$, and let $\Delta_1 \in \pi^{\mathfrak{s}}$ share
a wall with $\Delta_2$. Then the level $\operatorname{lev}\Delta_1$ of $\Delta_1$, in the sense of
Notation~\ref{5.4:not-zones-cells}, satisfies $\operatorname{lev}\Delta_1 \le k_{a''}$, and
\begin{equation*}
\rho(\Delta_1) \ge 2\check{R}_{k_{a''}} - 1 .
\end{equation*}
\end{lemma}

\begin{proof}
Throughout, $\check{R} := \check{R}_{k_{a''}}$, and distances and diameters with the index
$k_{a''}$ are measured in the units of level $k_{a''}$.

We first bound the level of $\Delta_1$. Since $\Delta_2 \subset W_{\mathfrak{a}''}$, the interior of
$\Delta_2$ is contained in $X''$. The collar-zone property in \eqref{5.4:eq-zones-cells-g}, applied
to the current cell $\Delta_2$, whose interior meets $X''$, gives
$\operatorname{lev}\Delta_2 \le k_{a''}$. Since $\Delta_1$ and $\Delta_2$ share a wall, the axiom in
\eqref{5.4:eq-zones-cells-b} on the levels of cells sharing a wall gives
$\operatorname{lev}\Delta_1 \le k_{a''} + 1$. Suppose that
$\operatorname{lev}\Delta_1 = k_{a''} + 1$. Then $\Delta_1 \subset \Omega^{\mathrm{cur}}_{k_{a''}+1}$,
and by the collar-zone property in \eqref{5.4:eq-zones-cells-g} the set
$\Omega^{\mathrm{cur}}_{k_{a''}+1}$ lies at positive distance from $X''$. This contradicts
$\Delta_1 \cap \Delta_2 \neq \emptyset$. Hence $\operatorname{lev}\Delta_1 \le k_{a''}$, which is the
first assertion, and consequently $\operatorname{diam}_{k_{a''}}\Delta_1 \le 1$.

Next we bound the distance of $\Delta_1$ from the complement of $X''$. Put
\begin{equation*}
f := \operatorname{dist}_{k_{a''}}\bigl(\,\cdot\,, (X'')^{c}\bigr).
\end{equation*}
It is $1$-Lipschitz for $\operatorname{dist}_{k_{a''}}$. By the window bound in
\eqref{5.4:eq-zones-cells-g}, $f \ge 2\check{R}$ on $W_{\mathfrak{a}''}$, hence on $\Delta_2$. Let
$y \in \Delta_1$ and pick $x \in \Delta_1 \cap \Delta_2$. Then
\begin{equation*}
f(y) \ge f(x) - \operatorname{dist}_{k_{a''}}(x,y) \ge 2\check{R} - \operatorname{diam}_{k_{a''}}\Delta_1
\ge 2\check{R} - 1 .
\end{equation*}
Thus $f \ge 2\check{R} - 1$ on $\Delta_1$.

Now let $\mathsf{y} \ni \Delta_1$ be a connected family containing a source cell with
$|\mathsf{y}| = \rho(\Delta_1)$. By the chain property in \eqref{5.4:eq-zones-cells-d}, $\mathsf{y}$
contains a chain $c_0 = \Delta_1, c_1, \dots, c_m$ of its members in which consecutive members share
a wall and $c_m$ is the first source cell of the chain. Source cells lie in $Z$, and
$Z \subset (X'')^{c}$; therefore $f \equiv 0$ on $c_m$.

Put $F := \bigcup_{i=0}^{m} c_i$. Each $c_i$ is connected and consecutive members share a wall, so
$F$ is connected. The function $f$ is continuous on $F$. On $c_0 = \Delta_1$ it takes values at
least $2\check{R} - 1$, and on $c_m$ it vanishes. Since $F$ is connected, the intermediate value
theorem gives
\begin{equation*}
f(F) \supseteq [\,0, 2\check{R} - 1\,].
\end{equation*}

Let $x \in F$ with $f(x) > 0$. Then $x$ has positive distance from $(X'')^{c}$, so
$x \in \operatorname{int} X''$, and $x \in c_i$ for some $i$. A neighbourhood of $x$ is therefore
contained in $X''$, and it meets $\operatorname{int} c_i$. Hence
$\operatorname{int} c_i \cap X'' \neq \emptyset$. The collar-zone property in
\eqref{5.4:eq-zones-cells-g} then gives $\operatorname{lev} c_i \le k_{a''}$, so that
$\operatorname{diam}_{k_{a''}} c_i \le 1$. Since $c_i$ is connected and $f$ is continuous and
$1$-Lipschitz, $f(c_i)$ is an interval of length at most $1$.

The intervals $f(c_i)$, taken over those members $c_i$ of the chain that contain a point where
$f > 0$, cover $(0, 2\check{R} - 1]$. Each has length at most $1$, so there are at least
$2\check{R} - 1$ distinct such members. They are members of $\mathsf{y}$, and by the property of
$\rho$ in \eqref{5.4:eq-zones-cells-d},
\begin{equation*}
\rho(\Delta_1) = |\mathsf{y}| \ge 2\check{R} - 1 .
\end{equation*}

The argument uses no hypothesis $k_{a''} < k$. No arrest nested in $\mathfrak{a}''$ or enclosing it
enters, because by the index convention in \eqref{5.4:eq-zones-cells-g} the collar and the radius
$\check{R}$ carry the index of the arrest $\mathfrak{a}''$ itself.
\end{proof}

\begin{corollary}[Absorption of a window weight by one collar factor]
Suppose the base bound of \eqref{5.4:eq-transport-constants-d} carries a window weight
$W^{\mathrm{win}}(\check{R}_{k_{a''}})$ at a raised, met cell or cube $\Delta_2 \not\subset Z$ with
$\Delta_2 \subset W_{\mathfrak{a}''}$ as in Lemma~\ref{5.4:lem-collar}, and let $\Delta_1$ be as
there. Write $\check{R} := \check{R}_{k_{a''}}$.
Then a single factor $e^{-\rho(\Delta_1)}$ suffices to absorb the weight:
\begin{equation}\label{5.4:eq-collar-1}
W^{\mathrm{win}}(\check{R})\,e^{-\rho(\Delta_1)} \le W^{\mathrm{win}}(\check{R})\,e^{-(2\check{R}-1)}
\le C^{\mathrm{win}}_{\mathrm{lg}} .
\end{equation}
For the weight $B^{\mathrm{win}}$ of \eqref{5.4:eq-transport-constants-d} the corresponding constant
is
\begin{equation}\label{5.4:eq-collar-a}
C^{\mathrm{win}}_{B'} := \sup_{r \ge 1} B^{\mathrm{win}}(r)\,e^{-(2r-1)}
\le 2e\,C_{\mathsf{b}}\,L^{d}\,C_{\uparrow}^{d}\,(2 + C_{\uparrow})^{p_{\mathsf{b}}}
\Bigl(\frac{q}{2e}\Bigr)^{q} .
\end{equation}
\end{corollary}

\begin{proof}
The first inequality of \eqref{5.4:eq-collar-1} follows from Lemma~\ref{5.4:lem-collar}, which
gives $\rho(\Delta_1) \ge 2\check{R} - 1$.
The second inequality follows from the definition of $C^{\mathrm{win}}_{\mathrm{lg}}$ in
\eqref{5.4:eq-transport-constants-d} together with
\begin{equation*}
-(2\check{R} - 1) \le 1 - \frac{2\check{R} - 2}{\mathsf{c}_d},
\end{equation*}
which holds because $\mathsf{c}_d \ge 1$. Indeed, the right-hand side minus the left-hand side
equals
\begin{equation*}
2 + (2\check{R} - 2)\Bigl(1 - \frac{1}{\mathsf{c}_d}\Bigr),
\end{equation*}
and this is nonnegative for $\check{R} \ge 0$ and $\mathsf{c}_d \ge 1$.

To obtain the bound in \eqref{5.4:eq-collar-a}, insert $\log_L(C_{\uparrow}r) \le C_{\uparrow}r$
into the definition of $B^{\mathrm{win}}$. This gives, for $r \ge 1$,
\begin{equation*}
B^{\mathrm{win}}(r)\,e^{-(2r-1)}
\le 2e\,C_{\mathsf{b}}\,L^{d}\,C_{\uparrow}^{d}\,(2 + C_{\uparrow})^{p_{\mathsf{b}}}\,
r^{q} e^{-2r} .
\end{equation*}
Finally, the function $r \mapsto r^{q} e^{-2r}$ attains its maximum over
$r > 0$ at $r = q/2$, where it equals $(q/(2e))^{q}$.
\end{proof}

\begin{lemma}[Two Cauchy estimates for interpolated derivatives]\label{5.4:lem-two-cauchy}
Let $T$ be a finite set, let $\kappa_1\ge 1$ and $\rho>0$, and put
\begin{equation*}
E:=\{\sigma\in\mathbb{C}:|\sigma|<e^{\kappa_1}\},\qquad
D_\rho:=\{\mathsf{t}\in\mathbb{C}:|\mathsf{t}|<\rho\}.
\end{equation*}
Let $g\colon D_\rho\times[0,1]\times\bar{E}^{T}\to\mathbb{C}$, $(\mathsf{t},t,\mathsf{s})\mapsto
g(\mathsf{t},t,\mathsf{s})$, be continuous and holomorphic in $\mathsf{t}\in D_\rho$, and for each
$\Delta\in T$ holomorphic in $\mathsf{s}(\Delta)\in E$ when all other arguments are held fixed.
Suppose that there are a number $\mathfrak{N}\ge 0$ and a function $g_0$ of $t$ alone such that
\begin{equation*}
|g(\mathsf{t},t,\mathsf{s})-g_0(t)|\le\mathfrak{N}\qquad
\text{for all }(\mathsf{t},t,\mathsf{s})\in D_\rho\times[0,1]\times\bar{E}^{T}.
\end{equation*}
For $\mathsf{y}\subset T$ put
\begin{equation}\label{5.4:eq-two-cauchy-1}
v_{\mathsf{y}}:=\Bigl[\prod_{\Delta\in\mathsf{y}}\int_0^1 d\mathsf{s}(\Delta)\,
\partial_{\mathsf{s}(\Delta)}\Bigr]\int_0^1 dt\,
\partial_{\mathsf{t}}g(\mathsf{t},t,\mathsf{s})\Big|_{\mathsf{t}=0}\,
\Big|_{\mathsf{s}(\Delta)=0\text{ for }\Delta\notin\mathsf{y}}.
\end{equation}
Then
\begin{equation}\label{5.4:eq-two-cauchy-2}
|v_{\mathsf{y}}|\le\rho^{-1}\mathfrak{N}\,(e^{\kappa_1}-1)^{-|\mathsf{y}|}
\le\rho^{-1}\mathfrak{N}\,e^{-(\kappa_1-1)|\mathsf{y}|}.
\end{equation}
The bound carries no factor $2^{|\mathsf{y}|}$ and does not depend on the cardinality of $T$.
\end{lemma}

\begin{proof}
We proceed in four steps.

\emph{Step one: the $\mathsf{t}$-derivative.} Put
$h(t,\mathsf{s}):=\partial_{\mathsf{t}}g(\mathsf{t},t,\mathsf{s})|_{\mathsf{t}=0}$ for
$(t,\mathsf{s})\in[0,1]\times\bar{E}^{T}$. For each fixed $(t,\mathsf{s})$ and each $0<\rho'<\rho$,
since $g$ is holomorphic in $\mathsf{t}\in D_\rho$, Cauchy's formula for the first derivative gives
the first equality below; the second holds because $\oint_{|\mathsf{t}|=\rho'}\mathsf{t}^{-2}\,d\mathsf{t}=0$
and $g_0(t)$ does not depend on $\mathsf{t}$:
\begin{equation}\label{5.4:eq-two-cauchy-3}
\begin{split}
h(t,\mathsf{s})
&=\frac{1}{2\pi i}\oint_{|\mathsf{t}|=\rho'}g(\mathsf{t},t,\mathsf{s})\,\mathsf{t}^{-2}\,d\mathsf{t}\\
&=\frac{1}{2\pi i}\oint_{|\mathsf{t}|=\rho'}\bigl(g(\mathsf{t},t,\mathsf{s})-g_0(t)\bigr)\,
\mathsf{t}^{-2}\,d\mathsf{t}.
\end{split}
\end{equation}
The contour has length $2\pi\rho'$ and the integrand is bounded by $\mathfrak{N}\rho'^{-2}$, so
$|h(t,\mathsf{s})|\le\mathfrak{N}/\rho'$. Letting $\rho'\uparrow\rho$ gives
$|h(t,\mathsf{s})|\le\mathfrak{N}/\rho$.

The set $\{|\mathsf{t}|=\rho'\}\times[0,1]\times\bar{E}^{T}$ is compact ($T$ is finite and $\bar{E}$
is a closed disc), so $g$ is uniformly continuous on it, and the first line of
\eqref{5.4:eq-two-cauchy-3} shows that $h$ is continuous in $(t,\mathsf{s})$. Fix $\Delta\in T$
and all arguments other than $\mathsf{s}(\Delta)$, and let $\Gamma$ be the boundary of a closed
triangle contained in $E$. The function $(\mathsf{t},\sigma)\mapsto g\,\mathsf{t}^{-2}$, with
$\sigma=\mathsf{s}(\Delta)$, is continuous on the compact set $\{|\mathsf{t}|=\rho'\}\times\Gamma$, so
Fubini's theorem allows us to exchange the two contour integrals:
\begin{equation*}
\oint_\Gamma h\,d\mathsf{s}(\Delta)
=\frac{1}{2\pi i}\oint_{|\mathsf{t}|=\rho'}\mathsf{t}^{-2}
\Bigl(\oint_\Gamma g\,d\mathsf{s}(\Delta)\Bigr)d\mathsf{t}=0,
\end{equation*}
the inner integral vanishing by the Cauchy--Goursat theorem, since $g$ is holomorphic in
$\mathsf{s}(\Delta)\in E$. By Morera's theorem, $h$ is holomorphic in each $\mathsf{s}(\Delta)\in E$.
Put $H(\mathsf{s}):=\int_0^1 h(t,\mathsf{s})\,dt$. Then $|H|\le\mathfrak{N}/\rho$, and the same two
arguments, applied with the compact interval $[0,1]$ in place of the circle, show that $H$ is
continuous and holomorphic in each $\mathsf{s}(\Delta)\in E$. By definition, the $t$-integral in
\eqref{5.4:eq-two-cauchy-1} is $H(\mathsf{s})$.

\emph{Step two: the interpolation integrals are differences of evaluations.} For $\Delta\in T$ and
$a\in\{0,1\}$ let $\varepsilon^{a}_{\Delta}$ denote evaluation at $\mathsf{s}(\Delta)=a$. Let $F$ be
a function of $\mathsf{s}(\Delta)$, and possibly of further variables held fixed, holomorphic in
$\mathsf{s}(\Delta)\in E$. Since $e^{\kappa_1}>1$ we have $[0,1]\subset E$. On $[0,1]$ the derivative
of $u\mapsto F|_{\mathsf{s}(\Delta)=u}$ is the complex derivative $\partial_{\mathsf{s}(\Delta)}F$,
which is continuous, so by the fundamental theorem of calculus
\begin{equation*}
\int_0^1 d\mathsf{s}(\Delta)\,\partial_{\mathsf{s}(\Delta)}F
=(\varepsilon^{1}_{\Delta}-\varepsilon^{0}_{\Delta})F.
\end{equation*}
Evaluations in distinct coordinates commute with one another and with setting $\mathsf{s}(\Delta)=0$
for $\Delta\notin\mathsf{y}$. Applying an evaluation $\varepsilon^a_{\Delta}$ to a function holomorphic
in each of its variables separately leaves a function holomorphic in each remaining variable, because
for every remaining $\Delta'$, with all other variables fixed, the result is the original function of
$\mathsf{s}(\Delta')$ with $\mathsf{s}(\Delta)$ frozen at $a$; the same holds for differences of
evaluations. Hence the identity above may be applied once for each $\Delta\in\mathsf{y}$, starting
from $H$, and
\begin{equation}\label{5.4:eq-two-cauchy-4}
v_{\mathsf{y}}=\Bigl[\prod_{\Delta\in\mathsf{y}}(\varepsilon^{1}_{\Delta}-\varepsilon^{0}_{\Delta})\Bigr]H
\,\Big|_{\mathsf{s}(\Delta)=0\text{ for }\Delta\notin\mathsf{y}}.
\end{equation}

\emph{Step three: induction over $\mathsf{y}$.} Write $\mathsf{y}=\{\Delta_1,\dots,\Delta_m\}$ and
$R:=e^{\kappa_1}-1>0$. Put $H_0:=H$ and, for $1\le i\le m$,
$H_i:=(\varepsilon^{1}_{\Delta_i}-\varepsilon^{0}_{\Delta_i})H_{i-1}$. By Step two, each $H_i$ is
holomorphic in each of its remaining variables. We claim that
\begin{equation}\label{5.4:eq-two-cauchy-5}
\sup|H_i|\le(\mathfrak{N}/\rho)\,R^{-i},
\end{equation}
the supremum being taken over $\bar{E}$ in each of the remaining coordinates
$\mathsf{s}(\Delta_{i+1}),\dots,\mathsf{s}(\Delta_m)$ and $\mathsf{s}(\Delta)$, $\Delta\notin\mathsf{y}$.
For $i=0$ this is Step one. Let $i\ge1$ and
assume \eqref{5.4:eq-two-cauchy-5} for $i-1$. Freeze all variables of $H_{i-1}$ except
$\mathsf{s}(\Delta_i)$ at points of $\bar{E}$, and let $\varphi(\sigma)$ be the value of $H_{i-1}$ at
$\mathsf{s}(\Delta_i)=\sigma$. Then $\varphi$ is holomorphic on $E$ and, writing $\sup|\varphi|$ for its
supremum over $E$, $\sup|\varphi|\le(\mathfrak{N}/\rho)R^{-(i-1)}$. Let $u\in[0,1]$ and $0<R'<R$.
Every $\zeta$ with $|\zeta-u|=R'$ satisfies
\begin{equation*}
|\zeta|\le u+R'\le 1+R'<1+R=e^{\kappa_1},
\end{equation*}
so the circle $|\zeta-u|=R'$ lies in $E$, and Cauchy's formula
\begin{equation*}
\varphi'(u)=\frac{1}{2\pi i}\oint_{|\zeta-u|=R'}\varphi(\zeta)(\zeta-u)^{-2}\,d\zeta
\end{equation*}
gives $|\varphi'(u)|\le\sup|\varphi|/R'$. Letting $R'\uparrow R$ gives
$|\varphi'(u)|\le\sup|\varphi|/R$. Therefore
\begin{equation*}
|\varphi(1)-\varphi(0)|\le\int_0^1|\varphi'(u)|\,du\le\sup|\varphi|/R
\le(\mathfrak{N}/\rho)R^{-i}.
\end{equation*}
Since $\varphi(1)-\varphi(0)$ is the value of $H_i$ at the frozen variables, and these were
arbitrary, \eqref{5.4:eq-two-cauchy-5} follows for $i$. By \eqref{5.4:eq-two-cauchy-4}, $v_{\mathsf{y}}$
is $H_m$ evaluated at $\mathsf{s}(\Delta)=0$ for $\Delta\notin\mathsf{y}$, so
$|v_{\mathsf{y}}|\le\rho^{-1}\mathfrak{N}(e^{\kappa_1}-1)^{-|\mathsf{y}|}$, which is the first
inequality in \eqref{5.4:eq-two-cauchy-2}. For $\mathsf{y}=\emptyset$ this is the case $i=0$.

\emph{Step four: the exponential form.} We have
\begin{equation*}
e^{\kappa_1}-1\ge e^{\kappa_1-1}
\iff e^{\kappa_1}(1-e^{-1})\ge 1
\iff \kappa_1\ge\log\bigl(e/(e-1)\bigr).
\end{equation*}
Since $e-1>1$, we have $e/(e-1)<e$, so $\log(e/(e-1))<1\le\kappa_1$. Hence
$(e^{\kappa_1}-1)^{-|\mathsf{y}|}\le e^{-(\kappa_1-1)|\mathsf{y}|}$, which is the second inequality
in \eqref{5.4:eq-two-cauchy-2}.

Only holomorphy in each variable separately was used, so no appeal to the theorems of Hartogs or
Osgood is needed. Neither $|T|$ nor any factor $2^{|\mathsf{y}|}$ enters the estimates.
\end{proof}

\begin{proposition}[The analytic size bound for long pieces]\label{5.4:prop-long-piece-size}
Let $\kappa_1\ge 1$. Assume the analyticity hypothesis \eqref{5.4:eq-age-base-bounds-e} of
Definition~\ref{5.4:def-age-base-bounds}, clauses (i) and (ii), which we recall:
for every long unit $v\in\mathfrak{L}^{\mathrm{lg}}$, every $s$, every block $\square$ and every
finite family $\mathsf{y}\subset T_v$ of derivative cells, with
$E:=\{\sigma\in\mathbb{C}:|\sigma|<e^{\kappa_1}\}$,
\begin{itemize}
\item[(i)] for $|\mathsf{t}|<\rho_A(v,p)$, $t\in[0,1]$ and $\mathsf{s}\in\bar{E}^{\mathsf{y}}$
(extended by $0$ off $\mathsf{y}$) the dilated chain
$\Phi_s^{t+\mathsf{t}\hat{\chi}_{\square}}(\mathsf{s})$ lies in the analyticity domain
$\mathfrak{D}_v$ of $f_v$, and the function
\begin{equation}\label{5.4:eq-long-piece-size-1}
g(\mathsf{t},t,\mathsf{s}):=f_v\bigl(\Phi_s^{t+\mathsf{t}\hat{\chi}_{\square}}(\mathsf{s})\bigr)
\end{equation}
is continuous, holomorphic in $\mathsf{t}$, and holomorphic in each $\mathsf{s}(\Delta)\in E$ with
the other arguments frozen;
\item[(ii)] on the same set, $|g-f_v(1)|\le 2\mathfrak{n}_v=\mathfrak{N}^{\flat}(v)$.
\end{itemize}
Then, for every piece $p$ of a long unit $v$ (Definition~\ref{5.4:def-pointed-units}), with $s$,
$\square$ and $\mathsf{y}$ the data of $p$ and with $m(p):=|\mathsf{y}|$ the number of derivative
cells of $p$, the piece $v_p$ satisfies
\begin{equation}\label{5.4:eq-long-piece-size-2}
|v_p|\le \mathsf{w}_L^{-1}(\eta_p\mathsf{f}_v)^{1/2}\,\mathfrak{N}^{\flat}(v)\,
e^{-(\kappa_1-1)m(p)}=\tfrac12\,\mathfrak{N}^{\mathrm{ex}}_{A}(p).
\end{equation}
\end{proposition}

\begin{proof}
We apply the two Cauchy estimates for interpolated derivatives, Lemma~\ref{5.4:lem-two-cauchy},
with the following data: the finite index set is $T:=\mathsf{y}$; the radius is
\begin{equation*}
\rho:=\rho_A(v,p)=\frac{\mathsf{w}_L}{(\eta_p\mathsf{f}_v)^{1/2}};
\end{equation*}
the function is $g$ of \eqref{5.4:eq-long-piece-size-1}; the comparison function, which depends on
$t$ alone (indeed it is constant), is $g_0:=f_v(1)$; and the bound is
$\mathfrak{N}:=\mathfrak{N}^{\flat}(v)$. The number $\kappa_1$ of the lemma is the $\kappa_1\ge 1$
of the statement, and the set $E$ of the lemma is the set $E$ above.

The hypotheses of Lemma~\ref{5.4:lem-two-cauchy} are then exactly hypotheses (i) and (ii). The
regularity required by the lemma --- $g$ continuous on
$D_\rho\times[0,1]\times\bar{E}^{T}$, where $D_\rho=\{|\mathsf{t}|<\rho\}$, holomorphic in
$\mathsf{t}\in D_\rho$, and, for each $\Delta\in T$, holomorphic in $\mathsf{s}(\Delta)\in E$ with all
other arguments frozen --- is the content of (i) for $|\mathsf{t}|<\rho_A(v,p)$, $t\in[0,1]$ and
$\mathsf{s}\in\bar{E}^{\mathsf{y}}$. The bound required by the lemma,
$|g(\mathsf{t},t,\mathsf{s})-g_0(t)|\le\mathfrak{N}$ on the same set, is (ii), since
$g_0=f_v(1)$ and $\mathfrak{N}=\mathfrak{N}^{\flat}(v)$.

The piece $v_p$ is the quantity $v_{\mathsf{y}}$ of the lemma for this choice of $g$, and the number
of its derivative cells is $|\mathsf{y}|=m(p)$. The lemma therefore gives
\begin{equation*}
|v_p|\le\rho^{-1}\mathfrak{N}\,(e^{\kappa_1}-1)^{-m(p)}
\le\rho^{-1}\mathfrak{N}\,e^{-(\kappa_1-1)m(p)}.
\end{equation*}
Inserting $\rho^{-1}=\mathsf{w}_L^{-1}(\eta_p\mathsf{f}_v)^{1/2}$ and
$\mathfrak{N}=\mathfrak{N}^{\flat}(v)$ yields the inequality in \eqref{5.4:eq-long-piece-size-2}.
The right-hand side is $\rho_A(v,p)^{-1}\mathfrak{N}^{\flat}(v)e^{-(\kappa_1-1)m(p)}$, which is
one half of the size
\begin{equation*}
\mathfrak{N}^{\mathrm{ex}}_{A}(p)=2\rho_A(v,p)^{-1}\,\mathfrak{N}^{\flat}(v)\,e^{-(\kappa_1-1)m(p)};
\end{equation*}
this is the equality in \eqref{5.4:eq-long-piece-size-2}.
\end{proof}

\begin{remark}[Three factors $2$]
Three factors $2$ occur in the size bookkeeping of long units, and they are distinct.
($\alpha$) The factor $2$ inside $\mathfrak{N}^{\flat}(v)=2\mathfrak{n}_v$ comes from the size bound
for a dilated term, $|f\circ\mathsf{R}-f(1)|\le 2\mathfrak{n}_v$. ($\beta$) The factor $2$ in the
bundle norm $\|\mathfrak{b}\|$ comes from the bundle having two members, each within $\mathfrak{N}$
of $f(1)$; here $\Phi_s^{0}:=\Phi_s^{\chi\equiv 0}$ does not depend on $\mathsf{t}$, but it does
depend on $\mathsf{s}$. ($\gamma$) The factor $2$ in the size
$\mathfrak{N}^{\mathrm{ex}}_{A}(p)$, which carries $2\rho_A(v,p)^{-1}$, is a spare factor $2$ by
\eqref{5.4:eq-long-piece-size-2}. For the long units $\mathfrak{L}^{\mathrm{lg}}_{c}$ of the class
introduced in Lemma~\ref{5.2:lem-exterior-pieces}, treated through the bundles of
Definition~\ref{5.2:def-boundary-bundles}, the present proposition is not needed: there the bundle
norm $\|\mathfrak{b}\|$ is estimated on the diagonal of the polydisc on which the bundles are
estimated, from the undilated containment alone.
\end{remark}

\begin{corollary}[Long units in the boundary-term sums]\label{5.4:cor-long-unit-sums}
Let $a:=\tfrac{31}{20}\kappa$ and let $Q_*\in\pi_k$. Assume the following statements.
\begin{enumerate}
\item[(a)] Summation of boundary terms at finer zones, Theorem~\ref{5.4:thm-finer-zone-sum}, which is
proved as an implication from the statements named among its hypotheses; we use it in its clause
for the cells of level $k$ not contained in $Z$ and within distance $2\check{R}_k$ of $\partial Z$. At such
a cell $\Delta_2$ at distance $0$ from $\partial Z$ it gives, for every $r'\le\tfrac85\kappa$, that the sum
$S_{r'}(\Delta_2;\mathfrak{L}^{\mathrm{lg}})$ of \eqref{5.4:eq-pointed-units-c} is at most
$\mathsf{K}^{\mathrm{lg}}$ under the first alternative of (b), and at most
$\mathsf{K}^{\mathrm{lg}}\,\mathfrak{W}$, with $\mathfrak{W}$ defined below, under the second.
Here $\mathsf{K}^{\mathrm{lg}}$ is the constant of
\eqref{5.4:eq-transport-constants-c}.
\item[(b)] The base bound for long units, with its age function $\mathsf{b}^{\mathrm{lg}}$, in the first
or in the second of its two alternatives listed in Theorem~\ref{5.4:thm-finer-zone-sum}.
\item[(c)] Parts (b), (c), (d) and (e) of the frame of the bundle bounds.
\item[(d)] The bound on the number $\#\partial^{+}\mathrm{St}^{j_v}_v$ of rim cells of a long unit $v$
of level $j_v$, stated together with the frame of the bundle bounds.
\item[(e)] The conversion corollary for $\theta'(g_k)$, with its exponents $r_{\mathrm{pr}}$ and $p_1$,
together with the conditions
\begin{equation*}
\check{R}_k\ge\check{\mathsf{T}}_k^{\,r_{\mathrm{pr}}},\qquad
\mathsf{T}_k\le\check{\mathsf{T}}_k+1,\qquad
\check{\mathsf{T}}_k\ge\mathsf{T}_{\gamma}\ge 1 .
\end{equation*}
\item[(f)] Under the second alternative of (b), the statement on the in-$Z$ weights of the long units,
which supplies the weight factors $W_A$, $W_{\mathrm{own}}$ and their exponent.
\end{enumerate}
Put $\mathfrak{W}:=1$ under the first alternative of (b). Under the second alternative put
$\mathfrak{W}:=W_A W_{\mathrm{own}}$, where the in-$Z$ weight factors satisfy
\begin{equation*}
W_A\in\bigl\{1,\tfrac32\cdot 3^{p_W}\bigr\},\qquad
W_{\mathrm{own}}\in\bigl\{1,\,1+C^{\mathrm{win}}_{\mathrm{lg}}\bigr\}.
\end{equation*}
Here $W_A$, $W_{\mathrm{own}}$ and the exponent $p_W$ are those of the statement in hypothesis (f).
Define
\begin{equation}\label{5.4:eq-long-unit-sums-a}
\begin{aligned}
\mathsf{C}^{\mathrm{lg},c}&:=2e\,S_d^{(1/2)}\,\mathfrak{K}\,\mathsf{b}^{\mathrm{lg},\natural}\,\mathfrak{W},\\
\mathsf{C}^{\Sigma,\mathrm{lg}}&:=80e\cdot d\,L^{d-1}S_d\,C^{0}_{\square}\,\mathsf{c}_d^{2}\,
\mathsf{c}_{\mathsf{f}}\,\bigl(1+C^{\mathrm{win}}_{\mathrm{lg}}\bigr)\,C_{\mathsf{b}}\,\mathfrak{s}_{\mathsf{b}}\,
\mathfrak{K}\,\mathsf{b}^{\mathrm{lg},\natural},\\
\mathsf{C}^{\mathrm{lg},d}_{\circ}&:=16e\cdot d\,L^{d-1}S_d\,\bigl(1+C^{\mathrm{win}}_{\mathrm{lg}}\bigr)\,
C_{\mathsf{b}}\,\mathfrak{s}_{\mathsf{b}}\,\mathfrak{K},\qquad
\mathsf{C}^{\mathrm{lg},d}:=\mathsf{C}^{\mathrm{lg},d}_{\circ}\,\mathsf{b}^{\mathrm{lg},\natural}.
\end{aligned}
\end{equation}
\begin{enumerate}
\item[(i)] For every $r\le a+7$, the sum over the bundles $\mathfrak{b}$ of the long singles
$v\in\mathfrak{L}^{\mathrm{lg}}_c$ of class (c) whose label contains $Q_*$ satisfies
\begin{equation}\label{5.4:eq-long-unit-sums-1}
\begin{aligned}
\sum_{\substack{\mathfrak{b}:\ v\in\mathfrak{L}^{\mathrm{lg}}_c\\ \mathrm{lab}(\mathfrak{b})\ni Q_*}}
\|\mathfrak{b}\|\,e^{r d_k(\mathrm{lab}\,\mathfrak{b})}
&\le e^{-\frac12(9\check{R}_k-2)}\cdot 2e\,S_d^{(1/2)}\cdot\mathfrak{K}\,\mathsf{b}^{\mathrm{lg}}(0)\cdot\mathfrak{W}\\
&\le \mathsf{c}_{\theta'}\,\theta'(g_k)\,\mathsf{C}^{\mathrm{lg},c}.
\end{aligned}
\end{equation}
If the in-$Z$ weights are treated by the alternative in which the supremum $\sigma^{\mathrm{lg}}$ of
\eqref{5.4:eq-transport-constants-d} enters, the constant $\mathsf{C}^{\mathrm{lg},c}$ is replaced by
$\mathsf{C}^{\mathrm{lg},c}\sigma^{\mathrm{lg}}$ and the first right-hand side of
\eqref{5.4:eq-long-unit-sums-1} acquires the factor $\sigma^{\mathrm{lg}}$, the weight in
$\check{R}_k$ being kept as a separate displayed factor and not absorbed into $\sigma^{\mathrm{lg}}$.
\item[(ii)] Assume the hypotheses of the piece-sum bound for long tails of class (d), proved in this
chapter. Then for every $r\le a+8$, the sum over the pieces $p$ of the long tails of class (d) whose label
$\mathrm{lab}^{\circ}_1(p)$ contains $Q_*$ satisfies
\begin{equation*}
\sum_{\substack{p\ \text{of}\ \mathfrak{L}^{\mathrm{lg}}_d\\ \mathrm{lab}^{\circ}_1(p)\ni Q_*}}
\mathfrak{N}^{\mathrm{ex}}_{A}(p)\,e^{r d_k(\mathrm{lab}^{\circ}_1(p))}
\le \theta'(g_k)\,\mathsf{C}^{\Sigma,\mathrm{lg}} .
\end{equation*}
\item[($\mathrm{ii}\flat$)] Assume the hypotheses of the bound on $\square$-bundles of long tails of
class (d), proved in this chapter; among them is the supplementary hypothesis on $\square$-bundles of long
tails of class (d) under which the bound on $\square$-bundles is formulated.
Then for every $r\le a+8$ the sum over the $\square$-bundles $\mathfrak{b}^{\circ}$ of
$v\in\mathfrak{L}^{\mathrm{lg}}_d$ whose label contains $Q_*$ satisfies
\begin{equation*}
\sum_{\substack{\mathfrak{b}^{\circ}:\ v\in\mathfrak{L}^{\mathrm{lg}}_d\\ \mathrm{lab}\ni Q_*}}
\|\mathfrak{b}^{\circ}\|\,e^{r d_k(\mathrm{lab})}
\le e^{-(9\check{R}_k-2)}\,\mathsf{C}^{\mathrm{lg},d}_{\circ}\,\mathsf{b}^{\mathrm{lg}}(0)
\le \mathsf{c}_{\theta'}\,\theta'(g_k)\,\mathsf{C}^{\mathrm{lg},d}.
\end{equation*}
This bound rests only on the following:
\begin{itemize}
\item the zeroth-order bound for $\square$-bundles of long tails of class (d);
\item the linearity statement for the $\square$-bundles;
\item a further input estimate of the bound on $\square$-bundles, fixed with that bound;
\item the geometric axiom \eqref{5.4:eq-zones-cells-e}, or the geometric statement that may replace it;
\item the base bound for long units.
\end{itemize}
In particular it uses neither the amplitude bound for long units, nor the first part of the fifth of the
statements on units of class (c), nor the cell-weight bound, nor the lemma on the blocks $\square$ that
supplies the constant $C^{0}_{\square}$.
\item[(iii)] The constant $\mathsf{C}^{\mathrm{ex}}_1$ has one value throughout the boundary-term sums;
that value is fixed where these sums are combined into the total.
\item[(iv)] No constant of this corollary, that is, none of the constants in
\eqref{5.4:eq-long-unit-sums-a} and $\mathsf{C}^{\mathrm{ex}}_1$, depends on any of the following:
\begin{itemize}
\item $K$, $k$, a level or an age;
\item the carrier sizes $N_j(X_v)$ of Lemma~\ref{5.4:lem-cells-carrier}, and a further size parameter
$N_0$;
\item a number of windows, of arrests or of hubs;
\item the volume.
\end{itemize}
The radius $\check{R}$ enters only in the following ways:
\begin{itemize}
\item inside the displayed suprema $\sigma^{\mathrm{lg}}$, $C^{\mathrm{win}}_{\mathrm{lg}}$
of \eqref{5.4:eq-transport-constants-d} and $C^{\mathrm{win}}_{B'}$ of \eqref{5.4:eq-collar-a};
\item through the lower bound on $\check{R}_k$ imposed where $\check{R}_k$ is fixed;
\item in the displayed factors $e^{-\frac12(9\check{R}_k-2)}$ and $e^{-(9\check{R}_k-2)}$;
\item in the inequality $\mathsf{b}^{\mathrm{lg}}(0)\le\mathsf{b}^{\mathrm{lg},\natural}\check{R}_k$.
\end{itemize}
\end{enumerate}
\end{corollary}

\begin{proof}
We prove clause (i) from hypotheses (a)--(f). Clause (ii) is the conclusion of the piece-sum bound for
long tails of class (d), and clause ($\mathrm{ii}\flat$) that of the bound on $\square$-bundles of long tails
of class (d), each under its own hypotheses; both are proved in this chapter and their proofs are not
repeated here. The value in clause (iii) is fixed where the boundary-term sums are combined into the total.
Fix $r\le a+7$ and $Q_*\in\pi_k$.

\emph{The root cell.} Let $v\in\mathfrak{L}^{\mathrm{lg}}_c$, let $X_v$ be its carrier, as in
Lemma~\ref{5.4:lem-cells-carrier}, and let $K_v$ be the family of $\pi_k$-cubes attached to $v$ in the
definition of pointed units accompanying \eqref{5.4:eq-pointed-units-c}. Since $K_v\not\subset Z$, there
is a $\pi_k$-cube $P'\in K_v$ with
$P'\not\subset Z$ which holds a cube $\square'$ of $X_v$. The interior of $P'$ does not meet $Z$, and
therefore $\square'\not\subset Z$. For a cube $\square$ we write $[\square]_k$ for the $\pi_k$-cube
containing it, as in \eqref{5.4:eq-zones-cells-a}. Two cases arise.

\emph{First case: some cube $\square$ of $X_v$ lies in a $\pi_k$-cube $P\subset Z$.} Take a path in $X_v$
from $\square$ to $\square'$ along which consecutive cubes share a wall. At its start the path has
$[\square]_k=P\subset Z$. At its end it has $[\square']_k=P'\not\subset Z$. Hence there is a first
consecutive pair $\square_1,\square_2$ on the path with
\begin{equation*}
[\square_1]_k\subset Z,\qquad [\square_2]_k\not\subset Z .
\end{equation*}
Taking the first such pair fixes the choice. Put $P_i:=[\square_i]_k$. Then $P_1\ne P_2$, since one
lies in $Z$ and the other does not. Moreover $P_1\cap P_2\supseteq\square_1\cap\square_2$, which is
$(d-1)$-dimensional, so $P_1$ and $P_2$ share a wall. This wall lies in $P_1\subset Z$. It also lies in
$P_2$, whose interior is contained in $Z^c$, so it lies in the closure of $Z^c$. Consequently the wall lies
in $Z\cap\operatorname{cl}Z^c\subset\partial Z$. We set $\Delta^{\partial}_v:=P_2$.

\emph{Second case: no cube of $X_v$ lies in a $\pi_k$-cube of $Z$, and $X_v\cap Z\neq\emptyset$.} Here
$X_v$ and $Z$ are in closed contact. Take a cube $\square$ of $X_v$ which touches $Z$, and set
$\Delta^{\partial}_v:=[\square]_k$. By the assumption of this case $\Delta^{\partial}_v\not\subset Z$, and
$\Delta^{\partial}_v$ touches $Z$ because it contains $\square$.

In both cases $\Delta^{\partial}_v$ is a $\pi_k$-cube, not contained in $Z$, at distance $0$ from
$\partial Z$. Hence, by \eqref{5.4:eq-zones-cells-f}, it is a current cell of level $k$ that has not been
raised. Further, $X_v$ meets $\Delta^{\partial}_v$, since $\Delta^{\partial}_v$ contains a cube of $X_v$,
and $\Delta^{\partial}_v\in K_v$. Since the interior of $\Delta^{\partial}_v$ is contained in $Z^c$, we
have
\begin{equation*}
\operatorname{dist}_k(\Delta^{\partial}_v,\Lambda^{\sim})\ge 9\check{R}_k .
\end{equation*}
Thus $\Delta^{\partial}_v\in\pi^{\mathfrak{s}}$, and $\Delta^{\partial}_v$ lies in no window.

\emph{The frame.} The bundles $\mathfrak{b}$ of $v$ are indexed by the bundle families $\tilde{\mathsf{y}}$.
We use the frame of the bundle bounds, whose parts make no use of the class of $v$.
\begin{itemize}
\item By part (b), $|\mathfrak{b}|\le\|\mathfrak{b}\|$.
\item By part (c), $\mathfrak{b}=0$ unless $\tilde{\mathsf{y}}\in\mathfrak{Y}_v$.
\item By part (e), $\mathrm{lab}(\mathfrak{b})\in\mathbf{D}_k$, and
\begin{equation*}
d_k(\mathrm{lab}\,\mathfrak{b})\le d_k(K_v)+c_l^{\circ}|\tilde{\mathsf{y}}|,\qquad
d_k(\mathrm{lab}\,\mathfrak{b})\ge 9\check{R}_k-2 .
\end{equation*}
\end{itemize}
Moreover $\mathrm{lab}(\mathfrak{b})\supseteq K_v\ni\Delta^{\partial}_v$, and in the sum of clause (i)
also $\mathrm{lab}(\mathfrak{b})\ni Q_*$. Hence \eqref{5.4:eq-zones-cells-c} gives
\begin{equation*}
d_k(\mathrm{lab}\,\mathfrak{b})\ge\operatorname{dist}_k(Q_*,\Delta^{\partial}_v).
\end{equation*}
Averaging the two lower bounds gives
\begin{equation}\label{5.4:eq-long-unit-sums-2}
d_k(\mathrm{lab}\,\mathfrak{b})\ge\tfrac12(9\check{R}_k-2)
+\tfrac12\operatorname{dist}_k(Q_*,\Delta^{\partial}_v).
\end{equation}
We write $e^{r d_k(\mathrm{lab}\,\mathfrak{b})}$ as the product of $e^{-d_k(\mathrm{lab}\,\mathfrak{b})}$
and $e^{(r+1)d_k(\mathrm{lab}\,\mathfrak{b})}$. We bound the first factor by
\eqref{5.4:eq-long-unit-sums-2}. We bound the second by the upper bound of part (e). This
gives
\begin{equation}\label{5.4:eq-long-unit-sums-3}
e^{r d_k(\mathrm{lab}\,\mathfrak{b})}
\le e^{-\frac12(9\check{R}_k-2)}\,e^{-\frac12\operatorname{dist}_k(Q_*,\Delta^{\partial}_v)}\,
e^{(r+1)d_k(K_v)}\,e^{c_l^{\circ}(r+1)|\tilde{\mathsf{y}}|}.
\end{equation}
Set $\lambda:=\kappa_1-1-c_l^{\circ}(r+1)$; then $\lambda\ge\lambda_{\Sigma}$. Let $j:=j_v$. Part (d) of
the frame gives, for the sum over the families $\tilde{\mathsf{y}}$ admitted there,
\begin{equation*}
\sum_{\tilde{\mathsf{y}}}e^{-\lambda|\tilde{\mathsf{y}}|}
\le e^{\epsilon_{\partial}\,\#\partial^{+}\mathrm{St}^j_v}\le e^{d_v+1}.
\end{equation*}
The second inequality follows from the bound on the number $\#\partial^{+}\mathrm{St}^j_v$ of rim cells,
hypothesis (d), and \eqref{5.4:eq-zones-cells-i}. From \eqref{5.4:eq-long-unit-sums-3} and this bound,
with the sum restricted to $\tilde{\mathsf{y}}\in\mathfrak{Y}_v$ by part (c) of the frame, we obtain
\begin{equation}\label{5.4:eq-long-unit-sums-4}
\begin{aligned}
\sum_{\substack{\tilde{\mathsf{y}}\in\mathfrak{Y}_v\\ \mathrm{lab}\ni Q_*}}
\|\mathfrak{b}\|\,e^{r d_k(\mathrm{lab}\,\mathfrak{b})}
&\le e^{-\frac12(9\check{R}_k-2)}\,e^{-\frac12\operatorname{dist}_k(Q_*,\Delta^{\partial}_v)}\\
&\quad\cdot 2\,\mathfrak{N}^{\flat}(v)\,e^{(r+1)d_k(K_v)}\,e^{d_v+1}.
\end{aligned}
\end{equation}

\emph{Grouping by the root.} We sum \eqref{5.4:eq-long-unit-sums-4} over $v\in\mathfrak{L}^{\mathrm{lg}}_c$
and group the terms according to $\Delta:=\Delta^{\partial}_v$. Each $v$ falls into exactly one group. All
terms are non-negative, and $X_v$ meets $\Delta^{\partial}_v$. Hence the group of $\Delta$ is a sub-sum of
the sum over all $v$ with $X_v$ meeting $\Delta$. In that sum, the weight $w_{r+1}(v)$ of the sum
$S_{r+1}$ of \eqref{5.4:eq-pointed-units-c} dominates $\mathfrak{N}^{\flat}(v)e^{(r+1)d_k(K_v)}$, so that
\begin{equation*}
2\,\mathfrak{N}^{\flat}(v)\,e^{(r+1)d_k(K_v)}\,e^{d_v+1}\le 2e\,w_{r+1}(v)\,e^{d_v}.
\end{equation*}
We get
\begin{equation}\label{5.4:eq-long-unit-sums-5}
\begin{aligned}
\sum_{v}\sum_{\substack{\tilde{\mathsf{y}}\in\mathfrak{Y}_v\\ \mathrm{lab}\ni Q_*}}
\|\mathfrak{b}\|\,e^{r d_k(\mathrm{lab}\,\mathfrak{b})}
&\le e^{-\frac12(9\check{R}_k-2)}\sum_{\Delta\in\pi_k}e^{-\frac12\operatorname{dist}_k(Q_*,\Delta)}
\cdot 2e\sum_{v:\ X_v\ \text{meets}\ \Delta}w_{r+1}(v)\,e^{d_v}\\
&\le e^{-\frac12(9\check{R}_k-2)}\cdot S_d^{(1/2)}\cdot 2e\cdot\mathsf{K}^{\mathrm{lg}}\cdot\mathfrak{W}.
\end{aligned}
\end{equation}
The second inequality in \eqref{5.4:eq-long-unit-sums-5} uses three facts.

First, only cells of the form $\Delta=\Delta^{\partial}_v$ carry a non-empty group. Such a cell is a
current cell of level $k$ not contained in $Z$. In the inner sum we bound $e^{d_v}$ by $e^{2d_v}$ and
apply Theorem~\ref{5.4:thm-finer-zone-sum} at $\Delta$ with $r':=r+1$. This choice is admissible:
$r+1\le a+8$, and
\begin{equation*}
a+8=\tfrac{31}{20}\kappa+8\le\tfrac{32}{20}\kappa=\tfrac85\kappa
\end{equation*}
holds if and only if $\kappa\ge160$, which is implied by the standing condition $\kappa\ge180$ on the
auxiliary parameter $\kappa$.

Second, the factor $\mathfrak{W}$ depends on the alternative of the base bound.
\begin{itemize}
\item Under the first alternative hypothesis (a) gives at $\Delta$ the bound $\mathsf{K}^{\mathrm{lg}}$;
it is the long-unit bound at a cell of the current
level, Corollary~\ref{5.4:cor-current-level-bound}. Since $\Delta\not\subset Z$, it involves only cubes
off $Z$ and no column in $Z$. Hence $\mathfrak{W}=1$.
\item Under the second alternative hypothesis (a) gives at $\Delta$ the bound
$\mathsf{K}^{\mathrm{lg}}\mathfrak{W}$ with $\mathfrak{W}=W_AW_{\mathrm{own}}$, the weight factors being
those supplied by hypothesis (f).
\end{itemize}

Third, the number of grid cubes at integer distance $x$ from $Q_*$ is at most $(2x+3)^d$. Hence
\begin{equation*}
\sum_{\Delta\in\pi_k}e^{-\frac12\operatorname{dist}_k(Q_*,\Delta)}\le S_d^{(1/2)},
\end{equation*}
with $S_d^{(1/2)}$ as in \eqref{5.4:eq-zones-cells-j}.

Since $\mathsf{K}^{\mathrm{lg}}=\mathfrak{K}\,\mathsf{b}^{\mathrm{lg}}(0)$ by
\eqref{5.4:eq-transport-constants-c}, the estimate \eqref{5.4:eq-long-unit-sums-5} is the first inequality
of \eqref{5.4:eq-long-unit-sums-1}. No window term is used in it.

\emph{Conversion into $\theta'(g_k)$.} For every real $x$ one has
\begin{equation*}
\tfrac12(9x-2)=(x-2)+\tfrac72x+1 .
\end{equation*}
We apply this at $x=\check{R}_k$ and distinguish two cases.

Suppose first that $\mathsf{b}^{\mathrm{lg}}(0)$ does not depend on $\check{R}$.
Then
\begin{equation*}
e^{-\frac12(9\check{R}_k-2)}\,\mathsf{K}^{\mathrm{lg}}\le e^{-(\check{R}_k-2)}\,\mathsf{K}^{\mathrm{lg}} .
\end{equation*}
The factor $e^{-(\frac72\check{R}_k+1)}$ is not used.

Suppose next that $\mathsf{b}^{\mathrm{lg}}(0)\le\mathsf{b}^{\mathrm{lg},\natural}\check{R}_k$, with the
constant $\mathsf{b}^{\mathrm{lg},\natural}$ of \eqref{5.4:eq-transport-constants-d}; this inequality is of
the same kind as the
bound of the age function of the frame of the bundle bounds by $\mathsf{b}^{\natural}\check{R}_k$. The
maximum of $x e^{-\frac72 x}$ over $x\ge0$ is attained at $x=\frac27$, so $x e^{-\frac72 x}\le\frac{2}{7e}$.
Therefore
\begin{equation*}
\begin{aligned}
e^{-\frac12(9\check{R}_k-2)}\,\check{R}_k
&=e^{-(\check{R}_k-2)}\cdot e^{-1}\cdot\bigl[\check{R}_k\,e^{-\frac72\check{R}_k}\bigr]\\
&\le e^{-(\check{R}_k-2)}\cdot\frac{2}{7e^{2}}\le e^{-(\check{R}_k-2)} .
\end{aligned}
\end{equation*}
In either case the left-hand side of \eqref{5.4:eq-long-unit-sums-1} is at most
\begin{equation*}
e^{-(\check{R}_k-2)}\cdot 2e\,S_d^{(1/2)}\,\mathfrak{K}\,\mathsf{b}^{\mathrm{lg},\natural}\,\mathfrak{W}
=e^{-(\check{R}_k-2)}\,\mathsf{C}^{\mathrm{lg},c}.
\end{equation*}
Finally, $\check{R}_k\ge\check{\mathsf{T}}_k^{\,r_{\mathrm{pr}}}$ gives the first inequality in
\begin{equation*}
e^{-(\check{R}_k-2)}\le e^{2}\,e^{-\check{\mathsf{T}}_k^{\,r_{\mathrm{pr}}}}
\le\mathsf{c}_{\theta'}\,C_{\theta'}\,(\check{\mathsf{T}}_k+1)^{p_1-q}
\le\mathsf{c}_{\theta'}\,\theta'(g_k).
\end{equation*}
The second and third inequalities are those of the conversion corollary for $\theta'(g_k)$, which apply
under the conditions $\mathsf{T}_k\le\check{\mathsf{T}}_k+1$ and
$\check{\mathsf{T}}_k\ge\mathsf{T}_{\gamma}\ge1$ of hypothesis (e). This proves the second inequality of
\eqref{5.4:eq-long-unit-sums-1}, and with it clause (i).

\emph{Particular configurations.} The argument covers the following configurations without change.
\begin{itemize}
\item If $\tilde{\mathsf{y}}=\emptyset$, that is, $X_v$ passes through a hub, the bundle is the term with
$|\tilde{\mathsf{y}}|=0$ of the sum in part (d) of the frame. Its norm is
$\|\mathfrak{b}\|=2\mathfrak{N}^{\flat}(v)$, and the smallness comes from
$d_k(K_v)\ge 9\check{R}_k-2$.
\item If $X_v$ is walled to several hubs, their number satisfies $\#\mathfrak{H}_v\le 2N_j(X_v)$, with
$N_j(X_v)$ as in Lemma~\ref{5.4:lem-cells-carrier}. It is
counted inside $\#\partial^{+}\mathrm{St}^j_v$ and controlled by \eqref{5.4:eq-zones-cells-i}.
\item A long single which runs along $\partial Z$ through many $\pi_k$-cubes has one root cell
$\Delta^{\partial}_v$ and is counted once. Its cubes cost $e^{(r+1)d_k(K_v)}$, which is carried by
$w_{r+1}(v)$.
\item The case $j_v=k$ needs no change.
\end{itemize}

\emph{Contact by touching.} Lemma~\ref{5.4:lem-cells-carrier} distinguishes the cells that meet $X_v$
from the cells that have a common point with $\bigcup X_v$. Suppose the contact of $X_v$ with a cell is
taken in the second sense, touching, rather than meeting. Then the argument goes through with the
following changes.
\begin{itemize}
\item The rim constant of part (d) of the frame depends on $d$ only, against
\eqref{5.4:eq-zones-cells-i}.
\item The label is formed with the attachment cubes $\mathrm{Att}^{T}$, whose label-length constant,
denoted $c_l^{T}$, satisfies $c_l^{T}\le c_l$.
\item The smallness factor becomes $e^{-\frac12(9\check{R}_k-3)}$.
\item Replacing \emph{$X_v$ meets $\Delta$} by \emph{$X_v$ touches $\Delta$} in the grouping costs the
factor $\hat{b}_t\,C_{\mathsf{b}}\,2^{p_{\mathsf{b}}}$.
\end{itemize}
The constants are otherwise unchanged, so that under this reading the first inequality of
\eqref{5.4:eq-long-unit-sums-1} takes the form
\begin{equation*}
\begin{aligned}
\sum_{\substack{\mathfrak{b}:\ v\in\mathfrak{L}^{\mathrm{lg}}_c\\ \mathrm{lab}(\mathfrak{b})\ni Q_*}}
\|\mathfrak{b}\|\,e^{r d_k(\mathrm{lab}\,\mathfrak{b})}
&\le e^{-\frac12(9\check{R}_k-3)}\cdot 2e\,S_d^{(1/2)}\cdot\mathfrak{K}\,\mathsf{b}^{\mathrm{lg}}(0)
\cdot\mathfrak{W}\\
&\quad\cdot\hat{b}_t\,C_{\mathsf{b}}\,2^{p_{\mathsf{b}}}.
\end{aligned}
\end{equation*}

Clause (iv) is read off
\eqref{5.4:eq-long-unit-sums-a}. Apart from the
numerical factors $2e$, $80e$, $16e$, the numbers $d$ and $L^{d-1}$, and the grid sums $S_d$ and
$S_d^{(1/2)}$ (the latter as in \eqref{5.4:eq-zones-cells-j}), each factor there is a constant fixed
earlier: $\mathfrak{K}$ and $\mathfrak{s}_{\mathsf{b}}$ in the constants of the two transports,
\eqref{5.4:eq-transport-constants-c};
$\mathsf{b}^{\mathrm{lg},\natural}$ and $C^{\mathrm{win}}_{\mathrm{lg}}$ in
\eqref{5.4:eq-transport-constants-d}; $\mathsf{c}_d$ in \eqref{5.4:eq-zones-cells-c}; and $C^{0}_{\square}$,
$\mathsf{c}_{\mathsf{f}}$, $C_{\mathsf{b}}$, $W_A$, $W_{\mathrm{own}}$ with the statements that supply
them. Each of these is, where it is fixed, independent of the quantities listed in clause (iv), and
$\check{R}$ enters these factors only through $C^{\mathrm{win}}_{\mathrm{lg}}$, through
$\sigma^{\mathrm{lg}}$ in the variant of clause (i), and through the lower bound on $\check{R}_k$ imposed
where $\check{R}_k$ is fixed; the supremum $C^{\mathrm{win}}_{B'}$ listed in clause (iv) does not occur in
\eqref{5.4:eq-long-unit-sums-a}.
Together with the two displayed exponentials and the inequality
$\mathsf{b}^{\mathrm{lg}}(0)\le\mathsf{b}^{\mathrm{lg},\natural}\check{R}_k$, this is clause (iv) for the
constants of \eqref{5.4:eq-long-unit-sums-a}; for $\mathsf{C}^{\mathrm{ex}}_1$ it is read off where its
value is fixed, as in clause (iii).
\end{proof}

\begin{proposition}[The pointed-piece sum for long units off the treated class]
\label{5.4:prf-off-class-piece-sum}
Assume the following statements:
\begin{itemize}
\item the analytic size bound for long pieces, Proposition~\ref{5.4:prop-long-piece-size};
\item the summation of boundary terms at finer zones, Theorem~\ref{5.4:thm-finer-zone-sum},
which is proved as an implication from the statements named in it,
in the form \eqref{5.4:eq-finer-zone-sum-a}, together with the standing data of the chosen
transport (the first or the second transport of Notation~\ref{5.4:not-transport-constants}):
the tree-length axiom (g4) among the
standing geometric axioms of levels, zones and cells and, when the base
bound is weighted at raised cells or cubes, the window remark for the first transport and the
weighted form of the second transport bound;
\item the cell-weight bound \eqref{5.4:eq-age-base-bounds-d};
\item the vanishing lemma for pieces off the star, the block-sum lemma, the exempt-family lemma
and the label-sum theorem, and, when the base bound carries a window weight besides the age
function, the label-sum theorem at arbitrary grain and dial.
\end{itemize}
Assume moreover that $\kappa \ge 180$ and $\lambda'_{\Sigma} \ge 5$, and, for the second
transport, that the
constant $c$ of the reserve factor $\mathfrak{r}_v = e^{-(c-2)d_v}$ satisfies $c \ge 3$.
Let $r \le a+8$ and $Q_* \in \pi_k$. Then the sum over the pieces $p$ of the long units
$v \in \mathfrak{L}^{\mathrm{lg}}_{d}$ whose label $\mathrm{lab}^{\circ}_1(p)$ contains $Q_*$
satisfies
\begin{equation}\label{5.4:eq-off-class-piece-sum-1}
\sum_{p:\ \mathrm{lab}^{\circ}_1(p) \ni Q_*} \mathfrak{N}^{\mathrm{ex}}_{A}(p)\,
e^{r d_k(\mathrm{lab}^{\circ}_1(p))}
\le \theta'(g_k)\,\mathsf{C}^{\Sigma,\mathrm{lg}},
\end{equation}
with the constant $\mathsf{C}^{\Sigma,\mathrm{lg}}$ of \eqref{5.4:eq-long-unit-sums-a}.
\end{proposition}

\begin{proof}
Throughout, $v \in \mathfrak{L}^{\mathrm{lg}}_{d}$, $j := j_v$, $\mathrm{St}_v$ is the family
of cells meeting $X_v$, $\partial\mathrm{St}_v$ is its rim and $S_v$ the union of its cells.

\smallskip\noindent
\textit{(C1) The star is exempt.} By Lemma~\ref{5.4:lem-cells-carrier}, parts (f) and (g),
\begin{equation*}
\#\mathrm{St}_v \le \mathsf{c}_d(d_v+1), \qquad [\mathrm{St}_v]_k \subset K_v, \qquad
\epsilon_{\partial}\,\#\partial\mathrm{St}_v \le \tfrac12(d_v+1),
\end{equation*}
where $K_v$ is the set attached to $v$ in Definition~\ref{5.4:def-pointed-units},
and $\mathrm{St}_v$ is wall-connected. The exempt-family lemma admits any subfamily of cells
as an exempt family; hence $\mathrm{St}_v$ is a legitimate exempt family.

\smallskip\noindent
\textit{(C2) The star carries no source.} The family $\mathrm{St}_v \cup \partial\mathrm{St}_v$
contains no source cell and does not meet $\Lambda^{\sim}$. The vanishing lemma for pieces off
the star is formulated for plain units, but its proof uses only the locality of the pieces in
the cells and the absence of source cells from the star, and both hold for the present $v$.
It therefore gives $v_p = 0$ unless the family $\mathsf{y} = \mathsf{y}(p)$ of derivative cells
of $p$ belongs to the set $\mathfrak{Y}(\mathrm{St}_v)$ of
\eqref{5.4:eq-pointed-units-d}. Only such pieces contribute to
\eqref{5.4:eq-off-class-piece-sum-1}.

\smallskip\noindent
\textit{(C3) Size and blocks.} Let $p$ be a piece with
$\mathsf{y} \in \mathfrak{Y}(\mathrm{St}_v)$. Proposition~\ref{5.4:prop-long-piece-size}
attaches to $p$ the size
\begin{equation*}
\mathfrak{N}^{\mathrm{ex}}_{A}(p)
= 2\,\mathsf{w}_L^{-1}\,(\eta_p\,\mathsf{f}_v)^{1/2}\,\mathfrak{N}^{\flat}(v)\,
e^{-(\kappa_1-1)m(p)},
\end{equation*}
where $m(p)$ is as in Proposition~\ref{5.4:prop-long-piece-size}.
Put
$\beta_{\square}(\mathrm{St}_v) := C^{0}_{\square}\,\#\mathrm{St}_v\,(\#\mathrm{St}_v+3)$.
The block-sum lemma gives
\begin{equation}\label{5.4:eq-off-class-piece-sum-2}
e^{-|\mathsf{y}|}\sum_{\square} e^{-\frac12 \delta_P\operatorname{dist}(\tilde{\square},S_v)}
\le \beta_{\square}(\mathrm{St}_v),
\end{equation}
and summing the size over the blocks $\square$ with \eqref{5.4:eq-off-class-piece-sum-2},
at the cost of one unit of the exponent $\kappa_1 - 1$, yields
\begin{equation}\label{5.4:eq-off-class-piece-sum-3}
\sum_{\square} \mathfrak{N}^{\mathrm{ex}}_{A}(p)
\le 2\,\mathsf{w}_L^{-1}\,\mu_v\, e^{-(\kappa_1-2)|\mathsf{y}|},
\qquad
\mu_v := \beta_{\square}(\mathrm{St}_v)\,\mathsf{f}_v^{1/2}\,\mathfrak{N}^{\flat}(v).
\end{equation}
We bound the three factors that will enter below. First, by (C1),
\begin{equation*}
\beta_{\square}(\mathrm{St}_v) \le C^{0}_{\square}\,\mathsf{c}_d^2\,(d_v+1)(d_v+4)
\le 5\,C^{0}_{\square}\,\mathsf{c}_d^2\, e^{d_v}.
\end{equation*}
The second inequality is $(t+1)(t+4) \le 5e^t$ for $t \ge 0$: the function
$\varphi(t) := 5e^t - t^2 - 5t - 4$ has $\varphi(0) = 1$, $\varphi'(0) = 0$ and
$\varphi''(t) = 5e^t - 2 > 0$, so $\varphi'$ is increasing, $\varphi' \ge 0$ on $[0,\infty)$,
and $\varphi \ge 1$ there. Second, $\mathsf{f}_v^{1/2} \le \mathsf{c}_{\mathsf{f}}$ by the
cell-weight bound \eqref{5.4:eq-age-base-bounds-d}. Third, by (C1),
$e^{\epsilon_{\partial}\#\partial\mathrm{St}_v} \le e^{\frac12(d_v+1)}$.

\smallskip\noindent
\textit{(C4) The label-sum theorem.} The triple
$\mathsf{u}_v := (\mathrm{St}_v, K_v, \mu_v)$ is a unit datum; let $\mathbb{D}^{\mathrm{lg}}_d$
be the family of these data. The conclusion of the label-sum theorem is a sum of four
brackets and carries the factor $4S_d\,\mathsf{M}^{\partial}_{r+1}[\mathbb{D}^{\mathrm{lg}}_d]$;
moreover $\mathsf{w}_L^{-1} = \theta'(g_k)$. Applying the theorem at $r \le a+8$ to the pieces
bounded by \eqref{5.4:eq-off-class-piece-sum-3}, and writing the prefactor
$2\mathsf{w}_L^{-1}$ as $2\theta'(g_k)$, we obtain
\begin{equation}\label{5.4:eq-off-class-piece-sum-4}
\sum_{p:\ \mathrm{lab}^{\circ}_1(p) \ni Q_*} \mathfrak{N}^{\mathrm{ex}}_{A}(p)\,
e^{r d_k(\mathrm{lab}^{\circ}_1(p))}
\le 2\theta'(g_k)\cdot 4\cdot S_d\,\mathsf{M}^{\partial}_{r+1}[\mathbb{D}^{\mathrm{lg}}_d].
\end{equation}
If a window weight is present besides the age function, we apply instead the label-sum theorem
at arbitrary grain, at grain $\mathfrak{g} = \pi^{\mathfrak{s}}$ and dial $\beta = 2$; the
proviso $\beta \le \lambda'_{\Sigma} - 2$ of that theorem holds because
$\lambda'_{\Sigma} \ge 5$. Its statement allows the grain
$\pi^{\mathfrak{s}}_j$ or $\pi^{\mathfrak{s}}$; the family of unit data is the same and so is
the sum $S_d$. Its rim supremum carries two factors $e^{-\rho}$: one is spent on the age
function, as below, and the other on the window, in (C6).

The rim supremum $\mathsf{M}^{\partial}_{r+1}[\mathbb{D}^{\mathrm{lg}}_d]$ of the label-sum
theorem is the supremum over $\Delta_1 \in \pi^{\mathfrak{s}}$ of the left member of
\eqref{5.4:eq-off-class-piece-sum-5} below; we bound it. Fix $\Delta_1 \in \pi^{\mathfrak{s}}$. If
$\Delta_1 \in \partial\mathrm{St}_v$, then by Lemma~\ref{5.4:lem-cells-carrier}(f) there is a
cell $\Delta_2 \in \mathrm{St}_v$, $\Delta_2 \not\subset Z$, sharing a wall with $\Delta_1$.
There are at most $2dL^{d-1}$ candidates for $\Delta_2$, and $X_v$ meets $\Delta_2$ because
$\Delta_2 \in \mathrm{St}_v$. Using the three bounds of (C3),
\begin{equation}\label{5.4:eq-off-class-piece-sum-5}
\begin{aligned}
&e^{-\rho(\Delta_1)} \sum_{v:\ \Delta_1 \in \partial\mathrm{St}_v}
\mu_v\, e^{(r+1)d_k(K_v)}\, e^{\epsilon_{\partial}\#\partial\mathrm{St}_v} \\
&\quad\le 5e^{1/2} C^{0}_{\square}\mathsf{c}_d^2\mathsf{c}_{\mathsf{f}}
\sum_{\Delta_2} e^{-\rho(\Delta_1)}
\sum_{v:\ X_v \text{ meets } \Delta_2} w_{r+1}(v)\, e^{2d_v}\, e^{-d_v/2} \\
&\quad\le 2dL^{d-1}\cdot 5e^{1/2} C^{0}_{\square}\mathsf{c}_d^2\mathsf{c}_{\mathsf{f}}
\cdot C_{\mathsf{b}}\,\mathfrak{s}_{\mathsf{b}}\,\mathsf{K}^{\mathrm{lg}},
\end{aligned}
\end{equation}
where $w_{r+1}(v)\,e^{2d_v}$ is the term of $v$ in the met-cell sum
$S_{r+1}(\Delta_2;\mathfrak{L}^{\mathrm{lg}})$ of \eqref{5.4:eq-pointed-units-c}. The last
inequality is obtained as follows. We discard $e^{-d_v/2} \le 1$, which is not used further.
We have $r+1 \le a+9$; since $a = \tfrac{31}{20}\kappa$, the inequality
$a+9 \le \tfrac85\kappa$ is equivalent to $\kappa \ge 180$, which is assumed. Thus
\eqref{5.4:eq-finer-zone-sum-a} at index $r+1$ and at the cell $\Delta_2$ gives
$S_{r+1}(\Delta_2;\mathfrak{L}^{\mathrm{lg}}) \le
\mathfrak{K}\,\mathsf{b}^{\mathrm{lg}}(k - \operatorname{lev}\Delta_2)$. Put
$x := k - \operatorname{lev}\Delta_2$. Lemma~\ref{5.4:lem-rim-trunk}(iii) gives
$\rho(\Delta_1) \ge x$, and Lemma~\ref{5.4:lem-poly-exp-sup} gives
\begin{equation*}
e^{-\rho(\Delta_1)}\,\mathsf{b}^{\mathrm{lg}}(x) \le e^{-x}\,\mathsf{b}^{\mathrm{lg}}(x)
\le C_{\mathsf{b}}\,\mathfrak{s}_{\mathsf{b}}\,\mathsf{b}^{\mathrm{lg}}(0).
\end{equation*}
Finally $\mathfrak{K}\,\mathsf{b}^{\mathrm{lg}}(0) = \mathsf{K}^{\mathrm{lg}}$ by
\eqref{5.4:eq-transport-constants-c}. Taking the supremum over $\Delta_1$ and replacing
$10e^{1/2}$ by $10e$,
\begin{equation}\label{5.4:eq-off-class-piece-sum-6}
\mathsf{M}^{\partial}_{r+1}[\mathbb{D}^{\mathrm{lg}}_d]
\le 10e\, dL^{d-1}\, C^{0}_{\square}\,\mathsf{c}_d^2\,\mathsf{c}_{\mathsf{f}}\,
C_{\mathsf{b}}\,\mathfrak{s}_{\mathsf{b}}\,\mathsf{K}^{\mathrm{lg}}.
\end{equation}

\smallskip\noindent
\textit{(C5) No unit of class (d) is filed in $Z$.} For the first transport there is nothing to
prove. For the second transport the alternative for units filed at cells contained in $Z$ is
void: every
$v \in \mathfrak{L}^{\mathrm{lg}}_{d}$ is filed at $\Delta'_v \supseteq q_v$ with
$q_v \not\subset Z$, so $\Delta'_v \not\subset Z$, and only cells not contained in $Z$ occur.

\smallskip\noindent
\textit{(C6) The window weight.} This step is needed only if the base bound is weighted at
raised cells or cubes; otherwise $C^{\mathrm{win}}_{\mathrm{lg}} := 0$. In what follows
$\mathfrak{a}''$ denotes a live arrest, $X'' := X_{\mathfrak{a}''}$ its component and
$k_{a''}$ its level, and we write $\check{R} := \check{R}_{k_{a''}}$.

For the first transport, the window remark shows that the weight appears only when
$\Delta_2 \subset W_{\mathfrak{a}''}$ for a live arrest $\mathfrak{a}''$. There the collar
lemma, Lemma~\ref{5.4:lem-collar}, gives $\rho(\Delta_1) \ge 2\check{R} - 1$. With the second
factor $e^{-\rho(\Delta_1)}$ of the dial $\beta = 2$, the weight is bounded by
$C^{\mathrm{win}}_{\mathrm{lg}}$ of \eqref{5.4:eq-transport-constants-d}.

For the second transport, where $c \ge 3$ by assumption, the second transport bound attaches to
the term of
$v$ the factor $W_v\,\mathfrak{r}_v$, with $\mathfrak{r}_v = e^{-(c-2)d_v} \le e^{-d_v}$. Here
$W_v := W^{\mathrm{win}}(\check{R})$ if $\Delta'_v$ is raised by a live arrest
$\mathfrak{a}''$, in the sense of \eqref{5.4:eq-zones-cells-g}, and $W_v := 1$ otherwise. By
(C5), $\Delta'_v \not\subset Z$, so such an $\mathfrak{a}''$ is foreign or a sibling. It
suffices to show
$e^{-\rho(\Delta_1)}\,W_v\,\mathfrak{r}_v \le 1 + C^{\mathrm{win}}_{\mathrm{lg}}$. If
$W_v = 1$ the bound is immediate, since $\rho(\Delta_1) \ge x \ge 0$ by
Lemma~\ref{5.4:lem-rim-trunk}(iii) and $\mathfrak{r}_v \le e^{-d_v} \le 1$. Suppose that
$\Delta'_v$ is raised by $\mathfrak{a}''$.
\begin{itemize}
\item[$(\alpha)$] If $\Delta_2 \subset W_{\mathfrak{a}''}$, the collar lemma,
Lemma~\ref{5.4:lem-collar}, gives $\rho(\Delta_1) \ge 2\check{R} - 1$.
\item[$(\beta)$] Otherwise let $\delta := \operatorname{dist}_{k_{a''}}(\Delta_2,(X'')^{c})$.
Then $\delta = 0$ if $\Delta_2$ does not meet $X''$; otherwise $\Delta_2 \subset X''$ and
$\Delta_2$ has diameter at most one, by the collar-zone axiom of
\eqref{5.4:eq-zones-cells-g}. The carrier $X_v$ contains $q_v$, and
$q_v \subset \Delta'_v \subset W_{\mathfrak{a}''}$, at distance at least $2\check{R}$ from
$(X'')^{c}$. It also contains a cube inside $\Delta_2$, every point of which lies within
$\delta+1$ of $(X'')^{c}$. Since $j_v \le \operatorname{lev}\Delta'_v \le k_{a''}$,
the tree-length axiom (g4), applied at level $k_{a''}$ to these
two cubes of $X_v$, gives
$d_v \ge (2\check{R} - \delta - 2)_+$. Moreover
$\rho(\Delta_1) \ge \max\{\delta-1,1\}$, by \eqref{5.4:eq-zones-cells-d} and, for
$\delta \ge 2$, by the chain of the proof of the collar lemma started from depth at least
$\delta - 1$.
\end{itemize}
For $\delta \ge 2$ the sum $(2\check{R} - \delta - 2)_+ + \max\{\delta-1,1\}$ is at least
$2\check{R} - 3$. For $\delta < 2$ it is at least $2\check{R} - 1 - \delta > 2\check{R} - 3$.
Hence
\begin{equation*}
\min_{\delta \ge 0}\bigl[(2\check{R} - \delta - 2)_+ + \max\{\delta-1,1\}\bigr]
\ge 2\check{R} - 3 \ge \frac{2\check{R}-2}{\mathsf{c}_d} - 1.
\end{equation*}
In both cases $\rho(\Delta_1) + d_v \ge (2\check{R}-2)/\mathsf{c}_d - 1$. With
$\mathfrak{r}_v \le e^{-d_v}$ this gives
\begin{equation*}
e^{-\rho(\Delta_1)}\,W_v\,\mathfrak{r}_v
\le W^{\mathrm{win}}(\check{R})\, e^{1 - (2\check{R}-2)/\mathsf{c}_d}
\le C^{\mathrm{win}}_{\mathrm{lg}}.
\end{equation*}
Nested arrests need no count: the inequality $\rho(\Delta_1) \ge \delta - 1$ holds
simultaneously for every arrest whose component contains $\Delta_2$.

\smallskip\noindent
\textit{(C7) The total.} In the weighted case each $v$-term of
\eqref{5.4:eq-off-class-piece-sum-5} carries in addition the factor
$e^{-\rho(\Delta_1)}W_v\mathfrak{r}_v$, which is at most $1 + C^{\mathrm{win}}_{\mathrm{lg}}$
by (C6); by (C5) no factor for units filed at cells contained in $Z$ occurs. Inserting
\eqref{5.4:eq-off-class-piece-sum-6}, with this factor, into
\eqref{5.4:eq-off-class-piece-sum-4}, the left member of
\eqref{5.4:eq-off-class-piece-sum-1} is at most
\begin{equation}\label{5.4:eq-off-class-piece-sum-7}
\theta'(g_k)\cdot 8S_d\cdot 10e\, dL^{d-1}\, C^{0}_{\square}\,\mathsf{c}_d^2\,
\mathsf{c}_{\mathsf{f}}\,(1 + C^{\mathrm{win}}_{\mathrm{lg}})\,
C_{\mathsf{b}}\,\mathfrak{s}_{\mathsf{b}}\cdot \mathfrak{K}\,\mathsf{b}^{\mathrm{lg}}(0).
\end{equation}
By \eqref{5.4:eq-transport-constants-d}, the $\check{R}$-free constant
$\mathsf{b}^{\mathrm{lg},\natural}$
satisfies $\mathsf{b}^{\mathrm{lg}}(0) \le \mathsf{b}^{\mathrm{lg},\natural}\,\check{R}_k$. The
dials used are $\beta \in \{1,2\}$, and the proviso $\beta \le \lambda'_{\Sigma} - 3$ of
Lemma~\ref{5.4:lem-rim-trunk}(v) holds for both because $\lambda'_{\Sigma} \ge 5$. So
Lemma~\ref{5.4:lem-rim-trunk}, parts (iv) and (v), applies: the bound
\eqref{5.4:eq-off-class-piece-sum-7} improves by the
prefactor $e^{-(9\check{R}_k - 2)}$, and no constant changes. The function $t e^{-9t}$ has
maximum $1/(9e)$ on $[0,\infty)$, attained at $t = 1/9$. Hence
\begin{equation*}
\check{R}_k\, e^{-(9\check{R}_k - 2)} = e^2\,\check{R}_k\, e^{-9\check{R}_k}
\le \frac{e}{9} \le 1,
\end{equation*}
with no lower bound on $\check{R}_k$, and therefore
$e^{-(9\check{R}_k - 2)}\,\mathsf{b}^{\mathrm{lg}}(0) \le \mathsf{b}^{\mathrm{lg},\natural}$.
The bound \eqref{5.4:eq-off-class-piece-sum-7} becomes
\begin{equation*}
\theta'(g_k)\cdot 80e\, dL^{d-1} S_d\, C^{0}_{\square}\,\mathsf{c}_d^2\,
\mathsf{c}_{\mathsf{f}}\,(1 + C^{\mathrm{win}}_{\mathrm{lg}})\,
C_{\mathsf{b}}\,\mathfrak{s}_{\mathsf{b}}\,\mathfrak{K}\,\mathsf{b}^{\mathrm{lg},\natural}
= \theta'(g_k)\,\mathsf{C}^{\Sigma,\mathrm{lg}},
\end{equation*}
with $\mathsf{C}^{\Sigma,\mathrm{lg}}$ as in \eqref{5.4:eq-long-unit-sums-a}. This is
\eqref{5.4:eq-off-class-piece-sum-1}.
\end{proof}

\begin{remark}[The filing column inside the treated class]\label{5.4:rem-in-class-column}
We work on the route in which the base bounds are the age-function bounds
\eqref{5.4:eq-age-base-bounds-c}, with $c\ge 2$, and in which the transport is the per-cell
transport (Proposition~\ref{5.4:prop-per-cell-transport}). Take
$\Delta_2:=\Delta^{\partial}$, the anchor cell of Theorem~\ref{5.4:thm-finer-zone-sum}: the
level-$k$ cell not contained in $Z$ and lying within
distance $2\check{R}_k$ of $\partial Z$ that carries the charge. At this cell the per-cell
transport sums the filing columns $\mathrm{Col}^{\mathrm{file}}(\Delta')$ over the filing
cells $\Delta'$ of the long units that meet $\Delta^{\partial}$. For
$v\in\mathfrak{L}^{\mathrm{lg}}_{c}$ the filing cell $y_v$ may lie inside $Z$, in one of the
following places:
\begin{itemize}
\item in a level-$k$ cell of $Z\setminus\Lambda^{\sim}$, which is not raised if it lies within
$2\check{R}_k$ of $\partial Z$;
\item in a hub;
\item in a core;
\item in the window of its own live arrest $\mathfrak{a}$.
\end{itemize}
In the last case $X_{\mathfrak{a}}\subset\mathsf{C}_C\subset\Lambda_C$.
The per-cell transport is independent of $Z$. The only question is therefore whether the base
bounds deliver the long and tail column at a filing cell in $Z$ in the form
\eqref{5.4:eq-age-base-bounds-c}, and with which weight factor $\mathfrak{W}$. Two
alternatives are considered.
\begin{itemize}
\item[(A)] The column in $Z$ is weight-free, or it carries only the un-raised cell weight
$\bar{\mathsf{w}}=3$; in the second case the column is bounded with the factor
$(3/2)3^{p}$, where $p$ is the exponent to which the cell weight $\bar{\mathsf{w}}$ is raised in
that bound. Then $\mathfrak{W}=W_A$ with $W_A\in\{1,(3/2)3^{p}\}$. If the base bounds for long
units carry no factor $\bar{\mathsf{w}}$, then $W_A=1$.
\item[(B)] The column in $Z$ carries the weights of the own windows. Then
$\mathfrak{W}=W_{\mathrm{own}}$, where $W_{\mathrm{own}}:=1+C^{\mathrm{win}}_{\mathrm{lg}}$,
and $c\ge 3$ is required.
\end{itemize}
Alternative (A) is taken as a hypothesis wherever it is invoked.
On the other route the question does not arise, by
Corollary~\ref{5.4:cor-current-level-bound}; see the end of the proof.
\end{remark}

\begin{proof}
The per-cell transport (Proposition~\ref{5.4:prop-per-cell-transport}) gives, for every cell
$\Delta_2$ of level $\ell\le k$ and every $r'\le(8/5)\kappa$,
\begin{equation*}
S_{r'}(\Delta_2;\mathfrak{L}^{\mathrm{lg}})\le
C^{\mathrm{gen}}_{\mathsf{b}}\,\mathsf{b}^{\mathrm{lg}}(k-\ell)\,
\mathsf{K}^{\flat}_{\mathrm{lg}} .
\end{equation*}
Its constant $C^{\mathrm{gen}}_{\mathsf{b}}$ depends only on $d$, $L$, $\mathsf{c}_d$,
$C_{\mathsf{b}}$ and $p_{\mathsf{b}}$, and not on $Z$. The proposition therefore does not
depend on $Z$. Applied at $\Delta_2=\Delta^{\partial}$, the only input that distinguishes filing cells
in $Z$ from filing cells off $Z$ is the column supplied by the base bounds at those cells.
This is the weight question stated above.

Two of the four places listed for $y_v$ call for comment.
\begin{itemize}
\item A level-$k$ cell of $Z\setminus\Lambda^{\sim}$ within $2\check{R}_k$ of $\partial Z$ is
not raised, by the geometric axiom \eqref{5.4:eq-zones-cells-f}.
\item If $y_v$ lies in the window of its own live arrest $\mathfrak{a}$, the component of that
arrest satisfies $X_{\mathfrak{a}}\subset\mathsf{C}_C\subset\Lambda_C$; this inclusion, for every
live arrest, is part of the representation and bound of the large-field terms.
\end{itemize}
Under (A) the weight attached to a filing cell in $Z$ is either $1$ or comes from
$\bar{\mathsf{w}}=3$. These are the two values of $W_A$. Under (B) the weight attached is
that of the own window, which enters through $1+C^{\mathrm{win}}_{\mathrm{lg}}$.

On the other route, the long-unit bound at a cell of the current level
(Corollary~\ref{5.4:cor-current-level-bound}) gives, at every level-$k$ cell $\Delta_2$,
\begin{equation*}
S_{r'}(\Delta_2;\mathfrak{L}^{\mathrm{lg}}_{c})\le
S_{r'}(\Delta_2;\mathfrak{L}^{\mathrm{lg}})\le\mathsf{K}^{\mathrm{lg}} .
\end{equation*}
This holds in particular at $\Delta_2=\Delta^{\partial}$. The bound requires no filing column
at a cell of $Z$, so no weight $\mathfrak{W}$ enters it.
\end{proof}

\begin{lemma}[Own windows of live arrests]\label{5.4:lem-own-windows}
Let $k$ be the current level and $C$ the treated component. Assume the following.
\begin{enumerate}
\item[(a)] The label changes only through the events listed in
\eqref{5.4:eq-zones-cells-h}.
\item[(b)] For every live arrest $\mathfrak{a}'$, every preparation attempt $b$ after step
$k_{a'}$ over a component containing $X_{\mathfrak{a}'}$ has range $]h_b,k_b]$ with
$h_b\ge k_{a'}$ (the rule on arrests stated in \eqref{5.4:eq-zones-cells-h}).
\item[(c)] The definition of the window and the collar of a live arrest in
\eqref{5.4:eq-zones-cells-g}; and, from the definition of an arrest, that for a live arrest
$\mathfrak{a}$ the set $X_{\mathfrak{a}}$ is at step $k_a$ a component of $(\Lambda_{k_a})^c$,
with $X_{\mathfrak{a}}\subset Z_{k_a}$ in the unprimed label. The window is
$W_{\mathfrak{a}}:=X_{\mathfrak{a}}\cap(\Omega^{\mathbf{T}}_{k_a})^c$, and these sets are
fixed at step $k_a$.
\item[(d)] The axioms (A0)--(A2) of \eqref{5.4:eq-zones-cells-b}, in the current label.
\item[(e)] For every live own arrest $\mathfrak{a}$ of $C$,
$X_{\mathfrak{a}}\subset\mathsf{C}_C\subset\Lambda_C$.
\item[(f)] The geometric axiom \eqref{5.4:eq-zones-cells-e}; the radius inequality
$\check{R}_{n-1}\le L\check{R}_n$ for every level $n$; and the separation of distinct
components of $\Lambda_n^c$ by at least $\check{R}_n$; the last two as in
\eqref{5.4:eq-zones-cells-g}.
\end{enumerate}
Let $\mathfrak{a}$ be a live own arrest of $C$ with $k_a<k$. Live means that no operation
$\mathbf{R}$ over a component containing $X:=X_{\mathfrak{a}}$ has been completed since
step $k_a$. Put $W:=W_{\mathfrak{a}}$. Then the following hold.
\begin{enumerate}
\item[(i)] $X\subset Z^{\mathrm{cur}}_{k_a}$, where
$Z^{\mathrm{cur}}_{k_a}:=(\Lambda^{\mathrm{cur}}_{k_a})^c$ as in the convention of
Lemma~\ref{5.2:lem-current-label-anchors}, and $X\cap\Omega^{\mathrm{cur}}_{k_a+1}=\emptyset$.
Every cell meeting $X$ lies in $X$ and has level $\le k_a$.
\item[(ii)] The window is separated from the complement of $X$:
\begin{equation}\label{5.4:eq-own-windows-1}
\operatorname{dist}_{k_a}(W,X^c)\ge 2\check{R}_{k_a}.
\end{equation}
This bound involves only the sets fixed at step $k_a$ and does not depend on the current
label.
\item[(iii)] $X\subset\Lambda^{\sim}_C$. Thus the own windows of earlier arrests are never
cut, and a filing cell there is an index, not a cut cell.
\end{enumerate}
\end{lemma}

\begin{proof}
\emph{Clause (i).} By hypothesis (c), at step $k_a$ the set $X$ is a component of
$(\Lambda_{k_a})^c$, and $X\subset Z_{k_a}$ in the unprimed label. The preparation of
$\mathfrak{a}$ re-labels only the sets $Z''_j$, $\Omega''_j$ of its prepared label, and only
inside its class. We
show that no event after the exit of $\mathfrak{a}$ changes the label at any level
$j\le k_a$ inside $X$. By the list of label-changing events (hypothesis (a)), the label
changes only through the events listed there, which we take in turn.
\begin{itemize}
\item A $\mathbf{T}$-step $m\to m+1$ defines the new top pair
$(\Omega_{m+1},\Lambda_{m+1})$ and leaves $(\Omega_i,\Lambda_i)_{i\le m}$ unchanged. Since
it occurs after step $k_a$, it touches only levels $>k_a$.
\item Let $b$ be a later attempt over a component containing $X$. By hypothesis (b), $b$
has $h_b\ge k_a$. Its preparation leaves the levels $j\le h_b$ untouched:
$\Omega''_j=\Omega_j$ for $j\le h_b$ by \cite[(1.10)--(1.12)]{B-CMP122a}, as listed in
that display, and $Z''_j=Z_j$ for $j\le h_b$. It therefore leaves
every level $\le k_a$ unchanged.
\item Attempts and completions on components disjoint from $X$ act inside those
components. For preparations this is the intersection with the class in
\cite[(1.12)]{B-CMP122a}; for completions it is the prescription on $Z$ in
\cite[(1.33)]{B-CMP122b}. Hence they do not change the label inside $X$.
\item A completed operation $\mathbf{R}$ over a component containing $X$ is excluded,
because $\mathfrak{a}$ is live.
\item Arrests, exclusions and freezing re-label no frozen level, by the last item of
\eqref{5.4:eq-zones-cells-h}; in particular they leave the levels $\le k_a$ inside $X$
unchanged.
\end{itemize}
Hence the set $\Lambda_{k_a}$ inside $X$ is that of step $k_a$: the preparation of
$\mathfrak{a}$ re-labels only $Z''_j$, $\Omega''_j$, and no later event re-labels a level
$\le k_a$ inside $X$. Since $X\subset(\Lambda_{k_a})^c$ at that step, we get
\begin{equation*}
\Lambda^{\mathrm{cur}}_{k_a}\cap X=\emptyset .
\end{equation*}
With the convention $Z_j=\Lambda_j^c$ of Lemma~\ref{5.2:lem-current-label-anchors}, taken
in the current label, this is the inclusion $X\subset Z^{\mathrm{cur}}_{k_a}$ of clause (i).
Axiom (A0) of \eqref{5.4:eq-zones-cells-b}, in the current label (hypothesis (d)), gives
$\Omega^{\mathrm{cur}}_{k_a+1}\subset\Lambda^{\mathrm{cur}}_{k_a}$. Therefore
$X\cap\Omega^{\mathrm{cur}}_{k_a+1}=\emptyset$.

This is the argument behind the statement in \eqref{5.4:eq-zones-cells-g} on live arrests
over components not in $Z$, transposed from such a component to $C$. In that argument the
hypothesis that the component does not lie in $Z$ served only to keep the running attempt
off the list of events that re-label levels $\le k_a$. Here hypothesis (b) does this.

Now consider the cells. A cell of level $m\ge k_a+1$ lies in
$\Omega^{\mathrm{cur}}_{k_a+1}$, which is the union of the current cells of level
$\ge k_a+1$. Since $X\cap\Omega^{\mathrm{cur}}_{k_a+1}=\emptyset$, such a cell does not
meet $X$. The set $X$ is a union of $\pi_{k_a}$-cubes. Hence, by nesting, a cell of level
$\le k_a$ that meets $X$ lies in $X$. So every cell meeting $X$ lies in $X$ and has level
$\le k_a$.

\emph{Clause (ii).} In this part we write $\Omega:=\Omega^{\mathbf{T}}_{k_a}$ and
$\Lambda:=\Lambda^{\mathbf{T}}_{k_a}\subset\Omega$ (the sets $\Omega_{k_a}$,
$\Lambda_{k_a}$ of the label at step $k_a$, which enter hypothesis (c)). Distances $|x-y|$
and $\operatorname{dist}(\cdot,\cdot)$ are measured as in $\operatorname{dist}_{k_a}$. By
\cite{B-CMP119}, $\Omega$ and $\Lambda$ are unions of $M\check{R}_{k_a}$-cubes with
\begin{equation}\label{5.4:eq-own-windows-2}
\operatorname{dist}_{k_a}(\partial\Omega,\partial\Lambda)\ge 2\check{R}_{k_a}.
\end{equation}

\emph{First step.} Let $w$ lie in the closure of $\Omega^c$ and let $\lambda\in\Lambda$. Let
$z_1$ be the last point of the segment $[w,\lambda]$ lying in the closure of $\Omega^c$.
Then $z_1\in\partial\Omega$.

We check that $z_1$ is not in the closure of $\Lambda$. Since $z_1$ is in the closure of
$\Omega^c$, it is not in the interior of $\Omega$. The interior of $\Lambda$ is contained
in the interior of $\Omega$, so $z_1$ is not in the interior of $\Lambda$. If $z_1$ were in
the closure of $\Lambda$, it would therefore lie in $\partial\Lambda\cap\partial\Omega$.
This contradicts \eqref{5.4:eq-own-windows-2}.

Let $z_2$ be the first point of $[z_1,\lambda]$ in the closure of $\Lambda$. Since $z_1$ is
not in that closure, $z_2\in\partial\Lambda$. The points $z_1,z_2$ lie on $[w,\lambda]$ in
this order. Hence, by \eqref{5.4:eq-own-windows-2},
\begin{equation*}
|w-\lambda|\ge|z_1-z_2|\ge 2\check{R}_{k_a}.
\end{equation*}
Thus $\operatorname{dist}(w,\Lambda)\ge 2\check{R}_{k_a}$ for every $w$ in the closure of
$\Omega^c$.

\emph{Second step.} Let $w\in W$ and $y\notin X$. By hypothesis (c),
$W=X\cap\Omega^c$, so $W$ lies in the closure of $\Omega^c$. Let $z$ be the first point of
$[w,y]$ outside the interior of $X$; then $z\in\partial X$. Points of $[w,y]$ just beyond
$z$ lie in $\Lambda$. Indeed, by hypothesis (c) $X$ is a component of $\Lambda^c$, and by
the separation in
hypothesis (f) any other component of $\Lambda^c$ is at distance
$\ge\check{R}_{k_a}$ from $X$. Hence $z$ lies in the closure of $\Lambda$. By the first
step,
\begin{equation*}
|w-y|\ge|w-z|\ge\operatorname{dist}(W,\Lambda)\ge 2\check{R}_{k_a}.
\end{equation*}
Taking the infimum over $w\in W$ and $y\notin X$ gives \eqref{5.4:eq-own-windows-1}. The
sets $X$, $W$, $\Omega$, $\Lambda$ are those of hypothesis (c), fixed at step $k_a$ and
never redefined. Hence the bound does not depend on the current label.

\emph{Clause (iii).} By hypothesis (e), $X\subset\mathsf{C}_C\subset\Lambda_C$. Since
$\Lambda_C\subset\Lambda^{\sim}_C$, it follows that $X\subset\Lambda^{\sim}_C$.
\end{proof}

\begin{proposition}[Own-window weights paid by the unit's own length]\label{5.4:prop-own-window-weights}
Let $C$ be the treated component, $Z$ the current large-field region, and $\Delta^{\partial}$ a cell of
level $k$ with $\Delta^{\partial}\not\subset Z$ which touches $Z$. Consider the long units
$v\in\mathfrak{L}^{\mathrm{lg}}_{c}$ with $X_v$ meeting $\Delta^{\partial}$ whose filing cell satisfies
$\Delta'_v\subset W_{\mathfrak{a}}$ for a live own arrest $\mathfrak{a}$ of $C$ of level $k_a<k$; here
$\mathfrak{a}$ is the arrest named by the weight at $\Delta'_v$ (nested arrests have disjoint ranges, and
the weight names one of them). Assume that at every such cell the filing column satisfies the weighted
bound of Lemma~\ref{5.4:lem-weighted-columns} with
\begin{equation*}
G(\Delta'_v)\le W^{\mathrm{win}}(\check{R}_{k_a}),
\end{equation*}
$W^{\mathrm{win}}$ as in \eqref{5.4:eq-transport-constants-d}, and that $c\ge 3$ in
\eqref{5.4:eq-age-base-bounds-c}. Then for each such $v$
\begin{equation}\label{5.4:eq-own-window-weights-1}
G(\Delta'_v)\,\mathfrak{r}_v\le C^{\mathrm{win}}_{\mathrm{lg}},
\end{equation}
where $\mathfrak{r}_v:=e^{-(c-2)d_v}$, and, for $r'\le(8/5)\kappa$ as in
Proposition~\ref{5.4:prop-per-cell-transport}, the met-cell sum
$S(\Delta^{\partial};\mathfrak{L}^{\mathrm{lg}}_{c})=S_{r'}(\Delta^{\partial};\mathfrak{L}^{\mathrm{lg}}_{c})$
of \eqref{5.4:eq-pointed-units-c} satisfies
\begin{equation}\label{5.4:eq-own-window-weights-2}
S(\Delta^{\partial};\mathfrak{L}^{\mathrm{lg}}_{c})
\le W_{\mathrm{own}}\,G^{\mathrm{oth}}\,
C^{\mathrm{gen}}_{\mathsf{b}}\,\mathsf{b}^{\mathrm{lg}}(0)\,\mathsf{K}^{\flat}_{\mathrm{lg}},
\end{equation}
where $W_{\mathrm{own}}:=1+C^{\mathrm{win}}_{\mathrm{lg}}$ and $G^{\mathrm{oth}}$ denotes the factor
contributed by the remaining weight factors carried by $G$ in Lemma~\ref{5.4:lem-weighted-columns},
other than the own-window weights bounded in \eqref{5.4:eq-own-window-weights-1}. The factor
$W_{\mathrm{own}}$ in \eqref{5.4:eq-own-window-weights-2} contains no age and no dependence on $k-k_a$.
\end{proposition}

\begin{proof}
Fix $v$ and $\mathfrak{a}$ as in the statement and write $X:=X_{\mathfrak{a}}$, $W:=W_{\mathfrak{a}}$,
$j_v$ for the level of $v$, $y_v$ for its point, and $q_v$ for the cube of $X_v$ with
$\Delta'_v=\Delta(q_v)$ (Lemma~\ref{5.4:lem-cells-carrier}\,(c)). For a level $m$,
$\operatorname{dist}_m$ denotes distance measured in level-$m$ units. The arrest $\mathfrak{a}$ is live,
so Lemma~\ref{5.4:lem-own-windows} applies to it; we use its clauses (i), (ii) and (iii).

\emph{Step 1: levels.} By Lemma~\ref{5.4:lem-own-windows}\,(ii),
$\operatorname{dist}_{k_a}(W,X^c)\ge 2\check{R}_{k_a}>0$, so $W\subset X$, and the cell
$\Delta'_v\subset W$ meets $X$. By Lemma~\ref{5.4:lem-own-windows}\,(i) every cell meeting $X$ has level
at most $k_a$, hence $\operatorname{lev}\Delta'_v\le k_a$. Since
$\operatorname{lev}\Delta'_v=\operatorname{lev}\Delta(q_v)\ge j_v$ by
Lemma~\ref{5.4:lem-cells-carrier}\,(c), also $j_v\le k_a$.

\emph{Step 2: the collar bound.} Put $g:=\operatorname{dist}_{k_a}(\cdot\,,X^c)$; this function is
$1$-Lipschitz for $\operatorname{dist}_{k_a}$. By Lemma~\ref{5.4:lem-own-windows}\,(ii) and
$q_v\subset\Delta'_v\subset W$ we have $g\ge 2\check{R}_{k_a}$ on $q_v$. Next, $X\subset Z$
while $\Delta^{\partial}\not\subset Z$, so
$\Delta^{\partial}\not\subset X$; since by Lemma~\ref{5.4:lem-own-windows}\,(i) a cell meeting $X$ lies
in $X$, the cell
$\Delta^{\partial}$ does not meet $X$, and $g\equiv 0$ on $\Delta^{\partial}$. As $X_v$ meets
$\Delta^{\partial}$, Lemma~\ref{5.4:lem-cells-carrier}\,(b) provides a cube $Q_{\partial}$ of $X_v$ with
$Q_{\partial}\subset\Delta^{\partial}$, and $g\equiv0$ on $Q_{\partial}$. The Lipschitz property gives
$\operatorname{dist}_{k_a}(q_v,Q_{\partial})\ge 2\check{R}_{k_a}$. Since $j_v\le k_a$ (Step~1),
passing to level-$j_v$ units multiplies distances by $L^{k_a-j_v}\ge 1$:
\begin{equation*}
\operatorname{dist}_{j_v}(q_v,Q_{\partial})\ge L^{k_a-j_v}\cdot 2\check{R}_{k_a}\ge 2\check{R}_{k_a}.
\end{equation*}
Both $q_v$ and $Q_{\partial}$ are cubes of $X_v$, so (g4) in \eqref{5.4:eq-zones-cells-c} yields
$d_v\ge 2\check{R}_{k_a}$. If only (g1) in
\eqref{5.4:eq-zones-cells-c} is used, the same distance bound yields
$\mathsf{c}_d(d_v+1)\ge 2\check{R}_{k_a}$.

\emph{Step 3: a bound free of the collar.} This step gives a second lower bound on $d_v$, which does not
use the bound $g\ge 2\check{R}_{k_a}$ on $q_v$ of Step~2; either lower bound suffices for Step~4, which
uses only $d_v\ge 2\check{R}_{k_a}$. By Lemma~\ref{5.4:lem-own-windows}\,(iii),
$y_v\in X\subset\Lambda^{\sim}_C$; moreover $\operatorname{dist}_k(\Lambda^{\sim}_C,Z^c)\ge 9\check{R}_k$.
Hence $\operatorname{dist}_k(q_v,Q_{\partial})\ge 9\check{R}_k-1$, and this is at least
$8\check{R}_k$ by the condition on the radii in \eqref{5.4:eq-zones-cells-g}. In level-$j_v$ units,
using $j_v\le k_a<k$ and then the comparison $\check{R}_kL^{k-k_a}\ge\check{R}_{k_a}$ of the radii in
\eqref{5.4:eq-zones-cells-g},
\begin{equation*}
\operatorname{dist}_{j_v}(q_v,Q_{\partial})\ge 8\check{R}_k L^{k-j_v}
\ge 8\check{R}_k L^{k-k_a}\ge 8\check{R}_{k_a},
\end{equation*}
and, as in Step~2, $d_v\ge 8\check{R}_{k_a}$ (through (g1) alone,
$\mathsf{c}_d(d_v+1)\ge 8\check{R}_{k_a}$).

\emph{Step 4: the weight.} Since $c\ge 3$, $\mathfrak{r}_v=e^{-(c-2)d_v}\le e^{-d_v}$. With the hypothesis
on $G(\Delta'_v)$ and $d_v\ge 2\check{R}_{k_a}$ from Step~2,
\begin{equation*}
G(\Delta'_v)\,\mathfrak{r}_v\le W^{\mathrm{win}}(\check{R}_{k_a})\,e^{-d_v}
\le W^{\mathrm{win}}(\check{R}_{k_a})\,e^{-2\check{R}_{k_a}}\le C^{\mathrm{win}}_{\mathrm{lg}},
\end{equation*}
the last inequality by the definition of $C^{\mathrm{win}}_{\mathrm{lg}}$ in
\eqref{5.4:eq-transport-constants-d}. This is \eqref{5.4:eq-own-window-weights-1}.
If only the comparison (g1) of \eqref{5.4:eq-zones-cells-c} is used, Step~2 gives
$\mathsf{c}_d(d_v+1)\ge 2\check{R}_{k_a}$, hence
$e^{-d_v}\le e^{1-2\check{R}_{k_a}/\mathsf{c}_d}$, the factor $e^{-2\check{R}_{k_a}}$ above is
replaced by $e^{1-2\check{R}_{k_a}/\mathsf{c}_d}$, and
\begin{equation*}
W^{\mathrm{win}}(\varrho)\,e^{1-2\varrho/\mathsf{c}_d}\big|_{\varrho=\check{R}_{k_a}}
\le C^{\mathrm{win}}_{\mathrm{lg}}
\end{equation*}
again by the definition of $C^{\mathrm{win}}_{\mathrm{lg}}$ in \eqref{5.4:eq-transport-constants-d}.

The arrest $\mathfrak{a}$ is determined by the cell $\Delta'_v$ (it is the one named by the weight
there), so the bound holds for every $v$ filed at $\Delta'_v$ with $X_v$ meeting $\Delta^{\partial}$,
and therefore, with $\mathfrak{r}(\Delta';\Delta_2)$ as in Lemma~\ref{5.4:lem-weighted-columns},
\begin{equation*}
G(\Delta'_v)\,\mathfrak{r}(\Delta'_v;\Delta^{\partial})\le C^{\mathrm{win}}_{\mathrm{lg}}
\le 1+C^{\mathrm{win}}_{\mathrm{lg}}.
\end{equation*}
With this bound on $G\,\mathfrak{r}$ at the filing cells lying in windows of live own arrests of $C$,
and with the remaining weight factors carried by $G$ contributing the factor $G^{\mathrm{oth}}$, the right
side of the bound of Lemma~\ref{5.4:lem-weighted-columns} for $\Delta_2=\Delta^{\partial}$ is at most
\begin{equation*}
W_{\mathrm{own}}\,G^{\mathrm{oth}}\,e^{2}\,\mathsf{K}^{\flat}_{\mathrm{lg}}
\sum_{\Delta'}\mathsf{b}^{\mathrm{lg}}(k-\operatorname{lev}\Delta')\,e^{-\theta m(\Delta')},
\end{equation*}
with $\theta$ as in \eqref{5.4:eq-transport-constants-b} and $m(\Delta')$ as in
Lemma~\ref{5.4:lem-weighted-columns}. The quantity
$e^{2}\mathsf{K}^{\flat}_{\mathrm{lg}}\sum_{\Delta'}\mathsf{b}^{\mathrm{lg}}(k-\operatorname{lev}\Delta')
e^{-\theta m(\Delta')}$ is bounded, exactly as in the proof of the per-cell transport
(Proposition~\ref{5.4:prop-per-cell-transport}) with $\Delta_2=\Delta^{\partial}$, a cell of level
$\ell=k$, by $C^{\mathrm{gen}}_{\mathsf{b}}\,\mathsf{b}^{\mathrm{lg}}(0)\,\mathsf{K}^{\flat}_{\mathrm{lg}}$.
This gives \eqref{5.4:eq-own-window-weights-2}.

Thus the length of the unit across the collar of that arrest pays the weight, in the arrest's own scale
$\check{R}_{k_a}$. No age and no $k-k_a$ enters: the factors $L^{k_a-j_v}$ and $L^{k-k_a}$ were used
only through $L^{k_a-j_v}\ge 1$ in Step~2 and through the comparison of the radii in Step~3, and were
then dropped.
Finally, completed earlier attempts carry no weight: they are healed and have $w=1$, since, by
clause (d) of the window-weight lemma of this chapter,
a cell weight exceeds $1$ only on windows of live arrests.
\end{proof}

\begin{remark}[The running-window weight and the residual case]\label{5.4:rem-running-window}
Let $C$ be the treated component, with core $\mathsf{C}_C$, let $k$ be the current level, and let
$\Delta^{\partial}$ be a level-$k$ cell with $\Delta^{\partial}\not\subset Z$ that touches $Z$. Write
$S(\Delta^{\partial};\mathfrak{L}^{\mathrm{lg}}_{c})$ for the sum of
Lemma~\ref{5.4:lem-weighted-columns} at the cell $\Delta^{\partial}$, restricted to the long singles
$\mathfrak{L}^{\mathrm{lg}}_{c}$ of class~(c) in the classification of long units.
\begin{enumerate}
\item[(a)] Consider the running preparation attempt. Suppose that the in-$Z$ filing column carries, at
the running-window cells, a weight at most $W^{\mathrm{win}}(\check{R}_k)$, with index $k$. Then the
long-unit member of the boundary-term sum of Corollary~\ref{5.4:cor-long-unit-sums}(i) is at most
\begin{equation}\label{5.4:eq-running-window-1}
\theta'(g_k)\,\mathsf{c}_{\theta'}\cdot 2e\,S_d^{(1/2)}\,W_{\mathrm{own}}\,\mathfrak{K}\,
\mathsf{b}^{\mathrm{lg},\natural}\,\sigma^{\mathrm{lg}},
\end{equation}
where $W_{\mathrm{own}}\in\{1,\,1+C^{\mathrm{win}}_{\mathrm{lg}}\}$ is as in that
corollary. The dependence on $\check{R}_k$ sits inside the single supremum $\sigma^{\mathrm{lg}}$ of
\eqref{5.4:eq-transport-constants-d} (read at rate $7/4$ when a top-step factor $\check{R}_k$ is
present), and no upper bound on $\check{R}_k$ is required.
\item[(b)] Apart from the own-window weights of Proposition~\ref{5.4:prop-own-window-weights} and the
running-window weight of~(a), no position-dependent weight occurs in an in-$Z$ long column.
\item[(c)] In the two corollaries that bound the sums of long units at a general rate $r$ (the
rate-$r$ long-unit bound and its flat counterpart), the case distinction that fixes the factor
$\mathfrak{W}$ in Corollary~\ref{5.4:cor-long-unit-sums}(i) is void.
\end{enumerate}
\end{remark}

\begin{proof}
(a) For the running attempt the arrest data are $k_a=k$ and $X_{\mathfrak{a}}=C$, and, by the
definition of the arrest data for the running attempt, a cell $\Delta$ with
$\Delta\not\subset\mathsf{C}_C$ belongs to the running window. Consequently every cut
level-$k$ cell of $C\setminus\Lambda^{\sim}$, the collar included, is a running-window cell. In
particular, the cells adjacent to $\Delta^{\partial}$ may carry the weight
$W^{\mathrm{win}}(\check{R}_k)$.

At those cells neither a distance nor a length is available to pay this weight. This is the
difference from the own-window weights of Proposition~\ref{5.4:prop-own-window-weights}, which are
paid by the distance to the complement of the arrest's component $X_{\mathfrak{a}}$ and by the
unit's own length. The weight must therefore be carried
along. Lemma~\ref{5.4:lem-weighted-columns}, applied with $G\le W^{\mathrm{win}}(\check{R}_k)$ and
summed with the constant $C^{\mathrm{gen}}_{\mathsf{b}}$ of \eqref{5.4:eq-transport-constants-b} and
with the in-$Z$ factor $W_{\mathrm{own}}$ of Corollary~\ref{5.4:cor-long-unit-sums}(i), gives
\begin{equation}\label{5.4:eq-running-window-2}
S(\Delta^{\partial};\mathfrak{L}^{\mathrm{lg}}_{c})
\le W^{\mathrm{win}}(\check{R}_k)\,W_{\mathrm{own}}\,C^{\mathrm{gen}}_{\mathsf{b}}\,
\mathsf{b}^{\mathrm{lg}}(0)\,\mathsf{K}^{\flat}_{\mathrm{lg}}.
\end{equation}

The argument proving Corollary~\ref{5.4:cor-long-unit-sums}(i) applies word for word with
\eqref{5.4:eq-running-window-2} in place of the unweighted column bound. It yields the bound of that
corollary with $W^{\mathrm{win}}(\check{R}_k)$ displayed as an additional factor outside the named
constant $\mathsf{C}^{\mathrm{lg},c}$.

The corollary supplying the small factor $\theta'(g_k)$, through which the proof of
Corollary~\ref{5.4:cor-long-unit-sums}(i) passes, also provides a factor
$e^{-((7/2)\check{R}_k+1)}$. When $\mathsf{b}^{\mathrm{lg}}(0)$ does not depend on $\check{R}$, this
factor is not used in the proof of Corollary~\ref{5.4:cor-long-unit-sums}(i), so it remains
available here. By the definition of $\sigma^{\mathrm{lg}}$ in
\eqref{5.4:eq-transport-constants-d},
\begin{equation}\label{5.4:eq-running-window-3}
W^{\mathrm{win}}(\check{R}_k)\,e^{-((7/2)\check{R}_k+1)}\le\sigma^{\mathrm{lg}}.
\end{equation}

It may happen that both a factor $\check{R}_k$ from the top step and the weight
$W^{\mathrm{win}}(\check{R}_k)$ are present. In that case one splits $e^{-(7/2)\check{R}_k}$ into two
halves $e^{-(7/4)\check{R}_k}\cdot e^{-(7/4)\check{R}_k}$. The first half absorbs the factor
$\check{R}_k$, since $\sup_{x\ge 0}x\,e^{-(7/4)x}=4/(7e)$, so that
\begin{equation*}
\check{R}_k\,e^{-(7/4)\check{R}_k}\le\frac{4}{7e}.
\end{equation*}
The second half is used in \eqref{5.4:eq-running-window-3} with $\sigma^{\mathrm{lg}}$ taken at rate
$7/4$ in place of $7/2$. Since $e^{-(7/4)x}\ge e^{-(7/2)x}$ for $x\ge 0$, the supremum at rate $7/4$
dominates the one at rate $7/2$, so reading $\sigma^{\mathrm{lg}}$ at rate $7/4$ covers both cases.

We combine these bounds. When $\mathsf{b}^{\mathrm{lg}}(0)$ carries the top-step factor
$\check{R}_k$, that factor has been absorbed by the first half with the constant $4/(7e)<1$; in
either case $\mathsf{b}^{\mathrm{lg}}(0)$ is thereby replaced by its $\check{R}$-free counterpart
$\mathsf{b}^{\mathrm{lg},\natural}$. Here the transport is the general one of
\eqref{5.4:eq-transport-constants-b}, for which
$\mathfrak{K}=C^{\mathrm{gen}}_{\mathsf{b}}\mathsf{K}^{\flat}_{\mathrm{lg}}$. The long-unit member is
then bounded by \eqref{5.4:eq-running-window-1}. The only $\check{R}$-dependence left is inside the
single supremum $\sigma^{\mathrm{lg}}$, and no cap on $\check{R}_k$ has been used.

(b) Every position-dependent weight that a long column can carry is one of the following.
\begin{itemize}
\item The weight $\bar{\mathsf{w}}=3w$ on $W_{\mathfrak{a}}$, where $w$ obeys the window-weight
bound
\begin{equation*}
w\le\mathsf{L}_0^{w2}\,L\,C_{\uparrow}\,\check{R}_{k_a},
\end{equation*}
with $\mathsf{L}_0^{w2}$ the constant of that bound. This involves the radius of a single arrest
only, and its supremum is contained in $C^{\mathrm{win}}_{\mathrm{lg}}$ or in $\sigma^{\mathrm{lg}}$:
according as the arrest $\mathfrak{a}$ is an own arrest of $C$ or the running attempt, this is the
weight of Proposition~\ref{5.4:prop-own-window-weights} or that of~(a).
\item The weight $w=1$ off the windows.
\end{itemize}
Moreover, by the clause to this effect in the statement of the bounds for long units, those bounds
read no factor $\bar{\mathsf{w}}$. Hence no weight of a third kind occurs, which proves~(b).

(c) On the general transport the factor $\mathfrak{W}$ of
Corollary~\ref{5.4:cor-long-unit-sums}(i) is $W_AW_{\mathrm{own}}$, the values of $W_A$ and
$W_{\mathrm{own}}$ being fixed there by a case distinction; on the other transport
$\mathfrak{W}=1$. The two corollaries of part~(c), the rate-$r$ long-unit bound and its flat
counterpart, are stated under a standing condition by which this case distinction does not apply to
their bounds; it has nothing to decide there, which is part~(c).
\end{proof}

\begin{definition}[Bundles of long units off the treated class]\label{5.4:def-long-bundles}
Let $\mathfrak{L}^{\mathrm{lg}}_{d}$ be the family of long pointed units of class (d), singles and
tails, as in Definition~\ref{5.4:def-pointed-units}. For $v\in\mathfrak{L}^{\mathrm{lg}}_{d}$ let
$\mathrm{St}_v$, $T_v$ and $\partial\mathrm{St}_v$ be as there: the family of cells meeting $X_v$,
the complement of its union, and its rim.

A \emph{$\square$-bundle} over $v$ is a triple $\mathfrak{b}^{\circ}=(v;s;\mathsf{y})$ in which $s$ is an
interpolation bracket, the index of the cover $\mathfrak{B}_s$ below, and $\mathsf{y}$ is a finite
family of cells contained in $T_v$. It stands for the sum of the pieces of $v$ attached to the
blocks of the cover $\mathfrak{B}_s$,
\begin{equation*}
\mathfrak{b}^{\circ}:=\sum_{\square}v_{(\square,\mathsf{y};s)},
\end{equation*}
one piece $v_{(\square,\mathsf{y};s)}$ for each block $\square$ of $\mathfrak{B}_s$; here the indexed
letter $v_{(\square,\mathsf{y};s)}$ denotes a piece of the unit $v$, not a pointed unit. We write
$|\mathsf{y}|$ for the
number of members of $\mathsf{y}$ and $K_v$ for the set attached to $v$ in
Definition~\ref{5.4:def-pointed-units}. The norm and the label of $\mathfrak{b}^{\circ}$ occupy the
first row of \eqref{5.4:eq-long-bundles-a}, and its support set the second. In the second row
the last union runs over the hubs walled to a member of $\mathsf{y}$. The third row of
\eqref{5.4:eq-long-bundles-a} is the weighted met-cell condition for the long units of class (d)
that these bundles serve; wherever it is used it enters as a hypothesis. It is stated for every
current cell $\Delta_2\not\subset Z$ and every
$r'\le\tfrac{8}{5}\kappa$, with $S_{r'}(\Delta_2;\cdot)$ as in \eqref{5.4:eq-pointed-units-c}:
\begin{equation}\label{5.4:eq-long-bundles-a}
\begin{aligned}
\|\mathfrak{b}^{\circ}\|&:=2\,\mathfrak{N}^{\flat}(v)\,e^{-(\kappa_1-1)|\mathsf{y}|},\qquad
\mathrm{lab}(\mathfrak{b}^{\circ}):=\mathrm{lab}(K_v,\mathsf{y}),\\
Y_0(\mathfrak{b}^{\circ})&:=\bigcup\mathrm{St}_v\;\cup\;\bigcup\mathsf{y}\;\cup\;
\{\text{hubs walled to a member of }\mathsf{y}\},\\
S_{r'}(\Delta_2;\mathfrak{L}^{\mathrm{lg}}_{d})
&\le\mathfrak{B}_{\mathrm{win}}(\Delta_2)\cdot\mathfrak{K}\,
\mathsf{b}^{\mathrm{lg}}(k-\mathrm{lev}\,\Delta_2).
\end{aligned}
\end{equation}

The window factor is defined as follows. Set
$\mathfrak{B}_{\mathrm{win}}(\Delta_2):=W^{\mathrm{win}}(\check{R}_{k_{a''}})$ if two
conditions hold. The first is that $\Delta_2$ is raised, in the sense of the statement on levels,
zones, cells and the standing geometric axioms; by \eqref{5.4:eq-zones-cells-g} it is
then raised by exactly one live arrest, denoted $\mathfrak{a}''$, of level $k_{a''}$. The second is
that the filing-column bound on which Theorem~\ref{5.4:thm-finer-zone-sum} rests carries this window
weight in the sense of Lemma~\ref{5.4:lem-weighted-columns}. Set
$\mathfrak{B}_{\mathrm{win}}(\Delta_2):=1$ otherwise.

The third row of \eqref{5.4:eq-long-bundles-a} is the bound of the form
\eqref{5.4:eq-finer-zone-sum-a} at a current cell. Here it is written for class (d) and carries the
window factor of the weighted filing columns of Lemma~\ref{5.4:lem-weighted-columns} and of the
remark on the window factor.
Class (d) is treated by route (ii) of the framework statement of the label-sum theorem;
along this route that theorem is applied once, at the grain $\pi^{\mathfrak{s}}$. At this grain the
bundles above are the general
bundles of the label-sum theorem, with $\mathfrak{g}:=\pi^{\mathfrak{s}}$ and with the identity as
transfer map; this is shown below.

Besides the third row of \eqref{5.4:eq-long-bundles-a}, the bundles are used under three
hypotheses, each taken by statement.

The first is an analyticity hypothesis at every $v\in\mathfrak{L}^{\mathrm{lg}}_{d}$. Here $f_v$
denotes the function of the field carried by the unit $v$, and $\mathfrak{P}$ is a polydisc in
the variables $m=(\zeta,\varpi,\vec{\mathsf{s}})$ (in this hypothesis only, $m$ denotes a point
of $\mathfrak{P}$) on which every component $\mathsf{s}^{(t')}(\Delta)$ of
$\vec{\mathsf{s}}$, indexed by a cell $\Delta$ and a second index $t'$, satisfies
$|\mathsf{s}^{(t')}(\Delta)|\le e^{\kappa_1}$. At the first link the interpolation in
$\vec{\mathsf{s}}$ is taken over the cells of $T_v$ only, the cells of $\mathrm{St}_v$ being
exempt; this is admissible
because the subfamily clause of the lemma on $\vec{\mathsf{s}}$ holds for any subfamily. The
hypothesis asks that, for each
interpolation bracket $s$ and each $t\in[0,1]$, three things hold.
\begin{itemize}
\item The first line of \eqref{5.4:eq-long-bundles-b} holds.
\item The map $\chi\mapsto f_v\circ\Phi^{\chi}_s$ is holomorphic in the finitely many values of
$\chi$, in a neighbourhood of the constant $\chi\equiv t$ (as a composite of holomorphic maps).
\item $f_v\circ\Phi$ depends on $\Phi$ only through its restriction to $X_v$.
\end{itemize}

The second requires that the partition of unity be exact on the set where the data $A'_s$, $X_s$
and the field $\mathbb{H}_s$ of the bracket $s$ are defined; this is the second line of
\eqref{5.4:eq-long-bundles-b}.

The third requires that clauses (ii) and (iii$'$) and the subfamily clause of the lemma on the
interpolation parameters $\vec{\mathsf{s}}$ hold; this is the third line of
\eqref{5.4:eq-long-bundles-b}.
\begin{equation}\label{5.4:eq-long-bundles-b}
\begin{gathered}
m\mapsto f_v\circ\Phi^{t}_s\ \text{is holomorphic on }\mathfrak{P},\qquad
\bigl|f_v\circ\Phi^{t}_s-f_v(1)\bigr|\le\mathfrak{N}^{\flat}(v),\\
\sum_{\square}\hat{\chi}_{\square}=1\quad\text{where }A'_s,\ X_s,\ \mathbb{H}_s\text{ are defined},\\
\text{the lemma on }\vec{\mathsf{s}}\text{: clauses (ii), (iii$'$) and the subfamily clause hold.}
\end{gathered}
\end{equation}

Together with these hypotheses the bundles are used with the following geometry: the standing
geometric axioms on levels, zones
and cells, and clause (f) of the geometric lemma accompanying them. For the window factor they are
used in addition with the axioms \eqref{5.4:eq-zones-cells-g}. From the framework of the label-sum
theorem the following enter:
\begin{itemize}
\item clauses (b), (c) and (e) of the framework statement of the label-sum theorem;
\item the label-sum theorem itself, at grain $\pi^{\mathfrak{s}}$ and with the parameter $\beta$
of the sums $S^{(\beta)}_{r'}$ of \eqref{5.4:eq-pointed-units-c} (the theorem is stated for the
grain $\pi^{\mathfrak{s}}_j$ or $\pi^{\mathfrak{s}}$);
\item two corollaries of the label-sum theorem, one of which is the corollary on the small factor
$\theta'(g_k)$.
\end{itemize}
\end{definition}

\begin{proof}
Two assertions in Definition~\ref{5.4:def-long-bundles} need justification.

The first is that the window factor is well defined. Let $\Delta_2\not\subset Z$ be a raised
current cell. By \eqref{5.4:eq-zones-cells-g} it is raised by exactly one live arrest
$\mathfrak{a}''$. Hence the level $k_{a''}$ of $\mathfrak{a}''$, and with it
$W^{\mathrm{win}}(\check{R}_{k_{a''}})$, is determined by $\Delta_2$. As $\Delta_2\not\subset Z$,
the cell $\Delta_2$ is foreign or sibling to $\mathfrak{a}''$, in the sense of the statement on
levels, zones, cells and the standing geometric axioms; the first alternative in the
definition of the window factor is thus applied only to cells of these two kinds.

The second is the identification with the general bundles of the label-sum theorem. Take the grain
$\mathfrak{g}:=\pi^{\mathfrak{s}}$ and the identity as transfer map. By clause (h) of
Lemma~\ref{5.4:lem-cells-carrier}, three facts hold at this grain:
\begin{itemize}
\item the star of $v$ at grain $\mathfrak{g}$ is $\mathrm{St}^{\mathfrak{g}}_v=\mathrm{St}_v$;
\item the set of hubs walled to the star is empty, $\mathfrak{H}_v=\emptyset$;
\item the set of attachment cubes is empty.
\end{itemize}
Substitute these three facts and the identity transfer map into the general bundle of the label-sum
theorem at grain
$\mathfrak{g}$. The star is replaced by $\mathrm{St}_v$. The hub contribution attached to the star
disappears, and so does the attachment part of the label. What remains is the sum over the blocks of
$\mathfrak{B}_s$ displayed before \eqref{5.4:eq-long-bundles-a}, together with the three definitions
in the first two rows of \eqref{5.4:eq-long-bundles-a}.
\end{proof}

\begin{lemma}[The bundle frame for long units off the treated class]\label{5.4:lem-bundle-frame}
Let $v$ be a long unit of the class $\mathfrak{L}^{\mathrm{lg}}_{d}$, read at the grain
$\mathfrak{g}=\pi^{\mathfrak{s}}$, and let $\mathfrak{b}^{\circ}(v;s;\mathsf{y})$ denote its
$\square$-bundles, with norm $\|\mathfrak{b}^{\circ}\|$, label $\mathrm{lab}(\mathfrak{b}^{\circ})$
and set $Y_0(\mathfrak{b}^{\circ})$ as in Definition~\ref{5.4:def-long-bundles} and
\eqref{5.4:eq-long-bundles-a}. Assume the following three hypotheses, in the form displayed in
\eqref{5.4:eq-long-bundles-b}:
\begin{itemize}
\item the undilated containment;
\item the linearity of the partition of unity $\{\hat{\chi}_{\square}\}$;
\item the support lemma.
\end{itemize}
Then the following hold. The clauses carry the letters of the corresponding clauses of the frame
lemma for the summation of boundary terms at finer zones; clauses (a) and (d) of that lemma have
no counterpart here.
\begin{enumerate}
\item[(b)] For every $s$ and every family $\mathsf{y}$ of derivative cells,
\begin{equation}\label{5.4:eq-bundle-frame-1}
\mathfrak{b}^{\circ}(v;s;\mathsf{y})
=\prod_{\Delta\in\mathsf{y}}\int_0^1 d\mathsf{s}(\Delta)\,\partial_{\mathsf{s}(\Delta)}\,
F_s(\mathsf{s})\Big|_{\mathsf{s}=0\text{ off }\mathsf{y}},
\end{equation}
where
\begin{equation}\label{5.4:eq-bundle-frame-2}
F_s:=\bigl[f_v\circ\Phi_s^{1}-f_v\circ\Phi_s^{0}\bigr](\mathsf{s}),
\qquad
\Phi_s^{0}:=\Phi_s^{\chi\equiv 0},
\end{equation}
and $\Phi_s^{0}$ does not depend on the dilation parameter $\mathsf{t}$, although it does depend
on $\mathsf{s}$. Moreover $|\mathfrak{b}^{\circ}|\le\|\mathfrak{b}^{\circ}\|$; the bundle
$\mathfrak{b}^{\circ}$ is linear in $f_v$, and it is holomorphic in every parameter in which
$f_v\circ\Phi$ is holomorphic.
\item[(c)] The bundle $\mathfrak{b}^{\circ}$ depends on the data and the background fields only
through their restrictions to $Y_0(\mathfrak{b}^{\circ})$; $\mathfrak{b}^{\circ}=0$ unless
$\mathsf{y}\in\mathfrak{Y}(\mathrm{St}_v)$, with $\mathfrak{Y}(\mathrm{St}_v)$ as in
\eqref{5.4:eq-pointed-units-d}; the bundle with $\mathsf{y}=\emptyset$ vanishes.
\item[(e)] Every piece lies in exactly one $\square$-bundle. All pieces of one bundle carry one and
the same pair $(K_p,Y_0)$ under the labelling rule of Definition~\ref{5.4:def-long-bundles},
taken at the identity value of its parameter and followed by the relabelling map $\mathcal{R}_1$.
For
$\mathsf{y}\in\mathfrak{Y}(\mathrm{St}_v)$ one has $\mathrm{lab}(\mathfrak{b}^{\circ})\in\mathbf{D}_k$
and
\begin{equation}\label{5.4:eq-bundle-frame-3}
\begin{gathered}
d_k(\mathrm{lab})\le d_k(K_v)+c_l^{\circ}|\mathsf{y}|,\qquad
d_{k,Z}\bigl(Y(\mathrm{lab})\bigr)\le d_k(\mathrm{lab}),\\
d_k(\mathrm{lab})\ge 9\check{R}_k-2 .
\end{gathered}
\end{equation}
\end{enumerate}
\end{lemma}

\begin{proof}
\emph{Clause} (b). Fix $s$, a point $\mathsf{s}$ of the sealed polydisc $\mathfrak{P}$, and
$t\in[0,1]$; write $\Phi_s^{t}:=\Phi_s^{\chi\equiv t}$. By the clause on $\chi$ of the hypotheses
displayed in \eqref{5.4:eq-long-bundles-b},
the differential of the map $\chi\mapsto f_v\circ\Phi_s^{\chi}(\mathsf{s})$ at the constant
function $\chi\equiv t$ is linear in the direction of differentiation. The functions
$\hat{\chi}_{\square}$ form a partition of unity, so by the linearity hypothesis the
derivatives in the directions $\hat{\chi}_{\square}$ add up to the derivative in the direction of
the constant function one, which is the derivative in $t$:
\begin{equation}\label{5.4:eq-bundle-frame-4}
\sum_{\square}\partial_{\mathsf{t}}\big|_{\mathsf{t}=0}\,
f_v\bigl(\Phi_s^{\,t+\mathsf{t}\hat{\chi}_{\square}}\bigr)(\mathsf{s})
=\frac{d}{dt}\,f_v\bigl(\Phi_s^{t}\bigr)(\mathsf{s}).
\end{equation}
This is the identity \cite[(3.4)]{B-CMP109}, read from its last member to its first; in other
words, the pieces localised at the blocks $\square$ sum to the bracket that they localise.
Integrating \eqref{5.4:eq-bundle-frame-4} over $t\in[0,1]$ and using the fundamental theorem of
calculus in $t$ gives
\begin{equation}\label{5.4:eq-bundle-frame-5}
\sum_{\square}\int_0^1 dt\,\partial_{\mathsf{t}}\big|_{\mathsf{t}=0}\,
f_v\bigl(\Phi_s^{\,t+\mathsf{t}\hat{\chi}_{\square}}\bigr)(\mathsf{s})
=f_v\circ\Phi_s^{1}(\mathsf{s})-f_v\circ\Phi_s^{0}(\mathsf{s})=F_s(\mathsf{s}).
\end{equation}
The bundle $\mathfrak{b}^{\circ}(v;s;\mathsf{y})$ is the sum over the blocks $\square$ of the
pieces obtained from the summands on the left of \eqref{5.4:eq-bundle-frame-5} by applying the
operators $\prod_{\Delta\in\mathsf{y}}\int_0^1 d\mathsf{s}(\Delta)\,\partial_{\mathsf{s}(\Delta)}$
and setting $\mathsf{s}=0$ off $\mathsf{y}$. The sum over $\square$ is finite, so it commutes with
these operators, and \eqref{5.4:eq-bundle-frame-5} gives \eqref{5.4:eq-bundle-frame-1}. Since
$\Phi_s^{0}$ is obtained at $\chi\equiv 0$, it carries no dilation parameter $\mathsf{t}$; it still
depends on $\mathsf{s}$.

For the bound: on the polydisc both members of $F_s$, namely $f_v\circ\Phi_s^{1}$ and
$f_v\circ\Phi_s^{0}$, lie within $\mathfrak{N}^{\flat}(v)$ of $f_v(1)$, hence
\begin{equation*}
|F_s(\mathsf{s})|\le 2\,\mathfrak{N}^{\flat}(v)\qquad(\mathsf{s}\in\mathfrak{P}).
\end{equation*}
In \eqref{5.4:eq-bundle-frame-1} each integration in $\mathsf{s}(\Delta)$, $\Delta\in\mathsf{y}$,
runs over an interval of length one, so the bundle is bounded by the supremum of the integrand
$\prod_{\Delta\in\mathsf{y}}\partial_{\mathsf{s}(\Delta)}F_s$ over $\mathsf{s}(\Delta)\in[0,1]$
for $\Delta\in\mathsf{y}$ and $\mathsf{s}=0$ off $\mathsf{y}$. Each derivative
$\partial_{\mathsf{s}(\Delta)}$ at such a point is estimated by Cauchy's formula in the single
variable $\mathsf{s}(\Delta)$, the others frozen, on the circle of radius $e^{\kappa_1}-1$ about
$\mathsf{s}(\Delta)$, which lies inside the polydisc; this is the Cauchy
step of the proof of the two Cauchy estimates (Lemma~\ref{5.4:lem-two-cauchy}), with $g$ replaced
by $F_s$ and $\mathfrak{N}$ by $2\,\mathfrak{N}^{\flat}(v)$. Each derivative
costs a factor $(e^{\kappa_1}-1)^{-1}$. Since $\kappa_1\ge 1$ (as in
Lemma~\ref{5.4:lem-two-cauchy}) we have $e^{\kappa_1}\ge e$ and
therefore
\begin{equation*}
e^{\kappa_1}-1-e^{\kappa_1-1}=e^{\kappa_1}\bigl(1-e^{-1}\bigr)-1\ge e-2>0,
\end{equation*}
that is, $e^{\kappa_1}-1\ge e^{\kappa_1-1}$. Altogether
\begin{equation}\label{5.4:eq-bundle-frame-6}
|\mathfrak{b}^{\circ}(v;s;\mathsf{y})|
\le 2\,\mathfrak{N}^{\flat}(v)\,(e^{\kappa_1}-1)^{-|\mathsf{y}|}
\le 2\,\mathfrak{N}^{\flat}(v)\,e^{-(\kappa_1-1)|\mathsf{y}|},
\end{equation}
which is the bound $|\mathfrak{b}^{\circ}|\le\|\mathfrak{b}^{\circ}\|$, with the norm of
\eqref{5.4:eq-long-bundles-a}. By \eqref{5.4:eq-bundle-frame-1}--\eqref{5.4:eq-bundle-frame-2},
$F_s$ is linear in $f_v$ and the operators applied to it are linear, so $\mathfrak{b}^{\circ}$ is
linear in $f_v$. The function $f_v\circ\Phi$ enters \eqref{5.4:eq-bundle-frame-1} only through
derivatives in the $\mathsf{s}(\Delta)$ and integrals over $[0,1]$, so $\mathfrak{b}^{\circ}$ is
holomorphic in every parameter in which $f_v\circ\Phi$ is.

\emph{Clause} (c). We run the argument of clause (c) of the frame lemma for the summation of
boundary terms at finer zones, with the grain star $\mathrm{St}^j_v$ replaced by $\mathrm{St}_v$
and with $\mathfrak{H}_v=\emptyset$; the latter is the case here because, by
Lemma~\ref{5.4:lem-cells-carrier}(f), no member of $\mathrm{St}_v$ is walled to a hub (a cell of
$\Lambda^{\sim}$). We write the argument out.

The function $F_s$ depends on the parameters $\mathsf{s}$ of the first link of the chain only
through the restriction to $X_v$ of the amplitude $\mathcal{A}(\mathsf{s})$, taken at amplitude
one, by the statement fixing the amplitude $\mathcal{A}(\mathsf{s})$ of the first link.
At $\mathsf{s}=0$ off $\mathsf{y}$ the support set $Y(\sigma)$ of the support lemma is
\begin{equation*}
Y(\sigma)=\bigcup\mathrm{St}_v\cup\Lambda^{\sim}\cup\bigcup\mathsf{y}.
\end{equation*}
The carrier $X_v$ is wall-connected, so it lies in one connected component of $Y(\sigma)$; call it
$Y^{v}$. By the support lemma, $F_s$ depends on $\mathsf{s}$, on the data and on the background
fields only through their restrictions to $Y^{v}$, and by \eqref{5.4:eq-bundle-frame-1} so does
$\mathfrak{b}^{\circ}$.

Now let $\mathsf{y}_a$ be a connected component of $\mathsf{y}$. If
$\bigcup\mathsf{y}_a\not\subset Y^{v}$, then $F_s$ does not depend on the parameters
$\mathsf{s}(\Delta)$, $\Delta\in\mathsf{y}_a$, and the derivatives $\partial_{\mathsf{s}(\Delta)}$
in \eqref{5.4:eq-bundle-frame-1} annihilate it, as in \cite{B-CMP116} and by the clause of the
statement on the first link of the chain asserting that $F_s$ does not depend on
$\mathsf{s}(\Delta)$ for cells $\Delta$ outside the component that is read. If
$\bigcup\mathsf{y}_a\subset Y^{v}$, then $\mathsf{y}_a$ is joined to $\mathrm{St}_v$ along the
wall-contact graph of $Y^{v}$. No member of $\mathrm{St}_v$ is walled to a hub
(Lemma~\ref{5.4:lem-cells-carrier}(f)) and $X_v$ meets no hub; hence a shortest such path leaves
$\mathsf{y}_a$ into $\mathrm{St}_v$ across a wall. Thus $\mathsf{y}_a$ has a member in
$\partial\mathrm{St}_v$, and $\mathrm{St}_v\cup\mathsf{y}_a$ is connected.

Suppose $Y^{v}$ contains no hub. The data are supported in the set $\Lambda$: for $s=4$ by the
support statement for the sourced data, for $s=3$ because $\mathsf{J}_3\subset\Lambda$, by the
clause of the statement on the first link of the chain placing $\mathsf{J}_3$ inside $\Lambda$,
and for $s=2$ and
$s=1$ by the defining clauses of those data. Hence the datum vanishes on $Y^{v}$, and there the
amplitude is $\mathcal{A}\equiv(0,0,1)$, by the statement fixing the amplitude of the first link.
Consequently $\Phi_s^{1}=\Phi_s^{0}$ on $X_v$, and
$F_s\equiv 0$. Suppose instead that $Y^{v}$ contains a hub. That hub is walled to a member of
$\mathsf{y}$, and that member is then a source cell.

Hence $\mathfrak{b}^{\circ}\neq 0$ forces $\mathsf{y}\in\mathfrak{Y}(\mathrm{St}_v)$. Components of
$\mathsf{y}$ joined to $\mathrm{St}_v$ but themselves walled to no hub of $Y^{v}$ are counted
here; they are not claimed to be annihilated. Because $\mathfrak{H}_v=\emptyset$,
however, in every configuration that survives $\mathrm{St}_v\cup\mathsf{y}$ is connected and
$\mathsf{y}$ contains a source cell. Finally, if $\mathsf{y}=\emptyset$ then
$Y^{v}\subset\bigcup\mathrm{St}_v$ contains
no hub, so $F_s\equiv 0$ by the case just treated, and $\mathfrak{b}^{\circ}=0$.

\emph{Clause} (e). A piece is indexed by $(v,s,\mathsf{y})$ together with a block $\square$, and the
bundle $\mathfrak{b}^{\circ}(v;s;\mathsf{y})$ is the sum over $\square$ of the pieces with these
indices. So every piece lies in exactly one $\square$-bundle. The block $\square$ does not enter
the pair $(\mathsf{y},\mathrm{St}_v)$ from which the labelling rule, at the identity value of its
parameter and followed by $\mathcal{R}_1$, computes $(K_p,Y_0)$. Hence all pieces of one bundle
carry the same pair.

Let $\mathsf{y}\in\mathfrak{Y}(\mathrm{St}_v)$. The wall-connectedness of the label,
$\mathrm{lab}\in\mathbf{D}_k$, and the bound $d_k(\mathrm{lab})\le d_k(K_v)+c_l^{\circ}|\mathsf{y}|$
follow by the argument of clause (e) of the frame lemma for the summation of boundary terms at
finer zones. That argument is applied here with $[\mathrm{St}_v]_k\subset K_v$, which holds by
Lemma~\ref{5.4:lem-cells-carrier}(f).

For the lower bound, the label contains two things. It contains a cube of $K_v$, and
$K_v\subset\operatorname{cl}Z^{c}$. It also contains the level-$k$ parent of a source cell, and
that parent lies at distance at least $9\check{R}_k-1$ from $Z^{c}$, by the geometric condition on
source cells. Since $\mathrm{lab}$ is wall-connected and contains both,
$d_k(\mathrm{lab})\ge 9\check{R}_k-2$.
\end{proof}

\begin{lemma}[Rim mass of long units, all levels at once]\label{5.4:lem-rim-mass}
Assume the bundle condition \eqref{5.4:eq-long-bundles-a}. Then for every
$r'\le \tfrac{8}{5}\kappa$ and every cell $\Delta_1\in\pi^{\mathfrak{s}}$,
\begin{equation}\label{5.4:eq-rim-mass-1}
\begin{split}
&e^{-3\rho(\Delta_1)}
\sum_{\substack{v\in\mathfrak{L}^{\mathrm{lg}}_{d}:\\ \Delta_1\in\partial\mathrm{St}_v}}
2\,\mathfrak{N}^{\flat}(v)\,e^{r' d_k(K_v)}\,e^{\epsilon_{\partial}\,\#\partial\mathrm{St}_v}\\
&\qquad\le e^{-(9\check{R}_k-2)}\cdot 4e\cdot dL^{d-1}\,
\bigl(1+C^{\mathrm{win}}_{\mathrm{lg}}\bigr)\,C_{\mathsf{b}}\,\mathfrak{s}_{\mathsf{b}}\,
\mathfrak{K}\,\mathsf{b}^{\mathrm{lg}}(0),
\end{split}
\end{equation}
where $\mathfrak{L}^{\mathrm{lg}}_{d}$ is the set of long units off the treated class
(Definition~\ref{5.4:def-long-bundles}), $\mathfrak{N}^{\flat}(v)$ is the size of the unit $v$,
$K_v$ is attached to the unit $v$ as in Definition~\ref{5.4:def-long-bundles}, and $\rho$,
$\mathfrak{K}$, $C^{\mathrm{win}}_{\mathrm{lg}}$ are as in \eqref{5.4:eq-zones-cells-d},
\eqref{5.4:eq-transport-constants-c} and \eqref{5.4:eq-transport-constants-d}.
\end{lemma}

\begin{proof}
If the sum on the left of \eqref{5.4:eq-rim-mass-1} is empty, there is nothing to prove,
the right-hand side being non-negative. Otherwise there is a unit
$v\in\mathfrak{L}^{\mathrm{lg}}_{d}$ with $\Delta_1\in\partial\mathrm{St}_v$, and part~(f) of the
lemma on cells against a carrier (Lemma~\ref{5.4:lem-cells-carrier}) gives
$\rho(\Delta_1)\ge 9\check{R}_k-2$. We write $e^{-3\rho(\Delta_1)}$ as the product of three
factors $e^{-\rho(\Delta_1)}$; the first of them is at most $e^{-(9\check{R}_k-2)}$.

Next we organise the sum over $v$. The relation $\Delta_1\in\partial\mathrm{St}_v$ means that
$\Delta_1$ shares a wall with some cell $\Delta_2\in\mathrm{St}_v$; such a cell is not contained
in $Z$, it is one of the at most $2dL^{d-1}$ cells sharing a wall with $\Delta_1$, and $X_v$
meets $\Delta_2$. Moreover
\begin{equation*}
e^{\epsilon_{\partial}\,\#\partial\mathrm{St}_v}\le e^{\frac12(d_v+1)}\le e\cdot e^{2d_v},
\end{equation*}
the second inequality because $\tfrac12(d_v+1)\le 1+2d_v$. Bounding the sum over $v$ by the
sum over the cells $\Delta_2$ sharing a wall with $\Delta_1$ of the sums over the units $v$
whose carrier $X_v$ meets $\Delta_2$, and applying the bundle condition
\eqref{5.4:eq-long-bundles-a} with $r'\le\tfrac85\kappa$, we obtain
\begin{equation*}
\sum_{\substack{v\in\mathfrak{L}^{\mathrm{lg}}_{d}:\\ \Delta_1\in\partial\mathrm{St}_v}}
2\,\mathfrak{N}^{\flat}(v)\,e^{r' d_k(K_v)}\,e^{\epsilon_{\partial}\,\#\partial\mathrm{St}_v}
\le 2e\sum_{\Delta_2}\mathfrak{B}_{\mathrm{win}}(\Delta_2)\,\mathfrak{K}\,
\mathsf{b}^{\mathrm{lg}}(k-\operatorname{lev}\Delta_2).
\end{equation*}

The two remaining factors $e^{-\rho(\Delta_1)}$ are spent on each term of this sum.

For the age factor, we use part~(iii) of the lemma on the rim and trunk of a long unit off the
treated class (Lemma~\ref{5.4:lem-rim-trunk}): it gives $\rho(\Delta_1)\ge k-\operatorname{lev}\Delta_2$,
since $\Delta_1\in\partial\mathrm{St}_v$ shares a wall with $\Delta_2\in\mathrm{St}_v$. Hence
$e^{-\rho(\Delta_1)}\le e^{-(k-\operatorname{lev}\Delta_2)}$, and the lemma on polynomial
against exponential (Lemma~\ref{5.4:lem-poly-exp-sup}), applied at
$x=k-\operatorname{lev}\Delta_2$, gives
\begin{equation*}
e^{-\rho(\Delta_1)}\,\mathsf{b}^{\mathrm{lg}}(k-\operatorname{lev}\Delta_2)
\le C_{\mathsf{b}}\,\mathfrak{s}_{\mathsf{b}}\,\mathsf{b}^{\mathrm{lg}}(0).
\end{equation*}

For the window factor, suppose first that $\Delta_2$ is raised, that is,
$\Delta_2\subset W_{\mathfrak{a}''}$
for a live arrest $\mathfrak{a}''$. Since $\Delta_2\not\subset Z$, the component
$X_{\mathfrak{a}''}$ is foreign or a sibling, so the collar lemma for foreign and sibling
windows (Lemma~\ref{5.4:lem-collar}) applies to $\Delta_2$ and to $\Delta_1$, which shares a
wall with it, and gives $\rho(\Delta_1)\ge 2\check{R}_{k_{a''}}-1$. Therefore
\begin{equation*}
\mathfrak{B}_{\mathrm{win}}(\Delta_2)\,e^{-\rho(\Delta_1)}
\le W^{\mathrm{win}}(\check{R}_{k_{a''}})\,e^{-(2\check{R}_{k_{a''}}-1)}
\le C^{\mathrm{win}}_{\mathrm{lg}},
\end{equation*}
with $W^{\mathrm{win}}$ and $C^{\mathrm{win}}_{\mathrm{lg}}$ as in
\eqref{5.4:eq-transport-constants-d}. If $\Delta_2$ is not raised, then
$\mathfrak{B}_{\mathrm{win}}(\Delta_2)\,e^{-\rho(\Delta_1)}\le 1$. In either case the
product is at most $1+C^{\mathrm{win}}_{\mathrm{lg}}$.

Combining the three factors, each term of the sum over $\Delta_2$, multiplied by
$e^{-3\rho(\Delta_1)}$, is at most
\begin{equation*}
e^{-(9\check{R}_k-2)}\,\bigl(1+C^{\mathrm{win}}_{\mathrm{lg}}\bigr)\,
C_{\mathsf{b}}\,\mathfrak{s}_{\mathsf{b}}\,\mathfrak{K}\,\mathsf{b}^{\mathrm{lg}}(0).
\end{equation*}
Summing over the at most $2dL^{d-1}$ cells $\Delta_2$ and restoring the prefactor $2e$ gives
the factor $4e\cdot dL^{d-1}$ and hence \eqref{5.4:eq-rim-mass-1}.
\end{proof}

\begin{lemma}[Placing the window weight]\label{5.4:lem-window-placement}
Let $v\in\mathfrak{L}^{\mathrm{lg}}_{d}$, let $\Delta_2\in\mathrm{St}_v$, and let
$\Delta_1\in\partial\mathrm{St}_v$ share a wall with $\Delta_2$. Let $\Delta''$ be a raised cell with
$\Delta''\cap\bigcup X_v\neq\emptyset$ and $\Delta''\subset W_{\mathfrak{a}''}$ for an arrest
$\mathfrak{a}''$. Put $X'':=X_{\mathfrak{a}''}$, the component of $\mathfrak{a}''$, and
$\varrho:=\check{R}_{k_{a''}}$. Then:
\begin{enumerate}
\item[(0)] $\mathfrak{a}''$ is a live arrest on a component not in $Z$;
\item[(i)] $j_v\le k_{a''}$;
\item[(ii)] $\rho(\Delta_1)+d_v\ge 2\varrho-3$;
\item[(iii)] for $c\ge 3$,
\begin{equation*}
W^{\mathrm{win}}(\varrho)\,e^{-\rho(\Delta_1)}\,e^{-(c-2)d_v}\le C^{\mathrm{win}}_{\mathrm{lg}}.
\end{equation*}
\end{enumerate}
\end{lemma}

\begin{proof}
\emph{Clause (0).} The cell $\Delta''$ has a common point with $\bigcup X_v$, and
$\bigcup X_v$ is contained in the closure of $Z^c$. Hence $\Delta''$ lies within distance $1$ of
the closure of $Z^c$. Suppose $\Delta''\subset Z$. Then $\Delta''$ would be a cell of $Z$ lying
within distance $2\check{R}_k$ of $\partial Z$, and by (\ref{5.4:eq-zones-cells-f}) such a cell is
not raised. This contradicts the hypothesis that $\Delta''$ is raised, so $\Delta''\not\subset Z$.
The running preparation attempt and the arrests belonging to $Z$ itself raise cells only inside
$Z$. Since $\Delta''$ is raised, lies in the window $W_{\mathfrak{a}''}$, and is not contained in
$Z$, the arrest $\mathfrak{a}''$ is neither of these. It is therefore foreign or a sibling, and
its component satisfies $X''\cap Z=\emptyset$.

\emph{Clause (i).} Fix a point $y\in\Delta''\cap\bigcup X_v$, and let $Q$ be a cube of $X_v$ with
$y\in Q$. By clause (c) of the lemma on cells against a carrier
(Lemma~\ref{5.4:lem-cells-carrier}), $Q$ lies in a cell $\Delta(Q)$, and
$\operatorname{lev}\Delta(Q)\ge j_v$. Since $y\in\Delta''\subset W_{\mathfrak{a}''}$, the
window-inclusion axiom in \eqref{5.4:eq-zones-cells-g} gives
\begin{equation*}
\{x:\ \mathrm{dist}_{k_{a''}}(x,y)<2\varrho\}\subset X''.
\end{equation*}
This set contains a neighbourhood of $y$, and $y\in Q\subset\Delta(Q)$, so it meets the interior
of $\Delta(Q)$. Thus $X''$ meets the interior of $\Delta(Q)$, and the zone--level axiom in
\eqref{5.4:eq-zones-cells-g} (a cell whose interior meets the component $X_{\mathfrak{a}}$ of a live
arrest $\mathfrak{a}$ has level at most $k_a$) gives $\operatorname{lev}\Delta(Q)\le k_{a''}$.
Together, $j_v\le\operatorname{lev}\Delta(Q)\le k_{a''}$.

\emph{Clause (ii).} Put
\begin{equation*}
g:=\mathrm{dist}_{k_{a''}}\bigl(\cdot\,,(X'')^{c}\bigr).
\end{equation*}
The inclusion displayed in the proof of clause (i) gives $g(y)\ge 2\varrho$. Define $\delta$ as
follows. If $\Delta_2$ does not meet $X''$, set $\delta:=0$. Then $\Delta_2\subset(X'')^c$, so
$g=0\le\delta+1$ on $\Delta_2$. If $\Delta_2$ meets $X''$, set $\delta:=\min_{\Delta_2}g$; in this
case $\max_{\Delta_2}g\le\delta+1$. In either case $g\le\delta+1$ on $\Delta_2$.

The carrier $X_v$ contains $y$, where $g\ge 2\varrho$. Since $\Delta_2\in\mathrm{St}_v$, it also
meets $\Delta_2$, where $g\le\delta+1$. Applied to $X_v$, the distance axiom in
\eqref{5.4:eq-zones-cells-c} therefore gives
\begin{equation*}
d_v\ge(2\varrho-\delta-2)_+ .
\end{equation*}

Next we bound $\rho(\Delta_1)$ from below. In all cases $\rho(\Delta_1)\ge 1$. Now let $\delta>1$.
The cells $\Delta_1$ and $\Delta_2$ share a wall. At a common point $g\ge\delta>0$, so that point
lies in $X''$; hence $\Delta_1$ meets $X''$. Moreover $\min_{\Delta_1}g\ge\delta-1$. The chain
count in the proof of the collar lemma (Lemma~\ref{5.4:lem-collar}), which bounds $\rho(\Delta_1)$
from below by the level-$k_{a''}$ distance from $\Delta_1$ to $(X'')^c$, that is by
$\min_{\Delta_1}g$, then gives $\rho(\Delta_1)\ge\delta-1$. In every case, therefore,
$\rho(\Delta_1)\ge\max\{1,\delta-1\}$.
Combining this with the bound on $d_v$,
\begin{equation}\label{5.4:eq-window-placement-1}
\rho(\Delta_1)+d_v\ge h(\delta)
:=(2\varrho-\delta-2)_+ +\max\{1,\delta-1\}.
\end{equation}

It remains to show that $h(\delta)\ge 2\varrho-3$ for all $\delta\ge0$. If $0\le\delta\le2$, then
\begin{equation*}
h(\delta)\ge(2\varrho-\delta-2)+1=2\varrho-\delta-1\ge 2\varrho-3 .
\end{equation*}
If $\delta>2$, then
\begin{equation*}
h(\delta)\ge(2\varrho-\delta-2)+(\delta-1)=2\varrho-3 .
\end{equation*}
With \eqref{5.4:eq-window-placement-1} this proves clause (ii).

\emph{Clause (iii).} Since $d_v\ge0$ and $c-2\ge1$, we have
$e^{-(c-2)d_v}\le e^{-d_v}$. Hence, by clause (ii),
\begin{equation*}
e^{-\rho(\Delta_1)}\,e^{-(c-2)d_v}\le e^{-(\rho(\Delta_1)+d_v)}\le e^{-(2\varrho-3)} .
\end{equation*}
The exponents satisfy
\begin{equation*}
-(2\varrho-3)-\Bigl(1-\frac{2\varrho-2}{\mathsf{c}_d}\Bigr)
=-(2\varrho-2)\bigl(1-\mathsf{c}_d^{-1}\bigr),
\end{equation*}
which is non-positive because $\check{R}_{k_{a''}}\ge1$ and $\mathsf{c}_d\ge1$, by the standing
bounds on the radii in \eqref{5.4:eq-zones-cells-g} and on $\mathsf{c}_d$ in
\eqref{5.4:eq-zones-cells-c}; and so
\begin{equation*}
e^{-(2\varrho-3)}\le e^{\,1-(2\varrho-2)/\mathsf{c}_d}.
\end{equation*}
By the definition of $C^{\mathrm{win}}_{\mathrm{lg}}$ among the constants of the two transports,
\eqref{5.4:eq-transport-constants-d}, we have
$W^{\mathrm{win}}(\varrho)\,e^{1-(2\varrho-2)/\mathsf{c}_d}\le C^{\mathrm{win}}_{\mathrm{lg}}$ for
every level. Multiplying the last two displays by $W^{\mathrm{win}}(\varrho)$ proves clause (iii).
\end{proof}

\begin{remark}
On route G of the two transports (Notation~\ref{5.4:not-transport-constants}) the lemma allows the
window weight to be placed at the cell $\Delta_2$ met by $X_v$, at the cell containing $y_v$, or at
any cell in closed contact with $X_v$; the available factors are then divided as follows. The factor
$e^{-3\rho(\Delta_1)}$ is spent once on each of the smallness, the age and the window. The factor
$e^{2d_v}$ splits as
\begin{equation*}
e^{2d_v}=e^{d_v/2}\cdot e^{d_v}\cdot e^{d_v/2}.
\end{equation*}
The first factor goes to the shape of the rim. The second is a reserve, spent only when the
weight is placed at a cell other than $\Delta_2$. The third is not used.
\end{remark}

\begin{theorem}[The bundle sum for long units off the treated class]\label{5.4:thm-bundle-sum}
Assume the following.
\begin{enumerate}
\item The three conditions displayed in \eqref{5.4:eq-long-bundles-b} hold for the long units
$\mathfrak{L}^{\mathrm{lg}}_d$.
\item The base bound for long units displayed in \eqref{5.4:eq-long-bundles-a} holds.
\item The label-sum theorem holds, for an arbitrary grain and every admissible value of its
parameter $\beta$.
\item The floor conditions on $\kappa_1$ hold (among them $\kappa_1\ge 1$), and $\kappa\ge 180$,
$L\ge 3$, together with the floor condition on $\check{R}_k$.
\end{enumerate}
Then for every $K$, every $k\le K$, every label, every $r\le a+8$ and every $Q_*\in\pi_k$,
\begin{equation}\label{5.4:eq-bundle-sum-1}
\sum_{\substack{\mathfrak{b}^{\circ}=(v;s;\mathsf{y}):\ v\in\mathfrak{L}^{\mathrm{lg}}_d,\ s\le 4,\\
\mathsf{y}\in\mathfrak{Y}(\mathrm{St}_v),\ \mathrm{lab}(\mathfrak{b}^{\circ})\ni Q_*}}
\|\mathfrak{b}^{\circ}\|\,e^{r d_k(\mathrm{lab}\,\mathfrak{b}^{\circ})}
\le e^{-(9\check{R}_k-2)}\,\mathsf{C}^{\mathrm{lg},d}_{\circ}\,\mathsf{b}^{\mathrm{lg}}(0),
\end{equation}
where $\mathfrak{Y}(\mathrm{St}_v)$ is as in \eqref{5.4:eq-pointed-units-d} and
$\mathsf{C}^{\mathrm{lg},d}_{\circ}$ as in \eqref{5.4:eq-long-unit-sums-a}. The sum runs over every
level $j_v\le k$, every zone, $\Omega_k$, the untreated shells, the foreign frozen windows, the four
brackets $s\le 4$ and all families $\mathsf{y}$. By clauses (b), (c) and (e) of
Lemma~\ref{5.4:lem-bundle-frame}, which is proved as an implication from the conditions of
\eqref{5.4:eq-long-bundles-b}, the left-hand side of \eqref{5.4:eq-bundle-sum-1} majorises, for
$Q_*\not\subset Z$, the quantity
\begin{equation}\label{5.4:eq-bundle-sum-2}
\sum_{Y\ni Q_*} e^{r d_{k,Z}(Y)}\,\Bigl|\sum_{p:\ Y(\mathrm{lab}_1 p)=Y} v_p\Bigr|,
\end{equation}
in which the inner sum runs over the pieces $v_p$ of the long units of $\mathfrak{L}^{\mathrm{lg}}_d$
whose label has region $Y$.
\end{theorem}

\begin{proof}
Fix the bracket index $s$. We apply the label-sum theorem at the grain
$\mathfrak{g}:=\pi^{\mathfrak{s}}$ to the family
\begin{equation}\label{5.4:eq-bundle-sum-3}
\mathbb{D}:=\bigl\{(\mathrm{St}_v,\,K_v,\,2\mathfrak{N}^{\flat}(v)) :
v\in\mathfrak{L}^{\mathrm{lg}}_d\bigr\}.
\end{equation}
Here $K_v$ is as in Lemma~\ref{5.4:lem-rim-mass}. The family $\mathbb{D}$ contains the units of all
levels $j_v\le k$ at once. Applying the label-sum
theorem separately to the units of each level $j\le k$ and summing over $j$ would give the bound
$(k+1)\mathsf{K}^{\mathrm{lg}}$, which grows with $k$; this is not done.

The hypotheses of the label-sum theorem hold for $\mathbb{D}$. By clause (f) of
Lemma~\ref{5.4:lem-cells-carrier}, for every $v\in\mathfrak{L}^{\mathrm{lg}}_d$ the set
$\mathrm{St}_v$ is finite, wall-connected and sourceless, and its coarse hull satisfies
$[\mathrm{St}_v]_k\subset K_v\in\mathbf{D}_k$. We take the parameter $\beta:=3$, which satisfies
$\beta\le\lambda'_{\Sigma}-2$. By clauses (c) and (e) of Lemma~\ref{5.4:lem-bundle-frame}, the
bundles with bracket index $s$ which do not vanish have $\mathsf{y}\in\mathfrak{Y}(\mathrm{St}_v)$,
and for them $\mathrm{lab}(\mathfrak{b}^{\circ})$ coincides with the label which the label-sum
theorem attaches to $(\mathrm{St}_v,K_v)$ and $\mathsf{y}$. Moreover, with
$\mu_v:=2\mathfrak{N}^{\flat}(v)$ and the norm $\|\mathfrak{b}^{\circ}\|$ of
\eqref{5.4:eq-long-bundles-a}, every such bundle satisfies
\begin{equation}\label{5.4:eq-bundle-sum-4}
\|\mathfrak{b}^{\circ}\|\le \mu_v\,e^{-(\kappa_1-2)|\mathsf{y}|}.
\end{equation}
These are the hypotheses of the label-sum theorem at grain $\pi^{\mathfrak{s}}$, with weights
$\mu_v$ and parameter $\beta=3$. Its conclusion bounds the slice of the left-hand side of
\eqref{5.4:eq-bundle-sum-1} at the fixed index $s$:
\begin{equation}\label{5.4:eq-bundle-sum-5}
\sum_{\substack{\mathfrak{b}^{\circ}=(v;s;\mathsf{y}):\ v\in\mathfrak{L}^{\mathrm{lg}}_d,\\
\mathsf{y}\in\mathfrak{Y}(\mathrm{St}_v),\ \mathrm{lab}(\mathfrak{b}^{\circ})\ni Q_*}}
\|\mathfrak{b}^{\circ}\|\,e^{r d_k(\mathrm{lab}\,\mathfrak{b}^{\circ})}
\le S_d\,\mathsf{M}^{\partial}_{r+1,3}[\mathbb{D}].
\end{equation}

We bound the rim supremum by Lemma~\ref{5.4:lem-rim-mass} at $r':=r+1$. That lemma requires
$r'\le\tfrac85\kappa$. Since $r\le a+8$ and $a=\tfrac{31}{20}\kappa$,
\begin{equation*}
r'\le a+9=\tfrac{31}{20}\kappa+9\le\tfrac{31}{20}\kappa+\tfrac{1}{20}\kappa=\tfrac85\kappa,
\end{equation*}
where the second inequality is $9\le\kappa/20$, that is, $\kappa\ge 180$. The left-hand side of
Lemma~\ref{5.4:lem-rim-mass} carries the rim factor $e^{-3\rho(\Delta_1)}$ of that lemma and the
weights
$2\mathfrak{N}^{\flat}(v)=\mu_v$, that is, $\beta=3$ and the weights of $\mathbb{D}$, and the lemma
holds for every $\Delta_1\in\pi^{\mathfrak{s}}$. It therefore gives
\begin{equation*}
S_d\,\mathsf{M}^{\partial}_{r+1,3}[\mathbb{D}]
\le S_d\,e^{-(9\check{R}_k-2)}\cdot 4e\,dL^{d-1}\bigl(1+C^{\mathrm{win}}_{\mathrm{lg}}\bigr)
C_{\mathsf{b}}\,\mathfrak{s}_{\mathsf{b}}\,\mathfrak{K}\,\mathsf{b}^{\mathrm{lg}}(0),
\end{equation*}
with $C^{\mathrm{win}}_{\mathrm{lg}}$ as in \eqref{5.4:eq-transport-constants-d} and
$C_{\mathsf{b}}$, $\mathfrak{s}_{\mathsf{b}}$, $\mathfrak{K}$ as in
\eqref{5.4:eq-transport-constants-c}.

Finally we sum over the four values of $s$. This multiplies the bound by $4$, and the constant
becomes
\begin{align*}
4\cdot S_d\cdot 4e\,dL^{d-1}\bigl(1+C^{\mathrm{win}}_{\mathrm{lg}}\bigr)
C_{\mathsf{b}}\,\mathfrak{s}_{\mathsf{b}}\,\mathfrak{K}
&=16e\,dL^{d-1}S_d\bigl(1+C^{\mathrm{win}}_{\mathrm{lg}}\bigr)
C_{\mathsf{b}}\,\mathfrak{s}_{\mathsf{b}}\,\mathfrak{K}\\
&=\mathsf{C}^{\mathrm{lg},d}_{\circ},
\end{align*}
the last equality being the definition of $\mathsf{C}^{\mathrm{lg},d}_{\circ}$ in
\eqref{5.4:eq-long-unit-sums-a}. This proves \eqref{5.4:eq-bundle-sum-1}.
\end{proof}

\begin{corollary}[Conversion of the bundle sum into the small factor]
Assume the hypotheses of Theorem~\ref{5.4:thm-bundle-sum}, the conversion corollary for the small
factor $\theta'(g_k)$, and
\begin{equation*}
\mathsf{b}^{\mathrm{lg}}(0)\le\mathsf{b}^{\mathrm{lg},\natural}\,\check{R}_k .
\end{equation*}
Then the left-hand side of \eqref{5.4:eq-bundle-sum-1} is bounded by
\begin{equation*}
e^{-(\check{R}_k-2)}\,\mathsf{C}^{\mathrm{lg},d}
\le \mathsf{c}_{\theta'}\,\theta'(g_k)\,\mathsf{C}^{\mathrm{lg},d},
\end{equation*}
with $\mathsf{C}^{\mathrm{lg},d}$ as in \eqref{5.4:eq-long-unit-sums-a}.
\end{corollary}

\begin{proof}
We have $e^{-(9\check{R}_k-2)}=e^{-8\check{R}_k}e^{-(\check{R}_k-2)}$. With the assumed bound on
$\mathsf{b}^{\mathrm{lg}}(0)$, the right-hand side of \eqref{5.4:eq-bundle-sum-1} is at most
\begin{equation*}
e^{-(\check{R}_k-2)}\,\bigl(\check{R}_k e^{-8\check{R}_k}\bigr)\,
\mathsf{C}^{\mathrm{lg},d}_{\circ}\,\mathsf{b}^{\mathrm{lg},\natural}.
\end{equation*}
For every real $x$ one has $x e^{-8x}\le 1/(8e)$, the maximum being attained at $x=1/8$. Hence
$\check{R}_k e^{-8\check{R}_k}\le 1/(8e)\le 1$. Since
$\mathsf{C}^{\mathrm{lg},d}=\mathsf{C}^{\mathrm{lg},d}_{\circ}\mathsf{b}^{\mathrm{lg},\natural}$ by
\eqref{5.4:eq-long-unit-sums-a}, the sum is at most $e^{-(\check{R}_k-2)}\mathsf{C}^{\mathrm{lg},d}$.
The conversion corollary for the small factor gives
$e^{-(\check{R}_k-2)}\le\mathsf{c}_{\theta'}\,\theta'(g_k)$, which is the second inequality.
\end{proof}

\begin{corollary}[Label reading for long units off the treated class]
Assume the hypotheses of Theorem~\ref{5.4:thm-bundle-sum} and the label-reading corollary. Then the
label-reading corollary holds for the family $\mathfrak{L}^{\mathrm{lg}}_d$. In it the triple
$(\mathrm{lab}\,\mathfrak{b}^{\circ},\,Y_0(\mathfrak{b}^{\circ}),\,\|\mathfrak{b}^{\circ}\|)$ enters
the linked-sum estimate at its second, third and fourth links, and the following condition (R)
holds.
\begin{enumerate}
\item[(R)] Let a member of $\mathsf{y}_t$ meet $Y_0(\mathfrak{b}^{\circ})\setminus\Lambda^{\sim}$.
Then the cube of $\pi_k$ containing it lies in $[\mathrm{St}_v]_k\cup[\mathsf{y}]_k$, and
\begin{equation*}
[\mathrm{St}_v]_k\cup[\mathsf{y}]_k\subset\mathrm{lab}(\mathfrak{b}^{\circ}).
\end{equation*}
\end{enumerate}
\end{corollary}

\begin{proof}
The label-reading corollary is applied to the family $\mathfrak{L}^{\mathrm{lg}}_d$. Its input at
the second, third and fourth links of the linked-sum estimate is supplied by the label
$\mathrm{lab}\,\mathfrak{b}^{\circ}$, the set $Y_0(\mathfrak{b}^{\circ})$ and the norm
$\|\mathfrak{b}^{\circ}\|$ of each bundle. The condition it requires beyond these is (R). A member
of $\mathsf{y}_t$ meeting $Y_0(\mathfrak{b}^{\circ})$ outside the hubs $\Lambda^{\sim}$ has its
level-$k$ parent in $[\mathrm{St}_v]_k\cup[\mathsf{y}]_k$. Here $[\mathrm{St}_v]_k\subset K_v$ by
clause (f) of Lemma~\ref{5.4:lem-cells-carrier}, and the label $\mathrm{lab}(\mathfrak{b}^{\circ})$
of \eqref{5.4:eq-long-bundles-a}, which is formed from $K_v$ and $\mathsf{y}$, contains
$[\mathrm{St}_v]_k\cup[\mathsf{y}]_k$. With these data the label-reading corollary gives its
conclusion for $\mathfrak{L}^{\mathrm{lg}}_d$.
\end{proof}

\begin{remark}
The pieces $v_p$ of the long units $v\in\mathfrak{L}^{\mathrm{lg}}_d$ enter the estimates above
only through the bundles $\mathfrak{b}^{\circ}$ in which they are grouped. The gain factor
$\mathsf{w}_L^{-1}(\eta_p\mathsf{f}_v)^{1/2}$ of the
individual bound of a piece is not used on these units. Accordingly, the dilation condition, the
block decay $\eta_p$, the cell weights $\mathsf{f}_v$ and the cell-cover lemma for the blocks
$\square$ are not hypotheses of Theorem~\ref{5.4:thm-bundle-sum}. Of the size conditions on the
pieces, only the size clause of the conditions displayed in \eqref{5.4:eq-long-bundles-b} is used.
\end{remark}

\begin{remark}[Foreign windows]\label{5.4:rem-foreign-windows}
Let $\mathfrak{a}'$ be a live arrest on a component not in $Z$, with window $W_{\mathfrak{a}'}$ and
level $k_{a'}$. This covers a sibling arrest, with $k_{a'}=k$, and arrests nested inside other windows.
Four facts hold about long singles and tails (the units of $\mathfrak{L}^{\mathrm{lg}}$) lying inside
$W_{\mathfrak{a}'}$.
\begin{enumerate}
\item[(1)] The per-cube transport (Proposition~\ref{5.4:prop-per-cube-transport}) is insensitive to the
window. Such a unit is counted at its own cube inside the cell it meets, against the factor
$e^{-c_{\mathrm{sp}}\mathfrak{m}_u}$ of that proposition's proof, and no factor
$(L\check{R}_{k_{a'}})^d$ arises.
\item[(2)] The per-cell transport (Proposition~\ref{5.4:prop-per-cell-transport}) is likewise insensitive
to the window. The count runs over filing cells $\Delta'$ (the cells at which units are filed), at the
factor $e^{-\theta m(\Delta')}$ of that proposition's proof, and never over cubes.
\item[(3)] Suppose the bound carrying the age function $\mathsf{b}^{\mathrm{lg}}$ is supplied with the
weight $W^{\mathrm{win}}(\check{R}_{k_{a'}})$ at raised cubes or cells (in the sense of
Lemma~\ref{5.4:lem-window-placement}), and let $\rho$ be the function of the collar lemma
(Lemma~\ref{5.4:lem-collar}). Then that weight is paid in one of the following ways:
\begin{itemize}
\item if $\Delta_2$ is a current cell (in the sense of the collar lemma) contained in
$W_{\mathfrak{a}'}$ and $\Delta_1\in\pi^{\mathfrak{s}}$ shares a wall with $\Delta_2$, by a single
factor $e^{-\rho(\Delta_1)}$ of the consuming estimate, taken across the collar of that arrest
(Lemma~\ref{5.4:lem-collar}); or
\item if the unit is $v\in\mathfrak{L}^{\mathrm{lg}}_{d}$, with region $X_v$ and length $d_v$,
$\Delta_2\in\mathrm{St}_v$ is a cell meeting $X_v$, $\Delta_1\in\partial\mathrm{St}_v$ is a rim cell
sharing a wall with $\Delta_2$, and the exponent $c$ of the collar hypothesis of
Proposition~\ref{5.4:prop-per-cell-transport} is at least $3$, by the unit's own length together
with the lower bound $\rho(\Delta_1)\ge\max\{\delta-1,1\}$, where $\delta$ is the distance parameter
of that bound, the product being absorbed into the displayed supremum
$C^{\mathrm{win}}_{\mathrm{lg}}$ (Lemma~\ref{5.4:lem-window-placement}); or
\item at belt cells (the cells of Lemma~\ref{5.4:lem-boundary-belt}), by the exit length, into the
constants $\mathfrak{G}_p$ and $\mathfrak{G}'_p$.
\end{itemize}
The weight of an own window inside a core is paid by the unit's length across the collar of that
window.
\item[(4)] No window radius enters any named constant, except inside the displayed suprema
$C^{\mathrm{win}}_{\mathrm{lg}}$, $C^{\mathrm{win}}_{B'}$, $\mathfrak{G}_p$, $\mathfrak{G}'_p$ and
$\sigma^{\mathrm{lg}}$. The running weight $W^{\mathrm{win}}(\check{R}_k)$, if present, stands outside
every named constant and is absorbed by the residue of the corollary on the small factor
$\theta'(g_k)$.
\end{enumerate}
\end{remark}

\begin{proof}
(1) In the proof of Proposition~\ref{5.4:prop-per-cube-transport}, hypothesis (A3) enters only through
its lower bound ``$\ge 1$''. This bound holds in the frozen strata of the window $W_{\mathfrak{a}'}$,
and it holds with equality on a prepared staircase of one layer. No other step of that proof reads
whether the cube lies inside a window. The count is therefore not taken over the cubes of the window,
and no factor $(L\check{R}_{k_{a'}})^d$ forms, for sibling arrests with $k_{a'}=k$ and for nested
arrests alike.

(2) The count in the proof of Proposition~\ref{5.4:prop-per-cell-transport} is per filing cell, so a
window contributes no count of its cubes, as in (1).

(3) Suppose the weight $W^{\mathrm{win}}(\check{R}_{k_{a'}})$ sits at a raised cube or cell. It is
paid in one of three ways, according to where it sits; own windows inside cores are a separate matter.

First, let $\Delta_2$ be a current cell with $\Delta_2\subset W_{\mathfrak{a}'}$ and let
$\Delta_1\in\pi^{\mathfrak{s}}$ share a wall
with $\Delta_2$. The collar lemma (Lemma~\ref{5.4:lem-collar}) gives
\begin{equation*}
\rho(\Delta_1)\ \ge\ 2\check{R}_{k_{a'}}-1 .
\end{equation*}
The collar lemma is stated for foreign and sibling arrests alike, so it applies when $k_{a'}=k$. Its
proof uses only the zoning of levels about the component $X_{\mathfrak{a}'}$ of the arrest, one
property of the level structure of the cells, and one lower bound for the function $\rho$; none of
these three refers to a window enclosing $W_{\mathfrak{a}'}$, so the bound holds inside
nested windows as well. One factor $e^{-\rho(\Delta_1)}$ of the consuming estimate thus pays the weight.

Second, let $v\in\mathfrak{L}^{\mathrm{lg}}_{d}$ be the unit, $X_v$ its region, $d_v$ its length,
$\Delta_2\in\mathrm{St}_v$, $\Delta_1\in\partial\mathrm{St}_v$ sharing a wall with $\Delta_2$, and
$c\ge 3$ the exponent $c$ of the collar hypothesis of
Proposition~\ref{5.4:prop-per-cell-transport}. In the setting of that proposition the weight may be
placed at a raised cell in closed contact with $X_v$ lying in $W_{\mathfrak{a}'}$. The window-placement
lemma (Lemma~\ref{5.4:lem-window-placement}), used together with the lower bound
$\rho(\Delta_1)\ge\max\{\delta-1,1\}$ for the rim cell $\Delta_1$, combines the unit's own length with
that bound and gives, for $c\ge 3$,
\begin{equation*}
W^{\mathrm{win}}(\check{R}_{k_{a'}})\,e^{-\rho(\Delta_1)}\,e^{-(c-2)d_v}\ \le\ C^{\mathrm{win}}_{\mathrm{lg}} .
\end{equation*}

Third, at belt cells the weight is paid by the exit length. This is clause (iv) of the belt lemma
(Lemma~\ref{5.4:lem-boundary-belt}), with the constants $\mathfrak{G}_p$ and $\mathfrak{G}'_p$.

Finally, own windows inside cores are a separate matter: the units concerned are those of
$\mathfrak{L}^{\mathrm{lg}}_{c}$ treated in Proposition~\ref{5.4:prop-own-window-weights}, where
the weight is paid by the unit's length across the collar of its own window.

(4) By (1) and (2), the constants $C^{LS}$ and $C^{\mathrm{gen}}_{\mathsf{b}}$ of the two transport
propositions read no window quantity. The only named constants carrying a window radius are the
displayed suprema $C^{\mathrm{win}}_{\mathrm{lg}}$, $C^{\mathrm{win}}_{B'}$, $\mathfrak{G}_p$,
$\mathfrak{G}'_p$, $\sigma^{\mathrm{lg}}$; of these, $C^{\mathrm{win}}_{\mathrm{lg}}$ is the supremum
of Lemma~\ref{5.4:lem-window-placement}, and $\mathfrak{G}_p$, $\mathfrak{G}'_p$ are the constants of
clause (iv) of Lemma~\ref{5.4:lem-boundary-belt}. The running weight $W^{\mathrm{win}}(\check{R}_k)$ of
Remark~\ref{5.4:rem-running-window}, when present, is not placed in any of these constants. It is
absorbed by the residue of the corollary on $\theta'(g_k)$.
\end{proof}

\begin{remark}[One value of the total constant and the order of choice of constants]
\label{5.4:rem-constant-order}
Let $\mathfrak{K}$ be the transport constant of \eqref{5.4:eq-transport-constants-c}: it is
$C^{LS}$ if the base estimate is taken from the first of the two transports, and
$C^{\mathrm{gen}}_{\mathsf{b}}\,\mathsf{K}^{\flat}_{\mathrm{lg}}$ if it is taken from the
second; one of the two is used, never both. Put
$\mathsf{K}^{\mathrm{lg}} := \mathfrak{K}\,\mathsf{b}^{\mathrm{lg}}(0)$.

In the first of the two ways of summing the long tails of class (d), namely by the pointed-piece
sum for long units off the treated class, the
total constant is
\begin{equation}\label{5.4:eq-constant-order-1}
\begin{aligned}
\mathsf{C}_1^{\mathrm{ex}} &:= \mathsf{C}^{\Sigma,4}+\mathsf{C}^{\Sigma}_{a+7}
  +\mathsf{C}^{\Sigma,\mathrm{lg}},\\
\mathsf{C}^{\Sigma,4} &:= \mathsf{c}_{\theta'}\bigl[(\mathsf{C}^{4,d}+\mathsf{C}^{4,c})
  \mathsf{b}^{\natural}+\mathsf{C}^{\mathrm{lg},c}\bigr],\\
\mathsf{C}^{\mathrm{lg},c} &:= 2e\,S_d^{(1/2)}\,\mathfrak{K}\,\mathsf{b}^{\mathrm{lg},\natural}
  \,\mathfrak{W},\\
\mathsf{C}^{\Sigma}_{a+7} &:= 8S_d\,\mathsf{M}^{\mathrm{sh}}_{a+8},\\
\mathsf{C}^{\Sigma,\mathrm{lg}} &:= 80e\,dL^{d-1}S_d\,C^{0}_{\square}\,\mathsf{c}_d^{2}\,
  \mathsf{c}_{\mathsf{f}}\,(1+C^{\mathrm{win}}_{\mathrm{lg}})\,C_{\mathsf{b}}\,
  \mathfrak{s}_{\mathsf{b}}\,\mathfrak{K}\,\mathsf{b}^{\mathrm{lg},\natural}.
\end{aligned}
\end{equation}
Here $\mathfrak{W}$ is the in-$Z$ weight factor of part (i) of the corollary on long units in
the boundary-term sums (Corollary~\ref{5.4:cor-long-unit-sums}), and
$\mathsf{C}^{\Sigma}_{a+7}$ is the constant of the corollary of the label-sum theorem that
bounds the label sum by means of the rim supremum $\mathsf{M}^{\mathrm{sh}}_{a+8}$. In
the alternative of
Corollary~\ref{5.4:cor-long-unit-sums}(i) in which the $\check{R}_k$-weight is displayed
outside the sum, $\mathsf{C}^{\mathrm{lg},c}$ carries the further factor $\sigma^{\mathrm{lg}}$;
it is still one value.

In the second way, by the bundle sum for
long units off the treated class (Theorem~\ref{5.4:thm-bundle-sum}), the constant
$\mathsf{C}^{\Sigma,\mathrm{lg}}$ does not occur and
\begin{equation}\label{5.4:eq-constant-order-2}
\begin{aligned}
\mathsf{C}_1^{\mathrm{ex}} &:= \mathsf{C}^{\Sigma,4}+\mathsf{C}^{\Sigma}_{a+7},\\
\mathsf{C}^{\Sigma,4} &:= \mathsf{c}_{\theta'}\bigl[(\mathsf{C}^{4,d}+\mathsf{C}^{4,c})
  \mathsf{b}^{\natural}+\mathsf{C}^{\mathrm{lg},c}+\mathsf{C}^{\mathrm{lg},d}\bigr],\\
\mathsf{C}^{\mathrm{lg},d} &:= 16e\,dL^{d-1}S_d\,(1+C^{\mathrm{win}}_{\mathrm{lg}})\,
  C_{\mathsf{b}}\,\mathfrak{s}_{\mathsf{b}}\,\mathfrak{K}\,\mathsf{b}^{\mathrm{lg},\natural},
\end{aligned}
\end{equation}
with $\mathsf{C}^{\mathrm{lg},c}$ and $\mathsf{C}^{\Sigma}_{a+7}$ as in
\eqref{5.4:eq-constant-order-1}.

The constants are chosen in the following order. First the numbers depending on $d$ only,
among them $\mathsf{c}_d$, $\mathfrak{a}_d$, $\mathfrak{e}_d$, $S_d$, $S_d^{(1/2)}$. Then $L$,
and the numbers depending on $(d,L)$: $C^{0}_{\square}$, $\hat{b}$, $\hat{b}_t$. Then
$\kappa_1$ and $\kappa$, and with them $\lambda_{\Sigma}$ and $\epsilon_{\partial}$. Then the
growth constants $C_{\mathsf{b}}$, $p_{\mathsf{b}}$ of the age function together with
$C_{\uparrow}$ and $\mathsf{L}_0^{w}$, and with them $C^{LS}$, $C^{\mathrm{gen}}_{\mathsf{b}}$,
$\mathfrak{s}_{\mathsf{b}}$, $S_{\mathsf{b}}$, $C^{\mathrm{win}}_{\mathrm{lg}}$,
$C^{\mathrm{win}}_{B'}$, $\mathfrak{G}_p$, $\mathfrak{G}'_p$, $\sigma^{\mathrm{lg}}$. Then $M$,
and with it $\mathsf{K}^{\flat}_{\mathrm{lg}}$, $\mathsf{b}^{\mathrm{lg}}(0)$,
$\mathsf{b}^{\mathrm{lg},\natural}$, $c_{\chi}$, $\mathsf{C}_{fw}$, $\mathsf{b}^{\natural}$,
$\mathsf{M}^{\mathrm{sh}}$. Then $C_{\theta'}$, and with it $\mathsf{c}_{\theta'}$ and
$\mathsf{C}^{\Sigma,4}$. Last $\mathsf{C}_1^{\mathrm{ex}}$. All of this precedes the choice of
$\gamma$; the coupling $g$ is chosen last and is read by none of these constants, entering only
through the factor $\theta'(g_k)$.
\end{remark}

\begin{proof}
In either case $\mathsf{C}_1^{\mathrm{ex}}$ is a finite sum of finite products of constants,
so it is one value as soon as each factor is. The factor $\mathfrak{K}$ is one number, because
exactly one of $C^{LS}$ and $C^{\mathrm{gen}}_{\mathsf{b}}\mathsf{K}^{\flat}_{\mathrm{lg}}$ is
taken; hence $\mathsf{K}^{\mathrm{lg}}$, $\mathsf{C}^{\mathrm{lg},c}$,
$\mathsf{C}^{\Sigma,\mathrm{lg}}$ and $\mathsf{C}^{\mathrm{lg},d}$ are single numbers (compare
\eqref{5.4:eq-transport-constants-c}, \eqref{5.4:eq-transport-constants-d} and
\eqref{5.4:eq-long-unit-sums-a}). Multiplying $\mathsf{C}^{\mathrm{lg},c}$ by
$\sigma^{\mathrm{lg}}$ in the outside-weight alternative replaces one number by another. The two
cases are not added: in \eqref{5.4:eq-constant-order-1} the long tails of class (d) contribute
$\mathsf{C}^{\Sigma,\mathrm{lg}}$ and $\mathsf{C}^{\mathrm{lg},d}$ does not appear, while in
\eqref{5.4:eq-constant-order-2} they contribute $\mathsf{C}^{\mathrm{lg},d}$ inside
$\mathsf{C}^{\Sigma,4}$ and $\mathsf{C}^{\Sigma,\mathrm{lg}}$ does not appear. Thus in each
case the total constant is one value.

For the order, each constant named in the order above is fixed at the stage at which the order
places it. The factor $\mathsf{c}_{\theta'}$ is fixed with $C_{\theta'}$, so
$\mathsf{C}^{\Sigma,4}$ is fixed at that stage, and $\mathsf{C}_1^{\mathrm{ex}}$, a sum of
$\mathsf{C}^{\Sigma,4}$ with constants already fixed, is fixed after it. No factor in
\eqref{5.4:eq-constant-order-1} or \eqref{5.4:eq-constant-order-2} involves $\gamma$ or $g$; the
coupling appears in the bounds only through the factor $\theta'(g_k)$ that multiplies these
constants.
\end{proof}

\begin{remark}[The touch reading of meeting]\label{5.4:rem-touch-reading}
Two readings of the word ``meet'' are possible for a cell, or a member of a grain, and the carrier
$X_v$ of a unit $v$. In
reading $(\mathrm{O})$ a cell meets $X_v$ if the two have a common interior point; in reading
$(\mathrm{T})$ it meets $X_v$ if the two touch, that is, have a common point. Let $K_v$ be the
family of $\pi_k$-cubes against which the label-sum theorem forms the label
$\mathrm{lab}(K_v,\cdot)$. Put $j:=j_v$, write $\mathrm{St}^{j}_v$, the \emph{star} of $v$, for the
family of members of the level-$j$ grain $\pi^{\mathfrak{s}}_j$ of the label-sum theorem that
meet $X_v$ under reading $(\mathrm{O})$, and write
$\mathrm{St}^{j,T}_v$ for the family of members of $\pi^{\mathfrak{s}}_j$ that touch $X_v$.
Throughout this section ``meet'' is taken in reading $(\mathrm{O})$.
\begin{enumerate}
\item Under reading $(\mathrm{O})$ the conclusions of the label-sum theorem and of its corollaries
hold, because
\begin{equation*}
[\mathrm{St}^{j}_v]_k\subset K_v .
\end{equation*}
\item Under reading $(\mathrm{T})$ the conclusion of the label-sum theorem does not hold in general
for admissible units: the label of a bundle family entering it need not lie in $\mathbf{D}_k$.
Precisely: let $v$ be an admissible unit with $j_v=k$ whose carrier is a single cube
$X_v=\{Q\}$ of the level-$j$ grain, let $e_1,e_2$ be the positive unit vectors of two distinct
coordinate directions of that grain, and
suppose that the coarse hull of the cube $G:=Q+e_1+e_2$ does not lie in $K_v$. Suppose further
that an admissible walk enters $\mathrm{St}^{j,T}_v$ through $G$ with members
$E,E+e_1,\dots,E+ne_1$ for some $n\ge 0$, where $E:=Q+2e_1+e_2$, that this walk is removed by no
clause of the filing rule other than clause $(c)$, and that no cube of the hub set
$\mathrm{Hub}(\tilde{\mathsf{y}})$ of the bundle family $\tilde{\mathsf{y}}$ built from it shares
a wall with $Q$. Then $\mathrm{lab}(K_v,\tilde{\mathsf{y}})\notin\mathbf{D}_k$.
\item Let $v$ be an admissible unit with $j_v<k$, and let $X_v=\{Q\}$, $e_1,e_2$, $G$, $E$, the
walk and $\tilde{\mathsf{y}}$ be as in the second item, the walk being removed by no clause of the
filing rule other than clause $(c)$. Suppose that $[G]_k\notin K_v$, and that none of the coarse
hulls $[E+ie_1]_k$, $0\le i\le n$, and $[H]_k$, $H\in\mathrm{Hub}(\tilde{\mathsf{y}})$, equals or
shares a wall with $[Q]_k$. Then $\mathrm{lab}(K_v,\tilde{\mathsf{y}})\notin\mathbf{D}_k$.
\end{enumerate}
Under reading $(\mathrm{T})$ two further assertions fail as stated: a piece that cuts far cells
of a partially exempt member lies in no bundle of Definition~\ref{5.4:def-long-bundles}, and the
assertion of the corollary of the label-sum theorem that the parent lies in the label does not
hold. For all three see the label-sum theorem under the touch reading.
\end{remark}

\begin{proof}
We construct the example for the second item. Let $v$ be an admissible unit as in the second
item, $j=j_v=k$, whose carrier is a
single cube $X_v=\{Q\}$ of the level-$j$ grain, and let $e_1,e_2$ be the positive unit vectors of
two distinct coordinate directions of
that grain, so that $Q+e_1$, $Q+e_1+e_2$ and $Q+2e_1+e_2$ are again cubes of the grain. Since
$j=k$, every cube of the level-$j$ grain is its own coarse hull $[\cdot]_k$, so the label
$\mathrm{lab}(K_v,\cdot)$ is formed from these cubes themselves.

Put
\begin{equation*}
G:=Q+e_1+e_2 .
\end{equation*}
The cubes $Q$ and $G$ differ by one step in each of two coordinate directions. Hence they have
common points but no common interior point, and they share no wall. Thus $G$ is a member of
$\mathrm{St}^{j,T}_v$, in diagonal position to $Q$. By the hypothesis of the second item, the
coarse hull of $G$ does not lie in $K_v$,
\begin{equation*}
[G]_k\notin K_v .
\end{equation*}
The cube $Q+e_1$ shares a wall with $Q$, so it touches $X_v$. It also shares a wall with $G$.
The cube $Q$ itself touches $X_v=\{Q\}$, so it is a member of $\mathrm{St}^{j,T}_v$.
Hence $Q$, $Q+e_1$ and $G$ lie in one wall-component of $\bigcup\mathrm{St}^{j,T}_v$.

Consider now the admissible walk of the second item, which enters $\mathrm{St}^{j,T}_v$ through
$G$ and whose first member is
\begin{equation*}
E=Q+2e_1+e_2 .
\end{equation*}
The cube $E=G+e_1$ shares a wall with $G$, so $E$ is walled to $G$. The offset of $E$ from $Q$
has coordinates $2$ and $1$ in the directions $e_1$ and $e_2$. Hence $E$ is at
$\infty$-distance $1$ from $Q$ and does not touch $Q$.

The entry member $G$ lies in the wall-component of $\bigcup\mathrm{St}^{j,T}_v$ that contains
$Q$ and $Q+e_1$. The walk therefore belongs to the family $\mathfrak{Y}_v$ of admissible walks
attached to $v$ by the filing rule, and it is not removed by clause $(c)$ of that rule; by the
hypothesis of the second item it is removed by no other clause. The bundle
family $\tilde{\mathsf{y}}$ built from this walk therefore
enters the label-sum theorem, and its label is
\begin{equation*}
\mathrm{lab}(K_v,\tilde{\mathsf{y}})
=\{Q\}\cup\{E,E+e_1,\dots,E+ne_1\}\cup\mathrm{Hub}(\tilde{\mathsf{y}}).
\end{equation*}
This family is not connected. For $0\le i\le n$ the offset of the member $E+ie_1$ from $Q$ has
coordinates $2+i$ and $1$ in the directions $e_1$ and $e_2$, so $E+ie_1$ is at
$\infty$-distance $1+i\ge 1$ from $Q$, does not touch $Q$, and in particular does not share a
wall with $Q$; by the hypothesis of the second item no cube of $\mathrm{Hub}(\tilde{\mathsf{y}})$
shares a wall with $Q$; and the
cube $G$, through which the walk entered $\mathrm{St}^{j,T}_v$, is not in the label because
$[G]_k\notin K_v$. Hence no cube of the label other than $Q$ itself shares a wall with $Q$, while
the label contains the cube $E\neq Q$. Consequently
\begin{equation*}
\mathrm{lab}(K_v,\tilde{\mathsf{y}})\notin\mathbf{D}_k ,
\end{equation*}
since $\mathbf{D}_k$ consists of wall-connected finite families of $\pi_k$-cubes.

For the third item, let $j=j_v<k$. The facts used above to place the walk in $\mathfrak{Y}_v$ ---
that $G$ touches $Q$, that $Q+e_1$ shares a wall with $Q$ and with $G$, that $E$ shares a wall
with $G$ and does not touch $Q$ --- are statements about cubes of the level-$j$ grain and hold
at every level $j$; with the hypothesis of the third item on the other clauses of the filing rule,
the bundle family $\tilde{\mathsf{y}}$ enters the label-sum theorem as before. Read through coarse
hulls, its label is
\begin{equation*}
\mathrm{lab}(K_v,\tilde{\mathsf{y}})
=\{[Q]_k\}\cup\{[E+ie_1]_k:0\le i\le n\}
\cup\{[H]_k:H\in\mathrm{Hub}(\tilde{\mathsf{y}})\}.
\end{equation*}
If $[G]_k$ is not one of the hulls $[E+ie_1]_k$, it does not occur in this family, since
$[G]_k\notin K_v$; if it is one of them, it is covered by the hypothesis of the third item on
those hulls. By the hypothesis of the third item no hull in
this family other than $[Q]_k$ equals or shares a wall with $[Q]_k$, while the family contains
$[E]_k\neq[Q]_k$. Hence $\mathrm{lab}(K_v,\tilde{\mathsf{y}})$ is not wall-connected, and
$\mathrm{lab}(K_v,\tilde{\mathsf{y}})\notin\mathbf{D}_k$.

The hypotheses of the second item are met at $j=k$ by a unit of class $(d)$ entering the
summation of boundary terms at finer zones, and those of the third item by a long unit of class
$(d)$, a member of $\mathfrak{L}^{\mathrm{lg}}_{d}$, at $j<k$. For such admissible units the
conclusion of the label-sum theorem does not hold under reading $(\mathrm{T})$.

Under reading $(\mathrm{O})$ the member $G$ does not meet $X_v$, since it has no common interior
point with $Q$, so $G$ is not a member of the star $\mathrm{St}^{j}_v$. Under that reading the
coarse hull of every member of the star lies in $K_v$, $[\mathrm{St}^{j}_v]_k\subset K_v$; thus
the situation of the second item, a member of the star $\mathrm{St}^{j}_v$ whose coarse hull
does not lie in $K_v$, does not arise, and the first item holds.
\end{proof}

\begin{theorem}[The label sum with attachments under the touch reading]\label{5.4:thm-touch-label-sum}
Fix the level $k$, the grain $\mathfrak{g}$ and the large-field region $Z$, and let $\mathbb{D}$ be a family
of data $\mathsf{u}=(\mathrm{St}_{\mathsf{u}},\mathrm{St}^{\mathrm{cut}}_{\mathsf{u}},A_{\mathsf{u}},\mu_{\mathsf{u}})$ with
$\mathrm{St}^{\mathrm{cut}}_{\mathsf{u}}\subset\mathrm{St}_{\mathsf{u}}$, where $\mathrm{St}_{\mathsf{u}}\subset\mathfrak{g}$ is a
finite family of cells holding no source cell, $A_{\mathsf{u}}$ is a family of $\pi_k$-cubes and
$\mu_{\mathsf{u}}$ is a weight. Assume the \emph{anchoring condition}:
\begin{itemize}
\item for every $\Delta\in\mathrm{St}_{\mathsf{u}}$ the cube $[\Delta]_k$ belongs to $A_{\mathsf{u}}$ or touches a
cube of $A_{\mathsf{u}}$.
\end{itemize}
(It holds with $A_{\mathsf{u}}=K_v$, the label set $K_v$ of a unit $v$, for any star of members touching the
region $X_v$ of $v$.) Put
$R_{\mathsf{u}}:=\partial\mathrm{St}_{\mathsf{u}}\cup\mathrm{St}^{\mathrm{cut}}_{\mathsf{u}}$, where
$\partial\mathrm{St}_{\mathsf{u}}$ is the rim of $\mathrm{St}_{\mathsf{u}}$ in $\mathfrak{g}$, and let
$\mathfrak{Y}^{\mathfrak{g}}_{T}(\mathsf{u})$ be the set of finite families
$\tilde{\mathsf{y}}\subset(\mathfrak{g}\setminus\mathrm{St}_{\mathsf{u}})\cup\mathrm{St}^{\mathrm{cut}}_{\mathsf{u}}$ which
hold a source cell and each component of which holds a member of $R_{\mathsf{u}}$.

For each component $\tilde{\mathsf{y}}_a$ of $\tilde{\mathsf{y}}$ fix by a rule a member
$\Delta_a\in\mathrm{St}_{\mathsf{u}}$ through which $\tilde{\mathsf{y}}_a$ enters, either by a wall or with
$\Delta_a\in\tilde{\mathsf{y}}_a\cap\mathrm{St}^{\mathrm{cut}}_{\mathsf{u}}$; put $P_a:=[\Delta_a]_k$, fix
$P'_a\in A_{\mathsf{u}}$ equal to or touching $P_a$, and let $\Gamma_a$ be a wall-chain of at most $d+1$
$\pi_k$-cubes from $P_a$ to $P'_a$, obtained by changing the at most $d$ differing coordinates one at a
time. Define
\begin{equation*}
\mathrm{Att}^{T}(\tilde{\mathsf{y}}):=\bigcup_a\bigl(\Gamma_a\setminus\{P'_a\}\bigr),
\end{equation*}
which has at most $d$ cubes per component, and
\begin{equation*}
\mathrm{lab}^{T}(A,\tilde{\mathsf{y}}):=A\cup[\tilde{\mathsf{y}}]_k\cup\mathrm{Hub}(\tilde{\mathsf{y}})
\cup\mathrm{Att}^{T}(\tilde{\mathsf{y}}).
\end{equation*}
Put $c_l^{T}:=c_l^{\circ}+d$, which satisfies $c_l^{T}\le c_l$, and
\begin{equation}\label{5.4:eq-touch-label-sum-1}
\lambda^{T}_{\Sigma}:=\kappa_1-2-c_l^{T}(r+1),\qquad
\lambda'^{T}_{\Sigma}:=\lambda^{T}_{\Sigma}-\log(e\hat{b}),\qquad
\epsilon^{T}_{\partial}:=2e^{-\lambda'^{T}_{\Sigma}},
\end{equation}
\begin{equation}\label{5.4:eq-touch-label-sum-2}
\mathsf{M}^{\partial,T}_{r,\beta}[\mathbb{D}]:=\sup_{\Delta_1}e^{-\beta\rho(\Delta_1)}
\sum_{\mathsf{u}\in\mathbb{D}:\,\Delta_1\in R_{\mathsf{u}}}\mu_{\mathsf{u}}\,e^{r d_k(A_{\mathsf{u}})}\,
e^{\epsilon^{T}_{\partial}\#R_{\mathsf{u}}},
\end{equation}
where $\rho(\Delta_1)$ is as in Notation~\ref{5.4:not-zones-cells}.
Then for $r\le a+9$, $1\le\beta\le\lambda'^{T}_{\Sigma}-2$ and $Q_*\in\pi_k$,
\begin{equation}\label{5.4:eq-touch-label-sum-3}
\sum_{\mathsf{u}\in\mathbb{D}}\;
\sum_{\substack{\tilde{\mathsf{y}}\in\mathfrak{Y}^{\mathfrak{g}}_{T}(\mathsf{u}):\\
\mathrm{lab}^{T}\ni Q_*}}
\mu_{\mathsf{u}}\,e^{-(\kappa_1-2)|\tilde{\mathsf{y}}|}\,e^{r d_k(\mathrm{lab}^{T})}
\le S_d\,\mathsf{M}^{\partial,T}_{r+1,\beta}[\mathbb{D}],
\end{equation}
where $A=A_{\mathsf{u}}$ and $\mathrm{lab}^{T}=\mathrm{lab}^{T}(A,\tilde{\mathsf{y}})$. Moreover:
\begin{enumerate}
\item[(i)] $\mathrm{lab}^{T}\in\mathbf{D}_k$;
\item[(ii)] $d_k(\mathrm{lab}^{T})\le d_k(A)+c_l^{T}|\tilde{\mathsf{y}}|$;
\item[(iii)] if $A$ holds a cube not contained in $Z$, then $d_k(\mathrm{lab}^{T})\ge 9\check{R}_k-2$.
\end{enumerate}
\end{theorem}

\begin{proof}
Write $A=A_{\mathsf{u}}$ throughout.

(i) We build $\mathrm{lab}^{T}$ from $A$ by attaching cubes in the following order, component by
component. First attach $\Gamma_a$, running from $P'_a\in A$ outward to $P_a$; consecutive cubes of
$\Gamma_a$ share a wall, so each new cube is wall-adjacent to one attached before. Next attach
$[E_a]_k$, where $E_a\in\tilde{\mathsf{y}}_a$ is the member walled to $\Delta_a$; the cube $[E_a]_k$ is equal
or wall-adjacent to $P_a=[\Delta_a]_k$. If instead $\Delta_a\in\tilde{\mathsf{y}}_a$, nothing is attached at
this stage, since $[\Delta_a]_k=P_a$ is already present. Finally attach the cubes of
$\tilde{\mathsf{y}}_a$ along a tree of $\tilde{\mathsf{y}}_a$ rooted at $E_a$ (respectively at $\Delta_a$):
across a wall-edge of the tree the coarse hull of the next cell is equal or wall-adjacent to the coarse
hull of its parent; across a hub-edge one first attaches the chain of $\mathrm{Hub}(\tilde{\mathsf{y}})$
joining the two cells, and then the coarse hull of the next cell. Starting from $A$, every cube attached
in this order is equal or wall-adjacent to a cube of $A$ or to one attached before, and the order exhausts
$A\cup[\tilde{\mathsf{y}}]_k\cup\mathrm{Hub}(\tilde{\mathsf{y}})\cup\mathrm{Att}^{T}(\tilde{\mathsf{y}})$. Hence
$\mathrm{lab}^{T}\in\mathbf{D}_k$.

(ii) Follow the order of attachment of (i) and apply, at each attached cube, the growth bound for
$d_k$ among the standing geometric axioms (Notation~\ref{5.4:not-zones-cells}). For a component
$\tilde{\mathsf{y}}_a$ with $n_a$ members the increase is at most
\begin{equation*}
d+n_a+n_a d(h_{\sharp}+1)=d+c_l^{\circ}n_a\le(c_l^{\circ}+d)\,n_a ,
\end{equation*}
where the middle equality is the definition of $c_l^{\circ}$, and
the last inequality uses $n_a\ge1$. Summing over the components, whose sizes add up to
$|\tilde{\mathsf{y}}|$, gives $d_k(\mathrm{lab}^{T})\le d_k(A)+c_l^{T}|\tilde{\mathsf{y}}|$.

(iii) If $A$ holds a cube not contained in $Z$, then $\mathrm{lab}^{T}$ holds a cube contained in the
closure of $Z^{c}$. Since $\tilde{\mathsf{y}}$ holds a source cell, $\mathrm{lab}^{T}$ also holds the coarse
hull of that source cell, which lies at distance at least $9\check{R}_k-1$ from $Z^{c}$. The lower bound
of $d_k$ by distances among the standing geometric axioms (Notation~\ref{5.4:not-zones-cells}) then
gives $d_k(\mathrm{lab}^{T})\ge 9\check{R}_k-2$.

(iv) The bound \eqref{5.4:eq-touch-label-sum-3}. We run again the three steps of the proof of the
label-sum theorem that bound the sum there, namely the rooting of the sum at a cell
$\Delta_{\mathfrak{F}}$, the bound of the label factor, and the summation over the components at the rim,
and record what changes. The sum is rooted at the cell $\Delta_{\mathfrak{F}}$ of that
proof, and the sum over the position of the root relative to $Q_*$ produces the factor $S_d$. For the
label factor write $r\,d_k(\mathrm{lab}^{T})=(r+1)\,d_k(\mathrm{lab}^{T})-d_k(\mathrm{lab}^{T})$; in the
second term bound $d_k(\mathrm{lab}^{T})$ from below by $\mathrm{dist}_k(Q_*,\Delta_{\mathfrak{F}})$, as in
that proof, and bound $(r+1)\,d_k(\mathrm{lab}^{T})$ from above by (ii). This gives
\begin{equation}\label{5.4:eq-touch-label-sum-4}
e^{r d_k(\mathrm{lab}^{T})}\le e^{-\mathrm{dist}_k(Q_*,\Delta_{\mathfrak{F}})}\,e^{(r+1)d_k(A)}\,
e^{c_l^{T}(r+1)|\tilde{\mathsf{y}}|}.
\end{equation}
Together with the factor $e^{-(\kappa_1-2)|\tilde{\mathsf{y}}|}$ of the summand, the last factor in
\eqref{5.4:eq-touch-label-sum-4} leaves $e^{-\lambda^{T}_{\Sigma}|\tilde{\mathsf{y}}|}$, by the definition
\eqref{5.4:eq-touch-label-sum-1} of $\lambda^{T}_{\Sigma}$. Thus the attachment cubes, at most $d$ per
component, are charged to $|\tilde{\mathsf{y}}|$ through $c_l^{T}$ rather than to the label $A$; this is the
only change in the rate, which is $\lambda^{T}_{\Sigma}$ in place of $\lambda_{\Sigma}$, and
correspondingly $\lambda'^{T}_{\Sigma}$ and $\epsilon^{T}_{\partial}$ in place of $\lambda'_{\Sigma}$ and
$\epsilon_{\partial}$.

The components of $\tilde{\mathsf{y}}$ other than the trunk are summed at the members of $R_{\mathsf{u}}$
they hold, and their total contribution is at most $e^{\epsilon^{T}_{\partial}\#R_{\mathsf{u}}}$, by the
third of the properties of $\rho$ among the standing geometric axioms
(Notation~\ref{5.4:not-zones-cells}), applied with $\lambda'^{T}_{\Sigma}$ and $\epsilon^{T}_{\partial}$ in
place of $\lambda'_{\Sigma}$ and $\epsilon_{\partial}$.

The trunk $\tilde{\mathsf{y}}_1$, the component holding the first source cell, satisfies
$|\tilde{\mathsf{y}}_1|\ge\rho(\Delta_1)$ for its member $\Delta_1\in R_{\mathsf{u}}$. Hence
$e^{\beta\rho(\Delta_1)}\le e^{\beta|\tilde{\mathsf{y}}_1|}$, and the factor $e^{\beta\rho(\Delta_1)}$ needed
to pass to the supremum \eqref{5.4:eq-touch-label-sum-2} is compensated in the sum over the trunk: since
$\lambda'^{T}_{\Sigma}-\beta\ge2$,
\begin{equation*}
\sum_{n\ge1}n\,e^{-(\lambda'^{T}_{\Sigma}-\beta)n}\le\sum_{n\ge1}n\,e^{-2n}
=\frac{e^{-2}}{(1-e^{-2})^{2}}<1 .
\end{equation*}
What remains, for fixed $\Delta_1$, is at most
\begin{equation*}
e^{-\beta\rho(\Delta_1)}\sum_{\mathsf{u}\in\mathbb{D}:\,\Delta_1\in R_{\mathsf{u}}}\mu_{\mathsf{u}}\,
e^{(r+1)d_k(A_{\mathsf{u}})}\,e^{\epsilon^{T}_{\partial}\#R_{\mathsf{u}}},
\end{equation*}
which is bounded by
$\mathsf{M}^{\partial,T}_{r+1,\beta}[\mathbb{D}]$; with the factor $S_d$ from the root this is
\eqref{5.4:eq-touch-label-sum-3}. At no point does the argument use the inclusion
$[\mathrm{St}_{\mathsf{u}}]_k\subset A_{\mathsf{u}}$; the anchoring condition enters only through the choice of
$P'_a$ in the construction of $\mathrm{Att}^{T}$.

(v) The range of the parameters. Let $\lambda_*$ denote the constant entering the standing lower bound on
$\kappa_1$; that bound gives $\kappa_1-2\ge c_l(2\kappa+4)-1+\lambda_*$. Since $r\le a+9$ and
$c_l^{T}\le c_l$ give $c_l^{T}(r+1)\le c_l(a+10)$,
\begin{equation*}
\lambda'^{T}_{\Sigma}\ge c_l(2\kappa+4)-c_l(a+10)-1+\lambda_*-\log(e\hat{b}).
\end{equation*}
By the standing condition $\lambda_*\ge\log(e\hat{b})$ and $a=\tfrac{31}{20}\kappa$, the right side is at
least
\begin{equation*}
c_l\bigl(2\kappa+4-\tfrac{31}{20}\kappa-10\bigr)-1=c_l\bigl(\tfrac{9}{20}\kappa-6\bigr)-1\ge 75c_l-1
\qquad\text{for }\kappa\ge180 .
\end{equation*}
The conditions of the label-sum theorem at general grain and with the parameter $\beta$ read
$c_l^{\circ}$ only through
$c_l^{\circ}\le c_l$, and $c_l^{T}\le c_l$ as well; so the same lower bound on $\kappa_1$ and the same
condition \eqref{5.4:eq-zones-cells-i} serve here unchanged, and $\beta=3$ is available. No lower bound
on $\kappa_1$ beyond those already imposed is needed.
\end{proof}

\begin{remark}[The touch reading carried through the frame]
\label{5.4:rem-touch-propagation}
Read the anchoring of stars by touching, as in the label sum with attachments under the touch
reading (Theorem~\ref{5.4:thm-touch-label-sum}). Then the following statements of the
summation of boundary terms at finer zones hold, under the same hypotheses, in the modified
forms below.
\begin{enumerate}
\item[(i)] In the counting lemma of that summation, the counts of the star
$\mathrm{St}^{j,T}_v$ of the unit $v$ and of its rim set $R^{+}_v$ become
\begin{equation}\label{5.4:eq-touch-propagation-1}
\begin{gathered}
\#\mathrm{St}^{j,T}_v\le 3^d N_j(X_v),\qquad
\#R^{+}_v\le 3^d(2\hat{b}+1)N_j,\\
\epsilon^{T}_{\partial}\,\#R^{+}_v\le d_j+1 ,
\end{gathered}
\end{equation}
where $N_j$ and $d_j$ abbreviate $N_j(X_v)$ and $d_j(X_v)$.
\item[(ii)] On the rim, the decay exponent $\rho_j$ satisfies $\rho_j\ge 9\check{R}_k-3$.
\item[(iii)] Clauses (c), (d) and (e) of the bundle-indexing statement hold with bundles indexed
by families $\tilde{\mathsf{y}}\subset\mathfrak{g}^{\mathrm{cut}}_v$ of cells of the cut grain
$\mathfrak{g}^{\mathrm{cut}}_v$ of $v$; these families host the pieces that cut
partially exempt members, and for a set $G$ one puts $T_G:=G\cap T_v$.
\item[(iv)] The proximity statement holds in the form ``within $2$ cells''.
\item[(v)] Clause~(c) of the main theorem of that summation and its clause for long units hold
with the factor $e^{-\frac12(9\check{R}_k-3)}$, the constants of their original statements being
multiplied by $\hat{b}_t C_{\mathsf{b}}2^{p_{\mathsf{b}}}$.
\item[(vi)] The corollary producing the small factor $\theta'(g_k)$ holds with an additional
factor $e$.
\item[(vii)] The clause of the parent-location corollary that locates the parent holds under the
touch reading in the following form: the parent lies in $\mathrm{lab}(\mathfrak{b})$ or in a
$\pi_k$-cube touching a cube of $K_v$, so that its distance is at most $1\le c_{\rho}\rho$.
\item[(viii)] For the long units of class (d), the bundle sum for long units off the treated class
(Theorem~\ref{5.4:thm-bundle-sum}) holds under the touch reading, the constant
$\mathsf{C}^{\mathrm{lg},d}$ of \eqref{5.4:eq-long-unit-sums-a} that it feeds being replaced by
\begin{equation}\label{5.4:eq-touch-propagation-2}
\mathsf{C}^{\mathrm{lg},d}_{T}:=e\,\hat{b}_t C_{\mathsf{b}}2^{p_{\mathsf{b}}}\,
\mathsf{C}^{\mathrm{lg},d}.
\end{equation}
\end{enumerate}
\end{remark}

\begin{proof}
For~(i), the third inequality of \eqref{5.4:eq-touch-propagation-1} holds by a comparison, with
constants depending on $d$ only, against the margin in the smallness condition on
$\epsilon_{\partial}$.

For~(v), the constants of the original clauses are multiplied by the factor
$\hat{b}_t C_{\mathsf{b}}2^{p_{\mathsf{b}}}$ of Remark~\ref{5.4:rem-per-cube-remarks}, and the
exponential factor is the one furnished by~(ii). For~(vi), note that
\begin{equation*}
e^{-\frac12(9\check{R}_k-3)}=e^{1/2}\,e^{-\frac12(9\check{R}_k-2)}\le e\,
e^{-\frac12(9\check{R}_k-2)},
\end{equation*}
so the exponential factor of~(v) exceeds the one used in the long units in the boundary-term
sums (Corollary~\ref{5.4:cor-long-unit-sums}) by at most the factor $e^{1/2}\le e$; the corollary
producing $\theta'(g_k)$ holds with an additional factor $e$, which in particular covers this
difference.

For~(vii), under the touch reading the parent of the unit lies in the label
$\mathrm{lab}(\mathfrak{b})$ of the bundle $\mathfrak{b}$, or in a $\pi_k$-cube that touches a
cube of the anchoring set $K_v$. In both cases the distance controlled by the parent-location
corollary is at most $1$, and $1\le c_{\rho}\rho$, where $c_{\rho}\rho$ is the parent-distance
allowance of that corollary; this is the bound that clause requires.

For~(viii), apply Theorem~\ref{5.4:thm-bundle-sum} under the touch reading. The window is entered at
$2\varrho-2$. The smallness available on the rim is $9\check{R}_k-3$ by~(ii), and the loss
relative to the original reading is absorbed by the spare factor $e^{-8\check{R}_k}$. The
constant $\mathsf{C}^{\mathrm{lg},d}$ is multiplied by the factor
$\hat{b}_t C_{\mathsf{b}}2^{p_{\mathsf{b}}}$, as in~(v),
and by a further factor $e$, which gives \eqref{5.4:eq-touch-propagation-2}.

One might instead enlarge the anchoring set of the unit $v$ to the coarse hull
$A:=\bigl[\bigcup_{\Delta\in\mathrm{St}^{T}_v}\Delta\bigr]_k$ of its star under the touch reading.
That choice is not used here. Indeed, $d_k(A)$ can be as large as
\begin{equation*}
\bigl(1+(3^d-1)\mathsf{c}_d\bigr)\,d_k(K_v),
\end{equation*}
and the resulting rate exceeds the rate $\frac85\kappa$ available for the long units. In the
argument above the attachment is charged to $|\tilde{\mathsf{y}}|$ and never to $d_k(K_v)$.

The estimates of the summation of boundary terms at finer zones therefore close, under the same
hypotheses, under both the original reading and the touch reading; under the touch
reading that summation is read throughout with the attachment cubes $\mathrm{Att}^{T}$ of
Theorem~\ref{5.4:thm-touch-label-sum}.
\end{proof}

\begin{proposition}[Long units anchored at their own scale]\label{5.4:prop-anchor-route}
Assume the following three statements.
\begin{enumerate}
\item[(a)] The base bound \eqref{5.4:eq-age-base-bounds-b} for long units, with the age function
$\mathsf{b}^{\mathrm{lg}}$ and its growth constants $C_{\mathsf{b}}$, $p_{\mathsf{b}}$ of
Definition~\ref{5.4:def-age-base-bounds}, and with its decay constant $c_{\mathrm{sp}}>0$.
\item[(b)] The own-scale anchoring property of the representation and bound of the large-field
terms, in the current label. There is $\delta_{\mathrm{anc}}\in\{0,1\}$ with the
following property. Every long unit $v\in\mathfrak{L}^{\mathrm{lg}}$, single or tail, that is live
at the current level has a cube $Q_1(v)\in X_v$ which, in the current label, meets the anchoring
set $Z^{\boxplus}_{j_v}$ of Lemma~\ref{5.2:lem-current-label-anchors} within one layer of
$\pi_{j_v}$-cubes, the width of this layer being
$\delta_{\mathrm{anc}}$, the anchoring margin. For a tail $v$ the cube $Q_1(v)$ is required to be a cube of the
tail's own carrier $X_v$. (Compare the condition written $X\cap Z^{\sim}_j\neq\emptyset$, in the
notation of that paper, in \cite[(2.41)]{B-CMP119}.)
\item[(c)] The cell dichotomy of the summation of boundary terms at finer zones, in the following form.
Let $\Delta$ be a cell of level $n$ such that $\Delta\not\subset Z$ or $\Delta$ lies within
distance $2\check{R}_k$ of $\partial Z$, and let $j<n$. Then either
\begin{equation*}
\text{(i)}\quad \operatorname{dist}_n\bigl(Z^{\boxplus}_j,\Delta\bigr)\ge 1,
\end{equation*}
or the alternative (ii) of that dichotomy holds. Of alternative (ii) only the following two
properties are used: it is reached only when $\Delta$ is raised, in the sense of the standing
geometric axioms \eqref{5.4:eq-zones-cells-g}, so that it never holds at an un-raised cell; and at
a raised cell of level $n$ with raiser $\mathfrak{a}''$ it holds only for levels $j$ with
$k_{0a''}\le j\le n-1$, where $k_{0a''}$ is the lower end of the range of levels which that
dichotomy attaches to the raiser $\mathfrak{a}''$, distinct from the level $k_{a''}$ of the
arrest $\mathfrak{a}''$.
\end{enumerate}
The star-exit condition \eqref{5.4:eq-pointed-units-a} is not assumed. Put $q:=p_{\mathsf{b}}+d$,
a letter local to this proposition, and
\begin{equation}\label{5.4:eq-anchor-route-1}
\begin{split}
C_{\mathrm{anc}}&:=C_{\mathsf{b}}\sum_{u\ge 0}(1+u)^{p_{\mathsf{b}}}L^{du}
\exp\bigl(-c_{\mathrm{sp}}(L^u-1-\delta_{\mathrm{anc}})\,\mathbf{1}[u\ge 1]\bigr)\\
&\le C_{\mathsf{b}}\Bigl[1+2^{p_{\mathsf{b}}}\Bigl(\frac{6q}{e\,c_{\mathrm{sp}}}\Bigr)^{q}
\bigl(e^{c_{\mathrm{sp}}/2}-1\bigr)^{-1}\Bigr].
\end{split}
\end{equation}
Then the following hold.
\begin{enumerate}
\item[(1)] Let $v$ be a long unit, $j:=j_v$, and let $\Delta$ be a cell of level $n>j$ that meets
$X_v$, and suppose that $\operatorname{dist}_n(Z^{\boxplus}_j,\Delta)\ge 1$, i.e.\ alternative
(i) of (c) holds at $(\Delta,j)$. Then $d_v\ge L^{n-j}-1-\delta_{\mathrm{anc}}$.
\item[(2)] Let $r'\le\frac{8}{5}\kappa$. At every un-raised cell $\Delta_2$ of level $\ell$ such
that either $\Delta_2\not\subset Z$ or $\Delta_2$ lies within distance $2\check{R}_k$ of
$\partial Z$ (hence in particular at the charge's cell $\Delta^{\partial}$, which lies within
distance $2\check{R}_k$ of $\partial Z$ and is un-raised), the met-cell sum obeys
\begin{equation}\label{5.4:eq-anchor-route-2}
S_{r'}(\Delta_2;\mathfrak{L}^{\mathrm{lg}})\le C_{\mathrm{anc}}\,\mathsf{b}^{\mathrm{lg}}(k-\ell).
\end{equation}
\item[(3)] Let $r'$ be as in (2). At a raised cell $\Delta_2\not\subset Z$ of level $n$, with
raiser $\mathfrak{a}''$ foreign or sibling,
\begin{equation}\label{5.4:eq-anchor-route-3}
S_{r'}(\Delta_2;\mathfrak{L}^{\mathrm{lg}})\le
\bigl(C_{\mathrm{anc}}+\mathfrak{B}^{+}_{\mathrm{win}}(\Delta_2)\bigr)\,\mathsf{b}^{\mathrm{lg}}(k-n),
\end{equation}
where $\mathfrak{B}^{+}_{\mathrm{win}}(\Delta_2):=B^{\mathrm{win}}(\check{R}_{k_{a''}})$ at raised
cells and $\mathfrak{B}^{+}_{\mathrm{win}}(\Delta_2):=0$ at all other cells. With this definition,
\eqref{5.4:eq-anchor-route-2} and \eqref{5.4:eq-anchor-route-3} together give, at every cell
$\Delta_2$ covered by (2) or by (3), of level $\ell$,
\begin{equation*}
S_{r'}(\Delta_2;\mathfrak{L}^{\mathrm{lg}})\le
\bigl(C_{\mathrm{anc}}+\mathfrak{B}^{+}_{\mathrm{win}}(\Delta_2)\bigr)\,\mathsf{b}^{\mathrm{lg}}(k-\ell).
\end{equation*}
If $\Delta_2$ is as in (3) and $\Delta_1\in\pi^{\mathfrak{s}}$ shares a wall with $\Delta_2$, then
$\rho(\Delta_1)\ge 2\check{R}_{k_{a''}}-1$ and
\begin{equation}\label{5.4:eq-anchor-route-4}
\begin{split}
C_{\mathrm{anc}}+B^{\mathrm{win}}(\check{R}_{k_{a''}})\,e^{-\rho(\Delta_1)}
&\le C_{\mathrm{anc}}+B^{\mathrm{win}}(\check{R}_{k_{a''}})\,e^{-(2\check{R}_{k_{a''}}-1)}\\
&\le C_{\mathrm{anc}}+C^{\mathrm{win}}_{B'} .
\end{split}
\end{equation}
\end{enumerate}
The value of $c_{\mathrm{sp}}>0$ is arbitrary, and the point $y_v$ of a unit is never used.
\end{proposition}

\begin{proof}
Part (1). Since $\Delta$ meets $X_v$, Lemma~\ref{5.4:lem-cells-carrier}~(b) gives a cube
$Q_2\in X_v$ with $Q_2\subset\Delta$. By hypothesis (b), the cube $Q_1:=Q_1(v)\in X_v$ meets
$Z^{\boxplus}_j$
within one layer of $\pi_j$-cubes of width $\delta_{\mathrm{anc}}$. Alternative (i) at
$(\Delta,j)$ says that $Z^{\boxplus}_j$ and $\Delta$ are at distance at least one in the units of
level $n$, that is, at distance at least $L^{n-j}$ in the units of level $j$. Subtracting the
side of the cube $Q_1$ and the anchoring margin $\delta_{\mathrm{anc}}$, every $x\in Q_2$ and every
$x'\in Q_1$ satisfy, in the units of level $j$,
\begin{equation*}
|x-x'|\ge L^{n-j}-(1+\delta_{\mathrm{anc}}).
\end{equation*}
The geometric axiom \eqref{5.4:eq-zones-cells-c}, applied to the two cubes $Q_1$, $Q_2$ of $X_v$,
turns this separation into $d_v\ge L^{n-j}-1-\delta_{\mathrm{anc}}$.

Part (2). Let $\Delta_2$ be as in (2). Because $\Delta_2$ is un-raised in the sense of
\eqref{5.4:eq-zones-cells-g}, by hypothesis (c) alternative
(ii) is void there: alternative (i) holds at $(\Delta_2,j)$ for every $j<\ell$. Let $v$ be a unit
counted in $S_{r'}(\Delta_2;\mathfrak{L}^{\mathrm{lg}})$, so that $X_v$ meets $\Delta_2$. By
Lemma~\ref{5.4:lem-cells-carrier}~(e), $j_v\le\ell$; put $u:=\ell-j_v\ge 0$. We group the counted
units by the value of $u$ and, for fixed $u$, by the cube of level $j_v$ inside $\Delta_2$ through
which $X_v$ meets $\Delta_2$; by Lemma~\ref{5.4:lem-cells-carrier}~(e) there are $L^{du}$ such
cubes. If $u\ge 1$, then $\ell>j_v$ and part (1) with $\Delta=\Delta_2$, $n=\ell$ gives
\begin{equation*}
d_v\ge L^u-1-\delta_{\mathrm{anc}} ;
\end{equation*}
if $u=0$ we use only $d_v\ge 0$. The base bound (a), applied with the index $r'$ at level $j_v$
with this lower bound on $d_v$, bounds the contribution of one group by
\begin{equation*}
C_{\mathsf{b}}(1+u)^{p_{\mathsf{b}}}
\exp\bigl(-c_{\mathrm{sp}}(L^u-1-\delta_{\mathrm{anc}})\mathbf{1}[u\ge 1]\bigr)\,
\mathsf{b}^{\mathrm{lg}}(k-\ell).
\end{equation*}
Multiplying by the $L^{du}$ cubes and summing over $u\ge 0$ gives \eqref{5.4:eq-anchor-route-2}
with the constant defined in \eqref{5.4:eq-anchor-route-1}. No property of $c_{\mathrm{sp}}$ beyond
positivity, and no property of $y_v$, has been used.

It remains to prove the inequality in \eqref{5.4:eq-anchor-route-1}. The term $u=0$ equals
$C_{\mathsf{b}}$. Let $u\ge 1$ and $y:=L^u$. Since $L$ is an odd integer larger than one, $y\ge 3$,
so $2\le 2y/3$ and $\delta_{\mathrm{anc}}\le 1$ give
\begin{equation*}
L^u-1-\delta_{\mathrm{anc}}\ge y-2\ge y/3 .
\end{equation*}
Moreover $1+u\le 2^u\le 2L^u$, so $(1+u)^{p_{\mathsf{b}}}L^{du}\le 2^{p_{\mathsf{b}}}y^{q}$ with
$q=p_{\mathsf{b}}+d$. Hence the term with index $u$ is at most
\begin{align*}
C_{\mathsf{b}}2^{p_{\mathsf{b}}}y^q e^{-c_{\mathrm{sp}}y/3}
&=C_{\mathsf{b}}2^{p_{\mathsf{b}}}\bigl(y^q e^{-c_{\mathrm{sp}}y/6}\bigr)e^{-c_{\mathrm{sp}}y/6}\\
&\le C_{\mathsf{b}}2^{p_{\mathsf{b}}}\Bigl(\frac{6q}{e\,c_{\mathrm{sp}}}\Bigr)^q
e^{-c_{\mathrm{sp}}L^u/6},
\end{align*}
because the function $y\mapsto y^qe^{-ay}$ on $(0,\infty)$, $a>0$, has maximum $(q/(ea))^q$,
attained at $y=q/a$; here $a=c_{\mathrm{sp}}/6$. Finally $L^u\ge 3^u\ge 3u$, so
\begin{equation*}
\sum_{u\ge 1}e^{-c_{\mathrm{sp}}L^u/6}\le\sum_{u\ge 1}e^{-c_{\mathrm{sp}}u/2}
=\bigl(e^{c_{\mathrm{sp}}/2}-1\bigr)^{-1},
\end{equation*}
and adding the term $u=0$ yields the right side of \eqref{5.4:eq-anchor-route-1}.

Part (3). Let $\Delta_2$ be as in (3): raised, $\Delta_2\not\subset Z$, of level $n$, with raiser
$\mathfrak{a}''$ foreign or sibling, and put
\begin{equation*}
J_{\mathrm{win}}:=\{\,j<n:\ \text{alternative (ii) of hypothesis (c) holds at }(\Delta_2,j)\,\}.
\end{equation*}
By hypothesis (c), $J_{\mathrm{win}}\subset[k_{0a''},\,n-1]$, and for these levels the radius bound in
\eqref{5.4:eq-zones-cells-g} gives $L^{n-j}\le L\,C_{\uparrow}\check{R}_{k_{a''}}$. Grouping the
units with $j_v\in J_{\mathrm{win}}$ as in part (2), but using only $d_v\ge 0$, the base bound (a)
bounds their contribution by
\begin{equation*}
\sum_{j\in J_{\mathrm{win}}}L^{d(n-j)}C_{\mathsf{b}}(1+n-j)^{p_{\mathsf{b}}}\,
\mathsf{b}^{\mathrm{lg}}(k-n)\le B^{\mathrm{win}}(\check{R}_{k_{a''}})\,\mathsf{b}^{\mathrm{lg}}(k-n),
\end{equation*}
where the last inequality is the definition of $B^{\mathrm{win}}$ in
\eqref{5.4:eq-transport-constants-d}, together with $\sum_{0\le u\le U}L^{du}\le 2L^{dU}$ for the
largest $U$ with $L^U\le L\,C_{\uparrow}\check{R}_{k_{a''}}$. By
Lemma~\ref{5.4:lem-cells-carrier}~(e), every counted unit has $j_v\le n$. For the remaining levels
$j\le n$ with $j\notin J_{\mathrm{win}}$ (for $j<n$ alternative (i) holds at $(\Delta_2,j)$ by (c)
and part (1) applies; for $j=n$ only $d_v\ge 0$ is used), the computation of part (2) bounds
their contribution by $C_{\mathrm{anc}}\mathsf{b}^{\mathrm{lg}}(k-n)$. Adding
the two contributions gives \eqref{5.4:eq-anchor-route-3}. At an un-raised cell of (2) one has
$\mathfrak{B}^{+}_{\mathrm{win}}(\Delta_2)=0$ and the combined bound is
\eqref{5.4:eq-anchor-route-2}; at a raised cell of (3) it is \eqref{5.4:eq-anchor-route-3}.

The collar lemma (Lemma~\ref{5.4:lem-collar}), applied to the raiser $\mathfrak{a}''$ and to each
cell $\Delta_1\in\pi^{\mathfrak{s}}$ sharing a wall with $\Delta_2$, gives
$\rho(\Delta_1)\ge 2\check{R}_{k_{a''}}-1$. Against one factor $e^{-\rho(\Delta_1)}$ of the
consuming estimate, the window term is therefore controlled as follows. Since $x\mapsto e^{-x}$ is
decreasing, the lower bound on $\rho(\Delta_1)$ gives the
first inequality in \eqref{5.4:eq-anchor-route-4}, and the second is the bound
$B^{\mathrm{win}}(\check{R}_{k_{a''}})e^{-(2\check{R}_{k_{a''}}-1)}\le C^{\mathrm{win}}_{B'}$ of
\eqref{5.4:eq-collar-a}. This proves \eqref{5.4:eq-anchor-route-4}.
\end{proof}

Consequently, the estimates that use the transport constant $\mathfrak{K}$ of
Notation~\ref{5.4:not-transport-constants} hold under the hypotheses of this proposition with
$C_{\mathrm{anc}}$ in place of $\mathfrak{K}$, with $C^{\mathrm{win}}_{B'}$ added in the estimate for the
long units of class (d).

Hypothesis (b) cannot simply be omitted. By Remark~\ref{5.4:rem-per-cube-remarks}, if neither the
star-exit condition \eqref{5.4:eq-pointed-units-a} nor an anchor is available, then units each
consisting of one $\pi_m$-cube inside $\Delta_2$, one for every level $m<\ell$, with $d_v=0$,
satisfy the base bound \eqref{5.4:eq-age-base-bounds-b} and make $S\ge\ell-1$. An anchor supplies
$d_v\ge L^u-1-\delta_{\mathrm{anc}}$, as in part (1); the star-exit condition supplies
$d_v\ge L^{u-1}-\delta$, with the integer $\delta\ge 0$ fixed in Remark~\ref{5.4:rem-per-cube-remarks}; either
lower bound suffices for the convergence of the sum over $u$. A lower bound of $d_v$ by a constant
independent of $u$ does not suffice.

\begin{remark}[Two organisations of the long-unit sum]\label{5.4:rem-two-organisations}
\leavevmode
\begin{enumerate}
\item[(a)] The sum over the long units of class (d) admits two organisations.

The first is the pointed-piece sum for long units off the treated class. It keeps the
organisation by pieces and uses:
\begin{itemize}
\item the gain lemma of the piece sum;
\item the block lemma, with its constant $C^{0}_{\square}$;
\item the geometric bound in its cell form.
\end{itemize}
It needs the long-unit hypothesis of the pointed-piece sum, which follows from clauses (i)
and (ii) of the condition on long units.

The second is the bundle sum for long units off the treated class
(Theorem~\ref{5.4:thm-bundle-sum}). It keeps the organisation of the summation of boundary
terms at finer zones:
\begin{itemize}
\item the terms are grouped into $\square$-bundles $\mathfrak{b}^{\circ}$;
\item the small factor $\theta'(g_k)$ is given up in exchange for $e^{-(9\check{R}_k-2)}$;
\item the label-sum theorem is applied once, on the grain $\pi^{\mathfrak{s}}$.
\end{itemize}
It needs only the undilated long-unit hypothesis of Theorem~\ref{5.4:thm-bundle-sum}.

The second organisation invokes strictly fewer auxiliary results, and it carries no floor
condition on the window weight and no accompanying cell condition. Its grouping of the terms
into bundles differs from the grouping of the pieces in Ba{\l}aban's construction. Either
organisation yields the constant $\mathsf{C}_1^{\mathrm{ex}}$ as one value.

\item[(b)] Suppose the long column inside $Z$ carries the running-window weight, whose shape is
fixed by $\check{R}_k$ (Remark~\ref{5.4:rem-running-window}). Then, along the route through the
series constants $\mathfrak{a}_d,\mathfrak{e}_d$, the long-column bound for class (c) at a
level-$k$ cell $\Delta^{\partial}\not\subset Z$ touching $Z$ (the cell of
Proposition~\ref{5.4:prop-own-window-weights}) is polynomial in $\check{R}_k$ as written. A
weight of this kind is
among the weights admitted in the bound that fixes the window constant, and it is paid by the
residue $\sigma^{\mathrm{lg}}$ of the corollary on the small factor $\theta'(g_k)$.

Own-window weights, indexed by the level $k_a<k$ of a live own arrest $\mathfrak{a}$, are not of
this kind. They are paid by the unit's own length
(Proposition~\ref{5.4:prop-own-window-weights}), never by the residue $\sigma^{\mathrm{lg}}$.
\end{enumerate}
\end{remark}

\begin{proof}
(a) Both organisations bound one and the same quantity, the sum over the long units of
class (d): the first through the pointed-piece sum for long units off the treated class, the
second through Theorem~\ref{5.4:thm-bundle-sum}. They are two groupings of the same terms and
can be compared input by input.

(b) If the long column inside $Z$ carries the running-window weight, then along the route
through $\mathfrak{a}_d,\mathfrak{e}_d$ the class-(c) long-column bound at $\Delta^{\partial}$
is, as written, a polynomial in $\check{R}_k$; a weight of this kind is among those admitted in
the bound that fixes the window constant, and so it is paid by the residue
$\sigma^{\mathrm{lg}}$.
An own-window weight is indexed by the level $k_a$ of a live own arrest, with $k_a<k$,
rather than by $k$. It is therefore not of the running-window kind, whose shape is fixed by
$\check{R}_k$, and the residue $\sigma^{\mathrm{lg}}$ is not used for it. By
Proposition~\ref{5.4:prop-own-window-weights} it is paid by the length of the unit itself.
\end{proof}

\begin{definition}[Letters, levels and weights for the weight lemma]\label{5.5:not-weight-notation}
The following letters are fixed for the weight lemma and its proof. Where one letter carries
several meanings, the meanings are listed together.

\emph{Stage letters and unit letters.} The letters $\mathsf{n}$ and $\mathsf{j}:=h+\mathsf{n}$
are the stage letters of the block sequences, with collar $R_{\mathsf{n}}$; here $h$ is the
offset of the stage letter $\mathsf{j}$ over $\mathsf{n}$, and not a history. The stage letters
are distinct from the letters $j$, $n$, $\nu$, $t$ of a unit in the sense of
Definition~\ref{5.2:def-sites-units}, which are the unit letters of the correlation theorem
(Theorem~\ref{5.1:thm-correlation-theorem}).

\emph{Letters with one meaning throughout.}
\begin{itemize}
\item $\mathfrak{K}(\square)$ is one and the same count-times-weight sum at every block
$\square$, fixed in the statement of the weight lemma. On the blocks with
$\tilde{\square}^5\cap\mathcal{W}=\emptyset$, where $\mathcal{W}$ is the set of cells defined
below, it is the sum appearing in item (a) of the tiling statement; on the remaining blocks it
is the sum appearing in the lemma that treats those blocks. The two are one sum.
\item $\mathsf{l}^2(\square)$ is defined once, at every block, by
\begin{equation}\label{5.5:eq-weight-notation-1}
\mathsf{l}^2(\square):=\max\bigl\{\mathsf{L}_0^{6s(\square)},\ \kappa_R c_W^3\,
\mathfrak{K}(\square)\bigr\},
\end{equation}
with $\mathsf{L}_0$ the weight base defined below, and $s(\square)$, $\kappa_R$ and $c_W$ as in
the weight lemma. Where $\mathsf{l}^2(\square)$ is
compared with the data factor $\mathfrak{d}(\square)$, the relation is the inequality
\begin{equation*}
\mathsf{l}^2(\square)\ge \mathsf{L}_0^{6s(\square)}\,\mathfrak{d}(\square),
\end{equation*}
not a definition; wherever the product $\mathsf{L}_0^{6s(\square)}\mathfrak{d}(\square)$ is
meant, it is written out.
\item $\mathfrak{r}_v$ is the near unit's own box at the site (P).
\item The moves of the production closure $\mathcal{K}^{rd}$ are
$\mathsf{M}1,\dots,\mathsf{M}8$, with $\mathsf{M}8:=\mathsf{M}^{\perp}$. The sans-serif letter
$\mathsf{M}$ denotes moves only; the italic $M$ is the auxiliary parameter $M=L^{m'}$.
\end{itemize}

\emph{The letter $t$.}
\begin{itemize}
\item $t\in[0,1]$ is the interpolation parameter of the source interpolation.
\item $t:=L^{j-n}$ is the scale ratio of a near unit, and $t_p:=L^{j-p}$.
\item $t(\Delta)$ is the step of the arrest that raised the cell $\Delta$.
\item $t_{\mathfrak{c}}$ is the step of the preparation $\mathfrak{c}$. Where $t:=t_{\mathfrak{c}}$
is written, $t$ is that step and never the scale ratio.
\end{itemize}
The sans-serif $\mathsf{t}$ is the complex source parameter, $\mathsf{t}\in\mathbb{D}_{\square}$.

\emph{The letter $A$.} In the displays of this paragraph the letters other than $A$,
$\mathbf{A}$, $t$ and $\mathsf{t}$ carry the meanings fixed where those displays occur.
\begin{itemize}
\item $\mathbf{A}:=(t\zeta'_{\square}+\mathsf{t}\zeta_{\square})\mathbb{H}^w$ is the fluctuation
carried by the source; here $t\in[0,1]$ is the interpolation parameter.
\item $A$ is the potential
\begin{equation*}
A:=\frac{1}{i\eta}\log\bigl(e^{i\eta\operatorname{Ad}_g\mathbf{A}}\,e^{i\eta H}\bigr),
\end{equation*}
and this is also the letter $A$ of the display $\mathfrak{b}=e^{i\eta A}$ of the statement in
which this potential is introduced.
\item $A(h_y,\cdot)$ and $A(c_v,\cdot)$ are the Wilson and counterterm plaquette functionals of
the class of Wilson-type pieces; $\hat{A}$ is their slot letter.
\item $A:=\check{\mathsf{T}}_k+1+c_{k-t}$ is a number, used only in the payment for depth.
\item $\mathsf{A}_w(\gamma)$ is the ceiling function.
\end{itemize}

\emph{The letter $\mathcal{K}$.} In the lemma at which the source's fluctuation $\mathbf{A}$ is
substituted into a slot written $\mathcal{K}$, the assignment $\mathcal{K}:=\mathbf{A}$ fixes that
slot letter, and it has this meaning there only. Everywhere else $\mathcal{K}:=\mathcal{K}^{rd}$,
the pair system.

\emph{The letter $w$.} The letters $w(x)$, $w(\Delta)$, $w_Q$, $w^{\vee}$, $w_{\mathfrak{a}}$
and $w^+$ are weights. The bare letter $w$ in
\begin{equation*}
\varrho(\square):=1+d^{\wedge}(\square,R_{\mathsf{n}})/w,
\end{equation*}
and in the width condition that accompanies the definition of $w$, is $w:=M/M_1$. This is the
side of the own-scale cube measured in units of $M_1$; it is never a weight. The index $w$ of
$\mathbb{H}^w$ is the complex source $w=\sum_i z_i f_i$ of the correlation theorem.

\emph{Levels and weights.}
\begin{itemize}
\item $\mathfrak{B}:=\mathbb{B}^{(\mathsf{n})}_k$ and
$\mathfrak{B}^+:=\mathbb{B}^{(\mathsf{n}+1)}_k$; both are frozen on live pieces.
\item $\ell(x)$ is the current level of $x$.
\item $w(x)$ is the actual weight at $x$, as fixed in the definition of the weights of the
large-field operations. On a cell raised during the present operation
$\mathbf{R}_k$, $w(x)\le\mathsf{L}_0^{2\mathsf{r}_0}$, where $\mathsf{r}_0$ is the largest raise
of this $\mathbf{R}_k$. On the region $W_{\mathfrak{a}}$ of a live arrest $\mathfrak{a}$, $w(x)$
is the frozen weight $w_{\mathfrak{a}}$, by the freezing clause of the arrest construction.
Elsewhere $w(x)=1$.
\item For a cell $\Delta$, the weight clause of the arrest construction gives
\begin{equation}\label{5.5:eq-weight-notation-2}
w(\Delta)\le\mathsf{L}_0^{2\mathsf{r}(\Delta)},\qquad
\mathsf{r}(\Delta):=\operatorname{lev}(\Delta)-b(\Delta),
\end{equation}
where $\operatorname{lev}(\Delta)$ and $b(\Delta)$ are the levels of $\Delta$ after and before
its raise in the arrest construction, so that $\mathsf{r}(\Delta)$ is the raise of $\Delta$.
\item $\mathcal{W}$ is the set of inherited raised cells of live arrests of earlier steps. The
data factor $\mathfrak{d}$ lives on this set. No other cells enter $\mathcal{W}$: after a
completed operation $\mathbf{R}_m$, $m$ an earlier step (in this item and the next, $m$ is a
running index, not the torus exponent), the partition $\mathbb{B}_m$ on $Z$ is $\{Z^{(m)}\}$
\cite[(1.33)]{B-CMP122b}. Hence no stratum of cells other than $Z^{(m)}$ remains on $Z$.
\item For a level $m$, $\ell_m:=L^m\eta$, and
\begin{equation}\label{5.5:eq-weight-notation-3}
r_2(m):=\tilde{\alpha}_{0,m}\,\eta^2\ell_m^{-2}=\tilde{\alpha}_{0,m}L^{-2m};
\end{equation}
this is the function $r_2$ of the arrest construction, with its constant
$\tilde{\alpha}_{0,m}$. The constant $\tilde{\alpha}_{0,m}$ satisfies
\begin{equation}\label{5.5:eq-weight-notation-4}
\tilde{\alpha}_{0,m}\le 2\alpha_{0,m},
\end{equation}
and this is the only relation between $\tilde{\alpha}_{0,m}$ and Ba{\l}aban's $\alpha_{0,m}$
used below.
\end{itemize}

\emph{The weight base.} The weight base $\mathsf{L}_0$, also written $L_{\circ}:=\mathsf{L}_0$,
is printed in sans-serif. It is not the constant $L_0=L_0(N)$. It satisfies the one-sided
condition $\mathsf{L}_0^2\le L/3$ of \cite{B-CMP122a}. Finally, the function $\theta_2$ on the
points $x$ above, and the form of this condition that is used in dimension $d=4$, are
\begin{equation}\label{5.5:eq-weight-notation-a}
\theta_2(x):=w(x)\,r_2(\ell(x)),\qquad
\mathsf{L}_0^6\le \frac{L^3}{27}.
\end{equation}
\end{definition}

\begin{proof}
The second expression for $r_2(m)$ in \eqref{5.5:eq-weight-notation-3} is obtained by
substituting $\ell_m=L^m\eta$: then $\eta^2\ell_m^{-2}=\eta^2L^{-2m}\eta^{-2}=L^{-2m}$.

For the second formula in \eqref{5.5:eq-weight-notation-a}, start from
$\mathsf{L}_0^2\le L/3$. Both sides are
non-negative, and $u\mapsto u^3$ is increasing on $[0,\infty)$. Cubing therefore preserves the
inequality and gives
\begin{equation*}
\mathsf{L}_0^6=\bigl(\mathsf{L}_0^2\bigr)^3\le\Bigl(\frac{L}{3}\Bigr)^3=\frac{L^3}{27}.
\end{equation*}
This is the second formula in \eqref{5.5:eq-weight-notation-a}.
\end{proof}

\begin{definition}[Windowed source, box pairs and weighted scaffold majorants]
\label{5.5:def-windowed-source}
The stage letters $(\mathsf{n},\mathsf{j})$ with the collar $R_{\mathsf{n}}$, the unit letters
$(j,n,\nu,t)$, the block sequences $\mathfrak{B}$ and $\mathfrak{B}^+$, the current level
$\ell(\cdot)$, the weights $w$ and the family $\mathcal{W}$ are those of
Notation~\ref{5.5:not-weight-notation}. The letter $t$ has two meanings below, both listed there:
in \textup{(a)} and \textup{(b)} it is the interpolation parameter $t\in[0,1]$, and in
\textup{(d)} and \textup{(e)} it is the scale ratio of a near unit.
The words zone, own unit, core, short single, comb gauge and graded rim used below carry the
meanings fixed with the table of units and with the scaffold bounds of near units.

\smallskip\noindent\emph{(a) The source and the window.}
Let $(\mathbb{U},\mathbb{J})$ be a configuration pair of the source of stage $\mathsf{n}+1$ over
the region $\hat{X}$ of that source, read with the block sequences frozen on the live pieces of
arrests, and assume $(\mathbb{U},\mathbb{J})\in\mathfrak{K}_P$, the admissible source class. Let
$b$ be a datum on the set $\mathcal{C}_{\mathsf{j}}$ of stage $\mathsf{j}$, and let $\sigma$ be a
function of the cells $\Delta$ with
$|\sigma(\Delta)|\le e^{\kappa_1}$. Let $\mathbb{H}$ be the fluctuation field given by the
theorem on the fluctuation field at the data $(\mathfrak{B};\mathbb{U},\mathbb{J};\sigma;(0,b))$.
Let $\square$ be a block, that is, a union of $2^d$ cubes of the partition $\pi_{n_0}$, at
zone $n_0$, and put
\begin{equation*}
\square_0 := \tilde{\square}^5 .
\end{equation*}
The \emph{windowed minimiser} $\mathbb{W}$ is
$U(\mathfrak{B}^+(\square_0),M^{\cdot}(\mathbb{U}))$ on $\square_0$, extended by $\mathbb{U}$
outside $\square_0$, as in \cite{B-CMP109}; here $M^{\cdot}$ denotes Ba{\l}aban's minimiser map
and $U(\cdot,\cdot)$ the configuration that his construction assigns to a block sequence and a
configuration. Let $\mathbb{H}^w$ be the field $\mathbb{H}$ at the
data $(\mathfrak{B};\mathbb{W},\mathbb{J})$, the remaining data being unchanged. For the
interpolation parameter $t\in[0,1]$ and the complex source parameter $\mathsf{t}$ put
\begin{equation*}
\mathbf{A} := \bigl(t\zeta'_{\square}+\mathsf{t}\,\zeta_{\square}\bigr)\mathbb{H}^w ,
\end{equation*}
where $\zeta_{\square}$ and $\zeta'_{\square}$ are the cut-off functions of the block $\square$,
built from the cut-off function $\zeta$ of \cite{B-CMP109}, which vanishes wherever its argument
has absolute value at least $\frac23$.
The normaliser of the block is the number $\vartheta(\square)$ defined in
\eqref{5.5:eq-windowed-source-a} below, in which $\delta_P$ is the rate of decay of the
normaliser in the distance to the collar and
$d^{\wedge}(\square,R_{\mathsf{n}})$ is the distance from $\square$ to the collar; the normaliser
$\vartheta(\square)$ therefore decays as $\square$ moves away from the collar $R_{\mathsf{n}}$.

\smallskip\noindent\emph{(b) The two box pairs.}
For $j\le k$ put
\begin{equation*}
\mathfrak{B}^*_j := \mathbb{B}_j(\square_0)\cup
\bigl\{\Gamma^{(\mathsf{n})}_{i'}\cap\square_0\bigr\}_{i'<j} ,
\qquad
\mathfrak{b} := e^{i\eta\mathbf{A}}\,\mathbb{W} ,
\end{equation*}
where $\Gamma^{(\mathsf{n})}_{i'}$ is the set of level $i'$ of stage $\mathsf{n}$ entering the
refined block sequence in the statement fixing $\mathfrak{B}^*_j$. Below,
$(U,J)(\cdot,\cdot)$ denotes the configuration pair of which $U(\cdot,\cdot)$ of part
\textup{(a)} is the first component.
The \emph{first box pair} is
\begin{equation*}
\mathsf{P}^0 := \mathfrak{b}'_j := (U,J)\bigl(\mathfrak{B}^*_j,M^{\cdot}(\mathfrak{b})\bigr),
\end{equation*}
the pair at which a constituent term $T$ of level $j$ of a near unit (see \textup{(d)}) is
evaluated. The \emph{second box pair}, which uses the map $M_*$ of the same construction in
place of $M^{\cdot}$, is
\begin{equation*}
\mathsf{P}^1(t,\mathsf{t}) := (U,J)\bigl(\mathfrak{B}^*_j,e^{i\mathsf{E}}M_*\mathfrak{b}\bigr),
\end{equation*}
where $\mathsf{E}$ is the stored-remainder phase, a function of cubes $c$, to each of which an
enlarged block $\tilde{B}(c)$ is attached. We recall the properties of $\mathsf{E}$ that
come with its definition. First, $\mathsf{E}$ vanishes on the cubes whose enlarged block lies in
$\tilde{\square}^3$:
\begin{equation}\label{5.5:eq-windowed-source-1}
\mathsf{E}(c)=0\quad\text{whenever }\tilde{B}(c)\subset\tilde{\square}^3 .
\end{equation}
Second, the set
\begin{equation*}
F_{\square} := \bigl\{c:\ \text{a block of } c \text{ meets the complement of }
\tilde{\square}^3\bigr\}
\end{equation*}
contains $\operatorname{supp}\mathsf{E}$. For a near unit $v$ with own box $\mathfrak{r}_v$
(defined in \textup{(e)}) put
\begin{equation*}
D_v := d_\star(\mathfrak{r}_v,F_{\square}),
\end{equation*}
a geometric distance, independent of the configurations. Third, $D_v$ obeys the lower bound
\begin{equation}\label{5.5:eq-windowed-source-2}
D_v\ \ge\ \tfrac15\,M L^{n_0-j}.
\end{equation}

\smallskip\noindent\emph{(c) The own-level family and the radii.}
For a term $T$ of level $\ell_T$ with domain $X_T$, the \emph{own-level family} is
$\mathcal{W}_T := \mathcal{W}^{\ell_T,\square}\restriction X_T$, the restriction to $X_T$ of the
family $\mathcal{W}^{\ell_T,\square}$ formed from $\mathcal{W}$ at the level $\ell_T$ and the
block $\square$ in the statement introducing the own-level families. Put
\begin{equation*}
K'' := 28K^+,\qquad \mathsf{w}_0 := \tfrac{1}{24},\qquad
\vartheta^{\natural} := \vartheta/\mathsf{w}_0,\qquad
K_A := 4\mathsf{K}_r e^{C_2\kappa_1},
\end{equation*}
where $C_2$ is the constant multiplying $\kappa_1$ in the factor $e^{-C_2\kappa_1}$ of the box
radius $r^{\sharp}_{\square}$, fixed with that radius in the statement introducing it.
The box radius $r^{\sharp}_{\square}$ of that statement is
\begin{equation*}
r^{\sharp}_{\square} = \frac{e^{-C_2\kappa_1}\,
e^{\frac12\delta d^{-}(\square)}}{2\mathsf{K}_r\vartheta},
\end{equation*}
where $d^{-}(\square)$ is that statement's distance of the block from the collar and $\delta$ its
decay rate, so that $\vartheta e^{-\frac12\delta d^{-}(\square)}$ is the normaliser of the block
in that statement. This normaliser is the normaliser $\vartheta(\square)$ of
\eqref{5.5:eq-windowed-source-a}, that is,
$e^{-\frac12\delta d^{-}(\square)} = e^{-\frac12\delta_P d^{\wedge}(\square,R_{\mathsf{n}})}$,
so that
\begin{equation*}
r^{\sharp}_{\square} = \frac{e^{-C_2\kappa_1}}{2\mathsf{K}_r\,\vartheta(\square)} .
\end{equation*}
In what follows $d_\ell$ is the distance measured at the level of the term and
$\square^{\zeta}$ denotes the set attached to the block's cut-off functions.
The normaliser, the own-level radius $r_T(\square)$, the factor $\eta_p(v)$ and the box radius
are related as follows:
\begin{equation}\label{5.5:eq-windowed-source-a}
\begin{gathered}
\vartheta(\square) := \vartheta\,e^{-\frac12\delta_P d^{\wedge}(\square,R_{\mathsf{n}})},
\qquad
r_T(\square) := \frac{e^{\frac14\delta d_\ell(X_T,\square^{\zeta})}}{14K^+\vartheta(\square)},
\\
\eta_p(v) := e^{-\frac14\delta d_\ell(X_T,\square^{\zeta})}\ \le\ 1,
\qquad
\frac{2}{r^{\sharp}_{\square}} = K_A\,\vartheta(\square).
\end{gathered}
\end{equation}

\smallskip\noindent\emph{(d) Near units.}
A \emph{near unit} at $\square$ is a unit $v$ of one of the first three rows or of the fifth row
of the table of units, with $X_v\subset\tilde{\square}^2$. Let $y$ be
its anchor and $j$ its level, and put
\begin{gather*}
n := \min\{\ell(y),n_0\},\qquad \nu := 1+n-j,\qquad t := L^{j-n},\\
w^{\vee}(y) := \max\bigl\{w(\Delta):\ \operatorname{dist}_n(\Delta,y)<1\bigr\}.
\end{gather*}
A \emph{unit $n$-cell} $\Delta'$ is either a cell of level at most $n_0$, or a sub-cube of level
$n_0$ of a cell of level $n_0+1$; $n(\Delta')$ denotes its level. For such $\Delta'$ put
\begin{equation*}
w^{\vee}(\Delta') := \max\bigl\{w(\Delta):\ \operatorname{dist}_{n(\Delta')}(\Delta,\Delta')<1\bigr\};
\end{equation*}
then $w^{\vee}(\Delta')\ge w^{\vee}(y)$ for $y\in\Delta'$. Let $\mathfrak{R}(\square)$ be the set
of the cells within such reach of $\tilde{\square}^2$, the reach being the one in the definition
of $w^{\vee}(\Delta')$; each of them meets
$\tilde{\square}^3$, and $\tilde{\square}^2$ lies in a cube of side six own units. Put
\begin{equation*}
\Theta := c_W\,w^{\vee}(y),
\end{equation*}
and note $\Theta\ge1$.

\smallskip\noindent\emph{(e) The own box and the weighted scaffold majorants.}
The \emph{own box} $\mathfrak{r}_v$ of a near unit $v$ is as follows. If $v$ is a core,
$\mathfrak{r}_v := \square_y^{\sim(\nu+1)}$, where $\square_y$ is the cube of $\pi_j$ containing
$y$ (not the block $\square$) and $\sim$ is the enlargement by cubes of $\pi_j$ in the sense of
\cite{B-CMP119}; this box carries one layer of $\pi_j$-cubes of margin. If $v$ is a short
single, $\mathfrak{r}_v$ is the box of $\hat{X}$ with a margin of $\frac52$ cubes of $\pi_j$.
Put $D := M(2\nu+3)$, and put $D_X := \operatorname{diam}_j X+5M$ for a single, where $X$ denotes
the domain of the unit in the table of units, and
$D_X := D$ for a core.
On $\mathfrak{r}_v$, in the comb gauge at $y$, the \emph{weighted scaffold majorants} of $v$ are
\begin{equation}\label{5.5:eq-windowed-source-b}
\begin{aligned}
|B^{ax}|_*(b) &:= 2d\,(1+|b-y|)\,\Theta\,\mathfrak{a}_n t^2,
&\qquad |F|_* &:= \Theta\,\mathfrak{a}_n t^2,\\
|\mathfrak{e}_\mu|_*(x) &:= c_1(1+|x-y|)^2\,\Theta\,\mathfrak{a}_n t^{2+\hat{\alpha}},
&\qquad |\mathfrak{e}^{rim}|_* &:= 4dD\,\Theta\,\mathfrak{a}_n t^2 ,
\end{aligned}
\end{equation}
where $b$ runs over the bonds and $x$ over the points of $\mathfrak{r}_v$ (this $b$ is not the
datum $b$ of part \textup{(a)}), and $\hat{\alpha}$ is the H\"older exponent and $c_1$ the
constant of the scaffold bound for the fields $\mathfrak{e}_\mu$, both fixed with the scaffold
bounds of near units.
The bounds of the scaffold fields of $v$ by these majorants are established with the scaffold
bounds of near units. On the graded rim of $\mathfrak{r}_v$ the pointwise value
$4d(1+|b-y|)\Theta\mathfrak{a}_n t^2/L$ is at most $4dD_X\Theta\mathfrak{a}_n t^2$, because
$1+|b-y|\le D_X$ there; for a core this is $|\mathfrak{e}^{rim}|_*$.
\end{definition}

\begin{proof}[Proof of the relations stated in Definition~\ref{5.5:def-windowed-source}]
The vanishing property \eqref{5.5:eq-windowed-source-1}, the inclusion
$\operatorname{supp}\mathsf{E}\subset F_{\square}$ and the lower
bound \eqref{5.5:eq-windowed-source-2} are those of the stored-remainder phase $\mathsf{E}$ that
accompany its definition; they are used here as stated with that definition.

The inequality in \eqref{5.5:eq-windowed-source-a} for $\eta_p(v)$: the quantity
$d_\ell(X_T,\square^{\zeta})$ is a distance, hence non-negative, and $\delta>0$; so the exponent
$-\frac14\delta d_\ell(X_T,\square^{\zeta})$ is non-positive and $\eta_p(v)\le1$.

The identity in \eqref{5.5:eq-windowed-source-a}: by \textup{(c)},
$r^{\sharp}_{\square}=e^{-C_2\kappa_1}/(2\mathsf{K}_r\vartheta(\square))$, the normaliser of the
statement fixing $r^{\sharp}_{\square}$ being $\vartheta(\square)$. Inverting this expression and
using $K_A=4\mathsf{K}_re^{C_2\kappa_1}$,
\begin{align*}
\frac{2}{r^{\sharp}_{\square}}
&= 2\cdot 2\mathsf{K}_r\, e^{C_2\kappa_1}\,\vartheta(\square)\\
&= 4\mathsf{K}_r e^{C_2\kappa_1}\,\vartheta(\square)
= K_A\,\vartheta(\square) .
\end{align*}

The inequality $w^{\vee}(\Delta')\ge w^{\vee}(y)$ for $y\in\Delta'$: we first show
$n(\Delta')=\min\{\ell(y),n_0\}=n$. If $\Delta'$ is a cell of level at most $n_0$, it is the
cell containing $y$, so its level is the current level $\ell(y)\le n_0$ and $n=\ell(y)=n(\Delta')$.
If $\Delta'$ is a sub-cube of level $n_0$ of a cell of level $n_0+1$, then $\ell(y)=n_0+1$ and
$n=n_0=n(\Delta')$. Since $y\in\Delta'$ and, by the definition of the distance
$\operatorname{dist}_n$ between sets, the distance of $\Delta$ to the set $\Delta'$ is the
infimum of its distances to the points of $\Delta'$,
\begin{equation*}
\operatorname{dist}_{n(\Delta')}(\Delta,\Delta') = \operatorname{dist}_n(\Delta,\Delta')
\le \operatorname{dist}_n(\Delta,y).
\end{equation*}
Hence every cell $\Delta$ with $\operatorname{dist}_n(\Delta,y)<1$ also satisfies
$\operatorname{dist}_{n(\Delta')}(\Delta,\Delta')<1$: the set over which the maximum defining
$w^{\vee}(y)$ is taken is contained in the set over which the maximum defining $w^{\vee}(\Delta')$
is taken, and $w^{\vee}(\Delta')\ge w^{\vee}(y)$.

The bound $\Theta\ge1$ holds by the choice of the constant $c_W$, which is fixed with the
scaffold bounds of near units; it is used here as stated with that choice.

The comparison on the graded rim: on the graded rim of $\mathfrak{r}_v$ one has
$1+|b-y|\le D_X$, a property of the own box established with the scaffold bounds of near units.
Since $L\ge1$,
\begin{equation*}
\frac{4d(1+|b-y|)\,\Theta\,\mathfrak{a}_n t^2}{L}
\ \le\ 4d(1+|b-y|)\,\Theta\,\mathfrak{a}_n t^2
\ \le\ 4dD_X\,\Theta\,\mathfrak{a}_n t^2 .
\end{equation*}
For a core, $D_X=D$ by \textup{(e)}, and the right-hand side is $|\mathfrak{e}^{rim}|_*$ of
\eqref{5.5:eq-windowed-source-b}.
\end{proof}

\begin{definition}[Three classes of near units and their discs]\label{5.5:def-near-unit-classes}
Fix a block $\square$. We use the windowed source, the box pairs $\mathfrak{b}'_j=\mathsf{P}^0$ and
$\mathsf{P}^1$, the own-level family $\mathcal{W}_T$, the normaliser $\vartheta(\square)$, the radii
$r_T(\square)$ and $r^{\sharp}_{\square}$, the factor $\eta_p(v)\le 1$, the constants $K_A$, $K^+$,
$K''=28K^+$ and $\mathsf{w}_0=1/24$, and the weighted scaffold majorants, all as fixed in
Definition~\ref{5.5:def-windowed-source}, in particular in \eqref{5.5:eq-windowed-source-a} and
\eqref{5.5:eq-windowed-source-b}. Let $v$ be a near unit of level $j$ in the sense of that
definition (a unit of rows~1--3 or~5 of the table of terms used there), with anchor $y$, capped
level $n=\min\{\ell(y),n_0\}$ of the anchor, index $\nu=1+n-j$, ratio $t=L^{j-n}$, weight factor
$\Theta=c_W w^{\vee}(y)\ge 1$, own box $\mathfrak{r}_v$ and length $D_X$.
The radii $a_j$ and $\alpha_{3,j}$ and the constant $c_3$ are those of
Definition~\ref{5.1:def-analyticity-spaces}, and $C_M$ and $\rho_{n,j}$ are those of
Notation~\ref{5.1:not-coupling-functions}; the constant $K^{rem}$ is the one of the bound for
stored remainders.

\emph{Readable units and the scaffold radius.} A near unit of rows~1--3 is called \emph{readable} if
$\nu\ge 3$ and it is a core member or a short single. Put $S^{\star}_v$ as in the first relation of
the following display; this number satisfies the second relation:
\begin{equation}\label{5.5:eq-near-unit-classes-a}
S^{\star}_v := \frac{\alpha_{3,j}}{\sup_{b\subset\mathfrak{r}_v}|B^{ax}|_*(b)},
\qquad
\bigl(\Theta C_M D_X t^2\rho_{n,j}\bigr)^{-1}\le S^{\star}_v .
\end{equation}
The number $S^{\star}_v$ is a scaffold number: it is defined through the weighted scaffold
majorant $|B^{ax}|_*$ of \eqref{5.5:eq-windowed-source-b}, not through the field. The second
relation in \eqref{5.5:eq-near-unit-classes-a} is the inverse-radius bound of the lemma on the
radius of the local chart of near units, whose step
\begin{equation*}
S^{-1}\le \frac{2dD_X\mathfrak{a}_n t^2}{c_3 a_j}
\end{equation*}
is applied here with $\mathfrak{a}_n$ replaced by $\Theta\mathfrak{a}_n$; the bound is used below in
this form. In the comb gauge at $y$, consider configurations on which the pointwise bounds
$|B^{ax}|(b)\le|B^{ax}|_*(b)$, $b\subset\mathfrak{r}_v$, hold, as given by the statement that
proves, in the comb gauge at $y$, the pointwise bounds of the axial field by the weighted scaffold
majorants of \eqref{5.5:eq-windowed-source-b} on the own box. For them the radius of that lemma,
\begin{equation*}
S = \frac{\alpha_{3,j}}{\sup_{b\subset\mathfrak{r}_v}|B^{ax}|(b)},
\end{equation*}
satisfies $S\ge S^{\star}_v$ (see the proof below).

\emph{The three classes.}
\begin{itemize}
\item[(a)] A near unit $v$ is \emph{box-read} if it is readable and $S^{\star}_v>1$. A box-read core
is then read as a whole on the one box pair $\mathfrak{b}'_j$.
\item[(b)] A near unit of rows~1--3 is \emph{own-read} if it is not box-read. This comprises the
tails, the long singles, the constituents of units with $\nu\le 2$, and the readable units with
$S^{\star}_v\le 1$, the last taken constituent by constituent. Own-read units are read on the
own-level family $\mathcal{W}_T$.
\item[(c)] The \emph{third class} consists of every near Wilson or counterterm piece: the pieces
$f_{\square}$ of row~5 and the pieces $\beta[F]A(h_y,\cdot)$ of a core, both in the notation of
the table of terms used in Definition~\ref{5.5:def-windowed-source}. In particular this piece
of a straddling core (its companion), that is, of a core with anchor $y\in\tilde{\square}^2$ and
own box $\mathfrak{r}_v\not\subset\tilde{\square}^2$, is included in the third class wherever it is
near, although the membership corollary for own-read units assigns it to none of its classes; it is
placed in the third class by the present definition. A piece of the third class is taken at
every $\nu$ and every value of $S^{\star}_v$, and is read on the box pair $\mathfrak{b}'_j$; it is
entire in the parameters, and no membership statement is required for it. Its size is bounded by
item (a2) of the box-pair estimate for near units of this section together with item (b) of the
bound for Wilson and counterterm pieces. The far pairs of the first family $\mathcal{T}_1$ of pairs
of terms in the resummation are read on $\mathcal{W}_T$.
\end{itemize}

\emph{Stored remainders.} Remainders of the comparison of the two box pairs exist exactly at the
pieces read on $\mathfrak{b}'_j$, that is, at the box-read units and at the pieces of the third
class.
For a box-read unit and a constituent term $T$ with domain $X\subset\mathfrak{r}_v$ the remainder
is stored as
\begin{equation*}
\operatorname{rem}(T,\square) := \int_0^1 dt\,\partial_{\mathsf{t}}\big|_{\mathsf{t}=0}
\bigl[T\bigl(X;\mathsf{P}^1(t,\mathsf{t})\bigr)-T\bigl(X;\mathsf{P}^0\bigr)\bigr],
\end{equation*}
where the integral is over the interpolation parameter $t\in[0,1]$ of the fluctuation
$\mathbf{A}$ of Definition~\ref{5.5:def-windowed-source}; in this display, and in
$\mathcal{W}_T(t,\tau)$ below, the letter $t$ denotes that parameter and not the ratio $L^{j-n}$.
The remainders of the pieces of the third class are those of the remainder statement for the
Wilson and counterterm pieces. There is no remainder at own-read units.

\emph{The three discs.} Every application of Lemma~\ref{5.2:lem-disc-bounds} to a piece of the
classes (a)--(c) or to a stored remainder, in the estimates that follow, is made on exactly one of
the following three discs: the box-pair disc
$\mathbb{D}_{\square}$ (for the pieces of classes (a) and (c)), the own-read disc in the parameter
$\tau$ of the own-level family $\mathcal{W}_T(t,\tau)$ (for class (b)), and the remainder disc (for
the stored remainders). They and their prefactors are
\begin{equation}\label{5.5:eq-near-unit-classes-b}
\begin{aligned}
&\mathbb{D}_{\square} := \{\mathsf{t}\in\mathbb{C}: |\mathsf{t}|\le r^{\sharp}_{\square}\},
&&\frac{2}{r^{\sharp}_{\square}} = K_A\vartheta(\square),\\
&\{\tau\in\mathbb{C}: |\tau|\le\rho_A(v)\},\quad \rho_A(v):=\mathsf{w}_0\, r_T(\square),
&&\frac{2}{\rho_A(v)} = 24K''\vartheta(\square)\eta_p(v),\\
&\{\mathsf{t}\in\mathbb{C}: |\mathsf{t}|\le r^{rem}_{\square}\},
&&\frac{2}{r^{rem}_{\square}} \le (K_A+28K^{rem})\vartheta(\square),
\end{aligned}
\end{equation}
where
\begin{equation*}
r^{rem}_{\square} := \min\Bigl\{r^{\sharp}_{\square},\ \frac{1}{14K^{rem}\vartheta(\square)}\Bigr\}.
\end{equation*}
These three discs are the disc factors of the parameter domains of the box read $\mathsf{M}2$, the
own-level read $\mathsf{M}3$ and the remainder read $\mathsf{M}4$ among the moves of the pair
system $\mathcal{K}$.
All three prefactors are proportional to $\vartheta(\square)$. The site-wide resummation weight
$\mathsf{w}_P$, which is proportional to $1/\vartheta$ and is used on these discs in the
resummation of the correlation theorem (Theorem~\ref{5.1:thm-correlation-theorem}) and in the rim
bounds, enters none of the definitions above; it satisfies
$\mathsf{w}_0\mathsf{w}_P/\eta_p(v)\le\rho_A(v)$ by the membership corollary for own-read units.
The predicates (a)--(c) and the three
discs are scaffold quantities: they are defined through the weighted scaffold majorants and the
radii of Definition~\ref{5.5:def-windowed-source}, not through the field. The factor
$24=1/\mathsf{w}_0$ in the prefactor of the own-read disc enters the constant $\kappa_R$ and, at
every block, the per-block constant $C_R$; each of the two is the least number for which its
bound holds, and both are fixed before $\gamma$ in the order of choice of the constants
(Definition~\ref{5.1:def-admitted-inputs}).
\end{definition}

\begin{proof}
We prove the relations stated in the definition other than the inverse-radius bound in
\eqref{5.5:eq-near-unit-classes-a}, which is the bound of the lemma named there, in the form
indicated.

\emph{The field radius.} For the configurations described in the definition, the axial field obeys,
in the comb gauge at $y$, the pointwise bound $|B^{ax}|(b)\le|B^{ax}|_*(b)$ for every
$b\subset\mathfrak{r}_v$, by the statement that proves, in the comb gauge at $y$, the pointwise
bounds of the axial field by the weighted scaffold majorants of \eqref{5.5:eq-windowed-source-b}
on the own box. Taking the supremum over $b\subset\mathfrak{r}_v$ gives
\begin{equation*}
\sup_{b\subset\mathfrak{r}_v}|B^{ax}|(b)\le\sup_{b\subset\mathfrak{r}_v}|B^{ax}|_*(b).
\end{equation*}
Both suprema are non-negative, and the one on the right is positive, since
$|B^{ax}|_*(b)=2d(1+|b-y|)\Theta\mathfrak{a}_nt^2$ with $\Theta\ge 1$, $\mathfrak{a}_n>0$ and
$t>0$. If the supremum on the left vanishes, then $S=+\infty\ge S^{\star}_v$. Otherwise both sides
are positive, and since $\alpha_{3,j}>0$, dividing $\alpha_{3,j}$ by the two sides reverses the
inequality and yields $S\ge S^{\star}_v$.

\emph{The box-pair disc.} The identity $2/r^{\sharp}_{\square}=K_A\vartheta(\square)$ is one of the
relations of \eqref{5.5:eq-windowed-source-a}.

\emph{The own-read disc.} By \eqref{5.5:eq-windowed-source-a},
$r_T(\square)=e^{\frac14\delta d_\ell(X_T,\square^\zeta)}/(14K^+\vartheta(\square))$ and
$\eta_p(v)=e^{-\frac14\delta d_\ell(X_T,\square^\zeta)}$. With $\mathsf{w}_0=1/24$ this gives
\begin{equation*}
\frac{2}{\rho_A(v)}=\frac{2}{\mathsf{w}_0\, r_T(\square)}
=2\cdot 24\cdot 14K^+\vartheta(\square)\,e^{-\frac14\delta d_\ell(X_T,\square^\zeta)}
=2\cdot 24\cdot 14K^+\vartheta(\square)\eta_p(v).
\end{equation*}
Since $K''=28K^+=2\cdot 14K^+$, the right-hand side equals $24K''\vartheta(\square)\eta_p(v)$.

\emph{The remainder disc.} Both numbers in the definition of $r^{rem}_{\square}$ are positive, so
\begin{equation*}
\frac{2}{r^{rem}_{\square}}
=\max\Bigl\{\frac{2}{r^{\sharp}_{\square}},\ 28K^{rem}\vartheta(\square)\Bigr\}
=\max\bigl\{K_A\vartheta(\square),\ 28K^{rem}\vartheta(\square)\bigr\},
\end{equation*}
where the box-pair identity was used in the second step. The maximum of two non-negative numbers is
at most their sum, whence $2/r^{rem}_{\square}\le(K_A+28K^{rem})\vartheta(\square)$.

\emph{Proportionality.} The three prefactors are $K_A\vartheta(\square)$,
$24K''\eta_p(v)\vartheta(\square)$ and $(K_A+28K^{rem})\vartheta(\square)$, each a multiple of
$\vartheta(\square)$.
\end{proof}

\begin{definition}[The data factor and its target bound]\label{5.5:def-data-factor}
Work in the setting of Notation~\ref{5.5:not-weight-notation} and
Definition~\ref{5.5:def-windowed-source}, with $\vartheta(\square)$, $\mathcal{W}$,
$\mathsf{L}_0=L_{\circ}$ and $\mathfrak{K}(\square)$ as fixed there. Thus $\square$ is a block
of zone $n_0$ and $\tilde{\square}^5=\square_0$ its enlargement. The collar $R_{\mathsf{n}}$ is
that of stage $\mathsf{n}$, with $\mathsf{j}=h+\mathsf{n}$. Let $\mathsf{T}_3(\square)$ denote the
total at $\square$ of the size sum of the weight lemma that is bounded in the
count-times-weight lemma.

\begin{enumerate}
\item[(a)] \emph{Base data.} With the constants $c_\alpha$, $C_R$, $E_0$, $C_2$, $C_N$ and the
raise bound $N_0$ of the weight lemma, put
\begin{equation*}
\mathfrak{n}_0 := c_\alpha C_R (E_0+1)\,e^{C_2\kappa_1},
\qquad
\mathfrak{N}_0(\square) := c_\alpha E_0 C_R\, e^{C_2\kappa_1}\,\vartheta(\square).
\end{equation*}

\item[(b)] \emph{The raise of the block.} Let $s(\square)$ be the largest raise
$\mathsf{r}_0$ of the present operation $\mathbf{R}_k$ over the cells meeting
$\tilde{\square}^5$. It satisfies $s(\square)\le N_0$.

\item[(c)] \emph{The data factor.} If $\tilde{\square}^5$ meets no cell of $\mathcal{W}$, put
$\mathfrak{d}(\square):=1$. Otherwise put
\begin{equation*}
\mathfrak{d}(\square):=\max\Bigl\{1,\
\frac{\mathsf{T}_3(\square)}{\mathsf{L}_0^{6s(\square)}\mathfrak{N}_0(\square)}\Bigr\}.
\end{equation*}
Thus $\mathfrak{d}(\square)\ge 1$ at every block. The \emph{data factor} of $\square$ is
$\mathsf{L}_0^{6s(\square)}\mathfrak{d}(\square)$.

\item[(d)] \emph{The size datum $\mathfrak{n}(\square)$.} With $\mathfrak{p}_{\mathsf{j}}$ the
collar datum of stage $\mathsf{n}$, put
\begin{equation}\label{5.5:eq-data-factor-1}
\mathfrak{n}(\square) := \mathsf{L}_0^{6s(\square)}\mathfrak{d}(\square)\cdot\mathfrak{n}_0
\cdot\bigl(\vartheta(\square)+\mathfrak{p}_{\mathsf{j}}\,
\mathbb{1}[\square\cap R_{\mathsf{n}}\neq\emptyset]\bigr).
\end{equation}
The one-maximum data functional $\mathsf{l}^2(\square)$ majorises the data factor at the
blocks where $\mathfrak{d}$ is not $1$ by definition: there, by the count-times-weight lemma,
\begin{equation*}
\mathsf{l}^2(\square)\ge \mathsf{L}_0^{6s(\square)}\mathfrak{d}(\square).
\end{equation*}

\item[(e)] \emph{Auxiliary constants.} For a block $\square$ put
\begin{equation*}
\varrho(\square) := 1+\frac{d^{\wedge}(\square,R_{\mathsf{n}})}{w},
\end{equation*}
where the denominator $w$ is the fixed length of the weight lemma by which the distance to
the collar is scaled; it is not the actual weight $w(x)$ of
Notation~\ref{5.5:not-weight-notation}. For $\varrho\ge 1$ put
\begin{equation*}
\mathsf{n}(\varrho) := C_N 2^{r_{pr}}(\varrho+4)+\varrho+4
\end{equation*}
(a function of $\varrho$, not the stage letter $\mathsf{n}$ of the collar $R_{\mathsf{n}}$).
Define
\begin{equation}\label{5.5:eq-data-factor-2}
C^{\mathrm{data}}_{0} := \sup_{\varrho\ge 1}\,
L^{12}\bigl(\varrho+7+c_{\mathsf{n}(\varrho)}\bigr)^{6r_{pr}}e^{2-\varrho}.
\end{equation}

\item[(f)] \emph{Properties of the blocks concerned.} The following properties are
established in the weight lemma and are used with this definition.
\begin{itemize}
\item[(f1)] $s=0$ on the old-block set $\mathfrak{O}^+$.
\item[(f2)] Suppose $\tilde{\square}^5$ meets $\mathcal{W}$. Then either
$\square\in\mathfrak{O}^+$, or $\square$ is exceptional in the sense of the weight lemma, that
is, a block of zone $k$ over the band of an exit of the last stage of the kind fixed there.
In both cases
\begin{equation*}
\varrho(\square)\ge \check{R}_k-4,
\qquad
k-n_0+1\le N_*+\varrho(\square),
\end{equation*}
where $N_*$ is the constant of the weight lemma that bounds the depth $k-n_0+1$ in this way.
\end{itemize}

\item[(g)] \emph{The target bound.} The \emph{target bound} for the data factor is the
statement that there is a constant $C^{\mathrm{data}}\ge C^{\mathrm{data}}_0$, chosen in the
order of choice of Definition~\ref{5.1:def-admitted-inputs} before $\gamma$, such that at
every block $\square$
\begin{equation}\label{5.5:eq-data-factor-a}
\mathfrak{d}(\square)\le 1+C^{\mathrm{data}}\,
e^{\frac18\delta_P d^{\wedge}(\square,R_{\mathsf{n}})},
\end{equation}
the bound being required both for the third and for the fourth item of the list of size
estimates in the weight lemma.

\item[(h)] \emph{Scope.}
\begin{itemize}
\item[(h1)] Let $\tilde{\square}^5$ meet no cell of $\mathcal{W}$. Then
$\mathfrak{d}(\square)=1$ by (c). Since $C^{\mathrm{data}}\ge C^{\mathrm{data}}_0\ge 0$, the
right member of \eqref{5.5:eq-data-factor-a} is at least $1$, so
\eqref{5.5:eq-data-factor-a} holds at such blocks by definition.
\item[(h2)] At the blocks of (h1), the sum $\mathfrak{K}(\square)$ is the sum of the tiling
lemma (cf.\ \cite{B-CMP122a}), and the tiling lemma bounds it there. At those blocks the
corollary to the tiling lemma uses the first clause of the count-times-weight lemma, namely
the chain of inequalities
\begin{equation*}
\mathsf{T}_3(\square)\le \kappa_R c_W^3\,\mathfrak{K}(\square)\,\mathfrak{N}_0(\square)
\le \mathsf{l}^2(\square)\,\mathfrak{N}_0(\square),
\end{equation*}
which holds at every block.
\item[(h3)] The count-times-weight lemma bounds $\mathfrak{K}(\square)$, and proves
\eqref{5.5:eq-data-factor-a}, only at the blocks with
$\tilde{\square}^5\cap\mathcal{W}\ne\emptyset$.
\end{itemize}

\item[(i)] \emph{The filing hull.} With $\pi^{\flat}$ the partition of the weight lemma, put
\begin{equation*}
\hat{\mathsf{K}}_{\square} := \bigcup\{G\in\pi^{\flat} : G\cap\tilde{\square}^5\ne\emptyset\}.
\end{equation*}
\end{enumerate}
\end{definition}

\begin{definition}[Constants of the weight lemma]\label{5.5:def-weight-constants}
The numbers fixed in this definition are chosen after $(d,L,M,\kappa,\kappa_1,\lambda,c_2)$ and
before $\gamma$. None of them depends on an age, on $N_0$, on a depth,
on $K$ or on a weight.

The following objects are those fixed earlier in this section:
\begin{itemize}
\item from Definitions~\ref{5.5:def-windowed-source} and~\ref{5.5:def-near-unit-classes}: the
blocks $\square$ and their enlargements $\tilde{\square}^2$, $\tilde{\square}^5$; the normaliser
$\vartheta(\square)$ of \eqref{5.5:eq-windowed-source-a}; the
filing of near units into three classes; the filed discs of \eqref{5.5:eq-near-unit-classes-b},
with their constants $K_A$, $K''$, $K^{+}$, $K^{rem}$, the factor $\eta_p\le 1$ and
$\mathsf{w}_0=1/24$;
\item from Definition~\ref{5.5:def-data-factor}: the family $\mathcal{W}$ of cells; the raise
$s(\square)$; the weight base $L_{\circ}=\mathsf{L}_0$; the exponent $r_{pr}$; the function
$\mathsf{n}(\varrho)$; the constant $C^{\mathrm{data}}_{pr}$; and the base size
\begin{equation*}
\mathfrak{N}_0(\square)=c_\alpha E_0C_Re^{C_2\kappa_1}\vartheta(\square).
\end{equation*}
\end{itemize}
Let $n_0$ be the zone of the block $\square$ (Definition~\ref{5.5:def-windowed-source}), and
$w(\Delta)$ the weight of a cell $\Delta$. A unit $n$-cell is a cell of level at most $n_0$, or a
sub-cube of level $n_0$ of a cell of level $n_0+1$. For a unit $n$-cell $\Delta'$ we write
$n(\Delta')$ for its level, and $w^{\vee}(\Delta')$ denotes the largest weight $w(\Delta)$ over
the cells $\Delta$ with $\operatorname{dist}_{n(\Delta')}(\Delta,\Delta')<1$.

\emph{Local constants.}
\begin{itemize}
\item $K_{ax}$ is the constant $O(1)$ of \cite[(3.26)]{B-CMP109} on a box of at most nine cubes
a side.
\item $C^{av}_M\ge 2$ is the constant of \cite[Proposition~2]{B-CMP98}, in the form extended to
complex data by the sixth-family lemma.
\item $K_h$ is the constant of the local bound from which the weighted scaffold majorants of
Definition~\ref{5.5:def-windowed-source} are obtained.
\item $c_1$ is the constant in the majorant of $\mathfrak{e}_\mu$ there.
\end{itemize}
These satisfy $K_h\ge1$, $K_{ax}M\ge1$ and $c_1\ge1$.

Let $\alpha_{i,n}$ be the analyticity radii at level $n$, $\alpha_{\cdot,n}$ the radius of this
family carrying the index $\cdot$ (the radius entering the sizes of the remainders of the third
class), and $\mathfrak{a}_n$ the level-$n$ scaffold smallness constant. Put
$c_{\cdot}:=\sup_n\alpha_{\cdot,n}/\alpha_{0,n}$.

Let $C'$ be the constant of the bound for the projection $\Pi^{\flat}$, by which the counterterm
companion of a core is estimated. Let $C_I^{\sharp}$ be the constant of the bound for the core
term $\mathfrak{i}_y$, with the rim contribution included. Let $C(\kappa)$ be a constant
depending only on $\kappa$, and $C_M$ the constant in the radius inequality of
Definition~\ref{5.5:def-near-unit-classes}. Put
\begin{equation}\label{5.5:eq-weight-constants-c}
\begin{aligned}
K_{loc}&:=4\bigl(c_{\cdot}+3K_hK_{ax}MC^{av}_M\bigr),\\
c_W&:=\max\Bigl\{1,\;2c_1K_{loc}L^3\sup_n\frac{\alpha_{0,n}}{\mathfrak{a}_n}\Bigr\},\\
C_I'&:=\max\Bigl\{C_I^{\sharp},\;8C(\kappa)(10C_MM)^{5/2},\;4C'(10C_MM)^{1/2}\Bigr\}.
\end{aligned}
\end{equation}
The one value of $K_{loc}$ in \eqref{5.5:eq-weight-constants-c} is the value used in $c_W$, in
the ceiling \eqref{5.5:eq-weight-constants-d} and in the constants and ceilings of the
remainder bounds. The factor $L^3$ in $c_W$ is the spare power of $L$ used in the remainder
bounds. $C_I''$ denotes the constant which replaces $C_I'$ in the bounds for differences of two
machines on one scaffold; it is obtained from $C_I'$ by a further maximisation, so that
$C_I''\ge C_I'$.

\emph{The size-free constant $\kappa_R$.} The total at $\square$ in question is the total at
$\square$ of Definition~\ref{5.5:def-data-factor}, the one whose ratio to
$L_{\circ}^{6s(\square)}\mathfrak{N}_0(\square)$ defines the data factor $\mathfrak{d}(\square)$.
For a block $\square$ and a unit $n$-cell $\Delta'$, the
\emph{contribution of $\Delta'$} to this total is the one made by the near units anchored in
$\Delta'$. It consists of:
\begin{itemize}
\item the pieces of all three classes;
\item the stored remainders of the box-pair-read pieces.
\end{itemize}
Each of these is formed by the Cauchy integral in the complex source parameter $\mathsf{t}$ on its
filed disc, in the form of the lemma of this chapter which bounds the pieces of near units by
Cauchy integrals in $\mathsf{t}$ over the filed discs. The prefactors are
\begin{equation*}
K_A\vartheta(\square),\qquad 24K''\vartheta(\square)\eta_p,\qquad
(K_A+28K^{rem})\vartheta(\square)
\end{equation*}
on the box-pair, own-read and remainder discs respectively. The pieces are then resummed from the
table $\mathfrak{N}^P(\cdot\,;1)$ of near-unit sizes of the weight lemma, with the scaffold
coefficients of the per-cell row table, that is, the table of per-cell bounds with rows 1--3,
$\partial$ and $5$; the near units sit in rows 1--3 and $5$.

The number $\kappa_R^{act}$ is the maximum of $1$ and the least number $q$ such that, for every
block $\square$ and every unit $n$-cell $\Delta'$, the contribution of $\Delta'$ at the actual
sizes is at most $q\,\mathfrak{N}_0(\square)$.

For a row $r$ of the per-cell row table, $c^{(r)}_{n,j}$ is the coefficient of row $r$ at levels
$(n,j)$; at sizes identically $1$ the per-cell row-$r$ sum at level $n$ is
$\sum_{j\le n}c^{(r)}_{n,j}$. Put
\begin{equation*}
K_r:=\sup_n\sum_{j\le n}c^{(r)}_{n,j},
\end{equation*}
the per-cell constant of row $r$ at size $1$. Explicitly:
\begin{itemize}
\item $K_1:=C_I''S_I$, $K_2:=CS_\omega$ with the constant $C$ of the second row, and $K_3:=1$,
where $S_I$, $S_\omega$ are as in Definition~\ref{5.2:def-comparison-constant};
\item $K_\partial$ is the per-cell constant of the row $\partial$;
\item for the fifth row, the remainder rows and the tail row, $K_r$ is the row's displayed
per-cell number at unit size, as transported, with the constants of the two machines included;
\item for the remainder row of the third class, $K_r$ is the per-cube number of the
$\circ$-slot bound, with the factor $C^{\perp}K^{W}$, where $C^{\perp}=720L^2$ is the constant
of the $\circ$-slot $\delta$-size bound and $K^{W}$ is the constant of the remainder bound of the
third class.
\end{itemize}

The number $\kappa_R^{\times}$ is the least number $q$ with the following property. For every
block $\square$, every unit $n$-cell $\Delta'$ and every filing, the contribution of $\Delta'$ is
formed as above but at arbitrary non-negative sizes whose per-cell row-$r$ sums are at most $K_r$
for every row $r$. The property is that every such contribution is at most
\begin{equation*}
q\,c_\alpha C_Re^{C_2\kappa_1}\vartheta(\square).
\end{equation*}
Put
\begin{equation}\label{5.5:eq-weight-constants-a}
\kappa_R:=\max\bigl\{\kappa_R^{act},\;\kappa_R^{\times}\bigr\}.
\end{equation}
The factor $24=1/\mathsf{w}_0$ of the own-read prefactor is contained in $\kappa_R$. The
quantities $C^{\mathrm{data}}$, $\mathsf{l}^2$ and $\mathfrak{d}^{\sharp}$ below are formed with
this $\kappa_R$.

\emph{Tiling and data constants; block quantities.} Put $c_s:=\log(1+3\lambda s+2c_2)$, the
function which also enters the constant $C^{\mathrm{data}}_{pr}$ of
Definition~\ref{5.5:def-data-factor} through $c_{\mathsf{n}(\varrho)}$, and
\begin{equation*}
y_*:=\sup\bigl\{y\ge1:\;y\le L\,(2c_{\log_Ly})^{r_{pr}}\bigr\}.
\end{equation*}
The constant $C^{\mathrm{data}}$ of the data-factor bound \eqref{5.5:eq-data-factor-a} is fixed
before $\gamma$, not in this definition, as the maximum
\begin{equation*}
\begin{split}
C^{\mathrm{data}}:=\max\Bigl\{&C^{\mathrm{data}}_{cnt},\;C^{\mathrm{data}}_{pr},\\
&\kappa_Rc_W^3\bigl[2(7L)^d(1+C^{\mathrm{data}}_{pr})+C^{tile}\bigr]\Bigr\},
\end{split}
\end{equation*}
in which $C^{\mathrm{data}}_{cnt}$ is the count constant of the correlation theorem
(Theorem~\ref{5.1:thm-correlation-theorem}), $C^{\mathrm{data}}_{pr}$ is the constant of
Definition~\ref{5.5:def-data-factor}, and the third member is formed with the constants of this
definition ($C^{tile}$ as in \eqref{5.5:eq-weight-constants-b}). In this section it is used only
through the inequalities $C^{\mathrm{data}}\ge C^{\mathrm{data}}_{pr}$ and the last line of
\eqref{5.5:eq-weight-constants-b}, both immediate from this maximum.

For every block $\square$, $\mathfrak{K}(\square)$ is the one sum displayed below. Where
$\tilde{\square}^5$ meets no cell of $\mathcal{W}$ it is bounded by the tiling bound. Where
$\tilde{\square}^5$ meets $\mathcal{W}$ it is bounded by the count-and-weight bound
\eqref{5.5:eq-weight-f}. There $\mathfrak{K}^{\sharp}(\square)$ denotes the right member of
\eqref{5.5:eq-weight-f} applicable at $\square$. We put the following, the last line recording
the inequality satisfied by $C^{\mathrm{data}}$:
\begin{equation}\label{5.5:eq-weight-constants-b}
\begin{aligned}
C^{tile}&:=3\cdot7^d\mathsf{L}_0^6L^d\Bigl[y_*^d\\
&\qquad\qquad+L^d2^{dr_{pr}}\sup_{\varrho\ge1}
\bigl(\varrho+5+c_{\mathsf{n}(\varrho)}\bigr)^{dr_{pr}}e^{1-\varrho}\Bigr],\\
\mathfrak{K}(\square)&:=\sum\bigl\{w^{\vee}(\Delta')^3:\;\Delta'\text{ a unit }n\text{-cell
inside }\tilde{\square}^2\bigr\},\\
\mathsf{l}^2(\square)&:=\max\bigl\{L_{\circ}^{6s(\square)},\;\kappa_Rc_W^3
\mathfrak{K}(\square)\bigr\},\\
\mathfrak{d}^{\sharp}(\square)&:=\kappa_Rc_W^3\,\mathfrak{K}^{\sharp}(\square)\big/
L_{\circ}^{6s(\square)},\\
C^{\mathrm{data}}&\ge\kappa_Rc_W^3\bigl[2(7L)^d(1+C^{\mathrm{data}}_{pr})+C^{tile}\bigr].
\end{aligned}
\end{equation}
The definitions of $\mathfrak{K}(\square)$ and $\mathsf{l}^2(\square)$ are the same at every
block. The quantity $\mathfrak{d}^{\sharp}(\square)$ is defined where $\tilde{\square}^5$ meets
$\mathcal{W}$.

\emph{Ceilings.} Put
\begin{equation*}
\mathsf{A}_w(\gamma):=\sup\{\mathsf{l}(\mathsf{x})\alpha_i(\mathsf{x}):\;
\mathsf{x}\ge\gamma^{-2}\},
\end{equation*}
where $\mathsf{l}(\mathsf{x})$ and $\alpha_i(\mathsf{x})$ are the functions of the coupling
variable $\mathsf{x}$, in the notation of Definition~\ref{5.1:not-coupling-functions}. By the
statement on the bounds for live arrests in which $\mathsf{A}_w$ is introduced,
$\mathsf{A}_w(\gamma)\downarrow0$ as $\gamma\downarrow0$. Let $c_h$ be half the least of the
absolute smallness constants
required in \cite[Proposition~2]{B-CMP98}, in \cite[Theorem~1, Proposition~9]{B-CMP102b} and in
the Baker--Campbell--Hausdorff bounds \cite[(3.30)--(3.32)]{B-CMP109}, each for the quantity it
bounds. Let $K_{min}$ be the window constant, the constant of the second and third conditions in
\eqref{5.5:eq-weight-constants-d}. The following one-sided conditions are imposed
with the ceiling $\gamma_J$, the last of them being a floor on $K_{min}$:
\begin{equation}\label{5.5:eq-weight-constants-d}
K_{loc}\,\mathsf{A}_w(\gamma)\le c_h,\qquad
K_{min}\vartheta\le\mathsf{w}_0,\quad K_{min}\ge28K^{+}.
\end{equation}
The second condition is the one read by the membership bound for own-read units.

Further ceilings are listed here, each one-sided:
\begin{itemize}
\item $17K^{+}\vartheta\le1$ and the smallness condition of the construction of
Definition~\ref{5.5:def-windowed-source};
\item the four ceilings of the remainder bounds, the last of which contains the threshold
condition of the refined remainder bound;
\item the three ceilings of the remainder theorem of the third class;
\item the ceiling of the $\circ$-slot $\delta$-size bound.
\end{itemize}

\emph{Floors.} The floors are:
\begin{itemize}
\item $L\ge 8$;
\item the floors under which the window construction of
Definition~\ref{5.5:def-windowed-source} is made, and the floors under which the data-factor
bound \eqref{5.5:eq-data-factor-a} is asked;
\item $\kappa\ge10\log L$ and $\min\{\delta_0M,\tfrac12\kappa\}\ge5\log L$, with $\delta_0$ the
constant of that name in the statement of the per-cell row table, under which that table is
obtained.
\end{itemize}
This definition imposes no floor beyond these.

In (iii) and (iv) below, the following per-cell row bounds at $w\equiv1$ are assumed, each taken
as stated elsewhere in this chapter:
\begin{itemize}
\item the rows of the per-cell row table for the core terms, the short singles and the
own-read units;
\item the collar row;
\item the block sum of the stored remainders;
\item the remainder row of the third class, from the remainder theorem of the third class and
its corollary;
\item the row of the near constituents of straddling cores, which is the row of the table of the
weight lemma.
\end{itemize}
With these choices the following hold.
\begin{enumerate}
\item[(i)] $K_{loc}\ge12C^{av}_M$.
\item[(ii)] $y_*<\infty$.
\item[(iii)] Under the assumptions above, $\kappa_R^{\times}$ is finite. Consider the size
vectors whose row sums, in the
normalisation used in the definition of the data functional $\mathsf{l}^J$ of the weight lemma,
are at most $1$. These lie in the set of size vectors admitted in the definition of
$\kappa_R^{\times}$. In particular the sizes divided by $\mathsf{l}^J$ in the transport step of
the weight lemma lie in that set.
\item[(iv)] Under the same assumptions, $\kappa_R$ is a number depending only on
$(d,L,M,\kappa,\kappa_1,\lambda,c_2,K^{+},K^{rem})$. It is uniform in $\square$, $\Delta'$, the
filing, the depth and the size vector.
\end{enumerate}
\end{definition}

\begin{proof}
(i) The radii $\alpha_{i,n}$ are positive, so $c_{\cdot}\ge0$. Together with $K_h\ge1$ and
$K_{ax}M\ge1$, the definition \eqref{5.5:eq-weight-constants-c} gives
\begin{equation*}
K_{loc}\ge12K_hK_{ax}MC^{av}_M\ge12C^{av}_M.
\end{equation*}

(ii) By the definition of $c_s$,
\begin{equation*}
c_{\log_Ly}=\log\bigl(1+3\lambda\log_Ly+2c_2\bigr).
\end{equation*}
Hence $L(2c_{\log_Ly})^{r_{pr}}/y\to0$ as $y\to\infty$. The inequality
$y\le L(2c_{\log_Ly})^{r_{pr}}$ therefore fails for all sufficiently large $y$. The set whose
supremum defines $y_*$ is thus bounded above, and $y_*<\infty$.

(iii) Fix $\square$, $\Delta'$, a filing, and non-negative sizes whose per-cell row-$r$ sums are
at most $K_r$ for every row $r$. Each piece and each stored remainder in the contribution of
$\Delta'$ is a Cauchy integral on its filed disc. Its bound is the prefactor of that disc, times
the numerical factor which the Cauchy-integral bound contributes, times the size. We write
$c_{cau}$ for this numerical factor of the lemma by which the Cauchy integrals are formed (a
convergent sum over its summation index), and we put
\begin{equation*}
K_{pre}:=\max\bigl\{K_A,\;24K'',\;K_A+28K^{rem}\bigr\}.
\end{equation*}

Consider first the pieces other than the near constituents of straddling cores. Since
$\eta_p\le1$, each of the three prefactors $K_A\vartheta(\square)$,
$24K''\vartheta(\square)\eta_p$ and $(K_A+28K^{rem})\vartheta(\square)$ is at most
$K_{pre}\vartheta(\square)$. The resummation over the anchors in
$\Delta'$ is carried out row by row with the scaffold coefficients. It therefore replaces the
sizes of row $r$ by their per-cell row-$r$ sum, which is at most $K_r$; there are finitely many
rows. These pieces therefore contribute at most
$K_{pre}\,c_{cau}\,\vartheta(\square)\sum_rK_r$.

Consider next the near constituents of straddling cores, at scale $j$ and anchored at level $n$,
with $\nu=1+n-j$ and $t=L^{j-n}$ as fixed for near units in this section. They are own-read units
of rows 1--3, formed on the own-read disc with prefactor $24K''\vartheta(\square)\eta_p$. By the
last of the row bounds assumed above, $\eta_p\le t^5$ at every $\nu$, and the number of their
anchors per unit $n$-cell is at most a constant times $(\nu+2)t^{-3}$. Since the sizes are
non-negative and their per-cell row-$r$ sum
is at most $K_r$, a level-$j$ size in row $r$ is at most $K_r/c^{(r)}_{n,j}$; since the
coefficient $c^{(r)}_{n,j}$ of the per-cell row table is at least a constant times
$t^{\hat{\alpha}}$, this is at most a constant times $t^{-\hat{\alpha}}$, where $\hat{\alpha}$ is
the H\"older exponent of the majorant of $\mathfrak{e}_\mu$. Let $c_{str}$ be the product of the
constants in these two bounds. The constituents at a given $\nu$ therefore contribute at most
$24K''c_{cau}\vartheta(\square)\,c_{str}(\nu+2)t^{2-\hat{\alpha}}$. Since $t=L^{1-\nu}$ with
$L\ge8$, and $2-\hat{\alpha}\ge1>0$ because $\hat{\alpha}\le1$ is a H\"older exponent,
\begin{equation*}
\sum_{\nu\ge1}(\nu+2)\,t^{2-\hat{\alpha}}
=\sum_{\nu\ge1}(\nu+2)\,L^{(1-\nu)(2-\hat{\alpha})}<\infty .
\end{equation*}
Moreover $24K''\le K_{pre}$.

Consequently the contribution of $\Delta'$ is at most $K_{pre}\,c_{cau}\,\vartheta(\square)$
times the sum of $\sum_rK_r$ and $c_{str}$ times the series just displayed. Comparing with the
bound $q\,c_\alpha C_Re^{C_2\kappa_1}\vartheta(\square)$ required in the definition of
$\kappa_R^{\times}$, the factor $\vartheta(\square)$ cancels, and
\begin{equation*}
\kappa_R^{\times}\le\frac{K_{pre}\,c_{cau}}{c_\alpha C_Re^{C_2\kappa_1}}
\Bigl(\sum_rK_r+c_{str}\sum_{\nu\ge1}(\nu+2)\,L^{(1-\nu)(2-\hat{\alpha})}\Bigr).
\end{equation*}
The right member is a number free of $\square$, $\Delta'$, the filing and the sizes. Hence
$\kappa_R^{\times}$ is finite, and it is size-free.

In the normalisation used in the definition of $\mathsf{l}^J$, a normalised row sum at most $1$
of row $r$ is a per-cell row-$r$ sum at most $K_r$. So the set of size vectors with normalised
row sums at most $1$ lies inside the set admitted in the definition of $\kappa_R^{\times}$. In
particular, the sizes which the transport step of the weight lemma divides by $\mathsf{l}^J$
lie in that set.

(iv) The three prefactors of the filed discs each carry the factor $\vartheta(\square)$ exactly
once. So do $\mathfrak{N}_0(\square)=c_\alpha E_0C_Re^{C_2\kappa_1}\vartheta(\square)$ and the
right member $c_\alpha C_Re^{C_2\kappa_1}\vartheta(\square)$ in the definition of
$\kappa_R^{\times}$. In the quotients defining $\kappa_R^{act}$ and $\kappa_R^{\times}$ the
factor $\vartheta(\square)$ therefore cancels.

What remains are the following, and none of them depends on $\square$, $\Delta'$, the filing or
the depth:
\begin{itemize}
\item the constants of the discs, including $24=1/\mathsf{w}_0$, and the bound $\eta_p\le1$;
\item the numerical factor $c_{cau}$ of the Cauchy-integral lemma;
\item the scaffold coefficients;
\item the per-cell row bounds at $w\equiv1$.
\end{itemize}
In the case of $\kappa_R^{\times}$ the supremum over the admitted sizes is included, by (iii).

The per-cell row bounds at $w\equiv1$ assumed above are taken as stated; only the last of them
needs comment here.

For the last of them, consider a constituent at scale $j$ anchored at level $n$, with
$\nu=1+n-j$ and $t=L^{j-n}$ as fixed for near units in this section, and let $X$ be its
localisation domain. Then $\eta_p\le t^5$, and the constituent has size at most
$2\mathfrak{n}_ve^{-\kappa d_j(X)}\le2\mathfrak{n}_v$, bounded at $w\equiv1$. The
number of anchors of such constituents per unit
$n$-cell is at most a constant times $(\nu+2)t^{-3}$. The per-cell sum is therefore at most a
constant times $t^5\cdot(\nu+2)t^{-3}=(\nu+2)t^2$. Since $t=L^{1-\nu}$ and $L\ge8$, this is
summable in $\nu$. For the
sizes admitted in the definition of $\kappa_R^{\times}$ the corresponding bound is the one
obtained in (iii), with $(\nu+2)t^{2-\hat{\alpha}}$ in place of $(\nu+2)t^2$.

Hence $\kappa_R$, as the maximum in \eqref{5.5:eq-weight-constants-a} of the two numbers, is a
number depending only on $(d,L,M,\kappa,\kappa_1,\lambda,c_2,K^{+},K^{rem})$. It is uniform in
$\square$, $\Delta'$, the filing, the depth and the size vector.
\end{proof}

\begin{definition}[Membership and remainder inputs of the weight lemma]\label{5.5:def-weight-inputs}
We work in the setting of Definition~\ref{5.5:def-windowed-source}. Let $\square$ be a block, with
decaying normaliser $\vartheta(\square)$, box radius $r^{\sharp}_{\square}$ and own-level radius
$r_T(\square)$ as in \eqref{5.5:eq-windowed-source-a}. Near units are box-read, own-read or of the
third class as in Definition~\ref{5.5:def-near-unit-classes}, and the Cauchy integrals below run on
the discs of \eqref{5.5:eq-near-unit-classes-b}:
\begin{equation*}
\mathbb{D}_{\square} = \{|\mathsf{t}|\le r^{\sharp}_{\square}\},\qquad
\rho_A(v) = \mathsf{w}_0\, r_T(\square),\qquad
\{|\mathsf{t}|\le r^{rem}_{\square}\}.
\end{equation*}
The constants are those of Definition~\ref{5.5:def-weight-constants}. Throughout, $t\in[0,1]$ and
$|\sigma|\le e^{\kappa_1}$. Two machines $\mathsf{A}$ and $\mathsf{B}$ are run on one scaffold. For a
quantity $F$ formed in both machines, $\delta F$ denotes the difference of its two values. The
letter $\delta$ in exponents such as $e^{-\delta D_v/8}$ is the margin parameter of
Definition~\ref{5.2:def-step-constants}. For a near unit $v$,
$\mathfrak{n}_v$ is its size in the table of near-unit sizes. For a box-read unit $v$ we put
$D_v := d_{\star}(\mathfrak{r}_v, F_{\square})$. This is the geometric distance from the own box
$\mathfrak{r}_v$ of $v$ to the support region $F_{\square}$ of the stored-remainder phase
$\mathsf{E}$.

The \emph{inputs of the weight lemma} are the five statements (i)--(v) below. The weight lemma
assumes each of them as a hypothesis.

\smallskip
\noindent\textup{(i)} \emph{Membership for own-read units.} Let $\square$ be a block and $v$ an
own-read near unit with term $T$ and domain $X_T$. Let
$\mathcal{W}_T := \mathcal{W}^{\ell_T,\square}\lceil X_T$ be its own-level family, where $\ell_T$ is
the own level of the term $T$. Then the following hold.
\begin{itemize}
\item[(a)] For $t\in[0,1]$, $|\tau|\le\rho_A(v)$ and $|\sigma|\le e^{\kappa_1}$ we have
$\mathcal{W}_T(t,\tau)\lceil X_T\in\mathfrak{U}_T=\mathfrak{D}_v$, holomorphically in the
complex parameters. On $|\tau|\le\rho_A(v)$ this follows from the analyticity statement for the
single term $T$ together with the window ceiling $K_{\min}\vartheta\le\mathsf{w}_0$ of
\eqref{5.5:eq-weight-constants-d}. The same containment holds on the disc
$|\tau|\le r_T(\square)$, which contains $|\tau|\le\rho_A(v)$.
\item[(b)] The jet in $\tau$ of $\mathcal{W}_T$ at the real point is the piece of $T$ that the
weight lemma assigns to $v$.
\item[(c)] Let $F$ be holomorphic on $\mathfrak{U}_T$ with $|F|\le\mathfrak{n}$ there. Then
$|F\circ\mathcal{W}_T|\le\mathfrak{n}$. Consequently the hypothesis of
Lemma~\ref{5.2:lem-disc-bounds} holds with $\mathfrak{N}(v):=2\mathfrak{n}_v$ on the disc
$|\tau|\le\rho_A(v)$. The coefficients $v_p$ of the piece of $v$ at the plaquettes $p$ satisfy
\begin{equation}\label{5.5:eq-weight-inputs-1}
|v_p|\le 2K''\,\vartheta^{\natural}(\square)\,\eta_p\,e^{-(\kappa_1-1)m(p)}\,\mathfrak{n}_v ,
\end{equation}
where $m(p)$ is the plaquette exponent of the own-read row of the near-unit sizes.
\item[(d)] The own-read pieces remain summands of the single near element of the block.
\item[(e)] For the two machines,
\begin{equation*}
|\delta(F\circ\mathcal{W}_T)|\le\sup_{\mathcal{K}}|\delta F| .
\end{equation*}
Here $\mathcal{K}$ is the pair system of (iii). This bound follows from the closure condition of
(iii) at the move $\mathsf{M}3$.
\end{itemize}
Box-read units need no membership input. Their membership on the chart $|\sigma'|<S^{\star}_v$,
where $\sigma'$ denotes the chart parameter,
comes from the box-pair clause of the weight lemma and from the chart clause of the transfer of the
analysis of \cite[\S 3]{B-CMP109} to the cover; that transfer rests on the sentence of
\cite[p.~276]{B-CMP119}, printed there without proof and cited here as printed. This chart contains
$\sigma'=1$, because $S^{\star}_v>1$ for box-read units (Definition~\ref{5.5:def-near-unit-classes}).
Pieces of the third class need no membership input, since they are entire.

\smallskip
\noindent\textup{(ii)} \emph{Remainders of box-read units.} Let $\square$ be a block and $v$ a
box-read near unit. Let $T$ be a constituent of $v$ with domain $X\subset\mathfrak{r}_v$, and take
any point of the source. Let $t\in[0,1]$, $|\mathsf{t}|\le r^{rem}_{\square}$ and
$|\sigma|\le e^{\kappa_1}$. Put
\begin{equation}\label{5.5:eq-weight-inputs-2}
G_v := T(X;\mathsf{P}^1)-T(X;\mathsf{P}^0),\qquad
\operatorname{rem}(T,\square) := \int_0^1 dt\,\partial_{\mathsf{t}}\big|_{\mathsf{t}=0}G_v .
\end{equation}
This is the stored remainder. At real parameters its jet is the sum of the expressions
\cite[(3.21), (3.22), p.~274]{B-CMP109}. Then the following hold.
\begin{itemize}
\item[(a)] The pairs $\mathsf{P}^0$, $\mathsf{P}^1$ and the segment $\mathsf{C}(\sigma_1)$ of box
pairs joining them, for $|\sigma_1|\le e^{\delta D_v/8}$, lie in $\mathfrak{U}_T$. They lie in the
own chart of $v$ at scale $j$, written $(U_j,J_j)(\mathfrak{r}_v, e^{i(B^0+\sigma_1\Delta)})$, with
\begin{equation*}
|B^0|\le|B^0|_* := 1.42\,L^{-1}\sup_b|B^{ax}|_*,\qquad
|\Delta|\le|\Delta|_* := \tfrac12\,\alpha_{3,j}\,e^{-\delta D_v/8}.
\end{equation*}
\item[(b)] These objects are holomorphic in the complex parameters.
\item[(c)] The size bounds read
\begin{equation}\label{5.5:eq-weight-inputs-3}
\begin{aligned}
|G_v| &\le 2\,\mathfrak{n}_v\,e^{-\delta D_v/8},\\
|\operatorname{rem}(T,\square)|
&\le \big[2/r^{\sharp}_{\square}+28K^{rem}\vartheta(\square)\big]\,
\mathfrak{n}_v\,e^{-\delta D_v/8}\\
&= (K_A+28K^{rem})\,\vartheta(\square)\,\mathfrak{n}_v\,e^{-\delta D_v/8}.
\end{aligned}
\end{equation}
The equality is the identity $2/r^{\sharp}_{\square}=K_A\vartheta(\square)$ of
\eqref{5.5:eq-windowed-source-a}.
\item[(d)] The block sum satisfies
\begin{equation}\label{5.5:eq-weight-inputs-4}
\sum\big\{\mathfrak{n}_v\,e^{-\delta D_v/8}:\ v \text{ box-read at } \square\big\}
\le c_{\Sigma}\,\mathfrak{s}^{rd}.
\end{equation}
Here $c_{\Sigma}$ is the block-sum constant of this remainder bound and $\mathfrak{s}^{rd}$ is the
size total of the read rows. The same bound holds for the sum over the box-read units of one unit
cell. The right side contains no weight and no age.
\item[(e)] For the two machines,
\begin{equation*}
|\delta G_v|\le 2\,\mathsf{S}\,e^{-\delta D_v/8}.
\end{equation*}
Here $\mathsf{S}$ is the $\delta$-size fixed in (iii). The closure conditions this bound requires
are those of (iii) at the move $\mathsf{M}4$, which reads $\mathsf{P}^1$, and at the fourth
splitting of $\mathsf{M}5$, the splitting $(B^0,\Delta)$ of (a).
\end{itemize}
Own-read units carry no remainder.

\smallskip
\noindent\textup{(iii)} \emph{Closure of the pair system.} With the following notation,
hypothesis (iii) is the statement under \emph{Closure conditions} and \emph{Construction} below.

\emph{Slots.} Slots carry one of two letters. The letter $\flat$ marks the rows one to three and the
boundary row. Their tubes are the $\mathfrak{U}_T$, and
$\mathfrak{U}(j,X;\flat):=\bigcap_T\mathfrak{U}_T$. The letter $\circ$ marks the Wilson and
counterterm pieces $\hat{A}$. These are entire and bounded by $\mathfrak{n}_{\hat{A}}$, the largest
of their size bounds over the moves landing in a $\circ$-slot.

\emph{The pair system.} The pair system is
\begin{equation*}
\mathcal{K}:=\mathcal{K}^{rd}:=\mathrm{cl}^{pr}(\mathfrak{M}^{rd})\subset\mathrm{cl}^{pr}_J .
\end{equation*}
It is the production closure of the seeds of the rim datum (Definition~\ref{5.2:def-rim-datum})
under the single move list $\mathfrak{M}^{rd}$, and it is contained in the production closure
$\mathrm{cl}^{pr}_J$ used for the correlation theorem. That list consists of the following moves.
\begin{itemize}
\item[] $\mathsf{M}1^M$: the response dilation $\Phi_t$ of
Corollary~\ref{5.2:cor-response-dilation}, taken at the radius of the parameter domain prescribed
for this move, with $\vartheta$ replaced by $\vartheta^{\natural}$.
\item[] $\mathsf{M}2$: the box read
$(\mathfrak{b}'^{\mathsf{B}}_j,\mathfrak{b}'^{\mathsf{A}}_j)\lceil X$ on
$(b,\sigma)\times[0,1]\times\mathbb{D}_{\square}$. It applies at box-read units (slot $\flat$) and
at every near Wilson or counterterm piece (slot $\circ$), with parent domain
$\hat{X}\supseteq\hat{\mathsf{K}}_{\square}$.
\item[] $\mathsf{M}3$: the own-level read $\mathcal{W}^{\times}_T$ on $\rho_A(v)$.
\item[] $\mathsf{M}4$: the remainder read $\mathsf{P}^{1\times}$ on $r^{rem}_{\square}$.
\item[] $\mathsf{M}5$: the paired charts through $\mathsf{M}2$ and $\mathsf{M}4$, at the weighted
majorants of Definition~\ref{5.5:def-windowed-source}. It uses four splittings: by $B^{ax}$; into
the $F$-part and $\mathfrak{e}$; by $B_0$; and into $(B^0,\Delta_c)$ at the radii
$1.42\,L^{-1}\sup_b|B^{ax}|_*$ and $\tfrac12\alpha_{3,j}e^{-\delta D_v/8}$. It also uses the
families from $\mathbb{1}$, including $\mathsf{h}[B]$ at $\|\mathsf{h}[B]\|_*$.
\item[] $\mathsf{M}6$: restriction.
\item[] $\mathsf{M}7$: the reads of the rim datum at the sites (T) and (L).
\item[] $\mathsf{M}8$: the comb-dilated remainder read. For every cell-unit $v=(c,i,\Delta')$ of the
third class, $\Delta'$ a cell, it takes a stage to $(i,\Delta';\circ)$ through the $\circ$-forms of
$\mathsf{M}2$ and $\mathsf{M}4$. It runs on the parameter domain
\begin{equation*}
\mathfrak{P}_{\perp}:=(b,\sigma)\times[0,1]\times\{|\mathsf{t}|\le r^{rem}_{\square}\}
\times\{|s|\le R^{\perp}_v\},\qquad R^{\perp}_v:=(32\kappa^*_v)^{-1},
\end{equation*}
where $\kappa^*_v$ is the constant of $v$ governing the comb dilation,
with the same construction in $\mathsf{A}$ and in $\mathsf{B}$. Its $\circ$-slot $\delta$-size is
$\mathsf{S}^{\perp}_v:=\mathsf{S}^{\circ}[A(c,\cdot)]$, where $A(c,\cdot)$ is the Wilson or
counterterm piece of the cell $c$ and $\mathsf{S}^{\circ}[\cdot]$ is the $\delta$-size in a
$\circ$-slot.
\end{itemize}

\emph{Sizes.} We put $\mathsf{S}:=\mathsf{S}^{rd}$ and $\mathsf{a}:=\mathsf{S}^{rd}[\hat{A}]$.

\emph{Closure conditions.} On $\mathcal{K}^{rd}$, by
construction, the closure conditions of $\mathcal{K}$ at the moves of $\mathfrak{M}^{rd}$ hold.
They yield the following.
\begin{itemize}
\item[] The closure at $\mathsf{M}6$ gives the stability of $\mathcal{K}$ under restriction.
\item[] The closure at $\mathsf{M}1^M$ gives the stability under the response dilation.
\item[] The closures at $\mathsf{M}2$, $\mathsf{M}3$ and $\mathsf{M}4$ give the closure under the
three reads.
\item[] The closure at $\mathsf{M}3$ gives (i)(e).
\item[] The closure at $\mathsf{M}5$ gives the closure under the paired charts.
\item[] The closures at $\mathsf{M}4$ and at the fourth splitting of $\mathsf{M}5$ give those used
in (ii)(e).
\item[] The closure at $\mathsf{M}8$ gives the one used in (v).
\end{itemize}

\emph{Construction.} There is one parameter domain $\mathfrak{P}$ per move. The construction is the
same in each machine and is jointly holomorphic; at $\mathsf{M}5$ it is jointly holomorphic in the
$(\sigma_i)$. Every element of $\mathcal{K}$ arises from the seeds by moves applied in production
order. Chains of moves stay in the product of the one-machine classes. Moreover
\begin{equation*}
\mathsf{S}^{rd}_j[F]\le 2\mathfrak{n}\ \text{ whenever } |F|\le\mathfrak{n},\qquad
\mathsf{a}\le 2\mathfrak{n}_{\hat{A}} .
\end{equation*}
The closure statement further comprises the product principle for the two machines, a dictionary
between the constructions in the two machines, and the list of the places at which the two-machine
clause of the weight lemma appeals to $\mathcal{K}$, the place at $\mathsf{M}8$ among them.
Every closure condition used in the two-machine clause of the weight lemma holds on $\mathcal{K}$;
this includes the closure at $\mathsf{M}8$.

\smallskip
\noindent\textup{(iv)} \emph{Remainders of the third class.} The remainders of the Wilson and
counterterm pieces of the third class obey a size bound and a block-sum bound of the form of (ii)(c)
and (ii)(d). The rate is $e^{-\delta D_v/4}$ in place of $e^{-\delta D_v/8}$, and the size bounds
carry the factor $\mathsf{s}_v\mathsf{v}_v$ of the third-class remainder sizes of $v$. These
remainders are entire. The sums are taken per unit $n$-cell, grouped as in the definition of
$\kappa_R$ in Definition~\ref{5.5:def-weight-constants}. A
block piece is assigned, plaquette by plaquette, to the level of its cell. Hypothesis (iv) also
comprises, on the chart of $v$, the two-machine clause of that statement, of the form of (ii)(e);
off the chart see (v).

\smallskip
\noindent\textup{(v)} \emph{Two-machine bound of the third class off the chart.} Let $v$ be any
unit of the third class. This covers the block cell-pieces and the companions
$b_j(\hat{w})A(h_y,\cdot)$ of cores, that is, the Wilson and counterterm pieces attached to the core
anchored at $y$, at every $\nu$ and every value of $S^{\star}_v$. Off the chart,
the two-machine differences of its remainders obey the bounds of (iv), with
$\mathsf{s}_v\mathsf{v}_v$ replaced by $C^{\perp}\mathsf{S}^{\perp}_v$, where
\begin{equation*}
C^{\perp}:=720\,L^2 .
\end{equation*}
This bound is used together with the closure of (iii) at $\mathsf{M}8$. It is taken under its own
one-sided smallness condition on $\gamma$, whose constant is fixed before $\gamma$ is chosen. For
the companions, (v) covers their remainder. It does not cover the chart factor of their item at
$\nu\le 2$.

\smallskip
Besides (i)--(v), the weight lemma uses, by statement, two further results. The first is the
per-cell row sums at weight identically one. The second is the analyticity statement for a single
term, which enters through (i)(a) and through the closure input (iii).

The following are not hypotheses of the weight lemma. The first is the closeness, in decaying
weights, of the responses at two cutoffs. This consists of the two-cutoff response statement at
zeroth order together with the energy-difference bound, at the dilations of radius
$e^{\delta D_v}$ of the moves of the production closure. These dilations are those of $\mathsf{M}8$
and of the fourth splitting of $\mathsf{M}5$ read through $\mathsf{M}4$. That closeness enters the
uniqueness theorem, not the weight lemma. The second is the low-order variation bound at the second
cube size, which would supply the chart factor at $\nu\le 2$ excluded in (v).

The two-machine data functional $\mathsf{l}^J$ used in the weight lemma is the one fixed with its
supplier rows, and on the production closure it is bounded by those rows, the row for the paired
read on $\mathrm{cl}^{pr}_J$ among them.
These rows are taken by statement and are not among the hypotheses (i)--(v). Finally,
the weight lemma is stated under the following conditions of the setting, which are not inputs of
the kind (i)--(v):
\begin{itemize}
\item[] the classification of Definition~\ref{5.5:def-near-unit-classes} and the prescriptions of
radii used at $\mathsf{M}1^M$ and at the fourth splitting of $\mathsf{M}5$;
\item[] the use, at the datum $Q_v$ attached to $v$, of the two bounds
\cite[(3.27), (3.32)]{B-CMP109};
\item[] the two ceilings on $\gamma$ displayed in \eqref{5.5:eq-weight-constants-d};
\item[] the smallness condition attached to (v);
\item[] a further listed condition on the weight sum $\mathfrak{K}(\square)$;
\item[] the tile bound at the blocks $\square$ for which $\tilde{\square}^5$ meets no cell of
$\mathcal{W}$;
\item[] the supplier rows of $\mathsf{l}^J$.
\end{itemize}
\end{definition}

\begin{lemma}[The weight lemma]\label{5.5:lem-weight}
Work in the setting of Definition~\ref{5.5:def-windowed-source}, with the letters, levels and
weights of Notation~\ref{5.5:not-weight-notation}, the near units, their three classes and
their filed discs as in Definition~\ref{5.5:def-near-unit-classes}, the data factor of
Definition~\ref{5.5:def-data-factor}, and the constants of
Definition~\ref{5.5:def-weight-constants}. Assume the floors and the one-sided
ceilings on $\gamma$ listed in Definition~\ref{5.5:def-weight-constants}, among them
\eqref{5.5:eq-weight-constants-d}. Assume moreover the following statements:
\begin{itemize}
\item the membership statement for own-read near units (Definition~\ref{5.5:def-weight-inputs});
\item the remainder statement for box-read near units (Definition~\ref{5.5:def-weight-inputs});
\item the closure statements for the pair system $\mathcal{K}^{rd}$
(Definition~\ref{5.5:def-weight-inputs}): its closure under
the read moves $\mathsf{M}2$, $\mathsf{M}3$, $\mathsf{M}4$, under the paired-chart move
$\mathsf{M}5$, under the restriction move $\mathsf{M}6$ and under the response-dilation move
$\mathsf{M}1$, and its closure under the comb-dilated remainder move $\mathsf{M}8$;
\item the remainder statement for the third class (size, block sum and filing of the
remainders of the near Wilson and counterterm pieces);
\item the two-run statement for the third class off the chart, namely the bound on the
two-machine differences of the third-class pieces, in the format of the size and block-sum
bounds of the remainder statement for the third class with $\mathsf{s}_v\mathsf{v}_v$ replaced
by $C^{\perp}\mathsf{S}^{\perp}_v$, at every third-class unit, the letters being those of
part~(e);
\item the per-cell row bounds of the near-unit table at unit weight, and the statement that
the own-level family of each own-read unit takes values in the tube $\mathfrak{U}_T$ of that
unit, on which the membership statement rests.
\end{itemize}
The two-run statement for the third class and the closure under $\mathsf{M}8$ enter only
part~(e), at the third class off the chart, that is, at the pieces at which the weighted
scaffold majorants of Definition~\ref{5.5:def-windowed-source}, on which the paired charts
rest, are not available. Then at every stage of every $\mathbf{R}_k$,
for every block $\square$, every $t\in[0,1]$ and every $\sigma$ with
$|\sigma|\le e^{\kappa_1}$ --- with $\mathsf{t}\in\mathbb{D}_{\square}$ in~(a) and in the rows
read on the box pair, $|\tau|\le\rho_A(v)$ in the own-read rows, and
$|\mathsf{t}|\le r^{rem}_{\square}$ in the remainder rows --- the following hold.
\begin{enumerate}
\item[(a)] \textup{(a1)} $|\partial\mathbb{W}(x)-1|<1.4\,\theta_2(x)$ for every cell $x$
meeting $\tilde{\square}^3$.

\textup{(a2)} For every cell $x$ meeting $\tilde{\square}^3$ and every $j\le k$,
\begin{equation}\label{5.5:eq-weight-a}
|\partial\mathfrak{b}'_j(x)-1|\le 4\,\theta_2(x).
\end{equation}

\textup{(a3)} Let $Q$ be a box of at most $9$ cubes of $\pi_p$ a side; write $Q^{(1)}$ for the
enlargement of $Q$ by one layer of $\pi_p$-cubes and $Q^{\sim-1}$, $Q^{\sim-2}$ for $Q$ shrunk
by one and by two such layers. Assume
$Q^{(1)}\subset\tilde{\square}^3$ and every cell meeting $Q^{(1)}$ of level at least $p$; let
$w_Q$ and $\alpha_Q$ be the largest values of $w(\Delta)$ and of $\alpha_{0,\ell(\Delta)}$
over those cells. In one gauge of the complexified gauge group $G^{\mathbb{C}}$ on
$Q^{\sim-1}$ one has $\mathfrak{b}=e^{i\eta A}$ with a potential $A$ satisfying, for
$0\le\beta\le\beta_0$ ($\beta_0$ the H\"older ceiling, $\|\cdot\|_{\beta}$ the H\"older
seminorm of exponent $\beta$, $\ell_p=L^p\eta$ the length of level $p$),
\begin{equation}\label{5.5:eq-weight-b}
\max\bigl\{\ell_p|A|,\ \ell_p^2|\nabla A|,\ \ell_p^{2+\beta}\|\nabla A\|_{\beta}\bigr\}
\le\mathsf{a}_Q:=K_{loc}\,w_Q\,\alpha_Q .
\end{equation}

\textup{(a4)} Let $j\le p$ and let $\square_0'\subset Q^{\sim-2}$ be a box of
$\pi_j$-cubes. Then $\mathfrak{b}'_j$ restricted to $\square_0'$ equals
$(U_j,J_j)(\square_0',\mathsf{V}')$ modulo gauge, where $\mathsf{V}'$ is
$M^{\cdot}(\mathfrak{b}'_j)$ restricted to $\mathfrak{B}_j(\square_0')$. Put $t_p:=L^{j-p}$.
In the comb gauge centred at any $y\in\square_0'$, one has
$|\partial\mathsf{V}'-1|\le c_1\mathsf{a}_Qt_p^2$ and, on all bonds $b$,
$|B^{ax}(b)|\le 2d(1+|b-y|)\,c_1\mathsf{a}_Qt_p^2$; moreover
\begin{equation}\label{5.5:eq-weight-c}
\begin{gathered}
B^{ax}_{\mu}(x)=\sum_{\iota<\mu}(x-y)_{\iota}F_{\iota\mu}(y)+\mathfrak{e}_{\mu}(x),
\qquad |F|\le c_1\mathsf{a}_Qt_p^2,\\
|\mathfrak{e}_{\mu}(x)|\le c_1(1+|x-y|)^2\mathsf{a}_Qt_p^{2+\beta}
\quad\text{on the level-}j\text{ bonds},\\
|\mathfrak{e}_{\mu}(b)|\le 4d(1+|b-y|)\,c_1\mathsf{a}_Qt_p^2
\quad\text{on the graded-rim bonds (levels }i<j).
\end{gathered}
\end{equation}

\item[(b)] Let $v$ be a near unit with anchor $y$, and let $n$, $j$, $\nu$ and the scale ratio
$t=L^{j-n}$ be its letters as in Definition~\ref{5.5:def-windowed-source} (see also
Notation~\ref{5.5:not-weight-notation} for the homonym $t$); this $t$ is the scale
ratio, not the interpolation parameter $t\in[0,1]$ of the preamble. Put
$\Theta:=c_Ww^{\vee}(y)$ and $z:=\Theta\rho_{n,j}$. On the filed discs the
near units satisfy the following bounds, which form the table $\mathfrak{N}^P(v;\Theta)$.
\begin{itemize}
\item A box-read core, taken whole together with its counterterm piece on the same box
pair: its term $\mathfrak{i}_y$ at the anchor $y$ satisfies
\begin{equation*}
|\mathfrak{i}_y|\le C_I'\,\mathfrak{n}\,\nu^3z^2(1+z)\,t^{4+\hat{\alpha}}.
\end{equation*}
\item A box-read short single: at most $2\mathfrak{n}_v(S^{\star}_v)^{-2}$.
\item An own-read unit: at most $2\mathfrak{n}_v$ each, by the membership statement.
\item The near constituents ($X\ni y$, $X\subset\tilde{\square}^2$) of a straddling core,
that is, of a core with $y\in\tilde{\square}^2$ and
$\mathfrak{r}_v=\square_y^{\sim\nu+1}\not\subset\tilde{\square}^2$, where $\square_y\in\pi_j$
is the cube of $y$, so that $v$ itself is not a near unit; this holds at every $\nu$. They are
own-read units on $\mathcal{W}_T$, each at most $2\mathfrak{n}_ve^{-\kappa d_j(X)}$ at the
prefactor $24K''\vartheta(\square)\eta_p$, with $\eta_p\le t^5$. For $\nu\ge2$ one has
$X\cap\tilde{\square}=\emptyset$, because $2\nu+3<L^{\nu-1}$ for $\nu\ge2$ (at $\nu=2$:
$7<8\le L$), and then $\eta_p\le e^{-\frac14\delta wL^{n_0-j}}\le t^5$. The first inequality,
which by the definition of $\eta_p$ amounts to $d_{\ell}(X_T,\square^{\zeta})\ge wL^{n_0-j}$,
follows from $X\cap\tilde{\square}=\emptyset$ by the distance bound for the own-level family;
the second follows from the level floor among the floors of
Definition~\ref{5.5:def-weight-constants}. Here $w=M/M_1$ is the side of the own-scale cube in
units of $M_1$ (Notation~\ref{5.5:not-weight-notation}), not a weight. At
$\nu=1$ one has $n=j$, so $t=1$, and $\eta_p\le t^5$ holds because $\eta_p\le1$; nothing is
claimed of $\mathfrak{r}_v\cap\tilde{\square}$.
For two machines each is at most $2\mathsf{S}e^{-\kappa d_j(X)}$. These bounds have
degree $0$ in $\Theta$. Their per-cell sums --- $t^5$ times the number of anchors of such
constituents per unit $n$-cell, a number proportional to $(\nu+2)t^{-3}$ --- are summable in
$\nu$ inside $\kappa_R$, by the per-cell row bounds at unit weight, and inside
$\kappa_R^{\times}$ as well: in the set where the normalised row sums are at most one, a
level-$j$ size is at most $K_r/c^{(r)}_{n,j}$, with the row constants $K_r$ and
$c^{(r)}_{n,j}$ of Definition~\ref{5.5:def-weight-constants}; this quotient is proportional
to $t^{-\hat{\alpha}}$, and $(\nu+2)t^{2-\hat{\alpha}}$ remains summable in $\nu$.
\item An own-read unit which is a read core ($S^{\star}_v\le1$, so that by
\eqref{5.5:eq-near-unit-classes-a} $C_MDt^2z\ge1>\frac12$): its constituents, taken singly,
together with its counterterm companion of the third class, are bounded inside
$C_I'\,\mathfrak{n}\,\nu^3z^2(1+z)t^{4+\hat{\alpha}}$.
\item At $\nu\le2$, where the units are own-read by the filing at every value of
$S^{\star}_v$: the constituents ($2\mathfrak{n}_v$ each) and the counterterm companion (at
most $C'\mathfrak{n}z^2t^4$, with the constant $C'$ of
Definition~\ref{5.5:def-weight-constants}) are separate units of the table.
\item The Wilson piece of the block (third class): at most $\Theta^2\mathfrak{N}_W(v)$.
\item The remainders of the box-pair-read pieces (slot $\flat$): at most
$(K_A+28K^{rem})\vartheta(\square)\mathfrak{n}_ve^{-\frac18\delta D_v}$, free of weights,
with block sum by the remainder statement for box-read near units; the remainders of the
third class by the remainder statement for the third class.
\end{itemize}
Every entry of the table $\mathfrak{N}^P(v;\Theta)$ is non-decreasing and of degree at most
$3$ in $\Theta$, so that
\begin{equation}\label{5.5:eq-weight-d}
\mathfrak{N}^P\bigl(v;c_Ww^{\vee}(y)\bigr)\le c_W^3\,w^{\vee}(y)^3\,\mathfrak{N}^P(v;1).
\end{equation}
The far pieces read no datum; the pieces of the fourth kind in the total of
Definition~\ref{5.5:def-data-factor} sit at blocks meeting
$R_{\mathsf{n}}$, which by~(c) read no cell of $\mathcal{W}$.

\item[(c)] If $\Delta\in\mathcal{W}$, $y\in\tilde{\square}^2$ and
$\operatorname{dist}_n(\Delta,y)<1$, with $n=\min\{\ell(y),n_0\}$ as in
Definition~\ref{5.5:def-windowed-source}, then, writing $t(\Delta)$ for the step of the arrest
that raised $\Delta$ (Notation~\ref{5.5:not-weight-notation}),
\begin{equation}\label{5.5:eq-weight-e}
t(\Delta)\ge\ell(y),\qquad t(\Delta)\ge n_0-1,\qquad \square\in\mathfrak{O}^+ .
\end{equation}
If $\tilde{\square}^2$ holds a point of level at most $n_0-2$, then
$\tilde{\square}^2\subset X_{\mathfrak{c}}$ crosses the thin strata of the sequence of block
partitions $\mathbb{B}$ of exactly one
preparation $\mathfrak{c}$ ($X_{\mathfrak{c}}$ the region of $\mathfrak{c}$), of step
$t_{\mathfrak{c}}\ge n_0$ ($t_{\mathfrak{c}}:=k_{\mathfrak{c}}$
for an arrest, $t_{\mathfrak{c}}:=k$ for the running or a sibling preparation of this
$\mathbf{R}_k$), with $k_0(\mathfrak{c})\le n_0-2$ for the level $k_0(\mathfrak{c})$ attached
to $\mathfrak{c}$; every cell inside $\tilde{\square}^2$ has
level in $[k_0(\mathfrak{c}),n_0+1]$, and every weighted cell within one own unit of
$\tilde{\square}^2$ belongs to $\mathfrak{c}$. No window block, the exceptional ones
included, reads a cell of $\mathcal{W}$ through $\mathfrak{R}(\square)$; the set
$\tilde{\square}^5$ of an exceptional block may nevertheless meet $\mathcal{W}$, and this case
is treated in~(c$'$).

\item[(c$'$)] \textup{(Count times weight.)} Let $\tilde{\square}^5$ meet $\mathcal{W}$, so
that $\square\in\mathfrak{O}^+$ or $\square$ is exceptional, and $\varrho=\varrho(\square)\ge
\check{R}_k-4$; let $\mathfrak{K}(\square)$ be the sum of
\eqref{5.5:eq-weight-constants-b}. In case \textup{(i)}, where $\tilde{\square}^2$ holds no
point of level at most $n_0-2$, the first line of \eqref{5.5:eq-weight-f} holds; in case
\textup{(ii)}, where it holds such a point, with $t:=t_{\mathfrak{c}}\ge n_0$ as
in~(c) ($t=k$ at an exceptional block), the second line holds (recall
$L_{\circ}=\mathsf{L}_0$, Notation~\ref{5.5:not-weight-notation}):
\begin{equation}\label{5.5:eq-weight-f}
\begin{aligned}
&\text{(i)}\quad \mathfrak{K}(\square)\le 2(7L)^d
\max\bigl\{L_{\circ}^{6s(\square)},\,C^{\mathtt{a}}_{\mathrm{pr}}e^{\varrho-2}\bigr\},\\
&\text{(ii)}\quad \mathfrak{K}(\square)\le 3\cdot7^d\mathsf{L}_0^6L^{dN_0(t)}
\le C^{tile}e^{\varrho-1}.
\end{aligned}
\end{equation}
The majorant $\mathfrak{K}^{\sharp}(\square)$ denotes the right-most member of the applicable
case.

\item[(d)] At every block, the near-unit total at $\square$ (the total whose excess over
$L_{\circ}^{6s(\square)}\mathfrak{N}_0(\square)$ defines $\mathfrak{d}(\square)$ in
Definition~\ref{5.5:def-data-factor}) is at most
$\kappa_Rc_W^3\mathfrak{K}(\square)\mathfrak{N}_0(\square)\le
\mathsf{l}^2(\square)\mathfrak{N}_0(\square)$. Hence, where $\mathfrak{d}$ is not $1$ by
definition,
\begin{equation}\label{5.5:eq-weight-g}
\begin{gathered}
L_{\circ}^{6s(\square)}\mathfrak{d}(\square)\le\mathsf{l}^2(\square)
\le L_{\circ}^{6s(\square)}\mathfrak{d}^{\sharp}(\square),\quad\text{that is,}\\
\mathfrak{d}(\square)\le\mathfrak{d}^{\sharp}(\square)
\le 1+C^{\mathtt{a}}e^{\frac18\delta_Pd^{\wedge}(\square,R_{\mathsf{n}})}
\qquad(\mathfrak{d}^{\sharp}\ge1);
\end{gathered}
\end{equation}
elsewhere $\mathfrak{d}=1$. Thus \eqref{5.5:eq-data-factor-a} holds at every block in its
one-machine form, with degree $0$ in $\Theta$, and $C^{\mathtt{a}}\ge C^{\mathtt{a}}_{\mathrm{pr}}$.

\item[(e)] Let $\mathsf{A}$, $\mathsf{B}$ be two machines on one scaffold, and let
$\mathcal{K}:=\mathcal{K}^{rd}$ be their pair system, on which the four closure properties
(closure under restriction, closure under the response dilation, closure under the read
moves, closure under the paired charts) hold by the closure statements assumed above, as
does the closure under $\mathsf{M}8$, by the same assumed closure statements. For a pair
$\mathsf{p}$ of $\mathcal{K}^{rd}$ with components $P^{\mathsf{A}}$ and $P^{\mathsf{B}}$ put
$\delta f(\mathsf{p}):=f^{\mathsf{B}}(P^{\mathsf{B}})-f^{\mathsf{A}}(P^{\mathsf{A}})$,
$\mathsf{S}:=\mathsf{S}^{rd}$ and $\mathsf{a}:=\mathsf{S}^{rd}[\hat{A}]$, where $\hat{A}$ is
the slot letter of the Wilson and counterterm plaquette functional. Let
$\mathsf{R}^{\times}_{\mathsf{p}}$ be the pair of the move $\mathsf{M}2$ at the pieces read
on the box pair and the pair of the move $\mathsf{M}3$ at own-read units. Then for every near
unit $v$ and every $\mathsf{p}\in\mathcal{K}^{rd}(k;\mathsf{n}+1,\hat{X})$ (the pairs of stage
$\mathsf{n}+1$ over the region $\hat{X}$ of the source, at level $k$), on the filed disc,
\begin{multline}\label{5.5:eq-weight-h}
\bigl|\delta(f_v\circ\mathsf{R}^{\times}_{\mathsf{p}})-\delta f_v(\mathbb{1})\bigr|\\
\le w^{\vee}(y)^3\,\mathfrak{N}^P(v;c_W)
\Big|_{\mathfrak{n}\mapsto\mathsf{S}+\mathfrak{n}\mathsf{a},\;C_I'\mapsto C_I'',\;
\mathsf{s}_v\mathsf{v}_v\mapsto C^{\perp}\mathsf{S}^{\perp}_v}.
\end{multline}
The last substitution is made at the row of (b) for the remainders of the third class:
$\mathsf{s}_v$, $\mathsf{v}_v$ are the letters of the bound
$\mathfrak{N}^W_v=K^W\mathsf{s}_v\mathsf{v}_ve^{-\frac14\delta D_v}\alpha_{\cdot,j}$ of the
remainder statement for the third class, with $K^W$ the constant of that bound and
$\alpha_{\cdot,j}$ the level-$j$ radius of Definition~\ref{5.5:def-weight-constants};
$\mathsf{S}^{\perp}_v$ is the $\delta$-size in the slot $\circ$ of the third-class cell-unit
$v$, bounded a priori by $2\max\{\mathfrak{n}^{\perp}_v,\mathfrak{n}_{\hat{A}}\}$ in the
two-run statement for the third class, $\mathfrak{n}^{\perp}_v$ being the size bound of $v$ in
the slot $\circ$ and $\mathfrak{n}_{\hat{A}}$ being the largest bound on
the Wilson and counterterm functional over the moves that land in the slot $\circ$;
$C^{\perp}:=720L^2$. At the third class (the Wilson pieces of the block and the counterterm
companions of the cores, with their remainders), \eqref{5.5:eq-weight-h} holds on the chart
as written wherever the weighted scaffold majorants of
Definition~\ref{5.5:def-windowed-source} exist on the unit's own box (the read cores); off the
chart --- the
Wilson pieces of the block at every $\nu$, the remainder rows of the companions at every
$\nu$, and the second member of the two-run bound of the third-class remainder statement ---
it holds by the two-run statement for the third class together with the closure under
$\mathsf{M}8$. For the item $\mathfrak{n}\mathsf{a}$ of the counterterm companions at
$\nu\le2$, where the radius of the chart in the slot $\hat{A}$ would exceed one, the value is
supplied by the closure under the move $\mathsf{M}2$ in the slot $\circ$, while the table
factor requires the low-order variation bound at the second cube size, which is not a
hypothesis of this lemma; for that item no claim is made. The two-cutoff response statement
at zeroth order and the flat energy-difference bound are not hypotheses of this lemma.

Consequently the two-machine differences of the pieces entering the one-machine form of the
data-factor bound of Definition~\ref{5.5:def-data-factor}, at the one-block set
$\{\square\}$, satisfy that bound in its two-machine form, with the bound
\begin{equation*}
\mathsf{l}^2(\square)\cdot\mathfrak{n}_0\big|_{E_0+1\mapsto\mathsf{l}^J}\cdot
\bigl(\vartheta(\square)+\mathfrak{p}_{\mathsf{j}}\mathbb{1}[\square\ \text{meets}\
R_{\mathsf{n}}]\bigr),
\end{equation*}
where $E_0$ and $\mathfrak{n}_0$ are as in Definition~\ref{5.5:def-data-factor}; the factor
$\mathsf{l}^2(\square)$ is the same as in~(d), and the ratio $C_I''/C_I'$ is absorbed into
$\mathfrak{n}_0$; the cost of this transport is carried by the size-free member
$\kappa_R^{\times}$ of $\kappa_R$ in \eqref{5.5:eq-weight-constants-a}. Here
\begin{equation*}
\begin{split}
\mathsf{l}^J:=\max\Bigl\{&\max_{j\le k,\,\mathcal{F}}\max\bigl\{\mathsf{S}_j[\mathcal{F}],\
\mathsf{a}_j,\
\mathsf{S}_j[\mathcal{F}]+\mathfrak{n}_{\mathcal{F}}\mathsf{a}_j,\
\mathsf{b}_j+b_+\mathsf{a}_j\bigr\},\
\sup_v\mathsf{S}^{\perp}_v,\\
&\sup_{n\le k}C\sum_{j\le n}\bar{\rho}^{\,2}_{n,j}\,\mathsf{S}_j[R^{(j)}]\Bigr\},
\end{split}
\end{equation*}
the constant $C$ being that of the third per-cell row, and
the supremum over $v$ running over the third-class cell-units. In the letters of the
definition of the two-machine data functional, $\mathsf{S}_j[\mathcal{F}]$ is the level-$j$
$\delta$-size of the family $\mathcal{F}$ and $\mathfrak{n}_{\mathcal{F}}$ its size bound,
$\mathsf{a}_j$ is the level-$j$ $\delta$-size in the slot $\hat{A}$, $\mathsf{b}_j$ and $b_+$
are the further $\delta$-size and coefficient of that definition, $\bar{\rho}_{n,j}$ is its
radius ratio of the levels $n$ and $j$, and $R^{(j)}$ is the level-$j$ family of the third
row. Described in the notation of
this section, $\mathsf{l}^J$ is
the maximum of the per-cell row sums of the near-unit table at the $\delta$-sizes, summed
over $j$ (the third row being the last member of the display), and of
$\mathsf{S}$, $\mathsf{a}$, $\mathsf{S}+\mathfrak{n}\mathsf{a}$,
$\mathsf{b}_j+b_+\mathsf{a}$ and $\sup_v\mathsf{S}^{\perp}_v$. It is read on
$\mathcal{K}^{rd}\subset\mathrm{cl}^{pr}_J$, the production class, so that the same
functional formed on
$\mathrm{cl}^{pr}_J$ dominates it; it is a function of the $\delta$-sizes, and it is one of
the members of the data functional $\mathsf{l}_{\mathsf{n}+1}$ of stage $\mathsf{n}+1$. Only
the definition of $\mathsf{l}^J$ enters
this lemma; bounds on $\mathsf{l}^J$ are not used in it.
\end{enumerate}
\end{lemma}

Parts (a1) and (a2) are proved below. The remaining parts are proved in the lemmas that follow
in this section; part~(c$'$) is proved in the lemma on the count times weight.

\begin{proof}[Proof of \textup{(a1)} and \textup{(a2)}]
\emph{(a1).} This clause is one of the defining conditions of the source class; no estimate
is involved. Since $\hat{X}\supseteq\hat{\mathsf{K}}_{\square}\supseteq\square_0$, the
triple $(k,\square_0,\mathfrak{B}^+)$ is one of the triples of the source over $\hat{X}$ at
which the conditions defining the admissible source class $\mathfrak{K}_P$ are imposed: the
definition of that class imposes its second and fourth entries for every triple
$(p,X',\mathfrak{s})$ among the current one and all its predecessors, and on the pieces
frozen by a live arrest $\mathfrak{a}$ it imposes them in the form attached to
$\mathfrak{a}$. By the definition of $\mathbb{W}=U(\mathfrak{B}^+(\square_0),M^{\cdot}(\mathbb{U}))$
in Definition~\ref{5.5:def-windowed-source}, the second entry of this triple at a cell $x$
is $|\partial\mathbb{W}(x)-1|$. In the reading with weights and arrest factor in which the
source class is taken (Definition~\ref{5.5:def-windowed-source}), the condition on this entry
states that the second entry at $x$ is less than
$F^{(k')}(x)\,w(x)\,r_2(\ell(x))$, taken at the levels and weights
$(\ell^+,w^+)$ of $\mathfrak{B}^+$; here $F^{(k')}$ is the arrest factor, $k'$ being the step
index it carries in that condition (the letter $F^{(k')}$ is used only in this paragraph). The
arrest factor satisfies $F^{(k')}(x)\le1+2\beta$, by
the bound on the arrest factor, where $\beta=\frac15$ is the parameter of that bound (not the
H\"older exponent of (a3)); thus $F^{(k')}(x)\le1.4$. At the data of $\mathfrak{B}^+$ one has
$w^+(x)r_2(\ell^+(x))\le\theta_2(x)$, by the comparison of the refined levels and weights with
the current ones. Hence
\begin{equation*}
|\partial\mathbb{W}(x)-1|<F^{(k')}(x)\,w^+(x)\,r_2(\ell^+(x))\le1.4\,\theta_2(x).
\end{equation*}
The condition is imposed on the set $X'^{\flat}$ attached to the triple, and
$X'^{\flat}\supset\tilde{\square}^3$; this gives (a1) for every cell $x$ meeting
$\tilde{\square}^3$. The window is read in the same way in \cite[(3.40)]{B-CMP109},
where the condition numbered (iv) of \cite[p.~278]{B-CMP109} is stated there to be a
consequence of its condition (iii).

\emph{(a2).} By Definition~\ref{5.5:def-windowed-source}, $\mathfrak{B}^*_j$ refines
$\mathfrak{B}^+(\square_0)$. The configuration $\mathbb{W}$ is minimal under the data of
$\mathfrak{B}^+(\square_0)$, and these data are averages of the data of $\mathfrak{B}^*_j$;
so $\mathbb{W}$ is critical under its own $\mathfrak{B}^*_j$-data, and
\begin{equation*}
W':=U(\mathfrak{B}^*_j,M^{\cdot}(\mathbb{W}))=\mathbb{W}^g
\end{equation*}
for a gauge transformation $g$, by \cite[(3.39)]{B-CMP109}; for complex $\mathbb{U}$ the
identity persists by the analytic continuation furnished in the setting of the perturbation
lemma for the windowed minimiser. The reason is that, by \cite[Theorem~1]{B-CMP102b}, the
minimal
orbit is an orbit of the group of gauge transformations which fix the constraints; these
transformations are the identity at the base points of the blocks, so $g=1$ there, and the
gauge factor is local to each block. For $x$ in the block with base point $y_0$, let
$\Gamma$ be the contour from $y_0$ to $x$; then
$g(x)=W'(\Gamma)^{-1}\mathbb{W}(\Gamma)$. Write $n(\cdot)$ for the size function on complex
group elements used in the defining condition of the admissible source class $\mathfrak{K}_P$
that controls the configuration along contours, $a_{\flat}$ for the radius in that condition,
$\mathcal{H}(B_0)$ for the field variable on the base cube $B_0$ of the block that enters the
contour bound of the perturbation lemma, and $\ell$
for the length scale entering the bound on it (not the level $\ell(x)$ of a cell). That
defining condition gives
$n(\mathbb{W}(\Gamma))\le e^{2da_{\flat}}$, and the bound on $W'$ along contours in the
setting of the perturbation lemma gives
$n(W'(\Gamma))\le2^{1/2}e^{2d\ell\sup|\mathcal{H}(B_0)|}$. These two bounds give
$n(g)\le2$ under the smallness condition of the perturbation lemma among the ceilings of
Definition~\ref{5.5:def-weight-constants}, since $a_{\flat}\le5\mathsf{A}_w(\gamma)$.
Combining $n(g)\le2$ with (a1), the second entry
$E_2^{\mathfrak{B}^*_j}(\mathbb{W})(x)$ of $\mathbb{W}$ relative to the data of
$\mathfrak{B}^*_j$ satisfies
\begin{equation*}
E_2^{\mathfrak{B}^*_j}(\mathbb{W})(x)\le2\cdot1.4\,\theta_2(x)=2.8\,\theta_2(x).
\end{equation*}

The perturbation lemma for the windowed minimiser holds in every frame. Apply it in the frame
$(\square_0,\mathfrak{B},\mathfrak{B}^*_j)$, with $\mathbb{W}$ as its configuration and
$\mathbf{A}$ in the slot of its fluctuation. Its hypotheses hold. Its defining condition on
the source holds because $(\mathbb{W},\mathbb{J})\in\mathfrak{K}_P$. Its two regularity
conditions and its smallness condition hold by the second clause of the corollary on the
parameters of the perturbation lemma, applied at the parameter value one half. That clause
bounds the first parameter of the perturbation lemma by one quarter of a radius which is
itself at most one, hence by one quarter, and it bounds the second parameter by the first.
With the first parameter at most $\frac14$, the perturbation lemma multiplies the entry bound
$2.8\,\theta_2(x)$ by at most $\frac54$ and adds at most $\frac14r_2(\ell(x))$; it thus gives
\begin{equation*}
|\partial\mathfrak{b}'_j(x)-1|\le2.8\,\theta_2(x)\cdot\tfrac54+\tfrac14r_2(\ell(x))
\le3.75\,\theta_2(x)\le4\,\theta_2(x),
\end{equation*}
where $2.8\cdot\frac54=3.5$ and $\frac14r_2(\ell(x))\le\frac14\theta_2(x)$, the weight
$w(x)$ being at least one by Notation~\ref{5.5:not-weight-notation}. The bound holds
pointwise in
$(t,\mathsf{t})\in[0,1]\times\mathbb{D}_{\square}$, which proves \eqref{5.5:eq-weight-a};
it is the entry bound used for the pieces of the third class in~(b).
\end{proof}

\begin{proof}[Proof of parts (a3) and (a4) of Lemma~\ref{5.5:lem-weight}: small-field regularity on a local box]\label{5.5:prf-box-regularity}
We work in the setting of the weight lemma (Lemma~\ref{5.5:lem-weight}). Its parts (a1) and (a2), the latter displayed in \eqref{5.5:eq-weight-a}, are used as inputs. The constants $K_{ax}$, $C_M^{av}\ge 2$, $K_h$, $c_1$, $c_{\cdot}$ and $K_{loc}$ are those of Definition~\ref{5.5:def-weight-constants}. In particular
\begin{equation*}
K_{loc} = 4\bigl(c_{\cdot}+3K_hK_{ax}MC_M^{av}\bigr),
\end{equation*}
see \eqref{5.5:eq-weight-constants-c}, and $K_h,\ K_{ax}M,\ c_1\ge 1$. The ceiling on $\gamma$ is \eqref{5.5:eq-weight-constants-d}, namely $K_{loc}\,\mathsf{A}_w(\gamma)\le c_h$, where $c_h$ is half the least of the absolute smallness constants of four published statements, each taken for the quantity it bounds: \cite[Proposition~2]{B-CMP98}, \cite[Theorem~1]{B-CMP102b}, \cite[Proposition~9]{B-CMP102b}, and the Baker--Campbell--Hausdorff step \cite[(3.30)--(3.32)]{B-CMP109}.

Let $Q$ be a box as in (a3). Thus $Q$ has at most $9$ cubes of $\pi_p$ a side, $Q^{(1)}\subset\tilde{\square}^3$, and every cell meeting $Q^{(1)}$ is of level at least $p$. Further, $w_Q$ and $\alpha_Q$ are the largest values of $w(\Delta)$ and of $\alpha_{0,\ell(\Delta)}$ over those cells. Sizes said to be taken \emph{in level-$p$ units} carry the powers of $\ell_p$ of \eqref{5.5:eq-weight-b}.

\emph{Part} (a3). \emph{The minimiser of $Q$.} Let $\mathfrak{B}_p(Q)$ be the maximal sequence of block partitions on $Q$ in the sense of \cite{B-CMP109}. Put
\begin{equation*}
\mathsf{V}_Q := M^{\cdot}(\mathbb{W})|_{\mathfrak{B}_p(Q)},
\end{equation*}
the family of averages of $\mathbb{W}$ over the blocks of $\mathfrak{B}_p(Q)$. We apply the locality part of the lemma on the constrained minimiser in a local box to $(\mathbb{W},Q,p)$, with $Q$ as its box and $\mathfrak{B}_p(Q)$ as its sequence. It yields two facts.

First: for any configuration $U$, every level-$p$ average $M^p(U,c)$ of $U$ is a function of the values of $U$ off the box and of the averages $M^{\cdot}(U)|_{\mathfrak{B}_p(Q)}$. Consequently the same holds for those data of the sequence $\mathfrak{B}^+(\square_0)$ whose blocks meet $Q$; all of these are of level at least $p$.

Second: let a variation of $\mathbb{W}$ be supported inside $Q$ and keep $\mathsf{V}_Q$ fixed. By the first fact it keeps those data of $\mathfrak{B}^+(\square_0)$ fixed, and it vanishes on the outer layer of $Q$. Now $\mathbb{W}$ is minimal under the data of $\mathfrak{B}^+(\square_0)$. Hence $\mathbb{W}|_Q$ is critical under the constraint $\mathsf{V}_Q$ in the smaller class of configurations on $Q$.

The critical point is unique by \cite[Theorem~1]{B-CMP102b}. The smallness that theorem requires is here
\begin{equation*}
1.4\,w(\Delta)\,\tilde{\alpha}_{0,\ell(\Delta)}\le 2.8\,\mathsf{A}_w(\gamma),
\end{equation*}
where $\tilde{\alpha}_{0,\ell(\Delta)}$ is the enlarged radius at level $\ell(\Delta)$, which satisfies $\tilde{\alpha}_{0,\ell(\Delta)}\le 2\alpha_{0,\ell(\Delta)}$, and $2.8\,\mathsf{A}_w(\gamma)$ is paid by the ceiling, as shown at the end of this part. Therefore
\begin{equation}\label{5.5:eq-box-regularity-1}
\mathbb{W}|_Q = U_p(Q,\mathsf{V}_Q)\quad\text{modulo gauge.}
\end{equation}
At complex data, \eqref{5.5:eq-box-regularity-1} follows from the real case by the identity theorem for holomorphic functions. The argument is the one used in the lemma on the constrained minimiser in a local box.

\emph{The data.} We apply \cite[Proposition~2]{B-CMP98} at complex data, in its extension to complex data given by the complex form of the averaging bound of this chapter, and feed it the bound (a1). For levels $\ell\ge p$, using $\tilde{\alpha}\le 2\alpha$, this gives in level-$p$ units
\begin{equation}\label{5.5:eq-box-regularity-2}
|\partial\mathsf{V}_Q-1|
\le 1.4\,C_M^{av}\,w_Q\,\tilde{\alpha}_{0,\ell}\,(\ell_p/\ell_\ell)^2
\le 2.8\,C_M^{av}\,w_Q\,\alpha_Q .
\end{equation}

\emph{The axial gauge on $Q$.} In the axial gauge of \cite[(3.26)]{B-CMP109} on a box of at most $9$ cubes a side, each bond variable is a product of at most $K_{ax}M$ conjugated plaquette variables of the box, with distortion at most $2$. Write $\mathsf{V}_Q^{v}$ for the data $\mathsf{V}_Q$ in this gauge. By \eqref{5.5:eq-box-regularity-2}, $\mathsf{V}_Q^{v}=e^{iB}$ with
\begin{equation}\label{5.5:eq-box-regularity-3}
|B|\le 5.6\,K_{ax}MC_M^{av}\,w_Q\,\alpha_Q\le K_{loc}\,\mathsf{A}_w(\gamma).
\end{equation}

\emph{Representation.} We apply \cite[Proposition~9]{B-CMP102b} at the identity configuration $\mathbb{1}$. The representation parts of the lemma on the constrained minimiser in a local box read the data through $B$ alone. They give
\begin{equation*}
U_p(Q,e^{iB}) = \bigl(e^{i\eta\mathcal{H}(B)}\bigr)^{u}
\end{equation*}
for a field $\mathcal{H}(B)$ on $Q$, determined by $B$, and a gauge transformation $u$. On $Q^{\sim-1}$, moreover,
\begin{equation*}
\max\bigl\{\ell_p|\mathcal{H}|,\ \ell_p^2|\nabla\mathcal{H}|,\
\ell_p^{2+\beta}\|\nabla\mathcal{H}\|_\beta\bigr\}\le K_h\sup|B| .
\end{equation*}
Combine this with \eqref{5.5:eq-box-regularity-1} and \eqref{5.5:eq-box-regularity-3}. There is a gauge $g$ on $Q^{\sim-1}$ in which $\mathbb{W}=e^{i\eta H}$, with
\begin{equation}\label{5.5:eq-box-regularity-4}
\max\bigl\{\ell_p|H|,\ \ell_p^2|\nabla H|,\ \ell_p^{2+\beta}\|\nabla H\|_\beta\bigr\}
\le 5.6\,K_hK_{ax}MC_M^{av}\,w_Q\,\alpha_Q .
\end{equation}
This is \cite[(3.27)]{B-CMP109}, with $Q$ in place of $\square_0$ and $w_Q\alpha_Q$ in place of the single $\alpha_0$.

\emph{The layer.} In the gauge $g$ we have $\mathfrak{b}^g=e^{i\eta\operatorname{Ad}_g\mathbf{A}}\,e^{i\eta H}$. By the bound on the source fluctuation of this chapter, on $\mathbb{D}_\square$
\begin{equation*}
N^{\Delta}_x(\mathbf{A})\le\alpha_{\cdot,\ell(x)},
\end{equation*}
where $N^{\Delta}_x$ is the local norm at $x$ on the cell $\Delta$; its H\"older member is included, and no gap between levels enters. This bound does not involve the weight. In level-$p$ units it is at most $c_{\cdot}\alpha_Q$, by the definition $c_{\cdot}=\sup_n\alpha_{\cdot,n}/\alpha_{0,n}$. The conjugation $\operatorname{Ad}_g$ costs a factor $2$. Put
\begin{equation*}
A:=\frac{1}{i\eta}\log\bigl(e^{i\eta\operatorname{Ad}_g\mathbf{A}}\,e^{i\eta H}\bigr).
\end{equation*}
The Baker--Campbell--Hausdorff bounds \cite[(3.30)--(3.32)]{B-CMP109} are local. With \eqref{5.5:eq-box-regularity-4} they give
\begin{equation*}
\begin{aligned}
\max\bigl\{\ell_p|A|,\ \ell_p^2|\nabla A|,\ \ell_p^{2+\beta}\|\nabla A\|_\beta\bigr\}
&\le 2\bigl(2c_{\cdot}\alpha_Q+5.6\,K_hK_{ax}MC_M^{av}\,w_Q\,\alpha_Q\bigr)\\
&\le K_{loc}\,w_Q\,\alpha_Q=\mathsf{a}_Q .
\end{aligned}
\end{equation*}
The last inequality uses $11.2\le 12$ against the definition of $K_{loc}$. This is \eqref{5.5:eq-weight-b}.

\emph{The smallness spends.} The ceiling \eqref{5.5:eq-weight-constants-d} is one-sided, and it pays all four smallness requirements used above.
\begin{itemize}
\item The quantity $2.8\,\mathsf{A}_w(\gamma)$ is used for \cite[Theorem~1]{B-CMP102b} and for \cite[Proposition~2]{B-CMP98}. Since $K_{loc}\ge 12C_M^{av}\ge 24$, it is at most $K_{loc}\mathsf{A}_w(\gamma)\le c_h$.
\item The quantity $K_{loc}\mathsf{A}_w(\gamma)$ is used for \cite[Proposition~9]{B-CMP102b}; it is at most $c_h$.
\item The quantity $2K_{loc}\mathsf{A}_w(\gamma)$ is used for the Baker--Campbell--Hausdorff step. It is at most $2c_h$, the least of the four smallness constants.
\end{itemize}

\emph{Part} (a4). \emph{The identity.} Let $j\le p$, and let $\square_0'\subset Q^{\sim-2}$ be a box of $\pi_j$-cubes. The configuration $\mathfrak{b}'_j$ is minimal under the data of $\mathfrak{B}^*_j$. Near $\square_0'$ these data are of level exactly $j$: the levels there are at least $p\ge j$, and the refinement $\mathfrak{B}^*_j$ caps them at $j$. We run the argument of the first paragraph of part (a3) at $(\mathfrak{b}'_j,\square_0',j)$ in place of $(\mathbb{W},Q,p)$. The regularity it needs is supplied by (a2) and by the bound on the fourth regularity entry of the admissible triple in the lemma on the bump construction of this chapter. It gives
\begin{equation*}
\mathfrak{b}'_j|_{\square_0'}=(U_j,J_j)(\square_0',\mathsf{V}')\quad\text{modulo gauge},
\end{equation*}
which is the identity \cite[(3.38)]{B-CMP109}.

\emph{The part of level $j$.} Block averages are covariant under gauge transformations, so
\begin{equation*}
\mathsf{V}'=M^j(\mathfrak{b})=M^j\bigl(e^{i\eta A}\bigr)^{g}.
\end{equation*}
We apply the averaging part of the lemma on the constrained minimiser in a local box, with $H$ replaced by $A$. This is legitimate because \cite[(4.16)--(4.18)]{B-CMP109} read the field only through the bounds \cite[(3.32)]{B-CMP109}, and part (a3) supplies these for $A$ with size $\mathsf{a}_Q$. With $t_p=L^{j-p}$ the result is
\begin{equation*}
|\partial\mathsf{V}'-1|\le c_1\mathsf{a}_Q t_p^2 .
\end{equation*}
It also gives the comb-gauge representation \eqref{5.5:eq-weight-c}, with $|F|\le c_1\mathsf{a}_Q t_p^2$ and, on the bonds of level $j$,
\begin{equation*}
|\mathfrak{e}_\mu(x)|\le c_1(1+|x-y|)^2\,\mathsf{a}_Q\,t_p^{2+\beta}.
\end{equation*}

\emph{The graded rim.} Near $\partial\square_0'$ there are levels $i<j$, and there $\mathsf{V}'=M^i(\mathfrak{b}'_j)$. Only supremum bounds are used at these levels. Part (a2), combined with \cite[Proposition~2]{B-CMP98}, bounds the deviation of a rim plaquette of level $i$ by
\begin{equation*}
4C_M^{av}\,w(x)\,\tilde{\alpha}_{0,\ell(x)}\,(\ell_i/\ell_{\ell(x)})^2
\le 8C_M^{av}\,w_Q\,\alpha_Q\,t_p^2\,L^{2(i-j)}
\le c_1\mathsf{a}_Q t_p^2 .
\end{equation*}
The last inequality holds because $L^{2(i-j)}\le 1$ for $i<j$, $K_{loc}\ge 8C_M^{av}$ and $c_1\ge 1$. Every plaquette deviation, of level $j$ or on the rim, is therefore at most $c_1\mathsf{a}_Qt_p^2$. Hence, on all bonds,
\begin{equation*}
|B^{ax}(b)|\le 2d(1+|b-y|)\,c_1\mathsf{a}_Q t_p^2 .
\end{equation*}
On the graded-rim bonds, $\mathfrak{e}^{rim}$ is $B^{ax}$ minus the linear term $\sum_{\kappa<\mu}(x-y)_\kappa F_{\kappa\mu}(y)$ of \eqref{5.5:eq-weight-c}. The bound on $|F|$ then gives
\begin{equation*}
|\mathfrak{e}^{rim}(b)|\le|B^{ax}(b)|+d\,|b-y|\,|F|\le 4d(1+|b-y|)\,c_1\mathsf{a}_Q t_p^2 .
\end{equation*}
Here we used $d|b-y|\le 2d(1+|b-y|)$. This completes the proof of (a4).

The argument above takes no supremum over $\square_0$ of any quantity, does not use the bounds \cite[(3.26)--(3.27)]{B-CMP109} on $\tilde{\square}^5$, and does not use the hypothesis of \cite[Lemma~4]{B-CMP109} on $\square_0$.
\end{proof}

\begin{proof}[Proof of Lemma~\ref{5.5:lem-weight}, part (b): the rows of the near-unit table]
\label{5.5:prf-near-unit-rows}
Throughout, $v$ is a near unit with anchor $y$ and $X_v\subset\tilde{\square}^2$. As in the setting
of Lemma~\ref{5.5:lem-weight}, $n=\min\{\ell(y),n_0\}$, $\nu=1+n-j$, $t=L^{j-n}$ and
$\Theta=c_W w^{\vee}(y)\ge 1$, and we put $z:=\Theta\rho_{n,j}$ as in the table
\eqref{5.5:eq-weight-d}. The weighted scaffold majorants $|B^{ax}|_*$, $|F|_*$,
$|\mathfrak{e}_\mu|_*$ and $|\mathfrak{e}^{rim}|_*$ are those of \eqref{5.5:eq-windowed-source-b}.
The radius $S^{\star}_v$ and the inequality it is to satisfy are those of
\eqref{5.5:eq-near-unit-classes-a}, and the filed discs are those of
\eqref{5.5:eq-near-unit-classes-b}. The filing into read, box-read, own-read and third-class units
is that of Definition~\ref{5.5:def-near-unit-classes}.

The analysis of the near units on the cover of blocks rests on the transfer of the analysis of
\cite[\S3]{B-CMP109} to that cover, printed without proof in \cite[p.~276]{B-CMP119}; cited as
printed. From the lemmas of that analysis, as they are stated in this book, we use the following
statements:
\begin{itemize}
\item its box conditions (c1), (c2) and (c3) on a box pair;
\item its short-single lemma;
\item its core lemma, which bounds the core term $\mathfrak{i}_y$ at the anchor $y$;
\item its bound for the Wilson pieces, clause (b);
\item its counterterm projection bound.
\end{itemize}
From the column bounds at the localisation site we use:
\begin{itemize}
\item the step-by-step bounds in the proof of the core lemma;
\item the field-free column bound.
\end{itemize}

\smallskip\noindent\emph{(i) Boxes; the weighted majorants.}
Let $v$ be a read unit, so that $\nu\ge 3$, and put $p:=n-1$. Let $Q_v$ be the union of the cubes of
$\pi_{n-1}$ meeting $\mathfrak{r}_v$, together with two further layers of cubes of $\pi_{n-1}$
around it. For a box $Q$ made of cubes of $\pi_{n-1}$ we write $Q^{(1)}$ for $Q$ enlarged by one
layer of such cubes, and $Q^{\sim-2}$ for $Q$ with its two outer layers removed.

Since $n-1-j=\nu-2$, a level-$j$ unit has length $L^{-(\nu-2)}$ in units of $\pi_{n-1}$. Hence,
measured in cubes of $\pi_{n-1}$, the sides of $\mathfrak{r}_v$ are at most
\begin{equation*}
\max\bigl\{(2\nu+3)L^{-(\nu-2)},\;1+5L^{-(\nu-2)}\bigr\}\le \tfrac{13}{8}<2 .
\end{equation*}
Indeed, $(2\nu+3)L^{-(\nu-2)}$ decreases in $\nu$ and equals $9/L\le 9/8$ at $\nu=3$. Likewise
$1+5L^{-(\nu-2)}\le 1+5/L\le 13/8$ for $\nu\ge 3$ and $L\ge 8$.

A segment of length less than $2$ meets at most $3$ unit intervals. So the union of the cubes of
$\pi_{n-1}$ meeting $\mathfrak{r}_v$ has at most $3$ cubes a side. Consequently $Q_v$ has at most
$7$ cubes a side and $Q_v^{(1)}$ at most $9$. Every point of $Q_v^{(1)}$ lies within $6$ cubes of
$\pi_{n-1}$ of $y$, that is within $6/L<1$ level-$n$ unit of $y$.

As $y\in\tilde{\square}^2$ and $n\le n_0$, this gives $Q_v^{(1)}\subset\tilde{\square}^3$. The cells
of $Q_v^{(1)}$ have level at least $n-1$, by the level bound for the cells meeting
$\tilde{\square}^3$. Since every cell meeting $Q_v$ is at level-$n$ distance at most $6/L<1$ from
$y$, the
definition of $w^{\vee}(y)$ gives $w_{Q_v}\le w^{\vee}(y)$. Further,
$\alpha_{Q_v}\le 2\alpha_{0,n}$ by the bound on the scaffold radii near an arrest, and
$\mathfrak{r}_v\subset Q_v^{\sim-2}$ by construction.

On $Q_v$ we read the small-field regularity on a local box, established in the part of the proof
of Lemma~\ref{5.5:lem-weight} under that heading. We apply the comb-gauge display
\eqref{5.5:eq-weight-c} at
$(Q,p,\square_0'):=(Q_v,n-1,\mathfrak{r}_v)$, for each of the two machines $\mathsf{A}$,
$\mathsf{B}$ by the per-machine hypothesis of the pair system $\mathcal{K}$, with the H\"older
exponent of \eqref{5.5:eq-weight-c} taken to be $\hat{\alpha}$, the exponent of condition (c3),
and with $t_p=L^{j-(n-1)}=Lt$. The definition of $c_W$ in
\eqref{5.5:eq-weight-constants-c} gives $\Theta\ge 2c_1K_{loc}L^3\alpha_{0,n}w^{\vee}(y)/\mathfrak{a}_n$.
Using also $\hat{\alpha}\le 1$, this yields
\begin{align*}
c_1K_{loc}\,w^{\vee}(y)\cdot 2\alpha_{0,n}L^2t^2
 &\le \Theta\,\mathfrak{a}_n t^2/L,\\
c_1K_{loc}\,w^{\vee}(y)\cdot 2\alpha_{0,n}L^{2+\hat{\alpha}}t^{2+\hat{\alpha}}
 &\le \Theta\,\mathfrak{a}_n t^{2+\hat{\alpha}} .
\end{align*}

The factor $L^{-1}$ in the first line is the one reserved in the power $L^3$ of $c_W$; it is the
spare factor used by the remainder hypothesis for box-read units of Lemma~\ref{5.5:lem-weight}. It
follows that conditions (c1) and (c3) hold for the box pair $\mathfrak{b}'_j$ on $\mathfrak{r}_v$ with
$(\mathfrak{a}_n,t)$ replaced by $(\Theta\mathfrak{a}_n,t)$. It also follows that the weighted
majorants dominate pointwise:
\begin{equation}\label{5.5:eq-near-unit-rows-1}
\begin{gathered}
|B^{ax}(b)|\le L^{-1}|B^{ax}|_*(b),\qquad |F|\le |F|_*,\\
|\mathfrak{e}_\mu|\le|\mathfrak{e}_\mu|_*\ \text{on level-$j$ bonds},\qquad
|\mathfrak{e}^{rim}|\le|\mathfrak{e}^{rim}|_*\ \text{on the graded rim of } \mathfrak{r}_v .
\end{gathered}
\end{equation}

Let $S=\alpha_{3,j}/\sup_b|B^{ax}(b)|$ be the radius of the field itself. Since
$|B^{ax}|\le|B^{ax}|_*$, comparison with the definition of $S^{\star}_v$ and with
\eqref{5.5:eq-near-unit-classes-a} gives
\begin{equation}\label{5.5:eq-near-unit-rows-2}
S\ \ge\ S^{\star}_v\ \ge\ \bigl(\Theta C_M D_X t^2\rho_{n,j}\bigr)^{-1}.
\end{equation}

\smallskip\noindent\emph{(ii) The rows.}

\emph{Box-read short single ($S^{\star}_v>1$).} The proof of the short-single lemma reads
$\mathfrak{b}'_j$ through condition (c1), which holds by step (i). It uses the chart
\begin{equation*}
\varphi(\sigma'):=F\bigl(X,U_j(\mathfrak{r}_v,e^{i\sigma'B^{ax}}),\dots\bigr).
\end{equation*}
For $|\sigma'|<S^{\star}_v$ we have $|\sigma'B^{ax}(b)|<S^{\star}_v|B^{ax}|_*(b)\le\alpha_{3,j}$.
So the chart stays in the domain given by condition (c2) and \cite[Lemma~4]{B-CMP109}, and $\varphi$
is holomorphic and bounded by $\mathfrak{n}$ on the disc $|\sigma'|<S^{\star}_v$. This disc contains
$1$ since $S^{\star}_v>1$. Moreover $\varphi'(0)=0$. The Schwarz lemma, in the form of
Lemma~\ref{5.2:lem-disc-bounds}, then gives the entry $2\mathfrak{n}_vS^{\star-2}_v$.

\emph{Box-read core, taken whole.} We use the steps of the proof of the core lemma rather than its
statement.
Each of its steps is a polynomial of degree at most $3$ in the data constant:
\begin{itemize}
\item the cubic Schwarz remainder on the chart;
\item the bilinear term $\varphi_X''(0)=D^2F(X;1,0)[\cdot,\cdot]$, through the bound
\begin{equation*}
|D^2F(X;1,0)[B_1,B_2]|\le 4\mathfrak{n}\,e^{-\kappa d_j(X)}\,|B_1|_\infty|B_2|_\infty\,
\alpha_{3,j}^{-2};
\end{equation*}
\item the main term of condition (c3);
\item the counterterm $A(h_y,\cdot)$, read on the same box pair $\mathfrak{b}'_j$, as the third class
files it.
\end{itemize}
Step (i) supplies the inputs with $\mathfrak{a}_n$ replaced by $\Theta\mathfrak{a}_n$. The
core lemma's bound therefore holds with $\rho_{n,j}$ replaced by $z$. Membership of the
arguments is given by condition (c2) on the chart. With the rim contribution of step (iii)
included, the constant is $C_I^{\sharp}\le C_I'$, and the entry is
$C_I'\mathfrak{n}\nu^3z^2(1+z)t^{4+\hat{\alpha}}$.

\emph{Own-read units.} By the own-read membership hypothesis of Lemma~\ref{5.5:lem-weight},
$|F(\mathcal{W}_T)|\le\mathfrak{n}$ for $|\tau|\le\rho_A(v)$. This gives the following entries.
\begin{itemize}
\item Tails, long singles and constituents with $\nu\le 2$: the entry $2\mathfrak{n}_v$.
\item The near constituents of a straddling core, at every $\nu$: such a core has
$y\in\tilde{\square}^2$ and $\mathfrak{r}_v\not\subset\tilde{\square}^2$. Its near constituents are
bounded by $2\mathfrak{n}_ve^{-\kappa d_j(X)}$ each, at the prefactor
$24K''\vartheta(\square)\eta_p$ with $\eta_p\le t^5$ by the own-read membership hypothesis, whose
two-machine clause gives the same row with $2\mathsf{S}$ in place of $2\mathfrak{n}_v$. This row
carries its own line in
\eqref{5.5:eq-weight-d}, together with its count per unit cell.
\item A read short single with $S^{\star}_v\le 1$: the entry
$2\mathfrak{n}_v=2\mathfrak{n}_v\min\{1,S^{\star-2}_v\}$, since $S^{\star-2}_v\ge 1$.
\item A read core with $S^{\star}_v\le 1$: treated in step (iv).
\end{itemize}

\emph{Third class.} The Wilson pieces $f_{\square}$ are entire. By the Wilson-piece bound, clause
(b), they are quadratic near the identity configuration. Each is supported within one level-$j$
unit of $y$, or in a block whose centre serves as $y$. Display \eqref{5.5:eq-weight-a}, applied
pointwise on $\mathfrak{b}'_j$, gives the entry $\Theta^2\mathfrak{N}_W(v)$.

For the counterterm piece of a core, the counterterm projection bound gives
$|\beta[F]|\le C\mathfrak{n}\alpha_{3,j}^{-2}$. Combined with the Wilson-piece bound, clause (b),
and \eqref{5.5:eq-weight-a}, this gives the entry $C'\mathfrak{n}z^2t^4$.

\smallskip\noindent\emph{(iii) The rim.}
The cover analysis has a step in which the insertion of condition (c3) into the core term costs
at most $C\mathfrak{n}e^{-\kappa d_j(X)}D^3\rho_{n,j}^2t^{4+\hat{\alpha}}$. That step uses the
H\"older member of (c3). In our setting, display \eqref{5.5:eq-weight-c} supplies the H\"older
member on level-$j$ bonds only, so we split $\mathfrak{e}=\mathfrak{e}^{in}+\mathfrak{e}^{rim}$.

The second derivative of the chart is
\begin{equation*}
\varphi_X''(0)=D^2F(X;1,0)\bigl[\mathsf{h}[B^{ax}],\mathsf{h}[B^{ax}]\bigr],
\qquad \mathsf{h}[B]:=\sum_b\mathsf{h}_bB(b).
\end{equation*}
Here, by \cite[(4.18)]{B-CMP109}, the family $\mathsf{h}_b$ decays exponentially, with the
constants $B_3$ and $\delta_0$ of \cite[(4.18)]{B-CMP109}:
\begin{equation*}
\max\bigl(|\mathsf{h}_b|,|\nabla\mathsf{h}_b|,\dots\bigr)\le
B_3e^{-\delta_0\operatorname{dist}_j(b,X)}=:\|\mathsf{h}_b\|_* .
\end{equation*}
For $\mathfrak{e}^{in}$, which lives on level-$j$ bonds, the insertion estimate of the cover
analysis applies unchanged.

By \eqref{5.5:eq-near-unit-rows-1} and the definition of $|\mathfrak{e}^{rim}|_*$, the remainder
$\mathfrak{e}^{rim}$ is bounded by $4dD\Theta\mathfrak{a}_nt^2$. It is supported on the outer
layer of $\mathfrak{r}_v$. Since $X\ni y$, that layer is at $\operatorname{dist}_j$ at least
$M(\nu-1-d_j(X))_+$ from $X$. Its insertion therefore costs at most
\begin{equation*}
C\mathfrak{n}D^{d+1}z^2t^4e^{-\kappa d_j(X)-\delta_0M(\nu-1-d_j(X))_+}
\ \le\ C''\mathfrak{n}\,e^{-\frac12\kappa d_j(X)}z^2t^{4+\hat{\alpha}} .
\end{equation*}

To justify the last inequality, put $\mu_0:=\min\{\tfrac12\kappa,\delta_0M\}$. Since
$d_j(X)+(\nu-1-d_j(X))_+\ge\nu-1$,
\begin{align*}
\kappa d_j(X)+\delta_0M(\nu-1-d_j(X))_+
 &\ge \tfrac12\kappa d_j(X)+\mu_0\bigl(d_j(X)+(\nu-1-d_j(X))_+\bigr)\\
 &\ge \tfrac12\kappa d_j(X)+\mu_0(\nu-1).
\end{align*}
The floor $\mu_0\ge 5\log L$ gives $e^{-\mu_0(\nu-1)}\le L^{-5(\nu-1)}$. We also have
$D=M(2\nu+3)\le 5M\nu$ and $t^{-\hat{\alpha}}=L^{\hat{\alpha}(\nu-1)}$. Hence
\begin{equation*}
D^{d+1}t^4e^{-\mu_0(\nu-1)}\ \le\ t^{4+\hat{\alpha}}\,(5M\nu)^{d+1}L^{(\hat{\alpha}-5)(\nu-1)} .
\end{equation*}
The last factor is bounded in $\nu$ because $\hat{\alpha}\le 1$. The constant $C''$ is one of the
contributions to $C_I^{\sharp}$ in Definition~\ref{5.5:def-weight-constants}.

\smallskip\noindent\emph{(iv) Near levels and the regime $S^{\star}_v\le 1$.}
Suppose $C_MDt^2z>\frac12$. Every read core with $S^{\star}_v\le 1$ is in this regime: for a core
$D_X=D$, so \eqref{5.5:eq-near-unit-rows-2} gives
\begin{equation*}
(C_MDt^2z)^{-1}\le S^{\star}_v\le 1<2 .
\end{equation*}

In this regime the core term $\mathfrak{i}_y$, read constituent by constituent together with its
counterterm companion, is trivially bounded by $2C(\kappa)\mathfrak{n}+C'\mathfrak{n}z^2t^4$. The
values are bounded as follows:
\begin{itemize}
\item at box-read units, by condition (c2) on the chart;
\item at own-read units, by the own-read membership hypothesis, constituent by constituent;
\item for the counterterm piece, by the counterterm projection bound, the Wilson-piece bound,
clause (b), and \eqref{5.5:eq-weight-a}.
\end{itemize}

Since $D\le 5M\nu$, the hypothesis $C_MDt^2z>\frac12$ gives $t^2>(10C_MM\nu z)^{-1}$. Hence
\begin{align*}
\nu^3z^2(1+z)t^{4+\hat{\alpha}}
 &\ge (10C_MM)^{-2-\hat{\alpha}/2}\,\nu^{1-\hat{\alpha}/2}\,z^{-\hat{\alpha}/2}(1+z)\\
 &\ge (10C_MM)^{-5/2}.
\end{align*}
In the second line we used $\nu\ge 1$ and $\hat{\alpha}\le 1$. We also used
$z^{-\hat{\alpha}/2}(1+z)\ge 1$: for $z\le 1$ the first factor is at least $1$, and for $z\ge 1$
the product is at least $z^{1-\hat{\alpha}/2}\ge 1$. By the member
$C_I'\ge 8C(\kappa)(10C_MM)^{5/2}$ of \eqref{5.5:eq-weight-constants-c},
\begin{equation*}
\tfrac12C_I'\mathfrak{n}\nu^3z^2(1+z)t^{4+\hat{\alpha}}\ \ge\ 4C(\kappa)\mathfrak{n}.
\end{equation*}
In the same way,
\begin{equation*}
\nu^3(1+z)t^{\hat{\alpha}}\ \ge\ (10C_MMz)^{-\hat{\alpha}/2}(1+z)\ \ge\ (10C_MM)^{-1/2}.
\end{equation*}
So by the member $C_I'\ge 4C'(10C_MM)^{1/2}$, the other half satisfies
\begin{equation*}
\tfrac12C_I'\mathfrak{n}z^2t^4\cdot\nu^3(1+z)t^{\hat{\alpha}}\ \ge\ 2C'\mathfrak{n}z^2t^4 .
\end{equation*}
Together, $C_I'\mathfrak{n}\nu^3z^2(1+z)t^{4+\hat{\alpha}}$ dominates the trivial bound. The
constant $C_I'$ does not depend on $\Theta$.

In the complementary regime $C_MDt^2z\le\frac12$, \eqref{5.5:eq-near-unit-rows-2} gives
$S^{\star}_v\ge 2>1$, so a read core is box-read there and its entry is that of step (ii), with
the rim of step (iii). The two regimes leave no gap.

The per-cell sums over own-read units do not depend on the field, by the field-free column bound.
They are therefore contained in the constant $\kappa_R$ of
Definition~\ref{5.5:def-weight-constants}.

\smallskip\noindent\emph{(v) Monotonicity in $\Theta$.}
Every row is non-decreasing in $\Theta$ and, since $\Theta\ge 1$, is at most $\Theta^a$ times its
value at $\Theta=1$, with $a\le 3$:
\begin{itemize}
\item Core: since $z=\Theta\rho_{n,j}$,
\begin{equation*}
z^2(1+z)=\Theta^2\rho_{n,j}^2(1+\Theta\rho_{n,j})\le\Theta^3\rho_{n,j}^2(1+\rho_{n,j}).
\end{equation*}
\item Short single: put
\begin{equation*}
x_v:=\Bigl(\sup_{b\subset\mathfrak{r}_v}2d(1+|b-y|)\mathfrak{a}_nt^2/\alpha_{3,j}\Bigr)^2 .
\end{equation*}
This is free of $\Theta$, and by \eqref{5.5:eq-windowed-source-b} we have
$S^{\star-2}_v=\Theta^2x_v$. Hence
\begin{equation*}
2\mathfrak{n}_v\min\{1,S^{\star-2}_v\}=2\mathfrak{n}_v\min\{1,\Theta^2x_v\}
\le\Theta^2\cdot2\mathfrak{n}_v\min\{1,x_v\}.
\end{equation*}
\item Wilson pieces: degree $2$.
\item Own-read units and remainders: degree $0$.
\end{itemize}
Comparing each row, under the filing that its own $\Theta$ prescribes, with its value at
$\Theta=1$ gives the degree bound of \eqref{5.5:eq-weight-d}. In particular,
\begin{equation*}
\mathfrak{N}^P(v;c_Ww)\le w^3\,\mathfrak{N}^P(v;c_W)\qquad (w\ge 1).
\end{equation*}
The Cauchy-integral bound on the filed discs is linear in $\mathfrak{N}(v)$, and its prefactors do
not depend on $\Theta$. The near terms at a block are grouped into the total by the grouping rule
of the sourced stage for the total at a block. By the own-read membership hypothesis, the own-read
pieces remain summands of the one near element.

\smallskip\noindent\emph{The other pieces.}
The far pieces are bounded by the far-piece clause of the sourced-stage estimate, with the far
terms of the second kind read on the configuration prescribed by that clause.

The remainders of clause (b$'$) of Lemma~\ref{5.5:lem-weight} are those named in
Definition~\ref{5.5:def-windowed-source}. For
them, membership, size and sum follow from two hypotheses of Lemma~\ref{5.5:lem-weight}: the
remainder hypothesis for box-read units, at the units of the rows $\flat$, and the remainder
hypothesis for the third class, at the third-class pieces. There are no remainders at own-read
units.

Finally, the pieces at blocks meeting the collar $R_{\mathsf{n}}$ read no cell of
$\mathcal{W}$, by clause (c) of Lemma~\ref{5.5:lem-weight}.
\end{proof}

\begin{proof}[Proof of \eqref{5.5:eq-weight-e} in Lemma~\ref{5.5:lem-weight}: weighted cells within
reach of a block]\label{5.5:prf-weighted-cells}
We work in the setting of Definition~\ref{5.5:def-windowed-source}, with the levels and weights of
Notation~\ref{5.5:not-weight-notation}. Thus $\ell(x)$ is the current level of a point or cell $x$ and
$w$ is the actual weight. The set $\mathcal{W}$ consists of the inherited raised cells of the live
arrests of earlier steps. The block $\square$ lies in zone $n_0$, its own cell $\Delta_0$ has level
$n_0$, and for $y\in\tilde{\square}^2$ we put $n:=\min\{\ell(y),n_0\}$. Finally, $\mathfrak{R}(\square)$
is the set of cells within one own unit of $\tilde{\square}^2$, each of which meets
$\tilde{\square}^3$.

For a level $m$, $\operatorname{dist}_m$ denotes distance measured in units of $\ell_m=L^m\eta$. It
does not increase when $m$ increases.

An arrest or preparation $\mathfrak{c}$ comes with the following data: its region $X_{\mathfrak{c}}$,
its steps $h_{\mathfrak{c}}$ and $k_{\mathfrak{c}}$, and its level $k_0(\mathfrak{c})$. Recall from
\eqref{5.5:eq-weight-e} in Lemma~\ref{5.5:lem-weight} that $t_{\mathfrak{c}}=k_{\mathfrak{c}}$ if
$\mathfrak{c}$ is an arrest and $t_{\mathfrak{c}}=k$ if $\mathfrak{c}$ is the running or a sibling
preparation of the present operation $\mathbf{R}_k$, and that $t(\Delta)=t_{\mathfrak{c}}$ for a
weighted cell $\Delta$, where $\mathfrak{c}$ is the arrest or preparation whose weight $\Delta$
carries. In particular $t(\Delta)=k_{\mathfrak{a}}$ if $\Delta$ is one
of the cells $W_{\mathfrak{a}}$ of a live arrest $\mathfrak{a}$, on which $w=w_{\mathfrak{a}}$.
We write $\check{R}_m$ for Ba{\l}aban's separation radius at level $m$, and $\check{R}$ for the
separation radius between disjoint regions; it satisfies $\check{R}\ge 38$.

\emph{Step one: the separation claim.} We first prove, for every point $y$,
\begin{equation}\label{5.5:eq-weighted-cells-a}
\mathfrak{a}\ \text{a live arrest},\quad m:=\ell(y)>k_{\mathfrak{a}}
\quad\Longrightarrow\quad \operatorname{dist}_m(y,X_{\mathfrak{a}})\ge 3 .
\end{equation}
The point $y$ lies in the current small-field region $\Omega_m$. The region of $\mathfrak{a}$ satisfies
$X_{\mathfrak{a}}\subset Z_{k_{\mathfrak{a}}}$, and $Z_{k_{\mathfrak{a}}}\subset Z_{m-1}$ by the
monotonicity of the large-field regions in the level, part of the level-and-zone statement of the
scaffold, since $k_{\mathfrak{a}}\le m-1$. We distinguish two cases according to how $\Omega_m$ near
$y$ was made.

If the operation $\mathbf{T}_m$ made $\Omega_m$ near $y$, then by Ba{\l}aban's construction of
$\mathbf{T}_m$ in \cite{B-CMP119} the point $y$ lies at level-$m$ distance at least $8\check{R}_m$ from
$Z_{m-1}\supset X_{\mathfrak{a}}$. In particular this distance is at least $3$.

Otherwise, by the statement on arrests of attempts, exactly one attempt $\mathfrak{c}$ re-made
$\Omega_m$ near $y$. Its steps satisfy
$h_{\mathfrak{c}}<m\le k_{\mathfrak{c}}$, and $y\in X_{\mathfrak{c}}$. We compare $X_{\mathfrak{a}}$ with
$X_{\mathfrak{c}}$ in three cases.
\begin{itemize}
\item Suppose $X_{\mathfrak{c}}\subsetneq X_{\mathfrak{a}}$. By the ordering statement for nested
regions, the region strictly inside the other belongs to the earlier step, so
$k_{\mathfrak{c}}<k_{\mathfrak{a}}$. Hence
$m\le k_{\mathfrak{c}}<k_{\mathfrak{a}}$, which contradicts $k_{\mathfrak{a}}<m$.
\item Suppose $X_{\mathfrak{a}}$ and $X_{\mathfrak{c}}$ are disjoint. By the separation statement for
disjoint regions, they lie at distance at
least $\check{R}$ from each other. Since $y\in X_{\mathfrak{c}}$, the distance from $y$ to
$X_{\mathfrak{a}}$ is at least $\check{R}\ge 3$.
\item Suppose $X_{\mathfrak{a}}\subset X_{\mathfrak{c}}$. By the ordering statement for nested regions,
the inner region was made no later than the outer one began, and the beginning step of a preparation
precedes its level $k_0$, by the definition of $N_0$ in \cite{B-CMP122a}. Therefore
\begin{equation*}
k_{\mathfrak{a}}\le h_{\mathfrak{c}}<k_0(\mathfrak{c}),
\end{equation*}
and $X_{\mathfrak{a}}\subset Z'_{k_0(\mathfrak{c})}$. The level-$m$ distance from $y$ to
$Z'_{k_0(\mathfrak{c})}$ is bounded below in three sub-cases:
\begin{itemize}
\item it is at least $3\check{R}_m$ if $m\le k_0(\mathfrak{c})$;
\item it is at least $L^{k_{\mathfrak{c}}-m}\ge L$ if $k_0(\mathfrak{c})<m<k_{\mathfrak{c}}$;
\item it is at least $3$ if $m=k_{\mathfrak{c}}$.
\end{itemize}
In the last sub-case the bound holds because $Z''_{k_{\mathfrak{c}}}$ is the enlargement
$Z'^{\sim 3}_{k_0(\mathfrak{c})}$ of $Z'_{k_0(\mathfrak{c})}$ by three layers \cite{B-CMP122a}. Since
$\check{R}_m\ge 1$ and $L>3$, the distance is at least $3$ in each sub-case.
\end{itemize}
This proves \eqref{5.5:eq-weighted-cells-a}.

We draw the first consequence. Let $\Delta\in\mathcal{W}$, $y\in\tilde{\square}^2$ and
$\operatorname{dist}_n(\Delta,y)<1$. Then $\Delta$ is one of the cells $W_{\mathfrak{a}}$ of a live
arrest $\mathfrak{a}$, and these lie in $X_{\mathfrak{a}}$. Since $n\le\ell(y)$, the bound
$\operatorname{dist}_n(\Delta,y)<1$ gives $\operatorname{dist}_{\ell(y)}(\Delta,y)<1$, hence
$\operatorname{dist}_{\ell(y)}(y,X_{\mathfrak{a}})<3$. By \eqref{5.5:eq-weighted-cells-a} this forces
$\ell(y)\le k_{\mathfrak{a}}$, that is,
\begin{equation}\label{5.5:eq-weighted-cells-1}
t(\Delta)=k_{\mathfrak{a}}\ge\ell(y).
\end{equation}
If $\ell(y)\ge n_0-1$, then \eqref{5.5:eq-weighted-cells-1} gives $t(\Delta)\ge n_0-1$. It remains to
treat blocks whose set $\tilde{\square}^2$ holds a point of level at most $n_0-2$.

\emph{Step two: the thin stratum.} Let $\tilde{\square}^2$ hold a point of level at most $n_0-2$. This
point lies within $3$ level-$n_0$ units of the own cell $\Delta_0$, which has level $n_0$. By the
level-and-zone statement of the scaffold, levels change
by at most one between neighbouring cells, so zone $n_0-1$ is crossed within $3$ such units.

By the level-and-zone statement, the zones of a preparation have widths, in level-$n_0$ units, at
least $8\check{R}_{n_0}$, at least $3\check{R}_{n_0}$, at least $L$ (a $k_0$-zone) and at least $1$
(a thin stratum), according to the kind of zone.
Of these, only a thin stratum can have width at most $3$. Hence the piece crossed is a thin stratum of
a preparation $\mathfrak{c}$, for which
\begin{equation*}
k_0(\mathfrak{c})\le n_0-2,\qquad t_{\mathfrak{c}}\ge n_0,\qquad
\Delta_0\subset X_{\mathfrak{c}}\setminus Z''^{(\mathfrak{c})}_{n_0}.
\end{equation*}
Moreover $\tilde{\square}^2\subset X_{\mathfrak{c}}$. Indeed, by the collar statement for
preparations, the strata lie behind the collar of
$X_{\mathfrak{c}}$ of width $2\check{R}_{t_{\mathfrak{c}}}$, while $\tilde{\square}^2$ reaches only $3$
units from $\Delta_0$, and $3<\check{R}$.

\emph{Step three: uniqueness.} Suppose two preparations have thin strata inside $\tilde{\square}^2$. By
step two, each of their regions contains $\Delta_0$, so the regions are nested. Let the inner one have
step $t_{\mathrm{inner}}$, and the outer one steps $h_{\mathrm{outer}}$ and $k_0(\mathrm{outer})$. Step
two gives $t_{\mathrm{inner}}\ge n_0$ and $k_0(\mathrm{outer})\le n_0-2$. By the ordering statement for
nested regions, the inner region was made no
later than the outer one began, and $h_{\mathrm{outer}}<k_0(\mathrm{outer})$ by the definition of $N_0$
in \cite{B-CMP122a}.
Together these give
\begin{equation*}
n_0-1<t_{\mathrm{inner}}\le h_{\mathrm{outer}}<k_0(\mathrm{outer})<n_0-1,
\end{equation*}
which is absurd. Hence $\mathfrak{c}$ is unique.

\emph{Step four: reach and the level window.} We measure from $\Delta_0$ in level-$n_0$ units.
\begin{itemize}
\item If $n_0<t_{\mathfrak{c}}$, the set $Z'_{k_0(\mathfrak{c})}$ lies
$L^{t_{\mathfrak{c}}-n_0}\ge L>3$ units away.
\item If $n_0=t_{\mathfrak{c}}$, it lies exactly $3$ units away, by the enlargement identity of
\cite{B-CMP122a} recalled in step one.
\end{itemize}
Hence $\tilde{\square}^2$ meets $Z'_{k_0(\mathfrak{c})}$ at most along its boundary. Cells of
$X_{\mathfrak{c}}$ off the interior of $Z'_{k_0(\mathfrak{c})}$ have level at least $k_0(\mathfrak{c})$
\cite[(1.11)]{B-CMP122a}. With the upper bound $n_0+1$ from the facts stated with
Definition~\ref{5.5:def-data-factor}, every cell inside $\tilde{\square}^2$ has level in
$[k_0(\mathfrak{c}),n_0+1]$.

Let $y\in\tilde{\square}^2$. Then $\ell(y)\ge k_0(\mathfrak{c})\ge h_{\mathfrak{c}}+1$. Let
$\mathfrak{a}$ be a live arrest with a cell within one unit of $y$. Then
\eqref{5.5:eq-weighted-cells-a}, applied as in the derivation of \eqref{5.5:eq-weighted-cells-1},
gives
\begin{equation}\label{5.5:eq-weighted-cells-2}
k_{\mathfrak{a}}\ge\ell(y)\ge h_{\mathfrak{c}}+1.
\end{equation}

\emph{Step five: a live arrest against the preparation.} Let $\mathfrak{a}$ be as at the end of step
four. The comparison below uses the following inputs:
\begin{itemize}
\item the comparison of a live arrest with another live arrest, if $\mathfrak{c}$ is a live arrest;
\item the separation statement for disjoint regions;
\item the geometry statements of the comparison lemma's region, if $\mathfrak{c}$ is the running or a
sibling preparation of the comparison lemma's region.
\end{itemize}
We consider three cases.
\begin{itemize}
\item Suppose $X_{\mathfrak{a}}$ and $X_{\mathfrak{c}}$ are disjoint. They then lie at least
$\check{R}\ge 38$ apart. But $\tilde{\square}^2\subset X_{\mathfrak{c}}$ has diameter at most $6$,
and $\mathfrak{a}$ has a cell within one unit of it. This is impossible.
\item Suppose $X_{\mathfrak{a}}\subsetneq X_{\mathfrak{c}}$. By the ordering statement for nested
regions, the inner region was made no later than
the outer one began, so $k_{\mathfrak{a}}\le h_{\mathfrak{c}}$. This contradicts
\eqref{5.5:eq-weighted-cells-2}.
\item Suppose $X_{\mathfrak{c}}\subsetneq X_{\mathfrak{a}}$. Then
$k_{\mathfrak{c}}\le h_{\mathfrak{a}}$. By the layer statement for raised cells, the raised cells of
$\mathfrak{a}$ lie at least three layers
of cubes of side $LM\check{R}_{h_{\mathfrak{a}}+1}$ away from
$Z_{h_{\mathfrak{a}}}\supset X_{\mathfrak{c}}$. If $\mathfrak{c}$ is the running or a sibling
preparation of the comparison lemma's region, we argue instead from the treated component $C$ of that
region. Let $\mathsf{C}_C$ be its core, $h$ the step at which the preparation treating $C$ began, and
$Y^{\sharp}$ the set fixed in the geometry statements of that region. By those statements the core
$\mathsf{C}_C$ contains $Y^{\sharp}$, and every collar $R_\nu$ of $C$ satisfies
\begin{equation*}
\operatorname{dist}_{h+1}(Y^{\sharp},R_\nu)\ge\tfrac52\check{R}_{h+1}.
\end{equation*}
In either case no raised cell of $\mathfrak{a}$ lies at distance less than $1$ from
$\tilde{\square}^2\subset X_{\mathfrak{c}}$, which is a contradiction.
\end{itemize}
Hence $\mathfrak{a}=\mathfrak{c}$. For $\Delta\in\mathcal{W}$ within one own unit of
$y\in\tilde{\square}^2$ this gives
\begin{equation}\label{5.5:eq-weighted-cells-3}
t(\Delta)=t_{\mathfrak{c}}\ge n_0 .
\end{equation}

The same conclusion holds for every weighted cell within one own unit of $\tilde{\square}^2$. By the
layer statement for raised cells, the cells raised during the present run of $\mathbf{R}_k$ belong to
the preparation of this run, and by the anchor statement at the current step this preparation is again
$\mathfrak{c}$, with $t_{\mathfrak{c}}=k$. A completed preparation's cells carry weight $1$
(Notation~\ref{5.5:not-weight-notation}; after a completed operation the block partition on $Z$ is that
of \cite[(1.33)]{B-CMP122b}), and so are not weighted.

Steps two to five give the second assertion of \eqref{5.5:eq-weight-e}. Together with
\eqref{5.5:eq-weighted-cells-1} and \eqref{5.5:eq-weighted-cells-3}, they give
$t(\Delta)\ge\ell(y)$ and $t(\Delta)\ge n_0-1$ for every $\Delta\in\mathcal{W}$ and
$y\in\tilde{\square}^2$ with $\operatorname{dist}_n(\Delta,y)<1$.

\emph{Step six: old and exceptional blocks.} Let $\Delta\in\mathcal{W}\cap\mathfrak{R}(\square)$. Then
$\Delta$ meets $\tilde{\square}^3\subset\tilde{\square}^5$. By the facts stated with
Definition~\ref{5.5:def-data-factor}, either $\square\in\mathfrak{O}^+$ or $\square$ is exceptional.

In the second case the own cell of $\square$ lies off
\begin{equation*}
Z''_k=Z\cap Z'^{\sim 3}_{k_0}.
\end{equation*}
Consequently every $y\in\tilde{\square}^2$ satisfies
\begin{equation*}
\ell(y)\ge k_0(k)>h\ge k_{\mathfrak{a}}
\end{equation*}
for every live arrest $\mathfrak{a}$ inside the component $C''$. Here $C''$ is the component treated
by the present preparation of $\mathbf{R}_k$, which began at step $h$ and has level $k_0(k)$. The last
inequality holds by the ordering statement for nested regions, since the region of $\mathfrak{a}$ lies
inside that of this preparation. Live arrests
in other components lie at least $\check{R}_k$ away. So \eqref{5.5:eq-weighted-cells-a} forbids any
contact within one unit between $\tilde{\square}^2$ and a cell of $\mathcal{W}$. Hence an exceptional
block reads no cell of $\mathcal{W}$ through $\mathfrak{R}(\square)$.

Finally, a residual cell in the sense of the facts stated with Definition~\ref{5.5:def-data-factor} has
$t\le n_0-3$, which contradicts $t(\Delta)\ge n_0-1$, proved above. Hence no window block, the
exceptional ones included, reads a cell
of $\mathcal{W}$ through $\mathfrak{R}(\square)$. A cell of $\mathcal{W}$ within one own unit of
$\tilde{\square}^2$ can therefore occur only at a block $\square\in\mathfrak{O}^+$. This completes the
proof of \eqref{5.5:eq-weight-e}.
\end{proof}

\begin{proof}[Proof of the count-times-weight bound and of the data-factor bound]
\label{5.5:prf-count-weight}
We prove the two parts of the count-times-weight bound \eqref{5.5:eq-weight-f} of
Lemma~\ref{5.5:lem-weight}, and then clause~(d) of that lemma, \eqref{5.5:eq-weight-g}. We
write $\mathsf{L}_0=L_{\circ}$ for the weight base of
Notation~\ref{5.5:not-weight-notation}. Recall from Definition~\ref{5.5:def-weight-constants}
that
\begin{equation*}
\mathfrak{K}(\square)=\sum\bigl\{w^{\vee}(\Delta')^3:\ \Delta'\text{ a unit $n$-cell inside }
\tilde{\square}^2\bigr\}
\end{equation*}
at every block $\square$, where $w^{\vee}(\Delta')$ is the largest weight of a cell within one
own unit of $\Delta'$. A cell $\Delta$ of raise $\mathsf{r}$ has weight
$w(\Delta)\le\mathsf{L}_0^{2\mathsf{r}}$,
and $\mathsf{L}_0^6\le L^3/27$ by \eqref{5.5:eq-weight-notation-a}; in particular
$\mathsf{L}_0^2\le L$. In the proof of the two parts of \eqref{5.5:eq-weight-f} the block
$\square$ is such that $\tilde{\square}^5$ meets $\mathcal{W}$; then $\square\in\mathfrak{O}^+$ or
$\square$ is exceptional, and $\varrho=\varrho(\square)\ge\check{R}_k-4$. The constant of the
target bound \eqref{5.5:eq-data-factor-a} of Definition~\ref{5.5:def-data-factor}, and the second
member of its defining maximum in \eqref{5.5:eq-weight-constants-b}, the constant introduced in
Definition~\ref{5.5:def-data-factor}, are written $C^{\mathrm{dat}}$ and $C^{\mathrm{dat}}_{pr}$
respectively; in particular $C^{\mathrm{dat}}\ge C^{\mathrm{dat}}_{pr}$.

\emph{Part~(i).} Suppose that $\tilde{\square}^2$ holds no point of level $\le n_0-2$. The unit
$n$-cells inside $\tilde{\square}^2$, which has side $6$, are then of three kinds: cells of level
$n_0-1$, of which there are at most $(6L)^d$; cells of level $n_0$, at most $6^d$; and level-$n_0$
sub-cubes of cells of level $n_0+1$, at most $6^d$. Their number is therefore at most
\begin{equation}\label{5.5:eq-count-weight-1}
(6L)^d+6^d+6^d\le 2(7L)^d .
\end{equation}
For each such $\Delta'$, the quantity $w^{\vee}(\Delta')$ is a maximum over cells within one own
unit of $\Delta'$, and two kinds of weighted cells can occur there.

A cell raised during the present operation $\mathbf{R}_k$ and lying within one own unit of
$\Delta'$ meets $\tilde{\square}^5$, so its raise $\mathsf{r}$ is at most $s(\square)$, which is,
by Definition~\ref{5.5:def-data-factor}, the largest raise of this $\mathbf{R}_k$ over cells
meeting $\tilde{\square}^5$; hence such a cell $\Delta$ has
$w(\Delta)^3\le\mathsf{L}_0^{6\mathsf{r}}\le L_{\circ}^{6s(\square)}$.

The other weighted cells there are inherited: let $\Delta\in\mathcal{W}$ lie within one own unit
of $\Delta'$. By the geometry clause \eqref{5.5:eq-weight-e} of
Lemma~\ref{5.5:lem-weight}, the step $t:=t(\Delta)$ satisfies $t\ge n_0-1$; by the bound on the
raises of inherited cells, $\mathsf{r}(\Delta)\le N_0(t)$. Hence
$w(\Delta)\le\mathsf{L}_0^{2N_0(t)}$ and,
as $\mathsf{L}_0^2\le L$, $w(\Delta)^3\le L^{3N_0(t)}$. The depth bound for the data factor gives
$L^{3N_0(t)}\le L^{12}(\check{\mathsf{T}}_k+3+c_{k-t})^{6r_{pr}}$. Since $\tilde{\square}^5$ meets
$\mathcal{W}$, the level estimate $k-n_0+1\le N_*+\varrho$ of
Definition~\ref{5.5:def-data-factor} holds, and the depth bound $N_*+\varrho\le\mathsf{n}(\varrho)$
together with the bound $\check{\mathsf{T}}_k\le\varrho+4$ on the depth parameter at old and
exceptional blocks give the two level estimates
\begin{equation*}
k-n_0+1\le N_*+\varrho\le\mathsf{n}(\varrho),\qquad \check{\mathsf{T}}_k\le\varrho+4 ;
\end{equation*}
and $t\ge n_0-1$ gives $k-t\le k-n_0+1\le\mathsf{n}(\varrho)$. Since
$c_j=\log(1+3\lambda j+2c_2)$ is increasing in $j$, $c_{k-t}\le c_{\mathsf{n}(\varrho)}$.
Therefore
\begin{equation}\label{5.5:eq-count-weight-2}
\begin{aligned}
w(\Delta)^3&\le L^{3N_0(t)}\le L^{12}\bigl(\check{\mathsf{T}}_k+3+c_{k-t}\bigr)^{6r_{pr}}
\le L^{12}\bigl(\varrho+7+c_{\mathsf{n}(\varrho)}\bigr)^{6r_{pr}}\\
&\le C^{\mathrm{dat}}_{pr}\,e^{\varrho-2},
\end{aligned}
\end{equation}
the last inequality by the definition
\begin{equation*}
C^{\mathrm{dat}}_{pr}=\sup_{\varrho'\ge1}L^{12}\bigl(\varrho'+7+c_{\mathsf{n}(\varrho')}\bigr)^{6r_{pr}}
e^{2-\varrho'}
\end{equation*}
of Definition~\ref{5.5:def-data-factor}, since $\varrho\ge1$. Consequently
\begin{equation*}
w^{\vee}(\Delta')^3\le\max\bigl\{L_{\circ}^{6s(\square)},\,C^{\mathrm{dat}}_{pr}e^{\varrho-2}\bigr\}
\end{equation*}
for every unit $n$-cell $\Delta'$ inside $\tilde{\square}^2$, and summing over the at most
$2(7L)^d$ such cells counted in \eqref{5.5:eq-count-weight-1} gives part~(i) of
\eqref{5.5:eq-weight-f}.

\emph{Part~(ii).} Suppose that $\tilde{\square}^2$ holds a point of level $\le n_0-2$, and let
$\mathfrak{c}$ be the one preparation given by the geometry clause \eqref{5.5:eq-weight-e}, of step
$t:=t_{\mathfrak{c}}\ge n_0$ ($t=k$ at an exceptional block). By that clause the cells inside
$\tilde{\square}^2$ have levels $m\in[k_0,n_0+1]$, where $k_0:=k_0(\mathfrak{c})=t-N_0(t)$ is the
lowest level of the preparation, $N_0(t)$ being its depth at step $t$ as in \cite{B-CMP122a}, and
every weighted cell within one own
unit of $\tilde{\square}^2$ is a cell of $\mathfrak{c}$. By the bound on the raises of the cells of
a preparation, such a cell $\Delta$ at level $i$ has raise at most $\max\{1,i-k_0\}$, and raise at
most $N_0(t)$, so that
\begin{equation*}
w(\Delta)\le\mathsf{L}_0^{2\max\{1,i-k_0\}}\le\mathsf{L}_0^{2N_0(t)} .
\end{equation*}
Cells within one level-$m$ unit of a level-$m$ cell have level at most $m+1$, by the level property
of the sequence of block partitions. Hence a unit cell $\Delta'$ of level $m\in[k_0,n_0]$ has
$w^{\vee}(\Delta')^3\le\mathsf{L}_0^{6(m+1-k_0)}$, and there are at most $(7L^{n_0-m})^d$ of them;
for the at most $7^d$ unit $n_0$-cells inside cells of level $n_0+1$,
$w^{\vee}\le\mathsf{L}_0^{2N_0(t)}$. Since $t\ge n_0$, we have $n_0-k_0=n_0-t+N_0(t)\le N_0(t)$.
Writing $m=k_0+q$ in the first sum, we obtain
\begin{equation}\label{5.5:eq-count-weight-a}
\begin{aligned}
\mathfrak{K}(\square)
&\le\sum_{m=k_0}^{n_0}7^dL^{d(n_0-m)}\mathsf{L}_0^{6(m+1-k_0)}+7^d\mathsf{L}_0^{6N_0(t)}\\
&\le 7^d\mathsf{L}_0^6L^{d(n_0-k_0)}\sum_{q\ge0}\bigl(\mathsf{L}_0^6L^{-d}\bigr)^q
+7^d\mathsf{L}_0^6\,\mathsf{L}_0^{6(N_0(t)-1)}\\
&\le 7^d\mathsf{L}_0^6L^{dN_0(t)}\Bigl[\bigl(1-(27L)^{-1}\bigr)^{-1}+1\Bigr]
\le 3\cdot7^d\mathsf{L}_0^6L^{dN_0(t)} .
\end{aligned}
\end{equation}
In the third line we used $n_0-k_0\le N_0(t)$; the inequality
$\mathsf{L}_0^6L^{-d}\le(27L)^{-1}$, which is \eqref{5.5:eq-weight-notation-a} with $d=4$,
so that the geometric series is at most $(1-(27L)^{-1})^{-1}$; and
\begin{equation*}
\mathsf{L}_0^{6(N_0(t)-1)}\le\bigl(L^3/27\bigr)^{N_0(t)-1}\le L^{dN_0(t)} .
\end{equation*}
In the last step, $(1-(27L)^{-1})^{-1}\le 27/26\le2$; the factor $3$ does not depend on $L$.

It remains to bound the factor $L^{dN_0(t)}$ by a constant times $e^{\varrho-1}$. Put
$y:=\check{R}_{k_0(t)+1}=L^{N_0(t)-1}$, so that $L^{dN_0(t)}=L^dy^d$. By
\cite[(2.5)]{B-CMP119}, $R(x)$ is the least power of $L$ which is $\ge(\log x)^{r_{pr}}$; hence
\begin{equation*}
y<L\max\bigl\{1,\check{\mathsf{T}}_{k_0(t)+1}^{\,r_{pr}}\bigr\}.
\end{equation*}
Since $k-k_0(t)-1=(k-t)+(N_0(t)-1)=(k-t)+\log_Ly$ and the sequence $c$ is sub-additive, the
comparison of the depth parameters at two levels gives
$\check{\mathsf{T}}_{k_0(t)+1}\le\check{\mathsf{T}}_k+1+c_{k-t}+c_{\log_Ly}$. With
$\mathsf{b}_t:=\check{\mathsf{T}}_k+1+c_{k-t}$ this reads $y<L(\mathsf{b}_t+c_{\log_Ly})^{r_{pr}}$. If
$c_{\log_Ly}\le\mathsf{b}_t$, then $y<L(2\mathsf{b}_t)^{r_{pr}}$. Otherwise
$y<L(2c_{\log_Ly})^{r_{pr}}$, and then
$y\le y_*$ by the definition $y_*=\sup\{y\ge1:\ y\le L(2c_{\log_Ly})^{r_{pr}}\}$ of
Definition~\ref{5.5:def-weight-constants}. In both cases
\begin{equation}\label{5.5:eq-count-weight-3}
y^d\le y_*^d+L^d2^{dr_{pr}}\mathsf{b}_t^{dr_{pr}} .
\end{equation}
At an old block, $t\ge n_0$ and the first level estimate give $k-t\le k-n_0+1\le\mathsf{n}(\varrho)$;
at an exceptional block $t=k$. In both cases $\check{\mathsf{T}}_k\le\varrho+4$ and
$c_{k-t}\le c_{\mathsf{n}(\varrho)}$, so that $\mathsf{b}_t\le\varrho+5+c_{\mathsf{n}(\varrho)}$.
Since $\varrho\ge1$, we have $y_*^d\le y_*^de^{\varrho-1}$ and
\begin{equation*}
\mathsf{b}_t^{dr_{pr}}\le\sup_{\varrho'\ge1}\bigl(\varrho'+5+c_{\mathsf{n}(\varrho')}\bigr)^{dr_{pr}}
e^{1-\varrho'}\cdot e^{\varrho-1}.
\end{equation*}
By \eqref{5.5:eq-count-weight-3} and the definition of $C^{tile}$ in
\eqref{5.5:eq-weight-constants-b}, therefore,
\begin{equation}\label{5.5:eq-count-weight-4}
\begin{aligned}
3\cdot7^d\mathsf{L}_0^6L^{dN_0(t)}&=3\cdot7^d\mathsf{L}_0^6L^dy^d
\le3\cdot7^d\mathsf{L}_0^6L^d\bigl[y_*^d+L^d2^{dr_{pr}}\mathsf{b}_t^{dr_{pr}}\bigr]\\
&\le C^{tile}e^{\varrho-1}.
\end{aligned}
\end{equation}
Together with \eqref{5.5:eq-count-weight-a} this is part~(ii) of \eqref{5.5:eq-weight-f}.

\emph{Clause~(d).} Let $\square$ be any block. The total at $\square$ which enters the definition
of $\mathfrak{d}(\square)$ in Definition~\ref{5.5:def-data-factor} is the resummation, by the
linear resummation lemma for the near-unit table, of the table of near-unit sizes. This
resummation is linear, its scaffold coefficients are non-negative, and the prefactors of the discs
\eqref{5.5:eq-near-unit-classes-b} are proportional to $\vartheta(\square)$. We group the near units
$v$ according to the unit $n$-cell $\Delta'$ which contains the anchor $y_v$. By the degree bound
of the table \eqref{5.5:eq-weight-d},
\begin{equation*}
\mathfrak{N}^P\bigl(v;c_Ww^{\vee}(y_v)\bigr)\le c_W^3\,w^{\vee}(\Delta')^3\,\mathfrak{N}^P(v;1).
\end{equation*}
By the definition of $\kappa_R^{act}$ in Definition~\ref{5.5:def-weight-constants} and
$\kappa_R\ge\kappa_R^{act}$, \eqref{5.5:eq-weight-constants-a}, the group of $\Delta'$, formed with
the table $\mathfrak{N}^P(\cdot\,;1)$, contributes at most $\kappa_R\mathfrak{N}_0(\square)$. By
linearity and non-negativity of the coefficients, the group of $\Delta'$ formed with the actual
sizes contributes at most $c_W^3w^{\vee}(\Delta')^3\kappa_R\mathfrak{N}_0(\square)$. Summing over
the unit $n$-cells $\Delta'$ inside $\tilde{\square}^2$, and using the definition
$\mathsf{l}^2(\square)=\max\{L_{\circ}^{6s(\square)},\kappa_Rc_W^3\mathfrak{K}(\square)\}$ of
\eqref{5.5:eq-weight-constants-b}, we obtain that at every block the total at $\square$ is at most
the left member of
\begin{equation}\label{5.5:eq-count-weight-5}
\kappa_Rc_W^3\,\mathfrak{K}(\square)\,\mathfrak{N}_0(\square)
\le\mathsf{l}^2(\square)\,\mathfrak{N}_0(\square).
\end{equation}

Now let $\tilde{\square}^5$ meet $\mathcal{W}$. By Definition~\ref{5.5:def-data-factor},
$\mathfrak{d}(\square)$ is the larger of $1$ and the quotient of the total at $\square$ by
$L_{\circ}^{6s(\square)}\mathfrak{N}_0(\square)$. Both are at most
$\mathsf{l}^2(\square)/L_{\circ}^{6s(\square)}$, the first because
$\mathsf{l}^2(\square)\ge L_{\circ}^{6s(\square)}$ and the second by \eqref{5.5:eq-count-weight-5};
that is, $L_{\circ}^{6s(\square)}\mathfrak{d}(\square)\le\mathsf{l}^2(\square)$. Next,
$\mathsf{l}^2(\square)\le\kappa_Rc_W^3\mathfrak{K}^{\sharp}(\square)
=L_{\circ}^{6s(\square)}\mathfrak{d}^{\sharp}(\square)$, which we check member by member. First,
$\mathfrak{K}\le\mathfrak{K}^{\sharp}$ by \eqref{5.5:eq-weight-f}. Second,
$L_{\circ}^{6s(\square)}\le\mathfrak{K}^{\sharp}(\square)\le\kappa_Rc_W^3\mathfrak{K}^{\sharp}(\square)$,
the second inequality because $\kappa_R\ge\kappa_R^{act}\ge1$ and $c_W\ge1$. The first holds in
case~(i) because
\begin{equation*}
\mathfrak{K}^{\sharp}(\square)=2(7L)^d\max\bigl\{L_{\circ}^{6s(\square)},\,
C^{\mathrm{dat}}_{pr}e^{\varrho-2}\bigr\}.
\end{equation*}
In case~(ii), $s(\square)=0$ when $\square\in\mathfrak{O}^+$ (Definition~\ref{5.5:def-data-factor}),
so that $L_{\circ}^{6s(\square)}=1\le\mathfrak{K}^{\sharp}(\square)$; at an exceptional block
$s(\square)\le N_0(k)$ and $t=k$, so that
\begin{equation*}
L_{\circ}^{6s(\square)}\le\mathsf{L}_0^6\,\mathsf{L}_0^{6(N_0(k)-1)}\le\mathsf{L}_0^6L^{dN_0(k)}
\le\mathfrak{K}^{\sharp}(\square).
\end{equation*}
Consequently $\mathfrak{d}^{\sharp}(\square)\ge1$ and
$\mathfrak{d}(\square)\le\mathfrak{d}^{\sharp}(\square)$.

By \eqref{5.5:eq-weight-f}, and since $L_{\circ}^{6s(\square)}\ge1$: in case~(i),
\begin{equation*}
\mathfrak{d}^{\sharp}(\square)\le2(7L)^d\kappa_Rc_W^3\max\bigl\{1,\,C^{\mathrm{dat}}_{pr}e^{\varrho-2}\bigr\};
\end{equation*}
in case~(ii), $\mathfrak{d}^{\sharp}(\square)\le\kappa_Rc_W^3C^{tile}e^{\varrho-1}$. By the floor
$\tfrac18\delta_Pw\ge2$ on the width $w$ which enters
$\varrho=1+d^{\wedge}(\square,R_{\mathsf{n}})/w$ in Definition~\ref{5.5:def-data-factor}, a floor
listed with the constants of Definition~\ref{5.5:def-weight-constants},
\begin{equation*}
e^{\varrho-2}\le e^{\varrho-1}=e^{d^{\wedge}(\square,R_{\mathsf{n}})/w}
\le e^{\frac18\delta_Pd^{\wedge}(\square,R_{\mathsf{n}})} .
\end{equation*}
Put $u:=\tfrac18\delta_Pd^{\wedge}(\square,R_{\mathsf{n}})\ge0$. Using $\max\{1,\xi\}\le1+\xi$, in
both cases $\mathfrak{d}^{\sharp}(\square)\le a+be^u$ with $a:=2(7L)^d\kappa_Rc_W^3\ge1$ and
$b:=\kappa_Rc_W^3[2(7L)^dC^{\mathrm{dat}}_{pr}+C^{tile}]$. Since $a\le ae^u$, we have
$a+be^u\le1+(a+b)e^u$ for $a\ge1$, $u\ge0$, and hence
\begin{equation}\label{5.5:eq-count-weight-6}
\begin{aligned}
\mathfrak{d}(\square)\le\mathfrak{d}^{\sharp}(\square)
&\le1+\kappa_Rc_W^3\bigl[2(7L)^d(1+C^{\mathrm{dat}}_{pr})+C^{tile}\bigr]
e^{\frac18\delta_Pd^{\wedge}(\square,R_{\mathsf{n}})}\\
&\le1+C^{\mathrm{dat}}e^{\frac18\delta_Pd^{\wedge}(\square,R_{\mathsf{n}})},
\end{aligned}
\end{equation}
the last inequality because $C^{\mathrm{dat}}$ is at least the third member of its defining maximum,
\eqref{5.5:eq-weight-constants-b}. Where $\tilde{\square}^5$ meets no cell of $\mathcal{W}$,
$\mathfrak{d}(\square)=1$ by Definition~\ref{5.5:def-data-factor}, and the same bound holds
trivially; this gives the target bound \eqref{5.5:eq-data-factor-a} at every block, with
$C^{\mathrm{dat}}\ge C^{\mathrm{dat}}_{pr}$.

All constants used in this proof are fixed before $\gamma$, and the bound on $\mathfrak{d}(\square)$
is of degree zero. The decay rate
$\tfrac18\delta_P$ spent here on $\mathfrak{d}$ and the rate $\tfrac38\delta_P$ spent in the other
decay estimate of Lemma~\ref{5.5:lem-weight} add up to $\tfrac12\delta_P$, the rate of
$\vartheta(\square)$, and the rate $\tfrac18\delta_P$ is spent only once.
\end{proof}

\begin{proof}[Proof of clause~(e) of Lemma~\ref{5.5:lem-weight}]\label{5.5:prf-machine-differences}
Let $\mathsf{A}$ and $\mathsf{B}$ be the two machines on one scaffold of clause~(e) of the weight lemma
(Lemma~\ref{5.5:lem-weight}). Let $\mathcal{K}=\mathcal{K}^{rd}$ be their pair system, with the four
closure properties assumed in that clause and the closure of $\mathcal{K}$ under the move
$\mathsf{M}8$, which holds by construction of $\mathcal{K}^{rd}$. As there, for a function $f$ read in
both machines and $p\in\mathcal{K}$ we write
\begin{equation*}
\delta f(p) := f^{\mathsf{B}}(P^{\mathsf{B}})-f^{\mathsf{A}}(P^{\mathsf{A}}),
\end{equation*}
$\mathsf{S}:=\mathsf{S}^{rd}$ for the $\delta$-size and $\mathsf{a}:=\mathsf{S}^{rd}[\hat{A}]$ for the
$\delta$-size of the Wilson and counterterm functional $\hat{A}$. The index $\mathsf{m}$ runs over
$\{\mathsf{A},\mathsf{B}\}$. Throughout we use the notation fixed in the proof of the rows of the
near-unit table, within the proof of Lemma~\ref{5.5:lem-weight}: the data constants $\rho$ and $z$ of
those rows, the quantity $D$ (with $D\le 5M\nu$), and the constants $C(\kappa)$, $C$ and $C'$; the
letters $E_0$, $\mathfrak{N}_0(\square)$ and the base data $\mathfrak{n}_0$ are those of clauses~(d)
and~(e) of Lemma~\ref{5.5:lem-weight}, and $c_{\alpha}$, $C_R$, $C_2$ are the constants entering the
definition of $\kappa_R^{\times}$ in Definition~\ref{5.5:def-weight-constants}. We prove
\eqref{5.5:eq-weight-h} in the steps below. Step~4$'$ treats
the cores with $\nu\le 2$ and is inserted between Step~4 and Step~5.

\emph{Step 1: one scaffold.} The scaffold is one and the same for both machines. Hence the following
objects coincide in $\mathsf{A}$ and in $\mathsf{B}$:
\begin{itemize}
\item the level function $\ell$, the weight $w$, the family $\mathcal{W}$, the largest weight $w^{\vee}$
and $\Theta$;
\item the scaffold radius $S^{\star}_v$ of \eqref{5.5:eq-near-unit-classes-a}, the assignment of near
units to their classes, and the discs of \eqref{5.5:eq-near-unit-classes-b};
\item the boxes $Q_v$ and $\mathfrak{r}_v$, and the weighted scaffold majorants of
\eqref{5.5:eq-windowed-source-b};
\item the count $\mathfrak{K}(\square)$, the constants $\kappa_R$ of \eqref{5.5:eq-weight-constants-a}
and $c_W$, and the data factors $\mathsf{l}^2(\square)$ and $\mathfrak{d}^{\sharp}(\square)$ of
\eqref{5.5:eq-weight-constants-b}.
\end{itemize}

\emph{Step 2: clause~(a) in each machine.} Clauses (a1)--(a4) of Lemma~\ref{5.5:lem-weight} hold for the
box pair $\mathfrak{b}^{\mathsf{m}}$ of each machine $\mathsf{m}$ with the same constants. So do the
statements of part~(i) of the proof of the rows of the near-unit table, within the proof of
Lemma~\ref{5.5:lem-weight}. This is the hypothesis on the paired data made in the closure theorem
for the pair system. The membership input and the remainder input of
Definition~\ref{5.5:def-weight-inputs} likewise hold in each machine separately. Consequently the two
requirements which the two-machine comparison at the rim places on the paired reads are met at the
common weighted majorants of Step~1. The closure of $\mathcal{K}$ under the paired charts follows from
the closure of $\mathcal{K}$ under the move $\mathsf{M}5$.

\emph{Step 3: the difference principle.} Let $\mu\in\mathfrak{M}^{rd}$ be a move and let
$p\in\mathcal{K}^{rd}$. On the parameter domain $\mathfrak{P}_{\mu}$ of $\mu$ put
\begin{equation}\label{5.5:eq-machine-differences-1}
\delta(\pi) := f^{\mathsf{B}}\bigl(\mu^{\mathsf{B}}(\pi;P^{\mathsf{B}})\bigr)
-f^{\mathsf{A}}\bigl(\mu^{\mathsf{A}}(\pi;P^{\mathsf{A}})\bigr),\qquad \pi\in\mathfrak{P}_{\mu}.
\end{equation}
By clause~(iii) of the closure theorem for the pair system, $\delta$ is holomorphic on
$\mathfrak{P}_{\mu}$. Moreover
\begin{equation*}
|\delta|\le\sup_{\mathcal{K}}|f^{\mathsf{B}}-f^{\mathsf{A}}|,
\end{equation*}
and the right-hand side is the $\delta$-size. At the slot $\circ$ of the Wilson and counterterm pieces,
$\delta$ is entire and $\mathsf{a}\le 2\mathfrak{n}_{\hat{A}}$.

Three tools are used in the one-machine proofs: the Cauchy estimates, the three Schwarz-lemma bounds
(a)--(c), and, for a function $g$ read in one machine, the vanishing $\partial^i g(\pi_0)=0$ of
prescribed derivatives at the base point $\pi_0$ of the parameter domain. All three pass to $\delta$.
The first two apply to every bounded holomorphic function, and $\delta$ is one. The identities hold in
each machine, hence hold for the difference $\delta$ by subtraction.

\emph{Step 4: the places at which $\mathcal{K}$ is read.} The pair system is read at six places; the
seventh place listed in clause~(v)(c) of the closure theorem for the pair system, the closure under the
move $\mathsf{M}8$, is read here inside the fourth of them, the place of the $\circ$-class. We take the
places in turn.

\begin{enumerate}
\item The Cauchy integrals of the Cauchy-integral lemma for near units are taken on the discs of
\eqref{5.5:eq-near-unit-classes-b}. The prefactor of the paired Cauchy row of the two-machine comparison
at the rim is read as the prefactors fixed with those discs. The difference is controlled in two cases:
\begin{itemize}
\item at box pairs, by the closure of $\mathcal{K}$ under $\mathsf{M}2$;
\item at own-read units, by the closure under $\mathsf{M}3$.
\end{itemize}
The values at near levels are treated in the same way.

\item The short-single lemma uses the disc in $\sigma'$. Its chart satisfies
$(\varphi^{\mathsf{m}})'(0)=0$ for each machine $\mathsf{m}$. The difference is controlled by the
closure under $\mathsf{M}5$ with its first splitting.

\item The core lemma has several items, among them the rim item and the families generated from
$\mathbb{1}$. Inside each machine the marginal part cancels at $y$. These items are controlled by the
closure under $\mathsf{M}5$.

\item The $\circ$-class.
\begin{itemize}
\item The $\delta$-size $\mathsf{a}$ is reached through $\mathsf{M}2^{\circ}$.
\item At read cores, that is at cores with $\nu\ge 3$ and at every value of $S^{\star}_v$, the chart of
$\hat{A}$ has radius of order one. This chart is the paired chart of item~(a) of the move $\mathsf{M}5$
at radius $\alpha_3(\circ)$, taken through $\mathsf{M}2^{\circ}$ on the common majorants of part~(i) of
the proof of the rows of the near-unit table. Part~(i) supplies these majorants in each machine. This is
the hypothesis on the $\circ$-slot paired data of the closure theorem.
\item Off that chart two groups of items remain: the block pieces of row~5 at every $\nu$, and the
remainder rows of the counterterm companions at every $\nu$. Both follow from the third-class off-chart
input of Lemma~\ref{5.5:lem-weight} together with the closure of $\mathcal{K}$ under $\mathsf{M}8$.
\item The item of the companions at $\nu\le 2$ whose size is $\mathfrak{n}\mathsf{a}$ is treated in
Step~4$'$.
\end{itemize}

\item The remainders follow from the closure under $\mathsf{M}4$ and under $\mathsf{M}5$ with its fourth
splitting. This is the remainder input of Definition~\ref{5.5:def-weight-inputs}.

The remainders of the third class follow from the third-class remainder input. Its first member holds
as stated there. Its second member follows from the third-class off-chart input together with the
closure under $\mathsf{M}8$. Concretely, the second member is the pair of bounds (d)--(e) of that input
with $\mathsf{s}_v\mathsf{v}_v$ replaced by $C^{\perp}\mathsf{S}^{\perp}_v$.

\item The far pieces are controlled through three moves: $\mathsf{M}3$ at the far pairs, $\mathsf{M}4$
at the remainders, and $\mathsf{M}1^{M}$ combined with $\mathsf{M}6$.
\end{enumerate}

Hence the paired rows of the two-machine comparison at the rim, from the Cauchy row to the Wilson row,
hold as stated there with $\rho$ replaced by $\Theta\rho$. In particular the core term satisfies
\begin{equation}\label{5.5:eq-machine-differences-2}
|\delta\mathfrak{i}_y|\le C_I''(\mathsf{S}+\mathfrak{n}\mathsf{a})\,\nu^3(\Theta\rho)^2(1+\Theta\rho)\,
t^{4+\hat{\alpha}}.
\end{equation}
The short singles satisfy
\begin{equation}\label{5.5:eq-machine-differences-3}
|\delta f_v|\le 2\mathsf{S}\,e^{-\kappa d_j(X)}\min\bigl\{1,(S^{\star}_v)^{-2}\bigr\}.
\end{equation}
The own-read units satisfy
\begin{equation}\label{5.5:eq-machine-differences-4}
|\delta f_v|\le 2\mathsf{S}\,e^{-\kappa d_j(X)},
\end{equation}
by the membership input of Definition~\ref{5.5:def-weight-inputs}.

The Wilson pieces are bounded by $\Theta^2$ times their $\delta$-size. Off the chart this bound is
reached through the third-class off-chart input together with the closure under $\mathsf{M}8$. One
exception is left: the factor of the companions' item at $\nu\le 2$ treated in Step~4$'$.

Each of these bounds is a polynomial of degree at most three in $\Theta=c_Ww^{\vee}(y)\ge 1$. The
fraction $\tfrac18$ of the decay rate is spent once. Hence the degree bound of
\eqref{5.5:eq-weight-d} gives the factor $w^{\vee}(y)^3\,\mathfrak{N}^P(v;c_W)$ of
\eqref{5.5:eq-weight-h}, and these are the per-unit bounds of \eqref{5.5:eq-weight-h}, with the
substitutions $\mathfrak{n}\mapsto\mathsf{S}+\mathfrak{n}\mathsf{a}$, $C_I'\mapsto C_I''$ and
$\mathsf{s}_v\mathsf{v}_v\mapsto C^{\perp}\mathsf{S}^{\perp}_v$ stated there, apart from the item
treated in Step~4$'$.

\emph{Step 4$'$: cores with $\nu\le 2$.} At $\nu\le 2$ a core is own-read, by the class assignment of
Definition~\ref{5.5:def-near-unit-classes}, at every value of $S^{\star}_v$. Part~(i) of the proof of the
rows is stated for $\nu\ge 3$ and $L\ge 8$. This proof therefore provides at $\nu\le 2$ none of the
following on $\mathfrak{r}_v$: the identity, the comb rows, or a chart. None of them is used.

For one machine the core consists of two parts.
\begin{itemize}
\item Its constituents are units on $\mathcal{W}_T$, each bounded by $2\mathfrak{n}_v$.
\item Its counterterm companion $\beta[F]\hat{A}$ is a unit of the third class on $\mathfrak{b}'_j$. It
is bounded by $C'\mathfrak{n}z^2t^4$, by three statements: the counterterm-coefficient bound, the bound
for Wilson pieces, and \eqref{5.5:eq-weight-a}. This is part~(ii) of the proof of the rows of the
near-unit table and the clause of (b) on counterterm companions.
\end{itemize}

For two machines the constituents follow from the closure under $\mathsf{M}3$. By the membership input of
Definition~\ref{5.5:def-weight-inputs} their differences are bounded by
$2\mathsf{S}e^{-\kappa d_j(X)}$.

For the companion, the functional $\beta[\cdot]$ is linear. Hence
\begin{align}
\delta\bigl(\beta[F]\hat{A}\bigr)
&=\beta[F^{\mathsf{B}}]\hat{A}^{\mathsf{B}}-\beta[F^{\mathsf{A}}]\hat{A}^{\mathsf{A}}\notag\\
&=\beta[\delta F]\cdot\hat{A}^{\mathsf{B}}+\beta[F^{\mathsf{A}}]\cdot\delta\hat{A}.
\label{5.5:eq-machine-differences-5}
\end{align}
This is the two-machine identity for counterterm companions. The three factors are bounded as follows.
\begin{itemize}
\item $|\beta[\delta F]|$ is at most $\mathsf{S}$ times the constant and the squared inverse radius of
the paired counterterm-coefficient row of the two-machine comparison at the rim; this is the content
of that row.
\item $|\hat{A}^{\mathsf{m}}|\le\mathfrak{n}_{\hat{A}}$ in each machine, by the bound for Wilson pieces
together with \eqref{5.5:eq-weight-a}.
\item $\delta\hat{A}$ is controlled by the closure under $\mathsf{M}2^{\circ}$. Apply Step~3 at the
$\circ$-slot, which is entire, to
\begin{equation*}
\delta(\pi):=\hat{A}\bigl(\mathfrak{b}'^{\mathsf{B}}_j(\pi)\bigr)
-\hat{A}\bigl(\mathfrak{b}'^{\mathsf{A}}_j(\pi)\bigr).
\end{equation*}
This function is holomorphic on the parameter domain
$\mathfrak{P}=(b,\sigma)\times[0,1]\times\mathbb{D}_{\square}$ of $\mathsf{M}2$, whose first factor is
the range of the bond and $\sigma$ variables of the box read, and
$|\delta\hat{A}|\le\mathsf{a}$. This is a bound on the value; no chart enters.
\end{itemize}

Consider first the item $|\beta[\delta F]|\,|\hat{A}^{\mathsf{B}}|$. Because $\beta[\cdot]$ is linear,
it is the item of the table with $\mathfrak{n}$ replaced by $\mathsf{S}$.

Consider next the item
\begin{equation*}
|\beta[F^{\mathsf{A}}]|\,|\delta\hat{A}|\le C\mathfrak{n}\alpha_{3,j}^{-2}\mathsf{a}.
\end{equation*}
The two-machine statement for counterterm companions, which accompanies this identity, asserts
\begin{equation*}
|\delta\hat{A}|\le 2\mathsf{a}\,(CD\Theta\mathfrak{a}_nt^2)^2 ,
\end{equation*}
on the ground that $\hat{A}$ is entire at radius of order one: by the second Schwarz-lemma bound where
the radius of the chart of $\hat{A}$ exceeds one, and trivially otherwise. In the factor
$CD\Theta\mathfrak{a}_nt^2$, here and below, $C$ denotes a constant of that statement, containing the
factor $2d/\alpha_3(\circ)$; it is not the constant $C$ of the item just displayed. At $\nu\le 2$ we
examine the two halves of this assertion separately. The radius in question is that of the chart of
$\hat{A}$ on $\mathfrak{r}_v$,
\begin{equation*}
\alpha_3(\circ)\Big/\sup_{b\subset\mathfrak{r}_v}|B^{ax}|_*(b).
\end{equation*}

\begin{itemize}
\item Suppose this radius would be at most one, that is, $CD\Theta\mathfrak{a}_nt^2\ge 1$. Then the
value already carries the factor:
\begin{equation*}
|\delta\hat{A}|\le\mathsf{a}\le 2\mathsf{a}\,(CD\Theta\mathfrak{a}_nt^2)^2 .
\end{equation*}
The item is therefore at most $C'\mathfrak{n}\mathsf{a}z^2t^4$. This is the item of the table with
$\mathfrak{n}$ replaced by $\mathfrak{n}\mathsf{a}$. In total the two-machine statement for counterterm
companions gives the bound
\begin{equation*}
2C(\kappa)\mathsf{S}+C'(\mathsf{S}+\mathfrak{n}\mathsf{a})z^2t^4 .
\end{equation*}
Here $D\le 5M\nu\le 10M$ at $\nu\le 2$, and the resulting factor depending on $D$ is absorbed into
$C_I''$.

\item Suppose this radius would exceed one. Then the factor $(CD\Theta\mathfrak{a}_nt^2)^2$ is the one
given by the second Schwarz-lemma bound on a chart, and this proof does not construct that chart at
$\nu\le 2$. The table's form of the item there is carried by the off-chart clause of the paired Wilson
row of the two-machine comparison at the rim. That clause bounds the difference
\begin{equation*}
A^{\mathsf{B}}(c^{\mathsf{B}},P^{\mathsf{B}})-A^{\mathsf{A}}(c^{\mathsf{B}},P^{\mathsf{A}}),
\end{equation*}
in its own notation, through the bound for Wilson pieces. At the pair points of $\mathsf{M}2^{\circ}$,
this factor is the content of the low-order variation bound at the second cube size.

That bound is not among the hypotheses of clause~(e), and clause~(e) does not assert the table's form
of this item at these points; as clause~(e) states, that form rests on the low-order
variation bound at the second cube size, a statement on which the uniqueness theorem depends.

The third-class off-chart input, together with the closure under $\mathsf{M}8$, covers the second member
of the companion's third-class remainder at every $\nu$. Apart from the factor just named, what holds at
these points from the value alone is the bound $C\mathfrak{n}\alpha_{3,j}^{-2}\mathsf{a}$, which is
linear in $\mathsf{a}$.
\end{itemize}

At read cores ($\nu\ge 3$, every value of $S^{\star}_v$) the situation is different. The identity, the
comb rows and the four pointwise rows of part~(i) of the proof of the rows hold in each machine on
$\mathfrak{r}_v$. This is the hypothesis on the $\circ$-slot paired data of the closure theorem. The chart
of $\hat{A}$ in the two-machine identity for counterterm companions has radius of order one. It is the
paired chart of item~(a) of the move $\mathsf{M}5$ at radius $\alpha_3(\circ)$, taken through
$\mathsf{M}2^{\circ}$ at the common majorants. This chart carries $\delta\hat{A}$, with
\begin{equation*}
|\delta\hat{A}|\le 2\mathsf{a}\,(CD\Theta\mathfrak{a}_nt^2)^2 .
\end{equation*}
This is the form of the table. It is the form of the remainder line of the paired core row of the
two-machine comparison at the rim, which bounds the paired remainder of $\hat{A}$ by $C_I\mathsf{a}$ times
the weight at radius of order one, $C_I$ being the constant of that line.

No new constant, floor or ceiling enters Step~4$'$. The regimes are those of part~(iv) of the proof of
the rows and of the two-machine statement for counterterm companions.

\emph{Step 5: the block.} Three ingredients are field-free: clause~(c), the block-count lemma, and the
grouping by unit $n$-cells used in the proof of clause~(d) of Lemma~\ref{5.5:lem-weight}. Hence the same
$\mathsf{l}^2(\square)$ multiplies $\mathfrak{n}_0|_{E_0+1\mapsto\mathsf{l}^J}$. The ratio
$C_I''/C_I'$ is absorbed into $\mathfrak{n}_0|_1$, because the resummed total is linear in the sizes.

The passage from unit sizes to the $\delta$-sizes is paid by $\kappa_R^{\times}$. Consider the
$\delta$-sizes of the group of near units anchored in one unit $n$-cell $\Delta'$. Their per-cell row
sums in the near-unit table satisfy, row by row,
\begin{equation}\label{5.5:eq-machine-differences-6}
\sum_{v\in\text{row }r,\ y_v\in\Delta'}(\text{the per-unit bound of }v)
\le K_r\,\mathsf{l}^J\times(\text{the common cell factors of row }r).
\end{equation}
This is the per-cell row bound for the two-machine data functional. It holds for rows~1--3 and
$\partial$, and also for row~5, the remainder rows and the tail row of the closure theorem for the pair
system; for each of them $K_r$ is the single constant of Definition~\ref{5.5:def-weight-constants}.
\begin{itemize}
\item For the extended rows, $K_r$ is the per-cell number of the row at unit size, in the form in which
it is transported; the two-run constants are included in it.
\item For the third-class remainder row, the bound is the per-cube line of that row in its per-cube
$\delta$-size $\mathfrak{s}^{\perp}_{\alpha}$, with the constant $C^{\perp}K^W$, where $K^W$ is the
constant of the third-class remainder row.
\end{itemize}
The bound \eqref{5.5:eq-machine-differences-6} is read in the normalisation of the definition of
$\mathsf{l}^J$. The condition that the per-cell row-$r$ sums be at most $K_r$, in the definition of
$\kappa_R^{\times}$ (Definition~\ref{5.5:def-weight-constants}), uses the same constants $K_r$.

The resummation of the Cauchy-integral lemma for near units, with its scaffold coefficients, is linear in
the size vector; see the proof of clause~(d) of Lemma~\ref{5.5:lem-weight}. Hence the $\Delta'$-group at
the $\delta$-sizes equals $\mathsf{l}^J$ times the $\Delta'$-group at the $\delta$-sizes divided by
$\mathsf{l}^J$. By \eqref{5.5:eq-machine-differences-6}, the divided sizes have per-cell row-$r$ sums at
most $K_r$. So they lie in the set of sizes over which $\kappa_R^{\times}$ is defined. By that definition,
and since $\kappa_R\ge\kappa_R^{\times}$, the $\Delta'$-group is bounded by
\begin{align*}
\mathsf{l}^J\,\kappa_R^{\times}c_{\alpha}C_Re^{C_2\kappa_1}\vartheta(\square)
&\le\kappa_R\,c_{\alpha}C_Re^{C_2\kappa_1}\vartheta(\square)\,\mathsf{l}^J\\
&=\frac{\mathfrak{N}_0(\square)}{E_0}\,\kappa_R\,\mathsf{l}^J .
\end{align*}
This is the display \eqref{5.5:eq-weight-g} of (d) with $E_0$ replaced by $\mathsf{l}^J$. The linearity
in the sizes is used row by row, and no operator bound is needed.

The per-cell sums of rows~1 and~3 are summed over $j$ at the $\delta$-sizes; for row~3, whose cell number
is summable only through the sizes, this $j$-sum is the second member of $\mathsf{l}^J$.

The far pieces and the remainders are linear in the function. The items of the third class are treated as
in Steps~4 and~4$'$. Off the chart they follow from the third-class off-chart input together with the
closure under $\mathsf{M}8$. The factor of Step~4$'$ at $\nu\le 2$ is the low-order variation bound at
the second cube size, which is named there and is not an input of clause~(e). This proves the bound on
the differences at the one-block argument $\{\square\}$ asserted in clause~(e).

Clause~(e) is used in three ways. In the bound \eqref{5.5:eq-data-factor-a} on the data factor, read at
unit normalisation, the factor $(1+C^{\mathtt{a}})$ is kept on differences. The data functional
$\mathsf{l}_{\mathsf{n}+1}$ of the next stage contains $\mathsf{l}^J$ in sum form. The difference
statement for the count sums is taken over the same pair system $\mathcal{K}$.
\end{proof}

\begin{definition}[Two cutoffs on one scaffold]\label{5.5:not-one-scaffold}
Throughout this section two runs of Ba{\l}aban's renormalisation construction are compared. They are
called the two cutoffs and are denoted $\mathsf{A}$ and $\mathsf{B}$; the run $\mathsf{A}$ has length
$n$ and the run $\mathsf{B}$ has length $n+1$. The letter $\mathsf{m}\in\{\mathsf{A},\mathsf{B}\}$
denotes either of them. The levels of both runs are indexed by the levels of $\mathsf{A}$, and the
labels of the two runs are aligned.

\smallskip
\emph{The scaffold.} The scaffold is that of Definition~\ref{5.1:def-source-free-scaffold}: the
partitions, the labels, the caps, the cover and its partition of unity, the enlargements, the radii,
the frozen weights and the associated geometric quantities fixed there, namely
\begin{equation}\label{5.5:eq-one-scaffold-a}
\begin{gathered}
\mathfrak{B}:=\mathbb{B}^{(\mathsf{n})}_k,\quad \mathfrak{B}^+,\quad \mathfrak{B}_\ell,\quad
\mathfrak{B}^*_j,\quad \{\square\},\quad \zeta_\square,\quad \zeta'_\square,\quad \square^\zeta,\quad
\square_0:=\tilde\square^5,\quad \hat{\mathsf{K}}_\square,\quad \mathcal{C}_{\mathsf{j}},\\
r_\square^\sharp,\quad R^\natural_\square,\quad r_T(\square),\quad r^{\mathrm{rem}}_\square,\quad
R^\perp_v,\quad \mathsf{w}_0=\tfrac{1}{24},\\
w,\quad \mathcal{W},\quad w^\vee,\quad \Theta:=c_W\,w^\vee(y),\quad S^*_v,\quad
\mathfrak{S}_\square,\quad D_v,
\end{gathered}
\end{equation}
together with every constant fixed there, among them $c_W$; the anchor $y$ entering $\Theta$ is that
of the near unit, as fixed there. Here $\mathfrak{B}$ and $\mathfrak{B}^+$ are the labels, $\mathfrak{B}_\ell$ and
$\mathfrak{B}^*_j$ are the caps, $\{\square\}$ is the cover, $\zeta_\square$ is the partition of unity
subordinate to it and $\zeta'_\square$ its companion function, as fixed there, and the radius
$R^\perp_v$ is the one carried by the entry, acting on cell-units, of the list of comparison maps
$\mathfrak{M}^{\mathrm{rd}}$ introduced below. The scaffold \eqref{5.5:eq-one-scaffold-a}, with
every constant, is one and the same for $\mathsf{A}$ and for $\mathsf{B}$: no object in
\eqref{5.5:eq-one-scaffold-a} depends on $\mathsf{m}$.
\end{definition}

\begin{remark}[Homonyms]
The following letters must be kept apart.
\begin{enumerate}
\item[(a)] The stage letters $\mathsf{n}$ and $\mathsf{j}$ (sans-serif) are distinct from the letters
$j$, $n$, $\nu$, $t$ of the correlation theorem (Theorem~\ref{5.1:thm-correlation-theorem}).
\item[(b)] The symbols of the form $(\mu\,\cdot)$ used in what follows are those of the rim
refinement (Definition~\ref{5.2:def-rim-refinement}), and not the symbols of the same shape in the
sourced preparatory step (Corollary~\ref{5.2:cor-sourced-preparatory}).
\item[(c)] The letter $\mathcal{K}$, where it enters through the corollary of the remainder bound
applied below, denotes a $1$-form; it is not the production closure $\mathcal{K}^{\mathrm{rd}}$ of
this section.
\item[(d)] In one comparison statement applied later in this section, identified at the place where
it is applied, the symbol $\star$ is written for the star $*$ of this section. The subscript $\star$
of $d_\star$ in \textup{(e)} below belongs to the letters of the entry described there and is not an
instance of this.
\item[(e)] In the entry, acting on cell-units, of the list of comparison maps
$\mathfrak{M}^{\mathrm{rd}}$ introduced below, $v=(c,i,\Delta)$ is a cell-unit of class $W$, and
\begin{equation*}
D_v:=d_\star(\Delta,F_\square),\qquad F_\square=\mathfrak{S}_\square ,
\end{equation*}
where $d_\star$ is the distance of that entry and $\mathfrak{S}_\square$ is the scaffold object of
\eqref{5.5:eq-one-scaffold-a}; in that entry $s$ denotes the dilation parameter, and $\kappa^*_v$,
$R^\perp_v$, $\mathsf{S}^\perp_v$ are the letters attached to this cell-unit. This $v$ and this $D_v$
are not the near unit $v$ of the scaffold and its $D_v$, and this $s$ is not the exponent $s$ in the
factor $ML^s$.
\end{enumerate}
\end{remark}

\begin{definition}[Stage slots, unit slots and their letters]\label{5.5:def-slots-letters}
Throughout, $\mathsf{A}$ and $\mathsf{B}$ are the two machines on the one scaffold of
Notation~\ref{5.5:not-one-scaffold}, $\mathsf{A}$ of length $n$ and $\mathsf{B}$ of length
$n+1$, with levels indexed by those of $\mathsf{A}$; $\mathsf{m}\in\{\mathsf{A},\mathsf{B}\}$
denotes either machine.

\smallskip
\noindent(a) \emph{Stage slots.} A \emph{stage slot} is a triple
$\mathfrak{s}=(k;\mathsf{n},\hat X)$ in which $k$ is a level and the slot belongs to the site
(P) of the large-field operation $\mathbf{R}_k$ (the sites are those of
Definition~\ref{5.1:def-expansion-rule}), $\mathsf{n}\in\{0,\dots,\mathsf{N}\}$ is a stage
of the expansion at that site, $\mathsf{N}$ denoting its last stage, and $\hat X$ is a
localisation domain of stage $\mathsf{n}$. Each of the terms evaluated at the site (P)
that enter the summation of boundary terms at finer zones, namely the terms $\mathbb{B}$ of
that summation (a symbol distinct from the ball $\mathbb{B}_k$ of complex histories) and
the chain terms $\mathbb{C}^{(i)}$, has a slot of its own, a \emph{boundary-term slot}. Each
boundary-term slot carries its own domain, the set
\begin{equation*}
\mathfrak{D}^\flat_n(X)\times\mathfrak{D}_{\mathcal{S}}(X),
\end{equation*}
where $X$ is the domain of the boundary term evaluated at the slot and $n$ is the stage
index of its inherited domain, the index which appears as $\mathsf{n}+1$ in part (c) below
(here $n$ is not the length of $\mathsf{A}$). This domain is the one fixed in the theorem on
inherited domains; its first factor is the inherited domain of part (a) of that theorem.

\smallskip
\noindent(b) \emph{Unit slots.} A \emph{unit slot} is a triple
$\mathfrak{u}=(j,X;\mathfrak{L})$ in which $(j,X)$ is a small-field unit, in the sense of
Definition~\ref{5.2:def-sites-units}, of the expansion on which the correlation theorem
(Theorem~\ref{5.1:thm-correlation-theorem}) rests, of level $j$ and with domain $X$, and
$\mathfrak{L}\in\{\flat,\circ\}$ is the \emph{letter} of the slot. For a term $T$ we write
$\ell_T$ for its level and $X_T$ for its domain.

\smallskip
\noindent(b$\flat$) \emph{The letter $\flat$.} The letter $\flat$ is given to the terms of
the first three families of the expansion at the cover and to the boundary terms,
that is, to the terms $T$ of the class $\mathcal{T}_1$ with $(\ell_T,X_T)=(j,X)$. At such a
slot we put
\begin{equation}\label{5.5:eq-slots-letters-a}
\mathfrak{U}^{\mathsf{m}}(j,X;\flat)
:=\bigcap\bigl\{\mathfrak{U}^{\mathsf{m}}_T:\ (\ell_T,X_T)=(j,X)\bigr\},
\end{equation}
the intersection of the analyticity tubes $\mathfrak{U}^{\mathsf{m}}_T$ of the terms $T$
stored at the slot; for a term $T$ attached to a unit $v$, the tube $\mathfrak{U}^{\mathsf{m}}_T$
is the set $\mathfrak{D}_v$. This is one set per slot; at near units, introduced later in this
section, it is the analyticity space $\mathcal{U}^\flat_j(X)$ of
Definition~\ref{5.1:def-analyticity-spaces}. Consequently every supremum over a unit slot of
letter $\flat$
taken below runs over one set. The requirements (m3), (m6) and (m7) imposed below on the
moves of the list $\mathfrak{M}^{\mathrm{rd}}$ of comparison maps are imposed for every term
$T$ stored at the slot; the map of the third entry of $\mathfrak{M}^{\mathrm{rd}}$ and its
parameter disc depend on $(\ell_T,X_T,\square)$ only. The set \eqref{5.5:eq-slots-letters-a}
is, like the spaces $U^c_j$ of \cite{B-CMP109}, an open union of $G^{\mathbb{C}}$-orbits,
where $G^{\mathbb{C}}$ is the
complexified gauge group acting by gauge transformations. The chart radii of the letter
$\flat$ are
\begin{equation*}
\alpha_3(\flat):=\alpha_{3,j}=c_3a_j,\qquad \mathfrak{r}(\flat):=\mathfrak{r}_j .
\end{equation*}

\smallskip
\noindent(b$\circ$) \emph{The letter $\circ$.} The letter $\circ$ is given to the Wilson
pieces of the expansion at the cover, to the counterterm pieces $\beta[F]\,A(h_y,\cdot)$ of
the cores,
and to the pointed Wilson family $\hat A$ itself. These functions are entire, and no
membership in an analyticity domain is required of their arguments. Their chart radii are
absolute:
\begin{equation*}
\alpha_3(\circ):=c_3a^\circ,\qquad a^\circ:=(\alpha_0,\alpha_1),\qquad
\mathfrak{r}(\circ):=\mathfrak{r}_\circ ,
\end{equation*}
where $(\alpha_0,\alpha_1)$ are the absolute constants of \cite{B-CMP109}. For the letter
$\circ$ there is no tube. The absolute radii $\alpha_3(\circ)$ and $\mathfrak{r}(\circ)$ are
chart radii, namely the radii of the parameter discs
\begin{equation*}
|B'|<\alpha_3(\circ)\quad\text{in (m5a)},\qquad
\sum_i|\sigma_i|\,\|(\mathsf{h}_i,\mathbf{j}_i)\|_*<\mathfrak{r}(\circ)\quad\text{in (m5b)},
\end{equation*}
and the condition ``$\in\mathfrak{U}^{\mathsf{m}}(j,X;\mathfrak{L})$'' at
$\mathfrak{L}=\circ$ in (m5a) expresses that the move lands at such a configuration with
its entry bound (c2); it does not express membership in an analyticity domain. Accordingly,
$\mathfrak{U}^{\mathsf{m}}(j,X;\circ)$ is the set of configurations at which the moves listed
below land, the $\circ$-families being bounded there by part (b) of the Wilson-family
lemma at the entry bound that the move carries. The entry bounds are the following.
\begin{itemize}
\item For $(\mathrm{m}2)^\circ$: the entry bound (a2).
\item For $(\mathrm{m}3)^\circ$, $(\mathrm{m}6)^\circ$, $(\mathrm{m}7)^\circ$: a scaffold
plaquette bound on the support of the piece, namely, cell by cell,
\begin{equation*}
|U(\partial p)-1|\le\mathsf{a}\,(L^{k-n}\eta)^2
\end{equation*}
for the plaquettes $p$ of the support, at a scaffold constant $\mathsf{a}$.
\item For $(\mathrm{m}4)^\circ$: the entry bound of the statement that supplies the move.
\item For $(\mathrm{m}\perp)^\circ$: for the entry $\mathsf{M}^\perp$ of
$\mathfrak{M}^{\mathrm{rd}}$, with its cell-unit, written $v=(c,i,\Delta)$ in the notation of
the lemma on the dilation move $\mathsf{M}^\perp$ and distinct from the unit $v$ of
part (b$\flat$): the plaquette bound of part (i) of that lemma on the dilated cell, for
dilation parameters $|s|\le R^\perp_v$, together with the bound of part (iii) of that lemma,
\begin{equation*}
|A(c,U_s)|\le\mathfrak{n}^\perp_v:=0.008\,\mathfrak{c}_{\mathrm{tr}}\,\mathfrak{n}_v ,
\end{equation*}
where $\mathfrak{c}_{\mathrm{tr}}$ is a constant of that lemma and $\mathfrak{n}_v$ is the
one-cutoff bound at the cell-unit $v$, the definition of $\mathfrak{n}^\perp_v$ being among
the constants fixed with that lemma;
this is the entry bound attached to the letter $\circ$ on the image of $\mathsf{M}^\perp$.
\item For $(\mathrm{m}5\mathrm{a})$, $(\mathrm{m}5\mathrm{b})$: the entry bound (c2), that
attached to the letter $\circ$ with its absolute radii in the lemma that supplies these two
moves.
\end{itemize}
We write $\mathfrak{n}_{\hat A}$ for the supremum, over the moves that land in the class of
letter $\circ$, of the bounds for $\hat A$ given by part (b) of the Wilson-family lemma.
What is asserted of $\mathfrak{n}_{\hat A}$ is only that $\mathfrak{n}_{\hat A}<\infty$, and
this is asserted given the entry bounds listed above.

\smallskip
\noindent(c) \emph{Source classes.} For a stage slot $\mathfrak{s}=(k;\mathsf{n},\hat X)$,
$\mathfrak{S}^{\mathsf{m}}(\mathfrak{s})$ is the \emph{source class} of machine $\mathsf{m}$
at $\mathfrak{s}$: the sources of stage $\mathsf{n}$ over $\hat X$, in the form in which the
arrest theorem reads them; each is a field on the whole domain of the label $\mathfrak{B}$,
with clauses imposed over $\hat X$. The terms $\mathbb{C}^{(\mathsf{n}+1)}_k(\hat X,\cdot)$ of
a stage at the site (P) are stored both on the stage-$(\mathsf{n}+1)$ domain
$\tilde U^{(\mathsf{n}+1)c,\Sigma}_k(\hat X)$ of the expansion at the site (P) and on the
inherited domain
$\mathfrak{D}^\flat_{\mathsf{n}+1}(\hat X)$ of part (a) of the theorem on inherited domains;
the latter is where their evaluations at the site (L) land. For the passage through such a
slot, treated in part (ii) of the first case of the statement on passages through slots,
the requirements (H2)(m1), (m6) and (H3) below are imposed on the elements
$P^{\mathsf{m}}\in\mathfrak{S}^{\mathsf{m}}(k;\mathsf{n}+1,\hat X)$, this class being read as
including that domain. This reading is taken by statement from part (a) of the theorem on
inherited domains, from the corollary of the Wilson-family lemma for machine $\mathsf{A}$,
and from the statement that fixes the source classes of machine $\mathsf{m}$.

\smallskip
\noindent(d) \emph{Product classes.} For a stage slot $\mathfrak{s}$ and a unit slot
$\mathfrak{u}$ we put
\begin{equation*}
\mathcal{U}^\times(\mathfrak{s}):=\mathfrak{S}^{\mathsf{B}}(\mathfrak{s})\times
\mathfrak{S}^{\mathsf{A}}(\mathfrak{s}),\qquad
\mathcal{U}^\times(\mathfrak{u}):=\mathfrak{U}^{\mathsf{B}}(\mathfrak{u})\times
\mathfrak{U}^{\mathsf{A}}(\mathfrak{u}),
\end{equation*}
and write their elements as $p=(P^{\mathsf{B}},P^{\mathsf{A}})$. The base point is
$\mathbb{1}:=((1,0),(1,0))$. For a quantity $F$ defined in both machines, with values
$F^{\mathsf{B}}$ and $F^{\mathsf{A}}$, the two-cutoff difference is
\begin{equation*}
\delta F(p):=F^{\mathsf{B}}(P^{\mathsf{B}})-F^{\mathsf{A}}(P^{\mathsf{A}}).
\end{equation*}

\smallskip
\noindent(e) The same pair $(j,X)$ with the two letters $\flat$ and $\circ$ gives two
distinct unit slots.
\end{definition}

\begin{definition}[Moves between slots]\label{5.5:def-moves}
We work on the one scaffold of Notation~\ref{5.5:not-one-scaffold}, with the two cutoffs
$\mathsf{A}$ and $\mathsf{B}$, and with the slots and one-cutoff classes of
Definition~\ref{5.5:def-slots-letters}. For a slot $\mathfrak{s}$ and a cutoff
$\mathsf{m}\in\{\mathsf{A},\mathsf{B}\}$ write $\mathfrak{C}^{\mathsf{m}}(\mathfrak{s})$ for the
one-cutoff class of $\mathfrak{s}$ in cutoff $\mathsf{m}$ fixed there, for every kind of slot
defined there; for instance, it is the source class $\mathfrak{S}^{\mathsf{m}}(\mathfrak{s})$ at
a stage slot, and the class $\mathfrak{U}^{\mathsf{m}}(j,X;\mathfrak{L})$ at a unit slot
$(j,X;\mathfrak{L})$.

A \emph{move} $\mu\colon\mathfrak{s}'\to\mathfrak{s}''$ from a slot $\mathfrak{s}'$ to a slot
$\mathfrak{s}''$ is a pair of maps $\mu^{\mathsf{m}}$, $\mathsf{m}\in\{\mathsf{A},\mathsf{B}\}$,
\begin{equation*}
\mu^{\mathsf{m}}\colon
\mathfrak{P}_\mu\times\mathfrak{C}^{\mathsf{m}}(\mathfrak{s}')
\longrightarrow
\{\text{configurations over the domain of }\mathfrak{s}''\},
\end{equation*}
subject to the following three requirements.
\begin{enumerate}
\item[(a)] Neither map depends on the complex coefficients $z_i$ of the source
$w=\sum_i z_i f_i$.
\item[(b)] $\mu^{\mathsf{A}}$ and $\mu^{\mathsf{B}}$ are given by the same construction, read
in cutoff $\mathsf{A}$ and in cutoff $\mathsf{B}$ respectively.
\item[(c)] The two maps have one and the same parameter domain $\mathfrak{P}_\mu$: the product of
a polydisc with a set of real supports, all radii being numbers of the scaffold of
Notation~\ref{5.5:not-one-scaffold}, the same for both cutoffs.
\end{enumerate}
The \emph{pair map} of $\mu$ is $\mu^\times:=(\mu^{\mathsf{B}},\mu^{\mathsf{A}})$. We evaluate its
two entries at one parameter $\pi\in\mathfrak{P}_\mu$: for a pair $p=(P^{\mathsf{B}},P^{\mathsf{A}})$
in the product class $\mathcal{U}^\times(\mathfrak{s}')$ of
Definition~\ref{5.5:def-slots-letters},
\begin{equation}\label{5.5:eq-moves-a}
\mu^\times(\pi,p):=\bigl(\mu^{\mathsf{B}}(\pi,P^{\mathsf{B}}),\,\mu^{\mathsf{A}}(\pi,P^{\mathsf{A}})\bigr).
\end{equation}

That $\mu^{\mathsf{m}}$ takes its values in the class $\mathfrak{C}^{\mathsf{m}}(\mathfrak{s}'')$
is not part of this definition; it is required of each move separately, one cutoff at a time,
as a hypothesis of the statements on moves that follow.
\end{definition}

\begin{definition}[Response maps, box pairs and weighted scaffold majorants]
\label{5.5:not-response-box-majorants}
The objects below live at the site (P) of the operation $\mathbf{R}_k$ in the sense of
Definition~\ref{5.2:def-sites-units}. They are defined for each
of the two cutoffs $\mathsf{A}$, $\mathsf{B}$ of Notation~\ref{5.5:not-one-scaffold}, on that
cutoff's own lattice. The labels $\mathfrak{B}$, $\mathfrak{B}^+$, $\mathfrak{B}_\ell$,
$\mathfrak{B}^*_j$, the cover $\{\square\}$ with $\zeta_\square$, $\zeta'_\square$,
$\square^\zeta$, $\square_0:=\tilde\square^5$, the set $\mathcal{C}_{\mathsf{j}}$, the radius
$r^\sharp_\square$, the constant $\mathsf{w}_0=1/24$, the frozen weights and every constant are
those of the scaffold of Notation~\ref{5.5:not-one-scaffold}. A term is said to be
\emph{read on} a pair (configuration, source) when it is evaluated at that pair.

\smallskip
\noindent(a) \emph{Parameters.} The source $b$ is a complex function on $\mathcal{C}_{\mathsf{j}}$
with $\|b\|_\infty\le 1.1\,C_1\delta'_{\mathsf{j}}$, where $C_1$ is a constant of the scaffold and
$\delta'_{\mathsf{j}}$ is the source ceiling of stage $\mathsf{j}$; $\sigma=(\sigma(\Delta))$ with
$|\sigma(\Delta)|\le e^{\kappa_1}$; $t\in[0,1]$; $(t_\square)_\square$ are block amplitudes.
The box amplitude $\mathsf{t}$ ranges over the disc $\mathbb{D}_\square$ of
\eqref{5.5:eq-response-box-majorants-c}, or over the disc
$\{|\mathsf{t}|\le r^{\mathrm{rem}}_\square\}$, which is contained in $\mathbb{D}_\square$
because $r^{\mathrm{rem}}_\square\le r^\sharp_\square$. For a term $T$ the own-level parameter
$\tau$ satisfies $|\tau|\le\mathsf{w}_0\,r_T(\square)$, with $r_T(\square)$ as in
\eqref{5.5:eq-response-box-majorants-c}. In (a)--(c) the bare letter $t$ is the interpolation
parameter in $[0,1]$, and $t_\square$ are the block amplitudes, on which $\chi_t$ and $\Phi_t$
depend; from (d) on $t$ denotes the ratio $L^{j-n}$ of (d).

\smallskip
\noindent(b) \emph{Kernel and response dilation.} $\mathbb{H}$ is the response kernel of the
theorem constructing the kernel and its box layers, taken at
$(\mathfrak{B};\mathbb{U},\mathbb{J};\sigma;(0,b))$, where $\mathbb{U}$ is the configuration and
$\mathbb{J}$ the source at which the response is taken. The response dilation $\Phi_t$ is given in
\eqref{5.5:eq-response-box-majorants-c}; its second component is local. There $\eta$ is the
lattice spacing and $\mathcal{J}^\sharp[1]$ is the source shift carried by the dilation.

\smallskip
\noindent(c) \emph{Window and box pairs.} The window $\mathbb{W}$ is the configuration
$U(\mathfrak{B}^+(\square_0),M^{\cdot}(\mathbb{U}))$ on $\square_0$, determined on
$\mathfrak{B}^+(\square_0)$ by the block averages $M^{\cdot}(\mathbb{U})$ of $\mathbb{U}$, extended
by $\mathbb{U}$ off $\square_0$; $\mathbb{H}^w$ is $\mathbb{H}$ at
$(\mathfrak{B};\mathbb{W},\mathbb{J})$. Put
\begin{equation*}
\mathbf{A}:=(t\zeta'_\square+\mathsf{t}\zeta_\square)\mathbb{H}^w,\qquad
\mathfrak{b}:=e^{i\eta\mathbf{A}}\mathbb{W}.
\end{equation*}
Let $\mathbf{H}_j(t,\mathsf{t})$ be the own-level box layer of $\mathbb{H}$ at $\mathfrak{B}_j$
given by the same theorem, and, with $\zeta'=\zeta'_\square$,
\begin{equation*}
\begin{aligned}
D^\sharp&:=\zeta'\cdot Q_*\bigl((t+\mathsf{t}\zeta_\square)\mathbb{H}^w\bigr)
+(1-\zeta')\cdot Q_*\bigl(\mathbf{H}_j(t,\mathsf{t})\bigr),\\
\mathsf{E}&:=\tfrac{1}{i}\log\bigl(e^{iD^\sharp}e^{-iQ_*(\mathbf{A})}\bigr).
\end{aligned}
\end{equation*}
The discs, the radii, the response dilation, the box pair $\mathsf{P}^0$ and the second box pair
$\mathsf{P}^1$ are
\begin{equation}\label{5.5:eq-response-box-majorants-c}
\begin{gathered}
\mathbb{D}_\square:=\{\mathsf{t}:|\mathsf{t}|\le r^\sharp_\square\},\qquad
r_T(\square):=\frac{e^{\delta d_\ell(X_T,\square^\zeta)/4}}{14K^+\vartheta(\square)},\\
r^{\mathrm{rem}}_\square:=\min\Bigl\{r^\sharp_\square,\,
\frac{1}{14K^{\mathrm{rem}}\vartheta(\square)}\Bigr\},\\
\Phi_t:=\bigl(e^{i\eta\chi_t\mathbb{H}}\mathbb{U},\ \mathbb{J}+\chi_t\mathcal{J}^\sharp[1]\bigr),
\qquad \chi_t:=\sum_\square t_\square\zeta_\square,\\
\mathsf{P}^0:=\mathfrak{b}'_j:=(U,J)(\mathfrak{B}^*_j,M_*\mathfrak{b}),\qquad
\mathsf{P}^1(t,\mathsf{t}):=(U,J)(\mathfrak{B}^*_j,e^{i\mathsf{E}}M_*\mathfrak{b}).
\end{gathered}
\end{equation}
Here $d_\ell$ is the level-$\ell$ distance of the scaffold, $K^+$ and $K^{\mathrm{rem}}$ are
constants of the scaffold, $\vartheta(\square)$ is the block weight of $\square$, $M_*$ is the
averaging map onto the label $\mathfrak{B}^*_j$, $Q_*$ its linearisation, and
$(U,J)(\mathfrak{B}^*_j,\cdot)$ is the configuration--source pair determined on
$\mathfrak{B}^*_j$ by the averaged data. $\mathsf{E}$ is a function on the bonds of
$\mathfrak{B}^*_j$.

\smallskip
\noindent(d) \emph{Own-level pair and near-unit geometry.} For a term $T$ of $\mathcal{T}_1$, the
class of stored terms each carrying a level $\ell_T$ and a domain $X_T$, the
own-level pair is $\mathcal{W}_T:=\mathcal{W}^{\ell_T,\square}\lceil X_T$, the restriction to
$X_T$ of the own-level pair $\mathcal{W}^{\ell_T,\square}$ of level $\ell_T$ at $\square$ given by
the scaffold. A near unit $v$ has domain $X_v\subset\tilde\square^2$, the second enlargement of
$\square$, and anchor $y$; with $\ell(y)$ the level of $y$ and $n_0$ the ceiling on the anchor's
level, put
\begin{equation*}
n:=\min\{\ell(y),n_0\},\qquad \nu:=1+n-j,\qquad t:=L^{j-n}.
\end{equation*}
In the majorants of \eqref{5.5:eq-response-box-majorants-a} and
\eqref{5.5:eq-response-box-majorants-1} the letter $t$ denotes this ratio. The near units are
divided, in the expansion at the site (P), into cores (with their members, constituents and
tails) and singles, short or long. The own box of $v$ is
$\mathfrak{r}_v$ (for a core, the enlargement $\square^{\sim\nu+1}$ of $\square$, with one layer
of $\pi_j$-blocks as margin),
and $D:=M(2\nu+3)$ is its side; $1+|b-y|\le D$ for every bond $b$ of $\mathfrak{r}_v$.

\smallskip
\noindent(e) \emph{Support of $\mathsf{E}$ and the geometric distance.} With $\tilde B(c)$ the
block of the bond $c$, $\tilde\square^3$ the third enlargement of $\square$ and $d_*$ the scaffold
distance between sets of bonds, the lemma on the second
box pair gives the first property below, which is produced by the cut-off $\zeta'$ in
$D^\sharp$ and not by the constraint on the configuration, and the lower bound (E1):
\begin{equation}\label{5.5:eq-response-box-majorants-b}
\begin{gathered}
\text{(E0)}\quad \mathsf{E}(c)=0\ \text{ whenever } \tilde B(c)\subset\tilde\square^3,
\ \text{ that is, }\ \operatorname{supp}\mathsf{E}\subset\mathfrak{S}_\square,\\
\mathfrak{S}_\square:=\{c\in\operatorname{bonds}(\mathfrak{B}^*_j):
\tilde B(c)\not\subset\tilde\square^3\},\qquad
D_v:=d_*(\mathfrak{r}_v,\mathfrak{S}_\square),\\
\text{(E1)}\quad D_v\ge\tfrac15\,ML^s.
\end{gathered}
\end{equation}
In (E1) the exponent $s$ is the one carried by that bound of the lemma on the second box pair;
it is not the depth below the cutoff.
The set $\mathfrak{S}_\square$ and the distance $D_v$ are determined by the scaffold alone, and
$D_v\le d_*(\mathfrak{r}_v,\operatorname{supp}\mathsf{E})$, the distance computed from the
field; accordingly clause (iii) of the remainder-decay lemma for the second box pair, the
remainder theorem for that pair and (E1) hold as stated with this $D_v$. The same $D_v$ is used
in clauses (a), (d), (e), (g) of the lemma on the second box pair, in the remainder-decay lemma,
in $|\Delta|_*$ below and in the estimate (S4) throughout.

\smallskip
\noindent(f) \emph{Weighted scaffold majorants.} On $\mathfrak{r}_v$, in the comb gauge anchored at
$y$, in which the configuration on $\mathfrak{r}_v$ is written through the comb row below, with
$\Theta=c_Ww^\vee(y)$ and with $c_1$, $\hat\alpha$ the constants of the comb-data bound,
\begin{equation}\label{5.5:eq-response-box-majorants-a}
\begin{aligned}
|B^{\mathrm{ax}}|_*(b)&:=2d\,(1+|b-y|)\,\Theta\,\mathfrak{a}_nt^2, &
|F|_*&:=\Theta\,\mathfrak{a}_nt^2,\\
|\mathfrak{e}_\mu|_*(x)&:=c_1(1+|x-y|)^2\,\Theta\,\mathfrak{a}_nt^{2+\hat\alpha}, &
|\mathfrak{e}^{\mathrm{rim}}|_*&:=4dD\,\Theta\,\mathfrak{a}_nt^2,\\
|B_0|_*&:=|B^{\mathrm{ax}}|_*, &
|B_F|_*(b)&:=d\,(1+|b-y|)\,\Theta\,\mathfrak{a}_nt^2.
\end{aligned}
\end{equation}
The constant $|\mathfrak{e}^{\mathrm{rim}}|_*$ is the majorant under which the rim family of the
comb data, the family of its rim terms, is estimated. It dominates on $\mathfrak{r}_v$ the
pointwise form of the rim majorant of clause (a4) of the weighted scaffold bound for the comb data:
that form depends on $b$ through the factor $1+|b-y|$ in the place of $D$, and $1+|b-y|\le D$ on
$\mathfrak{r}_v$. The comb row is
\begin{equation*}
B^{\mathrm{ax}}_\mu(x)=\sum_{\kappa<\mu}(x-y)_\kappa F_{\kappa\mu}(y)+\mathfrak{e}_\mu(x),
\end{equation*}
and $|B_F|_*$ is the weighted scaffold majorant of the $F$-part of (S2), the bond function
$B_F(b):=\sum_{\kappa<\mu}(x-y)_\kappa F_{\kappa\mu}(y)$ at the bond $b$ from $x$ in direction
$\mu$, the first member of the comb row; its pointwise bound is obtained from that of $F$. One has
\begin{equation}\label{5.5:eq-response-box-majorants-2}
|B_F|_*(b)=d\,(1+|b-y|)\,|F|_*\le\tfrac12|B^{\mathrm{ax}}|_*(b).
\end{equation}
For the estimate (S4), and for the classical directions $\mathsf{h}$, with $B_3$ and $\delta_0$
the constants of the bound on the classical directions, put
\begin{equation}\label{5.5:eq-response-box-majorants-1}
\begin{aligned}
|B^0|_*&:=1.42\,L^{-1}\sup_{b\in\mathfrak{r}_v}|B^{\mathrm{ax}}|_*(b), &
|\Delta|_*&:=\tfrac12\,\alpha_{3,j}\,e^{-\delta D_v/8},\\
\|\mathsf{h}_b\|_*&:=B_3\,e^{-\delta_0\operatorname{dist}_j(b,X)}, &
\|\mathsf{h}[B]\|_*&:=\sum_b\|\mathsf{h}_b\|_*\,|B|_*(b).
\end{aligned}
\end{equation}
Here $X$ is the domain of the unit slot $(j,X;\mathfrak{L})$ of
Definition~\ref{5.5:def-slots-letters} at which $\mathsf{h}$ is read.
The constant $|B^0|_*$ equals $1.42\,\alpha_{3,j}/(LS^*_v)$, with $S^*_v$ as in
\eqref{5.5:eq-response-box-majorants-d}, and dominates the pointwise bound
$1.2\,L^{-1}|B^{\mathrm{ax}}|_*(b)$ of clause (ii) of the remainder-decay lemma at every bond $b$
of $\mathfrak{r}_v$, $1.2/L$ being the constant of the working gauge there; $|\Delta|_*$ is the
bound of clause (iii) of that lemma.

\smallskip
\noindent(g) \emph{Assignment of the pairs on which near units are read.} A near unit of the kinds
to which
Definition~\ref{5.5:def-slots-letters} gives the letter $\flat$, other than the terms of
$\mathcal{T}_1$, is \emph{eligible} if $\nu\ge3$ and it is a short single or a member of a core.
With $\alpha_{3,j}$ the chart radius of Definition~\ref{5.5:def-slots-letters},
\begin{equation}\label{5.5:eq-response-box-majorants-d}
S^*_v:=\frac{\alpha_{3,j}}{\sup_{b\in\mathfrak{r}_v}|B^{\mathrm{ax}}|_*(b)},\qquad
v\in\text{box class}\iff v\ \text{is eligible and}\ S^*_v>1.
\end{equation}
A unit of the box class is read on the box pair; a core is then read as a whole, on a single box
pair. The \emph{own-level class} consists of the other near units of those kinds: tails, long
singles, constituents with $\nu\le2$, and eligible units with $S^*_v\le1$, the latter read
constituent by constituent. A unit of the own-level class is read at its own level, not on the
box pair. The \emph{Wilson class} consists of every near Wilson or counterterm
piece (the block Wilson pieces, taken plaquette by plaquette at the level of their cell, and a
core's counterterm piece $\beta[F]A(h_y,\cdot)$); at every $\nu$ and every value of $S^*_v$ it is
read on the box pair $\mathfrak{b}'_j$. Far pairs of $\mathcal{T}_1$, those with
$X_T\not\subset\tilde\square^2$, Wilson pieces included, are read on $\mathcal{W}_T$.

Every symbol of the predicates in (g) and of the majorants in (f) belongs to the scaffold of
Notation~\ref{5.5:not-one-scaffold}, the weight $w^\vee$ in $\Theta$ being a maximum of frozen
label weights; the assignment (g) and the radii in \eqref{5.5:eq-response-box-majorants-c} are
the same in $\mathsf{A}$ and $\mathsf{B}$. The entries M2--M4 of the list
$\mathfrak{M}^{\mathrm{rd}}$ of comparison maps are indexed by exactly the units as assigned in (g),
and its entry M8 by the Wilson class taken cell unit by cell unit; there a cell unit is written
$(c,i,\Delta)$, and the letter $v$, when used for it in that entry, does not denote the near unit
of (d). The list makes no assignment of its own.
\end{definition}

\begin{proof}
\emph{The support of $\mathsf{E}$ and $D_v$.} The function $\mathsf{E}$ is defined on the bonds of
$\mathfrak{B}^*_j$. If $c\in\operatorname{supp}\mathsf{E}$, then $\mathsf{E}(c)\ne0$, so by (E0)
the inclusion $\tilde B(c)\subset\tilde\square^3$ fails and $c\in\mathfrak{S}_\square$.
Conversely, if $\operatorname{supp}\mathsf{E}\subset\mathfrak{S}_\square$ and
$\tilde B(c)\subset\tilde\square^3$, then $c\notin\mathfrak{S}_\square$, hence
$c\notin\operatorname{supp}\mathsf{E}$ and $\mathsf{E}(c)=0$; the two forms of (E0) in
\eqref{5.5:eq-response-box-majorants-b} are therefore equivalent. Since
$\operatorname{supp}\mathsf{E}\subset\mathfrak{S}_\square$, the infimum defining
$d_*(\mathfrak{r}_v,\mathfrak{S}_\square)$ runs over a set of pairs containing the set over which
the infimum defining $d_*(\mathfrak{r}_v,\operatorname{supp}\mathsf{E})$ runs (the latter being
$+\infty$ when $\operatorname{supp}\mathsf{E}$ is empty); hence
$D_v\le d_*(\mathfrak{r}_v,\operatorname{supp}\mathsf{E})$. The bonds of $\mathfrak{B}^*_j$, the
sets $\tilde B(c)$ and $\tilde\square^3$ and the own box $\mathfrak{r}_v$ are scaffold objects of
Notation~\ref{5.5:not-one-scaffold}, so $\mathfrak{S}_\square$ and $D_v$ are the same for both
cutoffs.

\emph{The majorant of the $F$-part.} By \eqref{5.5:eq-response-box-majorants-a},
\begin{equation*}
|B_F|_*(b)=d\,(1+|b-y|)\,\Theta\,\mathfrak{a}_nt^2=d\,(1+|b-y|)\,|F|_*
=\tfrac12\cdot 2d\,(1+|b-y|)\,\Theta\,\mathfrak{a}_nt^2=\tfrac12|B^{\mathrm{ax}}|_*(b),
\end{equation*}
which gives \eqref{5.5:eq-response-box-majorants-2}, with equality.

\emph{The rim majorant.} For a bond $b$ of $\mathfrak{r}_v$ one has $1+|b-y|\le D$ by (d).
The pointwise form carries the factor $1+|b-y|$ in the place of the factor $D$ of
$|\mathfrak{e}^{\mathrm{rim}}|_*$; multiplying by the non-negative number
$4d\,\Theta\,\mathfrak{a}_nt^2$ gives
\begin{equation*}
4d\,(1+|b-y|)\,\Theta\,\mathfrak{a}_nt^2\le 4dD\,\Theta\,\mathfrak{a}_nt^2
=|\mathfrak{e}^{\mathrm{rim}}|_*,
\end{equation*}
so the constant $|\mathfrak{e}^{\mathrm{rim}}|_*$ dominates on $\mathfrak{r}_v$ the pointwise form
in which $D$ stands as $1+|b-y|$.

\emph{The constant $|B^0|_*$.} The own box $\mathfrak{r}_v$ has finitely many bonds, so the
supremum in \eqref{5.5:eq-response-box-majorants-d} is attained, and the definition of $S^*_v$
gives $\sup_{b\in\mathfrak{r}_v}|B^{\mathrm{ax}}|_*(b)=\alpha_{3,j}/S^*_v$. Substituting this in
\eqref{5.5:eq-response-box-majorants-1} yields $|B^0|_*=1.42\,\alpha_{3,j}/(LS^*_v)$. For every
bond $b$ of $\mathfrak{r}_v$,
\begin{equation*}
1.2\,L^{-1}|B^{\mathrm{ax}}|_*(b)\le 1.2\,L^{-1}\sup_{b'\in\mathfrak{r}_v}|B^{\mathrm{ax}}|_*(b')
\le 1.42\,L^{-1}\sup_{b'\in\mathfrak{r}_v}|B^{\mathrm{ax}}|_*(b')=|B^0|_*,
\end{equation*}
so $|B^0|_*$ dominates the pointwise bound of clause (ii) of the remainder-decay lemma.

\emph{Independence of the cutoff.} The predicates of (g) involve the level $j$, the anchor $y$
and its level $\ell(y)$, $n_0$, the index $\nu$, the letter of
Definition~\ref{5.5:def-slots-letters}, the division of near units into cores, tails and singles,
the own box $\mathfrak{r}_v$, the radius $\alpha_{3,j}$ and, through
\eqref{5.5:eq-response-box-majorants-a}, $\Theta=c_Ww^\vee(y)$, $\mathfrak{a}_n$, $t$ and $d$;
the radii of \eqref{5.5:eq-response-box-majorants-c} involve $r^\sharp_\square$, $\delta$,
$d_\ell(X_T,\square^\zeta)$, $K^+$, $K^{\mathrm{rem}}$, $\vartheta(\square)$, and the bound on
$\tau$ involves $\mathsf{w}_0$. All of these are objects or constants of the scaffold of
Notation~\ref{5.5:not-one-scaffold}, with $w^\vee$ a maximum of frozen label weights, and that
scaffold is one and the same for $\mathsf{A}$ and $\mathsf{B}$. Hence the assignment (g) and the
radii coincide for the two cutoffs.
\end{proof}

\begin{definition}[The list of comparison moves]\label{5.5:def-move-list}
We work at the site (P) of the operation $\mathbf{R}_k$ (Definition~\ref{5.1:def-expansion-rule}),
with the stage slots, leaf
slots, unit slots, letters and one-cutoff classes of Definition~\ref{5.5:def-slots-letters}, the
moves of Definition~\ref{5.5:def-moves}, and the common parameters, kernel, response dilation,
box pairs, weighted scaffold majorants and filing of
Notation~\ref{5.5:not-response-box-majorants}, in particular
\eqref{5.5:eq-response-box-majorants-a}, \eqref{5.5:eq-response-box-majorants-c} and
\eqref{5.5:eq-response-box-majorants-d}. Every object below is built in each cutoff
$\mathsf{m}\in\{\mathsf{A},\mathsf{B}\}$ on that cutoff's own lattice, and $p=(P^{\mathsf{B}},P^{\mathsf{A}})$
denotes a pair in the product class of the parent slot.

We use three conventions throughout. In a parameter domain, the factor $(b,\sigma)$ stands for
the range of the common parameters $b$ and $\sigma$ of
Notation~\ref{5.5:not-response-box-majorants}. For a configuration $P$ and a domain $X$,
$P\lceil X$ denotes the restriction of $P$ to $X$, and $p\lceil X:=(P^{\mathsf{B}}\lceil X,
P^{\mathsf{A}}\lceil X)$. For a slot $\mathfrak{s}$, $\mathcal{K}^{\mathrm{rd}}(\mathfrak{s})$ denotes the
production closure at $\mathfrak{s}$, which is defined after this definition from the list
$\mathfrak{M}^{\mathrm{rd}}$; below it enters only in $\mathrm{M5}$(b) and in $\mathrm{M8}$ (the
description of $\mathsf{S}^\perp_v$ and the last two lines of \eqref{5.5:eq-move-list-e}).

The \emph{list of comparison moves} $\mathfrak{M}^{\mathrm{rd}}$ consists of the moves
$\mathrm{M1}^{M}$, $\mathrm{M2}$, $\mathrm{M3}$, $\mathrm{M4}$, $\mathrm{M5}$, $\mathrm{M6}$, $\mathrm{M7}$ and
$\mathrm{M8}:=\mathsf{M}^\perp$ described below. As in Definition~\ref{5.5:def-moves}, each is one
construction applied in both cutoffs, does not depend on the complex sources $z$, and has a
single parameter domain. That its image lies in the class of the target slot is not part of
the definition.

\medskip
\noindent$\mathrm{M1}^{M}$, the all-blocks response dilation (stage to stage). For every stage slot
$(k;\mathsf{n}+1,\hat X)$ this move goes to $(k;\mathsf{n},\hat X)$. It is given by
\begin{equation}\label{5.5:eq-move-list-a}
\begin{gathered}
\bigl((b,\sigma,(t_\square)_\square);p\bigr)\longmapsto
\bigl(\Phi^{\mathsf{B}}_t(P^{\mathsf{B}}),\Phi^{\mathsf{A}}_t(P^{\mathsf{A}})\bigr),
\qquad (k;\mathsf{n}+1,\hat X)\to(k;\mathsf{n},\hat X),\\
\mathfrak{P}_{\mathrm{M1}}:=\Bigl\{(t_\square)_\square:\ |t_\square|\le R^\natural_\square
=\mathsf{w}_0\,r^\sharp_\square\ \text{if}\ \tilde\square\subset\hat X,\quad
t_\square:=0\ \text{otherwise}\Bigr\}.
\end{gathered}
\end{equation}
Here $\Phi_t$ is the response dilation of \eqref{5.5:eq-response-box-majorants-c} with
$\chi_t=\sum_\square t_\square\zeta_\square$, and $\mathfrak{P}_{\mathrm{M1}}$ is the polydisc of
Corollary~\ref{5.2:cor-response-dilation}. The stage classes are the source classes
$\mathfrak{S}^{\mathsf{m}}(\mathfrak{s})$ of Definition~\ref{5.5:def-slots-letters}. One and the same
parameter $(b,\sigma,(t_\square))$ is used in both cutoffs.

\medskip
\noindent$\mathrm{M2}$, the box read (stage to unit). Fix a block $\square$ and a parent stage slot
$(k;\mathsf{n}+1,\hat X)$ with $\hat X\supseteq\hat{\mathsf{K}}_\square\supseteq\square_0$. The
parameter domain is $(b,\sigma)\times[0,1]\times\mathbb{D}_\square$. The map $\rho^{\mathrm{box}}$
is the first line of \eqref{5.5:eq-move-list-b}; here
$\mathfrak{b}'_j=\mathsf{P}^0$ is the box pair of \eqref{5.5:eq-response-box-majorants-c}, at
interpolation parameter $t$ and box amplitude $\mathsf{t}$. The targets are of two kinds.
\begin{itemize}
\item[$(\flat)$] Let $v$ be a near unit of one of the classes that carry the letter $\flat$ in
Definition~\ref{5.5:def-slots-letters}, other than the $\partial$-terms, that is box-read at
$\square$ in the filing
\eqref{5.5:eq-response-box-majorants-d}. For each constituent $X\subset\mathfrak{r}_v$ of
$v$ there is one move, with target $(j,X;\flat)$. If $v$ is a core, every member lies on the one
box.
\item[$(\circ)$] For every near Wilson or counterterm piece at $\square$, at every $\nu$ and every
value of $S^*_v$, there is one move, with target $(j,X;\circ)$.
\end{itemize}

\medskip
\noindent$\mathrm{M3}$, the own-level read (stage to unit). Fix a block $\square$ and a pair
$(T,\square)$ filed own-read or far in \eqref{5.5:eq-response-box-majorants-d}. The parent is
$(k;\mathsf{n}+1,\hat X)$ with $\hat X\supseteq X_T$. The parameter domain is
$(b,\sigma)\times[0,1]\times\{|\tau|\le\mathsf{w}_0 r_T(\square)\}$. The map
$\rho^{\mathrm{own}}_{T,\square}$ is the second line of \eqref{5.5:eq-move-list-b}, and the target
is the slot $(\ell_T,X_T;\mathfrak{L}_T)$ of the term $T$.

\medskip
\noindent$\mathrm{M4}$, the remainder read (stage to unit). Fix a block $\square$ and a piece filed
on the box pair, that is, an element of the index set of $\mathrm{M2}$ with either letter. The
parent is the same as for $\mathrm{M2}$. The parameter domain is
$(b,\sigma)\times[0,1]\times\{|\mathsf{t}|\le r^{\mathrm{rem}}_\square\}$. The map $\rho^{\mathrm{rem}}$ is
the third line of \eqref{5.5:eq-move-list-b}, with $\mathsf{P}^1(t,\mathsf{t})$ the second box pair of
\eqref{5.5:eq-response-box-majorants-c}, and the target is $(j,X;\mathfrak{L})$. This move is the
first move of the lemma on the remainders of box-pair reads, by which every piece read
on the box pair has a remainder. For such a piece $T$, the stored remainder is
\begin{equation*}
\int_0^1 dt\,\partial_{\mathsf{t}}\big|_{\mathsf{t}=0}\bigl[T(\mathsf{P}^1)-T(\mathsf{P}^0)\bigr].
\end{equation*}
Here $T(\mathsf{P}^0)$ is read through $\mathrm{M2}$ and $T(\mathsf{P}^1)$ through $\mathrm{M4}$. The
expressions \cite[(3.21), (3.22)]{B-CMP109} are the jet of this remainder. The three reads are
\begin{equation}\label{5.5:eq-move-list-b}
\begin{aligned}
\rho^{\mathrm{box}}&:\ \bigl((b,\sigma,t,\mathsf{t});p\bigr)\longmapsto
\bigl(\mathfrak{b}'^{\mathsf{B}}_j\lceil X,\ \mathfrak{b}'^{\mathsf{A}}_j\lceil X\bigr),\\
\rho^{\mathrm{own}}_{T,\square}&:\ \bigl((b,\sigma,t,\tau);p\bigr)\longmapsto
\bigl(\mathcal{W}^{\mathsf{B}}_T(t,\tau),\ \mathcal{W}^{\mathsf{A}}_T(t,\tau)\bigr),\\
\rho^{\mathrm{rem}}&:\ \bigl((b,\sigma,t,\mathsf{t});p\bigr)\longmapsto
\bigl(\mathsf{P}^{1,\mathsf{B}}(t,\mathsf{t})\lceil X,\ \mathsf{P}^{1,\mathsf{A}}(t,\mathsf{t})\lceil X\bigr).
\end{aligned}
\end{equation}

\medskip
\noindent$\mathrm{M5}$, paired charts (stage to unit). This move has two parts.

(a) \emph{Comb dilations} through a box-pair read $\rho\in\{\mathrm{M2},\mathrm{M4}\}$, at the unit's
own box $\square_0'=\mathfrak{r}_v$. Let $(B_i)$ be one of the splittings (S1)--(S4) below of the
read's comb datum; the chart parameters $(\sigma_i)$ attached to it are not to be confused with
the $\sigma$-fields of the common parameters. Each $B_i$ comes with its scaffold majorant $|B_i|_*$ of
\eqref{5.5:eq-response-box-majorants-a}, and the move asks, in each cutoff, for the pointwise bound
$|B_i^{\mathsf{m}}|\le|B_i|_*$; as in Definition~\ref{5.5:def-moves}, this is a one-cutoff statement
asked of the move and is not part of this definition. With $\pi_\rho$ the parameters of $\rho$, the
move goes from the stage parent of $\rho$ to the unit slot $(j,X;\mathfrak{L})$ of the read; it is
the map (a) of \eqref{5.5:eq-move-list-c}, and its parameter domain $\mathfrak{P}$ is the line below
that map. Here $(U_j,J_j)(\square_0',V)$ denotes the level-$j$ map $(U,J)(\mathfrak{B}^*_j,\cdot)$ of
\eqref{5.5:eq-response-box-majorants-c}, by which $\mathsf{P}^0$ and $\mathsf{P}^1$ are defined,
applied over the box $\square_0'$ to a configuration $V$. The splittings are
the following.
\begin{itemize}
\item[(S1)] $B^{\mathrm{ax}}$ alone, with the single chart parameter $\sigma_1$. At the letter
$\flat$, where $\alpha_3(\flat)=\alpha_{3,j}$, the
parameter domain
\begin{equation*}
|\sigma_1|\,\sup_b|B^{\mathrm{ax}}|_*(b)<\alpha_3(\flat)
\end{equation*}
is, by the definition of $S^*_v$ in \eqref{5.5:eq-response-box-majorants-d}, the disc
$|\sigma_1|<S^*_v$ of the lemma on the $\sigma$-chart of $B^{\mathrm{ax}}$; on it the splitting is
taken through $\mathrm{M2}$. At the letter $\circ$ the
splitting is taken through $\mathrm{M2}$ at every read core, on the disc
$|\sigma_1|<\alpha_3(\circ)/\sup_b|B^{\mathrm{ax}}|_*(b)$; this is the chart of the pointed Wilson
family $\hat A$.
\item[(S2)] The F-part $B_F$ together with the terms $\mathfrak{e}^{\mathrm{in}}$ and
$\mathfrak{e}^{\mathrm{rim}}$ on the graded rim, as these are defined in the lemma on the rim
family. The F-part is read at the majorant
$|B_F|_*$ of \eqref{5.5:eq-response-box-majorants-a}.
\item[(S3)] The pointed family's $B_0$, at $|B_0|_*$.
\item[(S4)] The canonical splitting $(B^0,\Delta_c)$ of the comb datum of $\mathsf{P}^1$ on
$\mathfrak{r}_v$, taken through $\rho:=\mathrm{M4}$. Its base is $\mathsf{P}^0$, the read of
$\mathrm{M2}$ at the same parameters $\pi$, identified with
$(U_j,J_j)(\mathfrak{r}_v,e^{iB^0})$. Here $B^0$ is the comb datum of
$\mathsf{P}^0\lceil\mathfrak{r}_v$ from $y$, and $B^1$ is the comb datum of
$\mathsf{P}^1\lceil\mathfrak{r}_v$ itself. The second member is $\Delta_c$, defined in the line (S4)
of \eqref{5.5:eq-move-list-c}. By the lemma giving the canonical splitting, the data $B^0$,
$B^1$ and $\Delta_c$ are functions of $(\mathsf{P}^0,\mathsf{P}^1,y)$ alone. The majorants are
$|B^0|_*$ and $|\Delta|_*$ of \eqref{5.5:eq-response-box-majorants-a}.
\end{itemize}

(b) \emph{Linear families from $\mathbb{1}$}. The move is the map (b) of
\eqref{5.5:eq-move-list-c}, with the parameter domain displayed in the line below it and target
$(j,X;\mathfrak{L})$. Here $\xi$ is the scaling factor with which the classical directions of the
correlation theorem (Theorem~\ref{5.1:thm-correlation-theorem}) are exponentiated in the first
component. This move is listed at every unit slot $(j,X;\mathfrak{L})$ of
Definition~\ref{5.5:def-slots-letters}: $X$ is any small-field domain of level $j$, whether or not
it is a constituent of a read's unit. It is this $X$ that enters the weight
$\|\mathsf{h}_b\|_*=B_3e^{-\delta_0\operatorname{dist}_j(b,X)}$ of
\eqref{5.5:eq-response-box-majorants-a}, so that the closure inclusion for $\mathrm{M5}$(b), the
analogue for this move of the last line of \eqref{5.5:eq-move-list-e}, places these families in
$\mathcal{K}^{\mathrm{rd}}(j,X;\mathfrak{L})$ at every domain $X$ of level $j$, non-constituent $X$
included. The $\mathsf{h}_i$ range over the classical directions of the correlation theorem
(Theorem~\ref{5.1:thm-correlation-theorem}): the box direction, the free direction and their
difference. They also range over the contractions $\mathsf{h}[B]$ of these with a box-pair read's
comb data
$B\in\{B^{\mathrm{ax}},B_0,\mathfrak{e}^{\mathrm{rim}},B^{\mathrm{ax}}+B_0\}$, taken at
$\|\mathsf{h}[B]\|_*$. In the case of a contraction, the move starts from the stage parent of
that read.
\begin{equation}\label{5.5:eq-move-list-c}
\begin{aligned}
&\text{(a)}\quad\bigl((\pi_\rho,(\sigma_i));p\bigr)\longmapsto\Bigl((U_j,J_j)^{\mathsf{m}}
\bigl(\square_0',e^{i\sum_i\sigma_iB_i^{\mathsf{m}}}\bigr)\lceil X\Bigr)_{\mathsf{m}\in\{\mathsf{B},\mathsf{A}\}},\\
&\qquad\mathfrak{P}:=\mathfrak{P}_\rho\times\Bigl\{(\sigma_i):\ \sum_i|\sigma_i|\,|B_i|_*
<\alpha_3(\mathfrak{L})\Bigr\},\\
&\text{(S4)}\quad\Delta_c:=B^1-B^0,\qquad\text{with majorants}\ |B^0|_*,\ |\Delta|_*,\\
&\text{(b)}\quad(\sigma_i)\longmapsto\Bigl(\bigl(e^{i\xi\sum_i\sigma_i\mathsf{h}^{\mathsf{m}}_i},\
\textstyle\sum_i\sigma_i\mathbf{j}^{\mathsf{m}}_i\bigr)\Bigr)_{\mathsf{m}\in\{\mathsf{B},\mathsf{A}\}},\\
&\qquad\sum_i|\sigma_i|\,\|(\mathsf{h}_i,\mathbf{j}_i)\|_*<\mathfrak{r}(\mathfrak{L}).
\end{aligned}
\end{equation}

\medskip
\noindent$\mathrm{M6}$, restriction. This is the move $p\mapsto p\lceil X''$ in the two cases of
\eqref{5.5:eq-move-list-d}.
\begin{itemize}
\item[(i)] Stage to stage: from $(k;\mathsf{n}',X')$ to $(k;\mathsf{n}'',X'')$, with
$\mathsf{n}''\le\mathsf{n}'$ and $X''\subset X'$.
\item[(ii)] Stage to leaf or unit: from $(k;\mathsf{n}',\hat X)$ to $(\ell_T,X_T;\mathfrak{L}_T)$, with
$X_T\subset\hat X$. The leaf slots of Definition~\ref{5.5:def-slots-letters} are included among the
targets.
\end{itemize}
Restriction from a unit slot to a unit slot is not an entry of the list. The restriction property
of the stage classes, that is, the statement that the stage classes are stable under restriction,
enters only stage to stage and stage to leaf.
\begin{equation}\label{5.5:eq-move-list-d}
\begin{aligned}
p\mapsto p\lceil X''&:\ (k;\mathsf{n}',X')\to(k;\mathsf{n}'',X''),
\qquad \mathsf{n}''\le\mathsf{n}',\ X''\subset X',\\
p\mapsto p\lceil X_T&:\ (k;\mathsf{n}',\hat X)\to(\ell_T,X_T;\mathfrak{L}_T),\qquad X_T\subset\hat X.
\end{aligned}
\end{equation}

\medskip
\noindent$\mathrm{M7}$, the reads and step maps at the sites (T) and (L) (see
Definition~\ref{5.1:def-expansion-rule}). These moves are the reads of the correlation theorem into
the domains of earlier units, in the order of production, which is well-founded. A unit slot or
leaf slot is the parent of the reads of its own production formula, evaluated at its points.
They enter the list by this description; their formulae are those of the production formulae
themselves. They are taken together with the two seed families of those reads, which likewise
enter the list by reference to their definition and are not written out here.

\medskip
\noindent$\mathrm{M8}:=\mathsf{M}^\perp$, the comb-dilated remainder read (stage to unit), in the format
of Definition~\ref{5.5:def-moves}. Fix a block $\square$ and a cell-unit $v=(c,i,\Delta)$ of the
third class of the filing \eqref{5.5:eq-response-box-majorants-d} at $\square$, that is, that class
taken cell-unit by cell-unit, with piece $A(c,\cdot)$. Here $v$ is a cell-unit in the sense of the
definition of the cell-units of this class: $c$ is the index of the piece $A(c,\cdot)$ of the unit
within the class, $i$ is a level and
$\Delta$ a cell of level $i$. The parent is $(k;\mathsf{n}+1,\hat X)$ with
$\hat X\supseteq\hat{\mathsf{K}}_\square$, as for $\mathrm{M2}$ and $\mathrm{M4}$, and the target slot
is $(i,\Delta;\circ)$. The parameter domain $\mathfrak{P}_\perp$ and the map are given in
\eqref{5.5:eq-move-list-e}.

In each cutoff $\mathsf{m}$, the configuration $U^{\mathsf{m}}_s$ is the comb dilation at amplitude
$s$, in the sense of the definition of the comb-dilated configuration of a cell-unit, applied to
the pair
\begin{equation*}
\bigl(W^{\mathsf{m}},W^{1,\mathsf{m}}\bigr)
=\bigl(U^{\mathsf{m}}_{\mathsf{P}^0}(t,\mathsf{t}),\,U^{\mathsf{m}}_{\mathsf{P}^1}(t,\mathsf{t})\bigr),
\end{equation*}
on the lattice of that cutoff, where $U^{\mathsf{m}}_{\mathsf{P}^0}(t,\mathsf{t})$ and
$U^{\mathsf{m}}_{\mathsf{P}^1}(t,\mathsf{t})$ denote the bond-variable components of the box pairs
$\mathsf{P}^0(t,\mathsf{t})$ and $\mathsf{P}^1(t,\mathsf{t})$ of
\eqref{5.5:eq-response-box-majorants-c} in the cutoff $\mathsf{m}$.
It is restricted to the enlarged cell $\Delta^+$ of that definition.
The construction is
the same in $\mathsf{A}$ and $\mathsf{B}$, because the lexicographic comb from the base point of
$\Delta$ is part of the scaffold. It does not depend on $z$, and it has one parameter domain. The
radius $R^\perp_v:=(32\kappa^*_v)^{-1}$ of \eqref{5.5:eq-move-list-e} is a scaffold number in the
sense of \eqref{5.5:eq-one-scaffold-a}: $\kappa^*_v$ is the scaffold bound, given among the
constants of the definition of the comb-dilated configuration of a cell-unit, of the constant
$\kappa_v$ of that construction, so that $\kappa_v\le\kappa^*_v$, and $\kappa^*_v$ is a number fixed
after $(d,L,\delta)$ and before $\gamma$.

The number $\mathsf{S}^\perp_v$ is the $\delta$-size of this piece at the $\circ$-slot. The class
$\mathcal{K}^{\mathrm{rd}}(i,\Delta;\circ)$ over which its supremum is taken receives the images of
$\mathrm{M2}$ and $\mathrm{M4}$ at the letter $\circ$ and of $\mathsf{M}^\perp$. The last line of
\eqref{5.5:eq-move-list-e} is called the \emph{closure inclusion for $\mathsf{M}^\perp$}.
\begin{equation}\label{5.5:eq-move-list-e}
\begin{aligned}
\mathfrak{P}_\perp&:=(b,\sigma)\times[0,1]\times\{|\mathsf{t}|\le r^{\mathrm{rem}}_\square\}
\times\{|s|\le R^\perp_v\},\\
\mathsf{M}^\perp&:\ \bigl((b,\sigma,t,\mathsf{t},s);p\bigr)\longmapsto
\bigl(U^{\mathsf{B}}_s\lceil\Delta^+,\ U^{\mathsf{A}}_s\lceil\Delta^+\bigr),\qquad
R^\perp_v:=(32\kappa^*_v)^{-1},\\
\mathsf{S}^\perp_v&:=\mathsf{S}^\circ[A(c,\cdot)]
=\sup_{p\in\mathcal{K}^{\mathrm{rd}}(i,\Delta;\circ)}|\delta A(c,p)|,\\
&\mathsf{M}^\perp\bigl(\mathfrak{P}_\perp\times\mathcal{K}^{\mathrm{rd}}(k;\mathsf{n}+1,\hat X)\bigr)
\subset\mathcal{K}^{\mathrm{rd}}(i,\Delta;\circ).
\end{aligned}
\end{equation}
\end{definition}

\begin{definition}[The production closure of the comparison moves and its sizes]\label{5.5:def-production-closure}
Slots (stage slots $\mathfrak{s}=(k;\mathsf{n},\hat X)$, leaf slots and unit slots
$\mathfrak{u}=(j,X;\mathfrak{L})$ with $\mathfrak{L}\in\{\flat,\circ\}$) are those of
Definition~\ref{5.5:def-slots-letters}. Moves $\mu\colon\mathfrak{s}'\to\mathfrak{s}''$, their parameter
domains $\mathfrak{P}_\mu$ and pair maps $\mu^\times=(\mu^{\mathsf{B}},\mu^{\mathsf{A}})$ are those of
Definition~\ref{5.5:def-moves}. The list $\mathfrak{M}^{\mathrm{rd}}$ of comparison moves and the seeds are
those of Definition~\ref{5.5:def-move-list}. A \emph{family of slot classes} assigns to every slot
$\mathfrak{s}$ a set $\mathcal{K}(\mathfrak{s})$ of pairs $p=(P^{\mathsf{B}},P^{\mathsf{A}})$ over the
domain of $\mathfrak{s}$. Families are ordered by inclusion at every slot.
In (k0), (k1) below and in the definition of $\mathsf{S}^\circ$ the letter $\mathfrak{s}$ ranges over all
slots, stage, leaf or unit.

The \emph{production closure} $\mathcal{K}^{\mathrm{rd}}:=\operatorname{cl}^{\mathrm{pr}}(\mathfrak{M}^{\mathrm{rd}})$
is the smallest family of slot classes with the properties (k0), (k1), (k2) of
\eqref{5.5:eq-production-closure-a}.
\begin{equation}\label{5.5:eq-production-closure-a}
\begin{aligned}
&\text{(k0)}\quad \mathbb{1}\in\mathcal{K}^{\mathrm{rd}}(\mathfrak{s})\ \text{for every slot }\mathfrak{s};\\
&\text{(k1)}\quad \mu^\times\bigl(\mathfrak{P}_\mu\times\mathcal{K}^{\mathrm{rd}}(\mathfrak{s}')\bigr)
\subset\mathcal{K}^{\mathrm{rd}}(\mathfrak{s}'')\ \text{for every }\mu\in\mathfrak{M}^{\mathrm{rd}},\
\mu\colon\mathfrak{s}'\to\mathfrak{s}'';\\
&\text{(k2)}\quad \text{every seed lies in the class of its slot};\\
&\mathsf{S}^{\mathrm{rd}}_j[F]:=\sup_X e^{\kappa d_j(X)}
\sup_{p\in\mathcal{K}^{\mathrm{rd}}(j,X;\flat)}\ \sup_{y,w,u}\,\lvert\delta F(X;p)\rvert,\\
&\mathsf{a}:=\mathsf{S}^{\mathrm{rd}}_j[\hat A]\ \text{(over }\mathcal{K}^{\mathrm{rd}}(j,X;\circ)\text{)},\qquad
\mathsf{S}^\circ[F]:=\sup_{p\in\mathcal{K}^{\mathrm{rd}}(\mathfrak{s})}\lvert\delta F(p)\rvert .
\end{aligned}
\end{equation}
Each move in (k1) is applied where its one-machine construction is
defined; by clause~(iii) of the statement on the domains of the comparison moves, a chain of moves never
leaves that domain.

For a $\flat$-family $F$ at level $j$, the size $\mathsf{S}^{\mathrm{rd}}_j[F]$ is defined in
\eqref{5.5:eq-production-closure-a}. Here $\kappa$ is Ba{\l}aban's decay constant, $d_j(X)$ is the
distance function of the level-$j$ domain $X$, $\delta F(X;p)$ is the two-cutoff difference of $F$ at
$p$, and the innermost supremum runs over the remaining variables $y,w,u$ of $\delta F(X;p)$. The size
$\mathsf{a}$ is $\mathsf{S}^{\mathrm{rd}}_j[\hat A]$ for the pointed Wilson family $\hat A$, with the
classes $\mathcal{K}^{\mathrm{rd}}(j,X;\circ)$ in place of $\mathcal{K}^{\mathrm{rd}}(j,X;\flat)$. For a
$\circ$-family $F$ at a slot $\mathfrak{s}$, the size $\mathsf{S}^\circ[F]$ is also defined in
\eqref{5.5:eq-production-closure-a}.

The size $\mathsf{S}^\perp_v$ that enters the data of the move $\mathsf{M}^\perp$ in
\eqref{5.5:eq-move-list-e} is therefore
$\mathsf{S}^\circ[A(c,\cdot)]$ at the slot $(i,\Delta;\circ)$ of the cell-unit $v=(c,i,\Delta)$. The
a priori bound
\begin{equation*}
\mathsf{S}^\perp_v\le 2\max\{\mathfrak{n}^\perp_v,\mathfrak{n}_{\hat A}\}
\end{equation*}
is clause (iii) of the lemma bounding the comb-dilated remainder read. It is not part of this
definition.

Let $\mathfrak{M}^{\mathrm{pr}}$ be the list of moves used in the proof of the correlation theorem, as fixed
in the theorem on its production closure, and let
$\operatorname{cl}^{\mathrm{pr}}_J:=\operatorname{cl}^{\mathrm{pr}}(\mathfrak{M}^{\mathrm{pr}})$
be that closure. Write $\mathsf{S}^{\mathrm{cl}}_j[F]$ and $\mathsf{a}^{\mathrm{cl}}$ for the sizes defined as in
\eqref{5.5:eq-production-closure-a} with $\operatorname{cl}^{\mathrm{pr}}_J$ in place of
$\mathcal{K}^{\mathrm{rd}}$. Then, at every slot $\mathfrak{s}$,
\begin{equation}\label{5.5:eq-production-closure-1}
\mathcal{K}^{\mathrm{rd}}(\mathfrak{s})\subset\operatorname{cl}^{\mathrm{pr}}_J(\mathfrak{s}),\qquad
\mathsf{S}^{\mathrm{rd}}_j[F]\le\mathsf{S}^{\mathrm{cl}}_j[F],\qquad \mathsf{a}\le\mathsf{a}^{\mathrm{cl}} .
\end{equation}
Equality in \eqref{5.5:eq-production-closure-1} is not asserted here.
\end{definition}

\begin{proof}
\emph{Description by chains.} A \emph{chain} ending at a slot $\mathfrak{s}$ consists of the following
data:
\begin{itemize}
\item a starting point $p_0$ at a slot $\mathfrak{s}_0$, where $p_0$ is either $\mathbb{1}$ or a seed of
$\mathfrak{s}_0$;
\item moves $\mu_i\colon\mathfrak{s}_{i-1}\to\mathfrak{s}_i$ in $\mathfrak{M}^{\mathrm{rd}}$ for
$1\le i\le m$, with $m\ge0$ and $\mathfrak{s}_m=\mathfrak{s}$;
\item parameters $\pi_i\in\mathfrak{P}_{\mu_i}$.
\end{itemize}
The chain determines the points $p_i:=\mu_i^\times(\pi_i,p_{i-1})$, and its endpoint is $p_m$. Each move
is applied where its one-machine construction is defined; by the clause~(iii) invoked in the definition,
the chain never leaves that domain.

Let $\mathcal{K}'(\mathfrak{s})$ be the set of endpoints of chains ending at $\mathfrak{s}$. The chains
with $m=0$ give (k0) and (k2) for $\mathcal{K}'$. Let $p\in\mathcal{K}'(\mathfrak{s}')$ be the endpoint
of a chain, let $\mu\colon\mathfrak{s}'\to\mathfrak{s}''$ be in $\mathfrak{M}^{\mathrm{rd}}$, and let
$\pi\in\mathfrak{P}_\mu$. Appending $(\mu,\pi)$ to that chain gives a chain ending at $\mathfrak{s}''$
with endpoint $\mu^\times(\pi,p)$. This is (k1) for $\mathcal{K}'$.

Conversely, let $\mathcal{K}$ be any family with (k0), (k1), (k2). We show by induction on $i$ that
$p_i\in\mathcal{K}(\mathfrak{s}_i)$. For $i=0$ this is (k0) or (k2). The step from $i-1$ to $i$ is (k1)
applied to $\mu_i$, $\pi_i$ and $p_{i-1}$. Hence $\mathcal{K}'\subset\mathcal{K}$.

Thus $\mathcal{K}'$ has (k0)--(k2) and lies below every family that has them. It is therefore the
smallest such family, and $\mathcal{K}^{\mathrm{rd}}=\mathcal{K}'$. In words, the points of
$\mathcal{K}^{\mathrm{rd}}$ are the points obtained from a seed, or from $\mathbb{1}$, by a finite
chain of moves of $\mathfrak{M}^{\mathrm{rd}}$ at parameters in their parameter domains.

\emph{The inclusion.} The chain description has the same form for any list of moves with given seeds.
Let $\mathfrak{M}\subset\mathfrak{M}'$ be two lists whose common moves have the same maps and the same
parameter domains, and suppose the two lists have the same seeds. Then every chain of $\mathfrak{M}$ is a
chain of $\mathfrak{M}'$, so $\operatorname{cl}^{\mathrm{pr}}(\mathfrak{M})(\mathfrak{s})$ is contained
in $\operatorname{cl}^{\mathrm{pr}}(\mathfrak{M}')(\mathfrak{s})$ at every slot. In this sense the
closure is monotone in the list.

The list $\mathfrak{M}^{\mathrm{pr}}$ contains every move of $\mathfrak{M}^{\mathrm{rd}}$, the
comb-dilated remainder read $\mathsf{M}^\perp$ included, with the same maps, parameter domains and seeds.
Monotonicity therefore gives the first statement of \eqref{5.5:eq-production-closure-1}. This inclusion is
also the first clause of the theorem on the production closure of the correlation theorem, which asserts
$\operatorname{cl}^{\mathrm{pr}}_J(\mathfrak{s})\supseteq\mathcal{K}^{\mathrm{rd}}(\mathfrak{s})$ at every
slot, and hence $\mathsf{S}^{\mathrm{rd}}_j[F]\le\mathsf{S}^{\mathrm{cl}}_j[F]$.

\emph{The inequalities.} For each $X$, the weight $e^{\kappa d_j(X)}$ and the integrand
$\lvert\delta F(X;p)\rvert$ are the same in $\mathsf{S}^{\mathrm{rd}}_j[F]$ and in
$\mathsf{S}^{\mathrm{cl}}_j[F]$. By the inclusion, the supremum over
$p\in\mathcal{K}^{\mathrm{rd}}(j,X;\flat)$ runs over a subset of
$\operatorname{cl}^{\mathrm{pr}}_J(j,X;\flat)$. It is therefore at most the supremum over the larger
set. Taking $\sup_X$ gives $\mathsf{S}^{\mathrm{rd}}_j[F]\le\mathsf{S}^{\mathrm{cl}}_j[F]$. The same
argument, applied to $\hat A$ with the letter $\circ$, gives $\mathsf{a}\le\mathsf{a}^{\mathrm{cl}}$.
\end{proof}

\begin{definition}[The one-cutoff membership hypotheses]\label{5.5:def-membership-hypotheses}
Let $\mathsf{m}\in\{\mathsf{A},\mathsf{B}\}$ denote either of the two machines of
display~\eqref{5.5:eq-one-scaffold-a}, each on its own lattice. Slots, letters and
one-machine classes are those of Definition~\ref{5.5:def-slots-letters}. Response maps,
box pairs, parameter domains and weighted scaffold majorants are those of
Notation~\ref{5.5:not-response-box-majorants}. The moves are those of
Definition~\ref{5.5:def-move-list}. The production closure $\mathcal{K}^{\mathrm{rd}}$
and its sizes are those of Definition~\ref{5.5:def-production-closure}. For a
configuration or map and a set $X$ in its domain, ${}\upharpoonright X$ denotes
restriction to $X$.

The closure statement of this section is proved under the four hypotheses (a)--(d)
below. Together they are called the \emph{one-cutoff membership hypotheses}. None of
them is proved in this section. Hypotheses (a)--(c) concern one machine at a time.
Hypothesis~(d) concerns the production closure of the moves of the correlation theorem,
Theorem~\ref{5.1:thm-correlation-theorem}.
\begin{enumerate}
\item[(a)] The conventions of display~\eqref{5.5:eq-one-scaffold-a} and of
display~\eqref{5.5:eq-response-box-majorants-d} are in force. The index of the move
$\mathsf{M}^\perp$ of display~\eqref{5.5:eq-move-list-e} runs over the cell-units of the
class $W$, written
\begin{equation}\label{5.5:eq-membership-hypotheses-a}
v=(c,i,\Delta),
\end{equation}
with $c$, $i$ and the cell $\Delta$ as in the definition of the class $W$ and its
instances.

\item[(b)] In each machine $\mathsf{m}$ the following four conditions hold.
\begin{enumerate}
\item[(s)] The seeds of Definition~\ref{5.5:def-production-closure} lie in the classes
of their slots, and $\mathbb{1}$ lies in every class.
\item[(m1)] Let $\Phi^{\mathsf{m}}_t$ be the response dilation of the move
$\mathrm{M1}^{M}$, and let $\mathfrak{P}_{\mathrm{M1}}$ be its polydisc, as in
display~\eqref{5.5:eq-move-list-a}. Then the first inclusion of
\eqref{5.5:eq-membership-hypotheses-b} below holds.
\item[(m6)] Whenever restriction from the stage slot $(k;\mathsf{n}',X')$ to the stage
slot $(k;\mathsf{n}'',X'')$ is defined, it maps $\mathfrak{S}^{\mathsf{m}}(k;\mathsf{n}',X')$
into $\mathfrak{S}^{\mathsf{m}}(k;\mathsf{n}'',X'')$. Let $T$ be any term read at the slot
$(k;\mathsf{n}',\hat X)$, with level $\ell_T$, domain $X_T$ and letter
$\mathfrak{L}_T$. For every such $T$ the second inclusion of
\eqref{5.5:eq-membership-hypotheses-b} holds. The target class carries the room parameter
of the anchor's level on the active zone and the radii of the term classes on their
region.

The same holds at those slots of Definition~\ref{5.5:def-slots-letters} that belong to
terms read at the site (P) with a letter of their own; there the inclusion is into that
letter.

If $\mathfrak{L}_T=\circ$, the image of the restriction obeys a scaffold plaquette bound
on the support of the piece.
\begin{equation}\label{5.5:eq-membership-hypotheses-b}
\begin{gathered}
\Phi^{\mathsf{m}}_t\bigl(\mathfrak{P}_{\mathrm{M1}}
\times\mathfrak{S}^{\mathsf{m}}(k;\mathsf{n}+1,\hat X)\bigr)
\subset\mathfrak{S}^{\mathsf{m}}(k;\mathsf{n},\hat X),\\
\mathfrak{S}^{\mathsf{m}}(k;\mathsf{n}',\hat X)\upharpoonright X_T
\subset\mathfrak{U}^{\mathsf{m}}(\ell_T,X_T;\mathfrak{L}_T).
\end{gathered}
\end{equation}
\item[(m7)] At every move $\mathrm{M7}$ of display~\eqref{5.5:eq-move-list-d}, the
letter of the target slot contains the image. Consider a target of letter $\circ$: the
pieces of the fifth row of the expansion of Definition~\ref{5.5:def-slots-letters} and
the counterterm pieces, read at the sites
(T) and (L). There the image obeys a scaffold plaquette bound on the support of the
piece, namely the entry bound of part (b) of the Wilson-family lemma of the small-field
analysis.
\end{enumerate}

\item[(c)] Let $\mathsf{m}$ be a machine, let
$P^{\mathsf{m}}\in\mathfrak{S}^{\mathsf{m}}(k;\mathsf{n}+1,\hat X)$, and let the
parameters lie in the domain $\mathfrak{P}$ of the move concerned. Then the following
conditions hold at the site (P).
\begin{enumerate}
\item[$(\mathrm{m2})^\flat$] Let $v$ be a box-read unit. On $[0,1]\times\mathbb{D}_\square$
the relations of the first line of \eqref{5.5:eq-membership-hypotheses-c} hold.
Pointwise, the bounds of its second line hold.
\begin{equation}\label{5.5:eq-membership-hypotheses-c}
\begin{gathered}
\mathfrak{b}'^{\mathsf{m}}_j\upharpoonright X\in\mathfrak{U}^{\mathsf{m}}(j,X;\flat),
\qquad
\mathfrak{b}'^{\mathsf{m}}_j\upharpoonright\mathfrak{r}_v
\cong(U_j,J_j)\bigl(\mathfrak{r}_v,e^{iB^{\mathrm{ax},\mathsf{m}}}\bigr),\\
|B^{\mathrm{ax},\mathsf{m}}|\le|B^{\mathrm{ax}}|_*,\quad
|F^{\mathsf{m}}|\le|F|_*,\quad
|\mathfrak{e}^{\mathsf{m}}_\mu|\le|\mathfrak{e}_\mu|_*,\quad
|\mathfrak{e}^{\mathrm{rim},\mathsf{m}}|\le|\mathfrak{e}^{\mathrm{rim}}|_*.
\end{gathered}
\end{equation}
Consequently the $F$-part of condition (S2) of display~\eqref{5.5:eq-move-list-c}
obeys, for every bond $b$ of $\mathfrak{r}_v$,
$|B_F^{\mathsf{m}}(b)|\le|B_F|_*(b)$.
\item[$(\mathrm{m2})^\circ$] $\mathfrak{b}'^{\mathsf{m}}_j$ is defined and holomorphic on
$[0,1]\times\mathbb{D}_\square$. On the cells meeting $\tilde\square^3$ it obeys the
entry bound
\begin{equation*}
\bigl|\partial\mathfrak{b}'^{\mathsf{m}}_j(x)-1\bigr|\le 4\theta_2(x).
\end{equation*}
Here $\theta_2$ denotes the entry-bound profile on the cells meeting $\tilde\square^3$.
Moreover $|\hat A^{\mathsf{m}}(\mathfrak{b}'^{\mathsf{m}}_j)|\le\mathfrak{n}_{\hat A}$.
The latter bound is the consequence of three facts: the entry bound on
$\tilde\square^3$, the regularity of the input configuration $P^{\mathsf{m}}$ off
$\square_0$, and parts (b) and (c) of the Wilson-family lemma of the small-field analysis.

Further conditions hold at every read core $v$ with $\nu\ge3$. This includes every value
of $S^*_v$, and both box-read and own-read cores. On its own box
$\mathfrak{r}_v=\tilde\square^{\,\nu+1}$, in the comb gauge at $y$:
\begin{itemize}
\item the identity of $(\mathrm{m2})^\flat$ holds;
\item the comb data are
\begin{equation*}
B^{\mathrm{ax},\mathsf{m}}_\mu(x)=\sum_{\lambda<\mu}(x-y)_\lambda
F^{\mathsf{m}}_{\lambda\mu}(y)+\mathfrak{e}^{\mathsf{m}}_\mu(x),
\qquad
\mathfrak{e}=\mathfrak{e}^{\mathrm{in}}+\mathfrak{e}^{\mathrm{rim}};
\end{equation*}
\item the four pointwise bounds of \eqref{5.5:eq-membership-hypotheses-c} hold.
\end{itemize}
So at the counterterm piece of a read core of own-read type, the chart of $\hat A$ is the
listed move $\mathrm{M5}(a)$ at the letter $\circ$ taken through $\mathrm{M2}$ at the letter
$\circ$; see displays~\eqref{5.5:eq-move-list-c} and~\eqref{5.5:eq-move-list-b}. At
$\nu\le2$ no comb identity and no chart is part of this hypothesis; there the value of
the corresponding item is obtained from the closure property (k1) of
Definition~\ref{5.5:def-production-closure} at the move $\mathrm{M2}$ at the letter
$\circ$, and its chart factor is assigned to the low-order variation bound at the second
cube size at $\mathrm{M2}$.
\item[(m3)] Let $(T,\square)$ be an own-read or far pair. For $t\in[0,1]$,
$|\tau|\le\mathsf{w}_0r_T(\square)$ and $|\sigma|\le e^{\kappa_1}$,
\begin{equation*}
\mathcal{W}^{\mathsf{m}}_T(t,\tau)\upharpoonright X_T
\in\mathfrak{U}^{\mathsf{m}}(\ell_T,X_T;\flat)
=\bigcap_{T'}\mathfrak{U}^{\mathsf{m}}_{T'} .
\end{equation*}
The configuration restricted here is the single configuration
$\mathcal{W}^{\ell_T,\square,\mathsf{m}}(t,\tau)\upharpoonright X_T$. Its membership is
required for every $T'$ at the slot, with the same map and the same disc.
\item[$(\mathrm{m3})^\circ$] At far pieces of the class $W$,
$\mathcal{W}^{\mathsf{m}}_T(t,\tau)$ is defined and holomorphic on the same parameter
domain. Its image obeys a scaffold plaquette bound on the support of the piece.
\item[(m4)] Let $v$ be a box-read unit. For
$(t,\mathsf{t})\in[0,1]\times\{|\mathsf{t}|\le r^{\mathrm{rem}}_\square\}$, holomorphically,
\begin{equation*}
\mathsf{P}^{1,\mathsf{m}}(t,\mathsf{t})\upharpoonright X\in\mathfrak{U}^{\mathsf{m}}(j,X;\flat).
\end{equation*}
At the canonical data the box pair has the charts
\begin{equation*}
\begin{gathered}
\mathsf{P}^{0,\mathsf{m}}\upharpoonright\mathfrak{r}_v
\cong(U_j,J_j)\bigl(\mathfrak{r}_v,e^{iB^{0,\mathsf{m}}}\bigr),\\
\mathsf{P}^{1,\mathsf{m}}\upharpoonright\mathfrak{r}_v
\cong(U_j,J_j)\bigl(\mathfrak{r}_v,e^{i(B^{0,\mathsf{m}}+\Delta_c^{\mathsf{m}})}\bigr).
\end{gathered}
\end{equation*}
For every bond $b$ the following bounds hold pointwise:
\begin{equation*}
|B^{0,\mathsf{m}}(b)|\le 1.2\,L^{-1}|B^{\mathrm{ax}}(b)|_*\le|B^0|_*,
\qquad
|\Delta_c^{\mathsf{m}}|\le|\Delta|_*.
\end{equation*}
In the weighted scaffold majorants $|B^0|_*$ and $|\Delta|_*$ of
Notation~\ref{5.5:not-response-box-majorants}, the distance $D_v$ is the geometric
distance. At the letter $\circ$ the requirement is instead that, at every unit of the
class $W$, $\mathsf{P}^{1,\mathsf{m}}$ is defined and holomorphic and $\hat A$ is bounded
on it.
\item[$(\mathrm{m5a})$] Let $B'$ be complex on $\square_0'$ with
$|B'|<\alpha_3(\mathfrak{L})$. Then, analytically,
\begin{equation*}
(U_j,J_j)^{\mathsf{m}}\bigl(\square_0',e^{iB'}\bigr)\upharpoonright X
\in\mathfrak{U}^{\mathsf{m}}(j,X;\mathfrak{L}).
\end{equation*}
At $\mathfrak{L}=\circ$ the letter is the absolute one. This membership comes with three
sets of bounds:
\begin{itemize}
\item the pointwise bounds of $(\mathrm{m2})^\flat$;
\item the read-core bounds of $(\mathrm{m2})^\circ$;
\item the bounds of (m4).
\end{itemize}
The first two serve conditions (S1)--(S3) of display~\eqref{5.5:eq-move-list-c}. The
bounds of (m4) serve condition (S4).
\item[$(\mathrm{m5b})$] At the radius $\mathfrak{r}(\mathfrak{L})$ the conclusion of the
lemma giving the chart of $(U_j,J_j)$ holds. In addition,
$\|\mathsf{h}^{\mathsf{m}}_b\|\le\|\mathsf{h}_b\|_*$.
\item[$(\mathrm{m}\perp)^\circ$] This condition concerns the move $\mathsf{M}^\perp$ at
every cell-unit $v=(c,i,\Delta)$ of the class $W$ at $\square$. The configuration
$U^{\mathsf{m}}_s\upharpoonright\Delta^+$ of display~\eqref{5.5:eq-move-list-e} is
defined, where $\Delta^+$ is the enlargement of $\Delta$ fixed with the move
$\mathsf{M}^\perp$ in display~\eqref{5.5:eq-move-list-e}. It is built from
$\bigl(U^{\mathsf{m}}_{\mathsf{P}^0}(t,\mathsf{t}),U^{\mathsf{m}}_{\mathsf{P}^1}(t,\mathsf{t})\bigr)$,
which is the box pair of the moves $\mathrm{M2}$ and $\mathrm{M4}$. It is holomorphic in
$(\mathbb{U},\mathbb{J},b,\sigma,\mathsf{t})$, entire in $s$ and continuous in $t$.

For $|s|\le R^\perp_v$, the image of $U^{\mathsf{m}}_s$ obeys the plaquette bound of
part (i) of the lemma on the move $\mathsf{M}^\perp$ on the plaquettes of $\Delta$, and
\begin{equation*}
|A(c,U^{\mathsf{m}}_s)|\le\mathfrak{n}^\perp_v .
\end{equation*}
This is the entry bound of the letter $\circ$. The letter $\circ$ carries no tube, so no
comparison with $\alpha_3(\circ)$ enters.
\end{enumerate}

\item[(d)] The following statements are taken by statement and are not proved in this
section. They concern the production closure $\operatorname{cl}^{\mathrm{pr}}_J$ of the
move list $\mathfrak{M}^{\mathrm{pr}}$ of the correlation theorem,
Theorem~\ref{5.1:thm-correlation-theorem}, which contains $\mathcal{K}^{\mathrm{rd}}$.
\begin{enumerate}
\item[(d-i)] At every slot $\mathfrak{s}$ the first relation of
\eqref{5.5:eq-membership-hypotheses-d} holds, the containment holding through the unit
slots. Hence, for every family $F$, the first inequality of the second relation holds;
the second inequality, $\mathsf{S}^{\mathrm{cl}}_j[F]\le2\mathfrak{n}$, is part of the
first of the statements taken here.
\begin{equation}\label{5.5:eq-membership-hypotheses-d}
\mathcal{K}^{\mathrm{rd}}(\mathfrak{s})\subset\operatorname{cl}^{\mathrm{pr}}_J(\mathfrak{s}),
\qquad
\mathsf{S}^{\mathrm{rd}}_j[F]\le\mathsf{S}^{\mathrm{cl}}_j[F]\le2\mathfrak{n}.
\end{equation}
The same statement includes a product principle and a diagonal test.
\item[(d-ii)] The product statement for the rim data of the two machines holds on
$\operatorname{cl}^{\mathrm{pr}}_J$.
\item[(d-iii)] The constant $\mathsf{l}^J$ is defined by a single formula as the maximum
of two sizes on $\operatorname{cl}^{\mathrm{pr}}_J$. The first of them includes
$\sup_v\mathsf{S}^\perp_v$ and the per-cube contributions of the move $\mathsf{M}^\perp$.
\item[(d-iv)] The closure statements on $\operatorname{cl}^{\mathrm{pr}}_J$ hold: closure
under products, the two-cutoff Lipschitz step, and heredity. Each holds under its listed
hypotheses, which include inputs read at the sites (T) and (P).
\end{enumerate}
The smallness of $\mathsf{l}^J$ is not part of (d). None of the statements on which the
correlation theorem rests uses it; it enters the uniqueness theorem.
\end{enumerate}
\end{definition}

\begin{theorem}[Closure of the two-cutoff comparison family under the listed moves]
\label{5.5:thm-family-closure}
Let $\mathfrak{M}^{\mathrm{rd}}$ be the list of comparison moves of
Definition~\ref{5.5:def-move-list}. Let
$\mathcal{K}^{\mathrm{rd}}=\operatorname{cl}^{\mathrm{pr}}(\mathfrak{M}^{\mathrm{rd}})$ be its
production closure, with the sizes $\mathsf{S}^{\mathrm{rd}}_j[F]$, $\mathsf{a}$ and
$\mathsf{S}^\circ[F]$ and the closure $\operatorname{cl}^{\mathrm{pr}}_J$ with its sizes
$\mathsf{S}^{\mathrm{cl}}_j[F]$, all as in Definition~\ref{5.5:def-production-closure}. Stage
slots $\mathfrak{s}=(k;\mathsf{n},\hat X)$, unit slots $\mathfrak{u}=(j,X;\mathfrak{L})$ and
their letters $\mathfrak{L}\in\{\flat,\circ\}$ are those of
Definition~\ref{5.5:def-slots-letters}. Moves are those of Definition~\ref{5.5:def-moves}, and
the site objects and radii are those of Notation~\ref{5.5:not-response-box-majorants}. Assume
hypotheses (H1), (H2) and (H3) of Definition~\ref{5.5:def-membership-hypotheses}. Then the
following hold.
\begin{enumerate}
\item[(i)] Closure. $\mathcal{K}^{\mathrm{rd}}$ is a correspondence for
$\mathfrak{M}^{\mathrm{rd}}$ in the sense of the two-cutoff comparison principle. That is,
$\mathcal{K}^{\mathrm{rd}}$ satisfies the conditions (k0), (k2) and (k1) of
\eqref{5.5:eq-production-closure-a} at every
$\mu\in\mathfrak{M}^{\mathrm{rd}}$. In particular the following closure properties hold; each is
an implication, not a definition:
\begin{equation}\label{5.5:eq-family-closure-a}
\begin{aligned}
&\text{restriction property of the stage classes}
&&\Longleftarrow\ \text{(k1) at M6},\\
&\text{closure under the response dilation}
&&\Longleftarrow\ \text{(k1) at }\mathrm{M1}^{M},\\
&\text{closure under the response maps at (P)}
&&\Longleftarrow\ \text{(k1) at M2, at M3 and at M4},\\
&\text{closure at own-read pairs}
&&\Longleftarrow\ \text{(k1) at M3 at near own-read pairs},\\
&\text{closure at the weighted majorants}
&&\Longleftarrow\ \text{(k1) at M5},\\
&\text{the two remainder closures}
&&\Longleftarrow\ \text{(k1) at M4 and}\\
&&&\qquad\text{(k1) at M5(a) with (S4)},\\
&\text{the closure hypothesis along }\mathsf{M}^\perp
&&=\ \text{(k1) at M8}.
\end{aligned}
\end{equation}
The lines of \eqref{5.5:eq-family-closure-a} mean the following.
\begin{itemize}
\item The restriction property of the stage classes is taken in the form in which it enters the
definition of a pair system at the stage sites and the arguments built on it. It follows from
(k1) at M6.
\item The closure of the stage classes under the response dilation is taken in the form of the
pair-system definition, which in the notation of clause (v)(e) reads
\begin{equation*}
\Phi^\times_t\bigl(b;\mathcal{K}^{\mathrm{rd}}(k;\mathsf{n}+1,\hat X)\bigr)
\subset\mathcal{K}^{\mathrm{rd}}(k;\mathsf{n},\hat X).
\end{equation*}
Here $b$ and the remaining stage parameters $\mathsf{s}$ of the pair-system definition lie in
the polydisc of the stage, and one and the same
$(b,t,\mathsf{s})$ is used in both cutoffs. The block amplitudes are taken at the radius of
$\mathfrak{P}_{\mathrm{M1}}$ in \eqref{5.5:eq-move-list-a}, in place of the radius
$|t_\square|\le r_\square$ of that definition. That is, we take
$|t_\square|\le R^\natural_\square=\mathsf{w}_0 r^\sharp_\square$ if
$\tilde\square\subset\hat X$, and $t_\square:=0$ otherwise. Accordingly $\vartheta$ is replaced
by $\vartheta^\natural$ in the two estimates of the pair-system argument that carry
$K_A\vartheta(\square)$. That is, $K_A\vartheta(\square)$ is replaced by
$K_A\vartheta^\natural(\square)=2/R^\natural_\square$; this costs a factor $1/\mathsf{w}_0$,
which is fixed before $\gamma$ is chosen. This closure follows from (k1) at $\mathrm{M1}^{M}$.
\item The closure under the response maps concerns every response map that the proof of the
correlation theorem (Theorem~\ref{5.1:thm-correlation-theorem}) runs at the site (P) on small
units. These maps are:
\begin{itemize}
\item the box pairs
$(\mathfrak{b}'^{\mathsf{B}}_j,\mathfrak{b}'^{\mathsf{A}}_j)\!\upharpoonright_{X_v}$ on
$[0,1]\times\mathbb{D}_\square$;
\item for far small units, the pairs $\mathcal{W}^\times_T$ with
$|\tau|\le\mathsf{w}_0 r_T(\square)$;
\item for the remainders of box-read units, the pairs $(\mathsf{P}^1)^\times$ on
$[0,1]\times\{|\mathsf{t}|\le r^{\mathrm{rem}}_\square\}$.
\end{itemize}
One and the same $(b,t,\mathsf{t},\sigma)$ is used in both cutoffs. The closure states that
$\mathcal{K}^{\mathrm{rd}}(k;\mathsf{n}+1,\hat X)$ is carried into the unit class over which the
size $\mathsf{S}^{\mathrm{rd}}$ is the supremum. It follows from (k1) at M2, at M3 and at M4.
\item The closure at own-read pairs is used in the two statements on near units read at their
own level. It follows from (k1) at M3 at the near own-read pairs.
\item The closure at the weighted majorants is taken in the form used in the weighted-majorant
lemma for box reads. There it is read at the weighted majorants of
Notation~\ref{5.5:not-response-box-majorants}, so that it is (k1) at M5. It covers the
$\sigma$-disc, the items along the families produced by the entries $(\mu2)$ and $(\mu3)$ of the
two-cutoff comparison principle, and the rim item from $\mathbb{1}$. The passage of $(\mu3)$
through a read is removed at (P). This closure follows from (k1) at M5.
\item The two closures used in the remainder bound for box-read units follow from (k1) at M4 and
from (k1) at M5(a) together with (S4) of \eqref{5.5:eq-move-list-c}.
\item The comparison theorem along the comb dilation $\mathsf{M}^\perp$ has one closure
hypothesis. It is (k1) at M8, where $\mathrm{M8}=\mathsf{M}^\perp$ as in
\eqref{5.5:eq-move-list-e}.
\end{itemize}

\item[(ii)] Typing. Every $\mu\in\{\mathrm{M1}^{M},\mathrm{M2},\dots,\mathrm{M6},\mathrm{M8}\}$
is a move in the sense of Definition~\ref{5.5:def-moves}:
\begin{itemize}
\item it has one parameter domain $\mathfrak{P}_\mu$, whose radii are fixed under (H1);
\item $\mu^{\mathsf{B}}$ and $\mu^{\mathsf{A}}$ are the same construction;
\item the pair map
\begin{equation}\label{5.5:eq-family-closure-b}
\mu^\times(\pi;p)=\bigl(\mu^{\mathsf{B}}(\pi;P^{\mathsf{B}}),\,
\mu^{\mathsf{A}}(\pi;P^{\mathsf{A}})\bigr),\qquad
\pi\in\mathfrak{P}_\mu,\quad p=(P^{\mathsf{B}},P^{\mathsf{A}}),
\end{equation}
is jointly holomorphic in $(P^{\mathsf{m}},\pi)$ for complex $\pi$, and continuous in the real
parameters.
\end{itemize}
At M5(a) and M5(b) the holomorphy is in the chart parameters $(\sigma_i)$ at fixed
$(\pi_\rho;p)$, where $\rho\in\{\mathrm{M2},\mathrm{M4}\}$ is the read of which M5 is a chart
and $\pi_\rho\in\mathfrak{P}_\rho$ its parameter, as in \eqref{5.5:eq-move-list-c}; the data of
the read $\rho$ enter through (m2) and (m4) of \eqref{5.5:eq-membership-hypotheses-c}. A unit
carries the same read in both cutoffs, because the assignment of reads is part of the source-free
scaffold. At real parameters and $\sigma\equiv1$, the reads assigned in each cutoff reproduce
that cutoff's own expansion of every piece. At M8 this holds at the end points $s=0$ and $s=1$.

\item[(iii)] Containment along chains; finiteness; the transfer principle.
\begin{itemize}
\item Every point of $\mathcal{K}^{\mathrm{rd}}(\mathfrak{s})$ lies in
$\mathcal{U}^\times(\mathfrak{s})$.
\item Every point of $\mathcal{K}^{\mathrm{rd}}(j,X;\flat)$ lies in
$\mathfrak{U}^{\mathsf{B}}(j,X;\flat)\times\mathfrak{U}^{\mathsf{A}}(j,X;\flat)$.
\item At every point of $\mathcal{K}^{\mathrm{rd}}(j,X;\circ)$ the $\circ$-families are defined
and holomorphic along the move that produces the point, and $\hat A$ is bounded by
$\mathfrak{n}_{\hat A}$. Every other $\circ$-family is defined and holomorphic there. Its bound
by the number supplied with it is not part of this theorem.
\item Every move along every chain is applied inside its domain of definition.
\end{itemize}
Let $F$ be a $\flat$-family with $\|F^{\mathsf{m}}\|_j\le\mathfrak{n}$ in each cutoff, in the
weighted norm $\|F\|_j:=\sup_X e^{\kappa d_j(X)}|F(X,\cdot)|$ on the tube of $F$. Then
\begin{equation}\label{5.5:eq-family-closure-c}
\mathsf{S}^{\mathrm{rd}}_j[F]\le 2\mathfrak{n}<\infty,
\qquad
\mathsf{a}\le 2\mathfrak{n}_{\hat A}.
\end{equation}
The transfer principle holds for every $\mu\in\mathfrak{M}^{\mathrm{rd}}\colon
\mathfrak{s}'\to\mathfrak{s}''$ at every $p\in\mathcal{K}^{\mathrm{rd}}(\mathfrak{s}')$. That
is, for every pair $f=(f^{\mathsf{B}},f^{\mathsf{A}})$ of one-cutoff functionals at
$\mathfrak{s}''$ to which the transfer principle of the two-cutoff comparison principle applies,
the function
\begin{equation*}
\delta(\pi):=f^{\mathsf{B}}\bigl(\mu^{\mathsf{B}}(\pi;P^{\mathsf{B}})\bigr)
-f^{\mathsf{A}}\bigl(\mu^{\mathsf{A}}(\pi;P^{\mathsf{A}})\bigr)
\end{equation*}
is holomorphic on $\mathfrak{P}_\mu$, and
\begin{equation*}
|\delta|\le\sup_{\mathcal{K}^{\mathrm{rd}}(\mathfrak{s}'')}|f^{\mathsf{B}}-f^{\mathsf{A}}|.
\end{equation*}
The holomorphy is meant as follows at three moves: at M5, in $(\sigma_i)$ at fixed $\pi_\rho$; at
M7, through the holomorphic maps of (m7) of \eqref{5.5:eq-membership-hypotheses-b}; at M8, in
$(b,\sigma,\mathsf{t})$, and $\delta$
is entire in $s$. Consequently the following carry over to $\delta$, with $\mathfrak{N}$
replaced by $\sup|\delta|$: every Cauchy inequality and every inequality of Schwarz's lemma that
holds for a one-cutoff function $g$ with bound $\mathfrak{N}$, and every identity
$\partial^i g(\pi_0)=0$ that holds in each cutoff.

\item[(iv)$'$] Structure; production order. The moves $\mathrm{M1}^{M}$, M2, \dots, M6 and M8
have as parents stage slots of the operation $\mathbf{R}_k$ that runs them. The exception is the
families of M5(b) that start from $\mathbb{1}$, which have no parent. Consequently each of the
closure under the response maps at (P), the closure at own-read pairs and the closure at the
weighted majorants in \eqref{5.5:eq-family-closure-a} is closure under one move. For
$\mu\in\{\mathrm{M2},\mathrm{M3},\mathrm{M4},\mathrm{M5},\mathrm{M8}\}$ with stage parent
$\mathfrak{s}$ and unit target $\mathfrak{u}$, this reads
\begin{equation}\label{5.5:eq-family-closure-d}
\mu^\times\bigl(\mathfrak{P}_\mu\times\mathcal{K}^{\mathrm{rd}}(\mathfrak{s})\bigr)
\subset\mathcal{K}^{\mathrm{rd}}(\mathfrak{u}).
\end{equation}
The only moves of $\mathfrak{M}^{\mathrm{rd}}$ out of a unit slot or a leaf slot are the reads of
M7 in that unit's own production formula. Chains are finite and visit slots in strictly
decreasing production order. Hence no chain that leaves a stage class
$\mathcal{K}^{\mathrm{rd}}(k;\mathsf{n}+1,\hat X)$ through one of M2, \dots, M5, M8 returns to a
stage class of the same $\mathbf{R}_k$. In other words, closing under M2, \dots, M5 and M8 does
not enlarge the stage classes.

\item[(v)] Identifications.
\begin{itemize}
\item[(a)] The inclusion below is an identification of the objects as they are read, not an
identity of the two constructions term by term. Take the class $\mathcal{K}^\natural$ of the
two-cutoff comparison principle in the form in which it equals
$\operatorname{cl}^{\mathrm{pr}}_J$. Read its entries at the site (P) as M1--M6, as follows.
\begin{itemize}
\item The entry $(\mu1)$ at (P) is given by the maps $\Phi_t$, $\mathfrak{b}'_j$,
$\mathcal{W}_T$, $\mathsf{P}^1$ of Notation~\ref{5.5:not-response-box-majorants}.
\item The a priori majorants of the read in $(\mu2)$ and $(\mu3)$ are the weighted majorants of
that notation.
\item The passage of $(\mu3)$ through a read is dropped at (P).
\item $(\mu4)$ is M6.
\item The move M8 has no entry there. It enters $\operatorname{cl}^{\mathrm{pr}}_J$ as one
further generating move, the dilation of M8.
\end{itemize}
With this reading,
\begin{equation}\label{5.5:eq-family-closure-e}
\mathcal{K}^{\mathrm{rd}}\subset\mathcal{K}^\natural=\operatorname{cl}^{\mathrm{pr}}_J,
\qquad
\mathsf{S}^{\mathrm{rd}}_j[F]\le\mathsf{S}^{\mathrm{cl}}_j[F].
\end{equation}
Hence the transfer principle, and the four further statements of the two-cutoff comparison
principle on $\mathcal{K}^\natural$ that hypothesis (H4) of
Definition~\ref{5.5:def-membership-hypotheses} lists on $\operatorname{cl}^{\mathrm{pr}}_J$,
are statements about
$\operatorname{cl}^{\mathrm{pr}}_J\supseteq\mathcal{K}^{\mathrm{rd}}$. The two-cutoff sizes they
bound dominate, in the sense
$\mathsf{S}^{\mathrm{rd}}\le\mathsf{S}^{\mathrm{cl}}$. These statements are supplied on
$\operatorname{cl}^{\mathrm{pr}}_J$ by hypothesis (H4) of
Definition~\ref{5.5:def-membership-hypotheses}, and not by this theorem.
\item[(b)] The stage classes of $\mathcal{K}^{\mathrm{rd}}$ form a pair system, in the following
sense. They have the restriction property in the form of the pair-system definition and in the
form in which the pair-system argument reads it. They have the closure under the response
dilation in the form
of the pair-system definition, at the radius of $\mathfrak{P}_{\mathrm{M1}}$; that is, the two
estimates that carry $K_A\vartheta$ are run with $K_A\vartheta$ replaced by
$K_A\vartheta^\natural$. In the step of the pair-system argument that treats leaves and near
groups in the same way with the constant $\mathsf{l}$, two inputs are used: the leaves of
$\mathcal{T}_2$ are treated by the leaf lemma for $\mathcal{T}_2$, and the units of
$\mathcal{T}_1$ by (k1) at M2, M3 and M4.
\item[(c)] Four statements are statements about $\mathcal{K}:=\mathcal{K}^{\mathrm{rd}}$,
$\mathsf{S}:=\mathsf{S}^{\mathrm{rd}}$, $\mathsf{a}:=\mathsf{S}^{\mathrm{rd}}[\hat A]$ and
$\mathsf{S}^\perp_v$:
\begin{itemize}
\item clause (e) of the weighted-majorant lemma for box reads;
\item the statement on near units read at their own level;
\item the remainder bound for box-read units;
\item the comparison theorem along $\mathsf{M}^\perp$, in which $\mathcal{K}:=\mathcal{K}^{\mathrm{rd}}$
and the closure under $\mathsf{M}^\perp$ is assumed.
\end{itemize}
They read $\mathcal{K}$ at seven places:
\begin{enumerate}
\item[(1)] (k1) at M2, for the Cauchy integrals;
\item[(2)] (k1) at M3, for the values at near levels;
\item[(3)] (k1) at M5(a), for the $\sigma$-disc;
\item[(4)] (k1) at M5, for the items along families, with the rim item from $\mathbb{1}$;
\item[(5)] (k1) at M4 with (S4), for the remainders;
\item[(6)] the $\circ$-class, for the term of size $\mathfrak{n}\mathsf{a}$. Here the charts of
radius $O(1)$ of $\hat A$ are M5(a) at radius $\alpha_3(\circ)$ through $\mathrm{M2}^\circ$; at
every read core they are given by the read-core part of (m2)$^\circ$ of
\eqref{5.5:eq-membership-hypotheses-c};
\item[(7)] (k1) at M8, the one closure hypothesis of the comparison theorem along
$\mathsf{M}^\perp$. Its part (a) applies the transfer principle on $\mathfrak{P}_\perp$. Its part
(b) applies Schwarz's lemma in $s$, with $g(0)=0$, on $|s|\le R^\perp_v$, at the comb-dilated
reads
$(U^{\mathsf{B}}_s,U^{\mathsf{A}}_s)\!\upharpoonright_{\Delta^+}$ of the cell-units of class $W$.
Clause (e) of the weighted-majorant lemma at class $W$ follows from the two-cutoff bound for the
cell-units of class $W$ together with (k1) at M8.
\end{enumerate}
The far pieces in the weighted-majorant lemma are read through $\mathrm{M1}^{M}$ and M6 (the
$\mathcal{T}_2$ terms on $\Phi_t$). The statements listed read $\mathcal{K}$ at no other place.
\item[(d)] The class $\mathcal{K}^\sharp$ of the definition of the two-cutoff class over the
field pair systems is built over the field pair systems of the sites,
each with the restriction property and the closure under the response dilation. The stage classes
of $\mathcal{K}^{\mathrm{rd}}$ qualify as such pair systems. The two-cutoff bound on
$\mathcal{K}^\sharp$ stated in the remark following that definition is unchanged.
\item[(e)] Dictionary.
\begin{itemize}
\item The stage classes $\mathcal{K}_{\mathsf{n}}(\hat X)$, $\mathsf{n}\ge1$, of the pair
systems are $\mathcal{K}^{\mathrm{rd}}(k;\mathsf{n},\hat X)$.
\item The leaf class of a term $T$ is $\mathcal{K}^{\mathrm{rd}}(\ell_T,X_T;\mathfrak{L}_T)$.
\item The classes $\mathcal{K}_j(X)$ of the two-cutoff comparison principle are
$\mathcal{K}^{\mathrm{rd}}(j,X;\mathfrak{L})$, with the same letter.
\item In clause (e) of the weighted-majorant lemma, ``$p\in\mathcal{K}_{\mathsf{n}+1}(\hat X)$''
refers to the stage class, and ``the class over which $\mathsf{S}$ is the supremum'' refers to
the unit class.
\item The moves at the site (P) are $\mathrm{M1}^{M}$, M2, \dots, M6, M8. Below a unit only M7
occurs.
\end{itemize}
\end{itemize}
\end{enumerate}
\end{theorem}

The clauses are proved one at a time, beginning with clause (iv)$'$.

\begin{proof}[Proof of clause (iv)$'$]
We read the parent slot of each move off Definition~\ref{5.5:def-move-list}.
\begin{itemize}
\item The move $\mathrm{M1}^{M}$ goes from a stage slot to a stage slot,
$(k;\mathsf{n}+1,\hat X)\to(k;\mathsf{n},\hat X)$, by \eqref{5.5:eq-move-list-a}.
\item The moves M2, M3 and M4 of \eqref{5.5:eq-move-list-b}, and the move M8 of
\eqref{5.5:eq-move-list-e}, go from a stage slot to a unit slot.
\item The move M5(a) of \eqref{5.5:eq-move-list-c} is a chart of a read
$\rho\in\{\mathrm{M2},\mathrm{M4}\}$. Its data $B_i^{\mathsf{m}}$ are functions of
$(\pi_\rho;P^{\mathsf{m}})$. Here $\pi_\rho\in\mathfrak{P}_\rho$ is the parameter of $\rho$, and
$P^{\mathsf{m}}$ is the configuration, in the cutoff $\mathsf{m}$, at the stage slot to which
$\rho$ is applied. Hence the parent of M5(a) is the stage parent of $\rho$.
\item The same holds for the contractions of M5(b). The bare families of M5(b) start from
$\mathbb{1}$ and have no parent slot.
\item The move M6 goes from a stage slot to a stage slot or to a leaf slot.
\end{itemize}
Thus $\mathrm{M1}^{M}$, M2, \dots, M6 and M8 have stage slots of the operation $\mathbf{R}_k$
that runs them as parents, except for the bare families of M5(b), which have none. Consequently
each of (k1) at M2, (k1) at M3, (k1) at M4, (k1) at M5 and (k1) at M8 is one inclusion of the
form \eqref{5.5:eq-family-closure-d}. Here $\mathfrak{s}$ is the stage parent of the move and
$\mathfrak{u}$ is its unit target. This is the first assertion of (iv)$'$.

Consider next a chain that leaves a stage class
$\mathcal{K}^{\mathrm{rd}}(k;\mathsf{n}+1,\hat X)$ through one of M2, \dots, M5, M8. By the
preceding paragraph the chain then stands at a unit slot $(j,X;\mathfrak{L})$. By
Definition~\ref{5.5:def-move-list} (see \eqref{5.5:eq-move-list-d}), the only moves of
$\mathfrak{M}^{\mathrm{rd}}$ out of a unit slot or a leaf slot are the moves M7. An M7 move out of
$(j,X;\mathfrak{L})$ is a read in the production formula of that unit. Its target is therefore
the class of an object produced before the term of $(j,X)$. A fortiori, that object is produced
before stage $\mathsf{n}+1$ of $\mathbf{R}_k$.

The production order is well-founded. Hence every chain is finite and visits slots in strictly
decreasing production order. No later point of a chain as above lies in a stage slot
$(k;\mathsf{n}',\cdot)$ with $\mathsf{n}'\ge\mathsf{n}+1$. Stage slots of sites produced earlier
in the production order than stage $\mathsf{n}+1$ of $\mathbf{R}_k$ may be entered by such a
chain. This is compatible with the assertion, which concerns the stage classes
of the same $\mathbf{R}_k$ only. Therefore closing under M2, \dots, M5 and M8 does not enlarge the
stage classes, which is the second assertion of (iv)$'$.
\end{proof}

The remaining clauses are proved after this, in the order (ii), (iii), (i), (v).

\begin{proof}[Proof of clause~(ii) of Theorem~\ref{5.5:thm-family-closure}]\label{5.5:prf-moves-typed}
Assume hypotheses (H1)--(H3) of Definition~\ref{5.5:def-membership-hypotheses}. The claim is
\eqref{5.5:eq-family-closure-b}. For every entry
\begin{equation*}
\mu\in\{\mathrm{M1}^{M},\mathrm{M2},\mathrm{M3},\mathrm{M4},\mathrm{M5},\mathrm{M6},\mathrm{M8}\}
\end{equation*}
of the list of comparison moves (Definition~\ref{5.5:def-move-list}), the claim has five parts:
\begin{itemize}
\item $\mu$ is a move in the sense of \eqref{5.5:eq-moves-a}, with one parameter domain
$\mathfrak{P}_\mu$ whose radii are numbers of the scaffold \eqref{5.5:eq-one-scaffold-a};
\item its two components $\mu^{\mathsf{B}}$ and $\mu^{\mathsf{A}}$ are the same construction;
\item $\mu$ is jointly holomorphic in $P^{\mathsf{m}}$ and the complex parameters $\pi$, and
continuous in the real parameters. At $\mathrm{M5}$(a) and $\mathrm{M5}$(b) the claim is instead
holomorphy in the chart parameters $(\sigma_i)$ at fixed $(\pi_\rho;p)$, the data of the read
$\rho$ entering through (m2) and (m4) of \eqref{5.5:eq-membership-hypotheses-c};
\item a unit carries the same read at both cutoffs;
\item at real parameters and $\sigma\equiv1$, the reads assigned at cutoff $\mathsf{m}$
reproduce that cutoff's own expansion of every piece.
\end{itemize}
We verify these in five steps.

\emph{Step 1: parameter domains.} Each listed move has the single parameter domain given for it in
Definition~\ref{5.5:def-move-list}. That domain serves both cutoffs. Its radii, where not read
off directly from \eqref{5.5:eq-move-list-b} and \eqref{5.5:eq-move-list-c} (the disc
$\mathbb{D}_\square$ of radius $r^\sharp_\square$ for $\mathrm{M2}$, the radius
$\mathfrak{r}(\mathfrak{L})$ for $\mathrm{M5}$(b); $\mathrm{M6}$ has no parameter), are the
following.

For $\mathrm{M1}^{M}$ the domain is the polydisc of \eqref{5.5:eq-move-list-a}, with radii
$R^\natural_\square=\mathsf{w}_0r_\square^\sharp$ and $\mathsf{w}_0=1/24$.

For $\mathrm{M3}$ the radius is $\mathsf{w}_0r_T(\square)$, with $r_T(\square)$ as in
\eqref{5.5:eq-response-box-majorants-c}.

For $\mathrm{M4}$ and $\mathrm{M8}$ the amplitude radius is the remainder radius
$r^{\mathrm{rem}}_\square$ of \eqref{5.5:eq-response-box-majorants-c}. The constant
$K^{\mathrm{rem}}$ entering it is a number depending on $(d,L,M,\kappa_1)$ only.

For $\mathrm{M8}$, at a cell-unit $v=(c,i,\Delta)$ of the third filing class as in
\eqref{5.5:eq-move-list-e}, the domain is the one displayed there:
\begin{equation*}
\mathfrak{P}_\perp=(b,\sigma)\times[0,1]\times\{|\mathsf{t}|\le r^{\mathrm{rem}}_\square\}
\times\{|s|\le R^\perp_v\},
\end{equation*}
where
\begin{equation}\label{5.5:eq-moves-typed-1}
R^\perp_v=(32\,\kappa^*_v)^{-1},
\end{equation}
and $\kappa^*_v$ is the scaffold constant fixed beside $R^\perp_v$ in
\eqref{5.5:eq-move-list-e}, which carries the factor $e^{-\frac14\delta D_v}$.
The constants in \eqref{5.5:eq-moves-typed-1} are numbers fixed after $(d,L,\delta)$ and before
$\gamma$. None of them depends on an age, on the parameter $N_0$ of the constants fixed with the
comb-dilated remainder read, on a depth, on $K$, on
$\nu$, on $M$ or on a weight. The distance $D_v$ is the geometric distance of
\eqref{5.5:eq-response-box-majorants-b},
which depends on no field configuration. Thus \eqref{5.5:eq-moves-typed-1} consists of scaffold
data.

For $\mathrm{M5}$ the chart radii are taken at the weighted scaffold majorants of
\eqref{5.5:eq-response-box-majorants-a}. They are built from
\begin{equation*}
(d,M,\nu,t,\Theta,\mathfrak{a}_n,c_1,B_3,\delta_0,L,D_v),
\end{equation*}
where $t=L^{j-n}$ is the level ratio of the near unit fixed in
Notation~\ref{5.5:not-response-box-majorants} (not the interpolation parameter), the constants
$c_1$, $B_3$, $\delta_0$ are those of the majorants \eqref{5.5:eq-response-box-majorants-a}, and
$D_v$ is geometric. These are label data and constants.

By \eqref{5.5:eq-one-scaffold-a}, label data and constants are the same for the two cutoffs.
Hence every listed $\mu$ has one domain $\mathfrak{P}_\mu$, and its radii are scaffold numbers.

\emph{Step 2: the filing.} Let $v$ be a near unit or piece. The predicate of the filing
\eqref{5.5:eq-response-box-majorants-d} that decides whether $v$ is box-read, own-read, or a piece
of the third filing class of \eqref{5.5:eq-response-box-majorants-d} (the near Wilson and
counterterm pieces, called class~W below) reads four things:
\begin{itemize}
\item the index $\nu$;
\item the row of the table of units, in the proposition that tabulates them, to which $v$
belongs, and whether $v$ is a core member or a short single;
\item the scale ratio $S^*_v$, evaluated at the scaffold majorant
$\sup_{b\in\mathfrak{r}_v}|B^{\mathrm{ax}}|_*(b)$;
\item whether $v$ is a Wilson piece or a counterterm piece.
\end{itemize}
Each of these is a scaffold datum. By \eqref{5.5:eq-one-scaffold-a}, each is therefore the same
for $\mathsf{A}$ and for $\mathsf{B}$. Consequently every unit carries the same read at both
cutoffs.

\emph{Step 3: the same construction.} We take the moves in turn.

$\mathrm{M1}^{M}$ is, at each cutoff, the response dilation $\Phi_t$ of
\eqref{5.5:eq-response-box-majorants-c}. This is one formula, evaluated on the cutoff's own data.

$\mathrm{M2}$ is, at each cutoff, the chain of \eqref{5.5:eq-response-box-majorants-c}
\begin{equation*}
\mathbb{U}\mapsto\mathbb{W}\mapsto\mathbb{H}^w\mapsto\mathbf{A}\mapsto\mathfrak{b}
\mapsto\mathfrak{b}'_j,
\end{equation*}
followed by restriction to $X$.

$\mathrm{M3}$ is, at each cutoff, the own-level pair $\mathcal{W}_T$ of
Notation~\ref{5.5:not-response-box-majorants}. That is, it is the pair
$\mathcal{W}^{\ell_T,\square}$ restricted to $X_T$.

$\mathrm{M4}$ is, at each cutoff, the chain
\begin{equation*}
\mathbb{W}\mapsto(\mathbb{H}^w,\mathbf{H}_j)\mapsto D^\sharp\mapsto\mathsf{E}\mapsto\mathsf{P}^1,
\end{equation*}
followed by restriction to $X$.

The charts of $\mathrm{M5}$ are the paired charts of \eqref{5.5:eq-move-list-c}. At each cutoff
they are fed that cutoff's own read data. In the splitting (S4), the pair
$(B^0,\Delta_c)^{\mathsf{m}}$ is the canonical datum, in the canonical-splitting lemma, of the
triple $(\mathsf{P}^{0,\mathsf{m}},\mathsf{P}^{1,\mathsf{m}},y)$.

$\mathrm{M6}$ is one restriction formula, applied on each cutoff's own lattice.

$\mathrm{M8}$ is, at each cutoff, the configuration $U_s$ of \eqref{5.5:eq-move-list-e}
restricted to $\Delta^+$. Here $U_s$ is built from the pair
\begin{equation*}
(W,W^1)=(U_{\mathsf{P}^0},U_{\mathsf{P}^1})(t,\mathsf{t}),
\end{equation*}
which is the box pair of $\mathrm{M2}$ and $\mathrm{M4}$, along the lexicographic comb from the
base point of $\Delta$, which is a scaffold datum. Each cutoff does this on its own lattice.

Finally, consider the sets
\begin{equation*}
\square_0,\quad\mathfrak{B}^+(\square_0),\quad\mathfrak{B}^*_j,\quad\mathfrak{B}_j,\quad
\mathfrak{r}_v,\quad\mathfrak{S}_\square.
\end{equation*}
These are paired scaffold sets, one for each cutoff. In particular, at each cutoff
$\operatorname{supp}\mathsf{E}\subset\mathfrak{S}_\square$ by (E0) of
\eqref{5.5:eq-response-box-majorants-b}. Hence $\mu^{\mathsf{B}}$ and $\mu^{\mathsf{A}}$ are the
same construction for every listed $\mu$.

\emph{Step 4: holomorphy.} We again take the moves in turn.

For $\mathrm{M1}^{M}$, holomorphy is the holomorphy of the response dilation $\Phi_t$, as stated
with its formula in the localised corollary on the response dilation at the
$r^\sharp_\square$-polydisc, which rests on the transfer of the analysis of
\cite[\S3]{B-CMP109} to the cover, printed without proof in \cite[p.~276]{B-CMP119}; cited as
printed.

For $\mathrm{M2}$, we combine three facts:
\begin{itemize}
\item the holomorphy statements (a1)--(a4) of the lemma on the window and its box pair;
\item clause (a) of the kernel theorem, which gives $\mathbb{H}$;
\item the form $\mathbf{A}=(t\zeta'_\square+\mathsf{t}\zeta_\square)\mathbb{H}^w$, which is affine
in $(t,\mathsf{t})$.
\end{itemize}

For $\mathrm{M3}$, holomorphy is clause (b) of the lemma on the own-level pair.

For $\mathrm{M4}$, holomorphy is clause (b) of the second-box-pair lemma, together with the
holomorphy clause of the refinement lemma for the remainder read, in which the move
$\mathrm{M4}$ and the splitting (S4) are defined.

For $\mathrm{M5}$, the box functions of $e^{i\sum_i\sigma_iB_i}$ are analytic in $(\sigma_i)$, and
the exponentials are entire. These are the claims of clause~(ii) at $\mathrm{M5}$; all of them are
made at fixed $(\pi_\rho;P^{\mathsf{m}})$. In $\pi_\rho$ the data are read through (m2) and (m4)
of \eqref{5.5:eq-membership-hypotheses-c}.

For $\mathrm{M6}$, the move is linear.

For $\mathrm{M8}$, we use clause (iv) of the lemma on the comb-dilated remainder read. It states
that $U_s$ and $A(c,U_s)$ are holomorphic in $(\mathbb{U},\mathbb{J},b,\sigma,\mathsf{t})$, entire
in $s$, and continuous in $t$. Here $A(c,\cdot)$ denotes the class-W piece at the cell $c$, as in
the definition of $\mathsf{M}^\perp$ in \eqref{5.5:eq-move-list-e}. The
pair $(W,W^1)$ enters through $(\mathrm{m2})^\circ$ and $(\mathrm{m4})^\circ$ of
\eqref{5.5:eq-membership-hypotheses-c}.

\emph{Step 5: the real point.} Take real parameters and $\sigma\equiv1$. We show that the reads
assigned at cutoff $\mathsf{m}$ reproduce that cutoff's expansion of every piece.

First, the fourth statement of the lemma on the own-level pair asserts this reproduction
whatever filing is used. It therefore applies to the filing
\eqref{5.5:eq-response-box-majorants-d}.

The jets are reproduced as follows:
\begin{itemize}
\item for the jet of the first-order response term in the statement on box reads, at box-read
pairs, by clause $(b')$ of that statement;
\item for the same jet at own-read pairs, by clause (c) of the own-read statement of the lemma on
the own-level pair;
\item for the remainder, whose real jet is the expression in \cite[(3.21)]{B-CMP109}, by clause
(c) of the second-box-pair lemma;
\item for the Wilson and counterterm pieces, by clause (c) of the lemma on the block pieces of
class W.
\end{itemize}

Put $\mathfrak{F}_v(s):=A(c,U_s)$. At the endpoints $s=0$ and $s=1$ of $\mathrm{M8}$, clause (ii)
of the lemma on the comb-dilated remainder read states
that $U_0=W$ and $U_1=(W^1)^u$ on the bonds of $\Delta^+$, where $u$ is the comb transport of the
definition of $\mathsf{M}^\perp$. It concludes that
\begin{equation}\label{5.5:eq-moves-typed-2}
\mathfrak{F}_v(0)=A(c,W),\qquad\mathfrak{F}_v(1)=A(c,W^1).
\end{equation}
So at $s=0$ and $s=1$ the dilated cell returns the piece's own values on the box pair.

Together, these statements give each cutoff's own expansion from the reads assigned at that
cutoff, by clause (d1) of the reproduction statement that accompanies the moves between slots
of \eqref{5.5:eq-moves-a}. The filing
\eqref{5.5:eq-response-box-majorants-d} only decides which read carries a piece, and hence which
extension of the piece is carried. It changes no value.

This completes the verification of \eqref{5.5:eq-family-closure-b}.
\end{proof}

\begin{proof}[Proof of clause~(iii) of Theorem~\ref{5.5:thm-family-closure}: containment along chains and
finiteness of the difference sizes]\label{5.5:prf-chain-containment}
Throughout, hypotheses (H1)--(H3) of Definition~\ref{5.5:def-membership-hypotheses} are assumed, in
the form displayed in \eqref{5.5:eq-membership-hypotheses-b} and
\eqref{5.5:eq-membership-hypotheses-c}. The moves are those of Definition~\ref{5.5:def-move-list},
the slots, letters and one-machine classes those of Definition~\ref{5.5:def-slots-letters}, and
$\mathcal{K}^{\mathrm{rd}}$, $\mathsf{S}^{\mathrm{rd}}_j[F]$ and $\mathsf{a}$ are as in
Definition~\ref{5.5:def-production-closure}. Clause~(iii), displayed in
\eqref{5.5:eq-family-closure-c}, has four parts: (a) the points of $\mathcal{K}^{\mathrm{rd}}$ lie in
the products of the one-machine classes, and at $\circ$-slots the families are defined and
holomorphic, with $\hat A$ bounded by $\mathfrak{n}_{\hat A}$; (b) every move along every chain is
applied inside its domain of definition; (c) the two finiteness bounds for
$\mathsf{S}^{\mathrm{rd}}_j[F]$ and for $\mathsf{a}$; (d) the principle for the difference
$\delta(\pi)$ of the two machines along a move.

\emph{Assertions (a) and (b): induction on the length of the chain.} By
Definition~\ref{5.5:def-production-closure}, a point of $\mathcal{K}^{\mathrm{rd}}$ is a seed, or
$\mathbb{1}$, followed by a finite chain of moves of $\mathfrak{M}^{\mathrm{rd}}$ at parameters in
their domains $\mathfrak{P}_\mu$. We prove, by induction on the number $n$ of moves of a chain, that
the point produced by a chain of length $n$ satisfies (a) at its slot, and that each move of the
chain is applied inside its domain of definition.

\emph{Base, $n=0$.} The point is a seed or $\mathbb{1}$. By clause (s) of (H2) in
\eqref{5.5:eq-membership-hypotheses-b}, in each machine the seeds lie in the classes of their slots
and $\mathbb{1}$ lies in every class. Hence (a) holds, and there is no move to check.

\emph{Step.} Let $p=(P^{\mathsf{B}},P^{\mathsf{A}})\in\mathcal{K}^{\mathrm{rd}}(\mathfrak{s}')$ satisfy the
claim at $\mathfrak{s}'$, let $\mu\colon\mathfrak{s}'\to\mathfrak{s}''$ be a move of
$\mathfrak{M}^{\mathrm{rd}}$, and let $\pi\in\mathfrak{P}_\mu$. By the induction hypothesis
$P^{\mathsf{m}}$ lies in the one-machine class of $\mathfrak{s}'$ in machine $\mathsf{m}$. The map
$\mu^{\mathsf{m}}$ is built from minimisers, block averages and kernels
(Notation~\ref{5.5:not-response-box-majorants}, Definition~\ref{5.5:def-move-list}), and these are
constructed at every configuration of that class. Hence $\mu^{\mathsf{m}}(\pi;P^{\mathsf{m}})$ is
defined, so the move is applied inside its domain of definition. It remains to show, in each
machine, that $\mu^{\mathsf{m}}(\pi;P^{\mathsf{m}})$ lies in the class of $\mathfrak{s}''$. We take the
moves in turn.

\emph{The move $\mathrm{M1}^{M}$.} Clause (m1) of (H2) in \eqref{5.5:eq-membership-hypotheses-b}
gives
\begin{equation*}
\Phi^{\mathsf{m}}_t\bigl(\mathfrak{P}_{\mathrm{M1}}\times
\mathfrak{S}^{\mathsf{m}}(k;\mathsf{n}+1,\hat X)\bigr)\subset\mathfrak{S}^{\mathsf{m}}(k;\mathsf{n},\hat X),
\end{equation*}
and the right-hand side is the class of the target slot, see \eqref{5.5:eq-move-list-a}.

\emph{The moves $\mathrm{M6}$ and $\mathrm{M7}$.} Clauses (m6) and (m7) of (H2) give the landing, with
the target read at the letter of the target slot. By (m6), restrictions map source classes into
source classes, and the restriction to $X_T$ maps into
$\mathfrak{U}^{\mathsf{m}}(\ell_T,X_T;\mathfrak{L}_T)$ for every term $T$ at the slot. By (m7), at
every $\mathrm{M7}$ move the class of the letter of the target contains the image.

\emph{The moves $\mathrm{M2}^{\flat}$, $\mathrm{M3}$ and $\mathrm{M4}^{\flat}$.} These land by the
clauses (m2)$^\flat$, (m3) and (m4) of (H3) in \eqref{5.5:eq-membership-hypotheses-c}. For
$\mathrm{M3}$ one more remark is needed. The configuration
$\mathcal{W}^{\ell,\square,\mathsf{m}}\restriction X$ is one and the same for every term $T'$ at the
slot, since the map and the disc depend on $(\ell_{T},X_{T},\square)$ only. Clause (m3) is required for
every such $T'$, so it places this configuration in $\mathfrak{U}^{\mathsf{m}}_{T'}$ for every $T'$.
Hence the configuration lies in the intersection
$\bigcap_{T'}\mathfrak{U}^{\mathsf{m}}_{T'}$, which is $\mathfrak{U}^{\mathsf{m}}(\ell,X;\flat)$ by
Definition~\ref{5.5:def-slots-letters}.

\emph{Landings at $\circ$-slots.} For $\mathrm{M2}^{\circ}$ and $\mathrm{M4}^{\circ}$ the clauses
(m2)$^\circ$ and (m4)$^\circ$ of (H3) apply. They state that the image is defined and holomorphic and
that $\hat A$ is bounded there, by $\mathfrak{n}_{\hat A}$ in the case of (m2)$^\circ$.

For the move $\mathrm{M8}=\mathsf{M}^\perp$ of \eqref{5.5:eq-move-list-e} the clause
$(\mathrm{m}\perp)^\circ$ of (H3) applies. It gives the following for $P^{\mathsf{m}}$ in the stage
class and $\pi\in\mathfrak{P}_\perp$: the dilated configuration $U^{\mathsf{m}}_s$, restricted to
$\Delta^+$, is defined and entire in $s$. Moreover, for $|s|\le R^\perp_v$ it obeys the plaquette bound
contained in that clause, together with
$|A(c,U^{\mathsf{m}}_s)|\le\mathfrak{n}^\perp_v$.

For $\mathrm{M3}$, $\mathrm{M6}$ and $\mathrm{M7}$ into a $\circ$-slot, the clauses (m3)$^\circ$ of
(H3) and the $\circ$-parts of (m6) and (m7) of (H2) apply. They state that the image is defined and
holomorphic and that it obeys the entry bound of the move, namely a scaffold plaquette bound on the
support of the piece. By Definition~\ref{5.5:def-slots-letters}, $\mathfrak{n}_{\hat A}$ is the
supremum, over the moves landing at $\circ$-slots, of the bounds on $\hat A$ furnished by the bound
for the pointed Wilson family at these entry bounds. Hence $\hat A$ is bounded there by
$\mathfrak{n}_{\hat A}$.

\emph{The move $\mathrm{M5}$\textup{(a)}.} Let $B^{\mathsf{m}}_i$ be the splitting data of the read
and $|B_i|_*$ the corresponding weighted scaffold majorants of
\eqref{5.5:eq-response-box-majorants-a} and \eqref{5.5:eq-move-list-c}. By the pointwise bounds of
(m2)$^\flat$ and the pointwise bounds that (m2)$^\circ$ supplies at read cores for the splittings
(S1)--(S3), and by (m4) for (S4), one has, in each machine,
\begin{equation}\label{5.5:eq-chain-containment-1}
|B^{\mathsf{m}}_i|\le|B_i|_*\quad\text{pointwise on }\square_0' .
\end{equation}
For the part $B_F$ of (S2), the majorant is $|B_F|_*$, and the bondwise bound
$|B^{\mathsf{m}}_F(b)|\le|B_F|_*(b)$ of (m2)$^\flat$ applies. Hence, for
$(\sigma_i)\in\mathfrak{P}_{\mathrm{M5(a)}}$,
\begin{equation}\label{5.5:eq-chain-containment-2}
\sup_{\square_0'}\Bigl|\sum_i\sigma_iB^{\mathsf{m}}_i\Bigr|
\le\sum_i|\sigma_i|\sup_{\square_0'}|B_i|_*<\alpha_3(\mathfrak{L}).
\end{equation}
The first inequality is the triangle inequality together with \eqref{5.5:eq-chain-containment-1}.
The second is the condition defining $\mathfrak{P}_{\mathrm{M5(a)}}$. Clause (m5a) of (H3), applied
with $B'=\sum_i\sigma_iB^{\mathsf{m}}_i$, then gives two conclusions. At the letter $\flat$ it places
\begin{equation*}
(U_j,J_j)^{\mathsf{m}}\bigl(\square_0',e^{i\sum_i\sigma_iB^{\mathsf{m}}_i}\bigr)\restriction X
\in\mathfrak{U}^{\mathsf{m}}(j,X;\mathfrak{L}) .
\end{equation*}
At the letter $\circ$ it defines this configuration with $\hat A$ bounded.

This step fails if a weighted point is compared with an unweighted majorant in place of the
weighted one. This is why the corresponding closure statement of clause~(i),
\eqref{5.5:eq-family-closure-a}, is formulated at the weighted majorants; there it holds by the
argument just given.

For (S4) the domain is
\begin{equation*}
|\sigma_0|\,|B^0|_*+|\sigma_1|\,|\Delta|_*<\alpha_{3,j},
\end{equation*}
and it contains the parameters $\sigma_0=1$, $|\sigma_1|\le e^{\frac18\delta D_v}$. Indeed, by
Notation~\ref{5.5:not-response-box-majorants} and the definition of $S^*_v$,
\begin{equation*}
|B^0|_*=1.42\,L^{-1}\sup_{b\in\mathfrak{r}_v}|B^{\mathrm{ax}}|_*(b)=\frac{1.42\,\alpha_{3,j}}{L\,S^*_v},
\qquad |\Delta|_*=\tfrac12\,\alpha_{3,j}\,e^{-\frac18\delta D_v},
\end{equation*}
with one and the same geometric distance $D_v$. The splitting (S4) arises only at a box-read unit
$v$, the units at which clause (m4) of (H3) is stated, so $S^*_v>1$ by the filing of
Notation~\ref{5.5:not-response-box-majorants}; and $L\ge8$ is the floor on the blocking factor under
which the scaffold is set up. Hence, at these parameters,
\begin{equation*}
|\sigma_0|\,|B^0|_*+|\sigma_1|\,|\Delta|_*
\le\alpha_{3,j}\Bigl(\frac{1.42}{L\,S^*_v}+\frac12\Bigr)
\le\alpha_{3,j}\Bigl(\frac{1.42}{8}+\frac12\Bigr)<\alpha_{3,j},
\end{equation*}
since $1.42/8+1/2<1$. This is the inequality that enters part (a) of the theorem on the second box
pair.

\emph{The move $\mathrm{M5}$\textup{(b)}.} Clause (m5b) of (H3) gives
$\|\mathsf{h}^{\mathsf{m}}_b\|\le\|\mathsf{h}_b\|_*$. Moreover, where a direction is a contraction
$\mathsf{h}[B]$ with a bond function $B$, the pointwise bound
$|B^{\mathsf{m}}(b)|\le|B|_*(b)$ of \eqref{5.5:eq-chain-containment-1} holds, so that, term by term,
\begin{equation*}
\sum_b\|\mathsf{h}^{\mathsf{m}}_b\|\,|B^{\mathsf{m}}(b)|\le\sum_b\|\mathsf{h}_b\|_*\,|B|_*(b)
=\|\mathsf{h}[B]\|_*,
\end{equation*}
the majorant of Notation~\ref{5.5:not-response-box-majorants}. These give
\begin{equation*}
\Bigl\|\sum_i\sigma_i(\mathsf{h}^{\mathsf{m}}_i,\mathbf{j}^{\mathsf{m}}_i)\Bigr\|
\le\sum_i|\sigma_i|\,\|(\mathsf{h}_i,\mathbf{j}_i)\|_*<\mathfrak{r}(\mathfrak{L}),
\end{equation*}
and the analyticity at the chart radius $\mathfrak{r}(\mathfrak{L})$ contained in (m5b) yields the
landing.

In each of the $\circ$-landings above, that is, at $\mathrm{M2}^\circ$, $\mathrm{M4}^\circ$,
$\mathrm{M8}$, at $\mathrm{M3}$, $\mathrm{M6}$, $\mathrm{M7}$ into a $\circ$-slot, and at
$\mathrm{M5}$\textup{(a)} and $\mathrm{M5}$\textup{(b)} at the letter $\circ$, the bound on $\hat A$ is
one of those over which $\mathfrak{n}_{\hat A}$ is the supremum
(Definition~\ref{5.5:def-slots-letters}); hence $\hat A$ is bounded by $\mathfrak{n}_{\hat A}$ at
every $\circ$-landing.

This completes the step. Hence no chain leaves the product of the one-machine classes, and (a) and (b)
are proved.

\emph{Assertion (c): finiteness.} Let $F$ be a $\flat$-family with $\|F^{\mathsf{m}}\|_j\le\mathfrak{n}$
for $\mathsf{m}\in\{\mathsf{A},\mathsf{B}\}$, in the weighted norm
\begin{equation*}
\|F\|_j:=\sup_X e^{\kappa d_j(X)}|F(X,\cdot)|,
\end{equation*}
where $|F(X,\cdot)|$ is the supremum over the analyticity tube of $F$, and let
$p\in\mathcal{K}^{\mathrm{rd}}(j,X;\flat)$. By (a), $P^{\mathsf{m}}$ lies in
$\mathfrak{U}^{\mathsf{m}}(j,X;\flat)=\bigcap_T\mathfrak{U}^{\mathsf{m}}_T$. This intersection is
contained in the analyticity tube of $F^{\mathsf{m}}$, on which $\|F^{\mathsf{m}}\|_j$ is the
supremum. Hence, for every value of the remaining arguments $y,w,u$,
\begin{equation*}
\begin{aligned}
e^{\kappa d_j(X)}|\delta F(X;p)|
&\le e^{\kappa d_j(X)}\bigl(|F^{\mathsf{B}}(X;P^{\mathsf{B}})|+|F^{\mathsf{A}}(X;P^{\mathsf{A}})|\bigr)\\
&\le\|F^{\mathsf{B}}\|_j+\|F^{\mathsf{A}}\|_j\le 2\mathfrak{n}.
\end{aligned}
\end{equation*}
Taking the supremum over $y,w,u$, over $p$ and over $X$ in the definition
\eqref{5.5:eq-production-closure-a} of $\mathsf{S}^{\mathrm{rd}}_j[F]$ gives
$\mathsf{S}^{\mathrm{rd}}_j[F]\le2\mathfrak{n}<\infty$. At a $\circ$-slot, (a) bounds $\hat A$ by
$\mathfrak{n}_{\hat A}$ in each machine at every point of $\mathcal{K}^{\mathrm{rd}}(j,X;\circ)$, so
$|\delta\hat A|\le2\mathfrak{n}_{\hat A}$, and the definition of $\mathsf{a}$ gives
$\mathsf{a}\le2\mathfrak{n}_{\hat A}$.

\emph{Assertion (d).} Let $\mu\colon\mathfrak{s}'\to\mathfrak{s}''$ be a move of
$\mathfrak{M}^{\mathrm{rd}}$, let $p=(P^{\mathsf{B}},P^{\mathsf{A}})\in\mathcal{K}^{\mathrm{rd}}(\mathfrak{s}')$,
and let $f$ be a family at $\mathfrak{s}''$. Put
\begin{equation*}
\delta(\pi):=f^{\mathsf{B}}\bigl(\mu^{\mathsf{B}}(\pi;P^{\mathsf{B}})\bigr)
-f^{\mathsf{A}}\bigl(\mu^{\mathsf{A}}(\pi;P^{\mathsf{A}})\bigr),\qquad\pi\in\mathfrak{P}_\mu .
\end{equation*}
We show that $\delta$ is holomorphic on $\mathfrak{P}_\mu$ (at $\mathrm{M5}$: in $(\sigma_i)$ at fixed
remaining parameters $\pi_\rho$; at $\mathrm{M7}$: in the sense of the holomorphic maps of (m7); at
$\mathrm{M8}$: in $(b,\sigma,\mathsf{t})$, and entire in $s$), that
\begin{equation*}
|\delta|\le\sup_{\mathcal{K}^{\mathrm{rd}}(\mathfrak{s}'')}\bigl|f^{\mathsf{B}}-f^{\mathsf{A}}\bigr|
\quad\text{on }\mathfrak{P}_\mu ,
\end{equation*}
and that every Cauchy or Schwarz inequality, and every identity on derivatives that holds in each
machine, passes to $\delta$. Write
$g^{\mathsf{m}}(\pi):=f^{\mathsf{m}}(\mu^{\mathsf{m}}(\pi;P^{\mathsf{m}}))$, so that
$\delta=g^{\mathsf{B}}-g^{\mathsf{A}}$. The proof consists of three points: composition, the closure
property (k1), and the linearity of the inequalities and identities used.

First, composition. By the step above, $\mu^{\mathsf{m}}(\pi;P^{\mathsf{m}})$ lies in the class of
$\mathfrak{s}''$, where $f^{\mathsf{m}}$ is holomorphic; at a $\circ$-slot the families are entire.
By clause~(ii), \eqref{5.5:eq-family-closure-b}, $\mu^{\mathsf{m}}$ is holomorphic in $\pi$. Hence
$g^{\mathsf{m}}$, and therefore $\delta$, is holomorphic on $\mathfrak{P}_\mu$.

At $\mathrm{M5}$ the holomorphy is in the chart parameters $(\sigma_i)$ at fixed $\pi_\rho$. It is
used through the Schwarz lemma in $\sigma$, in the form in which the Schwarz lemma is applied to the
chart move in this chapter. The dependence on the box amplitude $\mathsf{t}$ enters through the
moves $\mathrm{M2}$ and $\mathrm{M4}$, and the Cauchy inequalities in $\mathsf{t}$ are applied there.
At $\mathrm{M8}$ the function is entire in $s$ by $(\mathrm{m}\perp)^\circ$. It is used through the
Schwarz lemma in $s$ for a function $g$ with $g(0)=0$, in the form in which the Schwarz lemma is
applied to the dilation move. At $\mathrm{M7}$ the holomorphy is that of the maps of (m7).

Second, the bound. By property (k1) of Definition~\ref{5.5:def-production-closure}, which is (k1)
of clause~(i), $\mu^\times(\pi;p)\in\mathcal{K}^{\mathrm{rd}}(\mathfrak{s}'')$. Hence
\begin{equation*}
|\delta(\pi)|=\bigl|\delta f\bigl(\mu^\times(\pi;p)\bigr)\bigr|
\le\sup_{\mathcal{K}^{\mathrm{rd}}(\mathfrak{s}'')}|\delta f| .
\end{equation*}

Third, linearity. A Cauchy or Schwarz inequality holds for every holomorphic function whose modulus
is bounded by a number $\mathfrak{N}$ on the relevant domain. Therefore it applies to $\delta$ with
$\mathfrak{N}=\sup|\delta|$. An identity $\partial^ig^{\mathsf{m}}(\pi_0)=0$ at a point
$\pi_0\in\mathfrak{P}_\mu$ that holds in each machine gives $\partial^i\delta(\pi_0)=0$, by linearity
of $\partial^i$.
\end{proof}

\begin{proof}[Proof of clauses (i) and (v) of Theorem~\ref{5.5:thm-family-closure}]
\label{5.5:prf-closure-clauses}
Throughout, hypotheses (H1)--(H3) of Definition~\ref{5.5:def-membership-hypotheses} are in force.
Hypothesis (H4), displayed in \eqref{5.5:eq-membership-hypotheses-d}, enters at one place, in
part (a) of clause (v). Stage slots $\mathfrak{s}=(k;\mathsf{n},\hat X)$, unit slots
$\mathfrak{u}=(j,X;\mathfrak{L})$ and the letters $\mathfrak{L}\in\{\flat,\circ\}$ are those of
\eqref{5.5:eq-slots-letters-a}. The moves $\mathrm{M1}^{M},\mathrm{M2},\dots,\mathrm{M8}$ and their
parameter domains $\mathfrak{P}_\mu$ are those of Definition~\ref{5.5:def-move-list}.

\emph{Clause (i).} By Definition~\ref{5.5:def-production-closure}, in the form
\eqref{5.5:eq-production-closure-a}, the family $\mathcal{K}^{\mathrm{rd}}$ contains $\mathbb{1}$ in
every class, which is (k0). It contains the seeds in their classes, which is (k2). It is closed
under every move $\mu\in\mathfrak{M}^{\mathrm{rd}}$, which is (k1). By clause (iii),
\eqref{5.5:eq-family-closure-c}, every point of $\mathcal{K}^{\mathrm{rd}}(\mathfrak{s})$ lies in
the product class $\mathcal{U}^\times(\mathfrak{s})$. Likewise every point of
$\mathcal{K}^{\mathrm{rd}}(j,X;\flat)$ lies in
$\mathfrak{U}^{\mathsf{B}}(j,X;\flat)\times\mathfrak{U}^{\mathsf{A}}(j,X;\flat)$. Hence
$\mathcal{K}^{\mathrm{rd}}$ is a family of subsets of the product classes. Together with (k0), (k2)
and (k1) at every move, this makes $\mathcal{K}^{\mathrm{rd}}$ a correspondence for
$\mathfrak{M}^{\mathrm{rd}}$, in the sense of the definition of correspondences attached to the
general move list $(\mu1)$--$(\mu4)$.

The named clauses of \eqref{5.5:eq-family-closure-a} are, in the order listed there, (k1) at
$\mathrm{M6}$; at $\mathrm{M1}^{M}$; at $\mathrm{M2}$, $\mathrm{M3}$ and $\mathrm{M4}$ together;
at $\mathrm{M3}$; at $\mathrm{M5}$; at $\mathrm{M4}$ together with $\mathrm{M5}$(a)(S4); and at
$\mathrm{M8}$. We check each identification in the letters in which the clause is quoted.

\emph{The first named clause.} This is the clause for moves from a stage slot to a stage slot or
to a unit slot. It is contained in (k1) at $\mathrm{M6}$, since $\mathrm{M6}$ is the move from a
stage slot to a stage slot or to a unit slot (Definition~\ref{5.5:def-move-list},
\eqref{5.5:eq-move-list-d}).

\emph{The dilation clause.} This clause is the inclusion
\begin{equation}\label{5.5:eq-closure-clauses-1}
  \Phi^\times_{n,t}\bigl(b;\mathcal{K}_{n+1}(Y)\bigr)\subset\mathcal{K}_n(Y),
\end{equation}
for $b$ and $\mathsf{s}$ in the stage's polydisc. It is written in the letters in which the clause
is consumed: by the dictionary of clause (v)(e), $\mathcal{K}_n(Y)=\mathcal{K}^{\mathrm{rd}}(k;n,Y)$
is the stage class at the stage domain $Y$, and $b$, $\mathsf{s}$ are the parameters of the
stage's polydisc. In the form in which it is used,
\eqref{5.5:eq-closure-clauses-1} holds with one parameter $(b,t,\mathsf{s})$ in both machines, and
it is read on the polydisc
\begin{equation}\label{5.5:eq-closure-clauses-2}
  |t_\square|\le R^\natural_\square=\mathsf{w}_0\,r^\sharp_\square
  \quad\text{if }\tilde\square\subset\hat X,
  \qquad t_\square:=0\quad\text{otherwise}.
\end{equation}
We show that this is (k1) at $\mathrm{M1}^{M}$ on $\mathfrak{P}_{\mathrm{M1}}$ itself. Three
observations give this.
\begin{itemize}
  \item The one-machine map of which $\Phi^\times_{n,t}$ in \eqref{5.5:eq-closure-clauses-1} is
    the pair, written $\Phi_{t,\mathsf{s}}(b)$ where the clause is consumed, is the response
    dilation $\Phi_t$ of Corollary~\ref{5.2:cor-response-dilation}.
  \item The one parameter $(b,t,\mathsf{s})$ common to both machines is the one parameter of
    $\mathrm{M1}^{M}$.
  \item The polydisc $\{|t_\square|\le r_\square\}$ with $r_\square:=c_M/\vartheta(\square)$,
    at which the clause is first stated, is replaced by \eqref{5.5:eq-closure-clauses-2}. By
    \eqref{5.5:eq-move-list-a}, the polydisc \eqref{5.5:eq-closure-clauses-2} is
    $\mathfrak{P}_{\mathrm{M1}}$.
\end{itemize}
The Cauchy estimates of the dilation step, which use \eqref{5.5:eq-closure-clauses-1}, then carry
$K_A\vartheta^\natural(\square)=K_A\vartheta(\square)/\mathsf{w}_0$ in place of
$K_A\vartheta(\square)$.

\emph{The response clause and the own-read clause.} The response clause asks that every response
map run by the correlation theorem (Theorem~\ref{5.1:thm-correlation-theorem}) at the site (P) of
Definition~\ref{5.2:def-sites-units} on
small units satisfy two requirements. First, it runs with one parameter $(b,t,\mathsf{t},\sigma)$
in both machines. Second, it sends the stage class $\mathcal{K}_{\mathsf{n}+1}(\hat X)$ into the
class over which $\mathsf{S}$ is the supremum. The response maps in question are of three kinds:
\begin{itemize}
  \item box reads $(\mathfrak{b}'^{\mathsf{B}}_j,\mathfrak{b}'^{\mathsf{A}}_j)$ restricted to
    $X_v$, on $[0,1]\times\mathbb{D}_\square$;
  \item the own-level pairs $\mathcal{W}^\times_T$ of far small units, with
    $|\tau|\le\mathsf{w}_0 r_T(\square)$;
  \item the stored remainders $(\mathsf{P}^1)^\times$, on
    $[0,1]\times\{|\mathsf{t}|\le r^{\mathrm{rem}}_\square\}$.
\end{itemize}
The three domains are those of $\mathrm{M2}$, $\mathrm{M3}$ and $\mathrm{M4}$, as displayed in
\eqref{5.5:eq-move-list-b} and \eqref{5.5:eq-response-box-majorants-c}, for the following reasons.
\begin{itemize}
  \item The disc $|\tau|\le\mathsf{w}_0 r_T(\square)$ of $\mathrm{M3}$ is the disc of the
    own-read radius of the own-read Cauchy lemma.
  \item The disc $\mathbb{D}_\square$ of $\mathrm{M2}$ contains the disc
    $\{|\mathsf{t}|\le\mathsf{w}_P\}$ on which box reads are run, where $\mathsf{w}_P$ is the
    box-read amplitude radius; this radius satisfies $\mathsf{w}_P\le r^\sharp_\square$ by the
    box-read radius statement.
  \item The disc of $\mathrm{M4}$ has the radius $r^{\mathrm{rem}}_\square$ of clause (d) of the
    remainder lemma.
\end{itemize}
Hence the response clause follows from (k1) at $\mathrm{M2}$, at $\mathrm{M3}$ and at $\mathrm{M4}$.
The own-read clause is (k1) at $\mathrm{M3}$, taken at near own-read pairs; this again uses the
first of the three identifications above.

\emph{The weighted closure clause.} In the form in which it is used, this clause is read at own
boxes and at weighted majorants. The own boxes are the boxes $\mathfrak{r}_v$ of $\mathrm{M5}$ in
\eqref{5.5:eq-move-list-c}. The weighted majorants are those of
\eqref{5.5:eq-response-box-majorants-a}. So the clause is (k1) at $\mathrm{M5}$.

\emph{The two closures of the remainder lemma.} The remainder lemma's clause (g) asks for (k1) at
two moves $\mathrm{M4}'$ and $\mathrm{M4}''$. Here $\mathrm{M4}'=\mathrm{M4}$. The move
$\mathrm{M4}''$ is $\mathrm{M5}$(a) under (S4), taken at the data $(B^0,\Delta_c)$ with majorants
$|B^0|_*$ and $|\Delta|_*$, as in \eqref{5.5:eq-move-list-c}. Hence the two closures follow from
(k1) at $\mathrm{M4}$ and (k1) at $\mathrm{M5}$(a)(S4).

\emph{The closure hypothesis of the orthogonal-move theorem.} That theorem has exactly one
closure hypothesis, namely (k1) at $\mathsf{M}^\perp$ on $\mathfrak{P}_\perp$. By
\eqref{5.5:eq-move-list-e}, $\mathrm{M8}:=\mathsf{M}^\perp$ on $\mathfrak{P}_\perp$, so the
hypothesis is (k1) at $\mathrm{M8}$. This proves clause (i).

\emph{Clause (v), part (a).} In the general move list, the move $(\mu1)$ at the site (P) is a
response map of (P), in the parameters $(\mathsf{t},\sigma)$, restricted to the set of the unit
$v$. The response maps at
(P) on small units, as fixed in \eqref{5.5:eq-response-box-majorants-c}, are the following:
\begin{itemize}
  \item $\Phi_t$ on the units of the class $\mathcal{T}_2$ of leaves;
  \item $\mathfrak{b}'_j$ on box-read units and on the pieces of class $W$;
  \item $\mathcal{W}_T$ on far own-read units;
  \item $\mathsf{P}^1$ on the stored remainder.
\end{itemize}
Definition~\ref{5.5:def-move-list} fixes only index sets and forms for these maps, and introduces
no map of its own. The move $\mathrm{M8}:=\mathsf{M}^\perp$ is the construction of the
orthogonal-move theorem, and it has no entry in the general move list. For the moves $(\mu2)$ and
$(\mu3)$, the a-priori majorant of the comb data of a box read is the weighted majorant of
\eqref{5.5:eq-response-box-majorants-a}; this is $\mathrm{M5}$(a). The move $(\mu3)$ from
$\mathbb{1}$ is $\mathrm{M5}$(b), and $(\mu4)$ is $\mathrm{M6}$. The seeds are the same, and so is
the typing.

Therefore $\mathcal{K}^{\mathrm{rd}}\subset\mathcal{K}^\natural$, where $\mathcal{K}^\natural$ is
the production closure of the general move list, which in the form typed here is
$\operatorname{cl}^{\mathrm{pr}}_J$.
Here $\mathrm{M8}$ enters $\operatorname{cl}^{\mathrm{pr}}_J$ through its dilation, as recorded in
Definition~\ref{5.5:def-production-closure}. Consider the statements of the general move list that
quantify over points of chains. These are statements about
$\operatorname{cl}^{\mathrm{pr}}_J\supseteq\mathcal{K}^{\mathrm{rd}}$. Their obligations at reads
of the site (P) are the entries collected in hypothesis (H4),
\eqref{5.5:eq-membership-hypotheses-d}.

\emph{Clause (v), part (b).} By the dilation clause above, (k1) at $\mathrm{M1}^{M}$ is
\eqref{5.5:eq-closure-clauses-1} on the polydisc \eqref{5.5:eq-closure-clauses-2}. By the first
named clause above, (k1) at $\mathrm{M6}$ contains the clause for moves from a stage slot to a
stage slot or to a unit slot.

The Cauchy estimates of the dilation step use the following ingredients:
\begin{itemize}
  \item the dilation clause;
  \item that stage clause;
  \item the product principle of clause (iii);
  \item the lemma on the units of $\mathcal{T}_2$;
  \item the one-machine sizes.
\end{itemize}
All of these hold on $\mathcal{K}^{\mathrm{rd}}$, by clause (i), by clause (iii),
\eqref{5.5:eq-family-closure-c}, and by clause (iv)$'$, \eqref{5.5:eq-family-closure-d}.

\emph{Clause (v), part (c).} The places listed in (c) are the enumeration of the places at which
the weighted closure clause reads $\mathcal{K}^{\mathrm{rd}}$, together with the $\circ$-class. At
the $\circ$-class, by clause (f5) of the bound on the pointed Wilson family, $\hat A$ is entire at
radius of order one, and the Schwarz bound applies where the radius of the chart exceeds one.
Through $\mathrm{M2}^\circ$, this is $\mathrm{M5}$(a) at $\alpha_3(\circ)/\sum_i|B_i|_*$. It holds
at every read core by hypothesis $(\mathrm{m}2)^\circ$ of (H3),
\eqref{5.5:eq-membership-hypotheses-c}.

The three further consumers read $\mathcal{K}^{\mathrm{rd}}$ as follows:
\begin{itemize}
  \item The own-read statement of near units reads only (k1) at $\mathrm{M3}$ and the product
    principle.
  \item Clause (g) of the remainder lemma reads (k1) at $\mathrm{M4}$, (S4) and the product
    principle.
  \item The orthogonal-move theorem reads (k1) at $\mathsf{M}^\perp$, that is (k1) at
    $\mathrm{M8}$, and, in its part (a), the product principle.
\end{itemize}
The orthogonal-move theorem is the seventh place, and it reads nothing else of
$\mathcal{K}^{\mathrm{rd}}$.

\emph{Clause (v), parts (d) and (e).} These consist of the lines quoted in their statement,
\eqref{5.5:eq-family-closure-e}.
\end{proof}

\begin{corollary}[The cited closure hypotheses hold on the production closure]
\label{5.5:cor-closure-hypotheses-met}
Assume hypotheses (H1)--(H3) of Definition~\ref{5.5:def-membership-hypotheses}, under which
Theorem~\ref{5.5:thm-family-closure} holds. Assume also hypothesis (H4) of the same definition.
Throughout the corollary and its proof, (H1)--(H4) denote the hypotheses of that definition. Take
$\mathcal{K}:=\mathcal{K}^{\mathrm{rd}}$, the production closure of
Definition~\ref{5.5:def-production-closure}, with $\mathsf{S}:=\mathsf{S}^{\mathrm{rd}}$ and
$\mathsf{a}:=\mathsf{S}^{\mathrm{rd}}[\hat A]$. Then the following hold.
\begin{enumerate}
\item[(a)] The comparison statement for the weighted scaffold majorants of
\eqref{5.5:eq-response-box-majorants-a} takes four closure hypotheses on $\mathcal{K}$:
\begin{itemize}
\item closure of the stage classes under the move $\mathrm{M6}$;
\item their closure under the move $\mathrm{M1}^{M}$ at the radius $R^\natural_\square$;
\item condition (k1) for the response maps on small units, that is, for the moves $\mathrm{M2}$,
$\mathrm{M3}$ and $\mathrm{M4}$;
\item the closure of the unit classes, with the a-priori majorants taken to be the weighted ones.
\end{itemize}
They are four of the properties displayed in \eqref{5.5:eq-family-closure-a}. All four hold for
$\mathcal{K}^{\mathrm{rd}}$. Among the items of the fourth closure, the one produced from the base
point $\mathbb{1}$ (the rim item) is available at the majorant
\begin{equation*}
4dD\Theta\mathfrak{a}_n t^2
\end{equation*}
at which that statement uses it. On the class $W$ of cell-units of hypothesis (H1), that
statement takes, besides the two-run comparison of the cell-units, the single closure hypothesis
(k1) at $\mathrm{M8}$; this hypothesis holds by (d).
\item[(b)] Condition (k1) holds at the own-read move $\mu^{\mathrm{own}}$, that is, at the move
$\mathrm{M3}$ at near own-read pairs. This is the closure hypothesis of the two statements on near
own-read pairs.
\item[(c)] In case (S4) of the move $\mathrm{M5}$ of \eqref{5.5:eq-move-list-c}, the two closure
hypotheses of the remainder statement hold, namely closure at the moves $\mathrm{M4}'$ and
$\mathrm{M4}''$ of the canonical splitting.
\item[(d)] Condition (k1) holds at $\mathrm{M8}=\mathsf{M}^\perp$ of \eqref{5.5:eq-move-list-e}.
This is the single closure hypothesis of the comparison theorem for the comb-dilated cell-units.
\item[(e)] The stage classes of $\mathcal{K}^{\mathrm{rd}}$ form a pair system: their points are
pairs $(P^{\mathsf{B}},P^{\mathsf{A}})$, one class for each stage slot, and they have the closures
of the first two items in (a). This is what the two statements built on the field pair systems of
the sites require.
\item[(f)] No bound, constant, disc, degree or rate of the statements in (a)--(e) changes, except
by the factor $\mathsf{w}_0=1/24$ of the radius $R^\natural_\square=\mathsf{w}_0 r_\square^\sharp$.
The discs of those statements are the parameter domains $\mathfrak{P}_\mu$ of the moves
$\mathrm{M2}$--$\mathrm{M5}$ and $\mathrm{M8}$ of \eqref{5.5:eq-move-list-b},
\eqref{5.5:eq-move-list-c} and \eqref{5.5:eq-move-list-e}. The radius $R^\natural_\square$
replaces $\vartheta$ by $\vartheta^\natural$, that is,
\begin{equation}\label{5.5:eq-closure-hypotheses-met-1}
K_A\vartheta(\square)\;\longmapsto\;K_A\vartheta^\natural(\square)=\frac{2}{R^\natural_\square}.
\end{equation}
The resulting factor $1/\mathsf{w}_0=24$ enters before $\gamma$ is chosen, both at the own-read
radii and in $K_A\vartheta^\natural$.
\item[(g)] The quantity $\mathsf{l}^J$ is, by the statement that defines it, the larger of two
expressions in the $\delta$-sizes over $\mathcal{K}^{\mathrm{rd}}$; it is thus a function of these
$\delta$-sizes, and it is a member of the family $\mathsf{l}_{\mathsf{n}+1}$.
Write $\operatorname{cl}^{\mathrm{pr}}_J\supseteq\mathcal{K}^{\mathrm{rd}}$
for the closure of part (v)(a) of Theorem~\ref{5.5:thm-family-closure}. Write
$\mathsf{S}^{\mathrm{cl}}$ for the $\delta$-size $\mathsf{S}$ with the supremum taken over
$\operatorname{cl}^{\mathrm{pr}}_J$. The members of $\mathsf{l}^J$ are bounded by the bounds of
hypothesis (H4) on $\operatorname{cl}^{\mathrm{pr}}_J$, through
$\mathsf{S}^{\mathrm{rd}}\le\mathsf{S}^{\mathrm{cl}}$.
\end{enumerate}
The corollary supplies none of the hypotheses (H2), (H3), (H4). In particular it supplies none of
the following:
\begin{itemize}
\item the items (m1), (m2), (m3), (m4), (m4)$^\circ$, (m$\perp$)$^\circ$, (m5a), (m5b), (m7) of
\eqref{5.5:eq-membership-hypotheses-b} and \eqref{5.5:eq-membership-hypotheses-c};
\item the item (m3)$^\circ$ of \eqref{5.5:eq-membership-hypotheses-c}, and the items
(m6)$^\circ$, (m7)$^\circ$ of Definition~\ref{5.5:def-membership-hypotheses};
\item the evaluations at the sites (T) and (P) that enter the product-closure statement.
\end{itemize}
Two statements are inputs of the comparison statement for the weighted scaffold majorants and of
the uniqueness theorem, and not items of (H3): the two-run comparison of the cell-units, and the
two-cutoff response statement. The latter consists of the two-cutoff response statement at zeroth
order together with a bound on the energy difference between the two cutoffs.
\end{corollary}

\begin{proof}
Part (i) of Theorem~\ref{5.5:thm-family-closure} states that $\mathcal{K}^{\mathrm{rd}}$ satisfies
(k0), (k2), and (k1) at every $\mu\in\mathfrak{M}^{\mathrm{rd}}$. The display
\eqref{5.5:eq-family-closure-a} records which instance of (k1) yields each closure property named
above:
\begin{itemize}
\item closure of the stage classes under $\mathrm{M6}$ follows from (k1) at $\mathrm{M6}$;
\item closure at the radius $R^\natural_\square$ follows from (k1) at $\mathrm{M1}^{M}$;
\item condition (k1) for the response maps on small units follows from (k1) at $\mathrm{M2}$,
$\mathrm{M3}$ and $\mathrm{M4}$;
\item (k1) at $\mu^{\mathrm{own}}$ follows from (k1) at $\mathrm{M3}$ at near own-read pairs;
\item the closure with the weighted majorants follows from (k1) at $\mathrm{M5}$;
\item the two closures of the remainder statement follow from (k1) at $\mathrm{M4}$ and at
$\mathrm{M5}$(a) in the case (S4);
\item the closure hypothesis of the comparison theorem for the comb-dilated cell-units is (k1) at
$\mathrm{M8}=\mathsf{M}^\perp$.
\end{itemize}
Each move named here lies in $\mathfrak{M}^{\mathrm{rd}}$, so part (i) gives every one of these
properties. This proves the closure assertions in (a)--(d).

The rim item in (a) is produced from $\mathbb{1}$. Condition (k0) puts $\mathbb{1}$ in every
class, and the item is then closed under $\mathrm{M5}$ by (k1). Part (v)(a) of the theorem takes
the a-priori majorants to be the weighted ones of \eqref{5.5:eq-response-box-majorants-a}; so the
rim item is available at the weighted majorants of \eqref{5.5:eq-response-box-majorants-a}, and in
particular at the majorant $4dD\Theta\mathfrak{a}_n t^2$ at which that statement uses it. The
clause on the class $W$ in (a) is (d), which is already proved.

Part (iii) of the theorem, displayed in \eqref{5.5:eq-family-closure-c}, gives the finiteness of
the sizes that these statements use. For every $\flat$-family $F$ with
$\|F^{\mathsf{m}}\|_j\le\mathfrak{n}$ in the weighted norm $\|\cdot\|_j$ of part (iii), one has
\begin{equation}\label{5.5:eq-closure-hypotheses-met-2}
\mathsf{S}^{\mathrm{rd}}_j[F]\le 2\mathfrak{n}<\infty;
\end{equation}
moreover
\begin{equation*}
\mathsf{a}\le 2\mathfrak{n}_{\hat A}<\infty .
\end{equation*}
Part (iii) also gives the product principle, in the form displayed in
\eqref{5.5:eq-family-closure-c}. Fix $\mu\in\mathfrak{M}^{\mathrm{rd}}$, a move from a slot
$\mathfrak{s}'$ to a slot $\mathfrak{s}''$, and a point $p=(P^{\mathsf{B}},P^{\mathsf{A}})$ of
$\mathcal{K}^{\mathrm{rd}}(\mathfrak{s}')$, the class at the slot from which $\mu$ starts. For
functions $f^{\mathsf{B}}$, $f^{\mathsf{A}}$ on the classes of the two machines at the slot
$\mathfrak{s}''$, put
\begin{equation*}
\delta(\pi):=f^{\mathsf{B}}\bigl(\mu^{\mathsf{B}}(\pi;P^{\mathsf{B}})\bigr)
-f^{\mathsf{A}}\bigl(\mu^{\mathsf{A}}(\pi;P^{\mathsf{A}})\bigr).
\end{equation*}
Then $\delta$ is holomorphic on $\mathfrak{P}_\mu$. At $\mathrm{M5}$ this holds in the chart
parameters $(\sigma_i)$ at fixed values $\pi_\rho$ of the remaining parameters of the move; at
$\mathrm{M7}$ it holds for the holomorphic maps of
item (m7) of \eqref{5.5:eq-membership-hypotheses-b}; and at $\mathrm{M8}$ it holds in
$(b,\sigma,\mathsf{t})$, with $\delta$ entire in the parameter $s$ of $\mathsf{M}^\perp$. Moreover
$|\delta|$ is bounded by the supremum of $|f^{\mathsf{B}}-f^{\mathsf{A}}|$ over
$\mathcal{K}^{\mathrm{rd}}(\mathfrak{s}'')$. Hence, for a quantity $\psi$ at one cutoff, every
Cauchy or Schwarz inequality and every identity $\partial^i\psi(\pi_0)=0$ at a point $\pi_0$ of
$\mathfrak{P}_\mu$ passes to $\delta$, with
the bound replaced by $\sup|\delta|$. This is the form in which the statements in (a)--(d) use the
two-cutoff difference. It includes the first clause of the comparison theorem for the comb-dilated
cell-units, which applies the principle on $\mathfrak{P}_\perp$.

Part (v)(c) of the theorem, displayed in \eqref{5.5:eq-family-closure-e}, shows that nothing else
of $\mathcal{K}$ is used. It states that the statements named in (a)--(d) concern
$\mathcal{K}=\mathcal{K}^{\mathrm{rd}}$, $\mathsf{S}=\mathsf{S}^{\mathrm{rd}}$,
$\mathsf{a}=\mathsf{S}^{\mathrm{rd}}[\hat A]$ and $\mathsf{S}^\perp_v$. It also states that they
use membership in $\mathcal{K}$ at exactly seven places:
\begin{itemize}
\item (k1) at $\mathrm{M2}$, for Cauchy integrals;
\item (k1) at $\mathrm{M3}$, for near-level values;
\item (k1) at $\mathrm{M5}$(a), for the $\sigma$-disc;
\item (k1) at $\mathrm{M5}$, for the items, the rim item coming from $\mathbb{1}$;
\item (k1) at $\mathrm{M4}$ with (S4), for the remainders;
\item the $\circ$-class, where $\hat A$ is bounded by $\mathfrak{n}_{\hat A}$ by part (iii);
\item (k1) at $\mathrm{M8}$.
\end{itemize}
The far pieces use $\mathrm{M1}^{M}$ and $\mathrm{M6}$. Every one of these places is covered by
the closure properties of part (i) listed above and by part (iii).

Assertion (e) is parts (v)(b) and (v)(d) of the theorem. Part (v)(b) makes the stage classes a
pair system, with closure under $\mathrm{M6}$ and closure under $\mathrm{M1}^{M}$ at the
radius $R^\natural_\square$. Part (v)(d) states that these stage classes qualify as field pair
systems.

Assertion (g) follows from part (v)(a) of the theorem, which gives
$\mathcal{K}^{\mathrm{rd}}\subset\operatorname{cl}^{\mathrm{pr}}_J$. A supremum over the smaller
class is at most the supremum over the larger, so each $\delta$-size over
$\mathcal{K}^{\mathrm{rd}}$ is at most the corresponding one over
$\operatorname{cl}^{\mathrm{pr}}_J$:
$\mathsf{S}^{\mathrm{rd}}\le\mathsf{S}^{\mathrm{cl}}$. The bounds of hypothesis (H4) bound the
$\delta$-sizes over $\operatorname{cl}^{\mathrm{pr}}_J$. The members of $\mathsf{l}^J$ are
expressions in the $\delta$-sizes over $\mathcal{K}^{\mathrm{rd}}$; through the inequality
$\mathsf{S}^{\mathrm{rd}}\le\mathsf{S}^{\mathrm{cl}}$ they are bounded by the same bounds.

Assertion (f) follows from part (ii) of the theorem, displayed in
\eqref{5.5:eq-family-closure-b}. That part gives each move one parameter domain $\mathfrak{P}_\mu$,
with the radii of hypothesis (H1), and the same construction at both cutoffs. The discs on which
the statements in (a)--(e) work are the domains $\mathfrak{P}_\mu$ of the moves $\mathrm{M2}$,
$\mathrm{M3}$, $\mathrm{M4}$, $\mathrm{M5}$ and $\mathrm{M8}$. Hence no value, disc or constant
changes, except through the radius $R^\natural_\square=\mathsf{w}_0 r^\sharp_\square$. This radius
appears in \eqref{5.5:eq-family-closure-a} and gives \eqref{5.5:eq-closure-hypotheses-met-1}.
\end{proof}

\begin{definition}[Standing notation at a large-field site]\label{5.5:not-site-notation}
The following notation is fixed for the rest of this section.

\emph{Site, stage and indices.} We work at the site (P) of the large-field operation
$\mathbf{R}_k$ (Definition~\ref{5.2:def-sites-units}). The stage is the pair
$(\mathsf{n},\mathsf{j})$; its number is $l:=\mathsf{n}+1$ and its level is
\begin{equation*}
\mathsf{j}:=h+\mathsf{n}=\ell_l .
\end{equation*}
Here $h$ is the level from which the stages of $\mathbf{R}_k$ at the site (P) are counted, and
$\ell_l$ is the level of stage $l$; in this section the letter $h$ does not denote a history.
The units called near below carry their own indices $(j,n,\nu,t)$, which are distinct from the
stage indices.

\begin{itemize}
\item \emph{Scaffold.} $\mathfrak{B}:=\mathbb{B}^{(\mathsf{n})}_k$ is the scaffold at stage
$\mathsf{n}$. It is frozen on live pieces, that is, on pieces still carrying their large-field
factor, by the zeroth clause of the arrest statement through which the source below is read.
For a point $y$, $m(y)$ is the level of $\mathfrak{B}$ at $y$. We write
$d^{\wedge}:=d^{\wedge}_{\mathfrak{B}}$ for the scaffold distance of $\mathfrak{B}$, with the
four properties of the scaffold distance stated in the stage estimate at site (P).
\item \emph{Source.} $(\mathbb{U},\mathbb{J})$ is the source field and current of stage $l$,
read through the arrest statement for the sourced stages over a region $\hat{X}$ containing the
block $\square_0$, with the properties stated for it in the sourced induction. It is a
field on the whole domain of $\mathfrak{B}$, and it lies in
\begin{equation*}
\mathfrak{K}_P:=\mathfrak{K}_{109}(\mathfrak{B};\gamma_0),
\end{equation*}
the source configuration space at the site, built on the spaces of \cite[p.~262]{B-CMP109}.
\item \emph{Complex data.} $b$ is a complex function on the set $\mathcal{C}_{\mathsf{j}}$ with
$\|b\|_\infty\le 1.1\,C_1\delta'_{\mathsf{j}}$, where $\delta'_{\mathsf{j}}$ is the parameter
of level $\mathsf{j}$ fixed in the stage estimate at site (P). The parameter $\sigma$ satisfies
$|\sigma(\Delta)|\le e^{\kappa_1}$ for every cell $\Delta$.
\item \emph{Response.} $(A',X,\mathbb{H})$ is the triple given by the response theorem at the
data $(\mathfrak{B};\mathbb{U},\mathbb{J};\sigma;(0,b))$.
\item \emph{Source profile.} For a point $y$,
\begin{equation*}
p(y):=\sup|b|\;e^{-\delta\, d^{\wedge}(y,\mathcal{C}_{\mathsf{j}})} .
\end{equation*}
Two further statements of the stage estimate at site (P) are used below by name.
\item \emph{Shell room fractions.} With $\beta:=\tfrac15$,
\begin{equation*}
\rho_m:=\beta\,2^{-|m-\mathsf{j}-1|}.
\end{equation*}
\item \emph{Blocks and the cover.} $\{\square\}$ is the cover of the site. For each block,
$\zeta_{\square}$ is its function in the partition of unity, and
$\square^{\zeta}:=\operatorname{supp}\zeta_{\square}\subset\square$. The block weight is
\begin{equation*}
\vartheta(\square):=\vartheta\, e^{-\frac12\delta d^{-}_{\square}} ,
\end{equation*}
where $\vartheta$ is the constant fixed in the stage estimate at site (P) and
$d^{-}_{\square}$ is the distance attached to the block $\square$ in that estimate.
\item \emph{Window fraction and room.} We put
\begin{equation*}
\mathsf{w}_0:=\tfrac1{24},\qquad \vartheta^{\natural}:=\vartheta/\mathsf{w}_0,\qquad
\iota^1_l(y):=\tfrac23\,\mathsf{w}_0\,\rho_{m(y)} .
\end{equation*}
The window condition of the threshold statement for the stage is
\begin{equation}\label{5.5:eq-site-notation-1}
K_{\min}\vartheta\le\mathsf{w}_0,\qquad K_{\min}\ge\max\{K_A,\,28K^{+}\}.
\end{equation}
\item \emph{Constants.} $K_A$, $K_J^{\sharp}$, $K_S^{b}$, $K_{\min}$ and $B_{\sharp}$ are the
constants fixed in the stage estimate at site (P) and in the threshold statement for the stage.
We put
\begin{equation*}
K^{+}:=16B_{\sharp}^2\max\{K_J^{\sharp},K_S^{b}\},\qquad K'':=28K^{+}.
\end{equation*}
The box disc radius $r_{\square}^{\sharp}$ of a block $\square$ is given by
\begin{equation*}
2/r_{\square}^{\sharp}=K_A\,\vartheta(\square).
\end{equation*}
\end{itemize}

\emph{Letters kept apart.}
\begin{itemize}
\item $S_v$ denotes only the read set of the stage estimate at site (P).
\item The chart radius of a unit $v$ is $S^{*}_v$. The smaller quantity of the weight lemma,
denoted $S^{\star}_v$, occurs in this section only as the left-hand member of an inequality.
\item $\mathsf{d}_v$ is the side of the unit's own box $\mathfrak{r}_v$. For a single unit with
domain $X_v$ it equals the side $D_{X_v}$ of the box in the transfer of the analysis of
\cite[\S3]{B-CMP109} to the cover, which is stated in \cite[p.~276]{B-CMP119} without proof;
for a core unit with enlargement index $\nu$ it equals $M(2\nu+3)$. It is a side length, not a
distance.
\item The word ``tail'', used alone, means a near constituent of a term of the class
$\mathcal{T}_1$ that lies outside the core, that is, a $\mathcal{T}_1$-piece. It is distinct
from the tail of the class $\mathcal{T}_2$ in the stage estimate at site (P).
\end{itemize}
\end{definition}

\begin{definition}[First-class terms, near units and their pieces]\label{5.5:def-near-units}
We work at the site (P) of the operation $\mathbf{R}_k$, at the stage and level fixed in
Notation~\ref{5.5:not-site-notation}, whose fields $\mathbb{U}$, $\mathbb{J}$, $\bar{U}$ and $A'$
are used below. The blocks $\square$ are those of the cover of
\cite[p.~275]{B-CMP119}. For a block $\square$, the sets $\tilde{\square}$ and $\tilde{\square}^2$
are the enlargements of $\square$ defined in \cite[p.~276]{B-CMP119}, taken at the scale of the
block.

\begin{enumerate}
\item[(a)] \emph{Terms.} Let $T\in\mathcal{T}_1$ be a term of level $\ell$. Here $\mathcal{T}_1$ is
the first class of effective-action and remainder terms of \cite{B-CMP119}, taken together with
their counterterm pieces and $W$-pieces, as fixed in the statement that defines the classes of
terms at the site (P). For every such $T$ the localisation domain $X_T$ lies in
$\tilde{\square}'^2$ for some block $\square'$. Let $\Lambda_\ell\subset\Omega_\ell$ be the regions
of Definition~\ref{5.1:def-expansion-rule}. By \cite[p.~259]{B-CMP119}, $X_T\subset\Lambda_\ell$,
and $\Lambda_\ell\subset\Omega_\ell$ by condition $(\nu3)$ of the statement that defines the classes
of terms. Thus
\begin{equation*}
X_T\subset\Lambda_\ell\subset\Omega_\ell .
\end{equation*}
The \emph{domain of $T$} is the analyticity domain $U^c_\ell(X_T,\cdot)$ of \cite{B-CMP109}
attached to $T$:
\begin{equation}\label{5.5:eq-near-units-1}
\mathfrak{U}_T:=U^c_\ell(X_T,\cdot)
=\bigl\{\,Q_a(y)<\theta_{T,a}(y)\ \text{for all}\ (y,a)\in\operatorname{Cl}(T)\,\bigr\}.
\end{equation}
In \eqref{5.5:eq-near-units-1}, $\operatorname{Cl}(T)$ is the clause set of $T$, the set of pairs
$(y,a)$ of a point $y$ and an entry index $a$ fixed with $T$ in the statement that defines the
classes of terms. The entries $a=1,\dots,7$ are as follows:
\begin{align*}
(Q_1,\dots,Q_4)&=\bigl(|\partial\mathbb{U}-1|,\ |\partial U_{p',X'}(M\bullet\mathbb{U})-1|,\
|\mathbb{J}|,\ |J_{p',X'}(M\bullet\mathbb{U})|\bigr),\\
(Q_5,Q_6,Q_7)&=\bigl(|\partial\bar{U}-1|,\ |A'|,\ |\nabla_{\bar{U}}A'|\bigr).
\end{align*}
The entry $a=8$ is the cube condition of the domain. The entries $a=2$ and $a=4$ are read only at
the triples $(p',X',M)$ that belong to $T$ itself in the sum over triples representing $T$ in the
statement that defines the classes of terms. We write $\mathbf{r}^T_a(y)$ for the radius of $T$
at $y$ in entry $a$. By the threshold statement at the site (P), the thresholds satisfy
\begin{equation*}
\theta_{T,a}(y)\le 2\,\mathbf{r}^T_a(y).
\end{equation*}
The set $\mathfrak{U}_T$ is open and is a union of orbits of the complexified gauge group
$G^{\mathbb{C}}$ \cite[p.~262]{B-CMP109}.

\item[(b)] \emph{Near pairs.} A pair $(T,\square)$ is \emph{near} if
$X_T\subset\tilde{\square}^2$.

\item[(c)] \emph{Near units.} Consider the table of units that accompanies the transfer of the
analysis of \cite[Sect.~3]{B-CMP109} to the cover $\{\square\}$ and the partition of unity
$\{\zeta_{\square}\}$, printed without proof in \cite[p.~276]{B-CMP119}; cited as printed. The
units of the second and third rows of that table are singles, and a unit of the first row has a
core $\mathfrak{i}^{(j)}_y$, in the sense of that table. A unit $v$ of the
first three rows or of the fifth row of that table is \emph{near} to $\square$ if its domain
satisfies $X_v\subset\tilde{\square}^2$; such a $v$ is called a \emph{near unit}. Among its indices
a near unit carries its level $j$ and the enlargement index $\nu$ of the cube
$\square_y^{\sim\nu}$ in item~(d).

\item[(d)] \emph{Member pieces.} The pieces of a near unit $v$ are of two sorts. The first sort
is the \emph{member pieces}, which are of three types:
\begin{enumerate}
\item[(i)] for a core $\mathfrak{i}^{(j)}_y$, each of its constituents $X$ with $y\in X$ and
$X\subset\square_y^{\sim\nu}$, where $\square_y$ is the cube of $y$;
\item[(ii)] a \emph{tail}, that is, a constituent of a near unit of the first row lying outside
its core;
\item[(iii)] a single of the second or the third row.
\end{enumerate}

\item[(e)] \emph{$W$-pieces.} The second sort is the \emph{$W$-pieces}, which are of two types,
in the notation of the table of units of item~(c):
\begin{enumerate}
\item[(i)] the counterterm companions $b_j(\hat{w})\,A(h_y,\cdot)$ of the near units of the first
row;
\item[(ii)] the block pieces $f_{\square'}$ of the fifth row with $\square'\subset\tilde{\square}^2$.
\end{enumerate}
Each block piece $f_{\square'}$ is assigned, plaquette by plaquette, to the level of the cell
containing the plaquette, in the sense of the lemma on the plaquette-wise filing of block pieces
and \cite[p.~275]{B-CMP119}.

\item[(f)] \emph{Excluded objects.} Two kinds of object are excluded from both sorts of piece.
\begin{enumerate}
\item[(i)] The second member $\zeta_k(p_y)\,G_y(B')$ of a unit of the fifth row is an object of
the site (T) of Definition~\ref{5.2:def-sites-units}, by the statement classifying the members of
the units of the fifth row. It is neither a $W$-piece nor a member piece.
\item[(ii)] The pieces of clause (S4) of the table of units of item~(c) are neither member pieces
nor $W$-pieces and are members of no near unit. They are entire, and they are bounded by
clause (b) of the bound on $W$-pieces in that table.
\end{enumerate}
\end{enumerate}
\end{definition}

\begin{definition}[Per-block pieces and their two extensions]\label{5.5:def-block-piece}
We work at the site of Notation~\ref{5.5:not-site-notation}, whose symbols are used without
change: the cover $\{\square\}$ with its partition of unity $\{\zeta_{\square}\}$, the block
weight $\vartheta(\square)$ and the response $\mathbb{H}$. Let $T\in\mathcal{T}_1$ be a
term of level $\ell$ with localisation domain $X_T$, as in Definition~\ref{5.5:def-near-units}.
For every block $\square$ of the cover, by one and the same formula, the \emph{piece} of $T$ at
$\square$ is
\begin{equation}\label{5.5:eq-block-piece-a}
\operatorname{piece}(T,\square)
:= \int_0^1 dt\;
\frac{\partial}{\partial t_{\square}}\bigg|_{t_{\square}=0}
T\Bigl(X_T;\,U_\ell\bigl(\exp iQ_\ell\bigl(\eta\,(t+t_{\square}\zeta_{\square})\,
\mathbb{H}\bigr)\,\bar{U}^\ell\bigr)\Bigr).
\end{equation}
Here $Q_\ell$ is the linearised block-averaging map of Ba{\l}aban's step at level $\ell$,
applied to the field $\eta\,(t+t_{\square}\zeta_{\square})\,\mathbb{H}$, and $U_\ell(\cdot)$
is the map of that step which assigns to a level-$\ell$ configuration the field on which terms
of level $\ell$ are evaluated. The configuration $\bar{U}^\ell$ is the starting point of the
interpolation in \eqref{5.5:eq-block-piece-a}: at $t=t_{\square}=0$ the argument of $T$ is
$U_\ell(\bar{U}^\ell)$; and $\eta$ is the scale factor of the response $\mathbb{H}$ along it,
the argument at $t=1$, $t_{\square}=0$ being $U_\ell(\exp iQ_\ell(\eta\,\mathbb{H})\,\bar{U}^\ell)$.
Write $\Phi^1$ and $\Phi^0$ for the two
real configurations of the summation identity for per-block pieces; by that identity, at the
real point, that is, before either of the holomorphic extensions below is taken, the pieces add
up to the full difference:
\begin{equation}\label{5.5:eq-block-piece-1}
\sum_{\square}\operatorname{piece}(T,\square) = T(\Phi^1)-T(\Phi^0).
\end{equation}

The pieces admit two holomorphic extensions.
\begin{itemize}
\item[(a)] \emph{The box extension}, that of the box-extension statement for per-block pieces.
Let $v$ be a unit at scale $j$, of the kind considered in Definition~\ref{5.5:def-near-units}
but not required to be near, with domain $X_v$. Let
$\mathfrak{b}'_j$ denote the box pair of the statement fixing the box pair and the box disc,
\begin{equation*}
\mathfrak{b}'_j := (U,J)\bigl(\mathfrak{B}^{*}_j,\,M^{\cdot}(\mathfrak{b})\bigr),
\end{equation*}
that is, the pair $(U,J)$ taken at the arguments $\mathfrak{B}^{*}_j$ and
$M^{\cdot}(\mathfrak{b})$, the symbols $\mathfrak{B}^{*}_j$, $M^{\cdot}(\cdot)$ and
$\mathfrak{b}$ having the meaning fixed in that statement; and put
$P := \mathfrak{b}'_j|_{X_v}$, its restriction to $X_v$. For $t\in[0,1]$ the extension
is in the complex variable $\mathsf{t}$ ranging over the box disc of the same statement,
\begin{equation}\label{5.5:eq-block-piece-2}
\mathbb{D}_{\square} := \{\,\mathsf{t}\in\mathbb{C} : |\mathsf{t}|\le r_{\square}^{\sharp}\,\},
\end{equation}
where $2/r_{\square}^{\sharp} = K_A\,\vartheta(\square)$ as fixed in
Notation~\ref{5.5:not-site-notation}. The box extension is taken at the datum $P$ and comes
together with the remainders of \cite[(3.21)--(3.22)]{B-CMP109}.
\item[(b)] \emph{The own extension}, that of the own-level extension statement for per-block
pieces. The term $T$ is evaluated at $\mathcal{W}_T(t,\tau)$, with $t$ the integration variable
of \eqref{5.5:eq-block-piece-a} and $\tau$ the
second variable of the own-level family; $\mathcal{W}_T$ is the restriction to $X_T$ of the
own-level family $\mathcal{W}^{\ell,\square}$, which is defined below.
\end{itemize}
The real points do not form a uniqueness set for these extensions: the values of the pieces at
the real point, and in particular the identity \eqref{5.5:eq-block-piece-1}, do not determine
which of the two extensions a given piece carries. That choice is therefore part of the datum
of the piece: it is made once and for all, and every later use of the piece refers to the
extension so chosen.
\end{definition}

\begin{definition}[The membership clause at own-level units]\label{5.5:def-membership-clause}
Let $\square$ be a block of the cover. Let $v$ be a near member unit in the sense of
Definition~\ref{5.5:def-near-units} that belongs to the own-read class, and let $T$ be the
term of $v$. Let $X_v$ be the localisation domain of $v$, $f_v$ its function and
\begin{equation*}
\mathfrak{D}_v=\mathfrak{U}_T
\end{equation*}
its domain, the analyticity domain of $T$ fixed in Definition~\ref{5.5:def-near-units}.
Let $\mathfrak{n}_v$ be the supremum of $|f_v|$ on $\mathfrak{D}_v$, and let $\rho_A(v)$ be
the disc radius of $v$; the own-read class and the radius $\rho_A(v)$ are both fixed in the
corollary of this chapter on own-read units. We write $p$ for the near piece of $v$ at
$\square$, that is, for the pair $(v,\square)$, which indexes the piece
$\operatorname{piece}(T,\square)$ of Definition~\ref{5.5:def-block-piece}.

We say that the \emph{membership clause} holds at $(v,\square)$ if there is a family
$\mathsf{R}_p(t,\mathsf{t},\sigma)$, $t\in[0,1]$, with the following five properties.
\begin{enumerate}
\item[(a)] For each $t\in[0,1]$ the family is holomorphic in
$(\mathbb{U},\mathbb{J},b,\sigma,\mathsf{t})$ on the product domain
\begin{equation*}
\mathfrak{K}_P\times\{\mathsf{t}:|\mathsf{t}|\le\rho_A(v)\}
\times\{\sigma:|\sigma(\Delta)|\le e^{\kappa_1}\ \text{for every}\ \Delta\}.
\end{equation*}
Here $(\mathbb{U},\mathbb{J},b)$ runs over the source configuration space $\mathfrak{K}_P$
at the site (P) of the operation $\mathbf{R}_k$ (Definition~\ref{5.2:def-sites-units}),
$\sigma=\{\sigma(\Delta)\}_{\Delta}$ is the family of complex parameters
$\sigma(\Delta)$, one for each index $\Delta$, and $\kappa_1$ is one of the auxiliary
parameters of Ba{\l}aban's construction.
\item[(b)] The family takes its values in $\mathfrak{D}_v=\mathfrak{U}_T$.
\item[(c)] The family depends on $\sigma$ only through the decoupled list of the parameters
$\sigma(\Delta)$ restricted to $X_v$.
\item[(d)] At the real point, at which the real formula \eqref{5.5:eq-block-piece-a} for the
piece is read, the family reproduces the piece of $v$ at $\square$ given by that formula:
\begin{equation}\label{5.5:eq-membership-clause-1}
\int_0^1 dt\,\partial_{\mathsf{t}}\big|_{\mathsf{t}=0}\,
f_v\big(\mathsf{R}_p(t,\mathsf{t},\sigma)\big)\Big|_{\mathrm{real}}
=\operatorname{piece}(T,\square),
\end{equation}
where $f_v$ is the function of the unit $v$ (Definition~\ref{5.2:def-sites-units}).
\item[(e)] On the whole domain in (a) the family satisfies
\begin{equation*}
\big|f_v\circ\mathsf{R}_p\big|\le\mathfrak{n}_v .
\end{equation*}
\end{enumerate}
The membership clause, without qualification, is the statement that it holds at
$(v,\square)$ for every block $\square$ of the cover and every near member unit $v$ of the
own-read class.
\end{definition}

\begin{definition}[The own-level family and its radius]\label{5.5:def-own-level-family}
We work at a large-field site in the standing notation of
Notation~\ref{5.5:not-site-notation}: $\mathfrak{B}$ is the frozen scaffold of the stage, $m(y)$
its level at $y$, $(\mathbb{U},\mathbb{J})$ the source field and current of the stage, $\sigma$
as in that notation, $\zeta_{\square}$ the partition-of-unity function of the block
$\square$ of the cover and $\square^{\zeta}$ its support, $\vartheta(\square)$ the block weight,
$K^{+}$ the fixed constant of that notation, and $\beta=1/5$. Fix a level $\ell\le k$ and a
block $\square$. Put:
\begin{itemize}
\item $\mathfrak{B}_\ell$ is the scaffold $\mathfrak{B}$ capped at level $\ell$, and
$d_\ell := d^{\wedge}_{\mathfrak{B}_\ell}$ is its scaffold distance;
\item $\varrho(y) := L^{-(m(\bar{y})-\ell)^{+}}$;
\item $\hat{p} := (\varrho p)^{[\delta/2]}$ and
$\hat{p}_{\square} := (\varrho p\,\mathbf{1}_{\square^{\zeta}})^{[\delta/2]}$, where $p$ is the
function of Notation~\ref{5.5:not-site-notation} and $\mathbf{1}_{\square^{\zeta}}$ is the
indicator of $\square^{\zeta}$;
\item $\mathsf{B}(t,\tau) := Q_{\mathfrak{B}_\ell}\bigl(\mathbb{U};\eta(t+\tau\zeta_{\square})
\mathbb{H}\bigr)$, with $\mathbb{H}$ the response field of the site;
\item $(A'_\ell,X_\ell,\mathbf{H}_\ell)$ is the triple furnished by the response-field theorem
at the data $\bigl(\mathfrak{B}_\ell;\mathbb{U},\mathbb{J};\sigma;(0,\mathsf{B}(t,0))\bigr)$;
\item $\delta(\cdot) := \partial_\tau|_{\tau=0}$, the first derivative in $\tau$ at
$\tau=0$; this operator is distinct from the decay constant $\delta$ which appears in the
exponents above and below;
\item $\mathbf{A} := \mathbf{H}_\ell + \tau\,\delta\mathbf{H}_\ell$.
\end{itemize}
The \emph{own-level family} is
\begin{equation}\label{5.5:eq-own-level-family-1}
\begin{split}
\mathcal{W}^{\ell,\square}(t,\tau) := \Bigl(\, & e^{i\eta\mathbf{A}}\mathbb{U},\;
\mathbb{J} + \bigl(\Delta_{\mathbb{U}} - P_1^{\ell}(\sigma) + \mathsf{K}\bigr)
\bigl(A'_\ell + \tau\,\delta A'_\ell\bigr) \\
&- \mathfrak{m}_H^{\ell}(\sigma)\bigl(X_\ell + \tau\,\delta X_\ell\bigr)
+ F(\mathbb{U},\mathbf{A})\Bigr).
\end{split}
\end{equation}
Here $\Delta_{\mathbb{U}}$, $P_1^{\ell}(\sigma)$, $\mathsf{K}$ and $\mathfrak{m}_H^{\ell}(\sigma)$
are the operators, and $F(\mathbb{U},\cdot)$ the map, entering the current component of the
family. These, the map $Q_{\mathfrak{B}_\ell}(\mathbb{U};\cdot)$, the operation
$(\cdot)^{[\delta/2]}$ and the point $\bar{y}$ in $\varrho(y)$ are fixed together with the
family $\mathcal{W}^{\ell,\square}$, which, with the data at which the response-field theorem
is applied above, is taken by statement from the construction of the own-level extension of
the terms. Its construction involves neither the analyticity domain nor the region of any
term; these enter only through the restriction below.

Now let $T$ be a term of the class $\mathcal{T}_1$ of Definition~\ref{5.5:def-near-units}, of
level $\ell_T$, with region $X_T$ and analyticity domain $\mathfrak{U}_T$, and write
$\ell=\ell_T$. Put $\mathcal{W}_T := \mathcal{W}^{\ell_T,\square}\restriction_{X_T}$, the
restriction of the own-level family to $X_T$. For an input entry
$E := Q_a(\mathbb{U},\mathbb{J})(y)$ in the sense of Definition~\ref{5.5:def-near-units},
with $(y,a)\in\mathrm{Cl}(T)$ and $\mathbf{r}^T_a(y)$ the radius of $T$ at entry $a$ and
point $y$, put the \emph{radius of $T$ at $\square$} and the \emph{entry coefficient}
\begin{equation}\label{5.5:eq-own-level-family-a}
\begin{gathered}
r_T(\square) := \frac{e^{\frac14\delta d_\ell(X_T,\square^{\zeta})}}{14K^{+}\vartheta(\square)},
\\
c_E := \tfrac12\Bigl(1+\frac{E}{\mathbf{r}^T_a(y)}\Bigr)\ \text{ for } a\in\{2,4\},
\qquad c_E := 1\ \text{ for } a\in\{1,3,6,7\}.
\end{gathered}
\end{equation}
Then on $\mathfrak{U}_T$ one has $c_E<3/2$, and $c_E\le 11/10$ wherever
$E\le(1+\beta)\mathbf{r}^T_a(y)$.
\end{definition}

\begin{proof}[Proof of the bounds on $c_E$]
For $a\in\{1,3,6,7\}$ we have $c_E=1$, which satisfies both bounds. Let $a\in\{2,4\}$. By
Definition~\ref{5.5:def-near-units}, $\mathfrak{U}_T$ is the set on which
$Q_a(y)<\theta_{T,a}(y)$ for the pairs $(y,a)$ of its defining clauses, and the threshold
satisfies $\theta_{T,a}(y)\le 2\mathbf{r}^T_a(y)$. Hence on $\mathfrak{U}_T$ we have
$E<2\mathbf{r}^T_a(y)$, and therefore
\begin{equation*}
c_E = \tfrac12\Bigl(1+\frac{E}{\mathbf{r}^T_a(y)}\Bigr) < \tfrac12(1+2) = \tfrac32 .
\end{equation*}
Where $E\le(1+\beta)\mathbf{r}^T_a(y)$, the same formula gives, with $\beta=1/5$,
\begin{equation*}
c_E \le \tfrac12\bigl(1+1+\beta\bigr) = 1+\frac{\beta}{2} = \frac{11}{10}. \qedhere
\end{equation*}
\end{proof}

\begin{lemma}[The margin condition on the input at a term]\label{5.5:lem-input-margin}
Work at the site (P) of the operation $\mathbf{R}_k$, at stage $l$ with stage level $\mathsf{j}$, with the frozen scaffold $\mathfrak{B}$ and its level $m(y)$ at a point $y$, and with the source field and current $(\mathbb{U},\mathbb{J})$ of that stage; fix a block $\square$ of the cover. Recall
\begin{equation*}
\rho_m=\beta\,2^{-|m-\mathsf{j}-1|},\qquad
\iota^1_l(y)=\tfrac{2}{3}\,\mathsf{w}_0\,\rho_{m(y)},\qquad
\beta=\tfrac15,\quad \mathsf{w}_0=\tfrac{1}{24}.
\end{equation*}
Let $T\in\mathcal{T}_1$ be a term of level $\ell$. Its domain $X_T$, its analyticity domain $\mathfrak{U}_T$, its set of clauses $\operatorname{Cl}(T)$, the entries $Q_a$, the thresholds $\theta_{T,a}$ and the radii $\mathbf{r}^T_a$ are those of Definition~\ref{5.5:def-near-units}. Every pair $(y,a)\in\operatorname{Cl}(T)$ with $a\notin\{5,8\}$ is called a location. The stage distinguishes two types of location: the $Z^{\mathrm{on}}$-locations, at which the margin of the input is read from statement (a) below, and the $\mathfrak{T}$-locations, at which it is read from statement (b) below. Assume the following four statements. Each is assumed for every term of Ba{\l}aban's classes of effective-action and remainder terms, not only for $T$, and none of them is restricted to terms $T$ for which $(T,\square)$ fails to be near in the sense of Definition~\ref{5.5:def-near-units}; we apply them to the fixed term $T$.
\begin{itemize}
\item[(a)] (The input margin statement of stage $l$, read at $Z^{\mathrm{on}}$-locations.) The restriction $(\mathbb{U},\mathbb{J})|_{X_T}$ lies in $\mathfrak{U}_T$. Moreover, at every $Z^{\mathrm{on}}$-location $(y,a)$ the unused margin in the entry $a$ on the shell of level $m(y)$ satisfies
\begin{equation*}
\theta_{T,a}(y)-Q_a(y)\ge\rho_{m(y)}\,\mathbf{r}^T_a(y).
\end{equation*}
On live pieces, that is, pieces still carrying their large-field factor at the level read, this hypothesis is supplied by the lemma on the arrested source at live pieces.
\item[(b)] (The interpolated threshold statement at stage $l$.) At every $\mathfrak{T}$-location $(y,a)$ we have $\theta_{T,a}(y)\ge\theta_{\mathrm{src},a}(y)$, where $\theta_{\mathrm{src},a}(y)$ is the threshold of entry $a$ at $y$ carried by the source $(\mathbb{U},\mathbb{J})$ of the stage, and, with
\begin{equation*}
\theta^{(l)}_{T,a}(y):=\theta_{T,a}(y)-\tfrac13\bigl(\theta_{T,a}(y)-\theta_{\mathrm{src},a}(y)\bigr)
-\sum_{l'\le l}\iota^{1,a}_{l'}\,\mathbf{r}^T_a(y),
\end{equation*}
where $\iota^{1,a}_{l'}$ is the room of entry $a$ spent at stage $l'$, the chain
\begin{equation*}
Q_a(y)<\theta^{(l)}_{T,a}(y)\le\theta_{T,a}(y)-\iota^1_l(y)\,\mathbf{r}^T_a(y).
\end{equation*}
\item[(b$'$)] (The threshold--radius ratio bound.) At every $\mathfrak{T}$-location $(y,a)$ with $a\in\{2,4\}$ the ratio of the threshold to the radius is at most $1+\beta$, that is, $\theta_{T,a}(y)\le(1+\beta)\,\mathbf{r}^T_a(y)$.
\item[(c)] (The partition of locations.) Every location is a $Z^{\mathrm{on}}$-location or a $\mathfrak{T}$-location.
\end{itemize}
Then $(\mathbb{U},\mathbb{J})|_{X_T}\in\mathfrak{U}_T$. Moreover, at every location $(y,a)$ at least one of the following two alternatives holds:
\begin{equation}\label{5.5:eq-input-margin-a}
\begin{aligned}
&[Z]\quad Q_a(y)\le\theta_{T,a}(y)-\rho_{m(y)}\,\mathbf{r}^T_a(y);\\
&[\mathfrak{T}]\quad Q_a(y)\le\theta_{T,a}(y)-\iota^1_l(y)\,\mathbf{r}^T_a(y),\ \text{and}\\
&\qquad\theta_{T,a}(y)\le(1+\beta)\,\mathbf{r}^T_a(y)\ \text{if}\ a\in\{2,4\}.
\end{aligned}
\end{equation}
\end{lemma}

\begin{proof}
Since the hypotheses (a), (b), (b$'$) and (c) are stated for every term, without any restriction to terms $T$ for which $(T,\square)$ fails to be near, they apply to the given $T\in\mathcal{T}_1$. This holds whether or not $(T,\square)$ is near.

The membership $(\mathbb{U},\mathbb{J})|_{X_T}\in\mathfrak{U}_T$ is the first assertion of hypothesis (a), applied to $T$.

Now let $(y,a)$ be a location, that is, $(y,a)\in\operatorname{Cl}(T)$ with $a\notin\{5,8\}$. By hypothesis (c) it is a $Z^{\mathrm{on}}$-location or a $\mathfrak{T}$-location. We treat the two cases in turn.

Suppose first that $(y,a)$ is a $Z^{\mathrm{on}}$-location. Hypothesis (a) gives $\theta_{T,a}(y)-Q_a(y)\ge\rho_{m(y)}\mathbf{r}^T_a(y)$. Rearranged, this is alternative $[Z]$ of \eqref{5.5:eq-input-margin-a}.

Suppose next that $(y,a)$ is a $\mathfrak{T}$-location. Hypothesis (b) gives
\begin{equation*}
Q_a(y)<\theta^{(l)}_{T,a}(y)\le\theta_{T,a}(y)-\iota^1_l(y)\,\mathbf{r}^T_a(y).
\end{equation*}
A fortiori the non-strict inequality holds, which is the first half of alternative $[\mathfrak{T}]$.

If in addition $a\in\{2,4\}$, hypothesis (b$'$) gives
\begin{equation*}
\theta_{T,a}(y)\le(1+\beta)\,\mathbf{r}^T_a(y).
\end{equation*}
This is the second half of alternative $[\mathfrak{T}]$.

Hence at every location one of the alternatives $[Z]$, $[\mathfrak{T}]$ of \eqref{5.5:eq-input-margin-a} holds, which proves the lemma. The derivation uses no hypothesis beyond (a), (b), (b$'$) and (c). In particular, reading the term $T$ through the own-level family of Definition~\ref{5.5:def-own-level-family} enlarges none of these hypotheses.
\end{proof}

We call the condition that alternative $[Z]$ of \eqref{5.5:eq-input-margin-a} holds at every pair $(y,a)$ the strong form of the margin condition at $T$; Lemma~\ref{5.5:lem-input-margin} does not assert it.

\begin{theorem}[Own-level membership at every block]\label{5.5:thm-own-level-membership}
Work in the setting of Notation~\ref{5.5:not-site-notation}. Assume the floors on $L$, on $M$ and on
Ba{\l}aban's auxiliary parameters, and the ceilings on $\gamma$, under which the own-level family of
Definition~\ref{5.5:def-own-level-family} is constructed. Among them are
\begin{equation*}
L\ge 8,\qquad M\ge 3L,\qquad 17K^{+}\vartheta\le 1 ,
\end{equation*}
and, listed with these ceilings, a separate smallness condition of the setting of
Notation~\ref{5.5:not-site-notation}, which is assumed as well.
For the clauses \textup{(a$'$)}, \textup{(d$'$)} and \textup{(f)} assume in addition the
\emph{window condition}
\begin{equation*}
K_{\min}\vartheta\le\mathsf{w}_0,\qquad K_{\min}\ge\max\{K_A,\,28K^{+}\},\qquad \mathsf{w}_0=\tfrac{1}{24}.
\end{equation*}
Here $K^{+}=16B_\sharp^2\max\{K_J^{\sharp},K_S^{b}\}$ and $K''=28K^{+}$. Assume the following
statements.
\begin{itemize}
\item[(A)] The own-level extension lemma with its clauses (i)--(v). Among these clauses are the
coarse-data comparison inequality and the bump bound for the displacement of the datum. Assume also:
\begin{itemize}
\item[--] the response theorem for the variational datum at the determining set $\mathfrak{B}$ and at
its cap $\mathfrak{B}_\ell$;
\item[--] the datum bound and the coarse-point bound;
\item[--] the first step of the comparison statement that accompanies the own-level extension lemma;
\item[--] the telescoping identity for the per-block pieces stated in
Definition~\ref{5.5:def-block-piece};
\item[--] the fourth standing condition on which the own-level extension lemma rests.
\end{itemize}
\item[(B)] The window statements:
\begin{itemize}
\item[--] the window condition above;
\item[--] the window fraction $\mathsf{w}_0$ and the spent room $\iota^1_l(y)=\tfrac23\mathsf{w}_0\rho_{m(y)}$;
\item[--] the interpolated threshold statement $\Im_l$;
\item[--] the input condition on the set $Z^{\mathrm{on}}$, which yields alternative $[Z]$ of
\eqref{5.5:eq-input-margin-a}, and the second auxiliary condition of the window statements.
\end{itemize}
\item[(C)] Lemma~\ref{5.5:lem-input-margin} (the margin condition on the input at a term), taken
by statement.
\item[(D)] For \textup{(f)} only: the first condition on pair systems, at the map $\mu^{\mathrm{own}}$
described in \textup{(f)}.
\end{itemize}
Let $T\in\mathcal{T}_1$ be a term of level $\ell$ in the sense of Definition~\ref{5.5:def-near-units}.
Let $\square$ be any block of the cover, near $T$ or not. Let $t\in[0,1]$, and let
$|\sigma(\Delta)|\le e^{\kappa_1}$. Let $(\mathbb{U},\mathbb{J})\in\mathfrak{K}_P$.
Let $\mathcal{W}_T$, $r_T(\square)$, $\hat p$, $\hat p_\square$ and
$c_E$ be as in \eqref{5.5:eq-own-level-family-a} and Definition~\ref{5.5:def-own-level-family}, and let
$\alpha_{\cdot,\ell}$ be the radius at level $\ell$ of the setting. Here
$E=Q_a(\mathbb{U},\mathbb{J})(y)$ is the input entry.
We say that $(\mathbb{U},\mathbb{J})$ satisfies the \emph{margin condition} if it satisfies
\eqref{5.5:eq-input-margin-a}. We say that it satisfies the \emph{strong margin condition} if it
satisfies \eqref{5.5:eq-input-margin-a} with alternative $[Z]$ chosen at every
$(y,a)\in\mathrm{Cl}(T)$ with $a\notin\{5,8\}$.
Then the following hold.
\begin{itemize}
\item[(o)] \emph{Increments.} Let $|\tau|\le r_T(\square)$. At $y\in X_T$, the entries $5$ and $8$ of
$\mathcal{W}_T(t,\tau)$ coincide
with those of the input $(\mathbb{U},\mathbb{J})$. For $a\notin\{5,8\}$,
\begin{equation}\label{5.5:eq-own-level-membership-a}
\begin{split}
Q_a(\mathcal{W}_T(t,\tau))(y)
&\le Q_a(\mathbb{U},\mathbb{J})(y)
 +c_E K^{+}\bigl[\hat p(y)+|\tau|\,\hat p_\square(y)\bigr]\alpha_{\cdot,\ell}^{-1}\,\mathbf{r}^T_a(y)\\
&\le Q_a(\mathbb{U},\mathbb{J})(y)
 +c_E K^{+}\Bigl[5.5\,\vartheta+4.4\,|\tau|\,\vartheta(\square)\,
 e^{-\frac14\delta d_\ell(y,\square^{\zeta})}\Bigr]\rho_{m(y)}\,\mathbf{r}^T_a(y).
\end{split}
\end{equation}
Consequently, suppose $(\mathbb{U},\mathbb{J})\lceil X_T\in\mathfrak{U}_T$, so that $c_E<1.5$. In units
of $\rho_{m(y)}\mathbf{r}^T_a(y)$, the increment
$Q_a(\mathcal{W}_T(t,\tau))(y)-Q_a(\mathbb{U},\mathbb{J})(y)$ is at most
\begin{equation*}
0.638\,c_E\le 0.957\qquad\text{for }|\tau|\le r_T(\square).
\end{equation*}
Under the window condition it is at most $0.511\,c_E\mathsf{w}_0$ for
$|\tau|\le\mathsf{w}_0r_T(\square)$. This is $<\iota^1_l(y)/\rho_{m(y)}$ if $c_E\le 1.1$, and it is
$\le 0.032$ in every case.
\item[(a)] Under the strong margin condition,
$\mathcal{W}_T(t,\tau)\lceil X_T\in\mathfrak{U}_T$ for $|\tau|\le r_T(\square)$. Moreover, take the
response field $\mathbf{H}_\ell$ of Definition~\ref{5.5:def-own-level-family} at the datum
$\mathsf{B}(t,0)$. It satisfies \cite[(3.14), p.~273]{B-CMP109} with $\tfrac13\alpha_2$, where
\begin{equation*}
\alpha_2(y):=\rho_{m(y)}\,\alpha_{\cdot,\ell}\cdot 4B_\sharp^2/K^{+}.
\end{equation*}
\item[(a$'$)] Under the margin condition and the window condition,
$\mathcal{W}_T(t,\tau)\lceil X_T\in\mathfrak{U}_T$ for $|\tau|\le\mathsf{w}_0r_T(\square)$.
\item[(b)] \emph{Holomorphy.} The map $(\mathbb{U},\mathbb{J},b,\sigma,\tau)\mapsto\mathcal{W}_T(t,\tau)$
is holomorphic on the product of the domain of $(\mathbb{U},\mathbb{J},b,\sigma)$ fixed in
Notation~\ref{5.5:not-site-notation} with the disc of \textup{(a)}, respectively of
\textup{(a$'$)}, and continuous in $t$. The variable $\tau$ enters only through $\mathbf{A}$, through
$A'_\ell+\tau\,\delta A'_\ell$ and through $X_\ell+\tau\,\delta X_\ell$.
\item[(c)] \emph{Real jet.} Take the real configuration of the third clause of the recursion statement
at the site (P), at which \eqref{5.5:eq-block-piece-a} is read, and $\sigma\equiv 1$. There
\begin{equation*}
\int_0^1 dt\,\partial_\tau\big|_{\tau=0}\,T\bigl(X_T;\mathcal{W}_T(t,\tau)\bigr)
=\operatorname{piece}(T,\square).
\end{equation*}
\item[(d), (d$'$)] Let $F$ be holomorphic on $\mathfrak{U}_T$ with $|F|\le\mathfrak{n}$. Then
$|F\circ\mathcal{W}_T|\le\mathfrak{n}$ on the polydisc of \textup{(a)}, respectively of
\textup{(a$'$)}. Under the hypotheses of \textup{(a)},
\begin{equation*}
\Bigl|\int_0^1dt\,\partial_\tau\big|_{\tau=0}F(\mathcal{W}_T)\Bigr|
\le\mathfrak{n}\,K''\vartheta(\square)\,e^{-\frac14\delta d_\ell(X_T,\square^{\zeta})}.
\end{equation*}
Under those of \textup{(a$'$)},
\begin{equation*}
\Bigl|\int_0^1dt\,\partial_\tau\big|_{\tau=0}F(\mathcal{W}_T)\Bigr|
\le\mathfrak{n}\,K''\vartheta^{\natural}(\square)\,e^{-\frac14\delta d_\ell(X_T,\square^{\zeta})},
\qquad K''\vartheta^{\natural}(\square)=24K''\vartheta(\square).
\end{equation*}
\item[(e)] \emph{Independence.} The statements \textup{(o)}--\textup{(d$'$)} depend on their data only as
follows:
\begin{itemize}
\item[--] on $\mathfrak{B}$ only through the four properties of its scaffold distance listed in
Notation~\ref{5.5:not-site-notation} and through its cap $\mathfrak{B}_\ell$;
\item[--] on $(\mathbb{U},\mathbb{J})$ only through membership in
$\mathfrak{K}_{109}(\mathfrak{B}_\ell;\gamma_0)$ and through the strong margin condition, respectively
the margin condition;
\item[--] on $b$ only through $p$;
\item[--] on $\sigma$ only through the decoupled list of the response theorem, at $\mathfrak{B}$ for
the datum and at $\mathfrak{B}_\ell$ for the responses.
\end{itemize}
None of them uses the block $\square_0$ of the site, the fields $\mathbb{U}^w_\square$ and
$\mathbb{H}^w$, an axial gauge, \cite[Lemma~4]{B-CMP109}, \cite[(3.26)--(3.32)]{B-CMP109}, or the two
comparison conditions on the constants of the site. None
uses the weight, the diameter, the hull or the shape of $X_T$, the zone of $\square$, the age of a
large-field region, the three auxiliary parameters of the large-field bookkeeping at the site, or the
depth. The position of $X_T$ enters only through
$d_\ell(X_T,\square^{\zeta})\ge 0$.
\item[(f)] \emph{Two runs.} Let $\mathsf{A}$ and $\mathsf{B}$ be two runs of the construction on one
scaffold. Let $\mathcal{K}$ be a pair system satisfying the first condition on pair systems at the map
\begin{equation*}
\mu^{\mathrm{own}}:\ \bigl((b,\sigma,t,\tau);p\bigr)\mapsto
\bigl(\mathcal{W}^{\mathsf{B}}_T,\mathcal{W}^{\mathsf{A}}_T\bigr)(t,\tau)
\end{equation*}
on the domain
\begin{equation*}
\mathfrak{P}:=(b,\sigma)\times[0,1]\times\{|\tau|\le\mathsf{w}_0r_T(\square)\}.
\end{equation*}
Let $F^{\mathsf{A}}$ and $F^{\mathsf{B}}$ be the functions of the two runs that $\mathcal{K}$ pairs,
and let $\sup_{\mathcal{K}}|\delta F|$ be the size of their difference in $\mathcal{K}$; in this
clause $\delta F$ denotes that difference and not a $\tau$-derivative.
Then on the polydisc of \textup{(a$'$)}
\begin{equation*}
\bigl|F^{\mathsf{B}}(\mathcal{W}^{\mathsf{B}}_T)-F^{\mathsf{A}}(\mathcal{W}^{\mathsf{A}}_T)\bigr|
\le\sup_{\mathcal{K}}|\delta F|.
\end{equation*}
Moreover \textup{(d$'$)} holds for the difference with $\mathfrak{n}$ replaced by
$\sup_{\mathcal{K}}|\delta F|$.
\end{itemize}
\end{theorem}

\begin{proof}[Proof of Theorem~\ref{5.5:thm-own-level-membership}: objects, responses and coarse data]
The proof follows the clauses of the own-level extension lemma at an arbitrary pair $(T,\square)$. At
each step we check that the position of $X_T$ relative to $\square$ is not used. The assumed statements
are used as stated. What is proved beyond them is:
\begin{itemize}
\item[--] that none of the steps reads the position of $X_T$ relative to $\square$;
\item[--] the role of the entry coefficient $c_E$;
\item[--] the numerical estimates;
\item[--] clause (b).
\end{itemize}

\emph{Step 0: the objects.} Consider every object entering Definition~\ref{5.5:def-own-level-family}:
\begin{itemize}
\item[--] the cap $\mathfrak{B}_\ell$ and its distance $d_\ell$;
\item[--] the damping $\varrho$ and the weights $\hat p$, $\hat p_\square$;
\item[--] the datum $\mathsf{B}(t,\tau)$;
\item[--] the triple $(A'_\ell,X_\ell,\mathbf{H}_\ell)$ and the field $\mathbf{A}$;
\item[--] the family $\mathcal{W}^{\ell,\square}(t,\tau)$.
\end{itemize}
Each of these is a function of $(\mathfrak{B},\ell,\square;\mathbb{U},\mathbb{J},b,\sigma;t,\tau)$ alone.
The set $X_T$ enters $\mathcal{W}_T=\mathcal{W}^{\ell_T,\square}\lceil X_T$ only as the set to which
$\mathcal{W}^{\ell_T,\square}$ is restricted. The response theorem for the variational datum holds at
any determining set, and in particular at the cap $\mathfrak{B}_\ell$. Hence the triple
$(A'_\ell,X_\ell,\mathbf{H}_\ell)$ is defined at $(\mathfrak{B}_\ell;\mathbb{U},\mathbb{J};\sigma;(0,\mathsf{B}(t,0)))$
for every block $\square$, whatever its position relative to $X_T$.

\emph{Step 1: clause (i) of the own-level extension lemma.} That clause states that a clause of the
analyticity conditions read at a coarser level implies the corresponding clause at the finer level; it
rests on the fourth of the properties of the scaffold distance listed in
Notation~\ref{5.5:not-site-notation}. The input satisfies
$(\mathbb{U},\mathbb{J})\in\mathfrak{K}_P=\mathfrak{K}_{109}(\mathfrak{B};\gamma_0)$. Applied to the cap
$\mathfrak{B}_\ell$, clause (i) gives $(\mathbb{U},\mathbb{J})\in\mathfrak{K}_{109}(\mathfrak{B}_\ell;\gamma_0)$.
One also has $\hat p,\hat p_\square\in\mathfrak{M}^{\wedge}_{\delta/2}(\mathfrak{B}_\ell)$.
Neither fact involves $X_T$.

\emph{Step 2: clause (ii), the responses.} Clause (ii) of the own-level extension lemma bounds the
responses at $\mathfrak{B}_\ell$ in three moves.
\begin{itemize}
\item[--] Its first bound is obtained with \cite[Proposition~7]{B-CMP98}. The datum bound yields the
estimate $\le B_\sharp\varrho p$ stated in that clause for the datum
$\mathsf{B}(t,0)=Q_{\mathfrak{B}_\ell}(\mathbb{U};\eta t\mathbb{H})$, in the weighted norm in which
the clause states it.
\item[--] Write $\delta\mathsf{B}:=\partial_\tau|_{\tau=0}\mathsf{B}(t,\tau)$. Here and below
$\delta(\cdot)$ denotes the first $\tau$-derivative at $\tau=0$, not the decay constant $\delta$. By
\cite[(134)--(135)]{B-CMP98} and the locality of \cite[Proposition~4]{B-CMP98}, $\delta\mathsf{B}$ lives
on bonds whose blocks meet $\operatorname{supp}\zeta_\square=\square^{\zeta}$. In the weighted norm
with weight $\hat p_\square$ it satisfies $\|\delta\mathsf{B}\|_{\hat p_\square}\le 4B_\sharp$.
\item[--] Clause (b) of the response theorem is then applied at $\mathfrak{B}_\ell$. The
$\tau$-derivatives are limits of difference quotients in $\tau$. The sup-norm bounds on the difference
quotients are uniform, so they persist in the limit.
\end{itemize}
The result holds at every $y$, not only at $y\in X_T$:
\begin{equation}\label{5.5:eq-own-level-membership-1}
N^{\Delta}_y(\mathbf{H}_\ell),\ N^{\Delta}_y(A'_\ell),\ |X_\ell|,\
\ell^3|P_1^{\ell}A'_\ell|,\ \ell^3|\mathfrak{m}_H^{\ell}X_\ell|\ \le\ 4B_\sharp^2\,\hat p(y).
\end{equation}
The corresponding $\delta$-objects satisfy
\begin{equation}\label{5.5:eq-own-level-membership-2}
N^{\Delta}_y(\delta\mathbf{H}_\ell),\ N^{\Delta}_y(\delta A'_\ell),\ |\delta X_\ell|,\
\ell^3|P_1^{\ell}\delta A'_\ell|,\ \ell^3|\mathfrak{m}_H^{\ell}\delta X_\ell|\ \le\ 4B_\sharp^2\,\hat p_\square(y).
\end{equation}
These bounds use only the following: \cite[Propositions~4 and~7]{B-CMP98}, the datum bound, and
clause (b) of the response theorem. The function $\zeta_\square$ enters only through its
values on bonds, and no gradient of $\zeta_\square$ is used. The set $X_T$ does not enter.

\emph{Step 3: clause (iii), the coarse-data comparison.} Clause (iii) of the own-level extension lemma
states the coarse-data comparison inequality. It holds for $y$ in $\Omega_\ell$, with $m:=m(\bar y)$,
where $\bar y$ is the point attached to $y$ in the damping $\varrho$ of
Definition~\ref{5.5:def-own-level-family}, and for every $y'$. Its consequences are
\begin{equation*}
\hat p(y)\le 5r\,\vartheta\,\rho_m\,\alpha_{\cdot,\ell}\,
e^{-\frac14\delta d^{\wedge}(\bar y,\mathcal{C}_{\mathsf{j}})}
\end{equation*}
and
\begin{equation*}
\hat p_\square(y)\le 4r\,\vartheta(\square)\,\rho_m\,\alpha_{\cdot,\ell}\,
e^{-\frac14\delta d_\ell(y,\square^{\zeta})}.
\end{equation*}
Here $r$ denotes the constant of the coarse-data comparison inequality (not the radius $r$ of
Definition~\ref{5.1:def-weight-classes}).
The inequality is derived in that clause from the following ingredients:
\begin{itemize}
\item[--] the second clause of the shell lemma at the shell $\ell$;
\item[--] the coarse-point bound at the point $\bar y'$ attached to $y'$ in the same way;
\item[--] the first three of the properties of the scaffold distance listed in
Notation~\ref{5.5:not-site-notation};
\item[--] the inequality $L^{-s}(1+s)\le 1$, valid for every integer $s\ge 0$ since $L\ge 2$, where
$s$ here denotes the level gap $(m(\bar y)-\ell)^{+}$. It is used
in the case of a wide box, where the factor $\varrho$ compensates the gap between the levels;
\item[--] the inequality $d^{\wedge}(\square^{\zeta},\mathcal{C}_{\mathsf{j}})\ge d^{\wedge}(\square,R_n)$,
where $R_n$ is the set against which that clause measures the distance of $\square$.
\end{itemize}
The one hypothesis of clause (iii) is $y\in\Omega_\ell$. It is met at every point of $X_T$, because
$X_T\subset\Lambda_\ell\subset\Omega_\ell$ for every term of $\mathcal{T}_1$; this is the inclusion
recorded in Definition~\ref{5.5:def-near-units}.

With $r\le 1.1$ one has $5r\le 5.5$ and $4r\le 4.4$. Moreover $d^{\wedge}\ge 0$ gives
$e^{-\frac14\delta d^{\wedge}(\bar y,\mathcal{C}_{\mathsf{j}})}\le 1$. Hence, at every $y\in X_T$ with
$m=m(\bar y)$,
\begin{equation}\label{5.5:eq-own-level-membership-3}
\hat p(y)\le 5.5\,\vartheta\,\rho_m\,\alpha_{\cdot,\ell},\qquad
\hat p_\square(y)\le 4.4\,\vartheta(\square)\,\rho_m\,\alpha_{\cdot,\ell}\,
e^{-\frac14\delta d_\ell(y,\square^{\zeta})}.
\end{equation}
Neither the derivation of \eqref{5.5:eq-own-level-membership-3} nor its statement uses the position of
$X_T$ relative to $\square$. The set $X_T$ enters only through the inclusion
$X_T\subset\Omega_\ell$, which holds for every term of $\mathcal{T}_1$.

The clauses (o)--(f) are derived from
\eqref{5.5:eq-own-level-membership-1}--\eqref{5.5:eq-own-level-membership-3} in the next two parts of
the proof.
\end{proof}

\begin{proof}[Proof of clause (o) of Theorem~\ref{5.5:thm-own-level-membership}]
\label{5.5:prf-increment-bound}
We use the objects of Definition~\ref{5.5:def-own-level-family}, in particular the scaffold
$\mathfrak{B}_\ell$ capped at $\ell=\ell_T$, its distance $d_\ell$, the weights $\hat{p}$ and
$\hat{p}_{\square}$, the field $\mathbf{A}=\mathbf{H}_\ell+\tau\delta\mathbf{H}_\ell$, and the
radius $r_T(\square)$ and the entry coefficient $c_E$ of \eqref{5.5:eq-own-level-family-a}. Let
$(\mathbb{U},\mathbb{J})\in\mathfrak{K}_P$, $t\in[0,1]$, $|\tau|\le r_T(\square)$ and $y\in X_T$;
for an entry index $a\notin\{5,8\}$, write $E:=Q_a(\mathbb{U},\mathbb{J})(y)$ for the input entry,
and write $\rho_m:=\rho_{m(y)}$. We treat
the entries in turn and then verify the arithmetic behind \eqref{5.5:eq-own-level-membership-a}.

\emph{Entries $5$ and $8$.} The first component of $\mathcal{W}_T(t,\tau)$ is
$e^{i\eta\mathbf{A}}e^{i\eta A'}\bar{U}$, where $\mathbb{U}=e^{i\eta A'}\bar{U}$ is the
decomposition of the input field fixed in the setting of
Theorem~\ref{5.5:thm-own-level-membership}. The configuration $\bar{U}$ and the gauges on cubes
are those of the input, for every $\tau$, every $t$ and every block $\square$: the gauges of
\cite[p.~262, (i)]{B-CMP109} are local, one for each cube of size $O(1)LM$. The
$\mathbb{U}$-part of these two entries is therefore kept, and entries $5$ and $8$ of
$\mathcal{W}_T(t,\tau)$ at $y$ are those of the input.

\emph{Entries $1$, $6$, $7$.} Parts (a) and (c) of the first-order field-entry bounds, written in
level-$\ell$ units, bound the increment of each of these entries at $y$ by $3N_y(\mathbf{A})$.

\emph{Entry $3$.} The first-order current bound, summed over the local norms bounded in clause (ii)
of Theorem~\ref{5.5:thm-own-level-membership}, bounds the increment of entry $3$ at $y$ by
$4K_J^{\sharp}$ times the sum of these local norms, each of which is at most
$4B_{\sharp}^2(\hat{p}(y)+|\tau|\hat{p}_{\square}(y))$.

\emph{Entries $2$ and $4$.} These are built from the structures of $T$ itself. We apply the bump
displacement bound in its second frame, namely the triple formed by $X_T$, $\mathfrak{B}_\ell$ and
one of the structures of $T$, at the datum
\begin{equation}\label{5.5:eq-increment-bound-1}
\varkappa:=4B_{\sharp}^2\bigl(\hat{p}+|\tau|\hat{p}_{\square}\bigr)
\in\mathfrak{M}^{\wedge}_{\delta/2}(\mathfrak{B}_\ell);
\end{equation}
here, and in this proof only, the letter $\varkappa$ denotes this datum.
Its hypothesis on the unit set is met by taking the unit set to be $\mathfrak{B}_\ell$ itself; its
hypothesis on the configuration and on the weights holds by clause (i), which gives
$(\mathbb{U},\mathbb{J})\in\mathfrak{K}_{109}(\mathfrak{B}_\ell;\gamma_0)$
and $\hat{p},\hat{p}_{\square}\in\mathfrak{M}^{\wedge}_{\delta/2}(\mathfrak{B}_\ell)$; its smallness
hypothesis $\sup_{X_T}\varkappa\le\alpha_{\cdot,\ell}$ is verified below. For large $|\tau|$ the
bump displacement bound uses the datum only on $X_T\cup B(x)$, where $B(x)$ is the set attached by
that bound to its point $x$, and there
$\varkappa\le c_B\sup_{X_T}\varkappa$. Let $\pi$, $\pi_\times$ be its displacement coefficients
at the datum $\varkappa$, and $r_a$ the radius it carries at entry $a$, so that $r_a\le\mathbf{r}^T_a$.
The bound moves the entry by at most $\pi(E+r_a)$; when $E\le\mathbf{r}^T_a$ this is at most
$2\pi\mathbf{r}^T_a$. In general the new entry is at most $(1+\pi_\times)E+\pi r_a$, and since
$\pi_\times\le\pi$ and $r_a\le\mathbf{r}^T_a$,
\begin{equation}\label{5.5:eq-increment-bound-2}
(1+\pi_\times)E+\pi r_a\le E+\pi\bigl(E+\mathbf{r}^T_a\bigr)
=E+2\cdot\tfrac12\Bigl(1+\frac{E}{\mathbf{r}^T_a}\Bigr)\pi\,\mathbf{r}^T_a
=E+2c_E\pi\,\mathbf{r}^T_a ,
\end{equation}
with $c_E=\frac12(1+E/\mathbf{r}^T_a)$ as in \eqref{5.5:eq-own-level-family-a} for $a\in\{2,4\}$.
Moreover, at the point $y$,
\begin{equation}\label{5.5:eq-increment-bound-3}
2\pi=\frac{8B_{\sharp}^2K_S^{b}\bigl(\hat{p}(y)+|\tau|\hat{p}_{\square}(y)\bigr)}
{\tilde{\alpha}_{0,\ell}}
\le\frac{K^{+}\bigl(\hat{p}(y)+|\tau|\hat{p}_{\square}(y)\bigr)}{\alpha_{\cdot,\ell}},
\end{equation}
because $K^{+}\ge16B_{\sharp}^2K_S^{b}$ and because the comparison of radii in the threshold
statement at the locations $[\mathfrak{T}]$ and $[Z]$ gives $\alpha_{\cdot,\ell}\le
2\tilde{\alpha}_{0,\ell}$, that is $1/\tilde{\alpha}_{0,\ell}\le2/\alpha_{\cdot,\ell}$.

\emph{The first inequality of \eqref{5.5:eq-own-level-membership-a}.} Each of the increments
above is at most
\begin{equation*}
K^{+}\bigl(\hat{p}(y)+|\tau|\hat{p}_{\square}(y)\bigr)\alpha_{\cdot,\ell}^{-1}\mathbf{r}^T_a(y),
\end{equation*}
multiplied by $c_E$. At $a\in\{2,4\}$ this follows from \eqref{5.5:eq-increment-bound-2} and
\eqref{5.5:eq-increment-bound-3}. At $a\in\{1,3,6,7\}$ the coefficient is $c_E=1$, as in
\eqref{5.5:eq-own-level-family-a}, and the bound $3N_y(\mathbf{A})$ for entries $1$, $6$, $7$ and
the bound by $4K_J^{\sharp}$ times the sum of the local norms of clause (ii) for entry $3$ are at
most the displayed quantity by the entry-by-entry bound for the own-level family, used here as
stated. For entries $1$, $6$, $7$ the factor $N_y(\mathbf{A})$ is
controlled by clause (ii): at every point, $N^{\Delta}_y(\mathbf{H}_\ell)\le4B_{\sharp}^2\hat{p}(y)$
and $N^{\Delta}_y(\delta\mathbf{H}_\ell)\le4B_{\sharp}^2\hat{p}_{\square}(y)$, so that, since
$\mathbf{A}=\mathbf{H}_\ell+\tau\delta\mathbf{H}_\ell$,
\begin{equation*}
N_y(\mathbf{A})\le4B_{\sharp}^2\bigl(\hat{p}(y)+|\tau|\hat{p}_{\square}(y)\bigr).
\end{equation*}
This is the first inequality of \eqref{5.5:eq-own-level-membership-a}.

\emph{The second inequality.} Every $\mathcal{T}_1$-term satisfies $X_T\subset\Lambda_\ell\subset
\Omega_\ell$, so clause (iii) applies on $X_T$ and gives there
\begin{equation*}
\hat{p}(y)\le5.5\,\vartheta\rho_m\alpha_{\cdot,\ell},\qquad
\hat{p}_{\square}(y)\le4.4\,\vartheta(\square)\rho_m\alpha_{\cdot,\ell}
e^{-\frac14\delta d_\ell(y,\square^{\zeta})}.
\end{equation*}
Inserting these in the first bound gives the second inequality of
\eqref{5.5:eq-own-level-membership-a}. By the arithmetic below, the common bound is at most
$0.638\,\rho_m\mathbf{r}^T_a\le\frac23\rho_m\mathbf{r}^T_a$ before the factor $c_E$.

\emph{Dependence on $\square$.} The bounds used above do not depend on the position of $X_T$
relative to $\square$. The bound $N_y(\mathbf{A})\le4B_{\sharp}^2(\hat{p}+|\tau|\hat{p}_{\square})(y)$
comes from clause (ii). The first-order field-entry bounds and the first-order current bound are
pointwise first-order bounds. The second frame of the bump displacement bound is the triple formed
by $X_T$, $\mathfrak{B}_\ell$ and one structure of $T$, and that bound holds in every frame. Its
constant $K_S^{b}$ is independent of $\square$, of $\tau$ and of the volume of the region to which
the bound is applied. Its comb gauge is a relative one, and so involves neither
$\square$ nor a hull. The only dependence on $\square$ is through $\hat{p}_{\square}$ and $|\tau|$.
Hence every estimate of this clause holds at every pair $(T,\square)$, whether or not $\square$
is near $T$.

\emph{The entry coefficient.} On $\mathfrak{U}_T$ one has $E<\theta_{T,a}(y)\le2\mathbf{r}^T_a(y)$,
hence $E/\mathbf{r}^T_a<2$ and $c_E<1.5$. At the locations of type $[\mathfrak{T}]$ with
$a\in\{2,4\}$, Lemma~\ref{5.5:lem-input-margin} gives $E<\theta_{T,a}(y)\le(1+\beta)\mathbf{r}^T_a(y)$,
and there one has
$c_E\le\frac12(2+\beta)=1.1$, since $\beta=\frac15$.

\emph{The hypothesis $\sup_{X_T}\varkappa\le\alpha_{\cdot,\ell}$.} For the radii
$|\tau|\le r_T(\square)$ the arithmetic below gives, on $X_T$,
$K^{+}(\hat{p}+|\tau|\hat{p}_{\square})\le0.638\,\rho_m\alpha_{\cdot,\ell}$. Since
$K^{+}\ge16B_{\sharp}^2K_S^{b}$, we obtain
\begin{equation*}
\sup_{X_T}\varkappa\le\frac{0.638}{4}\,\frac{\rho_m\alpha_{\cdot,\ell}}{K_S^{b}}
\le\alpha_{\cdot,\ell}.
\end{equation*}
The case of large $|\tau|$, in which the datum is used on $X_T\cup B(x)$, is paid by the constant
$c_B$. That constant is contained in $K_S^{b}$ as fixed in the bump displacement bound. This case
is void when $d_\ell(X_T,\square^{\zeta})=0$.

\emph{Arithmetic.} For $y\in X_T$ we have $d_\ell(y,\square^{\zeta})\ge d_\ell(X_T,\square^{\zeta})$,
with both sides possibly zero, the case $d_\ell(X_T,\square^{\zeta})=0$ included. By the definition
of $r_T(\square)$ in
\eqref{5.5:eq-own-level-family-a}, for $|\tau|\le r_T(\square)$,
\begin{equation}\label{5.5:eq-increment-bound-4}
\begin{aligned}
K^{+}\cdot4.4\,|\tau|\vartheta(\square)e^{-\frac14\delta d_\ell(y,\square^{\zeta})}
&\le\frac{4.4}{14}\,
e^{\frac14\delta(d_\ell(X_T,\square^{\zeta})-d_\ell(y,\square^{\zeta}))}\\
&\le\frac{11}{35}=0.3143.
\end{aligned}
\end{equation}
By the hypothesis $17K^{+}\vartheta\le1$ of Theorem~\ref{5.5:thm-own-level-membership},
\begin{equation*}
K^{+}\cdot5.5\,\vartheta\le\frac{5.5}{17}=\frac{11}{34}=0.3235 .
\end{equation*}
The sum is
\begin{equation}\label{5.5:eq-increment-bound-5}
\frac{11}{35}+\frac{11}{34}=\frac{759}{1190}=0.638,\qquad 1.5\times0.638=0.957<1 .
\end{equation}
For $|\tau|\le\mathsf{w}_0r_T(\square)$ under the window condition $28K^{+}\vartheta\le\mathsf{w}_0$,
the same computation gives
\begin{equation*}
K^{+}\cdot4.4\,|\tau|\vartheta(\square)e^{-\frac14\delta d_\ell(y,\square^{\zeta})}
\le\frac{11\mathsf{w}_0}{35},\qquad
K^{+}\cdot5.5\,\vartheta\le\frac{11\mathsf{w}_0}{56}.
\end{equation*}
These add to
\begin{equation}\label{5.5:eq-increment-bound-6}
\frac{11\mathsf{w}_0}{56}+\frac{11\mathsf{w}_0}{35}=\frac{1001\,\mathsf{w}_0}{1960}
=0.5107\,\mathsf{w}_0 .
\end{equation}
Multiplying by $1.1$ gives $0.5618\,\mathsf{w}_0<\frac23\mathsf{w}_0=\iota^1_l/\rho_m$; multiplying
by $1.5$ gives $0.766\,\mathsf{w}_0=0.0319$, since $\mathsf{w}_0=\frac1{24}$.

\emph{Conclusion.} If the restriction of $(\mathbb{U},\mathbb{J})$ to $X_T$ lies in
$\mathfrak{U}_T$, then $c_E<1.5$. By
\eqref{5.5:eq-own-level-membership-a} and \eqref{5.5:eq-increment-bound-5}, the increment in
units of $\rho_m\mathbf{r}^T_a$ is at most $0.638\,c_E\le0.957$ for $|\tau|\le r_T(\square)$. By
\eqref{5.5:eq-increment-bound-6}, under the window condition it is at most
$0.511\,c_E\mathsf{w}_0$ for $|\tau|\le\mathsf{w}_0r_T(\square)$. This quantity is smaller than
$\iota^1_l/\rho_m$ where $c_E\le1.1$, and at most $0.032$ in every case. This proves clause (o).
\end{proof}

\begin{proof}[Proof of Theorem~\ref{5.5:thm-own-level-membership}, clauses (a)--(f), continued]
\label{5.5:prf-holomorphy-jet}
We keep the setting of the theorem: $T\in\mathcal{T}_1$ has level $\ell$, $\square$ is any block, $t\in[0,1]$,
$(\mathbb{U},\mathbb{J})\in\mathfrak{K}_P$, and $\mathcal{W}_T=\mathcal{W}^{\ell_T,\square}$ restricted to $X_T$ is the
own-level family of Definition~\ref{5.5:def-own-level-family}, with its radius $r_T(\square)$ fixed there. Points
$y$ lie in $X_T$, and $m:=m(y)$. For an entry $a$ we write $Q_a$, $\theta_{T,a}$, $\mathbf{r}^T_a$ for the values at
$y$, $Q_a$ without further argument being the input's value and $Q_a(\mathcal{W}_T)$ the value at
$\mathcal{W}_T(t,\tau)$; $\mathrm{Cl}(T)$ is the set of pairs $(y,a)$ at which the strict inequalities
$Q_a<\theta_{T,a}$ defining $\mathfrak{U}_T$ are imposed. Three letters $\delta$ occur below: $\delta(\cdot)$ is the
derivative $\partial_\tau|_{\tau=0}$ of Definition~\ref{5.5:def-own-level-family}, $\delta$ in an exponent
$e^{-\frac14\delta d_\ell}$ is the decay constant, and $\delta F$ in the proof of clause (f) is the difference of
the two runs' functions defined there.

\emph{Clauses (a) and (a$'$).} Consider first the entries $a\in\{5,8\}$. By clause (o), the entries $5$ and $8$ of
$\mathcal{W}_T(t,\tau)$ at points of $X_T$ are those of the input, so the clauses of $\mathrm{Cl}(T)$ at these
entries are the input's clauses.

Let now $a\notin\{5,8\}$, and let $|\tau|\le r_T(\square)$. The proof of clause (o) gives, in the form
\eqref{5.5:eq-own-level-membership-a}, an increment of $Q_a$ at most $c_E\cdot 0.638\,\rho_m\mathbf{r}^T_a$, and
there $c_E<1.5$ (from $E<\theta_{T,a}\le 2\mathbf{r}^T_a$, where $E$ is the quantity of the proof of clause (o)
through which $c_E=\tfrac12(1+E/\mathbf{r}^T_a)$ is defined). Since $1.5\cdot 0.638=0.957$,
\begin{equation}\label{5.5:eq-holomorphy-jet-1}
Q_a(\mathcal{W}_T)\le Q_a+0.957\,\rho_m\mathbf{r}^T_a ,\qquad |\tau|\le r_T(\square).
\end{equation}
Under the strong margin condition $(\mathrm{I3})^{\mathrm{seal}}_T$ of Lemma~\ref{5.5:lem-input-margin}, see
\eqref{5.5:eq-input-margin-a}, alternative $[Z]$ holds at every $(y,a)$, that is
$Q_a\le\theta_{T,a}-\rho_m\mathbf{r}^T_a$. With \eqref{5.5:eq-holomorphy-jet-1} this gives
\begin{equation*}
Q_a(\mathcal{W}_T)\le\theta_{T,a}-0.043\,\rho_m\mathbf{r}^T_a<\theta_{T,a}
\end{equation*}
on $|\tau|\le r_T(\square)$, which is the membership asserted in (a). The bound $\tfrac23\rho_m\mathbf{r}^T_a$ for
the increment in clause (o) is the case $c_E=1$ of this estimate, since
$0.638<\tfrac23$.

For the second assertion of (a), put $\alpha_2(y):=\rho_m\alpha_{\cdot,\ell}\,4B_\sharp^2/K^{+}$, as in clause (a),
where $\alpha_{\cdot,\ell}$ is the level-$\ell$ constant of the proof of clause (o).
The proof of clause (o) gives $N_y(\mathbf{A})\le 4B_\sharp^2[\hat p+|\tau|\hat p_\square](y)$, and at $\tau=0$ one
has $\mathbf{A}=\mathbf{H}_\ell$. Together with the bound on $\hat p$ that enters the arithmetic of that proof, and
with the hypothesis $17K^{+}\vartheta\le 1$ of the theorem, this yields
\begin{equation*}
N_y(\mathbf{H}_\ell)\le 4B_\sharp^2\hat p(y)\le 5.5\,K^{+}\vartheta\,\alpha_2(y)
\le\tfrac{5.5}{17}\,\alpha_2(y)\le\tfrac13\,\alpha_2(y),
\end{equation*}
because $5.5/17<0.3236<\tfrac13$. Thus $\mathbf{H}_\ell$ satisfies \cite[(3.14)]{B-CMP109} with $\tfrac13\alpha_2$,
the constant in the form stated in \cite[p.~273]{B-CMP109}.

Let now $|\tau|\le\mathsf{w}_0r_T(\square)$ and assume the margin condition $(\mathrm{I5})_T$ of
Lemma~\ref{5.5:lem-input-margin} and the window condition $28K^{+}\vartheta\le\mathsf{w}_0$. The proof of clause (o)
gives, on this smaller disc, an increment of $Q_a$ at most $c_E\cdot\tfrac{1001}{1960}\,\mathsf{w}_0\rho_m\mathbf{r}^T_a$. By
\eqref{5.5:eq-input-margin-a}, at $(y,a)$ one of the alternatives $[\mathfrak{T}]$, $[Z]$ holds.

At a $[\mathfrak{T}]$-clause, $Q_a\le\theta_{T,a}-\iota^1_l\mathbf{r}^T_a$ with $\iota^1_l=\tfrac23\mathsf{w}_0\rho_m$.
If $a\in\{1,3,6,7\}$, then $c_E=1$. If $a\in\{2,4\}$, then $c_E=\tfrac12(1+E/\mathbf{r}^T_a)$ and
$E<\theta_{T,a}\le(1+\beta)\mathbf{r}^T_a$, whence $c_E\le 1+\beta/2=1.1$, as $\beta=1/5$. In both cases
$c_E\le 1.1$, and $1.1\cdot\tfrac{1001}{1960}<0.5618$, so
\begin{equation*}
Q_a(\mathcal{W}_T)\le Q_a+0.5618\,\mathsf{w}_0\rho_m\mathbf{r}^T_a
<Q_a+\tfrac23\mathsf{w}_0\rho_m\mathbf{r}^T_a=Q_a+\iota^1_l\mathbf{r}^T_a\le\theta_{T,a}.
\end{equation*}
Without the condition $\theta_{T,a}\le(1+\beta)\mathbf{r}^T_a$ one still has $c_E<1.5$. This gives an increment at
most $0.7661\,\mathsf{w}_0\rho_m\mathbf{r}^T_a$, and, as $\mathsf{w}_0=1/24$ and
$\rho_m=\beta2^{-|m-\mathsf{j}-1|}\le\beta$,
\begin{equation*}
0.7661\,\mathsf{w}_0\rho_m\mathbf{r}^T_a\le 0.032\,\rho_m\mathbf{r}^T_a\le 0.0064\,\mathbf{r}^T_a .
\end{equation*}
This lies below the unspent one-third reserve of the interpolated threshold statement $\Im_l$, which is at least
$\tfrac13\beta\hat c\,\mathbf{r}^T_a=\mathbf{r}^T_a/30$. That reserve is used elsewhere only at pairs evaluated
through the box extension of Definition~\ref{5.5:def-block-piece}, by the threshold proposition at
$[\mathfrak{T}]$-locations, and at most once for each pair. Hence $Q_a(\mathcal{W}_T)<\theta_{T,a}$ at a
$[\mathfrak{T}]$-clause also without the condition $\theta_{T,a}\le(1+\beta)\mathbf{r}^T_a$, the excess of the
increment over $\iota^1_l\mathbf{r}^T_a$ being absorbed by that reserve.

At a $[Z]$-clause, $Q_a\le\theta_{T,a}-\rho_m\mathbf{r}^T_a$. With $c_E<1.5$, the same computation gives
\begin{equation*}
Q_a(\mathcal{W}_T)\le Q_a+0.032\,\rho_m\mathbf{r}^T_a<\theta_{T,a}.
\end{equation*}
This is the membership asserted in (a$'$).

For (a) and (a$'$) alike, $\mathfrak{U}_T$ is defined by the inequalities indexed by $\mathrm{Cl}(T)$. It is open,
and its defining
inequalities are strict. The bounds above hold in supremum form over the closed discs $|\tau|\le r_T(\square)$,
respectively $|\tau|\le\mathsf{w}_0r_T(\square)$, over $t\in[0,1]$ and over $(\mathbb{U},\mathbb{J})\in\mathfrak{K}_P$.
Membership therefore holds uniformly on the closed polydiscs. The position of the block enters the argument only
through the inequality $d_\ell(y,\square^{\zeta})\ge d_\ell(X_T,\square^{\zeta})$ for $y\in X_T$ (equality $0=0$
allowed) and, for (a$'$), through the window condition. Clause (a$'$) thus holds at every pair $(T,\square)$, with no
condition that $X_T$ be far from $\square$.

\emph{Clause (b).} By clause (a) of the theorem defining the response $\mathbb{H}$, applied at $\mathfrak{B}$,
$\mathbb{H}$ is holomorphic in $(\mathbb{U},\mathbb{J},b,\sigma)$. By \cite[Proposition~4]{B-CMP98}, with
\cite[(135)]{B-CMP98}, the map $(\mathbb{U},A)\mapsto Q_{\mathfrak{B}_\ell}(\mathbb{U};\eta A)$ is analytic bond by
bond. Writing $U_0$ for the configuration at which it is expanded, it has the form
\begin{equation*}
Q=Q(U_0)A+C(U_0,A),\qquad |C(U_0,A)|\le C_2|A|^2,
\end{equation*}
with a constant $C_2$ as in \cite[(135)]{B-CMP98},
and it is not affine in $A$. Since $A=(t+\tau\zeta_\square)\mathbb{H}$ is affine in $\tau$, the datum
$\tau\mapsto\mathsf{B}(t,\tau)=Q_{\mathfrak{B}_\ell}(\mathbb{U};\eta(t+\tau\zeta_\square)\mathbb{H})$ is holomorphic
near the real segment and continuous in $t$.

By clause (a) of the theorem defining $\mathbb{H}$ and clause (b) of the statement on the response at the capped
scaffold, both applied at $\mathfrak{B}_\ell$, the response at $\mathfrak{B}_\ell$ is holomorphic in its datum.
Hence $(A'_\ell,X_\ell,\mathbf{H}_\ell)$, taken at the datum $\mathsf{B}(t,0)$, are holomorphic in
$(\mathbb{U},\mathbb{J},b,\sigma)$ and continuous in $t$. So are their derivatives
$\delta A'_\ell,\delta X_\ell,\delta\mathbf{H}_\ell$ at $\tau=0$ of the responses at the datum $\mathsf{B}(t,\tau)$.

By Definition~\ref{5.5:def-own-level-family}, the quantities
\begin{equation*}
\mathbf{A}=\mathbf{H}_\ell+\tau\delta\mathbf{H}_\ell,\qquad A'_\ell+\tau\delta A'_\ell,\qquad X_\ell+\tau\delta X_\ell
\end{equation*}
are affine in $\tau$, and $\mathcal{W}_T$ is a holomorphic function of them. Its first component is
$e^{i\eta\mathbf{A}}\mathbb{U}$. Its second component is affine in $A'_\ell+\tau\delta A'_\ell$ and
$X_\ell+\tau\delta X_\ell$, plus $F(\mathbb{U},\mathbf{A})$. The function $F(\mathbb{U},\mathbf{A})$ is analytic for
$\mathbf{A}$ in the range of \cite[(3.14)]{B-CMP109}, and clause (o) keeps $\mathbf{A}$ in that range on the discs of
(a) and (a$'$). Hence $(\mathbb{U},\mathbb{J},b,\sigma,\tau)\mapsto\mathcal{W}_T(t,\tau)$ is holomorphic on the
source space times the discs of (a) and (a$'$), and continuous in $t$; this is (b).

The $\tau$-derivatives $\delta A'_\ell,\delta X_\ell,\delta\mathbf{H}_\ell$ are those of the responses at the exact
datum $\mathsf{B}(t,\tau)$; since $Q$ is not affine, $\mathsf{B}$ is not replaced by the linearised datum
$Q_\ell(\eta\zeta_\square\mathbb{H})$.

\emph{Clauses (d) and (d$'$).} Let $F$ be holomorphic on $\mathfrak{U}_T$ with $|F|\le\mathfrak{n}$. By (a),
respectively (a$'$), the restriction of $\mathcal{W}_T(t,\tau)$ to $X_T$ lies in $\mathfrak{U}_T$ on the closed
polydisc of that clause, so $|F\circ\mathcal{W}_T|\le\mathfrak{n}$ there. By (b), for fixed source and $t$, the map
$\tau\mapsto F(\mathcal{W}_T(t,\tau))$ is holomorphic on the disc of radius $r$, with $r=r_T(\square)$ for (d) and
$r=\mathsf{w}_0r_T(\square)$ for (d$'$). The Cauchy estimate for the first derivative at the centre
(Lemma~\ref{5.2:lem-disc-bounds}) gives the bound $\mathfrak{n}/r$, which we use in the weaker form $2\mathfrak{n}/r$
matching the constant $K''=28K^{+}$ below:
\begin{equation*}
\bigl|\partial_\tau|_{\tau=0}F(\mathcal{W}_T(t,\tau))\bigr|\le\mathfrak{n}/r\le 2\mathfrak{n}/r
\end{equation*}
for every $t\in[0,1]$. The integrand being continuous in $t$ by (b), the same bound holds for
$\bigl|\int_0^1dt\,\partial_\tau|_{\tau=0}F(\mathcal{W}_T(t,\tau))\bigr|$. By the definition
$r_T(\square)=e^{\frac14\delta d_\ell(X_T,\square^{\zeta})}/(14K^{+}\vartheta(\square))$ and by $K''=28K^{+}$,
$\mathsf{w}_0=1/24$ and $\vartheta^{\natural}=\vartheta/\mathsf{w}_0$, we have
\begin{align*}
\frac{2}{r_T(\square)}&=K''\vartheta(\square)\,e^{-\frac14\delta d_\ell(X_T,\square^{\zeta})},\\
\frac{2}{\mathsf{w}_0r_T(\square)}&=24K''\vartheta(\square)\,e^{-\frac14\delta d_\ell(X_T,\square^{\zeta})}
=K''\vartheta^{\natural}(\square)\,e^{-\frac14\delta d_\ell(X_T,\square^{\zeta})}.
\end{align*}
These are the bounds of (d) and (d$'$).

\emph{Clause (c).} At the real point of clause (c), with $\sigma\equiv1$, the configuration is critical for
$\mathfrak{B}_\ell$ at its own averages, by part (i) of the lemma on critical configurations for nested classes.
Clauses (d) and (e) of the theorem defining $\mathbb{H}$ and clause (d) of the statement on the response at the
capped scaffold, applied at $\mathfrak{B}_\ell$, then give the relations \cite[(3.6)--(3.7)]{B-CMP109}.

Consequently, at the real point, $\mathcal{W}_T$ agrees to first order in $\tau$, together with the current it
carries, with the configuration
\begin{equation*}
U_\ell\bigl(\exp iQ_\ell(\eta(t+\tau\zeta_\square)\mathbb{H})\,\bar U^{\ell}\bigr)
\end{equation*}
of \cite[p.~273]{B-CMP109}. This is the first-jet device \cite[(3.13)--(3.15)]{B-CMP109}: $\delta$ is the exact
derivative at $\tau=0$, and the chain rule gives the $\tau$-jet. The identity is derived on
\cite[p.~273]{B-CMP109} from the gauge invariance \cite[(3.5)--(3.6), p.~271]{B-CMP109}, and holds at the level of
configurations modulo gauge. It
therefore applies equally to the pointed constituents $v$ of $T$, whose clause sets satisfy
$\mathrm{Cl}(v)=\mathrm{Cl}(T)$. The position of $X_T$ (the domain $X$ of \cite[(3.9)]{B-CMP109}) enters only that
bound.

Since $\partial_\tau|_{\tau=0}T(X_T;\cdot)$ depends only on the first-order jet in $\tau$, we obtain
\begin{equation*}
\begin{split}
&\int_0^1dt\,\partial_\tau|_{\tau=0}T(X_T;\mathcal{W}_T(t,\tau))\\
&\quad=\int_0^1dt\,\partial_{t_\square}|_{0}\,T\bigl(X_T;U_\ell(\exp iQ_\ell(\eta(t+t_\square\zeta_\square)\mathbb{H})
\bar U^{\ell})\bigr),
\end{split}
\end{equation*}
which is $\operatorname{piece}(T,\square)$ by Definition~\ref{5.5:def-block-piece}. This is (c).

\emph{The summation statement for the own-level family.} This statement asserts, for a pair with
$X_T\cap\tilde{\square}\neq\emptyset$, that $\operatorname{diam}X_T\ge 2Ms_i$, so that the pair lies in the class
$(\alpha)$ of the far-pair summation statement, and a second assertion covers the case
$X_T\cap\tilde{\square}=\emptyset$; here the length scale $s_i$ of the diameter bound, and the constants $K_S$ and
$E_0$ of the summation bound in the display below, are those fixed in the far-pair summation statement. The
statement is used for one purpose only, the bound
\begin{equation*}
\sum_{X_T\not\subset\tilde{\square}^2}\|T\|_\infty\,e^{-\frac14\delta d_\ell(X_T,\square^{\zeta})}\le K_S(E_0+1),
\end{equation*}
a sum over far pairs $X_T\not\subset\tilde{\square}^2$, for which it holds. Near pairs, for which a domain
$X_T\subset\tilde{\square}^2$ of small diameter may meet $\tilde{\square}$ without the diameter bound, are not summed
in this way; they are counted by clause (3) of the corollary counting near pairs. Apart from this summation statement,
$X_T$ enters
Definition~\ref{5.5:def-own-level-family} and the proof of clauses (o)--(d$'$) only as the set on which sizes are
evaluated. This completes the proof of clauses (a)--(d$'$).

\emph{Clause (e).} Clause (e) is the inventory of the objects on which the proof of clauses (o)--(d$'$) depends. It
is obtained by going through that proof, from its first step through the proof of clause (c) above, and recording
at each step the objects used there. The scaffold $\mathfrak{B}$ enters through the conditions on the scaffold in
the setting of Theorem~\ref{5.5:thm-own-level-membership} and
through its cap $\mathfrak{B}_\ell$; the pair $(\mathbb{U},\mathbb{J})$ through the space
$\mathfrak{K}_{109}(\mathfrak{B}_\ell;\gamma_0)$ and through the margin condition $(\mathrm{I3})^{\mathrm{seal}}_T$,
respectively $(\mathrm{I5})_T$; the function $b$ through $p$; and $\sigma$ through the decoupled list of the theorem
defining $\mathbb{H}$, at $\mathfrak{B}$ for the datum and at $\mathfrak{B}_\ell$ for the responses. The position of
$X_T$ enters only through $d_\ell(X_T,\square^{\zeta})\ge0$; no weight, diameter, hull or shape of $X_T$, zone of
$\square$, or age is used.

\emph{Clause (f).} For two runs $\mathsf{A}$, $\mathsf{B}$ on one scaffold, the family $\mathcal{W}^{\mathsf{m}}_T$,
$\mathsf{m}\in\{\mathsf{A},\mathsf{B}\}$, is in each run the same construction of
Definition~\ref{5.5:def-own-level-family}, with one parameter domain $\mathfrak{P}$. For a pair
$(\Phi^{\mathsf{B}},\Phi^{\mathsf{A}})$ of the pair system $\mathcal{K}$ put
\begin{equation*}
\delta F(\Phi^{\mathsf{B}},\Phi^{\mathsf{A}}):=F^{\mathsf{B}}(\Phi^{\mathsf{B}})-F^{\mathsf{A}}(\Phi^{\mathsf{A}}),
\end{equation*}
the difference of the two runs' functions at the pair, and write $\sup_{\mathcal{K}}|\delta F|$ for its supremum
over $\mathcal{K}$. By the
two-run principle, the difference
\begin{equation*}
\varpi\longmapsto F^{\mathsf{B}}\circ\mathcal{W}^{\mathsf{B}}_T-F^{\mathsf{A}}\circ\mathcal{W}^{\mathsf{A}}_T,
\qquad \varpi=((b,\sigma,t,\tau);p),
\end{equation*}
is holomorphic in $(b,\sigma,\tau)$ and continuous in $t$ on $\mathfrak{P}$. It is bounded in modulus by
$\sup_{\mathcal{K}}|\delta F|$ by hypothesis (k1) of the
pair system $\mathcal{K}$ at the map $\mu^{\mathrm{own}}$. The Cauchy estimate used for (d$'$) is linear in the
function to which it is applied. Hence (d$'$) holds for the difference, with $\mathfrak{n}$ replaced by
$\sup_{\mathcal{K}}|\delta F|$. This proves (f) and completes the proof of the theorem.
\end{proof}

\begin{remark}\label{5.5:cor-three-class-filing}
The statement of this corollary is not written out here; parts of its proof are printed below.
\end{remark}

\begin{proof}[Proof of Corollary~\ref{5.5:cor-three-class-filing}, continued: clauses (3)--(6), and the exponentially small weight of straddling cores]
\label{5.5:prf-straddling-cores}
Throughout, $v$, $X$, $j$, $y$, $n$, $\nu$, $t$, $\Theta$ and $\mathfrak{n}_v$ are as in the statement of
Corollary~\ref{5.5:cor-three-class-filing}, and we write $\mathfrak{n}:=\mathfrak{n}_v$.

\emph{Clause (3).} The bounds listed in clause (3) are exactly those supplied by clauses (iv), (a2)
and (b) of the weight lemma. We check that the two halves of the near-level clause (iv) lie inside
the common bound with the constant $C_I'$ of clause (3). Let $v$ be a read core with $S^{*}_v\le 1$,
and put $z:=\Theta\rho_{n,j}$. The inequality
$(\Theta C_M\mathsf{d}_v t^2\rho_{n,j})^{-1}\le S^{*}_v$, stated in
Corollary~\ref{5.5:cor-three-class-filing} with $S^{*}_v$ as in
the corresponding display, together with $S^{*}_v\le 1$ gives
\begin{equation}\label{5.5:eq-straddling-cores-1}
C_M\,\mathsf{d}_v\,t^2 z\ \ge\ 1 .
\end{equation}
For a core the side of the own box is $\mathsf{d}_v=M(2\nu+3)$, and $2\nu+3\le 5\nu$ because
$\nu\ge 1$; thus $\mathsf{d}_v\le 5M\nu$, and \eqref{5.5:eq-straddling-cores-1} yields
\begin{equation*}
t^2\ \ge\ (5C_M M\nu z)^{-1}\ \ge\ (10C_M M\nu z)^{-1}.
\end{equation*}
Since $\hat{\alpha}\ge 0$, we may raise this to the powers $2+\hat{\alpha}/2$ and $\hat{\alpha}/2$;
multiplying by $\nu^3z^2(1+z)$ and by $\nu^3(1+z)$ respectively, we obtain
\begin{align*}
\nu^3z^2(1+z)\,t^{4+\hat{\alpha}}
 &\ \ge\ \nu^{1-\hat{\alpha}/2}\,z^{-\hat{\alpha}/2}(1+z)\,(10C_MM)^{-2-\hat{\alpha}/2},\\
\nu^3(1+z)\,t^{\hat{\alpha}}
 &\ \ge\ \nu^{3-\hat{\alpha}/2}\,z^{-\hat{\alpha}/2}(1+z)\,(10C_MM)^{-\hat{\alpha}/2}.
\end{align*}
Here $\nu\ge 1$ and $\hat{\alpha}\le 1$, so both powers of $\nu$ are at least one; and
$z^{-\hat{\alpha}/2}(1+z)\ge 1$, since for $z\le 1$ the first factor is at least one, while for
$z\ge 1$ we have $1+z\ge z\ge z^{\hat{\alpha}/2}$. Since moreover $10C_MM\ge 1$, and
$2+\hat{\alpha}/2\le 5/2$ and $\hat{\alpha}/2\le 1/2$ by $\hat{\alpha}\le 1$, the two halves of
clause (iv) become
\begin{equation}\label{5.5:eq-straddling-cores-2}
\nu^3z^2(1+z)\,t^{4+\hat{\alpha}}\ \ge\ (10C_MM)^{-5/2},
\qquad
\nu^3(1+z)\,t^{\hat{\alpha}}\ \ge\ (10C_MM)^{-1/2}.
\end{equation}
By its definition in clause (3), $C_I'\ge 8C(\kappa)(10C_MM)^{5/2}$ and
$C_I'\ge 4C'(10C_MM)^{1/2}$. The first half of \eqref{5.5:eq-straddling-cores-2} therefore gives
\begin{equation*}
\tfrac12 C_I'\,\mathfrak{n}\,\nu^3z^2(1+z)\,t^{4+\hat{\alpha}}
\ \ge\ 4C(\kappa)(10C_MM)^{5/2}\,\mathfrak{n}\,(10C_MM)^{-5/2}
\ =\ 4C(\kappa)\,\mathfrak{n},
\end{equation*}
and the second half gives
\begin{align*}
\tfrac12 C_I'\,\mathfrak{n}\,\nu^3z^2(1+z)\,t^{4+\hat{\alpha}}
&=\tfrac12 C_I'\,\mathfrak{n}\,z^2t^4\cdot\nu^3(1+z)\,t^{\hat{\alpha}}\\
&\ge 2C'(10C_MM)^{1/2}\,\mathfrak{n}\,z^2t^4\,(10C_MM)^{-1/2}
\ =\ 2C'\,\mathfrak{n}\,z^2t^4 .
\end{align*}
Adding the two, $C_I'\mathfrak{n}\nu^3z^2(1+z)t^{4+\hat{\alpha}}$ is at least
$4C(\kappa)\mathfrak{n}+2C'\mathfrak{n}z^2t^4$, which is twice the sum of the bound
$2C(\kappa)\mathfrak{n}$ on the constituents, taken singly on $\mathcal{W}_T$, and the bound
$C'\mathfrak{n}z^2t^4$ on the counterterm in clause (3). The counterterm keeps the data form of
the bound in clause (a2) of the weight lemma because it belongs to class W and stays on the box
pair $\mathfrak{b}'_j$. The bound on the pieces $v_p$ in clause (2) of
Corollary~\ref{5.5:cor-three-class-filing} is linear in $\mathfrak{N}(v)$; the tabulated values
$\mathfrak{N}(v)=2\mathfrak{n}_v$ are therefore used as they stand, and at own-read units only
the prefactor of that bound changes, namely to the one displayed in clause (2).

\emph{Straddling cores.} We show that the near constituents of a straddling core carry the weight
$\eta_p\le t^5$ at every $\nu$. Let $v$ be a core with $y\in\tilde{\square}^2$ and own box
$\mathfrak{r}_v=\square_y^{\sim\nu+1}\not\subset\tilde{\square}^2$, and let $X\ni y$ be a
constituent with $X\subset\tilde{\square}^2$. Then $(X,\square)$ is a near pair whose term lies
in $\mathcal{T}_1$ (Definition~\ref{5.5:def-near-units}), and
$X\subset\square_y^{\sim\nu}\subset\mathfrak{r}_v$.

Let first $\nu\ge 2$. In units of level $j$ the side of $\mathfrak{r}_v$ is
$\mathsf{d}_v=M(2\nu+3)$, and
\begin{equation*}
M(2\nu+3)\ <\ ML^{\nu-1}\ =\ ML^{n-j}\ \le\ ML^{n_0-j}.
\end{equation*}
The first inequality is $2\nu+3<L^{\nu-1}$, valid for $\nu\ge 2$ and $L\ge 8$: at $\nu=2$ it
reads $7<8\le L$, and passing from $\nu$ to $\nu+1$ adds $2$ to the left member and multiplies
the right member by $L\ge 8$. The equality is $\nu-1=n-j$, and the last inequality is $n\le n_0$,
which holds as $n=\min\{\ell(y),n_0\}$. Now $ML^{n_0-j}$ is the side of one cube of $\pi_{n_0}$
in units of level $j$, and $\tilde{\square}^2\setminus\tilde{\square}$ is a full layer of cubes of
$\pi_{n_0}$. A set of side smaller than the thickness of that layer which is not contained in
$\tilde{\square}^2$ cannot meet $\tilde{\square}$. Hence $\mathfrak{r}_v\cap\tilde{\square}=\emptyset$,
and therefore $X\cap\tilde{\square}=\emptyset$.

We now use the second case of the block-sum separation bound, which belongs to the own holomorphic
extension of the pieces in Definition~\ref{5.5:def-block-piece}: if
$X_T\cap\tilde{\square}=\emptyset$, every contour from $X_T$ to the support
$\square^{\zeta}$ of $\zeta_{\square}$ crosses a cube of $\pi_i$ in units of level at most
$\ell$, so that $d_\ell\ge c_{\mathrm{sep}}L^{(i-\ell)^+}$, where $c_{\mathrm{sep}}>0$ denotes the
separation constant of that bound. This rests on the lower bound, stated in
\cite[p.~276]{B-CMP119}, for the distance from $\square$ of a localisation domain disjoint from
$\tilde{\square}$. We apply it with $i=n_0$ and $\ell=j$ and obtain
$d_j(X,\square^{\zeta})\ge c_{\mathrm{sep}}L^{n_0-j}$. Consequently, with $s:=n_0-j\ge 0$,
\begin{equation}\label{5.5:eq-straddling-cores-3}
\eta_p\ =\ e^{-\frac14\delta d_j(X,\square^{\zeta})}
\ \le\ e^{-\frac14\delta c_{\mathrm{sep}}L^{s}}
\ \le\ L^{-5s}
\ \le\ L^{-5(n-j)}\ =\ t^5 .
\end{equation}
The middle inequality uses the floor $\delta c_{\mathrm{sep}}\ge 16\ln L+4$ in the setting of
Theorem~\ref{5.5:thm-own-level-membership}: it gives
\begin{equation*}
\tfrac14\delta c_{\mathrm{sep}}L^{s}\ \ge\ (4\ln L+1)L^{s}\ \ge\ 5s\ln L
\qquad\text{for every } s\ge 0,
\end{equation*}
the second step because $4L^s\ge 5s$. The next inequality in \eqref{5.5:eq-straddling-cores-3}
is $n\le n_0$, and the final equality is $t=L^{j-n}$.

Let now $\nu=1$. Then $n=j$, so $t=L^{j-n}=1$; and $\eta_p\le 1$ because
$d_j(X,\square^{\zeta})\ge 0$, so that $\eta_p\le t^5$. In this case nothing is claimed about
$\mathfrak{r}_v\cap\tilde{\square}$. Together with \eqref{5.5:eq-straddling-cores-3} this gives
$\eta_p\le t^5$ at every $\nu$.

Clause (a$'$) of Theorem~\ref{5.5:thm-own-level-membership} applies to each such $X$, with $T$ the
term of $X$ (every such $T$ belongs to $\mathcal{T}_1$), and the bound on the pieces of $X$ in
clause (2) of Corollary~\ref{5.5:cor-three-class-filing} carries the factor $\eta_p\le t^5$. The
bound on the number and weight of these constituents in the data of the weight lemma is supplied
by the weight lemma and is not part of the present proof.

\emph{Clause (4).} The formula \eqref{5.5:eq-block-piece-a} of
Definition~\ref{5.5:def-block-piece} defines the piece of $T$ at $\square$ by one and the same
formula for every block $\square$. Clause (c) of Theorem~\ref{5.5:thm-own-level-membership}, the
real-point statement of the box extension in Definition~\ref{5.5:def-block-piece}, and clause (c)
of the remainder statement for the box construction identify each piece as the real jet of its
own holomorphic extension. The sum identity of Definition~\ref{5.5:def-block-piece},
\begin{equation*}
\sum_{\square}\operatorname{piece}(T,\square)\ =\ T(\Phi^1)-T(\Phi^0)
\qquad\text{at the real point},
\end{equation*}
sums these real pieces, and therefore holds whichever class each piece is placed in.

It remains to check that the own-read pieces lie in the class required of summands of the near
element in clause (4), that is, that they are holomorphic and $\mathsf{s}$-local. Holomorphy is
clause (b) of Theorem~\ref{5.5:thm-own-level-membership}. For locality: the response field
$\mathbf{H}_\ell$ at the cap $\mathfrak{B}_\ell$, its first $\tau$-derivative
$\delta\mathbf{H}_\ell$ at $\tau=0$, and the two further fields that the response theorem
furnishes at $\mathfrak{B}_\ell$ are the responses, in the sense of the response theorem, to a
datum built from $\mathbb{H}$ at $\mathfrak{B}$. All of them belong to the list of decoupled fields
by the first step of the decoupling argument, which holds at $\mathfrak{B}$ and at each
$\mathfrak{B}_\ell$. Restricted to $X_v\subset\tilde{\square}^4$ they are $\mathsf{s}$-local by the
locality clause (f-ii) of the two-run comparison statement. Hence the own-read pieces are
summands of the near element, filed among the summands assigned to it in clause (4), as clause (4)
asserts.

\emph{Clause (5).} The bound in clause (5) is clause (f) of
Theorem~\ref{5.5:thm-own-level-membership} applied with $F:=f_v$ at the map
$\mu^{\mathrm{own}}$; its hypothesis there is the input (k1) read at $\mu^{\mathrm{own}}$.

\emph{Clause (6).} Clause (6) assigns each bound on the pieces of class W to the statement that
supplies it and asserts nothing beyond these assignments. Its entry comparing the two runs
$\mathsf{A}$ and $\mathsf{B}$ rests on the input (k1) read at the two-run comparison statement,
at the W-pieces of every near unit.
\end{proof}

\begin{remark}[Own-level filing against the printed filing]\label{5.5:rem-own-level-filing}
Let $\square$ be a block, and let $T\in\mathcal{T}_1$ be a term of level $\ell$ with domain $X_T$, in the sense of Definition~\ref{5.5:def-near-units}. For a near unit $v$ at $\square$ let $S^{*}_v$ be the chart radius, and let $\nu$ and $\mathfrak{n}_v$ be as in the three-class filing of near-unit pieces (Corollary~\ref{5.5:cor-three-class-filing}). We use two extensions of the pieces:
\begin{itemize}
\item the box extension of Definition~\ref{5.5:def-block-piece}, on the box pair $\mathfrak{b}'_j$, with $t\in[0,1]$ and $\mathsf{t}\in\mathbb{D}_{\square}$;
\item the own-level family $\mathcal{W}_T$ with radius $r_T(\square)$ of Definition~\ref{5.5:def-own-level-family}.
\end{itemize}
As noted in Definition~\ref{5.5:def-block-piece}, the real points of the pieces do not determine which of these two extensions a piece carries. That choice is therefore a convention, fixed once for the whole construction. We call it the \emph{filing} of the piece, and fix it as follows.

\begin{enumerate}
\item[(a)] The \emph{box units} are the member units $v$ with $\nu\ge 3$ that are a core or a short single and satisfy $S^{*}_v>1$. Their member pieces carry the box extension on $\mathfrak{b}'_j$; we call these units and their member pieces \emph{box-read}.
\item[(b)] Every other pair $(T,\square)$ with $T\in\mathcal{T}_1$ is \emph{own-read}. Such a term is evaluated on $\mathcal{W}_T(t,\tau)$ with
\begin{equation*}
|\tau|\le \mathsf{w}_0\, r_T(\square),
\end{equation*}
where $\mathsf{w}_0$ is the window fraction; when the pair belongs to a unit $v$, the bound used for its term is $\mathfrak{N}(v)=2\mathfrak{n}_v$. Own-read pieces remain summands of the one near element attached to $\square$ in the summation of the near-unit pieces.
\item[(c)] There is one exception to (b): the W-pieces of class W of Definition~\ref{5.5:def-near-units} and Corollary~\ref{5.5:cor-three-class-filing}, that is, the counterterm companions of the near units and the block pieces over blocks $\square'\subset\tilde{\square}^2$, a block piece being filed plaquette-wise at the level of its cell. They carry the box extension on $\mathfrak{b}'_j$ at every near unit, for every value of $S^{*}_v$ and for every $\nu$, and never the own-level family $\mathcal{W}_T$.
\end{enumerate}

The following statements, used below by statement, are read with this filing.
\begin{itemize}
\item The box extension on $\mathfrak{b}'_j$ carries the hypothesis $X_T\subset\tilde{\square}^2$. This hypothesis, and the requirement that all box units near $\square$ be read on one background, concern the box-read units and the W-pieces only.
\item The own-level extension carries the hypothesis $X_T\not\subset\tilde{\square}^2$. That hypothesis is read as describing which pieces are summed there; this is legitimate because the family $\mathcal{W}^{\ell,\square}$ of Definition~\ref{5.5:def-own-level-family} is defined without reference to any domain $X_T$, the condition $X_T\not\subset\tilde{\square}^2$ entering only as the label of the sum. Its statements cover the own-read member pieces as well. Each such piece is an own-level layer in the sense of Definition~\ref{5.5:def-own-level-family}, sized by the layer bound of the own-level extension with the coarse-data damping $\varrho$, and it passes the layer test of the own-level extension by the argument of clause (iv) of that extension, as the pieces with $X_T\not\subset\tilde{\square}^2$ do.
\item Two further clauses stated for the box extension, and the part of the attribution line that refers to the box extension, are asked of box-read units and W-pieces alone.
\item In the radius statement, in which the radii $\mathbf{r}^T_a$ and the interpolated threshold statement $\Im_l$ are fixed, own-read pairs take the first radius rule and are served by its clauses (v) and $(\gamma)$ and by part (b) of its auxiliary hypothesis. The second radius rule, clause $(\delta)$, part (c) of the auxiliary hypothesis, the proposition on the locations of type $[\mathfrak{T}]$, one further lemma bound of that statement, and clause $(c')$ are asked of box-read units and W-pieces alone.
\item Own-read pieces stay in the total of the near element in the summation of the near-unit pieces, and the definition of the quantity $\mathfrak{d}$ entering that summation is unchanged.
\end{itemize}

In \cite[(3.5), (3.18)--(3.25)]{B-CMP109} and \cite[(1.59)]{B-CMP122b}, every domain $X\subset\tilde{\square}^2$ is expanded on the box. Domains $X\not\subset\tilde{\square}^2$ are treated by the own-level device or the raw layer, as in \cite[(3.6)--(3.15)]{B-CMP109} and \cite[(1.60)]{B-CMP122b}.

The filing (a)--(c) differs from that assignment in one respect only. The own-read near member pieces have $X\subset\tilde{\square}^2$ but carry the own-level family. The filing changes which extension these pieces carry, and no inequality of \cite{B-CMP109} or \cite{B-CMP122b} is changed.
\end{remark}

\begin{proof}
\emph{Each piece is assigned exactly one extension.} By Corollary~\ref{5.5:cor-three-class-filing}, every piece of the term of a near unit $v$ at $\square$ lies in exactly one of three classes. If it is a member piece of a box unit, clause (a) assigns it the box pair $\mathfrak{b}'_j$. If it is a W-piece of class W, clause (c) assigns it $\mathfrak{b}'_j$, whatever $S^{*}_v$ and $\nu$ are. Every remaining piece is own-read and carries $\mathcal{W}_T$ by (b). A pair $(T,\square)$ with $T\in\mathcal{T}_1$ that is not near is own-read by (b) as well. Hence each piece is assigned exactly one of the two extensions.

\emph{The box-read pieces use the inequalities of \cite{B-CMP109} and \cite{B-CMP122b} unchanged.} The box-read member pieces and the W-pieces have $X\subset\tilde{\square}^2$ and carry the box extension. This is the assignment of \cite[(3.5), (3.18)--(3.25)]{B-CMP109} and \cite[(1.59)]{B-CMP122b}.

Consider a box unit that is a core. By (a), all of its constituents lie on the one box pair $\mathfrak{b}'_j$. Its counterterm is a W-piece, so by (c) it lies on the same $\mathfrak{b}'_j$. Hence the core, all its constituents and its counterterm sit on the one pair $\mathfrak{b}'_j$, and the requirement that all box units near $\square$ be read on one background is kept. For these pieces, every inequality of \cite{B-CMP109} and \cite{B-CMP122b} is used in the form stated there.

\emph{The own-read pieces use only the own-level inequalities.} Two facts about Ba{\l}aban's construction are used.

First, the own-level device \cite[(3.6)--(3.15)]{B-CMP109} is applied on the basis of sizes \cite[p.~271]{B-CMP109}, \cite[p.~372]{B-CMP122b}. It is applied at pairs that are not near, and these include large domains $X$ that meet $\tilde{\square}$. Its inequalities are therefore used at pairs whose domains meet the block.

Second, the box at domains $X\subset\tilde{\square}^2$ serves the renormalisation cancellation \cite[p.~273]{B-CMP109}. There, the constituents of a core and its counterterm are expanded on one background, which is what the cancellation between the core and its counterterm requires.

The own-read units use no such cancellation. By Corollary~\ref{5.5:cor-three-class-filing}, an own-read core, namely one with $\nu\le 2$ or with $S^{*}_v\le 1$, is read constituent by constituent, and no cancellation between its constituents and its counterterm is read. The counterterm itself remains on $\mathfrak{b}'_j$ by (c).

Consequently, reading an own-read near member piece on $\mathcal{W}_T$, with $|\tau|\le\mathsf{w}_0 r_T(\square)$ and bound $\mathfrak{N}(v)=2\mathfrak{n}_v$, uses the own-level inequalities only. These are the inequalities already used at non-near pairs whose domains meet the block. The reason for the box, the cancellation, does not arise for these pieces.

\emph{The total is unchanged.} The own-read pieces remain summands of the near element at $\square$, and its total is defined as before. This proves the last assertion of the remark.
\end{proof}

\begin{remark}[The box pair and the factor $24$]\label{5.5:rem-box-pair}
Throughout, own-read and box-read units are the classes of Corollary~\ref{5.5:cor-three-class-filing}
(three-class filing of near-unit pieces), $j$ is the level of the unit, $\mathfrak{b}'_j$ is the box
pair at that level, $\mathcal{W}_T$ is the
own-level family of Definition~\ref{5.5:def-own-level-family} with the radius $r_T(\square)$ of
\eqref{5.5:eq-own-level-family-a}, and $\mathsf{w}_0 = 1/24$ is the window fraction. The
\emph{threshold induction} is the induction on the thresholds $\theta_{T,a}$ of the terms
$T\in\mathcal{T}_1$, with its two location types $[\mathfrak{T}]$ and $[Z]$ and its interpolated
threshold statement $\Im_l$. The \emph{piece lemma} is the lemma which, from its hypothesis at the
constant $\mathfrak{N}(v)$ on the disc of radius $\mathsf{w}/\eta_p$, bounds each piece $v_p$ of a
unit $v$ by $(2/\mathsf{w})\,\eta_p\,e^{-(\kappa_1-1)m(p)}\,\mathfrak{N}(v)$; here $m(p)$ is the level
of the cell of the piece $p$.

\begin{enumerate}
\item[(a)] \emph{Membership of the box pair at own-read units.} No statement of this section asserts
or uses that the box pair $\mathfrak{b}'_j$ is a member of the domain of an own-read unit. This
membership is not obtained in general by applying, on the own box of the domain $X$ of the unit, the
proposition of the threshold induction which, at box-read pairs, spends the unspent reserve of the
interpolated threshold statement $\Im_l$: there are a near box $X\subset\tilde{\square}^2$ which is
not read in the sense of Corollary~\ref{5.5:cor-three-class-filing} and lies in the slab of a tile
block, in the position fixed by the statement used in the proof below, and a real point of the source
configuration space at which the box pair $P:=\mathfrak{b}'_j$ has the following property. For
$g\in G^{\mathbb{C}}$ write $P^g$ for the gauge transform of $P$ by $g$; every $g\in G^{\mathbb{C}}$
for which $P^g$ is of the exponential form required by condition (ii) of \cite[p.~262]{B-CMP109},
written $e^{i\xi\phi}$ with $\phi$ the field and $\xi$ the parameter entering that form, satisfies
\begin{equation*}
  \sup_X |\phi| \ge \alpha_{1,j} .
\end{equation*}
Hence no representative of $P$ whose $G$-part is identically $1$ satisfies condition (ii) of
\cite[p.~262]{B-CMP109} on $X$, although the restriction of $P$ to $X$ is a member of the domain of
the unit on $X$. The
same source point satisfies the conclusion of Theorem~\ref{5.5:thm-own-level-membership} (own-level
membership at every block). The example contradicts no statement of this book.

\item[(b)] \emph{The threshold induction.} The increments on the right of
\eqref{5.5:eq-own-level-membership-a} are proportional to $\rho_{m(y)}$, so the own-level membership
draws on no fixed reserve of the thresholds. The factor $1/\mathsf{w}_0$ by which the $\tau$-disc of
clause (a$'$) of Theorem~\ref{5.5:thm-own-level-membership} is smaller than that of clause (a) is the
cost of the $[\mathfrak{T}]$-clauses. This is the same factor that the threshold induction charges
through $\vartheta^{\natural} = \vartheta/\mathsf{w}_0$. The clause of that induction which treats
near pairs covers the own-read near pairs. Three further clauses of that induction, and the
proposition spending the unspent reserve of $\Im_l$, concern box-read units.

\item[(c)] \emph{The factor $24$.} For a box-read unit the radius $r_{\square}^{\sharp}$ of the box
disc $\mathbb{D}_{\square}$ satisfies $2/r_{\square}^{\sharp} = K_A\vartheta(\square)$. For an
own-read unit the prefactor of the own-read piece bound in clause (2) of
Corollary~\ref{5.5:cor-three-class-filing} is $24K''\vartheta(\square)\eta_p$. The ratio of the second
to the first is at most $24K''/K_A$. No inequality between $K''$ and $K_A$ is available; the threshold
induction keeps both under $K_{\min}$, so the ratio is carried as $24K''/K_A$.

The constant $C_R$ and the constant $\kappa_R$ are least numbers, fixed before $\gamma$. Here $C_R$
is the constant in the per-block bounds $\mathfrak{n}_0$ and $\mathfrak{N}_0$; it serves every block,
and the bound of the piece lemma at constant weight equal to one lies inside it. The constant
$\kappa_R$ is the least number of the weight lemma for the bounds which the piece lemma forms on the
disc assigned to each unit; it serves the blocks at which $\mathcal{W}_T$ is used. The maxima
defining $C_R$ and $\kappa_R$ include, from the bounds at own-read units, the term
$\max\{1,\ 24K''/K_A\}$; in particular
\begin{equation}\label{5.5:eq-box-pair-1}
  C_R \ge \max\{1,\ 24K''/K_A\}, \qquad \kappa_R \ge \max\{1,\ 24K''/K_A\}.
\end{equation}
These two, and likewise the weight lemma's own least constant and its least number
$\mathfrak{d}^{\sharp}$, enter with degree $0$ and in no exponent.

The threshold induction carries the same number in the form $K''\vartheta^{\natural}(\square)$, where
$\vartheta^{\natural} = \vartheta/\mathsf{w}_0$. With $C_R$ taken in this way, no inequality fixed
earlier in the order of choice of the constants changes.
\end{enumerate}
\end{remark}

\begin{proof}
\emph{Clause (a).} The box $X$, the real source point and the lower bound on $\sup_X|\phi|$ are those
of the statement exhibiting a near box, not read in the sense of
Corollary~\ref{5.5:cor-three-class-filing}, and a real source point at which box-pair membership is
not obtained through the proposition spending the unspent reserve of $\Im_l$; that statement is used
here as stated.

Since the displayed lower bound holds for every $g\in G^{\mathbb{C}}$ for which $P^g$ has the
exponential form of condition (ii) of \cite[p.~262]{B-CMP109}, no representative of $P$ whose $G$-part
is identically $1$ satisfies that condition on $X$. This is the failure asserted in (a).

In that statement the family $\mathcal{W}_T$ keeps the entry $\bar U$ of the real source point, and
condition (i) of \cite[p.~262]{B-CMP109} is a local condition; by these two facts the
point satisfies the conclusion of Theorem~\ref{5.5:thm-own-level-membership}. Since no statement of
this section uses box-pair membership at own-read units, the example contradicts none of them.

\emph{Clause (b).} That the reduction of the $\tau$-disc by the factor $1/\mathsf{w}_0$ is what the
$[\mathfrak{T}]$-clauses require is the content of the closing step of
the proof of Theorem~\ref{5.5:thm-own-level-membership}, which establishes its clauses (a) and (a$'$).
At a $[\mathfrak{T}]$-clause one has $c_E\le 1.1$, and, under the window condition of clause (a$'$),
by clause (o) of that theorem the increment on $|\tau|\le\mathsf{w}_0 r_T(\square)$ is at most
$0.511c_E\mathsf{w}_0\rho_{m(y)}\mathbf{r}^T_a(y)$.
Since $0.511\cdot 1.1 < 2/3$ and $\iota^1_l(y) = \tfrac23\mathsf{w}_0\rho_{m(y)}$, this gives
\begin{equation*}
  0.511\,c_E\,\mathsf{w}_0\,\rho_{m(y)}\mathbf{r}^T_a(y) < \iota^1_l(y)\,\mathbf{r}^T_a(y).
\end{equation*}
On the larger disc $|\tau|\le r_T(\square)$ the bound of clause (o) is
$0.957\rho_{m(y)}\mathbf{r}^T_a(y)$, which exceeds $\iota^1_l(y)\mathbf{r}^T_a(y)$. In the same closing
step the unspent reserve of $\Im_l$ is spent at box-read pairs only, once per pair, and the clause of
the threshold induction which treats near pairs is applied to the own-read near pairs.

\emph{Clause (c).} By Corollary~\ref{5.5:cor-three-class-filing}, $\eta_p$ is an exponential
$e^{-\frac14\delta d_j(X_v,\square^{\zeta})}$ of a non-positive number, since the distance is
non-negative. Hence $\eta_p \le 1$, and
\begin{equation}\label{5.5:eq-box-pair-2}
  \frac{24K''\vartheta(\square)\eta_p}{K_A\vartheta(\square)} = \frac{24K''}{K_A}\,\eta_p
  \le \frac{24K''}{K_A}.
\end{equation}

Since $C_R$ and $\kappa_R$ are least numbers fixed before $\gamma$ whose defining maxima contain the
term $\max\{1,\ 24K''/K_A\}$, they satisfy \eqref{5.5:eq-box-pair-1}.

Finally, $\vartheta^{\natural}(\square) = \vartheta(\square)/\mathsf{w}_0$ and $\mathsf{w}_0 = 1/24$.
Therefore
\begin{equation*}
  K''\vartheta^{\natural}(\square) = 24K''\vartheta(\square).
\end{equation*}
\end{proof}

\begin{definition}[Two expansions at equal depth: objects, stored domains and uses]
\label{5.5:not-paired-expansions}
The following notation compares two runs of the expansion at equal depth, with two cutoffs. It
has three parts: \textup{(1)} the common frame, \textup{(2)} the objects and their stored
domains, \textup{(3)} the uses.

\smallskip
\noindent\textup{(1)} \emph{The frame.} The letters $\mathsf{A}$ and $\mathsf{B}$ denote the two
runs compared, $\mathsf{A}$ with $n$ and $\mathsf{B}$ with $n+1$ renormalisation steps. The letters
$\mathsf{m}$, $\mathsf{m}'$ range over $\{\mathsf{A},\mathsf{B}\}$. Levels are indexed as in
$\mathsf{A}$, and the labels of the two expansions are aligned. At equal depth the two expansions
have one and the same scaffold, namely:
\begin{itemize}
\item the partitions, the labels, the covers, the boxes and the caps;
\item the radii $r^{\sharp}_{\square}$, $r_T(\square)$, $\mathsf{w}_0$ and $\eta_p$;
\item the frozen weights $w$, $w^{\vee}$ and the label data $\Theta$, $S_v$, $Q_v$,
$\mathfrak{r}_v$;
\item the predicate which assigns each term to the box class, to the own-level class or to the
Wilson class;
\item every constant.
\end{itemize}
The sans-serif letters $(\mathsf{n},\mathsf{j})$ index the stages of the preparatory part (P) of
the large-field operation $\mathbf{R}_k$. They are distinct from the unit letters $j$, $n$, $\nu$
and $t=L^{j-n}$ of the correlation theorem (Theorem~\ref{5.1:thm-correlation-theorem}). The
letter $\mu$ carrying the subscript of a use denotes the maps of part~\textup{(3)}.

\smallskip
\noindent\textup{(2)} \emph{Objects and stored domains.} An \emph{object} $o$ is any family which
the construction in the proof of the correlation theorem creates, stores and later reads. The
objects are of four kinds.
\begin{itemize}
\item[(u)] A \emph{small-field unit}. This kind comprises:
\begin{itemize}
\item the families of the first three entries of the list of stored families in the proof of the
correlation theorem;
\item the pointed $E$-families with tails and the families $D^{(j)}$ and $R^{(j)}$;
\item the families $f_{\square}$ and $\zeta_k G_y$ of the fifth entry of that list, which are
independent of the source variables $z$;
\item the small terms $T$ produced in the proof of the correlation theorem, stored on the domains
$\mathfrak{U}_T$.
\end{itemize}
Such a unit has the \emph{unit slot} $\sigma(o)=(j,X;\mathfrak{L})$. Here $j$ is the scale and
$X$ the localisation domain of $o$, and the index $\mathfrak{L}$ of the stored domain on which
$o$ is stored is part of the slot.
\item[(s)] An \emph{object of a stage of} (P) of $\mathbf{R}_k$. This kind comprises the terms
$\mathbb{C}^{(\mathsf{n}+1)}_k$ of the stage $\mathsf{n}$ over $\hat X$ at the complex source $w$
(not the frozen weight of part~\textup{(1)}), and the source of the stage. Such an object
has the \emph{stage slot} $(k;\mathsf{n},\hat X)$.
\item[(r)] An object of the fourth entry of that list other than those of the kinds
\textup{(u)} and \textup{(s)}. This kind comprises the terms of the operation $\mathbf{B}$
created at the sites (T) and (L) of Definition~\ref{5.2:def-sites-units}, the arranged terms
$\mathcal{C}_{\mathrm{arr}}$, and the terms of one further family of the fourth entry, fixed in
the proof of the correlation theorem. Such an object has the slot $(n,Y)$.
\item[(b)] A step class seeded by the comparison at the rim.
\end{itemize}
For an object $o$ and $\mathsf{m}\in\{\mathsf{A},\mathsf{B}\}$, the \emph{stored domain}
$\mathfrak{D}^{\mathsf{m}}(o)$ is the class of the expansion $\mathsf{m}$ alone on which that
expansion stores $o$. It is given kind by kind as follows.
\begin{itemize}
\item[(u)] $\mathfrak{D}^{\mathsf{m}}(o)$ is the domain on which, at the time of the read, the
property (f1) of $o$ and the bound of $o$ are asserted in the proof of the correlation theorem.
The property (f1) is stated on $\mathcal{U}^{\flat}_j(X)$ or on $\mathcal{U}_j(X)$, and for
the small terms $T$ on $\mathfrak{U}_T$; the stored domain at the time of the read is as follows.
\begin{itemize}
\item While $o$ is fresh, that is, when it is read by the parts (P) and (L) of $\mathbf{R}_k$,
the stored domain is the output tube, at the radii $\hat\alpha_{k-1}$, of the statement which
produces the fresh units.
\item From the next step on, the stored domain is $\mathcal{U}^{\flat}_j(X)$, the domain on which
the proof of the correlation theorem states (f1).
\item For the small terms $T$ it is $\mathfrak{U}_T$, the domain denoted $\mathfrak{D}_v$ in the
proof of the correlation theorem; the index $v$ there is a unit index, not the read object of
part~\textup{(3)}.
\end{itemize}
At age zero, that is, at the step which creates $o$, the norm $\|\cdot\|_j$ in which item
\textup{(iii)} of the lemma bounding the stored small-field units in the proof of the correlation
theorem states the bound of $o$ is the fresh size which the list of stored families in that proof
already carries at $j=k$.
\item[(s)] $\mathfrak{D}^{\mathsf{m}}(o)=\mathfrak{S}^{\mathsf{m}}(k;\mathsf{n},\hat X)$. This is
the source of stage $\mathsf{n}$ over $\hat X$ which is read by the arrangement step. It is
contained in $\mathfrak{K}_P$ and is taken together with the conditions on its triples and the
margin conditions attached to it by the corollary of the construction which states the margins
of the sources of the stages of (P).
\item[(r)] The stored domain is the domain $\tilde U^{(n+1)c,\Sigma}_k(X)$ together with the
inherited domain of the inheritance theorem. For the arranged terms $\mathcal{C}_{\mathrm{arr}}$
it is $\tilde U^{\mathrm{arr}}$.
\item[(b)] The stored domains are the level-$k$ tubes.
\end{itemize}
Consider the statement that the formula producing $o$ is defined on a set containing
$\mathfrak{D}^{\mathsf{m}}(o)$, so that the reads of $o$ are defined on all of
$\mathfrak{D}^{\mathsf{m}}(o)$. This statement is not part of the definition of the stored domain:
it is a separate hypothesis on production domains of the comparison, used by statement. For a
unit of kind \textup{(u)} it holds with equality at age zero. At later ages it is supplied by the
proof of the correlation theorem. The \emph{product domain} of $o$ is
\begin{equation}\label{5.5:eq-paired-expansions-1}
\mathfrak{D}^{\times}(o):=\mathfrak{D}^{\mathsf{B}}(o)\times\mathfrak{D}^{\mathsf{A}}(o).
\end{equation}
Its points are written $p=(P^{\mathsf{B}},P^{\mathsf{A}})$. The diagonal base point is
\begin{equation}\label{5.5:eq-paired-expansions-2}
\mathbf{1}:=\bigl((1,0),(1,0)\bigr).
\end{equation}
The two-expansion difference of a family $F$ is
\begin{equation}\label{5.5:eq-paired-expansions-3}
\delta F(p):=F^{\mathsf{B}}(P^{\mathsf{B}})-F^{\mathsf{A}}(P^{\mathsf{A}}).
\end{equation}

\smallskip
\noindent\textup{(3)} \emph{Uses.} Let $o$ and $v$ be objects such that $v\prec o$ in the order
of production of the construction; this order is well-founded. A \emph{use}
$u\colon\sigma(o)\to\sigma(v)$ (``$o$ reads $v$'') is a pair of maps, one for each
$\mathsf{m}\in\{\mathsf{A},\mathsf{B}\}$,
\begin{equation}\label{5.5:eq-paired-expansions-4}
\begin{split}
\mu_u^{\mathsf{m}}\colon\mathfrak{P}_u\times\mathfrak{D}^{\mathsf{m}}(o)\longrightarrow{}&
\{\text{configurations of the expansion }\mathsf{m}\\
&\quad\text{over the domain of }v\},
\end{split}
\end{equation}
subject to the following requirements.
\begin{itemize}
\item The maps are independent of $z$.
\item They are given by the same construction, performed in each of the two expansions.
\item They share one parameter domain $\mathfrak{P}_u$, the product of a polydisc and of the real
supports of the common integration variables.
\item The radii of the polydisc are numbers of the common scaffold of part~\textup{(1)}.
\end{itemize}
At the uses at the sites (T) and (L), the response parameter $\mathsf{t}$, a coordinate of the
polydisc factor of $\mathfrak{P}_u$, ranges over a response disc taken separately for each object
$v$ that $o$ reads, namely
\begin{equation}\label{5.5:eq-paired-expansions-a}
|\mathsf{t}|<\rho_A(v),
\end{equation}
where $\rho_A$ is the radius in the notation of the response lemma in the proof of the
correlation theorem, namely
\begin{equation}\label{5.5:eq-paired-expansions-5}
\rho_A(v)=
\begin{cases}
\mathsf{w}/\eta_p, & \text{first three, fifth entries; pointed families},\\
\mathsf{w}/\eta_p^{1/2}, & \text{fourth entry},\\
\mathsf{w}_L/(\eta_p\,\mathsf{f}_v)^{1/2}, & \text{pointed pieces at (L)}.
\end{cases}
\end{equation}
The first radius is used at the sites (T) and (P) for the units of the first three entries and of
the fifth entry of the list of stored families and for the pointed families. The second radius is
the one of the half-power response bound and is used for every unit of the fourth entry. The third
is the one of the pointed version of that bound and is used at (L) for every pointed piece of
items \textup{(c)} and \textup{(d)} of the classification of the pieces at the site (L) in the
proof of the correlation theorem.
The cube parameters $\sigma(\Delta)$ among the coordinates of $\mathfrak{P}_u$ (not to be
confused with the slots $\sigma(o)$) satisfy $|\sigma(\Delta)|\le e^{\kappa_1}$. Two remarks on
\eqref{5.5:eq-paired-expansions-5} apply:
\begin{itemize}
\item at (L) the radius $\mathsf{w}/\eta_p$ is not used for the pointed pieces of the third case;
\item the radius of the disc \eqref{5.5:eq-paired-expansions-a} is $\rho_A(v)$, not
$\mathsf{w}_{\mathrm{site}}/\eta_p$.
\end{itemize}
The pair of maps of the use is
\begin{equation}\label{5.5:eq-paired-expansions-6}
\mu_u^{\times}:=(\mu_u^{\mathsf{B}},\mu_u^{\mathsf{A}}).
\end{equation}
The statement that $\mu_u^{\mathsf{m}}$ takes its values in $\mathfrak{D}^{\mathsf{m}}(v)$ is not
part of the definition of a use. It is a hypothesis stated separately for each use, among the
hypotheses of the comparison.
\end{definition}

\begin{definition}[The production list of uses and its dilations]\label{5.5:def-production-list}
Throughout, the two machines $\mathsf{A}$ and $\mathsf{B}$, the objects $o$, their slots $\sigma(o)$,
their stored domains $\mathfrak{D}^{\mathsf{m}}(o)$ and the uses $u$ with their maps
$\mu_u^{\mathsf{m}}$ ($\mathsf{m}\in\{\mathsf{A},\mathsf{B}\}$) and their pairs $\mu_u^{\times}$ are
those of Notation~\ref{5.5:not-paired-expansions}; the list $(\mu1)$--$(\mu4)$ of uses fixed together
with the response maps, and the parameter domains $\mathfrak{P}_u$, are those of
Definition~\ref{5.2:def-sites-units}. A slot is a
unit slot $(j,X;\mathfrak{L})$, a stage slot $(k;\mathsf{n},\hat X)$, an off-chart target
$(i,\Delta;\circ)$, or the slot of a boundary term, carried by a term of $\mathbf{B}$. The sites (P),
(T), (L) of Ba{\l}aban's large-field operation $\mathbf{R}_k$ are
those of Definition~\ref{5.2:def-sites-units}. We write $\mathbf{1}$ for the diagonal base point
$((1,0),(1,0))$.

The \emph{production list} is
\begin{equation}\label{5.5:eq-production-list-a}
\begin{split}
\mathfrak{M}^{\mathrm{pr}}
:=\bigl\{u:\ &u \text{ a use by which the production formula}\\
&\text{of an object } o \text{ reads another object } v\bigr\},
\end{split}
\end{equation}
and each use is typed with the slot of the reader as its
parent, that is,
\begin{equation}\label{5.5:eq-production-list-1}
\mu_u^{\mathsf{m}}\colon \sigma(o)\longrightarrow \sigma(v)\qquad(\mathsf{m}\in\{\mathsf{A},\mathsf{B}\}),
\end{equation}
with $\sigma(o)$ the slot of the object whose production performs the read; the pair $\mu_u^{\times}$ is
typed in the same way. Explicitly,
$\mathfrak{M}^{\mathrm{pr}}$ consists of the uses listed in (a)--(d) below.

\smallskip
\noindent(a) \emph{Uses at the site \textup{(P)} of $\mathbf{R}_k$.} The parent is a stage slot. The
entries are the following.
\begin{itemize}
\item $(\mathrm{M}1^{M})$\enspace The chart map $\Phi_t$ of the transfer of the small-field analysis of
\cite[\S3]{B-CMP109} to the cover $\{\square\}$ (printed without proof in \cite[p.~276]{B-CMP119};
cited as printed), applied at all cubes $\square$ of the cover at
once, on the polydisc
\begin{equation*}
|t_\square|\le \mathsf{w}_0\, r^{\sharp}_\square \qquad\text{for every cube } \square ,
\end{equation*}
where $t_\square$ is the parameter of $\Phi_t$ at the cube $\square$ and $r^{\sharp}_\square$ is the
enlarged radius of the cube $\square$
(Definition~\ref{5.1:def-analyticity-spaces}). Its entry condition is condition (k1) of the
paired-expansion theorem, taken at this entry; it is supplied by the clause of the localised remainder
corollary for the site (P).
\item $(\mathrm{M}1)$\enspace The chart map $\Phi^{\chi}$ on the polydisc of condition (C2) of the
statement fixing the chart map $\Phi^{\chi}$. This is the one-machine use
that enters condition (C5) of that statement; its one-machine bound is clause $(\mathrm{m}1)$ of
hypothesis (H2) of the paired-expansion theorem, where it is assumed. No two-machine estimate in the
proof of the correlation theorem
(Theorem~\ref{5.1:thm-correlation-theorem}) reads it, and it is not an entry of the read list
$\mathfrak{M}^{\mathrm{rd}}$, the smaller list of uses read by the two-machine estimates; it is an entry
of $\mathfrak{M}^{\mathrm{pr}}$.
\item $(\mathrm{M}2)$\enspace The box read, that is, the pair $(\mathfrak{b}'^{\mathsf{B}}_j,
\mathfrak{b}'^{\mathsf{A}}_j)$ restricted to the domain $X_v$ of the unit $v$, at the box-read units $v$
of Definition~\ref{5.2:def-sites-units}. Its parameter domain is
\begin{equation*}
(b,\sigma)\times[0,1]\times\mathbb{D}_\square .
\end{equation*}
In this domain $b$ is the box-read datum, $[0,1]$ is the range of one further real parameter and
$\mathbb{D}_\square$ is the disc of the cube $\square$.
Here, and in the response maps $\mathsf{R}_p(\mathsf{t},\sigma)$ of item~(b), the letter $\sigma$ is a
parameter of the read, not the slot map $\sigma(\cdot)$.
The Wilson pieces and the counterterm pieces are box-read on $\mathfrak{b}'_j$, by clause (a2) of the
statement in which the box read $\mathfrak{b}'_j$ is fixed.
\item $(\mathrm{M}3)$\enspace The own-level read, that is, the pair $\mathcal{W}^{\times}_T(t,\tau)$ at
the own-read units of Definition~\ref{5.2:def-sites-units}, near and far, with
\begin{equation*}
|\tau|\le \mathsf{w}_0\, r_T(\square),
\end{equation*}
where $r_T(\square)$ is the radius of the own-level read at the cube $\square$.
\item $(\mathrm{M}4)$\enspace The two moves of the remainder theorem at the two cutoffs and of its
corollary, together with the stored remainder
\begin{equation}\label{5.5:eq-production-list-2}
\int dt\;\partial_{\mathsf{t}}\big|_{0}\,\bigl[T(\mathsf{P}^{1})-T(\mathsf{P}^{0})\bigr],
\end{equation}
where $\mathsf{P}^{0}$ and $\mathsf{P}^{1}$ are the two points of the box pair, and $T(\cdot)$ is as in
the remainder theorem at the two cutoffs.
\item $(\mathrm{M}5)$\enspace The dilations at (P): (a) the comb dilations through box reads, taken at
the weighted per-unit scaffold majorants $\mathfrak{N}^{P}(v;c_W)$; (b) the linear dilations from
$\mathbf{1}$, contracted. A dilation through a read, form (ch3) of item~(d) below, does not
occur at (P).
\item $(\mathrm{M}6)$\enspace Restriction, from a stage slot to a stage slot and from a stage slot to
the slot of a terminal object of the expansion.
\item $(\mathrm{M}8)$\enspace The comparison move $\mathsf{M}^{\perp}$: the comb-dilated remainder read,
from the stage slot to the off-chart target $(i,\Delta;\circ)$, through the off-chart versions
$(\mathrm{M}2)^{\circ}$ and $(\mathrm{M}4)^{\circ}$ of the box read and of the two moves, for every
cell-unit $v=(c,i,\Delta)$ of the class to which the comparison move applies, namely the block
cell-pieces and their companions at every index $\nu$. Its
parameter domain is $\mathfrak{P}_{\perp}$, with the
dilation radius
\begin{equation}\label{5.5:eq-production-list-3}
R^{\perp}_v := \bigl(32\,\kappa^{*}_v\bigr)^{-1},
\end{equation}
where $\kappa^{*}_v$ is the scaffold majorant of $v$. Its dilation in the parameter $\mathsf{s}$ is the
form (ch4) of item~(d). Its one-machine bound is clause $(\mathrm{m}^{\perp})^{\circ}$ of
hypothesis (H3) of the paired-expansion theorem, that is, clause (iv) of the kernel lemma of the
comparison move together with clauses (i) and (iii) of the difference lemma for
$\mathsf{M}^{\perp}$. Its entry condition, condition (k1) of the paired-expansion theorem taken at this
entry, holds by the construction of $\mathsf{M}^{\perp}$ once $\mathsf{M}^{\perp}$ is listed, by the
corollary of the paired-expansion theorem.
\end{itemize}
The seventh entry $(\mathrm{M}7)$ of the list is formed by the response-map reads of item~(b) below
together with their dilations (ch1)--(ch3) of item~(d).

\smallskip
\noindent(b) \emph{Uses at the sites \textup{(T)} and \textup{(L)}.} The parent is the slot of the
reader. The entry is the read $(\mu1)$, in the numbering $(\mu1)$--$(\mu4)$ of the list of uses fixed
together with the response maps (Definition~\ref{5.2:def-sites-units}): a response map
$\mathsf{R}_p(\mathsf{t},\sigma)$ of the site
(T), or of the site (L), restricted to $S_v$. At (L) the response map is the chain of
\cite[(1.51), (1.54), (1.57), (1.58)]{B-CMP122b} taken together with the layer representations
\cite[(1.61)--(1.63)]{B-CMP122b} and with the cut of the analysis at the site (L); the rims within this
cut are characterised in form (ch3) of item~(d). The
parameters run over the response disc of the
response maps. The read is typed as in \eqref{5.5:eq-production-list-1}, $\sigma(o)$ being the slot of
the object whose production performs the read: $\sigma(o)$ is a unit slot $(j,X;\mathfrak{L})$ when $o$
is a small-field unit created at the site (T) of $\mathbf{R}_k$, and $\sigma(o)$ is the slot of a
boundary term when $o$ is a term of $\mathbf{B}$ created at (L) or at (T). In particular every unit slot
is also the parent of the reads performed by its own production.

\smallskip
\noindent(c) \emph{Restriction.} The entry $(\mu4)$ is the restriction $(\mathrm{M}6)$ of item~(a):
from a stage slot to a stage slot, and from a stage slot to the slot of a terminal object. It is the
same entry as $(\mathrm{M}6)$, listed here in its role as the identification between the stored
domain of a stage and the source of the read performed by the next stage or by a terminal object. The
list contains no restriction from a unit slot to a unit slot.

\smallskip
\noindent(d) \emph{Dilations.} Dilations enter in four forms.
\begin{itemize}
\item[(ch1)] The comb dilation $(\mu2)$, through reads of the reader's own boxes, taken at the
(weighted) scaffold majorants. At (P) this is $(\mathrm{M}5)$(a).
\item[(ch2)] The linear dilation $(\mu3)$ from $\mathbf{1}$, contracted, at the radius
$\|\mathsf{h}_b\|_{\star}$, the a priori majorant, in the norm $\|\cdot\|_{\star}$, of the datum
$\mathsf{h}_b$ of the lemma that supplies this dilation. At (P) this is
$(\mathrm{M}5)$(b).
\item[(ch3)] The dilation $(\mu3)$ through a read, and only in the following two cases: on the
reader's own boxes, at the radius
\begin{equation*}
C_{\mathrm{hsup}}\,|B^{\mathrm{ax}}|_{\star},
\end{equation*}
the a priori majorant, in the norm $\|\cdot\|_{\star}$, of the field pair read on those boxes, with the
constant $C_{\mathrm{hsup}}$ and the field $B^{\mathrm{ax}}$ of the one-machine bound for own-box reads;
and at the site (L), at rims within the cut of item~(b), that is, rims whose diameter, measured in units
of the lattice of level $n$, $n$ the level of the rim, is at most $2M$, the diameter bound of the cubes
of \cite[(1.69), p.~191]{B-CMP122a}, which enters as a clause of the stored domain at (L) and is used
here as stated there. At rims not within that cut there is no
dilation, by clause (ii) of the lemma on dilations at the rims. The form (ch3) has no entry at (P).
\item[(ch4)] The dilation of $(\mathrm{M}8)$: the dilation of the comparison move $\mathsf{M}^{\perp}$
in the parameter $\mathsf{s}$ on $\mathfrak{P}_{\perp}$, of radius
\begin{equation*}
|\mathsf{s}|\le R^{\perp}_v=\bigl(32\,\kappa^{*}_v\bigr)^{-1},
\end{equation*}
for every cell-unit $v=(c,i,\Delta)$ of the class to which the comparison move applies, namely the
block cell-pieces and their companions at every index $\nu$. This form occurs at (P) only,
where it is the entry $(\mathrm{M}8)$, with
a stage slot as parent and the off-chart target $(i,\Delta;\circ)$ as target.
\end{itemize}
For the forms (ch1)--(ch3) the radii are the a priori majorants of the read, never its actual size;
the radius of (ch4) is $R^{\perp}_v$ of \eqref{5.5:eq-production-list-3}. The rim case of (ch3)
depends on the hypothesis of a second lemma at the rims, distinct from the lemma on dilations at the
rims cited above; that hypothesis is carried as a named hypothesis of the statement in which the rim
case is used, and it is not derived in the proof of the correlation theorem.

\smallskip
\noindent(e) \emph{The list as a whole, and the read list.} Thus $\mathfrak{M}^{\mathrm{pr}}$ is the list
of uses fixed together with the response maps, with its entry at the site (P) written out as in
item~(a), including the comparison move $(\mathrm{M}8)=\mathsf{M}^{\perp}$, and with every entry
carrying the slot of the reader as its parent. The read list is contained in it,
\begin{equation}\label{5.5:eq-production-list-4}
\mathfrak{M}^{\mathrm{rd}}\subset\mathfrak{M}^{\mathrm{pr}},
\end{equation}
because the entries of $\mathfrak{M}^{\mathrm{rd}}$ are the same maps with the parameter domains fixed
together with the response maps, and the part of $\mathfrak{M}^{\mathrm{rd}}$ formed by the reads of
response maps and their dilations is the union, with the typing of item~(b), of the entries of
items~(b) and~(d) other than the form (ch4), which enters instead as the entry $(\mathrm{M}8)$ of
item~(a). The entries of $\mathfrak{M}^{\mathrm{rd}}$ whose parent is a unit slot, namely the entries
of items~(b) and~(d) in which $\sigma(o)$ is a unit slot, are treated in the statement of this section
on reads with a unit slot as parent.
\end{definition}

\begin{definition}[Seeds, the production closure and the difference sizes]
\label{5.5:def-production-closure-bis}
We work with the slots, the aligned labels, the backgrounds $U^{\mathsf{B}}_k$, $U^{\mathsf{A}}_k$
and their common variables, the pair points and the uses of
Notation~\ref{5.5:not-paired-expansions},
and with the production list $\mathfrak{M}^{\mathrm{pr}}$ of
Definition~\ref{5.5:def-production-list}, display~\eqref{5.5:eq-production-list-a}.
Write $\mathbf{1}$ for the diagonal base point $((1,0),(1,0))$, and for a pair point $p$ with
components $P^{\mathsf{B}}$, $P^{\mathsf{A}}$ and a family $F$ put
$\delta F(X;p) := F^{\mathsf{B}}(X;P^{\mathsf{B}}) - F^{\mathsf{A}}(X;P^{\mathsf{A}})$, the
remaining arguments $y$, $w$, $u$ of $F$ being suppressed from the notation; the letter $u$ in
this role is an argument of $F$ and is distinct from the uses $u$ of $\mathfrak{M}^{\mathrm{pr}}$.

\emph{Seeds.} The \emph{$J$-side seeds} are the point $\mathbf{1}$ at every slot and
\begin{itemize}
\item[(s1)] for every aligned label and for the whole lattice, the pairs of backgrounds
$(U^{\mathsf{B}}_k, U^{\mathsf{A}}_k)$ at equal values of the common variables, complex data in the
discs of hypothesis (H2) included.
\end{itemize}
The \emph{$U$-side seeds} are the $J$-side seeds together with
\begin{itemize}
\item[(s2)] the competitor seeds on the region $X^{+}$ of the proof of the uniqueness theorem;
they enter only the lemma on $X^{+}$ and its corollary in that proof.
\end{itemize}
Together, with the notion of a seed \emph{at} a slot,
\begin{equation}\label{5.5:eq-production-closure-a-bis}
\begin{aligned}
\text{$J$-side seeds} &:= \{\mathbf{1} \text{ at every slot}\}
  \cup \bigl\{(U^{\mathsf{B}}_k, U^{\mathsf{A}}_k)
  \text{ at equal values of the common variables,}\\
&\qquad\quad \text{for every aligned label and the whole lattice,}\\
&\qquad\quad \text{complex data in the discs of (H2) included}\bigr\},\\
\text{$U$-side seeds} &:= \text{$J$-side seeds} \cup \{\text{competitor seeds on } X^{+}\},\\
\text{seed at a slot } \sigma &:= \text{a seed restricted to the region of } \sigma.
\end{aligned}
\end{equation}

\emph{Production closure.} Let $\bullet \in \{J, U\}$. The production closure is the least
slot-indexed family $\sigma \mapsto \mathrm{cl}^{\mathrm{pr}}_{\bullet}(\sigma)$ of sets of pair
points such that, for every slot $\sigma$, the inclusion on the first two lines of
\eqref{5.5:eq-production-closure-b} holds; the last two lines of the same display define the
difference sizes. The set
$\mathrm{cl}^{\mathrm{pr}}_{\bullet}(\sigma)$ is also written
$\mathrm{cl}^{\mathrm{pr}}_{\bullet}(\mathfrak{M}^{\mathrm{rd}})(\sigma)$, the read list
$\mathfrak{M}^{\mathrm{rd}}$ being displayed as argument:
\begin{equation}\label{5.5:eq-production-closure-b}
\begin{aligned}
&\mathrm{cl}^{\mathrm{pr}}_{\bullet}(\sigma) \supset
  \{\text{$\bullet$-side seeds at } \sigma\} \cup \{\mathbf{1}\}\\
&\qquad \cup \bigl\{\mu^{\times}_u(\pi; p) : u\colon \sigma' \to \sigma
  \text{ a use of } \mathfrak{M}^{\mathrm{pr}},\
  \pi \in \mathfrak{P}_u,\ p \in \mathrm{cl}^{\mathrm{pr}}_{\bullet}(\sigma')\bigr\},\\[3pt]
&\mathsf{S}^{\mathrm{cl}_{\bullet}}_j[F] := \sup_X e^{\kappa d_j(X)}
  \sup_{p \in \mathrm{cl}^{\mathrm{pr}}_{\bullet}(j,X;\mathfrak{L}_F)}\ \sup_{y,w,u}
  \bigl|\delta F(X;p)\bigr|,\\
&\mathsf{a}^{\mathrm{cl}_{\bullet}}_j := \mathsf{S}^{\mathrm{cl}_{\bullet}}_j[\hat A].
\end{aligned}
\end{equation}
Thus $\mathrm{cl}^{\mathrm{pr}}_{\bullet}$ is the least family closed under the uses
\eqref{5.5:eq-production-list-a} and containing the seeds \eqref{5.5:eq-production-closure-a-bis}.
Equivalently, every point of $\mathrm{cl}^{\mathrm{pr}}_{\bullet}(\sigma)$ is a seed or
$\mathbf{1}$ followed by a finite chain of uses of $\mathfrak{M}^{\mathrm{pr}}$, each use
applied where its one-machine maps $\mu^{\mathsf{A}}_u$, $\mu^{\mathsf{B}}_u$ are defined. Clause~(ii)
of the theorem on the one-machine maps along the production list, later in this section, shows
that this proviso is always met.

The unadorned symbol $\mathrm{cl}^{\mathrm{pr}}$ stands for either closure, and a statement
written with it holds for both. In the two-machine recursion theorem, rows
I, II, III and V are read with $p$ ranging over $\mathrm{cl}^{\mathrm{pr}}_{J}$ and row IV with
$p$ ranging over $\mathrm{cl}^{\mathrm{pr}}_{U}$.

\emph{Difference sizes.} In \eqref{5.5:eq-production-closure-b}, $\kappa$ is the decay
constant and $d_j(X)$ Ba{\l}aban's size of the domain $X$ at scale $j$. The quantity
$\mathsf{S}^{\mathrm{cl}_{\bullet}}_j[F]$ is the \emph{difference size} of the family $F$ at
level $j$ over the closure. Here $\mathfrak{L}_F$ is the label of $F$, and the inner supremum
runs over the remaining arguments $y$, $w$, $u$ of $\delta F$. Finally,
$\mathsf{a}^{\mathrm{cl}_{\bullet}}_j$ is the difference size of the pointed Wilson family
$\hat A = A(h_y,\cdot)$.

For every slot $\sigma$, every family $F$ and every level $j$,
\begin{equation}\label{5.5:eq-production-closure-1-bis}
\mathrm{cl}^{\mathrm{pr}}_{J}(\sigma) \subset \mathrm{cl}^{\mathrm{pr}}_{U}(\sigma),
\qquad
\mathsf{S}^{\mathrm{cl}_J}_j[F] \le \mathsf{S}^{\mathrm{cl}_U}_j[F].
\end{equation}
\end{definition}

\begin{proof}
\emph{The inclusion of the closures.} Consider the family
$\sigma \mapsto \mathrm{cl}^{\mathrm{pr}}_{U}(\sigma)$. At each slot it contains $\mathbf{1}$.
By \eqref{5.5:eq-production-closure-a-bis}, every $J$-side seed is a $U$-side seed, so this family
also contains every $J$-side seed at $\sigma$. It is moreover closed under every use
$u\colon\sigma'\to\sigma$ of $\mathfrak{M}^{\mathrm{pr}}$ with $\pi\in\mathfrak{P}_u$, because the
third set in the union on the first two lines of \eqref{5.5:eq-production-closure-b}, read for
$\bullet = U$, is exactly this closedness.

The family $\mathrm{cl}^{\mathrm{pr}}_{U}$ therefore satisfies every defining condition of
$\mathrm{cl}^{\mathrm{pr}}_{J}$. Since $\mathrm{cl}^{\mathrm{pr}}_{J}$ is the least family
satisfying these conditions, $\mathrm{cl}^{\mathrm{pr}}_{J}(\sigma) \subset
\mathrm{cl}^{\mathrm{pr}}_{U}(\sigma)$ for every slot $\sigma$.

\emph{The inequality of the sizes.} Fix $j$ and $X$. By the inclusion just proved at the slot
$(j,X;\mathfrak{L}_F)$, the supremum over $p$ in \eqref{5.5:eq-production-closure-b} for
$\bullet = J$ runs over a subset of the set used for $\bullet = U$. The quantity
$\sup_{y,w,u}|\delta F(X;p)|$ under the supremum over $p$ and the weight $e^{\kappa d_j(X)}$ do
not depend on $\bullet$.
The inner supremum is therefore no larger for $J$ than for $U$, and taking $\sup_X$ preserves
the inequality. This gives $\mathsf{S}^{\mathrm{cl}_J}_j[F] \le \mathsf{S}^{\mathrm{cl}_U}_j[F]$;
applied to $F = \hat A$ it gives
$\mathsf{a}^{\mathrm{cl}_J}_j \le \mathsf{a}^{\mathrm{cl}_U}_j$.
\end{proof}

\begin{definition}[Standing hypotheses on the single expansions]\label{5.5:def-standing-hypotheses}
Throughout, the machine index $\mathsf{m}$ runs over $\{\mathsf{A},\mathsf{B}\}$. We use the objects $o$,
their slots $\sigma(o)$ and their stored domains $\mathfrak{D}^{\mathsf{m}}(o)$ of
Notation~\ref{5.5:not-paired-expansions}, and the production list $\mathfrak{M}^{\mathrm{pr}}$ of
Definition~\ref{5.5:def-production-list} (see \eqref{5.5:eq-production-list-a}). There each use
$u\colon\sigma(o)\to\sigma(v)$ carries a parameter domain $\mathfrak{P}_u$ and, for each machine, a
map $\mu_u^{\mathsf{m}}$ of the pair $(\pi,P^{\mathsf{m}})$, $\pi\in\mathfrak{P}_u$. The production
closures $\mathrm{cl}^{\mathrm{pr}}_{J}$, $\mathrm{cl}^{\mathrm{pr}}_{U}$ are those of
Definition~\ref{5.5:def-production-closure-bis}, and $\mathbf{1}$ is the diagonal base point. For an object
$o$ and a machine $\mathsf{m}$, the \emph{production source} $\mathfrak{D}^{\mathsf{m}}_{\mathrm{src}}(o)$
is the one-machine set on which $o$ is defined and on which the production formula of $o$ reads other
objects.

Hypotheses \textup{(S2)} and \textup{(S4)} are statements about one machine at a time. The
\emph{standing hypotheses on the single expansions} are \textup{(S1)}--\textup{(S4)}.
Hypotheses \textup{(S1)}--\textup{(S3)} are the hypotheses of the correspondence theorem for the
production closure; \textup{(S4)} is the input of the filing theorem at the rims and of the
correspondence lemma, read at the uses whose parent is a unit slot, a slot of the summation of
boundary terms at finer zones, or a step slot; none of them is a hypothesis of a new kind.

\smallskip
\noindent\textup{(S1)} The single expansions satisfy the first standing condition of the construction
and the filing predicate.

\smallskip
\noindent\textup{(S2)} (Containment in the one-machine systems.)
\begin{itemize}
\item[(sd)] The seeds, restricted seeds included, lie in their stored domains:
  \begin{itemize}
  \item[(s1)] on the $J$-side, the seeds and the base point $\mathbf{1}$;
  \item[(s2)] on the $U$-side only, the seeds of the closure $\mathrm{cl}^{\mathrm{pr}}_{U}$. This
    containment is a statement of the filing theorem at the rims, and it follows from three results:
    the statement fixing the domains of the $U$-side seeds; clause (c2) of the containment statement
    in the proof of the correlation theorem (Theorem~\ref{5.1:thm-correlation-theorem}), taken at the
    enlarged region; and the statement on the seed sets. The first of these three results is used by
    no statement about $\mathrm{cl}^{\mathrm{pr}}_{J}$.
  \end{itemize}
  No hypothesis of this kind is imposed at the seeds of $\mathrm{cl}^{\mathrm{pr}}_{J}$ beyond (s1).
\item[(m1)] The move of the production list \eqref{5.5:eq-production-list-a} that reads through the
  chart maps $\Phi_t$, $\Phi^{\chi}$ maps stored domains into stored domains. In its
  unlocalised form this is the containment clause of the statement on these chart maps; in its
  localised form it is the localisation corollary for that move, which is the
  row at the stage $(P)$ of the filing corollary and its claim.
\item[(m6)] Restrictions map stored domains into stored domains, for restrictions from a stage slot to
  a stage slot and from a stage slot to a unit slot. This is the one-machine restriction statement of
  the paired expansions; at the site $(T)$ it follows from the filing corollary and its claim.
\end{itemize}

\smallskip
\noindent\textup{(S3)} (Memberships at the units of the stage $(P)$.) The values produced at the stage
$(P)$ lie in the stored domains of their targets. This holds in each of the following cases.
\begin{itemize}
\item[(m2)] The members of the Wilson family, by clauses (a4), (b)(i) and (c2) of the statement on
  the Wilson family; the Wilson pieces, by its clause (a2). These rest also on an admitted input of
  Definition~\ref{5.1:def-admitted-inputs}.
\item[(m3)] The own-level members, by clauses (a) and (a$'$) of the statement on own-level members,
  which rest on the filing corollary with its claim and on the window statement; the far members, by
  clause (iv) of the statement on far members.
\item[(m4)] The members given by a theorem on members of the stage $(P)$ and its corollary.
\item[(m5)] The values taken at weighted majorants, or starting from the base point $\mathbf{1}$,
  contracted. These are given by clauses (c2) and $(\ell 0)$ of the containment statement in the proof
  of the correlation theorem (Theorem~\ref{5.1:thm-correlation-theorem}) together with its lemma on
  containment at weighted majorants; no such value passes through a read at the stage $(P)$.
\item[$(\mathrm{m}^{\perp})$] For the comparison move $\mathsf{M}^{\perp}$ of the production list
  \eqref{5.5:eq-production-list-a}, and for its dilation, two properties hold, by clause (iv) of the
  comparison-move lemma and clauses (i), (iii) of the comparison-move difference lemma:
  \begin{itemize}
  \item the dilated family is holomorphic on $\mathfrak{P}_{\perp}$ and entire in the dilation
    parameter $s$;
  \item on the image of the dilated family, the entry bound $\mathfrak{n}^{\perp}_v$ of the off-chart
    target slot $(i,\Delta;\circ)$ holds.
  \end{itemize}
  The smallness condition on which these two properties rest is among the conditions that fix
  $\gamma_J$.
\end{itemize}

\smallskip
\noindent\textup{(S4)} (Production containment at the stored domain of the reading object.) This
hypothesis concerns the uses of $\mathfrak{M}^{\mathrm{pr}}$ whose parent is a unit slot, a slot of the
summation of boundary terms at finer zones, or a step slot, one machine
at a time. It is the conjunction of \textup{(S4a)} and \textup{(S4b)}: for every object $o$, every
machine $\mathsf{m}$ and every such use $u\colon\sigma(o)\to\sigma(v)$, taken at the site $(T)$, at the
site $(L)$ or among the chart uses,
\begin{equation}\label{5.5:eq-standing-hypotheses-a}
\begin{gathered}
\textup{(S4b)}\qquad \mathfrak{D}^{\mathsf{m}}(o)\subset\mathfrak{D}^{\mathsf{m}}_{\mathrm{src}}(o),\\
\textup{(S4a)}\qquad
\begin{aligned}[t]
&\mu_u^{\mathsf{m}}\ \text{is defined and holomorphic on}\ \mathfrak{P}_u\times
  \mathfrak{D}^{\mathsf{m}}_{\mathrm{src}}(o),\\
&\mu_u^{\mathsf{m}}\bigl(\mathfrak{P}_u\times\mathfrak{D}^{\mathsf{m}}_{\mathrm{src}}(o)\bigr)
  \subset\mathfrak{D}^{\mathsf{m}}(v).
\end{aligned}
\end{gathered}
\end{equation}
Here ``holomorphic'' means holomorphic in $P^{\mathsf{m}}$ and in complex $\pi$, and continuous in real
$\pi$. By \textup{(S4a)}, every later read into $v$ lands in the stored domain of $v$. Under
\textup{(S4a)} and \textup{(S4b)}, the map $\mu_u^{\mathsf{m}}$ is defined and holomorphic on
$\mathfrak{P}_u\times\mathfrak{D}^{\mathsf{m}}(o)$ with values in $\mathfrak{D}^{\mathsf{m}}(v)$; this
one-clause form of \textup{(S4)} is proved below.

The production source in \textup{(S4b)} is the following set, for each kind of object in the
classification of Notation~\ref{5.5:not-paired-expansions}.
\begin{itemize}
\item For the objects of kind (u): the set of radii $\hat\alpha_k:=(1+2\beta_{\mathrm{s}})\alpha_{\cdot}(g_k)$,
  $\cdot\in\{0,1\}$, on which the families of the correlation theorem are defined; here
  $\beta_{\mathrm{s}}$ is the enlargement fraction of Lemma~\ref{5.2:lem-enlargement-fractions} and
  $\alpha_{\cdot}(g)$ are the radii of Notation~\ref{5.1:not-coupling-functions}.
\item For the objects of kinds (s) and (r): the source of the reads of the arrangement theorem. It
  consists of the rooms given by the filing corollary and its claim, the inherited domain of clause
  (a) of the inheritance statement, and the space $\tilde U^{\mathrm{arr}}$ on which the arranged
  terms $\mathcal{C}_{\mathrm{arr}}$ are defined by the arrangement theorem. For these kinds
  \textup{(S4b)} is the statement of the filing corollary with its claim, of the inheritance
  statement and of the arrangement theorem; it rests also on an admitted input of
  Definition~\ref{5.1:def-admitted-inputs}.
\item For the objects of kind (b): the tubes.
\end{itemize}
For the objects of kind (u), the sets given by the radii $\hat\alpha_k$ contain Ba{\l}aban's spaces
enlarged by the factor $1+\beta_{\mathrm{pr}}$ after the operation $\mathbf{T}$. These in turn contain
the stored spaces $\mathcal{U}^{\flat}$. The first containment rests on the chain
\begin{equation}\label{5.5:eq-standing-hypotheses-1}
(1+\beta_{\mathrm{pr}})\,\alpha_{\cdot,k+1}
\le(1+\beta_{\mathrm{s}})\,\alpha_{\cdot,k+1}
\le(1+2\beta_{\mathrm{s}})\,\alpha_{\cdot}(g_k)=\hat\alpha_k .
\end{equation}
Its first inequality is $\beta_{\mathrm{pr}}\le\beta_{\mathrm{s}}$ for the enlargement fractions of
Lemma~\ref{5.2:lem-enlargement-fractions}. Its second is the second clause of the lemma comparing the
enlarged radii at consecutive levels. By \eqref{5.5:eq-standing-hypotheses-1}, the outputs on
$\hat\alpha_k$ restrict to Ba{\l}aban's spaces enlarged after $\mathbf{T}$. Every containment into
Ba{\l}aban's enlarged spaces that is used in this chapter therefore holds a fortiori into the sets of
radii $\hat\alpha_k$.
\end{definition}

\begin{proof}[Proof of the one-clause form of \textup{(S4)}]
We claim that, under \textup{(S4)}, for every object $o$, machine $\mathsf{m}$ and use
$u\colon\sigma(o)\to\sigma(v)$ as in \eqref{5.5:eq-standing-hypotheses-a}, the map $\mu_u^{\mathsf{m}}$
is defined and holomorphic on $\mathfrak{P}_u\times\mathfrak{D}^{\mathsf{m}}(o)$ and takes values in
$\mathfrak{D}^{\mathsf{m}}(v)$.

By \textup{(S4b)}, $\mathfrak{D}^{\mathsf{m}}(o)\subset\mathfrak{D}^{\mathsf{m}}_{\mathrm{src}}(o)$.
Hence
\begin{equation*}
\mathfrak{P}_u\times\mathfrak{D}^{\mathsf{m}}(o)\subset\mathfrak{P}_u\times
\mathfrak{D}^{\mathsf{m}}_{\mathrm{src}}(o).
\end{equation*}
By the first line of \textup{(S4a)}, $\mu_u^{\mathsf{m}}$ is defined on the larger set and is
holomorphic there, in the sense of the definition. Holomorphy in $P^{\mathsf{m}}$ and in complex $\pi$,
and continuity in real $\pi$, are preserved under restriction to a subset. So $\mu_u^{\mathsf{m}}$ is
defined and holomorphic on $\mathfrak{P}_u\times\mathfrak{D}^{\mathsf{m}}(o)$.

For the values we use the monotonicity of images and the second line of \textup{(S4a)}:
\begin{equation*}
\mu_u^{\mathsf{m}}\bigl(\mathfrak{P}_u\times\mathfrak{D}^{\mathsf{m}}(o)\bigr)
\subset\mu_u^{\mathsf{m}}\bigl(\mathfrak{P}_u\times\mathfrak{D}^{\mathsf{m}}_{\mathrm{src}}(o)\bigr)
\subset\mathfrak{D}^{\mathsf{m}}(v).
\end{equation*}
This proves the claim. The claim is the form of \textup{(S4)} used in the proof of the lemma on the
one-machine maps of the paired expansions.
\end{proof}

\begin{remark}[Sources of the hypotheses \textup{(S4)}]
The containment \textup{(S4a)} is given,
statement by statement, by clauses (c1)--(c3) and $(\ell 0)$ of the containment statement in the proof
of the correlation theorem (Theorem~\ref{5.1:thm-correlation-theorem}). By these clauses the following
configurations, in Ba{\l}aban's notation, are admissible backgrounds at scale $n$:
\begin{itemize}
\item $U_k$;
\item $U_k(e^{iB'}V^{(k)})$, with its deformations in the parameters $(t,t_{\square},\mathsf{s})$
  on Ba{\l}aban's circles;
\item $U^{(n')}_k$ of \cite{B-CMP122a};
\item the chain \cite[(1.51)--(1.58)]{B-CMP122b};
\item the $(\mathbf{U},\mathbf{J})$-extensions of all of these on Ba{\l}aban's tubes.
\end{itemize}
For the chains at the stage $(P)$ and at the site $(L)$, the values lie in the space
$\tilde U''^{c}_k(Y_1,\tilde\alpha_0,\tilde\alpha_1)$, in the notation of the proof of the correlation
theorem. That space is contained in the analyticity domain of the term $E(X)$ there, and hence in the
analyticity domains of all terms. For the
$\mathbb{C}$-units of the summation of boundary terms at finer zones, the containment is the filing
corollary with its claim; for the arranged terms $\mathcal{C}_{\mathrm{arr}}$ it is the arrangement
theorem.

At the site $(T)$, the response map $\mathsf{R}_p$ restricted to $S_v$ lands in admissible
backgrounds; this is the content of a lemma on the containment clauses in the proof of the correlation
theorem. Containment in the stored domain of the target holds one machine at a time, as in the filing
theorem at the rims. For the move of
\eqref{5.5:eq-production-list-a} listed immediately before $\mathsf{M}^{\perp}$, the containment of its
values in the stored domain of its target is given by that move's row in the correspondence statement
for the read list.

For the chart uses the containment is given as follows:
\begin{itemize}
\item[(ch1)] by the same clauses, taken at the weighted a priori majorants of the read;
\item[(ch2)] by the lemma on containment at weighted majorants;
\item[(ch3)] at the own boxes of the read, by the statement supplying the own-box clause of the
  correspondence lemma; at the rims of the site $(L)$, by the cube clause of
  \cite[(1.69), p.~191]{B-CMP122a} together with the lemma on containment at weighted majorants;
\item[(ch4)] this case has no instance whose parent is a unit slot, a slot of the summation of
  boundary terms at finer zones, or a step slot. The parent of
  $\mathsf{M}^{\perp}$ is a stage slot, and that case is the content of $(\mathrm{m}^{\perp})$ in
  \textup{(S3)}.
\end{itemize}

In Ba{\l}aban's construction, the new expressions are defined on slightly larger spaces
\cite[p.~262]{B-CMP119}. The spaces carrying the terms form a descending sequence
\cite[p.~191]{B-CMP122a}. The space $\tilde U^{c}_k(Y_1,\alpha_0,\alpha_1)$ is contained in the
analyticity domains of all terms \cite[(1.70), p.~375]{B-CMP122b}; see also
\cite[(1.12), p.~262]{B-CMP109}.

\emph{Why (S4) is posed at the stored domain of the reading object.} In the correlation theorem a
unit is stored on the domain on which its analyticity property (f1) holds. For a new family, (f1) is
obtained from two facts:
\begin{itemize}
\item the production reads of that family are contained in the older stored domains;
\item the relevant series converge uniformly, by the convergence statements in the proof of the
  correlation theorem.
\end{itemize}
Thus \textup{(S4a)} is the input of that induction at the production of $o$, read on
$\mathfrak{D}^{\mathsf{m}}_{\mathrm{src}}(o)\supset\mathfrak{D}^{\mathsf{m}}(o)$. For objects of kind
(u), condition \textup{(S4b)} is the a fortiori containment following
\eqref{5.5:eq-standing-hypotheses-1}; for the other kinds it is the statement of the results named
for them after \eqref{5.5:eq-standing-hypotheses-a}. The
norm $\|\cdot\|_j$ of the correlation theorem is a supremum over the stored domain. It is therefore the
norm read in the third clause of the lemma on the one-machine maps of the paired expansions.
\end{remark}

\begin{theorem}[Minimality, containment and finiteness of the production closure]
\label{5.5:thm-closure-theorem}
Assume the standing hypotheses (H1)--(H4) of Definition~\ref{5.5:def-standing-hypotheses},
displayed in \eqref{5.5:eq-standing-hypotheses-a}, including their parts (H4a) and (H4b).
For the statements concerning $\mathrm{cl}^{\mathrm{pr}}_{U}$ assume in addition the $U$-side
clause (s2) of (H2)(s).

\emph{Notation.} We write $\mathrm{cl}^{\mathrm{pr}}_{\bullet}$, $\bullet\in\{J,U\}$, for the
production closures of Definition~\ref{5.5:def-production-closure-bis}
(see \eqref{5.5:eq-production-closure-b}), and $\mathrm{cl}^{\mathrm{pr}}$ for either of them.
$\mathfrak{M}^{\mathrm{pr}}$ is the production list of Definition~\ref{5.5:def-production-list},
and $\mathfrak{M}^{\mathrm{rd}}\subset\mathfrak{M}^{\mathrm{pr}}$ is the read list.

We write $\mathfrak{M}^{\mathrm{pr}}_{P}$, $\mathfrak{M}^{\mathrm{pr}}_{T}$ and
$\mathfrak{M}^{\mathrm{pr}}_{L}$ for the entries of \eqref{5.5:eq-production-list-a} at the sites
$(P)$, $(T)$ and $(L)$, $\mathfrak{M}^{\mathrm{pr}}_{\mathrm{ch}}$ for the entry whose containment
clauses are (ch1)--(ch4), and $\mathfrak{M}^{\mathrm{pr}}_{\mathrm{io}}$ for the entry
$\mathsf{M}_6$; the moves $\mathsf{M}_1$, $\mathsf{M}_1^{M}$, $\mathsf{M}_2,\dots,\mathsf{M}_6$ and
the comparison move $\mathsf{M}^{\perp}$ are those of Definition~\ref{5.5:def-production-list}.
A point of a slot is a pair $p=(P^{\mathsf{B}},P^{\mathsf{A}})$. For an object $o$ we put
$\mathfrak{D}^{\times}(o)=\mathfrak{D}^{\mathsf{B}}(o)\times\mathfrak{D}^{\mathsf{A}}(o)$.
For a family $F$ at the unit slot $(j,X;\mathfrak{L}_F)$ we put
$\delta F(X;p):=F^{\mathsf{B}}(X;P^{\mathsf{B}})-F^{\mathsf{A}}(X;P^{\mathsf{A}})$.
Then the following hold.
\begin{enumerate}
\item[(i)] \emph{Correspondence and minimality.} $\mathrm{cl}^{\mathrm{pr}}_{\bullet}$ contains
$\mathbf{1}$ and the $\bullet$-side seeds. It is closed under every use of
$\mathfrak{M}^{\mathrm{pr}}$; that is, it has the properties (k0)--(k2) of a correspondence class
of the two-machine correspondence scheme, taken for the list $\mathfrak{M}^{\mathrm{pr}}$. It is the
least family with these properties.

Split the read class alike: $\mathcal{K}^{\mathrm{rd}}_J$ is generated from $\mathbf{1}$ and the
seeds (s1), and $\mathcal{K}^{\mathrm{rd}}$ from $\mathbf{1}$ and the seeds (s1), (s2). Then at
every slot $\sigma$ at which the read class is defined
\begin{equation}\label{5.5:eq-closure-theorem-1}
\mathrm{cl}^{\mathrm{pr}}_{J}(\sigma)\supseteq\mathcal{K}^{\mathrm{rd}}_J(\sigma),
\qquad
\mathrm{cl}^{\mathrm{pr}}_{U}(\sigma)\supseteq\mathcal{K}^{\mathrm{rd}}(\sigma),
\end{equation}
and hence, for every level $j$ and every family $F$,
\begin{equation}\label{5.5:eq-closure-theorem-2}
\mathsf{S}^{\mathrm{rd}_J}_j[F]\le\mathsf{S}^{\mathrm{cl}_J}_j[F],
\qquad
\mathsf{S}^{\mathrm{rd}}_j[F]\le\mathsf{S}^{\mathrm{cl}_U}_j[F].
\end{equation}

Let $\mathcal{K}^{\natural}[\mathfrak{M}]$ denote the class generated by finite chains of uses of
a list $\mathfrak{M}$ from $\mathbf{1}$ and the seeds of the two-machine correspondence scheme. Then
$\mathrm{cl}^{\mathrm{pr}}_{U}=\mathcal{K}^{\natural}[\mathfrak{M}^{\mathrm{pr}}]$.

The class $\mathcal{K}^{\natural}$ of the two-machine correspondence scheme is formed for that
scheme's list of moves $(\mu1)$--$(\mu4)$, $(T)$, $(L)$, with its $(P)$-entry written out as
the uses of $\mathfrak{M}^{\mathrm{pr}}_{P}$. Under the interpretation of its moves given in the
proof of (i) it is identified with $\mathrm{cl}^{\mathrm{pr}}_{U}$; this is an identification, not
an equality of the defining formulas. The response discs of the
two-machine correspondence scheme are written with the radii $\mathsf{w}_{\mathrm{site}}$ and
$\eta_p$, whereas those of \eqref{5.5:eq-paired-expansions-a} in
Notation~\ref{5.5:not-paired-expansions} are written with
$\mathbb{D}_{\square}$, $\mathsf{w}_0r_T(\square)$ and $\rho_A(v)$. The move $(\mu1)$ at the
site $(P)$ is interpreted as the maps by which the production formulas are defined.

\item[(ii)] \emph{Containment through unit slots.} For every object $o$, every point of
$\mathrm{cl}^{\mathrm{pr}}(\sigma(o))$ lies in $\mathfrak{D}^{\times}(o)$. This holds whether
$\sigma(o)$ is a stage slot $(k;\mathsf{n},\hat X)$, a unit slot $(j,X;\mathfrak{L})$, a slot of
the summation of boundary terms at finer zones, or a step slot. Moreover, every use along every
chain is applied inside its domain of definition.

In particular, consider a chain that starts at a seed and enters the class of a unit through a
later use into that unit's slot. If it continues with that unit's own production read out of the
slot, the chain never leaves the product of the stored domains.

\item[(iii)] \emph{Finiteness.} For a level-$j$ family $F$ with
$\|F^{\mathsf{m}}\|_j\le\mathfrak{n}_F^{\mathsf{m}}$, $\mathsf{m}\in\{\mathsf{A},\mathsf{B}\}$,
\begin{equation}\label{5.5:eq-closure-theorem-a}
\mathsf{S}^{\mathrm{cl}}_j[F]\le\mathfrak{n}_F^{\mathsf{B}}+\mathfrak{n}_F^{\mathsf{A}}
\le 2\mathfrak{n}_F<\infty .
\end{equation}

\item[(iv)] \emph{The two-machine difference principle at every use.} Let
$u:\sigma(o)\to\sigma(v)$ be a use in $\mathfrak{M}^{\mathrm{pr}}$, let
$p\in\mathrm{cl}^{\mathrm{pr}}(\sigma(o))$, and let $f^{\mathsf{m}}$ be holomorphic on
$\mathfrak{D}^{\mathsf{m}}(v)$. Put
\begin{equation}\label{5.5:eq-closure-theorem-3}
\delta(\pi):=f^{\mathsf{B}}\bigl(\mu_u^{\mathsf{B}}(\pi;P^{\mathsf{B}})\bigr)
-f^{\mathsf{A}}\bigl(\mu_u^{\mathsf{A}}(\pi;P^{\mathsf{A}})\bigr).
\end{equation}
Then $\delta$ is holomorphic on $\mathfrak{P}_u$, and
\begin{equation}\label{5.5:eq-closure-theorem-4}
|\delta|\le\sup_{\mathrm{cl}^{\mathrm{pr}}(\sigma(v))}|f^{\mathsf{B}}-f^{\mathsf{A}}| .
\end{equation}
Every Cauchy inequality and every inequality (a)--(c) of Lemma~\ref{5.2:lem-disc-bounds} carries
over to $\delta$, with
the majorant $\mathfrak{N}$ replaced by $\sup|\delta|$. The same holds for every identity
$\partial^ig^{\mathsf{m}}(\pi_0)=0$, $\pi_0\in\mathfrak{P}_u$, satisfied in each machine by
$g^{\mathsf{m}}:=f^{\mathsf{m}}\circ\mu_u^{\mathsf{m}}(\cdot;P^{\mathsf{m}})$: then
$\partial^i\delta(\pi_0)=0$.

\item[(v)] \emph{Diagonal test.} Suppose $\mathsf{B}=\mathsf{A}$, that is, the same lattice, the
same operators and the same stored terms. Suppose also that every seed is diagonal, that is, of the
form $(P,P)$. Then every $\mathrm{cl}^{\mathrm{pr}}(\sigma)$ is diagonal and $\delta F\equiv0$ on
it. Hence $\mathsf{S}^{\mathrm{cl}}_j[F]=0$ for all $F$.

\item[(vi)] \emph{The theorem on the read class is kept.} Clauses (i)--(iii) and (v)(b)--(e) of
the theorem on the read class, and its corollary, hold with
$(\mathcal{K}^{\mathrm{rd}},\mathsf{S}^{\mathrm{rd}})$ replaced by
$(\mathrm{cl}^{\mathrm{pr}},\mathsf{S}^{\mathrm{cl}})$, for either closure.

In particular the following hold on $\mathrm{cl}^{\mathrm{pr}}$:
\begin{itemize}
\item the two properties (RM) and (MC)$^{M}$ of field pair systems required in that theorem;
\item property (k1) for the read list $\mathfrak{M}^{\mathrm{rd}}$;
\item property (k1) at the own production maps $\mu^{\mathrm{own}}$;
\item property (k1) at $\mathsf{M}_6$, which is the closure property of the class;
\item property (k1) at $\mathsf{M}_1^{M}$, at $\mathsf{M}_2\oplus\mathsf{M}_3\oplus\mathsf{M}_4$,
at $\mathsf{M}_3$ and at $\mathsf{M}_5$.
\end{itemize}
Property (k1) at $\mathsf{M}^{\perp}$ also holds, by (i), since
$\mathsf{M}^{\perp}\in\mathfrak{M}^{\mathrm{rd}}\subset\mathfrak{M}^{\mathrm{pr}}$.

Clause (e) of the statement on the pointed Wilson family $\hat A$, clause (f) of the statement on
the own production maps $\mu^{\mathrm{own}}$, and the fifth clause of the corollary on the
difference sizes are statements about
$(\mathrm{cl}^{\mathrm{pr}}_J,\mathsf{S}^{\mathrm{cl}_J},\mathsf{a}^{\mathrm{cl}_J})$. In the three
statements that require a field pair system with the properties (RM) and (MC)$^{M}$, the
classes of $\mathrm{cl}^{\mathrm{pr}}$ may be taken for that system.

In place of clause (iv) of the theorem on the read class, $\mathrm{cl}^{\mathrm{pr}}$ satisfies
the following:
\begin{itemize}
\item out of a unit slot, or a slot of the summation of boundary terms at finer zones, the uses
are the object's own production reads, those of
$\mathfrak{M}^{\mathrm{pr}}_{T}\cup\mathfrak{M}^{\mathrm{pr}}_{L}\cup\mathfrak{M}^{\mathrm{pr}}_{\mathrm{ch}}$;
\item the use $\mathsf{M}_6$, which constitutes $\mathfrak{M}^{\mathrm{pr}}_{\mathrm{io}}$, has
stage parents;
\item likewise $\mathsf{M}^{\perp}$ has stage parents; it leads from a stage slot to an off-chart
target $(i,\Delta;\circ)$.
\end{itemize}
\end{enumerate}
\end{theorem}

\begin{proof}
\emph{(i).} By Definition~\ref{5.5:def-production-closure-bis}, a point of
$\mathrm{cl}^{\mathrm{pr}}_{\bullet}(\sigma)$ is one of two kinds. Either it is $\mathbf{1}$ or a
$\bullet$-side seed at $\sigma$, or it is obtained from $\mathbf{1}$ or a $\bullet$-side seed at
some slot $\sigma'$ by a finite chain of uses of $\mathfrak{M}^{\mathrm{pr}}$ ending at $\sigma$,
each with its parameter in the corresponding domain $\mathfrak{P}_u$.
In particular $\mathrm{cl}^{\mathrm{pr}}_{\bullet}$ contains $\mathbf{1}$ and the $\bullet$-side
seeds.

Closure: let $p\in\mathrm{cl}^{\mathrm{pr}}_{\bullet}(\sigma')$, let $u:\sigma'\to\sigma$ be a use
of $\mathfrak{M}^{\mathrm{pr}}$, and let $\pi\in\mathfrak{P}_u$. The chain producing $p$, extended
by the one use $u$ with parameter $\pi$, is again a chain. Hence
$\mu_u^{\times}(\pi;p)\in\mathrm{cl}^{\mathrm{pr}}_{\bullet}(\sigma)$.

Minimality: let $\mathcal{G}$ be any family that contains $\mathbf{1}$ and the $\bullet$-side seeds
and is closed under the uses of \eqref{5.5:eq-production-list-a} as in
\eqref{5.5:eq-production-closure-a-bis}. We show by induction on the length of the producing chain that
every point of $\mathrm{cl}^{\mathrm{pr}}_{\bullet}$ lies in $\mathcal{G}$.
\begin{itemize}
\item A chain of length zero ends at $\mathbf{1}$ or at a seed, and these lie in $\mathcal{G}$.
\item Suppose every point produced by a chain of length $\ell$ lies in $\mathcal{G}$. A point
produced by a chain of length $\ell+1$ is $\mu_u^{\times}(\pi;p)$, where $p$ is produced by a chain
of length $\ell$. Then $p\in\mathcal{G}$, and closedness of $\mathcal{G}$ gives
$\mu_u^{\times}(\pi;p)\in\mathcal{G}$.
\end{itemize}
So $\mathrm{cl}^{\mathrm{pr}}_{\bullet}\subseteq\mathcal{G}$, and $\mathrm{cl}^{\mathrm{pr}}_{\bullet}$
is the least such family.

For \eqref{5.5:eq-closure-theorem-1}, note the following facts.
\begin{itemize}
\item $\mathfrak{M}^{\mathrm{rd}}\subset\mathfrak{M}^{\mathrm{pr}}$, with the same maps
$\mu_u^{\times}$ and the parameter domains $\mathfrak{P}_u$ of
\eqref{5.5:eq-paired-expansions-a}.
\item The seeds of $\mathcal{K}^{\mathrm{rd}}_J$ are those of $\mathrm{cl}^{\mathrm{pr}}_J$, namely
$\mathbf{1}$ and (s1).
\item The seeds of $\mathcal{K}^{\mathrm{rd}}$ are those of $\mathrm{cl}^{\mathrm{pr}}_U$, namely
$\mathbf{1}$, (s1) and (s2).
\end{itemize}
By the definition of the read class, a point of $\mathcal{K}^{\mathrm{rd}}_J(\sigma)$ (resp.\
$\mathcal{K}^{\mathrm{rd}}(\sigma)$) is a seed or $\mathbf{1}$ followed by a finite chain of uses
of $\mathfrak{M}^{\mathrm{rd}}$. The same chain is a chain of uses of $\mathfrak{M}^{\mathrm{pr}}$
from a seed of $\mathrm{cl}^{\mathrm{pr}}_J$ (resp.\ $\mathrm{cl}^{\mathrm{pr}}_U$). This proves
\eqref{5.5:eq-closure-theorem-1}. The difference sizes in \eqref{5.5:eq-closure-theorem-2} are
suprema over these classes, and a supremum over a subset is not larger. This gives
\eqref{5.5:eq-closure-theorem-2}.

The class $\mathcal{K}^{\natural}[\mathfrak{M}^{\mathrm{pr}}]$ is generated from $\mathbf{1}$ and
the seeds of the two-machine correspondence scheme, which are the seeds (s1) and (s2). It is by
construction $\mathrm{cl}^{\mathrm{pr}}_U$: both consist of the points reached from these seeds or
from $\mathbf{1}$ by finite chains of uses of $\mathfrak{M}^{\mathrm{pr}}$.

The class $\mathcal{K}^{\natural}$ of the two-machine correspondence scheme consists of every point
obtained as a seed followed by a finite chain of moves of that scheme's list. We interpret these
moves as uses of $\mathfrak{M}^{\mathrm{pr}}$ as follows.
\begin{itemize}
\item The $(P)$-entry of $(\mu1)$ is interpreted, as the production formulas define it, as a
composite of three pieces: the seed of the label at the stage slot
(see \eqref{5.5:eq-production-closure-a-bis}); then the chain of moves $\mathsf{M}_1$,
$\mathsf{M}_1^{M}$; then one of $\mathsf{M}_2$, $\mathsf{M}_3$, $\mathsf{M}_4$. This composite is
a chain of $\mathfrak{M}^{\mathrm{pr}}$.
\item The moves $(\mu2)$ and $(\mu3)$ are interpreted as uses of
$\mathfrak{M}^{\mathrm{pr}}_{\mathrm{ch}}$, which contains $\mathsf{M}_5$.
\item The move $(\mu4)$ is the use $\mathsf{M}_6$ of $\mathfrak{M}^{\mathrm{pr}}_{\mathrm{io}}$.
\item The entries $(T)$ and $(L)$ are the uses of $\mathfrak{M}^{\mathrm{pr}}_T$ and
$\mathfrak{M}^{\mathrm{pr}}_L$ on the response discs of \eqref{5.5:eq-paired-expansions-a}.
\end{itemize}
Under this interpretation every chain producing a point of $\mathcal{K}^{\natural}$ is a chain
producing a point of $\mathrm{cl}^{\mathrm{pr}}_U$. This gives
$\mathcal{K}^{\natural}\subseteq\mathrm{cl}^{\mathrm{pr}}_U$. The equality stated in (i) is the
identification under this interpretation and not an equality of the defining formulas.

One may instead interpret $(\mu1)$ at $(P)$ as a single direct map. In that case it yields unit
classes contained in those of $\mathrm{cl}^{\mathrm{pr}}$. The statements of this section that
follow the present theorem are formulated on $\mathrm{cl}^{\mathrm{pr}}$, whose definition does not
depend on which interpretation is chosen.

\emph{(ii).} We argue by induction on the length of the chain producing the point, one machine
$\mathsf{m}$ at a time.

\emph{Base.} A chain of length zero ends at $\mathbf{1}$ or at a seed. Clause (s) of (H2) places
$\mathbf{1}$ and the seeds, restricted seeds included, in the stored domains. On the $U$-side this
is clause (s2),
which rests on the inputs listed with it in Definition~\ref{5.5:def-standing-hypotheses}.

\emph{Step.} Let $p\in\mathrm{cl}^{\mathrm{pr}}(\sigma(o))$ with
$P^{\mathsf{m}}\in\mathfrak{D}^{\mathsf{m}}(o)$; this is the induction hypothesis. Let
$u:\sigma(o)\to\sigma(v)$ be in $\mathfrak{M}^{\mathrm{pr}}$ and $\pi\in\mathfrak{P}_u$. We show
$\mu_u^{\mathsf{m}}(\pi;P^{\mathsf{m}})\in\mathfrak{D}^{\mathsf{m}}(v)$.

\emph{Case 1: $\sigma(o)$ is a stage slot and $u\in\mathfrak{M}^{\mathrm{pr}}_{P}$.} This is the
containment step for the read class, taken move by move.
\begin{itemize}
\item For $\mathsf{M}_1$, clause (m1) of (H2) applies.
\item For $\mathsf{M}_1^{M}$, the localisation corollary for that move applies.
\item For $\mathsf{M}_2$, clause (m2) of (H3) applies. The Wilson and counterterm pieces read by
$\mathsf{M}_2$ are entire and obey the a priori bound on these pieces.
\item For $\mathsf{M}_3$, clause (m3) applies.
\item For $\mathsf{M}_4$, clause (m4) applies.
\item For $\mathsf{M}_5$, clause (m5) applies, in two parts. Part (a): the data $B_i^{\mathsf{m}}$
entering the use obey the weighted majorant in each machine. Hence, with the coefficients
$\sigma_i$ and the threshold $\alpha_{3,j}$ of clause (m5) of
Definition~\ref{5.5:def-standing-hypotheses}, pointwise
\begin{equation}\label{5.5:eq-closure-theorem-5}
\Bigl|\sum_i\sigma_iB_i^{\mathsf{m}}\Bigr|
\le\sum_i|\sigma_i|\,\sup|B_i^{\mathsf{m}}|_{*}<\alpha_{3,j},
\end{equation}
after which the containment and smallness clauses that turn this pointwise bound into containment
in the stored domain apply. Part (b) is given by the lemma on the ball of backgrounds in
the proof of the correlation theorem (Theorem~\ref{5.1:thm-correlation-theorem}).
\item For $\mathsf{M}^{\perp}$, clause (m$\perp$)$^{\circ}$ of (H3) applies, that is, (ch4) at
$(P)$. It consists of two statements on the comparison move, clause (iv) of the lemma on the
dilated comparison family and clauses (i), (iii) of the lemma on its entry bound. First, the family
dilated in $\mathsf{s}$ is holomorphic on $\mathfrak{P}_{\perp}$ and entire in $\mathsf{s}$.
Second, on $|\mathsf{s}|\le R^{\perp}_v$ it obeys the entry bound $\mathfrak{n}^{\perp}_v$ of the
slot $(i,\Delta;\circ)$. Property (k1) at $\mathsf{M}^{\perp}$ holds by construction.
\item For $\mathsf{M}_6$, the use of $\mathfrak{M}^{\mathrm{pr}}_{\mathrm{io}}$, clause (m6) of (H2)
applies.
\end{itemize}
In each case $\mu_u^{\mathsf{m}}(\pi;P^{\mathsf{m}})\in\mathfrak{D}^{\mathsf{m}}(v)$.

\emph{Case 2:
$u\in\mathfrak{M}^{\mathrm{pr}}_{T}\cup\mathfrak{M}^{\mathrm{pr}}_{L}\cup\mathfrak{M}^{\mathrm{pr}}_{\mathrm{ch}}$.}
This case includes every use out of a unit slot, a slot of the summation of boundary terms at finer
zones, or a step slot. Hypothesis (H4b) gives
$P^{\mathsf{m}}\in\mathfrak{D}^{\mathsf{m}}(o)\subset\mathfrak{D}^{\mathsf{m}}_{\mathrm{src}}(o)$,
and $u$ is run on $\mathfrak{D}^{\mathsf{m}}_{\mathrm{src}}(o)$. Hypothesis (H4a) for $u$ is
required on the whole of $\mathfrak{P}_u\times\mathfrak{D}^{\mathsf{m}}_{\mathrm{src}}(o)$, and it
gives $\mu_u^{\mathsf{m}}(\pi;P^{\mathsf{m}})\in\mathfrak{D}^{\mathsf{m}}(v)$.

For the uses of $\mathfrak{M}^{\mathrm{pr}}_{\mathrm{ch}}$ at the sites $(T)$ and $(L)$, (H4a)
consists of four clauses.
\begin{itemize}
\item In (ch1), the data entering the use obey the weighted a priori majorants pointwise; then the
same containment and smallness clauses apply, as in part (a) of (m5).
\item (ch2) is given by the lemma on the ball of backgrounds.
\item In (ch3), on the object's own boxes the statement supplying the box reads is used together
with the lemma on the ball of backgrounds. At the site $(L)$, at rims within the cut of the theorem
on the read class, one works in the cube gauge of \cite[(1.69)]{B-CMP122a} together with the ball
of that lemma.
\item (ch4) has no instance out of a unit slot, a slot of the boundary summation or a step slot. It
is the dilation of $\mathsf{M}^{\perp}$, whose parent is a stage slot, and it was treated in
Case~1.
\end{itemize}

Hence no chain leaves the stored domains. Moreover, each map $\mu_u^{\mathsf{m}}$ is constructed on
$\mathfrak{D}^{\mathsf{m}}_{\mathrm{src}}(o)\supseteq\mathfrak{D}^{\mathsf{m}}(o)$; the minimisers,
averages and kernels entering it are defined there, as in the constructions of the chart maps
$\Phi^{\chi}$, of the pointed Wilson family $\hat A$ and of the own production maps
$\mu^{\mathrm{own}}$, and in the proof of Theorem~\ref{5.1:thm-correlation-theorem}. So each use is
applied inside its domain of definition.

\emph{Passage through a unit slot.} Let $u'$ be a later use into $\sigma(o)$, and let
$q:=\mu_{u'}^{\times}(\pi';p')\in\mathrm{cl}^{\mathrm{pr}}(\sigma(o))$. Then
$Q^{\mathsf{m}}\in\mathfrak{D}^{\mathsf{m}}(o)$ by the step for $u'$, because the target of (H4a)
for $u'$ is the stored domain of $o$.

The object's own production read $u$ out of $\sigma(o)$ at $q$ is the step for $u$. It is
legitimate because (H4b) places all of $\mathfrak{D}^{\mathsf{m}}(o)$ inside
$\mathfrak{D}^{\mathsf{m}}_{\mathrm{src}}(o)$, where (H4a) for $u$ holds. The containment is needed
on all of $\mathfrak{D}^{\mathsf{m}}(o)$, not merely at the seeds or at the points at which $o$ is
first produced.

This is the form in which the statements required on the stored domains are formulated in the proof
of Theorem~\ref{5.1:thm-correlation-theorem}: on all admissible backgrounds. In
\cite[p.~280, Lemma~4]{B-CMP109} the condition is that $\alpha_0,\alpha_1,\alpha_2,\alpha_3$ be
sufficiently small. The objects are defined on these domains in \cite[p.~262]{B-CMP119} and
\cite[p.~375]{B-CMP122b}.

\emph{(iii).} Let $o$ be the object at the unit slot $(j,X;\mathfrak{L}_F)$ and
$p\in\mathrm{cl}^{\mathrm{pr}}(j,X;\mathfrak{L}_F)$. By (ii),
$P^{\mathsf{m}}\in\mathfrak{D}^{\mathsf{m}}(o)$ for both machines. The norm $\|\cdot\|_j$ is a
supremum over the stored domain $\mathfrak{D}^{\mathsf{m}}(o)$ with the weight
$e^{\kappa d_j(X)}$. Hence
\begin{equation}\label{5.5:eq-closure-theorem-6}
\begin{aligned}
|\delta F(X;p)|
&\le|F^{\mathsf{B}}(X;P^{\mathsf{B}})|+|F^{\mathsf{A}}(X;P^{\mathsf{A}})|\\
&\le e^{-\kappa d_j(X)}\bigl(\|F^{\mathsf{B}}\|_j+\|F^{\mathsf{A}}\|_j\bigr)
\le e^{-\kappa d_j(X)}\bigl(\mathfrak{n}_F^{\mathsf{B}}+\mathfrak{n}_F^{\mathsf{A}}\bigr)
\end{aligned}
\end{equation}
at every point over which $\mathsf{S}^{\mathrm{cl}}_j[F]$ of \eqref{5.5:eq-production-closure-b} is
formed. This gives \eqref{5.5:eq-closure-theorem-a}.

\emph{(iv).} By (ii), $\mu_u^{\mathsf{m}}(\pi;P^{\mathsf{m}})\in\mathfrak{D}^{\mathsf{m}}(v)$, where
$f^{\mathsf{m}}$ is holomorphic. So the composition in \eqref{5.5:eq-closure-theorem-3} is
legitimate. Holomorphy in $\pi$ is that of the maps in Notation~\ref{5.5:not-paired-expansions},
and it is supplied together with the containment by (H2)--(H4). Hence $\delta$ is holomorphic on
$\mathfrak{P}_u$.

Property (k1) is (i): $\mu_u^{\times}(\pi;p)\in\mathrm{cl}^{\mathrm{pr}}(\sigma(v))$. Writing
$\delta f:=f^{\mathsf{B}}-f^{\mathsf{A}}$ on pairs, we obtain
\begin{equation*}
|\delta(\pi)|=\bigl|\delta f\bigl(\mu_u^{\times}(\pi;p)\bigr)\bigr|
\le\sup_{\mathrm{cl}^{\mathrm{pr}}(\sigma(v))}|\delta f|,
\end{equation*}
which is \eqref{5.5:eq-closure-theorem-4}.

The Cauchy inequalities, the inequalities (a)--(c) of Lemma~\ref{5.2:lem-disc-bounds} and the
identities
$\partial^ig^{\mathsf{m}}(\pi_0)=0$ in each machine are linear in the function. Therefore they hold
for the difference $\delta=g^{\mathsf{B}}-g^{\mathsf{A}}$ of
$g^{\mathsf{m}}:=f^{\mathsf{m}}\circ\mu_u^{\mathsf{m}}(\cdot;P^{\mathsf{m}})$, with $\mathfrak{N}$
replaced by $\sup|\delta|$.

\emph{(v).} We argue by induction on the chain. The seeds are diagonal. At
$\mathsf{B}=\mathsf{A}$ the two components of a use are the same map on the same lattice, fed the
same $\pi$. So $\mu_u^{\times}(\pi;(P,P))=(Q,Q)$, and a diagonal point is carried to a diagonal
point. Since $F^{\mathsf{B}}=F^{\mathsf{A}}$, we get $\delta F\equiv0$ on every
$\mathrm{cl}^{\mathrm{pr}}(\sigma)$, and $\mathsf{S}^{\mathrm{cl}}_j[F]=0$.

\emph{(vi).} The statements named in (vi) use only the following four properties of the read class.
\begin{itemize}
\item Property (k1) at the moves $\mathsf{M}_1$--$\mathsf{M}_6$, and at $\mathsf{M}^{\perp}$ as
asserted in (vi). These hold on $\mathrm{cl}^{\mathrm{pr}}$ by (i).
\item The two-machine difference principle, which holds on $\mathrm{cl}^{\mathrm{pr}}$ by (iv).
\item Finiteness of the difference sizes, which holds by (iii).
\item The dictionary of clause (v)(e) of the theorem on the read class.
\end{itemize}
\end{proof}

\begin{lemma}[Pair bounds on the production closure]\label{5.5:lem-pair-bounds}
Assume the following three statements.
\begin{itemize}
\item[(a)] Hypotheses (H1)--(H4) of Definition~\ref{5.5:def-standing-hypotheses}, displayed in
\eqref{5.5:eq-standing-hypotheses-a}.
\item[(b)] The lemma stating the pair bounds $(\mathrm{A}^{\times})$, $(\mathrm{M}^{\times})$,
$(\Pi^{\times})$, $(\mathrm{I}^{\times})$, $(\partial^{\times})$, $(\mathrm{F}^{\times})$,
$(\mathrm{W}^{\times})$ for a correspondence class and the two difference sizes over it (of the list
families and of the pointed Wilson family), in the form in which they hold for the class
$\mathcal{K}^{\natural}$ and, with the same constants, for every correspondence class $\mathcal{K}$
that contains $\mathbf{1}$ and is closed under the use at each read (the uses (U-T), (U-L) and the
box-read and own-read pairs of hypothesis (H2) of Definition~\ref{5.5:def-standing-hypotheses}),
under the use (U-ch) and under the use of item (m5) of hypothesis (H2), with the difference sizes
over $\mathcal{K}$ in place of those over $\mathcal{K}^{\natural}$.
\item[(c)] Clause (e) of the lemma stating the weighted pair bound for a triple
$(\mathcal{K},\mathsf{S},\mathsf{a})$ (a correspondence class and the two difference sizes over it),
which holds whenever $(\alpha)$ every move applied in the weighted bound has one parameter domain
common to both machines $\mathsf{A}$ and $\mathsf{B}$, and $(\beta)$ $\mathcal{K}$ is closed under
every such move.
\end{itemize}
In the notation of Notation~\ref{5.5:not-paired-expansions} and
Definition~\ref{5.5:def-standing-hypotheses}, let $F^{\mathsf{m}}$, $\mathsf{m}\in\{\mathsf{A},\mathsf{B}\}$,
be level-$j$ list families with
$\|F^{\mathsf{m}}\|_j\le\mathfrak{n}$, put $\mathsf{S}:=\mathsf{S}^{\mathrm{cl}}_j[F]$ and
$\mathsf{a}:=\mathsf{a}^{\mathrm{cl}}_j$ as in \eqref{5.5:eq-production-closure-b}, let
$p\in\mathrm{cl}^{\mathrm{pr}}(\sigma)$ be admissible at scale $n$, and put $t:=L^{j-n}$. Then the seven
bounds of \textup{(b)} hold with the correspondence class replaced by the production closure
$\mathrm{cl}^{\mathrm{pr}}$ of Definition~\ref{5.5:def-production-closure-bis} and the difference sizes
replaced by $\mathsf{S}$ and $\mathsf{a}$, with no other change except the following readings.
\begin{itemize}
\item[$(\mathrm{A}^{\times})$] The bound holds, for the unit $v$ that reads $p$, on the disc
$|\mathsf{t}|<\rho_A(v)$ of
\eqref{5.5:eq-paired-expansions-a}, with the prefactor of the proof of the correlation theorem
(Theorem~\ref{5.1:thm-correlation-theorem}) taken per child, namely
$2\mathsf{w}_{\mathrm{site}}^{-1}\eta_p$; in the half-power case it is
$2\mathsf{w}_{\mathrm{site}}^{-1}\eta_p^{1/2}$, and in the pointed half-power case it is
$\tfrac{81}{64}\cdot 2\mathsf{w}_L^{-1}(\eta_p\mathsf{f}_v)^{1/2}$. These replace the single prefactor
$2\mathsf{w}_{\mathrm{site}}^{-1}\eta_p$ of \textup{(b)}.
\item[$(\mathrm{M}^{\times})$] For a read $q=\mu_u^{\times}(\pi;p)$, where $u$ is an own-box read carrying
a family of type (ch1) (the box read at the stage (P); the first-class reads at the sites (T) and
(L)), the first bound in \eqref{5.5:eq-pair-bounds-a} holds.
\item[$(\mathrm{I}^{\times})$] The second bound in \eqref{5.5:eq-pair-bounds-a} holds.
\end{itemize}
Moreover, at the stage (P) the weighted form holds: for every near unit $v$ and every
$p\in\mathrm{cl}^{\mathrm{pr}}(k;\mathsf{n}+1,\hat X)$, on the common disc, let
$\mathsf{R}^{\times}_p$ be the pair of response maps of the box read of hypothesis (H2) of
Definition~\ref{5.5:def-standing-hypotheses} at box-read units and
the pair of the own read of that hypothesis at own-read units. Then the third bound in
\eqref{5.5:eq-pair-bounds-a} holds, with $\mathsf{S}$ and $\mathsf{a}$ read as the $J$-side sizes
$\mathsf{S}^{\mathrm{cl}_J}_j[F]$ and $\mathsf{a}^{\mathrm{cl}_J}_j$; this is the weighted pair bound
\textup{(c)} with
$(\mathcal{K},\mathsf{S},\mathsf{a}):=(\mathrm{cl}^{\mathrm{pr}}_{J},\mathsf{S}^{\mathrm{cl}_J}_j[F],
\mathsf{a}^{\mathrm{cl}_J}_j)$.
The three bounds are
\begin{equation}\label{5.5:eq-pair-bounds-a}
\begin{aligned}
&|\delta F(X;q)-\delta F(X;\mathbf{1})|
\le 2\mathsf{S}\,e^{-\kappa d_j(X)}\bigl(C_M D_X t^2\rho_{n,j}\bigr)^2,\\
&|\mathfrak{i}_y^{\mathsf{B}}(\mathfrak{b}^{\mathsf{B}})-\mathfrak{i}_y^{\mathsf{A}}(\mathfrak{b}^{\mathsf{A}})|
\le C_I'(\mathsf{S}+\mathfrak{n}\mathsf{a})\,\nu^3\rho^2(1+\rho)\,t^{4+\hat\alpha},\\
&|\delta(f_v\circ\mathsf{R}^{\times}_p)-\delta f_v(\mathbf{1})|
\le w^{\vee}(y)^3\,\mathfrak{N}^{P}(v;c_W)\big|_{\mathfrak{n}\mapsto\mathsf{S}+\mathfrak{n}\mathsf{a},
\ C_I'\mapsto C_I''}.
\end{aligned}
\end{equation}
The clause of $(\mathrm{M}^{\times})$ for reads of any class, rims included, with $C_M$ replaced by
$C_\partial$, is not part of this lemma; it is the statement of the lemma giving the marginal pair
bound at rims, under the hypothesis named in that lemma.
\end{lemma}

\begin{proof}
The seven bounds of \textup{(b)} rest on the following inputs; for each we say what it requires of
the class.
\begin{itemize}
\item[(1)] The two Cauchy integrals, in the response parameter $\mathsf{t}$ and in the deformation
parameter, of the first-stage bound (Lemma~\ref{5.2:lem-first-stage-bound}), applied to
$\delta-\delta(\mathbf{1})$ on the disc fixed in the proof of the correlation theorem, taken per
child. This is the input of $(\mathrm{A}^{\times})$. It requires
$\mathsf{R}^{\times}_p(\mathfrak{P}\times\{p\})$ to lie in the class at the slot $\sigma(v)$, that is,
closure of the class under the use at the read; and it requires the clause of the principle for
pairs of machines by which identities holding in each machine separately hold for the pair.
\item[(2)] The families of type (ch1) entering through an own-box read and the families of type (ch2)
entering from $\mathbf{1}$, with $\varphi^{\mathsf{m}'}(0)=0$ in each machine $\mathsf{m}'$. These are the
inputs of the items of $(\mathrm{M}^{\times})$, $(\Pi^{\times})$ and $(\mathrm{I}^{\times})$: values,
cubic Schwarz remainders, and mixed second derivatives at $0$ along the $(\mu2)$- and
$(\mu3)$-families. They require that the class contain $\mathbf{1}$ and be closed under the use
(U-ch) and under the use of
item (m5) of hypothesis (H2) of Definition~\ref{5.5:def-standing-hypotheses}, together with the
clause of the principle for pairs of machines by which the identities valid in each machine pass
to the pair.
\item[(3)] The trivial bound $2\mathsf{S}e^{-\kappa d}$, which is a further clause of the principle for
pairs of machines.
\item[(4)] The bounds for the families $F^{\mathsf{m}}$ and for the pointed Wilson family $\hat A$
supplied by Theorem~\ref{5.1:thm-correlation-theorem} in each machine separately; these are the
inputs of $(\mathrm{F}^{\times})$ and $(\mathrm{W}^{\times})$.
\end{itemize}

Each of these inputs is available on the production closure. For (1): by
Theorem~\ref{5.5:thm-closure-theorem}\,(i), $\mathrm{cl}^{\mathrm{pr}}$ is closed under every use of
$\mathfrak{M}^{\mathrm{pr}}$, in particular under the use at the read, and this holds alike for the
uses (U-T), (U-L) and for the box-read and own-read pairs of hypothesis (H2) of
Definition~\ref{5.5:def-standing-hypotheses}; hence
$\mathsf{R}^{\times}_p(\mathfrak{P}\times\{p\})\subset\mathrm{cl}^{\mathrm{pr}}(\sigma(v))$ for
$p\in\mathrm{cl}^{\mathrm{pr}}(\sigma)$, and the two Cauchy integrals are taken on the disc
$|\mathsf{t}|<\rho_A(v)$ of \eqref{5.5:eq-paired-expansions-a}, which carries the prefactor per
child listed in the statement. For (2): by Theorem~\ref{5.5:thm-closure-theorem}\,(i),
$\mathrm{cl}^{\mathrm{pr}}$ contains $\mathbf{1}$, and by the same closure property, applied to the
use (U-ch) and to the use of item (m5)
of hypothesis (H2), the points reached through an own-box read and from $\mathbf{1}$ lie in
$\mathrm{cl}^{\mathrm{pr}}$; the per-machine identities $\varphi^{\mathsf{m}'}(0)=0$ do not depend on the
class and pass to the pair by the clause of that principle named in (2). For (3): the trivial bound
involves the class only through
the difference size, here $\mathsf{S}=\mathsf{S}^{\mathrm{cl}}_j[F]$. For (4): these are bounds in each
machine separately and involve no class. Thus $\mathrm{cl}^{\mathrm{pr}}$ contains $\mathbf{1}$ and has
every closure property that \textup{(b)} requires of a correspondence class, so \textup{(b)} applies
to $\mathcal{K}=\mathrm{cl}^{\mathrm{pr}}$ with the difference sizes $\mathsf{S}$ and $\mathsf{a}$ and
gives $(\mathrm{A}^{\times})$, $(\mathrm{M}^{\times})$,
$(\Pi^{\times})$, $(\mathrm{I}^{\times})$, $(\partial^{\times})$, $(\mathrm{F}^{\times})$,
$(\mathrm{W}^{\times})$ on $\mathrm{cl}^{\mathrm{pr}}$, among them the first two bounds of
\eqref{5.5:eq-pair-bounds-a}. In every one of the bounds the difference size enters linearly.
In particular, for a level-$j$ family $T$ entering through a read at
$q$, the first bound of \eqref{5.5:eq-pair-bounds-a} applied to $T$ gives
\begin{equation*}
|\delta T(X;q)-\delta T(X;\mathbf{1})|\le 2\mathsf{S}^{\mathrm{cl}}_j[T]\,e^{-\kappa d_j(X)}
\bigl(C_M D_X t^2\rho_{n,j}\bigr)^2,
\end{equation*}
so that the difference constant of $T$ is $2\mathsf{S}^{\mathrm{cl}}_j[T]$, linear in the difference
size.

For the weighted form we apply \textup{(c)}. Its proof is the proof of the pair bounds \textup{(b)}
with $\rho$ replaced by $\Theta\rho$ throughout, and it applies to a triple
$(\mathcal{K},\mathsf{S},\mathsf{a})$ provided $(\alpha)$ and
$(\beta)$ hold. Take
$(\mathcal{K},\mathsf{S},\mathsf{a})=(\mathrm{cl}^{\mathrm{pr}}_{J},\mathsf{S}^{\mathrm{cl}_J}_j[F],
\mathsf{a}^{\mathrm{cl}_J}_j)$. Condition $(\alpha)$ holds by \eqref{5.5:eq-paired-expansions-a} and
hypothesis (H1) of Definition~\ref{5.5:def-standing-hypotheses}: each move has one parameter domain
common to both machines, and the weights $\Theta$, the label sets $S_v$, the filing, the discs
and the majorants are label data, the same in both machines. Condition $(\beta)$ is
Theorem~\ref{5.5:thm-closure-theorem}\,(i). Hence \textup{(c)} gives the third bound of
\eqref{5.5:eq-pair-bounds-a}, with $\mathfrak{n}$ replaced by $\mathsf{S}+\mathfrak{n}\mathsf{a}$ and
$C_I'$ by $C_I''$.
\end{proof}

\begin{definition}[The difference number of the two expansions]\label{5.5:def-difference-number}
We work at Ba{\l}aban's large-field operation $\mathbf{R}_k$, at any stage of its preparatory
part (P) and at the sites (T) and (L) of Definition~\ref{5.2:def-sites-units} alike. The two
machines $\mathsf{A}$ and $\mathsf{B}$ are the expansions at the two cutoffs. A point $p$ of a
production closure has the components $P^{\mathsf{A}}$ and $P^{\mathsf{B}}$, and for a family
$F$ we write
\begin{equation*}
\delta F(p) := F^{\mathsf{B}}(P^{\mathsf{B}}) - F^{\mathsf{A}}(P^{\mathsf{A}})
\end{equation*}
for the two-machine difference. The difference sizes $\mathsf{S}^{\mathrm{cl}}_j[F]$ and
$\mathsf{a}^{\mathrm{cl}}_j$ are those of \eqref{5.5:eq-production-closure-b}, read on the
$J$-side production closure $\mathrm{cl}^{\mathrm{pr}}_J$ of
Definition~\ref{5.5:def-production-closure-bis}. The per-cell bounds used below
are those of the table of per-cell bounds in the proof of the correlation theorem
(Theorem~\ref{5.1:thm-correlation-theorem}); the \emph{small-field families} are those of the
lines $r\in\{1,2,3,\partial\}$ of that table, which carry respectively the pointed $E$-families
with their tails, the families $D^{(j)}$, the families $R^{(j)}$, and the boundary terms, together
with the small terms of the sourced preparatory step
(Corollary~\ref{5.2:cor-sourced-preparatory}). The fifth line of that table carries the block
pieces $f_\square$ and the terms $\zeta_kG_y$ used in (c) below. The per-cell scales $S_I$ and
$S_\omega$ are those of Definition~\ref{5.2:def-comparison-constant}. The letters $\nu$,
$\hat\alpha$, $\rho$, $\rho_{n,j}$, $C$, $C_I$, $C_B$, $\omega_{k,j}$, $C_\partial'$,
$\mathfrak{a}_n$ and $\kappa_0$ below are those of that table; the table writes $\rho$ without
indices in lines $1$ and $\partial$ and $\rho_{n,j}$ in line $3$.

\smallskip
\emph{(a) The off-chart difference size.} Let $v$ range over the class-$W$ cell-units of the
off-chart comparison that occur at $\mathbf{R}_k$, namely the cells $v=(c_\Delta,i,\Delta)$ of
the following, at every $\nu$:
\begin{itemize}
\item a block piece $f_{\square'}$ of the fifth line of the table;
\item a companion term $b_j(\hat w)A(h_y,\cdot)$ of the fifth line.
\end{itemize}
We write such a cell-unit as $v=(c,i,\Delta)$; its target slot is the off-chart slot
$(i,\Delta;\circ)$. Its \emph{off-chart difference size} $\mathsf{S}^{\perp}_v$ is the supremum
of $|\delta A(c,p)|$ over the points $p$ of the class $\mathrm{cl}^{\mathrm{pr}}_J(i,\Delta;\circ)$
(display \eqref{5.5:eq-difference-number-a} below), where $A(c,\cdot)$ is the pointed Wilson
family read at the coefficient function $c$ of the cell-unit. This class contains the images at
that slot of three uses of $\mathfrak{M}^{\mathrm{pr}}$: the off-chart forms of two of its moves,
and the move whose target slot is the slot $(i,\Delta;\circ)$ of $v$. It also
contains the class $\mathcal{K}^{\mathrm{rd}}_J(i,\Delta;\circ)$, by clause (i) of
Theorem~\ref{5.5:thm-closure-theorem}.

\smallskip
\emph{(b) The coefficient difference.} The number $\mathsf{b}_j$ is the modulus of the
two-machine difference of the counterterm coefficient $b_j(\hat w)$ of line $\partial$. This is
the difference size that enters the pair bound (W$^\times$) of
Lemma~\ref{5.5:lem-pair-bounds}: there the one-machine bound $b_++\tfrac{2}{3}r$ of the
coefficient is replaced by $\mathsf{b}_j+b_+\mathsf{a}^{\mathrm{cl}}_j$. Here $b_+$ and $r$ are
the constants of the one-machine bound on the counterterm coefficient in the proof of
Theorem~\ref{5.1:thm-correlation-theorem}; this constant $r$ is distinct from the index $r$ of
a line of the table.

\smallskip
\emph{(c) The $z$-free pieces.} The block pieces $f_\square$ and the terms $\zeta_kG_y$ of the
fifth line are $z$-free. Their two-machine differences are $z$-free difference pieces,
denoted by the letter $\mathsf{z}$, and they are not members of $\mathfrak{D}^J$. There is one
exception. The class-$W$ cell pieces of the block pieces $f_{\square'}$ and of the
companion terms $b_j(\hat w)A(h_y,\cdot)$, read off the chart at the points of the class at
$(i,\Delta;\circ)$, carry the
off-chart difference size $\mathsf{S}^{\perp}_v$ of (a), and this size is a member of
$\mathfrak{D}^J$. At these pieces the off-chart comparison statement is applied at the
coefficient
\begin{equation*}
c_\Delta=\zeta_k\hat\chi_{\square'}1_\Delta ,
\end{equation*}
where $1_\Delta$ is the indicator function of $\Delta$, $\hat\chi_{\square'}$ is the function
attached to the cube $\square'$ in the block piece $f_{\square'}$, and $\zeta_k$ is $z$-free. At
all other pieces the split of the pair bound (W$^\times$) is used. The difference of the
coefficient,
\begin{equation*}
\mathsf{z}_\zeta:=|\zeta^{\mathsf{B}}_k-\zeta^{\mathsf{A}}_k| ,
\end{equation*}
remains a $z$-free difference piece.

\smallskip
\emph{(d) Per-cell weights.} For each line $r\in\{1,2,3,\partial\}$ and each unit $n$-cell, we
write the last column of the table (per unit $n$-cell, summed over $j$) as
\begin{equation*}
\sum_{j\le n}c^{(r)}_{n,j}\,\mathfrak{n}^{(r)}_j .
\end{equation*}
Here $\mathfrak{n}^{(r)}_j$ is the size factor of line $r$: it is the size $E_0$ of the family
$E$ in lines $1$, $2$ and $\partial$, and $g_j^{\kappa_0}$ in line $3$. The number
$c^{(r)}_{n,j}\ge 0$ is the per-cell weight of line $r$ with the size divided out:
\begin{itemize}
\item in line $1$ it is $C_I\nu^3\rho^2(1+\rho)L^{-\hat\alpha(\nu-1)}$;
\item in line $2$ they are the weights carrying $S_\omega$, with the factor $C_B\omega_{k,j}$
inside;
\item in line $3$ it is $C\rho^2_{n,j}$;
\item in line $\partial$ it is $27dC_\partial'M^3(\nu+2)\rho^2L^{-(\nu-1)}$, and for the
counterterm entry, whose entry size in the table is $b_++r$ with $r$ the constant of (b), it
is $27d\cdot 21M\mathfrak{a}_n^2$.
\end{itemize}
These are numbers depending on $(\nu,n-j,L,M,d)$ and on the couplings of the label. Where a
weight reads the coupling of one machine, the larger of the two machines' values is taken:
\begin{equation*}
\bar\rho_{n,j}:=\max_{\mathsf{m}}\rho^{\mathsf{m}}_{n,j} .
\end{equation*}
Whether read at one machine's coupling or at the larger of the two, each weight is a number
fixed before $\gamma$.

We let $K_r$ be the per-cell constant of line $r$ at size one, computed with the constants
fixed in Definition~\ref{5.1:def-admitted-inputs}. In line $1$ the constant $C_I$ is replaced
by $C_I''$, the pair constant of the bound (I$^\times$) of Lemma~\ref{5.5:lem-pair-bounds}. The
same replacement $C_I\mapsto C_I''$ is made in the core terms of line $\partial$; this gives the
constant $K_\partial$.

\smallskip
\emph{(e) The difference data.} With the conventions (a)--(d) we put
\begin{equation}\label{5.5:eq-difference-number-a}
\begin{aligned}
\mathfrak{D}^J &:= \bigl\{\mathsf{S}^{\mathrm{cl}}_j[F] : j\le k,\\
&\qquad\quad F \text{ a small-field family of level } j\bigr\}\\
&\qquad \cup\bigl\{\mathsf{a}^{\mathrm{cl}}_j : j\le k\bigr\}
  \cup\bigl\{\mathsf{b}_j : j\le k\bigr\}
  \cup\bigl\{\mathsf{S}^{\perp}_v : v \text{ a class-}W\text{ cell-unit}\bigr\},\\
\mathsf{S}^{\perp}_v &:= \sup_{p\in\mathrm{cl}^{\mathrm{pr}}_J(i,\Delta;\circ)}|\delta A(c,p)|,
  \qquad v=(c,i,\Delta),\\
\mathsf{b}_j &:= \bigl|b^{\mathsf{B}}_j(\hat w)-b^{\mathsf{A}}_j(\hat w)\bigr|,\\
K_r &:= \sup_n\sum_{j\le n}c^{(r)}_{n,j}\\
&\qquad(\text{with } C_I \text{ replaced by } C_I'' \text{ in line } 1
  \text{ and in the core terms of line } \partial),\\
K_1 &= C_I''S_I,\\
d^{(1)}_j &= d^{(\partial)}_j := \mathsf{S}^{\mathrm{cl}}_j[E]+E_0\,\mathsf{a}^{\mathrm{cl}}_j,\\
d^{(2)}_j &:= \mathsf{S}^{\mathrm{cl}}_j[E],\qquad
d^{(3)}_j := \mathsf{S}^{\mathrm{cl}}_j[R^{(j)}].
\end{aligned}
\end{equation}
Here the small-field families are those fixed above: the families of lines $1$, $2$, $3$,
$\partial$ of the table and the small terms of
Corollary~\ref{5.2:cor-sourced-preparatory}. The number $d^{(r)}_j$ is the
\emph{difference datum of line $r$ at level $j$}. The
counterterm entry of line $\partial$ carries the datum
$\mathsf{b}_j+b_+\mathsf{a}^{\mathrm{cl}}_j$, by (W$^\times$). The datum of lines $1$ and
$\partial$ comes from (I$^\times$). The pair constant $C_I''$ of (I$^\times$) enters $K_1$ and
does not enter the number $\mathsf{l}^J$ below. The datum of line $3$ comes from the bound
(M$^\times$) of Lemma~\ref{5.5:lem-pair-bounds}. In line $2$ the size of $D^{(j)}$ is
$C_BE_0\omega_{k,j}$.

The family $\mathfrak{D}^J$ is a finite family of non-negative numbers, each defined by
\eqref{5.5:eq-production-closure-b} and \eqref{5.5:eq-difference-number-a} and finite. A
priori,
\begin{equation}\label{5.5:eq-difference-number-1}
\mathsf{S}^{\perp}_v\le 2\max\{\mathfrak{n}^{\perp}_v,\mathfrak{n}_{\hat A}\},
\end{equation}
where $\mathfrak{n}^{\perp}_v$ is the entry bound of the off-chart comparison for the
cell-unit $v$ and $\mathfrak{n}_{\hat A}$ is the one-machine size of the pointed Wilson family
$\hat A$, and
\begin{equation}\label{5.5:eq-difference-number-2}
\mathsf{b}_j\le C_\beta\mathfrak{t}_j^{-2}\,\mathsf{S},\qquad
\mathsf{b}_j\le 2\bigl(b_++\tfrac{2}{3}r\bigr),
\end{equation}
where $\mathsf{S}=\mathsf{S}^{\mathrm{cl}}_j[F]$ for the family $F$ to which the pair bound
(\(\Pi^\times\)) of Lemma~\ref{5.5:lem-pair-bounds} is applied at the counterterm entry, and
$\mathfrak{t}_j$ is the level-$j$ constant of that pair bound.

\smallskip
\emph{(f) The difference number.} The \emph{difference number of the two expansions} is
\begin{equation}\label{5.5:eq-difference-number-b}
\begin{aligned}
\mathsf{l}^J &:= \mathsf{l}^J[\mathrm{cl}^{\mathrm{pr}}_J] := \max\bigl\{(\alpha),\,(\beta)_3\bigr\},\\
(\alpha) &:= \max_{j\le k,\,F}\Bigl\{\mathsf{S}^{\mathrm{cl}}_j[F],\ \mathsf{a}^{\mathrm{cl}}_j,\
  \mathsf{S}^{\mathrm{cl}}_j[F]+\mathfrak{n}_F\,\mathsf{a}^{\mathrm{cl}}_j,\\
&\qquad\quad \sup_v\mathsf{S}^{\perp}_v,\ \mathsf{b}_j+b_+\mathsf{a}^{\mathrm{cl}}_j\Bigr\},\\
(\beta)_3 &:= \sup_{n\le k}\ C\sum_{j\le n}\bar\rho^{\,2}_{n,j}\,
  \mathsf{S}^{\mathrm{cl}}_j[R^{(j)}].
\end{aligned}
\end{equation}
Here $F$ runs over the small-field families of \eqref{5.5:eq-difference-number-a}, and
$\mathfrak{n}_F$ is the one-machine size of $F$. The supremum $\sup_v\mathsf{S}^{\perp}_v$ is
taken over the class-$W$ cell-units $v$ and enters with coefficient $1$.

The member $(\beta)_3$ is the weighted sum itself; it is not normalised. By the bound on the
flow sums $\mathfrak{m}_n$ in the proof of
Theorem~\ref{5.1:thm-correlation-theorem}, the one-machine $j$-sum of line $3$,
\begin{equation*}
C\sum_{j\le n}\rho^2_{n,j}\,g_j^{\kappa_0}\le C\mathfrak{m}_n ,
\end{equation*}
converges through the size factor $g_j^{\kappa_0}$, and so its two-machine version is carried
as the sum itself.

The same functional, read on the $U$-side production closure $\mathrm{cl}^{\mathrm{pr}}_U$ in
place of $\mathrm{cl}^{\mathrm{pr}}_J$, is written $\mathsf{l}^J[\mathrm{cl}^{\mathrm{pr}}_U]$. It
satisfies
\begin{equation*}
\mathsf{l}^J[\mathrm{cl}^{\mathrm{pr}}_U]\ge\mathsf{l}^J ,
\end{equation*}
and it is this quantity that the bound on the difference number estimates.
\end{definition}

\begin{proof}[Proof of the claims in Definition~\ref{5.5:def-difference-number}]
\emph{The family $\mathfrak{D}^J$.} The operation $\mathbf{R}_k$ acts on a lattice of a torus,
which is finite. There are finitely many levels $j\le k$, finitely many small-field families of
each level in the four lines of the table and among the small terms of
Corollary~\ref{5.2:cor-sourced-preparatory}, and finitely many cells, so finitely many
class-$W$ cell-units. Hence $\mathfrak{D}^J$ is a finite family. We next show that each of its
members is a well-defined non-negative real number.

\emph{The sizes $\mathsf{S}^{\mathrm{cl}}_j[F]$ and $\mathsf{a}^{\mathrm{cl}}_j$.} These are
defined by \eqref{5.5:eq-production-closure-b} on $\mathrm{cl}^{\mathrm{pr}}_J$. They are
non-negative because they are sizes of moduli of differences. They are finite by clause (iii)
of Theorem~\ref{5.5:thm-closure-theorem}.

\emph{The sizes $\mathsf{S}^{\perp}_v$.} Each $\mathsf{S}^{\perp}_v$ is a supremum of moduli,
hence non-negative. Its finiteness follows from the a priori bound
\eqref{5.5:eq-difference-number-1}, which is obtained by the mechanism of the finiteness
clause (iii) of Theorem~\ref{5.5:thm-closure-theorem}, applied at the off-chart slot. Let $p$
be a point of $\mathrm{cl}^{\mathrm{pr}}_J(i,\Delta;\circ)$. By the triangle inequality,
\begin{equation*}
|\delta A(c,p)|\le |A^{\mathsf{B}}(c,P^{\mathsf{B}})|+|A^{\mathsf{A}}(c,P^{\mathsf{A}})|.
\end{equation*}
We bound each one-machine value by $\max\{\mathfrak{n}^{\perp}_v,\mathfrak{n}_{\hat A}\}$, in two
cases.
\begin{itemize}
\item At the points of the image of the comparison move $\mathsf{M}^{\perp}$, the entry bound
of the off-chart comparison on that image bounds each one-machine value by
$\mathfrak{n}^{\perp}_v$. This bound
holds under the hypotheses of the comparison move, whose smallness condition is met for
couplings in the range fixed by $\gamma_J$. Hence the supremum of $|\delta A(c,p)|$ over the
image of $\mathsf{M}^{\perp}$ alone is at most $2\mathfrak{n}^{\perp}_v<\infty$.
\item At the remaining points of the class, each one-machine value of the pointed Wilson family
is bounded by its one-machine size $\mathfrak{n}_{\hat A}$, the one-machine entry bound for
$\hat A$ on the production closure.
\end{itemize}
Taking the supremum over $p$ gives \eqref{5.5:eq-difference-number-1}.

\emph{The coefficient differences $\mathsf{b}_j$.} Each $\mathsf{b}_j$ is a modulus, hence
non-negative. It satisfies two bounds. First, the pair bound (\(\Pi^\times\)) of
Lemma~\ref{5.5:lem-pair-bounds}, in the form displayed in \eqref{5.5:eq-pair-bounds-a}, gives
\begin{equation*}
\mathsf{b}_j\le C_\beta\mathfrak{t}_j^{-2}\,\mathsf{S},
\end{equation*}
with $\mathsf{S}$ as in \eqref{5.5:eq-difference-number-2}. Second, the proof of the
correlation theorem (Theorem~\ref{5.1:thm-correlation-theorem}) gives for each machine
$\mathsf{m}$ the one-machine bound
\begin{equation*}
|b^{\mathsf{m}}_j(\hat w)|\le b_++\tfrac{2}{3}r .
\end{equation*}
The triangle inequality, applied to the two machines' values, then gives
\begin{equation*}
\mathsf{b}_j\le |b^{\mathsf{B}}_j(\hat w)|+|b^{\mathsf{A}}_j(\hat w)|
\le 2\bigl(b_++\tfrac{2}{3}r\bigr).
\end{equation*}
This proves \eqref{5.5:eq-difference-number-2}; in particular $\mathsf{b}_j$ is finite.

\emph{The datum of line $2$.} With the size of $D^{(j)}$ as in (e), we apply (M$^\times$) of
Lemma~\ref{5.5:lem-pair-bounds} and then (F$^\times$) of
the same lemma, at $t=L^{j-n}$; the letters $\omega$, $C_M$, $D_X$, $\rho$, $\kappa$, $d_j$ and
$\delta$ below are those of these two bounds. This bounds the two-machine difference of line
$2$ by
\begin{equation*}
2C_B\,\mathsf{S}^{\mathrm{cl}}_j[E]\,\omega\bigl(C_MD_Xt^2\rho\bigr)^2
e^{-(1-\frac14\delta)\kappa d_j}.
\end{equation*}
The difference size entering this bound is $\mathsf{S}^{\mathrm{cl}}_j[E]$. This is the datum
$d^{(2)}_j$ of \eqref{5.5:eq-difference-number-a}. The data $d^{(1)}_j=d^{(\partial)}_j$ and
$d^{(3)}_j$ are the difference sizes entering (I$^\times$) and (M$^\times$) of the same lemma,
respectively.

\emph{Finiteness of $\mathsf{l}^J$.} The quantity $(\alpha)$ is a maximum over finitely many
$j$ and $F$ of finitely many numbers. Each of these numbers is finite by the preceding steps,
because $\mathfrak{n}_F$ and $b_+$ are finite constants. In particular $\sup_v\mathsf{S}^{\perp}_v$
is a maximum over the finitely many class-$W$ cell-units $v$, each term being finite by
\eqref{5.5:eq-difference-number-1}. The quantity $(\beta)_3$ is a supremum over the finitely
many $n\le k$ of finite sums of products of the numbers $C\bar\rho^{\,2}_{n,j}$ with finite
sizes. Hence $\mathsf{l}^J$ is a finite non-negative number.

\emph{The inequality for the $U$-side closure.} Every member of $(\alpha)$ and of $(\beta)_3$ is
non-decreasing when the family of points over which the difference sizes are taken is enlarged
slotwise. For the sizes $\mathsf{S}^{\mathrm{cl}}_j[F]$, $\mathsf{a}^{\mathrm{cl}}_j$ and
$\mathsf{S}^{\perp}_v$ this holds because a supremum over a larger set is not smaller. The
number $\mathsf{b}_j$, being the modulus of the difference of the two coefficients, does not
depend on the family of points. The members
$\mathsf{S}^{\mathrm{cl}}_j[F]+\mathfrak{n}_F\mathsf{a}^{\mathrm{cl}}_j$,
$\mathsf{b}_j+b_+\mathsf{a}^{\mathrm{cl}}_j$ and $(\beta)_3$ are combinations of these numbers
with non-negative coefficients, so they inherit the monotonicity. A maximum of non-decreasing
quantities is non-decreasing. The inequality
$\mathsf{l}^J[\mathrm{cl}^{\mathrm{pr}}_U]\ge\mathsf{l}^J$ is this monotonicity combined with
the slotwise inclusion
$\mathrm{cl}^{\mathrm{pr}}_J(\sigma)\subseteq\mathrm{cl}^{\mathrm{pr}}_U(\sigma)$ at every
slot $\sigma$. This inclusion holds because every $J$-side seed is a $U$-side seed: the family
$\mathrm{cl}^{\mathrm{pr}}_U$ therefore contains $\mathbf{1}$ and the $J$-side seeds and, by
clause (i) of Theorem~\ref{5.5:thm-closure-theorem}, is closed under every use of
$\mathfrak{M}^{\mathrm{pr}}$, while by the same clause $\mathrm{cl}^{\mathrm{pr}}_J$ is the
least family with these properties.
\end{proof}

\begin{proposition}[The difference number as a monotone linear functional]
\label{5.5:prop-difference-linear}
Assume the standing hypotheses (H1)--(H4) on the single expansions, displayed in
\eqref{5.5:eq-standing-hypotheses-a}. Let $\mathfrak{D}^{J}$ be the finite family of difference
data of \eqref{5.5:eq-difference-number-a}, formed over the production closure
$\mathrm{cl}^{\mathrm{pr}}_{J}$ of Definition~\ref{5.5:def-production-closure-bis}. Let
\begin{equation*}
\mathsf{l}^{J} = \mathsf{l}^{J}[\mathrm{cl}^{\mathrm{pr}}_{J}] = \max\{(\alpha),(\beta)_3\}
\end{equation*}
be the difference number of Definition~\ref{5.5:def-difference-number}, with $(\alpha)$ and
$(\beta)_3$ as in \eqref{5.5:eq-difference-number-b}. For a correspondence $\mathcal{K}$, that is,
an assignment of a family of pair points to every slot, $\mathsf{l}^{J}[\mathcal{K}]$ denotes the
same expression with every difference size taken over $\mathcal{K}$ in place of
$\mathrm{cl}^{\mathrm{pr}}_{J}$. The \emph{off-chart cell-units} are the cell-units
$v=(c,i,\Delta)$ of the block cell-pieces of $\mathbf{R}_k$ counted in the fifth row of the
per-cell line of Theorem~\ref{5.1:thm-correlation-theorem}, and their companions,
which lie off the chart; they are the units $v$ for which the difference size
$\mathsf{S}^{\perp}_v$ of \eqref{5.5:eq-difference-number-a} is defined. Write $\mathsf{A}$,
$\mathsf{B}$ for the two expansions compared in Definition~\ref{5.5:def-difference-number}. Then
the following hold.
\begin{enumerate}
\item[(a)] \emph{Linearity.} $\mathsf{l}^{J}$ is the maximum of finitely many forms
\begin{equation*}
\Lambda_i(\mathfrak{D}^{J}) = \sum_{d\in\mathfrak{D}^{J}} \lambda_{i,d}\, d ,
\qquad \lambda_{i,d}\ge 0 ,
\end{equation*}
whose coefficients $\lambda_{i,d}$ are numbers common to both expansions, namely $1$,
$\mathfrak{n}_F$, $E_0$, $b_+$, $C\bar\rho^2_{n,j}$, and $1$ on each $\mathsf{S}^{\perp}_v$; here
$\sup_v\mathsf{S}^{\perp}_v$ is the maximum, over the off-chart cell-units $v$ of
$\mathbf{R}_k$, of the coordinate forms $d\mapsto \mathsf{S}^{\perp}_v$. Consequently
$\mathsf{l}^{J}$ is positively homogeneous, subadditive and monotone in each datum, and
$\mathsf{l}^{J}=0$ when every datum is $0$.
\item[(b)] \emph{Floor.} For every $j\le k$ and every family $F$,
\begin{equation*}
\mathsf{l}^{J} \ge \max\bigl\{\mathsf{S}^{\mathrm{cl}}_j[F],\ \mathsf{a}^{\mathrm{cl}}_j,\
\mathsf{S}^{\mathrm{cl}}_j[F]+\mathfrak{n}_F\mathsf{a}^{\mathrm{cl}}_j,\
\sup_v\mathsf{S}^{\perp}_v,\ \mathsf{b}_j+b_+\mathsf{a}^{\mathrm{cl}}_j\bigr\};
\end{equation*}
moreover $\mathsf{l}^{J}\ge(\beta)_3$, and $\mathsf{l}^{J}\ge d^{(r)}_j$ for the difference datum
$d^{(r)}_j$ of every row $r$.
\item[(c)] \emph{Finiteness; no count of levels.} Let $\mathfrak{n}_F$, $\mathfrak{n}_{\hat A}$,
$g_j^{\kappa_0}$ and $\mathfrak{n}^{\perp}_v$ be the one-expansion sizes, each the larger of the
values in the two expansions; here $\mathfrak{n}^{\perp}_v$ is the entry bound, at the off-chart
slot $(i,\Delta;\circ)$, of the off-chart cell-unit $v$, given by the entry-bound lemma for the
comparison move $\mathsf{M}^{\perp}$. Then
\begin{multline*}
\mathsf{l}^{J} \le \max\Bigl\{ 2\max_F\mathfrak{n}_F(1+\mathfrak{n}_{\hat A}),\
2\mathfrak{n}_{\hat A},\ 2\max\bigl\{\sup_v\mathfrak{n}^{\perp}_v,\mathfrak{n}_{\hat A}\bigr\},\\
2\bigl(b_+ + \tfrac23 r\bigr) + 2b_+\mathfrak{n}_{\hat A},\ 2C\bar{\mathfrak{m}}\Bigr\} < \infty ,
\end{multline*}
where $r$ is the radius of Definition~\ref{5.1:def-weight-classes}, the third member is the
supremum over $v$ of the a-priori bound
$\mathsf{S}^{\perp}_v\le 2\max\{\mathfrak{n}^{\perp}_v,\mathfrak{n}_{\hat A}\}$, and
\begin{equation}\label{5.5:eq-difference-linear-a}
\bar{\mathfrak{m}} := 2\Bigl[(1+2c_2\gamma_J^2)\Bigl(1+\frac{1}{\lambda(\kappa_0/2-1)}\Bigr)
+ 3\gamma_J^2\Bigl(1+\frac{1}{\lambda(\kappa_0/2-2)}\Bigr)\Bigr].
\end{equation}
The constant $\bar{\mathfrak{m}}$ is the series bound of the proof of the correlation theorem
(Theorem~\ref{5.1:thm-correlation-theorem}) taken at the two exponents $a=\kappa_0/2$ and
$a=\kappa_0/2-1$, both larger than $1$ because $\kappa_0>4$ (the standing choice gives
$\kappa_0>6$). It is applied to the cross terms $\bar\rho^2_{n,j}(g_j^{\mathsf{m}'})^{\kappa_0}$ of
$(\beta)_3$, in which $\bar\rho_{n,j}$ may be the weight of the other expansion, and not only in
the one-expansion form in which it enters that proof. No member of the right-hand side depends on
the number of levels, on $K$, on the parameters denoted $N$, $N_0$ and $\check R$ in the
construction underlying the correlation theorem, or on the age of a large-field region (the number
of steps since the step that created it), and
none requires an upper bound on $L$; the coupling enters only through its ceiling $\gamma_J$.
\item[(d)] \emph{Monotonicity in the correspondence.} If $\mathcal{K}\subset\mathcal{K}'$ slotwise,
then $\mathsf{l}^{J}[\mathcal{K}]\le\mathsf{l}^{J}[\mathcal{K}']$. In particular
\begin{equation*}
\mathsf{l}^{J}[\mathcal{K}^{\mathrm{rd}}_J] \le \mathsf{l}^{J}[\mathrm{cl}^{\mathrm{pr}}_J]
= \mathsf{l}^{J} \le \mathsf{l}^{J}[\mathrm{cl}^{\mathrm{pr}}_U]
\quad\text{and}\quad
\mathsf{l}^{J}[\mathcal{K}^{\mathrm{rd}}] \le \mathsf{l}^{J}[\mathrm{cl}^{\mathrm{pr}}_U].
\end{equation*}
\item[(e)] \emph{Diagonal.} At $\mathsf{B}=\mathsf{A}$ with diagonal seeds, $\mathsf{l}^{J}=0$.
\item[(f)] \emph{The per-cell line with the size replaced by $\mathsf{l}^{J}$.} Let
$r\in\{1,2,3,\partial\}$ be a row of the per-cell line of the correlation theorem and
$\Delta'$ a unit cell of scale $n$; the row index, written $r$ here and as a sub- or superscript
in $K_r$, $d^{(r)}_j$ and $c^{(r)}_{n,j}$, is never the radius $r$ of (c). Suppose that every
unit $v$ of row $r$ anchored in $\Delta'$ is bounded, in the two expansions, by its row-by-row
pair bound built from the pair bounds of Lemma~\ref{5.5:lem-pair-bounds}: the cores of rows $1$
and $\partial$ by $(\mathrm{I}^{\times})$ with the pair constant $C_I''$, the units of rows $2$ and
$3$ and the tails by $(\mathrm{M}^{\times})$, $(\partial^{\times})$, $(\mathrm{F}^{\times})$, and
the counterterm entry of row $\partial$ by $(\mathrm{W}^{\times})$. Then
\begin{equation*}
\sum_{v\in\text{row } r,\ \Delta'} (\text{pair bound of } v)
\le K_r\,\mathsf{l}^{J}\times(\text{the common cell factors of row } r),
\end{equation*}
where $K_r$ is the per-cell constant of row $r$ with its size factor replaced by $1$:
\begin{gather*}
K_1 := C_I''S_I,\qquad K_2 := CS_{\omega},\qquad K_3 := 1,\\
K_{\partial} := 27d\,C_{\partial}'M^3S_{\partial}^{\sharp}\big|_{C_I\mapsto C_I''}
+ 27d\cdot 21M\mathfrak{a}_n^2 .
\end{gather*}
In $K_1$ the pair constant $C_I''$ of $(\mathrm{I}^{\times})$ stands in place of the
one-expansion constant $C_I'$; this replacement enters through $K_1$ and not through
$\mathsf{l}^{J}$. In $K_3$ the row constant $C$ is carried by $(\beta)_3$. In $K_{\partial}$, where
$d=4$ is the dimension, the
first term is the core constant with $C_I$ replaced by $C_I''$, and the second is the counterterm
entry's constant with its size factor replaced by $1$. Thus the one-expansion per-cell line of the
correlation theorem holds, row by row, for the differences of the two expansions with the size
factor replaced by $\mathsf{l}^{J}$, and hence for their sum over the rows. In this form the line
is the input of the per-cell summation bound of the two-expansion comparison, in which the size
constant $E_0+1$ is replaced by $\mathsf{l}^{J}$, and of the stored lines of that comparison, each
of which holds for the difference weights with its size constant replaced by the number
$\mathsf{l}_{\mathsf{n}+1}$ forwarded at length $\mathsf{n}+1$, of which $\mathsf{l}^{J}$ is a
member.

For the off-chart cell-units $v=(c,i,\Delta)$, whose two-run difference bound is given by the
off-chart comparison theorem for $\mathsf{M}^{\perp}$ and its corollary (the bound of the fifth
row with $\mathsf{s}_v\mathsf{v}_v$ replaced by $C^{\perp}\mathsf{S}^{\perp}_v$, which is linear in
$\mathsf{S}^{\perp}_v$; $C^{\perp}$ is the comparison constant of that theorem), the line with the
size replaced by $\mathsf{l}^{J}$ reads per unit
\begin{equation*}
\mathsf{S}^{\perp}_v \le \sup_v\mathsf{S}^{\perp}_v \le (\alpha) \le \mathsf{l}^{J},
\end{equation*}
and per cube as follows. Write $\Delta_v$ for the cell $\Delta$ of $v=(c,i,\Delta)$. The cubes
$q$, their families $\pi_i$ indexed by $i$, the weights $\alpha_{\cdot,i}$ and the distance $D_v$
of $v$ in the exponential factor below are those of the off-chart comparison theorem for
$\mathsf{M}^{\perp}$ and its corollary. For a cube $q$ let $V(q)$ be the set of cell-units $v$ of
the block pieces with $\Delta_v\subset q$, and put
\begin{equation}\label{5.5:eq-difference-linear-b}
\mathfrak{s}^{\perp}_{\alpha} := \sup_{i,\ q\in\pi_i}\ \sum_{v\in V(q)}\alpha_{\cdot,i}\,
\mathsf{S}^{\perp}_v ,
\end{equation}
a per-cube sum of difference sizes, finite a priori by clause (iii) of the entry-bound lemma for
$\mathsf{M}^{\perp}$ and by clause (c) of the corollary of the off-chart comparison theorem. The
per-cube line is the block sum of clause (b) of that corollary,
\begin{equation*}
\sum_{v \text{ at } \square} \alpha_{\cdot,i}\, e^{-\frac14\delta D_v}\,\mathsf{S}^{\perp}_v
\le c^{W}_{\Sigma}\, \mathfrak{s}^{\perp}_{\alpha},
\end{equation*}
where $c^{W}_{\Sigma}$ is the block-sum constant of that clause; it is the per-cube bound entering
the weighted per-unit majorant at
$\mathsf{s}_v\mathsf{v}_v\mapsto C^{\perp}\mathsf{S}^{\perp}_v$.
\end{enumerate}
\end{proposition}

\begin{proof}
We prove the six parts in order.

\smallskip\noindent
(a) We read off Definition~\ref{5.5:def-difference-number}. By
\eqref{5.5:eq-difference-number-b}, the members of the maximum $(\alpha)$ are of three kinds. The
first kind consists of single data: $\mathsf{S}^{\mathrm{cl}}_j[F]$, $\mathsf{a}^{\mathrm{cl}}_j$,
and the data $\mathsf{S}^{\perp}_v$, which enter with coefficient $1$ each and are maximised over
$v$ in $\sup_v\mathsf{S}^{\perp}_v$. The second kind is
$\mathsf{S}^{\mathrm{cl}}_j[F]+\mathfrak{n}_F\mathsf{a}^{\mathrm{cl}}_j$, with coefficients $1$
and $\mathfrak{n}_F$. The third kind is $\mathsf{b}_j+b_+\mathsf{a}^{\mathrm{cl}}_j$, with
coefficients $1$ and $b_+$. The quantity $(\beta)_3$ is, for each $n$, the sum
$\sum_{j\le n} C\bar\rho^2_{n,j}\,d^{(3)}_j$, with coefficients $C\bar\rho^2_{n,j}\ge 0$, and it
is maximised over $n$. The coefficient $E_0$ enters through the row data
$d^{(1)}_j=d^{(\partial)}_j=\mathsf{S}^{\mathrm{cl}}_j[E]+E_0\mathsf{a}^{\mathrm{cl}}_j$. All
these coefficients are nonnegative numbers that do not depend on the expansion.
Hence $\mathsf{l}^{J}=\max\{(\alpha),(\beta)_3\}$ is the maximum of finitely many linear forms
$\Lambda_i$ with nonnegative coefficients.

Each such form is positively homogeneous: $\Lambda_i(t\mathfrak{D}^{J})=t\Lambda_i(\mathfrak{D}^{J})$
for $t\ge0$. It is additive, hence subadditive. It is nondecreasing in each datum, because its
coefficients are nonnegative. A maximum of finitely many such forms keeps the three properties.
First, $\max_i\Lambda_i(t\mathfrak{D}^{J})=t\max_i\Lambda_i(\mathfrak{D}^{J})$ for $t\ge 0$.
Second,
\begin{equation*}
\max_i\Lambda_i(\mathfrak{D}^{J}+\mathfrak{D}'^{J})
\le \max_i\Lambda_i(\mathfrak{D}^{J})+\max_i\Lambda_i(\mathfrak{D}'^{J}).
\end{equation*}
Third, a maximum of nondecreasing functions is nondecreasing. Finally, every $\Lambda_i$ vanishes
when all data are $0$, and so does their maximum.

\smallskip\noindent
(b) By \eqref{5.5:eq-difference-number-b}, $(\alpha)$ and $(\beta)_3$ are the two members of the
maximum defining $\mathsf{l}^{J}$, so $\mathsf{l}^{J}\ge(\alpha)$ and
$\mathsf{l}^{J}\ge(\beta)_3$. The displayed inequality of (b) is $\mathsf{l}^{J}\ge(\alpha)$,
read member by member for each $j\le k$ and each $F$. Each row datum is a member of $(\alpha)$.
For rows $1$ and $\partial$,
\begin{equation*}
d^{(1)}_j = d^{(\partial)}_j = \mathsf{S}^{\mathrm{cl}}_j[E]+\mathfrak{n}_E\mathsf{a}^{\mathrm{cl}}_j,
\qquad \mathfrak{n}_E=E_0,
\end{equation*}
which is the member $\mathsf{S}^{\mathrm{cl}}_j[F]+\mathfrak{n}_F\mathsf{a}^{\mathrm{cl}}_j$ at
$F=E$. For row $2$, $d^{(2)}_j=\mathsf{S}^{\mathrm{cl}}_j[E]$, which is the member
$\mathsf{S}^{\mathrm{cl}}_j[F]$ at $F=E$. For row $3$,
$d^{(3)}_j=\mathsf{S}^{\mathrm{cl}}_j[R^{(j)}]$, the difference size of the row-three family
$R^{(j)}$, which is the member $\mathsf{S}^{\mathrm{cl}}_j[F]$ at $F=R^{(j)}$. Hence, for every
row $r$,
\begin{equation*}
\mathsf{l}^{J}\ge(\alpha)\ge d^{(r)}_j .
\end{equation*}

\smallskip\noindent
(c) We bound the members of $(\alpha)$ first. Clause (iii) of
Theorem~\ref{5.5:thm-closure-theorem}, displayed in \eqref{5.5:eq-closure-theorem-a}, gives
$\mathsf{S}^{\mathrm{cl}}_j[F]\le\mathfrak{n}_F^{\mathsf{B}}+\mathfrak{n}_F^{\mathsf{A}}$ and
$\mathsf{a}^{\mathrm{cl}}_j\le\mathfrak{n}_{\hat A}^{\mathsf{B}}+\mathfrak{n}_{\hat A}^{\mathsf{A}}$.
For the second inequality, $\mathfrak{n}_{\hat A}$ is the one-expansion bound on the pointed Wilson
family $\hat A$ established, for the fifth row, in the proof of the correlation theorem. Since
$\mathfrak{n}_F$ and $\mathfrak{n}_{\hat A}$ denote the larger of the two expansions' values,
\begin{equation*}
\mathsf{S}^{\mathrm{cl}}_j[F]\le 2\mathfrak{n}_F,\qquad
\mathsf{a}^{\mathrm{cl}}_j\le 2\mathfrak{n}_{\hat A},\qquad
\mathsf{S}^{\mathrm{cl}}_j[F]+\mathfrak{n}_F\mathsf{a}^{\mathrm{cl}}_j
\le 2\mathfrak{n}_F+2\mathfrak{n}_F\mathfrak{n}_{\hat A} = 2\mathfrak{n}_F(1+\mathfrak{n}_{\hat A}).
\end{equation*}
The proof of the correlation theorem bounds the counterterm coefficient in either expansion by
$|b_j^{\mathsf{m}}(\hat w)|\le b_+ + \tfrac23 r$, with $r$ the radius of
Definition~\ref{5.1:def-weight-classes}. Therefore
\begin{equation*}
\mathsf{b}_j \le |b_j^{\mathsf{B}}(\hat w)|+|b_j^{\mathsf{A}}(\hat w)|
\le 2\bigl(b_+ + \tfrac23 r\bigr),
\qquad
\mathsf{b}_j+b_+\mathsf{a}^{\mathrm{cl}}_j \le 2\bigl(b_+ + \tfrac23 r\bigr)+2b_+\mathfrak{n}_{\hat A}.
\end{equation*}

For the off-chart cell-units we use the a-priori bound
$\mathsf{S}^{\perp}_v\le 2\max\{\mathfrak{n}^{\perp}_v,\mathfrak{n}_{\hat A}\}$ for every such $v$.
It comes from clause (iii) of the entry-bound lemma for $\mathsf{M}^{\perp}$. That clause gives the
entry bound $\mathfrak{n}^{\perp}_v$ of the slot $(i,\Delta;\circ)$ on the whole image of the
comparison move, including its values at $\mathsf{s}=0,1$, and in particular
$\mathsf{S}^{\perp}_v\le 2\mathfrak{n}^{\perp}_v<\infty$. The class on which this is read is
$\mathrm{cl}^{\mathrm{pr}}_J$, by clause (ii) of Theorem~\ref{5.5:thm-closure-theorem}. Taking the
supremum over $v$ gives
$\sup_v\mathsf{S}^{\perp}_v\le 2\max\{\sup_v\mathfrak{n}^{\perp}_v,\mathfrak{n}_{\hat A}\}$.
Inserting these five bounds into $(\alpha)$ yields the first four members of (c).

We now bound $(\beta)_3$. The row-three datum is the difference size
$d^{(3)}_j=\mathsf{S}^{\mathrm{cl}}_j[R^{(j)}]$ of the row-three family $R^{(j)}$ of the
correlation theorem. Its one-expansion size in that theorem is $(g_j^{\mathsf{m}})^{\kappa_0}$, so
clause (iii) of Theorem~\ref{5.5:thm-closure-theorem} gives
\begin{equation*}
\mathsf{S}^{\mathrm{cl}}_j[R^{(j)}] \le (g_j^{\mathsf{B}})^{\kappa_0}+(g_j^{\mathsf{A}})^{\kappa_0},
\end{equation*}
and therefore
\begin{equation}\label{5.5:eq-difference-linear-1}
(\beta)_3 \le C\,\sup_n \sum_{\mathsf{m}'\in\{\mathsf{A},\mathsf{B}\}}\ \sum_{j\le n}
\bar\rho^2_{n,j}\,(g_j^{\mathsf{m}'})^{\kappa_0}.
\end{equation}
Fix $n$ and $\mathsf{m}'$. In this paragraph we write $m:=n-j\ge 0$; distinct $j\le n$ give
distinct $m$.

The proof of the correlation theorem bounds the row weight in either expansion $\mathsf{m}$, for
$g\le\gamma_J$, by
\begin{equation*}
(\rho^{\mathsf{m}}_{n,j})^2\le(1+o(1))\bigl(1+(3\lambda m+2c_2)\gamma_J^2\bigr).
\end{equation*}
Since $\bar\rho_{n,j}$ is the larger of $\rho^{\mathsf{A}}_{n,j}$ and $\rho^{\mathsf{B}}_{n,j}$,
the same bound holds for $\bar\rho^2_{n,j}$. In that proof the factor $1+o(1)$ is at most $2$
under $g\le\gamma_J$. Hence
\begin{equation*}
\bar\rho^2_{n,j}\le 2(1+2c_2\gamma_J^2)+6\gamma_J^2\,\lambda m .
\end{equation*}

The proof of the correlation theorem also gives a lower bound for the inverse squared coupling in
expansion $\mathsf{m}'$: $x_j^{\mathsf{m}'}\ge y+\lambda m$, with $y:=y_n^{\mathsf{m}'}\ge1$. As
$x_j^{\mathsf{m}'}=(g_j^{\mathsf{m}'})^{-2}$, this gives
$(g_j^{\mathsf{m}'})^{\kappa_0}\le(y+\lambda m)^{-\kappa_0/2}$. Multiplying the two bounds, using
$\lambda m\le y+\lambda m$ in the second term, and extending the sum to all $m\ge0$ (the terms are
nonnegative), we obtain
\begin{align*}
\sum_{j\le n}\bar\rho^2_{n,j}(g_j^{\mathsf{m}'})^{\kappa_0}
&\le \sum_{m\ge0}\bigl[2(1+2c_2\gamma_J^2)+6\gamma_J^2\lambda m\bigr](y+\lambda m)^{-\kappa_0/2}\\
&\le 2\Bigl[(1+2c_2\gamma_J^2)\sum_{m\ge0}(y+\lambda m)^{-\kappa_0/2}
+3\gamma_J^2\sum_{m\ge0}(y+\lambda m)^{-(\kappa_0/2-1)}\Bigr].
\end{align*}

We apply the series bound used in the proof of the correlation theorem: for $a>1$,
\begin{equation*}
\sum_{m\ge0}(y+\lambda m)^{-a}\le y^{-a}+\frac{y^{1-a}}{\lambda(a-1)} .
\end{equation*}
This holds because the term $m=0$ is $y^{-a}$, and for $m\ge1$ the term $(y+\lambda m)^{-a}$ is at
most $\int_{m-1}^{m}(y+\lambda s)^{-a}\,ds$, the integrand being decreasing in $s$. Moreover
$\int_0^\infty(y+\lambda s)^{-a}\,ds=y^{1-a}/(\lambda(a-1))$. We use it at $a=\kappa_0/2$ and at
$a=\kappa_0/2-1$; both exceed $1$ because $\kappa_0>4$. This gives
\begin{multline*}
\sum_{j\le n}\bar\rho^2_{n,j}(g_j^{\mathsf{m}'})^{\kappa_0}
\le 2\Bigl[(1+2c_2\gamma_J^2)\Bigl(y^{-\kappa_0/2}+\frac{y^{1-\kappa_0/2}}{\lambda(\kappa_0/2-1)}\Bigr)\\
+3\gamma_J^2\Bigl(y^{1-\kappa_0/2}+\frac{y^{2-\kappa_0/2}}{\lambda(\kappa_0/2-2)}\Bigr)\Bigr]
\le\bar{\mathfrak{m}} .
\end{multline*}
The last inequality holds because $y\ge1$ and the exponents $-\kappa_0/2$, $1-\kappa_0/2$ and
$2-\kappa_0/2$ are at most $0$ when $\kappa_0>4$, so each power of $y$ is at most $1$. Comparing
with \eqref{5.5:eq-difference-linear-a} gives the bound $\bar{\mathfrak{m}}$. There are two choices
of $\mathsf{m}'$, so \eqref{5.5:eq-difference-linear-1} gives $(\beta)_3\le 2C\bar{\mathfrak{m}}$.
This holds uniformly in $n$, $k$ and $K$, and whichever expansion carries $\bar\rho$. With the
bounds on the members of $(\alpha)$, this is the displayed inequality of (c).

It remains to prove the last assertion of (c). The right-hand side is built from the one-expansion
sizes, from $b_+$, the radius $r$ and $C$, and from $\bar{\mathfrak{m}}$, which involves only
$c_2$, $\lambda$, $\kappa_0$ and $\gamma_J$. The one-expansion sizes are those of the proof of the
correlation theorem. There the per-cell sums $S_I$, $S_{\omega}$, $S_{\partial}^{\sharp}$ are
independent of $K$ and bounded uniformly in $g\le\gamma_J$. The only sum over levels in that proof
without a weight, $\sum_j\rho^2g_j^{\kappa_0}$, is the one controlled by the series bound just
used. Hence no member depends on the number of levels, on $K$, on the parameters $N$, $N_0$,
$\check R$ of the statement or on the age of a large-field region,
or requires an upper bound on $L$, and the coupling enters only through $\gamma_J$.

\smallskip\noindent
(d) Each difference size of $\mathfrak{D}^{J}$, computed over a correspondence $\mathcal{K}$, is a
supremum over the points that $\mathcal{K}$ assigns to the relevant slots. If
$\mathcal{K}\subset\mathcal{K}'$ slotwise, each datum over $\mathcal{K}$ is therefore at most the
same datum over $\mathcal{K}'$. The monotonicity in each datum from (a) then gives
$\mathsf{l}^{J}[\mathcal{K}]\le\mathsf{l}^{J}[\mathcal{K}']$. Clause (i) of
Theorem~\ref{5.5:thm-closure-theorem} supplies the slotwise inclusions
\begin{equation*}
\mathrm{cl}^{\mathrm{pr}}_J\subset\mathrm{cl}^{\mathrm{pr}}_U,\qquad
\mathcal{K}^{\mathrm{rd}}_J\subset\mathrm{cl}^{\mathrm{pr}}_J,\qquad
\mathcal{K}^{\mathrm{rd}}\subset\mathrm{cl}^{\mathrm{pr}}_U ,
\end{equation*}
which give the three particular inequalities. The identity
$\mathsf{l}^{J}[\mathrm{cl}^{\mathrm{pr}}_J]=\mathsf{l}^{J}$ is
\eqref{5.5:eq-difference-number-b}.

\smallskip\noindent
(e) At $\mathsf{B}=\mathsf{A}$ with diagonal seeds, clause (v) of
Theorem~\ref{5.5:thm-closure-theorem} gives that every datum of $\mathfrak{D}^{J}$ vanishes. By
the last assertion of (a), $\mathsf{l}^{J}=0$.

\smallskip\noindent
(f) Fix $r$ and a cell $\Delta'$ of zone $n$, and for a unit of level $j$ write $t=L^{j-n}$.
Level by level, the units of row $r$ anchored in
$\Delta'$ are those counted in the proof of the correlation theorem. In rows $1$--$3$ and $5$,
there are $t^{-4}$ anchor points (respectively at most $t^{-4}$ cubes) per cell, and this cancels
the factor $t^4$; in row $\partial$ the counts are those of the boundary lemma of that proof. Both
expansions are built on the same source-free scaffold
(Definition~\ref{5.1:def-source-free-scaffold}), so the index sets and the cell factors are the
same in both.

The row-by-row pair bounds, built from the pair bounds of Lemma~\ref{5.5:lem-pair-bounds}
displayed in \eqref{5.5:eq-pair-bounds-a}, bound
the difference of each unit by its one-expansion bound with the size replaced by the difference
datum. The factors $\nu$, $\rho$, $\hat\alpha$, $\kappa$, $d_j$, $\omega$, $C_M$,
$D_X$, $C_B$, $\delta$ below are those of the one-expansion bounds in the proof of the correlation
theorem.

For row $1$, $(\mathrm{I}^{\times})$ bounds the core and counterterm group by
\begin{equation*}
C_I''\bigl(\mathsf{S}^{\mathrm{cl}}_j[E]+E_0\mathsf{a}^{\mathrm{cl}}_j\bigr)\nu^3\rho^2(1+\rho)
t^{4+\hat\alpha} = C_I''\,d^{(1)}_j\,\nu^3\rho^2(1+\rho)t^{4+\hat\alpha}.
\end{equation*}
The tails are bounded by
$2\mathsf{S}^{\mathrm{cl}}_j[E]e^{-\kappa d_j}\le 2d^{(1)}_je^{-\kappa d_j}$, since
$E_0\mathsf{a}^{\mathrm{cl}}_j\ge0$.

For row $2$, the composition of $(\mathrm{F}^{\times})$ with $(\mathrm{M}^{\times})$ gives
\begin{equation*}
2C_B\,\mathsf{S}^{\mathrm{cl}}_j[E]\,\omega\,(C_MD_Xt^2\rho)^2e^{-(1-\frac14\delta)\kappa d_j}
= 2C_B\,d^{(2)}_j\,\omega\,(C_MD_Xt^2\rho)^2e^{-(1-\frac14\delta)\kappa d_j}.
\end{equation*}

For row $3$, the bound is $2\mathsf{S}^{\mathrm{cl}}_j[R^{(j)}](C_MD_Xt^2\rho)^2e^{-\kappa d_j}$.

For row $\partial$, each unit's difference is bounded by the pair forms of clauses (i)--(iii) of
its one-expansion bound and by $(\mathrm{I}^{\times})$, all taken at the datum $d^{(\partial)}_j$
and with $C_I$ replaced by $C_I''$. The counterterm entry is bounded by
$(\mathrm{W}^{\times})$ as $(\mathsf{b}_j+b_+\mathsf{a}^{\mathrm{cl}}_j)\cdot4\mathfrak{a}_n^2t^4$.

Sum first over the level-$j$ units in $\Delta'$ and then over $j\le n$. This repeats the
computation in the last column of the per-cell line of the correlation theorem, with the row size
$\mathfrak{n}^{(r)}_j$ replaced by the difference datum. Hence the two-expansion total of row $r$
over $\Delta'$, the left-hand side of the display in (f), satisfies
\begin{multline*}
\sum_{v\in\text{row } r,\ \Delta'} (\text{pair bound of } v)\\
\le \sum_{j\le n} c^{(r)}_{n,j}\,(\text{datum})_j\times(\text{common cell factors}).
\end{multline*}
For $r\ne3$ every datum is at most $(\alpha)\le\mathsf{l}^{J}$, by (b). So the right-hand side is
at most $\mathsf{l}^{J}\sum_jc^{(r)}_{n,j}\times(\text{common cell factors})$, which is at most
$\mathsf{l}^{J}K_r\times(\text{common cell factors})$. The sums $S_I$, $S_{\omega}$,
$S_{\partial}^{\sharp}$ in $K_r$ stay independent of $K$ when $\rho$ is replaced by $\bar\rho$,
because the bound on $(\rho^{\mathsf{m}}_{n,j})^2$ used in (c) holds in either expansion. For
$r=3$ the weighted sum is at most $(\beta)_3\times(\text{common cell factors})$, and $(\beta)_3\le
\mathsf{l}^{J}\cdot1$ by (b); thus $K_3=1$. This proves the per-cell line for each row, and summing
over the rows gives it for the sum.

The clause on the off-chart cell-units is the off-chart comparison theorem for $\mathsf{M}^{\perp}$
and clause (b) of its corollary, read as follows. Per unit, we use the chain displayed in (f),
which holds by (b).
Per cube, we sum in $\mathfrak{s}^{\perp}_{\alpha}$ of \eqref{5.5:eq-difference-linear-b}. Nothing
further is to be proved.
\end{proof}

\begin{remark}[Why the third-row sum enters unnormalised]\label{5.5:rem-row-three-sum}
In the difference number $\mathsf{l}^{J}$ of Definition~\ref{5.5:def-difference-number}, the
maximum \eqref{5.5:eq-difference-number-b} contains the member $(\beta)_3$ beside $(\alpha)$.
If $\mathsf{l}^{J}$ were $(\alpha)$ alone, the per-cell constant $K_3$ of the third row in
\eqref{5.5:eq-difference-number-a} would contain a count of levels. With $(\beta)_3$ in the
maximum, $K_3=1$, and no member of $\mathsf{l}^{J}$ contains a count of levels.
\end{remark}

\begin{proof}
For a level $j\le n$ let $R^{(j)}$ denote the family of the third row at level $j$, and let
$\mathsf{S}^{\mathrm{cl}}_j[R^{(j)}]$ be its difference size. The third row contributes the sum
over $j\le n$ of $\bar\rho_{n,j}^{\,2}\,\mathsf{S}^{\mathrm{cl}}_j[R^{(j)}]$.

Suppose first that only $(\alpha)$ is available. Each $\mathsf{S}^{\mathrm{cl}}_j[R^{(j)}]$ is at
most $(\alpha)$, and the sum can only be bounded as follows:
\begin{equation*}
\sum_{j\le n}\bar\rho_{n,j}^{\,2}\,\mathsf{S}^{\mathrm{cl}}_j[R^{(j)}]
\le (\alpha)\sum_{j\le n}\bar\rho_{n,j}^{\,2}.
\end{equation*}
The row weights satisfy the lower bound $\rho_{n,j}\ge\tfrac12$, a floor fixed with the
constants of the correlation theorem. The weight $\bar\rho_{n,j}$ is the larger of the two
machines' weights,
so $\bar\rho_{n,j}^{\,2}\ge\tfrac14$. Hence
\begin{equation*}
(\alpha)\sum_{j\le n}\bar\rho_{n,j}^{\,2}\ \ge\ (\alpha)\cdot\tfrac14\,\#\{j: j\le n\}.
\end{equation*}
The only constant this route gives for $K_3$ is therefore at least a quarter of the number of
levels up to $n$. That is a count of levels, not a constant independent of $n$.
This count arises only from bounding every summand by the single number $(\alpha)$; it is a
defect of that way of bounding the sum, not a bound that the third-row sum itself requires.

Now take $(\beta)_3$ in the maximum. Then $(\beta)_3$ bounds the third-row sum with
coefficient one, so $K_3=1$.

In a single machine the sum over $j$ converges because the size carries the factor
$g_j^{\kappa_0}$, where $\kappa_0>0$ is the exponent of the coupling in the one-machine size
bound. For differences this role is taken by whatever bounds
$\mathsf{S}^{\mathrm{cl}}_j[R^{(j)}]$ in $j$. That is either the recursion in $j$ for these
difference sizes that belongs to the production closure $\mathrm{cl}^{\mathrm{pr}}_{U}$, or,
a priori, clause~(c) of
Proposition~\ref{5.5:prop-difference-linear}.

Consequently no member of the maximum defining $\mathsf{l}^{J}$ contains a count of levels.
\end{proof}

\begin{theorem}[The recursion for difference sizes on the larger closure]
\label{5.5:thm-difference-recursion}
Put $\mathcal{K} := \mathrm{cl}^{\mathrm{pr}}_{U}$, the production closure generated by the $U$-side
seeds (Definition~\ref{5.5:def-production-closure-bis}, display \eqref{5.5:eq-production-closure-b}).
Write $\mathsf{S}^{\mathrm{cl}_U}_j[F]$ and $\mathsf{a}^{\mathrm{cl}_U}_j$ for the difference sizes
$\mathsf{S}^{\mathrm{cl}}_j[F]$ and $\mathsf{a}^{\mathrm{cl}}_j$ of
\eqref{5.5:eq-production-closure-b} formed over $\mathrm{cl}^{\mathrm{pr}}_{U}$. Assume the
following.
\begin{enumerate}
\item[(a)] The standing hypotheses (H1)--(H4) of Definition~\ref{5.5:def-standing-hypotheses},
displayed in \eqref{5.5:eq-standing-hypotheses-a}, including the $U$-side clause of the seed
hypothesis (H2)(s).
\item[(b)] On $\mathrm{cl}^{\mathrm{pr}}_{U}$, the items of the proposition on closure of the pair
format beyond its first four items.
\item[(c)] The closeness input $(\mathcal{C}^{\mathcal{K}})$ on $\mathrm{cl}^{\mathrm{pr}}_{U}$,
namely clauses $(\mathcal{C}2)$--$(\mathcal{C}6)$ on the correspondence at $C\vartheta^n$,
for every use of $\mathfrak{M}^{\mathrm{pr}}$, where $n$ is the scale at which the correspondence
is taken and $C$ is a constant. Equivalently, the table of pair conditions for
$\mathfrak{M}^{\mathrm{pr}}$ holds when read in local-index form. In that form the hereditary size
$\mathsf{m}$ (the size bounded in the hereditary lemma; in this theorem the letter does not denote
the machine index) and the clauses (g1)--(g4) are taken at the label stratum $l$ (the stratum of
the label at level $l$), in units of that
stratum, with bounds of the form $\le C_*\vartheta^l$, where $C_*$ is the constant of the table;
the hereditary lemma is then read stratum by stratum (see (S3)). The table is used only in this
local-index form: the single-index form, in which (g1) is bounded by $c_0\vartheta^j$ on seeds,
$c_0$ being the constant of that form, does not hold where regions meet strata created fewer than
the prescribed number of steps earlier.
\item[(d)] The remaining inputs of the pair-format statements invoked below (the proposition on
closure of the pair format, the two-machine recursion theorem and the hereditary lemma), namely:
\begin{itemize}
\item[(d1)] the small-piece statement on the enlarged region $X^{+}$;
\item[(d2)] the first of the statements from which the $U$-side seed clause (H2)(s2) is derived
in \eqref{5.5:eq-standing-hypotheses-a}, the one listed there before the small-piece statement;
\item[(d3)] the input on the box-pair points $\mathsf{P}^0$, $\mathsf{P}^1$;
\item[(d4)] $g\le\gamma_R$, where $\gamma_R$ is the coupling ceiling of the two-machine
recursion theorem.
\end{itemize}
\item[(e)] Clause (g2) of the table of pair conditions at the comparison move
$\mathsf{M}^{\perp}$, the move at the slot $(i,\Delta;\circ)$, and at the move of the table at
which the same clause is asked again for the fourth conclusion of the remainder bound. At both
moves this clause is supplied by the two-cutoff response
statement at zeroth order taken together with the flat energy-difference bound.
\item[(f)] At the rims of the localisation site (L), the hypothesis of the boundary-move lemma.
\end{enumerate}

Then the following hold.
\begin{enumerate}
\item[(S1)] The first four items of the proposition on closure of the pair format hold on
$\mathrm{cl}^{\mathrm{pr}}_{U}$. Together with hypothesis (b), every item of that proposition
holds on $\mathrm{cl}^{\mathrm{pr}}_{U}$.
\item[(S2)] The two-machine recursion theorem holds on $\mathrm{cl}^{\mathrm{pr}}_{U}$; granted
(b) and (c), its proof is algebra on the difference sizes. These sizes are taken to be
\begin{equation*}
(\mathsf{p}^{\flat},\mathsf{q},\mathsf{r},\mathsf{e})_j
:= \mathsf{S}^{\mathrm{cl}_U}_j[P^{\flat},Q,R,E],
\end{equation*}
where $P^{\flat}$, $Q$, $R$, $E$ are the level-$j$ families of the two-machine recursion theorem
(in this theorem $R$ denotes that family, not the ball radius), and its combined difference
$[\![\delta]\!]^{\times}$ is built on $\mathsf{S}^{\mathrm{cl}_U}$.
In particular
\begin{equation}\label{5.5:eq-difference-recursion-1}
\begin{split}
\mathsf{S}^{\mathrm{cl}_U}_k[R']
\le{}& 4E_0^{-1}e^{-p_0(g_k)}\Sigma^q_k
+4e^{-p_0(g_k)}q^{k-n_c}\\
&+2\bigl(9C_\chi^{\sharp}/\vartheta+2\bigr)e^{-p_0(g_k)}\vartheta^k ,
\end{split}
\end{equation}
where $R'$ is the level-$k$ family of the two-machine recursion theorem so denoted, $p_0(\cdot)$ is
the degree function, and $E_0$, $n_c$, $C_\chi^{\sharp}$ are the constants, and $\Sigma^q_k$ the
flow sum, of the two-machine recursion theorem.
\item[(S3)] Let $j$ be a level and $X$ a domain. If (g1)--(g3) hold for
$\mathfrak{M}^{\mathrm{pr}}$ at label stratum $l$, then the hereditary size $\mathsf{m}$ of (c)
satisfies $\mathsf{m}\le C_*\vartheta^l$ on $\mathrm{cl}^{\mathrm{pr}}_{U}(j,X)$, stratum by
stratum. This is the arithmetic of the hereditary lemma, re-indexed to local-index form.
\item[(S4)] The following inequality holds:
\begin{equation*}
\mathsf{l}^{J}\le\mathsf{l}^{J}[\mathrm{cl}^{\mathrm{pr}}_{U}].
\end{equation*}
Moreover, $\mathsf{l}^{J}$ is bounded by the outputs of the same recursion. Specifically:
\begin{itemize}
\item[--] $\mathsf{S}^{\mathrm{cl}_U}_j[F]$ is bounded by (S2);
\item[--] $\mathsf{a}^{\mathrm{cl}_U}_j$ is bounded through the bracket of $(W^{\times})$ in
\eqref{5.5:eq-pair-bounds-a} under clause $(\mathcal{C}2)$ of the closeness input;
\item[--] $\mathsf{b}_j$ is bounded by $(\Pi^{\times})$ of Lemma~\ref{5.5:lem-pair-bounds};
\item[--] $\sup_v\mathsf{S}^{\perp}_v$ is bounded by the same recursion at the slot
$(i,\Delta;\circ)$.
\end{itemize}
Thus $\sup_v\mathsf{S}^{\perp}_v$ is bounded like the other members of $\mathsf{l}^{J}$: by the
difference sizes of the families from which the retained members of $\mathsf{l}^{J}$ are built.
\end{enumerate}
\end{theorem}

\begin{proof}
The argument transfers three proofs to the class $\mathcal{K}=\mathrm{cl}^{\mathrm{pr}}_{U}$:
those of the proposition on closure of the pair format, of the two-machine recursion theorem, and
of the hereditary lemma. These proofs use the correspondence class $\mathcal{K}$ only through
five properties:
\begin{enumerate}
\item[(1)] the pair principle of Theorem~\ref{5.5:thm-closure-theorem}~(iv);
\item[(2)] the closure property (k1) at the uses which the production formulas perform;
\item[(3)] the closure property (k0);
\item[(4)] the one-machine sizes;
\item[(5)] the items of the proposition on closure of the pair format beyond its first four, and
the closeness input $(\mathcal{C}^{\mathcal{K}})$, equivalent to the table of pair conditions
read in local-index form at the moves of the list.
\end{enumerate}
In addition, the arithmetic of the hereditary lemma is an induction on chain length over the
uses of the list, carried out stratum by stratum.

We check that each of (1)--(5) holds for $\mathcal{K}=\mathrm{cl}^{\mathrm{pr}}_{U}$. By
hypothesis (a), the standing hypotheses hold with the $U$-side clause of (H2)(s). These are the
hypotheses under which Theorem~\ref{5.5:thm-closure-theorem} is stated for
$\mathrm{cl}^{\mathrm{pr}}_{U}$, so that theorem applies.
\begin{itemize}
\item[--] Property (1) is Theorem~\ref{5.5:thm-closure-theorem}~(iv).
\item[--] Properties (2) and (3) are the closure properties (k1) and (k0) of
Theorem~\ref{5.5:thm-closure-theorem}~(i). By that theorem, $\mathrm{cl}^{\mathrm{pr}}_{U}$ is
closed under every use of $\mathfrak{M}^{\mathrm{pr}}$. This includes the uses out of unit slots
and the uses out of the slots of the fourth item of the proposition on closure of the pair
format. Those uses are exactly the uses that the fourth item and the step ``fix $X$ and
$p\in\mathcal{K}(X)$'' of the two-machine recursion theorem require.
\item[--] Property (4) is Theorem~\ref{5.5:thm-closure-theorem}~(iii).
\item[--] Property (5) consists of hypotheses (b) and (c), indexed by the uses of
$\mathfrak{M}^{\mathrm{pr}}$.
\end{itemize}
The remaining inputs of these three proofs are hypotheses (d1)--(d4), and, at the rims of the
site (L), hypothesis (f).
Hence the three proofs hold verbatim with $\mathcal{K}=\mathrm{cl}^{\mathrm{pr}}_{U}$. We read off
the four conclusions.

\emph{Conclusion (S1).} The first three items follow from the transferred proof. Consider the
fourth item, in which the terms $\mathbf{B}$, $\mathbb{C}^{(\mathsf{n}+1)}_k$ and $V$ are read
again at later steps. By the production list \eqref{5.5:eq-production-list-a}, the slot of a
derived or stored object is the parent of its production reads. Let $o$ be the object read and
$\sigma(v)$ the slot of the reader, and let $u$ be a use of $\mathfrak{M}^{\mathrm{pr}}$ from
$\sigma(o)$ to $\sigma(v)$, with $\pi\in\mathfrak{P}_u$ and
$p\in\mathrm{cl}^{\mathrm{pr}}_{U}(\sigma(o))$. By (k1),
\begin{equation*}
\mu_u^{\times}(\pi;p)\in\mathrm{cl}^{\mathrm{pr}}_{U}(\sigma(v)).
\end{equation*}
The induction hypothesis of the fourth item is therefore read at the point
$\mu_u^{\times}(\pi;p)$, which lies in the class. The bound by twice the supremum norm on
$\mathfrak{P}_u\times\mathbb{D}^{\times}$, used there, where
$\mathbb{D}^{\times}=\mathbb{D}\times\mathbb{D}$ is the product of the two machines' copies of the
parameter disc $\mathbb{D}$, is
supplied by Theorem~\ref{5.5:thm-closure-theorem}~(iii) and~(iv). So the induction of the fourth
item goes through on $\mathrm{cl}^{\mathrm{pr}}_{U}$. The items beyond the fourth are hypothesis
(b).

\emph{Conclusion (S2).} Given (b) and (c), the proof of the two-machine recursion theorem is
algebra on the difference sizes. Take these sizes to be
$\mathsf{S}^{\mathrm{cl}_U}_j[P^{\flat},Q,R,E]$ and build $[\![\delta]\!]^{\times}$ on
$\mathsf{S}^{\mathrm{cl}_U}$; the algebra is then unchanged.

The size $\mathsf{S}^{\mathrm{cl}_U}_k[R']$ is a supremum over a class of points. The step ``fix
$X$ and $p\in\mathcal{K}(X)$'' of that proof needs an object at every point of this class. Such
an object exists, precisely because the class is the parent of the production reads of $R'$
(display \eqref{5.5:eq-production-list-a} and Definition~\ref{5.5:def-production-closure-bis}). On the
smaller class $\mathcal{K}^{\mathrm{rd}}$ of Theorem~\ref{5.5:thm-closure-theorem}~(i), by
contrast, taken under the closure condition with which that class was defined, the step finds no
such object.

The bound \eqref{5.5:eq-difference-recursion-1} is the particular case of the recursion theorem's
conclusion for the family $R'$ at level $k$, with these sizes.

\emph{Conclusion (S3).} By Definition~\ref{5.5:def-production-closure-bis}, every point of
$\mathrm{cl}^{\mathrm{pr}}_{U}(j,X)$ is a seed or $\mathbf{1}$, followed by a finite chain of uses
of $\mathfrak{M}^{\mathrm{pr}}$. The arithmetic of the hereditary lemma is an induction on the
length of this chain, taken over the uses of the list. Carry it out stratum by stratum, in the
local-index form of (c), with (g1)--(g3) holding at label stratum $l$. The induction is then
unchanged, and it gives $\mathsf{m}\le C_*\vartheta^l$ on $\mathrm{cl}^{\mathrm{pr}}_{U}(j,X)$,
stratum by stratum.

\emph{Conclusion (S4).} By Proposition~\ref{5.5:prop-difference-linear}~(a) and~(d),
\begin{equation*}
\mathsf{l}^{J}\le\mathsf{l}^{J}[\mathrm{cl}^{\mathrm{pr}}_{U}].
\end{equation*}
By (a) of that proposition, this quantity is the maximum of finitely many linear forms in the
difference data $\mathfrak{D}^{J}$ of \eqref{5.5:eq-difference-number-a}. The coefficients of these
forms are non-negative, so the forms are monotone in each datum. Hence each member of the maximum
\eqref{5.5:eq-difference-number-b} is bounded once each datum entering it is bounded. We bound the
data in turn.
\begin{itemize}
\item[--] $\mathsf{S}^{\mathrm{cl}_U}_j[F]$ is bounded by (S2).
\item[--] $\mathsf{a}^{\mathrm{cl}_U}_j$ is bounded through the bracket of $(W^{\times})$ in
\eqref{5.5:eq-pair-bounds-a} (Lemma~\ref{5.5:lem-pair-bounds}). That bracket is controlled under
clause $(\mathcal{C}2)$ of the closeness input (c).
\item[--] $\mathsf{b}_j$ is bounded by $(\Pi^{\times})$ of Lemma~\ref{5.5:lem-pair-bounds}.
\item[--] Consider $\sup_v\mathsf{S}^{\perp}_v$, the supremum over the class-$W$ cell-units $v$ of
the off-chart difference sizes of \eqref{5.5:eq-difference-number-a}; its per-cube form is
$\mathfrak{s}^{\perp}_{\alpha}$ of \eqref{5.5:eq-difference-linear-b}. It is bounded by the same
recursion, now run at the slot $(i,\Delta;\circ)$. The clause (g2) at the comparison move
$\mathsf{M}^{\perp}$ needed there is hypothesis (e).
\end{itemize}
Every datum of $\mathsf{l}^{J}[\mathrm{cl}^{\mathrm{pr}}_{U}]$ is thus bounded by outputs of the
same recursion, by the difference sizes of the families from which the retained members of
$\mathsf{l}^{J}$ are built. This is (S4).
\end{proof}

\begin{remark}[Closeness conditions use by use and gauges at rims]\label{5.5:rem-closeness-conditions}
Let the uses of the production list $\mathfrak{M}^{\mathrm{pr}}$ be as in
Definition~\ref{5.5:def-production-list}, each typed with the reader's slot as parent, and consider
condition (g2) among the closeness conditions (g1)--(g4) on the uses of $\mathfrak{M}^{\mathrm{pr}}$
that enter the hypotheses of Theorem~\ref{5.5:thm-difference-recursion}. For the uses of
$\mathfrak{M}^{\mathrm{pr}}$, condition (g2) is given by the table of closeness conditions for
the rim pair format, re-indexed use by use; no new kind of use is added. The uses fall into the
following classes.
\begin{enumerate}
\item[(a)] \emph{Uses for which the table gives the form of (g2).} These are restriction and
refinement, for which the table gives (g2) in exact form; change of units, for which it gives (g2)
with a factor at most $\tfrac14$; and the cut of the layer use together with its layers at strata
of index $p\le k-q$, for which it gives (g2) when the parents are real and tame, and otherwise for
the linear part only. The table records only the form that (g2) takes at these uses; (g2) at them
is not proved, in the table or here, the closeness conditions on the uses of
$\mathfrak{M}^{\mathrm{pr}}$ being not proved.
\item[(b)] \emph{Uses at which the table gives no form of (g2) applicable to the use.} The final
dilation: the table's entry for it
covers only the dilation whose radius is the natural majorant, and does not cover the final
dilation of $\mathfrak{M}^{\mathrm{pr}}$; (g2) there is not proved, and is among the inputs of the
uniqueness theorem that are not proved, beside the box, layer, tube and stored-jet reads. The
further uses of this class are the following: the use at the site (T) and the
layer use at the top strata $k-q<p\le k$, that is, the layers there and the reads at (T) and (P);
the box and layer reads, which replace the pair $(\mathfrak{b},\mathbf{A})$ by $(\text{child},0)$,
where $\mathfrak{b}$ is the box read datum and $\mathbf{A}$ the component to which tube reads add,
including the companions
with $\nu\le2$ of the low-order variation bound at the second cube size, whose chart factor
\begin{equation*}
\bigl(CD\Theta\,\mathfrak{a}_n t^2\bigr)^2,
\end{equation*}
with $t$ as in Lemma~\ref{5.5:lem-pair-bounds},
enters at the pair points of the box read through the off-chart clause of the pair bound
$(W^{\times})$ in~\eqref{5.5:eq-pair-bounds-a}, that is, through clause $(\mathcal{C}2)$ of (g4);
the tube reads, which add to $\mathbf{A}$; the stored jet on the box pair, which is a tube read;
and the comparison move $\mathsf{M}^{\perp}$, the dilation in the parameter $\mathsf{s}$ of radius
$e^{\delta D}$ on the parameter domain $\mathfrak{P}_{\perp}$.
\item[(c)] \emph{The dilations of radius $e^{\delta D}$.} At every dilation of radius $e^{\delta D}$
among the uses of $\mathfrak{M}^{\mathrm{pr}}$ --- the comparison move $\mathsf{M}^{\perp}$, and the
dilation of radius $e^{\delta D}$ inside the stored-jet use --- condition (g2) is the conjunction of
the two-cutoff response statement at zeroth order and the flat energy-difference bound, which
together assert the closeness of the responses at the two cutoffs in decaying weights. This
conjunction is not proved.
\item[(d)] \emph{Rims.} At a rim the use is made through a read in the gauge centred at a point
$y$; as restated in~\eqref{5.5:eq-pair-bounds-a}, $(M^{\times})$ holds for every class, rims
included, with $C_M$ replaced by $C_{\partial}$. The one-machine step behind $(M^{\times})$ at a
rim is the one-machine rim bound, whose hypothesis reads: \emph{$P=(\mathbf{U},\mathbf{J})$ is
given on $X$,
and on $\square\cap\operatorname{dom}\supset X$, where $\square$ is a cube of side $D_X$ in the units
of level $j$ about $y$, it is $G^{c}$-gauge equivalent to $e^{iA}$ with}
\begin{equation*}
|A|,\ |\nabla A|,\ |\mathbf{J}|\le\mathfrak{a}
\end{equation*}
\emph{in the units of a level $n\ge j$} --- the clauses (i)--(iii) of the definition of
$U^{c}_j$ in \cite[p.~262]{B-CMP109}, together with \cite[(1.65), (1.67), (1.69)]{B-CMP122a}. This
hypothesis is required only on the cubes of $n$-diameter at most $2M$ of
\cite[(1.69)]{B-CMP122a}; outside them the trivial bound
$2\mathsf{S}e^{-\kappa d}$ of clause (ii) of the one-machine bound for the class $(M)$ is used
instead, $d$ denoting there a distance and not the dimension. So restricted, it is an input of the
uniqueness theorem. Ba{\l}aban's construction provides
such gauges only on cubes of a size $O(1)LM$ \cite[(1.12), p.~262]{B-CMP109}. On the side of
Theorem~\ref{5.1:thm-correlation-theorem} (the correlation theorem) the hypothesis is not
supplied, and no statement on that side uses it.
\item[(e)] \emph{The neutral count at (P).} The neutral count of (g3) at the site (P) is the count
along the chain of successive tube reads at (P), a chain already present within the closeness
conditions; passing to the production closure neither creates it nor removes it.
\end{enumerate}
\end{remark}

\begin{proof}
By Definition~\ref{5.5:def-production-list}, each use of $\mathfrak{M}^{\mathrm{pr}}$ is typed with
the reader's slot as parent. The table of closeness conditions
for the rim pair format is indexed by the kind of use; assigning to each use of
$\mathfrak{M}^{\mathrm{pr}}$ the entry of its kind therefore gives condition (g2) use by use, and
since the kind of each use of $\mathfrak{M}^{\mathrm{pr}}$ is one of the kinds of the table, no new
kind is added. The table itself lists (g2) as not proved for the layers at strata $k-q<p\le k$ and
for the reads at (T) and (P); the uses of (a) are those whose line of the table gives the form of
(g2) at them, a form that is not proved, and
the uses of (b) are those at which (g2) is listed as not proved, together with the final dilation,
whose entry covers only the dilation whose radius is the natural majorant.

For (c), the comparison move $\mathsf{M}^{\perp}$ and the dilation inside the stored-jet use have
radius $e^{\delta D}$ and not the natural majorant, so the table's entry for the final dilation
does not cover them; (g2) at each of them is the closeness of the two-cutoff responses in decaying
weights, which is the conjunction named in (c). These two are the only dilations of radius
$e^{\delta D}$ among the uses of $\mathfrak{M}^{\mathrm{pr}}$; hence at no dilation of radius
$e^{\delta D}$ is (g2) supplied by the table.

For the companions with $\nu\le2$ in (b), the low-order variation bound at the second cube size
carries the chart factor $(CD\Theta\,\mathfrak{a}_n t^2)^2$; at the pair points of the box read it
enters through the off-chart clause of $(W^{\times})$ in~\eqref{5.5:eq-pair-bounds-a}, which is
clause $(\mathcal{C}2)$ of (g4).

For (d), the pair bound $(M^{\times})$ at a rim is obtained, as restated
in~\eqref{5.5:eq-pair-bounds-a}, for every class with $C_M$ replaced by $C_{\partial}$, from the
one-machine rim bound applied to each of the two machines $\mathsf{A}$ and $\mathsf{B}$; the read at
the rim is made in the gauge centred at $y$. Hence the hypothesis of the one-machine rim bound is
needed wherever $(M^{\times})$ is used at a rim, and it is needed only on the cubes of
$n$-diameter at most $2M$ of \cite[(1.69)]{B-CMP122a}, since outside them the trivial bound
$2\mathsf{S}e^{-\kappa d}$ of clause (ii) of the one-machine bound for the class $(M)$ uses no
gauge. The hypothesis consists of the gauge conditions (i)--(iii) of the definition of $U^{c}_j$ in
\cite[p.~262]{B-CMP109} and of \cite[(1.65), (1.67), (1.69)]{B-CMP122a}, required on these cubes;
the gauges of \cite[(1.12), p.~262]{B-CMP109} are provided only on cubes of a size $O(1)LM$. The
hypothesis enters the uniqueness theorem as an input. On the side of
Theorem~\ref{5.1:thm-correlation-theorem} the hypothesis is not supplied; and no statement there
uses $(M^{\times})$ at a rim through this step.

For (e), the production closure only re-indexes the existing kinds of use by their readers' slots
and adds no new kind; it therefore leaves the chain of successive tube reads at (P), along which
the neutral count of (g3) at (P) is taken, as it is, neither creating it nor removing it.
\end{proof}

\begin{definition}[Cut-offs and the interpolated box field at a large-field site]
\label{5.5:not-interpolated-box-field}
The following notation is fixed at the site (P), in the sense of
Definition~\ref{5.2:def-sites-units}, of Ba{\l}aban's large-field operation $\mathbf{R}_k$, at its
stage $\mathsf{n}+1$.

\emph{Cut-off sets and the source.} Put $\mathfrak{B} := \mathbb{B}^{(\mathsf{n})}_k$ and
$\mathfrak{B}^{+} := \mathbb{B}^{(\mathsf{n}+1)}_k$, the cut-off sets of the stages $\mathsf{n}$ and
$\mathsf{n}+1$. The datum $(\mathbb{U},\mathbb{J},b,\sigma)$ is taken in the source over a region
$\hat X\supset\square_0$ (the region $\square_0$ is fixed below), as read through the arrangement
statement at the site (P), namely in the source class
$\tilde{U}^{(\mathsf{n}+1)c,\Sigma\flat}_k(\hat X)$ of stage $\mathsf{n}+1$ over $\hat X$; hence
\begin{equation*}
(\mathbb{U},\mathbb{J},b,\sigma)\in\mathfrak{G}(\mathfrak{B}^{+})\subset\mathfrak{G}(\mathfrak{B})
\subset\mathfrak{K}_P ,
\end{equation*}
where $\mathfrak{K}_P$ denotes the class of configurations of the type of \cite{B-CMP109} over
$\mathfrak{B}$ at the radius $\gamma_0$, the smallness radius of that class. Moreover
$a_{\flat}(\mathbb{U})\le\gamma_0^{(k)}$, with $a_{\flat}$ computed in the unit of
$\mathfrak{B}^{+}$; here $\gamma_0^{\top} := \bar\gamma_0(\mathsf{t}_{\gamma})$, the value at
$\mathsf{t}_{\gamma}$ of the window-smallness radius $\bar\gamma_0(\cdot)$ of the window statement at
the site (P), and $\gamma_0^{\top}\ge\gamma_0^{(k)}$. Finally $|\sigma(\Delta)|\le e^{\kappa_1}$ for
every cell $\Delta$ on which $\sigma$ is defined, and
$\|b\|_\infty\le 1.1\,C_1\delta'_{\mathsf{j}}$, where $\delta'_{\mathsf{j}}$ is the parameter at the
index $\mathsf{j}$ with which the source is described.

\emph{Blocks and cut-off functions.} Let $\square$ be a block of zone $n_0$, so that
$\square\subset\Omega_{n_0}$, with $\Omega_{n_0}$ and $\ell_{n_0}$ as defined in the paragraph on
levels below. For $a\ge 1$, $\tilde{\square}^{a}$ is $\square$ enlarged by $a$ layers
of $\pi_{n_0}$-cubes, which have side $M\ell_{n_0}$. Put $\square_0 := \tilde{\square}^{5}$; the region
$\square_0^{\flat}$ satisfies $\square_0^{\flat}\supset\tilde{\square}^{3}$. Besides the cut-off
$\zeta_{\square}$ of the block $\square$, let $\zeta'_{\square}$ be a cut-off with
$\zeta'_{\square}=1$ on $\tilde{\square}^{3}$ and $\zeta'_{\square}=0$ off $\tilde{\square}^{4}$.

\emph{The box field and the interpolated layer.} The box field is
$\mathbb{W} := U(\mathfrak{B}^{+}(\square_0),M^{\cdot}(\mathbb{U}))$ on $\square_0$, extended by
$\mathbb{U}$ off $\square_0$; then $(\mathbb{W},\mathbb{J})\in\mathfrak{K}_P$. The field
$\mathbb{H}^{w}$ and the map $Q_{\star}$ are taken at the kernel $(\mathfrak{B};\mathbb{W},\mathbb{J})$,
with the two properties stated where they are constructed at the site (P). For $t\in[0,1]$ and
$\mathsf{t}\in\mathbb{D}_{\square} := \{|\mathsf{t}|\le r^{\sharp}_{\square}\}$ put
\begin{equation*}
\mathbf{A}(t,\mathsf{t}) := (t\zeta'_{\square}+\mathsf{t}\zeta_{\square})\,\mathbb{H}^{w},
\qquad
\mathfrak{b} := e^{i\eta\mathbf{A}}\,\mathbb{W}.
\end{equation*}

\emph{The one-scale set inside the box.} For a level $j$ put
\begin{equation*}
\mathfrak{B}^{\star}_j := \mathbb{B}_j(\square_0)\cup
\bigl\{\Gamma^{(\mathsf{n})}_i\cap\square_0\bigr\}_{i<j},
\end{equation*}
where $\Gamma^{(\mathsf{n})}_i$, $i<j$, are the sets of stage $\mathsf{n}$ of the operation
$\mathbf{R}_k$ at level $i$;
let $M_{\star}$ denote its averages, and define $Q_{\star}(A)$ by
$M_{\star}(e^{i\eta A}\mathbb{W}) =: e^{iQ_{\star}(A)}M_{\star}\mathbb{W}$. For a data bond $c$,
$\tilde B(c)$ denotes the two $\mathfrak{B}^{\star}_j$-blocks of $c$, as in \cite{B-CMP98}. The distance
$d_{\star} := d^{\wedge}_{\mathfrak{B}^{\star}_j}$ is taken in the frame
$(\square_0,\mathfrak{B},\mathfrak{B}^{\star}_j)$, lengths inside $\square_0$ being measured in units
at most $\ell_j$. The pair fed to a near term $T$ of level $j$ is
\begin{equation*}
\mathsf{P}^{0}(t,\mathsf{t}) := \mathfrak{b}'_j :=
(U,J)\bigl(\mathfrak{B}^{\star}_j,M_{\star}\mathfrak{b}\bigr).
\end{equation*}

\emph{Kernel quantities capped at $j$.} Let $\mathfrak{B}_j$ be $\mathfrak{B}$ capped at $j$. The
profile $\varrho$ and the profiles $\hat p,\hat p_{\square}$ given by the weight-profile statement at
the site (P) belong to $\mathfrak{M}^{\wedge}_{\delta/2}(\mathfrak{B}_j)$, with $\delta$ the parameter
of Definition~\ref{5.2:def-step-constants}. Put
\begin{equation*}
\mathsf{B}(t,\mathsf{t}) :=
Q_{\mathfrak{B}_j}\bigl(\mathbb{W};\eta(t+\mathsf{t}\zeta_{\square})\mathbb{H}^{w}\bigr),
\end{equation*}
and let $\mathbf{H}_j(t,\mathsf{t})$ be the third member of the triple of fields given by the
theorem which provides $\mathbb{H}^{w}$, taken at the data
$(\mathfrak{B}_j;\mathbb{W},\mathbb{J};\sigma;(0,\mathsf{B}(t,\mathsf{t})))$.

\emph{Levels, weights and floors.} $\ell(x)$ and $w(x)$ are the level and weight functions,
$\ell_m := L^m\eta$, $\Omega_m := \{\ell\ge m\}$,
$\alpha_{\cdot,m} := \min\{\alpha_{0,m},\alpha_{1,m}\}$, and
\begin{equation*}
\theta_2(x) := w(x)\,\tilde{\alpha}_{0,\ell(x)}\,\eta^2\,\ell_{\ell(x)}^{-2}.
\end{equation*}
The constant $w := M/M_1$ satisfies $w\le M$ (the weight function is always written $w(x)$, with
its argument); with $\varrho_0$ the floor of Corollary~\ref{5.2:cor-response-dilation}, the
following two floors are in force:
\begin{equation}\label{5.5:eq-interpolated-box-field-a}
\delta w\ge 16\ln L+4,
\qquad
w\ge 3L(\varrho_0+2L).
\end{equation}

\emph{Non-unitarity number.} For $g\in G^{c}$ put $n(g) := \|g\|\,\|g^{-1}\|$; then
$n(g)\ge\|\operatorname{Ad}_g\|$.

\emph{Letters kept apart.} The response number is written $\nu_S(x)$, distinct from the depth
number $\nu$ of a unit; the interpolation parameter of Ba{\l}aban's construction is written
$\mathsf{s}$, distinct from $s := n_0-j$; the loop increment is written $\Upsilon$, while
$\Delta'$ remains the unit $n$-cell.
\end{definition}

\begin{proof}[Proof of the assertions in Notation~\ref{5.5:not-interpolated-box-field}]
The datum $(\mathbb{U},\mathbb{J},b,\sigma)$ is the source
$\tilde{U}^{(\mathsf{n}+1)c,\Sigma\flat}_k(\hat X)$ over $\hat X\supset\square_0$ described at the
site (P). Its membership in $\mathfrak{G}(\mathfrak{B}^{+})$ and the two inclusions
$\mathfrak{G}(\mathfrak{B}^{+})\subset\mathfrak{G}(\mathfrak{B})\subset\mathfrak{K}_P$ are given by the
description of that source at the site (P), by the arrangement statement through which it is read,
and by the window statement at the site (P). The bound
$a_{\flat}(\mathbb{U})\le\gamma_0^{(k)}$ in the unit of $\mathfrak{B}^{+}$ is the window statement
applied at $\mathfrak{B}^{+}$ to the source pairs. The inequality
$\gamma_0^{\top}=\bar\gamma_0(\mathsf{t}_{\gamma})\ge\gamma_0^{(k)}$ is part of the window statement.
The bounds on $|\sigma(\Delta)|$ and $\|b\|_\infty$ are part of the description of the source at the
site (P). The membership $(\mathbb{W},\mathbb{J})\in\mathfrak{K}_P$ of the box field is the one
established at the site (P) for the configuration
$U(\mathfrak{B}^{+}(\square_0),M^{\cdot}(\mathbb{U}))$ extended by $\mathbb{U}$ off $\square_0$.
The membership of the profiles $\varrho,\hat p,\hat p_{\square}$ in
$\mathfrak{M}^{\wedge}_{\delta/2}(\mathfrak{B}_j)$ is the conclusion of the weight-profile statement at
the site (P).

The only assertion proved here is the last inequality. Let $X$ lie in $\mathfrak{g}^{c}$; here
$\|\cdot\|$ is the operator norm on matrices, which is submultiplicative, and
$\|\operatorname{Ad}_g\|$ is the operator norm of $\operatorname{Ad}_g$ on $\mathfrak{g}^{c}$ with
respect to it. Then
\begin{equation*}
\|\operatorname{Ad}_g X\| = \|gXg^{-1}\|\le\|g\|\,\|X\|\,\|g^{-1}\| = n(g)\,\|X\| .
\end{equation*}
Taking the supremum over $\|X\|\le 1$ gives $\|\operatorname{Ad}_g\|\le n(g)$.
\end{proof}

\begin{definition}[Box-read near units and their own boxes]\label{5.5:def-box-read-units}
We work at the site (P) of $\mathbf{R}_k$, stage $\mathsf{n}+1$, in the setting and notation of
Notation~\ref{5.5:not-interpolated-box-field}. In particular $\square$ is a block of zone $n_0$,
$\tilde{\square}^{a}$ is $\square$ enlarged by $a$ layers of $\pi_{n_0}$-cubes, $\ell(x)$ and $w(x)$ are
the level and weight functions, and $\alpha_{0,m}$ are the radii of that notation; $\alpha_{3,j}$ denotes
the radius of index $3$ at level $j$ in the same family of radii.

\emph{(a) Near units and their depth numbers.} A \emph{near unit} is a pair $v=(T,\square)$ in which
$T$ is a constituent of level $j$ of one of the first three kinds in the table of constituents of the
localised terms, that is, a core or a short single, with localisation domain
$X:=X_T\subset\tilde{\square}^{2}$. The \emph{anchor} $y$ of $v$ is the level-$j$ lattice point of $T$,
and $|\cdot|$ denotes the sup norm measured in level-$j$ units. We put
\begin{equation*}
n:=\min\{\ell(y),n_0\},\qquad \nu:=1+n-j,\qquad t:=L^{j-n},\qquad s:=n_0-j,
\qquad \Theta:=c_W\,w^{\vee}(y),
\end{equation*}
where $w^{\vee}(y)$ denotes the weight read at the anchor $y$ in the sense of the statement on levels
and weights, and
\begin{equation}\label{5.5:eq-box-read-units-1}
c_W:=\max\Bigl\{1,\ 2c_1K_{\mathrm{loc}}L^{3}\,\sup_{n}\frac{\alpha_{0,n}}{\mathfrak{a}_n}\Bigr\}.
\end{equation}
In this definition and in what follows, $s$ and the number $p$ of (c) are these depth numbers; they are
not the depth $s$ below the cutoff nor a plaquette $p$ of the general conventions. Likewise $t$ and the
number $t_p$ of (c) denote the scale ratios $L^{j-n}$ and $L^{j-p}$ of the near unit, not the
interpolation parameter $t\in[0,1]$ of the layer $\mathbf{A}(t,\mathsf{t})$ of
Notation~\ref{5.5:not-interpolated-box-field}.
The \emph{scaffold majorant} and the \emph{box ratio} of $v$ are
\begin{equation}\label{5.5:eq-box-read-units-2}
|B^{ax}(b)|_{\star}:=2d\,(1+|b-y|)\,\Theta\,\mathfrak{a}_n\,t^{2},\qquad
S^{\star}_v:=\frac{\alpha_{3,j}}{\sup_{b}|B^{ax}(b)|_{\star}},
\end{equation}
where $b$ runs over the arguments at which the axial majorant $B^{ax}$ is read; this $b$ is not
the field $b$ of Notation~\ref{5.5:not-interpolated-box-field}.

\emph{(b) The box-read predicate.} The near unit $v$ is \emph{box-read} if and only if
\begin{equation}\label{5.5:eq-box-read-units-3}
\nu\ge 3\qquad\text{and}\qquad S^{\star}_v>1.
\end{equation}
This is a predicate of the source-free scaffold of Definition~\ref{5.1:def-source-free-scaffold}.

\emph{(c) The own box and its cubes.} For a box-read near unit $v$ put $p:=n-1$ and $t_p:=L^{j-p}$.
The \emph{own box} $\mathfrak{r}_v\supset X$ of $v$ is $\square_T^{\sim\nu+1}$ if $T$ is a core, where
$\square_T$ denotes the cube of the core $T$ in the table of constituents of the localised terms (not
the block $\square$ of zone $n_0$) and $(\cdot)^{\sim a}$ denotes the enlargement of index $a$ used in
that table, the index $-2$ in (iii) below being taken in the sense of the statement on the own box at a
box-read unit; and $\mathfrak{r}_v$ is the enlargement of $X$ by the
margin of $5/2$ cubes prescribed for short singles in that table if $T$ is a short single. We put
\begin{equation*}
Q_v:=[\mathfrak{r}_v]_p^{(2)},\qquad
\mathsf{a}_{Q_v}:=K_{\mathrm{loc}}\,w_{Q_v}\,\alpha_{Q_v},
\end{equation*}
where $[\,\cdot\,]^{(2)}_p$ is the hull of order $2$ in cubes of level $p$, $Q_v^{(1)}$ is the
companion set of $Q_v$ in the statement on the own box at a box-read unit, and $w_{Q_v}$, $\alpha_{Q_v}$
are the weight and the radius attached to the set $Q_v$ in the sense of the statement on levels and
weights.

\emph{(d) The block set of the own box.} $\mathfrak{B}_j(\mathfrak{r}_v)$ denotes the block set attached
to the box $\mathfrak{r}_v$ by the statement on block sets of a box, and $M_{\cdot}$ its
averages: it consists of the level-$j$ blocks inside $\mathfrak{r}_v$ and of a graded rim of blocks of
levels $i<j$ near $\partial\mathfrak{r}_v$, a rim bond $c$ of level $l(c)<j$ having length
$\ell_{l(c)}$; the anchor $y$ serves as a base point.

\emph{(e) Pieces placed at every near unit, and remainders.} Besides the constituents of (a), the
Wilson and counterterm pieces near $\square$ are placed on the pair of box data $\mathsf{P}^{0}$,
$\mathsf{P}^{1}$ at every near unit, box-read or not, as described in the remark on the Wilson and
counterterm pieces below. A box-read near unit carries no remainder term
$\operatorname{rem}(T,\square)$.

For a box-read near unit $v$ the following hold.
\begin{enumerate}
\item[(i)] $p\ge j+1$ and $t_p=Lt$.
\item[(ii)] The membership domain $\mathfrak{D}_v$ (a set of configurations, not the class
$\mathfrak{D}_j$ of localisation domains) is
$\mathfrak{D}_v=\mathfrak{U}_T=\mathcal{U}^{c}_j(X,a_j)$, the analyticity space of
Definition~\ref{5.1:def-analyticity-spaces}; on it $|T|\le\mathfrak{n}_v$; and the clause thresholds
satisfy $\theta_a\ge(1-\beta)\,\mathbf{r}^T_a$.
\item[(iii)] The own box $\mathfrak{r}_v$ has sides of at most $2M/t_p$ level-$j$ units;
$Q_v^{(1)}\subset\tilde{\square}^{3}$; every cell meeting $Q_v^{(1)}$ has level at least $p$;
$\alpha_{Q_v}\le 2\alpha_{0,n}$; and $\mathfrak{r}_v\subset Q_v^{\sim-2}$.
\item[(iv)] $\mathsf{a}_{Q_v}\le 2K_{\mathrm{loc}}\,\mathsf{A}_w(\gamma)$.
\end{enumerate}
\end{definition}

\begin{proof}
Let $v$ be box-read. By the first condition in \eqref{5.5:eq-box-read-units-3},
$1+n-j=\nu\ge 3$, that is $n\ge j+2$; hence $p=n-1\ge j+1$. Moreover
\begin{equation*}
t_p=L^{j-p}=L^{j-n+1}=L\cdot L^{j-n}=Lt,
\end{equation*}
which is (i).

Assertion (ii) is the description of the membership domain, of the bound $\mathfrak{n}_v$ and of the
clause thresholds given, at a unit satisfying \eqref{5.5:eq-box-read-units-3}, by the statement on
membership domains and thresholds of units of the source-free scaffold of
Definition~\ref{5.1:def-source-free-scaffold}, where the domain is $\mathcal{U}^{c}_j(X,a_j)$; the table
of constituents of the localised terms enters only through the shape of $T$, a core or a short single.

Assertion (iii) is the geometry of the own box: for a core, $\mathfrak{r}_v=\square_T^{\sim\nu+1}$, and
for a short single, $\mathfrak{r}_v$ is the enlargement of $X$ by $5/2$ cubes of margin, as in the
table of constituents of the localised terms; with $p=n-1\ge j+1$ and $t_p=Lt$ from (i), the statement
on the own box at a box-read unit gives the side bound $2M/t_p$ in level-$j$ units, the inclusion
$Q_v^{(1)}\subset\tilde{\square}^{3}$, the lower bound $p$ on the levels of the cells meeting
$Q_v^{(1)}$, the bound $\alpha_{Q_v}\le 2\alpha_{0,n}$ and the inclusion
$\mathfrak{r}_v\subset Q_v^{\sim-2}$.

Assertion (iv) is the bound on the localisation constant $\mathsf{a}_{Q_v}=K_{\mathrm{loc}}w_{Q_v}
\alpha_{Q_v}$ of the set $Q_v$ by $2K_{\mathrm{loc}}\mathsf{A}_w(\gamma)$, from the statement on
localisation constants of cubes of level at least $p$.
\end{proof}

\begin{definition}[The stored remainder of a box-read unit]\label{5.5:def-stored-remainder}
We work at the site fixed in Notation~\ref{5.5:not-interpolated-box-field} and use the objects
introduced there:
\begin{itemize}
\item the box field $\mathbb{W}$;
\item the field $\mathbb{H}^{w}$ and the map $Q_{\star}$ at the kernel;
\item the cut-offs $\zeta_{\square}$ and $\zeta'_{\square}$;
\item the interpolated layer
$\mathbf{A}(t,\mathsf{t}) = (t\zeta'_{\square}+\mathsf{t}\zeta_{\square})\mathbb{H}^{w}$,
where $t\in[0,1]$ and $\mathsf{t}\in\mathbb{D}_{\square}$;
\item the layered box field $\mathfrak{b} = e^{i\eta\mathbf{A}}\mathbb{W}$;
\item the set $\mathfrak{B}^{\star}_j$ and its averages $M_{\star}$;
\item the pair $\mathsf{P}^{0}(t,\mathsf{t}) = (U,J)(\mathfrak{B}^{\star}_j, M_{\star}\mathfrak{b})$;
\item the field $\mathbf{H}_j(t,\mathsf{t})$.
\end{itemize}
In this definition $t$ always denotes the real interpolation parameter. It is not the number
$L^{j-n}$ attached to a unit in Definition~\ref{5.5:def-box-read-units}.

Let $v = (T,\square)$ be a box-read unit in the sense of Definition~\ref{5.5:def-box-read-units},
where $T$ has level $j$, and put $X := X_T$. For a bond $c$ put
$\zeta'(c) := \zeta'_{\square}(c_-)$, where $c_-$ denotes the initial point of $c$. On the bonds
of $\mathfrak{B}^{\star}_j$ define
\begin{equation}\label{5.5:eq-stored-remainder-1}
\begin{aligned}
D^{\sharp} &:= \zeta'\cdot Q_{\star}\bigl((t+\mathsf{t}\zeta_{\square})\mathbb{H}^{w}\bigr)
 + (1-\zeta')\cdot Q_{\star}\bigl(\mathbf{H}_j(t,\mathsf{t})\bigr),\\
\tilde{\mathsf{B}} &:= Q_{\star}\bigl(\mathbf{A}(t,\mathsf{t})\bigr),\qquad
\mathsf{E} := \frac{1}{i}\log\bigl(e^{iD^{\sharp}}e^{-i\tilde{\mathsf{B}}}\bigr).
\end{aligned}
\end{equation}
Further put
\begin{equation}\label{5.5:eq-stored-remainder-2}
\mathsf{P}^{1}(t,\mathsf{t}) := (U,J)\bigl(\mathfrak{B}^{\star}_j, e^{i\mathsf{E}}M_{\star}\mathfrak{b}\bigr)
 = (U,J)\bigl(\mathfrak{B}^{\star}_j, e^{iD^{\sharp}}M_{\star}\mathbb{W}\bigr),
\end{equation}
and
\begin{equation}\label{5.5:eq-stored-remainder-3}
G_v(t,\mathsf{t}) := T(X;\mathsf{P}^{1}) - T(X;\mathsf{P}^{0}),\qquad
\operatorname{rem}(T,\square) := \int_0^1 dt\,
 \partial_{\mathsf{t}}\big|_{\mathsf{t}=0}\, G_v(t,\mathsf{t}).
\end{equation}
The quantity $\operatorname{rem}(T,\square)$ is called the \emph{stored remainder} of $v$.
\end{definition}

\begin{proof}[Proof of the equality in \eqref{5.5:eq-stored-remainder-2}]
In Notation~\ref{5.5:not-interpolated-box-field}, the map $Q_{\star}$ is defined, for a $1$-form
$A$, by
\begin{equation*}
M_{\star}\bigl(e^{i\eta A}\mathbb{W}\bigr) = e^{iQ_{\star}(A)}M_{\star}\mathbb{W}.
\end{equation*}
Take $A = \mathbf{A}(t,\mathsf{t})$. Since $\mathfrak{b} = e^{i\eta\mathbf{A}}\mathbb{W}$ and
$\tilde{\mathsf{B}} = Q_{\star}(\mathbf{A}(t,\mathsf{t}))$, this gives
$M_{\star}\mathfrak{b} = e^{i\tilde{\mathsf{B}}}M_{\star}\mathbb{W}$ bond by bond on
$\mathfrak{B}^{\star}_j$.

By the definition of $\mathsf{E}$ in \eqref{5.5:eq-stored-remainder-1}, and because the
exponential of a logarithm returns its argument, we have
$e^{i\mathsf{E}} = e^{iD^{\sharp}}e^{-i\tilde{\mathsf{B}}}$ on every bond of $\mathfrak{B}^{\star}_j$.
Multiplying bond by bond therefore gives
\begin{equation*}
e^{i\mathsf{E}}M_{\star}\mathfrak{b}
 = e^{iD^{\sharp}}e^{-i\tilde{\mathsf{B}}}e^{i\tilde{\mathsf{B}}}M_{\star}\mathbb{W}
 = e^{iD^{\sharp}}M_{\star}\mathbb{W}.
\end{equation*}
Applying $(U,J)(\mathfrak{B}^{\star}_j,\cdot)$ to both sides yields the equality in
\eqref{5.5:eq-stored-remainder-2}.
\end{proof}

\begin{lemma}[Support of the correction form and its distance]\label{5.5:lem-correction-support}
Let $v=(T,\square)$ be a box-read near unit in the sense of Definition~\ref{5.5:def-box-read-units},
$T$ of level $j$, with anchor $y$, $n=\min\{\ell(y),n_0\}$, depth number $s=n_0-j$, own box
$\mathfrak{r}_v$ and cube $Q_v^{(1)}$, in the setting of
Notation~\ref{5.5:not-interpolated-box-field}. Let $D^{\sharp}$, $\tilde{\mathsf{B}}$ and $\mathsf{E}$ be
the spliced data, the averaged layer and the correction form of
Definition~\ref{5.5:def-stored-remainder}, defined on the bonds of $\mathfrak{B}^{\star}_j$. For a bond
$c$ of $\mathfrak{B}^{\star}_j$ let $\tilde B(c)$ be the union of the two blocks of $c$.
\begin{enumerate}
\item[(a)] For every bond $c$ of $\mathfrak{B}^{\star}_j$ with $\tilde B(c)\subset\tilde{\square}^3$,
\begin{equation}\label{5.5:eq-correction-support-a}
\mathsf{E}(c)=0 ,
\end{equation}
for every $\sigma$ and every point of the complex source. Consequently, if
\begin{equation*}
F_{\square}:=\bigl\{c:\ \text{a block of } c \text{ meets } (\tilde{\square}^3)^{\complement}\bigr\},
\end{equation*}
then $\operatorname{supp}\mathsf{E}\subset F_{\square}$.
\item[(b)] Put $D_v:=d_{\star}(\mathfrak{r}_v,F_{\square})$.
Then $D_v$ is a geometric quantity, determined by $\square$, $\mathfrak{r}_v$ and the blocks of
$\mathfrak{B}^{\star}_j$ alone, and
\begin{equation}\label{5.5:eq-correction-support-b}
D_v\ge \tfrac15 ML^{s},\qquad 1+|b-y|\le 1+2D_v\quad\text{for every bond } b \text{ of } \mathfrak{r}_v .
\end{equation}
\end{enumerate}
\end{lemma}

\begin{proof}
(a) Let $c$ be a bond with $\tilde B(c)\subset\tilde{\square}^3$. On both blocks of $c$ the cut-off
satisfies $\zeta'_{\square}\equiv 1$, so on $\tilde B(c)$ the interpolated layer is
\begin{equation*}
\mathbf{A}(t,\mathsf{t})=(t\zeta'_{\square}+\mathsf{t}\zeta_{\square})\mathbb{H}^{w}
=(t+\mathsf{t}\zeta_{\square})\mathbb{H}^{w}.
\end{equation*}
By \cite[Proposition~4, (P8)]{B-CMP98}, the value $Q_{\star}(\cdot)(c)$ of the linearised map at $c$ reads its
$1$-form on $\tilde B(c)$ only. Hence
\begin{equation*}
\tilde{\mathsf{B}}(c)=Q_{\star}(\mathbf{A}(t,\mathsf{t}))(c)
=Q_{\star}\bigl((t+\mathsf{t}\zeta_{\square})\mathbb{H}^{w}\bigr)(c).
\end{equation*}
On the other hand $\zeta'(c)=\zeta'_{\square}(c_-)$, and the point $c_-$ of $c$ lies in $\tilde B(c)$, so
$\zeta'(c)=1$; the definition of $D^{\sharp}$ then gives
$D^{\sharp}(c)=Q_{\star}((t+\mathsf{t}\zeta_{\square})\mathbb{H}^{w})(c)$, because the term carrying the
factor $1-\zeta'(c)$ vanishes. Thus $D^{\sharp}(c)=\tilde{\mathsf{B}}(c)$, whence
$e^{iD^{\sharp}(c)}e^{-i\tilde{\mathsf{B}}(c)}=1$ and
\begin{equation*}
\mathsf{E}(c)=\tfrac1i\log\bigl(e^{iD^{\sharp}(c)}e^{-i\tilde{\mathsf{B}}(c)}\bigr)=\tfrac1i\log 1=0 .
\end{equation*}
No property of $\sigma$ or of the source point entered: the identity comes from the cut-off
$\zeta'_{\square}$ alone and uses no constraint on the source data. If $c\notin F_{\square}$, then
neither block of $c$ meets $(\tilde{\square}^3)^{\complement}$, that is $\tilde B(c)\subset\tilde{\square}^3$, and
\eqref{5.5:eq-correction-support-a} gives $\mathsf{E}(c)=0$; so $\operatorname{supp}\mathsf{E}\subset
F_{\square}$.

(b) The set $F_{\square}$ is determined by $\square$ and the blocks of $\mathfrak{B}^{\star}_j$ alone;
hence $D_v=d_{\star}(\mathfrak{r}_v,F_{\square})$ is determined by $\square$, $\mathfrak{r}_v$ and the
blocks of $\mathfrak{B}^{\star}_j$ alone.

By Definition~\ref{5.5:def-box-read-units}, the cube $Q_v^{(1)}$ lies within $6/L$ level-$n$ units of the
anchor $y\in\tilde{\square}^2$, and $n\le n_0$. Since $\tilde{\square}^3$ adds to $\tilde{\square}^2$ one
layer of $\pi_{n_0}$-cubes of side $M\ell_{n_0}$
(Notation~\ref{5.5:not-interpolated-box-field}), this gives
\begin{equation*}
\operatorname{dist}\bigl(\mathfrak{r}_v,(\tilde{\square}^3)^{\complement}\bigr)\ge (1-6/L)M\ell_{n_0}
\ge \tfrac14 M\ell_{n_0}.
\end{equation*}
A bond of $F_{\square}$ has a block meeting $(\tilde{\square}^3)^{\complement}$ and so lies within
$2\ell_j$ of $(\tilde{\square}^3)^{\complement}$. The distance $d_{\star}$ counts length in units of
size at most $\ell_j$, and
$M\ell_{n_0}/\ell_j=ML^{n_0-j}=ML^{s}$; therefore
\begin{equation*}
D_v\ge \tfrac14 ML^{s}-2\ge \tfrac15 ML^{s},
\end{equation*}
the last inequality because $ML^{s}\ge 40$ (it is equivalent to $\tfrac1{20}ML^{s}\ge 2$).
This is the first inequality of \eqref{5.5:eq-correction-support-b}.

For a bond $b$ of $\mathfrak{r}_v$, with $|\cdot|$ the sup norm in level-$j$ units,
\begin{equation*}
|b-y|\le \frac{2M}{t_p}\le 2ML^{s-1}=\frac{2}{L}\,ML^{s}\le \frac{10D_v}{L}\le 2D_v ,
\end{equation*}
where the third inequality uses $ML^{s}\le 5D_v$, just proved. Adding $1$ to both sides gives the second
inequality of \eqref{5.5:eq-correction-support-b}.
\end{proof}

\begin{definition}[Constants and smallness conditions for the remainder bound]
\label{5.5:def-remainder-constants}
We work at the large-field site in the setting and notation of
Notation~\ref{5.5:not-interpolated-box-field} and Definition~\ref{5.5:def-box-read-units}.

\smallskip
\emph{(a) Constants.} The following numbers are fixed after $d$, $L$, $M$, $\delta$ and $\kappa_1$,
and before the coupling ceiling $\gamma$. None of them depends on $K$, on a depth, on
the numbers $\nu$ and $s$ of Definition~\ref{5.5:def-box-read-units}, on the numbers denoted
$N$, $N_0$ and $\check{R}$ in the proof of the remainder bound, on the age of a large-field
region, or on the weight function $w(x)$. The input
constants are:
\begin{itemize}
\item $B_H$, $b_H$: the constants so denoted in the sourced preparatory bounds;
\item $c_A$, $c_1^{b}$, $K_S^{b}$: the constants so denoted in the smallness conditions of the
sourced preparatory bounds;
\item $r_V$, $C_V$: the radius and the constant of the bound on the averaged layer used in the
proof of the remainder bound;
\item $B_{\sharp}$, $K_J^{\sharp}$, $\varepsilon_{\sharp}$: the constants of the theorem
constructing the fields $\mathbb{H}^{w}$ at the kernel, $\varepsilon_{\sharp}$ being the radius
of the ball of data on which that theorem holds;
\item $c_4$: the constant of \cite[Proposition~4]{B-CMP98};
\item $B_3$: the constant of \cite[Theorem~1, Proposition~9]{B-CMP102b}, which depends on
$(d,L)$ only;
\item $c_1$, $K_{\mathrm{loc}}$, $c_W$: the constants of Definition~\ref{5.5:def-box-read-units}.
\end{itemize}
From them we put
\begin{align*}
K_w &:= \max\{4B_3,\ 2B_3+8LB_Hc_Ac_1^{b}+1\}, \qquad
K_1 := 270L^3B_HB_{\sharp}^2c_1^{b},\\
c_{\delta} &:= \sup_{x\ge0}\,(1+2x)\,e^{-\frac18\delta x}, \qquad
c_3' := \inf_j\frac{\alpha_{3,j}}{\alpha_{\cdot,j}}>0,\\
c_A' &:= c_A+\frac{1}{2dL^3r_V}, \qquad
K^{\mathrm{rem}} := \frac{26\,d\,\beta\,c_A'\,c_{\delta}\,K_1}{c_3'}, \qquad
K_N := 12\,dLM\,B_Hc_A^2c_1^{b},\\
r^{\mathrm{rem}}_{\square} &:= \min\Bigl\{r^{\sharp}_{\square},\
\frac{1}{14K^{\mathrm{rem}}\vartheta(\square)}\Bigr\}, \qquad
c_{\Sigma} := \sum_{i\ge2}\bigl(6L^i\bigr)^d e^{-(\delta w/40)L^i},\\
K_{\mathrm{inc}} &:= 4\max\{K_J^{\sharp},K_S^{b}\}.
\end{align*}
Here $w$ in $c_{\Sigma}$ denotes a positive number fixed together with $d$, $L$, $M$, $\delta$
and $\kappa_1$, so that $c_{\Sigma}$ is one of the numbers of this list; it is not the weight
function $w(x)$. Further, $c_3'$ is a ratio of the constants $\alpha_{3,j}$, $\alpha_{\cdot,j}$
of the cover used at
the localisation site; $\vartheta$,
$\vartheta(\square)$ and $\beta$ are the response numbers and the constant of the comparison of
responses at the site; $r^{\sharp}_{\square}$ is the radius of the disc $\mathbb{D}_{\square}$ of
the complex parameter $\mathsf{t}$.

\smallskip
\emph{(b) Ceilings.} The following conditions are smallness conditions on $\gamma$ of the same
kind as the smallness conditions of the sourced preparatory bounds; they enter
the ceiling $\gamma_J$. Each is imposed in supremum form over the parameter
$\mathsf{t}\ge\mathsf{t}_{\gamma}$, with $\mathsf{t}_{\gamma}$ and the function $\bar\gamma_0$ as in
Notation~\ref{5.5:not-interpolated-box-field}:
the window smallness enters through
$\gamma_0^{\top}=\bar\gamma_0(\mathsf{t}_{\gamma})$, which dominates $\gamma_0^{(k)}$ at every
level, and the constants $\alpha_{3,j}$, $\alpha_{\cdot,m}$ enter through their suprema over
levels. The constant $M$
enters \eqref{5.5:eq-remainder-constants-b}(i) through $K_N$ and
\eqref{5.5:eq-remainder-constants-b}(ii) explicitly, in the same way as $M$ enters
$K_{\mathrm{loc}}$ in Definition~\ref{5.5:def-box-read-units}.
\begin{equation}\label{5.5:eq-remainder-constants-a}
17K^{\mathrm{rem}}\vartheta\le1.
\end{equation}
\begin{equation}\label{5.5:eq-remainder-constants-b}
\begin{aligned}
&\text{(i)}\quad 100K_N(K_w+1)\gamma_0^{\top}\le1,\\
&\text{(ii)}\quad 220\,d^2M^2B_3c_1K_{\mathrm{loc}}\mathsf{A}_w(\gamma)\le1,\\
&\text{(iii)}\quad 100\,C_VB_3\sup_j\alpha_{3,j}\le1.
\end{aligned}
\end{equation}
Next, writing $(a_{\flat},\kappa_{\max})$ for the pair of parameters in which the smallness
conditions of the sourced preparatory bounds are written, we impose
\begin{equation}\label{5.5:eq-remainder-constants-c}
\begin{gathered}
\text{the smallness conditions of the sourced preparatory bounds hold with}\\
(a_{\flat},\kappa_{\max})=\Bigl((K_w+1)\gamma_0^{\top},\ 34B_{\sharp}^2\sup_m\alpha_{\cdot,m}\Bigr),
\end{gathered}
\end{equation}
together with
\begin{equation}\label{5.5:eq-remainder-constants-d}
\max\{1,2B_3\}\,\varepsilon_1^{\top}\le a_1,
\end{equation}
where $a_0$, $a_1$ are the two radii of \cite[Theorem~1]{B-CMP102b} and
\begin{equation*}
\varepsilon_1^{\top} := 2(K_w+1)\gamma_0^{\top}+5r^{\top},\qquad
r^{\top} := 2\Bigl(c_A(K_w+1)\gamma_0^{\top}+33B_{\sharp}^2\sup_m\alpha_{\cdot,m}\Bigr).
\end{equation*}
The four radius conditions imposed on the radius $r\le r^{\top}$ of the fourth step of the
proof of the remainder bound are implied by condition \eqref{5.5:eq-remainder-constants-c},
through the constant $C^{b}$ of the smallness conditions of the sourced preparatory bounds at the
substituted pair.
Finally,
\begin{equation}\label{5.5:eq-remainder-constants-e}
4B_{\sharp}^2\Bigl(1.1\,C_1\delta'_{\mathsf{j}}+\sup_m\alpha_{\cdot,m}\Bigr)
\le\min\{\varepsilon_{\sharp},c_4\}.
\end{equation}
Condition \eqref{5.5:eq-remainder-constants-e} ensures that the datum
$(0,\mathsf{B}(t,\mathsf{t}))$ lies in the ball of the theorem constructing $\mathbb{H}^{w}$, and
that the layer $\mathbf{H}_j$ lies in the domain of \cite[Proposition~4]{B-CMP98}, for every
$\mathsf{t}$ on the whole circle on which the parameter is read; this circle does not depend on
the region $X$.

\smallskip
\emph{(c) Floors.} No new floor condition is imposed. The floor conditions in force are the
floors on $L$ and on $M$ (among them $M\ge 3L$) fixed in
Notation~\ref{5.5:not-interpolated-box-field} and Definition~\ref{5.5:def-box-read-units}, and
those of the comparison of responses at the site.

\smallskip
\emph{(d) A consequence.} Let $\varepsilon_1'$ be the number of the smallness hypothesis of
\cite[Theorem~1]{B-CMP102b} evaluated at the data $\Phi(B)$ on $\mathfrak{B}_j(\mathfrak{r}_v)$ of
the configuration $Z(B)$ in the fourth step of the proof of the remainder bound, and assume
$\varepsilon_1'\le 2B_3\varepsilon_1^{\top}$, the bound established in the proof of the
remainder bound. Under
\eqref{5.5:eq-remainder-constants-d},
\begin{equation}\label{5.5:eq-remainder-constants-1}
B_3\,\varepsilon_1'\le a_0.
\end{equation}
\end{definition}

\begin{proof}[Proof of \eqref{5.5:eq-remainder-constants-1}]
Since $\max\{1,2B_3\}\ge 2B_3$, condition \eqref{5.5:eq-remainder-constants-d} gives
$2B_3\varepsilon_1^{\top}\le a_1$. Combining this with the hypothesis
$\varepsilon_1'\le 2B_3\varepsilon_1^{\top}$ of (d), we obtain
\begin{equation*}
\varepsilon_1'\le 2B_3\varepsilon_1^{\top}\le a_1 .
\end{equation*}
Multiplying by $B_3>0$ gives $B_3\varepsilon_1'\le B_3a_1$. The radii of
\cite[Theorem~1]{B-CMP102b} satisfy $B_3a_1\le a_0$, as stated in
\cite[p.~279]{B-CMP102b}. Hence $B_3\varepsilon_1'\le B_3a_1\le a_0$, which is
\eqref{5.5:eq-remainder-constants-1}.
\end{proof}

\begin{remark}\label{5.5:lem-near-unitary-transports}
The statement of this lemma is not written out here; parts of its proof are printed below.
\end{remark}

\begin{proof}[Proof of clauses (i) and (ii) of Lemma~\ref{5.5:lem-near-unitary-transports}]
\label{5.5:prf-near-unitary-proof}
The proof derives clauses (i) and (ii) from the following statements. The first is part (0) of
Lemma~\ref{5.5:lem-near-unitary-transports}, that is, the corresponding representation with the bound
$a_{\mathfrak{b}}\le (K_w+1)\gamma_0^{(k)}$. Part (0) gives this bound in two forms, since a
regularity number at one level holds with the same number at every finer level:
($\alpha$) in the length unit of $\mathfrak{B}$ on $\tilde{\square}^{4}$, and ($\beta$) in the
length unit of $\mathfrak{B}^{\star}_j$ on all of $\square_0$. The remaining inputs are these:
\begin{itemize}
\item the representation of averaged fields in a frame, together with its reality condition;
\item the axial-regauging bound;
\item the averaging representation and the averaging bound;
\item the exponential sum bound for the distance $d_{\star}$;
\item the block locality of the averaging operations, \cite[Proposition~6]{B-CMP98};
\item condition (i) of \eqref{5.5:eq-remainder-constants-b}.
\end{itemize}
The data $\mathfrak{r}_v$, $Q_v$, $Q_v^{(1)}$, $Q_v^{\sim-2}$, $p$, $t_p$ and
$[\mathfrak{r}_v]_p$ of the near unit $v$ are those of Definition~\ref{5.5:def-box-read-units},
and the staircase of $\square_0$ is the one of part (0), which lies in
$\tilde{\square}^{5}\setminus\tilde{\square}^{4}$. The constants $B_H$, $c_A$, $c_1^{b}$, $K_w$
and $K_N$ are those of Definition~\ref{5.5:def-remainder-constants}. In addition, the proof uses
the inequalities $\delta M\ge1$, $c_A\ge1$, $L\ge1$ and $M/t_p\ge2$ between these numbers.
We write $\ell_m:=L^m\eta$. For a cell
$c$ of $\mathfrak{B}^{\star}_j$ we write $l(c)$ for its level and $\tilde{B}(c)$ for its blocks.
Throughout we work in the gauge $u:=v_w^{-1}$ of part (0); in the native gauge the two end
factors $v_w$ contribute at most $\mathsf{n}_w^{2}$, where $\mathsf{n}_w:=2^{1/4}$, as stated in
clause (ii).

\emph{Clause} (i). Apply the representation of averaged fields in the frame
$(\square_0,\mathfrak{B},\mathfrak{B}^{\star}_j)$ to the field--current pair
$(\mathfrak{b}^{u},0)$, that is, to $\mathfrak{b}^{u}$ at zero current. By part (0),
$\mathfrak{b}^{u}=e^{i\eta A'_{\mathfrak{b}}}U_{00}^{w}$ with $U_{00}^{w}$ $G$-valued. The
representation therefore gives
\begin{equation}\label{5.5:eq-near-unitary-proof-1}
M_{\star}\mathfrak{b}^{u}=e^{iB_0}V_{00},\qquad V_{00}:=M_{\star}U_{00}^{w}.
\end{equation}
Here $V_{00}$ satisfies the reality condition of the representation. This follows from form
($\beta$) of part (0) together with the block locality \cite[Proposition~6]{B-CMP98}. Put
$U_{00}:=U(\mathfrak{B}^{\star}_j,V_{00})\in G$ and $\mathcal{H}_0:=\mathcal{H}(B_0)$. The output
of the representation is then
\begin{equation}\label{5.5:eq-near-unitary-proof-2}
W=U(\mathfrak{B}^{\star}_j,M_{\star}\mathfrak{b}^{u})
=\bigl[e^{i\eta\mathcal{H}_0}U_{00}\bigr]^{v_0},\qquad n(v_0)\le 2^{1/4}.
\end{equation}
The first equality is the definition of $W$ in clause (i).

We first bound $B_0$. The averaging bound, combined with block locality, gives for every cell $c$
\begin{equation}\label{5.5:eq-near-unitary-proof-3}
|B_0(c)|\le c_A\,\ell_{l(c)}\sup_{\tilde{B}(c)}|A'_{\mathfrak{b}}|\le c_A\,a_{\mathfrak{b}},
\end{equation}
where the second inequality is form ($\beta$) of part (0). Suppose moreover that
$\tilde{B}(c)\subset Q_v^{(1)}$. The cells of $\mathfrak{B}$ there have level at least $p$, so
form ($\alpha$) gives $|A'_{\mathfrak{b}}|\le a_{\mathfrak{b}}/\ell_p$ on $\tilde{B}(c)$.
Since $l(c)\le j$, it follows that
\begin{equation}\label{5.5:eq-near-unitary-proof-4}
|B_0(c)|\le c_A\,t_p\,a_{\mathfrak{b}}\qquad\text{whenever }\tilde{B}(c)\subset Q_v^{(1)}.
\end{equation}

We next bound $\mathcal{H}_0$. Let $P$ be a point of $\mathfrak{r}_v$. We have
$\mathfrak{r}_v\subset Q_v^{\sim-2}\subset\tilde{\square}^{3}$. On this set the cells of
$\mathfrak{B}$ have levels at least $p>j$, the set $\mathfrak{B}^{\star}_j$ is capped at $j$, and
there is no staircase; hence $l(P)=j$. The axial-regauging bound, applied on the segment $[0,B_0]$
at $P$, gives
\begin{equation}\label{5.5:eq-near-unitary-proof-5}
\ell_j|\mathcal{H}_0(P)|\le B_H\sum_{c}e^{-\delta d_{\star}(P,c)}|B_0(c)|.
\end{equation}
We split the sum into two groups.

The first group consists of the cells $c$ with $\tilde{B}(c)\subset Q_v^{(1)}$. By
\eqref{5.5:eq-near-unitary-proof-4} and the exponential sum bound, these cells contribute at most
$B_Hc_Ac_1^{b}t_pa_{\mathfrak{b}}$.

The second group consists of all other cells. For such a cell $c$, some block of $c$ meets
$(Q_v^{(1)})^{c}$. Now $Q_v^{(1)}$ contains three layers of $\pi_p$-cubes, each of side
$(M/t_p)\ell_j$, around $[\mathfrak{r}_v]_p\supseteq\mathfrak{r}_v\ni P$. Therefore
\begin{equation*}
d_{\star}(P,c)\ge \frac{3M}{t_p}-2\ge\frac{2M}{t_p}.
\end{equation*}
Write
$e^{-\delta d_{\star}}=e^{-\frac34\delta d_{\star}}e^{-\frac14\delta d_{\star}}$ and use
\eqref{5.5:eq-near-unitary-proof-3}. These cells then contribute at most
\begin{equation}\label{5.5:eq-near-unitary-proof-6}
B_Hc_A\,a_{\mathfrak{b}}\,e^{-\frac32\delta M/t_p}\sum_{c}e^{-\frac14\delta d_{\star}(P,c)}
\le B_Hc_Ac_1^{b}\,t_p\,a_{\mathfrak{b}},
\end{equation}
with the remaining sum over $c$ bounded by $c_1^{b}$. For the exponential factor, put
$x:=\frac32\delta M/t_p$. Since $\delta M\ge1$, we have $x\ge1/t_p$, and hence
$e^{-x}\le 1/x\le t_p$.

Adding the two groups gives
$\ell_j|\mathcal{H}_0(P)|\le 2B_Hc_Ac_1^{b}t_pa_{\mathfrak{b}}$. Since $L\ge1$, this yields
\begin{equation}\label{5.5:eq-near-unitary-proof-7}
\ell_j|\mathcal{H}_0|\le 2LB_Hc_Ac_1^{b}\,t_p\,a_{\mathfrak{b}}\qquad\text{on }\mathfrak{r}_v;
\end{equation}
the factor $L$ is retained only to match the form of clause (i).
Together with \eqref{5.5:eq-near-unitary-proof-2}, this is clause (i).

\emph{Clause} (ii). The averaging representation and the averaging bound rest on block locality
alone, \cite[Proposition~6]{B-CMP98}. They therefore hold for the averages $M_{\cdot}$ of
$\mathfrak{B}_j(\mathfrak{r}_v)$. Apply them to \eqref{5.5:eq-near-unitary-proof-2}. For every
bond $c$ of $\mathfrak{B}_j(\mathfrak{r}_v)$, with initial point $c_-$ and final point $c_+$, this
gives
\begin{equation*}
\mathsf{V}'_u(c)=v_0(c_-)\,e^{i\mathfrak{h}(c)}M_{\cdot}(U_{00})(c)\,v_0(c_+)^{-1},
\qquad M_{\cdot}(U_{00})(c)\in G.
\end{equation*}
The $1$-form piece $\mathfrak{h}(c)$ is bounded, using clause (i) in the form
\eqref{5.5:eq-near-unitary-proof-7}, by
\begin{equation}\label{5.5:eq-near-unitary-proof-8}
|\mathfrak{h}(c)|\le c_A\,\ell_{l(c)}\sup_{\text{blocks of }c}|\mathcal{H}_0|
\le 2Lc_AB_Hc_Ac_1^{b}\,t_p\,a_{\mathfrak{b}}\cdot\frac{\ell_{l(c)}}{\ell_j}.
\end{equation}

Let $\Gamma$ be a path as in clause (ii), with consecutive bonds $c(1),\dots,c(\mathsf{q})$. Put
$\mathsf{g}_i:=M_{\cdot}(U_{00})(c(i))\in G$ and $\mathsf{p}_i:=\mathsf{g}_1\cdots\mathsf{g}_i$,
with $\mathsf{p}_0:=1$. Along $\Gamma$ the factors $v_0$ at interior points
telescope. The $G$-valued factors can then be moved to the right, using
$\mathsf{g}\,e^{iY}=e^{i\operatorname{Ad}_{\mathsf{g}}Y}\,\mathsf{g}$ for $\mathsf{g}\in G$ and
$Y\in\mathfrak{g}^{c}$. This gives
\begin{equation*}
\mathsf{V}'_u(\Gamma)=v_0(c(1)_-)
\Bigl[\prod_{i=1}^{\mathsf{q}}e^{i\operatorname{Ad}_{\mathsf{p}_{i-1}}\mathfrak{h}(c(i))}\Bigr]
\mathsf{p}_{\mathsf{q}}\,v_0(c(\mathsf{q})_+)^{-1}.
\end{equation*}
Conjugation by an element of $G$ preserves $|\cdot|$. Moreover, the non-unitarity number $n$ is
submultiplicative, satisfies $n(g^{-1})=n(g)$, equals $1$ on $G$, and obeys
$n(e^{iY})\le e^{2|Y|}$. Using in addition $n(v_0)\le2^{1/4}$ and
\eqref{5.5:eq-near-unitary-proof-8}, we obtain
\begin{equation*}
\begin{split}
n(\mathsf{V}'_u(\Gamma))
&\le n(v_0)^2\exp\Bigl(2\sum_{c\in\Gamma}|\mathfrak{h}(c)|\Bigr)\\
&\le 2^{1/2}\exp\bigl(2\lambda_v\cdot 2Lc_AB_Hc_Ac_1^{b}\,t_p\,a_{\mathfrak{b}}\bigr)
\le 2^{1/2}e^{2K_Na_{\mathfrak{b}}}.
\end{split}
\end{equation*}
In the second line we used that $\sum_{c\in\Gamma}\ell_{l(c)}/\ell_j$ is the level-$j$ length of
$\Gamma$, which is at most $\lambda_v$. In the last step the factor $t_p$ cancels, because
$\lambda_vt_p=6dM$:
\begin{equation*}
2\lambda_v\cdot 2Lc_AB_Hc_Ac_1^{b}\,t_p\,a_{\mathfrak{b}}
=24\,dLM\,B_Hc_A^{2}c_1^{b}\,a_{\mathfrak{b}}\le 2K_Na_{\mathfrak{b}},
\end{equation*}
the inequality being the definition of $K_N$ in Definition~\ref{5.5:def-remainder-constants}.

Let $\gamma\subset\mathfrak{r}_v$ be a $T_\eta$-path of level-$j$ length at most $\lambda_v$. It
has at most $\lambda_v\ell_j/\eta$ fine bonds $b$, and on each of them
\begin{equation*}
W(b)=v_0(b_-)\,e^{i\eta\mathcal{H}_0(b)}U_{00}(b)\,v_0(b_+)^{-1}.
\end{equation*}
The same telescoping and conjugation, now with the $G$-valued factors $U_{00}(b)$, give
\begin{equation*}
\begin{split}
n(W(\gamma))
&\le 2^{1/2}\exp\Bigl(2\,\frac{\lambda_v\ell_j}{\eta}\,\eta\sup_{\mathfrak{r}_v}|\mathcal{H}_0|\Bigr)
=2^{1/2}\exp\Bigl(2\lambda_v\ell_j\sup_{\mathfrak{r}_v}|\mathcal{H}_0|\Bigr)\\
&\le 2^{1/2}\exp\bigl(2\lambda_v\cdot 2LB_Hc_Ac_1^{b}\,t_p\,a_{\mathfrak{b}}\bigr)
\le 2^{1/2}e^{2K_Na_{\mathfrak{b}}}.
\end{split}
\end{equation*}
Here the second line uses \eqref{5.5:eq-near-unitary-proof-7}, and the last step uses $c_A\ge1$
and the previous estimate.

Finally, combine part (0), the inequality $\gamma_0^{(k)}\le\gamma_0^{\top}$, and condition (i)
of \eqref{5.5:eq-remainder-constants-b}:
\begin{equation*}
2K_Na_{\mathfrak{b}}\le 2K_N(K_w+1)\gamma_0^{\top}\le 0.02.
\end{equation*}
Since $2^{1/2}e^{0.02}\le 1.45$, both transports satisfy the bound of clause (ii).
\end{proof}

\begin{definition}[The working gauge and transfer of estimates]\label{5.5:not-working-gauge}
Let $\square_0 := \tilde{\square}^{5}$ be the block $\square$ enlarged by five layers of cubes
of the partition $\pi_{n_0}$, and let $\mathbb{W}$ be the box field on $\square_0$. For
$g\in G^{c}$ write $n(g):=\|g\|\,\|g^{-1}\|$.

\begin{enumerate}
\item[(a)] \emph{The working gauge.} By clause (0) of the lemma on near-unitary transports on
the own box (Lemma~\ref{5.5:lem-near-unitary-transports}), there is a block-local
$G^{c}$-valued gauge transformation $v_w$ on $\square_0$ such that
\begin{equation}\label{5.5:eq-working-gauge-1}
  \mathbb{W} = \bigl[e^{i\eta\mathcal{H}^{w}}U_{00}^{w}\bigr]^{v_w},
  \qquad
  n(v_w(x)) \le \mathsf{n}_w := 2^{1/4},
\end{equation}
with $U_{00}^{w}$ $G$-valued. We put $u := v_w^{-1}$, and we let $\bar u$ be the restriction
of $u$ to the base points of the blocks, extended to a block-constant function. Objects
obtained by applying the gauge transformation $u$ are called \emph{working} objects and are
marked by a superscript or subscript $u$.

\item[(b)] \emph{Transfer.} Every sup-norm estimate in the proofs of this section is derived
for the working objects and transferred to the native objects as follows: $1$-forms and the
linearised maps $Q_{\star}$ are moved by $\operatorname{Ad}_{u(\cdot)}$, in accordance with
the transformation rule for $1$-forms and linearised maps under gauge transformations and with
\cite{B-CMP98}; loop data at a point $y$ are moved by $\operatorname{Ad}_{\bar u(y)}$. Each
such transfer costs a factor at most $1.2$ on $|\cdot|$. These factors are included in the
constants $K_1$ and $K^{\mathrm{rem}}$ and in the conditions
\eqref{5.5:eq-remainder-constants-c}.

\item[(c)] \emph{Gauge-blind objects.} The term $T$, membership in the set $\mathfrak{U}_T$
(a union of orbits of all $G^{c}$-valued gauge transformations), and the orbits of
$(U_j,J_j)(\mathfrak{r}_v,e^{iB})$ under constant conjugation of $B$ do not depend on the gauge.

\item[(d)] \emph{Native objects.} The loop data $B^{0}$, $B^{1}$ and the loop increment
$\Delta$, the difference $G_v$ of the near term at $\mathsf{P}^{1}$ and at $\mathsf{P}^{0}$,
the box fields $U_{\mathsf{P}^{\iota}}$, and their averages $\mathsf{V}'$ and $\mathsf{V}'^{1}$
are the native (canonical) objects. No superscript or subscript $u$ is attached to them.
\end{enumerate}
\end{definition}

\begin{proof}
We verify the two assertions contained in (b) and (c).

For (b), let $g\in G^{c}$. Then
$n(g^{-1}) = \|g^{-1}\|\,\|g\| = n(g)$. By \eqref{5.5:eq-working-gauge-1} we therefore have
$n(u(x)) = n(v_w(x)) \le 2^{1/4}$ at every point $x$ of $\square_0$. The same bound holds for
$\bar u(y)$, since $\bar u(y)$ is a value of $u$ at a base point.

Next, $\operatorname{Ad}_g X = gXg^{-1}$. Since $|gXg^{-1}| \le \|g\|\,|X|\,\|g^{-1}\|$ for
the norm $|\cdot|$ on $\mathfrak{g}^{c}$-valued data and the norm $\|\cdot\|$ on $G^{c}$ used
in the definition of $n(g)$, we have $\|\operatorname{Ad}_g\|\le n(g)$, whence
\begin{equation}\label{5.5:eq-working-gauge-2}
  |\operatorname{Ad}_g X| \le n(g)\,|X| .
\end{equation}
Applying \eqref{5.5:eq-working-gauge-2} with $g=u(x)$ or $g=\bar u(y)$ shows that each
transfer multiplies $|\cdot|$ by at most $2^{1/4}$. Finally, $2^{1/4} \le 1.2$, because
$1.2^{4} = 2.0736 \ge 2$.

For (c), membership in $\mathfrak{U}_T$ is the same for a configuration and for its transform
by $u$, because $\mathfrak{U}_T$ is a union of orbits of all $G^{c}$-valued gauge
transformations; this follows from the definition of the configuration classes in
\cite[(1.10)]{B-CMP109}. That the term $T$ takes the same value on a configuration and on its
transform by $u$, and that the orbits of $(U_j,J_j)(\mathfrak{r}_v,e^{iB})$ under constant
conjugation of $B$ are unchanged, follow from the gauge transformation formula
\cite[(181)]{B-CMP102b}.
\end{proof}

\begin{remark}\label{5.5:lem-loop-data}
The statement of this lemma is not written out here; parts of its proof are printed below.
\end{remark}

\begin{proof}[Proof of Lemma~\ref{5.5:lem-loop-data}, continued]
\label{5.5:prf-correction-response}
\emph{Response of the box field to the correction.}
This part of the proof is an implication from two statements of this chapter, which enter it as
hypotheses and are used as stated:
\begin{itemize}
\item[(a)] the response estimate of the closure step, which bounds the response of a box field to a
correction form and represents the corrected field as a gauge transform of the uncorrected one;
\item[(b)] the transport lemma whose constants $r_V$, $C_V$ enter
Definition~\ref{5.5:def-remainder-constants}: its clause on the regauging attached to block data,
together with the bound on the regauging obtained in the proof of that clause.
\end{itemize}
The response estimate (a) takes as datum a pair consisting of a box field and a correction form; we
write $\mathfrak{h}^{\natural}$ for the correction form. Its proof reads $\mathfrak{h}^{\natural}$ only through
$\sup|\mathfrak{h}^{\natural}|$ and through the kernel sum $\mathsf{T}$. We apply (a) with the
datum $(\mathbb{W},\mathfrak{h}^{\natural})$ replaced by $(\mathfrak{b}^{u},\mathsf{E}^{u})$, in the
frame $(\square_0,\mathfrak{B},\mathfrak{B}^{\star}_j)$. At this datum the kernel sum is the sum
$\mathsf{T}$ of Lemma~\ref{5.5:lem-loop-data}.

We check the hypotheses of (a) at this datum. The decomposition hypothesis of (a) asks that the
field be written as $e^{i\eta A'}\bar U$ with $\bar U$ $G$-valued. For $\mathfrak{b}^{u}$ this is
the representation $\mathfrak{b}^{u}=e^{i\eta A'_{\mathfrak{b}}}U^{w}_{00}$, with $U^{w}_{00}$
$G$-valued, of the corresponding display in
Lemma~\ref{5.5:lem-near-unitary-transports}. The segment condition of (a) at this datum reads
\begin{equation}\label{5.5:eq-correction-response-1}
|B_0|+1.1\,\sup|\mathsf{E}^{u}|\le c_A a_{\mathfrak{b}}+0.23\,b_H<b_H .
\end{equation}
Here $a_{\mathfrak{b}}$ is bounded by $(K_w+1)\gamma_0^{(k)}$ in
the corresponding display, and $\sup|\mathsf{E}^{u}|$ is bounded in
the corresponding display; the conditions under which \eqref{5.5:eq-correction-response-1} holds,
and the hypothesis of the regauging estimate (b), are among the conditions
\eqref{5.5:eq-remainder-constants-c} of Definition~\ref{5.5:def-remainder-constants}.

To state the output, put
\begin{equation}\label{5.5:eq-correction-response-2}
U^{u}_{\mathsf{P}^{1}}:=U\bigl(\mathfrak{B}^{\star}_j,\,e^{i\mathsf{E}^{u}}M_{\star}\mathfrak{b}^{u}\bigr).
\end{equation}
By the transformation rules of $1$-forms and averages under gauge transformations recorded in
Notation~\ref{5.5:not-working-gauge}, and by \cite[(181)]{B-CMP102b}, this field equals
$(U_{\mathsf{P}^{1}})^{\bar u}$. Here $U_{\mathsf{P}^{1}}$ is the $U$-component of
$\mathsf{P}^{1}$, taken in the native gauge. We further put
\begin{equation}\label{5.5:eq-correction-response-5}
\mathsf{V}'^{1}_u:=M_{\cdot}(U^{u}_{\mathsf{P}^{1}})\restriction\mathfrak{B}_j(\mathfrak{r}_v),
\end{equation}
which equals $(\mathsf{V}'^{1})^{\bar u}$, where $\bar u$ is restricted to
$\mathfrak{B}_j(\mathfrak{r}_v)$. Following the definitions of (a), we also put
\begin{equation}\label{5.5:eq-correction-response-6}
e^{iB_1}:=e^{i\mathsf{E}^{u}}e^{iB_0},\qquad \tilde v:=v(B_1)\,v(B_0)^{-1}.
\end{equation}
Here $v(B_0)$ and $v(B_1)$ are the regaugings attached to the block data $B_0$ and $B_1$.

Let $W$ be the box field on $\mathfrak{B}^{\star}_j$ of
Lemma~\ref{5.5:lem-near-unitary-transports}. Let $x\in\mathfrak{r}_v$, so that the level of $x$ in
$\mathfrak{B}^{\star}_j$ is the cap $j$, and let
$y_0(x)\in\mathfrak{r}_v$ be the base point of the block of $x$. Then (a) yields a $1$-form
$\mathcal{H}_S$ such that the following hold, the last line being proved below:
\begin{equation}\label{5.5:eq-correction-response-a}
\begin{aligned}
U^{u}_{\mathsf{P}^{1}}&=\bigl[e^{i\eta\mathcal{H}_S}W\bigr]^{\tilde v},\\
\ell_j|\mathcal{H}_S|,\ \ell_j^{2}|\nabla_W\mathcal{H}_S|,\
\ell_j^{3}|D^{\star}_WD_W\mathcal{H}_S|
&\le 8L^{3}B_H\,\mathsf{T}(y_0(x))\le h_v,\\
n(\tilde v(x))-1&\le\nu_S(x):=\frac{18B_H\,\mathsf{T}(y_0(x))}{r_V}\le\frac{9h_v}{4L^{3}r_V},\\
\|\operatorname{Ad}_{\tilde v(x)}-1\|&\le\nu_S(x).
\end{aligned}
\end{equation}
In the second line, the bound by $h_v$ holds because $y_0(x)\in\mathfrak{r}_v$ and
$h_v=8L^{3}B_H\sup_{\mathfrak{r}_v}\mathsf{T}$ by Lemma~\ref{5.5:lem-loop-data}. In the third line,
the last inequality is the same bound rewritten: $18B_H\mathsf{T}(y_0(x))/r_V$ equals
$\frac{9}{4L^{3}r_V}\cdot 8L^{3}B_H\mathsf{T}(y_0(x))$, which is at most $9h_v/(4L^{3}r_V)$.

We now prove the last line of \eqref{5.5:eq-correction-response-a}. The proof of the regauging
estimate (b) gives
\begin{equation}\label{5.5:eq-correction-response-3}
|\tilde v(x)-1|\le\varepsilon_x
:=\frac{20}{3r_V}\,\ell_j\sup_{\mathfrak{Q}_x}|\mathcal{H}(B_1)-\mathcal{H}(B_0)|\le 0.09 ,
\end{equation}
where $\mathfrak{Q}_x$ denotes the set over which the proof of that clause takes this supremum at
the point $x$.

Let $\|\cdot\|$ be the operator norm. For every matrix $Y$ we have the identity
\begin{equation*}
\tilde v(x)Y\tilde v(x)^{-1}-Y=(\tilde v(x)-1)Y\tilde v(x)^{-1}+Y(\tilde v(x)^{-1}-1).
\end{equation*}
Moreover $\tilde v(x)^{-1}-1=\tilde v(x)^{-1}(1-\tilde v(x))$, so
\begin{equation*}
\|\tilde v(x)^{-1}-1\|\le\|\tilde v(x)^{-1}\|\,\|\tilde v(x)-1\|\le\varepsilon_x\|\tilde v(x)^{-1}\|.
\end{equation*}
Hence, by the triangle inequality, $\|\tilde v(x)^{-1}\|\le 1+\varepsilon_x\|\tilde v(x)^{-1}\|$.
Since $\varepsilon_x<1$, this gives $\|\tilde v(x)^{-1}\|\le(1-\varepsilon_x)^{-1}$, and then
$\|\tilde v(x)^{-1}-1\|\le\varepsilon_x/(1-\varepsilon_x)$. Inserting these two bounds into the
identity gives
\begin{equation*}
\begin{aligned}
\|\tilde v(x)Y\tilde v(x)^{-1}-Y\|
&\le\frac{\varepsilon_x}{1-\varepsilon_x}\|Y\|+\frac{\varepsilon_x}{1-\varepsilon_x}\|Y\|\\
&=\frac{2\varepsilon_x}{1-\varepsilon_x}\|Y\|.
\end{aligned}
\end{equation*}
By $\varepsilon_x\le 0.09$ we have $2/(1-\varepsilon_x)\le 2/0.91\le 2.2$. Also
$2.2\cdot\frac{20}{3}\le 16$. Therefore
\begin{equation}\label{5.5:eq-correction-response-4}
\begin{aligned}
\|\operatorname{Ad}_{\tilde v(x)}-1\|
&\le\frac{2\varepsilon_x}{1-\varepsilon_x}\le 2.2\,\varepsilon_x\\
&\le\frac{16}{r_V}\,\ell_j\sup_{\mathfrak{Q}_x}|\mathcal{H}(B_1)-\mathcal{H}(B_0)|\le\nu_S(x).
\end{aligned}
\end{equation}
The last inequality in \eqref{5.5:eq-correction-response-4} is the closing inequality of the
response estimate (a). This proves the last line of \eqref{5.5:eq-correction-response-a}.

The form $\mathcal{H}_S$ and the gauge $\tilde v$ depend on the decomposition
$\mathbb{U}=e^{i\eta A'}\bar U$. They are used only at a fixed value of the parameters, as bounds,
and no holomorphy in the parameters is required of them.
\end{proof}

\begin{proof}[Proof of Lemma~\ref{5.5:lem-loop-data}, continued: charts of the two box fields by
analytic continuation]
\label{5.5:prf-box-field-charts}
This part of the proof establishes clause~(i) of Lemma~\ref{5.5:lem-loop-data} for both box pairs
$\mathsf{P}^{0}$ and $\mathsf{P}^{1}$, and the pointwise inequality of
its clause~(ii). It is an implication from the following statements, each used as stated:
\begin{itemize}
\item the statement on the chart and the comb datum of the first box field on an own box;
\item clauses (H), (A2), (A3), $(\gamma2)$ and the response bound of the closure estimates, and
the list $C^{b}$ of those estimates;
\item clause (V3) of the transport statement, the statement whose clause (V2) supplies the
constants $r_V$, $C_V$ of Definition~\ref{5.5:def-remainder-constants};
\item the transfer of the analysis of \cite[\S3]{B-CMP109} to the cover, printed without proof in
\cite[p.~276]{B-CMP119}; cited as printed.
\end{itemize}
The constants and smallness conditions are those of Definition~\ref{5.5:def-remainder-constants},
in particular \eqref{5.5:eq-remainder-constants-c} and \eqref{5.5:eq-remainder-constants-d}; the
data $y$, $n$, $t$, $\Theta$, $c_W$ and $|B^{ax}(b)|_{\star}$ of the box-read unit $v$ are those of
Definition~\ref{5.5:def-box-read-units}.

\emph{The first box pair.} We apply the chart statement for the first box field with its region,
its level and its box taken to be $Q_v$, $p$ and $\mathfrak{r}_v$. The proof of that statement does
not use the inductive bound on the remainders, so that the appeal to it here is not circular. It
gives two conclusions. First,
$\mathsf{P}^{0}\lceil\mathfrak{r}_v=(U_j,J_j)(\mathfrak{r}_v,\mathsf{V}')$ modulo gauge, the
identification being provided by the transfer of the analysis of \cite[\S3]{B-CMP109} to the
cover. Second, in the comb gauge centred at $y$ the datum $B^{ax}$ of $\mathsf{V}'$ obeys, on every
bond $b$,
\begin{equation}\label{5.5:eq-box-field-charts-1}
|B^{ax}(b)|\le 2d\,(1+|b-y|)\,c_1\,\mathsf{a}_{Q_v}\,t_p^{2}.
\end{equation}
This comb datum is $e^{iB^{0}}$. Indeed, let $\varpi$ be the comb gauge of $\mathsf{V}'$ from $y$,
defined on the base points by $\varpi(x):=\mathsf{V}'(\Gamma_{y,x})^{-1}$; then $\varpi(y)=1$, since
$\Gamma_{y,y}$ is the trivial path.
The transformed field at a bond $b$ combines $\mathsf{V}'(b)$ with the transports of $\mathsf{V}'$
along the comb paths $\Gamma_{y,b_-}$ and $\Gamma_{b_+,y}$, that is, with the transport along the
loop $\lambda_b=\Gamma_{y,b_-}\,b\,\Gamma_{b_+,y}$; by the definition of $B^{0}$ in
Lemma~\ref{5.5:lem-loop-data},
\begin{equation*}
\mathsf{V}'^{\,\varpi}(b)=\mathsf{V}'(\lambda_b)=e^{iB^{0}(b)} .
\end{equation*}
By \cite[(181)]{B-CMP102b},
$U_j(\mathfrak{r}_v,\mathsf{V}'^{\,\varpi})=U_j(\mathfrak{r}_v,\mathsf{V}')^{\bar\varpi}$, where
$\bar\varpi$ is the block-constant extension of $\varpi$ from the base points to the fine lattice;
the currents
transform by the representation $R(\bar\varpi)$, \cite[(1.10)]{B-CMP109}. Hence
$(U_j,J_j)(\mathfrak{r}_v,e^{iB^{0}})$ lies in the $G^{c}$-orbit of
$(U_j,J_j)(\mathfrak{r}_v,\mathsf{V}')$, and therefore in that of
$\mathsf{P}^{0}\lceil\mathfrak{r}_v$: this is clause~(i) for $\mathsf{P}^{0}$. Moreover
$B^{ax}=B^{0}$, so \eqref{5.5:eq-box-field-charts-1} is the pointwise inequality of clause~(ii) for
$B^{0}$ with the factor $1$ in place of $1.2$. Both conclusions are thus exactly those of the chart
statement for the first box field.

\emph{The pointwise inequality.} We use the two relations $\mathsf{a}_{Q_v}\le 2K_{\mathrm{loc}}
w^{\vee}(y)\alpha_{0,n}$ and $t_p=Lt$ between the data of the box-read unit $v$ of
Definition~\ref{5.5:def-box-read-units}. They give
\begin{equation*}
\begin{split}
2d\,(1+|b-y|)\,c_1\,\mathsf{a}_{Q_v}\,t_p^{2}
&\le 2d\,(1+|b-y|)\cdot 2c_1K_{\mathrm{loc}}L^{2}w^{\vee}(y)\,\alpha_{0,n}\,t^{2}\\
&\le L^{-1}|B^{ax}(b)|_{\star}.
\end{split}
\end{equation*}
For the second inequality: the definition of $c_W$ in Definition~\ref{5.5:def-box-read-units}
gives $c_W\ge 2c_1K_{\mathrm{loc}}L^{3}\alpha_{0,n}/\mathfrak{a}_n$, hence
\begin{equation*}
2c_1K_{\mathrm{loc}}L^{2}w^{\vee}(y)\,\alpha_{0,n}\le c_W w^{\vee}(y)\,\mathfrak{a}_n/L
=\Theta\,\mathfrak{a}_n/L,
\end{equation*}
and $|B^{ax}(b)|_{\star}=2d\,(1+|b-y|)\,\Theta\,\mathfrak{a}_n t^{2}$. Together with
\eqref{5.5:eq-box-field-charts-1} and $B^{ax}=B^{0}$ this gives
\begin{equation*}
|B^{0}(b)|\le 2d\,(1+|b-y|)\,c_1\,\mathsf{a}_{Q_v}\,t_p^{2}\le L^{-1}|B^{ax}(b)|_{\star}.
\end{equation*}
The working
datum is $B^{0}_u=\operatorname{Ad}_{\bar v_w(y)}^{-1}B^{0}$, and passing from the native to the
working datum costs at most a factor $1.2$ (Notation~\ref{5.5:not-working-gauge}). Hence both
$B^{0}$ and $B^{0}_u$ satisfy the pointwise line of the corresponding display.

The factor $1.2$ may also be placed on the other side. If the display of the chart statement is
read in the gauge, call it $\mathsf{g}$, in which the chart statement is derived and in which
$\operatorname{Ad}_{\mathsf{g}}$ costs at most a
factor $2$, then the canonical $B^{0}$ is the $\operatorname{Ad}_{\mathsf{g}(y)}$-conjugate of the
datum in
that display. The factor $1.2$ then falls on the native datum $B^{0}$, while the working datum
$B^{0}_u$ obeys the bound with factor $1$, by the rim bound in the chart statement and the chain of
factors $1.1\cdot1.45\cdot1.2\le2$ in part~(ii) of the lemma on near-unitary transports on the own
box. Since the corresponding display carries the factor $1.2$ for both $B^{0}$ and $B^{0}_u$, and
$|B^{0}|_{\star}$ bounds both, it holds under either reading.

\emph{The second box pair: the box field as a function of its data.} We work in the working gauge.
By the proof of part~(i) of the lemma on near-unitary transports on the own box
(proof~\ref{5.5:prf-near-unitary-proof}),
$M_{\star}\mathfrak{b}^{u}=e^{iB_0}V_{00}$. In the derivation of
\eqref{5.5:eq-correction-response-a} we put $e^{iB_1}:=e^{i\mathsf{E}^{u}}e^{iB_0}$ and
$U^{u}_{\mathsf{P}^{1}}:=U(\mathfrak{B}^{\star}_j,e^{i\mathsf{E}^{u}}M_{\star}\mathfrak{b}^{u})$.
Put $r_{\mathcal{O}}:=2\max\{\sup|B_0|,\sup|B_1|\}$, and let $\mathcal{O}$ be the set of
$\mathfrak{g}^{c}$-valued
$B$ on the bonds carrying the data of $\mathfrak{B}^{\star}_j$ with $|B|_{\infty}<r_{\mathcal{O}}$;
thus
$B_0,B_1\in\mathcal{O}$. For $B\in\mathcal{O}$ we put
\begin{equation}\label{5.5:eq-box-field-charts-a}
\begin{gathered}
Y(B):=U\bigl(\mathfrak{B}^{\star}_j,e^{iB}V_{00}\bigr),\qquad
\Phi(B):=M_{\cdot}\bigl(Y(B)\bigr)\lceil\mathfrak{B}_j(\mathfrak{r}_v),\qquad
Z(B):=U_j\bigl(\mathfrak{r}_v,\Phi(B)\bigr),\\
g_2(B)(x):=\text{the quotient of the transports } Z(B)(\gamma_{z,x})
\text{ and } Y(B)(\gamma_{z,x}),
\end{gathered}
\end{equation}
where in the last line $x$ lies in the block of the base point $z$, $\gamma_{z,x}$ is the path
from $z$ to $x$ in that block, and the quotient is taken in the order which makes $g_2(B)$ the
gauge transformation with $g_2(B)=1$ at base points and
$Z(B)^{g_2(B)}(\gamma_{z,x})=Y(B)(\gamma_{z,x})$. That $Z(B)$ and $g_2(B)$ are well defined on
$\mathcal{O}$ is shown in the paragraph on analyticity below. Since the product is taken
bondwise,
\begin{equation*}
\begin{gathered}
W=U(\mathfrak{B}^{\star}_j,e^{iB_0}V_{00})=Y(B_0),\\
U^{u}_{\mathsf{P}^{1}}=U(\mathfrak{B}^{\star}_j,e^{i\mathsf{E}^{u}}e^{iB_0}V_{00})
=U(\mathfrak{B}^{\star}_j,e^{iB_1}V_{00})=Y(B_1).
\end{gathered}
\end{equation*}

The radius satisfies $r_{\mathcal{O}}\le r^{\top}$, the radius of
\eqref{5.5:eq-remainder-constants-d}. Indeed,
$\sup|B_0|\le c_A a_{\mathfrak{b}}$ by clause (A3) of the closure estimates and
the corresponding display; and, by the response bound of the closure estimates
and the corresponding display,
\begin{equation*}
\sup|B_1-B_0|\le 1.1\,\sup|\mathsf{E}^{u}|\le 33B_{\sharp}^{2}\sup_m\alpha_{\cdot,m}.
\end{equation*}
Since $\sup|B_1|\le\sup|B_0|+\sup|B_1-B_0|$, it follows that
\begin{equation*}
r_{\mathcal{O}}\le 2\bigl(c_A a_{\mathfrak{b}}+33B_{\sharp}^{2}\sup_m\alpha_{\cdot,m}\bigr)
\le r^{\top},
\end{equation*}
the last inequality being part of \eqref{5.5:eq-remainder-constants-d}.

\emph{Analyticity on $\mathcal{O}$.} On $\mathcal{O}$ the box field has the representation
$Y(B)=\bigl[e^{i\eta\mathcal{H}(B)}U_{00}\bigr]^{v(B)}$; here $\mathcal{H}(B)$ and $v(B)$ (a
function of $B$, not to be confused with the unit $v$) are the small one-form and the regauging in
the representation of the box field at data $e^{iB}V_{00}$ used in the proof of part~(i) of the
lemma on near-unitary transports on the own box, so that $\mathcal{H}(B_0)=\mathcal{H}_0$ and
$v(B_0)=v_0$. The map $Y$ is analytic on $\mathcal{O}$,
by \cite[Proposition~9]{B-CMP102b}, clause (H) of the closure estimates and
clause (V3) of the transport statement. By clauses (A2), (A3), (H) and $(\gamma2)$ of the closure
estimates,
\begin{equation*}
\Phi(B)=\bigl(e^{i\mathfrak{h}_B}M_{\cdot}U_{00}\bigr)^{v(B)|},\qquad
|\mathfrak{h}_B|\le c_A B_H c_1^{b}\, r_{\mathcal{O}}<b_H,
\end{equation*}
where $v(B)|$ is $v(B)$ on the base points of $\mathfrak{B}_j(\mathfrak{r}_v)$; the last inequality
follows from $r_{\mathcal{O}}\le r^{\top}$ and the smallness condition
\eqref{5.5:eq-remainder-constants-c}, which uses from the list $C^{b}$ of the closure estimates
the inequality $c_AC_{\mathrm{sph}}\ge 2c_A^{2}L$ between their constants $c_A$ and
$C_{\mathrm{sph}}$. Thus
$\Phi(B)$ lies on a $G^{c}$-orbit contained in the domain of the box function
$U_j(\mathfrak{r}_v,\cdot)$ \cite[(181)]{B-CMP102b}.
Consequently $Z(B)$ and $g_2(B)$ are analytic on $\mathcal{O}$.

\emph{The identity at $\mathfrak{g}$-valued data.} Let $B\in\mathcal{O}$ be $\mathfrak{g}$-valued.
Then the data $e^{iB}V_{00}$ are $G$-valued, $V_{00}$ being $G$-valued by part~(i) of the lemma on
near-unitary transports on the own box, and $Y(B)$ is the minimiser itself \cite{B-CMP102b}, in the
space of \cite[(8)]{B-CMP102b}. A variation supported inside $\mathfrak{r}_v$ which keeps
$\Phi(B)$ keeps the data of $\mathfrak{B}^{\star}_j$, by part $(\beta)$ of the chart statement.
Hence $Y(B)\lceil\mathfrak{r}_v$ is critical under its own
$\mathfrak{B}_j(\mathfrak{r}_v)$-data $\Phi(B)$. Now \cite[Theorem~1]{B-CMP102b} itself, for
$G$-valued configurations and data in its small-field domain, the smallness being provided by
\eqref{5.5:eq-remainder-constants-d}, gives
\begin{equation}\label{5.5:eq-box-field-charts-2}
Y(B)\lceil\mathfrak{r}_v=Z(B)^{g_2(B)}.
\end{equation}

\emph{Continuation.} The ball $\mathcal{O}$ is convex, hence connected, and its
$\mathfrak{g}$-valued points form a totally real slice of maximal dimension, since
$\mathfrak{g}^{c}=\mathfrak{g}\otimes\mathbb{C}$. Both members of \eqref{5.5:eq-box-field-charts-2}
are analytic on $\mathcal{O}$ and agree on this slice; by the identity theorem for holomorphic
functions of several variables they agree on all of $\mathcal{O}$, in particular at $B_1$ (and at
$B_0$). Since $\mathsf{V}'^{1}_u=M_{\cdot}(U^{u}_{\mathsf{P}^{1}})\lceil
\mathfrak{B}_j(\mathfrak{r}_v)=\Phi(B_1)$, this reads
\begin{equation*}
U^{u}_{\mathsf{P}^{1}}\lceil\mathfrak{r}_v
=U_j\bigl(\mathfrak{r}_v,\mathsf{V}'^{1}_u\bigr)^{g_2(B_1)} .
\end{equation*}
The corresponding identity for the currents follows from \cite[(1.10)]{B-CMP109}. In the
derivation of \eqref{5.5:eq-correction-response-a},
$U^{u}_{\mathsf{P}^{1}}=(U_{\mathsf{P}^{1}})^{\bar u}$ and
$\mathsf{V}'^{1}_u=(\mathsf{V}'^{1})^{\bar u|}$; applying \cite[(181)]{B-CMP102b} with
$\bar v_w=\bar u^{-1}$ returns to the native objects. Hence
$\mathsf{P}^{1}\lceil\mathfrak{r}_v=(U_j,J_j)(\mathfrak{r}_v,\mathsf{V}'^{1})$ modulo gauge.

Finally, let $\varpi^{1}(x):=\mathsf{V}'^{1}(\Gamma_{y,x})^{-1}$ be the comb gauge of
$\mathsf{V}'^{1}$
from $y$. It is $G^{c}$-valued, and, as for $\mathsf{P}^{0}$,
$(\mathsf{V}'^{1})^{\varpi^{1}}(b)=\mathsf{V}'^{1}(\lambda_b)=e^{iB^{1}(b)}$ with $|B^{1}|<b_H$.
The
identity \cite[(181)]{B-CMP102b} is stated in \cite{B-CMP102b} for every gauge transformation for
which the transformed data lie in the analyticity domain; since
$(\mathsf{V}'^{1})^{\varpi^{1}}=e^{iB^{1}}$ with $|B^{1}|<b_H$ lies in that domain, it applies to
$\varpi^{1}$, the currents again
following by \cite[(1.10)]{B-CMP109}. Therefore $(U_j,J_j)(\mathfrak{r}_v,e^{iB^{1}})$ lies in the
$G^{c}$-orbit of $\mathsf{P}^{1}\lceil\mathfrak{r}_v$, which is clause~(i) for
$\mathsf{P}^{1}$.
\end{proof}

\begin{proof}[Proof of Lemma~\ref{5.5:lem-loop-data}, continued: bound on the difference of loop data]
\label{5.5:prf-loop-data-difference}
This part of the proof derives clause~(iii) of Lemma~\ref{5.5:lem-loop-data} and the holomorphy
assertion of that lemma. It uses four inputs besides the parts of the proof already given:
\begin{itemize}
\item the covariance of the block averages under gauge transformations;
\item the averaging bound at a regular $G$-based field, by which a small $1$-form
  $\mathcal{X}$ in the exponent changes the averages bond by bond by factors $e^{i\mathfrak{h}}$;
  it is used below in the form of the bound on $|\mathfrak{h}_S(c)|$ displayed after
  \eqref{5.5:eq-loop-data-difference-1};
\item the support and distance bound \eqref{5.5:eq-correction-support-b};
\item the smallness conditions \eqref{5.5:eq-remainder-constants-a}.
\end{itemize}
Throughout, for a bond $c$ of $\mathfrak{B}_j(\mathfrak{r}_v)$ we write $l(c)$ for its level,
$\ell_{l(c)}=L^{l(c)}\eta$ for its length, and $c_-$, $c_+$ for its initial and final points.

\emph{The working gauge.} We use the response \eqref{5.5:eq-correction-response-a}, in which
$U^{u}_{\mathsf{P}^{1}}=[e^{i\eta\mathcal{H}_S}W]^{\tilde v}$. We apply the covariance of the block
averages and the averaging bound at the regular $G$-based field $e^{i\eta\mathcal{H}_0}U_{00}$,
with the $1$-form $\mathcal{X}:=\operatorname{Ad}_{v_0^{-1}}\mathcal{H}_S$, which satisfies
$|\mathcal{X}|\le 2^{1/4}|\mathcal{H}_S|$. This gives, for every bond $c$ of
$\mathfrak{B}_j(\mathfrak{r}_v)$,
\begin{equation}\label{5.5:eq-loop-data-difference-1}
\mathsf{V}'^{1}_u(c)=\tilde v(c_-)\,e^{i\mathfrak{h}_S(c)}\,\mathsf{V}'_u(c)\,\tilde v(c_+)^{-1}
\end{equation}
with a $\mathfrak{g}^{c}$-valued $\mathfrak{h}_S(c)$ satisfying
\begin{equation*}
|\mathfrak{h}_S(c)|\le 2^{1/2}c_A\,\ell_{l(c)}\sup_{\text{blocks of }c}|\mathcal{H}_S|
\le 2c_A\,\frac{\ell_{l(c)}}{\ell_j}\,h_v .
\end{equation*}
The second inequality holds because $\ell_j|\mathcal{H}_S|\le h_v$ on $\mathfrak{r}_v$ by
\eqref{5.5:eq-correction-response-a}, and because $2^{1/2}\le 2$.

\emph{Telescoping round the loop.} Let $c_1,c_2,\dots$ be the successive bonds of the loop
$\lambda_b$. Let $\Pi_k:=\mathsf{V}'_u(c_1)\cdots\mathsf{V}'_u(c_k)$ be the partial products of
$\mathsf{V}'_u$ along $\lambda_b$, with $\Pi_0:=1$. By
the corresponding display, $n(\Pi_k)\le 1.45$. The level-$j$ length of the loop
satisfies
\begin{equation*}
\operatorname{len}_j(\lambda_b)\le 2d|b-y|+1\le\min\{2d(1+|b-y|),\lambda_v\}.
\end{equation*}

Multiply the identities \eqref{5.5:eq-loop-data-difference-1} along $\lambda_b$. For consecutive
bonds, $c_{k,+}=c_{k+1,-}$, so the factors $\tilde v(c_{k,+})^{-1}\tilde v(c_{k+1,-})$ equal
$1$. Only the factors $\tilde v(y)$ and $\tilde v(y)^{-1}$ at the base point $y$ remain.
Next, move each $e^{i\mathfrak{h}_S(c_k)}$ to the left through $\Pi_{k-1}$, using
$\Pi_{k-1}e^{iX}=e^{i\operatorname{Ad}_{\Pi_{k-1}}X}\Pi_{k-1}$. Since $\mathsf{V}'_u(\lambda_b)$ is
the full product and equals $e^{iB^{0}_u(b)}$, this yields
\begin{equation}\label{5.5:eq-loop-data-difference-2}
\tilde v(y)^{-1}\,\mathsf{V}'^{1}_u(\lambda_b)\,\tilde v(y)
=\Bigl[\prod_k e^{i\operatorname{Ad}_{\Pi_{k-1}}\mathfrak{h}_S(c_k)}\Bigr]\,e^{iB^{0}_u(b)}
=:e^{i(B^{0}_u+\Upsilon)(b)} .
\end{equation}

\emph{Bound on the increment $\Upsilon$.} Since $\|\operatorname{Ad}_g\|\le n(g)$, each exponent
in \eqref{5.5:eq-loop-data-difference-2} has norm at most $1.45\,|\mathfrak{h}_S(c_k)|$. Put
$\Sigma:=1.45\sum_k|\mathfrak{h}_S(c_k)|$. The lengths add up along the loop:
$\sum_k\ell_{l(c_k)}=\ell_j\operatorname{len}_j(\lambda_b)$. Hence
\begin{equation*}
\Sigma\le 1.45\cdot 2c_A\operatorname{len}_j(\lambda_b)\,h_v
\le 1.45\cdot 2c_A\cdot 2d(1+|b-y|)\,h_v .
\end{equation*}
By \eqref{5.5:eq-correction-support-b}, $1+|b-y|\le 1+2D_v$. By the definition of $c_\delta$ in
Definition~\ref{5.5:def-remainder-constants}, $1+2D_v\le c_\delta e^{\frac18\delta D_v}$. The
bound on $h_v$ in the corresponding display, established earlier in this proof, then gives
\begin{equation*}
\Sigma\le 5.8\,d\,c_Ac_\delta e^{\frac18\delta D_v}\cdot K_1\beta e^{-\frac14\delta D_v}
\bigl[5.5\vartheta+4.4|\mathsf{t}|\vartheta(\square)\bigr]\alpha_{\cdot,j}.
\end{equation*}
By the definition of $K^{\mathrm{rem}}$ in Definition~\ref{5.5:def-remainder-constants},
$\tfrac12K^{\mathrm{rem}}c_3'=13d\beta c_A'c_\delta K_1$, so that
$\tfrac14K^{\mathrm{rem}}c_3'=6.5\,d\beta c_A'c_\delta K_1$. This is at least
$5.8\,d\beta c_Ac_\delta K_1 e^{-\frac18\delta D_v}$, since $c_A'\ge c_A$ and
$e^{-\frac18\delta D_v}\le 1$. Next, $c_3'\alpha_{\cdot,j}\le\alpha_{3,j}$ by the definition of
$c_3'$. The conditions \eqref{5.5:eq-remainder-constants-a}, together with
$|\mathsf{t}|\le r^{\mathrm{rem}}_{\square}$, give
$K^{\mathrm{rem}}[5.5\vartheta+4.4|\mathsf{t}|\vartheta(\square)]\le 5.5/17+4.4/14$. Finally,
$\alpha_{3,j}\le 1/16$, one of the smallness conditions on $\gamma$ imposed at the beginning of
this section. Together these give
\begin{align*}
\Sigma&\le\frac{K^{\mathrm{rem}}c_3'}{4}\bigl[5.5\vartheta+4.4|\mathsf{t}|\vartheta(\square)\bigr]
\alpha_{\cdot,j}
\le\frac14\Bigl(\frac{5.5}{17}+\frac{4.4}{14}\Bigr)\alpha_{3,j}\\
&\le 0.16\,\alpha_{3,j}\le 0.01 .
\end{align*}
Consequently the product in \eqref{5.5:eq-loop-data-difference-2} satisfies
$|\prod-1|\le e^{\Sigma}-1\le 1.02\,\Sigma$.

Now apply the Baker--Campbell--Hausdorff formula at $1$. We use
$|B^{0}_u(b)|\le 1.2\,\alpha_{3,j}/L$, which follows from the corresponding display together
with $\sup_b|B^{ax}(b)|_{\star}<\alpha_{3,j}$ from clause~(ii) of
Lemma~\ref{5.5:lem-loop-data}. This gives
\begin{equation}\label{5.5:eq-loop-data-difference-3}
|\Upsilon(b)|\le 1.17\,\Sigma\le 6.8\,d\,c_A(1+|b-y|)\,h_v .
\end{equation}
Here the last inequality combines $1.17\cdot 1.45\cdot 4\le 6.8$ with the first bound on $\Sigma$
above. In particular $|\Upsilon|_\infty\le 1.17\cdot 0.16\,\alpha_{3,j}$.

\emph{The working comb data.} We have $\log(gXg^{-1})=\operatorname{Ad}_g\log X$ near
$1$. Together with \eqref{5.5:eq-loop-data-difference-2}, this shows
\begin{equation*}
B^{1}_u:=\tfrac1i\log\mathsf{V}'^{1}_u(\lambda_b)=\operatorname{Ad}_{\tilde v(y)}(B^{0}_u+\Upsilon).
\end{equation*}
Therefore
\begin{equation*}
\Delta_u:=B^{1}_u-B^{0}_u=\Upsilon+(\operatorname{Ad}_{\tilde v(y)}-1)(B^{0}_u+\Upsilon).
\end{equation*}
By \eqref{5.5:eq-correction-response-a}, $\|\operatorname{Ad}_{\tilde v(y)}-1\|\le\nu_S(y)$ and
$\nu_S(y)\le 9h_v/(4L^3r_V)$. The blocking factor $L$ is odd and larger than one, so $L\ge 3$, and
hence
\begin{equation*}
|B^{0}_u+\Upsilon|_\infty\le 1.2\,\alpha_{3,j}/L+1.17\cdot 0.16\,\alpha_{3,j}\le\alpha_{3,j}.
\end{equation*}
With \eqref{5.5:eq-loop-data-difference-3}, and with $\alpha_{3,j}\le 1\le 1+|b-y|$, we obtain
\begin{align*}
|\Delta_u(b)|&\le 6.8\,d\,c_A(1+|b-y|)h_v+\nu_S(y)\,\alpha_{3,j}
\le\Bigl[6.8\,d\,c_A+\frac{9}{4L^3r_V}\Bigr](1+|b-y|)h_v\\
&\le 6.8\,d\,c_A'(1+|b-y|)h_v .
\end{align*}
The last step uses the definition of $c_A'$ in Definition~\ref{5.5:def-remainder-constants}, by
which $6.8d(c_A'-c_A)=3.4/(L^3r_V)\ge 9/(4L^3r_V)$.

\emph{The native data.} By the covariance of the block averages and \cite[(181)]{B-CMP102b}, both
applied at the block-local gauge $v_w$ with $n(v_w(x))\le 2^{1/4}$, we have
$U_{\mathsf{P}^{\iota}}=(U^{u}_{\mathsf{P}^{\iota}})^{\bar v_w}$ for $\iota=0,1$, where
$U^{u}_{\mathsf{P}^{0}}:=W$ and $\bar v_w$ denotes the block-constant restriction of $v_w$ to base
points, so that $\bar v_w=\bar u^{-1}$ (Notation~\ref{5.5:not-working-gauge}). Hence the native
loop data at $y$ are the images under
$\operatorname{Ad}_{\bar v_w(y)}$ of the working ones:
\begin{equation*}
B^{0}=\operatorname{Ad}_{\bar v_w(y)}B^{0}_u,\qquad B^{1}=\operatorname{Ad}_{\bar v_w(y)}B^{1}_u,
\qquad\Delta=\operatorname{Ad}_{\bar v_w(y)}\Delta_u .
\end{equation*}
Since
\begin{equation*}
\|\operatorname{Ad}_{\bar v_w(y)}\|\le n(\bar v_w(y))\le 2^{1/4}\le 1.2,
\end{equation*}
it follows that
\begin{equation*}
|\Delta(b)|\le 1.2\cdot 6.8\,d\,c_A'(1+|b-y|)h_v\le 13\,d\,c_A'(1+|b-y|)h_v .
\end{equation*}
This is the first bound of clause~(iii) of Lemma~\ref{5.5:lem-loop-data}. Note that $B^{1}$ is
the comb datum of $\mathsf{P}^{1}$ itself, determined by $\mathsf{P}^{1}$ and $y$. No normalisation
$\tilde v(y)=1$ has been used.

\emph{The sup bound.} Combine the last bound with \eqref{5.5:eq-correction-support-b}, the bound on
$h_v$ in the corresponding display, and the identity
$\tfrac12K^{\mathrm{rem}}c_3'=13d\beta c_A'c_\delta K_1$. Use also
$c_3'\alpha_{\cdot,j}\le\alpha_{3,j}$, the conditions \eqref{5.5:eq-remainder-constants-a}, and
$|\mathsf{t}|\le 1/(14K^{\mathrm{rem}}\vartheta(\square))$. The result is
\begin{equation}\label{5.5:eq-loop-data-difference-4}
\begin{split}
|\Delta|_\infty&\le 13\,d\,c_A'c_\delta K_1\beta\bigl[5.5\vartheta+4.4|\mathsf{t}|\vartheta(\square)\bigr]
\alpha_{\cdot,j}e^{-\frac18\delta D_v}\\
&\le\tfrac12K^{\mathrm{rem}}\bigl[5.5\vartheta+4.4|\mathsf{t}|\vartheta(\square)\bigr]
\alpha_{3,j}e^{-\frac18\delta D_v}\\
&\le\tfrac12\Bigl(\frac{5.5}{17}+\frac{4.4}{14}\Bigr)\alpha_{3,j}e^{-\frac18\delta D_v}
\le\tfrac12\,\alpha_{3,j}e^{-\frac18\delta D_v}=|\Delta|_{\star}.
\end{split}
\end{equation}

\emph{Holomorphy of $B^{0}$, $B^{1}$, $\Delta$.} The averages $\mathsf{V}'$ and
$\mathsf{V}'^{1}$ are holomorphic functions of $(\mathbb{U},\mathbb{J},b,\sigma,\mathsf{t})$. This
follows from the holomorphy of the pairs $\mathsf{P}^{0},\mathsf{P}^{1}$ in these variables,
stated in the first sentence of part~(b) of the theorem constructing them, and from the
analyticity of the averaging operations \cite{B-CMP98}. Products along loops are analytic, and so
is the logarithm near $1$. Hence $B^{0}$, $B^{1}$ and $\Delta$ are holomorphic in
$(\mathbb{U},\mathbb{J},b,\sigma,\mathsf{t})$. This completes the proof of
Lemma~\ref{5.5:lem-loop-data}.
\end{proof}

\begin{theorem}[Membership, holomorphy and bound of the remainder]\label{5.5:thm-membership-remainder}
Work in the setting of Notation~\ref{5.5:not-interpolated-box-field}, under its floors, among them
\eqref{5.5:eq-interpolated-box-field-a}, under the smallness condition on $\gamma$ that
accompanies that setting for the sourced box field, and under the conditions
\eqref{5.5:eq-remainder-constants-a}, clause~(i) of \eqref{5.5:eq-remainder-constants-b},
\eqref{5.5:eq-remainder-constants-c} and \eqref{5.5:eq-remainder-constants-e} of
Definition~\ref{5.5:def-remainder-constants}. Let $v=(T,\square)$ be box-read in the sense of
Definition~\ref{5.5:def-box-read-units}, with own box $\mathfrak{r}_v$ and $X:=X_T$. Let
$\mathsf{P}^{0}$, $\mathsf{P}^{1}$, $G_v$ and $\operatorname{rem}(T,\square)$ be as in
Definition~\ref{5.5:def-stored-remainder}, let $D_v$ be as in
Lemma~\ref{5.5:lem-correction-support}, and let $B^{0}$, $B^{1}$, $\Delta=B^{1}-B^{0}$ be the
canonical loop data of Lemma~\ref{5.5:lem-loop-data}. Then for every source point, every
$t\in[0,1]$, every $\mathsf{t}$ with $|\mathsf{t}|\le r^{\mathrm{rem}}_{\square}$ and every
$\sigma$ with $|\sigma|\le e^{\kappa_1}$ the following hold.
\begin{itemize}
\item[(a)] \emph{Membership.} For $|\sigma_1|\le e^{\frac18\delta D_v}$ put
\begin{equation*}
\mathsf{C}(\sigma_1):=(U_j,J_j)\bigl(\mathfrak{r}_v,e^{i(B^{0}+\sigma_1\Delta)}\bigr)\restriction X .
\end{equation*}
Then $\mathsf{C}(\sigma_1)\in\mathfrak{U}_T$ for all such $\sigma_1$. Moreover
$\mathsf{C}(0)\cong\mathsf{P}^{0}\restriction X$ and $\mathsf{C}(1)\cong\mathsf{P}^{1}\restriction X$,
where $\cong$ denotes equality modulo $G^{c}$-gauge transformations, and both pairs
$\mathsf{P}^{0}\restriction X$, $\mathsf{P}^{1}\restriction X$ lie in $\mathfrak{U}_T$.
\item[(b)] \emph{Holomorphy.} The map $(\mathbb{U},\mathbb{J},b,\sigma,\mathsf{t})\mapsto G_v$ is
holomorphic on the product of the set of source points with
$\{|\sigma|\le e^{\kappa_1}\}$ and $\{|\mathsf{t}|\le r^{\mathrm{rem}}_{\square}\}$, and continuous
in $t$; the variable $\sigma$ enters only through $\mathbb{H}^{w}(\sigma)$ and
$\mathbf{H}_j(\sigma)$. The resulting term belongs to the near group, with domain
$\mathsf{d}=\tilde{\square}^{5}\cup\mathsf{y}$.
\item[(c)] \emph{The real jet.} At the real point of the construction, with $\sigma\equiv 1$
and $\mathsf{t}$ real and small, $T(X;\mathsf{P}^{1}(t,\mathsf{t}))$ is the integrand of the
interpolation formula by which the near term is filed, at the value $t_{\square}=\mathsf{t}$ of its
interpolation parameter. Consequently
\begin{equation*}
\int_0^1 dt\,\partial_{\mathsf{t}}\big|_{0}T(X;\mathsf{P}^{0})+\operatorname{rem}(T,\square)
=\operatorname{piece}(T,\square),
\end{equation*}
the first summand being the term \cite[(3.24)--(3.25)]{B-CMP109} in the form in which the filing
of the near terms stores it, and $\operatorname{rem}(T,\square)$ is the $t_{\square}$-jet of the sum
of the two terms \cite[(3.21), (3.22)]{B-CMP109} removed in Ba{\l}aban's construction. The
remaining clause of that filing is not affected.
\item[(d)] \emph{Bound.}
\begin{equation}\label{5.5:eq-membership-remainder-a}
|G_v|\le 2\,\mathfrak{n}_v\,e^{-\frac18\delta D_v},\qquad
|\operatorname{rem}(T,\square)|\le\Bigl[\frac{2}{r^{\sharp}_{\square}}
+28K^{\mathrm{rem}}\vartheta(\square)\Bigr]\mathfrak{n}_v\,e^{-\frac18\delta D_v}.
\end{equation}
The conclusion of the Cauchy-formula lemma attached to the transfer of the analysis of
\cite[\S3]{B-CMP109} to the cover holds for the function $G_v$ at the triple
$(\mathfrak{N},\rho_A,\eta_p):=(2\mathfrak{n}_v e^{-\frac18\delta D_v},r^{\mathrm{rem}}_{\square},1)$.
\item[(e)] \emph{Sum.} Put
\begin{equation*}
\mathfrak{s}^{\mathrm{rd}}:=\sup_{j,\;q\in\pi_j}\sum\{\mathfrak{n}_v:\ v\ \text{anchored in}\ q\},
\end{equation*}
a number independent of $K$. Then
\begin{equation*}
\sum\{\mathfrak{n}_v e^{-\frac18\delta D_v}:\ v\ \text{box-read at}\ \square\}
\le c_{\Sigma}\,\mathfrak{s}^{\mathrm{rd}},\qquad c_{\Sigma}\le 2\cdot 6^{d},
\end{equation*}
and likewise for the box-read units attached to one unit $n$-cell $\Delta'$.
\item[(f)] The inputs used by clauses (a)--(e) are listed in the remark following the proof.
\item[(g)] \emph{Two runs.} Let $\mathsf{A}$, $\mathsf{B}$ be two runs of the construction on one and
the same source-free scaffold (Definition~\ref{5.1:def-source-free-scaffold}), and let
$\mathcal{K}^{\mathrm{rd}}$ be the scaffold class. For $\mathsf{m}\in\{\mathsf{B},\mathsf{A}\}$
write $\mathsf{P}^{\iota,\mathsf{m}}$, $T^{\mathsf{m}}$, $\mathfrak{U}^{\mathsf{m}}_T$, $G_v^{\mathsf{m}}$
for the objects above in the run $\mathsf{m}$, and for any such quantity put
$\delta(\cdot):=(\cdot)^{\mathsf{B}}-(\cdot)^{\mathsf{A}}$. Define two moves.
The \emph{$\mathsf{P}^{1}$-move} sends $((b,\sigma,t,\mathsf{t});p)$, with $b$ and $\sigma$ ranging
as in the first move of the two-run comparison lemma, $t\in[0,1]$ and
$|\mathsf{t}|\le r^{\mathrm{rem}}_{\square}$, to the pair
$(\mathsf{P}^{1,\mathsf{B}}\restriction X,\mathsf{P}^{1,\mathsf{A}}\restriction X)$; it goes from
$(k;\mathsf{n}+1,\hat X)$ to $(j,X)$. The \emph{splitting move}, taken at the canonical splitting
$(B^{0},\Delta)$, sends $((\pi,\sigma_0,\sigma_1);p)$ to
\begin{equation*}
\bigl(\mu^{\mathsf{m}}(\pi)\bigr)_{\mathsf{m}=\mathsf{B},\mathsf{A}},\qquad
\mu^{\mathsf{m}}(\pi):=(U_j,J_j)^{\mathsf{m}}\bigl(\mathfrak{r}_v,
e^{i(\sigma_0B^{0,\mathsf{m}}+\sigma_1\Delta^{\mathsf{m}})}\bigr)\restriction X,
\end{equation*}
where $B^{0,\mathsf{m}}$, $\Delta^{\mathsf{m}}$ are the data of Lemma~\ref{5.5:lem-loop-data} of
$(\mathsf{P}^{0,\mathsf{m}},\mathsf{P}^{1,\mathsf{m}})$ at the parameters $\pi$ of the
$\mathsf{P}^{1}$-move, on the domain
\begin{equation*}
|\sigma_0|\,|B^{0}|_{\star}+|\sigma_1|\,|\Delta|_{\star}<\alpha_{3,j},
\end{equation*}
with $|B^{0}|_{\star}$, $|\Delta|_{\star}$ the scaffold numbers of Lemma~\ref{5.5:lem-loop-data}.
If the hypothesis on moves of the two-run comparison lemma holds for the $\mathsf{P}^{1}$-move and
for the splitting move on $\mathcal{K}^{\mathrm{rd}}$, then for
$p\in\mathcal{K}^{\mathrm{rd}}(k;\mathsf{n}+1,\hat X)$
\begin{equation*}
|\delta G_v|\le 2\,\mathsf{S}\,e^{-\frac18\delta D_v},
\end{equation*}
and \textup{(d)} and \textup{(e)} hold for the differences $\delta$ with $\mathfrak{n}_v$ replaced by
$\mathsf{S}:=\mathsf{S}^{\mathrm{rd}}$.
\end{itemize}
\end{theorem}

Stated; the proof below is incomplete at the step marked $(\ast)$.

\begin{proof}
\emph{Clause (a).} By Lemma~\ref{5.5:lem-loop-data}, the corresponding display gives
$|B^{0}|_{\infty}\le|B^{0}|_{\star}<1.42\,\alpha_{3,j}/L$, and
the corresponding display gives
$|\Delta|_{\infty}\le\tfrac12\alpha_{3,j}e^{-\frac18\delta D_v}=|\Delta|_{\star}$. Hence, for
$|\sigma_1|\le e^{\frac18\delta D_v}$,
\begin{equation*}
|B^{0}+\sigma_1\Delta|_{\infty}\le|B^{0}|_{\star}+\tfrac12\alpha_{3,j}
<\Bigl(\frac{1.42}{L}+\frac12\Bigr)\alpha_{3,j}<\alpha_{3,j},
\end{equation*}
where the last inequality uses the floor $L\ge 8$, which gives $1.42/L\le 0.18$, and
$0.18+0.5<1$. The own box $\mathfrak{r}_v$ is a box about $X$ with the margin required by
the margin statement of the transfer of the analysis of \cite[\S3]{B-CMP109} to the cover. That
transfer is stated in the sentence
\begin{quotation}
We repeat the analysis of Sect.\ 3 [I] using the above cover $\{\square\}$, and the partition of
unity $\{\zeta_{\square}\}$.
\end{quotation}
printed without proof in \cite[p.~276]{B-CMP119}; cited as printed. There the reference [I] is
\cite{B-CMP109}. For a core the margin is one layer of $\pi_j$-cubes, as in \cite{B-CMP119}; for a
short single it is at least two layers. The analyticity statement of that transfer applies to
complex configurations; on $\mathfrak{r}_v$, with the bound just displayed, it gives
$\mathsf{C}(\sigma_1)\in U^{c}_j(X,a_j)=\mathfrak{U}_T$, where $U^{c}_j(X,a_j)$ is the analyticity
space of Definition~\ref{5.1:def-analyticity-spaces}. By \cite[(1.10), (1.19)]{B-CMP109},
$\mathfrak{U}_T$ is a union of $G^{c}$-orbits. Clause (i) of Lemma~\ref{5.5:lem-loop-data} gives
$\mathsf{P}^{0}\restriction\mathfrak{r}_v\cong(U_j,J_j)(\mathfrak{r}_v,e^{iB^{0}})$ and
$\mathsf{P}^{1}\restriction\mathfrak{r}_v\cong(U_j,J_j)(\mathfrak{r}_v,e^{iB^{1}})$; since
$B^{1}=B^{0}+\Delta$ and $X\subset\mathfrak{r}_v$, restriction to $X$ gives
$\mathsf{C}(0)\cong\mathsf{P}^{0}\restriction X$ and $\mathsf{C}(1)\cong\mathsf{P}^{1}\restriction X$.
As $\mathsf{C}(0),\mathsf{C}(1)\in\mathfrak{U}_T$ and $\mathfrak{U}_T$ is a union of orbits, both
pairs $\mathsf{P}^{0}\restriction X$, $\mathsf{P}^{1}\restriction X$ lie in $\mathfrak{U}_T$.

\emph{Clause (b).} The map $V\mapsto(U,J)(\mathfrak{B}^{\star}_j,V)$ is analytic on the domain of
\cite[Proposition~9]{B-CMP102b}. By Definition~\ref{5.5:def-stored-remainder},
\begin{equation*}
\mathsf{P}^{1}=(U,J)(\mathfrak{B}^{\star}_j,e^{i\mathsf{E}}M_{\star}\mathfrak{b})
=(U,J)(\mathfrak{B}^{\star}_j,e^{iD^{\sharp}}M_{\star}\mathbb{W}),
\end{equation*}
and the argument
$e^{i\mathsf{E}}M_{\star}\mathfrak{b}$ does not leave that domain, by the corresponding bound on the correction form. The argument $e^{iD^{\sharp}}M_{\star}\mathbb{W}$ is
holomorphic in $(\mathbb{U},\mathbb{J},b,\sigma,\mathsf{t})$: the fields $\mathbb{H}^{w}$ and
$\mathbf{H}_j$ are holomorphic by clause (a) of the theorem on the fields $\mathbb{H}^{w}$,
$\mathbf{H}_j$, applied at $\mathfrak{B}$ and at $\mathfrak{B}_j$, whose hypotheses are met under
\eqref{5.5:eq-remainder-constants-e}; the linearised map $Q$ is analytic by
\cite[Proposition~4]{B-CMP98}; and $\mathbb{W}$ is analytic in $\mathbb{U}$. The near term $T$ is
analytic on $\mathfrak{U}_T$, which contains $\mathsf{P}^{1}\restriction X$ and
$\mathsf{P}^{0}\restriction X$ by clause (a). Continuity in $t$ holds because every ingredient is
continuous in $t$. That $\sigma$ enters only through $\mathbb{H}^{w}(\sigma)$ and $\mathbf{H}_j(\sigma)$
is the decoupled dependence on $\sigma$ of the theorem on the fields $\mathbb{H}^{w}$, $\mathbf{H}_j$,
at $\mathfrak{B}$ and at $\mathfrak{B}_j$, as used in Corollary~\ref{5.2:cor-sourced-preparatory}.

\emph{Clause (c).} We work at the real point. There the box field $\mathbb{W}$ equals
$U:=U^{(\mathsf{n}+1)}_k$ modulo $G$-gauge, by the real-point identification of the box field;
$\mathbb{H}^{w}=\mathbb{H}$ by clause (c) of the theorem on the fields $\mathbb{H}^{w}$,
$\mathbf{H}_j$; and clause (d) of that theorem, applied at $\mathfrak{B}_j$, gives
\begin{equation*}
e^{i\eta\mathbf{H}_j(t,\mathsf{t})}\mathbb{W}
=U_j\bigl(e^{i\mathsf{B}(t,\mathsf{t})}\bar{\mathbb{W}}^{j}\bigr)\quad\text{modulo gauge},
\end{equation*}
where $\bar{\mathbb{W}}^{j}$ and $\bar U^{j}$ denote the block averages at level $j$ of $\mathbb{W}$ and
of $U$, $Q_j$ the linearised averaging map at level $j$, and $U_j(\cdot)$ the configuration
function with prescribed block averages, in the notation of \cite[(3.19)]{B-CMP109}. Hence
\begin{equation*}
Q_{\mathfrak{B}_j}(\mathbb{W};\eta\mathbf{H}_j)=\mathsf{B}
=Q_{\mathfrak{B}_j}\bigl(\mathbb{W};\eta(t+\mathsf{t}\zeta_{\square})\mathbb{H}^{w}\bigr),
\end{equation*}
the second equality being the definition of $\mathsf{B}$. This is the sourced form of
\cite[(3.20)]{B-CMP109}, which states, in the notation of \cite{B-CMP109},
\begin{equation*}
Q(\xi\mathbf{H}_j(B_{\square}))=Q_j(\xi\mathbf{H}_j(B_{\square}))=B_{\square}
\quad\text{on a neighbourhood of}\ \tilde{\square}^{4}.
\end{equation*}
By \cite{B-CMP109}, $\mathfrak{B}^{\star}_j=\mathfrak{B}_j$ on $\tilde{\square}^{4}$, and
$\tilde{\square}^{4}\supset\operatorname{supp}\zeta'$. On the bonds where $\zeta'\neq0$ we therefore
have $Q_{\star}((t+\mathsf{t}\zeta_{\square})\mathbb{H}^{w})=Q_{\star}(\mathbf{H}_j)$, by the sourced
form of \cite[(3.20)]{B-CMP109} displayed above and the coincidence of $\mathfrak{B}^{\star}_j$ with
$\mathfrak{B}_j$ on $\tilde{\square}^{4}$, so that the definition of $D^{\sharp}$ gives
\begin{equation*}
D^{\sharp}=\zeta'\,Q_{\star}((t+\mathsf{t}\zeta_{\square})\mathbb{H}^{w})
+(1-\zeta')\,Q_{\star}(\mathbf{H}_j(t,\mathsf{t}))=Q_{\star}(\mathbf{H}_j)
\quad\text{on all bonds}.
\end{equation*}
Consequently
$e^{iD^{\sharp}}M_{\star}\mathbb{W}=M_{\star}(e^{i\eta\mathbf{H}_j}\mathbb{W})$, and a minimiser is the
box minimiser of its own data, as in \cite[(3.19)]{B-CMP109}, which reads, in the notation of
\cite{B-CMP109},
\begin{equation*}
E^{(j)}\bigl(X,U_j(\exp iB_{\square}\,\bar U^{j}_{k+1})\bigr)
=E^{(j)}\Bigl(X,U_j\bigl(\square_0,\exp iQ(\xi\mathbf{H}_j(B_{\square}))\,
M_{\cdot}(U_{k+1})\bigr)\Bigr).
\end{equation*}
Therefore
\begin{equation*}
\mathsf{P}^{1}(t,\mathsf{t})\cong U_j\bigl(\exp iQ_j(\eta(t+\mathsf{t}\zeta_{\square})\mathbb{H})
\,\bar U^{j}\bigr)\quad\text{on}\ \square_0^{\flat}\supset X,
\end{equation*}
which is the configuration of the interpolation formula at $t_{\square}=\mathsf{t}$; this proves the
first assertion of (c).

We now identify the two removed terms of \cite[p.~274]{B-CMP109}. The integral term of
\cite[(3.21)]{B-CMP109} is
\begin{multline*}
\int_0^1 d\tilde t\,\frac{d}{d\tilde t}\,E\Bigl(X,\exp i\xi\mathbf{H}_j\bigl(\square_0,
(\tilde t(1-\tilde\zeta_{\square})+\tilde\zeta_{\square})\,Q(\xi\mathbf{H}_j(B_{\square}))\bigr)
U_{k+1}\Bigr)\\
=\phi(1)-\phi(0),
\end{multline*}
where $\phi(\tilde t)$ denotes $E$ at the box minimiser with datum
$(\tilde t(1-\tilde\zeta_{\square})+\tilde\zeta_{\square})Q(\xi\mathbf{H}_j(B_{\square}))$. Here
$\phi(1)$ is the right member of \cite[(3.19)]{B-CMP109}, and, by \cite[(3.20)]{B-CMP109},
\begin{equation*}
\phi(0)=E\bigl(X,U_j(\square_0,\exp i\tilde\zeta_{\square}B_{\square}\,M_{\cdot}(U_{k+1}))\bigr).
\end{equation*}
Next, \cite[(3.22)]{B-CMP109} writes
\begin{equation*}
\tilde\zeta_{\square}B_{\square}
=\bigl[\tilde\zeta_{\square}Q_j(\eta(t+t_{\square}\zeta_{\square})\mathbf{H}_k(B'))
-Q_j(\eta(t\tilde\zeta_{\square}+t_{\square}\zeta_{\square})\mathbf{H}_k(B'))\bigr]
+\tilde B_{\square},
\end{equation*}
and Ba{\l}aban removes the expression in the square brackets by repeating the procedure, ending at
\cite[(3.23)]{B-CMP109},
\begin{equation*}
E\Bigl(X,U_j\bigl(\square_0,\exp i\tilde B_{\square}\,
M_{\cdot}(U_{k+1}(\square_0,M_{\cdot}(U_{k+1})))\bigr)\Bigr).
\end{equation*}
Each of the two removals is the difference of one functional on $X$ at two box minimisers, and the
two differences telescope exactly to $E(X;\mathsf{P}^{1})-E(X;\mathsf{P}^{0})$, where
$\mathsf{P}^{0}=(U,J)(\mathfrak{B}^{\star}_j,e^{i\tilde{\mathsf{B}}}M_{\star}\mathbb{W})$ is the pair
of \cite[(3.23)]{B-CMP109} on the capped set. This $\mathsf{P}^{0}$ coincides verbatim with the
stored form of the term \cite[(3.24)--(3.25)]{B-CMP109} in the filing of the near terms; both differ
from the expressions of \cite{B-CMP109} only in that $\mathfrak{B}^{\star}_j$ replaces the one-scale
set and $(\mathbb{U},\mathbb{J})$ are kept as variables. Passing to the
constituent $T$ of $E(X;\cdot)$ that the unit $v$ carries, the two removed terms contribute
$T(X;\mathsf{P}^{1})-T(X;\mathsf{P}^{0})=G_v$ to it, that is,
$T(X;\mathsf{P}^{1})=T(X;\mathsf{P}^{0})+G_v$. Applying $\partial_{\mathsf{t}}$ at $\mathsf{t}=0$,
integrating over $t\in[0,1]$, using the definition of $\operatorname{rem}(T,\square)$ in
Definition~\ref{5.5:def-stored-remainder} and the identification of $T(X;\mathsf{P}^{1})$ with the
integrand of the interpolation formula, we obtain the identity of (c), with
$\operatorname{rem}(T,\square)$ the $t_{\square}$-jet of the sum of the two removed terms. The
remaining clause of the filing of the near terms is not affected.

\emph{Clause (d).} Put $F(\sigma_1):=T(X;\mathsf{C}(\sigma_1))$. By clause (a), and because the
complex-analyticity statement of the transfer to the cover gives
analyticity of $T$ for complex data with $|T|\le\mathfrak{n}_v$ on $\mathfrak{U}_T$, the function $F$
is holomorphic with $|F|\le\mathfrak{n}_v$ on $|\sigma_1|<e^{\frac18\delta D_v}$. By gauge invariance
\cite[(1.19)]{B-CMP109} and clause (i) of Lemma~\ref{5.5:lem-loop-data}, $G_v=F(1)-F(0)$. The function
$F-F(0)$ is bounded by $2\mathfrak{n}_v$ and vanishes at $0$; the Schwarz lemma on the disc of radius
$e^{\frac18\delta D_v}$ gives
\begin{equation*}
|G_v|=|F(1)-F(0)|\le 2\,\mathfrak{n}_v\,e^{-\frac18\delta D_v}.
\end{equation*}
By clause (b), $G_v$ is holomorphic in $\mathsf{t}$ on $|\mathsf{t}|\le r^{\mathrm{rem}}_{\square}$ at
each $t$, and it obeys the bound just proved there. Cauchy's estimate on that disc, with
$r^{\mathrm{rem}}_{\square}=\min\{r^{\sharp}_{\square},1/(14K^{\mathrm{rem}}\vartheta(\square))\}$ from
Definition~\ref{5.5:def-remainder-constants}, gives
\begin{equation*}
\begin{split}
\bigl|\partial_{\mathsf{t}}\big|_{0}G_v\bigr|
&\le 2\,\mathfrak{n}_v\,e^{-\frac18\delta D_v}
\max\Bigl\{\frac{1}{r^{\sharp}_{\square}},14K^{\mathrm{rem}}\vartheta(\square)\Bigr\}\\
&\le\Bigl[\frac{2}{r^{\sharp}_{\square}}+28K^{\mathrm{rem}}\vartheta(\square)\Bigr]
\mathfrak{n}_v\,e^{-\frac18\delta D_v}.
\end{split}
\end{equation*}
Since $\operatorname{rem}(T,\square)=\int_0^1dt\,\partial_{\mathsf{t}}|_{0}G_v(t,\mathsf{t})$, the same
bound holds for $|\operatorname{rem}(T,\square)|$, which proves
\eqref{5.5:eq-membership-remainder-a}. Finally, the Cauchy-formula lemma attached to the transfer to
the cover is Cauchy's formula in $\mathsf{t}$ and in the values of $\sigma$ on cells for a bounded
holomorphic function, linear in that function; by clause
(b) and the first bound of \eqref{5.5:eq-membership-remainder-a} it applies to $G_v$ at the triple
$(\mathfrak{N},\rho_A,\eta_p)=(2\mathfrak{n}_ve^{-\frac18\delta D_v},r^{\mathrm{rem}}_{\square},1)$.

\emph{Clause (e).} Let $v$ be box-read at $\square$. Then $\nu\ge3$; since $\nu=1+n-j$ and $n\le n_0$,
we get $s=n_0-j\ge n-j=\nu-1\ge2$. The enlarged block $\tilde{\square}^{2}$ lies in a cube of side six
cubes of $\pi_{n_0}$ (Definition~\ref{5.5:def-box-read-units}); hence, for given $s$, at most
$(6L^{s})^{d}$ cubes of $\pi_j$ can carry the anchor of a box-read unit at $\square$. By
\eqref{5.5:eq-correction-support-b}, $D_v\ge\frac15ML^{s}$, and with $M\ge w$ this gives
$\frac18\delta D_v\ge(\delta w/40)L^{s}$. Grouping the units by $s$ and by the cube $q$ of $\pi_j$ in
which they are anchored, and bounding each group by $\mathfrak{s}^{\mathrm{rd}}$, we obtain
\begin{equation*}
\sum_v\mathfrak{n}_v\,e^{-\frac18\delta D_v}
\le\mathfrak{s}^{\mathrm{rd}}\sum_{s\ge2}(6L^{s})^{d}e^{-(\delta w/40)L^{s}}
=c_{\Sigma}\,\mathfrak{s}^{\mathrm{rd}}.
\end{equation*}
By the floor \eqref{5.5:eq-interpolated-box-field-a}, $\delta w/40\ge0.4\ln L+0.1$; and
$0.4L^{s}\ge ds$ for $L\ge8$, $d=4$, $s\ge2$. Hence
\begin{equation*}
(6L^{s})^{d}e^{-(\delta w/40)L^{s}}
\le6^{d}\exp\bigl(ds\ln L-0.4L^{s}\ln L-0.1L^{s}\bigr)\le6^{d}e^{-L^{s}/10},
\end{equation*}
and since $L^{s}/10\ge s$ for $L\ge8$, $s\ge2$, the sum over $s\ge2$ is at most
$6^{d}\sum_{s\ge2}e^{-s}\le2\cdot6^{d}$; thus $c_{\Sigma}\le2\cdot6^{d}$. The number
$\mathfrak{s}^{\mathrm{rd}}$ does not depend on $K$: it is formed from the per-cube numbers of the
format of the correlation theorem, Theorem~\ref{5.1:thm-correlation-theorem}, multiplied by the
factor $E_0$ of that format, and the independence of these numbers of $K$ is part of the statement
of the transfer to the cover.

\emph{Clause (g).} The splitting move has one parameter domain for both runs: it is defined through
the scaffold majorants $|B^{0}|_{\star}$, $|\Delta|_{\star}$ of Lemma~\ref{5.5:lem-loop-data}, which are
numbers of the scaffold. The construction is the same in each run, since $B^{0,\mathsf{m}}$,
$\Delta^{\mathsf{m}}$ are functions of $(\mathsf{P}^{0,\mathsf{m}},\mathsf{P}^{1,\mathsf{m}},y)$ alone,
and it is jointly holomorphic. Lemma~\ref{5.5:lem-loop-data} holds in each run, so that on the domain
of the splitting move
\begin{equation*}
|\sigma_0B^{0,\mathsf{m}}+\sigma_1\Delta^{\mathsf{m}}|_{\infty}
\le|\sigma_0|\,|B^{0}|_{\star}+|\sigma_1|\,|\Delta|_{\star}<\alpha_{3,j},
\end{equation*}
and $\mu^{\mathsf{m}}(\pi)\in\mathfrak{U}^{\mathsf{m}}_T$ by the complex-analyticity statement of the
transfer to the cover. Given the hypothesis on moves of the two-run comparison lemma at the splitting
move, the two-run principle gives that
\begin{equation*}
\mathcal{D}(\sigma_0,\sigma_1):=T^{\mathsf{B}}(\mu^{\mathsf{B}})-T^{\mathsf{A}}(\mu^{\mathsf{A}})
\end{equation*}
is holomorphic on the domain with $|\mathcal{D}|\le\mathsf{S}$. By gauge invariance and clause (a) in
each run, $\delta G_v=\mathcal{D}(1,1)-\mathcal{D}(1,0)$. The slice $\sigma_0=1$ of the domain contains
the disc $|\sigma_1|<(\alpha_{3,j}-|B^{0}|_{\star})/|\Delta|_{\star}$, and, by the values of
$|B^{0}|_{\star}$, $|\Delta|_{\star}$ in Lemma~\ref{5.5:lem-loop-data} and $L\ge8$,
\begin{equation*}
\frac{\alpha_{3,j}-|B^{0}|_{\star}}{|\Delta|_{\star}}
\ge2\Bigl(1-\frac{1.42}{L}\Bigr)e^{\frac18\delta D_v}
\ge1.64\,e^{\frac18\delta D_v}\ge e^{\frac18\delta D_v}.
\end{equation*}
$(\ast)$ The two-run principle is applied here on the slice $\sigma_0=1$ out to this radius, that is,
at the dilation of the canonical splitting $(B^{0},\Delta)$. At this dilation the cross-cutoff clause
of the hypothesis on moves at the splitting move is the conjunction of the two-cutoff response
statement at zeroth order and the flat energy-difference bound, whose proofs are outstanding. Under
this clause, the Schwarz lemma applied to
$\sigma_1\mapsto\mathcal{D}(1,\sigma_1)-\mathcal{D}(1,0)$, which is bounded by $2\mathsf{S}$ and
vanishes at $0$, gives $|\delta G_v|\le2\mathsf{S}e^{-\frac18\delta D_v}$ at every
$(b,\sigma,t,\mathsf{t})$, and Cauchy's estimate in $\mathsf{t}$, as in clause (d), gives the bound on
$\delta\operatorname{rem}(T,\square)$ with $\mathfrak{n}_v$ replaced by $\mathsf{S}$. Given the
hypothesis on moves of the two-run comparison lemma at the $\mathsf{P}^{1}$-move, the
two-run form of the Cauchy-formula lemma places the pairs of the remainder in
$\mathcal{K}^{\mathrm{rd}}(j,X)$, and clauses (d) and (e) follow for the differences $\delta$ with
$\mathfrak{n}_v$ replaced by $\mathsf{S}$.
\end{proof}

\begin{remark}[Clause (f): the inputs of clauses (a)--(e)]
The proofs of clauses (a)--(e) of Theorem~\ref{5.5:thm-membership-remainder} use only the following:
the inequalities $S^{\star}_v>1$ and $\nu\ge3$ of Definition~\ref{5.5:def-box-read-units}; the two
scaffold conditions on which that definition rests, together with the scaffold majorant
$|B^{ax}(b)|_{\star}$ defined there; the margin statement and the complex-analyticity statement of
the transfer of the analysis of \cite[\S3]{B-CMP109} to the cover, together with the remaining
condition of that transfer on the own box, applied on $\mathfrak{r}_v$; the response numbers
$\vartheta$, $\vartheta(\square)$ and the constant $\beta$ that enter
Definition~\ref{5.5:def-remainder-constants}; the four properties of the source field $\mathbb{U}$
recalled in Notation~\ref{5.5:not-interpolated-box-field}; and the three bounds of the variational
statement from which the constants $r_V$, $C_V$ of Definition~\ref{5.5:def-remainder-constants} are
taken. They use none of the displays \cite[(3.26)--(3.27)]{B-CMP109}, not the constant
$\mathfrak{a}_0$, not the quantity $G_{\square}$, not the second-derivative room of $T$, no weight
other than $\Theta$ inside $S^{\star}_v$ and $w$ inside the bound
$\mathsf{a}_{Q_v}\le2K_{\mathrm{loc}}\mathsf{A}_w$, no count of the cells of $\mathfrak{B}$, no age of
a large-field region, no depth, and none of the numbers $N$, $N_0$, $\check R$.
\end{remark}

\begin{remark}[Remainders of plaquette and counterterm pieces without membership]
\label{5.5:rem-counterterm-remainders}
Let $\square$ be a block and let $v$ be a near unit of $\square$ in the sense of
Definition~\ref{5.5:def-box-read-units}. No condition is placed on $S^{\star}_v$ or on $\nu$,
so $v$ need not be box-read. A \emph{plaquette or counterterm piece} at $v$ is a term
\begin{equation*}
T_v = \kappa_v\,A(c_v,\cdot),
\end{equation*}
where $c_v$ is a complex plaquette weight, $A(c_v,\cdot)$ is the Wilson action with the weight
$c_v$ in the sense of Lemma~\ref{5.2:lem-complex-weights}, and $\kappa_v$ is the coefficient of
the piece. By Definition~\ref{5.5:def-box-read-units}, such a term is filed on the two box pairs
$\mathsf{P}^{0}$ and $\mathsf{P}^{1}(t,\mathsf{t})$ of Definition~\ref{5.5:def-stored-remainder}
at every near unit of $\square$; the box pairs there are built from the block $\square$ and the
level $j$, not from the box-read condition. Put
\begin{equation}\label{5.5:eq-counterterm-remainders-1}
\operatorname{rem}^{W}(v,\square)
:= \int_0^1 dt\;\partial_{\mathsf{t}}\big|_{\mathsf{t}=0}
\Big[T_v\big(\mathsf{P}^{1}(t,\mathsf{t})\big)-T_v\big(\mathsf{P}^{0}\big)\Big].
\end{equation}
Assume the following two statements:
\begin{itemize}
\item[(i)] the theorem on plaquette and counterterm pieces at near units, clauses (a)--(f),
with the constant $K^{W}$ of its per-unit bound and the two numbers $\mathsf{s}_v$,
$\mathsf{v}_v$ that this bound attaches to the unit $v$;
\item[(ii)] its corollary on block sums, with its block-sum constant $c^{W}_{\Sigma}$ and its
size number $\mathfrak{s}^{W}$, and under the one smallness condition listed there.
\end{itemize}
Then the following hold for every near unit $v$ of $\square$ and every plaquette or
counterterm piece $T_v$ at $v$.
\begin{itemize}
\item[(1)] $\operatorname{rem}^{W}(v,\square)$ is defined without any membership
condition on the box pairs. This is in contrast with clause (a) of
Theorem~\ref{5.5:thm-membership-remainder}, which is stated with its proof incomplete at one
step.
\item[(2)] The difference $T_v(\mathsf{P}^{1})-T_v(\mathsf{P}^{0})$ is holomorphic in the sense
of clause (b) of that theorem.
\item[(3)] $\operatorname{rem}^{W}(v,\square)$ has the real jet of clause (c) of that
theorem.
\item[(4)] For each run of the construction separately, the per-unit bound
\eqref{5.5:eq-counterterm-remainders-2} holds.
\item[(5)] For each run of the construction separately, the block sum
\eqref{5.5:eq-counterterm-remainders-3} holds.
\end{itemize}
The two bounds in (4) and (5) are
\begin{equation}\label{5.5:eq-counterterm-remainders-2}
\big|T_v(\mathsf{P}^{1})-T_v(\mathsf{P}^{0})\big|
\le K^{W}\,\mathsf{s}_v\,\mathsf{v}_v\,e^{-\frac14\delta D_v}\,\alpha_{\cdot,j},
\end{equation}
\begin{equation}\label{5.5:eq-counterterm-remainders-3}
\sum_{v}\big|\operatorname{rem}^{W}(v,\square)\big|
\le \Big[\frac{1}{r^{\sharp}_{\square}}+14\,K^{\mathrm{rem}}\,\vartheta(\square)\Big]
K^{W}\,c^{W}_{\Sigma}\,\mathfrak{s}^{W}\,
\sup\alpha_{\cdot,j}.
\end{equation}
The sum in \eqref{5.5:eq-counterterm-remainders-3} runs over the near units of $\square$, and
the supremum is the one taken in the corollary of hypothesis (ii). These
two bounds have the format of clauses (d) and (e) of
Theorem~\ref{5.5:thm-membership-remainder}, compare \eqref{5.5:eq-membership-remainder-a},
with $c^{W}_{\Sigma}$ and $\mathfrak{s}^{W}$ in the places of $c_{\Sigma}$ and
$\mathfrak{s}^{\mathrm{rd}}$, but the rate is $e^{-\frac14\delta D_v}$.

Four ingredients lead to \eqref{5.5:eq-counterterm-remainders-2}. The first is the two
relations of the sealing statement, taken at the layered box field $\mathfrak{b}$ and the
correction form $\mathsf{E}$. The
second is the invariance of the trace under conjugation, which removes the regauging $\tilde v$.
The third is an estimate on each plaquette: the difference of the two plaquette
variables is bounded by the curvature times the layer plus the square of the layer. The fourth
is clause (a2) of the near-unit estimate, used pointwise.

This remark makes no assertion about the difference of these remainders between two runs of
the construction on one scaffold, that is, about the analogue for plaquette and counterterm
pieces of clause (g) of Theorem~\ref{5.5:thm-membership-remainder}. That difference is governed
by the two-run bound for this class, which is proved as an implication from a response
statement one half of which is not proved: the combination of the two-cutoff response statement
at zeroth order with the energy-difference bound. Nothing from it is used here.
\end{remark}

\begin{proof}
We derive (1)--(5) from hypotheses (i) and (ii).

For (1), the term $T_v=\kappa_v A(c_v,\cdot)$ is $\kappa_v$ times an entire function of the bond
variables. Hence $T_v(\mathsf{P}^{1}(t,\mathsf{t}))$ and $T_v(\mathsf{P}^{0})$ are defined
wherever the box pairs are defined. No membership of the restricted pair in a domain of the form
$\mathfrak{U}_T$ is needed, and the integrand of \eqref{5.5:eq-counterterm-remainders-1} is
defined as soon as the $\mathsf{t}$-derivative at $\mathsf{t}=0$ exists.

That derivative exists by the holomorphy in (2). Statement (2) is the holomorphy clause of
hypothesis (i), statement (3) its real-jet clause, and the bound
\eqref{5.5:eq-counterterm-remainders-2} its bound clause. These clauses are stated for every
near unit of $\square$, whatever $S^{\star}_v$ and $\nu$ are, and apply to each piece filed on
the box pairs by Definition~\ref{5.5:def-box-read-units}.

The block sum \eqref{5.5:eq-counterterm-remainders-3} is the conclusion of hypothesis (ii),
under its one listed smallness condition.
\end{proof}

\begin{corollary}[Membership of thin layers on the box field]\label{5.5:cor-thin-layers}
Work in the setting of Lemma~\ref{5.5:lem-loop-data}: $v=(T,\square)$ is the near unit, $\mathfrak{r}_v$
its own box, $y$ the base point from which the comb paths of $\mathfrak{B}_j(\mathfrak{r}_v)$ start,
and $\mathsf{P}^{0}$ the first of the two box pairs, with field component $U_{\mathsf{P}^{0}}$ and
current component $J_{\mathsf{P}^{0}}$. Let $X$ be the domain on which the members of
$\mathfrak{U}_T$ are defined. Assume the following statements.
\begin{itemize}
\item[(a)] Lemma~\ref{5.5:lem-near-unitary-transports}, with the window representation, the box field $W$ of its clause~(i), and the transport
bound.
\item[(b)] Lemma~\ref{5.5:lem-loop-data}, with the corresponding bounds on the
canonical loop datum $B^{0}$ and on its working form $B^{0}_u$.
\item[(c)] The charts of the two box fields by analytic continuation,
\eqref{5.5:eq-box-field-charts-a}.
\item[(d)] The regauging bound of the transport lemma with constants $r_V$, $C_V$.
\item[(e)] The bump lemma for the stored remainder in its second frame, together with its window
condition and its containment condition.
\item[(f)] The identity for $\Xi$ used in the analysis of the stored remainder, and the
Baker--Campbell--Hausdorff bound used there.
\item[(g)] Clause~(iv) of the increment bound for weight profiles in
$\mathfrak{M}^{\wedge}_{\delta/2}(\mathfrak{B}_j)$.
\item[(h)] Conditions (ii) and (iii) of \eqref{5.5:eq-remainder-constants-b}.
\item[(i)] The clause bounds for the representative of \cite[Lemma~4]{B-CMP109}: let $\xi$ be the
parameter with which the configuration of that lemma is written as $e^{i\xi\mathbf{H}_j(\cdot,B)}$;
for every $\mathfrak{g}^{c}$-valued datum $B$ on the bonds of $\mathfrak{B}_j(\mathfrak{r}_v)$ with
$|B(b)|<|B^{ax}(b)|_{\star}$ at every bond $b$, the configuration
\begin{equation*}
R_B:=\bigl(e^{i\xi\mathbf{H}_j(\cdot,B)},\ J_{R_B}\bigr),
\end{equation*}
with $J_{R_B}$ its own current, satisfies $Q_a(R_B)<\theta_a$ for $a\in\{1,2,3,4,6,7\}$.
\end{itemize}
Let $\mathcal{K}$ be a $\mathfrak{g}^{c}$-valued $1$-form and $\mathbf{j}$ a
$\mathfrak{g}^{c}$-valued bond function on $X$, both in the native gauge of $\mathsf{P}^{0}$, and
let $\varkappa\in\mathfrak{M}^{\wedge}_{\delta/2}(\mathfrak{B}_j)$ satisfy, for every point $y'$,
with $\tilde{\Delta}(y')$ the enlarged block of $y'$ and $\nabla_W$ the covariant derivative in the
box field, taken in the gauge in which $\mathcal{K}$ is given,
\begin{equation}\label{5.5:eq-thin-layers-1}
\sup_{\tilde{\Delta}(y')}\max\bigl\{\ell_j|\mathcal{K}|,\ \ell_j^{2}|\nabla_W\mathcal{K}|\bigr\}
\le\varkappa(y'),\qquad
\ell_j^{3}|\mathbf{j}|\le K_{\mathrm{inc}}\,\varkappa(y'),
\end{equation}
and
\begin{equation}\label{5.5:eq-thin-layers-2}
12K_{\mathrm{inc}}\sup_X\varkappa\le\alpha_{\cdot,j}.
\end{equation}
Then
\begin{equation}\label{5.5:eq-thin-layers-3}
\bigl(e^{i\eta\mathcal{K}}U_{\mathsf{P}^{0}},\ J_{\mathsf{P}^{0}}+\mathbf{j}\bigr)\restriction X
\in\mathfrak{U}_T .
\end{equation}
In particular, the layered extensions of $\mathsf{P}^{0}$ that enter the near terms, including the
exact layered configuration of the remainder, are members of $\mathfrak{U}_T$ as soon as their
layers obey the profile \eqref{5.5:eq-thin-layers-1}--\eqref{5.5:eq-thin-layers-2} in the native
gauge of $\mathsf{P}^{0}$; that they do is not shown here.
\end{corollary}

\begin{proof}
We work in the working gauge of Notation~\ref{5.5:not-working-gauge}, with $u=v_w^{-1}$. Put
$\mathcal{K}_u:=\operatorname{Ad}_u\mathcal{K}$ and $\mathbf{j}_u:=R(u)\mathbf{j}$. By the transfer
of estimates in Notation~\ref{5.5:not-working-gauge}, each such move costs at most a factor $1.2$
on $|\cdot|$, so $\mathcal{K}_u$ and $\mathbf{j}_u$ obey \eqref{5.5:eq-thin-layers-1} with
$\varkappa$ replaced by $1.2\varkappa$. Throughout, $L\ge 8$ is the floor on $L$ subject to which the
constants of Definition~\ref{5.5:def-remainder-constants} are chosen.

\smallskip\noindent\textit{Step 1 (the gauge).}
By the charts \eqref{5.5:eq-box-field-charts-a}, $W=Y(B_0)$, and $W\restriction\mathfrak{r}_v$ is
critical under its own $\mathfrak{B}_j(\mathfrak{r}_v)$-data
$\mathsf{V}'_u=M_{\cdot}(W)\restriction\mathfrak{B}_j(\mathfrak{r}_v)=\Phi(B_0)$ and is small.
Hence, by \cite[Theorem~1]{B-CMP102b}, it lies in the critical orbit under the group of
\cite[(4)]{B-CMP102b}:
\begin{equation*}
W\restriction\mathfrak{r}_v=U_j(\mathfrak{r}_v,\mathsf{V}'_u)^{g_2},
\qquad g_2=g_2(B_0)=1\ \text{at base points}.
\end{equation*}
Let $\bar v$ be the gauge transformation of \cite[(181)]{B-CMP102b} determined by the comb gauge of
$\mathsf{V}'_u$ from $y$, that is, by $x\mapsto\mathsf{V}'_u(\Gamma_{y,x})^{-1}$ on base points; in
particular $\bar v(y)=1$. By \cite[(181)]{B-CMP102b},
\begin{equation*}
U_j(\mathfrak{r}_v,\mathsf{V}'_u)=U_j(\mathfrak{r}_v,e^{iB^{0}_u})^{\bar v^{-1}} .
\end{equation*}
Put $H:=\mathbf{H}_j(\mathfrak{r}_v,B^{0}_u)$, with $\xi$ the parameter of hypothesis~(i) of the
corollary, and let $u_B$ be the regauging of \cite[Proposition~9]{B-CMP102b},
\cite[(3.37), (3.50)]{B-CMP109}; then
$U_j(\mathfrak{r}_v,e^{iB^{0}_u})=(e^{i\xi H})^{u_B}$. By the pointwise bound in
the corresponding display we have $|B^{0}_u(b)|\le 1.2L^{-1}|B^{ax}(b)|_{\star}$, and the last
inequality there, $1.42L^{-1}\sup_b|B^{ax}(b)|_{\star}<1.42\alpha_{3,j}/L$, says
$\sup_b|B^{ax}(b)|_{\star}<\alpha_{3,j}$; thus $|B^{0}_u|_\infty\le 1.2\alpha_{3,j}/L$. With
$\mathsf{b}$ ranging over the blocks of $\mathfrak{r}_v$ at scale $j$, the regauging bound (d), the
bound $\ell_j\sup_{\mathsf{b}}|\mathcal{H}(B)|\le B_3|B|_\infty$ of \cite[(3.38)]{B-CMP109}, and
condition~(iii)
of \eqref{5.5:eq-remainder-constants-b}, which gives $C_VB_3\alpha_{3,j}\le 0.01$, yield
\begin{equation}\label{5.5:eq-thin-layers-4}
|u_B(x)-1|\le C_V\ell_j\sup_{\mathsf{b}}|\mathcal{H}(B^{0}_u)|\le C_VB_3|B^{0}_u|_\infty
\le \frac{1.2\,C_VB_3\alpha_{3,j}}{L}\le\frac{0.012}{L}\le 0.002 .
\end{equation}
Consequently
\begin{equation*}
W^{g_v}=e^{i\xi H},\qquad g_v:=u_B^{-1}\,\bar v\,g_2^{-1}.
\end{equation*}
Since $\bar v(y)=1$ and $g_2(y)=1$, we have $n(g_v(y))=n(u_B(y))$. By
\eqref{5.5:eq-thin-layers-4} and the Neumann series,
$\|u_B(y)\|\le 1.002$ and $\|u_B(y)^{-1}\|\le(1-0.002)^{-1}$, so $n(g_v(y))\le 1.01$. Only this
bound enters below; the normalisation $g_v(y)=1$ is not used.

For $x\in\mathfrak{r}_v$ let $\gamma_x\subset\mathfrak{r}_v^{\flat}$ be the comb path from $y$ to
$x$; its level-$j$ length satisfies $\operatorname{len}_j(\gamma_x)\le 2dM/t_p\le\lambda_v$. The
relation $W^{g_v}=e^{i\xi H}$, read on transports along $\gamma_x$, gives
\begin{equation*}
g_v(x)=e^{i\xi H}(\gamma_x)^{-1}\,g_v(y)\,W(\gamma_x).
\end{equation*}
Since $\|abc\|\le\|a\|\|b\|\|c\|$ for the operator norm, $n$ is submultiplicative and
$n(a^{-1})=n(a)$; with the bound $n(e^{i\xi H}(\gamma_x))\le e^{2\int_{\gamma_x}|H|}$ and
the corresponding display, which bounds $n(W(\gamma_x))$ by $1.45$ for paths of
length at most $\lambda_v$,
\begin{equation*}
n(g_v(x))\le e^{2\int_{\gamma_x}|H|}\cdot 1.01\cdot 1.45 .
\end{equation*}
For every bond $b$ of $\mathfrak{B}_j(\mathfrak{r}_v)$ one has $1+|b-y|\le 3M/t_p$, so the
pointwise bound in the corresponding display,
$|B^{0}_u(b)|\le 1.2\cdot 2d(1+|b-y|)c_1\mathsf{a}_{Q_v}t_p^{2}$, gives
$|B^{0}_u|_\infty\le 1.2\cdot 6dMc_1\mathsf{a}_{Q_v}t_p$; with \cite[(3.38)]{B-CMP109} this gives
\begin{equation*}
|H|\le B_3|B^{0}_u|_\infty\le 1.2\,B_3\cdot 6dMc_1\mathsf{a}_{Q_v}t_p .
\end{equation*}
Multiplying by the length $2dM/t_p$ and using the bound
$\mathsf{a}_{Q_v}\le 2K_{\mathrm{loc}}\mathsf{A}_w(\gamma)$ of the weight $\mathsf{a}_{Q_v}$ by the
scaffold majorant $\mathsf{A}_w(\gamma)$ gives
\begin{equation*}
\int_{\gamma_x}|H|\le 29\,d^{2}M^{2}B_3c_1K_{\mathrm{loc}}\mathsf{A}_w(\gamma),
\end{equation*}
which is at most $29/220\le 0.14$ by condition~(ii) of \eqref{5.5:eq-remainder-constants-b}.
Hence, for every $x\in\mathfrak{r}_v$,
\begin{equation}\label{5.5:eq-thin-layers-5}
n(g_v(x))\le e^{0.28}\cdot 1.01\cdot 1.45\le 2 .
\end{equation}

\smallskip\noindent\textit{Step 2 (the representative without layer, by the Schwarz lemma).}
The proof of \cite[Lemma~4]{B-CMP109} is carried out on the representative $R_B$ of
hypothesis~(i) of the corollary; of this configuration it states: ``We will prove that this
configuration satisfies all the conditions (i)--(iv)''.
Write $J_R:=J_{R_{B^{0}_u}}$. For each clause functional $Q_a$
of $\mathfrak{U}_T$ with $a\in\{1,2,3,4,6,7\}$, write $Q_a(R_{zB^{0}_u})=\|F_a(z)\|$, where $F_a$
is holomorphic in $z$, since $\mathbf{H}_j$ is holomorphic in its datum. The function $F_a$
vanishes at $z=0$, because $R_0=(\mathbf{1},0)$. For $|z|<L/1.2$ the pointwise bound in
the corresponding display gives $|zB^{0}_u(b)|<|B^{ax}(b)|_{\star}$ at every bond $b$, and for
such data hypothesis~(i) of the corollary gives $\|F_a(z)\|<\theta_a$. The Schwarz
lemma, applied to $\varphi\circ F_a$ for every bounded linear functional $\varphi$ of norm at most
one on the target of $F_a$, then gives at $z=1$
\begin{equation}\label{5.5:eq-thin-layers-6}
Q_a(R_{B^{0}_u})\le\frac{1.2\,\theta_a}{L}\le 0.15\,\theta_a ,
\end{equation}
the last step by $L\ge 8$. Moreover $Q_5(R_{B^{0}_u})=Q_8(R_{B^{0}_u})=0$, and the part of
condition~(iv) of \cite[Lemma~4]{B-CMP109} that concerns the field is void, since the $G$-valued
part of the representative is identically $1$.

\smallskip\noindent\textit{Step 3 (increments).}
Put $\mathcal{K}':=\operatorname{Ad}_{g_v}\mathcal{K}_u$ and $\mathbf{j}':=R(g_v)\mathbf{j}_u$.
Since $\nabla_{W^{g}}\operatorname{Ad}_g=\operatorname{Ad}_g\nabla_W$, and since
$\|\operatorname{Ad}_g\|\le n(g)$ and $\|R(g)\|\le n(g)$, the bound
\eqref{5.5:eq-thin-layers-5} and the first paragraph of the proof show that $\mathcal{K}'$,
$\nabla_{e^{i\xi H}}\mathcal{K}'$ and $\mathbf{j}'$ obey \eqref{5.5:eq-thin-layers-1} with
$\varkappa$ replaced by $2\cdot 1.2\varkappa=2.4\varkappa$. We bound, clause by clause, the
increment of $Q_a$ from $R_{B^{0}_u}$ to
$(e^{i\eta\mathcal{K}'}e^{i\xi H},\ J_R+\mathbf{j}')\restriction X$.

\smallskip\noindent\emph{Entry~$1$.} This entry is controlled by the identity for $\Xi$ in (f),
which gives
\begin{equation*}
|\Xi-1|\le 4\xi^{2}|\nabla\mathcal{K}'|+8\xi\,|\partial e^{i\xi H}-1|\,|\mathcal{K}'|
+64\xi^{2}|\mathcal{K}'|^{2}.
\end{equation*}

\smallskip\noindent\emph{Entries~$6$ and~$7$.} These are bounded by the Baker--Campbell--Hausdorff
bound in (f), by at most $3\cdot 2.4\varkappa$.

\smallskip\noindent\emph{Entry~$3$.} The increment of this entry is $|\mathbf{j}'|$.

\smallskip\noindent\emph{Entries~$2$ and~$4$.} These entries, at the structures of $T$ itself, are
treated by the bump lemma (e) in its second frame, applied with $(\mathbb{W},\mathcal{K},\varkappa)$
replaced by $(e^{i\xi H},\mathcal{K}',2.4\varkappa)$. Its window condition holds in the form
$a_{\flat}\le 1.2B_3\alpha_{3,j}\le\alpha_{1,j}$, and its containment condition in the form
$B(x)\subset X$. These entries are moved by at most $\pi(E+r)\le 2\pi\mathbf{r}^T_a$, with $E$ and
$r$ the two terms of the displacement bound in the second frame of the bump lemma.

\smallskip\noindent
As in clause~(iv) of the increment bound (g), each of these increments is at most
\begin{equation}\label{5.5:eq-thin-layers-7}
K_{\mathrm{inc}}\cdot 2.4\,\varkappa(y)\,\alpha_{\cdot,j}^{-1}\,\mathbf{r}^T_a
\le 0.2\,\mathbf{r}^T_a\le 0.25\,\theta_a ,
\end{equation}
where the first inequality is \eqref{5.5:eq-thin-layers-2} and the second uses
$\theta_a\ge(1-\beta)\mathbf{r}^T_a$.

\smallskip\noindent\textit{Step 4 (conclusion).}
By \eqref{5.5:eq-thin-layers-6} and \eqref{5.5:eq-thin-layers-7},
$0.15\,\theta_a+0.25\,\theta_a<\theta_a$ for every clause, so
\begin{equation*}
\bigl(e^{i\eta\mathcal{K}'}e^{i\xi H},\ J_R+\mathbf{j}'\bigr)\restriction X\in\mathfrak{U}_T .
\end{equation*}
By the gauge transformation formulas \cite[(1.10)]{B-CMP109} for fields and currents, the
configuration $(e^{i\eta\mathcal{K}}U_{\mathsf{P}^{0}},\ J_{\mathsf{P}^{0}}+\mathbf{j})\restriction
X$ is the transform of this member by $(g_vu)^{-1}$: it lies in the same $G^{c}$-orbit. This is
\eqref{5.5:eq-thin-layers-3}.
\end{proof}

\begin{remark}[Ba{\l}aban's removed terms and an alternative route]\label{5.5:rem-removed-terms}
Let $v=(T,\square)$ be a box-read near unit, and let $\operatorname{rem}(T,\square)$ be its stored
remainder in the sense of Definition~\ref{5.5:def-stored-remainder}, built from the difference
$G_v=T(X;\mathsf{P}^{1})-T(X;\mathsf{P}^{0})$ of the near term at the two box pairs.
\begin{enumerate}
\item[(0)] At the real point, where $\sigma\equiv 1$ and $\mathsf{t}$ is real and small,
$\operatorname{rem}(T,\square)$ is the $t_{\square}$-jet (the jet in the interpolation parameter
$t_{\square}$ of \cite[\S 3]{B-CMP109}, evaluated at $t_{\square}=\mathsf{t}$), in the sense of
clause~(c) of
the theorem on membership, holomorphy and bound of the remainder
(Theorem~\ref{5.5:thm-membership-remainder}, whose proof is incomplete at the step marked $(\ast)$;
every appeal to it in this remark and its proof is subject to this), of the sum of the two terms
removed at
\cite[(3.21), (3.22)]{B-CMP109}. Other holomorphic extensions with the same real jet exist, among
them the layered configurations named at the end of Corollary~\ref{5.5:cor-thin-layers}. The
extension used here is the one given by $G_v$, and at the places where such an extension
enters, the membership family $\mathsf{C}(\sigma_1)$ of clause~(a) of
Theorem~\ref{5.5:thm-membership-remainder} is the one used.
The treatment here therefore differs from that of \cite[\S 3]{B-CMP109} in the choice of
extension and agrees with it in the real jet.
\item[(1)] In \cite[(3.21)]{B-CMP109} the second term is a function of $U_{k+1}$, and it is
treated, in the words of \cite[p.~274]{B-CMP109}, ``in exactly the same way as the terms (3.7)''.
Transcribed to the present setting, that treatment extends the term on the input field $U_{k+1}$
of the step, and its bound rests on the theorem furnishing the fields $\mathbb{H}^{w}$ and
$\mathbf{H}_j(t,\mathsf{t})$, applied on the set attached to $\square_0$ on which that theorem is
stated. That treatment is not used
here; the difference $G_v$ is used instead.
\item[(2)] The corresponding display of the canonical loop data of the two box fields
(Lemma~\ref{5.5:lem-loop-data}) rests on the constant $c_W$ carrying the factor $L^{3}$, as in its
definition. Were $c_W$ taken with the factor $L^{2}$ instead, the membership statement of clause~(a) of
Theorem~\ref{5.5:thm-membership-remainder} and Corollary~\ref{5.5:cor-thin-layers} would require the
margin furnished by a lemma of the present construction to be available on the ball of radius
$\tfrac32\alpha_{3,j}$. That estimate on the ball of radius $\tfrac32\alpha_{3,j}$ is not
proved here, and this second route is not taken. In neither case is the condition $S^{\star}_v>1$
changed.
\item[(3)] Theorem~\ref{5.5:thm-membership-remainder} is stated, and is applied, in one and the same
form at every place of use, for the partition function with sources and for
test functions of the class $\mathfrak{T}(\mathsf{d};\vartheta)$ alike; its proof
is incomplete at the step marked $(\ast)$ in both cases.
\end{enumerate}
\end{remark}

\begin{proof}
Only item~(0) calls for an argument; items~(1)--(3) record an alternative treatment, an alternative
route and the scope of Theorem~\ref{5.5:thm-membership-remainder}. At the real point, the
identification of $\operatorname{rem}(T,\square)$ with the
$t_{\square}$-jet of the sum of the two terms removed at \cite[(3.21), (3.22)]{B-CMP109} is clause~(c)
of Theorem~\ref{5.5:thm-membership-remainder}. By
Definition~\ref{5.5:def-stored-remainder},
\begin{equation*}
\operatorname{rem}(T,\square)=\int_0^1 dt\,\partial_{\mathsf{t}}\big|_{0}\,G_v(t,\mathsf{t}),
\end{equation*}
so the real jet only records $\partial_{\mathsf{t}}G_v$ at $\mathsf{t}=0$ and $\sigma\equiv 1$. This
real jet does not determine a holomorphic function on the domain of clause~(b) of
Theorem~\ref{5.5:thm-membership-remainder}: the layered configurations named at the end of
Corollary~\ref{5.5:cor-thin-layers} have the same real jet. The extension taken here is $G_v$.
By clause~(a) of the same theorem,
the pairs $\mathsf{C}(0)\cong\mathsf{P}^{0}\restriction X$ and
$\mathsf{C}(1)\cong\mathsf{P}^{1}\restriction X$ lie in $\mathfrak{U}_T$. Thus the family
$\mathsf{C}(\sigma_1)$, $|\sigma_1|\le e^{\frac18\delta D_v}$, supplies the configurations at which the
near term is evaluated wherever an extension of the real jet enters.
\end{proof}

\begin{definition}[Cells, adjacency and linear sizes of cell sets]\label{5.6:def-linear-sizes}
Throughout this section $d\ge 1$ is the dimension. The case of interest is $d=4$, but nothing
below uses this value.

\emph{Lattice and cells.} Fix a level $k$. Here $M$ is the cube side of the glossary, and of it
only the property that $M$ is an integer with $M\ge 1$ is used. The \emph{level-$k$ lattice
unit} is the unit of length in which a level-$k$ cell has integer side $M$. In this unit the
lattice is a box or a torus in $\mathbb{Z}^d$. Its \emph{sites} are its points and its
\emph{bonds} are its nearest-neighbour pairs of sites. The \emph{cells} are the $M$-cubes of the
level-$k$ partition, in the sense of the $M$-cubes of \cite[(1.67)]{B-CMP122b}:
\begin{equation}\label{5.6:eq-linear-sizes-1}
  \Delta_a=\prod_{\mu=1}^{d}\bigl([Ma_\mu,\,Ma_\mu+M-1]\cap\mathbb{Z}\bigr),
  \qquad a\in\mathbb{Z}^d .
\end{equation}
These cells form a grid that partitions the set of sites. For a site $x$ we write $a(x)$ for
the index of the cell containing $x$.

Some statements of this section use cells only as labels; none of this geometry is used there.
The grid geometry enters only through the adjacency relation and the linear sizes defined
below, and only where these are invoked.

\emph{Adjacency and connectedness.} Let $\Delta_a\neq\Delta_{a'}$ be two cells, and let
$e_1,\dots,e_d$ be the coordinate vectors.
\begin{itemize}
\item They are \emph{face-adjacent} if $a-a'=\pm e_\nu$ for one coordinate vector $e_\nu$. Then
  they share a $(d-1)$-dimensional face.
\item They are \emph{closure-touching} if $|a_\mu-a'_\mu|\le 1$ for every $\mu$.
\end{itemize}
The symbol $\sim$ denotes either relation. Each statement that uses it fixes the choice once and
for all. We write $D_d$ for the number of $\sim$-neighbours of a cell of the grid indexed by
$\mathbb{Z}^d$: thus $D_d=2d$ for face-adjacency and $D_d=3^d-1$ for closure-touching.

A finite set $Y$ of cells is \emph{$\sim$-connected} if the graph on $Y$ whose edges are the
pairs $\Delta\sim\Delta'$ is connected. For $Y_1\subset Y$, the phrase ``every $\sim$-component
of $Y$ meets $Y_1$'' means the following: every cell of $Y$ is joined to some cell of $Y_1$ by a
$\sim$-path of cells of $Y$. This condition holds automatically when $Y$ is $\sim$-connected and
$Y_1\neq\emptyset$. The \emph{outer boundary} of a set $Y_1$ of cells is
\begin{equation}\label{5.6:eq-linear-sizes-2}
  \partial^{+}Y_1=\{\Delta\notin Y_1:\ \Delta\sim\Delta' \text{ for some } \Delta'\in Y_1\}.
\end{equation}
It satisfies $\#\partial^{+}Y_1\le D_d\,\#Y_1$.

\emph{Tree graphs and lengths.} A \emph{tree graph} $T$ is a finite connected subgraph of the
lattice, that is, a set of sites together with bonds between them; a single site is allowed.
We write $|T|$ for the number of bonds of $T$; this is its length in the $1$-norm, in level-$k$
units. We say that $T$ \emph{meets} a cell $\Delta$ if some site of $T$ lies in $\Delta$. For a
finite nonempty set $Y$ of cells we define the \emph{free tree length}
\begin{equation}\label{5.6:eq-linear-sizes-3}
  d_k^{\mathrm{fr}}(Y)=M^{-1}\min\{|T|:\ T \text{ meets every cell of } Y\},
\end{equation}
and the \emph{contained tree length}
\begin{equation}\label{5.6:eq-linear-sizes-4}
  d_k^{\subset}(Y)=M^{-1}\min\Bigl\{|T|:\ T \text{ meets every cell of } Y
  \text{ and every site of } T \text{ lies in } \textstyle\bigcup Y\Bigr\}.
\end{equation}
We set $d_k^{\subset}(Y)=+\infty$ if no such $T$ exists. The relative linear size of
\cite[(1.67)]{B-CMP122b} is of the contained kind; there it is taken relative to a large-field
region $Z$.

\emph{Linear sizes.} Ba{\l}aban's lengths $d_k(\cdot)$ and $d_j(\cdot)$, which appear in
\cite[(1.66)]{B-CMP122b} and are defined in \cite[(0.26)]{B-CMP109}, do not enter the formulation.
Instead, the statements of this section that involve a length are formulated for an arbitrary
\emph{linear size}. This is a map
\[
  \ell:\ \{\text{finite sets of cells}\}\to[0,\infty]
\]
with the following property $(\mathrm{TL})_{c'}$, for a constant $c'$. Whenever
$Y_1\subset Y$ are finite sets of cells such that every $\sim$-component of $Y$ meets $Y_1$,
\begin{equation}\label{5.6:eq-linear-sizes-a}
  \ell(Y)\le \ell(Y_1)+c'\,\#(Y\setminus Y_1).
\end{equation}
A statement that applies the results of this section supplies its own $\ell$, its own $c'$ and
its own choice of $\sim$. Property \eqref{5.6:eq-linear-sizes-a} is established, for two
choices, in the lemma on the adjunction property of the tree lengths:
\begin{itemize}
\item $\ell=d_k^{\mathrm{fr}}$ with $c'=d$, for either adjacency;
\item $\ell=d_k^{\subset}$ with $c'=1$, for face-adjacency.
\end{itemize}
\end{definition}

\begin{proof}
We verify the values of $D_d$, the remark on the connected case, and the bound on
$\#\partial^{+}Y_1$.

\emph{The values of $D_d$.} The cells face-adjacent to $\Delta_a$ are the $\Delta_{a\pm e_\nu}$
with $\nu=1,\dots,d$. These are $2d$ cells. The cells closure-touching $\Delta_a$ are the
$\Delta_{a+b}$ with $b\in\{-1,0,1\}^d$ and $b\neq 0$. These are $3^d-1$ cells. In a box or a
torus the neighbours of a cell are among these index translates, so every cell there has at
most $D_d$ neighbours.

\emph{The connected case.} Suppose $Y$ is $\sim$-connected and $Y_1\neq\emptyset$. Choose a cell
$\Delta'\in Y_1\subset Y$. Since the graph on $Y$ is connected, every cell of $Y$ is joined to
$\Delta'$ by a $\sim$-path of cells of $Y$. Hence every $\sim$-component of $Y$ meets $Y_1$.

\emph{The size of the outer boundary.} By \eqref{5.6:eq-linear-sizes-2}, every cell of
$\partial^{+}Y_1$ is a $\sim$-neighbour of some cell of $Y_1$. Therefore
\[
  \partial^{+}Y_1\subset\bigcup_{\Delta'\in Y_1}\{\Delta:\ \Delta\sim\Delta'\}.
\]
Each set in this union has at most $D_d$ elements, and the union has $\#Y_1$ members. Hence
\[
  \#\partial^{+}Y_1\le D_d\,\#Y_1. \qedhere
\]
\end{proof}

\begin{definition}[Decoupling variables, polydiscs and the composite of a term]
\label{5.6:def-decoupling-variables}
We fix the variables, the domains, the data and the function on which the hypothesis
$(\mathrm{H}_{\circ})$, stated below in this section, bears.
\begin{enumerate}
\item[(a)] \emph{Decoupling variables.} A term $t$, in the sense of item~(d) below, comes with a
finite set $\mathcal{V} = \mathcal{V}(t)$ of indices and with complex variables
$\sigma_v$, $v \in \mathcal{V}$, called the \emph{decoupling variables} of $t$. They are the
parameters $s(\Delta)$ of \cite[(1.7)--(1.11)]{B-CMP116}, where they enter as the parameters
$s(Y_0)$. Along the chain of domains $Y_1, Y_2, \dots$ of \cite[(1.61)--(1.63)]{B-CMP122b},
whose first member is the exempt set $Y_1(t)$ of item~(d), each link $i$ of the chain
carries its own family $s_i(\Delta) = s(Y_i)(\Delta)$. When one variable
per cell is meant, $\mathcal{V} =
\mathcal{N}$ is a finite set of cells in the sense of
Definition~\ref{5.6:def-linear-sizes}, and $\sigma_v = \sigma(\Delta)$ for the cell
$\Delta = v$.

\item[(b)] \emph{Polydiscs and holomorphy.} Throughout, $\kappa_1 \ge 1$, so that
$e^{\kappa_1} \ge e > 1$. For a finite set $W$, for $\rho > 0$ and for radii
$r = (r_w)_{w \in W}$ with $r_w > 0$, put
\begin{align}
P_W(\rho) &:= \{\sigma \in \mathbb{C}^W : |\sigma_w| < \rho \text{ for all } w \in W\},
\label{5.6:eq-decoupling-variables-1}\\
\mathbb{D}_W(\rho) &:= \{\sigma \in \mathbb{C}^W : |\sigma_w| \le \rho \text{ for all }
w \in W\},
\label{5.6:eq-decoupling-variables-2}\\
T_W(r) &:= \prod_{w \in W} \{\sigma_w \in \mathbb{C} : |\sigma_w| = r_w\}.
\label{5.6:eq-decoupling-variables-3}
\end{align}
Thus $P_W(\rho)$ is the open polydisc, $\mathbb{D}_W(\rho)$ the closed polydisc and $T_W(r)$
the torus with radii $r$. A function $G : P_W(\rho) \to \mathbb{C}$ is called
\emph{holomorphic} if it is continuous and if, for each $w \in W$ and each fixed value of the
coordinates $\sigma_{w'}$, $w' \neq w$, it is holomorphic in $\sigma_w$ in the sense of one
complex variable. This is all that is used in what follows. A jointly complex-differentiable
function, for instance a composite of holomorphic maps, is holomorphic in this sense.

Since $e^{\kappa_1} > 1$, the real unit cube $[0,1]^W$ is contained in
$P_W(e^{\kappa_1})$. At a point of $[0,1]^W$ the partial derivative
$\partial_w := \partial/\partial\sigma_w$ taken along the real direction coincides with the
complex derivative in $\sigma_w$. The \emph{mixed first derivative} is
\begin{equation}
\partial^W := \prod_{w \in W} \partial_w;
\label{5.6:eq-decoupling-variables-4}
\end{equation}
that the order of its factors is immaterial is shown in the proof of the lemma on
Cauchy estimates in this section.

\item[(c)] \emph{Data.} The symbol $\mathfrak{d}$ denotes a \emph{datum}, ranging over a set
$\mathcal{D}_b$ of data of size at most $b$. In the applications, a datum is the complex
displacement of the data of a configuration of local form, its size being measured by a norm
on these data at the scale of the piece, or the datum attached to one
localised piece of the expansion; the letter $b$ denotes the data size, and the letter
$\mathfrak{b}$ is reserved for the supremum of the composite of item~(d) on a polydisc.
Nothing about $\mathcal{D}_b$ is used in what follows except that the holomorphy hypothesis
$(\mathrm{H}_{\circ})$, stated below in this section, holds for
each $\mathfrak{d} \in \mathcal{D}_b$. We write
\begin{equation*}
\sup_{\mathrm{size} \le b} (\,\cdot\,) := \sup_{\mathfrak{d} \in \mathcal{D}_b} (\,\cdot\,).
\end{equation*}

\item[(d)] \emph{The term, its exempt set, the composite.} A \emph{term} $t$ is one analytic
functional of the effective action, with a localisation domain $X$ and
an \emph{exempt set}
$Y_1 = Y_1(t)$: a finite nonempty set of cells that carry no decoupling variable of $t$.
In \cite[(1.61)--(1.65)]{B-CMP122b} the exempt set is built from $Y_0(X)$, the smallest domain
of the class $\mathbf{D}_k$ of localisation domains containing $X$ \cite[p.~375]{B-CMP122b}.
Outside the large-field region $Z$ it coincides with that domain \cite{B-CMP122b}, that is,
$Y_1 \setminus Z = Y_0(X) \setminus Z$.

Let $\mathsf{R}(\sigma;\mathfrak{d})$ denote the $\sigma$-interpolated, localised configuration
with data $\mathfrak{d}$; this is the configuration of \cite[(1.61)--(1.63)]{B-CMP122b},
taken with the parameters $s(Y_i)$. The \emph{composite} of
the term $t$ is the value of $t$ at this configuration,
\begin{equation}
F_t(\sigma;\mathfrak{d}) := t \circ \mathsf{R}(\sigma;\mathfrak{d}).
\label{5.6:eq-decoupling-variables-5}
\end{equation}
It is a function of $(\sigma,\mathfrak{d})$, with $\sigma$ in its domain in
$\mathbb{C}^{\mathcal{V}}$ and $\mathfrak{d} \in \mathcal{D}_b$. What is assumed of it is
exactly the hypothesis $(\mathrm{H}_{\circ})$; in particular nothing is assumed of the
dependence of $F_t$ on $\mathfrak{d}$ beyond what $(\mathrm{H}_{\circ})$ states for each
fixed $\mathfrak{d}$.
\end{enumerate}
\end{definition}

\begin{definition}[Pieces of a term and their counting forms]\label{5.6:def-pieces}
Let $t$ be a term with finite set $\mathcal{V}=\mathcal{V}(t)$ of decoupling variables and composite
$F_t(\sigma;\mathfrak{d})=t\circ\mathsf{R}(\sigma;\mathfrak{d})$, as in
Definition~\ref{5.6:def-decoupling-variables}.

\emph{Pieces.} A \emph{piece} of $t$ with \emph{variable set} $W\subset\mathcal{V}$ is a function of the
datum $\mathfrak{d}$ of the form
\begin{equation}\label{5.6:eq-pieces-a}
p(\mathfrak{d})=\int_{[0,1]^{\mathcal{V}}}\mu(d\sigma)\,(\partial^W F_t)(\sigma;\mathfrak{d}),
\end{equation}
where $\mu$ is a Borel probability measure on the real unit cube $[0,1]^{\mathcal{V}}$. The set $W$ and
the measure $\mu$ are part of the data of the piece and do not depend on $\mathfrak{d}$.

\emph{Product form.} A piece is of \emph{product form} if $\mu$ is the product of Lebesgue measure on
$[0,1]^W$ with the unit point mass $\delta_s$ at a fixed base point $s\in[0,1]^{\mathcal{V}\setminus W}$;
then
\begin{equation}\label{5.6:eq-pieces-1}
p(\mathfrak{d})=\int_{[0,1]^W}\Bigl(\prod_{w\in W}d\sigma_w\,\partial_{\sigma_w}\Bigr)
F_t(\sigma_W,s;\mathfrak{d}).
\end{equation}
The interpolations $\int ds\,\partial_s$ of \cite[(1.61)--(1.66)]{B-CMP122b} are of this form, with base
point $s\equiv1$ on the cells of the earlier links and $s\equiv0$ off the final domain (links and final
domains are introduced in the chain form below). The definition
\eqref{5.6:eq-pieces-a} also admits interpolations in which undifferentiated variables vary.

\emph{Chain form.} Fix an integer $\lambda\ge1$, the number of \emph{links}, indexed by
$i\in I=\{2,\dots,\lambda+1\}$; in \cite[(1.61)--(1.66)]{B-CMP122b} one has $I=\{2,3,4\}$, $\lambda=3$.
For a \emph{final domain} $Y\supseteq Y_1(t)$ let $\mathcal{N}'(t;Y)$ be the set of cells of
$Y\setminus Y_1(t)$ and $n=n(t;Y)=\#\mathcal{N}'(t;Y)$. A chain
\begin{equation*}
Y_1(t)=Y_1\subset Y_2\subset\cdots\subset Y_{1+\lambda}=Y
\end{equation*}
has increments $Y'_i=Y_i\setminus Y_{i-1}$, $i\in I$, which form an ordered partition of
$\mathcal{N}'(t;Y)$ into $\lambda$ possibly empty parts; the admissible chains are a subset of all such
chains. Let $\mathcal{P}(t;Y)$ be the set of pieces of $t$ with final domain $Y$, the final domain of a
piece being the last member $Y$ of its chain; in \cite{B-CMP122b} this is the localization domain of the
term, in the sense of the sentence preceding \cite[(1.72)]{B-CMP122b}. The family
$\mathcal{P}(t;Y)$ is said to be \emph{of chain form} if it has the three properties
\begin{equation}\label{5.6:eq-pieces-b}
\begin{aligned}
&(\pi 1)\quad&&\text{every $p\in\mathcal{P}(t;Y)$ is a piece in the sense of
\eqref{5.6:eq-pieces-a} (in \cite{B-CMP122b}, of product form),}\\
&&&\text{whose variable set $W(p)$ consists of one variable $\sigma_i(\Delta)$ for each cell}\\
&&&\text{$\Delta\in Y'_i$ and each $i\in I$;}\\
&(\pi 2)&&\text{for fixed $(t,Y)$ the map $p\mapsto(Y'_i)_{i\in I}$ is injective from
$\mathcal{P}(t;Y)$}\\
&&&\text{into the ordered partitions of $\mathcal{N}'(t;Y)$ into $\lambda$ parts;}\\
&(\pi 3)&&\text{$\mathcal{P}(t;Y)=\emptyset$ unless $Y\supseteq Y_1(t)$ and every
$\sim$-component of $Y$ meets $Y_1(t)$,}\\
&&&\text{with $\sim$ the cell adjacency of Definition~\ref{5.6:def-linear-sizes}.}
\end{aligned}
\end{equation}
These three properties are abstracted from \cite[(1.61)--(1.66)]{B-CMP122b}, in which the factors are
one per $M$-cube of each $Y'_i$ and the final domains are connected domains containing a prescribed
smallest domain; $(\pi 3)$ takes $Y_1(t)$ in the role of that domain. They are not proved here and
enter as hypotheses wherever they are used. In counting form they give
\begin{equation}\label{5.6:eq-pieces-2}
\begin{gathered}
|W(p)|\ge n(t;Y)\quad\text{for every }p\in\mathcal{P}(t;Y),\qquad
\#\mathcal{P}(t;Y)\le\lambda^{\,n(t;Y)},\\
\mathcal{P}(t;Y)=\emptyset\ \text{unless $Y\supseteq Y_1(t)$ and every $\sim$-component of $Y$
meets $Y_1(t)$,}
\end{gathered}
\end{equation}
the two bounds in the first line following from $(\pi 1)$ and $(\pi 2)$ as shown below, the second
line being $(\pi 3)$. In what follows, only the
two counting consequences in the first line of \eqref{5.6:eq-pieces-2} are used from $(\pi 1)$ and
$(\pi 2)$, and $(\pi 3)$ is used only where it is invoked explicitly.

\emph{Subset form.} Let $\mathcal{V}=\mathcal{N}$, one variable $\sigma(\Delta)$ per cell $\Delta$, and
write $F=F_t$. For $\mathsf{y}\subset\mathcal{N}$ put
\begin{equation}\label{5.6:eq-pieces-c}
F_{\mathsf{y}}(\mathfrak{d})=\int_{[0,1]^{\mathsf{y}}}
\Bigl(\prod_{\Delta\in\mathsf{y}}d\sigma(\Delta)\,\partial_{\sigma(\Delta)}\Bigr)
F(\sigma_{\mathsf{y}},0_{\mathcal{N}\setminus\mathsf{y}};\mathfrak{d}),
\qquad F_{\emptyset}(\mathfrak{d})=F(0;\mathfrak{d}).
\end{equation}
This is the piece of product form \eqref{5.6:eq-pieces-1} with $W=\mathsf{y}$ and base point $0$ off
$\mathsf{y}$.

That $\partial^W F_t(\cdot;\mathfrak{d})$ is continuous on $[0,1]^{\mathcal{V}}$ under the holomorphy
hypothesis $(\mathrm{H}_{\circ})$ of Definition~\ref{5.6:def-decoupling-variables}, so that the
integrals \eqref{5.6:eq-pieces-a},
\eqref{5.6:eq-pieces-1} and \eqref{5.6:eq-pieces-c} exist, follows from the Cauchy estimates for the
composite on the decoupling polydisc proved later in this section; it is used in the proof of clause (a)
of the weight lemma.
\end{definition}

\begin{proof}[Proof of the counting consequences in \eqref{5.6:eq-pieces-2}]
Let $p\in\mathcal{P}(t;Y)$ with chain increments $(Y'_i)_{i\in I}$. By $(\pi 1)$, $W(p)$ contains exactly
one variable for each pair $(i,\Delta)$ with $i\in I$ and $\Delta\in Y'_i$, so
$|W(p)|=\sum_{i\in I}\#Y'_i$. Since the $Y'_i$ are pairwise disjoint with union $\mathcal{N}'(t;Y)$, this
sum equals $n(t;Y)$; in particular $|W(p)|\ge n(t;Y)$.

For the second bound, an ordered partition of $\mathcal{N}'(t;Y)$ into $\lambda$ possibly empty parts
$(Y'_i)_{i\in I}$ is the same as a map $\mathcal{N}'(t;Y)\to I$, namely the map sending a cell to the
index of the part containing it; hence there are exactly $\lambda^{n(t;Y)}$ such partitions. Restricting
to admissible chains only removes partitions. By $(\pi 2)$ the map $p\mapsto(Y'_i)_{i\in I}$ is
injective from $\mathcal{P}(t;Y)$ into this set, so $\#\mathcal{P}(t;Y)\le\lambda^{n(t;Y)}$.

The last line of \eqref{5.6:eq-pieces-2} is $(\pi 3)$ restated.
\end{proof}

\begin{remark}[The exponential factors of Ba{\l}aban's bound in the present notation
{\cite[(1.66)]{B-CMP122b}}]
\label{5.6:rem-exponential-factors}
This is \cite[(1.66)]{B-CMP122b}, read in the notation of this section; it serves only to
identify Ba{\l}aban's symbols with those of this section and is not used as a step in any
proof. Let $t$ be a term of level $j<k$ with
domain $X$, considered at a chain of increments $Y'_2,Y'_3,Y'_4$ with final domain $Y_4$
(the chain form of the pieces, Definition~\ref{5.6:def-pieces}). For a cell set $Y$ write
$\#Y$ for the number of its cells, that is, of the $M$-cubes it contains. Write $Y_0(X)$
for the smallest domain of Ba{\l}aban's class of level-$k$ domains \cite{B-CMP122b} that
contains $X$, and put
$a=(1+3\beta)\kappa$. The bound of \cite[(1.66)]{B-CMP122b} carries three kinds of
exponential factor.

The inductive factor is
\begin{equation*}
  \exp\bigl(-\beta\kappa\, d_j(X)\bigr)\,\exp\bigl(-a\, d_k(Y_0(X))\bigr)
  \qquad\text{for } j<k,
\end{equation*}
and for a term of level $j=k$ it is instead $\exp\bigl(-(1+4\beta)\kappa\, d_k(X)\bigr)$.

For each $i=2,3,4$ there is a per-cell factor
\begin{equation*}
  \exp\bigl(-(\kappa_1-1)\,\#Y'_i\bigr).
\end{equation*}

When $Y'_2\cup Y'_3\cup Y'_4\neq\emptyset$ there is a third factor
\begin{equation*}
  \exp\Bigl(-\delta\, d\bigl((Y_1^0)^c\cap Y_1\cap Y'_{i_0},\,X\bigr)\Bigr),
\end{equation*}
where $i_0$ is the first index with $Y'_{i_0}\neq\emptyset$, $d(\cdot,\cdot)$ is
Ba{\l}aban's scaled distance, and $Y_1^0$, $Y_1$ are the domains so denoted in
\cite{B-CMP122b}, whose definition is not needed here; the letter $Y_1$ in this factor is
Ba{\l}aban's and is not identified here with the exempt set $Y_1(t)$ of $t$ that appears in
the last display below. This third factor is not used in what follows.

In this section the bound $\mathfrak{b}$ is the supremum of the composite $F_t$ of $t$ on
the decoupling polydisc (Definition~\ref{5.6:def-decoupling-variables}). It therefore
contains whatever inductive decay the term $t$ carries. That decay is made explicit below
in the form
\begin{equation*}
  \mathfrak{b}\le \mathfrak{b}_0\, e^{-a\,\ell(Y_1(t))},
\end{equation*}
with $\mathfrak{b}_0$ the part of the bound free of the inductive decay.
\end{remark}

\begin{lemma}[An elementary inequality for the Cauchy factor]
\label{5.6:lem-cauchy-factor}
For a real number $\kappa_1 \ge 1$ put
\begin{equation*}
r(\kappa_1) = e^{\kappa_1}\bigl(e^{\kappa_1}-1\bigr)^{-2}.
\end{equation*}
Then for every $\kappa_1 \ge 1$
\begin{equation}\label{5.6:eq-cauchy-factor-1}
r(\kappa_1) \le e^{-(\kappa_1-1)},
\end{equation}
hence $r(\kappa_1) \le 1$. Moreover
\begin{equation}\label{5.6:eq-cauchy-factor-2}
\bigl(e^{\kappa_1}-1\bigr)^{-1} \le e^{-(\kappa_1-1)}.
\end{equation}
\end{lemma}

\begin{proof}
Put $x = e^{\kappa_1}$. Since $\kappa_1 \ge 1$ we have $x \ge e > 1$, so $x - 1 > 0$, and
$e^{-(\kappa_1-1)} = e/x$.

We prove \eqref{5.6:eq-cauchy-factor-1}. In terms of $x$ it reads $x/(x-1)^2 \le e/x$. Multiplying by
the positive number $x(x-1)^2$, this is equivalent to $x^2 \le e(x-1)^2$. Both $x$ and
$\sqrt{e}\,(x-1)$ are positive, so taking positive square roots, this is equivalent to
$x \le \sqrt{e}\,(x-1)$, that is, to
\begin{equation*}
x\bigl(\sqrt{e}-1\bigr) \ge \sqrt{e}.
\end{equation*}
Since $\sqrt{e} - 1 > 0$, the left side increases with $x$; as $x \ge e$, it therefore suffices to
verify the inequality at $x = e$, which corresponds to $\kappa_1 = 1$. At $x = e$ it reads
$e(\sqrt{e}-1) \ge \sqrt{e}$. Multiplying out, this is $e\sqrt{e} - e \ge \sqrt{e}$, which after
moving terms is $e\sqrt{e} - \sqrt{e} \ge e$, that is, $\sqrt{e}\,(e-1) \ge e$. Dividing by
$\sqrt{e} > 0$, this is equivalent to $e - 1 \ge \sqrt{e}$. Both sides are positive, so squaring gives
the equivalent inequality $(e-1)^2 \ge e$, that is,
\begin{equation*}
e^2 - 3e + 1 \ge 0 .
\end{equation*}
The polynomial $y^2 - 3y + 1$ has the roots $(3 \pm \sqrt{5})/2$ and is non-negative exactly for $y$
at most the smaller root or at least the larger one. The smaller root satisfies
$(3-\sqrt{5})/2 < 1$ because $\sqrt{5} > 1$, while $e > 3/2 > 1$. Hence $e$ lies beyond the smaller
root, and the last inequality is equivalent to $e \ge (3+\sqrt{5})/2$.

It remains to check this. First, $\sqrt{5} < 9/4$, because both numbers are positive and
$(9/4)^2 = 81/16 > 80/16 = 5$. Therefore
\begin{equation*}
\frac{3+\sqrt{5}}{2} < \frac{3 + 9/4}{2} = \frac{21}{8}.
\end{equation*}
Second, all terms of the series $e = \sum_{n \ge 0} 1/n!$ are positive, so $e$ exceeds its partial sum
$1 + 1 + \tfrac12 + \tfrac16 = 8/3$; and $8/3 = 64/24 > 63/24 = 21/8$. Together,
\begin{equation*}
e > \frac{8}{3} > \frac{21}{8} > \frac{3+\sqrt{5}}{2}.
\end{equation*}
Hence $e \ge (3+\sqrt{5})/2$, and by the chain of equivalences above the inequality
\eqref{5.6:eq-cauchy-factor-1} holds for all $\kappa_1 \ge 1$. Since $\kappa_1 - 1 \ge 0$, we have
$e^{-(\kappa_1-1)} \le 1$, and therefore $r(\kappa_1) \le e^{-(\kappa_1-1)} \le 1$.

We prove \eqref{5.6:eq-cauchy-factor-2}. In terms of $x$ it reads $(x-1)^{-1} \le e/x$. Multiplying by
the positive number $x(x-1)$, this is equivalent to $x \le e(x-1)$, that is, to $x(e-1) \ge e$. Since
$e - 1 > 0$, the left side increases with $x$, and as $x \ge e$ it suffices to check $x = e$. There the
inequality reads $e(e-1) \ge e$, which after division by $e > 0$ is $e \ge 2$; this holds since
$e > 8/3 > 2$.

Finally, the chain of equivalences in the proof of \eqref{5.6:eq-cauchy-factor-1} shows that this
inequality holds exactly when $x \ge \sqrt{e}/(\sqrt{e}-1)$, that is, when
$\kappa_1 \ge \log\bigl(\sqrt{e}/(\sqrt{e}-1)\bigr)$; since the inequality $x(\sqrt{e}-1) \ge \sqrt{e}$
was shown to hold strictly at $x = e$, this threshold is less than $1$.
\end{proof}

\begin{lemma}[Cauchy estimates for mixed first derivatives on a polydisc]
\label{5.6:lem-cauchy-estimates}
Let $W$ be a finite set, let $R>1$, and let $G\colon P_W(R)\to\mathbb{C}$ be continuous and separately
holomorphic, with $|G|\le B$ on $P_W(R)$. Then $G$ is $C^\infty$ on $P_W(R)$, $\partial^W G$ is
continuous there, and for every $\sigma\in P_W(R)$
\begin{equation}\label{5.6:eq-cauchy-estimates-a}
|\partial^W G(\sigma)|\le B\prod_{w\in W} R\,(R-|\sigma_w|)^{-2};
\end{equation}
in particular
\begin{equation*}
|\partial^W G(\sigma)|\le B\,\bigl[R(R-1)^{-2}\bigr]^{|W|}\qquad\text{for }\sigma\in[0,1]^W.
\end{equation*}
Moreover, let $W\neq\emptyset$ and define the mixed difference of $G$ over the corners of the unit
cube by
\begin{equation*}
\Delta^W G:=\sum_{\mathsf{z}\subset W}(-1)^{|W\setminus\mathsf{z}|}\,
G(1_{\mathsf{z}},0_{W\setminus\mathsf{z}}),
\end{equation*}
where $\mathsf{z}$ runs over all subsets of $W$, the empty set and $W$ included, and
$(1_{\mathsf{z}},0_{W\setminus\mathsf{z}})\in\{0,1\}^W$ is the point whose coordinates are
$1$ on $\mathsf{z}$ and $0$ on $W\setminus\mathsf{z}$. Then
\begin{equation}\label{5.6:eq-cauchy-estimates-b}
|\Delta^W G|\le B\,(R-1)^{-|W|}.
\end{equation}
\end{lemma}

Here $P_W(\rho)$, $T_W(r)$ and $\partial^W=\prod_{w\in W}\partial_w$ are as in
Definition~\ref{5.6:def-decoupling-variables}; $\partial_w$ is the complex derivative in the variable
$\sigma_w$, which at real points coincides with the derivative along the real direction. For radii
$r=(r_w)_{w\in W}$ and a continuous function $\varphi$ on $T_W(r)$ we write
$\int_{T_W(r)}\varphi(\zeta)\,d\zeta$, with $d\zeta:=\prod_{w\in W}d\zeta_w$, for the contour
integral over the torus formed by the positively oriented circles $|\zeta_w|=r_w$, that is, with
$d\theta:=\prod_{w\in W}d\theta_w$,
\begin{equation*}
\int_{T_W(r)}\varphi(\zeta)\,d\zeta
=\int_{[0,2\pi]^W}\varphi\bigl((r_we^{i\theta_w})_{w\in W}\bigr)
\prod_{w\in W}\bigl(i\,r_we^{i\theta_w}\bigr)\,d\theta ,
\end{equation*}
so that
\begin{equation*}
\Bigl|\int_{T_W(r)}\varphi\,d\zeta\Bigr|\le\prod_{w\in W}(2\pi r_w)\cdot\sup_{T_W(r)}|\varphi| .
\end{equation*}

\begin{proof}
If $W=\emptyset$, then $P_W(R)$ is a single point, $G$ is a constant of modulus at most $B$,
$\partial^W G=G$ and the product in \eqref{5.6:eq-cauchy-estimates-a} is empty; there is nothing to
prove. Let $W\neq\emptyset$ and fix radii $r=(r_w)_{w\in W}$ with $1<r_w<R$ for every $w$; when
\eqref{5.6:eq-cauchy-estimates-a} is to be proved at a given $\sigma\in P_W(R)$ we require in addition
$r_w>|\sigma_w|$ for every $w$. Such radii exist because $\max\{1,|\sigma_w|\}<R$.

\emph{Step 1: representation.} We show that for every $\tau\in P_W(r)$
\begin{equation}\label{5.6:eq-cauchy-estimates-1}
G(\tau)=(2\pi i)^{-|W|}\int_{T_W(r)}G(\zeta)\prod_{w\in W}(\zeta_w-\tau_w)^{-1}\,d\zeta .
\end{equation}
The integrand is continuous on the compact torus $T_W(r)$, since $G$ is continuous and
$|\zeta_w-\tau_w|\ge r_w-|\tau_w|>0$ there. We argue by induction on $|W|$, the statement being
\eqref{5.6:eq-cauchy-estimates-1} for all $G$ satisfying the hypotheses of the lemma on $P_W(R)$, all
radii $r$ as above and all $\tau\in P_W(r)$.

If $|W|=1$, $G$ is a holomorphic function of one variable on the open disc $\{|\zeta|<R\}$, which
contains the closed disc $\{|\zeta|\le r_w\}$, and \eqref{5.6:eq-cauchy-estimates-1} is Cauchy's
integral formula in one complex variable.

Let $|W|\ge2$, choose $w\in W$ and put $W':=W\setminus\{w\}$. Fix $\tau\in P_W(r)$. The slice
$\zeta_w\mapsto G(\zeta_w,\tau_{W'})$ is holomorphic on $|\zeta_w|<R$ by separate holomorphy, and
$|\tau_w|<r_w<R$, so Cauchy's integral formula in the variable $\zeta_w$ gives
\begin{equation*}
G(\tau)=(2\pi i)^{-1}\oint_{|\zeta_w|=r_w}G(\zeta_w,\tau_{W'})\,(\zeta_w-\tau_w)^{-1}\,d\zeta_w .
\end{equation*}
For each $\zeta_w$ on this circle, the function $\tau_{W'}\mapsto G(\zeta_w,\tau_{W'})$ is continuous,
separately holomorphic and bounded by $B$ on $P_{W'}(R)$, and $\tau_{W'}\in P_{W'}(r_{W'})$ with
$1<r_{w'}<R$ for $w'\in W'$. By the induction hypothesis it equals its $(|W|-1)$-fold Cauchy integral
over $T_{W'}(r_{W'})$:
\begin{equation*}
G(\zeta_w,\tau_{W'})=(2\pi i)^{-(|W|-1)}\int_{T_{W'}(r_{W'})}G(\zeta_w,\zeta_{W'})
\prod_{w'\in W'}(\zeta_{w'}-\tau_{w'})^{-1}\,d\zeta_{W'} .
\end{equation*}
Substituting this into the previous display expresses $G(\tau)$ as $(2\pi i)^{-|W|}$ times an iterated
integral, first over $T_{W'}(r_{W'})$ and then over the circle $|\zeta_w|=r_w$, of the function
$G(\zeta)\prod_{v\in W}(\zeta_v-\tau_v)^{-1}$. This function is continuous on the compact torus
$T_W(r)$, so by Fubini's theorem the iterated integral equals the integral over $T_W(r)$, which is
\eqref{5.6:eq-cauchy-estimates-1}.

\emph{Step 2: differentiation.} For $\zeta\in T_W(r)$ the kernel
$\tau\mapsto\prod_{w\in W}(\zeta_w-\tau_w)^{-1}$ is a $C^\infty$ function on $P_W(r)$, since
$|\tau_w|<r_w=|\zeta_w|$ for every $w$; it is holomorphic in each $\tau_w$, and all its
$\tau$-derivatives are continuous in $(\zeta,\tau)$ on $T_W(r)\times K$, hence bounded there, for every
compact $K\subset P_W(r)$. The function $G$ is continuous and bounded on $T_W(r)$. Differentiating
under the integral sign in \eqref{5.6:eq-cauchy-estimates-1} therefore shows that
$G\in C^\infty(P_W(r))$, and, applying $\partial_w$ once for each $w\in W$ and using
$\tfrac{d}{d\tau}(\zeta-\tau)^{-1}=(\zeta-\tau)^{-2}$,
\begin{equation}\label{5.6:eq-cauchy-estimates-2}
\partial^W G(\tau)=(2\pi i)^{-|W|}\int_{T_W(r)}G(\zeta)\prod_{w\in W}(\zeta_w-\tau_w)^{-2}\,d\zeta,
\qquad \tau\in P_W(r).
\end{equation}
Each $\partial_w$ acts only on the factor $(\zeta_w-\tau_w)^{-1}$ of the product kernel, so the order
in which the $\partial_w$ are applied is immaterial. The right side of
\eqref{5.6:eq-cauchy-estimates-2} is continuous in $\tau$, because the integrand is continuous on
$T_W(r)\times K$ for every compact $K\subset P_W(r)$. Every point of $P_W(R)$ lies in $P_W(r)$ for
some admissible choice of radii (take $\max\{1,|\sigma_w|\}<r_w<R$), so $G$ is $C^\infty$ and
$\partial^W G$ is continuous on $P_W(R)$.

\emph{Step 3: bounds.} Let $\sigma\in P_W(R)$ and let $r$ satisfy, in addition, $r_w>|\sigma_w|$ for
all $w$. On the circle $|\zeta_w|=r_w$ we have $|\zeta_w-\sigma_w|\ge r_w-|\sigma_w|>0$, and its
length is $2\pi r_w$; moreover $|G|\le B$ on $T_W(r)\subset P_W(R)$. Hence
\eqref{5.6:eq-cauchy-estimates-2} at $\tau=\sigma$ gives
\begin{equation*}
|\partial^W G(\sigma)|\le(2\pi)^{-|W|}\prod_{w\in W}\bigl[2\pi r_w\,(r_w-|\sigma_w|)^{-2}\bigr]\cdot B
=B\prod_{w\in W}r_w\,(r_w-|\sigma_w|)^{-2}.
\end{equation*}
The left side does not depend on $r$. Letting $r_w\uparrow R$ for every $w$ and using the continuity
of $\rho\mapsto\rho\,(\rho-|\sigma_w|)^{-2}$ on $\rho>|\sigma_w|$, we obtain
\eqref{5.6:eq-cauchy-estimates-a}. If $\sigma\in[0,1]^W$, then $|\sigma_w|\le1<R$, so
$R(R-|\sigma_w|)^{-2}\le R(R-1)^{-2}$ for every $w$, and the bound for $\sigma\in[0,1]^W$ follows.

\emph{Proof of \eqref{5.6:eq-cauchy-estimates-b}.} Let $r$ be any radii with $1<r_w<R$. The corners
$(1_{\mathsf{z}},0_{W\setminus\mathsf{z}})\in\{0,1\}^W$ lie in $P_W(r)$ because $r_w>1$, so
\eqref{5.6:eq-cauchy-estimates-1} applies at each of them:
\begin{equation*}
G(1_{\mathsf{z}},0_{W\setminus\mathsf{z}})=(2\pi i)^{-|W|}\int_{T_W(r)}G(\zeta)
\prod_{w\in\mathsf{z}}(\zeta_w-1)^{-1}\prod_{w\in W\setminus\mathsf{z}}(\zeta_w-0)^{-1}\,d\zeta .
\end{equation*}
For complex numbers $a_w,b_w$ one has, by expanding the product,
\begin{equation*}
\sum_{\mathsf{z}\subset W}(-1)^{|W\setminus\mathsf{z}|}\prod_{w\in\mathsf{z}}a_w
\prod_{w\in W\setminus\mathsf{z}}b_w=\prod_{w\in W}(a_w-b_w).
\end{equation*}
Applying this under the integral with
$a_w=(\zeta_w-1)^{-1}$ and $b_w=\zeta_w^{-1}$ yields
\begin{equation*}
\Delta^W G=(2\pi i)^{-|W|}\int_{T_W(r)}G(\zeta)\prod_{w\in W}
\bigl[(\zeta_w-1)^{-1}-\zeta_w^{-1}\bigr]\,d\zeta .
\end{equation*}
On $|\zeta_w|=r_w$ we have $|\zeta_w-1|\ge r_w-1>0$, hence
\begin{equation*}
\bigl|(\zeta_w-1)^{-1}-\zeta_w^{-1}\bigr|=|\zeta_w(\zeta_w-1)|^{-1}\le\bigl(r_w(r_w-1)\bigr)^{-1}.
\end{equation*}
Together with $|G|\le B$ on $T_W(r)$ and the length $2\pi r_w$ of each circle this gives
\begin{equation*}
|\Delta^W G|\le\prod_{w\in W}\bigl[(2\pi)^{-1}\cdot2\pi r_w\cdot(r_w(r_w-1))^{-1}\bigr]\cdot B
=B\prod_{w\in W}(r_w-1)^{-1}.
\end{equation*}
The left side does not depend on $r$; letting $r_w\uparrow R$ for every $w$, the right side tends to
$B(R-1)^{-|W|}$, which is \eqref{5.6:eq-cauchy-estimates-b}.
\end{proof}

\begin{lemma}[The weight lemma]\label{5.6:lem-weight}
Let $t$ be a term with exempt set $Y_1(t)$, set of decoupling variables $\mathcal{V}=\mathcal{V}(t)$,
composite $F_t(\sigma;\mathfrak{d})=t\circ\mathsf{R}(\sigma;\mathfrak{d})$ and set $\mathcal{D}_b$ of data of
size at most $b$, as in Definition~\ref{5.6:def-decoupling-variables}, and put
\begin{equation*}
r(\kappa_1)=e^{\kappa_1}\bigl(e^{\kappa_1}-1\bigr)^{-2}.
\end{equation*}
Assume hypothesis $(\mathrm{H}_{\circ})$:
\begin{equation}\label{5.6:eq-weight-a}
\begin{gathered}
\kappa_1\ge 1,\quad\text{and for every }\mathfrak{d}\in\mathcal{D}_b\text{ the function }
\sigma\mapsto F_t(\sigma;\mathfrak{d})\\
\text{is holomorphic on the open polydisc }P_{\mathcal{V}}(e^{\kappa_1}),\\
\text{with }|F_t(\sigma;\mathfrak{d})|\le\mathfrak{b}\text{ for all }\sigma\in P_{\mathcal{V}}(e^{\kappa_1}).
\end{gathered}
\end{equation}
Here holomorphic means, as in Definition~\ref{5.6:def-decoupling-variables}, continuous and separately
holomorphic in each variable. Holomorphy is required in $\sigma$ for each $\mathfrak{d}$ separately;
nothing (no measurability, continuity or analyticity) is assumed of the dependence on
$\mathfrak{d}$. If $F_t(\cdot;\mathfrak{d})$ is holomorphic on a neighbourhood of the closed polydisc
$\mathbb{D}_{\mathcal{V}}(e^{\kappa_1})$ with $|F_t|\le\mathfrak{b}$ on
$\mathbb{D}_{\mathcal{V}}(e^{\kappa_1})$, then $(\mathrm{H}_{\circ})$ holds by restriction. Since
$\partial^W$ annihilates every function of $\mathfrak{d}$ alone when $W\neq\emptyset$, replacing $F_t$
by $F_t-c(\mathfrak{d})$, for any function $c$ of the datum, changes no piece with variable set
$W\neq\emptyset$; for such pieces $(\mathrm{H}_{\circ})$ may therefore be read with $\mathfrak{b}$ the
supremum of $|F_t-c|$ over $\sigma\in P_{\mathcal{V}}(e^{\kappa_1})$ and $\mathfrak{d}\in\mathcal{D}_b$.

Pieces, final domains $Y$, the sets $\mathcal{P}(t;Y)$, the numbers $n=n(t;Y)$ of cells of
$Y\setminus Y_1(t)$, the number $\lambda$ of links of the chain, and the counting readings
$(\pi 1)\text{--}(\pi 3)$ are those of Definition~\ref{5.6:def-pieces} and \eqref{5.6:eq-pieces-b}.
Then the following hold.
\begin{enumerate}
\item[(a)] (Per-piece factor.) Every piece $p$ of $t$ with variable set $W$, $|W|=m$, obeys
\begin{equation}\label{5.6:eq-weight-b}
\sup_{\mathfrak{d}\in\mathcal{D}_b}|p(\mathfrak{d})|\le\mathfrak{b}\,r(\kappa_1)^m
\le\mathfrak{b}\,e^{-(\kappa_1-1)m}.
\end{equation}
\item[(b)] (Count at a fixed final domain.) Assume $(\pi 1)$ and $(\pi 2)$ of \eqref{5.6:eq-pieces-b}.
Then for every final domain $Y$, with $n=n(t;Y)$,
\begin{equation}\label{5.6:eq-weight-c}
\sum_{p\in\mathcal{P}(t;Y)}\sup_{\mathfrak{d}\in\mathcal{D}_b}|p(\mathfrak{d})|
\le\mathfrak{b}\,\bigl(\lambda\,r(\kappa_1)\bigr)^n
\le\mathfrak{b}\,\bigl(\lambda e^{-(\kappa_1-1)}\bigr)^n .
\end{equation}
In particular the left side is at most $\mathfrak{b}$ when $n=0$ (the exceptional term of
\cite[text following (1.66)]{B-CMP122b}, in which $Y=Y_1(t)$ and the contribution is that of the term
itself); for $n\ge 1$ it is at
most $\mathfrak{b}$ when $\kappa_1\ge 1+\log\lambda$, and at most $\mathfrak{b}\,2^{-n}$ when
$\kappa_1\ge 1+\log(2\lambda)$. In the chain form of \cite[(1.61)--(1.63)]{B-CMP122b}, where
$\lambda=3$, this reads: at most $\mathfrak{b}$ for $\kappa_1\ge 1+\log 3$, and at most
$\mathfrak{b}\,2^{-n}$ for $\kappa_1\ge 1+\log 6$.
\item[(c)] (Decay transfer to the final domain.) Assume $(\pi 1)\text{--}(\pi 3)$ of
\eqref{5.6:eq-pieces-b}, and let $\ell$ be a linear size on cell sets obeying $(\mathrm{TL})_{c'}$,
\eqref{5.6:eq-linear-sizes-a}. Suppose $\mathfrak{b}\le\mathfrak{b}_0e^{-a\ell(Y_1(t))}$ with
$a\ge 0$ (the inductive factor extracted from the exponential of \cite[(1.66)]{B-CMP122b} decays at
the rate $a=(1+3\beta)\kappa$). Then for
every final domain $Y$ with $\mathcal{P}(t;Y)\neq\emptyset$,
\begin{equation}\label{5.6:eq-weight-d}
\sum_{p\in\mathcal{P}(t;Y)}\sup_{\mathfrak{d}\in\mathcal{D}_b}|p(\mathfrak{d})|
\le\mathfrak{b}_0e^{-a\ell(Y)}\bigl(\lambda e^{-(\kappa_1-1-c'a)}\bigr)^n
\le\mathfrak{b}_0e^{-a\ell(Y)},
\end{equation}
the second inequality under the one-sided floor $(\mathrm{F}\text{-}\kappa_1)^{P}$,
\begin{equation}\label{5.6:eq-weight-e}
\kappa_1\ge 1+\log\lambda+c'a,
\end{equation}
imposed on $\kappa_1$ after $\kappa$ has been chosen. At $n=0$ no floor is needed, since then
$Y=Y_1(t)$. With $\lambda=3$, $a\le 2\kappa$ and $c'=d=4$ (free tree length), the floor
\eqref{5.6:eq-weight-e} is implied by $\kappa_1\ge 1+\log 3+8\kappa$; with $c'=1$ (contained tree
length, face adjacency) it is implied by $\kappa_1\ge 1+\log 3+2\kappa$. The lower bound
$\kappa_1\ge 16\kappa+1$ of Definition~\ref{5.1:def-admitted-inputs} implies both as soon as
$8\kappa\ge\log 3$, that is $\kappa\ge(\log 3)/8$, for instance whenever $\kappa\ge 1$; where it would
not, the lower bound on $\kappa_1$ is raised to \eqref{5.6:eq-weight-e}, and raising this lower bound
is admissible. No value is assigned to $\kappa_1$.
\item[(d)] (All final domains.) Assume $(\pi 1)\text{--}(\pi 3)$ of \eqref{5.6:eq-pieces-b}. For
$n\ge 0$ let $N_t(n)$ be the number of final domains $Y$ with $\mathcal{P}(t;Y)\neq\emptyset$ and
$n(t;Y)=n$. Then
\begin{equation}\label{5.6:eq-weight-f}
\begin{gathered}
N_t(0)\le 1,\qquad N_t(n)\le 2^{D_d\cdot\#Y_1(t)}\,(8D_d)^n\quad(n\ge 1),\\
\sum_{Y}\sum_{p\in\mathcal{P}(t;Y)}\sup_{\mathfrak{d}\in\mathcal{D}_b}|p(\mathfrak{d})|
\le 2^{D_d\cdot\#Y_1(t)}\,\mathfrak{b}\sum_{n\ge 0}\bigl(8\lambda D_d\,r(\kappa_1)\bigr)^n\\
\qquad\le 2\cdot 2^{D_d\cdot\#Y_1(t)}\,\mathfrak{b},
\end{gathered}
\end{equation}
the last inequality for $\kappa_1\ge 1+\log(16\lambda D_d)$; for $\lambda=3$ the condition reads
$\kappa_1\ge 1+\log(48D_d)$. If the minimal tree length of $Y_1(t)$ is finite (for
$d_k^{\subset}$, for instance, when $Y_1(t)$ is face-connected; for $d_k^{\mathrm{fr}}$ always), then,
by the count of the cells of a cell set in terms of its tree length,
\begin{equation*}
\#Y_1(t)\le 3^d\bigl(2d_k^{\mathrm{fr}}(Y_1(t))+1\bigr)\le 3^d\bigl(2d_k^{\subset}(Y_1(t))+1\bigr),
\end{equation*}
so that, with $c''_d=3^dD_d\log 2$ and for $\ell\in\{d_k^{\mathrm{fr}},d_k^{\subset}\}$,
\begin{equation*}
2^{D_d\cdot\#Y_1(t)}\le e^{c''_d}\,e^{2c''_d\,\ell(Y_1(t))}.
\end{equation*}
Hence, if $\mathfrak{b}\le\mathfrak{b}_0e^{-a\ell(Y_1(t))}$, the final bound
$2\cdot 2^{D_d\cdot\#Y_1(t)}\,\mathfrak{b}$ of \eqref{5.6:eq-weight-f} satisfies
\begin{equation*}
2\cdot 2^{D_d\cdot\#Y_1(t)}\,\mathfrak{b}
\le 2e^{c''_d}\,\mathfrak{b}_0\,e^{-(a-2c''_d)\ell(Y_1(t))}
\le 2e^{c''_d}\,\mathfrak{b}_0\,e^{-(1-\varepsilon)a\ell(Y_1(t))},
\end{equation*}
the last inequality under the one-sided floor $(\mathrm{F}\text{-}\kappa)^{(d)}$, depending only on
$d$, $\beta$ and $\varepsilon$,
\begin{equation}\label{5.6:eq-weight-g}
\varepsilon a\ge 2c''_d,\quad\text{which for }a=(1+3\beta)\kappa\text{ reads}\quad
\kappa\ge\frac{2\cdot 3^dD_d\log 2}{\varepsilon(1+3\beta)},
\end{equation}
for any fixed $\varepsilon\in\,]0,1[$. Otherwise the prefactor $2^{D_d\cdot\#Y_1(t)}$ is kept
displayed. In either case no bound with a constant uniform in the term $t$ is asserted unless either
the floor \eqref{5.6:eq-weight-g} holds or the prefactor $2^{D_d\cdot\#Y_1(t)}$ is carried explicitly.
Clause (d) is not
needed for the order of summation of \cite[text following (1.66)]{B-CMP122b}, in which the final
domain is fixed first, nor for the per-final-domain bound of the correlation theorem.
\item[(e)] (Subset form.) Let $\mathcal{N}$ be a finite set of cells and let $F$ satisfy
$(\mathrm{H}_{\circ})$ with $\mathcal{V}$ replaced by $\mathcal{N}$ and $\mathfrak{b}$ denoted
$\mathfrak{N}(F)$, or with $\mathfrak{N}(F)$ the supremum of
$|F(\sigma;\mathfrak{d})-F(\mathbb{1};\mathfrak{d})|$ over $\sigma\in P_{\mathcal{N}}(e^{\kappa_1})$
and $\mathfrak{d}\in\mathcal{D}_b$, in the sense of the remark following \eqref{5.6:eq-weight-a},
where $\mathbb{1}$ is the point with all coordinates equal to $1$. Let $F_{\mathsf{y}}$ be the
subset-form pieces of \eqref{5.6:eq-pieces-c}. Then the bound in the first line of the following
display holds for every $\mathsf{y}\subset\mathcal{N}$, $\mathsf{y}\neq\emptyset$, and the
telescoping identity in its second line holds for every $\mathfrak{d}\in\mathcal{D}_b$:
\begin{equation}\label{5.6:eq-weight-h}
\begin{gathered}
\sup_{\mathfrak{d}\in\mathcal{D}_b}|F_{\mathsf{y}}(\mathfrak{d})|
\le\mathfrak{N}(F)\,r(\kappa_1)^{|\mathsf{y}|}\le\mathfrak{N}(F)\,e^{-(\kappa_1-1)|\mathsf{y}|},\\
F(\mathbb{1};\mathfrak{d})-F(0;\mathfrak{d})
=\sum_{\mathsf{y}\subset\mathcal{N},\ \mathsf{y}\neq\emptyset}F_{\mathsf{y}}(\mathfrak{d}).
\end{gathered}
\end{equation}
No disc in any variable other than the decoupling variables enters.
\item[(f)] (Dependence of the constants.) The constants in (a)--(e) are the numbers $r(\kappa_1)$,
$\lambda^n$, $2$, $4$, $8$, $16$, and numbers depending only on $d$ and the adjacency $\sim$, namely
\begin{equation}\label{5.6:eq-weight-i}
c'\in\{1,d\},\qquad D_d\in\{2d,\;3^d-1\},\qquad 3^d,\qquad c''_d=3^dD_d\log 2 .
\end{equation}
The conditions are one-sided lower bounds on $\kappa_1$ and, only where the absorption in (d) is wanted,
on $\kappa$. None of them depends on $K$, $k$, $N$, an age, a nesting depth, a position, the volume,
the size of a large-field region, the number of terms, or the coupling. The blocking factor $L$ enters
nowhere, and no power of $L$ is written or used; $M$ enters only through the definition of the lengths,
where it cancels, the increments in $(\mathrm{TL})_{c'}$ being $c'\cdot M/M$. No property of the group
$SU(N)$ is used.
\end{enumerate}
\end{lemma}

\begin{proof}[Proof of (a)]
Fix $\mathfrak{d}\in\mathcal{D}_b$ and a piece $p$ with data $(W,\mu)$, $|W|=m$; by
Definition~\ref{5.6:def-pieces},
\begin{equation*}
p(\mathfrak{d})=\int_{[0,1]^{\mathcal{V}}}\mu(d\sigma)\,(\partial^WF_t)(\sigma;\mathfrak{d}),
\end{equation*}
with $\mu$ a Borel probability measure on $[0,1]^{\mathcal{V}}$. Since
$\kappa_1\ge 1$ by $(\mathrm{H}_{\circ})$, $e^{\kappa_1}\ge e>1$.

For $s\in[0,1]^{\mathcal{V}\setminus W}$ define the slice
\begin{equation*}
G_s(\sigma_W)=F_t(\sigma_W,s;\mathfrak{d}),\qquad\sigma_W\in P_W(e^{\kappa_1}),
\end{equation*}
where $(\sigma_W,s)$ is the point of $\mathbb{C}^{\mathcal{V}}$ with $W$-coordinates $\sigma_W$ and
$(\mathcal{V}\setminus W)$-coordinates $s$. Every coordinate of $s$ has modulus at most
$1<e^{\kappa_1}$, so
$[0,1]^{\mathcal{V}\setminus W}\times P_W(e^{\kappa_1})\subset P_{\mathcal{V}}(e^{\kappa_1})$. By
$(\mathrm{H}_{\circ})$, $F_t(\cdot;\mathfrak{d})$ is continuous and separately holomorphic on
$P_{\mathcal{V}}(e^{\kappa_1})$ and bounded there by $\mathfrak{b}$; restricting to the set just
described, $G_s$ is continuous and separately holomorphic on $P_W(e^{\kappa_1})$ with
$|G_s|\le\mathfrak{b}$. Since
$\partial^W$ differentiates in the $W$-coordinates only,
\begin{equation*}
\partial^WF_t(\sigma;\mathfrak{d})=\partial^WG_{\sigma_{\mathcal{V}\setminus W}}(\sigma_W)
\qquad\text{for }\sigma\in[0,1]^{\mathcal{V}}.
\end{equation*}
Lemma~\ref{5.6:lem-cauchy-estimates}, applied to $G=G_{\sigma_{\mathcal{V}\setminus W}}$ with
radius $e^{\kappa_1}$ and bound $\mathfrak{b}$ in place of the lemma's $R$ and $B$, gives by
\eqref{5.6:eq-cauchy-estimates-a}, at the point
$\sigma_W\in[0,1]^W$ where $|\sigma_w|\le 1$ for each $w\in W$,
\begin{align*}
|\partial^WF_t(\sigma;\mathfrak{d})|
&\le\mathfrak{b}\prod_{w\in W}e^{\kappa_1}\bigl(e^{\kappa_1}-|\sigma_w|\bigr)^{-2}\\
&\le\mathfrak{b}\,\bigl[e^{\kappa_1}(e^{\kappa_1}-1)^{-2}\bigr]^m=\mathfrak{b}\,r(\kappa_1)^m
\qquad\text{for every }\sigma\in[0,1]^{\mathcal{V}}.
\end{align*}
(For $m=0$ one has $\partial^{\emptyset}F_t=F_t$, and the bound by $\mathfrak{b}=\mathfrak{b}\,r(\kappa_1)^0$
is $(\mathrm{H}_{\circ})$ itself, since $[0,1]^{\mathcal{V}}\subset P_{\mathcal{V}}(e^{\kappa_1})$.)

Lemma~\ref{5.6:lem-cauchy-estimates}, applied to $F_t(\cdot;\mathfrak{d})$ on
$P_{\mathcal{V}}(e^{\kappa_1})$ (the finite set there being $\mathcal{V}$, the radius $e^{\kappa_1}$
and the bound $\mathfrak{b}$), shows that $F_t(\cdot;\mathfrak{d})$
is $C^\infty$ on $P_{\mathcal{V}}(e^{\kappa_1})$. Hence $\partial^WF_t(\cdot;\mathfrak{d})$ is
continuous on the
compact cube $[0,1]^{\mathcal{V}}$, in particular Borel measurable and bounded, and the integral
defining $p(\mathfrak{d})$ exists. Since $\mu$ is a probability measure,
\begin{equation*}
|p(\mathfrak{d})|\le\int_{[0,1]^{\mathcal{V}}}\mu(d\sigma)\,|\partial^WF_t(\sigma;\mathfrak{d})|
\le\mathfrak{b}\,r(\kappa_1)^m .
\end{equation*}
This holds for each $\mathfrak{d}\in\mathcal{D}_b$ separately, with a right side independent of
$\mathfrak{d}$; taking the supremum over $\mathfrak{d}\in\mathcal{D}_b$ gives the first inequality of
\eqref{5.6:eq-weight-b}, and no property of the dependence on $\mathfrak{d}$ is needed. Finally, by
Lemma~\ref{5.6:lem-cauchy-factor}, $0\le r(\kappa_1)\le e^{-(\kappa_1-1)}$ for $\kappa_1\ge 1$; raising
to the power $m\ge 0$ gives $r(\kappa_1)^m\le e^{-(\kappa_1-1)m}$, the second inequality of
\eqref{5.6:eq-weight-b}.

Thus each differentiated variable costs the factor
\begin{equation*}
e^{\kappa_1}\bigl(e^{\kappa_1}-|\sigma_w|\bigr)^{-2}\le e^{\kappa_1}\bigl(e^{\kappa_1}-1\bigr)^{-2}
=r(\kappa_1),
\end{equation*}
the undifferentiated variables being held at their
real values in $[0,1]$, inside the polydisc where the bound $|F_t|\le\mathfrak{b}$ still holds, and
the averages over $[0,1]$ cost the factor $1$.
\end{proof}

Clauses (b)--(f) are proved below in this section.

\begin{remark}[The sharper factor for product-form pieces]\label{5.6:rem-sharper-factor}
Let $t$ be a term with decoupling variables $\mathcal{V}=\mathcal{V}(t)$ and composite
$F_t(\sigma;\mathfrak{d})$, and assume hypothesis $(\mathrm{H}_{\circ})$ of the weight lemma,
\eqref{5.6:eq-weight-a}. Let $p$ be a piece of $t$ with variable set $W\subset\mathcal{V}$, in the model
(product) form of Definition~\ref{5.6:def-pieces}, whose base point $s\in[0,1]^{\mathcal{V}\setminus W}$
has all its coordinates in $\{0,1\}$. Put $m:=|W|$ and, for $\mathfrak{d}\in\mathcal{D}_b$,
$G_s(\sigma_W):=F_t(\sigma_W,s;\mathfrak{d})$. Then $p(\mathfrak{d})=\Delta^W G_s$, the mixed difference
of $G_s$ over the corners of the unit cube, and
\begin{equation}\label{5.6:eq-sharper-factor-1}
\sup_{\mathfrak{d}\in\mathcal{D}_b}|p(\mathfrak{d})|\le\mathfrak{b}\,(e^{\kappa_1}-1)^{-m}.
\end{equation}
This is smaller than $\mathfrak{b}\,r(\kappa_1)^m$ by the factor $[(e^{\kappa_1}-1)/e^{\kappa_1}]^m$, and
$(e^{\kappa_1}-1)^{-1}\le e^{-(\kappa_1-1)}$. The remark is not used in what follows.
\end{remark}

\begin{proof}
Fix $\mathfrak{d}\in\mathcal{D}_b$ and write $R=e^{\kappa_1}>1$. If $W=\emptyset$, then
$p(\mathfrak{d})=F_t(s;\mathfrak{d})$, and $(\mathrm{H}_{\circ})$ gives $|p(\mathfrak{d})|\le\mathfrak{b}$,
which is \eqref{5.6:eq-sharper-factor-1} with $m=0$. Let $W\ne\emptyset$. Since
$s\in[0,1]^{\mathcal{V}\setminus W}$ lies in the polydisc $P_{\mathcal{V}\setminus W}(R)$,
$(\mathrm{H}_{\circ})$ shows that $G_s$ is continuous and separately holomorphic on $P_W(R)$, with
$|G_s|\le\mathfrak{b}$ there. By the Cauchy estimates on a polydisc
(Lemma~\ref{5.6:lem-cauchy-estimates}), $G_s$ is $C^\infty$ on $P_W(R)$, an open neighbourhood of
$[0,1]^W$.

By Definition~\ref{5.6:def-pieces},
$p(\mathfrak{d})=\int_{[0,1]^W}\partial^W G_s(\sigma_W)\,d\sigma_W$, because the derivatives act on the
$W$-variables only. The integrand is continuous on the compact cube, so by Fubini's theorem the integral
may be computed one variable at a time. For $w\in W$, on the real segment the complex derivative
$\partial_{\sigma_w}$ is the derivative in the real variable $\sigma_w$. Hence the fundamental theorem of
calculus replaces $\int_0^1 d\sigma_w\,\partial_{\sigma_w}$ by the difference of the values at
$\sigma_w=1$ and $\sigma_w=0$; the mixed derivatives commute because $G_s$ is smooth. After all $m$
integrations,
\begin{equation}\label{5.6:eq-sharper-factor-2}
p(\mathfrak{d})=\sum_{\mathsf{z}\subset W}(-1)^{|W\setminus\mathsf{z}|}\,
G_s(\mathbb{1}_{\mathsf{z}},0_{W\setminus\mathsf{z}})=\Delta^W G_s.
\end{equation}
The bound \eqref{5.6:eq-cauchy-estimates-b}, applied with $B=\mathfrak{b}$ and $R=e^{\kappa_1}$, gives
$|p(\mathfrak{d})|\le\mathfrak{b}(e^{\kappa_1}-1)^{-m}$. As $\mathfrak{d}$ was arbitrary, this proves
\eqref{5.6:eq-sharper-factor-1}.

For the comparison, recall $r(\kappa_1)=e^{\kappa_1}(e^{\kappa_1}-1)^{-2}$. Then
\begin{equation*}
r(\kappa_1)^m\Big[\frac{e^{\kappa_1}-1}{e^{\kappa_1}}\Big]^m
=e^{m\kappa_1}(e^{\kappa_1}-1)^{-2m}(e^{\kappa_1}-1)^m e^{-m\kappa_1}=(e^{\kappa_1}-1)^{-m},
\end{equation*}
and the bracket is less than $1$. Here $\mathfrak{b}\,r(\kappa_1)^m$ is the bound on the integrand
$|\partial^W G|$ at real points given by Lemma~\ref{5.6:lem-cauchy-estimates}.

Finally, since $\kappa_1\ge1$, the elementary inequality for the Cauchy factor
(Lemma~\ref{5.6:lem-cauchy-factor}) gives both $(e^{\kappa_1}-1)^{-1}\le e^{-(\kappa_1-1)}$ and
$r(\kappa_1)\le e^{-(\kappa_1-1)}$. The majorant $e^{-(\kappa_1-1)}$ per cell is therefore the same
whichever of the two Cauchy estimates is used.
\end{proof}

\begin{proof}[Proof of the count of pieces at a fixed final domain]
\label{5.6:prf-fixed-domain-count}
We prove clause \eqref{5.6:eq-weight-c} of the weight lemma (Lemma~\ref{5.6:lem-weight}). Its
hypotheses are \eqref{5.6:eq-weight-a}, that is, $(\mathrm{H}_{\circ})$, which includes
$\kappa_1\ge 1$, and the first two counting readings $(\pi 1)$--$(\pi 2)$ of
\eqref{5.6:eq-pieces-b} for the pieces of Definition~\ref{5.6:def-pieces}.

Fix a final domain $Y$. Let $\mathcal{N}'(t;Y)$ be the set of cells of $Y\setminus Y_1(t)$, and
let $n=n(t;Y)$ be their number. By $(\pi 1)$, every $p\in\mathcal{P}(t;Y)$ has $|W(p)|\ge n$. By
$(\pi 2)$, the map $p\mapsto(Y'_i)_{i\in I}$ from $\mathcal{P}(t;Y)$ to the ordered
$\lambda$-partitions of $\mathcal{N}'(t;Y)$ is injective, so
$\#\mathcal{P}(t;Y)\le\lambda^n$. The reason for this count is as follows. Each of the $n$ cells
of $Y\setminus Y_1(t)$ enters the chain of domains at exactly one of the $\lambda$ links. Sending
each cell to the index $i\in I$ of the increment $Y'_i$ that contains it identifies an ordered
$\lambda$-partition $(Y'_i)_{i\in I}$ with a map $\mathcal{N}'(t;Y)\to I$, that is, with a choice
function assigning to each cell its link. There are $\lambda^n$ such maps, and admissibility of a
partition only removes some of them.

Since $\kappa_1\ge 1$, the elementary inequality for the Cauchy factor
(Lemma~\ref{5.6:lem-cauchy-factor}) gives
\[
0<r(\kappa_1)\le e^{-(\kappa_1-1)}\le 1 .
\]
By clause \eqref{5.6:eq-weight-b}, each $p\in\mathcal{P}(t;Y)$ obeys
$\sup_{\mathfrak{d}\in\mathcal{D}_b}|p(\mathfrak{d})|\le\mathfrak{b}\,r(\kappa_1)^{|W(p)|}$.
Because $r(\kappa_1)\le 1$ and $|W(p)|\ge n$, this is at most $\mathfrak{b}\,r(\kappa_1)^n$.
Summing over the at most $\lambda^n$ pieces and applying the bound on $r(\kappa_1)$ once more, we
obtain
\begin{equation}\label{5.6:eq-fixed-domain-count-1}
\sum_{p\in\mathcal{P}(t;Y)}\ \sup_{\mathfrak{d}\in\mathcal{D}_b}|p(\mathfrak{d})|
\le\lambda^n\,\mathfrak{b}\,r(\kappa_1)^n
\le\mathfrak{b}\,\bigl(\lambda e^{-(\kappa_1-1)}\bigr)^n .
\end{equation}

If $n=0$, the right-hand side of \eqref{5.6:eq-fixed-domain-count-1} equals $\mathfrak{b}$.

Now let $n\ge 1$. Taking logarithms,
\[
\lambda e^{-(\kappa_1-1)}\le 1
\iff \log\lambda-(\kappa_1-1)\le 0
\iff \kappa_1\ge 1+\log\lambda .
\]
Under this condition, \eqref{5.6:eq-fixed-domain-count-1} is at most $\mathfrak{b}$. In the same
way,
\[
\lambda e^{-(\kappa_1-1)}\le\tfrac12
\iff \log(2\lambda)-(\kappa_1-1)\le 0
\iff \kappa_1\ge 1+\log 2\lambda .
\]
Under this condition, \eqref{5.6:eq-fixed-domain-count-1} is at most $\mathfrak{b}\,2^{-n}$.
These are the two clauses. For $\lambda=3$ the two conditions read $\kappa_1\ge 1+\log 3$ and
$\kappa_1\ge 1+\log 6$, where $1+\log 6=1+\log(2\cdot 3)$.
\end{proof}

\begin{lemma}[Tree lengths under adjunction of cells]\label{5.6:lem-tree-length-adjunction}
Cells, the two adjacency relations, tree graphs and the tree lengths are as in
Definition~\ref{5.6:def-linear-sizes}. Let $\sim$ be one of the two adjacency relations
(closure-touching or face-adjacent). Let $Y_1 \subset Y$ be finite sets of cells with
$Y_1 \neq \emptyset$, such that every $\sim$-component of $Y$ meets $Y_1$, and put
$n := \#(Y\setminus Y_1)$. Then:
\begin{enumerate}
\item[(i)] for $\sim$ either closure-touching or face adjacency,
\begin{equation*}
d_k^{\mathrm{fr}}(Y) \le d_k^{\mathrm{fr}}(Y_1) + d\, n ;
\end{equation*}
\item[(ii)] for $\sim$ face adjacency,
\begin{equation*}
d_k^{\subset}(Y) \le d_k^{\subset}(Y_1) + n .
\end{equation*}
\end{enumerate}
\end{lemma}

\begin{proof}
\emph{Ordering of the adjoined cells.} For a cell of $Y$, consider its distance from $Y_1$ in
the $\sim$-graph of $Y$, that is, the least number of $\sim$-steps in a $\sim$-path of cells
of $Y$ joining it to a cell of $Y_1$. Since every $\sim$-component of $Y$ meets $Y_1$, this
distance is finite for every cell of $Y$, and it is at least $1$ for the cells of
$Y\setminus Y_1$. List the cells of $Y\setminus Y_1$ as $\Delta_1,\dots,\Delta_n$ in
non-decreasing order of this distance (a breadth-first search from $Y_1$). A cell at distance
$\delta \ge 1$ is $\sim$-adjacent to a cell of $Y$ at distance $\delta-1$, namely the
next-to-last cell of a shortest path. If $\delta = 1$ that cell lies in $Y_1$. If
$\delta \ge 2$ it lies in $Y\setminus Y_1$ at strictly smaller distance, and so it was listed
earlier. Hence, for each $j = 1,\dots,n$, the cell $\Delta_j$ is $\sim$-adjacent to some cell
\begin{equation*}
\Delta'_j \in Y_1 \cup \{\Delta_1,\dots,\Delta_{j-1}\}.
\end{equation*}

\emph{Proof of} (i). The free tree length of a finite nonempty set of cells is finite, since a
tree graph passing through one site of each of its cells exists. Let $T_0$ be a tree graph
meeting every cell of $Y_1$ with $|T_0| = M\, d_k^{\mathrm{fr}}(Y_1)$.

We construct $T_1,\dots,T_n$ inductively. Suppose that $T_{j-1} \supseteq T_0$ is connected
and meets every cell of $Y_1\cup\{\Delta_1,\dots,\Delta_{j-1}\}$. Since $\Delta'_j$ belongs to
this set, we may pick a site $x'$ of $T_{j-1}$ in $\Delta'_j$. Write $\Delta'_j = \Delta_{a'}$
and $\Delta_j = \Delta_a$. Under either adjacency relation, $|a_\mu - a'_\mu| \le 1$ for all
$\mu$. Fix $\mu$. Then
\begin{equation*}
x'_\mu \in [Ma'_\mu,\, Ma'_\mu + M - 1],
\end{equation*}
and the distance from $x'_\mu$ to the interval $[Ma_\mu,\, Ma_\mu + M - 1]$ is at most $M$.
Let $x$ be the site of $\Delta_j$ nearest to $x'$; its $\mu$-th coordinate is the point of
$[Ma_\mu, Ma_\mu+M-1]\cap\mathbb{Z}$ nearest to $x'_\mu$. Summing over the $d$ coordinates
gives $\|x - x'\|_1 \le dM$.

Join $x'$ to $x$ by a coordinate-monotone lattice path, which has $\|x-x'\|_1$ bonds, and put
$T_j := T_{j-1} \cup (\text{that path})$. Then $T_j$ is connected, because the path meets
$T_{j-1}$ at $x'$, and $T_j$ meets $\Delta_j$ at $x$. Moreover
\begin{equation*}
|T_j| \le |T_{j-1}| + \|x-x'\|_1 \le |T_{j-1}| + dM .
\end{equation*}
Thus $T_j \supseteq T_0$ is connected and meets every cell of
$Y_1\cup\{\Delta_1,\dots,\Delta_j\}$, and the induction proceeds.

After $n$ steps, $T_n$ is a tree graph meeting every cell of
$Y = Y_1\cup\{\Delta_1,\dots,\Delta_n\}$, with $|T_n| \le |T_0| + n\, dM$. By the definition of
the free tree length,
\begin{equation*}
d_k^{\mathrm{fr}}(Y) \le M^{-1}|T_n| \le M^{-1}|T_0| + dn = d_k^{\mathrm{fr}}(Y_1) + dn .
\end{equation*}

\emph{Proof of} (ii). Now $\sim$ is face adjacency. If $d_k^{\subset}(Y_1) = \infty$ there is
nothing to prove. Otherwise let $T_0$ be an optimal tree graph: it meets every cell of $Y_1$,
every site of $T_0$ lies in $\bigcup Y_1$, and $|T_0| = M\, d_k^{\subset}(Y_1)$.

Suppose that $T_{j-1}$ is a tree graph meeting every cell of
$Y_1\cup\{\Delta_1,\dots,\Delta_{j-1}\}$, all of whose sites lie in the union of these cells.
Let $x'$ be a site of $\Delta'_j\cap T_{j-1}$. Write $\Delta'_j = \Delta_{a'}$. Since
$\Delta_j$ is face-adjacent to $\Delta'_j$, we have $\Delta_j = \Delta_{a'\pm e_\nu}$ for some
$\nu$ and some sign, and the two cells share the face across direction $\nu$.

From $x'$, move along $\pm e_\nu$ inside $\Delta'_j$ to the site $y'$ of $\Delta'_j$ that has
the same coordinates as $x'$ in the directions other than $\nu$ and the extreme $\nu$-coordinate
on the side of $\Delta_j$. This is $Ma'_\nu + M - 1$ for the sign $+$ and $Ma'_\nu$ for the
sign $-$. The move uses at most $M-1$ bonds, all inside the discrete box $\Delta'_j$. Then take
the one bond from $y'$ to $y := y' \pm e_\nu$, which lies in $\Delta_j$.

The added path has at most $M$ bonds and lies in
\begin{equation*}
\Delta'_j \cup \Delta_j \subset \bigcup\bigl(Y_1\cup\{\Delta_1,\dots,\Delta_j\}\bigr).
\end{equation*}
Put $T_j := T_{j-1}\cup(\text{that path})$. It is a tree graph in this union, connected through
$x'$ and meeting $\Delta_j$ at $y$, and $|T_j| \le |T_{j-1}| + M$. It therefore meets every cell
of $Y_1\cup\{\Delta_1,\dots,\Delta_j\}$, and the induction proceeds.

After $n$ steps, $T_n$ meets every cell of $Y$, all its sites lie in $\bigcup Y$, and
$|T_n| \le |T_0| + nM$. By the definition of the contained tree length,
\begin{equation*}
d_k^{\subset}(Y) \le M^{-1}|T_n| \le M^{-1}|T_0| + n = d_k^{\subset}(Y_1) + n .
\qedhere
\end{equation*}
\end{proof}

Consequently, the adjunction property \eqref{5.6:eq-linear-sizes-a} holds with $c' = d$ for the
free tree length $d_k^{\mathrm{fr}}$ under either adjacency relation. It also holds with
$c' = 1$ for the contained tree length $d_k^{\subset}$ under face adjacency. For the contained
tree length under closure-touching adjacency no claim is made: two cells that touch only at a
corner admit no lattice path inside their union.

\begin{proof}[Proof of part (c) of the weight lemma: transfer of the decay to the final domain]
\label{5.6:prf-decay-transfer}
We argue from three hypotheses: the holomorphy hypothesis $(\mathrm{H}_{\circ})$, \eqref{5.6:eq-weight-a}; the
counting forms $(\pi 1)\text{--}(\pi 3)$ of the pieces, \eqref{5.6:eq-pieces-b}; and the adjunction property
$(\mathrm{TL})_{c'}$ of the linear size $\ell$, \eqref{5.6:eq-linear-sizes-a}. We also use parts (a) and (b)
of the weight lemma (Lemma~\ref{5.6:lem-weight}), \eqref{5.6:eq-weight-b} and \eqref{5.6:eq-weight-c}.
Let $a \ge 0$, and assume $\mathfrak{b} \le \mathfrak{b}_0 e^{-a\ell(Y_1(t))}$.

Let $Y$ be a final domain with $\mathcal{P}(t;Y) \neq \emptyset$, and put $n = n(t;Y) = \#(Y \setminus Y_1(t))$.
By $(\pi 3)$, $Y \supseteq Y_1(t)$ and every $\sim$-component of $Y$ meets $Y_1(t)$. These are exactly the
hypotheses of $(\mathrm{TL})_{c'}$ with $Y_1 = Y_1(t)$, so
\begin{equation}\label{5.6:eq-decay-transfer-1}
\ell(Y) \le \ell(Y_1(t)) + c' n .
\end{equation}
Suppose first that $\ell(Y_1(t)) = \infty$. Then the assumption on $\mathfrak{b}$ forces $\mathfrak{b} = 0$, and
by (a) all pieces vanish, so there is nothing to prove. Suppose now that $\ell(Y_1(t)) < \infty$. Then
$\ell(Y) < \infty$ by \eqref{5.6:eq-decay-transfer-1}. Multiplying \eqref{5.6:eq-decay-transfer-1} by
$a \ge 0$ and exponentiating gives
\begin{equation*}
e^{-a\ell(Y_1(t))} \le e^{-a\ell(Y)}\, e^{a c' n}.
\end{equation*}

The bound (b), combined with $\mathfrak{b} \le \mathfrak{b}_0 e^{-a\ell(Y_1(t))}$ and then with the last
inequality, gives
\begin{align*}
\sum_{p \in \mathcal{P}(t;Y)} \sup |p|
 &\le \mathfrak{b}_0\, e^{-a\ell(Y_1(t))} \bigl(\lambda e^{-(\kappa_1 - 1)}\bigr)^{n}
 \le \mathfrak{b}_0\, e^{-a\ell(Y)} \bigl(\lambda e^{-(\kappa_1 - 1) + c' a}\bigr)^{n} \\
 &= \mathfrak{b}_0\, e^{-a\ell(Y)} \bigl(\lambda e^{-(\kappa_1 - 1 - c' a)}\bigr)^{n}.
\end{align*}
This is the first inequality of \eqref{5.6:eq-weight-d}. Taking logarithms, we have
$\lambda e^{-(\kappa_1 - 1 - c'a)} \le 1$ if and only if $\kappa_1 \ge 1 + \log\lambda + c'a$. This is the
one-sided floor $(\mathrm{F}\text{-}\kappa_1)^{P}$, \eqref{5.6:eq-weight-e}. Under it the $n$-th power is at
most $1$, which gives the second inequality of \eqref{5.6:eq-weight-d}. For $n = 0$ the power equals $1$,
and no floor is used.

We now read off the size of $c'a$. By the lemma on tree lengths under adjunction of cells
(Lemma~\ref{5.6:lem-tree-length-adjunction}), property $(\mathrm{TL})_{c'}$ holds in two cases:
\begin{itemize}
\item with $c' = d$ for the free tree length $d_k^{\mathrm{fr}}$;
\item with $c' = 1$ for the contained tree length $d_k^{\subset}$ under face adjacency.
\end{itemize}
The decay rate of Ba{\l}aban's inductive factor is $a = (1+3\beta)\kappa$, and $a \le 2\kappa$ for
$\beta \le \tfrac13$. Hence at $d = 4$ and $c' = d$ we have $c'a \le d \cdot 2\kappa = 8\kappa$, and at
$c' = 1$ we have $c'a \le 2\kappa$. The number of links of the chain is $\lambda = 3$. Thus
$(\mathrm{F}\text{-}\kappa_1)^{P}$ holds as soon as $\kappa_1 \ge 1 + \log 3 + 8\kappa$ in the first case,
and as soon as $\kappa_1 \ge 1 + \log 3 + 2\kappa$ in the second.

Finally, compare this with the one-sided condition $\kappa_1 \ge 16\kappa + 1$ on the auxiliary parameters.
We have
\begin{equation*}
16\kappa + 1 \ge 1 + \log 3 + 8\kappa
\iff 8\kappa \ge \log 3 .
\end{equation*}

The transfer proved here is in the plain linear size $\ell$. Ba{\l}aban's own extraction, from the exponential
of \cite[(1.66)]{B-CMP122b}, is made in a linear size relative to the large-field region, defined in
\cite[(1.67)]{B-CMP122b}; the argument above does not use it. Suppose the decay available where the bound is used
sits on a set $Y_0$ with $Y_0 \subsetneq Y_1(t)$. Then the comparison of $\ell(Y_1(t))$ with $\ell(Y_0)$ is
not supplied by part (c), and has to be made at the place of use.
\end{proof}

\begin{lemma}[Counting rooted connected sets of cells]\label{5.6:lem-rooted-sets}
Consider a graph whose vertices are cells, in which every cell has at most $D$ neighbours,
where $D \ge 1$. For every cell $\Delta_0$ and every integer $m \ge 1$ put
\begin{equation}\label{5.6:eq-rooted-sets-1}
a(m) := \#\{C : \Delta_0 \in C,\ \#C = m,\ C \text{ connected}\},
\end{equation}
where $C$ runs over sets of cells and connectedness refers to the graph induced on $C$. Then
\begin{equation*}
a(m) \le (4D)^{m-1}.
\end{equation*}
\end{lemma}

\begin{proof}
Fix a total order on the set of cells. It induces a total order on the set of neighbours of
each cell, which has at most $D$ elements. For a neighbour $\Delta'$ of a cell $\Delta$, the
\emph{rank} of $\Delta'$ at $\Delta$ is its position in the increasing enumeration of the
neighbours of $\Delta$. It is an element of $\{1,\dots,D\}$.

Let $C$ be a set counted in \eqref{5.6:eq-rooted-sets-1}. The graph induced on $C$ is finite
and connected, so it has a spanning tree. Fix one such tree $T_C$ for each $C$ and root it at
$\Delta_0$. It has $m$ vertices and $m-1$ edges, and each of its edges is an edge of the given
graph.

Consider the depth-first traversal of $T_C$. It starts at $\Delta_0$, visits the children of
each vertex in the fixed order, and returns to the parent when the subtree below the current
vertex is exhausted. This is a closed walk from $\Delta_0$ that crosses each of the $m-1$ tree
edges exactly twice: once away from the root, when the child is first entered, and once
towards the root, when the child's subtree is exhausted. It makes no other moves, so it has
exactly $2(m-1)$ steps. Every vertex of the rooted tree is reached, so the set of cells the
walk visits is $C$.

We encode the walk by a pair consisting of two pieces of data.
\begin{enumerate}
\item The first is the word in $\{\downarrow,\uparrow\}^{2(m-1)}$ that records, step by step,
whether the walk moves away from the root ($\downarrow$) or towards it ($\uparrow$). There are
at most $2^{2(m-1)} = 4^{m-1}$ such words.
\item The second records, for each of the $m-1$ steps marked $\downarrow$, in their order of
occurrence, the rank at the current cell of the cell moved to. This cell is a neighbour of the
current cell, because the step crosses an edge of $T_C$ and hence an edge of the graph. There
are at most $D^{m-1}$ such sequences.
\end{enumerate}

The code determines the walk, by induction on the number of steps. The walk begins at
$\Delta_0$. Suppose its first steps have been reconstructed. If the next letter is
$\downarrow$, the next cell is the neighbour of the current cell whose rank is the next entry
of the rank sequence. If the next letter is $\uparrow$, the walk returns to the parent of the
current cell in $T_C$. Since every cell other than $\Delta_0$ is first entered by a
$\downarrow$-step from its parent, this parent is the cell visited immediately before the
first visit to the current cell, and it is therefore determined by the part of the walk
already reconstructed.

Since the walk determines its set of visited cells, which is $C$, the code determines $C$.
Hence the map $C \mapsto$ code is injective. Counting the possible codes gives
\begin{equation*}
a(m) \le 4^{m-1} D^{m-1} = (4D)^{m-1}. \qedhere
\end{equation*}
\end{proof}

\begin{lemma}[Cells met by a tree graph]\label{5.6:lem-tree-cells}
Let $\ell \ge 0$ be a real number with $\ell M \in \mathbb{N}$, and let $T$ be a tree graph with
$|T| = \ell M$ bonds, in the sense of Definition~\ref{5.6:def-linear-sizes}. Then $T$ meets at most
$3^d(2\ell+1)$ cells. Consequently, for every finite cell set $Y_1$ with
$d_k^{\mathrm{fr}}(Y_1) < \infty$,
\begin{equation*}
\# Y_1 \le 3^d\bigl(2\,d_k^{\mathrm{fr}}(Y_1)+1\bigr) \le 3^d\bigl(2\,d_k^{\subset}(Y_1)+1\bigr).
\end{equation*}
\end{lemma}

In this lemma $\ell$ denotes a number, the length of $T$ in units of $M$, and not a linear size on
cell sets.

\begin{proof}
If $|T| = 0$, then $T$ consists of a single site, which lies in exactly one cell, so $T$ meets one cell,
and $1 \le 3^d(2\ell+1)$.

Let $|T| = \ell M \ge 1$. Fix a site of $T$ and traverse $T$ depth-first from it. Each bond of $T$ is
traversed once in each direction, so the traversal is a closed walk of $2\ell M$ steps along bonds of
$T$, and it passes through every site of $T$. Number its steps by the step-times
$q = 0, 1, \dots, 2\ell M$, and let $x(q)$ be the site occupied at step-time $q$. For
$j = 0, 1, \dots, \lfloor 2\ell \rfloor$ call $y_j := x(jM)$ the $j$-th post; each $jM$ lies in
$\{0,\dots,2\ell M\}$ because $j \le 2\ell$. There are $\lfloor 2\ell\rfloor + 1 \le 2\ell+1$
posts, and they occupy at most $2\ell+1$ sites.

Let $x$ be a site of $T$. It occurs on the walk at some step-time $q$ with $0 \le q \le 2\ell M$. Put
$j := \lfloor q/M \rfloor$. Then $jM \le q < (j+1)M$, and $j \le \lfloor 2\ell \rfloor$ because
$q/M \le 2\ell$. Hence $x$ is reached from the post $y_j$ in $q - jM \le M-1$ steps of the walk. Each
step moves along a bond and so changes exactly one coordinate by one. Therefore
\begin{equation*}
\|x - y_j\|_\infty \le \|x - y_j\|_1 \le M-1 .
\end{equation*}

We claim that for integers $x_\mu, y_\mu$ with $|x_\mu - y_\mu| \le M-1$,
\begin{equation}\label{5.6:eq-tree-cells-1}
\bigl|\lfloor x_\mu/M\rfloor - \lfloor y_\mu/M\rfloor\bigr| \le 1 .
\end{equation}
Suppose that $\lfloor x_\mu/M\rfloor \ge \lfloor y_\mu/M\rfloor + 2$. Since
$M\lfloor y_\mu/M\rfloor \ge y_\mu - (M-1)$, we obtain
\begin{equation*}
x_\mu \ge M\lfloor x_\mu/M\rfloor \ge M\lfloor y_\mu/M\rfloor + 2M
\ge y_\mu - (M-1) + 2M > y_\mu + M - 1,
\end{equation*}
which contradicts $|x_\mu - y_\mu| \le M-1$. The case
$\lfloor y_\mu/M\rfloor \ge \lfloor x_\mu/M\rfloor + 2$ is excluded in the same way with the roles of
$x_\mu$ and $y_\mu$ exchanged. This proves \eqref{5.6:eq-tree-cells-1}.

By the definition of the cells
$\Delta_a = \prod_{\mu=1}^{d}\bigl([Ma_\mu, Ma_\mu+M-1]\cap\mathbb{Z}\bigr)$, a site $x$ lies in $\Delta_a$
exactly when $a_\mu = \lfloor x_\mu/M\rfloor$ for every $\mu$. Thus
$a(x)_\mu = \lfloor x_\mu/M\rfloor$. Apply \eqref{5.6:eq-tree-cells-1} to $x$ and $y_j$ in each of
the $d$ coordinates. This gives $a(x) - a(y_j) \in \{-1,0,1\}^d$, so the cell of $x$ is one of the $3^d$
cells $\Delta_{a(y_j)+\epsilon}$ with $\epsilon \in \{-1,0,1\}^d$.

Every site of $T$ therefore lies in one of the cells $\Delta_{a(y_j)+\epsilon}$, where $j$ runs over the
at most $2\ell+1$ post indices and $\epsilon$ over $\{-1,0,1\}^d$. Hence $T$ meets at most
$3^d(2\ell+1)$ cells.

For the consequence, suppose first that $Y_1$ is nonempty. The quantity $d_k^{\mathrm{fr}}(Y_1)$ is
$M^{-1}$ times a minimum of numbers of bonds, taken over the tree graphs that meet every cell of $Y_1$.
That set of tree graphs is nonempty because $d_k^{\mathrm{fr}}(Y_1) < \infty$, so the minimum is
attained. Choose such an optimal free tree $T$. Then $|T| = \ell M$ with $\ell = d_k^{\mathrm{fr}}(Y_1)$
and $\ell M \in \mathbb{N}$. Every cell of $Y_1$ is met by $T$, so the first part gives
$\# Y_1 \le 3^d(2 d_k^{\mathrm{fr}}(Y_1)+1)$. If $Y_1$ is empty, this inequality holds trivially.

Finally, $d_k^{\mathrm{fr}}(Y_1) \le d_k^{\subset}(Y_1)$. Indeed, the tree graphs admitted in the
definition of $d_k^{\subset}(Y_1)$ are among those admitted in the definition of
$d_k^{\mathrm{fr}}(Y_1)$, and the inequality is trivial when $d_k^{\subset}(Y_1) = +\infty$. This gives
the second inequality.
\end{proof}

\begin{proof}[Proof of part (d) of Lemma~\ref{5.6:lem-weight}: sum over all final domains with prefactor]
\label{5.6:prf-all-final-domains}
Assume hypothesis $(\mathrm{H}_{\circ})$ of \eqref{5.6:eq-weight-a} and the counting forms
$(\pi 1)\text{--}(\pi 3)$ of \eqref{5.6:eq-pieces-b}. Write $Y_1:=Y_1(t)$, and recall that
$\lambda$ is the number of links of the chain and that
\begin{equation*}
N_t(n)=\#\{Y:\ \mathcal{P}(t;Y)\neq\emptyset,\ n(t;Y)=n\}.
\end{equation*}
The argument has three steps: a count of the final domains, a summation over them, and the
treatment of the prefactor.

\emph{Count.} Fix $n\ge 1$ and a finite set of cells $Y$ with $\mathcal{P}(t;Y)\neq\emptyset$ and
$n(t;Y)=n$. By $(\pi 3)$, $Y\supseteq Y_1$ and every $\sim$-component of $Y$ meets $Y_1$; that
is, every cell of $Y$ is joined to a cell of $Y_1$ by a $\sim$-path of cells of $Y$. Put
$A:=Y\setminus Y_1$, so that $\#A=n$, and let $C_1,\dots,C_q$ be the $\sim$-components of $A$,
that is, the components of the $\sim$-graph of $A$ itself.

Each $C_i$ is $\sim$-adjacent to $Y_1$. Indeed, take $\Delta\in C_i$ and a $\sim$-path of cells
of $Y$ from $\Delta$ to a cell of $Y_1$. The path starts in $C_i$ and ends in $Y_1$, which is
disjoint from $A\supseteq C_i$, so it has a first cell outside $C_i$. That cell is
$\sim$-adjacent to its predecessor on the path, which lies in $C_i$. It cannot lie in
$A\setminus C_i$, for then it would be joined to $C_i$ in the $\sim$-graph of $A$ and would
belong to $C_i$. Hence it lies in $Y_1$.

Fix once and for all a total order on the cells. Root $C_i$ at its least cell among those
$\sim$-adjacent to $Y_1$. Since $C_i\cap Y_1=\emptyset$, these are exactly the cells of
$C_i\cap\partial^{+}Y_1$, a set which is non-empty by the preceding paragraph. The components are
pairwise disjoint, so distinct components have distinct roots. Let $S\subset\partial^{+}Y_1$ be
the set of roots, so that $\#S=q$. For $s\in S$ let $m_s$ be the number of cells of the component
rooted at $s$. Assign to $Y$ the code
\begin{equation*}
\Bigl(S,\ (m_s)_{s\in S},\ \text{the components, as connected cell sets containing their roots}
\Bigr).
\end{equation*}
The map $Y\mapsto$ code is injective, because $Y=Y_1\cup\bigcup_i C_i$ and $Y_1$ depends on $t$
only.

We count the codes in three stages.
\begin{itemize}
\item The set $S$ is a non-empty subset of $\partial^{+}Y_1$. Each cell has $D_d$ neighbours,
so $\#\partial^{+}Y_1\le D_d\,\#Y_1$. This leaves at most
$2^{\#\partial^{+}Y_1}\le 2^{D_d\#Y_1}$ choices for $S$.
\item Given $S$ with $\#S=q\in[1,n]$, the sizes $(m_s)_{s\in S}$ satisfy $m_s\ge1$ and
$\sum_s m_s=n$. There are $\binom{n-1}{q-1}\le 2^{n-1}$ such choices.
\item Given $S$ and the sizes, the component rooted at $s$ is a $\sim$-connected set of $m_s$
cells containing $s$. Every cell has $D_d\ge1$ neighbours. The rooted-sets lemma
(Lemma~\ref{5.6:lem-rooted-sets}) therefore applies with $D=D_d$. It bounds the number of
choices for the components by
\begin{equation*}
\prod_{s\in S}a(m_s)\le\prod_{s\in S}(4D_d)^{m_s-1}\le(4D_d)^{n}.
\end{equation*}
\end{itemize}
Multiplying the three counts,
\begin{equation}\label{5.6:eq-all-final-domains-1}
N_t(n)\le 2^{D_d\#Y_1}\,2^{n-1}\,(4D_d)^n\le 2^{D_d\#Y_1}(8D_d)^n\qquad(n\ge1).
\end{equation}
Moreover $N_t(0)\le1$, since $n(t;Y)=0$ forces $A=\emptyset$, that is, $Y=Y_1$.

\emph{Sum.} By (b) of the weight lemma, \eqref{5.6:eq-weight-c}, for each $Y$ with
$n(t;Y)=n$,
\begin{equation*}
\sum_{p\in\mathcal{P}(t;Y)}\sup_{\mathcal{D}_b}|p|\le\mathfrak{b}\,(\lambda\,r(\kappa_1))^{n}.
\end{equation*}
Here $\sup_{\mathcal{D}_b}|p|$ is the supremum of $|p|$ over data of size at most $b$, as in (b).
All terms are non-negative, so the sum over all final domains may be regrouped by the value of
$n(t;Y)$. Using $N_t(0)\le1$, \eqref{5.6:eq-all-final-domains-1} and
$1\le 2^{D_d\#Y_1}$, we obtain
\begin{equation}\label{5.6:eq-all-final-domains-2}
\begin{split}
\sum_{Y}\sum_{p\in\mathcal{P}(t;Y)}\sup_{\mathcal{D}_b}|p|
&=\sum_{n\ge0}\ \sum_{Y:\,n(t;Y)=n}\ \sum_{p\in\mathcal{P}(t;Y)}\sup_{\mathcal{D}_b}|p|\\
&\le\mathfrak{b}\Bigl[1+\sum_{n\ge1}2^{D_d\#Y_1}(8D_d)^n(\lambda\,r(\kappa_1))^n\Bigr]\\
&\le 2^{D_d\#Y_1}\,\mathfrak{b}\sum_{n\ge0}\bigl(8\lambda D_d\,r(\kappa_1)\bigr)^n.
\end{split}
\end{equation}
By $(\mathrm{H}_{\circ})$ we have $\kappa_1\ge1$. Lemma~\ref{5.6:lem-cauchy-factor} then gives
$r(\kappa_1)\le e^{-(\kappa_1-1)}$. Consequently $8\lambda D_d\,r(\kappa_1)\le\tfrac12$ as soon
as $8\lambda D_d\,e^{-(\kappa_1-1)}\le\tfrac12$, and
\begin{equation*}
8\lambda D_d\,e^{-(\kappa_1-1)}\le\tfrac12
\quad\Longleftrightarrow\quad
\kappa_1\ge1+\log(16\lambda D_d).
\end{equation*}
Under this condition the geometric series in \eqref{5.6:eq-all-final-domains-2} is at most $2$.
The right side of \eqref{5.6:eq-all-final-domains-2} is then at most
$2\cdot2^{D_d\#Y_1}\,\mathfrak{b}$.

\emph{Prefactor.} Let $\ell\in\{d_k^{\mathrm{fr}},d_k^{\subset}\}$, and suppose
$d_k^{\mathrm{fr}}(Y_1)<\infty$, which is the hypothesis of the lemma on cells met by a tree
graph (Lemma~\ref{5.6:lem-tree-cells}). That lemma gives
$\#Y_1\le 3^d(2\ell(Y_1)+1)$ for both choices of $\ell$. Multiplying by $D_d\log 2$ and recalling
$c''_d=3^dD_d\log2$,
\begin{equation*}
D_d\,\#Y_1\log2\le D_d\,3^d\bigl(2\ell(Y_1)+1\bigr)\log2=c''_d\bigl(2\ell(Y_1)+1\bigr).
\end{equation*}
Hence $2^{D_d\#Y_1}\le e^{c''_d}\,e^{2c''_d\ell(Y_1)}$. This yields the absorption arithmetic
displayed in (d), \eqref{5.6:eq-weight-f}, with $a=(1+3\beta)\kappa$:
\begin{equation*}
e^{2c''_d\ell(Y_1)}\,e^{-a\ell(Y_1)}=e^{-(a-2c''_d)\ell(Y_1)}.
\end{equation*}
When $\ell(Y_1)>0$, the right side satisfies
\begin{equation*}
e^{-(a-2c''_d)\ell(Y_1)}\le e^{-(1-\varepsilon)a\ell(Y_1)}
\quad\Longleftrightarrow\quad
\varepsilon a\ge 2c''_d.
\end{equation*}
When $\ell(Y_1)=0$, both sides equal $1$. The condition $\varepsilon a\ge2c''_d$ is the one-sided
floor $(\mathrm{F}\text{-}\kappa)^{(d)}$ of \eqref{5.6:eq-weight-g}.
\end{proof}

\begin{remark}[No uniform bound without a decaying prefactor]\label{5.6:rem-no-uniform-bound}
Let $t$ be a term with exempt set $Y_1(t)$, variables $\mathcal{V}=\mathcal{V}(t)$ and composite
$F_t$ as in the weight lemma (Lemma~\ref{5.6:lem-weight}), with pieces as in
Definition~\ref{5.6:def-pieces}, and write $e^{\kappa_1}$ for the radius of the decoupling
polydisc and $r(\kappa_1)=e^{\kappa_1}(e^{\kappa_1}-1)^{-2}$.
\begin{enumerate}
\item A constant uniform in $t$ for the sum over all final domains cannot be derived from
$(\mathrm{H}_{\circ})$ and the counting forms $(\pi 1)$--$(\pi 3)$ of \eqref{5.6:eq-pieces-b}
by the majorants \eqref{5.6:eq-weight-b}; the same holds in the subset form
\eqref{5.6:eq-pieces-c}.
\item In the subset form no uniform constant exists: for every integer $J\ge 2$ and every
$\kappa_1\ge 1$ there is a polynomial $F$ in $J$ variables satisfying $(\mathrm{H}_{\circ})$
with $\mathfrak{N}(F)=\mathfrak{b}$ and
\begin{equation*}
\sum_{|\mathsf{y}|=2}|F_{\mathsf{y}}|\;\ge\;\tfrac12\sqrt{J}\,e^{-2\kappa_1}\mathfrak{b},
\end{equation*}
although each $|F_{\mathsf{y}}|$ obeys \eqref{5.6:eq-weight-b}.
\item In the chain form the sum over the final domains
$Y=Y_1(t)\cup\{\Delta,\Delta'\}$ is, for the same $F$, equally large with
$J=\#\partial^{+}Y_1(t)$, unless $\mathfrak{b}$ decays with the size of $Y_1(t)$.
\end{enumerate}
Consequently, at a fixed final domain the constant is $1$ (under $\kappa_1\ge 1+\log\lambda$),
while over all final domains a constant uniform in $t$ exists only with the decaying
prefactor displayed in \eqref{5.6:eq-weight-f}.
\end{remark}

\begin{proof}
The weight lemma, to which we appeal throughout, is proved as an implication from the
statements named in it.

(1) The majorant sum over all final domains obtained from \eqref{5.6:eq-weight-b} is
\begin{equation*}
\sum_{n\ge 0}N_t(n)\,(\lambda r(\kappa_1))^n\,\mathfrak{b}.
\end{equation*}
Every single-cell superset $Y=Y_1(t)\cup\{\Delta\}$ with $\Delta\in\partial^{+}Y_1(t)$ is
admissible for $(\pi 3)$, so $(\pi 1)$--$(\pi 3)$ do not exclude that each such $Y$ carries a
piece, and the majorants must allow
\begin{equation*}
N_t(1)\;\ge\;\#\partial^{+}Y_1(t).
\end{equation*}
Since $\lambda\ge 1$, the layer $n=1$ alone is at least
$\#\partial^{+}Y_1(t)\,r(\kappa_1)\,\mathfrak{b}$, which is unbounded in $\#Y_1(t)$ at
fixed $\mathfrak{b}$. In the subset form, with one variable per cell of $\mathcal{N}$, the
majorants $\mathfrak{N}(F)\,r(\kappa_1)^{|\mathsf{y}|}$ sum, by the binomial theorem, to
\begin{equation*}
\mathfrak{N}(F)\sum_{\emptyset\ne\mathsf{y}\subset\mathcal{N}}r(\kappa_1)^{|\mathsf{y}|}
=\mathfrak{N}(F)\bigl((1+r(\kappa_1))^{\#\mathcal{N}}-1\bigr),
\end{equation*}
which tends to infinity as $\#\mathcal{N}\to\infty$.

(2) The counterexample shows that in the subset form a uniform constant does not exist at
all, and not merely that none follows from the majorants. Let $J\ge 2$ be an integer and
$\mathcal{N}=\{1,\dots,J\}$. Put $\omega=e^{2\pi i/J}$ and
$\Phi=(\omega^{uv})_{u,v\in\mathcal{N}}$, and define
\begin{equation}\label{5.6:eq-no-uniform-bound-1}
F(\sigma)=\mathfrak{b}\,J^{-3/2}e^{-2\kappa_1}\sum_{u,v\in\mathcal{N}}\omega^{uv}\sigma_u\sigma_v .
\end{equation}
This is a polynomial and hence entire; it does not depend on the datum $\mathfrak{d}$. For
$u,u'\in\mathcal{N}$,
\begin{equation*}
\sum_{v\in\mathcal{N}}\omega^{uv}\,\overline{\omega^{u'v}}
=\sum_{v\in\mathcal{N}}\omega^{(u-u')v}=J\,\delta_{uu'},
\end{equation*}
because $\omega^{u-u'}\ne 1$ is a $J$-th root of unity when $u\ne u'$. Thus the rows of
$\Phi$ are orthogonal with squared norm $J$, $J^{-1/2}\Phi$ is unitary, and
$\|\Phi\sigma\|_2\le\sqrt{J}\,\|\sigma\|_2$. Writing
$\langle a,b\rangle=\sum_u\bar a_u b_u$, the Cauchy--Schwarz inequality gives, on
$\mathbb{D}_{\mathcal{N}}(e^{\kappa_1})$, where $\|\sigma\|_2^2\le Je^{2\kappa_1}$,
\begin{equation*}
\Bigl|\sum_{u,v}\omega^{uv}\sigma_u\sigma_v\Bigr|
=|\langle\bar\sigma,\Phi\sigma\rangle|
\le\|\sigma\|_2\,\sqrt{J}\,\|\sigma\|_2
\le\sqrt{J}\cdot J\,e^{2\kappa_1}.
\end{equation*}
Hence $|F|\le\mathfrak{b}$ on the closed, and so on the open, polydisc of radius
$e^{\kappa_1}$: $(\mathrm{H}_{\circ})$ holds with $\mathfrak{N}(F)=\mathfrak{b}$.

Let $\mathsf{y}=\{u,v\}$ with $u\ne v$. The monomial $\sigma_u\sigma_v$ occurs in
\eqref{5.6:eq-no-uniform-bound-1} with the coefficients of $(u,v)$ and of $(v,u)$, so
\begin{equation*}
\partial_u\partial_v F\equiv\mathfrak{b}J^{-3/2}e^{-2\kappa_1}(\omega^{uv}+\omega^{vu})
=2\mathfrak{b}J^{-3/2}e^{-2\kappa_1}\omega^{uv},
\end{equation*}
a constant. Integrating over $[0,1]^{\mathsf{y}}$ as in \eqref{5.6:eq-pieces-c} gives
$F_{\mathsf{y}}=2\mathfrak{b}J^{-3/2}e^{-2\kappa_1}\omega^{uv}$. There are $\tfrac12J(J-1)$
such $\mathsf{y}$, and $J-1\ge J/2$ for $J\ge 2$, so
\begin{equation*}
\sum_{|\mathsf{y}|=2}|F_{\mathsf{y}}|
=\tfrac12J(J-1)\cdot 2J^{-3/2}e^{-2\kappa_1}\mathfrak{b}
=(J-1)J^{-1/2}e^{-2\kappa_1}\mathfrak{b}
\ge\tfrac12\sqrt{J}\,e^{-2\kappa_1}\mathfrak{b}.
\end{equation*}
This is unbounded as $J\to\infty$ at fixed $\kappa_1$. Each single term is nevertheless
consistent with \eqref{5.6:eq-weight-b}. Indeed, $r(\kappa_1)\ge e^{-\kappa_1}$ and
$2J^{-3/2}\le 2^{-1/2}<1$, so
\begin{equation*}
|F_{\mathsf{y}}|=2J^{-3/2}e^{-2\kappa_1}\mathfrak{b}\le\mathfrak{b}\,r(\kappa_1)^2 .
\end{equation*}

(3) Let $J=\#\partial^{+}Y_1(t)$ and enumerate $\partial^{+}Y_1(t)=\{\Delta_1,\dots,\Delta_J\}$.
Read $F$ of \eqref{5.6:eq-no-uniform-bound-1} as the composite's dependence on the variables
of one link of the chain attached to the cells $\Delta_u$ only, the variable $\sigma_u$
belonging to $\Delta_u$. For $u\ne v$, a piece with variable set $\{u,v\}$ has final domain
$Y_1(t)\cup\{\Delta_u,\Delta_v\}$, and distinct pairs give distinct final domains, since
$\Delta_u\notin Y_1(t)$. Because $\partial_u\partial_vF$ is constant, such a piece equals
$2\mathfrak{b}J^{-3/2}e^{-2\kappa_1}\omega^{uv}$ whatever its probability measure. The
chain-form sum over these final domains is therefore at least
$\tfrac12\sqrt{J}\,e^{-2\kappa_1}\mathfrak{b}$, as in (2).

This quantity stays bounded uniformly in $t$ only if $\mathfrak{b}$ decays with the size of
$Y_1(t)$. That decay is exactly the prefactor paid in \eqref{5.6:eq-weight-f}, which rests on
\begin{equation*}
\#\partial^{+}Y_1\le D_d\,\#Y_1\le D_d\,3^d\bigl(2\ell(Y_1)+1\bigr).
\end{equation*}
Ba{\l}aban's construction, by contrast, sums at a fixed final domain: the domain is fixed
first \cite[text after (1.67)]{B-CMP122b}.

At a fixed final domain the weight lemma gives the constant $1$ under
$\kappa_1\ge 1+\log\lambda$. Over all final domains, (1)--(3) show that a constant uniform
in $t$ exists only with the decaying prefactor.
\end{proof}

We prove clauses (e) and (f) of the weight lemma (Lemma~\ref{5.6:lem-weight}). Hypothesis
$(\mathrm{H}_{\circ})$ of that lemma, displayed in \eqref{5.6:eq-weight-a}, is assumed throughout, in
the case $\mathcal{V}=\mathcal{N}$: one decoupling variable $\sigma(\Delta)$ is attached to each cell
$\Delta$ of the finite cell set $\mathcal{N}$. We write $F(\sigma;\mathfrak{d})$ for the composite, and
$F_{\mathsf{y}}$ for the subset-form piece \eqref{5.6:eq-pieces-c}; in particular
$F_{\emptyset}=F(0;\mathfrak{d})$. The quantity $\mathfrak{N}(F)$ is one of the following two numbers:
\begin{itemize}
\item the bound $\mathfrak{b}$ of $(\mathrm{H}_{\circ})$;
\item the supremum of $|F(\sigma;\mathfrak{d})-F(\mathbb{1};\mathfrak{d})|$ over
$\sigma\in P_{\mathcal{N}}(e^{\kappa_1})$ and $\mathfrak{d}\in\mathcal{D}_b$.
\end{itemize}
Here $\mathbb{1}$ and $0$ denote the points of $[0,1]^{\mathcal{N}}$ with all coordinates equal to $1$,
respectively to $0$.

\begin{proof}[Proof of clauses (e) and (f) of Lemma~\ref{5.6:lem-weight}]
\label{5.6:prf-telescoping-identity}
\emph{Proof of (e): the bound.} Let $\mathsf{y}\subset\mathcal{N}$ with $\mathsf{y}\neq\emptyset$.
By \eqref{5.6:eq-pieces-c}, $F_{\mathsf{y}}$ is a piece of the model form with variable set
$W=\mathsf{y}$ and base point $0$ on $\mathcal{N}\setminus\mathsf{y}$. Clause (a) of the weight lemma,
\eqref{5.6:eq-weight-b}, therefore applies to it with the bound $\mathfrak{b}$ of $(\mathrm{H}_{\circ})$,
which dominates $\sup|F|$ and equals $\mathfrak{N}(F)$ with the first choice of $\mathfrak{N}(F)$.

With the second choice of $\mathfrak{N}(F)$ we apply (a) to $F-F(\mathbb{1})$ instead. For fixed
$\mathfrak{d}$, the function $\sigma\mapsto F(\sigma;\mathfrak{d})-F(\mathbb{1};\mathfrak{d})$ is
holomorphic on $P_{\mathcal{N}}(e^{\kappa_1})$ and bounded there by $\mathfrak{N}(F)$. Since
$\mathsf{y}\neq\emptyset$, the operator $\partial^{\mathsf{y}}$ annihilates the constant
$F(\mathbb{1};\mathfrak{d})$. Hence $F-F(\mathbb{1})$ has the same piece $F_{\mathsf{y}}$.

In both cases, (a) gives the first inequality below. Lemma~\ref{5.6:lem-cauchy-factor} gives
$r(\kappa_1)\le e^{-(\kappa_1-1)}$ for $\kappa_1\ge 1$, and raising this to the power $|\mathsf{y}|$
gives the second:
\begin{equation*}
\sup_{\mathfrak{d}\in\mathcal{D}_b}|F_{\mathsf{y}}(\mathfrak{d})|
\le\mathfrak{N}(F)\,r(\kappa_1)^{|\mathsf{y}|}
\le\mathfrak{N}(F)\,e^{-(\kappa_1-1)|\mathsf{y}|}.
\end{equation*}
This is the bound in \eqref{5.6:eq-weight-h}.

\emph{Proof of (e): the telescoping identity.} Fix $\mathfrak{d}\in\mathcal{D}_b$ and write
$F(\sigma)$ for $F(\sigma;\mathfrak{d})$. By $(\mathrm{H}_{\circ})$, $F$ is continuous, separately
holomorphic and bounded on $P_{\mathcal{N}}(e^{\kappa_1})$, and $e^{\kappa_1}>1$ because
$\kappa_1\ge 1$. The Cauchy estimates (Lemma~\ref{5.6:lem-cauchy-estimates}) therefore show that $F$
is $C^\infty$ on $P_{\mathcal{N}}(e^{\kappa_1})$. Hence its restriction to the real points, an open
neighbourhood of $[0,1]^{\mathcal{N}}$ in $\mathbb{R}^{\mathcal{N}}$, is $C^\infty$ in the real
variables.

Let $\mathcal{C}=C^\infty([0,1]^{\mathcal{N}})$ be the space of functions $H$ on
$[0,1]^{\mathcal{N}}$ with continuous partial derivatives of all orders on the closed cube. By the
preceding paragraph, the restriction of $F$ to $[0,1]^{\mathcal{N}}$ lies in $\mathcal{C}$.
For $\Delta\in\mathcal{N}$,
$c\in\{0,1\}$ and $H\in\mathcal{C}$ set
\begin{gather*}
(E_\Delta^{c}H)(\sigma):=H(\sigma\text{ with }\sigma(\Delta):=c),\\
(I_\Delta H)(\sigma):=\int_0^1(\partial_{\sigma(\Delta)}H)(\sigma\text{ with }\sigma(\Delta):=u)\,du .
\end{gather*}
Both operators map $\mathcal{C}$ into itself; for $I_\Delta$ this is differentiation under the integral
sign, the integrand having continuous derivatives on the compact cube. The fundamental theorem of
calculus in the variable $\sigma(\Delta)$ gives
\begin{equation}\label{5.6:eq-telescoping-identity-1}
E_\Delta^{1}H=E_\Delta^{0}H+I_\Delta H\qquad(H\in\mathcal{C}).
\end{equation}

Operators attached to distinct cells $\Delta\neq\Delta'$ commute on $\mathcal{C}$.
\begin{itemize}
\item Two evaluations fix two different coordinates, so they commute.
\item $E_\Delta^{c}$ and $I_{\Delta'}$ commute, because fixing $\sigma(\Delta)$ commutes with
differentiating in $\sigma(\Delta')$ and integrating over $\sigma(\Delta')$.
\item For $I_\Delta$ and $I_{\Delta'}$, differentiation under the integral sign gives
$I_\Delta I_{\Delta'}H$ as the integral over $[0,1]^2$ of
$\partial_{\sigma(\Delta)}\partial_{\sigma(\Delta')}H$, evaluated at
$(\sigma(\Delta),\sigma(\Delta'))=(u,u')$. The same computation gives $I_{\Delta'}I_\Delta H$ as the
integral of $\partial_{\sigma(\Delta')}\partial_{\sigma(\Delta)}H$ taken in the other order. These
agree by Clairaut's theorem, since $H$ is $C^2$, and by Fubini's theorem on $[0,1]^2$.
\end{itemize}

Let $\prod$ denote composition of operators in any order; by the commutation just shown, the order is
immaterial. Apply \eqref{5.6:eq-telescoping-identity-1} once for each cell. Distributing the
composition of the $|\mathcal{N}|$ sums, and reordering each resulting term by commutativity, we
obtain, at every $\sigma$,
\begin{align*}
F(\mathbb{1})
&=\Bigl(\prod_{\Delta\in\mathcal{N}}E_\Delta^{1}\Bigr)F
=\Bigl(\prod_{\Delta\in\mathcal{N}}\bigl(E_\Delta^{0}+I_\Delta\bigr)\Bigr)F\\
&=\sum_{\mathsf{y}\subset\mathcal{N}}
\Bigl(\prod_{\Delta\in\mathsf{y}}I_\Delta\prod_{\Delta\notin\mathsf{y}}E_\Delta^{0}\Bigr)F .
\end{align*}

To evaluate the $\mathsf{y}$-summand, first apply the $E_\Delta^{0}$ with $\Delta\notin\mathsf{y}$;
they set the coordinates off $\mathsf{y}$ equal to $0$. Then apply the $I_\Delta$ with
$\Delta\in\mathsf{y}$ one after another, and use Fubini's theorem on $[0,1]^{\mathsf{y}}$. Writing
$u_{\mathsf{y}}\in[0,1]^{\mathsf{y}}$ for the integration variable, the summand equals
\begin{equation*}
\int_{[0,1]^{\mathsf{y}}}\partial^{\mathsf{y}}F(u_{\mathsf{y}},0_{\mathcal{N}\setminus\mathsf{y}})\,
du_{\mathsf{y}}
=F_{\mathsf{y}}
\end{equation*}
by \eqref{5.6:eq-pieces-c}. For $\mathsf{y}=\emptyset$ it is $F(0)=F_\emptyset$. Subtracting $F(0)$
from both sides gives
\begin{equation*}
F(\mathbb{1};\mathfrak{d})-F(0;\mathfrak{d})
=\sum_{\mathsf{y}\subset\mathcal{N},\ \mathsf{y}\neq\emptyset}F_{\mathsf{y}}(\mathfrak{d}),
\end{equation*}
which is the telescoping identity of \eqref{5.6:eq-weight-h}.

\emph{Proof of (f).} Clause (f), \eqref{5.6:eq-weight-i}, asserts that the constants and conditions in
(a)--(e) depend on no parameter beyond those listed below. Consider the constants and conditions in
clauses (a)--(e), and in the lemmas on which these clauses rest:
Lemma~\ref{5.6:lem-cauchy-factor}, Lemma~\ref{5.6:lem-cauchy-estimates},
Lemma~\ref{5.6:lem-tree-length-adjunction}, Lemma~\ref{5.6:lem-rooted-sets} and
Lemma~\ref{5.6:lem-tree-cells}. They involve the following parameters and no others:
\begin{itemize}
\item $\kappa_1$, through the factor $r(\kappa_1)$ and through the one-sided floors imposed on
$\kappa_1$;
\item the number $\lambda$ of links of the chain;
\item the dimension $d$ and the adjacency $\sim$, through the constants $c'$, $D_d$ and $3^d$;
\item in clause (c), and in the optional absorption of clause (d) only, the decay rate
$a=(1+3\beta)\kappa$ of the statement to which the weight lemma is applied (hence $\kappa$ and
$\beta$), together with $\varepsilon$.
\end{itemize}
None of these is $K$, $k$, $N$, $L$, the volume, a position, the number of terms or the coupling.
The cube side $M$ enters only through the definitions of the tree lengths $d_k^{\mathrm{fr}}$ and
$d_k^{\subset}$, which are measured in units of $M$; there it cancels, so that no constant or
condition depends on $M$. This is the assertion of (f).
\end{proof}

\begin{proposition}[Holomorphy of localised terms on the decoupling polydisc]
\label{5.6:prop-polydisc-holomorphy}
Let $\kappa_1 \ge 1$, let $b$ be fixed, and let $\mathcal{D}_b$ be the set of data of size at most $b$.
Assume the floor $\delta_1 M \ge \kappa_1$, where $\delta_1$ is the decay constant of
\cite{B-CMP116} and $M$ is chosen after $\kappa_1$; and assume that $\gamma$ does not exceed a
ceiling which depends on $\kappa_1$, among other parameters, and is fixed after $\kappa_1$. Both
conditions are one-sided.
Let $t$ be a term of the expansion, with decoupling variables $\mathcal{V}=\mathcal{V}(t)$, exempt set
$Y_1(t)$ and composite
\begin{equation*}
F_t(\sigma;\mathfrak{d}) = t\circ\mathsf{R}(\sigma;\mathfrak{d}),
\end{equation*}
the value of $t$ at the $\sigma$-interpolated localised configuration with datum $\mathfrak{d}$, as in
Definition~\ref{5.6:def-decoupling-variables}. Then the hypothesis $(\mathrm{H}_{\circ})$ of the weight
lemma (Lemma~\ref{5.6:lem-weight}, display \eqref{5.6:eq-weight-a}) holds for $t$. That is,
for every $\mathfrak{d} \in \mathcal{D}_b$ the function
$\sigma \mapsto F_t(\sigma;\mathfrak{d})$ is holomorphic on the open polydisc
$P_{\mathcal{V}}(e^{\kappa_1})$. Here holomorphic means continuous and separately holomorphic in each
$\sigma_v$, $v \in \mathcal{V}$. Moreover, with $\mathfrak{b}$ the supremum of $|t|$ on its domain of
analyticity,
\begin{equation*}
|F_t(\sigma;\mathfrak{d})| \le \mathfrak{b}
\qquad \text{for all } \sigma \in P_{\mathcal{V}}(e^{\kappa_1}),\ \mathfrak{d} \in \mathcal{D}_b.
\end{equation*}
Holomorphy in $\sigma$ is asserted separately for each $\mathfrak{d}$. Nothing is asserted about the
dependence on $\mathfrak{d}$: no measurability, continuity or analyticity. As in
Lemma~\ref{5.6:lem-weight}, the closed-polydisc form of the assertion, holomorphy of
$F_t(\cdot;\mathfrak{d})$ on a neighbourhood of $\mathbb{D}_{\mathcal{V}}(e^{\kappa_1})$ with
$|F_t| \le \mathfrak{b}$ on $\mathbb{D}_{\mathcal{V}}(e^{\kappa_1})$, implies it by restriction; and,
since replacing $F_t$ by $F_t - c(\mathfrak{d})$, for a constant $c(\mathfrak{d})$ depending on
$\mathfrak{d}$ alone, changes no piece $p$ with $W(p) \neq \emptyset$ in Lemma~\ref{5.6:lem-weight},
the assertion may equally be read for $F_t - c(\mathfrak{d})$, with $\mathfrak{b}$ replaced by the
supremum of $|F_t(\sigma;\mathfrak{d}) - c(\mathfrak{d})|$ over
$\sigma \in P_{\mathcal{V}}(e^{\kappa_1})$ and $\mathfrak{d} \in \mathcal{D}_b$.
\end{proposition}

Proof outstanding.

\emph{Route.} A proof would impose two one-sided conditions on the parameters and show that under
them the interpolated localised configuration $\mathsf{R}(\sigma;\mathfrak{d})$ lies inside the
analyticity domain of $t$ for every
$\sigma \in P_{\mathcal{V}}(e^{\kappa_1})$. The bound then holds with $\mathfrak{b}$ the supremum of
$|t|$ on that domain.

The first condition is the floor $\delta_1 M \ge \kappa_1$ of \cite{B-CMP116}. Here $\delta_1$ is the
decay constant of that paper, and $M$ is chosen after $\kappa_1$. Under this floor, on the circles
$|\sigma(\Delta)| = e^{\kappa_1}$ the interpolated operators of the random-walk and localisation
expansions keep their exponential decay. There they are bounded by $e^{16\kappa_1}$
\cite[(1.7), (1.11)]{B-CMP116}.

The configurations of \cite[(1.61)--(1.63)]{B-CMP122b} are built from these localised operators acting
on the field arguments of the large-field operation $\mathbf{R}$. On the circles, the images of those
arguments are therefore bounded only by a factor of order $e^{16\kappa_1}$, times a constant of
Ba{\l}aban's construction, times the size of the arguments. Containment of
$\mathsf{R}(\sigma;\mathfrak{d})$ in the analyticity domain of $t$
then follows once $e^{16\kappa_1}$ times the size of these field arguments is at most the radius of
that domain.

The second condition secures this. In \cite{B-CMP116}, at the step of the large-field operation
$\mathbf{T}$, it is expressed
through the radius $\mathsf{w}_T$ of the circle used at that step in
\cite[(1.17)--(1.22)]{B-CMP116}:
\begin{equation*}
\mathsf{w}_T = \frac{\alpha_2}{4 B_0 C_1 e^{16\kappa_1} \delta_k}.
\end{equation*}
Here $\alpha_2$ and $B_0$ are constants of \cite{B-CMP116} entering the radius of
\cite[(1.22)]{B-CMP116}. The field arguments are
small in the coupling, and
\begin{equation*}
\delta_k = \frac{A_1}{A_0}\,\varepsilon_k,
\qquad
\varepsilon_k = g_k\, p_0(g_k),
\end{equation*}
with $p_0(\cdot)$ the degree function. This $\delta_k$, a smallness parameter at level $k$, is
unrelated to the decay constant $\delta_1$ of \cite{B-CMP116} above. The requirement is therefore a
smallness condition on the
coupling: an upper bound on $\gamma$ that depends on $\kappa_1$, among other parameters, and is fixed
after $\kappa_1$. It is not a lower bound on $g$, and it places no upper bound on $L$, $M$, $K$, $k$,
the number of steps since a large-field region was created, or $\kappa_1$.

A proof must derive from \cite[(1.7), (1.11), (1.17)--(1.22)]{B-CMP116} and
\cite[(1.61)--(1.63)]{B-CMP122b} three things: first, that under $\delta_1 M \ge \kappa_1$, with
$\delta_1$ as in \cite{B-CMP116}, the
interpolated operators keep their exponential decay on $|\sigma(\Delta)| = e^{\kappa_1}$ and are
bounded there by $e^{16\kappa_1}$; second, the resulting bound on the images of the field arguments;
third, that under the upper bound on $\gamma$ these images lie in the analyticity domain of $t$.

\begin{definition}[Holomorphy in finitely many complex variables: the facts used]
\label{5.6:not-holomorphy-facts}
The facts collected here, all classical, are the only mathematics taken as known in the
verification of the holomorphy hypothesis of the weight lemma. No Banach-space theory and no
theorem of Hartogs is used there.
\begin{enumerate}
\item[(a)] \emph{Vector spaces.} For $N\ge 2$, $\mathfrak{g}^{\mathbb{C}}=\mathfrak{sl}(N,\mathbb{C})$
is a complex vector space of dimension $N^2-1$. A $\mathfrak{g}^{\mathbb{C}}$-valued function
on a finite set of bonds (or sites, or plaquettes) is a point of
$\mathfrak{F}:=(\mathfrak{g}^{\mathbb{C}})^{\#\mathrm{bonds}}\cong\mathbb{C}^{n}$, where
$n=(N^2-1)\,\#\mathrm{bonds}$. All norms on $\mathbb{C}^{n}$ are equivalent and continuous.
A seminorm $\mathcal{N}$ on $\mathbb{C}^{n}$ satisfies $\mathcal{N}(cA)=|c|\,\mathcal{N}(A)$
for $c\in\mathbb{C}$ and $\mathcal{N}(A_1+A_2)\le\mathcal{N}(A_1)+\mathcal{N}(A_2)$, and is
continuous on $\mathbb{C}^{n}$. Only this vector-space structure of
$\mathfrak{g}^{\mathbb{C}}$ enters that verification: no Lie bracket, group law, adjoint
action, Haar measure or character.
\item[(b)] \emph{Holomorphy in finitely many variables (Osgood).} Let $U\subset\mathbb{C}^{n}$
be open and $G\colon U\to\mathbb{C}^{m}$. The following are equivalent, and $G$ is called
\emph{holomorphic} if it has any of these properties:
\begin{itemize}
\item[(h1)] $G$ is continuous and separately holomorphic, that is, for each coordinate, with
the other coordinates held fixed, $G$ is holomorphic in the one-variable sense; this is the
sense in which the word is used in the weight lemma;
\item[(h2)] near each point of $U$, $G$ is the sum of a power series converging uniformly on
a small closed polydisc (power series in several variables being summed as absolutely
convergent multiple series);
\item[(h3)] $G$ is $C^{\infty}$ and $\mathbb{C}$-differentiable:
$G(z+v)=G(z)+DG(z)v+o(|v|)$ with $DG(z)$ $\mathbb{C}$-linear.
\end{itemize}
Continuity is part of (h1); Hartogs' theorem, which would remove that hypothesis, is not used.
\item[(c)] \emph{Closure facts.} Affine maps and polynomials are holomorphic; a map into
$\mathbb{C}^{m}$ is holomorphic if and only if its components are; a continuous
$\mathbb{C}$-linear map composed with a holomorphic map is holomorphic; the composite of
holomorphic maps between open sets is holomorphic, with
\begin{equation}\label{5.6:eq-holomorphy-facts-1}
D(\mathcal{T}\circ A)(z)=D\mathcal{T}(A(z))\circ DA(z);
\end{equation}
the restriction of a holomorphic map to an affine complex line or to a coordinate slice is
holomorphic; a holomorphic map is continuous. \emph{One-variable chain rule:} if
$\rho>0$, $\gamma\colon\{|w|<\rho\}\to U$ is a complex-differentiable curve (that is,
$\gamma'(w)$ exists in $\mathfrak{F}$, componentwise) and $\mathcal{T}\colon U\to\mathbb{C}$
satisfies (h3), then $w\mapsto\mathcal{T}(\gamma(w))$ is complex-differentiable with
derivative $D\mathcal{T}(\gamma(w))[\gamma'(w)]$.
\item[(d)] \emph{One variable.} Cauchy's integral formula holds for $\varphi$ continuous on
$\{|\tau|\le R'\}$ and holomorphic in $\{|\tau|<R'\}$:
\begin{equation}\label{5.6:eq-holomorphy-facts-2}
\varphi(\tau)=\frac{1}{2\pi i}\oint_{|\zeta|=R'}\frac{\varphi(\zeta)}{\zeta-\tau}\,d\zeta,
\qquad |\tau|<R'.
\end{equation}
For $\varphi$ holomorphic on $\{|\tau|<\rho\}$ and $0<R'<\rho$,
\begin{equation}\label{5.6:eq-holomorphy-facts-3}
\varphi'(0)=\frac{1}{2\pi i}\oint_{|\tau|=R'}\varphi(\tau)\,\tau^{-2}\,d\tau,
\qquad
|\varphi'(0)|\le\frac{1}{R'}\max_{|\tau|=R'}|\varphi|.
\end{equation}
\item[(e)] \emph{Parameter integrals.} Let $\mathsf{K}$ be a circle with arclength measure,
$U\subset\mathbb{C}^{n}$ open, and $h\colon U\times\mathsf{K}\to\mathbb{C}$ continuous with
$h(\cdot,\xi)$ holomorphic on $U$ for each $\xi\in\mathsf{K}$. Then
$\mathsf{H}(z):=\int_{\mathsf{K}}h(z,\xi)\,d\xi$ is continuous and separately holomorphic on $U$,
hence holomorphic.
\end{enumerate}
\end{definition}

\begin{proof}
\emph{Item (a).} The dimension count, the identification of $\mathfrak{F}$ with
$\mathbb{C}^{n}$ and the equivalence and continuity of norms are standard finite-dimensional
linear algebra. For the continuity of a seminorm, let $e_1,\dots,e_n$ be the standard basis
and $A_1=\sum_i a_ie_i$, $A_2=\sum_i b_ie_i$. The two defining properties give
$\mathcal{N}(A_1)\le\mathcal{N}(A_2)+\mathcal{N}(A_1-A_2)$ and the same with $A_1,A_2$
exchanged, so
\begin{equation*}
|\mathcal{N}(A_1)-\mathcal{N}(A_2)|\le\mathcal{N}(A_1-A_2)
\le\sum_{i=1}^{n}|a_i-b_i|\,\mathcal{N}(e_i),
\end{equation*}
which tends to zero as $A_2\to A_1$.

\emph{Item (b), (h1)$\Rightarrow$(h2).} It suffices to treat each component of $G$. Fix
$a\in U$ and radii $r_w>0$ with
$\mathbb{D}:=\{\zeta:|\zeta_w-a_w|\le r_w\ \forall w\}\subset U$. For $z$ in the open polydisc
and with $z_2,\dots,z_n$ fixed, the function $\zeta_1\mapsto G(\zeta_1,z_2,\dots,z_n)$ is
holomorphic on a neighbourhood of the closed disc $|\zeta_1-a_1|\le r_1$, because the closed
polydisc lies in the open set $U$; apply \eqref{5.6:eq-holomorphy-facts-2} to it, then,
for each fixed $\zeta_1$ on the circle $|\zeta_1-a_1|=r_1$, apply it in the second coordinate,
and so on. Since $G$ is continuous on the compact product of circles
$\Gamma_r:=\{\zeta:|\zeta_w-a_w|=r_w\ \text{for all } w\}$, the iterated integral is an
integral over $\Gamma_r$, and
\begin{equation}\label{5.6:eq-holomorphy-facts-4}
G(z)=\frac{1}{(2\pi i)^{n}}\oint_{\Gamma_r}G(\zeta)\prod_{w=1}^{n}(\zeta_w-z_w)^{-1}\,d\zeta ;
\end{equation}
this is the representation used in the proof of the weight lemma.
Fix $0<r'_w<r_w$. For $|z_w-a_w|\le r'_w$ and $\zeta\in\Gamma_r$ expand in geometric
series:
\begin{equation*}
(\zeta_w-z_w)^{-1}=\sum_{j\ge 0}\frac{(z_w-a_w)^{j}}{(\zeta_w-a_w)^{j+1}},
\end{equation*}
with terms bounded by $r_w^{-1}(r'_w/r_w)^{j}$, so the product over $w$ converges
absolutely and uniformly in $\zeta\in\Gamma_r$ and in $z$ in the closed polydisc of radii
$r'_w$. As $G$
is bounded on $\Gamma_r$, summation and integration in \eqref{5.6:eq-holomorphy-facts-4}
may be exchanged. For multi-indices $\alpha=(\alpha_w)\in\{0,1,2,\dots\}^{n}$ write
$(z-a)^{\alpha}:=\prod_w(z_w-a_w)^{\alpha_w}$, $r^{\alpha}:=\prod_w r_w^{\alpha_w}$ and
$\partial^{\beta}:=\prod_w\partial_{z_w}^{\beta_w}$; after applying $\partial^{\beta}$, the
terms with some $\alpha_w<\beta_w$ are absent. Hence on the closed polydisc of radii $r'_w$
the function $G$ is the sum of the absolutely and uniformly convergent power series
$\sum_{\alpha}c_\alpha(z-a)^{\alpha}$, with
\begin{equation*}
c_\alpha=\frac{1}{(2\pi i)^{n}}\oint_{\Gamma_r}G(\zeta)\prod_w(\zeta_w-a_w)^{-\alpha_w-1}\,d\zeta .
\end{equation*}

\emph{(h2)$\Rightarrow$(h3).} Let $G(z)=\sum_\alpha c_\alpha(z-a)^{\alpha}$ converge uniformly on
the closed polydisc of radii $r_w$. At $z_w=a_w+r_w$ the series converges absolutely, so its
terms are bounded: $|c_\alpha|r^{\alpha}\le C$. For $r'_w<r_w$ and every multi-index $\beta$,
the series of $\partial^{\beta}$-derivatives of the terms is then dominated on the polydisc of
radii $r'_w$ by $C\sum_\alpha\alpha^{\beta}r^{-\beta}(r'/r)^{\alpha-\beta}<\infty$ (with
$\alpha^{\beta}=\prod_w\alpha_w^{\beta_w}$), hence converges uniformly there. Differentiating
termwise, $G$ is $C^{\infty}$ near $a$. Each term is a polynomial in $z$, whose real
differential is $\mathbb{C}$-linear; the real differential of $G$ is the limit of these,
hence $\mathbb{C}$-linear, and this is the expansion of (h3).

\emph{(h3)$\Rightarrow$(h1).} A $C^{\infty}$ map is continuous. Taking $v=\eta e_w$ with
$\eta\in\mathbb{C}$ in (h3) gives $G(z+\eta e_w)=G(z)+\eta\,DG(z)e_w+o(|\eta|)$, so $G$ is
complex-differentiable in $z_w$ with the other coordinates fixed.

\emph{Item (c).} Affine maps and polynomials are continuous and holomorphic in each variable
separately, hence holomorphic by (h1). Property (h1) holds for a map into $\mathbb{C}^{m}$ if
and only if it holds for each component. For the composite, let $A\colon U\to V$ and
$\mathcal{T}\colon V\to\mathbb{C}^{m_1}$ be holomorphic, $V\subset\mathbb{C}^{m}$ open, in the
form (h3). Then
$\mathcal{T}\circ A$ is $C^{\infty}$ by the real chain rule, and inserting
$A(z+v)=A(z)+DA(z)v+o(|v|)$ into the expansion of $\mathcal{T}$ at $A(z)$ gives
\eqref{5.6:eq-holomorphy-facts-1}, a composite of $\mathbb{C}$-linear maps. A continuous
$\mathbb{C}$-linear map is holomorphic, so its composite with a holomorphic map is a special
case; so is the restriction to an affine complex line or coordinate slice, the composite with
an affine parametrisation. Continuity is part of (h1). For the one-variable chain rule, put
$\Delta(\eta):=\gamma(w+\eta)-\gamma(w)$; then $|\Delta(\eta)|=O(|\eta|)$ and, by (h3),
\begin{equation*}
\frac{\mathcal{T}(\gamma(w+\eta))-\mathcal{T}(\gamma(w))}{\eta}
=D\mathcal{T}(\gamma(w))\Bigl[\frac{\Delta(\eta)}{\eta}\Bigr]+\frac{o(|\Delta(\eta)|)}{\eta},
\end{equation*}
which tends to $D\mathcal{T}(\gamma(w))[\gamma'(w)]$ as $\eta\to 0$, by continuity of the
linear map $D\mathcal{T}(\gamma(w))$.

\emph{Item (d).} Formula \eqref{5.6:eq-holomorphy-facts-2} and the first formula in
\eqref{5.6:eq-holomorphy-facts-3} are Cauchy's integral formulas for $\varphi$ and for
$\varphi'$ in one complex variable; the estimate in \eqref{5.6:eq-holomorphy-facts-3}
follows from the latter by bounding the integrand by
$\max_{|\tau|=R'}|\varphi|/R'^{2}$ on a circle of length $2\pi R'$.

\emph{Item (e).} Let $z_0\in U$ and let $\bar B\subset U$ be a closed ball about $z_0$. Since
$\bar B\times\mathsf{K}$ is compact, $h$ is uniformly continuous there. Hence
$|\mathsf{H}(z)-\mathsf{H}(z_0)|$ is at most the length of $\mathsf{K}$ times
$\sup_{\xi\in\mathsf{K}}|h(z,\xi)-h(z_0,\xi)|$, which tends to zero as $z\to z_0$.
For separate holomorphy fix a coordinate, the others at fixed values, and an open disc
$O$ in that coordinate such that the points with that coordinate in the closure of $O$ and
the others at their fixed values lie in $U$. For every closed triangle $\triangle\subset O$,
Fubini's theorem for the continuous integrand on $\partial\triangle\times\mathsf{K}$ gives
\begin{equation*}
\oint_{\partial\triangle}\int_{\mathsf{K}}h\,d\xi\,d\zeta
=\int_{\mathsf{K}}\oint_{\partial\triangle}h\,d\zeta\,d\xi=0,
\end{equation*}
since each $h(\cdot,\xi)$ restricted to the coordinate disc is holomorphic by (c), so Cauchy's
theorem for the triangle applies. By Morera's theorem $\mathsf{H}$ is holomorphic in that
coordinate on $O$. Thus $\mathsf{H}$ satisfies (h1) and is holomorphic by (b).
\end{proof}

\begin{definition}[Analyticity on a closed domain]\label{5.6:def-closed-analyticity}
The following standing setting and reading of analyticity are fixed for the rest of this chapter.

\begin{enumerate}
\item[(a)] \emph{Finitely many variables.} At each renormalisation step the lattice is finite: it is
a finite torus or box. The field configurations denoted $\mathbf{U}$, $\mathbf{J}$, $A$, $A'$, $B$,
$B'$, \dots{} in \cite[\S1]{B-CMP116} are therefore functions on finitely many bonds,
sites or plaquettes; wherever that section is followed below, these letters keep the meaning
they have there. Those with values in
$\mathfrak{g}^{\mathbb{C}}=\mathfrak{sl}(N,\mathbb{C})$ are points of the finite-dimensional complex
vector space $\mathfrak{F}\cong\mathbb{C}^n$ of Definition~\ref{5.6:not-holomorphy-facts}.
Accordingly, the phrases ``an analytic function of $A$'', ``of $s(Y_0)$'', ``of $B'$'' in
\cite[\S1]{B-CMP116} are read as holomorphy in finitely many complex variables, in the sense of
Definition~\ref{5.6:not-holomorphy-facts} (on the closed domains of \cite[\S1]{B-CMP116} this is
made precise in part (b)). The admissible norms, the norms of
\cite[Theorems~3.1--3.10]{B-CMP99b} and the quantities of \cite[(3.31)]{B-CMP109} are seminorms or
norms on $\mathfrak{F}$.

\item[(b)] \emph{Analyticity on a set with a closed polydisc factor.} Let $Y_0$ be a connected
localisation domain
containing $\tilde{\square}^{4}$. Let $\mathcal{V}$ be the set of $M$-cubes
$\Delta\subset Y_0\setminus\tilde{\square}^{4}$, and write
$s(Y_0)=(s(\Delta))_{\Delta\in\mathcal{V}}$ for the family of complex parameters attached to them.
Put
\begin{equation*}
\mathbb{D}:=\bigl\{s(Y_0)\in\mathbb{C}^{\mathcal{V}}:\ |s(\Delta)|\le e^{\kappa_1}
\text{ for all }\Delta\in\mathcal{V}\bigr\},
\end{equation*}
the closed polydisc, whose interior is the open polydisc $P_{\mathcal{V}}(e^{\kappa_1})$.

Let $S=\mathbb{D}$, or $S=\mathbb{D}\times S'$, where $S'$ is the domain in the remaining variables
($A'$, $B'$ or $B$) of the statement of \cite[\S1]{B-CMP116} under consideration; the factor
$\mathbb{D}$ is closed, and no closedness is required of $S'$. Write $\operatorname{int}S$ for the
interior of $S$ in the ambient complex space. A function $G$ on $S$ is \emph{analytic on $S$} if
\begin{enumerate}
\item[(b1)] $G$ is continuous on $S$, and
\item[(b2)] $G$ is holomorphic, in the sense of Definition~\ref{5.6:not-holomorphy-facts}, at every
point of $\operatorname{int}S$.
\end{enumerate}
For $\mathfrak{F}$-valued $G$ both conditions are understood componentwise. Since all norms on
$\mathfrak{F}\cong\mathbb{C}^n$ are equivalent and continuous
(Definition~\ref{5.6:not-holomorphy-facts}), they may equivalently be understood in any norm on
$\mathfrak{F}$.
\end{enumerate}

We read in the sense of part (b) the phrases ``an analytic function of $s(Y_0)$'',
``of $A'$, $s(Y_0)$'', ``of $B'$, $s(Y_0)$'' and ``of $s(Y_0)$, $B$'', each qualified in
\cite[\S1, text at (1.11), (1.14), (1.16)--(1.18), and the sentence after (1.20)]{B-CMP116} by
``on the domain $|s(Y_0)|\le e^{\kappa_1}$'', that is on $\mathbb{D}$, or by the domains
considered there.

The reading (b) is the weakest of the usual ones: each of the following two properties implies it:
\begin{itemize}
\item $G$ is holomorphic on an open neighbourhood of $S$;
\item on $S$, $G$ is the sum of a series of functions, each continuous on $S$ and holomorphic at
every point of $\operatorname{int}S$, converging uniformly on $S$.
\end{itemize}
In the first case $G$ is holomorphic, hence continuous, at every point of the neighbourhood, in
particular on $S$ and at every point of $\operatorname{int}S$. In the second case (b1) holds
because a uniform limit of continuous functions is continuous. For (b2), on every closed polydisc
contained in $\operatorname{int}S$ each term equals its Cauchy integral over the distinguished
boundary of that polydisc (Definition~\ref{5.6:not-holomorphy-facts}); uniform convergence carries
this representation over to $G$, and a Cauchy integral of a function continuous on the
distinguished boundary is holomorphic inside the polydisc, since the Cauchy kernel expands in a
geometric series converging uniformly on compact subsets of the open polydisc.

In the representation \cite[text at (1.23)]{B-CMP116}, derivatives are represented by Cauchy
integrals in
the variables $s(\Delta)$, $\Delta\in\mathcal{V}$, written $\sigma(\Delta)$ as integration
variables in \cite{B-CMP116}; the $\sigma(\Delta)$-integrations run over the circles
$|\sigma(\Delta)|=e^{\kappa_1}$. The product of these circles lies in the boundary of
$\mathbb{D}$, not in its interior. On these circles the Cauchy integral formula of
Definition~\ref{5.6:not-holomorphy-facts} applies to functions continuous on the
closed disc and holomorphic inside it. This representation therefore uses at least the two
properties (b1) and (b2).

Conditions (b1), (b2) define analyticity on an arbitrary subset $S$ of the ambient complex space.
When $S$ is open, $\operatorname{int}S=S$, (b1) follows from (b2), and analyticity on $S$ is
holomorphy on $S$ in the sense of Definition~\ref{5.6:not-holomorphy-facts}.
\end{definition}

\begin{definition}[The term functional and the substituted field]\label{5.6:def-substituted-field}
Throughout, $k$ is the level, $g_k$ the coupling at level $k$, and $\kappa_1$ the exponent of
the decoupling radius $e^{\kappa_1}$ of \cite[\S1]{B-CMP116}. The letter $\kappa_1$ denotes the
same exponent in the weight lemma. The space $\mathfrak{F}$ of
$\mathfrak{g}^{\mathbb{C}}$-valued fields on a finite set of bonds, with
$\mathfrak{g}^{\mathbb{C}} = \mathfrak{sl}(N,\mathbb{C})$, is that of
Notation~\ref{5.6:not-holomorphy-facts} (holomorphy in finitely many complex variables).

\smallskip
\noindent\emph{The term functional.} Let
$\mathcal{T} = \mathcal{T}(X;\mathbf{U},\mathbf{J};A)$ be one of the functions $E(\,\cdot\,)$ of
the expansion \cite[(1.10)]{B-CMP116}. That is, $\mathcal{T}$ is a term
$E(X;\mathbf{U},\mathbf{J};A)$ of that expansion, which in \cite[\S1]{B-CMP116} arises by
substituting a nonlocal expression for the argument $A$. As a function of
$(\mathbf{U},\mathbf{J},A)$ it is analytic for $(\mathbf{U},\mathbf{J})$ in the domain of
\cite[\S1]{B-CMP116} and for $A$ satisfying the bounds \cite[(3.31)]{B-CMP109}. We write
$\mathfrak{A} \subset \mathfrak{F}$ for the open set of fields $A$ satisfying the bounds
\cite[(3.31)]{B-CMP109} with strict inequalities. The letter
$X$ is the localisation domain of the term; it is distinct from the test function on which one of
the operators depending on $s(Y_0)$ acts in \cite[\S1]{B-CMP116}.

The running example of \cite[\S1]{B-CMP116} is the last term of \cite[(3.34)]{B-CMP109}, the
fifth-order term bounded by \cite[Lemma~4, (3.54)]{B-CMP109}. The internal structure of
$\mathcal{T}$ is not used in this section. This structure comprises the box representation
\cite[(3.6)--(3.8)]{B-CMP109}, the gauges and the potential \cite[(3.26)--(3.32)]{B-CMP109}, and
the fifth-order remainder. The statements of this section are made for any $\mathcal{T}$
satisfying the analyticity hypothesis imposed on the term functional in this section, and
$\mathcal{T}$ enters them only through its values.

\smallskip
\noindent\emph{Decoupling variables.} In \cite[(1.10)]{B-CMP116} the sum runs over connected
domains $Y_0$ containing the fixed cube-neighbourhood $\tilde{\square}^{4}$ of the block
$\square$, with parameters $s = 0$ on $Y_0^{c}$; one writes $s = s(Y_0)$. For the term indexed by
$Y_0$ put
\begin{equation*}
\mathcal{V} := \{\Delta : \Delta \text{ an } M\text{-cube of } Y_0 \setminus \tilde{\square}^{4}\},
\qquad \sigma_v := s(\Delta) \quad (v = \Delta \in \mathcal{V}).
\end{equation*}
The exempt set $Y_1$ is the set of $M$-cubes of $\tilde{\square}^{4}$ (recall
$Y_0 \supseteq \tilde{\square}^{4}$). The values $s = 0$ on $Y_0^{c}$ are not variables of the
term indexed by $Y_0$. Put
\begin{equation*}
\mathbb{D} := \mathbb{D}_{\mathcal{V}}(e^{\kappa_1})
:= \{\sigma \in \mathbb{C}^{\mathcal{V}} : |\sigma_v| \le e^{\kappa_1} \text{ for all }
v \in \mathcal{V}\}.
\end{equation*}
This closed polydisc is the domain $|s(Y_0)| \le e^{\kappa_1}$ of \cite[\S1]{B-CMP116}. Its
interior $P := P_{\mathcal{V}}(e^{\kappa_1})$ is the open polydisc of the weight lemma.

\smallskip
\noindent\emph{Data.} A datum is a quadruple
$\mathfrak{d} := (\mathbf{U},\mathbf{J},B,\mathfrak{t}) \in \mathcal{D}$, where
\begin{equation*}
\mathcal{D} := \mathcal{D}^{UJ} \times \mathcal{D}^{B} \times [0,1].
\end{equation*}
The factor $\mathcal{D}^{UJ}$ is the intersection of two domains:
\begin{itemize}
\item the domain of configurations $(\mathbf{U},\mathbf{J})$ given by the conditions (i)--(iii)
of \cite{B-CMP109} with some constants $\alpha_0'$, $\alpha_1'$, as recalled in
\cite[text after (1.18)]{B-CMP116}, on
which all the functions and operators of \cite[(1.10)]{B-CMP116} are analytic and all their
bounds are uniform;
\item the domain $U^{c}_{k+1}(\square_0,(1+2\beta)\alpha_0,(1+2\beta)\alpha_1,\alpha_0)$ of
the term, as in \cite[\S1]{B-CMP116}: the space of configurations at level $k+1$ attached to
the block $\square_0$, with the parameters displayed, on which the term is analytic in
\cite[\S1]{B-CMP116}.
\end{itemize}
The factor $\mathcal{D}^{B}$ is the set of fields $B$ satisfying
\eqref{5.6:eq-auxiliary-constants-c}, that is, $g_k|B| < \varepsilon_1$. Here $B$ is the
fluctuation field restricted to $Y_0$, and $|B|$ is its norm in \cite[\S1]{B-CMP116}. Complex $B$
is allowed, as in the statement of \cite[\S1]{B-CMP116} that $H_k$ is an analytic function of
$s(Y_0)$ and $B$.

The last component $\mathfrak{t}$ is the real interpolation parameter of
\cite[(3.4), (3.6)]{B-CMP109}. It ranges over $[0,1]$ and is held fixed as part of the datum. It
is the first, unsubscripted $t$ of \cite[(1.1)]{B-CMP116}, and is written $\mathfrak{t}$ to keep
the letter $t$ free for the analytic functional of the weight lemma, whose composed function is
$F_t$. For version (a) below a fixed value
$t_{\square 0}$ is adjoined to the datum.

The weight lemma's data set $\mathcal{D}_b$ is matched with $\mathcal{D}$ by letting $b$ stand for
the triple $(\alpha_0',\alpha_1',\varepsilon_1)$. This matching is a naming only; no property of
$\mathcal{D}_b$ is used.

\smallskip
\noindent\emph{The substituted field and the composed term.} Let $H_k(s(Y_0),B')$ be the function
\cite[(1.17)]{B-CMP116}. It collects all the operators of \cite[(1.10)]{B-CMP116} depending on
$s(Y_0)$, together with the solution of the separated equations \cite[(1.4)]{B-CMP116}. Let
$B' = B'(B)$ be the function
\cite[(3.2)]{B-CMP109}, given by \cite[(1.19)]{B-CMP116}; it is a local function of $B$ and does
not depend on $s$. In \cite[text after (1.21)]{B-CMP116}, the function $H_k$ of
\cite[(1.21)]{B-CMP116}, multiplied by $\mathfrak{t}\zeta_{\ast} + t_{\square}\zeta_{\square}$,
is required to satisfy the bounds \cite[(3.31)]{B-CMP109}.

\emph{Standing form of the term.} We assume that the $s(Y_0)$-dependence of the term runs only
through the field substituted for $A$, as read from \cite[(1.1), text after (1.21)]{B-CMP116};
the case of further $s(Y_0)$-dependent arguments is treated in the paragraph on further arguments
below. Accordingly, for $s \in \mathbb{D}$, $t_\square \in \mathbb{C}$ and
$\mathfrak{d} \in \mathcal{D}$ we put
\begin{equation}\label{5.6:eq-substituted-field-a}
\begin{aligned}
A(s,t_\square;\mathfrak{d}) &:= \bigl(\mathfrak{t}\zeta_{\ast}
  + t_\square\zeta_{\square}\bigr)\, H_k\bigl(s(Y_0),B'(B)\bigr),\\
F(s,t_\square;\mathfrak{d}) &:= \mathcal{T}\bigl(X;\mathbf{U},\mathbf{J};
  A(s,t_\square;\mathfrak{d})\bigr).
\end{aligned}
\end{equation}
The function $F$ is defined at every $(s,t_\square)$ at which $A(s,t_\square;\mathfrak{d})$ lies
in $\mathfrak{A}$. In the notation of the weight lemma, the substitution map is
$\mathsf{R}(\sigma;\mathfrak{d}) := A(\sigma,t_\square;\mathfrak{d})$, and the function $F_t$ is
one of the two versions of $F$ given below.

\smallskip
\noindent\emph{The cut-off functions.} The functions $\zeta_{\ast}$ and
$\zeta_{\square}$ are the two fixed cut-off functions of \cite[(1.1)]{B-CMP116}, where the first
is printed as $\tilde{\zeta}_{\square}$, and
$\zeta_\square$ is the function of \cite[(3.4)]{B-CMP109}. Neither depends on $s$, $t_\square$,
$B$ or $(\mathbf{U},\mathbf{J})$, and both act on fields by pointwise multiplication. They
therefore define $\mathbb{C}$-linear maps
$\mathsf{M}_{\zeta_{\ast}},\ \mathsf{M}_{\zeta_{\square}} : \mathfrak{F} \to \mathfrak{F}$.
No property of $\zeta_{\ast}$ beyond this is used.

Writing $H := H_k(s(Y_0),B'(B))$, the first line of \eqref{5.6:eq-substituted-field-a} reads
\begin{equation}\label{5.6:eq-substituted-field-1}
A(s,t_\square;\mathfrak{d}) = \mathfrak{t}\,\mathsf{M}_{\zeta_{\ast}} H
  + t_\square\,\mathsf{M}_{\zeta_{\square}} H .
\end{equation}
Hence $A(s,t_\square;\mathfrak{d})$ is $\mathbb{C}$-linear in the vector $H \in \mathfrak{F}$ and
affine in $t_\square$. Of the parameter $\mathfrak{t}$ only two things are used: that it is fixed,
and that $|\mathfrak{t}| \le 1$. Two kinds of input about $A$ are used below: the linear structure
\eqref{5.6:eq-substituted-field-1}, together with $|\mathfrak{t}| \le 1$; and, where a bound is
taken directly from \cite{B-CMP116}, only the three statements of
\cite[text after (1.21)]{B-CMP116} about the product with
$\mathfrak{t}\zeta_{\ast}$, about the product with $\zeta_{\square}$, and about the
analytic extension in $t_\square$ made there.

\smallskip
\noindent\emph{Further arguments of the term.} Suppose $\mathcal{T}$ carries further arguments
depending on $s(Y_0)$, namely operators of \cite[(1.10)]{B-CMP116} not routed through $H_k$. Each
of them is stated in \cite[\S1]{B-CMP116} to be analytic on the domain
$|s(Y_0)| \le e^{\kappa_1}$, in the sense of Definition~\ref{5.6:def-closed-analyticity}
(analyticity on a closed domain), with bounds carrying the factor $e^{16\kappa_1}$. The joint
analytic extension is made in \cite[text after (1.21) and after (1.22)]{B-CMP116}. The proofs
below then treat each such $s(Y_0)$-dependent argument of $\mathcal{T}$ in turn, by composition,
under the analyticity hypothesis on $\mathcal{T}$ widened to include these arguments.

\smallskip
\noindent\emph{Constants and the $t_\square$-radius.} The constants are those of
\cite[\S1]{B-CMP116}:
\begin{itemize}
\item $B_0$ is the propagator constant;
\item $C_1$ is the absolute constant of \cite[(1.20)]{B-CMP116};
\item $\varepsilon_1$ is the constant of \eqref{5.6:eq-auxiliary-constants-c}.
\end{itemize}
The constant $\alpha_2$ is that of \cite[(3.13)--(3.14), (3.31)]{B-CMP109}. Put
\begin{equation}\label{5.6:eq-substituted-field-b}
\nu(\mathfrak{d}) := 4B_0C_1e^{16\kappa_1}g_k|B| .
\end{equation}
This is the bound stated in \cite[text after (1.21)]{B-CMP116} for the norms of
\cite[(3.31)]{B-CMP109} of the product with $\zeta_\square$. Since $g_k|B| < \varepsilon_1$ on
$\mathcal{D}$,
\begin{equation*}
\nu(\mathfrak{d}) < 4B_0C_1e^{16\kappa_1}\varepsilon_1 .
\end{equation*}
The proper fractions of $\alpha_2$ in \cite[text after (1.21)]{B-CMP116} are denoted as follows.
\begin{itemize}
\item $\mathfrak{c}_1 \in \left]0,1\right]$ is the fraction in the assumption
$4B_0C_1e^{16\kappa_1}\varepsilon_1 \le \mathfrak{c}_1\alpha_2$.
\item $\mathfrak{c}_2 \in \left]0,1\right[$ is the fraction with which the product with
$\mathfrak{t}\zeta_{\ast}$ satisfies \cite[(3.31)]{B-CMP109}.
\item $\mathfrak{c}_3 \in \left]0,1\right]$ is the fraction in the condition
$|t_\square|\,\nu(\mathfrak{d}) \le \mathfrak{c}_3\alpha_2$, under which the expression
\cite[(1.10)]{B-CMP116} is extended analytically in $t_\square$.
\end{itemize}
We write $\theta$ for the threshold of the bounds \cite[(3.31)]{B-CMP109}, formed with $\alpha_2$,
and $\theta_c$ for the threshold of the same bounds formed with $\mathfrak{c}_2\alpha_2$.

The $t_\square$-radius $R_\square(\mathfrak{d})$, which for $\nu(\mathfrak{d}) > 0$ is the radius
of the circle \cite[(1.22)]{B-CMP116}, on which equality holds in the condition above, and its
infimum $R_{\square,\ast}$ over $\mathcal{D}$ are those fixed in the definition of the response
radius.
The $t_\square$-radius $R_\square(\mathfrak{d})$ is distinct from the radius $e^{\kappa_1}$ of
$\mathbb{D}$, which is the radius of the polydisc of the weight lemma.

\smallskip
\noindent\emph{The two versions of $F_t$.} For $\sigma \in \mathbb{D}$ and
$\mathfrak{d} \in \mathcal{D}$ the weight lemma's function is taken in one of two versions.
\begin{itemize}
\item[(a)] $F_t(\sigma;\mathfrak{d}) := F(\sigma,t_{\square 0};\mathfrak{d})$, at a fixed
$t_{\square 0} \in \mathbb{C}$ with $|t_{\square 0}| < R_{\square,\ast}$; since
$R_{\square,\ast} \le R_\square(\mathfrak{d})$, such a point lies in the disc
$|t_\square| \le R_\square(\mathfrak{d})$ of \cite[(1.22)]{B-CMP116} for every
$\mathfrak{d} \in \mathcal{D}$.
\item[(b)] $F_t(\sigma;\mathfrak{d})$ is the derivative of $F$ in $t_\square$ at the origin:
\begin{equation*}
F_t(\sigma;\mathfrak{d}) := \partial_{t_\square}F(\sigma,t_\square;\mathfrak{d})
  \big|_{t_\square = 0} .
\end{equation*}
This is the derivative taken in \cite[(1.23)]{B-CMP116}. There the expansion is differentiated
with respect to $t_\square$ at $t_\square = 0$, and all derivatives are represented by the Cauchy
formula.
\end{itemize}

\smallskip
\noindent\emph{The bound letter.} Put
\begin{equation*}
\mathfrak{b}_{\mathcal{T}} := \sup\bigl\{\,|\mathcal{T}(X;\mathbf{U},\mathbf{J};A)| :
A \in \mathfrak{A},\ (\mathbf{U},\mathbf{J}) \in \mathcal{D}^{UJ}\bigr\}.
\end{equation*}
For the running example, \cite[(3.54)]{B-CMP109} bounds the term at the substituted small
potential. In \cite[\S1]{B-CMP116} the Cauchy-represented term \cite[(1.23)]{B-CMP116} is
estimated by means of \cite[(3.54)]{B-CMP109}, with the result \cite[(1.24)]{B-CMP116}.
That result bounds the represented term, and so carries $\mathfrak{b}_{\mathcal{T}}$ together with
the Cauchy factors: the factor $1/R_\square(\mathfrak{d})$ of the $t_\square$-circle and the
factors of the circles $|\sigma_v| = e^{\kappa_1}$, $v \in \mathcal{V}$. The supremum
$\mathfrak{b}_{\mathcal{T}}$ itself is the quantity that the analyticity hypothesis on
$\mathcal{T}$ takes to be finite.

For the function obtained by differentiation in $t_\square$, the bound
\cite[(3.17)]{B-CMP109} is a constant times $\exp(-\kappa d_j(X))$ times the size
$|\zeta_\square H_k(B')|$; for a big domain $X$ the exponential factor is absent. No value
of $\mathfrak{b}_{\mathcal{T}}$ is used in this section, and the weight lemma requires none.
\end{definition}

\begin{definition}[Conditions on the auxiliary constants and their order of choice]
\label{5.6:def-auxiliary-constants}
Let the data $\mathfrak{d}$, the cube-neighbourhood $\tilde{\square}^{4}$ of the block $\square$, the
connected localisation domain $Y_{0}\supset\tilde{\square}^{4}$, the complex parameters
$s(Y_{0})=(s(\Delta))_{\Delta\subset Y_{0}\setminus\tilde{\square}^{4}}$, the field $B'=B'(B)$, the
function $H_{k}(s(Y_{0}),B')$, the parameter $\mathfrak{t}\in[0,1]$, the cut-offs $\tilde{\zeta}_{\square}$,
$\zeta_{\square}$ and the quantity $\nu(\mathfrak{d})$ of \eqref{5.6:eq-substituted-field-b} be as in
Definition~\ref{5.6:def-substituted-field}. The constants
$B_{0}$, $\delta_{0}$, $C_{1}$, $C_{2}$, $C_{4}$ are those of \cite[\S1]{B-CMP116}: $B_{0}$ is the
prefactor and $\delta_{0}$ the decay constant of the bound \cite[(1.11)]{B-CMP116}, $C_{2}$ and $C_{4}$
are the constants of \cite[(1.13)--(1.16)]{B-CMP116}, and $C_{1}$ is the absolute constant of
\cite[(1.20)]{B-CMP116}. The constants $\alpha_{2}$ and $\alpha_{3}$ are those of \cite[\S3]{B-CMP109}:
$\alpha_{2}$ enters the bounds \cite[(3.31)]{B-CMP109} on the field $A$ on $\tilde{\square}^{4}$, and
$\alpha_{3}$ is the bound on $B'$ in the hypotheses of \cite[Lemma~4]{B-CMP109}. The numbers
$\mathfrak{c}_{1},\mathfrak{c}_{3}\in\left]0,1\right]$ and $\mathfrak{c}_{2}\in\left]0,1\right[$ are the
fractions of $\alpha_{2}$ in the text of \cite[after (1.21)]{B-CMP116}: $\mathfrak{c}_{1}$ in the
condition reproduced in (C3), $\mathfrak{c}_{2}$ in the bound for the product with
$\mathfrak{t}\tilde{\zeta}_{\square}$, and $\mathfrak{c}_{3}$ in the condition
$|t_{\square}|\,\nu(\mathfrak{d})\le\mathfrak{c}_{3}\alpha_{2}$ on the complex parameter $t_{\square}$
under which the term is extended analytically there.

The auxiliary constants $\kappa_{1}$, $M$, $\varepsilon_{1}$, $\varepsilon_{2}$, $\varepsilon_{3}$ are
chosen in the following order: first $\kappa_{1}$, which may be taken as large as desired; then $M$,
subject to a lower bound depending on $\kappa_{1}$; then $\varepsilon_{3}$, $\varepsilon_{2}$,
$\varepsilon_{1}$, each small depending on $\kappa_{1}$; last the coupling, small enough that at every step
the range of $g_{k}|B|$ allowed by the small-field conditions lies below $\varepsilon_{1}$. No later choice
imposes a condition on $\kappa_{1}$, on $M$ or on $L$. The conditions are the following; every one of them
is one-sided.
\begin{itemize}
\item[(C0)] $\kappa_{1}$ is a sufficiently big positive number \cite[before (1.11)]{B-CMP116}. This is a
lower bound on $\kappa_{1}$, used for the gains of the later Cauchy estimates. The analyticity statements
of this section for the term functional of Definition~\ref{5.6:def-substituted-field} and for
$H_{k}(s(Y_{0}),B'(B))$ hold for every $\kappa_{1}$ for which the statements of
\cite[(1.11)--(1.20) and the sentence after (1.20)]{B-CMP116} hold; any further lower bound on
$\kappa_{1}$ imposed where these statements are used is a condition of the same kind.
\item[(C1)] With $\delta_{1}>0$ a constant depending on $\delta_{0}$ only, as in
\cite[after (1.11)]{B-CMP116},
\begin{equation}\label{5.6:eq-auxiliary-constants-a}
\delta_{1}M\ \ge\ \kappa_{1}.
\end{equation}
This is a lower bound on $M$, namely $M\ge\kappa_{1}/\delta_{1}$, with $M$ chosen after $\kappa_{1}$; it is
not an upper bound on $\kappa_{1}$. It enters as follows. In the bound \cite[(1.11)]{B-CMP116} of the
$s$-dependent propagator by a supremum over random walks, the expression under the supremum is at most
$1$, under \eqref{5.6:eq-auxiliary-constants-a}, for the walks with which more than $2^{4}$ of the
parameters $s$ are connected, and at most $e^{16\kappa_{1}}$ for the remaining walks. In this way
\eqref{5.6:eq-auxiliary-constants-a} yields the bound of \cite[after (1.11)]{B-CMP116} by
$B_{0}e^{16\kappa_{1}}$ times the size of the function to which the propagator is applied, in all the
norms used there, on the whole closed polydisc
\begin{equation}\label{5.6:eq-auxiliary-constants-1}
|s(\Delta)|\le e^{\kappa_{1}}\qquad\text{for all }\Delta\subset Y_{0}\setminus\tilde{\square}^{4},
\end{equation}
and hence the bounds \cite[(1.14)--(1.21)]{B-CMP116}, each derived from this
$B_{0}e^{16\kappa_{1}}$-bound, on that
polydisc. Thus \eqref{5.6:eq-auxiliary-constants-a} is a hypothesis of the analyticity of
$H_{k}(s(Y_{0}),B'(B))$ on \eqref{5.6:eq-auxiliary-constants-1} stated in
\cite[after (1.20)]{B-CMP116}, and, through \cite[(1.21)]{B-CMP116}, of the statement that the
substituted field of Definition~\ref{5.6:def-substituted-field} satisfies \cite[(3.31)]{B-CMP109}.
\item[(C2)] With $c_{\mathrm{abs}}$ the absolute constant of \cite[after (1.21)]{B-CMP116},
\begin{equation}\label{5.6:eq-auxiliary-constants-2}
e^{32\kappa_{1}}\varepsilon_{1}\ \le\ c_{\mathrm{abs}} .
\end{equation}
This is an upper bound on $\varepsilon_{1}$, which is chosen after $\kappa_{1}$. In
\cite[after (1.21)]{B-CMP116} the constant $c_{\mathrm{abs}}$ is obtained by inspection of the conditions
preceding it, which are:
\begin{itemize}
\item[(C2a)] $4C_{2}B_{0}e^{16\kappa_{1}}\varepsilon_{2}\le 1$ \cite[between (1.13) and (1.14)]{B-CMP116};
\item[(C2b)] $4C_{4}B_{0}^{2}e^{32\kappa_{1}}\varepsilon_{3}\le 1$ \cite[after (1.15)]{B-CMP116};
\item[(C2c)] $2B_{0}e^{16\kappa_{1}}\varepsilon_{3}\le\varepsilon_{2}$
\cite[between (1.15) and (1.16)]{B-CMP116};
\end{itemize}
together with the domain condition $|A'|<\varepsilon_{2}$ on the field $A'$ of the change of variables
\cite[(1.12)]{B-CMP116}, imposed before \cite[(1.13)]{B-CMP116}, the domain condition
$|B'|<\varepsilon_{3}$ \cite[after (1.15)]{B-CMP116}, and the bound
\begin{equation}\label{5.6:eq-auxiliary-constants-3}
|B'|\ \le\ C_{1}g_{k}|B|\ <\ C_{1}\varepsilon_{1}
\end{equation}
of \cite[(1.20)]{B-CMP116}, which tie $\varepsilon_{3}$ and $\varepsilon_{2}$ to $\varepsilon_{1}$; with the
choice of $\varepsilon_{3}$ and $\varepsilon_{2}$ made in the proof below, (C2c) holds identically and
(C2a), (C2b) become upper bounds on $e^{32\kappa_{1}}\varepsilon_{1}$, which is how
\eqref{5.6:eq-auxiliary-constants-2} gathers them in the sentence of \cite[after (1.21)]{B-CMP116}. The
constants $\varepsilon_{1}$, $\varepsilon_{2}$, $\varepsilon_{3}$ are the auxiliary constants of
\cite[\S1]{B-CMP116}, chosen after $\kappa_{1}$; the one that binds is $\varepsilon_{1}$, the bound of
$g_{k}|B|$ in (C4), which is further bounded above in (C3) and (C5).
\item[(C3)] As in \cite[after (1.21)]{B-CMP116},
\begin{equation}\label{5.6:eq-auxiliary-constants-b}
4B_{0}C_{1}e^{16\kappa_{1}}\varepsilon_{1}\ \le\ \mathfrak{c}_{1}\alpha_{2}.
\end{equation}
This is an upper bound on $\varepsilon_{1}$, chosen after $\kappa_{1}$ and after $\alpha_{2}$. It enters as
follows: \cite[after (1.21)]{B-CMP116} states that it implies that the product of $H_{k}$ with
$\mathfrak{t}\tilde{\zeta}_{\square}$ satisfies \cite[(3.31)]{B-CMP109} with
$\mathfrak{c}_{2}\alpha_{2}$; and, since by \eqref{5.6:eq-substituted-field-b}
and (C4)
\begin{equation}\label{5.6:eq-auxiliary-constants-4}
\nu(\mathfrak{d})=4B_{0}C_{1}e^{16\kappa_{1}}g_{k}|B|<4B_{0}C_{1}e^{16\kappa_{1}}\varepsilon_{1}
\le\mathfrak{c}_{1}\alpha_{2},
\end{equation}
it bounds $\nu(\mathfrak{d})$ uniformly in the data. Hence it is an input of the statement that the
substituted field of Definition~\ref{5.6:def-substituted-field} satisfies \cite[(3.31)]{B-CMP109}, and,
through the uniform bound \eqref{5.6:eq-auxiliary-constants-4}, of the positive lower bound, uniform in
the data, on the radius of the disc in $t_{\square}$ on which the term is extended analytically; that
radius and its lower bound are fixed in the definition of the response radius in this section.
\item[(C4)] As in \cite[(1.20)]{B-CMP116},
\begin{equation}\label{5.6:eq-auxiliary-constants-c}
g_{k}|B|\ <\ \varepsilon_{1}.
\end{equation}
This is the range of the datum $B$; it defines the factor $\mathcal{D}^{B}$ of the set $\mathcal{D}$ of
admissible data of Definition~\ref{5.6:def-substituted-field}, and no bound on $g_{k}|B|$ beyond
\eqref{5.6:eq-auxiliary-constants-c} is used. The
quantity $g_{k}|B|$ is the coupling times the size of the fluctuation field on the
region read by the term; since the coupling is chosen last, (C2)--(C4) are upper bounds on the coupling
fixed after $\kappa_{1}$.
\item[(C5)] Condition \eqref{5.6:eq-auxiliary-constants-5} places the bound
\eqref{5.6:eq-auxiliary-constants-3} of \cite[(1.20)]{B-CMP116} inside the hypothesis $|B'|<\alpha_{3}$ of
\cite[Lemma~4]{B-CMP109}; it is carried among the hypotheses of Theorem~0 as (H5):
\begin{equation}\label{5.6:eq-auxiliary-constants-5}
C_{1}\varepsilon_{1}\ \le\ \alpha_{3}.
\end{equation}
This is an upper bound on $\varepsilon_{1}$, chosen after $\kappa_{1}$, of the same kind and at the same
place in the order as (C4). Its origin is the hypothesis $|B'|<\alpha_{3}$ of \cite[Lemma~4]{B-CMP109},
under which the functions of \cite[(3.53)]{B-CMP109} are analytic on the spaces of configurations there,
read together with \eqref{5.6:eq-auxiliary-constants-3}. No value of $\alpha_{3}$ is used.
\end{itemize}
In this order, for every $\kappa_{1}$ satisfying (C0) there is an $M$ satisfying (C1), and then, for these
$\kappa_{1}$ and $M$, every sufficiently small $\varepsilon_{1}>0$ satisfies (C2), (C3) and (C5), and
(C2a)--(C2c) hold with $\varepsilon_{3}:=C_{1}\varepsilon_{1}$ and
$\varepsilon_{2}:=2B_{0}e^{16\kappa_{1}}\varepsilon_{3}$; none of (C2), (C3), (C5) bounds $\kappa_{1}$
from above, since a condition $f(\kappa_{1})\,\varepsilon_{1}\le c$ with $\varepsilon_{1}$ chosen after
$\kappa_{1}$ is an upper bound on $\varepsilon_{1}$, not on $\kappa_{1}$. Moreover, under (C4) and
(C5) the field $B'$ satisfies the hypothesis $|B'|<\alpha_{3}$ of \cite[Lemma~4]{B-CMP109}.
\end{definition}

\begin{proof}
Let $\kappa_{1}$ satisfy (C0). Since $\delta_{1}>0$ depends on $\delta_{0}$ only, and $M$ is chosen after
$\kappa_{1}$, any $M\ge\kappa_{1}/\delta_{1}$ satisfies \eqref{5.6:eq-auxiliary-constants-a}.

Now fix $\kappa_{1}$ and $M$. The conditions \eqref{5.6:eq-auxiliary-constants-2},
\eqref{5.6:eq-auxiliary-constants-b} and \eqref{5.6:eq-auxiliary-constants-5}, together with the two
inequalities
\begin{equation*}
8C_{1}C_{2}B_{0}^{2}e^{32\kappa_{1}}\varepsilon_{1}\le 1,
\qquad
4C_{1}C_{4}B_{0}^{2}e^{32\kappa_{1}}\varepsilon_{1}\le 1,
\end{equation*}
hold simultaneously for every $\varepsilon_{1}$ with
\begin{equation*}
\begin{split}
0<\varepsilon_{1}\le\min\bigl\{&c_{\mathrm{abs}}e^{-32\kappa_{1}},\
\mathfrak{c}_{1}\alpha_{2}\,(4B_{0}C_{1})^{-1}e^{-16\kappa_{1}},\ \alpha_{3}C_{1}^{-1},\\
&(8C_{1}C_{2}B_{0}^{2})^{-1}e^{-32\kappa_{1}},\ (4C_{1}C_{4}B_{0}^{2})^{-1}e^{-32\kappa_{1}}\bigr\}.
\end{split}
\end{equation*}
For such $\varepsilon_{1}$ put $\varepsilon_{3}:=C_{1}\varepsilon_{1}$ and
$\varepsilon_{2}:=2B_{0}e^{16\kappa_{1}}\varepsilon_{3}=2B_{0}C_{1}e^{16\kappa_{1}}\varepsilon_{1}$. For
data satisfying \eqref{5.6:eq-auxiliary-constants-c}, the bound \eqref{5.6:eq-auxiliary-constants-3} gives
the domain condition $|B'|<\varepsilon_{3}$; (C2c) holds with
equality; (C2a) reads $4C_{2}B_{0}e^{16\kappa_{1}}\cdot 2B_{0}C_{1}e^{16\kappa_{1}}\varepsilon_{1}\le1$,
that is the first of the two inequalities displayed above; and (C2b) reads
$4C_{4}B_{0}^{2}e^{32\kappa_{1}}\cdot C_{1}\varepsilon_{1}\le 1$, that is the second. Both are upper
bounds on $e^{32\kappa_{1}}\varepsilon_{1}$ by constants depending on $B_{0}$, $C_{1}$, $C_{2}$, $C_{4}$
only, of the form \eqref{5.6:eq-auxiliary-constants-2}; they are among the conditions by whose inspection
the constant $c_{\mathrm{abs}}$ is obtained in \cite[after (1.21)]{B-CMP116}.

Finally, assume (C4) and (C5). By \eqref{5.6:eq-auxiliary-constants-3}, which is the bound of
\cite[(1.20)]{B-CMP116} and rests on \eqref{5.6:eq-auxiliary-constants-c}, and by
\eqref{5.6:eq-auxiliary-constants-5},
\begin{equation*}
|B'|\ \le\ C_{1}g_{k}|B|\ <\ C_{1}\varepsilon_{1}\ \le\ \alpha_{3},
\end{equation*}
which is the hypothesis on $B'$ in \cite[Lemma~4]{B-CMP109}.
\end{proof}

\begin{definition}[Analyticity and boundedness of the term functional]\label{5.6:def-term-analyticity}
Work in the setting of Notation~\ref{5.6:not-holomorphy-facts} and of Definitions~\ref{5.6:def-closed-analyticity} and~\ref{5.6:def-substituted-field}, with the conventions on the cut-off functions and on the interpolation parameter $\mathfrak{t}$ that enter the substituted field of Definition~\ref{5.6:def-substituted-field}, and with the conditions of Definition~\ref{5.6:def-auxiliary-constants} in force throughout. Fix the step $k$, the block $\square$, and a connected localisation domain $Y_{0}$ containing $\tilde{\square}^{4}$. Let $\mathcal{T}$ be the term functional of Definition~\ref{5.6:def-substituted-field}, with its localisation domain $X$. Let $\mathcal{D}$ be the set of admissible data $\mathfrak{d}=(\mathbf{U},\mathbf{J},B,\mathfrak{t})$, and $\mathcal{D}^{UJ}$ its $(\mathbf{U},\mathbf{J})$-factor. Let $\mathfrak{F}$ be the finite-dimensional space of $\mathfrak{g}^{\mathbb{C}}$-valued fields of Notation~\ref{5.6:not-holomorphy-facts}.

Let $\mathcal{N}_{1}$, $\mathcal{N}_{2}$, $\mathcal{N}_{3}$ be the three quantities of a field $A$ on $\tilde{\square}^{4}$ that enter \cite[(3.31)]{B-CMP109}, and let $\theta$ be the threshold there. Put
\begin{equation}\label{5.6:eq-term-analyticity-1}
\mathfrak{A}:=\bigl\{A\in\mathfrak{F}:\ \mathcal{N}_{i}(A)<\theta\ \text{ for } i=1,2,3\bigr\}.
\end{equation}
This is the space of fields $A$ satisfying \cite[(3.31)]{B-CMP109}, on which \cite[text preceding Lemma~4]{B-CMP109} considers the functions in \cite[(3.28), (3.34)]{B-CMP109}, with all three inequalities strict.

We say that $\mathcal{T}$ is \emph{analytic and bounded} if:
\begin{enumerate}
\item[(i)] $\mathfrak{A}$ is an open subset of $\mathfrak{F}$, and, for every $(\mathbf{U},\mathbf{J})\in\mathcal{D}^{UJ}$, the map $A\mapsto\mathcal{T}(X;\mathbf{U},\mathbf{J};A)$ is holomorphic on $\mathfrak{A}$ in the sense of Notation~\ref{5.6:not-holomorphy-facts};
\item[(ii)] there is a constant $\mathfrak{b}_{\mathcal{T}}<\infty$ such that, for every $(\mathbf{U},\mathbf{J})\in\mathcal{D}^{UJ}$,
\begin{equation}\label{5.6:eq-term-analyticity-2}
\sup_{A\in\mathfrak{A}}\bigl|\mathcal{T}(X;\mathbf{U},\mathbf{J};A)\bigr|\le\mathfrak{b}_{\mathcal{T}}<\infty .
\end{equation}
\end{enumerate}
Throughout this section, $\mathcal{T}$ is assumed to be analytic and bounded.
\end{definition}

For the typical term of the expansion, the last term on the right-hand side of \cite[(3.34)]{B-CMP109}, clause~(i) is Ba{\l}aban's statement in \cite[\S1, p.~2]{B-CMP116}:
\begin{quote}
We have proved that it is an analytic function of $\mathbf{U}$, $\mathbf{J}$, $\mathbf{A}$, for $(\mathbf{U},\mathbf{J})\in U^{c}_{k+1}(\square_{0},(1+2\beta)\alpha_{0},(1+2\beta)\alpha_{1},\alpha_{0})$ and $\mathbf{A}$ satisfying (I.3.31).
\end{quote}
Here (I.3.31) is \cite[(3.31)]{B-CMP109}, so that the condition on $\mathbf{A}$ is $A\in\mathfrak{A}$.
In that statement, analytic is read as in Definition~\ref{5.6:def-closed-analyticity}. Since $\mathfrak{A}$ is open, every point of it is interior, so that reading amounts to holomorphy on $\mathfrak{A}$.

For the other terms of the list in \cite[\S1]{B-CMP116}, clause~(i) is the form in which the structure stated there is used below: there is a function, written $E(X,\mathbf{U},\mathbf{J},A)$ in \cite[\S1]{B-CMP116} and denoted $\mathcal{T}$ here, that is analytic and localised to $X$ in $\mathbf{U}$, $\mathbf{J}$, $A$, and a given term is obtained by substituting a proper nonlocal expression in place of $A$.

\cite[Lemma~4, (3.53)]{B-CMP109} states that the functions displayed there are analytic on the spaces above. Under the conditions of Definition~\ref{5.6:def-auxiliary-constants}, that lemma places the substituted configurations in the spaces on which the term is analytic. It is used here for this purpose only; it is not a statement about $\mathcal{T}$.

The openness of $\mathfrak{A}$ belongs to clause~(i). The inequalities in \cite[(3.31)]{B-CMP109} are strict. If $\mathcal{N}_{1}$, $\mathcal{N}_{2}$, $\mathcal{N}_{3}$ are seminorms on $\mathfrak{F}$, they are continuous by Notation~\ref{5.6:not-holomorphy-facts}. Then $\mathfrak{A}$ is the intersection of the three preimages of the open interval $(-\infty,\theta)$, and hence is open. Otherwise, $\mathfrak{A}$ is open because \cite{B-CMP109} treats the space of fields satisfying \cite[(3.31)]{B-CMP109} as a domain of analyticity.

The only condition that \cite[text between (1.21) and (1.22)]{B-CMP116} places on the substituted field, for the analytic extension in the response parameter, is that it satisfy the bounds \cite[(3.31)]{B-CMP109}. This is the condition $A\in\mathfrak{A}$.

Clause~(ii) is taken here as a hypothesis in the form of the supremum \eqref{5.6:eq-term-analyticity-2}. For the typical term, \cite[(3.54)]{B-CMP109} states the required bound on the spaces above, and \cite[text preceding (1.24)]{B-CMP116} estimates the Cauchy representation \cite[(1.23)]{B-CMP116} of the term by means of \cite[(3.54)]{B-CMP109}. The supremum \eqref{5.6:eq-term-analyticity-2}, uniform in $(\mathbf{U},\mathbf{J})\in\mathcal{D}^{UJ}$, is not a display of \cite{B-CMP109} or of \cite{B-CMP116}; clause~(ii) is a hypothesis of this section. For the other terms, Ba{\l}aban states bounds of the type \cite[(1.18)]{B-CMP109} on their domains of analyticity \cite[\S1]{B-CMP116}. Only the existence of the finite constant $\mathfrak{b}_{\mathcal{T}}$, uniform over $\mathcal{D}^{UJ}$, is used; no value of $\mathfrak{b}_{\mathcal{T}}$ enters any argument.

\begin{proposition}[Analyticity of the localized field in the decoupling parameters {\cite[\S1, (1.17)--(1.21)]{B-CMP116}}]
\label{5.6:prop-localized-field}
Let the auxiliary constants satisfy the conditions of
Definition~\ref{5.6:def-auxiliary-constants}. In particular:
\begin{itemize}
\item $\kappa_1$ is a sufficiently big positive number;
\item $\delta_1 M\ge\kappa_1$, condition~\eqref{5.6:eq-auxiliary-constants-a};
\item $e^{32\kappa_1}\varepsilon_1$ is smaller than the absolute constant $c_{\mathrm{abs}}$.
\end{itemize}

Write $\mathcal{D}^{UJ}$ for the $(\mathbf{U},\mathbf{J})$-factor of the set $\mathcal{D}$
of admissible data of Definition~\ref{5.6:def-substituted-field}. This is the domain of
configurations given by the conditions recalled in the text following
\cite[(1.18)]{B-CMP116}; on it all the functions and operators of the construction are
analytic in $\mathbf{U},\mathbf{J}$, with bounds uniform on the domain.

Write $\mathcal{D}^{B}$ for the $B$-factor of $\mathcal{D}$. On $\mathcal{D}^{B}$ one has
$g_k|B|<\varepsilon_1$, condition~\eqref{5.6:eq-auxiliary-constants-c}.

Let $Y_0$ be the connected localization domain containing $\tilde{\square}^{4}$, and let
$\mathcal{V}$ be the set of $M$-cubes $\Delta$ of $Y_0\setminus\tilde{\square}^{4}$. Put
\begin{equation*}
\mathbb{D}:=\bigl\{s=(s(\Delta))_{\Delta\in\mathcal{V}}:\ |s(\Delta)|\le e^{\kappa_1}
\ \text{for all }\Delta\in\mathcal{V}\bigr\}.
\end{equation*}
This is a closed polydisc, with interior $P=P_{\mathcal{V}}(e^{\kappa_1})$.

Let $\mathfrak{F}$ be the finite-dimensional space of
$\mathfrak{g}^{\mathbb{C}}=\mathfrak{sl}(N,\mathbb{C})$-valued lattice fields. Let $H_k$ be
the function of \cite[(1.17)]{B-CMP116}, and let $B'=B'(B)$ be the field of
\cite[(1.19)]{B-CMP116}, as in Definition~\ref{5.6:def-substituted-field}.

Then for every $(\mathbf{U},\mathbf{J})\in\mathcal{D}^{UJ}$ and every $B\in\mathcal{D}^{B}$
the following hold.
\begin{enumerate}
\item[(a)] The $\mathfrak{F}$-valued map $s\mapsto H_k(s,B'(B))$ is analytic on
$\mathbb{D}$ in the sense of Definition~\ref{5.6:def-closed-analyticity}. That is, it is
continuous on $\mathbb{D}$ and holomorphic on $P$.
\item[(b)] For all $s\in\mathbb{D}$, and in each of the norms called admissible in
\cite[(1.18), (1.21)]{B-CMP116},
\begin{equation}\label{5.6:eq-localized-field-a}
\bigl\|H_k(s,B'(B))\bigr\|\le 4B_0C_1e^{16\kappa_1}g_k|B|
<4B_0C_1e^{16\kappa_1}\varepsilon_1 .
\end{equation}
\item[(c)] The bound \eqref{5.6:eq-localized-field-a} is uniform in
$(\mathbf{U},\mathbf{J})\in\mathcal{D}^{UJ}$.
\end{enumerate}
\end{proposition}

This is \cite[\S1, (1.17)--(1.21)]{B-CMP116}. Statement~(a) is the sentence following
\cite[(1.20)]{B-CMP116}. The bound~\eqref{5.6:eq-localized-field-a} is
\cite[(1.21)]{B-CMP116}; its first inequality is the composition of the bound
\cite[(1.18)]{B-CMP116} with the bound \cite[(1.20)]{B-CMP116}. The uniformity in
$(\mathbf{U},\mathbf{J})$ is stated in the text following \cite[(1.18)]{B-CMP116}.

The proof is Ba{\l}aban's and is the argument of \cite[\S1, (1.11)--(1.20)]{B-CMP116}.
It combines four ingredients:
\begin{itemize}
\item the walk bound \cite[(1.11)]{B-CMP116}, used under the condition $\delta_1M\ge\kappa_1$;
\item the contraction arguments of \cite[(1.12)--(1.16)]{B-CMP116}, which give the fixed points
$D$, $V$ and $A_0$ as analytic functions on the domains considered there;
\item the formula \cite[(1.17)]{B-CMP116} with the bound \cite[(1.18)]{B-CMP116};
\item the bound \cite[(1.20)]{B-CMP116} on $B'=B'(B)$.
\end{itemize}

\begin{definition}[Membership of the substituted field in the analyticity domain; the response radius]
\label{5.6:def-response-radius}
Let $\mathcal{D}$ be the set of admissible data
$\mathfrak{d}=(\mathbf{U},\mathbf{J},B,\mathfrak{t})$, let $\mathbb{D}$ be the closed polydisc
$|s(\Delta)|\le e^{\kappa_1}$, $\Delta\in\mathcal{V}$, and let $A(s,t_{\square};\mathfrak{d})$ and
$\nu(\mathfrak{d})$ be the substituted field and the number of
Definition~\ref{5.6:def-substituted-field}, displays \eqref{5.6:eq-substituted-field-a} and
\eqref{5.6:eq-substituted-field-b}. Here $H_k(s,B'(B))$ is the localized field of
Proposition~\ref{5.6:prop-localized-field}, display \eqref{5.6:eq-localized-field-a}, and
$\zeta_{\ast}$, $\zeta_{\square}$ are the two fixed cut-off functions of \cite[(1.1)]{B-CMP116},
acting on $\mathfrak{F}$ by pointwise multiplication through the $\mathbb{C}$-linear maps
$\mathsf{M}_{\zeta_{\ast}}$, $\mathsf{M}_{\zeta_{\square}}$. Write
$\mathcal{N}_1,\mathcal{N}_2,\mathcal{N}_3$ for the three quantities of a field $A\in\mathfrak{F}$
on $\tilde{\square}^{4}$ which enter the strict inequalities \cite[(3.31)]{B-CMP109}, $\theta$ for
their common threshold there, and $\mathfrak{A}$ for the open set of fields $A\in\mathfrak{F}$ with
$\mathcal{N}_i(A)<\theta$ for $i=1,2,3$. Let $\mathfrak{c}_1\in\left]0,1\right]$,
$\mathfrak{c}_2\in\left]0,1\right[$ and $\mathfrak{c}_3\in\left]0,1\right]$ be the fractions of
$\alpha_2$ used in \cite[text between (1.21) and (1.22)]{B-CMP116}, $\mathfrak{c}_1$ being the one
in \eqref{5.6:eq-auxiliary-constants-b}; only these ranges of $\mathfrak{c}_1$, $\mathfrak{c}_2$,
$\mathfrak{c}_3$ are used, not their numerical values.

The \emph{membership hypothesis} is either one of the following two forms.

\begin{itemize}
\item[(a)] For every $\mathfrak{d}\in\mathcal{D}$, every $s\in\mathbb{D}$ and every
$t_{\square}\in\mathbb{C}$ with $|t_{\square}|\,\nu(\mathfrak{d})\le\mathfrak{c}_3\alpha_2$ one has
$A(s,t_{\square};\mathfrak{d})\in\mathfrak{A}$.

\item[(b)] The following three statements hold.
\begin{itemize}
\item[(i)] Under \eqref{5.6:eq-auxiliary-constants-b}, for every $\mathfrak{d}\in\mathcal{D}$ and
every $s\in\mathbb{D}$,
\begin{equation*}
\mathcal{N}_i\bigl(\mathfrak{t}\,\mathsf{M}_{\zeta_{\ast}}H_k(s,B'(B))\bigr)\le\theta_{c},
\qquad i=1,2,3,
\end{equation*}
where $\theta_{c}$ is the threshold of \cite[(3.31)]{B-CMP109} with $\alpha_2$ replaced by
$\mathfrak{c}_2\alpha_2$, and $\theta-\theta_{c}>0$. This is the statement of
\cite[text between (1.21) and (1.22)]{B-CMP116} that the product with $\mathfrak{t}\zeta_{\ast}$
satisfies \cite[(3.31)]{B-CMP109} with $\mathfrak{c}_2\alpha_2$. If that statement is read as
$\mathcal{N}_i<\mathfrak{c}_2\alpha_2$, everything in what follows holds with $\theta_{c}$
replaced by $\mathfrak{c}_2\alpha_2$.
\item[(ii)] For every $\mathfrak{d}\in\mathcal{D}$ and every $s\in\mathbb{D}$,
\begin{equation*}
\mathcal{N}_i\bigl(\mathsf{M}_{\zeta_{\square}}H_k(s,B'(B))\bigr)\le\nu(\mathfrak{d}),
\qquad i=1,2,3.
\end{equation*}
By \eqref{5.6:eq-auxiliary-constants-c}, $\nu(\mathfrak{d})<4B_0C_1e^{16\kappa_1}\varepsilon_1$.
This is the statement of \cite[text between (1.21) and (1.22)]{B-CMP116} on the norms of the
product with $\zeta_{\square}$.
\item[(iii)] At fixed $\mathbf{U}$ each $\mathcal{N}_i$ is a seminorm on $\mathfrak{F}$: it is
absolutely homogeneous, subadditive and continuous. Here the three quantities are taken, as in
\cite[(1.13), (3.31)]{B-CMP109}, to be the supremum of $|A|$, the supremum of the covariant
difference $|\nabla A|$, and the H\"older seminorm of $\nabla A$. The covariant differences are
$\mathbb{C}$-linear in $A$ at fixed $\mathbf{U}$.
\end{itemize}
\end{itemize}

If, as in the display of \cite[(3.31)--(3.32)]{B-CMP109}, $\theta$ is $2\alpha_2$ plus a
non-negative quantity not involving $\alpha_2$, then under the first reading in (i)
$\theta-\theta_{c}=2(1-\mathfrak{c}_2)\alpha_2>0$, and under the second reading, in which
$\theta_{c}$ is replaced by $\mathfrak{c}_2\alpha_2$,
$\theta-\mathfrak{c}_2\alpha_2\ge(2-\mathfrak{c}_2)\alpha_2>0$.

Form (a) is the membership used by the analytic extension in \cite[text between (1.21) and
(1.22)]{B-CMP116}. There the expression \cite[(1.10)]{B-CMP116} is extended analytically in
$t_{\square}$ for $|t_{\square}|\,4B_0C_1e^{16\kappa_1}g_k|B|\le\mathfrak{c}_3\alpha_2$, by the
analyticity of $\mathcal{T}$ on the space of fields satisfying \cite[(3.31)]{B-CMP109}. The circle
$|t_{\square}|=\mathfrak{c}_3\alpha_2/\nu(\mathfrak{d})$ is then used in \cite[(1.22)]{B-CMP116}.
Both steps require $A(s,t_{\square};\mathfrak{d})\in\mathfrak{A}$ for all $s\in\mathbb{D}$ and all
$|t_{\square}|\le\mathfrak{c}_3\alpha_2/\nu(\mathfrak{d})$. Under form (b) the membership is not
assumed here; it is derived from (i)--(iii) in the proof of the theorem of this section on the
analyticity of the composed term in $(s,t_{\square})$, in the step treating form (b).

With the convention that a positive number divided by $0$ is $+\infty$, the \emph{response radius}
of a datum is
\begin{equation}\label{5.6:eq-response-radius-a}
\begin{gathered}
R_{\square}(\mathfrak{d}):=
\begin{cases}
\mathfrak{c}_3\alpha_2/\nu(\mathfrak{d}) & \text{under form (a)},\\
(\theta-\theta_{c})/\nu(\mathfrak{d}) & \text{under form (b)},
\end{cases}
\qquad R_{\square}(\mathfrak{d})\in\left]0,+\infty\right],\\
R_{\square,\ast}:=\inf_{\mathfrak{d}\in\mathcal{D}}R_{\square}(\mathfrak{d})>0 .
\end{gathered}
\end{equation}
More precisely, by \eqref{5.6:eq-auxiliary-constants-c},
\begin{equation}\label{5.6:eq-response-radius-1}
R_{\square,\ast}\ge
\begin{cases}
\dfrac{\mathfrak{c}_3\alpha_2}{4B_0C_1e^{16\kappa_1}\varepsilon_1} & \text{under form (a)},\\[2ex]
\dfrac{\theta-\theta_{c}}{4B_0C_1e^{16\kappa_1}\varepsilon_1} & \text{under form (b)}.
\end{cases}
\end{equation}
If in addition \eqref{5.6:eq-auxiliary-constants-b} holds, then
\begin{equation}\label{5.6:eq-response-radius-2}
R_{\square,\ast}\ge
\begin{cases}
\dfrac{\mathfrak{c}_3}{\mathfrak{c}_1} & \text{under form (a)},\\[2ex]
\dfrac{\theta-\theta_{c}}{\mathfrak{c}_1\alpha_2} & \text{under form (b)}.
\end{cases}
\end{equation}
Taking $|t_{\square}|$ equal to $R_{\square}(\mathfrak{d})$ under form (a) is the choice made in
\cite[(1.22)]{B-CMP116}.
\end{definition}

\begin{proof}[Proof that $R_{\square,\ast}>0$ and of \eqref{5.6:eq-response-radius-1},
\eqref{5.6:eq-response-radius-2}]
Fix $\mathfrak{d}\in\mathcal{D}$, and write $\mathfrak{m}$ for $\mathfrak{c}_3\alpha_2$ under
form (a) and for $\theta-\theta_{c}$ under form (b). In either case $\mathfrak{m}>0$:
$\mathfrak{c}_3>0$ and $\alpha_2>0$, and $\theta-\theta_{c}>0$ is part of (i).

Suppose first that $B\ne 0$. Then $\nu(\mathfrak{d})=4B_0C_1e^{16\kappa_1}g_k|B|>0$. Condition
\eqref{5.6:eq-auxiliary-constants-c} gives $g_k|B|<\varepsilon_1$, hence
\begin{equation*}
0<\nu(\mathfrak{d})<4B_0C_1e^{16\kappa_1}\varepsilon_1 .
\end{equation*}
Dividing the positive number $\mathfrak{m}$ by these quantities gives
\begin{equation*}
R_{\square}(\mathfrak{d})=\frac{\mathfrak{m}}{\nu(\mathfrak{d})}
>\frac{\mathfrak{m}}{4B_0C_1e^{16\kappa_1}\varepsilon_1}.
\end{equation*}
If $B=0$, then $R_{\square}(\mathfrak{d})=+\infty$ and the same inequality holds. Taking the infimum
over $\mathfrak{d}\in\mathcal{D}$ gives \eqref{5.6:eq-response-radius-1}. Its right-hand sides are
positive because $B_0,C_1,\varepsilon_1>0$, and this proves $R_{\square,\ast}>0$ in
\eqref{5.6:eq-response-radius-a}.

Under \eqref{5.6:eq-auxiliary-constants-b} we have
$0<4B_0C_1e^{16\kappa_1}\varepsilon_1\le\mathfrak{c}_1\alpha_2$, so
\begin{equation*}
\frac{1}{4B_0C_1e^{16\kappa_1}\varepsilon_1}\ge\frac{1}{\mathfrak{c}_1\alpha_2}.
\end{equation*}
Multiplying by $\mathfrak{m}$ turns \eqref{5.6:eq-response-radius-1} into
\eqref{5.6:eq-response-radius-2}, since $\mathfrak{c}_3\alpha_2/(\mathfrak{c}_1\alpha_2)
=\mathfrak{c}_3/\mathfrak{c}_1$.

Finally, consider again a datum with $B=0$; as noted, $\nu(\mathfrak{d})=0$ and
$R_{\square}(\mathfrak{d})=+\infty$. By \cite[(1.20), (1.21)]{B-CMP116},
$H_k(s,B'(0))=0$, so
\begin{equation*}
A(s,t_{\square};\mathfrak{d})
=(\mathfrak{t}\zeta_{\ast}+t_{\square}\zeta_{\square})\cdot 0=0
\quad\text{for all } s \text{ and } t_{\square}.
\end{equation*}
The field $0$ lies in $\mathfrak{A}$ under either form.
\begin{itemize}
\item Under form (a), the condition $|t_{\square}|\cdot 0\le\mathfrak{c}_3\alpha_2$ holds for every
$t_{\square}$, so the hypothesis gives membership at every $t_{\square}$.
\item Under form (b), $\mathcal{N}_i(0)=0$ by absolute homogeneity. Moreover $\theta_{c}\ge 0$,
since by (i) it bounds values of a seminorm, and $\theta-\theta_{c}>0$. Hence
$\mathcal{N}_i(0)=0<\theta$.
\end{itemize}
Consequently the composed term
$F(s,t_{\square};\mathfrak{d})=\mathcal{T}(X;\mathbf{U},\mathbf{J};A(s,t_{\square};\mathfrak{d}))$
of \eqref{5.6:eq-substituted-field-a}, with $X$ the domain fixed there, equals
$\mathcal{T}(X;\mathbf{U},\mathbf{J};0)$ and is constant in $(s,t_{\square})$. Thus
$\partial_{t_{\square}}F\equiv 0$, and any bound of the form
$|\partial_{t_{\square}}F|\le\mathfrak{b}_{\mathcal{T}}/R_{\square}(\mathfrak{d})$ used in what
follows holds at this datum with the convention $\mathfrak{b}_{\mathcal{T}}/(+\infty):=0$, where
$\mathfrak{b}_{\mathcal{T}}$ is the bound of Definition~\ref{5.6:def-term-analyticity}.
\end{proof}

\begin{lemma}[From the closed to the open polydisc of the same radius]
\label{5.6:lem-open-polydisc}
Let $W$ be a finite set and $\rho>0$. Write
\begin{equation*}
\begin{gathered}
\mathbb{D}_{W}(\rho)=\{\zeta\in\mathbb{C}^{W}:|\zeta_w|\le\rho\ \text{for all } w\in W\},\\
P_{W}(\rho)=\{\zeta\in\mathbb{C}^{W}:|\zeta_w|<\rho\ \text{for all } w\in W\}
\end{gathered}
\end{equation*}
for the closed and the open polydisc of radius $\rho$.
Let $E$ be a metric space, let $\|\cdot\|$ be a norm on $\mathbb{C}^m$, and let
$G:\mathbb{D}_{W}(\rho)\times E\to\mathbb{C}^m$ be such that, for each $e\in E$, the function
$G(\cdot,e)$ is analytic on $\mathbb{D}_{W}(\rho)$ in the sense of
Definition~\ref{5.6:def-closed-analyticity}, that is, continuous on $\mathbb{D}_{W}(\rho)$ and
holomorphic at every interior point of $\mathbb{D}_{W}(\rho)$, and such that, for some constant
$C(e)<\infty$, $\|G(\zeta,e)\|\le C(e)$ for all $\zeta\in\mathbb{D}_{W}(\rho)$. Then for each
$e\in E$ the restriction of $G(\cdot,e)$ to the open polydisc $P_{W}(\rho)$ of the same radius
$\rho$ is continuous, holomorphic (in particular separately holomorphic in each variable
$\zeta_w$), and bounded in norm by $C(e)$.
\end{lemma}

\begin{proof}
Fix $e\in E$. For $\zeta\in\mathbb{C}^{W}$ and $r>0$ write
$P_{W}(\zeta;r):=\{\zeta'\in\mathbb{C}^{W}:|\zeta'_w-\zeta_w|<r\ \text{for all } w\in W\}$, the
open polydisc of radius $r$ centred at $\zeta$; it is an open subset of $\mathbb{C}^{W}$.

Every point of $P_{W}(\rho)$ is an interior point of $\mathbb{D}_{W}(\rho)$. Let $\zeta\in P_{W}(\rho)$
and put $r:=\min_{w\in W}(\rho-|\zeta_w|)$. Each number $\rho-|\zeta_w|$ is positive, and the
minimum is taken over the finite set $W$, so $r>0$. For $\zeta'\in P_{W}(\zeta;r)$ and every
$w\in W$,
\begin{equation*}
|\zeta'_w|\le|\zeta_w|+|\zeta'_w-\zeta_w|<|\zeta_w|+r\le|\zeta_w|+(\rho-|\zeta_w|)=\rho ,
\end{equation*}
hence $P_{W}(\zeta;r)\subset\mathbb{D}_{W}(\rho)$, and $\zeta$ is an interior point of
$\mathbb{D}_{W}(\rho)$.

Consequently every $\zeta\in P_{W}(\rho)$ is an interior point of $\mathbb{D}_{W}(\rho)$, and
Definition~\ref{5.6:def-closed-analyticity} gives that $G(\cdot,e)$ is holomorphic at $\zeta$,
in the sense fixed in Notation~\ref{5.6:not-holomorphy-facts} (componentwise for the
$\mathbb{C}^m$-valued function). Thus the restriction to the open set $P_{W}(\rho)$ is
holomorphic at each of its points, and in particular holomorphic in each variable $\zeta_w$
separately, the others being held fixed. Finally $P_{W}(\rho)\subset\mathbb{D}_{W}(\rho)$, so the
restriction of the continuous function $G(\cdot,e)$ is continuous, and the bound
$\|G(\zeta,e)\|\le C(e)$, valid on $\mathbb{D}_{W}(\rho)$, holds in particular on $P_{W}(\rho)$.
\end{proof}

Holomorphy of $G(\cdot,e)$ on an open neighbourhood of $\mathbb{D}_{W}(\rho)$, or its
representation as the sum of an expansion convergent on $\mathbb{D}_{W}(\rho)$, implies
analyticity in the sense of Definition~\ref{5.6:def-closed-analyticity}, so that
Lemma~\ref{5.6:lem-open-polydisc} applies; if, at the radius $\rho=e^{\kappa_1}$, holomorphy is
available only on a smaller open polydisc $P_{W}(e^{\kappa_1-\eta})$ with $\eta>0$, the lemmas
that use Lemma~\ref{5.6:lem-open-polydisc} remain valid, since they are stated for an abstract
radius; they hold with $\kappa_1-\eta$ in place of $\kappa_1$ under one-sided conditions on the
constants.

\begin{theorem}[Analyticity of the cluster-expansion terms on the full decoupling polydisc]
\label{5.6:thm-full-polydisc}
Work at step $k$, at the block $\square$ and at a connected localisation domain $Y_0$ containing
$\tilde{\square}^{4}$, in the expansion \cite[(1.10)]{B-CMP116}. The term functional
$\mathcal{T}$ is as in Definition~\ref{5.6:def-substituted-field}, and so are the following objects:
\begin{itemize}
\item the finite index set $\mathcal{V}$ of the decoupling variables $s(\Delta)$;
\item the set $\mathcal{D}$ of data $\mathfrak{d}=(\mathbf{U},\mathbf{J},B,\mathfrak{t})$;
\item the substituted field $A(s,t_{\square};\mathfrak{d})$ and the composed term
$F(s,t_{\square};\mathfrak{d})$ of \eqref{5.6:eq-substituted-field-a}, with $X$ the domain to
which the term is localised, as there;
\item the number $\nu(\mathfrak{d})$ of \eqref{5.6:eq-substituted-field-b}.
\end{itemize}
In particular, the space $\mathfrak{F}$ of $\mathfrak{g}^{\mathbb{C}}$-valued fields on the finitely
many bonds involved is finite-dimensional. The cut-off functions $\zeta_{\ast}$, $\zeta_{\square}$
are fixed and act on fields by pointwise multiplication. The parameter $\mathfrak{t}\in[0,1]$ is a
fixed real number. Analyticity on a closed set is meant in the sense of
Definition~\ref{5.6:def-closed-analyticity}. Let $\mathfrak{A}\subset\mathfrak{F}$ be the set of
fields $A$ satisfying the bounds \cite[(3.31)]{B-CMP109}. Assume:
\begin{itemize}
\item[(a)] the analyticity and boundedness of the term functional,
Definition~\ref{5.6:def-term-analyticity}, with the bound $\mathfrak{b}_{\mathcal{T}}$ fixed there;
\item[(b)] the analyticity of the localized field in the decoupling parameters,
Proposition~\ref{5.6:prop-localized-field};
\item[(c)] the conditions (C0)--(C5) of Definition~\ref{5.6:def-auxiliary-constants};
\item[(d)] the membership hypothesis of Definition~\ref{5.6:def-response-radius}, in its first form
or in its second form, with $R_{\square}(\mathfrak{d})$ and $R_{\square,\ast}$ as in
\eqref{5.6:eq-response-radius-a} for that form.
\end{itemize}
For $\mathfrak{d}\in\mathcal{D}$ put
\begin{equation}\label{5.6:eq-full-polydisc-a}
\Omega_{\square}(\mathfrak{d}) := P_{\mathcal{V}}(e^{\kappa_1})\times
\{t_{\square}\in\mathbb{C} : |t_{\square}|<R_{\square}(\mathfrak{d})\}.
\end{equation}
Here $P_{\mathcal{V}}(e^{\kappa_1})$ is the open polydisc in $\mathbb{C}^{\mathcal{V}}$ consisting of
all $s$ with $|s(\Delta)|<e^{\kappa_1}$ for all $\Delta\in\mathcal{V}$. Its radius is exactly
$e^{\kappa_1}$ in all the variables $s(\Delta)$, $\Delta\in\mathcal{V}$. We use the convention
$\mathfrak{b}_{\mathcal{T}}/(+\infty):=0$. Then for every datum
$\mathfrak{d}=(\mathbf{U},\mathbf{J},B,\mathfrak{t})\in\mathcal{D}$ the following hold.
\begin{itemize}
\item[(i)] $A(s,t_{\square};\mathfrak{d})\in\mathfrak{A}$ for all $(s,t_{\square})\in
\Omega_{\square}(\mathfrak{d})$. The function
$F(\cdot;\mathfrak{d})\colon\Omega_{\square}(\mathfrak{d})\to\mathbb{C}$, given by
$F(s,t_{\square};\mathfrak{d})=\mathcal{T}(X;\mathbf{U},\mathbf{J};A(s,t_{\square};\mathfrak{d}))$,
has the following properties:
\begin{itemize}
\item it is continuous;
\item it is separately holomorphic in each variable $s(\Delta)$, $\Delta\in\mathcal{V}$, and in
$t_{\square}$;
\item hence it is jointly holomorphic on $\Omega_{\square}(\mathfrak{d})$;
\item $|F(s,t_{\square};\mathfrak{d})|\le\mathfrak{b}_{\mathcal{T}}$ there.
\end{itemize}
\item[(ii)] Let $t_{\square 0}\in\mathbb{C}$ with $|t_{\square 0}|<R_{\square,\ast}$. Then
$\sigma\mapsto F(\sigma,t_{\square 0};\mathfrak{d})$ is continuous and separately holomorphic on
$P_{\mathcal{V}}(e^{\kappa_1})$, with modulus at most $\mathfrak{b}_{\mathcal{T}}$. Hence the
holomorphy hypothesis of the weight lemma holds with the function of that lemma taken to be
$(\sigma;\mathfrak{d})\mapsto F(\sigma,t_{\square 0};\mathfrak{d})$, and with the following choices:
\begin{itemize}
\item the set of variables $\mathcal{V}(t):=\mathcal{V}$;
\item the data set $\mathcal{D}_{b}:=\mathcal{D}$;
\item the bound $\mathfrak{b}:=\mathfrak{b}_{\mathcal{T}}$.
\end{itemize}
That hypothesis asks for continuity, separate holomorphy and the bound by $\mathfrak{b}$ on the open
polydisc $P_{\mathcal{V}}(e^{\kappa_1})$, for each datum separately, together with $\kappa_1\ge 1$.
The condition $\kappa_1\ge 1$ is a one-sided floor on $\kappa_1$ comprised in (C0), by which
$\kappa_1$ is a sufficiently big positive number.
\item[(iii)] Put
\begin{equation*}
F_{t}(\sigma;\mathfrak{d}) := \partial_{t_{\square}}F(\sigma,t_{\square};\mathfrak{d})
\big|_{t_{\square}=0}.
\end{equation*}
This is the derivative with respect to $t_{\square}$ at $t_{\square}=0$ taken in
\cite[text before (1.23)]{B-CMP116}. It is
continuous and separately, hence jointly, holomorphic on $P_{\mathcal{V}}(e^{\kappa_1})$, and
\begin{equation*}
|F_{t}(\sigma;\mathfrak{d})| \le \mathfrak{b}_{\mathcal{T}}/R_{\square}(\mathfrak{d})
\le \mathfrak{b}_{\mathcal{T}}/R_{\square,\ast}
\qquad\text{for all }\sigma\in P_{\mathcal{V}}(e^{\kappa_1}).
\end{equation*}
Hence the holomorphy hypothesis of the weight lemma holds for this $F_{t}$, with
$\mathcal{V}(t):=\mathcal{V}$, $\mathcal{D}_{b}:=\mathcal{D}$ and
$\mathfrak{b}:=\mathfrak{b}_{\mathcal{T}}/R_{\square,\ast}$. Datum by datum the bound is
$\mathfrak{b}(\mathfrak{d})=\mathfrak{b}_{\mathcal{T}}/R_{\square}(\mathfrak{d})$. The bound in that
hypothesis is one constant for all data, and this is why the infimum $R_{\square,\ast}$ enters.
\item[(iv)] No radius is lost. The polydisc $P_{\mathcal{V}}(e^{\kappa_1})$ is, by
Lemma~\ref{5.6:lem-open-polydisc}, the interior of the closed polydisc
$\mathbb{D}_{\mathcal{V}}(e^{\kappa_1})$, where $|s(\Delta)|\le e^{\kappa_1}$, on which
\cite[(1.21)]{B-CMP116} is stated.
In \cite[text before (1.23)]{B-CMP116} the derivatives are represented by Cauchy integrals on the
boundary circles $|\sigma(\Delta)|=e^{\kappa_1}$ and on the circle $|t_{\square}|=R_{\square}(\mathfrak{d})$
of \cite[(1.22)]{B-CMP116}. These representations use, beyond \textup{(i)}--\textup{(iii)}, the
continuity of $F$ on the closed sets. That continuity is provided by the continuity on closed sets in
Definition~\ref{5.6:def-closed-analyticity}, together with the first form of the membership
hypothesis. The holomorphy hypothesis of the weight lemma uses only the open polydisc.
\end{itemize}
\end{theorem}

\begin{proof}
Fix $\mathfrak{d}=(\mathbf{U},\mathbf{J},B,\mathfrak{t})\in\mathcal{D}$. At the fixed
$(\mathbf{U},\mathbf{J})\in\mathcal{D}^{UJ}$ write $H(s):=H_k(s,B'(B))\in\mathfrak{F}$. Write
$\mathsf{M}_{\zeta_{\ast}},\mathsf{M}_{\zeta_{\square}}\colon\mathfrak{F}\to\mathfrak{F}$ for the
operators of pointwise multiplication by $\zeta_{\ast}$ and by $\zeta_{\square}$.

\emph{First step: the setting.} The space $\mathfrak{F}$ consists of
$\mathfrak{g}^{\mathbb{C}}$-valued functions on a finite set of bonds. By
Notation~\ref{5.6:not-holomorphy-facts} it is therefore a finite-dimensional complex vector space,
$\mathfrak{F}\cong\mathbb{C}^{n}$. The variables are $s=(s(\Delta))_{\Delta\in\mathcal{V}}\in
\mathbb{C}^{\mathcal{V}}$, with $\mathcal{V}$ finite, and $t_{\square}\in\mathbb{C}$. The set
$P_{\mathcal{V}}(e^{\kappa_1})$ is open in $\mathbb{C}^{\mathcal{V}}$, and
$\Omega_{\square}(\mathfrak{d})$ is open in $\mathbb{C}^{\mathcal{V}}\times\mathbb{C}$. The set
$\mathfrak{A}$ is open in $\mathfrak{F}$ by clause (i) of Definition~\ref{5.6:def-term-analyticity}.
All maps below are either between open subsets of finite-dimensional complex vector spaces or defined
on closed polydiscs. For these, ``holomorphic'' is meant in the sense of
Notation~\ref{5.6:not-holomorphy-facts}; by the statement recorded there, a continuous function on an
open set of $\mathbb{C}^{n}$ that is holomorphic in each variable separately is holomorphic.

\emph{Second step: the localized field on the open polydisc of full radius.} By
Proposition~\ref{5.6:prop-localized-field}, the map $s\mapsto H(s)$ is analytic on
$\mathbb{D}_{\mathcal{V}}(e^{\kappa_1})$ in the sense of Definition~\ref{5.6:def-closed-analyticity}.
It obeys \eqref{5.6:eq-localized-field-a}, so that in all admissible norms
\begin{equation*}
\sup_{s\in\mathbb{D}_{\mathcal{V}}(e^{\kappa_1})}\|H(s)\| < 4B_0C_1e^{16\kappa_1}\varepsilon_1 .
\end{equation*}
We apply Lemma~\ref{5.6:lem-open-polydisc} with $W:=\mathcal{V}$, $\rho:=e^{\kappa_1}$,
$E:=\mathcal{D}$ and $G:=H_k(\cdot,B'(\cdot))$. Its hypotheses are exactly the analyticity and the
bound just recalled. It gives that $s\mapsto H(s)$ is continuous on $\mathbb{D}_{\mathcal{V}}(e^{\kappa_1})$
and holomorphic on $P_{\mathcal{V}}(e^{\kappa_1})$, with the same bound on
$P_{\mathcal{V}}(e^{\kappa_1})$. This holds at every interior point of the domain of
\cite[(1.21)]{B-CMP116}, with radius exactly $e^{\kappa_1}$. Of the norms in which
\eqref{5.6:eq-localized-field-a} is stated, nothing is used beyond their being norms on $\mathfrak{F}$.
By Notation~\ref{5.6:not-holomorphy-facts} all such norms are equivalent and continuous.

\emph{Third step: the substituted field is jointly holomorphic in $(s,t_{\square})$.} The cut-offs act
by pointwise multiplication, and $\mathfrak{t}$ is fixed (Definition~\ref{5.6:def-substituted-field}).
By \eqref{5.6:eq-substituted-field-a} we therefore have
\begin{equation*}
A(s,t_{\square};\mathfrak{d}) = \mathfrak{t}\,\mathsf{M}_{\zeta_{\ast}}H(s)
+ t_{\square}\,\mathsf{M}_{\zeta_{\square}}H(s).
\end{equation*}
Here $\mathsf{M}_{\zeta_{\ast}}$ and $\mathsf{M}_{\zeta_{\square}}$ are fixed $\mathbb{C}$-linear maps
of $\mathfrak{F}$. They are continuous, because linear maps of a finite-dimensional space are
(Notation~\ref{5.6:not-holomorphy-facts}). The number $\mathfrak{t}$ is a fixed real number.

Fix $t_{\square}$. Then $s\mapsto A(s,t_{\square};\mathfrak{d})$ is a sum of two maps of the form
(continuous linear map)$\circ$(holomorphic map). It is therefore holomorphic on
$P_{\mathcal{V}}(e^{\kappa_1})$, by the second step and the composition rule of
Notation~\ref{5.6:not-holomorphy-facts}.

Fix $s\in\mathbb{D}_{\mathcal{V}}(e^{\kappa_1})$. Then $t_{\square}\mapsto A(s,t_{\square};\mathfrak{d})$
is affine on $\mathbb{C}$, hence holomorphic.

The map $(s,t_{\square})\mapsto A(s,t_{\square};\mathfrak{d})$ is continuous on
$\mathbb{D}_{\mathcal{V}}(e^{\kappa_1})\times\mathbb{C}$. Indeed, $H$ is continuous on
$\mathbb{D}_{\mathcal{V}}(e^{\kappa_1})$, the two linear maps are continuous, and multiplication by the
scalar $t_{\square}$ is jointly continuous.

Hence $A(\cdot,\cdot;\mathfrak{d})$ is continuous and separately holomorphic on
$P_{\mathcal{V}}(e^{\kappa_1})\times\mathbb{C}$, with values in $\mathfrak{F}$. By
Notation~\ref{5.6:not-holomorphy-facts} it is jointly holomorphic there.

\emph{Fourth step: membership on $\Omega_{\square}(\mathfrak{d})$.} We treat the two forms of the
membership hypothesis in turn.

Under the first form of the membership hypothesis of Definition~\ref{5.6:def-response-radius}, we have
$A(s,t_{\square};\mathfrak{d})\in\mathfrak{A}$ for all $s\in\mathbb{D}_{\mathcal{V}}(e^{\kappa_1})$
and all $t_{\square}$ with $|t_{\square}|\nu(\mathfrak{d})\le\mathfrak{c}_3\alpha_2$. This means
\begin{equation*}
|t_{\square}|\le \mathfrak{c}_3\alpha_2/\nu(\mathfrak{d}) = R_{\square}(\mathfrak{d}).
\end{equation*}
Since $\mathbb{D}_{\mathcal{V}}(e^{\kappa_1})\supset P_{\mathcal{V}}(e^{\kappa_1})$, this holds a
fortiori on $\Omega_{\square}(\mathfrak{d})$. Only the open disc
$|t_{\square}|<R_{\square}(\mathfrak{d})$ is used; the boundary circle
$|t_{\square}|=R_{\square}(\mathfrak{d})$ of \cite[(1.22)]{B-CMP116} plays no role here.

Under the second form of the membership hypothesis, $\mathcal{N}_1,\mathcal{N}_2,\mathcal{N}_3$ denote
the three quantities of \cite[(3.31)]{B-CMP109}, which under that form are seminorms on
$\mathfrak{F}$ at fixed $\mathbf{U}$ (Definition~\ref{5.6:def-response-radius}): absolutely
homogeneous, subadditive and continuous. The threshold $\theta$ and the reduced threshold $\theta_{c}$ satisfy
$\theta-\theta_{c}>0$, and $R_{\square}(\mathfrak{d})=(\theta-\theta_{c})/\nu(\mathfrak{d})$. That form
supplies, for every $s\in\mathbb{D}_{\mathcal{V}}(e^{\kappa_1})$ and $i=1,2,3$, the two bounds
\begin{equation*}
\mathcal{N}_i(\mathfrak{t}\,\mathsf{M}_{\zeta_{\ast}}H(s))\le\theta_{c},
\qquad
\mathcal{N}_i(\mathsf{M}_{\zeta_{\square}}H(s))\le\nu(\mathfrak{d}).
\end{equation*}
Let $s\in\mathbb{D}_{\mathcal{V}}(e^{\kappa_1})$ and $|t_{\square}|<R_{\square}(\mathfrak{d})$. Then
$|t_{\square}|\nu(\mathfrak{d})<\theta-\theta_{c}$. When $\nu(\mathfrak{d})>0$ this is the
hypothesis on $t_{\square}$ multiplied by $\nu(\mathfrak{d})$. When $\nu(\mathfrak{d})=0$ the left side
vanishes, and $\theta-\theta_{c}>0$. Subadditivity and absolute homogeneity of $\mathcal{N}_i$ then
give
\begin{align*}
\mathcal{N}_i(A(s,t_{\square};\mathfrak{d}))
&\le \mathcal{N}_i(\mathfrak{t}\,\mathsf{M}_{\zeta_{\ast}}H(s))
+ |t_{\square}|\,\mathcal{N}_i(\mathsf{M}_{\zeta_{\square}}H(s))\\
&\le \theta_{c} + |t_{\square}|\,\nu(\mathfrak{d})
< \theta_{c}+(\theta-\theta_{c}) = \theta .
\end{align*}
So $A(s,t_{\square};\mathfrak{d})$ satisfies the bounds \cite[(3.31)]{B-CMP109}, that is,
$A(s,t_{\square};\mathfrak{d})\in\mathfrak{A}$.

In both forms the radius $R_{\square}(\mathfrak{d})$ is of order $1/(e^{16\kappa_1}g_k|B|)$. The
conditions (C3) and (C4) of Definition~\ref{5.6:def-auxiliary-constants} bound it below by
$R_{\square,\ast}>0$, as in \eqref{5.6:eq-response-radius-a}. These conditions are a smallness of the
coupling imposed after $\kappa_1$ is chosen; in particular (C3) is
\eqref{5.6:eq-auxiliary-constants-b}. They are one-sided.

\emph{Fifth step: $F$ is continuous and separately holomorphic on $\Omega_{\square}(\mathfrak{d})$,
hence jointly holomorphic.} The composite
$F=\mathcal{T}(X;\mathbf{U},\mathbf{J};\cdot)\circ A(\cdot,\cdot;\mathfrak{d})$ is defined on
$\Omega_{\square}(\mathfrak{d})$ by the fourth step.

Continuity. By clause (i) of Definition~\ref{5.6:def-term-analyticity}, $\mathcal{T}$ is holomorphic
on $\mathfrak{A}$, hence continuous there (Notation~\ref{5.6:not-holomorphy-facts}). The map $A$ is
continuous by the third step. So $F$ is continuous.

Separate holomorphy. Fix all variables but one, and call the free one
$w\in\{s(\Delta):\Delta\in\mathcal{V}\}\cup\{t_{\square}\}$. It ranges over an open disc: either
$|s(\Delta)|<e^{\kappa_1}$ or $|t_{\square}|<R_{\square}(\mathfrak{d})$. On that disc, the curve
\begin{equation*}
w\mapsto\gamma(w):=A(\dots,w,\dots;\mathfrak{d})
\end{equation*}
is a complex-differentiable $\mathfrak{F}$-valued curve by the third step. It takes values in
$\mathfrak{A}$ by the fourth step. The functional $\mathcal{T}(X;\mathbf{U},\mathbf{J};\cdot)$ is
complex-differentiable on the open set $\mathfrak{A}$. By the one-variable chain rule of
Notation~\ref{5.6:not-holomorphy-facts}, the function $w\mapsto\mathcal{T}(X;\mathbf{U},\mathbf{J};\gamma(w))$
is complex-differentiable, with derivative the complex derivative of
$\mathcal{T}(X;\mathbf{U},\mathbf{J};\cdot)$ at $\gamma(w)$ applied to $\gamma'(w)$, in the notation of
Notation~\ref{5.6:not-holomorphy-facts}.

Hence $F$ is continuous and separately holomorphic on $\Omega_{\square}(\mathfrak{d})$. By
Notation~\ref{5.6:not-holomorphy-facts} it is therefore jointly holomorphic there. Equivalently, $F$ is
the composite of two holomorphic maps, $A(\cdot,\cdot;\mathfrak{d})\colon\Omega_{\square}(\mathfrak{d})\to\mathfrak{A}$
and $\mathcal{T}(X;\mathbf{U},\mathbf{J};\cdot)\colon\mathfrak{A}\to\mathbb{C}$, and the composition rule
of Notation~\ref{5.6:not-holomorphy-facts} applies.

\emph{Sixth step: the bound, and clauses \textup{(i)} and \textup{(ii)}.} Let
$(s,t_{\square})\in\Omega_{\square}(\mathfrak{d})$. Then $A(s,t_{\square};\mathfrak{d})\in\mathfrak{A}$
by the fourth step, and $(\mathbf{U},\mathbf{J})\in\mathcal{D}^{UJ}$. Clause (ii) of
Definition~\ref{5.6:def-term-analyticity} therefore gives
\begin{equation*}
|F(s,t_{\square};\mathfrak{d})|
= |\mathcal{T}(X;\mathbf{U},\mathbf{J};A(s,t_{\square};\mathfrak{d}))|
\le \sup_{A\in\mathfrak{A}}|\mathcal{T}(X;\mathbf{U},\mathbf{J};A)|
\le \mathfrak{b}_{\mathcal{T}} .
\end{equation*}
Together with the fourth and fifth steps, this proves clause (i).

For clause (ii), let $|t_{\square 0}|<R_{\square,\ast}$. Then $|t_{\square 0}|<R_{\square}(\mathfrak{d})$
for every $\mathfrak{d}\in\mathcal{D}$. So $\sigma\mapsto F(\sigma,t_{\square 0};\mathfrak{d})$ is the
restriction of $F(\cdot;\mathfrak{d})$ to the coordinate slice
\begin{equation*}
\{t_{\square}=t_{\square 0}\}\cap\Omega_{\square}(\mathfrak{d})
= P_{\mathcal{V}}(e^{\kappa_1})\times\{t_{\square 0}\}.
\end{equation*}
On this slice it is continuous, separately holomorphic and bounded in modulus by
$\mathfrak{b}_{\mathcal{T}}$, by clause (i). These are the three requirements of the holomorphy
hypothesis of the weight lemma, at radius $e^{\kappa_1}$ and with $\mathfrak{b}=\mathfrak{b}_{\mathcal{T}}$.
They hold for each $\mathfrak{d}$ separately, and that hypothesis imposes nothing on the dependence on
$\mathfrak{d}$.

\emph{Seventh step: the $t_{\square}$-derivative, clause \textup{(iii)}.} Fix
$\sigma\in P_{\mathcal{V}}(e^{\kappa_1})$ and put $\varphi_{\sigma}(\tau):=F(\sigma,\tau;\mathfrak{d})$.
This function is holomorphic on $\{|\tau|<R_{\square}(\mathfrak{d})\}$, because it is the restriction of
$F$ to a complex line (fifth step and Notation~\ref{5.6:not-holomorphy-facts}). Let
$0<R'<R_{\square}(\mathfrak{d})$. By the Cauchy integral formula of
Notation~\ref{5.6:not-holomorphy-facts},
\begin{equation*}
F_{t}(\sigma;\mathfrak{d}) = \varphi_{\sigma}'(0)
= \frac{1}{2\pi i}\oint_{|\tau|=R'}F(\sigma,\tau;\mathfrak{d})\,\tau^{-2}\,d\tau .
\end{equation*}
The circle has length $2\pi R'$, the factor $\tau^{-2}$ has modulus $R'^{-2}$ on it, and
$|\varphi_{\sigma}|\le\mathfrak{b}_{\mathcal{T}}$ there by the sixth step. Hence
\begin{equation*}
|F_{t}(\sigma;\mathfrak{d})| \le \frac{1}{2\pi}\cdot 2\pi R'\cdot\frac{1}{R'^{2}}
\max_{|\tau|=R'}|\varphi_{\sigma}(\tau)|
\le \frac{\mathfrak{b}_{\mathcal{T}}}{R'} .
\end{equation*}
The left side does not depend on $R'$. Letting $R'\uparrow R_{\square}(\mathfrak{d})$ gives
\begin{equation*}
|F_{t}(\sigma;\mathfrak{d})|\le\mathfrak{b}_{\mathcal{T}}/R_{\square}(\mathfrak{d})
\le\mathfrak{b}_{\mathcal{T}}/R_{\square,\ast}.
\end{equation*}
If $R_{\square}(\mathfrak{d})=+\infty$, letting $R'\to\infty$ gives $F_{t}(\sigma;\mathfrak{d})=0$, in
accordance with the convention $\mathfrak{b}_{\mathcal{T}}/(+\infty)=0$.

Holomorphy and continuity in $\sigma$. Fix $R'<R_{\square}(\mathfrak{d})$ and consider the integrand
$h(\sigma,\tau):=F(\sigma,\tau;\mathfrak{d})\tau^{-2}$. It is continuous on
$P_{\mathcal{V}}(e^{\kappa_1})\times\{|\tau|=R'\}$, and holomorphic in $\sigma$ for each $\tau$, by the
fifth step. By the statement on integrals depending holomorphically on a parameter in
Notation~\ref{5.6:not-holomorphy-facts}, the function $\sigma\mapsto F_{t}(\sigma;\mathfrak{d})$ is
continuous and separately holomorphic on $P_{\mathcal{V}}(e^{\kappa_1})$, hence holomorphic there
(Notation~\ref{5.6:not-holomorphy-facts}).

This is the holomorphy hypothesis of the weight lemma for the function of clause (iii), with
$\mathfrak{b}=\mathfrak{b}_{\mathcal{T}}/R_{\square,\ast}$. This proves clause (iii).

\emph{Eighth step: exactness of the radius, clause \textup{(iv)}.} The polydisc in the variables $s$
in (i)--(iii) is $P_{\mathcal{V}}(e^{\kappa_1})$. By Lemma~\ref{5.6:lem-open-polydisc} it is the
interior of $\mathbb{D}_{\mathcal{V}}(e^{\kappa_1})$, the closed polydisc $|s(\Delta)|\le e^{\kappa_1}$
of \cite[(1.21)]{B-CMP116}. Its radius is exactly $e^{\kappa_1}$. This is the radius below which the
lemma inside the proof of the weight lemma takes its Cauchy integrals: on tori whose radii are
strictly smaller than $e^{\kappa_1}$, after which the radii are let increase to $e^{\kappa_1}$.

In \cite[text before (1.23)]{B-CMP116} the integrations are on the circles
$|\sigma(\Delta)|=e^{\kappa_1}$ and on the circle $|t_{\square}|=R_{\square}(\mathfrak{d})$ of
\cite[(1.22)]{B-CMP116}. This requires $F$ to be continuous on the closed product and holomorphic
inside. Three facts provide this:
\begin{itemize}
\item the continuity of $H_k$ on $\mathbb{D}_{\mathcal{V}}(e^{\kappa_1})$, by
Definition~\ref{5.6:def-closed-analyticity} and Proposition~\ref{5.6:prop-localized-field};
\item the membership $A(s,t_{\square};\mathfrak{d})\in\mathfrak{A}$ for
$s\in\mathbb{D}_{\mathcal{V}}(e^{\kappa_1})$ and $|t_{\square}|\le R_{\square}(\mathfrak{d})$, by the
first form of the membership hypothesis;
\item the Cauchy integral formula for functions continuous on the closed disc and holomorphic inside,
Notation~\ref{5.6:not-holomorphy-facts}.
\end{itemize}
The holomorphy hypothesis of the weight lemma requires less, namely the open polydisc only, and so
(ii) and (iii) meet it with room to spare.
\end{proof}

\begin{remark}[A smaller decoupling radius at one-sided cost]\label{5.6:rem-smaller-radius}
This remark is not used in the sequel. Let $0\le\epsilon<\kappa_1-1$ and put $\kappa_1':=\kappa_1-\epsilon$.
Suppose that, in the setting of Theorem~\ref{5.6:thm-full-polydisc}, the holomorphy of the function
$s\mapsto H(s)$ of the proof of Theorem~\ref{5.6:thm-full-polydisc} is known only on the open
polydisc $P_{\mathcal{V}}(e^{\kappa_1'})$.
Then the following hold.
\begin{enumerate}
\item[(a)] The proof of Theorem~\ref{5.6:thm-full-polydisc}, with $P:=P_{\mathcal{V}}(e^{\kappa_1'})$,
yields the hypothesis of the weight lemma with $\kappa_1'$ in place of $\kappa_1$: for every datum
$\mathfrak{d}$ the function $F_{t}(\cdot\,;\mathfrak{d})$ is continuous, separately holomorphic and
bounded by $\mathfrak{b}$ on $P_{\mathcal{V}}(e^{\kappa_1'})$.
\item[(b)] The weight lemma applies with radius $e^{\kappa_1'}$. Its per-variable Cauchy factor
satisfies
\begin{equation*}
r(\kappa_1')\le e^{-(\kappa_1'-1)},
\end{equation*}
and each of its lower bounds on $\kappa_1$ is replaced by the same lower bound increased by $\epsilon$.
\end{enumerate}
These are again lower bounds on $\kappa_1$ only, with $\kappa_1$ chosen after $\kappa$; no upper bound is
imposed on any constant. For $\epsilon=0$ part (a) is Theorem~\ref{5.6:thm-full-polydisc} and part (b)
is the weight lemma as stated.
\end{remark}

\begin{proof}
(a) The steps of the proof of Theorem~\ref{5.6:thm-full-polydisc} that start from the holomorphy
of $H$ are carried out without change with $P:=P_{\mathcal{V}}(e^{\kappa_1'})$ in place of
$P_{\mathcal{V}}(e^{\kappa_1})$. They give the hypothesis of the weight lemma with $\kappa_1'$ in
place of $\kappa_1$: for every datum $\mathfrak{d}$ the function $F_{t}(\cdot\,;\mathfrak{d})$ is
continuous and separately holomorphic on the open polydisc $P_{\mathcal{V}}(e^{\kappa_1'})$, with the
sup bound $\mathfrak{b}$ there.

(b) The weight lemma is stated on an open polydisc of abstract radius $R=e^{\kappa_1}>1$. Its
estimate of the per-variable Cauchy factor
\begin{equation*}
r(\kappa_1)=e^{\kappa_1}(e^{\kappa_1}-1)^{-2}
\end{equation*}
is stated for this abstract radius. So are its representation of derivatives by Cauchy integrals
over tori whose radii lie strictly between $1$ and $R$, and its remaining statements. Its conditions
on $\kappa_1$ are one-sided: each is a lower bound $\kappa_1\ge c$ whose right-hand side does not
involve $\kappa_1$; one of them is $\kappa_1\ge1$, and the others have right-hand sides formed from
constants of the weight lemma.

We now replace $\kappa_1$ by $\kappa_1'$ throughout. Since $\epsilon<\kappa_1-1$, we have $\kappa_1'>1$.
The estimate $r(\kappa_1)\le e^{-(\kappa_1-1)}$, valid for $\kappa_1\ge1$, therefore gives
$r(\kappa_1')\le e^{-(\kappa_1'-1)}$.

Each condition of the form $\kappa_1'\ge c$ is, by the definition $\kappa_1'=\kappa_1-\epsilon$, the same
as $\kappa_1\ge c+\epsilon$. This is again a lower bound on $\kappa_1$ alone, to be met when $\kappa_1$ is
chosen after $\kappa$, and no upper bound arises.

Finally, $\epsilon=0$ gives $\kappa_1'=\kappa_1$; part (a) is then Theorem~\ref{5.6:thm-full-polydisc}
and every statement of part (b) is the weight lemma as stated.
\end{proof}

\begin{definition}[A domain term as one piece of the weight lemma's form]
\label{5.6:def-domain-term}
We use the datum $\mathfrak{d}$, the response parameter $t_{\square}$, the block $\square$ and
the composed term $F$ of Definition~\ref{5.6:def-substituted-field}, and we write
$\tilde{\square}^{4}$ for the fixed cube-neighbourhood of the block $\square$, as in
\cite[text before (1.10)]{B-CMP116}.

In \cite[text before (1.10)]{B-CMP116} the following is stated. In a term of the sum
considered there, the region carrying the parameters $s$ may have several components; the
derivatives with respect to the parameters $s$ restricted to a component not containing
$\tilde{\square}^{4}$ make the term vanish. Hence the
sum of \cite[(1.10)]{B-CMP116} runs over connected domains $Y_{0}$ containing
$\tilde{\square}^{4}$, with the parameters set to $s=0$ on $Y_{0}^{c}$, and the remaining
parameters are denoted $s(Y_{0})$.

Fix a connected domain $Y_{0}\supseteq\tilde{\square}^{4}$, and let $\mathcal{V}$ be the set
of $M$-cubes of $Y_{0}\setminus\tilde{\square}^{4}$.
In \cite[text before (1.23)]{B-CMP116} this expansion is differentiated with respect to
$t_{\square}$ at $t_{\square}=0$, and all derivatives are represented by the Cauchy
formula, with the $\sigma(\Delta)$-integrations over the circles $|\sigma(\Delta)|=e^{\kappa_1}$.
There is one such circle for each $\Delta\in\mathcal{V}$. Each variable
$\sigma(\Delta)=s(\Delta)$, $\Delta\in\mathcal{V}$, is therefore differentiated exactly once.

Let $t_{Y_{0}}$ denote the term of
\cite[(1.10)]{B-CMP116} corresponding to the domain $Y_{0}$ (the letter $t_{Y_{0}}$ names a
term and is unrelated to the complex parameter $t_{\square}$), and put
$\mathcal{V}(t_{Y_{0}}):=\mathcal{V}$. We regard $t_{Y_{0}}$ as one piece of the weight
lemma's form, with the following data.
\begin{enumerate}
\item[(a)] The set of differentiated variables is $W:=\mathcal{V}$, each differentiated once
as noted above, and the integration is against Lebesgue measure on $[0,1]^{\mathcal{V}}$.
\item[(b)] The function of the piece is the first $t_{\square}$-derivative of the composed term
at $t_{\square}=0$, as in \cite[text before (1.23)]{B-CMP116}:
\begin{equation}\label{5.6:eq-domain-term-1}
F_{t_{Y_{0}}}(\sigma;\mathfrak{d})
:=\partial_{t_{\square}}F(\sigma,t_{\square};\mathfrak{d})\big|_{t_{\square}=0},
\qquad \sigma=(\sigma(\Delta))_{\Delta\in\mathcal{V}}.
\end{equation}
\item[(c)] In the chain notation of the weight lemma's form:
\begin{itemize}
\item the chain of the piece has one link, $\lambda=1$;
\item its final domain is $Y_{0}$, and $Y_{0}\supseteq Y_{1}$;
\item the number of its variables is $n=\#\mathcal{V}$;
\item $\#\mathcal{P}(t_{Y_{0}};Y_{0})=1$.
\end{itemize}
\end{enumerate}
Only the single term $t_{Y_{0}}$, for a fixed $Y_{0}$, is so regarded; the sum over $Y_{0}$
in \cite[(1.10)]{B-CMP116} is not.
\end{definition}

\begin{corollary}[The weight lemma's hypothesis and bound for the cluster-expansion terms]
\label{5.6:cor-weight-hypothesis}
Assume the following:
\begin{itemize}
\item[(1)] the setting of Theorem~\ref{5.6:thm-full-polydisc}, and the conclusions of that
theorem (analyticity of the cluster-expansion terms on the full decoupling polydisc);
\item[(2)] the analyticity and boundedness of the term functional $\mathcal{T}$ of
Definition~\ref{5.6:def-term-analyticity};
\item[(3)] the analyticity of the localized field in the decoupling parameters,
Proposition~\ref{5.6:prop-localized-field};
\item[(4)] the membership condition of Definition~\ref{5.6:def-response-radius}, in either of
its two forms, with the response radius $R_{\square}(\mathfrak{d})$ and its infimum
$R_{\square,\ast}>0$ as in~\eqref{5.6:eq-response-radius-a};
\item[(5)] the conditions on the auxiliary constants of
Definition~\ref{5.6:def-auxiliary-constants};
\item[(6)] the identification of a domain term with one piece of the weight lemma's form,
Definition~\ref{5.6:def-domain-term};
\item[(7)] clauses (a)--(e) of the weight lemma, together with the bound on the Cauchy factor
stated with it,
\begin{equation*}
r(\kappa_1) := e^{\kappa_1}(e^{\kappa_1}-1)^{-2} \le e^{-(\kappa_1-1)}
\qquad\text{for } \kappa_1\ge 1;
\end{equation*}
\item[(8)] $\kappa_1 \ge 1$.
\end{itemize}
Let $\square$ be a block and $Y_0$ a connected localisation domain containing
$\tilde{\square}^{4}$, and let $t := t_{Y_0}$ be the term of \cite[(1.10)]{B-CMP116}
corresponding to $Y_0$. Then the analyticity hypothesis of the weight lemma holds for the
term $t$ with index set $\mathcal{V}(t) := \mathcal{V}$, the cubes of
$Y_0\setminus\tilde{\square}^{4}$, with exempt set the cubes of $\tilde{\square}^{4}$, with
$\mathcal{D}_{b} := \mathcal{D}$, and with
\begin{equation*}
F_{t}(\sigma;\mathfrak{d}) := \partial_{t_{\square}} F(\sigma,t_{\square};\mathfrak{d})
\big|_{t_{\square}=0},
\qquad
\mathfrak{b} := \mathfrak{b}_{\mathcal{T}}/R_{\square,\ast},
\end{equation*}
for $\sigma\in P=P_{\mathcal{V}}(e^{\kappa_1})$ and $\mathfrak{d}\in\mathcal{D}$, where $F$ is
the composed term of~\eqref{5.6:eq-substituted-field-a}. The same hypothesis also holds with
$F_{t} := F(\cdot,t_{\square 0};\cdot)$ for any fixed $t_{\square 0}\in\mathbb{C}$ with
$|t_{\square 0}|<R_{\square,\ast}$ and $\mathfrak{b} := \mathfrak{b}_{\mathcal{T}}$.
Consequently, for the first choice of $F_t$, denote by $p_{Y_0}(\mathfrak{d})$ the piece of
$t_{Y_0}$ in the sense of Definition~\ref{5.6:def-domain-term}, that is, the $Y_0$-term read
as the single piece of the weight lemma's form in which every variable $s(\Delta)$,
$\Delta\in\mathcal{V}$, is differentiated once; then
\begin{equation}\label{5.6:eq-weight-hypothesis-1}
\sup_{\mathfrak{d}\in\mathcal{D}} \bigl|p_{Y_0}(\mathfrak{d})\bigr|
\le \frac{\mathfrak{b}_{\mathcal{T}}}{R_{\square,\ast}}\, r(\kappa_1)^{\#\mathcal{V}}
\le \frac{\mathfrak{b}_{\mathcal{T}}}{R_{\square,\ast}}\,
e^{-(\kappa_1-1)\,\#(\text{cubes of } Y_0\setminus\tilde{\square}^{4})},
\end{equation}
and clauses (b)--(e) of the weight lemma hold as stated there, with one link and a single
piece in the weight lemma's count, under the further hypotheses that the weight lemma imposes
for those clauses, where these are invoked.
\end{corollary}

\begin{proof}
The index set $\mathcal{V}$ of decoupling variables of
Theorem~\ref{5.6:thm-full-polydisc} is the set of cubes of
$Y_0\setminus\tilde{\square}^{4}$, and the polydisc $P=P_{\mathcal{V}}(e^{\kappa_1})$ of
that theorem is the open polydisc of radius $e^{\kappa_1}$ in all the variables
$s(\Delta)$, $\Delta\in\mathcal{V}$. By hypothesis~(1) the conclusions of
Theorem~\ref{5.6:thm-full-polydisc} are available, its own hypotheses being (2)--(5) together
with the setting of (1); its clause~(iii) applies: the function
$F_{t}(\sigma;\mathfrak{d}) = \partial_{t_{\square}}F(\sigma,t_{\square};\mathfrak{d})
|_{t_{\square}=0}$ is continuous and separately holomorphic on $P$ for every
$\mathfrak{d}\in\mathcal{D}$, and for all $\sigma\in P$
\begin{equation*}
|F_{t}(\sigma;\mathfrak{d})| \le \mathfrak{b}_{\mathcal{T}}/R_{\square}(\mathfrak{d})
\le \mathfrak{b}_{\mathcal{T}}/R_{\square,\ast}.
\end{equation*}
Together with $\kappa_1\ge 1$, hypothesis~(8), this is the analyticity hypothesis of the
weight lemma for the term $t$, with $\mathcal{V}(t)=\mathcal{V}$,
$\mathcal{D}_{b}=\mathcal{D}$ and the datum-uniform bound
$\mathfrak{b}=\mathfrak{b}_{\mathcal{T}}/R_{\square,\ast}$, as clause~(iii) of
Theorem~\ref{5.6:thm-full-polydisc} records. In the same way clause~(ii) of that theorem
gives, for each fixed $t_{\square 0}$ with $|t_{\square 0}|<R_{\square,\ast}$, that
$\sigma\mapsto F(\sigma,t_{\square 0};\mathfrak{d})$ is continuous and separately
holomorphic on $P$ with modulus at most $\mathfrak{b}_{\mathcal{T}}$, which is the same
hypothesis with $F_{t}=F(\cdot,t_{\square 0};\cdot)$ and
$\mathfrak{b}=\mathfrak{b}_{\mathcal{T}}$. The cubes of $\tilde{\square}^{4}$ carry no
decoupling variable and form the exempt set.

By Definition~\ref{5.6:def-domain-term} (hypothesis~(6)), the $Y_0$-term is one piece of the
weight lemma's form in which the set of differentiated variables is $\mathcal{V}$: every
variable $s(\Delta)$, $\Delta\in\mathcal{V}$, is differentiated once and integrated against
Lebesgue measure on $[0,1]^{\mathcal{V}}$, with $F_{t}$ the function of clause~(iii). In the
weight lemma's notation, in which $W$ denotes the set of differentiated variables, $\lambda$
the number of links, $n$ the number of variables and $\mathcal{P}(t;Y_0)$ the set of pieces of
$t$ with final domain $Y_0$, this reads: $W=\mathcal{V}$, $\lambda=1$, the final domain is
$Y_0$, $n=\#\mathcal{V}$, and $\#\mathcal{P}(t;Y_0)=1$.
Since the hypothesis of the weight lemma holds with
$\mathfrak{b}=\mathfrak{b}_{\mathcal{T}}/R_{\square,\ast}$, its clause~(a) bounds
$p_{Y_0}(\mathfrak{d})$, uniformly in $\mathfrak{d}\in\mathcal{D}$, by
$(\mathfrak{b}_{\mathcal{T}}/R_{\square,\ast})\,r(\kappa_1)^{\#\mathcal{V}}$; this is the
first inequality in~\eqref{5.6:eq-weight-hypothesis-1}. For the second, the bound of
hypothesis~(7) gives $r(\kappa_1)\le e^{-(\kappa_1-1)}$, since $\kappa_1\ge 1$ by
hypothesis~(8). Raising to the power $\#\mathcal{V}$, which is the number of cubes of
$Y_0\setminus\tilde{\square}^{4}$, gives the second inequality. The resulting factor
$e^{-(\kappa_1-1)}$ per cube of $Y_0\setminus\tilde{\square}^{4}$ is of the same kind as
the factor carried by the exponent in $|Y_0\setminus\tilde{\square}^{4}|$ of
\cite[(1.24), (1.25)]{B-CMP116}. Finally, clauses (b)--(e) of the weight lemma require only
its hypothesis, established above, and its own further hypotheses where these are invoked;
with $\lambda=1$ and $\#\mathcal{P}(t;Y_0)=1$ they hold as stated there.

For the terms of \cite[(1.10)]{B-CMP116} the weight lemma's hypothesis thus follows from
Definition~\ref{5.6:def-term-analyticity}, Proposition~\ref{5.6:prop-localized-field} and
Definition~\ref{5.6:def-response-radius}, under the conditions of
Definition~\ref{5.6:def-auxiliary-constants} and the identification of
Definition~\ref{5.6:def-domain-term}; the sum over all terms having a common localisation
domain, with a common domain of analyticity for all of them, is the next step of
\cite{B-CMP116} after \cite[(1.25)]{B-CMP116} and is not part of
Corollary~\ref{5.6:cor-weight-hypothesis}.
\end{proof}

\begin{remark}[Scope: the terms of the large-field operation]\label{5.6:rem-large-field-terms}
Theorem~\ref{5.6:thm-full-polydisc} concerns the term functionals $\mathcal{T}(X;\mathbf{U},\mathbf{J};A)$
of the first expansion \cite[(1.10)]{B-CMP116} of the fluctuation integral in the renormalisation
transformation of \cite[\S1]{B-CMP116}. These functionals are
evaluated at the substituted field \eqref{5.6:eq-substituted-field-a}, built from
$H_k(s(Y_0),B'(B))$. The substituted field contains one family of decoupling variables,
$s(Y_0)=(s(\Delta))_{\Delta\in\mathcal{V}}$. These are the terms of the expansion at site (T) of
Definition~\ref{5.2:def-sites-units}.

The terms at the localisation site (L) of the operation $\mathbf{R}_k$
(Definition~\ref{5.2:def-sites-units}) are different objects. They
are the pieces \cite[(1.61)--(1.66)]{B-CMP122b} of the term $V(\tilde{X})$ in the definition
\cite[(1.71)]{B-CMP122b}, and $(\mathbf{U},\mathbf{J})$
enters them only through the composite configuration \cite[(1.61)--(1.63)]{B-CMP122b}. The
interpolated functions in this composite are localised functions defined in
\cite[(1.59)--(1.65)]{B-CMP122b} by localisation operations connected with the domains $Y_i$,
$i=2,3,4$; the operation connected with $Y_i$ carries its own family of decoupling parameters
$s(Y_i)$, so that the composite carries three such families.
The construction of \cite[\S1]{B-CMP116} belongs to the renormalisation transformation whose terms
are those of the first paragraph, not to the operation $\mathbf{R}_k$.

For the terms at site (L), the conclusion of Theorem~\ref{5.6:thm-full-polydisc}, with the
composite configuration in place of the substituted field and the three families of decoupling
parameters in place of the variables $s(\Delta)$, $\Delta\in\mathcal{V}$, is
Theorem~\ref{5.6:thm-link-holomorphy}. It is proved there as an implication from the hypotheses on
the localised functions of $\mathbf{R}_k$ stated in Definition~\ref{5.6:def-link-readings} and from
the analyticity and boundedness of the term, clause~(e) of
Definition~\ref{5.6:def-link-hypotheses}; the hypotheses of
Definition~\ref{5.6:def-link-readings} are not proved here.
\end{remark}

\begin{definition}[Holomorphy conventions and the complexified configuration manifold]
\label{5.6:not-holomorphy-conventions}
The following classical facts and conventions are fixed for every statement of this chapter
in which a function of complex configurations is called holomorphic or analytic. Throughout,
$N\ge 2$. The classical facts on holomorphic maps between open subsets of complex coordinate
spaces are those recorded in this chapter's list of classical facts on holomorphic maps: the
finite-dimensionality of the spaces of fields and the equivalence of norms, Osgood's lemma and the
meaning of the word \emph{holomorphic}, the closure properties of holomorphic maps, Cauchy's
integral formula and estimate in one variable, and the holomorphy of parameter integrals over a
circle. The same place fixes the notion of a function analytic on a closed polydisc in the
decoupling parameters times a set of $\mathfrak{g}^{\mathbb{C}}$-valued fields; item
\textup{(g)} below states its analogue for configurations with values in $G^{\mathbb{C}}$. The
following items concern complex configurations with values in $G^{\mathbb{C}}$.
\begin{itemize}
\item[\textup{(f)}] $G^{\mathbb{C}}:=\mathrm{SL}(N,\mathbb{C})\subset\mathbb{C}^{N\times N}$ is a closed
complex submanifold (a complex Lie group). The matrix exponential
$\exp\colon\mathbb{C}^{N\times N}\to\mathbb{C}^{N\times N}$ is an entire map taking
$\mathfrak{g}^{\mathbb{C}}$ into $G^{\mathbb{C}}$. Matrix multiplication and inversion on
$\mathrm{GL}(N,\mathbb{C})$ are holomorphic. Let $\mathcal{B}$ and $\mathcal{B}'$ be the finite sets
of bonds of a finite lattice on which $\mathbf{U}$ and $\mathbf{J}$ respectively are defined. A
configuration $(\mathbf{U},\mathbf{J})$, with $\mathbf{U}$ bond-wise $G^{\mathbb{C}}$-valued and
$\mathbf{J}$ bond-wise $\mathfrak{g}^{\mathbb{C}}$-valued, is a point of
\begin{equation*}
\mathfrak{M}:=(G^{\mathbb{C}})^{\mathcal{B}}\times(\mathfrak{g}^{\mathbb{C}})^{\mathcal{B}'}
\subset\mathbb{C}^{n'},\qquad n':=N^2\bigl(\#\mathcal{B}+\#\mathcal{B}'\bigr),
\end{equation*}
where $\mathfrak{g}^{\mathbb{C}}$ is realised as the trace-zero subspace of
$\mathbb{C}^{N\times N}$; $\mathfrak{M}$ is a closed complex submanifold of $\mathbb{C}^{n'}$. A
map into $\mathfrak{M}$ from an open set of $\mathbb{C}^m$ is holomorphic as a map into the
complex submanifold $\mathfrak{M}$ if and only if it is holomorphic as a $\mathbb{C}^{n'}$-valued
map; this is the sense used below. These facts use no Lie bracket, no adjoint action,
no Haar measure, no character and no commutativity, and hold verbatim for every $N\ge 2$.
\item[\textup{(g)}] For a finite set $W$ and $\rho>0$ let $P_W(\rho)$ and $\mathbb{D}_W(\rho)$ be the
open and the closed polydisc of radius $\rho$ in $\mathbb{C}^W$, in the polydisc notation of this
chapter.
Let $S$ be either
\textup{(1)} a relatively open subset of $\mathfrak{M}$, or
\textup{(2)} a set $S=\mathbb{D}_W(\rho)\times S'$ with $S'$ relatively open in $\mathfrak{M}$.
Case \textup{(1)} is the sense in which a set of configurations is called an analyticity domain
in Ba{\l}aban's construction; case \textup{(2)} is the sense in which \cite{B-CMP116} speaks of a
function analytic in the decoupling parameters $s(\cdot)$ of the cluster expansion,
$\sigma=(s(\Delta))_{\Delta\in W}\in\mathbb{D}_W(\rho)$,
and in the configuration. In case \textup{(1)} call $\mathbb{C}^{n'}$ the ambient space and put
$S^\circ:=S$; in case \textup{(2)} call $\mathbb{C}^W\times\mathbb{C}^{n'}$ the ambient space and
put $S^\circ:=P_W(\rho)\times S'$. A function $\Psi$ on $S$ is
\emph{analytic on $S$} if
\begin{equation}\label{5.6:eq-holomorphy-conventions-a}
\begin{gathered}
\Psi \text{ is continuous on } S, \text{ and every } q\in S^\circ \text{ has an open neighbourhood }
V\\
\text{in the ambient space and a holomorphic } \tilde{\Psi}\colon V\to\mathbb{C}\\
\text{with } \tilde{\Psi}=\Psi \text{ on } V\cap S .
\end{gathered}
\end{equation}
In case \textup{(1)} this says equivalently that $\Psi$ is holomorphic on the complex submanifold
$S$. Of the standard senses of analyticity on sets of this form, this is the weakest.
\item[\textup{(h)}] If $\Omega\subset\mathbb{C}^m$ is open, $\Phi\colon \Omega\to\mathfrak{M}$ is
holomorphic with $\Phi(\Omega)\subset S$, and $\Psi$ is analytic on the relatively open set
$S\subset\mathfrak{M}$ in the sense of \eqref{5.6:eq-holomorphy-conventions-a}, then
$\Psi\circ\Phi$ is holomorphic on $\Omega$.
\end{itemize}
\end{definition}

\begin{proof}
For \textup{(f)}: $G^{\mathbb{C}}$ is the zero set of the polynomial $\det-1$ on
$\mathbb{C}^{N\times N}$, hence closed, since polynomials are holomorphic and therefore
continuous; that
it is a complex submanifold and a complex Lie group is classical. The matrix exponential is the
sum of a power series converging on all of $\mathbb{C}^{N\times N}$, hence entire, holomorphy
being equivalent by Osgood's lemma to being locally the sum of a convergent power series; that
it maps $\mathfrak{g}^{\mathbb{C}}$ into $G^{\mathbb{C}}$ is classical. The
entries of a product of two matrices are polynomials in the entries of the factors, so
multiplication is holomorphic, polynomial maps being holomorphic; the entries of $A^{-1}$ are
rational functions of
the entries of $A$ whose denominator $\det A$ does not vanish on $\mathrm{GL}(N,\mathbb{C})$, so
inversion is holomorphic there. The set $\mathfrak{M}$ is a finite product of copies of the
closed complex submanifold $G^{\mathbb{C}}$ and of the linear subspace $\mathfrak{g}^{\mathbb{C}}$
of $\mathbb{C}^{N\times N}$, hence a closed complex submanifold of $\mathbb{C}^{n'}$.

For \textup{(h)}: fix $a_0\in\Omega$ and put $q_0:=\Phi(a_0)\in S$. By
\eqref{5.6:eq-holomorphy-conventions-a} in case \textup{(1)} there are an open
$V\subset\mathbb{C}^{n'}$ containing $q_0$ and a holomorphic $\tilde{\Psi}$ on $V$ with
$\tilde{\Psi}=\Psi$ on $V\cap S$. The map $\Phi$, regarded as $\mathbb{C}^{n'}$-valued, is
holomorphic, hence continuous, so $\Omega':=\Phi^{-1}(V)$ is an open neighbourhood of $a_0$ in
$\Omega$. On $\Omega'$ the composite $\tilde{\Psi}\circ\Phi$ is holomorphic by the chain rule for
composites of holomorphic maps between open sets, applied to the holomorphic maps
$\Phi\colon \Omega'\to V$ and $\tilde{\Psi}\colon V\to\mathbb{C}$. Since
$\Phi(\Omega')\subset V\cap S$ and $\tilde{\Psi}=\Psi$ there, we have
$\tilde{\Psi}\circ\Phi=\Psi\circ\Phi$ on $\Omega'$. This identification does not depend on the
choice of the local extension: any two extensions agree with $\Psi$ on $S\supseteq\Phi(\Omega)$,
so both composites
equal $\Psi\circ\Phi$ where both are defined. Thus $\Psi\circ\Phi$ coincides near every point of
$\Omega$ with a holomorphic function, and, holomorphy being a local property by its
characterisation in Osgood's lemma as a local sum of a convergent power series, $\Psi\circ\Phi$
is holomorphic on $\Omega$.
\end{proof}

\begin{remark}[The large-field papers' references to Part I read as Part II]
\label{5.6:rem-reference-reading}
In \cite{B-CMP122b}, and in \cite{B-CMP119} and \cite{B-CMP122a} wherever the same numbers occur,
the cross-references with section numeral $6$ or $7$ under the label $[\mathrm{I}]$ are read in
this book as references to \cite{B-CMP116}, according to the rule
\begin{equation}\label{5.6:eq-reference-reading-a}
\begin{aligned}
&(6.n)\,[\mathrm{I}] \longmapsto (1.n) \text{ of \cite{B-CMP116}},
&&(7.n)\,[\mathrm{I}] \longmapsto (2.n) \text{ of \cite{B-CMP116}},\\
&\text{Sect.~6 } [\mathrm{I}] \longmapsto \S1 \text{ of \cite{B-CMP116}},
&&\text{Sects.~6, 7 } [\mathrm{I}] \longmapsto \S\S1, 2 \text{ of \cite{B-CMP116}},
\end{aligned}
\end{equation}
where $n$ is the number of a display. In this reading the two-part work \emph{Renormalization group
approach to lattice gauge field theories} is numbered through: Part~I is \cite{B-CMP109}, with
its \S\S0--5, and Part~II is \cite{B-CMP116}, whose \S\S1, 2 correspond to \S\S6, 7. No sentence
of \cite{B-CMP109}, \cite{B-CMP116} or \cite{B-CMP122b} states this through-numbering. The reading
is determined by the table of contents of \cite{B-CMP109} together with the content of the
passages cited, and it is consistent with the list of references of \cite{B-CMP122b} and with the
title page of \cite{B-CMP116}. The
cross-references ``Sects.~3--5 $[\mathrm{I}]$'' and ``(3.3), (3.18), (3.19) $[\mathrm{I}]$'' denote
the sections and displays of \cite{B-CMP109} itself. ``Sects.~3--6 $[\mathrm{I}]$'' denotes \S\S3--5
of \cite{B-CMP109} together with \S1 of \cite{B-CMP116}. In \cite{B-CMP122b} the labels
$[\mathrm{III}]$ and $[\mathrm{IV}]$ denote \cite{B-CMP119} and \cite{B-CMP122a}. The reading
identifies numerals only.
It involves no constant, no threshold and no condition on the coupling.
\end{remark}

\begin{proof}
The reading rests on four facts about the published texts, set out in the next four paragraphs.
The second and the fourth determine it; the first and the third are consistent with it. It is
then confirmed item by item against the displays and sections of \cite{B-CMP116}.

\emph{The list of references.} The list of references of \cite{B-CMP122b} contains the entries
$[\mathrm{I}]$--$[\mathrm{IV}]$ as follows.
Entry $[\mathrm{I}]$ is \emph{Renormalization group approach to lattice gauge field theories.~I},
Commun.\ Math.\ Phys.\ 109, 249--301 (1987). Entry $[\mathrm{II}]$ is the same title with~II,
Commun.\ Math.\ Phys.\ 116, 1--22 (1988). Entry $[\mathrm{III}]$ is \emph{Convergent renormalization
expansions for lattice gauge theories}, Commun.\ Math.\ Phys.\ 119, 243--285 (1988). Entry
$[\mathrm{IV}]$ is \emph{Large field renormalization.~I. The basic step of the R operation},
Commun.\ Math.\ Phys.\ 122, 175--202.

\emph{The sections of \cite{B-CMP109}.} The table of contents of \cite{B-CMP109} lists \S\S0--5
and the references, and no other section:
\begin{itemize}
\item \S0, Introduction. Formulations of Results, p.~249;
\item \S1, An Inductive Description of the Effective Actions, p.~260;
\item \S2, Fluctuation Fields Integral in $(k+1)$-st Renormalization Transformation, p.~265;
\item \S3, An Expansion of Terms in Fluctuation Fields Integral, and Preliminary Analytic
Extensions, p.~269;
\item \S4, Ward--Takahashi Identities and Their Consequences, p.~281;
\item \S5, whose title begins The Analysis of the Vacuum Polarization Tensor, p.~292;
\item References, p.~298.
\end{itemize}
The paper occupies pp.~249--301. Hence \cite{B-CMP109} has no \S6 and no \S7, and a numeral $6$
or $7$ under $[\mathrm{I}]$ cannot refer to a section or display of \cite{B-CMP109}.

\emph{The title page of \cite{B-CMP116}.} The title of \cite{B-CMP116} is \emph{Renormalization
Group Approach to Lattice Gauge Field Theories~II. Cluster Expansions}. Its abstract opens: ``The
fluctuation field integral,
constructed in Part~I, is represented by the exponentiated cluster expansion.'' It refers to
\cite{B-CMP109} by writing I before the numbers, as in (I.3.34) and (I.3.54). Its own sections
are numbered~1, Localizations and Bounds of Terms in the Fluctuation Field Action, and~2, A Final
Localization by a Cluster Expansion. Analyticity Properties and Bounds for Terms in the Effective
Action. Its displays are numbered $(1.n)$ and $(2.n)$.

\emph{The content.} In the passage following \cite[(1.59)]{B-CMP122b}, Ba{\l}aban applies ``all
the transformations and estimates of Sects.~3--5 $[\mathrm{I}]$''.
He then applies ``the expansion (6.9) $[\mathrm{I}]$'', with $\sigma_0$ the set of
cubes $\Delta$ disjoint with $\square_0=\tilde{\square}^5$ and with the enlargement of~$\Lambda$.
He writes this expansion ``in the form (6.10) $[\mathrm{I}]$''. There the summation is over
domains $Y_0\in\mathbf{D}_k$ that contain $\square_0$ and at least one component of that
enlargement. The non-vanishing terms are estimated ``as in (6.24)--(6.29) $[\mathrm{I}]$''.
The sentence introducing \cite[(1.10)]{B-CMP116} describes the terms of the sum on the
right-hand side of (1.9) there as indexed by sets $\sigma$ of cubes. It reduces that sum to a sum
over
connected domains $Y_0$ containing $\tilde{\square}^4$, with the parameters $s(\Delta)=0$ off
$Y_0$. This is the same construction, with the protected neighbourhood $\tilde{\square}^4$ in
place of $\square_0$. In the passage of \cite{B-CMP122b} that carries these cross-references the
label $[\mathrm{II}]$ does not occur. The content of \cite{B-CMP116} is reached there through the
numerals $6$ and $7$ under $[\mathrm{I}]$.

The cross-references covered by \eqref{5.6:eq-reference-reading-a}, together with two further
cross-references, match as follows. For each item we say whether
the match rests on the content of the cited passage or on the numbering alone.
\begin{enumerate}
\item[(i)] ``(6.9) $[\mathrm{I}]$'' is \cite[(1.9)]{B-CMP116}. The match is by numbering and by
content. The content match rests on the sentence introducing \cite[(1.10)]{B-CMP116}: it describes
(1.9) as an expansion over sets $\sigma$ of cubes, as in the use with $\sigma_0$ above. The display
(1.9) itself is not used for the match.
\item[(ii)] ``(6.10) $[\mathrm{I}]$'' is \cite[(1.10)]{B-CMP116}. The match is by numbering and by
content: a sum over connected domains containing the protected neighbourhood, with $s=0$ off the
domain. This holds at each occurrence in \cite{B-CMP122b}, including the sums over domains $Y_i'$
each component of which belongs to $\mathbf{D}_k$ and contains a component of $\Lambda$.
\item[(iii)] ``(6.24)--(6.29) $[\mathrm{I}]$'' are \cite[(1.24)--(1.29)]{B-CMP116}. Display
\cite[(1.24)]{B-CMP116} estimates the term (1.23) of the expansion (1.10). That term is represented
by Cauchy integrals in the decoupling variables $s(\Delta)$ over circles of radius $e^{\kappa_1}$,
and the bound carries the factor
$\exp(-\kappa d_j(X))$. Display \cite[(1.25)]{B-CMP116} continues these estimates. For (1.24) and
(1.25) the match is by numbering and by content. For (1.26)--(1.29) it rests on the numbering alone.
\item[(iv)] ``Sect.~6 $[\mathrm{I}]$'', in the sentence ``This resummation is almost exactly the
same as the one discussed in Sect.~6 $[\mathrm{I}]$'' of \cite{B-CMP122b}, is \cite[\S1]{B-CMP116}.
There all terms
(1.23) having a common localization domain $Y$ are resummed, with a common domain of analyticity
for all terms in it. The match is by numbering and by content: in both places the terms with a
fixed final domain are resummed.
\item[(v)] ``Sects.~3--6 $[\mathrm{I}]$'', in ``We repeat all the considerations and
constructions of Sect.~3, $[\mathrm{III}]$, and Sects.~3--6 $[\mathrm{I}]$'', is \S\S3--5 of
\cite{B-CMP109} together with \S1 of \cite{B-CMP116}. This follows from the absence of a \S6 in
\cite{B-CMP109} and from the identification of Sect.~6 $[\mathrm{I}]$ in item~(iv).
\item[(vi)] ``Sects.~6,7 $[\mathrm{I}]$'', in ``All the above summations were discussed in
Sects.~6,7 $[\mathrm{I}]$, and in Sects.~2,3 $[\mathrm{III}]$'' of \cite{B-CMP122b}, is
\cite[\S\S1, 2]{B-CMP116}.
The match is by numbering. The resummations (2.1)--(2.13) of the cluster expansion in \S2 of
\cite{B-CMP116} are summations of the kind referred to. The label $[\mathrm{III}]$ there is
\cite{B-CMP119}.
\item[(vii)] ``the convergent series (7.13) $[\mathrm{I}]$'' of \cite{B-CMP122b} is
\cite[(2.13)]{B-CMP116}, the exponentiated series of the cluster expansion of \S2 of that paper.
The match is by numbering and by the kind of object. The space (I.6.34) of \cite{B-CMP119}, which
\cite[(3.43)]{B-CMP119} replaces, is the analyticity space (1.34) of \S1 of \cite{B-CMP116}. This
match rests on the numbering.
\item[(viii)] Of the displays named above, (1.26)--(1.29) and (1.34) of \cite{B-CMP116} are
identified by numbering alone.
\end{enumerate}
Together these give the reading \eqref{5.6:eq-reference-reading-a}. The references ``Sects.~3--5
$[\mathrm{I}]$'' and ``(3.3), (3.18), (3.19) $[\mathrm{I}]$'' have numerals at most~$5$, so they
remain references to \cite{B-CMP109}, which has those sections.
\end{proof}

\begin{definition}[The three localization links of a first-part term and their link maps]
\label{5.6:def-localization-links}
We work at the localisation site of the operation $\mathbf{R}_k$ of \cite{B-CMP122b}. There
$\Lambda$ is the integration domain of $\mathbf{R}_k$, $\tilde{\Lambda}$ the enlargement of
$\Lambda$ that \cite{B-CMP122b} defines in terms of $M$-cubes and writes $\Lambda^{\sim}$ (here and
in $\tilde{\square}^2$, $\tilde{\square}^5$ below, the tilde denotes this enlargement operation of
\cite{B-CMP122b}, not an adjacency relation between cells),
$Z$ the large-field region,
$\mathbf{D}_k$ the family of localization domains at scale $k$, and $\mathbf{B}_k$, $\mathbf{B}_1$,
$\mathbf{B}(Y)$ are determining sets. References of \cite{B-CMP122b} to the paper it cites
as [I] are read as in Remark~\ref{5.6:rem-reference-reading}, display
\eqref{5.6:eq-reference-reading-a}. In particular, ``Sects.~3--5 [I]'' are
\cite[\S\S3--5]{B-CMP109}. The expansion ``(6.9) [I]'', written ``in the form of the expansion
(6.10) [I]'', is the localization expansion \cite[(1.9)]{B-CMP116} written in the form
\cite[(1.10)]{B-CMP116}. Every localization expansion in this definition is of this kind, and we
call it a \emph{localization expansion in the form} \cite[(1.10)]{B-CMP116}.

\begin{enumerate}
\item[(a)] \emph{The chain of backgrounds.} Five background configurations occur, in this order:
\begin{itemize}
\item $U_k'$;
\item $U_k^0$;
\item $U_{k,Z}''$, the configuration given by \cite[(1.53)]{B-CMP122b};
\item $U_{k,\Lambda}$, given by \cite[(1.56)]{B-CMP122b};
\item $U_k$.
\end{itemize}
The step from $U_k'$ to $U_k^0$ is made by \cite[(1.59)]{B-CMP122b} or \cite[(1.60)]{B-CMP122b}.
The step from $U_k^0$ to $U_{k,Z}''$ is made by the representation \cite[(1.61)]{B-CMP122b},
which is built from \cite[(1.51), (1.54)]{B-CMP122b}. The step from $U_{k,Z}''$ to
$U_{k,\Lambda}$ is made by \cite[(1.62)]{B-CMP122b}, which uses the expansion
\cite[(1.57)]{B-CMP122b}. The step from $U_{k,\Lambda}$ to $U_k$ is made by
\cite[(1.63)]{B-CMP122b}, which uses the expansion \cite[(1.58)]{B-CMP122b}.

\item[(b)] \emph{The first operation.} Let $E(X;U_k')$ be a term of the first part of the effective
action in \cite[(1.49)]{B-CMP122b} other than its first two terms, or one of the new terms with
localization domains $Y_0$ constructed in \cite{B-CMP122b} for the second part. In
\cite{B-CMP122b} such a term also carries a point $z$, written as its last argument, with
$z\in\Omega_n\setminus\Omega_{n+1}$ for a scale $n\ge j$; here $j$ is the scale of the term and
$\Omega_n$ is the region of scale $n$ in the sequence of regions of \cite{B-CMP122b}. Let
$\square$ be the cube of the partition $\pi_n$ of scale $n$ that contains $z$, and let
$\tilde{\square}^2$ and $\square_0=\tilde{\square}^5$ be its cube neighbourhoods. In the case
$X\subset\tilde{\square}^2$ one applies the identity \cite[(1.59)]{B-CMP122b}, which passes
through a box representation over $\square_0$, of the kind of \cite[(3.6)]{B-CMP109}, at the
shifted background that appears in the square bracket of
\eqref{5.6:eq-localization-links-1}. Otherwise
one applies the identity \cite[(1.60)]{B-CMP122b}, which has the structure
\begin{equation}\label{5.6:eq-localization-links-1}
\begin{split}
E(X;U_k') = E(X;U_k^0)
 &+ \bigl[E\bigl(X;\exp\bigl(i\eta H_k'(g_kB)\bigr)\cdot U_k^0\bigr)\\
 &\qquad - E(X;U_k^0)\bigr].
\end{split}
\end{equation}
Here $\eta$ is the lattice spacing and $g_kB$ the small fluctuation field. To the square bracket one
applies ``all the transformations and estimates of Sects.~3--5 [I] (with $\eta$ replaced by
$L^{-n}$)'', that is, those of \cite[\S\S3--5]{B-CMP109}. One then applies a localization expansion
in the form \cite[(1.10)]{B-CMP116}. It is determined by the set $\sigma_0$ of cubes $\Delta$
disjoint with ``$\square_0$, or $X$, and with $\tilde{\Lambda}$''.
Here and in (d), the letters $\sigma_0$ and $\sigma_0^{(i)}$ denote sets of cubes, as in
\cite{B-CMP122b}; they are not to be confused with $\sigma^{(i)}$ and $\sigma$ of (f), which are
families of complex parameters.

This expansion decouples over domains $Y_0\in\mathbf{D}_k$. It is the first localization link of
Ba{\l}aban's construction. The terms having a common domain $Y_0$ are resummed into ``some new
terms''. This resummation is ``almost exactly the same as the one discussed in Sect.~6 [I]'', that
is, the resummation of \cite[\S1]{B-CMP116}. It takes place before a single term is taken in (c).
The inductive factor of \cite[(1.66)]{B-CMP122b} carries the domain $Y_0$, ``the
smallest domain from $\mathbf{D}_k$ containing $X$''. This first link is not among the three links
of (d).

\item[(c)] \emph{One term.} After the first operation all the terms depend on the background
$U_k^0$ of the chain (a). In \cite{B-CMP122b} the term chosen below is denoted $E(X,U_k)$,
``although it may be any term in the sum''. We write $E(X;U_k^0)$, the background being the one
that the first operation produces. Fix one such term with $X\cap Z\neq\emptyset$. Let
$Y_1\in\mathbf{D}_k$ be the smallest localization domain containing $X$ and the components of $Z$
that intersect $X$.

\item[(d)] \emph{The three links.} For $i=2,3,4$, link $i$ has three constituents: a determining
set, a representation of a background of the chain (a), and one localization expansion in the form
\cite[(1.10)]{B-CMP116}. In each link the expansion is determined by the set $\sigma_0^{(i)}$ of
cubes $\Delta$ disjoint with $\Lambda$.
\begin{itemize}
\item[] \emph{Link 2.} It uses the determining set $\mathbf{B}_k(Y_1)$ and the representation
\cite[(1.61)]{B-CMP122b} of the background on $Y_1$. The gauge transformation $u_k$ is kept in
this representation. Here $Y_1^0$ denotes the domain on which $\mathbf{B}_k(Y_1)$ and $\mathbf{B}_k$
coincide. The localization expansion is applied to the function $\mathbf{H}^{(2)}$ of (e). The sum
runs over domains $Y_2'$ such that each component of $Y_2'$ belongs to $\mathbf{D}_k$, contains at
least one component of $\Lambda$, and intersects $Y_1\setminus Y_1^0$. A term of the expansion
depends on $U_{k,Z}''$ restricted to $Y_2:=Y_1\cup Y_2'$.
\item[] \emph{Link 3.} It uses the determining set
$\mathbf{B}_1(Y_2)=\mathbf{B}(Y_2)\cup\mathbf{B}_1$ and the representation
\cite[(1.62)]{B-CMP122b} of $U_{k,Z}''$ on $Y_2$, where the expansion \cite[(1.57)]{B-CMP122b} is
used. The function on its right-hand side ``depends analytically on the field $U_{k,\Lambda}$''.
Here $Y_2^0$ denotes the domain on which $\mathbf{B}_1(Y_2)$ and $\mathbf{B}_1$ coincide. The
localization expansion is applied to the function $H_B$. The sum runs over domains $Y_3'$ such
that each component belongs to $\mathbf{D}_k$, contains a component of $\Lambda$, and intersects
$Y_2\setminus Y_2^0$. A term depends on $U_{k,\Lambda}$ restricted to $Y_3:=Y_2\cup Y_3'$.
\item[] \emph{Link 4.} It uses the determining set
$\mathbf{B}_k(Y_3)=\mathbf{B}(Y_3)\cup\mathbf{B}_k$ and the representation
\cite[(1.63)]{B-CMP122b} of $U_{k,\Lambda}$ on $Y_3$, where the expansion
\cite[(1.58)]{B-CMP122b} is used and ``the function $H_k$ is given by''
\cite[(1.64)]{B-CMP122b}. Here $Y_3^0$ denotes the
domain on which $\mathbf{B}_k(Y_3)$ and $\mathbf{B}_k$ coincide. The localization expansion is
applied to the function $H_k$. The sum runs over domains $Y_4'$ such that each component belongs
to $\mathbf{D}_k$, contains a component of $\Lambda$, and intersects $Y_3\setminus Y_3^0$. A term
depends on $U_k$ restricted to $Y_4:=Y_3\cup Y_4'$. ``This is the final localization.''
\end{itemize}

\item[(e)] \emph{The link functions.}
\begin{itemize}
\item[] $\mathbf{H}^{(4)}:=H_k$, the function given by \cite[(1.64)]{B-CMP122b}. It is the
$H$-function of the representation \cite[(1.63)]{B-CMP122b} and carries the name of the function
$H_k$ of \cite[(1.58)]{B-CMP122b}; the identity of the two functions so named is not used below.
In its discussion of the analyticity properties of the terms of this expansion,
\cite{B-CMP122b} names the localized function as a function of its link's parameters:
``The function $H_k(s(Y_4))$ on $(Y_3^0)^c\cap Y_3$ is also very small, because of the localization
of the argument, and the exponential decay''. Here $s(Y_4)$ denotes the family of parameters
written $\sigma^{(4)}$ in (f); \cite{B-CMP122b} names them at this place without fixing their
range.
\item[] $\mathbf{H}^{(3)}:=H_B$, the function of the expansion \cite[(1.57)]{B-CMP122b}. Of it
Ba{\l}aban states in \cite{B-CMP122b} that ``The field in the argument of the function $H_B$ has
again a right localization'', that the function is ``decaying exponentially off'' a domain written
there with the letter $Z$, and that ``The function $H_B$ is very small on
$Y_2\setminus Y_2^0$''. For it,
\cite{B-CMP122b} constructs ``again the localization expansion'' in the form
\cite[(1.10)]{B-CMP116}, as it does for $H_k$. The dependence of the localized function on the
parameters $\sigma^{(3)}$ of (f) is that of this localization expansion.
\item[] $\mathbf{H}^{(2)}$ is the function in the representation \cite[(1.61)]{B-CMP122b} to which
the localization operation of link~2 is applied. The letter $\mathbf{H}^{(2)}$ is a name fixed here.
Only its role as this function of the representation \cite[(1.61)]{B-CMP122b}, which is built from
\cite[(1.51), (1.54)]{B-CMP122b}, is used below. Ba{\l}aban states in \cite{B-CMP122b} that
$u_k(x)$ is an analytic function of the restriction of an $H$-function to the block of the set
$\mathbf{B}_k$ to which the point $x$ belongs, and that this function is given by an explicit
formula. In what follows this $H$-function is taken to be
$\mathbf{H}^{(2)}$. Thus
$u_k$ is one more map composed inside the representation of link~2, and it forms part of the link
map $\mathcal{E}_2$ of (f).
\end{itemize}

\item[(f)] \emph{The link maps.} For $i=2,3,4$, the localization expansion of link $i$ introduces
parameters $s(\Delta)$, indexed by a finite set $\mathcal{V}_i$ of cubes. We write
$\sigma^{(i)}=(s(\Delta))_{\Delta\in\mathcal{V}_i}$ for them, and
$\mathbf{H}^{(i)}(\sigma^{(i)})$ for the localized link function. Writing $t:=E(X;U_k^0)$ for the
fixed term of (c), we also put
\begin{equation*}
\mathcal{V}(t):=\mathcal{V}_2\sqcup\mathcal{V}_3\sqcup\mathcal{V}_4, \qquad
\sigma:=(\sigma^{(2)},\sigma^{(3)},\sigma^{(4)}).
\end{equation*}
The \emph{link map} $\mathcal{E}_i$ is the right-hand side of \cite[(1.61)]{B-CMP122b},
\cite[(1.62)]{B-CMP122b} and \cite[(1.63)]{B-CMP122b} for $i=2,3,4$ respectively, with
$\mathbf{H}^{(i)}$ replaced by $\mathbf{H}^{(i)}(\sigma^{(i)})$ and with the complex background
configuration variables $(\mathbf{U},\mathbf{J})$ introduced as described in \cite{B-CMP122b}.
It is regarded as a function of
two arguments: the values of the localized link function, and the background that the
representation uses. That background is $U_{k,Z}''$ for $i=2$, $U_{k,\Lambda}$ for $i=3$, and
$U_k$ for $i=4$. The value of $\mathcal{E}_i$ is a configuration on $Y_1$, $Y_2$, $Y_3$
respectively. By the chain (a), the output of $\mathcal{E}_{i+1}$ is the background argument of
$\mathcal{E}_i$ for $i=2,3$. The domain of the extension, the analyticity properties of the link
maps, and the domains on which they are used, are not part of this definition.

\item[(g)] \emph{The piece index.} The pieces of the term $E(X;U_k^0)$ are indexed by the triples
$(Y_2',Y_3',Y_4')$ ``satisfying all the conditions described above'', that is, those of (d).
Summing over these triples represents $E(X;U_k^0)$ as a sum of terms. The term of a given triple
depends analytically on $U_k$ restricted to $Y_4=Y_1\cup Y_2'\cup Y_3'\cup Y_4'$. The number of
links is $\lambda=3$.
\end{enumerate}
\end{definition}

\begin{definition}[The per-link analyticity readings]\label{5.6:def-link-readings}
We work at the localisation site (L) of the operation $\mathbf{R}_k$, with a first-part term
$t=E(X;\cdot)$, its three localisation links $i\in\{2,3,4\}$, their link domains $Y_i'$, the
nested domains $Y_1\subset Y_2\subset Y_3\subset Y_4$ and the link maps $\mathcal{E}_i$, all as in
Definition~\ref{5.6:def-localization-links}. We write
$\Lambda$ for the integration domain of $\mathbf{R}_k$.

These readings introduce no constant, threshold or smallness condition of their own. The only
conditions that occur in them are those of \cite[\S1]{B-CMP116}: $\kappa_1$ is sufficiently big,
and $\delta_1M\ge\kappa_1$, where $\delta_1$ is the constant in the walk bound
\cite[(1.11)]{B-CMP116}. Each of the analyticity readings at the links $i=2,3,4$ and the
outer-parameter reading below is a hypothesis. Wherever one of them is used, it appears as an
explicit hypothesis of the statement proved.

\emph{Polydiscs.} For a finite set $W$ and $\rho>0$ put
\begin{align*}
P_W(\rho)&:=\{\sigma=(s(\Delta))_{\Delta\in W}\in\mathbb{C}^W:\
|s(\Delta)|<\rho\ \text{for all }\Delta\in W\},\\
\mathbb{D}_W(\rho)&:=\{\sigma=(s(\Delta))_{\Delta\in W}\in\mathbb{C}^W:\
|s(\Delta)|\le\rho\ \text{for all }\Delta\in W\}.
\end{align*}
Thus $P_W(\rho)$ is the open polydisc and $\mathbb{D}_W(\rho)$ the closed polydisc of radius
$\rho$ in $\mathbb{C}^W$.

\emph{The decoupling variables of link $i$.} Let $i\in\{2,3,4\}$. In \cite{B-CMP122b}
link~$i$ is expanded in the form there called ``(6.10)~[I]''.
Under the reading
\eqref{5.6:eq-reference-reading-a} of Remark~\ref{5.6:rem-reference-reading}, this expansion is
the expansion \cite[(1.10)]{B-CMP116}.

The $Y_i'$-term of this expansion carries finitely many decoupling variables
$\sigma^{(i)}=(s(\Delta))_{\Delta\in\mathcal{V}_i}$. Here $\mathcal{V}_i$, fixed in
\eqref{5.6:eq-link-readings-a} below, is a finite set of $M$-cubes of the scale of link~$i$
contained in $Y_i'$. Since $M^{-d}|Y_i'|$ is the number of $M$-cubes in $Y_i'$, we have in
particular $\#\mathcal{V}_i\le M^{-d}|Y_i'|$.

We read $\mathcal{V}_i$ on the pattern of \cite[(1.10)]{B-CMP116}. There the sum runs over
connected domains $Y_0$, the parameters are $s=0$ on $Y_0^c$, and the parameters satisfy
$|s(\Delta)|\le e^{\kappa_1}$ for $\Delta\subset Y_0\setminus\tilde{\square}^4$. Accordingly we
set
\begin{equation}\label{5.6:eq-link-readings-a}
\begin{gathered}
\sigma_0^{(i)}:=\{\Delta\ \text{an $M$-cube of the scale of link }i:\ \Delta\cap\Lambda=\emptyset\},\\
\mathcal{V}_i:=\{\Delta\in\sigma_0^{(i)}:\ \Delta\subset Y_i'\},\\
s(\Delta)=0\quad\text{for }\Delta\in\sigma_0^{(i)}\setminus\mathcal{V}_i .
\end{gathered}
\end{equation}
We further put:
\begin{itemize}
\item $\mathcal{V}(t):=\mathcal{V}_2\sqcup\mathcal{V}_3\sqcup\mathcal{V}_4$, the set indexing
the variables $s(\Delta)$ of the three links of the term $t$ at the site (L);
\item $\mathbb{D}_i:=\mathbb{D}_{\mathcal{V}_i}(e^{\kappa_1})$ (closed) and
$P_i:=P_{\mathcal{V}_i}(e^{\kappa_1})$ (open);
\item $\mathbb{D}:=\mathbb{D}_2\times\mathbb{D}_3\times\mathbb{D}_4$ and
$P:=P_2\times P_3\times P_4$.
\end{itemize}
Under the identification of $\mathbb{C}^{\mathcal{V}(t)}$ with
$\mathbb{C}^{\mathcal{V}_2}\times\mathbb{C}^{\mathcal{V}_3}\times\mathbb{C}^{\mathcal{V}_4}$ we
have $P=P_{\mathcal{V}(t)}(e^{\kappa_1})$. Of the sets $\mathcal{V}_i$, only their finiteness is
used in what follows.

\emph{The parameter domain.} In \cite{B-CMP122b} the localised function $H_k$ of the
representation \cite[(1.63)]{B-CMP122b} is expanded in the form of \cite[(1.10)]{B-CMP116};
this again uses \eqref{5.6:eq-reference-reading-a}. The reference fixes the form of the
expansion. The construction of \cite[\S1]{B-CMP116} which produces this form estimates its
terms on complex parameters with $|s(\Delta)|\le e^{\kappa_1}$, with $\kappa_1$ sufficiently
big and under $\delta_1M\ge\kappa_1$, by Cauchy integrals over the circles
$|s(\Delta)|=e^{\kappa_1}$. As part of the readings, the parameter domain at link~4 is taken
to be the corresponding domain
\begin{equation*}
\mathbb{D}_4=\{\sigma^{(4)}:\ |s(\Delta)|\le e^{\kappa_1}\ \text{for all }\Delta\in\mathcal{V}_4\}
=\mathbb{D}_{\mathcal{V}_4}(e^{\kappa_1}),
\end{equation*}
with the estimates obtained by Cauchy integrals over the circles $|s(\Delta)|=e^{\kappa_1}$,
$\Delta\in\mathcal{V}_4$; this identification is inferred from the form of the expansion. The
same reference to \cite[(1.10)]{B-CMP116} is made at links 3 and 2, for the representations
\cite[(1.62)]{B-CMP122b} and \cite[(1.61)]{B-CMP122b}, and $\mathbb{D}_3$ and $\mathbb{D}_2$ are
taken in the same way as parameter domains of that construction. The readings therefore
consist of this identification of the domains $\mathbb{D}_i$ with the domains of
\cite[\S1]{B-CMP116}, together with what they assert about the functions of
\cite[(1.61)--(1.63)]{B-CMP122b} on these domains.

\emph{The localised functions.} We write:
\begin{itemize}
\item $\mathbf{H}^{(4)}$ for the localised function $H_k(s(Y_4))$ of the representation
\cite[(1.63)]{B-CMP122b}, given by \cite[(1.64)]{B-CMP122b};
\item $\mathbf{H}^{(3)}$ for the localised function $H_B$ of the representation
\cite[(1.62)]{B-CMP122b}, defined at \cite[(1.57)]{B-CMP122b};
\item $\mathbf{H}^{(2)}$ for the localised function of the representation
\cite[(1.61)]{B-CMP122b}, defined at \cite[(1.51), (1.54)]{B-CMP122b}.
\end{itemize}

The arguments and backgrounds of $\mathbf{H}^{(i)}$ are as follows:
\begin{itemize}
\item for $i=4$: the argument consists of the fields on which $H_k$ of \cite[(1.64)]{B-CMP122b}
acts, and the background is $(\mathbf{U},\mathbf{J})$ restricted to $Y_4$;
\item for $i=3$: the argument is the extended right-hand side of \cite[(1.63)]{B-CMP122b}, and
the background consists of the fields of \cite[(1.62)]{B-CMP122b};
\item for $i=2$: the argument is the extended right-hand side of \cite[(1.62)]{B-CMP122b}, and
the background consists of the other expressions of \cite[(1.61)]{B-CMP122b}, among them the
gauge transformation $u_k$.
It enters through the statement of \cite{B-CMP122b} that $u_k(x)$ is an analytic function of
the localised $H$-function restricted to the block in question in \cite{B-CMP122b}, here taken
at $\mathbf{H}^{(2)}$.
\end{itemize}
We write $\mathcal{O}_i$ for the relatively open domain of these arguments and backgrounds on
which \cite{B-CMP122b} considers $\mathbf{H}^{(i)}$.

\emph{The analyticity reading at link $i$, for $i=4,3,2$.} Write $\xi$ for the argument and
$\psi$ for the background of $\mathbf{H}^{(i)}$, as listed above. The analyticity reading at
link $i$ is the following hypothesis.
\begin{enumerate}
\item[(1)] The dependence of the function $\mathbf{H}^{(i)}$ on $\sigma^{(i)}$ enters through
operators with generalised random walk expansions, into which the parameters are introduced as
in \cite[\S1]{B-CMP116} and which obey the walk bound \cite[(1.11)]{B-CMP116}, and through
contraction fixed points of small arguments, of the type of \cite[(1.12)--(1.17)]{B-CMP116}.
\item[(2)] There is a constant $C$, of the kind of the constants in
\cite[(1.11), (1.18), (1.21)]{B-CMP116} and independent of
$(\sigma^{(i)},\xi,\psi)\in\mathbb{D}_i\times\mathcal{O}_i$, such that
\begin{equation}\label{5.6:eq-link-readings-b}
\begin{gathered}
(\sigma^{(i)},\xi,\psi)\longmapsto\mathbf{H}^{(i)}(\sigma^{(i)};\xi,\psi)\ \text{ is analytic on }
\mathbb{D}_i\times\mathcal{O}_i,\ \text{ and}\\
\|\mathbf{H}^{(i)}(\sigma^{(i)};\xi,\psi)\|\le C\,e^{16\kappa_1}\,\|\xi\|
\qquad\bigl((\sigma^{(i)},\xi,\psi)\in\mathbb{D}_i\times\mathcal{O}_i\bigr),
\end{gathered}
\end{equation}
where, in the bound, $\|\cdot\|$ on the left is any one of the admissible norms of
\cite[\S1]{B-CMP116}, namely the norms in which the bound following \cite[(1.11)]{B-CMP116} is
stated, and $\|\xi\|$ is the size of the argument $\xi$ measured as in those bounds. The bound
holds under $\delta_1M\ge\kappa_1$, uniformly in the background $\psi$, in the same way
as the bounds of \cite[\S1]{B-CMP116}.
\item[(3)] The values of $\mathbf{H}^{(i)}$ lie in the argument domain of the map applied next
in the composition of Definition~\ref{5.6:def-localization-links}.
\end{enumerate}
Analyticity on $\mathbb{D}_i\times\mathcal{O}_i$, which is closed in the variables
$\sigma^{(i)}$ only, is meant in the sense of \eqref{5.6:eq-holomorphy-conventions-a} of
Notation~\ref{5.6:not-holomorphy-conventions}: continuous on this set and holomorphic at its
interior points.

Items (1)--(3) say that $\mathbf{H}^{(i)}$ is a function to which the analysis of
\cite[\S1]{B-CMP116} applies. This is the analysis that leads from the bound
\cite[(1.11)]{B-CMP116} to the analyticity of $H_k(s(Y_0),B')$ of \cite[(1.17)]{B-CMP116} (in
the notation of that paper) in $s(Y_0)$ and $B'$. In \cite[\S1]{B-CMP116} the structure of
item (1) is what the analysis uses, and statements of the form of item (2) are what it yields
for the functions treated there; for $\mathbf{H}^{(i)}$ all three items are assumed.

For the functions of \cite[\S1]{B-CMP116} themselves, written in the notation of that paper,
this is proved there:
\begin{itemize}
\item for $H(s(Y_0))X$, and for the propagators $G(s(Y_0))$ and $H_0(s(Y_0))$, which have the
bound \cite[(1.11)]{B-CMP116};
\item for the objects denoted there $D$, $V$, $A_0$;
\item for $H_k(s(Y_0),B')$ of \cite[(1.17)]{B-CMP116}.
\end{itemize}
This is the analogue of \eqref{5.6:eq-link-readings-b} at the site (T) of
Definition~\ref{5.2:def-sites-units}, proved in
\cite[\S1]{B-CMP116}. For the functions $\mathbf{H}^{(i)}$ of the operation $\mathbf{R}_k$,
items (1)--(3), stated on the domains $\mathbb{D}_i$ so identified, are the content of the
analyticity reading at link $i$.

\emph{The outer-parameter reading.} The first operation of $\mathbf{R}_k$ (at
\cite[(1.59)]{B-CMP122b} or \cite[(1.60)]{B-CMP122b}) applies the transformations and estimates
of \cite[\S\S3--5]{B-CMP109}, among them the device of \cite[(3.4)]{B-CMP109} with parameters
$t_{\square}$ and cut-off functions $\zeta_{\square}$. Let a piece of the expansion at the site
(L) be a derivative with respect to an outer complex parameter $t_{\square}$ of this device
(a decoupling parameter of the first operation, not among the variables $s(\Delta)$).

The outer-parameter reading is the following hypothesis. The parameter $t_{\square}$ enters
such a piece through the outer factor
\begin{equation}\label{5.6:eq-link-readings-c}
\exp\Bigl(i\eta\bigl(\mathfrak{t}\,\zeta_{*}+t_{\square}\,\zeta_{\square}\bigr)\cdot
H_k'(g_kB)\Bigr).
\end{equation}
Here $H_k'(g_kB)$ is taken at its localised value, and \eqref{5.6:eq-link-readings-c} is the
shift inside the bracket of \cite[(1.59)]{B-CMP122b} and \cite[(1.60)]{B-CMP122b}. The cut-off
functions $\zeta_{*}$, $\zeta_{\square}$ are those of \cite[(3.4)]{B-CMP109}, and $\mathfrak{t}$
denotes the real interpolation parameters of \cite[\S\S3--5]{B-CMP109}, held fixed. The
hypothesis further provides, for each datum $\mathfrak{d}$, that is, each choice of the
background, the auxiliary fields and the parameters $\mathfrak{t}$ held fixed, a radius
$R(\mathfrak{d})>0$, the membership radius, of a disc $\{|t_{\square}|<R(\mathfrak{d})\}$ of
values of the outer parameter on which the piece is considered. This radius is of the kind of
the radius of the circle
\cite[(1.22)]{B-CMP116}, over which the $t_{\square}$-integration of the construction of
\cite[\S1]{B-CMP116} runs.

The outer-parameter reading concerns only the pieces just described.
\end{definition}

\begin{definition}[Data, argument and value spaces of a term]\label{5.6:not-term-spaces}
Let $Y_1\subset Y_2\subset Y_3\subset Y_4$ be the nested domains of the three localisation links of
Definition~\ref{5.6:def-localization-links}, let $X\subset Y_1$ be a localisation domain, and let
$t=E(X;\cdot)$ be a term of the effective action, regarded as a function of a configuration on $Y_1$.
Throughout the analysis of the link maps the letter $t$ denotes such a term.

\emph{Data.} A \emph{datum} $\mathfrak{d}$ consists of the complex background configuration
$(\mathbf{U},\mathbf{J})$ of the large-field step, together with the following objects, all held fixed:
the auxiliary real and complex fields of \cite[(1.51)--(1.63)]{B-CMP122b}, namely the fluctuation field
$g_kB$, localised in the integration domain $\Lambda$ of $\mathbf{R}_k$, the fields $V''$, $V'_\Lambda$,
$V'_k$, and the data of \cite[(1.53)]{B-CMP122b} and \cite[(1.56)]{B-CMP122b}; and the real
interpolation parameters $\mathfrak{t}$ of \cite[Sections~3--5]{B-CMP109}. The datum ranges over
\begin{equation}\label{5.6:eq-term-spaces-1}
\begin{split}
\mathcal{D} := {}&\{(\mathbf{U},\mathbf{J})\in\mathcal{U}_{\mathrm{arg}}(Y_4)\}\\
&\times(\text{the small-field ranges of the auxiliary fields in \cite{B-CMP122b}})\\
&\times(\text{the ranges of the interpolation parameters } \mathfrak{t}).
\end{split}
\end{equation}

\emph{Argument and value spaces.} The argument space $\mathcal{U}_{\mathrm{arg}}(Y_4)$ is the space
$\tilde{U}^c_k(Y_4,\alpha_0,\alpha_1)$ of \cite[(1.65)]{B-CMP122b}, the space of the variables
$(\mathbf{U},\mathbf{J})$ there. The value space $\mathcal{U}_{\mathrm{val}}(Y_1)$ is the space of values of
the function $U_k$ in \cite[(1.65)]{B-CMP122b}: the space on $Y_1$ defined by
\cite[(1.64)--(1.69), pp.~190--191]{B-CMP122a} with the index $n$ there set equal to the number of
stages of the construction of \cite{B-CMP122a} (this number is written $N$ in \cite{B-CMP122a} and in
\cite{B-CMP122b}; that letter is unrelated to the rank of $\mathrm{SU}(N)$).

\emph{Analyticity domain.} We write $\mathcal{A}_t$ for the analyticity domain of the term
$t=E(X;\cdot)$ as a function of a configuration on $Y_1\supseteq X$. In the sense of
\eqref{5.6:eq-holomorphy-conventions-a},
$\mathcal{A}_t$ is relatively open in $\mathfrak{M}(Y_1)$, and $t$ is continuous on $\mathcal{A}_t$ and
holomorphic there. \cite[(1.65) and the paragraph following it]{B-CMP122b} states that the space of
values of the function in \cite[(1.65)]{B-CMP122b} is contained in the analyticity domain of the term
$E(X)$; we use this inclusion in the form
\begin{equation}\label{5.6:eq-term-spaces-2}
\mathcal{U}_{\mathrm{val}}(Y_1)\subset\mathcal{A}_t .
\end{equation}

\emph{Auxiliary data and the link maps.} In \cite{B-CMP122b}, after $(\mathbf{U},\mathbf{J})$ has been
introduced, the remaining fields in \cite[(1.62)]{B-CMP122b} are replaced by functions of $\mathbf{U}$.
Those auxiliary data of $\mathfrak{d}$ with values in the group $\mathrm{SU}(N)$, or with real values,
that are not so replaced are carried as a fixed
parameter of each link map, outside the holomorphy. For $i\in\{2,3,4\}$ the link map of link $i$ is thus
written
\begin{equation}\label{5.6:eq-term-spaces-3}
\mathcal{E}_i(\cdot,\cdot;\mathfrak{d})\colon \mathbb{D}_i\times\mathcal{O}_i\to\mathcal{O}_{i-1},
\qquad
\mathcal{O}_i\subset\mathfrak{M}(Y_i)\times\mathbb{C}^{m_i}\ \text{relatively open}.
\end{equation}
Here $\mathbb{D}_i$ is the closed polydisc of radius $e^{\kappa_1}$ in the decoupling variables
$\sigma^{(i)}$ of link $i$, as in Definition~\ref{5.6:def-link-readings}, and $m_i\ge 0$ is the number of
additional complex coordinates, for auxiliary data, carried by the link map of link $i$; $m_i=0$ if there
are none. The link maps are considered in this form: the per-link analyticity readings of
Definition~\ref{5.6:def-link-readings} include the hypothesis that $\mathcal{E}_i$ is jointly holomorphic
in $\sigma^{(i)}$ and in the configuration and complex components at fixed $\mathfrak{d}$, the datum
entering only as a parameter, outside the holomorphy.

With this convention, and with $\mathfrak{M}(Y_i)$ a closed complex submanifold of a coordinate space
in the sense of \eqref{5.6:eq-holomorphy-conventions-a}, the product
$\mathfrak{M}(Y_i)\times\mathbb{C}^{m_i}$ is again a closed complex submanifold of a coordinate space,
and each $\mathcal{O}_i$ is literally of the form required by the hypothesis of the
composition lemma for the link maps that its argument
domains be relatively open in a closed complex submanifold. A factor consisting of a real group is never
treated as a complex submanifold. The convention is a choice of notation: the per-link analyticity
readings of Definition~\ref{5.6:def-link-readings} are hypotheses, and the convention changes none of them.

\emph{Homonyms.} The following letters are kept apart.
\begin{itemize}
\item The term $t$ is distinct from the interpolation parameters $\mathfrak{t}$ and from the outer complex
parameter $t_{\square}$ of the device of \cite[(3.4)]{B-CMP109}.
\item The subscripted letters $\mathcal{A}_t$ and $\mathcal{O}_i$ are unrelated to the free lattice field
$\mathcal{A}$ and to the observables $\mathcal{O}(f)$ of the glossary.
\item The function $H_k$ of \cite[(1.64)]{B-CMP122b} is not the function $H_k$ of
\cite[(1.17)]{B-CMP116}: it belongs to a different step and has a different argument. Where confusion
could arise, the former is written $\mathbf{H}^{(4)}$.
\item The decoupling parameters written $s$, with components $s(\Delta)$, in \cite{B-CMP116} and
\cite{B-CMP122b} are the components of $\sigma^{(i)}$.
\item The letter $\delta_1$ denotes the constant of the random-walk bound following
\cite[(1.11)]{B-CMP116}, the one in the condition $\delta_1 M\ge\kappa_1$ there.
\item The domain $\Lambda$ of integration of $\mathbf{R}_k$ is distinct from the number $\lambda$ of links
in the composition lemma.
\end{itemize}

\emph{The reference {[I]}.} Where \cite{B-CMP122b} cites Sections~0--5 of its reference [I], the
reference is to the sections of the same numbers in \cite{B-CMP109}. Sections~6 and~7 of [I] are
Sections~1 and~2 of \cite{B-CMP116}, and an equation cited in \cite{B-CMP122b} as $(6.n)$ [I] or
$(7.n)$ [I] is equation $(1.n)$ or $(2.n)$ of \cite{B-CMP116}, respectively.
\end{definition}

\begin{definition}[Setting and hypotheses for the link maps]\label{5.6:def-link-hypotheses}
The conventions of Notation~\ref{5.6:not-holomorphy-conventions} are in force, in particular its
finiteness convention, its analyticity convention, and the notion
of analyticity on configuration manifolds in \eqref{5.6:eq-holomorphy-conventions-a}. So are the reading
\eqref{5.6:eq-reference-reading-a} of the references in \cite{B-CMP122b}, fixed in
Remark~\ref{5.6:rem-reference-reading}, and the reading \eqref{5.6:eq-link-readings-a} of the decoupling
variables, fixed in Definition~\ref{5.6:def-link-readings}.

We work at step $k$ of the operation $\mathbf{R}_k$. Let $t=E(X;\cdot)$ be a term of the first part, in the
sense of Definition~\ref{5.6:def-localization-links}, with $X\cap Z\neq\emptyset$. Let
$Y_1\subset Y_2\subset Y_3\subset Y_4$ be its domains and $(Y_2',Y_3',Y_4')$ a fixed triple of link domains,
as in that definition. The polydiscs $\mathbb{D}_i$, $P_i$ ($i=2,3,4$), $\mathbb{D}$, $P$ are those of
Definition~\ref{5.6:def-link-readings}, with the index sets $\mathcal{V}_i$ fixed there. The datum space
$\mathcal{D}$, the data $\mathfrak{d}$, the argument space $\mathcal{U}_{\mathrm{arg}}(Y_4)$, the value
space $\mathcal{U}_{\mathrm{val}}(Y_1)$ and the analyticity domain $\mathcal{A}_t$ of $t$ are those of
Definition~\ref{5.6:not-term-spaces}.

For a set $Y$, $\mathfrak{M}(Y)$ denotes the manifold of complex configurations on $Y$ of
Notation~\ref{5.6:not-holomorphy-conventions}. An \emph{open set of $\mathfrak{M}(Y)$} means a relatively
open subset of $\mathfrak{M}(Y)$. We write $m_3,m_2\ge 0$ for the numbers of complex auxiliary arguments read
by the link maps at levels three and two, and $\mathbb{C}^{n_0}$ for the ambient coordinate space of
$\mathfrak{M}(Y_1)$. The letter $w$ denotes a point of the argument domain of a link map.

The setting for the link maps consists of the following five hypotheses.
\begin{enumerate}
\item[(a)] \emph{The link map at level four.} There are two sets with the following properties.
First, a relatively open set $\mathcal{O}_4$: it consists of configurations on $Y_4$ together with the
complex auxiliary arguments read by the representation \cite[(1.63)]{B-CMP122b}, and it contains the image
of the datum space $\mathcal{D}$, each datum $\mathfrak{d}$ being placed in these slots as in
Definition~\ref{5.6:not-term-spaces}. Second, a relatively open set
$\mathcal{O}_3\subset\mathfrak{M}(Y_3)\times\mathbb{C}^{m_3}$. The map $\mathcal{E}_4$ assigns to
$(\sigma^{(4)},w)$ the extended right-hand side of \cite[(1.63)]{B-CMP122b} on $Y_3$, built with the
localised function $\mathbf{H}^{(4)}(\sigma^{(4)};w)$. It has the property in the first line of
\eqref{5.6:eq-link-hypotheses-a} below: it is analytic in the sense of
\eqref{5.6:eq-holomorphy-conventions-a}, that is, continuous on $\mathbb{D}_4\times\mathcal{O}_4$,
holomorphic at every point of $P_4\times\mathcal{O}_4$, and $\mathfrak{M}(Y_3)$-valued in its configuration
components. Moreover, under the condition $\delta_1 M\ge\kappa_1$, the localised function
$\mathbf{H}^{(4)}$ obeys bounds of the type of \cite[(1.11), (1.18), (1.21)]{B-CMP116}. These bounds enter
the composition of the link maps only through the containment
$\mathcal{E}_4(\mathbb{D}_4\times\mathcal{O}_4)\subset\mathcal{O}_3$, which is part of (a). Their type is
used in the smallness conditions on the constants at the link maps.
\item[(b)] \emph{The link map at level three.} There is a relatively open set
$\mathcal{O}_2\subset\mathfrak{M}(Y_2)\times\mathbb{C}^{m_2}$ with the following property. Consider the
map $\mathcal{E}_3$ that assigns to $(\sigma^{(3)},w)$ the extended right-hand side of
\cite[(1.62)]{B-CMP122b} on $Y_2$, built as follows: the localised function
$\mathbf{H}^{(3)}(\sigma^{(3)};w)=H_B$, the localised $H$-function of the representation
\cite[(1.62)]{B-CMP122b}, is inserted at $w$ both as the background field and in the argument,
and the remaining fields in \cite[(1.62)]{B-CMP122b} are replaced by functions of $\mathbf{U}$. This map has
the property in the second line of \eqref{5.6:eq-link-hypotheses-a}: it is analytic in the sense of
\eqref{5.6:eq-holomorphy-conventions-a} with values in $\mathcal{O}_2$. Its localised function obeys bounds
of the same type as in (a).
\item[(c)] \emph{The link map at level two.} Consider the map $\mathcal{E}_2$ that assigns to
$(\sigma^{(2)},w)$ the right-hand side of \cite[(1.61)]{B-CMP122b} on $Y_1$, built with three ingredients:
\begin{itemize}
\item the localised function $\mathbf{H}^{(2)}(\sigma^{(2)};w)$;
\item the kept gauge transformation $u_k$, evaluated on the restriction of $\mathbf{H}^{(2)}$ to the block
of each point;
\item the variables $(\mathbf{U},\mathbf{J})$ in the other expressions of \cite[(1.61)]{B-CMP122b}.
\end{itemize}
This map has the property in the third line of \eqref{5.6:eq-link-hypotheses-a}: it is analytic in the sense
of \eqref{5.6:eq-holomorphy-conventions-a} as an $\mathfrak{M}(Y_1)$-valued map, that is, as a
$\mathbb{C}^{n_0}$-valued map. Here \cite[text preceding (1.61)]{B-CMP122b} states that $u_k(x)$ is
an analytic function of the restriction to the block of $x$ of the $H$-function of that representation.
That function is $\mathbf{H}^{(2)}$ by the identification made in
Definition~\ref{5.6:def-link-readings}. The requirement that the nested composite take its values in
$\mathcal{A}_t$ is stated separately as (d).
\end{enumerate}
In summary, (a)--(c) require
\begin{equation}\label{5.6:eq-link-hypotheses-a}
\begin{aligned}
&\mathcal{E}_4:\mathbb{D}_4\times\mathcal{O}_4\to\mathcal{O}_3,
  &(\sigma^{(4)},w)&\mapsto\mathcal{E}_4(\sigma^{(4)};w),\\
&\mathcal{E}_3:\mathbb{D}_3\times\mathcal{O}_3\to\mathcal{O}_2,
  &(\sigma^{(3)},w)&\mapsto\mathcal{E}_3(\sigma^{(3)};w),\\
&\mathcal{E}_2:\mathbb{D}_2\times\mathcal{O}_2\to\mathfrak{M}(Y_1),
  &(\sigma^{(2)},w)&\mapsto\mathcal{E}_2(\sigma^{(2)};w),
\end{aligned}
\end{equation}
each map being analytic in the sense of \eqref{5.6:eq-holomorphy-conventions-a}.
\begin{enumerate}
\item[(d)] \emph{Membership of the nested composite.} For every $\mathfrak{d}\in\mathcal{D}$ and every
$\sigma=(\sigma^{(2)},\sigma^{(3)},\sigma^{(4)})\in\mathbb{D}$,
\begin{equation}\label{5.6:eq-link-hypotheses-b}
\mathsf{R}(\sigma;\mathfrak{d})
:=\mathcal{E}_2\bigl(\sigma^{(2)};\,\mathcal{E}_3\bigl(\sigma^{(3)};\,
\mathcal{E}_4(\sigma^{(4)};\mathfrak{d})\bigr)\bigr)\in\mathcal{A}_t .
\end{equation}
\item[(e)] \emph{Analyticity and boundedness of the term.}
\begin{itemize}
\item[(i)] $\mathcal{A}_t$ is a relatively open subset of $\mathfrak{M}(Y_1)$, and
$\mathcal{U}_{\mathrm{val}}(Y_1)\subset\mathcal{A}_t$. The map $t=E(X;\cdot):\mathcal{A}_t\to\mathbb{C}$
is analytic: it is continuous, and locally the restriction of a holomorphic function of the ambient
coordinates.
\item[(ii)] There is a relatively open set $\mathcal{A}_t'$ such that
\begin{equation}\label{5.6:eq-link-hypotheses-c}
\mathcal{U}_{\mathrm{val}}(Y_1)\subset\mathcal{A}_t'\subset\mathcal{A}_t\subset\mathfrak{M}(Y_1),
\qquad
\mathfrak{b}_t:=\sup_{\mathcal{A}_t'}\,\lvert E(X;\cdot)\rvert<\infty .
\end{equation}
\end{itemize}
\end{enumerate}

\emph{Relation of (a)--(c) to Ba{\l}aban's papers.} For the expansion of the localised function
$\mathbf{H}^{(4)}$, which is the function $H_k$ of \cite[(1.63), (1.64)]{B-CMP122b},
\cite{B-CMP122b} refers to the expansion \cite[(1.10)]{B-CMP116},
read through \eqref{5.6:eq-reference-reading-a}. The construction of \cite[\S1]{B-CMP116} estimates the
terms of that expansion on the parameter domain $\lvert s(\Delta)\rvert\le e^{\kappa_1}$. The reference
names the form of the expansion; the domain $\mathbb{D}_4$ of the decoupling variables at this link is
inferred from it, and is itself a reading. The analyticity with bounds of $H_k$,
given by \cite[(1.64)]{B-CMP122b}, on $\mathbb{D}_4$ is the reading \eqref{5.6:eq-link-readings-b}. For the
function written $H_k(s(Y_0),B')$ in \cite[(1.17)]{B-CMP116}, the corresponding statement is established in
\cite[\S1]{B-CMP116}.

Likewise, at the representations \cite[(1.62)]{B-CMP122b} and \cite[(1.61)]{B-CMP122b} the domains
$\mathbb{D}_3$ and $\mathbb{D}_2$ are inferred in the same way from the reference to
\cite[(1.10)]{B-CMP116}, and are readings. The analyticity of $H_B$
on $\mathbb{D}_3$ and that of $\mathbf{H}^{(2)}$ on $\mathbb{D}_2$ are the readings
\eqref{5.6:eq-link-readings-b}.

\emph{Relation of (d) to Ba{\l}aban's papers.} \cite[(1.65)]{B-CMP122b} states the following. The extended
function $U_k$, that is, the composite of \cite[(1.61)--(1.63)]{B-CMP122b}
regarded as a function of $(\mathbf{U},\mathbf{J})$, is analytic in the space
$\mathcal{U}_{\mathrm{arg}}(Y_4)$ with values in the space $\mathcal{U}_{\mathrm{val}}(Y_1)$. The latter is
the space defined by \cite[(1.64)--(1.69)]{B-CMP122a} at the top value of the index $n$ of those spaces, as
specified in \cite[text following (1.65)]{B-CMP122b}. The same passage also states
that this space of values is contained in the analyticity domain of the term $E(X)$.
Hence $\mathsf{R}\in\mathcal{U}_{\mathrm{val}}(Y_1)\subset\mathcal{A}_t$ at the values of the decoupling
parameters at which \cite[(1.65)]{B-CMP122b} is stated; the range of these parameters is not displayed
there. Condition (d) requires this membership for every
$\sigma\in\mathbb{D}$ and every $\mathfrak{d}\in\mathcal{D}$; in this form it is a hypothesis, of the same
kind as the readings of Definition~\ref{5.6:def-link-readings}.

For the construction of \cite[\S1]{B-CMP116}, whose expansion \cite[(1.10)]{B-CMP116} is extended
analytically in the decoupling parameters, the corresponding membership is the requirement of
\cite[\S1]{B-CMP116} that the
function under consideration there, multiplied by $\mathfrak{t}\zeta_{*}+t_{\square}\zeta_{\square}$,
where $\mathfrak{t}$ denotes the real interpolation parameters of \cite[\S\S3--5]{B-CMP109}, $t_{\square}$
the outer complex parameter and $\zeta_{*}$, $\zeta_{\square}$ the cut-off functions of the device of
\cite[(3.4)]{B-CMP109},
satisfy the bounds \cite[(3.31)]{B-CMP109}; the expression \cite[(1.10)]{B-CMP116} is then extended
analytically. There the
membership is obtained by smallness conditions on constants chosen after $\kappa_1$; these conditions
absorb the factor $e^{16\kappa_1}$ in \cite[(1.21)]{B-CMP116}. For the link maps, the smallness conditions
on the constants chosen after $\kappa_1$ that concern the parameter domain and the circles of the Cauchy
integrals are stated in the same form among the conditions on the constants at the link maps; the
membership (d) itself remains a hypothesis, and no form of conditions from which it would follow is stated.

\emph{The extension in the outer parameter.} Where the outer complex parameter $t_{\square}$ of the device
of \cite[(3.4)]{B-CMP109} is used, (d) is assumed in the following strengthened form. Let
$\mathsf{R}(\sigma,t_{\square};\mathfrak{d})$ denote the nested composite extended in $t_{\square}$ as in
\eqref{5.6:eq-link-readings-c}. Then
\begin{equation}\label{5.6:eq-link-hypotheses-1}
\mathsf{R}(\sigma,t_{\square};\mathfrak{d})\in\mathcal{A}_t
\quad\text{for } \mathfrak{d}\in\mathcal{D},\ \sigma\in\mathbb{D},\
\lvert t_{\square}\rvert\le R(\mathfrak{d}).
\end{equation}
Here $R(\mathfrak{d})>0$ denotes the membership radius in the outer parameter, of the type of
\cite[(1.22)]{B-CMP116}, and
$R_{*}:=\inf_{\mathfrak{d}\in\mathcal{D}}R(\mathfrak{d})>0$ under the smallness conditions on the
constants chosen after $\kappa_1$. This strengthened membership and the positivity of $R_*$ are assumed
in addition to the reading \eqref{5.6:eq-link-readings-c}; like (d), they are hypotheses of the same kind
as the readings of Definition~\ref{5.6:def-link-readings}.

\emph{Relation of (e) to Ba{\l}aban's papers.} Clause (i) is stated in
\cite[text following (1.65)]{B-CMP122b}. The space named in that sentence is read as the value space
$\mathcal{U}_{\mathrm{val}}(Y_1)$ of \cite[(1.65)]{B-CMP122b} (its printed decorations differ from those of
the value space in \cite[(1.65)]{B-CMP122b}; the sentence continues the discussion of the smaller space of
values of the function in \cite[(1.65)]{B-CMP122b}, which fixes the sense), and analyticity is understood
in the sense of
\eqref{5.6:eq-holomorphy-conventions-a}. \cite{B-CMP122b} states the containment of the space of values in
the analyticity domain of the term $E(X)$ for three classes of terms:
\begin{itemize}
\item for the terms of the family written $\mathbf{B}''_k$ in \cite[text following (1.65)]{B-CMP122b},
from the inductive assumptions connected with
\cite[(1.70)]{B-CMP122a};
\item for the new terms created by the previous localisation operations;
\item for all the old terms.
\end{itemize}

Clause (ii) is of the type of the bounds that Ba{\l}aban's construction carries for the terms on their
domains of analyticity. These are the inductive bounds of the type of \cite[(1.68)--(1.70)]{B-CMP122a} and
\cite[(1.18)]{B-CMP109}. \cite{B-CMP122b} states that these bounds are satisfied for the analytically extended
expressions. It also states that the other factors coming from the renormalisation procedure, such as powers
of $L^j\eta$ or of $\alpha_{0,k}$, were treated earlier in the construction. The existence of a finite bound
on a set $\mathcal{A}_t'$ as in \eqref{5.6:eq-link-hypotheses-c} is the form in which this is assumed here.

Wherever $\mathcal{A}_t$ occurs in (d), in \eqref{5.6:eq-link-hypotheses-1} and in the statements proved
below from these hypotheses, it stands for $\mathcal{A}_t'$. If the bound
holds on the whole of $\mathcal{A}_t$, one takes
$\mathcal{A}_t':=\mathcal{A}_t$. No value of $\mathfrak{b}_t$ is used.
\end{definition}

\begin{remark}[Order of choice of $\kappa_1$, $M$ and the link smallness constants]
\label{5.6:rem-choice-order}
We work with the construction of \cite[\S1]{B-CMP116}. In that construction decoupling parameters
$s(\Delta)$ are introduced into the operators. The terms of the expansion \cite[(1.10)]{B-CMP116}
are then estimated by the Cauchy formula on the closed polydisc $|s(\Delta)|\le e^{\kappa_1}$,
$\Delta$ running over the cubes that carry the decoupling variables (at the links, over the sets
$\mathcal{V}_i$ of Definition~\ref{5.6:def-link-readings}).
Under the reading \eqref{5.6:eq-reference-reading-a}, \cite{B-CMP122b} applies this construction at
each link of the composite of link maps fixed in Definition~\ref{5.6:def-link-hypotheses}. The
readings \eqref{5.6:eq-link-readings-b}, the hypotheses
\eqref{5.6:eq-link-hypotheses-a} and the membership condition \eqref{5.6:eq-link-hypotheses-b} are
stated under the following conditions. Each of them is a condition of \cite[\S1]{B-CMP116}, and each
is one-sided:
\begin{equation}\label{5.6:eq-choice-order-a}
\begin{aligned}
&\text{(c0)}\quad \kappa_1 \text{ is a sufficiently big positive number};\\
&\text{(c1)}\quad \delta_1 M\ge \kappa_1;\\
&\text{(c2)}\quad 4C_2B_0e^{16\kappa_1}\varepsilon_2\le 1,\qquad
  4C_4B_0^2e^{32\kappa_1}\varepsilon_3\le 1,\qquad 2B_0e^{16\kappa_1}\varepsilon_3\le\varepsilon_2,\\
&\phantom{\text{(c2)}\quad} e^{32\kappa_1}\varepsilon_1 \text{ is smaller than an absolute constant};\\
&\text{(c3)}\quad 4B_0C_1e^{16\kappa_1}\varepsilon_1\le \mathfrak{c}_1\alpha_2;\\
&\text{(c4)}\quad g_k|B|<\varepsilon_1 .
\end{aligned}
\end{equation}

Here $\delta_1=\delta_1(\delta_0)$ is the positive constant of the walk bound
\cite[(1.11)]{B-CMP116}, and it depends only on the decay constant $\delta_0$ of the bound
\cite[(1.7)]{B-CMP116}, from which the walk bound is derived.
The numbers $\varepsilon_1,\varepsilon_2,\varepsilon_3$ are the auxiliary smallness constants of
\cite[\S1]{B-CMP116}, which bound the arguments of the functions of that construction, and
$B_0,C_1,C_2,C_4$ are the constants of those bounds. The numbers $\alpha_2$ and $\mathfrak{c}_1$ are
further auxiliary constants of \cite[\S1]{B-CMP116}, and condition (c3) is the membership condition
assumed there.

Part (c3) carries with it the analytic extension of \cite[\S1]{B-CMP116} with respect to $t_\square$
on the set
\begin{equation*}
|t_\square|\cdot 4B_0C_1e^{16\kappa_1}g_k|B|\le\alpha_2 .
\end{equation*}
Part (c3) is the membership condition entering \eqref{5.6:eq-link-hypotheses-b}.

The letter $\kappa_1$ of (c0) is the one that appears in the factor
$\exp(-(\kappa_1-1)M^{-d}|Y_i'|)$ of \cite{B-CMP122b}. Under the reading \eqref{5.6:eq-reference-reading-a},
condition (c0) is in force at every link. The lower bounds
$\kappa_1\ge 1$ and $\kappa_1\ge 16\kappa+1$ used in the proof of the correlation theorem
(Theorem~\ref{5.1:thm-correlation-theorem}) are conditions of the same kind as (c0).

The constants are chosen in the following order:
\begin{enumerate}
\item $\kappa_1$, subject only to lower bounds;
\item $M$, subject to (c1) given $\kappa_1$;
\item the smallness constants $\varepsilon_1,\varepsilon_2,\varepsilon_3$ of (c2), which govern the
domain of the parameters $s(\Delta)$ and the Cauchy circles of \cite[\S1]{B-CMP116}, given
$\kappa_1$; condition (c3) fixes no further constant, the membership condition
\eqref{5.6:eq-link-hypotheses-b} being a hypothesis;
\item last, the coupling.
\end{enumerate}
Condition (c4) concerns the ranges of the data: the small-field ranges of the auxiliary fields, on
which the estimates of \cite[\S1]{B-CMP116} are made. No size of these ranges is derived here.

This order is compatible with the order of choice of the constants in
Definition~\ref{5.1:def-admitted-inputs}, where $\kappa$ precedes $\kappa_1$, $\kappa_1$ precedes
$M_1$ and $M$, and the coupling comes last. No later choice imposes a condition on $\kappa_1$, on $M$
or on $L$.
\end{remark}

\begin{proof}
We take the conditions in the order of choice. At each stage we show that the condition can be met
once the earlier constants are fixed, and that it constrains no earlier constant.

\emph{Condition (c0).} This is a lower bound on $\kappa_1$ and nothing else; we fix $\kappa_1$
first. Any finite family of lower bounds of the same kind, among them $\kappa_1\ge1$ and
$\kappa_1\ge16\kappa+1$, is met by taking $\kappa_1$ at least their maximum. This is possible because
$\kappa$ is fixed before $\kappa_1$.

\emph{Condition (c1).} With $\kappa_1$ fixed, and with $\delta_1>0$ depending only on $\delta_0$,
condition (c1) reads $M\ge\kappa_1/\delta_1$. This is a lower bound on $M$ chosen after $\kappa_1$,
and it places no upper bound on $\kappa_1$.

Condition (c1) enters the construction as follows. In \cite[\S1]{B-CMP116}, when
$\delta_1M\ge\kappa_1$, the walk bound \cite[(1.11)]{B-CMP116} of the localised operators is at most
a constant times $e^{16\kappa_1}$ on the whole closed
polydisc $|s(\Delta)|\le e^{\kappa_1}$. Thus (c1) is a condition for the analyticity, with this bound,
on the closed polydisc. At the links it is a condition of the hypotheses
\eqref{5.6:eq-link-hypotheses-a}, entering by way of the readings \eqref{5.6:eq-link-readings-b}: that
the operators at the links obey the walk bound \cite[(1.11)]{B-CMP116} is part of those readings.

\emph{Conditions (c2).} With $\kappa_1$ and the constants $B_0,C_2,C_4$ fixed, each inequality in
(c2) has the form $f(\kappa_1)\,\varepsilon\le c$, with $f(\kappa_1)$ equal to $e^{16\kappa_1}$ or
$e^{32\kappa_1}$ times fixed constants and $c$ fixed before. Such an inequality is met by taking
$\varepsilon$ small enough, and it is an upper bound on $\varepsilon$. We choose them in turn:
\begin{itemize}
\item first $\varepsilon_2\le(4C_2B_0e^{16\kappa_1})^{-1}$;
\item then $\varepsilon_3$ with
\begin{equation*}
\varepsilon_3\le\min\bigl\{(4C_4B_0^2e^{32\kappa_1})^{-1},\ \varepsilon_2(2B_0e^{16\kappa_1})^{-1}\bigr\};
\end{equation*}
\item and $\varepsilon_1$ so small that $e^{32\kappa_1}\varepsilon_1$ lies below the absolute constant
of \cite[\S1]{B-CMP116}.
\end{itemize}

At the links, each of the localised functions $\mathbf{H}^{(4)},\mathbf{H}^{(3)},\mathbf{H}^{(2)}$ is a
function of the fixed-point type of \cite[\S1]{B-CMP116} evaluated at a small argument. According to
\cite{B-CMP122b}, these functions are very small because of the localisation of the argument, and
$H_B$ is very small on $Y_2\setminus Y_2^0$. The field $V'_\Lambda$ is in the axial gauge in
$\Lambda$, hence small inside $\Lambda$, so that the condition \cite[(1.84)]{B-CMP122a} is satisfied.
Their analyticity on $|s(\Delta)|\le e^{\kappa_1}$, with bounds inflated by $e^{16\kappa_1}$, requires,
as in \cite[\S1]{B-CMP116}, that $e^{16\kappa_1}$ or $e^{32\kappa_1}$ times the size of the argument
stays below the thresholds of the fixed-point construction. This is smallness of the arguments,
chosen after $\kappa_1$. Which quantities of the operation $\mathbf{R}_k$ play the roles of
$\varepsilon_1,\varepsilon_2,\varepsilon_3$ and of $g_k|B|$ at the links is part of the readings
\eqref{5.6:eq-link-readings-b} and of the membership condition \eqref{5.6:eq-link-hypotheses-b}.

\emph{Condition (c3).} With $\kappa_1$, $B_0$, $C_1$ and the threshold $\mathfrak{c}_1\alpha_2$
fixed, (c3) reads, in \cite[\S1]{B-CMP116},
\begin{equation*}
\varepsilon_1\le\mathfrak{c}_1\alpha_2\,(4B_0C_1e^{16\kappa_1})^{-1}.
\end{equation*}
There it is an upper bound on $\varepsilon_1$. At the links, (c3) fixes no constant in the order of
choice: the membership condition \eqref{5.6:eq-link-hypotheses-b} is a hypothesis, and it is not
deduced here from the smallness of the arguments.

\emph{Condition (c4).} The ranges of the data are the small-field ranges of the auxiliary fields.
They are supplied by the small-field characteristic functions of the integrations in the operation
$\mathbf{R}_k$, and they are bounded in terms of the coupling. The conditions (c2), read with the
ranges (c4), therefore take the form
\begin{multline*}
\bigl(\text{a function of }\kappa_1\bigr)\times\bigl(\text{a quantity tending to zero with the
coupling}\bigr)\\
\le\bigl(\text{a constant fixed before}\bigr).
\end{multline*}
For $\kappa_1$ and the earlier constants fixed, these inequalities are satisfied once the coupling
is small enough. They are upper bounds on the coupling, chosen last. No inequality of this form is
derived for the membership condition \eqref{5.6:eq-link-hypotheses-b}.

\emph{No condition bounds $\kappa_1$ from above.} Consider an inequality of the form
$f(\kappa_1)\,\varepsilon\le c$ among (c2)--(c4), with $f(\kappa_1)$ as large as $e^{32\kappa_1}$
times a constant. Here $\kappa_1$ is fixed first and may be taken as
large as the lower bounds require. Then $M\ge\kappa_1/\delta_1$ is chosen. Then the smallness
constants, and with them the admissible range of the coupling, are taken small enough. This is
possible for every given $\kappa_1$, because they are chosen after $\kappa_1$. Hence a condition
$f(\kappa_1)\,\varepsilon\le c$, with $\varepsilon$ chosen after $\kappa_1$, is an upper bound on
$\varepsilon$, that is, on the coupling. It is not an upper bound on $\kappa_1$.

Three nested links compound the inflation factor to at most $e^{48\kappa_1}$, so the quantity to be
kept below the thresholds is at most $e^{48\kappa_1}$ times the size of the argument. This is still an
upper bound on the coupling, chosen last. Every choice in the order above
is constrained only by constants fixed before it, so none returns to $\kappa_1$, to $M$ or to $L$.
\end{proof}

\begin{lemma}[From the closed to the open polydisc without loss of radius]\label{5.6:lem-open-polydisc-bis}
For a finite set $W$ and $\rho>0$ write
\begin{equation*}
P_W(\rho):=\{\zeta\in\mathbb{C}^W:\ |\zeta_w|<\rho\ \text{for all } w\in W\},\qquad
\mathbb{D}_W(\rho):=\{\zeta\in\mathbb{C}^W:\ |\zeta_w|\le\rho\ \text{for all } w\in W\}
\end{equation*}
for the open and the closed polydisc of radius $\rho$. Let $W$ be finite, $\rho>0$, $E$ a topological
space, $V\subset\mathbb{C}^{n'}$ any set, and let $G:\mathbb{D}_W(\rho)\times E\to V$ be analytic in
the following sense: $G$ is continuous on $\mathbb{D}_W(\rho)\times E$ for the product topology and,
for each $e\in E$, the map $G(\cdot,e)$ is holomorphic, as a $\mathbb{C}^{n'}$-valued map, at every
point of $P_W(\rho)$, holomorphy being understood as in
Notation~\ref{5.6:not-holomorphy-conventions}. Then $P_W(\rho)$ is the interior of
$\mathbb{D}_W(\rho)$, and for each $e\in E$ the restriction of $G(\cdot,e)$
to $P_W(\rho)$ is continuous and holomorphic on the open polydisc $P_W(\rho)$ of the same radius
$\rho$, and every bound that $G$ obeys on $\mathbb{D}_W(\rho)$ it obeys on $P_W(\rho)$.
\end{lemma}

\begin{proof}
For $\zeta\in\mathbb{C}^W$ and $r>0$ write
$P_W(\zeta;r):=\{\xi\in\mathbb{C}^W:\ |\xi_w-\zeta_w|<r\ \text{for all } w\in W\}$ for the open
polydisc of radius $r$ centred at $\zeta$. We first show that $P_W(\rho)$ is the interior of
$\mathbb{D}_W(\rho)$. Let $\zeta\in P_W(\rho)$ and put $r:=\min_{w\in W}(\rho-|\zeta_w|)$; since $W$ is
finite and $|\zeta_w|<\rho$ for every $w$, we have $r>0$. For $\xi\in P_W(\zeta;r)$ and every $w\in W$,
\begin{equation*}
|\xi_w|\le|\zeta_w|+|\xi_w-\zeta_w|<|\zeta_w|+r\le\rho ,
\end{equation*}
so $P_W(\zeta;r)\subset\mathbb{D}_W(\rho)$, and $P_W(\zeta;r)$ is an open neighbourhood of $\zeta$.
Hence every point of $P_W(\rho)$ is an interior point of $\mathbb{D}_W(\rho)$. Conversely, if
$\zeta\in\mathbb{D}_W(\rho)$ has $|\zeta_{w_0}|=\rho$ for some $w_0$, then $\zeta_{w_0}\neq0$, and for
every $\varepsilon>0$ the point obtained from $\zeta$ by replacing $\zeta_{w_0}$ with
$(1+\varepsilon)\zeta_{w_0}$ lies outside $\mathbb{D}_W(\rho)$, so $\zeta$ is not interior.

Fix $e\in E$ and $\zeta\in P_W(\rho)$. Since $\zeta$ is an interior point of $\mathbb{D}_W(\rho)$ and
$P_W(\rho)$ is open, the hypothesis that $G(\cdot,e)$ is holomorphic at $\zeta$ gives holomorphy of the
restriction of $G(\cdot,e)$ to $P_W(\rho)$ at $\zeta$; as $\zeta$ was arbitrary, this restriction is
holomorphic on $P_W(\rho)$. The map $\zeta\mapsto(\zeta,e)$ from $\mathbb{D}_W(\rho)$ to
$\mathbb{D}_W(\rho)\times E$ is continuous for the product topology, so $G(\cdot,e)$ is continuous on
$\mathbb{D}_W(\rho)$, and its restriction to the subset $P_W(\rho)$ is continuous. Finally, a bound
that holds at every point of $\mathbb{D}_W(\rho)$ holds in particular at every point of the subset
$P_W(\rho)$.
\end{proof}

\begin{lemma}[Nested composition of polydisc-analytic maps]\label{5.6:lem-nested-composition}
Let $\lambda\ge 1$, let $W_1,\dots,W_\lambda$ be finite, pairwise disjoint index sets, and let
$\rho>0$. For $i=1,\dots,\lambda$ put $\mathbb{D}^{(i)}:=\mathbb{D}_{W_i}(\rho)$ and
$P^{(i)}:=P_{W_i}(\rho)$. These are the closed and the open polydisc of radius $\rho$ in
$\mathbb{C}^{W_i}$, in the notation of Notation~\ref{5.6:not-holomorphy-conventions}.

Let $\mathcal{O}_\lambda,\mathcal{O}_{\lambda-1},\dots,\mathcal{O}_1,\mathcal{O}_0$ be sets with
the following properties. For $i=0,\dots,\lambda-1$, the set $\mathcal{O}_i$ is relatively open in
a closed complex submanifold $\mathfrak{M}_i$ of some $\mathbb{C}^{n_i}$. The set
$\mathcal{O}_\lambda$ may be any topological space; it is the space of data.

For $i=\lambda,\lambda-1,\dots,1$ let
\begin{equation*}
\mathcal{E}_i\colon \mathbb{D}^{(i)}\times\mathcal{O}_i\to\mathcal{O}_{i-1}
\end{equation*}
be analytic in the sense of \eqref{5.6:eq-holomorphy-conventions-a}. This means:
\begin{itemize}
\item[(a)] $\mathcal{E}_i$ is continuous on $\mathbb{D}^{(i)}\times\mathcal{O}_i$;
\item[(b)] for each $w\in\mathcal{O}_i$, the map $\zeta\mapsto\mathcal{E}_i(\zeta;w)$ is
holomorphic on $P^{(i)}$ as a $\mathbb{C}^{n_{i-1}}$-valued map;
\item[(c)] for $i\le\lambda-1$, where $\mathcal{O}_i\subset\mathfrak{M}_i\subset\mathbb{C}^{n_i}$,
the map $\mathcal{E}_i$ is jointly holomorphic on $P^{(i)}\times\mathcal{O}_i$ in the following
sense. Near each point $(\zeta,w)\in P^{(i)}\times\mathcal{O}_i$, $\mathcal{E}_i$ is the
restriction of a holomorphic map defined on the product of a neighbourhood of $\zeta$ in
$P^{(i)}$ and a neighbourhood of $w$ in $\mathbb{C}^{n_i}$.
\end{itemize}

Let $g\colon\mathcal{O}_0\to\mathbb{C}$ be analytic in the sense of
\eqref{5.6:eq-holomorphy-conventions-a}, and suppose
$\sup_{\mathcal{O}_0}|g|\le\mathfrak{b}$.

For $w\in\mathcal{O}_\lambda$ and
$\zeta=(\zeta_1,\dots,\zeta_\lambda)\in\prod_i\mathbb{D}^{(i)}$ put
\begin{equation}\label{5.6:eq-nested-composition-1}
\begin{gathered}
\mathsf{R}(\zeta_1,\dots,\zeta_\lambda;w)
:=\mathcal{E}_1\bigl(\zeta_1;\mathcal{E}_2\bigl(\zeta_2;\dots\mathcal{E}_\lambda(\zeta_\lambda;w)
\dots\bigr)\bigr)\in\mathcal{O}_0,\\
F(\zeta;w):=g\bigl(\mathsf{R}(\zeta;w)\bigr).
\end{gathered}
\end{equation}

Then, for every $w\in\mathcal{O}_\lambda$, the following hold.
\begin{itemize}
\item[(i)] $F(\cdot\,;w)$ is continuous on $\prod_i\mathbb{D}^{(i)}$.
\item[(ii)] $F(\cdot\,;w)$ is continuous and separately holomorphic, hence jointly holomorphic,
on the open polydisc
\begin{equation*}
\prod_{i=1}^{\lambda}P^{(i)}=P_{W_1\sqcup\dots\sqcup W_\lambda}(\rho).
\end{equation*}
This polydisc has radius exactly $\rho$ in all variables.
\item[(iii)] $|F(\zeta;w)|\le\mathfrak{b}$ for all $\zeta\in\prod_i\mathbb{D}^{(i)}$.
\end{itemize}
\end{lemma}

\begin{proof}
Fix $w\in\mathcal{O}_\lambda$ throughout.

\emph{Step 1 (well-definedness and the bound).} By hypothesis the ranges satisfy
\begin{equation*}
\mathcal{E}_\lambda\bigl(\mathbb{D}^{(\lambda)}\times\mathcal{O}_\lambda\bigr)
\subset\mathcal{O}_{\lambda-1},\ \dots,\
\mathcal{E}_1\bigl(\mathbb{D}^{(1)}\times\mathcal{O}_1\bigr)\subset\mathcal{O}_0 .
\end{equation*}
Let $\zeta\in\prod_i\mathbb{D}^{(i)}$. Then $\mathcal{E}_\lambda(\zeta_\lambda;w)$ is defined and
lies in $\mathcal{O}_{\lambda-1}$. Hence $\mathcal{E}_{\lambda-1}(\zeta_{\lambda-1};\cdot)$ may be
applied to it, and the value lies in $\mathcal{O}_{\lambda-2}$. Continuing down to the map
$\mathcal{E}_1$, we see that the nested value $\mathsf{R}(\zeta;w)$ of
\eqref{5.6:eq-nested-composition-1} is defined and lies in $\mathcal{O}_0$. Consequently
\begin{equation*}
|F(\zeta;w)|=|g(\mathsf{R}(\zeta;w))|\le\sup_{\mathcal{O}_0}|g|\le\mathfrak{b}
\qquad\text{for all } \zeta\in\prod_i\mathbb{D}^{(i)} .
\end{equation*}
This is assertion (iii).

\emph{Step 2 (continuity).} Put
$\mathsf{R}_\lambda(\zeta_\lambda;w):=\mathcal{E}_\lambda(\zeta_\lambda;w)$, and for
$i=\lambda-1,\dots,1$ put
\begin{equation}\label{5.6:eq-nested-composition-2}
\mathsf{R}_i(\zeta_i,\dots,\zeta_\lambda;w)
:=\mathcal{E}_i\bigl(\zeta_i;\mathsf{R}_{i+1}(\zeta_{i+1},\dots,\zeta_\lambda;w)\bigr),
\end{equation}
for $(\zeta_i,\dots,\zeta_\lambda)\in\mathbb{D}^{(i)}\times\dots\times\mathbb{D}^{(\lambda)}$.
By the argument of Step~1 these nested values are defined, $\mathsf{R}_i(\cdot\,;w)$ takes its
values in $\mathcal{O}_{i-1}$, and $\mathsf{R}=\mathsf{R}_1$.

Each $\mathcal{E}_i$ is continuous on
$\mathbb{D}^{(i)}\times\mathcal{O}_i$ by (a). The function $g$ is continuous on $\mathcal{O}_0$,
because continuity is part of analyticity in the sense of
\eqref{5.6:eq-holomorphy-conventions-a}.

The map $\zeta_\lambda\mapsto\mathsf{R}_\lambda(\zeta_\lambda;w)=\mathcal{E}_\lambda(\zeta_\lambda;w)$
is continuous on $\mathbb{D}^{(\lambda)}$, being the restriction of a continuous map to the slice
$\mathbb{D}^{(\lambda)}\times\{w\}$. Suppose, for some $i\le\lambda-1$, that
$\mathsf{R}_{i+1}(\cdot\,;w)$ is continuous on $\mathbb{D}^{(i+1)}\times\dots\times\mathbb{D}^{(\lambda)}$
with values in $\mathcal{O}_i$, and write $\zeta'=(\zeta_{i+1},\dots,\zeta_\lambda)$. Then the map
\begin{equation*}
(\zeta_i,\zeta')\mapsto
\bigl(\zeta_i,\mathsf{R}_{i+1}(\zeta';w)\bigr)
\end{equation*}
into $\mathbb{D}^{(i)}\times\mathcal{O}_i$ is continuous. Composing it with $\mathcal{E}_i$ gives
$\mathsf{R}_i(\cdot\,;w)$, which is therefore continuous. Descending to $i=1$ and composing with
$g$, we obtain that the nested composite $F(\cdot\,;w)=g\circ\mathsf{R}_1(\cdot\,;w)$ is continuous
on $\prod_i\mathbb{D}^{(i)}$. This is assertion (i).

\emph{Step 3 (induction on the nesting depth; the first link).} Recall that
$\mathsf{R}=\mathsf{R}_1$, with $\mathsf{R}_i$ as in \eqref{5.6:eq-nested-composition-2}. We prove
the following claim by descending induction on $i$.

\emph{Claim.} For $i=\lambda,\dots,1$, the map $\mathsf{R}_i(\cdot\,;w)$ is holomorphic on the
open set $P^{(i)}\times\dots\times P^{(\lambda)}$, as a $\mathbb{C}^{n_{i-1}}$-valued map, and it
takes its values in $\mathcal{O}_{i-1}$.

For $i=\lambda$ we apply Lemma~\ref{5.6:lem-open-polydisc-bis} with the data
\begin{equation*}
W:=W_\lambda,\quad E:=\mathcal{O}_\lambda,\quad V:=\mathcal{O}_{\lambda-1}\subset
\mathbb{C}^{n_{\lambda-1}},\quad G:=\mathcal{E}_\lambda .
\end{equation*}
Its hypotheses hold. By (a), the map $G$ is continuous on $\mathbb{D}_{W_\lambda}(\rho)\times E$
for the product topology. By (b), for each $e\in E$ the map $G(\cdot\,,e)$ is holomorphic at every
point of $P_{W_\lambda}(\rho)$ as a $\mathbb{C}^{n_{\lambda-1}}$-valued map.

Lemma~\ref{5.6:lem-open-polydisc-bis} therefore shows that
$\mathcal{E}_\lambda(\cdot\,;w)$, restricted to $P^{(\lambda)}$, is continuous and holomorphic on
the open polydisc $P^{(\lambda)}$. This polydisc has the same radius $\rho$. The values lie in
$\mathcal{O}_{\lambda-1}$ by Step~1. This proves the claim for $i=\lambda$.

\emph{Step 4 (from $i+1$ to $i$, for $i\le\lambda-1$).} Assume the claim for $i+1$. Write
$\zeta'=(\zeta_{i+1},\dots,\zeta_\lambda)$. On the open set
$P^{(i)}\times\bigl(P^{(i+1)}\times\dots\times P^{(\lambda)}\bigr)$ consider the map
\begin{equation*}
\Phi(\zeta_i,\zeta'):=\bigl(\zeta_i,\mathsf{R}_{i+1}(\zeta';w)\bigr).
\end{equation*}
Its values lie in $P^{(i)}\times\mathcal{O}_i\subset P^{(i)}\times\mathbb{C}^{n_i}$. The map
$\Phi$ is holomorphic, because its components are the identity in $\zeta_i$ and the map
$\mathsf{R}_{i+1}(\cdot\,;w)$. The latter is holomorphic as a $\mathbb{C}^{n_i}$-valued map by the
induction hypothesis.

Fix a point $(\zeta_i^0,\zeta'^0)$ of the domain, and put
$w_i^0:=\mathsf{R}_{i+1}(\zeta'^0;w)\in\mathcal{O}_i$. By the joint holomorphy hypothesis (c),
there are a neighbourhood $\mathcal{N}$ of $\zeta_i^0$ in $P^{(i)}$, a neighbourhood
$\mathcal{N}'$ of $w_i^0$ in $\mathbb{C}^{n_i}$, and a holomorphic map $\tilde{\mathcal{E}}_i$ on
$\mathcal{N}\times\mathcal{N}'$ that agrees with $\mathcal{E}_i$ on
$(\mathcal{N}\times\mathcal{N}')\cap(P^{(i)}\times\mathcal{O}_i)$.

Since $\Phi$ is continuous, $\Phi^{-1}(\mathcal{N}\times\mathcal{N}')$ is an open neighbourhood of
$(\zeta_i^0,\zeta'^0)$. On this neighbourhood $\tilde{\mathcal{E}}_i\circ\Phi$ is a composite of
holomorphic maps between open sets, and is therefore holomorphic by the composition rule for
holomorphic maps of Notation~\ref{5.6:not-holomorphy-conventions}.

This composite equals $\mathsf{R}_i(\cdot\,;w)$ there. Indeed, $\Phi$ takes its values in
$P^{(i)}\times\mathcal{O}_i$, where $\tilde{\mathcal{E}}_i$ agrees with $\mathcal{E}_i$. So
$\tilde{\mathcal{E}}_i\circ\Phi=\mathcal{E}_i\circ\Phi=\mathsf{R}_i(\cdot\,;w)$ by
\eqref{5.6:eq-nested-composition-2}.

Holomorphy is a local property, so $\mathsf{R}_i(\cdot\,;w)$ is holomorphic on
$P^{(i)}\times\dots\times P^{(\lambda)}$ as a $\mathbb{C}^{n_{i-1}}$-valued map. Its values lie in
$\mathcal{O}_{i-1}$ by Step~1. This proves the claim for $i$, and completes the induction.

\emph{Step 5 (conclusion).} By the claim for $i=1$, the map $\mathsf{R}=\mathsf{R}_1$ is
holomorphic on $\prod_iP^{(i)}$ with values in $\mathcal{O}_0\subset\mathfrak{M}_0$. The function
$g$ is analytic on the relatively open set $\mathcal{O}_0$ in the sense of
\eqref{5.6:eq-holomorphy-conventions-a}. So near each point of $\mathcal{O}_0$ it is the
restriction of a holomorphic function on an open subset of $\mathbb{C}^{n_0}$.

By the composition clause for analytic functions on the complex configuration manifold in
Notation~\ref{5.6:not-holomorphy-conventions}, the map $g\circ\mathsf{R}=F(\cdot\,;w)$ is
holomorphic on $\prod_iP^{(i)}$. That clause is applied to these local ambient extensions of $g$.
It is well defined because the extensions agree with $g$ on $\mathcal{O}_0$, which contains
$\mathsf{R}\bigl(\prod_iP^{(i)};w\bigr)$.

By the equivalences of Osgood's lemma in
Notation~\ref{5.6:not-holomorphy-conventions}, a holomorphic map is continuous and separately
holomorphic. This gives assertion (ii) on $\prod_iP^{(i)}$.

The product $\prod_iP^{(i)}$ of the open polydiscs of radius $\rho$ in the disjoint index sets
$W_i$ is, by definition, the open polydisc $P_{W_1\sqcup\dots\sqcup W_\lambda}(\rho)$ of radius
$\rho$ in every variable. Lemma~\ref{5.6:lem-open-polydisc-bis} produced the open polydisc of the same
radius $\rho$ at the first link, and Steps~4 and~5 did not shrink any domain. Hence the radius in
(ii) is exactly $\rho$ in all variables. Together with Steps~1 and~2, this proves the lemma.
\end{proof}

\begin{lemma}[Nested composition with the range condition on the composite]
\label{5.6:lem-composition-variant}
Let $\lambda\ge 1$, $W_1,\dots,W_\lambda$, $\rho$,
$\mathbb{D}^{(i)}=\mathbb{D}_{W_i}(\rho)$, $P^{(i)}=P_{W_i}(\rho)$, the sets
$\mathcal{O}_\lambda,\dots,\mathcal{O}_0$, the closed complex submanifolds
$\mathfrak{M}_i\subset\mathbb{C}^{n_i}$, the maps $\mathcal{E}_\lambda,\dots,\mathcal{E}_2$, the
function $g\colon\mathcal{O}_0\to\mathbb{C}$ and the bound $\mathfrak{b}$ be as in
Lemma~\ref{5.6:lem-nested-composition}. For $i=1$, replace the range hypothesis
$\mathcal{E}_1\colon\mathbb{D}^{(1)}\times\mathcal{O}_1\to\mathcal{O}_0$ of that lemma by the
following: $\mathcal{E}_1\colon\mathbb{D}^{(1)}\times\mathcal{O}_1\to\mathfrak{M}_0$ is continuous,
$\mathbb{C}^{n_0}$-valued, for each $w\in\mathcal{O}_1$ holomorphic on $P^{(1)}$, and, when
$\lambda\ge 2$, jointly holomorphic on $P^{(1)}\times\mathcal{O}_1$ in the sense of
Lemma~\ref{5.6:lem-nested-composition}; and,
with
\begin{equation*}
\mathsf{R}(\zeta_1,\dots,\zeta_\lambda;w)
:=\mathcal{E}_1\bigl(\zeta_1;\mathcal{E}_2(\zeta_2;\dots\mathcal{E}_\lambda(\zeta_\lambda;w)\dots)\bigr),
\end{equation*}
the composite satisfies, for every $w\in\mathcal{O}_\lambda$,
\begin{equation}\label{5.6:eq-composition-variant-a}
\mathsf{R}(\zeta;w)\in\mathcal{O}_0\qquad\text{for all }\zeta\in\textstyle\prod_i\mathbb{D}^{(i)}.
\end{equation}
Then the conclusion of Lemma~\ref{5.6:lem-nested-composition} holds unchanged for
$F(\zeta;w):=g(\mathsf{R}(\zeta;w))$: for every $w\in\mathcal{O}_\lambda$, $F(\cdot;w)$ is
continuous on $\prod_i\mathbb{D}^{(i)}$, continuous and separately holomorphic, hence jointly
holomorphic, on the open polydisc
$\prod_iP^{(i)}=P_{W_1\sqcup\dots\sqcup W_\lambda}(\rho)$ of radius exactly $\rho$ in all
variables, and $|F(\zeta;w)|\le\mathfrak{b}$ for all $\zeta\in\prod_i\mathbb{D}^{(i)}$.
\end{lemma}

\begin{proof}
We go through the proof of Lemma~\ref{5.6:lem-nested-composition} step by step and note
what each step uses of the range of $\mathcal{E}_1$. The ranges of
$\mathcal{E}_\lambda,\dots,\mathcal{E}_2$ are unchanged, so the inner composite
$\mathcal{E}_2(\zeta_2;\dots\mathcal{E}_\lambda(\zeta_\lambda;w)\dots)$ lies in $\mathcal{O}_1$,
and $\mathcal{E}_1$ can be applied to it.

The step giving well-definedness and the bound uses the range of $\mathcal{E}_1$ only to place
$\mathsf{R}(\zeta;w)$ in $\mathcal{O}_0$ for $\zeta\in\prod_i\mathbb{D}^{(i)}$. This is exactly
\eqref{5.6:eq-composition-variant-a}. Hence $F(\zeta;w)=g(\mathsf{R}(\zeta;w))$ is defined there,
and $|F|\le\sup_{\mathcal{O}_0}|g|\le\mathfrak{b}$.

The continuity step uses only that each $\mathcal{E}_i$ is continuous on
$\mathbb{D}^{(i)}\times\mathcal{O}_i$ and that $g$ is continuous on $\mathcal{O}_0$. The
composition with $g$ is legitimate by \eqref{5.6:eq-composition-variant-a}, and nothing about
the range of $\mathcal{E}_1$ enters.

The inductive holomorphy steps define
\begin{align*}
\mathsf{R}_\lambda(\zeta_\lambda;w)&:=\mathcal{E}_\lambda(\zeta_\lambda;w),\\
\mathsf{R}_i(\zeta_i,\dots,\zeta_\lambda;w)
&:=\mathcal{E}_i\bigl(\zeta_i;\mathsf{R}_{i+1}(\zeta_{i+1},\dots,\zeta_\lambda;w)\bigr),
\qquad i=\lambda-1,\dots,1,
\end{align*}
so that $\mathsf{R}=\mathsf{R}_1$. For $i\ge 2$ they are word for word as before. When
$\lambda=1$ the map in question is $\zeta_1\mapsto\mathcal{E}_1(\zeta_1;w)$, and the base case of
the induction uses only that $\mathcal{E}_1(\cdot;w)$ is holomorphic on $P^{(1)}$ as a
$\mathbb{C}^{n_0}$-valued map. When $\lambda\ge 2$, the step for $i=1$ asserts holomorphy of
$(\zeta_1,\zeta')\mapsto\mathcal{E}_1(\zeta_1;\mathsf{R}_2(\zeta';w))$ on $\prod_iP^{(i)}$ as a
$\mathbb{C}^{n_0}$-valued map, where
$\zeta'=(\zeta_2,\dots,\zeta_\lambda)\in P^{(2)}\times\dots\times P^{(\lambda)}$; this map is a
composite of holomorphic maps:
$(\zeta_1,\zeta')\mapsto(\zeta_1,\mathsf{R}_2(\zeta';w))$ maps holomorphically into
$P^{(1)}\times\mathcal{O}_1$, and $\mathcal{E}_1$ is, near each point of $P^{(1)}\times\mathcal{O}_1$,
the restriction of a holomorphic map on an open subset of $P^{(1)}\times\mathbb{C}^{n_1}$. In
neither case is the range of $\mathcal{E}_1$ used. The location of the values of $\mathsf{R}_1$
in $\mathcal{O}_0$ is taken, as in that proof, from the step giving well-definedness, hence from
\eqref{5.6:eq-composition-variant-a}.

The concluding step composes the holomorphic map $\mathsf{R}$ on $\prod_iP^{(i)}$ with $g$ by
means of local ambient holomorphic extensions of $g$. For this it needs
$\mathsf{R}(\prod_iP^{(i)};w)\subset\mathcal{O}_0\subset\mathfrak{M}_0$. This follows from
\eqref{5.6:eq-composition-variant-a}, since $\prod_iP^{(i)}\subset\prod_i\mathbb{D}^{(i)}$. Thus
$g\circ\mathsf{R}(\cdot;w)$ is holomorphic on $\prod_iP^{(i)}$, hence continuous and separately
holomorphic there. Together with the bound and the continuity established above, this is the
assertion, at radius $\rho$ exactly in every variable.
\end{proof}

\begin{remark}[Joint holomorphy in parameter and argument at each link]\label{5.6:rem-joint-holomorphy}
Lemma~\ref{5.6:lem-nested-composition} (nested composition of polydisc-analytic maps) assumes
that each link map $\mathcal{E}_i$ is jointly holomorphic on $P^{(i)}\times\mathcal{O}_i$, that
is, holomorphic in the decoupling variables $\sigma^{(i)}$ and in the argument $w\in\mathcal{O}_i$
together. This hypothesis is part of what the per-link analyticity readings
\eqref{5.6:eq-link-readings-b} of Definition~\ref{5.6:def-link-readings} assert. Suppose instead
that $\mathcal{E}_i$ were only holomorphic in $\sigma^{(i)}$ for each fixed $w$ and continuous
jointly. That would not suffice for the step of the proof of
Lemma~\ref{5.6:lem-nested-composition} at which the links are composed, unless theorems of
Hartogs type were invoked. The joint form is therefore assumed, and the readings state it as such.
\end{remark}

\begin{proof}
We first compare the readings with the hypothesis. Each of the readings
\eqref{5.6:eq-link-readings-b}, for $i\in\{2,3,4\}$, concerns the function of link $i$. It is
regarded as a function of the decoupling variables $\sigma^{(i)}=(s(\Delta))_{\Delta\in\mathcal{V}_i}$
and of its argument and background fields. The reading states that this function is analytic
jointly in $\sigma^{(i)}\in\mathbb{D}_i$ and in that argument on the open argument domain. This
is the hypothesis of Lemma~\ref{5.6:lem-nested-composition} for the map $\mathcal{E}_i$.

The same joint form is the form in which the terms arise in the construction of
\cite[\S1]{B-CMP116}, the construction whose analysis the readings apply. There the terms are
analytic functions of the parameters $s(Y_0)$ together with a field argument. In the notation of
\cite[\S1]{B-CMP116} that field argument is $B'$, or $B$, or $A'$. The analyticity is thus joint
in the parameter and in the field argument.

At the site (T) of Definition~\ref{5.2:def-sites-units}, joint holomorphy in parameter and
field needs no separate hypothesis. In the corresponding step of the analysis at that site, the
field argument enters complex-linearly once the function $H_k$ has been formed.

At the site (L) the situation is different. For $i<\lambda$, the argument of $\mathcal{E}_i$ is
the whole extended configuration produced by the previous link. That configuration is the value
\begin{equation*}
\mathcal{E}_{i+1}\bigl(\sigma^{(i+1)},w'\bigr)\in\mathcal{O}_i,
\qquad w'\in\mathcal{O}_{i+1},
\end{equation*}
and it depends on the variables $\sigma^{(i+1)}$ of link $i+1$. The nested composite contains
\begin{equation*}
\mathcal{E}_i\bigl(\sigma^{(i)},\mathcal{E}_{i+1}(\sigma^{(i+1)},w')\bigr).
\end{equation*}
In this expression the variables $\sigma^{(i+1)}$ enter $\mathcal{E}_i$ only through its argument
slot. For the composite to be holomorphic in $(\sigma^{(i)},\sigma^{(i+1)})$, the map
$\mathcal{E}_i$ must therefore be holomorphic in $(\sigma^{(i)},w)$ jointly.

Holomorphy of $\mathcal{E}_i$ in $\sigma^{(i)}$ for each fixed $w$ controls the first slot only.
Continuity in $w$ adds no complex differentiability in the second slot. To pass from separate
properties of this kind to joint holomorphy one would need theorems of Hartogs type, and no such
theorem is used in the proof of Lemma~\ref{5.6:lem-nested-composition}. Hence the joint form
enters as a hypothesis, and it is carried by the readings \eqref{5.6:eq-link-readings-b}.
\end{proof}

\begin{theorem}[Holomorphy of the link pieces in the decoupling parameters]\label{5.6:thm-link-holomorphy}
Work at step $k$ of $\mathbf{R}_k$ in the setting of Definition~\ref{5.6:def-link-hypotheses}. Consider a first-part term $t=E(X;\cdot)$ with $X\cap Z\neq\emptyset$, its nested domains $Y_1\subset Y_2\subset Y_3\subset Y_4$ and the fixed link domains $(Y_2',Y_3',Y_4')$. The large-field papers' references of the form ``(6.$n$) [I]'' in \cite{B-CMP122a,B-CMP122b} are read as the items (1.$n$) of \cite[\S1]{B-CMP116}, as in \eqref{5.6:eq-reference-reading-a} of Remark~\ref{5.6:rem-reference-reading}. For $i\in\{2,3,4\}$ let $\sigma^{(i)}=(s(\Delta))_{\Delta\in\mathcal{V}_i}$ be the decoupling variables of link $i$, and put
\begin{equation*}
\mathcal{V}(t):=\mathcal{V}_2\sqcup\mathcal{V}_3\sqcup\mathcal{V}_4,\qquad
\mathbb{D}:=\mathbb{D}_2\times\mathbb{D}_3\times\mathbb{D}_4=\mathbb{D}_{\mathcal{V}(t)}(e^{\kappa_1}),\qquad
P:=P_{\mathcal{V}(t)}(e^{\kappa_1}).
\end{equation*}
Here $P_W(\rho)=\{\zeta\in\mathbb{C}^W:|\zeta_w|<\rho\text{ for all }w\}$ is the open polydisc and $\mathbb{D}_W(\rho)$ is the corresponding closed polydisc. For $\sigma=(\sigma^{(2)},\sigma^{(3)},\sigma^{(4)})\in\mathbb{D}$ and a datum $\mathfrak{d}\in\mathcal{D}$ put
\begin{equation*}
\mathsf{R}(\sigma;\mathfrak{d}):=
\mathcal{E}_2\bigl(\sigma^{(2)};\mathcal{E}_3\bigl(\sigma^{(3)};\mathcal{E}_4(\sigma^{(4)};\mathfrak{d})\bigr)\bigr),
\qquad
F_t(\sigma;\mathfrak{d}):=E\bigl(X;\mathsf{R}(\sigma;\mathfrak{d})\bigr).
\end{equation*}
Assume the following.
\begin{itemize}
\item[(a)] The hypotheses \eqref{5.6:eq-link-hypotheses-a} on the link maps $\mathcal{E}_4$, $\mathcal{E}_3$, $\mathcal{E}_2$ hold. These are the per-link analyticity readings \eqref{5.6:eq-link-readings-b} of Definition~\ref{5.6:def-link-readings}, taken with the joint holomorphy in parameter and argument of Remark~\ref{5.6:rem-joint-holomorphy}.
\item[(b)] The membership hypothesis \eqref{5.6:eq-link-hypotheses-b} holds.
\item[(c)] The hypothesis \eqref{5.6:eq-link-hypotheses-c} on the analyticity of the term $E(X;\cdot)$ on its analyticity domain $\mathcal{A}_t$, and on its bound $\mathfrak{b}_t$ on the sub-domain $\mathcal{A}_t'\subset\mathcal{A}_t$ of \eqref{5.6:eq-link-hypotheses-c}, holds.
\item[(d)] The first four conditions displayed in \eqref{5.6:eq-choice-order-a} (Remark~\ref{5.6:rem-choice-order}) hold: the floor on $\kappa_1$ (which includes $\kappa_1\ge 1$), the floor $\delta_1M\ge\kappa_1$ on $M$, the smallness ceilings on the arguments of the link maps, and the membership ceiling; the fifth, on the ranges of the data, is not used. These are the conditions under which (a) and (b) are stated.
\end{itemize}
Then for every datum $\mathfrak{d}\in\mathcal{D}$ the following hold.
\begin{itemize}
\item[(i)] $\mathsf{R}(\sigma;\mathfrak{d})\in\mathcal{A}_t$ for all $\sigma\in\mathbb{D}$. The function $F_t(\cdot;\mathfrak{d})$ is continuous on $\mathbb{D}$. On the open polydisc $P$, of radius exactly $e^{\kappa_1}$ in all the decoupling variables $s(\Delta)$, $\Delta\in\mathcal{V}(t)$, of the term, it is continuous and separately holomorphic, hence jointly holomorphic. Moreover $|F_t(\sigma;\mathfrak{d})|\le\mathfrak{b}_t$ for all $\sigma\in\mathbb{D}$.
\item[(ii)] Consequently the following condition holds for the term $t$ at the site (L): $\kappa_1\ge1$, and for every $\mathfrak{d}\in\mathcal{D}_b$,
\begin{equation}\label{5.6:eq-link-holomorphy-1}
\begin{gathered}
\sigma\mapsto F_t(\sigma;\mathfrak{d})\ \text{is continuous and separately holomorphic on }
P_{\mathcal{V}(t)}(e^{\kappa_1}),\\
\sup_{\sigma\in P_{\mathcal{V}(t)}(e^{\kappa_1})}|F_t(\sigma;\mathfrak{d})|\le\mathfrak{b}.
\end{gathered}
\end{equation}
This holds with $\mathcal{D}_b:=\mathcal{D}$, $F_t=t\circ\mathsf{R}$ and $\mathfrak{b}:=\mathfrak{b}_t$. Nothing is required in \eqref{5.6:eq-link-holomorphy-1} of the dependence on $\mathfrak{d}$.
\item[(iii)] Assume in addition the reading \eqref{5.6:eq-link-readings-c} of Definition~\ref{5.6:def-link-readings}. Assume also that the datum carries one further complex parameter $t_{\square}$, entering through the outermost factor of \eqref{5.6:eq-link-readings-c}, which is holomorphic jointly in $t_{\square}$ and in the remaining variables on $\{|t_{\square}|<R(\mathfrak{d})\}$ times open sets; in that factor $t_{\square}$ enters through the quantity $(\mathfrak{t}\zeta_{*}+t_{\square}\zeta_{\square})$ times a localised H-function, with $\mathfrak{t}$, $\zeta_{*}$, $\zeta_{\square}$ held fixed, and this quantity is affine in $t_{\square}$. Let $\mathsf{R}(\sigma,t_{\square};\mathfrak{d})$ denote the composite so extended, and assume the membership
\begin{equation*}
\mathsf{R}(\sigma,t_{\square};\mathfrak{d})\in\mathcal{A}_t
\quad\text{for all }\sigma\in\mathbb{D},\ |t_{\square}|\le R(\mathfrak{d});
\end{equation*}
here $R(\mathfrak{d})>0$ is a membership radius of the kind of \cite[(1.22)]{B-CMP116}, and it is part of the assumption that $R_{*}:=\inf_{\mathcal{D}}R(\mathfrak{d})>0$. Then the function
\begin{equation*}
F(\sigma,t_{\square};\mathfrak{d}):=E\bigl(X;\mathsf{R}(\sigma,t_{\square};\mathfrak{d})\bigr)
\end{equation*}
is jointly holomorphic on $P\times\{|t_{\square}|<R(\mathfrak{d})\}$, with $|F|\le\mathfrak{b}_t$ there. The derivative piece
\begin{equation}\label{5.6:eq-link-holomorphy-2}
F_t^{\square}(\sigma;\mathfrak{d}):=\partial_{t_{\square}}F(\sigma,t_{\square};\mathfrak{d})\big|_{t_{\square}=0},
\end{equation}
a piece of the kind occurring in the construction of \cite[\S1]{B-CMP116}, satisfies the condition of clause (ii), with $F_t$ replaced by $F_t^{\square}$ in \eqref{5.6:eq-link-holomorphy-1}, with $\mathcal{D}_b:=\mathcal{D}$ and with $\mathfrak{b}:=\mathfrak{b}_t/R_{*}$. For each datum separately the bound is $\mathfrak{b}_t/R(\mathfrak{d})$.
\item[(iv)] The radius of the open polydisc in (i)--(iii) is the radius $e^{\kappa_1}$ of the closed polydiscs of the hypotheses, so no radius is lost. The conclusions use no value of $\mathfrak{b}_t$, of $R(\mathfrak{d})$, of any threshold in \cite[(1.64)--(1.69)]{B-CMP122a}, or of any other constant.
\end{itemize}
\end{theorem}

\begin{proof}
\emph{Clauses \textup{(i)} and \textup{(ii)}.} We apply the nested composition lemma, Lemma~\ref{5.6:lem-nested-composition}, in the form of Lemma~\ref{5.6:lem-composition-variant}. In that form the range condition on the lemma's first map (here $\mathcal{E}_2$, the outermost link) is replaced by the condition \eqref{5.6:eq-composition-variant-a} that the composite lies in the set at index zero. The objects of the lemma are identified as follows.
\begin{itemize}
\item[$\cdot$] Number of links and radius. There are three links, and the common radius is $\rho:=e^{\kappa_1}$.
\item[$\cdot$] Index sets. The lemma's index sets, in its order $1,2,3$, are $\mathcal{V}_2,\mathcal{V}_3,\mathcal{V}_4$. Thus the lemma's link index $1,2,3$ corresponds to links $2,3,4$ here. Its closed polydiscs are $\mathbb{D}_2,\mathbb{D}_3,\mathbb{D}_4$ and its open polydiscs are $P_2,P_3,P_4$. The product of its closed polydiscs is $\mathbb{D}$ and the product of its open ones is $P$.
\item[$\cdot$] Maps. The lemma's first, second and third maps are $\mathcal{E}_2$, $\mathcal{E}_3$ and $\mathcal{E}_4$ of hypothesis (a).
\item[$\cdot$] Domains. The lemma's datum space, at index three, is $\mathcal{D}$ with the discrete topology; one may equally take the single point $\{\mathfrak{d}\}$. On it $\mathcal{E}_4$ is evaluated through the data's image in its argument domain. The lemma's set at index two is the relatively open target set of $\mathcal{E}_4$ in \eqref{5.6:eq-link-hypotheses-a}, which is the argument domain of $\mathcal{E}_3$. Its set at index one is the relatively open target set of $\mathcal{E}_3$, which is the argument domain of $\mathcal{E}_2$. Its set at index zero is $\mathcal{A}_t$, inside the ambient manifold $\mathfrak{M}(Y_1)$.
\item[$\cdot$] Function and bound. The function composed last is $E(X;\cdot)$, with bound $\mathfrak{b}_t$.
\end{itemize}
The bound $\mathfrak{b}_t$ of (c) holds on the relatively open subdomain $\mathcal{A}_t'$ with $\mathcal{U}_{\mathrm{val}}(Y_1)\subset\mathcal{A}_t'\subset\mathcal{A}_t$, with $\mathcal{A}_t':=\mathcal{A}_t$ when the bound holds on all of $\mathcal{A}_t$; the set at index zero is taken to be $\mathcal{A}_t'$. Accordingly membership in $\mathcal{A}_t$, in (b) and in clause (i), means membership in $\mathcal{A}_t'$, as is stipulated in \eqref{5.6:eq-link-hypotheses-c}.

We now verify the hypotheses of Lemma~\ref{5.6:lem-composition-variant} under these identifications.

\emph{Continuity of the third map.} Since $\mathcal{D}$ carries the discrete topology, continuity on $\mathbb{D}_4\times\mathcal{D}$ is the same as continuity of $\sigma^{(4)}\mapsto\mathcal{E}_4(\sigma^{(4)};\mathfrak{d})$ on $\mathbb{D}_4$ for each fixed $\mathfrak{d}$. This follows from the continuity of $\mathcal{E}_4$ on $\mathbb{D}_4$ times its argument domain, which is part of \eqref{5.6:eq-link-hypotheses-a}. Condition \eqref{5.6:eq-link-holomorphy-1} asks nothing of the dependence on $\mathfrak{d}$, so this choice of topology costs nothing in (ii).

\emph{Ranges of the second and third maps.} The hypotheses in \eqref{5.6:eq-link-hypotheses-a} state that $\mathcal{E}_4$ maps $\mathbb{D}_4$ times its argument domain into the argument domain of $\mathcal{E}_3$. They also state that $\mathcal{E}_3$ maps $\mathbb{D}_3$ times that domain into the argument domain of $\mathcal{E}_2$. These are the range hypotheses of the lemma at indices three and two.

\emph{The first map.} At index one, Lemma~\ref{5.6:lem-composition-variant} requires two things. First, $\mathcal{E}_2$ must be continuous and $\mathfrak{M}(Y_1)$-valued, holomorphic on $P_2$ for each argument, and jointly holomorphic on $P_2$ times its argument domain, holomorphy being meant in the sense of \eqref{5.6:eq-holomorphy-conventions-a}. This is the hypothesis on $\mathcal{E}_2$ in \eqref{5.6:eq-link-hypotheses-a}, in the joint form of Remark~\ref{5.6:rem-joint-holomorphy}. Second, it requires $\mathsf{R}(\sigma;\mathfrak{d})\in\mathcal{A}_t$ for all $\sigma\in\mathbb{D}$. This is, word for word, the membership hypothesis \eqref{5.6:eq-link-hypotheses-b} of (b).

\emph{Analyticity of the three maps.} The analyticity hypotheses of the lemma are those of \eqref{5.6:eq-link-hypotheses-a}, read in the joint form of Remark~\ref{5.6:rem-joint-holomorphy}. Each map is continuous on its closed polydisc times its argument domain, and holomorphic at every point of its open polydisc times that domain, jointly in parameter and argument.

\emph{The function composed last.} By \eqref{5.6:eq-link-hypotheses-c}, $\mathcal{A}_t$ is a relatively open subset of $\mathfrak{M}(Y_1)$ and $E(X;\cdot)$ is analytic on it. Moreover $\sup|E(X;\cdot)|=\mathfrak{b}_t<\infty$ on $\mathcal{A}_t'$.

All hypotheses of Lemma~\ref{5.6:lem-composition-variant} are therefore met. Its conclusion, read in the identifications above, is clause (i) word for word. The passage from separate to joint holomorphy on $P$ for a continuous function is part of that conclusion.

For clause (ii), compare (i) with \eqref{5.6:eq-link-holomorphy-1} for $\mathcal{D}_b:=\mathcal{D}$ and $\mathfrak{b}:=\mathfrak{b}_t$. The floor $\kappa_1\ge1$ is part of (d). For each $\mathfrak{d}\in\mathcal{D}$, clause (i) gives that $\sigma\mapsto F_t(\sigma;\mathfrak{d})=(t\circ\mathsf{R})(\sigma;\mathfrak{d})$ is continuous and separately holomorphic on the open polydisc $P=P_{\mathcal{V}(t)}(e^{\kappa_1})$. Since $P\subset\mathbb{D}$, the bound $|F_t(\sigma;\mathfrak{d})|\le\mathfrak{b}_t$ on $\mathbb{D}$ gives the bound on $P$. Everything in \eqref{5.6:eq-link-holomorphy-1} is a statement for each $\mathfrak{d}$ separately, so (ii) follows.

\emph{Clause \textup{(iii)}.} The argument uses only the following four facts.

\emph{First fact: joint holomorphy with bound.} Apply Lemma~\ref{5.6:lem-composition-variant} as above, with one further outermost link: the factor of \eqref{5.6:eq-link-readings-c} through which $t_{\square}$ enters, holomorphic jointly in $t_{\square}$ and the remaining variables, with $t_{\square}$ entering affinely through the quantity $(\mathfrak{t}\zeta_{*}+t_{\square}\zeta_{\square})$ times a localised H-function. In place of (b), use the membership of the extended composite $\mathsf{R}(\sigma,t_{\square};\mathfrak{d})$ in $\mathcal{A}_t$ for $\sigma\in\mathbb{D}$ and $|t_{\square}|\le R(\mathfrak{d})$; the reading \eqref{5.6:eq-link-readings-c} is assumed. The lemma then gives that $F(\sigma,t_{\square};\mathfrak{d})$ is jointly holomorphic on $P\times\{|t_{\square}|<R(\mathfrak{d})\}$. It also gives the bound $|F(\sigma,t_{\square};\mathfrak{d})|\le\mathfrak{b}_t$ there, since $E(X;\cdot)$ is bounded by $\mathfrak{b}_t$ at every point of its domain reached by the composite.

\emph{Second fact: the Cauchy estimate.} Fix $\sigma\in P$ and $0<R'<R(\mathfrak{d})$. The function $\tau\mapsto F(\sigma,\tau;\mathfrak{d})$ is holomorphic on the disc $\{|\tau|<R(\mathfrak{d})\}$. Cauchy's integral formula on the circle $|\tau|=R'$ gives
\begin{equation}\label{5.6:eq-link-holomorphy-3}
F_t^{\square}(\sigma;\mathfrak{d})
=\frac{1}{2\pi i}\oint_{|\tau|=R'}\frac{F(\sigma,\tau;\mathfrak{d})}{\tau^{2}}\,d\tau .
\end{equation}
The one-variable Cauchy estimate follows:
\begin{equation*}
|F_t^{\square}(\sigma;\mathfrak{d})|
\le\frac{1}{2\pi}\cdot 2\pi R'\cdot\frac{\mathfrak{b}_t}{R'^{2}}
=\frac{\mathfrak{b}_t}{R'} .
\end{equation*}
The left side does not depend on $R'$. Letting $R'\uparrow R(\mathfrak{d})$ gives $|F_t^{\square}(\sigma;\mathfrak{d})|\le\mathfrak{b}_t/R(\mathfrak{d})$ for all $\sigma\in P$.

\emph{Third fact: holomorphy of the parameter integral.} Fix $R'$ as above. By the first fact the integrand $F(\sigma,\tau;\mathfrak{d})\tau^{-2}$ in \eqref{5.6:eq-link-holomorphy-3} is continuous on $P\times\{|\tau|=R'\}$, and for each $\tau$ it is holomorphic in $\sigma$ on $P$. The circle is compact, so the integrand is uniformly continuous on $\Gamma\times\{|\tau|=R'\}$ for every compact set $\Gamma\subset P$. Hence the integral \eqref{5.6:eq-link-holomorphy-3} is continuous in $\sigma$ on $P$.

For each variable $s(\Delta)$, with the others fixed, consider a closed triangle in the disc of that variable. By Fubini's theorem the contour integral of \eqref{5.6:eq-link-holomorphy-3} around the triangle equals the $\tau$-integral of the contour integrals of the integrand. Each of the latter vanishes by Cauchy's theorem. By Morera's theorem, $F_t^{\square}(\cdot;\mathfrak{d})$ is therefore holomorphic in each variable $s(\Delta)$ separately. Thus $F_t^{\square}(\cdot;\mathfrak{d})$ is continuous and separately holomorphic on $P$.

\emph{Fourth fact: uniformity in the datum.} Since $R(\mathfrak{d})\ge R_{*}>0$ for every $\mathfrak{d}\in\mathcal{D}$, the second fact gives
\begin{equation*}
|F_t^{\square}(\sigma;\mathfrak{d})|\le\frac{\mathfrak{b}_t}{R(\mathfrak{d})}\le\frac{\mathfrak{b}_t}{R_{*}}
\qquad(\sigma\in P,\ \mathfrak{d}\in\mathcal{D}).
\end{equation*}

Together with $\kappa_1\ge1$ from (d), these facts are the condition of clause (ii) for the derivative piece \eqref{5.6:eq-link-holomorphy-2}, that is, \eqref{5.6:eq-link-holomorphy-1} with $F_t$ replaced by $F_t^{\square}$, with $\mathcal{D}_b:=\mathcal{D}$ and $\mathfrak{b}:=\mathfrak{b}_t/R_{*}$. Datum by datum, the bound is $\mathfrak{b}_t/R(\mathfrak{d})$.

\emph{Clause \textup{(iv)}.} The passage from the closed polydiscs $\mathbb{D}_i$ of the hypotheses to the open polydisc $P$ in the conclusions keeps the radius $e^{\kappa_1}$, with every bound valid on the closed polydisc carried over to the open one. This is the radius statement of Lemma~\ref{5.6:lem-open-polydisc-bis} (from the closed to the open polydisc without loss of radius). In (iii) the circles of radius $R'$ exhaust the disc of radius $R(\mathfrak{d})$ as $R'\uparrow R(\mathfrak{d})$, so the bound involves the full radius $R(\mathfrak{d})$.

The steps of the proof are the following: the application of Lemma~\ref{5.6:lem-composition-variant}, the comparison with \eqref{5.6:eq-link-holomorphy-1}, Cauchy's formula and estimate, the holomorphy of the parameter integral, and the infimum $R_{*}$. In these steps the quantities $\mathfrak{b}_t$ and $R(\mathfrak{d})$ enter only as the symbols carried from the hypotheses into the conclusions. No threshold of \cite[(1.64)--(1.69)]{B-CMP122a} and no value of any constant is invoked.
\end{proof}

\begin{corollary}[The weight lemma for the link pieces]\label{5.6:cor-link-weight}
Fix $K$, a level $k\le K$, and a term $t=E(X;\cdot)$ of the first part of the expansion at the
localisation site $(L)$ of $\mathbf{R}_k$ treated in Theorem~\ref{5.6:thm-link-holomorphy},
with decoupling variables indexed by
\begin{equation*}
\mathcal{V}(t)=\mathcal{V}_2\sqcup\mathcal{V}_3\sqcup\mathcal{V}_4 .
\end{equation*}
Assume the following.
\begin{itemize}
\item[(0)] The setting and the hypotheses of Theorem~\ref{5.6:thm-link-holomorphy}.
\item[(1)] The per-link analyticity readings \eqref{5.6:eq-link-readings-b}, for the links
$i=2,3,4$.
\item[(2)] The membership hypothesis \eqref{5.6:eq-link-hypotheses-b}, and the analyticity
hypothesis \eqref{5.6:eq-link-hypotheses-c} on the term $E(X;\cdot)$, with its bound on
$\mathcal{A}_t'$.
\item[(3)] The conditions \eqref{5.6:eq-choice-order-a} on the order of choice of $\kappa_1$,
$M$ and the link smallness constants.
\item[(4)] The reading \eqref{5.6:eq-reference-reading-a} of the references marked [I] in
\cite{B-CMP122a} and \cite{B-CMP122b}.
\item[(5)] The published statements of \cite{B-CMP122a}, \cite{B-CMP122b}, \cite{B-CMP109},
\cite{B-CMP116} and \cite{B-CMP119} on which Theorem~\ref{5.6:thm-link-holomorphy} rests,
used as stated there.
\item[(6)] The weight lemma, parts (a)--(e), with the readings under which it is stated,
taken unchanged, with $\lambda=3$ links, and the lemma
on the per-cube factor
\begin{equation*}
r(\kappa_1)=e^{\kappa_1}\bigl(e^{\kappa_1}-1\bigr)^{-2}\le e^{-(\kappa_1-1)} .
\end{equation*}
\end{itemize}
Then the holomorphy hypothesis of the weight lemma holds for $t$: $\kappa_1\ge 1$, and
for every datum $\mathfrak{d}\in\mathcal{D}$ the function
$\sigma\mapsto F_t(\sigma;\mathfrak{d})$ is continuous and separately holomorphic on the open
polydisc $P_{\mathcal{V}(t)}(e^{\kappa_1})$, with
\begin{equation}\label{5.6:eq-link-weight-1}
|F_t(\sigma;\mathfrak{d})|\le\mathfrak{b}
\qquad\text{for all }\sigma\in P_{\mathcal{V}(t)}(e^{\kappa_1}) ,
\end{equation}
where $\mathfrak{b}=\mathfrak{b}_t$ is the uniform bound of $t$ on $\mathcal{A}_t'$ from
Theorem~\ref{5.6:thm-link-holomorphy}.
Consequently parts (a)--(e) of the weight lemma hold for the pieces $p\in\mathcal{P}(t)$ of $t$
produced by the links $2$, $3$, $4$, in the sense of the weight lemma: each such piece is
obtained by differentiating $F_t(\cdot;\mathfrak{d})$ in a set $W(p)\subset\mathcal{V}(t)$ of
decoupling variables and integrating the derivative over $[0,1]^{\mathcal{V}(t)}$ against the
measure of the weight lemma. In particular, part (a) gives
\begin{equation}\label{5.6:eq-link-weight-2}
\sup_{\mathfrak{d}\in\mathcal{D}}|p(\mathfrak{d})|
\le\mathfrak{b}\,r(\kappa_1)^{\#W(p)}
\le\mathfrak{b}\,e^{-(\kappa_1-1)\#W(p)} .
\end{equation}
\end{corollary}

\begin{proof}
We verify the holomorphy hypothesis of the weight lemma with $\mathcal{D}_b:=\mathcal{D}$. The
weight lemma uses nothing about $\mathcal{D}_b$ beyond the validity of that hypothesis for each
of its elements, so this choice is admissible.

By (0), Theorem~\ref{5.6:thm-link-holomorphy} applies; (1)--(5) are the statements from which
its hypotheses are read, so that what follows is an implication from them. Fix
$\mathfrak{d}\in\mathcal{D}$. By
part~(i) of the theorem, $\mathsf{R}(\sigma;\mathfrak{d})\in\mathcal{A}_t$ for every
$\sigma\in\mathbb{D}$. Moreover $F_t(\sigma;\mathfrak{d})=E(X;\mathsf{R}(\sigma;\mathfrak{d}))$
is continuous on $\mathbb{D}$, and it is jointly holomorphic on the open polydisc
$P=P_{\mathcal{V}(t)}(e^{\kappa_1})$ of radius exactly $e^{\kappa_1}$ in all the decoupling
variables $s(\Delta)$, $\Delta\in\mathcal{V}(t)$, of the three links. A jointly holomorphic
function is continuous and separately holomorphic, which is the sense of holomorphy that
the hypothesis requires. The bound \eqref{5.6:eq-link-weight-1} on $P$, with $\mathfrak{b}$ the
uniform bound of the term $t$ on $\mathcal{A}_t'$, is likewise part of the conclusion of
Theorem~\ref{5.6:thm-link-holomorphy}. The inequality $\kappa_1\ge 1$ follows from the
conditions \eqref{5.6:eq-choice-order-a} on $\kappa_1$. This proves the holomorphy hypothesis for
$t$ for each $\mathfrak{d}\in\mathcal{D}$, separately.

The decoupling variables form the disjoint union of the three families $\mathcal{V}_i$,
$i=2,3,4$. The configuration of the weight lemma is therefore the one with $\lambda=3$ links,
which is the case assumed in (6). The readings under which the weight lemma is stated enter
unchanged. All hypotheses of the weight lemma now hold for $t$, so its parts (a)--(e) hold for
every piece $p\in\mathcal{P}(t)$. Part (a) is the first inequality in
\eqref{5.6:eq-link-weight-2}.

For the second inequality, $\kappa_1\ge1$ and the lemma on the per-cube factor give
$0<r(\kappa_1)\le e^{-(\kappa_1-1)}$. Raising this to the power $\#W(p)\ge 0$ and multiplying
by $\mathfrak{b}\ge 0$ yields it.
\end{proof}

\begin{remark}
The bound \eqref{5.6:eq-link-weight-2} is of the same exponential form, $e^{-(\kappa_1-1)}$ per
cube, as two factors of Ba{\l}aban's construction.
\begin{itemize}
\item The factor $\exp(-(\kappa_1-1)M^{-d}|Y_i'|)$ that \cite{B-CMP122b} obtains for each link
$i=2,3,4$ after the resummations connected with a fixed domain $Y_i'$.
\item The factor of \cite[(1.24)]{B-CMP116} whose exponent carries
$|Y_0\setminus\tilde{\square}^4|$, where $Y_0$ is the domain so denoted in that display and
$\tilde{\square}^4$ is a protected cube neighbourhood.
\end{itemize}
The bound \eqref{5.6:eq-link-weight-2} is a statement about a single piece. The summations
\begin{itemize}
\item over $(Y_2',Y_3',Y_4')$ at a fixed final domain,
\item over the final domains,
\end{itemize}
are those of parts (b)--(d) of the weight lemma, which enter above by implication and are not
re-derived here, and those of \cite{B-CMP122b}. In \cite{B-CMP122b} they are referred to
``Sects.~6,7 [I]'' and ``Sects.~2,3 [III]''; the former is, by
\eqref{5.6:eq-reference-reading-a}, \cite[\S\S1--2]{B-CMP116}, and the latter is
\cite[\S\S2--3]{B-CMP119}.
\end{remark}

\begin{remark}[Keeping the first operation's decoupling parameters]\label{5.6:rem-first-operation}
In Theorem~\ref{5.6:thm-link-holomorphy} the decoupling variables are those of the three links of
the nested composite $\mathsf{R}=\mathcal{E}_2\circ\mathcal{E}_3\circ\mathcal{E}_4$; the
operation that precedes them in \cite{B-CMP122b}, called here the first operation, localises the
bracket in \cite[(1.59)]{B-CMP122b} (the case $X\subset\tilde{\square}^2$) or in
\cite[(1.60)]{B-CMP122b} (the other case) by a sum over domains $Y_0\in\mathbf{D}_k$ containing
$\square_0=\tilde{\square}^5$ (respectively $X$) and at least one of the components of
$\Lambda^{\sim}$, the operation $\sim$ being that of \cite{B-CMP122b}; \cite{B-CMP122b} resums
this sum into the new terms before one term is singled out.

Suppose that these variables are kept unsummed. Let $\mathcal{V}_1$ be the finite set of cubes
indexing the parameters $s$ of this localisation, and put
$\sigma^{(1)}=(s(\Delta))_{\Delta\in\mathcal{V}_1}$. Write
$\sigma=(\sigma^{(1)},\sigma^{(2)},\sigma^{(3)},\sigma^{(4)})$ for the full set of decoupling
variables. Let $\mathbb{D}_1$ and $P_1$ be the closed
and the open polydisc of radius $e^{\kappa_1}$ in $\mathbb{C}^{\mathcal{V}_1}$. Let
$\mathcal{E}_1$ be the link map of the first operation: the configuration at which $E(X;\cdot)$ is
evaluated inside the bracket of \cite[(1.59)]{B-CMP122b} (respectively \cite[(1.60)]{B-CMP122b}),
namely the background $U_k^0$ shifted by $\exp(i\eta H_k'(g_kB))$, with the localised function
$H_k'(g_kB)$ carrying the parameters $\sigma^{(1)}$, regarded as a function of $\sigma^{(1)}$ and
of the configuration supplied by $\mathcal{E}_2$.
Its relatively open argument domain is $\mathcal{O}_1$. Assume:
\begin{itemize}
\item[(a)] all hypotheses of Theorem~\ref{5.6:thm-link-holomorphy}, in particular the hypotheses
on the link maps $\mathcal{E}_2,\mathcal{E}_3,\mathcal{E}_4$ displayed in
\eqref{5.6:eq-link-hypotheses-a};
\item[(b)] the following hypotheses on the first operation. The map $\mathcal{E}_1$, placed
outermost in a composite of four links with $\mathcal{E}_2,\mathcal{E}_3,\mathcal{E}_4$ as the
three inner links, satisfies the hypotheses which Lemma~\ref{5.6:lem-composition-variant} places
on the outermost link. These are taken in the form \eqref{5.6:eq-composition-variant-a}, with the
range condition imposed on the four-link composite with target set $\mathcal{A}_t$:
\begin{gather*}
(\mathcal{E}_1\circ\mathcal{E}_2\circ\mathcal{E}_3\circ\mathcal{E}_4)(\sigma;\mathfrak{d})
\in\mathcal{A}_t\\
\text{for all }\sigma\in\mathbb{D}_1\times\mathbb{D}_2\times\mathbb{D}_3\times\mathbb{D}_4
\text{ and all }\mathfrak{d}\in\mathcal{D}.
\end{gather*}
It is assumed moreover that $\mathcal{E}_2(\mathbb{D}_2\times\mathcal{O}_2)\subset\mathcal{O}_1$;
this is the range hypothesis which Lemma~\ref{5.6:lem-composition-variant} places on
$\mathcal{E}_2$ as an inner link.
\end{itemize}

Then, for every datum $\mathfrak{d}\in\mathcal{D}$, the conclusions of
Theorem~\ref{5.6:thm-link-holomorphy} hold for the four-link composite
\begin{equation}\label{5.6:eq-first-operation-1}
\mathcal{E}_1\circ\mathcal{E}_2\circ\mathcal{E}_3\circ\mathcal{E}_4 ,
\end{equation}
and for $E(X;\cdot)$ evaluated at this composite,
\begin{equation*}
E\bigl(X;(\mathcal{E}_1\circ\mathcal{E}_2\circ\mathcal{E}_3\circ\mathcal{E}_4)(\sigma;\mathfrak{d})\bigr),
\end{equation*}
which takes the place of the function $F_t$ of that theorem. In these conclusions $\mathbb{D}$ is
replaced by
$\mathbb{D}_1\times\mathbb{D}_2\times\mathbb{D}_3\times\mathbb{D}_4$, and $P$ by
$P_1\times P_2\times P_3\times P_4$.

The holomorphy part of (b), namely the holomorphy of $\mathcal{E}_1$ in $\sigma^{(1)}$, is the
assumption that the analysis of \cite{B-CMP116} --- the expansion \cite[(1.9)]{B-CMP116} in the
form \cite[(1.10)]{B-CMP116}, with the estimates \cite[(1.24)--(1.29)]{B-CMP116}, which
\cite[following (1.59)]{B-CMP122b} invokes in the words ``the expansion (6.9) [I] \ldots\
(6.10) [I] \ldots\ estimated as in (6.24)--(6.29) [I]'' --- carried out there for the
$H$-function of \cite[(1.17)]{B-CMP116}, a function
of the parameters $s$ on $Y_0$ and of the auxiliary field written $B'(B)$ there, applies in the
same way to the localised function $H_k'(g_kB)$ restricted to $\square_0$.

The qualitative content of (b) is stated in \cite[following (1.59)]{B-CMP122b}: the expressions
obtained there after the first operation are described as ``depending analytically on
$H_k'(g_kB)$ restricted to the cube $\square_0$''; the field $B$ is localised in $\Lambda$, and
``on the domains with nonzero parameters $s$ the field $B$ is equal to $0$''. A further sentence,
in the discussion following \cite[(1.60)]{B-CMP122b}, reads: ``\ldots\ restricted only by the
analyticity properties of the function $H_k'$''.

This statement is not used for the pieces indexed by the triples of domains of links $2$, $3$,
$4$, nor at any place where Lemma~\ref{5.6:lem-composition-variant} is applied with three links.
It is stated here for completeness, under hypotheses (a) and (b).
\end{remark}

\begin{proof}
Apply Lemma~\ref{5.6:lem-composition-variant}, which is stated for an arbitrary finite number
$\lambda$ of links, with $\lambda=4$ and $\mathcal{E}_1$ outermost, the variables of link $i$
ranging over $\mathbb{D}_i$ and $P_i$. Its hypotheses on $\mathcal{E}_1$, including the range
condition on the composite in the form \eqref{5.6:eq-composition-variant-a}, and the inner-link
range hypothesis on $\mathcal{E}_2$ are (b); those on $\mathcal{E}_2,\mathcal{E}_3,\mathcal{E}_4$
are \eqref{5.6:eq-link-hypotheses-a}, which hold by (a). The lemma therefore gives its conclusion
for the composite \eqref{5.6:eq-first-operation-1}.

The conclusions of Theorem~\ref{5.6:thm-link-holomorphy} are derived in its proof from the
conclusion of Lemma~\ref{5.6:lem-composition-variant} with $\lambda=3$ and the theorem's remaining
hypotheses. The same derivation, with (a) supplying those hypotheses and the range condition in (b)
supplying the membership of the composite in $\mathcal{A}_t$, applies with $\lambda=4$ and gives
the conclusions of Theorem~\ref{5.6:thm-link-holomorphy} on $P_1\times P_2\times P_3\times P_4$,
as asserted.
\end{proof}

\begin{remark}[What the link holomorphy theorem does not assert]\label{5.6:rem-link-scope}
Theorem~\ref{5.6:thm-link-holomorphy} is an implication, and it does not prove its inputs. It
does not prove the per-link analyticity readings \eqref{5.6:eq-link-readings-b}, one for each
link $i\in\{2,3,4\}$. It does not prove the membership condition
\eqref{5.6:eq-link-hypotheses-b}. It does not prove clause~(ii) of the analyticity condition
\eqref{5.6:eq-link-hypotheses-c}.

What it does establish is that nothing beyond these inputs is needed for holomorphy of the link
pieces in the decoupling parameters at the localisation site (L) of $\mathbf{R}_k$. This rests
on four points.
\begin{itemize}
\item[(a)] The nested composition of the link maps $\mathcal{E}_4$, $\mathcal{E}_3$,
$\mathcal{E}_2$ is the content of Lemma~\ref{5.6:lem-nested-composition} and of its variant,
Lemma~\ref{5.6:lem-composition-variant}. In the variant the range condition is imposed on the
composite, \eqref{5.6:eq-composition-variant-a}. Both lemmas are proved, with their hypotheses
displayed.
\item[(b)] Each link map $\mathcal{E}_i$ is required to be holomorphic jointly in its
decoupling variables $\sigma^{(i)}$ and in its argument. This requirement is stated explicitly
in Remark~\ref{5.6:rem-joint-holomorphy}.
\item[(c)] The nesting is $G^{\mathbb{C}}$-valued, and hence non-linear. It is handled by
taking the argument domains relatively open in the closed complex submanifolds
$\mathfrak{M}(Y)$ recalled in Definition~\ref{5.6:def-link-hypotheses}.
\item[(d)] The open polydisc on which holomorphy is obtained has radius exactly
$e^{\kappa_1}$ in every decoupling variable $s(\Delta)$, $\Delta\in\mathcal{V}(t)$. This is the
radius of the decoupling parameters in the construction of \cite{B-CMP116}.
\end{itemize}
\end{remark}

\begin{proof}
The readings \eqref{5.6:eq-link-readings-b}, the membership condition
\eqref{5.6:eq-link-hypotheses-b} and clause~(ii) of the analyticity condition
\eqref{5.6:eq-link-hypotheses-c} are hypotheses of Theorem~\ref{5.6:thm-link-holomorphy}. They
are therefore assumed there, not proved.

The conclusion of that theorem is then obtained from Lemma~\ref{5.6:lem-composition-variant},
applied to the link maps with $\rho=e^{\kappa_1}$, whose remaining hypotheses are items
(b)--(d) above, each fixed in Remark~\ref{5.6:rem-joint-holomorphy} and
Definition~\ref{5.6:def-link-hypotheses}. Hence, within the setting of
Definition~\ref{5.6:def-link-hypotheses}, with the joint holomorphy of
Remark~\ref{5.6:rem-joint-holomorphy} and under the conditions in which the readings are
stated, no unproved input other than the three displayed hypotheses enters.
\end{proof}

\begin{lemma}[Vanishing of all two-cutoff differences on one lattice. Stated; proof incomplete at the step marked $(\ast)$]
\label{5.7:lem-one-lattice-vanishing}
Let the two machines $\mathsf{A}$ and $\mathsf{B}$ of the two-cutoff comparison have the same lattice
and the same scaffold. They may differ in their stored terms, in the arguments of their couplings and
in their sources. Then every point $P=(P^{\mathsf{B}},P^{\mathsf{A}})$ of the closure
$\mathrm{cl}^{\mathrm{pr}}_J(\sigma)$ is diagonal, that is,
\begin{equation*}
P^{\mathsf{B}}=P^{\mathsf{A}}\qquad\text{at every slot.}
\end{equation*}
Consequently, for every $z$-free functional $\mathsf{z}$ of the construction, that is, every functional
that does not depend on the sources $z$,
\begin{equation*}
\delta\mathsf{z}:=\mathsf{z}^{\mathsf{B}}-\mathsf{z}^{\mathsf{A}}\equiv 0
\qquad\text{on }\mathrm{cl}^{\mathrm{pr}}_J .
\end{equation*}
In particular, let $\mathcal{H}_S$ be the sealed increment and $\mathsf{E}$ the energy-difference datum;
let $\Delta$, $\mathcal{K}^{\perp}$, $\hat{A}$ and $\mathsf{S}^{\perp}_v$ be the $z$-free functions of the
point of the closure and of the parameter $\pi$ from the common parameter domain of the machines that
enter the energy-difference datum and the layers it reads, the dilated Wilson family and the
multiscale frames; and let $\mathsf{a}$ be the size of the differences $\delta(\cdot)$ over the class.
Then
\begin{equation}
\label{5.7:eq-one-lattice-vanishing-1}
\begin{gathered}
\delta\mathcal{H}_S\equiv 0,\qquad \delta\mathsf{E}\equiv 0,\qquad \delta\Delta\equiv 0,\qquad
\delta\mathcal{K}^{\perp}\equiv 0,\qquad \delta\hat{A}\equiv 0,\\
\mathsf{a}=0,\qquad \mathsf{S}^{\perp}_v=0 .
\end{gathered}
\end{equation}
Hence the left-hand sides of the two-cutoff response statement at zeroth order, of the
energy-difference bound and of the low-order variation bound at the second cube size vanish. Each of
these bounds then states only that $0$ is at most its right-hand side, which is non-negative.
The hypothesis $M^{\partial}$ of the construction of the closure is a hypothesis on a single machine;
it is not affected by the number
of machines. The three bounds named above are inputs of clause (iv) of Theorem~0 alone; they are not
used for clauses (i), (iii-b) and (iii-c), in which the number $K$ of renormalisation steps is a
single number.
\end{lemma}

\begin{proof}
$(\ast)$ The argument is an induction over the construction of the closure
$\mathrm{cl}^{\mathrm{pr}}_J(\sigma)$.
It is the same induction by which the closure is shown to be built from seeds by uses. That induction
uses the equality of the stored terms of the two machines only for its last clause, $\delta F\equiv 0$,
the vanishing of the difference of the stored terms $F$.
The present lemma does not assume that the stored terms are equal and does not assert that clause.
The description of the seeds and of the uses of the closure invoked below is taken by statement from
the construction of the closure; this is the step marked $(\ast)$.

\emph{Seeds.} The seeds of the closure are the trivial configuration $\mathbb{1}$ and the pairs of
backgrounds $(U_k^{\mathsf{B}},U_k^{\mathsf{A}})$ taken at equal values of the common variables. The
point $\mathbb{1}$ is diagonal. For a pair of backgrounds, the two machines have one lattice, so
$U_k^{\mathsf{B}}$ and $U_k^{\mathsf{A}}$ are one function on one lattice evaluated at one argument.
Hence the pair is diagonal.

\emph{Step.} A use $u$ is a pair of $z$-free maps $\mu_u^{\mathsf{m}}$, $\mathsf{m}\in\{\mathsf{A},
\mathsf{B}\}$. Each is the same construction performed in the respective machine, and the two share one
parameter domain. On one lattice with one scaffold the two components $\mu_u^{\mathsf{A}}$ and
$\mu_u^{\mathsf{B}}$ are one map. This map is fed one parameter $\pi$ from the common domain, and, by the
induction hypothesis, one diagonal point. Its output is therefore diagonal. Thus every point of
$\mathrm{cl}^{\mathrm{pr}}_J(\sigma)$ is diagonal at every slot.

\emph{The listed objects.} By their definitions, each of the five differences in the first line of
\eqref{5.7:eq-one-lattice-vanishing-1} is the difference of the two machines' values of one $z$-free
function of the point of the closure and of the parameter $\pi$. These definitions are
those of the energy-difference datum and the layers it reads, of the dilated Wilson family, and of the
multiscale frames. Evaluated at a diagonal point and at one $\pi$, each such function takes the same
value in both machines, so each of these five differences vanishes identically. The two remaining
statements of \eqref{5.7:eq-one-lattice-vanishing-1}, $\mathsf{a}=0$ and $\mathsf{S}^{\perp}_v=0$, are
consequences of $\delta\mathsf{z}\equiv 0$ on the closure, as listed in the lemma. The assertion
about the three bounds follows at once.
\end{proof}

\begin{lemma}[Interpolation between two analytic maps]\label{5.7:lem-two-map-interpolation}
Let $\mathfrak{c}$ be a finite set and let $a>0$. A field on $\mathfrak{c}$ is a family
$B=(B(c))_{c\in\mathfrak{c}}$
with $B(c)\in\mathfrak{g}^{c}$, where the complexified Lie algebra $\mathfrak{g}^{c}$ carries the
operator norm $|\cdot|$. Write
\begin{equation*}
\sup|B|:=\sup_{c\in\mathfrak{c}}|B(c)|,\qquad
\mathfrak{D}(a):=\bigl\{B:\ \sup_{c}|B(c)|<a\bigr\}.
\end{equation*}
A field $B$ is \emph{real} if it is $\mathfrak{g}$-valued. Every field is $B=B'+iB''$ with $B'$, $B''$
real and $|B'(c)|\le|B(c)|$, $|B''(c)|\le|B(c)|$ for every $c$ (the decomposition into Hermitian and
anti-Hermitian parts). Let $E$ be a finite-dimensional complex normed space, let
$k\colon\mathfrak{c}\to[0,\infty[$, and let $\mathsf{K}\ge\sum_{c}k(c)$. For a field $b$ on
$\mathfrak{c}$ and $0\le\theta<1$ put
\begin{equation*}
T^{(1-\theta)}[b]:=\sum_{c\in\mathfrak{c}}k(c)^{1-\theta}\mathsf{K}^{\theta}|b(c)|,
\end{equation*}
so that $T^{(1)}[b]=\sum_{c}k(c)|b(c)|\le T^{(1-\theta)}[b]$. Assume, for
$\mathsf{m}\in\{\mathsf{A},\mathsf{B}\}$,
\begin{equation}\label{5.7:eq-two-map-interpolation-a}
\begin{aligned}
&(\Lambda 0)\quad \Phi^{\mathsf{B}},\Phi^{\mathsf{A}}\colon\mathfrak{D}(a)\to E
\ \text{are holomorphic, and}\ \delta\Phi:=\Phi^{\mathsf{B}}-\Phi^{\mathsf{A}};\\
&(\Lambda 1)\quad |\Phi^{\mathsf{m}}(\tilde B)-\Phi^{\mathsf{m}}(B)|
\le\sum_{c}k(c)\,|\tilde B(c)-B(c)|
\quad\text{for } B,\tilde B\in\mathfrak{D}(a);\\
&(\Lambda 2)\quad |\delta\Phi(B)|\le\varepsilon_{\star}a\mathsf{K}
\quad\text{for real } B\in\mathfrak{D}(a),\ \text{where } 0\le\varepsilon_{\star}\le1.
\end{aligned}
\end{equation}
Then the following hold.
\begin{enumerate}
\item[(a)] $|\delta\Phi|\le 3a\mathsf{K}$ on $\mathfrak{D}(a)$.
\item[(b)] $|\delta\Phi|\le(3\varepsilon_{\star})^{1/2}a\mathsf{K}$ on $\mathfrak{D}(a/4)$.
\item[(c)] If $\sup|B|+\sup|b|\le a/8$ and $0\le\theta\le1$, then
\begin{equation*}
|\delta\Phi(B+b)-\delta\Phi(B)|\le 28\,\varepsilon_{\star}^{\theta/2}\,T^{(1-\theta)}[b].
\end{equation*}
\item[(d)] (One map.) If $\sup|B|+\sup|b|\le a/2$ and $\sup|e|\le a/4$, then for
$\mathsf{m}\in\{\mathsf{A},\mathsf{B}\}$
\begin{equation*}
\bigl|\Phi^{\mathsf{m}}(B+e+b)-\Phi^{\mathsf{m}}(B+e)-\Phi^{\mathsf{m}}(B+b)+\Phi^{\mathsf{m}}(B)\bigr|
\le\frac{4\sup|e|}{a}\,T^{(1)}[b].
\end{equation*}
\item[(e)] (Assembly.) If $\sup|B_0|+\sup|b^{\mathsf{A}}|\le a/8$, $\sup|e_0|\le a/4$ and
$\sup|B_0+e_0|+\sup|b^{\mathsf{B}}|<a$, then for $0\le\theta\le1$ the mixed difference
$\Delta^{\times}$ below obeys
\begin{equation}\label{5.7:eq-two-map-interpolation-b}
\begin{aligned}
\Delta^{\times}&:=\bigl[\Phi^{\mathsf{B}}(B_0+e_0+b^{\mathsf{B}})-\Phi^{\mathsf{B}}(B_0+e_0)\bigr]
-\bigl[\Phi^{\mathsf{A}}(B_0+b^{\mathsf{A}})-\Phi^{\mathsf{A}}(B_0)\bigr],\\
|\Delta^{\times}|&\le 28\,\varepsilon_{\star}^{\theta/2}\,T^{(1-\theta)}[b^{\mathsf{A}}]
+\frac{4\sup|e_0|}{a}\,T^{(1)}[b^{\mathsf{A}}]+T^{(1)}[b^{\mathsf{B}}-b^{\mathsf{A}}].
\end{aligned}
\end{equation}
\item[(f)] (Absolute form.) If $\sup|B_0|<a/4$ and $e_0$ is a field with
$B_0+e_0\in\mathfrak{D}(a)$, so that $\Phi^{\mathsf{B}}(B_0+e_0)$ is defined, then
\begin{equation*}
|\Phi^{\mathsf{B}}(B_0+e_0)-\Phi^{\mathsf{A}}(B_0)|\le(3\varepsilon_{\star})^{1/2}a\mathsf{K}
+T^{(1)}[e_0].
\end{equation*}
\end{enumerate}
\end{lemma}

\begin{proof}
Two elementary facts from one complex variable are used. The first is the two-constants lemma on one
complex line, in the following form: if $\varphi$ is holomorphic on the disc $\{|z|<3\}$,
$|\varphi|\le\bar\mu$ there and $|\varphi|\le\mu\le\bar\mu$ on the real segment $\left]-3,3\right[$,
then
$|\varphi(i)|\le(\mu\bar\mu)^{1/2}$. The second is Schwarz's lemma: if $h$ is holomorphic on
$\{|t|<r\}$, $h(0)=0$ and $|h|\le C$ there, then $|h(t)|\le C|t|/r$ for $|t|<r$.

\emph{(a).} For $B\in\mathfrak{D}(a)$ we have the identity
\begin{equation*}
\delta\Phi(B)=\delta\Phi(0)+\bigl[\Phi^{\mathsf{B}}(B)-\Phi^{\mathsf{B}}(0)\bigr]
-\bigl[\Phi^{\mathsf{A}}(B)-\Phi^{\mathsf{A}}(0)\bigr].
\end{equation*}
The zero field is real and lies in $\mathfrak{D}(a)$, so $|\delta\Phi(0)|\le\varepsilon_{\star}a\mathsf{K}$
by $(\Lambda 2)$. By $(\Lambda 1)$, applied in each of the two maps to $\tilde B=B$ and the zero field,
\begin{equation*}
|\Phi^{\mathsf{m}}(B)-\Phi^{\mathsf{m}}(0)|\le\sum_{c}k(c)|B(c)|\le a\sum_{c}k(c)\le a\mathsf{K}.
\end{equation*}
Hence $|\delta\Phi(B)|\le\varepsilon_{\star}a\mathsf{K}+2a\mathsf{K}\le3a\mathsf{K}$, since
$\varepsilon_{\star}\le1$.

\emph{(b).} Let $B=B'+iB''\in\mathfrak{D}(a/4)$, let $\mathsf{l}$ be a linear functional on $E$ of norm
$1$, and put $\varphi(z):=\mathsf{l}\bigl(\delta\Phi(B'+zB'')\bigr)$. For $|z|<3$ and every $c$,
\begin{equation*}
|B'(c)+zB''(c)|\le|B'(c)|+|z|\,|B''(c)|<\tfrac{a}{4}+\tfrac{3a}{4}=a,
\end{equation*}
so $B'+zB''\in\mathfrak{D}(a)$. Thus $\varphi$ is holomorphic on $\{|z|<3\}$, being the composition of
the holomorphic map $\delta\Phi$ with an affine map and a linear functional, and $|\varphi|\le3a\mathsf{K}$
there by (a). For real $z\in\left]-3,3\right[$ the field $B'+zB''$ is real and lies in
$\mathfrak{D}(a)$, so
$|\varphi(z)|\le\varepsilon_{\star}a\mathsf{K}$ by $(\Lambda 2)$; and
$\varepsilon_{\star}a\mathsf{K}\le3a\mathsf{K}$ because $\varepsilon_{\star}\le1$. The two-constants
lemma with $\mu=\varepsilon_{\star}a\mathsf{K}$ and $\bar\mu=3a\mathsf{K}$ gives
\begin{equation*}
|\mathsf{l}(\delta\Phi(B))|=|\varphi(i)|\le(3\varepsilon_{\star})^{1/2}a\mathsf{K}.
\end{equation*}
Since the norm of a vector of $E$ is the supremum of $|\mathsf{l}(\cdot)|$ over the functionals
$\mathsf{l}$ of norm $1$, (b) follows.

\emph{(c).} Enumerate $\mathfrak{c}=\{c_1,\dots,c_{|\mathfrak{c}|}\}$. For $c\in\mathfrak{c}$ and
$v\in\mathfrak{g}^{c}$ write $v[c]$ for the field equal to $v$ at $c$ and to $0$ elsewhere,
and put $B^{(j)}:=B+\sum_{i\le j}b(c_i)[c_i]$ for $0\le j\le|\mathfrak{c}|$, so that
$B^{(0)}=B$ and $B^{(|\mathfrak{c}|)}=B+b$. At each $c$ the value of $B^{(j)}$ is $B(c)$ or
$B(c)+b(c)$, so $\sup|B^{(j)}|\le\sup|B|+\sup|b|\le a/8$. Telescoping,
\begin{equation*}
\delta\Phi(B+b)-\delta\Phi(B)=\sum_{j=1}^{|\mathfrak{c}|}
\bigl[\delta\Phi(B^{(j)})-\delta\Phi(B^{(j-1)})\bigr].
\end{equation*}
The $j$-th term vanishes if $b(c_j)=0$. If $b(c_j)\neq0$, put
\begin{equation*}
f(t):=\delta\Phi\bigl(B^{(j-1)}+t\,b(c_j)[c_j]\bigr),\qquad
|t|<r_j:=\frac{a}{8|b(c_j)|},
\end{equation*}
and note $r_j\ge1$ because $|b(c_j)|\le\sup|b|\le a/8$. The argument of $f$ coincides with
$B^{(j-1)}$ off $c_j$, where its norm is at most $a/8$, and at $c_j$ it equals
$B(c_j)+t\,b(c_j)$, of norm at most $|B(c_j)|+|t|\,|b(c_j)|<a/8+a/8$; so it stays in
$\mathfrak{D}(a/4)$. Hence $f$ is holomorphic on $\{|t|<r_j\}$ and
$|f|\le(3\varepsilon_{\star})^{1/2}a\mathsf{K}$ there by (b), so
$|f-f(0)|\le2(3\varepsilon_{\star})^{1/2}a\mathsf{K}$. Schwarz's lemma for $f-f(0)$ gives
\begin{equation*}
|f(1)-f(0)|\le\frac{2(3\varepsilon_{\star})^{1/2}a\mathsf{K}}{r_j}
=16(3\varepsilon_{\star})^{1/2}\mathsf{K}\,|b(c_j)|=:X_j.
\end{equation*}
On the other hand $f(1)-f(0)$ is the difference of the increments of $\Phi^{\mathsf{B}}$ and of
$\Phi^{\mathsf{A}}$ between two fields in $\mathfrak{D}(a/4)$ that differ only at $c_j$, by $b(c_j)$;
applying $(\Lambda 1)$ in both maps,
\begin{equation*}
|f(1)-f(0)|\le2k(c_j)|b(c_j)|=:Y_j.
\end{equation*}
For non-negative numbers $\min\{X_j,Y_j\}\le X_j^{\theta}Y_j^{1-\theta}$,
and
\begin{equation*}
X_j^{\theta}Y_j^{1-\theta}
=(16\cdot3^{1/2})^{\theta}\,2^{1-\theta}\,\varepsilon_{\star}^{\theta/2}\,
\mathsf{K}^{\theta}k(c_j)^{1-\theta}|b(c_j)|.
\end{equation*}
The factor $(16\cdot3^{1/2})^{\theta}2^{1-\theta}$ is log-linear in $\theta$, hence at most
$\max\{2,16\cdot3^{1/2}\}\le28$. Since $f(0)=\delta\Phi(B^{(j-1)})$ and $f(1)=\delta\Phi(B^{(j)})$,
we obtain, for every $j$ (trivially when $b(c_j)=0$),
\begin{equation*}
|\delta\Phi(B^{(j)})-\delta\Phi(B^{(j-1)})|
\le28\,\varepsilon_{\star}^{\theta/2}\mathsf{K}^{\theta}k(c_j)^{1-\theta}|b(c_j)|.
\end{equation*}
The right side is $28\,\varepsilon_{\star}^{\theta/2}$ times the $c_j$-summand of $T^{(1-\theta)}[b]$,
so summing over $j$ gives (c).

\emph{(d).} If $e=0$ the left side vanishes. Otherwise put
$g(s):=\Phi^{\mathsf{m}}(B+se+b)-\Phi^{\mathsf{m}}(B+se)$ for $|s|<\rho:=a/(2\sup|e|)$; then $\rho\ge2$
because $\sup|e|\le a/4$. For $|s|<\rho$ and every $c$,
\begin{equation*}
|B(c)+se(c)+b(c)|\le\sup|B|+\sup|b|+|s|\sup|e|<\tfrac{a}{2}+\tfrac{a}{2}=a,
\end{equation*}
and likewise $|B(c)+se(c)|<a$; so both arguments lie in $\mathfrak{D}(a)$ and $g$ is holomorphic on
$\{|s|<\rho\}$. The two arguments differ by $b$, so $|g|\le\sum_{c}k(c)|b(c)|=T^{(1)}[b]$ by
$(\Lambda 1)$, and $|g-g(0)|\le2T^{(1)}[b]$. Schwarz's lemma for $g-g(0)$, at $s=1<\rho$, gives
\begin{equation*}
|g(1)-g(0)|\le\frac{2T^{(1)}[b]}{\rho}=\frac{4\sup|e|}{a}\,T^{(1)}[b],
\end{equation*}
and $g(1)-g(0)$ is the left side of (d).

\emph{(e).} Adding and subtracting $\Phi^{\mathsf{B}}(B_0+b^{\mathsf{A}})$, $\Phi^{\mathsf{B}}(B_0)$ and
$\Phi^{\mathsf{B}}(B_0+e_0+b^{\mathsf{A}})$ gives
\begin{align*}
\Delta^{\times}
&=\bigl[\delta\Phi(B_0+b^{\mathsf{A}})-\delta\Phi(B_0)\bigr]\\
&\quad+\bigl[\Phi^{\mathsf{B}}(B_0+e_0+b^{\mathsf{A}})-\Phi^{\mathsf{B}}(B_0+e_0)
-\Phi^{\mathsf{B}}(B_0+b^{\mathsf{A}})+\Phi^{\mathsf{B}}(B_0)\bigr]\\
&\quad+\bigl[\Phi^{\mathsf{B}}(B_0+e_0+b^{\mathsf{B}})-\Phi^{\mathsf{B}}(B_0+e_0+b^{\mathsf{A}})\bigr].
\end{align*}
The first bracket is bounded by (c) with $B=B_0$, $b=b^{\mathsf{A}}$, since
$\sup|B_0|+\sup|b^{\mathsf{A}}|\le a/8$. The second is bounded by (d) in the map $\Phi^{\mathsf{B}}$
with $B=B_0$, $e=e_0$, $b=b^{\mathsf{A}}$, since $\sup|B_0|+\sup|b^{\mathsf{A}}|\le a/8\le a/2$ and
$\sup|e_0|\le a/4$. For the third, $B_0+e_0+b^{\mathsf{B}}\in\mathfrak{D}(a)$ because
$\sup|B_0+e_0|+\sup|b^{\mathsf{B}}|<a$, and $B_0+e_0+b^{\mathsf{A}}\in\mathfrak{D}(a)$ because its
norm at each $c$ is at most $a/8+a/4$; the two fields differ by $b^{\mathsf{B}}-b^{\mathsf{A}}$, so
$(\Lambda 1)$ in $\Phi^{\mathsf{B}}$ bounds the third bracket by $T^{(1)}[b^{\mathsf{B}}-b^{\mathsf{A}}]$.
Adding the three bounds gives \eqref{5.7:eq-two-map-interpolation-b}.

\emph{(f).} Write
\begin{equation*}
\Phi^{\mathsf{B}}(B_0+e_0)-\Phi^{\mathsf{A}}(B_0)
=\delta\Phi(B_0)+\bigl[\Phi^{\mathsf{B}}(B_0+e_0)-\Phi^{\mathsf{B}}(B_0)\bigr].
\end{equation*}
Since
$B_0\in\mathfrak{D}(a/4)$, (b) gives $|\delta\Phi(B_0)|\le(3\varepsilon_{\star})^{1/2}a\mathsf{K}$. Since
$B_0$ and $B_0+e_0$ lie in $\mathfrak{D}(a)$ and differ by $e_0$, $(\Lambda 1)$ in $\Phi^{\mathsf{B}}$
bounds the bracket by $T^{(1)}[e_0]$.
\end{proof}

\begin{remark}[Where the fractional power falls]\label{5.7:rem-fractional-power}
The following observations concern Lemma~\ref{5.7:lem-two-map-interpolation} and its hypotheses
\eqref{5.7:eq-two-map-interpolation-a}.

In part~(c) of
Lemma~\ref{5.7:lem-two-map-interpolation}, which asserts that, whenever
$\sup|B|+\sup|b|\le a/8$ and $0\le\theta\le 1$,
\begin{equation*}
  \bigl|\delta\Phi(B+b)-\delta\Phi(B)\bigr|
  \le 28\,\varepsilon_{\star}^{\theta/2}\,T^{(1-\theta)}[b],
\end{equation*}
the fractional power falls on $\varepsilon_{\star}$ alone. This is the index factor of the
real-slice closeness bound among the hypotheses \eqref{5.7:eq-two-map-interpolation-a}. It is a
classical number, introduced afresh, and no recursion feeds it back. The inherited mismatch
$e_0$ and the data mismatch $b^{\mathsf{B}}-b^{\mathsf{A}}$ enter linearly, at the full kernel
$k$, and only through bounds for one of the two maps at a time, as in parts~(d) and~(e).

The device is the elementary inequality
\begin{equation}\label{5.7:eq-fractional-power-1}
  \min\{u,v\}\le u^{\theta}v^{1-\theta},\qquad u,v\ge 0,\ 0\le\theta\le 1,
\end{equation}
a minimum bounded by a geometric mean with unequal weights. The exponent is arranged to fall on a
fresh term only. It is combined with the splitting of the mixed difference $\Delta^{\times}$ in
second-difference form, as in part~(e). No root is taken of a size that the recursion
feeds back; on a classical number such as $\varepsilon_{\star}$ a square root costs nothing.

\emph{No volume.} No constant of Lemma~\ref{5.7:lem-two-map-interpolation} depends on the number
of data bonds, that is, on the volume. The constant $\mathsf{K}$ is an upper
bound for the kernel sum $\sum_{c}k(c)$, not for the number of its terms.

\emph{One complex line.} In the proof of Lemma~\ref{5.7:lem-two-map-interpolation}, the estimate
invoked there for analytic functions is applied on a single complex line only, the line
$t\mapsto B+tb$.

\emph{Variables of one map only.} Variables on which only $\Phi^{\mathsf{B}}$ depends are
admitted. Let $\mathfrak{c}_{\mathsf{B}}\subset\mathfrak{c}$ be the set of indices $c$ on whose
variables $B(c)$ only $\Phi^{\mathsf{B}}$ depends, and
let $\Phi^{\mathsf{A}}$ not depend on them:
\begin{equation}\label{5.7:eq-fractional-power-a}
  \Phi^{\mathsf{A}}(B+b)=\Phi^{\mathsf{A}}(B)
  \quad\text{whenever } b(c)=0 \text{ for all } c\notin\mathfrak{c}_{\mathsf{B}} .
\end{equation}
Then the first-difference bound among \eqref{5.7:eq-two-map-interpolation-a} holds for
$\Phi^{\mathsf{A}}$ with any non-negative values $k(c)$, $c\in\mathfrak{c}_{\mathsf{B}}$, for
which $\sum_{c}k(c)\le\mathsf{K}$ still holds.

\emph{Weighted-sup variant.} Let $p\colon\mathfrak{c}\to(0,\infty)$ be a weight, let
$x$ be a fixed point, and let $p_x>0$ be the value of the weight at $x$.
Suppose that, instead of the first-difference bound among
\eqref{5.7:eq-two-map-interpolation-a}, one has
\begin{equation}\label{5.7:eq-fractional-power-2}
  \bigl|\Phi^{\mathsf{m}}(B+b)-\Phi^{\mathsf{m}}(B)\bigr|
  \le \mathsf{K}\,p_x\,\|b\|_p,
  \qquad \|b\|_p:=\sup_{c}\,|b(c)|/p(c) ,
\end{equation}
for $\mathsf{m}\in\{\mathsf{A},\mathsf{B}\}$. Then, under the hypotheses of part~(c),
\begin{equation}\label{5.7:eq-fractional-power-3}
  \bigl|\delta\Phi(B+b)-\delta\Phi(B)\bigr|
  \le 28\,\varepsilon_{\star}^{\theta/2}\,\mathsf{K}\,
  \bigl(\sup|b|\bigr)^{\theta}\bigl(p_x\|b\|_p\bigr)^{1-\theta}.
\end{equation}
Under the hypotheses of part~(d), that part holds with $T^{(1)}[b]$ replaced by
$\mathsf{K}\,p_x\|b\|_p$.
\end{remark}

\begin{proof}
Only the assertion on variables of one map and the weighted-sup variant require proof.

For the assertion on variables of one map, let $b$ be given and let $b'$ be equal to $b$
on $\mathfrak{c}\setminus\mathfrak{c}_{\mathsf{B}}$ and to $0$ on $\mathfrak{c}_{\mathsf{B}}$;
then $\sup|b'|\le\sup|b|$. Since $\Phi^{\mathsf{A}}$ does not depend on the variables indexed by
$\mathfrak{c}_{\mathsf{B}}$, \eqref{5.7:eq-fractional-power-a} gives
$\Phi^{\mathsf{A}}(B+b)-\Phi^{\mathsf{A}}(B)=\Phi^{\mathsf{A}}(B+b')-\Phi^{\mathsf{A}}(B)$. The
first-difference bound applied to the pair $(B,b')$ has right-hand side
$T^{(1)}[b']=\sum_{c}k(c)|b'(c)|$, in which the values $k(c)$, $c\in\mathfrak{c}_{\mathsf{B}}$,
are multiplied by $|b'(c)|=0$. Hence it does not change when these values are replaced by any
non-negative ones, and it is at most $\sum_{c}k(c)|b(c)|=T^{(1)}[b]$ for every such choice.

For \eqref{5.7:eq-fractional-power-3}, repeat the proof of part~(c) of
Lemma~\ref{5.7:lem-two-map-interpolation} along the single line $t\mapsto B+tb$. At each place
where that proof applies the first-difference bound among
\eqref{5.7:eq-two-map-interpolation-a}, apply \eqref{5.7:eq-fractional-power-2} instead. This
yields \eqref{5.7:eq-fractional-power-3} in place of the bound of part~(c).

For part~(d), repeat its proof in the same way, with \eqref{5.7:eq-fractional-power-2} in place
of the first-difference bound. This yields part~(d) with $T^{(1)}[b]$ replaced by
$\mathsf{K}\,p_x\|b\|_p$.
\end{proof}

Clause~(c) of the interpolation between two analytic maps (Lemma~\ref{5.7:lem-two-map-interpolation})
returns the new term of the two-map difference at the rate $(1-\theta)\delta$. For $\theta=\tfrac14$
this rate is $\tfrac34\delta$.

The full-rate majorant of the profile sum $\mathsf{T}(x)$ of the energy-difference construction spends the whole of $\delta$:
\begin{itemize}
\item $\tfrac12\delta$ for the class move;
\item $\tfrac14\delta$ for the summation bound over all levels of the frame;
\item $\tfrac14\delta$ kept in the final exponent.
\end{itemize}

The two-cutoff difference bound is stated for a fixed $\delta'\in(0,\delta]$, and in its last
estimate it bounds the profile sum at the rate $\delta'$ by the majorant of $\mathsf{T}$. For
$\delta'<\delta$ the kept rate would then fall below $\tfrac14\delta$, while the dilation radii $R_u$,
$R^{\perp}_v$ and the factor $e^{\delta D_v/8}$ stay as they are. The product of a radius with the norm
of the layer difference would then grow like $e^{(\delta-\delta')D_v}$, with $D_v$ the distance of the
unit region. The following lemma shows that at $\delta'=\tfrac34\delta$ no such loss occurs.

\begin{lemma}[Majorants at three quarters of the kernel rate]
\label{5.7:lem-three-quarter-rate}
Work in the setting of the energy-difference construction at level $j$, with block $\square$, the set
$F_{\square}$ attached to $\square$ in that setting and the distance $D_v$ of the unit region
$\mathfrak{r}_v$. Alternatively, work in the setting
of the dilated Wilson family at level $i$ with distance $D_v$.
Let $\vartheta_{\mathrm{p}}$ be the small index base of the pointwise weighted profile bound, and let
$\vartheta_{\mathrm{p}}(\square)$ be the index that this bound attaches to the block $\square$.
On the frame $\mathfrak{B}_j$ put $f:=\hat p+|\mathsf{t}|\,\hat p_{\square}$. For $c\in F_{\square}$ put
\begin{equation*}
\mathfrak{h}_{\ast}(c):=15B_{\sharp}^{2}f(\bar c),\qquad
T_{\ast}(x):=\sum_{c\in F_{\square}}e^{-\frac34\delta d_{\star}(x,c)}\,\mathfrak{h}_{\ast}(c).
\end{equation*}
Then, with $r=1.1$ the constant of the energy-difference construction, for every element $x$ of the
region $\Omega_j$ at level $j$ and every cell $z$ of $\mathfrak{B}_j$, where $\bar c$ and $\bar x$
denote the centres of $c$ and of $x$ and $m(\bar x)$ is the index that the pointwise weighted
profile bound attaches to $\bar x$,
\begin{equation}\label{5.7:eq-three-quarter-rate-a}
\begin{aligned}
f(z)&\le 4r\,[\vartheta_{\mathrm{p}}+|\mathsf{t}|\vartheta_{\mathrm{p}}(\square)]\,
\rho_{m(\bar x)}\,\alpha_{\cdot,j}\,
e^{\frac14\delta d_j(x,z)},\\
T_{\ast}(x)&\le 79B_{\sharp}^{2}c_1^{b}\beta\,
[\vartheta_{\mathrm{p}}+|\mathsf{t}|\vartheta_{\mathrm{p}}(\square)]\,\alpha_{\cdot,j}\,
e^{-\frac14\delta d_{\star}(x,F_{\square})},\\
8L^{3}B_H\,T_{\ast}(x)&\le K_1\beta e^{-\frac14\delta D_v}\,
[5.5\vartheta_{\mathrm{p}}+4.4|\mathsf{t}|\vartheta_{\mathrm{p}}(\square)]\,\alpha_{\cdot,j}
\qquad (x\in\mathfrak{r}_v).
\end{aligned}
\end{equation}
The first two lines have the following content.
\begin{enumerate}
\item[(i)] The first line states that the profile $f$ grows away from $x$ at the rate $\tfrac14\delta$, not $\tfrac12\delta$.
\item[(ii)] The second line is below the majorant
\begin{equation}\label{5.7:eq-three-quarter-rate-1}
30B_{\sharp}^{2}c_1^{b}\beta\,
[5.5\vartheta_{\mathrm{p}}+4.4|\mathsf{t}|\vartheta_{\mathrm{p}}(\square)]\,\alpha_{\cdot,j}\,
e^{-\frac14\delta D_v}
\end{equation}
of $\mathsf{T}(x)$ in the energy-difference construction.
\end{enumerate}
Consequently the third line holds on $\mathfrak{r}_v$. In the setting of the dilated Wilson family, $8L^3B_HT_{\ast}\le\hat h_v$ holds on the
region $\Delta$ of that construction. Every majorant and radius derived from the corresponding bound for
$\mathsf{T}$ holds for $T_{\ast}$ as it does for $\mathsf{T}$. These are the majorants and radii of clause (iii) of the radius lemma of the
energy-difference construction, the constants of the dilated Wilson family, and the radii of the two-cutoff difference bound.
\end{lemma}

\begin{proof}
\emph{Proof of (i).} We use the pointwise weighted profile bound for the profile $p$ and its
smoothing weight $\varrho$. For every level $\ell$, every element $y$ of the region $\Omega_\ell$ at
level $\ell$ and every $y'$, that bound, multiplied by $e^{\frac12\delta d_\ell(y,y')}$, reads
\begin{equation*}
\varrho(y')p(y')\le 4r\vartheta_{\mathrm{p}}\rho_m\alpha_{\cdot,\ell}\,
e^{-\frac12\delta d^{\wedge}(\bar y',\mathcal{C}_j)}\,e^{+\frac14\delta d_\ell(y,y')},
\end{equation*}
where $\mathcal{C}_j$ is the reference set of that bound and $m=m(\bar y)$ is the index it attaches
to the centre $\bar y$ of $y$. Take $(\ell,y):=(j,x)$, so that $m=m(\bar x)$.

By definition,
\begin{equation*}
\hat p(z)=\sup_{y'}\varrho(y')p(y')e^{-\frac12\delta d_j(z,y')}.
\end{equation*}
The triangle inequality for the frame distances gives $d_j(x,y')\le d_j(x,z)+d_j(z,y')$. Hence
\begin{equation*}
e^{\frac14\delta d_j(x,y')}\,e^{-\frac12\delta d_j(z,y')}\le e^{\frac14\delta d_j(x,z)}\,e^{-\frac14\delta d_j(z,y')}.
\end{equation*}
Inserting this into the displayed bound, and using $e^{-\frac12\delta d^{\wedge}(\bar y',\mathcal{C}_j)}\le 1$, we obtain
\begin{align*}
\hat p(z)&\le 4r\vartheta_{\mathrm{p}}\rho_{m(\bar x)}\alpha_{\cdot,j}\,e^{\frac14\delta d_j(x,z)}
\sup_{y'}e^{-\frac14\delta d_j(z,y')}\\
&\le 4r\vartheta_{\mathrm{p}}\rho_{m(\bar x)}\alpha_{\cdot,j}\,e^{\frac14\delta d_j(x,z)}.
\end{align*}

For $\hat p_{\square}$ the supremum runs over $y'\in\square^{\zeta}$. The proof of the pointwise
weighted profile bound shows $d^{\wedge}(\square^{\zeta},\mathcal{C}_j)\ge d^{\wedge}(\square,R_n)$,
where $R_n$ is the set attached to $\square$ there. It follows that
$\vartheta_{\mathrm{p}}e^{-\frac12\delta d^{\wedge}(\bar y',\mathcal{C}_j)}\le\vartheta_{\mathrm{p}}(\square)$
for every $y'\in\square^{\zeta}$. The same computation therefore gives
\begin{equation*}
\hat p_{\square}(z)\le 4r\vartheta_{\mathrm{p}}(\square)\rho_{m(\bar x)}\alpha_{\cdot,j}\,
e^{\frac14\delta d_j(x,z)}.
\end{equation*}
Adding this bound, multiplied by $|\mathsf{t}|$, to the bound for $\hat p(z)$ gives the first line of
\eqref{5.7:eq-three-quarter-rate-a}.

\emph{Proof of (ii).} We use three facts:
\begin{itemize}
\item $r=1.1$ and $\rho_m\le\beta$, both from the energy-difference construction, so that $4r\rho_{m(\bar x)}\le 4.4\beta$;
\item by the comparison of frame distances with cell centres, $d_j(x,\bar c)$ exceeds
$d_{\star}(x,c)$ by at most an additive constant whose exponential at the rate $\delta$ is at
most $2$, so that
\begin{equation*}
e^{\frac14\delta d_j(x,\bar c)}\le 2^{1/4}e^{\frac14\delta d_{\star}(x,c)};
\end{equation*}
\item the first line of \eqref{5.7:eq-three-quarter-rate-a}, applied at $z=\bar c$.
\end{itemize}
Together they give, for $c\in F_{\square}$,
\begin{equation*}
\mathfrak{h}_{\ast}(c)\le 15\cdot 4.4\cdot 2^{1/4}\,B_{\sharp}^{2}\beta\,
[\vartheta_{\mathrm{p}}+|\mathsf{t}|\vartheta_{\mathrm{p}}(\square)]\,
\alpha_{\cdot,j}\,e^{\frac14\delta d_{\star}(x,c)}.
\end{equation*}
Since $2^{1/4}<1.19$, we have $15\cdot 4.4\cdot 2^{1/4}<66\cdot1.19<79$. Combining $e^{-\frac34\delta d_{\star}(x,c)}$ with $e^{\frac14\delta d_{\star}(x,c)}$ therefore gives
\begin{equation*}
T_{\ast}(x)\le 79B_{\sharp}^{2}\beta\,
[\vartheta_{\mathrm{p}}+|\mathsf{t}|\vartheta_{\mathrm{p}}(\square)]\,\alpha_{\cdot,j}
\sum_{c\in F_{\square}}e^{-\frac12\delta d_{\star}(x,c)}.
\end{equation*}
For $c\in F_{\square}$ we have $d_{\star}(x,c)\ge d_{\star}(x,F_{\square})$. Hence
\begin{equation*}
e^{-\frac12\delta d_{\star}(x,c)}\le e^{-\frac14\delta d_{\star}(x,F_{\square})}\,e^{-\frac14\delta d_{\star}(x,c)}.
\end{equation*}
The summation bound over all levels of the frame bounds $\sum_{c}e^{-\frac14\delta d_{\star}(x,c)}$ by $c_1^{b}$; every level of the frame enters that sum. This proves the second line of \eqref{5.7:eq-three-quarter-rate-a}.

Finally, $79\le 30\cdot 4.4<30\cdot 5.5$. So the coefficients of $\vartheta_{\mathrm{p}}$ and of
$|\mathsf{t}|\vartheta_{\mathrm{p}}(\square)$ in the second line are at most those in
\eqref{5.7:eq-three-quarter-rate-1}, the majorant of $\mathsf{T}(x)$. That majorant yields the third
line of \eqref{5.7:eq-three-quarter-rate-a} for $\mathsf{T}$ on $\mathfrak{r}_v$, and the bound by
$\hat h_v$ on $\Delta$; hence both hold for $T_{\ast}$.
\end{proof}

\begin{remark}[Distribution of the rate]
At the rate $\tfrac34\delta$ the exponent is distributed as follows:
\begin{itemize}
\item $\tfrac14\delta$ is used in (i);
\item $\tfrac14\delta$ is used in the summation bound over all levels of the frame;
\item $\tfrac14\delta$ is kept in $e^{-\frac14\delta d_{\star}(x,F_{\square})}$.
\end{itemize}
The quarter made available is the one that the pointwise weighted profile bound itself keeps. The class move of the full-rate majorant, $[\,\cdot\,](\bar c)\le 2e^{\frac12\delta d_{\star}(x,c)}[\,\cdot\,](x)$, does not use it.
\end{remark}

\begin{definition}[The paired frame and the data mismatch]\label{5.7:def-data-mismatch}
The following objects are fixed in the proof of the two-cutoff response statement at zeroth
order. The two cutoffs are the lattice theory $\mathsf{A}$ on $T_{\eta}$ and the lattice
theory $\mathsf{B}$ on $T_{\eta'}$, with $\eta'=\eta/L$.

\emph{Frame.} Let $(X',\mathfrak{A},\mathfrak{B}^{\ast})$ be a multiscale frame of the frame
construction, realised in $\mathsf{A}$ on $T_{\eta}$ and in $\mathsf{B}$ on $T_{\eta'}$, the two
realisations being exactly paired in the sense of the pairing statement for the two cutoffs;
write $\mathfrak{B}^{\ast\prime}$ for the realisation in $\mathsf{B}$. Let $\mathfrak{c}$ be the
set of data bonds of $\mathsf{A}$, which is also the set of paired data bonds of
$\mathsf{B}$, and let $\mathfrak{c}'_0$ be the set of collar bonds of $\mathsf{B}$, which
lie at level $-1$ of $\mathsf{A}$. Put $d_{\star}:=d^{\wedge}_{\mathfrak{B}^{\ast}}$, the
scaled distance of the realisation in $\mathsf{A}$. For a point $x\in X'^{\flat}$, that is,
a point of $X'$ off the staircase, $l_x$ denotes its level and $\ell_x:=\ell_{l_x}$.

\emph{Base.} Let $D$ be real on $\mathfrak{c}$ and $c_{\circ}$ real on $\mathfrak{c}'_0$,
such that $D\oplus c_{\circ}$ satisfies the regularity condition on real configurations of
the frame construction with constant $\varepsilon_D$, where
$\varepsilon_D\le\frac12\min\{a_H,a_A\}$; here $a_H$ is the radius of the relative response
statement, which bounds the regauged carried response minus its principal part, and is
distinct from the radius $b_H$ of the response function below, and $a_A$ is the radius of the
background-field representation. Put
$U:=U^{\mathsf{A}}(\mathfrak{B}^{\ast},D)$ and
$U':=U^{\mathsf{B}}(\mathfrak{B}^{\ast\prime},D\oplus c_{\circ})$, the latter rotated once
into the carried frame, so that $\mathcal{M}(U')=e^{i\eta A^{\natural}}U$, where
$\mathcal{M}$ denotes the one-step block-averaging map from $T_{\eta'}$ to $T_{\eta}$, and put
$\mathsf{w}_{\natural}:=\mathsf{w}[A^{\natural};U]$. Let $\mathcal{H}^{\mathsf{A}}(B)$ be
the response function about $U$ and $\mathcal{H}^{\mathsf{B}}(B\oplus B_c)$ the response
function about $U'$, for $B$ on $\mathfrak{c}$ and $B_c$ on $\mathfrak{c}'_0$. With the
carried form
\begin{equation*}
  Y[h;V]:=(i\eta)^{-1}\log\bigl[\mathcal{M}(e^{i\eta' h}V)\,\mathcal{M}(V)^{-1}\bigr],
\end{equation*}
define
\begin{equation}\label{5.7:eq-data-mismatch-c}
\begin{gathered}
  \mathsf{Y}(B,B_c):=Y\bigl[\mathcal{H}^{\mathsf{B}}(\operatorname{Ad}_{\mathsf{w}_{\natural}}B
  \oplus B_c);U'\bigr],\\
  a_{\Lambda}:=\tfrac14\min\bigl\{b_H,\,a_H,\,a_A/C_L,\,1/K_N^{\Lambda}\bigr\},\qquad
  K_N^{\Lambda}:=2400\,d\,L^4B_Hc_1^{b}\,(1+1/r_V),
\end{gathered}
\end{equation}
where $C_L\ge1$ is the constant for which the background-field representation holds both at
a real configuration $V$ satisfying the regularity condition with constant $\varepsilon$ and
at $e^{iB}V$, the latter under the regularity condition with constant
$\varepsilon+C_L\mathsf{b}$, $\mathsf{b}:=\sup|B|$, and $r_V$ is the constant of the transport
estimate for the dilated Wilson family whose arithmetic fixes $K_N^{\Lambda}$ (see the
verification below).

\emph{Paired data.} Let $(\mathfrak{b}^{\mathsf{B}},\mathfrak{b}^{\mathsf{A}})$ be a pair of
exactly paired fields of the zeroth data format of the frame construction, and let
$\mathcal{M}_{\ast}$ and $\mathcal{M}'_{\ast}$ denote the averaging maps of $\mathsf{A}$ and
$\mathsf{B}$ to the data level. The averaged field of $\mathsf{A}$ is
$V^{\mathsf{A}}:=\mathcal{M}_{\ast}\mathfrak{b}^{\mathsf{A}}=e^{iB_0^{\mathsf{A}}}D$, whose real part
$D:=\mathcal{M}_{\ast}\bar U^{\mathsf{A}}$, with $\bar U^{\mathsf{A}}$ the real configuration
attached to $\mathfrak{b}^{\mathsf{A}}$, satisfies the regularity condition and serves as the base
above, and $|B_0^{\mathsf{A}}|\le c_ALa_{\flat}$; both facts are supplied by the averaging
statement of the frame construction, and $c_A$, $a_{\flat}$ are the constants of that statement.
Likewise, by the same averaging statement applied to $\mathsf{B}$,
$\mathcal{M}'_{\ast}\mathfrak{b}^{\mathsf{B}}=e^{iB_0^{\mathsf{B}}}D^{\mathsf{B}}$ on
$\mathfrak{c}$ with $D^{\mathsf{B}}$ real; $D^{\mathsf{B}}$ and $B_0^{\mathsf{B}}$ are the real
part and the phase of the data of $\mathsf{B}$ in its own frame.
The data of $\mathsf{B}$, read in the carried frame and after a data gauge $g$, are
\begin{equation*}
  (\mathcal{M}'_{\ast}\mathfrak{b}^{\mathsf{B}})^{g}
  =e^{i\operatorname{Ad}_{\mathsf{w}_{\natural}}(B_0^{\mathsf{A}}+e_0)}D
  \ \text{on }\mathfrak{c},\qquad
  e^{iB_{c,0}}c_{\circ}\ \text{on }\mathfrak{c}'_0,
\end{equation*}
where $c_{\circ}$ is the part in $\mathrm{SU}(N)$ of the collar data of $\mathsf{B}$; this
defines $e_0$ on $\mathfrak{c}$, depending on $g$, and $B_{c,0}$ on $\mathfrak{c}'_0$. The
\emph{data mismatch} of the pair is
\begin{equation}\label{5.7:eq-data-mismatch-a}
  \mathsf{m}^{\mathrm{dat}}(\mathfrak{b}^{\mathsf{B}},\mathfrak{b}^{\mathsf{A}})
  :=\frac{4}{a_{\Lambda}}\,\inf_{g}\,\sup_{c\in\mathfrak{c}}|e_0(c)|,
\end{equation}
the infimum being taken over data gauges $g$.

\emph{Far data.} For $\mathsf{m}\in\{\mathsf{A},\mathsf{B}\}$ let
$e^{iB_1^{\mathsf{m}}}:=e^{i\mathfrak{h}^{\mathsf{m}}}e^{iB_0^{\mathsf{m}}}$, with
$\mathfrak{h}^{\mathsf{m}}$ the far data of $\mathsf{m}$ in the sense of the
response-and-regauging statement of the frame construction, and
$b^{\mathsf{m}}:=B_1^{\mathsf{m}}-B_0^{\mathsf{m}}$. The collar increment of $\mathsf{B}$ on
$\mathfrak{c}'_0$ is denoted $b_c$.

\emph{Standing condition.} The following one-sided condition is added to the conditions in
the order of choice of Definition~\ref{5.1:def-admitted-inputs}; once the auxiliary parameters
are chosen, it restricts $\gamma$. The outer supremum is taken over the classes of pairs to
which the closure condition on classes of pairs applies:
\begin{equation}\label{5.7:eq-data-mismatch-b}
\begin{gathered}
  \sup\Bigl(\sup|B_0^{\mathsf{m}}|+\sup|B_{c,0}|+1.1\,\sup|\mathfrak{h}^{\mathsf{m}}|\Bigr)
  \le\frac{a_{\Lambda}}{8},\\
  \varepsilon_D\le\tfrac12\min\{a_H,a_A\}.
\end{gathered}
\end{equation}
In the two settings of Lemma~\ref{5.7:lem-three-quarter-rate}, with $c_A$ as above and the
constants $K_w$, $\gamma_0^{\top}$ of the remainder bound of the energy-difference statement in
those settings, it reads
\begin{equation*}
\begin{gathered}
  c_A(K_w+1)\gamma_0^{\top}+33B_{\sharp}^2\sup_{j}\alpha_{\cdot,j}\le\frac{a_{\Lambda}}{8},\\
  (2+9c_A)(K_w+1)\gamma_0^{\top}\le\tfrac15\min\{a_H,a_A\}.
\end{gathered}
\end{equation*}
\end{definition}

\begin{proof}[Verification of the assertions accompanying the definition]
\emph{The frame.} The scaled distance of the realisation in $\mathsf{B}$ is no smaller than
$d_{\star}$, because the multiscale structure of $\mathsf{B}$ refines that of $\mathsf{A}$.
By the geometric lemma on the frames, $d_{\star}(x,c)\ge w\,l_x$ for every
$x\in X'^{\flat}$ and every $c\in\mathfrak{c}'_0$, where $w$ is the width constant of the
frames, the letter used for it in that lemma.

\emph{The carried response.} Multiplying the definition of the carried form by $i\eta$,
exponentiating and multiplying on the right by $\mathcal{M}(V)$ gives
$e^{i\eta Y[h;V]}\mathcal{M}(V)=\mathcal{M}(e^{i\eta'h}V)$. Applied with $V=U'$ and
$h=\mathcal{H}^{\mathsf{B}}(\operatorname{Ad}_{\mathsf{w}_{\natural}}B\oplus B_c)$, this is
\begin{equation*}
  e^{i\eta\mathsf{Y}}\mathcal{M}(U')
  =\mathcal{M}\bigl(e^{i\eta'\mathcal{H}^{\mathsf{B}}(\operatorname{Ad}_{\mathsf{w}_{\natural}}B
  \oplus B_c)}U'\bigr).
\end{equation*}

\emph{The radius.} The constant $K_N^{\Lambda}$ in \eqref{5.7:eq-data-mismatch-c} is chosen
so that on the polydisc $\mathfrak{D}(a_{\Lambda})$ the parallel transports of the native
canonical axial representatives $Y^{\mathsf{m}}(B)$ of the two cutoffs,
$\mathsf{m}\in\{\mathsf{A},\mathsf{B}\}$, along contours of at most $4d$ units of the
lattice's own spacing remain of norm at most $1.5$; the transport bounds then apply.
The arithmetic, through which $r_V$ enters $K_N^{\Lambda}$, is that of the corresponding
estimate in the setting of the dilated Wilson family, with the radius $c_Aa_{\mathfrak{b}}$ of
that setting replaced by $a_{\Lambda}$.

\emph{The base at the collar.} At the plaquettes joining $\mathfrak{c}'_0$ to
$\mathfrak{c}$,
\begin{equation*}
  |\partial(D\oplus c_{\circ})-1|\le\varepsilon^{\mathsf{B}}
  +4\,|D^{\mathsf{B}}D^{-1}-1|,
\end{equation*}
where $\varepsilon^{\mathsf{B}}$ is the constant of the regularity condition satisfied by the
data of $\mathsf{B}$ and $D^{\mathsf{B}}$ is the real part of those data introduced above.
Moreover
\begin{equation*}
  |D^{\mathsf{B}}D^{-1}-1|\le 1.1\bigl(\sup|B_0^{\mathsf{B}}|+\sup|B_0^{\mathsf{A}}|
  +\sup|e_0|\bigr).
\end{equation*}
For pairs whose data mismatch satisfies $\mathsf{m}^{\mathrm{dat}}\le\varpi$ with $\varpi<1$,
\eqref{5.7:eq-data-mismatch-a} provides a data gauge $g$ with
$\sup_{\mathfrak{c}}|e_0|<a_{\Lambda}/4$; for it the right-hand side of the first inequality
is, by the second, at most a number that is small with $\gamma$ plus $1.1a_{\Lambda}$. By
\eqref{5.7:eq-data-mismatch-b} and by $a_{\Lambda}\le\frac14\min\{a_H,a_A\}$, which follows
from \eqref{5.7:eq-data-mismatch-c} since $C_L\ge1$, the configuration $D\oplus c_{\circ}$
therefore satisfies the regularity condition also at these plaquettes with a constant
$\varepsilon_D\le\frac12\min\{a_H,a_A\}$.

\emph{Data gauges.} A data gauge moves the background fields $W^{\mathsf{B}}$ and
$W^{1,\mathsf{B}}$ by one common gauge transformation \cite[(181)]{B-CMP102b}. Every function
read in what follows, other than $e_0$ itself, is invariant under such a common transformation,
by the gauge-invariance clause of the carried frame and by clause (i) of the invariance lemma
for the dilated Wilson family; $e_0$ depends on $g$, which is why
\eqref{5.7:eq-data-mismatch-a} takes the infimum over data gauges.

\emph{The mismatch.} The quantity $\mathsf{m}^{\mathrm{dat}}$ vanishes at $\mathbb{1}$ and on
diagonal pairs. The closure condition on classes of pairs in which
$\mathsf{m}^{\mathrm{dat}}$ enters is of existential form in $\mathsf{m}^{\mathrm{dat}}$.

\emph{Far data.} For each $\mathsf{m}$, $|b^{\mathsf{m}}|\le1.1\,|\mathfrak{h}^{\mathsf{m}}|$.
Let $\mathfrak{h}_{\ast}$ be the common profile of Lemma~\ref{5.7:lem-three-quarter-rate}, and
suppose that $|\mathfrak{h}^{\mathsf{m}}(c)|\le\mathfrak{h}_{\ast}(c)$ for both $\mathsf{m}$
and that
$|\mathfrak{h}^{\mathsf{B}}-\mathfrak{h}^{\mathsf{A}}|(c)\le\varpi'\mathfrak{h}_{\ast}(c)$. The
Baker--Campbell--Hausdorff formula at small arguments then gives
\begin{equation*}
  |b^{\mathsf{B}}-b^{\mathsf{A}}|(c)\le\bigl(1.1\,\varpi'+2\,|e_0(c)|\bigr)
  \mathfrak{h}_{\ast}(c).
\end{equation*}
The collar increment satisfies $|b_c(c)|\le1.1\,\mathfrak{h}_{\ast}(c)$ on
$\mathfrak{c}'_0$ and has no partner in $\mathsf{A}$.

\emph{The standing condition.} The set of $\gamma$ for which the concrete form of
\eqref{5.7:eq-data-mismatch-b} holds is a down-set.
\end{proof}

\begin{proposition}[The two-cutoff response statement at zeroth order]
\label{5.7:prop-response-zeroth-order}
Let $\mathsf{A}$ and $\mathsf{B}$ be the two cutoffs, the lattices of spacing $\eta$ and $\eta'=\eta/L$,
and let $\mathsf{m}\in\{\mathsf{A},\mathsf{B}\}$ denote either of them. Assume the following.
\begin{enumerate}
\item[(a)] The frames of $\mathsf{A}$ and $\mathsf{B}$ are exactly paired, as in the paired frame of
Definition~\ref{5.7:def-data-mismatch}.
\item[(b)] In each cutoff the one-cutoff response bound for the increment $\mathcal{H}^{\mathsf{m}}$
about its base shifted by a field $e$ holds, with Lipschitz constant $8L^3B_H$, the stencil factor $8L^3$
times $B_H$. The paired frame satisfies the following conditions, for every $x\in X'^{\flat}$.
\begin{enumerate}
\item[(b1)] The kernel-sum bound:
\begin{equation*}
\sum_{c\in\mathfrak{c}\cup\mathfrak{c}'_0}e^{-\frac14\delta d_{\star}(x,c)}\le c_1^{b}.
\end{equation*}
\item[(b2)] The adjacent-cell bound: $e^{\delta d_{\star}(x,y)}\le2$ for every point $y$ in the cell
of $x$ or in a cell adjacent to it.
\item[(b3)] The collar-distance bound: $d_{\star}(x,c)\ge w\,l_x$ for every $c\in\mathfrak{c}'_0$,
where $w$ is the collar-depth constant of the frame.
\item[(b4)] The level condition: $e^{-\frac34\delta w\,l_x}\le L^{-12l_x}$.
\item[(b5)] The frame is in the setting of Lemma~\ref{5.7:lem-three-quarter-rate}, so that the
majorants at three quarters of the kernel rate hold.
\end{enumerate}
\item[(c)] The first clause of the relative response statement holds, in its field form.
\item[(d)] The two-constants lemma holds for a function holomorphic in a disc that is bounded on the
disc and bounded more sharply on its real diameter; part~(b) of
Lemma~\ref{5.7:lem-two-map-interpolation} is derived from it.
\item[(e)] The radius $a_{\Lambda}$ of \eqref{5.7:eq-data-mismatch-c} obeys the one-sided ceiling
\eqref{5.7:eq-data-mismatch-b}.
\item[(f)] The Lipschitz bound for the carried form $Y[\cdot\,;\cdot]$ holds: it reduces a difference
of the carried form at a bond to differences of its first argument at the fine points of the pair of
blocks of that bond.
\end{enumerate}
Put
\begin{gather*}
K_R := 136d\cdot 8L^3B_H,\qquad \vartheta := L^{-\theta_1/8},\qquad \delta' := \tfrac34\delta,\\
T_{\ast}(x) := \sum_{c\in\mathfrak{c}\cup\mathfrak{c}'_0}
e^{-\frac34\delta d_{\star}(x,c)}\,\mathfrak{h}_{\ast}(c),\\
c_1 := 28\cdot 1.1\,(c_1^{b})^{1/4}C_{\star}^{1/8},\qquad
C_{\star} := 1+\frac{2C_H(\varepsilon_D+a_{\Lambda})}{a_{\Lambda}K_R\,c_1^{b}}.
\end{gather*}
Here $\mathfrak{c}$ is the finite set of data bonds, $\mathfrak{c}'_0$ the set of collar bonds of
$\mathsf{B}$, $d_{\star}$ the scaled distance of the frame, $\mathfrak{h}_{\ast}$ the common profile of
the far data of the two cutoffs, $\mathsf{Y}$ the carried response of $\mathsf{B}$ of
\eqref{5.7:eq-data-mismatch-c}, defined there through the carried form $Y[\cdot\,;\cdot]$,
$B_H$ the response constant of hypothesis~(b), $c_1^{b}$ the kernel-sum constant of~(b1), and $C_H$,
$\varepsilon_D$, $\theta_1$ the constants of the relative response statement.
The sum $T_{\ast}(x)$ has the form of the sum in part~(ii) of
Lemma~\ref{5.7:lem-three-quarter-rate}, taken here over the index set
$\mathfrak{c}\cup\mathfrak{c}'_0$ and with the profile $\mathfrak{h}_{\ast}$ of the present frame;
throughout this proposition and its proof $T_{\ast}(x)$ denotes the sum displayed above.

Let $\mathfrak{b}^{\mathsf{A}},\mathfrak{b}^{\mathsf{B}}$ be the paired fields of the two cutoffs and let
$\varpi\ge\mathsf{m}^{\mathrm{dat}}(\mathfrak{b}^{\mathsf{B}},\mathfrak{b}^{\mathsf{A}})$, the data mismatch of
\eqref{5.7:eq-data-mismatch-a}. Assume moreover $\varpi<1$ (for $\varpi\ge1$ the bounds below follow by
subtracting the one-cutoff bounds), so that $\sup|e_0|<a_{\Lambda}/4$, where
$e_0$ is the mismatch of the base data of the two cutoffs.

With each cutoff carrying its own lattice spacing in the exponent, put
\begin{equation*}
W^{L,\mathsf{A}} := e^{i\eta\mathcal{H}^{\mathsf{A}}(B_0^{\mathsf{A}})}U,\qquad
W^{L,\mathsf{B}} := e^{i\eta'\mathcal{H}^{\mathsf{B}}(B_0^{\mathsf{B}})}U',
\end{equation*}
where $U$ and $U'$ are the base configurations of $\mathsf{A}$ and $\mathsf{B}$,
and define $\mathcal{X}^{\mathsf{A}}$ and $\mathcal{X}^{\mathsf{B}}$ by
\begin{align*}
e^{i\eta\mathcal{X}^{\mathsf{A}}}
&:= e^{i\eta\mathcal{H}^{\mathsf{A}}(B_1^{\mathsf{A}})}\,e^{-i\eta\mathcal{H}^{\mathsf{A}}(B_0^{\mathsf{A}})},\\
e^{i\eta'\mathcal{X}^{\mathsf{B}}}
&:= e^{i\eta'\mathcal{H}^{\mathsf{B}}(B_1^{\mathsf{B}})}\,e^{-i\eta'\mathcal{H}^{\mathsf{B}}(B_0^{\mathsf{B}})}.
\end{align*}
Thus $\mathcal{X}^{\mathsf{m}}$ is the sealed increment $\mathcal{H}_S^{\mathsf{m}}$ written in the
Landau frame, $\mathcal{H}_S=\operatorname{Ad}_{v(B_0)}\mathcal{X}$, where $\operatorname{Ad}_{v(B_0)}$
is the adjoint action of the regauging $v(B_0)$ that relates the Landau frame of the base $B_0$ to
the sealed frame.

For a coarse plaquette $q$ put
\begin{align*}
\mathsf{F}^{\mathsf{A}}(B)(q) &:= \eta^{-2}\bigl[\partial\bigl(e^{i\eta\mathcal{H}^{\mathsf{A}}(B)}U\bigr)(q)-1\bigr],\\
\mathsf{F}^{\mathsf{B}}(B)(q) &:= \eta^{-2}\bigl[\partial\mathcal{M}\bigl(e^{i\eta'\mathcal{H}^{\mathsf{B}}(B)}U'\bigr)(q)-1\bigr],
\end{align*}
where $\partial V(q)$ denotes the plaquette variable of a configuration $V$ at $q$, and where, here
and in the proof below, $\mathcal{M}$ denotes the one-step averaging map of
\cite[Propositions~6--7]{B-CMP98}, which carries configurations of $\mathsf{B}$ to configurations
of $\mathsf{A}$.

Let $x\in X'^{\flat}$, with level $l_x$ and length scale $\ell_x$. Consider all complex data
$(B_0^{\mathsf{m}},B_1^{\mathsf{m}})$ of the two cutoffs admitted in
Definition~\ref{5.7:def-data-mismatch}, and let $\varpi'\ge0$ be such that their far data
$\mathfrak{h}^{\mathsf{m}}$ obey, for every $c\in\mathfrak{c}\cup\mathfrak{c}'_0$,
\begin{equation*}
|\mathfrak{h}^{\mathsf{m}}(c)|\le\mathfrak{h}_{\ast}(c),\qquad
|\mathfrak{h}^{\mathsf{B}}(c)-\mathfrak{h}^{\mathsf{A}}(c)|\le\varpi'\,\mathfrak{h}_{\ast}(c).
\end{equation*}
Then the following three bounds hold, the third for every coarse plaquette $q$ based at $x$:
\begin{equation}\label{5.7:eq-response-zeroth-order-a}
\begin{aligned}
&\ell_x\bigl|[\mathsf{Y}(B_1^{\mathsf{B}})-\mathsf{Y}(B_0^{\mathsf{B}})]
-[\mathcal{H}^{\mathsf{A}}(B_1^{\mathsf{A}})-\mathcal{H}^{\mathsf{A}}(B_0^{\mathsf{A}})]\bigr|(x)\\
&\qquad\le K_R\bigl(2\varpi+1.1\varpi'+c_1\vartheta^{l_x}\bigr)T_{\ast}(x),\\
&\ell_x\bigl|Y[\mathcal{X}^{\mathsf{B}};W^{L,\mathsf{B}}]-\mathcal{X}^{\mathsf{A}}\bigr|(x)
\le 2K_R\bigl(2\varpi+1.1\varpi'+c_1\vartheta^{l_x}\bigr)T_{\ast}(x),\\
&\ell_x^2\bigl|[\mathsf{F}^{\mathsf{B}}(B_1^{\mathsf{B}})-\mathsf{F}^{\mathsf{B}}(B_0^{\mathsf{B}})]
-[\mathsf{F}^{\mathsf{A}}(B_1^{\mathsf{A}})-\mathsf{F}^{\mathsf{A}}(B_0^{\mathsf{A}})]\bigr|(q)\\
&\qquad\le K_F\bigl(2\varpi+1.1\varpi'+c_1'\vartheta^{l_x}\bigr)T_{\ast}(x).
\end{aligned}
\end{equation}
We call the three lines of this display the sup member, the averaged member and the curl member.
In the curl member
\begin{equation*}
K_F := C_M^{\mathrm{av}}L^2\cdot 6\cdot 8L^3B_H\cdot 2,
\end{equation*}
with $C_M^{\mathrm{av}}$ the constant of the averaging map $\mathcal{M}$,
and
\begin{equation*}
c_1' := 28\cdot1.1\,(c_1^{b})^{1/4}(C_{\star}')^{1/8},\qquad
C_{\star}' := 1+\frac{2C_H'(\varepsilon_D+a_{\Lambda})}{a_{\Lambda}K_F\,c_1^{b}},
\end{equation*}
where $C_H'$ is the constant that replaces $C_H$ in the curl member, fixed in the lemma on the
averaged and curl members of the two-cutoff response, in which that member is proved.
\end{proposition}

In \eqref{5.7:eq-response-zeroth-order-a} three parameters that could be left free, a rate
$\delta'\in(0,\delta]$, an index base and a mismatch parameter, are fixed as $\tfrac34\delta$,
$L^{-\theta_1/8}$ and $\mathsf{m}^{\mathrm{dat}}$. One-step averaging
is realised by the carried form $Y[\cdot\,;\cdot]$. The prefactor $8L^3B_H$ carries the number
$136d$ of \cite[(159)--(164)]{B-CMP98}. The sup member and the curl member together bound the pair
$\{\ell_x|\cdot|,\ \ell_x^2|\mathrm{curl}|\}$, which plays the role of a $C^1$ norm at scale
$\ell_x$ for the two-cutoff difference.

\begin{proof}[Proof of the sup member]
Fix $x\in X'^{\flat}$. We apply Lemma~\ref{5.7:lem-two-map-interpolation} on the index set
$\mathfrak{c}\cup\mathfrak{c}'_0$. The collar bonds $\mathfrak{c}'_0$ carry variables of $\mathsf{B}$ only.
This is admitted by \eqref{5.7:eq-fractional-power-a} of Remark~\ref{5.7:rem-fractional-power}:
$\Phi^{\mathsf{A}}$ does not depend on these
variables, so the one-cutoff Lipschitz hypothesis holds for $\Phi^{\mathsf{A}}$ there with any
kernel. We take $a:=a_{\Lambda}$. We take for $E$ the space of $\mathfrak{g}^{c}$-valued functions on
the bonds based at $x$, a finite-dimensional complex normed space. Finally we put
\begin{equation}\label{5.7:eq-response-zeroth-order-b}
\begin{gathered}
\Phi^{\mathsf{A}}_x(B,B_c) := \ell_x\mathcal{H}^{\mathsf{A}}(B)(x),\qquad
\Phi^{\mathsf{B}}_x(B,B_c) := \ell_x\mathsf{Y}(B,B_c)(x),\\
k(c) := K_Re^{-\delta d_{\star}(x,c)},\qquad \mathsf{K} := K_Rc_1^{b},\qquad
\varepsilon_{\star} := \min\{1,\,C_{\star}L^{-\theta_1l_x}\}.
\end{gathered}
\end{equation}
Write $\delta\Phi_x:=\Phi^{\mathsf{B}}_x-\Phi^{\mathsf{A}}_x$.

\emph{The kernel sum.} Since $e^{-\delta d_{\star}(x,c)}\le e^{-\frac14\delta d_{\star}(x,c)}$, the
kernel-sum bound of the frame gives
\begin{equation*}
\sum_c k(c)=K_R\sum_c e^{-\delta d_{\star}(x,c)}\le K_R\sum_c e^{-\frac14\delta d_{\star}(x,c)}
\le K_Rc_1^{b}=\mathsf{K}.
\end{equation*}
So $\mathsf{K}$ is an upper bound of the kernel sum, as Lemma~\ref{5.7:lem-two-map-interpolation}
requires.

\emph{The holomorphy hypothesis} of \eqref{5.7:eq-two-map-interpolation-a}. On each lattice,
$\Phi^{\mathsf{A}}_x$ and the response entering $\Phi^{\mathsf{B}}_x$ are holomorphic on
$\mathfrak{D}(a_{\Lambda})$ by \cite[Proposition~9]{B-CMP102b}. The averaging map $\mathcal{M}$ entering the
carried form preserves holomorphy by \cite[Propositions~6--7]{B-CMP98}. Hence $\Phi^{\mathsf{A}}_x$
and $\Phi^{\mathsf{B}}_x$ are holomorphic maps $\mathfrak{D}(a_{\Lambda})\to E$.

\emph{The one-cutoff Lipschitz hypothesis} of \eqref{5.7:eq-two-map-interpolation-a}.

For $\mathsf{A}$, this is the one-cutoff response bound of hypothesis~(b) at zero shift, $e=0$, with
Lipschitz constant $8L^3B_H$.

For $\mathsf{B}$, we proceed in four moves.
\begin{itemize}
\item The Lipschitz bound for the carried form, hypothesis~(f), reduces a difference of
$\Phi^{\mathsf{B}}_x$ to differences of $\mathcal{H}^{\mathsf{B}}$ at the fine points of the pair of
blocks of the bonds based at $x$.
\item The response bound of hypothesis~(b) in $\mathsf{B}$ is applied at these fine points.
\item These points lie in the cell of $x$ or in an adjacent cell. By the adjacent-cell bound, moving
the decay factors from them to $x$ costs at most a factor $e^{\delta d_{\star}(x,y)}\le2$ at such a
point $y$.
\item The frame regauging has norm $|\operatorname{Ad}_{\mathsf{w}_{\natural}}|=1$, and the scaled
distance of the frame of $\mathsf{B}$ dominates $d_{\star}$, $d_{\star}^{\mathsf{B}}\ge d_{\star}$. So the
decay factors of $\mathsf{B}$ are at most $e^{-\delta d_{\star}(x,c)}$.
\end{itemize}
This gives the one-cutoff Lipschitz hypothesis for $\mathsf{B}$ with the kernel $k$.

\emph{The real-slice closeness hypothesis} of \eqref{5.7:eq-two-map-interpolation-a}. Let
$(B,B_c)\in\mathfrak{D}(a_{\Lambda})$ be real. Then
\begin{equation*}
\delta\Phi_x(B,B_c)=\ell_x\mathsf{D}(B)(x)+\ell_x\bigl[\mathsf{Y}(B,B_c)-\mathsf{Y}(B,0)\bigr](x),
\end{equation*}
where $\mathsf{D}(B)$ is the relative response difference of the relative response statement.

The first clause of that statement applies, since its condition $\varepsilon_D+\mathsf{b}\le a_H$
holds, where $\mathsf{b}$ and $a_H$ are the quantities so named in that statement. It bounds the
first term by $2C_H(\varepsilon_D+a_{\Lambda})L^{-\theta_1l_x}$.

The second term changes only the collar variables, from $0$ to $B_c$ with
$\sup|B_c|<a_{\Lambda}$. We apply the one-cutoff Lipschitz bound for $\mathsf{B}$ just established,
then the collar-distance bound of the frame, $d_{\star}(x,c)\ge wl_x$ for $c\in\mathfrak{c}'_0$,
together with the kernel-sum bound, and finally the level condition:
\begin{align*}
\ell_x\bigl|\mathsf{Y}(B,B_c)-\mathsf{Y}(B,0)\bigr|(x)
&\le a_{\Lambda}K_R\sum_{c\in\mathfrak{c}'_0}e^{-\delta d_{\star}(x,c)}\\
&\le a_{\Lambda}K_Rc_1^{b}\,e^{-\frac34\delta wl_x}
\le a_{\Lambda}\mathsf{K}L^{-12l_x}.
\end{align*}
By the definition of $C_{\star}$,
\begin{equation*}
2C_H(\varepsilon_D+a_{\Lambda})=(C_{\star}-1)\,a_{\Lambda}K_Rc_1^{b}=(C_{\star}-1)\,a_{\Lambda}\mathsf{K}.
\end{equation*}
The two bounds therefore add up to
\begin{equation*}
|\delta\Phi_x(B,B_c)|\le a_{\Lambda}\mathsf{K}\bigl[(C_{\star}-1)L^{-\theta_1l_x}+L^{-12l_x}\bigr].
\end{equation*}
Hence the real-slice closeness hypothesis holds with the constant $\varepsilon_{\star}$ of
\eqref{5.7:eq-response-zeroth-order-b}, and $0\le\varepsilon_{\star}\le1$.

\emph{Assembly.} In the notation of the data, write on $\mathfrak{c}$
\begin{equation*}
B_0^{\mathsf{B}}=B_0^{\mathsf{A}}+e_0,\qquad B_1^{\mathsf{m}}=B_0^{\mathsf{m}}+b^{\mathsf{m}}.
\end{equation*}
Let $B_{c,0}$ and $B_{c,0}+b_c$ be the collar components of $B_0^{\mathsf{B}}$ and $B_1^{\mathsf{B}}$. We
apply part~(e) of Lemma~\ref{5.7:lem-two-map-interpolation}, displayed in
\eqref{5.7:eq-two-map-interpolation-b}, with $\theta:=\tfrac14$ and
\begin{equation*}
\hat B_0:=(B_0^{\mathsf{A}},B_{c,0}),\quad \hat e_0:=(e_0,0),\quad
\hat b^{\mathsf{A}}:=(b^{\mathsf{A}},b_c),\quad \hat b^{\mathsf{B}}:=(b^{\mathsf{B}},b_c).
\end{equation*}
Its condition on $\hat e_0$, namely $\sup|\hat e_0|=\sup|e_0|<a_{\Lambda}/4$, holds because
$\varpi<1$. With these choices,
\begin{gather*}
\Phi^{\mathsf{B}}_x(\hat B_0+\hat e_0+\hat b^{\mathsf{B}})=\ell_x\mathsf{Y}(B_1^{\mathsf{B}})(x),\qquad
\Phi^{\mathsf{B}}_x(\hat B_0+\hat e_0)=\ell_x\mathsf{Y}(B_0^{\mathsf{B}})(x),\\
\Phi^{\mathsf{A}}_x(\hat B_0+\hat b^{\mathsf{A}})=\ell_x\mathcal{H}^{\mathsf{A}}(B_1^{\mathsf{A}})(x),\qquad
\Phi^{\mathsf{A}}_x(\hat B_0)=\ell_x\mathcal{H}^{\mathsf{A}}(B_0^{\mathsf{A}})(x).
\end{gather*}
So $\Delta^{\times}$ is the quantity whose modulus stands on the left of the sup member, and
\begin{equation}\label{5.7:eq-response-zeroth-order-1}
|\Delta^{\times}|\le 28\,\varepsilon_{\star}^{1/8}\,T^{(3/4)}[\hat b^{\mathsf{A}}]
+\frac{4\sup|e_0|}{a_{\Lambda}}\,T^{(1)}[\hat b^{\mathsf{A}}]
+T^{(1)}[\hat b^{\mathsf{B}}-\hat b^{\mathsf{A}}].
\end{equation}

We bound the three weighted sums using $k(c)=K_Re^{-\delta d_{\star}(x,c)}$ and
$\mathsf{K}=K_Rc_1^{b}$, together with the bounds on the data.

For the first sum,
\begin{equation*}
T^{(3/4)}[\hat b^{\mathsf{A}}]
=K_R(c_1^{b})^{1/4}\sum_{c}e^{-\frac34\delta d_{\star}(x,c)}|\hat b^{\mathsf{A}}(c)|
\le 1.1\,K_R(c_1^{b})^{1/4}\,T_{\ast}(x).
\end{equation*}
For the second, since $e^{-\delta d_{\star}(x,c)}\le e^{-\frac34\delta d_{\star}(x,c)}$,
\begin{equation*}
T^{(1)}[\hat b^{\mathsf{A}}]\le 1.1\,K_R\,T_{\ast}(x).
\end{equation*}
For the third,
\begin{equation*}
T^{(1)}[\hat b^{\mathsf{B}}-\hat b^{\mathsf{A}}]\le K_R\bigl(1.1\varpi'+2\sup|e_0|\bigr)T_{\ast}(x).
\end{equation*}

Since $\varpi\ge\mathsf{m}^{\mathrm{dat}}(\mathfrak{b}^{\mathsf{B}},\mathfrak{b}^{\mathsf{A}})$, taking the infimum
over gauge transformations in the definition \eqref{5.7:eq-data-mismatch-a} of the data mismatch
gives $4\sup|e_0|/a_{\Lambda}\le\varpi$. Hence
\begin{equation*}
2\sup|e_0|\le\tfrac12a_{\Lambda}\varpi\le\tfrac12\varpi,
\end{equation*}
where the last inequality uses $a_{\Lambda}\le1$.
Inserting these bounds into \eqref{5.7:eq-response-zeroth-order-1} gives
\begin{equation*}
|\Delta^{\times}|\le K_R\Bigl[28\cdot1.1\,(c_1^{b})^{1/4}\varepsilon_{\star}^{1/8}
+1.1\varpi+1.1\varpi'+\tfrac12\varpi\Bigr]T_{\ast}(x).
\end{equation*}

It remains to compare with the right side of the sup member. First, $1.1\varpi+\tfrac12\varpi\le2\varpi$.
Second, $\varepsilon_{\star}\le C_{\star}L^{-\theta_1l_x}$, so
\begin{equation*}
28\cdot1.1\,(c_1^{b})^{1/4}\varepsilon_{\star}^{1/8}
\le 28\cdot1.1\,(c_1^{b})^{1/4}C_{\star}^{1/8}L^{-\theta_1l_x/8}=c_1\vartheta^{l_x}.
\end{equation*}
This is the sup member of \eqref{5.7:eq-response-zeroth-order-a}.
\end{proof}

The averaged member and the curl member of \eqref{5.7:eq-response-zeroth-order-a} are proved in
the lemma on the averaged and curl members of the two-cutoff response, which also fixes the
constant $C_H'$.

\begin{proof}[Proof of the averaged member and the curl member]
\label{5.7:prf-averaged-curl-members}
We prove the averaged member (S$'$) and the curl member (C) of the two-cutoff response statement
at zeroth order, Proposition~\ref{5.7:prop-response-zeroth-order}, displayed in
\eqref{5.7:eq-response-zeroth-order-a}. The hypotheses are those of that proposition, and we use
its sup member (S), whose proof precedes this one. We use the interpolation between two analytic
maps, Lemma~\ref{5.7:lem-two-map-interpolation}, whose three hypotheses are displayed in
\eqref{5.7:eq-two-map-interpolation-a}: analyticity on the polydisc, the first-order bound
with a kernel $k$, and the closeness on the real points of the polydisc with constant
$\varepsilon_{\star}$.

Fix $x\in X'^{\flat}$ and complex data as in the proposition. For $\iota=0,1$ write
$\mathsf{Y}_{\iota}:=\mathsf{Y}(B_{\iota}^{\mathsf{B}})$ and
$\mathcal{H}^{\mathsf{A}}_{\iota}:=\mathcal{H}^{\mathsf{A}}(B_{\iota}^{\mathsf{A}})$. As in the
definition of $\mathsf{Y}$ recalled in \eqref{5.7:eq-data-mismatch-c}, write
$\mathcal{H}^{\mathsf{B}}(B_{\iota})$ for the response of $\mathsf{B}$ that enters
$\mathsf{Y}_{\iota}$, so that $\mathsf{Y}_{\iota}=Y[\mathcal{H}^{\mathsf{B}}(B_{\iota});U']$, where
\begin{equation*}
Y[h;V]=(i\eta)^{-1}\log\bigl[M(e^{i\eta' h}V)\,M(V)^{-1}\bigr].
\end{equation*}

\emph{The averaged member} (S$'$). At the cutoff $\mathsf{B}$, whose lattice has spacing $\eta'$
(so that in $\mathsf{B}$ the exponent $\eta$ of (S$'$) is the spacing $\eta'$ of its lattice),
the definitions of $W^{L,\mathsf{B}}$ and $\mathcal{X}^{\mathsf{B}}$ in (S$'$) give
$e^{i\eta'\mathcal{X}^{\mathsf{B}}}W^{L,\mathsf{B}}=e^{i\eta'\mathcal{H}^{\mathsf{B}}(B_1)}U'$.
The definition of $Y$ then gives
\begin{equation}\label{5.7:eq-averaged-curl-members-1}
e^{i\eta Y[\mathcal{X}^{\mathsf{B}};W^{L,\mathsf{B}}]}
= M\bigl(e^{i\eta'\mathcal{H}^{\mathsf{B}}(B_1)}U'\bigr)\,
M\bigl(e^{i\eta'\mathcal{H}^{\mathsf{B}}(B_0)}U'\bigr)^{-1}
= e^{i\eta\mathsf{Y}_1}e^{-i\eta\mathsf{Y}_0}.
\end{equation}
The second equality holds because
\begin{equation*}
e^{i\eta\mathsf{Y}_{\iota}}=M\bigl(e^{i\eta'\mathcal{H}^{\mathsf{B}}(B_{\iota})}U'\bigr)\,M(U')^{-1}
\qquad(\iota=0,1),
\end{equation*}
so that the product $e^{i\eta\mathsf{Y}_1}e^{-i\eta\mathsf{Y}_0}$ contains the factor
$M(U')^{-1}M(U')=\mathbb{1}$.

Take logarithms in \eqref{5.7:eq-averaged-curl-members-1} and apply the Baker--Campbell--Hausdorff
formula. This gives
\begin{equation}\label{5.7:eq-averaged-curl-members-2}
Y[\mathcal{X}^{\mathsf{B}};W^{L,\mathsf{B}}]
=(i\eta)^{-1}\log\bigl(e^{i\eta\mathsf{Y}_1}e^{-i\eta\mathsf{Y}_0}\bigr)
=(\mathsf{Y}_1-\mathsf{Y}_0)+\mathsf{c}[\mathsf{Y}_1-\mathsf{Y}_0,\mathsf{Y}_0].
\end{equation}
Here $\mathsf{c}$ is the sum of the brackets of order at least two. Write each entry $\mathsf{Y}_1$
as $(\mathsf{Y}_1-\mathsf{Y}_0)+\mathsf{Y}_0$ and expand. A bracket all of whose entries equal
$\mathsf{Y}_0$ vanishes; for instance $[\mathsf{Y}_1,\mathsf{Y}_0]=[\mathsf{Y}_1-\mathsf{Y}_0,\mathsf{Y}_0]$.
Hence every bracket of $\mathsf{c}$ carries $\mathsf{Y}_1-\mathsf{Y}_0$ as an entry. After division
by $i\eta$, it carries besides that entry a factor $\eta\mathsf{Y}_0$ or
$\eta(\mathsf{Y}_1-\mathsf{Y}_0)$.

In the same way, the definition of $\mathcal{X}^{\mathsf{A}}$ gives
\begin{equation}\label{5.7:eq-averaged-curl-members-3}
\mathcal{X}^{\mathsf{A}}
=(\mathcal{H}^{\mathsf{A}}_1-\mathcal{H}^{\mathsf{A}}_0)
+\mathsf{c}[\mathcal{H}^{\mathsf{A}}_1-\mathcal{H}^{\mathsf{A}}_0,\mathcal{H}^{\mathsf{A}}_0],
\end{equation}
with the same series $\mathsf{c}$. Subtract \eqref{5.7:eq-averaged-curl-members-3} from
\eqref{5.7:eq-averaged-curl-members-2}. At $x$, the difference
$(\mathsf{Y}_1-\mathsf{Y}_0)-(\mathcal{H}^{\mathsf{A}}_1-\mathcal{H}^{\mathsf{A}}_0)$, multiplied by
$\ell_x$, is the left member of (S). By (S) it is at most
$K_R(2\varpi+1.1\varpi'+c_1\vartheta^{l_x})T_{\ast}(x)$.

It remains to bound the difference of the two series. Each bracket is multilinear in its entries.
Replace the entries of the first series one at a time by those of the second, and sum the series.
At $x$ this gives
\begin{equation}\label{5.7:eq-averaged-curl-members-4}
\begin{split}
&\bigl|\mathsf{c}[\mathsf{Y}_1-\mathsf{Y}_0,\mathsf{Y}_0]
-\mathsf{c}[\mathcal{H}^{\mathsf{A}}_1-\mathcal{H}^{\mathsf{A}}_0,\mathcal{H}^{\mathsf{A}}_0]\bigr|\\
&\quad\le 2\eta\bigl(|\mathsf{Y}_0|+|\mathcal{H}^{\mathsf{A}}_0|+\cdots\bigr)\,
\bigl|(\mathsf{Y}_1-\mathsf{Y}_0)-(\mathcal{H}^{\mathsf{A}}_1-\mathcal{H}^{\mathsf{A}}_0)\bigr|\\
&\qquad+2\eta\,|\mathsf{Y}_0-\mathcal{H}^{\mathsf{A}}_0|\,
|\mathcal{H}^{\mathsf{A}}_1-\mathcal{H}^{\mathsf{A}}_0|.
\end{split}
\end{equation}
Here the dots stand for the contributions of the brackets of higher order; by the preceding
paragraph, each of them carries further factors $\eta\mathsf{Y}_0$ or
$\eta(\mathsf{Y}_1-\mathsf{Y}_0)$, respectively $\eta\mathcal{H}^{\mathsf{A}}_0$ or
$\eta(\mathcal{H}^{\mathsf{A}}_1-\mathcal{H}^{\mathsf{A}}_0)$.
Four bounds enter the right-hand side, all at $x$:
\begin{itemize}
\item $\eta|\mathsf{Y}_0|$ and $\eta|\mathcal{H}^{\mathsf{A}}_0|$ are each at most $L^{-l_x}$ times
a number not exceeding one;
\item Lemma~\ref{5.7:lem-two-map-interpolation}(a), together with the first-order bound, gives
$\eta|\mathsf{Y}_0-\mathcal{H}^{\mathsf{A}}_0|\le L^{-l_x}\cdot 3a_{\Lambda}\mathsf{K}$;
\item as in the proof of (S), the first-order bound for $\mathsf{A}$ and the estimate
$T^{(1)}[b^{\mathsf{A}}]\le 1.1K_RT_{\ast}(x)$ give
$\ell_x|\mathcal{H}^{\mathsf{A}}_1-\mathcal{H}^{\mathsf{A}}_0|\le 1.1K_RT_{\ast}(x)$.
\end{itemize}
By the first item, the displayed factors $\eta|\mathsf{Y}_0|$ and $\eta|\mathcal{H}^{\mathsf{A}}_0|$
of the prefactor $2\eta(|\mathsf{Y}_0|+|\mathcal{H}^{\mathsf{A}}_0|+\cdots)$ are at most
$L^{-l_x}$. The first line on the right of \eqref{5.7:eq-averaged-curl-members-4}, multiplied by
$\ell_x$, is this prefactor times the left member of (S); by (S) it is therefore at most the
prefactor times $K_R(2\varpi+1.1\varpi'+c_1\vartheta^{l_x})T_{\ast}(x)$. By the second and third
items, the last line, multiplied by $\ell_x$, is at most
\begin{align*}
2\cdot 3a_{\Lambda}\mathsf{K}L^{-l_x}\cdot 1.1K_RT_{\ast}(x)
&=6.6\,a_{\Lambda}\mathsf{K}\,L^{-l_x}K_RT_{\ast}(x)\\
&\le 6.6\,a_{\Lambda}\mathsf{K}\,\vartheta^{l_x}K_RT_{\ast}(x),
\end{align*}
where $L^{-l_x}\le\vartheta^{l_x}$ was used. Both contributions are contained in
$K_R(2\varpi+1.1\varpi'+c_1\vartheta^{l_x})T_{\ast}(x)$; written out, the inequality used is
\begin{equation*}
\begin{split}
&2\eta\bigl(|\mathsf{Y}_0|+|\mathcal{H}^{\mathsf{A}}_0|+\cdots\bigr)\,
K_R(2\varpi+1.1\varpi'+c_1\vartheta^{l_x})T_{\ast}(x)
+6.6\,a_{\Lambda}\mathsf{K}\,\vartheta^{l_x}K_RT_{\ast}(x)\\
&\quad\le K_R(2\varpi+1.1\varpi'+c_1\vartheta^{l_x})T_{\ast}(x),
\end{split}
\end{equation*}
in which the factors of the prefactor are bounded by the first item and the second term is placed
in the term $c_1\vartheta^{l_x}$. Hence $\ell_x$ times the difference of the two series is at most
$K_R(2\varpi+1.1\varpi'+c_1\vartheta^{l_x})T_{\ast}(x)$. Adding the bound (S) for the first
difference gives
\begin{equation*}
\ell_x\bigl|Y[\mathcal{X}^{\mathsf{B}};W^{L,\mathsf{B}}]-\mathcal{X}^{\mathsf{A}}\bigr|(x)
\le 2K_R(2\varpi+1.1\varpi'+c_1\vartheta^{l_x})T_{\ast}(x),
\end{equation*}
which is (S$'$): the bound of (S) with $K_R$ replaced by $2K_R$.

\emph{The curl member} (C). Fix a coarse plaquette $q$ based at $x$. Apply
Lemma~\ref{5.7:lem-two-map-interpolation} with
\begin{equation*}
\Phi^{\mathsf{m}}:=\ell_x^2\,\mathsf{F}^{\mathsf{m}}(\cdot)(q)\quad(\mathsf{m}\in\{\mathsf{A},\mathsf{B}\}),
\qquad k(c):=K_F\,e^{-\delta d_{\star}(x,c)},
\end{equation*}
on the index set $\mathfrak{c}\cup\mathfrak{c}'_0$ and the polydisc $\mathfrak{D}(a_{\Lambda})$,
with $\mathsf{K}:=K_Fc_1^{b}$. This value of $\mathsf{K}$ is the value entering $c_1'$.

First-order bound, cutoff $\mathsf{A}$. Let $\tilde B$ be a second datum and consider the
segment from $B$ to $\tilde B$. Define
\begin{equation*}
h_q:=8L^3B_H\sum_{c\in\mathfrak{c}\cup\mathfrak{c}'_0}e^{-\delta d_{\star}(x,c)}|\tilde B(c)-B(c)|,
\end{equation*}
and let $\Xi_q$ be the ratio of the plaquette variables
$\partial(e^{i\eta\mathcal{H}^{\mathsf{A}}(\tilde B)}U)(q)$ and
$\partial(e^{i\eta\mathcal{H}^{\mathsf{A}}(B)}U)(q)$. Apply the plaquette-ratio identity, which
expresses this ratio along the segment, and hypothesis~(b) of
Proposition~\ref{5.7:prop-response-zeroth-order} at the two endpoints of the segment, that is,
with its shift $e$ equal to $0$ and to $\tilde B-B$. This gives
\begin{equation*}
\begin{split}
|\Xi_q-1|
&\le 5h_q\eta^2\ell_x^{-2}
+8h_q\bigl|\partial(e^{i\eta\mathcal{H}^{\mathsf{A}}(B)}U)(q)-1\bigr|\\
&\le 6h_q\eta^2\ell_x^{-2}.
\end{split}
\end{equation*}
This is the first-order bound for $\Phi^{\mathsf{A}}$ with the kernel $k$; the factor
$6\cdot 8L^3B_H$ is a factor of $K_F$.

First-order bound, cutoff $\mathsf{B}$. The plaquette variable of a one-step average is an analytic
local function of the fine plaquette variables and of the transports. Each averaging rescales
plaquette bounds by $L^2$, up to the constant $C_M^{\mathrm{av}}$. This is
\cite[Proposition~2 and (47)--(51)]{B-CMP98}, applied to complex data, which is where the constant
$C_M^{\mathrm{av}}$ is fixed, together with the locality statement for averaged plaquette variables.
The argument given for $\mathsf{A}$ is then repeated in $\mathsf{B}$ on the fine plaquettes. The factor
$C_M^{\mathrm{av}}L^2$ is a factor of $K_F$.

Closeness on real points. Let $(B,0)$ be real and put $\mathsf{b}:=\sup|B|$. The regauging
statement for the once-averaged response rests on an exact identity, which in the carried frame
reads
\begin{equation*}
M(e^{i\eta'\mathbb{H}'}U')=(e^{i\eta\tilde A}e^{i\eta\mathbb{H}}U)^{w}.
\end{equation*}
In the notation of that statement, $\mathbb{H}$ and $\mathbb{H}'$ are the responses of
$\mathsf{A}$ and $\mathsf{B}$ at the real datum $B$, so that the plaquette variables of
$e^{i\eta\mathbb{H}}U$ and of the left member are those entering $\mathsf{F}^{\mathsf{A}}(B)(q)$
and $\mathsf{F}^{\mathsf{B}}(B)(q)$; further $\tilde A:=R(a^{-1})A_B$ is the potential $A_B$ of
that statement rotated by its frame element $a$, and $w=e^{i\lambda}$ is the regauging of
that statement. The plaquette variable of a regauged configuration is conjugated by the value of the
gauge function at the base point, so
\begin{equation*}
\partial M(e^{i\eta'\mathbb{H}'}U')(q)
=\operatorname{Ad}_{w(x)}\partial(e^{i\eta\tilde A}\mathsf{X})(q),
\qquad \mathsf{X}:=e^{i\eta\mathbb{H}}U.
\end{equation*}
For $P:=\partial(e^{i\eta\tilde A}\mathsf{X})(q)$ and unitary $w(x)$ (the data being real),
\begin{equation*}
w(x)P-Pw(x)=(w(x)-1)(P-1)-(P-1)(w(x)-1).
\end{equation*}
Hence $|\operatorname{Ad}_{w(x)}P-P|\le 2|w(x)-1|\,|P-1|$. Also
$|P-1|\le|\partial\mathsf{X}(q)-1|+|P-\partial\mathsf{X}(q)|$. Therefore
\begin{equation*}
\begin{split}
\eta^{-2}\bigl|\partial M(e^{i\eta'\mathbb{H}'}U')(q)-\partial\mathsf{X}(q)\bigr|
&\le 2|w(x)-1|\,\eta^{-2}\bigl(|\partial\mathsf{X}(q)-1|
+|P-\partial\mathsf{X}(q)|\bigr)\\
&\quad+\eta^{-2}\bigl|\partial(e^{i\eta\tilde A}\mathsf{X})(q)-\partial\mathsf{X}(q)\bigr|.
\end{split}
\end{equation*}
Three bounds enter the right-hand side.
\begin{itemize}
\item By the regauging statement,
\begin{equation*}
|w(x)-1|\le 1.1|\lambda|\le 1.1C_H(\varepsilon_D+\mathsf{b})L^{-\theta_1l_x}.
\end{equation*}
\item $\eta^{-2}|\partial\mathsf{X}(q)-1|\le C(\varepsilon_D+\mathsf{b})\ell_x^{-2}$, with a
constant $C$.
\item Apply the plaquette-ratio identity, together with the bound
$\|A\|_{\theta}\le C_A\varepsilon\eta^{\theta}$ on the potential $A$ of the regauged configuration
relative to the block configuration, at the datum $e^{iB}D$; here $\theta$ is the regularity
exponent of that bound and of the regauging statement. That datum satisfies the hypothesis of the
bound with $\varepsilon_D+C_L\mathsf{b}$ in place of its parameter $\varepsilon$, where $C_L$ is a
constant, and it is the
gradient member of $\|A_B\|_{\theta}$, $\ell_x^2|\nabla A_B|$, that is used here. This gives
\begin{equation*}
\eta^{-2}\bigl|\partial(e^{i\eta\tilde A}\mathsf{X})(q)-\partial\mathsf{X}(q)\bigr|
\le 5\ell_x^{-2}C_A(\varepsilon_D+C_L\mathsf{b})L^{-\theta l_x}.
\end{equation*}
\end{itemize}
So the closeness on real points holds with $C_H':=C(d,L)(C_H+C_A)$ in place of $C_H$. Adjoin
collar data $B_c$ as in the proof of (S): let $\mathsf{F}^{\mathsf{B}}(B,B_c)(q)$ be formed as
$\mathsf{F}^{\mathsf{B}}(B)(q)$, with the response of $\mathsf{B}$ taken at $(B,B_c)$ as in
$\mathsf{Y}(B,B_c)$, and put $\Phi^{\mathsf{B}}(B,B_c):=\ell_x^2\mathsf{F}^{\mathsf{B}}(B,B_c)(q)$.
The collar contribution $\Phi^{\mathsf{B}}(B,B_c)-\Phi^{\mathsf{B}}(B,0)$ is bounded by
the first-order bound on the collar bonds, as in the proof of (S) with the data of
\eqref{5.7:eq-response-zeroth-order-b}.

Finally, apply Lemma~\ref{5.7:lem-two-map-interpolation}(e), with the interpolation exponent of its
parts (c) and (e) equal to $\tfrac14$ and with
$B_0$, $e_0$, $b^{\mathsf{A}}$, $b^{\mathsf{B}}$ chosen as in the proof of (S). It yields (C) with
the constant $c_1'$, which is $c_1$ with $(K_R,C_H)$ replaced by $(K_F,C_H')$.
\end{proof}

\begin{remark}[The derivative read through plaquettes only]\label{5.7:rem-plaquette-reading}
In the relative response statement, with $U$ the base configuration, $D_U$ the covariant difference
in $U$, $\lambda$ the gauge function of the regauging, $\mathsf{e}$ the remainder term of the
regauged response in that statement, and $\eta$ the spacing of the lattice on which that
statement is read, the relative response difference is
\begin{equation}\label{5.7:eq-plaquette-reading-1}
\mathsf{D}=\mathsf{e}-D_U\lambda .
\end{equation}
\begin{enumerate}
\item[(a)] The symmetric part of $\nabla\mathsf{D}$ contains $\nabla D_U\lambda$, which is a second
difference of the gauge function $\lambda$. The relative response statement bounds this term only
crudely, $|\nabla D_U\lambda|\le 2\eta^{-1}|\nabla\lambda|$; in its proof the term occurs once,
inside the commutator $[\lambda,\nabla D_U\lambda]$, bounded there by
$2\eta^{-1}|\lambda|\,|\nabla\lambda|$. At a point $x$ off the staircase, of level $l_x$ and
length scale $\ell_x$ as in Proposition~\ref{5.7:prop-response-zeroth-order}, neither bound supplies
a power $\vartheta^{l_x}$ of the index base $\vartheta$ at the weight $\ell_x^{-2}$, and the term
is seen by no gauge-invariant object.
\item[(b)] The gradient $\nabla\mathcal{H}_S$ of the sealed increment enters the later estimates
only through the plaquette variables $\Xi_q$ of the sealed increment, one for each plaquette $q$.
First-order control at the scale $\ell_x$ is asserted in \eqref{5.7:eq-response-zeroth-order-a}
of the two-cutoff response statement at zeroth order
(Proposition~\ref{5.7:prop-response-zeroth-order}) only as the pair $\ell_x|\cdot|$ (the sup
member) and $\ell_x^{2}|\mathrm{curl}\,\cdot|$ (the curl member). The analogue of the curl member
with $\ell_x^{2}|\nabla\,\cdot|$ in place of $\ell_x^{2}|\mathrm{curl}\,\cdot|$ is not asserted.
\item[(c)] No auxiliary gauge is used to compare the two cutoffs $\mathsf{A}$ and $\mathsf{B}$.
The comparison uses only canonical objects built from the native canonical axial representatives
$Y^{\mathsf{m}}(B)$: holonomies along loops based at a point $y_0$ common to the two lattices,
determined modulo one constant gauge transformation, and class functions of them.
\end{enumerate}
\end{remark}

\begin{proof}
(a) From \eqref{5.7:eq-plaquette-reading-1} we have
\begin{equation*}
\nabla\mathsf{D}=\nabla\mathsf{e}-\nabla D_U\lambda .
\end{equation*}
The second term is a difference of the covariant difference $D_U\lambda$, that is, a second
difference of $\lambda$. Its symmetric part contributes to that of $\nabla\mathsf{D}$.
The relative response statement bounds this term only crudely:
\begin{equation*}
|\nabla D_U\lambda|\le 2\eta^{-1}|\nabla\lambda| .
\end{equation*}

In the proof of the relative response statement the one second difference that occurs is
$[\lambda,\nabla D_U\lambda]$, and it is bounded there only crudely, by
\begin{equation*}
|[\lambda,\nabla D_U\lambda]|\le 2\eta^{-1}|\lambda|\,|\nabla\lambda| .
\end{equation*}
Neither bound supplies a power $\vartheta^{l_x}$ of the index base $\vartheta$ at the weight
$\ell_x^{-2}$. Finally, $-D_U\lambda$ is the term of $\mathsf{D}$ produced by the regauging with
gauge function $\lambda$ in the relative response statement, and no gauge-invariant object sees it.

(b) The statements that use $\nabla\mathcal{H}_S$ depend on it only through the variables
$\Xi_q$, that is, through plaquettes. This is the case in the clause of the statement on the
dilated Wilson family $\mathfrak{F}_v(s)$ that involves $\nabla\mathcal{H}_S$, and in clause (iii)
of the sixth-family lemma, where the quantity estimated is expressed through the variables
$\Xi_q$, one for each plaquette $q$ of a family of plaquettes.

The step of the proof of the flat energy-difference bound that bounds the datum $\mathsf{E}$, and
the corresponding step of the proof of the difference bound from which the curl member is taken,
use $\mathcal{H}_S$ only through $\sup|\mathcal{H}_S|$. They do so through the estimate
\begin{equation*}
|\mathfrak{h}_S(c)|\le 2c_A\,\ell_{l(c)}\sup|\mathcal{H}_S|
\end{equation*}
for the far datum $\mathfrak{h}_S(c)$ of the sealed increment at a data bond $c$ of level $l(c)$,
where $c_A$ is a fixed constant.
Thus every occurrence is covered by $\ell_x|\cdot|$ and $\ell_x^{2}|\mathrm{curl}\,\cdot|$. In the
proof of the averaged member and the curl member of
Proposition~\ref{5.7:prop-response-zeroth-order},
first-order control at the scale $\ell_x$ is used in exactly this sense.

(c) The construction of $Y^{\mathsf{m}}(B)$ yields, as its outputs, holonomies along loops based
at $y_0$, determined modulo one constant gauge transformation, and class functions of them. These
objects are defined without reference to any further choice of gauge.
\end{proof}

\begin{definition}[Canonical loops and class functions on two lattices]\label{5.7:not-canonical-loops}
We work in the frame at level $j$ with block $\square$ and set $F_{\square}$, and in the frame at
level $i$, that is, in the setting of the majorants at three quarters of the kernel rate
(Lemma~\ref{5.7:lem-three-quarter-rate}). For $\mathsf{m}\in\{\mathsf{A},\mathsf{B}\}$, the two
cutoffs, $\mathsf{A}$ on the lattice $T_{\eta}$ and $\mathsf{B}$ on the lattice $T_{\eta'}$ with
$\eta'=\eta/L$, which we call the two \emph{machines}, the
\emph{native function} of machine $\mathsf{m}$ is
\begin{equation}\label{5.7:eq-canonical-loops-1}
Y^{\mathsf{m}}(B) := U^{\mathsf{m}}\bigl(\mathfrak{B}^{\star}, e^{iB}D\bigr),
\end{equation}
Ba{\l}aban's canonical axial representative, where $D$ is the real base data of
Definition~\ref{5.7:def-data-mismatch}. For $\mathsf{B}$ the argument is the pair
$(B,B_c)$ of data and collar variables. In the frame at level $j$ the one-machine form of
\eqref{5.7:eq-canonical-loops-1} reads
\begin{equation}\label{5.7:eq-canonical-loops-2}
W = Y(B_0), \qquad U^{u}_{\mathsf{P}^1} = Y(B_1), \qquad
Y(B) := U\bigl(\mathfrak{B}^{\star}_j, e^{iB}V_{00}\bigr),
\end{equation}
where $V_{00}$ is the real base data of the frame at level $j$, $W$ is its base configuration
and $U^{u}_{\mathsf{P}^1}$ is its far-data configuration. Under the hypothesis of
Proposition~\ref{5.7:prop-response-zeroth-order} on the response functions
$\mathcal{H}^{\mathsf{A}}$ and $\mathcal{H}^{\mathsf{B}}$ of the two machines, $Y$ is analytic
in $B$ on a ball containing $B_0$ and $B_1$.

Proposition~\ref{5.7:prop-dilated-action-family} and Theorem~\ref{5.7:thm-remainder-value} use
no carried frame, and in them common data means equal data in Ba{\l}aban's own frames of the two
machines; Corollary~\ref{5.7:cor-chart-data-control}, drawn from
Proposition~\ref{5.7:prop-response-zeroth-order}, by contrast, is stated in the frame of that
proposition, in which the frame regauging $\mathsf{w}_{\natural}$ is one more data gauge.

In Ba{\l}aban's own frames of the two machines, after one data gauge, we put
$e_0 := B_0^{\mathsf{B}} - B_0^{\mathsf{A}}$ and let $\varpi$ be any
number with
\begin{equation}\label{5.7:eq-canonical-loops-3}
\varpi \ge \frac{4}{a_{\Lambda}}\,\sup|e_0| .
\end{equation}
The far data $\mathfrak{h}^{\mathsf{m}} := \mathsf{E}^{u,\mathsf{m}}$ determine
$B_1^{\mathsf{m}}$ through $e^{iB_1^{\mathsf{m}}} = e^{i\mathfrak{h}^{\mathsf{m}}}e^{iB_0^{\mathsf{m}}}$,
and we put $b^{\mathsf{m}} := B_1^{\mathsf{m}} - B_0^{\mathsf{m}}$. On $F_{\square}$ we put
$\mathfrak{h}_{\ast}(c) := 15B_{\sharp}^2 f(\bar c)$, where $f$ is the profile of
Lemma~\ref{5.7:lem-three-quarter-rate}, and, writing
$\delta\Phi := \Phi^{\mathsf{B}}-\Phi^{\mathsf{A}}$ for the two-machine difference of a
functional $\Phi$, we let $\varpi'$ be any number with
\begin{equation}\label{5.7:eq-canonical-loops-4}
|\delta\mathsf{E}^{u}(c)| \le \varpi'\,\mathfrak{h}_{\ast}(c), \qquad c \in F_{\square}.
\end{equation}
The sum $T_{\ast}$ is the profile sum of Lemma~\ref{5.7:lem-three-quarter-rate} at rate
$\tfrac34\delta$, with the collar bonds included among the cells summed over.

All functions of the two machines considered below are of one of two kinds. Either they are loops
based at one lattice point common to the two lattices, or they are class functions. In both
cases they are canonical, in the sense fixed in the canonicity lemma for the chart datum
$\Delta^{\natural}$.
\end{definition}

\begin{proof}
Two assertions above need proof: the analyticity of $Y$ asserted after
\eqref{5.7:eq-canonical-loops-2}, and the canonicity asserted in the last paragraph.

\emph{Analyticity.} The function $Y$ is Ba{\l}aban's canonical axial representative evaluated at
$e^{iB}V_{00}$, and $V_{00}$ is real. Its analyticity in $B$ on a ball containing
$B_0$ and $B_1$ is \cite[Proposition~9]{B-CMP102b}, applied under the hypothesis of
Proposition~\ref{5.7:prop-response-zeroth-order} on the response functions
$\mathcal{H}^{\mathsf{A}}$ and $\mathcal{H}^{\mathsf{B}}$ of the two machines, together with
clause~(V3) of the transport lemma.

\emph{Canonicity.} That loops based at one lattice point common to the two lattices, and class
functions, are canonical is the content of the canonicity lemma for the chart datum
$\Delta^{\natural}$ and of clause~(i) of the gauge-blindness lemma of the frame at level $i$.
\end{proof}

\begin{lemma}[Closeness of native loops and plaquettes at real data]\label{5.7:lem-native-loop-closeness}
In this lemma $\mathrm{i}$ denotes the imaginary unit and $i$ a level. Let $(B,B_c)$ be real and
in $\mathfrak{D}(a_{\Lambda})$, with $a_{\Lambda}$ as in \eqref{5.7:eq-data-mismatch-c}. Let
$\varepsilon_D$, $C_L$, $a_A$ and the real base datum $D$ be those of the paired frame of
Definition~\ref{5.7:def-data-mismatch}, and put $\varepsilon:=\varepsilon_D+C_La_{\Lambda}$. Assume
$\varepsilon\le a_A$. Let $\theta_A$, with $3/5<\theta_A\le1$, be the exponent of the weighted norm in the axial
comparison theorem, and put $\theta_1:=(5\theta_A-3)/2$, as in the response statement.
Throughout, $C$ denotes constants depending only on $d$, $L$ and $\theta_A$, possibly different at
each occurrence; in (a), (b) and (d) they are the constants of the statements named there.

Let $\Delta$ be a cell of level $i$ of the frame $\mathfrak{B}^{\star}$, with length scale
$\ell_i=L^{i}\eta$, base point $y_0$ and
enlargement $\Delta^{+}$. Let $\gamma_z$ be the lexicographic combs on $T_{\eta}$ from $y_0$, in the
sense of the comb-product lemma recalled in (c) below, and let
$\lambda_b:=\gamma_{b_-}\,b\,\gamma_{b_+}^{-1}$ be the comb loops. On
$T_{\eta'}$, where
$\eta'=\eta/L$, let $\gamma'_{z'}$ be the lexicographic combs from the same base point $y_0$, and let
$\lambda'_{b'}$ be the corresponding comb loops.

Write $Y^{\mathsf{A}}=Y^{\mathsf{A}}(B)$ and $Y^{\mathsf{B}}=Y^{\mathsf{B}}(B,B_c)$ for the native
canonical axial representatives on $T_{\eta}$ and on $T_{\eta'}$. For a plaquette $P$ based at $x$,
$\partial Y(P)$ denotes the holonomy of $Y$ around $P$. For a configuration $Y$ on a lattice of
spacing $\eta$ and a plaquette $q$ of it, put
\begin{equation*}
\mathfrak{F}_Y(q):=\eta^{-2}\,\tfrac{1}{\mathrm{i}}\log\partial Y(q),
\qquad
\mathfrak{F}^{\rho}_Y(q):=\eta^{-2}\,\tfrac{1}{\mathrm{i}}\log\operatorname{Ad}_{Y(\rho)}\partial Y(q),
\end{equation*}
where $\rho$ is a path; on $T_{\eta'}$ the spacing $\eta'$ replaces $\eta$. Let $M$ be the
block-averaging map, and put $\tilde W:=M(Y^{\mathsf{B}})$. The bonds and plaquettes below all lie
in $\Delta^{+}$.

Assume the following statements.
\begin{itemize}
\item[(a)] The axial comparison theorem, in the following form for the present data. There are a
gauge transformation $\mathsf{u}$ and a Lie-algebra valued field $A$ such that
\begin{equation*}
\tilde W^{\mathsf{u}}=e^{\mathrm{i}\eta A}Y^{\mathsf{A}},\qquad
\ell_i|A|\le C_A\varepsilon L^{-\theta_A i},\qquad |\mathsf{u}-1|\le C\varepsilon L^{-\theta_A i}.
\end{equation*}
Here $C_A$ is the constant of that theorem.
\item[(b)] Part (ii) of its corollary. Let $P'\subset T_{\eta'}$ and $P\subset T_{\eta}$ be parallel
plaquettes based at $x'$ and $x$, with $|x'-x|_{\infty}\le\eta$. Then
\begin{equation*}
\bigl|R\,\mathfrak{F}_{Y^{\mathsf{B}}}(P')-\mathfrak{F}_{Y^{\mathsf{A}}}(P)\bigr|
\le C\varepsilon\ell_i^{-2}L^{-\theta_A i},
\end{equation*}
where $R$ is the adjoint action of $\mathsf{u}(x)$ times the transport from $x$ to $x'$.
\item[(c)] The comb-product lemma. For a bond $b$ of direction $\mu$, $Y^{\mathsf{A}}(\lambda_b)$ is
an ordered product of at most $(d-1)r_b/\eta$ conjugated plaquette variables
$\operatorname{Ad}_{Y^{\mathsf{A}}(\rho_l)}(\partial Y^{\mathsf{A}}(q_l)^{\pm1})$, where $r_b$ is the
length of the comb segments along which these plaquettes run; these segments lie in $\Delta^{+}$
and have length at most $\ell_i$. Here, for each
direction $\kappa$ that follows $\mu$ in the lexicographic order, the plaquettes $q_l$ run along
the segment of the comb in direction
$\kappa$. The same holds on $T_{\eta'}$ with step $\eta'$.
\item[(d)] The regauging identity of the response statement. Write
$Y^{\mathsf{m}}(B)=[e^{\mathrm{i}\eta\mathcal{H}^{\mathsf{m}}(B)}U^{(\prime)}]^{v^{\mathsf{m}}(B)}$
for the representation through the response function $\mathcal{H}^{\mathsf{m}}$ about the base
($U$ for $\mathsf{m}=\mathsf{A}$, $U'$ for $\mathsf{m}=\mathsf{B}$), read in the carried frame. Put
$a:=v^{\mathsf{A}}(B)$ and $a':=v^{\mathsf{B}}(B)$. Let $\mathsf{u}_B$ be the gauge transformation
that the axial comparison theorem produces at the shifted configuration $e^{\mathrm{i}B}D$, and let
$a'{\lceil}$ denote $a'$ read at the points of $T_{\eta}$. Then
\begin{equation*}
w:=\mathsf{u}(a'{\lceil})^{-1}\mathsf{u}_B^{-1}a=e^{\mathrm{i}\lambda},\qquad
|\lambda|\le C_H(\varepsilon_D+\mathsf{b})L^{-\theta_1 i},\qquad \mathsf{b}:=\sup|B|,
\end{equation*}
with the constant $C_H$ of the response statement.
Moreover, $\mathsf{u}$ and $\mathsf{u}_B$ are $C\varepsilon\eta^{\theta_A}$-close to $1$.
\item[(e)] The regularity hypothesis: the configuration $e^{\mathrm{i}B}D$ satisfies the regularity
condition of the paired frame with constant $\varepsilon$.
\end{itemize}
Then there is $C'=C'(d,L,\theta_A)$ with the following property. Let
$b'\parallel b$ with $|b'_--b_-|_{\infty}\le\eta$, and let $P'\parallel P$ be based at $x'$ and $x$
with $|x'-x|_{\infty}\le\eta$, all in $\Delta^{+}$. Then
\begin{equation}\label{5.7:eq-native-loop-closeness-a}
\begin{aligned}
&\text{(k)}\quad
\Bigl|\frac{\ell_i}{\eta'}\frac{1}{\mathrm{i}}\log Y^{\mathsf{B}}(\lambda'_{b'})
-\frac{\ell_i}{\eta}\frac{1}{\mathrm{i}}\log Y^{\mathsf{A}}(\lambda_b)\Bigr|
\le C'\varepsilon L^{-\theta_A i},\\
&\text{(F)}\quad
\ell_i^{2}\Bigl|\eta'^{-2}\log\operatorname{Ad}_{Y^{\mathsf{B}}(\gamma'_{x'})}\partial Y^{\mathsf{B}}(P')
-\eta^{-2}\log\operatorname{Ad}_{Y^{\mathsf{A}}(\gamma_x)}\partial Y^{\mathsf{A}}(P)\Bigr|\\
&\qquad\qquad\le C'\varepsilon L^{-\theta_A i},\\
&\text{(v)}\quad
\bigl|v^{\mathsf{B}}(B)(y_0)\,v^{\mathsf{A}}(B)(y_0)^{-1}-1\bigr|
\le C'\varepsilon L^{-\theta_1 i}.
\end{aligned}
\end{equation}
\end{lemma}

\begin{proof}
Since $(B,B_c)$ is real, all fields below are $\mathrm{SU}(N)$-valued. By (e), the configuration
$e^{\mathrm{i}B}D$
satisfies the regularity condition of the paired frame with constant $\varepsilon$. In this proof
we write $Y:=Y^{\mathsf{A}}$ and $Y':=Y^{\mathsf{B}}$, so that $\tilde W=M(Y')$, and (a) reads
$\tilde W^{\mathsf{u}}=e^{\mathrm{i}\eta A}Y$.

The regularity bound of \cite[(9)]{B-CMP102b} applies on each of the two lattices. It gives
\begin{equation*}
|\mathfrak{F}_Y|\le C\varepsilon\ell_i^{-2},\qquad |\mathfrak{F}_{Y'}|\le C\varepsilon\ell_i^{-2}.
\end{equation*}
It also bounds the H\"older quotients of $\mathfrak{F}_Y$ and $\mathfrak{F}_{Y'}$ by
$C\varepsilon\ell_i^{-2-\beta}$, where $\beta$ is the H\"older exponent of \cite[(9)]{B-CMP102b}.

\emph{First step: a fine straight segment against the average.} Let $b$ be a bond of $T_{\eta}$,
and let $\hat b$ be its subdivision into bonds of $T_{\eta'}$. By \cite[(42)]{B-CMP98}, $M(Y')(b)$
is the mean, in the chart about $Y'(\hat b)$, of the transports of $Y'$ along contours homotopic
to $\hat b$. Each such contour encloses, together with $\hat b$, an area at most $d\eta^2$, and by
the regularity bound the curvature of $Y'$ is at most $C\varepsilon\ell_i^{-2}$. Hence
\begin{equation*}
Y'(\hat b)=e^{\mathrm{i}\eta\mathfrak{z}(b)}\tilde W(b),\qquad
\eta|\mathfrak{z}(b)|\le C\varepsilon\eta^2\ell_i^{-2}.
\end{equation*}
Since $\eta/\ell_i=L^{-i}$, this is the same as $\ell_i|\mathfrak{z}(b)|\le C\varepsilon L^{-i}$.

\emph{Second step: transports.} Let $x$ be a point of $T_{\eta}$ in $\Delta^{+}$. The comb
$\gamma'_x$ is the comb $\gamma_x$ subdivided. The comb $\gamma_x$ has at most $d\ell_i/\eta$ bonds,
and by the first step each of them contributes a discrepancy at most $\eta|\mathfrak{z}|$. Hence
\begin{equation*}
\bigl|Y'(\gamma'_x)\tilde W(\gamma_x)^{-1}-1\bigr|\le\frac{d\ell_i}{\eta}\,\eta|\mathfrak{z}|
\le Cd\varepsilon L^{-i}.
\end{equation*}
By (a),
\begin{equation*}
\tilde W(\gamma_x)=\mathsf{u}(y_0)^{-1}\,(e^{\mathrm{i}\eta A}Y)(\gamma_x)\,\mathsf{u}(x).
\end{equation*}
Again $\gamma_x$ has at most $d\ell_i/\eta$ bonds, and each of them carries the factor
$e^{\mathrm{i}\eta A}$. Hence
\begin{equation*}
\bigl|(e^{\mathrm{i}\eta A}Y)(\gamma_x)Y(\gamma_x)^{-1}-1\bigr|\le d\ell_i\sup|A|
\le dC_A\varepsilon L^{-\theta_A i}.
\end{equation*}
Now let $x'$ be within $\eta$ of $x$. The contours $\gamma'_{x'}$ and $\gamma'_x\cdot(x\to x')$
bound an area at most $d^2\ell_i\eta$. The curvature of $Y'$ is at most $C\varepsilon\ell_i^{-2}$,
so the holonomies along the two contours differ by at most
$d^2\ell_i\eta\cdot C\varepsilon\ell_i^{-2}=Cd^2\varepsilon L^{-i}$.

\emph{Third step: the bound} (F). We apply (b) to the pair $(P',P)$. There $R$ is the adjoint
action of $\mathsf{u}(x)$ times the short transport $x\to x'$.

First, the two logarithms in (F) are curvatures. Since the logarithm near $1$ commutes with
conjugation, $\log\operatorname{Ad}_g h=\operatorname{Ad}_g\log h$, we have
\begin{equation*}
\eta'^{-2}\log\operatorname{Ad}_{Y'(\gamma'_{x'})}\partial Y'(P')
=\mathrm{i}\operatorname{Ad}_{Y'(\gamma'_{x'})}\mathfrak{F}_{Y'}(P'),
\end{equation*}
and likewise on $T_{\eta}$. Both plaquette variables lie near $1$, since
$\eta^2|\mathfrak{F}|\le C\varepsilon\eta^2\ell_i^{-2}$.

Next, we compare the conjugating elements. By the second step, $Y'(\gamma'_{x'})$ agrees with
$Y'(\gamma'_x)$ times the transport from $x$ to $x'$ within $Cd^2\varepsilon L^{-i}$. In turn,
$Y'(\gamma'_x)$ agrees with $\tilde W(\gamma_x)$ within $Cd\varepsilon L^{-i}$. Finally,
$\tilde W(\gamma_x)$ equals $\mathsf{u}(y_0)^{-1}(e^{\mathrm{i}\eta A}Y)(\gamma_x)\mathsf{u}(x)$,
where $(e^{\mathrm{i}\eta A}Y)(\gamma_x)$ lies within $dC_A\varepsilon L^{-\theta_A i}$ of
$Y(\gamma_x)$, and $|\mathsf{u}(y_0)-1|\le C\varepsilon L^{-\theta_A i}$ by (a). Thus
$\operatorname{Ad}_{Y'(\gamma'_{x'})}$ equals $\operatorname{Ad}_{Y(\gamma_x)}R$ up to these
discrepancies.

We now conjugate the inequality of (b) by $\operatorname{Ad}_{Y(\gamma_x)}$, which preserves the
norm. Each replacement of a conjugating element by one within distance $t$ of it changes the
conjugated curvature by at most a multiple of $t|\mathfrak{F}|\le Ct\varepsilon\ell_i^{-2}$. After
multiplication by $\ell_i^2$, the contribution of (b) is $C\varepsilon L^{-\theta_A i}$, and each
discrepancy contributes at most $C\varepsilon\cdot C\varepsilon L^{-\theta_A i}$. For the latter we
use $\varepsilon\le a_A$ and $L^{-i}\le L^{-\theta_A i}$, which holds since $\theta_A\le1$. This
gives (F).

\emph{Fourth step: the bound} (k). By (c), $Y(\lambda_b)$ is an ordered product of at most
$(d-1)r_b/\eta$ factors $\operatorname{Ad}_{Y(\rho_l)}(\partial Y(q_l)^{\pm1})$. For each direction
$\kappa$ that follows $\mu$ in the lexicographic order, the $q_l$ run along the segment of the comb
in direction $\kappa$. Each factor
satisfies
\begin{equation*}
\operatorname{Ad}_{Y(\rho_l)}(\partial Y(q_l)^{\pm1})
=1\pm\mathrm{i}\eta^2\mathfrak{F}^{\rho_l}_Y(q_l)+O(\eta^4|\mathfrak{F}|^2).
\end{equation*}
Hence the product is
\begin{equation*}
Y(\lambda_b)=1+\mathrm{i}\eta^2\sum_l\pm\mathfrak{F}^{\rho_l}_Y(q_l)+O((\eta\ell_i|\mathfrak{F}|)^2).
\end{equation*}
Taking the logarithm and writing the signs into the components $\mathfrak{F}_{\mu\kappa}$, we get
\begin{equation*}
\frac{\ell_i}{\eta}\frac{1}{\mathrm{i}}\log Y(\lambda_b)
=\ell_i\sum_{\kappa}\sum_l\eta\,\mathfrak{F}^{\rho_l}_{Y,\mu\kappa}(q_l)+O(\varepsilon^2L^{-i}).
\end{equation*}
The error term is justified by
\begin{equation*}
\frac{\ell_i}{\eta}(\eta\ell_i|\mathfrak{F}|)^2
\le\frac{\ell_i}{\eta}\bigl(C\varepsilon\eta\ell_i^{-1}\bigr)^2=C^2\varepsilon^2L^{-i}.
\end{equation*}
The main term is a Riemann sum of step $\eta$. By the fine-lattice part of (c), the corresponding
quantity for $Y'(\lambda'_{b'})$ is, with the same error, the Riemann sum of step $\eta'$ of
$\mathfrak{F}^{\rho'}_{Y'}$. Its segments have ends that differ from those of the coarse segments
by at most $\eta$.

To each fine ladder plaquette we assign the coarse one whose running coordinate is the floor, on
$T_{\eta}$, of its own. Then every coordinate differs by at most $\eta$. By the second and third
steps, the integrands therefore agree within $C\varepsilon L^{-\theta_A i}\ell_i^{-2}$. The $L$ fine
plaquettes assigned to one coarse plaquette carry total weight $L\eta'=\eta$. The end pieces, of
length at most $d\eta$, cost at most $d\eta\cdot C\varepsilon\ell_i^{-2}$ each.

It remains to multiply by $\ell_i\cdot(d-1)\ell_i$: the prefactor $\ell_i$, and at most $d-1$
segments of length at most $\ell_i$. The integrand differences then contribute at most
$C(d-1)\varepsilon L^{-\theta_A i}$. The end pieces, multiplied by $\ell_i(d-1)$, contribute at most
$Cd(d-1)\varepsilon L^{-i}$. Together with the two errors $O(\varepsilon^2L^{-i})$, this gives (k).

\emph{Fifth step: the bound} (v). We use (d), with $v^{\mathsf{A}}=a$ and $v^{\mathsf{B}}=a'$. The
element $w=\mathsf{u}(a'{\lceil})^{-1}\mathsf{u}_B^{-1}a=e^{\mathrm{i}\lambda}$ satisfies
$|\lambda|\le C_H(\varepsilon_D+\mathsf{b})L^{-\theta_1 i}$, and $\mathsf{u}$, $\mathsf{u}_B$ are
$C\varepsilon\eta^{\theta_A}$-close to $1$. All four of $a$, $a'$, $\mathsf{u}$, $\mathsf{u}_B$ are
$\mathrm{SU}(N)$-valued. Since $B\in\mathfrak{D}(a_{\Lambda})$, we have
$\mathsf{b}\le a_{\Lambda}\le C_La_{\Lambda}$, because $C_L\ge1$ in the paired frame; hence
$\varepsilon_D+\mathsf{b}\le\varepsilon$ and $|\lambda|\le C_H\varepsilon L^{-\theta_1 i}$. Since
$\eta=\ell_iL^{-i}$ and $\ell_i\le1$, we have $\eta^{\theta_A}\le L^{-\theta_A i}$; moreover
$\theta_1\le\theta_A$, because $\theta_A\le1$, so that $\eta^{\theta_A}\le L^{-\theta_1 i}$.
Evaluating at $y_0$, these bounds on $\lambda$, $\mathsf{u}$ and $\mathsf{u}_B$ give the bound (v).

The collar field enters only through the statements (a)--(e), which hold for every real collar
datum because the hypothesis of the axial comparison theorem admits any collar; hence for
$(B,B_c)$ with $B_c\ne0$ real nothing changes.
\end{proof}

\begin{lemma}[Reduction of the finer averaged datum to one lattice]\label{5.7:lem-datum-reduction}
Let $\mathsf{A}$ act on the lattice of spacing $\eta$ and $\mathsf{B}$ on the lattice of spacing
$\eta'=\eta/L$. Let $M$ be the block-averaging map from the lattice of $\mathsf{B}$ to that of
$\mathsf{A}$, let $M_{\star}$ be the averaging map of $\mathsf{A}$ onto its data bonds $c$, and let
$M'_{\star}$ be that of $\mathsf{B}$. For a layer $A$ over a base $V$ of $\mathsf{A}$ recall, from the
definition of the energy-difference datum, that
\begin{equation*}
  Q_{\star}(A;V)(c)=\frac1i\log\bigl[M_{\star}(e^{i\eta A}V)(c)\,M_{\star}(V)(c)^{-1}\bigr],
\end{equation*}
and that $Q'_{\star}(A';V')(c)$ is defined in the same way with $\eta'$ and $M'_{\star}$ in place of
$\eta$ and $M_{\star}$. Let $A'$ be a layer over a base $\mathbb{W}'$ of $\mathsf{B}$, and let $\mathbb{W}$
be the base of $\mathsf{A}$ paired with $\mathbb{W}'$ in the paired frame of
Definition~\ref{5.7:def-data-mismatch}, with frame potential $A_W$ defined by
$M\mathbb{W}'=e^{i\eta A_W}\mathbb{W}$. Let $A$ be a layer of $\mathsf{A}$ over $\mathbb{W}$, and put
$\delta Q_{\star}(c):=Q'_{\star}(A';\mathbb{W}')(c)-Q_{\star}(A;\mathbb{W})(c)$. Here $l(c)$ is the level
of the data bond $c$, $\ell_{l(c)}$ is the length scale of that level (written $\ell$ when the level
is not displayed), and $\tilde B(c)$ is the neighbourhood of $c$ through which $Q_{\star}(\cdot)(c)$
depends on its argument. Then:
\begin{enumerate}
\item[(i)] for every data bond $c$,
\begin{equation}\label{5.7:eq-datum-reduction-1}
  Q'_{\star}(A';\mathbb{W}')(c)=Q_{\star}\bigl(Y[A';\mathbb{W}'];\,e^{i\eta A_W}\mathbb{W}\bigr)(c);
\end{equation}
\item[(ii)] if $\ell\sup_{\tilde B(c)}|A|$ and $\ell\sup_{\tilde B(c)}|Y[A';\mathbb{W}']|$ are at most
$c_4/11$ and $\ell\sup_{\tilde B(c)}|A_W|$ is less than $c_4$, the radius of analyticity of
\cite[Propositions~6--7]{B-CMP98}, then
\begin{equation}\label{5.7:eq-datum-reduction-a}
\begin{split}
  |\delta Q_{\star}(c)|\le{}& 2.2\,\ell_{l(c)}\sup_{\tilde B(c)}\bigl|Y[A';\mathbb{W}']-A\bigr|\\
  &+\frac{4}{c_4}\Bigl(\ell\sup_{\tilde B(c)}|A_W|\Bigr)\cdot 2\ell_{l(c)}\sup_{\tilde B(c)}|A|;
\end{split}
\end{equation}
\item[(iii)] the scaffold cut-offs $\zeta_{\square}$ of $\mathsf{A}$ and $\zeta'_{\square}$ of $\mathsf{B}$
are one and the same function $\zeta$ on the torus, and
\begin{equation}\label{5.7:eq-datum-reduction-2}
\begin{split}
  \bigl|Y[\zeta A';\mathbb{W}']-\zeta\,Y[A';\mathbb{W}']\bigr|
  &\le 34d\cdot\eta\sup|\nabla\zeta|\cdot\sup|A'|\\
  &\le CL^{-l}L^{-m'}\sup|A'|,
\end{split}
\end{equation}
where $l$ is the level of the block $\square$ and $C$ is the constant of the commutation bound of
the carried-form lemma.
\end{enumerate}
\end{lemma}

\begin{proof}
Recall the three properties of $Q_{\star}$, from the definition of the energy-difference datum,
used below: $Q_{\star}(A)(c)$ depends on $A$ only through its restriction to the set $\tilde B(c)$;
it is analytic; and $|Q_{\star}(A)(c)|\le 2\ell_{l(c)}\sup_{\tilde B(c)}|A|$.

(i) By \cite[(43)]{B-CMP98}, $M'_{\star}=M_{\star}\circ M$. Hence
\begin{equation*}
  Q'_{\star}(A';\mathbb{W}')(c)=\frac1i\log\bigl[M_{\star}\bigl(M(e^{i\eta'A'}\mathbb{W}')\bigr)(c)
  \cdot M_{\star}(M\mathbb{W}')(c)^{-1}\bigr].
\end{equation*}
The carried form is
\begin{equation*}
  Y[h;V]=(i\eta)^{-1}\log\bigl[M(e^{i\eta'h}V)\,M(V)^{-1}\bigr].
\end{equation*}
Exponentiating its definition with $h=A'$ and $V=\mathbb{W}'$ therefore gives
\begin{equation*}
  M(e^{i\eta'A'}\mathbb{W}')=e^{i\eta Y[A';\mathbb{W}']}M\mathbb{W}'.
\end{equation*}
Substituting this into the previous display and comparing with the definition of $Q_{\star}$ at the
layer $Y[A';\mathbb{W}']$ over the base $M\mathbb{W}'$ yields
\begin{equation*}
  Q'_{\star}(A';\mathbb{W}')(c)=Q_{\star}\bigl(Y[A';\mathbb{W}'];M\mathbb{W}'\bigr)(c).
\end{equation*}
Thus the datum of
$\mathsf{B}$ is the functional $Q_{\star}$ of $\mathsf{A}$, read at the averaged layer over the averaged
base. Inserting the frame relation $M\mathbb{W}'=e^{i\eta A_W}\mathbb{W}$ gives
\eqref{5.7:eq-datum-reduction-1}.

(ii) By \cite[Propositions~6--7]{B-CMP98}, the map $(A,a)\mapsto Q_{\star}(A;e^{i\eta a}\mathbb{W})(c)$
is analytic on the ball $\ell|A|,\ \ell|a|<c_4$. It vanishes at $A=0$ by definition, since the
logarithm of the identity is zero. It is bounded by $2\ell_{l(c)}\sup_{\tilde B(c)}|A|$. By
\eqref{5.7:eq-datum-reduction-1},
\begin{equation*}
\begin{split}
  \delta Q_{\star}(c)={}&\Bigl[Q_{\star}\bigl(Y[A';\mathbb{W}'];e^{i\eta A_W}\mathbb{W}\bigr)(c)
  -Q_{\star}\bigl(A;e^{i\eta A_W}\mathbb{W}\bigr)(c)\Bigr]\\
  &+\Bigl[Q_{\star}\bigl(A;e^{i\eta A_W}\mathbb{W}\bigr)(c)-Q_{\star}(A;\mathbb{W})(c)\Bigr].
\end{split}
\end{equation*}
The first bracket is the difference of $Q_{\star}$ at two layers over the same base
$V:=e^{i\eta A_W}\mathbb{W}$, which lies in the domain of analyticity because
$\ell\sup_{\tilde B(c)}|A_W|<c_4$. Put $F(B):=Q_{\star}(B;V)(c)$. By locality $F$ depends on the
layer $B$ only through its restriction to $\tilde B(c)$, so $F$ is analytic on the set of layers with
$\ell\sup_{\tilde B(c)}|B|<c_4$, and there $|F(B)|\le 2\ell_{l(c)}\sup_{\tilde B(c)}|B|<2c_4$.
Let $B$ satisfy $\ell\sup_{\tilde B(c)}|B|\le c_4/11$, and let $h$ be a layer with
$\sup_{\tilde B(c)}|h|=1$. For complex $z$ with
\begin{equation*}
  |z|<\rho:=\frac{c_4-\ell\sup_{\tilde B(c)}|B|}{\ell}
\end{equation*}
one has $\ell\sup_{\tilde B(c)}|B+zh|<c_4$, so $z\mapsto F(B+zh)$ is analytic on the disc of radius
$\rho$ and bounded there by $2c_4$. Cauchy's estimate for the derivative at the centre gives
\begin{equation*}
  \Bigl|\frac{d}{dz}F(B+zh)\Big|_{z=0}\Bigr|\le\frac{2c_4}{\rho}
  =\frac{2c_4\,\ell}{c_4-\ell\sup_{\tilde B(c)}|B|}
  \le\frac{2c_4\,\ell}{10c_4/11}=2.2\,\ell_{l(c)}.
\end{equation*}
If $Y[A';\mathbb{W}']-A$ vanishes on $\tilde B(c)$, the first bracket is zero by locality. Otherwise
put $s:=\sup_{\tilde B(c)}|Y[A';\mathbb{W}']-A|$, $h:=(Y[A';\mathbb{W}']-A)/s$ and
$B_t:=A+t(Y[A';\mathbb{W}']-A)$ for $0\le t\le 1$. By the hypothesis of (ii) and the convexity of
$B\mapsto\sup_{\tilde B(c)}|B|$, every $B_t$ satisfies $\ell\sup_{\tilde B(c)}|B_t|\le c_4/11$. Since
$\frac{d}{dt}F(B_t)=s\,\frac{d}{dz}F(B_t+zh)|_{z=0}$, integrating over $t\in[0,1]$ shows that the
first bracket is at most $2.2\,\ell_{l(c)}\sup_{\tilde B(c)}|Y[A';\mathbb{W}']-A|$.

For the second bracket, fix $A$ and consider
\begin{equation*}
  g(a):=Q_{\star}(A;e^{i\eta a}\mathbb{W})(c)-Q_{\star}(A;\mathbb{W})(c).
\end{equation*}
This function is analytic in
$a$ on $\ell|a|<c_4$ and vanishes at $a=0$. By the triangle inequality and the bound on $Q_{\star}$,
it is bounded there by $2\cdot 2\ell_{l(c)}\sup_{\tilde B(c)}|A|$. By locality, $g$ depends on $a$
only through its restriction to $\tilde B(c)$; the Schwarz lemma on the ball
$\ell\sup_{\tilde B(c)}|a|<c_4$, evaluated at $a=A_W$, gives
\begin{equation*}
  |g(A_W)|\le\frac{4}{c_4}\Bigl(\ell\sup_{\tilde B(c)}|A_W|\Bigr)\cdot
  \ell_{l(c)}\sup_{\tilde B(c)}|A|,
\end{equation*}
which is at most the second term on the right of \eqref{5.7:eq-datum-reduction-a}.
Adding the two bounds gives \eqref{5.7:eq-datum-reduction-a}.

(iii) Since the cut-offs $\zeta_{\square}$ and $\zeta'_{\square}$ are one function $\zeta$ on the torus,
the same $\zeta$ multiplies the layer $A'$ of $\mathsf{B}$ and the carried layer
$Y[A';\mathbb{W}']$ read by $\mathsf{A}$. Both inequalities in \eqref{5.7:eq-datum-reduction-2} are
the commutation bound of the carried-form lemma, applied to the layer $A'$ over the base
$\mathbb{W}'$ and the function $\zeta$.
\end{proof}

\begin{corollary}[Linear control of the children's chart data]\label{5.7:cor-chart-data-control}
Assume the following statements:
\begin{enumerate}
\item the hypotheses of the two-cutoff response statement at zeroth order,
Proposition~\ref{5.7:prop-response-zeroth-order}, at both cutoffs $\mathsf{A}$ and $\mathsf{B}$;
\item the second part of the transport lemma for comb loops, which bounds at the actual data,
at each cutoff, the transports along the comb loops $\lambda_b$;
\item the canonical-data lemma, including its part (iii);
\item the distance bound $1+|b-y|\le 1+2D_v$ for the bonds $b$ of the unit region $\mathfrak{r}_v$;
\item the regauging bound at a base point;
\item the averaging theorem for the canonical representative, in the notation below: the
composition identity $M'_{\star}=M_{\star}\circ M$ for the averaging maps of the two lattices, the
conjugation identity
\begin{equation*}
(M_{\star}Y)(\lambda_b)=v(y)\,\mathsf{M}(\lambda_b)\,v(y)^{-1},
\end{equation*}
and the bound $\ell_j|A^{\natural}|\le C\vartheta^{4j}$ on $\mathfrak{r}_v$ for the carried-frame
potential $A^{\natural}$.
\end{enumerate}
At each cutoff work in the Landau frame about the real base. For $\iota\in\{0,1\}$ let $B_\iota$
be the two data of the canonical-data lemma, $\mathcal{H}(B_\iota)$ the response about the base, and
put $W^{L}_{\iota}:=e^{i\eta\mathcal{H}(B_\iota)}U$ in $\mathsf{A}$ and
$W^{L}_{\iota}:=e^{i\eta'\mathcal{H}(B_\iota)}U'$ in $\mathsf{B}$, each cutoff with its own
lattice spacing; put $\mathsf{M}_{\iota}:=M_{\star}W^{L}_{\iota}$ on the bonds of
$\mathfrak{B}_j(\mathfrak{r}_v)$, the averaging $M_{\star}$ being taken to the level $l(\kappa)$ of
each bond $\kappa$ (in $\mathsf{B}$ the averaging map is $M'_{\star}=M_{\star}\circ M$), and
\begin{equation*}
G_{\iota}(b):=\tfrac{1}{i}\log\mathsf{M}_{\iota}(\lambda_b),\qquad
v_{\iota}:=v(B_\iota)(y),\qquad \tilde v':=v_0^{-1}v_1 ,
\end{equation*}
where $y$ is the common base point of the comb loops $\lambda_b$ and $v(B)(y)$ is the regauging
element at $y$. Then the canonical data of the canonical-data lemma are
\begin{equation}\label{5.7:eq-chart-data-control-1}
B^{0}=\operatorname{Ad}_{v_0}G_0,\qquad \Delta=\operatorname{Ad}_{v_0}\Delta^{\natural},\qquad
\Delta^{\natural}:=\operatorname{Ad}_{\tilde v'}G_1-G_0 ,
\end{equation}
and the common factor $\operatorname{Ad}_{v_0}$ is a constant gauge of the chart
(\cite[(181)]{B-CMP102b}), to which the children of the chart move of the canonical-data lemma are
covariant and to which every function read is insensitive. Let
$\varpi\ge\mathsf{m}^{\mathrm{dat}}(\mathfrak{b}^{\mathsf{B}},\mathfrak{b}^{\mathsf{A}})$ be the data
mismatch and $\varpi'$ the mismatch as in Proposition~\ref{5.7:prop-response-zeroth-order}, and put
$\hat\varpi:=2\varpi+1.1\varpi'+c_1\vartheta^{j}$. There are constants $C_{\mathrm{ch}}$ and
$C'_{\mathrm{ch}}$, independent of $b$ and of $D_v$, such that for every bond $b$
\begin{equation}\label{5.7:eq-chart-data-control-2}
|\delta\Delta^{\natural}(b)|\le C_{\mathrm{ch}}(1+|b-y|)^{2}\,\hat\varpi\,
\sup_{\mathfrak{r}_v}8L^{3}B_H T_{\ast},
\end{equation}
hence
\begin{equation}\label{5.7:eq-chart-data-control-3}
e^{\frac18\delta D_v}\sup_{b}|\delta\Delta^{\natural}(b)|\le C'_{\mathrm{ch}}\,\hat\varpi\,
\alpha_{3,j},
\end{equation}
and
\begin{equation}\label{5.7:eq-chart-data-control-4}
|\delta G_0(b)|\le C_{\mathrm{ch}}(1+|b-y|)\bigl(\varpi+\vartheta^{4j}\bigr).
\end{equation}
Consequently, on the whole domain
$|\sigma_0|\,|B^{0}|_{\star}+|\sigma_1|\,|\Delta|_{\star}<\alpha_{3,j}$ of the chart move, where
$\sigma_0,\sigma_1$ are the complex chart parameters of the chart move and $|\cdot|_{\star}$ is its
weighted supremum norm (canonical-data lemma), the chart
data of the children at the two cutoffs differ, modulo one constant gauge, by at most
\begin{equation*}
|\sigma_0|\,|\delta G_0|+2C'_{\mathrm{ch}}\,\hat\varpi\,\alpha_{3,j}:
\end{equation*}
linearly in the parent's data mismatch $\varpi$ and in the mismatch $\varpi'$ of the
energy-difference bound, uniformly in $D_v$. The same holds, with the same proof, for the
analogous control of the chart data of the unit region $\mathfrak{r}_u$ at level $m$, with
$(j,\mathfrak{r}_v,D_v)$ replaced by $(m,\mathfrak{r}_u,D_u)$ and the factor $e^{\frac18\delta D_v}$
replaced by the dilation radius $R_u$, which is proportional to $e^{\frac14\delta D_u}$; there the
polynomial factor in $|b-y|$ is bounded by a number depending only on $d$ and $L$.
\end{corollary}

\begin{proof}
We first record why the canonical data have the form \eqref{5.7:eq-chart-data-control-1}. By the
conjugation identity of the averaging theorem, $(M_{\star}Y)(\lambda_b)=v(y)\,\mathsf{M}(\lambda_b)\,v(y)^{-1}$,
where $Y$ is
the native canonical axial representative, the averaged transport of the canonical representative
along $\lambda_b$ is the averaged Landau transport conjugated by the regauging element at $y$.
Taking logarithms at $\iota=0$ gives $B^{0}=\operatorname{Ad}_{v_0}G_0$; at $\iota=1$ it gives
$\operatorname{Ad}_{v_1}G_1=\operatorname{Ad}_{v_0}\operatorname{Ad}_{\tilde v'}G_1$, and subtracting
the first gives $\Delta=\operatorname{Ad}_{v_0}(\operatorname{Ad}_{\tilde v'}G_1-G_0)$.

The interpolation theorem for two analytic maps is applied below to local functions only: the
fields at a point, and the regauging element at a base point. The loops are composed afterwards
at the actual data, where the second part of the transport lemma for comb loops bounds the
transports at each cutoff. On the uniform polydisc $\mathfrak{D}(a_{\Lambda})$ the chart loops
themselves are not admissible functions $\Phi$ for the interpolation theorem: their length at
level $j$ reaches $6dM/t_p$, where $M=L^{m'}$ is the auxiliary parameter of that name and not the
averaging map, and a complex potential of size $a_{\Lambda}/\ell_j$ along them
gives transports of norm $e^{Ca_{\Lambda}M/t_p}$, so that part (\(\Lambda\)1) of
\eqref{5.7:eq-two-map-interpolation-a} fails. The gain $t_p$ of the first part of the transport
lemma belongs to the actual data only.

Throughout, $\ell_{l(\kappa)}$ is the length scale of the level $l(\kappa)$ of a bond $\kappa$,
$w_\kappa:=\ell_{l(\kappa)}/\ell_j$, and $n(\cdot)$ is the matrix norm of the transport lemma.

\emph{Step 1: the base.} Put $\mathsf{Y}_0:=\mathsf{Y}(B_0^{\mathsf{B}})$. Averaging the base
of the cutoff $\mathsf{B}$ once gives
\begin{equation*}
M\bigl(W_0^{L,\mathsf{B}}\bigr)=e^{i\eta\mathsf{Y}_0}e^{i\eta A^{\natural}}U
=:e^{i\eta a_W}W_0^{L,\mathsf{A}} ,
\end{equation*}
which defines $a_W$. On $\mathfrak{r}_v$, with $\ell=\ell_j$ here and below,
\begin{equation}\label{5.7:eq-chart-data-control-5}
\ell|a_W|\le 1.1\Bigl(\ell|A^{\natural}|
+\ell\bigl|\mathsf{Y}(B_0^{\mathsf{B}})-\mathcal{H}^{\mathsf{A}}(B_0^{\mathsf{A}})\bigr|\Bigr)
\le C_{\mathrm{base}}\bigl(\vartheta^{4j}+\varpi\bigr),
\end{equation}
by the bound on the carried-frame potential $A^{\natural}$ in the averaging theorem for the first
term, and for the second
term by part (f) of the interpolation theorem for two analytic maps, applied with the function
$\Phi_x$ of
\eqref{5.7:eq-response-zeroth-order-b}. By the composition identity,
$\mathsf{M}_0^{\mathsf{B}}=M_{\star}(MW_0^{L,\mathsf{B}})=M_{\star}(e^{i\eta a_W}W_0^{L,\mathsf{A}})$,
hence for every bond $\kappa$
\begin{equation*}
\mathsf{M}_0^{\mathsf{B}}(\kappa)=e^{iq(\kappa)}\mathsf{M}_0^{\mathsf{A}}(\kappa),\qquad
|q(\kappa)|\le 2\ell_{l(\kappa)}\sup|a_W|\le 2w_\kappa C_{\mathrm{base}}(\vartheta^{4j}+\varpi),
\end{equation*}
where $q(\kappa)$ is the averaged datum of Lemma~\ref{5.7:lem-datum-reduction} read at the layer
$a_W$ over the base $W_0^{L,\mathsf{A}}$, bounded there by $2\ell_{l(\kappa)}\sup|a_W|$, and the
last inequality is \eqref{5.7:eq-chart-data-control-5}.

Let $\Pi_k^{\mathsf{m}}$ be the partial products of $\mathsf{M}_0^{\mathsf{m}}$ along $\lambda_b$.
By the second part of the transport lemma, together with $n(v_0)\le 2^{1/4}$, one has
$n(\Pi_k^{\mathsf{m}})\le 1.45\cdot 2^{1/2}$ at both cutoffs. The length of $\lambda_b$ at level
$j$ is $\operatorname{len}_j(\lambda_b):=\sum_{\kappa\in\lambda_b}w_\kappa$, and
$\operatorname{len}_j(\lambda_b)\le 2d(1+|b-y|)$ by the comb-loop length bound stated with the
telescoping identity of the canonical-data lemma. Put
$E:=4.3\sum_{\kappa\in\lambda_b}|q(\kappa)|$; then
\begin{equation*}
E\le 4.3\cdot 2C_{\mathrm{base}}(\vartheta^{4j}+\varpi)\cdot 2d(1+|b-y|)
\le 18dC_{\mathrm{base}}(1+|b-y|)(\vartheta^{4j}+\varpi),
\end{equation*}
and
\begin{equation}\label{5.7:eq-chart-data-control-6}
\bigl|\Pi_k^{\mathsf{B}}(\Pi_k^{\mathsf{A}})^{-1}-1\bigr|\le 6E .
\end{equation}
Indeed, writing $\lambda_b=(\kappa_1,\kappa_2,\dots)$, the relation
$\mathsf{M}_0^{\mathsf{B}}(\kappa_l)=e^{iq(\kappa_l)}\mathsf{M}_0^{\mathsf{A}}(\kappa_l)$ gives
$\Pi_k^{\mathsf{B}}(\Pi_k^{\mathsf{A}})^{-1}$ as the ordered product over $l\le k$ of the factors
$\exp\bigl(i\operatorname{Ad}_{\Pi^{\mathsf{A}}_{l-1}}q(\kappa_l)\bigr)$; since
$(1.45\cdot2^{1/2})^{2}\le 4.3$, each factor differs from $1$ by at most $e^{4.3|q(\kappa_l)|}-1$,
and therefore
\begin{equation*}
\bigl|\Pi_k^{\mathsf{B}}(\Pi_k^{\mathsf{A}})^{-1}-1\bigr|
\le\prod_{l\le k}e^{4.3|q(\kappa_l)|}-1\le e^{E}-1 .
\end{equation*}
If $E\le1$, then $e^{E}-1\le 2E\le 6E$; if $E>1$ the left side is at
most $1+(1.45\cdot2^{1/2})^{2}\le 6\le 6E$. Applied to the complete product along $\lambda_b$, which
is $\mathsf{M}_0^{\mathsf{m}}(\lambda_b)$, \eqref{5.7:eq-chart-data-control-6} gives
\eqref{5.7:eq-chart-data-control-4}.

\emph{Step 2: the increments bond by bond.} Let $\mathcal{X}^{\mathsf{m}}$ be the increment in
the Landau frame and define
\begin{equation*}
\mathfrak{h}^{L}(\kappa):=\tfrac{1}{i}\log\bigl[\mathsf{M}_1(\kappa)\mathsf{M}_0(\kappa)^{-1}\bigr]
=Q_{\star}\bigl(\mathcal{X};W_0^{L}\bigr)(\kappa).
\end{equation*}
By the exact reduction of Lemma~\ref{5.7:lem-datum-reduction} and Step 1,
\begin{equation*}
\mathfrak{h}^{L,\mathsf{B}}(\kappa)=Q_{\star}\Bigl(Y\bigl[\mathcal{X}^{\mathsf{B}};W_0^{L,\mathsf{B}}\bigr];
e^{i\eta a_W}W_0^{L,\mathsf{A}}\Bigr)(\kappa).
\end{equation*}
By the analyticity of $Q_{\star}(A;e^{i\eta a}\mathbb{W})$ in $(A,a)$ on the ball
$\ell|A|,\ell|a|<c_4$ and the estimate \eqref{5.7:eq-datum-reduction-a},
\begin{align*}
|\delta\mathfrak{h}^{L}(\kappa)|
&\le 2.2\,\ell_{l(\kappa)}\sup\bigl|Y[\mathcal{X}^{\mathsf{B}}]-\mathcal{X}^{\mathsf{A}}\bigr|
+\frac{8}{c_4}\bigl(\ell\sup|a_W|\bigr)\,\ell_{l(\kappa)}\sup|\mathcal{X}^{\mathsf{A}}|\\
&\le w_\kappa K_1'\,\hat\varpi\,T_{\ast}(\kappa).
\end{align*}
The second inequality uses the bound (S$'$) of \eqref{5.7:eq-response-zeroth-order-a} for the first
term, which applies with $l_x=j$ on $\mathfrak{r}_v$, the bonds of its rim included; and, for the
second term, \eqref{5.7:eq-chart-data-control-5} together with
$\ell_j|\mathcal{X}^{\mathsf{A}}|\le 1.2K_R T_{\ast}$. Here $K_1'$ denotes a constant independent of
$\kappa$, $b$ and $D_v$, determined by the constant of (S$'$), by $C_{\mathrm{base}}$, $K_R$ and
$c_4$. At a single cutoff the same argument gives
\begin{equation*}
|\mathfrak{h}^{L}(\kappa)|\le 2.2\,w_\kappa K_R\,T_{\ast}(\kappa).
\end{equation*}

\emph{Step 3: the loops.} Moving, bond by bond along $\lambda_b=(\kappa_1,\kappa_2,\dots)$, each
factor $e^{\pm i\mathfrak{h}^{L}(\kappa_k)}$ to the left past the partial product $\Pi_{k-1}$ of
$\mathsf{M}_0$ gives the telescoping identity
\begin{equation*}
\mathsf{M}_1(\lambda_b)=\Bigl[\prod_k e^{\pm i\operatorname{Ad}_{\Pi_{k-1}}\mathfrak{h}^{L}(\kappa_k)}\Bigr]
\mathsf{M}_0(\lambda_b)=:e^{i\Upsilon_b}\mathsf{M}_0(\lambda_b).
\end{equation*}
For each factor,
\begin{equation*}
\delta\bigl(\operatorname{Ad}_{\Pi}\mathfrak{h}^{L}\bigr)
=\operatorname{Ad}_{\Pi^{\mathsf{B}}}\delta\mathfrak{h}^{L}
+\bigl(\operatorname{Ad}_{\Pi^{\mathsf{B}}}-\operatorname{Ad}_{\Pi^{\mathsf{A}}}\bigr)
\mathfrak{h}^{L,\mathsf{A}} .
\end{equation*}
The first term is bounded by Step 2 and the bound on $n(\Pi_k^{\mathsf{m}})$; the second by
\eqref{5.7:eq-chart-data-control-6} and the single-cutoff bound of Step 2. Summing over the
bonds of $\lambda_b$, whose number of bonds weighted by $w_\kappa$ is at most $2d(1+|b-y|)$, we
obtain, with constants $C$ independent of $b$ and $D_v$,
\begin{equation*}
|\delta\Upsilon_b|\le C(1+|b-y|)K_1'\hat\varpi\sup T_{\ast}
+C(1+|b-y|)^{2}(\vartheta^{4j}+\varpi)K_R\sup T_{\ast}.
\end{equation*}
Next, $G_1-G_0=\frac1i\log\bigl(e^{i\Upsilon}e^{iG_0}\bigr)-G_0$ vanishes at $\Upsilon=0$, and
therefore
\begin{equation*}
|\delta(G_1-G_0)|\le 1.1\,|\delta\Upsilon_b|+2\,|\Upsilon_b|\,|\delta G_0| .
\end{equation*}

\emph{Step 4: the regauging at $y$.} Apply the interpolation theorem for two analytic maps with
$\Phi^{\mathsf{m}}(B):=v^{\mathsf{m}}(B)(y)$. Condition \eqref{5.7:eq-two-map-interpolation-a}, in
its part (\(\Lambda\)1), holds by the regauging bound
\begin{equation*}
\bigl|v(\tilde B)(y)v(B)(y)^{-1}-1\bigr|\le \frac{20}{3r_V}\,\ell_j
\sup_{\Delta_y}\bigl|\mathcal{H}(\tilde B)-\mathcal{H}(B)\bigr|,
\end{equation*}
where $r_V$ is the constant of the regauging bound and $\Delta_y$ is the cell of
$\mathfrak{B}^{\star}$ with base point $y$ (not to be confused with the datum $\Delta$),
with kernel $\frac{20}{3r_V}\cdot 8L^{3}B_He^{-\delta d_{\star}(y,c)}$; its part (\(\Lambda\)2) is
part (v) of \eqref{5.7:eq-native-loop-closeness-a}, the closeness of native loops and plaquettes
at real data (Lemma~\ref{5.7:lem-native-loop-closeness}). Hence, with a constant $K_v$ given by the
interpolation theorem for this kernel, and $C$ as above,
\begin{equation*}
|\delta(v_1-v_0)|\le K_v\hat\varpi\,T_{\ast}(y),\qquad
|\delta v_0|\le C(\vartheta^{4j}+\varpi),\qquad |\tilde v'-1|\le K_vT_{\ast}(y).
\end{equation*}
Write $\Delta^{\natural}=(\operatorname{Ad}_{\tilde v'}-1)G_1+(G_1-G_0)$. For the first term,
\begin{equation*}
\delta\bigl[(\operatorname{Ad}_{\tilde v'}-1)G_1\bigr]
=(\delta\operatorname{Ad}_{\tilde v'})G_1+(\operatorname{Ad}_{\tilde v'}-1)\delta G_1,
\qquad |G_1|\le 1 .
\end{equation*}
The second term was bounded in Step 3. Adding the bounds of Steps 3 and 4, and using
$\delta G_1=\delta(G_1-G_0)+\delta G_0$, \eqref{5.7:eq-chart-data-control-4}, $\varpi\le\hat\varpi/2$
and $\vartheta^{4j}\le\vartheta^{j}$ (as $\vartheta=L^{-\theta_1/8}<1$), gives
\eqref{5.7:eq-chart-data-control-2}.

\emph{Step 5: the second display.} By Lemma~\ref{5.7:lem-three-quarter-rate}, part (ii), in the
form \eqref{5.7:eq-three-quarter-rate-a},
\begin{equation*}
8L^{3}B_HT_{\ast}\le K_1\beta e^{-\frac14\delta D_v}
\bigl[5.5\vartheta_p+4.4|\mathsf{t}|\vartheta_p(\square)\bigr]\alpha_{\cdot,j}
\quad\text{on }\mathfrak{r}_v ,
\end{equation*}
where $\vartheta_p$ and $\vartheta_p(\square)$ are the indices of the profiles in
Lemma~\ref{5.7:lem-three-quarter-rate}, distinct from the index base $\vartheta=L^{-\theta_1/8}$.
By the distance bound for the bonds of $\mathfrak{r}_v$, $1+|b-y|\le 1+2D_v$. Inserting both into
\eqref{5.7:eq-chart-data-control-2} and multiplying by $e^{\frac18\delta D_v}$ leaves the factor
$(1+2D_v)^{2}e^{-\frac18\delta D_v}$, which is bounded by
$\sup_{x\ge0}(1+2x)^{2}e^{-\frac18\delta x}<\infty$, a number depending only on $\delta$. Finally, the
last line of the majorant budget of the canonical-data lemma and the arithmetic of its part (iii)
turn $K_1\beta[5.5\vartheta_p+4.4|\mathsf{t}|\vartheta_p(\square)]\alpha_{\cdot,j}$ into a multiple of
$\alpha_{3,j}$. This gives \eqref{5.7:eq-chart-data-control-3}, uniformly in $D_v$.

In \eqref{5.7:eq-chart-data-control-4} the coefficient of the parent's data mismatch $\varpi$ is
$C_{\mathrm{ch}}(1+|b-y|)$, proportional to the length of the loop $\lambda_b$.
\end{proof}

\begin{proposition}[Two-cutoff difference of the dilated action family]
\label{5.7:prop-dilated-action-family}
Let $d=4$, and work in the setting of the dilated Wilson family at level $i$. Here $v=(c,\Delta)$ is a
unit, formed by a function $c$ and a cell $\Delta$ of level $i$ with base point $y_0$, lexicographic
combs $\gamma_z$ from $y_0$ and comb loops $\lambda_b$. The family $\mathfrak{F}_v(s):=A(c,U_s)$ is the
dilated Wilson family of $v$; here $A(c,U)$ is the sum over plaquettes $P$ of $c(P)$ times the
one-plaquette action $a$ of the plaquette variable of $U$ at $P$. The family is defined for
$|s|\le R^{\perp}_v$, where $R^{\perp}_v:=(32\kappa^{\ast}_v)^{-1}$
is proportional to $e^{\delta D_v/4}$, the constant $\kappa^{\ast}_v$ and the distance $D_v$ of the
unit region being fixed in that setting. The cutoff $\mathsf{A}$ lives on $T_{\eta}$ and the cutoff
$\mathsf{B}$ on $T_{\eta'}$, with $\eta'=\eta/L$. We write $\ell_i$ for the length scale of level $i$
and put $u_i:=(\eta/\ell_i)^2$; in the dilated setting $\eta/\ell_i$ is at most a constant multiple
of $L^{-i}$. For a quantity $\Phi$ defined at both cutoffs,
$\delta\Phi:=\Phi^{\mathsf{B}}-\Phi^{\mathsf{A}}$. Assume the following.
\begin{itemize}
\item The hypotheses of the two-cutoff response statement at zeroth order,
Proposition~\ref{5.7:prop-response-zeroth-order}. In what follows $\varpi$ is the bound on the data
mismatch assumed there, $\varpi'$ is the far-data mismatch of the energy-difference bound, and
$c_1$, $\vartheta$ and the profile sum $T_{\ast}$ are as in that proposition.
\item Clause (iv) of the comb lemma of the dilated setting, together with four standing bounds of
that setting: a bound (a) used for comb loops only, a bound (b) used both for comb loops and for
plaquettes, and two bounds (c) and (d) used for plaquettes only.
\item The plaquette-logarithm bound of the dilated setting.
\item The pairing corollary for plaquettes of the two lattices, which supplies the pairing
$P'\sim P$ used below.
\end{itemize}
Then for every $s$ with $|s|\le R^{\perp}_v$,
\begin{equation}\label{5.7:eq-dilated-action-family-1}
|\delta\mathfrak{F}_v(s)|\le C_{\mathfrak{F}}\,\mathfrak{c}_{\mathrm{tr}}
\bigl(\mathfrak{n}_v+\mathsf{p}_d\mathsf{s}_v\bigr)
\bigl(2\varpi+1.1\varpi'+c_1'''\vartheta^{i}\bigr).
\end{equation}
Here $\mathsf{s}_v:=\sup|c|+\ell_i\sup|\nabla c|$, $c_1'''$ is a constant which contains $c_1$
together with the contributions carrying a factor $L^{-i}$ that are collected in the proof, and
$C_{\mathfrak{F}}=C_{\mathfrak{F}}(d,L)$. The constants $\mathfrak{c}_{\mathrm{tr}}$, $\mathfrak{n}_v$
and $\mathsf{p}_d$ are those of the dilated setting; $\mathfrak{c}_{\mathrm{tr}}$ enters through the
trace in the one-plaquette action, and $\mathsf{p}_d$ is the sum of $u_i^2$ over the plaquettes of
the cell (both appear in the last step of the proof). Thus, at the points dilated by
$e^{\delta D_v/4}$, the
two-cutoff difference of the Wilson family is bounded by three quantities, uniformly in $D_v$:
\begin{itemize}
\item the data mismatch $\varpi$ of the response statement;
\item the far-data mismatch $\varpi'$ of the energy-difference bound;
\item the term $c_1'''\vartheta^{i}$ in the level $i$.
\end{itemize}
\end{proposition}

\begin{proof}
Throughout, $C$ denotes a constant whose value may change from one occurrence to the next.

We first pass to the comb gauge. Fix one of the two cutoffs and suppress its label until the two
are compared. Let $W$ and $W^1$ be the native configurations at the two data $B_0$ and $B_1$. The
comb loop of a bond $b$ is $\lambda_b=\gamma_{b_-}b\gamma_{b_+}^{-1}$. Hence
$W(\lambda_b)=W(\gamma_{b_-})W(b)W(\gamma_{b_+})^{-1}$, or equivalently
\begin{equation*}
W(\gamma_{b_+})W(b)^{-1}=W(\lambda_b)^{-1}W(\gamma_{b_-}).
\end{equation*}
The regauging field through which the dilated family $U_s$ is defined is
\begin{equation*}
\mathfrak{E}(b)=\operatorname{Ad}_{W(\gamma_{b_-})^{-1}}\bigl(W^1(\lambda_b)W(\lambda_b)^{-1}\bigr).
\end{equation*}
Perform the gauge transformation $g(z):=W(\gamma_z)$. By the definition of $U_s$ through
$\mathfrak{E}$ and the preceding identity, in this gauge
\begin{equation*}
(U_s)^{g}(b)=\exp\bigl(s\log[W^1(\lambda_b)W(\lambda_b)^{-1}]\bigr)\,W(\lambda_b).
\end{equation*}
The family $\mathfrak{F}_v$ is a class function of the plaquette variables, and a plaquette variable
changes by conjugation under a gauge transformation. So $\mathfrak{F}_v$ is gauge invariant and may
be evaluated in this gauge. Define $\mathfrak{x}_b$ and $\mathfrak{y}_b$ by
\begin{equation*}
e^{\mathfrak{x}_b}:=W(\lambda_b),\qquad e^{\mathfrak{y}_b}e^{\mathfrak{x}_b}:=W^1(\lambda_b).
\end{equation*}
The loops $W(\lambda_b)$ and $W^1(\lambda_b)$ are native loops based at $y_0$. In $1/i$, $i\eta$ and
$e^{iB}$ the letter $i$ denotes the
imaginary unit; as a subscript, or as the exponent of $\vartheta$ or of $L$, it denotes the level.
In the notation of the dilated setting, $\mathfrak{x}_b=i\eta a_0(b)$. For real or complex data $B$, put
\begin{align*}
\Phi_b(B)&:=\frac{\ell_i}{\eta}\,\frac{1}{i}\,\log Y(B)(\lambda_b),\\
\Phi_P(B)&:=\ell_i^{2}\eta^{-2}\log\operatorname{Ad}_{Y(B)(\gamma_x)}\partial Y(B)(P),
\end{align*}
for a plaquette $P$ based at $x$. Then $W(\lambda_b)$ and $W^1(\lambda_b)$ are the native loops whose
rescaled logarithms are the values $\Phi_b(B_0)$ and $\Phi_b(B_1)$. For a single cutoff, clause
(iv) of the comb lemma and the standing bounds (a) and (b) give
\begin{equation*}
|\mathfrak{x}_b|\le Cu_i^{1/2},\qquad |s\mathfrak{y}_b|\le u_i^{1/2}/14\quad\text{for }|s|\le R^{\perp}_v.
\end{equation*}

Next we prove an exact interpolation identity. Let $P=b_1b_2b_3^{-1}b_4^{-1}$, and let $G(s)$ be the
plaquette product of the factors $e^{s\mathfrak{y}_b}e^{\mathfrak{x}_b}$:
\begin{equation*}
G(s):=e^{s\mathfrak{y}_{b_1}}e^{\mathfrak{x}_{b_1}}\,e^{s\mathfrak{y}_{b_2}}e^{\mathfrak{x}_{b_2}}\,
\bigl(e^{s\mathfrak{y}_{b_3}}e^{\mathfrak{x}_{b_3}}\bigr)^{-1}
\bigl(e^{s\mathfrak{y}_{b_4}}e^{\mathfrak{x}_{b_4}}\bigr)^{-1}.
\end{equation*}
Write $\sum_{\pm}$ for the sum over $j=1,\dots,4$ with sign $+$ for $j=1,2$ and $-$ for $j=3,4$.
For bond-indexed families $Z=(Z_{b_j})$, $Z'=(Z'_{b_j})$ and $Z''=(Z''_{b_j})$, put
\begin{align*}
\mathsf{L}_Z&:=\sum_{\pm}Z_{b_j}, &
\mathsf{D}&:=\tfrac12\sum_{\pm}[\mathfrak{y}_{b_j},\mathfrak{x}_{b_j}],\\
\mathsf{C}(Z',Z'')&:=\tfrac12\sum_{j<k}[\pm Z'_{b_j},\pm Z''_{b_k}].
\end{align*}
The letter $\mathsf{L}_Z$ is unrelated to the blocking factor $L$, and $\mathsf{D}$ here is the
commutator sum just defined, not a response difference. Let $R_3(s)$ denote the terms of
order at least three. The Baker--Campbell--Hausdorff formula to second order gives
\begin{equation*}
\log G(s)=s\mathsf{L}_{\mathfrak{y}}+\mathsf{L}_{\mathfrak{x}}+s\mathsf{D}
+\mathsf{C}(s\mathfrak{y}+\mathfrak{x},s\mathfrak{y}+\mathfrak{x})+R_3(s).
\end{equation*}
By bilinearity of $\mathsf{C}$,
\begin{equation*}
\mathsf{C}(s\mathfrak{y}+\mathfrak{x},s\mathfrak{y}+\mathfrak{x})
=s^2\mathsf{C}(\mathfrak{y},\mathfrak{y})
+s[\mathsf{C}(\mathfrak{y},\mathfrak{x})+\mathsf{C}(\mathfrak{x},\mathfrak{y})]
+\mathsf{C}(\mathfrak{x},\mathfrak{x}).
\end{equation*}
At $s=0$ this yields $\log G(0)=\mathsf{L}_{\mathfrak{x}}+\mathsf{C}(\mathfrak{x},\mathfrak{x})+R_3(0)$.
At $s=1$ it yields
\begin{equation*}
\begin{split}
\log G(1)-\log G(0)&=\mathsf{L}_{\mathfrak{y}}+\mathsf{D}+\mathsf{C}(\mathfrak{y},\mathfrak{x})
+\mathsf{C}(\mathfrak{x},\mathfrak{y})+\mathsf{C}(\mathfrak{y},\mathfrak{y})\\
&\quad+R_3(1)-R_3(0).
\end{split}
\end{equation*}
Eliminate $\mathsf{L}_{\mathfrak{x}}+\mathsf{C}(\mathfrak{x},\mathfrak{x})$ and
$\mathsf{L}_{\mathfrak{y}}+\mathsf{D}+\mathsf{C}(\mathfrak{y},\mathfrak{x})+\mathsf{C}(\mathfrak{x},\mathfrak{y})$
between these two relations and the expansion. This gives, for every $s$,
\begin{equation}\label{5.7:eq-dilated-action-family-a}
\begin{split}
\log G(s)&=\log G(0)+s\bigl[\log G(1)-\log G(0)\bigr]
+(s^2-s)\,\mathsf{C}(\mathfrak{y},\mathfrak{y})+\mathsf{r}(s),\\
\mathsf{r}(s)&:=R_3(s)-(1-s)R_3(0)-sR_3(1).
\end{split}
\end{equation}
By the one-cutoff bounds above,
$|R_3(s)|\le C(|\mathfrak{x}|+|s\mathfrak{y}|)^3\le Cu_i^{3/2}$. Every term of
$R_3(1)-R_3(0)$ contains a factor $\mathfrak{y}$, so
\begin{equation*}
|s|\,|R_3(1)-R_3(0)|\le C|s\mathfrak{y}|(|\mathfrak{x}|+|\mathfrak{y}|)^2\le Cu_i^{3/2}.
\end{equation*}
Since $\mathsf{r}(s)=R_3(s)-R_3(0)-s[R_3(1)-R_3(0)]$, it follows that $|\mathsf{r}(s)|\le Cu_i^{3/2}$.
In their own comb gauges, $G(0)=\operatorname{Ad}_{W(\gamma_x)}\partial W(P)$ and
$G(1)=\operatorname{Ad}_{W^1(\gamma_x)}\partial W^1(P)$. Divide \eqref{5.7:eq-dilated-action-family-a}
by $u_i$ and express everything through $\Phi_P$ and the rescaled increments
\begin{equation*}
\mathsf{K}_b:=\frac{\ell_i}{\eta}\,\frac{1}{i}\,\mathfrak{y}_b.
\end{equation*}
This gives
\begin{equation}\label{5.7:eq-dilated-action-family-2}
\ell_i^2\eta^{-2}\log G(s)=\Phi_P(B_0)+s\bigl[\Phi_P(B_1)-\Phi_P(B_0)\bigr]
-(s^2-s)\,\mathsf{C}(\mathsf{K},\mathsf{K})+O(L^{-i}).
\end{equation}
The Baker--Campbell--Hausdorff formula applied to $e^{\mathfrak{y}_b}e^{\mathfrak{x}_b}=W^1(\lambda_b)$ shows
\begin{equation*}
\mathsf{K}_b=\Phi_b(B_1)-\Phi_b(B_0)+O\bigl(L^{-i}|\Phi_b(B_1)-\Phi_b(B_0)|\bigr).
\end{equation*}
The letter $\mathsf{K}_b$ denotes a loop increment. It is unrelated to the kernel-sum bound
$\mathsf{K}$ of Lemma~\ref{5.7:lem-two-map-interpolation}.

We now compare the two cutoffs. Make the same construction on $T_{\eta'}$, with $\eta'$ in place of
$\eta$ and with the fine combs $\gamma'_{z'}$ from the same base point $y_0$. Apply it at a fine
plaquette $P'\parallel P$ based at a point $x'$ of the $\eta$-cell of $x$, whose bonds
$b'_j\parallel b_j$ lie within distance $\eta$ of the $b_j$. The fine functions are
\begin{align*}
\Phi^{\mathsf{B}}_{b'}(B)&:=\frac{\ell_i}{\eta'}\,\frac{1}{i}\,\log Y^{\mathsf{B}}(B)(\lambda'_{b'}),\\
\Phi^{\mathsf{B}}_{P'}(B)&:=\ell_i^{2}\eta'^{-2}\log\operatorname{Ad}_{Y^{\mathsf{B}}(B)(\gamma'_{x'})}
\partial Y^{\mathsf{B}}(B)(P').
\end{align*}
Identities \eqref{5.7:eq-dilated-action-family-a} and \eqref{5.7:eq-dilated-action-family-2} hold for
them verbatim. Put $Z_s:=\eta^{-2}\log G(s)$ at the cutoff $\mathsf{A}$, and define $Z_s$ at
$\mathsf{B}$ in the same way with $\eta'$.

Apply Lemma~\ref{5.7:lem-two-map-interpolation}, the interpolation between two analytic maps, to the
pairs $(\Phi^{\mathsf{B}}_{b'},\Phi^{\mathsf{A}}_b)$ and $(\Phi^{\mathsf{B}}_{P'},\Phi^{\mathsf{A}}_P)$.
For the first pair, take the kernel $k(q):=K_{\kappa}e^{-\delta d_{\star}(y_0,q)}$ on data bonds
$q\in\mathfrak{c}\cup\mathfrak{c}'_0$, with $K_{\kappa}:=2.2\cdot 210dL^2\cdot 8L^3B_H$. The kernel
hypothesis of that lemma (among \eqref{5.7:eq-two-map-interpolation-a}) is furnished by the proof of
clause (iv) of the comb lemma, run at a pair $(B,\tilde B)$. That proof bounds $|\mathfrak{E}-1|$
through the ladder products it uses. For the second pair, take the kernel of the same form with a
constant $K_F'$, supplied by the standing bounds (b), (c) and (d). The real-slice closeness hypothesis
of the lemma is, for both pairs, clauses (k) and (F) of the closeness of native loops and plaquettes
at real data, Lemma~\ref{5.7:lem-native-loop-closeness}, \eqref{5.7:eq-native-loop-closeness-a}.

Read in the setting of Proposition~\ref{5.7:prop-response-zeroth-order}, the interpolation lemma gives
three outputs. First,
\begin{equation*}
\bigl|\delta\bigl(\Phi_b(B_1)-\Phi_b(B_0)\bigr)\bigr|
\le K_{\kappa}\bigl(2\varpi+1.1\varpi'+c_1\vartheta^{i}\bigr)T_{\ast}(y_0).
\end{equation*}
Second, the same bound holds for $\Phi_P$.
Third, clause (f) of the interpolation lemma gives
\begin{equation*}
|\delta\Phi_P(B_0)|\le C\bigl(\vartheta^{4i}+a_{\Lambda}\varpi\bigr).
\end{equation*}
By Lemma~\ref{5.7:lem-three-quarter-rate}, the majorants at three quarters of the kernel rate
\eqref{5.7:eq-three-quarter-rate-a}, we have
\begin{equation*}
R^{\perp}_v\cdot 210dL^2\cdot 8L^3B_H\,T_{\ast}\le 1/32.
\end{equation*}
Since $|s|\le R^{\perp}_v$, it follows that
\begin{equation*}
|s|\,\bigl|\delta\bigl(\Phi_P(B_1)-\Phi_P(B_0)\bigr)\bigr|\le C\bigl(2\varpi+1.1\varpi'+c_1\vartheta^{i}\bigr).
\end{equation*}

For the quadratic term, put $S:=\max\{1,|s|\}\le R^{\perp}_v$. Each of the six commutators in
$\mathsf{C}$ is bounded by twice the product of the norms of its entries, so
$|\mathsf{C}(Z',Z'')|\le 6\sup|Z'|\sup|Z''|$. By bilinearity,
\begin{equation*}
\delta\mathsf{C}(\mathsf{K},\mathsf{K})=\mathsf{C}(\delta\mathsf{K},\mathsf{K}^{\mathsf{B}})
+\mathsf{C}(\mathsf{K}^{\mathsf{A}},\delta\mathsf{K}).
\end{equation*}
Together with $|s^2-s|\le|s|^2+|s|\le 2S^2$, this gives
\begin{equation*}
\begin{split}
|s^2-s|\,|\delta\mathsf{C}(\mathsf{K},\mathsf{K})|
&\le 2S^2\cdot 12\sup|\mathsf{K}|\sup|\delta\mathsf{K}|
=24\bigl(S\sup|\mathsf{K}|\bigr)\bigl(S\sup|\delta\mathsf{K}|\bigr)\\
&\le 24\cdot\tfrac{1}{14}\cdot C\bigl(2\varpi+1.1\varpi'+c_1\vartheta^{i}\bigr).
\end{split}
\end{equation*}
The justifications for the last line are as follows.
\begin{itemize}
\item The one-cutoff bound on $s\mathfrak{y}_b$ from the comb-gauge step gives $S\sup|\mathsf{K}|\le 1/14$.
\item The first output above, multiplied by $S\le R^{\perp}_v$ and combined with the bound from
Lemma~\ref{5.7:lem-three-quarter-rate}, controls $S\sup|\delta\mathsf{K}|$.
\item The difference $\delta$ of the Baker--Campbell--Hausdorff correction in $\mathsf{K}_b$ is
bilinear and carries a factor $L^{-i}$; it is accounted for in the term $c_1'''\vartheta^{i}$.
\end{itemize}
Taking the difference between the two cutoffs of \eqref{5.7:eq-dilated-action-family-2}, and
collecting the bounds on $\delta\Phi_P(B_0)$, on the linear term and on the quadratic term, gives,
with a constant $C_Z$,
\begin{equation}\label{5.7:eq-dilated-action-family-3}
|\delta(\ell_i^2Z_s)|\le C_Z\bigl(2\varpi+1.1\varpi'+c_1'''\vartheta^{i}\bigr).
\end{equation}

It remains to sum over plaquettes. For the one-plaquette action,
\begin{equation*}
a(e^{\eta^2Z})=-\operatorname{tr}\bigl(\cosh(\eta^2Z)-1\bigr)
=-\tfrac12\eta^4\operatorname{tr}Z^2\bigl(1+O(u_i^2)\bigr).
\end{equation*}
Because $d=4$, $L^4\eta'^4=\eta^4$. At the fine cutoff,
\begin{equation*}
\mathfrak{F}^{\mathsf{B}}_v(s)=\sum_{P}\sum_{P'\sim P}c(P')\,a'(P'),
\end{equation*}
where $a'$ is the one-plaquette action on $T_{\eta'}$ and there are $L^4$ fine plaquettes $P'$ for
each $P$. The family $\mathfrak{F}^{\mathsf{A}}_v(s)$ is the sum over $P$ of $c(P)a(P)$. Write
$\mathcal{P}_{\Delta}$ for the set of plaquettes $P$ over which these sums run. Since each $P$ has
exactly $L^4$ partners $P'\sim P$,
\begin{equation*}
\delta\mathfrak{F}_v(s)=\sum_{P}L^{-4}\sum_{P'\sim P}\bigl[L^4c(P')a'(P')-c(P)a(P)\bigr],
\end{equation*}
and each bracket splits as
\begin{equation*}
\begin{split}
L^4c(P')a'(P')-c(P)a(P)
&=\bigl(c(P')-c(P)\bigr)L^4a'(P')\\
&\quad+c(P)\bigl(L^4a'(P')-a(P)\bigr).
\end{split}
\end{equation*}
Put $\hat z:=\sup\ell_i^2|Z_s|$, and use the expansion of $a$ and of $a'$ together with
$L^4\eta'^4=\eta^4$. This gives
\begin{equation}\label{5.7:eq-dilated-action-family-b}
\begin{split}
|\delta\mathfrak{F}_v(s)|\le\sum_{P}L^{-4}\sum_{P'\sim P}\Bigl\{&|c(P')-c(P)|\,
\mathfrak{c}_{\mathrm{tr}}u_i^2\hat z^{2}\\
&\quad+|c(P)|\,\mathfrak{c}_{\mathrm{tr}}u_i^2\hat z\,|\delta(\ell_i^2Z_s)|\Bigr\}+O(u_i).
\end{split}
\end{equation}
Three further facts enter the estimate of this right-hand side.
\begin{itemize}
\item The plaquette-logarithm bound of the dilated setting gives $\hat z\le 1$.
\item Since $P'$ is based within distance $\eta$ of the base point of $P$, we have
$|c(P')-c(P)|\le\eta\sup|\nabla c|$.
\item Summing over the plaquettes gives $\sum_{P\in\mathcal{P}_{\Delta}}u_i^2=\mathsf{p}_d$.
\end{itemize}
Insert these three facts and the bound \eqref{5.7:eq-dilated-action-family-3} on
$\delta(\ell_i^2Z_s)$ into \eqref{5.7:eq-dilated-action-family-b}; the inner sum over the $L^4$
partners $P'$ is cancelled by the factor $L^{-4}$, and $\eta=u_i^{1/2}\ell_i$. This gives
\begin{equation*}
\begin{split}
|\delta\mathfrak{F}_v(s)|&\le\mathfrak{c}_{\mathrm{tr}}\mathsf{p}_d\Bigl[u_i^{1/2}\ell_i\sup|\nabla c|\\
&\qquad+C_Z\sup|c|\,\bigl(2\varpi+1.1\varpi'+c_1'''\vartheta^{i}\bigr)\Bigr]+O(u_i).
\end{split}
\end{equation*}
With $\sup|c|\le\mathsf{s}_v$ and $\ell_i\sup|\nabla c|\le\mathsf{s}_v$, this right-hand side is
estimated by the right-hand side of \eqref{5.7:eq-dilated-action-family-1}.
\end{proof}

\begin{theorem}[The two-cutoff value of the remainder without dilation]
\label{5.7:thm-remainder-value}
Work in the setting of Proposition~\ref{5.7:prop-dilated-action-family}: a unit $v=(c,\Delta)$ at
level $i$, with distance $D_v$, the two machines $\mathsf{A}$ (spacing $\eta$) and $\mathsf{B}$
(spacing $\eta'=\eta/L$), and a base point $y_0$ common to the two lattices. For
$\mathsf{m}\in\{\mathsf{A},\mathsf{B}\}$ let $Y^{\mathsf{m}}(B)$ be the canonical axial
representative of machine $\mathsf{m}$ as a function of its data $B$, let $B^{\mathsf{m}}_0$ and
$B^{\mathsf{m}}_1$ be the base datum and the changed datum of machine $\mathsf{m}$, as in
Proposition~\ref{5.7:prop-response-zeroth-order}, let
$W:=Y^{\mathsf{m}}(B^{\mathsf{m}}_0)$ and $W^1:=Y^{\mathsf{m}}(B^{\mathsf{m}}_1)$ in each machine,
and put
\begin{equation*}
G_v:=A(c,W^1)-A(c,W),\qquad \Phi^{\mathsf{m}}(B):=A^{\mathsf{m}}\bigl(c,Y^{\mathsf{m}}(B)\bigr).
\end{equation*}
Let $\varpi\ge\mathsf{m}^{\mathrm{dat}}(\mathfrak{b}^{\mathsf{B}},\mathfrak{b}^{\mathsf{A}})$ and
$\varpi'$ be the data and far-data majorants and $c_1$, $\vartheta$ the constants of
Proposition~\ref{5.7:prop-response-zeroth-order}, and $T_{\ast}$ the profile sum at rate
$\tfrac34\delta$ of Lemma~\ref{5.7:lem-three-quarter-rate}. Set
\begin{equation*}
K_G:=36\,\mathfrak{c}_{\mathrm{tr}}\cdot 8L^3B_H,\qquad
k(c'):=K_G\,\mathfrak{n}_v\,e^{-\delta d_{\star}(y_0,c')}.
\end{equation*}
Assume:
\begin{enumerate}
\item[(a)] the hypotheses of the two-cutoff response statement at zeroth order
(Proposition~\ref{5.7:prop-response-zeroth-order});
\item[(b)] clauses (ii) and (iii) of the corollary on the closeness of the one-plaquette actions on
the two lattices, with the closeness constant $\varepsilon$ and the exponent $\theta$ of that
corollary: (ii) for plaquettes $P$ of the coarser and $P'$ of the finer lattice paired as in that
corollary, with $a'$ the one-plaquette action of the finer lattice,
\begin{equation*}
\bigl|L^4a'(P')-a(P)\bigr|\le C\varepsilon^2u_i^2L^{-\theta i};
\end{equation*}
(iii) the sum of (ii) over a cell, the relative closeness of the actions;
\item[(c)] the plaquette-factor bound of the dilated-action setting, $|\Xi-\mathbb{1}|\le 6hu_i$,
where, as in that setting, $V$ and $\Xi V$ are the plaquette variables entering
$Y^{\mathsf{m}}(B)$ and $Y^{\mathsf{m}}(\tilde B)$ at two data $B$, $\tilde B$, and $h$ is the
parameter of that bound.
\end{enumerate}
Then
\begin{equation}\label{5.7:eq-remainder-value-1}
\begin{aligned}
|\delta G_v|&=\bigl|\delta[A(c,W^1)-A(c,W)]\bigr|
\le K_G\,\mathfrak{n}_v\,(2\varpi+1.1\varpi'+c_1\vartheta^i)\,T_{\ast}(y_0)\\
&\le C\,\mathfrak{n}_v\,(2\varpi+1.1\varpi'+c_1\vartheta^i)\,e^{-\frac14\delta D_v}\,
\alpha_{\cdot,i}.
\end{aligned}
\end{equation}
\end{theorem}

\begin{proof}
Write $B_0:=B^{\mathsf{A}}_0$, $e_0:=B^{\mathsf{B}}_0-B^{\mathsf{A}}_0$ after one data gauge, as
in Proposition~\ref{5.7:prop-response-zeroth-order}, and
$b^{\mathsf{m}}:=B^{\mathsf{m}}_1-B^{\mathsf{m}}_0$. By the definitions of $W$, $W^1$ and
$\Phi^{\mathsf{m}}$, in machine $\mathsf{m}$ one has
$G_v=\Phi^{\mathsf{m}}(B^{\mathsf{m}}_0+b^{\mathsf{m}})-\Phi^{\mathsf{m}}(B^{\mathsf{m}}_0)$, so that
\begin{equation*}
\delta G_v
=\bigl[\Phi^{\mathsf{B}}(B_0+e_0+b^{\mathsf{B}})-\Phi^{\mathsf{B}}(B_0+e_0)\bigr]
-\bigl[\Phi^{\mathsf{A}}(B_0+b^{\mathsf{A}})-\Phi^{\mathsf{A}}(B_0)\bigr],
\end{equation*}
which is the mixed second difference $\Delta^{\times}$ of
\eqref{5.7:eq-two-map-interpolation-b}, formed with the maps $\Phi^{\mathsf{A}}$,
$\Phi^{\mathsf{B}}$. We verify the hypotheses of \eqref{5.7:eq-two-map-interpolation-a} with the
kernel $k$ defined above.

\emph{The one-machine Lipschitz hypothesis.} For the plaquette variables $V$ and $\Xi V$ entering
$Y^{\mathsf{m}}(B)$ and $Y^{\mathsf{m}}(\tilde B)$, the one-plaquette action $a$ satisfies,
plaquette by plaquette,
\begin{equation*}
\bigl|a(\Xi V)-a(V)\bigr|\le \tfrac32\,\mathfrak{c}_{\mathrm{tr}}
\bigl(|V-\mathbb{1}|+|\Xi-\mathbb{1}|\bigr)|\Xi-\mathbb{1}|,
\end{equation*}
and the factor $|\Xi-\mathbb{1}|$ is bounded by $6hu_i$ by hypothesis~(c).

\emph{The real-slice hypothesis.} For real common data, the two-cutoff difference is estimated
plaquette by plaquette by clause~(ii) of hypothesis~(b),
$|L^4a'(P')-a(P)|\le C\varepsilon^2u_i^2L^{-\theta i}$, whose sum over a cell is clause~(iii) of
hypothesis~(b), the relative closeness of the actions; the oscillation of the plaquette weights is
controlled as in \eqref{5.7:eq-dilated-action-family-b} in the proof of
Proposition~\ref{5.7:prop-dilated-action-family}.

With these two hypotheses in place, the assembly clause~(e) of
Lemma~\ref{5.7:lem-two-map-interpolation}, applied to $\Delta^{\times}=\delta G_v$, gives the
first inequality of \eqref{5.7:eq-remainder-value-1}; the second is the bound of clause~(ii) of
Lemma~\ref{5.7:lem-three-quarter-rate} on $T_{\ast}$, at $x=y_0$.
\end{proof}

Thus the two-cutoff value of the remainder $G_v$ is small at the rate $e^{-\frac14\delta D_v}$,
times the sum of the mismatch $2\varpi+1.1\varpi'$ and the index term $c_1\vartheta^i$. The
argument needs neither the passage to the closure of the class of pairs nor the dilation of
Proposition~\ref{5.7:prop-dilated-action-family}.

\begin{definition}[The energy-difference datum and its layers]\label{5.7:not-energy-datum-layers}
Fix the frame $\mathfrak{B}_j$, the base $\mathbb{W}$, the configuration $\mathbb{J}$, the parameter $\sigma$ and the
interpolation parameters $t,\mathsf{t}$. Put
\begin{equation}\label{5.7:eq-energy-datum-layers-1}
\mathsf{B}(t,\mathsf{t}) := Q_{\mathfrak{B}_j}\bigl(\mathbb{W};\,
\eta\,(t+\mathsf{t}\,\zeta_{\square})\,\mathbb{H}^{w}\bigr),
\end{equation}
where $Q_{\mathfrak{B}_j}(\mathbb{W};\cdot)$ is the averaging map of the frame $\mathfrak{B}_j$
about the base $\mathbb{W}$.
Let $\mathbf{H}_j(t,\mathsf{t})$ be the third member furnished by the response-layer theorem at
the data
\begin{equation*}
\bigl(\mathfrak{B}_j;\,\mathbb{W},\mathbb{J};\,\sigma;\,(0,\mathsf{B}(t,\mathsf{t}))\bigr).
\end{equation*}
Set
\begin{align}
D^{\sharp} &:= \zeta'\, Q_{\star}\bigl((t+\mathsf{t}\,\zeta_{\square})\,\mathbb{H}^{w}\bigr)
 + (1-\zeta')\, Q_{\star}\bigl(\mathbf{H}_j(t,\mathsf{t})\bigr),
 \label{5.7:eq-energy-datum-layers-2}\\
\tilde{\mathsf{B}} &:= Q_{\star}\bigl(\mathbf{A}(t,\mathsf{t})\bigr),
 \label{5.7:eq-energy-datum-layers-3}
\end{align}
where $\mathbf{A}(t,\mathsf{t})$ is the field of the energy-difference construction at the
interpolation parameters $t,\mathsf{t}$.
Define the \emph{energy-difference datum}
\begin{equation}\label{5.7:eq-energy-datum-layers-4}
\mathsf{E} := \frac{1}{i}\,\log\bigl(e^{iD^{\sharp}}\,e^{-i\tilde{\mathsf{B}}}\bigr).
\end{equation}
On the set $F_{\square}$ of the energy-difference construction, which contains
$\operatorname{supp}\mathsf{E}$, the \emph{three layers} of $\mathsf{E}$ are
\begin{equation}\label{5.7:eq-energy-datum-layers-5}
t\,\mathbb{H}^{w}, \qquad t\,\zeta'_{\square}\,\mathbb{H}^{w}, \qquad \mathbf{H}_j(t,\mathsf{t}).
\end{equation}
The map $Q_{\star}$ is local and analytic in its argument, and for every bond $c$ it
satisfies
\begin{equation}\label{5.7:eq-energy-datum-layers-6}
\bigl|Q_{\star}(A)(c)\bigr| \le 2\,\ell_{l(c)}\,\sup_{\tilde B(c)}|A| .
\end{equation}
Here $l(c)$ is the level of $c$, $\ell_{l(c)}$ is the length scale of that level, and
$\tilde B(c)$ is a set attached to $c$ (not to be confused with the field $\tilde{\mathsf{B}}$
of \eqref{5.7:eq-energy-datum-layers-3}).
Consequently every estimate of $\mathsf{E}(c)$ obtained through
\eqref{5.7:eq-energy-datum-layers-6} involves the layers only through their suprema over the
sets $\tilde B(c)$; no other norm of the layers is used.
\end{definition}

\begin{proof}[Proof of the last assertion]
By \eqref{5.7:eq-energy-datum-layers-4}, $\mathsf{E}(c)$ is a function of the pair
$\bigl(D^{\sharp}(c),\tilde{\mathsf{B}}(c)\bigr)$ alone.

By \eqref{5.7:eq-energy-datum-layers-2}, $D^{\sharp}(c)$ is formed from
$Q_{\star}\bigl((t+\mathsf{t}\zeta_{\square})\mathbb{H}^{w}\bigr)(c)$ and
$Q_{\star}\bigl(\mathbf{H}_j(t,\mathsf{t})\bigr)(c)$, with the weights $\zeta'$ and $1-\zeta'$.
By \eqref{5.7:eq-energy-datum-layers-3}, $\tilde{\mathsf{B}}(c)$ is
$Q_{\star}\bigl(\mathbf{A}(t,\mathsf{t})\bigr)(c)$.

Each of these three values is a local analytic functional of the corresponding field, and by
\eqref{5.7:eq-energy-datum-layers-6} each is bounded by $2\ell_{l(c)}$ times the supremum of
that field over $\tilde B(c)$.

Hence $\mathsf{E}(c)$ is a function of the three values of $Q_{\star}$ at $c$, and the bounds on
$\mathsf{E}(c)$ obtained through \eqref{5.7:eq-energy-datum-layers-6} therefore
involve the fields $(t+\mathsf{t}\zeta_{\square})\mathbb{H}^{w}$, $\mathbf{H}_j(t,\mathsf{t})$ and
$\mathbf{A}(t,\mathsf{t})$ only through those suprema.
\end{proof}

\begin{lemma}[The flat profile at five eighths of the rate]\label{5.7:lem-flat-profile}
Let $p$ be the profile of the datum $b$ at level $j$, and let $\hat p$, $\hat p_{\square}$ and
$f := \hat p + |\mathsf{t}|\,\hat p_{\square}$ be its smoothings, as fixed in the definition of the
smoothed profiles. Define the \emph{flat profile} and its smoothings by
\begin{equation}\label{5.7:eq-flat-profile-a}
\begin{gathered}
p^{\flat}(y) := \sup|b|\; e^{-\frac{5}{8}\delta\, d^{\wedge}(y,\mathcal{C}_j)},\\
\hat p^{\flat} := \bigl(\varrho\, p^{\flat}\bigr)^{[\delta/2]},\qquad
\hat p^{\flat}_{\square} := \bigl(\varrho\, p^{\flat}\,\mathbb{1}_{\square^{\zeta}}\bigr)^{[\delta/2]},\\
f^{\flat} := \hat p^{\flat} + |\mathsf{t}|\,\hat p^{\flat}_{\square},
\end{gathered}
\end{equation}
where $\mathcal{C}_j$ is the set entering the exponent of $p$, fixed together with $p$ in the
definition of the smoothed profiles; one has $p^{\flat} \ge p$ and $f^{\flat} \ge f$. Then:
\begin{enumerate}
\item[(a)] the pointwise profile bound holds with $p$ replaced by $p^{\flat}$ and with the same
right member;
\item[(b)] consequently, the two displays of clause (iii) of the lemma on smoothed profiles, and
the bounds (i) and (ii) of Lemma~\ref{5.7:lem-three-quarter-rate} (majorants at three quarters
of the kernel rate), displayed in \eqref{5.7:eq-three-quarter-rate-a}, hold with $f$ replaced by
$f^{\flat}$.
\end{enumerate}
\end{lemma}

\begin{proof}
We first record the two comparisons stated with the definitions. The profile $p$ is defined as
$p^{\flat}$ is, with the same prefactor $\sup|b|$ and with the rate $\delta$ in place of
$\frac58\delta$ in the exponent:
\begin{equation*}
p(y) = \sup|b|\; e^{-\delta\, d^{\wedge}(y,\mathcal{C}_j)}.
\end{equation*}
Since $\frac58\delta < \delta$ and $d^{\wedge} \ge 0$, this gives $p^{\flat} \ge p$ pointwise. The
smoothing
$(\cdot)^{[\delta/2]}$ is monotone, and so is multiplication by $\varrho$ and by the indicator
$\mathbb{1}_{\square^{\zeta}}$; hence $\hat p^{\flat} \ge \hat p$,
$\hat p^{\flat}_{\square} \ge \hat p_{\square}$, and $f^{\flat} \ge f$.

We turn to (a). In the proof of the pointwise profile bound the exponent of $p$ is used at
exactly one place: at the point $\bar y'$, where the proof of the absorption estimate converts the
factor
\begin{equation}\label{5.7:eq-flat-profile-1}
3C_1\lambda\,\bigl(1+|m'-j|\bigr)/\rho_{m'}
\end{equation}
(notation of that proof) into $\vartheta$, at the price of a factor
$e^{-\frac12\delta\, d^{\wedge}(\bar y',\mathcal{C}_j)}$.
The other half of the rate, $\frac12\delta$, is kept in the right member, where it appears as the
factor $e^{-\frac12\delta\, d^{\wedge}(\bar y',\mathcal{C}_j)}$. It therefore suffices to show that
the conversion of \eqref{5.7:eq-flat-profile-1} into $\vartheta$ can be achieved at the price
$e^{-\frac18\delta\, d^{\wedge}(\bar y',\mathcal{C}_j)}$; for then the exponent of $p^{\flat}$,
which is $\frac58\delta = \frac18\delta + \frac12\delta$, pays the price with its first part and
leaves the second part, $\frac12\delta$, in the right member exactly as before, so that the right
member is unchanged.

The proof of the absorption estimate consists of the inequality
\begin{equation}\label{5.7:eq-flat-profile-2}
v_{\rho}(a) := 3(1+a)\,2^{a+1}\,e^{-\rho w (a-1)^{+}} \;\le\; 24
\qquad \text{for every integer } a \ge 0,
\end{equation}
at the rate $\rho = \frac14\delta$; here $v_{\rho}(a)$ is notation introduced for the left member,
and $w$ and the integer $a$ are the quantities of the absorption estimate.
We verify \eqref{5.7:eq-flat-profile-2} at $\rho = \frac18\delta$.

For $a = 0$ and $a = 1$ one has $(a-1)^{+} = 0$, and the left member equals
$3\cdot 1\cdot 2 = 6$ and $3\cdot 2\cdot 4 = 24$ respectively; both are at most $24$.

Let $a \ge 2$. The level-spacing condition $\delta w \ge 16\ln L + 4$ gives
$\frac18\delta w \ge 2\ln L + \frac12$, that is,
\begin{equation}\label{5.7:eq-flat-profile-3}
e^{-\frac18\delta w} \le L^{-2}e^{-1/2}.
\end{equation}
At $a = 2$ the left member of \eqref{5.7:eq-flat-profile-2} is $3\cdot 3\cdot 8\,e^{-\frac18\delta w}$,
and by \eqref{5.7:eq-flat-profile-3}
\begin{equation*}
v_{\delta/8}(2) = 72\, e^{-\frac18\delta w} \le 72\,L^{-2}e^{-1/2} < 24 .
\end{equation*}
For $a \ge 2$ one has $(a-1)^{+} = a-1$ and $(a)^{+} = a$, so the ratio of consecutive values is
\begin{equation*}
\frac{v_{\delta/8}(a+1)}{v_{\delta/8}(a)}
= \frac{2+a}{1+a}\cdot 2\cdot e^{-\frac18\delta w}
\le \frac43\cdot 2\cdot L^{-2}e^{-1/2}
= \frac83\,L^{-2}e^{-1/2},
\end{equation*}
where we used $(2+a)/(1+a) \le \frac43$ for $a \ge 2$ and \eqref{5.7:eq-flat-profile-3}. The
inequalities $72\,L^{-2}e^{-1/2} < 24$ and $\frac83 L^{-2}e^{-1/2} < 1$ both hold whenever
$L \ge 2$, since then $L^{-2}e^{-1/2} \le \frac14 e^{-1/2} < \frac14$; in particular they hold
for the blocking factor $L$. Hence, for $a \ge 2$, the values $v_{\delta/8}(a)$ form a
non-increasing sequence beginning with a value below $24$, and are therefore all at most $24$.
This proves \eqref{5.7:eq-flat-profile-2} at $\rho = \frac18\delta$ with the same constant $24$.

So the price $e^{-\frac18\delta\, d^{\wedge}(\bar y',\mathcal{C}_j)}$ suffices for the conversion
of \eqref{5.7:eq-flat-profile-1} into $\vartheta$. Since $\frac58\delta = \frac18\delta+\frac12\delta$,
the argument above shows that the pointwise profile bound holds for $p^{\flat}$ with the same
right member. This is (a).

For (b): the right member being the same, the two displays of clause (iii) of the lemma on
smoothed profiles and the bounds (i), (ii) of Lemma~\ref{5.7:lem-three-quarter-rate}, which are
drawn from the pointwise profile bound, hold with $p$, $\hat p$, $\hat p_{\square}$ replaced by
$p^{\flat}$, $\hat p^{\flat}$, $\hat p^{\flat}_{\square}$, that is, with $f$ replaced by
$f^{\flat}$.
\end{proof}

\begin{lemma}[The flat energy-difference bound]\label{5.7:lem-flat-energy-bound}
Let $\mathsf{A}$ and $\mathsf{B}$ be the two cutoffs, on $T_{\eta}$ and on $T_{\eta'}$ with $\eta'=\eta/L$.
For a functional $\Phi$ write $\delta\Phi:=\Phi^{\mathsf{B}}-\Phi^{\mathsf{A}}$, and read every
layer of $\mathsf{B}$ in the carried frame through $Y[\cdot\,;\cdot]$. Let $p^{\flat}$,
$\hat p^{\flat}$ and $f^{\flat}$ be the flat profile and its smoothings of
\eqref{5.7:eq-flat-profile-a}. Assume the relative closeness, in sup norm, of the layers
$\mathbb{H}^{w}$ and $\mathbf{H}_j$ across the two cutoffs, that is, the following. At a pair of
the stage class, on the parameter domains of the construction and in the carried frame,
\begin{equation}\label{5.7:eq-flat-energy-bound-b}
\begin{aligned}
&\text{(B1)}\quad \ell_{\mathfrak{B}}(y)\,
\bigl|Y[\mathbb{H}^{w,\mathsf{B}};\mathbb{W}^{\mathsf{B}}]-\mathbb{H}^{w,\mathsf{A}}\bigr|(y)
\le \varpi_{\mathbb{H}}\,B_{\sharp}\,p^{\flat}(y)\qquad\text{for all } y,\\
&\text{(B2)}\quad \ell_{\mathfrak{B}_j}(y)\,
\bigl|Y[\mathbf{H}_j^{\mathsf{B}}(\mathsf{d}^{\mathsf{B}})]
-\mathbf{H}_j^{\mathsf{A}}(\mathsf{d}^{\mathsf{A}})\bigr|(y)
\le 2B_{\sharp}(\varpi_{\mathbb{H}}+\varpi_{\mathsf{d}})\,\mathsf{d}_{\ast}(y),\\
&\text{(B3)}\quad \ell\,\sup_{\square_0}|A_W|\le \varpi_W .
\end{aligned}
\end{equation}
Here (B2) is required for all data $\mathsf{d}^{\mathsf{A}},\mathsf{d}^{\mathsf{B}}$ on
$\mathfrak{B}_j$ that satisfy $|\mathsf{d}^{\mathsf{m}}|\le\mathsf{d}_{\ast}$ for
$\mathsf{m}\in\{\mathsf{A},\mathsf{B}\}$ and $|\delta\mathsf{d}|\le\varpi_{\mathsf{d}}\mathsf{d}_{\ast}$,
where $\mathsf{d}_{\ast}\in\mathfrak{M}^{\wedge}_{\delta/2}(\mathfrak{B}_j)$.
Then the energy-difference datum $\mathsf{E}^{u}$, in the notation fixed
in~\ref{5.7:not-energy-datum-layers}, satisfies
\begin{equation}\label{5.7:eq-flat-energy-bound-a}
|\delta\mathsf{E}^{u}(c)|\le \varpi'\cdot 15\,B_{\sharp}^{2}\,f^{\flat}(\bar c)
\qquad\text{on } F_{\square},
\end{equation}
with
\begin{equation*}
\varpi':=C_E\bigl(\varpi_{\mathbb{H}}+\varpi_W+L^{-l(c)}\bigr),
\end{equation*}
where the constant $C_E=C_E(d,c_4)$ depends only on $d$ and $c_4$.
\end{lemma}

\begin{proof}
We recall the datum in the notation fixed in~\ref{5.7:not-energy-datum-layers}. With
\begin{equation*}
D^{\sharp}:=\zeta'_{\square}\,Q_{\star}\bigl((t+\mathsf{t}\zeta_{\square})\mathbb{H}^{w}\bigr)
+(1-\zeta'_{\square})\,Q_{\star}\bigl(\mathbf{H}_j(t,\mathsf{t})\bigr)
\quad\text{and}\quad
\tilde{\mathsf{B}}:=Q_{\star}\bigl(\mathbf{A}(t,\mathsf{t})\bigr),
\end{equation*}
it is defined by $\mathsf{E}:=\tfrac1i\log\bigl(e^{iD^{\sharp}}e^{-i\tilde{\mathsf{B}}}\bigr)$.
On $F_{\square}$, which contains the support of $\mathsf{E}$, the datum reads exactly three layers,
namely $t\mathbb{H}^{w}$, $t\zeta'_{\square}\mathbb{H}^{w}$ and $\mathbf{H}_j(t,\mathsf{t})$. Each
$Q_{\star}(A)(c)$ satisfies $|Q_{\star}(A)(c)|\le 2\ell_{l(c)}\sup_{\tilde B(c)}|A|$, so that only
sup norms of the layers enter.

We therefore apply the reduction of the finer averaged datum to one lattice,
\eqref{5.7:eq-datum-reduction-a} of Lemma~\ref{5.7:lem-datum-reduction}, at each of the three
layers. That bound has two members. The first is proportional to
$\ell_{l(c)}\sup_{\tilde B(c)}|Y[A';\mathbb{W}']-A|$, the distance between the averaged finer
layer and the coarser layer. The second is proportional to the product of
$\ell\sup_{\tilde B(c)}|A_W|$ with $\ell_{l(c)}\sup_{\tilde B(c)}|A|$.

\emph{The layers $t\mathbb{H}^{w}$ and $t\zeta'_{\square}\mathbb{H}^{w}$.}
Three bounds enter here. First, (B1) of \eqref{5.7:eq-flat-energy-bound-b} bounds the first member
by $\varpi_{\mathbb{H}}B_{\sharp}p^{\flat}$. Second, (B3) bounds the factor
$\ell\sup|A_W|$ of the second member by $\varpi_W$. Third, the remaining factor
$\ell_{l(c)}\sup|\mathbb{H}^{w}|$ of the second member is controlled by the size bound of the
first member of the layer theorem, $\ell_{l(c)}|\mathbb{H}^{w}|\le B_{\sharp}\varrho p$,
continued as follows:
\begin{equation*}
\ell_{l(c)}|\mathbb{H}^{w}|\le B_{\sharp}\,\varrho\, p
\le e^{\delta/2}B_{\sharp}\,\hat p^{\flat}(\bar c).
\end{equation*}
In the layer $t\zeta'_{\square}\mathbb{H}^{w}$ the scaffold cut-off is moved through $Y$; by the
commutation bound for the scaffold cut-offs recorded with Lemma~\ref{5.7:lem-datum-reduction},
$|Y[\zeta A']-\zeta Y[A']|$ is bounded by $CL^{-l}M^{-1}\sup|A'|$, which is the origin of the
term $L^{-l}$ below.

\emph{The datum $\mathsf{B}(t,\mathsf{t})$.}
The third layer is the third member of the layer theorem evaluated at the datum
$\mathsf{B}(t,\mathsf{t}):=Q_{\mathfrak{B}_j}\bigl(\mathbb{W};\eta(t+\mathsf{t}\zeta_{\square})
\mathbb{H}^{w}\bigr)$. We apply the reduction \eqref{5.7:eq-datum-reduction-a} on
$\mathfrak{B}_j$, with the layer $(t+\mathsf{t}\zeta_{\square})\mathbb{H}^{w}$. This gives
\begin{equation}\label{5.7:eq-flat-energy-bound-1}
|\delta\mathsf{B}|\le C\bigl(\varpi_{\mathbb{H}}+\varpi_W+L^{-l}\bigr)\cdot 4B_{\sharp}f^{\flat}
=:\varpi_{\mathsf{d}}\,\mathsf{d}_{\ast},
\qquad \mathsf{d}_{\ast}:=4B_{\sharp}f^{\flat}.
\end{equation}
The majorant $\mathsf{d}_{\ast}$ belongs to $\mathfrak{M}^{\wedge}_{\delta/2}(\mathfrak{B}_j)$ by
clause (g4) of the majorant lemma. Recall that
$f^{\flat}=\hat p^{\flat}+|\mathsf{t}|\hat p^{\flat}_{\square}$. The rate of $p^{\flat}$ is used
only in the member $|\mathsf{t}|\hat p^{\flat}_{\square}$, which lives on $\square^{\zeta}$.

\emph{The layer $\mathbf{H}_j(t,\mathsf{t})$.}
We apply (B2) of \eqref{5.7:eq-flat-energy-bound-b} with the data
$\mathsf{d}^{\mathsf{m}}:=\mathsf{B}^{\mathsf{m}}(t,\mathsf{t})$, and with $\varpi_{\mathsf{d}}$
and $\mathsf{d}_{\ast}$ taken from \eqref{5.7:eq-flat-energy-bound-1}. It bounds the first member
of the reduction at this layer by
\begin{equation*}
2B_{\sharp}(\varpi_{\mathbb{H}}+\varpi_{\mathsf{d}})\cdot 4B_{\sharp}f^{\flat}.
\end{equation*}
Since $\varpi_{\mathsf{d}}=C(\varpi_{\mathbb{H}}+\varpi_W+L^{-l})$, we have
$\varpi_{\mathbb{H}}+\varpi_{\mathsf{d}}\le(1+C)(\varpi_{\mathbb{H}}+\varpi_W+L^{-l})$.

\emph{Assembly.}
On $F_{\square}$ the quantity $D^{\sharp}$ is a convex combination, with the scaffold weights
$\zeta'_{\square}$ and $1-\zeta'_{\square}$, of the values of $Q_{\star}$ at the layers just
treated. Hence $\delta D^{\sharp}$ is bounded by the corresponding convex combination of their
bounds.

For $\mathsf{E}=\tfrac1i\log(e^{iD^{\sharp}}e^{-i\tilde{\mathsf{B}}})$ we use the
Baker--Campbell--Hausdorff series at small arguments. It consists of $D^{\sharp}-\tilde{\mathsf{B}}$
and of terms at least bilinear in $(D^{\sharp},\tilde{\mathsf{B}})$. The difference between the two
cutoffs is therefore estimated bilinearly, by the differences of the arguments times their sizes.

Collecting the bounds of the three layers gives \eqref{5.7:eq-flat-energy-bound-a}, with
$\varpi'=C_E(\varpi_{\mathbb{H}}+\varpi_W+L^{-l(c)})$ and $C_E$ depending only on $d$ and $c_4$.
\end{proof}

Wherever the two-cutoff response statement at zeroth order, and the corollary, the proposition and
the theorem that read the energy-difference bound in the same way, use the energy-difference bound
with the full-rate profile $\hat p+|\mathsf{t}|\hat p_{\square}$, the datum enters only through the
interpolated weighted sum
$T^{(1)}[b^{\mathsf{B}}-b^{\mathsf{A}}]$, that is, against the full kernel $e^{-\delta d_{\star}}$.
We take the common profile of the far data to be
$\mathfrak{h}_{\ast}:=15B_{\sharp}^{2}f^{\flat}(\bar c)$ throughout. With this choice these four
statements, and every majorant they use, remain valid by Lemma~\ref{5.7:lem-flat-profile}. The bound
\eqref{5.7:eq-flat-energy-bound-a} may then be used wherever the energy-difference bound with the
full-rate profile is required.

\begin{proposition}[Layer closeness at critical parents]\label{5.7:prop-critical-layer-closeness}
Assume the following statements.
\begin{enumerate}
\item[(A1)] The refinement statement for critical configurations and its clause on own averages:
a critical configuration is critical for every refining set of blocks at its own averages.
\item[(A2)] Clauses (a) and (d) of the layer theorem for the layers $\mathbb{H}^{w}$ and
$\mathbf{H}_j$.
\item[(A3)] Clause (i) of the relative real-slice response bound, on the pair of frames
$(\mathfrak{B}',\mathfrak{B})$.
\item[(A4)] The decay statement for the scaled distances $d$, $d^{\wedge}$ and $d_j$ of the
multiscale frames, by which a loss of $\tfrac14\delta$ in the decay rate $\delta>0$ of the kernels
is absorbed; the letter $\delta$ denotes this rate throughout.
\item[(A5)] Clause (0) of the carried-frame theorem.
\item[(A6)] The data-gauge identity $b^{\mathsf{B}}-b^{\mathsf{A}}=(\operatorname{Ad}_g-1)b$ for a
common variable $b$ read after a data gauge $g$, $b^{\mathsf{A}}$ and $b^{\mathsf{B}}$ denoting
its readings in the two constructions.
\end{enumerate}
Let the datum $\mathsf{s}$ be trivial, $\mathsf{s}\equiv\mathbb{1}$, and let $\mathfrak{q}$ be a
critical pair. That is, for each of the two constructions
$\mathsf{m}\in\{\mathsf{A},\mathsf{B}\}$, at the cutoffs $\eta$ and $\eta'=\eta/L$ respectively,
\begin{equation*}
\mathbb{U}^{\mathsf{m}}=U^{\mathsf{m}}\bigl(\mathfrak{B}^{+},e^{i\beta^{\mathsf{m}}}V_{\circ}\bigr),
\end{equation*}
each with its own current. Here $V_{\circ}$ is real and satisfies the small-field condition on
real configurations, $\beta^{\mathsf{m}}$ is complex and lies in
$\mathfrak{D}(a_{\Lambda}/8)$, and the labels of the two constructions are exactly paired.

Then the three bounds \eqref{5.7:eq-flat-energy-bound-b} hold on the parameter domains: complex
$b$ with $\|b\|_\infty\le 1.1\,C_1 r'_{\mathsf{j}}$, $t\in[0,1]$ and
$|\mathsf{t}|\le r^{\mathrm{rem}}_{\square}$, $t$ and $\mathsf{t}$ being the two interpolation
parameters of the layer $\mathsf{B}(t,\mathsf{t})$ read in Lemma~\ref{5.7:lem-flat-energy-bound};
here $1.1\,C_1 r'_{\mathsf{j}}$ is the radius of
the disc for $b$ at the index $\mathsf{j}$, and $r^{\mathrm{rem}}_{\square}$ is the radius of the
disc for $\mathsf{t}$ attached to the block $\square$. They hold with
\begin{equation}\label{5.7:eq-critical-layer-closeness-1}
\varpi_{\mathbb{H}}:=C\bigl(\mathsf{m}^{\mathrm{dat}}(\mathfrak{q})+c_1L^{-\theta_1 m(y)/16}\bigr),
\qquad
\varpi_W:=C\bigl(\mathsf{m}^{\mathrm{dat}}(\mathfrak{q})+L^{-\theta_1 m/2}\bigr),
\end{equation}
where $C$ is a constant, $\mathsf{m}^{\mathrm{dat}}(\mathfrak{q})$ is the data mismatch
\eqref{5.7:eq-data-mismatch-a} of the pair $\mathfrak{q}$, $\theta_1$ is the exponent in the index
base $L^{-\theta_1/8}$ of the two-cutoff response statement at zeroth order, and $m(y)$ and $m$ are
level indices, $m(y)$ depending on the point $y$ at which the first of the bounds
\eqref{5.7:eq-flat-energy-bound-b} is read; this $m$ is not the torus index of
$\mathbb{T}_m$.

This class of pairs is the class $\mathcal{K}^{\mathrm{crit}}$ of critical pairs. It contains the
seed pairs with which the construction of the classes of pairs starts, and it consists of the pairs
obtained from slice pairs by restriction, by vertex moves without dilation and by complex
increments of the data.
\end{proposition}

\begin{proof}
\emph{The layers as increments of the response function.} By (A1), each configuration
$\mathbb{U}^{\mathsf{m}}$ of the pair $\mathfrak{q}$ is critical for every refining set of blocks
at its own averages. Hence, at real
$\beta^{\mathsf{m}}$ and real $b$, the base $\mathbb{W}$ of the layer theorem (A2) coincides with
$\mathbb{U}$ modulo a gauge transformation on the base block $\square_0$.

By clause (d) of the layer theorem (A2), $\mathbb{H}^{w}(b)$ is then the Landau-gauge function of
Ba{\l}aban's construction on the frame $\mathfrak{B}$ about $\mathbb{W}$. Likewise
$\mathbf{H}_j(\mathsf{d})$ is the corresponding function on $\mathfrak{B}_j$. Equivalently,
with $M_{\mathfrak{B}}$ the block average over the frame $\mathfrak{B}$,
\begin{equation}\label{5.7:eq-critical-layer-closeness-2}
e^{i\eta\mathbb{H}^{w}}\mathbb{W}\cong U\bigl(\mathfrak{B},e^{ib}M_{\mathfrak{B}}\mathbb{W}\bigr),
\end{equation}
where $\cong$ denotes equality modulo gauge.

Both sides of \eqref{5.7:eq-critical-layer-closeness-2} are holomorphic in $(\beta,b)$. For the
left side this is clause (a) of the layer theorem (A2); for the right side it is
\cite[Proposition~9]{B-CMP102b}. Two holomorphic functions on a ball centred at a real point that
agree at its real points agree on the whole ball. The identity
\eqref{5.7:eq-critical-layer-closeness-2} therefore holds on the complex ball.

Consequently the layers are increments of the response function $\mathcal{H}$ on the frames
$\mathfrak{B}$ and $\mathfrak{B}_j$. The base datum is $B_0$, defined as the datum of
$M_{\mathfrak{B}}\mathbb{U}(\beta)$ relative to $M_{\mathfrak{B}}U(\mathfrak{B}^{+},V_{\circ})$,
and $e_0$ is the increment of the base datum of $\mathsf{B}$ over that of $\mathsf{A}$, read after
a data gauge $g$, as in \eqref{5.7:eq-data-mismatch-a}.

\emph{The first bound.} We apply the assembly clause (e) of the interpolation between two
analytic maps (Lemma~\ref{5.7:lem-two-map-interpolation}; see
\eqref{5.7:eq-two-map-interpolation-b}) to these increments. We take $\theta:=\tfrac18$ and the
kernel
\begin{equation*}
k(c'):=K_R\,e^{-\delta d(y,c')},
\end{equation*}
where $d(y,c')$ is the scaled distance of (A4) between the point $y$ and the index $c'$.
Hypothesis ($\Lambda2$) of \eqref{5.7:eq-two-map-interpolation-a} on the pair of frames
$(\mathfrak{B}',\mathfrak{B})$ is (A3). The three members of the bound on $\Delta^{\times}$ are
treated in turn.

The first member, $28\varepsilon_{\star}^{\theta/2}T^{(1-\theta)}[b^{\mathsf{A}}]$, is the term
in which the two constructions are compared at common data. We take as the bound of the kernel sum
$\mathsf{K}:=K_R c_1^{b}$, $c_1^{b}$ being the constant that bounds
$\sum_{c'}e^{-\delta d(y,c')}$; then $\mathsf{K}^{1/8}$ contributes the factor
$K_R^{1/8}(c_1^b)^{1/8}$. Further
$28\varepsilon_{\star}^{\theta/2}=28\varepsilon_{\star}^{1/16}$. Since $b$ is a common variable of
the two constructions, $b^{\mathsf{A}}=b$, and
\begin{equation}\label{5.7:eq-critical-layer-closeness-3}
\begin{aligned}
28\varepsilon_{\star}^{1/16}\,T^{(7/8)}[b]
&\le 28\varepsilon_{\star}^{1/16}K_R(c_1^{b})^{1/8}
\sum_{c'\in\mathcal{C}_{\mathsf{j}}}e^{-\frac78\delta d(y,c')}\,|b(c')|\\
&\le C\varepsilon_{\star}^{1/16}\,\sup|b|\;e^{-\frac58\delta d^{\wedge}(y,\mathcal{C}_{\mathsf{j}})},
\end{aligned}
\end{equation}
where $\mathcal{C}_{\mathsf{j}}$ is the index set of the variable $b$.
The last inequality is (A4), which takes $\tfrac14\delta$ of the rate $\tfrac78\delta$. The
resulting rate $\tfrac58\delta$ is the rate of the flat profile $p^{\flat}$ of
\eqref{5.7:eq-flat-profile-a}.

The third member is $T^{(1)}[b^{\mathsf{B}}-b^{\mathsf{A}}]$. Since $b$ is a common variable of
the two constructions, (A6) gives $b^{\mathsf{B}}-b^{\mathsf{A}}=(\operatorname{Ad}_g-1)b$, where
$g$ is the data gauge.
Hence
\begin{equation*}
T^{(1)}[b^{\mathsf{B}}-b^{\mathsf{A}}]\le\sup|\operatorname{Ad}_g-1|\;T^{(1)}[b],
\end{equation*}
so this member is of the size of the relative data mismatch and is absorbed into it; it therefore
contributes to the term $C\,\mathsf{m}^{\mathrm{dat}}(\mathfrak{q})$ of $\varpi_{\mathbb{H}}$.

The second member is $(4\sup|e_0|/a)\,T^{(1)}[b^{\mathsf{A}}]$, with $a=a_{\Lambda}$; taking the
infimum over data gauges, its prefactor is the data mismatch $\mathsf{m}^{\mathrm{dat}}(\mathfrak{q})$
of \eqref{5.7:eq-data-mismatch-a}. It is linear in $\sup|e_0|$ and
carries the kernel at the full rate $\delta$.

Together these give the first bound of \eqref{5.7:eq-flat-energy-bound-b} with
$\varpi_{\mathbb{H}}$ as in \eqref{5.7:eq-critical-layer-closeness-1}.

\emph{The second bound.} We apply clause (e) of
Lemma~\ref{5.7:lem-two-map-interpolation} at $\theta:=\tfrac14$ on the frame $\mathfrak{B}_j$.
The weighted kernel sum satisfies
\begin{equation}\label{5.7:eq-critical-layer-closeness-4}
\sum_{c}e^{-\frac34\delta d_j(y,c)}\,\mathsf{d}_{\ast}(c)\le 2c_1^{b}\,\mathsf{d}_{\ast}(y).
\end{equation}
This follows from the temperedness of
$\mathsf{d}_{\ast}\in\mathfrak{M}^{\wedge}_{\delta/2}(\mathfrak{B}_j)$ at rate $\tfrac12\delta$
together with (A4). This gives the second bound of
\eqref{5.7:eq-flat-energy-bound-b}.

\emph{The third bound.} It follows from the absolute form, clause (f) of
Lemma~\ref{5.7:lem-two-map-interpolation}, applied to the frame potential $A_W$ of the third
bound. Alternatively it is (A5).
\end{proof}

Proposition~\ref{5.7:prop-critical-layer-closeness} supplies the hypothesis
\eqref{5.7:eq-flat-energy-bound-b} of the flat energy-difference bound
(Lemma~\ref{5.7:lem-flat-energy-bound}) for pairs of sources in $\mathcal{K}^{\mathrm{crit}}$.
Hence, by Lemma~\ref{5.7:lem-flat-energy-bound}, the flat energy-difference bound holds for such
pairs, with $\varpi_{\mathbb{H}}$ and $\varpi_W$ given by \eqref{5.7:eq-critical-layer-closeness-1},
at the radii fixed above.

Let $p$ be a pair that is a chain of the dilations $\Phi_t$, with
$\chi_t=\sum_{\square}t_{\square}\zeta_{\square}\not\equiv 1$. We call such a pair a
\emph{rippled parent}. At a rippled parent the base is not critical, and the layers are critical
points in the presence of a source term (the sourced critical points) of the relative layer theory.

At real pairs presented in the format of the relative layer theory, with $\mathsf{s}=\mathbb{1}$,
the relative layer bound gives
\begin{equation}\label{5.7:eq-rippled-layer-closeness-1}
\bigl|Y[\mathbb{H}';\mathbf{U}']-\mathbb{H}\bigr|_x
\le C_R'\bigl\{\mathsf{P}_b(x)\Psi_p(x)+\mathsf{P}_{\delta b}(x)\bigr\},
\qquad
\Psi_p=\mathsf{T}_p^{\sim}+\mathfrak{r}_0L^{-\theta\lambda},
\end{equation}
with $\mathsf{P}_b(x)=\sum_y e^{-\delta_1 d(x,y)}|b(y)|$, and with $\mathfrak{r}_0L^{-\theta\lambda}$
the member of $\Psi_p$ that decays in the level function $\lambda$ at the exponent $\theta$. The
exchange of the base
$\mathbb{U}$ for $\mathbb{W}$ is supplied by the hybrid window lemma. The bound
\eqref{5.7:eq-rippled-layer-closeness-1} is linear in the inherited discrepancy
$\mathsf{T}_p$, as the recursion requires. Its rate, however, is too small.

The smear $\mathsf{P}_b$ is taken at the rate
$\delta_1=\delta_0/32=\tfrac14\delta$. By the summation bound for the scaled distance, which
sums at half the rate, it returns on $\square^{\zeta}$ a quantity of order
$\sup|b|\,e^{-(\delta/4)d^{-}(\square)/2}$, where $d^{-}(\square)$ is the scaled distance of the
block $\square$ that enters the disc radius. The interpolation disc, on the other hand, allows
$|\mathsf{t}|$ up to $1/(14K^{\mathrm{rem}}\vartheta(\square))$, where $K^{\mathrm{rem}}$ is the
constant in the radius of these discs; this is a
quantity proportional to
$e^{\frac12\delta d^{-}(\square)}$. Since $\tfrac12\delta-\tfrac18\delta=\tfrac38\delta$, the
product of $|\mathsf{t}|$ with the bound can grow like $e^{\frac38\delta d^{-}(\square)}$. This
is the rate obstruction of the relative layer theory, and it has two exits:
\begin{itemize}
\item[(i)] A lower bound on $M_1$. This is valid if Ba{\l}aban's scaled distance counts
cubes of side $M_1$. It is not available here, since the scaled distance used in the present
frames is $d^{\wedge}$ in each frame's own units.
\item[(ii)] For pair differences only, a smaller circle $|\mathsf{t}|\le r^{\bullet}_{\square}$.
\end{itemize}
On the smaller circle the relative layer bound suffices, as the following remark shows.

\begin{remark}[Pair differences on the smaller circle]
Let $\sigma\equiv\mathbb{1}$, that is, let $\sigma$ be identically the trivial element
$\mathbb{1}$, let the parameters be real, and let $|\mathsf{t}|\le
r^{\bullet}_{\square}$. Assume the following three statements:
\begin{itemize}
\item[(a)] the relative layer bound \eqref{5.7:eq-rippled-layer-closeness-1} on the frame
$\mathfrak{B}$;
\item[(b)] the reduction of the finer averaged datum to one lattice,
Lemma~\ref{5.7:lem-datum-reduction};
\item[(c)] the relative layer bound \eqref{5.7:eq-rippled-layer-closeness-1} on the frame
$\mathfrak{B}_j$, with $(b,\delta b)$ replaced by $(\mathsf{B},\delta\mathsf{B})$.
\end{itemize}
Put
\begin{equation*}
\mathsf{h}:=t\,\mathsf{P}_b+\mathsf{H},
\qquad
\mathsf{H}:=|\mathsf{t}|\,\mathsf{P}_b\mathbb{1}_{\square^{\zeta}},
\end{equation*}
and let $\mathsf{h}^{\sim}$ be the smear of $\mathsf{h}$ at rate $\delta_1$. Then the
energy-difference bound holds with the profile $\mathsf{h}$, linearly in $\Psi_p$:
\begin{equation}\label{5.7:eq-rippled-layer-closeness-2}
|\delta\mathsf{E}(c)|\le C\,\Psi_p(c)\,\mathsf{h}^{\sim}(c).
\end{equation}
Moreover, with $F_{\square}$ the set of plaquettes $c$ over which the energy-difference datum at
the block $\square$ is summed,
\begin{equation}\label{5.7:eq-rippled-layer-closeness-3}
\sum_{c\in F_{\square}}e^{-4\delta_1 d_{\star}(x,c)}|\delta\mathsf{E}(c)|
\le C L c_1^{b}\,\Psi_p(x)\,\mathsf{h}^{\sim}(x)\,e^{-\frac32\delta_1 D_v}.
\end{equation}
\end{remark}

\begin{proof}
The datum $\delta\mathsf{E}(c)$ reads three layers:
\begin{itemize}
\item the layers $t\mathbb{H}^{w}$ on $\mathfrak{B}$;
\item the datum $\mathsf{B}(t,\mathsf{t})$ on $\mathfrak{B}_j$;
\item the layer $\mathbf{H}_j$ on $\mathfrak{B}_j$, read at that datum.
\end{itemize}
We bound the three differences in turn.

The two-construction difference of the layers $t\mathbb{H}^{w}$ is bounded by hypothesis (a). Its
right member is $C_R'\mathsf{P}_b\Psi_p$, which is linear in $\Psi_p$, together with the member
$C_R'\mathsf{P}_{\delta b}$ in the difference $\delta b$ of the two constructions' variables. The
steps below bound the members linear in $\Psi_p$, which make up the right member of
\eqref{5.7:eq-rippled-layer-closeness-2}.

The datum $\mathsf{B}(t,\mathsf{t})$ and its mismatch are handled by hypothesis (b), applied on
$\mathfrak{B}_j$ with the layer $(t+\mathsf{t}\zeta_{\square})\mathbb{H}^{w}$. That lemma writes
the finer construction's datum as the coarser construction's functional $Q_{\star}$, read at the
averaged layer over the averaged base. The mismatch $\delta\mathsf{B}$ is therefore bounded by
\eqref{5.7:eq-datum-reduction-a} in terms of the two-construction difference of that layer, which
hypothesis (a) bounds. The member of the layer in $\mathsf{t}$ is carried on $\square^{\zeta}$;
it is the origin of the member $\mathsf{H}$ of the profile $\mathsf{h}$. Hence the mismatch too
is linear in $\Psi_p$.

The layer $\mathbf{H}_j$ is bounded by hypothesis (c), with $(b,\delta b)$ replaced by
$(\mathsf{B},\delta\mathsf{B})$. Its right member consists of a member linear in $\Psi_p$ and
the smear of $\delta\mathsf{B}$. By the preceding step, the latter is again linear in $\Psi_p$.

The profile carried by these members is that of the layers, $t\mathsf{P}_b$, together with the
interpolation member on $\square^{\zeta}$, $\mathsf{H}$. On the circle
$|\mathsf{t}|\le r^{\bullet}_{\square}$ the member $\mathsf{H}$ is bounded by the constant
$C_{\flat}(c_M\mathsf{b}_j/\vartheta_P^{0})$ of the smaller circle, by exit (ii) of the rate
obstruction.
The datum $D^{\sharp}$ is a convex combination of these layers with scaffold weights, and
$\mathsf{E}$ is formed from them by the Baker--Campbell--Hausdorff formula at small arguments,
bilinearly, as in the proof of the flat energy-difference bound
(Lemma~\ref{5.7:lem-flat-energy-bound}). Hence the three bounds combine to
\eqref{5.7:eq-rippled-layer-closeness-2}, with the members in $\Psi_p$ entering linearly.

For \eqref{5.7:eq-rippled-layer-closeness-3}, insert \eqref{5.7:eq-rippled-layer-closeness-2}
into the sum. The quantities $\Psi_p$ and $\mathsf{h}^{\sim}$ are tempered at rate $\delta_1$
each. For $\Psi_p$ this reads
\begin{equation*}
\Psi_p(c)\le L\,e^{\delta_1 d(x,c)}\Psi_p(x),
\end{equation*}
which produces the factor $L$. Replacing $\Psi_p(c)\mathsf{h}^{\sim}(c)$ by
$\Psi_p(x)\mathsf{h}^{\sim}(x)$ thus costs the rate $2\delta_1$ out of the kernel's
$4\delta_1$. The summation bound for the scaled distance takes a further
$\tfrac12\delta_1$ and contributes the constant $c_1^{b}$. The remaining rate is
\begin{equation*}
4\delta_1-\delta_1-\delta_1-\tfrac12\delta_1=\tfrac32\delta_1,
\end{equation*}
and it yields the factor $e^{-\frac32\delta_1 D_v}$ of
\eqref{5.7:eq-rippled-layer-closeness-3}. Since
$\delta_1=\frac14\delta$, we have
\begin{equation*}
e^{-\frac32\delta_1 D_v}=e^{-\frac38\delta D_v}.
\end{equation*}
This factor decays in $D_v$ at a rate exceeding $\frac14\delta$.
\end{proof}

On the full interpolation discs $|\mathsf{t}|\le r^{\mathrm{rem}}_{\square}$ of the correlation
theorem the remark is not
available. An interpolation of \eqref{5.7:eq-rippled-layer-closeness-1} against clause (b) of
the layer theorem would put a fractional power on $\mathsf{T}_p$. Such a power is admissible
only in a bound on which no further step of the recursion depends; in a bound that feeds the
recursion it is not, since a recursion of the form
$\mathsf{T}\mapsto C\mathsf{T}^{1/3}$ does not close. What is required is the following.

\begin{proposition}[Linear layer closeness at rippled parents; proof outstanding]
\label{5.7:prop-rippled-layer-closeness}
Let $p$ be a rippled parent, that is, a pair which is a chain of the dilations $\Phi_t$ with
$\chi_t=\sum_{\square}t_{\square}\zeta_{\square}\not\equiv 1$. On the parameter
domains of hypothesis \eqref{5.7:eq-flat-energy-bound-b}, the interpolation parameter ranging
over the full discs $|\mathsf{t}|\le r^{\mathrm{rem}}_{\square}$, put
\begin{equation*}
\mathsf{p}:=\sup|b|\,e^{-\frac52\delta_1 d^{\wedge}(\cdot,\mathcal{C}_{\mathsf{j}})},
\end{equation*}
where $\mathcal{C}_{\mathsf{j}}$ is the set of cells on which the data $b$ are given;
this is a profile decaying at the rate $\frac52\delta_1=\frac58\delta$ of the profile $p^{\flat}$ of
\eqref{5.7:eq-flat-profile-a}. Then, at $p$, the first inequality of
\eqref{5.7:eq-flat-energy-bound-b} holds at the rate $\frac58\delta$ of $p^{\flat}$, with a
right member linear in $\mathsf{T}_p$ through $\Psi_p$: in the bound
\eqref{5.7:eq-rippled-layer-closeness-1}, the member $C_R'\mathsf{P}_b(x)\Psi_p(x)$, which
carries the profile, is replaced by $C\,c_1^{b}\,\mathsf{p}(x)\Psi_p(x)$.
\end{proposition}

Proof outstanding.

\medskip\noindent\emph{Route.}\
The bound is to be obtained by reading the proof of the relative layer bound
\eqref{5.7:eq-rippled-layer-closeness-1} at a general profile. In that proof the profile
$\mathsf{P}$ enters through the regularity corollary read in the finer construction, namely, for
the quantities $X'$ and $\mu'$ of that corollary and its constant $C_s$,
\begin{equation*}
X'\in C^{1,\beta}(C_s\mathsf{P}),\qquad |\mu'|\le C_s\mathsf{P}.
\end{equation*}
It leaves through the stability estimate, whose kernel is
$e^{-4\delta_1 d}=e^{-\delta d}$.

The proof splits the difference into six residuals. Four of them are local in the profile. The
profile is smeared only in two: the Lipschitz residual of the scalar operator $R$, and the
residual collecting the $\varsigma$-lemma in the finer construction, the nesting bound for $R$ and
the divergence identity. There the smearing is through the scalar operators $R$, $G'$, $P$. It is
taken at $\delta_1$ by statement, namely through the uniform smear of the weight calculus, while
the kernel estimates \cite[(3.42)--(3.49)]{B-CMP99b} give the kernels' own rate,
$e^{-8\delta_1 d}$. The weight calculus states that all these kernels decay at rate at least
$4\delta_1$ in $d$, and that the calculus uses, of the weight, only its slowness.

If the two smears are read at $4\delta_1$, every step is to hold for any profile $\mathsf{p}$
that is tempered at $\frac52\delta_1$ and majorises the $C^{1,\beta}$-sizes of $\mathbb{H}'$ and
$\mathbb{H}$. The rate budget then closes exactly. Consider a residual with density
$|\mathfrak{a}|_y$ in $y$. Since $\mathsf{p}$ is tempered at $\frac52\delta_1$, and
$4\delta_1-\frac52\delta_1=\frac32\delta_1$, the density of the residual satisfies
\begin{equation*}
C\sum_y e^{-4\delta_1 d(x,y)}|\mathfrak{a}|_y\,\mathsf{p}(y)
\le C\,\mathsf{p}(x)\sum_y e^{-\frac32\delta_1 d(x,y)}|\mathfrak{a}|_y .
\end{equation*}
The stability estimate smears once more, at $\frac32\delta_1$. The triangle inequality gives
\begin{equation*}
e^{-\frac32\delta_1[d(x_0,x)+d(x,y)]}
\le e^{-\delta_1 d(x_0,y)}\,e^{-\frac12\delta_1 d(x_0,x)},
\end{equation*}
and the summation bound for the scaled distance, which sums at $\frac12\delta_1$, then yields
\begin{equation}\label{5.7:eq-rippled-layer-closeness-4}
\sum_x e^{-\frac32\delta_1[d(x_0,x)+d(x,y)]}\le c_1^{b}\,e^{-\delta_1 d(x_0,y)}.
\end{equation}
The budget is therefore
\begin{equation*}
4\delta_1=\tfrac52\delta_1+\delta_1+\tfrac12\delta_1 .
\end{equation*}
The first term is taken by the profile, the second by the smear of $\mathsf{T}_p^{\sim}$ in
$\Psi_p$, and the third by the summation in \eqref{5.7:eq-rippled-layer-closeness-4}. The
stability estimate then returns $C c_1^{b}\,\mathsf{p}(x_0)\Psi_p(x_0)$.

For the profile $\mathsf{p}$ of the statement, which decays at the rate
$\frac52\delta_1=\frac58\delta$ of $p^{\flat}$, regularity
at this weight is clause (b) of the layer theorem, which holds for every profile in
$\mathfrak{M}_{\delta}$. The returned bound is then the first inequality of
\eqref{5.7:eq-flat-energy-bound-b} at the rate of $p^{\flat}$, linear in $\mathsf{T}_p$.

The same profile, inserted in the companion clause of the relative layer bound, makes the
member $\mathsf{H}=\bar{\tau}\,\mathsf{p}$ of that clause, where $\bar{\tau}$ is the factor
through which the interpolation parameter $\mathsf{t}$ enters that member, bounded
on the full interpolation discs. This follows from the inequality which
Lemma~\ref{5.7:lem-flat-profile} shows to hold for the profile $p^{\flat}$.
The real recursion along the chains of dilations needs this as well.

What has to be established is therefore that the following five arguments hold verbatim at a
general weight tempered at $\frac52\delta_1$, with the two smears read at $4\delta_1$:
\begin{itemize}
\item the Lipschitz bound for $R$;
\item the nesting bound;
\item the $\varsigma$-lemma;
\item the penalty bound;
\item clause (c) of the statement on the response function $\mathcal{H}$.
\end{itemize}
The decay statement of the weight calculus is itself an input taken by statement. Moreover, the
slowness used in the relative layer theory,
\begin{equation*}
\Psi(y)\le L\,e^{\delta_1 d(x,y)}\Psi(x),
\end{equation*}
differs from the slowness of the weight calculus, which allows the factor
$e^{\delta_1 d(x,y)/4}$. The fifth residual, though local in the profile, smears $\Psi_p$; it
therefore has to be checked in the same way.

\begin{proposition}[Layer closeness at real rippled parents with a fractional power]
\label{5.7:prop-fractional-layer-closeness}
Let $p$ be a real presented pair of the class of rippled parents, the pairs reached by chains of
dilations, whose discrepancy profile $\mathsf{T}_p$ obeys $\sup\mathsf{T}_p\le\varepsilon_D$, where
$\varepsilon_D$ is the constant of the paired frame
(Definition~\ref{5.7:def-data-mismatch}), and with $\sigma$ the trivial
configuration, $\sigma\equiv\mathbb{1}$. Let $\Psi_p$ be the inherited discrepancy of the pair,
formed from the smear $\mathsf{T}_p^{\sim}$ of $\mathsf{T}_p$, let $g$ be the gauge element of the
pair, which obeys $|g-1|\le C\Psi_p$, and let $\delta_1=\tfrac14\delta$. For a point $y$ write
\begin{equation*}
\Phi^{\mathsf{A}}_y(b):=\ell\,\mathbb{H}(b;P)(y),\qquad
\Phi^{\mathsf{B}}_y(b):=\ell\,Y[\mathbb{H}'(\operatorname{Ad}_g b\oplus 0;P');\mathbf{U}'](y),
\end{equation*}
for the layers of the two machines (the two cutoffs) at $p$, read at $y$, where $\ell$ is the
length-scale prefactor of the layer bounds and $P$, $P'$ are the arguments at which the pair forms
the layers of the machines $\mathsf{A}$ and $\mathsf{B}$; and write
$\delta\Phi_y:=\Phi^{\mathsf{B}}_y-\Phi^{\mathsf{A}}_y$. Assume:
\begin{enumerate}
\item[(i)] the relative layer bound for this class, read in sup form: for real $b$ with
$\sup_{c'}|b(c')|<a_{R'}$, where $a_{R'}$ is the radius of this bound,
\begin{equation*}
|\delta\Phi_y(b)|\le C_{R'}\{\mathsf{P}_b\Psi_p+\mathsf{P}_{\delta b}\}(y),
\qquad \delta b:=(\operatorname{Ad}_g-1)b,
\end{equation*}
where $C_{R'}$ is the constant of this bound and $\mathsf{P}_b$ is the smear of $b$ at rate
$\delta_1$; and the same bound on the frame $\mathfrak{B}_j$ for the layer $\mathbf{H}_j$ of the
two machines as a function of a real datum $\mathsf{d}$ with $\sup_{c}|\mathsf{d}(c)|<a_{R'}$,
with $(b,\delta b)$ replaced by $(\mathsf{d},\delta\mathsf{d})$;
\item[(ii)] the regularity of the layers:
complex data are allowed, the layers depend analytically on them, and they obey the decay
bound with kernel $e^{-4\delta_1 d(y,c')}=e^{-\delta d(y,c')}$ in the scaled distance $d(y,c')$
of the frame;
\item[(iii)] the bound for the frame potential $A_W$ of the passage from the base $\mathbb{U}$
to the base $\mathbb{W}$ at $p$:
$\ell\sup_{\square_0}|A_W|\le C\sup_{\square_0}\Psi_p$;
\item[(iv)] the profile bound of which Lemma~\ref{5.7:lem-flat-profile} is the flat form, under its
smallness condition, by which the datum $\mathsf{d}=\mathsf{B}(t,\mathsf{t})$ is
small on the domains below, with smallness measured by the parameter $\gamma$ of that condition.
\end{enumerate}
Then on the domains of the last stage's own parameters, namely complex $b$ with
$\|b\|_\infty\le 1.1\,C_1\delta'_{\mathsf{j}}$, the radius fixed for $b$ at this stage,
$t\in[0,1]$, and $|\mathsf{t}|\le r^{\mathrm{rem}}_{\square}$, the remainder radius of the block
$\square$, the inequality (B1) of
\eqref{5.7:eq-flat-energy-bound-b} holds with
$\varpi_{\mathbb{H}}(y):=C\Psi_p(y)^{1/16}$, the inequality (B2) of
\eqref{5.7:eq-flat-energy-bound-b} holds with $\varpi_{\mathbb{H}}$ replaced by
$C\Psi_p(y)^{1/8}$ and $\varpi_{\mathsf{d}}$ unchanged, and the inequality (B3) of
\eqref{5.7:eq-flat-energy-bound-b} holds with
$\varpi_W:=C\sup_{\square_0}\Psi_p$.
\end{proposition}

\begin{proof}
Fix $y$. We apply part~(c) of the interpolation between two analytic maps
(Lemma~\ref{5.7:lem-two-map-interpolation}) to the pair $\Phi^{\mathsf{A}}_y$,
$\Phi^{\mathsf{B}}_y$, with base $B:=0$ and radius $a:=a_{R'}$, the radius of the
relative layer bound fixed in hypothesis~(i), the condition
\eqref{5.7:eq-data-mismatch-b} of the paired frame (Definition~\ref{5.7:def-data-mismatch}) being
read with $a_{R'}$ in place of its radius. At $B=0$ the condition of part~(c),
$\sup|B|+\sup|b|\le a/8$, reads $\sup_{c'}|b(c')|\le a_{R'}/8$, and for (B2)
$\sup_{c}|\mathsf{d}(c)|\le a_{R'}/8$; it is this condition that
\eqref{5.7:eq-data-mismatch-b}, read with $a_{R'}$, supplies on the domains of the statement.
For (B1) the variables are the values $b(c')$,
$c'\in\mathcal{C}_{\mathsf{j}}$, the cells on which $b$ is defined; for (B2) they are the values
$\mathsf{d}(c)$, $c\in\mathfrak{B}_j$, of a datum on the frame $\mathfrak{B}_j$, and the two maps
are the layers $\mathbf{H}_j$ of the two machines, read at $y$. For (B1) the variable is
common to the two maps, the second reading it through $\operatorname{Ad}_g$.

We check the hypotheses \eqref{5.7:eq-two-map-interpolation-a} of that lemma. The first two
ask for analyticity of the two maps on the polydisc
$\mathfrak{D}(a)$ in complex data and for the kernel bound; both are hypothesis~(ii), which allows
complex data, gives analytic dependence on them, and supplies the decay bound with kernel
$e^{-4\delta_1 d(y,c')}=e^{-\delta d(y,c')}$ in the scaled distance $d(y,c')$. The third is the
real-slice closeness of the two
maps. By hypothesis~(i), on the real slice of $\mathfrak{D}(a)$,
\begin{equation*}
|\delta\Phi_y(b)|\le C_{R'}\{\mathsf{P}_b\Psi_p+\mathsf{P}_{\delta b}\}(y)\le C c_1 a\,\Psi_p(y),
\end{equation*}
where the second inequality uses $\delta b=(\operatorname{Ad}_g-1)b$ and $|g-1|\le C\Psi_p$. For
(B2) the second form of hypothesis~(i) is read in the same way, with
$(\mathsf{d},\delta\mathsf{d})$ in place of $(b,\delta b)$. Thus
the third hypothesis holds at $y$ with a closeness constant $\varepsilon_{\star}$ bounded by a
constant multiple of $\Psi_p(y)$. No decay in the distance is asked of this bound.

The parameter $\mathsf{t}$ is not one of the variables of the application: it enters only through
the datum $\mathsf{d}=\mathsf{B}(t,\mathsf{t})$, which is small on the domains of the
statement, with smallness measured by $\gamma$, by hypothesis~(iv).

Part~(c) of Lemma~\ref{5.7:lem-two-map-interpolation} now gives, for $0\le\theta\le1$,
\begin{equation}\label{5.7:eq-fractional-layer-closeness-1}
|\delta\Phi_y(b)-\delta\Phi_y(0)|\le 28\,\varepsilon_{\star}^{\theta/2}\,T^{(1-\theta)}[b].
\end{equation}
We take $\theta:=\tfrac18$ for (B1) and $\theta:=\tfrac14$ for (B2). For (B1) the kernel is
$k(c')=K_R e^{-\delta d(y,c')}$, so that $T^{(7/8)}[b]$ is at most a constant multiple of
\begin{equation*}
\sum_{c'\in\mathcal{C}_{\mathsf{j}}}e^{-\frac78\delta d(y,c')}|b(c')|
\le C\sup|b|\,e^{-\frac58\delta d^{\wedge}(y,\mathcal{C}_{\mathsf{j}})},
\end{equation*}
the summation over $c'$ costing $\tfrac14\delta$ of the rate; the right member decays at the rate
$\tfrac58\delta$ of the flat profile \eqref{5.7:eq-flat-profile-a}. For (B2) the kernel on the
frame $\mathfrak{B}_j$ is $K_R e^{-\delta d_j(y,c)}$, and for data bounded by
$\mathsf{d}_{\ast}$ the sum $T^{(3/4)}$ is at most a constant multiple of
\begin{equation*}
\sum_{c\in\mathfrak{B}_j}e^{-\frac34\delta d_j(y,c)}\,\mathsf{d}_{\ast}(c)
\le 2c_1^{b}\,\mathsf{d}_{\ast}(y),
\end{equation*}
since $\mathsf{d}_{\ast}\in\mathfrak{M}^{\wedge}_{\delta/2}(\mathfrak{B}_j)$ is tempered at rate
$\tfrac12\delta$ and the summation again costs $\tfrac14\delta$.
Since $\varepsilon_{\star}\le C\Psi_p(y)$, the factor
$\varepsilon_{\star}^{\theta/2}$ in \eqref{5.7:eq-fractional-layer-closeness-1} is at most
$C\Psi_p(y)^{1/16}$ for $\theta=\tfrac18$ and at most $C\Psi_p(y)^{1/8}$ for $\theta=\tfrac14$. This
is (B1) with $\varpi_{\mathbb{H}}(y)=C\Psi_p(y)^{1/16}$ and (B2) with $C\Psi_p(y)^{1/8}$ in the place
of $\varpi_{\mathbb{H}}$. In (B2) only this factor changes: the member
$2B_{\sharp}\varpi_{\mathsf{d}}\mathsf{d}_{\ast}$ accounts for the difference of the two data; it is
the member $T^{(1)}[\mathsf{d}^{\mathsf{B}}-\mathsf{d}^{\mathsf{A}}]$ of part~(e) of
Lemma~\ref{5.7:lem-two-map-interpolation}, which does not contain $\varepsilon_{\star}$, and it is
therefore unchanged.

Finally, (B3) with $\varpi_W=C\sup_{\square_0}\Psi_p$ is hypothesis~(iii).
\end{proof}

\begin{remark}
At real rippled parents, that is, at real presented pairs as in
Proposition~\ref{5.7:prop-fractional-layer-closeness}, the flat energy-difference bound
\eqref{5.7:eq-flat-energy-bound-a} therefore holds with $\varpi'=C\Psi_p^{1/16}$. This is
Lemma~\ref{5.7:lem-flat-energy-bound} applied to (B1)--(B3) in the form just obtained. The data
mismatch of the paired frame (Definition~\ref{5.7:def-data-mismatch}) obeys the same bound,
through the reduction \eqref{5.7:eq-datum-reduction-a} at the layer $\mathbf{A}$ together with
Lemma~\ref{5.7:lem-flat-profile}. Theorem~\ref{5.7:thm-remainder-value} then gives, with $i$ the
level index of that theorem,
\begin{equation*}
|\delta G_v|\le C\,\mathfrak{n}_v\bigl(\Psi_p^{1/16}+c_1\vartheta^{i}\bigr)
e^{-\frac14\delta D_v}\,\alpha_{\cdot,i}.
\end{equation*}
A fractional power of $\Psi_p$ is harmless where the bound is used terminally, for a value that is
read once and not fed into a further move. It is not admissible where the
pair produced from $p$ by a move is moved again. The chains of moves producing pairs run through
all levels, and the recursion of data mismatches
$\mathsf{m}^{\mathrm{dat}}\mapsto C(\mathsf{m}^{\mathrm{dat}})^{1/16}$ from one level to the next
has no small fixed point. There
a bound linear in $\Psi_p$
is needed, which Proposition~\ref{5.7:prop-fractional-layer-closeness} does not provide.
\end{remark}

\begin{remark}[No common object off the vertices of the decoupling parameters. Proof outstanding]
\label{5.7:rem-decoupling-vertices}
Take both machines $\mathsf{A}$ and $\mathsf{B}$ with the decoupling parameters of Ba{\l}aban's own
random-walk expansions. Write $\sigma(\Delta)$ for the decoupling parameter attached to a box
$\Delta$; its vertex values are the endpoints of the real interval over which Ba{\l}aban's
random-walk expansion interpolates $\sigma(\Delta)$, and $\sigma$ is at a vertex when every
$\sigma(\Delta)$ takes one of them. Let $\sigma$ range over the closed polydisc
\begin{equation*}
  \{\sigma : |\sigma(\Delta)| \le e^{\kappa_1} \text{ for every box } \Delta\} .
\end{equation*}
Write $\mathsf{R}^{\mathsf{m}}(\mathsf{t},\sigma)$, $\mathsf{m}\in\{\mathsf{A},\mathsf{B}\}$, for the
random-walk expansion of machine $\mathsf{m}$ at interpolation parameter $\mathsf{t}$ and decoupling
parameters $\sigma$.

On this polydisc, establishing the bounds \eqref{5.7:eq-flat-energy-bound-b}, which together form the
hypothesis of the flat energy-difference bound (Lemma~\ref{5.7:lem-flat-energy-bound}), requires
comparing $\mathsf{R}^{\mathsf{B}}(\mathsf{t},\sigma)$ with $\mathsf{R}^{\mathsf{A}}(\mathsf{t},\sigma)$
on the random-walk layers.
These are walk expansions on two different lattices, and they have no common walk space.

The comparison in the uniqueness theorem is carried out by natural interpolations through the
vertex values of the decoupling parameters. In that comparison the passage from the trivial smearing
profile $\mathsf{s}=\mathbb{1}$ to a profile $\mathsf{s}\neq\mathbb{1}$ is not a variational step, and
nothing is claimed for it.

The two-cutoff response statement at zeroth order (Proposition~\ref{5.7:prop-response-zeroth-order})
is itself free of $\sigma$, as its statement shows. The configurations from which its responses are
computed are, box by box, minimisers whose only input is the data. The decoupling parameters therefore
enter the bounds \eqref{5.7:eq-response-zeroth-order-a} only through the data, that is, through the
mismatches $\varpi$ and $\varpi'$.

The statement whose proof is outstanding concerns the decoupling parameters off their vertex values:
the bounds \eqref{5.7:eq-flat-energy-bound-b} hold for every $\sigma$ in the closed polydisc above,
and not only at those $\sigma$ whose values $\sigma(\Delta)$ are all vertex values.
Proof outstanding.

\emph{Route.} No route is given. The missing ingredient is a space of walks common to the two
random-walk expansions, on which $\mathsf{R}^{\mathsf{B}}(\mathsf{t},\sigma)$ and
$\mathsf{R}^{\mathsf{A}}(\mathsf{t},\sigma)$ could be compared term by term, and none is available.
\end{remark}

\begin{definition}[The comb dilation from the trivial configuration]\label{5.7:def-comb-dilation}
Let $\mathsf{A}$ and $\mathsf{B}$ be the two cutoffs, on the lattices $T_{\eta}$ and $T_{\eta'}$ respectively, with
$\eta'=\eta/L$, and let $\mathsf{m}\in\{\mathsf{A},\mathsf{B}\}$ denote either of them. Fix a block $\square$, a
level $j$, an anchor $y\in T_{\ell_j}\cap\tilde{\square}^{2}$, and a near companion $\hat A=A(h_y,\cdot)$ of
level $j$. Here $\ell_j$ is the length scale of level $j$, with $\ell_{j+1}=L\ell_j$, and $T_{\ell_j}$ is the
lattice of spacing $\ell_j$. The function $h_y$ is the tent function of \cite{B-CMP119} with anchor $y$, built from
$h(t)=\max\{1-|t|,0\}$ in each coordinate, in units of $\ell_j$. Put
\begin{equation*}
n:=\min\{l(y),n_0\},\qquad \nu:=1+n-j,\qquad t:=L^{\,j-n},\qquad u_j:=(\eta/\ell_j)^2 ,
\end{equation*}
where $l(y)$ denotes the largest level $i$ with $y\in T_{\ell_i}$, and $n_0$ is a fixed level capping $n$,
so that $n\le n_0$.
Let $w(\Delta)$ be the block weight of the $n$-cell $\Delta$, with the constant $c_W$, as fixed with the block
weights, let $\operatorname{dist}_n(\Delta,\cdot)$ be the sup-distance from the $n$-cell $\Delta$ to a point or a
cell, in units of $\ell_n$, the distance with which the block weights are defined, and put
\begin{equation*}
w^{\vee}(y):=\max\{w(\Delta):\operatorname{dist}_n(\Delta,y)<1\},\qquad
w^{\vee}(\Delta'):=\max\{w(\Delta):\operatorname{dist}_n(\Delta,\Delta')<1\}.
\end{equation*}
Let $\Delta'_y$ be the unit $n$-cell with lower corner $y$ (a point of $T_{\ell_n}$, since $n\le l(y)$).

\begin{enumerate}
\item \emph{Support cube and natural sizes.} Define
\begin{equation}\label{5.7:eq-comb-dilation-1}
\begin{aligned}
C_y^{\sharp}&:=\{z\in T_{\eta}:\ y_\kappa-\ell_j+\eta\le z_\kappa\le y_\kappa+\ell_j
\ \text{for every }\kappa\},\\
w^{\sharp}(y)&:=\max\{w(\Delta):\Delta\cap C_y^{\sharp}\neq\emptyset\},\qquad
\Theta^{\sharp}:=c_W\,w^{\sharp}(y),\\
\mathsf{f}_{\ast}&:=\Theta^{\sharp}\,\mathfrak{a}_n\,t^{2},\qquad
\mathsf{b}_{\ast}(b):=2d\,(1+|b-y|)\,\mathsf{f}_{\ast},\\
a^{\mathrm{ax}}&:=\min\{\alpha_3(\circ),1/16\},\qquad
S^{\mathrm{ax}}_v:=\frac{a^{\mathrm{ax}}}{4d\,\mathsf{f}_{\ast}} .
\end{aligned}
\end{equation}
Here $|b-y|$ is the sup-distance from the bond $b$ to $y$ in units of $\ell_j$, and $\alpha_3(\circ)$ is the
small-field radius $\alpha_3$ read at the slot mark $\circ$ of the target slot below, and the subscript $v$
denotes the unit carrying the companion $\hat A$. The dilation radius
$S^{\mathrm{ax}}_v$ is a number in the sense of the first of the three numerical statements accompanying the
analysis of Ba{\l}aban's step.
The constant $\mathfrak{a}_n$ in $\mathsf{f}_{\ast}$ is the level-$n$ coefficient of the natural size data of
companions, fixed with those data; it is a constant distinct from the $\delta$-size $\mathsf{a}$ of the move below.
The cube $C_y^{\sharp}$ contains every corner of every plaquette $P$ with $h_y(x_P)\neq0$, where
$x_P\in T_{\eta}$ is the lattice point labelling $P$, its lower corner (not the barycentre $x(p)$ of a
plaquette). The number
$\mathsf{f}_{\ast}$ is the natural curvature size $|F|_{\ast}$ at $\Theta^{\sharp}$, and $\mathsf{b}_{\ast}(b)$ is
the natural size $|B^{\mathrm{ax}}|_{\ast}(b)$ of the axial potential at $\Theta^{\sharp}$. Moreover
\begin{equation}\label{5.7:eq-comb-dilation-2}
\sup_{C_y^{\sharp}}\mathsf{b}_{\ast}=4d\,\mathsf{f}_{\ast},
\qquad
w^{\vee}(y)\le w^{\sharp}(y)\le w^{\vee}(\Delta'_y)\quad(\nu\in\{1,2\}).
\end{equation}

\item \emph{The comb.} Let $W$ be the link-variable component $U_{\mathsf{P}^0}$ of the box field
$\mathfrak{b}'_j$ at level $j$, and $W^{\mathsf{m}}$ its version for cutoff $\mathsf{m}$.

For the cutoff $\mathsf{m}$, the cube $C_y^{\sharp}$ is read in the comb and in the move below on
$\mathsf{m}$'s own lattice: it is the set of points of $T_{\eta}$ for $\mathsf{A}$, respectively of
$T_{\eta'}$ for $\mathsf{B}$, lying in the coordinate box of \eqref{5.7:eq-comb-dilation-1}. The contours,
bonds and loops below are taken in that lattice.

For $z\in C_y^{\sharp}$, let $\gamma_z$ be the lexicographic contour from $y$ to $z$ in $C_y^{\sharp}$, along
which each coordinate is monotone.

For a bond $b$ from $b_-$ to $b_+$, the comb loop based at $y$ is
$\lambda_b:=\gamma_{b_-}\,b\,\gamma_{b_+}^{-1}$.

\item \emph{The move.} Consider every block $\square$ and every near pointed family or counterterm
companion of level $j$ anchored at $y\in\tilde{\square}^{2}$, for every value of $\nu$. Its parent datum is
$(k;\mathsf{n}+1,\hat X)$ with $\hat X\supseteq\hat{\mathsf{K}}_{\square}$, as for the second listed move;
here $\mathsf{n}+1$ is the index of the parent slot and $\hat{\mathsf{K}}_{\square}$ the region attached to the
block $\square$ in that move. Set
\begin{equation}\label{5.7:eq-comb-dilation-3}
\mathfrak{P}_{\mathrm{ax}}:=(b,\sigma)\times[0,1]\times\mathbb{D}_{\square}
\times\{\varsigma:\ |\varsigma|<S^{\mathrm{ax}}_v\},
\end{equation}
where the first factor stands for the set of index pairs $(b,\sigma)$ of the listed-move format, taken over
unchanged; and,
for every bond $b'$ of the lattice of cutoff $\mathsf{m}$ with both endpoints in $C_y^{\sharp}$, so that the
comb loop $\lambda_{b'}$ is defined,
\begin{equation}\label{5.7:eq-comb-dilation-4}
U^{\mathsf{m}}_{\varsigma}(b'):=\exp\bigl(\varsigma\,\log W^{\mathsf{m}}(\lambda_{b'})\bigr),
\end{equation}
where $W^{\mathsf{m}}(\lambda_{b'})$ is the value of $W^{\mathsf{m}}$ on the loop $\lambda_{b'}$, in the
notation of Definition~\ref{5.7:not-canonical-loops} for loops and class functions on the two lattices.
The comb dilation of the box read from the trivial configuration $\mathbb{1}$ is the map
\begin{equation*}
\mathsf{M}^{\mathrm{ax}}:\ \bigl((b,\sigma,t,\mathsf{t},\varsigma);p\bigr)\longmapsto
\bigl(U^{\mathsf{B}}_{\varsigma}\big|_{C_y^{\sharp}},\ U^{\mathsf{A}}_{\varsigma}\big|_{C_y^{\sharp}}\bigr),
\qquad p\in\mathcal{K}^{\mathrm{rd}}(k;\mathsf{n}+1,\hat X),
\end{equation*}
with target slot $(j,C_y;\circ)$, $C_y$ denoting the region of that slot; in the listed-move format it is a
move from the stage to the unit.
Here $t$ is the coordinate of the factor $[0,1]$ of $\mathfrak{P}_{\mathrm{ax}}$, not the ratio $L^{j-n}$ of
\eqref{5.7:eq-comb-dilation-1}, $\mathsf{t}$ is the coordinate of $\mathbb{D}_{\square}$, $b$ is the index of
the listed-move format, not the bond argument of $\mathsf{b}_{\ast}$ and of the comb loops, and $p$ is a member
of the parent class $\mathcal{K}^{\mathrm{rd}}(k;\mathsf{n}+1,\hat X)$, not a plaquette. The map is given the
format of the listed moves and is adjoined to the list
$\mathfrak{M}^{\mathrm{rd}}$ of moves that generates the reduced classes $\mathcal{K}^{\mathrm{rd}}$.

The construction is the same for $\mathsf{A}$ and for $\mathsf{B}$: the comb from $y$ is an object of the
source-free scaffold of Definition~\ref{5.1:def-source-free-scaffold}; the map does not depend on the complex
sources $z_i$; and both components are defined on the one
parameter domain $\mathfrak{P}_{\mathrm{ax}}$. In this sense it is the orthogonal dilation read at the pair
$(\mathbb{1},W)$ in place of the pair $(W,W^{1})$ at which that dilation is defined; it is also the dilation of
the chart-based move, with the chart replaced by the comb formed at the lattice's own level; and it makes
canonical the $\sigma$-slot of the interpolating exponential in $\sigma$ and $\sigma'$ of the
interpolation construction. With this move listed,
\begin{equation}\label{5.7:eq-comb-dilation-a}
\mathsf{M}^{\mathrm{ax}}\bigl(\mathfrak{P}_{\mathrm{ax}}\times\mathcal{K}^{\mathrm{rd}}(k;\mathsf{n}+1,\hat X)\bigr)
\subset\mathcal{K}^{\mathrm{rd}}(j,C_y;\circ),
\end{equation}
and $\mathsf{a}$ denotes the $\delta$-size over the class so enlarged.
\end{enumerate}
\end{definition}

\begin{proof}
We first prove the identity in \eqref{5.7:eq-comb-dilation-2}. Every point of $C_y^{\sharp}$ satisfies
$y_\kappa-\ell_j+\eta\le z_\kappa\le y_\kappa+\ell_j$ in every coordinate. Hence it lies at sup-distance at
most $\ell_j$ from $y$, that is, at distance at most one in units of $\ell_j$. Hence, over bonds $b$ with both
endpoints in $C_y^{\sharp}$, one has $\sup(1+|b-y|)=2$, the value being attained by a bond lying in an upper
face $z_\kappa=y_\kappa+\ell_j$, and
\begin{equation*}
\sup_{C_y^{\sharp}}\mathsf{b}_{\ast}=2d\cdot2\,\mathsf{f}_{\ast}=4d\,\mathsf{f}_{\ast}.
\end{equation*}

We now prove the two bounds on $w^{\sharp}(y)$ in \eqref{5.7:eq-comb-dilation-2}, for $\nu\in\{1,2\}$.

\emph{Lower bound.} Since $n=\min\{l(y),n_0\}\le l(y)$, and $y$ is a point of $T_{\ell_{l(y)}}$, which is
contained in $T_{\ell_n}$ because $\ell_{l(y)}$ is an integer power of $L$ times $\ell_n$, the anchor $y$ is a
point of $T_{\ell_n}$. The $n$-cells have their corners on $T_{\ell_n}$, so the number
$\operatorname{dist}_n(\Delta,y)$ is a non-negative integer for
every $n$-cell $\Delta$, and the condition $\operatorname{dist}_n(\Delta,y)<1$ in the definition of $w^{\vee}(y)$
means $y\in\Delta$. Since $\eta\le\ell_j$ and $y\in T_{\ell_j}\subset T_{\eta}$, the point $y$ satisfies
$y_\kappa-\ell_j+\eta\le y_\kappa\le y_\kappa+\ell_j$ for every $\kappa$, so $y\in C_y^{\sharp}$. Hence every cell
entering the maximum defining $w^{\vee}(y)$ meets $C_y^{\sharp}$ and enters the maximum defining $w^{\sharp}(y)$,
which gives $w^{\vee}(y)\le w^{\sharp}(y)$.

\emph{Upper bound, case $\nu=2$.} Here $n=j+1$. The cube $C_y^{\sharp}$ has half-width $\ell_j$ about $y$,
which in units of $\ell_n$ is $\ell_j/\ell_n=1/L$. A cell $\Delta$ meeting $C_y^{\sharp}$ therefore satisfies
\begin{equation*}
\operatorname{dist}_n(\Delta,y)\le 1/L<1,
\end{equation*}
and so $w^{\sharp}(y)\le w^{\vee}(y)$. Moreover $y\in\Delta'_y$, so for every cell $\Delta$
\begin{equation*}
\operatorname{dist}_n(\Delta,\Delta'_y)\le\operatorname{dist}_n(\Delta,y).
\end{equation*}
Hence every cell entering $w^{\vee}(y)$ also enters $w^{\vee}(\Delta'_y)$, and
$w^{\vee}(y)\le w^{\vee}(\Delta'_y)$. Together, $w^{\sharp}(y)\le w^{\vee}(\Delta'_y)$.

\emph{Upper bound, case $\nu=1$.} Here $n=j$ and $\ell_n=\ell_j$. The upper faces $z_\kappa=y_\kappa+\ell_n$
of $C_y^{\sharp}$ lie in cells at $\operatorname{dist}_n=1$ from $y$. They are, however, within $\eta$ of the
cell $\Delta'_y$, whose lower corner is $y$. In the other direction the coordinates of $C_y^{\sharp}$ stop at
$y_\kappa-\ell_j+\eta$. Hence every point of $C_y^{\sharp}$ lies at sup-distance strictly less than $\ell_n$
from $\Delta'_y$. A cell $\Delta$ meeting $C_y^{\sharp}$ therefore has $\operatorname{dist}_n(\Delta,\Delta'_y)<1$,
so it enters the maximum defining $w^{\vee}(\Delta'_y)$. This gives $w^{\sharp}(y)\le w^{\vee}(\Delta'_y)$.

In the sums over blocks in which this weight is used, the weight is read as $w^{\vee}(\Delta')$ with
$\Delta'\ni y$. For $y\in\Delta'$ one has $w^{\vee}(\Delta')\ge w^{\vee}(y)$: every cell with
$\operatorname{dist}_n(\Delta,y)<1$ satisfies
\begin{equation*}
\operatorname{dist}_n(\Delta,\Delta')\le\operatorname{dist}_n(\Delta,y)<1 .
\end{equation*}
By this inequality and by the clause of the block-weight statement on which those sums rest, the upper bound
by $w^{\vee}(\Delta'_y)$ does not change those sums.

Finally, \eqref{5.7:eq-comb-dilation-a} holds by construction. The reduced classes $\mathcal{K}^{\mathrm{rd}}$
are defined as the classes generated by the moves of the list $\mathfrak{M}^{\mathrm{rd}}$. Once
$\mathsf{M}^{\mathrm{ax}}$ is on that list, its image of $\mathfrak{P}_{\mathrm{ax}}$ times the parent class
$\mathcal{K}^{\mathrm{rd}}(k;\mathsf{n}+1,\hat X)$ lies in the class of its target slot $(j,C_y;\circ)$.
\end{proof}

\begin{lemma}[Bounds for the comb dilation on one lattice]\label{5.7:lem-comb-dilation-bounds}
Let a block $\square$ with its enlarged blocks
$\tilde{\square}^{2}\subset\tilde{\square}^{3}\subset\tilde{\square}^{4}$, a near companion
$\hat A=A(h_y,\cdot)$ of level $j$ anchored at
$y\in\tilde{\square}^{2}$, and the quantities $n_0$, $n=\min\{\ell(y),n_0\}$, $t$,
$u_j=(\eta/\ell_j)^2$, $C_y^{\sharp}$,
$w^{\sharp}(y)$, $\Theta^{\sharp}=c_Ww^{\sharp}(y)$, $\mathsf{f}_{\ast}=\Theta^{\sharp}\mathfrak{a}_nt^2$,
$\mathsf{b}_{\ast}(b)=2d(1+|b-y|)\mathsf{f}_{\ast}$, $W$, $\gamma_z$, $\lambda_b$, $a^{\mathrm{ax}}$,
$S^{\mathrm{ax}}_v$ and $U_{\varsigma}(b)=\exp(\varsigma\log W(\lambda_b))$ be as in
Definition~\ref{5.7:def-comb-dilation}, all taken on one cutoff $\mathsf{m}\in\{\mathsf{A},\mathsf{B}\}$ on
its own lattice, of spacing $\eta$ for $\mathsf{A}$ and $\eta'$ for $\mathsf{B}$ (we write $\eta$ for the
spacing throughout), in the working gauge. In the combinations $\mathfrak{a}_nt^2$ and $t^2u_j$ the symbol $t$
denotes the level ratio $L^{j-n}$ of Definition~\ref{5.7:def-comb-dilation}; the interpolation parameter
$t\in[0,1]$ of the parameter domain enters only through the continuity statement in (v). For a plaquette
$q$ we write $\partial W(q)$ for the holonomy of $W$ around $q$. The assertions below hold at every value
of the complex source, every $t\in[0,1]$, every $\mathsf{t}\in\mathbb{D}_{\square}$ and every $\sigma$
with $|\sigma|\le e^{\kappa_1}$, under the following hypotheses.
\begin{enumerate}
\item[(a)] (Curvature bound for the box field.) At every cell $x$ meeting $\tilde{\square}^{3}$ and at
every level $j'$ up to the current level $k$ of the parent configuration of
Definition~\ref{5.7:def-comb-dilation}, in particular at the companion's level $j$,
\begin{equation*}
|\partial\mathfrak{b}'_{j'}(x)-1|\le 4\theta_2(x),\qquad
\theta_2(x)=w(x)\,\tilde\alpha_{0,\ell(x)}\,\eta^2\ell_{\ell(x)}^{-2},
\end{equation*}
in the gauge of the box field $\mathfrak{b}'_{j'}$, where $\ell(x)$ is the level of $x$, $w(x)$ its
weight, $\alpha_{0,i}$ the small-field radius at level $i$ and $\tilde\alpha_{0,i}$ its enlarged variant
entering this bound; in the working gauge of Definition~\ref{5.7:def-comb-dilation} the left side is larger
by at most a factor $2^{1/2}$; the constants satisfy $\tilde\alpha_{0,i}\le 2\alpha_{0,i}$ at every
level $i$, $c_{\mathrm{comb}}\ge 1$ for the constant $c_{\mathrm{comb}}$ of the order of choice that enters
the choice of $c_W$ below, $K_{\mathrm{loc}}\ge 12C_M^{\mathrm{av}}\ge 24$, and, in the order of choice
of the constants, $K_{\mathrm{loc}}L\ge 192$ and
$c_W\ge 2c_{\mathrm{comb}}K_{\mathrm{loc}}L^3\alpha_{0,n}/\mathfrak{a}_n$.
\item[(b)] (Levels near the anchor.) Every cell meeting $y+[-4,4]^d\ell_n$ has level at least $n-1$.
(This is the second clause of the level lemma for the multiscale frames, whose proof is run with its
index equal to $n$; it applies because $\ell(y)\ge n$.)
\item[(c)] (Comparison of radii.) For the cells $x$ of (b), $\alpha_{0,\ell(x)}\le 2\alpha_{0,n}$. (This
is the comparison of small-field radii over the cells of one unit region, taken by statement from the
construction of the energy-difference datum.)
\item[(d)] (Ladder representation and holonomy bound.) For every bond $b\subset C_y^{\sharp}$, $W(\lambda_b)$
is an ordered product of at most $(d-1)r_b/\eta+1$ factors
$\operatorname{Ad}_{W(\rho_l)}(\partial W(q_l)^{\pm1})$ with plaquettes $q_l\subset C_y^{\sharp}$ and
comb contours $\rho_l\subset C_y^{\sharp}$ of at most $4d\ell_j/\eta$ bonds, $r_b$ denoting the
sup-distance of $b$ from $y$, so that $r_b=\ell_j|b-y|$ with $|b-y|\le 1$ on $C_y^{\sharp}$; the proof of
this representation uses only that the comb $\{\gamma_z\}$ is lexicographic. Moreover, the box field
satisfies on $\tilde{\square}^{4}$ the curvature ceiling on which the contour-holonomy bound rests, and,
for the norm $n(\cdot)$ of the contour-holonomy bound, $n(W(\rho))\le 1.5$ for every contour
$\rho\subset C_y^{\sharp}$ of at most $4d\ell_j/\eta$ bonds through cells of level at least $j-1$, given
the ceilings of (f); this is the contour-holonomy bound
with its level and region taken to be $j$ and $C_y^{\sharp}$.
\item[(e)] (Holomorphy of the base.) $W$ is holomorphic in the data $(\mathbb{U},\mathbb{J},b,\sigma,\mathsf{t})$
of Definition~\ref{5.7:def-comb-dilation}, $b$ here denoting the source datum and not a bond, as asserted
in clause (b) of the theorem on the energy-difference datum.
\item[(f)] (Ceilings.) The configurations at which the lemma is applied satisfy the a priori ceilings of
the first clause of the remainder-ceiling statement; apart from the inputs stated in (a) and (d), no
further a priori ceilings of the construction are assumed.
\item[(g)] (Continuity of the base.) $W$ is continuous in $t$.
\end{enumerate}
Assume $4d\mathsf{f}_{\ast}<a^{\mathrm{ax}}$, so that $S^{\mathrm{ax}}_v>1$. Put
$c_F:=12/(c_{\mathrm{comb}}K_{\mathrm{loc}}L)$, $\mathfrak{n}^{\mathrm{ax}}:=0.1\,\mathfrak{c}_{\mathrm{tr}}
(a^{\mathrm{ax}})^2\mathsf{p}_d$ $(d=4)$, $g(z):=W(\gamma_z)$ and $\varphi(\varsigma):=\hat A(U_{\varsigma})$.
Then $c_F\le 1/16$ and, for every plaquette $q\subset C_y^{\sharp}$, every bond $b\subset C_y^{\sharp}$,
every $|\varsigma|<S^{\mathrm{ax}}_v$ and every plaquette $P\subset C_y^{\sharp}$,
\begin{equation}\label{5.7:eq-comb-dilation-bounds-a}
\begin{aligned}
&\text{(i)} && (\ell_j/\eta)^2\,|\partial W(q)-1|\le c_F\,\mathsf{f}_{\ast},\\
&\text{(ii)} && |\log W(\lambda_b)|\le \tfrac18\,\mathsf{b}_{\ast}(b)\,u_j^{1/2},\\
&\text{(iii)} && |\partial U_{\varsigma}(P)-1|\le 0.3\,a^{\mathrm{ax}}u_j,\qquad
|\hat A(U_{\varsigma})|\le\mathfrak{n}^{\mathrm{ax}},\\
&\text{(iv)} && U_0=\mathbb{1},\qquad U_1=W^{g},\qquad \hat A(U_1)=\hat A(\mathfrak{b}'_j),\\
&\text{(v)} && \varphi(0)=\varphi'(0)=0;
\end{aligned}
\end{equation}
moreover $\varphi$ is entire in $\varsigma$, holomorphic in the data
$(\mathbb{U},\mathbb{J},b,\sigma,\mathsf{t})$ of Definition~\ref{5.7:def-comb-dilation}, $b$ here denoting
the source datum and not a bond, and continuous in $t$.
\end{lemma}

\begin{proof}
(i) By the definition of $C_y^{\sharp}$, every $z\in C_y^{\sharp}$ satisfies
$y_\kappa-\ell_j+\eta\le z_\kappa\le y_\kappa+\ell_j$ for all $\kappa$, so
$C_y^{\sharp}\subset y+[-1,1]^d\ell_j$; since $y\in\tilde{\square}^{2}$ and
$\ell_j\le\ell_{n_0}\le M\ell_{n_0}$, this cube lies in $\tilde{\square}^{3}$. Hence (a) applies at every
cell $x$ meeting $C_y^{\sharp}$. Such a cell meets $y+[-4,4]^d\ell_n$, because $n\ge j$ and so
$\ell_j\le\ell_n$; by (b) its level satisfies $\ell(x)\ge n-1$ (the argument behind (b) rests on the
level-gradient bound and the flatness condition of the frames), so
$\ell_{\ell(x)}\ge\ell_{n-1}=L^{-1}\ell_n$ and
\begin{equation*}
\eta^2\ell_{\ell(x)}^{-2}\le L^2\eta^2\ell_n^{-2}=L^2t^2u_j,
\end{equation*}
the equality because $\ell_n=L^{n-j}\ell_j=t^{-1}\ell_j$ and $u_j=(\eta/\ell_j)^2$. By (a),
$\tilde\alpha_{0,\ell(x)}\le 2\alpha_{0,\ell(x)}$; by (c), $\alpha_{0,\ell(x)}\le 2\alpha_{0,n}$; and
$w(x)\le w^{\sharp}(y)$, since $x$ meets $C_y^{\sharp}$. Therefore
\begin{align*}
4\theta_2(x)&\le 4\,w^{\sharp}(y)\cdot 4\alpha_{0,n}\cdot L^2t^2u_j
=16L^2w^{\sharp}(y)\alpha_{0,n}t^2u_j\\
&\le \frac{8}{c_{\mathrm{comb}}K_{\mathrm{loc}}L}\,\Theta^{\sharp}\mathfrak{a}_nt^2u_j
=\frac{8}{c_{\mathrm{comb}}K_{\mathrm{loc}}L}\,\mathsf{f}_{\ast}u_j,
\end{align*}
where the second inequality is
$c_W\ge 2c_{\mathrm{comb}}K_{\mathrm{loc}}L^3\alpha_{0,n}/\mathfrak{a}_n$ multiplied by $w^{\sharp}(y)$,
i.e.\ $16L^2w^{\sharp}(y)\alpha_{0,n}\le 8\Theta^{\sharp}\mathfrak{a}_n/(c_{\mathrm{comb}}K_{\mathrm{loc}}L)$.
Passing to the
working gauge costs at most a factor $2^{1/2}$ by (a), and $8\cdot 2^{1/2}<12$; multiplying by
$(\ell_j/\eta)^2=u_j^{-1}$ gives (i) with $c_F=12/(c_{\mathrm{comb}}K_{\mathrm{loc}}L)$. Since
$c_{\mathrm{comb}}\ge 1$ and
$K_{\mathrm{loc}}L\ge 192$ by (a), we get $c_F\le 12/192=1/16$.

(ii) We use the ladder representation of (d) on $C_y^{\sharp}$. Every contour $\rho_l$ lies in
$C_y^{\sharp}$ and has at most $4d\ell_j/\eta$ bonds, and the cells it passes through have level at least
$n-1\ge j-1$ by (b); these cells meet $C_y^{\sharp}\subset\tilde{\square}^{3}\subset\tilde{\square}^{4}$,
so they are covered by the curvature ceiling of (d). Hence the holonomy bound of (d) applies, with the
ceilings of (f), and gives
$n(W(\rho_l))\le 1.5$ for the norm $n(\cdot)$ of the contour-holonomy bound. The argument behind that bound
uses only the curvature ceiling of the box field on $\tilde{\square}^{4}$, the lower bound $j-1$ on the
levels of the cells, the bound $4d\ell_j/\eta$ on the number of bonds of the contours, and the ceilings
of (f). By (i), each factor
$\operatorname{Ad}_{W(\rho_l)}(\partial W(q_l)^{\pm1})$ differs from $1$ by at most
$(1.5)^2\cdot 1.1\,c_F\mathsf{f}_{\ast}u_j=2.25\cdot 1.1\,c_F\mathsf{f}_{\ast}u_j$, the factor $1.1$
covering the inverse plaquette holonomies: as $|\partial W(q_l)-1|\le c_F\mathsf{f}_{\ast}u_j<10^{-2}$,
\begin{equation*}
|\partial W(q_l)^{-1}-1|\le\frac{|\partial W(q_l)-1|}{1-|\partial W(q_l)-1|}\le 1.1\,|\partial W(q_l)-1|.
\end{equation*}
Since the number of
factors is at most $(d-1)r_b/\eta+1$, and $r_bu_j/\eta=|b-y|\,u_j^{1/2}$, $u_j=u_j^{1/2}\,u_j^{1/2}$,
the sum of these deviations is at most
\begin{equation*}
s_b:=2.25\cdot 1.1\,\bigl[(d-1)|b-y|+u_j^{1/2}\bigr]\,u_j^{1/2}\,c_F\mathsf{f}_{\ast}.
\end{equation*}
For an ordered product of $V_l$ one has
\begin{equation*}
\Bigl|\prod_lV_l-1\Bigr|\le\prod_l\bigl(1+|V_l-1|\bigr)-1\le\exp\Bigl(\sum_l|V_l-1|\Bigr)-1,
\end{equation*}
hence
$|W(\lambda_b)-1|\le e^{s_b}-1$. As $d-1\le d$ and $u_j^{1/2}\le 1\le d$,
$s_b\le 2.475\,d(1+|b-y|)c_F\mathsf{f}_{\ast}u_j^{1/2}$. Since $|b-y|\le 1$ for $b\subset C_y^{\sharp}$
(so that $\sup_{C_y^{\sharp}}\mathsf{b}_{\ast}=4d\mathsf{f}_{\ast}$), and
$4d\mathsf{f}_{\ast}<a^{\mathrm{ax}}\le 1/16$, $c_F\le 1/16$, we get $s_b<10^{-2}$; then
$e^{s_b}-1\le 1.01\,s_b$ and
\begin{equation*}
|W(\lambda_b)-1|\le e^{s_b}-1\le 2.6\,d(1+|b-y|)\,c_F\mathsf{f}_{\ast}u_j^{1/2}.
\end{equation*}
Since $|\log V|\le 1.1|V-1|$ for $|V-1|\le 10^{-2}$, and
$1.1\cdot 2.6\,d\,c_F\le\tfrac{3}{16}d=\tfrac18\cdot 2d\cdot\tfrac34$, we obtain
$|\log W(\lambda_b)|\le\tfrac34\cdot\tfrac18\,\mathsf{b}_{\ast}(b)u_j^{1/2}$, which implies (ii).

(iii) Put $X_b:=\log W(\lambda_b)$, let $b_1,\dots,b_4$ be the bonds of $P$ and
$\bar x:=\max_i|X_{b_i}|$. Since $\mathsf{b}_{\ast}\le 4d\mathsf{f}_{\ast}$ on $C_y^{\sharp}$, (ii) gives
$\bar x\le\tfrac18\cdot 4d\mathsf{f}_{\ast}u_j^{1/2}$. The holonomy $\partial U_{\varsigma}(P)$ is the
ordered product of the $e^{\pm\varsigma X_{b_i}}$ around $P$; by the Baker--Campbell--Hausdorff formula,
\begin{equation*}
Z(\varsigma):=\log\partial U_{\varsigma}(P)=\varsigma L_X+\varsigma^2\mathsf{C}(X,X)+R_3(\varsigma),
\end{equation*}
where $L_X$ is the signed sum of the $X_{b_i}$ around $P$, $\mathsf{C}$ is the bilinear term and $R_3$ the
third-order remainder, with $|\mathsf{C}(X,X)|\le 6\bar x^2$ and $|R_3(\varsigma)|\le C_3(|\varsigma|\bar
x)^3$ for $|\varsigma|\bar x\le 1/16$, where $C_3:=90$. For the last bound put
$\xi:=\sum_i|\varsigma X_{b_i}|\le 4|\varsigma|\bar x\le\tfrac14$. Since
$|\partial U_{\varsigma}(P)-1|\le e^{\xi}-1<1$, the logarithm series
$\sum_{m\ge1}(-1)^{m+1}(\partial U_{\varsigma}(P)-1)^m/m$ converges; expanding each
$e^{\pm\varsigma X_{b_i}}$ in its power series, the sum of the norms of the terms of total degree $m$ in
the $\varsigma X_{b_i}$ is at most the coefficient of $\xi^m$ in
\begin{equation*}
\sum_{m\ge1}\frac{(e^{\xi}-1)^m}{m}=-\log(2-e^{\xi})=\xi+\xi^2+\xi^3+\cdots,
\end{equation*}
a series with nonnegative coefficients. Its part of degree at least three, divided by $\xi^3$, increases
with $\xi$, and at $\xi=\tfrac14$ that part equals
\begin{equation*}
-\log\bigl(2-e^{1/4}\bigr)-\tfrac14-\tfrac1{16}<0.0217<1.4\cdot4^{-3};
\end{equation*}
so it is at most $1.4\,\xi^3$ for $\xi\le\tfrac14$. As $R_3(\varsigma)$ is the sum of the terms of degree
at least three,
$|R_3(\varsigma)|\le 1.4\,\xi^3\le 1.4\cdot 64\,(|\varsigma|\bar x)^3\le C_3(|\varsigma|\bar x)^3$.
At $\varsigma=1$,
$\partial U_1(P)=g\,\partial W(P)\,g^{-1}$ by the
telescoping in (iv), the conjugating holonomy being that of a comb contour, of norm at most $1.5$ by (d);
with (i) and $|\log V|\le 1.1|V-1|$ this gives $|Z(1)|\le 1.1\cdot 2.25\,c_F\mathsf{f}_{\ast}u_j$. Since
$\bar x\le\tfrac12 d\mathsf{f}_{\ast}<1/16$,
\begin{equation*}
|L_X|\le|Z(1)|+|\mathsf{C}(X,X)|+|R_3(1)|\le 2.5\,c_F\mathsf{f}_{\ast}u_j+6\bar x^2+C_3\bar x^3.
\end{equation*}
Let $|\varsigma|<S^{\mathrm{ax}}_v=a^{\mathrm{ax}}/(4d\mathsf{f}_{\ast})$. Then
$|\varsigma|\bar x\le\tfrac18a^{\mathrm{ax}}u_j^{1/2}\le 1/128$, so the expansion applies, and
\begin{align*}
|Z(\varsigma)|&\le|\varsigma|\,|L_X|+6|\varsigma|^2\bar x^2+C_3(|\varsigma|\bar x)^3\\
&\le\frac{a^{\mathrm{ax}}}{4d\mathsf{f}_{\ast}}\bigl(2.5\,c_F\mathsf{f}_{\ast}u_j+6\bar x^2
+C_3\bar x^3\bigr)\\
&\qquad+6\bigl(\tfrac18a^{\mathrm{ax}}u_j^{1/2}\bigr)^2+C_3\bigl(\tfrac18a^{\mathrm{ax}}u_j^{1/2}\bigr)^3\\
&\le u_ja^{\mathrm{ax}}\Bigl[\frac{2.5\,c_F}{4d}+\frac{6}{64}\,4d\mathsf{f}_{\ast}
+\frac{6}{64}\,a^{\mathrm{ax}}\\
&\qquad+\frac{C_3d^2\mathsf{f}_{\ast}^2u_j^{1/2}}{32}+\frac{C_3(a^{\mathrm{ax}})^2u_j^{1/2}}{512}\Bigr]
\le\tfrac14\,a^{\mathrm{ax}}u_j,
\end{align*}
where we used $\bar x^2\le\tfrac1{64}(4d\mathsf{f}_{\ast})^2u_j$ and
$\bar x^3\le(\tfrac12d\mathsf{f}_{\ast})^3u_j^{3/2}$; the first three terms of the bracket are at most
$0.0098$, $0.0059$, $0.0059$ by $c_F\le 1/16$, $4d\mathsf{f}_{\ast}<a^{\mathrm{ax}}\le 1/16$ and $d=4$, and
the last two carry the factor $2^{-17}C_3$, so that, as $C_3=90$, the bracket is at most
$0.0216+2^{-16}\cdot 90<0.0236\le\tfrac14$. Hence,
since $|Z(\varsigma)|\le\tfrac14a^{\mathrm{ax}}u_j\le 1/64$,
\begin{equation*}
|\partial U_{\varsigma}(P)-1|=|e^{Z(\varsigma)}-1|\le|Z(\varsigma)|\,e^{|Z(\varsigma)|}
\le\tfrac14 e^{1/64}a^{\mathrm{ax}}u_j\le 0.3\,a^{\mathrm{ax}}u_j.
\end{equation*}
For the action, $a(V)=-\tfrac12\operatorname{tr}(V^{-1}(V-1)^2)$ satisfies
$|a(V)|\le\mathfrak{c}_{\mathrm{tr}}|V-1|^2$ for $|V-1|\le\tfrac12$, which holds for
$V=\partial U_{\varsigma}(P)$. Every plaquette $P$ with $h_y(x(P))\ne0$, $x(P)$ its barycentre, has all its
corners in $C_y^{\sharp}$, and $\sum_Ph_y(x(P))\,u_j^2=\mathsf{p}_d(\ell_j/\eta)^du_j^2=\mathsf{p}_d$ at
$d=4$; hence
\begin{equation*}
|\hat A(U_{\varsigma})|\le\sum_Ph_y(x(P))\,\mathfrak{c}_{\mathrm{tr}}(0.3\,a^{\mathrm{ax}}u_j)^2
=0.09\,\mathfrak{c}_{\mathrm{tr}}(a^{\mathrm{ax}})^2\mathsf{p}_d\le\mathfrak{n}^{\mathrm{ax}}.
\end{equation*}

(iv) $U_0(b)=\exp(0)=\mathbb{1}$. At $\varsigma=1$, $U_1(b)=W(\lambda_b)$, and since
$\lambda_b=\gamma_{b_-}b\gamma_{b_+}^{-1}$,
\begin{equation*}
U_1(b)=W(\gamma_{b_-})\,W(b)\,W(\gamma_{b_+})^{-1}=g(b_-)\,W(b)\,g(b_+)^{-1}=W^{g}(b).
\end{equation*}
Around a plaquette the factors $g$ telescope, so $\partial U_1(P)=g\,\partial W(P)\,g^{-1}$ with $g$
evaluated at the base corner of $P$; since $a$ is a class function,
$a(\partial U_1(P))=a(\partial W(P))$ for every $P$, and $\hat A(U_1)=\hat A(W)=\hat A(\mathfrak{b}'_j)$,
$W$ being the bond-variable component of $\mathfrak{b}'_j$.

(v) By (e), $W$ is holomorphic in $(\mathbb{U},\mathbb{J},b,\sigma,\mathsf{t})$, and by (g) it is
continuous in $t$; by the proof of (ii) (the display preceding its last step),
$|W(\lambda_b)-1|<10^{-2}$, so $X_b=\log W(\lambda_b)$, given by the convergent logarithm series,
has the same regularity. For fixed parameters, $\varsigma\mapsto e^{\pm\varsigma X_b}$ is entire, hence so
are $\partial U_{\varsigma}(P)$ and its inverse, and $a(\partial U_{\varsigma}(P))$, a polynomial in these;
$\varphi$ is a finite sum of such terms with weights $h_y(x(P))$, which proves the regularity assertions.
Finally $Z(\varsigma)=O(\varsigma)$ by the expansion in (iii), and since
$e^{-Z}(e^Z-1)^2=(e^{Z/2}-e^{-Z/2})^2$ contains only even powers of $Z$,
$a(e^Z)=-\tfrac12\operatorname{tr}Z^2+O(Z^4)$. Hence $\varphi(\varsigma)=O(\varsigma^2)$, that is,
$\varphi(0)=\varphi'(0)=0$.

For the cutoff $\mathsf{B}$ the statement and its proof are the same, with $\eta$ replaced by $\eta'$.
\end{proof}

\begin{theorem}[The low-order variation bound at the second cube size]
\label{5.7:thm-low-order-variation}
Let $\mathsf{A}$, $\mathsf{B}$ be two cutoffs on one scaffold. Assume the following hypotheses.
\begin{itemize}
\item Hypothesis \eqref{5.7:eq-comb-dilation-a} holds at the listed move $\mathsf{M}^{\mathrm{ax}}$. This move
is the comb dilation from the trivial configuration (Definition~\ref{5.7:def-comb-dilation}).
\item The estimate for the companions holds.
\item The bound on the energy-difference datum $\mathsf{E}$ holds.
\item The bound at the flat profile rate holds, together with the lemma in which it is stated.
\end{itemize}
Let $\hat A$ be the functional of clause~(iii) of Lemma~\ref{5.7:lem-comb-dilation-bounds}, there
bounded by $\mathfrak{n}^{\mathrm{ax}}$ along the comb dilation, and put
$\mathsf{a}:=\mathsf{S}^{\mathrm{rd}}_j[\hat A]$, the $\delta$-size of $\hat A$ over the reduced class
$\mathcal{K}^{\mathrm{rd}}(j,C_y;\circ)$.

Then the following holds for every point of the reduced class $\mathcal{K}^{\mathrm{rd}}(k;\mathsf{n}+1,\hat X)$
and for all parameters at the second cube size. Wherever $4d\,\mathsf{f}_{\ast}<a^{\mathrm{ax}}$,
\begin{equation}\label{5.7:eq-low-order-variation-1}
\bigl|\delta\hat A(\mathfrak{b}'_j)\bigr|
\le \mathsf{a}\,\Bigl(\frac{4d\,\Theta^{\sharp}\mathfrak{a}_n t^2}{a^{\mathrm{ax}}}\Bigr)^{2}
\le \mathsf{a}\,\bigl(C D\,\Theta^{\sharp}\mathfrak{a}_n t^2\bigr)^{2}.
\end{equation}
Here $D:=M(2\nu+3)$, so that $D\ge 5M$, and $C$ is a constant with $C\ge 4d/(5Ma^{\mathrm{ax}})$.
The bound holds for every $\nu$; the case $\nu\le 2$ gives the statement its name.
\end{theorem}

\begin{proof}
Fix a point of $\mathcal{K}^{\mathrm{rd}}(k;\mathsf{n}+1,\hat X)$ and the parameters at the second cube size, and assume
$4d\,\mathsf{f}_{\ast}<a^{\mathrm{ax}}$, so that $S^{\mathrm{ax}}_v>1$, as in the hypothesis of
Lemma~\ref{5.7:lem-comb-dilation-bounds}.

For $\mathsf{m}\in\{\mathsf{A},\mathsf{B}\}$ let $\varphi^{\mathsf{m}}$ be the function of clause~(v) of
Lemma~\ref{5.7:lem-comb-dilation-bounds} for the cutoff $\mathsf{m}$. It is defined along the comb dilation
$U^{\mathsf{m}}_{\varsigma}$ for $|\varsigma|<S^{\mathrm{ax}}_v$. Set
\begin{equation*}
\delta\varphi(\varsigma):=\varphi^{\mathsf{B}}(\varsigma)-\varphi^{\mathsf{A}}(\varsigma),
\qquad |\varsigma|<S^{\mathrm{ax}}_v .
\end{equation*}
The function $\delta\varphi$ is holomorphic on the disc $|\varsigma|<S^{\mathrm{ax}}_v$.

It satisfies $|\delta\varphi(\varsigma)|\le\mathsf{a}$ on that disc. This follows from the principle
bounding two-cutoff differences along a move, applied with hypothesis \eqref{5.7:eq-comb-dilation-a}
at $\mathsf{M}^{\mathrm{ax}}$: along the comb dilation, the difference of the values of $\hat A$ at
$U^{\mathsf{B}}_{\varsigma}$ and at $U^{\mathsf{A}}_{\varsigma}$ is bounded by $\mathsf{a}$, the
$\delta$-size of $\hat A$ over the reduced class $\mathcal{K}^{\mathrm{rd}}(j,C_y;\circ)$.

Clause~(v) of Lemma~\ref{5.7:lem-comb-dilation-bounds}, applied for each cutoff separately, gives
$\varphi^{\mathsf{m}}(0)=0$ and $(\varphi^{\mathsf{m}})'(0)=0$ for $\mathsf{m}=\mathsf{A}$ and for
$\mathsf{m}=\mathsf{B}$. Subtracting, $\delta\varphi(0)=0$ and $\delta\varphi'(0)=0$.

We now apply the Schwarz lemma for a function with a double zero at the origin. Since $\delta\varphi$ has a double
zero at $0$, the quotient $\delta\varphi(\varsigma)/\varsigma^{2}$ is holomorphic on $|\varsigma|<S^{\mathrm{ax}}_v$.
For $0<r<S^{\mathrm{ax}}_v$, the maximum modulus principle on $|\varsigma|\le r$ gives
\begin{equation*}
\Bigl|\frac{\delta\varphi(\varsigma)}{\varsigma^{2}}\Bigr|\le\frac{\mathsf{a}}{r^{2}}
\qquad\text{for } |\varsigma|\le r .
\end{equation*}
Letting $r\uparrow S^{\mathrm{ax}}_v$ yields
$|\delta\varphi(\varsigma)|\le\mathsf{a}\,|\varsigma|^{2}(S^{\mathrm{ax}}_v)^{-2}$ on the whole disc. The point
$\varsigma=1$ lies in the disc because $S^{\mathrm{ax}}_v>1$. Hence
\begin{equation*}
|\delta\varphi(1)|\le \mathsf{a}\,(S^{\mathrm{ax}}_v)^{-2}.
\end{equation*}

Since $\varphi^{\mathsf{m}}(\varsigma)=\hat A(U^{\mathsf{m}}_{\varsigma})$ by clause~(v) of
Lemma~\ref{5.7:lem-comb-dilation-bounds}, we have $\varphi^{\mathsf{m}}(1)=\hat A(U^{\mathsf{m}}_1)$, and
clause~(iv) of that lemma gives $\hat A(U^{\mathsf{m}}_1)=\hat A(\mathfrak{b}'_j)$ for each cutoff. Subtracting,
$\delta\varphi(1)=\delta\hat A(\mathfrak{b}'_j)$. By Definition~\ref{5.7:def-comb-dilation}, in which the radius
$S^{\mathrm{ax}}_v$ is defined,
\begin{equation*}
(S^{\mathrm{ax}}_v)^{-1}\le\frac{4d\,\Theta^{\sharp}\mathfrak{a}_n t^2}{a^{\mathrm{ax}}}.
\end{equation*}
Combining $|\delta\varphi(1)|\le\mathsf{a}(S^{\mathrm{ax}}_v)^{-2}$ with
$\delta\varphi(1)=\delta\hat A(\mathfrak{b}'_j)$, and then squaring the last inequality, both sides of which are
non-negative, and multiplying by $\mathsf{a}\ge 0$, we obtain
\begin{align*}
\bigl|\delta\hat A(\mathfrak{b}'_j)\bigr|
&\le\mathsf{a}\,(S^{\mathrm{ax}}_v)^{-2}\\
&\le\mathsf{a}\,\Bigl(\frac{4d\,\Theta^{\sharp}\mathfrak{a}_n t^2}{a^{\mathrm{ax}}}\Bigr)^{2},
\end{align*}
which is the first inequality of \eqref{5.7:eq-low-order-variation-1}.

For the second inequality, since $C\ge 4d/(5Ma^{\mathrm{ax}})$ and $D\ge 5M$,
\begin{equation*}
CD\ge\frac{4d}{5Ma^{\mathrm{ax}}}\cdot 5M=\frac{4d}{a^{\mathrm{ax}}}.
\end{equation*}
Multiplying by $\Theta^{\sharp}\mathfrak{a}_n t^2\ge 0$ and squaring both sides, which are non-negative, gives
\begin{equation*}
\Bigl(\frac{4d\,\Theta^{\sharp}\mathfrak{a}_n t^2}{a^{\mathrm{ax}}}\Bigr)^{2}
\le\bigl(CD\,\Theta^{\sharp}\mathfrak{a}_n t^2\bigr)^{2}.
\end{equation*}
Multiplying by $\mathsf{a}\ge 0$ gives the second inequality. When the two cutoffs have the same lattice,
Lemma~\ref{5.7:lem-one-lattice-vanishing}, whose proof is incomplete at one step, gives $\delta\hat A\equiv 0$ and
$\mathsf{a}=0$; in that case all three members of \eqref{5.7:eq-low-order-variation-1} vanish.
\end{proof}

\begin{proposition}[The relative form of the variation bound]\label{5.7:prop-relative-variation}
Let $y$, $C_y^{\sharp}$, $\mathsf{f}_{\ast}$, $a^{\mathrm{ax}}$, $S^{\mathrm{ax}}_v$, the comb loops
$\lambda_b$, the combs $\gamma_z$, the tent $h_y$, the companion $\hat A=A(h_y,\cdot)$, the length
$\ell_j$, the ratio $u_j=(\eta/\ell_j)^2$ and the configuration $W^{\mathsf{m}}$, the $U$-component
of the box field $\mathfrak{b}'_j$, be as in Definition~\ref{5.7:def-comb-dilation}, on $T_{\eta}$
for $\mathsf{m}=\mathsf{A}$ and on $T_{\eta'}$, $\eta'=\eta/L$, for $\mathsf{m}=\mathsf{B}$. Let
$U^{\mathsf{m}}_{\varsigma}(b)=\exp(\varsigma\log W^{\mathsf{m}}(\lambda_b))$,
$\mathsf{m}\in\{\mathsf{A},\mathsf{B}\}$, be the comb dilation of
Definition~\ref{5.7:def-comb-dilation} on machine $\mathsf{m}$'s own lattice; we write $W$ and
$U_{\varsigma}$ for $W^{\mathsf{A}}$ and $U^{\mathsf{A}}_{\varsigma}$. Assume
$4d\mathsf{f}_{\ast}<a^{\mathrm{ax}}$, as in Lemma~\ref{5.7:lem-comb-dilation-bounds}. For a bond $b$ and
a plaquette $P$ of $T_{\eta}$ in $C_y^{\sharp}$, with base point $x_P$ (the corner at which its loop
starts), put, with $g(z):=W(\gamma_z)$,
\begin{equation*}
\Phi_b:=\frac{\ell_j}{\eta}\,\frac{1}{i}\,\log W(\lambda_b),\qquad
\Phi_P:=\eta^{-2}\log\partial W^{g}(P),
\end{equation*}
the comb potential and the comb curvature of $W$; on $T_{\eta'}$ the same with $\eta\mapsto\eta'$.
Here $\Phi_P$ is normalised without the factor $\ell_j^{2}$ that the comb curvature carries in the
proof of Proposition~\ref{5.7:prop-dilated-action-family}. Write
$\delta\Phi_b:=\Phi^{\mathsf{B}}_{b'}-\Phi^{\mathsf{A}}_b$ and $\delta\Phi_P:=\Phi^{\mathsf{B}}_{P'}-\Phi^{\mathsf{A}}_P$
for every bond $b$ and every plaquette $P\subset C_y^{\sharp}$ of $T_{\eta}$, every fine plaquette
$P'\parallel P$ of $T_{\eta'}$ based in the $\eta$-cell of $x_P$ (we write $P'\sim P$ for these $L^4$
fine plaquettes), and every bond $b'\parallel b$ of
$T_{\eta'}$ within $\eta$ of $b$. Let $\varpi^{\mathrm{nat}}\ge0$ be a number such that, for all these
$b$, $b'$, $P$, $P'$,
\begin{equation}\label{5.7:eq-relative-variation-1}
\ell_j^{2}\,|\delta\Phi_P|\le\varpi^{\mathrm{nat}}\,\mathsf{f}_{\ast},\qquad
|\delta\Phi_b|\le\varpi^{\mathrm{nat}}\,\mathsf{f}_{\ast}
\end{equation}
(the relative closeness, at rate $\varpi^{\mathrm{nat}}$, of the comb curvature and the comb potential
of the two cutoffs on $C_y^{\sharp}$).
Then, with $\mathfrak{n}^{\mathrm{ax}}$ as in Lemma~\ref{5.7:lem-comb-dilation-bounds}\,(iii) and a constant $C$,
\begin{equation*}
|\delta\hat A(U_{\varsigma})|\le C\,(\varpi^{\mathrm{nat}}+L^{-j})\,\mathfrak{n}^{\mathrm{ax}}
\qquad\text{for } |\varsigma|<S^{\mathrm{ax}}_v ,
\end{equation*}
where $\delta\hat A(U_{\varsigma}):=\hat A(U^{\mathsf{B}}_{\varsigma})-\hat A(U^{\mathsf{A}}_{\varsigma})$.
\end{proposition}

\begin{proof}
The proof transposes the two-machine step and the summation step of the proof of
Proposition~\ref{5.7:prop-dilated-action-family}.

\emph{The identity.} Fix $P=b_1b_2b_3^{-1}b_4^{-1}\subset C_y^{\sharp}$ and put $X_b:=\log W(\lambda_b)$. By
the definition of the comb dilation, $\partial U_{\varsigma}(P)$ is the plaquette product of the
$e^{\varsigma X_b}$, that is, the product $G$ of \eqref{5.7:eq-dilated-action-family-a} with the
exponents of the fixed factors set to $0$ and the exponents multiplied by the parameter set to $X_b$.
Then $G(0)=\mathbb{1}$, so $\log G(0)=0$,
and $R_3(0)=0$, since every term of order at
least three vanishes at $0$. Hence \eqref{5.7:eq-dilated-action-family-a} reads
\begin{equation}\label{5.7:eq-relative-variation-2}
\log\partial U_{\varsigma}(P)=\varsigma\log\partial U_1(P)+(\varsigma^2-\varsigma)\,\mathsf{C}(X,X)
+\bigl[R_3(\varsigma)-\varsigma R_3(1)\bigr].
\end{equation}
By Lemma~\ref{5.7:lem-comb-dilation-bounds}\,(iv), $U_1=W^{g}$, so that $\log\partial U_1(P)=\eta^2\Phi_P$.
Moreover $X_b=i\eta\ell_j^{-1}\Phi_b$, whence $\ell_j^2\eta^{-2}\mathsf{C}(X,X)=-\mathsf{C}(\Phi,\Phi)$ by
bilinearity. With $Z_{\varsigma}(P):=\eta^{-2}\log\partial U_{\varsigma}(P)$, identity
\eqref{5.7:eq-relative-variation-2} becomes
\begin{equation*}
\ell_j^2 Z_{\varsigma}(P)=\varsigma\,\ell_j^2\Phi_P-(\varsigma^2-\varsigma)\,\mathsf{C}(\Phi,\Phi)
+\ell_j^2\eta^{-2}\bigl[R_3(\varsigma)-\varsigma R_3(1)\bigr].
\end{equation*}
The same identity holds on $T_{\eta'}$ at $P'$, with bonds $b'_i\parallel b_i$ within $\eta$.

\emph{Sizes.} By Lemma~\ref{5.7:lem-comb-dilation-bounds}\,(ii) and
$\sup_{C_y^{\sharp}}\mathsf{b}_{\ast}=4d\mathsf{f}_{\ast}$ (Definition~\ref{5.7:def-comb-dilation}), we have
$|X_b|\le\bar x:=\tfrac12 d\mathsf{f}_{\ast}u_j^{1/2}$ and $|\Phi_b|\le\tfrac12 d\mathsf{f}_{\ast}$.
Since $|\varsigma|<S^{\mathrm{ax}}_v=a^{\mathrm{ax}}/(4d\mathsf{f}_{\ast})$, it follows that
$|\varsigma|\bar x<\tfrac18a^{\mathrm{ax}}u_j^{1/2}$; also $\bar x\le\tfrac18a^{\mathrm{ax}}u_j^{1/2}\le 1/16$.
The third-order bound $|R_3(\varsigma)|\le C_3(|\varsigma|\bar x)^3$ for $|\varsigma|\bar x\le1/16$, with
the constant $C_3$ used in the proof of Lemma~\ref{5.7:lem-comb-dilation-bounds}\,(iii), gives
\begin{equation*}
\ell_j^2\eta^{-2}\bigl(|R_3(\varsigma)|+|\varsigma||R_3(1)|\bigr)
\le u_j^{-1}C_3\bigl((|\varsigma|\bar x)^3+|\varsigma|\bar x\,\bar x^2\bigr)
\le 2C_3\,(a^{\mathrm{ax}}/8)^3\,u_j^{1/2}.
\end{equation*}
In $\mathsf{B}$ the same holds with $u_j$ replaced by $(\eta'/\ell_j)^2\le u_j$.

\emph{Two machines.} By \eqref{5.7:eq-relative-variation-1},
$|\varsigma|\,\ell_j^2|\delta\Phi_P|\le S^{\mathrm{ax}}_v\varpi^{\mathrm{nat}}\mathsf{f}_{\ast}
=\varpi^{\mathrm{nat}}a^{\mathrm{ax}}/(4d)$. Put $S:=\max\{1,|\varsigma|\}$; then $S<S^{\mathrm{ax}}_v$,
because $S^{\mathrm{ax}}_v>1$, and $|\varsigma^2-\varsigma|\le 2S^2$. We have
$\delta\mathsf{C}(\Phi,\Phi)=\mathsf{C}(\delta\Phi,\Phi^{\mathsf{B}})+\mathsf{C}(\Phi^{\mathsf{A}},\delta\Phi)$
and $|\mathsf{C}(A,B)|\le6\sup|A|\sup|B|$. Hence
\begin{equation*}
|\varsigma^2-\varsigma|\,|\delta\mathsf{C}(\Phi,\Phi)|
\le 2S^2\cdot 6\cdot\varpi^{\mathrm{nat}}\mathsf{f}_{\ast}\cdot d\mathsf{f}_{\ast}
\le \frac{3}{4d}\,\varpi^{\mathrm{nat}}(a^{\mathrm{ax}})^2 .
\end{equation*}
Adding the remainders of both machines yields, with a constant $C_Z$,
$|\delta(\ell_j^2Z_{\varsigma})|\le C_Z\,a^{\mathrm{ax}}(\varpi^{\mathrm{nat}}+u_j^{1/2})$.

\emph{Sum.} Since $e^{-Z}(e^{Z}-\mathbb{1})^2=e^{Z}-2\cdot\mathbb{1}+e^{-Z}$, we have
$a(e^{Z})=-\operatorname{tr}(\cosh Z-\mathbb{1})=-\tfrac12\operatorname{tr}Z^2+O(|Z|^4)$. Moreover
$L^4\eta'^4=\eta^4$ (here $d=4$). Write $a(P):=a(\partial U^{\mathsf{A}}_{\varsigma}(P))$ and
$a(P'):=a(\partial U^{\mathsf{B}}_{\varsigma}(P'))$, so that
\begin{equation*}
\hat A(U^{\mathsf{A}}_{\varsigma})=\sum_P h_y(x_P)\,a(P),\qquad
\hat A(U^{\mathsf{B}}_{\varsigma})=\sum_P\sum_{P'\sim P}h_y(x_{P'})\,a(P');
\end{equation*}
the inner sum runs, for each $P$, over the $L^4$ fine plaquettes $P'\sim P$ of the hypothesis.
Writing each term as
\begin{equation*}
h_y(x_{P'})a(P')-L^{-4}h_y(x_P)a(P)
=\bigl(h_y(x_{P'})-h_y(x_P)\bigr)a(P')+h_y(x_P)\bigl(a(P')-L^{-4}a(P)\bigr),
\end{equation*}
we obtain
\begin{equation*}
|\delta\hat A(U_{\varsigma})|\le\sum_P L^{-4}\sum_{P'\sim P}
C\,\mathfrak{c}_{\mathrm{tr}}\bigl\{|h_y(x_{P'})-h_y(x_P)|\,u_j^2\hat z^2
+h_y(x_P)\,u_j^2\,\hat z\,|\delta(\ell_j^2Z_{\varsigma})|\bigr\}+Q,
\end{equation*}
with $\hat z:=\sup\ell_j^2|Z_{\varsigma}|$ over both machines and $Q$ the sum of the quartic remainders
of both machines. By
Lemma~\ref{5.7:lem-comb-dilation-bounds}\,(iii), $|\partial U_{\varsigma}(P)-\mathbb{1}|\le0.3a^{\mathrm{ax}}u_j\le\tfrac12$,
and $|\log(\mathbb{1}+A)|\le2|A|$ for $|A|\le\tfrac12$; hence $\hat z\le0.6\,a^{\mathrm{ax}}$. Since $h_y$ is the
tent in $\ell_j$-units, $|h_y(x_{P'})-h_y(x_P)|\le\eta\sup|\nabla h_y|\le Cu_j^{1/2}$. The number of
plaquettes in $C_y^{\sharp}$ is at most $Cu_j^{-2}$, and $\sum_Ph_y(x_P)u_j^2=\mathsf{p}_d$. Since
$|\eta^2Z_{\varsigma}(P)|\le u_j\hat z$ and $|\eta'^2Z^{\mathsf{B}}_{\varsigma}(P')|\le L^{-2}u_j\hat z$, the
quartic remainder of a plaquette of $T_{\eta}$ is at most $C_4u_j^4\hat z^4$, and that of a plaquette of
$T_{\eta'}$ at most $C_4L^{-8}u_j^4\hat z^4$, for a constant $C_4$; with $h_y\le1$, $L^4$ fine plaquettes
per $P$ and at most $Cu_j^{-2}$ plaquettes $P$, this gives
\begin{equation*}
Q\le Cu_j^{2}\hat z^{4}\le Cu_j^{2}(a^{\mathrm{ax}})^4\le Cu_j^{1/2}(a^{\mathrm{ax}})^2,
\end{equation*}
because $u_j\le1$ and $a^{\mathrm{ax}}\le1$. Therefore
\begin{equation*}
\begin{split}
|\delta\hat A(U_{\varsigma})|
&\le C\,\mathfrak{c}_{\mathrm{tr}}(a^{\mathrm{ax}})^2\bigl(\varpi^{\mathrm{nat}}+u_j^{1/2}\bigr)(\mathsf{p}_d+1)
+Cu_j^{1/2}(a^{\mathrm{ax}})^2\\
&\le C\bigl(\varpi^{\mathrm{nat}}+u_j^{1/2}\bigr)\mathfrak{n}^{\mathrm{ax}},
\end{split}
\end{equation*}
since $\mathfrak{n}^{\mathrm{ax}}=0.1\,\mathfrak{c}_{\mathrm{tr}}(a^{\mathrm{ax}})^2\mathsf{p}_d$ and
$\mathfrak{c}_{\mathrm{tr}}$, $\mathsf{p}_d$ are fixed positive constants. Since $\ell_j=L^j\eta$, we have
$u_j^{1/2}=\eta/\ell_j=L^{-j}$, and this is the bound of the statement; the term $u_j^{1/2}$ is the same
term that the two-machine step of Proposition~\ref{5.7:prop-dilated-action-family} writes $O(L^{-i})$,
$i$ there denoting the level.
\end{proof}

\begin{remark}
The hypothesis \eqref{5.7:eq-relative-variation-1}, the relative closeness of the comb curvature and
the comb potential of the two cutoffs on $C_y^{\sharp}$, is not proved here; nothing above supplies it.
The hypothesis, with its number $\varpi^{\mathrm{nat}}$, belongs to the setting whose parameters enter
Theorem~\ref{5.7:thm-low-order-variation}, and the present proposition leaves it unchanged.
Its right side is measured against the natural size $\mathsf{f}_{\ast}=\Theta^{\sharp}\mathfrak{a}_nt^2$,
the majorant proportional to $\mathfrak{a}_nt^2$, whereas clause~(f) of
Lemma~\ref{5.7:lem-two-map-interpolation} bounds the mismatch only relative to the radius $a_{\Lambda}$.
It has the form of clause~(a) of the multiscale closeness bound with $\mathfrak{a}_n$ replaced by
$C_{\mathrm{cl}}\vartheta^{j}\mathfrak{a}_n$, where $C_{\mathrm{cl}}$ is a constant of that bound; since
$\mathsf{f}_{\ast}=\Theta^{\sharp}\mathfrak{a}_nt^2$, this is \eqref{5.7:eq-relative-variation-1} with
$\varpi^{\mathrm{nat}}=C_{\mathrm{cl}}\vartheta^{j}$.
At real data the relative form is the localised regularity bound for the potential $A^{\natural}$ of the
carried frame,
\begin{equation*}
|A^{\natural}|_y\le C_A\,\mathsf{S}(y)\,L^{-\theta\lambda(y)},
\end{equation*}
where $\mathsf{S}$ is the smeared regularity profile and $C_A$ and $\theta\lambda(y)$ are the constant and
the exponent of that bound. At complex data the relative form rests instead on homogeneity, which gives
the index factor $\vartheta^{j}$ in full; that argument is not carried out here.
\end{remark}

In the lemma below and its proof, $G=SU(N)$ and $\mathfrak{g}^{c}$ is its complexified Lie algebra. The lemma applies to two configuration spaces defined by Ba{\l}aban, which we now recall. In both, Ba{\l}aban writes the complex factor of a configuration as $U'$; we write it $U_{\mathrm{c}}$.

\emph{First space.} This is the space $U^{c}_j(X,\alpha_0,\alpha_1)$ of \cite[(1.11)--(1.14)]{B-CMP109}, on a lattice of spacing $\xi$. Its configurations are pairs $(\mathbf{U},\mathbf{J})$ with the following properties.
\begin{itemize}
\item $\mathbf{U}=U_{\mathrm{c}}U$, where $U$ is $G$-valued with $|\partial U-1|<\alpha_0\xi^2$ on $X$.
\item For each cube $\square\subset X$ of size $O(1)LM$ there is a $G$-valued gauge transformation $u$ on $\square$ with $U^u=\exp i\xi A$ and $|A|,|\nabla^{\xi}A|<O(1)LMB\alpha_0$ on $\square$.
\item $U_{\mathrm{c}}=\exp i\xi A'$, where $A'$ is $\mathfrak{g}^{c}$-valued with $|A'|,|\nabla^{\xi}_{U}A'|<\alpha_1$ on $X$.
\item $\mathbf{U}$ and $\mathbf{J}$ obey $|\partial\mathbf{U}-1|<\alpha_0\xi^2$ and a bound on $|\mathbf{J}|$ on $X$.
\end{itemize}

\emph{Second space.} This is the space $\tilde{U}^{(n)c}_k(X,\tilde{\alpha}_0,\tilde{\alpha}_1)$ of \cite[(1.64)--(1.69), pp.~190--191]{B-CMP122a}, on a lattice of spacing $\eta$. Write $\ell_m:=L^m\eta$ for the length unit of level $m$. Its configurations $(\mathbb{U},\mathbb{J})$ have the following properties.
\begin{itemize}
\item $\mathbb{U}=U_{\mathrm{c}}U$, where $U$ is $G$-valued and $U_{\mathrm{c}}=\exp i\eta A'$.
\item On the stratum of level $m$ one has $\ell_m^3|\mathbb{J}|<\mathsf{w}_m\alpha_{0,m}$, $\ell_m|A'|<\mathsf{w}_m\alpha_{1,m}$ and $\ell_m^2|\nabla^{\eta}_{U}A'|<\mathsf{w}_m\alpha_{1,m}$ \cite[(1.64)]{B-CMP122a}.
\item For every cube $\square$ of the family of cubes of the definition there is a gauge transformation $u$ defined on $\square\cap X$ such that $U^u=\exp i\eta A$.
\item On a cube $\square$ of side $CML^m\eta$ one has $\ell_m|A|,\ \ell_m^2|\nabla^{\eta}A|<\mathsf{w}_mBCM\alpha_{0,m}$ \cite[(1.65), (1.67), (1.69)]{B-CMP122a}.
\end{itemize}
Here $\mathsf{w}_m$ is the numerical factor printed in the respective clause of \cite[(1.64)--(1.69)]{B-CMP122a}, a power of one of the constants of that paper. It equals $1$ in \cite[(1.69)]{B-CMP122a}, the clause for cubes $\square\subset\Omega_m$ with $\square\cap\Omega^c_{m+1}\neq\emptyset$ or $\square\subset\Omega_k$, where $\Omega_m$, $\Omega_{m+1}$ and $\Omega_k$ are the regions so denoted in \cite[pp.~190--191]{B-CMP122a}. As in \cite[pp.~190--191]{B-CMP122a}, one may assume $C\le 2$; there $B=B_3$, the constant so named in that paper.

For the first space the statement below is read with $\eta=\xi$, with $\mathsf{w}_n=1$, $\alpha_{0,n}=\alpha_0$ and $\alpha_{1,n}=\alpha_1$, and with the constant $O(1)LMB$ of \cite[(1.12)]{B-CMP109} in place of $BCM$; its cubes are of size $O(1)LM$ in $\xi$-units. For a bond $b$ from $b_-$ to $b_+$, gauge transformations act by $U^u(b)=u(b_-)U(b)u(b_+)^{-1}$. We write $R(u)\mathbb{J}$ for the transform of $\mathbb{J}$ under $u$.

\begin{lemma}[Axial gauge on the cubes of Ba{\l}aban's configuration spaces]
\label{5.7:lem-axial-gauge-cubes}
Let $(\mathbb{U},\mathbb{J})$ belong to one of the two spaces above over a region $Y$ (that is, with $X=Y$). Let $\square$ be a cube of its family at level $n$, of side at most $CM\le 2M$ in units of level $n$ (for the first space, of size $O(1)LM$ in $\xi$-units). Set
\begin{equation}\label{5.7:eq-axial-gauge-cubes-a}
\begin{gathered}
\mathsf{v}_n:=\mathsf{w}_n\bigl(BCM\alpha_{0,n}+\alpha_{1,n}+\alpha_{0,n}\bigr),\qquad
\mathfrak{a}^{\partial}_n:=\mathsf{v}_n(1+4\mathsf{v}_n),\\
\text{and assume}\qquad \sup_n\mathsf{v}_n\le\tfrac14 .
\end{gathered}
\end{equation}
Then on $\square\cap Y$ the pair $(\mathbb{U},\mathbb{J})$ is gauge equivalent, under a $G$-valued gauge transformation $u$, to $(e^{i\eta A_{\square}},R(u)\mathbb{J})$. Here $A_{\square}$ is $\mathfrak{g}^{c}$-valued and satisfies
\begin{equation}\label{5.7:eq-axial-gauge-cubes-1}
\ell_n|A_{\square}|,\quad \ell_n^2|\nabla^{\eta}A_{\square}|,\quad \ell_n^3|R(u)\mathbb{J}|
\ \le\ \mathfrak{a}^{\partial}_n .
\end{equation}
\end{lemma}

\begin{proof}
Let $u$ be the $G$-valued gauge transformation that the definition of the space attaches to $\square$, and let $A$ be the potential with $U^u=e^{i\eta A}$ on $\square\cap Y$.

\emph{The bounds at level $n$.} By the defining clauses at level $n$, on $\square\cap Y$ we have
\begin{equation}\label{5.7:eq-axial-gauge-cubes-2}
\begin{gathered}
\ell_n|A'|,\ \ell_n^2|\nabla^{\eta}_{U}A'|<\mathsf{w}_n\alpha_{1,n},\qquad
\ell_n^3|\mathbb{J}|<\mathsf{w}_n\alpha_{0,n},\\
\ell_n|A|,\ \ell_n^2|\nabla^{\eta}A|<\mathsf{w}_nBCM\alpha_{0,n}.
\end{gathered}
\end{equation}
Each right-hand side is at most $\mathsf{v}_n$. Since $\eta\le\ell_n$, it follows that $\eta|A'|,\eta|A|\le\mathsf{v}_n\le\frac14$ by \eqref{5.7:eq-axial-gauge-cubes-a}. The inequality $xy\le\frac14(x+y)^2$, applied with $x=\mathsf{w}_n\alpha_{1,n}$ and $y=\mathsf{w}_nBCM\alpha_{0,n}$, gives
\begin{equation}\label{5.7:eq-axial-gauge-cubes-3}
\ell_n^2|A'|\,|A|\le\tfrac14\mathsf{v}_n^2 .
\end{equation}

\emph{The gauge-transformed link.} For a bond $b$ in $\square\cap Y$ we have
\begin{align*}
\mathbb{U}^u(b)
&=u(b_-)U_{\mathrm{c}}(b)U(b)u(b_+)^{-1}\\
&=u(b_-)e^{i\eta A'(b)}u(b_-)^{-1}\cdot U^u(b)
=e^{i\eta\operatorname{Ad}_{u(b_-)}A'(b)}\,e^{i\eta A(b)}
=:e^{i\eta A_{\square}(b)} .
\end{align*}
Because $u$ is $G$-valued, $|\operatorname{Ad}_u|=1$, and so $|\operatorname{Ad}_uA'|=|A'|$.

\emph{The Baker--Campbell--Hausdorff series.} The two exponents have norm at most $\frac14$ by \eqref{5.7:eq-axial-gauge-cubes-a}, so the series for the logarithm of the product converges. It gives
\begin{equation*}
A_{\square}=\operatorname{Ad}_uA'+A+\mathsf{c},\qquad
|\mathsf{c}|\le\eta|A'|\,|A|\,(1+\cdots),
\end{equation*}
where the bracket collects the commutators of order three and higher. By $\eta\le\ell_n$ and \eqref{5.7:eq-axial-gauge-cubes-3}, $\eta|A'||A|\le\frac14\ell_n^{-1}\mathsf{v}_n^2$. With this, the bound reads $|\mathsf{c}|\le\ell_n^{-1}\mathsf{v}_n^2$. Hence, by \eqref{5.7:eq-axial-gauge-cubes-2},
\begin{equation*}
\ell_n|A_{\square}|\le\mathsf{w}_n\alpha_{1,n}+\mathsf{w}_nBCM\alpha_{0,n}+\mathsf{v}_n^2
\le\mathsf{v}_n+\mathsf{v}_n^2\le\mathfrak{a}^{\partial}_n .
\end{equation*}

\emph{The gradient.} The lattice derivative is covariant:
\begin{equation*}
\operatorname{Ad}_u\nabla^{\eta}_{U}A'=\nabla^{\eta}_{U^u}(\operatorname{Ad}_uA')
=\nabla^{\eta}(\operatorname{Ad}_uA')+O(|A|\,|A'|).
\end{equation*}
The second equality uses $U^u=e^{i\eta A}$ on $\square\cap Y$ and $|\operatorname{Ad}_uA'|=|A'|$. Since $|\operatorname{Ad}_u|=1$, \eqref{5.7:eq-axial-gauge-cubes-2} and \eqref{5.7:eq-axial-gauge-cubes-3} give
\begin{equation*}
\ell_n^2|\nabla^{\eta}(\operatorname{Ad}_uA')|\le\mathsf{w}_n\alpha_{1,n}+2\mathsf{v}_n^2 .
\end{equation*}
In the same way, $\mathsf{c}$ is bilinear in $A'$ and $A$ up to the higher commutators. Its difference quotient is therefore controlled by \eqref{5.7:eq-axial-gauge-cubes-2} and \eqref{5.7:eq-axial-gauge-cubes-3}, which gives $\ell_n^2|\nabla^{\eta}\mathsf{c}|\le2\mathsf{v}_n^2$. Together with $\ell_n^2|\nabla^{\eta}A|<\mathsf{w}_nBCM\alpha_{0,n}$ this yields
\begin{equation*}
\ell_n^2|\nabla^{\eta}A_{\square}|
\le\mathsf{w}_n\alpha_{1,n}+2\mathsf{v}_n^2+\mathsf{w}_nBCM\alpha_{0,n}+2\mathsf{v}_n^2
\le\mathsf{v}_n+4\mathsf{v}_n^2=\mathfrak{a}^{\partial}_n .
\end{equation*}

\emph{The current.} Since $u$ is $G$-valued, $|R(u)\mathbb{J}|=|\mathbb{J}|$. By \eqref{5.7:eq-axial-gauge-cubes-2},
\begin{equation*}
\ell_n^3|R(u)\mathbb{J}|<\mathsf{w}_n\alpha_{0,n}\le\mathsf{v}_n\le\mathfrak{a}^{\partial}_n .
\end{equation*}
This proves \eqref{5.7:eq-axial-gauge-cubes-1}.
\end{proof}

\begin{corollary}[The boundary hypothesis as a membership clause]\label{5.7:cor-boundary-membership}
Read clauses (i)--(iii) of \cite[(1.11)--(1.14)]{B-CMP109} and the clauses
\cite[(1.64)--(1.69)]{B-CMP122a} as the definitions of Ba{\l}aban's configuration spaces
and of their families of cubes. Assume the condition
\eqref{5.7:eq-axial-gauge-cubes-a} of Lemma~\ref{5.7:lem-axial-gauge-cubes}, and assume:
\begin{enumerate}
\item[(a)] the closure-membership clause: every pair $p\in\mathrm{cl}^{\mathrm{pr}}$ at which the
two cutoffs $\mathsf{A}$ and $\mathsf{B}$ are compared, at the image of a read at the site $(L)$
or at a source at the site $(P)$ (Definition~\ref{5.2:def-sites-units}), has each of its two
components in the configuration space of its own cutoff;
\item[(b)] the containment hypothesis: the configurations of the two cutoffs at the pairs
compared in (a) take values in the space
$\tilde{U}''^{c}_k(Y_1,\tilde{\alpha}_0,\tilde{\alpha}_1)$, with $Y_1$ the
region and $k$ the level at which the pair is read; the statements supplying this hypothesis
include the cube clause \cite[(1.69)]{B-CMP122a} together with the lemma bounding the
configurations on the rims read at the site $(L)$;
\item[(c)] the boundary step lemma, in the following form. Let $G=\mathrm{SU}(N)$ be the
structure group and $G^{c}$ its complexification. Let $j\le n$ be levels, $t:=L^{j-n}$, let $X$
be a set and $\square_X$ a cube about a point $y$, of $j$-side $D_X$, with $X$ contained in
$\square_X$ intersected with the domain of the configurations; let $F(X;\cdot)$ be the
functional the lemma estimates, $(1,0)$ the trivial point, and $\mathfrak{n}$, $\mathfrak{r}_j$
its constants at level $j$. If, for some $\mathfrak{a}>0$, a point $P$ on $X$ is, on $\square_X$
intersected with its domain, $G^{c}$-gauge equivalent to a configuration $(e^{i\eta A},\mathbb{J}')$
with
\begin{equation*}
\ell_n|A|,\quad \ell_n^{2}|\nabla A|,\quad \ell_n^{3}|\mathbb{J}'|\;\le\;\mathfrak{a},
\end{equation*}
then
\begin{equation*}
\bigl|F(X;P)-F(X;1,0)\bigr|
\le 2\mathfrak{n}\Bigl(\frac{2D_X t^{2}\,\mathfrak{a}}{\mathfrak{r}_j}\Bigr)^{2};
\end{equation*}
\item[(d)] clause (ii) of the one-cutoff step lemma, of which the boundary step lemma of (c) is
the variant at the rim: it bounds $|F(X;P)-F(X;1,0)|$ by twice a size constant times an
exponential decay factor with Ba{\l}aban's decay constant $\kappa$ as rate (the trivial bound).
\end{enumerate}
Let $(\mathbb{U},\mathbb{J})$ be any point of one of the configuration spaces of \cite{B-CMP109},
\cite{B-CMP122a} just read as definitions, over a region $Y$. Let $j\le n$ be two levels and
$t:=L^{j-n}$, as in Lemma~\ref{5.2:lem-local-gauge-bound}, so that a length $D$ in the units of
level $j$ is $Dt$ in the units of level $n$. For a set $X$, let $\square_X$ be the cube about a
point $y$ on which hypothesis (c) is read for $X$, let $D_X$ be the
$j$-side of $\square_X$, so that $D_Xt$ is its side in the units of level $n$, and suppose
$X\subset\square_X\cap Y$. Let $\square$ be a cube of the family of $(\mathbb{U},\mathbb{J})$ at
level $n$ about $y$, and let $C\le 2$ be the constant of the cube family of
\cite[pp.~190--191]{B-CMP122a}. We say that $X$ lies \emph{within the cut} of the boundary
comparison (relative to $\square$) if $\square_X\subset\square$, and \emph{outside the cut}
otherwise. If $X$ lies within the cut, then $X\subset\square\cap Y$ and
\begin{equation}\label{5.7:eq-boundary-membership-1}
D_X t\le CM\le 2M .
\end{equation}
Moreover, for $X$ within the cut, the hypothesis of (c) holds at
$P:=(\mathbb{U},\mathbb{J})$ on $X$ with $(n,\mathfrak{a}):=(n,\mathfrak{a}^{\partial}_n)$, where
$\mathfrak{a}^{\partial}_n=\mathsf{a}_n(1+4\mathsf{a}_n)$ is the majorant of
Lemma~\ref{5.7:lem-axial-gauge-cubes}: on
$\square_X\cap Y\supset X$ the point $P$ is $G^{c}$-gauge equivalent to a configuration
$(e^{i\eta A},\mathbb{J}')$ with
\begin{equation*}
\ell_n|A|,\quad \ell_n^{2}|\nabla A|,\quad \ell_n^{3}|\mathbb{J}'|\;\le\;\mathfrak{a}^{\partial}_n ,
\end{equation*}
that is, the potential, its lattice gradient and the current are bounded by
$\mathfrak{a}^{\partial}_n$ in the units of the level $n\ge j$. Consequently, by (c) at
$\mathfrak{a}=\mathfrak{a}^{\partial}_n$, with $F(X;\cdot)$, $(1,0)$, $\mathfrak{n}$ and
$\mathfrak{r}_j$ as there,
\begin{equation}\label{5.7:eq-boundary-membership-2}
\bigl|F(X;P)-F(X;1,0)\bigr|
\le 2\mathfrak{n}\Bigl(\frac{2D_X t^{2}\,\mathfrak{a}^{\partial}_n}{\mathfrak{r}_j}\Bigr)^{2}
\end{equation}
for every such $\square$ and every $X$ within the cut.
For $X$ outside the cut, the trivial bound of hypothesis (d) applies. Both conclusions
hold for each component of every pair $p$ as in (a), each in the space of its own cutoff.
\end{corollary}

\begin{proof}
The condition \eqref{5.7:eq-boundary-membership-1} is the cube cut of the boundary comparison:
cubes of $n$-diameter at most $2M$. Its second inequality uses Ba{\l}aban's normalisation
$C\le 2$ of the cube size. He states this in \cite[pp.~190--191]{B-CMP122a},
where the cubes of the family have side $CML^{n}\eta$ at level $n$.

Fix $(\mathbb{U},\mathbb{J})$, $\square$ and $X$ as in the statement, with $X$ within the cut.
Since $\square_X$ and $\square$ are both cubes about $y$, the inclusion $\square_X\subset\square$
says that the $n$-side $D_Xt$ of $\square_X$ is at most the $n$-side of $\square$. For the spaces
of \cite{B-CMP122a}, this $n$-side is $CM\le 2M$ by \cite[(1.69)]{B-CMP122a}. For the spaces of
\cite{B-CMP109}, the cubes of the family are the cubes of \cite[(1.12)]{B-CMP109}, of side
$O(1)LM$ in the units of the lattice spacing $\xi$; these are the cubes to which
Lemma~\ref{5.7:lem-axial-gauge-cubes} applies in those spaces, its hypothesis reading their
$n$-side as at most $CM\le 2M$. In either case
$D_X t\le CM\le 2M$, which is \eqref{5.7:eq-boundary-membership-1}. Moreover
$X\subset\square_X\cap Y\subset\square\cap Y$. Together with the
membership of $(\mathbb{U},\mathbb{J})$ in one of the spaces over $Y$ and the condition
\eqref{5.7:eq-axial-gauge-cubes-a}, these are precisely the hypotheses of
Lemma~\ref{5.7:lem-axial-gauge-cubes}.

That lemma therefore gives a $G$-valued gauge transformation $u$ on $\square\cap Y$. It carries
$(\mathbb{U},\mathbb{J})$ to $(e^{i\eta A_{\square}},R(u)\mathbb{J})$, and
\begin{equation*}
\ell_n|A_{\square}|,\quad \ell_n^{2}|\nabla A_{\square}|,\quad \ell_n^{3}|R(u)\mathbb{J}|
\;\le\;\mathfrak{a}^{\partial}_n \qquad\text{on } \square\cap Y .
\end{equation*}
Since $G\subset G^{c}$, a $G$-gauge equivalence is in particular a $G^{c}$-gauge equivalence.

Restricting $u$ and these bounds to $\square_X\cap Y\subset\square\cap Y$, which contains $X$, we
see that on $\square_X\cap Y$ the configuration is $G^{c}$-gauge equivalent to an exponential
of a potential. The
potential, its lattice gradient and the current are bounded by $\mathfrak{a}^{\partial}_n$ in
the units of the level $n\ge j$. This is the hypothesis of (c)
with $(n,\mathfrak{a}):=(n,\mathfrak{a}^{\partial}_n)$.

The conclusion of (c), at $\mathfrak{a}=\mathfrak{a}^{\partial}_n$, is the bound
\eqref{5.7:eq-boundary-membership-2}. For $X$ outside the cut,
neither Lemma~\ref{5.7:lem-axial-gauge-cubes} nor hypothesis (c) is invoked; hypothesis (d)
supplies the trivial bound instead.

Now let $p\in\mathrm{cl}^{\mathrm{pr}}$ be a pair at which the two cutoffs are compared, at the
image of a read at the site $(L)$ or at a source at the site $(P)$. By the
closure-membership clause (a), which clause (ii) of the closure lemma for the two cutoffs
derives from the containment hypothesis (b), each component of $p$ lies in the configuration
space of its
own cutoff. The argument of the two preceding paragraphs therefore applies to each component
separately, in its own space, at every cube of its family and every $X$ within the cut. This
gives the hypothesis of (c) with
$(n,\mathfrak{a}):=(n,\mathfrak{a}^{\partial}_n)$, the bound
\eqref{5.7:eq-boundary-membership-2} within the cut, and the trivial bound outside it.

Hence the hypothesis of (c) is not an additional
estimate. It is the membership of the pair in the spaces of its two cutoffs, read at one of the
clauses that define those spaces: the cube clause \cite[(1.69)]{B-CMP122a}.
\end{proof}

\begin{remark}[Asserted and constructed gauges]\label{5.7:rem-asserted-gauges}
The gauge $u$ of Lemma~\ref{5.7:lem-axial-gauge-cubes} exists, but it is asserted, not constructed:
the lemma does not single out a canonical one. Consider the
two cutoffs $\mathsf{A}$, of spacing $\eta$, and $\mathsf{B}$, of spacing $\eta'=\eta/L$, with
$\mathsf{m}\in\{\mathsf{A},\mathsf{B}\}$ denoting either of them. The one-parameter family of pairs
of configurations, one configuration for each cutoff $\mathsf{m}$, that is built in the gauge $u$,
with parameter $\sigma$, therefore depends on a choice made
separately for each cutoff $\mathsf{m}$ and at each point.

This dependence does not enter the two places where the family is used. The interpolation
principle in $\sigma$ for paired configurations, applied to this family, has its hypothesis
verified in Corollary~\ref{5.7:cor-boundary-membership}; it requires
holomorphy in $\sigma$ only, at one fixed pair of configurations of the two cutoffs and for fixed
choices. The closure condition for the move built on $u$ requires that
the move map the reduced class of its parent into the reduced class $\mathcal{K}^{\mathrm{rd}}$ of its
target slot. This condition holds by construction, provided the move is listed over all admissible
couples $(u,A)$ of a gauge and its potential; here $A$ is the axial potential $A_{\square}$ of
Lemma~\ref{5.7:lem-axial-gauge-cubes} in the gauge $u$.

For families of class functions, which read the plaquette variables only, the choice can be
avoided altogether: the comb gauge of the comb dilation (Definition~\ref{5.7:def-comb-dilation}),
which is read from the trivial configuration, needs only the clause of the configuration space
that bounds the plaquette variables, that is the curvature, and not a clause on the gradient.

For families of objects that read the full gradient, the choice cannot be avoided. The ball that
enters the bounds for these families reads the full gradient, and the comb gauge does not bound the
full gradient in terms of the curvature.
\end{remark}

\begin{definition}[Two cutoffs on exactly paired sets: geometry and floors]
\label{5.8:not-paired-sets}
The following notation is fixed for the comparison of two cutoffs, and is used without further
comment below.
Throughout, $p$ denotes a pair of points, one of cutoff $\mathsf{B}$ and one of cutoff
$\mathsf{A}$ (in particular a presented pair, as below), except in the plaquette commutator below,
where it is a plaquette; $\Delta$ without argument is the operator of the quadratic part of
$\mathcal{W}_{\mathbf{U}}$, distinct from the blocks $\Delta(y)$; $A(\cdot)$ with a group-valued
argument is the Wilson action, distinct from the fields $A$; $M(\cdot)$ is the block averaging,
distinct from the constants $M$ and $M_{1}$; $\kappa$, as the argument of the format
$\sharp(\kappa)$, is a profile on sites, not the auxiliary parameter $\kappa$; $d\Sigma$ and
$d\mathbf{Q}$ denote differentials; and the regions $\Omega_{j}$ are unrelated to the
vacuum $\Omega$.

\emph{The two cutoffs and the paired sets.} The coarser cutoff $\mathsf{A}$ has lattice spacing
$\eta$, and the finer cutoff $\mathsf{B}$ has spacing $\eta' = \eta/L$. A quantity attached to
one of them carries the superscript $\mathsf{A}$ or $\mathsf{B}$ where this is needed.
Let $\mathfrak{S}$ be an admissible set of $\mathsf{A}$ and $\mathfrak{S}'$ an admissible set
of $\mathsf{B}$. Write $\Omega_j$ and $\Omega'_j$ for the regions that make up $\mathfrak{S}$
and $\mathfrak{S}'$ respectively. The pair $(\mathfrak{S}',\mathfrak{S})$ is \emph{exactly
paired} if $\Omega'_{j+1} = \Omega_j$ for every $j$. The collar $\Lambda'_{0}$ of
$\mathsf{B}$ is left free. The layers of $\mathfrak{S}$ are $\Lambda_{j}$, and the label is
$\mathfrak{B} = \bigcup_j \Lambda_{j}$.

For a site $x$, $\lambda(x)$ is its level, counted in the indexing of $\mathsf{A}$, and we set
$\lambda(x) := -1$ on $\Lambda'_{0}$. The scaled length is $\ell := L^{\lambda}\eta$, and
$\ell_y := \ell(y) = L^{\lambda(y)}\eta$.

\emph{Distance and blocks.} The distance $d(\cdot,\cdot)$ is Ba{\l}aban's distance
\cite[(2.46)]{B-CMP96}; the bare letter $d = 4$ remains the dimension. By the statement fixing
the geometry of exactly paired sets, on sites
of exactly paired sets there is one distance function $d(\cdot,\cdot)$ and one level function
$\lambda$, common to both cutoffs. The blocks $\Delta(y)$ and $\tilde{\Delta}(y)$ are those of
\cite[p.~397]{B-CMP99b}. For a field $A$, its scaled local size at $y$ is
\begin{equation*}
|A|_y := \ell_y \sup_{\tilde{\Delta}(y)} |A| .
\end{equation*}
With $\delta_{0}$ the fixed decay rate of Ba{\l}aban's construction (so that
$\delta = \delta_{0}/8$), set $\delta_{1} := \delta_{0}/32$.
By the same statement, on exactly paired sets the distance of \cite[(2.46)]{B-CMP96} has the
following two properties.
For all paired sites $x,y$,
\begin{equation}\label{5.8:eq-paired-sets-a}
d(x,y) \ge R M_{1}\bigl(|\lambda(x)-\lambda(y)| - 1\bigr)_{+} ,
\end{equation}
and for every site $x$,
\begin{equation}\label{5.8:eq-paired-sets-b}
\sum_y e^{-\delta_{1} d(x,y)/2} \le c_{1} ,
\end{equation}
with $c_{1}$ a constant independent of $x$.
The constants satisfy the floor
\begin{equation}\label{5.8:eq-paired-sets-c}
\delta_{1} R M_{1} \ge 16 \log L .
\end{equation}
This floor is \cite[(2.59)]{B-CMP96} taken at $\alpha = 1/64$, for $L \ge 5$.
Here $M_{1}$ is Ba{\l}aban's auxiliary parameter, and $R$ is Ba{\l}aban's constant entering
\cite[(2.46), (2.59)]{B-CMP96}, not the ball radius of the observables.

\emph{H\"older indices and the smear.} Fix $\theta \in (3/5, 1)$, and set
\begin{equation*}
\beta := \frac{1+\theta}{2}, \qquad \varepsilon := \frac{1-\theta}{2}, \qquad
\theta_{2} := \frac{5\theta - 3}{8} .
\end{equation*}
For a function $f$ on sites, its $\delta_{1}$-smear is
\begin{equation*}
f^{\sim} := \sum_y e^{-\delta_{1} d(\cdot,y)} f(y) .
\end{equation*}

\emph{The mismatch potential and the blocked-increment map.} Write $M(\cdot)$ for the block
averaging that carries a background of $\mathsf{B}$ to one of $\mathsf{A}$. Let $p$ be a pair
whose point of $\mathsf{B}$ has background $\mathbf{U}'$ and whose point of $\mathsf{A}$ has
background $\mathbf{U}$. Its mismatch
potential is
\begin{equation*}
A_{p} := (i\eta)^{-1} \log\bigl[M(\mathbf{U}')\,\mathbf{U}^{-1}\bigr] ,
\end{equation*}
so that, exponentiating, $M(\mathbf{U}') = e^{i\eta A_{p}}\mathbf{U}$.
It is read in the frame of the pair $p$, that is, in a gauge satisfying the three conditions of
the lemma fixing the frame of a presented pair.
For a $\mathsf{B}$-field $h$ and a $\mathsf{B}$-background $V$, the blocked-increment map is
\begin{equation*}
Y[h;V] := (i\eta)^{-1}
\log\bigl[M(e^{i\eta' h}V)\,M(V)^{-1}\bigr] .
\end{equation*}
By the lemma on the blocked-increment map, this map is local and analytic and, with $d = 4$ the
dimension, it satisfies
\begin{equation*}
|Y[h;V]| \le 34\, d\, \sup|h| ,
\end{equation*}
together with the Lipschitz estimate for $Y[\cdot\,;V]$ stated in that lemma.
The lemma on the averaged background $W := M(\mathbf{U}')$, in particular its second assertion,
is used as stated there.

\emph{The format $\sharp(\kappa)$.} Let $\kappa$ be a function on sites. A linear functional
$\mathfrak{r}$ of coarse fields is of format $\sharp(\kappa)$ if it can be written as
\begin{equation*}
\mathfrak{r}[u] = \langle \nabla u, s_{1}\rangle + \langle u, s_{0}\rangle ,
\end{equation*}
with $\ell^{2}|s_{1}| \le \kappa(y)$ and $\ell^{3}|s_{0}| \le \kappa(y)$ on
$\tilde{\Delta}(y)$, for every $y$. In dimension four the pairings of data with coarse fields
are dimensionless.

The next three paragraphs recall, without proof, the objects and conclusions of the
constrained-critical-point statement, of its chart, and of the definition of the regularity
classes.

\emph{The slot.} For a background $\mathbf{U}$ of
$\mathsf{A}$, the nonlinear averaging
constraint $\mathbf{Q}_{\mathbf{U}}$ splits as
\begin{equation*}
\mathbf{Q}_{\mathbf{U}}(\eta A) = (\ell Q) A + C_{\mathbf{U}}(A) ,
\end{equation*}
where $C_{\mathbf{U}}$ is its nonlinear part. The action functional is
\begin{align*}
\mathcal{W}_{\mathbf{U}}(X)
&:= A(e^{i\eta X}\mathbf{U}) - A(\mathbf{U})
 - \langle D_{\mathbf{U}}X, \mathfrak{F}_{\mathbf{U}}\rangle \\
&= \tfrac12 \langle X, \Delta X\rangle + V_{0}(X) ,
\end{align*}
with $\mathfrak{F}_{\mathbf{U}}$ the curvature of $\mathbf{U}$. Here $\Delta$ is the operator of
the quadratic part, and $V_{0}$ is the remainder. Set $T := \Delta + D R D^{*}$.
Here $D_{\mathbf{U}}$ is the covariant derivative in the background $\mathbf{U}$; $D$ and
$D^{*}$ are the derivative operator of the constrained-critical-point statement, written there
without a background subscript, and its
adjoint; the operator $R$ is that
of the gauge component $R D^{*} X$, distinct from the constant $R$ of
\eqref{5.8:eq-paired-sets-a}.

For a complex data vector $\mathsf{m}$, arbitrary and not required to come from a current, put
\begin{equation*}
\Sigma_{\mathbf{U},\mathsf{m}}(X) := \mathcal{W}_{\mathbf{U}}(X)
 + \tfrac12 \|R D^{*} X\|^{2} - \langle C_{\mathbf{U}}(X), \mathsf{m}\rangle .
\end{equation*}
For a prescribed value $B$ of the averaging constraint (the italic $B$ is unrelated to the
cutoff $\mathsf{B}$), $\mathfrak{X}(B;\mathbf{U},\mathsf{m})$ denotes the one critical point of
$\Sigma_{\mathbf{U},\mathsf{m}}$ on the set $\{\mathbf{Q}_{\mathbf{U}}(\eta X) = B\}$ that lies
in the ball $\ell|X| \le \varepsilon_{3}$, $\ell^{2}|\nabla X| \le \varepsilon_{3}$, with
$\varepsilon_{3}$ the radius of the constrained-critical-point statement. At this
point the Euler--Lagrange equation holds with a Lagrange multiplier $\mu_{B}$:
\begin{equation*}
d\Sigma(X)[u] = \langle d\mathbf{Q}(\eta X) u, \mu_{B}\rangle
\quad\text{for all } u .
\end{equation*}
For a point $P = (\mathbf{U},\mathbf{J})$, that is, a background of $\mathsf{A}$ together with a
current, the response on the label is
\begin{equation*}
\mathbb{H}(B;P) = \mathfrak{X}(B;\mathbf{U},\mathsf{m}_{P}),
\qquad \mathsf{m}_{P} := H_{\mathbf{U}}^{\dagger}\mathbf{J} ,
\end{equation*}
where $H_{\mathbf{U}}$ is the operator of the constrained-critical-point statement.

\emph{The chart.} The chart is
\begin{equation*}
\mathsf{A}_{0}(A') := A' - H_{0} D_{0}(A') ,
\end{equation*}
where the constraint function $D_{0}$ is determined by $D_{0}(A') = C(\mathsf{A}_{0}A')$, with
$C$ the block function of the averaging constraint. The
chart satisfies
\begin{equation*}
\mathbf{Q}(\eta\mathsf{A}_{0}A') = (\ell Q)A', \qquad
T H_{0} = (\ell Q)^{\dagger}\Theta ,
\end{equation*}
with $H_{0}$ and $\Theta$ Ba{\l}aban's operators of the chart. In the chart,
\begin{equation*}
\Sigma(\mathsf{A}_{0}A') = \tfrac12\langle A', T A'\rangle + \mathcal{N}(A') ,
\end{equation*}
where the nonlinear part is
\begin{equation*}
\begin{split}
\mathcal{N} := {}& -\langle (\ell Q)A', \Theta D_{0}\rangle
 + \tfrac12 \langle D_{0}, \Theta D_{0}\rangle \\
&- \langle D_{0}, \mathsf{m}\rangle + V_{0}(\mathsf{A}_{0}A') .
\end{split}
\end{equation*}
The remainder $V_{0}$ splits as $V_{0}(X) = \langle DX, \mathfrak{c}(X)\rangle + V_{0}'(X)$,
with $V_{0}'$ the further remainder of the constrained-critical-point statement.
Here $\mathfrak{c}$ is the plaquette commutator
\begin{equation*}
\mathfrak{c}(X)(p) := \tfrac12 \sum_{b_{1}<b_{2}} i\,[X'(b_{1}), X'(b_{2})] ,
\end{equation*}
with $p$ a plaquette, $b_{1} < b_{2}$ its bonds, and $X'$ as in \cite[(39)--(40)]{B-CMP102b}.
The constrained Green operator is $\mathcal{G}_{0} := G_{0}P_{0}^{*}$, with $G_{0}$ and $P_{0}$
the operators of the constrained-critical-point statement. It satisfies
$\mathcal{G}_{0}\mathfrak{f} \in \ker(\ell Q)$ and
\begin{equation*}
\langle v, T\mathcal{G}_{0}\mathfrak{f}\rangle = \mathfrak{f}[v]
\quad\text{for } v \in \ker(\ell Q) .
\end{equation*}

\emph{Classes and tallies.} Let $\mathfrak{r} > 0$.
The regularity class
$\mathsf{C}_{\beta}(\mathfrak{r})$ consists of the real points $(\mathbf{U},\mathbf{J})$ with the
following two properties. First, on every $\tilde{\Delta}(y)$ some gauge gives
$\mathbf{U} = e^{i\eta A}$ with
\begin{equation*}
\ell|A| \le \mathfrak{r}, \qquad \ell^{2}|\nabla A| \le \mathfrak{r}, \qquad
\ell^{2+\beta}[\nabla A]_{\beta} \le \mathfrak{r} ,
\end{equation*}
where $[\cdot]_{\beta}$ is the H\"older seminorm of exponent $\beta$. Second,
$\ell^{3}|\mathbf{J}| \le \mathfrak{r}$. The increment class
$\operatorname{Inc}_{\beta}(\mathsf{r})$, with $\mathsf{r}$ its size parameter, and the
presented points and pairs are used as defined
with these classes.

For a presented pair $p$, let $p_{0}, p_{1}, \dots$ be the successive pairs of its presentation,
as defined with the presented pairs in the definition of the regularity classes.
Its mismatch tally is
\begin{equation*}
\mathsf{T}_{p}(x) := |A_{p_{0}}|_{x} + \sum_{s} |A_{p_{s}} - A_{p_{s-1}}|_{x} .
\end{equation*}
Its smeared currencies are
\begin{equation*}
\Psi_{p} := \mathsf{T}_{p}^{\sim} + \mathfrak{r}_{0} L^{-\theta\lambda}, \qquad
\Psi^{B}_{p} := \mathsf{T}_{p}^{\sim} + L^{-\theta\lambda} ,
\end{equation*}
with $\mathfrak{r}_{0}$ the fixed constant carried in the definition of the smeared currencies
in the definition of the regularity classes.
Entries and their moduli are used as defined with the presented pairs.
\end{definition}

\begin{lemma}[Linearised difference of the two kernels]\label{5.8:lem-kernel-difference}
Let $P=(\mathbf{U},\mathsf{m})$ be a real base of the constrained variational problem of the sourced
slot, with background field $\mathbf{U}$, data vector $\mathsf{m}$ and functional
$\Sigma_{\mathbf{U},\mathsf{m}}$. Write $T_{\mathbf{U}}$ for the operator on fields,
$C^{(2)}_{\mathbf{U}}$ for the quadratic map from fields to data, and $\ell Q_{\mathbf{U}}$ for the
linear constraint map from fields to block data, which the sourced slot attaches to the
background $\mathbf{U}$. Define $\hat{T}_{P}:=T_{\mathbf{U}}-\Delta^{(2)}_{\mathsf{m}}$, where
$\Delta^{(2)}_{\mathsf{m}}$ is given by
\begin{equation}\label{5.8:eq-kernel-difference-1}
\langle X,\Delta^{(2)}_{\mathsf{m}}X\rangle:=2\langle C^{(2)}_{\mathbf{U}}(X),\mathsf{m}\rangle
\end{equation}
(the quadratic part of $\Sigma_{\mathbf{U},\mathsf{m}}$).
For a datum $b$, let $K_{P}b$ be the stationary point of
$\tfrac12\langle X,\hat{T}_{P}X\rangle$ on the set $\{X:(\ell Q_{\mathbf{U}})X=b\}$; to first order
in the datum the response is $\mathbb{H}(b;P)=K_{P}b$. Let $\mathcal{G}_{P}$ be the constrained
Green operator: for a linear functional $\mathfrak{f}$ on fields, $\mathcal{G}_{P}\mathfrak{f}$ is the
element of $\ker(\ell Q_{\mathbf{U}})$ such that
\begin{equation}\label{5.8:eq-kernel-difference-2}
\langle v,\hat{T}_{P}\mathcal{G}_{P}\mathfrak{f}\rangle=\mathfrak{f}[v]
\qquad\text{for all } v\in\ker(\ell Q_{\mathbf{U}}).
\end{equation}
Let $\mu_{P}b$ be the multiplier, defined by
\begin{equation}\label{5.8:eq-kernel-difference-3}
\langle v,\hat{T}_{P}K_{P}b\rangle=\langle(\ell Q_{\mathbf{U}})v,\mu_{P}b\rangle
\qquad\text{for all fields } v .
\end{equation}
Let $\tilde{P}=(W,\mathsf{m}')$ be a second base on the same lattice, and let
$\hat{T}_{\tilde{P}}$, $K_{\tilde{P}}$, $\mathcal{G}_{\tilde{P}}$, $\mu_{\tilde{P}}$ be defined in the
same way with $\mathbf{U},\mathsf{m}$ replaced by $W,\mathsf{m}'$. Then, identically,
\begin{equation}\label{5.8:eq-kernel-difference-a}
\begin{split}
K_{\tilde{P}}-K_{P}
&=-K_{\tilde{P}}\bigl(\ell Q_{W}-\ell Q_{\mathbf{U}}\bigr)K_{P}\\
&\quad-\mathcal{G}_{\tilde{P}}\Bigl[\bigl(\hat{T}_{\tilde{P}}-\hat{T}_{P}\bigr)K_{P}
-\bigl(\ell Q_{W}-\ell Q_{\mathbf{U}}\bigr)^{\dagger}\mu_{P}\Bigr],
\end{split}
\end{equation}
where $\dagger$ is the adjoint with respect to the pairings $\langle\cdot,\cdot\rangle$ on fields and
on block data, and a field $Y$ inside the square bracket is read as the functional
$v\mapsto\langle v,Y\rangle$.
\end{lemma}

\begin{proof}
Fix a datum $b$ and put $X:=K_{P}b$, $\tilde{X}:=K_{\tilde{P}}b$ and $E:=\tilde{X}-X$. By the
definition of the two kernels, $(\ell Q_{\mathbf{U}})X=b$ and $(\ell Q_{W})\tilde{X}=b$. Hence
\begin{equation}\label{5.8:eq-kernel-difference-4}
(\ell Q_{W})E=b-(\ell Q_{W})X=(\ell Q_{\mathbf{U}})X-(\ell Q_{W})X
=-\bigl(\ell Q_{W}-\ell Q_{\mathbf{U}}\bigr)X .
\end{equation}

Now let $v\in\ker(\ell Q_{W})$. By \eqref{5.8:eq-kernel-difference-3} written for $\tilde{P}$,
\begin{equation*}
\langle v,\hat{T}_{\tilde{P}}\tilde{X}\rangle=\langle(\ell Q_{W})v,\mu_{\tilde{P}}b\rangle=0 .
\end{equation*}
Writing $\hat{T}_{\tilde{P}}X=(\hat{T}_{\tilde{P}}-\hat{T}_{P})X+\hat{T}_{P}X$ we obtain
\begin{equation*}
\langle v,\hat{T}_{\tilde{P}}E\rangle
=0-\langle v,(\hat{T}_{\tilde{P}}-\hat{T}_{P})X\rangle-\langle v,\hat{T}_{P}X\rangle .
\end{equation*}
By \eqref{5.8:eq-kernel-difference-3} for $P$, and since $(\ell Q_{W})v=0$,
\begin{equation*}
\langle v,\hat{T}_{P}X\rangle
=\langle(\ell Q_{\mathbf{U}})v,\mu_{P}b\rangle
=-\langle(\ell Q_{W}-\ell Q_{\mathbf{U}})v,\mu_{P}b\rangle
=-\langle v,(\ell Q_{W}-\ell Q_{\mathbf{U}})^{\dagger}\mu_{P}b\rangle .
\end{equation*}
Therefore, with the functional
\begin{equation*}
\mathfrak{f}[v]:=\bigl\langle v,\bigl[(\hat{T}_{\tilde{P}}-\hat{T}_{P})K_{P}
-(\ell Q_{W}-\ell Q_{\mathbf{U}})^{\dagger}\mu_{P}\bigr]b\bigr\rangle ,
\end{equation*}
we have
\begin{equation}\label{5.8:eq-kernel-difference-5}
\langle v,\hat{T}_{\tilde{P}}E\rangle=-\mathfrak{f}[v]
\qquad\text{for all } v\in\ker(\ell Q_{W}).
\end{equation}

We split $E=E_{1}+E_{2}$ with
$E_{1}:=-K_{\tilde{P}}(\ell Q_{W}-\ell Q_{\mathbf{U}})X$. Since $K_{\tilde{P}}c$ satisfies the
constraint $(\ell Q_{W})K_{\tilde{P}}c=c$ for every datum $c$,
\eqref{5.8:eq-kernel-difference-4} gives $(\ell Q_{W})E_{1}=(\ell Q_{W})E$, so
$E_{2}\in\ker(\ell Q_{W})$. For $v\in\ker(\ell Q_{W})$,
\eqref{5.8:eq-kernel-difference-3} for $\tilde{P}$ gives
\begin{equation*}
\langle v,\hat{T}_{\tilde{P}}E_{1}\rangle
=-\langle(\ell Q_{W})v,\mu_{\tilde{P}}(\ell Q_{W}-\ell Q_{\mathbf{U}})X\rangle=0 ,
\end{equation*}
and so
\eqref{5.8:eq-kernel-difference-5} holds with $E_{2}$ in place of $E$. Since
$E_{2}\in\ker(\ell Q_{W})$, the defining property \eqref{5.8:eq-kernel-difference-2} of
$\mathcal{G}_{\tilde{P}}$, applied to the functional $-\mathfrak{f}$, identifies
$E_{2}=-\mathcal{G}_{\tilde{P}}\mathfrak{f}$. Hence
\begin{equation*}
K_{\tilde{P}}b-K_{P}b=E_{1}+E_{2}
=-K_{\tilde{P}}(\ell Q_{W}-\ell Q_{\mathbf{U}})K_{P}b-\mathcal{G}_{\tilde{P}}\mathfrak{f}.
\end{equation*}
Since $b$ is arbitrary, this is \eqref{5.8:eq-kernel-difference-a}.
\end{proof}

\begin{remark}
In the application, $\tilde{P}$ is the base of the finer cutoff $\mathsf{B}$ after one exact
blocking, $W=M(\mathbf{U}')=e^{i\eta\mathfrak{a}}\mathbf{U}$, where $M(\mathbf{U}')$ is the image of
$\mathbf{U}'$ under one exact blocking. The three members on the right of
\eqref{5.8:eq-kernel-difference-a} are the first-order parts, respectively, of the first step of the
two-cutoff comparison of the sourced slot, of those residuals of that comparison whose first-order
parts carry the difference $\hat{T}_{\tilde{P}}-\hat{T}_{P}$, and of its residual whose first-order
part carries the multiplier $\mu_{P}$. Each middle factor, $\ell Q_{W}-\ell Q_{\mathbf{U}}$,
$\hat{T}_{\tilde{P}}-\hat{T}_{P}$ and $(\ell Q_{W}-\ell Q_{\mathbf{U}})^{\dagger}$, is local (or
local up to a scalar projection, which is itself a kernel) and is bounded by the mismatch at
the place. Thus the inherited member of the linearised difference has the form
(kernel)$\cdot$(multiplication by the mismatch)$\cdot$(kernel). The remaining, fresh member, which
compares $\mathsf{B}$ with its own blocking (the residuals of the comparison with index
$L^{-\theta\lambda}$), has the same shape with $\mathfrak{a}(y)$ replaced by $L^{-\theta\lambda(y)}$.
\end{remark}

\begin{proposition}[Sharp rate window for transported mismatches]\label{5.8:prop-rate-window}
Let $(\mathfrak{B},d)$ be a finite metric set and $\nu_{\mathrm{sum}}\ge 0$ a rate such that
\begin{equation*}
\sum_{y\in\mathfrak{B}}e^{-\nu_{\mathrm{sum}}d(x,y)}\le c_{\mathrm{sum}}\qquad\text{for every }x\in\mathfrak{B}.
\end{equation*}
Let $\mathcal{G}$, $\tilde{K}$, $K$, $\mu$ be operators on functions on $\mathfrak{B}$ whose kernels are
bounded in absolute value by $Be^{-\kappa_{\mathrm{ker}}d}$; for instance
$|\mathcal{G}(x,y)|\le Be^{-\kappa_{\mathrm{ker}}d(x,y)}$.
Let $V$ be a middle factor whose size is controlled by the mismatch profile $\mathfrak{a}$, a
nonnegative function:
\begin{equation*}
|(VX)(y)|\le\mathfrak{a}(y)\sup_{d(y,y')\le 1}|X(y')|.
\end{equation*}
Let $\mathsf{p}$ and $\Psi$ be nonnegative profiles with rates
$\nu_{\mathsf{p}},\nu_{\Psi}\ge 0$, that is, $\mathsf{p}(y)\le e^{\nu_{\mathsf{p}}d(y,y')}\mathsf{p}(y')$ and
$\Psi(y)\le e^{\nu_{\Psi}d(y,y')}\Psi(y')$ for all $y,y'$. Let $b$ be a datum with $|b|\le\mathsf{p}$, and
assume $\mathfrak{a}\le\Psi$.
\begin{enumerate}
\item[(a)] If $\nu_{\mathsf{p}}+\nu_{\Psi}+\nu_{\mathrm{sum}}\le\kappa_{\mathrm{ker}}$, then for every
$x\in\mathfrak{B}$
\begin{equation}\label{5.8:eq-rate-window-1}
|(\mathcal{G}VKb)(x)|\le e^{\nu_{\mathsf{p}}}B^{2}c_{\mathrm{sum}}^{2}\,\Psi(x)\,\mathsf{p}(x).
\end{equation}
The bound is linear in the mismatch, local in it, and holds at the profile's own rate.
\item[(b)] The window in (a) is sharp. On the chain $\mathbb{Z}_{\ge 0}$ let $T:=-\Delta+m^{2}$, where
$m>0$ is a mass parameter and $-\Delta$ is the nearest-neighbour Laplacian of the chain, so that
\begin{equation*}
(TX)(x)=-X(x+1)+2X(x)-X(x-1)+m^{2}X(x);
\end{equation*}
the datum is placed at $0$, and $\kappa_{\star}>0$ is defined by
$2\cosh\kappa_{\star}=2+m^{2}$. Then the response and the Green operator are
\begin{equation*}
K(x)=e^{-\kappa_{\star}x},\qquad
\mathcal{G}(x,z)=\frac{e^{-\kappa_{\star}|x-z|}-e^{-\kappa_{\star}(x+z)}}{2\sinh\kappa_{\star}}.
\end{equation*}
Write $\delta K$ for the change of $K$ when $T$ is replaced by $T+V$.
\begin{enumerate}
\item[(b1)] For $V=\varepsilon\mathbf{1}_{\{z\}}$, to first order in $\varepsilon$ and for $x\ge z$,
\begin{equation}\label{5.8:eq-rate-window-2}
\delta K(x)=-\varepsilon\,(1-e^{-2\kappa_{\star}z})\,(2\sinh\kappa_{\star})^{-1}e^{-\kappa_{\star}x},
\end{equation}
which does not depend on $x-z$. With the local weight $\Psi(x):=\varepsilon e^{-\nu_{\Psi}|x-z|}$ and
$z=1$, the first-order term satisfies $|\delta K(x)|\le C\Psi(x)e^{-\nu_{\mathsf{p}}x}$ for all $x$, with
some constant $C$, if and only if $\nu_{\mathsf{p}}+\nu_{\Psi}\le\kappa_{\star}$.
\item[(b2)] For $V=\varepsilon\mathbf{1}_{[1,x]}$, to first order in $\varepsilon$,
\begin{equation}\label{5.8:eq-rate-window-3}
\delta K(x)=-\varepsilon K(x)\,\frac{x+O(1)}{2\sinh\kappa_{\star}}.
\end{equation}
Relative to the true size $K(x)$, the transported mismatch $\delta K(x)$ therefore grows linearly in
the distance $x$. With the sup
weight ($\nu_{\Psi}=0$, $\Psi\equiv\varepsilon$) the first-order term satisfies
$|\delta K(x)|\le C\Psi(x)e^{-\nu_{\mathsf{p}}x}$ for all $x$
at every $\nu_{\mathsf{p}}<\kappa_{\star}$, with a constant $C$ proportional to
$(\kappa_{\star}-\nu_{\mathsf{p}})^{-1}$,
and this fails at $\nu_{\mathsf{p}}=\kappa_{\star}$.
\end{enumerate}
\end{enumerate}
\end{proposition}

\begin{proof}
(a) We first bound $Kb$ pointwise. For $y'\in\mathfrak{B}$, the kernel bound on $K$ and $|b|\le\mathsf{p}$
give the first inequality below. The second inequality uses the profile property
$\mathsf{p}(c)\le e^{\nu_{\mathsf{p}}d(c,y')}\mathsf{p}(y')$. The third uses
$\kappa_{\mathrm{ker}}-\nu_{\mathsf{p}}\ge\nu_{\mathrm{sum}}+\nu_{\Psi}\ge\nu_{\mathrm{sum}}$,
which gives $e^{-(\kappa_{\mathrm{ker}}-\nu_{\mathsf{p}})d}\le e^{-\nu_{\mathrm{sum}}d}$, together with the
defining bound for $c_{\mathrm{sum}}$:
\begin{equation}\label{5.8:eq-rate-window-4}
\begin{split}
|(Kb)(y')|&\le B\sum_{c}e^{-\kappa_{\mathrm{ker}}d(y',c)}\mathsf{p}(c)\\
&\le B\,\mathsf{p}(y')\sum_{c}e^{-(\kappa_{\mathrm{ker}}-\nu_{\mathsf{p}})d(y',c)}
\le Bc_{\mathrm{sum}}\,\mathsf{p}(y').
\end{split}
\end{equation}
Next we apply $V$. By the hypothesis on $V$, by $\mathfrak{a}\le\Psi$ and by \eqref{5.8:eq-rate-window-4},
\begin{equation*}
|(VKb)(y)|\le\Psi(y)\,Bc_{\mathrm{sum}}\max_{d(y,y')\le 1}\mathsf{p}(y').
\end{equation*}
For $d(y,y')\le 1$ we have
$\mathsf{p}(y')\le e^{\nu_{\mathsf{p}}d(y',y)}\mathsf{p}(y)\le e^{\nu_{\mathsf{p}}}\mathsf{p}(y)$,
since $\nu_{\mathsf{p}}\ge 0$. Hence
\begin{equation}\label{5.8:eq-rate-window-5}
|(VKb)(y)|\le e^{\nu_{\mathsf{p}}}Bc_{\mathrm{sum}}\,\Psi(y)\,\mathsf{p}(y).
\end{equation}
Finally we apply $\mathcal{G}$, using its kernel bound and \eqref{5.8:eq-rate-window-5}. We then use the
profile properties $\Psi(y)\le e^{\nu_{\Psi}d(y,x)}\Psi(x)$ and
$\mathsf{p}(y)\le e^{\nu_{\mathsf{p}}d(y,x)}\mathsf{p}(x)$:
\begin{align*}
|(\mathcal{G}VKb)(x)|
&\le e^{\nu_{\mathsf{p}}}B^{2}c_{\mathrm{sum}}\sum_{y}e^{-\kappa_{\mathrm{ker}}d(x,y)}\Psi(y)\,\mathsf{p}(y)\\
&\le e^{\nu_{\mathsf{p}}}B^{2}c_{\mathrm{sum}}\,\Psi(x)\,\mathsf{p}(x)
\sum_{y}e^{-(\kappa_{\mathrm{ker}}-\nu_{\mathsf{p}}-\nu_{\Psi})d(x,y)}.
\end{align*}
Since $\kappa_{\mathrm{ker}}-\nu_{\mathsf{p}}-\nu_{\Psi}\ge\nu_{\mathrm{sum}}$, the last sum is at most
$c_{\mathrm{sum}}$. This proves
\eqref{5.8:eq-rate-window-1}.

(b) We first check the two kernels. On the chain,
\begin{align*}
(TK)(x)&=-K(x+1)+2K(x)-K(x-1)+m^{2}K(x)\\
&=e^{-\kappa_{\star}x}(2+m^{2}-2\cosh\kappa_{\star})=0,
\end{align*}
and $K(0)=1$; thus $K$ is the decaying solution of $TK=0$ on $x\ge 1$ with $K(0)=1$. The function
$e^{-\kappa_{\star}|x|}/(2\sinh\kappa_{\star})$ of $x\in\mathbb{Z}$ is mapped by $-\Delta+m^{2}$ to
$\mathbf{1}_{\{0\}}$:
away from $0$ the computation is the one just made, and at $0$ one obtains
$(2+m^{2}-2e^{-\kappa_{\star}})/(2\sinh\kappa_{\star})=1$. Let $z\ge 1$. The image term
$e^{-\kappa_{\star}(x+z)}/(2\sinh\kappa_{\star})$ is, as a function of $x$, a multiple of $K$; it has no
kink for $x\ge 1$ and is annihilated by $T$ there. Subtracting it therefore gives a function
$\mathcal{G}(\cdot,z)$ with $T\mathcal{G}(\cdot,z)=\mathbf{1}_{\{z\}}$ on $x\ge 1$ and $\mathcal{G}(0,z)=0$.

We now apply the linearised difference of the two kernels (Lemma~\ref{5.8:lem-kernel-difference}),
that is, \eqref{5.8:eq-kernel-difference-a}, with the averaging unchanged. The two terms carrying the
difference of the linearised averagings vanish, one of them through its adjoint acting on the
multiplier. The Hessian difference is $V$, and $V=O(\varepsilon)$. In the notation of
Lemma~\ref{5.8:lem-kernel-difference}, with $P$ standing for the problem with Hessian $T$ and
$\tilde{P}$ for the problem with Hessian $T+V$, what remains is
$\delta K=-\mathcal{G}_{\tilde{P}}VK$, where $\mathcal{G}_{\tilde{P}}$ is the Green operator of $T+V$.
Replacing $\mathcal{G}_{\tilde{P}}$ by $\mathcal{G}$ costs $O(\varepsilon^{2})$, since
$\mathcal{G}_{\tilde{P}}-\mathcal{G}=-\mathcal{G}_{\tilde{P}}V\mathcal{G}=O(\varepsilon)$ by the
resolvent identity and $V=O(\varepsilon)$, so that
\begin{equation*}
\delta K=-\mathcal{G}VK+O(\varepsilon^{2}).
\end{equation*}
We insert the two explicit kernels into this formula.

(b1) Let $V=\varepsilon\mathbf{1}_{\{z\}}$. Then $\delta K(x)=-\varepsilon\,\mathcal{G}(x,z)K(z)+O(\varepsilon^{2})$.
For $x\ge z$,
\begin{equation*}
\mathcal{G}(x,z)K(z)
=\frac{e^{-\kappa_{\star}(x-z)}-e^{-\kappa_{\star}(x+z)}}{2\sinh\kappa_{\star}}\,e^{-\kappa_{\star}z}
=\frac{(1-e^{-2\kappa_{\star}z})\,e^{-\kappa_{\star}x}}{2\sinh\kappa_{\star}}.
\end{equation*}
This is \eqref{5.8:eq-rate-window-2}, and it is free of $x-z$.

Now take $z=1$. The coefficient $(1-e^{-2\kappa_{\star}})/(2\sinh\kappa_{\star})$ is positive. At $x=0$ we
have $\mathcal{G}(0,1)=0$, so the first-order term vanishes there. For $x\ge 1$ we have
$\Psi(x)=\varepsilon e^{\nu_{\Psi}}e^{-\nu_{\Psi}x}$, so the inequality
$|\delta K(x)|\le C\varepsilon e^{-\nu_{\Psi}(x-1)}e^{-\nu_{\mathsf{p}}x}$ reads
\begin{equation*}
\frac{1-e^{-2\kappa_{\star}}}{2\sinh\kappa_{\star}}\,e^{-\nu_{\Psi}}
e^{-(\kappa_{\star}-\nu_{\mathsf{p}}-\nu_{\Psi})x}\le C.
\end{equation*}
It holds for all $x\ge 1$ with some constant $C$ if and only if
$\kappa_{\star}-\nu_{\mathsf{p}}-\nu_{\Psi}\ge 0$. When
this holds, $C=(1-e^{-2\kappa_{\star}})e^{-\nu_{\Psi}}/(2\sinh\kappa_{\star})$ suffices. When it fails,
the left side is unbounded in $x$, since its coefficient is positive.

(b2) Let $V=\varepsilon\mathbf{1}_{[1,x]}$. Every $z$ in the support of $V$ satisfies $z\le x$, so the
formula of (b1) applies to each of them:
\begin{equation*}
\delta K(x)=-\varepsilon\,\frac{e^{-\kappa_{\star}x}}{2\sinh\kappa_{\star}}\sum_{z=1}^{x}(1-e^{-2\kappa_{\star}z})
+O(\varepsilon^{2}).
\end{equation*}
Put $\xi_{x}:=\sum_{z=1}^{x}e^{-2\kappa_{\star}z}$. Then $0\le\xi_{x}\le(e^{2\kappa_{\star}}-1)^{-1}$
uniformly in $x$, and
\begin{equation*}
\sum_{z=1}^{x}(1-e^{-2\kappa_{\star}z})=x-\xi_{x}=x+O(1).
\end{equation*}
This gives \eqref{5.8:eq-rate-window-3}.

The same leading term appears if $V$ is uniform. A uniform $\varepsilon$ shifts the decay rate
to $\tilde{\kappa}$, defined by $2\cosh\tilde{\kappa}=2+m^{2}+\varepsilon$. Then
\begin{equation*}
\tilde{\kappa}=\kappa_{\star}+\frac{\varepsilon}{2\sinh\kappa_{\star}}+O(\varepsilon^{2}),
\qquad
e^{-\tilde{\kappa}x}-e^{-\kappa_{\star}x}=-(\tilde{\kappa}-\kappa_{\star})\,x\,e^{-\kappa_{\star}x}+\dots,
\end{equation*}
whose first-order term is $-\varepsilon K(x)\,x/(2\sinh\kappa_{\star})$.

With the sup weight, $\nu_{\Psi}=0$ and $\Psi\equiv\varepsilon$, the inequality at $x\ge 1$ reads
\begin{equation*}
(x-\xi_{x})\,e^{-(\kappa_{\star}-\nu_{\mathsf{p}})x}\le 2C\sinh\kappa_{\star},
\end{equation*}
and at $x=0$ the first-order term vanishes. Suppose $\nu_{\mathsf{p}}<\kappa_{\star}$. The function
$t\mapsto te^{-at}$ on $t\ge 0$ attains its maximum $1/(ea)$ at $t=1/a$. Using also
$x-\xi_{x}\le x$, we get
\begin{equation*}
(x-\xi_{x})\,e^{-(\kappa_{\star}-\nu_{\mathsf{p}})x}\le x\,e^{-(\kappa_{\star}-\nu_{\mathsf{p}})x}
\le\frac{1}{e(\kappa_{\star}-\nu_{\mathsf{p}})}.
\end{equation*}
So the inequality holds with
$C=\bigl(2e\sinh\kappa_{\star}\,(\kappa_{\star}-\nu_{\mathsf{p}})\bigr)^{-1}$, which is
proportional to $(\kappa_{\star}-\nu_{\mathsf{p}})^{-1}$. Suppose instead $\nu_{\mathsf{p}}=\kappa_{\star}$.
Then the left side
equals $x-\xi_{x}\ge x-(e^{2\kappa_{\star}}-1)^{-1}$, which is unbounded in $x$, and the inequality fails.
\end{proof}

\begin{remark}[Three readings of the comparison that fail]\label{5.8:rem-failing-readings}
Let $\mathsf{A}$ and $\mathsf{B}$ be the machines with the coarser cutoff $\eta$ and the finer cutoff
$\eta' = \eta/L$, and let $p$ be the profile at rate $\delta$ on which the two-cutoff comparison of
this section is stated. Write $\kappa_{\star}$ for the threshold rate of the sharp rate window
(Proposition~\ref{5.8:prop-rate-window}) beyond which the sum of a profile rate and a tempering rate
is excluded. Three readings of the comparison are
false in the generality in which they might be applied, and none of them is the comparison itself:
\begin{enumerate}
\item[(f1)] the layers of $\mathsf{A}$ and of $\mathsf{B}$ agree to within their common exponential
smallness;
\item[(f2)] a local currency tempered at rate $\tau$ controls the transported mismatch also for
profiles with $\rho + \tau > \kappa_{\star}$;
\item[(f3)] at carried slots, a currency in the size $|A_{p}|$ of the pair's mismatch potential
alone controls the mismatch.
\end{enumerate}
The amplitude and the non-linear terms give no further reading of this kind. Let $\mathsf{E}$ be
the far datum, $\delta\mathsf{E}$ its two-cutoff mismatch, $\mathbb{H}^{w}$ the response on the cap
layer and $\zeta_{\square}$ the cutoff functions of the box $\square$, all as defined for the flat
energy-difference bound of this section. The complex amplitude
$\mathsf{t}$ enters $\delta\mathsf{E}$ linearly, through the
product $\mathsf{t}\cdot\delta Q(\zeta_{\square}\mathbb{H}^{w})$. Its only requirement is that $\rho$
be at least $\tfrac12\delta$ plus the rate consumed by the absorption estimate of the comparison. The
non-linear terms are local products, and in weighted norms each is dominated by the supremum of its
datum times its profile.
\end{remark}

\begin{proof}
(f1) The true sizes of the two layers are governed by two different decay lengths, since the two
cutoffs have two decay lengths; this is the content of the sharpness part (b) of the sharp rate window
(Proposition~\ref{5.8:prop-rate-window}). The layers therefore do not agree to within a common
exponential smallness of their true sizes. What is true is that the two layers agree to within the
smallness of one common majorant whose rate has room to spare. The comparison of this section is
stated on such a majorant, the profile $p$, and not on the layers' true sizes, so reading (f1) is not
used.

(f2) Let a mismatch sit at a source point $y$, at distance $d = d(x,y)$ from $x$. By the sharpness
part (b) of Proposition~\ref{5.8:prop-rate-window}, the mismatch is transported from $y$ to $x$
without damping relative to the kernel at $x$ of the transport operator $K$ of
Proposition~\ref{5.8:prop-rate-window}. A currency tempered at rate $\tau$ reads the mismatch locally
at $x$,
and that local reading has decayed by the factor $e^{-\tau d}$. The transported mismatch is
therefore not controlled by its local reading once $\rho + \tau > \kappa_{\star}$. Hence a local
currency tempered at $\tau$ excludes the profiles with $\rho + \tau > \kappa_{\star}$, which is the
failure of (f2).

(f3) By the memory statement for carried slots (a carried slot retains each entry it has received,
so that at carried slots it is the mismatch tally $\mathsf{T}_{p}$, and not $|A_{p}|$, that controls
the mismatch), a currency in $|A_{p}|$ alone has the wrong shape there, and (f3) fails.

It remains to show that the amplitude and the non-linear terms give nothing further. In the
representation of $\delta\mathsf{E}$, the amplitude $\mathsf{t}$ occurs only in the product
$\mathsf{t}\cdot\delta Q(\zeta_{\square}\mathbb{H}^{w})$, which is linear in $\mathsf{t}$. The
contribution of this product is therefore $|\mathsf{t}|$ times the bound for
$\delta Q(\zeta_{\square}\mathbb{H}^{w})$, and the only condition that bound imposes is on the
profile rate, namely $\rho \ge \tfrac12\delta$ plus the rate consumed by the absorption estimate of
the comparison. The
non-linear terms are local products. The weighted-norm estimates for such products dominate each of
them by the supremum of its datum times its profile, so they impose no condition beyond those
already stated.
\end{proof}

\begin{remark}[Available kernel rates against required rates]\label{5.8:rem-kernel-rates}
Recall that the profile decay rates are fixed as $\delta=\tfrac18\delta_{0}$ and
$\delta_{1}=\delta_{0}/32=\tfrac14\delta$.

The estimates that use Proposition~\ref{5.8:prop-rate-window} require the following profile rates:
\begin{itemize}
\item $\rho\ge\tfrac58\delta$ on $\mathfrak{B}$, as required by the lemma on the flat profile
$p^{\flat}$, the profile of rate $\tfrac58\delta$;
\item $\rho=\tfrac12\delta$ on $\mathfrak{B}_{j}$;
\item rate $\tau=\delta_{1}=\tfrac14\delta$ for the smeared currency $\Psi_{p}$ in which the
mismatch tally $\mathsf{T}_{p}$ of a presented pair is measured.
\end{itemize}

Every kernel entering the linearised difference \eqref{5.8:eq-kernel-difference-a} decays at a rate $\kappa\ge\tfrac12\delta_{0}=4\delta$. Consequently
\begin{equation}\label{5.8:eq-kernel-rates-1}
\rho+\tau=\tfrac58\delta+\tfrac14\delta=\tfrac78\delta<4\delta\le\kappa .
\end{equation}
Even the rate $\rho=\delta$ of the profile $p(y)$ gives
\begin{equation*}
\rho+\tau=\tfrac54\delta<4\delta .
\end{equation*}

Hence, in the linearised model, in which the difference of the two kernels is given by
\eqref{5.8:eq-kernel-difference-a}, no admissible rippled background violates, at
$\rho=\tfrac58\delta$, the estimate that bounds the resulting change of the kernel, acting on
the datum $b$, by a constant multiple of $\Psi(y)\mathsf{p}(y)$ (the form of the conclusion of
Proposition~\ref{5.8:prop-rate-window}(a)); we call it the transported-mismatch estimate.
Here a rippled background is one obtained from the unrippled background $U_{c}^{\mathsf{m}}$
by inserting one or more factors $e^{i\eta\mathsf{t}_{i}\zeta_{\square_{i}}\mathbb{H}_{i}^{\mathsf{m}}}$,
the ripples, where $\zeta_{\square_{i}}$ is the cutoff function carried by the box $\square_{i}$;
it is admissible if each complex amplitude $\mathsf{t}_{i}$ lies in the disc in which it ranges
in the ripple construction.
\end{remark}

\begin{proof}
\emph{The rates.} The kernels in the chain \eqref{5.8:eq-kernel-difference-a} fall into two groups:
\begin{itemize}
\item Ba{\l}aban's operators $G'$, $(Q'G'^{2}Q'^{*})^{-1}$ and $G_{0}$ of \cite{B-CMP99b}, which decay at rate $\delta_{0}$;
\item $P$, $H_{0}$, $(I+\mathfrak{R})^{-1}$ and $\mathfrak{D}$, which decay at rate $\tfrac12\delta_{0}$, by the weighted-calculus bounds for these operators.
\end{itemize}
So every such kernel has $\kappa\ge\tfrac12\delta_{0}=4\delta$. Inserting
$\delta=\tfrac18\delta_{0}$ and $\delta_{1}=\tfrac14\delta$ gives \eqref{5.8:eq-kernel-rates-1},
and likewise $\delta+\tfrac14\delta=\tfrac54\delta<4\delta$. Thus $\rho+\tau\le\kappa$ holds with
room to spare:
\begin{equation*}
\kappa-\rho-\tau\ge4\delta-\tfrac78\delta=\tfrac{25}{8}\delta .
\end{equation*}
In particular, the hypothesis $\rho+\tau+\varsigma\le\kappa$ of
Proposition~\ref{5.8:prop-rate-window}(a) holds for every summability rate
$\varsigma\le\tfrac{25}{8}\delta$.

\emph{Uniformity under rippling.} For an admissible rippled background, $\mathbf{U}$ stays inside the class \cite[(3.35)]{B-CMP99b}, and on that class Ba{\l}aban's rates are uniform. Indeed, \cite[Theorem~3.1]{B-CMP99b} is stated for an arbitrary configuration satisfying the regularity condition \cite[(3.35)]{B-CMP99b}. Moreover, \cite[p.~399]{B-CMP99b} states that the constants in the formulations of both theorems there do not depend on the sequence of regions $\{\Omega_j\}$ of that paper. The data vector $\mathsf{m}$ stays inside $|\mathsf{m}|\le\mathfrak{r}_{0}$. Hence the rate $\kappa$ and the constant $B$ in the kernel bounds $Be^{-\kappa d}$ of Proposition~\ref{5.8:prop-rate-window} are the same for every admissible rippled background.

Among the weighted-calculus bounds for this family of operators, the one member whose decay is
printed only at the lower rate $\tfrac18\delta_{0}=4\delta_{1}=\delta$ is the derivative of
$\mathcal{H}$ in \cite[(190)]{B-CMP102b} with respect to the field variable of that formula; it is
printed there without proof. It does not occur in the chain \eqref{5.8:eq-kernel-difference-a},
and it is used neither here nor in the bound for the smeared datum majorant $\mathsf{P}_{b}$.

\emph{The candidate counterexample.} The most unfavourable admissible background for the
transported-mismatch estimate is one ripple of natural size placed on the way:
\begin{equation*}
\mathbf{U}^{\mathsf{m}}
=e^{i\eta\mathsf{t}_{1}\zeta_{\square_{1}}\mathbb{H}_{1}^{\mathsf{m}}}U_{c}^{\mathsf{m}} .
\end{equation*}
Here $|\mathsf{t}_{1}|$ is equal to the radius of the disc in which the complex amplitude
$\mathsf{t}_{1}$ ranges, and $\square_{1}$ lies between the set $\mathcal{C}_{\mathsf{j}}$
carrying the datum $b$ and the box $\square$, or on $\mathcal{C}_{\mathsf{j}}$ itself.

Its two-cutoff mismatch is
\begin{equation*}
\mathfrak{a}=\mathsf{t}_{1}\zeta_{\square_{1}}\delta\mathbb{H}_{1}+\dots ,
\end{equation*}
where $\delta\mathbb{H}_{1}$ denotes the two-cutoff difference of the response $\mathbb{H}_{1}$
(not a multiple of the rate $\delta$); on $\square_{1}$ it has size $\varpi_{1}\,\mathsf{H}^{\sharp}$
and no smaller, where $\varpi_{1}$ is the relative-mismatch factor of the ripple and
$\mathsf{H}^{\sharp}$ is the natural unit of size of the two-cutoff mismatch. By
\eqref{5.8:eq-kernel-difference-a}, for $y\in\square^{\zeta}$ and $c\in\mathcal{C}_{\mathsf{j}}$
the kernel $K(y,c)$ changes by the
term $-\mathcal{G}(y,\square_{1})\,V_{\mathfrak{a}}\,K(\square_{1},c)$. Here, in the notation of
Proposition~\ref{5.8:prop-rate-window}, $\mathcal{G}$ and $K$ stand for the outer kernels of a
term of \eqref{5.8:eq-kernel-difference-a}, each bounded by $Be^{-\kappa d}$, and
$V_{\mathfrak{a}}$ for its middle factor, which carries the mismatch $\mathfrak{a}$ and is of
size $|\mathfrak{a}|$ on $\square_{1}$. Hence this term is at most
\begin{equation*}
B^{2}|\mathfrak{a}|\,e^{-\kappa[d(y,\square_{1})+d(\square_{1},c)]}
\le B^{2}|\mathfrak{a}|\,e^{-\kappa d(y,c)} .
\end{equation*}

For $c\in\mathcal{C}_{\mathsf{j}}$, the right member $\Psi(y)\mathsf{p}(y)$ of the
transported-mismatch estimate satisfies
\begin{equation*}
\Psi(y)\mathsf{p}(y)\ge|\mathfrak{a}|\,e^{-\tau d(y,\square_{1})}\,
\sup|b|\,e^{-\rho d(y,\mathcal{C}_{\mathsf{j}})} .
\end{equation*}
To compare the two, choose $y_{1},y_{2}\in\square_{1}$ and $y_{3}\in\mathcal{C}_{\mathsf{j}}$ with
$d(y,y_{1})=d(y,\square_{1})$ and $d(y_{2},y_{3})=d(\square_{1},\mathcal{C}_{\mathsf{j}})$. The
triangle inequality then gives
\begin{equation*}
d(y,\mathcal{C}_{\mathsf{j}})
\le d(y,\square_{1})+\operatorname{diam}\square_{1}+d(\square_{1},\mathcal{C}_{\mathsf{j}}) .
\end{equation*}
Also $d(\square_{1},c)\ge d(\square_{1},\mathcal{C}_{\mathsf{j}})$, because
$c\in\mathcal{C}_{\mathsf{j}}$. Multiplying the kernel bound above, in its first form, by
$|b(c)|\le\sup|b|$ and dividing by the lower bound for $\Psi(y)\mathsf{p}(y)$, we find that the
quotient is at most
\begin{equation}\label{5.8:eq-kernel-rates-2}
B^{2}e^{\rho\operatorname{diam}\square_{1}}
\exp\bigl\{-(\kappa-\rho-\tau)\,d(y,\square_{1})
-(\kappa-\rho)\,d(\square_{1},\mathcal{C}_{\mathsf{j}})\bigr\}.
\end{equation}

By \eqref{5.8:eq-kernel-rates-1}, $\rho+\tau\le\kappa$, so both exponents in \eqref{5.8:eq-kernel-rates-2} are non-positive. The quotient is therefore bounded, whatever the values of $|\mathsf{t}_{1}|$ and of the amplitude $|\mathsf{t}|$ carried by the box $\square$. It is also bounded whatever the number and the places of the ripples: their contributions add, and the sum over $y$ in Proposition~\ref{5.8:prop-rate-window}(a) controls them.

To violate the transported-mismatch estimate one would need a ripple that slows the kernels below
$\rho+\tau=\tfrac78\delta$. By the uniformity established above, this is impossible inside the
class \cite[(3.35)]{B-CMP99b}.
\end{proof}

\begin{definition}[Tempered weights and weighted supremum norms]\label{5.8:def-tempered-weights}
The admissible sets $\mathfrak{S}$, $\mathfrak{S}'$, the set of sites $\mathfrak{B}=\bigcup_{j}\Lambda_{j}$ with
its layers $\Lambda_{j}$ and its caps $\mathfrak{B}_{j}$, the distance $d$, the level function $\lambda$, the scaled length
$\ell=L^{\lambda}\eta$, the blocks $\Delta(y)$, $\tilde{\Delta}(y)$ and the scaled size $|A|_{y}$ of a
one-form are as fixed in Definition~\ref{5.8:not-paired-sets} (two cutoffs on exactly
paired sets: geometry and floors).
\begin{enumerate}
\item For $\sigma\ge 0$ and $A\ge 1$ let
\begin{equation*}
\mathfrak{M}^{A}_{\sigma}(\mathfrak{S}):=\bigl\{\,w:\mathfrak{B}\to(0,\infty)\;:\;
w(y)\le A\,e^{\sigma d(y,y')}\,w(y')\ \text{for all } y,y'\in\mathfrak{B}\,\bigr\},
\end{equation*}
and $\mathfrak{M}_{\sigma}:=\mathfrak{M}^{1}_{\sigma}$. The elements of $\mathfrak{M}^{A}_{\sigma}(\mathfrak{S})$
are the \emph{tempered weights} of rate $\sigma$ and constant $A$; since they are functions on
$\mathfrak{B}$, we also write $\mathfrak{M}^{A}_{\sigma}(\mathfrak{B})$ for the same class. For any set $Y$
of sites on which $d$ is defined, in particular $\mathfrak{B}_{j}$ and the set of sites of
$\mathfrak{S}'$, the class $\mathfrak{M}^{A}_{\sigma}(Y)$ of functions $w:Y\to(0,\infty)$ is defined by
the same inequality for $y,y'\in Y$; $\mathfrak{M}^{A}_{\sigma}(\mathfrak{S}')$ denotes the class formed
on the set of sites of $\mathfrak{S}'$.
\item The \emph{$\sigma$-hull} of $w:\mathfrak{B}\to(0,\infty)$ is
$w^{[\sigma]}(y):=\sup_{y'}w(y')\,e^{-\sigma d(y,y')}$.
\item The \emph{size at $y$} of an object is the following quantity.
\begin{itemize}
\item A one-form $A$ (the letter $A$ denoting in this item a one-form, not the constant of item 1):
$|A|_{y}$; moreover
\begin{align*}
N_{y}(A)&:=\sup_{\tilde{\Delta}(y)}\max\{\ell|A|,\ \ell^{2}|\nabla A|\},\\
N^{\beta}_{y}(A)&:=\max\{N_{y}(A),\ \ell^{2+\beta}[\nabla A]_{\beta;\tilde{\Delta}(y)}\},\\
N^{\Delta}_{y}(A)&:=\max\{N^{\beta}_{y}(A),\ \sup\ell^{3}|\Delta_{\mathbf{U}}A|\},
\end{align*}
where $\nabla$ is the derivative of one-forms, $[\cdot]_{\beta;\tilde{\Delta}(y)}$ is the H\"older
seminorm of index $\beta$ on $\tilde{\Delta}(y)$, and $\Delta_{\mathbf{U}}$ is the Laplacian of
one-forms in the background field $\mathbf{U}$; the superscript $\Delta$ of $N^{\Delta}_{y}$ refers to
this Laplacian, not to a block.
\item A source $f$: $\sup\ell^{3}|f|$.
\item A scalar of weight $a\in\{0,2,4\}$: $\sup\ell^{a}|\cdot|$, together with
$\ell^{a+1}|D\cdot|$ and $\ell^{a+1+\theta}[D\cdot]_{\theta}$ where this is said; here $D$ is the
derivative of scalars and $[\cdot]_{\theta}$ the H\"older seminorm of index $\theta$.
\item A datum $B$, a function of cells $c$: $|B(c)|$, with $c\ni y$.
\item A linear functional in the format $\sharp(\kappa)$: $\kappa(y)$.
\end{itemize}
\item For a weight $w$, $\|\cdot\|_{w}:=\sup_{y}(\text{size at }y)/w(y)$.
\end{enumerate}
\end{definition}

The weight classes have the following properties, for $\sigma,\sigma',\rho\ge0$, $A,A'\ge1$ and real $a$,
proved below:
\begin{equation}\label{5.8:eq-tempered-weights-a}
\begin{aligned}
&\text{(w1)}\quad \mathfrak{M}^{A}_{\sigma}\cdot\mathfrak{M}^{A'}_{\sigma'}\subset
\mathfrak{M}^{AA'}_{\sigma+\sigma'},\ \text{and sums and suprema stay in } \mathfrak{M}^{A}_{\sigma};\\
&\text{(w2)}\quad w\in\mathfrak{M}^{A}_{\sigma}\ \Longrightarrow\ w\le w^{[\sigma]}\le Aw\ \text{ and }\
w^{[\sigma]}\in\mathfrak{M}_{\sigma};\\
&\text{(w3)}\quad L^{a\lambda}\in\mathfrak{M}^{L^{|a|}}_{|a|\delta_{1}/16};\\
&\text{(w4)}\quad f^{\sim}\in\mathfrak{M}_{\delta_{1}},\qquad
\Psi_{p},\ \Psi^{B}_{p}\in\mathfrak{M}^{L}_{\delta_{1}};\\
&\text{(w5)}\quad \mathfrak{M}^{\wedge}_{\rho}(\mathfrak{B})\subset\mathfrak{M}_{\rho}(\mathfrak{B});\\
&\hphantom{\text{(w5)}}\quad p\in\mathfrak{M}_{\delta}(\mathfrak{B})=\mathfrak{M}_{4\delta_{1}}(\mathfrak{B}),\
p^{\flat}\in\mathfrak{M}_{5\delta_{1}/2}(\mathfrak{B}),\
\hat{p},\hat{p}_{\square}\in\mathfrak{M}_{2\delta_{1}}(\mathfrak{B}_{j});\\
&\text{(w6)}\quad w\in\mathfrak{M}_{\sigma}(\mathfrak{S})\ \text{extends to}\
w'\in\mathfrak{M}_{\sigma}(\mathfrak{S}'),\ \ w'=w\ \text{on the paired sites}.
\end{aligned}
\end{equation}
In (w6), which concerns the collar $\Lambda'_{0}$ of the finer cutoff $\mathsf{B}$, fixed in
Definition~\ref{5.8:not-paired-sets},
$w'(x):=\inf_{y\in\mathfrak{B}}w(y)e^{\sigma d(x,y)}$ for $x$ a site of $\mathfrak{S}'$, that is, a site
of $\mathfrak{B}$ or of $\Lambda'_{0}$. In
(w4), $f^{\sim}$ is the $\delta_{1}$-smear of a profile $f$, and $\Psi_{p}$, $\Psi^{B}_{p}$ are the
smeared currencies of a presented pair; in (w5), $p$ is the profile at rate $\delta$, $\hat{p}$ and
$\hat{p}_{\square}$ are its versions on a cap $\mathfrak{B}_{j}$, $p^{\flat}$ is the flat profile at
rate $\tfrac58\delta$, and
$\mathfrak{M}^{\wedge}_{\rho}$ is the weight class of rate $\rho$ formed with the distance $d^{\wedge}$,
the second distance on $\mathfrak{B}$ with respect to which the profiles $p$, $p^{\flat}$, $\hat{p}$,
$\hat{p}_{\square}$ are tempered.
The cost rule attached to (w2) is this: the constants of the weighted bounds that follow are stated
for weights with $A=1$, and a weight of
$\mathfrak{M}^{A}_{\sigma}$ costs the factor $A$ once,
in the final bound, never in an absorption. This is possible by (w2), since $w\le w^{[\sigma]}$ and
$w^{[\sigma]}\le Aw$ give
\begin{equation*}
\|\cdot\|_{w^{[\sigma]}}\le\|\cdot\|_{w}\le A\|\cdot\|_{w^{[\sigma]}},
\end{equation*}
with $w^{[\sigma]}\in\mathfrak{M}_{\sigma}$.

\begin{proof}[Proof of \eqref{5.8:eq-tempered-weights-a}]
\textit{(w1).} Let $w\in\mathfrak{M}^{A}_{\sigma}$, $v\in\mathfrak{M}^{A'}_{\sigma'}$. Multiplying the two
defining inequalities,
\begin{equation*}
w(y)v(y)\le A e^{\sigma d(y,y')}w(y')\cdot A' e^{\sigma' d(y,y')}v(y')
=AA'e^{(\sigma+\sigma')d(y,y')}\,w(y')v(y').
\end{equation*}
If $w_{1},w_{2}\in\mathfrak{M}^{A}_{\sigma}$, adding their inequalities gives the inequality for
$w_{1}+w_{2}$. If $(w_{i})$ is a family in $\mathfrak{M}^{A}_{\sigma}$ whose supremum is finite at one
site $y_{0}$, then $w_{i}(y)\le Ae^{\sigma d(y,y_{0})}w_{i}(y_{0})$ shows that it is finite everywhere,
and it is positive; from $w_{i}(y)\le Ae^{\sigma d(y,y')}\sup_{k}w_{k}(y')$, taking the supremum
over $i$ gives the inequality for $\sup_{i}w_{i}$.

\textit{(w2).} The choice $y'=y$ in the supremum gives $w\le w^{[\sigma]}$. By the defining inequality
with the roles of $y,y'$ exchanged and the symmetry of $d$,
$w(y')e^{-\sigma d(y,y')}\le Aw(y)$ for every $y'$; hence $w^{[\sigma]}\le Aw$, and
$w^{[\sigma]}$ takes values in $(0,\infty)$. By the triangle inequality for $d$,
\cite[(2.54)]{B-CMP96}, $-d(y,y'')\le d(y,y')-d(y',y'')$, so
\begin{align*}
w^{[\sigma]}(y)&=\sup_{y''}w(y'')e^{-\sigma d(y,y'')}\\
&\le e^{\sigma d(y,y')}\sup_{y''}w(y'')e^{-\sigma d(y',y'')}=e^{\sigma d(y,y')}w^{[\sigma]}(y'),
\end{align*}
that is, $w^{[\sigma]}\in\mathfrak{M}_{\sigma}$.

\textit{(w3).} By \eqref{5.8:eq-paired-sets-a},
$d(y,y')\ge RM_{1}(|\lambda(y)-\lambda(y')|-1)_{+}$, whence
$|\lambda(y)-\lambda(y')|\le 1+d(y,y')/(RM_{1})$: this is trivial if $|\lambda(y)-\lambda(y')|<1$,
and otherwise it is the inequality itself. By \eqref{5.8:eq-paired-sets-c},
$\log L/(RM_{1})\le\delta_{1}/16$. Therefore
\begin{align*}
\frac{L^{a\lambda(y)}}{L^{a\lambda(y')}}
\le L^{|a|\,|\lambda(y)-\lambda(y')|}
&\le L^{|a|}\exp\Bigl(|a|\,\frac{\log L}{RM_{1}}\,d(y,y')\Bigr)\\
&\le L^{|a|}\,e^{|a|\delta_{1}d(y,y')/16}.
\end{align*}

\textit{(w4).} The three memberships of (w4) are the statement that the $\delta_{1}$-smear $f^{\sim}$
and the smeared currencies $\Psi_{p}$, $\Psi^{B}_{p}$ vary slowly, at rate $\delta_{1}$ with the
constants $1$ and $L$ respectively; this statement is used here as it is stated where these objects
are defined.

\textit{(w5).} By the inequality $d^{\wedge}\le d$ between the two distances, stated where $d^{\wedge}$
is introduced, and
since $\rho\ge0$, $w(y)\le e^{\rho d^{\wedge}(y,y')}w(y')$ implies $w(y)\le e^{\rho d(y,y')}w(y')$,
which is the inclusion. Here $\delta_{1}=\delta_{0}/32$ by Definition~\ref{5.8:not-paired-sets},
and the decay rate of the profiles is $\delta=\delta_{0}/8$, as fixed where the profiles are
defined, so $\delta=4\delta_{1}$ and $\tfrac58\delta=\tfrac52\delta_{1}$. The inclusion is applied to
the memberships
\begin{equation*}
p\in\mathfrak{M}^{\wedge}_{\delta}(\mathfrak{B}),\qquad
p^{\flat}\in\mathfrak{M}^{\wedge}_{5\delta/8}(\mathfrak{B}),\qquad
\hat{p},\ \hat{p}_{\square}\in\mathfrak{M}^{\wedge}_{2\delta_{1}}(\mathfrak{B}_{j}),
\end{equation*}
which are an input of this proof, taken as stated for the profiles in the distance $d^{\wedge}$;
it gives the memberships listed in \eqref{5.8:eq-tempered-weights-a}.

\textit{(w6).} There is one distance $d$ on paired sites, shared by $\mathfrak{S}$ and $\mathfrak{S}'$,
by Definition~\ref{5.8:not-paired-sets}. Fix $y_{0}\in\mathfrak{B}$ and a site $x$ of
$\mathfrak{S}'$. Then
$w'(x)\le w(y_{0})e^{\sigma d(x,y_{0})}<\infty$, and, since
$w(y)\ge w(y_{0})e^{-\sigma d(y,y_{0})}$ and $d(x,y)-d(y,y_{0})\ge -d(x,y_{0})$,
\begin{equation*}
w(y)e^{\sigma d(x,y)}\ge w(y_{0})e^{\sigma(d(x,y)-d(y,y_{0}))}\ge w(y_{0})e^{-\sigma d(x,y_{0})}
\quad\text{for all }y\in\mathfrak{B},
\end{equation*}
so $w'(x)>0$. For sites $x,x'$ of $\mathfrak{S}'$, $d(x,y)\le d(x,x')+d(x',y)$ gives
\begin{equation*}
w'(x)\le\inf_{y\in\mathfrak{B}}w(y)e^{\sigma(d(x,x')+d(x',y))}=e^{\sigma d(x,x')}w'(x'),
\end{equation*}
so $w'\in\mathfrak{M}_{\sigma}(\mathfrak{S}')$. At a paired site $x\in\mathfrak{B}$, the choice $y=x$
gives $w'(x)\le w(x)$, and $w(x)\le w(y)e^{\sigma d(x,y)}$ for all $y\in\mathfrak{B}$ gives
$w(x)\le w'(x)$; thus $w'=w$ there.
\end{proof}

\begin{proposition}[Ba{\l}aban's local decay of propagators and projections
{\cite[(3.42)--(3.45), (3.48), (3.49), Theorems~3.1, 3.3, (3.132)--(3.133)]{B-CMP99b}},
{\cite[(71), (73)]{B-CMP102b}}]
\label{5.8:prop-local-decay}
Use the notation of Notation~\ref{5.8:not-paired-sets}: the cutoffs $\mathsf{A}$ (spacing
$\eta$) and $\mathsf{B}$
(spacing $\eta' = \eta/L$), the level function $\lambda(x)$, the scaled length
$\ell_{y} = L^{\lambda(y)}\eta$, the blocks $\Delta(y)$ and their neighbourhoods
$\tilde{\Delta}(y)$, the distance $d$, and $\delta_{1} = \delta_{0}/32$. Let the background be
one of $\mathbf{U}$, $W$ (cutoff $\mathsf{A}$) or $\mathbf{U}'$ (cutoff $\mathsf{B}$), and assume
that it satisfies \cite[(3.35)]{B-CMP99b} with $M\alpha_{0} \le a_{0}$, where $M$, $\alpha_{0}$
and $a_{0}$ are the parameters denoted by these letters in \cite{B-CMP99b}. Consider the list
\begin{align*}
\mathfrak{L} := \bigl\{ & G',\ \nabla G',\ G'\nabla^{*},\ \Delta_{U}G';\
\nabla G'\nabla^{*};\ C';\ P,\ DP,\ PD^{*},\ DPD^{*}; \\
& G_{0},\ \nabla G_{0},\ G_{0}\nabla^{*},\ \Delta_{U}G_{0},\
\nabla G_{0}\nabla^{*};\ C_{Q};\ H_{0},\ \nabla H_{0};\ \Theta;\
(I+\mathfrak{R})^{-1} \bigr\},
\end{align*}
where
\begin{gather*}
C' := (Q'G'^{2}Q'^{*})^{-1}, \qquad C_{Q} := (QG_{0}Q^{*})^{-1}, \qquad
H_{0} := G_{0}Q^{*}C_{Q}\ell^{-1}, \\
\Theta := \ell^{-1}(C_{Q} - a)\ell^{-1}, \qquad \mathfrak{R} := dC(X)H_{0}.
\end{gather*}
Here $G'$, $Q'$, $Q$, $P$, $D$, $\nabla$, $\nabla^{*}$ and $\Delta_{U}$ are the
objects denoted by the same letters in \cite{B-CMP99b}, and $G_{0}$ is the operator $G(U)$ of
\cite[Theorem~3.3]{B-CMP99b}, all built on the given background, which
is denoted $U$ there and is one of $\mathbf{U}$, $W$, $\mathbf{U}'$; $\ell^{-1}$ is
multiplication by the function $y \mapsto \ell_{y}^{-1}$; $a$ is the quantity
denoted by this letter in \cite{B-CMP99b}; and $dC(X)$ is the derivative of the
block function $C$ of
\cite{B-CMP102b}, taken at the variable denoted $X$ there (this $X$ is not a large-field
region). Then each operator of $\mathfrak{L}$ obeys a local inequality with the exponential
factor $e^{-8\delta_{1}d(y,y')}$, as follows.

For $G'$, $\nabla G'$, $G'\nabla^{*}$, $\Delta_{U}G'$, $\nabla G'\nabla^{*}$, $C'$, $P$, $DP$,
$PD^{*}$, $DPD^{*}$, and, through \cite[Theorem~3.3]{B-CMP99b}, for $G_{0}$, $\nabla G_{0}$,
$G_{0}\nabla^{*}$, $\Delta_{U}G_{0}$, $\nabla G_{0}\nabla^{*}$, it is the local inequality
printed for the operator in \cite[(3.42)--(3.45), (3.48), (3.49)]{B-CMP99b}. In each
inequality of \cite{B-CMP99b} so cited the source is
supported in $\Delta(y')$ or in $\tilde{\Delta}(y')$, and the value is read at points of
$\Delta(y)$; the inequality keeps the powers of $\ell_{y}$ and $\ell_{y'}$ and the norms of
the printed one --- the supremum norm, and the H\"older norms of
\cite[(3.43)--(3.45)]{B-CMP99b} --- and only its exponential factor is replaced by
$e^{-8\delta_{1}d(y,y')}$.

For $C_{Q}$ and $H_{0}$ the local inequality is the block bound of the lemma on the local
decay of $C_{Q}$ and $H_{0}$, whose derivation uses \cite[Theorem~3.3]{B-CMP99b} and
\cite[(3.132)]{B-CMP99b}; $\nabla H_{0}$ and $\Theta$ obey the local bound that results from
these bounds, from the bounds above and from the definitions of $H_{0}$ and $\Theta$, with the
resulting powers of $\ell_{y}$, $\ell_{y'}$ and the factor
$e^{-8\delta_{1}d(y,y')}$. The operator $(I+\mathfrak{R})^{-1}$ obeys
\cite[(71)]{B-CMP102b} with the kernel denoted $h$ in \cite[(69)--(71)]{B-CMP102b} taken to be
$H_{0}$, the kernel bound of $H_{0}$ being taken from the lemma on the local decay of $C_{Q}$
and $H_{0}$, and with the points $c_{-}$, $c'_{-}$ attached to the cubes $c$, $c'$ of
\cite{B-CMP102b} in the roles of $y$, $y'$.
\end{proposition}

This is \cite[(3.42)--(3.45), (3.48), (3.49), Theorems~3.1, 3.3, (3.132)--(3.133)]{B-CMP99b}
together with \cite[(71), (73)]{B-CMP102b}, and, for $C_{Q}$ and $H_{0}$, the lemma on the
local decay of $C_{Q}$ and $H_{0}$, which supplies their decay between blocks from
\cite[Theorem~3.3]{B-CMP99b} and \cite[(3.132)]{B-CMP99b}, from which the bounds for
$\nabla H_{0}$ and $\Theta$ follow, and which furnishes the kernel bound of $H_{0}$ used in
\cite[(71)]{B-CMP102b}; we record the printed inequalities in the present notation.

Throughout, $y, y'$ are sites of the paired sets, and $\varphi$ is a form supported in
$\Delta(y')$ unless $\tilde{\Delta}(y')$ is stated. The point $x$ ranges over $\Delta(y)$,
and $\psi$ is a function. The letters $\varphi$ and $\psi$ are written $\lambda$ and $\zeta$
in \cite{B-CMP99b}; they are renamed here because $\lambda$ denotes the level function and
$\zeta$ the cutoff functions $\zeta_{\square}$. The symbol $|\cdot|$ is the supremum norm, and
$\|\cdot\|^{\xi}_{\beta}$, $\|\cdot\|^{\xi'}_{\varepsilon}$ are the H\"older norms of
\cite[(3.43)--(3.45)]{B-CMP99b}, with the indices $\xi, \xi', \beta, \beta_{0},
\varepsilon$, the H\"older norm $\|\cdot\|_{\beta}$ and the constants $B_{0}$,
$B_{0}(\beta_{0})$, $B_{0}'(\varepsilon)$, $B_{0}'(\varepsilon,\beta_{0})$ and $O(1)$ of that
paper. Ba{\l}aban's scaled length at the level of $y$ --- $L^{j}\eta$ for the cutoff
$\mathsf{A}$, and $L^{j}\eta'$ for the cutoff $\mathsf{B}$, with $j$ the level in that
cutoff's own indexing --- equals $\ell_{y}$ by the single level function of
Notation~\ref{5.8:not-paired-sets}, and likewise at $y'$. Since the dimension is four, his
scaled length at $y'$ raised to minus the dimension is written $\ell_{y'}^{-4}$, and at
\cite[(73)]{B-CMP102b} the power one minus the dimension equals $-3$. The identification with
$\ell_{y}$, $\ell_{y'}$ is used in the displays taken from \cite{B-CMP99b}; in those taken
from \cite{B-CMP102b} the scaled lengths are written as printed.
With these conventions, and with the exponential factor printed in \cite{B-CMP99b} weakened
to $e^{-8\delta_{1}d(y,y')}$ as justified after the displays, the inequalities give the
following bounds.

By \cite[(3.42)]{B-CMP99b},
\begin{equation*}
|(G'\varphi)(x)|,\ |(\nabla G'\varphi)(x)|,\ |(G'\nabla^{*}\varphi)(x)|,\
|(\Delta_{U}G'\varphi)(x)|
\end{equation*}
are bounded, respectively, by $B_{0}\ell_{y}^{2}$, $B_{0}\ell_{y}$, $B_{0}\ell_{y}$ and
$B_{0}$, each times $e^{-8\delta_{1}d(y,y')}|\varphi|$.

By \cite[(3.43)]{B-CMP99b}, both $\|\psi\nabla G'\varphi\|_{\beta}$ and
$\|\psi G'\nabla^{*}\varphi\|_{\beta}$ are bounded by
\begin{equation*}
B_{0}(\beta_{0})\,\ell_{y}^{1-\beta}\bigl(\|\psi\|^{\xi}_{\beta}+|\psi|\bigr)
e^{-8\delta_{1}d(y,y')}|\varphi| .
\end{equation*}

By \cite[(3.44)]{B-CMP99b}, for $\operatorname{supp}\varphi \subset \tilde{\Delta}(y')$,
\begin{equation*}
|(\nabla G'\nabla^{*}\varphi)(x)|
\le B_{0}'(\varepsilon)\,e^{-8\delta_{1}d(y,y')}
\bigl(\|\varphi\|^{\xi'}_{\varepsilon}+|\varphi|\bigr).
\end{equation*}

By \cite[(3.45)]{B-CMP99b},
\begin{equation*}
\|\psi\nabla G'\nabla^{*}\varphi\|_{\beta}
\le B_{0}'(\varepsilon,\beta_{0})\,\ell_{y}^{-\beta}
\bigl(\|\psi\|^{\xi}_{\beta}+|\psi|\bigr)e^{-8\delta_{1}d(y,y')}
\bigl(\|\varphi\|^{\xi'}_{\beta+\varepsilon}+|\varphi|\bigr).
\end{equation*}

By \cite[(3.48)]{B-CMP99b},
\begin{equation*}
|C'(y,y')| \le B_{0}\,\ell_{y}^{-4}\ell_{y'}^{-4}e^{-8\delta_{1}d(y,y')} .
\end{equation*}

By \cite[(3.49)]{B-CMP99b}, for $x \in \Delta(y)$ and $x' \in \Delta(y')$, the kernels
$P$, $(DP)_{\mu}$, $(PD^{*})_{\nu}$ and $(DPD^{*})_{\mu\nu}$ satisfy
\begin{align*}
|P(x,x')| &\le O(1)\,\ell_{y'}^{-4}e^{-8\delta_{1}d(y,y')}, \\
|(DP)_{\mu}(x,x')| &\le O(1)\,\ell_{y}^{-1}\ell_{y'}^{-4}e^{-8\delta_{1}d(y,y')}, \\
|(PD^{*})_{\nu}(x,x')| &\le O(1)\,\ell_{y}^{-1}\ell_{y'}^{-4}e^{-8\delta_{1}d(y,y')}, \\
|(DPD^{*})_{\mu\nu}(x,x')| &\le O(1)\,\ell_{y}^{-2}\ell_{y'}^{-4}e^{-8\delta_{1}d(y,y')}.
\end{align*}
As stated in
\cite[p.~399]{B-CMP99b}, the corresponding bounds hold for H\"older norms of the kernel
$(DPD^{*})_{\mu\nu}(x,x')$.

By \cite[Theorem~3.3]{B-CMP99b}, the operator $G_{0}$ (denoted $G(U)$ there) satisfies
the inequalities \cite[(3.42)--(3.47)]{B-CMP99b}. This gives the bounds displayed above
for $G_{0}$, $\nabla G_{0}$, $G_{0}\nabla^{*}$, $\Delta_{U}G_{0}$ and
$\nabla G_{0}\nabla^{*}$.

In each of the inequalities \cite[(3.42)--(3.45), (3.48), (3.49)]{B-CMP99b} displayed above
the printed exponent is $-\delta_{0}d(y,y')$ or
$-\tfrac12\delta_{0}d(y,y')$. Here $\delta_{0}$ is the rate fixed in \cite{B-CMP99b} by
the sentence ``Let us denote a common, best possible, decay rate for all these operators
by $\delta_{0}$.'' Since $\delta_{1} = \delta_{0}/32$, we have
$\tfrac12\delta_{0} = 16\delta_{1}$. Since $d(y,y') \ge 0$, it follows that
\begin{equation*}
e^{-\delta_{0}d(y,y')} \le e^{-\frac12\delta_{0}d(y,y')} = e^{-16\delta_{1}d(y,y')}
\le e^{-8\delta_{1}d(y,y')} .
\end{equation*}
Hence each printed inequality implies the one displayed above with the factor
$e^{-8\delta_{1}d(y,y')}$.

The members $C_{Q}$, $H_{0}$, $\nabla H_{0}$ and $\Theta$ rest on
\cite[(3.132)--(3.133)]{B-CMP99b}. For $y$, $y'$ in one layer $\Lambda_{j}$ the exponential
factor in the bound of \cite[(3.132)--(3.133)]{B-CMP99b} for $C_{Q}$ does not depend on
$d(y,y')$; decay between such blocks is the content of the lemma on the local decay of
$C_{Q}$ and $H_{0}$, which supplies the local bound for $C_{Q}$ at rate $\delta_{0}$ and for
$H_{0}$ with the factor $e^{-\frac12\delta_{0}d(y,y')}$. Its derivation uses only
\cite[Theorem~3.3]{B-CMP99b} and \cite[(3.132)]{B-CMP99b}, and therefore holds for
$H_{0} = G_{0}Q^{*}C_{Q}\ell^{-1}$ under \cite[(3.35)]{B-CMP99b} alone. The member
$\Theta = \ell^{-1}(C_{Q}-a)\ell^{-1}$ is built from $C_{Q}$ by its definition above; since
$\ell^{-1}$ is multiplication by $y \mapsto \ell_{y}^{-1}$, its kernel is
\begin{equation*}
\Theta(y,y') = \ell_{y}^{-1}\,(C_{Q}-a)(y,y')\,\ell_{y'}^{-1},
\end{equation*}
so that its powers of $\ell$ are those of the bound for $C_{Q}-a$ with the two further
factors $\ell_{y}^{-1}\ell_{y'}^{-1}$. By the definition of $H_{0}$, the member
$\nabla H_{0}$ is the product $(\nabla G_{0})Q^{*}C_{Q}\ell^{-1}$.
By the display above, the factors $e^{-\delta_{0}d(y,y')}$ and
$e^{-\frac12\delta_{0}d(y,y')}$ in these bounds are at most $e^{-8\delta_{1}d(y,y')}$.

For the last member, the inequalities \cite[(69)--(71)]{B-CMP102b} are stated for an
operator built from a kernel $h$, and they use only the kernel bound of $h$. Taking
$h = H_{0}$, so that the operator there is $\mathfrak{R} = dC(X)H_{0}$,
\cite[(71)]{B-CMP102b} gives
\begin{equation*}
|(I+\mathfrak{R})^{-1}(c,c')|
\le 2\,(L^{j'}\eta)^{-4}e^{-\frac12\delta_{0}d(c_{-},c'_{-})}
\le 2\,(L^{j'}\eta)^{-4}e^{-8\delta_{1}d(c_{-},c'_{-})} .
\end{equation*}
Here $c, c'$ are the cubes of \cite{B-CMP102b}, $c_{-}, c'_{-}$ are the points attached
to them there, and $L^{j'}\eta$ is the scaled length at the level $j'$ of $c'$ (with
$\eta'$ in place of $\eta$ for the cutoff $\mathsf{B}$, here and in the next display). In the
same way, \cite[(73)]{B-CMP102b} bounds the kernel of
$\mathfrak{D} := (I+\mathfrak{R})^{-1}dC(X)$; it reads
\begin{align*}
|\mathfrak{D}(A';c,b)|
&\le O(1)\,C_{3}\varepsilon_{3}(L^{j}\eta)^{-3}e^{-\frac12\delta_{0}d(c_{-},y)} \\
&\le O(1)\,C_{3}\varepsilon_{3}(L^{j}\eta)^{-3}e^{-8\delta_{1}d(c_{-},y)} ,
\end{align*}
where $A'$, $b$, the level $j$ and the constants $C_{3}$, $\varepsilon_{3}$ are the
objects of \cite{B-CMP102b}, and $y$ is the site that enters \cite[(73)]{B-CMP102b}. The
operator $\mathfrak{D}$ is not a member of
$\mathfrak{L}$; its bound \cite[(73)]{B-CMP102b} is recorded here beside the bound for
$(I+\mathfrak{R})^{-1}$, from which $\mathfrak{D}$ is obtained by composition with $dC(X)$.
Nor is the operator of \cite[(190)]{B-CMP102b} a member of
$\mathfrak{L}$.

\begin{lemma}[Weighted calculus: products spend no decay rate]\label{5.8:lem-weighted-calculus}
Assume the geometric bounds \eqref{5.8:eq-paired-sets-a} and \eqref{5.8:eq-paired-sets-b}, the floor
\eqref{5.8:eq-paired-sets-c}, and Ba{\l}aban's local decay of propagators and projections,
Proposition~\ref{5.8:prop-local-decay}, for the operators of its list $\mathfrak{L}$. Let
$\sigma\le 6\delta_{1}$ and let $w\in\mathfrak{M}_{\sigma}$, on $\mathfrak{S}$, or on $\mathfrak{S}'$ by
the last of the properties \eqref{5.8:eq-tempered-weights-a} of tempered weights, which carries them to
$\mathfrak{S}'$.
\begin{enumerate}
\item[(i)] Every $\mathcal{O}\in\mathfrak{L}$ is bounded from its source norm to its value norm, both
weighted by $w$:
\begin{equation}\label{5.8:eq-weighted-calculus-1}
\|\mathcal{O}f\|_{w}\le C_{wt}\,\|f\|_{w},\qquad C_{wt}:=2d\,\nu_{d}\,B_{\star}L^{4}c_{1},
\end{equation}
where $d=4$ is the dimension (the letter $d$ without arguments denotes the dimension, while
$d(\cdot,\cdot)$ is the distance of \eqref{5.8:eq-paired-sets-a}), $B_{\star}$ is the largest of the
constants in the local inequalities of
Proposition~\ref{5.8:prop-local-decay}, $\nu_{d}$ is the overlap number of the family
$\{\tilde{\Delta}(y)\}$, and $c_{1}$ is the constant of \eqref{5.8:eq-paired-sets-b}.
\item[(ii)] Every finite product of members of $\mathfrak{L}$ and of local maps whose value at $y$
reads only cells $y'$ with $d(y,y')\le 2$ is likewise bounded from the $w$-weighted source norm to the
$w$-weighted value norm, with constant the product of the constants of its factors: each member of
$\mathfrak{L}$ contributes $C_{wt}$, and each such local map contributes, besides its own norm, a factor
$e^{2\sigma}\le e^{12\delta_{1}}$. In particular this holds for $R=1-P$, $RG'$, $RG'D^{*}$, $DRG'D^{*}$,
$PD^{*}$, and for
\begin{equation*}
\mathcal{G}_{0}=G_{0}-G_{0}Q^{*}C_{Q}QG_{0}
\end{equation*}
on sources $f$, measured by $N_{y}(\mathcal{G}_{0}f)$ and $N^{\Delta}_{y}(\mathcal{G}_{0}f)$, and on
$\sharp$-functionals $\mathfrak{f}$, measured by $|\mathcal{G}_{0}\mathfrak{f}|_{y}$.
\item[(iii)] (Transposes.) Let $\delta_{c}$ denote the datum concentrated on the cell $c$. The map
$H_{0}^{\top}$ from sources and $\sharp$-functionals to data, defined by
$(H_{0}^{\top}\mathfrak{f})(c):=\mathfrak{f}[H_{0}\delta_{c}]$, and the maps
$((I+\mathfrak{R})^{-1})^{\top}$ and $\Theta^{\top}=\Theta$ on data obey the same bound.
\end{enumerate}
\end{lemma}

\begin{proof}
(i) Let $\{h_{y'}\}$ be a partition of unity subordinate to the cover $\{\tilde{\Delta}(y')\}$ with
$\ell\,|\nabla h_{y'}|\le C$ for a constant $C$; for sup norms, take instead the indicator functions of the
blocks $\Delta(y')$. Split $f=\sum_{y'}h_{y'}f$. Since $\ell|\nabla h_{y'}|\le C$, the local source norm
of $h_{y'}f$ is at most $C$ times the local source norm of $f$ at $y'$, and by the definition of the
weighted norm the latter is at most $w(y')\|f\|_{w}$.

Apply Proposition~\ref{5.8:prop-local-decay} to each piece $\mathcal{O}(h_{y'}f)$: its source lies in
$\tilde{\Delta}(y')$ (in $\Delta(y')$ for sup norms), and the local inequality bounds its value on
$\Delta(y)$ with Ba{\l}aban's powers of $\ell_{y}$ and $\ell_{y'}$ and the decay
$e^{-8\delta_{1}d(y,y')}$. Writing these powers at the two ends in the scaled sizes at $y$ and at $y'$
produces at most the factor $L^{4|\lambda(y)-\lambda(y')|}$. Indeed, each factor of the local
inequalities of Proposition~\ref{5.8:prop-local-decay} shifts at most four powers of $\ell$ between $y$
and $y'$ when its two ends are written in the scaled sizes at $y$ and at $y'$. Since $\ell=L^{\lambda}\eta$, we have
\begin{equation*}
\ell_{y'}=L^{\lambda(y')-\lambda(y)}\,\ell_{y},
\end{equation*}
so that moving one power of $\ell$ from $y'$ to $y$, or from $y$ to $y'$, costs at most the factor
$L^{|\lambda(y)-\lambda(y')|}$, and four such moves cost at most $L^{4|\lambda(y)-\lambda(y')|}$.
Write $\operatorname{size}_{y}(\mathcal{O}f)$ for the scaled size at $y$, in the
sense of Definition~\ref{5.8:def-tempered-weights}, of the value $\mathcal{O}f$ (the scaled size
appropriate to the kind of value that $\mathcal{O}$ produces). Hence
\begin{equation*}
\operatorname{size}_{y}(\mathcal{O}f)\le B_{\star}\sum_{y'}L^{4|\lambda(y)-\lambda(y')|}
e^{-8\delta_{1}d(y,y')}\,w(y')\,\|f\|_{w}.
\end{equation*}
Since $w\in\mathfrak{M}_{\sigma}$, we have $w(y')\le e^{\sigma d(y,y')}w(y)$, and $\sigma\le 6\delta_{1}$
gives $e^{-8\delta_{1}d(y,y')}w(y')\le e^{-2\delta_{1}d(y,y')}w(y)$. Therefore
\begin{equation}\label{5.8:eq-weighted-calculus-2}
\operatorname{size}_{y}(\mathcal{O}f)\le B_{\star}\,w(y)\,\|f\|_{w}\sum_{y'}
L^{4|\lambda(y)-\lambda(y')|}e^{-2\delta_{1}d(y,y')}.
\end{equation}
By \eqref{5.8:eq-paired-sets-a} together with the floor \eqref{5.8:eq-paired-sets-c},
\begin{equation*}
L^{4|\lambda(y)-\lambda(y')|}e^{-\delta_{1}d(y,y')/4}\le L^{4},
\end{equation*}
so that each summand in \eqref{5.8:eq-weighted-calculus-2} is at most
\begin{equation*}
L^{4}e^{-7\delta_{1}d(y,y')/4}\le L^{4}e^{-\delta_{1}d(y,y')/2},
\end{equation*}
and by \eqref{5.8:eq-paired-sets-b} the sum $\sum_{y'}e^{-\delta_{1}d(y,y')/2}$ is at most $c_{1}$.
Dividing by $w(y)$ and taking the supremum over $y$ gives \eqref{5.8:eq-weighted-calculus-1}.

(ii) Operator norms multiply: the weighted norm of a finite product is at most the product of the weighted
norms of its factors. A factor from $\mathfrak{L}$ contributes $C_{wt}$ by (i). A local map whose value at
$y$ reads only cells $y'$ with $d(y,y')\le 2$ contributes, besides its own norm, the factor $e^{2\sigma}$,
because $w(y')\le e^{\sigma d(y,y')}w(y)\le e^{2\sigma}w(y)$ for such $y'$; and
$e^{2\sigma}\le e^{12\delta_{1}}$ since $\sigma\le 6\delta_{1}$. The operators listed in (ii) are such
products.

(iii) Let $\mathfrak{f}$ be a $\sharp$-functional with local density profile $\kappa$. Pairing
$\mathfrak{f}$ with the one-form $H_{0}\delta_{c}$ block by block, through the first differences and
through the values of $H_{0}\delta_{c}$ on the bonds of each $\Delta(y)$, gives
\begin{align*}
|\mathfrak{f}[H_{0}\delta_{c}]|
&\le\sum_{y}\Bigl(\sum_{b\in\Delta(y)}\eta^{4}\Bigr)
\Bigl[\ell^{-2}N_{y}(H_{0}\delta_{c})\,\ell^{-2}\kappa(y)
+\ell^{-1}N_{y}(H_{0}\delta_{c})\,\ell^{-3}\kappa(y)\Bigr]\\
&\le 2d\sum_{y}N_{y}(H_{0}\delta_{c})\,\kappa(y),\qquad d=4.
\end{align*}
Since $H_{0}\in\mathfrak{L}$, Proposition~\ref{5.8:prop-local-decay}, written in scaled sizes as in (i),
gives
\begin{equation*}
N_{y}(H_{0}\delta_{c})\le B_{\star}L^{4|\lambda(y)-\lambda(c)|}e^{-8\delta_{1}d(y,c)},
\end{equation*}
with $\lambda$ and $d$ evaluated at a site of $c$. Bounding $\kappa(y)$ by $w(y)$ times the weighted norm
of $\mathfrak{f}$, the sum over $y$ is carried out exactly as in (i): the weight is moved from $y$ to $c$
at the cost of $6\delta_{1}$ of the rate, the powers of $L$ are absorbed by \eqref{5.8:eq-paired-sets-a}
and \eqref{5.8:eq-paired-sets-c}, and the remaining sum is bounded by \eqref{5.8:eq-paired-sets-b}. This
is the bound for $H_{0}^{\top}$. The kernels of \cite[(71)]{B-CMP102b} and the kernel of $\Theta$ are bounded
symmetrically in their two arguments up to powers of $\ell$; hence the argument of (i), with the roles of
the two arguments exchanged, applies to $((I+\mathfrak{R})^{-1})^{\top}$ and to $\Theta^{\top}=\Theta$ on
data.
\end{proof}

\begin{remark}
The rate used from Proposition~\ref{5.8:prop-local-decay} is $8\delta_{1}$, while that proposition provides
at least $16\delta_{1}$; the margin $2\delta_{1}=\delta_{0}/16$ left after the weight is what the proof
above uses. The rate $8\delta_{1}$ is spent as
\begin{equation*}
8\delta_{1}=6\delta_{1}+\tfrac{1}{4}\delta_{1}+\tfrac{1}{2}\delta_{1}+\tfrac{5}{4}\delta_{1},
\end{equation*}
namely $6\delta_{1}$ on the weight, $\tfrac14\delta_{1}$ on the powers of $\ell$, $\tfrac12\delta_{1}$ on
the sum \eqref{5.8:eq-paired-sets-b}, and $\tfrac54\delta_{1}$ unused. The weights to which the lemma is
applied in what follows are: the tempered profile majorant $\mathsf{p}$, of some rate $\rho\le 4\delta_{1}$;
$\mathsf{p}L^{-a\lambda}$ with $a\le 1$; and $\mathsf{p}\mathsf{M}$ with $\mathsf{M}$ a tempered majorant
of some rate $\tau$, $\mathsf{M}\in\mathfrak{M}_{\tau}$, where the two rates satisfy
$\rho+\tau\le 5\delta_{1}$. All of them lie inside the rate $6\delta_{1}$ allowed by the lemma.
\end{remark}

Throughout the rest of this section the following standing conventions are in force:
\begin{itemize}
\item the background field $\mathbf{U}$ is real and satisfies the standing regularity condition imposed on background fields;
\item $\mathsf{m}$ is a complex data vector with $|\mathsf{m}| \le \mathfrak{r}_{0}$, and the data $B$ are complex;
\item $\sigma \le 6\delta_{1}$ and $w \in \mathfrak{M}_{\sigma}$, with $\mathfrak{M}_{\sigma}$ as in Definition~\ref{5.8:def-tempered-weights}.
\end{itemize}

The block function $C(\cdot)$ is the nonlinear averaging correction; the letter $C$ without an index always denotes it. The constants $C_{2}$, $C_{3}$ below, and the generic constants written $C$ with an index in this section, depend only on the dimension, on $L$, on $\theta$ and on $M_{1}$; $C_{wt}$ is the constant in the calculus of the tempered weights of Definition~\ref{5.8:def-tempered-weights}, and $C_{W}$ is the constant by which sizes of fields bound sizes of data; the two are distinct. The derivative of the block function $C(\cdot)$ is written $dC(X)$, and its kernel $\mathsf{c}_{X}$. Neither is to be confused with the plaquette commutator $\mathfrak{c}(X)$. For a site $x$, $B(x)$ denotes the block attached to $x$; it is not the datum $B(c)$ of a bond $c$, nor the complex data $B$ above.

For a data bond $c$ with endpoints $c_{-}$, $c_{+}$ put
\[P(c) := B(c_{-}) \cup B(c_{+}),\]
the union of the two blocks attached to the endpoints of $c$. For a field $X$ put
\[|X|_{c} := \ell \sup_{P(c)} |X|.\]
On the restrictions to $P(c)$, the function $X \mapsto |X|_{c}$ is a norm. Hence every ball $\{|X|_{c} < r\}$ is a convex and balanced open subset of a finite-dimensional complex space.

\begin{lemma}[Local analyticity of the nonlinear averaging correction]
\label{5.8:lem-averaging-correction}
Let $c$ be a data bond. Then $C(X)(c)$ is an analytic function of $X \restriction P(c)$ on the domain $|X|_{c} < c_{4}$, and on that domain
\begin{equation}\label{5.8:eq-averaging-correction-1}
|C(X)(c)| \le C_{2}\, |X|_{c}^{2}.
\end{equation}
Here $c_{4} > 0$ and $C_{2}$ are constants.

Suppose $2\varepsilon_{3} \le \tfrac14 c_{4}$. Then there is a constant $C_{3}$ with the following property. For fields $X$, $X_{1}$, $X_{2}$ with $|X|_{c} \le 2\varepsilon_{3}$ and $|X_{i}|_{c} \le 2\varepsilon_{3}$, for every field $v$, and, in the two bounds on the kernel, for every bond $b \subset P(c)$,
\begin{equation}\label{5.8:eq-averaging-correction-a}
\begin{gathered}
|C(X_{1})(c) - C(X_{2})(c)| \le 9 C_{2}\,\varepsilon_{3}\, |X_{1} - X_{2}|_{c},\\
dC(X)v(c) = \sum_{b \subset P(c)} \eta^{4}\, \mathsf{c}_{X}(c,b)\cdot v(b),
\qquad \ell^{3}|\mathsf{c}_{X}(c,b)| \le C_{3}\, |X|_{c},\\
\ell^{3}|\mathsf{c}_{X_{1}}(c,b) - \mathsf{c}_{X_{2}}(c,b)| \le C_{3}\, |X_{1} - X_{2}|_{c}.
\end{gathered}
\end{equation}
The kernel $\mathsf{c}_{X}(c,b)$ is analytic in $X \restriction P(c)$ on $|X|_{c} < c_{4}$ and vanishes at $X = 0$. The set $P(c)$ lies in cells at distance $d \le 2$ from $c$, where $d$ is the distance of Notation~\ref{5.8:not-paired-sets}.
\end{lemma}

\begin{proof}
The analyticity of $C(X)(c)$ in $X \restriction P(c)$ on $|X|_{c} < c_{4}$, and the bound \eqref{5.8:eq-averaging-correction-1}, are \cite[Propositions~4, 5 and~7]{B-CMP98}, in the form in which they are used in \cite[(44), (54)]{B-CMP102b}.

\emph{The first line of \eqref{5.8:eq-averaging-correction-a}.}
Let $|X|_{c} \le 2\varepsilon_{3}$, and let $Y$ be a field with $|Y|_{c} = 1$. For complex $\zeta$ with $|\zeta| < c_{4} - 2\varepsilon_{3}$ we have
\[|X + \zeta Y|_{c} \le 2\varepsilon_{3} + |\zeta| < c_{4}.\]
Hence $g(\zeta) := C(X + \zeta Y)(c)$ is analytic on that disc. Since $4\varepsilon_{3} \le \tfrac12 c_{4} < c_{4}$, the disc contains the closed disc $|\zeta| \le 2\varepsilon_{3}$. On the circle $|\zeta| = 2\varepsilon_{3}$ we have $|X + \zeta Y|_{c} \le 4\varepsilon_{3}$, so \eqref{5.8:eq-averaging-correction-1} gives $|g(\zeta)| \le 16 C_{2}\varepsilon_{3}^{2}$. The Cauchy estimate on this circle then gives
\[|g'(0)| \le \frac{16 C_{2}\varepsilon_{3}^{2}}{2\varepsilon_{3}} = 8 C_{2}\varepsilon_{3}.\]
Since $g'(0) = dC(X)Y(c)$, and $dC(X)Y(c)$ is linear in $Y$, it follows that for every field $Y$
\begin{equation}\label{5.8:eq-averaging-correction-2}
|dC(X)Y(c)| \le 8 C_{2}\,\varepsilon_{3}\, |Y|_{c}.
\end{equation}

Now let $|X_{i}|_{c} \le 2\varepsilon_{3}$ and put $X_{t} := X_{2} + t(X_{1} - X_{2})$ for $0 \le t \le 1$. By convexity of the norm,
\[|X_{t}|_{c} \le (1-t)|X_{2}|_{c} + t|X_{1}|_{c} \le 2\varepsilon_{3}.\]
By the fundamental theorem of calculus along the segment,
\[C(X_{1})(c) - C(X_{2})(c) = \int_{0}^{1} dC(X_{t})(X_{1} - X_{2})(c)\, dt.\]
Applying \eqref{5.8:eq-averaging-correction-2} at each $X_{t}$ bounds the right side by $8 C_{2}\varepsilon_{3}|X_{1} - X_{2}|_{c}$. This is at most the right side of the first line of \eqref{5.8:eq-averaging-correction-a}.

\emph{The second line.}
Since $C(X)(c)$ depends only on $X \restriction P(c)$, the number $dC(X)v(c)$ depends complex-linearly on $v \restriction P(c)$ alone. It can therefore be written uniquely as the sum over the bonds $b \subset P(c)$ displayed in \eqref{5.8:eq-averaging-correction-a}. In that sum $\eta^{4}\mathsf{c}_{X}(c,b)$ is the partial derivative of $C(X)(c)$ with respect to $X(b)$, paired with $v(b)$. This is the definition of the kernel $\mathsf{c}_{X}$.

As a partial derivative of an analytic function of finitely many complex variables, $\mathsf{c}_{X}(c,b)$ is analytic in $X \restriction P(c)$ on $|X|_{c} < c_{4}$.

The kernel vanishes at $X = 0$. Indeed, for $|Y|_{c} = 1$ and $0 < r < c_{4}$, the function $\zeta \mapsto C(\zeta Y)(c)$ is bounded by $C_{2} r^{2}$ on $|\zeta| = r$, by \eqref{5.8:eq-averaging-correction-1}. The Cauchy estimate gives $|dC(0)Y(c)| \le C_{2} r$. Letting $r \to 0$ gives $dC(0)Y(c) = 0$ for every $Y$, hence $\mathsf{c}_{0}(c,b) = 0$ for every $b$.

In the range $|X|_{c} \le 2\varepsilon_{3}$, the representation together with the bound $\ell^{3}|\mathsf{c}_{X}(c,b)| \le C_{3}|X|_{c}$ is \cite[(72)]{B-CMP102b}, where it is obtained from \cite[Propositions~4, 5 and~7]{B-CMP98}; we take this bound from there rather than deriving it here, with the constant $C_{3}$ fixed below. The bound is linear in $|X|_{c}$, in accordance with the vanishing of the kernel at $X = 0$ (Schwarz's lemma).

\emph{The third line.}
By the convexity bound above, every point of the segment $X_{t}$, $0 \le t \le 1$, joining $X_{2}$ to $X_{1}$ satisfies $|X_{t}|_{c} \le 2\varepsilon_{3} \le \tfrac14 c_{4}$, so the segment lies in the domain $|X|_{c} < c_{4}$ on which $X \mapsto \mathsf{c}_{X}(c,b)$ is analytic. The third line is the Cauchy estimate for this analytic function of the restriction of $X$ to the blocks of $P(c)$, along that segment; it rests on a majorant of the kernel on a neighbourhood of the segment, as follows. We use the bound
\[
\ell^{3}|\mathsf{c}_{X'}(c,b)| \le C'\,|X'|_{c}
\qquad\text{for all fields } X' \text{ with } |X'|_{c} \le \tfrac12 c_{4},
\]
with a constant $C'$ depending only on the dimension, on $L$, on $\theta$ and on $M_{1}$. This is a bound for the kernel of the derivative of the block function on the closed ball $|X'|_{c} \le \tfrac12 c_{4}$, which lies inside its analyticity domain $|X'|_{c} < c_{4}$, and we take it from \cite[Propositions~4, 5 and~7]{B-CMP98}; the bound \cite[(72)]{B-CMP102b} is used above only in the range $|X|_{c} \le 2\varepsilon_{3}$ of the second line. Put $C_{3} := 2C'$. Since $|X|_{c} \le 2\varepsilon_{3} \le \tfrac12 c_{4}$ in the second line of \eqref{5.8:eq-averaging-correction-a}, that line holds with this $C_{3}$.

If $X_{1} = X_{2}$ there is nothing to prove. Otherwise put $Y := (X_{1} - X_{2})/|X_{1} - X_{2}|_{c}$, so that $|Y|_{c} = 1$, fix $t \in [0,1]$, and put $h(\zeta) := \ell^{3}\mathsf{c}_{X_{t} + \zeta Y}(c,b)$. For $|\zeta| < \tfrac34 c_{4}$ we have $|X_{t} + \zeta Y|_{c} \le \tfrac14 c_{4} + |\zeta| < c_{4}$, so $h$ is analytic on that disc, which contains the closed disc $|\zeta| \le \tfrac14 c_{4}$. On the circle $|\zeta| = \tfrac14 c_{4}$ we have $|X_{t} + \zeta Y|_{c} \le \tfrac12 c_{4}$, so the bound just quoted gives $|h(\zeta)| \le \tfrac12 c_{4} C'$. The Cauchy estimate on this circle, applied to $h$ as an analytic function with values in a finite-dimensional normed space, gives
\[
|h'(0)| \le \frac{\tfrac12 c_{4} C'}{\tfrac14 c_{4}} = 2C' = C_{3}.
\]
Since $\frac{d}{dt}X_{t} = X_{1} - X_{2} = |X_{1} - X_{2}|_{c}\, Y$, the derivative of $t \mapsto \ell^{3}\mathsf{c}_{X_{t}}(c,b)$ at $t$ equals $|X_{1} - X_{2}|_{c}\, h'(0)$, and is therefore bounded in norm by $C_{3}|X_{1} - X_{2}|_{c}$. Integrating over $t \in [0,1]$ gives
\[
\ell^{3}|\mathsf{c}_{X_{1}}(c,b) - \mathsf{c}_{X_{2}}(c,b)|
\le \int_{0}^{1} C_{3}\,|X_{1} - X_{2}|_{c}\, dt
= C_{3}\,|X_{1} - X_{2}|_{c},
\]
which is the third line of \eqref{5.8:eq-averaging-correction-a}.

Finally, the set $P(c)$, on which all the quantities in \eqref{5.8:eq-averaging-correction-a} depend, lies in cells at distance $d \le 2$ from $c$, where $d$ is the distance of Notation~\ref{5.8:not-paired-sets}.
\end{proof}

\begin{lemma}[The chart in weighted norms]\label{5.8:lem-weighted-chart}
Assume the following.
\begin{itemize}
\item The hypotheses of the weighted calculus, Lemma~\ref{5.8:lem-weighted-calculus}, hold.
Let $\sigma\le 6\delta_{1}$ and $w\in\mathfrak{M}_{\sigma}$ be as in that lemma. All weighted
norms $\|\cdot\|_{w}$ below are those of Definition~\ref{5.8:def-tempered-weights}.
\item The chart $\mathsf{A}_{0}$, $D_{0}$ obeys the following algebra, taken by statement.
For $A'_{1}$, $A'_{2}$ in its domain, put $X_{i}:=\mathsf{A}_{0}A'_{i}$ and
$E:=A'_{1}-A'_{2}$. Then
\begin{equation*}
D_{0}(A'_{i})=C(X_{i}),\qquad X_{1}-X_{2}=E-H_{0}\bigl(D_{0}(A'_{1})-D_{0}(A'_{2})\bigr),
\end{equation*}
and
\begin{equation*}
dD_{0}(A'_{i})=(I+\mathfrak{R}_{i})^{-1}dC(X_{i}),\qquad \mathfrak{R}_{i}:=dC(X_{i})H_{0}.
\end{equation*}
Here $C$ is the block function of the nonlinear averaging correction of
Lemma~\ref{5.8:lem-averaging-correction}.
\end{itemize}
Let $A'_{1}$, $A'_{2}$ satisfy $\sup_{y}N_{y}(A'_{i})\le\varepsilon_{3}$ for $i=1,2$. Let
$X_{i}$, $E$ and $\mathfrak{R}_{i}$ be as above, and put
$\mathfrak{D}_{i}:=dD_{0}(A'_{i})=(I+\mathfrak{R}_{i})^{-1}dC(X_{i})$. Suppose that
\begin{equation}\label{5.8:eq-weighted-chart-a}
40\,(C_{2}+C_{3})\,C_{wt}\,e^{12\delta_{1}}\,\varepsilon_{3}\le 1 .
\end{equation}
Then the following hold. In each item, $C$ without an argument denotes a constant; it is not the block function.
\begin{enumerate}
\item[(a)] With the sizes $|\cdot|_{y}$ we have $\|X_{1}-X_{2}\|_{w}\le 2\|E\|_{w}$. The same
bound holds with the sizes $N_{y}$. Moreover
$\|D_{0}(A'_{1})-D_{0}(A'_{2})\|_{w}\le C\varepsilon_{3}\|E\|_{w}$.
\item[(b)] The operator $(I+\mathfrak{R}_{i})^{-1}$ and its transpose have norm at most $2$ on data
measured in $\|\cdot\|_{w}$.
\item[(c)] For a data vector $Z$, the source $\mathfrak{D}_{i}^{\top}Z$ obeys
\begin{equation*}
\|\mathfrak{D}_{i}^{\top}Z\|_{w}\le C\min\bigl\{\varepsilon_{3}\|Z\|_{w},\ \|X_{i}\|_{w}\sup|Z|\bigr\}.
\end{equation*}
\item[(d)] For a data vector $Z$,
$\|(\mathfrak{D}_{1}-\mathfrak{D}_{2})^{\top}Z\|_{w}\le C\|E\|_{w}\sup|Z|$.
\end{enumerate}
\end{lemma}

\begin{proof}
Write $D_{0,i}:=D_{0}(A'_{i})$. By the choice $\sigma\le6\delta_{1}$ we have
$e^{2\sigma}\le e^{12\delta_{1}}$. This is used throughout to compare with
\eqref{5.8:eq-weighted-chart-a}.

(a) By the algebra of the chart, $D_{0,i}=C(X_{i})$. Apply the first of the local bounds in
\eqref{5.8:eq-averaging-correction-a} at each data bond $c$, with $X_{1}$, $X_{2}$ in place of the
two arguments. This gives $|D_{0,1}(c)-D_{0,2}(c)|\le 9C_{2}\varepsilon_{3}|X_{1}-X_{2}|_{c}$.
The size $|\cdot|_{c}$ reads the blocks of $P(c)$. Passing from it to the weighted norm costs
the factor $e^{2\sigma}$, and we obtain
\begin{equation}\label{5.8:eq-weighted-chart-1}
\|D_{0,1}-D_{0,2}\|_{w}\le 9C_{2}\varepsilon_{3}e^{2\sigma}\|X_{1}-X_{2}\|_{w}.
\end{equation}
By the algebra of the chart, $X_{1}-X_{2}=E-H_{0}(D_{0,1}-D_{0,2})$.
Lemma~\ref{5.8:lem-weighted-calculus} applies to $H_{0}$, and with the sizes $N_{y}$ also to
$\nabla H_{0}$. It bounds the second term by $C_{wt}\|D_{0,1}-D_{0,2}\|_{w}$. Together with
\eqref{5.8:eq-weighted-chart-1} this gives
\begin{equation*}
\|X_{1}-X_{2}\|_{w}\le\|E\|_{w}+9C_{2}C_{wt}e^{2\sigma}\varepsilon_{3}\|X_{1}-X_{2}\|_{w}.
\end{equation*}
By \eqref{5.8:eq-weighted-chart-a}, the coefficient $9C_{2}C_{wt}e^{2\sigma}\varepsilon_{3}$ is at
most $9/40\le\frac12$. Absorbing the last term into the left-hand side gives
$\|X_{1}-X_{2}\|_{w}\le2\|E\|_{w}$. The argument is the same with the sizes $N_{y}$. Inserting
this bound into \eqref{5.8:eq-weighted-chart-1} yields
\begin{equation*}
\|D_{0,1}-D_{0,2}\|_{w}\le 18C_{2}e^{12\delta_{1}}\varepsilon_{3}\|E\|_{w}.
\end{equation*}

(b) We have $\mathfrak{R}_{i}Z=dC(X_{i})(H_{0}Z)$. By the second of the local bounds in
\eqref{5.8:eq-averaging-correction-a}, $dC(X_{i})$ is a sum over the bonds $b\subset P(c)$ with
kernel $\ell^{3}|\mathsf{c}_{X_{i}}(c,b)|\le C_{3}|X_{i}|_{c}$, where $|X_{i}|_{c}\le2\varepsilon_{3}$.
By part~(i) of Lemma~\ref{5.8:lem-weighted-calculus}, $\|H_{0}Z\|_{w}\le C_{wt}\|Z\|_{w}$. Hence
\begin{equation*}
\|\mathfrak{R}_{i}Z\|_{w}\le 4dC_{3}\cdot2\varepsilon_{3}\,e^{2\sigma}C_{wt}\|Z\|_{w}
\le\tfrac12\|Z\|_{w}.
\end{equation*}
The transpose is $\mathfrak{R}_{i}^{\top}=H_{0}^{\top}dC(X_{i})^{\top}$. It is bounded in the same
way, by the second local bound and part~(iii) of Lemma~\ref{5.8:lem-weighted-calculus}.
For $T$ equal to $\mathfrak{R}_{i}$ or $\mathfrak{R}_{i}^{\top}$, the Neumann series
$\sum_{n\ge0}(-T)^{n}$ converges in the norm $\|\cdot\|_{w}$. Its sum has norm at most
$\sum_{n\ge0}2^{-n}=2$, and it represents $(I+\mathfrak{R}_{i})^{-1}$, respectively its
transpose.

(c) Transposing the product $\mathfrak{D}_{i}=(I+\mathfrak{R}_{i})^{-1}dC(X_{i})$ gives
\begin{equation*}
\mathfrak{D}_{i}^{\top}Z=dC(X_{i})^{\top}\tilde{Z},\qquad
\tilde{Z}:=\bigl((I+\mathfrak{R}_{i})^{-1}\bigr)^{\top}Z.
\end{equation*}
By the second local bound in \eqref{5.8:eq-averaging-correction-a}, for every bond $b$,
\begin{equation*}
\ell^{3}\bigl|dC(X_{i})^{\top}\tilde{Z}\bigr|(b)\le\sum_{c:\,b\subset P(c)}C_{3}|X_{i}|_{c}\,|\tilde{Z}(c)|,
\end{equation*}
and at most $4d$ bonds $c$ occur in the sum.

For the first alternative, use $|X_{i}|_{c}\le2\varepsilon_{3}$ and apply (b) at the weight $w$,
which gives $\|\tilde{Z}\|_{w}\le2\|Z\|_{w}$. The passage to the weighted norm costs $e^{2\sigma}$,
so
\begin{equation*}
\|\mathfrak{D}_{i}^{\top}Z\|_{w}\le 16dC_{3}e^{12\delta_{1}}\varepsilon_{3}\|Z\|_{w}.
\end{equation*}
For the second alternative, keep $|X_{i}|_{c}$ in the weighted norm. Apply (b) at the weight
identically equal to one. This weight lies in $\mathfrak{M}_{\sigma}$ by
Definition~\ref{5.8:def-tempered-weights}, and its weighted norm is $\sup|\cdot|$. Thus
$\sup|\tilde{Z}|\le2\sup|Z|$, and
\begin{equation*}
\|\mathfrak{D}_{i}^{\top}Z\|_{w}\le 8dC_{3}e^{12\delta_{1}}\|X_{i}\|_{w}\sup|Z|.
\end{equation*}
The two bounds together give (c).

(d) The definition of $\mathfrak{D}_{2}$ gives $(I+\mathfrak{R}_{2})\mathfrak{D}_{2}=dC(X_{2})$,
that is, $\mathfrak{D}_{2}=dC(X_{2})-dC(X_{2})H_{0}\mathfrak{D}_{2}$. Using this,
\begin{align*}
(I+\mathfrak{R}_{1})(\mathfrak{D}_{1}-\mathfrak{D}_{2})
&=dC(X_{1})-\mathfrak{D}_{2}-dC(X_{1})H_{0}\mathfrak{D}_{2}\\
&=\bigl(dC(X_{1})-dC(X_{2})\bigr)\bigl(I-H_{0}\mathfrak{D}_{2}\bigr).
\end{align*}
Therefore
\begin{equation}\label{5.8:eq-weighted-chart-2}
\mathfrak{D}_{1}-\mathfrak{D}_{2}
=(I+\mathfrak{R}_{1})^{-1}\bigl(dC(X_{1})-dC(X_{2})\bigr)\bigl(I-H_{0}\mathfrak{D}_{2}\bigr).
\end{equation}
Transposing \eqref{5.8:eq-weighted-chart-2} gives
\begin{equation*}
(\mathfrak{D}_{1}-\mathfrak{D}_{2})^{\top}Z
=\bigl(I-\mathfrak{D}_{2}^{\top}H_{0}^{\top}\bigr)\mathsf{g},\qquad
\mathsf{g}:=\bigl(dC(X_{1})-dC(X_{2})\bigr)^{\top}\tilde{Z}_{1},
\end{equation*}
with $\tilde{Z}_{1}:=((I+\mathfrak{R}_{1})^{-1})^{\top}Z$.

The third local bound in \eqref{5.8:eq-averaging-correction-a} controls the kernel of
$dC(X_{1})-dC(X_{2})$ by $C_{3}|X_{1}-X_{2}|_{c}$. Part~(b), at the weight one, gives
$\sup|\tilde{Z}_{1}|\le2\sup|Z|$. As in (c), these yield
\begin{equation*}
\|\mathsf{g}\|_{w}\le 4dC_{3}e^{2\sigma}\|X_{1}-X_{2}\|_{w}\cdot2\sup|Z|.
\end{equation*}
Next we bound $\mathfrak{D}_{2}^{\top}H_{0}^{\top}\mathsf{g}$. Part~(iii) of
Lemma~\ref{5.8:lem-weighted-calculus} bounds $\|H_{0}^{\top}\mathsf{g}\|_{w}$ by a constant times
$\|\mathsf{g}\|_{w}$. The first alternative of (c), applied with $H_{0}^{\top}\mathsf{g}$ in place
of $Z$, bounds $\|\mathfrak{D}_{2}^{\top}H_{0}^{\top}\mathsf{g}\|_{w}$ by $C\varepsilon_{3}$ times
that. Hence $\|(\mathfrak{D}_{1}-\mathfrak{D}_{2})^{\top}Z\|_{w}\le C\|\mathsf{g}\|_{w}$. Finally,
(a) gives $\|X_{1}-X_{2}\|_{w}\le2\|E\|_{w}$, and (d) follows.
\end{proof}

\begin{definition}[Smallness ceilings for the weighted estimates]\label{5.8:not-smallness-ceilings}
Let $C_{wt}$ be the constant of the weighted calculus (Lemma~\ref{5.8:lem-weighted-calculus}).
Let $C_{ch}$ denote the constant written $C$ in the chart lemma (Lemma~\ref{5.8:lem-weighted-chart}),
and let $C_{adj}$ and $C_{s}^{\sharp}$ be the constants bearing these names that are produced by two
weighted estimates of this section. Each of these constants is fixed after $L$, in the statement that
produces it. None of $C_{wt}$, $C_{ch}$, $C_{adj}$, $C_{s}^{\sharp}$ depends on $\varepsilon_{3}$,
$\mathfrak{r}_{0}$ or $a_{R}^{\sharp}$, since all four precede these ceilings in the order of choice,
and the conditions below are imposed only after all of them have been fixed. Let
$\mathfrak{r}_{0}$, $\varepsilon_{D}$, $a_{R}'$ be ceilings, depending on $(d,L,\theta,M_{1})$
only, as given by the statement on the existence of the ceilings $\mathfrak{r}_{0}$,
$\varepsilon_{D}$, $a_{R}'$. Let $\varepsilon_{3}$ be the ceiling on $\sup_{y}N_{y}(A'_{i})$ in
the chart lemma (Lemma~\ref{5.8:lem-weighted-chart}), subject to the one-sided condition imposed
on it there.

The \emph{smallness ceilings for the weighted estimates} are the numbers $\varepsilon_{3}$,
$\mathfrak{r}_{0}$ and $a_{R}^{\sharp}$ (the ceiling $\varepsilon_{D}$ is not changed here), subject
to the one-sided conditions
\begin{equation}\label{5.8:eq-smallness-ceilings-a}
\begin{gathered}
2C_{wt}\,(C_{ch}+C_{adj})\,(\varepsilon_{3}+\mathfrak{r}_{0}+a_{R}^{\sharp}) \le 1,\\
a_{R}^{\sharp} \le a_{R}', \qquad 70\,d\,C_{s}^{\sharp}\,a_{R}^{\sharp} \le \varepsilon_{3}.
\end{gathered}
\end{equation}

They are fixed by one further minimisation, made after $L$ and after $C_{ch}$, $C_{adj}$,
$C_{s}^{\sharp}$ have been fixed: $\mathfrak{r}_{0}$ and $\varepsilon_{3}$ are replaced by
smaller values, and $a_{R}^{\sharp}\le a_{R}'$ is chosen, so that
\eqref{5.8:eq-smallness-ceilings-a} holds. The re-minimised values depend only on $d$, on
$(d,L,\theta,M_{1})$ through the ceilings $\mathfrak{r}_{0}$, $a_{R}'$ given above, on the value of
$\varepsilon_{3}$ given by Lemma~\ref{5.8:lem-weighted-chart}, and on the parameters on which
$C_{wt}$, $C_{ch}$, $C_{adj}$, $C_{s}^{\sharp}$ depend; from here on $\mathfrak{r}_{0}$ and
$\varepsilon_{3}$ denote the re-minimised values. This minimisation is made before $\gamma$, which
is chosen last. Such a choice exists, and \eqref{5.8:eq-smallness-ceilings-a} is preserved under any
further decrease of $\mathfrak{r}_{0}$ and $a_{R}^{\sharp}$; both facts are proved below.
\end{definition}

\begin{proof}
We first show that the conditions are preserved when $\mathfrak{r}_{0}$ and $a_{R}^{\sharp}$
are decreased with $\varepsilon_{3}$ held fixed.
\begin{itemize}
\item The left side of the first inequality is a positive multiple of
$\varepsilon_{3}+\mathfrak{r}_{0}+a_{R}^{\sharp}$. It therefore does not increase when
$\mathfrak{r}_{0}$ or $a_{R}^{\sharp}$ decreases (nor when $\varepsilon_{3}$ decreases).
\item The left sides of the second and third inequalities are nondecreasing in $a_{R}^{\sharp}$.
\end{itemize}

We next show that the conditions can be met by a single decrease. Since $C_{wt}$, $C_{ch}$ and
$C_{adj}$ do not depend on $\varepsilon_{3}$, we may decrease the ceiling
$\varepsilon_{3}$ of Lemma~\ref{5.8:lem-weighted-chart}, if necessary, so that
\begin{equation*}
6\,C_{wt}\,(C_{ch}+C_{adj})\,\varepsilon_{3} \le 1 .
\end{equation*}
The condition imposed on $\varepsilon_{3}$ in Lemma~\ref{5.8:lem-weighted-chart} bounds a positive
multiple of $\varepsilon_{3}$ by one, and therefore continues to hold after this decrease.
The ceiling $\mathfrak{r}_{0}$ is provided by the statement on the existence of the ceilings
$\mathfrak{r}_{0}$, $\varepsilon_{D}$, $a_{R}'$, and the set of values of $\mathfrak{r}_{0}$ for which
the conclusions of that statement hold is a down-set; so they remain valid after the following
replacement.
Replace $\mathfrak{r}_{0}$ by $\min\{\mathfrak{r}_{0},\varepsilon_{3}\}$, and put
\begin{equation*}
a_{R}^{\sharp} := \min\Bigl\{a_{R}',\ \varepsilon_{3},\
\frac{\varepsilon_{3}}{70\,d\,\max\{1,C_{s}^{\sharp}\}}\Bigr\}.
\end{equation*}
The second inequality of \eqref{5.8:eq-smallness-ceilings-a} then holds by the definition of
$a_{R}^{\sharp}$, and so does the third, since $a_{R}^{\sharp}\ge 0$ and
\begin{equation*}
70\,d\,C_{s}^{\sharp}\,a_{R}^{\sharp} \le 70\,d\,\max\{1,C_{s}^{\sharp}\}\,a_{R}^{\sharp}
\le \varepsilon_{3}.
\end{equation*}
For the first, we have
$\varepsilon_{3}+\mathfrak{r}_{0}+a_{R}^{\sharp} \le 3\varepsilon_{3}$, and hence
\begin{equation*}
2C_{wt}(C_{ch}+C_{adj})(\varepsilon_{3}+\mathfrak{r}_{0}+a_{R}^{\sharp})
\le 6C_{wt}(C_{ch}+C_{adj})\,\varepsilon_{3} \le 1. \qedhere
\end{equation*}
\end{proof}

\begin{lemma}[Critical-point size at a tempered profile]\label{5.8:lem-critical-size}
Assume the weighted calculus, Lemma~\ref{5.8:lem-weighted-calculus}, together with its hypotheses.
Assume also that the constrained critical point $\mathfrak{X}(B;\mathbf{U},\mathsf{m})$ exists, is
unique and is analytic in $(B,\mathsf{m})$, and that on the constraint set $d\mathcal{N}$ satisfies
the field identity \eqref{5.8:eq-critical-size-2} below. Then there are a constant $C_{s}^{\sharp}$ and
ceilings such that the following holds. If $\mathsf{p}\in\mathfrak{M}_{\sigma}$,
$|B|\le\mathsf{p}$ and $\sup\mathsf{p}\le a_{R}^{\sharp}$, then the constrained critical point
$X:=\mathfrak{X}(B;\mathbf{U},\mathsf{m})$ and its Lagrange multiplier $\mu_{B}$ obey, at every
$y$,
\begin{equation}\label{5.8:eq-critical-size-1}
N^{\Delta}_{y}(X),\quad \ell^{2}|D^{*}X|,\quad \ell^{2+\beta}[D^{*}X]_{\beta},\quad
|\mu_{B}(c)|_{c\ni y}\ \le\ C_{s}^{\sharp}\,\mathsf{p}(y).
\end{equation}
\end{lemma}

\begin{proof}
Existence and uniqueness of $X$, and its analyticity in $(B,\mathsf{m})$, are among the
hypotheses. Write $A'$ for the chart
coordinate of $X$, so that $X=\mathsf{A}_{0}A'$. The coordinate $A'$ lies on the constraint set
$\{(\ell Q)A'=B\}$.

By hypothesis the field identity holds on this set, read as an identity between fields
tested on $\ker(\ell Q)$:
\begin{equation}\label{5.8:eq-critical-size-2}
d\mathcal{N}(A')=\bigl(1-\mathfrak{D}^{\top}H_{0}^{\top}\bigr)\mathsf{v}(X)
+\mathfrak{D}^{\top}Z(A')=:\mathsf{g}(A').
\end{equation}
Here
\begin{align*}
\mathsf{v}(X)&:=D^{*}\mathfrak{c}(X)+d\mathfrak{c}(X)^{\dagger}DX+\nabla V_{0}'(X),\\
Z(A')&:=\Theta\bigl(D_{0}(A')-B\bigr)-\mathsf{m},
\end{align*}
with $V_{0}'$ as in \cite[(40), (90)--(96)]{B-CMP102b}.
The force $\mathsf{v}$ is local. By \cite[(40), (90)--(96)]{B-CMP102b} it obeys
\begin{equation}\label{5.8:eq-critical-size-3}
\ell^{3}|\mathsf{v}(X)|(y)\le C_{\mathsf{v}}N_{y}(X)^{2}\le 2C_{\mathsf{v}}\varepsilon_{3}N_{y}(X),
\end{equation}
with a constant $C_{\mathsf{v}}$ depending only on the constants of
\cite[(40), (90)--(96)]{B-CMP102b}.
The vector $Z$ satisfies
\begin{equation*}
\sup|Z|\le C\bigl(\varepsilon_{3}^{2}+\sup|B|\bigr)+\mathfrak{r}_{0}.
\end{equation*}

Write $\|\cdot\|_{\mathsf{p}}$ for the weighted supremum norm of
Definition~\ref{5.8:def-tempered-weights} with weight $w=\mathsf{p}$. Write
$\|\cdot\|^{N}_{\mathsf{p}}$ for the same norm built on the local sizes $N_{y}$. We bound the two
terms of \eqref{5.8:eq-critical-size-2} using three results:
\begin{itemize}
\item part (a) of the chart lemma, Lemma~\ref{5.8:lem-weighted-chart};
\item part (c) of the same lemma in its second alternative, $\|\mathfrak{D}^{\top}Z\|_{w}\le
C\|X\|_{w}\sup|Z|$;
\item part (iii) of Lemma~\ref{5.8:lem-weighted-calculus}.
\end{itemize}
Together with \eqref{5.8:eq-critical-size-3} and the bound on $\sup|Z|$, they give
\begin{equation}\label{5.8:eq-critical-size-4}
\|\mathsf{g}(A')\|_{\mathsf{p}}\le C\bigl(\varepsilon_{3}+\mathfrak{r}_{0}+\sup|B|\bigr)\,
\|A'\|^{N}_{\mathsf{p}}.
\end{equation}

The coordinate $A'$ solves the fixed-point equation
\begin{equation*}
A'=H_{0}B-\mathcal{G}_{0}\,\mathsf{g}(A').
\end{equation*}
Apply Lemma~\ref{5.8:lem-weighted-calculus} to $H_{0}$ and to $\mathcal{G}_{0}$. The hypothesis
$|B|\le\mathsf{p}$ bounds the first term. The chain $\sup|B|\le\sup\mathsf{p}\le a_{R}^{\sharp}$,
used in \eqref{5.8:eq-critical-size-4}, bounds the second. Hence
\begin{equation}\label{5.8:eq-critical-size-5}
\|A'\|^{N}_{\mathsf{p}}\le C_{wt}+C_{wt}\,C\bigl(\varepsilon_{3}+\mathfrak{r}_{0}
+a_{R}^{\sharp}\bigr)\,\|A'\|^{N}_{\mathsf{p}}.
\end{equation}
The lattice is finite and $\mathsf{p}>0$, so every norm in \eqref{5.8:eq-critical-size-5} is
finite. Under the smallness ceiling \eqref{5.8:eq-smallness-ceilings-a} the coefficient of
$\|A'\|^{N}_{\mathsf{p}}$ on the right is at most $\tfrac12$. The last term can therefore be
absorbed into the left-hand side, and
\begin{equation*}
\|A'\|^{N}_{\mathsf{p}}\le 2C_{wt}.
\end{equation*}

It remains to treat the H\"older and Laplacian members of $N^{\Delta}_{y}$ for the two terms
$H_{0}B$ and $\mathcal{G}_{0}\mathsf{g}(A')$. These are given by part (ii) of
Lemma~\ref{5.8:lem-weighted-calculus}. For the propagator $G_{0}$ of \cite{B-CMP99b} that enters
$H_{0}$ and $\mathcal{G}_{0}$, sources in the supremum norm suffice, by \cite[(3.43)]{B-CMP99b}
and the fourth member of \cite[(3.42)]{B-CMP99b}.

The chart reads
\begin{equation*}
X=\mathsf{A}_{0}A'=A'-H_{0}D_{0}(A'),\qquad
\|D_{0}(A')\|_{\mathsf{p}}\le C\varepsilon_{3}\|A'\|_{\mathsf{p}}.
\end{equation*}
Together with the bounds on $A'$ just obtained, this gives the bound on $N^{\Delta}_{y}(X)$ in
\eqref{5.8:eq-critical-size-1}. The field $D^{*}X$ is read off from $\nabla X$. Its two members
$\ell^{2}|D^{*}X|$ and $\ell^{2+\beta}[D^{*}X]_{\beta}$ are therefore bounded in the same way.

For the multiplier, recall that the derivative of the nonlinear averaging constraint along the
chart is the linear one: $d\mathbf{Q}_{\mathbf{U}}(\eta X)\,d\mathsf{A}_{0}=\ell Q$. The critical
point therefore satisfies, for all $v$ (not only for $v\in\ker(\ell Q)$),
\begin{equation*}
\langle v,TA'\rangle+d\mathcal{N}(A')[v]=\langle(\ell Q)v,\mu_{B}\rangle,
\end{equation*}
where $T$ is the operator of the quadratic part of the functional in the chart, so that the left
side is the derivative of the functional at $A'$ in the direction $v$.
Let $\delta_{c}$ be the datum equal to one at the cell $c$ and zero elsewhere.
Taking $v=H_{0}\delta_{c}$ gives
\begin{equation}\label{5.8:eq-critical-size-6}
\mu_{B}(c)=\bigl(\Theta(B-D_{0})\bigr)(c)+\bigl(H_{0}^{\top}\mathsf{g}\bigr)(c).
\end{equation}
At $c\ni y$ the right side of \eqref{5.8:eq-critical-size-6} is estimated by three inputs:
\begin{itemize}
\item the hypothesis $|B|\le\mathsf{p}$;
\item the bound on $D_{0}(A')$ above;
\item Lemma~\ref{5.8:lem-weighted-calculus} applied to $H_{0}^{\top}$ and
\eqref{5.8:eq-critical-size-4}.
\end{itemize}
Collecting the constants of all these steps into $C_{s}^{\sharp}$ gives
\eqref{5.8:eq-critical-size-1}.
\end{proof}

For the slot in which $X$ is the response $\mathbb{H}$, and with $\mathsf{p}:=p$, the profile at
rate $\delta$, \eqref{5.8:eq-critical-size-1} is the response-size statement
$N^{\Delta}_{y}(\mathbb{H})\le C_{s}^{\sharp}p(y)$: the size of the response is controlled at the
full rate $\delta$.

\begin{lemma}[Weighted stability of the critical point]\label{5.8:lem-critical-stability}
Assume the weighted calculus, Lemma~\ref{5.8:lem-weighted-calculus}, the smallness ceilings
\eqref{5.8:eq-smallness-ceilings-a}, and the algebra of the
stability lemma for the constrained critical point in unweighted norms. This algebra consists of
two parts: the difference equation satisfied by the critical point and a second one-form, namely
the constraint equation for the difference and the variational identity on the kernel of the
linearised constraint; and the expression of the difference
of the differentials of the functional $V_{0}$ of that lemma as a sum of three terms, each
carrying the plaquette commutator $\mathfrak{c}$ or its derivative. Let $B$ be a datum with
$\sup|B|\le a_{R}^{\sharp}$, and let $X:=\mathfrak{X}(B;\mathbf{U},\mathsf{m})$. Let
$\tilde{X}$ be a one-form with $\ell|\tilde{X}|\le\varepsilon_{3}$ and
$\ell^{2}|\nabla\tilde{X}|\le\varepsilon_{3}$. Let $\tilde{\beta}$ be a datum and $\kappa$ a
profile such that
\begin{equation*}
\mathbf{Q}_{\mathbf{U}}(\eta\tilde{X})=B+\tilde{\beta},\qquad
d\Sigma_{\mathbf{U},\mathsf{m}}(\tilde{X})\in\sharp(\kappa)\ \text{ on }\
\ker d\mathbf{Q}_{\mathbf{U}}(\eta\tilde{X}).
\end{equation*}
Then for every $\sigma\le 6\delta_{1}$, every $w\in\mathfrak{M}_{\sigma}$ and every point $x$,
\begin{equation*}
|\tilde{X}-X|_{x}\le C_{S}^{\sharp}\,w(x)\,\bigl\||\tilde{\beta}|+\kappa\bigr\|_{w}.
\end{equation*}
In particular, let $B^{1}$, $B^{2}$ be data with $\sup|B^{i}|\le a_{R}^{\sharp}$ for $i=1,2$,
and write $\mathfrak{X}(B^{i})$ for $\mathfrak{X}(B^{i};\mathbf{U},\mathsf{m})$. Then
\begin{equation}\label{5.8:eq-critical-stability-b}
\|\mathfrak{X}(B^{2})-\mathfrak{X}(B^{1})\|_{w}\le C_{S}^{\sharp}\,\|B^{2}-B^{1}\|_{w}.
\end{equation}
\end{lemma}

\begin{proof}
We write $\Sigma$ and $\mathbf{Q}$ for $\Sigma_{\mathbf{U},\mathsf{m}}$ and $\mathbf{Q}_{\mathbf{U}}$.
The letter $C$ denotes a constant whose value may change from one occurrence to the next.

\emph{The difference equation.} Put
\begin{equation*}
\tilde{A}':=\mathsf{A}_{0}^{-1}\tilde{X},\qquad A':=\mathsf{A}_{0}^{-1}X,\qquad
E:=\tilde{A}'-A'.
\end{equation*}
Write $\tilde{\mathfrak{D}}$ for the operator $\mathfrak{D}_{i}$ of the chart in weighted norms,
Lemma~\ref{5.8:lem-weighted-chart}, taken at $\tilde{A}'$, that is,
$\tilde{\mathfrak{D}}=(I+dC(\tilde{X})H_{0})^{-1}dC(\tilde{X})$. The difference equation of the
unweighted stability lemma consists of two statements. First, $(\ell Q)E=\tilde{\beta}$. Second,
for every $v\in\ker(\ell Q)$,
\begin{equation}\label{5.8:eq-critical-stability-1}
\langle v,TE\rangle+\bigl(d\mathcal{N}(\tilde{A}')-d\mathcal{N}(A')\bigr)[v]
=d\Sigma(\tilde{X})[t_{v}],\qquad
t_{v}:=v-H_{0}\tilde{\mathfrak{D}}v\in\ker d\mathbf{Q}(\eta\tilde{X}).
\end{equation}
Here $T$ is the linear operator of the difference equation of the unweighted stability lemma,
taken by statement. Write $d\Sigma(\tilde{X})\circ t$
for the functional $v\mapsto d\Sigma(\tilde{X})[t_{v}]$ on $\ker(\ell Q)$. Combining
$(\ell Q)E=\tilde{\beta}$ with \eqref{5.8:eq-critical-stability-1} gives the representation
\begin{equation}\label{5.8:eq-critical-stability-2}
E=H_{0}\tilde{\beta}+\mathcal{G}_{0}\bigl\{d\Sigma(\tilde{X})\circ t
-\bigl(d\mathcal{N}(\tilde{A}')-d\mathcal{N}(A')\bigr)\bigr\}.
\end{equation}
Steps (i) and (ii) place the two functionals in the braces in formats $\sharp(\cdot)$ with profiles
proportional to $w$. Step (iii) then reads off $E$.

\emph{Step (i): the force.} By hypothesis, the restriction of $d\Sigma(\tilde{X})$ to
$\ker d\mathbf{Q}(\eta\tilde{X})$ is given by a functional
\begin{equation*}
\mathfrak{f}:=\langle\nabla\cdot,s_{1}\rangle+\langle\cdot,s_{0}\rangle\in\sharp(\kappa),
\end{equation*}
whose densities $s_{1}$, $s_{0}$ are controlled by the profile $\kappa$ as in the format
$\sharp(\kappa)$. The formula for $\mathfrak{f}$ defines it on all vectors. Since
$t_{v}\in\ker d\mathbf{Q}(\eta\tilde{X})$ and $t_{v}=v-H_{0}\tilde{\mathfrak{D}}v$,
\begin{equation*}
d\Sigma(\tilde{X})[t_{v}]=\mathfrak{f}[v]-\mathfrak{f}[H_{0}\tilde{\mathfrak{D}}v]
=\mathfrak{f}[v]-\langle v,\tilde{\mathfrak{D}}^{\top}H_{0}^{\top}\mathfrak{f}\rangle .
\end{equation*}
We bound the two terms separately. For the first, the definition of $\|\cdot\|_{w}$ gives
$\kappa\le\|\kappa\|_{w}\,w$ pointwise, so $\mathfrak{f}\in\sharp(\|\kappa\|_{w}w)$. For the
second, we apply Lemma~\ref{5.8:lem-weighted-calculus}(iii) to the data vector
$H_{0}^{\top}\mathfrak{f}$. We then apply the chart lemma, Lemma~\ref{5.8:lem-weighted-chart}(c),
to the source $\tilde{\mathfrak{D}}^{\top}H_{0}^{\top}\mathfrak{f}$; its first alternative
supplies a factor $\varepsilon_{3}$. Together,
\begin{equation}\label{5.8:eq-critical-stability-3}
d\Sigma(\tilde{X})\circ t\in\sharp\bigl((1+C\varepsilon_{3})\|\kappa\|_{w}\,w\bigr).
\end{equation}

\emph{Step (ii): the nonlinear part.} Put $A'_{1}:=\tilde{A}'$ and $A'_{2}:=A'$. We claim
\begin{equation}\label{5.8:eq-critical-stability-a}
d\mathcal{N}(A'_{1})-d\mathcal{N}(A'_{2})\in
\sharp\bigl(C_{adj}(\varepsilon_{3}+\mathfrak{r}_{0}+a_{R}^{\sharp})\|E\|_{w}\,w\bigr)
\quad\text{on }\ker(\ell Q),
\end{equation}
and that only the sizes $|E|_{y}$ of $E$ are read in it. We use the notation of
Lemma~\ref{5.8:lem-weighted-chart}:
$X_{i}:=\mathsf{A}_{0}A'_{i}$, so that $X_{1}=\tilde{X}$ and $X_{2}=X$, and
$\mathfrak{D}_{i}:=dD_{0}(A'_{i})$, so that $\mathfrak{D}_{1}=\tilde{\mathfrak{D}}$. For a data
bond $c$, let $\delta_{c}$ denote the data vector supported at the bond $c$ used in the unweighted
stability lemma. For
$v\in\ker(\ell Q)$ write
\begin{equation}\label{5.8:eq-critical-stability-4}
\begin{gathered}
d\mathcal{N}(A')[v]=dV_{0}(X)[v]+\langle\mathfrak{D}v,\mathsf{Z}(A')\rangle,\\
\mathsf{Z}(A'):=\Theta\bigl(D_{0}(A')-(\ell Q)A'\bigr)-\mathsf{m}-\mathsf{h}(X),\qquad
\mathsf{h}(X)(c):=dV_{0}(X)[H_{0}\delta_{c}].
\end{gathered}
\end{equation}
Here $X=\mathsf{A}_{0}A'$ and $\mathfrak{D}=dD_{0}(A')$. Put $\mathsf{Z}_{i}:=\mathsf{Z}(A'_{i})$
and $E_{X}:=X_{1}-X_{2}$. Subtract \eqref{5.8:eq-critical-stability-4} at $A'_{2}$ from the same
identity at $A'_{1}$, and add and subtract $\langle\mathfrak{D}_{1}v,\mathsf{Z}_{2}\rangle$. This
gives
\begin{equation}\label{5.8:eq-critical-stability-5}
\begin{aligned}
\bigl(d\mathcal{N}(A'_{1})-d\mathcal{N}(A'_{2})\bigr)[v]
&=\bigl(dV_{0}(X_{1})-dV_{0}(X_{2})\bigr)[v]
+\langle\mathfrak{D}_{1}v,\mathsf{Z}_{1}-\mathsf{Z}_{2}\rangle\\
&\quad+\langle(\mathfrak{D}_{1}-\mathfrak{D}_{2})v,\mathsf{Z}_{2}\rangle .
\end{aligned}
\end{equation}
We treat the three terms in turn.

$(\alpha)$ \emph{The difference of $dV_{0}$.} We show that
$dV_{0}(X_{1})-dV_{0}(X_{2})\in\sharp(C\varepsilon_{3}|E_{X}|_{y})$ locally. That is, the
densities representing it at a point $y$ are bounded by $C\varepsilon_{3}$ times the size of
$E_{X}$ at $y$. Write $u$ for the vector on which the differentials act. By the algebra of the
unweighted stability lemma, the difference is the sum of three terms.
\begin{itemize}
\item The first is the term carrying $\mathfrak{c}(X_{1})-\mathfrak{c}(X_{2})$. By Cauchy's
inequality applied to \cite[(40)]{B-CMP102b}, it satisfies
$|\mathfrak{c}(X_{1})-\mathfrak{c}(X_{2})|\le C|X||E_{X}|$; it is paired with a gradient of $u$,
so it contributes a density of the type $s_{1}$.
\item The second is $\langle DX_{2},(d\mathfrak{c}(X_{1})-d\mathfrak{c}(X_{2}))u\rangle$. By
Cauchy's inequality applied to \cite[(40)]{B-CMP102b}, the difference
$d\mathfrak{c}(X_{1})-d\mathfrak{c}(X_{2})$ is controlled by the size of $E_{X}$; the term is
paired with $u$ itself and contributes a density of the type $s_{0}$.
\item The third is $\langle DE_{X},d\mathfrak{c}(X_{1})u\rangle$. It equals
$\langle E_{X},D^{*}(d\mathfrak{c}(X_{1})u)\rangle$, with
\begin{equation*}
|D^{*}(d\mathfrak{c}(X_{1})u)|\le C\bigl(|\nabla X_{1}||u|+|X_{1}||\nabla u|\bigr).
\end{equation*}
So it contributes an $s_{0}$-density bounded by $C|\nabla X_{1}||E_{X}|$ and an $s_{1}$-density
bounded by $C|X_{1}||E_{X}|$.
\end{itemize}
In each term, $E_{X}$ is paired with a factor controlled by a scaled size of $X_{1}$ or of
$X_{2}$: $|X|$ in the first, the derivative $DX_{2}$ in the second, $|\nabla X_{1}|$ and
$|X_{1}|$ in the third. The scaled sizes of $X_{1}=\tilde{X}$ are at most $\varepsilon_{3}$ by
hypothesis. For $X_{2}=X$, the constant profile $\mathsf{p}\equiv a_{R}^{\sharp}$ belongs to
$\mathfrak{M}_{\sigma}$ and majorises $|B|$, so the critical-point size lemma,
Lemma~\ref{5.8:lem-critical-size}, gives $\ell|X|,\ \ell^{2}|\nabla X|\le C_{s}^{\sharp}a_{R}^{\sharp}$,
and this is at most $\varepsilon_{3}$ by the ceiling $70dC_{s}^{\sharp}a_{R}^{\sharp}\le\varepsilon_{3}$
of \eqref{5.8:eq-smallness-ceilings-a}. This gives the profile $C\varepsilon_{3}|E_{X}|_{y}$.

$(\beta)$ \emph{The difference of the $\mathsf{Z}$.} We have
$\langle\mathfrak{D}_{1}v,\mathsf{Z}_{1}-\mathsf{Z}_{2}\rangle
=\langle v,\mathfrak{D}_{1}^{\top}(\mathsf{Z}_{1}-\mathsf{Z}_{2})\rangle$. The data vector
$\mathsf{Z}_{1}-\mathsf{Z}_{2}$ obeys $\|\mathsf{Z}_{1}-\mathsf{Z}_{2}\|_{w}\le C\|E\|_{w}$. Its
parts are bounded as follows. The term $\mathsf{m}$ is the same in $\mathsf{Z}_{1}$ and
$\mathsf{Z}_{2}$ and cancels. The operator $\Theta$ is bounded in the weighted norms by
Lemma~\ref{5.8:lem-weighted-calculus}. The difference $D_{0}(A'_{1})-D_{0}(A'_{2})$ is bounded by
Lemma~\ref{5.8:lem-weighted-chart}(a). The term $(\ell Q)E$ obeys
$|(\ell Q)E(c)|\le C|E|_{c}$. Finally, by the definition of $\mathsf{h}$,
$\mathsf{h}(X_{1})-\mathsf{h}(X_{2})$ is $H_{0}^{\top}$ applied to the functional of $(\alpha)$.
The chart lemma, Lemma~\ref{5.8:lem-weighted-chart}(c), first alternative, applied to
$Z:=\mathsf{Z}_{1}-\mathsf{Z}_{2}$, then gives
\begin{equation*}
\|\mathfrak{D}_{1}^{\top}(\mathsf{Z}_{1}-\mathsf{Z}_{2})\|_{w}\le C\varepsilon_{3}\|E\|_{w}.
\end{equation*}

$(\gamma)$ \emph{The difference of the $\mathfrak{D}$.} We have
$\langle(\mathfrak{D}_{1}-\mathfrak{D}_{2})v,\mathsf{Z}_{2}\rangle
=\langle v,(\mathfrak{D}_{1}-\mathfrak{D}_{2})^{\top}\mathsf{Z}_{2}\rangle$.
Lemma~\ref{5.8:lem-weighted-chart}(d) gives
\begin{equation*}
\|(\mathfrak{D}_{1}-\mathfrak{D}_{2})^{\top}\mathsf{Z}_{2}\|_{w}
\le C\|E\|_{w}\sup|\mathsf{Z}_{2}|,
\end{equation*}
and for the data vector $\mathsf{Z}_{2}$, defined by \eqref{5.8:eq-critical-stability-4} at
$A'_{2}$, we use the bound
\begin{equation*}
\sup|\mathsf{Z}_{2}|\le C(\varepsilon_{3}+a_{R}^{\sharp})+\mathfrak{r}_{0}.
\end{equation*}

It remains to combine $(\alpha)$, $(\beta)$ and $(\gamma)$ in \eqref{5.8:eq-critical-stability-5}.
By the definition of the weighted norm, $|E_{X}|_{y}\le\|E_{X}\|_{w}w(y)$. By
Lemma~\ref{5.8:lem-weighted-chart}(a), $\|E_{X}\|_{w}\le 2\|E\|_{w}$. So the profile of
$(\alpha)$ is at most $2C\varepsilon_{3}\|E\|_{w}w$. The profiles of $(\beta)$ and $(\gamma)$ are
$C\varepsilon_{3}\|E\|_{w}w$ and
$C\bigl(C(\varepsilon_{3}+a_{R}^{\sharp})+\mathfrak{r}_{0}\bigr)\|E\|_{w}w$. Their sum is at most
$C_{adj}(\varepsilon_{3}+\mathfrak{r}_{0}+a_{R}^{\sharp})\|E\|_{w}w$, with $C_{adj}$ the
resulting constant. This is \eqref{5.8:eq-critical-stability-a}. In all three terms $E$ enters
only through $|E_{X}|_{y}$, $|E|_{c}$ and $\|E\|_{w}$, that is, through its sizes.

\emph{Step (iii): absorption.} Apply Lemma~\ref{5.8:lem-weighted-calculus}(ii) to the operators
$H_{0}$ and $\mathcal{G}_{0}$ in \eqref{5.8:eq-critical-stability-2}, using
\eqref{5.8:eq-critical-stability-3} and \eqref{5.8:eq-critical-stability-a}. This gives
\begin{equation*}
\|E\|_{w}\le C_{wt}\|\tilde{\beta}\|_{w}+C_{wt}(1+C\varepsilon_{3})\|\kappa\|_{w}
+C_{wt}C_{adj}(\varepsilon_{3}+\mathfrak{r}_{0}+a_{R}^{\sharp})\|E\|_{w}.
\end{equation*}
The first inequality of \eqref{5.8:eq-smallness-ceilings-a}, in which the constant added to
$C_{adj}$ is non-negative, gives
\begin{equation*}
C_{wt}C_{adj}(\varepsilon_{3}+\mathfrak{r}_{0}+a_{R}^{\sharp})\le\tfrac{1}{2}.
\end{equation*}
On a finite lattice, with $w>0$, the quantity $\|E\|_{w}$ is finite. So the last term can be
absorbed into the left side:
\begin{equation*}
\|E\|_{w}\le 2C_{wt}\|\tilde{\beta}\|_{w}+2C_{wt}(1+C\varepsilon_{3})\|\kappa\|_{w}.
\end{equation*}
Now $\tilde{X}-X=X_{1}-X_{2}=E_{X}$, and Lemma~\ref{5.8:lem-weighted-chart}(a) gives
$\|\tilde{X}-X\|_{w}\le 2\|E\|_{w}$. Since $\kappa\ge 0$, both $\|\tilde{\beta}\|_{w}$ and
$\|\kappa\|_{w}$ are at most $\bigl\||\tilde{\beta}|+\kappa\bigr\|_{w}$. Finally,
$|\tilde{X}-X|_{x}\le\|\tilde{X}-X\|_{w}\,w(x)$. This yields the first assertion, with a constant
$C_{S}^{\sharp}$ determined by $C_{wt}$ and the constants $C$ above.

For \eqref{5.8:eq-critical-stability-b}, take $B:=B^{1}$, so $X=\mathfrak{X}(B^{1})$, and
$\tilde{X}:=\mathfrak{X}(B^{2})$. As for $X$ in step (ii)$(\alpha)$, the constant profile
$\mathsf{p}\equiv a_{R}^{\sharp}$ majorises $|B^{2}|$, so Lemma~\ref{5.8:lem-critical-size} and the
ceiling $70dC_{s}^{\sharp}a_{R}^{\sharp}\le\varepsilon_{3}$ of
\eqref{5.8:eq-smallness-ceilings-a} give
$\ell|\tilde{X}|,\ \ell^{2}|\nabla\tilde{X}|\le C_{s}^{\sharp}a_{R}^{\sharp}\le\varepsilon_{3}$. Then
$\mathbf{Q}(\eta\tilde{X})=B^{2}=B+\tilde{\beta}$ with $\tilde{\beta}=B^{2}-B^{1}$. Since
$\tilde{X}$ is the constrained critical point, $d\Sigma(\tilde{X})$ vanishes on
$\ker d\mathbf{Q}(\eta\tilde{X})$, so we may take $\kappa=0$. The first assertion then gives,
at every $x$,
\begin{equation*}
|\mathfrak{X}(B^{2})-\mathfrak{X}(B^{1})|_{x}\le C_{S}^{\sharp}\,w(x)\,\|B^{2}-B^{1}\|_{w}.
\end{equation*}
Dividing by $w(x)$ and taking the supremum over $x$ gives
\eqref{5.8:eq-critical-stability-b}.

The proof nowhere uses that $B$, $X$, $\tilde{X}$, $\mathsf{m}$ or $u$ are real. The operators
are taken at the real background field $\mathbf{U}$. Every other step is either an identity
between holomorphic functions or a supremum estimate.
\end{proof}

\begin{lemma}[Gauge component of the critical point]\label{5.8:lem-gauge-component}
Assume the following three statements:
\begin{enumerate}
\item[(a)] the weighted calculus lemma, Lemma~\ref{5.8:lem-weighted-calculus};
\item[(b)] the gauge-component identity of the critical-point algebra, taken by statement. For
$X$ as below it asserts that for $\varsigma=RD^{*}X$ the identity
\eqref{5.8:eq-gauge-component-1} below holds, holomorphically in $B$. Its sources $\rho_{0}$,
$\rho_{1}$ are local: $\rho_{1}$ is built from products of $X$ with
$(\mathfrak{F}_{\mathbf{U}},\varsigma)$, and $\rho_{0}$ from products of $\nabla X$ with
$(\mathfrak{F}_{\mathbf{U}},\varsigma)$, together with the data terms;
\item[(c)] the density bound for the data terms of the critical-point algebra, taken by statement.
It asserts that at every $y$ the data terms have density at most
$C\ell^{-4}|X|_{y}(|\mathsf{m}|+|\mu_{B}|)$.
\end{enumerate}
Let $(\mathbf{U},\cdot)\in\mathsf{C}_{\beta}(\mathfrak{r}_{0})$. Let $\mathsf{p}$, $B$ and
$\mathsf{m}$, the critical point $X=\mathfrak{X}(B;\mathbf{U},\mathsf{m})$ and its multiplier
$\mu_{B}$ be as in the critical-point size lemma, Lemma~\ref{5.8:lem-critical-size}. Then, by (b),
\begin{equation}\label{5.8:eq-gauge-component-1}
\varsigma:=RD^{*}X=-RG'\bigl(\rho_{0}+D^{*}\rho_{1}\bigr),
\end{equation}
where, locally at every $y$,
\begin{equation}\label{5.8:eq-gauge-component-2}
\ell^{4}|\rho_{0}|,\quad \ell^{3}|\rho_{1}|,\quad \ell^{3+\beta}[\rho_{1}]_{\beta}
\;\le\; C\,\mathsf{p}(y)\bigl(\mathfrak{r}_{0}+\sup\mathsf{p}\bigr).
\end{equation}
Moreover
\begin{align}
\ell^{2}|\varsigma|,\quad \ell^{3}|D\varsigma|,\quad \ell^{3+\theta}[D\varsigma]_{\theta}
&\;\le\; C_{\varsigma}^{\sharp}\,\mathsf{p}(y)\bigl(\mathfrak{r}_{0}+\sup\mathsf{p}\bigr),
\label{5.8:eq-gauge-component-3}\\
\ell^{3}|DD^{*}X|,\quad \ell^{3+\theta}[DD^{*}X]_{\theta}
&\;\le\; C_{\varsigma}^{\sharp}\,\mathsf{p}(y).
\label{5.8:eq-gauge-component-4}
\end{align}
In these statements:
\begin{itemize}
\item $G'$ and $P$ are the operators of \cite{B-CMP99b} entering
Lemma~\ref{5.8:lem-weighted-calculus}, and $R:=1-P$;
\item $D$, $D^{*}$ and $\nabla$, $\nabla^{*}$ are the derivative operators of \cite{B-CMP99b}
and their adjoints, as they appear in its estimates \cite[(3.42)--(3.45)]{B-CMP99b};
\item $[\,\cdot\,]_{\beta}$ denotes the H\"older seminorm of index $\beta$, and
$\beta=\theta+\varepsilon$.
\end{itemize}
\end{lemma}

\begin{proof}
\emph{The identity.} The identity \eqref{5.8:eq-gauge-component-1} is statement (b). It comes from
the gauge flow together with the Euler--Lagrange equation of the critical point evaluated at the
variation $u=\dot{X}$, using $D^{*}J_{\mathbf{U}}=0$, where $J_{\mathbf{U}}$ is the current of the
background $\mathbf{U}$, and it is holomorphic in $B$. The sources of
\eqref{5.8:eq-gauge-component-1} are local:
\begin{itemize}
\item $\rho_{1}$ is of the form $X\cdot(\mathfrak{F}_{\mathbf{U}},\varsigma)$;
\item $\rho_{0}$ is of the form $\nabla X\cdot(\mathfrak{F}_{\mathbf{U}},\varsigma)$;
\item the data terms have, by statement (c), density at most
$C\ell^{-4}|X|_{y}(|\mathsf{m}|+|\mu_{B}|)$.
\end{itemize}

\emph{Sizes of the factors at $y$.} The hypotheses of Lemma~\ref{5.8:lem-critical-size} hold by
assumption. That lemma therefore gives
\begin{equation*}
N^{\beta}_{y}(X),\ |\mu_{B}|\le C_{s}^{\sharp}\,\mathsf{p}(y).
\end{equation*}
Since $(\mathbf{U},\cdot)\in\mathsf{C}_{\beta}(\mathfrak{r}_{0})$, the curvature
$\mathfrak{F}_{\mathbf{U}}$ is of class $C^{\beta}$ at size $\mathfrak{r}_{0}$; this is the
definition of the regularity class. For the gauge component we write, since $R=1-P$,
\begin{equation*}
\varsigma=D^{*}X-PD^{*}X .
\end{equation*}
The first term obeys $\ell^{2}|D^{*}X|,\ \ell^{2+\beta}[D^{*}X]_{\beta}\le C_{s}^{\sharp}\mathsf{p}(y)$
by Lemma~\ref{5.8:lem-critical-size}. For the second term we use
Lemma~\ref{5.8:lem-weighted-calculus}, applied to $P$ and $DP$ at the weight $\mathsf{p}$. This is
admissible because $\mathsf{p}\in\mathfrak{M}_{\sigma}$ is a hypothesis of
Lemma~\ref{5.8:lem-critical-size}. Together these give
\begin{equation}\label{5.8:eq-gauge-component-5}
\ell^{2}|\varsigma|,\quad \ell^{2+\beta}[\varsigma]_{\beta}\;\le\; C\,\mathsf{p}(y).
\end{equation}
We insert these sizes into the local forms of the sources. The factors $X$ and $\nabla X$ are
bounded through $N^{\beta}_{y}(X)$. The factor $\mathfrak{F}_{\mathbf{U}}$ contributes
$\mathfrak{r}_{0}$. The factor $\varsigma$ contributes, by \eqref{5.8:eq-gauge-component-5}, at most
$C\mathsf{p}\le C\sup\mathsf{p}$. The data terms are bounded by statement (c). This yields
\eqref{5.8:eq-gauge-component-2}.

The bound \eqref{5.8:eq-gauge-component-5} serves only as the a priori input to
\eqref{5.8:eq-gauge-component-2}; it does not carry the factor
$\mathfrak{r}_{0}+\sup\mathsf{p}$. The bound on $\ell^{2}|\varsigma|$ in
\eqref{5.8:eq-gauge-component-3} is obtained from the identity \eqref{5.8:eq-gauge-component-1}
together with \eqref{5.8:eq-gauge-component-2}. The bounds of \cite{B-CMP99b} for $G'$ on the local
source $\rho_{0}$ and for $G'D^{*}$ on $\rho_{1}$ are carried to the weight $\mathsf{p}$ by
Lemma~\ref{5.8:lem-weighted-calculus}: part~(i) for the single operators, and part~(ii) for the
products $RG'=(1-P)G'$ and $RG'D^{*}$. Applied to the sources bounded in
\eqref{5.8:eq-gauge-component-2}, they give
\begin{equation*}
\ell^{2}|\varsigma|\;\le\;C_{\varsigma}^{\sharp}\,\mathsf{p}(y)\bigl(\mathfrak{r}_{0}+\sup\mathsf{p}\bigr).
\end{equation*}

\emph{The derivative of $\varsigma$.} Applying $D$ to the identity and using $R=1-P$ gives
\begin{equation}\label{5.8:eq-gauge-component-6}
D\varsigma=-D(1-P)G'\bigl(\rho_{0}+D^{*}\rho_{1}\bigr).
\end{equation}
We estimate the three contributions to the right-hand side separately.
\begin{enumerate}
\item[(1)] The $\rho_{1}$ part is the operator $\nabla G'\nabla^{*}$ acting on the source
$\rho_{1}$, which is of class $C^{\beta}$ by \eqref{5.8:eq-gauge-component-2}. It is bounded in
sup norm by \cite[(3.44)]{B-CMP99b}. It is bounded in $C^{\theta}$ by \cite[(3.45)]{B-CMP99b},
which applies because $\beta=\theta+\varepsilon$.
\item[(2)] The $\rho_{0}$ part is $\nabla G'\rho_{0}$, which is bounded by
\cite[(3.42), (3.43)]{B-CMP99b}.
\item[(3)] The part containing $P$ is controlled through $DP$. The sup bound is
\cite[(3.49)]{B-CMP99b}, and the H\"older bound is its H\"older companion.
\end{enumerate}
In each of these applications the unweighted estimates of \cite{B-CMP99b} are carried to the weight
$\mathsf{p}$ by Lemma~\ref{5.8:lem-weighted-calculus}: part~(i) for single operators of the list,
and part~(ii) for their finite products, such as $D(1-P)G'$. Inserting
\eqref{5.8:eq-gauge-component-2} on the right of \eqref{5.8:eq-gauge-component-6} gives the bounds
on $\ell^{3}|D\varsigma|$ and $\ell^{3+\theta}[D\varsigma]_{\theta}$ in
\eqref{5.8:eq-gauge-component-3}.

\emph{The derivative of $D^{*}X$.} By the definition of $\varsigma$,
\begin{equation*}
D^{*}X=\varsigma+PD^{*}X,
\end{equation*}
and hence
\begin{equation*}
DD^{*}X=D\varsigma+DP\,D^{*}X .
\end{equation*}
The first term is bounded by \eqref{5.8:eq-gauge-component-3}. Since
$\sup\mathsf{p}\le a_{R}^{\sharp}$ by the hypotheses of Lemma~\ref{5.8:lem-critical-size}, the
factor $\mathfrak{r}_{0}+\sup\mathsf{p}\le\mathfrak{r}_{0}+a_{R}^{\sharp}$ is bounded by a constant,
which is absorbed into $C_{\varsigma}^{\sharp}$. For the second term,
\cite[(3.49)]{B-CMP99b} and its H\"older companion, carried to the weight $\mathsf{p}$ by
Lemma~\ref{5.8:lem-weighted-calculus}~(i) and~(ii), are applied to $D^{*}X$. The input is the bound
$\ell^{2}|D^{*}X|,\ \ell^{2+\beta}[D^{*}X]_{\beta}\le C_{s}^{\sharp}\mathsf{p}(y)$ of
Lemma~\ref{5.8:lem-critical-size}. Together these give \eqref{5.8:eq-gauge-component-4}.
\end{proof}

\begin{lemma}[Lipschitz dependence of the gauge projection]\label{5.8:lem-projection-lipschitz}
Let $V$ be a background, let $\mathfrak{a}$ be a one-form, and let $V_1 = e^{i\eta\mathfrak{a}}V$, where $V$
and $V_1$ satisfy \cite[(3.35)]{B-CMP99b}. For $V' \in \{V, V_1\}$ write $R_{V'}$ for the gauge projection
at $V'$. Write $D_{V'}$ for the covariant derivative and $D_{V'}^*$ for its adjoint. Write
$\Delta_{V'} = D_{V'}^* D_{V'}$ for the covariant Laplacian, $Q'(V')$ for the averaging operator,
$Q'_j(V')$ for the averaging operators to which the averaging-difference bound below refers, $G'_{V'}$ for the
Green operator and $C'_{V'}$ for the operator on block functions, all as in \cite{B-CMP99b}. Let $P_{V'}$ be
the operator
\begin{equation}\label{5.8:eq-projection-lipschitz-1}
P_{V'} g = G'_{V'} Q'(V')^* C'_{V'} Q'(V') G'_{V'} g .
\end{equation}
Assume the following.
\begin{enumerate}
\item[(1)] The hypotheses of the weighted calculus (Lemma~\ref{5.8:lem-weighted-calculus}) hold.
\item[(2)] The projections at the two backgrounds satisfy the projection-difference identity
\begin{equation}\label{5.8:eq-projection-lipschitz-2}
R_V - R_{V_1} = R_V P_{V_1} - P_V R_{V_1} .
\end{equation}
\item[(3)] The background-change formula holds: for all $\varphi$ and $\nu$ with
$\Delta_{V_1}\varphi = Q'(V_1)^*\nu$, one has
\begin{equation}\label{5.8:eq-projection-lipschitz-3}
\begin{split}
R_V\varphi = R_V G'_V\bigl[\,& D_V^*\bigl((D_V - D_{V_1})\varphi\bigr)
+ (D_V^* - D_{V_1}^*)D_{V_1}\varphi \\
&+ \bigl(Q'(V_1) - Q'(V)\bigr)^*\nu\,\bigr].
\end{split}
\end{equation}
\item[(4)] The averaging-difference bound holds: for every $j$ and every $y$, the difference
$Q'_j(V_1) - Q'_j(V)$, read at $y$, has size at most $16d\,|\mathfrak{a}|_y$, where $d = 4$ is the
dimension.
\end{enumerate}
Let $\rho, \tau \ge 0$ with $\rho + \tau \le 6\delta_{1}$. Let $\mathsf{g} \in \mathfrak{M}_{\rho}$ and
$\mathsf{M} \in \mathfrak{M}_{\tau}$, and assume that $|\mathfrak{a}|_y \le \mathsf{M}(y) \le 1$ for all $y$.
Then, with a constant $C$, the following hold.
\begin{enumerate}
\item[(a)] If $g$ is a scalar function with $\ell^{2}|g|(y) \le \mathsf{g}(y)$ for all $y$, then for
every $x$
\begin{equation}\label{5.8:eq-projection-lipschitz-4}
\ell^{2}\bigl|(R_{V_1} - R_V)g\bigr|(x) \le C\,\mathsf{M}(x)\,\mathsf{g}(x).
\end{equation}
\item[(b)] If $A$ is a one-form with $\ell|A|(y) \le \mathsf{g}(y)$ and
$\ell^{2}|\nabla A|(y) \le \mathsf{g}(y)$ for all $y$, then for
every $x$
\begin{equation}\label{5.8:eq-projection-lipschitz-5}
\ell^{2}\bigl|(R_{V_1}D_{V_1}^* - R_V D_V^*)A\bigr|(x) \le C\,\mathsf{M}(x)\,\mathsf{g}(x).
\end{equation}
\end{enumerate}
\end{lemma}

\begin{proof}
We use the following abbreviations for the operators at $V_1$:
\begin{equation*}
R_1 := R_{V_1},\quad P_1 := P_{V_1},\quad G'_1 := G'_{V_1},
\end{equation*}
\begin{equation*}
C'_1 := C'_{V_1},\quad D_1 := D_{V_1},\quad \Delta_1 := \Delta_{V_1}.
\end{equation*}
The letter $C$ denotes a constant that may change from line to line.

We first record which weights are admissible in Lemma~\ref{5.8:lem-weighted-calculus}. Since
$\mathsf{g} \in \mathfrak{M}_{\rho}$ and
$\mathsf{M} \in \mathfrak{M}_{\tau}$, the definition of tempered weights
(Definition~\ref{5.8:def-tempered-weights}) gives, for all $y, y'$,
\begin{equation*}
\mathsf{M}(y)\mathsf{g}(y) \le e^{\tau d(y,y')}\mathsf{M}(y')\,e^{\rho d(y,y')}\mathsf{g}(y')
= e^{(\rho+\tau)d(y,y')}\mathsf{M}(y')\mathsf{g}(y').
\end{equation*}
Hence $\mathsf{M}\mathsf{g} \in \mathfrak{M}_{\rho+\tau}$. Since $\rho \le \rho+\tau \le 6\delta_{1}$,
Lemma~\ref{5.8:lem-weighted-calculus} applies both at the weight $\mathsf{g}$ and at the weight
$\mathsf{M}\mathsf{g}$.

\emph{Proof of (a).} By \eqref{5.8:eq-projection-lipschitz-2},
\begin{equation*}
(R_{V_1} - R_V)g = P_V R_1 g - R_V P_1 g .
\end{equation*}
It therefore suffices to bound $\ell^{2}|R_V P_1 g|$ and $\ell^{2}|P_V R_1 g|$ by
$C\mathsf{M}\mathsf{g}$, pointwise.

\emph{First half: the term $R_V P_1 g$.} By \eqref{5.8:eq-projection-lipschitz-1},
\begin{equation*}
\varphi_1 := P_1 g = G'_1 Q'(V_1)^{*} C'_1 Q'(V_1) G'_1 g ,
\end{equation*}
and $\varphi_1$ satisfies
\begin{equation*}
\Delta_1\varphi_1 = Q'(V_1)^*\nu_1,\qquad \nu_1 := C'_1 Q'(V_1) G'_1 g - a\,Q'(V_1)\varphi_1 ,
\end{equation*}
where $a$ denotes the constant of the averaging term in \cite{B-CMP99b}. The hypothesis
$\ell^{2}|g| \le \mathsf{g}$ says
that the $\mathsf{g}$-weighted norm of $g$, read as a scalar of weight two, is at most one.
The functions $\varphi_1$, $D_1\varphi_1$ and $\nu_1$ are obtained from $g$ by finite products of the
operators $G'_1$, $Q'(V_1)$, $Q'(V_1)^{*}$, $C'_1$, $D_1$. Lemma~\ref{5.8:lem-weighted-calculus}, applied at
the weight $\mathsf{g}$, therefore gives for every $y$
\begin{equation*}
\ell^{2}|\varphi_1|,\quad \ell^{3}|D_1\varphi_1|,\quad \ell^{4}|\nu_1| \;\le\; C\,\mathsf{g}(y).
\end{equation*}
Since $\Delta_1\varphi_1 = Q'(V_1)^*\nu_1$, the background-change formula
\eqref{5.8:eq-projection-lipschitz-3} applies with $\varphi = \varphi_1$ and $\nu = \nu_1$. It gives
\begin{equation*}
R_V\varphi_1 = R_V G'_V D_V^* s + R_V G'_V (s_0 + s_1),
\end{equation*}
with the three sources
\begin{equation*}
s := (D_V - D_1)\varphi_1,\qquad s_0 := (D_V^* - D_1^*)D_1\varphi_1,\qquad
s_1 := \bigl(Q'(V_1) - Q'(V)\bigr)^*\nu_1 .
\end{equation*}
Since $V_1 = e^{i\eta\mathfrak{a}}V$, the difference $D_V - D_1$ is an operator of order zero and size
$|\mathfrak{a}|$, and so, by taking adjoints, is $D_V^* - D_1^*$. By the averaging-difference bound,
for every $j$ the difference $Q'_j(V_1) - Q'_j(V)$ has size at most $16d\,|\mathfrak{a}|_y$ at $y$.
Combining these facts
with the bounds just obtained for
$\varphi_1$, $D_1\varphi_1$, $\nu_1$, and with $|\mathfrak{a}|_y \le \mathsf{M}(y)$, the three sources
satisfy locally, that is at each $y$,
\begin{equation*}
\ell^{3}|s|,\quad \ell^{4}|s_0|,\quad \ell^{4}|s_1| \;\le\; C\,\mathsf{M}(y)\,\mathsf{g}(y).
\end{equation*}
Thus the $\mathsf{M}\mathsf{g}$-weighted norms of the three sources are bounded by $C$.
Lemma~\ref{5.8:lem-weighted-calculus}~(ii), applied at the weight
$\mathsf{M}\mathsf{g} \in \mathfrak{M}_{\rho+\tau}$ to the operators $R_V G'_V D_V^*$ and $R_V G'_V$, yields
$\ell^{2}|R_V P_1 g|(x) \le C\mathsf{M}(x)\mathsf{g}(x)$.

\emph{Second half: the term $P_V R_1 g$.} Put $\chi_1 := G'_1 R_1 g$. Then
$\chi_1 \in \ker Q'(V_1)$ and $R_1 g = \Delta_1\chi_1$. By
Lemma~\ref{5.8:lem-weighted-calculus} at the weight $\mathsf{g}$, as in the first half,
\begin{equation*}
|\chi_1|,\quad \ell|D_1\chi_1| \;\le\; C\,\mathsf{g}(y).
\end{equation*}
We write
\begin{equation*}
P_V R_1 g = P_V\Delta_1\chi_1 = P_V(\Delta_1 - \Delta_V)\chi_1 + P_V\Delta_V\chi_1 .
\end{equation*}
Since $\Delta_{V'} = D_{V'}^* D_{V'}$ for $V' \in \{V, V_1\}$, adding and subtracting $D_V^* D_1\chi_1$ gives
\begin{equation*}
(\Delta_1 - \Delta_V)\chi_1 = D_V^* t + t_0,\qquad t := (D_1 - D_V)\chi_1,\qquad
t_0 := (D_1^* - D_V^*)D_1\chi_1 .
\end{equation*}
For the second term we have
\begin{equation*}
P_V\Delta_V\chi_1
= G'_V Q'(V)^* C'_V\bigl(1 - Q'(V)G'_V Q'(V)^* a\bigr)Q'(V)\chi_1 .
\end{equation*}
This term reads $\chi_1$ only through $Q'(V)\chi_1$. Since $Q'(V_1)\chi_1 = 0$,
\begin{equation*}
Q'(V)\chi_1 = t_1,\qquad t_1 := \bigl(Q'(V) - Q'(V_1)\bigr)\chi_1 .
\end{equation*}
The three sources $t$, $t_0$ and $t_1$ are bounded locally, in their scaled sizes, by
$C\mathsf{M}\mathsf{g}$. For $t$ and $t_0$ this follows from the order-zero difference $D_V - D_1$ of size
$|\mathfrak{a}|$. For $t_1$ it follows from the averaging-difference bound. In all cases one also uses
the bounds on $\chi_1$, $D_1\chi_1$ and $|\mathfrak{a}|_y \le \mathsf{M}(y)$. The operators acting on the
sources $t$ and $t_0$ are $P_V D_V^*$ and $P_V$, and $t_1$ enters through
$G'_V Q'(V)^* C'_V\bigl(1 - Q'(V)G'_V Q'(V)^* a\bigr)$, a finite product of members of $\mathfrak{L}$.
Lemma~\ref{5.8:lem-weighted-calculus}~(ii) bounds
them at the weight $\mathsf{M}\mathsf{g}$, so that
$\ell^{2}|P_V R_1 g|(x) \le C\mathsf{M}(x)\mathsf{g}(x)$.

The two halves and the triangle inequality give \eqref{5.8:eq-projection-lipschitz-4}.

\emph{Proof of (b).} Adding and subtracting $R_V D_1^* A$ gives
\begin{equation*}
(R_{V_1}D_{V_1}^* - R_V D_V^*)A = (R_{V_1} - R_V)D_1^* A + R_V(D_1^* - D_V^*)A .
\end{equation*}
Consider the first term. Since $D_1^*$ is a first-order operator, the bounds $\ell|A| \le \mathsf{g}$ and
$\ell^{2}|\nabla A| \le \mathsf{g}$ at the neighbouring points it reads give
$\ell^{2}|D_1^* A| \le C\mathsf{g}$, by the tempering of $\mathsf{g}$. Part (a), applied to
$D_1^* A / C$, then bounds the first term by $C\mathsf{M}\mathsf{g}$.

Consider the second term. The operator $D_1^* - D_V^*$ is of order zero and size $|\mathfrak{a}|$. Hence,
locally,
\begin{equation*}
\ell^{2}|(D_1^* - D_V^*)A| \le C\mathsf{M}\,\ell|A| \le C\mathsf{M}\mathsf{g}.
\end{equation*}
Lemma~\ref{5.8:lem-weighted-calculus}~(ii), applied to $R_V$ at the weight $\mathsf{M}\mathsf{g}$, bounds
the second term by $C\mathsf{M}\mathsf{g}$. This proves \eqref{5.8:eq-projection-lipschitz-5}.
\end{proof}

\begin{remark}
In the unweighted form of the bound \eqref{5.8:eq-projection-lipschitz-4}, the right-hand side carries
the $\delta_{1}$-smear $(\ell^{2}|g|)^{\sim}$ in place of $\mathsf{g}(x)$; this
smear is produced by the operators $P_1$, $D_1 P_1$ and $C'_1 Q'(V_1) G'_1$, whose kernels decay at rates
at least $16\delta_{1}$.
\end{remark}

\begin{lemma}[Nesting of fine and coarse gauge projections]\label{5.8:lem-projection-nesting}
Let $\mathbf{U}'$ satisfy the field clauses of the regularity class
$\mathsf{C}'_{\beta}(\mathfrak{r}_{0})$, put $W:=M(\mathbf{U}')$, and let $f'$ be a fine scalar.
Here $R'$ and $P'$ denote the fine projections on fine scalars in the construction $\mathsf{B}$ at
the finer cutoff, of lattice spacing $\eta'=\eta/L$, at the background $\mathbf{U}'$, and $R_{W}$,
$P_{W}$ the coarse projections at the
background $W$; $Q'_{1}$ is the averaging operator from fine to coarse scalars, and
$N(Q'(\mathbf{U}'))$ is the null space of the fine averaging operator $Q'=Q'(\mathbf{U}')$;
$G'$, $C'$ and the covariant derivative $D'$ are operators of $\mathsf{B}$ at $\mathbf{U}'$, and
$G'_{W}$, $D_{W}^{*}$ the corresponding operators at $W$; $[\,\cdot\,]_{\beta}$ is the H\"older
seminorm of index $\beta$ and $\operatorname{osc}_{2d\eta}$ the oscillation over distances
$2d\eta$; $f^{s}[\cdot]$ and $f^{s}_{0}[\cdot]$ are the source terms of the background-field
representation, $\alpha'$ its field-size parameter, and $C_{L}$ a constant depending on $L$.
Put $\varphi':=P'f'$ and $\chi':=G'R'f'$. Assume the following.
\begin{enumerate}
\item The weighted calculus, Lemma~\ref{5.8:lem-weighted-calculus}, in the finer-cutoff
construction $\mathsf{B}$ (at the background $\mathbf{U}'$) and at the coarse background $W$.
\item The nesting identities for the fine and coarse projections:
\begin{equation*}
Q'_{1}R'f'-R_{W}Q'_{1}f' = P_{W}Q'_{1}R'f'-R_{W}Q'_{1}P'f' .
\end{equation*}
\item The representation of the fine projection through the Green operator of $\mathsf{B}$:
$P'f'=G'Q'^{*}\mu'$ with $\mu':=C'Q'G'f'$.
\item The background-field representation of $R_{W}Q'_{1}$: one has
$R_{W}Q'_{1}\varphi'=R_{W}G'_{W}[D_{W}^{*}f^{s}[\varphi']+f^{s}_{0}[\varphi']]$, where, for
$\psi\in\{\varphi',\chi'\}$,
\begin{align*}
|f^{s}[\psi]| &\le \operatorname{osc}_{2d\eta}(D'\psi)+C_{L}\eta^{-1}\alpha'|\psi| ,\\
|f^{s}_{0}[\varphi']| &\le C_{L}\eta^{-1}\alpha'|D'\varphi'| ,
\qquad \alpha'\le \mathfrak{r}_{0}\eta'^{2}\ell^{-2}.
\end{align*}
\item The dual form of the nesting identities:
$P_{W}Q'_{1}R'f'=-P_{W}D_{W}^{*}f^{s}[\chi']$.
\item The null-space property of the fine Green operator: $\chi'\in N(Q'(\mathbf{U}'))$.
\item The operators $R_{W}G'_{W}D_{W}^{*}$ and $R_{W}G'_{W}$ at the background $W$ fall under
clause (ii) of Lemma~\ref{5.8:lem-weighted-calculus}.
\end{enumerate}
Let $\mathsf{f}\in\mathfrak{M}_{\sigma}$ with $\sigma\le 5\delta_{1}$, and suppose
$\ell^{2}|f'|\le\mathsf{f}$. Then
\begin{equation}\label{5.8:eq-projection-nesting-1}
\ell^{2}\,\bigl|Q'_{1}R'f'-R_{W}Q'_{1}f'\bigr|(x)\le C\,\mathsf{f}(x)\,L^{-\beta\lambda(x)} .
\end{equation}
\end{lemma}

\begin{proof}
Throughout, $C$ denotes a constant whose value may change from one inequality to the next.
By the nesting identities,
\begin{equation*}
Q'_{1}R'f'-R_{W}Q'_{1}f' = P_{W}Q'_{1}R'f'-R_{W}Q'_{1}P'f' ,
\end{equation*}
and we bound the two terms on the right separately, the second one first. Since
$\ell^{2}|\cdot|$ is subadditive, \eqref{5.8:eq-projection-nesting-1} follows once each term
satisfies the same bound.

\emph{The second term.} Recall $\varphi'=P'f'$. By the representation of the fine projection,
$\varphi'=G'Q'^{*}\mu'$ with
$\mu':=C'Q'G'f'$, all operators being those of the finer-cutoff construction $\mathsf{B}$. We apply
Lemma~\ref{5.8:lem-weighted-calculus} in $\mathsf{B}$ at the weight $\mathsf{f}$; this is
admissible because $\mathsf{f}\in\mathfrak{M}_{\sigma}$ with
$\sigma\le 5\delta_{1}\le 6\delta_{1}$, and the bounds on the operators entering the product
$G'Q'^{*}C'Q'G'$ are those of \cite[(3.42), (3.43), (3.48)]{B-CMP99b}. Starting from
$\ell^{2}|f'|\le\mathsf{f}$, it gives, at every point $y$,
\begin{equation}\label{5.8:eq-projection-nesting-2}
\ell^{2}|\varphi'|,\quad \ell^{3}|D'\varphi'|,\quad
\ell^{3+\beta}[D'\varphi']_{\beta}\ \le\ C\,\mathsf{f}(y) .
\end{equation}
By the background-field representation, with $f^{s}:=f^{s}[\varphi']$ and
$f^{s}_{0}:=f^{s}_{0}[\varphi']$,
\begin{equation*}
R_{W}Q'_{1}\varphi'=R_{W}G'_{W}\bigl[D_{W}^{*}f^{s}+f^{s}_{0}\bigr],
\end{equation*}
with $f^{s}$ bounded by the oscillation of $D'\varphi'$ over distances $2d\eta$ plus
$C_{L}\eta^{-1}\alpha'|\varphi'|$, and $f^{s}_{0}$ bounded by
$C_{L}\eta^{-1}\alpha'|D'\varphi'|$, where $\alpha'\le\mathfrak{r}_{0}\eta'^{2}\ell^{-2}$.
Inserting \eqref{5.8:eq-projection-nesting-2} into these bounds, the oscillation through the
H\"older bound on $D'\varphi'$ and the remaining terms through the bounds on $|\varphi'|$,
$|D'\varphi'|$ and on $\alpha'$, gives locally, at every point $y$,
\begin{equation}\label{5.8:eq-projection-nesting-3}
\ell^{3}|f^{s}|,\quad \ell^{4}|f^{s}_{0}|\ \le\ C\,\mathsf{f}(y)\,L^{-\beta\lambda(y)} .
\end{equation}
By the weight properties \eqref{5.8:eq-tempered-weights-a} of
Definition~\ref{5.8:def-tempered-weights}, the weight $\mathsf{f}L^{-\beta\lambda}$ lies in
$\mathfrak{M}^{L}_{\sigma+\delta_{1}/16}$; since $\sigma\le 5\delta_{1}$, its rate satisfies
$\sigma+\delta_{1}/16\le 6\delta_{1}$. By assumption, the operators $R_{W}G'_{W}D_{W}^{*}$
and $R_{W}G'_{W}$ at the background $W$ fall under clause (ii) of
Lemma~\ref{5.8:lem-weighted-calculus}; applied with the weight $\mathsf{f}L^{-\beta\lambda}$ to
the sources $f^{s}$ and $f^{s}_{0}$, whose weighted sizes are controlled by
\eqref{5.8:eq-projection-nesting-3}, it gives
\begin{equation*}
\ell^{2}\,|R_{W}Q'_{1}P'f'|(x)=\ell^{2}\,|R_{W}Q'_{1}\varphi'|(x)
\le C\,\mathsf{f}(x)\,L^{-\beta\lambda(x)} .
\end{equation*}

\emph{The first term.} Recall $\chi'=G'R'f'$; by the null-space property of the fine Green
operator, $\chi'\in N(Q'(\mathbf{U}'))$. By
Lemma~\ref{5.8:lem-weighted-calculus} in $\mathsf{B}$ at the weight $\mathsf{f}$, applied as in
the second term, at every point $y$,
\begin{equation*}
\ell^{1+\beta}[D'\chi']_{\beta},\quad \ell|D'\chi'|,\quad |\chi'|\ \le\ C\,\mathsf{f}(y) .
\end{equation*}
By the dual form of the nesting identities,
$P_{W}Q'_{1}R'f'=-P_{W}D_{W}^{*}f^{s}[\chi']$. The bounds just obtained on $\chi'$, inserted
in the bound on $f^{s}[\chi']$ of the background-field representation, give,
locally,
\begin{equation*}
\ell\,|f^{s}[\chi']|\le C\,\mathsf{f}\,L^{-\beta\lambda} .
\end{equation*}
Applying Lemma~\ref{5.8:lem-weighted-calculus} to $P_{W}D_{W}^{*}$ at the weight
$\mathsf{f}L^{-\beta\lambda}$, which, as shown above, lies in
$\mathfrak{M}^{L}_{\sigma+\delta_{1}/16}$ with $\sigma+\delta_{1}/16\le 6\delta_{1}$, we obtain
\begin{equation*}
\ell^{2}\,|P_{W}Q'_{1}R'f'|(x)=\ell^{2}\,|P_{W}D_{W}^{*}f^{s}[\chi']|(x)
\le C\,\mathsf{f}(x)\,L^{-\beta\lambda(x)} .
\end{equation*}
Adding the bounds for the two terms gives \eqref{5.8:eq-projection-nesting-1}.
\end{proof}

\begin{theorem}[Two-cutoff comparison of responses at tempered profiles]
\label{5.8:thm-response-comparison}
Assume the following statements.
\begin{itemize}
\item The weighted calculus, Lemma~\ref{5.8:lem-weighted-calculus}.
\item From the comparison of responses at smeared majorants:
  \begin{itemize}
  \item its constraint step;
  \item its residual identity, which writes $d\Sigma_{\mathbf{U},\mathsf{m}}(Y)[u]$ as the sum of the six
    members (I)--(VI) listed in the proof below;
  \item its bound for block means of a $C^{1}$ field;
  \item the local statements used in its proof, namely:
    \begin{itemize}
    \item the background-change bound for $d\mathcal{W}$;
    \item the nesting bound for $d\mathcal{W}$;
    \item the extension map $\hat{E}$ with its transfer identity;
    \item the divergence comparison;
    \item the constraint comparison;
    \item the condition \cite[(3.35)]{B-CMP99b} for $W$.
    \end{itemize}
  \end{itemize}
\item Parts (a) and (b) of the background lemma for the paired backgrounds.
\item The blocked-increment lemma and the Lipschitz bound for the blocked-increment map
  $Y[\,\cdot\,;\mathbf{U}']$.
\item The statements of the second item hold, as stated, when the data
  $(b,\delta b,b_{c},\mathsf{m}',\mathsf{m})$ are complex and the backgrounds $\mathbf{U}$, $\mathbf{U}'$,
  $W$ are real.
\end{itemize}

Let $(\mathfrak{S}',\mathfrak{S})$ be exactly paired admissible sets
(Notation~\ref{5.8:not-paired-sets}). Let $\mathbf{U}'$ be real and satisfy the field clauses of
$\mathsf{C}'_{\beta}(\mathfrak{r}_{0})$. Let $\mathbf{U}$ be real and satisfy
\cite[(3.35)]{B-CMP99b}. Put $W:=M(\mathbf{U}')=e^{i\eta\mathfrak{a}}\mathbf{U}$ and assume
$\sup_{y}|\mathfrak{a}|_{y}\le\varepsilon_{D}$. Let $\mathsf{m}',\mathsf{m}$ be data vectors with
$|\mathsf{m}'|,|\mathsf{m}|\le\mathfrak{r}_{0}$.

Let the profiles satisfy
\begin{equation*}
\mathsf{p}\in\mathfrak{M}_{\rho},\qquad \mathsf{M}\in\mathfrak{M}^{A}_{\tau},\qquad
\mathsf{p}_{\delta}\in\mathfrak{M}_{\rho+\tau},\qquad \rho+\tau\le 5\delta_{1},
\end{equation*}
together with $10\sup\mathsf{p}\le a_{R}^{\sharp}$, $\mathsf{p}_{\delta}\le 4\mathsf{p}$ and $\mathsf{M}\le 1$,
and
\begin{equation}\label{5.8:eq-response-comparison-1}
|\mathfrak{a}|_{y}\le\mathsf{M}(y),\qquad |\mathsf{m}'(c)-\mathsf{m}(c)|\le\mathsf{M}(c),\qquad
L^{-\theta\lambda(y)}\le\mathsf{M}(y).
\end{equation}

Let $b$ and $\delta b$ be complex on $\mathfrak{B}$, and let $b_{c}$ be complex on the collar
$\Lambda'_{0}$ of $\mathsf{B}$. Assume $|b|\le\mathsf{p}$, $|\delta b|\le\mathsf{p}_{\delta}$ and
$|b_{c}|\le\mathsf{p}'$, where $\mathsf{p}'$ is the extension of $\mathsf{p}$ to $\mathfrak{S}'$ of
\eqref{5.8:eq-tempered-weights-a}. Put
\begin{equation*}
X:=\mathfrak{X}_{\mathfrak{S}}(b;\mathbf{U},\mathsf{m}),\qquad
X':=\mathfrak{X}_{\mathfrak{S}'}\bigl((b+\delta b)\oplus b_{c};\mathbf{U}',\mathsf{m}'\bigr).
\end{equation*}
Assume \eqref{5.8:eq-paired-sets-a}, \eqref{5.8:eq-paired-sets-b} and
\eqref{5.8:eq-paired-sets-c}, the fourth hypothesis of Lemma~\ref{5.8:lem-weighted-calculus} (the
condition on weights assumed there besides \eqref{5.8:eq-paired-sets-a},
\eqref{5.8:eq-paired-sets-b} and \eqref{5.8:eq-paired-sets-c}), and the smallness conditions
\eqref{5.8:eq-weighted-chart-a} and \eqref{5.8:eq-smallness-ceilings-a}. Then
\begin{equation}\label{5.8:eq-response-comparison-2}
\bigl|Y[X';\mathbf{U}']-X\bigr|_{x}\le C_{B}^{\sharp}\bigl\{A\,\mathsf{p}(x)\mathsf{M}(x)
+\mathsf{p}_{\delta}(x)\bigr\}\qquad\text{for every }x,
\end{equation}
with $C_{B}^{\sharp}=C_{B}^{\sharp}(d,L,\theta,M_{1})$, where $d$ is the dimension.
\end{theorem}

\begin{proof}
Throughout the proof $C$ denotes a constant depending only on $(d,L,\theta,M_{1})$. Its value may
change from one occurrence to the next.

A linear functional of the variation $u$ is said to lie in $\sharp(\kappa)$ if it has a local
density with profile $\kappa$. A bound is called local at $y$ if it holds at the site $y$ with the
profiles evaluated at $y$.

\emph{Reduction to $A=1$.} Replace $\mathsf{M}$ by $\min\{1,\mathsf{M}^{[\tau]}\}$. This function lies in
$\mathfrak{M}_{\tau}$ by \eqref{5.8:eq-tempered-weights-a}. It also lies between $\mathsf{M}$ and
$A\mathsf{M}$, for two reasons:
\begin{itemize}
\item taking $y'=y$ in $\mathsf{M}^{[\tau]}(y)=\sup_{y'}\mathsf{M}(y')e^{-\tau d(y,y')}$ gives
  $\mathsf{M}^{[\tau]}\ge\mathsf{M}$, and since $\mathsf{M}\le 1$ the minimum with $1$ is still
  $\ge\mathsf{M}$;
\item the defining inequality $\mathsf{M}(y')\le Ae^{\tau d(y,y')}\mathsf{M}(y)$ of
  $\mathfrak{M}^{A}_{\tau}$ (Definition~\ref{5.8:def-tempered-weights}) gives
  $\mathsf{M}^{[\tau]}\le A\mathsf{M}$.
\end{itemize}
Every hypothesis on $\mathsf{M}$ is one of three kinds: the bound $\le 1$, a lower bound
\eqref{5.8:eq-response-comparison-1}, or membership in a class of tempered weights. Each therefore
persists after the replacement. The conclusion \eqref{5.8:eq-response-comparison-2} for the replaced
profile with $A=1$ implies it for $\mathsf{M}$ with the factor $A$. We therefore take $A=1$.

\emph{The collar.} Put
\begin{equation*}
X'_{0}:=\mathfrak{X}_{\mathfrak{S}'}\bigl((b+\delta b)\oplus 0;\mathbf{U}',\mathsf{m}'\bigr),
\end{equation*}
and on $\mathfrak{B}'$ the weight
\begin{equation*}
\tilde{w}:=\bigl(\mathsf{p}'L^{-\theta\lambda}\bigr)^{[\rho+\delta_{1}/16]},
\end{equation*}
with $\lambda=-1$ on $\Lambda'_{0}$. We record three of its properties.
\begin{itemize}
\item It is tempered at the rate $\rho+\delta_{1}/16$, by \eqref{5.8:eq-tempered-weights-a} read on
  $\mathfrak{S}'$. This rate is $\le 6\delta_{1}$, because $\rho\le\rho+\tau\le 5\delta_{1}$.
\item On $\Lambda'_{0}$ the hull is at least the function itself, so that
  \begin{equation*}
  \tilde{w}\ge\mathsf{p}'L^{\theta}\ge\mathsf{p}'\ge|b_{c}|\qquad\text{on }\Lambda'_{0}.
  \end{equation*}
\item On $\mathfrak{B}$, by the same properties read on $\mathfrak{S}'$,
  $\tilde{w}\le L\mathsf{p}L^{-\theta\lambda}$.
\end{itemize}
The data $(b+\delta b)\oplus b_{c}$ and $(b+\delta b)\oplus 0$ differ by $b_{c}$, which is supported on
$\Lambda'_{0}$. Their difference therefore has $\|\cdot\|_{\tilde{w}}$-norm at most $1$. The weighted
stability \eqref{5.8:eq-critical-stability-b}, applied at the finer cutoff $\mathsf{B}$, gives
$|X'-X'_{0}|_{x}\le C_{S}^{\sharp}\tilde{w}(x)$. The Lipschitz bound for the blocked-increment map then
gives
\begin{equation*}
\bigl|Y[X';\mathbf{U}']-Y[X'_{0};\mathbf{U}']\bigr|_{x}\le C\mathsf{p}(x)L^{-\theta\lambda(x)}
\le C\mathsf{p}(x)\mathsf{M}(x),
\end{equation*}
where the last inequality is the third bound in \eqref{5.8:eq-response-comparison-1}. The right-hand
side is of the form of \eqref{5.8:eq-response-comparison-2}. Moreover $X'_{0}$ satisfies all hypotheses
of the theorem, with $b_{c}=0$. By the triangle inequality it therefore suffices to prove the theorem
with $b_{c}=0$, which we assume from now on.

\emph{Sizes.} The datum of $X'$ obeys $|b+\delta b|\le\mathsf{p}+\mathsf{p}_{\delta}\le 5\mathsf{p}$. The
weight $5\mathsf{p}$ lies in $\mathfrak{M}_{\rho}$, and $\sup 5\mathsf{p}\le a_{R}^{\sharp}/2$.

We apply the critical-point size at a tempered profile (Lemma~\ref{5.8:lem-critical-size}) at the
finer cutoff $\mathsf{B}$ at the weight $5\mathsf{p}$. It gives, locally at every $y$,
\begin{equation*}
N^{\beta}_{y}(X')\le N^{\Delta}_{y}(X')\le C\mathsf{p}(y),\qquad |\mu'|\le C\mathsf{p}(y),
\end{equation*}
where $\mu'$ is the Lagrange multiplier of $X'$.

We apply the gauge component of the critical point (Lemma~\ref{5.8:lem-gauge-component}) in
$\mathsf{B}$ at the same weight. It gives, for $\varsigma':=R'D'^{*}X'$ and locally at $y$,
\begin{equation*}
\ell^{2}|\varsigma'|,\ \ \ell^{3}|D'\varsigma'|,\ \ \ell^{3+\theta}[D'\varsigma']_{\theta}\le
C_{\varsigma}^{\sharp}\,5\mathsf{p}(y)\,(\mathfrak{r}_{0}+5\sup\mathsf{p})\le C\mathsf{p}(y).
\end{equation*}

Put $Y:=Y[X';\mathbf{U}']$. The blocked-increment lemma and the bound for block means of a $C^{1}$
field give, locally at $y$,
\begin{equation*}
\ell|Y|,\ \ \ell^{2}|\nabla Y|\le C\mathsf{p}(y)\le\varepsilon_{3}.
\end{equation*}
The last inequality holds because $\sup\mathsf{p}\le a_{R}^{\sharp}/10$ and by the ceilings
\eqref{5.8:eq-smallness-ceilings-a}.

Finally, $\lambda\ge 0$ on $\mathfrak{B}$ (Notation~\ref{5.8:not-paired-sets}) and
$\theta\le\beta\le 1$. Hence on $\mathfrak{B}$, using
\eqref{5.8:eq-response-comparison-1},
\begin{equation}\label{5.8:eq-response-comparison-3}
L^{-\lambda}\le L^{-\beta\lambda}\le L^{-\theta\lambda}\le\mathsf{M}.
\end{equation}

\emph{The constraint (local).} The constraint comparison gives
\begin{equation*}
\mathbf{Q}_{W}(\eta Y)=(b+\delta b)+\beta_{1},\qquad |\beta_{1}|\le CL^{-\lambda}\mathsf{p}.
\end{equation*}
The difference $\mathbf{Q}_{\mathbf{U}}(\eta Y)-\mathbf{Q}_{W}(\eta Y)$ is analytic in
$(\mathfrak{a},Y)$. It vanishes at $Y=0$, and at $\mathfrak{a}=0$ (where $W=\mathbf{U}$). It is therefore
bounded by $C|\mathfrak{a}|_{y}|Y|_{y}$.

Hence $\mathbf{Q}_{\mathbf{U}}(\eta Y)=b+\tilde{\beta}$ with
\begin{equation*}
\tilde{\beta}:=\delta b+\beta_{1}+\bigl(\mathbf{Q}_{\mathbf{U}}(\eta Y)-\mathbf{Q}_{W}(\eta Y)\bigr).
\end{equation*}
Its three terms are bounded as follows:
\begin{itemize}
\item $|\delta b|\le\mathsf{p}_{\delta}$;
\item $|\beta_{1}|\le C\mathsf{p}\mathsf{M}$, by \eqref{5.8:eq-response-comparison-3};
\item the difference of constraints is $\le C|\mathfrak{a}|_{y}|Y|_{y}\le C\mathsf{M}\mathsf{p}$, by
  \eqref{5.8:eq-response-comparison-1} and the size of $Y$.
\end{itemize}
Altogether,
\begin{equation}\label{5.8:eq-response-comparison-4}
|\tilde{\beta}|\le\mathsf{p}_{\delta}+C\mathsf{p}\mathsf{M}.
\end{equation}

\emph{The residual.} Let $u$ be a coarse variation with $u=0$ on $\Lambda_{0}$, and let $u''$ be the
fine variation attached to $u$ in the nesting bound for $d\mathcal{W}$. The residual identity of the
comparison of responses at smeared majorants is one of the statements assumed in the last item of the
hypotheses. It holds at the present complex data and gives
\begin{equation*}
d\Sigma_{\mathbf{U},\mathsf{m}}(Y)[u]=\mathrm{(I)}+\mathrm{(II)}+\mathrm{(III)}+\mathrm{(IV)}
+\mathrm{(V)}+\mathrm{(VI)}.
\end{equation*}
We bound the six members one at a time.

\begin{itemize}
\item[(I)] The member $\{d\mathcal{W}_{\mathbf{U}}(Y)-d\mathcal{W}_{W}(Y)\}[u]$ lies in
  $\sharp(C\mathsf{M}\mathsf{p})$. The background-change bound applies as stated and locally, with the
  local size $C\mathsf{p}(y)$ of $Y$. It gives the factor $|\mathfrak{a}|_{y}$, and
  $|\mathfrak{a}|_{y}\le\mathsf{M}(y)$ by \eqref{5.8:eq-response-comparison-1}.
\item[(II)] The member $d\mathcal{W}_{W}(Y)[u]-d\mathcal{W}_{\mathbf{U}'}(X')[u'']$ lies in
  $\sharp(C\mathsf{p}L^{-\beta\lambda})$. This is the nesting bound for $d\mathcal{W}$, applied as
  stated with its size parameter equal to $C\mathsf{p}$. It is local, and it holds uniformly for
  complex $X'$.
\item[(III)] The member is
  \begin{equation*}
  \langle(D_{\mathbf{U}}^{*}-D_{W}^{*})u,R_{\mathbf{U}}D_{\mathbf{U}}^{*}Y\rangle
  +\langle D_{W}^{*}u,(R_{\mathbf{U}}D_{\mathbf{U}}^{*}-R_{W}D_{W}^{*})Y\rangle,
  \end{equation*}
  and it lies in $\sharp(C\mathsf{M}\mathsf{p})$.
  \begin{itemize}
  \item In the first pairing, $\ell^{2}|R_{\mathbf{U}}D_{\mathbf{U}}^{*}Y|\le C\mathsf{p}$, by the
    weighted calculus (Lemma~\ref{5.8:lem-weighted-calculus}) at the weight $\mathsf{p}$ and the size
    of $Y$. The backgrounds $\mathbf{U}$ and $W=e^{i\eta\mathfrak{a}}\mathbf{U}$ differ by the
    one-form $\mathfrak{a}$, and $|\mathfrak{a}|_{y}\le\mathsf{M}(y)$.
  \item For the second pairing we apply Lipschitz dependence of the gauge projection
    (Lemma~\ref{5.8:lem-projection-lipschitz}(b)) with $(\mathsf{g},\mathsf{M}):=(C\mathsf{p},\mathsf{M})$
    and $a:=Y$. Its hypotheses hold: $C\mathsf{p}\in\mathfrak{M}_{\rho}$;
    $\mathsf{M}\in\mathfrak{M}_{\tau}$; $\rho+\tau\le 5\delta_{1}\le 6\delta_{1}$;
    $|\mathfrak{a}|_{y}\le\mathsf{M}(y)\le 1$; $\ell|Y|,\ell^{2}|\nabla Y|\le C\mathsf{p}$; and
    $\mathbf{U}$ and $W$ satisfy \cite[(3.35)]{B-CMP99b}. The lemma gives
    \begin{equation*}
    \ell^{2}\bigl|(R_{W}D_{W}^{*}-R_{\mathbf{U}}D_{\mathbf{U}}^{*})Y\bigr|(x)
    \le C\mathsf{M}(x)\mathsf{p}(x).
    \end{equation*}
  \end{itemize}
\item[(IV)] The member
  \begin{equation*}
  \langle D_{W}^{*}u,R_{W}D_{W}^{*}Y\rangle-\langle D'^{*}u'',\varsigma'\rangle
  \end{equation*}
  lies in $\sharp(C\mathsf{p}L^{-\theta\lambda})$. The proof is given below.
\item[(V)] The member
  \begin{equation*}
  \langle dC'(X')u'',\mathsf{m}'\rangle-\langle dC_{\mathbf{U}}(Y)u,\mathsf{m}\rangle
  \end{equation*}
  lies in $\sharp(C\mathsf{p}\mathsf{M})$. We write it as
  \begin{equation*}
  \bigl\langle dC'(X')u''-dC_{\mathbf{U}}(Y)u,\mathsf{m}'\bigr\rangle
  +\bigl\langle dC_{\mathbf{U}}(Y)u,\mathsf{m}'-\mathsf{m}\bigr\rangle .
  \end{equation*}
  The first term carries $|\mathsf{m}'|\le\mathfrak{r}_{0}$ and the mismatch
  $dC'(X')u''-dC_{\mathbf{U}}(Y)u$ of the two derivatives. The constraint comparison, differentiated,
  bounds this mismatch by $C\mathsf{p}$ times a mismatch factor, and this factor is at most
  $C\mathsf{M}$ by \eqref{5.8:eq-response-comparison-1} and \eqref{5.8:eq-response-comparison-3}.
  The second term carries the factor $C\mathsf{p}$ together with
  $|\mathsf{m}'(c)-\mathsf{m}(c)|\le\mathsf{M}(c)$. Both terms are bounded by the Schwarz inequality.
  They are block-local: they read only cells at distance at most $2$. So, by
  Lemma~\ref{5.8:lem-weighted-calculus}(ii), no decay rate is spent beyond a factor
  $e^{12\delta_{1}}$.
\item[(VI)] The member $\langle d\mathbf{Q}'(\eta'X')u'',\mu'\rangle$, for $u$ in
  $\ker d\mathbf{Q}_{\mathbf{U}}(\eta Y)$, lies in $\sharp(C\mathsf{p}\mathsf{M})$. It arises from the
  Euler--Lagrange equation of $X'$ in $\mathsf{B}$. Write $\epsilon[u]$ for the defect
  $d\mathbf{Q}'(\eta'X')u''$ of $u''$ with respect to the finer constraint, at such $u$. The estimate
  of this defect used at this member in the comparison of responses at smeared majorants gives
  \begin{equation*}
  |\epsilon[u]|(c)\le C\bigl(L^{1-\lambda}+|\mathfrak{a}|_{c}\bigr)\,\ell\sup|u|,
  \end{equation*}
  and $|\mu'|\le C\mathsf{p}$ by Lemma~\ref{5.8:lem-critical-size}. Moreover
  $L^{1-\lambda}=L\cdot L^{-\lambda}\le L\mathsf{M}$ by \eqref{5.8:eq-response-comparison-3}, and
  $|\mathfrak{a}|\le\mathsf{M}$.
\end{itemize}

We now prove the bound for (IV). As in the comparison of responses at smeared majorants, the
pairing splits as
\begin{equation*}
\langle D'^{*}u'',\varsigma'\rangle=\langle\hat{E}u,D'\varsigma'\rangle+\langle k'',D'\varsigma'\rangle,
\end{equation*}
where $\hat{E}u$ is the extension of $u$ supplied by the extension map and $k'':=u''-\hat{E}u$ is the
remainder of $u''$ with respect to it.

\emph{The term with $k''$.} It obeys
\begin{equation*}
|\langle k'',D'\varsigma'\rangle|\le\sum_{c}\eta^{d}\,C\eta|X'|\,|D'\varsigma'|\,|u|.
\end{equation*}
By the sizes of $X'$ and $\varsigma'$ this lies in $\sharp(C\mathsf{p}^{2}L^{-\lambda})$. Since
$\mathsf{p}\le a_{R}^{\sharp}/10$, this is contained in $\sharp(C\mathsf{p}L^{-\lambda})$.

\emph{The transfer identity.} The transfer identity for $\hat{E}$ gives
\begin{equation*}
\langle\hat{E}u,D'\varsigma'\rangle=\langle D_{W}^{*}\tilde{u},Q'_{1}\varsigma'\rangle
-\langle\tilde{u},f^{s}[\varsigma']\rangle,
\end{equation*}
where $\tilde{u}:=(1+\mathsf{k})^{-1}u$, with $|\mathsf{k}|,\ \ell|\nabla\mathsf{k}|\le
C\mathfrak{r}_{0}L^{-\lambda}$, and where the source term $f^{s}[\varsigma']$ of the transfer identity
obeys
\begin{equation*}
\ell^{3}|f^{s}[\varsigma']|\le\ell^{3}\operatorname{osc}_{2d\eta}(D'\varsigma')
+C\mathfrak{r}_{0}L^{-\lambda}\ell^{2}|\varsigma'|.
\end{equation*}
Since $\ell=L^{\lambda}\eta$, the oscillation term is bounded by
\begin{equation*}
\ell^{3}(2d\eta)^{\theta}[D'\varsigma']_{\theta}=(2d)^{\theta}L^{-\theta\lambda}
\ell^{3+\theta}[D'\varsigma']_{\theta}\le CL^{-\theta\lambda}\mathsf{p}.
\end{equation*}
The second term is $\le C\mathfrak{r}_{0}L^{-\lambda}\mathsf{p}$. Both bounds come from the sizes of
$\varsigma'$, that is, from Lemma~\ref{5.8:lem-gauge-component} in $\mathsf{B}$, and they hold locally.
Thus $\ell^{3}|f^{s}[\varsigma']|\le C\mathsf{p}L^{-\theta\lambda}$ locally.

\emph{The projection difference.} Put
$\varphi:=R_{W}D_{W}^{*}Y-Q'_{1}\varsigma'$. Since $\varsigma'=R'D'^{*}X'$,
\begin{equation*}
\varphi=R_{W}\bigl(D_{W}^{*}Y-Q'_{1}D'^{*}X'\bigr)+\bigl(R_{W}Q'_{1}-Q'_{1}R'\bigr)D'^{*}X'.
\end{equation*}
We bound the two terms separately.
\begin{itemize}
\item In the first term, the divergence comparison bounds $D_{W}^{*}Y-Q'_{1}D'^{*}X'$ locally by
  $C\mathsf{p}L^{-\beta\lambda}$. Lemma~\ref{5.8:lem-weighted-calculus}, applied to $R_{W}$ at the
  weight $\mathsf{p}L^{-\beta\lambda}$, carries this bound through $R_{W}$.
\item For the second term we apply nesting of fine and coarse gauge projections
  (Lemma~\ref{5.8:lem-projection-nesting}) with $\mathsf{f}:=C\mathsf{p}$ and $f':=D'^{*}X'$. Its
  hypotheses hold: $\mathbf{U}'$ satisfies the field clauses; $C\mathsf{p}\in\mathfrak{M}_{\rho}$ with
  $\rho\le 5\delta_{1}$; and $\ell^{2}|D'^{*}X'|\le C\mathsf{p}$ by
  Lemma~\ref{5.8:lem-critical-size}. The lemma bounds the second term by
  $C\mathsf{p}L^{-\beta\lambda}$.
\end{itemize}
Hence $\ell^{2}|\varphi|\le C\mathsf{p}L^{-\beta\lambda}$.

\emph{Assembly.} Using $R_{W}D_{W}^{*}Y=\varphi+Q'_{1}\varsigma'$, the member (IV) equals
\begin{equation*}
\langle D_{W}^{*}u,\varphi\rangle+\langle D_{W}^{*}(u-\tilde{u}),Q'_{1}\varsigma'\rangle
+\langle\tilde{u},f^{s}[\varsigma']\rangle-\langle k'',D'\varsigma'\rangle.
\end{equation*}
The four terms are bounded as follows.
\begin{itemize}
\item The first lies in $\sharp(C\mathsf{p}L^{-\beta\lambda})$, by the bound on $\varphi$.
\item The second carries the factor $\mathsf{k}$, through $u-\tilde{u}=(1+\mathsf{k})^{-1}\mathsf{k}u$,
  and lies in $\sharp(C\mathsf{p}L^{-\lambda})$ by the bounds on $\mathsf{k}$ and $\varsigma'$.
\item The third lies in $\sharp(C\mathsf{p}L^{-\theta\lambda})$, by the bound on $f^{s}[\varsigma']$.
\item The fourth lies in $\sharp(C\mathsf{p}L^{-\lambda})$, as shown above.
\end{itemize}
By \eqref{5.8:eq-response-comparison-3}, (IV) lies in $\sharp(C\mathsf{p}L^{-\theta\lambda})$.

Collecting the six members and using \eqref{5.8:eq-response-comparison-3} once more, we obtain
\begin{equation}\label{5.8:eq-response-comparison-5}
d\Sigma_{\mathbf{U},\mathsf{m}}(Y)\in\sharp(C\mathsf{p}\mathsf{M})\qquad\text{on }
\ker d\mathbf{Q}_{\mathbf{U}}(\eta Y).
\end{equation}

\emph{Conclusion.} Put $w:=\mathsf{p}\mathsf{M}+\mathsf{p}_{\delta}$. This weight lies in
$\mathfrak{M}_{\rho+\tau}$, for two reasons:
\begin{itemize}
\item multiplying the defining inequalities of $\mathsf{p}\in\mathfrak{M}_{\rho}$ and
  $\mathsf{M}\in\mathfrak{M}_{\tau}$ gives $\mathsf{p}\mathsf{M}\in\mathfrak{M}_{\rho+\tau}$;
\item a sum of two members of $\mathfrak{M}_{\rho+\tau}$ lies in $\mathfrak{M}_{\rho+\tau}$.
\end{itemize}
Its rate satisfies $\rho+\tau\le 5\delta_{1}\le 6\delta_{1}$.

We apply weighted stability of the critical point (Lemma~\ref{5.8:lem-critical-stability}) at
$(\mathbf{U},\mathsf{m})$, with $B:=b$, $\tilde{X}:=Y$ and the weight $w$. Its hypotheses hold:
\begin{itemize}
\item $X=\mathfrak{X}(b;\mathbf{U},\mathsf{m})$ with $\sup|b|\le\sup\mathsf{p}\le a_{R}^{\sharp}$;
\item $\ell|Y|,\ell^{2}|\nabla Y|\le\varepsilon_{3}$, by the sizes established above;
\item $\mathbf{Q}_{\mathbf{U}}(\eta Y)=b+\tilde{\beta}$, by the constraint step;
\item $d\Sigma_{\mathbf{U},\mathsf{m}}(Y)\in\sharp(\kappa)$ on $\ker d\mathbf{Q}_{\mathbf{U}}(\eta Y)$
  with $\kappa:=C\mathsf{p}\mathsf{M}$, by \eqref{5.8:eq-response-comparison-5}.
\end{itemize}
By \eqref{5.8:eq-response-comparison-4},
\begin{equation*}
\bigl\||\tilde{\beta}|+\kappa\bigr\|_{w}\le\sup\frac{\mathsf{p}_{\delta}+C\mathsf{p}\mathsf{M}}
{\mathsf{p}\mathsf{M}+\mathsf{p}_{\delta}}\le C.
\end{equation*}
Hence
\begin{equation*}
|Y-X|_{x}\le C_{S}^{\sharp}w(x)\bigl\||\tilde{\beta}|+\kappa\bigr\|_{w}\le C_{S}^{\sharp}C\,w(x)
=C_{S}^{\sharp}C\bigl(\mathsf{p}(x)\mathsf{M}(x)+\mathsf{p}_{\delta}(x)\bigr).
\end{equation*}
Together with the collar estimate and the reduction to $A=1$, this is
\eqref{5.8:eq-response-comparison-2}.
\end{proof}

\begin{remark}
We add four observations on the proof above.

First, in the comparison of responses at smeared majorants, exactly five of the lemmas used carry the
$\delta_{1}$-smear in their statements. They are the lemmas on:
\begin{itemize}
\item the size of the critical point;
\item its gauge component;
\item the Lipschitz dependence of the gauge projection;
\item the nesting of gauge projections;
\item the stability of the critical point.
\end{itemize}
Their versions at tempered profiles are Lemmas~\ref{5.8:lem-critical-size},
\ref{5.8:lem-gauge-component}, \ref{5.8:lem-projection-lipschitz},
\ref{5.8:lem-projection-nesting} and~\ref{5.8:lem-critical-stability}. The remaining lemmas of that
proof serve only its datum-free part, which enters unchanged.

Second, no decay rate is exchanged between $\mathsf{p}$ and $\mathsf{M}$ beyond the condition
$\rho+\tau\le 5\delta_{1}$. Nothing is interpolated, and no root of a profile is taken.

Third, linearised in $b$, the estimate \eqref{5.8:eq-response-comparison-2} is the combination of the
linearised difference of the two kernels, \eqref{5.8:eq-kernel-difference-a}, with part (a) of the
sharp rate window (Proposition~\ref{5.8:prop-rate-window}).

Fourth, reality is required of the backgrounds only. Ba{\l}aban's \cite[Theorems~3.1--3.3]{B-CMP99b}
are applied at the real backgrounds $\mathbf{U}$, $\mathbf{U}'$, $W$, while
$(b,\delta b,b_{c},\mathsf{m}',\mathsf{m})$ are complex. The comparison of responses at smeared
majorants is stated for real $b,\delta b$, but its proof uses the reality of $b$ at no point:
\begin{itemize}
\item the critical-point size bound allows complex data;
\item the nesting bound for $d\mathcal{W}$ holds uniformly for complex $X'$;
\item the Euler--Lagrange equation, the identity for the gauge component, the constraint comparison
  and the residual identity are holomorphic in $(b,\delta b)$ and hold at real points.
\end{itemize}
This reading is the last hypothesis of the theorem. It is taken by statement.
\end{remark}

\begin{corollary}[Response comparison in local, uniform and framed forms]
\label{5.8:cor-comparison-forms}
Assume the following three statements.
\begin{enumerate}
\item[(1)] The two-cutoff comparison of responses at tempered profiles,
Theorem~\ref{5.8:thm-response-comparison}.
\item[(2)] Part (A) of the statement on real presented pairs, the base-mismatch bound, used
as stated. It gives, for the base mismatch one-form $\mathfrak{a}$ of such a pair $p$, the
bound $|\mathfrak{a}|_{y} \le \mathsf{T}_{p}(y)$ for every $y$. For a
presentation over an entry, part (b) of the corresponding statement for presentations over
an entry is used in its place, also as stated.
\item[(3)] The frame bound. A datum prescribed in Ba{\l}aban's frames in both machines
$\mathsf{A}$ and $\mathsf{B}$ has two prescriptions that differ by conjugation with a
gauge element $g$ satisfying
\begin{equation*}
|g-\mathbf{1}| \le C_{N}\,\mathsf{S}\,L^{-\theta_{2}\lambda},
\end{equation*}
where $C_{N}$ and $\mathsf{S}$ are the constants of that bound.
\end{enumerate}
Let $p$ be a real presented pair, or a presentation over an entry, with mismatch tally
satisfying $\sup \mathsf{T}_{p} \le \varepsilon_{D}$. Let $P$ and $P'$ be the presentation
data of $p$ in the machines $\mathsf{A}$ and $\mathsf{B}$ respectively (the primed letters
refer to $\mathsf{B}$, as for $\mathfrak{S}'$). Let $(\mathfrak{S}',\mathfrak{S})$ and
$\mathbf{U}'$, $\mathbf{U}$ be the exactly paired sets and the backgrounds of $p$, in the
setting of Theorem~\ref{5.8:thm-response-comparison}. Set
\begin{equation*}
\mathbb{H} := \mathbb{H}_{\mathfrak{S}}(b;P), \qquad
\mathbb{H}' := \mathbb{H}_{\mathfrak{S}'}\bigl((b+\delta b)\oplus b_{c};P'\bigr).
\end{equation*}
Put $\mathsf{t}_{p} := \sup_{x} \mathsf{T}_{p}(x)$. The data $b$, $\delta b$, $b_{c}$, the
profile $\mathsf{p}_{\delta}$, the data vectors $\mathsf{m}'$, $\mathsf{m}$ of $p$ in the two
machines, and the smallness conditions are those of
Theorem~\ref{5.8:thm-response-comparison}, taken at the mismatch majorant
\begin{equation*}
\mathsf{M} := C_{\mathsf{m}}\Psi^{B}_{p} \wedge 1 \quad\text{in (a)}, \qquad
\mathsf{M} := C'\bigl(\mathsf{t}_{p} + L^{-\theta\lambda}\bigr) \quad\text{in (b)}.
\end{equation*}
Here $C_{\mathsf{m}}$ is the constant with which the smeared currency $\Psi^{B}_{p}$
majorises the data-vector mismatch $|\mathsf{m}'-\mathsf{m}|$ of a presented pair, and
$C' \ge c_{1}$ is a constant, with $c_{1}$ the
constant of \eqref{5.8:eq-paired-sets-b}. Then the following hold.
\begin{enumerate}
\item[(a)] (Local form; this is the form used in the uniqueness theorem.) For every
$\mathsf{p} \in \mathfrak{M}_{\rho}$ with
$\rho \le 4\delta_{1} = \delta$,
\begin{equation}\label{5.8:eq-comparison-forms-1}
\bigl|Y[\mathbb{H}';\mathbf{U}'] - \mathbb{H}\bigr|_{x}
\le C_{B}^{\sharp}\,L\,C_{\mathsf{m}}
\bigl\{\mathsf{p}\,\Psi^{B}_{p} + \mathsf{p}_{\delta}\bigr\}(x).
\end{equation}
This is part (B) of the statement on real presented pairs whose part (A) is hypothesis (2),
with its smeared datum majorant $\mathsf{P}_{b}$ replaced by $\mathsf{p}$.
\item[(b)] (Uniform form with the local index.) For every $\mathsf{p}$ as in (a), with a
constant $C$ (a different constant from $C'$),
\begin{equation}\label{5.8:eq-comparison-forms-2}
\bigl|Y[\mathbb{H}';\mathbf{U}'] - \mathbb{H}\bigr|_{x}
\le C\bigl\{\mathsf{p}(x)\bigl[\mathsf{t}_{p} + L^{-\theta\lambda(x)}\bigr]
+ \mathsf{p}_{\delta}(x)\bigr\}.
\end{equation}
\item[(c)] (Framed form.) A datum prescribed in Ba{\l}aban's frames in both machines is read
as the pair $(b+\delta b,\,b)$ with $\delta b = (\operatorname{Ad}_{g}-1)b$, where $g$ is
as in (3). In this case one may take
\begin{equation}\label{5.8:eq-comparison-forms-3}
\mathsf{p}_{\delta} := 2C_{N}\,\mathsf{S}\,\mathsf{p}\,L^{-\theta_{2}\lambda}.
\end{equation}
\end{enumerate}
\end{corollary}

\begin{proof}
(a) We apply Theorem~\ref{5.8:thm-response-comparison} with the mismatch majorant
\begin{equation*}
\mathsf{M} := C_{\mathsf{m}}\Psi^{B}_{p} \wedge 1 \in \mathfrak{M}^{L}_{\delta_{1}} .
\end{equation*}
This $\mathsf{M}$ lies in $\mathfrak{M}^{L}_{\tau}$ with $\tau = \delta_{1}$. Since
$\rho \le 4\delta_{1}$, we have $\rho + \tau \le 5\delta_{1}$, which is the condition on the
rates in that theorem. Because of the truncation at $1$, we also have $\mathsf{M} \le 1$.

The base mismatch one-form $\mathfrak{a}$ of $p$ satisfies
\begin{equation*}
|\mathfrak{a}|_{y} \le \mathsf{T}_{p}(y) \le \mathsf{T}_{p}^{\sim}(y)
\end{equation*}
for every $y$. The first inequality is hypothesis (2); the second is the domination of a
profile by its $\delta_{1}$-smear $f^{\sim}$, fixed with the conventions on the smear. In
particular, $\sup_{y}|\mathfrak{a}|_{y} \le \sup \mathsf{T}_{p} \le \varepsilon_{D}$, as the
theorem requires. By hypothesis, the data vectors satisfy the condition of that theorem at
this $\mathsf{M}$, namely
$|\mathsf{m}'(c)-\mathsf{m}(c)| \le \mathsf{M}(c)$ for every $c$. With this $\mathsf{M}$, the
conclusion of Theorem~\ref{5.8:thm-response-comparison} is
\eqref{5.8:eq-comparison-forms-1}.

(b) By the summability bound \eqref{5.8:eq-paired-sets-b}, the $\delta_{1}$-smear of the tally
is controlled by its supremum:
\begin{equation*}
\mathsf{T}_{p}^{\sim} \le c_{1}\,\mathsf{t}_{p}.
\end{equation*}
We now apply Theorem~\ref{5.8:thm-response-comparison} with
\begin{equation*}
\mathsf{M} := C'\bigl(\mathsf{t}_{p} + L^{-\theta\lambda}\bigr)
\in \mathfrak{M}^{L}_{\delta_{1}/16} .
\end{equation*}
Combining the chain of part (a) with the last inequality and with $C' \ge c_{1}$, we get
\begin{equation*}
|\mathfrak{a}|_{y} \le \mathsf{T}_{p}^{\sim}(y) \le c_{1}\,\mathsf{t}_{p} \le \mathsf{M}(y)
\end{equation*}
for every $y$; and since $\rho \le 4\delta_{1}$, the rates satisfy
$\rho + \delta_{1}/16 \le 5\delta_{1}$. By hypothesis, the data, the data vectors and the
smallness conditions are those of that theorem at this $\mathsf{M}$; its conclusion is
\eqref{5.8:eq-comparison-forms-2}.

(c) For a datum given in Ba{\l}aban's frames in both machines, we have
$\delta b = (\operatorname{Ad}_{g}-1)b$. The factor $2$ arises as follows. Write
\begin{equation*}
\operatorname{Ad}_{g} b - b = g\,b\,g^{-1} - b = (g-\mathbf{1})\,b\,g^{-1} + b\,(g^{-1}-\mathbf{1}).
\end{equation*}
Since $g$ is unitary, $|g^{-1}| = 1$ and
$|g^{-1}-\mathbf{1}| = |(g-\mathbf{1})^{\dagger}| = |g-\mathbf{1}|$. Hence
\begin{equation*}
|\delta b| \le 2\,|g-\mathbf{1}|\,|b|
\le 2C_{N}\,\mathsf{S}\,L^{-\theta_{2}\lambda}\,\mathsf{p},
\end{equation*}
where the second inequality uses the frame bound (3) and $|b| \le \mathsf{p}$. That the
profile \eqref{5.8:eq-comparison-forms-3} belongs to the weight class required of
$\mathsf{p}_{\delta}$ in Theorem~\ref{5.8:thm-response-comparison} follows from rule (w3)
among the rules for tempered weights collected in \eqref{5.8:eq-tempered-weights-a}.
\end{proof}

\begin{corollary}[Increment bound for a move by a cut response]\label{5.8:cor-cut-move}
Assume the following three statements:
\begin{enumerate}
\item the response comparison in its local form, Corollary~\ref{5.8:cor-comparison-forms}(a);
\item the decomposition of the blocked-increment map: for a form $h$ and a background $V$,
\begin{equation*}
Y[h;V] = Q_{1}h + N_{2}[h],
\end{equation*}
where $Q_{1} = Q_{1}(V)$ is a linear operator depending on the background and $N_{2}[h]$ is the
nonlinear remainder, and there is a
constant $C_{W}$ such that
\begin{equation*}
|N_{2}[h]| \le C_{W}\,\eta\,\sup|h|^{2};
\end{equation*}
moreover, for a scalar $\tau$, the commutator $[Q_{1},\tau]h' := Q_{1}(\tau h') - \tau Q_{1}h'$
obeys
\begin{equation*}
|[Q_{1},\tau]h'| \le C\eta\,|\nabla\tau|\,|h'|;
\end{equation*}
\item the increment dictionary.
\end{enumerate}
Work in the setting of Corollary~\ref{5.8:cor-comparison-forms}(a), with $p$ (a real presented
pair, or a presentation over an entry), the responses $\mathbb{H}$ and $\mathbb{H}'$, the
background $\mathbf{U}'$, the profiles $\mathsf{p}$, $\mathsf{p}_{\delta}$ and the currency
$\Psi^{B}_{p}$ as there. Assume moreover $\sup\mathsf{p} \le a_{R}^{\sharp}$, the ceiling of
Lemma~\ref{5.8:lem-critical-size}. Let $\tau$ be a
scalar, the same for the coarser and for the finer cutoff, and put
\begin{equation*}
\bar{\tau}(y) := \sup_{\tilde{\Delta}(y)}\max\{|\tau|,\ \ell|\nabla\tau|,\ \ell^{2}|\nabla^{2}\tau|\},
\qquad
\mathsf{H}^{\sharp} := C_{s}^{\sharp}\,\bar{\tau}\,\mathsf{p},
\end{equation*}
with $C_{s}^{\sharp}$ the constant of Lemma~\ref{5.8:lem-critical-size}; $\bar{\tau}(y)$ is the
function denoted $a(y)$ in the bound for the cutoff scalar at a sharp
profile. Assume $\sup\mathsf{H}^{\sharp} \le a_{R}^{\sharp}$. Then the following hold.

\smallskip
\noindent The mismatch $e := Y[\tau\mathbb{H}';\mathbf{U}'] - \tau\mathbb{H}$ obeys, at every $x$,
the first inequality of \eqref{5.8:eq-cut-move-a} below. For real data $b$ and real $\tau$, the
pair $p^{+}$ obtained by extending $p$ by $(\tau\mathbb{H}',\tau\mathbb{H})$ satisfies, at every
$x$, the second:
\begin{equation}\label{5.8:eq-cut-move-a}
\begin{split}
|e|_{x} &\le \bar{\tau}(x)\cdot C\{\mathsf{p}\Psi^{B}_{p} + \mathsf{p}_{\delta}\}(x)
+ C L^{-\lambda(x)}\,\mathsf{H}^{\sharp}(x)\,[1 + \mathsf{H}^{\sharp}(x)],\\
\mathsf{T}_{p^{+}}(x) - \mathsf{T}_{p}(x) &= |A_{p^{+}} - A_{p}|_{x}\\
&\le \bigl(1 + C L^{-\lambda(x)}\mathsf{H}^{\sharp}(x)\bigr)|e|_{x}
+ C L^{-\lambda(x)}\mathsf{H}^{\sharp}(x)\,\mathsf{T}_{p}(x).
\end{split}
\end{equation}
In the first inequality $\tau$ may be complex. If moreover $\ell|\nabla\tau| \le \bar{\tau}/M$ on
$\tilde{\Delta}(x)$, the part of the term $1$ in the bracket that comes from the derivative of
$\tau$ becomes $M^{-1}$, and the first inequality holds in the form
\begin{equation*}
\begin{split}
|e|_{x} &\le \bar{\tau}(x)\cdot C\{\mathsf{p}\Psi^{B}_{p} + \mathsf{p}_{\delta}\}(x)\\
&\quad + C L^{-\lambda(x)}\,\mathsf{H}^{\sharp}(x)\,
\bigl[M^{-1} + C_{s}^{\sharp}\mathsf{p}(x) + \mathsf{H}^{\sharp}(x)\bigr].
\end{split}
\end{equation*}

\smallskip
\noindent The cut responses satisfy
\begin{equation}\label{5.8:eq-cut-move-b}
\begin{gathered}
\tau\mathbb{H} \in \operatorname{Inc}_{\beta}(C\mathsf{H}^{\sharp}), \qquad
\ell^{3}|DD^{*}(\tau\mathbb{H})| \le C\mathsf{H}^{\sharp},\\
\tau\mathbb{H}' \in \operatorname{Inc}_{\beta}(C\mathsf{H}^{\sharp}), \qquad
\ell^{3}|DD^{*}(\tau\mathbb{H}')| \le C\mathsf{H}^{\sharp}.
\end{gathered}
\end{equation}
Consequently $p^{+}$ is presented, and $\mathfrak{r}$ grows by $C\sup\mathsf{H}^{\sharp}$.

\smallskip
\noindent No slowness is asked of $\bar{\tau}$: no condition on the variation of $\bar{\tau}$
enters beyond its definition.
\end{corollary}

\begin{proof}
Throughout, $C$ denotes a constant that may change from one occurrence to the next, and the
background in $Y$, $Q_{1}$ and $N_{2}$ is $\mathbf{U}'$.

\emph{The decomposition of $e$.} By the decomposition of the blocked-increment map, applied to
$h = \tau\mathbb{H}'$ and to $h = \mathbb{H}'$, and by the definition of the commutator,
\begin{align*}
Y[\tau\mathbb{H}';\mathbf{U}'] &= \tau Q_{1}\mathbb{H}' + [Q_{1},\tau]\mathbb{H}'
+ N_{2}[\tau\mathbb{H}'],\\
\tau\, Y[\mathbb{H}';\mathbf{U}'] &= \tau Q_{1}\mathbb{H}' + \tau N_{2}[\mathbb{H}'].
\end{align*}
Subtracting the second line from the first and then subtracting $\tau\mathbb{H}$ gives
\begin{equation}\label{5.8:eq-cut-move-2}
e = \tau\bigl(Y[\mathbb{H}';\mathbf{U}'] - \mathbb{H}\bigr) + [Q_{1},\tau]\mathbb{H}'
+ N_{2}[\tau\mathbb{H}'] - \tau N_{2}[\mathbb{H}'].
\end{equation}
The identity uses no reality of $\tau$, so what follows holds for complex $\tau$.

\emph{Sizes.} By the hypothesis $\sup\mathsf{p} \le a_{R}^{\sharp}$, the
critical-point size lemma (Lemma~\ref{5.8:lem-critical-size}) applies to $\mathbb{H}$ and to
$\mathbb{H}'$: the scaled local sizes of $\mathbb{H}$ and $\mathbb{H}'$ at $y$ are at most
$C_{s}^{\sharp}\mathsf{p}(y)$.
All sizes
below are read locally at $x$, on $\tilde{\Delta}(x)$, where $|\tau| \le \bar{\tau}(x)$ and
$\ell|\nabla\tau| \le \bar{\tau}(x)$. In particular the local sizes of $\tau\mathbb{H}$ and
$\tau\mathbb{H}'$ at $x$ are at most $\bar{\tau}C_{s}^{\sharp}\mathsf{p} = \mathsf{H}^{\sharp}$.
By the definition of the scaled length, $\ell = L^{\lambda(x)}\eta$ on $\tilde{\Delta}(x)$, so
$\eta/\ell = L^{-\lambda(x)}$: measured in $\ell$ at $x$, the factor $\eta$ of the two bounds of
the blocked-increment map is the factor $L^{-\lambda(x)}$ of \eqref{5.8:eq-cut-move-a}, and
$\eta|\nabla\tau| = (\eta/\ell)\,\ell|\nabla\tau|$.

We bound the four terms of \eqref{5.8:eq-cut-move-2} in turn.

\emph{First term.} By Corollary~\ref{5.8:cor-comparison-forms}(a) and $|\tau| \le \bar{\tau}(x)$,
\begin{equation*}
\bigl|\tau\bigl(Y[\mathbb{H}';\mathbf{U}'] - \mathbb{H}\bigr)\bigr|_{x}
\le \bar{\tau}(x)\cdot C\{\mathsf{p}\Psi^{B}_{p} + \mathsf{p}_{\delta}\}(x),
\end{equation*}
with $C$ the constant of that corollary; this is the first term of the first inequality of
\eqref{5.8:eq-cut-move-a}.

\emph{Commutator.} By the commutator bound and the size of $\mathbb{H}'$,
\begin{equation*}
|[Q_{1},\tau]\mathbb{H}'|_{x} \le C\,\frac{\eta}{\ell}\,\ell|\nabla\tau|\,C_{s}^{\sharp}\mathsf{p}
\le C L^{-\lambda(x)}\,\bar{\tau}\,C_{s}^{\sharp}\mathsf{p}
= C L^{-\lambda(x)}\mathsf{H}^{\sharp}(x).
\end{equation*}

\emph{Quadratic remainders.} By the bound on $N_{2}$, measured in $\ell$, and the size of
$\tau\mathbb{H}'$,
\begin{equation*}
|N_{2}[\tau\mathbb{H}']|_{x} \le C_{W}\,\frac{\eta}{\ell}\,(\mathsf{H}^{\sharp})^{2}
= C_{W} L^{-\lambda(x)}\mathsf{H}^{\sharp}(x)^{2}.
\end{equation*}
By the same bound, the size of $\mathbb{H}'$ and $|\tau| \le \bar{\tau}$,
\begin{equation*}
|\tau N_{2}[\mathbb{H}']|_{x} \le \bar{\tau}\,C_{W}\,\frac{\eta}{\ell}\,(C_{s}^{\sharp}\mathsf{p})^{2}
= C_{W} L^{-\lambda(x)}\,\mathsf{H}^{\sharp}\,C_{s}^{\sharp}\mathsf{p}
\le C L^{-\lambda(x)}\mathsf{H}^{\sharp}(x),
\end{equation*}
where $C_{s}^{\sharp}\mathsf{p} \le C_{s}^{\sharp}a_{R}^{\sharp}$ was used.

Adding the four bounds gives the first inequality of \eqref{5.8:eq-cut-move-a}.

\emph{The case $\ell|\nabla\tau| \le \bar{\tau}/M$.} The commutator is the only term of
\eqref{5.8:eq-cut-move-2} in which a derivative of $\tau$ appears, and its bound above used
$\ell|\nabla\tau| \le \bar{\tau}(x)$. If $\ell|\nabla\tau| \le \bar{\tau}/M$ on
$\tilde{\Delta}(x)$, the same computation gives
\begin{equation*}
|[Q_{1},\tau]\mathbb{H}'|_{x} \le C L^{-\lambda(x)}M^{-1}\mathsf{H}^{\sharp}(x),
\end{equation*}
so that the share of the term $1$ contributed by the commutator is $M^{-1}$.
The bounds of the first term and of the two quadratic remainders involve no derivative of
$\tau$ and are unchanged. Adding to this bound the bound of the first term, the bound
$C_{W}L^{-\lambda(x)}\mathsf{H}^{\sharp}(x)^{2}$ of $N_{2}[\tau\mathbb{H}']$ and the bound
$C_{W}L^{-\lambda(x)}\mathsf{H}^{\sharp}(x)\,C_{s}^{\sharp}\mathsf{p}(x)$ of
$\tau N_{2}[\mathbb{H}']$, taken before the last step of its estimate, gives the form of the
first inequality with the bracket
$[M^{-1} + C_{s}^{\sharp}\mathsf{p}(x) + \mathsf{H}^{\sharp}(x)]$ stated in the corollary.

\emph{The extended pair.} Let $b$ and $\tau$ be real, so that $\tau\mathbb{H}$ and
$\tau\mathbb{H}'$ are real and $p^{+}$ is a real pair. The increment dictionary expresses the
change of the mismatch tally under the extension of $p$ by $(\tau\mathbb{H}',\tau\mathbb{H})$ as
$\mathsf{T}_{p^{+}}(x) - \mathsf{T}_{p}(x) = |A_{p^{+}} - A_{p}|_{x}$ at every $x$, and bounds it
by the blocked
mismatch $e$ of the added pair and the sizes of its members. Inserting the local sizes of
$\tau\mathbb{H}$ and $\tau\mathbb{H}'$, at most $\mathsf{H}^{\sharp}$, with the factor
$L^{-\lambda(x)}$ of the scaled reading of $\eta$, gives the second inequality of
\eqref{5.8:eq-cut-move-a}.

\emph{Proof of \eqref{5.8:eq-cut-move-b}.} Let $X$ be $\mathbb{H}$ or $\mathbb{H}'$. The
members of the critical-point size lemma (Lemma~\ref{5.8:lem-critical-size}) bound
$N^{\Delta}_{y}(X)$, $\ell^{2}|D^{*}X|$ and $\ell^{2+\beta}[D^{*}X]_{\beta}$ by
$C_{s}^{\sharp}\mathsf{p}(y)$. The last display of the lemma on the gauge component of the
critical point (Lemma~\ref{5.8:lem-gauge-component}), applied in the same setting, bounds
$\ell^{3}|DD^{*}X|$ and $\ell^{3+\theta}[DD^{*}X]_{\theta}$ by $C_{\varsigma}^{\sharp}\mathsf{p}(y)$.
By the Leibniz rule, $D^{*}(\tau X)$, its $\beta$-H\"older difference quotients and
$DD^{*}(\tau X)$, and the derivatives entering $N^{\Delta}_{y}(\tau X)$, are each a sum of
$\tau$ times the corresponding expression for $X$ and of terms in which one or two derivatives
fall on $\tau$. The first kind is bounded by $\bar{\tau}$ times the bounds just listed. The
second kind is bounded through $\ell|\nabla\tau|$ and
$\ell^{2}|\nabla^{2}\tau|$, hence by $\bar{\tau}$, times members of
Lemma~\ref{5.8:lem-critical-size} of lower order. Since
$\bar{\tau}C_{s}^{\sharp}\mathsf{p} = \mathsf{H}^{\sharp}$ and
$\bar{\tau}C_{\varsigma}^{\sharp}\mathsf{p} \le C\mathsf{H}^{\sharp}$, this gives
$\tau X \in \operatorname{Inc}_{\beta}(C\mathsf{H}^{\sharp})$ and
$\ell^{3}|DD^{*}(\tau X)| \le C\mathsf{H}^{\sharp}$, which is \eqref{5.8:eq-cut-move-b}.
These are the bounds asked of an added pair; by the hypothesis
$\sup\mathsf{H}^{\sharp} \le a_{R}^{\sharp}$ the sizes of the added members are at most
$Ca_{R}^{\sharp}$ everywhere, so $p^{+}$ is presented, and $\mathfrak{r}$ grows by
$C\sup\mathsf{H}^{\sharp} \le Ca_{R}^{\sharp}$.

Only $\bar{\tau}$ itself, through $|\tau|$, $\ell|\nabla\tau|$ and $\ell^{2}|\nabla^{2}\tau|$
on $\tilde{\Delta}(x)$, has entered; no property of the variation of $\bar{\tau}$ was used.
\end{proof}

\begin{corollary}[Two-cutoff comparison of the multipliers]\label{5.8:cor-multiplier-comparison}
Assume the following:
\begin{itemize}
\item the two-cutoff comparison of responses at tempered profiles
(Theorem~\ref{5.8:thm-response-comparison}) in its setting, together with the identity for the
derivative of $\Sigma_{\mathbf{U},\mathsf{m}}$ established in its proof, which holds for every
coarse direction vanishing on $\Lambda_{0}$;
\item the bound \eqref{5.8:eq-critical-stability-a} of the weighted stability of the critical
point (Lemma~\ref{5.8:lem-critical-stability});
\item the single-cutoff multiplier algebra of the chart, namely the single-cutoff multiplier
corollary and the chart identities $(\ell Q)H_{0} = 1$ and
$T H_{0} = (\ell Q)^{\dagger}\Theta$, where $T$ denotes the linear part of the stationarity
equation of the chart $\mathsf{A}_{0}$, and $H_{0}$, $\Theta$ are the operators of the chart
entering Lemma~\ref{5.8:lem-weighted-chart}.
\end{itemize}
In the setting of Theorem~\ref{5.8:thm-response-comparison}, with
$b_{c} = 0$, let $\mu'$ and $\mu$ be the Lagrange multipliers $\mu_{B}$ of the two constrained
critical points compared there. Then at every paired data site $c$
\begin{equation}\label{5.8:eq-multiplier-comparison-1}
|\mu'(c) - \mu(c)| \le C_{\mu}^{\sharp}\,\{A\mathsf{p}\mathsf{M} + \mathsf{p}_{\delta}\}(c).
\end{equation}
\end{corollary}

\begin{proof}
In the proof of Theorem~\ref{5.8:thm-response-comparison}, the identity for the derivative
of the functional $\Sigma_{\mathbf{U},\mathsf{m}}$ at the configuration $Y$ fixed in that proof
(not to be confused with the blocked-increment map $Y[h;V]$) is
applied to a direction $u$. The tangency of $u$ is used there for one purpose only: to drop
the term $\langle d\mathbf{Q}_{\mathbf{U}}(\eta Y)u,\mu'\rangle$. By the first hypothesis, the
identity itself holds for every coarse direction $u$ vanishing on $\Lambda_{0}$:
\begin{equation}\label{5.8:eq-multiplier-comparison-2}
d\Sigma_{\mathbf{U},\mathsf{m}}(Y)[u]
= \langle d\mathbf{Q}_{\mathbf{U}}(\eta Y)u,\mu'\rangle + \mathfrak{r}[u],
\qquad \mathfrak{r} \in \sharp(C\mathsf{p}\mathsf{M}).
\end{equation}

We write \eqref{5.8:eq-multiplier-comparison-2} in the chart, and denote by $\tilde{A}'$ and
$A'$ the chart coordinates of the primed and of the unprimed critical point. For a direction
$v$, let $\tilde{t}_{v}$ denote the coarse direction associated with the chart direction $v$,
as in the display defining $t_{v}$ in the proof of Lemma~\ref{5.8:lem-critical-stability}.
Then for every
$v$
\begin{equation}\label{5.8:eq-multiplier-comparison-3}
\langle v, T\tilde{A}'\rangle + d\mathcal{N}(\tilde{A}')[v]
= \langle (\ell Q)v,\mu'\rangle + \mathfrak{r}[\tilde{t}_{v}].
\end{equation}
The same identity holds at $A'$, with the pair $(\mu',\mathfrak{r})$ replaced by $(\mu,0)$;
this is the identity of the single-cutoff multiplier corollary:
\begin{equation}\label{5.8:eq-multiplier-comparison-4}
\langle v, TA'\rangle + d\mathcal{N}(A')[v] = \langle (\ell Q)v,\mu\rangle .
\end{equation}

Put $E := \tilde{A}' - A'$. Subtract \eqref{5.8:eq-multiplier-comparison-4} from
\eqref{5.8:eq-multiplier-comparison-3} and evaluate at $v = H_{0}\delta_{c}$, $\delta_{c}$
being the data vector equal to $1$ at $c$ and $0$ elsewhere. By
$(\ell Q)H_{0} = 1$ we have $(\ell Q)H_{0}\delta_{c} = \delta_{c}$, so the right-hand side
contributes $(\mu'-\mu)(c) + \mathfrak{r}[\tilde{t}_{H_{0}\delta_{c}}]$. By
$TH_{0} = (\ell Q)^{\dagger}\Theta$, the term
$\langle H_{0}\delta_{c}, TE\rangle$ equals $(\Theta(\ell Q)E)(c)$. Hence
\begin{equation}\label{5.8:eq-multiplier-comparison-5}
\begin{split}
(\mu'-\mu)(c) &= (\Theta(\ell Q)E)(c)
+ \bigl(d\mathcal{N}(\tilde{A}') - d\mathcal{N}(A')\bigr)[H_{0}\delta_{c}]\\
&\quad - \mathfrak{r}[\tilde{t}_{H_{0}\delta_{c}}].
\end{split}
\end{equation}

Off $\ker(\ell Q)$ the derivative $d\mathcal{N}[v]$ contains one further member,
$-\langle (\ell Q)v, \Theta D_{0}\rangle$. The direction $v = H_{0}\delta_{c}$ does not lie in
$\ker(\ell Q)$, so this member is present in the second term of
\eqref{5.8:eq-multiplier-comparison-5}, where it contributes
\begin{equation*}
-\bigl(\Theta(D_{0}(\tilde{A}') - D_{0}(A'))\bigr)(c).
\end{equation*}
Part~(a) of the chart in weighted norms (Lemma~\ref{5.8:lem-weighted-chart}) bounds the
difference $D_{0}(\tilde{A}') - D_{0}(A')$: its $\|\cdot\|_{w}$-norm is at most
$C\varepsilon_{3}\|E\|_{w}$.

The right-hand side of \eqref{5.8:eq-multiplier-comparison-5}, this member included,
is estimated with three inputs, the last term $\mathfrak{r}[\tilde{t}_{H_{0}\delta_{c}}]$
entering through the membership $\mathfrak{r}\in\sharp(C\mathsf{p}\mathsf{M})$ of
\eqref{5.8:eq-multiplier-comparison-2}:
\begin{itemize}
\item the bound \eqref{5.8:eq-critical-stability-a};
\item part~(iii) of the weighted calculus (Lemma~\ref{5.8:lem-weighted-calculus});
\item the bound $\|E\|_{w} \le C$ at the weight $w = \mathsf{p}\mathsf{M} + \mathsf{p}_{\delta}$,
established for the chart difference in the concluding step of the proof of the two-cutoff
comparison of responses (Theorem~\ref{5.8:thm-response-comparison}).
\end{itemize}
This gives \eqref{5.8:eq-multiplier-comparison-1}.
\end{proof}

\begin{definition}[The locally tempered majorant under refinement]\label{5.8:def-tempered-majorant}
Let $\mathsf{T}\ge 0$ be a profile on the label $\mathfrak{B}$, that is, a non-negative function on
$\mathfrak{B}$. Its \emph{locally tempered majorant} is the function
\begin{equation}\label{5.8:eq-tempered-majorant-1}
\mathsf{T}^{\wedge}(y) := \sup_{y'\in\mathfrak{B}}
\mathsf{T}(y')\,e^{-\frac{1}{2}\delta_{1}d^{\wedge}_{\mathfrak{B}}(y,y')},
\qquad y\in\mathfrak{B},
\end{equation}
where $d^{\wedge}_{\mathfrak{B}}$ is the distance $d^{\wedge}$ on $\mathfrak{B}$, as fixed in this
chapter; it satisfies $d^{\wedge}_{\mathfrak{B}}\le d$.

The statement uses the following objects, fixed earlier in this chapter.
\begin{itemize}
\item For $\sigma\ge 0$ and $A\ge 1$, $\mathfrak{M}^{\wedge,A}_{\sigma}(\mathfrak{B})$ is the class
$\mathfrak{M}^{A}_{\sigma}$ of Definition~\ref{5.8:def-tempered-weights},
with $d$ replaced by $d^{\wedge}_{\mathfrak{B}}$, and
$\mathfrak{M}^{\wedge}_{\sigma}(\mathfrak{B}) := \mathfrak{M}^{\wedge,1}_{\sigma}(\mathfrak{B})$.
\item The $\delta_{1}$-smear of a profile $f$ is
$f^{\sim}(y) = \sum_{y'} e^{-\delta_{1}d(y,y')}f(y')$.
\item For a refinement $\mathfrak{B}^{(1)}$ of $\mathfrak{B}$ and $y\in\mathfrak{B}^{(1)}$,
$\bar{y}\in\mathfrak{B}$
is the image of $y$ under the cell map of the refinement, in the sense fixed in this chapter, and
$d_{\mathfrak{B}^{(1)}}$ is the distance $d$ on $\mathfrak{B}^{(1)}$.
\end{itemize}
Then
\begin{equation}\label{5.8:eq-tempered-majorant-a}
\begin{aligned}
&\text{(t1)}\quad \mathsf{T}\le\mathsf{T}^{\wedge}\in\mathfrak{M}^{\wedge}_{\delta_{1}/2}(\mathfrak{B})
 \subset\mathfrak{M}_{\delta_{1}/2}(\mathfrak{B});\\
&\text{(t2)}\quad \mathsf{T}^{\sim}\le c_{1}\,\mathsf{T}^{\wedge};\\
&\text{(t3)}\quad \text{on every refinement $\mathfrak{B}^{(1)}$ of $\mathfrak{B}$, the function }
 y\mapsto\mathsf{T}^{\wedge}(\bar{y})\text{ lies in }
 \mathfrak{M}^{\wedge,2}_{\delta_{1}/2}(\mathfrak{B}^{(1)}),\\
&\phantom{\text{(t3)}\quad}\text{and its $\delta_{1}$-smear in $d_{\mathfrak{B}^{(1)}}$ is at most }
 2c_{1}\,\mathsf{T}^{\wedge}(\bar{y}),
\end{aligned}
\end{equation}
where $c_{1}$ is the constant of \eqref{5.8:eq-paired-sets-b}.
\end{definition}

\begin{proof}[Proof of \eqref{5.8:eq-tempered-majorant-a}]
(t1). Since $d^{\wedge}_{\mathfrak{B}}(y,y)=0$, the choice $y'=y$ in \eqref{5.8:eq-tempered-majorant-1}
gives $\mathsf{T}(y)\le\mathsf{T}^{\wedge}(y)$. Let $y,y',y''\in\mathfrak{B}$. The triangle inequality
$d^{\wedge}_{\mathfrak{B}}(y,y'')\ge d^{\wedge}_{\mathfrak{B}}(y',y'')-d^{\wedge}_{\mathfrak{B}}(y,y')$ gives
\begin{equation*}
\mathsf{T}(y'')\,e^{-\frac{1}{2}\delta_{1}d^{\wedge}_{\mathfrak{B}}(y,y'')}
\le e^{\frac{1}{2}\delta_{1}d^{\wedge}_{\mathfrak{B}}(y,y')}\,
\mathsf{T}(y'')\,e^{-\frac{1}{2}\delta_{1}d^{\wedge}_{\mathfrak{B}}(y',y'')}
\le e^{\frac{1}{2}\delta_{1}d^{\wedge}_{\mathfrak{B}}(y,y')}\,\mathsf{T}^{\wedge}(y').
\end{equation*}
Taking the supremum over $y''$ yields
$\mathsf{T}^{\wedge}(y)\le e^{\frac{1}{2}\delta_{1}d^{\wedge}_{\mathfrak{B}}(y,y')}\mathsf{T}^{\wedge}(y')$.
This is membership in $\mathfrak{M}^{\wedge}_{\delta_{1}/2}(\mathfrak{B})$. The inclusion in
$\mathfrak{M}_{\delta_{1}/2}(\mathfrak{B})$ holds because $d^{\wedge}_{\mathfrak{B}}\le d$, so that
$e^{\sigma d^{\wedge}_{\mathfrak{B}}}\le e^{\sigma d}$.

(t2). By \eqref{5.8:eq-tempered-majorant-1} and $d^{\wedge}_{\mathfrak{B}}\le d$, for all $y,y'$,
\begin{equation*}
\mathsf{T}(y')\le e^{\frac{1}{2}\delta_{1}d^{\wedge}_{\mathfrak{B}}(y,y')}\,\mathsf{T}^{\wedge}(y)
\le e^{\frac{1}{2}\delta_{1}d(y,y')}\,\mathsf{T}^{\wedge}(y).
\end{equation*}
Hence
\begin{equation*}
\mathsf{T}^{\sim}(y)=\sum_{y'}e^{-\delta_{1}d(y,y')}\mathsf{T}(y')
\le\mathsf{T}^{\wedge}(y)\sum_{y'}e^{-\frac{1}{2}\delta_{1}d(y,y')}
\le c_{1}\,\mathsf{T}^{\wedge}(y),
\end{equation*}
where the last step is \eqref{5.8:eq-paired-sets-b}.

(t3). Put $F(y):=\mathsf{T}^{\wedge}(\bar{y})$ for $y\in\mathfrak{B}^{(1)}$, and fix
$y,y'\in\mathfrak{B}^{(1)}$.
Applying (t1) at the points $\bar{y},\bar{y}'$ gives
\begin{equation*}
F(y)\le e^{\frac{1}{2}\delta_{1}d^{\wedge}_{\mathfrak{B}}(\bar{y},\bar{y}')}F(y').
\end{equation*}
By the lemma comparing the distance $d^{\wedge}$ on a label with the distance $d^{\wedge}$ on its
refinements, there is an additive constant, which we denote $d_{+}$, such that
\begin{equation*}
d^{\wedge}_{\mathfrak{B}}(\bar{y},\bar{y}')\le d^{\wedge}_{\mathfrak{B}^{(1)}}(y,y')+d_{+}.
\end{equation*}
Moreover $e^{d_{+}\delta}\le 2$ for the rate $\delta$ fixed in this chapter. Since
$\delta_{1}=\delta_{0}/32$ and $\delta=\delta_{0}/8$, we have $\delta_{1}/2\le\delta$, so that also
$e^{\frac{1}{2}\delta_{1}d_{+}}\le 2$. Combining these,
\begin{equation*}
F(y)\le 2\,e^{\frac{1}{2}\delta_{1}d^{\wedge}_{\mathfrak{B}^{(1)}}(y,y')}\,F(y'),
\end{equation*}
which is membership of $F$ in $\mathfrak{M}^{\wedge,2}_{\delta_{1}/2}(\mathfrak{B}^{(1)})$.

Exchanging $y$ and $y'$ and using $d^{\wedge}_{\mathfrak{B}^{(1)}}\le d_{\mathfrak{B}^{(1)}}$, we get
\begin{equation*}
F(y')\le 2\,e^{\frac{1}{2}\delta_{1}d_{\mathfrak{B}^{(1)}}(y,y')}F(y).
\end{equation*}
Therefore, by \eqref{5.8:eq-paired-sets-b} on $\mathfrak{B}^{(1)}$,
\begin{equation*}
\sum_{y'\in\mathfrak{B}^{(1)}}e^{-\delta_{1}d_{\mathfrak{B}^{(1)}}(y,y')}F(y')
\le 2F(y)\sum_{y'\in\mathfrak{B}^{(1)}}e^{-\frac{1}{2}\delta_{1}d_{\mathfrak{B}^{(1)}}(y,y')}
\le 2c_{1}\,\mathsf{T}^{\wedge}(\bar{y}).
\end{equation*}
By (t3) the majorant $\mathsf{T}^{\wedge}$ reproduces itself on every refinement up to the factor $2$,
which is what allows it to be carried through successive refinements in place of kernel-sum bounds on
smeared profiles.
\end{proof}

\begin{lemma}[Moves with the local second slot remain entries]\label{5.8:lem-local-slot}
Assume the following statements:
\begin{itemize}
\item[(a)] the two-cutoff comparison of responses at tempered profiles
(Theorem~\ref{5.8:thm-response-comparison}), the two-cutoff comparison of the multipliers
(Corollary~\ref{5.8:cor-multiplier-comparison}) and the increment bound for a move by a cut
response (Corollary~\ref{5.8:cor-cut-move});
\item[(b)] the Euler--Lagrange equation of the constrained critical point
$\mathfrak{X}(B;\mathbf{U},\mathsf{m})$, whose Lagrange multiplier is $\mu_B$;
\item[(c)] the identity \cite[(3.11)]{B-CMP109} for the current, displayed as
\eqref{5.8:eq-local-slot-3} below, and the first-variation identity
\eqref{5.8:eq-local-slot-5} for the action functional $\mathcal{W}$;
\item[(d)] the exact identity between the linearised averagings of the two machines and the
comparison identity for the averaging correction $C$ of the two machines;
\item[(e)] the definition of entries of a given modulus on a label, with its clauses
(e1)--(e3) and its level condition; the entry theorem, clauses (b) and (c), of which (b) bounds the
mismatch of the data vectors over an entry; and the hybrid lemma, read over entries;
\item[(f)] the pointwise summation bound for the profiles of the cut moves.
\end{itemize}
Let $p_0$ be an entry of modulus $(\mathcal{E}_0,C_e)$ on the label $\mathfrak{B}$; here and
below the first member of the modulus of an entry, a profile on the label, is written
$\mathcal{E}$. Let $r\ge 1$
and, for $s=1,\dots,r$, let the following be given:
\begin{enumerate}
\item[(1)] exactly paired admissible sets $(\mathfrak{S}'_s,\mathfrak{S}_s)$, in the sense of
Definition~\ref{5.8:not-paired-sets}, refining the label and satisfying the level condition of
the definition of entries;
\item[(2)] a real datum $b_s$ with $|b_s|\le\mathsf{p}_s$, where
$\mathsf{p}_s\in\mathfrak{M}_{4\delta_1}(\mathfrak{S}_s)$ and $10\sup\mathsf{p}_s\le a_R^{\sharp}$;
the datum $b_s$ is a single datum, prescribed in the frames of Ba{\l}aban's construction in both
machines and read there as in Corollary~\ref{5.8:cor-comparison-forms}(c);
\item[(3)] a real scalar $\tau_s$, the same in both machines $\mathsf{A}$ and $\mathsf{B}$, whose
support is disjoint from $\Lambda_0(\mathfrak{S}_s)\cup\Lambda'_0(\mathfrak{S}'_s)$; this is the
case for the cut-off scalars of the interpolating move along the line of sourced points, which
vanish on every box whose enlargement is not contained in the enlarged region carrying the
cut-offs, and for the cut-off scalars of the construction of the sourced step whose cut-off
scalars are supported in the interior of its region.
\end{enumerate}
Write $p_0=(P'_0,P_0)$, and for $s\ge1$ let $p_s=(P'_s,P_s)$ be the pair obtained from
$p_{s-1}$ by the $s$-th move defined now, in which $P_{s-1}$ carries the background field and
current $(\mathbf{U}_{s-1},\mathbf{J}_{s-1})$ of the machine $\mathsf{A}$ and $P'_{s-1}$ carries
those, $(\mathbf{U}'_{s-1},\mathbf{J}'_{s-1})$, of the machine $\mathsf{B}$. Put
$\mathbb{H}_s:=\mathbb{H}_{\mathfrak{S}_s}(b_s;P_{s-1})$, likewise $\mathbb{H}'_s$ in the machine
$\mathsf{B}$, and
\begin{equation}\label{5.8:eq-local-slot-1}
\begin{gathered}
\mathbf{U}_s:=e^{i\eta\tau_s\mathbb{H}_s}\mathbf{U}_{s-1},\qquad
\mathbf{J}_s:=\mathbf{J}_{s-1}+\tau_s\mathfrak{L}_s,\\
\mathfrak{L}_s:=\mathcal{J}_{\mathfrak{S}_s}(\mathbb{H}_s;\mathbf{U}_{s-1})
=(D^*D+DR_sD^*)\mathbb{H}_s+F(\mathbf{U}_{s-1},\mathbb{H}_s),
\end{gathered}
\end{equation}
with $R_s$ the operator $R$ of Lemma~\ref{5.8:lem-gauge-component} at $\mathfrak{S}_s$; the
objects $\mathbf{U}'_s$, $\mathbf{J}'_s$, $\mathfrak{L}'_s$ are defined by the same formulas in
the machine $\mathsf{B}$. The move \eqref{5.8:eq-local-slot-1} is the interpolating move
\begin{equation*}
\Phi_t:=\bigl(e^{i\eta\chi_t\mathbb{H}}\mathbb{U},\ \mathbb{J}+\chi_t\mathcal{J}^{\sharp}[1]\bigr),
\qquad \mathcal{J}^{\sharp}[1]=\mathfrak{s}(\mathbb{U},\mathbb{H})+DRD^*\mathbb{H},
\end{equation*}
along the line of sourced points $(\mathbb{U},\mathbb{J})$, taken at the identity vertex
$\mathbf{1}$ of the decoupling parameters and with $\tau=\chi_t$, the identity for
$\mathcal{J}^{\sharp}[1]$ being
clause (i) of the description of the sourced increment along that line; in a presented pair the
second slot is moved instead by $\mathcal{J}_{\mathfrak{S}_s}(\tau_s\mathbb{H}_s;\mathbf{U}_{s-1})$.

Let $\bar{\tau}_s$ be the size of $\tau_s$ defined in Corollary~\ref{5.8:cor-cut-move}, let
$C_s^{\sharp}$ be the constant of Lemma~\ref{5.8:lem-critical-size} (the subscript in
$C_s^{\sharp}$ is part of the name of that constant and does not refer to the index of the
move), and put
\begin{equation}\label{5.8:eq-local-slot-2}
\begin{gathered}
\mathsf{H}^{\sharp}_s:=C_s^{\sharp}\bar{\tau}_s\mathsf{p}_s,\qquad
\bar h:=\sup_y\sum_s\mathsf{H}^{\sharp}_s(y),\qquad \mathsf{i}:=L^{-\theta_2\lambda},\\
\mathsf{T}_{p_s}:=\mathcal{E}_0+\sum_{q\le s}\bigl|A_{p_q}-A_{p_{q-1}}\bigr|_{\cdot}
\quad(0\le s\le r),
\end{gathered}
\end{equation}
and let $\mathfrak{r}_e$ be the radius of the definition of entries. Since clause (e2) of the
definition of entries for $p_0$ gives $|A_{p_0}|\le\mathcal{E}_0$, the profile
$\mathsf{T}_{p_s}$ of \eqref{5.8:eq-local-slot-2} majorises the mismatch tally of the pair $p_s$
of Definition~\ref{5.8:not-paired-sets}, formed with $|A_{p_0}|$ in place of $\mathcal{E}_0$; in
this lemma and its proof $\mathsf{T}_{p_s}$ always denotes the profile
\eqref{5.8:eq-local-slot-2}. There are constants
$C_{sl}$, $C_{\star\star}$, depending on $(d,L,\theta,M_1,C_e)$ only and in particular free of
$r$, such that, with $\mathsf{T}^{\wedge}$ as in Definition~\ref{5.8:def-tempered-majorant} and
\begin{equation*}
\mathsf{M}_s:=C_{\star\star}\bigl[\mathsf{T}^{\wedge}_{p_s}+\mathsf{i}\bigr]\qquad(0\le s\le r),
\end{equation*}
under
\begin{equation}\label{5.8:eq-local-slot-b}
C_{sl}\bar h\le 1,\qquad \mathfrak{r}_e+C_{sl}\bar h\le\mathfrak{r}_0,\qquad
\sup\mathsf{T}_{p_r}\le\varepsilon_D,
\end{equation}
the following hold for every $s\le r$.
\begin{enumerate}
\item[(i)] $p_s$ is an entry of modulus $(\mathcal{E}_s,2C_e+1)$, where
\begin{equation*}
\mathcal{E}_s:=\mathsf{T}_{p_s}+C_{sl}\sum_{q\le s}\mathsf{H}^{\sharp}_q\mathsf{M}_{q-1}.
\end{equation*}
\item[(ii)] At every data bond $c$,
\begin{equation*}
|A_{p_s}|_c,\quad |\mathsf{m}_{P'_s}(c)-\mathsf{m}_{P_s}(c)|,\quad \mathsf{i}(c)
\ \le\ \mathsf{M}_s(c)=C_{\star\star}\bigl[\mathsf{T}^{\wedge}_{p_s}+\mathsf{i}\bigr](c).
\end{equation*}
\item[(iii)] Theorem~\ref{5.8:thm-response-comparison},
Corollary~\ref{5.8:cor-multiplier-comparison} and Corollary~\ref{5.8:cor-cut-move} hold at
$p_s$ with $\mathsf{M}:=\mathsf{M}_s$ and with the constant $A$ of the class
$\mathfrak{M}^{A}_{\tau}$ in Theorem~\ref{5.8:thm-response-comparison} equal to $2L$, and
clauses (i)--(iii)$_1$ of the hybrid
lemma hold with $\mathsf{T}_p$ replaced by $\mathcal{E}_s$.
\end{enumerate}
\end{lemma}

\begin{proof}
Steps~1--3 prove (i) at $s$ from (ii) at $s-1$; Step~4 proves (ii) and (iii) at $s$ from (i)
at $s$.

\emph{Step 1: in one machine, the slot increment is an image of the adjoint of the averaging,
by the Euler--Lagrange equation.} Fix $s$ and write $X:=\mathbb{H}_s$,
$\mathbf{U}:=\mathbf{U}_{s-1}$, $\mathbf{U}_1:=e^{i\eta X}\mathbf{U}$,
$\varsigma:=R_sD^*X$, $\mathsf{p}:=\mathsf{p}_s$ and $\mathfrak{L}:=\mathfrak{L}_s$. Let
$J(V)$ denote the current of a field $V$. By \cite[(3.11)]{B-CMP109}, for every field $h$,
\begin{equation}\label{5.8:eq-local-slot-3}
J(e^{i\eta h}\mathbf{U})-J(\mathbf{U})=D^*Dh+F(\mathbf{U},h).
\end{equation}
Taking $h=X$ and comparing with \eqref{5.8:eq-local-slot-1}, in which
$DR_sD^*X=D\varsigma$, we obtain
\begin{equation*}
\mathfrak{L}=J(\mathbf{U}_1)-J(\mathbf{U})+D\varsigma .
\end{equation*}
This identity is read weakly, the current being paired with a test field $\psi$ through
\begin{equation}\label{5.8:eq-local-slot-4}
\langle\psi,J(V)\rangle:=\langle D_V\psi,\mathfrak{F}_V\rangle ,
\end{equation}
where $\mathfrak{F}_V$ is the curvature of $V$. Let $\psi$ be a $C^1$ field vanishing on
$\Lambda_0$ and put $u:=\mathsf{g}_X^{-1}\psi$, where $\mathsf{g}_{\xi}$ is the linear map of
fields that enters the first-variation identity of the action functional,
\begin{equation}\label{5.8:eq-local-slot-5}
d\mathcal{W}_V(\xi)[z]=\langle\mathsf{g}_{\xi}z,J(e^{i\eta\xi}V)\rangle-\langle z,J(V)\rangle .
\end{equation}
Applying \eqref{5.8:eq-local-slot-5} with $V=\mathbf{U}$, $\xi=X$, $z=u$, and using
$\mathsf{g}_Xu=\psi$, we get
\begin{equation*}
\langle\psi,J(\mathbf{U}_1)\rangle=d\mathcal{W}_{\mathbf{U}}(X)[u]+\langle u,J(\mathbf{U})\rangle .
\end{equation*}
Since $\langle\psi,D\varsigma\rangle=\langle D^*\psi,\varsigma\rangle$ and
$D^*\psi=D^*u+D^*(\psi-u)$, it follows that
\begin{equation*}
\langle\psi,\mathfrak{L}\rangle
=d\mathcal{W}_{\mathbf{U}}(X)[u]+\langle D^*u,\varsigma\rangle
+\langle u-\psi,J(\mathbf{U})\rangle+\langle D^*(\psi-u),\varsigma\rangle .
\end{equation*}
Now $X$ is the constrained critical point with data vector $\mathsf{m}:=\mathsf{m}_{P_{s-1}}$;
let $\mu$ be its Lagrange multiplier. The Euler--Lagrange equation of
hypothesis~(b) states that the sum of the first two terms on the right equals
\begin{equation*}
\langle d\mathbf{Q}(\eta X)u,\mu\rangle+\langle dC(X)u,\mathsf{m}\rangle ,
\end{equation*}
where $dC(X)$ is the derivative of the averaging correction $C$ of
Lemma~\ref{5.8:lem-averaging-correction}. Hence
\begin{equation*}
\langle\psi,\mathfrak{L}\rangle
=\langle d\mathbf{Q}(\eta X)u,\mu\rangle+\langle dC(X)u,\mathsf{m}\rangle
+\langle u-\psi,J(\mathbf{U})\rangle+\langle D^*(\psi-u),\varsigma\rangle .
\end{equation*}
We replace $u$ by $\psi$ in the first two terms and collect what remains in
\begin{equation*}
\begin{split}
\epsilon[\psi]:={}&\langle d\mathbf{Q}(\eta X)(u-\psi),\mu\rangle
+\langle dC(X)(u-\psi),\mathsf{m}\rangle\\
&+\langle u-\psi,J(\mathbf{U})\rangle+\langle D^*(\psi-u),\varsigma\rangle .
\end{split}
\end{equation*}
The difference $u-\psi$ obeys
\begin{equation*}
|u-\psi|\le C\eta|X||\psi|,\qquad
|\nabla(u-\psi)|\le C\eta\bigl(|\nabla X||\psi|+|X||\nabla\psi|\bigr).
\end{equation*}
At a site $y$ the factor $\eta$ in these bounds equals $L^{-\lambda(y)}\ell_y$, since
$\ell=L^{\lambda}\eta$. By Lemma~\ref{5.8:lem-critical-size}, $\ell|X|$ and $\ell^2|\nabla X|$
are at most $N^{\Delta}_y(X)\le C_s^{\sharp}\mathsf{p}(y)$, and the multiplier obeys
$|\mu|\le C_s^{\sharp}\mathsf{p}$ at the data bonds containing $y$; by
Lemma~\ref{5.8:lem-gauge-component}, $\ell^2|\varsigma|\le
C_{\varsigma}^{\sharp}\mathsf{p}(y)(\mathfrak{r}_0+\sup\mathsf{p})$; the kernel of $dC(X)$ is
bounded by the kernel bound $(\ell2)$ of \eqref{5.8:eq-averaging-correction-a}; and the current
$J(\mathbf{U})$ and the
data vector $\mathsf{m}$ are bounded through the regularity class $\mathsf{C}_{\beta}(\mathfrak{r}_0)$.
Inserting these bounds into the four terms of $\epsilon[\psi]$, we obtain
\begin{equation}\label{5.8:eq-local-slot-a}
\begin{gathered}
\langle\psi,\mathfrak{L}\rangle=\langle d\mathbf{Q}(\eta X)\psi,\mu\rangle
+\langle dC(X)\psi,\mathsf{m}\rangle+\epsilon[\psi],\\
|\epsilon[\psi]|\le C\sum_yL^{-\lambda(y)}N_y(\psi)\mathsf{p}(y).
\end{gathered}
\end{equation}
No non-local operator acts on $\psi$ in \eqref{5.8:eq-local-slot-a}; consequently, for
$\psi=\tau\phi$ with a scalar $\tau$, the amplitude of $\psi$ meets the profile $\mathsf{p}$ at
the data sites of the support of $\tau$. Moreover, $\ell^3|\mathfrak{L}|\le C\mathsf{p}$
pointwise: this follows from the Laplace member $\ell^3|\Delta_{\mathbf{U}}X|$ of the bound
$N^{\Delta}_y(X)\le C_s^{\sharp}\mathsf{p}(y)$ of Lemma~\ref{5.8:lem-critical-size}, from
\cite[(3.49)]{B-CMP99b}, and from the pointwise bound on the slot increment of presented pairs.

\emph{Step 2: two machines, one stage.} Keep $s$ fixed and the notation of Step~1 in the
machine $\mathsf{A}$; in the machine $\mathsf{B}$ write $X':=\mathbb{H}'_s$,
$\mathbf{U}':=\mathbf{U}'_{s-1}$, $\mu'$ for the multiplier of $X'$,
$\mathsf{m}':=\mathsf{m}_{P'_{s-1}}$ and $\mathfrak{L}':=\mathfrak{L}'_s$, and write
$\tau:=\tau_s$. Let $\phi'$ be a test field of the machine $\mathsf{B}$, and let $w$ be a weight
bounding the five local sizes of $\phi'$ that enter clause (e3) of the definition of entries.
Let $Q_1$ denote the averaging map from fields of the machine $\mathsf{B}$ to fields of the
machine $\mathsf{A}$, taken at the background $\mathbf{U}'_{s-1}$, and put
$\psi':=\tau\phi'$, $\bar\psi:=Q_1\psi'$. Then $N_y(\psi')$, $N_y(\bar\psi)\le
C\bar{\tau}(y)w(y)$ at every $y$.

Let $W:=M(\mathbf{U}')=e^{i\eta\mathfrak{a}}\mathbf{U}$ be as in
Theorem~\ref{5.8:thm-response-comparison}, let $Y:=Y[X';\mathbf{U}']$, and write
$\mathbf{Q}^{\mathrm{lin}}_V$ for the linear averaging at the background $V$, with
$\mathbf{Q}'^{\mathrm{lin}}_{\mathbf{U}'}$ the one of the machine $\mathsf{B}$; the first terms
of \eqref{5.8:eq-local-slot-a} in the two machines are compared through these maps. By the exact
identity of hypothesis~(d),
\begin{equation*}
\mathbf{Q}'^{\mathrm{lin}}_{\mathbf{U}'}\psi'=\mathbf{Q}^{\mathrm{lin}}_W\bar\psi ,
\qquad
\bigl|(\mathbf{Q}^{\mathrm{lin}}_W-\mathbf{Q}^{\mathrm{lin}}_{\mathbf{U}})\bar\psi\bigr|(c)
\le C|\mathfrak{a}|_c|\bar\psi|_c .
\end{equation*}
Write $dC_V$ for the derivative of the averaging correction at the background $V$ and $dC'$ for
that of the machine $\mathsf{B}$. Differentiating the comparison identity of hypothesis~(d) at
$X'$ in the direction $\psi'$ shows that $dC'(X')\psi'$ is the sum of
\begin{equation*}
dC_W(Y)\bigl[dY[X']\psi'\bigr]
+\mathbf{Q}^{\mathrm{lin}}_W\bigl(dY[X']\psi'-Q_1\psi'\bigr)
\end{equation*}
and of the terms obtained by differentiating the remaining member of the comparison identity of
hypothesis~(d) along $\psi'$,
where $dY[X']\psi'=\bar\psi+O(\eta|X'||\psi'|)$. By the Schwarz inequality and the Lipschitz
bound $(\ell3)$ for the kernel in \eqref{5.8:eq-averaging-correction-a},
\begin{equation*}
\bigl|(dC_W(Y)-dC_{\mathbf{U}}(X))\bar\psi\bigr|(c)
\le C\bigl(|\mathfrak{a}|_c|Y|_c+|Y-X|_c\bigr)|\bar\psi|_c .
\end{equation*}

We now insert \eqref{5.8:eq-local-slot-a} in both machines, with $\psi'$ in $\mathsf{B}$ and
$\bar\psi$ in $\mathsf{A}$, and compare term by term. The smallness condition
$\sup\mathsf{T}_{p_{s-1}}\le\varepsilon_D$ holds by the third condition of
\eqref{5.8:eq-local-slot-b}, since $\mathsf{T}_{p_q}$ is non-decreasing in $q$. By (ii) at
$s-1$, the hypotheses of Theorem~\ref{5.8:thm-response-comparison} hold with
$\mathsf{M}:=\mathsf{M}_{s-1}$; that theorem gives $|Y-X|_c\le C\mathsf{p}\mathsf{M}_{s-1}(c)$,
the profile $\mathsf{p}_{\delta}$ of the two-cutoff difference of the datum being supplied by
Corollary~\ref{5.8:cor-comparison-forms}(c), since $b_s$ is prescribed in the frames of
Ba{\l}aban's construction in both machines. Corollary~\ref{5.8:cor-multiplier-comparison} gives
$|\mu'-\mu|\le C\mathsf{p}\mathsf{M}_{s-1}$, and Lemma~\ref{5.8:lem-critical-size} gives
$|\mu|,|\mu'|\le C\mathsf{p}$. Combining these with the three bounds displayed above and with the
error bounds of \eqref{5.8:eq-local-slot-a} in both machines, we obtain
\begin{equation*}
\bigl|\langle\psi',\mathfrak{L}'\rangle_{\eta'}-\langle\bar\psi,\mathfrak{L}\rangle_{\eta}\bigr|
\le C\sum_c\bar{\tau}w\mathsf{p}\mathsf{M}_{s-1}(c),
\end{equation*}
where the subscripts indicate the lattice spacings of the pairings. Finally,
$\langle\phi',\tau\mathfrak{L}'\rangle=\langle\psi',\mathfrak{L}'\rangle$ and
\begin{equation*}
\langle Q_1\phi',\tau\mathfrak{L}\rangle
=\langle\bar\psi,\mathfrak{L}\rangle+\langle\tau Q_1\phi'-Q_1(\tau\phi'),\mathfrak{L}\rangle ,
\end{equation*}
and the last term is controlled by the commutator bound
$|\tau Q_1\phi'-Q_1(\tau\phi')|\le C\eta|\nabla\tau||\phi'|$ together with
$\ell^3|\mathfrak{L}|\le C\mathsf{p}$ from Step~1. Recalling
$\mathsf{H}^{\sharp}_s=C_s^{\sharp}\bar{\tau}_s\mathsf{p}_s$, we conclude
\begin{equation}\label{5.8:eq-local-slot-6}
\bigl|\langle\phi',\tau_s\mathfrak{L}'_s\rangle-\langle Q_1\phi',\tau_s\mathfrak{L}_s\rangle\bigr|
\le C_{sl}\sum_yw(y)\mathsf{H}^{\sharp}_s(y)\mathsf{M}_{s-1}(y).
\end{equation}

\emph{Step 3: item (i).} Write $Q_1^{(q)}$ for the map $Q_1$ taken at the background
$\mathbf{U}'_q$, so that Step~2 at the stage $q$ concerns $Q_1^{(q-1)}$. Since
$\mathbf{J}_q-\mathbf{J}_{q-1}=\tau_q\mathfrak{L}_q$ and
$\mathbf{J}'_q-\mathbf{J}'_{q-1}=\tau_q\mathfrak{L}'_q$, telescoping gives
\begin{equation*}
\begin{split}
\langle\phi',\mathbf{J}'_s\rangle-\langle Q_1^{(s)}\phi',\mathbf{J}_s\rangle
={}&\bigl\{\langle\phi',\mathbf{J}'_0\rangle-\langle Q_1^{(0)}\phi',\mathbf{J}_0\rangle\bigr\}\\
&+\sum_{q\le s}\bigl\{\langle\phi',\tau_q\mathfrak{L}'_q\rangle
-\langle Q_1^{(q-1)}\phi',\tau_q\mathfrak{L}_q\rangle\bigr\}\\
&-\sum_{q\le s}\bigl\langle\bigl(Q_1^{(q)}-Q_1^{(q-1)}\bigr)\phi',\mathbf{J}_q\bigr\rangle ,
\end{split}
\end{equation*}
the last sum being the drift of $Q_1$ along the moves. The first brace is bounded by clause
(e3) for $p_0$, the five local sizes of $\phi'$ at $\mathbf{U}'_0$ being at most
$(1+C\mathfrak{r}_0)w$; the drift is at most
$C\sum_yw(y)\mathfrak{r}_0^2L^{-\lambda(y)}$; both bounds are obtained as in the proof of the
entry theorem. Each term of the middle sum is bounded by \eqref{5.8:eq-local-slot-6} at the
stage $q$, which uses (ii) at $q-1\le s-1$; since $w\le w^{\sim}$, all three bounds may be
written with $w^{\sim}$ on the right, as clause (e3) requires. Together these bounds give
clause (e3) for $p_s$ at the modulus $(\mathcal{E}_s,2C_e+1)$.

Clause (e1) for the fields follows from the bound \eqref{5.8:eq-cut-move-b} of
Corollary~\ref{5.8:cor-cut-move} applied to each move. For the currents,
$\ell^3|\mathbf{J}_s|\le\ell^3|\mathbf{J}_0|+\sum_{q\le s}|\tau_q|\,\ell^3|\mathfrak{L}_q|$, and
by (e1) for $p_0$, $|\tau_q|\le\bar{\tau}_q$ and $\ell^3|\mathfrak{L}_q|\le C\mathsf{p}_q$ from
Step~1,
\begin{equation*}
\ell^3|\mathbf{J}_s|\le\mathfrak{r}_e+C\sum_{q\le s}\mathsf{H}^{\sharp}_q ,
\end{equation*}
which is at most $\mathfrak{r}_0$ by the second condition of \eqref{5.8:eq-local-slot-b} once
$C_{sl}$ is at least the constant $C$ here. For clause (e2): by (e2) for $p_0$, which bounds
$|A_{p_0}|$ by $\mathcal{E}_0$, and by the triangle inequality,
\begin{equation*}
|A_{p_s}|\le|A_{p_0}|+\sum_{q\le s}|A_{p_q}-A_{p_{q-1}}|\le\mathsf{T}_{p_s}\le\mathcal{E}_s .
\end{equation*}
This proves (i) at $s$.

\emph{Step 4: item (ii), by induction and without compounding of constants; item (iii).}
Write $\delta\mathsf{m}_s:=\mathsf{m}_{P'_s}-\mathsf{m}_{P_s}$. Clause (b) of the entry
theorem, applied over the entry $p_s$ with empty presentation, gives
\begin{equation*}
|\delta\mathsf{m}_s(c)|\le C_{\mathsf{m}}\bigl[\mathcal{E}_s^{\sim}
+\mathfrak{r}_0L^{-\theta\lambda}\bigr](c),\qquad C_{\mathsf{m}}=C_{\mathsf{m}}(2C_e+1).
\end{equation*}
At $s=0$ this applies to $p_0$ itself, with $\mathcal{E}_0$ as in (i); this starts the induction.
The sum defining $\mathsf{T}_{p_q}$ has non-negative terms, so $\mathsf{T}_{p_q}$ is
non-decreasing in $q$; since $\mathsf{T}^{\wedge}$ is a supremum of $\mathsf{T}$ against a fixed
kernel, $\mathsf{T}^{\wedge}_{p_q}$ is non-decreasing in $q$ as well, and therefore, for $q\le s$,
\begin{equation*}
\mathsf{M}_{q-1}\le C_{\star\star}\bigl[\mathsf{T}^{\wedge}_{p_s}+\mathsf{i}\bigr].
\end{equation*}
With $\mathsf{T}_{p_s}\le\mathsf{T}^{\wedge}_{p_s}$ (property (t1) of
\eqref{5.8:eq-tempered-majorant-a}) and $\sum_q\mathsf{H}^{\sharp}_q\le\bar h$, this gives
\begin{equation*}
\mathcal{E}_s\le\bigl(1+C_{sl}\bar hC_{\star\star}\bigr)\bigl[\mathsf{T}^{\wedge}_{p_s}
+\mathsf{i}\bigr].
\end{equation*}
By property (t3) of \eqref{5.8:eq-tempered-majorant-a}, which bounds the $\delta_1$-smear of
$\mathsf{T}^{\wedge}$ on every refinement of the label by $2c_1\mathsf{T}^{\wedge}$, and by
clause (w3) of the calculus of tempered weights, the bound for the $\delta_1$-smear of the level
weight $\mathsf{i}$,
\begin{equation*}
\mathcal{E}_s^{\sim}\le 2Lc_1\bigl(1+C_{sl}\bar hC_{\star\star}\bigr)
\bigl[\mathsf{T}^{\wedge}_{p_s}+\mathsf{i}\bigr].
\end{equation*}
Put $C_{\star\star}:=8LC_{\mathsf{m}}c_1+C_{\mathsf{i}}+1$, where $C_{\mathsf{i}}$ is a constant
depending on $(d,L,\theta,M_1,C_e)$ only, the member of $C_{\star\star}$ by which the term
$C_{\mathsf{m}}\mathfrak{r}_0L^{-\theta\lambda}$ of the first display of this step is carried
into the member $C_{\star\star}\mathsf{i}$ of the bound: it is taken so large that
\begin{equation*}
C_{\mathsf{m}}\mathfrak{r}_0L^{-\theta\lambda}\le\tfrac34C_{\mathsf{i}}\,\mathsf{i}
\quad\text{on the label};
\end{equation*}
and read the first condition of
\eqref{5.8:eq-local-slot-b} as $8LC_{\mathsf{m}}c_1C_{sl}\bar h\le1$. Then the last two displays
give
\begin{equation*}
|\delta\mathsf{m}_s|\le C_{\star\star}\bigl[\mathsf{T}^{\wedge}_{p_s}+\mathsf{i}\bigr]
=\mathsf{M}_s ,
\end{equation*}
with the same constant $C_{\star\star}$ at every $s$, so that no constant compounds along the
moves. The bound $\mathsf{i}\le\mathsf{M}_s$ holds because $C_{\star\star}\ge1$ and
$\mathsf{T}^{\wedge}\ge0$; the bound $|A_{p_s}|_c\le\mathsf{M}_s(c)$ follows from
$|A_{p_s}|\le\mathsf{T}_{p_s}$ of Step~3, from (t1) and from $C_{\star\star}\ge1$. This proves
(ii) at $s$.

For (iii): Theorem~\ref{5.8:thm-response-comparison} is stated in free form, its hypotheses
concerning only the fields, the data and the majorant $\mathsf{M}$, so that it applies at $p_s$
with the majorant $\mathsf{M}_s$ of (ii); Corollaries~\ref{5.8:cor-multiplier-comparison}
and~\ref{5.8:cor-cut-move} are then available in the same setting. The proof of the hybrid lemma
reads the pair $p$ only through its fields, for its clause (i), and, for its clause (ii), only
through the fact that $p$ is then itself an entry, to which clause (c) of the entry theorem is
applied with the slots of $p$ itself. By (i) of the present lemma, $p_s$ is an entry of modulus
$(\mathcal{E}_s,2C_e+1)$, so clauses (i)--(iii)$_1$ of the hybrid lemma hold with $\mathsf{T}_p$
replaced by $\mathcal{E}_s$.
\end{proof}

\begin{remark}
(r1) The number $r$ of moves enters nowhere: the sum $\sum_q\mathsf{H}^{\sharp}_q$ is bounded
pointwise by the pointwise summation bound of hypothesis~(f), so that
\eqref{5.8:eq-local-slot-b} is a smallness condition on $\gamma$.

(r2) With the smeared datum majorant $\mathsf{P}_b$ in place of $\mathsf{p}$ the lemma would be
void on the product of polydiscs $\mathfrak{P}$ of the amplitudes, where
$\mathsf{H}=\bar{\tau}\mathsf{P}_b$ is unbounded.

(r3) The identity \eqref{5.8:eq-local-slot-a} is the closed form of the source term obtained
for the sourced step, bare sources plus adjoints of the averaging acting on block fields, read
on the second slot.

(r4) The slot rule of the construction of the sourced step whose cut-off scalars are supported
in the interior of its region,
$\tau\mathfrak{L}_Z+[D^*D,\tau]\mathbb{H}_Z+F(\mathbf{U},\tau\mathbb{H}_Z)$ with $Z$ the region
of that construction, differs from the rule \eqref{5.8:eq-local-slot-1} by the term
$-\tau F(\mathbf{U},\mathbb{H})$ and by the two local first-order members
$[D^*D,\tau]\mathbb{H}$ and $F(\mathbf{U},\tau\mathbb{H})$. Weakly these are
\begin{equation*}
\langle D\phi,D(\tau\mathbb{H})\rangle-\langle D(\tau\phi),D\mathbb{H}\rangle
\quad\text{and}\quad
\langle D_{\mathbf{U}_{\tau}}\phi,\mathfrak{F}_{\mathbf{U}_{\tau}}\rangle
-\langle D_{\mathbf{U}}\phi,\mathfrak{F}_{\mathbf{U}}\rangle
-\langle D\phi,D(\tau\mathbb{H})\rangle ,
\end{equation*}
with $\mathbf{U}_{\tau}$ the field moved by $\tau\mathbb{H}$; they are products of $C^{\beta}$
fields, which would nest by the two nesting statements for products of $C^{\beta}$ fields and
the nesting statement (b) of the comparison of the two machines, whose clause (c) is the exact
identity of hypothesis~(d), in the way this is done for the first clause of the lemma on
presented pairs. This nesting is not carried out here: for that slot rule the
statement of Lemma~\ref{5.8:lem-local-slot} is not proved.
\end{remark}

\begin{corollary}[Natural size of the cut response on the contour]\label{5.8:cor-contour-size}
Assume the following.
\begin{itemize}
\item The increment bound for a move by a cut response, Corollary~\ref{5.8:cor-cut-move}, in
particular its natural unit $\mathsf{H}^{\sharp} = C_{s}^{\sharp}\bar{\tau}\mathsf{p}$.
\item The radius setting of the rate bound for the cut response. The datum $b$ lives on the live
sites of the region $R_{\nu}$, with $\sup|b| \le \mathsf{b}_{j}$, where $\mathsf{b}_{j}$ is the
bound on the datum. The radius is $r_{\square} = c_{M}/\vartheta_{P}(\square)$, with $c_{M}$ a
constant and with the radius parameter $\vartheta_{P}(\square)$ of the box $\square$ given, in
terms of its reference value $\vartheta_{P}^{0}$ and its rate $\delta_{P}$, by
\begin{equation*}
\vartheta_{P}(\square) = \vartheta_{P}^{0}\,e^{-\delta_{P} d_{sc}(\square,R_{\nu})},
\qquad \delta_{P} = \tfrac12\delta = 2\delta_{1},
\end{equation*}
where $d_{sc}$ is the distance described next.
\item In that setting $\delta$ denotes $\tfrac18\delta_{0}$, and $d_{sc}$ denotes the distance
$d^{\wedge}$ of $\mathfrak{B}$; it satisfies $d_{sc} \le d$, where
$d$ is the distance in which $\mathsf{p}$ below and $\operatorname{diam}_{d}$ are measured.
\end{itemize}
Put $\mathsf{p} := \mathsf{b}_{j}\,e^{-4\delta_{1} d(\cdot,R_{\nu})} \in \mathfrak{M}_{4\delta_{1}}$.
Then, for $|\tau| \le r_{\square}$ on the enlarged box $\tilde{\square} \supseteq \square$ and every
$x \in \tilde{\square}$,
\begin{equation}\label{5.8:eq-contour-size-1}
\mathsf{H}^{\sharp}(x) \le C_{s}^{\sharp}\,e^{2\delta_{1}\operatorname{diam}_{d}\tilde{\square}}
\,\frac{c_{M}\mathsf{b}_{j}}{\vartheta_{P}^{0}}\,e^{-2\delta_{1} d(x,R_{\nu})}.
\end{equation}
The bound \eqref{5.8:eq-contour-size-1} holds on the whole closed disc $|\tau| \le r_{\square}$,
in particular on the contour $|\tau| = r_{\square}$. Neither of the two
alternative conditions under which the rate bound for the cut response is established is
assumed.
\end{corollary}

\begin{proof}
The two expressions for $\delta_{P}$ agree. Since $\delta = \tfrac18\delta_{0}$ and
$\delta_{1} = \tfrac{1}{32}\delta_{0}$, we have $\tfrac12\delta = \tfrac1{16}\delta_{0} = 2\delta_{1}$.
Hence $\delta_{P} = 2\delta_{1}$.

By Corollary~\ref{5.8:cor-cut-move}, $\mathsf{H}^{\sharp}(x) = C_{s}^{\sharp}\bar{\tau}(x)\mathsf{p}(x)$.
The hypothesis $|\tau| \le r_{\square}$ on $\tilde{\square}$ enters through the size $\bar{\tau}$ of
Corollary~\ref{5.8:cor-cut-move} and is used in the form $\bar{\tau}(x) \le r_{\square}$ for
$x \in \tilde{\square}$; when $\tau$ is constant on $\tilde{\Delta}(x)$ this is the hypothesis
itself, since then $\nabla\tau$ and $\nabla^{2}\tau$ vanish there and $\bar{\tau}(x) = |\tau|$.
Insert
\begin{equation*}
r_{\square} = \frac{c_{M}}{\vartheta_{P}^{0}}\,e^{2\delta_{1} d_{sc}(\square,R_{\nu})}
\end{equation*}
and the definition of $\mathsf{p}$. This gives
\begin{equation*}
\mathsf{H}^{\sharp}(x) \le C_{s}^{\sharp}\,\frac{c_{M}\mathsf{b}_{j}}{\vartheta_{P}^{0}}\,
\exp\bigl(2\delta_{1} d_{sc}(\square,R_{\nu}) - 4\delta_{1} d(x,R_{\nu})\bigr).
\end{equation*}
So \eqref{5.8:eq-contour-size-1} follows from the inequality between exponents
\begin{equation}\label{5.8:eq-contour-size-2}
\begin{split}
2\delta_{1} d_{sc}(\square,R_{\nu}) - 4\delta_{1} d(x,R_{\nu})
&\le 2\delta_{1}\bigl[d(x,R_{\nu}) + \operatorname{diam}_{d}\tilde{\square}\bigr]
- 4\delta_{1} d(x,R_{\nu})\\
&= 2\delta_{1}\operatorname{diam}_{d}\tilde{\square} - 2\delta_{1} d(x,R_{\nu}).
\end{split}
\end{equation}
The first line of \eqref{5.8:eq-contour-size-2} holds because $d_{sc} \le d$. It also uses, for
$x \in \tilde{\square}$, the triangle inequality
\begin{equation*}
d(\square,R_{\nu}) \le d(x,R_{\nu}) + \operatorname{diam}_{d}\tilde{\square},
\end{equation*}
which applies since $\square \subset \tilde{\square}$. The second line is algebra.
\end{proof}

In particular, the display of the rate bound for the cut response, with decay rate
$\tfrac38\delta_{1}$, holds a fortiori for $|\tau| \le r_{\square}$, since
$\tfrac38\delta_{1} \le 2\delta_{1}$. Consequently the decay rate of that display requires neither
the pair radii $r^{\bullet}_{\square}$ nor the condition
$\tfrac18\delta_{1}M \ge (1+3\beta_{I})\kappa + c_{d}$ of that setting, in which $\beta_{I}$ and
$c_{d}$ are constants of that setting.

\begin{definition}[The far datum and its smallness ceiling]\label{5.8:def-far-datum}
We use the two cutoffs $\mathsf{A}$, at spacing $\eta$, and $\mathsf{B}$, at spacing $\eta'=\eta/L$, on
exactly paired admissible sets, in the notation for two cutoffs on exactly paired sets fixed in
Notation~\ref{5.8:not-paired-sets}.

The far datum $\mathsf{E}$ is defined below at a point $(\mathbb{U},\mathbb{J})$, with window field
$\mathbb{W}$ and layers $\mathbb{H}^{w}$, $\mathbf{H}_{j}(t,\mathsf{t})$. The identifications listed
below express these objects, and the two-cutoff mismatches $\delta\mathsf{E}$, $\delta^{0}\mathsf{E}$
of $\mathsf{E}$, in the notation of the two-cutoff comparison. Each identification is stated
together with the statement on which it rests.

We recall the far datum from the statement in which it is constructed.
Here $t\in[0,1]$, the complex amplitude is $\mathsf{t}$, and
\begin{align*}
D^{\sharp}&:=\zeta'\,Q_{\star}\bigl((t+\mathsf{t}\zeta_{\square})\mathbb{H}^{w}\bigr)
+(1-\zeta')\,Q_{\star}\bigl(\mathbf{H}_{j}(t,\mathsf{t})\bigr),\\
\tilde{\mathsf{B}}&:=Q_{\star}\bigl(\mathbf{A}(t,\mathsf{t})\bigr),\qquad
\mathbf{A}(t,\mathsf{t}):=(t\zeta'_{\square}+\mathsf{t}\zeta_{\square})\mathbb{H}^{w},\\
\mathsf{B}(t,\mathsf{t})&:=Q_{\mathfrak{B}_{j}}\bigl(\mathbb{W};\eta(t+\mathsf{t}\zeta_{\square})
\mathbb{H}^{w}\bigr),\\
\mathsf{E}&:=\frac{1}{i}\log\bigl(e^{iD^{\sharp}}e^{-i\tilde{\mathsf{B}}}\bigr),
\end{align*}
where $Q_{\mathfrak{B}_{j}}(\mathbb{W};\cdot)$ is the averaging map of that construction on the cap
$\mathfrak{B}_{j}$ in the background field $\mathbb{W}$; its arguments are written in this order
throughout, the background field first and the one-form being averaged second, so that
$Q_{\mathfrak{B}_{j}}(\mathbb{W};u)$ is the average of the one-form $u$ in the background
$\mathbb{W}$, and in $\mathsf{B}(t,\mathsf{t})$ the factor $\eta$ is part of the averaged one-form.
Further, $\mathbf{H}_{j}(t,\mathsf{t})$ is
the response on the cap $\mathfrak{B}_{j}$ provided by the response theorem (the third of the
responses it constructs), at the window field $\mathbb{W}$, the current $\mathbb{J}$, the vertex
$\sigma$ and the datum $(0,\mathsf{B}(t,\mathsf{t}))$.
The map $Q_{\star}$ is the averaging map of that construction taken at $(\mathfrak{B};\mathbb{W},
\mathbb{J})$, the same triple at which $\mathbb{H}^{w}$ is taken; $\zeta'$, $\zeta_{\square}$ and
$\zeta'_{\square}$ are the cutoff functions of that construction, carried by the box $\square$.
For a bond $c$, $\tilde{B}(c)$ is the neighbourhood formed of the blocks of $c$, and
$\tilde{\square}^{3}$ is the neighbourhood of $\square$, as fixed in that construction; we put
\[
F_{\square}:=\{c:\ \text{some block of } c \text{ meets the complement of } \tilde{\square}^{3}\}.
\]
In that construction $\mathsf{E}(c)=0$ whenever $\tilde{B}(c)\subset\tilde{\square}^{3}$, and
$F_{\square}$ contains the support of $\mathsf{E}$.
The datum $\mathsf{B}(t,\mathsf{t})$ on the cap always carries its arguments $(t,\mathsf{t})$, and
every decorated form of the letter (with arguments or a tilde) denotes a datum; the
undecorated letter $\mathsf{B}$ denotes only the finer cutoff, in the phrase ``the cutoff
$\mathsf{B}$'' and as a superscript, $\mathsf{x}^{\mathsf{B}}$.

\begin{enumerate}
\item \emph{Sets.} The admissible set $\mathfrak{S}$ is the label $\mathfrak{B}$, which is the set of
level $k$ on which the far datum is constructed.
Its cap $\mathfrak{B}_{j}$ is an initial segment of the admissible sequence of $\mathfrak{B}$, and is
therefore admissible, by the statement on caps and their profiles. The set $\mathfrak{B}^{+}$ is the
one used for the window. These sets are exactly paired with the corresponding sets of the cutoff
$\mathsf{B}$, in the sense of the notation for two cutoffs on exactly paired sets fixed in
Notation~\ref{5.8:not-paired-sets}, by the pairing statement. The caps of paired labels are again
paired, by the statement on the pairing of caps.

\item \emph{Distance and rates.} The distance $d$ is one and the same on both sides of the
correspondence. It is the distance of \cite[(2.46)]{B-CMP96}, as fixed in
Notation~\ref{5.8:not-paired-sets}. The decay rate is
\[
\delta=4\delta_{1}.
\]

\item \emph{Point.} The point $(\mathbb{U},\mathbb{J})$ corresponds to the point
$P=(\mathbf{U},\mathbf{J})$ formed by the background field $\mathbf{U}$ and the current $\mathbf{J}$.

\item \emph{Window.} The window field $\mathbb{W}$ corresponds to the hybrid-gauge representative
$\mathbf{U}^{\natural}$ of $\mathbf{U}$, at the box $\square_{0}$, with shells over
$\mathfrak{B}^{+}$. This correspondence holds under the
hybrid representation. The axial representative of Ba{\l}aban's construction differs from
$\mathbf{U}^{\natural}$ by a rotation; this rotation is fixed in the statement on the rotation of
the hybrid-gauge representative.

\item \emph{Layers.} Fix the identity vertex $\sigma\equiv\mathbf{1}$, and take the response on each
side by Ba{\l}aban's defining formula for it, not by its representation. Under these, the layers
correspond as follows:
\begin{align*}
\mathbb{H}^{w}&\ \longleftrightarrow\ \mathbb{H}_{\mathfrak{B}}(b;\mathbf{U}^{\natural},\mathbf{J}),\\
\mathbf{H}_{j}(t,\mathsf{t})&\ \longleftrightarrow\
\mathbb{H}_{\mathfrak{B}_{j}}\bigl(\mathsf{B}(t,\mathsf{t});\mathbf{U}^{\natural},\mathbf{J}\bigr),
\end{align*}
where $b$ is the datum on the label $\mathfrak{B}$.
At a real critical point with its own current, both sides of each correspondence are Ba{\l}aban's Landau function. This is clause~(d) of the response theorem, together with clause~(c) of the naturality statement under tilting. Off shell, the correspondence is the one given by the off-shell pair extension.

\item \emph{Gauge.} We write $\psi$ for the gauge function of the hybrid gauge; the gauge
transformation is $g:=e^{i\psi}$, the one of the hybrid gauge. It is not the
frame gauge. The frame gauge enters separately, in the estimates that follow, through a term
written $\delta b_{0}$ at the place where it is introduced.

\item \emph{Mismatches.} Let $\mathsf{x}$ be a quantity defined at the bonds for both cutoffs, with values $\mathsf{x}^{\mathsf{A}}$ and $\mathsf{x}^{\mathsf{B}}$. Its mismatch in the carried frame at a paired bond $c$, with initial site $c_{-}$, is
\begin{equation}\label{5.8:eq-far-datum-1}
\delta\mathsf{x}(c):=\mathsf{x}^{\mathsf{B}}(c)-\operatorname{Ad}_{g(c_{-})}\mathsf{x}^{\mathsf{A}}(c).
\end{equation}
Its unframed mismatch is
\[
\delta^{0}\mathsf{x}:=\mathsf{x}^{\mathsf{B}}-\mathsf{x}^{\mathsf{A}}.
\]
For $\mathsf{x}=\mathsf{E}$ these are $\delta\mathsf{E}$ and $\delta^{0}\mathsf{E}$.

\item \emph{Constant.} Let $C_{s}^{\sharp}$ be the constant of the critical-point size lemma (Lemma~\ref{5.8:lem-critical-size}). Put
\[
K_{U}:=\max\{1,C_{s}^{\sharp}\}.
\]
The bound on the far datum in the statement in which it is constructed carries a constant written
$B_{\sharp}$, with a subscript; it is distinct from the constant $B^{\sharp}$ bounded by
$C_{\beta}+1$. In the bounds on the far datum stated in this chapter the constant $K_{U}$ plays the
role of that constant $B_{\sharp}$.

\item \emph{Profile and cells.} For a bond $c$ of $\mathfrak{B}^{\star}_{j}$, $\bar{c}$ denotes the
$\mathfrak{B}_{j}$-cell of $c$. For a $\mathfrak{B}_{j}$-cell $y$, $\bar{y}$ denotes the
$\mathfrak{B}$-cell of $y$. On $\mathfrak{B}_{j}$ the majorant is
\[
f:=\hat{p}+|\mathsf{t}|\,\hat{p}_{\square},
\]
built from the cap profiles $\hat{p}$ and $\hat{p}_{\square}$ of the statement on caps and their
profiles. These are
\[
\hat{p}:=(\varrho\,p)^{[\delta/2]},\qquad
\hat{p}_{\square}:=(\varrho\,p\,\mathbf{1}_{\square^{\zeta}})^{[\delta/2]},
\]
where $(\cdot)^{[\delta/2]}$ is the $(\delta/2)$-hull of the statement on caps and their profiles,
where the profile $p$ at rate $\delta$ is
\[
p(y):=\sup|b|\cdot e^{-\delta d^{\wedge}(y,\mathcal{C}_{\mathsf{j}})},
\]
and the level ratio $\varrho(y):=\ell_{\mathfrak{B}_{j}}(y)/\ell_{\mathfrak{B}}(\bar{y})$ is the scaled
length of the cap at $y$ over the scaled length of the label at $\bar{y}$; it is a non-positive
power of $L$.
\end{enumerate}

\emph{The smallness ceiling.} The far datum is considered under the following condition, in
supremum form over $m$. Here $\alpha_{\cdot,m}$ are the parameters of the smallness ceiling of the
construction of the far datum. They carry two indices: the condition below is imposed for each value
of the first index, written as a dot, and the supremum is taken over the second index $m$.
The constant $c_{4}$ is the one of that ceiling, and $\beta$ is the exponent of that ceiling.
Further, $a_{R}^{\sharp}$ is the ceiling on tempered profiles of
Lemma~\ref{5.8:lem-critical-size}. The condition reads
\begin{equation}\label{5.8:eq-far-datum-a}
200\,K_{U}^{2}\,(1+\beta)\,\sup_{m}\alpha_{\cdot,m}
\ \le\ \min\Bigl\{a_{R}^{\sharp},\ \frac{c_{4}}{11}\Bigr\}.
\end{equation}
The condition \eqref{5.8:eq-far-datum-a} is one-sided. It bounds the left-hand side from above and
imposes nothing else. It is a smallness condition additional to the smallness ceiling of the
construction of the far datum, and of the same kind; it is a hypothesis of every statement below
that concerns the far datum.
\end{definition}

\begin{lemma}[One-cutoff sizes of the layers and far datum]\label{5.8:lem-one-cutoff-sizes}
Work in either of the two machines (cutoffs $\eta$ and $\eta'=\eta/L$), with the weights
and scaled sizes of
Definition~\ref{5.8:def-tempered-weights}. Let $b$ be the data, $\mathbb{H}^{w}$ the
response on the label $\mathfrak{B}$ and $\mathbf{H}_{j}$ the response on the cap layer
$\mathfrak{B}_{j}$, $\mathsf{B}(t,\mathsf{t})$ the datum of the cap layer, and
$\mathsf{E}$ the far datum of Definition~\ref{5.8:def-far-datum}. Let $p$ be the profile
at rate $\delta$, $\hat{p}$ and $\hat{p}_{\square}$ its cap versions, $\varrho$ the level
ratio, $\ell_{\mathfrak{B}_{j}}$ the scaled length of the cap layer $\mathfrak{B}_{j}$, and
$f:=\hat{p}+|\mathsf{t}|\hat{p}_{\square}$. Let $C_{1}$,
$\delta'_{\mathsf{j}}$, $\rho_{m}$, $\alpha_{\cdot,m}$ and $K^{rem}$ be the constants of
the correlation theorem, and let $m(y)$ be the index $m$ of the sequences $\rho_{m}$,
$\alpha_{\cdot,m}$ attached to the point $y$ in the correlation theorem
(Theorem~\ref{5.1:thm-correlation-theorem}).
Let $\vartheta(\square)$ be the box-dependent parameter, obtained from $\vartheta$, that
bounds the radius $r^{rem}_{\square}$ in (F4) below; it is defined by
\begin{equation*}
\vartheta(\square):=\vartheta\, e^{-\frac12\delta\, d_{sc}(\square,R_{n})},
\end{equation*}
where $d_{sc}(\square,R_{n})$ is the scaled distance between the box $\square$ and the
region $R_{n}$ of the construction of the radius $r^{rem}_{\square}$. Put
\begin{equation*}
\mathcal{T}:=\{t,\ t\zeta'_{\square},\ t+\mathsf{t}\zeta_{\square},\
t\zeta'_{\square}+\mathsf{t}\zeta_{\square}\};
\end{equation*}
for $\tau\in\mathcal{T}$, $\bar{\tau}$ denotes the $C^{2}$-size of $\tau$, and
$\mathsf{H}^{\sharp}:=C_{s}^{\sharp}\bar{\tau}p$ is the natural unit, with $C_{s}^{\sharp}$
the constant of (F1) below.
Assume the following.
\begin{enumerate}
\item[(F1)] The critical-point size at a tempered profile
(Lemma~\ref{5.8:lem-critical-size}), with constant $C_{s}^{\sharp}$; put
$K_{U}:=\max\{1,C_{s}^{\sharp}\}$.
\item[(F2)] The cutoff-size bound: a scalar built from the cutoff functions of the boxes
$\square'$ with coefficients $c_{\square'}$ has $C^{2}$-size $a$ obeying
\begin{equation*}
a(y)\le 5\sum_{\square':\ {\square'}^{\zeta}\cap\tilde{\Delta}(y)\neq\emptyset}
|c_{\square'}|.
\end{equation*}
\item[(F3)] The normalisation of the cutoff functions $\zeta_{\square}$,
$\zeta'_{\square}$ of the one-machine construction, in the form in which it is used: both
are scalars of the kind in (F2); at every $y$ the sum of the moduli of the coefficients of
$\zeta'_{\square}$ over the boxes $\square'$ with
${\square'}^{\zeta}\cap\tilde{\Delta}(y)\neq\emptyset$ is at most $1$, and the
corresponding sum for $\zeta_{\square}$ is at most
\begin{equation*}
\mathbf{1}[\square^{\zeta}\cap\tilde{\Delta}(y)\neq\emptyset],
\end{equation*}
$\square$ being the box that carries the amplitude $\mathsf{t}$.
\item[(F4)] The radius $r^{rem}_{\square}$ obeys
$r^{rem}_{\square}\le 1/(14K^{rem}\vartheta(\square))$.
\item[(F5)] For $y\in\square^{\zeta}$, the comparison of distances
\begin{equation*}
d_{sc}(\square,R_{n})\le d^{\wedge}(y,\mathcal{C}_{\mathsf{j}})
\end{equation*}
and the absorption bound
\begin{equation*}
C_{1}\delta'_{\mathsf{j}}\,e^{-\frac12\delta d^{\wedge}(y,\mathcal{C}_{\mathsf{j}})}
\le\vartheta\,\rho_{m(y)}\alpha_{\cdot,m(y)}.
\end{equation*}
\item[(F6)] The level-ratio comparison of the profiles, in the form in which it is used:
on the two blocks of a cell,
\begin{equation*}
\varrho\, p\le e^{\delta/2}\hat{p},\qquad
\varrho\, p\,\mathbf{1}[\square^{\zeta}\cap\tilde{\Delta}(y)\neq\emptyset]
\le e^{\delta/2}\hat{p}_{\square};
\end{equation*}
the factor $e^{\delta/2}$, multiplied by the exponential factor $e^{\delta d(y,y')}$ of the
passage between points $y$, $y'$ of the two blocks ($d$ the distance of
Definition~\ref{5.8:def-tempered-weights}), is at most $2$; the points of
$\tilde{\Delta}(y)$ lie at distance at most $2$ from $y$; and the rate $\delta$ obeys
$e^{2\delta}\le 2$.
\item[(F7)] The absolute bounds, in the form in which they are used: at every $y$,
$p(y)\le\sup_{m}\alpha_{\cdot,m}$; for every $m$,
\begin{equation*}
2.2\,\rho_{m}\alpha_{\cdot,m}\le 14K^{rem}\beta\,\sup_{m'}\alpha_{\cdot,m'};
\end{equation*}
and $\sup f\le(1+\beta)\sup_{m}\alpha_{\cdot,m}$. The last of these is the absolute bound
of the one-machine construction; the first two are stated here in the form needed in
item (b).
\item[(F8)] The profile majorant of the data: for data with
$\|b\|_{\infty}\le 1.1C_{1}\delta'_{\mathsf{j}}$ one has
$|b|\le p$ pointwise, and the profile obeys, at every $y$,
\begin{equation*}
p(y)\le 1.1C_{1}\delta'_{\mathsf{j}}\,e^{-\delta d^{\wedge}(y,\mathcal{C}_{\mathsf{j}})};
\end{equation*}
the cap versions $\hat{p}$ and $\hat{p}_{\square}$ belong to
$\mathfrak{M}_{2\delta_{1}}(\mathfrak{B}_{j})$.
\item[(F9)] The bound of the far datum by the response of the cap layer, in the form in
which it is used: if $N^{\Delta}_{y}(\mathbf{H}_{j})\le 20K_{U}^{2}f(y)$ at every $y$ of
the cap layer, then $|\mathsf{E}(c)|\le 100K_{U}^{2}f(\bar{c})$ at every cell $c$.
\item[(F10)] The bound of the datum of the cap layer by the response on the label, in the
form in which it is used: at every cell $c$, with $\tau\in\mathcal{T}$ the cutoff scalar
entering $\mathsf{B}(t,\mathsf{t})$,
\begin{equation*}
|\mathsf{B}(t,\mathsf{t})(c)|\le 2\ell_{\mathfrak{B}_{j}}\sup_{c}|\tau\mathbb{H}^{w}|,
\end{equation*}
where $\sup_{c}$ is the supremum over the points of the two blocks of $c$; and
$\ell_{\mathfrak{B}_{j}}=\varrho\,\ell$, with $\ell$ the scaled length on the label
$\mathfrak{B}$.
\end{enumerate}
Then for $\|b\|_{\infty}\le 1.1C_{1}\delta'_{\mathsf{j}}$, $t\in[0,1]$ and
$|\mathsf{t}|\le r^{rem}_{\square}$, in each machine, at every point $y$, every cell $c$
and every $\tau\in\mathcal{T}$ (in (c), $\tau$ the cutoff scalar entering
$\mathsf{B}(t,\mathsf{t})$, and $\sup_{c}$ as in (F10)):
\begin{equation}\label{5.8:eq-one-cutoff-sizes-a}
\begin{aligned}
&\text{(a)}\quad N^{\Delta}_{y}(\mathbb{H}^{w})\le K_{U}\,p(y);\\
&\text{(b)}\quad \bar{\tau}(y)\le 5\bigl(1+|\mathsf{t}|\,
\mathbf{1}[\square^{\zeta}\cap\tilde{\Delta}(y)\neq\emptyset]\bigr),\\
&\phantom{\text{(b)}\quad}\mathsf{H}^{\sharp}=C_{s}^{\sharp}\bar{\tau}p\le K_{U}\bar{\tau}p
\le 5K_{U}(1+\beta)\sup_{m}\alpha_{\cdot,m},\\
&\phantom{\text{(b)}\quad}|\mathsf{t}|\,p(y)
\le\frac{1.1C_{1}\delta'_{\mathsf{j}}\,
e^{-\frac12\delta d^{\wedge}(y,\mathcal{C}_{\mathsf{j}})}}{14K^{rem}\vartheta}
\le\frac{1.1\rho_{m(y)}\alpha_{\cdot,m(y)}}{14K^{rem}}
\quad(y\in\square^{\zeta});\\
&\text{(c)}\quad |\mathsf{B}(t,\mathsf{t})(c)|
\le 2\ell_{\mathfrak{B}_{j}}\sup_{c}|\tau\mathbb{H}^{w}|\le 20K_{U}f(c);\\
&\text{(d)}\quad N^{\Delta}_{y}(\mathbf{H}_{j})\le 20K_{U}^{2}f(y);\\
&\text{(e)}\quad |\mathsf{E}(c)|\le 100K_{U}^{2}f(\bar{c}).
\end{aligned}
\end{equation}
\end{lemma}

\begin{proof}
Fix a machine, $b$ with $\|b\|_{\infty}\le 1.1C_{1}\delta'_{\mathsf{j}}$, $t\in[0,1]$ and
$\mathsf{t}$ with $|\mathsf{t}|\le r^{rem}_{\square}$. The five items are proved in order.

\emph{Item (a).} Apply (F1), that is Lemma~\ref{5.8:lem-critical-size}, on the label
$\mathfrak{B}$ with the profile $\mathsf{p}:=p$, which belongs to
$\mathfrak{M}_{4\delta_{1}}(\mathfrak{B})$ by the relations
\eqref{5.8:eq-tempered-weights-a}. The other two hypotheses of that lemma hold: $|b|\le p$
by (F8), and, by the first bound of (F7), by $K_{U}\ge 1$, $\beta\ge 0$ and the smallness
ceiling of Definition~\ref{5.8:def-far-datum},
\begin{equation*}
\sup p\le\sup_{m}\alpha_{\cdot,m}\le 200K_{U}^{2}(1+\beta)\sup_{m}\alpha_{\cdot,m}
\le a_{R}^{\sharp}.
\end{equation*}
The lemma bounds the scaled size of the critical point
by $C_{s}^{\sharp}p(y)$, so $N^{\Delta}_{y}(\mathbb{H}^{w})\le C_{s}^{\sharp}p(y)\le
K_{U}p(y)$, since $K_{U}\ge C_{s}^{\sharp}$ by definition.

\emph{Item (b), the cutoff scalars.} By (F3), $\zeta'_{\square}$ and $\zeta_{\square}$ are
scalars of the kind in (F2), and inserting the coefficient sums of (F3) into the
cutoff-size bound (F2) shows that, at every $y$, the $C^{2}$-size of $\zeta'_{\square}$ is
at most $5$ and that of $\zeta_{\square}$ is at most
$5\,\mathbf{1}[\square^{\zeta}\cap\tilde{\Delta}(y)\neq\emptyset]$. A constant has
vanishing derivatives, so the $C^{2}$-size of the constant $t$ is $t\le 1$. The
$C^{2}$-size is built from suprema of the moduli of a scalar and of its derivatives up to
order two; it is therefore subadditive and absolutely homogeneous. For $\tau=t$ and
$\tau=t\zeta'_{\square}$ this gives $\bar{\tau}(y)\le 5t\le 5$; for
$\tau=t+\mathsf{t}\zeta_{\square}$ and $\tau=t\zeta'_{\square}+\mathsf{t}\zeta_{\square}$
it gives
\begin{equation*}
\bar{\tau}(y)\le 5t+5|\mathsf{t}|\,
\mathbf{1}[\square^{\zeta}\cap\tilde{\Delta}(y)\neq\emptyset].
\end{equation*}
Since $t\le 1$, hence
\begin{equation*}
\bar{\tau}(y)\le 5\bigl(1+|\mathsf{t}|\,
\mathbf{1}[\square^{\zeta}\cap\tilde{\Delta}(y)\neq\emptyset]\bigr)
\qquad(\tau\in\mathcal{T}).
\end{equation*}

\emph{Item (b), the natural unit.} The natural unit is
$\mathsf{H}^{\sharp}=C_{s}^{\sharp}\bar{\tau}p$, and $C_{s}^{\sharp}\le K_{U}$ by the
definition of $K_{U}$, so $\mathsf{H}^{\sharp}\le K_{U}\bar{\tau}p$.
Its bound rests on the estimate of the amplitude term on $\square^{\zeta}$, obtained as
follows. By the hypothesis $|\mathsf{t}|\le r^{rem}_{\square}$ and (F4),
\begin{equation*}
|\mathsf{t}|\le\frac{1}{14K^{rem}\vartheta(\square)}.
\end{equation*}
For $y\in\square^{\zeta}$, the comparison of distances in (F5), the monotonicity of the
exponential and the definition of $\vartheta(\square)$ give
\begin{equation*}
\vartheta e^{-\frac12\delta d^{\wedge}(y,\mathcal{C}_{\mathsf{j}})}
\le\vartheta e^{-\frac12\delta d_{sc}(\square,R_{n})}=\vartheta(\square),
\end{equation*}
that is,
$1/\vartheta(\square)\le e^{\frac12\delta d^{\wedge}(y,\mathcal{C}_{\mathsf{j}})}/\vartheta$
there. Since $\|b\|_{\infty}\le 1.1C_{1}\delta'_{\mathsf{j}}$, (F8) gives
$p(y)\le 1.1C_{1}\delta'_{\mathsf{j}}e^{-\delta d^{\wedge}(y,\mathcal{C}_{\mathsf{j}})}$.
Multiplying these three bounds, the factor
$e^{\frac12\delta d^{\wedge}(y,\mathcal{C}_{\mathsf{j}})}$ coming from
$1/\vartheta(\square)$ is spent against half of the decay of $p$; this is the first
inequality below. The second is the absorption bound of (F5), divided by $\vartheta$ and
multiplied by $1.1/(14K^{rem})$. Thus, for $y\in\square^{\zeta}$,
\begin{equation}\label{5.8:eq-one-cutoff-sizes-1}
|\mathsf{t}|\,p(y)
\le\frac{1.1C_{1}\delta'_{\mathsf{j}}\,
e^{-\frac12\delta d^{\wedge}(y,\mathcal{C}_{\mathsf{j}})}}{14K^{rem}\vartheta}
\le\frac{1.1\rho_{m(y)}\alpha_{\cdot,m(y)}}{14K^{rem}}.
\end{equation}
At the neighbours of $\square^{\zeta}$, the points $y\notin\square^{\zeta}$ with
$\square^{\zeta}\cap\tilde{\Delta}(y)\neq\emptyset$, pick
$y'\in\square^{\zeta}\cap\tilde{\Delta}(y)$. Since $\delta=\delta_{0}/8$ and
$\delta_{1}=\delta_{0}/32$, one has $4\delta_{1}=\delta$, so
$p\in\mathfrak{M}_{\delta}(\mathfrak{B})$; as $d(y,y')\le 2$ by (F6), the definition of
this class and $e^{2\delta}\le 2$ from (F6) give $p(y)\le e^{2\delta}p(y')\le 2p(y')$.
Hence, by \eqref{5.8:eq-one-cutoff-sizes-1} at $y'$,
\begin{equation*}
|\mathsf{t}|\,p(y)\le\frac{2.2\,\rho_{m(y')}\alpha_{\cdot,m(y')}}{14K^{rem}}.
\end{equation*}
By the first part of item (b), at
every $y$,
\begin{equation*}
\mathsf{H}^{\sharp}=C_{s}^{\sharp}\bar{\tau}p\le K_{U}\bar{\tau}p
\le 5K_{U}\bigl(p(y)+|\mathsf{t}|\,p(y)\,
\mathbf{1}[\square^{\zeta}\cap\tilde{\Delta}(y)\neq\emptyset]\bigr).
\end{equation*}
The first term is at most $\sup_{m}\alpha_{\cdot,m}$ by the first bound of (F7). The
second term vanishes unless $\square^{\zeta}\cap\tilde{\Delta}(y)\neq\emptyset$; then $y$
lies in $\square^{\zeta}$, where \eqref{5.8:eq-one-cutoff-sizes-1} applies with the index
$m_{*}:=m(y)$, or is a neighbour, where the display above applies with $m_{*}:=m(y')$. In
either case
\begin{equation*}
|\mathsf{t}|\,p(y)\le\frac{2.2\,\rho_{m_{*}}\alpha_{\cdot,m_{*}}}{14K^{rem}}
\le\beta\sup_{m}\alpha_{\cdot,m},
\end{equation*}
the last step by the second bound of (F7), which holds for every index. Adding the two
terms gives
\begin{equation*}
\mathsf{H}^{\sharp}=C_{s}^{\sharp}\bar{\tau}p\le K_{U}\bar{\tau}p
\le 5K_{U}(1+\beta)\sup_{m}\alpha_{\cdot,m}.
\end{equation*}

\emph{Item (c).} By (F10), the datum of the cap layer at a cell $c$ satisfies
\begin{equation*}
|\mathsf{B}(t,\mathsf{t})(c)|\le 2\ell_{\mathfrak{B}_{j}}\sup_{c}|\tau\mathbb{H}^{w}|
=2\varrho\sup_{c}\ell|\tau\mathbb{H}^{w}|,
\end{equation*}
and the bound of the right-hand side by $20K_{U}f(c)$ is the product of the following
factors.
First, $|\tau\mathbb{H}^{w}|\le\bar{\tau}\,|\mathbb{H}^{w}|$, with
$\bar{\tau}$ bounded by the first line of item (b), and the scaled size
$\ell|\mathbb{H}^{w}|$ is at most
$N^{\Delta}_{y}(\mathbb{H}^{w})\le K_{U}p(y)$ by the definition of $N^{\Delta}_{y}$ and
item (a). Next, the level ratio $\varrho$ is absorbed through (F6): on the two blocks of
the cell $c$,
\begin{equation*}
\varrho\, p\le e^{\delta/2}\hat{p},\qquad
\varrho\, p\,\mathbf{1}[\square^{\zeta}\cap\tilde{\Delta}(y)\neq\emptyset]
\le e^{\delta/2}\hat{p}_{\square},
\end{equation*}
which turn $p$ into $\hat{p}$ and the amplitude term into $|\mathsf{t}|\hat{p}_{\square}$;
the factor $e^{\delta/2}$, together with the exponential factor $e^{\delta d(y,y')}$ of the
passage between points $y$, $y'$ of the two blocks, is at most $2$ by (F6).
Hence
\begin{align*}
2\ell_{\mathfrak{B}_{j}}\sup_{c}|\tau\mathbb{H}^{w}|
&\le 2\cdot 5\cdot K_{U}\cdot 2\,\bigl(\hat{p}+|\mathsf{t}|\hat{p}_{\square}\bigr)(c)\\
&=20K_{U}f(c),
\end{align*}
by the definition $f=\hat{p}+|\mathsf{t}|\hat{p}_{\square}$.

\emph{Item (d).} Apply (F1), that is Lemma~\ref{5.8:lem-critical-size}, on the cap layer
$\mathfrak{B}_{j}$ with the profile $\mathsf{p}:=20K_{U}f$, which belongs to
$\mathfrak{M}_{2\delta_{1}}(\mathfrak{B}_{j})$: by (F8), $\hat{p}$ and
$\hat{p}_{\square}$ belong to this class, and by its definition in
Definition~\ref{5.8:def-tempered-weights} the class is closed under sums and under
multiplication by positive numbers (if $w_{1}(y)\le e^{2\delta_{1}d(y,y')}w_{1}(y')$ and
$w_{2}(y)\le e^{2\delta_{1}d(y,y')}w_{2}(y')$ for all $y$, $y'$, the same holds for
$a_{1}w_{1}+a_{2}w_{2}$ with $a_{1},a_{2}\ge 0$ not both zero). By item (c) the profile
majorises the datum
$\mathsf{B}(t,\mathsf{t})$. It stays below the ceiling of that lemma: by (F7), by
$K_{U}\ge 1$ and by the smallness ceiling of Definition~\ref{5.8:def-far-datum},
\begin{equation*}
\sup 20K_{U}f\le 20K_{U}(1+\beta)\sup_{m}\alpha_{\cdot,m}
\le 200K_{U}^{2}(1+\beta)\sup_{m}\alpha_{\cdot,m}\le a_{R}^{\sharp}.
\end{equation*}
The lemma gives
\begin{equation*}
N^{\Delta}_{y}(\mathbf{H}_{j})\le C_{s}^{\sharp}\cdot 20K_{U}f(y)\le 20K_{U}^{2}f(y),
\end{equation*}
again because $C_{s}^{\sharp}\le K_{U}$.

\emph{Item (e).} By item (d), $N^{\Delta}_{y}(\mathbf{H}_{j})\le 20K_{U}^{2}f(y)$ at every
$y$ of the cap layer, which is the hypothesis of (F9). Hence (F9) gives
\begin{equation*}
|\mathsf{E}(c)|\le 100K_{U}^{2}f(\bar{c})
\end{equation*}
at every cell $c$, which completes \eqref{5.8:eq-one-cutoff-sizes-a}.
\end{proof}

\begin{lemma}[The flat energy-difference bound]\label{5.8:lem-energy-difference}
Let the coarser cutoff $\mathsf{A}$ (spacing $\eta$) and the finer cutoff $\mathsf{B}$ (spacing
$\eta'=\eta/L$) act on the exactly paired admissible sets $(\mathfrak{S}',\mathfrak{S})$ of
Notation~\ref{5.8:not-paired-sets}, with label $\mathfrak{B}$, cap $\mathfrak{B}_{j}$ (and
$\mathfrak{B}'_{j+1}$ in the finer cutoff) and frame set $\mathfrak{B}^{\star}_{j}$. Quantities of
the finer cutoff carry a prime. The bases of the presented pair $p=(P',P)$ of the statement below
in the two cutoffs are $V'$, $V$, and its currents are $\mathbf{J}'$, $\mathbf{J}$;
$g=e^{i\lambda}$ is the hybrid gauge, with rotation
$R(g)$ and mismatch one-form $A^{\natural}$. The exponent $\lambda$ of $g$ is a Lie-algebra valued
function, distinct from the level function $\lambda(\cdot)$ of the paired sets: where both occur,
$|\lambda|$ and $|\lambda(y)|$ denote the former, while every power $L^{-a\lambda}$,
$L^{-a\lambda(y)}$, $L^{-a\lambda_{j}}$ with $a>0$, as well as $\lambda_{j}$ and $\ell$, are formed
from the latter.
For a Lie-algebra valued quantity $x$ at a bond $c$, present in both cutoffs, put
\begin{equation*}
\delta x(c):=x'(c)-\operatorname{Ad}_{g(c_{-})}x(c),\qquad \delta^{0}x(c):=x'(c)-x(c),
\end{equation*}
where $c_{-}$ is the initial site of $c$. In particular $\delta\mathsf{m}^{\mathfrak{B}}$ and
$\delta\mathsf{m}^{\mathfrak{B}_{j}}$ are the mismatches of the slots' data vectors on
$\mathfrak{B}$ and on $\mathfrak{B}_{j}$. Assume the following.
\begin{enumerate}
\item \emph{Exact pairing.} The admissible sets $(\mathfrak{S}',\mathfrak{S})$ are exactly paired
in the sense of Notation~\ref{5.8:not-paired-sets}, the floors \eqref{5.8:eq-paired-sets-c} hold,
and $\theta>3/5$.
\item \emph{The response.} In each cutoff the response $\mathbb{H}$ is taken to be Ba{\l}aban's
definition of it, through the constrained critical point and its Landau function, and not any
representation formula for it.
\item \emph{The hybrid gauge.} The presented pair is written in the hybrid gauge $g$; in this
representation there are constants $C_{h}$ and $C_{h}'$ with
\begin{equation*}
|A^{\natural}|_{x}\le C_{h}\Psi_{p}(x),\qquad |\lambda|\le C_{h}\mathsf{T}_{p}^{\sim},\qquad
|\delta\mathsf{m}^{\mathfrak{B}}|\le C_{h}'\Psi_{p}.
\end{equation*}
The pair $p^{\natural}$ on $\mathfrak{B}_{j}$ obtained in this gauge is an entry of modulus
$C_{h}\Psi_{p}$.
\item \emph{Healing.} $V'\in\mathsf{C}'_{\beta^{\flat}}(C_{h}\mathfrak{r}_{0})$, where
$\beta^{\flat}$ is the H\"older index of the healed class: $\mathsf{C}'_{\beta^{\flat}}$ is the
regularity class $\mathsf{C}'_{\beta}$ of Theorem~\ref{5.8:thm-response-comparison} with the indices
$(\beta,\varepsilon)$ replaced by $(\beta^{\flat},\varepsilon/2)$ and $\theta$ unchanged.
Theorem~\ref{5.8:thm-response-comparison} applies with these replaced indices, and with
$C_{h}\mathfrak{r}_{0}$ in place of $\mathfrak{r}_{0}$ in the field clauses of the base.
\item \emph{Entries.} The pair $(\mathfrak{B}'_{j+1},\mathfrak{B}_{j})$ is exactly paired and its
levels do not exceed those of the label. The entry theorem, part (b), gives on it the response
comparison at entries. At the entry $p^{\natural}$ this reads, with a constant $C_{\mathsf{m}}$,
\begin{equation*}
|\delta\mathsf{m}^{\mathfrak{B}_{j}}(c)|\le C_{\mathsf{m}}\big[(C_{h}\Psi_{p})^{\sim}
+\mathfrak{r}_{0}L^{-\theta\lambda_{j}}\big](c),
\end{equation*}
where the smear is taken in $d_{\mathfrak{B}_{j}}$.
\item \emph{Covariance and frames.} The response is covariant, so that
$\mathbb{H}(\operatorname{Ad}_{g}b;(V,\mathbf{J})^{g})=R(g)\mathbb{H}(b;V,\mathbf{J})$.
The datum in the finer cutoff is $b+\delta b_{0}$ with
$|\delta b_{0}|\le 2C_{N}\mathsf{S}L^{-\theta_{2}\lambda}|b|$, where $C_{N}$ and $\mathsf{S}$ are
the constants of the bound for the change of a datum prescribed in Ba{\l}aban's frames.
\item \emph{Levels and decay.} The hypothesis on the levels of the paired sets under which the
profiles $\hat{p}$, $\hat{p}_{\square}$ below are formed holds, and
$e^{d\delta}\le 2$, where $d=4$ is the dimension, which is also the additive constant in the
comparison of the distances on a
refinement with those on the label, on which the third of the properties
\eqref{5.8:eq-tempered-majorant-a} rests. Ba{\l}aban's local decay of propagators and projections
(Proposition~\ref{5.8:prop-local-decay}) holds.
\item \emph{Ceilings.} The ceilings \eqref{5.8:eq-weighted-chart-a},
\eqref{5.8:eq-smallness-ceilings-a} and \eqref{5.8:eq-far-datum-a} hold.
\end{enumerate}
Let $p=(P',P)$ be either a real presented pair or a real presentation over a diagonal entry, on
the label, with regularity radii $\mathfrak{r}_{P'},\mathfrak{r}_{P}\le\mathfrak{r}_{0}$ of the two
members of the pair and
$\sup\mathsf{T}_{p}\le\varepsilon_{D}$. (By the corollary on diagonal entries, every entry of the
second kind at real dilations and natural vertices is diagonal.) Let the vertex of the presentation
be the identity, $\sigma\equiv\mathbf{1}$. Let $b$ be
complex on $\mathcal{C}_{\mathsf{j}}$ with
$\|b\|_{\infty}\le 1.1C_{1}\delta'_{\mathsf{j}}$, where $\delta'_{\mathsf{j}}$ is the size parameter
of the source cube $\mathcal{C}_{\mathsf{j}}$, the same $b$ being prescribed in Ba{\l}aban's
frames in both cutoffs. Let $t\in[0,1]$, and let $\mathsf{t}$ be complex with
$|\mathsf{t}|\le r^{rem}_{\square}$. Let
$p(y):=\sup|b|\,e^{-\delta d^{\wedge}(y,\mathcal{C}_{\mathsf{j}})}$ be the profile of the datum; it
is not to be confused with the pair $p=(P',P)$, which occurs only as a subscript and in the symbol
$p^{\natural}$ for the pair in the hybrid gauge. Let
\begin{equation*}
\hat{p}:=(\varrho p)^{[\delta/2]},\qquad
\hat{p}_{\square}:=(\varrho p\,\mathbf{1}_{\square^{\zeta}})^{[\delta/2]}
\end{equation*}
be its cap versions, where $\varrho\le 1$ is the ratio of the scaled length on the cap to that on
the label. Put
$f:=\hat{p}+|\mathsf{t}|\hat{p}_{\square}$. Let $A\ge 1$, and let
$\mathsf{M}^{\mathfrak{B}}\in\mathfrak{M}^{A}_{\delta_{1}}(\mathfrak{B})$ and
$\mathsf{M}\in\mathfrak{M}^{A}_{\delta_{1}}(\mathfrak{B}_{j})$, with
$\mathsf{M}^{\mathfrak{B}}(\bar{y})\le\mathsf{M}(y)\le 1$, be majorants, in the sense that
\begin{equation}\label{5.8:eq-energy-difference-a}
\begin{gathered}
|A^{\natural}|_{y},\ |\lambda(y)|,\ |\delta\mathsf{m}^{\mathfrak{B}}(c)|_{c\ni y},\
C_{N}\mathsf{S}L^{-\theta_{2}\lambda(y)}\le\mathsf{M}^{\mathfrak{B}}(y)\quad\text{on }\mathfrak{B},\\
|\delta\mathsf{m}^{\mathfrak{B}_{j}}(c)|_{c\ni y},\ L^{-\theta_{2}\lambda_{j}(y)}\le\mathsf{M}(y)
\quad\text{on }\mathfrak{B}_{j},\qquad \lambda_{j}:=\min\{\lambda,j\}.
\end{gathered}
\end{equation}
Then at every paired bond $c$ of $\mathfrak{B}^{\star}_{j}$
\begin{equation}\label{5.8:eq-energy-difference-1}
|\delta\mathsf{E}(c)|\le C_{E}\,A\,\mathsf{M}(\bar{c})\,
\big[\hat{p}+|\mathsf{t}|\hat{p}_{\square}\big](\bar{c}),\qquad C_{E}=C_{E}(d,L,\theta,M_{1},c_{4}).
\end{equation}
The unframed mismatch $|\delta^{0}\mathsf{E}(c)|$ obeys the same bound with $C_{E}$ replaced by
$C_{E}+200K_{U}^{2}$. Here $c_{4}$ is the analyticity radius of the averaging operations of
\cite[Propositions~4, 6, 7]{B-CMP98} recalled in the proof.

The condition \eqref{5.8:eq-energy-difference-a} is met in each of the following two cases. Here
$\mathsf{t}_{p}:=\sup\mathsf{T}_{p}$, $\mathsf{T}^{\wedge}$ is the locally tempered majorant of
Definition~\ref{5.8:def-tempered-majorant}, and $C_{\star}$ is a constant fixed in the proof.
\begin{enumerate}
\item[(a)] (Supremum of the mismatch, local index.) Take $A=L$ and
\begin{equation*}
\mathsf{M}^{\mathfrak{B}}:=1\wedge C_{\star}\big[\mathsf{t}_{p}+L^{-\theta_{2}\lambda}\big],\qquad
\mathsf{M}:=1\wedge C_{\star}\big[\mathsf{t}_{p}+L^{-\theta_{2}\lambda_{j}}\big].
\end{equation*}
\item[(b)] (Local mismatch, on both sets.) Take $A=2L$ and
\begin{equation*}
\mathsf{M}^{\mathfrak{B}}:=1\wedge C_{\star}\big[\mathsf{T}_{p}^{\wedge}+L^{-\theta_{2}\lambda}\big],
\qquad
\mathsf{M}(y):=1\wedge C_{\star}\big[\mathsf{T}_{p}^{\wedge}(\bar{y})+L^{-\theta_{2}\lambda_{j}(y)}\big].
\end{equation*}
\end{enumerate}
Consequently the energy-difference hypothesis of the two-cutoff response statement at zeroth
order holds, in the form
\begin{equation*}
|\delta\mathsf{E}(c)|\le\varpi''(c)\cdot 14B_{\sharp}^{2}
\big[\hat{p}+|\mathsf{t}|\hat{p}_{\square}\big](\bar{c}),
\end{equation*}
with
\begin{equation*}
\varpi''(c):=\frac{2C_{E}LC_{\star}}{14B_{\sharp}^{2}}
\big[\mathsf{T}_{p}^{\wedge}(\bar{c})+L^{-\theta_{2}l(c)}\big],
\end{equation*}
where $l(c)$ is the level of the cell of $\mathfrak{B}^{\star}_{j}$ carrying the bond $c$.
This is at most the same expression with $\mathsf{t}_{p}$ in place of $\mathsf{T}_{p}^{\wedge}$.
The bound is linear in the inherited mismatch: no power of the mismatch appears, and the index is
local. It holds uniformly in $|\mathsf{t}|$, $\square$, $j$, $k$ and the length of the
presentation.
Since the flat profile
\begin{equation*}
p^{\flat}(y):=\sup|b|\,e^{-\frac{5}{8}\delta d^{\wedge}(y,\mathcal{C}_{\mathsf{j}})}
\end{equation*}
satisfies $p^{\flat}\ge p$, and the hull $w\mapsto w^{[\delta/2]}$ is monotone, the bound
\eqref{5.8:eq-energy-difference-1}, its analogue for $\delta^{0}\mathsf{E}$ and the bound with
$\varpi''$ remain true with $\hat{p}$,
$\hat{p}_{\square}$ formed from $p^{\flat}$ in place of $p$; this gives the flat energy-difference
bound at the rate $\tfrac{5}{8}\delta$.
\end{lemma}

\begin{proof}
\emph{The majorants.} We first show that the majorants of (a) and (b) satisfy
\eqref{5.8:eq-energy-difference-a}.

On $\mathfrak{B}$ the third hypothesis bounds $|A^{\natural}|_{x}$, $|\lambda|$ and
$|\delta\mathsf{m}^{\mathfrak{B}}|$ by $C_{h}\Psi_{p}(x)$, $C_{h}\mathsf{T}_{p}^{\sim}$ and
$C_{h}'\Psi_{p}$ respectively.
In case (a) we use $\mathsf{T}_{p}^{\sim}\le c_{1}\mathsf{t}_{p}$. In case (b) we use
$\mathsf{T}_{p}^{\sim}\le c_{1}\mathsf{T}_{p}^{\wedge}$, which is the second of the properties
\eqref{5.8:eq-tempered-majorant-a}. The same property gives
\begin{equation*}
\Psi_{p}\le c_{1}\mathsf{T}_{p}^{\wedge}+\mathfrak{r}_{0}L^{-\theta\lambda}.
\end{equation*}

On $\mathfrak{B}_{j}$ the pair $p^{\natural}$ is an entry of modulus $C_{h}\Psi_{p}$ (third
hypothesis). The pair $(\mathfrak{B}'_{j+1},\mathfrak{B}_{j})$ is exactly paired, and its levels do
not exceed the label's (fifth hypothesis). So part (b) of the entry theorem gives there the bound on
$|\delta\mathsf{m}^{\mathfrak{B}_{j}}(c)|$ displayed in the fifth hypothesis, with the smear taken
in $d_{\mathfrak{B}_{j}}$. A larger modulus is again a modulus, so the bound
on $\Psi_{p}$ just recalled may be inserted. The third of the properties
\eqref{5.8:eq-tempered-majorant-a} and the smear bound for tempered weights on $\mathfrak{B}_{j}$,
used together with $L^{-\theta\lambda}\le L^{-\theta\lambda_{j}}$, then give
\begin{equation*}
(C_{h}\Psi_{p})^{\sim}(c)\le C\big[\mathsf{T}_{p}^{\wedge}(\bar{c})+L^{-\theta\lambda_{j}(c)}\big]
\le C\big[\mathsf{t}_{p}+L^{-\theta\lambda_{j}}\big].
\end{equation*}
Since $\theta_{2}<\theta$, each term on the left of \eqref{5.8:eq-energy-difference-a} is thereby
bounded by a constant times $\mathsf{t}_{p}+L^{-\theta_{2}\lambda}$ on $\mathfrak{B}$, or
$\mathsf{t}_{p}+L^{-\theta_{2}\lambda_{j}}$ on $\mathfrak{B}_{j}$, and by a constant times the
corresponding expression with $\mathsf{T}_{p}^{\wedge}$ in place of $\mathsf{t}_{p}$. We fix
$C_{\star}$ as the largest of $1$ and of these constants. With this choice each term on the left of
\eqref{5.8:eq-energy-difference-a}
is bounded by the majorant of case (a), and by that of case (b).

That the majorants lie in the classes $\mathfrak{M}^{A}_{\delta_{1}}$ with the stated $A$ follows
from the third of the weight properties \eqref{5.8:eq-tempered-weights-a}, together with the first
and third of the properties \eqref{5.8:eq-tempered-majorant-a}.

When the presented pair is a real child of a move with the local second slot, the same verification
holds with $\mathsf{T}_{p}$ replaced by the entry modulus $\mathsf{E}_{s}$. This is parts (ii) and
(iii) of Lemma~\ref{5.8:lem-local-slot}.

\smallskip
\emph{The response on the label.} We apply Theorem~\ref{5.8:thm-response-comparison} on
$(\mathfrak{B}',\mathfrak{B})$ with bases $(V',V^{g})$. Put $\mathfrak{a}:=R(g)A^{\natural}$. Then
$M(V')=e^{i\eta\mathfrak{a}}V^{g}$, and $|\mathfrak{a}|_{y}=|A^{\natural}|_{y}$, since the rotation
$R(g)$ preserves the norm. By the healing hypothesis,
$V'\in\mathsf{C}'_{\beta^{\flat}}(C_{h}\mathfrak{r}_{0})$, and the theorem is applied at the indices
$(\beta^{\flat},\varepsilon/2)$, with the same $\theta$ and with $C_{h}\mathfrak{r}_{0}$ in place of
$\mathfrak{r}_{0}$, as that hypothesis permits. The slots are
$(\mathsf{m}_{(V',\mathbf{J}')},\operatorname{Ad}_{g}\mathsf{m}_{(V,\mathbf{J})})$. We take
$\mathsf{p}:=p$, which lies in $\mathfrak{M}_{\rho}$ with $\rho=4\delta_{1}$, and
$\mathsf{M}:=\mathsf{M}^{\mathfrak{B}}$, which lies in the class $\mathfrak{M}^{A}_{\delta_{1}}$
(the theorem's second rate taken equal to $\delta_{1}$).
The theorem's hypotheses $|\mathfrak{a}|_{y}\le\mathsf{M}^{\mathfrak{B}}(y)$ and
$|\mathsf{m}_{(V',\mathbf{J}')}-\operatorname{Ad}_{g}\mathsf{m}_{(V,\mathbf{J})}|(c)
\le\mathsf{M}^{\mathfrak{B}}(c)$ are the first and third entries of
\eqref{5.8:eq-energy-difference-a}.

By the covariance in the sixth hypothesis, the response in the coarser cutoff at the rotated datum
$\operatorname{Ad}_{g}b$ and the rotated base is $R(g)\mathbb{H}$; in the finer cutoff the datum is
$b+\delta b_{0}$, with $\delta b_{0}$ bounded as in the sixth hypothesis. Hence the two data differ by
$\delta b:=b+\delta b_{0}-\operatorname{Ad}_{g}b$. Since $g=e^{i\lambda}$, we have
$|\operatorname{Ad}_{g}b-b|\le 2|\lambda||b|$. On $\mathcal{C}_{\mathsf{j}}$ we also have
$|b|\le\sup|b|=p$. Using \eqref{5.8:eq-energy-difference-a},
\begin{equation*}
|\delta b|\le 2\big(C_{N}\mathsf{S}L^{-\theta_{2}\lambda}+|\lambda|\big)|b|
\le 4\mathsf{M}^{\mathfrak{B}}p.
\end{equation*}
Put
\begin{equation*}
\mathsf{p}_{\delta}:=4p\cdot\min\{1,(\mathsf{M}^{\mathfrak{B}})^{[\delta_{1}]}\}.
\end{equation*}
This lies in $\mathfrak{M}_{5\delta_{1}}$, since $p\in\mathfrak{M}_{4\delta_{1}}$ and the
$\delta_{1}$-hull lies in $\mathfrak{M}_{\delta_{1}}$. It dominates $|\delta b|$, because
$\mathsf{M}^{\mathfrak{B}}\le(\mathsf{M}^{\mathfrak{B}})^{[\delta_{1}]}$. It satisfies
$\mathsf{p}_{\delta}\le 4\mathsf{p}$. By the second of the weight properties
\eqref{5.8:eq-tempered-weights-a}, it is at most $4A\mathsf{M}^{\mathfrak{B}}p$.

Theorem~\ref{5.8:thm-response-comparison}, applied with these data, therefore gives, for all $y$,
\begin{equation}\label{5.8:eq-energy-difference-2}
|Y[\mathbb{H}'^{w};V']-R(g)\mathbb{H}^{w}|_{y}\le 5C_{B}^{\sharp}A\,\mathsf{M}^{\mathfrak{B}}(y)p(y).
\end{equation}

\smallskip
\emph{Cut and dilated layers.} Let $\tau$ lie in
\begin{equation*}
\mathcal{T}:=\{t,\ t\zeta'_{\square},\ t+\mathsf{t}\zeta_{\square},\
t\zeta'_{\square}+\mathsf{t}\zeta_{\square}\}.
\end{equation*}
As in Corollary~\ref{5.8:cor-cut-move}, write $Y[\cdot\,;V']=Q_{1}[\cdot]+N_{2}[\cdot]$, where
$Q_{1}$ is the linear part and $N_{2}$ the remainder. Since $Q_{1}$ is linear and the scalar $\tau$
commutes with $R(g)$,
\begin{multline*}
Y[\tau\mathbb{H}'^{w};V']-R(g)\tau\mathbb{H}^{w}
=\tau\big(Y[\mathbb{H}'^{w};V']-R(g)\mathbb{H}^{w}\big)\\
+[Q_{1},\tau]\mathbb{H}'^{w}+N_{2}[\tau\mathbb{H}'^{w}]-\tau N_{2}[\mathbb{H}'^{w}].
\end{multline*}
The first term is bounded by \eqref{5.8:eq-energy-difference-2}. The commutator and the difference
of nonlinear remainders are bounded by parts (a) and (b) of
Lemma~\ref{5.8:lem-one-cutoff-sizes}, together with $L^{-\lambda}\le\mathsf{M}^{\mathfrak{B}}$.
Together these give, with a constant $C_{2}^{E}$,
\begin{equation}\label{5.8:eq-energy-difference-3}
|Y[\tau\mathbb{H}'^{w};V']-R(g)\tau\mathbb{H}^{w}|_{y}\le C_{2}^{E}A\,\bar{\tau}(y)
\mathsf{M}^{\mathfrak{B}}(y)p(y).
\end{equation}
The amplitude $\mathsf{t}$ enters only through $\bar{\tau}$, which multiplies the bound, so no
decay rate is required of it; the product $|\mathsf{t}|p$ is bounded uniformly in $\square$ over
the disc $|\mathsf{t}|\le r^{rem}_{\square}$, by part (b) of Lemma~\ref{5.8:lem-one-cutoff-sizes}.

\smallskip
\emph{Reduction of the averaged increment to one cutoff.} Throughout this step and the next,
$\mathsf{u}$ and $\mathsf{u}'$ denote one-forms in the coarser and in the finer cutoff; the letter
$A$ is kept for the constant of the weight classes. Let
$\mathfrak{Q}\in\{\mathfrak{B}^{\star}_{j},\mathfrak{B}_{j}\}$. By
\cite[(43)]{B-CMP98} and exact pairing, $M'_{\mathfrak{Q}}=M_{\mathfrak{Q}}\circ M$ on paired
bonds. Recall the definition of $Q$: $M_{\star}(e^{i\eta\mathsf{u}}\mathbb{W})=:e^{iQ_{\star}(\mathsf{u})}
M_{\star}\mathbb{W}$. Recall also that of $Y$: the averaging $M$ carries the increment $\mathsf{u}'$
of $V'$ into the increment $Y[\mathsf{u}';V']$ of $M(V')$. These two definitions give
\begin{align*}
Q'_{\mathfrak{Q}}(\mathsf{u}';V')(c)
&=\tfrac{1}{i}\log\Big[M_{\mathfrak{Q}}\big(e^{i\eta Y[\mathsf{u}';V']}M(V')\big)(c)\cdot
M_{\mathfrak{Q}}\big(M(V')\big)(c)^{-1}\Big]\\
&=Q_{\mathfrak{Q}}\big(Y[\mathsf{u}';V'];M(V')\big)(c),
\end{align*}
and here $M(V')=e^{i\eta\mathfrak{a}}V^{g}$.

For a bond $c$ of $\mathfrak{Q}$ let $l(c)$ be the level of the cell of $\mathfrak{Q}$ containing
$c$ (for $\mathfrak{Q}=\mathfrak{B}^{\star}_{j}$ this is the $l(c)$ of the statement), and
$\ell_{l(c)}:=L^{l(c)}\eta$.
By \cite[Propositions~4, 6, 7]{B-CMP98}, $Q_{\mathfrak{Q}}(\mathsf{u};e^{i\eta a}V^{g})(c)$ has three
properties. It is analytic in $(\mathsf{u},a)\restriction\tilde{B}(c)$ on the region
$\ell_{l(c)}|\mathsf{u}|,\ \ell_{l(c)}|a|<c_{4}$. It vanishes at $\mathsf{u}=0$. It is at most
$2\ell_{l(c)}\sup|\mathsf{u}|$. Moreover, by \cite{B-CMP98},
\begin{equation*}
Q_{\mathfrak{Q}}(R(g)\mathsf{u};V^{g})(c)=\operatorname{Ad}_{g(c_{-})}Q_{\mathfrak{Q}}(\mathsf{u};V)(c).
\end{equation*}
Hence
\begin{multline*}
\delta Q_{\mathfrak{Q}}(c)
=\Big[Q_{\mathfrak{Q}}\big(Y[\mathsf{u}';V'];e^{i\eta\mathfrak{a}}V^{g}\big)
-Q_{\mathfrak{Q}}\big(R(g)\mathsf{u};e^{i\eta\mathfrak{a}}V^{g}\big)\Big](c)\\
+\Big[Q_{\mathfrak{Q}}\big(R(g)\mathsf{u};e^{i\eta\mathfrak{a}}V^{g}\big)
-Q_{\mathfrak{Q}}\big(R(g)\mathsf{u};V^{g}\big)\Big](c).
\end{multline*}
We bound the first bracket by Cauchy's inequality for the derivative in $\mathsf{u}$ on the ball of
radius $c_{4}/11$. The second bracket vanishes at $\mathfrak{a}=0$, and as a function of
$\mathsf{u}$ it vanishes at $\mathsf{u}=0$. We bound it by the Schwarz lemma in $\mathfrak{a}$, using
$|\mathfrak{a}|=|A^{\natural}|$. The result is
\begin{multline}\label{5.8:eq-energy-difference-b}
|\delta Q_{\mathfrak{Q}}(c)|\le 2.2\,\ell_{l(c)}\sup_{\tilde{B}(c)}|Y[\mathsf{u}';V']-R(g)\mathsf{u}|\\
+\frac{8}{c_{4}}\Big(\ell_{l(c)}\sup_{\tilde{B}(c)}|A^{\natural}|\Big)\cdot
\ell_{l(c)}\sup_{\tilde{B}(c)}|\mathsf{u}|.
\end{multline}
On cells of $\mathfrak{B}_{j}$ one has $l(c)\le\lambda_{j}$. Hence
$\ell_{l(c)}|\cdot|\le\varrho\cdot\ell_{\mathfrak{B}}|\cdot|$ there.

\smallskip
\emph{The datum of $\mathbf{H}_{j}$.} The datum of $\mathbf{H}_{j}$ is
$\mathsf{B}(t,\mathsf{t})=Q_{\mathfrak{B}_{j}}\big((t+\mathsf{t}\zeta_{\square})\mathbb{H}^{w};V\big)$
in the coarser cutoff, in the notation of the preceding step, $V$ being the base on which
$\mathbb{H}^{w}$ is formed. We apply \eqref{5.8:eq-energy-difference-b} at $\mathfrak{Q}=\mathfrak{B}_{j}$,
with $\mathsf{u}=\tau\mathbb{H}^{w}$ and $\tau=t+\mathsf{t}\zeta_{\square}$.

In the first term we insert \eqref{5.8:eq-energy-difference-3}. In both terms the scale
$\ell_{l(c)}$ is converted into $\varrho\,\ell_{\mathfrak{B}}$ by the last remark. In the second
term the factor $\ell_{\mathfrak{B}}\sup_{\tilde{B}(c)}|A^{\natural}|$ is then bounded by
$|A^{\natural}|_{y}\le\mathsf{M}^{\mathfrak{B}}(y)$, the first entry of
\eqref{5.8:eq-energy-difference-a}, and the factor
$\ell_{\mathfrak{B}}\sup_{\tilde{B}(c)}|\tau\mathbb{H}^{w}|$ by part (a) of
Lemma~\ref{5.8:lem-one-cutoff-sizes}. This gives, with constants $C$ and $C_{4}^{E}$,
\begin{equation*}
|\delta\mathsf{B}(c)|\le CA\,\mathsf{M}^{\mathfrak{B}}(\bar{c})\cdot
\sup_{\tilde{B}(c)}\varrho\bar{\tau}p\le C_{4}^{E}A\,\mathsf{M}(c)f(c).
\end{equation*}

The second inequality needs three facts. First, by the construction of the caps, the two blocks of
$c$ lie within $d_{\mathfrak{B}_{j}}$-distance $1$ of $c$. On them
\begin{equation*}
\varrho p\le e^{\delta/2}\hat{p},\qquad
\varrho p\mathbf{1}_{\square^{\zeta}}\le e^{\delta/2}\hat{p}_{\square},
\end{equation*}
by the definition of $\hat{p}$ and $\hat{p}_{\square}$ as $\delta/2$-hulls. Second, part (b) of
Lemma~\ref{5.8:lem-one-cutoff-sizes} gives
$\bar{\tau}\le 5(1+|\mathsf{t}|\mathbf{1}[\square^{\zeta}\cap\tilde{\Delta}(y)\ne\emptyset])$.
Third, $\mathsf{M}^{\mathfrak{B}}(\bar{y})\le\mathsf{M}(y)$.

On the finer cutoff's own collar bonds of $\mathfrak{B}'_{j+1}$, if there are any,
$|\mathsf{B}'|\le 20K_{U}f'$. Here $f'$ is the profile $f$ formed in the finer cutoff.

\smallskip
\emph{$\mathbf{H}_{j}$ on the cap.} We apply Theorem~\ref{5.8:thm-response-comparison} on
$(\mathfrak{B}'_{j+1},\mathfrak{B}_{j})$ with the same bases. This is legitimate because the
regularity classes weaken as the levels drop. The slots are $\mathsf{m}^{\mathfrak{B}_{j}}$. We
take $\mathsf{p}:=20K_{U}f$, which lies in $\mathfrak{M}_{2\delta_{1}}(\mathfrak{B}_{j})$, so
$\rho=2\delta_{1}$. We take $\mathsf{M}$, in the class $\mathfrak{M}^{A}_{\delta_{1}}(\mathfrak{B}_{j})$
(the theorem's second rate taken equal to $\delta_{1}$), and data
$(\mathsf{B}',\operatorname{Ad}_{g}\mathsf{B})$. We take
\begin{equation*}
\mathsf{p}_{\delta}:=\min\{2\mathsf{p},\ C_{4}^{E}A\,\mathsf{M}^{[\delta_{1}]}f\}
\in\mathfrak{M}_{3\delta_{1}}(\mathfrak{B}_{j}),
\end{equation*}
and the collar datum $b_{c}$, the restriction of $\mathsf{B}'$ to the finer cutoff's own collar
bonds of $\mathfrak{B}'_{j+1}$. The theorem gives, with a constant $C_{5}^{E}$,
\begin{equation}\label{5.8:eq-energy-difference-4}
\ell_{\mathfrak{B}_{j}}|Y[\mathbf{H}'_{j};V']-R(g)\mathbf{H}_{j}|(y)\le C_{5}^{E}A\,\mathsf{M}(y)f(y).
\end{equation}

\smallskip
\emph{The far datum.} Recall the definitions of the far datum and its constituents
\begin{equation*}
D^{\sharp}=\zeta'\,Q_{\star}\big((t+\mathsf{t}\zeta_{\square})\mathbb{H}^{w}\big)
+(1-\zeta')\,Q_{\star}\big(\mathbf{H}_{j}(t,\mathsf{t})\big),\qquad
\tilde{\mathsf{B}}=Q_{\star}\big(\mathbf{A}(t,\mathsf{t})\big),
\end{equation*}
with $\mathbf{A}(t,\mathsf{t})=(t\zeta'_{\square}+\mathsf{t}\zeta_{\square})\mathbb{H}^{w}$, and
$\mathsf{E}=\tfrac{1}{i}\log(e^{iD^{\sharp}}e^{-i\tilde{\mathsf{B}}})$.

If $\tilde{B}(c)\subset\tilde{\square}^{3}$, then $\mathsf{E}(c)=0$ in both cutoffs, and so
$\delta\mathsf{E}(c)=0$. Otherwise $c\in F_{\square}$. There $\zeta_{\square}=0$ on $\tilde{B}(c)$,
and the three layers are $t\mathbb{H}^{w}$, $t\zeta'_{\square}\mathbb{H}^{w}$ and $\mathbf{H}_{j}$.

We collect the three comparisons used:
\begin{equation}\label{5.8:eq-energy-difference-c}
\begin{gathered}
|Y[\mathbb{H}'^{w};V']-R(g)\mathbb{H}^{w}|_{y}\le 5C_{B}^{\sharp}A\,\mathsf{M}^{\mathfrak{B}}(y)p(y)
\quad\text{for all }y,\\
\ell_{\mathfrak{B}_{j}}|Y[\mathbf{H}'_{j};V']-R(g)\mathbf{H}_{j}|(y)\le C_{5}^{E}A\,\mathsf{M}(y)f(y),\\
|A^{\natural}|_{y}\le\mathsf{M}^{\mathfrak{B}}(y).
\end{gathered}
\end{equation}
We apply \eqref{5.8:eq-energy-difference-b} at $\mathfrak{Q}=\mathfrak{B}^{\star}_{j}$. For the
layers $t\mathbb{H}^{w}$ and $t\zeta'_{\square}\mathbb{H}^{w}$ we insert
\eqref{5.8:eq-energy-difference-3}. For $\mathbf{H}_{j}$ we insert the second line of
\eqref{5.8:eq-energy-difference-c}. The third line bounds the term in $A^{\natural}$. This gives
\begin{equation*}
|\delta Q_{\star}(t\mathbb{H}^{w})|(c),\ |\delta Q_{\star}(t\zeta'_{\square}\mathbb{H}^{w})|(c),\
|\delta Q_{\star}(\mathbf{H}_{j})|(c)\le CA\,\mathsf{M}(\bar{c})f(\bar{c}).
\end{equation*}
The function $\zeta'(c)$ is a scaffold, the same in both cutoffs. Hence
\begin{equation*}
\delta D^{\sharp}=\zeta'\,\delta Q_{\star}(t\mathbb{H}^{w})+(1-\zeta')\,\delta Q_{\star}(\mathbf{H}_{j}),
\qquad \delta\tilde{\mathsf{B}}=\delta Q_{\star}(t\zeta'_{\square}\mathbb{H}^{w}).
\end{equation*}
The map $(D,B)\mapsto\frac{1}{i}\log(e^{iD}e^{-iB})$ is analytic and
$\operatorname{Ad}$-covariant. At arguments of size at most $c_{4}/11$ it is $1.1$-Lipschitz in
each argument. By covariance, $\delta\mathsf{E}$ is the difference of its values at
$(D^{\sharp\prime},\tilde{\mathsf{B}}')$ and at
$(\operatorname{Ad}_{g}D^{\sharp},\operatorname{Ad}_{g}\tilde{\mathsf{B}})$. Therefore
\begin{equation*}
|\delta\mathsf{E}(c)|\le 1.1\big(|\delta D^{\sharp}|+|\delta\tilde{\mathsf{B}}|\big)(c).
\end{equation*}
Inserting the three bounds above yields \eqref{5.8:eq-energy-difference-1}.

Finally, $\delta^{0}\mathsf{E}=\delta\mathsf{E}+(\operatorname{Ad}_{g}-1)\mathsf{E}$, with
$|\operatorname{Ad}_{g}-1|\le 2|\lambda|\le 2\mathsf{M}^{\mathfrak{B}}$. The size of the far datum
$\mathsf{E}$ in one cutoff is bounded by part (e) of Lemma~\ref{5.8:lem-one-cutoff-sizes}. This gives
the bound for $\delta^{0}\mathsf{E}$ with $C_{E}+200K_{U}^{2}$.

The consequence stated for $\varpi''$ follows by inserting the majorant of case (b), with $A=2L$,
into \eqref{5.8:eq-energy-difference-1}. The comparison with $\mathsf{t}_{p}$ follows from
$\mathsf{T}_{p}^{\wedge}\le\sup\mathsf{T}_{p}=\mathsf{t}_{p}$.
\end{proof}

\begin{corollary}[The three-member comparison at the flat profile]\label{5.8:cor-three-member}
Assume the hypotheses of the flat energy-difference bound
(Lemma~\ref{5.8:lem-energy-difference}), and let $(P',P)$ be a pair at which that
lemma applies. Let the complex data $b$ and the complex amplitude $\mathsf{t}$ range over the
domains fixed in its hypotheses. Then the three estimates
\eqref{5.8:eq-energy-difference-c} hold at this pair on these domains. They hold linearly,
in the sense displayed by the factors below. In (i) and (iii) the factor is a constant
multiple of the majorant $\mathsf{M}^{\mathfrak{B}}$. In (ii) the right-hand side is
$2B_{\sharp}(\varpi_{\mathbb{H}}+\varpi_{\mathsf{d}})\,\mathsf{d}_{*}$, with the product
$\varpi_{\mathbb{H}}B_{\sharp}$ fixed in (i). The factors are the following; the letter $A$ in
\eqref{5.8:eq-three-member-1} denotes a constant.
\begin{enumerate}
\item[(i)] The first estimate of \eqref{5.8:eq-energy-difference-c} holds with the profile
$p(\cdot)$ itself, and with the factor
\begin{equation}\label{5.8:eq-three-member-1}
\varpi_{\mathbb{H}}B_{\sharp} := 5C_{B}^{\sharp}A\,\mathsf{M}^{\mathfrak{B}}(y).
\end{equation}
Hence it holds also with the flat profile $p^{\flat}$ of Lemma~\ref{5.8:lem-energy-difference}
in place of $p(\cdot)$, since $p^{\flat} \ge p$ pointwise.
\item[(ii)] Let $\mathsf{d}_{*} \in \mathfrak{M}_{2\delta_{1}}(\mathfrak{B}_{j})$ be arbitrary
subject to $10\sup\mathsf{d}_{*} \le a_{R}^{\sharp}$. Let data
$\mathsf{d}^{\mathsf{A}}$, $\mathsf{d}^{\mathsf{B}}$ on $\mathfrak{B}_{j}$ be given in the two
machines, and let them satisfy
$|\mathsf{d}^{\mathsf{A}}| \le \mathsf{d}_{*}$ and $|\mathsf{d}^{\mathsf{B}}| \le \mathsf{d}_{*}$;
let their two-cutoff mismatch $\delta\mathsf{d}$ satisfy
$|\delta\mathsf{d}| \le \varpi_{\mathsf{d}}\mathsf{d}_{*}$ for some $\varpi_{\mathsf{d}}\ge 0$.
Then, with the product
$\varpi_{\mathbb{H}}B_{\sharp}$ of \eqref{5.8:eq-three-member-1}, for all $y$
\begin{equation}\label{5.8:eq-three-member-2}
\ell_{\mathfrak{B}_{j}}(y)\,
\bigl|Y[\mathbf{H}_{j}^{\mathsf{B}}(\mathsf{d}^{\mathsf{B}})]
-\mathbf{H}_{j}^{\mathsf{A}}(\mathsf{d}^{\mathsf{A}})\bigr|(y)
\le 2B_{\sharp}(\varpi_{\mathbb{H}}+\varpi_{\mathsf{d}})\,\mathsf{d}_{*}(y).
\end{equation}
This is the second estimate of \eqref{5.8:eq-energy-difference-c}, now for an arbitrary
majorant $\mathsf{d}_{*}$ of the stated class.
\item[(iii)] The third estimate of \eqref{5.8:eq-energy-difference-c} holds with the factor
$\varpi_{W} := \sup_{\square_{0}}\mathsf{M}^{\mathfrak{B}}$; it is the bound on the mismatch
assumed in \eqref{5.8:eq-energy-difference-a}, read on the window box $\square_{0}$.
\end{enumerate}
\end{corollary}

\begin{proof}
Fix a pair $(P',P)$ at which
Lemma~\ref{5.8:lem-energy-difference} applies. Fix data $(b,\mathsf{t})$ in the complex domains
prescribed there.

The proof of Lemma~\ref{5.8:lem-energy-difference} proceeds in steps. Three of them
establish, under exactly the present hypotheses and on the same domains, the three
comparison estimates we need.

The opening step of that proof reads the hypothesis \eqref{5.8:eq-energy-difference-a}
of Lemma~\ref{5.8:lem-energy-difference}, which bounds the mismatch pointwise by the
majorant $\mathsf{M}^{\mathfrak{B}}$, on the window box $\square_{0}$, and yields the third
estimate of \eqref{5.8:eq-energy-difference-c} with
$\varpi_{W} = \sup_{\square_{0}}\mathsf{M}^{\mathfrak{B}}$. This is (iii).

The step immediately after it compares the responses of the two machines on the label.
It gives the first estimate of \eqref{5.8:eq-energy-difference-c}, which is
(i) with the profile $p(\cdot)$ and the factor \eqref{5.8:eq-three-member-1}.

The step of the same proof in which the responses $\mathbf{H}_{j}^{\mathsf{B}}$ and
$\mathbf{H}_{j}^{\mathsf{A}}$ are compared
on the cap layer $\mathfrak{B}_{j}$ does so for data dominated by an arbitrary
$\mathsf{d}_{*} \in \mathfrak{M}_{2\delta_{1}}(\mathfrak{B}_{j})$ with
$10\sup\mathsf{d}_{*} \le a_{R}^{\sharp}$, and the result is \eqref{5.8:eq-three-member-2},
the second estimate of \eqref{5.8:eq-energy-difference-c} for every such $\mathsf{d}_{*}$.
This is (ii).

In (i) and (iii) the factor is, by its definition, a constant multiple of
$\mathsf{M}^{\mathfrak{B}}$; in (ii) the right-hand side is
$2B_{\sharp}(\varpi_{\mathbb{H}}+\varpi_{\mathsf{d}})\,\mathsf{d}_{*}$, with the product
$\varpi_{\mathbb{H}}B_{\sharp}$ of (i).

It remains to pass from $p(\cdot)$ to $p^{\flat}$ in (i). The factor
$5C_{B}^{\sharp}A\,\mathsf{M}^{\mathfrak{B}}(y)$ is non-negative, since the constants
$C_{B}^{\sharp}$ and $A$ are non-negative and
$\mathsf{M}^{\mathfrak{B}}$ is a majorant. Moreover $p(y) \le p^{\flat}(y)$ for every $y$.
The right-hand side of the first estimate of \eqref{5.8:eq-energy-difference-c} is this
factor times the profile at $y$; hence replacing $p(y)$ by $p^{\flat}(y)$ only enlarges it.
This is (i) with $p^{\flat}$, and the factor \eqref{5.8:eq-three-member-1} is unchanged.
\end{proof}

\begin{corollary}[The local mismatch factor in the zeroth-order response]
\label{5.8:cor-mismatch-response}
In the setting and under the hypotheses of Lemma~\ref{5.8:lem-energy-difference}, in particular
$e^{d\delta}\le 2$, assume the following statements.
\begin{itemize}
\item The flat energy-difference bound (Lemma~\ref{5.8:lem-energy-difference}). It bounds the
two-cutoff mismatch $\delta\mathsf{E}$ of the far datum at $c$ by the local factor $\varpi''$
times a constant multiple of the profile majorant $f := \hat{p}+|\mathsf{t}|\hat{p}_{\square}$,
evaluated at $\bar{c}$.
\item Clause (e) of the theorem that represents the third member of the two-cutoff response
statement at zeroth order as a kernel sum.
\item The lemma that frees one quarter of the decay rate $\delta$ in that statement.
\item Clauses (g2) and (g3) of the lemma on labels and their refinements.
\item The level-function hypothesis of the flat energy-difference bound.
\end{itemize}
Here $x$ is the point at which the third member of the two-cutoff response statement at zeroth
order is evaluated, and $l_{x}$ is its level. In the kernel sum \eqref{5.8:eq-mismatch-response-1}
below, $T^{(1)}[b^{\mathsf{B}}-b^{\mathsf{A}}]$ denotes that third member, and the letter $b$ is
the generic datum of clause (e): in the present application it is the far datum
$\mathfrak{h}:=\mathsf{E}$, not the cap-layer data $b$ of Lemma~\ref{5.8:lem-energy-difference},
and $b^{\mathsf{A}}$, $b^{\mathsf{B}}$ are its values at the cutoffs $\mathsf{A}$ and
$\mathsf{B}$, whose difference is the two-cutoff mismatch $\delta\mathsf{E}$. The distance
$d_{\star}$ is the one in which the kernel $k$ below decays, and $K_{R}$ is the constant of that
kernel. For a rate $\rho>0$ let $\mathfrak{M}^{\wedge,L}_{\rho}(\mathfrak{B}_{j})$ and
$\mathfrak{M}^{\wedge,2}_{\rho}(\mathfrak{B}_{j})$ be the classes, in the sense of
Definition~\ref{5.8:def-tempered-majorant} applied to the label $\mathfrak{B}_{j}$ and its
distance $d^{\wedge}_{\mathfrak{B}_{j}}$ (the distance in which clause (g2) and property (t3)
temper), of profiles $\mathsf{T}\ge 0$ on $\mathfrak{B}_{j}$ with
\begin{align*}
\mathsf{T}(y)&\le L\,e^{\rho\, d^{\wedge}_{\mathfrak{B}_{j}}(y,y')}\,\mathsf{T}(y'),
&&\text{respectively}\\
\mathsf{T}(y)&\le 2\,e^{\rho\, d^{\wedge}_{\mathfrak{B}_{j}}(y,y')}\,\mathsf{T}(y'),
&&\text{for all } y,y'.
\end{align*}
Property (t3) of Definition~\ref{5.8:def-tempered-majorant} places
$y\mapsto\mathsf{T}^{\wedge}(\bar{y})$ in $\mathfrak{M}^{\wedge,2}_{\delta_{1}/2}(\mathfrak{B}_{j})$
when $\mathfrak{B}_{j}$ refines the label of $\mathsf{T}$.
The letter $C$ denotes a constant independent of $x$ and $c$. The local factor $\varpi''$ of the
flat energy-difference bound enters the third member at $c$, up to the constant $C$, through the
value $\mathsf{M}(\bar{c})$ of the tempered majorant $\mathsf{M}$ of the inherited mismatch. On
the label $\mathfrak{B}_{j}$ this majorant is the profile
\begin{equation*}
c\longmapsto\mathsf{M}(\bar{c})=\mathsf{T}^{\wedge}_{p}(\bar{c})+L^{-\theta_{2}\lambda_{j}(c)},
\qquad c\in\mathfrak{B}_{j},
\end{equation*}
and $\mathsf{M}(x)$ denotes the value of this profile at the point $x$. In that value the first
summand $\mathsf{T}^{\wedge}_{p}(\bar{x})$ is written $\mathsf{T}^{\wedge}_{p}(x)$, and the level
$\lambda_{j}(x)$ is $l_{x}$; thus
\begin{equation*}
\mathsf{M}(x)=\mathsf{T}^{\wedge}_{p}(x)+L^{-\theta_{2}l_{x}}.
\end{equation*}
Then the following hold.
\begin{enumerate}
\item[(a)] Every majorant and every radius of the two-cutoff response statement at zeroth order
stands unchanged. The profile entering it is its own $\mathfrak{h}_{*}$, a constant multiple of
$f$, and no flattened profile or flattened rate is needed.
\item[(b)] By clause (e) of the theorem in the second hypothesis, the third member of that
statement reads
\begin{equation}\label{5.8:eq-mismatch-response-1}
T^{(1)}[b^{\mathsf{B}}-b^{\mathsf{A}}]=\sum_{c}k(c)\,|b^{\mathsf{B}}-b^{\mathsf{A}}|(c),
\qquad k(c)=K_{R}\,e^{-\delta d_{\star}(x,c)} .
\end{equation}
In it the local factor $\varpi''$ costs no more than the quarter of $\delta$ freed by the lemma
of the third hypothesis. Precisely,
\begin{equation}\label{5.8:eq-mismatch-response-2}
\begin{split}
\sum_{c}e^{-\delta d_{\star}(x,c)}\,\mathsf{M}(\bar{c})\,\mathfrak{h}_{*}(c)
&\le 4L\,\mathsf{M}(x)\sum_{c}e^{-\frac{7}{8}\delta d_{\star}(x,c)}\,\mathfrak{h}_{*}(c)\\
&\le 4L\,\mathsf{M}(x)\,T_{*}(x),
\end{split}
\end{equation}
where $T_{*}$ is the majorant of the two-cutoff response statement at zeroth order that bounds
the kernel sum $\sum_{c}e^{-\delta' d_{\star}(x,c)}\,\mathfrak{h}_{*}(c)$, taken at the rate
$\delta'=\tfrac{3}{4}\delta$. Consequently the display of the
two-cutoff response statement at zeroth order is supplied with
\begin{equation*}
\varpi'\ \longmapsto\ C\bigl[\mathsf{T}^{\wedge}_{p}(x)+L^{-\theta_{2}l_{x}}\bigr].
\end{equation*}
This lies within that display's own format, in the notation of that display,
$\varpi+\varpi'+c_{1}\vartheta^{l_{x}}$.
\item[(c)] The far datum in the working gauge, $\mathsf{E}^{u}=\operatorname{Ad}_{u}\mathsf{E}$,
where $u$ is the working-gauge transformation,
differs from $\mathsf{E}$ by one further block-local gauge transformation of the data at each of
the two cutoffs $\mathsf{A}$ and $\mathsf{B}$. This transformation is $\mathrm{SU}(N)$-valued at
real parents. It is absorbed in the same way as the data gauges in clause (h5) of the theorem named
in the second hypothesis.
\end{enumerate}
\end{corollary}

\begin{proof}
(a) The two-cutoff response statement at zeroth order takes the mismatch of the far datum as a
hypothesis of the form $|\delta\mathfrak{h}(c)|\le\varpi'\,\mathfrak{h}_{*}(c)$, with
$\mathfrak{h}:=\mathsf{E}$. The flat energy-difference bound
(Lemma~\ref{5.8:lem-energy-difference}) bounds $|\delta\mathsf{E}|$ by $\varpi''$ times a multiple
of $f=\hat{p}+|\mathsf{t}|\hat{p}_{\square}$. The profile it supplies is therefore the
statement's own $\mathfrak{h}_{*}$, proportional to $f$. No majorant and no radius of the
statement has to be altered, and no flattened profile or rate enters.

(b) Fix $x$. By clause (e), the third member is the kernel sum
\eqref{5.8:eq-mismatch-response-1}. In it the generic datum $b$ of clause (e) is the far datum,
$b=\mathfrak{h}=\mathsf{E}$, so that $|b^{\mathsf{B}}-b^{\mathsf{A}}|(c)=|\delta\mathsf{E}(c)|$.
The bound of the first hypothesis, in the form recorded in the proof of (a) and with the local
factor evaluated at $\bar{c}$, gives
\begin{equation*}
|b^{\mathsf{B}}-b^{\mathsf{A}}|(c)\le C\,\mathsf{M}(\bar{c})\,\mathfrak{h}_{*}(c).
\end{equation*}
Inserting this in \eqref{5.8:eq-mismatch-response-1} and dividing by $CK_{R}$, the quantity to
be bounded is the left-hand side of \eqref{5.8:eq-mismatch-response-2}. Three inputs are used.

First, clause (g2) of the lemma on labels, together with the level-function hypothesis, gives
\begin{equation*}
L^{-\theta_{2}\lambda_{j}}\in\mathfrak{M}^{\wedge,L}_{\delta/16}(\mathfrak{B}_{j}).
\end{equation*}

Second, $\delta=\delta_{0}/8$ and $\delta_{1}=\delta_{0}/32$, so
$\delta_{1}/2=\delta_{0}/64=\delta/8$. Property (t3) of the locally tempered majorant
(Definition~\ref{5.8:def-tempered-majorant}, display~\eqref{5.8:eq-tempered-majorant-a}) is
applied to the refinement $\mathfrak{B}_{j}$. It gives that $y\mapsto\mathsf{T}^{\wedge}_{p}(\bar{y})$
lies in $\mathfrak{M}^{\wedge,2}_{\delta/8}(\mathfrak{B}_{j})$.

Third, clause (g3) of the lemma on labels gives $d^{\wedge}_{\mathfrak{B}_{j}}\le d_{\star}+d$.

By the first two inputs, the two summands of $c\mapsto\mathsf{M}(\bar{c})$ lie in
$\mathfrak{M}^{\wedge,L}_{\delta/16}(\mathfrak{B}_{j})$ and
$\mathfrak{M}^{\wedge,2}_{\delta/8}(\mathfrak{B}_{j})$. Since $\delta/16\le\delta/8$, both
satisfy the tempering inequality at rate $\delta/8$, with constants $L$ and $2$ respectively, and
their sum therefore satisfies, for every $c$,
\begin{equation*}
\mathsf{M}(\bar{c})\le\max\{L,2\}\,e^{\frac{1}{8}\delta d^{\wedge}_{\mathfrak{B}_{j}}(x,c)}\,\mathsf{M}(x)
\le 2L\,e^{\frac{1}{8}\delta d^{\wedge}_{\mathfrak{B}_{j}}(x,c)}\,\mathsf{M}(x).
\end{equation*}
By the third input, $d^{\wedge}_{\mathfrak{B}_{j}}(x,c)\le d_{\star}(x,c)+d$, and
$e^{\frac{1}{8}\delta d}\le e^{\delta d}\le 2$
by the condition $e^{d\delta}\le 2$ among the hypotheses of the flat energy-difference bound.
Hence
\begin{equation*}
\mathsf{M}(\bar{c})\le 4L\,e^{\frac{1}{8}\delta d_{\star}(x,c)}\,\mathsf{M}(x),
\end{equation*}
the constant $4L$ being the product of the tempering constant $2L$ and the factor $2$. Since
$e^{-\delta d_{\star}}e^{\frac{1}{8}\delta d_{\star}}=e^{-\frac{7}{8}\delta d_{\star}}$, this is
the first inequality of \eqref{5.8:eq-mismatch-response-2}.

For the second inequality, note that $\tfrac{7}{8}\delta\ge\delta'=\tfrac{3}{4}\delta$. Hence
$e^{-\frac{7}{8}\delta d_{\star}(x,c)}\le e^{-\delta' d_{\star}(x,c)}$ for every $c$. The remaining
sum is therefore bounded by $T_{*}(x)$, taken at the rate $\delta'$ at which the lemma of the
third hypothesis leaves it. The decay consumed by the local factor is $\delta/8$, which is less
than the quarter $\delta/4$ that this lemma frees.

By the convention fixed in the statement,
$\mathsf{M}(x)=\mathsf{T}^{\wedge}_{p}(x)+L^{-\theta_{2}l_{x}}$. By
\eqref{5.8:eq-mismatch-response-2}, the
third member is therefore bounded in the form that the display of the two-cutoff response
statement at zeroth order requires. The only change is that $\varpi'$ is replaced by
$C[\mathsf{T}^{\wedge}_{p}(x)+L^{-\theta_{2}l_{x}}]$. This is a term of the display's own format,
in the notation of that display, $\varpi+\varpi'+c_{1}\vartheta^{l_{x}}$.

(c) Write $\mathsf{E}^{u}=\operatorname{Ad}_{u}\mathsf{E}$ for the far datum in the working gauge.
At each of the two cutoffs it differs from $\mathsf{E}$ by one more block-local gauge
transformation of the data, which is $\mathrm{SU}(N)$-valued at real parents. Such a transformation
is absorbed exactly as the data gauges are absorbed in clause (h5) of the theorem whose clause (e)
was used in (b).
\end{proof}

\begin{corollary}[Absolute mismatch bound at complex ancestor parameters]
\label{5.8:cor-complex-ancestors}
Let $\Pi \mapsto p(\Pi) = (P'(\Pi),P(\Pi))$ be a chain of moves and restrictions at natural vertices,
starting from an entry. Here $\Pi$, consisting of the entry data together with the parameters
$(b_s,t_s)_s$ of the successive stages, denotes the ancestors' parameters, and one and the same
$\Pi$ is used for both cutoffs $\mathsf{A}$ and $\mathsf{B}$. Let $\mathfrak{P}$ be the product of
the prescribed polydiscs of these parameters. For the move whose prescribed polydisc has radius
$R^{\natural}_{\square}$ one has $R^{\natural}_{\square} = \mathsf{w}_{0}\,r^{\sharp}_{\square}$ with
the ratio $\mathsf{w}_{0} = 1/24$, so that $4R^{\natural}_{\square} = r^{\sharp}_{\square}/6$; here
$r^{\sharp}_{\square}$ denotes the radius of reference fixed for that move. Fix the data
$(b,t,\mathsf{t})$ of the last stage in the prescribed domain of the last move: $b$ complex with
$\|b\|_{\infty} \le 1.1C_{1}\delta'_{\mathsf{j}}$, $t \in [0,1]$, and $\mathsf{t}$ complex with
$|\mathsf{t}| \le r^{rem}_{\square}$, as in Lemma~\ref{5.8:lem-energy-difference}. Write
$\delta^{0}\mathsf{E}(\Pi;b,t,\mathsf{t})$ for the
unframed two-cutoff mismatch of the far datum produced by the chain at $\Pi$ and at these last-stage
data. Assume the following.
\begin{enumerate}
\item[(a)] On $4\mathfrak{P}$ all compositions of moves and restrictions, for every box, are
defined, are holomorphic, and stay in the stage spaces. This includes the one-cutoff size (e)
of Lemma~\ref{5.8:lem-one-cutoff-sizes}, in the display
\eqref{5.8:eq-one-cutoff-sizes-a}, at the points of $4\mathfrak{P}$.
\item[(b)] On the real trace of $4\mathfrak{P}$ the pairs $p(\Pi)$ are real presentations, and
there is $\mathsf{t}_{\star}$ with
\begin{equation*}
\sup \mathsf{T}_{p(\Pi)} \le \mathsf{t}_{\star} \le \varepsilon_{D}.
\end{equation*}
\item[(c)] The flat energy-difference bound, Lemma~\ref{5.8:lem-energy-difference}, in the form
in which it applies to a real presented pair with complex last-stage data $b$ and $\mathsf{t}$.
\item[(d)] The two-constants lemma for a function holomorphic on a disc which is bounded on the
disc and bounded on a diameter of the disc.
\end{enumerate}
Then for every $\Pi \in \mathfrak{P}$ and every paired bond $c$,
\begin{equation}\label{5.8:eq-complex-ancestors-1}
|\delta^{0}\mathsf{E}(\Pi;b,t,\mathsf{t})(c)|
\le f(\bar{c})\cdot\Bigl\{200K_{U}^{2}\,\bigl(C_{E}+200K_{U}^{2}\bigr)\,L\,C_{\star}
\bigl[\mathsf{t}_{\star}+L^{-\theta_{2}l(c)}\bigr]\Bigr\}^{1/2},
\end{equation}
where $f = \hat{p} + |\mathsf{t}|\hat{p}_{\square}$ is the majorant and $l(c)$ is defined as in the
flat energy-difference bound, Lemma~\ref{5.8:lem-energy-difference}.
\end{corollary}

\begin{proof}
Fix $\Pi \in \mathfrak{P}$ and a paired bond $c$. Write $\Pi = \Pi_{1} + i\Pi_{2}$ with $\Pi_{1}$ and
$\Pi_{2}$ real. Since each component of $\Pi_{1}$ and of $\Pi_{2}$ is the real or the imaginary
part of the corresponding component of $\Pi$, we have $|\Pi_{1}| \le |\Pi|$ and
$|\Pi_{2}| \le |\Pi|$ componentwise. The map $z \mapsto \Pi_{1} + z\Pi_{2}$ therefore sends the
disc $D(0,3)$ into $4\mathfrak{P}$, and it sends the real interval $]{-3},3[$ into the real trace of
$4\mathfrak{P}$, because for real $z$ the point $\Pi_{1} + z\Pi_{2}$ is real.

Let $\mathsf{l}$ be a linear functional of norm one, and set
\begin{equation*}
\varphi(z) := \mathsf{l}\bigl(\delta^{0}\mathsf{E}(\Pi_{1}+z\Pi_{2};b,t,\mathsf{t})(c)\bigr),
\qquad z \in D(0,3).
\end{equation*}

\emph{Holomorphy.} By hypothesis (a), every composition of moves and restrictions is defined and
holomorphic at the points $\Pi_{1}+z\Pi_{2}$, $z \in D(0,3)$, and stays in the stage spaces. The
response and the far datum are analytic in the block data, the slot data and the background
field, $(B,\mathsf{m},\mathbf{U})$, by the regularity statement for the response and by part (b)
of the healing statement, which gives the analyticity in $(B,\mathsf{m},\mathbf{U})$; these rest
in turn on the analyticity of the averaging operations,
\cite[Propositions~6--7]{B-CMP98}. Composing these holomorphic dependences, $\varphi$ is
holomorphic on $D(0,3)$.

\emph{Bound on the disc.} At every point of $4\mathfrak{P}$ the one-cutoff size (e) of
Lemma~\ref{5.8:lem-one-cutoff-sizes} is available by hypothesis (a). It gives, at every
$z \in D(0,3)$, a bound for the mismatch, which is a difference of two one-cutoff quantities:
\begin{equation}\label{5.8:eq-complex-ancestors-2}
|\varphi(z)| \le 2\cdot 100K_{U}^{2}\, f(\bar{c}), \qquad z \in D(0,3).
\end{equation}

\emph{Bound on the diameter.} Let $z \in\, ]{-3},3[$. Then $\Pi_{1}+z\Pi_{2}$ lies in the real trace
of $4\mathfrak{P}$. By hypothesis (b) the parent pair $p(\Pi_{1}+z\Pi_{2})$ is a real presentation
with $\sup \mathsf{T}_{p} \le \mathsf{t}_{\star} \le \varepsilon_{D}$. The flat energy-difference
bound, Lemma~\ref{5.8:lem-energy-difference}, admits a real presented pair together with complex
last-stage data $b$ and $\mathsf{t}$; it therefore applies at the complex $(b,\mathsf{t})$ fixed
above. Replacing $\sup \mathsf{T}_{p}$ by its majorant $\mathsf{t}_{\star}$, it gives
\begin{equation}\label{5.8:eq-complex-ancestors-3}
|\varphi(z)| \le \bigl(C_{E}+200K_{U}^{2}\bigr)\,L\,C_{\star}
\bigl[\mathsf{t}_{\star}+L^{-\theta_{2}l(c)}\bigr] f(\bar{c}),
\qquad z \in\, ]{-3},3[.
\end{equation}

\emph{Interpolation.} We now apply the two-constants lemma to the holomorphic function
$\varphi$ on $D(0,3)$, with the bound \eqref{5.8:eq-complex-ancestors-2} on the disc and the bound
\eqref{5.8:eq-complex-ancestors-3} on the diameter $]{-3},3[$, and evaluate at $z = i$. At this point
$\Pi_{1}+i\Pi_{2} = \Pi$, and the result is the geometric mean of the two bounds:
\begin{equation*}
|\varphi(i)|
\le \Bigl\{2\cdot 100K_{U}^{2} f(\bar{c})\cdot
\bigl(C_{E}+200K_{U}^{2}\bigr)L\,C_{\star}
\bigl[\mathsf{t}_{\star}+L^{-\theta_{2}l(c)}\bigr] f(\bar{c})\Bigr\}^{1/2}.
\end{equation*}
The right-hand side equals the right-hand side of \eqref{5.8:eq-complex-ancestors-1}. It does not
depend on $\mathsf{l}$. The norm of $\delta^{0}\mathsf{E}(\Pi;b,t,\mathsf{t})(c)$ is the supremum of
$|\mathsf{l}(\cdot)|$ over functionals of norm one, so \eqref{5.8:eq-complex-ancestors-1} follows.
\end{proof}

The square root in \eqref{5.8:eq-complex-ancestors-1} falls on
$\mathsf{t}_{\star}+L^{-\theta_{2}l(c)}$. This is an absolute bound for the real closure. The
two-constants lemma is applied once, to the total mismatch, and not move by move. Nothing is fed
back into the hypotheses, and the disc of the last stage is not shrunk. The absolute form
\eqref{5.8:eq-complex-ancestors-1} is the form in which this bound is required where the corollary
is applied.

Hypothesis (b) is assumed in the corollary and is not derived here for the polydiscs
$\mathfrak{P}$. On the circles of the uniqueness theorem, hypothesis (b) is supplied by the local
growth statement together with its corollary, and the local growth argument uses only a one-move
bound and the accompanying identity. On the polydiscs $\mathfrak{P}$ the one-move input is
supplied by the increment bound for a move by a cut response
(Corollary~\ref{5.8:cor-cut-move}) and the natural size of the cut response on the contour
(Corollary~\ref{5.8:cor-contour-size}). The bound \eqref{5.8:eq-cut-move-a} has the form of the
one-move bound used in the local growth argument, with $\mathsf{H}$ replaced by
$\mathsf{H}^{\sharp}$, and $\mathsf{H}^{\sharp}$
is a natural unit by clause (b) of Lemma~\ref{5.8:lem-one-cutoff-sizes}. The geometric part of the
argument for a chain of moves, consisting of the statements on the move geometry among which is
the increment statement, is not carried through here in the unit $\mathsf{H}^{\sharp}$; on the
polydiscs $\mathfrak{P}$ hypothesis (b) is therefore taken by statement.
The corresponding linear bound at complex ancestor parameters is not asserted. It would require
the analysis of the response on complex backgrounds.

\begin{proposition}[The bound off the natural decoupling vertices]\label{5.8:prop-off-vertex}
Let $\mathsf{s}(\Delta)$ denote the decoupling parameters of a family of presented pairs and
$\sigma(\Delta)$ the decoupling parameters of the interpolated family $K(\sigma)$ displayed in
\textup{(c)}. A \emph{vertex} is a parameter point at which every decoupling parameter takes the
value $0$ or $1$. The \emph{natural vertices} are the vertices singled out in the lemma on the
natural vertices; for the family $K(\sigma)$ the natural vertex is the point described in
\textup{(c)}. The vertex function of $\delta\mathsf{E}$ is the function on the vertices whose
value at $\mathsf{s}$ is the value of $\delta\mathsf{E}$ at the pair of the family with parameter
$\mathsf{s}$. For a finite index set $\mathsf{B}$ of decoupling parameters,
$(\delta\mathsf{E})^{\wedge}(\mathsf{B})$
denotes its M\"obius coefficient with index set $\mathsf{B}$, that is, the mixed difference, in
the parameters $\mathsf{s}(\Delta)$, $\Delta\in\mathsf{B}$, of its values at the vertices.
Here $\mathsf{B}$ is an index set and not the finer cutoff. The argument $c$ of
$\delta\mathsf{E}$, the point $\bar{c}$ and the majorant $f$ are as in
Lemma~\ref{5.8:lem-energy-difference}. Assume the following statements:
\begin{itemize}
\item the M\"obius-coefficient lemma for families of pairs parametrised by the decoupling
parameters, in its clause \textup{(a)};
\item the proposition bounding M\"obius coefficients through vertex values and bounds at one
cutoff alone, in its clause \textup{(b)};
\item the relative small-field machine, in its clause \textup{(s1)}, which supplies the
operators $K_{\omega}$ of \textup{(c)}, together with the random-walk bound quoted there;
\item the lemma on the natural vertices;
\item the representation of the interpolated family $K(\sigma)$ in one lattice geometry;
\item clauses \textup{(d)} and \textup{(e)} of the statement on natural vertices for tilted
presentations;
\item the statement that the quantities carried at either cutoff read the decoupling parameters,
of either family, only at natural vertices.
\end{itemize}
Then the following hold.
\begin{enumerate}
\item[(a)] Suppose that the conclusion of Lemma~\ref{5.8:lem-energy-difference}, in its form for
presented pairs, holds at every pair of some family
of pairs of the kind to which the M\"obius-coefficient lemma applies, parametrised by the
decoupling parameters on a polydisc $|\mathsf{s}(\Delta)| \le \rho$ and passing through the
pairs at the natural vertices; here $\rho$ is the radius of the polydisc and not the contraction
constant of clause \textup{(v-d)}. Then for every index set $\mathsf{B}$
\begin{equation}\label{5.8:eq-off-vertex-1}
\bigl|(\delta\mathsf{E})^{\wedge}(\mathsf{B})(c)\bigr|
\le \Bigl(\frac{2}{\rho-1}\Bigr)^{|\mathsf{B}|}\,\varpi''\cdot 14\,B_{\sharp}^{2}\,f(\bar{c}).
\end{equation}
The right-hand side of \eqref{5.8:eq-off-vertex-1} is small and, when $\rho>3$, decays
geometrically in $|\mathsf{B}|$.
\item[(b)] The bounds at the vertices, namely the flat energy-difference bound of
Lemma~\ref{5.8:lem-energy-difference} applied at every box $Z$ to which clauses \textup{(d)}
and \textup{(e)} of the statement on natural vertices for tilted presentations attach a pair
$(\mathbf{\Omega}'(Z),\mathbf{\Omega}(Z))$ (here $Z$ denotes a box, not a
large-field region), together with bounds at one cutoff alone, give for the M\"obius
coefficients only
\begin{equation}\label{5.8:eq-off-vertex-2}
\bigl|(\delta\mathsf{E})^{\wedge}(\mathsf{B})(c)\bigr|
\le \min\bigl\{2^{|\mathsf{B}|}\beta,\; q^{|\mathsf{B}|}G\bigr\},
\end{equation}
with $\beta$, $q$ and $G$ the constants of clause \textup{(b)} of the proposition bounding
M\"obius coefficients through vertex values and bounds at one cutoff alone; these letters are
local to this item and are not the
H\"older index $\beta$ or the auxiliary parameter $q$. The bound \eqref{5.8:eq-off-vertex-2} is
sharp. The bound \eqref{5.8:eq-off-vertex-1} of \textup{(a)} is not implied by these inputs.
\item[(c)] The interpolated family is
\begin{equation}\label{5.8:eq-off-vertex-3}
K(\sigma) := \sum_{\omega}\Bigl(\prod_{\Delta:\,\Delta\cap\operatorname{dom}\omega\neq\emptyset}
\sigma(\Delta)\Bigr)K_{\omega},
\end{equation}
where $\omega$ runs over the walks of the representation, $\operatorname{dom}\omega$ is the
domain of the walk $\omega$, and all operators are on one lattice geometry. Its only natural
vertex is $\sigma \equiv \mathbf{1}$. For a set $Y$ of indices $\Delta$, let $\mathbf{1}_{Y}$
be the parameter point with $\sigma(\Delta)=1$ for $\Delta\in Y$ and $\sigma(\Delta)=0$
otherwise. At $\sigma = \mathbf{1}_{Y} \neq \mathbf{1}$ its operators are sums
over walks confined to $Y$, and no bound comparing such sums between the lattices of the two
cutoffs $\eta$ and $\eta'$ is proved.
\item[(d)] The proof of the uniqueness theorem reads the decoupling parameters, of either
family, only at natural vertices.
\end{enumerate}
\end{proposition}

\begin{proof}
(a) By hypothesis, the flat energy-difference bound holds along the family on the polydisc
$|\mathsf{s}(\Delta)| \le \rho$, and the family passes through the pairs at the natural
vertices. These are the hypotheses of clause (a) of the M\"obius-coefficient lemma. Applying that
clause to the family of values $\delta\mathsf{E}(c)$ yields \eqref{5.8:eq-off-vertex-1}.

(b) By clauses (d) and (e) of the statement on natural vertices for tilted presentations, the
hypotheses of Lemma~\ref{5.8:lem-energy-difference} are met at every box $Z$, so that
Lemma~\ref{5.8:lem-energy-difference} bounds the vertex values at every box $Z$. The
bounds at one cutoff alone supply the remaining input. Clause (b) of the proposition bounding
M\"obius coefficients through vertex values and bounds at one cutoff alone converts these two
inputs into
\eqref{5.8:eq-off-vertex-2}. By the same clause the bound is sharp. Sharpness is witnessed by the
family
\begin{equation*}
\mathsf{s} \longmapsto a\prod_{\Delta\in\mathsf{B}}\bigl(2\mathsf{s}(\Delta)-1\bigr),
\end{equation*}
where $a$ is a constant and $\mathsf{s}(\Delta)$, $\Delta \in \mathsf{B}$, are the decoupling
parameters indexed by $\mathsf{B}$: its values at the vertices have modulus $|a|$, while its
M\"obius coefficient at $\mathsf{B}$ is $2^{|\mathsf{B}|}a$.

(c) By the representation of the interpolated family in one lattice geometry, $K(\sigma)$ has the
form \eqref{5.8:eq-off-vertex-3}, with all its operators on one lattice geometry; the summation
runs over walks $\omega$, and the operators $K_{\omega}$ are those given by clause (s1) of the
relative small-field machine, with the random-walk bound quoted there.
Its only natural vertex is $\sigma \equiv \mathbf{1}$. Now take $\sigma = \mathbf{1}_{Y}$. For a
walk $\omega$, the product in \eqref{5.8:eq-off-vertex-3} equals $1$ if every $\Delta$ meeting
$\operatorname{dom}\omega$ lies in $Y$, and $0$ otherwise. Hence $K(\mathbf{1}_{Y})$ is the sum
of $K_{\omega}$ over the walks $\omega$ with that property, which is a walk-confined sum. By the
lemma on the natural vertices, off the natural vertex both the constraint and the stationarity
equation acquire walk-sum defects. Comparing these defects between the two cutoffs is exactly
the two-lattice comparison of walk-confined sums named in (c), and this comparison is not
proved.

(d) By the statement that the quantities carried at either cutoff read the decoupling parameters
only at natural vertices, these quantities depend on the decoupling parameters, of either
family, only through their values at natural vertices. Since the proof of the uniqueness theorem
reads the flat energy-difference bound only through these quantities, it reads the decoupling
parameters only at natural vertices.
\end{proof}

\begin{remark}[The bound off the natural vertices is not proved]
The flat energy-difference bound off the natural vertices, on the polydisc
$|\sigma(\Delta)| \le e^{\kappa_{1}}$ in the decoupling parameters $\sigma(\Delta)$ of the
family \eqref{5.8:eq-off-vertex-3}, is not proved. Clause (a) of
Proposition~\ref{5.8:prop-off-vertex}, by contrast, is a conditional statement about a family in
the parameters $\mathsf{s}(\Delta)$, with the radius $\rho$ among its hypotheses. An argument
for the bound on $|\sigma(\Delta)| \le e^{\kappa_{1}}$
would have to supply the two-cutoff comparison of the walk-sum defects described in clause (c)
of Proposition~\ref{5.8:prop-off-vertex}.

By clauses (d) and (e) of the statement on natural vertices for tilted presentations,
Lemma~\ref{5.8:lem-energy-difference} applies without change to the pair
$(\mathbf{\Omega}'(Z),\mathbf{\Omega}(Z))$ attached to the box $Z$ by that statement, for every
box $Z$. It is the bound at the natural vertices that clause (b) of
Proposition~\ref{5.8:prop-off-vertex} uses and that the proof of the uniqueness theorem, as
described in clause (d), uses.
\end{remark}

\begin{remark}[The two off-shell extensions of the response]\label{5.8:rem-off-shell}
The flat energy-difference bound (Lemma~\ref{5.8:lem-energy-difference}) is a statement about
two objects fixed among its hypotheses. The first is the response $\mathbb{H}$, built from
Ba{\l}aban's operator $H$ taken by its definition and not by the representation
\cite[(3.120)]{B-CMP99b}. The second is the window field $\mathbb{W}$ of the hybrid
representation. A second convention for the response is the one in force in the theorem on the
response of the off-shell pair extension. It extends the response of the presented pair
$p = (P',P)$ of Lemma~\ref{5.8:lem-energy-difference} off shell by replacing the current
$\mathbf{J}$ with $\mathbb{J}$ in every propagator and every operator,
Ba{\l}aban's operator $H$ of \cite[(3.126)]{B-CMP99b} included, so that $\mathbb{J}$ enters the
response through that operator.

On shell the two conventions agree. Off shell they are two holomorphic extensions of the
response.

In the second convention the response is the critical point of the functional
\begin{equation}\label{5.8:eq-off-shell-1}
\mathcal{W} - \langle C, \mathsf{m}_{P}\rangle
+ \tfrac{1}{2}\,\bigl\| R D^{*} X + \mathsf{h}_{P}\, C(X) \bigr\|^{2} .
\end{equation}
Here $\mathcal{W}$ is the action functional. The block function is $C$, evaluated at $X$ in the
last term, and
$RD^{*}X$ is the gauge component of $X$. The coefficient with which $C(X)$ enters that term is
$\mathsf{h}_{P}$. The operator $H_{P}$ enters the data vector $\mathsf{m}_{P}$.
In this remark $H$, $R$ and $D^{*}$ are operators of Ba{\l}aban's construction and $X$ is the
variable of the variational problem; they are not the Hamiltonian, the ball radius or the
large-field regions denoted by the same letters elsewhere.

The second convention also calls for a comparison of the critical point of
\eqref{5.8:eq-off-shell-1} between the coarser cutoff $\eta$ and the finer cutoff
$\eta' = \eta/L$. When the mismatch between the two cutoffs is measured in a supremum norm, this
comparison meets an operator of order zero acting on quantities
$e_{s}$ that are small only in the supremum norm. The tempered weights used in the estimates of
this section improve decay rates, not the order of an operator. They therefore do not remove
this obstruction. This comparison is the statement addressed by the route below.

In the proof of the uniqueness theorem, the extension at a single cutoff is re-based on the
convention of Lemma~\ref{5.8:lem-energy-difference}, at a price paid inside that proof.
In that proof the flat energy-difference bound speaks of the proof's own objects.

Proof outstanding.

\emph{Route.} A proof has to compare the critical points of \eqref{5.8:eq-off-shell-1} at the
cutoffs $\eta$ and $\eta'$. The route is to carry the mismatch between the two cutoffs in a
H\"older norm of index $\varepsilon$ instead of the supremum norm. In that norm the term of
order zero is bounded linearly in the mismatch by \cite[(3.44), (3.45)]{B-CMP99b}.

What then has to be established is that the H\"older norm is preserved under one step of the
construction. Three facts are needed:
\begin{itemize}
\item the seeds are small in the class $C^{1}$, since the norm $\|\cdot\|_{\theta}$ of
the theorem on the initial data of the iteration controls the gradient;
\item the fresh terms are small in the H\"older class $C^{\varepsilon}$ at the index
$\theta(1-\varepsilon)$, by interpolation, at the cost of replacing the small number that bounds
a fresh term by a fractional power of that number;
\item the stability estimate of the step, with forces of format $\sharp^{+}$, returns increments
that are small in the class $C^{1}$.
\end{itemize}
Preservation of the H\"older norm under the step, together with the linear bound on the term of
order zero, is what the comparison requires.
\end{remark}

\begin{remark}[The three slot rules at a cut]\label{5.8:rem-slot-rules}
The flat energy-difference bound (Lemma~\ref{5.8:lem-energy-difference}) reads the second slot of
the parent pair only through property (A) and clause (ii) of the hybrid representation, that is,
through clauses (e1) and (e3) of the definition of an entry. At a cut, that is, where
$\nabla\tau\neq 0$, three rules for moving the second slot are in use:
\begin{align}
\mathbf{J}_{s} &:= \mathbf{J}_{s-1}
  + \mathcal{J}_{\mathfrak{S}_{s}}(h_{s};\mathbf{U}_{s-1}),
  \qquad h_{s} := \tau\mathbb{H}, \label{5.8:eq-slot-rules-1}\\
\mathcal{J}_{\mu} &:= \tau_{t}\mathfrak{L}_{Z} + [D^{*}D,\tau_{t}]\,\mathbb{H}_{Z}
  + F(\mathbf{U},\tau_{t}\mathbb{H}_{Z}), \label{5.8:eq-slot-rules-2}\\
\Phi_{t} &:= \bigl(e^{i\eta\chi_{t}\mathbb{H}}\,\mathbb{U},\;
  \mathbb{J} + \chi_{t}\,\mathcal{J}^{\sharp}[1]\bigr). \label{5.8:eq-slot-rules-3}
\end{align}
Here $\tau$, $\tau_{t}$ and $\chi_{t}$ are the cutoff scalars of the three rules and $h_{s}$ is the
cut field; $\mathfrak{L}_{Z}$, $\mathbb{H}_{Z}$, $F$, $D$, the index $\mu$ and the region $Z$ are
those of the statement in which the cutoff scalar $\tau_{t}$ is introduced.
The first is the slot rule with which property (A) is proved; the second is the
slot rule of the statement in which the cutoff scalar $\tau_{t}$ is introduced; the third is the
local second slot of Lemma~\ref{5.8:lem-local-slot} (see \eqref{5.8:eq-local-slot-1}). The three
rules agree at $\tau\equiv 1$. The relative small-field machine treats the children of the moves
with the local second slot as presented pairs;
what follows justifies this. The real children of the moves with the local second slot (chains of
every
length, the amplitudes $\mathsf{t}$ admitted in Lemma~\ref{5.8:lem-local-slot},
$\sigma\equiv\mathbf{1}$) and their restrictions are entries satisfying (A) in the local currency
$\mathsf{M}_{s}$, so the flat energy-difference bound applies to them as stated. The parents at
the sites (T) and (L) of the expansion rule (Definition~\ref{5.1:def-expansion-rule}) are entries
by the corollary on entries at real dilations and natural vertices.
\end{remark}

\begin{proof}
In the proof of the flat energy-difference bound (Lemma~\ref{5.8:lem-energy-difference}), the
second slot of the parent enters only through property (A) and clause (ii) of the hybrid
representation. The proof of property (A) reads the presentation
of the parent only through the pairing condition on the label sets and through clause (b) of the
second defining condition of a presented pair; accordingly, reading the second slot through (A)
and clause
(ii) of the hybrid representation amounts to reading it through clauses (e1) and (e3) of the
definition of an entry.

We now describe the three rules. In the first rule \eqref{5.8:eq-slot-rules-1}, the slot
$\mathbf{J}$ moves by the increment of \cite[(3.13)]{B-CMP109} evaluated at the cut field
$h_{s}=\tau\mathbb{H}$, the same cut field as in the move of the first slot. In the second rule
\eqref{5.8:eq-slot-rules-2}, the cutoff scalar $\tau_{t}$ multiplies $\mathfrak{L}_{Z}$ and also
enters inside the commutator with $D^{*}D$ and inside the term $F$. In the third rule
\eqref{5.8:eq-slot-rules-3}, the cutoff scalar
$\chi_{t}$ multiplies the response in the exponent of the background field and multiplies the
increment $\mathcal{J}^{\sharp}[1]$ in the current.

By the lemma on moves with the local second slot (Lemma~\ref{5.8:lem-local-slot}), the following
are entries, of the modulus fixed in that lemma, and satisfy (A) in the local currency
$\mathsf{M}_{s}$:
\begin{itemize}
\item the real children of the moves with the local second slot \eqref{5.8:eq-slot-rules-3}, of
every length, at the amplitudes $\mathsf{t}$ admitted in that lemma, with
$\sigma\equiv\mathbf{1}$;
\item their restrictions under the restricting move of the chain construction, which is applied
to pairs.
\end{itemize}
By the first paragraph, clauses (e1) and (e3) and property (A) are all that the flat
energy-difference bound reads of the parent's second slot. What that bound requires of the
parent's second slot therefore holds for these children, and under its remaining hypotheses it
applies to them as stated.
\end{proof}

\begin{definition}[Low-depth cores and their counterterm companions]\label{5.9:not-low-depth-cores}
Throughout, the two machines $\mathsf{A}$ and $\mathsf{B}$, the stage class $\mathcal{K}^{\mathrm{rd}}$ with its list of moves, among them the box-read move $\mathrm{M}2$, and the $\delta$-sizes are those of the definition of the two-machine stage classes and their $\delta$-sizes. Fix a level $k$ and work at the site $(\mathrm{P})$ of the operation $\mathbf{R}_k$ (Definition~\ref{5.2:def-sites-units}). At this site fix a block $\square$ of zone $n_0$, with its enlargement $\tilde{\square}^2$. Let $v$ be a near core at $\square$ of one of the first three kinds in the classification of near cores, with core functional $F$, level $j$, localisation domain $X\in\mathfrak{D}_j$ (Definition~\ref{5.1:def-analyticity-spaces}) and anchor $y\in\tilde{\square}^2$. On the machine $\mathsf{m}\in\{\mathsf{A},\mathsf{B}\}$ the core functional is written $F^{\mathsf{m}}$. For a functional $f$ read on both machines we write $\delta f:=f^{\mathsf{B}}-f^{\mathsf{A}}$.

\smallskip\noindent\emph{Depth.} We set
\begin{equation}\label{5.9:eq-low-depth-cores-1}
n:=\min\{\ell(y),n_0\},\qquad \nu:=1+n-j.
\end{equation}
We call $v$ a \emph{low-depth core} if $\nu\le 2$.

\smallskip\noindent\emph{The Wilson functional.} Let $\mathbf{U}$ be a bond configuration with values in $G^{c}$, and let $\partial\mathbf{U}(P)$ be its plaquette variables. For a function $c$ on plaquettes we set
\begin{equation}\label{5.9:eq-low-depth-cores-2}
A(c,\mathbf{U}):=\sum_P c(P)\,a(\partial\mathbf{U}(P)),\qquad
a(V):=\mathfrak{c}_0-\tfrac12\operatorname{tr}\bigl(V+V^{-1}\bigr)
=-\tfrac12\operatorname{tr}\bigl(V^{-1}(V-1)^2\bigr).
\end{equation}
The \emph{counterterm companion} of $v$ is $\beta[F]\,\hat{A}$. Here $\beta[F]$ is the counterterm coefficient attached to the core $F$ in the definition of the counterterm coefficients of cores, and it depends linearly on $F$; further,
\begin{equation*}
\hat{A}:=A(h_y,\cdot).
\end{equation*}
The function $h_y$ is the tent of \cite[(3.65)]{B-CMP119}, and
\begin{equation*}
\operatorname{supp}h_y\subset y+(-1,1)^d\,\ell_j .
\end{equation*}

\smallskip\noindent\emph{The two machines.} The functional $\hat{A}$ is read on the box pair $\mathfrak{b}'_j=\mathsf{P}^0$ through the companion-slot form $\mathrm{M}2^{\circ}$ of the box-read move $\mathrm{M}2$. The read lands in the slot $\mathfrak{u}^{\circ}_v:=(j,X;\circ)$. On machine $\mathsf{m}$ the read gives $\hat{A}^{\mathsf{m}}$. The companion then splits across the machines as
\begin{equation}\label{5.9:eq-low-depth-cores-a}
\delta\bigl(\beta[F]\hat{A}\bigr)=\beta[\delta F]\,\hat{A}^{\mathsf{B}}+\beta[F^{\mathsf{A}}]\,\delta\hat{A}.
\end{equation}
Let $\mathcal{K}^{\mathrm{rd}}(\mathfrak{u}^{\circ}_v)$ be the stage class at the slot. Its $\circ$-slot $\delta$-size is
\begin{equation}\label{5.9:eq-low-depth-cores-c}
\mathsf{S}^{\circ}[\hat{A}]:=\sup\bigl\{|\delta\hat{A}(p')|:\ p'\in\mathcal{K}^{\mathrm{rd}}(\mathfrak{u}^{\circ}_v)\bigr\}
\ \le\ \mathsf{a}:=\mathsf{S}^{\mathrm{rd}}_j[\hat{A}].
\end{equation}

\smallskip\noindent\emph{Weights and sizes near the anchor.} Let $w(\Delta)$ be the frozen cell weights. Further, $\mathfrak{a}_n$ denotes the field size at level $n$, $\operatorname{dist}_n$ denotes distance in units of the level-$n$ spacing, and for a bond $b$, $|b-y|$ denotes the distance of its base point from $y$ in level-$n$ units. We set
\begin{equation}\label{5.9:eq-low-depth-cores-b}
\begin{gathered}
\Theta:=c_W\,w^{\vee}(y),\qquad
w^{\vee}(y):=\max\{w(\Delta):\ \operatorname{dist}_n(\Delta,y)<1\},\qquad
D:=M(2\nu+3),\\
t:=L^{j-n},\qquad
|B^{\mathrm{ax}}|_{*}(b):=2d\,(1+|b-y|)\,\Theta\,\mathfrak{a}_n\,t^2 .
\end{gathered}
\end{equation}
For a cell $x$ of the scaffold we set
\begin{equation}\label{5.9:eq-low-depth-cores-3}
\theta_2(x):=w(x)\,r_2(\ell(x)),\qquad
r_2(i):=\tilde\alpha_{0,i}\,\eta^2\,\ell_i^{-2}.
\end{equation}
For the near cores considered here $\nu\ge 1$; moreover $\Theta\ge 1$, and $\tilde\alpha_{0,i}\le 2\alpha_{0,i}$ for every level $i$.

\smallskip\noindent\emph{The class-$\mathrm{W}$ unit and its chart.} To the companion we attach the class-$\mathrm{W}$ unit
\begin{equation*}
u:=(1,h_y,j,y),\qquad T_u=\hat{A}.
\end{equation*}
Its base-point set $S_u$ is the set of base points of the plaquettes of $h_y$, and
\begin{equation*}
S_u\subset y+[-1,1]^d\,\ell_j .
\end{equation*}
Let $a_M$ be the least $a\ge 1$ with $L^a\ge 4M$, and set $m_u:=\max\{j-a_M,0\}$. The shifted own box is $\mathfrak{r}^{\circ}_u:=[S_u]^{(2)}_{m_u}$, the level-$m_u$ box built on $S_u$, with index $2$ in the bracket notation $[\,\cdot\,]^{(2)}_{m}$ for boxes, and $\mathfrak{r}^{\circ,(1)}_u$ denotes its one-layer collar. The chart of $u$ is the absolute chart at level $m_u$,
\begin{equation*}
(U_{m_u},J_{m_u})(\mathfrak{r}^{\circ}_u,e^{iB'}),\qquad |B'|_{\infty}<\alpha_3(\circ).
\end{equation*}
Here $\alpha_3(\circ):=c_3a^{\circ}$ is the radius of the absolute chart; it is an absolute constant, $c_3$ and $a^{\circ}$ being the constants fixed with that chart. Finally, $B^0$ is the comb datum of $\mathsf{P}^0$ restricted to $\mathfrak{r}^{\circ}_u$, taken from $y$.
\end{definition}

\begin{proof}[Proof of the identities \eqref{5.9:eq-low-depth-cores-2}, \eqref{5.9:eq-low-depth-cores-a} and the bound \eqref{5.9:eq-low-depth-cores-c}]
\emph{The two forms of $a(V)$.} For invertible $V$,
\begin{equation*}
V^{-1}(V-1)^2=V^{-1}\bigl(V^2-2V+1\bigr)=V-2\cdot 1+V^{-1}.
\end{equation*}
Taking $-\tfrac12\operatorname{tr}$ of this gives
\begin{equation*}
-\tfrac12\operatorname{tr}\bigl(V^{-1}(V-1)^2\bigr)
=\operatorname{tr}1-\tfrac12\operatorname{tr}\bigl(V+V^{-1}\bigr).
\end{equation*}
Hence the two expressions for $a(V)$ in \eqref{5.9:eq-low-depth-cores-2} coincide exactly when the normalising constant is $\mathfrak{c}_0=\operatorname{tr}1=1$, the trace being normalised so that the identity has trace $1$; this is the value of $\mathfrak{c}_0$.

\emph{The split \eqref{5.9:eq-low-depth-cores-a}.} By the definition of $\delta$, the left-hand side is
\begin{equation*}
\delta\bigl(\beta[F]\hat{A}\bigr)=\beta[F^{\mathsf{B}}]\hat{A}^{\mathsf{B}}-\beta[F^{\mathsf{A}}]\hat{A}^{\mathsf{A}}.
\end{equation*}
Add and subtract $\beta[F^{\mathsf{A}}]\hat{A}^{\mathsf{B}}$ on the right. This gives
\begin{equation*}
\bigl(\beta[F^{\mathsf{B}}]-\beta[F^{\mathsf{A}}]\bigr)\hat{A}^{\mathsf{B}}
+\beta[F^{\mathsf{A}}]\bigl(\hat{A}^{\mathsf{B}}-\hat{A}^{\mathsf{A}}\bigr).
\end{equation*}
By linearity of $\beta[\cdot]$, the first bracket is $\beta[F^{\mathsf{B}}-F^{\mathsf{A}}]=\beta[\delta F]$; the second bracket is $\delta\hat{A}$ by definition. This is the right-hand side of \eqref{5.9:eq-low-depth-cores-a}.

\emph{The bound in \eqref{5.9:eq-low-depth-cores-c}.} By the definition of the $\delta$-size $\mathsf{S}^{\mathrm{rd}}_j$, the quantity $\mathsf{a}=\mathsf{S}^{\mathrm{rd}}_j[\hat{A}]$ dominates $e^{\kappa d_j(X)}|\delta\hat{A}(p')|$ for every $p'\in\mathcal{K}^{\mathrm{rd}}(\mathfrak{u}^{\circ}_v)$. The weight satisfies $e^{\kappa d_j(X)}\ge 1$, because $\kappa d_j(X)\ge 0$. Hence $|\delta\hat{A}(p')|\le\mathsf{a}$ for every such $p'$, and the unweighted supremum $\mathsf{S}^{\circ}[\hat{A}]$ is at most $\mathsf{a}$.
\end{proof}

\begin{remark}[Why the absolute chart is level-shifted]\label{5.9:rem-shifted-chart}
Let $v$ be a near core of level $j$ with anchor $y$, as in the notation for low-depth cores and their
counterterm companions, Definition~\ref{5.9:not-low-depth-cores}. Put $n := \min\{\ell(y),n_0\}$,
$\nu := 1+n-j$ and
$t := L^{j-n}$, and assume $\nu \le 2$. Let $\hat{A} := A(h_y,\cdot)$ be the Wilson functional of the
counterterm companion of $v$. It is read on the box pair $\mathfrak{b}'_j = \mathsf{P}^0$ through
$\mathrm{M}2^{\circ}$ into the companion slot $\mathfrak{u}^{\circ}_v$, and
$\delta\hat{A} := \hat{A}^{\mathsf{B}} - \hat{A}^{\mathsf{A}}$. Then the following hold.
\begin{enumerate}
\item The low-order variation bound at the second cube size is a bound on $\delta\hat{A}$ at the base
point $\mathsf{P}^0$ itself. It is not a Taylor coefficient at $\mathrm{M}2$ of the dilation families
at which the two-cutoff response statement at zeroth order and its energy-difference companion are
invoked.
\item The dial that carries its factor $(CD\Theta\mathfrak{a}_n t^2)^2$ is the comb-chart parameter
$\sigma_0$ of the base. It runs from the flat configuration ($\sigma_0 = 0$) to $\mathsf{P}^0$
($\sigma_0 = 1$).
\item A chart for that dial on the own box $\mathfrak{r}_v = \tilde{\square}^{\nu+1}$ at the native level
is not available when $\nu \le 2$. The chart is therefore taken at the shifted level
$m := \max\{j-a_M,0\}$, where it exists at every class-W unit, the companions at every $\nu$ included,
by the lemmas that construct the level-shifted absolute chart.
\end{enumerate}
\end{remark}

\begin{proof}
For the first item, consider the dilation families at which the two-cutoff response statement at
zeroth order and its energy-difference companion are invoked. They are the dilation in a complex
parameter $s$ of the relative layer, and the dilation in a
complex parameter $\sigma_1$ of the increment $\Delta_c$. Both interpolate from the configuration
$U_{\mathsf{P}^0}$ of the base point to the configuration $U_{\mathsf{P}^1}$ of the dilated point, and
both keep the base $\mathsf{P}^0$ of the box pair fixed. Every Taylor coefficient in $s$ or in
$\sigma_1$ is therefore a statement about the response to the far datum $\mathsf{E}$.

The low-order variation bound at the second cube size concerns $\delta\hat{A}$ at $\mathsf{P}^0$
itself, that is, the value of the family at the origin of the parameter. A family that interpolates
from $U_{\mathsf{P}^0}$ to $U_{\mathsf{P}^1}$ with $U_{\mathsf{P}^0}$ fixed as base gives no
information on that value beyond what the maximum-modulus principle gives, namely the value bound
\begin{equation*}
|\delta\hat{A}(\mathsf{P}^0)| \le \mathsf{S}^{\circ}[\hat{A}] \le \mathsf{a},
\end{equation*}
where $\mathsf{a} := \mathsf{S}^{\mathrm{rd}}_j[\hat{A}]$. The required bound thus has no object among the
Taylor coefficients of these families.

For the second item, the factor $(CD\Theta\mathfrak{a}_n t^2)^2$ in the bound of the form
\begin{equation*}
|\delta\hat{A}| \le 2\,\mathsf{a}\,(CD\Theta\mathfrak{a}_n t^2)^2
\end{equation*}
that the low-order variation bound at the second cube size asks for
is the one produced by the Schwarz lemma for a holomorphic function with a zero of second order at the
origin. That lemma yields it from a second-order zero of a family that dilates the own curvature of
$\mathfrak{b}'_j$, from the
flat configuration to the natural radius of the chart of $\hat{A}$. The parameter of such a family is
the comb-chart parameter $\sigma_0$ of the base, running from the flat configuration ($\sigma_0 = 0$)
to $\mathsf{P}^0$ ($\sigma_0 = 1$).

For the third item, let $\pi_{n-1}$ be the partition into level-$(n-1)$ cubes. A chart for
$\sigma_0$ on the own box $\mathfrak{r}_v = \tilde{\square}^{\nu+1}$ at the native level is available
when $\nu \ge 3$, for then, for $L \ge L_0$, the sides of $\mathfrak{r}_v$ measured in cubes of
$\pi_{n-1}$ are less than $2$. Hence the region $Q^{(1)}_v$ attached to $v$ in the construction of
that chart stays within one level-$n$ unit of $y$.

When $\nu \le 2$ this fails. The box $\tilde{\square}^{\nu+1}$ is $5L$ cubes of $\pi_{n-1}$ wide when
$\nu = 1$ and $7$ cubes wide when $\nu = 2$, so it leaves that level-$n$ unit. A chart on the own box
at the native, unshifted level therefore is not available for $\nu \le 2$. This is why the chart is
taken at the level $m = \max\{j-a_M,0\}$, shifted down from $j$ by $a_M$ (and equal to $0$ when
$j < a_M$), where $a_M$ is the least $a \ge 1$ with $L^a \ge 4M$.

In the notation of Definition~\ref{5.9:not-low-depth-cores}, the class-W unit $u = (1,h_y,j,y)$ of the
companion
has functional $T_u = \hat{A}$ and shifted own box $\mathfrak{r}^{\circ}_u = [S_u]^{(2)}_m$. The
absolute chart $(U_m,J_m)(\mathfrak{r}^{\circ}_u, e^{iB'})$, for $|B'|_{\infty} < \alpha_3(\circ)$,
exists at every class-W unit, the companions at every $\nu$ included and in particular those with
$\nu \le 2$, by the lemmas that construct the level-shifted absolute chart; on it the
comb datum $B^0$ of $\mathsf{P}^0$ restricted to $\mathfrak{r}^{\circ}_u$ is taken from $y$.
\end{proof}

\begin{definition}[hypotheses: the absolute chart and its comb dilation]
\label{5.9:def-variation-hypotheses}
We work in the setting of Definition~\ref{5.9:not-low-depth-cores}
(low-depth cores and their counterterm companions), with the letters $\Theta$,
$w^{\vee}(y)$, $D$, $t$ of
\eqref{5.9:eq-low-depth-cores-b}. The near core $v$ with anchor $y$, the numbers $n$ and
$\nu$, the functional $\hat A$ read through $\mathrm{M}2^{\circ}$ into the slot
$\mathfrak{u}^{\circ}_v$, the unit $u := (1,h_y,j,y)$ with $T_u = \hat A$, its base-point
set $S_u$, its chart level $m$, its own box $\mathfrak{r}^{\circ}_u$ with one-layer collar
$\mathfrak{r}^{\circ,(1)}_u$, the absolute chart
$B' \mapsto (U_m,J_m)(\mathfrak{r}^{\circ}_u,e^{iB'})$ for $|B'|_\infty < \alpha_3(\circ)$,
and the comb datum $B^{0,\mathsf{m}}$ of the base point $\mathsf{P}^{0,\mathsf{m}}$ in the
machine $\mathsf{m} \in \{\mathsf{A},\mathsf{B}\}$ are as there. For a configuration
$\mathsf{P}$ (a point or a chart value) and a set $\mathfrak{r}$ of bonds or of base points,
$\mathsf{P}|_{\mathfrak{r}}$ denotes its restriction to $\mathfrak{r}$. By
$\operatorname{dist}_n$ we denote the distance between cells measured at level $n$.

The hypotheses of the two-machine variation estimate for the counterterm companion at
$\nu \le 2$ are (A1), (A2), (A3) and, at $\nu = 1$, in addition (A3)$'$.

\smallskip\noindent
(A1) \emph{The absolute chart, in one machine.} The following holds in each machine
$\mathsf{m}$, under the floor conditions of the box-inclusion and absolute-chart lemmas (a lower
bound on $L$, the condition $M \ge 3L$, and two further floor conditions of those lemmas). It
holds for every point of the stage class
$\mathcal{K}^{\mathrm{rd}}(\mathfrak{u}^{\circ}_v)$ and for every parameter of the box read
$\mathrm{M}2$: interpolation parameter in $[0,1]$, dial $\mathsf{t} \in \mathbb{D}_{\square}$,
and $|\sigma| \le e^{\kappa_1}$.
Here and below, $\mathfrak{A}_w(\gamma)$ denotes the bound on the weight-radius products
given by the cell-wise bound of the weight arrangement, which holds throughout the scaffold:
$w(\Delta)\,\alpha_{0,\ell(\Delta)} \le \mathfrak{A}_w(\gamma)$ for every cell $\Delta$ of
$\mathfrak{B}$.
\begin{itemize}
\item[(o)] One has
$\mathfrak{r}^{\circ,(1)}_u \subset y+[-4,4]^d\ell_j \subset \tilde{\square}^3$. Every
$\mathfrak{B}$-cell meeting $y+[-4,4]^d\ell_j$ has level at least $\max\{j-1,0\}$, and
$\max\{j-1,0\} \ge m$.
\item[(i)] The base point restricted to the own box is represented by the chart at its
comb datum:
\begin{equation*}
\mathsf{P}^{0,\mathsf{m}}|_{\mathfrak{r}^{\circ}_u} \cong
(U_m,J_m)(\mathfrak{r}^{\circ}_u,e^{iB^{0,\mathsf{m}}}).
\end{equation*}
\item[(ii)] The comb datum obeys the local bound
\begin{equation}\label{5.9:eq-variation-hypotheses-a}
|B^{0,\mathsf{m}}(b)| \le \mathsf{b}_u := 16\,d\,C^{\mathrm{av}}(1+D^{\circ})\,\omega_u
\qquad\text{for every bond } b \text{ of } \mathfrak{B}_m(\mathfrak{r}^{\circ}_u),
\end{equation}
where
\begin{equation}\label{5.9:eq-variation-hypotheses-1}
\omega_u := \max\bigl\{ w(x)\,\alpha_{0,\ell(x)} :
x \text{ a } \mathfrak{B}\text{-cell meeting } \mathfrak{r}^{\circ,(1)}_u \bigr\}
\le \mathfrak{A}_w(\gamma).
\end{equation}
\item[(iii)] For $|B'|_\infty < \alpha_3(\circ)$, the chart
$(U_m,J_m)(\mathfrak{r}^{\circ}_u,e^{iB'})$ is analytic in $B'$. Moreover $|T_u(\cdot)|$,
evaluated on it, is bounded by $\mathfrak{n}^{\circ}_u$.
\end{itemize}
Item (o) is the box-inclusion lemma, and items (i) and (iii) are statements of the
absolute-chart lemma; item (ii) is the comb-datum bound of the absolute-chart lemma in local
form, obtained from that lemma's proof by the one modification described below. The
box-inclusion lemma, statements (i) and (iii) of the absolute-chart lemma, the cell-wise bound
of the weight arrangement, the regime facts of the near-core analysis in (A3)
and the closure theorem for the stage class used below are taken by statement; only their
content as displayed in this definition is used. The two
one-sided ceilings on $\gamma$ are in force; both are among the
conditions that fix $\gamma_J$:
\begin{equation}\label{5.9:eq-variation-hypotheses-2}
\begin{gathered}
(\mathrm{c\text{-}DR0})\quad 150\,d\,(1+D^{\circ})K'\gamma_0^{\top} \le 1,\\
(\mathrm{c\text{-}DR1})\quad 64\,d\,C^{\mathrm{av}}(1+D^{\circ})\,\mathfrak{A}_w(\gamma)
\le \min\{\alpha_3(\circ),\tfrac14\}.
\end{gathered}
\end{equation}

The chart rests on the following results of Ba{\l}aban, used by statement:
\cite[Proposition~2]{B-CMP98}, \cite[Theorem~1 and (181)]{B-CMP102b} and
\cite[Lemma~4]{B-CMP109}. Bound (ii) is the comb-datum bound of the absolute-chart lemma
with the local product $\omega_u$ kept in place of its global majorant. The step of
that lemma's proof which estimates the comb datum level by level bounds the contribution of
each level $i$ at a cell $x$ meeting the box
by
\begin{equation*}
1.4\,C^{\mathrm{av}}_M\, w(x)\,\tilde\alpha_{0,\ell(x)}\,(\ell_i/\ell_{\ell(x)})^2
\le 2.8\,C^{\mathrm{av}}\, w(x)\,\alpha_{0,\ell(x)}.
\end{equation*}
Only after this does the proof majorise
$w\alpha$ by $\mathfrak{A}_w(\gamma)$. Keeping $w(x)\alpha_{0,\ell(x)} \le \omega_u$ instead
leaves the proof unchanged. Each smallness step remains met, because
$\omega_u \le \mathfrak{A}_w(\gamma)$. These steps are $(\mathrm{c\text{-}DR0})$ for the
transports of the chart that precede the level-by-level estimate, and
$(\mathrm{c\text{-}DR1})$ for the inequality
$1.6\,N\varepsilon^{\circ} \le 0.05$ within that estimate.

\smallskip\noindent
(A2) \emph{The comb dilation is a listed move.} The move list $\mathfrak{M}^{\mathrm{rd}}$
contains the move $\mathrm{M}5(\mathrm{a})(S^0)$. This is the comb dilation through
$\mathrm{M}2^{\circ}$ at the own box $\mathfrak{r}^{\circ}_u$, for the single splitting
$(B^0)$, which we label $(S^0)$, with target $\mathfrak{u}^{\circ}_v$:
\begin{equation}\label{5.9:eq-variation-hypotheses-b}
\begin{gathered}
\bigl((\pi_{\mathrm{M}2},\sigma_0);p\bigr) \longmapsto
\Bigl((U_m,J_m)^{\mathsf{m}}\bigl(\mathfrak{r}^{\circ}_u,e^{i\sigma_0 B^{0,\mathsf{m}}}\bigr)
\big|_{S_u}\Bigr)_{\mathsf{m}},\\
\text{on}\quad \mathfrak{P}_{(S^0)} := \mathfrak{P}_{\mathrm{M}2} \times
\bigl\{ |\sigma_0| < R_0(u) \bigr\},
\qquad R_0(u) := \alpha_3(\circ)/\mathsf{b}_u .
\end{gathered}
\end{equation}
Here $\pi_{\mathrm{M}2} \in \mathfrak{P}_{\mathrm{M}2}$ and $p$ is a point of the stage
class. The majorant $\mathsf{b}_u$ is a scaffold number. It is built from $w$,
$\alpha_{0,\cdot}$, $\ell(\cdot)$, $\mathfrak{B}$ and the box, all of which belong to the
scaffold, so it is the same in both machines. Its pointwise bound is (A1)(ii). The
one-machine bounds attached to the move, $(\mathrm{m}5\mathrm{a})^{\mathrm{sh}}$, are the
three statements (A1)(i), (ii), (iii). The entry bound
$\mathfrak{n}^{\circ}_u$ is included in the supremum defining $\mathfrak{n}_{\hat A}$. The
move has the form of the comb dilations $\mathrm{M}5(\mathrm{a})$, with these choices: the
read $\mathrm{M}2^{\circ}$ (with parameter polydisc $\mathfrak{P}_{\mathrm{M}2}$); the single
splitting $(B^0)$ with majorant $\mathsf{b}_u$; the absolute
radius $\alpha_3(\circ)$; the chart level $m \le j$ in place of $j$; and the smaller own box
$\mathfrak{r}^{\circ}_u$.
Under (A1), (A2) and $(\mathrm{c\text{-}DR1})$ one has
$\mathsf{b}_u \le \tfrac14\alpha_3(\circ)$, hence $R_0(u) \ge 4$; the dilation reads the
chart within $|B'|_\infty < \alpha_3(\circ)$; and $(\mathrm{k}1)$ holds at
$\mathrm{M}5(\mathrm{a})(S^0)$. These are proved below.

\smallskip\noindent
(A3) \emph{Regime near the anchor.} These are the three regime facts of the near-core
analysis. Let $x$ be a $\mathfrak{B}$-cell meeting
$\tilde{\square}^3$ with $\operatorname{dist}_n(x,y) < 1$. Then
$\ell(x) \in [n-1,n+1]$, $w(x) \le w^{\vee}(y)$ and $\alpha_{0,\ell(x)} \le 2\alpha_{0,n}$.
Moreover, by the definition of $c_W$ in the near-core analysis,
\begin{equation}\label{5.9:eq-variation-hypotheses-c}
2\,c_1 K_{\mathrm{loc}} L^3\, w^{\vee}(y)\,\alpha_{0,n} \le \Theta\,\mathfrak{a}_n .
\end{equation}

\smallskip\noindent
(A3)$'$ \emph{Regime at reach three; $\nu = 1$ only.} Set
\begin{equation*}
\begin{gathered}
w^{\vee,3}(y) := \max\{ w(\Delta) : \Delta \text{ a } \mathfrak{B}\text{-cell},\
\operatorname{dist}_n(\Delta,y) < 3 \},\\
c_{\alpha} := 2(1+\beta_0)^4, \qquad \Theta^{(3)} := c_W\, w^{\vee,3}(y).
\end{gathered}
\end{equation*}
Let $x$ be a
$\mathfrak{B}$-cell meeting $\tilde{\square}^3$ with $\operatorname{dist}_n(x,y) < 3$.
Then $\ell(x) \in [n-1,n+3]$ and $\ell(x) \le n_0+1$. Further,
$w(x) \le w^{\vee,3}(y)$ and $\alpha_{0,\ell(x)} \le c_{\alpha}\,\alpha_{0,n}$, and
\begin{equation}\label{5.9:eq-variation-hypotheses-d}
2\,c_1 K_{\mathrm{loc}} L^3\, w^{\vee,3}(y)\,\alpha_{0,n} \le \Theta^{(3)}\,\mathfrak{a}_n .
\end{equation}
The constant $c_{\alpha}$ comes from two further steps of the growth law $\chi_g$,
\begin{equation*}
\alpha_{\cdot,m'} \le (1+\beta_0)^2 (m'-\mu)^{1/2}\,\alpha_{\cdot,\mu}, \qquad
m'-\mu \in \{1,2\},
\end{equation*}
used together with the one-unit level law of the scaffold, from which the range
$\ell(x) \in [n-1,n+1]$ in (A3) is taken, applied twice. Of the clauses of (A3)$'$, the
inequality $w(x) \le w^{\vee,3}(y)$ holds by the definition of $w^{\vee,3}(y)$, since $x$
is among the cells over which that maximum is taken; the inequality
\eqref{5.9:eq-variation-hypotheses-d} is \eqref{5.9:eq-variation-hypotheses-c} with
$w^{\vee}(y)$ replaced by $w^{\vee,3}(y)$ and $\Theta$ by $\Theta^{(3)}$; the level range and
the bound $\alpha_{0,\ell(x)} \le c_{\alpha}\,\alpha_{0,n}$ are those obtained from the
growth law and the level law as just described. Unlike (A3), the statement
(A3)$'$ as a whole is not taken from the near-core analysis; it is assumed here, at
$\nu = 1$.
The near-core analysis, from which (A3) is taken, also contains a transport estimate; it
counts, for each unit $n$-cell $\Delta'$, the factor $w^{\vee}(\Delta')^3$, the cube of the
local weight maximum $w^{\vee}$ taken at the cell $\Delta'$ in place of the anchor $y$, at
reach one. We
do not extend it to reach three here. Such an extension would gain a number of neighbouring
cells depending on
$(d,L)$ only, and the resulting bound would remain of degree two, hence at most three, in
$\Theta$.
\end{definition}

\begin{proof}
We prove the consequences stated at the end of (A2).

First, $\omega_u \le \mathfrak{A}_w(\gamma)$. Indeed, $\omega_u$ is a maximum of products
$w(x)\alpha_{0,\ell(x)}$ over cells $x$. Each such product is at most
$\mathfrak{A}_w(\gamma)$ by the cell-wise bound stated in (A1). Multiplying by the positive
constant $16\,d\,C^{\mathrm{av}}(1+D^{\circ})$ gives
\begin{equation*}
\mathsf{b}_u \le |B^0|^{\circ}_{*} := 16\,d\,C^{\mathrm{av}}(1+D^{\circ})\,\mathfrak{A}_w(\gamma).
\end{equation*}
By $(\mathrm{c\text{-}DR1})$ in \eqref{5.9:eq-variation-hypotheses-2},
$64\,d\,C^{\mathrm{av}}(1+D^{\circ})\mathfrak{A}_w(\gamma) \le \alpha_3(\circ)$. Dividing by
four,
\begin{equation*}
\mathsf{b}_u \le |B^0|^{\circ}_{*} \le \tfrac14\,\alpha_3(\circ).
\end{equation*}
Hence $\alpha_3(\circ) \ge 4\mathsf{b}_u$, that is, $R_0(u) = \alpha_3(\circ)/\mathsf{b}_u \ge 4$.

Second, the move \eqref{5.9:eq-variation-hypotheses-b} reads the chart inside its domain
of analyticity. Let $(\pi_{\mathrm{M}2},\sigma_0) \in \mathfrak{P}_{(S^0)}$, and let $b$ be
a bond of $\mathfrak{B}_m(\mathfrak{r}^{\circ}_u)$. By
\eqref{5.9:eq-variation-hypotheses-a},
\begin{equation*}
|\sigma_0 B^{0,\mathsf{m}}(b)| \le |\sigma_0|\,\mathsf{b}_u < R_0(u)\,\mathsf{b}_u
= \alpha_3(\circ).
\end{equation*}
So $|\sigma_0 B^{0,\mathsf{m}}|_\infty < \alpha_3(\circ)$, and (A1)(iii) applies to
$B' = \sigma_0 B^{0,\mathsf{m}}$ in each machine. Equivalently, the $\sigma_0$-factor of
$\mathfrak{P}_{(S^0)}$ is the disc $\{|\sigma_0|\mathsf{b}_u < \alpha_3(\circ)\}$, which is
the domain prescribed for a comb dilation $\mathrm{M}5(\mathrm{a})$ with one splitting,
majorant $\mathsf{b}_u$ and radius $\alpha_3(\circ)$, as stated in (A2).

Third, the closure condition $(\mathrm{k}1)$ holds at $\mathrm{M}5(\mathrm{a})(S^0)$ by
construction. The stage class is
$\mathcal{K}^{\mathrm{rd}} := \operatorname{cl}^{\mathrm{pr}}(\mathfrak{M}^{\mathrm{rd}})$,
the closure under the moves of the list. By the closure theorem for the stage class,
$(\mathrm{k}1)$ holds at every move of $\mathfrak{M}^{\mathrm{rd}}$. By (A2),
$\mathrm{M}5(\mathrm{a})(S^0)$ is such a move. What the listing requires in addition is
the one-machine bound $(\mathrm{m}5\mathrm{a})^{\mathrm{sh}}$, which is (A1)(i)--(iii).
\end{proof}

\begin{theorem}[The low-order variation bound at the second cube size]\label{5.9:thm-low-order-variation}
Let the site, the block $\square$, the near core $v$ at level $j$ with anchor $y$, the numbers
$n$, $\nu$, $t$, $D$ and the weight $\Theta$, the counterterm companion $\beta[F]\hat{A}$ with
$\hat{A}=A(h_y,\cdot)$, the companion slot $\mathfrak{u}^{\circ}_v$, the stage class
$\mathcal{K}^{\mathrm{rd}}(k;\mathsf{n}+1,\hat{X})$ of source points at the site and the class
$\mathcal{K}^{\mathrm{rd}}(\mathfrak{u}^{\circ}_v)$ of the companion slot, the class-W unit
$u=(1,h_y,j,y)$ with its base-point set $S_u$, chart level $m$, own box $\mathfrak{r}^{\circ}_u$,
chart $(U_m,J_m)(\mathfrak{r}^{\circ}_u,e^{iB'})$ and comb datum $B^0$, and the $\delta$-size
$\mathsf{S}^{\circ}[\hat{A}]$ be as in Definition~\ref{5.9:not-low-depth-cores},
\eqref{5.9:eq-low-depth-cores-b} and \eqref{5.9:eq-low-depth-cores-c}. Assume hypotheses (A1) and
(A2) of Definition~\ref{5.9:def-variation-hypotheses}. Let
$p\in\mathcal{K}^{\mathrm{rd}}(k;\mathsf{n}+1,\hat{X})$ and
$\pi=(b,\sigma,t,\mathsf{t})\in\mathfrak{P}_{\mathrm{M}2}$ (here the first component $b$ and the
third component $t\in[0,1]$ of $\pi$ are parameters of the box-read move; this $t$ is not the scale
ratio $t=L^{j-n}$, which is the meaning of $t$ everywhere outside $\pi$), and let
\begin{equation*}
q:=\mathrm{M}2^{\circ,\times}(\pi;p)
=\bigl(\mathfrak{b}'^{\mathsf{B}}_j\restriction X,\ \mathfrak{b}'^{\mathsf{A}}_j\restriction X\bigr)
\end{equation*}
be the box-pair point. Then the following hold.
\begin{enumerate}
\item[(a)] One has
\begin{equation}\label{5.9:eq-low-order-variation-a}
|\delta\hat{A}(q)|\le\mathsf{S}^{\circ}[\hat{A}]\cdot
\min\Bigl\{1,\ 2\Bigl(\frac{\varkappa_u}{\alpha_3(\circ)}\Bigr)^{2}\Bigr\}.
\end{equation}
\item[(b)] If moreover hypothesis (A3) of Definition~\ref{5.9:def-variation-hypotheses} holds and
$\nu=2$, then
\begin{equation}\label{5.9:eq-low-order-variation-b}
\begin{gathered}
\Bigl(\frac{\varkappa_u}{\alpha_3(\circ)}\Bigr)^{2}
\le\bigl(C_A^{\times}D\Theta\mathfrak{a}_nt^{2}\bigr)^{2},\\
C_A^{\circ}:=\frac{16dC^{\mathrm{av}}(1+D^{\circ})}{c_1K_{\mathrm{loc}}L^{3}\alpha_3(\circ)},
\qquad
C_A^{\times}:=\frac{C_A^{\circ}L^{2}}{5M}.
\end{gathered}
\end{equation}
The constants $C_A^{\circ}$, $C_A^{\times}$ are numbers depending on $(d,L,M)$, given the constants
$c_1$, $K_{\mathrm{loc}}$, $C^{\mathrm{av}}$, $\alpha_3(\circ)$ and $D^{\circ}\le 16ML$; they are
fixed before $\gamma$, and they depend on none of the age parameters, the depth and $K$, on $\nu$
only through $\nu\le 2$, and on no weight and no coupling. If instead hypothesis (A3)$'$ of
Definition~\ref{5.9:def-variation-hypotheses} holds and $\nu=1$, the same bound holds with
$\Theta$ replaced by $\Theta^{(3)}$ and $C_A^{\circ}$ replaced by $C_A^{\circ}c_{\alpha}/2$, that
is, with $C_A^{\times}$ replaced by $C_A^{\circ}c_{\alpha}L^{2}/(10M)$.
\item[(c)] Consequently, under the hypotheses of (b) in the respective case $\nu=2$ or $\nu=1$,
writing $\Theta^{(\cdot)}$ for $\Theta$ when $\nu=2$ and for
$\Theta^{(3)}$ when $\nu=1$, and $C_A^{\times}$ for the corresponding constant of (b), and assuming
the core's counterterm coefficient bound $|\beta[F]|\le C\mathfrak{n}\alpha_{3,j}^{-2}$, with
$C$, $\mathfrak{n}$ and $\alpha_{3,j}$ the constant, the one-machine size bound of the core and
Ba{\l}aban's small-field radius at level $j$ in that bound, the second
term $\beta[F^{\mathsf{A}}]\cdot\delta\hat{A}$ of the splitting \eqref{5.9:eq-low-depth-cores-a}
obeys
\begin{equation}\label{5.9:eq-low-order-variation-c}
|\beta[F^{\mathsf{A}}]|\cdot|\delta\hat{A}(q)|
\le C\mathfrak{n}\alpha_{3,j}^{-2}\cdot\mathsf{S}^{\circ}[\hat{A}]\cdot
\min\bigl\{1,\ 2\bigl(C_A^{\times}D\Theta^{(\cdot)}\mathfrak{a}_nt^{2}\bigr)^{2}\bigr\}.
\end{equation}
\end{enumerate}
Parts (a) and (b) give, with a two-machine constant, the dichotomy for $\delta\hat{A}$ obtained
from part (b) of Lemma~\ref{5.2:lem-disc-bounds} on a chart of radius of order one, which bounds
$|\delta\hat{A}|$ by $2\mathsf{a}$ times the square of a constant multiple of
$D\Theta\mathfrak{a}_nt^{2}$ where the radius of the chart exceeds $1$, and by the value
otherwise: the constant of that dichotomy, which contains the factor $2d/\alpha_3(\circ)$, is
replaced in \eqref{5.9:eq-low-order-variation-a} and \eqref{5.9:eq-low-order-variation-b} by the
pair constant $C_A^{\times}$, and the $\delta$-datum $\mathsf{a}$ by
$\mathsf{S}^{\circ}[\hat{A}]\le\mathsf{a}$. Accordingly the bound
\eqref{5.9:eq-low-order-variation-c}, which carries in addition the coefficient factor
$C\mathfrak{n}\alpha_{3,j}^{-2}$, has the form of the corresponding one-machine bound with
$\mathfrak{n}$ replaced by $\mathfrak{n}\mathsf{a}$, the constant being a two-machine constant.
\end{theorem}

\begin{proof}
\emph{Part (a).} Fix a machine $\mathsf{m}\in\{\mathsf{A},\mathsf{B}\}$ and recall
$R_0(u)=\alpha_3(\circ)/\varkappa_u$.

\emph{Definedness.} Let $\sigma_0\in\mathbb{C}$ with $|\sigma_0|<R_0(u)$. By the pointwise bound
\eqref{5.9:eq-variation-hypotheses-a} of (A1), for every bond $b'$ of
$\mathfrak{B}_m(\mathfrak{r}^{\circ}_u)$,
\begin{equation*}
|\sigma_0B^{0,\mathsf{m}}(b')|\le|\sigma_0|\,\varkappa_u<R_0(u)\varkappa_u=\alpha_3(\circ),
\end{equation*}
so $\sup_{b'}|\sigma_0B^{0,\mathsf{m}}(b')|<\alpha_3(\circ)$. Hence, by the chart clause of (A1), the
function
\begin{equation*}
\varphi^{\mathsf{m}}(\sigma_0):=T_u^{\mathsf{m}}\Bigl((U_m,J_m)^{\mathsf{m}}
\bigl(\mathfrak{r}^{\circ}_u,e^{i\sigma_0B^{0,\mathsf{m}}}\bigr)\restriction S_u\Bigr)
\end{equation*}
is defined on the disc $|\sigma_0|<R_0(u)$. It is holomorphic there as a composition of analytic
maps: $\sigma_0\mapsto\sigma_0B^{0,\mathsf{m}}$ is linear, the chart is analytic in $B'$ for
$|B'|_{\infty}<\alpha_3(\circ)$ by (A1), and $T_u=\hat{A}=\sum_Ph_y(P)\,a(\partial\mathbf{U}(P))$
with $a(V)=\mathfrak{c}_0-\tfrac12\operatorname{tr}(V+V^{-1})$ is a polynomial in the plaquette
variables and their inverses. By the same clause of (A1),
$|\varphi^{\mathsf{m}}(\sigma_0)|\le\mathfrak{n}^{\circ}_u$.

\emph{The value at $\sigma_0=1$.} By the identification clause of (A1),
$\mathsf{P}^{0,\mathsf{m}}\restriction
\mathfrak{r}^{\circ}_u\cong(U_m,J_m)(\mathfrak{r}^{\circ}_u,e^{iB^{0,\mathsf{m}}})$, the
identification holding up to a $G^{c}$-valued gauge map. The functional $T_u$ is invariant under
such maps: a gauge map acts on each plaquette variable $\partial\mathbf{U}(P)$ by conjugation with
an element of $G^{c}$, and $a(V)$ is unchanged when $V$ is replaced by a conjugate of $V$ in
$G^{c}$, because the trace is invariant under conjugation,
so that $a$ is a class function (see~\cite{B-CMP109}). Therefore
$\varphi^{\mathsf{m}}(1)=\hat{A}^{\mathsf{m}}(\mathsf{P}^{0,\mathsf{m}})$. Moreover $\hat{A}$
depends only on the plaquette variables based in $S_u\subset X\cup(y+[-1,1]^d\ell_j)$, and the
box-read move $\mathrm{M}2^{\circ}$ evaluates exactly the restriction of the box pair to these
plaquettes; hence $\varphi^{\mathsf{m}}(1)$ is the value of $\hat{A}^{\mathsf{m}}$ at the
$\mathsf{m}$-component of $q$.

\emph{The value at $\sigma_0=0$.} The configuration $(U_m,J_m)(\mathfrak{r}^{\circ}_u,1)$ is flat:
the minimiser with all block averages equal to $1$ is $\mathbf{U}\equiv 1$, $\mathbf{J}\equiv 0$,
the unique critical orbit of the flat data being flat
(\cite[Proposition~9, (181)]{B-CMP102b}); for $m=0$ this holds by the convention
$(U_0,J_0)(\mathfrak{r},V):=V$. The constant $\mathfrak{c}_0$ in the definition of $a$
in Definition~\ref{5.9:not-low-depth-cores} is normalised so that $a(1)=0$, that is,
$\mathfrak{c}_0=\operatorname{tr}1$, the trace of the identity. Hence
\begin{equation*}
\varphi^{\mathsf{m}}(0)=\sum_Ph_y(P)\,a(1)=0 .
\end{equation*}

\emph{Second-order zero.} The functional $A(c,\cdot)$ depends on the plaquette variables
$\partial\mathbf{U}$ only and not on the current $\mathbf{J}$; so it does not matter that
$J_m$ need not vanish along the family. Put
\begin{equation*}
\partial U_{\sigma_0}(P):=\partial\bigl[U_m\bigl(\mathfrak{r}^{\circ}_u,
e^{i\sigma_0B^{0,\mathsf{m}}}\bigr)\bigr](P).
\end{equation*}
The map $\sigma_0\mapsto\partial U_{\sigma_0}(P)$ is analytic, since
$U_m(\mathfrak{r}^{\circ}_u,e^{i\sigma_0B^{0,\mathsf{m}}})$
is analytic in $\sigma_0$ by the chart clause of (A1), as in the definedness step, and
$\partial U_0(P)=1$ by the
flatness just shown. Hence $\partial U_{\sigma_0}(P)-1=\sigma_0E_P(\sigma_0)$ with $E_P$ analytic on
$|\sigma_0|<R_0(u)$. The identity $V+V^{-1}-2=V^{-1}(V-1)^2$, verified by expanding
$V^{-1}(V^2-2V+1)$, together with $\mathfrak{c}_0=\operatorname{tr}1$, gives
\begin{equation*}
a(V)=-\tfrac12\operatorname{tr}\bigl(V+V^{-1}-2\bigr)
=-\tfrac12\operatorname{tr}\bigl(V^{-1}(V-1)^2\bigr),
\end{equation*}
and therefore
\begin{equation*}
a\bigl(\partial U_{\sigma_0}(P)\bigr)
=-\tfrac12\sigma_0^{2}\operatorname{tr}\bigl(\partial U_{\sigma_0}(P)^{-1}E_P(\sigma_0)^{2}\bigr).
\end{equation*}
Summing against $h_y(P)$,
\begin{equation}\label{5.9:eq-low-order-variation-1}
\varphi^{\mathsf{m}}(\sigma_0)=\sigma_0^{2}\psi^{\mathsf{m}}(\sigma_0),
\qquad
\psi^{\mathsf{m}}(\sigma_0):=-\tfrac12\sum_Ph_y(P)\operatorname{tr}
\bigl(\partial U_{\sigma_0}(P)^{-1}E_P(\sigma_0)^{2}\bigr),
\end{equation}
with $\psi^{\mathsf{m}}$ analytic on $|\sigma_0|<R_0(u)$. In particular
$\varphi^{\mathsf{m}}(0)=\varphi^{\mathsf{m}\prime}(0)=0$ in each machine. No invariance argument
is needed here: the double zero is the double zero of $a$ at $1$.

\emph{Two machines.} By (A2) (see \eqref{5.9:eq-variation-hypotheses-b}), for every
$|\sigma_0|<R_0(u)$ the pair
\begin{equation*}
\Bigl((U_m,J_m)^{\mathsf{B}}\bigl(\mathfrak{r}^{\circ}_u,e^{i\sigma_0B^{0,\mathsf{B}}}\bigr)
\restriction S_u,\ (U_m,J_m)^{\mathsf{A}}\bigl(\mathfrak{r}^{\circ}_u,
e^{i\sigma_0B^{0,\mathsf{A}}}\bigr)\restriction S_u\Bigr)
=\mathrm{M}5(\mathrm{a})(S^0)^{\times}\bigl((\pi,\sigma_0);p\bigr)
\end{equation*}
lies in $\mathcal{K}^{\mathrm{rd}}(\mathfrak{u}^{\circ}_v)$: the move $\mathrm{M}5(\mathrm{a})(S^0)$
is listed by (A2), and by the closure statement for the stage class (the class
$\mathcal{K}^{\mathrm{rd}}$ is by definition the closure under the listed moves, so that the
closure condition at a listed move is met by construction) the pair belongs to
this class. We use the two-machine transfer principle, which states: for a function $f$
holomorphic in each machine on the one-machine class of the companion slot
$\mathfrak{u}^{\circ}_v$, a listed move $\mu$ and a point of the stage class, the difference
between the values of $f^{\mathsf{B}}$ and $f^{\mathsf{A}}$ along the move is holomorphic on the
parameter polydisc $\mathfrak{P}_{\mu}$ of the move and bounded there by the supremum of
$|\delta f|$ over the class $\mathcal{K}^{\mathrm{rd}}(\mathfrak{u}^{\circ}_v)$ of the slot, and
every vanishing of a derivative at a parameter point that holds in each machine holds for the
difference. The principle applies here: the
difference $\delta\varphi:=\varphi^{\mathsf{B}}-\varphi^{\mathsf{A}}$ is holomorphic on the disc
$|\sigma_0|<R_0(u)$ with $|\delta\varphi|\le\mathsf{S}^{\circ}[\hat{A}]$ there, and the per-machine
identities pass to the difference, so that $\delta\varphi(0)=0$ and $\delta\varphi'(0)=0$.

\emph{The Schwarz-lemma bound.} Part (b) of Lemma~\ref{5.2:lem-disc-bounds}, the bounds for a
bounded analytic function on a disc, taken with centre $0$, point
$1$ and a radius $R>1$, states: if $f$ is analytic on the disc $|z|<R$ with
$|f|\le\mathsf{S}^{\circ}[\hat{A}]$ there and $f'(0)=0$, then
$|f(1)-f(0)|\le 2\mathsf{S}^{\circ}[\hat{A}]/R^{2}$. Here $R_0(u)>1$: under (A1) one has
$\varkappa_u\le\frac14\alpha_3(\circ)$ by the ceiling condition on $\gamma$ listed in (A1) of
Definition~\ref{5.9:def-variation-hypotheses}, so
$R_0(u)\ge 4$. Apply the device to $f=\delta\varphi$ with $R=R_0(u)$; since $\delta\varphi(0)=0$,
\begin{equation*}
|\delta\varphi(1)|\le\frac{2\mathsf{S}^{\circ}[\hat{A}]}{R_0(u)^{2}}
=2\mathsf{S}^{\circ}[\hat{A}]\Bigl(\frac{\varkappa_u}{\alpha_3(\circ)}\Bigr)^{2}.
\end{equation*}
By the value at $\sigma_0=1$ established above, $\delta\varphi(1)=\delta\hat{A}(q)$.

\emph{The member $1$ of the minimum.} The move $\mathrm{M}2^{\circ}$ is listed, so by the closure
statement for the stage class the closure condition at $\mathrm{M}2^{\circ}$ is met by construction,
and
$q\in\mathcal{K}^{\mathrm{rd}}(\mathfrak{u}^{\circ}_v)$; by the definition
\eqref{5.9:eq-low-depth-cores-c} of $\mathsf{S}^{\circ}[\hat{A}]$ as a supremum over this class,
$|\delta\hat{A}(q)|\le\mathsf{S}^{\circ}[\hat{A}]$. Both bounds hold, hence their minimum, which
is \eqref{5.9:eq-low-order-variation-a}.

\emph{Part (b).} First, $\mathfrak{r}^{\circ,(1)}_u\subset y+[-2,2]^d\ell_j$. Indeed
$S_u\subset y+[-1,1]^d\ell_j$, and the widths add up as follows:
the passage from $S_u$ to
the level-$m$ set $[S_u]_m$, of which $\mathfrak{r}^{\circ}_u=[S_u]^{(2)}_m$ is the enlargement by
two layers, adds at most $M\ell_m\le\frac14\ell_j$ (by the choice of $a_M$), and the two layers
together with the collar add at most $3M\ell_m\le\frac34\ell_j$, in total at most $\ell_j$; when
$m=0$ the additions are at most $3\eta\le\ell_j$ for $j\ge 1$, and at $j=0$ everything lies in one
cell. Consequently a $\mathfrak{B}$-cell $x$ meeting $\mathfrak{r}^{\circ,(1)}_u$ satisfies
\begin{equation*}
\operatorname{dist}_n(x,y)\le\frac{2\ell_j}{\ell_n}=2L^{1-\nu},
\end{equation*}
where $\operatorname{dist}_n$ is the distance measured in units of $\ell_n$, as in (A3);
this is $<1$ at $\nu=2$, the blocking factor $L$ being an odd integer larger than $2$, and $<3$ at
$\nu=1$; and such a cell meets
$\tilde{\square}^3$, since $\mathfrak{r}^{\circ,(1)}_u\subset\tilde{\square}^3$ by the box clause of
(A1).

\emph{The case $\nu=2$.} Every cell $x$ entering the maximum that defines $\omega_u$ meets
$\tilde{\square}^3$ and has $\operatorname{dist}_n(x,y)<1$, so (A3) (see
\eqref{5.9:eq-variation-hypotheses-c}) gives $w(x)\le w^{\vee}(y)$ and
$\alpha_{0,\ell(x)}\le 2\alpha_{0,n}$, together with
$2c_1K_{\mathrm{loc}}L^{3}w^{\vee}(y)\alpha_{0,n}\le\Theta\mathfrak{a}_n$. Hence
\begin{equation*}
\omega_u\le w^{\vee}(y)\cdot 2\alpha_{0,n}
\le\frac{\Theta\mathfrak{a}_n}{c_1K_{\mathrm{loc}}L^{3}},
\end{equation*}
and, with $\varkappa_u=16dC^{\mathrm{av}}(1+D^{\circ})\omega_u$ from
\eqref{5.9:eq-variation-hypotheses-a},
\begin{equation*}
\frac{\varkappa_u}{\alpha_3(\circ)}
\le\frac{16dC^{\mathrm{av}}(1+D^{\circ})\Theta\mathfrak{a}_n}
{c_1K_{\mathrm{loc}}L^{3}\alpha_3(\circ)}
=C_A^{\circ}\Theta\mathfrak{a}_n .
\end{equation*}
At $\nu=2$ one has $D=M(2\nu+3)=7M$ and $t=L^{j-n}=L^{1-\nu}=L^{-1}$, so $Dt^{2}=7M/L^{2}$ and
\begin{equation*}
\Theta\mathfrak{a}_n=\frac{D\Theta\mathfrak{a}_nt^{2}}{Dt^{2}}
=D\Theta\mathfrak{a}_nt^{2}\cdot\frac{L^{2}}{7M}
\le\frac{L^{2}}{5M}\cdot D\Theta\mathfrak{a}_nt^{2}.
\end{equation*}
Therefore
\begin{equation*}
\frac{\varkappa_u}{\alpha_3(\circ)}\le\frac{C_A^{\circ}L^{2}}{5M}\,D\Theta\mathfrak{a}_nt^{2}
=C_A^{\times}D\Theta\mathfrak{a}_nt^{2};
\end{equation*}
both sides being non-negative, squaring gives
\eqref{5.9:eq-low-order-variation-b}.

\emph{The case $\nu=1$.} Here $D=5M$ and $t=1$, and every cell entering $\omega_u$ meets
$\tilde{\square}^3$ and has $\operatorname{dist}_n(x,y)<3$, so (A3)$'$ (see
\eqref{5.9:eq-variation-hypotheses-d}) gives $w(x)\le w^{\vee,3}(y)$,
$\alpha_{0,\ell(x)}\le c_{\alpha}\alpha_{0,n}$ and
$2c_1K_{\mathrm{loc}}L^{3}w^{\vee,3}(y)\alpha_{0,n}\le\Theta^{(3)}\mathfrak{a}_n$. Hence
\begin{equation*}
\omega_u\le w^{\vee,3}(y)c_{\alpha}\alpha_{0,n}
\le\frac{c_{\alpha}}{2}\cdot\frac{\Theta^{(3)}\mathfrak{a}_n}{c_1K_{\mathrm{loc}}L^{3}},
\end{equation*}
so $\varkappa_u/\alpha_3(\circ)\le(C_A^{\circ}c_{\alpha}/2)\Theta^{(3)}\mathfrak{a}_n$ exactly as in
the case $\nu=2$, and, because $L^{2}\ge 1$,
\begin{equation*}
\Theta^{(3)}\mathfrak{a}_n=\frac{D\Theta^{(3)}\mathfrak{a}_nt^{2}}{5M}
\le\frac{L^{2}}{5M}\,D\Theta^{(3)}\mathfrak{a}_nt^{2}.
\end{equation*}
Squaring gives the bound of (b) with
$\Theta^{(3)}$ in place of $\Theta$ and $C_A^{\circ}c_{\alpha}/2$ in place of $C_A^{\circ}$.

\emph{Part (c).} By the core's counterterm coefficient bound, assumed in (c), the counterterm
coefficient of each core $F$ considered here satisfies $|\beta[F]|\le C\mathfrak{n}\alpha_{3,j}^{-2}$,
in particular $|\beta[F^{\mathsf{A}}]|\le C\mathfrak{n}\alpha_{3,j}^{-2}$. Multiply
\eqref{5.9:eq-low-order-variation-a} by $|\beta[F^{\mathsf{A}}]|$ and insert part (b) inside the
minimum, the map $s\mapsto\min\{1,2s\}$ being non-decreasing on $[0,\infty)$; this gives
\eqref{5.9:eq-low-order-variation-c}.
\end{proof}

\begin{remark}[What the chart factor cancels]\label{5.9:rem-chart-factor}
Let $\nu$, $n$, $t$, $\Theta$ and $D$ be as in \eqref{5.9:eq-low-depth-cores-b}, with $\nu\le 2$, and
let $q$ be a box-pair point as in the low-order variation bound at the second cube size
(Theorem~\ref{5.9:thm-low-order-variation}).
\begin{enumerate}
\item In the second term $\beta[F^{\mathsf{A}}]\cdot\delta\hat A$ of the split
\eqref{5.9:eq-low-depth-cores-a}, the factor $(C_A^{\times}D\Theta^{(\cdot)}\mathfrak{a}_n t^2)^2$ of
\eqref{5.9:eq-low-order-variation-c}, where $\Theta^{(\cdot)}$ stands for $\Theta$ at $\nu=2$ and for
$\Theta^{(3)}$ at $\nu=1$, cancels the factor $\alpha_{3,j}^{-2}$ of the bound on the
counterterm coefficient against $(\Theta\mathfrak{a}_n)^2$, the ratio $\mathfrak{a}_n/\alpha_{3,j}$ being
of the order of $\rho_{n,j}$. The result has the form of the one-machine bound for the counterterm
companion, $C'\mathfrak{n}z^2t^4$ with $z:=\Theta\rho_{n,j}$, with $\mathfrak{n}$ replaced by
$\mathfrak{n}\,\mathsf{S}^{\circ}[\hat A]$ and with a two-machine constant. The value bound
$|\delta\hat A(q)|\le\mathsf{S}^{\circ}[\hat A]$ alone misses this form by $\alpha_{3,j}^{-2}$. That
factor reads the level-$j$ radius from below; it depends on the coupling and is not a constant.
Theorem~\ref{5.9:thm-low-order-variation} supplies the cancellation from a chart, as the bound obtained
from the analyticity of $\hat A$ at radius of order one does at relative depth $\nu\ge 3$.
\item The comb dilation $\mathrm{M}5(\mathrm{a})(S^0)$, whose closure condition is hypothesis (A2),
\eqref{5.9:eq-variation-hypotheses-b}, is a final comb dilation whose radius is fixed by the natural
majorant $\varkappa_u$ of the comb datum of the base point. It is not a dilation at a radius inflated by
a factor exponential in $D$. Adding it to the move list therefore brings in no two-cutoff response
hypothesis. Theorem~\ref{5.9:thm-low-order-variation} does not assert the smallness of
$\mathsf{S}^{\circ}[\hat A]$, for which
\begin{equation*}
\mathsf{S}^{\circ}[\hat A]\le\mathsf{a}\le\mathsf{l}^{J},
\end{equation*}
where $\mathsf{l}^{J}$ denotes the bound on two-machine difference sizes read in the uniqueness theorem.
That smallness is left to be established within the proof of the uniqueness theorem, as is the
smallness of every two-machine difference size.
\item On the stage class $\mathcal{K}^{\mathrm{rd}}$ with $\mathrm{M}5(\mathrm{a})(S^0)$ added to its
move list $\mathfrak{M}^{\mathrm{rd}}$, $\mathsf{a}$ is still finite a priori: $\mathsf{a}\le
2\mathfrak{n}_{\hat A}$. Here the supremum defining $\mathfrak{n}_{\hat A}$ also runs over the
configurations produced by the dilation, and at these the entry bound
\begin{equation}\label{5.9:eq-chart-factor-1}
\mathfrak{n}^{\circ}_u\le 25\,\mathfrak{c}_{\mathrm{tr}}
\max\{\alpha_0^{\circ},\alpha_3(\circ)\}^2\,\|h_y\|_1\,\xi_m^2
\end{equation}
holds, where $\xi_m$ denotes the positive number attached to the chart level $m$ in the entry bound of
the chart construction; it is not related to the unit $u$.
\item The straddling case, $y\in\tilde{\square}^2$ with $\mathfrak{r}_v\not\subset\tilde{\square}^2$,
needs no separate treatment, because $\mathfrak{r}^{\circ}_u\subset\tilde{\square}^3$ in every case.
\end{enumerate}
\end{remark}

\begin{proof}
(1) Part (c) of Theorem~\ref{5.9:thm-low-order-variation} multiplies parts (a) and (b) by the
bound $|\beta[F^{\mathsf{A}}]|\le C\mathfrak{n}\alpha_{3,j}^{-2}$ on the core's counterterm piece. It
gives
\begin{equation*}
|\beta[F^{\mathsf{A}}]|\cdot|\delta\hat A(q)|\le C\mathfrak{n}\,\alpha_{3,j}^{-2}\,
\mathsf{S}^{\circ}[\hat A]\,\min\bigl\{1,\,2(C_A^{\times}D\Theta^{(\cdot)}\mathfrak{a}_nt^2)^2\bigr\}.
\end{equation*}
In the second member of the minimum, the powers of $\alpha_{3,j}$ and $\mathfrak{a}_n$ combine by the
identity
\begin{equation*}
\alpha_{3,j}^{-2}\bigl(C_A^{\times}D\Theta^{(\cdot)}\mathfrak{a}_nt^2\bigr)^2
=(C_A^{\times}D)^2\,t^4\,(\Theta^{(\cdot)})^2\Bigl(\frac{\mathfrak{a}_n}{\alpha_{3,j}}\Bigr)^2 .
\end{equation*}
In the one-machine bound for the counterterm companion, $\mathfrak{a}_n/\alpha_{3,j}$ is of the order of
$\rho_{n,j}$. At $\nu=2$, where $\Theta^{(\cdot)}=\Theta$, $D=7M$ and $t=L^{-1}$, the right-hand side is
therefore of the order of $(C_A^{\times}D)^2z^2t^4$; at $\nu=1$, where $D=5M$ and $t=1$, the same holds
with $\Theta^{(3)}$ and $C_A^{\circ}c_{\alpha}/2$ in place of $\Theta$ and $C_A^{\circ}$. Here
$C_A^{\times}$ is one of the $(d,L,M)$-numbers of \eqref{5.9:eq-low-order-variation-b}.

If only the first member of the minimum is used, which is the value bound, the right-hand side is
$C\mathfrak{n}\alpha_{3,j}^{-2}\mathsf{S}^{\circ}[\hat A]$. Compared with $C'\mathfrak{n}z^2t^4$, the
factor $\alpha_{3,j}^{-2}$ is left uncompensated.

The second member is produced from a chart. In the proof of
Theorem~\ref{5.9:thm-low-order-variation}(a), the Schwarz-lemma device is applied to the two-machine
difference of the dilated family on the disc of radius $R_0(u)=\alpha_3(\circ)/\varkappa_u$ of the
absolute chart. This difference has a double zero at the origin, which yields the factor
$2(\varkappa_u/\alpha_3(\circ))^2$ of \eqref{5.9:eq-low-order-variation-a}. Part (b),
\eqref{5.9:eq-low-order-variation-b}, bounds $(\varkappa_u/\alpha_3(\circ))^2$ by
$(C_A^{\times}D\Theta\mathfrak{a}_nt^2)^2$, so the second member of the minimum is at most
$2(C_A^{\times}D\Theta\mathfrak{a}_nt^2)^2$.

(2) In the proof of Theorem~\ref{5.9:thm-low-order-variation}(a), the dilation is
\begin{equation*}
\sigma_0\mapsto(U_m,J_m)^{\mathsf{m}}\bigl(\mathfrak{r}^{\circ}_u,e^{i\sigma_0B^{0,\mathsf{m}}}\bigr),
\qquad |\sigma_0|<R_0(u),
\end{equation*}
and by (A1)(ii), \eqref{5.9:eq-variation-hypotheses-a},
\begin{equation*}
\sup_b|\sigma_0B^{0,\mathsf{m}}(b)|<R_0(u)\varkappa_u=\alpha_3(\circ).
\end{equation*}
The radius of the dilation is thus determined by the natural majorant $\varkappa_u$ and the chart
radius $\alpha_3(\circ)$ alone. No factor exponential in $D$ enters. By the statement classifying final
comb dilations, a final comb dilation whose radius is the natural majorant of the comb datum of its base
point is of the kind that requires no two-cutoff response estimate, in contrast with dilations at an
inflated radius; hence adding $\mathrm{M}5(\mathrm{a})(S^0)$ to the move list brings in no two-cutoff
response hypothesis. Theorem~\ref{5.9:thm-low-order-variation} bounds $|\delta\hat A(q)|$ only through
$\mathsf{S}^{\circ}[\hat A]\le\mathsf{a}$, and it says nothing about the smallness of that quantity.

(3) By part (iii) of the theorem on the stage class $\mathcal{K}^{\mathrm{rd}}$, the bound
$\mathsf{a}\le 2\mathfrak{n}_{\hat A}$ holds on the class closed under any list of moves satisfying the
closure condition $(\mathrm{k}1)$, and $(\mathrm{k}1)$ holds at $\mathrm{M}5(\mathrm{a})(S^0)$ by
hypothesis (A2). Adding $\mathrm{M}5(\mathrm{a})(S^0)$ to $\mathfrak{M}^{\mathrm{rd}}$ enlarges the set of
configurations over which the supremum defining $\mathfrak{n}_{\hat A}$ is taken. At the new
configurations, the entry bound \eqref{5.9:eq-chart-factor-1} of the chart construction applies, so
$\mathfrak{n}_{\hat A}$, and with it $\mathsf{a}$, stays finite.

(4) The inclusion $\mathfrak{r}^{\circ}_u\subset\tilde{\square}^3$ is part of hypothesis (A1). It holds
without any condition on the position of $\mathfrak{r}_v$ relative to $\tilde{\square}^2$. The proof of
Theorem~\ref{5.9:thm-low-order-variation} reads only $\mathfrak{r}^{\circ}_u$ and the cells meeting its
one-layer collar $\mathfrak{r}^{\circ,(1)}_u$, and each of these cells meets $\tilde{\square}^3$ by the
same part of (A1).
\end{proof}

\begin{remark}[The boundary companion's hypothesis holds by definition]\label{5.9:rem-boundary-companion}
The hypothesis of the marginal-weight lemma without a box, the boundary companion of the marginal
bound on a domain's own box (Lemma~\ref{5.2:lem-marginal-bound}), is not a statement of its own:
at every rim evaluation at the sites (L) and (T) of Definition~\ref{5.2:def-sites-units} at which
that lemma is used, it follows from two hypotheses of the uniqueness theorem, taken as stated (that
each map producing an evaluation point is defined and holomorphic on the product of its complex
parameter polydisc with the domain of the object from which it issues, and takes this product into
the domain of the object it evaluates), together with \cite[(1.64)--(1.69)]{B-CMP122a} and
\cite[clauses (i)--(iii) and (1.12)]{B-CMP109}, read at the evaluation point as the definition of
those domains; the bound $\mathfrak{a}^{\partial}$ required there of the gauge potential, of its
gradient and of the current is the right member of the defining clause at the cube of the family of
\cite[(1.64)--(1.69)]{B-CMP122a} that contains the localisation domain, and no numerical value is
assigned to it; when the localisation domain meets several cubes of that family, a single gauge is
first assembled cube by cube from the pointwise plaquette clauses \cite[(1.66), (1.68)]{B-CMP122a},
by the lemma that builds one gauge on a union of cubes of that family from these clauses, and
$\mathfrak{a}^{\partial}$ is then the
maximum of those right members over the levels met; no step of the proof of the correlation theorem
uses this hypothesis, which enters only the proof of the uniqueness theorem.
\end{remark}

\begin{remark}[Ba{\l}aban's bounds for old terms beside a new large-field region (printed without
proof in {\cite[p.~375]{B-CMP122b}}; cited as printed)]\label{5.10:rem-old-terms-beside}
The old terms of the effective action standing beside a newly created large-field region keep
their analyticity, and with it their inductive bounds, on the space of values of the new
determining set. This is the space
$\tilde U''^c_k(Y_1,\tilde\alpha_0,\tilde\alpha_1)$ of values of the extended function $U^0_k$ of
\cite[(1.65)]{B-CMP122b}. According to that display, $U^0_k$ is analytic in
$(\mathbf{U},\mathbf{J})\in\tilde U^c_k(Y_4,\tilde\alpha_0,\tilde\alpha_1)$. On
\cite[p.~375]{B-CMP122b} this space of values is stated to lie in the analyticity domain of the
term $\mathbf{E}(X)$ considered there, and hence in the analyticity domains of all terms of the
effective action. For the old terms, that statement is the sentence quoted below. Composed with
$U^0_k$, each old term is therefore analytic in $(\mathbf{U},\mathbf{J})$ on
$\tilde U^c_k(Y_4,\tilde\alpha_0,\tilde\alpha_1)$. Together with the same containment for the other
terms of the effective action, this gives the analyticity of $\mathbf{E}(X,U^0_k)$ with which the
quoted passage closes. The analyticity of the old terms is an input of the weight lemma and of
the treatment of boundary terms beside large-field regions in the correlation theorem
(Theorem~\ref{5.1:thm-correlation-theorem}). It is printed without proof in
\cite[p.~375]{B-CMP122b} and is cited as printed:
\begin{quotation}
It is also clear for all new terms created by the previous localization operations. Finally, it
is simple to see for all the old terms, because they are connected with the old determining set,
for which the regularity conditions on the crucial domain
$\Omega''_{k_0+1}\setminus\Omega_{k_0+1}$ are weaker than the regularity conditions for the new
determining set. This implies the basic fact that the function $\mathbf{E}(X,U^0_k)$ is an
analytic function of the variables $(\mathbf{U},\mathbf{J})$ on the space
$\tilde U^c_k(Y_4,\tilde\alpha_0,\tilde\alpha_1)$.
\end{quotation}
The same step also uses the independence from $M$ of the kernel constants and of the decay rate
of the random walk. This is not a printed sentence; it is
Proposition~\ref{5.10:prop-m-freeness} below (stated; proof incomplete). The printed block
\cite[(1.23)--(1.48)]{B-CMP122b} contains the preparatory estimates of the large-field analysis.
It is likewise cited as printed wherever this chapter uses it. It is not among the four sentences
of Ba{\l}aban printed without proof that are listed in the index closing Part~III.
\end{remark}

\begin{proposition}[$M$-freeness of the random-walk constants, companion of the bounds for old terms beside a new region]\label{5.10:prop-m-freeness}
Consider backgrounds given only as local data on the cells of the one family of cells, at the
decomposition of the density at which the following bounds are applied, and let $M$ range over its
values above its floors. For such backgrounds, the kernel $\mathcal{K}$ of the conditional
fluctuation operator satisfies
\begin{equation}\label{5.10:eq-m-freeness-1}
|\mathcal{K}(x,y)|\le B_W e^{-\delta_W|x-y|},
\end{equation}
and the companion kernel of the random-walk expansion of that operator satisfies its exponential
bound with decay rate $\mu_0$. Here $(B_W,\delta_W)$ and $\mu_0$ are the constants of Stage~5 of the
order of choice of Chapter~\ref{ch:2}, that is, the $M$-free constants of the same two bounds for
backgrounds in the global regularity class.
\end{proposition}

Stated; the proof below is incomplete at the step marked $(\ast)$.

Proposition~\ref{2.4:prop-m-freeness} is the same statement at the place where the order of choice
of Chapter~\ref{ch:2} reads it; its proof is not repeated there, and the proof below, with its step
marked $(\ast)$, is given only here.

\begin{proof}
We first treat backgrounds in the global regularity class. For these, at the same decomposition of
the density, the kernel $\mathcal{K}$ of the conditional fluctuation operator satisfies the bound
\eqref{5.10:eq-m-freeness-1} with the constants $(B_W,\delta_W)$. Moreover, the companion kernel of
the random-walk expansion satisfies its exponential bound with decay rate $\mu_0$. Both bounds hold
by the operator theorem at free current, and the constants $(B_W,\delta_W)$ and $\mu_0$ obtained
there do not depend on $M$.

For the pair $(B_W,\delta_W)$ the passage to backgrounds given only as local data is supplied by the
operator theorem at free current read at arbitrary sub-families of the cells of the one family:
read in this way, it gives the bound \eqref{5.10:eq-m-freeness-1} with the constants
$(B_W,\delta_W)$ for such backgrounds, for every $M$ above its floors.

$(\ast)$ It remains to show that the companion kernel of the random-walk expansion satisfies its
exponential bound with the same decay rate $\mu_0$ for backgrounds given only as local data on the
cells of the one family, for every $M$ above its floors.
\end{proof}

The proposition is needed in the order of choice of Chapter~\ref{ch:2}: the floors of the width
parameters and the three floors on $M$ of Stage~6 are written in the constants of Stage~5, whereas
the bounds they serve are applied at local data, and if the constants of these bounds at local data
depended on $M$, the three floors on $M$ would compare $M$ with an unspecified function of $M$ and
their side would not be determined.

\begin{remark}\label{5.11:lem-sixth-family}
The statement and proof of this lemma are not written out here.
\end{remark}

\begin{remark}[A question on ejected components]\label{5.11:rem-ejected-components}
One question about Ba{\l}aban's large-field operation $\mathbf{R}$ is left undecided in
this chapter, and no statement of this book depends on its answer.
Suppose the operation $\mathbf{R}$ at a level $k \ge 1$ returns a component $Y$ of its
domain to the large-field region of the new action. Such a component is one of the second
class in \cite[p.~378]{B-CMP122b}, that is, one on which the large-field characteristic
function of the decomposition of unity is selected. The question is whether the values of
the block field $V_k$ that the operation holds fixed on the cells of $Y$ are regular at
their own scale, that is, lie within a fixed multiple of the small-field threshold
$\varepsilon_k$ of level $k$, or whether they may be rough.
On the reading of \cite[p.~378]{B-CMP122b} together with
\cite[(1.74)--(1.75), Proposition~1]{B-CMP122a}, the characteristic functions that the
operation keeps on such a component force $|V_k(\partial p) - 1| < C\varepsilon_k$ at every
plaquette $p$ of $Y$, with a constant $C$ depending only on $d$, $L$ and $M$. The data are
then regular at their own scale, and no further hypothesis is needed to answer the
question; on this reading the affirmative answer is proved as an implication from the
printed sentences just cited, read by statement.
On the stricter reading, the displayed sentences just cited do not by themselves settle the
question, and the affirmative answer is not proved.
Neither reading is preferred here. Since no statement of this book uses the answer, the
question is stated only so that the list in Chapter~\ref{ch:13} is complete.
\end{remark}

\begin{proof}[Proof of the torus bound]\label{5.12:prf-torus-bound-proof}
The torus bound of the correlation theorem (Theorem~\ref{5.1:thm-correlation-theorem}) is derived
here as an implication. The statements
entering as hypotheses are the following.
\begin{enumerate}
\item The induction format with sources, Definition~\ref{5.2:def-induction-format}, at every
level $k\le K'$ of the local run (the renormalisation steps $k\le K'$, performed on the
physical torus), in particular its clauses (c), (d) and (f). We also use the
reference bound on the running inverse coupling, from the reference recursion of
Definition~\ref{5.1:def-reference-recursion},
$x_{K'}\le g^{-2}+b_++3\lambda m+2c_2$.
\item The sourced versions, on real and on complex backgrounds, of the bounds
\cite[(2.43)--(2.46)]{B-CMP119}, given by the last column of the budget at each site
(Definition~\ref{5.2:def-comparison-constant}), whose constant is $C_\sharp$, and the sup-norm
bounds \cite[(2.47)--(2.48)]{B-CMP119}.
\item Clause (c) of the lemma on complex plaquette weights in the expansions
(Lemma~\ref{5.2:lem-complex-weights}), with its constant
$\theta_W$, the ratio of complex to real weight.
\item The per-zone-cell value bound, with its constant $\mathsf{K}_{\mathrm{cell}}(k,n)$, and
the volume ceiling on $\gamma$ that bounds this constant.
\item Clause (c) of the modification bound for arbitrary regions on the flat domain
class (the class of localisation domains attached to the flat, coupling-dependent radii), that
is, the modification bound with the additional term
$-\kappa_{\mathfrak{d}}d_k(X)$.
\item The bound for the gas of live large-field components on the flat domain class, together
with the polymer bound $\Xi\le\exp(e^{-p_0(g_k)}|\mathbb{L}_k|)$, valid under the two ceilings
on $\gamma$. Here $\Xi$ is the partition function of that gas at level $k$, $\mathbb{L}_k$ is
the set of level-$k$ cubes over which the live components range, and $p_0(\cdot)$ is the degree
function.
\item The counting lemma, which controls the number $\#\bigl((\Lambda^{*})^{c}\bigr)$ of elements
of the complement of the set $\Lambda^{*}$ entering the indexing of the histories by the labels
$I^*$.
\item The floor for the unsourced integral at the level $K'$.
\item Clause (c) of the lemma indexing the histories, and the payment per $I$-cube with its
half factors.
\end{enumerate}

The local run is stopped at the level $K'=K-m$ of Definition~\ref{5.2:def-auxiliary-torus}, on
the physical torus. Its lattice $T^{(K')}$
has $N_{\mathrm{side}}$ sites per side; the letter $N$ is reserved for the group
$\mathrm{SU}(N)$. Throughout, $w=\sum_i z_if_i$ with $z$ in the ball
$\mathbb{D}$ of
admissible sources; $\mathbb{D}$, the number $r$ and the constant $\mathsf{Q}$ are those of the
torus bound in Theorem~\ref{5.1:thm-correlation-theorem}.

We begin by separating off the following terms. For $j\le K'$ and a localisation
domain $X$ at level $j$, put
\begin{equation}\label{5.12:eq-torus-bound-proof-1}
e^{(j)}(X;w):=\sum_{y\in X}E^{(j)}(X,1,y;w)+R^{(j)}(X,1;w),
\end{equation}
to which, when $X$ is the cube of a point $y$ at level $j$, the further term
$\log[\mathfrak{z}(a_j(y))/\mathfrak{z}(x_j)]$ is added.
The first two members are built from the $E$- and $R$-terms of the induction format at source $w$
(Definition~\ref{5.2:def-induction-format}); the further term on the cube of $y$ comes from the
one-bond Haar integrals
$\mathfrak{z}$ of the gauge-fixing move; here $a_j(y)$ is the argument at which the one-bond
integral $\mathfrak{z}$ is evaluated on the cube of $y$ at level $j$, and $x_j$ is the running
inverse coupling of level $j$, with which it is compared. Define
\begin{equation}\label{5.12:eq-torus-bound-proof-2}
e_j(X;z):=e^{(j)}(X;w)-e^{(j)}(X;0).
\end{equation}
By definition $e_j(X;0)=0$. Clauses (c), (d) and (f) of the induction format with sources
(Definition~\ref{5.2:def-induction-format}) give three further properties:
\begin{itemize}
\item $e_j(X;\cdot)$ is analytic on $\mathbb{D}$;
\item $e_j(X;\cdot)$ does not depend on $z_i$ unless $X^+\cap\operatorname{supp}f_i\neq\emptyset$,
where $X^+$ is the enlargement of $X$;
\item $e_j(X;\cdot)$ obeys, with the constant $C_e$ of the bounds of the induction format with
sources (Definition~\ref{5.2:def-induction-format}), the bound
\begin{equation}\label{5.12:eq-torus-bound-proof-3}
|e_j(X;z)|\le C_e\,e^{-(1-\delta)\kappa d_j(X)},\qquad z\in\mathbb{D}.
\end{equation}
\end{itemize}
Put $\mathcal{E}(z):=\sum_{j\le K'}\sum_Xe_j(X;z)$; the sum is finite on the torus. Define
$\tilde Z$ by $Z=e^{\mathcal{E}}\tilde Z$.

We next represent $\tilde Z$. By the induction format with sources at $k=K'$, $Z$ is the
integral over the variables $V_{K'}$ of $\rho^w_{K'}$. Here $\rho^w_{K'}$ is the sum over
histories $\{\Omega_j,\Lambda_j,I_j\}$ of $\chi_{K'}\mathbf{T}^w_{K'}\exp A^w_{K'}$. Hence
\begin{equation}\label{5.12:eq-torus-bound-proof-4}
\tilde Z(z)=\int dV_{K'}\sum\chi_{K'}\mathbf{T}^w_{K'}\exp\bigl[A^w_{K'}-\mathcal{E}\bigr].
\end{equation}
In a given history, the constant term of Ba{\l}aban's expansion appears together with
$-\sum_j\sum_{X\not\subset\Lambda_j}e_j(X;z)$, and this additional sum has modulus at most
$O(1)\sum_j|Z_j|$, where $Z_j$ are the large-field regions of the history.

The histories are indexed, at each step, by a pair consisting of a region and a label $I^*$;
this is clause (c) of the lemma indexing the histories. The additional
sum over the labels $I^*$ at step $j$ is paid for by the half factors reserved in the payment
made per $I$-cube.
These half factors lie inside the factor $\exp O(1)(MR_j)^{-d}|Z_j|$ of
\cite{B-CMP122b}, where $R_j$ denotes Ba{\l}aban's parameter $R$ at step $j$. The count
$\#\bigl((\Lambda^{*})^{c}\bigr)$ is controlled by the counting lemma. Together,
these two payments change the constants $E_+^{\sharp}$ and $C_{\mathrm{low}}$ below only by
numbers depending on $d$ and $L$.

We now bound the exponent. The sourced versions of \cite[(2.43)--(2.46)]{B-CMP119} on real and
complex backgrounds, in the last column of the budget at each site
(Definition~\ref{5.2:def-comparison-constant}), the sup-norm bounds
\cite[(2.47)--(2.48)]{B-CMP119} and clause (c) of
Lemma~\ref{5.2:lem-complex-weights} together give
\begin{equation}\label{5.12:eq-torus-bound-proof-a}
\begin{split}
\operatorname{Re}A^w_{K'}\le{}&-(1-\theta_W)\,A(x_{K'}(\cdot),U_{K'})
+(\text{logarithmic terms})\\
&+\sum_j\mathsf{K}_{\mathrm{cell}}(K',j)\,|\Gamma_j|
+\sum_{\mathfrak{a}\ \text{live at}\ K'}\mathsf{P}(\mathfrak{a}).
\end{split}
\end{equation}
Here $\mathsf{K}_{\mathrm{cell}}(K',j)$ is the per-zone-cell constant of the per-zone-cell value
bound, and the $\Gamma_j$ are the regions of \cite[(2.49)]{B-CMP119}, the zones. The constant
$\mathsf{K}_{\mathrm{cell}}(K',j)$ covers the
classes of terms listed in the per-zone-cell value bound under its first, second, third and
fifth items, its item of boundary-type terms and its item of the boundary terms $\mathbf{B}$,
together with the chain terms of the last step. By the volume ceiling on $\gamma$,
\begin{equation}\label{5.12:eq-torus-bound-proof-5}
\mathsf{K}_{\mathrm{cell}}(K',j)\le 2^{-d}\bigl(C_{\sharp}O(1)+1\bigr)\log g_j^{-2}
\qquad\text{per cell.}
\end{equation}
This is one more logarithmic term.

For the zones $j<K'$, which carry large-field factors, this term is controlled together with
the other logarithmic terms. In \cite[(2.49)]{B-CMP119} the corresponding member is
$O(1)\sum_j|\Gamma_j|$, and the first term of that bound also controls the logarithmic terms
there. The constants $E_\pm$ of \cite[Corollary~3]{B-CMP119} depend on $g_k$, and so do the
corresponding constants in \cite{B-CMP122b}.

The zone $j=K'$ has $|\Gamma_{K'}|$ of order $N_{\mathrm{side}}^4$ and no large-field factor.
There the term
joins $E_+^{\sharp}$ as the member $2^{-d}(C_\sharp O(1)+1)\mathsf{T}_{K'}$. This member reads
$(g,m)$ through
\begin{equation}\label{5.12:eq-torus-bound-proof-6}
\mathsf{T}_{K'}=\log g_{K'}^{-2}\le\log(g^{-2}+b_++3\lambda m+2c_2).
\end{equation}
This inequality is the reference bound on $x_{K'}=g_{K'}^{-2}$, read through the logarithm,
exactly as for $C_{\mathrm{low}}$ below.

In the last sum of \eqref{5.12:eq-torus-bound-proof-a}, $\mathfrak{a}$ runs over the arrests,
that is, the arrested chains of the large-field expansion, that are still live at the level
$K'$, that is, still carry their large-field factor at that level; for such an arrest
$\mathfrak{a}$, $\mathsf{P}(\mathfrak{a})$ denotes its frozen pile, the contribution of the
frozen left-over terms of $\mathfrak{a}$. Clause $(\beta')$ of the per-zone-cell value
bound describes it. It is extensive, each of its terms being counted at that term's own scale.
It is absorbed in the large-field volume charge $O(1)M^dR_j^{d+1}d'_j(Z_j)$ of the gas of live
components, where $d'_j(Z_j)$ is the size of $Z_j$ measured in cubes of side $MR_j$, with the
constant $2c_{\mathsf{P}}M^{-d}$; here $c_{\mathsf{P}}$ is the constant of the frozen pile in the
per-zone-cell value bound, and $2c_{\mathsf{P}}M^{-d}$ is a number once $M$ is fixed.
Accordingly, the charge of the bound for the gas of live large-field components contains it, by
the charge clause of that bound. A
live arrest lies in a large-field region of the history contained in $\Lambda^c_{K'}$, never in
the top zone. Therefore the upper bound
on $|\tilde Z|$ below holds with $E_+^{\sharp}$ changed only by the one member reading $(g,m)$
named above.

We turn to the upper bound. We use clause (c) of the modification bound for arbitrary regions
on the flat domain class, that is, the modification bound with the additional term
$-\kappa_{\mathfrak{d}}d_k(X)$; here $\kappa_{\mathfrak{d}}$ is the coefficient of the additional
term of \cite{B-CMP122b}, denoted there $\kappa_1$.

Ba{\l}aban proves the upper member of \cite[(2.50)]{B-CMP119} by bounding the moduli of the
terms. We follow that proof, with the moduli integrated over the variables $V_{K'}$ by the device
for integrating bounds on moduli over these variables. On the
flat domain class the gas of live components is controlled by the bound for the gas of live
large-field components, together with the polymer bound for $\Xi$, which holds under
the two ceilings on $\gamma$. Combining this with \eqref{5.12:eq-torus-bound-proof-4},
\eqref{5.12:eq-torus-bound-proof-a} and the history bound above gives
\begin{equation}\label{5.12:eq-torus-bound-proof-7}
|\tilde Z(z)|\le e^{E_+^{\sharp}N_{\mathrm{side}}^4},\qquad z\in\mathbb{D}.
\end{equation}

For the floor, note that $\mathcal{E}(0)=0$, since every $e_j(X;0)$ vanishes. The floor for the
unsourced integral then gives
\begin{equation}\label{5.12:eq-torus-bound-proof-8}
\tilde Z(0)=\int dV_{K'}\,\rho^0_{K'}\ge e^{-C_{\mathrm{low}}N_{\mathrm{side}}^4}.
\end{equation}
The constant $C_{\mathrm{low}}$ reads the coupling through the reference bound
$x_{K'}\le g^{-2}+b_++3\lambda m+2c_2$.

We now estimate the derivatives. Put $\mathsf{h}:=\tilde Z/\tilde Z(0)$. Dividing
\eqref{5.12:eq-torus-bound-proof-7} by \eqref{5.12:eq-torus-bound-proof-8} gives
\begin{equation*}
|\mathsf{h}|\le e^{(E_+^{\sharp}+C_{\mathrm{low}})N_{\mathrm{side}}^4}\le\mathsf{Q}
\qquad\text{on }\mathbb{D},
\end{equation*}
the constant $\mathsf{Q}$ of the torus bound in Theorem~\ref{5.1:thm-correlation-theorem}
dominating $e^{(E_+^{\sharp}+C_{\mathrm{low}})N_{\mathrm{side}}^4}$.

Since $Z=e^{\mathcal{E}}\tilde Z(0)\,\mathsf{h}$, wherever $\operatorname{Log}\mathsf{h}$ is
defined the logarithm of $Z$ equals $\mathcal{E}+\log\tilde Z(0)+\operatorname{Log}\mathsf{h}$.
Its mixed derivative $\partial_{z_1}\cdots\partial_{z_n}$ at $z=0$ is therefore the sum of the
corresponding derivatives of $\operatorname{Log}\mathsf{h}$ and of $\mathcal{E}$. We estimate
them in turn.

($\alpha$) The derivatives of $\operatorname{Log}\mathsf{h}$. The function $\mathsf{h}-1$
vanishes at $0$ and is bounded by $\mathsf{Q}+1$ on $\mathbb{D}$. Fix $z\in\mathbb{D}$. Since
$\mathbb{D}$ is a ball about $0$, the map $\varsigma\mapsto\mathsf{h}(\varsigma z)-1$ is analytic
for complex $\varsigma$ of modulus less than one, vanishes at $\varsigma=0$ and is bounded by
$\mathsf{Q}+1$. The Schwarz lemma on
this complex line gives $|\mathsf{h}(\varsigma z)-1|\le(\mathsf{Q}+1)|\varsigma|$. Hence
\begin{equation}\label{5.12:eq-torus-bound-proof-9}
|\mathsf{h}-1|\le\tfrac12\quad\text{on }(2\mathsf{Q}+2)^{-1}\mathbb{D}.
\end{equation}
There $\operatorname{Log}\mathsf{h}$, the principal branch, is analytic. Its power series in
$\mathsf{h}-1$ gives
\begin{equation}\label{5.12:eq-torus-bound-proof-10}
|\operatorname{Log}\mathsf{h}|
\le\sum_{\ell\ge1}\frac{|\mathsf{h}-1|^\ell}{\ell}
\le\sum_{\ell\ge1}\frac{2^{-\ell}}{\ell}
=\log2 .
\end{equation}

We apply the Cauchy estimate on the circles
$|z_i|=(1-\upsilon)r/((2\mathsf{Q}+2)n\|f_i\|_\sharp)$, $i=1,\dots,n$, with $0<\upsilon<1$. For
$\upsilon>0$ these circles lie in $(2\mathsf{Q}+2)^{-1}\mathbb{D}$; at $\upsilon=0$ the polyradii
lie on its boundary. The estimate gives
\begin{equation}\label{5.12:eq-torus-bound-proof-11}
\bigl|\partial_{z_1}\cdots\partial_{z_n}\operatorname{Log}\mathsf{h}(0)\bigr|
\le\log2\,\prod_{i=1}^n\frac{(2\mathsf{Q}+2)\,n\,\|f_i\|_\sharp}{(1-\upsilon)\,r}.
\end{equation}
Letting $\upsilon\downarrow0$ removes the factors $(1-\upsilon)^{-1}$. Alternatively, the
closure of $(2\mathsf{Q}+2)^{-1}\mathbb{D}$ lies in $\mathbb{D}$, and the set where
$|\mathsf{h}-1|\le\frac12$ is closed. So \eqref{5.12:eq-torus-bound-proof-9} holds on that
closure, and one may take $\upsilon=0$ directly. The constant is the same either way.

($\beta$) The derivatives of $\mathcal{E}$. Since $e_j(X;\cdot)$ does not depend on $z_i$
unless $X^+$ meets $\operatorname{supp}f_i$, the derivative
$\partial_{z_1}\cdots\partial_{z_n}e_j(X;0)$ vanishes unless $X^+$ meets every
$\operatorname{supp}f_i$. In that case
$4M\mathsf{t}_j(d_j(X)+3)\ge\mathsf{d}$, where $\mathsf{t}_j:=L^j\varepsilon$ and $\varepsilon$
is the lattice spacing. Since $d_j(X)\ge0$, this means
$d_j(X)\ge[\mathsf{d}/(4M\mathsf{t}_j)-3]_+$.

For each such $X$ we apply the Cauchy estimate to $e_j(X;\cdot)$ on the circles
$|z_i|=(1-\upsilon)r/(n\|f_i\|_\sharp)$, using the bound \eqref{5.12:eq-torus-bound-proof-3},
and let $\upsilon\downarrow0$. This gives
\begin{equation}\label{5.12:eq-torus-bound-proof-12}
\bigl|\partial_{z_1}\cdots\partial_{z_n}e_j(X;0)\bigr|
\le C_e\,e^{-(1-\delta)\kappa d_j(X)}\prod_{i=1}^n\frac{n\|f_i\|_\sharp}{r}.
\end{equation}
We anchor each contributing $X$ at one of the at most $1+(M\mathsf{t}_j)^{-4}$ cubes of level
$j$ that meet $(\operatorname{supp}f_1)^+$. We then sum over $X$ with
\cite[(1.26)]{B-CMP116}, under the constraint
$d_j(X)\ge[\mathsf{d}/(4M\mathsf{t}_j)-3]_+$. After summation over $j\le K'$ this gives
\begin{equation}\label{5.12:eq-torus-bound-proof-13}
\bigl|\partial_{z_1}\cdots\partial_{z_n}\mathcal{E}(0)\bigr|
\le C_{\mathrm{UV}}(\mathsf{d})\prod_{i=1}^n\frac{n\|f_i\|_\sharp}{r},
\end{equation}
where
\begin{equation}\label{5.12:eq-torus-bound-proof-14}
C_{\mathrm{UV}}(\mathsf{d}):=C_e\,O(1)\sum_{j\le K'}\bigl(1+(M\mathsf{t}_j)^{-4}\bigr)
e^{-\frac14\kappa[\mathsf{d}/(4M\mathsf{t}_j)-3]_+}.
\end{equation}

It remains to bound $C_{\mathrm{UV}}(\mathsf{d})$ independently of $K$. The numbers
$\mathsf{t}_j$ form a
geometric sequence of ratio $L$. The summands of \eqref{5.12:eq-torus-bound-proof-14} with
$M\mathsf{t}_j\le\mathsf{d}/24$, prefactor $C_e\,O(1)$ included, sum to at most
$C\mathsf{d}^{-4}$, with a constant $C$ that does
not depend on $K$. Write
$\log_L^+x:=\max\{\log_Lx,0\}$. The remaining terms, those with $M\mathsf{t}_j>\mathsf{d}/24$,
number at most $1+\log_L^+(24M/(\mathsf{d}N_{\mathrm{side}}))$. Each of them is at most
$C_e\,O(1)(1+(24/\mathsf{d})^4)$, since for them $(M\mathsf{t}_j)^{-4}<(24/\mathsf{d})^4$ and
the exponential factor is at most one. Hence
\begin{equation}\label{5.12:eq-torus-bound-proof-15}
C_{\mathrm{UV}}(\mathsf{d})\le C\mathsf{d}^{-4}
+C_e\,O(1)\bigl(1+(24/\mathsf{d})^{4}\bigr)\bigl(1+\log_L^+(24M/(\mathsf{d}N_{\mathrm{side}}))\bigr),
\end{equation}
and the right-hand side does not depend on $K$.

Finally, the derivative of the constant $\log\tilde Z(0)$ vanishes, so by the decomposition of
the logarithm of $Z$ above, \eqref{5.12:eq-torus-bound-proof-11} with $\upsilon\downarrow0$ and
\eqref{5.12:eq-torus-bound-proof-13} give
\begin{equation}\label{5.12:eq-torus-bound-proof-16}
\bigl|\partial_{z_1}\cdots\partial_{z_n}\log Z(0)\bigr|
\le\Bigl[(2\mathsf{Q}+2)^n\log2+C_{\mathrm{UV}}(\mathsf{d})\Bigr]
\prod_{i=1}^n\frac{n\|f_i\|_\sharp}{r},
\end{equation}
with $C_{\mathrm{UV}}(\mathsf{d})$ bounded independently of $K$ by
\eqref{5.12:eq-torus-bound-proof-15}.
\end{proof}

\begin{proposition}[Subsequential continuum limits of the Schwinger functions]
\label{5.12:prop-subsequential-limits}
Fix $(g,m,N_T)$, where $N_T$ is the positive
integer for which $N_{\mathrm{side}}:=N_TL^m$ is the number of sites per side of the unit lattice
$T^{(K')}$ after $K'=K-m$ renormalisation steps (the letter $N$ being reserved for
$\mathrm{SU}(N)$); $N_T$ is not to be
confused with the torus floor function
$N_{\mathbb{T}}(g,m)$ of Definition~\ref{5.1:def-torus-floor}. In this proposition the index $m$ of
$S^T_{n;K,m}$ and of its limit $S^T_{n;m}$ is this number of renormalisation steps withheld, the run
being stopped at level $K'$.
For test functions $f_1,\dots,f_n$ of class $C^1$ on the physical torus $\mathbb{T}^4$ of
Definition~\ref{5.2:def-auxiliary-torus} and the complex source
$w=\sum_i z_i f_i$, with
$Z(z)$ the sourced partition function and
$\mathcal{O}(f)=\sum_p f(x(p))\,[1-\operatorname{Re}\operatorname{tr}U(\partial p)]$, set
\begin{equation}\label{5.12:eq-subsequential-limits-1}
S^T_{n;K,m}(f_1,\dots,f_n)
:=\partial_{z_1}\cdots\partial_{z_n}\log Z(z)\big|_{z=0}
=\langle \mathcal{O}(f_1);\dots;\mathcal{O}(f_n)\rangle^{\mathrm{conn}}_K .
\end{equation}
At fixed $N_T$, the left side of \eqref{5.12:eq-subsequential-limits-1} does not depend on $m$.
Let $(g_0^{(i)})_{i\ge 1}$ be a
sequence of bare couplings, let $K_i:=K(g_0^{(i)},g)$ be the first-passage levels of
Definition~\ref{5.1:def-reference-recursion}, with $x_0^{(i)}=(g_0^{(i)})^{-2}$, and assume:
\begin{enumerate}
\item the torus bound of the correlation theorem (Theorem~\ref{5.1:thm-correlation-theorem}), which
bounds $S^T_{n;K,m}$ uniformly in $K$ at pairwise separated supports;
\item the compactness lemma for equicontinuous multilinear functionals, in the following form. Let
$n\ge2$, let $B_1,\dots,B_n\subset\mathbb{T}^4$ be closed sets pairwise at positive distance, and
let $C^1_{B_j}$ be the linear space of test functions supported in $B_j$. Let $(S_i)_{i\ge1}$ be a
sequence of multilinear functionals on $C^1_{B_1}\times\dots\times C^1_{B_n}$, equicontinuous on
$\|\cdot\|_\sharp$-bounded subsets and convergent on a countable $\|\cdot\|_\sharp$-dense subset.
Then $(S_i)$ converges at every tuple of $C^1_{B_1}\times\dots\times C^1_{B_n}$ to a multilinear
limit. If, moreover, $(S_i)$ is defined on all tuples of test functions with pairwise disjoint
supports and these hypotheses hold on every such product, then the limits so obtained on all these
products are the restrictions of one distribution of finite order on the open set
\begin{equation}\label{5.12:eq-subsequential-limits-2}
\{(x_1,\dots,x_n)\in(\mathbb{T}^4)^n : x_i\ne x_j \text{ for } i\ne j\};
\end{equation}
\item $K_i\to\infty$ as $i\to\infty$;
\item the torus-compatibility condition of Definition~\ref{5.1:def-torus-floor} on $(g,m,N_T)$, with
the constant $C_{\mathbb{T}}=L^{11}M$.
\end{enumerate}
Then there is a subsequence $(K_{i_l})_{l\ge1}$ such that, for every $n\ge 2$ and every tuple
$f_1,\dots,f_n$ of test functions with pairwise disjoint supports,
$S^T_{n;K_{i_l},m}(f_1,\dots,f_n)$
converges; the limit, denoted $S^T_{n;m}(f_1,\dots,f_n)$ and depending on the subsequence, is
multilinear, and $S^T_{n;m}$ is the restriction of one distribution of finite order on the open set
\eqref{5.12:eq-subsequential-limits-2}. Each $S^T_{n;m}$ is symmetric in its $n$
arguments and is
covariant under the lattice symmetries common to all members of the subsequence. Nothing is asserted for
$n=1$ or at coinciding points.
\end{proposition}

\begin{proof}
The first equality in \eqref{5.12:eq-subsequential-limits-1} is the definition. The second is the
definition of the connected correlation (joint cumulant) as the mixed derivative at the origin of the
logarithm of the generating function. Indeed, $Z(z)$ is the integral of the density sourced by
$w=\sum_i z_if_i$, and $w$ enters the weight through
$\sum_i z_i\,\mathcal{O}(f_i)$. Neither side refers to $m$ at fixed $N_T$: the observable
$\mathcal{O}(f)$ and the bare sourced density $\rho_0^w$ live on the finest lattice of the physical
torus of Definition~\ref{5.2:def-auxiliary-torus}, whose number of sites per side,
$N_{\mathrm{side}}L^{K'}=N_TL^mL^{K-m}=N_TL^K$, is fixed by $K$ and $N_T$ alone, and they make no
reference to $m$.

Fix $(g,m,N_T)$ obeying the torus-compatibility condition, under which the torus bound of
assumption (1) is in force, and let $K_i\to\infty$. Since
$N_{\mathrm{side}}$ is fixed, the torus bound, in the form \eqref{5.12:eq-torus-bound-proof-a} of
its proof, applies, with the same value of $N_{\mathrm{side}}$, to
every member of the sequence with $K_i$ sufficiently large (in particular $K_i\ge m$, so that
$K'=K_i-m\ge0$); by the third assumption this excludes only finitely many members, and discarding them
we assume that it applies to every member of the sequence.
It bounds $S^T_{n;K_i,m}(f_1,\dots,f_n)$ uniformly in $i$ in terms of the
norms $\|f_j\|_\sharp$ and the separation of the supports.

Fix $n\ge2$ and closed sets $B_1,\dots,B_n\subset\mathbb{T}^4$ pairwise at distance at least
$\mathsf{d}>0$, with $C^1_{B_j}$ as in assumption (2).
For tuples $f=(f_1,\dots,f_n)$ and $f'=(f'_1,\dots,f'_n)$ in
$C^1_{B_1}\times\dots\times C^1_{B_n}$, multilinearity gives the telescoping identity
\begin{equation*}
S^T_{n;K_i,m}(f)-S^T_{n;K_i,m}(f')
=\sum_{j=1}^n S^T_{n;K_i,m}(f'_1,\dots,f'_{j-1},f_j-f'_j,f_{j+1},\dots,f_n).
\end{equation*}
Each tuple on the right again lies in $C^1_{B_1}\times\dots\times C^1_{B_n}$, because each
$C^1_{B_j}$ is a linear space; hence its supports are pairwise at distance at least $\mathsf{d}$, and
the torus bound estimates each term by a constant, independent of $i$, depending on $n$, $\mathsf{d}$
and the fixed data $(g,m,N_T)$, times the product of the $\|\cdot\|_\sharp$-norms of its entries.
Calling this constant $C_n(\mathsf{d})$, we obtain
\begin{equation*}
\bigl|S^T_{n;K_i,m}(f)-S^T_{n;K_i,m}(f')\bigr|
\le C_n(\mathsf{d})\sum_{j=1}^n\|f_j-f'_j\|_\sharp\prod_{r\ne j}\max\{\|f_r\|_\sharp,\|f'_r\|_\sharp\}
\end{equation*}
uniformly in $i$. Thus $\{S^T_{n;K_i,m}\}_{i\ge1}$ is equicontinuous on $\|\cdot\|_\sharp$-bounded
subsets of $C^1_{B_1}\times\dots\times C^1_{B_n}$.

Every tuple of test functions with pairwise disjoint supports lies in such a product for a choice in
which each $B_j$ is a finite union of closed balls with rational centres and rational radii. Indeed,
the supports are compact and pairwise disjoint, hence pairwise at distance at least $3\mathsf{d}$ for
some $\mathsf{d}>0$. Each point of $\operatorname{supp}f_j$ lies in a closed ball with rational centre
within $\mathsf{d}/4$ of it and rational radius strictly between $\mathsf{d}/4$ and $\mathsf{d}/2$.
By compactness
finitely many of these balls cover $\operatorname{supp}f_j$; let $B_j$ be their union. Every point of
$B_j$ lies within $\mathsf{d}$ of $\operatorname{supp}f_j$, so the sets $B_j$ are pairwise at distance
at least $\mathsf{d}$. There are countably many such choices of $(B_1,\dots,B_n)$.

On $C^1$ functions on the compact torus the norm $\|\cdot\|_\sharp$ is dominated by a multiple of the
$C^1$ norm, and the space of $C^1$ functions is separable for the $C^1$ norm; hence each
$C^1_{B_1}\times\dots\times C^1_{B_n}$ contains a countable $\|\cdot\|_\sharp$-dense subset. Let
$\mathcal{D}$ be the union of these subsets over all $n\ge2$ and all countably many choices of
$(B_1,\dots,B_n)$ above; $\mathcal{D}$ is countable. For each tuple in $\mathcal{D}$, the numbers
$S^T_{n;K_i,m}(f_1,\dots,f_n)$ form a bounded sequence, bounded by the same estimate, so it has a
convergent subsequence. Enumerate $\mathcal{D}$ and extract successively nested subsequences, one for
each of its elements. The diagonal sequence $(K_{i_l})$ then gives convergence of
$S^T_{n;K_{i_l},m}$ on the whole of $\mathcal{D}$, for every $n\ge 2$.

Equicontinuity on each of the products of the countable family above and convergence on a dense
subset of it are exactly the hypotheses of assumption (2). Its first conclusion gives, on each of
these products, convergence of $S^T_{n;K_{i_l},m}$ at every tuple to a multilinear limit. Since every
tuple with pairwise disjoint supports lies in one of these products, and at a tuple lying in several
of them the limit is the limit of one sequence of numbers, hence the same, this defines, for each
$n\ge 2$, a multilinear limit $S^T_{n;m}$ of $S^T_{n;K_{i_l},m}$ on all tuples with pairwise disjoint
supports. The hypotheses of assumption (2) now hold on every product
$C^1_{B_1}\times\dots\times C^1_{B_n}$ with closed sets $B_j$ pairwise at positive distance:
equicontinuity was shown above for every such product, and convergence holds at every tuple of it,
in particular on the countable $\|\cdot\|_\sharp$-dense subset that it contains by the separability
argument above, which applies to every such product. The second conclusion of assumption (2)
therefore shows that the limits on all these products are the restrictions of one distribution of
finite order on the open set \eqref{5.12:eq-subsequential-limits-2}. This open set is covered by the
sets $\operatorname{int}B_1\times\dots\times\operatorname{int}B_n$ of the countable family: if the
coordinates of $(x_1,\dots,x_n)$ are pairwise at distance at least $3\mathsf{d}$ for some
$\mathsf{d}>0$, a closed ball $B_j$ with rational centre within $\mathsf{d}/4$ of $x_j$ and rational
radius strictly between $\mathsf{d}/4$ and $\mathsf{d}/2$ contains $x_j$ in its interior, and these
balls are pairwise at distance at least $\mathsf{d}$.

For symmetry, note that $S^T_{n;K,m}$ is, by \eqref{5.12:eq-subsequential-limits-1}, a mixed partial
derivative of the function $z\mapsto\log Z(z)$. The lattice is finite and the integrand defining
$Z(z)$ is the exponential of a function affine in $z$, so $Z$ is holomorphic in $z$; at $z=0$ the
integrand is positive, so $Z(0)>0$. Hence $\log Z$ is holomorphic near $z=0$ and its mixed partial
derivatives commute, so the value of $S^T_{n;K,m}(f_1,\dots,f_n)$ is unchanged under every
permutation of its arguments. This identity holds for every member of the subsequence and therefore
for the limit.

For covariance, let $\iota$ be an isometry of $\mathbb{T}^4$ that maps the finest lattice at cutoff
$K_{i_l}$ onto itself for every $l$, and for a configuration $U$ let $\iota_*U$ be the
configuration with $(\iota_*U)(\iota b)=U(b)$ on every oriented bond $b$. The product of the Haar
measures over the bonds and the Wilson action are invariant under $U\mapsto\iota_*U$. Relabelling the
plaquettes by $p'=\iota p$, using $x(\iota p)=\iota\, x(p)$, and using that
$(\iota_*U)(\partial(\iota p))$ equals $U(\partial p)$ up to inversion and conjugation, neither of
which changes $\operatorname{Re}\operatorname{tr}$, one has
\begin{equation*}
\mathcal{O}(f)(\iota_*U)
=\sum_{p}f(\iota\, x(p))\,[1-\operatorname{Re}\operatorname{tr}U(\partial p)]
=\mathcal{O}(f\circ\iota)(U).
\end{equation*}
Since $w$ enters the weight only through $\sum_i z_i\,\mathcal{O}(f_i)$, the change of variables
$U\mapsto\iota_*U$ in the integral defining $Z(z)$ shows that $Z(z)$ is unchanged when every $f_i$ is
replaced by $f_i\circ\iota$. By \eqref{5.12:eq-subsequential-limits-1},
\begin{equation*}
S^T_{n;K_{i_l},m}(f_1\circ\iota,\dots,f_n\circ\iota)=S^T_{n;K_{i_l},m}(f_1,\dots,f_n)
\end{equation*}
for every $l$, and letting $l\to\infty$ gives the same identity for $S^T_{n;m}$.
\end{proof}

\begin{lemma}[Reflection positivity of the limits at disjoint supports]
\label{5.12:lem-limit-positivity}
Assume the lemma on reflection positivity of the lattice measure \cite{OSe78}. For a point $x$
of the torus write
$x^{0}$ for its time coordinate, and let $\theta$ denote the time reflection $x^{0}\mapsto -x^{0}$ for
which that lemma holds; assume further that the expectation
$\langle\cdot\rangle_K$ is invariant under $\theta$. The reflection $\theta$ acts on test functions by
$(\theta f)(x):=f(\theta x)$, on bond variables by $(\theta U)(b):=U(\theta b)$ for an oriented
bond $b$, and on functions of the bond variables by $(\theta F)(U):=F(\theta U)$. Let
$(K_{i_l})_{l\ge 1}$ be the subsequence of the sequence $(K_i)$ along which
Proposition~\ref{5.12:prop-subsequential-limits} provides the limits $S^T_{n;m}$, $n\ge 2$, of the
connected Schwinger functions $S^T_{n;K,m}$. In this lemma $g$ and $m$ are fixed as in
Proposition~\ref{5.12:prop-subsequential-limits}, and we abbreviate $S^K_n:=S^T_{n;K,m}$ and
$S_n:=S^T_{n;m}$.

Let $(f_i)_{i\in I}$, with $I$ a finite index set, be test functions whose supports are pairwise
disjoint and contained in $\{0<x^{0}<\tfrac12\}$. For subsets $\alpha\subset I$ let $c_\alpha$ be
complex numbers, and put
\begin{equation*}
F:=\sum_{\alpha}c_\alpha\prod_{i\in\alpha}\bigl(\mathcal{O}(f_i)-\langle\mathcal{O}(f_i)\rangle_K\bigr).
\end{equation*}
Then $\langle\overline{\theta(F)}\,F\rangle_K\ge 0$ for every $K$. This quantity is a
polynomial $P_K(F)$ in the connected functions $S^K_n$, $n\ge 2$, evaluated at test functions with
pairwise disjoint supports; it depends only on the $c_\alpha$ and the $f_i$, not on the centring
constants. As $l\to\infty$, $P_{K_{i_l}}(F)$ converges to the same polynomial $P(F)$ in the
limits $S_n$, and
\begin{equation}\label{5.12:eq-limit-positivity-1}
P(F)=\lim_{l\to\infty}\bigl\langle\overline{\theta(F)}\,F\bigr\rangle_{K_{i_l}}\ \ge\ 0 .
\end{equation}
Nothing is asserted for the one-point functions ($n=1$) or at coinciding points.
\end{lemma}

\begin{proof}
For each $K$, $F$ is a function of the bond variables of the plaquettes whose
barycentres lie in the union of the supports of the $f_i$. This union is contained in
$\{0<x^{0}<\tfrac12\}$. Reflection positivity of the lattice measure, applied to $F$, therefore gives
$\langle\overline{\theta(F)}F\rangle_K\ge 0$ (cf.\ \cite{OSe78}).

We next write the left side in terms of connected functions. The reflection maps plaquettes to
plaquettes and barycentres to reflected barycentres, and $1-\operatorname{Re}\operatorname{tr}U(\partial p)$ is
unchanged under reversal of orientation. Hence $\theta(\mathcal{O}(f))=\mathcal{O}(\theta f)$. By the
hypothesis that the expectation is invariant under $\theta$,
$\langle\mathcal{O}(\theta f)\rangle_K=\langle\mathcal{O}(f)\rangle_K$. Since $\mathcal{O}(f)$ is linear in
$f$ with real plaquette weights, $\overline{\mathcal{O}(f)}=\mathcal{O}(\bar f)$. Together these give
\begin{equation*}
\overline{\theta(F)}=\sum_{\alpha}\bar c_\alpha\prod_{i\in\alpha}
\bigl(\mathcal{O}(\theta\bar f_i)-\langle\mathcal{O}(\theta\bar f_i)\rangle_K\bigr).
\end{equation*}
For a pair $(\alpha,\beta)$ let $\alpha'=\{i':i\in\alpha\}$ be a disjoint copy of $\alpha$. Put
$h_{i'}:=\theta\bar f_i$ for $i\in\alpha$ and $h_j:=f_j$ for $j\in\beta$. Then
$\langle\overline{\theta(F)}F\rangle_K$ is the sum over $(\alpha,\beta)$ of $\bar c_\alpha c_\beta$ times
the moment of the centred variables $\mathcal{O}(h_j)-\langle\mathcal{O}(h_j)\rangle_K$,
$j\in\alpha'\sqcup\beta$.

By the moment--cumulant formula, a moment is the sum over partitions $\pi$ of the index set of the
products over blocks $B\in\pi$ of joint cumulants. Cumulants of order one of centred variables vanish.
Cumulants of order at least two do not change when constants are added to the variables. By the
definition in Proposition~\ref{5.12:prop-subsequential-limits}, the joint cumulant of
$\mathcal{O}(h_j)$, $j\in B$, is
\begin{equation*}
S^K_{|B|}\bigl((h_j)_{j\in B}\bigr)
=\langle\mathcal{O}(h_{j_1});\dots;\mathcal{O}(h_{j_{|B|}})\rangle^{\mathrm{conn}}_K,
\end{equation*}
where $j_1,\dots,j_{|B|}$ enumerate $B$; the order is immaterial, since this function is symmetric
in its arguments. Write $\Pi_{\ge 2}(J)$ for the set of partitions of a finite set
$J$ into blocks of size at least two; for $J=\emptyset$ it consists of the empty partition, with empty
product $1$. Then $\langle\overline{\theta(F)}F\rangle_K=P_K(F)$, where
\begin{equation}\label{5.12:eq-limit-positivity-2}
P_K(F):=\sum_{\alpha,\beta}\bar c_\alpha c_\beta
\sum_{\pi\in\Pi_{\ge2}(\alpha'\sqcup\beta)}\ \prod_{B\in\pi}S^K_{|B|}\bigl((h_j)_{j\in B}\bigr).
\end{equation}
Let $P(F)$ be the same expression with $S^K_{|B|}$ replaced by $S_{|B|}$. Both $P_K(F)$ and $P(F)$
depend only on the coefficients $c_\alpha$ and the test functions $f_i$, not on the centring
constants.

In each term of \eqref{5.12:eq-limit-positivity-2}, every $f_i$ and every $\theta\bar f_i$ occurs at
most once. The supports of the $h_j$, $j\in\beta$, are pairwise disjoint by hypothesis. Those of the
$h_{i'}$, $i\in\alpha$, are pairwise disjoint because $\theta$ is a bijection and conjugation preserves
supports. The two families are disjoint from each other because they lie in
$\{0<x^{0}<\tfrac12\}$ and $\{-\tfrac12<x^{0}<0\}$ respectively. Hence each tensor product
$\bigotimes_{j\in B}h_j$ is supported in $\{x_i\ne x_j\}$. On that set the limits of
Proposition~\ref{5.12:prop-subsequential-limits} are defined and attained along $(K_{i_l})$, so
$S^{K_{i_l}}_{|B|}((h_j)_{j\in B})\to S_{|B|}((h_j)_{j\in B})$ as $l\to\infty$.

The sum \eqref{5.12:eq-limit-positivity-2} is finite and each term is a finite product. Therefore
$P_{K_{i_l}}(F)\to P(F)$. Since every $P_{K_{i_l}}(F)=\langle\overline{\theta(F)}F\rangle_{K_{i_l}}$ is
non-negative, so is the limit, which is \eqref{5.12:eq-limit-positivity-1}. Only cumulants of order at
least two, at pairwise disjoint supports, entered. This is why nothing is asserted for the one-point
functions ($n=1$) or at coinciding points.
\end{proof}

\begin{proposition}[The correlation theorem for $\mathrm{SU}(N)$, $N\ge3$]
\label{5.13:prop-sun-correlation}
Let $N\ge3$, and assume hypotheses (H1)--(H7) of Theorem~0 with this value of $N$. Then the
statements (iii-a), (iii-b) and (iii-c) of clause (iii) of Theorem~0 hold as stated there,
with one difference: the constants $C_n(\mathsf{d})$ and $C^\infty_n(\mathsf{d};R)$, and the
thresholds on the coupling, depend on $N$. They depend on it through the existential constant
$L_0(N)$, through the one-loop constant $c_0=(11N^2/12\pi^2)\log L$, and through the constants
of the Lie algebra $\mathfrak{su}(N)$. Apart from this dependence on $N$ they read the same
data as at $N=2$.
\end{proposition}

Proof outstanding.

\paragraph{Route.}
A proof is a reading, at general $N$, of the proof of the correlation theorem of this chapter
and of the proof of the per-ball volume-free bound in Chapter~\ref{ch:6}. Those proofs are
given for $N=2$. The rank enters through the following quantities:
\begin{itemize}
\item the one-loop constant $c_0$, through the lower bound
$\lambda_\sharp\ge\underline{\lambda}=c_0/4$ on the slope of the flow of the coupling;
\item the existential constant $L_0(N)$;
\item the constants of $\mathfrak{su}(N)$ that enter the stability constants $E_\pm$, and
hence the thresholds $\gamma_H$ and $\gamma_J$.
\end{itemize}
A proof must establish two things. First, no step of the argument at $N=2$ uses $N=2$ other
than through these quantities. Second, the stability constants $E_\pm$ of clause (i) have an
$\mathrm{SU}(N)$ form, and a proof must supply it, since the argument rests on it; this is the
subject of the two statements of Chapter~\ref{ch:4} that concern $\mathrm{SU}(N)$. With these
two points in hand, the torus bound (iii-a), the per-ball volume-free bound (iii-b) and the
running-shift statement (iii-c) would be obtained by the same arguments. In those arguments
$L_0(N)$ remains a single constant bound existentially, every condition on $L$ remains a lower
bound, and the dependence on $N$ is carried through $E_\pm$ into the thresholds.

Clause (iv) of Theorem~0 at $N\ge3$ would rest on this proposition. At $N\ge3$, Theorem~0
asserts nothing beyond clauses (0) and (i).

\paragraph{Inputs used by statement}
Besides what is proved in it, this chapter uses published results of Ba{\l}aban, as
stated and with the reference given, and statements of this book whose
own proofs are incomplete or outstanding. Both are listed below, clause by clause of
Theorem~0. Ba{\l}aban's definitions and notation are cited where they are introduced
and are not repeated here. A hypothesis that a statement of this chapter displays is
stated where it is used. Two published results of other authors are used as stated and
are cited at their use: the estimate of Combes and Thomas \cite{CT73}, in
Lemma~\ref{5.2:lem-covariance-root}, and reflection positivity of lattice gauge theories
\cite{OSe78}, in Lemma~\ref{5.12:lem-limit-positivity}.

\begin{description}
\item[(0)] This chapter proves no part of clause~(0).
\item[(i)] This chapter proves no part of clause~(i). The proof of
Theorem~\ref{5.4:thm-finer-zone-sum} (summation of boundary terms at finer zones) takes as
hypotheses the containment statements of clause~(i), at the standing at which the chapter
on clauses~(0) and~(i) establishes them; those among these statements and their inputs
whose proofs are incomplete or outstanding are named, with that word, in the corresponding
list of that chapter.
\item[(iii)] The proofs of items (a), (c) and (d) of clause~(iii) use the following.
\begin{itemize}
\item[] \emph{Published statements cited by statement.}
\begin{itemize}
\item \cite[Proposition~2]{B-CMP98};
\item \cite[(11), (97), (148)]{B-CMP98}, on the
averaging operations;
\item \cite[(1.3)--(1.6)]{B-CMP99a}, on regular gauge field configurations;
\item \cite[Theorems~3.1--3.11]{B-CMP99b};
\item \cite[Theorem~3.15]{B-CMP99b};
\item \cite[(3.19), (3.35)--(3.36), (3.108), (3.113)--(3.114), (3.128)--(3.129), (3.138),
(3.154)--(3.156), (3.158), (3.161), (3.168)--(3.169), (3.185), (3.187)]{B-CMP99b}, on
propagators in a background field;
\item \cite[Theorem~1]{B-CMP102b};
\item \cite[Proposition~9]{B-CMP102b};
\item \cite[(7), (9)--(10), (170)--(171), (174), (178), (181), (190)]{B-CMP102b}, on the
variational problem and background fields;
\item the analysis of \cite[Sect.~F]{B-CMP102b};
\item the analysis of \cite[Sect.~G]{B-CMP102b};
\item \cite[(2.1)--(2.12)]{B-CMP109}, and the analysis of \cite[Sect.~2]{B-CMP109};
\item \cite[Lemma~4, pp.~276--277]{B-CMP109};
\item the analysis of \cite[Sects.~3--5, (3.2)--(3.54), (4.1)--(4.36), (5.10)]{B-CMP109};
\item \cite[Lemma~1, p.~9]{B-CMP116};
\item \cite[Lemma~2, p.~11]{B-CMP116};
\item \cite[Lemma~3]{B-CMP116};
\item \cite[\S1, (1.9)--(1.29)]{B-CMP116}, the localised terms of the cluster expansion,
their domains of analyticity and their bounds;
\item \cite[(1.42)--(1.43), (2.1)--(2.42)]{B-CMP116}, the cluster expansion and its bounds;
\item \cite[Theorem~1, Lemma~4, Corollary~3, (2.19)--(2.22), (2.27), (2.30)--(2.50)]{B-CMP119},
the inductive hypotheses and bounds of the renormalization expansions;
\item the analysis of \cite[Sect.~3, (3.2)--(3.68)]{B-CMP119};
\item the cube conditions \cite[p.~190, (1.64)]{B-CMP122a}, with Ba{\l}aban's enlargement
fraction one fifth, used in Lemma~\ref{5.2:lem-enlargement-fractions};
\item \cite[(1.60)--(1.63), (1.65)--(1.70), (1.84), pp.~188--189]{B-CMP122a}, on the basic
step of the operation $\mathbf{R}$;
\item \cite[(1.49)--(1.65)]{B-CMP122b}, the localisation of the operation $\mathbf{R}$ and the
analyticity of its terms;
\item \cite[(1.66)--(1.72), (1.74)--(1.99), (1.103)--(1.104), (2.21)]{B-CMP122b}, bounds
for the operation $\mathbf{R}$;
\item \cite[Theorem~1]{B-CMP122b};
\item Ba{\l}aban's bounds for old terms beside a new large-field region,
\cite[p.~375]{B-CMP122b} --- printed without proof; cited as printed;
\item the assertion \cite[p.~368, after (1.46)--(1.48)]{B-CMP122b} that the sum of the
localised remainders of the localisation and exponentiation
\cite[(1.11), (1.23)--(1.48)]{B-CMP122b} is small, used with the preceding bounds ---
printed without proof; cited as printed.
\end{itemize}
\item[] \emph{Statements of this book whose proofs are incomplete or outstanding.}
\begin{itemize}
\item Proposition~\ref{5.6:prop-polydisc-holomorphy} (holomorphy of the localised terms on
the decoupling polydisc) --- proof outstanding.
\end{itemize}
\end{itemize}
This chapter also prints Proposition~\ref{5.13:prop-sun-correlation} (the correlation
theorem for $\mathrm{SU}(N)$, $N\ge 3$), with its proof outstanding; it is not used in the
proofs above.
\item[(ii)] Proposition~\ref{5.12:prop-subsequential-limits} and
Lemma~\ref{5.12:lem-limit-positivity}, which concern clause~(ii), rest on
Theorem~\ref{5.1:thm-correlation-theorem} (the correlation theorem) and on reflection
positivity of the lattice measure, proved in \cite{OSe78}. Beyond the
inputs listed under~(iii), they use no published statement of Ba{\l}aban by statement
and no statement of this book whose proof is incomplete or outstanding.
\item[(iv)] This chapter proves no part of clause~(iv) other than the sentence on the
running of the coupling stated with it: the one-loop law of the coupling at separation
$\mathsf{d}$, Corollary~\ref{5.14:cor-separation-one-loop}, and, jointly with the two-point
law of clause~(iv), the statement that the charge read off the two-point function runs at
the one-loop rate within fixed factors, Corollary~\ref{5.14:cor-the-charge-read-off-the-two};
the inputs of these corollaries of the two kinds listed here are among those listed
under~(iii), and those of the two-point law are listed with the chapter on clause~(iv).
\item[(v)] This chapter proves no part of clause~(v). The statements of this chapter
that compare two cutoffs are used in the proof of clause~(v) and not in that of
clause~(iii); besides published statements listed under~(iii), they use the following.
\begin{itemize}
\item[] \emph{Published statements cited by statement.}
\begin{itemize}
\item \cite[(2.46), (2.54), (2.59)]{B-CMP96};
\item \cite[Propositions~4, 5, 6 and~7]{B-CMP98};
\item \cite[(42), (43), (47)--(51), (159)--(164)]{B-CMP98}, on the averaging operations;
\item \cite[(3.42)--(3.49), (3.120), (3.126), (3.132)--(3.133)]{B-CMP99b}, on propagators
in a background field;
\item \cite[(39)--(40), (44), (54), (69)--(73), (90)--(96)]{B-CMP102b}, on the variational
problem and background fields.
\end{itemize}
\item[] \emph{Statements of this book whose proofs are incomplete or outstanding.}
\begin{itemize}
\item the statement that all two-cutoff differences on one lattice vanish --- stated;
proof incomplete.
\end{itemize}
\end{itemize}
\item[(vi)] This chapter proves no part of clause~(vi).
\end{description}

The published statements listed above depend on nothing proved in this book. The
constants of this chapter are chosen in the order fixed in
Definition~\ref{5.1:def-admitted-inputs}, each after every constant on which it depends.
Two dependences among the statements used here are to be noted. First, at the site
labelled (P) in Proposition~\ref{5.2:prop-three-sites}, part of the hypothesis of
Lemma~\ref{5.6:lem-weight} (the weight lemma) is supplied by
Corollary~\ref{5.2:cor-created-bound-induction} (the created-term bound at every step), a
statement of the joint induction by which Theorem~\ref{5.1:thm-correlation-theorem} is
proved; there the weight lemma is proved as an implication from that corollary. No cycle
results: in the order of the joint induction, the created-term bound is obtained at every
step already performed before the weight lemma is used at the site (P) of the next step,
so the dependence runs along the induction. Second,
Corollary~\ref{5.14:cor-the-charge-read-off-the-two} uses the two-point law of
clause~(iv), which the chapter on clause~(iv) deduces from the correlation theorem and not
from any corollary of the one-loop law.

%% ---- end inlined file: chapters/c05 ----

\clearpage
\setcounter{section}{5}
\refstepcounter{section}
\section*{Chapter~\thesection\quad A volume-free bound for each ball: clause (iii-b)\addcontentsline{toc}{section}{\protect\numberline{\thesection}A volume-free bound for each ball: clause (iii-b)}\label{ch:6}}
%% ---- begin inlined file: chapters/c06 ----
\begin{definition}[Last-lattice units, plaquette functionals and sources]\label{6.1:not-units-sources}
In this chapter lengths are measured in the units of the last lattice $T^{(K')}$ of the correlation theorem, where $K'$ is the index of that theorem's last level.

\begin{enumerate}
\item[(a)] \emph{Units.} We use \emph{last-lattice units}: $T^{(K')}$ has spacing $1$ and $\mathcal{N}$ sites per side, where $\mathcal{N}:=2L^mL^b$ and the integer $b\ge0$ is specified below. In the units of the correlation theorem, in which the torus has side $1$, the same lattice has spacing $1/\mathcal{N}$. When the last level of the correlation theorem is identified with the unit lattice on which Theorem~0 is stated, one has $K'=K$ and $b=0$, so that $\mathcal{N}=2L^m$. When it is instead identified, as in the lemma on uniformity of the last-scale constants in the torus exponent (Lemma~\ref{6.3:lem-last-scale-constants}), with the lattice of the correlation theorem's torus index $m_J$, the exponent $b$ is the unit-conversion exponent fixed there, and
\begin{equation*}
0\le b\le b_*(g),
\end{equation*}
where $b$ and its bound $b_*(g)$, also fixed there, are independent of $K$ and of $m$; since $b$ is an integer in $[0,b_*(g)]$, the side $\mathcal{N}$ takes finitely many values once $g$, $m$ and the torus exponent $m_1$ of Definition~\ref{6.3:def-source-classes} are fixed. In that case physical lengths are last-lattice lengths multiplied by $L^{-b}$. Every object below lives on the bare lattice of spacing $\varepsilon=L^{-K'}$ of item~(b), inside the torus of side $\mathcal{N}$, with all lengths expressed in last-lattice units. The conversion to the units of Theorem~0 multiplies the radius $R$ and the separation $\mathsf{d}$ of item~(e) below by $L^b$: a radius $R$ and a separation $\mathsf{d}$ in the units of Theorem~0 are a radius $RL^b$ and a separation $\mathsf{d}L^b$ in last-lattice units. Every constant remains independent of $K$ and of $m$; the conversion is carried out at the end of this chapter.

\item[(b)] \emph{Lattice and torus.} Let $\varepsilon:=L^{-K'}$ be the bare spacing. The bare sites are the points $x\in\varepsilon(\mathbb{Z}+\tfrac12)^4$ with $|x_\mu|<\mathcal{N}/2$ for $\mu=1,\dots,4$; this is the lattice of centres of \cite[(0.1)]{B-CMP109}. Let $\mathbb{T}:=\mathbb{R}^4/\mathcal{N}\mathbb{Z}^4$ be the continuum torus. For a plaquette $p$ of the bare lattice, $x_p\in\mathbb{T}$ denotes its barycentre, as in the definition of the observables $\mathcal{O}(f)$.

\item[(c)] \emph{Plaquette functionals.} For a bare plaquette $p$ and a configuration $U$ put
\begin{equation}\label{6.1:eq-units-sources-1}
a_p(U):=1-\operatorname{Re}\operatorname{tr}U(\partial p),
\end{equation}
as in \cite[(0.2)]{B-CMP109}. For $v\colon\mathbb{T}\to\mathbb{C}$ put
\begin{equation}\label{6.1:eq-units-sources-2}
A(v,U):=\sum_p v(x_p)\,a_p(U),
\end{equation}
which is linear in $v$. Let $\mu_{K,m}$ be the bare Wilson probability measure, proportional to $e^{-x_0A(U)}\,dU$ with $x_0=g_0^{-2}$, and let $\langle\cdot\rangle$ denote its expectation. Here $dU$ is the product over the bare bonds of the normalised Haar measure of the gauge group $\mathrm{SU}(N)$. The sourced partition function is
\begin{equation}\label{6.1:eq-units-sources-3}
Z(v):=\int\exp\bigl[-A(x_0-v,U)-E\bigr]\,dU,
\end{equation}
where $x_0-v$ is the function $x\mapsto x_0-v(x)$ on $\mathbb{T}$, and $E$ is the constant, independent of $v$ and of $U$, that the correlation theorem carries in its partition function. This is the partition function $Z$ of clause (iii)(b) of Theorem~0.

\item[(d)] \emph{Scales.} Let $M=L^{m'}$ be as in the statement of Theorem~0. Let $\pi_j$ be the scale-$j$ partition into $M$-cubes of \cite[(0.24)--(0.26)]{B-CMP109}. In last-lattice units its cubes have side
\begin{equation*}
s_j:=ML^{j-K'}.
\end{equation*}
In the units of the correlation theorem these cubes have side $M\mathsf{t}_j$, where $\mathsf{t}_j$ is the scale of that theorem at level $j$; that is, $s_j=M\mathsf{t}_j\mathcal{N}$, so that $\mathsf{t}_j=L^{j-K'}/\mathcal{N}$. The class $D_j$ of localisation domains consists of the finite unions of $\pi_j$-cubes that are wall-connected in the sense of \cite[(0.24)--(0.26)]{B-CMP109}.

For $X\in D_j$, let $d_j(X)$ be the length of a shortest tree graph in $X$ meeting all cubes of $X$, divided by $s_j$, so that the cubes of $\pi_j$ become unit cubes \cite[\S0]{B-CMP109}. Let $X^+$ be $X$ together with one layer of $\pi_j$-cubes around it. For a set $Y$, let $\operatorname{plaq}(Y)$ be the set of bare plaquettes that intersect $Y$ \cite[\S0]{B-CMP109}.

For a Lipschitz function $f$ on $\mathbb{T}$ put
\begin{equation*}
\|f\|_{\sharp}:=\max\{\|f\|_\infty,\ \operatorname{Lip}f/\lambda_*\},\qquad \lambda_*:=(40M)^{-1},
\end{equation*}
with the Lipschitz constant taken in last-lattice units.

\item[(e)] \emph{Sources.} Let $n\ge2$. Let $f_1,\dots,f_n\colon\mathbb{T}\to\mathbb{R}$ be Lipschitz, with
\begin{equation*}
\operatorname{supp}f_i\subset B_i:=\overline{B}(x_i,R),
\end{equation*}
the closed Euclidean ball of radius $R$ about $x_i$. Suppose that $\operatorname{dist}(\operatorname{supp}f_i,\operatorname{supp}f_k)\ge\mathsf{d}>0$ for $i\ne k$. Let $r>0$ be the source radius of the correlation theorem; it is distinct from the extraction depth $r(s)$ of clause (iv). Put
\begin{equation}\label{6.1:eq-units-sources-4}
\mathbb{D}:=\Bigl\{z\in\mathbb{C}^n:\ \rho(z):=r^{-1}\sum_{i=1}^n|z_i|\,\|f_i\|_{\sharp}<1\Bigr\},
\qquad
w_z:=\sum_{i=1}^n z_if_i .
\end{equation}
Define
\begin{equation}\label{6.1:eq-units-sources-5}
S^T_n(\vec f):=\partial_{z_1}\cdots\partial_{z_n}\log Z(w_z)\big|_{z=0}.
\end{equation}
\end{enumerate}

Then $a_p(U)$ is real, satisfies $a_p(U)\ge0$, and does not depend on the orientation of $p$. For every $v$ one has $Z(v)/Z(0)=\langle e^{A(v,\cdot)}\rangle$; in particular the constant $E$ and the normalisation of $\mu_{K,m}$ cancel. Finally, $S^T_n(\vec f)$ is well defined and equals the joint cumulant of $\mathcal{O}(f_1),\dots,\mathcal{O}(f_n)$ under $\mu_{K,m}$.
\end{definition}

\begin{proof}
\emph{Properties of $a_p$.} Since $U(\partial p)$ is unitary, its eigenvalues have modulus one; as $\operatorname{tr}$ is normalised by $\operatorname{tr}\mathbf{1}=1$, we have $|\operatorname{tr}U(\partial p)|\le1$. Hence $a_p(U)$ is real and lies in $[0,2]$.

Changing the starting point of $\partial p$ replaces $U(\partial p)$ by a conjugate, which has the same trace. Reversing the orientation replaces $U(\partial p)$ by $U(\partial p)^{-1}=U(\partial p)^*$, whose trace is the complex conjugate. Neither operation changes $\operatorname{Re}\operatorname{tr}U(\partial p)$.

\emph{The ratio $Z(v)/Z(0)$.} By linearity of \eqref{6.1:eq-units-sources-2} in $v$,
\begin{equation*}
A(x_0-v,U)=x_0A(1,U)-A(v,U).
\end{equation*}
Moreover $A(1,U)=\sum_pa_p(U)=A(U)$. Therefore \eqref{6.1:eq-units-sources-3} gives
\begin{equation*}
Z(v)=e^{-E}\int e^{-x_0A(U)}\,e^{A(v,U)}\,dU,\qquad Z(0)=e^{-E}\int e^{-x_0A(U)}\,dU.
\end{equation*}
The integrand of $Z(0)$ is continuous and strictly positive, and $dU$ is a probability measure, so $Z(0)>0$. Dividing the first identity by the second, the factor $e^{-E}$ cancels. What remains is the expectation of $e^{A(v,\cdot)}$ under the probability measure proportional to $e^{-x_0A(U)}\,dU$, that is, under $\mu_{K,m}$. This proves $Z(v)/Z(0)=\langle e^{A(v,\cdot)}\rangle$.

\emph{The cumulant.} For $z\in\mathbb{C}^n$, linearity of \eqref{6.1:eq-units-sources-2} and the definition of $\mathcal{O}(f)$ give
\begin{equation*}
A(w_z,U)=\sum_{i=1}^nz_iA(f_i,U)=\sum_{i=1}^nz_i\,\mathcal{O}(f_i).
\end{equation*}
Hence, by the identity just proved,
\begin{equation*}
Z(w_z)=Z(0)\,\Bigl\langle\exp\sum_{i=1}^nz_i\,\mathcal{O}(f_i)\Bigr\rangle .
\end{equation*}

There are finitely many bare plaquettes, and $0\le a_p\le2$. So each $\mathcal{O}(f_i)$ is bounded in modulus by $2\|f_i\|_\infty$ times the number of plaquettes. The expectation on the right is therefore an entire function of $z$, obtained by differentiating under the integral sign over a compact group, and it equals $1$ at $z=0$. Consequently $Z(w_z)\ne0$ in a neighbourhood of $z=0$. Using the principal logarithm of the expectation there,
\begin{equation*}
\log Z(w_z)=\log Z(0)+\log\Bigl\langle\exp\sum_{i=1}^nz_i\,\mathcal{O}(f_i)\Bigr\rangle,
\end{equation*}
so \eqref{6.1:eq-units-sources-5} is well defined. Since $n\ge2$, the mixed derivative $\partial_{z_1}\cdots\partial_{z_n}$ annihilates the constant $\log Z(0)$. Therefore $S^T_n(\vec f)$ equals $\partial_{z_1}\cdots\partial_{z_n}$ at $z=0$ of the cumulant generating function of $(\mathcal{O}(f_1),\dots,\mathcal{O}(f_n))$ under $\mu_{K,m}$, which by definition is their joint cumulant.
\end{proof}

\begin{lemma}[Reflection positivity of the lattice measure]\label{6.1:lem-reflection-positivity}
Work in the last-lattice units of Definition~\ref{6.1:not-units-sources}, in which the bare sites of
$\mathbb{T}=\mathbb{R}^4/N\mathbb{Z}^4$ have their coordinates in $\varepsilon(\mathbb{Z}+\tfrac12)$. Let
$c\in\varepsilon\mathbb{Z}$, let $\mu$ be a coordinate direction, and put
\begin{equation*}
\Pi:=\{x_\mu=c\}\cup\{x_\mu=c+N/2\}\subset\mathbb{T}.
\end{equation*}
The set $\Pi$ is a pair of antipodal link planes: no bare site lies on it, since $c/\varepsilon\in\mathbb{Z}$,
$\mathbb{Z}\cap(\mathbb{Z}+\tfrac12)=\emptyset$, and $N/2\in\mathbb{Z}\subset\varepsilon\mathbb{Z}$. Let
$\theta\colon x_\mu\mapsto 2c-x_\mu$ (the other coordinates unchanged) be the reflection, and let $H_\pm$ be
the two open components of $\mathbb{T}\setminus\Pi$. For a bare bond $b\subset H_+$ put
\begin{equation*}
(\theta\cdot U)(b):=
\begin{cases}
U(\theta b) & \text{if $\theta$ preserves the orientation of $b$ (bonds $\perp e_\mu$)},\\
U(\theta b)^{-1} & \text{otherwise (bonds $\parallel e_\mu$)}.
\end{cases}
\end{equation*}
Let $\langle\cdot\rangle$ denote expectation with respect to the bare Wilson measure $\mu_{K,m}$. Then for
every bounded measurable complex-valued function $F$ of the bond variables in $H_+$:
\begin{enumerate}
\item[(i)] $\langle F(\theta\cdot U)\rangle=\langle F(U)\rangle$; indeed $\mu_{K,m}$ is invariant under
every isometry of $\mathbb{T}$ preserving the bare lattice;
\item[(ii)] $\langle\overline{F(\theta\cdot U)}\,F(U)\rangle\ge 0$.
\end{enumerate}
\end{lemma}

\begin{proof}
Recall that $\mu_{K,m}$ has density proportional to $e^{-x_0A(U)}$, $x_0=g_0^{-2}$, with respect to the
product over bare bonds of the normalised Haar measure $dU$ of $\mathrm{SU}(N_{\mathrm{col}})$, and that
$A(U)=\sum_p a_p(U)$ with $a_p(U)=1-\operatorname{Re}\operatorname{tr}U(\partial p)$
\cite[(0.1), (0.2)]{B-CMP109}.

\emph{Part (i).} Let $g$ be an isometry of $\mathbb{T}$ mapping the bare lattice onto itself, and let
$U\circ g$ be the configuration obtained by relabelling the bonds by $g$, inverting the variable where $g$
reverses the orientation of the bond: $(U\circ g)(b):=U(gb)$ if $g$ carries the orientation of $b$ to
that of $gb$, and $(U\circ g)(b):=U(gb)^{-1}$ otherwise. The map $U\mapsto U\circ g$ preserves product Haar
measure: the relabelling permutes the factors of the product, and Haar measure on
$\mathrm{SU}(N_{\mathrm{col}})$ is invariant under inversion. It also preserves $A$: $g$ permutes the
plaquettes, and $a_p(U\circ g)=a_{gp}(U)$, because $\operatorname{Re}\operatorname{tr}$ of a holonomy is
invariant under cyclic relabelling of the loop and under inversion of the holonomy. Hence
$\mu_{K,m}$ is invariant under $U\mapsto U\circ g$. The reflection $\theta$ is such an isometry: it maps
$\varepsilon(n+\tfrac12)$ to $\varepsilon(2c/\varepsilon-n-\tfrac12)\in\varepsilon(\mathbb{Z}+\tfrac12)$.
On the bonds of $H_+$ the configuration $\theta\cdot U$ coincides with $U\circ\theta$, so
$\langle F(\theta\cdot U)\rangle=\langle F(U\circ\theta)\rangle=\langle F(U)\rangle$.

\emph{Part (ii): classification of plaquettes.} Since $\theta$ fixes each plane of $\Pi$ modulo $N$
($2c-(c+N/2)=c-N/2$), it exchanges $H_+$ and $H_-$. A bare bond either lies in $H_+$, or lies in $H_-$, or
crosses $\Pi$; the bonds crossing $\Pi$ are the bonds parallel to $e_\mu$ with endpoints at
$x_\mu=c\pm\varepsilon/2$, respectively at $x_\mu=c+N/2\pm\varepsilon/2$. We call them \emph{crossing links}
and write $\ell_x$ for the crossing link through the point $x\in\Pi$; the points $x$ form two
three-dimensional lattices of spacing $\varepsilon$, one in each plane of $\Pi$. Each $\ell_x$ is oriented
from its endpoint in $H_-$ to its endpoint in $H_+$.

A plaquette containing no crossing link lies entirely in $H_+$ or entirely in $H_-$: its boundary is
connected through its sites, each of which lies in one open half, and none of its bonds crosses $\Pi$.
A plaquette containing a crossing link $\ell_x$ lies in a $(\mu,\nu)$-plane with $\nu\neq\mu$ and is of the
form $p_{x,\nu}$ (for the appropriate $x$), the plaquette with boundary links
$\ell_x$, $b^+_{x,\nu}$, $\ell_{x+\varepsilon e_\nu}$, $b^-_{x,\nu}$. Here $b^\pm_{x,\nu}\subset H_\pm$ joins
the endpoints in $H_\pm$ of $\ell_x$ and $\ell_{x+\varepsilon e_\nu}$ and is oriented from the side of $x$
to the side of $x+\varepsilon e_\nu$. Since $\nu\perp\mu$, $\theta b^+_{x,\nu}=b^-_{x,\nu}$ with the same
orientation. Denote by $U_\pm$ the bond variables in $H_\pm$ and by $\theta\cdot U_-$ the configuration
$b\mapsto(\theta\cdot U)(b)$ on the bonds of $H_+$, which depends on $U_-$ only. Then
\begin{equation}\label{6.1:eq-reflection-positivity-1}
A(U)=A_+(U_+)+A_-(U_-)+A_\times(U),
\qquad
A_\times(U):=\sum_{x,\nu}a_{p_{x,\nu}}(U),
\end{equation}
where $A_\pm(U):=\sum_{p\subset H_\pm}a_p(U)$, and $A_-(U_-)=A_+(\theta\cdot U_-)$, because $\theta$ maps
the plaquettes of $H_+$ onto those of $H_-$ and $a_p$ does not depend on the orientation of the loop.

\emph{The crossing kernel.} Let
\begin{equation*}
W_{x,\nu}:=U(\ell_x)\,U(b^+_{x,\nu})\,U(\ell_{x+\varepsilon e_\nu})^{-1}\,U(b^-_{x,\nu})^{-1}
\end{equation*}
be the holonomy of $p_{x,\nu}$ from the endpoint of $\ell_x$ in $H_-$. Then
$e^{-x_0a_{p_{x,\nu}}(U)}=e^{-x_0}\varphi(W_{x,\nu})$ with
$\varphi(W):=e^{x_0\operatorname{Re}\operatorname{tr}W}$. The function $\varphi$ is a $C^\infty$
positive-definite class function on $\mathrm{SU}(N_{\mathrm{col}})$: $\operatorname{Re}\operatorname{tr}W$
is a positive multiple of $\tfrac12(\chi_{\mathrm{f}}+\chi_{\bar{\mathrm{f}}})(W)$, where $\chi_{\mathrm{f}}$,
$\chi_{\bar{\mathrm{f}}}$ are the characters of the defining representation and of its conjugate, and is
therefore positive-definite, since every character is; sums, products (by
Schur's product theorem) and uniform limits of positive-definite functions are positive-definite, and
$\varphi$ is the uniform limit of the partial sums of its exponential series. Consequently
\begin{equation*}
\varphi=\sum_\lambda a_\lambda\chi_\lambda,\qquad a_\lambda\ge0,
\end{equation*}
with $\lambda$ running over the irreducible unitary representations $D^\lambda$ of
$\mathrm{SU}(N_{\mathrm{col}})$, $\chi_\lambda=\operatorname{Tr}D^\lambda$, $d_\lambda:=\dim D^\lambda$ (a
continuous positive-definite class function on a compact group is a series in the irreducible characters
with non-negative coefficients),
and $\sum_\lambda a_\lambda d_\lambda^k<\infty$ for every $k$, because the Fourier coefficients of a
$C^\infty$ function decay faster than any power of the Casimir eigenvalue and $d_\lambda$ is polynomial in
it. Since all matrix entries of unitary representations are bounded by $1$, every rearrangement below is
absolutely convergent, and Fubini's theorem applies on the configuration space, which is a finite product
of copies of the compact group $\mathrm{SU}(N_{\mathrm{col}})$.

By unitarity, for $\sigma_1,\dots,\sigma_4\in\mathrm{SU}(N_{\mathrm{col}})$,
\begin{equation*}
\chi_\lambda(\sigma_1\sigma_2\sigma_3^{-1}\sigma_4^{-1})
=\sum_{i,j,k,l}D^\lambda_{ij}(\sigma_1)\,D^\lambda_{jk}(\sigma_2)\,
\overline{D^\lambda_{lk}(\sigma_3)}\;\overline{D^\lambda_{il}(\sigma_4)}.
\end{equation*}
The functions $F(U)$ and $\overline{F(\theta\cdot U)}$ do not depend on the crossing links, since $F$ reads
bonds in the open half $H_+$. We therefore integrate the crossing links first. Expanding each
$\varphi(W_{x,\nu})$ with its own representation $\lambda_{x,\nu}$ and indices
$i_{x,\nu},j_{x,\nu},k_{x,\nu},l_{x,\nu}$, we get
\begin{equation}\label{6.1:eq-reflection-positivity-2}
\begin{split}
\mathcal{J}(U_+,U_-)&:=\int\prod_x dU(\ell_x)\prod_{x,\nu}\varphi(W_{x,\nu})\\
&=\sum_{\lambda(\cdot)}\prod_{x,\nu}a_{\lambda_{x,\nu}}\sum_{i,j,k,l}\Bigl[\prod_xI_x\Bigr]
\prod_{x,\nu}D^{\lambda_{x,\nu}}_{j_{x,\nu}k_{x,\nu}}(U(b^+_{x,\nu}))\,
\overline{D^{\lambda_{x,\nu}}_{i_{x,\nu}l_{x,\nu}}(U(b^-_{x,\nu}))}.
\end{split}
\end{equation}
Here the crossing link $\ell_x$ occurs as $\sigma_1$ in $p_{x,\nu}$ and as $\sigma_3$ in
$p_{x-\varepsilon e_\nu,\nu}$, for each of the three directions $\nu\neq\mu$ (both plaquettes exist on the
torus), so that
\begin{equation*}
I_x:=\int d\sigma\prod_{\nu\neq\mu}D^{\lambda_{x,\nu}}_{i_{x,\nu}j_{x,\nu}}(\sigma)\,
\overline{D^{\lambda_{x-\varepsilon e_\nu,\nu}}_{l_{x-\varepsilon e_\nu,\nu}k_{x-\varepsilon e_\nu,\nu}}(\sigma)}.
\end{equation*}

We now factor $I_x$. By invariance of Haar measure,
$\int\psi(\sigma)\,d\sigma=\iint\psi(\sigma'\sigma''^{-1})\,d\sigma'\,d\sigma''$ for integrable $\psi$, and
by unitarity
\begin{equation*}
D(\sigma'\sigma''^{-1})_{ij}=\sum_mD(\sigma')_{im}\overline{D(\sigma'')_{jm}},\qquad
\overline{D(\sigma'\sigma''^{-1})_{lk}}=\sum_{m'}\overline{D(\sigma')_{lm'}}\,D(\sigma'')_{km'}.
\end{equation*}
Inserting these into $I_x$, with summation indices $m=(m_\nu)_{\nu\neq\mu}$, $m'=(m'_\nu)_{\nu\neq\mu}$, the
integral splits into a $\sigma'$-integral and a $\sigma''$-integral, and
\begin{equation*}
I_x=\sum_{m,m'}u_x^{m,m'}(i,l)\,\overline{u_x^{m,m'}(j,k)},
\end{equation*}
\begin{equation*}
u_x^{m,m'}(i,l):=\int d\sigma\prod_{\nu\neq\mu}D^{\lambda_{x,\nu}}_{i_{x,\nu}m_\nu}(\sigma)\,
\overline{D^{\lambda_{x-\varepsilon e_\nu,\nu}}_{l_{x-\varepsilon e_\nu,\nu}m'_\nu}(\sigma)};
\end{equation*}
indeed the $\sigma''$-integral is the complex conjugate of the same expression with $(j,k)$ in place of
$(i,l)$. Put $\alpha:=(\lambda(\cdot),m(\cdot),m'(\cdot))$, $c_\alpha:=\prod_{x,\nu}a_{\lambda_{x,\nu}}\ge0$
and
\begin{equation*}
G_\alpha(U_+):=\sum_{j,k}\prod_x\overline{u_x^{m,m'}(j,k)}
\prod_{x,\nu}D^{\lambda_{x,\nu}}_{j_{x,\nu}k_{x,\nu}}(U(b^+_{x,\nu})),
\end{equation*}
a bounded function of the bond variables in $H_+$ (for fixed $\alpha$ the sum is finite and every factor is
bounded by $1$). In \eqref{6.1:eq-reflection-positivity-2} the sum over $(i,l)$ is
\begin{equation*}
\begin{split}
&\sum_{i,l}\prod_xu_x^{m,m'}(i,l)
\prod_{x,\nu}\overline{D^{\lambda_{x,\nu}}_{i_{x,\nu}l_{x,\nu}}(U(b^-_{x,\nu}))}\\
&\qquad=\overline{\sum_{i,l}\prod_x\overline{u_x^{m,m'}(i,l)}
\prod_{x,\nu}D^{\lambda_{x,\nu}}_{i_{x,\nu}l_{x,\nu}}(U(b^-_{x,\nu}))},
\end{split}
\end{equation*}
and $U(b^-_{x,\nu})=(\theta\cdot U)(b^+_{x,\nu})$, because $\theta b^+_{x,\nu}=b^-_{x,\nu}$ with the same
orientation. Renaming the summation indices $(i,l)$ as $(j,k)$, the right side is
$\overline{G_\alpha(\theta\cdot U_-)}$. Therefore
\begin{equation}\label{6.1:eq-reflection-positivity-3}
\mathcal{J}(U_+,U_-)=\sum_\alpha c_\alpha\,G_\alpha(U_+)\,\overline{G_\alpha(\theta\cdot U_-)},
\qquad c_\alpha\ge0.
\end{equation}
The crossing kernels at the two planes of $\Pi$ multiply, since they involve disjoint sets of crossing links
and of crossing plaquettes (the planes are $N/2\ge1>\varepsilon$ apart); the product of two sums of the form
\eqref{6.1:eq-reflection-positivity-3}, indexed by pairs of indices, is again of the form
$\sum_\alpha c_\alpha G_\alpha\overline{G_\alpha\circ\theta}$ with $c_\alpha\ge0$.

\emph{Conclusion.} Let $\#_\times$ be the number of crossing plaquettes. By
\eqref{6.1:eq-reflection-positivity-1}, $A_-(U_-)=A_+(\theta\cdot U_-)$, the reality of $A_+$,
and \eqref{6.1:eq-reflection-positivity-3},
\begin{equation*}
\begin{split}
\int\overline{F(\theta\cdot U)}\,F(U)\,e^{-x_0A(U)}\,dU
&=e^{-x_0\#_\times}\sum_\alpha c_\alpha\int dU_+\,dU_-\\
&\qquad\times\bigl[FG_\alpha e^{-x_0A_+}\bigr](U_+)\,
\overline{\bigl[FG_\alpha e^{-x_0A_+}\bigr](\theta\cdot U_-)}\\
&=e^{-x_0\#_\times}\sum_\alpha c_\alpha\Bigl|\int dU_+\,F\,G_\alpha\,e^{-x_0A_+}\Bigr|^2\ \ge\ 0.
\end{split}
\end{equation*}
In the second equality we used, in the integral over $U_-$, the change of variables
$U_-\mapsto\theta\cdot U_-$, which is a bijection of the bond variables in $H_-$ onto the bond variables in
$H_+$ preserving product Haar measure (it is a relabelling combined with inversions). Dividing by the
normalising constant of $\mu_{K,m}$, which is positive, gives (ii).
\end{proof}

\begin{remark}
The argument is the link reflection positivity of Osterwalder and Seiler~\cite{OSe78}. It uses of the gauge
group only that it is compact, and of the plaquette weight only that it is positive-definite, so it holds for
every compact gauge group and every positive-definite plaquette weight. Since it applies to every bounded
measurable function of the bond variables in $H_+$, it applies in particular to observables built from block
variables of bonds in $H_+$ and to polynomials in the bond variables of $H_+$.
\end{remark}

\begin{corollary}[The reflection form is positive semidefinite]\label{6.1:cor-reflection-form}
Let $\Pi$, $\theta$, $H_\pm$ and the reflected configuration $\theta\cdot U$ be as in
Lemma~\ref{6.1:lem-reflection-positivity}, and let $\langle\cdot\rangle$ denote expectation
with respect to $\mu_{K,m}$. For bounded measurable complex functions $F,G$ of the bond variables
in $H_+$ put
\begin{equation}\label{6.1:eq-reflection-form-1}
(F,G)_\theta := \bigl\langle \overline{F(\theta\cdot U)}\,G(U)\bigr\rangle .
\end{equation}
Then $(\cdot,\cdot)_\theta$ is a Hermitian positive-semidefinite sesquilinear form on the space
of such functions. Consequently
\begin{equation}\label{6.1:eq-reflection-form-2}
|(F,G)_\theta|^2 \le (F,F)_\theta\,(G,G)_\theta .
\end{equation}
Moreover, define $\theta\cdot U(b)$ for bonds $b\subset H_-$ by the same rule as for bonds in
$H_+$. If $F$ is a bounded measurable function of the bond variables in $H_-$, then
\begin{equation*}
\bigl\langle \overline{F(\theta\cdot U)}\,F(U)\bigr\rangle \ge 0 .
\end{equation*}
\end{corollary}

\begin{proof}
All expectations below exist because the integrands are bounded. The expression
\eqref{6.1:eq-reflection-form-1} is conjugate-linear in $F$ and linear in $G$, since
$F\mapsto\overline{F(\theta\cdot U)}$ is conjugate-linear and the expectation is linear. It is
therefore sesquilinear.

The form is Hermitian. Complex conjugation commutes with the expectation, so
\begin{equation*}
\overline{(F,G)_\theta} = \bigl\langle F(\theta\cdot U)\,\overline{G(U)}\bigr\rangle .
\end{equation*}
We now replace $U$ by $\theta\cdot U$. This substitution does not change the expectation, by
clause~(i) of Lemma~\ref{6.1:lem-reflection-positivity}, namely by the invariance of
$\mu_{K,m}$ under the isometry $\theta$ recorded there. We also use that the map
$U\mapsto\theta\cdot U$ is an involution. To see this, note that if $\theta$ preserves the
orientation of $b$, then it preserves that of $\theta b$, so that
$\theta\cdot(\theta\cdot U)(b)=U(\theta\theta b)=U(b)$. Otherwise both applications invert, and
$\theta\cdot(\theta\cdot U)(b)=(U(b)^{-1})^{-1}=U(b)$. Hence
\begin{equation*}
\bigl\langle F(\theta\cdot U)\,\overline{G(U)}\bigr\rangle
= \bigl\langle F(U)\,\overline{G(\theta\cdot U)}\bigr\rangle = (G,F)_\theta .
\end{equation*}

The form is positive semidefinite, because $(F,F)_\theta\ge 0$ for every admissible $F$ by
clause~(ii) of Lemma~\ref{6.1:lem-reflection-positivity}. The inequality
\eqref{6.1:eq-reflection-form-2} is the Cauchy--Schwarz inequality, which holds for every
Hermitian positive-semidefinite sesquilinear form.

Finally, let $F$ read only bonds of $H_-$, and put $G(U):=\overline{F(\theta\cdot U)}$. Since
$\theta$ maps bonds of $H_+$ to bonds of $H_-$, the function $G$ reads only bonds of $H_+$, so
the lemma applies to it. By the involution property, $\overline{G(\theta\cdot U)}=F(U)$.
Therefore
\begin{equation*}
\bigl\langle \overline{F(\theta\cdot U)}\,F(U)\bigr\rangle
= \bigl\langle G(U)\,\overline{G(\theta\cdot U)}\bigr\rangle = (G,G)_\theta \ge 0 .
\qedhere
\end{equation*}
\end{proof}

\begin{definition}[Four properties of the localised terms]\label{6.2:def-localised-terms}
Work under the standing hypotheses of the correlation theorem, in the last-lattice units and
with the plaquette functionals and sources fixed in Definition~\ref{6.1:not-units-sources}.
Here $K'$ is the number of renormalisation steps of the correlation theorem, so that the bare
spacing is $\varepsilon=L^{-K'}$ in these units and the levels run over $0\le k\le K'$. For a
complex weight $v$, that is, a complex function on $\mathbb{T}$ of which only the values $v(x_p)$
at the barycentres $x_p$ of the bare plaquettes $p$ enter, for $j\le K'$ and $X\in D_j$, let
$e^{(j)}(X;v)$ denote the
vacuum energy localised in $X$ that the correlation theorem produces at scale $j$ at the weight
$v$, and let
\begin{equation}\label{6.2:eq-localised-terms-1}
e_j(X;v):=e^{(j)}(X;v)-e^{(j)}(X;0)
\end{equation}
be the sourced vacuum-energy difference. Put
\begin{equation}\label{6.2:eq-localised-terms-2}
\mathcal{E}(v):=\sum_{j\le K'}\ \sum_{X\in D_j}e_j(X;v),
\qquad
\widetilde{Z}(v):=e^{-\mathcal{E}(v)}Z(v).
\end{equation}
For $j\le K'$ and $X\in D_j$ the following are the four properties of the localised terms,
together with the floor property (F); the constants $C_e$, $\delta$, $\kappa$, $E_+^{\sharp}$ and
$C_{\mathrm{low}}$ in them are those of the correlation theorem.

(L) \emph{Locality.} $e_j(X;v)$ is a function of $(v(x_p))_{p\in\operatorname{plaq}(X^+)}$ alone. In
particular, if $v(x_p)=v'(x_p)$ for every $p\in\operatorname{plaq}(X^+)$, then
\begin{equation}\label{6.2:eq-localised-terms-a}
e_j(X;v)=e_j(X;v'),
\end{equation}
and $e_j(X;v)=0$ if $v$ vanishes at the barycentres of $\operatorname{plaq}(X^+)$.

(A) \emph{Analyticity.}
\begin{equation}\label{6.2:eq-localised-terms-b}
v\longmapsto e_j(X;v)\ \text{ is holomorphic where defined.}
\end{equation}

(B) \emph{Bound.}
\begin{equation}\label{6.2:eq-localised-terms-c}
|e_j(X;v)|\le C_e\,e^{-(1-\delta)\kappa\, d_j(X)}.
\end{equation}

(S) \emph{Stability.}
\begin{equation}\label{6.2:eq-localised-terms-d}
|\widetilde{Z}(v)|=\bigl|e^{-\mathcal{E}(v)}Z(v)\bigr|\le e^{E_+^{\sharp}N^4}.
\end{equation}

(F) \emph{Floor.}
\begin{equation}\label{6.2:eq-localised-terms-e}
\widetilde{Z}(0)\ge e^{-C_{\mathrm{low}}N^4}.
\end{equation}
Moreover $\widetilde{Z}(0)=Z(0)$ and $E_+^{\sharp}+C_{\mathrm{low}}\ge 0$.

The domain of $v$ in (A), (B), (S) is taken in one of two readings. In the \emph{class reading},
(A), (B), (S) hold for every $v\in\mathcal{W}_{K'}$, the weight class of the correlation theorem
at the last level. In the \emph{literal reading}, (A), (B), (S) hold for
$v=\sum_l\zeta_lh_l$, where $(h_l)$ is any finite system of real Lipschitz functions on
$\mathbb{T}$ and $\zeta$ lies in that system's polydisc
$\{\sum_l|\zeta_l|\,\|h_l\|_{\sharp}<r\}$ (the coefficients $\zeta_l$ are complex source
coefficients, unrelated to the running shift $\zeta_k$). An index $\mathfrak{r}\in\{1,2\}$ records
the reading:
$\mathfrak{r}=1$ under the class reading, $\mathfrak{r}=2$ under the literal reading, and
\begin{equation}\label{6.2:eq-localised-terms-f}
\mathbb{D}_{\mathfrak{r}}:=\{z:\ \rho(z)<1/\mathfrak{r}\},
\qquad
\rho(z):=r^{-1}\sum_i|z_i|\,\|f_i\|_{\sharp},
\end{equation}
where $\rho(z)$ is the gauge of the source $z=(z_i)$ for the test functions $f_i$ of
Definition~\ref{6.1:not-units-sources} (this $\rho$ is unrelated to the effective densities
$\rho_k$), and $\mathbb{D}:=\mathbb{D}_1$ is the source polydisc.
All bounds of this chapter that rest on (A), (B), (S) at the weight $w_z$, or at weights built
from it, are asserted for $z\in\mathbb{D}_{\mathfrak{r}}$.
\end{definition}

\begin{proof}[Proof that the properties hold]
We take the properties in turn; each is read off from the correlation theorem.

(L). The correlation theorem states that for $k\le K'$ and every weight $v\in\mathcal{W}_k$ the
sourced effective density at level $k$ and weight $v$ has the inductive format
$(\mathrm{IH}^v)_k$, with terms
analytic in $v$ that depend on $v$ only through its values at the barycentres of
$\operatorname{plaq}(X^+)$. Hence $e^{(j)}(X;v)$ depends on $v$ only through
$(v(x_p))_{p\in\operatorname{plaq}(X^+)}$, and so does the difference
\eqref{6.2:eq-localised-terms-1}, since its second term does not depend on $v$. If $v$ and $v'$
agree at the barycentres of $\operatorname{plaq}(X^+)$, then \eqref{6.2:eq-localised-terms-a}
follows. If $v$ vanishes at these barycentres, then $v$ agrees there with the weight $0$, and by
\eqref{6.2:eq-localised-terms-a} and \eqref{6.2:eq-localised-terms-1}
\begin{equation*}
e_j(X;v)=e_j(X;0)=e^{(j)}(X;0)-e^{(j)}(X;0)=0.
\end{equation*}
For $v=w_z$ the torus bound states (L) in the form: $e_j(X;w_z)$ is independent of $z_i$ unless
$X^+\cap\operatorname{supp}f_i\neq\emptyset$, and vanishes at $z=0$.

(A). The correlation theorem states that the local terms of the inductive format at scale $j$
are analytic in the weight, uniformly on $\mathcal{W}_j(X)$, where they are defined by their
production formula, and that they depend on the weight only through its values at the
barycentres of $\operatorname{plaq}(X^+)$; and that the large-field
operations $\mathbf{R}$ and $\mathbf{B}$ of the sourced step are analytic in the weight,
uniformly on $\mathcal{W}_j(X)$, and depend on it only through its values on
$\operatorname{plaq}(X^+)$.
Since $\mathcal{W}_{K'}\subset\mathcal{W}_j$ for $j\le K'$, this gives
\eqref{6.2:eq-localised-terms-b} for $v\in\mathcal{W}_{K'}$, that is, in the class reading. The
torus bound states, as a consequence, that $e_j(X;\cdot)$ is analytic on $\mathbb{D}$; applied to
the system $(h_l)$, this gives \eqref{6.2:eq-localised-terms-b} in the literal reading.

(B). The torus bound states $|e_j|\le C_e\,e^{-(1-\delta)\kappa d_j(X)}$, which is
\eqref{6.2:eq-localised-terms-c}.

(S). The equality in \eqref{6.2:eq-localised-terms-d} is the definition of $\widetilde{Z}$ in
\eqref{6.2:eq-localised-terms-2}. In the torus bound, $\mathcal{E}$ is the double sum in
\eqref{6.2:eq-localised-terms-2} and $Z=e^{\mathcal{E}}\widetilde{Z}$, so its $\widetilde{Z}$ is
the same object; the torus bound states $|\widetilde{Z}|\le e^{E_+^{\sharp}N^4}$ on $\mathbb{D}$,
which is the inequality in \eqref{6.2:eq-localised-terms-d}.

(F). This is the floor statement of the correlation theorem: $\widetilde{Z}(0)$, the integral of
the effective density at level $K'$ and weight $0$, satisfies
$\widetilde{Z}(0)\ge e^{-C_{\mathrm{low}}N^4}$; in particular $\widetilde{Z}(0)$ is
a positive real number, so $|\widetilde{Z}(0)|=\widetilde{Z}(0)$. By (L), $e_j(X;0)=0$ for every
$j\le K'$ and $X\in D_j$, so $\mathcal{E}(0)=0$ by \eqref{6.2:eq-localised-terms-2}, and
$\widetilde{Z}(0)=e^{0}Z(0)=Z(0)$. The weight $0$ lies in the domain of (S) in either reading:
$0\in\mathcal{W}_{K'}$, and $z=0$ lies in $\mathbb{D}_{\mathfrak{r}}$. Hence (S) at $v=0$,
together with the floor, gives
\begin{equation*}
e^{-C_{\mathrm{low}}N^4}\le\widetilde{Z}(0)\le e^{E_+^{\sharp}N^4},
\end{equation*}
and, since $N^4>0$, taking logarithms gives $E_+^{\sharp}+C_{\mathrm{low}}\ge 0$.

Domains. That (A), (B), (S) hold on the domains of the class reading and of the literal reading,
and hence that all bounds resting on them hold for $z\in\mathbb{D}_{\mathfrak{r}}$ with
$\mathfrak{r}$ as in \eqref{6.2:eq-localised-terms-f}, is stated and proved in the remark on the
two readings of the domain of the weight, later in this section.
\end{proof}

\begin{definition}[Cells, the cell-reflection group and transport]\label{6.2:def-cell-reflections}
We work on the torus $\mathbb{T}=\mathbb{R}^4/N\mathbb{Z}^4$ of the last lattice $T^{(K')}$, in
last-lattice units, with the plaquette functionals and sources of
Notation~\ref{6.1:not-units-sources}. The bare spacing is $\varepsilon=L^{-K'}$ with
$K'\ge 0$, and the bare lattice is $\varepsilon(\mathbb{Z}+\tfrac12)^4$. Let $\ell$ be a power of $L$
such that $\ell$ divides $N/2$. Since $\ell$ and $L^m$ are powers of $L$ and $L^m$ divides $N/2$,
this holds as soon as $\ell\le L^m$, that is, as soon as $m\ge\log_L\ell$. Put $n_c:=N/\ell$. Since
$n_c=2\cdot(N/2)/\ell$, the number $n_c$ is an even integer with $n_c\ge 2$.

\smallskip
\emph{Cells.} For $k\in(\mathbb{Z}/n_c\mathbb{Z})^4$ put $c_k:=\ell k+[0,\ell]^4\subset\mathbb{T}$.
Let $\mathcal{C}$ be the set of these cells, so that $|\mathcal{C}|=n_c^4=N^4/\ell^4$. For
$c\in\mathcal{C}$ let $c^\circ$ be its interior, the open cell. The faces of the cells lie on the planes
$\{x_\mu\in\ell\mathbb{Z}\}$, and $\ell\mathbb{Z}\subset\mathbb{Z}\subset\varepsilon\mathbb{Z}$,
because $\ell a/\varepsilon=\ell aL^{K'}\in\mathbb{Z}$ for $a\in\mathbb{Z}$ and every $K'\ge0$. Bare
sites have coordinates in $\varepsilon(\mathbb{Z}+\frac12)$. These planes are therefore link planes:
no bare site lies on them. Let $\mathfrak{A}_c$ be the algebra of bounded measurable functions of the
bond variables $U(b)$, where $b$ ranges over the bare bonds $b\subset c^\circ$, that is, the bonds
with both endpoints in the open cell.

\smallskip
\emph{Face reflections.} For a coordinate direction $\mu$ and $a\in\mathbb{Z}/n_c\mathbb{Z}$ put
\begin{equation*}
\Pi_{\mu,a}:=\{x_\mu=a\ell\}\cup\{x_\mu=a\ell+N/2\}.
\end{equation*}
Let $\theta_{\mu,a}$ be the reflection of $\mathbb{T}$ in $\Pi_{\mu,a}$. It sends $x_\mu$ to
$2a\ell-x_\mu$ and leaves the coordinates $x_\nu$, $\nu\ne\mu$, fixed. Since
$a\ell+N/2=(a+n_c/2)\ell$, both planes of $\Pi_{\mu,a}$ are planes of cell faces, and
$\theta_{\mu,a}=\theta_{\mu,a+n_c/2}$. On cell indices, $\theta_{\mu,a}c_k=c_{k'}$ with
$k'_\mu=2a-1-k_\mu$ and $k'_\nu=k_\nu$ for $\nu\ne\mu$. The maps $\theta_{\mu,a}$ are called
\emph{face reflections}. The \emph{cell-reflection group} $\mathfrak{G}_\ell$ is the group of
isometries of $\mathbb{T}$ generated by all face reflections. For a face reflection
$\theta=\theta_{\mu,a}$, let $H_\pm(\theta)$ be the two open components of
$\mathbb{T}\setminus\Pi_{\mu,a}$, and let $\mathfrak{A}(H_\pm(\theta))$ be the bounded measurable
functions of the bond variables of the bonds in $H_\pm(\theta)$.

A \emph{lattice isometry} is an isometry of $\mathbb{T}$ that maps $\varepsilon(\mathbb{Z}+\frac12)^4$
onto itself. The following hold.
\begin{enumerate}
\item[(i)] For all $\mu$ and $a,a'$,
\begin{equation}\label{6.2:eq-cell-reflections-a}
\theta_{\mu,a}\theta_{\mu,a'}\,x=x+2(a-a')\ell\,e_\mu\qquad(x\in\mathbb{T}).
\end{equation}
The group $\mathfrak{G}_\ell$ is the direct product of four dihedral groups of order $n_c$, one for
each direction. It acts simply transitively on $\mathcal{C}$: for $c,c'\in\mathcal{C}$ there is a
unique $g_{c\to c'}\in\mathfrak{G}_\ell$ with $g_{c\to c'}c=c'$.
\item[(ii)] There is a well-defined homomorphism $\operatorname{sgn}\colon\mathfrak{G}_\ell\to\{\pm1\}$,
the \emph{parity}, with
$\operatorname{sgn}(g)=(-1)^r$ whenever $g$ is a product of $r$ face reflections.
For all $k,k'$,
\begin{equation}\label{6.2:eq-cell-reflections-b}
\operatorname{sgn}(g_{c_k\to c_{k'}})=(-1)^{\sum_\mu(k'_\mu-k_\mu)}.
\end{equation}
\item[(iii)] Every $g\in\mathfrak{G}_\ell$ is a lattice isometry. It maps cells to cells, $c^\circ$
onto $(gc)^\circ$, bonds in $c^\circ$ to bonds in $(gc)^\circ$, and plaquettes to plaquettes.
\item[(iv)] For every $c\in\mathcal{C}$,
\begin{equation}\label{6.2:eq-cell-reflections-c}
\#\{c'\in\mathcal{C}:\operatorname{sgn}(g_{c\to c'})=+1\}
=\#\{c'\in\mathcal{C}:\operatorname{sgn}(g_{c\to c'})=-1\}=\tfrac12|\mathcal{C}|.
\end{equation}
\end{enumerate}

\smallskip
\emph{Transport.} Each bare bond $b$ carries a fixed orientation, its positive orientation, along
which $U(b)$ is read; the same bond traversed in the opposite direction carries $U(b)^{-1}$. Let $g$
be a lattice
isometry. Put $s_g(b):=+1$ if $g$ carries the positive orientation of $b$ to that of the bond $gb$,
and $s_g(b):=-1$ otherwise. Define
\begin{equation*}
(U\circ g)(b):=U(gb)^{s_g(b)},\qquad T_gF(U):=F(U\circ g).
\end{equation*}
For a face reflection $\theta$ put $\Theta_\theta(F):=\overline{T_\theta F}$. Then the following hold.
\begin{enumerate}
\item[(v)] For lattice isometries $g,h$ and every bare bond $b$, $s_{gh}(b)=s_g(hb)s_h(b)$,
$(U\circ g)\circ h=U\circ(gh)$ and $T_gT_h=T_{gh}$. For $g\in\mathfrak{G}_\ell$ the map $T_g$ sends
$\mathfrak{A}_c$ into $\mathfrak{A}_{gc}$. Moreover $T_g$ commutes with complex conjugation,
$\|T_gF\|_\infty=\|F\|_\infty$, and $\langle T_gF\rangle=\langle F\rangle$.
\item[(vi)] $\Theta_\theta$ is antilinear, multiplicative and involutive, and it maps
$\mathfrak{A}_c$ into $\mathfrak{A}_{\theta c}$. Moreover $\langle\Theta_\theta(F)\cdot F\rangle\ge0$
for every $F\in\mathfrak{A}(H_+(\theta))$ and for every $F\in\mathfrak{A}(H_-(\theta))$.
\end{enumerate}

\smallskip
\emph{Chessboard copies.} For a complex number or function $\xi$ put $\xi^{(+)}:=\xi$ and
$\xi^{(-)}:=\overline{\xi}$, and read $\xi^{(\pm1)}$ as $\xi^{(\pm)}$. Let $c\in\mathcal{C}$ and
$F\in\mathfrak{A}_c$, regarded as based at $c$. For $c'\in\mathcal{C}$ put
\begin{equation}\label{6.2:eq-cell-reflections-1}
F^{[c']}:=(T_gF)^{(\operatorname{sgn}(g))},\qquad g:=g_{c\to c'}.
\end{equation}
Thus $F^{[c']}=T_gF$ if $\operatorname{sgn}(g)=+1$, and $F^{[c']}=\overline{T_gF}$ if
$\operatorname{sgn}(g)=-1$.
Then $F^{[c']}\in\mathfrak{A}_{c'}$ and $F^{[c]}=F$. Let $c',c''\in\mathcal{C}$, let $\theta$ be a
face reflection, let $F,G\in\mathfrak{A}_c$ and let $\lambda\in\mathbb{C}$. Copies of $F$, $G$,
$\overline{F}$, $\lambda F$, $FG$ and $1$ are based at $c$, and copies of $F^{[c']}$ are based at
$c'$. Then
\begin{gather}
\Theta_\theta\bigl(F^{[c']}\bigr)=F^{[\theta c']},\label{6.2:eq-cell-reflections-d}\\
\bigl(F^{[c']}\bigr)^{[c'']}=F^{[c'']},\label{6.2:eq-cell-reflections-e}\\
\begin{gathered}
(\overline{F})^{[c']}=\overline{F^{[c']}},\qquad
(\lambda F)^{[c']}=\lambda^{(\operatorname{sgn}(g_{c\to c'}))}F^{[c']},\\
(FG)^{[c']}=F^{[c']}G^{[c']},\qquad 1^{[c']}=1.
\end{gathered}\label{6.2:eq-cell-reflections-f}
\end{gather}
\end{definition}

\begin{proof}
We first prove (i). Fix $\mu$ and let $\mathfrak{G}^{(\mu)}_\ell$ be the subgroup generated by the
$\theta_{\mu,a}$, $a\in\mathbb{Z}/n_c\mathbb{Z}$. Its elements change only the coordinate $x_\mu$.
The map $\theta_{\mu,a'}$ sends $x_\mu$ to $2a'\ell-x_\mu$, and $\theta_{\mu,a}$ then sends this to
$2a\ell-2a'\ell+x_\mu$. This proves \eqref{6.2:eq-cell-reflections-a}.

Consider the translations $x_\mu\mapsto x_\mu+2t\ell$, $t\in\mathbb{Z}$, together with the
reflections $\theta_{\mu,a}$. This set is closed under composition. Two translations compose to a
translation. Two reflections compose to a translation by \eqref{6.2:eq-cell-reflections-a}. Composing
$\theta_{\mu,a}$ with the translation by $2t\ell$, in either order, gives the map
$x_\mu\mapsto 2(a\mp t)\ell-x_\mu$, that is, $\theta_{\mu,a-t}$ or $\theta_{\mu,a+t}$. The set
contains the inverse of each of its elements, and it contains the generators. Each translation by
$2t\ell$ equals $\theta_{\mu,t}\theta_{\mu,0}$, again by \eqref{6.2:eq-cell-reflections-a}. Hence
$\mathfrak{G}^{(\mu)}_\ell$ is exactly this set.

Translation by $2t\ell$ modulo $N=n_c\ell$ depends only on $t$ modulo $n_c/2$, so there are $n_c/2$
translations. The map $\theta_{\mu,a}$ depends only on $2a\ell$ modulo $N$, that is, on the residue
$2a$ modulo $n_c$, which takes $n_c/2$ values because $n_c$ is even; so there are $n_c/2$
distinct reflections. No translation $x_\mu\mapsto x_\mu+u$ equals a reflection
$x_\mu\mapsto v-x_\mu$, since equality would require $2x_\mu\equiv v-u$ modulo $N$ for every
$x_\mu$. Hence $|\mathfrak{G}^{(\mu)}_\ell|=n_c$. The translations form a cyclic group of order
$n_c/2$, generated by the translation by $2\ell$. Moreover $\theta_{\mu,0}$ conjugates the translation
by $2t\ell$ into the translation by $-2t\ell$. So $\mathfrak{G}^{(\mu)}_\ell$ is dihedral of order
$n_c$.

On the index $k_\mu\in\mathbb{Z}/n_c\mathbb{Z}$, the translation acts by $k_\mu\mapsto k_\mu+2t$ and
$\theta_{\mu,a}$ acts by $k_\mu\mapsto 2a-1-k_\mu$. The orbit of $0$ therefore contains every even
residue $2t$ and every odd residue $2a-1$. Since $n_c$ is even, these exhaust $\mathbb{Z}/n_c\mathbb{Z}$,
so there is one orbit, of size $n_c$. By the orbit--stabiliser relation the stabilisers have order
$n_c/n_c=1$. Hence $\mathfrak{G}^{(\mu)}_\ell$ acts simply transitively on $\mathbb{Z}/n_c\mathbb{Z}$.

Face reflections in different directions act on different coordinates and therefore commute.
Hence every $g\in\mathfrak{G}_\ell$ is a product $g_1g_2g_3g_4$ with
$g_\mu\in\mathfrak{G}^{(\mu)}_\ell$. Each factor $g_\mu$ is determined by $g$, as the action of $g$
on $x_\mu$. So $\mathfrak{G}_\ell$ is the direct product of the four groups
$\mathfrak{G}^{(\mu)}_\ell$, and it acts coordinatewise on $(\mathbb{Z}/n_c\mathbb{Z})^4$. A direct
product of simply transitive actions is simply transitive. Since $k\mapsto c_k$ is a bijection onto
$\mathcal{C}$, the action on $\mathcal{C}$ is simply transitive. This gives the unique
$g_{c\to c'}$.

Next we prove (ii). On $\mathfrak{G}^{(\mu)}_\ell$ put $\operatorname{sgn}_\mu:=+1$ on the
translations and
$\operatorname{sgn}_\mu:=-1$ on the reflections. By the description of $\mathfrak{G}^{(\mu)}_\ell$
obtained in
the proof of (i), these two sets are disjoint and exhaust the group,
so $\operatorname{sgn}_\mu$ is well defined. Every element of $\mathfrak{G}^{(\mu)}_\ell$ is an affine map
$x_\mu\mapsto\sigma x_\mu+\mathrm{const}$ of $\mathbb{R}/N\mathbb{Z}$, and $\operatorname{sgn}_\mu$
is its
coefficient $\sigma$. The coefficient of a composite of affine maps is the product of the
coefficients, so $\operatorname{sgn}_\mu$ is a homomorphism. Put
$\operatorname{sgn}(g):=\prod_\mu\operatorname{sgn}_\mu(g_\mu)$; this is a homomorphism on the
direct product. A
face reflection has one factor in the reflections and all other factors equal to the identity, so
its parity is $-1$. Hence $\operatorname{sgn}(g)=(-1)^r$ for every product of $r$ face reflections,
and this
value does not depend on the word chosen.

For \eqref{6.2:eq-cell-reflections-b}, let $g=g_{c_k\to c_{k'}}$, so that $g_\mu k_\mu=k'_\mu$ for
each $\mu$. Parity of residues modulo $n_c$ is meaningful because $n_c$ is even. If $g_\mu$ is a
translation, then $k'_\mu-k_\mu=2t$ is even. If $g_\mu=\theta_{\mu,a}$, then
$k'_\mu-k_\mu=2a-1-2k_\mu$ is odd. Thus $\operatorname{sgn}_\mu(g_\mu)=(-1)^{k'_\mu-k_\mu}$, and
taking the
product over $\mu$ gives \eqref{6.2:eq-cell-reflections-b}. Here the evenness of the cell count is
used: for odd $n_c$ the stabiliser of a cell index in one direction would have order greater than
one, and the parity of $k'_\mu-k_\mu$ modulo $n_c$ would not be defined.

Now (iii). Under $\theta_{\mu,a}$, the bare coordinate $x_\mu=(k+\frac12)\varepsilon$, $k\in\mathbb{Z}$,
goes to $2a\ell-x_\mu$, which equals $(k'+\frac12)\varepsilon$ with
\begin{equation*}
k'=2a\ell/\varepsilon-k-1\in\mathbb{Z},
\end{equation*}
because $a\ell\in\mathbb{Z}\subset\varepsilon\mathbb{Z}$. The other coordinates are unchanged, and
$N\in\varepsilon\mathbb{Z}$, so the reduction modulo $N$ is compatible. Hence $\theta_{\mu,a}$ maps
the bare lattice into itself, and onto itself since it is an involution. Compositions of lattice
isometries are lattice isometries, so every element of $\mathfrak{G}_\ell$ is one.

The map $\theta_{\mu,a}$ sends $[\ell k_\mu,\ell(k_\mu+1)]$ onto
$[\ell(2a-1-k_\mu),\ell(2a-k_\mu)]$, so face reflections, and hence all $g\in\mathfrak{G}_\ell$, map
cells onto cells. Being homeomorphisms, they map $c^\circ$ onto $(gc)^\circ$. A lattice isometry maps
pairs of bare sites at distance $\varepsilon$ to such pairs, and elementary squares of side
$\varepsilon$ of the bare lattice, that is, bare plaquettes, to such squares. Hence it maps bonds to
bonds and plaquettes to plaquettes. A bond with both endpoints
in $c^\circ$ goes to a bond with both endpoints in $(gc)^\circ$.

For (iv), let $c=c_k$. By \eqref{6.2:eq-cell-reflections-b}, the parity of $g_{c\to c_{k'}}$ is
$+1$ or $-1$ according as $\sum_\mu(k'_\mu-k_\mu)$ is even or odd; this is meaningful since $n_c$ is
even. Let $e_1$ denote the first unit vector of $(\mathbb{Z}/n_c\mathbb{Z})^4$. The map
$k'\mapsto k'+e_1$ is a bijection of $(\mathbb{Z}/n_c\mathbb{Z})^4$ that changes this parity. So
$c_{k'}\mapsto c_{k'+e_1}$ is a bijection between the two sets in
\eqref{6.2:eq-cell-reflections-c}. Their union is $\mathcal{C}$, so each has $\frac12|\mathcal{C}|$
elements.

Now (v). The lattice isometry $h$ carries the positive orientation of $b$ to $s_h(b)$ times the
positive orientation of $hb$. Then $g$ carries the positive orientation of $hb$ to $s_g(hb)$ times
that of $ghb$. Composing, $s_{gh}(b)=s_g(hb)s_h(b)$. Consequently
\begin{equation*}
((U\circ g)\circ h)(b)=(U\circ g)(hb)^{s_h(b)}
=U(ghb)^{s_g(hb)s_h(b)}=(U\circ(gh))(b).
\end{equation*}
Since $s_{\mathrm{id}}\equiv1$, we have $U\circ\mathrm{id}=U$. Hence $U\mapsto U\circ g$ is a
bijection of the configurations, with inverse $U\mapsto U\circ g^{-1}$. It is continuous, since it
permutes the coordinates and inverts some of them.

Let $F\in\mathfrak{A}_c$. Then $T_gF$ is bounded and measurable. It depends only on $(U\circ g)(b)$
for $b\subset c^\circ$, that is, on $U(gb)$, and by (iii) these bonds $gb$ lie in $(gc)^\circ$. So
$T_gF\in\mathfrak{A}_{gc}$. Next,
\begin{equation*}
T_g(T_hF)(U)=(T_hF)(U\circ g)=F((U\circ g)\circ h)=F(U\circ(gh))=T_{gh}F(U).
\end{equation*}
Conjugation commutes with $T_g$, because $T_g$ is composition with a map. The supremum norm is
preserved, because $U\mapsto U\circ g$ is onto. Finally, $\langle T_gF\rangle=\langle F\rangle$
because the bare Wilson probability measure, whose expectation is $\langle\cdot\rangle$, is
invariant under every isometry of $\mathbb{T}$ preserving the bare lattice
(Lemma~\ref{1.1:lem-invariance}; the invariance is recalled with clause (i) of
Lemma~\ref{6.1:lem-reflection-positivity}), and $g$ is such an isometry by (iii).

For (vi), $T_\theta$ is linear and multiplicative, while conjugation is antilinear and
multiplicative; hence $\Theta_\theta$ is antilinear and multiplicative. By
\eqref{6.2:eq-cell-reflections-a} with $a=a'$ we have $\theta^2=\mathrm{id}$, so by (v)
\begin{equation*}
\Theta_\theta(\Theta_\theta F)=\overline{T_\theta\overline{T_\theta F}}
=T_\theta T_\theta F=T_{\theta^2}F=F.
\end{equation*}
Also $\Theta_\theta$ maps $\mathfrak{A}_c$ into $\mathfrak{A}_{\theta c}$ by (v).

For the positivity, let $\theta=\theta_{\mu,a}$. Then $\Pi_{\mu,a}$ is the antipodal pair of link
planes of Lemma~\ref{6.1:lem-reflection-positivity}, with first plane at $x_\mu=a\ell\in\varepsilon\mathbb{Z}$,
and $\theta$ is the reflection of that lemma. For a bare bond $b$ contained in $H_+(\theta)$ or in
$H_-(\theta)$, the definition of $s_\theta$ gives $(U\circ\theta)(b)=U(\theta b)$ if $\theta$ preserves
the orientation of $b$, and $(U\circ\theta)(b)=U(\theta b)^{-1}$ otherwise. This is the orientation
rule defining $\theta\cdot U(b)$ in Lemma~\ref{6.1:lem-reflection-positivity}, used there for bonds
of $H_+(\theta)$ and in Corollary~\ref{6.1:cor-reflection-form} also for bonds of $H_-(\theta)$.
Hence $\Theta_\theta(F)(U)=\overline{F(\theta\cdot U)}$ for $F\in\mathfrak{A}(H_+(\theta))$ and for
$F\in\mathfrak{A}(H_-(\theta))$. The inequality
$\langle\Theta_\theta(F)\cdot F\rangle\ge0$ is then clause (ii) of
Lemma~\ref{6.1:lem-reflection-positivity}. For $F\in\mathfrak{A}(H_-(\theta))$ it is the last
assertion of Corollary~\ref{6.1:cor-reflection-form}.

Next the copies. By (i), $g=g_{c\to c'}$ exists and is unique. By (ii), $\operatorname{sgn}(g)$ is
defined.
By (v), $T_gF\in\mathfrak{A}_{gc}=\mathfrak{A}_{c'}$, and conjugation preserves $\mathfrak{A}_{c'}$.
So \eqref{6.2:eq-cell-reflections-1} is well defined and $F^{[c']}\in\mathfrak{A}_{c'}$. By
uniqueness $g_{c\to c}=\mathrm{id}$, whose parity is $+1$, and $T_{\mathrm{id}}F=F$; so $F^{[c]}=F$.
By \eqref{6.2:eq-cell-reflections-a}, the reflections in two adjacent parallel faces compose to a
translation by $2\ell$: $\theta_{\mu,a+1}\theta_{\mu,a}$ is the translation by $2\ell e_\mu$, and
$\theta_{\mu,a}\theta_{\mu,a+1}$ the translation by $-2\ell e_\mu$. By
\eqref{6.2:eq-cell-reflections-b}, $F^{[c_{k'}]}$ with $c=c_k$ is a transported copy of $F$ or of
$\overline{F}$ according as $(-1)^{\sum_\mu(k'_\mu-k_\mu)}$ equals $+1$ or $-1$.

We prove \eqref{6.2:eq-cell-reflections-d}. Let $g=g_{c\to c'}$. Then $\theta g\in\mathfrak{G}_\ell$
maps $c$ to $\theta c'$, so $\theta g=g_{c\to\theta c'}$ by uniqueness in (i). By (ii),
$\operatorname{sgn}(\theta g)=-\operatorname{sgn}(g)$. If $\operatorname{sgn}(g)=+1$, then by (v)
\begin{equation*}
\Theta_\theta(T_gF)=\overline{T_\theta T_gF}=\overline{T_{\theta g}F}=F^{[\theta c']},
\end{equation*}
since $\operatorname{sgn}(\theta g)=-1$. If $\operatorname{sgn}(g)=-1$, then, since $T_\theta$
commutes with
conjugation,
\begin{equation*}
\Theta_\theta\bigl(\overline{T_gF}\bigr)=T_\theta T_gF=T_{\theta g}F=F^{[\theta c']},
\end{equation*}
since $\operatorname{sgn}(\theta g)=+1$.

We prove \eqref{6.2:eq-cell-reflections-e}. Let $g=g_{c\to c'}$ and $h=g_{c'\to c''}$. Then $hg$
maps $c$ to $c''$, so $hg=g_{c\to c''}$, and
$\operatorname{sgn}(hg)=\operatorname{sgn}(h)\operatorname{sgn}(g)$. For
$\sigma,\tau\in\{\pm1\}$ we have $(\xi^{(\sigma)})^{(\tau)}=\xi^{(\sigma\tau)}$, and by (v)
$T_h(\xi^{(\sigma)})=(T_h\xi)^{(\sigma)}$. Therefore
\begin{align*}
\bigl(F^{[c']}\bigr)^{[c'']}
&=\bigl(T_h(T_gF)^{(\operatorname{sgn}(g))}\bigr)^{(\operatorname{sgn}(h))}
=(T_{hg}F)^{(\operatorname{sgn}(g)\operatorname{sgn}(h))}\\
&=(T_{hg}F)^{(\operatorname{sgn}(hg))}=F^{[c'']}.
\end{align*}
Thus the family of copies of $F^{[c']}$, based at $c'$, is exactly the family of copies of $F$.

Finally \eqref{6.2:eq-cell-reflections-f}. Let $g=g_{c\to c'}$ and $\sigma=\operatorname{sgn}(g)$. The
operator $T_g$ is linear and multiplicative, commutes with conjugation, and satisfies $T_g1=1$;
conjugation is multiplicative and fixes $1$. Hence
\begin{equation*}
(T_g\overline{F})^{(\sigma)}=\overline{(T_gF)^{(\sigma)}},\qquad
(\lambda T_gF)^{(\sigma)}=\lambda^{(\sigma)}(T_gF)^{(\sigma)},
\end{equation*}
and
\begin{equation*}
(T_g(FG))^{(\sigma)}=(T_gF)^{(\sigma)}(T_gG)^{(\sigma)},\qquad (T_g1)^{(\sigma)}=1.
\end{equation*}
These are the four identities of \eqref{6.2:eq-cell-reflections-f}.
\end{proof}

\begin{lemma}[Positivity of the chessboard value]\label{6.2:lem-chessboard-value}
Let $c_0\in\mathcal{C}$, let $F\in\mathfrak{A}_{c_0}$, and let $F^{[c']}$, $c'\in\mathcal{C}$, be its chessboard
copies (Definition~\ref{6.2:def-cell-reflections}). Put
\begin{equation}\label{6.2:eq-chessboard-value-1}
\mathsf{c}(F):=\Bigl\langle\prod_{c'\in\mathcal{C}}F^{[c']}\Bigr\rangle .
\end{equation}
Then $\mathsf{c}(F)$ is real and $\mathsf{c}(F)\ge 0$. Consequently
$\|F\|_{\mathrm{ch}}:=\mathsf{c}(F)^{1/|\mathcal{C}|}$, the non-negative root, is defined. Moreover, for
every $c'\in\mathcal{C}$ and every $\lambda\in\mathbb{C}$,
\begin{equation}\label{6.2:eq-chessboard-value-2}
\|F^{[c']}\|_{\mathrm{ch}}=\|F\|_{\mathrm{ch}},\qquad
\|\overline{F}\|_{\mathrm{ch}}=\|F\|_{\mathrm{ch}},\qquad
\|\lambda F\|_{\mathrm{ch}}=|\lambda|\,\|F\|_{\mathrm{ch}},\qquad
\|1\|_{\mathrm{ch}}=1 .
\end{equation}
\end{lemma}

\begin{proof}
Fix a face reflection $\theta=\theta_{\mu,a}$, the reflection of $\mathbb{T}$ in
$\Pi_{\mu,a}=\{x_\mu=a\ell\}\cup\{x_\mu=a\ell+N/2\}$. Since $a\ell\in\mathbb{Z}\subset\varepsilon\mathbb{Z}$,
$\Pi_{\mu,a}$ is a pair of antipodal link planes of the kind considered in
Lemma~\ref{6.1:lem-reflection-positivity}. Let $H_\pm(\theta)$ be the two open components of
$\mathbb{T}\setminus\Pi_{\mu,a}$.

\emph{Every cell lies in exactly one half.} Both planes of $\Pi_{\mu,a}$ are face planes, that is,
lie in $\{x_\mu\in\ell\mathbb{Z}\}$. On the interior $c_k^{\circ}$ of a cell the coordinate $x_\mu$ ranges over
the open interval $(\ell k_\mu,\ell k_\mu+\ell)$, which contains no point of $\ell\mathbb{Z}$. Hence
$c_k^{\circ}$ does not meet $\Pi_{\mu,a}$; being connected, it lies in exactly one of $H_+(\theta)$,
$H_-(\theta)$. Let $\mathcal{C}_\pm(\theta)$ be the set of cells whose interior lies in $H_\pm(\theta)$. Then
$\mathcal{C}$ is the disjoint union of $\mathcal{C}_+(\theta)$ and $\mathcal{C}_-(\theta)$. The reflection
$\theta$ fixes $\Pi_{\mu,a}$ setwise and reverses the coordinate $x_\mu-a\ell$, so it exchanges
$H_+(\theta)$ and $H_-(\theta)$. It maps cells to cells and open cells to open cells, and it is an
involution. Therefore $c'\mapsto\theta c'$ is a bijection of $\mathcal{C}_+(\theta)$ onto
$\mathcal{C}_-(\theta)$.

\emph{The splitting.} Define
\begin{equation*}
G_\pm:=\prod_{c'\in\mathcal{C}_\pm(\theta)}F^{[c']} .
\end{equation*}
Each $F^{[c']}$ is $T_gF$ or $\overline{T_gF}$ with $g=g_{c_0\to c'}$, hence lies in $\mathfrak{A}_{c'}$. Its
bonds lie in $c'^{\circ}$, and for $c'\in\mathcal{C}_\pm(\theta)$ these bonds lie in $H_\pm(\theta)$. A
finite product of bounded measurable functions is bounded and measurable. Hence $G_\pm$ belongs to
$\mathfrak{A}(H_\pm(\theta))$, the algebra of bounded measurable functions of the bonds in
$H_\pm(\theta)$. Since $\mathcal{C}$ is the disjoint union of the two halves, we have
$\prod_{c'\in\mathcal{C}}F^{[c']}=G_-G_+$. The operation $\Theta_\theta$ is multiplicative. Using in
turn \eqref{6.2:eq-cell-reflections-d} and the bijection $c'\mapsto\theta c'$, we obtain
\begin{equation*}
\Theta_\theta(G_+)=\prod_{c'\in\mathcal{C}_+(\theta)}\Theta_\theta\bigl(F^{[c']}\bigr)
=\prod_{c'\in\mathcal{C}_+(\theta)}F^{[\theta c']}
=\prod_{c''\in\mathcal{C}_-(\theta)}F^{[c'']}=G_- .
\end{equation*}
Therefore $\mathsf{c}(F)=\langle\Theta_\theta(G_+)\,G_+\rangle$. By reflection positivity of the
lattice measure (Lemma~\ref{6.1:lem-reflection-positivity}), in the form
$\langle\Theta_\theta(G)\,G\rangle\ge 0$ for $G\in\mathfrak{A}(H_+(\theta))$, this quantity is real
and non-negative.

\emph{Invariance under copies.} We have $F^{[c']}\in\mathfrak{A}_{c'}$, and its chessboard copies are
taken with base cell $c'$. By \eqref{6.2:eq-cell-reflections-e} they are
$(F^{[c']})^{[c'']}=F^{[c'']}$ for every $c''\in\mathcal{C}$. Hence
$\mathsf{c}(F^{[c']})=\mathsf{c}(F)$, and the first identity in \eqref{6.2:eq-chessboard-value-2}
follows.

\emph{Conjugation and scalars.} Let $g=g_{c_0\to c'}$. Since $T_g$ commutes with complex conjugation,
the definition of the copies gives $(\overline{F})^{[c']}=\overline{T_gF}=\overline{F^{[c']}}$ when
$\varepsilon(g)=+1$, and $(\overline{F})^{[c']}=T_gF=\overline{F^{[c']}}$ when $\varepsilon(g)=-1$.
Similarly, since $T_g$ is linear, $(\lambda F)^{[c']}=\lambda F^{[c']}$ when $\varepsilon(g)=+1$, and
$(\lambda F)^{[c']}=\overline{\lambda}\,F^{[c']}$ when $\varepsilon(g)=-1$. For conjugation this gives
\begin{equation*}
\mathsf{c}(\overline{F})=\Bigl\langle\,\overline{\textstyle\prod_{c'}F^{[c']}}\,\Bigr\rangle
=\overline{\mathsf{c}(F)}=\mathsf{c}(F),
\end{equation*}
because $\mathsf{c}(F)$ is real. For scalars, by \eqref{6.2:eq-cell-reflections-c} exactly
$|\mathcal{C}|/2$ cells $c'$ have $\varepsilon(g_{c_0\to c'})=+1$ and the other $|\mathcal{C}|/2$ have
$\varepsilon(g_{c_0\to c'})=-1$. Hence
\begin{equation*}
\mathsf{c}(\lambda F)=\lambda^{|\mathcal{C}|/2}\,\overline{\lambda}^{\,|\mathcal{C}|/2}\,\mathsf{c}(F)
=|\lambda|^{|\mathcal{C}|}\,\mathsf{c}(F).
\end{equation*}
Taking non-negative $|\mathcal{C}|$-th roots gives the second and third identities in
\eqref{6.2:eq-chessboard-value-2}.

\emph{The constant function.} We have $T_g1=1$ and $\overline{1}=1$, so $1^{[c']}=1$ for every $c'$.
Since $\mu_{K,m}$ is a probability measure, $\mathsf{c}(1)=\langle 1\rangle=1$, and hence
$\|1\|_{\mathrm{ch}}=1$.
\end{proof}

\begin{lemma}[The chessboard estimate for complex cell observables]\label{6.2:lem-chessboard-estimate}
For every family $(F_c)_{c\in\mathcal{C}}$ with $F_c\in\mathfrak{A}_c$ for each $c\in\mathcal{C}$,
\begin{equation*}
\Bigl|\Bigl\langle\prod_{c\in\mathcal{C}}F_c\Bigr\rangle\Bigr|\le\prod_{c\in\mathcal{C}}\|F_c\|_{\mathrm{ch}}.
\end{equation*}
In particular, if $c_0\in\mathcal{C}$ and $F\in\mathfrak{A}_{c_0}$, then $|\langle F\rangle|\le\|F\|_{\mathrm{ch}}$.
\end{lemma}

\begin{proof}
\emph{Notation.} A \emph{cell array} is a family $\mathsf{X}=(\mathsf{X}_c)_{c\in\mathcal{C}}$ with
$\mathsf{X}_c\in\mathfrak{A}_c$ for every $c\in\mathcal{C}$; we put
\begin{equation*}
\mathsf{a}(\mathsf{X}):=\Bigl\langle\prod_{c\in\mathcal{C}}\mathsf{X}_c\Bigr\rangle .
\end{equation*}
Let $\theta=\theta_{\mu,a}$ be a face reflection and let $H_\pm(\theta)$ be the two open components of
$\mathbb{T}\setminus\Pi_{\mu,a}$. The set $\Pi_{\mu,a}=\{x_\mu=a\ell\}\cup\{x_\mu=a\ell+N/2\}$ is an antipodal pair
of link planes, since $a\ell\in\mathbb{Z}\subset\varepsilon\mathbb{Z}$, so Lemma~\ref{6.1:lem-reflection-positivity}
and Corollary~\ref{6.1:cor-reflection-form} apply to the reflection $\theta$. The faces of the cells lie on the
planes $x_\mu\in\ell\mathbb{Z}$, and the interior
of a cell meets none of them, so the interior of every cell lies in exactly one of $H_+(\theta)$, $H_-(\theta)$.
Let $\mathcal{C}_\pm(\theta)$ be the set of cells whose interior lies in $H_\pm(\theta)$, and let
$\mathfrak{A}(H_\pm(\theta))$ be the bounded measurable functions of the bond variables of bonds in $H_\pm(\theta)$.
Since $\theta$ is an isometry of $\mathbb{T}$ that maps cells to cells and exchanges the two halves, it maps
$\mathcal{C}_+(\theta)$ bijectively onto $\mathcal{C}_-(\theta)$. We use the following properties of the antilinear
reflection, all contained in Definition~\ref{6.2:def-cell-reflections}: $\Theta_\theta$ is antilinear, multiplicative
and involutive, preserves the sup norm, and maps $\mathfrak{A}_c$ into $\mathfrak{A}_{\theta c}$ and
$\mathfrak{A}(H_\mp(\theta))$ into $\mathfrak{A}(H_\pm(\theta))$. Unwinding the definitions of the transport $T_\theta$
and of the reflected configuration $\theta\cdot U$ of Lemma~\ref{6.1:lem-reflection-positivity}, one has
$\Theta_\theta(F)(U)=\overline{F(\theta\cdot U)}$, so that the reflection form of
Corollary~\ref{6.1:cor-reflection-form} reads
\begin{equation}\label{6.2:eq-chessboard-estimate-1}
(F,G)_\theta=\langle\Theta_\theta(F)\,G\rangle,\qquad F,G\in\mathfrak{A}(H_+(\theta)).
\end{equation}
For a cell array $\mathsf{X}$ define the two \emph{mirrored cell arrays} $\mathsf{X}^{\theta,+}$ and
$\mathsf{X}^{\theta,-}$ by
\begin{equation*}
\mathsf{X}^{\theta,\pm}_c:=
\begin{cases}
\mathsf{X}_c, & c\in\mathcal{C}_\pm(\theta),\\
\Theta_\theta(\mathsf{X}_{\theta c}), & c\in\mathcal{C}_\mp(\theta).
\end{cases}
\end{equation*}
These are cell arrays: for $c\in\mathcal{C}_\mp(\theta)$ we have $\mathsf{X}_{\theta c}\in\mathfrak{A}_{\theta c}$, hence
$\Theta_\theta(\mathsf{X}_{\theta c})\in\mathfrak{A}_{\theta\theta c}=\mathfrak{A}_c$.

\emph{The Schwarz step.} We show that for every cell array $\mathsf{X}$ and every face reflection $\theta$ the numbers
$\mathsf{a}(\mathsf{X}^{\theta,+})$ and $\mathsf{a}(\mathsf{X}^{\theta,-})$ are real and non-negative and
\begin{equation}\label{6.2:eq-chessboard-estimate-2}
|\mathsf{a}(\mathsf{X})|^2\le\mathsf{a}(\mathsf{X}^{\theta,+})\,\mathsf{a}(\mathsf{X}^{\theta,-}).
\end{equation}
Put $G_\pm:=\prod_{c\in\mathcal{C}_\pm(\theta)}\mathsf{X}_c$. Each factor reads only bonds in the open cell $c$, which
lies in $H_\pm(\theta)$, so $G_\pm\in\mathfrak{A}(H_\pm(\theta))$; consequently
$\Theta_\theta(G_-)\in\mathfrak{A}(H_+(\theta))$.
Since every cell lies in exactly one of $\mathcal{C}_\pm(\theta)$ and $\Theta_\theta$ is involutive,
\begin{equation*}
\mathsf{a}(\mathsf{X})=\langle G_-G_+\rangle
=\bigl\langle\Theta_\theta\bigl(\Theta_\theta(G_-)\bigr)\,G_+\bigr\rangle
=\bigl(\Theta_\theta(G_-),G_+\bigr)_\theta ,
\end{equation*}
the last equality by \eqref{6.2:eq-chessboard-estimate-1}. The Cauchy--Schwarz inequality of
Corollary~\ref{6.1:cor-reflection-form} therefore gives
\begin{align*}
|\mathsf{a}(\mathsf{X})|^2
&\le\bigl(\Theta_\theta(G_-),\Theta_\theta(G_-)\bigr)_\theta\,(G_+,G_+)_\theta\\
&=\bigl\langle G_-\,\Theta_\theta(G_-)\bigr\rangle\,\bigl\langle\Theta_\theta(G_+)\,G_+\bigr\rangle ,
\end{align*}
where we used \eqref{6.2:eq-chessboard-estimate-1} once more together with
$\Theta_\theta(\Theta_\theta(G_-))=G_-$. Both factors on the right are values of the positive-semidefinite form
$(\cdot,\cdot)_\theta$ on the diagonal, hence real and non-negative by Lemma~\ref{6.1:lem-reflection-positivity} and
Corollary~\ref{6.1:cor-reflection-form}. It remains to identify them. Since $c\mapsto\theta c$ is a bijection of
$\mathcal{C}_+(\theta)$ onto $\mathcal{C}_-(\theta)$ and $\Theta_\theta$ is multiplicative,
\begin{equation*}
\prod_{c\in\mathcal{C}_+(\theta)}\Theta_\theta(\mathsf{X}_{\theta c})
=\Theta_\theta\Bigl(\prod_{c''\in\mathcal{C}_-(\theta)}\mathsf{X}_{c''}\Bigr)=\Theta_\theta(G_-),
\end{equation*}
so $\prod_{c\in\mathcal{C}}\mathsf{X}^{\theta,-}_c=G_-\,\Theta_\theta(G_-)$ and
$\langle G_-\Theta_\theta(G_-)\rangle=\mathsf{a}(\mathsf{X}^{\theta,-})$. In the same way
$\prod_{c\in\mathcal{C}_-(\theta)}\Theta_\theta(\mathsf{X}_{\theta c})=\Theta_\theta(G_+)$, whence
$\langle\Theta_\theta(G_+)G_+\rangle=\mathsf{a}(\mathsf{X}^{\theta,+})$. This proves
\eqref{6.2:eq-chessboard-estimate-2}.

\emph{The spreading sequence.} Let $c_0\in\mathcal{C}$ and $F\in\mathfrak{A}_{c_0}$, and let $Y^0$ be a cell array
with $Y^0_{c_0}=F=F^{[c_0]}$. We construct a finite sequence $Y^0,Y^1,\dots,Y^T$ of cell arrays, with
$Y^{i+1}=(Y^i)^{\theta_i,\sigma_i}$ for a face reflection $\theta_i$ and a sign $\sigma_i\in\{+,-\}$, such that
$Y^T_{c'}=F^{[c']}$ for every $c'\in\mathcal{C}$; that is, $Y^T$ is the chessboard array of $F$ and
$\mathsf{a}(Y^T)=\mathsf{c}(F)$. Write $c_0=c_{k^0}$. For $\mu\in\{1,2,3,4\}$ and $t\in\mathbb{Z}/n_c\mathbb{Z}$ let
$B^{(\mu)}_t$ be the set of cells $c_k$ with $k_\mu=t$ and $k_\nu=k^0_\nu$ for all $\nu>\mu$ (the indices $k_\nu$ with
$\nu<\mu$ being arbitrary); thus $B^{(1)}_t$ is a single cell of the direction-one row of $c_0$, for $\mu\ge2$ one has
$B^{(\mu)}_{k^0_\mu}=\bigcup_tB^{(\mu-1)}_t$, and $\bigcup_tB^{(4)}_t=\mathcal{C}$. We say that a set of cells
\emph{carries copies}
in a cell array $Y$ if $Y_{c'}=F^{[c']}$ for every cell $c'$ in it.

We work through the directions $\mu=1,2,3,4$ in this order. At the start of direction $\mu$ the block
$B^{(\mu)}_{k^0_\mu}$ carries copies (for $\mu=1$ because $Y^0_{c_0}=F^{[c_0]}$; for $\mu\ge2$ because
$B^{(\mu)}_{k^0_\mu}=\bigcup_tB^{(\mu-1)}_t$ is the union filled in direction $\mu-1$). Let
$S\subset\mathbb{Z}/n_c\mathbb{Z}$ be the set of indices $t$ such that $B^{(\mu)}_t$ carries copies; initially
$S\supseteq\{k^0_\mu\}$, and we follow an arc $\{b,b+1,\dots,b+s-1\}\subseteq S$ of the cycle
$\mathbb{Z}/n_c\mathbb{Z}$ (of even length $n_c$), starting with $s=1$, $b=k^0_\mu$. Recall from
Definition~\ref{6.2:def-cell-reflections} that $\theta_{\mu,a}$ acts on cell indices by
$k_\mu\mapsto2a-1-k_\mu$, fixing $k_\nu$ for $\nu\neq\mu$. Since $c_k=\ell k+[0,\ell]^4$ and
$\Pi_{\mu,a}=\{x_\mu=a\ell\}\cup\{x_\mu=(a+n_c/2)\ell\}$, the cells in the two halves of $\theta_{\mu,a}$ are those
with $k_\mu\in\{a,\dots,a+n_c/2-1\}$ and those with $k_\mu\in\{a-n_c/2,\dots,a-1\}$. Such a reflection maps
$B^{(\mu)}_t$ onto
$B^{(\mu)}_{2a-1-t}$, and by the key identity \eqref{6.2:eq-cell-reflections-d},
$\Theta_\theta(F^{[c']})=F^{[\theta c']}$, the mirror of a block carrying copies carries copies.
\begin{itemize}
\item If $s\le n_c/2$, take $\theta=\theta_{\mu,b}$; the arc $\{b,\dots,b+s-1\}$ lies in the half
$\{b,\dots,b+n_c/2-1\}$, which has $n_c/2\ge s$ consecutive indices. Take the mirrored cell array keeping that
half. The blocks with index in the arc keep their entries, and the blocks with index in the reflected arc
$\{b-s,\dots,b-1\}$ receive $\Theta_\theta(F^{[\theta c']})=F^{[c']}$ by \eqref{6.2:eq-cell-reflections-d}. Hence the
arc $\{b-s,\dots,b+s-1\}$, of length $2s$ (the whole cycle if $2s=n_c$), indexes blocks carrying copies.
\item If $n_c/2<s<n_c$, the arc contains $\{b,\dots,b+n_c/2-1\}$, which is a half of $\theta_{\mu,b}$. Mirroring while
keeping that half, the blocks with index in the other half $\{b-n_c/2,\dots,b-1\}$ receive $F^{[c']}$ again by
\eqref{6.2:eq-cell-reflections-d}; so all $n_c$ blocks carry copies.
\end{itemize}
Each step strictly increases $s$ until $s=n_c$, so after finitely many steps all blocks $B^{(\mu)}_t$ carry copies,
and we pass to direction $\mu+1$. After direction $4$ every cell of $\mathcal{C}$ carries copies; every other entry of
$Y^0$ has been overwritten, and the terminal cell array is the chessboard array of $F$, with value
$\mathsf{a}(Y^T)=\mathsf{c}(F)$.

\emph{Vanishing chessboard norms.} Suppose $\|F_{c_0}\|_{\mathrm{ch}}=0$ for some $c_0\in\mathcal{C}$, and put
$F:=F_{c_0}$ and
\begin{equation*}
\mathsf{b}:=\max\Bigl(1,\max_{c\in\mathcal{C}}\|F_c\|_\infty\Bigr)^{|\mathcal{C}|}.
\end{equation*}
Apply the spreading sequence with $Y^0:=(F_c)_{c\in\mathcal{C}}$. Every entry of every cell array met (including
the branches not followed) is obtained from some $F_c$ by transports and complex conjugations, which preserve the sup
norm; so each of the $|\mathcal{C}|$ factors has sup norm at most $\max(1,\max_c\|F_c\|_\infty)$ and $|\mathsf{a}|\le\mathsf{b}$
on all these cell arrays. By \eqref{6.2:eq-chessboard-estimate-2}, with the non-followed branch bounded by
$\mathsf{b}$ and non-negative,
\begin{equation*}
|\mathsf{a}(Y^i)|^2\le\mathsf{a}(Y^{i+1})\,\mathsf{a}\bigl((Y^i)^{\theta_i,-\sigma_i}\bigr)
\le\mathsf{b}\,\mathsf{a}(Y^{i+1}),\qquad 0\le i<T .
\end{equation*}
By induction on $i$, $|\mathsf{a}(Y^0)|^{2^i}\le\mathsf{b}^{2^i-1}|\mathsf{a}(Y^i)|$: the case $i=0$ is trivial, and
\begin{equation*}
|\mathsf{a}(Y^0)|^{2^{i+1}}\le\mathsf{b}^{2^{i+1}-2}|\mathsf{a}(Y^i)|^2
\le\mathsf{b}^{2^{i+1}-1}\mathsf{a}(Y^{i+1}).
\end{equation*}
For $i=T$ this gives
\begin{equation*}
|\mathsf{a}(Y^0)|^{2^T}\le\mathsf{b}^{2^T-1}\,\mathsf{c}(F)=\mathsf{b}^{2^T-1}\,\|F\|_{\mathrm{ch}}^{|\mathcal{C}|}=0 .
\end{equation*}
Hence $\mathsf{a}(Y^0)=0\le\prod_c\|F_c\|_{\mathrm{ch}}$, and the lemma holds. From now on we assume
$\|F_c\|_{\mathrm{ch}}>0$ for all $c\in\mathcal{C}$.

\emph{A maximiser.} Let
\begin{equation*}
\mathcal{S}:=\bigl\{F_c^{[c']}:\ c,c'\in\mathcal{C}\bigr\}.
\end{equation*}
This set is finite. It is closed under chessboard copying, since $(F_c^{[c']})^{[c'']}=F_c^{[c'']}$ by
\eqref{6.2:eq-cell-reflections-e}, and under every $\Theta_\theta$, since $\Theta_\theta(F_c^{[c']})=F_c^{[\theta c']}$
by \eqref{6.2:eq-cell-reflections-d}. By Lemma~\ref{6.2:lem-chessboard-value},
$\|F_c^{[c']}\|_{\mathrm{ch}}=\|F_c\|_{\mathrm{ch}}$, so every element of $\mathcal{S}$ has chessboard norm in
$\{\|F_c\|_{\mathrm{ch}}:c\in\mathcal{C}\}\subset\,]0,\infty[$. Call a cell array $Y$ \emph{admissible} if
$Y_{c'}\in\mathcal{S}\cap\mathfrak{A}_{c'}$ for every $c'$; there are finitely many, and
$\mathsf{X}:=(F_c)_{c\in\mathcal{C}}$ is admissible because $F_c=F_c^{[c]}$. For admissible $Y$ put
\begin{equation*}
R(Y):=\frac{|\mathsf{a}(Y)|}{\prod_{c'\in\mathcal{C}}\|Y_{c'}\|_{\mathrm{ch}}},\qquad
R_{\max}:=\max_{Y\ \mathrm{admissible}}R(Y)\ge R(\mathsf{X}).
\end{equation*}
Let $\theta$ be a face reflection and $Y$ admissible. Then $Y^{\theta,\pm}$ are admissible, because
$\Theta_\theta(Y_{\theta c})\in\mathcal{S}\cap\mathfrak{A}_c$. Moreover $\|\Theta_\theta Z\|_{\mathrm{ch}}=\|Z\|_{\mathrm{ch}}$
for $Z\in\mathcal{S}$: if $Z=F_{c_1}^{[c']}$ then $\Theta_\theta Z=F_{c_1}^{[\theta c']}$ by
\eqref{6.2:eq-cell-reflections-d}, and both have norm $\|F_{c_1}\|_{\mathrm{ch}}$ by \eqref{6.2:eq-cell-reflections-e}
and Lemma~\ref{6.2:lem-chessboard-value}. Using that $\theta$ maps $\mathcal{C}_\mp(\theta)$ bijectively onto
$\mathcal{C}_\pm(\theta)$, we obtain
\begin{equation*}
\prod_{c'}\|Y^{\theta,\pm}_{c'}\|_{\mathrm{ch}}=\Bigl(\prod_{c\in\mathcal{C}_\pm(\theta)}\|Y_c\|_{\mathrm{ch}}\Bigr)^2,
\end{equation*}
and hence
\begin{equation*}
\prod_{c'}\|Y^{\theta,+}_{c'}\|_{\mathrm{ch}}\cdot\prod_{c'}\|Y^{\theta,-}_{c'}\|_{\mathrm{ch}}
=\Bigl(\prod_{c'\in\mathcal{C}}\|Y_{c'}\|_{\mathrm{ch}}\Bigr)^2 .
\end{equation*}
Dividing \eqref{6.2:eq-chessboard-estimate-2} by this identity gives
\begin{equation}\label{6.2:eq-chessboard-estimate-3}
R(Y)^2\le R(Y^{\theta,+})\,R(Y^{\theta,-})\le R_{\max}^2 .
\end{equation}
We have $R_{\max}>0$: the chessboard array $(F_{c_0}^{[c']})_{c'}$ of any $F_{c_0}$ is admissible, and since
$\mathsf{c}(F_{c_0})\ge0$ by Lemma~\ref{6.2:lem-chessboard-value}, its $R$-value is
$\mathsf{c}(F_{c_0})/\|F_{c_0}\|_{\mathrm{ch}}^{|\mathcal{C}|}=1$. Consequently, if $R(Y)=R_{\max}$, then
\eqref{6.2:eq-chessboard-estimate-3} gives $R(Y^{\theta,+})R(Y^{\theta,-})\ge R_{\max}^2$ with both factors at most
$R_{\max}>0$, so $R(Y^{\theta,+})=R(Y^{\theta,-})=R_{\max}$: mirrored cell arrays of a maximiser are maximisers.

\emph{Spreading a maximiser.} Take an admissible $Y$ with $R(Y)=R_{\max}$, fix $c_0\in\mathcal{C}$ and put
$F:=Y_{c_0}\in\mathcal{S}\cap\mathfrak{A}_{c_0}$. Run the spreading sequence with $Y^0:=Y$, choosing at each step the
mirrored cell array that spreads the copies of $F$. By the preceding paragraph and induction on the step, each
cell array of the sequence is a maximiser. The terminal one is admissible, being an iterated mirrored
cell array of the admissible $Y$, and it is the chessboard array $(F^{[c']})_{c'}$ of $F$ based at $c_0$;
since $\|F\|_{\mathrm{ch}}>0$, its $R$-value is, by
Lemma~\ref{6.2:lem-chessboard-value},
\begin{equation*}
\frac{\mathsf{c}(F)}{\prod_{c'}\|F^{[c']}\|_{\mathrm{ch}}}
=\frac{\|F\|_{\mathrm{ch}}^{|\mathcal{C}|}}{\|F\|_{\mathrm{ch}}^{|\mathcal{C}|}}=1 .
\end{equation*}
Hence $R_{\max}=1$, and $R(\mathsf{X})\le1$ is the inequality
$|\mathsf{a}(\mathsf{X})|\le\prod_c\|F_c\|_{\mathrm{ch}}$.

For the last assertion take $F_{c_0}:=F$ and $F_c:=1\in\mathfrak{A}_c$ for $c\neq c_0$, and use
$\|1\|_{\mathrm{ch}}=1$ from Lemma~\ref{6.2:lem-chessboard-value}.

The proof uses only the following: reflection positivity (Lemma~\ref{6.1:lem-reflection-positivity} and
Corollary~\ref{6.1:cor-reflection-form}) at every pair of face planes $\Pi_{\mu,a}$, in all four directions $\mu$,
which Lemma~\ref{6.1:lem-reflection-positivity} provides for every direction and every link plane; the evenness of
the cell count $n_c$, through \eqref{6.2:eq-cell-reflections-a}, \eqref{6.2:eq-cell-reflections-b} and
\eqref{6.2:eq-cell-reflections-c}; and the fact that each $F_c$ reads only bonds of the open cell $c$. The lemma is
the chessboard estimate of Fr\"ohlich, Israel, Lieb and Simon \cite[Theorem~4.1]{FILS78} in the generality
needed here, namely complex observables, an antilinear reflection $\Theta_\theta$, gauge link variables and open
cells; it is proved above rather than cited.
\end{proof}

\begin{lemma}[Placing the supports inside cells by translation]\label{6.2:lem-support-placement}
Let the sources $f_1,\dots,f_n$, the balls $B_i=\bar B(x_i,R)\supset\operatorname{supp}f_i$ and the weight
$w_z=\sum_i z_if_i$ be as in Notation~\ref{6.1:not-units-sources}, and let the cells be as in
Definition~\ref{6.2:def-cell-reflections}. For $a\in\mathbb{R}^4$ put $\tau_af:=f(\cdot-a)$ and
$\tau_a\vec f:=(\tau_af_1,\dots,\tau_af_n)$. Let
\begin{equation*}
\ell\ \ge\ 8n(R+1).
\end{equation*}
Then there is $a\in\varepsilon\mathbb{Z}^4$ such that every translated ball $B_i+a$ lies in the
interior of a cell, at Euclidean distance at least $R+2$ from all face planes. Moreover
$S^T_n(\tau_a\vec f)=S^T_n(\vec f)$; indeed
\begin{equation*}
Z\Bigl(\sum_i z_i\tau_af_i\Bigr)=Z\Bigl(\sum_i z_if_i\Bigr)\qquad\text{for all } z\in\mathbb{C}^n.
\end{equation*}
\end{lemma}

\begin{proof}
Write $e_\mu$ for the unit vector in direction $\mu$. The translation $a$ is constructed one
direction at a time. Fix $\mu$ and $i$. For $t\in\mathbb{R}$, the $\mu$-th coordinates of the
points of $B_i+te_\mu$ fill the interval $[x_{i,\mu}+t-R,\,x_{i,\mu}+t+R]$, and the Euclidean
distance of a point $y$ to the union of the face planes $\{x_\mu\in\ell\mathbb{Z}\}$ is
$\operatorname{dist}(y_\mu,\ell\mathbb{Z})$ (this union is $N\mathbb{Z}$-periodic, since $\ell\mid N$).
Hence $B_i+te_\mu$ comes within distance $R+2$ of the face planes $\{x_\mu\in\ell\mathbb{Z}\}$
if and only if
\begin{equation*}
\operatorname{dist}(x_{i,\mu}+t,\ \ell\mathbb{Z})<2R+2 .
\end{equation*}
This condition depends on $t$ only modulo $\ell$. Since $n\ge 2$, we have
\begin{equation*}
\ell\ \ge\ 8n(R+1)\ \ge\ 16(R+1)\ >\ 4R+4,
\end{equation*}
so the condition says that $t$ lies in an open arc of length $4R+4$ of the
circle $\mathbb{R}/\ell\mathbb{Z}$.

The union over $i\le n$ of these $n$ open arcs has measure at most $n(4R+4)$, so its complement
has measure at least
\begin{equation*}
\ell-n(4R+4)\ \ge\ 8n(R+1)-4n(R+1)\ =\ 4n(R+1).
\end{equation*}
A union of $n$ open arcs of the circle has at most $n$ connected components, so its complement
consists of at most $n$ closed arcs. One of them therefore has length at least $4(R+1)\ge 4>\varepsilon$,
since $\varepsilon=L^{-K'}\le 1$. Because $\ell\in\varepsilon\mathbb{Z}$, the set
$\varepsilon\mathbb{Z}/\ell\mathbb{Z}$ is a well-defined subset of $\mathbb{R}/\ell\mathbb{Z}$ whose
consecutive points are at distance $\varepsilon$; hence a closed arc of length at least
$\varepsilon$ contains a point of it. Choose $t_\mu\in\varepsilon\mathbb{Z}$ representing such a
point, do this for every $\mu$, and put $a:=(t_\mu)_\mu\in\varepsilon\mathbb{Z}^4$.

By construction, for every $i$ and every $\mu$ we have
$\operatorname{dist}(x_{i,\mu}+t_\mu,\ell\mathbb{Z})\ge 2R+2$. Hence every point $y$ of $B_i+a$
satisfies $\operatorname{dist}(y_\mu,\ell\mathbb{Z})\ge R+2$ for every $\mu$. Thus, in each
direction $\mu$, the $\mu$-th coordinates of $B_i+a$ lie strictly between two consecutive multiples of
$\ell$. Consequently $B_i+a$ lies in the interior of a cell, at Euclidean distance at least $R+2$
from every face plane. Since $\operatorname{supp}\tau_af_i=\operatorname{supp}f_i+a\subset B_i+a$,
the translated sources are placed as claimed.

Invariance. The translation $x\mapsto x+a$ with $a\in\varepsilon\mathbb{Z}^4$ maps the bare lattice
onto itself and bare plaquettes bijectively onto bare plaquettes, $p'\mapsto p'+a$, with
$x_{p'+a}=x_{p'}+a$. It preserves the orientation of every bond. Let $U\circ\tau_a$ be the
transported configuration of Definition~\ref{6.2:def-cell-reflections} for this isometry, namely
$(U\circ\tau_a)(b)=U(b+a)$, so that $a_{p'+a}(U)=a_{p'}(U\circ\tau_a)$. Substituting $p=p'+a$
gives
\begin{align*}
A\Bigl(\sum_i z_i\tau_af_i,\,U\Bigr)
&=\sum_p w_z(x_p-a)\,a_p(U)
=\sum_{p'}w_z(x_{p'})\,a_{p'+a}(U)\\
&=\sum_{p'}w_z(x_{p'})\,a_{p'}(U\circ\tau_a)
=A(w_z,\,U\circ\tau_a).
\end{align*}

By Notation~\ref{6.1:not-units-sources}, $Z(v)=\int\exp[-A(x_0-v,U)-E]\,dU$. Since $A(\cdot,U)$ is
linear in its weight, $A(x_0,U)=x_0A(U)$, and $E$ does not depend on $U$, we have
\begin{equation*}
Z(v)=\int e^{-x_0A(U)-E}\,e^{A(v,U)}\,dU=Z(0)\,\bigl\langle e^{A(v,U)}\bigr\rangle,
\end{equation*}
because $\mu_{K,m}$ has density $e^{-x_0A(U)-E}/Z(0)$ with respect to $dU$. Therefore
\begin{equation*}
Z\Bigl(\sum_i z_i\tau_af_i\Bigr)=Z(0)\bigl\langle e^{A(w_z,U\circ\tau_a)}\bigr\rangle
=Z(0)\bigl\langle e^{A(w_z,U)}\bigr\rangle=Z(w_z).
\end{equation*}
The middle equality is the invariance of $\mu_{K,m}$ under every isometry of the lattice, the clause
contained in part (i) of the lemma on reflection positivity of the lattice measure
(Lemma~\ref{6.1:lem-reflection-positivity}), applied to the translation $x\mapsto x+a$ and to the
function $U\mapsto e^{A(w_z,U)}$, which is bounded and measurable, because it depends on finitely many bond
variables and $0\le a_p(U)\le 2$.

Finally, $S^T_n(\vec f)=\partial_{z_1}\cdots\partial_{z_n}\log Z(w_z)|_{z=0}$, and $S^T_n(\tau_a\vec f)$
is given by the same formula with $w_z$ replaced by $\sum_iz_i\tau_af_i$. The two functions of $z$
have just been shown to coincide. Hence $S^T_n(\tau_a\vec f)=S^T_n(\vec f)$.
\end{proof}

In the remainder of this chapter the sources are assumed so placed. That is, $f_i$ is replaced by
$\tau_af_i$ and $B_i$ by $B_i+a$; the cells themselves are fixed by the lattice, and only the
sources are moved. For a cell $c$ with interior $c^\circ$ put
\begin{equation*}
I_c:=\{\,i: B_i\subset c^\circ\,\},\qquad q:=\#\{\,c: I_c\neq\emptyset\,\}.
\end{equation*}
By Lemma~\ref{6.2:lem-support-placement}, each ball lies in the interior of a cell. The interiors
of distinct cells are disjoint. Hence each index $i$ belongs to exactly one $I_c$, so that
$\sum_c|I_c|=n$ and $q\le n$. For the same reason, two balls lying in different cells are disjoint.

\begin{lemma}[The periodised weight and its class]\label{6.2:lem-periodised-weight}
Let the sources be placed as in Lemma~\ref{6.2:lem-support-placement}, with $I_c$ as there, so that
every ball $B_i$ with $i\in I_c$ lies in the open cell $c^\circ$, at distance at least $R+2$ from
all face planes, and a cell $c$ is occupied when $I_c\neq\emptyset$. Let $\xi^{(\pm)}$ be as in
Definition~\ref{6.2:def-cell-reflections}. For an occupied cell $c$ and $z\in\mathbb{C}^n$
put
\begin{equation*}
w_{c,z}:=\sum_{i\in I_c}z_if_i,\qquad F_{c,z}:=e^{A(w_{c,z},\cdot)},
\end{equation*}
so that $w_z=\sum_c w_{c,z}$, the sum running over the occupied cells. Put, for each
$g\in\mathfrak{G}_\ell$ in the first line and for each $i\in I_c$ in the second line (where
$g$ is a summation index),
\begin{equation}\label{6.2:eq-periodised-weight-a}
\begin{aligned}
u_{c,g,z}&:=w_{c,z}\circ g^{-1}=\sum_{i\in I_c}z_i\,(f_i\circ g^{-1}),\\
w^{\mathrm{ch}}_{c,z}&:=\sum_{g\in\mathfrak{G}_\ell}u_{c,g,z}^{(\varepsilon(g))},\qquad
h^{\pm}_{c,i}:=\sum_{g\in\mathfrak{G}_\ell:\ \varepsilon(g)=\pm1}f_i\circ g^{-1}.
\end{aligned}
\end{equation}
With the polydisc of a finite system of real Lipschitz functions on $\mathbb{T}$ as in
Definition~\ref{6.2:def-localised-terms}, the following hold.
\begin{enumerate}
\item[(a)] $F_{c,z}\in\mathfrak{A}_c$, and $F_{c,z}$ is bounded at finite $K$. Put
$F_{c,z}:=1$ for unoccupied cells $c$. Then $\prod_{c\in\mathcal{C}}F_{c,z}=e^{A(w_z,\cdot)}$,
and so $Z(w_z)/Z(0)=\langle\prod_cF_{c,z}\rangle$.
\item[(b)] For every $g\in\mathfrak{G}_\ell$ the chessboard copy of $F_{c,z}$ based at $c$
satisfies $F_{c,z}^{[gc]}=e^{A(u_{c,g,z}^{(\varepsilon(g))},\cdot)}$. Hence
$\prod_{c'\in\mathcal{C}}F_{c,z}^{[c']}=e^{A(w^{\mathrm{ch}}_{c,z},\cdot)}$, the summands in
the definition of $w^{\mathrm{ch}}_{c,z}$ have pairwise disjoint supports, and
\begin{equation*}
w^{\mathrm{ch}}_{c,z}=\sum_{i\in I_c}\bigl(z_ih^+_{c,i}+\bar z_ih^-_{c,i}\bigr).
\end{equation*}
Each $h^\pm_{c,i}$ is a real Lipschitz function on $\mathbb{T}$ with
$\|h^\pm_{c,i}\|_{\sharp}=\|f_i\|_{\sharp}$. Finally
$\|F_{c,z}\|_{\mathrm{ch}}^{|\mathcal{C}|}=Z(w^{\mathrm{ch}}_{c,z})/Z(0)\in[0,\infty[$.
\item[(c)] For $z\in\mathbb{D}$, $w^{\mathrm{ch}}_{c,z}\in\mathcal{W}_{K'}$, the weight class
of the correlation theorem. For $z\in\tfrac12\mathbb{D}=\{z:\rho(z)<\tfrac12\}$, the pair
$(\zeta,\zeta')=((z_i)_{i\in I_c},(\bar z_i)_{i\in I_c})$ lies in the polydisc
\begin{equation*}
\Bigl\{(\zeta,\zeta'):\ \sum_{i\in I_c}|\zeta_i|\,\|h^+_{c,i}\|_{\sharp}
+\sum_{i\in I_c}|\zeta'_i|\,\|h^-_{c,i}\|_{\sharp}<r\Bigr\}
\end{equation*}
of the system $(h^+_{c,i},h^-_{c,i})_{i\in I_c}$ of $2|I_c|$ sources. Moreover, for every
$z\in\mathbb{D}$ and every $g$, $u_{c,g,z}$ is the weight of the transported system
$(f_i\circ g^{-1})_{i\in I_c}$ at the same parameter $z$, and $z$ lies in the polydisc of
that system.
\end{enumerate}
\end{lemma}

\begin{proof}
\emph{(a).} Let $p$ be a bare plaquette with $w_{c,z}(x_p)\neq0$. Since
$\operatorname{supp}f_i\subset B_i$, we have $x_p\in\bigcup_{i\in I_c}B_i$. Every point of
$p$ lies within distance $2\varepsilon$ of $x_p$. Hence $p$ is contained in the closed
$2\varepsilon$-neighbourhood of $\bigcup_{i\in I_c}B_i$. Each $B_i$ with $i\in I_c$ lies in
$c^\circ$ at distance at least $R+2$ from the face planes, and $R+2>2\varepsilon$ because
$\varepsilon\le1$. So this neighbourhood lies in $c^\circ$. Every bond of such a plaquette
therefore has both endpoints in $c^\circ$. Consequently
$A(w_{c,z},U)=\sum_pw_{c,z}(x_p)a_p(U)$ depends only on the variables $U(b)$ with
$b\subset c^\circ$, and it is continuous in them.

Since $\operatorname{tr}\mathbf{1}=1$, we have $|\operatorname{Re}\operatorname{tr}V|\le1$ for
$V\in SU(N_{\mathrm{col}})$, and hence $0\le a_p(U)\le2$. Let $Y$ be the closed
$2$-neighbourhood of $\bigcup_{i\in I_c}B_i$. It follows that
\begin{equation*}
|A(w_{c,z},U)|\le2\,\sup|w_{c,z}|\cdot\#\operatorname{plaq}(Y),
\end{equation*}
which is finite because at finite $K$ the lattice is finite. Thus
$|F_{c,z}|\le e^{|A(w_{c,z},\cdot)|}$ is bounded, and $F_{c,z}\in\mathfrak{A}_c$.

Every index $i$ belongs to exactly one $I_c$, since $B_i$ lies in exactly one open cell and
distinct open cells are disjoint. Hence $w_z=\sum_cw_{c,z}$. The weighted plaquette sum is
additive in the weight, $A(v+v',U)=A(v,U)+A(v',U)$, so
$\prod_cF_{c,z}=e^{A(w_z,\cdot)}$; unoccupied cells contribute $e^0=1$.

By additivity, $A(x_0-v,U)=x_0A(U)-A(v,U)$. Therefore, in the definition of $Z(v)$, the
integrand is $e^{A(v,U)}$ times the integrand of $Z(0)$, while $\mu_{K,m}$ is the probability
measure obtained by dividing the integrand of $Z(0)$ by $Z(0)$. Hence, for every complex
weight $v$,
\begin{equation}\label{6.2:eq-periodised-weight-1}
\frac{Z(v)}{Z(0)}=\bigl\langle e^{A(v,\cdot)}\bigr\rangle .
\end{equation}
With $v=w_z$ this is the last assertion of (a).

\emph{(b).} Let $v$ be a weight and $g\in\mathfrak{G}_\ell$. By definition
$(U\circ g)(b)=U(gb)^{s_g(b)}$. So the ordered product of $U\circ g$ around $\partial p$ is
the ordered product of $U$ around $g(\partial p)=\partial(gp)$, traversed in one of its two
orientations and from one of its corners. A change of corner conjugates this product and a
change of orientation inverts it. On $SU(N_{\mathrm{col}})$,
$\operatorname{Re}\operatorname{tr}$ is invariant under both, because
$\operatorname{tr}V^{-1}=\overline{\operatorname{tr}V}$. Hence $a_p(U\circ g)=a_{gp}(U)$.

By the isometry property of the cell-reflection group (Definition~\ref{6.2:def-cell-reflections}),
$g$ maps plaquettes bijectively onto plaquettes, and $x_{gp}=g\,x_p$. Substituting $p'=gp$
gives
\begin{equation}\label{6.2:eq-periodised-weight-2}
\begin{aligned}
T_g\bigl(e^{A(v,\cdot)}\bigr)(U)
&=\exp\sum_pv(x_p)\,a_p(U\circ g)=\exp\sum_pv(x_p)\,a_{gp}(U)\\
&=\exp\sum_{p'}v(g^{-1}x_{p'})\,a_{p'}(U)=e^{A(v\circ g^{-1},U)}.
\end{aligned}
\end{equation}
Since the $a_p$ are real, $\overline{e^{A(v,U)}}=e^{A(\bar v,U)}$.

Now take $v=w_{c,z}$. Then $F_{c,z}\in\mathfrak{A}_c$ by (a). By \eqref{6.2:eq-cell-reflections-a}
the group $\mathfrak{G}_\ell$ acts simply transitively on $\mathcal{C}$, so $g_{c\to gc}=g$.
By the definition of chessboard copies, $F_{c,z}^{[gc]}$ equals $T_gF_{c,z}$ if
$\varepsilon(g)=+1$ and $\overline{T_gF_{c,z}}$ if $\varepsilon(g)=-1$, with $\varepsilon$ the
parity homomorphism of \eqref{6.2:eq-cell-reflections-b}. Applying
\eqref{6.2:eq-periodised-weight-2} and the conjugation rule gives
$F_{c,z}^{[gc]}=e^{A(u_{c,g,z}^{(\varepsilon(g))},\cdot)}$.

The map $g\mapsto gc$ is a bijection $\mathfrak{G}_\ell\to\mathcal{C}$, again by simple
transitivity. Additivity of $A$ in the weight then gives
$\prod_{c'}F^{[c']}_{c,z}=e^{A(w^{\mathrm{ch}}_{c,z},\cdot)}$.

For the supports: $\operatorname{supp}u_{c,g,z}\subset g(\bigcup_{i\in I_c}B_i)$, and this set
lies in $g(c^\circ)=(gc)^\circ$ because $g$ maps open cells onto open cells
(Definition~\ref{6.2:def-cell-reflections}). Distinct $g$ give distinct cells, and the open
cells of distinct cells are disjoint. So the summands have pairwise disjoint supports.

The $f_i$ are real, so
\begin{equation*}
\bigl(z_i\,f_i\circ g^{-1}\bigr)^{(\varepsilon(g))}=z_i^{(\varepsilon(g))}\,f_i\circ g^{-1}.
\end{equation*}
Grouping the $g$ by parity yields
$w^{\mathrm{ch}}_{c,z}=\sum_{i\in I_c}(z_ih^+_{c,i}+\bar z_ih^-_{c,i})$.

The functions $h^\pm_{c,i}$ are real. Their summands $f_i\circ g^{-1}$ are supported in the
pairwise disjoint sets $g(B_i)\subset(gc)^\circ$. Each parity class is non-empty: the identity
has parity $+1$ and a single face reflection has parity $-1$. Each summand has supremum norm
$\|f_i\|_\infty$, since $g$ is a bijective isometry of $\mathbb{T}$. Hence
$\|h^\pm_{c,i}\|_\infty=\|f_i\|_\infty$. The proof of (c) below shows that
$\operatorname{Lip}h^\pm_{c,i}=\operatorname{Lip}f_i$. Together these give
$\|h^\pm_{c,i}\|_{\sharp}=\|f_i\|_{\sharp}$.

Finally, the chessboard value of $F_{c,z}$ is
\begin{equation*}
\mathsf{c}(F_{c,z})=\Bigl\langle\prod_{c'}F^{[c']}_{c,z}\Bigr\rangle
=\bigl\langle e^{A(w^{\mathrm{ch}}_{c,z},\cdot)}\bigr\rangle
=\frac{Z(w^{\mathrm{ch}}_{c,z})}{Z(0)},
\end{equation*}
by \eqref{6.2:eq-periodised-weight-1}. It is finite because the integrand is bounded. It is
real and non-negative by positivity of the chessboard value
(Lemma~\ref{6.2:lem-chessboard-value}). By definition
$\|F_{c,z}\|_{\mathrm{ch}}^{|\mathcal{C}|}=\mathsf{c}(F_{c,z})$.

\emph{(c).} Let $z\in\mathbb{D}$ and put $\Lambda:=\sum_{i\in I_c}|z_i|\operatorname{Lip}f_i$.
Since $\rho(z)<1$, and since $\|f\|_{\sharp}\ge\|f\|_\infty$ and
$\|f\|_{\sharp}\ge\operatorname{Lip}f/\lambda_*$, the definition of the gauge $\rho$ gives
\begin{equation}\label{6.2:eq-periodised-weight-3}
\sum_{i\in I_c}|z_i|\,\|f_i\|_\infty\le r\rho(z)<r,\qquad
\Lambda\le\lambda_*r\rho(z)<\lambda_*r .
\end{equation}

\emph{Supremum.} By (b), $w^{\mathrm{ch}}_{c,z}$ is a sum of disjointly supported copies
$u_{c,g,z}^{(\varepsilon(g))}$, each with the supremum of $|w_{c,z}|$. Hence, using
\eqref{6.2:eq-periodised-weight-3},
\begin{equation*}
\sup|w^{\mathrm{ch}}_{c,z}|=\sup|w_{c,z}|\le\sum_{i\in I_c}|z_i|\,\|f_i\|_\infty<r\le r_{K'}.
\end{equation*}

\emph{Lipschitz bound in one cell.} Fix $g$. Every summand $u_{c,g',z}$ with $g'\neq g$ is
supported in $(g'c)^\circ$. This set does not meet the closed cell $gc$, because distinct
cells meet only in points of their faces. So $w^{\mathrm{ch}}_{c,z}=u^{(\varepsilon(g))}_{c,g,z}$
on $gc$. For all $x',y'$,
\begin{equation*}
|w_{c,z}(x')-w_{c,z}(y')|\le\sum_{i\in I_c}|z_i|\,|f_i(x')-f_i(y')|\le\Lambda|x'-y'| .
\end{equation*}
Also $g^{-1}$ is an isometry, and conjugation preserves moduli of differences. Hence
$|w^{\mathrm{ch}}_{c,z}(x)-w^{\mathrm{ch}}_{c,z}(y)|\le\Lambda|x-y|$ whenever $x,y$ lie in a
common cell.

\emph{Vanishing on faces.} Each $B_i$ lies at distance at least $R+2$ from the face planes.
Each $g$ is an isometry mapping cells onto cells, and hence the union of the face planes onto
itself. So every copy $g(B_i)$ lies at distance at least $R+2$ from all faces, and
$w^{\mathrm{ch}}_{c,z}$ vanishes on the faces of every cell.

\emph{Gluing.} Let $x\in gc$ and $y\in g'c$ lie in no common cell. Let
$\gamma\colon[0,1]\to\mathbb{T}$ be a shortest segment from $x$ to $y$, so that
$|\gamma(t)-\gamma(s)|=|t-s|\,|x-y|$. Put
\begin{equation*}
t_1:=\sup\{t:\gamma([0,t])\subset gc\},\qquad t_2:=\inf\{t:\gamma([t,1])\subset g'c\},
\end{equation*}
and $x^*:=\gamma(t_1)$, $y^*:=\gamma(t_2)$. Since $y\notin gc$ and $x\notin g'c$, the point
$x^*$ lies on the boundary of $gc$ and $y^*$ on the boundary of $g'c$. Both lie on faces, so
$w^{\mathrm{ch}}_{c,z}(x^*)=w^{\mathrm{ch}}_{c,z}(y^*)=0$.

If $t_1\le t_2$, the points $x,x^*,y^*,y$ lie on $\gamma$ in this order, and the bound in one
cell gives
\begin{equation*}
|w^{\mathrm{ch}}_{c,z}(x)-w^{\mathrm{ch}}_{c,z}(y)|
\le\Lambda\bigl(|x-x^*|+|y^*-y|\bigr)\le\Lambda|x-y| .
\end{equation*}
If $t_1>t_2$, then $x^*$ lies in both $gc$ and $g'c$ and on $\gamma$. The same bound gives
$\Lambda(|x-x^*|+|x^*-y|)=\Lambda|x-y|$.

Every point of $\mathbb{T}$ lies in some cell. Hence, with \eqref{6.2:eq-periodised-weight-3},
\begin{equation*}
\operatorname{Lip}_{K'}w^{\mathrm{ch}}_{c,z}\le\Lambda<r\lambda_*=\mathfrak{l}_{K'} .
\end{equation*}
Together with the supremum bound, this shows $w^{\mathrm{ch}}_{c,z}\in\mathcal{W}_{K'}$.

\emph{The functions $h^\pm_{c,i}$.} The same three steps apply to $h^\pm_{c,i}$. On the cell
$gc$, $h^\pm_{c,i}$ equals $f_i\circ g^{-1}$ if $\varepsilon(g)=\pm1$ and vanishes otherwise.
It is Lipschitz on each cell with constant at most $\operatorname{Lip}(f_i|_{c})$, and it
vanishes on all faces. Gluing gives
$\operatorname{Lip}h^\pm_{c,i}\le\operatorname{Lip}(f_i|_c)\le\operatorname{Lip}f_i$.
Conversely, on a cell $gc$ with $\varepsilon(g)=\pm1$ the function $h^\pm_{c,i}$ is an
isometric copy of $f_i|_c$, so $\operatorname{Lip}h^\pm_{c,i}\ge\operatorname{Lip}(f_i|_c)$.
The same gluing argument applied to $f_i$, which equals $f_i|_c$ on $c$ and vanishes on every
other cell and on all faces, gives $\operatorname{Lip}f_i\le\operatorname{Lip}(f_i|_c)$.
Hence $\operatorname{Lip}h^\pm_{c,i}=\operatorname{Lip}f_i$, which completes (b).

\emph{Polydiscs.} Let $z\in\tfrac12\mathbb{D}$. By (b),
$\|h^\pm_{c,i}\|_{\sharp}=\|f_i\|_{\sharp}$, so
\begin{equation*}
\sum_{i\in I_c}|z_i|\,\|h^+_{c,i}\|_{\sharp}+\sum_{i\in I_c}|\bar z_i|\,\|h^-_{c,i}\|_{\sharp}
=2\sum_{i\in I_c}|z_i|\,\|f_i\|_{\sharp}\le2r\rho(z)<r .
\end{equation*}
By (b), $w^{\mathrm{ch}}_{c,z}$ is the weight of this system at $(\zeta,\zeta')$. Now let
$z\in\mathbb{D}$ and $g\in\mathfrak{G}_\ell$. Each $f_i\circ g^{-1}$ is real Lipschitz with
$\|f_i\circ g^{-1}\|_{\sharp}=\|f_i\|_{\sharp}$, because $g$ is an isometry of $\mathbb{T}$. So
$u_{c,g,z}=\sum_{i\in I_c}z_i(f_i\circ g^{-1})$ is the weight of the transported system at
$z$, and
\begin{equation*}
\sum_{i\in I_c}|z_i|\,\|f_i\circ g^{-1}\|_{\sharp}\le r\rho(z)<r .
\end{equation*}
\end{proof}

\begin{remark}[The bounds at transported and periodised weights]
\label{6.2:rem-transported-weights}
\begin{enumerate}
\item[(i)] The analyticity \eqref{6.2:eq-localised-terms-b}, the bound
\eqref{6.2:eq-localised-terms-c} and the stability bound \eqref{6.2:eq-localised-terms-d} of
Definition~\ref{6.2:def-localised-terms} hold for every weight $v\in\mathcal{W}_{K'}$, and not
only at the weights $v=w_z$, $z\in\mathbb{D}$.
\item[(ii)] Suppose the correlation theorem is used only for weights that are complex linear
combinations of a finite system $(h_l)$ of real Lipschitz functions on $\mathbb{T}$, with
coefficients in that system's polydisc as in Definition~\ref{6.2:def-localised-terms}. Then for
every occupied cell $c$ and every
$g\in\mathfrak{G}_\ell$ the following hold, with $u_{c,g,z}$ and $w^{\mathrm{ch}}_{c,z}$ as in
\eqref{6.2:eq-periodised-weight-a}:
\begin{itemize}
\item the properties \eqref{6.2:eq-localised-terms-b} and \eqref{6.2:eq-localised-terms-c} hold at
$v=u_{c,g,z}$ for every $z\in\mathbb{D}$;
\item the properties \eqref{6.2:eq-localised-terms-c} and \eqref{6.2:eq-localised-terms-d} hold at
$v=w^{\mathrm{ch}}_{c,z}$ for every $z\in\mathbb{D}_2=\{\rho(z)<\tfrac12\}$.
\end{itemize}
At these weights $\mathcal{E}(v)$ and $\widetilde{Z}(v)$ are given by the formulas that accompany
\eqref{6.2:eq-localised-terms-d}.
Here $e_j(X;v)=e^{(j)}(X;v)-e^{(j)}(X;0)$, and each localised vacuum-energy term
$e^{(j)}(X;\cdot)$ is defined on $\mathcal{W}_j(X)$ by its production formula. The only loss is
that the bounds are asserted on $\mathbb{D}_2$ instead
of $\mathbb{D}$. Accordingly the radii in the estimates drawn from these bounds halve,
and the factor $(n/r)^n$ in the resulting bounds becomes $(2n/r)^n$.
\end{enumerate}
Part~(i) is drawn from the proof of the correlation theorem, part~(ii) from its statement alone;
both are proved below. In what follows all bounds are used on the polydisc
$\mathbb{D}_{\mathfrak{r}}$ of \eqref{6.2:eq-localised-terms-f}, with $\mathfrak{r}=1$ when the
bounds are taken from (i) and $\mathfrak{r}=2$ when they are taken from (ii) alone; the later
estimates are written with $\mathfrak{r}$ and hold for either value.
\end{remark}

\begin{proof}
(i) We begin with the stability bound $|\widetilde{Z}(v)|\le e^{E_+^{\sharp}N^4}$. In the torus
bound of the correlation theorem it is derived from the following five ingredients.

The first is the representation of $\widetilde{Z}$ as the $V_{K'}$-integral of the large-field
expansion of the sourced effective density at level $K'$, with the counterterm $\mathcal{E}$
subtracted in the exponent. This representation rests on the inductive format that the proof of
the correlation theorem establishes for the sourced effective density at level $K'$. The inductive
statement of that proof supplies this format for every weight in $\mathcal{W}_{K'}$.

The second is the family of bounds on the individual terms of the effective action. These are
stated for all sites and for every weight in $\mathcal{W}_k$, or for a constant running shift.

The third is the bound for complex weights of the correlation theorem's proof. It applies to every
complex weight whose modulus is at most $\vartheta$ times the real reference value fixed there.
In the proof it is applied to the running shift $\zeta'_{K'}$ at level $K'$, for which it suffices
that $\sup|\zeta'_{K'}|\le\tfrac83 r$. For $v\in\mathcal{W}_{K'}$ one has $\sup|v|<r_{K'}$. By the
running-shift statement,
\begin{equation*}
\sup|\zeta'_{K'}|\le\tfrac43\sup|v|<\tfrac43 r_{K'}<\tfrac83 r .
\end{equation*}
Hence this ingredient uses only $\sup|v|<2r$.

The fourth is the family of bounds on the moduli of the terms. These are stated for every weight
in $\mathcal{W}_k$, or constant.

The fifth is Ba{\l}aban's proof of the upper bound in \cite[(2.50)]{B-CMP119}, which proceeds
from \cite[(2.49)]{B-CMP119}. That proof
estimates the integral $\int dV_k\rho_k$ by a sum of terms and bounds each term by its modulus;
here it is integrated over $V_{K'}$.

None of these five ingredients uses the presentation $v=\sum_iz_if_i$. The torus bound states the
stability bound on $\mathbb{D}$, and this is its specialisation along $z\mapsto w_z$, which is
legitimate because $w_z\in\mathcal{W}_{K'}$ for $z\in\mathbb{D}$. Hence
\eqref{6.2:eq-localised-terms-d} holds for every $v\in\mathcal{W}_{K'}$.

For $j\le K'$ and a localisation domain $X$, write $e^{(j)}(X;v)$ for the localised vacuum-energy
term of the correlation theorem's construction. Then $e_j(X;v)=e^{(j)}(X;v)-e^{(j)}(X;0)$. The
analyticity and the bound of $e_j(X;v)$ in the torus bound come from clauses~(c), (d) and~(f) of
the inductive format of the correlation theorem's proof: clause~(c) bounds the local
vacuum-energy terms, and clause~(d) bounds the large-field operations $\mathbf{R}$, $\mathbf{B}$
of the sourced step. These bounds hold uniformly on $\mathcal{W}_j(X)$, and
the terms are analytic
there and read the weight on $X^+$ only. Hence \eqref{6.2:eq-localised-terms-b} and
\eqref{6.2:eq-localised-terms-c} hold for every $v\in\mathcal{W}_{K'}$. The statement of the
torus bound that $e_j(X;\cdot)$ is independent of $z_i$ unless
$X^+\cap\operatorname{supp}f_i\ne\emptyset$ is the locality property of
Definition~\ref{6.2:def-localised-terms}, specialised to $v=w_z$.

(ii) The correlation theorem is stated for every value of the reference coupling not exceeding
$\gamma_J$, every $x_0$, every $K$ and $m$
and every torus, and for every choice of sources $f_i$. It places no restriction on the number of
sources, their supports or their sizes. Moreover $K$, $k$, $m$, the torus and the position enter
none of its constants, the number of sites $N$ enters only through $\|\cdot\|_{\sharp}$ and
$\mathsf{Q}$, and the sources enter only through the norms $\|f_i\|_{\sharp}$.
It may therefore be applied to any finite system of real Lipschitz test functions, one system at
a time, with the same constants.

\emph{The transported weight.} Fix an occupied cell $c$ and $g\in\mathfrak{G}_\ell$. By
clause~(c) of Lemma~\ref{6.2:lem-periodised-weight},
$u_{c,g,z}=\sum_{i\in I_c}z_i(f_i\circ g^{-1})$ is the weight of the transported system
$(f_i\circ g^{-1})_{i\in I_c}$ at the same parameter $z$. For $z\in\mathbb{D}$ that parameter
lies inside the system's polydisc, because $\|f_i\circ g^{-1}\|_{\sharp}=\|f_i\|_{\sharp}$.
Applying the correlation theorem to this system gives \eqref{6.2:eq-localised-terms-b} and
\eqref{6.2:eq-localised-terms-c} at $u_{c,g,z}$ for $z\in\mathbb{D}$.

\emph{The periodised weight.} Also by clause~(c) of Lemma~\ref{6.2:lem-periodised-weight},
$w^{\mathrm{ch}}_{c,z}$ is the weight of the system $(h^+_{c,i},h^-_{c,i})_{i\in I_c}$, which
has $2|I_c|$ sources. It is taken at the parameter point
$((z_i)_{i\in I_c},(\bar z_i)_{i\in I_c})$. For $z\in\mathbb{D}_2$ this point lies
in the polydisc of that system, that is,
\begin{equation*}
\sum_{i\in I_c}|z_i|\,\|h^+_{c,i}\|_{\sharp}
+\sum_{i\in I_c}|\bar z_i|\,\|h^-_{c,i}\|_{\sharp}<r .
\end{equation*}
Applying the correlation theorem to this system gives
\eqref{6.2:eq-localised-terms-c} and \eqref{6.2:eq-localised-terms-d} at $w^{\mathrm{ch}}_{c,z}$
for $z\in\mathbb{D}_2$.

It remains to check that the quantities in \eqref{6.2:eq-localised-terms-d} are the ones
attached to the weight itself, namely those given by the formulas that accompany
\eqref{6.2:eq-localised-terms-d}, for which
$Z=e^{\mathcal{E}}\widetilde{Z}$. Each $e^{(j)}(X;\cdot)$ is defined on
$\mathcal{W}_j(X)$ by its production formula. That formula belongs to the construction written
directly in terms of the weight, without reference to the parameters $z$ or to the system of
test functions presenting it. Hence the values obtained from the transported or
periodised system coincide with those the definition assigns to $v$.

\emph{The loss.} The bounds at the periodised weight are asserted on $\mathbb{D}_2$ rather than
on $\mathbb{D}$. Consequently the radii in the estimates drawn from these bounds halve,
and $(n/r)^n$ is replaced by $(2n/r)^n$. This is recorded by carrying $\mathfrak{r}\in\{1,2\}$
in \eqref{6.2:eq-localised-terms-f}.
\end{proof}

\begin{lemma}[Diameter of an enlarged localisation domain]\label{6.2:lem-domain-diameter}
We work in the last-lattice units of Notation~\ref{6.1:not-units-sources}. A \emph{face plane} is a
plane $\{x_\mu=a\ell\}$, $a\in\mathbb{Z}$, on which faces of the cells of
Definition~\ref{6.2:def-cell-reflections} lie. Let $X\in D_j$, write $s_j:=ML^{j-K'}$ for the side
length, in last-lattice units, of the cubes of the scale-$j$ partition $\pi_j$, and let
$x_1,x_2\in X^+$ be points of the torus $\mathbb{T}$. Then
\begin{equation}\label{6.2:eq-domain-diameter-1}
|x_1-x_2|_{\mathbb{T}}\le s_j\bigl(d_j(X)+8\bigr)\le 4s_j\bigl(d_j(X)+3\bigr).
\end{equation}
Consequently the following hold.
\begin{enumerate}
\item[(a)] Let $v_1,v_2\colon\mathbb{T}\to\mathbb{C}$ and $D>0$. Suppose that for $i=1,2$ there is
$p_i\in\operatorname{plaq}(X^+)$ with $v_i(x_{p_i})\neq 0$. Suppose also that the set of
barycentres at which $v_1$ does not vanish and the corresponding set for $v_2$ are at distance at
least $D$. Then
\begin{equation}\label{6.2:eq-domain-diameter-a}
d_j(X)\ge\frac{D-4\varepsilon}{4s_j}-3\ge\frac{D\,L^{K'-j}}{4M}-4.
\end{equation}
\item[(b)] Let $R>0$, and let $B\subset\mathbb{T}$ be a set at distance at least $R+2$ from every
face plane; in particular, $B$ may be one of the ball copies of a cell, placed at distance at least
$R+2$ from the faces of that cell. Let $v\colon\mathbb{T}\to\mathbb{C}$ vanish outside $B$. Suppose
that $X^+$ contains a point of a face plane and that $v(x_p)\neq0$ for some
$p\in\operatorname{plaq}(X^+)$. Then $X^+$
contains two points at distance at least $R+2-2\varepsilon\ge R$, and
\begin{equation}\label{6.2:eq-domain-diameter-b}
d_j(X)\ge\frac{R\,L^{K'-j}}{4M}-3.
\end{equation}
\item[(c)] For a cell $c'\in\mathcal{C}$ put
\begin{equation}\label{6.2:eq-domain-diameter-c}
\mathcal{X}_{c'}:=\{X\in D_j:\ X^+\subset(c')^\circ\},\qquad
\mathcal{X}_\partial:=D_j\setminus\bigsqcup_{c'\in\mathcal{C}}\mathcal{X}_{c'}.
\end{equation}
The union is disjoint. If $X\in\mathcal{X}_\partial$, then $X^+$ contains a point of a face plane.
\item[(d)] Let $R>0$. For each $c'\in\mathcal{C}$ let $B_{c'}\subset c'$ be a set at distance at
least $R+2$ from every face plane; in particular, $B_{c'}$ may be a ball copy sitting in
$c'$ at distance at least $R+2$ from the faces of $c'$. Let $v=\sum_{c'\in\mathcal{C}}v_{c'}$, where
$v_{c'}$ vanishes outside $B_{c'}$. If $g\in\mathfrak{G}_\ell$, $c\in\mathcal{C}$ and
$X\in\mathcal{X}_{gc}$, then
\begin{equation}\label{6.2:eq-domain-diameter-d}
e_j(X;v)=e_j(X;v_{gc}).
\end{equation}
\end{enumerate}
\end{lemma}

\begin{proof}
\emph{The bound \eqref{6.2:eq-domain-diameter-1}.}
By definition, $X^+$ is $X$ together with one layer of $\pi_j$-cubes. Hence every $\pi_j$-cube of
$X^+$ either is a cube of $X$ or shares a point with a cube of $X$. Consequently, for $i=1,2$ there
are $\pi_j$-cubes $\square_i'\subset X^+$ and $\square_i\subset X$ such that $x_i\in\square_i'$ and
$\square_i'\cap\square_i\neq\emptyset$. (If $\square_i'\subset X$, take $\square_i=\square_i'$.)

By \cite[(0.24)--(0.26)]{B-CMP109}, $d_j(X)$ is the length of a shortest tree graph in $X$ meeting
all cubes of $X$, divided by the side of the cubes of $\pi_j$. In last-lattice units that side is
$s_j$, so a shortest such tree $T\subset X$ has length $s_j d_j(X)$. The tree $T$ meets
$\square_1$ and $\square_2$; choose $t_i\in T\cap\square_i$. Since $T$ is connected, $t_1$ and $t_2$
are joined by a path inside $T$. This path has length at most the total length $s_jd_j(X)$ of $T$,
so $|t_1-t_2|_{\mathbb{T}}\le s_jd_j(X)$.

Two cubes of side $s_j$ that share a point lie in a cube of side $2s_j$. In dimension $d=4$ the
Euclidean diameter of that cube is $2s_j\sqrt4=4s_j$. Since $x_i,t_i\in\square_i'\cup\square_i$,
\begin{equation*}
|x_i-t_i|_{\mathbb{T}}\le\operatorname{diam}(\square_i'\cup\square_i)\le 4s_j .
\end{equation*}
The triangle inequality gives
\begin{equation*}
|x_1-x_2|_{\mathbb{T}}\le|x_1-t_1|_{\mathbb{T}}+|t_1-t_2|_{\mathbb{T}}+|t_2-x_2|_{\mathbb{T}}
\le 4s_j+s_jd_j(X)+4s_j ,
\end{equation*}
which is the first inequality of \eqref{6.2:eq-domain-diameter-1}. The second inequality is
equivalent to $d_j(X)+8\le4d_j(X)+12$, that is, to $3d_j(X)+4\ge0$, which holds because
$d_j(X)\ge0$. The weaker form $4s_j(d_j(X)+3)$ is the one in which the bound enters the correlation
theorem.

\emph{Part (a).} Recall that $\operatorname{plaq}(X^+)$ is the set of bare plaquettes intersecting
$X^+$. For $i=1,2$ choose $y_i\in p_i\cap X^+$. The barycentre $x_{p_i}$ lies in $p_i$, and
$\operatorname{diam}p_i\le2\varepsilon$, so $|x_{p_i}-y_i|_{\mathbb{T}}\le2\varepsilon$. Since
$v_i(x_{p_i})\neq0$, the hypothesis on the two barycentre sets, followed by
\eqref{6.2:eq-domain-diameter-1} applied to $y_1,y_2\in X^+$, gives
\begin{equation*}
D\le|x_{p_1}-x_{p_2}|_{\mathbb{T}}\le|y_1-y_2|_{\mathbb{T}}+4\varepsilon
\le4s_j\bigl(d_j(X)+3\bigr)+4\varepsilon .
\end{equation*}
Rearranging yields the first inequality of \eqref{6.2:eq-domain-diameter-a}.

For the second inequality, note that $s_j=ML^{j-K'}$, so $D/(4s_j)=DL^{K'-j}/(4M)$ and
$4\varepsilon/(4s_j)=\varepsilon/s_j$. Since $j\ge0$,
\begin{equation*}
\varepsilon=L^{-K'}\le L^{j-K'}=\frac{s_j}{M},
\end{equation*}
hence $\varepsilon/s_j\le1/M\le1$. Therefore
\begin{equation*}
\frac{D-4\varepsilon}{4s_j}-3\ge\frac{DL^{K'-j}}{4M}-1-3 ,
\end{equation*}
which is the second inequality.

\emph{Part (b).} Let $x_1\in X^+$ lie on a face plane. Let $p\in\operatorname{plaq}(X^+)$ satisfy
$v(x_p)\neq0$; then $x_p\in B$, because $v$ vanishes outside $B$. Choose $x_2\in p\cap X^+$. The
point $x_p$ is at distance at least $R+2$ from every face plane, so
$|x_p-x_1|_{\mathbb{T}}\ge R+2$. Moreover $|x_2-x_p|_{\mathbb{T}}\le\operatorname{diam}p\le2\varepsilon$.
Hence
\begin{equation*}
|x_1-x_2|_{\mathbb{T}}\ge R+2-2\varepsilon\ge R,
\end{equation*}
where the last step uses $\varepsilon=L^{-K'}\le1$. Combining this with
\eqref{6.2:eq-domain-diameter-1} gives $R\le4s_j(d_j(X)+3)$. Therefore
\begin{equation*}
d_j(X)\ge\frac{R}{4s_j}-3=\frac{RL^{K'-j}}{4M}-3,
\end{equation*}
which is \eqref{6.2:eq-domain-diameter-b}.

\emph{Part (c).} The open cells are pairwise disjoint and $X^+\neq\emptyset$, so no $X$ belongs to
two of the families $\mathcal{X}_{c'}$; the union is therefore disjoint.

The cells are the closed cubes $\ell k+[0,\ell]^4$. The boundary of each cell lies in the union of
the face planes. The complement of that union is the disjoint union of the open cells, and each open
cell is a connected component of the complement.

Now $X\in D_j$ is a wall-connected union of closed $\pi_j$-cubes. Each cube added to form $X^+$
shares a point with $X$. Hence $X^+$ is a connected union of closed cubes.

Let $X\in\mathcal{X}_\partial$. Suppose first that $X^+\subset c'$ for some cell $c'$. Since $X^+$
is not contained in $(c')^\circ$, it contains a point of $\partial c'$, which lies on a face plane.

Suppose next that $X^+$ is contained in no cell. Then either $X^+$ contains a point of a face plane,
or $X^+$ lies in the complement of the face planes. In the latter case $X^+$ meets two different
open cells, since otherwise it would lie in one closed cell. This is impossible for a connected set
in the complement, because the open cells are the connected components of the complement. Hence
$X^+$ contains a point of a face plane.

The argument uses only that $X^+$ is a connected union of closed cubes; it uses no alignment of the
cubes of $\pi_j$ with the cells.

\emph{Part (d).} By Definition~\ref{6.2:def-cell-reflections}, $gc$ is a cell. Let
$X\in\mathcal{X}_{gc}$ and $p\in\operatorname{plaq}(X^+)$. Then $p$ meets $X^+\subset(gc)^\circ$, and
$\operatorname{diam}p\le2\varepsilon$. Hence every point $y$ of $p$, in particular $x_p$, lies within
$2\varepsilon$ of some point $x'\in(gc)^\circ$.

Let $c'\neq gc$ and $x''\in B_{c'}$. The open cell $(gc)^\circ$ is disjoint from the closed cell
$c'$, so $x'\notin c'$. A shortest segment from $x'$ to $x''$ therefore meets $\partial c'$ at a
point $y'\neq x'$. Since $y'$ lies on a shortest segment from $x'$ to $x''$, and $y'$ lies on a face
plane,
\begin{equation*}
|x'-x''|_{\mathbb{T}}=|x'-y'|_{\mathbb{T}}+|y'-x''|_{\mathbb{T}}>|y'-x''|_{\mathbb{T}}\ge R+2 .
\end{equation*}
It follows that
\begin{equation*}
|y-x''|_{\mathbb{T}}>R+2-2\varepsilon\ge R>0 .
\end{equation*}
So $p$ lies at distance greater than $R+2-2\varepsilon\ge R$ from $B_{c'}$ for every cell
$c'\neq gc$. In particular $x_p\notin B_{c'}$, and so $v_{c'}(x_p)=0$ for $c'\neq gc$.

Thus $v$ and $v_{gc}$ agree at the barycentres of all plaquettes of $\operatorname{plaq}(X^+)$. The
locality property \eqref{6.2:eq-localised-terms-a} of the localised terms
(Definition~\ref{6.2:def-localised-terms}) gives \eqref{6.2:eq-domain-diameter-d}.
\end{proof}

\begin{lemma}[Anchor count and tree-graph sum with threshold]\label{6.2:lem-tree-sum}
Fix a scale $j$, a point $x\in\mathbb{T}$ and $R\ge 0$.
Here $|\cdot|_{\mathbb{T}}$ is the Euclidean distance of the torus, $\bar B(x,r)$ is the closed ball of
radius $r$ about $x$ for this distance, and two $\pi_j$-cubes are \emph{adjacent} if they share a point,
so that every cube is adjacent to itself.
\begin{enumerate}
\item[(a)] Let $X\in D_j$. The number of $\pi_j$-cubes of $X$ that are adjacent to some $\pi_j$-cube meeting
$\bar B(x,R+2)$ is at most
\begin{equation}\label{6.2:eq-tree-sum-a}
\mathfrak{a}(R,K'-j):=\Bigl(\frac{2(R+2)L^{K'-j}}{M}+6\Bigr)^4 .
\end{equation}
Let $p\in\operatorname{plaq}(X^+)$ satisfy $x_p\in\bar B(x,R)$. Then $p$ meets a $\pi_j$-cube of $X^+$ that
meets $\bar B(x,R+2)$, and every cube of $X$ adjacent to that cube is among the cubes counted above.
\item[(b)] Let $C_e,\delta,\kappa$ be the constants of the bound \eqref{6.2:eq-localised-terms-c}.
Let $c_\kappa>0$ be such that the rate
\[
\kappa_{\mathrm{adm}}:=(1-\delta)\kappa-c_\kappa
\]
is admissible for \cite[(1.26)]{B-CMP116}. This means there is a constant $C_{(1.26)}$, independent of $j$,
of the cube $\square'$ and of the volume, with
\[
\sum_{X\in D_j,\ X\supset\square'} e^{-\kappa_{\mathrm{adm}} d_j(X)}\le C_{(1.26)}
\qquad\text{for every $\pi_j$-cube $\square'$.}
\]
By \cite[(1.26)]{B-CMP116}, every sufficiently large rate is admissible in this sense.
For the constants $\delta,\kappa$ of \eqref{6.2:eq-localised-terms-c}, which are those of the correlation
theorem, such a $c_\kappa$ exists: the proof of that theorem already writes $(1-\delta)\kappa$ as an
admissible rate plus a retained positive part.
Let $D\ge 0$, and let $\mathcal{F}$ be a set of $\pi_j$-cubes with $\#\mathcal{F}\le\mathfrak{a}(R,K'-j)$.
Let $\mathcal{Y}\subset D_j$ be a family such that every $X\in\mathcal{Y}$ contains one of the cubes
of $\mathcal{F}$ and has $d_j(X)\ge D$. Then
\begin{equation}\label{6.2:eq-tree-sum-b}
\sum_{X\in\mathcal{Y}}C_e\,e^{-(1-\delta)\kappa\, d_j(X)}
\le C_e\,C_{(1.26)}\,\mathfrak{a}(R,K'-j)\,e^{-c_\kappa D}.
\end{equation}
\end{enumerate}
\end{lemma}

\begin{proof}
\emph{(a), the count.} The cubes of $\pi_j$ have side $s_j=ML^{j-K'}$, so $1/s_j=L^{K'-j}/M$.

Fix a lift $\tilde x\in\mathbb{R}^4$ of $x$, and regard $\pi_j$ as the image of the grid of cubes of side
$s_j$ in $\mathbb{R}^4$. Let $\square$ be a cube of the kind counted, adjacent to a cube $\square''$ that
meets $\bar B(x,R+2)$ at a point $y$. Choose a lift $\tilde y$ of $y$ with $|\tilde y-\tilde x|\le R+2$.
Choose a grid cube $\tilde\square''$ containing $\tilde y$ and projecting onto $\square''$. Finally choose a
grid cube $\tilde\square$ that shares a point with $\tilde\square''$ and projects onto $\square$. Distinct
cubes $\square$ receive distinct lifts $\tilde\square$, so it suffices to count the grid cubes
$\tilde\square$ so obtained.

The Euclidean ball of radius $R+2$ about $\tilde x$ lies in the $\ell^\infty$-ball of the same radius.
The cube $\tilde\square''$ has side $s_j$ and contains $\tilde y$, so it lies in the $\ell^\infty$-ball of
radius $R+2+s_j$ about $\tilde x$. The cube $\tilde\square$ shares a point with $\tilde\square''$ and has side
$s_j$, so it lies in the $\ell^\infty$-ball of radius $\eta:=R+2+2s_j$ about $\tilde x$.

Now count grid cubes meeting such a ball, one coordinate at a time. An interval $[ks_j,(k+1)s_j]$,
$k\in\mathbb{Z}$, meets an interval $[t-\eta,t+\eta]$, $t\in\mathbb{R}$, only if
\[
(t-\eta)/s_j-1\le k\le (t+\eta)/s_j .
\]
There are at most $2\eta/s_j+2$ such $k$. Hence at most $(2\eta/s_j+2)^4$ grid cubes of side $s_j$ meet the
$\ell^\infty$-ball of radius $\eta$. With $\eta=R+2+2s_j$,
\[
\frac{2\eta}{s_j}+2=\frac{2(R+2)}{s_j}+6=\frac{2(R+2)L^{K'-j}}{M}+6,
\]
which gives the bound \eqref{6.2:eq-tree-sum-a}.

\emph{(a), plaquettes.} Let $p\in\operatorname{plaq}(X^+)$ with $x_p\in\bar B(x,R)$. By definition of
$\operatorname{plaq}(X^+)$, the plaquette $p$ meets a $\pi_j$-cube $\square''$ of $X^+$ at a point $y$. The
plaquette $p$ is a square of side $\varepsilon=L^{-K'}\le 1$, so its Euclidean diameter is
$\sqrt2\,\varepsilon<2$. Hence
\[
|y-x|_{\mathbb{T}}\le|y-x_p|_{\mathbb{T}}+|x_p-x|_{\mathbb{T}}<2+R,
\]
and $\square''$ meets $\bar B(x,R+2)$. Since $X^+$ is the one-layer enlargement of $X$, the cube $\square''$
is adjacent to at least one cube of $X$. Every cube of $X$ adjacent to $\square''$ is, by definition, among
the cubes counted in the first assertion of (a), so there are at most $\mathfrak{a}(R,K'-j)$ of them.

\emph{(b).} Let $X\in\mathcal{Y}$. Since $(1-\delta)\kappa=c_\kappa+\kappa_{\mathrm{adm}}$, $c_\kappa>0$ and
$d_j(X)\ge D$,
\[
e^{-(1-\delta)\kappa\, d_j(X)}=e^{-c_\kappa d_j(X)}e^{-\kappa_{\mathrm{adm}} d_j(X)}
\le e^{-c_\kappa D}\,e^{-\kappa_{\mathrm{adm}} d_j(X)} .
\]
Assign to each $X\in\mathcal{Y}$ one cube $\square'_X\in\mathcal{F}$ with $X\supset\square'_X$; such a
cube exists by hypothesis. All terms are non-negative, and each $X\in\mathcal{Y}$ occurs in the inner sum
below for $\square'=\square'_X$. Therefore
\begin{align*}
\sum_{X\in\mathcal{Y}}C_e\,e^{-(1-\delta)\kappa\, d_j(X)}
&\le C_e\,e^{-c_\kappa D}\sum_{\square'\in\mathcal{F}}\ \sum_{X\in D_j,\,X\supset\square'}
e^{-\kappa_{\mathrm{adm}} d_j(X)}\\
&\le C_e\,e^{-c_\kappa D}\,(\#\mathcal{F})\,C_{(1.26)} ,
\end{align*}
where the second inequality is the admissibility of $\kappa_{\mathrm{adm}}$ for \cite[(1.26)]{B-CMP116},
applied to each cube $\square'\in\mathcal{F}$. Since $\#\mathcal{F}\le\mathfrak{a}(R,K'-j)$, the bound
\eqref{6.2:eq-tree-sum-b} follows.
\end{proof}

\begin{lemma}[The ultraviolet bound for separated ball sources]\label{6.2:lem-ultraviolet-bound}
Let $h_1,\dots,h_n$ be real Lipschitz functions on $\mathbb{T}$. Suppose that each $\operatorname{supp}h_l$
lies in a ball of radius $R$ and that the supports are pairwise at distance at least $d$. For
$0\le j\le K'$ let $\mathcal{Y}_j\subset D_j$ be an arbitrary sub-family of localisation domains.
For $\zeta=(\zeta_1,\dots,\zeta_n)$ in
\begin{equation*}
\mathbb{D}_h:=\Bigl\{\zeta\in\mathbb{C}^n:\ \sum_{l=1}^n|\zeta_l|\,\|h_l\|_{\sharp}<r\Bigr\}
\end{equation*}
put
\begin{equation}\label{6.2:eq-ultraviolet-bound-1}
\Phi(\zeta):=\sum_{j=0}^{K'}\ \sum_{X\in\mathcal{Y}_j}e_j\Bigl(X;\sum_{l=1}^n\zeta_l h_l\Bigr).
\end{equation}
Then $\Phi$ is holomorphic on $\mathbb{D}_h$ and $\Phi(0)=0$. Moreover
\begin{equation}\label{6.2:eq-ultraviolet-bound-2}
\bigl|\partial_{\zeta_1}\cdots\partial_{\zeta_n}\Phi(0)\bigr|
\le C_{\mathrm{UV}}(d;R)\Bigl(\frac{n}{r}\Bigr)^{n}\prod_{l=1}^n\|h_l\|_{\sharp},
\end{equation}
where
\begin{equation}\label{6.2:eq-ultraviolet-bound-a}
C_{\mathrm{UV}}(d;R):=C_eC_{(1.26)}\sum_{i\ge0}\mathfrak{a}(R,i)\,
e^{-c_\kappa\left[dL^i/(4M)-4\right]_+}<\infty .
\end{equation}
The constant $C_{\mathrm{UV}}(d;R)$ does not depend on $K$, $g_0$, $m$, $N$ or on the positions of
the supports, and it is non-increasing in $d$.
\end{lemma}

Here $[x]_+:=\max\{x,0\}$, and $\mathfrak{a}(R,i)$ is the anchor count of
\eqref{6.2:eq-tree-sum-a}. The constants $C_e,\delta,\kappa$ are those of the bound
\eqref{6.2:eq-localised-terms-c}, and $c_\kappa>0$ and $C_{(1.26)}$ are those of the tree-graph sum
\eqref{6.2:eq-tree-sum-b}. The latter rests on \cite[(1.26)]{B-CMP116}.

\begin{proof}
For each $l$ fix a point $y_l\in\mathbb{T}$ with $\operatorname{supp}h_l\subset\overline{B}(y_l,R)$.
Here $\overline{B}(y,\cdot)$ denotes the closed ball in $\mathbb{T}$ of the indicated radius about $y$.

We first prove holomorphy and the vanishing at the origin. For $\zeta\in\mathbb{D}_h$ the weight
$\sum_l\zeta_lh_l$ is the complex source $w_\zeta$ of the family $(h_l)$, and its gauge is
$\rho(\zeta)=r^{-1}\sum_l|\zeta_l|\,\|h_l\|_{\sharp}$. Hence $\mathbb{D}_h$ is exactly the source
polydisc $\{\rho<1\}$ for this family. The family $(h_l)_{l\le n}$, consisting of real Lipschitz
functions on $\mathbb{T}$ with supports in balls of radius $R$ pairwise at distance at least $d$, is
a system of test functions to which the correlation theorem applies as it is stated. By the
correlation theorem the localised terms $e_j(X;w)$ are therefore defined for every source $w$ of
this system with gauge $\rho<1$, that is, for $w=\sum_l\zeta_lh_l$ with $\zeta\in\mathbb{D}_h$. By
property \eqref{6.2:eq-localised-terms-b} of the localised terms
(Definition~\ref{6.2:def-localised-terms}), each map $v\mapsto e_j(X;v)$ is holomorphic where
defined, hence on these sources.
The map $\zeta\mapsto\sum_l\zeta_lh_l$ is linear, so each summand of
\eqref{6.2:eq-ultraviolet-bound-1} is holomorphic on $\mathbb{D}_h$. The sum is finite: $j$ ranges
over $0\le j\le K'$, and each $D_j$ consists of unions of cubes of the partition $\pi_j$ of the
finite torus. Hence $\Phi$ is holomorphic on $\mathbb{D}_h$. At $\zeta=0$ the weight vanishes
identically. By \eqref{6.2:eq-localised-terms-a} every $e_j(X;0)$ is then zero, and so $\Phi(0)=0$.

We next show that most domains do not contribute to the mixed derivative. Fix $j$ and
$X\in\mathcal{Y}_j$. Say that $X^+$ meets $\operatorname{supp}h_l$ if some plaquette
$p\in\operatorname{plaq}(X^+)$ has $h_l(x_p)\ne0$. Suppose that $X^+$ does not meet
$\operatorname{supp}h_l$. Then changing $\zeta_l$ does not change the values of $\sum_k\zeta_kh_k$ at
the barycentres of $\operatorname{plaq}(X^+)$. By the locality property
\eqref{6.2:eq-localised-terms-a}, the function $\zeta\mapsto e_j(X;\sum_k\zeta_kh_k)$ is therefore
independent of $\zeta_l$, and its derivative in $\zeta_l$ vanishes identically. Consequently
\begin{equation*}
\partial_{\zeta_1}\cdots\partial_{\zeta_n}\,e_j\Bigl(X;\sum_k\zeta_kh_k\Bigr)\Big|_{\zeta=0}=0
\end{equation*}
unless $X^+$ meets every $\operatorname{supp}h_l$, $1\le l\le n$. Let $\mathcal{Y}_j'$ be the set of
those $X\in\mathcal{Y}_j$ for which $X^+$ meets every $\operatorname{supp}h_l$. In particular, if
$\mathcal{Y}_j'\neq\emptyset$ then no $h_l$ vanishes identically, so $\|h_l\|_{\sharp}>0$ for every $l$.

For $X\in\mathcal{Y}_j'$ the enlarged domain $X^+$ meets $\operatorname{supp}h_1$ and
$\operatorname{supp}h_2$ in the sense just described. The sets of barycentres at which $h_1$ and
$h_2$ are non-zero are at distance at least $d$. The lower bound
\eqref{6.2:eq-domain-diameter-a}, a consequence of the diameter lemma
(Lemma~\ref{6.2:lem-domain-diameter}), applies with separation $d$. Since $d_j(X)\ge0$, it gives
\begin{equation}\label{6.2:eq-ultraviolet-bound-3}
d_j(X)\ \ge\ \Delta_j:=\Bigl[\frac{dL^{K'-j}}{4M}-4\Bigr]_+ .
\end{equation}

We now estimate each contributing term by Cauchy's inequality. Fix $X\in\mathcal{Y}_j'$ and
$0<\varepsilon'<1$, and set $\sigma_l:=(1-\varepsilon')r/(n\|h_l\|_{\sharp})$. On the closed polydisc
$P:=\prod_l\{|\zeta_l|\le\sigma_l\}$ we have
\begin{equation*}
\sum_l|\zeta_l|\,\|h_l\|_{\sharp}\le(1-\varepsilon')r<r,
\end{equation*}
so $P\subset\mathbb{D}_h$. Write $F(\zeta):=e_j(X;\sum_l\zeta_lh_l)$. Cauchy's integral formula on the
distinguished boundary of $P$ reads
\begin{equation*}
\partial_{\zeta_1}\cdots\partial_{\zeta_n}F(0)=\frac{1}{(2\pi i)^n}\oint_{|\zeta_1|=\sigma_1}\cdots
\oint_{|\zeta_n|=\sigma_n}\frac{F(\zeta)}{\zeta_1^2\cdots\zeta_n^2}\,d\zeta_1\cdots d\zeta_n .
\end{equation*}
It gives $|\partial_{\zeta_1}\cdots\partial_{\zeta_n}F(0)|\le\sup_P|F|\prod_l\sigma_l^{-1}$.
Bounding $\sup_P|F|$ by \eqref{6.2:eq-localised-terms-c} yields
\begin{equation*}
\bigl|\partial_{\zeta_1}\cdots\partial_{\zeta_n}F(0)\bigr|\le C_ee^{-(1-\delta)\kappa d_j(X)}
\prod_{l=1}^n\frac{n\|h_l\|_{\sharp}}{(1-\varepsilon')r}.
\end{equation*}
Letting $\varepsilon'\downarrow0$ we obtain
\begin{equation}\label{6.2:eq-ultraviolet-bound-4}
\bigl|\partial_{\zeta_1}\cdots\partial_{\zeta_n}F(0)\bigr|\le C_ee^{-(1-\delta)\kappa d_j(X)}
\Bigl(\frac{n}{r}\Bigr)^{n}\prod_{l=1}^n\|h_l\|_{\sharp}.
\end{equation}

Next we count the contributing domains and sum over them. Let $X\in\mathcal{Y}_j'$. There is a
plaquette $p\in\operatorname{plaq}(X^+)$ with $h_1(x_p)\ne0$, hence with
$x_p\in\overline{B}(y_1,R)$. By the parenthetical observation accompanying the anchor count in the
tree-graph lemma (Lemma~\ref{6.2:lem-tree-sum}), $X$ contains a $\pi_j$-cube lying within one layer
of a cube that meets $\overline{B}(y_1,R+2)$. By \eqref{6.2:eq-tree-sum-a}, at most
$\mathfrak{a}(R,K'-j)$ such cubes exist, and they are determined by $y_1$, $R$ and $j$ alone. By
\eqref{6.2:eq-ultraviolet-bound-3}, every $X\in\mathcal{Y}_j'$ also satisfies $d_j(X)\ge\Delta_j$.
The tree-graph sum \eqref{6.2:eq-tree-sum-b}, applied to the family $\mathcal{Y}_j'$ with threshold
$\Delta_j$, therefore gives
\begin{equation}\label{6.2:eq-ultraviolet-bound-5}
\sum_{X\in\mathcal{Y}_j'}C_ee^{-(1-\delta)\kappa d_j(X)}\le C_eC_{(1.26)}\,\mathfrak{a}(R,K'-j)\,
e^{-c_\kappa\Delta_j}.
\end{equation}

We now sum over $j$. The derivative of the finite sum \eqref{6.2:eq-ultraviolet-bound-1} is the sum
of the derivatives, and only $X\in\mathcal{Y}_j'$ contribute. Combining
\eqref{6.2:eq-ultraviolet-bound-4} and \eqref{6.2:eq-ultraviolet-bound-5} we obtain
\begin{equation*}
\bigl|\partial_{\zeta_1}\cdots\partial_{\zeta_n}\Phi(0)\bigr|
\le\Bigl(\frac{n}{r}\Bigr)^{n}\prod_{l=1}^n\|h_l\|_{\sharp}\
\sum_{j=0}^{K'}C_eC_{(1.26)}\,\mathfrak{a}(R,K'-j)\,e^{-c_\kappa\Delta_j}.
\end{equation*}
Substitute $i:=K'-j\in[0,K']$. Then $\Delta_j=[dL^i/(4M)-4]_+$, and the sum over $j$ becomes a
partial sum of the series of non-negative terms in \eqref{6.2:eq-ultraviolet-bound-a}. It is
therefore majorised by the full sum over $i\ge0$, which proves \eqref{6.2:eq-ultraviolet-bound-2}.

It remains to prove that the series converges. Put $\eta_i:=dL^i/(4M)$, so that
$2(R+2)L^i/M=8(R+2)\eta_i/d$ and $\mathfrak{a}(R,i)=(8(R+2)\eta_i/d+6)^4$.

Consider first the indices with $\eta_i\ge8$. For these $[\eta_i-4]_+\ge\eta_i/2$. Since $6\le\eta_i$,
we also have $\mathfrak{a}(R,i)\le(1+8(R+2)/d)^4\eta_i^4$. Hence the $i$-th summand of
\eqref{6.2:eq-ultraviolet-bound-a}, without the prefactor $C_eC_{(1.26)}$, is at most
\begin{equation*}
\Bigl(1+\frac{8(R+2)}{d}\Bigr)^4\eta_i^4e^{-c_\kappa\eta_i/2}
\le\Bigl(1+\frac{8(R+2)}{d}\Bigr)^4\Bigl(\frac{16}{c_\kappa}\Bigr)^4e^{-4}\,e^{-c_\kappa\eta_i/4}.
\end{equation*}
The second inequality holds because $\max_{x\ge0}x^4e^{-c_\kappa x/4}=(16/c_\kappa)^4e^{-4}$. Let
$i_0$ be the least index with $\eta_{i_0}\ge8$. Then $\eta_{i_0+k}=\eta_{i_0}L^k\ge8L^k$ for all
$k\ge0$. Since $L>1$, we have $L^k\ge1+k(L-1)$, so $\sum_{k\ge0}e^{-2c_\kappa L^k}$ is dominated by a
convergent geometric series. Its value depends on $c_\kappa$ and $L$ only. Thus this part of the
series converges.

The remaining indices satisfy $L^i<32M/d$. There are at most $1+\log^+_L(32M/d)$ of them, where
$\log^+_L:=\max\{\log_L,0\}$. For each of them $\mathfrak{a}(R,i)\le(2(R+2)\cdot32/d+6)^4$, and the
exponential factor is at most one. So each such summand, including the prefactor, is at most
$C_eC_{(1.26)}(64(R+2)/d+6)^4$.

Using the inequalities
\begin{equation*}
\Bigl(1+\frac{8(R+2)}{d}\Bigr)^4\le8^4\Bigl(1+\frac{R+2}{d}\Bigr)^4,\qquad
\Bigl(\frac{64(R+2)}{d}+6\Bigr)^4\le64^4\Bigl(1+\frac{R+2}{d}\Bigr)^4,
\end{equation*}
we conclude that
\begin{equation*}
C_{\mathrm{UV}}(d;R)\le C(L,\mathfrak{p})\Bigl(1+\frac{R+2}{d}\Bigr)^4
\Bigl(1+\log^+_L\frac{32M}{d}\Bigr)<\infty .
\end{equation*}
Here $C(L,\mathfrak{p})$ depends only on $L$, $C_e$, $C_{(1.26)}$ and $c_\kappa$.

Finally, the right side of \eqref{6.2:eq-ultraviolet-bound-a} is built from $C_e$, $C_{(1.26)}$,
$c_\kappa$, $L$, $M$, $R$ and $d$ alone. The first three are the constants of
\eqref{6.2:eq-localised-terms-c} and \eqref{6.2:eq-tree-sum-b}, and neither $K$, $g_0$, $m$, $N$ nor
the positions of the supports enter the expression. For each $i$ the quantity
$[dL^i/(4M)-4]_+$ is non-decreasing in $d$, and $c_\kappa>0$. Hence every summand, and with it
$C_{\mathrm{UV}}(d;R)$, is non-increasing in $d$.
\end{proof}

\begin{lemma}[The bound on domains meeting a cell face]\label{6.2:lem-face-estimate}
Let $R>0$ and let $v$ be a weight on $\mathbb{T}$ of the form $v=\sum_{\iota\in\mathsf{K}}v_\iota$, where
$\mathsf{K}$ is a finite index set and, for each $\iota\in\mathsf{K}$, the weight $v_\iota$ is supported,
at the barycentres, in a ball copy $B'_\iota$ of radius $R$: that is, $v_\iota(x_p)=0$ unless
$x_p\in B'_\iota$. Suppose that each $B'_\iota$ lies at distance at least $R+2$ from every face plane,
and that the bound \eqref{6.2:eq-localised-terms-c} holds at $v$. For each $j$ let
$\mathcal{X}_\partial\subset D_j$ be the family of face-meeting domains of
\eqref{6.2:eq-domain-diameter-c}. Then
\begin{equation}\label{6.2:eq-face-estimate-1}
\sum_{j\le K'}\ \sum_{X\in\mathcal{X}_\partial}\bigl|e_j(X;v)\bigr|
\le |\mathsf{K}|\cdot C_{\mathrm{face}}(R),
\end{equation}
where
\begin{equation}\label{6.2:eq-face-estimate-a}
C_{\mathrm{face}}(R):=C_e\,C_{(1.26)}\sum_{i\ge 0}\mathfrak{a}(R,i)\,
e^{-c_\kappa\left[RL^i/(4M)-3\right]_+}<\infty ,
\end{equation}
with $\mathfrak{a}(R,i)$ the anchor count of \eqref{6.2:eq-tree-sum-a} and $C_{(1.26)}$, $c_\kappa$ the
constant and the retained rate of \eqref{6.2:eq-tree-sum-b}. The constant $C_{\mathrm{face}}(R)$ is free
of $K$, $g_0$, $m$, $N$, $\ell$ and of the positions of the ball copies.
\end{lemma}

\begin{proof}
For each $\iota\in\mathsf{K}$ let $x_\iota$ be the centre of $B'_\iota$, so that $B'_\iota$ is the ball
of radius $R$ about $x_\iota$.

\emph{Nonzero terms.} Fix $j\le K'$ and $X\in\mathcal{X}_\partial$ with $e_j(X;v)\neq 0$. By the
locality property \eqref{6.2:eq-localised-terms-a} of Definition~\ref{6.2:def-localised-terms},
$e_j(X;v)=0$ if $v$ vanishes at the barycentres of $\operatorname{plaq}(X^+)$; hence there is a
plaquette $p\in\operatorname{plaq}(X^+)$ with $v(x_p)\neq 0$. Since $v(x_p)=\sum_{\iota}v_\iota(x_p)$,
there is an index $\iota\in\mathsf{K}$ with $v_\iota(x_p)\neq 0$, and then $x_p\in B'_\iota$ by the
support hypothesis. Thus $X^+$ meets the weight $v_\iota$, supported in $B'_\iota$, in the sense of
\eqref{6.2:eq-localised-terms-a}. Moreover, since $X\in\mathcal{X}_\partial$, the enlarged domain $X^+$
contains a point of a face plane, by \eqref{6.2:eq-domain-diameter-c} of the diameter lemma
(Lemma~\ref{6.2:lem-domain-diameter}).

\emph{Lower bound on the size.} $X^+$ thus contains a point of a face plane and meets a weight supported
in a ball copy at distance at least $R+2$ from the faces. By \eqref{6.2:eq-domain-diameter-b} of
Lemma~\ref{6.2:lem-domain-diameter},
\begin{equation*}
d_j(X)\ge \frac{RL^{K'-j}}{4M}-3 .
\end{equation*}
Since $d_j(X)$ is the length of a tree measured in units of $s_j$, it is nonnegative, and therefore
\begin{equation}\label{6.2:eq-face-estimate-2}
d_j(X)\ge \Bigl[\frac{RL^{K'-j}}{4M}-3\Bigr]_+ .
\end{equation}

\emph{Grouping.} For $j\le K'$ and $\iota\in\mathsf{K}$ let $\mathcal{Y}_{j,\iota}$ be the family of
those $X\in\mathcal{X}_\partial$ for which some $p\in\operatorname{plaq}(X^+)$ satisfies
$v_\iota(x_p)\neq 0$. By the first step, every $X\in\mathcal{X}_\partial$ with $e_j(X;v)\neq0$ belongs
to at least one $\mathcal{Y}_{j,\iota}$. A domain may belong to several of these families, but all terms
are nonnegative, so counting it more than once only increases the sum:
\begin{equation}\label{6.2:eq-face-estimate-3}
\sum_{X\in\mathcal{X}_\partial}\bigl|e_j(X;v)\bigr|
\le\sum_{\iota\in\mathsf{K}}\ \sum_{X\in\mathcal{Y}_{j,\iota}}\bigl|e_j(X;v)\bigr| .
\end{equation}

\emph{Anchors.} Let $X\in\mathcal{Y}_{j,\iota}$ and let $p\in\operatorname{plaq}(X^+)$ have
$v_\iota(x_p)\neq0$, so that $x_p\in B'_\iota$. By the parenthetical observation accompanying
\eqref{6.2:eq-tree-sum-a} in Lemma~\ref{6.2:lem-tree-sum}, the cube of $X^+$ that $p$ meets meets the
closed ball of radius $R+2$ about $x_\iota$. The cube of $X$ adjacent to it is therefore one of the
$\pi_j$-cubes lying within one layer of a cube meeting that ball. By \eqref{6.2:eq-tree-sum-a} there are
at most $\mathfrak{a}(R,K'-j)$ such cubes, and they depend only on $x_\iota$, $R$ and $j$. Hence every
domain of $\mathcal{Y}_{j,\iota}$ contains one of at most $\mathfrak{a}(R,K'-j)$ given cubes, and by
\eqref{6.2:eq-face-estimate-2} each satisfies
$d_j(X)\ge [RL^{K'-j}/(4M)-3]_+$.

\emph{Tree-graph sum.} The bound \eqref{6.2:eq-localised-terms-c} holds at $v$, so
$|e_j(X;v)|\le C_e e^{-(1-\delta)\kappa d_j(X)}$. The tree-graph sum with threshold
\eqref{6.2:eq-tree-sum-b} of Lemma~\ref{6.2:lem-tree-sum}, which rests on \cite[(1.26)]{B-CMP116},
applies to $\mathcal{Y}=\mathcal{Y}_{j,\iota}$ with $\mathfrak{a}=\mathfrak{a}(R,K'-j)$ and
$D=[RL^{K'-j}/(4M)-3]_+$. It gives
\begin{equation}\label{6.2:eq-face-estimate-4}
\begin{split}
\sum_{X\in\mathcal{Y}_{j,\iota}}\bigl|e_j(X;v)\bigr|
&\le\sum_{X\in\mathcal{Y}_{j,\iota}}C_e\,e^{-(1-\delta)\kappa d_j(X)}\\
&\le C_e\,C_{(1.26)}\,\mathfrak{a}(R,K'-j)\,e^{-c_\kappa[RL^{K'-j}/(4M)-3]_+}.
\end{split}
\end{equation}

\emph{Summation.} The right side of \eqref{6.2:eq-face-estimate-4} does not depend on $\iota$. Inserting
it in \eqref{6.2:eq-face-estimate-3} therefore produces the factor $|\mathsf{K}|$. We then sum over
$j\le K'$. The index $i=K'-j$ runs over distinct nonnegative integers, and all terms are nonnegative, so
the sum over $j$ is at most the sum over all $i\ge0$:
\begin{equation*}
\sum_{j\le K'}\ \sum_{X\in\mathcal{X}_\partial}\bigl|e_j(X;v)\bigr|
\le|\mathsf{K}|\,C_e\,C_{(1.26)}\sum_{i\ge0}\mathfrak{a}(R,i)\,e^{-c_\kappa[RL^i/(4M)-3]_+}
=|\mathsf{K}|\,C_{\mathrm{face}}(R).
\end{equation*}
This is \eqref{6.2:eq-face-estimate-1}.

\emph{Convergence.} Since $[y-3]_+\ge y-3$, each summand of \eqref{6.2:eq-face-estimate-a} satisfies
\begin{equation*}
\mathfrak{a}(R,i)\,e^{-c_\kappa[RL^i/(4M)-3]_+}
\le e^{3c_\kappa}\Bigl(\frac{2(R+2)L^i}{M}+6\Bigr)^4 e^{-c_\kappa RL^i/(4M)} .
\end{equation*}
Because $c_\kappa R>0$, the function $t\mapsto(2(R+2)t+6)^4e^{-c_\kappa Rt/8}$ is bounded on
$[0,\infty)$. Taking $t=L^i/M$, the summand is therefore at most a constant, depending on $R$, $c_\kappa$
and $M$, times $e^{-c_\kappa RL^i/(8M)}$. The blocking factor satisfies $L\ge2$, so $L^i\ge 2^i\ge i$, and
hence $e^{-c_\kappa RL^i/(8M)}\le e^{-c_\kappa Ri/(8M)}$, which is the general term of a convergent
geometric series. Thus $C_{\mathrm{face}}(R)<\infty$. Finally, the right side of
\eqref{6.2:eq-face-estimate-a} involves only $C_e$, $C_{(1.26)}$, $c_\kappa$, $L$, $M$ and $R$. It
contains none of $K$, $g_0$, $m$, $N$, $\ell$, and no centre $x_\iota$.
\end{proof}

\begin{lemma}[The logarithm of a bounded holomorphic function]\label{6.2:lem-bounded-logarithm}
Let $f_1,\dots,f_n$ be the sources of~\ref{6.1:not-units-sources}. Let
$\rho(z)=r^{-1}\sum_{i=1}^n|z_i|\,\|f_i\|_{\sharp}$ be their gauge on $\mathbb{C}^n$, and let
$\mathfrak{r}\in\{1,2\}$, with $\mathbb{D}_{\mathfrak{r}}=\{z\in\mathbb{C}^n:\rho(z)<1/\mathfrak{r}\}$ the
polydisc of~\eqref{6.2:eq-localised-terms-f}. Let $\mathsf{h}$ be holomorphic on
$\mathbb{D}_{\mathfrak{r}}$, with $\mathsf{h}(0)=1$ and $|\mathsf{h}|\le Q$ on $\mathbb{D}_{\mathfrak{r}}$
for some constant $Q\ge1$. Let $\operatorname{Log}$ denote the principal logarithm. Then:
\begin{enumerate}
\item for every $z\in\mathbb{D}_{\mathfrak{r}}$,
\begin{equation}\label{6.2:eq-bounded-logarithm-1}
|\mathsf{h}(z)-1|\le(Q+1)\,\mathfrak{r}\,\rho(z);
\end{equation}
\item on the set $\{\rho\le 1/(2\mathfrak{r}(Q+1))\}$ one has $|\mathsf{h}-1|\le\tfrac12$; moreover
$\operatorname{Log}\mathsf{h}$ is holomorphic there and satisfies $|\operatorname{Log}\mathsf{h}|\le\log 2$;
\item the mixed derivative at the origin satisfies
\begin{equation}\label{6.2:eq-bounded-logarithm-2}
\bigl|\partial_{z_1}\cdots\partial_{z_n}\operatorname{Log}\mathsf{h}(0)\bigr|
\le\log 2\cdot(2Q+2)^n\,\mathfrak{r}^n\Bigl(\frac{n}{r}\Bigr)^n\prod_{i=1}^n\|f_i\|_{\sharp}.
\end{equation}
\end{enumerate}
\end{lemma}

\begin{proof}
\emph{The bound \eqref{6.2:eq-bounded-logarithm-1}.} At $z=0$ both sides vanish, because
$\mathsf{h}(0)=1$. Now let $z\neq0$ lie in $\mathbb{D}_{\mathfrak{r}}$. For $\lambda\in\mathbb{C}$ with
$|\lambda|<1$ put
\begin{equation*}
\varphi(\lambda):=\mathsf{h}\Bigl(\frac{\lambda z}{\mathfrak{r}\rho(z)}\Bigr)-1 .
\end{equation*}
The gauge $\rho$ is absolutely homogeneous, that is, $\rho(cz)=|c|\rho(z)$ for $c\in\mathbb{C}$. Hence
\begin{equation*}
\rho\Bigl(\frac{\lambda z}{\mathfrak{r}\rho(z)}\Bigr)=\frac{|\lambda|}{\mathfrak{r}}<\frac1{\mathfrak{r}},
\end{equation*}
so the point $\lambda z/(\mathfrak{r}\rho(z))$ lies in $\mathbb{D}_{\mathfrak{r}}$ whenever $|\lambda|<1$.
Consequently $\varphi$ is holomorphic on the open unit disc, as the composition of $\mathsf{h}$ with an
affine map of $\lambda$. The hypothesis $|\mathsf{h}|\le Q$ on $\mathbb{D}_{\mathfrak{r}}$ gives
$|\varphi|\le Q+1$ there, and $\varphi(0)=\mathsf{h}(0)-1=0$.

Apply Schwarz's lemma to $\varphi/(Q+1)$, which maps the open unit disc into the closed unit disc and
vanishes at $0$. This gives
\begin{equation*}
|\varphi(\lambda)|\le(Q+1)|\lambda|\qquad(|\lambda|<1).
\end{equation*}
Since $z\in\mathbb{D}_{\mathfrak{r}}$, the number $\lambda=\mathfrak{r}\rho(z)$ satisfies $|\lambda|<1$.
For this $\lambda$ we have $\lambda z/(\mathfrak{r}\rho(z))=z$, so $\varphi(\lambda)=\mathsf{h}(z)-1$, and
the last inequality becomes $|\mathsf{h}(z)-1|\le(Q+1)\mathfrak{r}\rho(z)$. This proves
\eqref{6.2:eq-bounded-logarithm-1}.

\emph{The bounds on $\operatorname{Log}\mathsf{h}$.} Let $\rho(z)\le 1/(2\mathfrak{r}(Q+1))$. Since
$Q+1>1$, this point lies in $\mathbb{D}_{\mathfrak{r}}$, and \eqref{6.2:eq-bounded-logarithm-1} gives
$|\mathsf{h}(z)-1|\le\tfrac12$.

For $u\in\mathbb{C}$ with $|u|\le\tfrac12$, expand the principal logarithm as
$\operatorname{Log}(1+u)=\sum_{k\ge1}(-1)^{k+1}u^k/k$. Hence
\begin{equation*}
|\operatorname{Log}(1+u)|\le\sum_{k\ge1}\frac{|u|^k}{k}=-\log(1-|u|)\le\log2 .
\end{equation*}
Applied with $u=\mathsf{h}(z)-1$, this gives $|\operatorname{Log}\mathsf{h}(z)|\le\log2$.

For holomorphy, consider the open set $\{z:(Q+1)\mathfrak{r}\rho(z)<1\}$. Because $Q+1>1$, it is
contained in $\mathbb{D}_{\mathfrak{r}}$. On it, \eqref{6.2:eq-bounded-logarithm-1} gives
$|\mathsf{h}-1|<1$, hence $\operatorname{Re}\mathsf{h}>0$. The principal logarithm is holomorphic on the
open right half-plane, so $\operatorname{Log}\mathsf{h}$ is holomorphic on this open set. In particular
it is holomorphic on the set $\{\rho\le 1/(2\mathfrak{r}(Q+1))\}$, which is contained in it.

\emph{The bound \eqref{6.2:eq-bounded-logarithm-2}.} Define radii and a closed polydisc by
\begin{equation*}
\rho_i:=\frac{r}{2\mathfrak{r}(Q+1)\,n\,\|f_i\|_{\sharp}},\qquad
P:=\prod_{i=1}^n\{z_i\in\mathbb{C}:|z_i|\le\rho_i\}.
\end{equation*}
For $z\in P$,
\begin{equation*}
\rho(z)\le\frac1r\sum_{i=1}^n\rho_i\|f_i\|_{\sharp}=\sum_{i=1}^n\frac{1}{2\mathfrak{r}(Q+1)n}
=\frac{1}{2\mathfrak{r}(Q+1)} .
\end{equation*}
Thus $P$ is contained in $\{\rho\le1/(2\mathfrak{r}(Q+1))\}$, and hence in the open set
$\{(Q+1)\mathfrak{r}\rho<1\}$ on which $\operatorname{Log}\mathsf{h}$ is holomorphic. By the preceding
step, $\sup_P|\operatorname{Log}\mathsf{h}|\le\log2$.

The derivative $\partial_{z_1}\cdots\partial_{z_n}\operatorname{Log}\mathsf{h}(0)$ is the coefficient of
the monomial $z_1\cdots z_n$ in the Taylor expansion of $\operatorname{Log}\mathsf{h}$ at $0$. By Cauchy's
integral formula on the distinguished boundary $\{|\zeta_i|=\rho_i,\ i=1,\dots,n\}$ of $P$,
\begin{equation*}
\partial_{z_1}\cdots\partial_{z_n}\operatorname{Log}\mathsf{h}(0)
=\frac{1}{(2\pi i)^n}\oint\cdots\oint
\frac{\operatorname{Log}\mathsf{h}(\zeta)}{\zeta_1^2\cdots\zeta_n^2}\,d\zeta_1\cdots d\zeta_n .
\end{equation*}
Bounding the integrand by $\sup_P|\operatorname{Log}\mathsf{h}|\,\prod_i\rho_i^{-2}$ and the length of
each circle by $2\pi\rho_i$ gives Cauchy's estimate
\begin{equation*}
\bigl|\partial_{z_1}\cdots\partial_{z_n}\operatorname{Log}\mathsf{h}(0)\bigr|
\le\frac{\sup_P|\operatorname{Log}\mathsf{h}|}{\prod_{i=1}^n\rho_i}
\le\log2\cdot\prod_{i=1}^n\frac{2\mathfrak{r}(Q+1)\,n\,\|f_i\|_{\sharp}}{r}.
\end{equation*}
The right-hand side equals
$\log2\cdot(2Q+2)^n\mathfrak{r}^n(n/r)^n\prod_i\|f_i\|_{\sharp}$, which is
\eqref{6.2:eq-bounded-logarithm-2}.
\end{proof}

\begin{lemma}[The chessboard inequality with bulk and face terms]\label{6.2:lem-chessboard-inequality}
Let $\mathfrak{r}\in\{1,2\}$ be the reading index and $z\in\mathbb{D}_{\mathfrak{r}}$, see
\eqref{6.2:eq-localised-terms-f}. Let $f_1,\dots,f_n$ be the test functions, $n$ being the number
of sources, and let the sources be placed so that every ball $B_i$ of radius $R$ carrying
$\operatorname{supp} f_i$ lies in the open cell $c^\circ$ of some cell $c\in\mathcal{C}$ and is at
distance at least $R+2$ from every face plane. For $c\in\mathcal{C}$ let $I_c$ be the set of
indices $i$ with $B_i\subset c^\circ$, so that $\sum_c|I_c|=n$, call $c$ occupied if
$I_c\neq\emptyset$, and let $q$ be the number of occupied cells. Then
\begin{equation}\label{6.2:eq-chessboard-inequality-a}
\begin{split}
\Bigl|\frac{\widetilde{Z}(w_z)}{\widetilde{Z}(0)}\Bigr|
&= e^{-\operatorname{Re}\mathcal{E}(w_z)}\,\Bigl|\frac{Z(w_z)}{Z(0)}\Bigr|\\
&\le \exp\Bigl\{\sum_{c\ \mathrm{occupied}}|\mathcal{C}|^{-1}
\operatorname{Re}\mathcal{E}(w^{\mathrm{ch}}_{c,z})-\operatorname{Re}\mathcal{E}(w_z)
+q\,(E_+^{\sharp}+C_{\mathrm{low}})\,\ell^4\Bigr\}.
\end{split}
\end{equation}
Moreover, let $\mathcal{X}_{c'}$ and $\mathcal{X}_\partial$ be the classes of localisation domains
of \eqref{6.2:eq-domain-diameter-c}, and put, for every cell $c'$ and every weight $v$ at which the
localised terms are defined,
\begin{equation}\label{6.2:eq-chessboard-inequality-b}
\begin{gathered}
\mathcal{E}^{\mathrm{bulk}}_{c'}(v):=\sum_{j}\sum_{X\in\mathcal{X}_{c'}}e_j(X;v),\qquad
\mathcal{E}^{\mathrm{bulk}}:=\sum_{c'}\mathcal{E}^{\mathrm{bulk}}_{c'},\\
\mathcal{E}^{\partial}(v):=\sum_{j}\sum_{X\in\mathcal{X}_\partial}e_j(X;v).
\end{gathered}
\end{equation}
Then $\mathcal{E}=\mathcal{E}^{\mathrm{bulk}}+\mathcal{E}^{\partial}$, and the exponent on the
right-hand side of \eqref{6.2:eq-chessboard-inequality-a} is at most
\begin{equation}\label{6.2:eq-chessboard-inequality-1}
\begin{split}
&\sum_{c\ \mathrm{occupied}}\Bigl\{|\mathcal{C}|^{-1}\sum_{g\in\mathfrak{G}_\ell}
\operatorname{Re}\mathcal{E}^{\mathrm{bulk}}_{gc}\bigl(u^{(\varepsilon(g))}_{c,g,z}\bigr)
-\operatorname{Re}\mathcal{E}^{\mathrm{bulk}}_{c}(w_{c,z})\Bigr\}\\
&\qquad+\Bigl(\sum_{c}|I_c|+n\Bigr)C_{\mathrm{face}}(R)
+q\,(E_+^{\sharp}+C_{\mathrm{low}})\,\ell^4 .
\end{split}
\end{equation}
\end{lemma}

\begin{proof}
Fix $z\in\mathbb{D}_{\mathfrak{r}}$ and the placed sources.

\emph{The per-cell bound.} Let $c$ be occupied. By Lemma~\ref{6.2:lem-periodised-weight}(b),
$\|F_{c,z}\|_{\mathrm{ch}}^{|\mathcal{C}|}=Z(w^{\mathrm{ch}}_{c,z})/Z(0)$, a non-negative real
number, with $w^{\mathrm{ch}}_{c,z}$ the periodised weight of
\eqref{6.2:eq-periodised-weight-a}. By the locality property \eqref{6.2:eq-localised-terms-a}
every $e_j(X;0)$ vanishes, so $\mathcal{E}(0)=0$ and $\widetilde{Z}(0)=Z(0)$; and $Z(0)>0$, being
the integral of an everywhere positive function against the Haar measure. By the definition
$\widetilde{Z}(v)=e^{-\mathcal{E}(v)}Z(v)$ in \eqref{6.2:eq-localised-terms-d}, and since the
non-negative number $Z(w^{\mathrm{ch}}_{c,z})/Z(0)$ equals its modulus,
\begin{equation*}
\begin{split}
\frac{Z(w^{\mathrm{ch}}_{c,z})}{Z(0)}
&=\frac{\bigl|e^{\mathcal{E}(w^{\mathrm{ch}}_{c,z})}\widetilde{Z}(w^{\mathrm{ch}}_{c,z})\bigr|}
{\widetilde{Z}(0)}
=e^{\operatorname{Re}\mathcal{E}(w^{\mathrm{ch}}_{c,z})}\,
\frac{|\widetilde{Z}(w^{\mathrm{ch}}_{c,z})|}{\widetilde{Z}(0)}\\
&\le e^{\operatorname{Re}\mathcal{E}(w^{\mathrm{ch}}_{c,z})}\,
e^{(E_+^{\sharp}+C_{\mathrm{low}})N^4}.
\end{split}
\end{equation*}
In the last step the stability bound \eqref{6.2:eq-localised-terms-d}, which holds at the
periodised weight for $z\in\mathbb{D}_{\mathfrak{r}}$ by Remark~\ref{6.2:rem-transported-weights},
gives $|\widetilde{Z}(w^{\mathrm{ch}}_{c,z})|\le e^{E_+^{\sharp}N^4}$, and the floor bound
\eqref{6.2:eq-localised-terms-e} bounds $1/\widetilde{Z}(0)$ by $e^{C_{\mathrm{low}}N^4}$. The
first two steps are identities and the third uses only the two extensive bounds
\eqref{6.2:eq-localised-terms-d} and \eqref{6.2:eq-localised-terms-e}; no bound on the modulus of
$\mathcal{E}$ alone is taken. Taking $|\mathcal{C}|$-th roots and using $N^4/|\mathcal{C}|=\ell^4$,
\begin{equation}\label{6.2:eq-chessboard-inequality-2}
\|F_{c,z}\|_{\mathrm{ch}}\le\exp\bigl\{|\mathcal{C}|^{-1}
\operatorname{Re}\mathcal{E}(w^{\mathrm{ch}}_{c,z})+(E_+^{\sharp}+C_{\mathrm{low}})\ell^4\bigr\}.
\end{equation}

\emph{Proof of \eqref{6.2:eq-chessboard-inequality-a}.} By
Lemma~\ref{6.2:lem-periodised-weight}(a), $Z(w_z)/Z(0)=\langle\prod_{c}F_{c,z}\rangle$ with
$F_{c,z}\in\mathfrak{A}_c$ and $F_{c,z}:=1$ on unoccupied cells. The chessboard estimate
(Lemma~\ref{6.2:lem-chessboard-estimate}) bounds the modulus of this expectation by
$\prod_c\|F_{c,z}\|_{\mathrm{ch}}$; the factor of an unoccupied cell is the chessboard norm of the
constant $1$, which equals $1$, since every chessboard copy of $1$ is $1$ and
$\langle 1\rangle=1$. Hence, by \eqref{6.2:eq-chessboard-inequality-2},
\begin{equation*}
\Bigl|\frac{Z(w_z)}{Z(0)}\Bigr|\le\prod_{c\ \mathrm{occupied}}\|F_{c,z}\|_{\mathrm{ch}}
\le\prod_{c\ \mathrm{occupied}}\exp\bigl\{|\mathcal{C}|^{-1}
\operatorname{Re}\mathcal{E}(w^{\mathrm{ch}}_{c,z})+(E_+^{\sharp}+C_{\mathrm{low}})\ell^4\bigr\}.
\end{equation*}
Since $\widetilde{Z}(w_z)=e^{-\mathcal{E}(w_z)}Z(w_z)$, $\widetilde{Z}(0)=Z(0)$ and
$|e^{-\mathcal{E}(w_z)}|=e^{-\operatorname{Re}\mathcal{E}(w_z)}$, the identity in
\eqref{6.2:eq-chessboard-inequality-a} holds; the product over the $q$ occupied cells then gives
the inequality.

\emph{The split of the counterterm.} By \eqref{6.2:eq-domain-diameter-c}, each $D_j$ is the
disjoint union of the classes $\mathcal{X}_{c'}$, $c'\in\mathcal{C}$, and of $\mathcal{X}_\partial$
(the classes $\mathcal{X}_{c'}$ are disjoint because the open cells are); hence
$\mathcal{E}(v)=\mathcal{E}^{\mathrm{bulk}}(v)+\mathcal{E}^{\partial}(v)$ with the definitions
\eqref{6.2:eq-chessboard-inequality-b}.

\emph{Bulk terms.} Let $c$ be occupied. The periodised weight
$w^{\mathrm{ch}}_{c,z}=\sum_{g'\in\mathfrak{G}_\ell}u^{(\varepsilon(g'))}_{c,g',z}$ is a sum over
cells of pieces, the piece in $g'c$ being $u^{(\varepsilon(g'))}_{c,g',z}$, supported in the
copies $g'(B_i)$, $i\in I_c$. Fix $g\in\mathfrak{G}_\ell$. For $X\in\mathcal{X}_{gc}$,
\eqref{6.2:eq-domain-diameter-d} together with the
locality property \eqref{6.2:eq-localised-terms-a} gives
$e_j(X;w^{\mathrm{ch}}_{c,z})=e_j(X;u^{(\varepsilon(g))}_{c,g,z})$, and summing,
\begin{equation*}
\mathcal{E}^{\mathrm{bulk}}_{gc}(w^{\mathrm{ch}}_{c,z})
=\mathcal{E}^{\mathrm{bulk}}_{gc}\bigl(u^{(\varepsilon(g))}_{c,g,z}\bigr).
\end{equation*}
Likewise $w_z=\sum_{c'}w_{c',z}$ with $w_{c',z}$ supported in the balls $B_i$, $i\in I_{c'}$, so
$\mathcal{E}^{\mathrm{bulk}}_{c'}(w_z)=\mathcal{E}^{\mathrm{bulk}}_{c'}(w_{c',z})$; this vanishes
if $c'$ is unoccupied, since then $w_{c',z}=0$ and $e_j(X;0)=0$.

\emph{Face terms.} The weight $w_z$ is the sum of the $n$ pieces $z_if_i$, each supported in a
ball $B_i$ of radius $R$ at distance at least $R+2$ from every face plane, and the bound
\eqref{6.2:eq-localised-terms-c} holds at $w_z$ for $z\in\mathbb{D}_{\mathfrak{r}}$ by the four
properties of the localised terms (Definition~\ref{6.2:def-localised-terms}). The bound on domains
meeting a cell face
(Lemma~\ref{6.2:lem-face-estimate}), with $C_{\mathrm{face}}(R)$ as in
\eqref{6.2:eq-face-estimate-a}, gives $|\mathcal{E}^{\partial}(w_z)|\le n\,C_{\mathrm{face}}(R)$.
The weight $w^{\mathrm{ch}}_{c,z}$ is the sum of $|\mathcal{C}|\cdot|I_c|$ pieces, the
$|\mathcal{C}|$ copies of the $|I_c|$ balls of the cell $c$; each copy $g(B_i)$ is again at
distance at least $R+2$ from every face plane, because the elements of $\mathfrak{G}_\ell$ are
isometries of $\mathbb{T}$ generated by reflections in face planes and so map the set of face
planes onto itself. The bound \eqref{6.2:eq-localised-terms-c} holds at $w^{\mathrm{ch}}_{c,z}$ by
Remark~\ref{6.2:rem-transported-weights}, and Lemma~\ref{6.2:lem-face-estimate} gives
$|\mathcal{E}^{\partial}(w^{\mathrm{ch}}_{c,z})|\le|\mathcal{C}|\cdot|I_c|\cdot
C_{\mathrm{face}}(R)$.

\emph{The reduced exponent.} For each cell $c'$ there is exactly one $g\in\mathfrak{G}_\ell$ with
$c'=gc$, namely $g_{c\to c'}$. Hence, by the bulk identities,
\begin{equation*}
\operatorname{Re}\mathcal{E}(w^{\mathrm{ch}}_{c,z})
=\sum_{g\in\mathfrak{G}_\ell}\operatorname{Re}\mathcal{E}^{\mathrm{bulk}}_{gc}
\bigl(u^{(\varepsilon(g))}_{c,g,z}\bigr)+\operatorname{Re}\mathcal{E}^{\partial}
(w^{\mathrm{ch}}_{c,z}),
\end{equation*}
\begin{equation*}
\operatorname{Re}\mathcal{E}(w_z)
=\sum_{c\ \mathrm{occupied}}\operatorname{Re}\mathcal{E}^{\mathrm{bulk}}_{c}(w_{c,z})
+\operatorname{Re}\mathcal{E}^{\partial}(w_z).
\end{equation*}
Insert both into the exponent of \eqref{6.2:eq-chessboard-inequality-a}. By the face bounds, the
terms $|\mathcal{C}|^{-1}\operatorname{Re}\mathcal{E}^{\partial}(w^{\mathrm{ch}}_{c,z})$, summed
over the occupied cells, contribute at most $\sum_c|I_c|\,C_{\mathrm{face}}(R)$ (for unoccupied
$c$, $|I_c|=0$), and the single term $-\operatorname{Re}\mathcal{E}^{\partial}(w_z)$ contributes at
most $n\,C_{\mathrm{face}}(R)$. This gives \eqref{6.2:eq-chessboard-inequality-1}.
\end{proof}

\begin{proposition}[The bound under cell covariance]\label{6.2:prop-cell-covariance}
Let $f_1,\dots,f_n$ be real Lipschitz functions on $\mathbb{T}$, placed in the cells as in
Lemma~\ref{6.2:lem-chessboard-inequality}, each supported in a ball of radius $R$, with pairwise
distance of supports at least $d$. Let $w_z=\sum_i z_if_i$. Let $w_{c,z}$ and
$u_{c,g,z}=w_{c,z}\circ g^{-1}$ be the cell weight and the transported weight of
\eqref{6.2:eq-periodised-weight-a}, and let $\mathfrak{r}$ and $\mathbb{D}_{\mathfrak{r}}$ be as in
\eqref{6.2:eq-localised-terms-f}. For $g\in\mathfrak{G}_\ell$ write $u^{(\varepsilon(g))}_{c,g,z}$
for $u_{c,g,z}$ if $g$ is even and for its complex conjugate $\bar u_{c,g,z}$ if $g$ is odd. For a
cell $c$ let $c^\circ$ denote the open cell. Assume either both of the following hypotheses:
\begin{enumerate}
\item[(a)] (cell covariance) for every occupied cell $c$, every $g\in\mathfrak{G}_\ell$, every
$j\le K'$ and every $z\in\mathbb{D}_{\mathfrak{r}}$,
\begin{equation}\label{6.2:eq-cell-covariance-1}
\sum_{X\in D_j:\,X^+\subset(gc)^\circ}e_j(X;u_{c,g,z})
=\sum_{Y\in D_j:\,Y^+\subset c^\circ}e_j(Y;w_{c,z});
\end{equation}
\item[(b)] (reality) for the same $c,g,j,z$, with $u=u_{c,g,z}$,
\begin{equation}\label{6.2:eq-cell-covariance-2}
\operatorname{Re}\sum_{X\in D_j:\,X^+\subset(gc)^\circ}e_j(X;\bar u)
=\operatorname{Re}\sum_{X\in D_j:\,X^+\subset(gc)^\circ}e_j(X;u);
\end{equation}
\end{enumerate}
or the following weaker hypothesis:
\begin{enumerate}
\item[(c)] (bounded covariance defect) there is a constant $C_{\mathrm{cov}}$, free of $K$,
$g_0$, $m$ and of the positions of the supports, such that
\begin{equation}\label{6.2:eq-cell-covariance-3}
\sup_{c,g,z}\Bigl|\sum_{j}\operatorname{Re}\!\!\sum_{X^+\subset(gc)^\circ}\!\!
e_j\bigl(X;u^{(\varepsilon(g))}_{c,g,z}\bigr)
-\sum_{j}\operatorname{Re}\!\!\sum_{Y^+\subset c^\circ}\!\!e_j(Y;w_{c,z})\Bigr|
\le C_{\mathrm{cov}},
\end{equation}
the supremum running over occupied cells $c$, over $g\in\mathfrak{G}_\ell$ and over
$z\in\mathbb{D}_{\mathfrak{r}}$, the sums over $j\le K'$ and over $X,Y\in D_j$.
\end{enumerate}
Under (a) and (b) put $C_{\mathrm{cov}}:=0$, and set
\begin{equation}\label{6.2:eq-cell-covariance-4}
Q_\infty:=\exp\bigl[(E_+^{\sharp}+C_{\mathrm{low}})\ell^4+2C_{\mathrm{face}}(R)
+C_{\mathrm{cov}}\bigr].
\end{equation}
Then for every $z\in\mathbb{D}_{\mathfrak{r}}$
\begin{equation}\label{6.2:eq-cell-covariance-5}
\Bigl|\frac{\widetilde{Z}(w_z)}{\widetilde{Z}(0)}\Bigr|
\le\exp\bigl\{2nC_{\mathrm{face}}(R)+n(E_+^{\sharp}+C_{\mathrm{low}})\ell^4
+nC_{\mathrm{cov}}\bigr\}\le Q_\infty^n ,
\end{equation}
and, with $C_{\mathrm{UV}}(d;R)$ as in \eqref{6.2:eq-ultraviolet-bound-a},
\begin{equation}\label{6.2:eq-cell-covariance-6}
\Bigl|\partial_{z_1}\cdots\partial_{z_n}\log Z(w_z)\big|_{z=0}\Bigr|
\le C^\infty_n(d;R)\prod_{i=1}^n\|f_i\|_{\sharp},
\end{equation}
\begin{equation}\label{6.2:eq-cell-covariance-7}
C^\infty_n(d;R):=\bigl[C_{\mathrm{UV}}(d;R)
+\log 2\cdot\mathfrak{r}^n(2Q_\infty^n+2)^n\bigr](n/r)^n .
\end{equation}
\end{proposition}

\begin{proof}
Fix $z\in\mathbb{D}_{\mathfrak{r}}$. For a cell $c'$, the family $\mathcal{X}_{c'}$ of the
splitting \eqref{6.2:eq-chessboard-inequality-b} consists, at scale $j$, of the $X\in D_j$ with
$X^+\subset c'^\circ$. Hence, for every weight $v$,
\begin{equation*}
\operatorname{Re}\mathcal{E}^{\mathrm{bulk}}_{c'}(v)
=\sum_{j\le K'}\operatorname{Re}\sum_{X\in D_j:\,X^+\subset c'^\circ}e_j(X;v).
\end{equation*}
In the proof of Lemma~\ref{6.2:lem-chessboard-inequality} the exponent on the right of
\eqref{6.2:eq-chessboard-inequality-a} is estimated. That estimate uses the splitting
\eqref{6.2:eq-chessboard-inequality-b}, the locality property of
Definition~\ref{6.2:def-localised-terms} and the face estimate, Lemma~\ref{6.2:lem-face-estimate}.
It gives
\begin{equation}\label{6.2:eq-cell-covariance-8}
\begin{split}
\log\Bigl|\frac{\widetilde{Z}(w_z)}{\widetilde{Z}(0)}\Bigr|
\le{}&\sum_{c\ \mathrm{occ}}\Bigl\{|\mathcal{C}|^{-1}\sum_{g\in\mathfrak{G}_\ell}
\operatorname{Re}\mathcal{E}^{\mathrm{bulk}}_{gc}\bigl(u^{(\varepsilon(g))}_{c,g,z}\bigr)
-\operatorname{Re}\mathcal{E}^{\mathrm{bulk}}_{c}(w_{c,z})\Bigr\}\\
&+\Bigl(\sum_c|I_c|+n\Bigr)C_{\mathrm{face}}(R)+q(E_+^{\sharp}+C_{\mathrm{low}})\ell^4 .
\end{split}
\end{equation}
Each cell $c'\in\mathcal{C}$ is the image of $c$ under exactly one element $g_{c\to c'}$ of
$\mathfrak{G}_\ell$ (Definition~\ref{6.2:def-cell-reflections}). So $g\mapsto gc$ is a bijection of
$\mathfrak{G}_\ell$ onto $\mathcal{C}$, and $|\mathfrak{G}_\ell|=|\mathcal{C}|$.

\emph{The braces under (a) and (b).} Let $c$ be occupied and $u=u_{c,g,z}$. If $g$ is even, then
$u^{(\varepsilon(g))}_{c,g,z}=u$. Taking real parts in \eqref{6.2:eq-cell-covariance-1} and
summing over $j\le K'$ gives
$\operatorname{Re}\mathcal{E}^{\mathrm{bulk}}_{gc}(u)=\operatorname{Re}\mathcal{E}^{\mathrm{bulk}}_{c}(w_{c,z})$.
If $g$ is odd, then $u^{(\varepsilon(g))}_{c,g,z}=\bar u$. For each $j$, \eqref{6.2:eq-cell-covariance-2}
first replaces $\bar u$ by $u$ in the real part of the sum over $X^+\subset(gc)^\circ$. Then the real
part of \eqref{6.2:eq-cell-covariance-1} replaces that sum by the sum over $Y^+\subset c^\circ$
at $w_{c,z}$. Summing over $j$ again gives
$\operatorname{Re}\mathcal{E}^{\mathrm{bulk}}_{gc}(\bar u)=\operatorname{Re}\mathcal{E}^{\mathrm{bulk}}_{c}(w_{c,z})$.
So every $g$-term in the brace of \eqref{6.2:eq-cell-covariance-8} equals
$\operatorname{Re}\mathcal{E}^{\mathrm{bulk}}_{c}(w_{c,z})$. Since $|\mathfrak{G}_\ell|=|\mathcal{C}|$,
each brace vanishes.

\emph{The braces under (c).} By $|\mathfrak{G}_\ell|=|\mathcal{C}|$, the brace of an occupied cell
$c$ equals
\begin{equation*}
|\mathcal{C}|^{-1}\sum_{g\in\mathfrak{G}_\ell}\Bigl[
\operatorname{Re}\mathcal{E}^{\mathrm{bulk}}_{gc}\bigl(u^{(\varepsilon(g))}_{c,g,z}\bigr)
-\operatorname{Re}\mathcal{E}^{\mathrm{bulk}}_{c}(w_{c,z})\Bigr].
\end{equation*}
By \eqref{6.2:eq-cell-covariance-3} each bracket is at most $C_{\mathrm{cov}}$, so the brace is at
most $C_{\mathrm{cov}}$.

\emph{The bound \eqref{6.2:eq-cell-covariance-5}.} In both cases the sum of the braces is at most
$qC_{\mathrm{cov}}$, where $C_{\mathrm{cov}}=0$ under (a) and (b). Every occupied cell contains at
least one of the $n$ supports, so $q\le n$. Each support lies in exactly one cell, so
$\sum_c|I_c|=n$. By the floor statement of the correlation theorem,
$E_+^{\sharp}+C_{\mathrm{low}}\ge0$, and therefore
$q(E_+^{\sharp}+C_{\mathrm{low}})\ell^4\le n(E_+^{\sharp}+C_{\mathrm{low}})\ell^4$. Inserting these
facts in \eqref{6.2:eq-cell-covariance-8} gives the first inequality of
\eqref{6.2:eq-cell-covariance-5}. By \eqref{6.2:eq-cell-covariance-4}, the middle member of
\eqref{6.2:eq-cell-covariance-5} equals $Q_\infty^n$.

\emph{Holomorphy.} Let $\mathsf{h}(z):=\widetilde{Z}(w_z)/\widetilde{Z}(0)$. At finite $K$, the map
$z\mapsto Z(w_z)$ is the integral, over a compact product of copies of
$SU(N_{\mathrm{col}})$, of a function that is entire in $z$ locally uniformly in $U$; hence it is
entire. The map $z\mapsto\mathcal{E}(w_z)$ is holomorphic on $\mathbb{D}$ by the analyticity
property \eqref{6.2:eq-localised-terms-b}, because $z\mapsto w_z$ is linear. By the locality
property of Definition~\ref{6.2:def-localised-terms}, $e_j(X;0)=0$; hence $\mathcal{E}(0)=0$ and
$\widetilde{Z}(0)=Z(0)$. Moreover $Z(0)>0$, being the integral of a positive function. Thus
$\mathsf{h}=e^{-\mathcal{E}(w_\cdot)}Z(w_\cdot)/Z(0)$ is holomorphic on $\mathbb{D}$, and so on
$\mathbb{D}_{\mathfrak{r}}\subset\mathbb{D}$, with $\mathsf{h}(0)=1$. By
\eqref{6.2:eq-cell-covariance-5}, $|\mathsf{h}|\le Q_\infty^n$ on $\mathbb{D}_{\mathfrak{r}}$. The
exponent in \eqref{6.2:eq-cell-covariance-4} is non-negative, because
$E_+^{\sharp}+C_{\mathrm{low}}\ge0$, $C_{\mathrm{face}}(R)\ge0$ and $C_{\mathrm{cov}}\ge0$; so
$Q_\infty^n\ge1$.

\emph{The two derivative bounds.} The logarithm lemma, Lemma~\ref{6.2:lem-bounded-logarithm},
applies with $Q=Q_\infty^n$. It shows that $\operatorname{Log}\mathsf{h}$ is holomorphic on
$\{\rho\le1/(2\mathfrak{r}(Q_\infty^n+1))\}$ and that
\begin{equation}\label{6.2:eq-cell-covariance-9}
\bigl|\partial_{z_1}\cdots\partial_{z_n}\operatorname{Log}\mathsf{h}(0)\bigr|
\le\log2\cdot(2Q_\infty^n+2)^n\mathfrak{r}^n(n/r)^n\prod_{i=1}^n\|f_i\|_{\sharp}.
\end{equation}
Next apply Lemma~\ref{6.2:lem-ultraviolet-bound} with $h_l=f_l$ and $\mathcal{Y}_j=D_j$ for every
$j$. The functions $f_l$ are real and Lipschitz on $\mathbb{T}$, supported in balls of radius $R$,
with pairwise support distance at least $d$. With these choices $\Phi(\zeta)=\mathcal{E}(w_\zeta)$.
The polydisc $\mathbb{D}_h=\{\sum_l|\zeta_l|\|f_l\|_{\sharp}<r\}$ is
$\{\rho<1\}=\mathbb{D}$. The lemma gives
\begin{equation}\label{6.2:eq-cell-covariance-10}
\bigl|\partial_{z_1}\cdots\partial_{z_n}\mathcal{E}(w_z)\big|_{z=0}\bigr|
\le C_{\mathrm{UV}}(d;R)(n/r)^n\prod_{i=1}^n\|f_i\|_{\sharp}.
\end{equation}

\emph{Conclusion.} Since $\widetilde{Z}(w_z)=e^{-\mathcal{E}(w_z)}Z(w_z)$, we have
$Z(w_z)=e^{\mathcal{E}(w_z)}\widetilde{Z}(0)\mathsf{h}(z)$. On the neighbourhood of $0$ where
$\operatorname{Log}\mathsf{h}$ is holomorphic, a logarithm of $Z(w_z)$ is therefore
\begin{equation*}
\log Z(w_z)=\mathcal{E}(w_z)+\log\widetilde{Z}(0)+\operatorname{Log}\mathsf{h}(z).
\end{equation*}
It is holomorphic, and at $z=0$ it takes the real value $\log Z(0)$, because $\mathcal{E}(0)=0$ and
$\operatorname{Log}1=0$. The middle term does not depend on $z$, so it is annihilated by
$\partial_{z_1}\cdots\partial_{z_n}$. Adding \eqref{6.2:eq-cell-covariance-9} and
\eqref{6.2:eq-cell-covariance-10} gives \eqref{6.2:eq-cell-covariance-6}, with the constant
\eqref{6.2:eq-cell-covariance-7}.
\end{proof}

Without any of the hypotheses (a), (b), (c) one still obtains an average bound whose constant does
not depend on $m$. Consider a single occupied cell $c$ carrying all $n$ supports, and the weights
$w^{(g)}:=u^{(\varepsilon(g))}_{c,g,z}$, $g\in\mathfrak{G}_\ell$. By
\eqref{6.2:eq-cell-reflections-e} they all have the same periodised weight $w^{\mathrm{ch}}_{c,z}$.
Apply \eqref{6.2:eq-chessboard-inequality-a} to each $w^{(g)}$ and average over $g$, using
$|\mathfrak{G}_\ell|=|\mathcal{C}|$ and the splitting \eqref{6.2:eq-chessboard-inequality-b}. The
locality identities used in the proof of Lemma~\ref{6.2:lem-chessboard-inequality} show that
$|\mathcal{C}|^{-1}\operatorname{Re}\mathcal{E}^{\mathrm{bulk}}(w^{\mathrm{ch}}_{c,z})$ and the
average of $\operatorname{Re}\mathcal{E}^{\mathrm{bulk}}(w^{(g)})$ over $g$ are the same sum
$|\mathcal{C}|^{-1}\sum_g\operatorname{Re}\mathcal{E}^{\mathrm{bulk}}_{gc}(w^{(g)})$, so the bulk
terms cancel. By Lemma~\ref{6.2:lem-face-estimate}, the face terms contribute at most
$nC_{\mathrm{face}}(R)$ each. Hence
\begin{equation*}
|\mathcal{C}|^{-1}\sum_{g\in\mathfrak{G}_\ell}
\log\Bigl|\frac{\widetilde{Z}(w^{(g)})}{\widetilde{Z}(0)}\Bigr|
\le2nC_{\mathrm{face}}(R)+(E_+^{\sharp}+C_{\mathrm{low}})\ell^4 .
\end{equation*}
Hypothesis (c) is exactly what turns this average over the orbit into a bound for each member of it.

\begin{definition}[Standing parameters and flow constants]\label{6.3:not-standing-parameters}
The following parameters and constants are fixed throughout this chapter.

\emph{Fixed parameters.} $L$ is an odd integer with $L \ge L_0$. After $L$, the following are fixed in this order, as in the conventions of the book:
\begin{itemize}
\item Ba{\l}aban's auxiliary parameter list $\mathfrak{p}$, which contains $M = L^{m'} \ge L^4$, $r \ge 2$, $\kappa_0$, \ldots;
\item the coupling window $\gamma$;
\item the terminal ceiling $\gamma_H$.
\end{itemize}
The gauge group is $G = SU(2)$, and $\operatorname{tr}$ is normalised by $\operatorname{tr}\mathbb{1} = 1$. The dimension is $d = 4$. The separation distance of supports used below is written $\mathsf{d}$.

\emph{Flow constants.} All of the following are functions of $(L,\mathfrak{p},\gamma)$ only. Put
\begin{equation}\label{6.3:eq-standing-parameters-1}
c_0 = \frac{11\,C_{\mathfrak{g}}}{24\pi^2}\,\log L > 0,
\qquad C_{\mathfrak{g}} = 2\cdot 2^2 .
\end{equation}
Let $c_2$ be the constant of the coupling-flow estimate that enters the flow bounds, and put $c_5 := 2c_2$. Let $\bar{\varsigma}$ be a number with
\begin{equation*}
0 \le \bar{\varsigma} \le \min\{c_0/4,\,1\}.
\end{equation*}
With these, define
\begin{equation*}
\lambda^{\sharp} := \frac{c_0}{2} - \bar{\varsigma},
\qquad
B^{\sharp} := C_\beta + \bar{\varsigma},
\qquad
g_{-}(g) := \bigl(g^{-2} + B^{\sharp}\bigr)^{-1/2}.
\end{equation*}
The flow letters of the correlation theorem are
\begin{equation*}
\lambda := \frac{c_0}{2},
\qquad
b_{+} := 3\lambda + 2c_2 .
\end{equation*}
Let $R(g)$ be the least power of $L$ that is at least $(\log g^{-2})^{r}$, and put
\begin{equation*}
m_{\mathbb{T}}(g) := 11 + m' + \log_L R(g).
\end{equation*}

With these definitions the following hold:
\begin{itemize}
\item $|\operatorname{tr} U| \le 1$ for every $U \in G$;
\item $0 < \lambda^{\sharp} \le \lambda$;
\item $m_{\mathbb{T}}(g)$ is an integer.
\end{itemize}
\end{definition}

\begin{proof}
\emph{The trace bound.} On $2\times 2$ matrices the normalisation $\operatorname{tr}\mathbb{1} = 1$ means that $\operatorname{tr}$ is one half of the matrix trace. An element $U \in SU(2)$ is unitary, so its two eigenvalues have modulus one. The matrix trace of $U$, being the sum of these eigenvalues, therefore has modulus at most $2$. Halving gives $|\operatorname{tr}U| \le 1$.

\emph{The bounds on $\lambda^{\sharp}$.} We have $c_0 > 0$ by \eqref{6.3:eq-standing-parameters-1}. The condition $\bar{\varsigma} \le c_0/4$ then gives
\begin{equation*}
\lambda^{\sharp} = \frac{c_0}{2} - \bar{\varsigma} \ge \frac{c_0}{4} > 0 .
\end{equation*}
The condition $\bar{\varsigma} \ge 0$ gives
\begin{equation*}
\lambda^{\sharp} = \frac{c_0}{2} - \bar{\varsigma} \le \frac{c_0}{2} = \lambda .
\end{equation*}

\emph{Integrality of $m_{\mathbb{T}}(g)$.} By definition $R(g)$ is a power of $L$, so $\log_L R(g)$ is an integer. Since $11$ and the exponent $m'$ in $M = L^{m'}$ are integers, $m_{\mathbb{T}}(g)$ is an integer.
\end{proof}

\begin{definition}[The torus floor and the least admissible block exponent]
\label{6.3:def-block-exponent-floor}
Let $L$ and $\mathfrak{p}$ be fixed, let $\lambda$, $b_{+}$ and $c_2$ be the flow constants of the
correlation theorem (see~\ref{6.3:not-standing-parameters}), let $M = L^{m'}$, and let $r > 0$ be
the exponent in the ball radius, so that $R(g')$ is the least power of $L$ that is
$\ge (\log g'^{-2})^{r}$.
\begin{enumerate}
\item For $x > 1$, let $R_J(x)$ be the least power of $L$ that is $\ge (\log x)^{r}$. This is the
function $R$ read at an inverse coupling, $R_J(x) = R(x^{-1/2})$, and it is non-decreasing in $x$.
\item Put $C_{\mathbb{T}} := L^{11} M$.
\item For $g \in \,]0,1[$ and an integer $m \ge 0$, the \emph{torus floor} of the correlation
theorem is
\begin{equation}\label{6.3:eq-block-exponent-floor-1}
N_{\mathbb{T}}(g,m) := C_{\mathbb{T}} \sup_{u \in \mathbb{Z},\, u \ge 0}
L^{-u} R_J\bigl(g^{-2} + b_{+} + 4c_2 + 3\lambda(m+u)\bigr).
\end{equation}
\item The \emph{least admissible block exponent} $m^{J}_{\mathbb{T}}(g)$ is the least integer
$m \ge 0$ with $L^{m} \ge N_{\mathbb{T}}(g,m)$.
\end{enumerate}
Then: $N_{\mathbb{T}}(g,m)$ is a power of $L$, the supremum in
\eqref{6.3:eq-block-exponent-floor-1} being attained; $m^{J}_{\mathbb{T}}(g)$ exists for every
$g \in \,]0,1[$; and $m^{J}_{\mathbb{T}}(g) \to \infty$ and
$m_{\mathbb{T}}(g_{-}(g)) \to \infty$ as $g \downarrow 0$.
\end{definition}

\begin{proof}
Fix $g \in \,]0,1[$ and write, for integers $m, u \ge 0$,
\begin{equation*}
x_{m,u} := g^{-2} + b_{+} + 4c_2 + 3\lambda(m+u),
\end{equation*}
the argument of $R_J$ in \eqref{6.3:eq-block-exponent-floor-1}.

We first bound $R_J$. Let $y > 0$. If $y \ge 1$ and $L^{j}$ is the least power of $L$ that is
$\ge y$, then $L^{j-1} < y$, so $L^{j} < Ly$; if $y < 1$, the least power of $L$ that is $\ge y$ is
at most $L^{0} = 1 \le L$. Applied to $y = (\log x)^{r}$ this gives
\begin{equation}\label{6.3:eq-block-exponent-floor-2}
R_J(x) \le L \max\{1, (\log x)^{r}\} \qquad (x > 1).
\end{equation}
Put $a_g := |g^{-2} + b_{+} + 4c_2| + 3\lambda + e$, so that $\log a_g \ge 1$. Since
$1 + m + u \le (1+m)(1+u)$,
\begin{equation*}
x_{m,u} \le |g^{-2} + b_{+} + 4c_2| + 3\lambda(m+u) \le a_g(1+m+u) \le a_g(1+m)(1+u).
\end{equation*}
Hence, when $x_{m,u} > 1$, $\log x_{m,u} \le \log a_g + \log(1+m) + \log(1+u)$. Each of the
three summands is at most $(\log 2)^{-2}\log a_g \log(2+m)\log(2+u)$, because $\log a_g \ge 1$,
the quotients $\log(2+m)/\log 2$ and $\log(2+u)/\log 2$ are $\ge 1$, and $\log 2 < 1$. The
resulting product bound is $\ge 3$, hence $\ge 1$, so
\begin{equation*}
\max\{1, (\log x_{m,u})^{r}\}
\le \bigl(3(\log 2)^{-2}\log a_g\bigr)^{r} (\log(2+m))^{r} (\log(2+u))^{r}.
\end{equation*}
With \eqref{6.3:eq-block-exponent-floor-2} this yields
\begin{equation}\label{6.3:eq-block-exponent-floor-3}
L^{-u} R_J(x_{m,u})
\le L \bigl(3(\log 2)^{-2}\log a_g\bigr)^{r} (\log(2+m))^{r} \, L^{-u}(\log(2+u))^{r}.
\end{equation}
Since $L > 1$, $L^{-u}(\log(2+u))^{r} \to 0$ as $u \to \infty$. Therefore
$c_{L,r} := \sup_{u \in \mathbb{Z},\, u \ge 0} L^{-u}(\log(2+u))^{r}$ is finite.

The supremum is attained. Every term $C_{\mathbb{T}} L^{-u} R_J(x_{m,u})$ of
\eqref{6.3:eq-block-exponent-floor-1} is a power of $L$, because $C_{\mathbb{T}} = L^{11+m'}$ and
$R_J$ takes values in powers of $L$. By \eqref{6.3:eq-block-exponent-floor-3} the terms tend to $0$
as $u \to \infty$. Consequently only finitely many terms are $\ge$ the term with $u = 0$, which is
positive. The largest of these finitely many terms is the supremum. Thus the supremum is attained,
and $N_{\mathbb{T}}(g,m)$ is a power of $L$.

Existence of $m^{J}_{\mathbb{T}}(g)$. Taking the supremum over $u$ in
\eqref{6.3:eq-block-exponent-floor-3} gives
\begin{equation*}
N_{\mathbb{T}}(g,m)
\le L C_{\mathbb{T}}\, c_{L,r} \bigl(3(\log 2)^{-2}\log a_g\bigr)^{r} (\log(2+m))^{r}.
\end{equation*}
For fixed $g$, then, $N_{\mathbb{T}}(g,m)$ grows at most like a power of $\log(2+m)$, while $L^{m}$
grows geometrically, and $L^{m}/(\log(2+m))^{r} \to \infty$. Hence $L^{m} \ge N_{\mathbb{T}}(g,m)$
for some integer $m \ge 0$, and the least such $m$ exists.

Divergence. Let $u_0$ be the least integer $\ge 0$ with $b_{+} + 4c_2 + 3\lambda u_0 \ge 0$. It
exists because $\lambda > 0$, and it is a number fixed by $(L, \mathfrak{p})$. For every $m \ge 0$
we have $x_{m,u_0} \ge g^{-2}$. Keeping only the term $u = u_0$ in
\eqref{6.3:eq-block-exponent-floor-1} and using that $R_J$ is non-decreasing, together with its
definition, gives
\begin{equation*}
N_{\mathbb{T}}(g,m) \ge C_{\mathbb{T}} L^{-u_0} R_J(g^{-2})
\ge C_{\mathbb{T}} L^{-u_0} (\log g^{-2})^{r}.
\end{equation*}
Taking $m = m^{J}_{\mathbb{T}}(g)$ and using its defining property,
\begin{equation*}
L^{m^{J}_{\mathbb{T}}(g)} \ge N_{\mathbb{T}}\bigl(g, m^{J}_{\mathbb{T}}(g)\bigr)
\ge C_{\mathbb{T}} L^{-u_0} (\log g^{-2})^{r},
\end{equation*}
and the right member tends to $\infty$ as $g \downarrow 0$. Hence
$m^{J}_{\mathbb{T}}(g) \to \infty$.

Likewise, $g_{-}(g) = (g^{-2} + B^{\sharp})^{-1/2}$ decreases to $0$ as $g \downarrow 0$. By the
definition of $R$, we have $R(g_{-}(g)) \ge (\log g_{-}(g)^{-2})^{r}$, which tends to $\infty$.
Since $m_{\mathbb{T}}(g') = 11 + m' + \log_L R(g')$, it follows that
$m_{\mathbb{T}}(g_{-}(g)) \to \infty$ as $g \downarrow 0$.
\end{proof}

\begin{definition}[Connected Schwinger functions as joint cumulants]\label{6.3:not-joint-cumulants}
The parameters $L$, $\mathfrak{p}$ (with $M=L^{m'}$) and the compact group $G$ in which the bond
variables take their values are those of Definition~\ref{6.3:not-standing-parameters};
$\operatorname{tr}$ is normalised by $\operatorname{tr}\mathbb{1}=1$.

The torus exponent is written $m_1$, and
\begin{equation*}
  \mathbb{T}_{m_1}:=\mathbb{R}^4/(2L^{m_1}\mathbb{Z})^4
\end{equation*}
is the continuum torus of side $2L^{m_1}$ in unit-lattice units, with its quotient metric. For a
cutoff $K\ge 0$ the bare lattice has spacing $\varepsilon:=L^{-K}$, hence $2L^{m_1+K}$ sites per
side. Its bonds carry variables $U$ in $G$, and $dU$ is the product Haar measure. For a plaquette $p$
of the bare lattice, $U(\partial p)$ is the holonomy around $p$, $x(p)\in\mathbb{T}_{m_1}$ is its
barycentre, and
\begin{equation*}
  a_p(U):=1-\operatorname{Re}\operatorname{tr}U(\partial p)\in[0,2],
  \qquad A(U):=\sum_p a_p(U).
\end{equation*}
For bounded $f\colon\mathbb{T}_{m_1}\to\mathbb{R}$ put
$\mathcal{O}(f):=\sum_p f(x(p))\,a_p(U)$, the sum running over the bare plaquettes. For Lipschitz $f$
put $\|f\|_{\sharp}:=\max\{\|f\|_\infty,\,40M\operatorname{Lip}f\}$.

Let $g_0>0$ be the bare coupling and $x_0:=g_0^{-2}$. For $n\ge1$ and bounded
$f_1,\dots,f_n\colon\mathbb{T}_{m_1}\to\mathbb{R}$ set
\begin{equation}\label{6.3:eq-joint-cumulants-1}
  Z(z):=\int\exp\Big[-x_0A(U)+\sum_{i=1}^n z_i\,\mathcal{O}(f_i)\Big]\,dU,
  \qquad z\in\mathbb{C}^n .
\end{equation}
Then $Z$ is an entire function of $z$ with $Z(0)>0$, and we define
\begin{equation*}
  S^T_n(f_1,\dots,f_n):=\partial_{z_1}\cdots\partial_{z_n}\log Z(z)\big|_{z=0}.
\end{equation*}
Here the cutoff $K$, the torus exponent $m_1$ and $x_0$ are held fixed and are suppressed from the
notation.
Let $\mu$ be the probability measure $Z(0)^{-1}e^{-x_0A(U)}\,dU$, and let $\langle\cdot\rangle$ denote
its expectation. Then $S^T_n(f_1,\dots,f_n)$ is the joint cumulant of the bounded random variables
$\mathcal{O}(f_1),\dots,\mathcal{O}(f_n)$ under $\mu$. It is also written as the truncated
expectation $\langle\mathcal{O}(f_1);\dots;\mathcal{O}(f_n)\rangle^{\mathrm{conn}}_K$, the subscript
recording the cutoff $K$.

The partition function used in the correlation theorem differs from the function $Z$ of
\eqref{6.3:eq-joint-cumulants-1} by a constant factor $e^{-E}$, with $E$ independent of $z$. This
factor does not affect derivatives of $\log Z$.
\end{definition}

\begin{proof}
The bare lattice has $2L^{m_1+K}$ sites per side, so it has finitely many bonds and finitely many
plaquettes. The integral in \eqref{6.3:eq-joint-cumulants-1} is therefore an integral over a finite
product of copies of the compact group $G$.

Since $0\le a_p(U)\le 2$, we have $A(U)\ge0$ and
\begin{equation*}
  |\mathcal{O}(f)|\le 2\,\|f\|_\infty\cdot\#\{p\}
\end{equation*}
for every $U$, where $\#\{p\}$ is the number of bare plaquettes. Hence each $\mathcal{O}(f_i)$ is a
bounded random variable. Moreover, since $x_0>0$ and $A(U)\ge0$, the modulus of the integrand is at
most $\exp|\sum_i z_i\,\mathcal{O}(f_i)|$, and
\begin{equation*}
  \Big|\sum_{i=1}^n z_i\,\mathcal{O}(f_i)\Big|
  \le 2\,\#\{p\}\,\max_i|z_i|\sum_{i=1}^n\|f_i\|_\infty .
\end{equation*}
Hence on every bounded set of $z\in\mathbb{C}^n$ the integrand is bounded uniformly in $U$.
For each fixed $U$ the integrand is entire in $z$, being the exponential of a polynomial of degree
one in $z$. It is also continuous in $(U,z)$. A parameter integral over a finite measure of a
function that is holomorphic in the parameter and locally uniformly bounded is holomorphic.
Therefore $Z$ is entire.

At real $z$ the integrand is a strictly positive continuous function on the compact group, and the
Haar measure has positive total mass. Hence $Z(z)>0$ for real $z$, and in particular $Z(0)>0$.
Since $Z$ is continuous and $Z(0)>0$, $Z$ does not vanish on a neighbourhood of $0$. On that
neighbourhood $\log Z$ (the branch that is real on real $z$) is holomorphic, so the derivative
defining $S^T_n$ exists.

Dividing \eqref{6.3:eq-joint-cumulants-1} by $Z(0)$ gives
\begin{equation*}
  \frac{Z(z)}{Z(0)}=\Big\langle\exp\sum_{i=1}^n z_i\,\mathcal{O}(f_i)\Big\rangle .
\end{equation*}
Thus $\log Z(z)-\log Z(0)$ is the joint cumulant generating function of
$\mathcal{O}(f_1),\dots,\mathcal{O}(f_n)$ under $\mu$. The constant $\log Z(0)$ is annihilated by
$\partial_{z_1}\cdots\partial_{z_n}$ for $n\ge1$. Consequently $S^T_n(f_1,\dots,f_n)$ is, by the
definition of the joint cumulant, the joint cumulant of these variables.

Finally, replacing $Z$ by $e^{-E}Z$ adds the constant $-E$ to $\log Z$. This leaves every derivative
$\partial_{z_1}\cdots\partial_{z_n}\log Z$ with $n\ge1$ unchanged.
\end{proof}

\begin{definition}[First-passage cutoff and last run index]\label{6.3:def-first-passage}
Let $x_k := g_k^{-2}$, where $(x_k)_{k \ge 0}$ is the sequence generated by the coupling recursion
of the effective-action theorem, started at $x_0 = g_0^{-2}$. Throughout this section we assume, as a
standing hypothesis, that the coefficient $b_{k+1}(0)$ of the correlation theorem equals the coefficient
$\beta_{k+1}$ of this recursion, so that the two sequences coincide as long as both are defined.
Hence $(x_k)$ is the sequence of the correlation theorem,
\begin{equation}\label{6.3:eq-first-passage-1}
  x_{k+1} = x_k - b_{k+1}(0).
\end{equation}
As a sequence of numbers, which do not depend on the torus, it is defined for every index reached
by the run of the correlation theorem. Put
\begin{equation*}
  X := g^{-2}.
\end{equation*}
The \emph{first-passage cutoff} is the cutoff of the effective-action theorem,
\begin{equation}\label{6.3:eq-first-passage-2}
  K := K(g_0, g) := \min\{\, k \ge 0 : x_k \le X + B^{\sharp} \,\}.
\end{equation}
The \emph{last run index} is the last index of the correlation theorem,
\begin{equation}\label{6.3:eq-first-passage-3}
  K_J := \max\{\, k \ge 0 : x_j > X \text{ for all } j \le k \,\}.
\end{equation}
In the statement of the correlation theorem the index \eqref{6.3:eq-first-passage-3} is written $K$.
Throughout this book it is written $K_J$, so that it is kept apart from the cutoff
\eqref{6.3:eq-first-passage-2}.
\end{definition}

\begin{definition}[The per-ball source classes]\label{6.3:def-source-classes}
Let $\mathbb{T}_{m_1}$ be the torus of side $2L^{m_1}$, where $m_1$ is the torus exponent, and
equip it with its torus metric. For $c \in \mathbb{T}_{m_1}$ and $R > 0$ write $B(c,R)$ for the
closed ball of centre $c$ and radius $R$ in this metric. For subsets $A, A' \subset \mathbb{T}_{m_1}$
write $\operatorname{dist}(A,A')$ for their distance in this metric.

For $n \ge 1$ and $R > 0$, the class $\mathfrak{S}(n,R)$ is the set of $n$-tuples
$(f_1,\dots,f_n)$ of real bounded functions on $\mathbb{T}_{m_1}$ with the following property: for
each $i$ there is a centre $c_i \in \mathbb{T}_{m_1}$ such that
\begin{equation*}
\operatorname{supp} f_i \subset B(c_i,R).
\end{equation*}

For $d > 0$, the class $\mathfrak{S}(n,R,d)$ is the set of those tuples
$(f_1,\dots,f_n) \in \mathfrak{S}(n,R)$ that satisfy two further conditions. First, every $f_i$ is
Lipschitz. Second, the supports are pairwise separated:
\begin{equation*}
\operatorname{dist}(\operatorname{supp} f_i, \operatorname{supp} f_j) \ge d
\qquad \text{for } i \ne j .
\end{equation*}
\end{definition}

\begin{definition}[The two families of data compared]\label{6.3:not-data-families}
In this section two families of data are matched against each other: the data and clause~(0) of
Theorem~0, and the data and the two conclusions of the correlation theorem that enter the
comparison. We record both below, each in the notation in which it is used. The flow constants
$\lambda^{\sharp}$, $B^{\sharp}$, $c_5$, $\lambda$, $b_{+}$, $c_2$ are those of
Definition~\ref{6.3:not-standing-parameters}. The torus floor $m_{\mathbb{T}}$ and the floor function
$N_{\mathbb{T}}(g,m)$ are those of Definition~\ref{6.3:def-block-exponent-floor}. The first-passage
index $K(g_0,g)$ and the last run index $K_J$ are those of Definition~\ref{6.3:def-first-passage}.
Throughout, $X=g^{-2}$.

\begin{enumerate}
\item[(a)] \emph{The family of Theorem~0.} The data are $(L,\mathfrak{p},\gamma,\gamma_H)$ and
a target coupling $g\in(0,\gamma_H]$. The torus exponent $m_1$ is subject to the floor
\begin{equation}\label{6.3:eq-data-families-1}
m_1\ \ge\ m_{\mathbb{T}}\bigl(g_{-}(g)\bigr),
\end{equation}
which is the only condition on the torus. The bare coupling is $g_0\in(0,g_{-}(g)]$. The cutoff is
$K=K(g_0,g)$, the index of first passage at $X+B^{\sharp}$. The bare lattice has $2L^{m_1+K}$
sites per side. The lattice at level $K$ has unit spacing and $2L^{m_1}$ sites per side.
Clause~(0) of Theorem~0 asserts the following:
\begin{itemize}
\item $K$ is well defined;
\item $x_k>X$ for every $k\le K$;
\item $x_K\in[X,X+B^{\sharp}]$, which is the first-passage property, and
$g_K\in[g_{-}(g),g)$;
\item for $0\le i\le j\le K$ the two-sided flow bounds hold:
\begin{equation}\label{6.3:eq-data-families-2}
\lambda^{\sharp}(j-i)-c_5\ \le\ x_i-x_j\ \le\ B^{\sharp}(j-i).
\end{equation}
\end{itemize}
The constants of clauses~(i) and~(ii) of Theorem~0 depend on $(L,\mathfrak{p},\gamma)$ and
on the gauge group. They depend on the level $k$ only through $g_k$, and never on
$K$, $m_1$, $g_0$ or position.

\item[(b)] \emph{The family of the correlation theorem.} The data are as follows.
\begin{itemize}
\item The four-dimensional torus of side one carries a bare lattice of spacing
$L^{-K_J}N_T^{-1}$, that is, with $L^{K_J}N_T$ sites per side. Here $N_T\ge 1$ is a free integer.
\item $m\ge 0$ is an integer independent of $K_J$, the number of omitted last steps.
\item $K':=K_J-m\ge 0$ is the last step of the local run.
\item $N:=N_T L^m$ is the number of sites per side of the last lattice, the lattice at level
$K'$. In this family $N$ counts sites; it is not the $N$ of $SU(N)$.
\item The block exponent obeys the floor $L^m\ge N_{\mathbb{T}}(g,m)$. This is a lower bound on
the block exponent $m$, never an upper bound on $K_J$.
\item The initial inverse squared coupling satisfies $x_0>X$.
\end{itemize}
For a function $f$ on the torus of side one, and with $\lambda_{*}:=(40M)^{-1}$, we set
\begin{equation}\label{6.3:eq-data-families-3}
\|f\|^{J}_{\sharp}:=\max\Bigl\{\|f\|_{\infty},\ \frac{\operatorname{Lip} f}{\lambda_{*}N}\Bigr\},
\qquad
\mathbb{D}:=\Bigl\{\,z:\ \sum_i |z_i|\,\|f_i\|^{J}_{\sharp}<\varrho\,\Bigr\}.
\end{equation}
Two conclusions of the correlation theorem enter the comparison. Here $(x_i)$ denotes the sequence
of inverse squared couplings of that theorem's run. The first is a pair of flow bounds: for
$0\le i<j\le K_J$,
\begin{equation}\label{6.3:eq-data-families-4}
\lambda(j-i)-2c_2\ \le\ x_i-x_j\ \le\ 3\lambda(j-i)+2c_2,
\end{equation}
and moreover $x_{K_J}\in(X,X+b_{+}]$. The second is the torus bound. Let $n\ge 2$, and let the
supports of $f_1,\dots,f_n$ lie pairwise at distance at least $d>0$. Then
\begin{equation}\label{6.3:eq-data-families-5}
\bigl|\partial_{z_1}\cdots\partial_{z_n}\log Z(0)\bigr|
\ \le\ \Bigl[C_{\mathrm{UV}}(d)+\log 2\cdot(2\mathsf{Q}+2)^{n}\Bigr]
\Bigl(\frac{n}{\varrho}\Bigr)^{n}\prod_{i=1}^{n}\|f_i\|^{J}_{\sharp},
\end{equation}
with
\begin{equation}\label{6.3:eq-data-families-6}
\mathsf{Q}=e^{(E_{+}^{\sharp}+C_{\mathrm{low}})N^{4}} .
\end{equation}
In~\eqref{6.3:eq-data-families-5} and~\eqref{6.3:eq-data-families-6}, $C_{\mathrm{UV}}(d)$ is the
correlation theorem's constant depending on the separation $d$. The constants $E_{+}^{\sharp}$ and
$C_{\mathrm{low}}$ are its last-scale per-site constants.
\end{enumerate}
\end{definition}

\begin{remark}[What the correlation theorem's constants depend on]\label{6.3:rem-constant-dependence}
We record which objects of the correlation theorem depend on which of its data, as its statement
asserts. Throughout, $m$ denotes the correlation theorem's number of omitted last steps, and the
torus exponent of this chapter is written $m_1$. Let $x_{K'} = g_{K'}^{-2}$ be the inverse coupling
at the last step $K'$ of the local run, $\mathsf{T}_{K'} = \log x_{K'}$, $X = g^{-2}$, and
$\lambda = c_0/2$, $b_{+} = 3\lambda + 2c_2$.
\begin{enumerate}
\item Its sentence on constants asserts that the indices $K$, $k$, $m$, the torus and the positions
occur in no constant. The number $N$ of sites per side of the last lattice occurs only in the norm
$\|\cdot\|^{J}_{\sharp}$, in $\mathsf{Q}$, and in the torus floor
condition of the correlation theorem, which is a condition on $m$. The objects produced by the
recursion are numbers free of the torus.
\item Its last-scale clause names two exceptions to the independence from $m$, and only these two.
Both enter through $x_{K'}$,
and both are bounded at once by the coupling bound of the correlation theorem,
\begin{equation}\label{6.3:eq-constant-dependence-1}
x_{K'} \le X + b_{+} + 3\lambda m + 2c_2 .
\end{equation}
First, the top zone, of index $j = K'$, contributes to $E_{+}^{\sharp}$ the member
$2^{-d}(C_{\sharp}O(1)+1)\,\mathsf{T}_{K'}$. This member reads $(g,m)$, in the same way as
$C_{\mathrm{low}}$ does, through
\begin{equation*}
\mathsf{T}_{K'} = \log g_{K'}^{-2} \le \log(X + b_{+} + 3\lambda m + 2c_2),
\end{equation*}
and the bound $|\tilde{Z}| \le e^{E_{+}^{\sharp}N^4}$ holds with $E_{+}^{\sharp}$ changed by this
one member. Here $\tilde{Z}$ denotes the sourced integral at the last step of the correlation
theorem, whose value at zero source is $\int \rho^{0}_{K'}$. Second, the floor reads
\begin{equation}\label{6.3:eq-constant-dependence-2}
\tilde{Z}(0) = \int \rho^{0}_{K'} \ge e^{-C_{\mathrm{low}}N^4},
\end{equation}
with $C_{\mathrm{low}}$ read at the bound \eqref{6.3:eq-constant-dependence-1} for $x_{K'}$.
Here $\rho^{0}_{K'}$ is the integrand of $\tilde{Z}(0)$ at step $K'$ in the correlation theorem.
Every other member of $E_{+}^{\sharp}$ is one of Ba{\l}aban's constants, changed by numbers that
depend on $(d,L)$ only.
\item The limit clause of the correlation theorem is stated along every sequence of cutoff indices
$K_i \to \infty$ at fixed $(g, m, N_T)$ obeying the torus floor condition.
\item The torus $\mathbb{T}_{m_1}$ of this chapter enters the correlation theorem only through
$N_T$.
\end{enumerate}
\end{remark}

\begin{proof}
The first three items are read directly from the statement of the correlation theorem. The first
item is its sentence on constants. The two exceptions of the second item are the two sentences of
that statement that name a member reading $(g,m)$. Each bounds that member through the inverse
coupling $x_{K'}$, and
\eqref{6.3:eq-constant-dependence-1} is the coupling bound that the correlation theorem supplies
for $x_{K'}$. The statement also says that the remaining members of $E_{+}^{\sharp}$ differ from
Ba{\l}aban's constants by $(d,L)$-numbers only. The third item is its limit clause.

The last item is the comparison of data in Notation~\ref{6.3:not-data-families}, under which the
torus $\mathbb{T}_{m_1}$ enters the correlation theorem only through $N_T$.
\end{proof}

\begin{definition}[One-sided floors on coupling and torus exponent]\label{6.3:def-one-sided-floors}
Let $L$, $\mathfrak{p}$, $\gamma$, $\varrho$ be the standing parameters and $\lambda$, $\lambda^{\sharp}$, $B^{\sharp}$,
$c_5$ the flow constants of Notation~\ref{6.3:not-standing-parameters}. Let $m_{\mathbb{T}}$ be the torus floor
and $m^{J}_{\mathbb{T}}$ the least admissible block exponent of
Definition~\ref{6.3:def-block-exponent-floor}, and let $g_{-}(g)=(g^{-2}+B^{\sharp})^{-1/2}$. Put
\begin{equation}\label{6.3:eq-one-sided-floors-1}
\mathcal{A}:=\{0\}\cup\{a\in\mathbb{Z}_{\ge 1} : \lambda a< B^{\sharp}+c_5\},
\qquad a^{\natural}:=\max\mathcal{A}.
\end{equation}
For $n\ge 2$ and $R>0$ let $\ell(n,R)$ be the least power of $L$ that is $\ge 8n(R+1)$, with $R$
measured in the units in which the torus $\mathbb{T}_{m_1}$ has side $2L^{m_1}$. Let $\mathcal{B}(g)$ be
the finite set of exponents $b=m^{J}_{\mathbb{T}}(g)-a$, $a\in\mathcal{A}$.

\smallskip
\noindent\textit{Floor on the coupling.} This is the following condition on $g$:
\begin{equation}\label{6.3:eq-one-sided-floors-a}
m_{\mathbb{T}}(g_{-}(g))\ge a^{\natural}
\quad\text{and}\quad
m^{J}_{\mathbb{T}}(g)\ge a^{\natural}.
\end{equation}

\noindent\textit{Floor on the torus exponent.} This is the following condition on $m_1$:
\begin{equation}\label{6.3:eq-one-sided-floors-b}
m_1\ge m_{\mathbb{T}}(g_{-}(g)).
\end{equation}

\noindent\textit{Floor for the per-ball bounds.} For part (c) of Lemma~\ref{6.3:lem-last-scale-constants}
and for its corollary only, $m_1$ is required in addition to satisfy
\begin{equation}\label{6.3:eq-one-sided-floors-c}
L^{m_1}\ge L^{m^{J}_{\mathbb{T}}(g)}\,\ell(n,R).
\end{equation}

\noindent These floors have the following properties.
\begin{enumerate}
\item[(1)] The number $a^{\natural}$ is fixed with $(L,\mathfrak{p},\gamma)$, before $g$. There is
$\gamma_D=\gamma_D(L,\mathfrak{p},\gamma)>0$ such that every $g\in\,]0,\gamma_D]$ satisfies
\eqref{6.3:eq-one-sided-floors-a}. Hence the hypothesis that $g\le\gamma_J(L,\mathfrak{p},\varrho)$ and
that \eqref{6.3:eq-one-sided-floors-a} holds is implied by
$g\le\min\{\gamma_J(L,\mathfrak{p},\varrho),\gamma_D(L,\mathfrak{p},\gamma)\}$, an upper ceiling on the
coupling of the same kind as $\gamma_J$, that is, a floor on $g^{-2}$. Further,
\eqref{6.3:eq-one-sided-floors-a} is implied by the two inequalities
$m_{\mathbb{T}}(g_{-}(g))\ge (B^{\sharp}+c_5)/\lambda^{\sharp}$ and
$m^{J}_{\mathbb{T}}(g)\ge (B^{\sharp}+c_5)/\lambda^{\sharp}$.
\item[(2)] The floors \eqref{6.3:eq-one-sided-floors-a} and \eqref{6.3:eq-one-sided-floors-b} together give
$m_1-a\ge 0$ for every $a\in\mathcal{A}$.
\item[(3)] Condition \eqref{6.3:eq-one-sided-floors-c} is equivalent to
$m_1\ge m^{J}_{\mathbb{T}}(g)+\log_L\ell(n,R)$; it is a one-sided floor on $m_1$ depending on
$(L,\mathfrak{p},g;n,R)$, and it implies the floor $m_1\ge\log_L\ell(n,RL^{b})$ of the per-ball
volume-free bound (Theorem~\ref{6.4:thm-per-ball-bound}) for every $b\in\mathcal{B}(g)$.
\item[(4)] The correlation theorem carries a floor on its own block exponent, which we denote $m_J$;
this floor is met by the choice $m_J=m^{J}_{\mathbb{T}}(g)$ and imposes nothing on the data
$(g,m_1,g_0)$. In particular it is a condition distinct from \eqref{6.3:eq-one-sided-floors-b}, which
bears on $m_1$.
\end{enumerate}
No floor on $L$ beyond $L\ge L_0$ is introduced.
\end{definition}

\begin{proof}
(1) The set $\mathcal{A}$ is defined from $(\lambda,B^{\sharp},c_5)$, which are functions of
$(L,\mathfrak{p},\gamma)$; it is finite since $\lambda>0$, so $a^{\natural}$ exists and is fixed with
$(L,\mathfrak{p},\gamma)$. By Definition~\ref{6.3:def-block-exponent-floor},
$m_{\mathbb{T}}(g_{-}(g))\to\infty$ and $m^{J}_{\mathbb{T}}(g)\to\infty$ as $g\downarrow 0$. Hence there is
$\gamma_D=\gamma_D(L,\mathfrak{p},\gamma)>0$ such that both quantities are $\ge a^{\natural}$ for every
$g\in\,]0,\gamma_D]$, which is \eqref{6.3:eq-one-sided-floors-a}. If
$g\le\min\{\gamma_J(L,\mathfrak{p},\varrho),\gamma_D(L,\mathfrak{p},\gamma)\}$, then $g\le\gamma_J$, and
$g\in\,]0,\gamma_D]$, so \eqref{6.3:eq-one-sided-floors-a} holds.

For the last assertion, let $a\in\mathcal{A}$ with $a\ge 1$. Then $\lambda a<B^{\sharp}+c_5$, so
$B^{\sharp}+c_5>\lambda>0$. Since $\lambda\ge\lambda^{\sharp}>0$ (Notation~\ref{6.3:not-standing-parameters}),
\begin{equation*}
a<\frac{B^{\sharp}+c_5}{\lambda}\le\frac{B^{\sharp}+c_5}{\lambda^{\sharp}}.
\end{equation*}
Hence every $a\in\mathcal{A}$ satisfies $a\le\max\{0,(B^{\sharp}+c_5)/\lambda^{\sharp}\}$, and so does
$a^{\natural}$. Since $m_{\mathbb{T}}(g_{-}(g))$ and $m^{J}_{\mathbb{T}}(g)$ are non-negative by
Definition~\ref{6.3:def-block-exponent-floor}, the two
inequalities with $(B^{\sharp}+c_5)/\lambda^{\sharp}$ imply that both are
$\ge\max\{0,(B^{\sharp}+c_5)/\lambda^{\sharp}\}\ge a^{\natural}$, which is
\eqref{6.3:eq-one-sided-floors-a}.

(2) For $a\in\mathcal{A}$, the first inequality of \eqref{6.3:eq-one-sided-floors-a} and
\eqref{6.3:eq-one-sided-floors-b} give
\begin{equation*}
a\le a^{\natural}\le m_{\mathbb{T}}(g_{-}(g))\le m_1 ,
\end{equation*}
so $m_1-a\ge 0$.

(3) Write $\ell(n,R)=L^{j}$ with $j$ an integer. Since $L>1$, \eqref{6.3:eq-one-sided-floors-c} holds
if and only if $m_1\ge m^{J}_{\mathbb{T}}(g)+j$, and $j=\log_L\ell(n,R)$. This condition bounds $m_1$ from below
only, and its right member depends on $(L,\mathfrak{p},g)$ through $m^{J}_{\mathbb{T}}(g)$, and on $L$ and
$(n,R)$ through $\ell(n,R)$.

Let $b\in\mathcal{B}(g)$, so $b=m^{J}_{\mathbb{T}}(g)-a$ with $a\in\mathcal{A}$. Since
$a\ge 0$, we have $b\le m^{J}_{\mathbb{T}}(g)$, so $L^{m^{J}_{\mathbb{T}}(g)}\ge L^{b}$. Since
$m^{J}_{\mathbb{T}}(g)\ge 0$ by Definition~\ref{6.3:def-block-exponent-floor}, we
have $L^{m^{J}_{\mathbb{T}}(g)}\ge 1$. Moreover $\ell(n,R)\ge 8n(R+1)$ and $R>0$. Hence
\begin{equation*}
\begin{split}
L^{m_1}&\ge L^{m^{J}_{\mathbb{T}}(g)}\,\ell(n,R)\ge 8n\,L^{m^{J}_{\mathbb{T}}(g)}(R+1)\\
&=8n\bigl(RL^{m^{J}_{\mathbb{T}}(g)}+L^{m^{J}_{\mathbb{T}}(g)}\bigr)\ge 8n\bigl(RL^{b}+1\bigr).
\end{split}
\end{equation*}
Thus $L^{m_1}$ is a power of $L$ that is $\ge 8n(RL^{b}+1)$. Since $\ell(n,RL^{b})$ is the least such
power, $L^{m_1}\ge\ell(n,RL^{b})$, that is, $m_1\ge\log_L\ell(n,RL^{b})$.

(4) The exponent $m_J=m^{J}_{\mathbb{T}}(g)$ is the least admissible block exponent of
Definition~\ref{6.3:def-block-exponent-floor}, which exists for every $g$ and depends on
$(L,\mathfrak{p})$ and $g$ only. The floor of the correlation theorem on $m_J$ is met by this choice,
since by Definition~\ref{6.3:def-block-exponent-floor} $m^{J}_{\mathbb{T}}(g)$ is the least integer
$m\ge 0$ with $L^{m}\ge N_{\mathbb{T}}(g,m)$, which is that floor. So it places no condition on $m_1$ or
on $g_0$, nor any condition on $g$ beyond the ceilings
already imposed. By contrast, \eqref{6.3:eq-one-sided-floors-b} is a lower bound on $m_1$. The two
conditions are therefore distinct.

No condition on $L$ appears in (1)--(4) beyond $L\ge L_0$, which is part of the standing parameters.
\end{proof}

\begin{remark}\label{6.3:lem-last-scale-constants}
The statement of this lemma is not written out here; parts of its proof are printed below.
\end{remark}

\begin{proof}[Proof of Lemma~\ref{6.3:lem-last-scale-constants}, part~(b)]\label{6.3:prf-last-scale-coupling}
Let $L$, $\mathfrak{p}$, $\gamma$, $\gamma_H$, $\varrho$, $g$, $m_1$, $g_0$, $K=K(g_0,g)$ and $K_J$ be as in
Lemma~\ref{6.3:lem-last-scale-constants}. Put $X:=g^{-2}$ and $x_0:=g_0^{-2}$. Let $a:=K_J-K$,
which lies in $\mathcal{A}$ by part~(a) of Lemma~\ref{6.3:lem-last-scale-constants}, and put
$a^{\natural}:=\max\mathcal{A}$. The argument uses the following antecedents.
\begin{itemize}
\item The least admissible block exponent $m^{J}_{\mathbb{T}}(g)$ of
Definition~\ref{6.3:def-block-exponent-floor}.
\item The floors \eqref{6.3:eq-one-sided-floors-a} and \eqref{6.3:eq-one-sided-floors-b} of
Definition~\ref{6.3:def-one-sided-floors}. The first gives
$m_{\mathbb{T}}(g_{-}(g))\ge a^{\natural}$ and $m^{J}_{\mathbb{T}}(g)\ge a^{\natural}$; the second gives
$m_1\ge m_{\mathbb{T}}(g_{-}(g))$.
\item $x_0>X$; $g\le\gamma_J(L,\mathfrak{p},\varrho)$; $g\le\gamma_H<1$; the gauge group is $SU(2)$.
\item The correlation theorem, in the following form. Let $N_T\ge 1$ be an integer, $m\ge 0$ an
integer not depending on $K_J$, with $K':=K_J-m\ge 0$, $L^{m}\ge N_{\mathbb{T}}(g,m)$, $x_0>g^{-2}$,
$g\le\gamma_J$ and gauge group $SU(2)$. The theory is set on the torus of side $1$ with bare lattice of
spacing $L^{-K_J}N_T^{-1}$, bare density
$\rho_0^{w}=\exp[-A(x_0-w,U)-E]$ with $A(c,U)=\sum_p c(p)\,a_p(U)$, $a_p(U):=1-\operatorname{Re}
\operatorname{tr}U(\partial p)$, where $E$ is a constant and $c(p)$ is the value of the weight $c$ at the
barycentre of $p$, and partition function $Z^{J}(z):=\int\rho_0^{w}$; its lattices $T^{(k)}$,
$0\le k\le K'$, have spacing $L^{k-K_J}N_T^{-1}$, so that the last one, $T^{(K')}$, has
$N:=N_TL^{m}$ sites per side; for $f$ on the side-one torus its norm is
$\|f\|^{J}_{\sharp}:=\max\{\|f\|_{\infty},\,40M\operatorname{Lip}f/N\}$, and its polydisc is
$\mathbb{D}:=\{z:\sum_i|z_i|\,\|f_i\|^{J}_{\sharp}<\varrho\}$. Its Schwinger functions are
$S^{K_J}_{n}=\partial_{z_1}\cdots\partial_{z_n}\log Z^{J}(0)$. Its inverse couplings $x_k$ obey the
reference bounds
\begin{equation}\label{6.3:eq-last-scale-coupling-1}
\lambda(j-i)-2c_2\le x_i-x_j\le 3\lambda(j-i)+2c_2\quad(0\le i<j\le K_J),\qquad
x_{K_J}\in\,]X,X+b_{+}],
\end{equation}
and its constants obey the following clause: the constants $E_{+}^{\sharp}$ and $C_{\mathrm{low}}$ of
the torus bound satisfy
\begin{equation}\label{6.3:eq-last-scale-coupling-a}
E_{+}^{\sharp}+C_{\mathrm{low}}\le E_{J}(L,\mathfrak{p},\varrho;g,m)
\end{equation}
for a function $E_J$ of the displayed arguments only: they do not depend on $K_J$, on $x_0$, on $N_T$,
on $N$ otherwise than through the factor $N^{4}$ in the exponent of
$\mathsf{Q}=e^{(E_{+}^{\sharp}+C_{\mathrm{low}})N^{4}}$ of the torus bound, or on positions; and they
depend on the pair $(g,m)$ only through the member $2^{-d}(C_{\sharp}O(1)+1)\mathsf{T}_{K'}$, with
$\mathsf{T}_{K'}=\log x_{K'}$, of $E_{+}^{\sharp}$ contributed by the last zone, and through
$C_{\mathrm{low}}$, each of which is bounded in terms of the bound $X+b_{+}+3\lambda m+2c_2$ for $x_{K'}$
furnished by \eqref{6.3:eq-last-scale-coupling-1} and its logarithm.
\end{itemize}
Assume $K_J\ge m^{J}_{\mathbb{T}}(g)$, and put $m_J:=m^{J}_{\mathbb{T}}(g)$, $K':=K_J-m_J$,
$N_T:=2L^{m_1-a}$ and $b:=m_J-a$. We show that $b\ge 0$, that the correlation theorem applies with
data $(g,m_J,N_T,x_0)$, that its bare lattice is that of Lemma~\ref{6.3:lem-last-scale-constants} under
the dilation $y\mapsto y/(2L^{m_1})$, that its last lattice $T^{(K')}$ is the block lattice at depth $b$
below the cutoff, with $N=2L^{m_1+b}$ sites per side, that the domain
$\{\sum_i|z_i|\,\|f_i\|_{\sharp}<\varrho\}$ is contained in $\mathbb{D}$, that
$S^{T}_{n;K,m_1}=S^{K_J}_{n}$, and that \eqref{6.3:eq-last-scale-coupling-b} and
\eqref{6.3:eq-last-scale-coupling-c} below hold, where neither $x^{*}(g)$ nor $E_{*}$ contains $K$,
$g_0$, $m_1$, $a$ or positions.

By the assumption $K_J\ge m_J$, $K'=K_J-m_J\ge 0$.

\emph{Admissibility of the data.} By the floor \eqref{6.3:eq-one-sided-floors-a},
$m_{\mathbb{T}}(g_{-}(g))\ge a^{\natural}$ and $m^{J}_{\mathbb{T}}(g)\ge a^{\natural}$. Since
$a\in\mathcal{A}$ we have $a\le a^{\natural}$, and by the floor \eqref{6.3:eq-one-sided-floors-b}
$m_1\ge m_{\mathbb{T}}(g_{-}(g))$; hence
\begin{equation*}
m_1-a\ge m_{\mathbb{T}}(g_{-}(g))-a^{\natural}\ge 0,\qquad b=m_J-a\ge m_J-a^{\natural}\ge 0 .
\end{equation*}
Consequently $N_T=2L^{m_1-a}$ is an integer $\ge 1$, and, since $K_J=K+a$,
\begin{equation}\label{6.3:eq-last-scale-coupling-2}
L^{K_J}N_T=L^{K+a}\cdot 2L^{m_1-a}=2L^{m_1+K}.
\end{equation}
The left member is the number of sites per side of the bare lattice of the correlation theorem
(spacing $L^{-K_J}N_T^{-1}$ on the side-one torus); the right member is the number of sites per side
of the bare lattice of spacing $L^{-K}$ on the torus $\mathbb{T}_{m_1}$ of side $2L^{m_1}$. The
dilation $y\mapsto y/(2L^{m_1})$ maps $\mathbb{T}_{m_1}$ onto the side-one torus and, by
\eqref{6.3:eq-last-scale-coupling-2}, maps the spacing $L^{-K}$ to $L^{-K_J}N_T^{-1}$; it therefore
identifies the two bare lattices, their plaquettes and their barycentres.

The integer $m=m_J$ does not depend on $K$: it is a function of $g$. It satisfies the floor
$L^{m_J}\ge N_{\mathbb{T}}(g,m_J)$ by the definition of $m^{J}_{\mathbb{T}}(g)$ in
Definition~\ref{6.3:def-block-exponent-floor}. Further $K'\ge 0$ by the assumption $K_J\ge m_J$,
$x_0>X$, $g\le\gamma_J$, and the gauge group is $SU(2)$. So the correlation theorem applies with data
$(g,m_J,N_T,x_0)$.

Its last lattice $T^{(K')}$ has
\begin{equation*}
N=N_TL^{m_J}=2L^{m_1-a+m_J}=2L^{m_1+b}
\end{equation*}
sites per side. Under the dilation, the spacing $L^{k-K_J}N_T^{-1}$ of $T^{(k)}$ becomes
$2L^{m_1}\cdot L^{k-K-a}/(2L^{m_1-a})=L^{k-K}$ in the units of $\mathbb{T}_{m_1}$, in which the bare
lattice has spacing $L^{-K}$ and the level-$K$ block lattice has spacing one. Since
$K'=K+a-m_J=K-b$, the lattice
$T^{(K')}=T^{(K-b)}$ has spacing $L^{-b}$: it is the block lattice at level $K-b$, at depth $b$ below
the cutoff, and lengths measured in last-lattice units (spacing $1$ on $T^{(K')}$) are the lengths on
$\mathbb{T}_{m_1}$ multiplied by $L^{b}$.

\emph{Norms.} For $f$ on $\mathbb{T}_{m_1}$ put $\tilde f(y):=f(2L^{m_1}y)$ on the side-one torus.
Then $\|\tilde f\|_{\infty}=\|f\|_{\infty}$ and $\operatorname{Lip}\tilde f=2L^{m_1}\operatorname{Lip}f$,
so, using $N=2L^{m_1+b}$,
\begin{equation*}
\|\tilde f\|^{J}_{\sharp}
=\max\Bigl\{\|f\|_{\infty},\,\frac{40M\cdot 2L^{m_1}\operatorname{Lip}f}{N}\Bigr\}
=\max\{\|f\|_{\infty},\,40ML^{-b}\operatorname{Lip}f\}\le\|f\|_{\sharp},
\end{equation*}
the last inequality because $b\ge 0$ gives $L^{-b}\le 1$ and
$\|f\|_{\sharp}=\max\{\|f\|_{\infty},40M\operatorname{Lip}f\}$. Consequently, if
$\sum_i|z_i|\,\|f_i\|_{\sharp}<\varrho$ then $\sum_i|z_i|\,\|\tilde f_i\|^{J}_{\sharp}<\varrho$, that is,
the domain $\{\sum_i|z_i|\,\|f_i\|_{\sharp}<\varrho\}$ is contained in $\mathbb{D}$; and any bound by
$\prod_i\|\tilde f_i\|^{J}_{\sharp}$ is a bound by $\prod_i\|f_i\|_{\sharp}$.

\emph{Partition functions.} Take $w=\sum_iz_i\tilde f_i$. Since $A(c,U)=\sum_pc(p)a_p(U)$ is linear in
the weight $c$, since $\tilde f_i$ at the barycentre of the image of $p$ equals $f_i(x(p))$, and since
$\sum_pa_p(U)=A(U)$ and $\sum_pf_i(x(p))a_p(U)=\mathcal{O}(f_i)$,
\begin{equation*}
A(x_0-w,U)=x_0A(U)-\sum_iz_i\,\mathcal{O}(f_i).
\end{equation*}
Hence
\begin{equation*}
\rho_0^{w}=e^{-E}\exp\Bigl[-x_0A(U)+\sum_i z_i\,\mathcal{O}(f_i)\Bigr],
\end{equation*}
the integrand of the partition
function $Z(z)$ of Lemma~\ref{6.3:lem-last-scale-constants} with bare inverse coupling $x_0=g_0^{-2}$
multiplied by $e^{-E}$, so $Z(z)=e^{E}Z^{J}(z)$. The factor $e^{E}$ does not depend on $z$, so
$\log Z$ and $\log Z^{J}$ differ by a constant, and their derivatives
$\partial_{z_1}\cdots\partial_{z_n}$ at $0$ coincide; by Notation~\ref{6.3:not-joint-cumulants} the
former is $S^{T}_{n;K,m_1}$, and by the correlation theorem the latter is $S^{K_J}_{n}$. Thus
$S^{T}_{n;K,m_1}=S^{K_J}_{n}$.

\emph{The last-scale coupling.} Since $K'\le K_J$, $x_{K'}>X$. If $m_J=0$, then $K'=K_J$ and
$x_{K'}=x_{K_J}\le X+b_{+}$ by the second assertion of \eqref{6.3:eq-last-scale-coupling-1}. If
$m_J\ge 1$, the upper member of \eqref{6.3:eq-last-scale-coupling-1} with $i=K'<j=K_J$, for which
$j-i=m_J$, together with $x_{K_J}\le X+b_{+}$, gives
\begin{equation*}
x_{K'}\le x_{K_J}+3\lambda m_J+2c_2\le X+b_{+}+3\lambda m_J+2c_2 .
\end{equation*}
In both cases
\begin{equation}\label{6.3:eq-last-scale-coupling-b}
X<x_{K'}\le x^{*}(g):=X+b_{+}+\max\{0,\,3\lambda\,m^{J}_{\mathbb{T}}(g)+2c_2\}.
\end{equation}
Since $g\le\gamma_H<1$ we have $X=g^{-2}>1$, hence $x_{K'}>1$, and the logarithm being increasing,
\begin{equation*}
\mathsf{T}_{K'}:=\log x_{K'}\le\log x^{*}(g).
\end{equation*}
The right member of \eqref{6.3:eq-last-scale-coupling-b} is a function of $(L,\mathfrak{p},g)$ alone;
in particular $m^{J}_{\mathbb{T}}(g)$ is a function of $g$.

\emph{The constants.} The clause of the correlation theorem used now is
\eqref{6.3:eq-last-scale-coupling-a}, and nothing else. Applying it with
$m=m_J=m^{J}_{\mathbb{T}}(g)$ gives
\begin{equation}\label{6.3:eq-last-scale-coupling-c}
E_{+}^{\sharp}+C_{\mathrm{low}}\le E_{*}(L,\mathfrak{p},g,\varrho)
:=E_{J}(L,\mathfrak{p},\varrho;g,m^{J}_{\mathbb{T}}(g)).
\end{equation}
The two members that read
the pair $(g,m)$ are, explicitly, at most
\begin{equation*}
2^{-4}(C_{\sharp}O(1)+1)\cdot\log\bigl(X+b_{+}+3\lambda m^{J}_{\mathbb{T}}(g)+2c_2\bigr)
\end{equation*}
and $C_{\mathrm{low}}$ read at $X+b_{+}+3\lambda m^{J}_{\mathbb{T}}(g)+2c_2$; these are expressions in
$(L,\mathfrak{p},g)$ alone, consistent with \eqref{6.3:eq-last-scale-coupling-b}. The right member of
\eqref{6.3:eq-last-scale-coupling-c} contains neither $K$, nor $g_0$, nor $m_1$ (which entered the
correlation theorem only through $N_T$), nor $a$ (which entered only through $N_T$ and $K_J$), nor
positions.
\end{proof}

\begin{remark}[The statements relied on for block-exponent dependence]
\label{6.3:rem-exponent-statements}
The assertion of part~(b) of the lemma on uniformity of the last-scale constants in the torus
exponent, that the places listed there are the only ones through which the block exponent $m$
of the correlation theorem enters the constants $E_{+}^{\sharp}$ and $C_{\mathrm{low}}$, rests on
exactly three assertions of the correlation theorem, and on no others:
\begin{enumerate}
\item its constants clause;
\item the assertion that the bound of the modulus of the sourced partition function by
$e^{E_{+}^{\sharp}N^{4}}$ on the polydisc $\mathbb{D}$ of sources continues to hold with
$E_{+}^{\sharp}$ changed by a single member, and that this single member of $E_{+}^{\sharp}$ is
the only one depending on the pair $(g,m)$, namely the member
$2^{-d}(C_{\sharp}O(1)+1)\mathsf{T}_{K'}$, with $d=4$ the dimension, contributed by the top zone
$j=K'$,
which depends on $(g,m)$ through $\mathsf{T}_{K'} = \log x_{K'}$;
\item the floor assertion, that the sourced partition function at zero source is bounded below
by $e^{-C_{\mathrm{low}}N^{4}}$ with a constant $C_{\mathrm{low}}$ that depends on the pair
$(g,m)$ only through the last-scale inverse coupling $x_{K'}$, which obeys
$x_{K'} \le X + b_{+} + 3\lambda m + 2c_{2}$.
\end{enumerate}
Read together, these three assertions are exactly what the display
\eqref{6.3:eq-last-scale-coupling-a}, in the statement on admissible data and the last-scale
coupling, isolates from the statement of the correlation theorem. The clause of
\cite[p.~355]{B-CMP122b} stating that $E_{-}$ and $E_{+}$ are independent of $k$, of
the quantity denoted $T_{\eta}$ there and of $U_{k}$ is not used here.
\end{remark}

\begin{proof}
The correlation theorem does not display $E_{+}^{\sharp}$ and $C_{\mathrm{low}}$ in closed form
as functions of $x_{K'}$. What it asserts nevertheless suffices for the bound of part~(b). By
the second and third assertions above, $E_{+}^{\sharp}$ depends on the pair $(g,m)$ only through
the member named in the second assertion, and $C_{\mathrm{low}}$ only through $x_{K'}$; both
dependences therefore pass through $x_{K'}$. By the first
assertion, no other dependence on $m$ enters. Finally, under the identification of the
correlation theorem's data with those of this chapter, the block exponent $m$ of the correlation
theorem is a function of $g$, namely $m = m^{J}_{\mathbb{T}}(g)$. A dependence on the pair $(g,m)$ is
therefore, at fixed $L$, $\mathfrak{p}$ and $\varrho$, a dependence on $g$.
\end{proof}

\begin{proof}[Proof of part~(c) of Lemma~\ref{6.3:lem-last-scale-constants}: the per-ball bound across unit rescalings]
\label{6.3:prf-unit-rescalings}
We deduce part~(c) from the per-ball volume-free bound, parts (i)--(iii), taken by statement
together with its tracing of constants and its change to last-lattice units, which we read
as follows. On the last lattice
$T^{(K')}$ of the correlation theorem's run, for sources supported in balls of radius $R$
(in last-lattice units) whose supports are pairwise at distance at least $d$, part (iii)
holds with
\begin{equation*}
C^{\infty}_{n}(d;R)=\Phi_{n,d,R}(E_{+}^{\sharp}+C_{\mathrm{low}}),
\end{equation*}
where $\Phi$ is the function \eqref{6.3:eq-unit-rescalings-a} below. Its coefficients
$C_{\mathrm{UV}}(d;R)$, $C_{\mathrm{face}}(R)$, $\mathfrak{r}$, $\ell(n,R)$ are functions of
$(L,\mathfrak{p},g,\varrho;n,d,R)$ that do not depend on $K$, $g_0$, $m_1$, $N$, $a$ or the
positions of the supports, and $C_{\mathrm{UV}}$ is non-increasing in $d$.

The remaining inputs of the lemma enter only through the identifications of Case~A made in
the earlier part of the proof of Lemma~\ref{6.3:lem-last-scale-constants} in which the
admissible data and the last-scale coupling are fixed, through the bound
\eqref{6.3:eq-last-scale-coupling-c} established there, and through the floor
\eqref{6.3:eq-one-sided-floors-c}.

We work in Case~A, with those identifications. Thus $a=K_J-K\in\mathcal{A}$ and
$b=m^{J}_{\mathbb{T}}(g)-a$, so that $0\le b\le m^{J}_{\mathbb{T}}(g)$, the integer $b$ ranging
over the finite set $\mathcal{B}(g)$ displayed below. Fix $n\ge 2$, $R>0$, $d>0$ and
$(f_1,\dots,f_n)\in\mathfrak{S}(n,R,d)$ (Definition~\ref{6.3:def-source-classes}).

On $T^{(K')}$, in last-lattice units, $\operatorname{supp} f_i$ lies in a ball of radius
$RL^{b}$, and the supports are pairwise at distance at least $dL^{b}$. For a function $f$
write $\operatorname{Lip}_{\mathrm{ll}}f:=L^{-b}\operatorname{Lip}f$ for its Lipschitz
constant in last-lattice units. The norm in which the per-ball volume-free bound is stated
on the last lattice is $\max\{\|f\|_{\infty},\,40M\operatorname{Lip}_{\mathrm{ll}}f\}$. The
sup norm does not change under the change of units. Since $b\ge 0$, we have
$\operatorname{Lip}_{\mathrm{ll}}f\le\operatorname{Lip}f$. Hence this norm is at most
$\|f\|_{\sharp}$.

We apply the per-ball volume-free bound on $T^{(K')}$ with $(R,d)$ replaced by
$(RL^{b},dL^{b})$. Read on the torus exponent $m_1$ of the lemma, that theorem requires the
floor
\begin{equation*}
m_1\ge\max\{m_{\mathbb{T}}(g_{-}(g)),\,\log_L\ell\},
\qquad \ell=\ell(n,RL^{b}),
\end{equation*}
where $\ell(n,R')$ denotes the least power of $L$ that is $\ge 8n(R'+1)$.

We first bound $\ell(n,RL^{b})$. Since $L^{b}\ge 1$,
\begin{equation*}
8n(RL^{b}+1)\le L^{b}\cdot 8n(R+1)\le L^{b}\ell(n,R).
\end{equation*}
Now $L^{b}\ell(n,R)$ is a power of $L$ that is $\ge 8n(RL^{b}+1)$, and $\ell(n,RL^{b})$ is
the least such power. Together with $b\le m^{J}_{\mathbb{T}}(g)$, this gives
\begin{equation*}
\ell(n,RL^{b})\le L^{b}\ell(n,R)\le L^{m^{J}_{\mathbb{T}}(g)}\ell(n,R).
\end{equation*}
The floor \eqref{6.3:eq-one-sided-floors-c} states
$L^{m_1}\ge L^{m^{J}_{\mathbb{T}}(g)}\ell(n,R)$. It therefore gives
$m_1\ge\log_L\ell(n,RL^{b})$. Together with the hypothesis $m_1\ge m_{\mathbb{T}}(g_{-}(g))$
of Lemma~\ref{6.3:lem-last-scale-constants}, this is the floor of the per-ball volume-free
bound, for every $b$ that occurs.

Part (iii) of the per-ball volume-free bound then gives
\begin{equation}\label{6.3:eq-unit-rescalings-1}
|S^{T}_{n}(f_1,\dots,f_n)|
\le\Phi_{n,dL^{b},RL^{b}}(E_{+}^{\sharp}+C_{\mathrm{low}})\prod_{i=1}^{n}\|f_i\|_{\sharp},
\end{equation}
where, for $t\in\mathbb{R}$, $d'>0$ and $R'>0$,
\begin{equation}\label{6.3:eq-unit-rescalings-a}
\begin{split}
\Phi_{n,d',R'}(t):={}&\Bigl[2C_{\mathrm{UV}}(d';R')
+\log 2\cdot\mathfrak{r}^{n}\Bigl(2e^{n(t\,\ell(n,R')^{4}+2C_{\mathrm{face}}(R'))}+2\Bigr)^{n}\Bigr]\\
&\times\Bigl(\frac{n}{\varrho}\Bigr)^{n}.
\end{split}
\end{equation}

Evaluated at $t=E_{+}^{\sharp}+C_{\mathrm{low}}$, the function
\eqref{6.3:eq-unit-rescalings-a} is the constant $C^{\infty}_{n}(d';R')$ of part (iii) of the
per-ball volume-free bound, in which the last-scale factor
\begin{equation*}
\exp\bigl[(E_{+}^{\sharp}+C_{\mathrm{low}})\,\ell(n,R')^{4}+2C_{\mathrm{face}}(R')\bigr]
\end{equation*}
is written with $t$ in place of $E_{+}^{\sharp}+C_{\mathrm{low}}$. By the tracing of constants
recalled above, the coefficients $C_{\mathrm{UV}}(d';R')$, $C_{\mathrm{face}}(R')$,
$\mathfrak{r}$ and $\ell(n,R')$ are functions of $(L,\mathfrak{p},g,\varrho;n,d',R')$. They
do not depend on $K$, $g_0$, $m_1$, $N$, $a$ or positions, and $C_{\mathrm{UV}}$ is
non-increasing in $d'$.

The function $t\mapsto\Phi_{n,d',R'}(t)$ is non-decreasing. Indeed, the exponential is
increasing in $t$, and all coefficients in \eqref{6.3:eq-unit-rescalings-a} are
non-negative. By \eqref{6.3:eq-last-scale-coupling-c}, $E_{+}^{\sharp}+C_{\mathrm{low}}$ is
majorised by $E_{*}(L,\mathfrak{p},g,\varrho)$. Hence
\begin{equation}\label{6.3:eq-unit-rescalings-2}
\Phi_{n,dL^{b},RL^{b}}(E_{+}^{\sharp}+C_{\mathrm{low}})
\le\Phi_{n,dL^{b},RL^{b}}\bigl(E_{*}(L,\mathfrak{p},g,\varrho)\bigr).
\end{equation}

The integer $b=m^{J}_{\mathbb{T}}(g)-a$ ranges over the finite set
\begin{equation*}
\mathcal{B}(g):=\{m^{J}_{\mathbb{T}}(g)-a' : a'\in\mathcal{A}\}\subset[0,m^{J}_{\mathbb{T}}(g)],
\end{equation*}
which is determined by $(L,\mathfrak{p},\gamma,g)$. Define
\begin{equation}\label{6.3:eq-unit-rescalings-b}
C^{*}_{n}(d;R):=\max_{b'\in\mathcal{B}(g)}
\Phi_{n,dL^{b'},RL^{b'}}\bigl(E_{*}(L,\mathfrak{p},g,\varrho)\bigr).
\end{equation}
Combining \eqref{6.3:eq-unit-rescalings-1}, \eqref{6.3:eq-unit-rescalings-2} and
$b\in\mathcal{B}(g)$, we obtain
\begin{equation*}
|S^{T}_{n}(f_1,\dots,f_n)|\le C^{*}_{n}(d;R)\prod_{i=1}^{n}\|f_i\|_{\sharp}.
\end{equation*}

The constant $C^{*}_{n}$ is a function of $(L,\mathfrak{p},g,\varrho;n,d,R)$ alone, for
three reasons. First, $K$, $g_0$ and $m_1$ occur in no ingredient of
\eqref{6.3:eq-unit-rescalings-b}. Second, the dependence on $a$ has been removed by
maximising over $\mathcal{A}$. Third, the positions of the supports do not occur, by the
tracing of constants in the per-ball volume-free bound: translation covariance removes the
placement of the balls.

Monotonicity in $d$ follows from two observations. For each $b'\in\mathcal{B}(g)$, the map
$d\mapsto\Phi_{n,dL^{b'},RL^{b'}}(E_{*})$ is non-increasing, because only $C_{\mathrm{UV}}$
depends on $d$ and it is non-increasing in its first argument. A maximum of finitely many
non-increasing functions is non-increasing. Hence $d\mapsto C^{*}_{n}(d;R)$ is
non-increasing. This proves part~(c).
\end{proof}

\begin{proof}[Proof of part (d): local plaquette count and cumulant bound]
\label{6.3:prf-plaquette-count}
We prove part (d) of Lemma~\ref{6.3:lem-last-scale-constants}. Let $K=K(g_0,g)$ be the
first-passage cutoff of Definition~\ref{6.3:def-first-passage}. Let $K_J$ be the last
index of the run of the correlation theorem, and let $m^{J}_{\mathbb{T}}(g)$ be the least
block exponent $m$ with $L^{m}\ge N_{\mathbb{T}}(g,m)$. Assume
$K_J<m^{J}_{\mathbb{T}}(g)$.

Earlier in the proof of Lemma~\ref{6.3:lem-last-scale-constants}, the inequality
$K\le K_J$ was derived from hypotheses (I1), (I2) and (I3) of that lemma. It is the only
use of these hypotheses in the present part. Since $K$ and $m^{J}_{\mathbb{T}}(g)$ are
integers, it gives $K\le m^{J}_{\mathbb{T}}(g)-1$. For the bare spacing
$\varepsilon=L^{-K}$ it follows that $\varepsilon\ge L^{1-m^{J}_{\mathbb{T}}(g)}$, and
$\varepsilon\le 1$ because $K\ge 0$.

Fix $n\ge 1$, $R>0$ and $(f_1,\dots,f_n)\in\mathfrak{S}(n,R)$, a member of the per-ball
source classes of Definition~\ref{6.3:def-source-classes}. For each $i$ let $c_i$ be the
centre of a ball of radius $R$ in $\mathbb{T}_{m_1}$ that contains
$\operatorname{supp} f_i$.

\emph{Local plaquette count.} A bare plaquette $p$ is determined by two data:
\begin{itemize}
\item a base site $y$, a point of the bare site lattice of spacing $\varepsilon$ on
$\mathbb{T}_{m_1}$;
\item an unordered pair of distinct coordinate directions, of which there are six at each
site.
\end{itemize}
Its barycentre is obtained from $y$ by adding $\varepsilon/2$ in each of the two
directions, so $|x(p)-y|_\infty=\varepsilon/2\le\varepsilon$. Here $|\cdot|_\infty$
denotes the sup-distance on the torus.

If $x(p)\in\operatorname{supp} f_i$, then $|x(p)-c_i|_\infty\le R$, and the triangle
inequality gives $|y-c_i|_\infty\le R+\varepsilon$. This condition on $y$ is imposed
coordinate by coordinate. In each coordinate the bare sites form a one-dimensional lattice
of spacing $\varepsilon$ on a circle of length $2L^{m_1}$. The points of that lattice
within distance $R+\varepsilon$ of a given point number at most
\begin{equation*}
\min\Bigl\{\frac{2L^{m_1}}{\varepsilon},\ \frac{2(R+\varepsilon)}{\varepsilon}+1\Bigr\}
\le\frac{2R}{\varepsilon}+3\le(2R+3)\,\varepsilon^{-1},
\end{equation*}
where the last step uses $\varepsilon\le 1$. Multiplying over the four coordinates and
over the six direction pairs, and using $\varepsilon^{-1}=L^{K}$ with
$K\le m^{J}_{\mathbb{T}}(g)-1$, we obtain
\begin{equation}\label{6.3:eq-plaquette-count-a}
\begin{split}
\#\{p:\ x(p)\in\operatorname{supp} f_i\}
&\le 6(2R+3)^4\varepsilon^{-4}=6(2R+3)^4L^{4K}\\
&\le 6(2R+3)^4L^{4m^{J}_{\mathbb{T}}(g)}=:P_R .
\end{split}
\end{equation}
The number $P_R$ does not depend on $m_1$. The torus entered only through the first
member of the minimum, and that member was discarded in favour of the larger second one.
Nor does $P_R$ depend on the positions of the supports.

Write $a_p(U)=1-\operatorname{Re}\operatorname{tr}U(\partial p)$. Then
$0\le a_p(U)\le 2$, and
$\mathcal{O}(f_i)=\sum_p f_i(x(p))\,a_p(U)$ receives contributions only from the
plaquettes counted in \eqref{6.3:eq-plaquette-count-a}. Hence
\begin{equation}\label{6.3:eq-plaquette-count-b}
|\mathcal{O}(f_i)|\le 2P_R\,\|f_i\|_\infty
\quad\text{everywhere on the configuration space.}
\end{equation}

\emph{Cumulant bound.} Put $Y_i:=\mathcal{O}(f_i)$. By
Notation~\ref{6.3:not-joint-cumulants}, $S^{T}_{n}(f_1,\dots,f_n)$ is the joint cumulant
$\langle Y_1;\dots;Y_n\rangle$ of the bounded random variables $Y_1,\dots,Y_n$ under the
probability measure $\mu=Z(0)^{-1}e^{-g_0^{-2}A(U)}\,dU$. Here $\langle\cdot\rangle$
denotes expectation under $\mu$.

Let $\operatorname{Part}(n)$ denote the set of partitions of $\{1,\dots,n\}$, and write
$|\pi|$ for the number of blocks of a partition $\pi$. The moment--cumulant formula is
Möbius inversion on the lattice of set partitions of $\{1,\dots,n\}$. On a partition $\pi$
the Möbius function of that lattice equals $(-1)^{|\pi|-1}(|\pi|-1)!$, and the formula
reads
\begin{equation}\label{6.3:eq-plaquette-count-1}
\langle Y_1;\dots;Y_n\rangle
=\sum_{\pi\in\operatorname{Part}(n)}(-1)^{|\pi|-1}(|\pi|-1)!
\prod_{B\in\pi}\Bigl\langle\prod_{i\in B}Y_i\Bigr\rangle .
\end{equation}
Since $\mu$ is a probability measure, each moment satisfies
$|\langle\prod_{i\in B}Y_i\rangle|\le\prod_{i\in B}\|Y_i\|_\infty$. The blocks of $\pi$
partition $\{1,\dots,n\}$, so for every $\pi$ the product over $B\in\pi$ of these bounds
equals $\prod_{i=1}^n\|Y_i\|_\infty$. Therefore
\begin{equation*}
|\langle Y_1;\dots;Y_n\rangle|\le c_n\prod_{i=1}^n\|Y_i\|_\infty,
\qquad c_n:=\sum_{\pi\in\operatorname{Part}(n)}(|\pi|-1)! .
\end{equation*}
Inserting \eqref{6.3:eq-plaquette-count-b} and the value of $P_R$ from
\eqref{6.3:eq-plaquette-count-a} gives
\begin{equation*}
\begin{split}
|S^{T}_{n}(f_1,\dots,f_n)|
&\le c_n(2P_R)^n\prod_{i=1}^n\|f_i\|_\infty
=c_n\bigl[12(2R+3)^4\bigr]^n L^{4n\,m^{J}_{\mathbb{T}}(g)}\prod_{i=1}^n\|f_i\|_\infty\\
&=C^{B}_{n}\prod_{i=1}^n\|f_i\|_\infty .
\end{split}
\end{equation*}

The constant $C^{B}_{n}$ depends only on $L$, $g$, $n$ and $R$. It is free of $m_1$,
$\mathsf{d}$, $K$, $g_0$ and the positions of the supports. Finally, for Lipschitz $f_i$
we have $\|f_i\|_\infty\le\|f_i\|_{\sharp}$ by the definition
$\|f\|_{\sharp}=\max\{\|f\|_\infty,40M\operatorname{Lip}f\}$.

The local count \eqref{6.3:eq-plaquette-count-a} is what makes this bound free of $m_1$. A
count over all plaquettes of $\mathbb{T}_{m_1}$, of which there are
$6(2L^{m_1+K})^4$, would depend on $m_1$.
\end{proof}

\begin{proof}[Proof of part (e) of Lemma~\ref{6.3:lem-last-scale-constants}: bare couplings tending to zero]
\label{6.3:prf-bare-coupling-limit}
The inputs of Lemma~\ref{6.3:lem-last-scale-constants} used below are the first, by which the first-passage
cutoff $K(g_0,g)$ tends to infinity as $g_0\downarrow 0$, and the third, the correlation theorem's upper bound on
the last-scale inverse coupling; beside them we use the limit clauses of the correlation theorem.

Let $g_0^{(i)}\downarrow 0$ at fixed $(g,m_1)$, and put $K_i:=K(g_0^{(i)},g)$. By the first input of
Lemma~\ref{6.3:lem-last-scale-constants},
$K(g_0,g)\to\infty$ as $g_0\downarrow 0$; hence $K_i\to\infty$. By part (a), the integer $a_i$ attached there
to $g_0^{(i)}$ lies in
$\mathcal{A}$ for every $i$. Write $K_J^{(i)}$ for the last index of the correlation theorem's run belonging to
$g_0^{(i)}$. Then
\begin{equation*}
K_J^{(i)}=K_i+a_i\ \ge\ K_i\ \longrightarrow\ \infty .
\end{equation*}
Hence $K_J^{(i)}\ge m^{J}_{\mathbb{T}}(g)$ for all but finitely many $i$. For every such $i$ we are in the case
$K_J\ge m^{J}_{\mathbb{T}}(g)$ of part (a).

The set $\mathcal{A}$ is finite, so some $a\in\mathcal{A}$ equals $a_i$ for infinitely many $i$. Pass to that
subsequence. Let $m_J$ denote the block exponent at which the correlation theorem is applied. Along the subsequence
the following quantities are constant:
\begin{equation*}
N_T=2L^{m_1-a},\qquad m_J=m^{J}_{\mathbb{T}}(g),\qquad N=2L^{m_1+b},
\end{equation*}
where $b\in\mathcal{B}(g)$ is the unit-rescaling exponent $b=m^{J}_{\mathbb{T}}(g)-a$. The limit clauses of the
correlation theorem are stated along cutoffs tending to infinity with the coupling $g$, the block exponent and the
integer $N_T$ held fixed, the block exponent being subject to that theorem's torus floor. Those clauses therefore
apply along this subsequence; it is at this step that the third input of
Lemma~\ref{6.3:lem-last-scale-constants}, the correlation theorem's upper bound on the last-scale inverse
coupling, is used.

It remains to treat the bounds. The bounds \eqref{6.3:eq-last-scale-coupling-c} and \eqref{6.3:eq-unit-rescalings-b},
together with those of part (d), hold for every $i$. Their constants do not depend on $a_i$. Consequently they hold
along the whole sequence $(g_0^{(i)})$, and no passage to a subsequence is needed for them. This proves part (e).
\end{proof}

\begin{corollary}[The per-ball volume-free bound in unit-lattice units]
\label{6.3:cor-volume-free-constant}
Let $L$, $\mathfrak{p}$, $\varrho$, $g$ and the torus exponent $m_1$ be fixed as in
Lemma~\ref{6.3:lem-last-scale-constants}, with $m_1$ satisfying in addition the floor
\eqref{6.3:eq-one-sided-floors-c}. For $g_0 \in \,]0, g_{-}(g)]$ put $K = K(g_0,g)$ and
$S^T_n = S^T_{n;K,m_1}$; let $K_J$ be the last index of the run of the correlation theorem and
$a = K_J - K$. Assume:
\begin{itemize}
\item parts (c) and (d) of Lemma~\ref{6.3:lem-last-scale-constants};
\item parts (i)--(iii) of the per-ball volume-free bound, with its tracing of constants and its
reading in units.
\end{itemize}
Then, under the hypotheses of part (c) of Lemma~\ref{6.3:lem-last-scale-constants}, part (iii)
of the per-ball volume-free bound holds in the units in which $\mathbb{T}_{m_1}$ has side
$2L^{m_1}$, with $C^{\infty}_n(\mathsf{d};R)$ replaced by $C^{*}_{n}(\mathsf{d};R)$. Together
with part (d) this gives the following. For every $n \ge 2$, $R > 0$, $\mathsf{d} > 0$, every
$g_0 \in \,]0, g_{-}(g)]$, and all real Lipschitz functions $f_1, \dots, f_n$ on
$\mathbb{T}_{m_1}$, each supported in a ball of radius $R$, with supports pairwise at distance
at least $\mathsf{d}$,
\begin{equation}\label{6.3:eq-volume-free-constant-1}
\bigl|S^T_n(f_1, \dots, f_n)\bigr|
\le \max\bigl\{C^{*}_{n}(\mathsf{d};R),\, C^{B}_{n}(n,R)\bigr\}
\prod_{i=1}^{n} \|f_i\|_{\sharp}.
\end{equation}
The constant in \eqref{6.3:eq-volume-free-constant-1} depends on
$(L,\mathfrak{p},g,\varrho;n,\mathsf{d},R)$ alone. It is free of $K$, $g_0$, $m_1$, $a$, the
positions of the supports, and the number $N$ of sites per side of the last lattice, and it is
non-increasing in $\mathsf{d}$.
\end{corollary}

\begin{proof}
Case B is the case $K_J < m^{J}_{\mathbb{T}}(g)$, and Case A is the complementary case. Every
$g_0 \in \,]0, g_{-}(g)]$ therefore falls in exactly one of the two cases.

\emph{Case A.} Here part (c) of Lemma~\ref{6.3:lem-last-scale-constants} applies under the
hypotheses listed. It gives
\begin{equation*}
|S^T_n(f_1,\dots,f_n)| \le C^{*}_{n}(\mathsf{d};R) \prod_{i=1}^n \|f_i\|_{\sharp}.
\end{equation*}
This bound is part (iii) of the per-ball volume-free bound, transported to the units in which
$\mathbb{T}_{m_1}$ has side $2L^{m_1}$, with the constant $C^{*}_{n}(\mathsf{d};R)$. This proves
the first assertion, and \eqref{6.3:eq-volume-free-constant-1} follows because
$C^{*}_{n}(\mathsf{d};R)$ does not exceed the maximum.

\emph{Case B.} Here $K_J < m^{J}_{\mathbb{T}}(g)$, so part (d) of
Lemma~\ref{6.3:lem-last-scale-constants} applies. Its hypotheses hold for the given functions:
\begin{itemize}
\item $n \ge 2 \ge 1$;
\item each $f_i$ is real and supported in a ball of radius $R$;
\item each $f_i$ is bounded, being Lipschitz on the compact torus.
\end{itemize}
Part (d) therefore gives
$|S^T_n(f_1,\dots,f_n)| \le C^{B}_{n}(n,R) \prod_{i=1}^n \|f_i\|_{\infty}$. By definition,
$\|f_i\|_{\sharp} = \max\{\|f_i\|_{\infty},\, 40M \operatorname{Lip} f_i\}$, so
$\|f_i\|_{\infty} \le \|f_i\|_{\sharp}$ for each $i$. Hence
\begin{equation*}
|S^T_n(f_1,\dots,f_n)|
\le C^{B}_{n}(n,R) \prod_{i=1}^n \|f_i\|_{\sharp}
\le \max\bigl\{C^{*}_{n}(\mathsf{d};R),\, C^{B}_{n}(n,R)\bigr\} \prod_{i=1}^n \|f_i\|_{\sharp}.
\end{equation*}

\emph{The constant.} By part (c), $C^{*}_{n}(\mathsf{d};R)$ is a function of
$(L,\mathfrak{p},g,\varrho;n,\mathsf{d},R)$. It is free of $K$, $g_0$, $m_1$, $a$ and the
positions of the supports, and it is non-increasing in $\mathsf{d}$. By part (d), $C^{B}_{n}$
depends on $(L,g;n,R)$ and is free of $m_1$, $\mathsf{d}$, $K$, $g_0$ and the positions. The
maximum of the two constants therefore depends on $(L,\mathfrak{p},g,\varrho;n,\mathsf{d},R)$
only, and in particular not on $N$ or $a$. It is non-increasing in $\mathsf{d}$, since it is
the maximum of a non-increasing function of $\mathsf{d}$ and of $C^{B}_{n}$, which is free of
$\mathsf{d}$.
\end{proof}

\begin{remark}[Last-scale constants in the per-ball bound]\label{6.3:rem-per-ball-reading}
In the per-ball volume-free bound, $E_{+}^{\sharp}$ and $C_{\mathrm{low}}$ are the correlation theorem's
per-site last-scale constants. They are free of $K$, of $N_T$, of position, and of $N$ apart from the
factor $N^4$ of $\mathsf{Q}$, which in the per-ball bound is replaced by $\ell^4$, where
$\ell=\ell(n,R)$ is the coefficient of that name in the per-ball bound. They depend on the
correlation theorem's pair $(g,m)$ only through the bound
\begin{equation}\label{6.3:eq-per-ball-reading-1}
x_{K'} \le X + b_{+} + 3\lambda m + 2c_2 ,
\end{equation}
namely through the top-zone member $2^{-4}\bigl(C_{\sharp}O(1)+1\bigr)\mathsf{T}_{K'}$ of $E_{+}^{\sharp}$
and through $C_{\mathrm{low}}$. The correlation theorem is applied with $m=m^{J}_{\mathbb{T}}(g)$, a
function of $(L,\mathfrak{p},g)$ alone; hence
\begin{equation}\label{6.3:eq-per-ball-reading-2}
E_{+}^{\sharp} + C_{\mathrm{low}} \le E_{*}(L,\mathfrak{p},g,\varrho)
\end{equation}
uniformly in $K$, $g_0$, the torus exponent $m_1$ and $a=K_J-K$, and $C^{\infty}_{n}(d;R)$ is to be
read as $C^{*}_{n}(d;R)$ of \eqref{6.3:eq-unit-rescalings-b}, the maximum over the finitely many unit
rescalings $b\in\mathcal{B}(g)$. For the finitely many cutoffs with $K_J<m^{J}_{\mathbb{T}}(g)$ the
bound is the local plaquette count.
\end{remark}

\begin{proof}
That $E_{+}^{\sharp}$ and $C_{\mathrm{low}}$ do not depend on $K$, $N_T$ or position, and depend on
$N$ only through the factor $N^4$ of $\mathsf{Q}$, is the content of the constants clause of
the correlation theorem, displayed in \eqref{6.3:eq-last-scale-coupling-a}. The last-scale analysis
in the proof of the correlation theorem names the two members of these constants that depend on
$(g,m)$, the top-zone member $2^{-4}(C_{\sharp}O(1)+1)\mathsf{T}_{K'}$ of $E_{+}^{\sharp}$
and $C_{\mathrm{low}}$, and bounds each through
\eqref{6.3:eq-per-ball-reading-1} only.

The exponent $m^{J}_{\mathbb{T}}(g)$ is the least block exponent $m$ with
$L^{m}\ge N_{\mathbb{T}}(g,m)$; in the first case of the choice of the block exponent the
correlation theorem is applied with this $m$, a function of $(L,\mathfrak{p},g)$ alone. With this
choice, the result on admissible data and the last-scale coupling gives
\eqref{6.3:eq-per-ball-reading-2}, with a right member containing none of $K$, $g_0$, $m_1$, $a$.

The per-ball constant is then evaluated at $E_{*}(L,\mathfrak{p},g,\varrho)$ in each unit rescaling
$b\in\mathcal{B}(g)$, and its maximum over this finite set is, by definition,
$C^{*}_{n}(d;R)$ of \eqref{6.3:eq-unit-rescalings-b}, as established in the result on the
per-ball bound across unit rescalings. The finitely many cutoffs with
$K_J<m^{J}_{\mathbb{T}}(g)$, to which the preceding bound does not apply, are covered by the local
plaquette count and cumulant bound.
\end{proof}

Consider the volume dependence in the argument that combines the chessboard inequality (Lemma~\ref{6.2:lem-chessboard-inequality}) with the splitting of the counterterm into bulk and face parts \eqref{6.2:eq-chessboard-inequality-b}, and which yields the bound under cell covariance (Proposition~\ref{6.2:prop-cell-covariance}). In that argument the covariance-bound hypothesis is used for one purpose only: to identify the real part of the periodised bulk vacuum energy,
$|\mathcal{C}|^{-1}\operatorname{Re}\mathcal{E}^{\mathrm{bulk}}(w^{\mathrm{ch}}_{c,z})$, with the real part
$\operatorname{Re}\mathcal{E}^{\mathrm{bulk}}(w_z)$ of the counterterm of the correlation theorem. The construction below avoids this hypothesis.

The counterterm, however, may be chosen freely. We have
$S^T_n(\vec f)=\partial_{z_1}\cdots\partial_{z_n}\log Z(w_z)|_{z=0}$. Write
$\log Z(w_z)=\mathcal{E}'(z)+\log\widetilde{Z}'(z)$, where $\mathcal{E}'$ is holomorphic on $\mathbb{D}_{\mathfrak{r}}$, $\mathcal{E}'(0)=0$, and $\widetilde{Z}'(z):=e^{-\mathcal{E}'(z)}Z(w_z)$. Any such splitting serves in the argument by which the correlation theorem estimates $S^T_n$, provided two conditions hold:
\begin{itemize}
\item the mixed derivative $\partial_{z_1}\cdots\partial_{z_n}\mathcal{E}'(0)$ is bounded independently of $m$;
\item $|\widetilde{Z}'/\widetilde{Z}'(0)|\le Q$ on $\mathbb{D}_{\mathfrak{r}}$, with $Q$ independent of $m$.
\end{itemize}
We take for $\mathcal{E}'$ the cell average of the periodised bulk energies themselves, together with the face part of the original counterterm. On the cells $gc$ with $\varepsilon(g)=-1$ the values are complex conjugated; this restores holomorphy and does not change real parts. With this choice the second condition holds because the bulk terms cancel identically. The first condition holds because each averaged piece is the bulk vacuum energy of an isometric copy of the source system. The locality of the correlation theorem therefore applies to each piece with the same $(R,d)$, in the form of the ultraviolet bound for separated ball sources (Lemma~\ref{6.2:lem-ultraviolet-bound}).

\begin{theorem}[The per-ball volume-free bound]\label{6.4:thm-per-ball-bound}
Let $n\ge2$, $R>0$ and $d>0$. Let $\ell=\ell(n,R)$ be the least power of $L$ with
$\ell\ge 8n(R+1)$. Let
\begin{equation*}
m\ge\max\{m_{\mathbb{T}}(g_-(g)),\ \log_L\ell\},
\end{equation*}
let $K'$ be arbitrary, and let $g\le\gamma_J(L,\mathfrak{p},r)$, the radius $r$ being fixed first. Let $f_1,\dots,f_n$ be sources as in Notation~\ref{6.1:not-units-sources}, placed inside the cells of side $\ell$ by Lemma~\ref{6.2:lem-support-placement}. Let
$\mathcal{E}^{\mathrm{bulk}}_{c'}$ and $\mathcal{E}^{\partial}$ be as in \eqref{6.2:eq-chessboard-inequality-b}. Let
$u_{c,g,z}=\sum_{i\in I_c}z_i\,(f_i\circ g^{-1})$ be as in \eqref{6.2:eq-periodised-weight-a}, and put $\xi^{(+)}=\xi$, $\xi^{(-)}=\overline{\xi}$. For $z\in\mathbb{D}_{\mathfrak{r}}$ define
\begin{equation}\label{6.4:eq-per-ball-bound-a}
\mathcal{E}''(z):=\sum_{c\ \text{occupied}}|\mathcal{C}|^{-1}
\sum_{g\in\mathfrak{G}_\ell}
\Bigl[\mathcal{E}^{\mathrm{bulk}}_{gc}\bigl(u^{(\varepsilon(g))}_{c,g,z}\bigr)\Bigr]^{(\varepsilon(g))}
+\mathcal{E}^{\partial}(w_z),
\end{equation}
and $\mathsf{h}''(z):=e^{-\mathcal{E}''(z)}Z(w_z)/Z(0)$. Then:
\begin{enumerate}
\item[(i)] $\mathcal{E}''$ is holomorphic on $\mathbb{D}_{\mathfrak{r}}$, $\mathcal{E}''(0)=0$, and
the first bound of \eqref{6.4:eq-per-ball-bound-b} holds;
\item[(ii)] $\mathsf{h}''$ is holomorphic on $\mathbb{D}_{\mathfrak{r}}$, $\mathsf{h}''(0)=1$, and
the second bound of \eqref{6.4:eq-per-ball-bound-b} holds on $\mathbb{D}_{\mathfrak{r}}$;
\item[(iii)] consequently the third bound of \eqref{6.4:eq-per-ball-bound-b} holds.
\end{enumerate}
Here
\begin{equation}\label{6.4:eq-per-ball-bound-b}
\begin{gathered}
|\partial_{z_1}\cdots\partial_{z_n}\mathcal{E}''(0)|
\le 2C_{\mathrm{UV}}(d;R)\,(n/r)^n\prod_{i=1}^n\|f_i\|_{\sharp},\\
|\mathsf{h}''(z)|\le Q_\infty^n\quad(z\in\mathbb{D}_{\mathfrak{r}}),\\
|S^T_n(\vec f)|\le C^\infty_n(d;R)\prod_{i=1}^n\|f_i\|_{\sharp},\\
Q_\infty:=\exp\bigl[(E_+^{\sharp}+C_{\mathrm{low}})\ell^4+2C_{\mathrm{face}}(R)\bigr],\\
C^\infty_n(d;R):=\bigl[2C_{\mathrm{UV}}(d;R)
+(\log 2)\,\mathfrak{r}^n(2Q_\infty^n+2)^n\bigr](n/r)^n.
\end{gathered}
\end{equation}
The constant $C^\infty_n(d;R)$ is free of $K$, $g_0$, $m$, $N$ and of the positions of the supports, and it is non-increasing in $d$.
\end{theorem}

Thus (iii) is the volume-free form, for sources in separated balls, of the torus bound of the correlation theorem.

\begin{proof}[Proof of (i)]
\emph{Holomorphy.} Since $K'$ is finite, each $\mathcal{E}^{\mathrm{bulk}}_{gc}$ and $\mathcal{E}^{\partial}$ is a finite sum over $j\le K'$ and over the domains $X$, and the sums over $c$ and $g$ in
\eqref{6.4:eq-per-ball-bound-a} are finite, so $\mathcal{E}''$ has finitely many summands. It suffices to show that each summand is holomorphic on $\mathbb{D}\supset\mathbb{D}_{\mathfrak{r}}$.

Fix an occupied cell $c$ and $g\in\mathfrak{G}_\ell$, and put $h_i^{(g)}:=f_i\circ g^{-1}$ for $i\in I_c$. Since $g$ is an isometry of $\mathbb{T}$, each $h_i^{(g)}$ is real and Lipschitz, with the same sup norm and Lipschitz constant as $f_i$. Hence $\|h_i^{(g)}\|_{\sharp}=\|f_i\|_{\sharp}$.

Suppose first $\varepsilon(g)=+1$. The summand is
\begin{equation*}
\mathcal{E}^{\mathrm{bulk}}_{gc}\Bigl(\sum_{i\in I_c}z_ih_i^{(g)}\Bigr)
=\sum_{j\le K'}\sum_{X\in\mathcal{X}_{gc}}e_j\Bigl(X;\sum_{i\in I_c}z_ih_i^{(g)}\Bigr).
\end{equation*}
This is a finite sum of functions $e_j(X;\cdot)$ composed with a $\mathbb{C}$-linear map of $z$. By part~(c) of Lemma~\ref{6.2:lem-periodised-weight}, for every $z\in\mathbb{D}$ the weight $u_{c,g,z}$ is the weight of the transported system $(h_i^{(g)})_{i\in I_c}$ at the same parameter $z$, and it lies inside that system's polydisc, the norms being equal. The transported system is a finite system of real Lipschitz functions on $\mathbb{T}$. The analyticity property \eqref{6.2:eq-localised-terms-b}, applied to this system at its own parameter, therefore shows that the summand is holomorphic on $\mathbb{D}$.

Suppose now $\varepsilon(g)=-1$. Since the $h_i^{(g)}$ are real,
\begin{equation*}
u^{(-)}_{c,g,z}=\overline{\sum_{i\in I_c}z_ih_i^{(g)}}=\sum_{i\in I_c}\overline{z_i}\,h_i^{(g)}.
\end{equation*}
The summand is therefore $\overline{\Phi_g(\overline{z})}$, where
\begin{equation*}
\Phi_g(\zeta):=\mathcal{E}^{\mathrm{bulk}}_{gc}\Bigl(\sum_{i\in I_c}\zeta_ih_i^{(g)}\Bigr).
\end{equation*}
By the argument of the case $\varepsilon(g)=+1$, $\Phi_g$ is holomorphic on $\mathbb{D}$. We now use a standard fact about complete Reinhardt domains invariant under conjugation. The domain $\mathbb{D}=\{\rho(z)<1\}$ is a complete Reinhardt domain invariant under $z\mapsto\overline{z}$. Hence a function $\Phi$ holomorphic on it is the sum of its power series about $0$, $\Phi(\zeta)=\sum_\alpha a_\alpha\zeta^\alpha$, on all of $\mathbb{D}$. It follows that
\begin{equation}\label{6.4:eq-per-ball-bound-1}
z\longmapsto\overline{\Phi(\overline{z})}=\sum_\alpha\overline{a_\alpha}\,z^\alpha
\end{equation}
is holomorphic on $\mathbb{D}$. Moreover, comparing the coefficients $\alpha!\,\overline{a_\alpha}$ and $\alpha!\,a_\alpha$ gives
\begin{equation}\label{6.4:eq-per-ball-bound-2}
\partial^\alpha\overline{\Phi(\overline{z})}\big|_{z=0}=\overline{\partial^\alpha\Phi(0)}.
\end{equation}
Applied to $\Phi=\Phi_g$, this shows that the summands with $\varepsilon(g)=-1$ are holomorphic on $\mathbb{D}$.

The face part is
\begin{equation*}
\mathcal{E}^{\partial}(w_z)=\sum_j\sum_{X\in\mathcal{X}_\partial}e_j\Bigl(X;\sum_{i=1}^nz_if_i\Bigr).
\end{equation*}
By \eqref{6.2:eq-localised-terms-b} applied to the system $(f_i)_{i\le n}$, it is holomorphic on $\mathbb{D}$.

\emph{Value at zero.} At $z=0$ every weight $u^{(\pm)}_{c,g,0}$ and $w_0$ vanishes. By the locality property \eqref{6.2:eq-localised-terms-a}, $e_j(X;0)=0$ for every $j$ and $X$. Hence every summand of \eqref{6.4:eq-per-ball-bound-a} vanishes at $0$, and $\mathcal{E}''(0)=0$.

\emph{Mixed derivative.} Let $q$ be the number of occupied cells. By Lemma~\ref{6.2:lem-support-placement} and the definition of $I_c$, the sets $I_c$ of the occupied cells are non-empty, pairwise disjoint, and have union $\{1,\dots,n\}$.

Suppose $q\ge2$. Then each $I_c$ is a proper subset of $\{1,\dots,n\}$. Each bulk summand of \eqref{6.4:eq-per-ball-bound-a} belonging to $c$ depends on $z$ only through $(z_i)_{i\in I_c}$, so it is constant in $z_i$ for $i\notin I_c$. The operator $\partial_{z_1}\cdots\partial_{z_n}$ therefore annihilates it, and only the face part contributes.

Suppose $q=1$, so that $I_c=\{1,\dots,n\}$ for the single occupied cell $c$. Take $g$ with $\varepsilon(g)=+1$. The supports of the system $(h_i^{(g)})_{i\le n}$ are the sets $g(\operatorname{supp}f_i)$. Since $g$ is an isometry, each lies in a ball of radius $R$, and they are pairwise at distance at least $d$. Apply Lemma~\ref{6.2:lem-ultraviolet-bound} to this system with the sub-families $\mathcal{Y}_j:=\mathcal{X}_{gc}$ of $D_j$, the domains of scale $j$ lying in the cell $gc$. Its polydisc is $\mathbb{D}$, because $\|h_i^{(g)}\|_{\sharp}=\|f_i\|_{\sharp}$. The lemma gives
\begin{equation*}
\Bigl|\partial_{\zeta_1}\cdots\partial_{\zeta_n}
\Bigl[\mathcal{E}^{\mathrm{bulk}}_{gc}\Bigl(\sum_{i=1}^n\zeta_ih_i^{(g)}\Bigr)\Bigr](0)\Bigr|
\le C_{\mathrm{UV}}(d;R)\,(n/r)^n\prod_{i=1}^n\|f_i\|_{\sharp}.
\end{equation*}
For $g$ with $\varepsilon(g)=-1$ the summand is $\overline{\Phi_g(\overline{z})}$. By \eqref{6.4:eq-per-ball-bound-2} its mixed derivative at $0$ is the complex conjugate of that of $\Phi_g$. Since $\Phi_g$ is the function just bounded, the same bound holds.

For each cell $c'$ there is exactly one element $g_{c\to c'}\in\mathfrak{G}_\ell$ mapping $c$ to $c'$, so $\mathfrak{G}_\ell$ has $|\mathcal{C}|$ elements. The bulk part of \eqref{6.4:eq-per-ball-bound-a} is thus the average, with weight $|\mathcal{C}|^{-1}$, of $|\mathcal{C}|$ terms, each obeying the same bound. Its mixed derivative at $0$ is therefore at most $C_{\mathrm{UV}}(d;R)\,(n/r)^n\prod_i\|f_i\|_{\sharp}$, with no volume factor.

For the face part, in either case, apply Lemma~\ref{6.2:lem-ultraviolet-bound} to the system $(f_i)_{i\le n}$ with the sub-families $\mathcal{Y}_j:=\mathcal{X}_\partial$. This gives
\begin{equation*}
\Bigl|\partial_{z_1}\cdots\partial_{z_n}\bigl[\mathcal{E}^{\partial}(w_z)\bigr](0)\Bigr|
\le C_{\mathrm{UV}}(d;R)\,(n/r)^n\prod_{i=1}^n\|f_i\|_{\sharp}.
\end{equation*}
Adding the bulk and face contributions gives
\begin{equation*}
|\partial_{z_1}\cdots\partial_{z_n}\mathcal{E}''(0)|
\le 2C_{\mathrm{UV}}(d;R)\,(n/r)^n\prod_{i=1}^n\|f_i\|_{\sharp},
\end{equation*}
which is (i).
\end{proof}

\begin{proof}[Proof of parts (ii) and (iii) of Theorem~\ref{6.4:thm-per-ball-bound}]
\label{6.4:prf-recentred-quotient}
Throughout, $\mathcal{E}''$ is the re-centred counterterm of \eqref{6.4:eq-per-ball-bound-a},
$\mathsf{h}''(z)=e^{-\mathcal{E}''(z)}Z(w_z)/Z(0)$, and $\mathcal{E}^{\mathrm{bulk}}_{c'}$,
$\mathcal{E}^{\mathrm{bulk}}$, $\mathcal{E}^{\partial}$ are the bulk and face parts of the counterterm
defined in \eqref{6.2:eq-chessboard-inequality-b}. We write
$\partial^n:=\partial_{z_1}\cdots\partial_{z_n}$.

\emph{Part (ii): holomorphy and normalisation.} The sourced partition function $Z(w_z)$ is the
integral, over the product of the normalised Haar measures of $SU(N_{\mathrm{col}})$ on the bonds of
$T^{(K')}$, of a $z$-independent factor times $e^{A(w_z,U)}$. The domain of integration is
compact, and the integrand is continuous in $(z,U)$ and, for each fixed $U$, an entire function of
$z\in\mathbb{C}^n$, because $A(w_z,U)=\sum_i z_i\,\mathcal{O}(f_i)(U)$ is linear in $z$. Hence
$z\mapsto Z(w_z)$ is entire. Moreover $w_0=0$, so $\mathcal{E}(0)=0$ by the locality property (L) of
\eqref{6.2:eq-localised-terms-a}, and therefore $Z(0)=\widetilde{Z}(0)$, which by (F) in
\eqref{6.2:eq-localised-terms-e} satisfies $\widetilde{Z}(0)\ge e^{-C_{\mathrm{low}}N^4}$; the
right-hand side is positive, so $Z(0)>0$. Since $\mathcal{E}''$ is holomorphic on
$\mathbb{D}_{\mathfrak{r}}$ by part (i), $\mathsf{h}''$ is holomorphic on $\mathbb{D}_{\mathfrak{r}}$.
At $z=0$ every transported weight $u_{c,g,0}$ vanishes, as does its complex conjugate, and $w_0=0$. By (L)
every localised term $e_j(X;\cdot)$ vanishes at the zero weight, so each summand of
\eqref{6.4:eq-per-ball-bound-a} vanishes at $z=0$. Thus $\mathcal{E}''(0)=0$ and
$\mathsf{h}''(0)=Z(0)/Z(0)=1$.

Fix now $z\in\mathbb{D}_{\mathfrak{r}}$.

\emph{The real part of the re-centred counterterm.} For every complex number $\xi$ we have
$\operatorname{Re}\xi^{(+)}=\operatorname{Re}\xi^{(-)}=\operatorname{Re}\xi$. Applying this to each
summand $[\mathcal{E}^{\mathrm{bulk}}_{gc}(u^{(\varepsilon(g))}_{c,g,z})]^{(\varepsilon(g))}$ of
\eqref{6.4:eq-per-ball-bound-a} gives the first equality in
\begin{equation}\label{6.4:eq-recentred-quotient-a}
\begin{split}
\operatorname{Re}\mathcal{E}''(z)
&=\sum_{c\ \mathrm{occ}}|\mathcal{C}|^{-1}\sum_{g\in\mathfrak{G}_\ell}
\operatorname{Re}\mathcal{E}^{\mathrm{bulk}}_{gc}\bigl(u^{(\varepsilon(g))}_{c,g,z}\bigr)
+\operatorname{Re}\mathcal{E}^{\partial}(w_z)\\
&=\sum_{c\ \mathrm{occ}}|\mathcal{C}|^{-1}
\operatorname{Re}\mathcal{E}^{\mathrm{bulk}}\bigl(w^{\mathrm{ch}}_{c,z}\bigr)
+\operatorname{Re}\mathcal{E}^{\partial}(w_z).
\end{split}
\end{equation}
The second equality rests on two facts.

First, for every occupied cell $c$ and every $g\in\mathfrak{G}_\ell$,
\begin{equation}\label{6.4:eq-recentred-quotient-1}
\mathcal{E}^{\mathrm{bulk}}_{gc}\bigl(w^{\mathrm{ch}}_{c,z}\bigr)
=\mathcal{E}^{\mathrm{bulk}}_{gc}\bigl(u^{(\varepsilon(g))}_{c,g,z}\bigr).
\end{equation}
To see this, let $X\in\mathcal{X}_{gc}$ and $p\in\operatorname{plaq}(X^+)$. By the last consequence
of Lemma~\ref{6.2:lem-domain-diameter}, \eqref{6.2:eq-domain-diameter-d}, the plaquette $p$ lies
within $2\varepsilon$ of the open cell $(gc)^\circ$. Consequently it lies at distance greater than
$R+2-2\varepsilon\ge R$ from every ball copy sitting in another cell. By part (b) of the lemma on the
periodised weight (Lemma~\ref{6.2:lem-periodised-weight}),
\[
w^{\mathrm{ch}}_{c,z}=\sum_{g'\in\mathfrak{G}_\ell}u^{(\varepsilon(g'))}_{c,g',z},
\]
and the summand indexed by $g'$ is supported, in the sense of its values at barycentres, in the ball
copies $g'(\bigcup_{i\in I_c}B_i)\subset(g'c)^\circ$, where $B_i$ denotes the ball of radius $R$
carrying the placed source $f_i$. Hence at the barycentre $x_p$ every summand with
$g'\ne g$ vanishes, so
$w^{\mathrm{ch}}_{c,z}(x_p)=u^{(\varepsilon(g))}_{c,g,z}(x_p)$ for all
$p\in\operatorname{plaq}(X^+)$. In words: both weights vanish off the copies, no copy belonging to
another cell comes within $2\varepsilon$ of $(gc)^\circ$, and inside $gc$ the two weights are equal by
definition. The locality property (L) of \eqref{6.2:eq-localised-terms-a} therefore gives
$e_j(X;w^{\mathrm{ch}}_{c,z})=e_j(X;u^{(\varepsilon(g))}_{c,g,z})$. Summing over $j\le K'$ and
$X\in\mathcal{X}_{gc}$ yields \eqref{6.4:eq-recentred-quotient-1}.

Second, by \eqref{6.2:eq-cell-reflections-a} the group $\mathfrak{G}_\ell$ acts simply
transitively on $\mathcal{C}$, so $g\mapsto gc$ is a bijection of $\mathfrak{G}_\ell$ onto
$\mathcal{C}$. Inserting \eqref{6.4:eq-recentred-quotient-1} and reindexing by $c'=gc$ gives
\[
\sum_{g\in\mathfrak{G}_\ell}\operatorname{Re}\mathcal{E}^{\mathrm{bulk}}_{gc}
\bigl(u^{(\varepsilon(g))}_{c,g,z}\bigr)
=\sum_{c'\in\mathcal{C}}\operatorname{Re}\mathcal{E}^{\mathrm{bulk}}_{c'}\bigl(w^{\mathrm{ch}}_{c,z}\bigr)
=\operatorname{Re}\mathcal{E}^{\mathrm{bulk}}\bigl(w^{\mathrm{ch}}_{c,z}\bigr).
\]
This is the second equality of \eqref{6.4:eq-recentred-quotient-a}.

\emph{The chessboard side.} We argue as in the proof of the chessboard inequality with bulk and face
terms (Lemma~\ref{6.2:lem-chessboard-inequality}). By part (a) of Lemma~\ref{6.2:lem-periodised-weight},
$Z(w_z)/Z(0)=\langle\prod_{c\in\mathcal{C}}F_{c,z}\rangle$, with $F_{c,z}\in\mathfrak{A}_c$ for
occupied $c$ and $F_{c,z}:=1$ on unoccupied cells. The chessboard estimate
(Lemma~\ref{6.2:lem-chessboard-estimate}) bounds this by $\prod_{c}\|F_{c,z}\|_{\mathrm{ch}}$. For the
constant function we have $\mathsf{c}(1)=\langle 1\rangle=1$, so $\|1\|_{\mathrm{ch}}=1$. Therefore
\[
\Bigl|\frac{Z(w_z)}{Z(0)}\Bigr|\le\prod_{c\ \mathrm{occ}}\|F_{c,z}\|_{\mathrm{ch}}.
\]

Now let $c$ be occupied. By part (b) of Lemma~\ref{6.2:lem-periodised-weight},
$\|F_{c,z}\|_{\mathrm{ch}}^{|\mathcal{C}|}=Z(w^{\mathrm{ch}}_{c,z})/Z(0)$, and this is a non-negative
real number. We have $Z=e^{\mathcal{E}}\widetilde{Z}$ at $w^{\mathrm{ch}}_{c,z}$, and
$Z(0)=\widetilde{Z}(0)>0$. A non-negative real number equals its modulus, so
\[
\frac{Z(w^{\mathrm{ch}}_{c,z})}{Z(0)}
=\frac{|e^{\mathcal{E}(w^{\mathrm{ch}}_{c,z})}\widetilde{Z}(w^{\mathrm{ch}}_{c,z})|}{\widetilde{Z}(0)}
=\frac{e^{\operatorname{Re}\mathcal{E}(w^{\mathrm{ch}}_{c,z})}\,|\widetilde{Z}(w^{\mathrm{ch}}_{c,z})|}
{\widetilde{Z}(0)}.
\]
The stability bound (S) of \eqref{6.2:eq-localised-terms-d} holds at the periodised weight
$w^{\mathrm{ch}}_{c,z}$ for $z\in\mathbb{D}_{\mathfrak{r}}$. For $\mathfrak{r}=1$ this follows from the
membership $w^{\mathrm{ch}}_{c,z}\in\mathcal{W}_{K'}$ of part (c) of
Lemma~\ref{6.2:lem-periodised-weight}. For $\mathfrak{r}=2$ it follows from the membership of the
parameters $((z_i)_{I_c},(\bar z_i)_{I_c})$ in the polydisc of the system
$(h^+_{c,i},h^-_{c,i})_{i\in I_c}$, see \eqref{6.2:eq-periodised-weight-a}.
Both cases are treated in Remark~\ref{6.2:rem-transported-weights}. The bound (S) gives
$|\widetilde{Z}(w^{\mathrm{ch}}_{c,z})|\le e^{E_+^{\sharp}N^4}$, and (F) gives
$\widetilde{Z}(0)\ge e^{-C_{\mathrm{low}}N^4}$. Splitting
$\mathcal{E}=\mathcal{E}^{\mathrm{bulk}}+\mathcal{E}^{\partial}$ according to the decomposition of each
$D_j$ into the families $\mathcal{X}_{c'}$ and $\mathcal{X}_\partial$, and using
$N^4/|\mathcal{C}|=\ell^4$, we obtain
\begin{equation}\label{6.4:eq-recentred-quotient-2}
\|F_{c,z}\|_{\mathrm{ch}}
\le\exp\Bigl\{|\mathcal{C}|^{-1}\operatorname{Re}\mathcal{E}^{\mathrm{bulk}}(w^{\mathrm{ch}}_{c,z})
+|\mathcal{C}|^{-1}\operatorname{Re}\mathcal{E}^{\partial}(w^{\mathrm{ch}}_{c,z})
+(E_+^{\sharp}+C_{\mathrm{low}})\ell^4\Bigr\}.
\end{equation}

\emph{Combination.} Since $|e^{-\mathcal{E}''(z)}|=e^{-\operatorname{Re}\mathcal{E}''(z)}$, we have
$|\mathsf{h}''(z)|=e^{-\operatorname{Re}\mathcal{E}''(z)}|Z(w_z)/Z(0)|$. We take the product of
\eqref{6.4:eq-recentred-quotient-2} over the $q$ occupied cells and insert
\eqref{6.4:eq-recentred-quotient-a}. This gives
\begin{equation}\label{6.4:eq-recentred-quotient-3}
\begin{split}
|\mathsf{h}''(z)|
\le\exp\Bigl\{&\sum_{c\ \mathrm{occ}}\Bigl[|\mathcal{C}|^{-1}
\operatorname{Re}\mathcal{E}^{\mathrm{bulk}}(w^{\mathrm{ch}}_{c,z})
-|\mathcal{C}|^{-1}\operatorname{Re}\mathcal{E}^{\mathrm{bulk}}(w^{\mathrm{ch}}_{c,z})\Bigr]\\
&+\sum_{c\ \mathrm{occ}}|\mathcal{C}|^{-1}\operatorname{Re}\mathcal{E}^{\partial}(w^{\mathrm{ch}}_{c,z})
-\operatorname{Re}\mathcal{E}^{\partial}(w_z)
+q(E_+^{\sharp}+C_{\mathrm{low}})\ell^4\Bigr\}.
\end{split}
\end{equation}
Each bracket in the first line of \eqref{6.4:eq-recentred-quotient-3} is identically zero, and this
uses no hypothesis beyond the definitions.

We next apply the bound on domains meeting a cell face (Lemma~\ref{6.2:lem-face-estimate}, with the
constant $C_{\mathrm{face}}(R)$ of \eqref{6.2:eq-face-estimate-a}) to the two remaining face terms.

The periodised weight $w^{\mathrm{ch}}_{c,z}$ is a sum over $g\in\mathfrak{G}_\ell$ and $i\in I_c$ of
pieces supported in the ball copies $g(B_i)$ of radius $R$. There are
$|\mathfrak{G}_\ell|\,|I_c|=|\mathcal{C}|\,|I_c|$ such copies, and each lies at distance at least $R+2$
from all face planes. The bound (B) of \eqref{6.2:eq-localised-terms-c} holds at
$w^{\mathrm{ch}}_{c,z}$ by Remark~\ref{6.2:rem-transported-weights}. The face lemma therefore gives
\[
|\mathcal{C}|^{-1}|\mathcal{E}^{\partial}(w^{\mathrm{ch}}_{c,z})|\le|I_c|\,C_{\mathrm{face}}(R).
\]
The weight $w_z=\sum_i z_if_i$ is a sum of $n$ pieces supported in the balls $B_i$, which satisfy the
same distance condition, and (B) of \eqref{6.2:eq-localised-terms-c} holds at $w_z$ for
$z\in\mathbb{D}_{\mathfrak{r}}\subset\mathbb{D}$ (Remark~\ref{6.2:rem-transported-weights}). The face
lemma gives
\[
|\mathcal{E}^{\partial}(w_z)|\le n\,C_{\mathrm{face}}(R).
\]

Every source lies in exactly one occupied cell, so $\sum_c|I_c|=n$. Every occupied cell carries at
least one source, so $q\le n$. Finally $E_+^{\sharp}+C_{\mathrm{low}}\ge 0$: applying (S) and (F) at
the zero weight gives $e^{-C_{\mathrm{low}}N^4}\le\widetilde{Z}(0)\le e^{E_+^{\sharp}N^4}$. Using
$|\operatorname{Re}\xi|\le|\xi|$ for both face terms in \eqref{6.4:eq-recentred-quotient-3}, we
conclude
\[
|\mathsf{h}''(z)|
\le\exp\bigl\{2n\,C_{\mathrm{face}}(R)+n(E_+^{\sharp}+C_{\mathrm{low}})\ell^4\bigr\}
=Q_\infty^{\,n},
\]
with $Q_\infty$ as in \eqref{6.4:eq-per-ball-bound-b}. This holds for every
$z\in\mathbb{D}_{\mathfrak{r}}$, and proves part (ii).

\emph{Part (iii): the cumulant bound.} We have $C_{\mathrm{face}}(R)\ge 0$, since it is a series of
non-negative terms, and $E_+^{\sharp}+C_{\mathrm{low}}\ge 0$. Hence $Q:=Q_\infty^{\,n}\ge 1$. By part
(ii), $\mathsf{h}''$ is holomorphic on $\mathbb{D}_{\mathfrak{r}}$, satisfies $\mathsf{h}''(0)=1$, and
is bounded there by $Q$. So the lemma on the logarithm of a bounded holomorphic function
(Lemma~\ref{6.2:lem-bounded-logarithm}) applies with $\mathsf{h}:=\mathsf{h}''$ and gives
\begin{equation}\label{6.4:eq-recentred-quotient-4}
\bigl|\partial^n\operatorname{Log}\mathsf{h}''(0)\bigr|
\le\log 2\,\bigl(2Q_\infty^{\,n}+2\bigr)^n\mathfrak{r}^n\Bigl(\frac{n}{r}\Bigr)^n
\prod_{i=1}^n\|f_i\|_{\sharp}.
\end{equation}
On the neighbourhood of $z=0$ where $|\mathsf{h}''-1|\le\tfrac12$, the definition of $\mathsf{h}''$
gives $Z(w_z)=Z(0)\,e^{\mathcal{E}''(z)}\,\mathsf{h}''(z)$, with $Z(0)>0$. Hence
\[
\log Z(w_z)=\mathcal{E}''(z)+\log Z(0)+\operatorname{Log}\mathsf{h}''(z).
\]
By part (a) of Lemma~\ref{6.2:lem-periodised-weight} and $A(w_z,U)=\sum_iz_i\,\mathcal{O}(f_i)(U)$,
the quotient $Z(w_z)/Z(0)$ equals $\langle e^{\sum_iz_i\mathcal{O}(f_i)}\rangle$. The joint cumulant
$S^T_n(\vec f)$ of $\mathcal{O}(f_1),\dots,\mathcal{O}(f_n)$ is therefore $\partial^n\log Z(w_z)$ at
$z=0$. The constant $\log Z(0)$ contributes nothing, so
\[
S^T_n(\vec f)=\partial^n\mathcal{E}''(0)+\partial^n\operatorname{Log}\mathsf{h}''(0).
\]
Adding the bound of part (i) on $\partial^n\mathcal{E}''(0)$ to \eqref{6.4:eq-recentred-quotient-4}
gives the bound of part (iii).

\emph{The constants.} The bounds of parts (i)--(iii) are built from the following constants: the
bound constants $C_e$, $\delta$, $\kappa$ of (B) in \eqref{6.2:eq-localised-terms-c}; the retained
rate $c_\kappa$, the anchor count $\mathfrak{a}(R,i)$ and the constant $C_{(1.26)}$ of the
anchor count and tree-graph sum with threshold (Lemma~\ref{6.2:lem-tree-sum}), where $C_{(1.26)}$ is
the constant of the tree-graph bound in the form stated in \cite[(1.26)]{B-CMP116}, on which that sum
rests; the parameter $M$; the per-site constants $E_+^{\sharp}$, $C_{\mathrm{low}}$ of (S) and (F);
the radius $r$ of the source polydisc; and the coupling ceiling $\gamma_J$. By the constants clause
of the correlation theorem these letters are free of $K$, $k$, the torus index $m$, the torus and the
positions of the sources. The constants $E_+^{\sharp}$ and
$C_{\mathrm{low}}$ are the per-site constants of \cite[(0.1)]{B-CMP122b}, which are independent of
the level $k$, of the torus and of the block field $U_k$. The value of $C_{\mathrm{low}}$ at the
torus index $m$ is bounded above by a quantity free of $m$, by the uniformity of the last-scale
constants in the torus exponent (Lemma~\ref{6.3:lem-last-scale-constants}). The cell side
$\ell=\ell(n,R,L)$ depends on $n$, $R$ and $L$ only. The quantities
$\mathfrak{a}$, $C_{\mathrm{UV}}$ (see \eqref{6.2:eq-ultraviolet-bound-a}) and $C_{\mathrm{face}}$ are
explicit series free of $K$, $m$, $N$ and the positions; they are free of $K$ because the sum over
$i\ge 0$ in their definitions majorises the sum over $i=K'-j$, $j\le K'$, that occurs in the
estimates. The translation invariance $S^T_n(\tau_a\vec f)=S^T_n(\vec f)$ of
Lemma~\ref{6.2:lem-support-placement} removes the placement of the sources.
Finally, the constant $C^\infty_n(d;R)$ of \eqref{6.4:eq-per-ball-bound-b} depends on $d$ only through
$C_{\mathrm{UV}}(d;R)$, which is non-increasing in $d$. This gives the monotonicity in $d$ asserted in
part (iii).
\end{proof}

\begin{remark}[The estimate holds at complex sources directly]\label{6.4:rem-complex-sources}
Write $A(v)$ for the weighted plaquette sum $A(v,U)=\sum_p v(x_p)\,a_p(U)$ of
Notation~\ref{6.1:not-units-sources}, where $a_p(U):=1-\operatorname{Re}\operatorname{tr}U(\partial p)$
and $x_p$ is the barycentre of $p$; it is not the Wilson action $A(U)$. We write $x_0:=g_0^{-2}$ for
the constant weight, so that $A(x_0,U)=g_0^{-2}A(U)$. The chessboard inequality
(Lemma~\ref{6.2:lem-chessboard-inequality}) and the
per-ball volume-free bound (Theorem~\ref{6.4:thm-per-ball-bound}) hold at complex sources $z$ directly.
In their proofs the bound is applied to $|\langle e^{A(w_z)}\rangle|$ itself, the vacuum energies compared
are evaluated at complex weights on both sides of each comparison, and the argument is carried out on the
polydisc itself. The case $\operatorname{Im} z=0$ is the special case of the chessboard
inequality in which complex conjugation acts trivially.

Clause (ii) of Theorem~\ref{6.4:thm-per-ball-bound} is a zero-free statement: $\mathsf{h}''$ has no zero
on the polydisc
\begin{equation}\label{6.4:eq-complex-sources-1}
  \Bigl\{\rho(z)\le \frac{1}{2\mathfrak{r}(Q_\infty^{\,n}+1)}\Bigr\},
\end{equation}
whose radius does not depend on $m$. This statement is obtained from reflection positivity without any
decoupling.
\end{remark}

\begin{proof}
The chessboard estimate (Lemma~\ref{6.2:lem-chessboard-estimate}) is proved for complex bounded cell
observables $F_c\in\mathfrak{A}_c$. The reflection $\Theta_\theta$ in it is antilinear. Hence the cell
observable $F_{c,z}=e^{A(w_{c,z})}$, with complex weight $w_{c,z}$, is disseminated as it stands.
On the even cells its copies are $e^{A(u)}$ and on the odd cells they are $e^{A(\bar u)}$, with $u$ the
transported weight $u_{c,g,z}$ of \eqref{6.2:eq-periodised-weight-a}. The conjugation on odd cells is the
only place where complex sources differ from real ones. When $\operatorname{Im} z=0$ it is the identity,
and the argument reduces to the one for real sources.

Next we identify the quantity whose modulus is taken. The action $A(v)$ is linear in $v$, so
$e^{-A(x_0-v)}=e^{-A(x_0)}e^{A(v)}$ and
\begin{equation*}
  \frac{Z(w_z)}{Z(0)}=\bigl\langle e^{A(w_z)}\bigr\rangle .
\end{equation*}
The modulus taken is that of $\langle e^{A(w_z)}\rangle$ itself. Lemma~\ref{6.2:lem-chessboard-estimate}
bounds it by the product of the chessboard norms of the cell observables. Each such norm is the
$|\mathcal{C}|$-th root of the positive number $Z(w^{\mathrm{ch}}_{c,z})/Z(0)$.

The vacuum energies compared in the two results are as follows. In the chessboard inequality
\eqref{6.2:eq-chessboard-inequality-b} the quantity is the real part of the bulk energy
$\mathcal{E}^{\mathrm{bulk}}$ at the complex periodised weight $w^{\mathrm{ch}}$. It is set against
$|\mathcal{C}|$ times the real part of $\mathcal{E}^{\mathrm{bulk}}$ at the same complex weight from which
$w^{\mathrm{ch}}$ is periodised. The two are equal under the covariance hypothesis with which that
inequality is used. In Theorem~\ref{6.4:thm-per-ball-bound} the same quantity is set against
$\operatorname{Re}\mathcal{E}''$ of \eqref{6.4:eq-per-ball-bound-a}. There the two are equal by
definition, since $\mathcal{E}''$ is the cell average of the periodised bulk energies plus the face part
$\mathcal{E}^{\partial}(w_z)$. On odd cells the values are conjugated, which leaves real parts unchanged;
this is the identity \eqref{6.4:eq-recentred-quotient-a}.

Finally, the zero-free statement. By clause (ii) of Theorem~\ref{6.4:thm-per-ball-bound}, $\mathsf{h}''$
is holomorphic on $\mathbb{D}_{\mathfrak{r}}$, $\mathsf{h}''(0)=1$, and $|\mathsf{h}''|\le Q_\infty^{\,n}$
there, with $Q_\infty^{\,n}\ge1$. These are the hypotheses of the lemma on the logarithm of a bounded
holomorphic function (Lemma~\ref{6.2:lem-bounded-logarithm}) with $Q=Q_\infty^{\,n}$. That lemma gives
$|\mathsf{h}''(z)-1|\le\tfrac12$ on the set \eqref{6.4:eq-complex-sources-1}, so $\mathsf{h}''$ has no zero
there. The constant $Q_\infty$ depends on $E_+^{\sharp}$, $C_{\mathrm{low}}$, $\ell$ and
$C_{\mathrm{face}}(R)$ only. Hence the radius in \eqref{6.4:eq-complex-sources-1} is independent of $m$.
The only inputs are the chessboard estimate, which rests on reflection positivity of the lattice measure,
and the cancellation of the bulk terms described above; no decoupling of the measure enters.
\end{proof}

\begin{remark}[Per-ball constants and uniformity under time translation]
\label{6.4:rem-per-ball-constant}
The constant $C^\infty_n(d;R)$ of the per-ball volume-free bound
(Theorem~\ref{6.4:thm-per-ball-bound}) is a per-ball constant. It depends on the common radius
$R$ of the individual balls carrying the supports of the test functions and on the minimal
pairwise distance $d$ between these supports. It does not depend on the mutual positions of
the balls, nor on the diameter of their union.

Let $\mathcal{T}$, $\mathsf{T}_t$ and $(\cdot,\cdot)_{\mathrm{OS}}$ denote the time-zero
algebra, the time translation and the Osterwalder--Schrader form of the following chapter.
The per-ball dependence gives, for every $F\in\mathcal{T}$,
\begin{equation}\label{6.4:eq-per-ball-constant-1}
\sup_{t\ge 0}\,(F,\mathsf{T}_tF)_{\mathrm{OS}}<\infty .
\end{equation}
This is the boundedness hypothesis of the proposition of the following chapter that
establishes the contraction property of the time translations.

A bound of the same shape whose constant depends on the diameter of the union of the supports
does not yield \eqref{6.4:eq-per-ball-constant-1} by the doubling iteration used in that
proposition.
\end{remark}

\begin{proof}
\emph{Per-ball dependence.} In Theorem~\ref{6.4:thm-per-ball-bound} the sources are placed in
cells of side $\ell=\ell(n,R)$, the least power of $L$ with $\ell\ge 8n(R+1)$. This is the
only requirement of the placement, and it involves only $n$ and $R$. The cells absorb
arbitrary mutual distances between the balls. The number $q$ of occupied cells satisfies
$q\le n$. Hence the constant $C^\infty_n(d;R)$
of clause~(iii) of Theorem~\ref{6.4:thm-per-ball-bound}, in the form
\eqref{6.4:eq-per-ball-bound-b}, is free of the positions of the supports. It is also free of
$K$, $g_0$, $m$ and $N$, and it is non-increasing in $d$.

\emph{Uniformity in $t$.} Let $F\in\mathcal{T}$ and $t\ge 0$. Let $e_0$ denote the unit vector
of the time direction. The quantity $(F,\mathsf{T}_tF)_{\mathrm{OS}}$ is a finite sum of
products of limits $S^T_{|B|;\infty}$, where, by the expression of moments through cumulants,
$B$ runs over the blocks of partitions of the set of smeared observables entering
$(F,\mathsf{T}_tF)_{\mathrm{OS}}$. Each limit is evaluated at test
functions supported in balls of a common radius $R_F$ depending only on $F$, of one of two
kinds. The reflected ones are
fixed. The shifted ones are translated by $te_0$, so their supports still lie in balls of the
same radius $R_F$.

Set
\[
d_0:=\min\bigl\{\text{the gaps within the sequence of time slabs carrying }F,\
2\min I_1\bigr\}>0,
\]
where $I_1$ is the time interval of the slab of $F$ nearest the reflection plane, as in the
definition of $\mathcal{T}$ in the following chapter. The pairwise distances between the
supports are at
least $d_0$, and $d_0$ is independent of $t$: the time translation only increases the distance
between the reflected group and the shifted group.

Moreover $\|\mathsf{T}_tf\|_{\sharp}=\|f\|_{\sharp}$ for every test function $f$. Clause~(iii)
of Theorem~\ref{6.4:thm-per-ball-bound}, in the form \eqref{6.4:eq-per-ball-bound-b}, has a
constant free of $K$, $g_0$, $m$ and of the positions of the supports. Since the limits
$S^T_{|B|;\infty}$ exist by the infinite-volume theorem of the following chapter, the bound
passes to them. Since $C^\infty_{|B|}(\cdot\,;R_F)$ is non-increasing in its first
argument, each factor is bounded by
\[
C^\infty_{|B|}(d_0;R_F)\prod_{i\in B}\|f_i\|_{\sharp},
\]
uniformly in $t$. Summing the finitely many terms gives \eqref{6.4:eq-per-ball-constant-1}.

\emph{Failure of a union-diameter constant.} Now replace the per-ball bound by its single-ball
form, in which all supports lie in one ball $B_R$. Under $\mathsf{T}_t$ the radius of that ball
must satisfy $R\gtrsim t$. Let $D(t)$ denote the diameter of the union of the supports of the
test functions entering $(F,\mathsf{T}_tF)_{\mathrm{OS}}$, so that $D(t)\gtrsim 2t$. The bound
then available, which we denote by $\Xi(t)$, satisfies
\begin{equation}\label{6.4:eq-per-ball-constant-2}
\Xi(t)\asymp\exp\bigl\{n\,(E_+^{\sharp}+C_{\mathrm{low}})\,D(t)^4\bigr\},
\end{equation}
with $n$ here the number of smeared observables entering $(F,\mathsf{T}_tF)_{\mathrm{OS}}$.
The doubling iteration in the proof of the contraction property bounds the quantity it
controls, at step $k$, by
\[
\bigl(\Xi(2^{k+1}t)/(F,F)_{\mathrm{OS}}\bigr)^{2^{-k-1}} .
\]
By \eqref{6.4:eq-per-ball-constant-2} with $D(2^{k+1}t)\gtrsim 2^{k+2}t$, and with a constant
$C>0$ independent of $k$, this is of order
\[
\bigl(e^{C16^kt^4}\bigr)^{2^{-k-1}}=e^{(C/2)\,8^kt^4}\longrightarrow\infty
\qquad(k\to\infty).
\]
So the iteration diverges, and the single-ball form does not yield the contraction property.
This does not show that the contraction property fails; it shows only that the single-ball
form is not the estimate to use there. The per-ball form is what the proof of the single-ball
form itself gives once the placement is carried out ball by ball, as in
Theorem~\ref{6.4:thm-per-ball-bound}.
\end{proof}

\begin{remark}[The bound at small cutoffs]\label{6.4:rem-small-cutoffs}
The per-ball volume-free bound, Theorem~\ref{6.4:thm-per-ball-bound}, holds at every $K'\ge 0$;
here $K'$ is the index of the last lattice of the correlation theorem, equal to the number $K$ of
renormalisation steps when that lattice is the unit lattice.
A separate treatment of finitely many small values of $K$, as in the second case of
Lemma~\ref{6.3:lem-last-scale-constants}, is therefore not needed. If one is used,
each observable must be bounded through the plaquettes of its ball and not through those of the
whole torus. Let $B_i$ be the ball of radius $R$ containing the support of $f_i$, and let
$B_i^{+1}$ be its enlargement by distance one. The required bound is
\begin{equation}\label{6.4:eq-small-cutoffs-1}
|\mathcal{O}(f_i)|
\le 2\,\#\operatorname{plaq}(B_i^{+1})\,\|f_i\|_\infty
\le 12(2R+4)^4L^{4K}\|f_i\|_\infty ,
\end{equation}
for the finitely many $K$ so treated, all of which satisfy $K<m_{\mathbb{T}}(g)$. Here
$m_{\mathbb{T}}(g)$ is the torus floor of the correlation theorem taken at the coupling $g$ itself;
it is to be distinguished from the floor $m_{\mathbb{T}}(g_-(g))$ in the hypothesis on $m$ of
Theorem~\ref{6.4:thm-per-ball-bound}, and it does not depend on $m$, so that the
right-hand side is bounded independently of $m$. With either treatment, the bound
\eqref{6.4:eq-per-ball-bound-b}
holds for every $K$ with a constant independent of $m$.
\end{remark}

\begin{proof}
The hypotheses of Theorem~\ref{6.4:thm-per-ball-bound} place no restriction on $K'$, which is
arbitrary there. Its conclusion therefore holds at every $K'\ge 0$, with the constant stated there,
and that constant does not depend on $m$. Consequently \eqref{6.4:eq-per-ball-bound-b} holds for
every $K$ with an $m$-free constant, and no separate treatment of small $K$ is required.

Suppose nevertheless that finitely many small $K$ are treated separately, by bounding each
observable $\mathcal{O}(f_i)$ directly, as in the second case of
Lemma~\ref{6.3:lem-last-scale-constants}. The bound of that case counts all plaquettes of the
torus and so produces a factor equal to
the number of those plaquettes, namely $6(2L^{m_1+m_{\mathbb{T}}(g)})^4$, with $m_1$ as in that
case. That number depends on $m$; its $m$-dependence is
an artefact of the choice of units, of the same kind as the one met in the ultraviolet locality
constant $C_{\mathrm{UV}}(d;R)$.

Instead, recall that $\mathcal{O}(f_i)=\sum_p f_i(x_p)\,a_p(U)$. Since $\operatorname{tr}$ is
normalised by $\operatorname{tr}\mathbf{1}=1$, we have $|\operatorname{Re}\operatorname{tr}U(\partial p)|\le 1$,
and hence $0\le a_p(U)\le 2$. Only plaquettes whose barycentre $x_p$ lies in the support of $f_i$
contribute, since for every other plaquette $f_i(x_p)=0$. The support of $f_i$ lies in $B_i$, so
every contributing plaquette meets $B_i$, and in particular meets $B_i^{+1}$. This gives the first
inequality in \eqref{6.4:eq-small-cutoffs-1}.

The bare plaquettes meeting $B_i^{+1}$ number at most $6(2R+4)^4L^{4K}$, by counting six plaquettes
per bare site. This gives the second inequality.

In the range of $K$ treated separately, $K$ lies below the torus floor $m_{\mathbb{T}}(g)$ of the
correlation theorem at the coupling $g$ (not at $g_-(g)$), and that floor does not depend on $m$.
Hence $L^{4K}$, and with it the right-hand
side of \eqref{6.4:eq-small-cutoffs-1}, is bounded independently of $m$.

Thus, whichever treatment of small $K$ is adopted, \eqref{6.4:eq-per-ball-bound-b} holds for every
$K$ with an $m$-free constant.
\end{proof}

\begin{remark}[Conversion to the reference lattice units]\label{6.4:rem-unit-conversion}
Let $n\ge 2$, $R>0$, $d>0$, and let $C^\infty_n(d;R)$ be the per-ball constant
\eqref{6.4:eq-per-ball-bound-b} in clause (iii) of the per-ball volume-free bound
(Theorem~\ref{6.4:thm-per-ball-bound}). That bound is written on the last lattice $T^{(K')}$, with
$N$ sites per side, in last-lattice units. Its conversion to the reference lattice units, the
units in which Theorem~\ref{3.1:thm-main} is stated, is as follows.
\begin{enumerate}
\item[(1)] Under the identification of the last lattice of the correlation theorem with the
reference unit lattice, $K'=K$, $T^{(K')}$ is the reference unit lattice and $N=2L^m$; the bound
then holds in the reference units without change.
\item[(2)] In the first case of the lemma fixing the torus index of the correlation theorem,
$N=2L^m\cdot L^b$, where $b$ is a non-negative integer with $b\le b_*(g)$, the bound $b_*(g)$ is
independent of $K$ and of $m$, and $b$ is eventually constant along the subsequence of clause
(ii-b) of Theorem~\ref{3.1:thm-main}. Then, under the hypotheses of the per-ball volume-free bound on $T^{(K')}$
with $(d,R)$ replaced by $(dL^b,RL^b)$, in particular with $\ell(n,RL^b)$ in place of
$\ell(n,R)$ in its lower bound on the torus index, for test
functions $f_1,\dots,f_n$ supported in balls of radius $R$ with pairwise support distance at
least $d$, radii, distances and norms $\|f_i\|_{\sharp}$ all taken in the reference units,
\begin{equation}\label{6.4:eq-unit-conversion-1}
|S^T_{n;K,m}(f_1,\dots,f_n)|\le \Bigl(\max_{\substack{0\le b\le b_*(g)\\ b\in\mathbb{Z}}}
C^\infty_n(dL^b;RL^b)\Bigr)\prod_{i=1}^n\|f_i\|_{\sharp}.
\end{equation}
The constant in \eqref{6.4:eq-unit-conversion-1} is free of $K$, $g_0$, $m$ and of the
positions of the supports.
\end{enumerate}
\end{remark}

\begin{proof}
In case~(1) the last lattice of the per-ball volume-free bound is the reference unit lattice,
so its radii, distances, norms and Schwinger functions are those of the reference units, and
nothing changes.

In case~(2) a physical length equals the corresponding last-lattice length times $L^{-b}$. For
a test function $f$ in the reference units put $f^{(b)}(y):=f(L^{-b}y)$, the same function
written in last-lattice coordinates. A barycentre at last-lattice position $x_p$ sits at
physical position $L^{-b}x_p$, so $\mathcal{O}(f)$ coincides with the observable built from
$f^{(b)}$ on $T^{(K')}$, and $S^T_{n;K,m}(f_1,\dots,f_n)$ with the connected function
$S^T_n(f^{(b)}_1,\dots,f^{(b)}_n)$ on $T^{(K')}$. A support in a ball of physical radius $R$ lies in
a ball of last-lattice radius $RL^b$, and a physical support distance at least $d$ is a
last-lattice distance at least $dL^b$. Moreover $\|f^{(b)}\|_\infty=\|f\|_\infty$ and
$\operatorname{Lip}f^{(b)}=L^{-b}\operatorname{Lip}f\le\operatorname{Lip}f$. Both norms are
$\max\{\|\cdot\|_\infty,40M\operatorname{Lip}\}$, so the last-lattice norm of $f^{(b)}$ is at
most the reference norm of $f$. Clause (iii) of the per-ball volume-free bound
(Theorem~\ref{6.4:thm-per-ball-bound}) on $T^{(K')}$ therefore gives
\begin{align*}
|S^T_{n;K,m}(f_1,\dots,f_n)|&=|S^T_n(f^{(b)}_1,\dots,f^{(b)}_n)|
\le C^\infty_n(dL^b;RL^b)\prod_{i=1}^n\|f^{(b)}_i\|_{\sharp}\\
&\le C^\infty_n(dL^b;RL^b)\prod_{i=1}^n\|f_i\|_{\sharp},
\end{align*}
and bounding the constant by its maximum over the finitely many values of $b$ gives
\eqref{6.4:eq-unit-conversion-1}.

Each $C^\infty_n(dL^b;RL^b)$ is finite and free of $K$, $g_0$, $m$ and of the positions, by
the per-ball volume-free bound. The admissible values of $b$ are finitely many, and their bound
$b_*(g)$ is free of $K$ and $m$. Hence the maximum has the same properties.

Its dependence on the radius can be described further. The proof of the ultraviolet bound for
separated ball sources (Lemma~\ref{6.2:lem-ultraviolet-bound}) shows that the constant
\eqref{6.2:eq-ultraviolet-bound-a} satisfies
\begin{equation*}
C_{\mathrm{UV}}(d;R)\le C(L,\mathfrak{p})\Bigl(1+\frac{R+2}{d}\Bigr)^4
\Bigl(1+\log^+_L\frac{32M}{d}\Bigr),
\end{equation*}
where $\log^+_L x:=\max\{\log_L x,0\}$,
with $C(L,\mathfrak{p})$ depending only on $L$ and $\mathfrak{p}$. Since $L^b\ge 1$, we have
$(RL^b+2)/(dL^b)\le (R+2)/d$. Since $\log^+_L$ is non-decreasing,
$\log^+_L(32M/(dL^b))\le\log^+_L(32M/d)$. Thus $C_{\mathrm{UV}}$ is scale-invariant in
$(d,R)$ up to its logarithmic tail, and at $(dL^b,RL^b)$ it is bounded by the same expression
as at $(d,R)$. The remaining radius-dependent quantities of the per-ball bound are the
constant $C_{\mathrm{face}}(RL^b)$ of the bound on domains meeting a cell face
(Lemma~\ref{6.2:lem-face-estimate}, display \eqref{6.2:eq-face-estimate-a}) and the cell side
$\ell(n,RL^b)$. Both are bounded over the finitely many values of $b$; for $\ell(n,RL^b)$ this
means that the lower bound on the torus index required in the per-ball volume-free bound can be
taken the same for all admissible $b$.
\end{proof}

\paragraph{Inputs used by statement}
For each clause of Theorem~\ref{3.1:thm-main} we list, in two categories, the inputs used in the
proofs of this chapter. The first comprises the published statements of Ba{\l}aban cited by statement,
each with its reference; the second, the statements of this book whose own proofs are
incomplete or outstanding. The proofs of this chapter bear on clause (ii), through reflection
positivity of the lattice measure (Lemma~\ref{6.1:lem-reflection-positivity}), and on clauses (ii)
and (iii)(b), through the per-ball volume-free bound (Theorem~\ref{6.4:thm-per-ball-bound}); the
published statements used in the proof of that theorem are listed once, under (iii).
The chapter also uses the correlation theorem by statement: Definition~\ref{6.2:def-localised-terms}
is stated under its standing hypotheses, and the dependence of its constants on the data is taken
from its statement. That theorem is proved as an implication from statements of Ba{\l}aban as cited
in its proof, with the readings declared there where his construction admits more than one reading,
and from the statements of this book on which clause (i) rests, which are listed with
their standing together with the proof of clause (i); its antecedents are therefore not repeated
here, and it is an item of neither category below.

\begin{description}
\item[(0)] This chapter proves no part of clause (0); there are no items under it.

\item[(i)] This chapter proves no part of clause (i); there are no items under it.

\item[(iii)] Published statements cited by statement:
\begin{itemize}
\item \cite[(0.24)--(0.26)]{B-CMP109}, used in the proof of
Lemma~\ref{6.2:lem-domain-diameter};
\item \cite[(1.26)]{B-CMP116}, used in the proofs of Lemmas~\ref{6.2:lem-tree-sum},
\ref{6.2:lem-ultraviolet-bound} and~\ref{6.2:lem-face-estimate};
\item \cite[(2.49), (2.50)]{B-CMP119}, used in the bounds at transported and periodised weights;
\item \cite[p.~355]{B-CMP122b}, cited among the statements on the dependence on the block
exponent that are assumed in Lemma~\ref{6.3:lem-last-scale-constants};
\item \cite[(0.1)]{B-CMP122b}, used for the re-centred quotient and the cumulant bound in the
proof of Theorem~\ref{6.4:thm-per-ball-bound}.
\end{itemize}
Statements of this book whose proofs are incomplete or outstanding: none.

\item[(ii)] Published statements cited by statement:
\begin{itemize}
\item \cite[(0.1), (0.2)]{B-CMP109}, used in the proof of
Lemma~\ref{6.1:lem-reflection-positivity}.
\end{itemize}
Statements of this book whose proofs are incomplete or outstanding: none.

\item[(iv)] This chapter proves no part of clause (iv); there are no items under it.

\item[(v)] This chapter proves no part of clause (v); there are no items under it.

\item[(vi)] This chapter proves no part of clause (vi); there are no items under it.
\end{description}

The standing parameters of Notation~\ref{6.3:not-standing-parameters} and the floors of
Definitions~\ref{6.3:def-one-sided-floors} and~\ref{6.3:def-block-exponent-floor} are fixed in the
order of Notations~\ref{2.3:not-stages-first}, \ref{2.3:not-stages-middle}
and~\ref{2.3:not-stages-last}. By Lemma~\ref{2.4:lem-no-cycle}, that order has no cycle and
imposes only one-sided conditions.

%% ---- end inlined file: chapters/c06 ----

\clearpage
\setcounter{section}{6}
\refstepcounter{section}
\section*{Chapter~\thesection\quad Limits, invariance, reflection positivity and reconstruction: clause (ii)\addcontentsline{toc}{section}{\protect\numberline{\thesection}Limits, invariance, reflection positivity and reconstruction: clause (ii)}\label{ch:7}}
%% ---- begin inlined file: chapters/c07 ----
\begin{definition}[Curvature cumulants and moments on the torus]\label{7.1:not-torus-cumulants}
Throughout, $N=2$, and the hypotheses (H1)--(H7) of Theorem~0 hold. The observable
$\mathcal{O}(f)$, the norm $\|f\|_\sharp$, the weight $w=\sum_i z_if_i$ and the partition function
with sources $Z(z)$ are as in the glossary of Chapter~0; $\omega_K$ denotes the Wilson state and
$\mu_{K,m}$ the Wilson measure on $\mathbb{T}_m$ at cutoff $K$.

Consider the torus $\mathbb{T}_m$ at cutoff $K$, with lattice spacing $\varepsilon=L^{-K}$. For
real Lipschitz functions $f_1,\dots,f_n$ on $\mathbb{T}_m$ set
\begin{equation}\label{7.1:eq-torus-cumulants-a}
S^T_{n;K,m}(f_1,\dots,f_n)
:=\partial_{z_1}\cdots\partial_{z_n}\log Z(0)
=\bigl\langle\mathcal{O}(f_1);\dots;\mathcal{O}(f_n)\bigr\rangle^{\mathrm{conn}}_{K},
\end{equation}
where $Z$ is the partition function with sources built from $f_1,\dots,f_n$. Thus
$S^T_{n;K,m}(f_1,\dots,f_n)$ is the joint cumulant of $\mathcal{O}(f_1),\dots,\mathcal{O}(f_n)$
under the Wilson measure $\mu_{K,m}$. It is multilinear in $(f_1,\dots,f_n)$, and we extend it
complex-multilinearly to complex test functions.

We use the following conventions.
\begin{itemize}
\item Distances on $\mathbb{T}_m$ are torus distances, and
\begin{equation*}
(\mathbb{T}_m)^n_{\neq}
:=\bigl\{(x_1,\dots,x_n)\in\mathbb{T}_m^n:\ x_i\neq x_j\ \text{for}\ i\neq j\bigr\}.
\end{equation*}
\item $\mathcal{C}^{(n)}_{\mathrm{sep}}(\mathbb{T}_m)$ is the set of $n$-tuples of real
$C^\infty$ functions on $\mathbb{T}_m$ whose supports are pairwise at positive distance.
\item For a $C^1$ function $f$ on $\mathbb{T}_m$ put
$\|f\|_{C^1}:=\max\{\|f\|_\infty,\|\nabla f\|_\infty\}$. With $c_\sharp$ the constant in the
comparison of the norm $\|\cdot\|_\sharp$ with the $C^1$ norm, put $c':=\max\{1,c_\sharp\}$;
then $\|f\|_\sharp\le c'\|f\|_{C^1}$.
\item $W_4$ is the hypercubic group (permutations and sign changes of the coordinates), acting on
$\mathbb{T}_m$ about the origin, and $\tau_af:=f(\cdot-a)$ for $a\in\mathbb{T}_m$.
\item Take coordinates in $(-L^m,L^m]$. Then $\theta(x_0,\vec x):=(-x_0,\vec x)$ is the reflection
of $\mathbb{T}_m$ in the pair of hyperplanes $\{x_0=0\}$ and $\{x_0=L^m\}$, and
$\mathbb{T}_m^+:=\{0<x_0<L^m\}$.
\end{itemize}

For $q\ge0$ and functions $h_1,\dots,h_q$ with pairwise disjoint compact supports put
\begin{equation}\label{7.1:eq-torus-cumulants-b}
S_{q;K,m}(h_1,\dots,h_q)
:=\sum_{\pi\in\Pi_{\ge2}[q]}\ \prod_{B\in\pi}S^T_{|B|;K,m}\bigl((h_l)_{l\in B}\bigr),
\end{equation}
where $\Pi_{\ge2}[q]$ is the set of partitions of $\{1,\dots,q\}$ into blocks of size at
least $2$. In particular $S_{0;K,m}=1$ and $S_{1;K,m}=0$. The same formula with the fixed-torus
limits $S^T_{|B|;m}$ in place of $S^T_{|B|;K,m}$ defines $S_{q;m}$.
\end{definition}

\begin{proof}
Multilinearity of \eqref{7.1:eq-torus-cumulants-a}. The map $f\mapsto\mathcal{O}(f)$ is linear,
since $\mathcal{O}(f)=\sum_p f(x(p))\bigl[1-\operatorname{Re}\operatorname{tr}U(\partial p)\bigr]$
depends linearly on the values of $f$. A joint cumulant of $n$ random variables is multilinear in
those variables. Composing the two maps shows that
$(f_1,\dots,f_n)\mapsto S^T_{n;K,m}(f_1,\dots,f_n)$ is multilinear over the reals. A real
multilinear form has a unique complex-multilinear extension to complex arguments, and this
extension is the one meant in the definition.

The case $q=0$ of \eqref{7.1:eq-torus-cumulants-b}. The set $\{1,\dots,q\}$ is empty, and
$\Pi_{\ge2}[0]$ consists of exactly one element, the empty partition. This partition has no blocks,
so the condition on block sizes holds vacuously. The corresponding product over blocks is empty and
equals $1$. Hence $S_{0;K,m}=1$.

The case $q=1$. A partition of $\{1\}$ has the single block $\{1\}$, which has size $1<2$. Hence
$\Pi_{\ge2}[1]$ is empty, the sum in \eqref{7.1:eq-torus-cumulants-b} has no terms, and
$S_{1;K,m}=0$.

The same two arguments apply verbatim to $S_{q;m}$, since they use only the index set of the sum
and not the factors.
\end{proof}

\begin{proposition}[Fixed-torus limits along first passage]\label{7.1:prop-torus-limits}
Fix $m\ge m_{\mathbb{T}}(g_-(g))$ and bare couplings $g_0^{(i)}\in(0,g_-(g)]$ with
$K_i:=K(g_0^{(i)},g)\to\infty$. Assume:
\begin{itemize}
\item[(B)] for every $n\ge2$ and $\mathsf{d}>0$ there is $C_n(\mathsf{d})<\infty$ such that
\begin{equation}\label{7.1:eq-torus-limits-1}
|S^T_{n;K_i,m}(f_1,\dots,f_n)|\le C_n(\mathsf{d})\prod_l\|f_l\|_\sharp
\end{equation}
for every $i$ and all real Lipschitz $f_1,\dots,f_n$ on $\mathbb{T}_m$ whose supports are
pairwise at distance at least $\mathsf{d}$.
\end{itemize}
Then there is a subsequence $(i_k)$ with the following properties.
\begin{itemize}
\item[(a)] For every $n\ge2$ and every $\vec f\in\mathcal{C}^{(n)}_{\mathrm{sep}}(\mathbb{T}_m)$,
the values $S^T_{n;K_{i_k},m}(\vec f)$ converge to a limit $S^T_{n;m}(\vec f)$. The limit is
multilinear. Moreover,
\[
|S^T_{n;m}(\vec f)|\le C_n(\mathsf{d})\prod_l\|f_l\|_\sharp
\]
whenever the supports are pairwise at distance at least $\mathsf{d}$.
\item[(b)] There is a unique distribution on $(\mathbb{T}_m)^n_{\neq}$, again denoted
$S^T_{n;m}$, with $S^T_{n;m}(f_1\otimes\cdots\otimes f_n)=S^T_{n;m}(\vec f)$ for all
$\vec f\in\mathcal{C}^{(n)}_{\mathrm{sep}}(\mathbb{T}_m)$. It is of locally finite order.
\item[(c)] $S^T_{n;m}$ is symmetric. It is invariant under $\tau_a$ for all
$a\in\mathbb{R}^4/2L^m\mathbb{Z}^4$ and under $W_4$.
\item[(d)] Let $r\ge1$, let $f_1,\dots,f_r$ be real $C^\infty$ functions with pairwise disjoint
compact supports in $\mathbb{T}_m^+$, and let $c_\alpha$, $\alpha\subset\{1,\dots,r\}$, be
complex numbers. Then, with $S_{q;m}$ the moments built from the limits $S^T_{|B|;m}$ as in
\eqref{7.1:eq-torus-cumulants-b},
\begin{equation}\label{7.1:eq-torus-limits-b}
\sum_{\alpha,\beta}\overline{c_\alpha}\,c_\beta\,
S_{|\alpha|+|\beta|;m}\bigl((f_i\circ\theta)_{i\in\alpha},(f_j)_{j\in\beta}\bigr)\ge0 .
\end{equation}
By (c), \eqref{7.1:eq-torus-limits-b} also holds with $\theta$ replaced by the reflection
$x_\mu\mapsto2c-x_\mu$ of $\mathbb{T}_m$, which fixes
$\{x_\mu=c\}\cup\{x_\mu=c+L^m\}$, for every $\mu$ and every $c\in\mathbb{R}$, provided
$f_1,\dots,f_r$ are real $C^\infty$ functions with pairwise disjoint compact supports in
$\{c<x_\mu<c+L^m\}$.
\end{itemize}
Nothing is claimed for $n=1$ or at coinciding points.
\end{proposition}

Hypothesis (B) is the torus bound, that is, the fixed-torus bound of clause (iii)(b) of
Theorem~0, which Theorem~0 does not assert.

\begin{proof}[Proof of Proposition~\ref{7.1:prop-torus-limits}: the subsequence and clause (a)]
Throughout, $S^T_{n;K,m}$ is the joint cumulant \eqref{7.1:eq-torus-cumulants-a} of the
curvature observables (Notation~\ref{7.1:not-torus-cumulants}). It is multilinear in
$(f_1,\dots,f_n)$ because $f\mapsto\mathcal{O}(f)$ is linear. We use the $C^1$ norm
$\|f\|_{C^1}=\max\{\|f\|_\infty,\|\nabla f\|_\infty\}$ and the comparison of norms
\begin{equation}\label{7.1:eq-torus-limits-2}
\|f\|_\sharp\le c'\|f\|_{C^1},
\end{equation}
where the norm $\|\cdot\|_{C^1}$ and the constant $c'$ are those fixed in
Notation~\ref{7.1:not-torus-cumulants}.

\emph{Step 1: a countable family of test functions.} Lipschitz functions do not form a separable
space in the Lipschitz norm, so we work in the $C^1$ norm instead. The real $C^\infty$ functions
on $\mathbb{T}_m$ form a subset of the separable space of real $C^1$ functions on the compact
torus $\mathbb{T}_m$, and they are dense in it for the $C^1$ norm (convolution with a smooth
approximate identity on $\mathbb{T}_m$). Let $D_0$ be a countable set of real $C^\infty$
functions on $\mathbb{T}_m$, dense in the real $C^1$ functions for the $C^1$ norm.

Let $Q$ be a finite union of closed balls with rational centres and radii, and let $\eta>0$ be
rational. There are countably many such pairs $(Q,\eta)$. For each of them fix a real
$\chi_{Q,\eta}\in C^\infty(\mathbb{T}_m)$ with $\chi_{Q,\eta}=1$ on $Q$ and
$\operatorname{supp}\chi_{Q,\eta}\subset Q^{+\eta}$, the closed $\eta$-neighbourhood of $Q$.
Set
\[
D:=\{\chi_{Q,\eta}h:\ h\in D_0,\ Q,\eta\text{ as above}\}.
\]
The set $D$ is countable, and all its elements are real $C^\infty$ functions, hence real and
Lipschitz. Suppose $\operatorname{supp}f\subset Q$ and $h_k\to f$ in $C^1$. For every $C^1$
function $u$, $\nabla(\chi_{Q,\eta}u)=u\nabla\chi_{Q,\eta}+\chi_{Q,\eta}\nabla u$. Hence
\[
\|\chi_{Q,\eta}u\|_{C^1}\le\bigl(\|\chi_{Q,\eta}\|_\infty+\|\nabla\chi_{Q,\eta}\|_\infty\bigr)
\|u\|_{C^1}.
\]
Applied to $u=h_k-f$, this gives $\chi_{Q,\eta}h_k\to\chi_{Q,\eta}f$ in $C^1$. Moreover
$\chi_{Q,\eta}f=f$, because $\chi_{Q,\eta}=1$ on $Q\supset\operatorname{supp}f$. All
approximants $\chi_{Q,\eta}h_k$ are supported in $Q^{+\eta}$.

\emph{Step 2: equicontinuity.} Let $\vec f=(f_1,\dots,f_n)$ and $\vec g=(g_1,\dots,g_n)$ be
tuples of real Lipschitz functions on $\mathbb{T}_m$. Suppose that $f_l$ and $g_l$ are supported
in $A_l$, where the sets $A_1,\dots,A_n$ are pairwise at distance at least $\mathsf{d}'>0$.
Write $S$ for any map that is multilinear in its $n$ arguments. For $0\le l\le n$ let
$\vec t^{\,(l)}:=(g_1,\dots,g_l,f_{l+1},\dots,f_n)$ be the hybrid tuple, so that
$\vec t^{\,(0)}=\vec f$ and $\vec t^{\,(n)}=\vec g$. The tuples $\vec t^{\,(l-1)}$ and
$\vec t^{\,(l)}$ differ only in the $l$-th entry, which is
$f_l$ in the first and $g_l$ in the second. Telescoping and linearity in the $l$-th entry
therefore give
\begin{align*}
S(\vec f)-S(\vec g)
&=\sum_{l=1}^n\bigl(S(\vec t^{\,(l-1)})-S(\vec t^{\,(l)})\bigr)\\
&=\sum_{l=1}^n S(g_1,\dots,g_{l-1},f_l-g_l,f_{l+1},\dots,f_n).
\end{align*}
We apply this with $S=S^T_{n;K_i,m}$.

In each hybrid tuple every entry is real and Lipschitz. The $l'$-th entry is $g_{l'}$, $f_{l'}$
or $f_{l'}-g_{l'}$, and in each case it is supported in $A_{l'}$. Hence the supports of the
entries are pairwise at distance at least $\mathsf{d}'$, and (B) applies with
$\mathsf{d}=\mathsf{d}'$. It bounds the $l$-th term by $C_n(\mathsf{d}')$ times
\[
\|f_l-g_l\|_\sharp\prod_{l'<l}\|g_{l'}\|_\sharp\prod_{l'>l}\|f_{l'}\|_\sharp .
\]
Each factor $\|g_{l'}\|_\sharp$ or $\|f_{l'}\|_\sharp$ is at most
$\|f_{l'}\|_\sharp+\|g_{l'}\|_\sharp$. Therefore, uniformly in $i$,
\begin{equation}\label{7.1:eq-torus-limits-a}
\begin{split}
\bigl|S^T_{n;K_i,m}(\vec f)&-S^T_{n;K_i,m}(\vec g)\bigr|\\
&\le C_n(\mathsf{d}')\sum_{l=1}^n\|f_l-g_l\|_\sharp
\prod_{l'\neq l}\bigl(\|f_{l'}\|_\sharp+\|g_{l'}\|_\sharp\bigr).
\end{split}
\end{equation}

\emph{Step 3: extraction; proof of (a).} For each $n\ge2$ the set $D^n$ is countable. Hence the
set of all tuples in $\bigcup_{n\ge2}D^n$ whose supports are pairwise at positive distance is
countable. Let $\vec g\in D^n$ be such a tuple, and let $\mathsf{d}_{\vec g}>0$ be the smallest
distance between two of its supports. Its entries are real and Lipschitz, so (B) gives
\[
|S^T_{n;K_i,m}(\vec g)|\le C_n(\mathsf{d}_{\vec g})\prod_l\|g_l\|_\sharp
\qquad\text{for every }i.
\]
Thus
the sequence $\bigl(S^T_{n;K_i,m}(\vec g)\bigr)_i$ is bounded. Enumerate these countably many
bounded sequences. Extract a convergent subsequence of the first, from it a convergent
subsequence of the second, and so on, and take the diagonal. This gives a subsequence $(i_k)$
along which all these sequences converge.

Now let $\vec f\in\mathcal{C}^{(n)}_{\mathrm{sep}}(\mathbb{T}_m)$ have supports pairwise at
distance at least $\mathsf{d}>0$.
\begin{itemize}
\item Each $\operatorname{supp}f_l$ is compact. Cover it by finitely many rational balls of
radius at most $\mathsf{d}/16$ whose centres lie within $\mathsf{d}/16$ of
$\operatorname{supp}f_l$, and let $Q_l$ be their union. Then
$\operatorname{supp}f_l\subset Q_l$.
\item Take a rational $\eta\in(0,\mathsf{d}/16]$. A point of $Q_l^{+\eta}$ lies within $\eta$
of a point of one of the balls. That point lies within $\mathsf{d}/16$ of the ball's centre,
and the centre lies within $\mathsf{d}/16$ of $\operatorname{supp}f_l$. By the triangle
inequality for the torus distance, $Q_l^{+\eta}\subset(\operatorname{supp}f_l)^{+3\mathsf{d}/16}$.
Hence, for $l\neq l'$, the distance between $Q_l^{+\eta}$ and $Q_{l'}^{+\eta}$ is at least
$\mathsf{d}-2\cdot3\mathsf{d}/16=5\mathsf{d}/8\ge\mathsf{d}/2$.
\item By the density of $D_0$ there are $h_l^{(j)}\in D_0$ with $h_l^{(j)}\to f_l$ in $C^1$
as $j\to\infty$. By Step~1, applied with $Q=Q_l$, the functions
$g^{(j)}_l:=\chi_{Q_l,\eta}h^{(j)}_l\in D$ are supported in $Q_l^{+\eta}$ and satisfy
$g^{(j)}_l\to f_l$ in $C^1$.
\end{itemize}
For each $j$ the tuple $\vec g^{(j)}:=(g^{(j)}_1,\dots,g^{(j)}_n)$ lies in $D^n$, and its
supports are pairwise at distance at least $\mathsf{d}/2>0$. So
$\bigl(S^T_{n;K_{i_k},m}(\vec g^{(j)})\bigr)_k$ converges by the choice of $(i_k)$.

Apply \eqref{7.1:eq-torus-limits-a} with $A_l=Q_l^{+\eta}$ and $\mathsf{d}'=\mathsf{d}/2$. This
is legitimate: $f_l$ is supported in $Q_l\subset Q_l^{+\eta}$, $g^{(j)}_l$ is supported in
$Q_l^{+\eta}$, and all these functions are real and $C^\infty$, hence Lipschitz. Then bound each
$\|\cdot\|_\sharp$ by \eqref{7.1:eq-torus-limits-2}. The result is that, for every $i$,
\[
\bigl|S^T_{n;K_i,m}(\vec f)-S^T_{n;K_i,m}(\vec g^{(j)})\bigr|\le\delta_j,
\]
where
\[
\delta_j:=C_n(\mathsf{d}/2)\,c'^{\,n}\sum_{l=1}^n\|f_l-g^{(j)}_l\|_{C^1}
\prod_{l'\neq l}\bigl(\|f_{l'}\|_{C^1}+\|g^{(j)}_{l'}\|_{C^1}\bigr).
\]
The norms $\|g^{(j)}_{l'}\|_{C^1}$ converge to $\|f_{l'}\|_{C^1}$ and are therefore bounded in
$j$, while $\|f_l-g^{(j)}_l\|_{C^1}\to0$. Hence $\delta_j\to0$.

Let $\epsilon>0$. Choose $j$ with $\delta_j<\epsilon/3$. Then choose $k_0$ such that
\[
\bigl|S^T_{n;K_{i_k},m}(\vec g^{(j)})-S^T_{n;K_{i_{k'}},m}(\vec g^{(j)})\bigr|<\epsilon/3
\qquad\text{for }k,k'\ge k_0 .
\]
For $k,k'\ge k_0$ we compare $\vec f$ with $\vec g^{(j)}$ at $i_k$, pass from $i_k$ to
$i_{k'}$ at $\vec g^{(j)}$, and compare back at $i_{k'}$. This gives
\[
\bigl|S^T_{n;K_{i_k},m}(\vec f)-S^T_{n;K_{i_{k'}},m}(\vec f)\bigr|
<\delta_j+\epsilon/3+\delta_j<\epsilon .
\]
So $\bigl(S^T_{n;K_{i_k},m}(\vec f)\bigr)_k$ is a Cauchy sequence, and we call its limit
$S^T_{n;m}(\vec f)$. This proves the convergence asserted in (a).

Let $\vec f,\vec f'\in\mathcal{C}^{(n)}_{\mathrm{sep}}(\mathbb{T}_m)$ differ only in the $l$-th
entry, and let $a,b$ be real. Then the support of $af_l+bf'_l$ lies in
$\operatorname{supp}f_l\cup\operatorname{supp}f'_l$, which is at positive distance from the
supports of the other entries. So the tuple with $l$-th entry $af_l+bf'_l$ again belongs to
$\mathcal{C}^{(n)}_{\mathrm{sep}}(\mathbb{T}_m)$. Its value under $S^T_{n;K_{i_k},m}$ is
$a$ times the value at $\vec f$ plus $b$ times the value at $\vec f'$. Letting $k\to\infty$
shows that $S^T_{n;m}$ is multilinear.

If the supports of $\vec f$ are pairwise at distance at least $\mathsf{d}$, then
\[
|S^T_{n;K_{i_k},m}(\vec f)|\le C_n(\mathsf{d})\prod_l\|f_l\|_\sharp
\qquad\text{for every }k
\]
by (B). The
bound in (a) follows in the limit. In the same way \eqref{7.1:eq-torus-limits-a}, which holds
along $(i_k)$ uniformly in $k$, passes to the limit. It holds for $S^T_{n;m}$ whenever $\vec f$
and $\vec g$ belong to $\mathcal{C}^{(n)}_{\mathrm{sep}}(\mathbb{T}_m)$ and satisfy the support
hypotheses of Step~2. This proves (a).

Clauses (b), (c) and (d) are proved below for the subsequence $(i_k)$ fixed in this proof.
\end{proof}

\emph{Distributions; proof of clause (b).}\label{7.1:prf-limit-distributions}
We continue the proof of Proposition~\ref{7.1:prop-torus-limits}, under hypothesis (B) and
along the subsequence $(i_k)$ for which clause (a) has been established above. Throughout,
$\|f\|_{C^1}=\max\{\|f\|_\infty,\|\nabla f\|_\infty\}$, and $c'=\max\{1,c_\sharp\}$ is the
constant with $\|f\|_\sharp\le c'\|f\|_{C^1}$ used above.

Let $U_1,\dots,U_n\subset\mathbb{T}_m$ be open sets of diameter less than $L^m/2$, pairwise at
distance at least $\mathsf{d}>0$. Each $U_l$ is isometric to an open subset of $\mathbb{R}^4$.
We extend $S^T_{n;m}$ to complex test functions by multilinearity. Let $f_l\in C^\infty_c(U_l)$
be complex and write $f_l=u_l+iv_l$ with $u_l,v_l$ real. Then $u_l,v_l\in C^\infty_c(U_l)$, and
\begin{equation*}
\|u_l\|_{C^1}\le\|f_l\|_{C^1},\qquad \|v_l\|_{C^1}\le\|f_l\|_{C^1}.
\end{equation*}
Expanding $S^T_{n;m}(f_1,\dots,f_n)$ by multilinearity gives $2^n$ terms. Each term is $\pm1$
or $\pm i$ times the value of $S^T_{n;m}$ on a tuple of real $C^\infty$ functions whose $l$-th
entry is supported in $U_l$. The supports of such a tuple are pairwise at distance at least
$\mathsf{d}$, so the tuple lies in $\mathcal{C}^{(n)}_{\mathrm{sep}}(\mathbb{T}_m)$. By the bound
in clause (a) and $\|\cdot\|_\sharp\le c'\|\cdot\|_{C^1}$, each term is at most
$c'^{\,n}C_n(\mathsf{d})\prod_l\|f_l\|_{C^1}$ in modulus. Summing the $2^n$ terms,
\begin{equation}\label{7.1:eq-limit-distributions-1}
|S^T_{n;m}(\vec f)|\le(2c')^nC_n(\mathsf{d})\prod_{l=1}^n\|f_l\|_{C^1},\qquad
f_l\in C^\infty_c(U_l).
\end{equation}
By \eqref{7.1:eq-limit-distributions-1}, the $n$-linear form $\vec f\mapsto S^T_{n;m}(\vec f)$
is continuous on $C^\infty_c(U_1)\times\cdots\times C^\infty_c(U_n)$. We apply the Schwartz
kernel theorem $n-1$ times, on the open subsets of $\mathbb{R}^4$ to which the $U_l$ are
isometric. This yields a unique distribution $T_{\vec U}$ on $\prod_lU_l$ with
\begin{equation*}
T_{\vec U}(f_1\otimes\cdots\otimes f_n)=S^T_{n;m}(\vec f),\qquad f_l\in C^\infty_c(U_l).
\end{equation*}

Finite sums of tensors $f_1\otimes\cdots\otimes f_n$ are dense in
$C^\infty_c(\prod_lU_l)$. Hence a distribution on such a product is determined by its values on
tensors.

Let $\vec U$ and $\vec U'$ be two families as above, with separations $\mathsf{d}$ and
$\mathsf{d}'$. Then
\begin{equation*}
\prod_lU_l\cap\prod_lU'_l=\prod_l(U_l\cap U'_l).
\end{equation*}
If this intersection is empty, there is nothing to check. Otherwise the sets $U_l\cap U'_l$ are
open, of diameter less than $L^m/2$, and pairwise at distance at least
$\max\{\mathsf{d},\mathsf{d}'\}$. So the intersection is a product of the same kind. On tensors
with $f_l\in C^\infty_c(U_l\cap U'_l)$, the restrictions of $T_{\vec U}$ and $T_{\vec U'}$ both
take the value $S^T_{n;m}(\vec f)$. By the density just noted, the restrictions coincide. Hence
the distributions $T_{\vec U}$ agree on overlaps.

These products cover $(\mathbb{T}_m)^n_{\neq}$. Indeed, let $x=(x_1,\dots,x_n)$ be a point with
\begin{equation*}
\delta:=\min_{i\neq j}|x_i-x_j|>0,
\end{equation*}
and take open balls about the $x_l$ of a common radius $r<\min\{\delta/3,L^m/4\}$. Each ball has
diameter at most $2r<L^m/2$. Two of the balls are at distance at least $\delta-2r>\delta/3>0$.
So their product is a product of the kind above, and it contains $x$.

The local distributions $T_{\vec U}$ agree on overlaps. Therefore they glue to a unique
distribution on $(\mathbb{T}_m)^n_{\neq}$, which we again denote $S^T_{n;m}$. Every
distribution has finite order on each relatively compact open set. This gives locally finite
order.

It remains to prove the tensor identity for an arbitrary
$\vec f\in\mathcal{C}^{(n)}_{\mathrm{sep}}(\mathbb{T}_m)$, whose supports may be large. Let the
supports be pairwise at distance at least $\mathsf{d}$.
\begin{itemize}
\item By a smooth real partition of unity, write each $f_l$ as a finite sum
$f_l=\sum_af_{l,a}$ of real $C^\infty$ functions.
\item Choose the pieces so that each $f_{l,a}$ lies in $C^\infty_c(U_{l,a})$. Here $U_{l,a}$ is
open, has diameter less than $\min\{\mathsf{d}/4,\,L^m/2\}$, and is contained in
$(\operatorname{supp}f_l)^{+\mathsf{d}/4}$.
\item Then for every choice $a_1,\dots,a_n$ the sets $U_{1,a_1},\dots,U_{n,a_n}$ are pairwise at
distance at least $\mathsf{d}/2$. Indeed, for $l\neq l'$,
\begin{equation*}
\operatorname{dist}(U_{l,a_l},U_{l',a_{l'}})
\ge\operatorname{dist}(\operatorname{supp}f_l,\operatorname{supp}f_{l'})
-\frac{\mathsf{d}}{4}-\frac{\mathsf{d}}{4}\ge\frac{\mathsf{d}}{2}.
\end{equation*}
\end{itemize}
Consequently each $\prod_lU_{l,a_l}$ is a product of the kind considered above, with separation
$\mathsf{d}/2$. On it the glued distribution $S^T_{n;m}$ coincides with $T_{\vec U}$ for
$\vec U=(U_{1,a_1},\dots,U_{n,a_n})$. Moreover,
\begin{equation*}
f_1\otimes\cdots\otimes f_n=\sum_{a_1,\dots,a_n}f_{1,a_1}\otimes\cdots\otimes f_{n,a_n}.
\end{equation*}
Each tuple $(f_{1,a_1},\dots,f_{n,a_n})$ consists of real $C^\infty$ functions with supports
pairwise at distance at least $\mathsf{d}/2$, so it lies in
$\mathcal{C}^{(n)}_{\mathrm{sep}}(\mathbb{T}_m)$. We now use three facts in turn: linearity of
the distribution; the construction of $T_{\vec U}$ on each product $\prod_lU_{l,a_l}$; and the
multilinearity in clause (a). Together they give
\begin{equation}\label{7.1:eq-limit-distributions-2}
\begin{split}
S^T_{n;m}(f_1\otimes\cdots\otimes f_n)
&=\sum_{a_1,\dots,a_n}S^T_{n;m}(f_{1,a_1},\dots,f_{n,a_n})\\
&=S^T_{n;m}(\vec f).
\end{split}
\end{equation}
Uniqueness of the distribution follows from the agreement with each $T_{\vec U}$ on the cover of
$(\mathbb{T}_m)^n_{\neq}$ by the products above, together with the uniqueness of each
$T_{\vec U}$. This proves clause (b).

\emph{Symmetry and invariance; proof of clause~(c) of
Proposition~\ref{7.1:prop-torus-limits}.}\label{7.1:prf-limit-invariance}
We continue the proof of Proposition~\ref{7.1:prop-torus-limits} under hypothesis~(B), along
the subsequence $(i_k)$ of clause~(a).

At finite $K$ the partition function $Z(z)$ with sources is entire in $z$, and $Z(0)>0$. Hence
$\log Z$ is analytic in a neighbourhood of $z=0$, and its mixed derivatives at $z=0$ commute.
The connected correlations $S^T_{n;K,m}$ of Notation~\ref{7.1:not-torus-cumulants} are these
mixed derivatives, so $S^T_{n;K,m}(f_1,\dots,f_n)$ is a symmetric function of
$f_1,\dots,f_n$.

Let $w\in W_4$. It maps the bare sites $L^{-K}(\mathbb{Z}+\tfrac12)^4$ onto themselves, and the
plaquettes onto themselves, with $x(wp)=w\,x(p)$. For a configuration $U$ write $U\circ w$ for the
transported configuration, and for an observable $F$ put $\alpha_wF:=F\circ(U\mapsto U\circ w)$.
By the transport identity for the plaquette function under $W_4$ we have
$a_p(U\circ w)=a_{wp}(U)$. Substituting $q=wp$ in the sum over plaquettes,
\begin{equation*}
\mathcal{O}(f)(U\circ w)=\sum_p f(x(p))\,a_{wp}(U)
=\sum_q f\bigl(w^{-1}x(q)\bigr)\,a_q(U)=\mathcal{O}(f\circ w^{-1})(U),
\end{equation*}
that is, $\mathcal{O}(f\circ w^{-1})=\alpha_w\mathcal{O}(f)$. The measure $\mu_{K,m}$ is invariant
under $U\mapsto U\circ w$. Therefore the family $(\mathcal{O}(f_l\circ w^{-1}))_{l}$ has under
$\mu_{K,m}$ the same joint law as $(\mathcal{O}(f_l))_{l}$, and, writing
$\vec f\circ w^{-1}:=(f_1\circ w^{-1},\dots,f_n\circ w^{-1})$,
\begin{equation*}
S^T_{n;K,m}(\vec f\circ w^{-1})=S^T_{n;K,m}(\vec f).
\end{equation*}
The same holds for the translations $\tau_a$ with $a\in L^{-K}\mathbb{Z}^4$, which map the bare
lattice at cutoff $K$ onto itself. Since $L$ is an integer, $L^{-K}\mathbb{Z}^4\subset
L^{-K'}\mathbb{Z}^4$ for every $K'\ge K$, so for such $a$ the identity
$S^T_{n;K',m}(\tau_a\vec f)=S^T_{n;K',m}(\vec f)$ holds at every cutoff $K'\ge K$.

All these identities pass to the limit. Indeed, $w$ and $\tau_a$ are isometries of
$\mathbb{T}_m$, so $\vec f\in\mathcal{C}^{(n)}_{\mathrm{sep}}(\mathbb{T}_m)$ implies that
$\vec f\circ w^{-1}$, $\tau_a\vec f$ and every permutation of $\vec f$ lie in
$\mathcal{C}^{(n)}_{\mathrm{sep}}(\mathbb{T}_m)$. Clause~(a) then gives convergence of both
sides along $(i_k)$. For a translation by $a\in L^{-K}\mathbb{Z}^4$ we use in addition that
$K_{i_k}\ge K$ for all large $k$, because $K_{i_k}\to\infty$. Thus $S^T_{n;m}$ is symmetric and
satisfies $S^T_{n;m}(\vec f\circ w^{-1})=S^T_{n;m}(\vec f)$ for $w\in W_4$.

Consider the group
\begin{equation*}
G_m:=\bigcup_k L^{-K_{i_k}}\mathbb{Z}^4/2L^m\mathbb{Z}^4 .
\end{equation*}
Any two members of this union are comparable by inclusion, by the inclusion
$L^{-K}\mathbb{Z}^4\subset L^{-K'}\mathbb{Z}^4$ for $K\le K'$. Hence the union is nested and
$G_m$ is a subgroup of $\mathbb{T}_m$. Since $K_{i_k}\to\infty$, the mesh $L^{-K_{i_k}}$ tends to
zero, and $G_m$ is dense in $\mathbb{T}_m$. Every $a\in G_m$ lies in
$L^{-K_{i_{k_0}}}\mathbb{Z}^4/2L^m\mathbb{Z}^4$ for some $k_0$. By the preceding paragraph,
$S^T_{n;m}(\tau_a\vec f)=S^T_{n;m}(\vec f)$ for $a\in G_m$.

Fix $\vec f\in\mathcal{C}^{(n)}_{\mathrm{sep}}(\mathbb{T}_m)$ with supports pairwise at distance
at least $\mathsf{d}>0$, and set $\varphi(b):=S^T_{n;m}(\tau_b\vec f)$ for
$b\in\mathbb{R}^4/2L^m\mathbb{Z}^4$. Applying the identity of the previous paragraph to the tuple
$\tau_b\vec f$ gives $\varphi(b+a)=\varphi(b)$ for all $a\in G_m$.

Let $|b-b'|<\mathsf{d}/4$. The functions $\tau_bf_l$ and $\tau_{b'}f_l$ are supported in
$A_l:=(\operatorname{supp}f_l+b)^{+|b-b'|}$, the closed $|b-b'|$-neighbourhood of
$\operatorname{supp}f_l+b$. Translation preserves torus distances, so these sets are pairwise at
distance at least $\mathsf{d}-2|b-b'|\ge\mathsf{d}/2$. Moreover, for each $f=f_l$,
\begin{equation*}
\|\tau_bf-\tau_{b'}f\|_\infty\le\|\nabla f\|_\infty|b-b'|,\qquad
\operatorname{Lip}(\tau_bf-\tau_{b'}f)\le\|\nabla^2f\|_\infty|b-b'|.
\end{equation*}
Hence
\begin{equation*}
\|\tau_bf_l-\tau_{b'}f_l\|_\sharp
\le\max\{\|\nabla f_l\|_\infty,\,40M\|\nabla^2f_l\|_\infty\}\,|b-b'| .
\end{equation*}
This is where test functions of class $C^2$ are used. Also $\|\tau_bf_l\|_\sharp=\|f_l\|_\sharp$.

Now apply the equicontinuity estimate \eqref{7.1:eq-torus-limits-a} with $\mathsf{d}'=\mathsf{d}/2$
to $\tau_b\vec f$ and $\tau_{b'}\vec f$. It holds uniformly in $i$, and both tuples lie in
$\mathcal{C}^{(n)}_{\mathrm{sep}}(\mathbb{T}_m)$, so its left side converges along $(i_k)$ by
clause~(a). In the limit,
\begin{equation*}
|\varphi(b)-\varphi(b')|\le 2^{n-1}C_n(\mathsf{d}/2)\sum_{l=1}^n
\max\{\|\nabla f_l\|_\infty,40M\|\nabla^2f_l\|_\infty\}
\prod_{l'\neq l}\|f_{l'}\|_\sharp\;|b-b'|.
\end{equation*}
This constant does not depend on $b,b'$. Joining two arbitrary points by a segment cut into pieces
of length less than $\mathsf{d}/4$ and summing shows that $\varphi$ is Lipschitz on $\mathbb{T}_m$.

A continuous function on $\mathbb{T}_m$ invariant under translation by a dense subgroup is
constant. Indeed, given $b$, choose $a_j\in G_m$ with $a_j\to b$; then
$\varphi(b)=\lim_j\varphi(a_j)=\varphi(0)$. Thus
$S^T_{n;m}(\tau_a\vec f)=S^T_{n;m}(\vec f)$ for every $a\in\mathbb{R}^4/2L^m\mathbb{Z}^4$.

The statements for the distribution of clause~(b) follow from the uniqueness asserted there. Let
$\sigma$ be a permutation of $\{1,\dots,n\}$, the diagonal action of $w\in W_4$, or the diagonal
translation by $a$. Each is a diffeomorphism of $(\mathbb{T}_m)^n_{\neq}$ onto itself. Composing
the distribution $S^T_{n;m}$ with it therefore gives again a distribution on
$(\mathbb{T}_m)^n_{\neq}$. On every product $f_1\otimes\cdots\otimes f_n$ with
$\vec f\in\mathcal{C}^{(n)}_{\mathrm{sep}}(\mathbb{T}_m)$ it takes, by the identities just proved,
the value $S^T_{n;m}(\vec f)$. By the uniqueness in clause~(b) it coincides with $S^T_{n;m}$. This
proves clause~(c).

\begin{proof}[Proof of Proposition~\ref{7.1:prop-torus-limits}, continued]
\emph{Reflection positivity; proof of (d).}\label{7.1:prf-limit-reflection}
Write $\varepsilon=L^{-K}$ for the bare spacing and $\langle\cdot\rangle_{K,m}$ for the expectation
with respect to the Wilson measure $\mu_{K,m}$ on $\mathbb{T}_m$. The time coordinates of the bare
sites lie in $(\mathbb{Z}+\frac12)\varepsilon$. Hence the hyperplanes $\{x_0=0\}$ and
$\{x_0=L^m\}$ contain no bare site.

Let $f$ be a function on $\mathbb{T}_m$ with $\operatorname{supp}f\subset\mathbb{T}_m^+$, and let
$p$ be a plaquette with $f(x(p))\neq0$. Then $x(p)\in\mathbb{T}_m^+$, and there are two cases.
\begin{itemize}
\item Suppose $p$ lies in a $(0,j)$-plane. Its corners have times $(k+\frac12)\varepsilon$ and
$(k+\frac32)\varepsilon$ for some integer $k$, and its barycentre has time
$(k+1)\varepsilon\in(0,L^m)$. Since $L^m/\varepsilon=L^{m+K}$ is an integer, $k+1$ is an integer
with $1\le k+1\le L^m/\varepsilon-1$. Hence all corner times of $p$ lie in
$[\varepsilon/2,\,L^m-\varepsilon/2]$.
\item Suppose $p$ is a spatial plaquette. Then all of $p$ lies at the single time
$(k+\frac12)\varepsilon$, which is the time of $x(p)$ and so lies in $(0,L^m)$.
\end{itemize}
In both cases every link of $\partial p$ lies in $\mathbb{T}_m^+$. Therefore $\mathcal{O}(f)$ is a
function of the links in $\mathbb{T}_m^+$, and $\mathcal{O}(f)\circ\theta=\mathcal{O}(f\circ\theta)$.
The Wilson measure is link reflection positive \cite{OSe78}. Hence
\begin{equation*}
\langle\overline{F\circ\theta}\,F\rangle_{K,m}\ge0
\end{equation*}
for every polynomial $F$ in bounded functions of links in $\mathbb{T}_m^+$.

Now let $r$, $f_1,\dots,f_r$ and $c_\alpha$ be as in (d). For a bounded function $X$ of the links
put ${:}X{:}:=X-\langle X\rangle_{K,m}$, and set
\begin{equation*}
F:=\sum_{\alpha\subset\{1,\dots,r\}}c_\alpha\prod_{i\in\alpha}{:}\mathcal{O}(f_i){:}.
\end{equation*}
Each ${:}\mathcal{O}(f_i){:}$ is a bounded function of links in $\mathbb{T}_m^+$ plus a constant,
by the preceding paragraph. So $F$ is a polynomial of the kind just described.

Since $f_i$ is real, $\mathcal{O}(f_i)$ and $\langle\mathcal{O}(f_i)\rangle_{K,m}$ are real. Since
$\mu_{K,m}$ is $\theta$-invariant,
\begin{equation*}
\langle\mathcal{O}(f_i\circ\theta)\rangle_{K,m}=\langle\mathcal{O}(f_i)\circ\theta\rangle_{K,m}
=\langle\mathcal{O}(f_i)\rangle_{K,m}.
\end{equation*}
Together these give
$\overline{{:}\mathcal{O}(f_i){:}\circ\theta}={:}\mathcal{O}(f_i\circ\theta){:}$. Hence, at every
cutoff $K=K_{i_k}$,
\begin{align*}
0&\le\langle\overline{F\circ\theta}\,F\rangle_{K,m}\\
&=\sum_{\alpha,\beta}\overline{c_\alpha}\,c_\beta\,\Bigl\langle\prod_{i\in\alpha}
{:}\mathcal{O}(f_i\circ\theta){:}\prod_{j\in\beta}{:}\mathcal{O}(f_j){:}\Bigr\rangle_{K,m}.
\end{align*}

We next read each term as a moment. The support of $f_i\circ\theta$ is $\theta(\operatorname{supp}f_i)$,
a compact subset of $\{-L^m<x_0<0\}$, while the supports of the $f_j$ are compact subsets of
$\mathbb{T}_m^+$. The $f_i$ have pairwise disjoint supports, and so do the $f_i\circ\theta$. Hence the
functions $f_i\circ\theta$ ($i\in\alpha$) and $f_j$ ($j\in\beta$) have pairwise disjoint compact
supports, and therefore pairwise positive distances.

Fix $\alpha,\beta$, put $q:=|\alpha|+|\beta|$, and let
$(h_1,\dots,h_q):=((f_i\circ\theta)_{i\in\alpha},(f_j)_{j\in\beta})$. We identify the corresponding
term with a moment of the form \eqref{7.1:eq-torus-cumulants-b}, in three steps.
\begin{itemize}
\item By the moment--cumulant formula, the moment
$\langle\prod_{l=1}^q{:}\mathcal{O}(h_l){:}\rangle_{K,m}$ is the sum, over all partitions $\pi$ of
$\{1,\dots,q\}$, of the products over $B\in\pi$ of the joint cumulants of the variables
${:}\mathcal{O}(h_l){:}$, $l\in B$.
\item Cumulants of order one of centred variables vanish. So only partitions in
$\Pi_{\ge2}[q]$ contribute.
\item Joint cumulants of order at least $2$ do not change when constants are added to the variables.
So for $|B|\ge2$ the joint cumulant of the ${:}\mathcal{O}(h_l){:}$, $l\in B$, equals
$S^T_{|B|;K,m}((h_l)_{l\in B})$.
\end{itemize}
Therefore
\begin{equation*}
\Bigl\langle\prod_{l=1}^q{:}\mathcal{O}(h_l){:}\Bigr\rangle_{K,m}
=\sum_{\pi\in\Pi_{\ge2}[q]}\ \prod_{B\in\pi}S^T_{|B|;K,m}\bigl((h_l)_{l\in B}\bigr)
=S_{q;K,m}(h_1,\dots,h_q).
\end{equation*}
So each term of the double sum is $S_{|\alpha|+|\beta|;K,m}$ at pairwise separated supports, as in
\eqref{7.1:eq-torus-cumulants-b}. It is a finite sum of products of cumulants of order at least $2$.
The arguments of each cumulant are real $C^\infty$ functions with pairwise positive distances, so
each cumulant converges along $(i_k)$ by (a). Its limit is the corresponding $S^T_{|B|;m}$, and
$S_{q;m}$ is defined by the same formula with $S^T_{|B|;m}$ in place of $S^T_{|B|;K,m}$.

Consequently, for each $k$ the double sum at $K=K_{i_k}$ is a nonnegative number. It is also a
finite sum of terms, each of which converges along $(i_k)$ to the corresponding term of
\eqref{7.1:eq-torus-limits-b}. The left side of \eqref{7.1:eq-torus-limits-b} is therefore a limit
of nonnegative numbers, hence nonnegative. This proves \eqref{7.1:eq-torus-limits-b}.

Finally, fix $\mu$ and $c\in\mathbb{R}$. Let $\psi$ be the map of $\mathbb{T}_m$ obtained by first
translating by $-c$ in the direction $\mu$ and then applying the element of $W_4$ that exchanges the
coordinate axes $0$ and $\mu$ (the identity if $\mu=0$). Then $\psi$ carries
$\{c<x_\mu<c+L^m\}$ onto $\mathbb{T}_m^+$.

We check that $\psi^{-1}\circ\theta\circ\psi$ is the reflection $x_\mu\mapsto2c-x_\mu$. The map
$\psi$ sends the coordinate $x_\mu$ to the new time coordinate $x_\mu-c$. Then $\theta$ changes its
sign, giving $c-x_\mu$. Finally $\psi^{-1}$ sends this to $c+(c-x_\mu)=2c-x_\mu$. The other
coordinates are left unchanged. Thus the reflection $x_\mu\mapsto2c-x_\mu$ is conjugate to $\theta$
by an element of $W_4$ and a translation.

By (c), every $S^T_{n;m}$ is invariant under both of these maps, and hence so is every $S_{q;m}$.
Now let $f_1,\dots,f_r$ be real $C^\infty$ functions with pairwise disjoint compact supports in
$\{c<x_\mu<c+L^m\}$. Apply \eqref{7.1:eq-torus-limits-b} to the functions $f_i\circ\psi^{-1}$, whose
supports are pairwise disjoint compact subsets of $\mathbb{T}_m^+$. Then use this invariance to
compose all arguments with $\psi$. The result is \eqref{7.1:eq-torus-limits-b} with $\theta$ replaced
by the reflection $x_\mu\mapsto2c-x_\mu$.
\end{proof}

\begin{remark}[Dependence of the limits on subsequence and torus]\label{7.1:rem-subsequence-dependence}
The subsequence $(i_k)$ of Proposition~\ref{7.1:prop-torus-limits} depends on $m$ and on the
sequence of bare couplings $g_0^{(i)}$.

Convergence without passing to subsequences at the tuned bare couplings
\begin{equation*}
g_0^{-2}=a_n-\hat w
\end{equation*}
is the subject of Theorem~\ref{7.3:thm-tuned-family}, which is clause~(ii-$\mathbb{T}$-a) of
Theorem~0 (\ref{3.1:thm-main}). For $\varrho\ge 6(B_\sharp+1)$ and $g$ below a further
threshold, the identification of every subsequential limit obtained in
Proposition~\ref{7.1:prop-torus-limits} with some $S^{(\hat w)}$ belongs to clause~(v) of
Theorem~0, the uniqueness clause. The condition $\varrho\ge 6(B_\sharp+1)$ is also a hypothesis
of Lemma~\ref{7.3:lem-first-passage-offsets} on the non-convergence of the first-passage offsets.
For a general sequence of bare couplings, no independence of the limit from the
subsequence is asserted.

The constant $C_n(\mathsf{d})$ of hypothesis~(B) of Proposition~\ref{7.1:prop-torus-limits}
contains the volume factor
\begin{equation*}
Q_m=\exp\bigl[(E^\sharp_++C_{\mathrm{low}})(2L^m)^4\bigr]
\end{equation*}
of the torus bound. Here $E^\sharp_+$ is the constant of the upper global stability bound in
the presence of sources, and $C_{\mathrm{low}}$ is the constant of the lower bound
$\tilde Z(0)\ge e^{-C_{\mathrm{low}}(2L^m)^4}$ at zero source. The function
$\tilde Z=e^{-E-\mathcal{E}}Z$ is the integral over the last lattice of the density that remains
when the normalization constant $E$ of Ba{\l}aban's bare density, which does not depend on the
sources, and the sum $\mathcal{E}$ of the source-dependent vacuum energies have been removed
(these are the objects entering the torus bound of the correlation theorem).
Because of this factor,
Proposition~\ref{7.1:prop-torus-limits} gives nothing uniform in $m$.

Hypothesis~(B) is the torus bound of the correlation theorem. The correlation theorem is proved
as an implication from the four hypotheses listed in its statement. Consequently, the limits
$S^T_{n;m}$ of Proposition~\ref{7.1:prop-torus-limits} are proved as an implication from the
same four statements.
\end{remark}

\begin{proposition}[Reflection positivity of the torus Schwinger functions]
\label{7.2:prop-torus-reflection}
Fix integers $K, m \ge 0$ and $g_0 > 0$, put $\varepsilon := L^{-K}$ and $x_0 := g_0^{-2}$, and let
$\langle\,\cdot\,\rangle$ be the expectation in the bare Wilson measure, the probability measure
proportional to $e^{-x_0 A(U)}\,dU$ on the bare lattice of $\mathbb{T}_m$, whose sites are
$\varepsilon(\mathbb{Z}+\tfrac12)^4$ modulo $2L^m$.
Let $c \in \varepsilon\mathbb{Z}$ and let $\mu$ be a coordinate direction.
Let $\theta$ be the reflection $x_\mu \mapsto 2c - x_\mu$, the other coordinates being fixed, and put
\begin{equation*}
\Pi := \{x_\mu = c\} \cup \{x_\mu = c + L^m\}, \qquad
H_{+} := \{c < x_\mu < c + L^m\}, \qquad H_{-} := \theta(H_{+}).
\end{equation*}
For every bare bond $b$ let $(\theta\cdot U)(b) := U(\theta b)$ if $\theta$ preserves the
orientation of $b$ (bonds perpendicular to $e_{\mu}$), and $(\theta\cdot U)(b) := U(\theta b)^{-1}$
otherwise (bonds parallel to $e_{\mu}$).

Assume clause (ii) of the lemma on reflection positivity of the lattice measure (in origin the
link reflection positivity of Osterwalder and Seiler \cite{OSe78}) for the $c$ and $\mu$ fixed
above, namely: for every bounded measurable function $F$ of the bond variables in $H_{+}$,
\begin{equation}\label{7.2:eq-torus-reflection-1}
\langle \overline{F(\theta\cdot U)}\, F(U) \rangle \ge 0 .
\end{equation}
Then for every $n \ge 1$, all real Lipschitz functions $h_1, \dots, h_n$ on $\mathbb{T}_m$ with
$\operatorname{supp} h_i \subset H_{+}$, and every polynomial
$P(y) = \sum_{\eta} \xi_\eta\, y^\eta$ in $y = (y_1, \dots, y_n)$ with complex coefficients $\xi_\eta$
(multi-indices $\eta \in \mathbb{N}^n$, $y^\eta := \prod_i y_i^{\eta_i}$),
\begin{equation}\label{7.2:eq-torus-reflection-2}
\sum_{\eta, \eta'} \overline{\xi_\eta}\, \xi_{\eta'}
\Big\langle \prod_{i=1}^n \mathcal{O}(h_i\circ\theta)^{\eta_i}
\prod_{j=1}^n \mathcal{O}(h_j)^{\eta'_j} \Big\rangle \ge 0 .
\end{equation}
\end{proposition}

\begin{proof}
\emph{Geometry of $\theta$.} Since $L^m = L^{K+m}\varepsilon \in \varepsilon\mathbb{Z}$ and
$c \in \varepsilon\mathbb{Z}$, the map $x_\mu \mapsto 2c - x_\mu$ maps $\varepsilon(\mathbb{Z}+\frac12)$
onto itself and is compatible with reduction modulo $2L^m$; so $\theta$ is an isometry of
$\mathbb{T}_m$ preserving the bare lattice, and it maps bonds to bonds and plaquettes to plaquettes.
It fixes both planes of $\Pi$: $2c - c = c$, and
\begin{equation*}
2c - (c + L^m) = c - L^m \equiv c + L^m \pmod{2L^m}.
\end{equation*}
It maps the open slab $H_{+}$ onto $\{c - L^m < x_\mu < c\}$, which modulo $2L^m$ is the other
open component of $\mathbb{T}_m \setminus \Pi$; thus $H_{-}$ is that component. No bare site lies on
$\Pi$, because the site coordinates lie in $\varepsilon(\mathbb{Z}+\frac12)$ while $c$ and $c + L^m$ lie
in $\varepsilon\mathbb{Z}$.

\emph{The observables read only bonds in $H_{+}$.} Let $h$ be real with
$\operatorname{supp} h \subset H_{+}$. In
\begin{equation*}
\mathcal{O}(h) = \sum_p h(x(p))\, a_p(U), \qquad
a_p(U) := 1 - \operatorname{Re}\operatorname{tr} U(\partial p) \in [0,2],
\end{equation*}
with $x(p)$ the barycentre of the plaquette $p$, only plaquettes with $x(p) \in H_{+}$ contribute. Such a
plaquette lies in $H_{+}$. If it is perpendicular to $e_{\mu}$, all its points have the
$\mu$-coordinate of its sites, which equals that of $x(p)$. If it lies in a $(\mu,\nu)$-plane, then
$a := x(p)_\mu \in \varepsilon\mathbb{Z}$, and its points have $\mu$-coordinates in
$[a - \varepsilon/2,\, a + \varepsilon/2]$. Since $c < a < c + L^m$ with all three numbers in
$\varepsilon\mathbb{Z}$, we have $c + \varepsilon \le a \le c + L^m - \varepsilon$, and the interval
lies in $H_{+}$. Hence $\mathcal{O}(h)$ is a function of the bond variables in $H_{+}$. It is bounded,
since $0 \le a_p \le 2$ and the lattice is finite, and it is continuous.

\emph{Reflection of the observables.} For a plaquette $p \subset H_{+}$, evaluate the holonomy
$U(\partial p)$ on $\theta\cdot U$. It is the ordered product of the variables $(\theta\cdot U)(b)$
over the bonds $b$ of $\partial p$. By the definition of $\theta\cdot U$, each factor is the variable of
the bond $\theta b$, traversed in the direction that is the image under $\theta$ of the direction of
$b$. So the product is the holonomy of $U$ around the boundary of $\theta p$, traversed in one of its
two orientations and from one of its corners. It therefore equals $U(\partial(\theta p))$ or its
inverse, up to cyclic conjugation. Now $\operatorname{Re}\operatorname{tr}$ is invariant under
conjugation, and also under inversion, because
$\operatorname{tr}(W^{-1}) = \operatorname{tr}(W^{*}) = \overline{\operatorname{tr} W}$ for unitary $W$.
Hence $a_p(\theta\cdot U) = a_{\theta p}(U)$. Since $x(\theta p) = \theta x(p)$ and $\theta$ is a
bijection of the plaquettes,
\begin{equation}\label{7.2:eq-torus-reflection-3}
\mathcal{O}(h)(\theta\cdot U)
= \sum_{p} h(x(p))\, a_{\theta p}(U)
= \sum_{q} h(\theta x(q))\, a_q(U)
= \mathcal{O}(h\circ\theta)(U),
\end{equation}
where $\theta^{-1} = \theta$ was used. Note that
$\operatorname{supp}(h\circ\theta) = \theta(\operatorname{supp} h) \subset H_{-}$.

\emph{Conclusion.} Put
\begin{equation*}
F(U) := P\big(\mathcal{O}(h_1)(U), \dots, \mathcal{O}(h_n)(U)\big).
\end{equation*}
By the second paragraph,
$F$ is a bounded continuous, hence measurable, function of the bond variables in $H_{+}$. By
\eqref{7.2:eq-torus-reflection-3},
\begin{equation*}
F(\theta\cdot U) = P\big(\mathcal{O}(h_1\circ\theta)(U), \dots, \mathcal{O}(h_n\circ\theta)(U)\big).
\end{equation*}
Because the $h_i$ are real, the functions $\mathcal{O}(h_i\circ\theta)$ are real, and so
\begin{equation*}
\overline{F(\theta\cdot U)}\, F(U)
= \sum_{\eta, \eta'} \overline{\xi_\eta}\, \xi_{\eta'}
\prod_{i} \mathcal{O}(h_i\circ\theta)(U)^{\eta_i}
\prod_{j} \mathcal{O}(h_j)(U)^{\eta'_j}.
\end{equation*}
The sum is finite. Taking expectations and applying the hypothesis
\eqref{7.2:eq-torus-reflection-1} gives \eqref{7.2:eq-torus-reflection-2}.
\end{proof}

The expectations in \eqref{7.2:eq-torus-reflection-2} are moments of the smeared curvature
observables $\mathcal{O}(h)$. By the moment--cumulant formula, each is a finite sum of products of
joint cumulants $\kappa_q$, $q \ge 1$, of observables among the $\mathcal{O}(h_i)$ and
$\mathcal{O}(h_i\circ\theta)$, a test function possibly occurring more than once; the cumulants of
order one are the expectations $\langle \mathcal{O}(h) \rangle$. A cumulant of order $q \ge 2$ whose
$q$ test functions have pairwise separated supports is the connected Schwinger function
$S^T_{q;K,m}$ of those test functions.

On the periodic torus $\mathbb{T}_m$, the reflection $x_\mu \mapsto 2c - x_\mu$ fixes the pair of planes
$\{x_\mu = c\}$ and $\{x_\mu = c + L^m\}$. Reflection positivity ``at the hyperplane
$x_\mu = c$'' on $\mathbb{T}_m$ refers to this reflection and to test functions supported in the
open slab $c < x_\mu < c + L^m$. Every
$c \in \varepsilon\mathbb{Z}$ and every direction $\mu$ are admitted. In particular, every
$c \in L^{-k}\mathbb{Z}$ with $0 \le k \le K$ is admitted, since $L^{-k}\mathbb{Z} \subset \varepsilon\mathbb{Z}$.

\begin{corollary}[The reflection form and its Cauchy--Schwarz inequality]\label{7.2:cor-reflection-form}
Work in the setting of the lemma on reflection positivity of the lattice measure. There $\theta$ is the
reflection of the torus $\mathbb{T}$ in a pair of antipodal link planes, $H_{+}$ and $H_{-}$ are the two
open components of their complement, and $\langle\,\cdot\,\rangle$ is the expectation in the bare
Wilson measure. For a bare bond $b$, put $(\theta\cdot U)(b):=U(\theta b)$ if $\theta$ preserves the
orientation of $b$, and $(\theta\cdot U)(b):=U(\theta b)^{-1}$ otherwise. For bounded measurable
functions $F,G$ of the bond variables in $H_{+}$ put
\begin{equation}\label{7.2:eq-reflection-form-1}
(F,G)_{\theta}:=\big\langle \overline{F(\theta\cdot U)}\,G(U)\big\rangle .
\end{equation}
Then $(\cdot,\cdot)_{\theta}$ is a Hermitian, positive-semidefinite sesquilinear form on these functions,
and
\begin{equation}\label{7.2:eq-reflection-form-2}
|(F,G)_{\theta}|^{2}\le (F,F)_{\theta}\,(G,G)_{\theta}.
\end{equation}
Moreover, if $F$ is a bounded measurable function of the bond variables in $H_{-}$, then
$\langle\overline{F(\theta\cdot U)}\,F(U)\rangle\ge 0$.
\end{corollary}

\begin{proof}
The expectation is linear. Hence \eqref{7.2:eq-reflection-form-1} is conjugate-linear in $F$ and
linear in $G$, that is, sesquilinear.

Hermitian symmetry. The complex conjugate of an expectation is the expectation of the complex conjugate,
so
\begin{equation*}
\overline{(F,G)_{\theta}}=\big\langle F(\theta\cdot U)\,\overline{G(U)}\big\rangle .
\end{equation*}
By clause (i) of the lemma on reflection positivity of the lattice measure, the measure is invariant
under every isometry of the torus that preserves the bare lattice. Apply this to $\theta$: the
substitution $U\mapsto\theta\cdot U$ preserves the measure. Since $\theta$ is an involution, two
inversions cancel, and $\theta\cdot(\theta\cdot U)=U$. The substitution therefore turns the right-hand
side into $\langle F(U)\,\overline{G(\theta\cdot U)}\rangle$, which is $(G,F)_{\theta}$ by
\eqref{7.2:eq-reflection-form-1}. In one line,
\begin{equation*}
\overline{(F,G)_{\theta}}
=\big\langle F(\theta\cdot U)\,\overline{G(U)}\big\rangle
=\big\langle F(U)\,\overline{G(\theta\cdot U)}\big\rangle
=(G,F)_{\theta}.
\end{equation*}

Positivity. The inequality $(F,F)_{\theta}\ge 0$ is clause (ii) of the same lemma.

Cauchy--Schwarz. A Hermitian, positive-semidefinite sesquilinear form satisfies the Cauchy--Schwarz
inequality. By the three properties just established, this gives \eqref{7.2:eq-reflection-form-2}.

Functions of $H_{-}$. Let $F$ read bonds of $H_{-}$ and put $G(U):=\overline{F(\theta\cdot U)}$.
For a bond $b$ of $H_{-}$, the variable $(\theta\cdot U)(b)$ is $U(\theta b)$ or its inverse, and
$\theta b$ lies in $H_{+}$. Hence $G$ is a bounded measurable function of the bond variables in $H_{+}$.
Using $\theta\cdot(\theta\cdot U)=U$ once more,
\begin{equation*}
\overline{G(\theta\cdot U)}=F\big(\theta\cdot(\theta\cdot U)\big)=F(U).
\end{equation*}
Therefore
\begin{equation*}
\big\langle\overline{F(\theta\cdot U)}\,F(U)\big\rangle
=\big\langle G(U)\,\overline{G(\theta\cdot U)}\big\rangle
=(G,G)_{\theta}\ge 0,
\end{equation*}
where the middle equality is \eqref{7.2:eq-reflection-form-1} with both arguments equal to $G$, and the
inequality is clause (ii) of the lemma applied to $G$.
\end{proof}

\begin{theorem}[Convergence of the curvature Schwinger functions along the tuned family on a fixed
torus (clause (ii-$\mathbb{T}$-a))]\label{7.3:thm-tuned-family}
This is clause (ii-$\mathbb{T}$-a) of Theorem~0 (\ref{3.1:thm-main}). The tuned bare inverse
couplings $a_n$, the complex offset $\hat w$, the width $\upsilon_w$, the reference length $\ell_g$
and the bare and renormalised labellings are those of clause (v) of Theorem~0, and they are
constructed in Chapter~\ref{ch:9}.

Assume the hypotheses of Theorem~0 and its quantifier prefix. Take the reference coupling $g$ in
the range in which clause (v) of Theorem~0 constructs the tuned couplings $a_n$. Fix $m_0\ge 0$.
Let $\mathbb{T}$ be the torus of side $2L^{m_0}\ell_g$, with hypothesis (H5) taken for this
torus. Put $\varrho_2:=\varrho/2$. Let $\mathsf{p}\ge 2$. Let $f_1,\dots,f_{\mathsf{p}}$ be real
Lipschitz functions on $\mathbb{T}$ whose supports are pairwise at distance at least $\mathsf{d}$,
for some $\mathsf{d}>0$. Write $\vec f=(f_1,\dots,f_{\mathsf{p}})$. For a bare inverse coupling
$x$, real or complex, denote by $S^n(x;\vec f)$ the connected $\mathsf{p}$-point function of
$\mathcal{O}(f_1),\dots,\mathcal{O}(f_{\mathsf{p}})$ at bare inverse coupling $x$. It is taken
on the bare lattice of $\mathbb{T}$ lying $n$ levels below the reference scale.

Let $\hat w\in\mathbb{C}$ satisfy
\begin{equation*}
|\operatorname{Re}\hat w|\le \varrho_2/2,\qquad |\operatorname{Im}\hat w|<\upsilon_w .
\end{equation*}
Then there is a limit $S^{(\hat w)}(\vec f)$ such that
\begin{equation}\label{7.3:eq-tuned-family-1}
S^n(a_n-\hat w;\vec f)\longrightarrow S^{(\hat w)}(\vec f)\qquad (n\to\infty).
\end{equation}
The convergence \eqref{7.3:eq-tuned-family-1} holds along the full sequence, without passage to
subsequences. It is uniform in $\hat w$ in the set just described. Its rate is
$n^{-1/2}(\log n)^{p_0+1}$
in the bare labelling, where $p_0$ is the fixed exponent of the degree function
$p_0(\cdot)$. In the renormalised labelling of clause (v-d) of Theorem~0 the convergence is
exponential.

For real $\hat w$ the limits $S^{(\hat w)}$ have the following properties.
\begin{itemize}
\item They are multilinear and symmetric in $f_1,\dots,f_{\mathsf{p}}$.
\item They are invariant under all translations of $\mathbb{T}$, for $C^1$ test functions.
\item They are invariant under the group $W_4$.
\item They are reflection positive with respect to every hyperplane $x_\mu=c$, on centred
tensors of test functions with separated supports.
\item They obey the bound of clause (iii-a) of Theorem~0.
\end{itemize}
\end{theorem}

The proof occupies Chapter~\ref{ch:9}, where this statement is the first conclusion of the
uniqueness theorem: the limit family and its rate.

\begin{lemma}[Non-convergence of the first-passage offsets. Proof outstanding]\label{7.3:lem-first-passage-offsets}
Assume the hypotheses of Theorem~0. Assume moreover $\varrho\ge 6(B_\sharp+1)$, and that $g$ lies
below the threshold of clause (v-c)(2) of Theorem~0, a threshold not exceeding $\gamma_U$. For
$g_0\in(0,g_-(g))$ let $K=K(g_0,g)\ge 1$ be the first-passage depth of
Definition~\ref{1.2:def-flow-constants}. Let $a_n$ be the tuned inverse couplings of
clause (v-a) of Theorem~0. Then the offset
\begin{equation*}
g_0\longmapsto a_{K(g_0,g)}-g_0^{-2}
\end{equation*}
takes all its values in $[-\varrho_2/2,\varrho_2/2]$, and it has no limit as $g_0\downarrow 0$: its
values range through an interval of positive length. In particular there is a sequence
$g_0^{(i)}\downarrow 0$ with $K_i=K(g_0^{(i)},g)$ along which the offsets
$a_{K_i}-(g_0^{(i)})^{-2}$ do not converge.

Consequently the family indexed by $g_0\downarrow 0$ and read at the first-passage depth
$K(g_0,g)$ (the first-passage family) is not a tuned family in the sense of
Theorem~\ref{7.3:thm-tuned-family}: it is not of the form $g_0^{-2}=a_n-\hat w$ with $\hat w$
fixed, and its offsets do not converge to any such $\hat w$. Its limits are asserted only along
subsequences, as in clause (ii-$\mathbb{T}$-b) of Theorem~0.
\end{lemma}

\medskip\noindent\emph{Route.} The route is the implication from three clauses of Theorem~0.

The first is clause (0) of Theorem~0. For the first-passage depth $K(g_0,g)$ of
Definition~\ref{1.2:def-flow-constants} it places the last inverse coupling $x_K$ of the run in the
half-open interval of length $B_\sharp$ above $g^{-2}$.

The second is clause (v-a) of Theorem~0. It supplies the tuned couplings $a_n$ with two-sided linear
bounds in the depth $n$, uniformly in the torus.

The third is clause (v-c)(2) of Theorem~0, the corollary of the uniqueness theorem on first-passage
runs. Under the two additional hypotheses of the lemma it states that every bare inverse coupling
whose first-passage run stops at depth $K$ satisfies $|g_0^{-2}-a_K|\le\varrho_2/2$. Hence every
offset lies in $[-\varrho_2/2,\varrho_2/2]$.

Since the offsets are bounded, it suffices to show that they are not a Cauchy family as
$g_0\downarrow 0$; boundedness then yields two distinct accumulation points. A proof would therefore
have to establish the following from clause (0) and clause (v-a). At a bare coupling at which the
first-passage depth passes from $K$ to $K+1$, the inverse coupling $g_0^{-2}$ is held fixed while
the tuned coupling read in the offset changes from $a_K$ to $a_{K+1}$, so that the offset changes by
$a_{K+1}-a_K$; and, by the two-sided linear bounds of clause (v-a), there is a $\delta>0$ with
$a_{K+1}-a_K\ge\delta$ for infinitely many such passages, and these passages occur along a sequence
of bare couplings tending to zero.

For the statement that the values range through an interval of positive length, a proof would
moreover have to establish that the bare couplings with a given first-passage depth form an
interval, on which the offset is continuous in $g_0$ because $a_K$ is fixed there, and that these
continuous pieces, between the changes at the passages, fill an interval of positive length.

\begin{definition}[Tori, bare lattices and curvature cumulants at finite cutoff]
\label{7.4:not-finite-cutoff}
Lengths are measured in physical units, which are the units of the unit lattice. Points of
$\mathbb{R}^4$ are written $x=(x^0,x^1,x^2,x^3)$, $x^0$ being the time coordinate; coordinate
indices $\mu$ run over $0,1,2,3$, and $e_0,\dots,e_3$ is the standard basis.

\emph{Torus.} For $m\ge 1$ let $\mathbb{T}_m:=\mathbb{R}^4/2L^m\mathbb{Z}^4$, the torus of
side $2L^m$ of \cite[(0.1)]{B-CMP109}. Let $q_m\colon\mathbb{R}^4\to\mathbb{T}_m$ be the
quotient map, and let $d_{\mathbb{T}}$ be the flat metric,
\begin{equation*}
d_{\mathbb{T}}\bigl(q_m(x),q_m(y)\bigr):=\min_{k\in\mathbb{Z}^4}\,|x-y-2L^m k| .
\end{equation*}

\emph{Bare lattice.} At cutoff $K$ the lattice spacing is $\varepsilon:=L^{-K}$. The sites are
the points of $q_m\bigl(\varepsilon(\mathbb{Z}+\tfrac12)^4\bigr)$. They are in bijection with
the lattice of \cite[(0.1)]{B-CMP109},
\begin{equation}\label{7.4:eq-finite-cutoff-1}
T_\varepsilon=\Bigl\{x\in\varepsilon\mathbb{Z}^4+\sum_{\mu=0}^{3}\tfrac12\varepsilon e_\mu:\
-L^m<x^\mu<L^m\ \text{for }\mu=0,\dots,3\Bigr\},
\end{equation}
through $q_m$. Bonds and plaquettes are those of \cite{B-CMP109}. For a plaquette $p$,
$\partial p$ is its boundary, and $x(p)\in\mathbb{T}_m$ is its barycentre.

\emph{Measure.} A configuration $U$ assigns an element of $SU(N)$ to each bond; a bond
traversed against its orientation carries the inverse element. $U(\partial p)$ is the ordered
product around $\partial p$, and $dU$ is the product of normalised Haar measures. Put
\begin{equation}\label{7.4:eq-finite-cutoff-2}
a_p(U):=1-\operatorname{Re}\operatorname{tr}U(\partial p),\qquad A(U):=\sum_p a_p(U),
\end{equation}
with $\operatorname{tr}\mathbb{1}=1$, as in \cite[(0.2)]{B-CMP109} at $d=4$. Here $a_p(U)$ is
real and lies in $[0,2]$. It is unchanged under a cyclic relabelling and under an inversion of
$\partial p$. Let $\mu_{K,m}$ be the probability measure on configurations with density
proportional to $e^{-g_0^{-2}A(U)}$ with respect to $dU$. Expectation under $\mu_{K,m}$ is
written $\langle\cdot\rangle_{K,m}$.

\emph{Observables.} For Lipschitz $h\colon\mathbb{T}_m\to\mathbb{C}$ put
$\mathcal{O}(h):=\sum_p h(x(p))\,a_p(U)$. At finite $K$ this is a bounded function of $U$.
Write $\operatorname{Lip}_{\mathbb{T}}h$ for the Lipschitz constant of $h$ with respect to
$d_{\mathbb{T}}$, and put
\begin{equation*}
\|h\|_{\sharp}:=\max\{\|h\|_\infty,\ 40M\operatorname{Lip}_{\mathbb{T}}h\}.
\end{equation*}
For $n\ge1$, Lipschitz $h_1,\dots,h_n$ and $z\in\mathbb{C}^n$ put
\begin{equation}\label{7.4:eq-finite-cutoff-3}
Z_{K,m}(z;\vec h):=\int\exp\Bigl[-g_0^{-2}A(U)+\sum_{i=1}^n z_i\,\mathcal{O}(h_i)\Bigr]\,dU .
\end{equation}
For $n\ge2$ the connected Schwinger function is
\begin{equation}\label{7.4:eq-finite-cutoff-4}
S^T_{n;K,m}(\vec h):=\partial_{z_1}\cdots\partial_{z_n}\log Z_{K,m}(0;\vec h)
=\operatorname{cum}_n\bigl(\mathcal{O}(h_1),\dots,\mathcal{O}(h_n)\bigr).
\end{equation}
Here $\operatorname{cum}_n$ is the joint cumulant under $\mu_{K,m}$. For bounded random variables
$X_1,\dots,X_n$ it is defined by
\begin{equation}\label{7.4:eq-finite-cutoff-5}
\operatorname{cum}_n(X_1,\dots,X_n):=\partial_{z_1}\cdots\partial_{z_n}
\log\Bigl\langle\exp\sum_{i=1}^n z_iX_i\Bigr\rangle_{K,m}\Big|_{z=0}.
\end{equation}

\emph{Test functions on $\mathbb{R}^4$.} These are the elements of
$C^1_c(\mathbb{R}^4;\mathbb{R})$, normed by
\begin{equation*}
\|f\|_{\sharp}:=\max\{\|f\|_\infty,\ 40M\|\nabla f\|_\infty\}.
\end{equation*}
On $\mathbb{R}^4$ one has $\operatorname{Lip}f=\|\nabla f\|_\infty$. For $n\ge2$ the
\emph{separated} $n$-tuples are the elements of
\begin{equation*}
\mathfrak{S}_n:=\bigl\{(f_1,\dots,f_n):\ f_i\in C^1_c(\mathbb{R}^4;\mathbb{R}),\
\text{supports pairwise disjoint}\bigr\}.
\end{equation*}
For such a tuple put
$\mathsf{d}(\vec f):=\min_{i\neq k}\operatorname{dist}(\operatorname{supp}f_i,
\operatorname{supp}f_k)$; then $\mathsf{d}(\vec f)>0$.

\emph{Symmetries.} For $a\in\mathbb{R}^4$ let $\tau_af:=f(\cdot-a)$. $W_4$ is the group of
signed coordinate permutations, and for $w\in W_4$ we put $(w\cdot f)(x):=f(w^{-1}x)$. The
time reflection $\theta(x^0,x^1,x^2,x^3):=(-x^0,x^1,x^2,x^3)$ is an element of $W_4$.
\end{definition}

\begin{proof}
\emph{Sites.} We have $2L^m=2L^{m+K}\varepsilon\in\varepsilon\mathbb{Z}$. Hence
$\varepsilon(\mathbb{Z}+\frac12)^4$ is invariant under translation by $2L^m\mathbb{Z}^4$, and
its image under $q_m$ is well defined. In one coordinate, the points $\varepsilon(j+\frac12)$,
$j\in\mathbb{Z}$, that lie in the open interval $(-L^m,L^m)=(-L^{m+K}\varepsilon,
L^{m+K}\varepsilon)$ are exactly those with $-L^{m+K}\le j\le L^{m+K}-1$. These are
$2L^{m+K}$ points, one in each residue class of $j$ modulo $2L^{m+K}$. So $q_m$ maps
\eqref{7.4:eq-finite-cutoff-1} bijectively onto the set of sites.

\emph{Properties of $a_p$.} $U(\partial p)$ is unitary, with eigenvalues
$e^{i\varphi_1},\dots,e^{i\varphi_N}$. Since $\operatorname{tr}\mathbb{1}=1$, we have
$\operatorname{tr}U(\partial p)=N^{-1}\sum_j e^{i\varphi_j}$, which has modulus at most $1$.
Hence $\operatorname{Re}\operatorname{tr}U(\partial p)\in[-1,1]$ and $a_p(U)\in[0,2]$;
$a_p(U)$ is real by definition.

A cyclic relabelling of $\partial p$ replaces $u_1u_2u_3u_4$ by
$u_2u_3u_4u_1=u_1^{-1}(u_1u_2u_3u_4)u_1$, which has the same trace. Inverting $\partial p$
replaces $U(\partial p)$ by $U(\partial p)^{-1}=U(\partial p)^{*}$. Since
$\operatorname{tr}U^*=\overline{\operatorname{tr}U}$, the real part is unchanged.

\emph{Boundedness of $\mathcal{O}(h)$.} At finite $K$ and $m$ there are finitely many
plaquettes, say $P$ of them. Then $|\mathcal{O}(h)|\le 2P\|h\|_\infty$.

\emph{The identity in \eqref{7.4:eq-finite-cutoff-4}.} The configuration space is compact and
the integrand of \eqref{7.4:eq-finite-cutoff-3} is continuous. For $z$ in a compact set the
integrand is bounded by $\exp\bigl(2P\sum_i|z_i|\,\|h_i\|_\infty\bigr)$. Hence
$Z_{K,m}(\cdot;\vec h)$ is entire, and $Z_{K,m}(0;\vec h)>0$. So $\log Z_{K,m}$ is
holomorphic on a neighbourhood of $z=0$ on which $\operatorname{Re}Z_{K,m}>0$, taking the
principal branch. By the definition of $\mu_{K,m}$,
\begin{equation*}
\log Z_{K,m}(z;\vec h)=\log Z_{K,m}(0;\vec h)
+\log\Bigl\langle\exp\sum_i z_i\,\mathcal{O}(h_i)\Bigr\rangle_{K,m}.
\end{equation*}
The derivative $\partial_{z_1}\cdots\partial_{z_n}$ annihilates the $z$-independent first
term. What remains is \eqref{7.4:eq-finite-cutoff-5} with $X_i=\mathcal{O}(h_i)$. For the same
reason, multiplying $Z_{K,m}$ by any positive factor independent of $z$ leaves $S^T_{n;K,m}$
unchanged. The sourced integral of the correlation theorem differs from
\eqref{7.4:eq-finite-cutoff-3} by such a factor, so the two define the same connected
functions.

The exponent in \eqref{7.4:eq-finite-cutoff-3} equals
\begin{equation*}
-\sum_p\Bigl(g_0^{-2}-\sum_i z_i\,h_i(x(p))\Bigr)a_p(U).
\end{equation*}
Thus the sources shift the bare inverse coupling additively, plaquette by plaquette.

\emph{$\operatorname{Lip}f=\|\nabla f\|_\infty$.} For $x,y\in\mathbb{R}^4$,
\begin{equation*}
f(x)-f(y)=\int_0^1\nabla f(y+t(x-y))\cdot(x-y)\,dt ,
\end{equation*}
so $|f(x)-f(y)|\le\|\nabla f\|_\infty|x-y|$. Conversely, if $\nabla f(x)\neq0$, put
$u:=\nabla f(x)/|\nabla f(x)|$. Then
$|\nabla f(x)|=\lim_{t\downarrow0}t^{-1}(f(x+tu)-f(x))\le\operatorname{Lip}f$.

\emph{Positivity of $\mathsf{d}(\vec f)$.} Each support is closed and bounded, hence compact.
For two non-empty disjoint compact sets, the continuous function $(x,y)\mapsto|x-y|$ attains
its minimum on their product. That minimum is not $0$, since otherwise the sets would share a
point. If a support is empty, its distance to any set is $\inf\emptyset=+\infty$. A minimum of
finitely many positive numbers is positive.

\emph{$\theta\in W_4$.} $\theta$ is the signed coordinate permutation with trivial permutation
and sign $-1$ on the coordinate $x^0$ only.
\end{proof}

\begin{lemma}[Transport of compactly supported test functions to the torus]
\label{7.4:lem-torus-transport}
Let $\rho>0$ and let $m$ satisfy $L^m\ge 2\rho+2$. Write $B_{\rho}$ for the open ball of radius
$\rho$ about $0$ in $\mathbb{R}^4$. For $f$ Lipschitz on $\mathbb{R}^4$ with $\operatorname{supp}f$
compact and contained in $B_{\rho}$, define $f^{(m)}$ on $\mathbb{T}_m$ by
\begin{equation}\label{7.4:eq-torus-transport-1}
f^{(m)}(q_m(x)):=f(x)\quad(x\in B_{\rho}),\qquad
f^{(m)}(\xi):=0\quad(\xi\in\mathbb{T}_m\setminus q_m(B_{\rho})).
\end{equation}
Then the following hold.
\begin{enumerate}
\item[(i)] $f^{(m)}$ is well defined and Lipschitz on $\mathbb{T}_m$, with
$\|f^{(m)}\|_\infty=\|f\|_\infty$ and $\operatorname{Lip}_{\mathbb{T}}f^{(m)}=\operatorname{Lip}f$;
hence $\|f^{(m)}\|_{\sharp}=\|f\|_{\sharp}$.
\item[(ii)] For $f$ and $g$ both supported as above,
\begin{equation*}
\operatorname{dist}_{\mathbb{T}}\bigl(\operatorname{supp}f^{(m)},\operatorname{supp}g^{(m)}\bigr)
=\operatorname{dist}(\operatorname{supp}f,\operatorname{supp}g),
\end{equation*}
and if $\operatorname{supp}f\subset B(x,R)$, then
$\operatorname{supp}f^{(m)}\subset B_{\mathbb{T}}(q_m(x),R)$.
\item[(iii)] If moreover $\operatorname{supp}\tau_af\subset B_{\rho}$, then
$(\tau_af)^{(m)}=\tau_{q_m(a)}(f^{(m)})$. For $w\in W_4$ one has $(w\cdot f)^{(m)}=w\cdot(f^{(m)})$;
here $W_4$ acts on $\mathbb{T}_m$ fixing $q_m(0)$, since each of its elements preserves
$2L^m\mathbb{Z}^4$. In particular $(f\circ\theta)^{(m)}=f^{(m)}\circ\theta$.
\end{enumerate}
Here $\operatorname{dist}_{\mathbb{T}}$ is the distance between subsets of $\mathbb{T}_m$ for
$d_{\mathbb{T}}$, $B_{\mathbb{T}}(\xi,R)$ is the ball of radius $R$ about $\xi$ for $d_{\mathbb{T}}$,
$\tau_{\xi}h:=h(\cdot-\xi)$ for $\xi\in\mathbb{T}_m$, and $(w\cdot h)(\xi):=h(w^{-1}\xi)$.
\end{lemma}

\begin{proof}
Write $|\cdot|$ for the Euclidean norm and $|k|_\infty$ for the maximum norm of $k\in\mathbb{Z}^4$.
The flat metric is
\begin{equation*}
d_{\mathbb{T}}(q_m(x),q_m(y))=\min_{k\in\mathbb{Z}^4}|x-y+2L^mk| .
\end{equation*}
Let $x,y\in B_{\rho}$ and $k\in\mathbb{Z}^4\setminus\{0\}$. Since $|2L^mk|\ge 2L^m|k|_\infty$,
$|k|_\infty\ge1$ and $|x-y|<2\rho$,
\begin{equation}\label{7.4:eq-torus-transport-2}
|x-y+2L^mk|\ge 2L^m|k|_\infty-|x-y|\ge 2L^m-2\rho\ge 2\rho+2>|x-y|,
\end{equation}
where the third inequality uses $L^m\ge2\rho+2$. Hence the minimum is attained at $k=0$ and
$d_{\mathbb{T}}(q_m(x),q_m(y))=|x-y|$: the map $q_m$ is isometric on $B_{\rho}$, in particular
injective there (as $2\rho<2L^m$). Moreover the choice $k=0$ gives
$d_{\mathbb{T}}(q_m(x),q_m(y))\le|x-y|$ for all $x,y\in\mathbb{R}^4$, so $q_m$ is $1$-Lipschitz.

(i) By injectivity of $q_m$ on $B_{\rho}$, \eqref{7.4:eq-torus-transport-1} defines $f^{(m)}$
unambiguously. Its values are those of $f$ on $B_{\rho}$ together with $0$, and $f$ vanishes off
$B_{\rho}$; so $\|f^{(m)}\|_\infty=\|f\|_\infty$. Since $\operatorname{supp}f$ is compact in the open
ball $B_{\rho}$, there is $\rho'<\rho$ with $\operatorname{supp}f\subset B_{\rho'}$ (if
$\operatorname{supp}f=\emptyset$ everything is trivial). Let $\xi,\eta\in\mathbb{T}_m$; there are three
cases.

If $\xi=q_m(x)$, $\eta=q_m(y)$ with $x,y\in B_{\rho}$, the isometry gives
\begin{equation*}
|f^{(m)}(\xi)-f^{(m)}(\eta)|=|f(x)-f(y)|\le\operatorname{Lip}f\,|x-y|
=\operatorname{Lip}f\,d_{\mathbb{T}}(\xi,\eta).
\end{equation*}

If (up to exchanging the two points) $\eta\notin q_m(B_{\rho})$, then $f^{(m)}(\eta)=0$ and
$\eta\notin q_m(\bar B_{\rho'})$. Suppose $\xi=q_m(x)$ with $x\in\operatorname{supp}f$. Choose
$\hat y\in q_m^{-1}(\eta)$ with $|x-\hat y|=d_{\mathbb{T}}(\xi,\eta)$; the segment from $x$ to
$\hat y$ is mapped by $q_m$ onto a $d_{\mathbb{T}}$-geodesic from $\xi$ to $\eta$. As $|x|<\rho'$ and
$|\hat y|>\rho'$, the geodesic leaves $q_m(\bar B_{\rho'})$ at $q_m(x^*)$ for a point $x^*$ of the
segment with $|x^*|=\rho'$; then $f(x^*)=0$, $|x-x^*|\le|x-\hat y|$, and, $x^*$ lying in $B_{\rho}$,
the isometry gives
\begin{equation*}
\begin{split}
|f^{(m)}(\xi)-f^{(m)}(\eta)|&=|f(x)-f(x^*)|\le\operatorname{Lip}f\,|x-x^*|\\
&=\operatorname{Lip}f\,d_{\mathbb{T}}(\xi,q_m(x^*))\le\operatorname{Lip}f\,d_{\mathbb{T}}(\xi,\eta).
\end{split}
\end{equation*}

In the remaining case both $f^{(m)}(\xi)$ and $f^{(m)}(\eta)$ vanish and the difference is $0$.
Hence $\operatorname{Lip}_{\mathbb{T}}f^{(m)}\le\operatorname{Lip}f$. For the reverse inequality,
the same segment argument in $\mathbb{R}^4$ shows that $\operatorname{Lip}f$ is the supremum of
$|f(x)-f(y)|/|x-y|$ over $x\ne y$ in $B_{\rho}$; for such pairs
$|f(x)-f(y)|=|f^{(m)}(q_m(x))-f^{(m)}(q_m(y))|$ and $|x-y|=d_{\mathbb{T}}(q_m(x),q_m(y))$ by the
isometry, so $\operatorname{Lip}f\le\operatorname{Lip}_{\mathbb{T}}f^{(m)}$. The equality of the
$\|\cdot\|_{\sharp}$ norms follows from the two equalities.

(ii) The set where $f^{(m)}\ne0$ is the image under $q_m$ of the set where $f\ne0$; since $q_m$ is
continuous and $q_m(\operatorname{supp}f)$ is compact, $\operatorname{supp}f^{(m)}=q_m(\operatorname{supp}f)$,
and likewise for $g$. For $a\in\operatorname{supp}f$ and $b\in\operatorname{supp}g$, the candidates
$|a-b+2L^mk|$ with $k\ne0$ (the wrapped ones) are at least $2\rho+2$ by
\eqref{7.4:eq-torus-transport-2}, which exceeds every distance inside $B_{\rho}$; so
$d_{\mathbb{T}}(q_m(a),q_m(b))=|a-b|$. Minimising over the two compact supports gives the equality of
distances. If $\operatorname{supp}f\subset B(x,R)$, then for $y\in\operatorname{supp}f$ one has
$d_{\mathbb{T}}(q_m(x),q_m(y))\le|x-y|$, since $q_m$ is $1$-Lipschitz, so
$q_m(\operatorname{supp}f)\subset B_{\mathbb{T}}(q_m(x),R)$.

(iii) Both sides vanish off $q_m(B_{\rho})$ and agree on it. For translations: the left side vanishes
off $q_m(B_{\rho})$ by definition; the right side is non-zero only on
$q_m(\{f\ne0\}+a)=q_m(\{\tau_af\ne0\})\subset q_m(B_{\rho})$. At $q_m(x)$ with $x\in B_{\rho}$ the left
side is $f(x-a)$ and the right side is $f^{(m)}(q_m(x-a))$. If $x-a\in B_{\rho}$ these agree by
\eqref{7.4:eq-torus-transport-1}. Otherwise $f(x-a)=0$, and if $q_m(x-a)=q_m(y)$ with $y\in B_{\rho}$
and $f(y)\ne0$, then $y+a\in\operatorname{supp}\tau_af\subset B_{\rho}$ and $q_m(y+a)=q_m(x)$ with
$y+a\ne x$, against injectivity on $B_{\rho}$; so the right side is $0$ as well. For $w\in W_4$: $w$ is
linear, preserves $2L^m\mathbb{Z}^4$ and the Euclidean norm, so $w(q_m(x))=q_m(wx)$ defines its action
on $\mathbb{T}_m$, fixing $q_m(0)$, and $w^{\pm1}$ maps $B_{\rho}$ and $q_m(B_{\rho})$ onto
themselves; also $\operatorname{supp}(w\cdot f)=w(\operatorname{supp}f)\subset B_{\rho}$. At $q_m(x)$,
$x\in B_{\rho}$, both sides equal $f(w^{-1}x)$, because
$f^{(m)}(w^{-1}q_m(x))=f^{(m)}(q_m(w^{-1}x))$ and $w^{-1}x\in B_{\rho}$; off $q_m(B_{\rho})$ both vanish,
since $w^{-1}\xi\notin q_m(B_{\rho})$ there. Finally $\theta\in W_4$ and $\theta^{-1}=\theta$, so
$\theta\cdot f=f\circ\theta$ and the last assertion is the case $w=\theta$.
\end{proof}

The additive $2$ in the hypothesis $L^m\ge2\rho+2$ is a margin which absorbs plaquette diameters
where the lemma is applied.

For $\vec f=(f_1,\dots,f_n)\in\mathfrak{S}_n$ (see~\ref{7.4:not-finite-cutoff}) with
$\bigcup_i\operatorname{supp}f_i\subset B_{\rho}$ and $L^m\ge2\rho+2$, each $f_i$ is Lipschitz with
compact support in $B_{\rho}$, and we write
\begin{equation}\label{7.4:eq-torus-transport-a}
S^T_{n;K,m}(\vec f):=S^T_{n;K,m}\bigl(f_1^{(m)},\dots,f_n^{(m)}\bigr),
\end{equation}
the right side being the connected Schwinger function on $\mathbb{T}_m$ of~\ref{7.4:not-finite-cutoff}.

\begin{remark}[Standard facts from analysis used below]\label{7.4:rem-standard-facts}
The following standard facts are used in what follows in exactly the form stated here. Their proofs are not reproduced. Where a short argument is indicated in parentheses, it is the argument by which the fact is proved.
\begin{enumerate}
\item[(a)] (Cantor's diagonal argument.) Let $\mathfrak{I}$ be a countable set. For each $s\in\mathfrak{I}$ let an integer $j(s)$ and a bounded sequence of complex numbers $(v_j(s))_{j\ge j(s)}$ be given. Then there is an infinite subset $\mathcal{J}\subset\mathbb{N}$ such that the sequence $(v_j(s))_{j\in\mathcal{J},\,j\ge j(s)}$ converges for every $s\in\mathfrak{I}$. (Enumerate $\mathfrak{I}$ and choose nested infinite subsets of $\mathbb{N}$, the $i$-th one giving convergence for the first $i$ elements of the enumeration. Then take the diagonal subsequence.)

\item[(b)] (Separability.) A subspace of a separable metric space is separable. For a closed box $Q'\subset\mathbb{R}^4$, the space $C^1(Q')$ with the norm
\begin{equation*}
\|f\|_{C^1}:=\max\{\|f\|_\infty,\|\nabla f\|_\infty\}
\end{equation*}
is separable. (The map $f\mapsto(f,\partial_1 f,\dots,\partial_4 f)$ is an isometric embedding of $C^1(Q')$ into $C(Q')^5$, which is separable.)

\item[(c)] (Moments and cumulants.) Let $\mathbb{E}$ denote expectation, and let $X_1,\dots,X_q$ be bounded random variables. For a nonempty $B\subset\{1,\dots,q\}$, let $\kappa_{|B|}(X_B)$ be the joint cumulant of $(X_l)_{l\in B}$. Let $\operatorname{Part}[q]$ be the set of partitions of $\{1,\dots,q\}$ into nonempty blocks. Then
\begin{equation*}
\mathbb{E}\Bigl[\prod_{l=1}^{q}X_l\Bigr]=\sum_{\pi\in\operatorname{Part}[q]}\ \prod_{B\in\pi}\kappa_{|B|}(X_B).
\end{equation*}
The cumulant of a single variable is its expectation, $\kappa_1=\mathbb{E}$. Hence, if $\mathbb{E}X_l=0$ for every $l$, the partitions with a block of size one contribute nothing, and the sum runs over $\operatorname{Part}_{\ge2}[q]$. For $|B|\ge2$, $\kappa_{|B|}(X_B)$ is invariant under $X_l\mapsto X_l+c_l$ with constants $c_l$. The cumulants $\kappa_n$ are multilinear and symmetric in their arguments. For the smeared curvature observables,
\begin{equation*}
\kappa_n\bigl(\mathcal{O}(h_1),\dots,\mathcal{O}(h_n)\bigr)
=\partial_{z_1}\cdots\partial_{z_n}
\log\mathbb{E}\Bigl[e^{\sum_{l=1}^{n}z_l\mathcal{O}(h_l)}\Bigr]\Big|_{z=0}.
\end{equation*}

\item[(d)] (Schwartz kernel theorem, $n$-fold form.) Let $U_1,\dots,U_n$ be open subsets of $\mathbb{R}^4$. Let
\begin{equation*}
\mathcal{B}\colon\mathcal{D}(U_1)\times\dots\times\mathcal{D}(U_n)\to\mathbb{C}
\end{equation*}
be $n$-linear and separately continuous. Then there is a unique $T\in\mathcal{D}'(U_1\times\dots\times U_n)$ such that
\begin{equation*}
T(\varphi_1\otimes\dots\otimes\varphi_n)=\mathcal{B}(\varphi_1,\dots,\varphi_n)
\end{equation*}
for all $\varphi_i\in\mathcal{D}(U_i)$. (The case $n=2$ is the classical kernel theorem. The case of $n$ factors follows by induction, applying the two-factor case to $\mathcal{D}(U_1\times\dots\times U_{n-1})\times\mathcal{D}(U_n)$. Uniqueness holds because finite sums of tensors $\varphi_1\otimes\dots\otimes\varphi_n$ are dense in $\mathcal{D}(U_1\times\dots\times U_n)$.)

\item[(e)] (Gluing of distributions.) Let $(V_\alpha)$ be an open cover of an open set $V$. Suppose that $T_\alpha\in\mathcal{D}'(V_\alpha)$ satisfy $T_\alpha=T_\beta$ on $V_\alpha\cap V_\beta$ for all $\alpha,\beta$. Then there is a unique $T\in\mathcal{D}'(V)$ whose restriction to $V_\alpha$ is $T_\alpha$ for every $\alpha$.

\item[(f)] (Fourier series.) Let $D\ge1$ be an integer, let $Q$ be an open box in $\mathbb{R}^D$, and let $\Psi\in\mathcal{D}(Q)$. Regard $\Psi$ as periodic with period box $\bar Q$, and let $\Lambda^*$ be the dual lattice of the lattice of periods. Let $\hat\Psi(\xi)$, $\xi\in\Lambda^*$, be the Fourier coefficients of $\Psi$. Then the Fourier series of $\Psi$ converges to $\Psi$ in $C^\infty(\bar Q)$. For every integer $\mathsf{n}\ge0$ there is a constant $C_{\mathsf{n},Q}$ such that
\begin{equation*}
|\hat\Psi(\xi)|\le C_{\mathsf{n},Q}\,(1+|\xi|)^{-\mathsf{n}}\,\|\Psi\|_{C^{\mathsf{n}}}
\qquad(\xi\in\Lambda^*).
\end{equation*}
Here $\|\Psi\|_{C^{\mathsf{n}}}$ is the largest sup norm of a derivative of $\Psi$ of order at most $\mathsf{n}$. Moreover, for real $\mathsf{s}$, $\sum_{\xi\in\Lambda^*}(1+|\xi|)^{-\mathsf{s}}<\infty$ if and only if $\mathsf{s}>D$.

\item[(g)] (Cauchy--Schwarz.) Let $(\cdot,\cdot)$ be a Hermitian positive-semidefinite sesquilinear form on a complex vector space. Then $|(F,G)|^2\le(F,F)(G,G)$ for all $F,G$. The null vectors, those $F$ with $(F,F)=0$, form a subspace.

\item[(h)] (Square roots.) A bounded positive operator on a Hilbert space has a unique positive square root.

\item[(i)] (Contraction semigroups.) Let $(T_t)_{t\ge0}$ be a semigroup of positive self-adjoint contractions on a Hilbert space, strongly continuous at $0$, with $T_0=1$. Let $E_1$ be the spectral measure of $T_1$. Then $\ker T_t=\{0\}$ for all $t\ge0$. Moreover, there is a unique self-adjoint operator $H\ge0$ with $T_t=e^{-tH}$ for all $t\ge0$, namely
\begin{equation*}
H=-\log T_1=\int_{]0,1]}(-\log\lambda)\,dE_1(\lambda).
\end{equation*}

\item[(j)] (Stone--Naimark--Ambrose--Godement, spectral calculus.) Let $\vec a\mapsto U(\vec a)$ be a strongly continuous unitary representation of $\mathbb{R}^3$. Then there is a unique projection-valued measure $E$ on $\mathbb{R}^3$ such that
\begin{equation*}
U(\vec a)=\int_{\mathbb{R}^3}e^{i\vec a\cdot\vec\eta}\,dE(\vec\eta).
\end{equation*}
One sets $P_l:=\int_{\mathbb{R}^3}\eta_l\,dE(\vec\eta)$ for $l=1,2,3$. A bounded operator commuting with all $U(\vec a)$ commutes with $E$. A bounded operator commuting with a bounded self-adjoint operator $A$ commutes with the spectral measure of $A$.

\item[(k)] (Uniform continuity.) A continuous function with compact support is uniformly continuous. In particular, if $f$ is continuously differentiable with compact support in $\mathbb{R}^4$, the modulus
\begin{equation*}
\omega_{\nabla f}(\delta):=\sup_{|x-y|\le\delta}|\nabla f(x)-\nabla f(y)|
\end{equation*}
tends to $0$ as $\delta\downarrow0$.
\end{enumerate}
\end{remark}

\begin{lemma}[Lattice symmetries of the curvature cumulants]\label{7.4:lem-lattice-symmetries}
Fix $K$ and $m$, and let $\varepsilon=L^{-K}$, with the torus $\mathbb{T}_m$, the bare lattice, the
measure $\mu_{K,m}$ and the observables $\mathcal{O}(h)$ as in
Notation~\ref{7.4:not-finite-cutoff}. Let
$b\in\varepsilon\mathbb{Z}^4$ \textup{(}taken modulo $2L^m$\textup{)}, let $w\in W_4$, and let
$h_1,\dots,h_n$ be Lipschitz functions on $\mathbb{T}_m$. Write
$\tau_b\vec h=(\tau_bh_1,\dots,\tau_bh_n)$ and $w\cdot\vec h=(w\cdot h_1,\dots,w\cdot h_n)$. For a bond
configuration $U$ write $w\cdot U$ for the configuration given on every positively oriented bond
$\ell$ by $(w\cdot U)(\ell):=U(\overline{w\ell})^{\pm1}$, where $\overline{w\ell}$ is the positively
oriented bond underlying $w\ell$ and the exponent is $+1$ or $-1$ according as $w$ preserves or
reverses the orientation of $\ell$; equivalently, with $U(\ell^{-1}):=U(\ell)^{-1}$ on
reversed bonds, $(w\cdot U)(\ell)=U(w\ell)$ for every oriented bond $\ell$. Then the
joint laws under $\mu_{K,m}$ of $(\mathcal{O}(\tau_bh_i))_{i}$, of $(\mathcal{O}(w\cdot h_i))_{i}$ and of
$(\mathcal{O}(h_i))_{i}$ coincide. Hence, when $n\ge 2$,
\begin{equation*}
S^T_{n;K,m}(\tau_b\vec h)=S^T_{n;K,m}(w\cdot\vec h)=S^T_{n;K,m}(\vec h),
\end{equation*}
and, for every Lipschitz $h$ on $\mathbb{T}_m$, with $\langle\cdot\rangle_{K,m}$ the expectation
under $\mu_{K,m}$,
\begin{equation*}
\langle\mathcal{O}(\tau_bh)\rangle_{K,m}=\langle\mathcal{O}(w\cdot h)\rangle_{K,m}
=\langle\mathcal{O}(h)\rangle_{K,m}.
\end{equation*}
\end{lemma}

\begin{proof}
\emph{Translations.} Since $b\in\varepsilon\mathbb{Z}^4$, the map $\tau_b\colon x\mapsto x+b$ sends
the site lattice $\varepsilon(\mathbb{Z}+\tfrac12)^4$ onto itself, and it sends $2L^m\mathbb{Z}^4$
onto itself, so it descends to $\mathbb{T}_m$. It therefore maps nearest-neighbour pairs of sites to
nearest-neighbour pairs, bonds to bonds and plaquettes to plaquettes; thus $p\mapsto p+b$ is a
permutation of the bare plaquettes, with inverse $p\mapsto p-b$, and $x(p+b)=x(p)+b$. For a bond
configuration $U$ and every bond $\ell$ put $(U\circ\tau_b)(\ell):=U(\ell+b)$. The holonomy of
$U\circ\tau_b$ around
$\partial p'$ is the product of the $U(\ell+b)$, $\ell\in\partial p'$, taken in order, that is
$U(\partial(p'+b))$; hence $a_{p'+b}(U)=a_{p'}(U\circ\tau_b)$. Substituting $p=p'+b$,
\begin{equation}\label{7.4:eq-lattice-symmetries-1}
\begin{aligned}
\mathcal{O}(\tau_bh)(U)&=\sum_p h(x(p)-b)\,a_p(U)=\sum_{p'}h(x(p'))\,a_{p'+b}(U)\\
&=\sum_{p'}h(x(p'))\,a_{p'}(U\circ\tau_b)=\mathcal{O}(h)(U\circ\tau_b).
\end{aligned}
\end{equation}
The measure $\mu_{K,m}$ is invariant under $U\mapsto U\circ\tau_b$: first,
\begin{equation*}
A(U\circ\tau_b)=\sum_p a_p(U\circ\tau_b)=\sum_p a_{p+b}(U)=A(U),
\end{equation*}
because the plaquettes are permuted; second, the product Haar measure $dU$ is carried to itself,
the map only relabelling the bonds. This is also the invariance part of the lemma on reflection
positivity of the lattice measure, part~(i), which asserts that $\mu_{K,m}$ is invariant under every
isometry of $\mathbb{T}_m$ preserving the bare lattice; the invariance of the plaquette terms under
relabelling is that of \cite[(0.2)]{B-CMP109}.

\emph{Signed permutations.} Let $w\in W_4$ be a signed permutation of the coordinates. Since
$-(k+\tfrac12)=(-k-1)+\tfrac12$ for $k\in\mathbb{Z}$, $w$ maps $\varepsilon(\mathbb{Z}+\tfrac12)^4$
onto itself, and it maps $2L^m\mathbb{Z}^4$ onto itself. Being linear and orthogonal, $w$ is
therefore an isometry of $\mathbb{T}_m$ preserving the bare lattice: it permutes the sites, the
unoriented bonds and the plaquettes, and $x(wp')=w\,x(p')$. Let $w\cdot U$ be the configuration
defined in the statement. The holonomy of $w\cdot U$ around $\partial p'$
is then the
holonomy of $U$ around $\partial(wp')$, possibly with its starting point moved cyclically and
possibly traversed in the opposite sense, i.e.\ inverted. Since $\operatorname{tr}(VV')=
\operatorname{tr}(V'V)$ for all $V,V'\in\mathrm{SU}(N)$ and, for $V\in\mathrm{SU}(N)$,
$\operatorname{Re}\operatorname{tr}V^{-1}=
\operatorname{Re}\overline{\operatorname{tr}V}=\operatorname{Re}\operatorname{tr}V$, the real
trace of a holonomy is unchanged by cyclic relabelling and by inversion
\cite[(0.2)]{B-CMP109}; hence $a_{wp'}(U)=a_{p'}(w\cdot U)$. Substituting $p=wp'$,
\begin{equation}\label{7.4:eq-lattice-symmetries-2}
\begin{aligned}
\mathcal{O}(w\cdot h)(U)&=\sum_p h(w^{-1}x(p))\,a_p(U)=\sum_{p'}h(x(p'))\,a_{wp'}(U)\\
&=\sum_{p'}h(x(p'))\,a_{p'}(w\cdot U)=\mathcal{O}(h)(w\cdot U).
\end{aligned}
\end{equation}
By part~(i) of the lemma on reflection positivity of the lattice measure, $\mu_{K,m}$ is invariant
under $U\mapsto w\cdot U$, since $w$ is an isometry of $\mathbb{T}_m$ preserving the bare lattice.

\emph{Conclusion.} Let $\Phi$ denote either $U\mapsto U\circ\tau_b$ or $U\mapsto w\cdot U$. By
\eqref{7.4:eq-lattice-symmetries-1} and \eqref{7.4:eq-lattice-symmetries-2}, applied to each $h_i$
with the same map $\Phi$, the vector $(\mathcal{O}(\tau_bh_i))_i$, respectively
$(\mathcal{O}(w\cdot h_i))_i$, equals $(\mathcal{O}(h_i))_i\circ\Phi$, and $\Phi$ preserves
$\mu_{K,m}$; so the three vectors have the same joint law under $\mu_{K,m}$. At finite $K$ each
$\mathcal{O}(h_i)$ is a bounded function, so its joint moments, and with them the joint cumulants,
are functionals of the joint law. Equal joint laws therefore give equal means and equal cumulants;
since $S^T_{n;K,m}(\vec h)=\kappa_n(\mathcal{O}(h_1),\dots,\mathcal{O}(h_n))$ for $n\ge2$, the two
displayed identities follow.
\end{proof}

\begin{lemma}[Half-space observables depend only on an open slab]\label{7.4:lem-half-space}
Let $\rho>0$ and $L^m\ge 2\rho+2$, let $B_{\rho}$ be the open ball of radius $\rho$ about $0$ in
$\mathbb{R}^4$, and let $h$ be Lipschitz on $\mathbb{R}^4$ with $\operatorname{supp}h$ compact and
contained in $B_{\rho}\cap\{x^0>0\}$. Let $\theta$ denote the time reflection $x^0\mapsto -x^0$, on
$\mathbb{R}^4$ and on $\mathbb{T}_m$, and put
\begin{equation*}
H_{+}:=q_m\bigl(\{x\in\mathbb{R}^4:\ 0<x^0<L^m\}\bigr)\subset\mathbb{T}_m .
\end{equation*}
Then every bare plaquette $p$ of $\mathbb{T}_m$ with $h^{(m)}(x(p))\neq 0$ has all four of its bonds in
the open slab $H_{+}$. In particular no such bond lies in or crosses the plane $x^0=0$, and every
such bond keeps time-distance at least $\rho+1$ from the antipodal plane $x^0=L^m$, which is
identified with $x^0=-L^m$.
Consequently $\mathcal{O}(h^{(m)})$ and the centred observable
\begin{equation*}
{:}\mathcal{O}(h^{(m)}){:}\ :=\ \mathcal{O}(h^{(m)})-\langle\mathcal{O}(h^{(m)})\rangle
\end{equation*}
are bounded measurable functions of the bond variables in $H_{+}$; here $\langle\cdot\rangle$ is the
expectation under $\mu_{K,m}$. Moreover
\begin{equation*}
\mathcal{O}(h^{(m)})(\theta\cdot U)=\mathcal{O}\bigl((h\circ\theta)^{(m)}\bigr)(U),
\qquad
\langle\mathcal{O}(h^{(m)})\rangle=\langle\mathcal{O}\bigl((h\circ\theta)^{(m)}\bigr)\rangle .
\end{equation*}
\end{lemma}

\begin{proof}
We use the tori, bare lattices, plaquette barycentres $x(p)$ and the quantities $a_p(U)$ fixed in
Notation~\ref{7.4:not-finite-cutoff}. The bare spacing is
$\varepsilon=L^{-K}$, and the sites have coordinates in $\varepsilon(\mathbb{Z}+\tfrac12)$, taken
modulo $2L^m$, by \cite[(0.1)]{B-CMP109}. In particular every site has time coordinate in
$\varepsilon(\mathbb{Z}+\tfrac12)$. Since $2L^m/\varepsilon=2L^{m+K}$ is an integer, the set
$\varepsilon(\mathbb{Z}+\tfrac12)^4$ is invariant under translation by $2L^m\mathbb{Z}^4$.

Let $p$ be a bare plaquette with $h^{(m)}(x(p))\neq0$. By definition $h^{(m)}$ vanishes off
$q_m(B_{\rho})$, and $h^{(m)}(q_m(y))=h(y)$ for $y\in B_{\rho}$. Hence $x(p)=q_m(y)$ for some
$y\in B_{\rho}$ with $h(y)\ne0$, so $y\in\operatorname{supp}h$. The hypothesis on
$\operatorname{supp}h$ and
$|y|<\rho$ give $0<y^0<\rho$.

Choose a square $\tilde p$ of the lattice $\varepsilon(\mathbb{Z}+\tfrac12)^4$ with
$q_m(\tilde p)=p$. Its barycentre is mapped by $q_m$ to $x(p)=q_m(y)$, so it is congruent to $y$
modulo $2L^m\mathbb{Z}^4$. After a translation by an element of $2L^m\mathbb{Z}^4$, which preserves
$\varepsilon(\mathbb{Z}+\tfrac12)^4$ by the invariance just noted and commutes with $q_m$, the square
$\tilde p$ has barycentre $y$, its corners lie in $\varepsilon(\mathbb{Z}+\tfrac12)^4$, and still
$q_m(\tilde p)=p$.

Suppose first that $p$ lies in a $(0,i)$-plane. Then the corners of $\tilde p$ have times
$(k+\tfrac12)\varepsilon$
and $(k+\tfrac32)\varepsilon$ for some $k\in\mathbb{Z}$, and its barycentre has time
$y^0=(k+1)\varepsilon$. The condition $0<y^0$ forces $k+1>0$, that is $k\ge0$. All its bonds
therefore have times in $[(k+\tfrac12)\varepsilon,(k+\tfrac32)\varepsilon]$, and this interval lies
in $[\varepsilon/2,\rho+\varepsilon]$. Indeed $(k+\tfrac12)\varepsilon\ge\varepsilon/2>0$, while
\begin{equation*}
(k+\tfrac32)\varepsilon=y^0+\tfrac{\varepsilon}{2}<\rho+\tfrac{\varepsilon}{2}\le\rho+1,
\end{equation*}
because $\varepsilon=L^{-K}\le1$.

Suppose next that $p$ lies in an $(i,j)$-plane with $i,j$ spatial. Then all four corners of
$\tilde p$ have the
same time $(k'+\tfrac12)\varepsilon$ for some $k'\in\mathbb{Z}$. This time equals $y^0$ and so lies
in $(0,\rho)$.

In both cases every bond of $\tilde p$ has all its times in $(0,\rho+1)$. From
$L^m\ge2\rho+2$ we get $\rho+1\le L^m-(\rho+1)$, hence $(0,\rho+1)\subset(0,L^m)$. Moreover
every such time is at distance at least $\rho+1$ from $L^m$. Since $q_m(\tilde p)=p$, applying
$q_m$ shows that all four bonds of $p$
lie in $H_{+}$. They keep distance at least $\rho+1$ in time from the antipodal plane, and they
neither lie in nor cross $x^0=0$.

Boundedness and measurability come next. The torus carries finitely many bare plaquettes, and
\begin{equation*}
\mathcal{O}(h^{(m)})(U)=\sum_p h^{(m)}(x(p))\,a_p(U).
\end{equation*}
Only plaquettes with $h^{(m)}(x(p))\neq0$ contribute. For each of them $a_p(U)$ is a continuous
function of the four bond variables of $\partial p$, and these lie in $H_{+}$. Since
$|\operatorname{Re}\operatorname{tr}V|\le1$ for $V\in\mathrm{SU}(N)$ with
$\operatorname{tr}\mathbb{1}=1$, we
have $a_p\in[0,2]$. Hence $\mathcal{O}(h^{(m)})$ is a bounded measurable function of the bond
variables in $H_{+}$, and its mean under the probability measure $\mu_{K,m}$ is finite. Subtracting
this constant, ${:}\mathcal{O}(h^{(m)}){:}$ has the same property.

Reflection. The map $\theta$ preserves $\varepsilon(\mathbb{Z}+\tfrac12)^4$ and $2L^m\mathbb{Z}^4$,
so it maps sites to sites, bonds to bonds and plaquettes to plaquettes. Being an involution, it
permutes the plaquettes. Being affine, it maps barycentres to barycentres, so $x(\theta p)=\theta x(p)$.

Recall, as in Lemma~\ref{1.1:lem-invariance}, that the reflected configuration $\theta\cdot U$
assigns to the oriented bond $b$ the variable $U(\theta b)$. The holonomy of $\theta\cdot U$ around
$\partial p$ is
therefore the holonomy of $U$ around $\partial(\theta p)$, possibly started at another corner and
possibly traversed in the opposite orientation. A change of starting corner conjugates the
holonomy, which leaves the trace unchanged. A reversal of orientation inverts it, and
\begin{equation*}
\operatorname{Re}\operatorname{tr}V^{-1}=\operatorname{Re}\operatorname{tr}V^{*}
=\operatorname{Re}\operatorname{tr}V
\end{equation*}
for unitary $V$. Hence $a_p(\theta\cdot U)=a_{\theta p}(U)$.

The function $h\circ\theta$ is Lipschitz, with $\operatorname{supp}(h\circ\theta)=\theta(\operatorname{supp}h)$
compact in $B_{\rho}$. Part (iii) of Lemma~\ref{7.4:lem-torus-transport} therefore gives
$(h\circ\theta)^{(m)}=h^{(m)}\circ\theta$. Writing $p'=\theta p$, we obtain
\begin{equation}\label{7.4:eq-half-space-1}
\begin{aligned}
\mathcal{O}(h^{(m)})(\theta\cdot U)
&=\sum_p h^{(m)}(x(p))\,a_{\theta p}(U)
=\sum_{p'} h^{(m)}(\theta x(p'))\,a_{p'}(U)\\
&=\sum_{p'} (h\circ\theta)^{(m)}(x(p'))\,a_{p'}(U)
=\mathcal{O}\bigl((h\circ\theta)^{(m)}\bigr)(U).
\end{aligned}
\end{equation}

For the means, apply part (i) of the lemma on reflection positivity of the lattice measure in the
time direction, with the plane pair $\{x^0=0\}\cup\{x^0=L^m\}$. The reflection in this pair is
$\theta$, and $H_{+}$ is one of the two open components of its complement. By the first part of the
proof, $F=\mathcal{O}(h^{(m)})$ is a bounded measurable function of the bond variables in $H_{+}$.
That lemma gives
$\langle\mathcal{O}(h^{(m)})\rangle=\langle\mathcal{O}(h^{(m)})(\theta\cdot U)\rangle$, and by
\eqref{7.4:eq-half-space-1} the right side equals $\langle\mathcal{O}((h\circ\theta)^{(m)})\rangle$.
\end{proof}

\begin{lemma}[The translation modulus in the sharp norm]\label{7.4:lem-translation-modulus}
For $a\in\mathbb{R}^4$ let $\tau_a f:=f(\cdot-a)$. For $f\in C^1_c(\mathbb{R}^4)$ write
$\|\nabla f\|_\infty:=\sup_{x}|\nabla f(x)|$, with the Euclidean norm on $\mathbb{R}^4$, and for
$\delta\ge 0$ let
\begin{equation*}
\omega_{\nabla f}(\delta):=\sup\bigl\{|\nabla f(x)-\nabla f(y)| : x,y\in\mathbb{R}^4,\ |x-y|\le\delta\bigr\}
\end{equation*}
be the modulus of uniform continuity of $\nabla f$. Let $f\in C^1_c(\mathbb{R}^4)$ and $a,b\in\mathbb{R}^4$.
\begin{enumerate}
\item We have $\|\tau_af-\tau_bf\|_\infty\le\|\nabla f\|_\infty|a-b|$ and
\begin{equation*}
\operatorname{Lip}(\tau_af-\tau_bf)=\|\nabla f(\cdot-a)-\nabla f(\cdot-b)\|_\infty
\le\omega_{\nabla f}(|a-b|).
\end{equation*}
\item Consequently
\begin{equation}\label{7.4:eq-translation-modulus-a}
\|\tau_af-\tau_bf\|_{\sharp}\le\varpi_f(|a-b|),
\qquad
\varpi_f(\delta):=\max\{\|\nabla f\|_\infty\delta,\ 40M\,\omega_{\nabla f}(\delta)\},
\end{equation}
and $\varpi_f(\delta)\to 0$ as $\delta\to 0$.
\item If $f\in C^2_c(\mathbb{R}^4)$, then $\omega_{\nabla f}(\delta)\le\|\nabla^2f\|_\infty\,\delta$, where
$\|\nabla^2 f\|_\infty$ is the supremum over $x$ of the operator norm of the Hessian
$\nabla^2f(x)$; this is an explicit rate.
\item $\|\tau_af\|_{\sharp}=\|f\|_{\sharp}$.
\item Let $\rho>0$ and $L^m\ge 2\rho+2$, and suppose that $\tau_af$ and $\tau_bf$ are supported in
$B_{\rho}$. Then the transported functions on $\mathbb{T}_m$ satisfy the same bounds:
$\|(\tau_af)^{(m)}-(\tau_bf)^{(m)}\|_\infty\le\|\nabla f\|_\infty|a-b|$,
\begin{equation*}
\operatorname{Lip}_{\mathbb{T}}\bigl((\tau_af)^{(m)}-(\tau_bf)^{(m)}\bigr)\le\omega_{\nabla f}(|a-b|),
\end{equation*}
hence $\|(\tau_af)^{(m)}-(\tau_bf)^{(m)}\|_{\sharp}\le\varpi_f(|a-b|)$, and
$\|(\tau_af)^{(m)}\|_{\sharp}=\|f\|_{\sharp}$. If moreover $\operatorname{supp}f\subset B_\rho$, then
$(\tau_af)^{(m)}=\tau_{q_m(a)}(f^{(m)})$ and $(\tau_bf)^{(m)}=\tau_{q_m(b)}(f^{(m)})$.
\item For merely Lipschitz $f$ the conclusion of \textup{(2)} fails in general. Let $e_1$ be a
coordinate unit vector, $x_1:=x\cdot e_1$, and $f(x)=|x_1|\chi(x)$ with $\chi\in C^1_c(\mathbb{R}^4)$,
$\chi\equiv 1$ near $0$. Then $\operatorname{Lip}(\tau_{se_1}f-f)\ge 2$ for all small $s\ne 0$: the
gradient jump $2e_1$ across $x_1=0$ is displaced, not removed. Hence
$\|\tau_{se_1}f-f\|_{\sharp}\ge 80M$, which does not tend to $0$.
\end{enumerate}
\end{lemma}

\begin{proof}
\emph{Sup part.} Fix $x$. The function $t\mapsto f(x-b-t(a-b))$ on $[0,1]$ is $C^1$ with
derivative $-\nabla f(x-b-t(a-b))\cdot(a-b)$. By the mean value theorem along the segment from
$x-b$ to $x-a$,
\begin{equation*}
|f(x-a)-f(x-b)|\le\|\nabla f\|_\infty|a-b|.
\end{equation*}

\emph{Lipschitz part.} Let $g\in C^1_c(\mathbb{R}^4)$. The same argument along the segment from
$y$ to $x$ gives $|g(x)-g(y)|\le\|\nabla g\|_\infty|x-y|$, so
$\operatorname{Lip}g\le\|\nabla g\|_\infty$. Conversely, if $\nabla g(x)\ne 0$, put
$u:=\nabla g(x)/|\nabla g(x)|$. Then $(g(x+hu)-g(x))/h\to\nabla g(x)\cdot u=|\nabla g(x)|$ as
$h\downarrow 0$, so $\operatorname{Lip}g\ge|\nabla g(x)|$ for every $x$. Thus the Lipschitz constant
of a $C^1$ function with compact support on $\mathbb{R}^4$ is the supremum of its gradient. Apply this
to $g:=\tau_af-\tau_bf\in C^1_c(\mathbb{R}^4)$, for which
$\nabla g(x)=\nabla f(x-a)-\nabla f(x-b)$. This gives the identity in (1). Since
$|(x-a)-(x-b)|=|a-b|$, each $|\nabla g(x)|$ is at most $\omega_{\nabla f}(|a-b|)$, which gives the
inequality.

\emph{(2).} By definition $\|g\|_{\sharp}=\max\{\|g\|_\infty,40M\operatorname{Lip}g\}$, so
\eqref{7.4:eq-translation-modulus-a} follows from (1). The gradient $\nabla f$ is continuous with
compact support, hence bounded and uniformly continuous. Therefore $\|\nabla f\|_\infty<\infty$ and
$\omega_{\nabla f}(\delta)\to0$ as $\delta\to 0$, so $\varpi_f(\delta)\to 0$.

\emph{(3).} For $f\in C^2_c(\mathbb{R}^4)$ and $x,y\in\mathbb{R}^4$, the mean value theorem gives
\begin{equation*}
\nabla f(x)-\nabla f(y)=\int_0^1\nabla^2f(y+t(x-y))\,(x-y)\,dt .
\end{equation*}
Hence $|\nabla f(x)-\nabla f(y)|\le\|\nabla^2f\|_\infty|x-y|$, and taking the supremum over
$|x-y|\le\delta$ gives the bound.

\emph{(4).} The translation $x\mapsto x-a$ is an isometry of $\mathbb{R}^4$ onto itself. It therefore
changes neither the supremum nor the Lipschitz constant.

\emph{(5).} The functions $\tau_af$, $\tau_bf$ and $g=\tau_af-\tau_bf$ are Lipschitz with compact
support in $B_\rho$, so their transports are defined by Lemma~\ref{7.4:lem-torus-transport}. By the
definition of the transport, $g^{(m)}$ takes the value $g(x)$ at $q_m(x)$ for $x\in B_\rho$ and
vanishes off $q_m(B_\rho)$. Hence $g^{(m)}=(\tau_af)^{(m)}-(\tau_bf)^{(m)}$. By
Lemma~\ref{7.4:lem-torus-transport}\,(i), $\|g^{(m)}\|_\infty=\|g\|_\infty$ and
$\operatorname{Lip}_{\mathbb{T}}g^{(m)}=\operatorname{Lip}g$, so the bounds of (1) and (2) carry over
unchanged. The same clause together with (4) gives
$\|(\tau_af)^{(m)}\|_{\sharp}=\|\tau_af\|_{\sharp}=\|f\|_{\sharp}$. If also
$\operatorname{supp}f\subset B_\rho$, then Lemma~\ref{7.4:lem-torus-transport}\,(iii) gives
$(\tau_af)^{(m)}=\tau_{q_m(a)}(f^{(m)})$, and likewise for $b$.

\emph{(6).} Choose $r>0$ with $\chi\equiv 1$ on the ball of radius $r$ about $0$, and $R>0$ with
$\operatorname{supp}\chi\subset B_R$. Then $f=\chi\cdot\min\{|x_1|,R\}$ is a product of bounded
Lipschitz functions, so $f$ is Lipschitz with compact support. Let $0<|s|<r$ and
$g_s:=\tau_{se_1}f-f$. The points $0$, $se_1$, $-se_1$ lie in the ball where $\chi\equiv1$. Hence
\begin{equation*}
g_s(0)=|{-s}|-0=|s|,\qquad g_s(se_1)=0-|s|=-|s|,
\end{equation*}
so $|g_s(0)-g_s(se_1)|=2|s|=2|0-se_1|$ and $\operatorname{Lip}g_s\ge 2$. Therefore
$\|g_s\|_{\sharp}\ge 40M\cdot2=80M$ for all such $s$, and this does not tend to $0$ as $s\to0$.
\end{proof}

The bound \eqref{7.4:eq-translation-modulus-a} is what makes translations usable at all.
Approximating a translation by translations from a dense subgroup and using continuity in the
translation parameter of the limit, and rounding a translation to the lattice at the same cutoff, both
require $\varpi_f(\delta)\to0$. By (6), this fails for Lipschitz data; it requires $f\in C^1$, or
$f\in C^2$ for a rate.

\begin{definition}[Centred moments as sums of cumulants]\label{7.4:def-centred-moments}
Let $q\ge 1$ and let $h_1,\dots,h_q$ be Lipschitz functions on $\mathbb{T}_m$ with pairwise
disjoint supports; write $\vec h=(h_1,\dots,h_q)$. For a non-empty $B\subset\{1,\dots,q\}$ let
$\vec h_B:=(h_l)_{l\in B}$, listed in increasing order of $l$. Let
$\operatorname{Part}_{\ge 2}[q]$ be the set of partitions $\pi$ of $\{1,\dots,q\}$ all of whose
blocks have at least two elements. Let $\langle\cdot\rangle_{K,m}$ be the expectation under
$\mu_{K,m}$, and ${:}\mathcal{O}(h){:}:=\mathcal{O}(h)-\langle\mathcal{O}(h)\rangle_{K,m}$. Then
\begin{equation}\label{7.4:eq-centred-moments-a}
\Bigl\langle\prod_{l\le q}{:}\mathcal{O}(h_l){:}\Bigr\rangle_{K,m}
=\sum_{\pi\in\operatorname{Part}_{\ge 2}[q]}\ \prod_{B\in\pi}S^T_{|B|;K,m}(\vec h_B)
=:S_{q;K,m}(\vec h),
\end{equation}
and each factor $S^T_{|B|;K,m}(\vec h_B)$ is a cumulant of observables with pairwise disjoint
supports. For $\vec f=(f_1,\dots,f_q)\in\mathfrak{S}_q$ with all supports in $B_{\rho}$ and
$L^m\ge 2\rho+2$ we put
$S_{q;K,m}(\vec f):=S_{q;K,m}(\vec f^{(m)})$, where
$\vec f^{(m)}:=(f_1^{(m)},\dots,f_q^{(m)})$ is the transport of
Lemma~\ref{7.4:lem-torus-transport}.
\end{definition}

\begin{proof}
By \ref{7.4:not-finite-cutoff}, $\mathcal{O}(h_l)=\sum_p h_l(x(p))\,a_p(U)$ is a finite sum over the
plaquettes of the bare lattice on $\mathbb{T}_m$. Since the normalised trace of a matrix in
$SU(N)$ is the mean of its eigenvalues, which lie on the unit circle, $0\le a_p(U)\le 2$; hence
$|\mathcal{O}(h_l)|\le 2\|h_l\|_\infty$ times the number of plaquettes, and the variables
$X_l:={:}\mathcal{O}(h_l){:}$ are bounded.

We apply the moment--cumulant facts of Remark~\ref{7.4:rem-standard-facts} to $X_1,\dots,X_q$.
First,
\begin{equation*}
\Bigl\langle\prod_{l\le q}X_l\Bigr\rangle_{K,m}
=\sum_{\pi}\ \prod_{B\in\pi}\kappa_{|B|}\bigl((X_l)_{l\in B}\bigr),
\end{equation*}
the sum running over all partitions $\pi$ of $\{1,\dots,q\}$. Since $\kappa_1=\langle\cdot\rangle_{K,m}$
and $\langle X_l\rangle_{K,m}=0$ for every $l$, each partition with a block of size one contributes
zero, and only $\pi\in\operatorname{Part}_{\ge 2}[q]$ remain. Second, for $|B|\ge 2$ the joint
cumulant is invariant under $X_l\mapsto X_l+c_l$ with constants $c_l$; taking
$c_l=\langle\mathcal{O}(h_l)\rangle_{K,m}$ gives
\begin{equation*}
\kappa_{|B|}\bigl(({:}\mathcal{O}(h_l){:})_{l\in B}\bigr)
=\kappa_{|B|}\bigl((\mathcal{O}(h_l))_{l\in B}\bigr)=S^T_{|B|;K,m}(\vec h_B),
\end{equation*}
the last equality being the definition of the connected Schwinger function as the joint cumulant.
This proves \eqref{7.4:eq-centred-moments-a}. The functions in $\vec h_B$ are among
$h_1,\dots,h_q$, so their supports are pairwise disjoint.

For the transported definition, Lemma~\ref{7.4:lem-torus-transport}(i) shows that each $f_l^{(m)}$
is Lipschitz on $\mathbb{T}_m$, and Lemma~\ref{7.4:lem-torus-transport}(ii) gives
$d_{\mathbb{T}}(\operatorname{supp}f_l^{(m)},\operatorname{supp}f_{l'}^{(m)})
=\operatorname{dist}(\operatorname{supp}f_l,\operatorname{supp}f_{l'})$ for $l\ne l'$. The
supports of $f_l,f_{l'}$ are compact and disjoint, so this distance is positive and the transported
supports are disjoint. Thus $\vec f^{(m)}$ satisfies the hypotheses of the first part, and
$S_{q;K,m}(\vec f)$ is defined. With the convention \eqref{7.4:eq-torus-transport-a},
$S^T_{|B|;K,m}(\vec f_B)=S^T_{|B|;K,m}(\vec f^{(m)}_B)$, so
\eqref{7.4:eq-centred-moments-a} holds with $\vec f$ in place of $\vec h$.
\end{proof}

\begin{lemma}[Linearity in the test function]\label{7.4:lem-linearity}
Let $\rho>0$, let $m$ satisfy $L^m\ge 2\rho+2$, let $K$ be a cutoff and let $q\ge 1$. Tori, bare lattices
and barycentres are as in Notation~\ref{7.4:not-finite-cutoff}, and $\langle\cdot\rangle_{K,m}$ is
the expectation with respect to $\mu_{K,m}$.

Let $\mathsf{F}_q(\rho)$ be the complex vector space of functions on $(\mathbb{R}^4)^q$ spanned by the
tensors $f_1\otimes\cdots\otimes f_q$ of Lipschitz functions $f_l$ on $\mathbb{R}^4$ with compact
supports in $B_{\rho}$. For a bare plaquette $p$ with $x(p)\in q_m(B_{\rho})$ let $\hat x(p)\in B_{\rho}$
be the lift of $x(p)$, that is, the point with $q_m(\hat x(p))=x(p)$. For $F\in\mathsf{F}_q(\rho)$
let $\Phi_{K,m}(F)$ and $\Psi_{K,m}(F)$ be defined by the first two lines of
\begin{equation}\label{7.4:eq-linearity-a}
\begin{aligned}
\Phi_{K,m}(F)&:=\sum_{p_1,\dots,p_q}F\bigl(\hat x(p_1),\dots,\hat x(p_q)\bigr)
\prod_{l\le q}\bigl(a_{p_l}(U)-\langle a_{p_l}\rangle_{K,m}\bigr),\\
\Psi_{K,m}(F)&:=\sum_{p_1,\dots,p_q}F\bigl(\hat x(p_1),\dots,\hat x(p_q)\bigr)\,
\kappa_q\bigl(a_{p_1},\dots,a_{p_q}\bigr),\\
\Phi_{K,m}(F)\,\Phi_{K,m}(G)&=\Phi_{K,m}(F\otimes G),
\end{aligned}
\end{equation}
the sums running over $q$-tuples of bare plaquettes whose barycentres lie in $q_m(B_{\rho})$. The same
formulas with $q$ replaced by $q'$ define $\Phi_{K,m}$ and $\Psi_{K,m}$ on $\mathsf{F}_{q'}(\rho)$ for
every $q'\ge 1$, the number of arguments being read off the function. Then:
\begin{enumerate}
\item $\Phi_{K,m}$ and $\Psi_{K,m}$ are linear in $F$;
\item on tensors,
\begin{gather*}
\Phi_{K,m}(f_1\otimes\cdots\otimes f_q)=\prod_{l\le q}{:}\mathcal{O}(f_l^{(m)}){:},\\
\Psi_{K,m}(f_1\otimes\cdots\otimes f_q)
=\kappa_q\bigl(\mathcal{O}(f_1^{(m)}),\dots,\mathcal{O}(f_q^{(m)})\bigr);
\end{gather*}
\item the third line of \eqref{7.4:eq-linearity-a} holds for $F\in\mathsf{F}_q(\rho)$ and
$G\in\mathsf{F}_{q'}(\rho)$.
\end{enumerate}
Let $\mathcal{T}^{\mathrm{sep}}_q(\rho)$ be the complex span of the separated real $C^1$ tensors
$f_1\otimes\cdots\otimes f_q$, that is, those with $f_l$ real, of class $C^1$, with compact supports in
$B_{\rho}$ and with pairwise disjoint supports. Consequently, for $F\in\mathcal{T}^{\mathrm{sep}}_q(\rho)$ the
expressions
\begin{equation*}
S_{q;K,m}(F):=\bigl\langle\Phi_{K,m}(F)\bigr\rangle_{K,m},\qquad
S^T_{q;K,m}(F):=\Psi_{K,m}(F)
\end{equation*}
are well-defined linear functionals of the function $F$. They extend, respectively, the centred
moments of Definition~\ref{7.4:def-centred-moments} and the cumulant, and they are independent of the
representation of $F$ as a sum of tensors.
\end{lemma}

\begin{proof}
The lift $\hat x(p)$ is unique: two points of $B_{\rho}$ differ in each coordinate by less than
$2\rho<2L^m$, so they cannot differ by a non-zero element of $2L^m\mathbb{Z}^4$. Hence $q_m$ is
injective on $B_{\rho}$.

Linearity in $F$ is visible from \eqref{7.4:eq-linearity-a}. The coefficients
$\prod_l(a_{p_l}(U)-\langle a_{p_l}\rangle_{K,m})$ and $\kappa_q(a_{p_1},\dots,a_{p_q})$ do not depend
on $F$. Moreover, $F$ enters only through its values at the tuples
$(\hat x(p_1),\dots,\hat x(p_q))$.

Now let $F=f_1\otimes\cdots\otimes f_q$. By the definition of the transport in
Lemma~\ref{7.4:lem-torus-transport}, $f_l^{(m)}(q_m(x))=f_l(x)$ for $x\in B_{\rho}$. Hence
\begin{equation*}
F\bigl(\hat x(p_1),\dots,\hat x(p_q)\bigr)=\prod_{l\le q}f_l^{(m)}\bigl(x(p_l)\bigr).
\end{equation*}
Since $f_l^{(m)}=0$ off $q_m(B_{\rho})$, each sum over $p_l$ may be extended to all bare plaquettes of
$\mathbb{T}_m$ without changing its value. The sum in $\Phi_{K,m}(F)$ is then a sum over independent
indices of a product, so it factorises:
\begin{equation*}
\Phi_{K,m}(F)=\prod_{l\le q}\,\sum_{p}f_l^{(m)}\bigl(x(p)\bigr)\bigl(a_p(U)-\langle a_p\rangle_{K,m}\bigr)
=\prod_{l\le q}{:}\mathcal{O}(f_l^{(m)}){:}.
\end{equation*}
For the last equality we used that, for every function $h$ on the torus,
\begin{equation*}
{:}\mathcal{O}(h){:}=\sum_p h\bigl(x(p)\bigr)\bigl(a_p(U)-\langle a_p\rangle_{K,m}\bigr),
\end{equation*}
which is linear in $h$. This identity follows by subtracting the expectation of
$\mathcal{O}(h)=\sum_p h(x(p))\,a_p(U)$ term by term.

For $\Psi_{K,m}$ we use that the random variables $a_p$ are bounded and that the joint cumulant
$\kappa_q$ of bounded random variables is multilinear, by the moment--cumulant relation recalled in
Remark~\ref{7.4:rem-standard-facts}. Applying multilinearity in each of the $q$ arguments gives
\begin{equation*}
\begin{split}
\Psi_{K,m}(F)&=\kappa_q\Bigl(\sum_{p_1}f_1^{(m)}\bigl(x(p_1)\bigr)a_{p_1},\dots,
\sum_{p_q}f_q^{(m)}\bigl(x(p_q)\bigr)a_{p_q}\Bigr)\\
&=\kappa_q\bigl(\mathcal{O}(f_1^{(m)}),\dots,\mathcal{O}(f_q^{(m)})\bigr).
\end{split}
\end{equation*}
This proves the tensor identities.

For multiplicativity, note first that $F\otimes G\in\mathsf{F}_{q+q'}(\rho)$, since a tensor of two
tensors of admissible functions is again such a tensor. Second, writing
$\hat x(\vec p):=(\hat x(p_1),\dots,\hat x(p_q))$ for a $q$-tuple $\vec p=(p_1,\dots,p_q)$, and
similarly for $q'$-tuples, we have
$(F\otimes G)(\hat x(\vec p),\hat x(\vec p\,'))=F(\hat x(\vec p))\,G(\hat x(\vec p\,'))$. The sum
defining $\Phi_{K,m}(F\otimes G)$ runs over pairs $(\vec p,\vec p\,')$ of a $q$-tuple and a $q'$-tuple.
Its summand is the product of the summand of $\Phi_{K,m}(F)$ at $\vec p$ and the summand of
$\Phi_{K,m}(G)$ at $\vec p\,'$. Hence the sum factorises into $\Phi_{K,m}(F)\,\Phi_{K,m}(G)$.

Finally, let $F\in\mathcal{T}^{\mathrm{sep}}_q(\rho)$. Compactly supported $C^1$ functions are Lipschitz, so
$\mathcal{T}^{\mathrm{sep}}_q(\rho)\subset\mathsf{F}_q(\rho)$. Suppose $F=\sum_k c_k\bigotimes_l f_{k,l}$.
By linearity and the tensor identities,
\begin{equation*}
\Phi_{K,m}(F)=\sum_k c_k\prod_{l\le q}{:}\mathcal{O}(f_{k,l}^{(m)}){:},
\end{equation*}
and similarly for $\Psi_{K,m}$. In particular, if $\sum_k c_k\bigotimes_l f_{k,l}=0$ as a function,
then $\Phi_{K,m}$ and $\Psi_{K,m}$ of it vanish, whatever the term structure of the sum, because they
read only the values of the function at tuples of lifted barycentres. Thus the right-hand sides depend
only on the function $F$, and $S_{q;K,m}$ and $S^T_{q;K,m}$ are well-defined linear functionals of $F$.

On a single separated tensor with supports in $B_{\rho}$, the identity
\begin{equation*}
\bigl\langle\Phi_{K,m}(f_1\otimes\cdots\otimes f_q)\bigr\rangle_{K,m}
=\Bigl\langle\prod_{l\le q}{:}\mathcal{O}(f_l^{(m)}){:}\Bigr\rangle_{K,m}
\end{equation*}
holds. The supports of the $f_l^{(m)}$ are pairwise disjoint: by
Lemma~\ref{7.4:lem-torus-transport}(ii) the torus distance between two of them equals the distance
between the corresponding supports of the $f_l$, which is positive since these are disjoint compact
sets. Hence the right-hand side is the centred moment
$S_{q;K,m}(\vec f^{(m)})=S_{q;K,m}(\vec f)$ of \eqref{7.4:eq-centred-moments-a}. Similarly,
\begin{equation*}
\Psi_{K,m}(f_1\otimes\cdots\otimes f_q)
=\kappa_q\bigl(\mathcal{O}(f_1^{(m)}),\dots,\mathcal{O}(f_q^{(m)})\bigr)
\end{equation*}
is the cumulant $S^T_{q;K,m}(\vec f)$ in the convention \eqref{7.4:eq-torus-transport-a}. Hence both
functionals extend the quantities previously defined on tensors.
\end{proof}

\begin{remark}
Every lattice quantity built from $S_{q;K,m}$ and $S^T_{q;K,m}$ in what follows is therefore a
functional of the test function itself. No appeal to the kernel theorem is needed for $C^1$ data.
\end{remark}

We first recall the countable family $\mathcal{T}^{\vee}$ of block and smeared observables on which the joint block/smeared limit state is built. This family is fixed, as a countable set, in the definition of the skeleton of block and smeared observables. By the uniform bound for the block and smeared observables, for each $A\in\mathcal{T}^{\vee}$ the sequence $(\langle A\rangle_j)_{j\ge j(A)}$ is bounded uniformly in $K$. Here $\langle\cdot\rangle_j$ is the expectation at the $j$-th cutoff pair $(K_j,m_j)$, and $j(A)$ is the index from which it is defined. The family defined below is adjoined to $\mathcal{T}^{\vee}$, so that a single diagonal argument over the countable index set
\[\mathfrak{I}:=\mathcal{T}^{\vee}\sqcup\mathfrak{F}\]
serves both.

\begin{definition}[A countable support-adapted family of test tuples]
\label{7.4:def-countable-family}
A \emph{closed rational box} is a set $\prod_{\mu=1}^{4}[a_\mu,b_\mu]\subset\mathbb{R}^4$ with rational
$a_\mu<b_\mu$. Let $\mathcal{Q}$ be the countable family of finite unions of closed rational boxes.

For $Q\in\mathcal{Q}$ put
\begin{equation*}
C^1_Q:=\{f\in C^1(\mathbb{R}^4;\mathbb{R}):\ f=0\ \text{on}\ \mathbb{R}^4\setminus Q\}.
\end{equation*}
We write
\[\|f\|_{C^1}:=\max\{\|f\|_\infty,\|\nabla f\|_\infty\},\qquad
\|\nabla f\|_\infty:=\sup_x|\nabla f(x)|.\]
The space $C^1_Q$ is separable for $\|\cdot\|_{C^1}$ (see the proof below). Fix a countable
$\|\cdot\|_{C^1}$-dense subset $D_Q\subset C^1_Q$, and define
\begin{equation}\label{7.4:eq-countable-family-a}
\begin{split}
\mathfrak{F}:=\bigcup_{n\ge 2}\bigl\{(\varphi_1,\dots,\varphi_n):{}&Q_1,\dots,Q_n\in\mathcal{Q}\
\text{pairwise disjoint},\\
&\varphi_i\in D_{Q_i}\ \text{for}\ 1\le i\le n\bigr\}.
\end{split}
\end{equation}

These objects have the following properties.
\begin{itemize}
\item[(i)] For $f\in C^1_Q$,
\begin{equation}\label{7.4:eq-countable-family-1}
\|f\|_{\sharp}=\max\{\|f\|_\infty,\,40M\|\nabla f\|_\infty\}\le 40M\|f\|_{C^1}.
\end{equation}
Consequently $D_Q$ is dense in $C^1_Q$ for $\|\cdot\|_{\sharp}$.
\item[(ii)] (Support adaptation.) Let $f\in C^1_c(\mathbb{R}^4;\mathbb{R})$ and let $W\supset\operatorname{supp}f$ be open. Then there is $Q\in\mathcal{Q}$ with $\operatorname{supp}f\subset Q\subset W$. Moreover $f\in C^1_Q$ is a $\|\cdot\|_{C^1}$-limit of elements of $D_Q$, all of which are supported in $Q\subset W$.
\end{itemize}
The family $\mathfrak{F}$ is countable. Each $\vec\varphi=(\varphi_1,\dots,\varphi_n)\in\mathfrak{F}$, built from pairwise disjoint $Q_1,\dots,Q_n$ as in \eqref{7.4:eq-countable-family-a}, lies in the set $\mathfrak{S}_n$ of separated $n$-tuples and satisfies
\begin{equation}\label{7.4:eq-countable-family-2}
\mathsf{d}(\vec\varphi)\ \ge\ \min_{i\ne k}\operatorname{dist}(Q_i,Q_k)\ >\ 0.
\end{equation}
\end{definition}

\begin{proof}
\emph{Separability of $C^1_Q$.} Since $Q$ is bounded, there is a closed box $Q'\supset Q$. Let $f\in C^1_Q$. Then $f$ vanishes on the open set $\mathbb{R}^4\setminus Q$, so $\nabla f$ vanishes there too. Hence the suprema of $|f|$ and $|\nabla f|$ over $\mathbb{R}^4$ equal those over $Q'$. The restriction map $f\mapsto f|_{Q'}$ is therefore a linear isometry from $(C^1_Q,\|\cdot\|_{C^1})$ into $C^1(Q')$ with the norm $\|\cdot\|_{C^1}$.

Its image is closed. Indeed, suppose $f_\ell|_{Q'}$ converges in $C^1(Q')$. Since $f_\ell$ and $\nabla f_\ell$ vanish off $Q$, the sequence $(f_\ell)_{\ell\ge1}$ is Cauchy for $\|\cdot\|_{C^1}$ on all of $\mathbb{R}^4$. Then $f_\ell$ and $\nabla f_\ell$ converge uniformly on $\mathbb{R}^4$, and the limit $f$ is of class $C^1$ with $\nabla f=\lim\nabla f_\ell$. It vanishes off $Q$, and its restriction is the given limit.

So $C^1_Q$ is isometric to a closed subspace of $C^1(Q')$. The space $C^1(Q')$ is separable, and a subspace of a separable metric space is separable, both by Remark~\ref{7.4:rem-standard-facts}. Hence $C^1_Q$ is separable, and the choice of $D_Q$ is possible.

\emph{(i).} For $f\in C^1_Q$ we have $\operatorname{Lip}f\le\|\nabla f\|_\infty$ by the mean value theorem along segments. Conversely, every directional derivative of $f$ is bounded by $\operatorname{Lip}f$. Hence $\operatorname{Lip}f=\|\nabla f\|_\infty$, and the equality in \eqref{7.4:eq-countable-family-1} is the definition of $\|\cdot\|_{\sharp}$.

Since $M=L^{m'}\ge L^4\ge 1$, we have $40M\ge 1$. Therefore
\[\|f\|_\infty\le 40M\|f\|_{C^1}\quad\text{and}\quad 40M\|\nabla f\|_\infty\le 40M\|f\|_{C^1},\]
which is the inequality in \eqref{7.4:eq-countable-family-1}.

Applied to $f-\varphi$ with $\varphi\in D_Q$, it gives
\[\|f-\varphi\|_{\sharp}\le 40M\|f-\varphi\|_{C^1}.\]
So the $\|\cdot\|_{C^1}$-density of $D_Q$ in $C^1_Q$ implies its $\|\cdot\|_{\sharp}$-density.

\emph{(ii).} Let $x\in\operatorname{supp}f$. Since $W$ is open, there is an open box with rational corners containing $x$ whose closure lies in $W$. These open boxes cover the compact set $\operatorname{supp}f$, so finitely many of them already cover it. Let $Q$ be the union of their closures. Then $Q\in\mathcal{Q}$ and $\operatorname{supp}f\subset Q\subset W$.

The function $f$ is real and of class $C^1$, and it vanishes off $\operatorname{supp}f$, hence off $Q$. Thus $f\in C^1_Q$. By the density of $D_Q$ it is a $\|\cdot\|_{C^1}$-limit of elements of $D_Q$. Each of these vanishes off the closed set $Q$, so its support lies in $Q\subset W$.

\emph{Countability of $\mathfrak{F}$.} In \eqref{7.4:eq-countable-family-a} there are countably many $n$. For each $n$ there are countably many tuples $(Q_1,\dots,Q_n)\in\mathcal{Q}^n$. For each such tuple, $D_{Q_1}\times\dots\times D_{Q_n}$ is countable. A countable union of countable sets is countable.

\emph{Separation.} Let $\vec\varphi\in\mathfrak{F}$ with $\varphi_i\in D_{Q_i}$ and $Q_1,\dots,Q_n$ pairwise disjoint. Each $\varphi_i$ is real, of class $C^1$, and supported in the compact set $Q_i$, so $\varphi_i\in C^1_c(\mathbb{R}^4;\mathbb{R})$. The supports satisfy $\operatorname{supp}\varphi_i\subset Q_i$, so they are pairwise disjoint, and
\[\operatorname{dist}(\operatorname{supp}\varphi_i,\operatorname{supp}\varphi_k)\ge\operatorname{dist}(Q_i,Q_k)\quad\text{for } i\ne k.\]
Taking the minimum over $i\ne k$ gives the first inequality in \eqref{7.4:eq-countable-family-2}. The second holds because the distance between two disjoint compact sets is positive and the minimum is over finitely many pairs. Hence $\vec\varphi\in\mathfrak{S}_n$.
\end{proof}

The family is made support-adapted, and not merely dense, because of how it is used in the approximation arguments that follow. There, the approximants of each entry $f_i$ of a separated tuple must keep their supports inside prescribed sets that are disjoint from the other supports and lie inside balls. In this way every tuple obtained by replacing some entries by their approximants is again a separated, ball-supported tuple, to which the per-ball volume-free bound applies. A plain dense sequence in $C^1_c(\mathbb{R}^4)$ does not control supports.

Lipschitz test functions are not admitted in this construction. The space of Lipschitz functions is not separable for the Lipschitz seminorm, and $C^1$ functions are not $\|\cdot\|_{\sharp}$-dense among the Lipschitz functions: a Lipschitz function with a kink is not such a limit (compare Lemma~\ref{7.4:lem-translation-modulus}). Nothing is lost by this restriction, since $C^1_c(\mathbb{R}^4)$ contains $\mathcal{D}(\mathbb{R}^4)$, which determines the distributions.

\begin{theorem}[The infinite-volume theorem]\label{7.4:thm-infinite-volume}
Let $L\ge L_0$ be odd, let $\mathfrak{p}$, $\gamma$ and $\gamma_H$ be fixed after $L$, let $r>0$ be
fixed first, and let
\begin{equation*}
0<g\le\min\{\gamma_H,\gamma_J(L,\mathfrak{p},r)\}.
\end{equation*}
Let $(g_0^{(j)})_{j\in\mathbb{N}}$ be a sequence in $(0,g_{-}(g)]$ decreasing to $0$, for instance the
subsequence of clause~(ii-b) of Theorem~\ref{3.1:thm-main}. Let $K_j:=K(g_0^{(j)},g)$ be the
first-passage cutoffs, so that $K_j\to\infty$ by clause~(0) of Theorem~\ref{3.1:thm-main}. Let
$(m_j)_{j\in\mathbb{N}}$ be integers with
$m_j\to\infty$ and $m_j\ge m_{\mathbb{T}}(g_{-}(g))$ for every $j$. Write
$\langle\cdot\rangle_j:=\langle\cdot\rangle_{K_j,m_j}$, $S^T_{n;j}:=S^T_{n;K_j,m_j}$ and
$S_{q;j}:=S_{q;K_j,m_j}$.

For $n\ge1$ let $\mathfrak{S}_n$ be the set of $n$-tuples $\vec f=(f_1,\dots,f_n)$ of real-valued
$C^1$ functions on $\mathbb{R}^4$ with compact, pairwise disjoint supports; for $n\ge2$ let
$\mathsf{d}(\vec f)>0$ be the least distance between two of these supports. On such tuples,
$S^T_{n;j}(\vec f)$ ($n\ge2$) and $S_{n;j}(\vec f)$ ($n\ge1$) are understood through the transport
to the torus, \eqref{7.4:eq-torus-transport-a}, and Definition~\ref{7.4:def-centred-moments}; they
are defined for every $j\ge j(\vec f)$, for some index $j(\vec f)$. For $a\in\mathbb{R}^4$ put
$\tau_af:=f(\cdot-a)$; for $w\in W_4$
put $(w\cdot f)(x):=f(w^{-1}x)$. Let $\theta$ be the time reflection $(x^0,\vec x)\mapsto(-x^0,\vec x)$
and $e_0$ the unit vector in the time direction.

Assume:
\begin{enumerate}
\item[(a)] the correlation theorem, in the standing setting of
Definition~\ref{7.1:not-torus-cumulants}, in which the gauge group is $\mathrm{SU}(2)$, that is
$N=2$, and the hypotheses \textup{(H1)}--\textup{(H7)} hold;
\item[(b)] reflection positivity of the lattice measure;
\item[(c)] the per-ball volume-free bound;
\item[(d)] the extraction of the limit state of block and smeared observables, in the following form.
Let $\mathcal{T}^{\vee}$ be the countable family of block and smeared observables on which that state
is constructed; for $A\in\mathcal{T}^{\vee}$ the expectation $\langle A\rangle_j$ is defined for every
$j\ge j(A)$, for some index $j(A)$. For every $A\in\mathcal{T}^{\vee}$ the sequence
$(\langle A\rangle_j)_{j\ge j(A)}$ is bounded uniformly in the cutoff. Along every infinite subset of
$\mathbb{N}$ along which $\langle A\rangle_j$ converges for every $A\in\mathcal{T}^{\vee}$, the joint
limit state $\omega_{\infty}$ exists on the quasilocal algebra
$\mathfrak{A}^{\tau}_{\mathrm{loc}}\otimes\mathfrak{A}^{\mathrm{sm}}_{\mathrm{loc}}$ of block and
smeared observables; its construction uses only this convergence together with the density of
$\mathcal{T}^{\vee}$ in that algebra, uniform in the cutoff.
\end{enumerate}
Then there is an infinite set $\mathcal{J}\subset\mathbb{N}$ with the following properties. It is
obtained by a single diagonal extraction over the countable set
$\mathcal{T}^{\vee}\sqcup\mathfrak{F}$, where $\mathfrak{F}$ is the family of
Definition~\ref{7.4:def-countable-family}. Along this same $\mathcal{J}$ the joint limit state
$\omega_{\infty}$ exists. Moreover, the following hold.
\begin{enumerate}
\item[(A)] \emph{Existence, bounds, structure.} For every $n\ge2$ and every $\vec f\in\mathfrak{S}_n$
the limit
\begin{equation*}
S^T_n(\vec f):=\lim_{j\in\mathcal{J}}S^T_{n;j}(\vec f)
\end{equation*}
exists, the terms being defined for $j\ge j(\vec f)$. The functional $S^T_n$ has the following
properties.
\begin{itemize}
\item It is multilinear on separated tuples, in each argument among the functions that keep the tuple
in $\mathfrak{S}_n$.
\item It is symmetric under permutations of its arguments.
\item It is real-valued.
\item Suppose each $f_i$ is supported in a ball of radius $R$ and the supports are pairwise at
distance at least $\mathsf{d}$. Then
\begin{equation}\label{7.4:eq-infinite-volume-a}
|S^T_n(\vec f)|\le C^{\infty}_n(\mathsf{d};R)\prod_{i=1}^n\|f_i\|_{\sharp},
\end{equation}
where $C^{\infty}_n(\mathsf{d};R)$ is the constant of the per-ball volume-free bound.
\end{itemize}
Let $V_1,\dots,V_n\subset\mathbb{R}^4$ be open sets pairwise at positive distance. There is a unique
$T_V\in\mathcal{D}'(V_1\times\cdots\times V_n)$ with
$T_V(\varphi_1\otimes\cdots\otimes\varphi_n)=S^T_n(\vec\varphi)$ for real-valued
$\varphi_i\in\mathcal{D}(V_i)$; the right side is extended complex-multilinearly to complex-valued
$\varphi_i$. These distributions glue to a unique $S^T_n\in\mathcal{D}'((\mathbb{R}^4)^n_{\neq})$, where
\begin{equation*}
(\mathbb{R}^4)^n_{\neq}:=\{(x_1,\dots,x_n)\in(\mathbb{R}^4)^n:\ x_i\neq x_k\ \text{for all}\ i\neq k\}.
\end{equation*}
This $S^T_n$ has order at most $5n+1$ on every product of pairwise separated balls, so its order is
locally finite. Convergence of $S^T_{n;j}(\vec f)$ for tuples of merely Lipschitz functions is not
asserted.
\item[(B)] \emph{Invariance.} For every $n$, every $\vec f\in\mathfrak{S}_n$, $a\in\mathbb{R}^4$,
$w\in W_4$ and every permutation $\sigma$ of $\{1,\dots,n\}$,
\begin{equation*}
\begin{aligned}
S^T_n(\tau_af_1,\dots,\tau_af_n)&=S^T_n(\vec f)=S^T_n(w\cdot f_1,\dots,w\cdot f_n)\\
&=S^T_n(f_{\sigma(1)},\dots,f_{\sigma(n)}).
\end{aligned}
\end{equation*}
The distributions $S^T_n$ are invariant under the diagonal action of $\mathbb{R}^4\rtimes W_4$ on
$(\mathbb{R}^4)^n_{\neq}$. Invariance under $O(4)$ is not asserted.
\item[(C)] \emph{Osterwalder--Schrader positivity for the centred field.} Put $S_0:=1$ and $S_1:=0$.
Put $\mathcal{T}^{\mathrm{sep}}_{0,+}:=\mathbb{C}$. For $q\ge1$ let
$\mathcal{T}^{\mathrm{sep}}_{q,+}$ be the complex span of the tensors $f_1\otimes\cdots\otimes f_q$ with
$f_a\in C^1_c(\mathbb{R}^4;\mathbb{R})$ whose supports are compact, pairwise disjoint and contained in
$\{x^0>0\}$. Let $\mathcal{P}_+:=\bigoplus_{q\ge0}\mathcal{T}^{\mathrm{sep}}_{q,+}$, whose elements
$F=(F_q)_q$ have finitely many non-zero components. Put
\begin{equation*}
(\Theta\bar F)_q(x_1,\dots,x_q):=\overline{F_q(\theta x_1,\dots,\theta x_q)}.
\end{equation*}
For $q\ge2$ and $\vec h\in\mathfrak{S}_q$, and for $F,G\in\mathcal{P}_+$, put
\begin{equation}\label{7.4:eq-infinite-volume-b}
\begin{gathered}
S_q(\vec h):=\sum_{\pi\in\operatorname{Part}_{\ge2}[q]}\ \prod_{B\in\pi}
S^T_{|B|}\bigl((h_l)_{l\in B}\bigr),\\
(F,G)_{\mathrm{OS}}:=\sum_{q,q'}S_{q+q'}\bigl((\Theta\bar F)_q\otimes G_{q'}\bigr).
\end{gathered}
\end{equation}
Let $\mathcal{T}^{\mathrm{sep}}_q$ be the complex span of the tensors $h_1\otimes\cdots\otimes h_q$
with $\vec h\in\mathfrak{S}_q$, regarded as functions on $(\mathbb{R}^4)^q$. By the linearity in the test
function (Lemma~\ref{7.4:lem-linearity}), $S_q$ extends linearly and unambiguously to
$\mathcal{T}^{\mathrm{sep}}_q$, and $S_{q;j}\to S_q$ there as $j\to\infty$ in $\mathcal{J}$.
Since the reflected factors are supported in $\{x^0<0\}$, $(\Theta\bar F)_q\otimes G_{q'}$ is a
separated tensor, that is, an element of $\mathcal{T}^{\mathrm{sep}}_{q+q'}$, so the second line of
\eqref{7.4:eq-infinite-volume-b} is defined.

Then $(\cdot,\cdot)_{\mathrm{OS}}$ is sesquilinear, Hermitian and positive semidefinite on
$\mathcal{P}_+$, and
\begin{equation*}
(F,G)_{\mathrm{OS}}=\lim_{j\in\mathcal{J}}
\bigl\langle\overline{\mathbb{F}_j(\theta\cdot U)}\,\mathbb{G}_j(U)\bigr\rangle_j .
\end{equation*}
Here $\mathbb{F}_j:=\sum_q\Phi_{K_j,m_j}(F_q)$ and $\mathbb{G}_j:=\sum_q\Phi_{K_j,m_j}(G_q)$ are the
centred lattice polynomials of Lemma~\ref{7.4:lem-linearity}, and $\theta\cdot U$ is the bond
configuration transformed by the time reflection of $\mathbb{T}_{m_j}$ in the plane pair
$\{x^0=0\}\cup\{x^0=L^{m_j}\}$, acting on bond variables as in Lemma~\ref{1.1:lem-invariance}.
Moreover, the positivity condition $(\mathrm{E}2)$ of \cite{OS73} holds on distributions. Put
\begin{equation*}
(\mathbb{R}^4)^q_{\neq,+}:=\{x\in(\mathbb{R}^4)^q_{\neq}:\ x_a^0>0\ \text{for all}\ a\}.
\end{equation*}
Then for every finite sequence $(F_q)_q$ with $F_q\in\mathcal{D}((\mathbb{R}^4)^q_{\neq,+})$,
\begin{equation*}
\sum_{q,q'}S_{q+q'}\bigl((\Theta\bar F)_q\otimes F_{q'}\bigr)\ge0,
\end{equation*}
where $S_{q+q'}\in\mathcal{D}'((\mathbb{R}^4)^{q+q'}_{\neq})$ now denotes the distribution given by
the first line of \eqref{7.4:eq-infinite-volume-b} with the distributions $S^T_{|B|}$ of (A),
tensored over the blocks and restricted to $(\mathbb{R}^4)^{q+q'}_{\neq}$.
\item[(D)] \emph{Osterwalder--Schrader reconstruction.} Put
\begin{equation}\label{7.4:eq-infinite-volume-c}
\begin{gathered}
\mathcal{N}:=\{F\in\mathcal{P}_+:\ (F,F)_{\mathrm{OS}}=0\},\qquad
\mathcal{H}^{\mathrm{curv}}:=\text{the completion of}\ \mathcal{P}_+/\mathcal{N},\\
\Omega:=[(1,0,0,\dots)],\\
(\mathsf{T}_tF)_q(x_1,\dots,x_q):=F_q(x_1-te_0,\dots,x_q-te_0)\qquad(t\ge0),\\
(\mathsf{S}_{\vec a}F)_q(x_1,\dots,x_q):=F_q\bigl(x_1-(0,\vec a),\dots,x_q-(0,\vec a)\bigr)
\qquad(\vec a\in\mathbb{R}^3).
\end{gathered}
\end{equation}
Then $\mathcal{H}^{\mathrm{curv}}$ is a separable Hilbert space and $\|\Omega\|=1$.
\begin{enumerate}
\item[(D1)] For real $t\ge0$, $\mathsf{T}_t$ induces a well-defined operator $T_t$ on
$\mathcal{H}^{\mathrm{curv}}$. The family $(T_t)_{t\ge0}$ is a strongly continuous semigroup of
positive self-adjoint contractions with $T_t\Omega=\Omega$ and $\ker T_t=\{0\}$. There is a unique
self-adjoint operator $H\ge0$, acting in the whole of $\mathcal{H}^{\mathrm{curv}}$, with
$T_t=e^{-tH}$ for $t\ge0$. Moreover $H=-\log T_1$ and $H\Omega=0$.
\item[(D2)] $\mathsf{S}_{\vec a}$ induces a strongly continuous unitary representation
$\vec a\mapsto\mathcal{U}(\vec a)$ of $\mathbb{R}^3$ with $\mathcal{U}(\vec a)\Omega=\Omega$ and
$\mathcal{U}(\vec a)T_t=T_t\mathcal{U}(\vec a)$. Its generators $\vec P=(P_1,P_2,P_3)$, obtained from
the projection-valued measure of the representation, are self-adjoint and strongly commute pairwise
and with $H$. The joint spectrum of $(H,\vec P)$ lies in $[0,\infty)\times\mathbb{R}^3$.
\item[(D3)] $W_3$, the group of signed permutations of $x^1,x^2,x^3$, acts by unitaries
$\mathcal{U}(w)$. These fix $\Omega$, commute with $T_t$ and $H$, and satisfy
$\mathcal{U}(w)\mathcal{U}(\vec a)\mathcal{U}(w)^*=\mathcal{U}(w\vec a)$. The joint
spectrum of $(H,\vec P)$ is $W_3$-symmetric.
\end{enumerate}
\item[(E)] \emph{Not asserted.} The theorem asserts nothing about the following:
\begin{itemize}
\item the case $n=1$, or anything at coinciding points;
\item temperedness (the extension of $S^T_n$ to test spaces of Schwartz type), or the linear growth
condition $(\mathrm{E}0')$ of \cite{OS75};
\item field operators or Wightman functions;
\item uniqueness of $\Omega$, clustering, or a mass gap;
\item invariance under $O(4)$, or the relativistic spectrum condition $H\ge|\vec P|$;
\item that $\dim\mathcal{H}^{\mathrm{curv}}>1$, or that $S^T_2\not\equiv0$ (both are consequences
stated in clause (iv) of Theorem~\ref{3.1:thm-main} and drawn in Chapter~\ref{ch:8});
\item uniqueness of $\mathcal{J}$, independence of $S^T_n$ from the choice of $(K_j,m_j)$, or
independence from the regularisation;
\item convergence on merely Lipschitz tuples;
\item any identification of $\mathcal{H}^{\mathrm{curv}}$ with the space
$\mathcal{H}^{\mathrm{sm}}$ reconstructed from $\omega_{\infty}$, or with any other Hilbert space
reconstructed in this book;
\item any mixed correlation of block and curvature observables.
\end{itemize}
\end{enumerate}
\end{theorem}

\begin{proof}[Proof of Theorem~\ref{7.4:thm-infinite-volume}: the choice of $\mathcal{J}$]
We first choose $\mathcal{J}$ by a single diagonal extraction. Let
\begin{equation*}
\mathfrak{I}:=\mathcal{T}^{\vee}\sqcup\mathfrak{F}
\end{equation*}
be the disjoint union. It is countable: $\mathcal{T}^{\vee}$ is countable by hypothesis~(d), and
$\mathfrak{F}$ is countable by Definition~\ref{7.4:def-countable-family}. For $s\in\mathfrak{I}$ put
\begin{equation*}
v_j(s):=
\begin{cases}
\langle A\rangle_j, & s=A\in\mathcal{T}^{\vee},\\
S^T_{n;j}(\vec g), & s=\vec g\in\mathfrak{F}\cap\mathfrak{S}_n ,
\end{cases}
\end{equation*}
defined for $j\ge j(s)$. For $s=A\in\mathcal{T}^{\vee}$, $j(A)$ is the index from hypothesis~(d).

Now let $s=\vec g\in\mathfrak{F}\cap\mathfrak{S}_n$. Fix $\rho>0$ such that
$\bigcup_i\operatorname{supp}g_i\subset B_{\rho}$; such a $\rho$ exists because the supports are
compact. Each $g_i$ is then supported in a ball of radius $R:=\rho$, namely $B_{\rho}$. By the
definition of $\mathfrak{F}$ the supports are pairwise at distance at least
$\mathsf{d}:=\mathsf{d}(\vec g)>0$. Choose $j(\vec g)$ such that for every $j\ge j(\vec g)$,
\begin{equation*}
L^{m_j}\ge2\rho+2,\qquad m_j\ge\max\{m_{\mathbb{T}}(g_{-}(g)),\ \log_L\ell(n,\rho,g)\}.
\end{equation*}
No condition on $K_j$ is needed, since the per-ball volume-free bound holds at every first-passage
cutoff $K(g_0,g)$ with $g_0\in(0,g_{-}(g)]$. Such a $j(\vec g)$ exists because $m_j\to\infty$.

Next, the sequences are bounded. Let $\vec g\in\mathfrak{F}\cap\mathfrak{S}_n$ and $j\ge j(\vec g)$.
Since $L^{m_j}\ge2\rho+2$, the transport lemma (Lemma~\ref{7.4:lem-torus-transport}) applies to each
$g_i$, and gives:
\begin{itemize}
\item $g_i^{(m_j)}$ is a real Lipschitz function on $\mathbb{T}_{m_j}$ with
$\|g_i^{(m_j)}\|_{\sharp}=\|g_i\|_{\sharp}$ (each $g_i$, being $C^1$ with compact support, is
Lipschitz);
\item $g_i^{(m_j)}$ is supported in the torus ball of radius $\rho$ about $q_{m_j}(0)$;
\item the supports of the $g_i^{(m_j)}$ are pairwise at $d_{\mathbb{T}}$-distance at least
$\mathsf{d}(\vec g)$.
\end{itemize}
The remaining hypotheses of the per-ball volume-free bound also hold: $g_0^{(j)}\in(0,g_{-}(g)]$,
$K_j=K(g_0^{(j)},g)$ is the first-passage cutoff, and
$m_j\ge\max\{m_{\mathbb{T}}(g_{-}(g)),\log_L\ell(n,\rho,g)\}$. By hypothesis~(c), applied to the
tuple $(g_i^{(m_j)})_{i\le n}$ with $R=\rho$ and $\mathsf{d}=\mathsf{d}(\vec g)$, and by the convention
\eqref{7.4:eq-torus-transport-a}, we get
\begin{equation*}
|v_j(\vec g)|=\bigl|S^T_{n;K_j,m_j}\bigl((g_i^{(m_j)})_{i\le n}\bigr)\bigr|
\le C^{\infty}_n(\mathsf{d}(\vec g);\rho)\prod_{i=1}^n\|g_i\|_{\sharp}
\qquad(j\ge j(\vec g)).
\end{equation*}
The constant is free of $K$, $g_0$, $m$ and of the torus side, so the right side does not depend on
$j$. For $A\in\mathcal{T}^{\vee}$ the sequence $(v_j(A))_{j\ge j(A)}$ is bounded by hypothesis~(d).

Every sequence $(v_j(s))_{j\ge j(s)}$, $s\in\mathfrak{I}$, is therefore a bounded complex sequence,
and $\mathfrak{I}$ is countable. The diagonal argument recalled in
Remark~\ref{7.4:rem-standard-facts} therefore yields an infinite set
\begin{equation}\label{7.4:eq-infinite-volume-d}
\mathcal{J}\subset\mathbb{N}\quad\text{such that}\quad
\bigl(v_j(s)\bigr)_{j\in\mathcal{J},\,j\ge j(s)}\ \text{converges for every}\ s\in\mathfrak{I}.
\end{equation}

The half of \eqref{7.4:eq-infinite-volume-d} indexed by $\mathcal{T}^{\vee}$ says that
$\langle A\rangle_j$ converges along $\mathcal{J}$ for every $A\in\mathcal{T}^{\vee}$. By
hypothesis~(d), the joint limit state $\omega_{\infty}$ on
$\mathfrak{A}^{\tau}_{\mathrm{loc}}\otimes\mathfrak{A}^{\mathrm{sm}}_{\mathrm{loc}}$ therefore exists
along this $\mathcal{J}$. Indeed, its construction uses only the convergence on $\mathcal{T}^{\vee}$
and the density of $\mathcal{T}^{\vee}$ uniform in the cutoff, and both are intact. Adjoining
$\mathfrak{F}$ to the index set only enlarges the family of sequences made convergent and changes
nothing in that construction.

Further arguments may pass to an infinite subset of $\mathcal{J}$. Every conclusion below persists
along such a subset, since subsequences of convergent sequences converge to the same limits. The half
of \eqref{7.4:eq-infinite-volume-d} indexed by $\mathfrak{F}$ is the input of the existence part of
(A), proved next. No mixed expectation
$\langle A\cdot\Phi_{K_j,m_j}(F)\rangle_j$ of an element $A\in\mathcal{T}^{\vee}$ and a centred lattice
polynomial is included in $\mathfrak{I}$; the reason why no such quantity can be adjoined is explained
in the remark of this chapter on mixed correlations of block and curvature observables.

This establishes $\mathcal{J}$ as one diagonal extraction over
$\mathcal{T}^{\vee}\sqcup\mathfrak{F}$, along which $\omega_{\infty}$ exists. Parts (A)--(D) are
proved from this choice in the lemmas that follow in this chapter, among them
Lemma~\ref{7.7:prf-reconstructed-space} for (D).
\end{proof}

\begin{lemma}[Equicontinuity and existence of the limits]\label{7.4:prf-equicontinuity}
Let $L$, $\mathfrak{p}$, $\gamma$, $\gamma_H$, $r$ and $g$ be as in the standing data of
Theorem~\ref{7.4:thm-infinite-volume}, and let $\mathcal{J}$ be the diagonal subsequence of
\eqref{7.4:eq-infinite-volume-d}.
\begin{itemize}
\item[(i)] Let $n\ge 2$, $R'>0$, $\mathsf{d}'>0$ and $\rho>0$, and let
$Q_1,\dots,Q_n\subset\mathbb{R}^4$ be compact sets such that:
\begin{itemize}
\item each $Q_i$ is contained in a ball of radius $R'$;
\item the $Q_i$ are pairwise at distance at least $\mathsf{d}'$;
\item $Q_1\cup\dots\cup Q_n\subset B_{\rho}$.
\end{itemize}
Let $f_i$, $\varphi_i$ ($i\le n$) be real Lipschitz functions on $\mathbb{R}^4$ supported in $Q_i$.
Let $(K,m)$ satisfy
\begin{equation*}
m\ge\max\{m_{\mathbb{T}}(g_{-}(g)),\ \log_L\ell(n,R',g)\},\qquad L^m\ge 2\rho+2,
\end{equation*}
and let $K=K(g_0,g)$ be the first-passage cutoff of a bare coupling $g_0\in\,]0,g_{-}(g)]$.
Then
\begin{equation}\label{7.4:eq-equicontinuity-a}
\bigl|S^T_{n;K,m}(\vec f)-S^T_{n;K,m}(\vec\varphi)\bigr|
\le C^{\infty}_n(\mathsf{d}';R')\sum_{i\le n}\|f_i-\varphi_i\|_{\sharp}
\prod_{l<i}\|\varphi_l\|_{\sharp}\prod_{l>i}\|f_l\|_{\sharp}.
\end{equation}
\item[(ii)] Let $n\ge 2$ and $\vec f\in\mathfrak{S}_n$, put $\mathsf{d}:=\mathsf{d}(\vec f)$, choose
$\rho>0$ with $\operatorname{supp}f_1\cup\dots\cup\operatorname{supp}f_n\subset B_{\rho}$, and put
$R':=\rho+\mathsf{d}/3$. For $\eta>0$ put
\begin{equation}\label{7.4:eq-equicontinuity-1}
\Delta_{\vec f}(\eta):=C^{\infty}_n(\mathsf{d}/3;R')\sum_{i\le n}40M\eta
\prod_{l\neq i}\bigl(\|f_l\|_{\sharp}+40M\eta\bigr).
\end{equation}
Then $\Delta_{\vec f}(\eta)\to 0$ as $\eta\downarrow 0$. For every $\eta>0$ there are
$\vec\varphi\in\mathfrak{F}$ with $\|f_i-\varphi_i\|_{\sharp}\le 40M\eta$ for $i\le n$, and an index
$j_0(\vec f,\eta)$, such that
\begin{equation*}
\bigl|S^T_{n;j}(\vec f)-S^T_{n;j}(\vec\varphi)\bigr|\le\Delta_{\vec f}(\eta)
\qquad\text{for all }j\in\mathcal{J}\text{ with }j\ge j_0(\vec f,\eta).
\end{equation*}
The limit
\begin{equation}\label{7.4:eq-equicontinuity-b}
S^T_n(\vec f):=\lim_{j\in\mathcal{J}}S^T_{n;j}(\vec f),\qquad\vec f\in\mathfrak{S}_n,\ n\ge 2,
\end{equation}
exists.
\end{itemize}
\end{lemma}

\begin{proof}
\emph{Proof of} (i).
Throughout, $S^T_{n;K,m}$ is evaluated on test functions on $\mathbb{R}^4$ through the convention
\eqref{7.4:eq-torus-transport-a}: the tuple is first transported to the torus $\mathbb{T}_m$ by
$f\mapsto f^{(m)}$. This is legitimate because all the functions involved have compact support in
$B_{\rho}$ and $L^m\ge 2\rho+2$.

For $i=1,\dots,n$ introduce the hybrid tuples
\begin{equation*}
\vec h^{(i)}:=(\varphi_1,\dots,\varphi_{i-1},\,f_i-\varphi_i,\,f_{i+1},\dots,f_n).
\end{equation*}
The transport $f\mapsto f^{(m)}$ is linear, because $f^{(m)}(q_m(x))=f(x)$ on $B_{\rho}$ and
$f^{(m)}=0$ off $q_m(B_{\rho})$. The observable $\mathcal{O}(h)$ is linear in $h$. The joint
cumulant $\kappa_n$ of bounded random variables is multilinear, by the standard facts recalled in
Remark~\ref{7.4:rem-standard-facts}; equivalently, by linearity in the test function
(Lemma~\ref{7.4:lem-linearity}), $S^T_{n;K,m}$ is linear in each entry. Hence, for each $i$,
\begin{equation*}
S^T_{n;K,m}(\varphi_1,\dots,\varphi_{i-1},f_i,\dots,f_n)
-S^T_{n;K,m}(\varphi_1,\dots,\varphi_i,f_{i+1},\dots,f_n)
=S^T_{n;K,m}(\vec h^{(i)}).
\end{equation*}
The left-hand sides telescope as $i$ runs from $1$ to $n$. The first tuple for $i=1$ is $\vec f$,
and the second tuple for $i=n$ is $\vec\varphi$. Therefore
\begin{equation*}
S^T_{n;K,m}(\vec f)-S^T_{n;K,m}(\vec\varphi)=\sum_{i\le n}S^T_{n;K,m}(\vec h^{(i)}).
\end{equation*}

Fix $i$. The $i$-th entry $f_i-\varphi_i$ of $\vec h^{(i)}$ is real, Lipschitz and supported in
$Q_i$. For $l<i$ the $l$-th entry is $\varphi_l$, supported in $Q_l$; for $l>i$ it is $f_l$,
supported in $Q_l$.
So the $l$-th entry of $\vec h^{(i)}$ is a real Lipschitz function supported in $Q_l$ for every $l$.

By the transport lemma (Lemma~\ref{7.4:lem-torus-transport}), applied with $\rho$ and $m$ as above:
\begin{itemize}
\item the transported entries are real Lipschitz functions on $\mathbb{T}_m$ with the same
$\|\cdot\|_{\sharp}$-norms;
\item each is supported in a torus ball of radius $R'$, because $Q_l$ lies in a ball of radius
$R'$;
\item their supports are pairwise at $d_{\mathbb{T}}$-distance at least $\mathsf{d}'$, because
torus distances between the transported supports equal the distances between the supports in
$\mathbb{R}^4$, which are at least $\mathsf{d}'$ since the supports lie in the $Q_l$.
\end{itemize}

The per-ball volume-free bound applies at every $m\ge\max\{m_{\mathbb{T}}(g_{-}(g)),
\log_L\ell(n,R',g)\}$ and every first-passage cutoff $K=K(g_0,g)$ with $g_0\in\,]0,g_{-}(g)]$, to
real Lipschitz tuples on $\mathbb{T}_m$ of this kind. Its constant $C^{\infty}_n(\mathsf{d}';R')$ is
free of $K$, $g_0$, $m$ and of the positions of the supports. It gives
\begin{equation*}
\bigl|S^T_{n;K,m}(\vec h^{(i)})\bigr|
\le C^{\infty}_n(\mathsf{d}';R')\,\|f_i-\varphi_i\|_{\sharp}
\prod_{l<i}\|\varphi_l\|_{\sharp}\prod_{l>i}\|f_l\|_{\sharp}.
\end{equation*}
Summing over $i$ and using the telescoped identity yields \eqref{7.4:eq-equicontinuity-a}.

\smallskip\noindent
\emph{Proof of} (ii).
We now construct the limits asserted in Theorem~\ref{7.4:thm-infinite-volume} along the diagonal
sequence $\mathcal{J}$ of \eqref{7.4:eq-infinite-volume-d}.
Recall that $\mathfrak{S}_n$ is the set of $n$-tuples of real compactly supported $C^1$ functions on
$\mathbb{R}^4$ with pairwise disjoint supports. Let $\vec f$, $\mathsf{d}$ and $\rho$ be as in (ii);
$\mathsf{d}$ is positive because the supports are disjoint compact sets.

Apply the support-adapted property of the family $\mathcal{Q}$ in
\eqref{7.4:eq-countable-family-a} (Definition~\ref{7.4:def-countable-family}) to $f_i$, with the
open $\mathsf{d}/3$-neighbourhood of
$\operatorname{supp}f_i$ as the open set. It gives $Q_i\in\mathcal{Q}$ with
$\operatorname{supp}f_i\subset Q_i$ and every point of $Q_i$ at distance less than $\mathsf{d}/3$
from $\operatorname{supp}f_i$. Consequently:
\begin{itemize}
\item For $i\neq k$, a point of $Q_i$ and a point of $Q_k$ lie within $\mathsf{d}/3$ of
$\operatorname{supp}f_i$ and of $\operatorname{supp}f_k$ respectively. These supports are at
distance at least $\mathsf{d}$, so $Q_1,\dots,Q_n$ are pairwise at distance at least
$\mathsf{d}/3$, and in particular pairwise disjoint.
\item Each $Q_i$ is contained in the ball $B_{\rho+\mathsf{d}/3}$, which has radius
$R'=\rho+\mathsf{d}/3$.
\item $Q_1\cup\dots\cup Q_n\subset B_{\rho'}$ with $\rho':=\rho+\mathsf{d}/3$.
\end{itemize}

Let $\eta>0$. Since $f_i$ is a real $C^1$ function vanishing off $Q_i$, we have
$f_i\in C^1_{Q_i}$. We may therefore pick $\varphi_i\in D_{Q_i}$ with
$\|f_i-\varphi_i\|_{C^1}\le\eta$.
The $Q_i$ are pairwise disjoint, so $\vec\varphi=(\varphi_1,\dots,\varphi_n)\in\mathfrak{F}$ by the
definition of $\mathfrak{F}$ in \eqref{7.4:eq-countable-family-a}. By the comparison
$\|f\|_{\sharp}\le 40M\|f\|_{C^1}$ recorded there,
$\|f_i-\varphi_i\|_{\sharp}\le 40M\eta$. Since $\|\cdot\|_{\sharp}$ satisfies the triangle
inequality, $\|\varphi_l\|_{\sharp}\le\|f_l\|_{\sharp}+40M\eta$ for every $l$. Each $\varphi_i$,
being $C^1$ with compact support, is real Lipschitz and supported in $Q_i$.

Let $\mathcal{J}$ be the diagonal subsequence of \eqref{7.4:eq-infinite-volume-d}, and write
$S^T_{n;j}:=S^T_{n;K_j,m_j}$. Choose $j_0(\vec f,\eta)$ such that every $j\in\mathcal{J}$ with
$j\ge j_0(\vec f,\eta)$ satisfies
\begin{equation*}
m_j\ge\max\{m_{\mathbb{T}}(g_{-}(g)),\ \log_L\ell(n,R',g)\},\qquad L^{m_j}\ge 2\rho'+2,
\end{equation*}
together with $j\ge j(\vec\varphi)$, the index from which the sequence attached to $\vec\varphi$ in
the diagonal argument \eqref{7.4:eq-infinite-volume-d} is defined. This is possible since
$m_j\to\infty$. For such $j$ the cutoff $K_j=K(g_0^{(j)},g)$ is a first-passage cutoff with
$g_0^{(j)}\in\,]0,g_{-}(g)]$. Hence part (i) applies with
$(n,R',\mathsf{d}/3,\rho')$ in place of $(n,R',\mathsf{d}',\rho)$, and gives, by the definition
\eqref{7.4:eq-equicontinuity-1},
\begin{align*}
\bigl|S^T_{n;j}(\vec f)-S^T_{n;j}(\vec\varphi)\bigr|
&\le C^{\infty}_n(\mathsf{d}/3;R')\sum_{i\le n}40M\eta
\prod_{l\neq i}\bigl(\|f_l\|_{\sharp}+40M\eta\bigr)
=\Delta_{\vec f}(\eta).
\end{align*}
The right-hand side does not depend on $j$, and $\Delta_{\vec f}(\eta)\to 0$ as $\eta\downarrow 0$.
Moreover $(S^T_{n;j}(\vec\varphi))_{j\in\mathcal{J}}$ converges, by the choice of $\mathcal{J}$ in
\eqref{7.4:eq-infinite-volume-d}, because $\vec\varphi\in\mathfrak{F}$.

The sequence $(S^T_{n;j}(\vec f))_{j\in\mathcal{J}}$ is real (see the reality property below).
The last estimate places $S^T_{n;j}(\vec f)$ within $\Delta_{\vec f}(\eta)$ of the convergent
sequence $(S^T_{n;j}(\vec\varphi))_{j\in\mathcal{J}}$.
Hence the real sequence $(S^T_{n;j}(\vec f))_{j\in\mathcal{J}}$ is bounded, and
\begin{equation*}
\limsup_{j\in\mathcal{J}}S^T_{n;j}(\vec f)-\liminf_{j\in\mathcal{J}}S^T_{n;j}(\vec f)
\le 2\Delta_{\vec f}(\eta).
\end{equation*}
This holds for every $\eta>0$, so upper and lower limits coincide. The sequence is therefore Cauchy
and convergent, and the limit \eqref{7.4:eq-equicontinuity-b} exists.
\end{proof}

\begin{proof}[Proof of Theorem~\ref{7.4:thm-infinite-volume}: reality, symmetry, multilinearity
and bound of the limits]
The functionals \eqref{7.4:eq-equicontinuity-b} have the following properties.

\smallskip\noindent
\emph{Reality.} Each $S^T_{n;j}(\vec f)$ is a cumulant of real bounded random variables, hence
real. So is its limit.

\smallskip\noindent
\emph{Symmetry.} The joint cumulant $\kappa_n$ is symmetric in its arguments
(Remark~\ref{7.4:rem-standard-facts}). Hence $S^T_{n;j}$ is symmetric under permutations of the
entries of $\vec f$, and so is the limit, $\mathfrak{S}_n$ being invariant under such permutations.

\smallskip\noindent
\emph{Multilinearity on separated tuples.} The cumulant $\kappa_n$ is multilinear, so a linear
relation in one entry among tuples of $\mathfrak{S}_n$ holds for every $S^T_{n;j}$. It therefore
persists in the limit, since limits of multilinear quantities are multilinear. By linearity in the
test function (Lemma~\ref{7.4:lem-linearity}), more is true. By the definition of $\Psi_{K,m}$ in
\eqref{7.4:eq-linearity-a} and the multilinearity of $\kappa_n$, we have
$\Psi_{K_j,m_j}(f_1\otimes\dots\otimes f_n)=S^T_{n;j}(\vec f)$. Let $F$ be a finite sum of such
tensors in the span $\mathcal{T}^{\mathrm{sep}}_n$ of separated real-$C^1$ tensors. Then
$\Psi_{K_j,m_j}(F)$ is the corresponding finite combination of convergent sequences. Consequently
$S^T_n$ extends linearly to $\mathcal{T}^{\mathrm{sep}}_n$ by
\begin{equation*}
S^T_n(F)=\lim_{j\in\mathcal{J}}\Psi_{K_j,m_j}(F).
\end{equation*}
This is well defined on functions, not merely on their representations as sums of tensors, because
$\Psi_{K,m}(F)$ is defined through the values of the function $F$.

\smallskip\noindent
\emph{Bound.} Suppose each $f_i$ is supported in a ball of radius $R$ and the supports are pairwise
at distance at least $\mathsf{d}$. The per-ball volume-free bound, applied to the transported tuple
(Lemma~\ref{7.4:lem-torus-transport}), gives
\begin{equation*}
\bigl|S^T_{n;j}(\vec f)\bigr|\le C^{\infty}_n(\mathsf{d};R)\prod_{i\le n}\|f_i\|_{\sharp}
\end{equation*}
for all $j\in\mathcal{J}$ large enough that
\begin{equation*}
m_j\ge\max\{m_{\mathbb{T}}(g_{-}(g)),\ \log_L\ell(n,R,g)\},\qquad L^{m_j}\ge 2\rho+2;
\end{equation*}
such $j$ exist since $m_j\to\infty$. Passing to
the limit gives the bound \eqref{7.4:eq-infinite-volume-a}.
\end{proof}

\begin{remark}[Every cluster point is a limit along a subsequence]
Let $\Xi$ be a pointwise cluster point, on
$\bigcup_{n\ge 2}\mathfrak{S}_n$, of the functionals $\vec f\mapsto S^T_{n;K_j,m_j}(\vec f)$ as
$(K_j,m_j)\to\infty$ along the standing data. Then $\Xi$ is determined by its values on the
countable family $\mathfrak{F}$, and there is a subsequence of the index sequence of the standing
data along which $S^T_{n;K_j,m_j}(\vec f)\to\Xi(\vec f)$ for every $n\ge 2$ and every
$\vec f\in\mathfrak{S}_n$.
\end{remark}

\begin{proof}
First, $\Xi$ is determined by its values on the countable family $\mathfrak{F}$. The right-hand
side of \eqref{7.4:eq-equicontinuity-a} does not depend on $(K,m)$, and the estimate holds for all
$(K,m)$ beyond its floors, so it passes to $\Xi$. With the approximants $\vec\varphi\in\mathfrak{F}$
constructed above, it gives $|\Xi(\vec f)-\Xi(\vec\varphi)|\le\Delta_{\vec f}(\eta)$ for every
$\eta>0$.

Second, $\Xi$ is a sequential limit of the kind constructed above. Tuple by tuple, the values are
bounded by the per-ball volume-free bound; in this sense precompactness is of the trivial kind.
Enumerate the countable family $\mathfrak{F}$ as $\vec\varphi_{[1]},\vec\varphi_{[2]},\dots$, and let
$n_l$ be the length of $\vec\varphi_{[l]}$. Since $\Xi$ is a pointwise cluster point, for every $k$
there are arbitrarily large indices $j$ of the standing data with
\begin{equation*}
\bigl|S^T_{n_l;K_j,m_j}(\vec\varphi_{[l]})-\Xi(\vec\varphi_{[l]})\bigr|\le 1/k
\qquad\text{for all }l\le k.
\end{equation*}
Choosing such indices $j_1<j_2<\cdots$, one for each $k$, extracts a subsequence of the index
sequence of the standing data along which the values converge to $\Xi$ at every element of
$\mathfrak{F}$; this is the nested-and-diagonal mechanism of the diagonal fact recalled in
Remark~\ref{7.4:rem-standard-facts}, by which $\mathcal{J}$ was obtained in
\eqref{7.4:eq-infinite-volume-d}. Along that
subsequence, the argument above yields convergence on every $\mathfrak{S}_n$, to values which agree
with $\Xi$ because both are determined by the same values on $\mathfrak{F}$. Hence every limit point
of the connected Schwinger functions on separated tuples is covered by the construction
\eqref{7.4:eq-equicontinuity-b}.
\end{proof}

\begin{remark}[Iterated limits are diagonal limits]\label{7.4:rem-iterated-limits}
Nothing in this section lets $K\to\infty$ at fixed $m$. Suppose one does, in the following form.
Let $n\ge 2$, and let $(K_i)_{i\in\mathbb{N}}$ be cutoffs $K_i=K(g_0^{(i)},g)$ with bare couplings
$g_0^{(i)}$. Suppose, as in Proposition~\ref{7.1:prop-torus-limits}, that for
each $m\ge m_{\mathbb{T}}(g_{-}(g))$ there is an infinite set $I'_m\subset\mathbb{N}$ with
$K_i\to\infty$ as $i\to\infty$ in $I'_m$, such that $S^T_{n;K_i,m}(\vec f)$ converges, as
$i\to\infty$ in $I'_m$, to a limit $S^{(m)}_n(\vec f)$. This is required for every
$\vec f\in\mathfrak{S}_n$ whose transport to $\mathbb{T}_m$ is defined, that is, whose supports lie
in a ball $B_{\rho}$ with $L^m\ge 2\rho+2$. Suppose moreover that the torus
limits $S^{(m)}_n$ obey the torus bound of the correlation theorem with its constant at fixed $m$,
and that each is invariant under the whole group $(\mathbb{R}/2L^m\mathbb{Z})^4\rtimes W_4$;
neither property is used below.

Let $m_j\to\infty$, with $m_j\ge m_{\mathbb{T}}(g_{-}(g))$, be a sequence along which $S^{(m_j)}_n$
converges on the $n$-tuples of $\mathfrak{F}$. Then there
are indices $i_j\in I'_{m_j}$, increasing in $j$ and with $K_{i_j}\to\infty$, that have the
following property. The diagonal sequence
$(K_{i_j},m_j)$ has limits on the $n$-tuples of $\mathfrak{F}$ and on all of $\mathfrak{S}_n$, and
these limits are
the iterated limits $\lim_j S^{(m_j)}_n$. Thus iterated limits are diagonal limits. The converse is
not claimed, and neither is uniqueness.
\end{remark}

\begin{proof}
The family $\mathfrak{F}$ of Definition~\ref{7.4:def-countable-family} is countable; enumerate its
$n$-tuples as $\vec g_{[1]},\vec g_{[2]},\dots$. The supports of $\vec g_{[l]}$ lie in a ball
$B_{\rho_l}$, so that $S^{(m)}_n(\vec g_{[l]})$ is defined once $L^m\ge 2\rho_l+2$.

We choose the indices inductively. Put $i_0:=0$. Given $i_{j-1}$, consider the finitely many
$l\le j$ with $L^{m_j}\ge 2\rho_l+2$. At each such $l$ the values $S^T_{n;K_i,m_j}(\vec g_{[l]})$
converge to $S^{(m_j)}_n(\vec g_{[l]})$ as $i\to\infty$ in the infinite set $I'_{m_j}$, and
$K_i\to\infty$ as $i\to\infty$ in $I'_{m_j}$. Hence there
is $i_j\in I'_{m_j}$ with $i_j>i_{j-1}$ and $K_{i_j}\ge j$ such that
\begin{equation}\label{7.4:eq-iterated-limits-1}
\bigl|S^T_{n;K_{i_j},m_j}(\vec g_{[l]})-S^{(m_j)}_n(\vec g_{[l]})\bigr|\le 2^{-j}
\end{equation}
for all these $l$. Since $m_j\to\infty$, every fixed $l$ is among them for all large $j$; and
$K_{i_j}\ge j$ gives $K_{i_j}\to\infty$.
Write $\sigma_j(\vec f):=S^T_{n;K_{i_j},m_j}(\vec f)$.

\emph{Limits on $\mathfrak{F}$.} Fix $l$. For all large $j$, \eqref{7.4:eq-iterated-limits-1} gives
$|\sigma_j(\vec g_{[l]})-S^{(m_j)}_n(\vec g_{[l]})|\le 2^{-j}$. By hypothesis
$S^{(m_j)}_n(\vec g_{[l]})$ converges as $j\to\infty$. Hence $\sigma_j(\vec g_{[l]})$ converges to
the same limit.

\emph{Limits on $\mathfrak{S}_n$.} Let $\vec f\in\mathfrak{S}_n$ and $\eta>0$. Choose
$\vec g\in\mathfrak{F}$ approximating $\vec f$ entry by entry to within $\eta$ in the
$\|\cdot\|_{C^1}$ norm, exactly as in the proof of
Lemma~\ref{7.4:prf-equicontinuity}, and let $\Delta_{\vec f}(\eta)$ be the bound
defined there; it tends
to $0$ as $\eta\downarrow 0$ with $\vec f$ fixed. Say $\vec g=\vec g_{[l]}$.

Let $j$ be so large that $m_j$ satisfies the floors of the equicontinuity estimate for the radii
fixed in that proof. Then the equicontinuity estimate \eqref{7.4:eq-equicontinuity-a} applies at
every $(K_i,m_j)$ with $i\in I'_{m_j}$, because each $K_i$ has the form $K(g_0^{(i)},g)$. It gives
\begin{equation*}
\bigl|S^T_{n;K_i,m_j}(\vec f)-S^T_{n;K_i,m_j}(\vec g)\bigr|\le\Delta_{\vec f}(\eta),
\qquad i\in I'_{m_j}.
\end{equation*}
For $i=i_j$ this reads
$|\sigma_j(\vec f)-\sigma_j(\vec g)|\le\Delta_{\vec f}(\eta)$. The constant $C^{\infty}_n$ depends on
neither $K$ nor $m$, so letting $i\to\infty$ in $I'_{m_j}$ yields
\begin{equation*}
\bigl|S^{(m_j)}_n(\vec f)-S^{(m_j)}_n(\vec g)\bigr|\le\Delta_{\vec f}(\eta).
\end{equation*}

For all sufficiently large such $j$ (so that \eqref{7.4:eq-iterated-limits-1} holds at $l$, as noted
after \eqref{7.4:eq-iterated-limits-1}), combining these two bounds with
\eqref{7.4:eq-iterated-limits-1} gives
\begin{equation*}
\bigl|\sigma_j(\vec f)-S^{(m_j)}_n(\vec f)\bigr|\le 2\Delta_{\vec f}(\eta)+2^{-j}.
\end{equation*}
For two such indices $j,j'$ we likewise get
\begin{equation*}
|\sigma_j(\vec f)-\sigma_{j'}(\vec f)|\le 2\Delta_{\vec f}(\eta)+|\sigma_j(\vec g)-\sigma_{j'}(\vec g)|.
\end{equation*}
The last term tends to $0$ by the first step. Hence
\begin{align*}
\limsup_{j}\sigma_j(\vec f)-\liminf_{j}\sigma_j(\vec f)&\le 2\Delta_{\vec f}(\eta),\\
\limsup_{j}\bigl|\sigma_j(\vec f)-S^{(m_j)}_n(\vec f)\bigr|&\le 2\Delta_{\vec f}(\eta).
\end{align*}

Since $\Delta_{\vec f}(\eta)\to 0$ as $\eta\downarrow 0$ with $\vec f$ fixed, $\sigma_j(\vec f)$
converges. Moreover $S^{(m_j)}_n(\vec f)$ converges to the same limit. So the diagonal limit at
$\vec f$ exists and equals the iterated limit.
\end{proof}

In what follows, for a function $\Phi$ on an open subset of $\mathbb{R}^D$ and an integer $\nu\ge 0$, $\|\Phi\|_{C^{\nu}}$ is the maximum of the sup norms of all partial derivatives of $\Phi$ of order at most $\nu$. The one exception is $\nu=1$ and $D=4$, where $\|f\|_{C^1}=\max\{\|f\|_\infty,\|\nabla f\|_\infty\}$ as in Remark~\ref{7.4:rem-standard-facts}. Here $\|\nabla f\|_\infty:=\sup_x|\nabla f(x)|$, and $|\nabla f(x)|$ is the Euclidean length of the gradient. This norm dominates the maximum of the sup norms of $f,\partial_1f,\dots,\partial_4f$.

\begin{lemma}[The limits as distributions off the diagonals]\label{7.5:prf-limit-distributions}
Assume the standing data of the infinite-volume theorem, Theorem~\ref{7.4:thm-infinite-volume}.
Let $n\ge 2$, and let $S^T_n$ be the limit functional on $\mathfrak{S}_n$ defined by
\eqref{7.4:eq-equicontinuity-b}. Put
\begin{equation*}
(\mathbb{R}^4)^n_{\neq}:=\{(x_1,\dots,x_n)\in(\mathbb{R}^4)^n:\ x_i\neq x_k\ \text{for all } i\neq k\}.
\end{equation*}
\begin{itemize}
\item[(a)] Let $R>0$ and $\mathsf{d}>0$, and let $G_1,\dots,G_n\subset\mathbb{R}^4$ be open sets, each
of diameter at most $2R$, pairwise at distance at least $\mathsf{d}$. There is a unique
distribution $T_G\in\mathcal{D}'(G_1\times\dots\times G_n)$ such that
$T_G(\varphi_1\otimes\dots\otimes\varphi_n)=S^T_n(\vec\varphi)$ for all real-valued
$\varphi_i\in\mathcal{D}(G_i)$, a space we denote $\mathcal{D}(G_i;\mathbb{R})$.
\item[(b)] There is a unique distribution on $(\mathbb{R}^4)^n_{\neq}$, again written $S^T_n$, whose
restriction to $G_1\times\dots\times G_n$ is $T_G$ for every family $(G_i)$ as in \textup{(a)}
(with any $R>0$ and $\mathsf{d}>0$).
\item[(c)] For $(G_i)$ as in \textup{(a)}, finite sums of tensors $\varphi_1\otimes\dots\otimes\varphi_n$
with $\varphi_i\in\mathcal{D}(G_i)$ are dense in $\mathcal{D}(\prod_iG_i)$. Moreover, if
$\chi_i\in\mathcal{D}(G_i;\mathbb{R})$ are given, there is a constant
$C=C(n,\mathsf{d},R,G,\chi)$ such that
\begin{equation}\label{7.5:eq-limit-distributions-a}
|T_G(\Phi)|\le C\,\|\Phi\|_{C^{5n+1}}
\end{equation}
for every $\Phi\in\mathcal{D}(\prod_iG_i)$ such that each $\chi_i$ is identically $1$ on a
neighbourhood of the projection of $\operatorname{supp}\Phi$ to the $i$-th factor. Thus $T_G$ has
order at most $5n+1$ on $\prod_iG_i$, and the distribution $S^T_n$ of \textup{(b)} has locally finite
order on $(\mathbb{R}^4)^n_{\neq}$.
\end{itemize}
\end{lemma}

\begin{proof}
\emph{(a) Local pieces.} Let $\varphi_i\in\mathcal{D}(G_i;\mathbb{R})$. The supports of the $\varphi_i$ are pairwise at distance at least $\mathsf{d}$, so $\vec\varphi\in\mathfrak{S}_n$. Each $G_i$ has diameter at most $2R$, so it lies in the ball of radius $2R$ about any of its points; any enclosing radius would serve equally. The bound \eqref{7.4:eq-infinite-volume-a} was established for the limits in the proof of Lemma~\ref{7.4:prf-equicontinuity} (equicontinuity and existence of the limits). Applied with radius $2R$ and separation $\mathsf{d}$, it gives
\begin{equation*}
|S^T_n(\vec\varphi)|\le C^{\infty}_n(\mathsf{d};2R)\prod_{i}\|\varphi_i\|_{\sharp}
\le C^{\infty}_n(\mathsf{d};2R)\,(40M)^n\prod_i\|\varphi_i\|_{C^1}.
\end{equation*}
The second inequality holds because $\operatorname{Lip}\varphi_i\le\|\nabla\varphi_i\|_\infty$. This follows from the mean value theorem along segments of $\mathbb{R}^4$, with $\varphi_i$ extended by zero and its gradient measured in Euclidean length, as in the definition above. Since $40M\ge1$ and $\|\cdot\|_\sharp$ is the larger of the sup norm and $40M$ times the Lipschitz constant, we get $\|\varphi_i\|_\sharp\le 40M\|\varphi_i\|_{C^1}$.

The limits are multilinear on separated tuples, as established in the same proof. Hence the map $\mathsf{B}(\vec\varphi):=S^T_n(\vec\varphi)$ is $n$-linear on
$\mathcal{D}(G_1;\mathbb{R})\times\dots\times\mathcal{D}(G_n;\mathbb{R})$. We extend $\mathsf{B}$
complex-multilinearly to $\mathcal{D}(G_1)\times\dots\times\mathcal{D}(G_n)$. Writing each complex
$\varphi_i$ as $\operatorname{Re}\varphi_i+\mathrm{i}\operatorname{Im}\varphi_i$ expresses
$\mathsf{B}(\vec\varphi)$ as a sum of $2^n$ real terms. Since
$\|\operatorname{Re}\varphi_i\|_{C^1}\le\|\varphi_i\|_{C^1}$ and
$\|\operatorname{Im}\varphi_i\|_{C^1}\le\|\varphi_i\|_{C^1}$, the bound above applies to each term, and
\begin{equation}\label{7.5:eq-limit-distributions-1}
|\mathsf{B}(\vec\varphi)|\le 2^n(40M)^nC^{\infty}_n(\mathsf{d};2R)\prod_i\|\varphi_i\|_{C^1}
\qquad(\varphi_i\in\mathcal{D}(G_i)).
\end{equation}

For compact sets $\Gamma_i\subset G_i$, let $\mathcal{D}_{\Gamma_i}(G_i)$ be the space of test functions supported in $\Gamma_i$, with its Fr\'echet topology. The norm $\|\cdot\|_{C^1}$ is a continuous seminorm on it. By \eqref{7.5:eq-limit-distributions-1}, $\mathsf{B}$ is jointly continuous on $\prod_i\mathcal{D}_{\Gamma_i}(G_i)$. A linear form on $\mathcal{D}(G_i)$ is continuous if and only if its restriction to every $\mathcal{D}_{\Gamma_i}(G_i)$ is continuous. Hence $\mathsf{B}$ is separately continuous for the $\mathcal{D}$-topologies. The kernel theorem recorded in Remark~\ref{7.4:rem-standard-facts} therefore yields a unique $T_G\in\mathcal{D}'(\prod_iG_i)$ with $T_G(\varphi_1\otimes\dots\otimes\varphi_n)=\mathsf{B}(\vec\varphi)$ for all complex $\varphi_i\in\mathcal{D}(G_i)$. By multilinearity, any distribution that takes the value $S^T_n(\vec\varphi)$ on all real tensors takes the value $\mathsf{B}(\vec\varphi)$ on all complex tensors. So the uniqueness in (a) holds as stated.

\emph{(b) Gluing.} Consider the family of all products $\prod_iG_i$ with $(G_i)$ as in (a), for some $R>0$ and $\mathsf{d}>0$. This family covers $(\mathbb{R}^4)^n_{\neq}$. Indeed, let $x\in(\mathbb{R}^4)^n_{\neq}$ and put $\delta:=\min_{i\neq k}|x_i-x_k|>0$. Take the open balls $G_i$ of radius $\eta<\delta/3$ about the $x_i$. They have diameter at most $2\eta$ and are pairwise at distance at least $\delta-2\eta>\delta/3$, and their product contains $x$.

Let $\prod_iG_i$ and $\prod_iG'_i$ be two members of the family. Their overlap is $\prod_iG_i\cap\prod_iG'_i=\prod_i(G_i\cap G'_i)$. If some $G_i\cap G'_i$ is empty, there is nothing to check. Otherwise the sets $G_i\cap G'_i$ are open, each of diameter at most that of $G_i$, and pairwise at distance at least that of the $G_i$. So $(G_i\cap G'_i)$ is again a family as in (a). For $\varphi_i\in\mathcal{D}(G_i\cap G'_i;\mathbb{R})$, the tensor $\varphi_1\otimes\dots\otimes\varphi_n$ lies in both $\mathcal{D}(\prod_iG_i)$ and $\mathcal{D}(\prod_iG'_i)$. The restrictions of $T_G$ and $T_{G'}$ to $\prod_i(G_i\cap G'_i)$ therefore both take the value $S^T_n(\vec\varphi)$ on it. By the uniqueness in (a), applied to $(G_i\cap G'_i)$, the two restrictions agree.

The gluing theorem for distributions recorded in Remark~\ref{7.4:rem-standard-facts} now gives a unique $S^T_n\in\mathcal{D}'((\mathbb{R}^4)^n_{\neq})$ whose restriction to every member $\prod_iG_i$ of the cover is $T_G$.

The same uniqueness transfers invariance statements from tensors to the distribution. Consider a symmetry of the kind appearing in clause (B) of Theorem~\ref{7.4:thm-infinite-volume}, and suppose that the values of $S^T_n$ on tensors are invariant under it. On each product $\prod_iG_i$, the distribution pulled back by it then takes on tensors the same values as $S^T_n$. Hence it coincides with $T_G$ on each $\prod_iG_i$, and therefore with $S^T_n$.

\emph{(c) Order at most $5n+1$, and density of tensors.} Let $\Phi\in\mathcal{D}(\prod_iG_i)$. For each $i$, the projection of $\operatorname{supp}\Phi$ to the $i$-th factor is a compact subset of $G_i$. Choose $\chi_i\in\mathcal{D}(G_i;\mathbb{R})$ identically $1$ on a neighbourhood of it. If the $\chi_i$ are prescribed in advance with this property, as in the statement, the argument below applies to them unchanged. Choose an open box $Y\subset\mathbb{R}^4$ containing all $\operatorname{supp}\chi_i$. Then $\Phi\in\mathcal{D}(Y^n)$, and $Y^n$ is an open box in $\mathbb{R}^{4n}$.

Let $\Lambda^*_Y\subset\mathbb{R}^4$ be the dual lattice of the period box $\bar Y$. Then $(\Lambda^*_Y)^n$ is the dual lattice of $\bar Y^n$. For $\xi=(\xi_1,\dots,\xi_n)\in(\Lambda^*_Y)^n$ put $e_\xi(x):=\prod_ie^{\mathrm{i}\xi_i\cdot x_i}$. By the Fourier-series facts recorded in Remark~\ref{7.4:rem-standard-facts}, in dimension $4n$,
$\Phi=\sum_{\xi\in(\Lambda^*_Y)^n}\hat\Phi(\xi)e_\xi$ with convergence in $C^\infty(\bar Y^n)$. For every integer $\nu\ge0$,
\begin{equation}\label{7.5:eq-limit-distributions-2}
|\hat\Phi(\xi)|\le C_{\nu,Y^n}(1+|\xi|)^{-\nu}\|\Phi\|_{C^{\nu}} .
\end{equation}
If $\Phi(x)\neq0$, then each $x_i$ lies in the projection of $\operatorname{supp}\Phi$, where $\chi_i=1$. Hence $\Phi=\Phi\cdot\prod_i\chi_i(x_i)$. For $\Xi>0$ put
\begin{equation}\label{7.5:eq-limit-distributions-3}
\Phi_\Xi:=\sum_{|\xi|\le\Xi}\hat\Phi(\xi)\,
(\chi_1e^{\mathrm{i}\xi_1\cdot})\otimes\dots\otimes(\chi_ne^{\mathrm{i}\xi_n\cdot}).
\end{equation}
This is a finite sum of tensors of functions in $\mathcal{D}(G_i)$, that is, an element of the span of separated tensors.

We show that $\Phi_\Xi\to\Phi$ in $\mathcal{D}(\prod_iG_i)$. All supports lie in the fixed compact set $\prod_i\operatorname{supp}\chi_i$. Let $\nu_1\ge0$ be an integer. By the Leibniz rule, for a constant $C_{\nu_1}$ depending only on $\nu_1$,
\begin{equation*}
\|\chi_ie^{\mathrm{i}\xi_i\cdot}\|_{C^{\nu_1}}\le C_{\nu_1}(1+|\xi_i|)^{\nu_1}\|\chi_i\|_{C^{\nu_1}}.
\end{equation*}
A partial derivative of order at most $\nu_1$ of a product of functions of the separate variables $x_1,\dots,x_n$ is a product of partial derivatives of the factors, each of order at most $\nu_1$. Using also $1+|\xi_i|\le1+|\xi|$, the $C^{\nu_1}$ norm of the $\xi$-th tensor in \eqref{7.5:eq-limit-distributions-3} is therefore at most $C^n_{\nu_1}(1+|\xi|)^{n\nu_1}\prod_i\|\chi_i\|_{C^{\nu_1}}$. On $Y^n$ we have $\Phi-\Phi_\Xi=\prod_i\chi_i\cdot\sum_{|\xi|>\Xi}\hat\Phi(\xi)e_\xi$, and both sides vanish outside $\prod_i\operatorname{supp}\chi_i$. Taking $\nu=n\nu_1+4n+1$ in \eqref{7.5:eq-limit-distributions-2}, we obtain
\begin{equation*}
\|\Phi-\Phi_\Xi\|_{C^{\nu_1}}\le C^n_{\nu_1}\prod_i\|\chi_i\|_{C^{\nu_1}}\,
C_{\nu,Y^n}\|\Phi\|_{C^{\nu}}\sum_{|\xi|>\Xi}(1+|\xi|)^{-4n-1}.
\end{equation*}
The right-hand side tends to $0$ as $\Xi\to\infty$, because $\sum_{\xi\in(\Lambda^*_Y)^n}(1+|\xi|)^{-4n-1}<\infty$: the exponent $4n+1$ exceeds $4n$ (Remark~\ref{7.4:rem-standard-facts}). This proves the convergence, and with it the density of finite sums of tensors asserted in (c). That density is the fact behind the uniqueness in the kernel theorem used in (a), and it is used again below in this chapter.

By the continuity of $T_G$ and the multilinear extension in (a),
\begin{equation}\label{7.5:eq-limit-distributions-4}
T_G(\Phi)=\lim_{\Xi\to\infty}T_G(\Phi_\Xi)
=\sum_{\xi\in(\Lambda^*_Y)^n}\hat\Phi(\xi)\,
\mathsf{B}\bigl((\chi_ie^{\mathrm{i}\xi_i\cdot})_{i\le n}\bigr).
\end{equation}
For a real $\chi\in\mathcal{D}(\mathbb{R}^4)$ we have $\nabla(\chi e^{\mathrm{i}\xi\cdot})=(\nabla\chi+\mathrm{i}\xi\chi)e^{\mathrm{i}\xi\cdot}$. Hence $\|\chi e^{\mathrm{i}\xi\cdot}\|_{C^1}\le(1+|\xi|)\|\chi\|_{C^1}$. The supports of the $\chi_i$ lie in the $G_i$, so \eqref{7.5:eq-limit-distributions-1} applies and gives
\begin{equation*}
\bigl|\mathsf{B}\bigl((\chi_ie^{\mathrm{i}\xi_i\cdot})_{i}\bigr)\bigr|
\le 2^nC^{\infty}_n(\mathsf{d};2R)\prod_i40M(1+|\xi_i|)\|\chi_i\|_{C^1}.
\end{equation*}
Inserting this and \eqref{7.5:eq-limit-distributions-2} into \eqref{7.5:eq-limit-distributions-4} yields
\begin{equation*}
\begin{split}
|T_G(\Phi)|&\le 2^n(40M)^nC^{\infty}_n(\mathsf{d};2R)\prod_i\|\chi_i\|_{C^1}\\
&\qquad\cdot C_{\nu,Y^n}\|\Phi\|_{C^{\nu}}\sum_{\xi}(1+|\xi|)^{-\nu}\prod_i(1+|\xi_i|).
\end{split}
\end{equation*}
Since $\prod_i(1+|\xi_i|)\le(1+|\xi|)^n$, the last sum is at most $\sum_\xi(1+|\xi|)^{-\nu+n}$. By the Fourier-series facts in dimension $D=4n$, this sum is finite if and only if $\nu-n>4n$. So $\nu=5n+1$ suffices, and \eqref{7.5:eq-limit-distributions-a} follows with
\begin{equation*}
\begin{split}
C(n,\mathsf{d},R,G,\chi)&:=2^n(40M)^nC^{\infty}_n(\mathsf{d};2R)\prod_i\|\chi_i\|_{C^1}\\
&\qquad\cdot C_{5n+1,Y^n}\sum_{\xi}(1+|\xi|)^{-4n-1}.
\end{split}
\end{equation*}
Here $Y$ is an open box containing the $\operatorname{supp}\chi_i$, chosen once for the given $\chi_i$, so the dependence of $C$ on $Y$ is part of its dependence on $\chi$.

Let $\Gamma\subset\prod_iG_i$ be compact. Choosing each $\chi_i$ identically $1$ on a neighbourhood of the projection of $\Gamma$ to the $i$-th factor gives one constant $C$ serving every $\Phi$ supported in $\Gamma$. Hence $T_G$ has order at most $5n+1$ on $\prod_iG_i$. Since the products $\prod_iG_i$ cover $(\mathbb{R}^4)^n_{\neq}$ by (b), $S^T_n$ has locally finite order there.

Finally, the series \eqref{7.5:eq-limit-distributions-4}, which converges absolutely by the estimate just proved, also exhibits $T_G$ explicitly. This completes the proof of clause (A) of the infinite-volume theorem.
\end{proof}

\begin{lemma}[Exact invariance under translations and signed permutations]\label{7.5:prf-exact-invariance}
Assume the standing data of the infinite-volume theorem, Theorem~\ref{7.4:thm-infinite-volume}.
Let $\mathcal{J}$ be the diagonal subsequence of \eqref{7.4:eq-infinite-volume-d}. For $n\ge 2$ let
$S^T_n$ be the limits along $\mathcal{J}$ given by \eqref{7.4:eq-equicontinuity-b}. Let
$\mathfrak{S}_n$ be the class of separated real $C^1$ tuples of Lemma~\ref{7.4:prf-equicontinuity},
on which these limits are defined; we recall that its elements are the $n$-tuples of real $C^1$
functions of compact support on $\mathbb{R}^4$ with pairwise disjoint supports. For
$a\in\mathbb{R}^4$ put $\tau_a f:=f(\cdot-a)$. For
a signed permutation of the coordinates $\eta\in W_4$ put $(\eta\cdot f)(x):=f(\eta^{-1}x)$. Both actions
are applied componentwise to tuples.
\begin{itemize}
\item[(a)] For every $n\ge 2$, every $\vec f=(f_1,\dots,f_n)\in\mathfrak{S}_n$, every
$a\in\mathbb{R}^4$, every $\eta\in W_4$ and every permutation $\sigma$ of $\{1,\dots,n\}$,
\begin{equation}\label{7.5:eq-exact-invariance-1}
S^T_n(\tau_a f_1,\dots,\tau_a f_n)=S^T_n(f_1,\dots,f_n)=S^T_n(\eta\cdot f_1,\dots,\eta\cdot f_n)
\end{equation}
and
\begin{equation*}
S^T_n(f_{\sigma(1)},\dots,f_{\sigma(n)})=S^T_n(f_1,\dots,f_n).
\end{equation*}
\item[(b)] Let $q\ge1$ and let $\mathcal{T}^{\mathrm{sep}}_q$ be the linear span of the tensors
$f_1\otimes\dots\otimes f_q$ with $(f_1,\dots,f_q)\in\mathfrak{S}_q$. Let
$\mathbb{R}^4\rtimes W_4$ act on such tensors factorwise through $f\mapsto\tau_a(\eta\cdot f)$. For
$F\in\mathcal{T}^{\mathrm{sep}}_q$ let $S_q(F)$ be the limit Schwinger function of $q$ arguments
of the infinite-volume theorem, Theorem~\ref{7.4:thm-infinite-volume}, that is, the limit along
$\mathcal{J}$ of the centred moments of Definition~\ref{7.4:def-centred-moments}, taken on the
transported factors; on
$f_1\otimes\dots\otimes f_q$ it equals the sum over partitions $\pi\in\operatorname{Part}_{\ge2}[q]$
of the products, over the blocks $B$ of $\pi$, of the cumulants $S^T_{|B|}$ of the subtuples
$(f_i)_{i\in B}$, which exist by \eqref{7.4:eq-equicontinuity-b}, and it is linear in $F$. Then
\begin{equation}\label{7.5:eq-exact-invariance-a}
S_q\bigl((a,\eta)\cdot F\bigr)=S_q(F)\qquad
\text{for all } F\in\mathcal{T}^{\mathrm{sep}}_q,\ (a,\eta)\in\mathbb{R}^4\rtimes W_4 .
\end{equation}
In particular $S_q$ is invariant under the reflection $\theta$ of the time coordinate and
under all time and space translations.
\end{itemize}
\end{lemma}

\begin{proof}
\emph{Signed permutations.} Let $\eta\in W_4$ and $\vec f\in\mathfrak{S}_n$, and choose $\rho>0$ with
$\bigcup_i\operatorname{supp}f_i\subset B_\rho$, the open ball of radius $\rho$ about $0$. A signed
permutation is an orthogonal map, so $\eta\cdot B_\rho=B_\rho$.
Let $j\in\mathcal{J}$ be so large that $L^{m_j}\ge 2\rho+2$.

By the convention \eqref{7.4:eq-torus-transport-a}, $S^T_{n;j}(\eta\cdot\vec f)$ is the torus
cumulant of the transported functions $(\eta\cdot f_i)^{(m_j)}$. By part~(iii) of the transport
lemma, Lemma~\ref{7.4:lem-torus-transport}, $(\eta\cdot f_i)^{(m_j)}=\eta\cdot(f_i^{(m_j)})$. By the
lemma on lattice symmetries of the curvature cumulants, Lemma~\ref{7.4:lem-lattice-symmetries},
applied with $h_i:=f_i^{(m_j)}$, the torus cumulant is unchanged by $\eta$. Hence
\begin{align*}
S^T_{n;j}(\eta\cdot\vec f)
&=S^T_{n;K_j,m_j}\bigl(((\eta\cdot f_i)^{(m_j)})_i\bigr)
=S^T_{n;K_j,m_j}\bigl((\eta\cdot(f_i^{(m_j)}))_i\bigr)\\
&=S^T_{n;K_j,m_j}\bigl((f_i^{(m_j)})_i\bigr)=S^T_{n;j}(\vec f),
\end{align*}
and this equality is exact at every such $j$.

Moreover $\eta\cdot\vec f\in\mathfrak{S}_n$. Indeed, each $\eta\cdot f_i$ is real, $C^1$ and
compactly supported, and the supports $\eta(\operatorname{supp}f_i)$ are pairwise disjoint. By
\eqref{7.4:eq-equicontinuity-b} both sides therefore converge along $\mathcal{J}$, and so
$S^T_n(\eta\cdot\vec f)=S^T_n(\vec f)$. The symmetry of $S^T_n$ under permutations of its $n$
arguments is established in Lemma~\ref{7.4:prf-equicontinuity}, together with the existence of
the limits.

\emph{Translations: the mechanism.} Let $a\in\mathbb{R}^4$. Neither of the following two devices
suffices alone, and both are used.
\begin{itemize}
\item[(i)] The first is the exact lattice identity at a rounded vector. Let $a_K$ be a point of
$L^{-K}\mathbb{Z}^4=\varepsilon\mathbb{Z}^4$ nearest to $a$. Then
$|a-a_K|_\infty\le\varepsilon/2$, hence $|a-a_K|\le\varepsilon$ in the Euclidean norm on
$\mathbb{R}^4$. For Lipschitz $\vec f$ such that $\vec f$ and $\tau_{a_K}\vec f$ are supported in
a ball $B_\rho$ with $L^m\ge 2\rho+2$, we have
\begin{equation}\label{7.5:eq-exact-invariance-2}
S^T_{n;K,m}(\tau_{a_K}\vec f)=S^T_{n;K,m}(\vec f).
\end{equation}
This follows from Lemma~\ref{7.4:lem-lattice-symmetries} with $b=a_K\in\varepsilon\mathbb{Z}^4$,
together with part~(iii) of Lemma~\ref{7.4:lem-torus-transport}, which carries the translation
$\tau_{a_K}$ through the transport to the torus.
\item[(ii)] The second is the $C^1$ translation modulus of
Lemma~\ref{7.4:lem-translation-modulus}, applied to the rounding error
$\tau_a f-\tau_{a_K}f$. By \eqref{7.4:eq-translation-modulus-a} its sharp norm is at most
$\varpi_f(\varepsilon)$, which tends to $0$ as $K\to\infty$. This is converted into a bound on
cumulants by the equicontinuity estimate \eqref{7.4:eq-equicontinuity-a}, whose constant does
not depend on $m$. This is the second place where the $m$-independent constant $C^\infty_n$ of
the per-ball volume-free bound is used, the first being the proof of existence of the limits
\eqref{7.4:eq-equicontinuity-b}.
\end{itemize}
With merely Lipschitz data step~(ii) is unavailable, since by
Lemma~\ref{7.4:lem-translation-modulus} the modulus need not tend to zero, and for the same
reason continuity of $a\mapsto S^T_n(\tau_a\vec f)$ is not available either.

With $C^1$ data one may alternatively argue as follows. First obtain exact invariance under the
dense subgroup
\begin{equation*}
\mathbb{Z}[1/L]^4=\bigcup_{j}L^{-K_j}\mathbb{Z}^4 .
\end{equation*}
Each $a\in\mathbb{Z}[1/L]^4$ lies in $L^{-K_{j_0}}\mathbb{Z}^4$ for some $j_0$, and then step~(i)
holds with $a_K=a$ at every $j$ with $K_j\ge K_{j_0}$. Then
use the continuity of $a\mapsto S^T_n(\tau_a\vec f)$ to reach all of $\mathbb{R}^4$; this
continuity follows from \eqref{7.4:eq-equicontinuity-a} in the limit together with
Lemma~\ref{7.4:lem-translation-modulus}. The proof below follows the first route.

\emph{Proof of $S^T_n(\tau_a\vec f)=S^T_n(\vec f)$ for every $a\in\mathbb{R}^4$: choice of the
data.} Fix $\vec f\in\mathfrak{S}_n$ and put $\mathsf{d}:=\mathsf{d}(\vec f)>0$. Choose
$\rho$ with $\bigcup_i\operatorname{supp}f_i\subset B_\rho$, and fix $a\in\mathbb{R}^4$. Then
$\tau_a\vec f\in\mathfrak{S}_n$, with supports in $B_{\rho+|a|}$. For $j\in\mathcal{J}$ put
$K:=K_j$, $m:=m_j$ and $\varepsilon:=L^{-K}$. Assume $j$ so large that the following hold:
\begin{itemize}
\item $L^m\ge 2(\rho+|a|+1)+2$;
\item $\varepsilon\le\min\{1,\mathsf{d}/4\}$;
\item the conditions on $m$ in \eqref{7.4:eq-equicontinuity-a} hold for $n$, $R':=\rho+|a|+1$,
$\mathsf{d}':=\mathsf{d}/2$ and the transport radius $\rho':=\rho+|a|+1$.
\end{itemize}
All of these hold for large $j$, since $m_j\to\infty$ and $K_j\to\infty$ by the standing data,
and $K_j=K(g_0^{(j)},g)$ is the first passage. Let $a_K$ be the rounded vector of step~(i).

\emph{The compact sets $Q_i$.} Put
\begin{equation*}
Q_i:=\operatorname{supp}f_i+a+\bar B_\varepsilon ,
\end{equation*}
the closed $\varepsilon$-neighbourhood of $\operatorname{supp}\tau_a f_i$. Each $Q_i$ is compact,
being the sum of two compact sets. It contains $\operatorname{supp}\tau_a f_i$, and it contains
$\operatorname{supp}\tau_{a_K}f_i=\operatorname{supp}f_i+a_K$ because $|a_K-a|\le\varepsilon$.
Its points have the form $x+a+y$ with $|x|<\rho$ and $|y|\le\varepsilon\le1$. Hence
$Q_i\subset B_{\rho+|a|+1}$, a ball of radius $R'$ about $0$.

For $i\ne l$, the distance between $Q_i$ and $Q_l$ is at least
$\mathsf{d}-2\varepsilon\ge\mathsf{d}/2$, because $\varepsilon\le\mathsf{d}/4$. The functions
$\tau_a f_i$ and
$\tau_{a_K}f_i$ are real and Lipschitz and are supported in $Q_i$.

\emph{The equicontinuity bound.} We apply \eqref{7.4:eq-equicontinuity-a} with $f_i$ replaced by
$\tau_a f_i$ and $g_i$ replaced by $\tau_{a_K}f_i$. By Lemma~\ref{7.4:lem-translation-modulus}
and its display \eqref{7.4:eq-translation-modulus-a},
\begin{align*}
\|\tau_a f_l\|_\sharp&=\|\tau_{a_K}f_l\|_\sharp=\|f_l\|_\sharp ,\\
\|\tau_a f_i-\tau_{a_K}f_i\|_\sharp&\le\varpi_{f_i}(|a-a_K|)\le\varpi_{f_i}(\varepsilon).
\end{align*}
The
last step uses that $\varpi_{f_i}$ is non-decreasing, since $\omega_{\nabla f_i}$ is. We obtain
\begin{equation}\label{7.5:eq-exact-invariance-3}
\bigl|S^T_{n;j}(\tau_a\vec f)-S^T_{n;j}(\tau_{a_K}\vec f)\bigr|
\le C^\infty_n\bigl(\mathsf{d}/2;\rho+|a|+1\bigr)\sum_{i\le n}\varpi_{f_i}(\varepsilon)
\prod_{l\ne i}\|f_l\|_\sharp .
\end{equation}

\emph{Passage to the limit.} The right side of \eqref{7.5:eq-exact-invariance-3} tends to $0$ as
$j\to\infty$ along $\mathcal{J}$. Indeed $\varepsilon=L^{-K_j}\to0$ and
$\varpi_{f_i}(\varepsilon)\to0$ by Lemma~\ref{7.4:lem-translation-modulus}, while the constant
does not depend on $j$: it is uniform in $m$, and uniform in $a$ on bounded sets.

By step~(i), in the form \eqref{7.5:eq-exact-invariance-2} with $m=m_j$ and transport radius
$\rho+|a|+1$, we have $S^T_{n;j}(\tau_{a_K}\vec f)=S^T_{n;j}(\vec f)$. Hence
\begin{equation*}
S^T_{n;j}(\tau_a\vec f)-S^T_{n;j}(\vec f)\to0\qquad\text{along }\mathcal{J}.
\end{equation*}
Both sequences converge along $\mathcal{J}$ by \eqref{7.4:eq-equicontinuity-b}, because
$\tau_a\vec f\in\mathfrak{S}_n$. This step relies on convergence on all separated real $C^1$
tuples, and not only on a countable family of them. So the limits agree, which together with
the first paragraph proves \eqref{7.5:eq-exact-invariance-1}.

\emph{Proof of part~(b).} Let $(f_1,\dots,f_q)\in\mathfrak{S}_q$ and
$(a,\eta)\in\mathbb{R}^4\rtimes W_4$. The tuple $(\tau_a(\eta\cdot f_i))_i$ again lies in
$\mathfrak{S}_q$, by the two membership observations above. Every subtuple of a tuple in
$\mathfrak{S}_q$ with at least two entries lies in the corresponding $\mathfrak{S}_{n}$.

By Definition~\ref{7.4:def-centred-moments}, at each $j$ the centred moment of the transported
functions $f_1^{(m_j)},\dots,f_q^{(m_j)}$ is the sum over partitions
$\pi\in\operatorname{Part}_{\ge2}[q]$ of the products, over the blocks $B$ of $\pi$, of the torus
cumulants of the corresponding subtuples. Letting $j\to\infty$ along $\mathcal{J}$ gives the same
expression for $S_q$, with each torus cumulant replaced by its limit $S^T_{|B|}$. Each factor
$S^T_{|B|}$ of each product is invariant under
$f\mapsto\tau_a(\eta\cdot f)$ by \eqref{7.5:eq-exact-invariance-1}, applied first with $\eta$ and then
with $a$. Hence each product is invariant, and so is $S_q(f_1\otimes\dots\otimes f_q)$.

Since $S_q$ is linear on $\mathcal{T}^{\mathrm{sep}}_q$ and the action is linear,
\eqref{7.5:eq-exact-invariance-a} follows. The reflection $\theta$ changes the sign of one
coordinate, so $\theta\in W_4$. Time and space translations are translations $\tau_a$.
\end{proof}

\begin{lemma}[Reflection positivity of the centred limit]
\label{7.6:prf-centred-positivity}
Let the standing data of the infinite-volume theorem (Theorem~\ref{7.4:thm-infinite-volume}) be in
force, let $\mathcal{J}$ be the diagonal subsequence of \eqref{7.4:eq-infinite-volume-d}, and write
$\langle\cdot\rangle_j:=\langle\cdot\rangle_{K_j,m_j}$ and $S_{q;j}:=S_{q;K_j,m_j}$. Put $S_0:=1$,
$S_1:=0$,
$\mathcal{T}^{\mathrm{sep}}_0:=\mathbb{C}$ and $\Phi_{K,m}(c):=c$ for $c\in\mathbb{C}$. Recall that
$\mathcal{P}_+=\bigoplus_{q\ge 0}\mathcal{T}^{\mathrm{sep}}_{q,+}$, where
$\mathcal{T}^{\mathrm{sep}}_{0,+}=\mathbb{C}$ and, for $q\ge 1$, $\mathcal{T}^{\mathrm{sep}}_{q,+}$ is
the complex span of the tensors $f_1\otimes\cdots\otimes f_q$ of real $f_a\in C^1_c(\mathbb{R}^4)$
with compact, pairwise disjoint supports contained in $\{x^0>0\}$. Then the following hold.
\begin{enumerate}
\item[(a)] Let $q\ge 2$ and $F\in\mathcal{T}^{\mathrm{sep}}_q$, all supports in some $B_{\rho}$. For
$j\in\mathcal{J}$ with $L^{m_j}\ge 2\rho+2$ the number
$S_{q;j}(F):=\langle\Phi_{K_j,m_j}(F)\rangle_j$ is well defined, the limit
\begin{equation*}
S_q(F):=\lim_{j\in\mathcal{J}}S_{q;j}(F)
\end{equation*}
exists, it is linear in the function $F$, and on a separated tensor $h_1\otimes\cdots\otimes h_q$ it
equals
\begin{equation}\label{7.6:eq-centred-positivity-1}
S_q(h_1\otimes\cdots\otimes h_q)
=\sum_{\pi\in\operatorname{Part}_{\ge 2}[q]}\ \prod_{B\in\pi}S^T_{|B|}\bigl((h_l)_{l\in B}\bigr).
\end{equation}
\item[(b)] The form $(\cdot,\cdot)_{\mathrm{OS}}$ of \eqref{7.4:eq-infinite-volume-b} is
sesquilinear, Hermitian and positive semidefinite on $\mathcal{P}_+$, and for
$F=(F_q)_q$, $G=(G_q)_q\in\mathcal{P}_+$
\begin{equation}\label{7.6:eq-centred-positivity-2}
(F,G)_{\mathrm{OS}}=\lim_{j\in\mathcal{J}}
\bigl\langle\overline{\mathbb{F}_j(\theta\cdot U)}\,\mathbb{G}_j(U)\bigr\rangle_j,
\qquad
\mathbb{F}_j:=\sum_q\Phi_{K_j,m_j}(F_q),\quad \mathbb{G}_j:=\sum_q\Phi_{K_j,m_j}(G_q).
\end{equation}
\item[(c)] Let $(\mathbb{R}^4)^q_{\neq,+}:=\{x\in(\mathbb{R}^4)^q_{\neq}:x_a^0>0\ \text{for all } a\}$.
For every $q_0$ and every finite sequence $F_q\in\mathcal{D}((\mathbb{R}^4)^q_{\neq,+})$, $q\le q_0$,
\begin{equation}\label{7.6:eq-centred-positivity-3}
\sum_{q,q'\le q_0}S_{q+q'}\bigl((\Theta\bar F)_q\otimes F_{q'}\bigr)\ \ge\ 0,
\end{equation}
where $S_{q+q'}$ is the distribution on $(\mathbb{R}^4)^{q+q'}_{\neq}$.
\end{enumerate}
\end{lemma}

\begin{proof}
Throughout, $\theta$ is the time reflection $x^0\mapsto -x^0$; for a function $F$ of $q$ points we
write $(F\circ\theta)(x_1,\dots,x_q):=F(\theta x_1,\dots,\theta x_q)$, $\bar F$ for the complex
conjugate function, and $\theta\vec f:=(f_1\circ\theta,\dots,f_q\circ\theta)$; thus
$(\Theta\bar F)_q=\overline{F_q\circ\theta}$, as in \eqref{7.4:eq-infinite-volume-b}.

\emph{Part (a).} Write $F=\sum_k c_k\,h_{k,1}\otimes\cdots\otimes h_{k,q}$, a finite combination of
separated real tensors with supports in $B_{\rho}$. For $L^{m_j}\ge 2\rho+2$, linearity in the test
function (Lemma~\ref{7.4:lem-linearity}) shows that $\Phi_{K_j,m_j}(F)$, hence
$S_{q;j}(F)=\langle\Phi_{K_j,m_j}(F)\rangle_j$, is a well-defined linear functional of the function
$F$, independent of the chosen tensor representation. By this linearity,
$S_{q;j}(F)=\sum_k c_k\,S_{q;j}(h_{k,1}\otimes\cdots\otimes h_{k,q})$, and on each separated tensor
$S_{q;j}$ is the sum over $\pi\in\operatorname{Part}_{\ge 2}[q]$ of the products over $B\in\pi$ of
$S^T_{|B|;j}((h_{k,l})_{l\in B})$, by \eqref{7.4:eq-centred-moments-a}. Each sub-tuple
$(h_{k,l})_{l\in B}$ is again separated, so each factor converges along $\mathcal{J}$ to
$S^T_{|B|}((h_{k,l})_{l\in B})$ by \eqref{7.4:eq-equicontinuity-b}. There are finitely many $k$,
partitions and blocks; hence $\lim_{j\in\mathcal{J}}S_{q;j}(F)$ exists and equals
$\sum_k c_k$ times the right side of \eqref{7.6:eq-centred-positivity-1} at
$(h_{k,1},\dots,h_{k,q})$. As a pointwise limit of linear functionals of the function $F$, $S_q$ is
linear in $F$ and depends on $F$ only; taking a single tensor gives
\eqref{7.6:eq-centred-positivity-1}. For $q=1$, $\langle\Phi_{K_j,m_j}(F)\rangle_j=0$, each
${:}\mathcal{O}(h^{(m_j)}){:}$ being centred, so the limit is $0=S_1$, which is also the sum over
the empty set $\operatorname{Part}_{\ge 2}[1]$; for $q=0$,
$\langle\Phi_{K_j,m_j}(c)\rangle_j=c=c\,S_0$.

\emph{Part (b).} Let $F=(F_q)_{q\le q_0}\in\mathcal{P}_+$. Only finitely many test functions occur,
all real, of class $C^1_c$, with compact supports in $B_{\rho}\cap\{x^0>0\}$ for some $\rho$. Take
$j\in\mathcal{J}$ so large that $L^{m_j}\ge 2\rho+2$ and that $m_j$ is at least each of the
finitely many floors required in the equicontinuity estimate in the proof of
Theorem~\ref{7.4:thm-infinite-volume}, for the
block sizes, radii and separations occurring, and put $\mathbb{F}_j:=\sum_q\Phi_{K_j,m_j}(F_q)$.

\emph{Claim 1: $\mathbb{F}_j$ is a bounded measurable function of the bond variables in the open slab
$H_+=\{0<x^0<L^{m_j}\}$.} By Lemma~\ref{7.4:lem-half-space}, each ${:}\mathcal{O}(f^{(m_j)}){:}$ with
$f$ as above is such a function. By Lemma~\ref{7.4:lem-linearity}, $\Phi_{K_j,m_j}(F_q)$ is a finite
linear combination of products of such centred observables, and products and sums of bounded
measurable functions of the bonds in $H_+$ are again such functions.

\emph{Claim 2: the complex conjugate of $\mathbb{F}_j$ at the reflected configuration is the lattice
polynomial of $\Theta\bar F$, that is,}
\begin{equation*}
\overline{\mathbb{F}_j(\theta\cdot U)}=\sum_q\Phi_{K_j,m_j}((\Theta\bar F)_q)(U).
\end{equation*}
On a tensor, Lemma~\ref{7.4:lem-linearity} and the reflection identity of
Lemma~\ref{7.4:lem-half-space} (for the observable and for its mean) give
\begin{equation*}
\Phi_{K_j,m_j}(f_1\otimes\cdots\otimes f_q)(\theta\cdot U)
=\prod_a{:}\mathcal{O}(f_a^{(m_j)}){:}(\theta\cdot U)
=\prod_a{:}\mathcal{O}((f_a\circ\theta)^{(m_j)}){:}(U),
\end{equation*}
and the right side is $\Phi_{K_j,m_j}((f_1\otimes\cdots\otimes f_q)\circ\theta)(U)$. By linearity,
$\Phi_{K_j,m_j}(F_q)(\theta\cdot U)=\Phi_{K_j,m_j}(F_q\circ\theta)(U)$. Since every $a_p(U)$ is real,
complex conjugation of $\Phi_{K_j,m_j}(G)$ acts on the values of $G$ only:
$\overline{\Phi_{K_j,m_j}(G)}=\Phi_{K_j,m_j}(\bar G)$. With $G=F_q\circ\theta$ this gives
$\overline{\Phi_{K_j,m_j}(F_q)(\theta\cdot U)}=\Phi_{K_j,m_j}((\Theta\bar F)_q)(U)$; sum over $q$.

\emph{Claim 3.} By Claim 2 and the multiplicativity in \eqref{7.4:eq-linearity-a},
\begin{equation}\label{7.6:eq-centred-positivity-4}
\begin{split}
\overline{\mathbb{F}_j(\theta\cdot U)}\,\mathbb{F}_j(U)
&=\sum_{q,q'}\Phi_{K_j,m_j}((\Theta\bar F)_q)\,\Phi_{K_j,m_j}(F_{q'})\\
&=\sum_{q,q'}\Phi_{K_j,m_j}\bigl((\Theta\bar F)_q\otimes F_{q'}\bigr),
\end{split}
\end{equation}
and $(\Theta\bar F)_q\otimes F_{q'}\in\mathcal{T}^{\mathrm{sep}}_{q+q'}$: in each of its tensor terms
the $q$ reflected functions have pairwise disjoint supports in $\{x^0<0\}$ and the $q'$ others
pairwise disjoint supports in $\{x^0>0\}$, so all $q+q'$ supports are pairwise disjoint, all in
$B_{\rho}$, whatever the overlaps between different tensor terms of $F$.

Taking expectations in \eqref{7.6:eq-centred-positivity-4} and using part (a), with finitely many
terms,
\begin{align*}
\bigl\langle\overline{\mathbb{F}_j(\theta\cdot U)}\,\mathbb{F}_j(U)\bigr\rangle_j
&=\sum_{q,q'}S_{q+q';j}\bigl((\Theta\bar F)_q\otimes F_{q'}\bigr)\\
&\longrightarrow\sum_{q,q'}S_{q+q'}\bigl((\Theta\bar F)_q\otimes F_{q'}\bigr)=(F,F)_{\mathrm{OS}}.
\end{align*}
On the other hand, the left side is non-negative for every such $j$. Indeed, $0$ and
$L^{m_j}=L^{m_j+K_j}L^{-K_j}$ lie in $\varepsilon\mathbb{Z}$, $\varepsilon=L^{-K_j}$, so
$\Pi:=\{x^0=0\}\cup\{x^0=L^{m_j}\}$ consists of two planes at times in $\varepsilon\mathbb{Z}$, which
contain no lattice sites and bisect the time-like links, whatever $K$
\cite[(0.1)]{B-CMP109}; the reflection $\theta_{\Pi}$ in it is $x^0\mapsto -x^0$, i.e.\ $\theta$, and
$H_+$ is one of the two open components of $\mathbb{T}_{m_j}\setminus\Pi$. By Claim 1, part (ii) of
reflection positivity of the lattice measure, applied to $\mathbb{F}_j$ at $\Pi$, gives
$\langle\overline{\mathbb{F}_j(\theta\cdot U)}\,\mathbb{F}_j(U)\rangle_j\ge 0$. A limit of
non-negative numbers is non-negative:
\begin{equation*}
(F,F)_{\mathrm{OS}}\ge 0 .
\end{equation*}
For $F,G\in\mathcal{P}_+$, Claims 1--3 with $\mathbb{G}_j$ in the second factor of
\eqref{7.6:eq-centred-positivity-4} give \eqref{7.6:eq-centred-positivity-2}.

\emph{Hermitian symmetry.} By \eqref{7.6:eq-centred-positivity-2}, then the invariance of
$\mu_{K_j,m_j}$ under the isometry $\theta$ of $\mathbb{T}_{m_j}$, which preserves the bare lattice
(part (i) of reflection positivity of the lattice measure), together with $\theta^2=1$,
\begin{align*}
\overline{(F,G)_{\mathrm{OS}}}
&=\lim_{j}\overline{\bigl\langle\overline{\mathbb{F}_j(\theta\cdot U)}\,
\mathbb{G}_j(U)\bigr\rangle_j}
=\lim_{j}\bigl\langle\mathbb{F}_j(\theta\cdot U)\,\overline{\mathbb{G}_j(U)}\bigr\rangle_j\\
&=\lim_{j}\bigl\langle\mathbb{F}_j(U)\,\overline{\mathbb{G}_j(\theta\cdot U)}\bigr\rangle_j
=(G,F)_{\mathrm{OS}} .
\end{align*}
Equivalently, on real separated tensors $\otimes f=f_1\otimes\cdots\otimes f_q$ and
$\otimes g=g_1\otimes\cdots\otimes g_{q'}$,
\begin{equation*}
(\otimes f,\otimes g)_{\mathrm{OS}}=S_{q+q'}(\theta\vec f,\vec g)
=S_{q+q'}(\vec f,\theta\vec g)=S_{q+q'}(\theta\vec g,\vec f)=(\otimes g,\otimes f)_{\mathrm{OS}},
\end{equation*}
by the invariance of $S_{q+q'}$ under $\theta\in W_4$ applied to all arguments,
\eqref{7.5:eq-exact-invariance-a}, with $\theta^2=1$, and then the symmetry of $S_{q+q'}$ in its
arguments (the sum over partitions and each joint cumulant are symmetric); the common value is
real, being a limit of expectations of real random variables. The general case follows by
sesquilinear extension to complex coefficients.

\emph{Sesquilinearity.} $G\mapsto(F,G)_{\mathrm{OS}}$ is linear since $\otimes$ and each
$S_{q+q'}$ are linear, and $F\mapsto(F,G)_{\mathrm{OS}}$ is conjugate linear since
$F_q\mapsto(\Theta\bar F)_q$ is.

\emph{Part (c).} Let $F_q\in\mathcal{D}((\mathbb{R}^4)^q_{\neq,+})$, $q\le q_0$, and
$\mathsf{K}_q:=\operatorname{supp}F_q$. For $x\in\mathsf{K}_q$ choose $r_x>0$ with
$r_x<\tfrac13\min_{a\neq b}|x_a-x_b|$ and $r_x<\min_a x_a^0$, and put
$W(x):=\prod_a B(x_a,r_x)$. It is an open neighbourhood of $x$ in $(\mathbb{R}^4)^q_{\neq,+}$ and a
product of pairwise separated open balls inside $\{x^0>0\}$. Finitely many of these,
$W_l:=W(x_l)$ with $x_l=(x_{l,1},\dots,x_{l,q})\in\mathsf{K}_q$ and $r_l:=r_{x_l}$,
$l=1,\dots,\nu$, cover the compact set $\mathsf{K}_q$; a smooth partition of unity
$\chi_1,\dots,\chi_{\nu}$
subordinate to them on a neighbourhood of $\mathsf{K}_q$ splits $F_q=\sum_l\chi_lF_q$ with
$\chi_lF_q\in\mathcal{D}(W_l)$. For each piece, choose, as in \eqref{7.5:eq-limit-distributions-a},
real one-variable cutoffs $\eta_{l,a}\in\mathcal{D}(B(x_{l,a},r_l))$, $a\le q$, equal to $1$ on a
neighbourhood of the $a$-th projection of $\operatorname{supp}(\chi_lF_q)$, and let
$\widehat{\chi_lF_q}(\xi)$ be the Fourier coefficients of $\chi_lF_q$ used there, indexed by the
frequency tuples $\xi=(\xi_1,\dots,\xi_q)$ of that display. The partial sums
\begin{equation*}
F^{(\Xi)}_{q,l}:=\sum_{|\xi|\le\Xi}\widehat{\chi_lF_q}(\xi)\,
\eta_{l,1}e^{i\xi_1\cdot}\otimes\cdots\otimes\eta_{l,q}e^{i\xi_q\cdot},\qquad \Xi>0,
\end{equation*}
lie in $\mathcal{T}^{\mathrm{sep}}_{q,+}$: their tensor factors are supported in the fixed compact
sets $\operatorname{supp}\eta_{l,a}$, contained in the pairwise separated balls $B(x_{l,a},r_l)$
inside $\{x^0>0\}$, and each complex factor is a combination of the real factors
$\eta_{l,a}\cos(\xi_a\cdot)$, $\eta_{l,a}\sin(\xi_a\cdot)$, with the same supports. By that display,
$F^{(\Xi)}_{q,l}\to\chi_lF_q$ in $\mathcal{D}(W_l)$ as $\Xi\to\infty$. Put
$F^{(\Xi)}:=(\sum_lF^{(\Xi)}_{q,l})_q\in\mathcal{P}_+$. By part (b),
$(F^{(\Xi)},F^{(\Xi)})_{\mathrm{OS}}\ge 0$.

For open sets $W$, $W'$ of the above kind, the map $(\varphi,\psi)\mapsto\Theta\bar\varphi\otimes\psi$
is continuous from $\mathcal{D}(W)\times\mathcal{D}(W')$ to $\mathcal{D}(\Theta W\times W')$, being a
product of smooth functions in separate variables with supports in a fixed compact set; and
$\Theta W\times W'\subset(\mathbb{R}^4)^{q+q'}_{\neq}$, since reflected points have negative times
and the others positive times. The distribution $S_{q+q'}$ obtained in the proof of
Theorem~\ref{7.4:thm-infinite-volume} by identifying the limits as distributions off the diagonals
(see \eqref{7.5:eq-limit-distributions-a}) is
continuous on $\mathcal{D}(\Theta W\times W')$ and agrees with the $S_{q+q'}$ of part (a) on
separated tensors, both being equal to the sum over partitions of products of $S^T$ there. Applying
this to each pair of pieces $(l,l')$ and summing,
\begin{align*}
\sum_{q,q'}S_{q+q'}\bigl((\Theta\bar F)_q\otimes F_{q'}\bigr)
&=\lim_{\Xi\to\infty}\sum_{q,q'}S_{q+q'}\bigl((\Theta\overline{F^{(\Xi)}})_q\otimes
F^{(\Xi)}_{q'}\bigr)\\
&=\lim_{\Xi\to\infty}(F^{(\Xi)},F^{(\Xi)})_{\mathrm{OS}}\ \ge\ 0,
\end{align*}
which is \eqref{7.6:eq-centred-positivity-3}.
\end{proof}

The argument uses Lemmas~\ref{7.4:lem-half-space} and~\ref{7.4:lem-linearity}, the moment--cumulant
formula (Remark~\ref{7.4:rem-standard-facts}), the equicontinuity estimate and the existence of the
limits (in the proof of Theorem~\ref{7.4:thm-infinite-volume}, \eqref{7.4:eq-equicontinuity-b}),
the limits as distributions off the diagonals
(\eqref{7.5:eq-limit-distributions-a}) and their exact invariance under translations and signed
permutations (\eqref{7.5:eq-exact-invariance-a}), and reflection positivity of
the lattice measure at the single plane pair $\{x^0=0\}\cup\{x^0=L^{m_j}\}$ only, no reflection
through sites and no other plane; the class $\mathcal{P}_+$ is preserved by time translations by
$t\ge 0$, by spatial translations and by $W_3$, since these preserve the property that supports are
compact, pairwise disjoint and contained in $\{x^0>0\}$.

\begin{lemma}[The reconstructed space and bounds uniform in time]\label{7.7:prf-reconstructed-space}
Assume the standing data of the infinite-volume theorem, Theorem~\ref{7.4:thm-infinite-volume}.
Let $\mathcal{P}_+$ and the Osterwalder--Schrader form $(\cdot,\cdot)_{\mathrm{OS}}$ be as in
\eqref{7.4:eq-infinite-volume-b}, let $\mathsf{T}_t$ be as in \eqref{7.4:eq-infinite-volume-c},
and write $\mathbb{1}:=(1,0,0,\dots)\in\mathcal{P}_+$.
\begin{enumerate}
\item[(a)] The set $\mathcal{N}:=\{F\in\mathcal{P}_+ : (F,F)_{\mathrm{OS}}=0\}$ is a linear subspace
of $\mathcal{P}_+$, and $(\cdot,\cdot)_{\mathrm{OS}}$ descends to an inner product on
$\mathcal{P}_+/\mathcal{N}$. Its completion $\mathcal{H}^{\mathrm{curv}}$ is a separable Hilbert
space. Writing $[F]$ for the class of $F$, the vector $\Omega:=[\mathbb{1}]$ satisfies
$\|\Omega\|=1$.
\item[(b)] Let $F\in\mathcal{P}_+$ be written with real tensors. Define three quantities:
\begin{itemize}
\item for each test function occurring in $F$ choose a closed ball containing its support, and let
$R_F$ be the largest of the radii of these balls; every occurring function is then supported in a
ball of radius $R_F$;
\item $\mathsf{d}_F$ is the least distance between two supports within one tensor term;
\item $\delta_F$ is the least distance of an occurring support from the plane $\{x^0=0\}$.
\end{itemize}
Here a minimum over the empty set is $+\infty$; if $F=F_0$ all claims below are trivial.
Then $\mathsf{d}_F>0$ and $\delta_F>0$. Moreover, for every $t\ge 0$, each connected function
arising when $(F,\mathsf{T}_tF)_{\mathrm{OS}}$ is expanded by \eqref{7.4:eq-infinite-volume-b}
into connected functions is evaluated at a separated real
$C^1$ tuple with supports in balls of radius $R_F$, at mutual distances at least $\mathsf{d}_0$,
where
\begin{equation}\label{7.7:eq-reconstructed-space-a}
\mathsf{d}_0:=\min\{\mathsf{d}_F,2\delta_F\}>0,\qquad
M_F:=\sup_{t\ge 0}\,(F,\mathsf{T}_tF)_{\mathrm{OS}}<\infty .
\end{equation}
The separation $\mathsf{d}_0$ does not depend on $t$.
\end{enumerate}
\end{lemma}

\begin{proof}
\emph{Notation.} Write $F=(F_q)_{q\ge 0}$, where only finitely many $F_q$ are non-zero, and
\begin{equation*}
F_q=\sum_k c_{q,k}\,f_{q,k,1}\otimes\cdots\otimes f_{q,k,q}
\end{equation*}
with finite sums, coefficients $c_{q,k}\in\mathbb{C}$, and real $C^1$ functions $f_{q,k,a}$ of
compact support. Within each tensor these functions have pairwise disjoint supports, all contained
in the open half-space $\{x^0>0\}$.

There are finitely many tensor terms, and within each the supports are pairwise disjoint compact
sets. Hence $\mathsf{d}_F>0$. The occurring supports are finitely many compact subsets of the open
half-space, so $\delta_F>0$. Every occurring support lies in $\{x^0\ge\delta_F\}$. Each occurring
$f$ is supported in a ball of radius at most $R_F$, hence in a ball of radius $R_F$.

\emph{(a) The space.} By part (C) of Theorem~\ref{7.4:thm-infinite-volume}, the form
$(\cdot,\cdot)_{\mathrm{OS}}$ is Hermitian and positive semidefinite on $\mathcal{P}_+$. The
Cauchy--Schwarz inequality for such forms gives
$|(F,G)_{\mathrm{OS}}|^2\le (F,F)_{\mathrm{OS}}(G,G)_{\mathrm{OS}}$. Hence $F\in\mathcal{N}$ if and
only if $(F,G)_{\mathrm{OS}}=0$ for all $G\in\mathcal{P}_+$. Therefore $\mathcal{N}$ is a linear
subspace, and $(\cdot,\cdot)_{\mathrm{OS}}$ descends to a well-defined inner product on
$\mathcal{P}_+/\mathcal{N}$. The space $\mathcal{H}^{\mathrm{curv}}$ is its completion, a Hilbert
space, and $\|\Omega\|^2=(\mathbb{1},\mathbb{1})_{\mathrm{OS}}=S_0=1$.

\emph{Separability.} Let $\mathcal{Q}$ and $D_Q$ ($Q\in\mathcal{Q}$) be as in
\eqref{7.4:eq-countable-family-a}. Let $\mathcal{P}_+^{\mathbb{Q}}\subset\mathcal{P}_+$ be the set
of finite $\mathbb{Q}(i)$-combinations of tensors $g_1\otimes\cdots\otimes g_q$ (for $q=0$, of
$\mathbb{1}$), where $Q_1,\dots,Q_q\in\mathcal{Q}$ are pairwise disjoint and contained in
$\{x^0>0\}$, and $g_a\in D_{Q_a}$. This set is countable: there are countably many $q$, countably
many tuples of sets in $\mathcal{Q}$, countably many elements of each $D_Q$, and countably many
coefficients in $\mathbb{Q}(i)$.

Fix $F$ as above and $0<\eta\le 1$. Choose $\delta>0$ with $\delta\le\mathsf{d}_F/4$ and
$\delta\le\delta_F/2$, and let $(\operatorname{supp}f)^{+\delta}$ denote the open
$\delta$-neighbourhood of $\operatorname{supp}f$. This choice has three consequences.
\begin{itemize}
\item Within each tensor, the $\delta$-neighbourhoods of two occurring supports are at distance at
least $\mathsf{d}_F-2\delta\ge\mathsf{d}_F/2$.
\item All these neighbourhoods lie in $\{x^0\ge\delta_F-\delta\}\subset\{x^0\ge\delta_F/2\}$.
\item Each neighbourhood lies in a ball of radius $R_F+\delta$.
\end{itemize}

For each occurring $f$, the support-adaptedness property (b2) in
\eqref{7.4:eq-countable-family-a} gives $Q(f)\in\mathcal{Q}$ with
\begin{equation*}
\operatorname{supp}f\subset Q(f)\subset(\operatorname{supp}f)^{+\delta}.
\end{equation*}
Then $f$ vanishes off $Q(f)$, so $f$ lies in the space of $C^1$ functions vanishing off $Q(f)$, in
which $D_{Q(f)}$ is dense for $\|\cdot\|_{C^1}$. Choose $g\in D_{Q(f)}$ with
$\|f-g\|_{C^1}\le\eta$, and coefficients $c'_{q,k}\in\mathbb{Q}(i)$ with
$|c_{q,k}-c'_{q,k}|\le\eta$. Replacing every $f_{q,k,a}$ by the chosen $g_{q,k,a}$ and every
$c_{q,k}$ by $c'_{q,k}$ defines $G$. Within each tensor the sets $Q(f_{q,k,a})$ are pairwise
disjoint and lie in $\{x^0>0\}$, so $G\in\mathcal{P}_+^{\mathbb{Q}}$.

In the tensor product we have the telescoping identity
\begin{equation*}
\bigotimes_{a\le q}f_a-\bigotimes_{a\le q}g_a
=\sum_{a\le q}g_1\otimes\cdots\otimes g_{a-1}\otimes(f_a-g_a)\otimes f_{a+1}\otimes\cdots\otimes f_q .
\end{equation*}
Applying it to each tensor term gives
\begin{equation}\label{7.7:eq-reconstructed-space-1}
F_q-G_q=\sum_k(c_{q,k}-c'_{q,k})\bigotimes_{a\le q}g_{q,k,a}
+\sum_k c_{q,k}\sum_{a\le q}g_{q,k,1}\otimes\cdots\otimes(f_{q,k,a}-g_{q,k,a})\otimes\cdots
\otimes f_{q,k,q}.
\end{equation}
So $F-G=\sum_j \xi_j\,\varphi_{j,1}\otimes\cdots\otimes\varphi_{j,q_j}$ is a finite sum of tensors.
The number
of terms depends only on the form of $F$. In each term, the factor in slot $a$ is supported in the
set $Q(f_{q,k,a})$. The tensors are therefore separated, with factors at mutual distance at least
$\mathsf{d}_F/2$, supported in $\{x^0\ge\delta_F/2\}$ and in balls of radius $R_F+\delta$.

Each term carries either a coefficient of modulus at most $\eta$, or one factor $f_a-g_a$. By
property (b1) in \eqref{7.4:eq-countable-family-a}, such a factor has
$\|f_a-g_a\|_{\sharp}\le 40M\|f_a-g_a\|_{C^1}\le 40M\eta$. Put
\begin{equation*}
\Gamma_F:=1+\max_{q,k}|c_{q,k}|,\qquad
\Lambda_F:=1+40M+\max\|f_{q,k,a}\|_{\sharp}.
\end{equation*}
Then every $g_{q,k,a}$ satisfies
$\|g_{q,k,a}\|_\sharp\le\|f_{q,k,a}\|_\sharp+40M\eta\le \Lambda_F$, and $|c'_{q,k}|\le \Gamma_F$.
Let
$\bar q$ be the largest $q$ with $F_q\ne 0$. Since $40M\le \Lambda_F$ and $\Lambda_F\ge 1$, every
term obeys
\begin{equation*}
|\xi_j|\prod_a\|\varphi_{j,a}\|_\sharp\le \Gamma_F\Lambda_F^{\bar q}\,\eta .
\end{equation*}

By sesquilinearity and the definition \eqref{7.4:eq-infinite-volume-b},
\begin{equation}\label{7.7:eq-reconstructed-space-2}
(F-G,F-G)_{\mathrm{OS}}=\sum_{j,j'}\overline{\xi_j}\,\xi_{j'}\,
S_{q_j+q_{j'}}\bigl((\varphi_{j,a}\circ\theta)_{a\le q_j},(\varphi_{j',b})_{b\le q_{j'}}\bigr).
\end{equation}
Each $S_{q_j+q_{j'}}$ is, by \eqref{7.4:eq-infinite-volume-b}, a sum over
$\pi\in\operatorname{Part}_{\ge2}[q_j+q_{j'}]$ of products of connected functions over the blocks
of $\pi$. A term with $q_j=q_{j'}=0$ equals $\overline{\xi_j}\xi_{j'}S_0$. For $q_j+q_{j'}=1$ the
set $\operatorname{Part}_{\ge2}[1]$ is empty.

The time reflection $\theta$ is an isometry. So the reflected factors are pairwise at distance at
least $\mathsf{d}_F/2$ and lie in $\{x^0\le-\delta_F/2\}$. A reflected factor and an unreflected
one are therefore at distance at least $\delta_F$. All factors are real $C^1$ functions supported in
balls of radius $R_F+\delta$, and $\|\varphi\circ\theta\|_\sharp=\|\varphi\|_\sharp$.

The bound \eqref{7.4:eq-infinite-volume-a}, applied to every block with separation
$\min\{\mathsf{d}_F/2,\delta_F\}$ and radius $R_F+\delta$, bounds each summand of
\eqref{7.7:eq-reconstructed-space-2} by
\begin{equation*}
\Bigl(\sum_{\pi}\prod_{B\in\pi}C^\infty_{|B|}\bigl(\min\{\mathsf{d}_F/2,\delta_F\};R_F+\delta\bigr)
\Bigr)\,(\Gamma_F\Lambda_F^{\bar q})^2\,\eta^2 .
\end{equation*}
The number of summands depends only on $F$. Hence
\begin{equation*}
\|[F]-[G]\|^2=(F-G,F-G)_{\mathrm{OS}}\le C_F\,\eta^2 ,
\end{equation*}
with $C_F$ independent of $\eta$, and this tends to $0$ as $\eta\to0$. The classes $[F]$,
$F\in\mathcal{P}_+$, are dense in $\mathcal{H}^{\mathrm{curv}}$ by construction. Therefore the
countable set $\{[G]:G\in\mathcal{P}_+^{\mathbb{Q}}\}$ is dense, and $\mathcal{H}^{\mathrm{curv}}$
is separable.

\emph{(b) Bounds uniform in time.} Let $t\ge0$. Unwinding the definitions
\eqref{7.4:eq-infinite-volume-b} and \eqref{7.4:eq-infinite-volume-c} (the functions are real, so
conjugation acts on the coefficients only) gives
\begin{equation}\label{7.7:eq-reconstructed-space-3}
(F,\mathsf{T}_tF)_{\mathrm{OS}}=\sum_{q,q'}\sum_{k,k'}\overline{c_{q,k}}\,c_{q',k'}\,
S_{q+q'}\bigl((f_{q,k,a}\circ\theta)_{a\le q},(\tau_{te_0}f_{q',k',b})_{b\le q'}\bigr).
\end{equation}
Each value $S_{q+q'}(\dots)$ here equals
\begin{equation*}
\sum_{\pi\in\operatorname{Part}_{\ge2}[q+q']}\ \prod_{B\in\pi}S^T_{|B|}(\vec\psi_B),
\end{equation*}
where $\vec\psi_B$, the sub-tuple of functions indexed by the block $B$, consists of $|B|\ge2$
functions from the set
\begin{equation*}
\{f_{q,k,a}\circ\theta\}\cup\{\tau_{te_0}f_{q',k',b}\}.
\end{equation*}
A term with $q=q'=0$ does not involve $t$, and for $q+q'=1$ the set $\operatorname{Part}_{\ge2}[1]$
is empty.

Each function in $\vec\psi_B$ is supported in a ball of radius $R_F$. Reflection and
translation carry balls to balls of the same radius, and this is the point at which the per-ball
form of the bound is used. The mutual distances are controlled in three cases.
\begin{itemize}
\item Two reflected functions are at distance at least $\mathsf{d}_F$, since $\theta$ is an
isometry.
\item Two translated functions are at distance at least $\mathsf{d}_F$, since translation is an
isometry.
\item A reflected function is supported in $\{x^0\le-\delta_F\}$ and a translated one in
$\{x^0\ge\delta_F+t\}$. Their supports are at distance at least $\delta_F+(\delta_F+t)\ge2\delta_F$,
because separation in the time coordinate bounds the distance from below.
\end{itemize}
So $\vec\psi_B$ is a separated real $C^1$ tuple whose supports lie in balls of radius $R_F$,
pairwise at distance at least $\mathsf{d}_0=\min\{\mathsf{d}_F,2\delta_F\}>0$. This separation does
not depend on $t$, since the cross distances only grow with $t$.

The bound \eqref{7.4:eq-infinite-volume-a} now applies directly, because the supports are pairwise
at distance at least $\mathsf{d}_0$; no monotonicity of the constant in $\mathsf{d}$ is needed. It
gives
\begin{equation*}
|S^T_{|B|}(\vec\psi_B)|\le C^\infty_{|B|}(\mathsf{d}_0;R_F)\prod_{\psi\in\vec\psi_B}
\|\psi\|_\sharp .
\end{equation*}
Here $\|f\circ\theta\|_\sharp=\|f\|_\sharp$, since $\theta$ is an isometry, and
$\|\tau_{te_0}f\|_\sharp=\|f\|_\sharp$ by Lemma~\ref{7.4:lem-translation-modulus}.

The number of terms in \eqref{7.7:eq-reconstructed-space-3} and in the partition expansions is
finite and does not depend on $t$. Therefore, for every $t\ge0$,
\begin{equation}\label{7.7:eq-reconstructed-space-4}
\begin{split}
|(F,\mathsf{T}_tF)_{\mathrm{OS}}|
&\le\sum_{q+q'\ge2}\sum_{k,k'}|c_{q,k}|\,|c_{q',k'}|
\sum_{\pi\in\operatorname{Part}_{\ge2}[q+q']}\ \prod_{B\in\pi}C^\infty_{|B|}(\mathsf{d}_0;R_F)\\
&\qquad\times\prod_{a\le q}\|f_{q,k,a}\|_\sharp\prod_{b\le q'}\|f_{q',k',b}\|_\sharp
\ +\ |F_0|^2S_0 ,
\end{split}
\end{equation}
where the last term is the summand with $q=q'=0$ (recall $S_0=1$), and the summands with $q+q'=1$
vanish. The right-hand side does not depend on $t$.
Hence the numbers $(F,\mathsf{T}_tF)_{\mathrm{OS}}$, $t\ge0$, are bounded in modulus uniformly in
$t$.

Each number $(F,\mathsf{T}_tF)_{\mathrm{OS}}$ is real. Indeed, by the infinite-volume theorem,
Theorem~\ref{7.4:thm-infinite-volume}, the connected functions at real tuples are limits of
connected lattice Schwinger functions, which are joint cumulants of real random variables; they
are therefore real, and so are the values $S_{q+q'}$ in \eqref{7.7:eq-reconstructed-space-3}, by
\eqref{7.4:eq-infinite-volume-b}. The affine map $x\mapsto\theta x+te_0$ is the composite of the
signed permutation $\theta\in W_4$ with the translation by $te_0$, and its induced action on test
functions carries $f\circ\theta$ to $\tau_{te_0}f$ and $\tau_{te_0}f$ to $f\circ\theta$. By the
invariance \eqref{7.5:eq-exact-invariance-a} and the symmetry of $S_{q+q'}$ under permutations of
its arguments (it is a sum over all partitions of products of connected functions, each symmetric
in its arguments), the value of $S_{q+q'}$ in the summand of
\eqref{7.7:eq-reconstructed-space-3} indexed by $(q,k),(q',k')$ equals the value in the summand
indexed by $(q',k'),(q,k)$. The coefficient matrix of the Hermitian expression
\eqref{7.7:eq-reconstructed-space-3} is therefore real and symmetric, so the expression is real.
Hence $M_F<\infty$ in \eqref{7.7:eq-reconstructed-space-a}.
\end{proof}

\begin{remark}
The bound \eqref{7.7:eq-reconstructed-space-a} is the bound uniform in time that enters the
doubling argument for the contraction property of the time translations in part (D) of
Theorem~\ref{7.4:thm-infinite-volume}. The per-ball form of the constant is what supplies it.

Consider a hypothesis of the form $\mathcal{C}(\mathsf{d},R)\le\mathcal{C}(\mathsf{d})(1+R)^{\sigma}$
for the $n$-point constant, with a fixed exponent $\sigma$ and with $R$ the
radius of one ball containing all the supports. The per-ball volume-free bound does not assert it:
the dependence of its constant on $R$, through the factor $\ell(n,R,g)^4$ in the exponent of the
constant $Q_\infty$, is super-polynomial. Under $\mathsf{T}_t$ the radius of a single ball
containing all supports of the tuples in \eqref{7.7:eq-reconstructed-space-3} is $\gtrsim t$.
Polynomial growth would be used at the $k$-th step of the doubling argument to obtain
\begin{equation*}
\bigl(\mathcal{C}(\mathsf{d})(1+R+2^{k+1}t)^{\sigma}\bigr)^{2^{-k-1}}\to1 .
\end{equation*}
The role of such a hypothesis is supplied by the per-ball form. There the relevant radius is the
individual radius $R_F$, which does not depend on $t$, so the doubling step runs with growth
exponent $0$.

A bound in terms of one ball containing all supports would not serve. Suppose all supports are
forced into one ball of radius $D(t)\gtrsim2t$, and write $\mathcal{M}(t)$ for the available bound
on $(F,\mathsf{T}_tF)_{\mathrm{OS}}$. For an $n$-point block this bound is of the type
$\exp\{n(E_+^\sharp+C_{\mathrm{low}})\,\ell(n,D(t),g)^4\}$. At the $k$-th doubling step it gives,
for some $c>0$ depending on $F$,
\begin{equation*}
\Bigl(\frac{\mathcal{M}(2^{k+1}t)}{(F,F)_{\mathrm{OS}}}\Bigr)^{2^{-k-1}}
\gtrsim\exp\bigl\{c\cdot16^{k+1}t^4\,2^{-k-1}\bigr\}\to\infty .
\end{equation*}
Such a bound neither proves nor refutes the contraction property.
\end{remark}

\begin{proof}[The transfer semigroup, the Hamiltonian and the momenta]
\label{7.7:prf-semigroup-hamiltonian}
We complete the proof of clause (D) of the infinite-volume theorem
(Theorem~\ref{7.4:thm-infinite-volume}).

\emph{Notation and standing facts.} The following objects are those fixed in the statement of
Theorem~\ref{7.4:thm-infinite-volume}, in \eqref{7.4:eq-infinite-volume-b} and
\eqref{7.4:eq-infinite-volume-c}:
\begin{itemize}
\item the space $\mathcal{P}_+$, the map $\Theta$ and the form $(\cdot,\cdot)_{\mathrm{OS}}$;
\item the null space $\mathcal{N}=\{F\in\mathcal{P}_+:(F,F)_{\mathrm{OS}}=0\}$ of the form, and
the space $\mathcal{H}^{\mathrm{curv}}$, the completion of $\mathcal{P}_+/\mathcal{N}$;
\item the vector $\Omega=[\mathbb{1}]$, where $[F]$ denotes the class of $F\in\mathcal{P}_+$;
\item the maps $\mathsf{T}_t$ and $\mathsf{S}_{\vec a}$.
\end{itemize}
On a tensor of $q+q'$ test functions, $\mathsf{T}_t$ translates every argument by $te_0$ and
$\mathsf{S}_{\vec a}$ translates every argument by $(0,\vec a)$. This is the action meant in
$\mathsf{T}_t[\,\cdot\otimes\cdot\,]$ and $\mathsf{S}_{\vec a}[\,\cdot\otimes\cdot\,]$ below.

By the reflection positivity of the centred limit, which is clause (C) of
Theorem~\ref{7.4:thm-infinite-volume}, the form
$(\cdot,\cdot)_{\mathrm{OS}}$ is Hermitian and positive semidefinite on $\mathcal{P}_+$. Hence the
Cauchy--Schwarz inequality applies to it, and its null vectors form a subspace
(Remark~\ref{7.4:rem-standard-facts}).

For $F\in\mathcal{P}_+$ written with real tensors, we use the radius $R_F$, the within-tensor
separation $\mathsf{d}_F$, the distance $\delta_F$ from the plane $\{x^0=0\}$, and the constant
$M_F$ of \eqref{7.7:eq-reconstructed-space-a}, all as fixed in the proof of the bounds uniform in
time that precedes this one. For $G\in\mathcal{P}_+$ the corresponding quantities,
defined in the same way, are denoted $R_G$, $\mathsf{d}_G$, $\delta_G$. We write
\begin{equation*}
F_q=\sum_k c_{q,k}\, f_{q,k,1}\otimes\cdots\otimes f_{q,k,q},
\end{equation*}
with complex $c_{q,k}$ and real $f_{q,k,a}$, and similarly $G_{q'}$ with coefficients
$c'_{q',k'}$ and functions $g_{q',k',b}$.

\emph{Time translations: well-defined positive self-adjoint contractions.}
For $s,t\ge 0$, $\mathsf{T}_t$ maps $\mathcal{P}_+$ into itself, and
$\mathsf{T}_s\mathsf{T}_t=\mathsf{T}_{s+t}$, $\mathsf{T}_t\mathbb{1}=\mathbb{1}$.

\emph{Symmetry.} Since $\theta(y+te_0)=\theta y-te_0$, we have
\begin{equation*}
(\Theta\overline{\mathsf{T}_tF})_q(x)
=\overline{F_q(\theta x_1-te_0,\dots,\theta x_q-te_0)}
=(\mathsf{T}_{-t}(\Theta\bar F))_q(x).
\end{equation*}
Each real tensor term of $\mathsf{T}_{-t}(\Theta\bar F)_q\otimes G_{q'}$ is separated: its first
$q$ supports lie in $\{x^0\le-\delta_F-t\}$ and are pairwise at distance at least $\mathsf{d}_F$,
and its last $q'$ supports lie in $\{x^0\ge\delta_G\}$ and are pairwise at distance at least
$\mathsf{d}_G$. So this tensor lies in $\mathcal{T}^{\mathrm{sep}}_{q+q'}$. Translating all
$q+q'$ of its arguments by $te_0$ leaves $S_{q+q'}$ unchanged, by the exact invariance
\eqref{7.5:eq-exact-invariance-a}. Therefore
\begin{align*}
(\mathsf{T}_tF,G)_{\mathrm{OS}}
&=\sum_{q,q'}S_{q+q'}\bigl(\mathsf{T}_{-t}(\Theta\bar F)_q\otimes G_{q'}\bigr)\\
&=\sum_{q,q'}S_{q+q'}\bigl(\mathsf{T}_t[\mathsf{T}_{-t}(\Theta\bar F)_q\otimes G_{q'}]\bigr)\\
&=\sum_{q,q'}S_{q+q'}\bigl((\Theta\bar F)_q\otimes\mathsf{T}_tG_{q'}\bigr)
=(F,\mathsf{T}_tG)_{\mathrm{OS}} .
\end{align*}

\emph{Null space.} By the symmetry, the semigroup law and the Cauchy--Schwarz inequality,
\begin{equation*}
(\mathsf{T}_tF,\mathsf{T}_tF)_{\mathrm{OS}}=(F,\mathsf{T}_{2t}F)_{\mathrm{OS}}
\le(F,F)_{\mathrm{OS}}^{1/2}(\mathsf{T}_{2t}F,\mathsf{T}_{2t}F)_{\mathrm{OS}}^{1/2}.
\end{equation*}
So $(F,F)_{\mathrm{OS}}=0$ implies $\mathsf{T}_tF\in\mathcal{N}$, and $\mathsf{T}_t$ descends to
$\mathcal{P}_+/\mathcal{N}$.

\emph{Contraction.} Fix $t\ge0$. For integers $k\ge0$ put
\begin{equation*}
\eta_k:=(\mathsf{T}_{2^kt}F,\mathsf{T}_{2^kt}F)_{\mathrm{OS}}^{1/2}.
\end{equation*}
By the symmetry, the semigroup law and the Cauchy--Schwarz inequality,
\begin{equation*}
\eta_k^2=(F,\mathsf{T}_{2^{k+1}t}F)_{\mathrm{OS}}\le(F,F)_{\mathrm{OS}}^{1/2}\,\eta_{k+1}.
\end{equation*}
By \eqref{7.7:eq-reconstructed-space-a} we also have $\eta_k^2\le M_F$.

If $(F,F)_{\mathrm{OS}}=0$, then $\eta_0=0$ by the null-space statement. Otherwise put
$\tilde\eta_k:=\eta_k/(F,F)_{\mathrm{OS}}^{1/2}$. Dividing the last inequality by
$(F,F)_{\mathrm{OS}}$ gives $\tilde\eta_k^2\le\tilde\eta_{k+1}$. Moreover
$\tilde\eta_k\le(M_F/(F,F)_{\mathrm{OS}})^{1/2}$. All $\tilde\eta_k$ are non-negative, so
induction on $k$ gives
\begin{equation*}
\tilde\eta_0\le\tilde\eta_1^{1/2}\le\cdots\le\tilde\eta_k^{2^{-k}}
\le\Bigl(\frac{M_F}{(F,F)_{\mathrm{OS}}}\Bigr)^{2^{-k-1}}\longrightarrow1
\qquad(k\to\infty).
\end{equation*}
Hence $\tilde\eta_0\le1$, that is,
\begin{equation}\label{7.7:eq-semigroup-hamiltonian-1}
(\mathsf{T}_tF,\mathsf{T}_tF)_{\mathrm{OS}}\le(F,F)_{\mathrm{OS}}\qquad(F\in\mathcal{P}_+,\ t\ge0).
\end{equation}

Thus $\mathsf{T}_t$ induces a contraction of $\mathcal{P}_+/\mathcal{N}$. We extend it by
continuity to a contraction $T_t$ of $\mathcal{H}^{\mathrm{curv}}$. The operator $T_t$ is symmetric
and bounded, hence self-adjoint. Further $T_sT_t=T_{s+t}$ and $T_0=1$, and $T_t\Omega=\Omega$
because $\mathsf{T}_t\mathbb{1}=\mathbb{1}$.

The operator $T_t$ is positive. Indeed, by the symmetry and the semigroup law,
$(v,T_tv)=(T_{t/2}v,T_{t/2}v)\ge0$ for $v$ in the dense quotient, hence for all
$v\in\mathcal{H}^{\mathrm{curv}}$ by continuity. No reflection positivity with respect to any plane
other than $\{x^0=0\}$ is used.

\emph{Strong continuity.} Let $F,G\in\mathcal{P}_+$ and $s,t\ge0$ with
$|t-s|\le\min\{\mathsf{d}_G/4,1\}$. Then
$(F,\mathsf{T}_tG)_{\mathrm{OS}}-(F,\mathsf{T}_sG)_{\mathrm{OS}}$ is the finite sum, over
$q,q',k,k'$, of the terms
\begin{equation*}
\overline{c_{q,k}}\,c'_{q',k'}\Bigl[S_{q+q'}\bigl(\theta\vec f,\tau_{te_0}\vec g\bigr)
-S_{q+q'}\bigl(\theta\vec f,\tau_{se_0}\vec g\bigr)\Bigr].
\end{equation*}
Here $\theta\vec f:=(f_{q,k,a}\circ\theta)_{a\le q}$ and
$\tau_{ue_0}\vec g:=(\tau_{ue_0}g_{q',k',b})_{b\le q'}$.

The functional $S_{q+q'}$ is multilinear in its arguments. We telescope in the $g$'s, writing
$g_b$ for $g_{q',k',b}$:
\begin{align*}
&S_{q+q'}(\theta\vec f,\tau_{te_0}\vec g)-S_{q+q'}(\theta\vec f,\tau_{se_0}\vec g)\\
&\quad=\sum_{b=1}^{q'}S_{q+q'}\bigl(\theta\vec f,\tau_{se_0}g_1,\dots,\tau_{se_0}g_{b-1},
\tau_{te_0}g_b-\tau_{se_0}g_b,\\
&\qquad\qquad\tau_{te_0}g_{b+1},\dots,\tau_{te_0}g_{q'}\bigr).
\end{align*}

Each tuple on the right is a separated tuple of real $C^1$ functions with the following
properties.
\begin{itemize}
\item The entry in position $q+l$, for $l\le q'$, whether translated or the difference entry, is
supported in $\operatorname{supp}g_l+[\min\{s,t\},\max\{s,t\}]e_0$. This set lies in a ball of
radius $R_G+1$ because $|t-s|\le1$.
\item Two such sets are at distance at least $\mathsf{d}_G-|t-s|\ge\mathsf{d}_G/2$.
\item Two reflected $f$'s are at distance at least $\mathsf{d}_F$.
\item A reflected $f$ is supported in $\{x^0\le-\delta_F\}$ and a translated $g$ in
$\{x^0\ge\delta_G\}$. Their distance is therefore at least $\delta_F+\delta_G$.
\end{itemize}
Hence all supports lie in balls of radius $R_{F,G}:=\max\{R_F,R_G+1\}$ and are pairwise at
distance at least
\begin{equation*}
\mathsf{d}_{F,G}:=\min\{\mathsf{d}_F,\mathsf{d}_G/2,\delta_F+\delta_G\}>0.
\end{equation*}
Both numbers are independent of $s,t$.

Each tuple has exactly one entry $\tau_{te_0}g_b-\tau_{se_0}g_b$. By the translation modulus in the
sharp norm (Lemma~\ref{7.4:lem-translation-modulus}, display
\eqref{7.4:eq-translation-modulus-a}), its sharp norm is at most $\varpi_{g_b}(|t-s|)$. The other
entries have sharp norms $\|f\circ\theta\|_{\sharp}=\|f\|_{\sharp}$ and
$\|\tau_{ue_0}g\|_{\sharp}=\|g\|_{\sharp}$.

We expand each $S_{q+q'}$ as $\sum_{\pi\in\operatorname{Part}_{\ge2}[q+q']}\prod_{B\in\pi}S^T_{|B|}$
by \eqref{7.4:eq-infinite-volume-b}. We bound every factor by the bound
\eqref{7.4:eq-infinite-volume-a} with the constant $C^{\infty}_{|B|}(\mathsf{d}_{F,G};R_{F,G})$.
This applies directly, since the supports are pairwise at distance at least $\mathsf{d}_{F,G}$ and
lie in balls of radius $R_{F,G}$. Exactly one factor contains the difference entry. This gives
\begin{equation*}
\bigl|(F,\mathsf{T}_tG)_{\mathrm{OS}}-(F,\mathsf{T}_sG)_{\mathrm{OS}}\bigr|
\le C_{F,G}\,\max_b\varpi_{g_b}(|t-s|)\longrightarrow0\qquad(|t-s|\to0).
\end{equation*}
Here the maximum runs over the finitely many functions occurring in $G$. The constant $C_{F,G}$
collects the finitely many coefficients, constants $C^{\infty}$ and sharp norms, and does not
depend on $s,t$. So $t\mapsto(F,\mathsf{T}_tG)_{\mathrm{OS}}$ is continuous on $[0,\infty[$.

By \eqref{7.7:eq-semigroup-hamiltonian-1},
\begin{align*}
\|T_t[F]-[F]\|^2
&=(\mathsf{T}_tF,\mathsf{T}_tF)_{\mathrm{OS}}-2\operatorname{Re}(F,\mathsf{T}_tF)_{\mathrm{OS}}
+(F,F)_{\mathrm{OS}}\\
&\le2(F,F)_{\mathrm{OS}}-2\operatorname{Re}(F,\mathsf{T}_tF)_{\mathrm{OS}}\longrightarrow0
\qquad(t\downarrow0),
\end{align*}
by the continuity just proved at $s=0$, where $\mathsf{T}_0$ is the identity.

Let $v\in\mathcal{H}^{\mathrm{curv}}$. For every $F\in\mathcal{P}_+$, since $\|T_t\|\le1$,
\begin{equation*}
\|T_tv-v\|\le\|T_t(v-[F])\|+\|T_t[F]-[F]\|+\|[F]-v\|\le2\|v-[F]\|+\|T_t[F]-[F]\|.
\end{equation*}
The classes $[F]$ are dense, so the first term on the right can be made smaller than any given
positive number; for that $F$ the second term is smaller than the same number for all small $t$, by
the preceding display. So $T_t\to1$ strongly as $t\downarrow0$. For
$t\ge0$ and $h>0$ we then have
\begin{equation*}
\|T_{t+h}v-T_tv\|=\|T_t(T_hv-v)\|\le\|T_hv-v\|.
\end{equation*}
When $t-h\ge0$, the same bound holds for $\|T_tv-T_{t-h}v\|$. Hence $t\mapsto T_t$ is strongly
continuous on $[0,\infty[$.

\emph{Kernel and Hamiltonian.} The family $(T_t)_{t\ge0}$ is a semigroup of positive self-adjoint
contractions. It is strongly continuous at $0$, and $T_0=1$. We show that $\ker T_t=\{0\}$ for
every $t\ge0$, and that $T_t=e^{-tH}$ with $H:=-\log T_1$ self-adjoint and non-negative on all of
$\mathcal{H}^{\mathrm{curv}}$. This is the content of the semigroup lemma of
Remark~\ref{7.4:rem-standard-facts}; we carry out both steps.

\emph{The kernel.} Suppose $T_tv=0$. Then
\begin{equation*}
\|T_{t/2}v\|^2=(v,T_tv)=0,
\end{equation*}
so $T_{t/2}v=0$. By induction, $T_{t2^{-k}}v=0$ for every $k$. Since
$T_{t2^{-k}}v\to v$ by strong continuity at $0$, we get $v=0$.

\emph{The identification.} The operator $T_{t/2}$ is positive and $T_{t/2}^2=T_t$. By uniqueness
of the positive square root of a bounded positive operator (Remark~\ref{7.4:rem-standard-facts}),
$T_{t/2}=T_t^{1/2}$. Hence $T_{2^{-k}}=T_1^{2^{-k}}$ by induction on $k$: if
$T_{2^{-k}}=T_1^{2^{-k}}$, then $T_{2^{-k-1}}$ and $T_1^{2^{-k-1}}$ are both positive square roots
of this operator, so they coincide by the same uniqueness. By the semigroup law, $T_u=T_1^u$ for
every dyadic rational $u\ge0$.

Let $E_1$ be the spectral measure of $T_1$. It is supported in $[0,1]$, since $T_1$ is a positive
contraction. Since $\ker T_1=\{0\}$, we have $E_1(\{0\})=0$. So
\begin{equation*}
H:=-\log T_1=\int_{]0,1]}(-\log\lambda)\,dE_1(\lambda)
\end{equation*}
is self-adjoint and non-negative on all of $\mathcal{H}^{\mathrm{curv}}$, and $T_1^u=e^{-uH}$. Both
$t\mapsto T_t$ and $t\mapsto e^{-tH}$ are strongly continuous, and the dyadic rationals are dense.
Therefore $T_t=e^{-tH}$ for all $t\ge0$. The operator $H$ is unique: if also $T_t=e^{-tH'}$ for all
$t\ge0$ with $H'$ self-adjoint, then $e^{-H'}=T_1$, and since $\lambda\mapsto e^{-\lambda}$ is
injective on $\mathbb{R}$ with inverse $-\log$ on $]0,\infty[$, the functional calculus gives
$H'=-\log T_1=H$.

Finally, $T_1\Omega=\Omega$ places $\Omega$ in the range of $E_1(\{1\})$, where $-\log\lambda$
vanishes. Hence $H\Omega=0$.

\emph{Space translations, momenta, joint spectrum.} Let $\vec a\in\mathbb{R}^3$. The map
$\mathsf{S}_{\vec a}$ maps $\mathcal{P}_+$ onto itself, with inverse $\mathsf{S}_{-\vec a}$. It
commutes with $\Theta$, because $\theta$ fixes the spatial coordinates $\vec x$. It also commutes
with $\mathsf{T}_t$. By the exact invariance \eqref{7.5:eq-exact-invariance-a},
\begin{equation*}
(\mathsf{S}_{\vec a}F,\mathsf{S}_{\vec a}G)_{\mathrm{OS}}
=\sum_{q,q'}S_{q+q'}\bigl(\mathsf{S}_{\vec a}[(\Theta\bar F)_q\otimes G_{q'}]\bigr)
=(F,G)_{\mathrm{OS}}.
\end{equation*}

So $\mathsf{S}_{\vec a}$ preserves $\mathcal{N}$ and induces a surjective isometry of
$\mathcal{P}_+/\mathcal{N}$, that is, a unitary $U(\vec a)$ on $\mathcal{H}^{\mathrm{curv}}$. It
satisfies:
\begin{itemize}
\item $U(\vec a)U(\vec b)=U(\vec a+\vec b)$ and $U(0)=1$;
\item $U(\vec a)T_t=T_tU(\vec a)$, first on the dense quotient and then by continuity;
\item $U(\vec a)\Omega=\Omega$.
\end{itemize}

\emph{Strong continuity of $U$.} We have
\begin{equation*}
\|U(\vec a)[F]-[F]\|^2=2(F,F)_{\mathrm{OS}}-2\operatorname{Re}(F,\mathsf{S}_{\vec a}F)_{\mathrm{OS}}.
\end{equation*}
This tends to $0$ as $\vec a\to0$, by the argument used for strong continuity of $T_t$, with
$\tau_{(0,\vec a)}$ in place of $\tau_{te_0}$. That argument uses the translation modulus of
Lemma~\ref{7.4:lem-translation-modulus} and the same separation margins. Since the classes $[F]$
are dense and $\|U(\vec a)\|=1$, $U$ is strongly continuous.

By the SNAG theorem (Remark~\ref{7.4:rem-standard-facts}),
\begin{equation*}
U(\vec a)=\int e^{i\vec a\cdot\vec k}\,dE(\vec k)
\end{equation*}
for a unique projection-valued measure $E$. The operators
$P_l:=\int k_l\,dE(\vec k)$, $l=1,2,3$, are self-adjoint and commute strongly pairwise.

\emph{Commutation with $H$.} Each $U(\vec a)$ commutes with $T_1$. By the last clause of the SNAG
theorem, it then commutes with the spectral measure of $T_1$. It therefore commutes with the
spectral measure $E_H$ of $H=-\log T_1$, which consists of the same projections. Each $E_H(\Delta)$
is bounded and commutes with every $U(\vec a)$, so it commutes with $E$. Hence $P_l$ and $H$
commute strongly. The pair $(H,\vec P)$ has a joint spectral measure on $\sigma(H)\times\mathbb{R}^3$,
and
\begin{equation*}
\sigma(H)\times\mathbb{R}^3\subset[0,\infty[\times\mathbb{R}^3.
\end{equation*}
Finally, $\vec P\Omega=0$, because $U(\vec a)\Omega=\Omega$ for all $\vec a$.

\emph{Signed permutations of the spatial coordinates.} Let $\varsigma\in W_3$ be a signed
permutation of $x_1,x_2,x_3$. As an element of $W_4$, $\varsigma$ has the following properties:
\begin{itemize}
\item it fixes $x^0$ and preserves $\{x^0>0\}$;
\item it commutes with $\theta$;
\item $\varsigma e_0=e_0$ and $\varsigma(0,\vec a)=(0,\varsigma\vec a)$.
\end{itemize}
Hence $F\mapsto \varsigma\cdot F$, with $\varsigma$ acting on every argument, maps $\mathcal{P}_+$
onto itself. It commutes with $\Theta$ and $\mathsf{T}_t$, and
$\varsigma\,\mathsf{S}_{\vec a}\,\varsigma^{-1}=\mathsf{S}_{\varsigma\vec a}$. It preserves
$(\cdot,\cdot)_{\mathrm{OS}}$, by the $W_4$-invariance in \eqref{7.5:eq-exact-invariance-a}.

As for the space translations, this yields unitaries $U(\varsigma)$ with the following properties:
\begin{itemize}
\item $U(\varsigma)\Omega=\Omega$;
\item $U(\varsigma)T_t=T_tU(\varsigma)$, so that $U(\varsigma)$ commutes with $H$, as in the
commutation argument for $U(\vec a)$;
\item $U(\varsigma)U(\vec a)U(\varsigma)^*=U(\varsigma\vec a)$;
\item $U(\varsigma)E(\Delta)U(\varsigma)^*=E(\varsigma\Delta)$.
\end{itemize}
The joint spectrum of $(H,\vec P)$ is therefore $W_3$-symmetric. The argument has used only the
bound \eqref{7.4:eq-infinite-volume-a}, in the form in which its constant depends on the radius of
the individual balls containing the supports (the per-ball volume-free bound), the exact invariance
\eqref{7.5:eq-exact-invariance-a}, the reflection positivity of the centred limit, the translation
modulus of Lemma~\ref{7.4:lem-translation-modulus}, and the Cauchy--Schwarz inequality, uniqueness
of positive square roots, the semigroup lemma and the SNAG theorem recalled in
Remark~\ref{7.4:rem-standard-facts}; no temperedness, no rotation invariance and no clustering
enters. This completes the proof of clause (D).
\end{proof}

\begin{remark}[Properties not obtained in infinite volume]\label{7.7:rem-not-obtained}
Let the standing data be those of the infinite-volume theorem (Theorem~\ref{7.4:thm-infinite-volume}).
Let $\mathcal{H}^{\mathrm{curv}}$, $\Omega$, $H$ and $\vec P=(P_1,P_2,P_3)$ be the objects constructed there,
in the construction of the reconstructed space (\ref{7.7:prf-reconstructed-space}) and in that of the
transfer semigroup, the Hamiltonian and the momenta (\ref{7.7:prf-semigroup-hamiltonian}). None of the
following is obtained from that theorem.
\begin{itemize}
\item[(a)] The growth condition (E0$'$) of Osterwalder and Schrader \cite{OS75}, which is needed beyond
reflection positivity and translation invariance to reconstruct fields and Wightman functions. It has
the shape
\begin{equation}\label{7.7:eq-not-obtained-1}
|S_q(\Phi)|\le \sigma_q\,\|\Phi\|_{(\nu q)},\qquad \sigma_q\le c_{\mathrm{OS}}\,(q!)^{\beta_{\mathrm{OS}}},
\end{equation}
for every $q$ and every Schwartz test function $\Phi$ of $q$ arguments, where $\|\cdot\|_{(\nu q)}$ is a
Schwartz norm of order $\nu q$, $\nu$ is a fixed degree per argument, and $c_{\mathrm{OS}}$,
$\beta_{\mathrm{OS}}$ are constants. The bound \eqref{7.4:eq-infinite-volume-a}, with
the constant $C^{\infty}_n(\mathsf{d};R)$ that the per-ball volume-free bound supplies, does not deliver
\eqref{7.7:eq-not-obtained-1}.
\item[(b)] Field operators, Wightman functions, and Lorentz covariance.
\item[(c)] Uniqueness of the vacuum $\Omega$.
\item[(d)] The relativistic spectrum condition $H\ge|\vec P|$.
\item[(e)] Any statement for $n=1$ or at coinciding points.
\item[(f)] The inequality $\dim\mathcal{H}^{\mathrm{curv}}>1$; it is a consequence stated in clause (iv) of
Theorem~\ref{3.1:thm-main} and drawn in Chapter~\ref{ch:8}, not one of the infinite-volume theorem.
\item[(g)] Convergence on tuples of test functions that are merely Lipschitz. Uniqueness of the limit
in $(K,m)$. Independence of the regularisation.
\end{itemize}
\end{remark}

\begin{proof}
\emph{Item} (a). Only the shape \eqref{7.7:eq-not-obtained-1} of the condition is used.

In the language of per-ball bounds, the condition requires constants $\widetilde C_n(\mathsf{d};R)$,
in place of $C^{\infty}_n(\mathsf{d};R)$ in \eqref{7.4:eq-infinite-volume-a}, satisfying the following,
with fixed $\bar c,\bar a,\bar b>0$:
\begin{equation}\label{7.7:eq-not-obtained-2}
\widetilde C_n(\mathsf{d};R)\le \bar c^{\,n}\,(n!)^{\bar b}\,\mathsf{d}^{-\bar a}\,
\bigl(1+\log^{+}\mathsf{d}^{-1}\bigr)\,(1+R)^{\bar b}.
\end{equation}
Here $\log^{+}x=\max\{\log x,0\}$. In words, \eqref{7.7:eq-not-obtained-2} allows polynomial blow-up
at the diagonals, polynomial growth at infinity, and factorial growth in $n$.

The per-ball volume-free bound gives its constant in the form
\begin{equation}\label{7.7:eq-not-obtained-3}
C^{\infty}_n(\mathsf{d};R)=\Bigl[2C_{\mathrm{UV}}(\mathsf{d};R)
+\log 2\cdot\mathfrak{r}^{\,n}\bigl(2Q_{\infty}^{\,n}+2\bigr)^{n}\Bigr]\Bigl(\frac{n}{r}\Bigr)^{n}.
\end{equation}
The quantity $Q_\infty$ is
\begin{equation*}
Q_{\infty}=\exp\bigl[(E_+^{\sharp}+C_{\mathrm{low}})\,\ell(n,R,g)^4+2C_{\mathrm{face}}(R)\bigr].
\end{equation*}
Here $E_+^{\sharp}$, $C_{\mathrm{low}}$, $\mathfrak{r}$, $C_{\mathrm{face}}(R)$ and
$C_{\mathrm{UV}}(\mathsf{d};R)$ are constants of the per-ball volume-free bound; their dependence on
$n$, $\mathsf{d}$ and $R$ is the one indicated by the arguments written, so that $E_+^{\sharp}$,
$C_{\mathrm{low}}$ and $\mathfrak{r}$ depend on none of $n$, $\mathsf{d}$, $R$.
Here $\ell(n,R,g)$ is the least power of $L$ that is at least $8n(R+1)$, and $r>0$ is the standing
parameter fixed first. We examine the three dependences of \eqref{7.7:eq-not-obtained-3} in turn.

First, the dependence on $\mathsf{d}$ enters only through $C_{\mathrm{UV}}(\mathsf{d};R)$, for which
\begin{equation*}
C_{\mathrm{UV}}(\mathsf{d};R)\le C(L,\mathfrak{p})\Bigl(1+\frac{R+2}{\mathsf{d}}\Bigr)^{4}
\Bigl(1+\log^{+}_{L}\frac{32M}{\mathsf{d}}\Bigr),
\end{equation*}
where $\log^{+}_L x=\max\{\log_L x,0\}$. This dependence is polynomial, which is what
\eqref{7.7:eq-not-obtained-2} allows.

Second, the constant does not depend on the positions of the supports. This is better than
\eqref{7.7:eq-not-obtained-2} requires at infinity.

Third, the dependence on $n$ contains the factor $(2Q_{\infty}^{\,n}+2)^n$. Since
$2Q_\infty^{\,n}+2\ge Q_\infty^{\,n}$, we have
\begin{equation*}
\bigl(2Q_{\infty}^{\,n}+2\bigr)^{n}\ge Q_{\infty}^{\,n^2}
=\exp\bigl[n^2\bigl((E_+^{\sharp}+C_{\mathrm{low}})\,\ell(n,R,g)^4+2C_{\mathrm{face}}(R)\bigr)\bigr].
\end{equation*}
Since $\ell(n,R,g)\ge 8n(R+1)$, we have $\ell(n,R,g)^4\ge 8^4\,n^4(R+1)^4$. Hence, when
$E_+^{\sharp}+C_{\mathrm{low}}>0$, the exponent is at least
\begin{equation*}
8^4\,(E_+^{\sharp}+C_{\mathrm{low}})\,n^6(R+1)^4+2n^2C_{\mathrm{face}}(R).
\end{equation*}
So, unless $E_+^{\sharp}+C_{\mathrm{low}}$ vanishes, the factor $(2Q_{\infty}^{\,n}+2)^n$ of the
constant \eqref{7.7:eq-not-obtained-3} is bounded below by
\begin{equation*}
\exp\bigl[8^4\,(E_+^{\sharp}+C_{\mathrm{low}})\,n^6(R+1)^4+2n^2C_{\mathrm{face}}(R)\bigr],
\end{equation*}
a lower bound that grows faster than factorially in $n$ and faster than polynomially in $R$.
Therefore \eqref{7.7:eq-not-obtained-2} is not met, and
\eqref{7.7:eq-not-obtained-1} is not obtained.

We do not address here whether, for each fixed $n$, the function $S^T_n$ extends continuously to
Schwartz functions vanishing to infinite order on the diagonals.

\emph{Item} (b). Field operators, Wightman functions and Lorentz covariance require two things:
\begin{itemize}
\item the growth condition of item (a), which is not obtained;
\item invariance under the full Euclidean group, including $O(4)$.
\end{itemize}
Only invariance under the hyperoctahedral group $W_4$ is available here. Invariance under $O(4)$ is
not among the standing data or the conclusions of the infinite-volume theorem: in Theorem~0
(\ref{3.1:thm-main}) it appears only in clause (vi), under the far-sphere flux hypothesis
$(\mathrm{IR\text{-}fl})_{\omega_0}$, which is not assumed in this chapter.

\emph{Item} (c). In the framework of Osterwalder and Schrader, uniqueness of the vacuum is equivalent
to the cluster property \cite[(E4)]{OS73}. No cluster property is proved or claimed here.

\emph{Item} (d). The condition $H\ge|\vec P|$ requires Euclidean rotations that mix $x^0$ with the
spatial coordinates, and these are not available. What is obtained in the construction of the transfer
semigroup, the Hamiltonian and the momenta (\ref{7.7:prf-semigroup-hamiltonian}) is weaker:
\begin{itemize}
\item $H\ge 0$;
\item the joint spectrum of $(H,\vec P)$ lies in $[0,\infty)\times\mathbb{R}^3$;
\item this joint spectrum is symmetric under $W_3$.
\end{itemize}

\emph{Item} (e). The correlation theorem claims nothing for $n=1$ or at coinciding points. Accordingly,
no bounds uniform in $j$ are obtained for $\langle\mathcal{O}(f)\rangle_j$ or for
$\langle{:}\mathcal{O}(f){:}^2\rangle_j$.

The choice $S_1:=0$ in the Osterwalder--Schrader data is a centring, not a computed one-point function.
It is legitimate for the following reason. A joint cumulant of order at least two is unchanged when
each of its variables is shifted by a constant. Hence the connected functions with $n\ge2$ do not
depend on the centring.

\emph{Item} (f). The inequality $\dim\mathcal{H}^{\mathrm{curv}}>1$ would need a real test function $f$
supported in the open half-space $\{x^0>0\}$ with $S^T_2(\theta f,f)\neq 0$. No such $f$ is exhibited
in the present setting of smeared curvature observables with separated supports. The case
$\mathcal{H}^{\mathrm{curv}}=\mathbb{C}\Omega$ is therefore not excluded by what is proved in this chapter; for
reference couplings below the ceiling of clause (iv) it is excluded by the consequences of that clause,
drawn in Chapter~\ref{ch:8} after Corollary~\ref{8.22:cor-reference-torus-kernel}.

\emph{Item} (g). Three further properties are not asserted by the infinite-volume theorem and are not
obtained:
\begin{itemize}
\item convergence on tuples that are Lipschitz but not $C^1$;
\item independence of the limit from the subsequence of $(K_j,m_j)$;
\item independence from the lattice regularisation.
\end{itemize}
\end{proof}

\begin{remark}[Curvature and block reconstructions are not identified]
\label{7.7:rem-spaces-not-identified}
Let $\mathcal{H}^{\mathrm{curv}}$ be the Hilbert space reconstructed from the Osterwalder--Schrader form
$(\cdot,\cdot)_{\mathrm{OS}}$ of \eqref{7.4:eq-infinite-volume-b}. Let $\mathcal{H}^{\mathrm{sm}}$ be the
space reconstructed from the joint block and smeared limit state $\omega_{\infty}$ on the quasilocal
algebra $\mathfrak{A}^{\tau}_{\mathrm{loc}}\otimes\mathfrak{A}^{\mathrm{sm}}_{\mathrm{loc}}$ of bounded
block and smeared observables. This state is obtained along the subsequence $\mathcal{J}$ of
\eqref{7.4:eq-infinite-volume-d}: the diagonal argument there also runs over a countable family of
bounded block and smeared observables, uniformly bounded in $K$, and $\omega_{\infty}$ is built from the
limits along $\mathcal{J}$ of their expectations. The spaces
$\mathcal{H}^{\mathrm{curv}}$ and $\mathcal{H}^{\mathrm{sm}}$ are not identified. No map between them is
asserted, no common Osterwalder--Schrader space containing both is asserted, and no mixed correlation
$\omega_{\infty}(G\cdot{:}\mathcal{O}(h_1){:}\cdots)$, with $G$ a bounded block observable and
$h_1,\dots$ test functions, is asserted.
Both spaces are reconstructions from limits of the same lattice states $\mu_{K_j,m_j}$ of
Theorem~\ref{7.4:thm-infinite-volume} along the same
subsequence $\mathcal{J}$ of \eqref{7.4:eq-infinite-volume-d}. Nor is $\mathcal{H}^{\mathrm{curv}}$
identified with the Hilbert space reconstructed from the bounded block observables alone.
\end{remark}

\begin{proof}
We explain why the two spaces are not identified.

\emph{What ties them.} Both are Osterwalder--Schrader spaces of limits of the Gibbs states
$\mu_{K_j,m_j}$, $j\in\mathcal{J}$, fixed in the infinite-volume theorem
(Theorem~\ref{7.4:thm-infinite-volume}); as there, we write $\langle\cdot\rangle_j$ for the
expectation in $\mu_{K_j,m_j}$.

The space $\mathcal{H}^{\mathrm{sm}}$ comes from $\omega_{\infty}$, a state on bounded observables.

The space $\mathcal{H}^{\mathrm{curv}}$ comes from the connected Schwinger functions $S^T_n$. These are
the cumulants of the observables $\mathcal{O}(f)$. Such observables are unbounded in the limit and are not
elements of any $C^{*}$-algebra carried to the continuum. The positivity of the form
$(\cdot,\cdot)_{\mathrm{OS}}$ from which $\mathcal{H}^{\mathrm{curv}}$ is built is the reflection
positivity of the centred limit proved in Lemma~\ref{7.6:prf-centred-positivity}.

No further relation between the two spaces is asserted.

\emph{Mixed terms.} For the bounded block and smeared observables, the algebra generated by those
localised at positive times has finite-cutoff realisations by bounded functions with sup bounds uniform
in $K$, by the uniform boundedness in $K$ of the countable family of block and smeared observables
entering the diagonal argument of \eqref{7.4:eq-infinite-volume-d}, and
$\omega_{\infty}$ is defined on all of it along $\mathcal{J}$. So the limit along $\mathcal{J}$ of the
Osterwalder--Schrader form pairing block with smeared observables exists on this algebra; it is computed
from $\omega_{\infty}$. For
$\mathcal{H}^{\mathrm{sm}}$ and $\mathcal{H}^{\mathrm{curv}}$ the analogous form, pairing bounded block
observables with centred curvature observables, would need the
limits, as $j\to\infty$ in $\mathcal{J}$, of the mixed terms, for a bounded block observable $G$, a test
function $\varphi$, and real test functions $h_1,\dots,h_n\in C^1_c(\mathbb{R}^4)$ with pairwise disjoint
supports, transported to the torus as $h_l^{(m_j)}$:
\begin{equation}\label{7.7:eq-spaces-not-identified-1}
\Bigl\langle (G,\varphi)_{K_j}\cdot\prod_{l}{:}\mathcal{O}\bigl(h_l^{(m_j)}\bigr){:}\Bigr\rangle_j .
\end{equation}
Here $(G,\varphi)_{K}$ is the realisation at cutoff $K$ of a smeared block observable, given in
\eqref{7.7:eq-spaces-not-identified-3} below.

At finite cutoff this form is positive. Suppose moreover that $\varphi$, the $h_l$ and $G$ are
localised in the open half-space $\{x^0>0\}$, and that $j$ is so large that the transported supports lie
in the open slab $H_{+}:=\{0<x^0<L^{m_j}\}$ of $\mathbb{T}_{m_j}$, one of the two open components of the
complement of the plane pair $\{x^0=0\}\cup\{x^0=L^{m_j}\}$. Then both factors in
\eqref{7.7:eq-spaces-not-identified-1} are
bounded measurable functions of the bond variables in $H_{+}$, for the curvature factor by
Lemma~\ref{7.4:lem-half-space}. Reflection positivity of the
lattice measure therefore applies to any finite combination of such products.

However, none of the bounds proved for the lattice measures controls the mixed terms
\eqref{7.7:eq-spaces-not-identified-1} uniformly in $K$, let alone uniformly in $m$, so the diagonal
argument of \eqref{7.4:eq-infinite-volume-d} cannot be applied to them to extract a convergent
subsequence. The generic estimate does not help. For a bounded measurable
function $B$ of the bond variables write $\|B\|_{\infty}$ for the supremum of $|B|$, and put
$Y=\prod_l{:}\mathcal{O}(h_l^{(m_j)}){:}$. The estimate
\begin{equation}\label{7.7:eq-spaces-not-identified-2}
\bigl|\langle B\cdot Y\rangle_j\bigr|
\le \|B\|_{\infty}\,\langle |Y|\rangle_j
\le \|B\|_{\infty}\,\langle Y^{2}\rangle_j^{1/2}
\end{equation}
leads to $\langle Y^{2}\rangle_j$. This is a quantity at coinciding points. It is ultraviolet divergent,
and no bound on it is asserted.

\emph{Smeared block observables.} The bounds of the correlation theorem that are uniform in $m$ come from
Ba{\l}aban's renormalisation steps carried out with the complex sources $w$. These steps are run on
Ba{\l}aban's single fixed hierarchy of $L$-adic
grids, made of cubes with corners at points of $L^{-k}\mathbb{Z}^4$, $0\le k\le K$
\cite[(0.1)]{B-CMP109}, with no insertion. The insertion identity behind clause (ii-a) of Theorem~0
(\ref{3.1:thm-main})
carries only a block observable on this fixed hierarchy: $G$ composed with the $(K-s)$-fold
block-averaging map, at $w=0$.

A smeared block observable, by contrast, is realised at cutoff $K$ as follows. Take a block observable
$G$ of depth $s$ and a test function $\varphi$. Let $\alpha_{\tau_v}$ denote the action on observables of
the translation $\tau_v$. Then
\begin{equation}\label{7.7:eq-spaces-not-identified-3}
(G,\varphi)_K=\sum_{v\in L^{-K}\mathbb{Z}^4}L^{-4K}\,\varphi(v)\,\bigl(\alpha_{\tau_v}G\bigr)_K .
\end{equation}
This is a Riemann sum over bare translates of block functions. For the $v$ not in $L^{-s}\mathbb{Z}^4$
these block functions live on grids shifted by $v$. Relative to the fixed hierarchy they are generic
bounded functions of the bare field, with no block structure that the renormalisation steps can carry.
A product of $q$ such smeared block observables is a $q$-fold sum, over independent translates
$(v_1,\dots,v_q)$, of functions living on several grids at once.
Consequently the renormalisation steps with complex sources cannot carry $(G,\varphi)_K$. Write $A(w)$
for the action weighted by the complex source $w=\sum_i z_i f_i$, as in
Definition~\ref{1.1:def-wilson-action}. Neither the correlation theorem nor the insertion identity
supplies a bound, uniform in $m$, on sourced expectations of the type
\begin{equation*}
\frac{\bigl\langle \prod_i (G_i,\varphi_i)_{K_j}\,e^{A(w)}\bigr\rangle_j}
{\bigl\langle e^{A(w)}\bigr\rangle_j},
\end{equation*}
through which the mixed terms \eqref{7.7:eq-spaces-not-identified-1} would have to be controlled.
\end{proof}

\paragraph{Inputs used by statement}
This list covers the proofs of this chapter. They use two kinds of input by statement: published statements of Ba{\l}aban, each given with its reference, and statements of this book whose own proof is incomplete or outstanding. Each clause of Theorem~0 (\ref{3.1:thm-main}) has one item. Only clause (ii) is proved in this chapter, and its item is divided by sub-clause, each entry naming the sub-clause it serves.

\begin{description}
\item[(0)] No statement of clause (0) is proved in this chapter.
\item[(i)] No statement of clause (i) is proved in this chapter.
\item[(iii)] No statement of clause (iii) is proved in this chapter.
\item[(ii)] Published statements cited by statement:
\begin{itemize}
\item \cite[(0.1)]{B-CMP109} --- quoted in Definition~\ref{7.4:not-finite-cutoff}; used for
clause (ii-$\mathbb{R}^4$) in Lemma~\ref{7.4:lem-half-space}, in the proof of reflection
positivity of the centred limit, and in the comparison of the curvature reconstruction with the
block reconstruction.
\item \cite[(0.2)]{B-CMP109} --- quoted in Definition~\ref{7.4:not-finite-cutoff}; used for
clause (ii-$\mathbb{R}^4$) in Lemma~\ref{7.4:lem-lattice-symmetries}.
\end{itemize}
No published statement of Ba{\l}aban is cited by statement in the proofs of clauses
(ii-$\mathbb{T}$-a) and (ii-$\mathbb{T}$-b).

Statements of this book whose proofs are incomplete or outstanding:
\begin{itemize}
\item Non-convergence of the first-passage offsets,
Lemma~\ref{7.3:lem-first-passage-offsets} --- used for clause (ii-$\mathbb{T}$-a)
(Theorem~\ref{7.3:thm-tuned-family}) --- proof outstanding.
\end{itemize}
No other statement of this book whose proof is incomplete or outstanding is used in the proofs
of (ii-$\mathbb{T}$-a), (ii-$\mathbb{T}$-b) and (ii-$\mathbb{R}^4$).
\item[(iv)] No statement of clause (iv) is proved in this chapter.
\item[(v)] No statement of clause (v) is proved in this chapter.
\item[(vi)] No statement of clause (vi) is proved in this chapter.
\end{description}

The proofs of clause (ii) appeal by statement to the two equations of \cite{B-CMP109} listed
above, to the correlation theorem, the per-ball volume-free bound, reflection positivity of the
lattice measure and uniformity of the last-scale constants in the torus exponent, and, for
(ii-$\mathbb{T}$-a), to the uniqueness theorem.
None of these statements is proved using a statement of clause (ii). Clauses (ii-$\mathbb{T}$-b) and (ii-$\mathbb{R}^4$) are deduced from clauses (0) and (iii) and from reflection positivity of the lattice measure. Clause (ii-$\mathbb{T}$-a) is deduced from the uniqueness theorem. In each case the dependence runs in one direction only. The constants entering these proofs are chosen in the order of
Definitions~\ref{2.3:not-stages-first}, \ref{2.3:not-stages-middle} and \ref{2.3:not-stages-last}. By Lemma~\ref{2.4:lem-no-cycle}, this order has no cycle.

%% ---- end inlined file: chapters/c07 ----

\clearpage
\setcounter{section}{7}
\refstepcounter{section}
\section*{Chapter~\thesection\quad A two-point witness of non-triviality: clause (iv)\addcontentsline{toc}{section}{\protect\numberline{\thesection}A two-point witness of non-triviality: clause (iv)}\label{ch:8}}
%% ---- begin inlined file: chapters/c08 ----
\begin{definition}[Separated test pairs, the room and the two-point kernel]
\label{8.1:not-test-pairs-kernel}
Throughout this chapter $\mathbb{T}_m$ is the torus of Definition~\ref{1.1:def-lattice},
$\mathbb{F}$ is the bare lattice of spacing $L^{-K}$ on it ($=\mathbb{T}^{(0)}$ in the notation
of Definition~\ref{1.5:def-block-average}), $\Lambda^2(\mathbb{F})$ is its set
of plaquettes, and $x(p)$ is the barycentre of a plaquette $p$. For a function $f$ on
$\mathbb{T}_m$, $\|f\|_\infty=\sup_{\mathbb{T}_m}|f|$ is its supremum norm and
$\operatorname{Lip} f$ its Lipschitz constant.
\begin{enumerate}
\item[(a)] \emph{Separated test pairs.} Let $\vartheta\in(0,1]$ be a constant fixing the shape of
the test functions, and let $\mathsf{d}>0$. The class $\mathfrak{T}(\mathsf{d};\vartheta)$ is
the set of pairs $(f_1,f_2)$ of Lipschitz functions $f_i\colon\mathbb{T}_m\to[0,\infty)$,
$f_i\not\equiv0$, for which there are points $x_1,x_2\in\mathbb{T}_m$ such that, $B(x,\mathsf{d})$
denoting the ball of radius $\mathsf{d}$ about $x$ for the distance of $\mathbb{T}_m$,
\begin{equation}\label{8.1:eq-test-pairs-kernel-1}
\begin{gathered}
\operatorname{supp} f_i\subset B(x_i,\mathsf{d}),\qquad
\int_{\mathbb{T}_m} f_i\ \ge\ \vartheta\,\|f_i\|_\infty\,\mathsf{d}^4,\\
\operatorname{Lip} f_i\ \le\ \frac{\|f_i\|_\infty}{\vartheta\,\mathsf{d}}
\quad(i=1,2),\\
\mathsf{d}\ \le\ \operatorname{dist}(\operatorname{supp} f_1,\operatorname{supp} f_2)
\ \le\ 6\,\mathsf{d}.
\end{gathered}
\end{equation}
The last line is a two-sided condition on the pair $(f_1,f_2)$ of test functions.
\item[(b)] \emph{The room.} The witness-scale envelope of the running inverse coupling is
\begin{equation}\label{8.1:eq-test-pairs-kernel-2}
\bar{x}_s\ :=\ g^{-2}+B^{\sharp}\,(s+1),\qquad s=0,1,2,\dots,
\end{equation}
with $g$ the reference coupling and $B^{\sharp}$ the flow constant of
Definition~\ref{1.2:def-flow-constants}; it does not involve the cutoff $K$. The \emph{room}
is a function $s\mapsto C_W(s)\ge1$ that depends on $s$ only through $\bar{x}_s$ and is
non-decreasing in $s$; in this chapter it is
\begin{equation}\label{8.1:eq-test-pairs-kernel-3}
C_W(s)\ =\ C_{W,0}\,\bigl[p_0\bigl(\bar{x}_s^{-1/2}\bigr)\bigr]^2
\ =\ C_{W,0}\,A_0^2\,(\log\bar{x}_s)^{2p_0},
\end{equation}
which is the room of \eqref{8.1:eq-coupling-envelope-a}; here $p_0(\cdot)$ is the degree
function, $A_0$ and the exponent $p_0$ are auxiliary parameters of $\mathfrak{p}$, and
$C_{W,0}$ is the constant of the room fixed together with \eqref{8.1:eq-coupling-envelope-a}
(it is not Ba{\l}aban's auxiliary parameter $C_0$); all three enter as in
\eqref{8.1:eq-coupling-envelope-a}. This form of the room satisfies
$C_W(0)\ge1$, as fixed together with \eqref{8.1:eq-coupling-envelope-a}.
\item[(c)] \emph{Witness depth and support radius.} For $\mathsf{d}>0$,
\begin{equation}\label{8.1:eq-test-pairs-kernel-4}
s(\mathsf{d})\ :=\ \min\bigl\{\,s\ge0:\ C_W(s)\,L^{-s}\le\mathsf{d}\,\bigr\},
\qquad \rho\ :=\ L^{s}\,\mathsf{d}\quad\text{with } s=s(\mathsf{d}).
\end{equation}
When $s(\mathsf{d})\le K$, the number $\rho$ is the radius $\mathsf{d}$ measured in cells of
the block lattice $\mathbb{T}^{(K-s)}$ of Definition~\ref{1.5:def-block-average}, whose
spacing is $L^{-s}$.
The depth $s(\mathsf{d})$ is well defined and does not depend on $K$; it is
non-increasing in $\mathsf{d}$ and tends to $\infty$ as $\mathsf{d}\downarrow0$; one has
$\rho\ge C_W(s)$, and $\rho<L\,C_W(s)$ whenever $s(\mathsf{d})\ge1$; for $s(\mathsf{d})=0$
one has $\rho=\mathsf{d}$.
\item[(d)] \emph{The two-point kernel.} Let $\mathcal{A}$ be the free lattice field on
$\mathbb{F}$, a $1$-cochain, $\gamma$ its Gaussian law, $\mathbb{E}_{\gamma}$ the
expectation in $\gamma$, and $Q_{K-s}$ the linearised block-averaging map to
level $K-s$, $0\le s\le K$; these objects are those of the Gaussian
extraction of this chapter. The background curvature, its covariance and the kernel are
\begin{equation}\label{8.1:eq-test-pairs-kernel-5}
\begin{aligned}
F_B\ &:=\ d_1\,\mathbb{E}_{\gamma}\bigl[\mathcal{A}\,\big|\,Q_{K-s}\mathcal{A}\bigr],
\qquad
C^{\mathrm{bg}}_{K,s}(p,q)\ :=\ \mathbb{E}_{\gamma}\bigl[F_B(p)\,F_B(q)\bigr],\\
\mathcal{K}_{K,s}(f_1,f_2)\ &:=\ \sum_{p,q\in\Lambda^2(\mathbb{F})}
f_1\bigl(x(p)\bigr)\,f_2\bigl(x(q)\bigr)\,\bigl[C^{\mathrm{bg}}_{K,s}(p,q)\bigr]^2,
\end{aligned}
\end{equation}
where $d_1$ is the lattice coboundary taking $1$-cochains to $2$-cochains.
\item[(e)] \emph{Normalisation.} We put
\begin{equation}\label{8.1:eq-test-pairs-kernel-6}
b_0\ :=\ \frac{N^2-1}{2}.
\end{equation}
\end{enumerate}
\end{definition}

\begin{proof}
We prove the second equality in \eqref{8.1:eq-test-pairs-kernel-3}, the monotonicity of
$C_W$ in that form and the bound $C_W(s)\ge1$ for all $s$ from its value at $s=0$, and the
assertions of (c).

\emph{The two forms of the room.} The degree function is
$p_0(g')=A_0(\log g'^{-2})^{p_0}$. For $g'=\bar{x}_s^{-1/2}$ one has $g'^{-2}=\bar{x}_s$,
hence $\bigl[p_0(\bar{x}_s^{-1/2})\bigr]^2=A_0^2(\log\bar{x}_s)^{2p_0}$; multiplying by $C_{W,0}$
gives the second equality in \eqref{8.1:eq-test-pairs-kernel-3}.

\emph{Positivity and monotonicity.} Since $B^{\sharp}>0$ and $g<1$
(Definition~\ref{1.2:def-flow-constants}, which fixes the flow constants and the reference
window), \eqref{8.1:eq-test-pairs-kernel-2} gives $\bar{x}_s\ge\bar{x}_0$ for every $s\ge0$,
where $\bar{x}_0=g^{-2}+B^{\sharp}>1$, and $\bar{x}_s$ is strictly increasing in $s$. Hence
$\log\bar{x}_s>0$, and
the right-hand side of \eqref{8.1:eq-test-pairs-kernel-3} is a well-defined real number. By (b),
$C_W(0)\ge1$, as fixed together with \eqref{8.1:eq-coupling-envelope-a}; since
$(\log\bar{x}_0)^{2p_0}>0$, it follows that $C_{W,0}A_0^2=C_W(0)/(\log\bar{x}_0)^{2p_0}>0$. As
$p_0>0$, the map $t\mapsto(\log t)^{2p_0}$ is increasing on $(1,\infty)$, so $C_W$ is
non-decreasing in $s$, and $C_W(s)\ge C_W(0)\ge1$ for every $s\ge0$.

\emph{Independence of $K$.} By \eqref{8.1:eq-test-pairs-kernel-2}, $\bar{x}_s$ depends only
on $g$, $B^{\sharp}$ and $s$, and by \eqref{8.1:eq-test-pairs-kernel-3} $C_W(s)$ depends on $s$
only through $\bar{x}_s$ and otherwise on $C_{W,0}$, $A_0$, $p_0$. Hence the set in
\eqref{8.1:eq-test-pairs-kernel-4}, and with it $s(\mathsf{d})$ and $\rho$, do not involve
$K$.

\emph{Well-definedness.} As $\bar{x}_s$ is an affine function of $s$ with $\bar{x}_s>1$,
$\log\bar{x}_s$ is bounded by a constant times $\log(s+2)$, so $C_W(s)$ is bounded by a
constant times $(\log(s+2))^{2p_0}$. Since $L>1$, the factor $L^{-s}$ decays geometrically, and
therefore $C_W(s)L^{-s}\to0$ as $s\to\infty$. For every $\mathsf{d}>0$ the set
$\{s\ge0:C_W(s)L^{-s}\le\mathsf{d}\}$ is thus a non-empty set of non-negative integers, and
its minimum $s(\mathsf{d})$ exists.

\emph{Monotonicity and divergence.} If $0<\mathsf{d}'\le\mathsf{d}$, every $s$ with
$C_W(s)L^{-s}\le\mathsf{d}'$ also satisfies $C_W(s)L^{-s}\le\mathsf{d}$, so the set defining
$s(\mathsf{d}')$ is contained in the set defining $s(\mathsf{d})$, and
$s(\mathsf{d}')\ge s(\mathsf{d})$. Let $S\ge1$ be an integer. Since $C_W\ge1$,
\begin{equation*}
\min_{0\le s<S} C_W(s)\,L^{-s}\ \ge\ L^{-(S-1)}\ >\ 0 .
\end{equation*}
If $\mathsf{d}<L^{-(S-1)}$, no $s<S$ satisfies $C_W(s)L^{-s}\le\mathsf{d}$, hence
$s(\mathsf{d})\ge S$. As $S$ is arbitrary, $s(\mathsf{d})\to\infty$ as $\mathsf{d}\downarrow0$.

\emph{The range of $\rho$.} Let $s=s(\mathsf{d})$. By the definition of $s(\mathsf{d})$,
$C_W(s)L^{-s}\le\mathsf{d}$; multiplying by $L^{s}$ gives $\rho=L^{s}\mathsf{d}\ge C_W(s)$.
If $s=0$, then $\rho=L^{0}\mathsf{d}=\mathsf{d}$.
If $s\ge1$, the minimality of $s$ gives $C_W(s-1)L^{-(s-1)}>\mathsf{d}$; multiplying by
$L^{s}$ and using that $C_W$ is non-decreasing,
\begin{equation*}
\rho\ =\ L^{s}\,\mathsf{d}\ <\ L\,C_W(s-1)\ \le\ L\,C_W(s).
\end{equation*}
This proves the assertions listed at the beginning of the proof.
\end{proof}

\begin{definition}[The curvature observable and the connected two-point function]
\label{8.1:not-curvature-observable}
We write $\mathbb{F}$ for the bare lattice of spacing $L^{-K}$ on the torus $\mathbb{T}_m$ and
$\Lambda^2(\mathbb{F})$ for its set of plaquettes (Definition~\ref{1.1:def-lattice}); $x(p)$ is
the barycentre of $p$, and $\operatorname{tr}$ is normalised by $\operatorname{tr}\mathbf{1}=1$.
For a configuration $U$ and $p\in\Lambda^2(\mathbb{F})$ we put
$a_p(U):=1-\operatorname{Re}\operatorname{tr}U(\partial p)$, and for $f\colon\mathbb{T}_m\to\mathbb{R}$
\begin{equation*}
\mathcal{O}(f):=\sum_{p\in\Lambda^2(\mathbb{F})}f(x(p))\,a_p(U),
\end{equation*}
the smeared curvature field of Definition~\ref{1.3:def-curvature-field}. For test functions
$f_1,f_2$ and $z=(z_1,z_2)\in\mathbb{C}^2$ we put $w:=z_1f_1+z_2f_2$ and
\begin{equation*}
\rho_0^w:=\exp\bigl[-A(x_0-w,U)-E\bigr],\qquad x_0=g_0^{-2},
\end{equation*}
where $A(y,U)=\sum_{p}y(x(p))\,a_p(U)$ is the weighted Wilson action of
Definition~\ref{1.1:def-wilson-action} (the constant $x_0$ read as a constant function) and $E$ is
the constant of the bare density there, which depends neither on $U$ nor on $z$. We put
$Z(z):=\int\rho_0^w\,dU$, with $dU$ the product Haar measure, let $\mu_K$ be the probability measure
$Z(0)^{-1}\rho_0^0\,dU$, and define the connected two-point function
\begin{equation}\label{8.1:eq-curvature-observable-1}
S^T_{2;K}(f_1,f_2):=\partial_{z_1}\partial_{z_2}\log Z(z)\big|_{z=0}
=\operatorname{Cov}_{\mu_K}\bigl(\mathcal{O}(f_1),\mathcal{O}(f_2)\bigr).
\end{equation}
Finally $\dim\mathfrak{su}(N)=N^2-1$, and for $Y\in\mathfrak{su}(N)$ we put
$|Y|^2:=-\operatorname{tr}(Y^2)$; then
\begin{equation}\label{8.1:eq-curvature-observable-2}
1-\operatorname{Re}\operatorname{tr}e^{Y}=\tfrac12|Y|^2-\tfrac1{24}\operatorname{tr}Y^4+O(|Y|^6),
\end{equation}
where $O(|Y|^6)$ is bounded by $N^3e^{\sqrt N}|Y|^6$ for $|Y|\le1$.
\end{definition}

\begin{proof}
\emph{The second equality in \eqref{8.1:eq-curvature-observable-1}.} By the definition of the
weighted action, $A(x_0-w,U)=x_0A(U)-z_1\mathcal{O}(f_1)-z_2\mathcal{O}(f_2)$, so
$\rho_0^w=\rho_0^0\exp[z_1\mathcal{O}(f_1)+z_2\mathcal{O}(f_2)]$ and
\begin{equation*}
Z(z)=Z(0)\int\exp\bigl[z_1\mathcal{O}(f_1)+z_2\mathcal{O}(f_2)\bigr]\,d\mu_K .
\end{equation*}
Here $Z(0)>0$ because $\rho_0^0>0$. The configuration space is a finite product of copies of
$\mathrm{SU}(N)$, hence compact, and $\mathcal{O}(f_i)$ is bounded and continuous on it; therefore
$Z$ is entire in $z$, differentiation under the integral is permitted, and $Z(z)\neq0$ near $z=0$.
Differentiating, $\partial_{z_1}\log Z=Z^{-1}\int\mathcal{O}(f_1)\rho_0^w\,dU$, and differentiating
this in $z_2$ and setting $z=0$ gives
$\int\mathcal{O}(f_1)\mathcal{O}(f_2)\,d\mu_K-\int\mathcal{O}(f_1)\,d\mu_K\int\mathcal{O}(f_2)\,d\mu_K$,
which is the covariance.

\emph{The expansion \eqref{8.1:eq-curvature-observable-2}.} Since $Y$ is anti-Hermitian its
eigenvalues are $i\lambda_1,\dots,i\lambda_N$ with $\lambda_j\in\mathbb{R}$, so
$\operatorname{tr}Y^k=N^{-1}\sum_j(i\lambda_j)^k$ is real for even $k$ and purely imaginary for odd
$k$; in particular $|Y|^2=N^{-1}\sum_j\lambda_j^2\ge0$. Hence
$\operatorname{Re}\operatorname{tr}e^Y=\sum_{k\ \mathrm{even}}\operatorname{tr}Y^k/k!$, whose terms
$k=0,2,4$ are $1$, $-\tfrac12|Y|^2$ and $\tfrac1{24}\operatorname{tr}Y^4$. For the rest,
$|\operatorname{tr}Y^k|\le\max_j|\lambda_j|^k$ and $\max_j|\lambda_j|\le\sqrt N|Y|$, so for
$|Y|\le1$
\begin{equation*}
\sum_{k\ge6}\frac{|\operatorname{tr}Y^k|}{k!}
\le N^3|Y|^6\sum_{k\ge6}\frac{N^{(k-6)/2}}{k!}
\le N^3e^{\sqrt N}|Y|^6 .
\end{equation*}
\end{proof}

\begin{proposition}[The correlation theorem at the witness depth]\label{8.1:prop-correlation-witness-depth}
Assume the correlation theorem in the following form: it holds for every value of its parameter $m$
with $m_{\mathbb{T}}(g)\le m\le K$ and for every value $N_T\ge 1$ of its torus parameter. Here and
throughout this proposition $K$ is the bare cutoff, so that the bare spacing is $L^{-K}$
(Definition~\ref{1.1:def-lattice}). This $m$
belongs to the statement of the correlation theorem and is set to $s$ below; the torus of the
two-point witness theorem enters through $N_T$. None of its
constants depends on $K$, on the level $k$, on $m$, on the torus or on the position, and $N$ enters
only through the norm $\|\cdot\|_{\sharp}$, through the constant $\mathsf{Q}$ and through the lower
bound $m\ge m_{\mathbb{T}}(g)$ on the torus exponent.

Let $s$ be an integer with $m_{\mathbb{T}}(g)\le s\le K$. Apply the correlation theorem with $m:=s$,
and with $N_T$ the size of the depth-zero torus of the two-point witness theorem. Put
$k_0:=K-s$, $\eta:=L^{-k_0}$ and $\rho:=L^{s}\mathsf{d}$. Then the following hold.
\begin{itemize}
\item[(a)] The parameter $K'$ of the correlation theorem, the number of blocking steps from
its bare lattice to its unit lattice, equals $K-s=k_0$, and its unit lattice is the block lattice
$\mathbb{T}^{(k_0)}$ of Definition~\ref{1.5:def-block-average}, the witness lattice. Lengths are
measured in cells of
$\mathbb{T}^{(k_0)}$; in these units the bare spacing is $\eta$ and the support scale $\mathsf{d}$
of the test functions is $\rho$. The constants of the correlation theorem depend on none of $s$,
$K$, $k_0$, the running level $k$, the torus or the position.
\item[(b)] At $m=s$ the notation of the correlation theorem specialises as follows. Write
$\operatorname{Lip}_j$ for the Lipschitz constant in units of cells of $\mathbb{T}^{(j)}$, as in the
correlation theorem, and let
$r_j$, $r$ and $\lambda_*$ be the radii and the constant of the correlation theorem. Then
\begin{align}
\|f\|_{\sharp}&=\max\{\|f\|_\infty,\ 40M\operatorname{Lip}_{k_0}f\},\label{8.1:eq-correlation-witness-depth-1}\\
\mathcal{W}_j(X)&=\{\,|w|<r_j,\ \operatorname{Lip}_jw<\mathfrak{l}_j\,\},\qquad
\mathfrak{l}_j=r\lambda_*L^{-(k_0-j)/2},\label{8.1:eq-correlation-witness-depth-2}\\
\mathbb{D}&=\Bigl\{\,\sum_i|z_i|\,\|f_i\|_{\sharp}<r\,\Bigr\}.\label{8.1:eq-correlation-witness-depth-3}
\end{align}
Moreover, $w=\sum_iz_if_i\in\mathcal{W}_{k_0}$ for $(z_i)\in\mathbb{D}$. The lower bound on the
torus exponent reads $s\ge m_{\mathbb{T}}(g)$. This is a one-sided floor on the depth, it depends
on $g$, and it is one of the conditions of the fixed-depth requirement of the two-point witness
theorem. For a level $j\le k_0$ we put $\nu:=k_0-j\ge0$. In this chapter a level index of the form
$k_0+i$ occurs only inside the Gaussian extraction. The Lipschitz radius $\mathfrak{l}_j$ of
\eqref{8.1:eq-correlation-witness-depth-2} is not the unindexed letter $\mathfrak{l}$, which later
in this chapter denotes a number of levels.
\item[(c)] Let $f$ be a member of a pair in $\mathfrak{T}(\mathsf{d};\vartheta)$
(Definition~\ref{8.1:not-test-pairs-kernel}). Then
$\operatorname{Lip}_{k_0}f\le\|f\|_\infty/(\vartheta\rho)$. Consequently
\begin{equation}\label{8.1:eq-correlation-witness-depth-4}
\|f\|_{\sharp}=\|f\|_\infty\qquad\text{whenever }\rho\ge 40M/\vartheta .
\end{equation}
The lower bound $\rho\ge40M/\vartheta$ is one of the conditions collected in
\eqref{8.1:eq-constants-ceiling-c}.
\item[(d)] The following items of the correlation theorem hold at $m=s$ as stated there, with the
same constants.
\begin{itemize}
\item[--] The running inverse couplings satisfy $x_{k+1}=x_k-b_{k+1}(0)$, and the correlation
theorem uses the number $\lambda:=c_0/2$; this $\lambda$ is a constant, not a gauge transformation
as in Lemma~\ref{1.1:lem-invariance}.
\item[--] The running-shift statement holds. With $d_k(\varrho)=\theta'_k+\theta''_k/\varrho$,
where $\theta'_k$ and $\theta''_k$ are the two level-$k$ constants of the running-shift statement,
as fixed in the correlation theorem,
\begin{equation}\label{8.1:eq-correlation-witness-depth-5}
|b_{k+1}(\zeta)-b_{k+1}(0)|\le d_k(\varrho)|\zeta|,\qquad \sum_{k<k_0}d_k\le\tfrac13,\qquad
N_{k_0}\in[\tfrac23,\tfrac43].
\end{equation}
\item[--] The coupling increment at shift $\zeta$ is $b_j(\zeta):=\beta[E^{(j)}_\zeta]$, where
$E^{(j)}_\zeta$ is the level-$j$ effective action at shift $\zeta$ and the functional
$\beta[\,\cdot\,]$ extracts the coupling increment, both as fixed in the correlation theorem; the
contribution of the operation $\tilde{\mathbf{R}}$ is included in $E^{(j)}_\zeta$.
\item[--] The sourced effective densities admit the representation
\begin{equation}\label{8.1:eq-correlation-witness-depth-6}
\rho^w_k=\sum_h\chi_k\,\mathbf{T}^w_k\exp A^w_k ,
\end{equation}
in which
\begin{equation}\label{8.1:eq-correlation-witness-depth-7}
A^w_k=-A(x_k(\cdot)-\zeta_k,U_k)+\mathbf{E}^w+\mathbf{D}^w+\mathbf{R}^w+\mathbf{B}^w-E^w ;
\end{equation}
here $E^w$, in italic type, is the term subtracted in the exponent as fixed in the correlation
theorem, and it is a different object from the term $\mathbf{E}^w$.
Here the running shift is
\begin{equation}\label{8.1:eq-correlation-witness-depth-8}
\zeta_k(p):=w(p)+\sum_{j\le k}\varphi_j(p)\sum_yh^{(j)}_y(p)\bigl[b_j(w(p_y))-b_j(0)\bigr],
\end{equation}
with $\varphi_j$, $h^{(j)}_y$ and $p_y$ as fixed in the correlation theorem;
and the further clauses of that representation hold as stated in the correlation theorem.
\item[--] Each $\rho^w_k$ has the format \eqref{8.1:eq-correlation-witness-depth-6} and is analytic
in $w$. A term indexed by a region $X$ depends on $w$ only through its values on the enlargement
$X^+$ of $X$ fixed in the correlation theorem. Finally,
$\int\rho^w_k=\int\rho^w_0$, with $\rho^w_0$ as in
Definition~\ref{8.1:not-curvature-observable}.
\item[--] The remaining items of the correlation theorem also hold. Among them are the torus bound
and the bounds of the correlation theorem's unit table on the per-cell sums over scales $S_I$,
$S_\omega$, $\mathfrak{m}$, together with its complete list of such sums; these bounds make the
sums over levels used later in this chapter independent of $K$. The unindexed
$\mathfrak{m}$ is a per-cell sum over scales and is not the running quantity $\mathfrak{m}_k$ of the
correlation theorem.
\end{itemize}
\end{itemize}
\end{proposition}

\begin{proof}
We first check that $m:=s$ is admissible. By hypothesis the correlation theorem holds for every
integer $m$ with $m_{\mathbb{T}}(g)\le m\le K$ and for every $N_T\ge1$. The integer $s$ satisfies
$m_{\mathbb{T}}(g)\le s\le K$, so $m:=s$ is admissible. The size of the depth-zero torus of the
two-point witness theorem is an admissible value of $N_T$. Hence the correlation theorem applies with
these choices.

Proof of (a). In the correlation theorem the parameter $K'$ is $K-m$; with $m=s$ it is
$K'=K-s=k_0$. Its unit
lattice is the block lattice obtained after
$k_0$ blocking steps, that is, $\mathbb{T}^{(k_0)}$ of Definition~\ref{1.5:def-block-average}.

We now convert lengths into cells of $\mathbb{T}^{(k_0)}$. The bare spacing is $L^{-K}$
(Definition~\ref{1.1:def-lattice}). A cell of $\mathbb{T}^{(k_0)}$ is a block of side $L^{k_0}$
bare bonds, so its physical side is $L^{k_0}L^{-K}=L^{-s}$. Measured in such cells, the bare spacing
is therefore $L^{-k_0}=\eta$, and a physical length $\mathsf{d}$ is $L^{s}\mathsf{d}=\rho$ cells.

The constants of the correlation theorem do not depend on $K$, $k$, $m$, the torus or the position.
In particular they do not depend on $m=s$, hence not on $k_0=K-s$, and not on $K$, on the running
level $k$, on the torus or on the position. This proves (a).

Proof of (b). The displays \eqref{8.1:eq-correlation-witness-depth-1}--\eqref{8.1:eq-correlation-witness-depth-3}
are the definitions of $\|\cdot\|_{\sharp}$, of the
analyticity domains $\mathcal{W}_j(X)$ with their Lipschitz radii $\mathfrak{l}_j$, and of the
polydisc $\mathbb{D}$ in the correlation theorem, written at $m=s$. The unit lattice is
$\mathbb{T}^{(k_0)}$ by (a), so the level difference in $\mathfrak{l}_j$ is $k_0-j$ and the
Lipschitz constant in $\|\cdot\|_{\sharp}$ is $\operatorname{Lip}_{k_0}$. The membership
$\sum_iz_if_i\in\mathcal{W}_{k_0}$ on $\mathbb{D}$ is part of the statement of the correlation
theorem. The lower bound $m\ge m_{\mathbb{T}}(g)$ on the torus exponent becomes, at $m=s$, the
condition $s\ge m_{\mathbb{T}}(g)$. This condition is an inequality in one direction only, and its
right-hand side depends on $g$.

Proof of (c). Let $f$ be a member of a pair in $\mathfrak{T}(\mathsf{d};\vartheta)$. By
Definition~\ref{8.1:not-test-pairs-kernel}, in physical units
$\operatorname{Lip}f\le\|f\|_\infty/(\vartheta\mathsf{d})$. Measuring distances in cells of
$\mathbb{T}^{(k_0)}$ multiplies every length by $L^{s}$, by (a), so the Lipschitz constant measured
in these units is at most $L^{-s}\operatorname{Lip}f$. Therefore
\begin{equation*}
\operatorname{Lip}_{k_0}f\le L^{-s}\operatorname{Lip}f\le\frac{\|f\|_\infty}{\vartheta L^{s}\mathsf{d}}
=\frac{\|f\|_\infty}{\vartheta\rho}.
\end{equation*}
If $\rho\ge40M/\vartheta$, it follows that
\begin{equation*}
40M\operatorname{Lip}_{k_0}f\le\frac{40M}{\vartheta\rho}\,\|f\|_\infty\le\|f\|_\infty .
\end{equation*}
Hence the maximum in \eqref{8.1:eq-correlation-witness-depth-1} equals $\|f\|_\infty$, which is
\eqref{8.1:eq-correlation-witness-depth-4}.

Proof of (d). Each item listed in (d) is part of the statement of the correlation theorem. By the
first paragraph, that theorem holds at $m=s$. By (a), its constants are the same as for every other
admissible $m$. Hence each item holds at $m=s$, as stated, with those constants. In the
running-shift statement the sum of the $d_k$ runs over the levels below the unit lattice of the
correlation theorem, and the normalisation is read at that lattice; by (a) the unit lattice is
reached after $K'=k_0$ blocking steps, so at $m=s$ these read $\sum_{k<k_0}d_k\le\frac13$ and
$N_{k_0}\in[\frac23,\frac43]$, as in \eqref{8.1:eq-correlation-witness-depth-5}.
\end{proof}

\begin{definition}[The witness test class and witness pairs]\label{8.1:def-witness-class}
We work on the witness lattice of the correlation theorem at the witness depth
(Proposition~\ref{8.1:prop-correlation-witness-depth}). Thus $k_0=K-s$, lengths are measured in
$k_0$-cells, that is, in units of the lattice spacing of the witness lattice, $\mathbb{F}$ is the
bare torus lattice, of spacing $\eta=L^{-k_0}$ in these units, and
$\rho:=L^s\mathsf{d}$. For a function $f\ge0$ on the torus, $\|f\|_\infty$ is its supremum,
$\operatorname{Lip}_{k_0}f$ is its Lipschitz constant for distances measured in $k_0$-cells, and
\begin{equation}\label{8.1:eq-witness-class-1}
\mathfrak{n}(f):=\eta^4\sum_{x\in\mathbb{F}}f(x).
\end{equation}
For $\delta>0$, the $\delta$-neighbourhood of a subset of the torus, the set of points at distance at
most $\delta$ from it, is denoted by the superscript $+\delta$. Let
$\vartheta$ be a shape constant as in Notation~\ref{8.1:not-test-pairs-kernel}.
\begin{itemize}
\item[(a)] $\mathfrak{F}^{+}(\mathsf{d},\vartheta)$ is the set of Lipschitz functions $f\ge0$,
$f\not\equiv0$, whose support lies in a ball of radius $\rho$ and which satisfy
\begin{equation}\label{8.1:eq-witness-class-2}
\mathfrak{n}(f)\ge\vartheta\,\|f\|_\infty\,\rho^4,\qquad
\operatorname{Lip}_{k_0}f\le\frac{\|f\|_\infty}{\vartheta\rho}.
\end{equation}
\item[(b)] A \emph{witness pair for the shape constant} $\vartheta$ is a pair
$f_1,f_2\in\mathfrak{F}^{+}(\mathsf{d},\vartheta)$ with
$\rho\le\operatorname{dist}(\operatorname{supp}f_1,\operatorname{supp}f_2)\le6\rho$, that is, with
separation between $\mathsf{d}$ and $6\mathsf{d}$ in physical units.
\item[(c)] For a witness pair, $\mathcal{N}_i:=(\operatorname{supp}f_i)^{+\rho/4}$ and
$\mathcal{N}_i^{++}:=(\operatorname{supp}f_i)^{+3\rho/8}$, $i=1,2$.
\item[(d)] $b_0:=n/2$, where $n=N^2-1$; $\mathcal{K}_{K,s}$ denotes the two-point kernel at finite
cutoff $K$ corresponding to the two-point kernel $\mathcal{K}_s$ of
Notation~\ref{8.1:not-test-pairs-kernel}; and, with the running
shift $\zeta_k$ and the increments $b_j$ of
Proposition~\ref{8.1:prop-correlation-witness-depth}, the normalisation at level $k$ is
\begin{equation}\label{8.1:eq-witness-class-3}
N_k:=\zeta_k'(0)=1+\sum_{j\le k}b_j'(0),
\end{equation}
the second equality because $\zeta_k$ arises from $\zeta$ by adding the increments $b_j$, $j\le k$;
$N_k$ is the quantity denoted $\zeta'_k$ in clause (iv) of
Theorem~\ref{3.1:thm-main}.
\end{itemize}
Then:
\begin{itemize}
\item[(e)] every $f\in\mathfrak{F}^{+}(\mathsf{d},\vartheta)$ satisfies
$\mathfrak{n}(f)\le10\rho^4\|f\|_\infty$ whenever $\eta\le\rho/6$;
\item[(f)] for a witness pair, $\rho\le\operatorname{dist}(x,y)\le10\rho$ for all
$x\in\operatorname{supp}f_1$, $y\in\operatorname{supp}f_2$, and
$\operatorname{dist}(\mathcal{N}_1^{++},\mathcal{N}_2^{++})\ge\rho/4$;
\item[(g)] sup-normalised dilates of two fixed bumps are witness pairs with one $\vartheta$:
let $\varphi_1,\varphi_2\colon\mathbb{R}^4\to[0,\infty[$ be Lipschitz, $\varphi_i\not\equiv0$,
$\operatorname{supp}\varphi_i\subset B(c_i,1)$, with
$1\le\operatorname{dist}(\operatorname{supp}\varphi_1,\operatorname{supp}\varphi_2)\le6$, and put
$\mathsf{r}_\varphi:=2+\max_i|c_i|$. There is $\vartheta_\varphi\in\,]0,1]$ depending only on
$\varphi_1,\varphi_2$ with the following property. Let $o$ be a point of the torus, and suppose that
$\rho=L^s\mathsf{d}$ is such that $4\mathsf{r}_\varphi\rho$ is smaller than the side of the torus in
$k_0$-cells and
\begin{equation*}
\eta\le\rho\,\min_{i=1,2}\frac{\int\varphi_i}{32\operatorname{Lip}\varphi_i}.
\end{equation*}
Define
\begin{equation*}
f_i(o+\rho z):=\frac{\varphi_i(z)}{\|\varphi_i\|_\infty}\quad(|z|\le\mathsf{r}_\varphi),\qquad
f_i:=0\ \text{elsewhere}.
\end{equation*}
Then $(f_1,f_2)$ is a witness pair for the shape constant $\vartheta_\varphi$. One may take
\begin{equation*}
\vartheta_\varphi:=\min_{i=1,2}\min\Bigl\{\frac{\int\varphi_i}{2\|\varphi_i\|_\infty},\
\frac{\|\varphi_i\|_\infty}{\operatorname{Lip}\varphi_i}\Bigr\}.
\end{equation*}
\item[(h)] the pairs of Notation~\ref{8.1:not-test-pairs-kernel}, read in $k_0$-cells, are witness
pairs for the shape constant $\vartheta/2$ once $\eta\le\vartheta^2\rho/32$: let
$(f_1,f_2)\in\mathfrak{T}(\mathsf{d};\vartheta)$, read with lengths measured in $k_0$-cells, so
that physical lengths are multiplied by $L^s$. If
\begin{equation*}
\eta\le\frac{\vartheta^2\rho}{32},
\end{equation*}
then $f_1,f_2\in\mathfrak{F}^{+}(\mathsf{d},\vartheta/2)$ and $(f_1,f_2)$ is a witness pair for
the shape constant $\vartheta/2$.
\end{itemize}
\end{definition}

\begin{proof}
Since $k_0=K-s$, at the witness depth
\begin{equation*}
\frac{\eta}{\rho}=\frac{L^{-k_0}}{L^s\mathsf{d}}=\frac{L^{-K}}{\mathsf{d}},
\end{equation*}
so each of the thresholds on $\eta$ in (e), (g) and (h) holds once $K$ is large enough at fixed
$\mathsf{d}$, $\vartheta$, $\varphi_1$, $\varphi_2$.

(e) Let $\operatorname{supp}f\subset B(c,\rho)$. The cubes of side $\eta$ centred at the points of
$\mathbb{F}\cap B(c,\rho)$ are disjoint, and each lies in $B(c,\rho+\eta)$, since the half-diagonal
of a four-dimensional cube of side $\eta$ is $\eta$. A ball of the torus has volume at most that of
the Euclidean ball of the same radius, being the image of it under the covering map. Since
$0\le f(x)\le\|f\|_\infty$,
\begin{equation*}
\mathfrak{n}(f)\le\|f\|_\infty\,\eta^4\,\#\bigl(\mathbb{F}\cap B(c,\rho)\bigr)
\le\tfrac{\pi^2}{2}(\rho+\eta)^4\|f\|_\infty\le10\rho^4\|f\|_\infty ,
\end{equation*}
the last step because $\rho+\eta\le\tfrac76\rho$ for $\eta\le\rho/6$ and
$\tfrac{\pi^2}{2}(7/6)^4<10$.

(f) The lower bound is the definition of the distance of two sets. The supports are compact; choose
$x'\in\operatorname{supp}f_1$, $y'\in\operatorname{supp}f_2$ with
$\operatorname{dist}(x',y')\le6\rho$. Each support lies in a ball of radius $\rho$, so
$\operatorname{dist}(x,x')\le2\rho$ and $\operatorname{dist}(y,y')\le2\rho$. Hence
$\operatorname{dist}(x,y)\le10\rho$. If $u\in\mathcal{N}_1^{++}$, $v\in\mathcal{N}_2^{++}$, there are
$x\in\operatorname{supp}f_1$, $y\in\operatorname{supp}f_2$ with
$\operatorname{dist}(u,x),\operatorname{dist}(v,y)\le3\rho/8$, whence
$\operatorname{dist}(u,v)\ge\rho-3\rho/4=\rho/4$.

(g) For a displacement $\xi$ with $|\xi|$ less than half the side of the torus, the torus
distance of $u$ and $u+\xi$ is $|\xi|$. Hence
$z\mapsto o+\rho z$ is isometric, up to the factor $\rho$, on $B(0,\mathsf{r}_\varphi)$, as
differences there have length at most $2\mathsf{r}_\varphi\rho$. Consequently $\|f_i\|_\infty=1$,
$\int f_i=\rho^4\int\varphi_i/\|\varphi_i\|_\infty$,
$\operatorname{supp}f_i\subset B(o+\rho c_i,\rho)$, and the distance of the supports is $\rho$ times
that of the $\operatorname{supp}\varphi_i$, hence lies in $[\rho,6\rho]$.

Since $\varphi_i$ is continuous and vanishes off a ball of radius $1$, it vanishes at some point at
distance at most $1$ from a maximiser, so $\|\varphi_i\|_\infty\le\operatorname{Lip}\varphi_i$; in
particular $0<\vartheta_\varphi\le1$. Let $u,v$ lie in the image of $B(0,\mathsf{r}_\varphi)$.
Then
\begin{equation*}
|f_i(u)-f_i(v)|\le\frac{\operatorname{Lip}\varphi_i}{\rho\|\varphi_i\|_\infty}\,
\operatorname{dist}(u,v).
\end{equation*}

Now let $f_i(u)>0$ and let $v$ lie outside the image. Then $\operatorname{dist}(u,v)\ge\rho$:
otherwise $v=u+\xi$ with $|\xi|<\rho$, and $u=o+\rho z$ with $\varphi_i(z)>0$, so
$|z-c_i|\le1$; then $v=o+\rho(z+\xi/\rho)$ with
\begin{equation*}
|z+\xi/\rho|\le|z-c_i|+|c_i|+|\xi|/\rho<|c_i|+2\le\mathsf{r}_\varphi,
\end{equation*}
and $v$ would lie in the image. Hence
\begin{equation*}
\frac{f_i(u)}{\operatorname{dist}(u,v)}\le\frac1\rho
\le\frac{\operatorname{Lip}\varphi_i}{\rho\|\varphi_i\|_\infty},
\end{equation*}
and in all cases
\begin{equation*}
\operatorname{Lip}_{k_0}f_i\le\frac{\operatorname{Lip}\varphi_i}{\rho\|\varphi_i\|_\infty}
\le\frac{1}{\vartheta_\varphi\rho}.
\end{equation*}

For the mass, let $\mathsf{c}_x$ be the cube of side $\eta$ centred at $x\in\mathbb{F}$; these cubes
tile the torus. We have
\begin{equation*}
\bigl|\eta^4f_i(x)-\textstyle\int_{\mathsf{c}_x}f_i\bigr|\le\eta^5\operatorname{Lip}_{k_0}f_i .
\end{equation*}
This is non-zero only if $x$ lies within $\eta$ of $\operatorname{supp}f_i$, and then
$\mathsf{c}_x\subset B(o+\rho c_i,\rho+2\eta)$. As in (e), there are at most
$\tfrac{\pi^2}{2}(\rho+2\eta)^4\eta^{-4}$ such $x$. The hypothesis on $\eta$ gives $\eta\le\rho/6$.
Indeed, by the support condition and the bound $\|\varphi_i\|_\infty\le\operatorname{Lip}\varphi_i$,
\begin{equation*}
\int\varphi_i\le\frac{\pi^2}{2}\|\varphi_i\|_\infty\le\frac{\pi^2}{2}\operatorname{Lip}\varphi_i ,
\end{equation*}
so that $\eta\le\pi^2\rho/64$, and $\pi^2/64<1/6$. Moreover $\tfrac{\pi^2}{2}(4/3)^4<16$. Therefore
\begin{equation*}
\Bigl|\mathfrak{n}(f_i)-\int f_i\Bigr|
\le16\,\frac{\eta}{\rho}\,\frac{\operatorname{Lip}\varphi_i}{\|\varphi_i\|_\infty}\,\rho^4
\le\frac{\rho^4\int\varphi_i}{2\|\varphi_i\|_\infty},
\end{equation*}
and so
\begin{equation*}
\mathfrak{n}(f_i)\ge\frac{\rho^4\int\varphi_i}{2\|\varphi_i\|_\infty}\ge\vartheta_\varphi\rho^4\|f_i\|_\infty .
\end{equation*}
The constant $\vartheta_\varphi$ depends on $\varphi_1,\varphi_2$ only.

(h) A $k_0$-cell has physical side $L^{-s}$, so a physical length $l$ is $L^sl$ in
$k_0$-cells, and $\rho=L^s\mathsf{d}$. Hence each $f_i$ is Lipschitz, $f_i\ge0$, $f_i\not\equiv0$,
its support lies in a ball of radius $\rho$, and
$\rho\le\operatorname{dist}(\operatorname{supp}f_1,\operatorname{supp}f_2)\le6\rho$. The Lipschitz
constant in $k_0$-cells is $L^{-s}$ times the physical one, so
\begin{equation*}
\operatorname{Lip}_{k_0}f_i\le\frac{L^{-s}\|f_i\|_\infty}{\vartheta\mathsf{d}}
=\frac{\|f_i\|_\infty}{\vartheta\rho}\le\frac{\|f_i\|_\infty}{(\vartheta/2)\rho}.
\end{equation*}
Volumes in $k_0$-cell units are $L^{4s}$ times physical volumes, so the integral of $f_i$ in these
units, written $\int f_i$ as in (g), satisfies $\int f_i\ge\vartheta\|f_i\|_\infty\rho^4$. Since
$\vartheta\le1$, the hypothesis gives $\eta\le\rho/32\le\rho/6$. With the cubes $\mathsf{c}_x$ of
the proof of (g), the cube estimate of that proof holds for $f_i$; the difference it bounds
is non-zero only if $x$ lies within $\eta$ of $\operatorname{supp}f_i$, and then $\mathsf{c}_x$ lies
in the ball of radius $\rho+2\eta$ with the same centre as the ball containing
$\operatorname{supp}f_i$. As in (e) there are at most
$\tfrac{\pi^2}{2}(\rho+2\eta)^4\eta^{-4}<16\rho^4\eta^{-4}$ such $x$, since $\eta\le\rho/6$ and
$\tfrac{\pi^2}{2}(4/3)^4<16$. Therefore
\begin{equation*}
\Bigl|\mathfrak{n}(f_i)-\int f_i\Bigr|\le16\,\eta\,\rho^4\operatorname{Lip}_{k_0}f_i
\le16\,\frac{\eta}{\vartheta\rho}\,\rho^4\|f_i\|_\infty
\le\frac{\vartheta}{2}\,\rho^4\|f_i\|_\infty ,
\end{equation*}
the last step by $\eta\le\vartheta^2\rho/32$. Hence
$\mathfrak{n}(f_i)\ge\tfrac{\vartheta}{2}\|f_i\|_\infty\rho^4$. Since $\vartheta/2\in\,]0,1]$, this
shows $f_i\in\mathfrak{F}^{+}(\mathsf{d},\vartheta/2)$, and with the separation above $(f_1,f_2)$ is a
witness pair for $\vartheta/2$.

The items (a)--(d) are
definitions.
\end{proof}

\begin{definition}[Envelope of the running coupling, depth floor and coupling ceiling]
\label{8.1:def-coupling-envelope}
Let $g>0$ be the reference coupling, put $X:=g^{-2}$, and let $x_k=g_k^{-2}$ be the running
inverse couplings of Definition~\ref{1.2:def-couplings}. Let $\lambda^{\sharp}$, $B^{\sharp}$ and
$c_5$ be the constants of Definition~\ref{1.2:def-flow-constants}. Let $p_0(\cdot)$ be the degree
function, $p_0(g')=A_0(\log g'^{-2})^{p_0}$. The room constant $C_{W,0}$, the depth offset $r_0$,
a constant $C_D$ and the level $k_W(L)$ are constants of the two-point witness theorem, fixed in
the choice of its constants; of these, $C_{W,0}$, $r_0$ and $k_W(L)$ are used below. They are
chosen from $L$, $M$, $\mathfrak{p}$, $N$, $\vartheta$ and the other fixed parameters only.

\emph{Envelope.} For integers $s\ge0$ and real $y$ put
\begin{equation}\label{8.1:eq-coupling-envelope-1}
\underline{x}_s:=X+\lambda^{\sharp}s-c_5,\qquad
\bar{x}_s:=X+B^{\sharp}(s+1),\qquad
\bar{x}(y):=\frac{B^{\sharp}}{\lambda^{\sharp}}\,(y+c_5)+B^{\sharp}.
\end{equation}
For every $K$ and every $0\le s\le K$ these satisfy
\begin{equation}\label{8.1:eq-coupling-envelope-f}
\underline{x}_s\;\le\;x_{K-s}\;\le\;\bar{x}_s\;\le\;\bar{x}(\underline{x}_s).
\end{equation}
The first two inequalities follow from clause~(0) of Theorem~\ref{3.1:thm-main}, read in the
reference window of clause~(ii-b). This window is the reference window of the correlation theorem,
written in that theorem's constants. The third is proved below. None of the three bounds depends
on $K$.
The affine function $\bar{x}(\cdot)$ is the only envelope function used in this chapter.

\emph{Room.} Put
\begin{equation}\label{8.1:eq-coupling-envelope-a}
C_W(s):=C_{W,0}\,p_0\bigl(\bar{x}_s^{-1/2}\bigr)^2=C_{W,0}A_0^2(\log\bar{x}_s)^{2p_0}
\;\ge\;C_{W,0}\,p_0(g_{K-s})^2 .
\end{equation}
The room is taken at the envelope $\bar{x}_s$ of the inverse coupling at the witness scale $K-s$,
so that it dominates the room at $g_{K-s}$ itself. The room at the witness scale is evaluated at
$\bar{x}_s$, not at $X=g^{-2}$.
Since $C_W$ is non-decreasing in $s$ (proved below), for every $s\ge0$
\begin{equation*}
C_W(s)\;\ge\;C_W(0)\;=\;C_{W,0}\,p_0\bigl(\bar{x}_0^{-1/2}\bigr)^2 ;
\end{equation*}
in particular $C_W(s)\ge C_{W,0}$ for every $s\ge0$ whenever $p_0(\bar{x}_0^{-1/2})\ge1$.

\emph{Witness depth.} For $\mathsf{d}>0$ put
\begin{equation}\label{8.1:eq-coupling-envelope-b}
s(\mathsf{d}):=\min\{s\ge0:\ L^{s}\mathsf{d}\ge C_W(s)\}.
\end{equation}
Let $\rho=L^{s(\mathsf{d})}\mathsf{d}$ be the support radius in cells of
Proposition~\ref{8.1:prop-correlation-witness-depth}. If $s(\mathsf{d})\ge1$, then
$C_W(s(\mathsf{d}))\le\rho<L\,C_W(s(\mathsf{d}))$.
In all cases $\rho$ is at least $C_W(s(\mathsf{d}))$, hence at least $C_W(0)$; so
$\rho\ge C_{W,0}$ at the witness depth whenever $p_0(\bar{x}_0^{-1/2})\ge1$.

\emph{Extraction depth.} Put
\begin{equation}\label{8.1:eq-coupling-envelope-c}
\begin{aligned}
r(s)&:=\bigl\lceil\log_L\bigl(20L\,C_W(s)\bigr)\bigr\rceil
+\bigl\lceil(p_0/2+1)\log_L\log\bar{x}_s\bigr\rceil+r_0,\\
r_1&:=\bigl\lceil\log_L(20L\rho)\bigr\rceil .
\end{aligned}
\end{equation}

\emph{Depth floor.} Put
\begin{equation}\label{8.1:eq-coupling-envelope-d}
s_W(g):=\min\bigl\{s\ge0:\ s-r(s)\ge s_*,\ \ s\ge m_{\mathbb{T}}(g),\ \
2B^{\sharp}\bigl(r(s)+1\bigr)\le\underline{x}_s\bigr\}.
\end{equation}
Statements at the witness depth are made under the hypothesis $s(\mathsf{d})\ge s_W(g)$,
equivalently $\mathsf{d}\le\mathsf{d}_W(g)$, for a threshold $\mathsf{d}_W(g)$ depending on $g$
only. It is a one-sided condition and depends on $g$. It is of the same kind as the condition of
clause~(ii-d) of Theorem~\ref{3.1:thm-main}, and milder than it. This hypothesis is retained in
every statement made at the witness depth; it is not discharged.

\emph{Ultraviolet condition and coupling ceiling.} The witness statements assume
\begin{equation}\label{8.1:eq-coupling-envelope-e}
K-s(\mathsf{d})\ge k_W(L),\qquad g\le\gamma_W\le\gamma_J .
\end{equation}
Here $\gamma_J$ is the coupling threshold of the correlation theorem, used at the witness depth
in Proposition~\ref{8.1:prop-correlation-witness-depth}. The threshold $\gamma_W$ is fixed last,
after every other constant. Each condition defining it has the form
$\sup_{y\ge X-c}\Phi(y)\le C$, where $c$ and $C$ are constants and $\Phi$ is non-increasing with
$\Phi(y)\to0$ as $y\to\infty$.

None of the conditions above bounds from above $L$, $K$, $m$, the torus parameter $N_T$ of the
correlation theorem, the number of steps, an age or a depth, and none bounds a coupling from
below. The torus exponent $m$ is subject to the torus clause of the two-point witness theorem and
to no further condition. The torus lower bound of the correlation theorem, applied in
Proposition~\ref{8.1:prop-correlation-witness-depth} with its torus exponent equal to $s$ and with
$N_T$ free, becomes the depth floor $s\ge m_{\mathbb{T}}(g)$.
\end{definition}

\begin{proof}
We prove the assertions of the definition that are not definitions.

\emph{The third inequality of \eqref{8.1:eq-coupling-envelope-f}.} The definitions
\eqref{8.1:eq-coupling-envelope-1} give
\begin{equation}\label{8.1:eq-coupling-envelope-2}
\bar{x}(\underline{x}_s)-\bar{x}_s
=\frac{B^{\sharp}}{\lambda^{\sharp}}(X+\lambda^{\sharp}s)+B^{\sharp}-X-B^{\sharp}s-B^{\sharp}
=\frac{B^{\sharp}-\lambda^{\sharp}}{\lambda^{\sharp}}\,X .
\end{equation}
Since $X>0$ and $\lambda^{\sharp}\ge c_0/4>0$, the third inequality is equivalent to
$\lambda^{\sharp}\le B^{\sharp}$. The first two inequalities of
\eqref{8.1:eq-coupling-envelope-f} hold for every $K$ and every $0\le s\le K$, with bounds
independent of $K$. Hence $X+\lambda^{\sharp}s-c_5\le X+B^{\sharp}(s+1)$ for every $s\ge0$.
Dividing by $s+1$ and letting $s\to\infty$ gives $\lambda^{\sharp}\le B^{\sharp}$.

\emph{The room.} By the definition of the degree function,
$p_0(\bar{x}_s^{-1/2})=A_0(\log\bar{x}_s)^{p_0}$. This gives the equality in
\eqref{8.1:eq-coupling-envelope-a}. Likewise $p_0(g_{K-s})=A_0(\log x_{K-s})^{p_0}$. The
function $t\mapsto(\log t)^{2p_0}$ is non-decreasing on $[1,\infty[$, and
$x_{K-s}\le\bar{x}_s$ by \eqref{8.1:eq-coupling-envelope-f}. Hence
$C_{W,0}p_0(g_{K-s})^2\le C_W(s)$ for $x_{K-s}\ge1$, the domain on which the degree
function is defined. Since $B^{\sharp}\ge\lambda^{\sharp}>0$, $\bar{x}_s$ is increasing in $s$. By
the same monotonicity, $C_W$ is non-decreasing in $s$.
Consequently $C_W(s)\ge C_W(0)$ for every $s\ge0$, and $C_W(0)=C_{W,0}\,p_0(\bar{x}_0^{-1/2})^2$
is at least $C_{W,0}$ when $p_0(\bar{x}_0^{-1/2})\ge1$.

\emph{The witness depth.} Let $\mathsf{d}>0$. Since $L\ge3$, $L^{s}\mathsf{d}$ grows
geometrically in $s$. On the other hand, $C_W(s)=C_{W,0}A_0^2(\log(X+B^{\sharp}(s+1)))^{2p_0}$
grows like a power of $\log s$. Hence $L^{s}\mathsf{d}\ge C_W(s)$ for all large $s$, the set in
\eqref{8.1:eq-coupling-envelope-b} is a non-empty set of non-negative integers, and
$s(\mathsf{d})$ exists. By definition, $\rho=L^{s(\mathsf{d})}\mathsf{d}\ge C_W(s(\mathsf{d}))$.
Together with $C_W(s(\mathsf{d}))\ge C_W(0)$, proved above, this gives $\rho\ge C_W(0)$, and
hence $\rho\ge C_{W,0}$ when $p_0(\bar{x}_0^{-1/2})\ge1$.
Suppose $s(\mathsf{d})\ge1$. Minimality gives
$L^{s(\mathsf{d})-1}\mathsf{d}<C_W(s(\mathsf{d})-1)$, and $C_W$ is non-decreasing, so
\begin{equation*}
L^{s(\mathsf{d})-1}\mathsf{d}<C_W(s(\mathsf{d})-1)\le C_W(s(\mathsf{d})).
\end{equation*}
Multiplying by $L$ gives $\rho<L\,C_W(s(\mathsf{d}))$.

\emph{Existence of the depth floor.} By the formula for $C_W$ and by
\eqref{8.1:eq-coupling-envelope-c},
\begin{align*}
\log_L\bigl(20LC_W(s)\bigr)&=\log_L\bigl(20LC_{W,0}A_0^2\bigr)+2p_0\log_L\log\bar{x}_s,\\
r(s)&\le\log_L\bigl(20LC_{W,0}A_0^2\bigr)+\Bigl(\frac{5p_0}{2}+1\Bigr)\log_L\log\bar{x}_s
+r_0+2 .
\end{align*}
Since $\bar{x}_s=X+B^{\sharp}(s+1)$ is affine in $s$, the right-hand side grows like
$\log\log s$. Thus $r(s)/s\to0$ as $s\to\infty$. We check the three conditions of
\eqref{8.1:eq-coupling-envelope-d} in turn.
\begin{itemize}
\item[(a)] We have $(s-r(s))/s\to1$, so $s-r(s)\ge s_*$ for all large $s$.
\item[(b)] The number $m_{\mathbb{T}}(g)$ does not depend on $s$, so $s\ge m_{\mathbb{T}}(g)$ for
all large $s$.
\item[(c)] Write
\begin{equation*}
\underline{x}_s-2B^{\sharp}\bigl(r(s)+1\bigr)=X-c_5-2B^{\sharp}+\lambda^{\sharp}s-2B^{\sharp}r(s).
\end{equation*}
Divided by $s$, this tends to $\lambda^{\sharp}>0$. Hence it is non-negative for all large $s$.
\end{itemize}
So the set in \eqref{8.1:eq-coupling-envelope-d} contains all large integers and is non-empty,
and $s_W(g)$ exists.

\emph{Uniformity of the coupling ceiling.} The condition $g\le\gamma_W$ is equivalent to
$X\ge\gamma_W^{-2}$. Hence each condition $\sup_{y\ge X-c}\Phi(y)\le C$ entering
\eqref{8.1:eq-coupling-envelope-e} follows from $\sup_{y\ge\gamma_W^{-2}-c}\Phi(y)\le C$. The
quantity on the left tends to $0$ as $\gamma_W\to0$, because $\Phi$ is non-increasing with limit
$0$. The quantities $c$, $C$ and $\Phi$ do not depend on $s$, $K$ or $m$. Hence these conditions
hold uniformly in $s$, $K$ and $m$.
\end{proof}

\begin{lemma}[Splitting off the spanning local terms]\label{8.1:lem-spanning-split}
Let $(f_1,f_2)$ be a witness pair in the sense of Definition~\ref{8.1:def-witness-class}, with
the neighbourhood $N_1$ of $\operatorname{supp}f_1$ defined there, and use the units of the
correlation theorem at the witness depth (Proposition~\ref{8.1:prop-correlation-witness-depth}):
lengths are measured in $k_0$-cells, and the supports are at distance at least $\rho$.
Let $w=z_1f_1+z_2f_2$, $z=(z_1,z_2)$, be the complex source and $Z=Z(z)$ the partition function
with source $w$, and write $\partial_i=\partial/\partial z_i$; we write $S^T_{2;K}$ for the
connected two-point function $S^T_{2;K,m}(f_1,f_2)$ on the torus of
Proposition~\ref{8.1:prop-correlation-witness-depth}, the torus index and the test functions
being suppressed, so that $S^T_{2;K}=\partial_1\partial_2\log Z(0)$
(Definition~\ref{1.3:def-curvature-field}). Assume the following, which the proof of the torus
bound in Proposition~\ref{8.1:prop-correlation-witness-depth} supplies. For $j\le k_0$ put
$\nu:=k_0-j$; for each large-field region $X$ of level $j$, with enlargement $X^{+}$, the local term
$e_j(X;z):=e^{(j)}(X;w)-e^{(j)}(X;0)$ is the difference of the level-$j$ local exponent terms
of the torus bound at source $w$ and at source $0$, and:
\begin{itemize}
\item[(a)] $e_j(X;\cdot)$ is analytic on the domain $\mathbb{D}$, which contains the polydisc
$\{|z_i|<r/(2\|f_i\|_{\sharp}),\ i=1,2\}$, where $r$ is the radius of the analyticity domains of
Proposition~\ref{8.1:prop-correlation-witness-depth}; $e_j(X;0)=0$;
$e_j(X;z)$ does not depend on $z_i$ unless $X^{+}\cap\operatorname{supp}f_i\neq\emptyset$;
and on $\mathbb{D}$, for constants $C_e>0$ and $\delta<\frac34$ of the torus bound,
\[
|e_j(X;z)|\le C_e\,e^{-(1-\delta)\kappa d_j(X)};
\]
\item[(b)] $Z=e^{\mathcal{E}}\tilde{Z}$, where
$\mathcal{E}(z)=\sum_{j\le k_0}\sum_{X}e_j(X;z)$;
\item[(c)] $\mathsf{h}:=\tilde{Z}/\tilde{Z}(0)$ satisfies $|\mathsf{h}|\le\mathsf{Q}$, where
$\mathsf{Q}$ is the constant of the torus bound of
Proposition~\ref{8.1:prop-correlation-witness-depth}; this is the bound to which the Schwarz lemma
is applied in the proof of the torus bound;
\item[(d)] in these units the unit length of level $j$ is $\mathsf{t}_j=L^{-\nu}$ $k_0$-cells,
and if $X^{+}$ meets both supports, then $4M\mathsf{t}_j(d_j(X)+3)\ge\rho$;
\item[(e)] if $X^{+}$ meets $\operatorname{supp}f_1$, then $X$ contains a level-$j$ cube meeting
$N_1$; the level-$j$ cubes at which the tree bound \cite[(1.26)]{B-CMP116} is anchored form a grid
of side at least $M\mathsf{t}_j/L$ $k_0$-cells; and the
decay rate $(\frac34-\delta)\kappa$ lies in the range in which \cite[(1.26)]{B-CMP116} bounds,
for each fixed level-$j$ cube $\square$, the sum of $e^{-(\frac34-\delta)\kappa d_j(X)}$ over the
regions $X$ of level $j$ containing $\square$ by a constant $O(1)$.
\end{itemize}
Then
\begin{equation}\label{8.1:eq-spanning-split-a}
\begin{aligned}
&S^T_{2;K}=T_{\mathcal{E}}+\partial_1\partial_2\log\tilde{Z}(0),\qquad
T_{\mathcal{E}}:=\sum_{j\le k_0}\ \sum_{X:\ X^{+}\text{ meets both supports}}
\partial_1\partial_2e_j(X;0),\\
&|T_{\mathcal{E}}|\le 4C_e\,r^{-2}\|f_1\|_{\sharp}\|f_2\|_{\sharp}\cdot O(1)
\sum_{\nu\ge0}\Bigl(\frac{4L\rho L^{\nu}}{M}+3\Bigr)^{4}e^{-\frac14\kappa\ell_\nu},\\
&\ell_\nu:=\Bigl[\frac{\rho L^{\nu}}{4M}-3\Bigr]_{+},
\end{aligned}
\end{equation}
where $O(1)$ is the constant of the bound \cite[(1.26)]{B-CMP116} used in the proof.
\end{lemma}

\begin{proof}
Item~(c) is not used in this proof.

\emph{The decomposition.} The torus is finite, so the sums defining $\mathcal{E}$ and
$T_{\mathcal{E}}$ have finitely many terms. By (a), $e_j(X;0)=0$ for all $j,X$, hence
$\mathcal{E}(0)=0$ and, by (b), $\tilde{Z}(0)=Z(0)$, which is the integral of a positive density
and so is positive. By continuity $Z$ and $\tilde{Z}$ do not vanish near $z=0$, and with the
branches of the logarithm that are real at $z=0$ we get from (b)
\[
\log Z(z)=\mathcal{E}(z)+\log\tilde{Z}(z)
\]
near $z=0$. The connected two-point function $S^T_{2;K}$ is the mixed derivative
$\partial_1\partial_2\log Z$ at $z=0$, so
\[
S^T_{2;K}=\partial_1\partial_2\mathcal{E}(0)+\partial_1\partial_2\log\tilde{Z}(0),\qquad
\partial_1\partial_2\mathcal{E}(0)=\sum_{j\le k_0}\sum_X\partial_1\partial_2e_j(X;0).
\]
If $X^{+}$ does not meet $\operatorname{supp}f_i$ for some $i$, then by (a) $e_j(X;\cdot)$ does
not depend on $z_i$, so $\partial_ie_j(X;\cdot)\equiv0$ and $\partial_1\partial_2e_j(X;0)=0$.
Only regions with $X^{+}$ meeting both supports remain, which is the first line of
\eqref{8.1:eq-spanning-split-a}.

\emph{Cauchy's estimate.} Put $r_i:=r/(2\|f_i\|_{\sharp})$. For $0<r_i'<r_i$, Cauchy's formula on
the product of the circles $|z_1|=r_1'$, $|z_2|=r_2'$, which lies in $\mathbb{D}$ by (a), gives
\[
|\partial_1\partial_2e_j(X;0)|\le\frac{1}{r_1'r_2'}\,\sup_{\mathbb{D}}|e_j(X;\cdot)|.
\]
Letting $r_i'\to r_i$ and using $(r_1r_2)^{-1}=4r^{-2}\|f_1\|_{\sharp}\|f_2\|_{\sharp}$ and the
bound in (a),
\[
|\partial_1\partial_2e_j(X;0)|\le4C_e\,r^{-2}\|f_1\|_{\sharp}\|f_2\|_{\sharp}\,
e^{-(1-\delta)\kappa d_j(X)}.
\]

\emph{Lower bound on $d_j(X)$.} Let $X^{+}$ meet both supports and $\nu=k_0-j$. By (d),
$4ML^{-\nu}(d_j(X)+3)\ge\rho$, that is, $d_j(X)\ge\rho L^{\nu}/(4M)-3$; since $d_j(X)\ge0$, this
gives $d_j(X)\ge\ell_\nu$.

\emph{Anchoring and the tree bound.} Fix $j$, hence $\nu$. Each region $X$ in the inner sum of
$T_{\mathcal{E}}$ has $X^{+}$ meeting $\operatorname{supp}f_1$, so by (e) $X$ contains a level-$j$
cube meeting $N_1$; we anchor $X$ at such a cube. By Definition~\ref{8.1:def-witness-class},
$\operatorname{supp}f_1$ lies in a ball of radius $\rho$, so $N_1$, its $\rho/4$-neighbourhood,
lies in a ball of radius $5\rho/4$ and hence in a cube of side $5\rho/2$. In each coordinate
direction an interval of length $5\rho/2$ meets at most $5\rho/(2a)+2$ intervals of a grid of
side $a$; with $a\ge M\mathsf{t}_j/L=ML^{-\nu-1}$ by (e) and (d), the number of level-$j$ cubes
meeting $N_1$ is at most
\[
\Bigl(\frac{5\rho L^{\nu+1}}{2M}+2\Bigr)^{4}\le\Bigl(\frac{4L\rho L^{\nu}}{M}+3\Bigr)^{4}.
\]
Write
\[
e^{-(1-\delta)\kappa d_j(X)}=e^{-\frac14\kappa d_j(X)}\,e^{-(\frac34-\delta)\kappa d_j(X)};
\]
the first
factor is at most $e^{-\frac14\kappa\ell_\nu}$ by the lower bound on $d_j(X)$. Since every region
of the inner sum contains at least one anchor cube, the sum over these regions is at most the sum
over the anchor cubes $\square$ of the sums over the regions $X$ of level $j$ containing
$\square$. For each fixed cube the sum of the second factor over these regions is bounded by a
constant $O(1)$ by the bound \cite[(1.26)]{B-CMP116} on sums over connected sets containing a
given cube, applied with the decay rate $(\frac34-\delta)\kappa$ as in (e). Hence, for each
$j\le k_0$,
\[
\sum_{X:\ X^{+}\text{ meets both supports}}|\partial_1\partial_2e_j(X;0)|
\le4C_e\,r^{-2}\|f_1\|_{\sharp}\|f_2\|_{\sharp}\cdot O(1)
\Bigl(\frac{4L\rho L^{\nu}}{M}+3\Bigr)^{4}e^{-\frac14\kappa\ell_\nu}.
\]

\emph{Summation over levels.} As $j$ runs over the levels $j\le k_0$, $\nu=k_0-j$ runs over
integers $\nu\ge0$, each at most once. Summing the last bound over $j$ and enlarging the sum to all
$\nu\ge0$ gives the second line of \eqref{8.1:eq-spanning-split-a}. The series converges, since
$\ell_\nu$ grows like $\rho L^{\nu}/(4M)$ while the prefactor grows like $L^{4\nu}$, and for
$\rho/M$ large enough, depending on $\kappa$ and $L$, its term $\nu=0$ is the largest.
\end{proof}

\begin{remark}[What the witness theorem uses of the torus bound]\label{8.1:rem-torus-bound-use}
Consider the torus bound of the correlation theorem in the case $n=2$, with the complex source
$w=z_1f_1+z_2f_2$. This bound has two parts, described in items (a) and (b); item (c) states where
the response of the healed large-field regions to the sources sits.
\begin{enumerate}
\item[(a)] Every cross term in $z_1z_2$ of the whole-lattice small-field families and of the
whole-lattice families produced by the operation $\mathbf{R}$
lies in one spanning local exponent term
$e_j(X;z)$, that is, in the spanning sum $T_{\mathcal{E}}$ of \eqref{8.1:eq-spanning-split-a}.
\item[(b)] All large-field and history content enters through the global quotient $\tilde{Z}$. The
control of this quotient is uniform in $K$, depends on the volume, and carries no factor $g^4$.
\item[(c)] For the large-field regions that have been healed, the marginal part of their response to
the sources is contained in the normalisation $N_{K-s}$. Their irrelevant part is contained in the
families of effective-action terms from which the coupling increment $\beta_{k+1}$ of
Definition~\ref{1.2:def-couplings} has been extracted, and in the third term of the
correlation theorem's bound at $n=2$.
\end{enumerate}
The two-point witness theorem does not use the control of the global quotient in item (b). Of the
torus bound it uses only the bounds on the local exponent terms $e_j(X;z)$, through the spanning sum
of item (a). It re-analyses the integral over the
fields remaining at the witness level $k_0$ on the torus $\mathbb{T}_m$ of that theorem, with the
correlation theorem applied at its parameter value $m=s$, by total covariance on
Ba{\l}aban's history space at level $k_0=K-s$.
The auxiliary object formed in the proof of the torus bound is never formed there.
\end{remark}

\begin{proof}
\emph{Part (a).} We apply Lemma~\ref{8.1:lem-spanning-split}, splitting off the spanning local
terms. By \eqref{8.1:eq-spanning-split-a}, the connected two-point function is the sum of two terms:
\begin{itemize}
\item the spanning sum $T_{\mathcal{E}}$, which collects the mixed derivatives
$\partial_1\partial_2e_j(X;0)$ of the local
exponent terms $e_j(X;z)$ with $j\le k_0$ whose enlarged regions $X^{+}$, in the sense of
Lemma~\ref{8.1:lem-spanning-split}, meet both supports;
\item the mixed derivative $\partial_1\partial_2\log\tilde{Z}(0)$.
\end{itemize}
By the construction in the proof of Lemma~\ref{8.1:lem-spanning-split}, every $z_1z_2$ cross term of
the whole-lattice small-field families and of the whole-lattice families produced by $\mathbf{R}$
lies in a single term $e_j(X;z)$ of the spanning sum.

The same lemma bounds the spanning sum by
\begin{align*}
|T_{\mathcal{E}}|
&\le 4C_e r^{-2}\|f_1\|\|f_2\|\cdot O(1)
\sum_{\nu\ge0}\Bigl(\frac{4L\rho L^{\nu}}{M}+3\Bigr)^{4}e^{-\frac14\kappa\ell_\nu},\\
\ell_\nu&=\Bigl[\frac{\rho L^{\nu}}{4M}-3\Bigr]_{+},
\end{align*}
with the constants $C_e$ and $r$ of that lemma and with $\rho$ the support radius of the $f_i$,
measured in cells of the witness level; the term $\nu=0$ dominates the sum. For
$\rho\ge 12M$ one has $\ell_0=\rho/(4M)-3$, so the summand with $\nu=0$ is
\begin{equation*}
\Bigl(\frac{4L\rho}{M}+3\Bigr)^{4}e^{-\frac{\kappa\rho}{16M}+\frac{3\kappa}{4}},
\end{equation*}
and its exponential factor lies below a fixed negative power of $x_{K-s}$ only when $\rho$ is at
least of order $(M/\kappa)\log x_{K-s}$. In this way the bound
forces the room condition that the support radius $\rho$, measured in cells, be at least of order
$(M/\kappa)\log x_{K-s}$.

\emph{Part (b).} All remaining content of the $n=2$ bound lies in the second term of the
decomposition, $\partial_1\partial_2\log\tilde{Z}(0)$. Every large-field and history contribution
therefore passes through the global quotient. The torus bound controls this quotient uniformly in
$K$, with a volume-dependent constant, and this control carries no factor $g^4$.

\emph{Part (c).} We compose four inputs:
\begin{itemize}
\item the torus bound's formula for the local exponent terms $e_j(X;z)$;
\item the running-shift statement, whose recursion contains, at each step from level $k$ to level
$k+1$, an additive term in which the coupling increment enters at level $k+1$;
\item the bounds of \cite{B-CMP119} on the convergent renormalization expansion;
\item the clause of the corollary to the large-field part of the coupling increment that concerns
the contribution of the healed large-field regions.
\end{itemize}
Together they place the marginal source response of the healed large-field regions inside
$N_{K-s}$. They place the irrelevant response inside the families of effective-action terms from
which the coupling increment $\beta_{k+1}$ has been extracted, and inside the third term of the
correlation theorem's bound at $n=2$.

\emph{The witness theorem.} The two-point witness theorem works on its own torus $\mathbb{T}_m$, with
the correlation theorem applied at its parameter value $m=s$, and treats the integral over the
fields remaining at the witness level $k_0$ directly. It decomposes
the covariance of the two observables by total
covariance with respect to Ba{\l}aban's histories at level $k_0$. The decomposition of its own
two-point function then has two halves:
\begin{itemize}
\item The small-field half is the same spanning sum $T_{\mathcal{E}}$ as in part (a); in the witness
theorem it is the localisation remainder $R^{\mathrm{loc}}$. The witness theorem bounds it through
Lemma~\ref{8.1:lem-spanning-split}, that is, from the torus bound's bounds on the local exponent
terms, and this forces the room condition of part (a).
\item The half carried by the history labels is treated by the large-field block bound, healed part.
This is combined with either the flat shell bound or Ba{\l}aban's summation of the $\mathbf{R}$
operation in \cite{B-CMP122b}.
\end{itemize}
No step of this analysis invokes the torus bound's control of the global quotient $\tilde{Z}$.
\end{proof}

\begin{lemma}[Invariant-theoretic facts for $\mathfrak{su}(N)$]\label{8.1:lem-invariant-facts}
Let $N\ge 2$. On $\mathfrak{su}(N)$ put $\langle X,Y\rangle:=-\operatorname{tr}(XY)$, where
$\operatorname{tr}\mathbf{1}=1$ as in Notation~\ref{8.1:not-curvature-observable}; then
$\langle X,X\rangle=|X|^2$ with $|\cdot|$ the norm of that notation, since the second-order term of
$1-\operatorname{Re}\operatorname{tr}e^{X}$ is $-\tfrac12\operatorname{tr}(X^2)$. Put
$\operatorname{Ad}(\Gamma)X:=\Gamma X\Gamma^{-1}$ for
$\Gamma\in\mathrm{SU}(N)$. For a finite set $I$, $\operatorname{Ad}(\Gamma)$ acts on $\mathfrak{su}(N)^I$
componentwise, and for $B,B'\in\mathfrak{su}(N)^I$ we put
$\langle B,B'\rangle:=\sum_{x\in I}\langle B_x,B'_x\rangle$.
\begin{enumerate}
\item[(i)] There is no non-zero $\operatorname{Ad}$-invariant vector in $\mathfrak{su}(N)$ and no non-zero
$\operatorname{Ad}$-invariant linear functional on $\mathfrak{su}(N)$. Every $\operatorname{Ad}$-invariant
symmetric bilinear form on $\mathfrak{su}(N)$ equals $c\langle\cdot,\cdot\rangle$ for some $c\in\mathbb{R}$.
The assertions on vectors and functionals fail for the Lie algebra $\mathfrak{u}(N)$ of $\mathrm{U}(N)$,
for every $N\ge1$, in particular for $\mathrm{U}(1)$; for $N\ge2$ the classification of invariant
symmetric bilinear forms fails for $\mathfrak{u}(N)$ as well.
\item[(ii)] Let $\mathcal{B}$ be a finite set of oriented bonds, let $\lambda$ range over gauge
transformations, acting by $(U^\lambda)_b=\lambda(x)U_b\lambda(y)^{-1}$ for $b$ from $x$ to $y$, and let
$\Psi$ be a real function of class $C^2$ on $\mathrm{SU}(N)^{\mathcal{B}}$ with $\Psi(U^\lambda)=\Psi(U)$ for
all $\lambda$. For $Y\in\mathfrak{su}(N)^{\mathcal{B}}$ put
$D\Psi(U)[Y]:=\tfrac{d}{dt}\Psi\bigl((e^{tY_b}U_b)_{b}\bigr)\big|_{t=0}$. Then $D\Psi(\mathbf{1}^\lambda)=0$ for
every $\lambda$, i.e.\ at every pure-gauge configuration. If $\psi$ is a real $C^2$ function on
$\mathrm{SU}(N)$ with $\psi(\Gamma U\Gamma^{-1})=\psi(U)$, there is a unique $c_\psi\in\mathbb{R}$ with
\begin{equation*}
\frac{d^2}{dt^2}\psi(e^{tY})\Big|_{t=0}=2c_\psi|Y|^2\qquad(Y\in\mathfrak{su}(N)).
\end{equation*}
For $\psi(V)=1-\operatorname{Re}\operatorname{tr}V$, so that $\psi(U(\partial p))=a_p(U)$, one has
$c_\psi=\tfrac12\neq0$.
\item[(iii)] Let $\mu$ be a probability measure on a measurable space $\mathcal{M}$ on which each
$\Gamma\in\mathrm{SU}(N)$ acts by a measurable map $\omega\mapsto\Gamma\cdot\omega$ preserving $\mu$. Let
$\Psi$ be real with $\Psi(\Gamma\cdot\omega)=\Psi(\omega)$, and $\xi$ be $\mathfrak{su}(N)$-valued with
$\xi(\Gamma\cdot\omega)=\operatorname{Ad}(\Gamma)\xi(\omega)$, and let $\Psi\xi$ be $\mu$-integrable. Then
$\int\Psi\xi\,d\mu=0$.
\item[(iv)] Let $C$ be a symmetric positive semi-definite operator on $\mathfrak{su}(N)^I$ commuting with every
$\operatorname{Ad}(\Gamma)$, let $\mu_C$ be the centred Gaussian measure with covariance $C$ (invariant
under $B\mapsto-B$), with expectation $\mathbb{E}_C$. For $m\ge3$ let $v_m$ be an $\operatorname{Ad}$-invariant
real polynomial on $\mathfrak{su}(N)^I$, homogeneous of degree $m$, and put
$\mathcal{V}:=\sum_{m\ge3}g^{m-2}v_m$, which is invariant under $(B,g)\mapsto(-B,-g)$. For $i=1,2$ let
$\Psi_i=\sum_{m\ge1}g^m\Psi_{i,m}$ with
$\Psi_{i,m}$ $\operatorname{Ad}$-invariant and homogeneous of degree $m$, and write
$\Psi_{i,2}(B)=\tfrac12\langle B,\mathsf{A}_iB\rangle$ with $\mathsf{A}_i$ symmetric. As formal power series
in $g$ put
$\mathbb{E}_g[F]:=\mathbb{E}_C[Fe^{-\mathcal{V}}]/\mathbb{E}_C[e^{-\mathcal{V}}]$ and
\begin{equation*}
\sum_{r\ge0}g^r\varsigma_r:=\mathbb{E}_g[\Psi_1\Psi_2]-\mathbb{E}_g[\Psi_1]\,\mathbb{E}_g[\Psi_2].
\end{equation*}
Then $\varsigma_r=0$ for $r\le3$ and for every odd $r$, and
\begin{equation}\label{8.1:eq-invariant-facts-1}
\varsigma_4=\mathbb{E}_C[\Psi_{1,2}\Psi_{2,2}]-\mathbb{E}_C[\Psi_{1,2}]\,\mathbb{E}_C[\Psi_{2,2}]
=\tfrac12\operatorname{Tr}(\mathsf{A}_1C\mathsf{A}_2C),
\end{equation}
where $\operatorname{Tr}$ denotes the ordinary (unnormalised) trace of linear operators on
$\mathfrak{su}(N)^I$; this is the bubble. Thus $g^2$ and odd orders are absent, the $g^4$ coefficient is
exactly the bubble, and all other terms are of order $g^6$ or higher.
\end{enumerate}
\end{lemma}

Assertion (iv) concerns formal power series in $g$; the weights $g^{m-2}$ of vertices of degree $m$ are
those of the expansion of the effective actions in \cite{B-CMP109}. The non-perturbative counterpart of
(iv) is the Gaussian extraction. In perturbative language, (iii) says that an invariant observable is
joined to the rest by no contribution consisting of a single propagator line, since such a line would
carry one adjoint index.

\begin{proof}
(i) For $X,Y\in\mathfrak{su}(N)$ we have $X^*=-X$, so $\langle X,Y\rangle=\operatorname{tr}(X^*Y)$; this is
real, positive definite, and $\operatorname{Ad}$-invariant because
$\operatorname{tr}(\Gamma X\Gamma^{-1}\Gamma Y\Gamma^{-1})=\operatorname{tr}(XY)$. If
$\operatorname{Ad}(e^{tY})X=X$ for all $t$ and $Y$, differentiation at $t=0$ gives $[Y,X]=0$ for all
$Y\in\mathfrak{su}(N)$. Every complex matrix is $S+iS'$ with $S,S'$ anti-Hermitian, and each of these is
an element of $\mathfrak{su}(N)$ plus a multiple of $\mathbf{1}$; hence $X$ commutes with all matrices,
so $X=c\mathbf{1}$, and $\operatorname{tr}X=0$ gives $X=0$. An invariant functional is
$\langle X_0,\cdot\rangle$ for a unique $X_0$, and invariance of $\langle\cdot,\cdot\rangle$ makes $X_0$
invariant, hence $X_0=0$.

Let $\beta$ be an invariant symmetric bilinear form and write $\beta(X,Y)=\langle TX,Y\rangle$ with $T$
symmetric; invariance gives $T\operatorname{Ad}(\Gamma)=\operatorname{Ad}(\Gamma)T$. Let $c$ be an
eigenvalue of $T$ (real) and $\mathfrak{k}:=\ker(T-c)\neq0$; $\mathfrak{k}$ is $\operatorname{Ad}$-invariant,
so differentiating, $[Y,\mathfrak{k}]\subset\mathfrak{k}$ for all $Y\in\mathfrak{su}(N)$. Since
$\mathfrak{sl}(N,\mathbb{C})=\mathfrak{su}(N)\oplus i\,\mathfrak{su}(N)$ over $\mathbb{R}$,
$\mathfrak{k}_{\mathbb{C}}:=\mathfrak{k}+i\mathfrak{k}$ is an ideal of $\mathfrak{sl}(N,\mathbb{C})$ with
$\dim_{\mathbb{R}}\mathfrak{k}_{\mathbb{C}}=2\dim\mathfrak{k}$. Take real $\delta_1,\dots,\delta_N$ with sum
zero and all $\delta_a-\delta_b$ ($a\neq b$) distinct and non-zero, and $D:=\operatorname{diag}(\delta_a)$.
Then $\operatorname{ad}_D$ is diagonalisable on $\mathfrak{sl}(N,\mathbb{C})$, with kernel the diagonal
matrices and eigenvectors the matrix units $e_{ab}$, $a\neq b$, for the eigenvalues $\delta_a-\delta_b$; the
invariant subspace $\mathfrak{k}_{\mathbb{C}}$ contains the eigencomponents of each of its elements. Let
$0\neq X'\in\mathfrak{k}_{\mathbb{C}}$ have entries $x'_{ab}$. If $x'_{ab}\neq0$ for some $a\neq b$, then
$e_{ab}$ lies in $\mathfrak{k}_{\mathbb{C}}$; if $X'$ is diagonal, some $x'_{aa}\neq x'_{bb}$ and
$[X',e_{ab}]=(x'_{aa}-x'_{bb})e_{ab}\in\mathfrak{k}_{\mathbb{C}}$. From one $e_{ab}\in\mathfrak{k}_{\mathbb{C}}$
the relations $[e_{a'a},e_{ab}]=e_{a'b}$ ($a'\notin\{a,b\}$), $[e_{ab},e_{bb'}]=e_{ab'}$
($b'\notin\{a,b\}$), $[e_{ab},e_{ba}]=e_{aa}-e_{bb}$ and $[e_{ba},e_{aa}-e_{bb}]=2e_{ba}$, applied
repeatedly, give all of $\mathfrak{sl}(N,\mathbb{C})$. Hence
$2\dim\mathfrak{k}=2(N^2-1)$, $\mathfrak{k}=\mathfrak{su}(N)$, $T=c$ and $\beta=c\langle\cdot,\cdot\rangle$.

In $\mathfrak{u}(N)$ the vector $i\mathbf{1}$ is invariant and $X\mapsto i\operatorname{tr}X$ is a real,
invariant functional taking the value $-1$ at $i\mathbf{1}$. For $N\ge2$, the form
$(X,Y)\mapsto\operatorname{tr}X\,\operatorname{tr}Y$ on $\mathfrak{u}(N)$ is real (both traces are
imaginary), symmetric and invariant; it vanishes on $\mathfrak{su}(N)\times\mathfrak{su}(N)$, where
$-\operatorname{tr}(XY)$ is positive definite and $\mathfrak{su}(N)\neq0$, and it takes the value $-1$ at
$(i\mathbf{1},i\mathbf{1})$, where $-\operatorname{tr}(XY)=1$; so it is not a multiple of
$-\operatorname{tr}(XY)$.

(ii) Let $b$ run from $x$ to $y$. Apply $\Psi(U)=\Psi(U^{\lambda^{-1}})$, where
$\lambda^{-1}(x):=\lambda(x)^{-1}$, to $U_b=e^{tY_b}\lambda(x)\lambda(y)^{-1}$; the $b$-component of
$U^{\lambda^{-1}}$ is $\lambda(x)^{-1}e^{tY_b}\lambda(x)=\exp(t\operatorname{Ad}(\lambda(x)^{-1})Y_b)$.
Hence $D\Psi(\mathbf{1}^\lambda)[Y]=D\Psi(\mathbf{1})[Y']$ with
$Y'_b=\operatorname{Ad}(\lambda(x)^{-1})Y_b$. With the constant transformation $\lambda\equiv\Gamma$,
$(e^{tY_b})^\Gamma=e^{t\operatorname{Ad}(\Gamma)Y_b}$, so the linear functional $D\Psi(\mathbf{1})$
satisfies $D\Psi(\mathbf{1})[\operatorname{Ad}(\Gamma)Y]=D\Psi(\mathbf{1})[Y]$. Write
$D\Psi(\mathbf{1})[Y]=\sum_b\ell_b(Y_b)$; taking $Y$ supported on a single bond shows each $\ell_b$
invariant, hence zero by (i). So $D\Psi(\mathbf{1}^\lambda)=0$.

For $\psi$ as stated, $\psi\circ\exp$ is $C^2$ near $0$ and
$\tfrac{d^2}{dt^2}\psi(e^{tY})|_{t=0}=D^2(\psi\circ\exp)(0)[Y,Y]$, a quadratic form; its polarisation is
invariant because $\psi(e^{t\operatorname{Ad}(\Gamma)Y})=\psi(\Gamma e^{tY}\Gamma^{-1})=\psi(e^{tY})$. By
(i) it equals $2c_\psi\langle\cdot,\cdot\rangle$, and $c_\psi$ is unique since $|Y|^2>0$ for $Y\neq0$.
For $\psi(V)=1-\operatorname{Re}\operatorname{tr}V$ the second derivative at $t=0$ is
$-\operatorname{Re}\operatorname{tr}(Y^2)=|Y|^2$, as $\operatorname{tr}(Y^2)=-|Y|^2$ is real; so
$c_\psi=\tfrac12$.

(iii) For each $\Gamma$, since $\operatorname{Ad}(\Gamma)$ is linear and $\mu$ is invariant,
\begin{equation*}
\operatorname{Ad}(\Gamma)\int\Psi\xi\,d\mu=\int\Psi(\Gamma\cdot\omega)\,\xi(\Gamma\cdot\omega)\,d\mu(\omega)
=\int\Psi\xi\,d\mu .
\end{equation*}
So $\int\Psi\xi\,d\mu$ is an invariant vector, hence zero by (i).

(iv) Both expectations are well defined: $\mathcal{V}=O(g)$, so the coefficient of $g^r$ in
$e^{-\mathcal{V}}=\sum_l(-\mathcal{V})^l/l!$ involves $l\le r$ only and is a polynomial, and
$\mathbb{E}_C[e^{-\mathcal{V}}]$ has constant term $1$, hence is invertible. Call a series
$\sum_rg^rp_r(B)$ balanced if every homogeneous part of $p_r$ has degree $\equiv r \pmod 2$. Products of
balanced series are balanced; $\mathcal{V}$, $e^{-\mathcal{V}}$ and $\Psi_i$ are balanced, since
$m-2\equiv m$. For balanced $F$ and odd $r$, the $g^r$ coefficient of $\mathbb{E}_C[F]$ is a sum of
expectations of odd homogeneous polynomials, which vanish by invariance of $\mu_C$ under $B\mapsto-B$.
So $\mathbb{E}_C[e^{-\mathcal{V}}]$, $\mathbb{E}_C[\Psi_ie^{-\mathcal{V}}]$ and
$\mathbb{E}_C[\Psi_1\Psi_2e^{-\mathcal{V}}]$ are even series; the inverse of an even series
$1+\phi$ is $\sum_l(-\phi)^l$, also even. Hence $\varsigma_r=0$ for odd $r$.

By the argument of (ii), $\Psi_{i,1}(B)=\sum_x\ell_x(B_x)$ with each $\ell_x$ invariant, so
$\Psi_{i,1}=0$ by (i), and $\Psi_i=g^2\Psi_{i,2}+O(g^3)$. Then
$\Psi_1\Psi_2e^{-\mathcal{V}}=g^4\Psi_{1,2}\Psi_{2,2}+O(g^5)$ and $\mathbb{E}_C[e^{-\mathcal{V}}]=1+O(g^2)$;
by evenness, $\mathbb{E}_g[\Psi_1\Psi_2]=g^4\mathbb{E}_C[\Psi_{1,2}\Psi_{2,2}]+O(g^6)$ and
$\mathbb{E}_g[\Psi_i]=g^2\mathbb{E}_C[\Psi_{i,2}]+O(g^4)$. This gives $\varsigma_r=0$ for $r\le3$, the first
equality in \eqref{8.1:eq-invariant-facts-1}, and order $\ge6$ for the rest. In an orthonormal basis,
with $\varpi_a$ the coordinates of $B$, Wick's formula
\begin{equation*}
\mathbb{E}_C[\varpi_a\varpi_b\varpi_c\varpi_d]=C_{ab}C_{cd}+C_{ac}C_{bd}+C_{ad}C_{bc}
\end{equation*}
gives
\begin{equation*}
\varsigma_4=\tfrac14\sum_{a,b,c,d}(\mathsf{A}_1)_{ab}(\mathsf{A}_2)_{cd}\,(C_{ac}C_{bd}+C_{ad}C_{bc})
=\tfrac12\operatorname{Tr}(\mathsf{A}_1C\mathsf{A}_2C),
\end{equation*}
the term $C_{ab}C_{cd}$ cancelling against the product of the means, and symmetry of $\mathsf{A}_i$, $C$
being used in the last step.
\end{proof}

\begin{lemma}[Covariance of two sums of Gaussian squares]\label{8.1:lem-gaussian-squares}
Let $I$ and $J$ be finite index sets, and let $(\xi_\alpha)_{\alpha\in I}$ and
$(\eta_\beta)_{\beta\in J}$ be real random variables such that the joint family
$(\xi_\alpha,\eta_\beta)_{\alpha\in I,\,\beta\in J}$ is Gaussian and centred. Write
$\mathbb{E}$ for the expectation and
$\operatorname{Cov}(X,Y)=\mathbb{E}[XY]-\mathbb{E}[X]\,\mathbb{E}[Y]$. Then
\begin{equation}\label{8.1:eq-gaussian-squares-1}
\operatorname{Cov}\Bigl(\sum_{\alpha\in I}\xi_\alpha^2,\ \sum_{\beta\in J}\eta_\beta^2\Bigr)
=2\sum_{\alpha\in I}\sum_{\beta\in J}\bigl(\mathbb{E}[\xi_\alpha\eta_\beta]\bigr)^2\ \ge\ 0 .
\end{equation}
\end{lemma}

\begin{proof}
Gaussian random variables have finite moments of all orders, so every expectation below is
finite. For real random variables $a,b,c,d$ that are jointly Gaussian and centred, the
fourth-moment case of the Gaussian pairing formula reads
\begin{equation*}
\mathbb{E}[abcd]=\mathbb{E}[ab]\,\mathbb{E}[cd]+\mathbb{E}[ac]\,\mathbb{E}[bd]
+\mathbb{E}[ad]\,\mathbb{E}[bc].
\end{equation*}
Fix $\alpha\in I$ and $\beta\in J$. We apply this formula with $a=b=\xi_\alpha$ and
$c=d=\eta_\beta$, which are jointly Gaussian and centred by hypothesis. The second and third
terms on the right are then both equal to $(\mathbb{E}[\xi_\alpha\eta_\beta])^2$, so
\begin{equation*}
\mathbb{E}[\xi_\alpha^2\eta_\beta^2]
=\mathbb{E}[\xi_\alpha^2]\,\mathbb{E}[\eta_\beta^2]+2\bigl(\mathbb{E}[\xi_\alpha\eta_\beta]\bigr)^2 .
\end{equation*}
By the definition of the covariance, this is the identity
$\operatorname{Cov}(\xi_\alpha^2,\eta_\beta^2)=2(\mathbb{E}[\xi_\alpha\eta_\beta])^2$.
The covariance is bilinear on random variables with finite second moments. The sums in
\eqref{8.1:eq-gaussian-squares-1} are finite and each $\xi_\alpha^2$, $\eta_\beta^2$ has finite
second moment. Hence
\begin{equation*}
\operatorname{Cov}\Bigl(\sum_{\alpha\in I}\xi_\alpha^2,\ \sum_{\beta\in J}\eta_\beta^2\Bigr)
=\sum_{\alpha\in I}\sum_{\beta\in J}\operatorname{Cov}(\xi_\alpha^2,\eta_\beta^2)
=2\sum_{\alpha\in I}\sum_{\beta\in J}\bigl(\mathbb{E}[\xi_\alpha\eta_\beta]\bigr)^2 .
\end{equation*}
Each summand on the right is the square of a real number and is therefore non-negative. This
gives the inequality in \eqref{8.1:eq-gaussian-squares-1}.
\end{proof}

Throughout the lemma below, $\mathbb{F}$ is the bare torus lattice of Notation~\ref{8.1:not-curvature-observable}. We write it in lattice units, so that its bonds have length one. It is a discrete torus with the same number of points in each of the four coordinate directions; $|\mathbb{F}|$ is its number of points and $e_1,\dots,e_4$ are the unit vectors.

For $f:\mathbb{F}\to\mathbb{R}$ we put
\begin{gather*}
\nabla^{+}_{\mu}f(z)=f(z+e_\mu)-f(z),\qquad \nabla^{-}_{\mu}f(z)=f(z)-f(z-e_\mu),\\
\nabla^2_\mu=\nabla^+_\mu\nabla^-_\mu,\qquad \Delta=\sum_{\mu=1}^{4}\nabla^2_\mu ,
\end{gather*}
and we write $(G*f)(z)=\sum_{z'\in\mathbb{F}}G(z-z')f(z')$ for convolution. We also put $\langle f,g\rangle=\sum_{z\in\mathbb{F}}f(z)g(z)$.

A $1$-cochain is a family $\xi=(\xi_\alpha)_{\alpha=1}^4$ of functions on $\mathbb{F}$, with $\langle \xi,\xi'\rangle=\sum_\alpha\langle \xi_\alpha,\xi'_\alpha\rangle$. A $2$-cochain is a function $\omega(x;\mu\nu)$ of $x\in\mathbb{F}$ and an ordered pair $\mu\neq\nu$, subject to $\omega(x;\nu\mu)=-\omega(x;\mu\nu)$. An element of $\Lambda^2(\mathbb{F})$ is written $p=(x;\mu\nu)$ with $\mu<\nu$, and $\langle\omega,\omega'\rangle=\sum_{x}\sum_{\mu<\nu}\omega(x;\mu\nu)\omega'(x;\mu\nu)$.

The coboundaries are
\begin{equation*}
(d_0\varphi)_\alpha=\nabla^+_\alpha\varphi,\qquad
(d_1\xi)(x;\mu\nu)=\nabla^+_\mu \xi_\nu(x)-\nabla^+_\nu \xi_\mu(x),
\end{equation*}
and we set $\Delta_1=d_1^*d_1+d_0d_0^*$. A superscript $+$ denotes the Moore--Penrose pseudoinverse. For an ordered pair $\mu\ne\nu$, $\delta_{(x;\mu\nu)}$ is the $2$-cochain with $\langle\delta_{(x;\mu\nu)},\omega\rangle=\omega(x;\mu\nu)$; thus $\delta_{(x;\nu\mu)}=-\delta_{(x;\mu\nu)}$. Finally, $\delta_x$ is the indicator of $\{x\}$. For an operator $P$ on $2$-cochains we write $P(q,q')=\langle\delta_q,P\delta_{q'}\rangle$. Within the lemma below and its proof, $\mu,\nu,\rho,\sigma,\alpha,a,b$ denote direction indices in $\{1,\dots,4\}$, and $q,q'$ denote ordered plaquettes $(x;\mu\nu)$ with $\mu\neq\nu$.

\begin{lemma}[The free lattice curvature covariance]\label{8.1:lem-free-curvature-covariance}
Let $n=N^2-1$. Let $\gamma_{\mathrm{ex}}$ be the standard Gaussian law on the orthogonal sum of $n$ copies of $\operatorname{im}d_1\subset\mathbb{R}^{\Lambda^2(\mathbb{F})}$, one copy for each colour; this is the free lattice Maxwell curvature in lattice units. Then under $\gamma_{\mathrm{ex}}$ distinct colour components are uncorrelated, and each colour component has covariance
\begin{equation*}
P_{\mathrm{ex}}=d_1\Delta_1^{+}d_1^{*}.
\end{equation*}
Put $G=(-\Delta)^+\delta_0$. For $x,y\in\mathbb{F}$ and ordered pairs $\mu\neq\nu$, $\rho\neq\sigma$,
\begin{equation}\label{8.1:eq-free-curvature-covariance-a}
\begin{aligned}
P_{\mathrm{ex}}\bigl((x;\mu\nu),(y;\rho\sigma)\bigr)
 &=\sum_{\alpha\in\{\mu,\nu\}\cap\{\rho,\sigma\}}
 \bigl\langle (d_1^*\delta_{(x;\mu\nu)})_\alpha,\,
 G*(d_1^*\delta_{(y;\rho\sigma)})_\alpha\bigr\rangle,\\
P_{\mathrm{ex}}\bigl((x;\mu\nu),(y;\mu\nu)\bigr)
 &=-\bigl(\nabla^2_\mu+\nabla^2_\nu\bigr)G(x-y),\\
P_{\mathrm{ex}}\bigl((x;\mu\nu),(y;\mu\sigma)\bigr)
 &=-\nabla^+_\nu\nabla^-_\sigma G(x-y)\qquad(\sigma\notin\{\mu,\nu\}).
\end{aligned}
\end{equation}
In particular the entry vanishes when $\{\mu,\nu\}\cap\{\rho,\sigma\}=\emptyset$. In $d=4$ the diagonal entries are
\begin{equation*}
P_{\mathrm{ex}}(p,p)=\tfrac12-\tfrac12|\mathbb{F}|^{-1}\qquad\text{for every } p\in\Lambda^2(\mathbb{F});
\end{equation*}
that is, the diagonal value is $\tfrac12$ up to the volume correction $-\tfrac12|\mathbb{F}|^{-1}$.
Every entry of $P_{\mathrm{ex}}$ is a linear combination of at most sixteen values of $G$.
\end{lemma}

\begin{proof}
\emph{Covariance.} The standard Gaussian law on a subspace $V$ of a Euclidean space has covariance $\langle u,\Pi_Vv\rangle$, where $\Pi_V$ is the orthogonal projection onto $V$. For the orthogonal sum of $n$ copies of $\operatorname{im}d_1$ this projection is the direct sum of $n$ copies of $\Pi_{\operatorname{im}d_1}$. Hence distinct colours are uncorrelated, and it suffices to show $P_{\mathrm{ex}}=\Pi_{\operatorname{im}d_1}$.

Since difference operators commute,
\begin{equation*}
(d_1d_0\varphi)(x;\mu\nu)=\nabla^+_\mu\nabla^+_\nu\varphi(x)-\nabla^+_\nu\nabla^+_\mu\varphi(x)=0.
\end{equation*}
Hence $d_0^*d_1^*=(d_1d_0)^*=0$, and so $\Delta_1=d_1^*d_1$ on $\operatorname{im}d_1^*$.

The operator $d_1^*d_1$ maps $\operatorname{im}d_1^*$ into itself and is injective there. Indeed, $d_1^*d_1u=0$ gives $\|d_1u\|^2=0$, hence $u\in\ker d_1\cap\operatorname{im}d_1^*=\{0\}$. Being injective on a finite-dimensional space, it is bijective on $\operatorname{im}d_1^*$. Therefore $\operatorname{im}d_1^*\subset\operatorname{im}\Delta_1$, and for $v\in\operatorname{im}d_1^*$ there is a unique $u\in\operatorname{im}d_1^*$ with $\Delta_1u=v$. Since $\Delta_1$ is self-adjoint, $\operatorname{im}\Delta_1=(\ker\Delta_1)^\perp$, and the solution $u\in\operatorname{im}d_1^*\subset\operatorname{im}\Delta_1=(\ker\Delta_1)^\perp$ of $\Delta_1u=v$ is $\Delta_1^+v$.

Write $\xi=\xi'+\xi''$ with $\xi'\in\operatorname{im}d_1^*$ and $\xi''\in\ker d_1$. Then $\Delta_1^+d_1^*d_1\xi=\Delta_1^+d_1^*d_1\xi'=\xi'$, so $P_{\mathrm{ex}}d_1\xi=d_1\xi'=d_1\xi$. Moreover $P_{\mathrm{ex}}$ vanishes on $\ker d_1^*=(\operatorname{im}d_1)^\perp$. Thus $P_{\mathrm{ex}}=\Pi_{\operatorname{im}d_1}$.

\emph{The Laplacians.} Summation by parts gives $\langle\nabla^+_\mu f,g\rangle=-\langle f,\nabla^-_\mu g\rangle$. Hence $(d_0^*\xi)=-\sum_\alpha\nabla^-_\alpha \xi_\alpha$. Extending the sum over $\mu<\nu$ antisymmetrically, we get $\langle\omega,d_1\xi\rangle=\sum_x\sum_{\mu\neq\nu}\omega(x;\mu\nu)\nabla^+_\mu \xi_\nu(x)$, so that $(d_1^*\omega)_\nu=-\sum_{\mu\neq\nu}\nabla^-_\mu\omega(\cdot;\mu\nu)$. Therefore
\begin{equation*}
(\Delta_1\xi)_\nu=-\Delta \xi_\nu+\sum_\mu\bigl(\nabla^-_\mu\nabla^+_\nu-\nabla^+_\nu\nabla^-_\mu\bigr)\xi_\mu
=-\Delta \xi_\nu .
\end{equation*}

Since $\langle f,-\Delta f\rangle=\sum_\mu\|\nabla^+_\mu f\|^2$ and the torus is connected, $\ker(-\Delta)$ consists of the constants. Consequently $\operatorname{im}(-\Delta)$ is the space of functions of mean zero. It follows that $G$ has mean zero and that $(-\Delta)G=\delta_0-|\mathbb{F}|^{-1}$.

Now let $v$ have mean zero. Then $G*v$ has mean zero and
\begin{equation*}
(-\Delta)(G*v)=v-|\mathbb{F}|^{-1}\textstyle\sum_{z'} v(z')=v,
\end{equation*}
so $(-\Delta)^+v=G*v$. Componentwise, therefore, $\Delta_1^+v=(G*v_\alpha)_\alpha$ whenever every component $v_\alpha$ has mean zero.

\emph{Proof of \eqref{8.1:eq-free-curvature-covariance-a}.} From $\langle d_1^*\delta_{(x;\mu\nu)},\xi\rangle=\nabla^+_\mu \xi_\nu(x)-\nabla^+_\nu \xi_\mu(x)$ and summation by parts we read off
\begin{equation*}
(d_1^*\delta_{(x;\mu\nu)})_\nu=-\nabla^-_\mu\delta_x,\qquad (d_1^*\delta_{(x;\mu\nu)})_\mu=\nabla^-_\nu\delta_x,
\end{equation*}
and all other components vanish. Each component has mean zero. By the previous paragraph,
\begin{equation*}
P_{\mathrm{ex}}(q,q')=\langle d_1^*\delta_q,\Delta_1^+d_1^*\delta_{q'}\rangle=\sum_\alpha\langle(d_1^*\delta_q)_\alpha,G*(d_1^*\delta_{q'})_\alpha\rangle .
\end{equation*}
Only $\alpha\in\{\mu,\nu\}\cap\{\rho,\sigma\}$ contribute, which is the first line of \eqref{8.1:eq-free-curvature-covariance-a}; the sum is empty for disjoint index pairs.

For directions $a,b$ we have $G*\nabla^-_b\delta_y=\nabla^-_b G(\cdot-y)$. Summation by parts then gives
\begin{equation*}
\langle\nabla^-_a\delta_x,G*\nabla^-_b\delta_y\rangle=-(\nabla^+_a\nabla^-_bG)(x-y).
\end{equation*}
For $q=(x;\mu\nu)$ and $q'=(y;\mu\nu)$, the term $\alpha=\nu$ pairs $-\nabla^-_\mu\delta_x$ with $-\nabla^-_\mu\delta_y$ and gives $-\nabla^2_\mu G(x-y)$. The term $\alpha=\mu$ pairs $\nabla^-_\nu\delta_x$ with $\nabla^-_\nu\delta_y$ and gives $-\nabla^2_\nu G(x-y)$. This proves the second line.

For $q'=(y;\mu\sigma)$ with $\sigma\notin\{\mu,\nu\}$, only $\alpha=\mu$ contributes. It pairs $\nabla^-_\nu\delta_x$ with $\nabla^-_\sigma\delta_y$ and gives $-\nabla^+_\nu\nabla^-_\sigma G(x-y)$, which is the third line.

\emph{Diagonal.} Both the equation $(-\Delta)G=\delta_0-|\mathbb{F}|^{-1}$ and the condition of mean zero are preserved by $z\mapsto-z$ and by permutations of the coordinates. Since these two conditions determine $G$ uniquely, $G$ is invariant under both operations. Hence
\begin{equation*}
\nabla^2_\mu G(0)=2\bigl(G(e_\mu)-G(0)\bigr)
\end{equation*}
does not depend on $\mu$. Consequently
\begin{equation*}
-\nabla^2_\mu G(0)=\tfrac14(-\Delta)G(0)=\tfrac14\bigl(1-|\mathbb{F}|^{-1}\bigr).
\end{equation*}
The second line of \eqref{8.1:eq-free-curvature-covariance-a} at $x=y$ therefore gives $P_{\mathrm{ex}}(p,p)=\tfrac12-\tfrac12|\mathbb{F}|^{-1}$, that is, $\tfrac12$ up to the volume correction $-\tfrac12|\mathbb{F}|^{-1}$.

\emph{Count.} In the first line of \eqref{8.1:eq-free-curvature-covariance-a} there are at most two indices $\alpha$. Each term has the form $-(\nabla^+_a\nabla^-_bG)(x-y)$, which involves the four values $G(x-y+e_a)$, $G(x-y)$, $G(x-y+e_a-e_b)$ and $G(x-y-e_b)$. Every entry is therefore a combination of at most eight, and hence at most sixteen, values of $G$.
\end{proof}

\begin{definition}[The background curvature and the two-point kernel]\label{8.1:def-background-kernel}
Let $\mathbb{F}$ be the bare lattice of spacing $L^{-K}$ on $\mathbb{T}_m$ of
Definition~\ref{1.1:def-lattice}, and $\Lambda^2(\mathbb{F})$ its set of plaquettes. Let $s$
be the depth below the cutoff and $k_0=K-s$ the corresponding level.
\begin{enumerate}
\item Let $\mathcal{A}$ be the centred massless lattice Gaussian $1$-cochain on $\mathbb{F}$, in
the sense of \cite[(1.3), (1.5)]{B-CMP95}. We write $\mathbb{E}$ for expectation with respect to
the law of $\mathcal{A}$.
\item The \emph{linearised block field} is $B:=Q_{k_0}\mathcal{A}$, where $Q_{k_0}$ is the linearised
block-averaging map of \cite[(1.18)]{B-CMP95}, taken at level $k_0$.
\item The \emph{background field} is the conditional expectation
$A_B:=\mathbb{E}[\mathcal{A}\mid B]$. It coincides with the minimiser of the linear
variational problem of \cite[(1.19)]{B-CMP95} with block datum $B$.
\item The \emph{background curvature} is $F_B:=d_1A_B$, a $2$-cochain on $\mathbb{F}$; here $d_1$
denotes the lattice coboundary from $1$-cochains to $2$-cochains, as in
Lemma~\ref{8.1:lem-free-curvature-covariance}.
\item The \emph{background covariance} and the \emph{fluctuation covariance} are the functions
on pairs of plaquettes $p,q\in\Lambda^2(\mathbb{F})$ given by
\begin{equation}\label{8.1:eq-background-kernel-1}
\begin{aligned}
C^{\mathrm{bg}}_{K,s}(p,q)&:=\mathbb{E}\bigl[F_B(p)\,F_B(q)\bigr],\\
C^{\mathrm{fl}}(p,q)&:=\mathbb{E}\bigl[\operatorname{Cov}\bigl(d_1\mathcal{A}(p),
d_1\mathcal{A}(q)\,\big|\,B\bigr)\bigr].
\end{aligned}
\end{equation}
\item For test functions $f_1,f_2$ the \emph{two-point kernel} is
\begin{equation}\label{8.1:eq-background-kernel-2}
\mathcal{K}_{K,s}(f_1,f_2):=\sum_{p}\sum_{q}f_1(x_p)\,f_2(x_q)\,
\bigl[C^{\mathrm{bg}}_{K,s}(p,q)\bigr]^2 .
\end{equation}
Here both sums run over $\Lambda^2(\mathbb{F})$. Every plaquette $p$ is written $p=(x;\mu\nu)$ with
$1\le\mu<\nu\le4$, where $x$ is the site of $\mathbb{F}$ from which $p$ is spanned by the positive
lattice directions $\mu$ and $\nu$; this site is the \emph{attachment point} $x_p:=x$ of $p$.
\item For sites $x,y$ of $\mathbb{F}$ the \emph{point kernel} is
\begin{equation}\label{8.1:eq-background-kernel-3}
k(x,y):=\sum_{p:\,x_p=x}\ \sum_{q:\,x_q=y}\bigl[C^{\mathrm{bg}}_{K,s}(p,q)\bigr]^2 .
\end{equation}
Here $p$ runs over the six plaquettes attached at $x$ and $q$ over the six attached at $y$.
\item With $n=N^2-1$ we put $b_0:=n/2$.
\end{enumerate}
\end{definition}

\emph{Decomposition of the free covariance.} The covariance of the curvature $d_1\mathcal{A}$
under the law of $\mathcal{A}$ is the free lattice curvature covariance $P_{\mathrm{ex}}$ of
Lemma~\ref{8.1:lem-free-curvature-covariance}, whose entries are given in
\eqref{8.1:eq-free-curvature-covariance-a}. It splits as
\begin{equation}\label{8.1:eq-background-kernel-4}
P_{\mathrm{ex}}(p,q)=C^{\mathrm{bg}}_{K,s}(p,q)+C^{\mathrm{fl}}(p,q),
\qquad p,q\in\Lambda^2(\mathbb{F}).
\end{equation}
Indeed, $d_1$ is linear and conditional expectation is linear, so
\begin{equation*}
\mathbb{E}[d_1\mathcal{A}\mid B]=d_1\,\mathbb{E}[\mathcal{A}\mid B]=d_1A_B=F_B .
\end{equation*}
Since $\mathcal{A}$ is centred, so are $d_1\mathcal{A}$ and $F_B$. The law of total covariance
then reads
\begin{equation*}
\operatorname{Cov}\bigl(d_1\mathcal{A}(p),d_1\mathcal{A}(q)\bigr)
=\mathbb{E}\bigl[\operatorname{Cov}\bigl(d_1\mathcal{A}(p),d_1\mathcal{A}(q)\,\big|\,B\bigr)\bigr]
+\mathbb{E}\bigl[F_B(p)\,F_B(q)\bigr].
\end{equation*}
By \eqref{8.1:eq-background-kernel-1}, this is \eqref{8.1:eq-background-kernel-4}.

\emph{The free kernel.} The kernel built from the full free covariance is
\begin{equation}\label{8.1:eq-background-kernel-5}
\mathcal{K}^{\mathrm{ex}}(f_1,f_2):=\sum_{p}\sum_{q}f_1(x_p)\,f_2(x_q)\,
\bigl[P_{\mathrm{ex}}(p,q)\bigr]^2 .
\end{equation}
Let $C_W=C_W(s)$ be the room function fixed in Notation~\ref{8.1:not-test-pairs-kernel}.
For the pairs $(f_1,f_2)$ of test functions (pairs of the class $\mathfrak{T}(\mathsf{d};\vartheta)$
of that notation) and the depths $s$ admitted by the two-point witness theorem, the comparison
of the free and background kernels, which is proved as an implication from the two hypotheses
named in its statement, gives
\begin{equation}\label{8.1:eq-background-kernel-6}
\bigl|\mathcal{K}^{\mathrm{ex}}(f_1,f_2)-\mathcal{K}_{K,s}(f_1,f_2)\bigr|
\le C\,C_W^{4}\,e^{-\delta C_W}\,\mathcal{K}_{K,s}(f_1,f_2),
\end{equation}
with the constants $C$ and $\delta>0$ of that statement. The right-hand side of
\eqref{8.1:eq-background-kernel-6} is of the size absorbed in the remainder $R_K$ of the
two-point witness theorem; the kernel of that theorem is
$\mathcal{K}_{K,s}$ of \eqref{8.1:eq-background-kernel-2}, not $\mathcal{K}^{\mathrm{ex}}$.

\begin{lemma}[The Gaussian two-point covariance and positivity of the kernel]
\label{8.1:lem-kernel-positivity}
Let $n=N^2-1$. Write $F^a(p)$, $a=1,\dots,n$, for the colour components at the plaquette $p$ of
$\mathbb{F}$ of a curvature field $F$, and let $F=(F^a(p))_{a,p}$ be a centred, jointly Gaussian
family with
\begin{equation}\label{8.1:eq-kernel-positivity-1}
\mathbb{E}\bigl[F^a(p)F^b(q)\bigr]=\delta^{ab}\,C(p,q)
\end{equation}
for a real function $C$ on pairs of plaquettes. For a real function $f$ on the plaquettes of
$\mathbb{F}$ put
\begin{equation*}
X^{(2)}(f):=\tfrac12\sum_{a=1}^{n}\sum_{p}f(p)\bigl(F^a(p)\bigr)^2 .
\end{equation*}
For real functions $f_1,f_2$ on the plaquettes we write $\mathcal{K}^{\mathrm{ex}}(f_1,f_2)$ and
$\mathcal{K}_{K,s}(f_1,f_2)$ for the kernels of Definition~\ref{8.1:def-background-kernel} with
$f_i(x(p))$ replaced by $f_i(p)$; in particular
\begin{equation*}
\mathcal{K}^{\mathrm{ex}}(f_1,f_2)=\sum_{p}\sum_{q}f_1(p)f_2(q)\,P_{\mathrm{ex}}(p,q)^2 .
\end{equation*}
For test functions on the torus the two readings agree with $f_i(p):=f_i(x(p))$.
\begin{itemize}
\item[(a)] For all real $f_1,f_2$,
\begin{equation}\label{8.1:eq-kernel-positivity-a}
\operatorname{Cov}\bigl(X^{(2)}(f_1),X^{(2)}(f_2)\bigr)
=b_0\sum_{p}\sum_{q}f_1(p)f_2(q)\,C(p,q)^2,
\end{equation}
where
\begin{equation*}
b_0=\bigl(\tfrac12\bigr)^2\cdot 2\cdot n=\frac{N^2-1}{2}\neq0 .
\end{equation*}
In particular, the covariance equals $b_0\,\mathcal{K}^{\mathrm{ex}}(f_1,f_2)$ when $F$ is
distributed according to $\gamma_{\mathrm{ex}}$, so that $C=P_{\mathrm{ex}}$. It equals
$b_0\,\mathcal{K}_{K,s}(f_1,f_2)$ when $F=F_B$ is the background curvature, so that
$C=C^{\mathrm{bg}}_{K,s}$.
\item[(b)] For every real function $f$ on the plaquettes with $f\not\equiv0$,
\begin{equation}\label{8.1:eq-kernel-positivity-b}
\mathcal{K}^{\mathrm{ex}}(f,f)>0 .
\end{equation}
\item[(c)] If $f_1\ge0$ and $f_2\ge0$, every term of the double sums defining
$\mathcal{K}^{\mathrm{ex}}(f_1,f_2)$ and $\mathcal{K}_{K,s}(f_1,f_2)$ is non-negative.
\item[(d)] With $\delta_p$ the $2$-cochain equal to $1$ on $p$, to $-1$ on $p$ with reversed
orientation, and to $0$ on every other plaquette, $d^*$ the adjoint
of $d_1$, $\langle\cdot,\cdot\rangle$ the standard inner product of $1$-cochains on $\mathbb{F}$, and
$(-\Delta)^+$ the pseudo-inverse of the lattice Laplacian acting on each component of a
$1$-cochain,
\begin{equation}\label{8.1:eq-kernel-positivity-2}
b_0\,\mathcal{K}^{\mathrm{ex}}(f_1,f_2)=\frac{N^2-1}{2}\sum_{p}\sum_{q}f_1(p)f_2(q)
\bigl[\langle d^*\delta_p,(-\Delta)^+d^*\delta_q\rangle\bigr]^2 .
\end{equation}
Each bracket is a sum of at most two second differences of $G$. The right side therefore
involves only $f_1$, $f_2$ and finitely many values of the Green's function
$G=(-\Delta)^+\delta_0$ of Lemma~\ref{8.1:lem-free-curvature-covariance}, and no quantity of the
renormalisation-group construction.
\end{itemize}
\end{lemma}

\begin{proof}
\emph{(a).} By bilinearity of the covariance,
\begin{equation*}
\operatorname{Cov}\bigl(X^{(2)}(f_1),X^{(2)}(f_2)\bigr)
=\frac14\sum_{a,b=1}^{n}\sum_{p,q}f_1(p)f_2(q)\,
\operatorname{Cov}\bigl(F^a(p)^2,F^b(q)^2\bigr).
\end{equation*}
The pair $F^a(p),F^b(q)$ is centred and jointly Gaussian. Lemma~\ref{8.1:lem-gaussian-squares}
applied to the one-element families $\xi=F^a(p)$ and $\eta=F^b(q)$ gives
$\operatorname{Cov}(F^a(p)^2,F^b(q)^2)=2(\mathbb{E}[F^a(p)F^b(q)])^2$. By
\eqref{8.1:eq-kernel-positivity-1} this equals $2\delta^{ab}C(p,q)^2$. The sum over $a,b$ then
contributes a factor $n$, so the prefactor is $\frac14\cdot2\cdot n=\frac n2=b_0$, which is
\eqref{8.1:eq-kernel-positivity-a}. Since $N\ge2$, we have $n\ge3$ and $b_0\ne0$.

The self-contractions of each square, $\mathbb{E}[F^a(p)^2]=C(p,p)$, enter only the mean
$\mathbb{E}X^{(2)}(f)=\frac n2\sum_pf(p)C(p,p)$. This is a constant and does not contribute to the
covariance.

For $F$ distributed according to $\gamma_{\mathrm{ex}}$, Lemma~\ref{8.1:lem-free-curvature-covariance}
states that the colour components of $F$ are independent centred Gaussians, each with covariance
$P_{\mathrm{ex}}$. Hence
\eqref{8.1:eq-kernel-positivity-1} holds with $C=P_{\mathrm{ex}}$, and the right side of
\eqref{8.1:eq-kernel-positivity-a} is $b_0\mathcal{K}^{\mathrm{ex}}(f_1,f_2)$.

For $F=F_B=dA_B$, Definition~\ref{8.1:def-background-kernel} exhibits $F_B$ as a linear image of the
centred Gaussian field $\mathcal{A}$, so $F_B$ is centred and jointly Gaussian. The colour
components of $\mathcal{A}$ are independent, and $Q_{k_0}$, hence the conditional expectation
$A_B$, acts on each colour component separately; so the cross-colour covariances of $F_B$ vanish,
and its covariance is $C^{\mathrm{bg}}_{K,s}$ in each colour. Thus
\eqref{8.1:eq-kernel-positivity-1} holds with $C=C^{\mathrm{bg}}_{K,s}$. Comparing with the
definition of $\mathcal{K}_{K,s}$ there gives $b_0\mathcal{K}_{K,s}(f_1,f_2)$.

\emph{(b).} By (a) with $f_1=f_2=f$ and $F$ distributed according to $\gamma_{\mathrm{ex}}$,
$b_0\mathcal{K}^{\mathrm{ex}}(f,f)=\operatorname{Var}X^{(2)}(f)\ge0$. Since $b_0>0$, this gives
$\mathcal{K}^{\mathrm{ex}}(f,f)\ge0$. Suppose $\mathcal{K}^{\mathrm{ex}}(f,f)=0$. Put
$\mathsf{q}_f(F'):=\sum_pf(p)F'(p)^2$ for $F'\in\operatorname{im}d_1$. Then
$X^{(2)}(f)=\frac12\sum_a\mathsf{q}_f(F^a)$ is a sum of independent random variables, so
\begin{equation*}
0=\operatorname{Var}X^{(2)}(f)=\sum_{a=1}^n\operatorname{Var}\bigl(\tfrac12\mathsf{q}_f(F^a)\bigr).
\end{equation*}
Every summand is therefore zero. Hence $\mathsf{q}_f(F^1)=c$ almost surely for some constant $c$,
where the colour component $F^1$ of $F$ has, by Lemma~\ref{8.1:lem-free-curvature-covariance}, the
law of the standard Gaussian on $\operatorname{im}d_1$.
This law has a strictly positive density on $\operatorname{im}d_1$, so every non-empty open subset
of $\operatorname{im}d_1$ has positive measure. The set where the continuous function
$\mathsf{q}_f-c$ does not vanish is open and has measure zero, hence it is empty. Thus
$\mathsf{q}_f\equiv c$ on $\operatorname{im}d_1$, and evaluating at $F'=0$ gives $c=0$. By
polarisation, for $F',F''\in\operatorname{im}d_1$,
\begin{equation*}
\sum_pf(p)F'(p)F''(p)=\tfrac14\bigl(\mathsf{q}_f(F'+F'')-\mathsf{q}_f(F'-F'')\bigr)=0 .
\end{equation*}

Fix a plaquette $\bar p$. Let it lie in the $(\mu,\nu)$-plane with lower corner $x$, and write
$e_\rho$ for the unit lattice vectors. Let $b$ and $b'$ be its two bonds in direction $\mu$,
oriented in the positive $\mu$-direction, from $x$ and from $x+e_\nu$. Let $\delta_b$ be the
$1$-cochain equal to $1$ on $b$, to $-1$ on $b$ reversed, and to $0$ elsewhere; define
$\delta_{b'}$ in the same way. Then $d_1\delta_b(p)=\pm1$ if $b$ lies in the boundary of $p$, and
$d_1\delta_b(p)=0$ otherwise; the same holds for $b'$.

A plaquette containing $b$ lies, for some $\rho\ne\mu$, in the $(\mu,\rho)$-plane and has lower
corner $x$ or $x-e_\rho$. A plaquette containing $b'$ lies, for some $\rho\ne\mu$, in the
$(\mu,\rho)$-plane and has lower corner $x+e_\nu$ or
$x+e_\nu-e_\rho$. Since every period of $\mathbb{F}$ is $2L^{m+K}\ge6$
(Definition~\ref{1.1:def-lattice}), these coincide only for
$\rho=\nu$ with lower corner $x$, that is, only for $\bar p$.

Take $F'=d_1\delta_b$ and $F''=d_1\delta_{b'}$, both in $\operatorname{im}d_1$, in the polarised
identity. By the previous paragraph only $p=\bar p$ contributes, so
\begin{equation*}
0=f(\bar p)\,d_1\delta_b(\bar p)\,d_1\delta_{b'}(\bar p)=\pm f(\bar p).
\end{equation*}
As $\bar p$ was arbitrary, $f\equiv0$, contrary to the hypothesis. Hence
\eqref{8.1:eq-kernel-positivity-b} holds.

\emph{(c).} If $f_1,f_2\ge0$, each term $f_1(p)f_2(q)C(p,q)^2$ is a product of non-negative
numbers. This holds for $C=P_{\mathrm{ex}}$ and for $C=C^{\mathrm{bg}}_{K,s}$.

\emph{(d).} Let $p$ lie in the $(\mu,\nu)$-plane and $q$ in the $(\rho,\sigma)$-plane. The
$1$-cochain $d^*\delta_p$ is supported on the four boundary bonds of $p$, so $(d^*\delta_p)_\alpha=0$
for $\alpha\notin\{\mu,\nu\}$; likewise $(d^*\delta_q)_\alpha=0$ for $\alpha\notin\{\rho,\sigma\}$.
The sum over $\alpha\in\{\mu,\nu\}\cap\{\rho,\sigma\}$ in
\eqref{8.1:eq-free-curvature-covariance-a} of Lemma~\ref{8.1:lem-free-curvature-covariance} may
therefore be extended to all four directions $\alpha$ without change.

Since $-\Delta$ commutes with the translations of the torus, its pseudo-inverse acts on each
component as convolution with $G=(-\Delta)^+\delta_0$. Hence
\begin{equation*}
P_{\mathrm{ex}}(p,q)=\sum_{\alpha}\langle(d^*\delta_p)_\alpha,G\ast(d^*\delta_q)_\alpha\rangle
=\langle d^*\delta_p,(-\Delta)^+d^*\delta_q\rangle .
\end{equation*}
Squaring, multiplying by $f_1(p)f_2(q)$, summing, and using $b_0=(N^2-1)/2$ from (a) gives
\eqref{8.1:eq-kernel-positivity-2}.

The sum over $\alpha$ has at most two non-zero terms. For each of them, $(d^*\delta_p)_\alpha$ and
$(d^*\delta_q)_\alpha$ are, up to sign, differences of the unit cochains on the two
$\alpha$-bonds of $p$, respectively of $q$; these two bonds are translates of each other by one
lattice step. The pairing is therefore a second difference of $G$, in agreement with the explicit
forms in Lemma~\ref{8.1:lem-free-curvature-covariance}. Hence each bracket involves finitely many
values of $G$ and nothing else.
\end{proof}

\begin{remark}
Under the lattice measure $\mu_K$, the Gaussian covariance \eqref{8.1:eq-kernel-positivity-a} with
$C=C^{\mathrm{bg}}_{K,s}$, that is, $b_0\,\mathcal{K}_{K,s}(f_1,f_2)$, is, after the normalisation
relating $\mathcal{O}(f)$ to $X^{(2)}(f)$, the leading term of the covariance of the curvature
observables $\mathcal{O}(f_1)$ and $\mathcal{O}(f_2)$ of Definition~\ref{1.3:def-curvature-field}:
the tadpoles, that is, the self-contractions $\mathbb{E}[F^a(p)^2]=C(p,p)$, are constants and enter
only the mean, as shown in the proof of (a), and the first correction to the covariance is of
order $g^6$.
\end{remark}

\begin{lemma}[The background covariance as an orthogonal tower]\label{8.1:lem-background-tower}
Let $\mathcal{A}$, $B=Q_{k_0}\mathcal{A}$, $A_B$, $F_B=dA_B$,
$C^{\mathrm{bg}}=C^{\mathrm{bg}}_{K,s}$ and the Gaussian law with its expectation $\mathbb{E}_\gamma$
be as in Definition~\ref{8.1:def-background-kernel}, with $k_0=K-s$. Let $i_*\ge0$ be the index of
the last block step of Ba{\l}aban's tower above level $k_0$. For $0\le i\le i_*$ let $Q^{(i)}$ be the
$i$-fold composition of the linearised one-step block-averaging map $Q$ of
Definition~\ref{1.5:def-block-average}, with $Q^{(0)}$ the identity, and put
$\mathcal{G}_i:=\sigma(Q^{(i)}B)$, the $\sigma$-algebra generated by the $i$-fold blocked field; for
$i>i_*$ let $\mathcal{G}_i$ be the trivial $\sigma$-algebra. Put
\begin{equation}\label{8.1:eq-background-tower-1}
F_{B,i}:=\mathbb{E}_\gamma[F_B\,|\,\mathcal{G}_i]-\mathbb{E}_\gamma[F_B\,|\,\mathcal{G}_{i+1}],
\qquad
C^{\mathrm{step}}_i(p,q):=\mathbb{E}_\gamma[F_{B,i}(p)\,F_{B,i}(q)].
\end{equation}
Then the $F_{B,i}$ are orthogonal martingale differences along the decreasing sequence of
$\sigma$-algebras $(\mathcal{G}_i)_{i\ge0}$, only finitely many of them are non-zero, and
\begin{equation}\label{8.1:eq-background-tower-2}
C^{\mathrm{bg}}(p,q)=\sum_{i\ge0}C^{\mathrm{step}}_i(p,q)
\qquad\text{for all plaquettes } p,q.
\end{equation}
The right-hand side of \eqref{8.1:eq-background-tower-2} is Ba{\l}aban's tower of renormalisation
steps above level $k_0$ with all couplings set equal to one, in the form stated in
\cite[(1.16)--(1.19)]{B-CMP95} and used in \cite{B-CMP109}.
\end{lemma}

\begin{proof}
For $i<i_*$ we have $Q^{(i+1)}B=Q\bigl(Q^{(i)}B\bigr)$, so the variable $Q^{(i+1)}B$ is a function of
$Q^{(i)}B$; hence $\mathcal{G}_{i+1}\subset\mathcal{G}_i$. For $i\ge i_*$ the inclusion
$\mathcal{G}_{i+1}\subset\mathcal{G}_i$ is trivial, since $\mathcal{G}_{i+1}$ is the trivial
$\sigma$-algebra.

The bare lattice is finite and $F_B$ is a linear function of the centred Gaussian field $\mathcal{A}$,
so $F_B(p)$ is centred and square integrable. It is $\sigma(B)$-measurable, and
$\sigma(B)=\mathcal{G}_0$, so $\mathbb{E}_\gamma[F_B\,|\,\mathcal{G}_0]=F_B$. The conditional expectation
given the trivial $\sigma$-algebra is the expectation, so
$\mathbb{E}_\gamma[F_B\,|\,\mathcal{G}_{i_*+1}]=\mathbb{E}_\gamma[F_B]=0$. Summing
\eqref{8.1:eq-background-tower-1} over $0\le i\le i_*$, the sum telescopes:
\begin{equation}\label{8.1:eq-background-tower-3}
F_B=\sum_{i=0}^{i_*}F_{B,i},
\qquad F_{B,i}=0\ \text{for } i>i_*.
\end{equation}

By the tower property of conditional expectation and $\mathcal{G}_{i+1}\subset\mathcal{G}_i$,
$\mathbb{E}_\gamma[F_{B,i}\,|\,\mathcal{G}_{i+1}]=0$, and $F_{B,i}$ is $\mathcal{G}_i$-measurable; so
the $F_{B,i}$ are martingale differences along $(\mathcal{G}_i)$. Let $i<j$. Then $F_{B,j}(q)$ is
$\mathcal{G}_j$-measurable and $\mathcal{G}_j\subset\mathcal{G}_{i+1}$, hence
\begin{equation*}
\mathbb{E}_\gamma[F_{B,i}(p)F_{B,j}(q)]
=\mathbb{E}_\gamma\bigl[\mathbb{E}_\gamma[F_{B,i}(p)\,|\,\mathcal{G}_{i+1}]\,F_{B,j}(q)\bigr]=0 .
\end{equation*}
Inserting \eqref{8.1:eq-background-tower-3} for both factors in
$C^{\mathrm{bg}}(p,q)=\mathbb{E}_\gamma[F_B(p)F_B(q)]$ and using this orthogonality, only the diagonal
terms $i=j$ remain, which are the $C^{\mathrm{step}}_i(p,q)$; this is
\eqref{8.1:eq-background-tower-2}.

By Definition~\ref{8.1:def-background-kernel}, $B=Q_{k_0}\mathcal{A}$ is the block field of
\cite[(1.18)]{B-CMP95} and $A_B$ is the linear minimiser of \cite[(1.19)]{B-CMP95}. By the
statement that the Gaussian integration of one renormalisation step at the flat background
coincides with the corresponding increment of the free field, a statement proved in this book as an
implication,
the covariances $C^{\mathrm{step}}_i$ are then the successive steps of the decomposition of
\cite[(1.16)--(1.19)]{B-CMP95} above level $k_0$ with all couplings set equal to one, as used in
\cite{B-CMP109}; this is the identification in the last sentence of the lemma.
\end{proof}

Throughout, $C^{\mathrm{bg}}=C^{\mathrm{bg}}_{K,s}$, $F_B$, $B$, $\mathcal{K}_{K,s}$ and the point kernel $k(x,y)$ are as
in Definition~\ref{8.1:def-background-kernel}. The tower steps $C^{\mathrm{step}}_i$ are as in
Lemma~\ref{8.1:lem-background-tower}. The class $\mathfrak{F}^{+}(\mathsf{d},\vartheta)$, the mass
$\mathfrak{n}(f)$, the radius $\rho$ in cells and the witness pairs are as in
Definition~\ref{8.1:def-witness-class}. The letter $\eta$ is the bare spacing in units of the cells of the
witness level $k_0=K-s$. Plaquettes are written $p=(x;\mu\nu)$, as in
Lemma~\ref{8.1:lem-free-curvature-covariance}, and $x_p$ denotes the point $x$ to which $p$ is attached.
In $\mathcal{K}_{K,s}$ and in $k(x,y)$ the plaquettes enter through these points, as in
Definition~\ref{8.1:def-background-kernel}. At each point there are six plaquettes $(x;\mu\nu)$, one
for each coordinate plane.

\begin{lemma}[Upper and lower bounds and marginal scaling of the kernel]\label{8.1:lem-kernel-bounds}
Assume the following.
\begin{itemize}
\item[(h1)] The bounds on the kernel of the linear minimiser hold; they are used here as stated in
\cite[Theorem~1, (9)--(10)]{B-CMP102b}. The operator theorem at free current also holds.
\item[(h2)] The flux identity holds with its constant $C_d$. For every plaquette $P$ of the block lattice
of level $k_0$ there are weights $\varphi_P$ on the plaquettes of $\mathbb{F}$ such that
\begin{equation}\label{8.1:eq-kernel-bounds-1}
\sum_q \varphi_P(q)\,F_B(q)=(dB)(P),\qquad \|\varphi_P\|_1=\sum_q|\varphi_P(q)|=\eta^{-2}.
\end{equation}
Moreover $\varphi_P$ is supported on plaquettes $q$ with $|x_q-x_P|\le C_d$. Here $x_P$ is the point to
which $P$ is attached, and distances are measured in cells. The lower bound on the flux variance also
holds: the variance $\mathfrak{g}_{k_0}$ of $(dB)(P)$ satisfies $\mathfrak{g}_{k_0}\ge c_d$, with a
constant $c_d>0$.
\item[(h3)] The continuum kernel holds in the following form. Let $C_0^{\mathrm{ker}}(L)$ be the constant
of \eqref{8.1:eq-constants-ceiling-c}. For $x\ne0$ let the continuum curvature two-point function be
\begin{equation}\label{8.1:eq-kernel-bounds-2}
\begin{gathered}
\langle F_{\mu\nu}(x)F_{\rho\sigma}(0)\rangle
:=\frac{1}{\pi^2|x|^4}\bigl(I_{\mu\rho}I_{\nu\sigma}-I_{\mu\sigma}I_{\nu\rho}\bigr),\\
I_{\mu\nu}=\delta_{\mu\nu}-2\hat{x}_\mu\hat{x}_\nu,\quad \hat{x}=x/|x|.
\end{gathered}
\end{equation}
In this lemma and its proof only, $\hat{x}$ denotes the unit vector $x/|x|$. The premises of the
continuum kernel, in the form used here, are the bounds \eqref{8.1:eq-kernel-bounds-a} of item~(a) below
and the operator theorem at free current; its conclusion is that, for $K-s\ge k_W(L)$ and all points
$x,y$ with $|x-y|\ge C_0^{\mathrm{ker}}(L)$,
\begin{equation*}
\eta^{-8}\,k(x,y)\ge\frac13\sum_{\mu<\nu,\ \rho<\sigma}\langle F_{\mu\nu}(x-y)F_{\rho\sigma}(0)\rangle^2 .
\end{equation*}
\end{itemize}
Then the following hold.
\begin{itemize}
\item[(a)] There are constants $C_u$ and $\delta_{\square}>0$ such that, for all $i\ge0$ and all
plaquettes $p,q$ of $\mathbb{F}$,
\begin{equation}\label{8.1:eq-kernel-bounds-a}
\begin{gathered}
|C^{\mathrm{bg}}(p,q)|\le C_u\,\eta^4\,(1+|x_p-x_q|)^{-4},\\
|C^{\mathrm{step}}_i(p,q)|\le C_u\,\eta^4 L^{-4i}\,e^{-\delta_{\square}|x_p-x_q|L^{-i}}.
\end{gathered}
\end{equation}
\item[(b)] For every plaquette $P$ of the block lattice of level $k_0$,
\begin{equation*}
\max_{q\in\operatorname{supp}\varphi_P} C^{\mathrm{bg}}(q,q)\ge c_d\,\eta^4 .
\end{equation*}
\item[(c)] For every $f\in\mathfrak{F}^{+}(\mathsf{d},\vartheta)$ with $\rho\ge 2(C_d+2)/\vartheta$, and
at every $K$,
\begin{equation*}
\mathcal{K}_{K,s}(f,f)\ge \tfrac14\,\|f\|_\infty^2\,c_d^2\,\eta^8>0 .
\end{equation*}
\item[(d)] Put $c_*:=2/\pi^4$. For $K-s\ge k_W(L)$ and $|x-y|\ge C_0^{\mathrm{ker}}(L)$,
\begin{equation*}
k(x,y)\ge c_*\,\eta^8\,|x-y|^{-8}.
\end{equation*}
The continuum two-point function \eqref{8.1:eq-kernel-bounds-2} satisfies
\begin{equation}\label{8.1:eq-kernel-bounds-3}
\sum_{\mu<\nu,\ \rho<\sigma} \langle F_{\mu\nu}(x)F_{\rho\sigma}(0)\rangle^2=\frac{6}{\pi^4|x|^8}.
\end{equation}
Its components change sign, and $\langle F_{12}(x)F_{34}(0)\rangle=0$ for every $x\ne0$.
\item[(e)] Put $c_{\mathcal{K}}:=2\cdot10^{-8}/\pi^4$ and $C_{\mathcal{K}}:=36\,C_u^2$. Let
$(f_1,f_2)$ be a witness pair. The upper bounds below hold for every such pair. The lower bounds hold
when in addition $K-s\ge k_W(L)$ and $\rho\ge C_0^{\mathrm{ker}}(L)$; these are the conditions under
which (d) applies to the point pairs of a witness pair. The bounds are
\begin{equation}\label{8.1:eq-kernel-bounds-b}
\begin{gathered}
c_{\mathcal{K}}\,\rho^{-8}\,\mathfrak{n}(f_1)\,\mathfrak{n}(f_2)\le\mathcal{K}_{K,s}(f_1,f_2)
\le C_{\mathcal{K}}\,\rho^{-8}\,\mathfrak{n}(f_1)\,\mathfrak{n}(f_2),\\
c_{\mathcal{K}}\,\vartheta^2\,\|f_1\|_\infty\|f_2\|_\infty\le\mathcal{K}_{K,s}(f_1,f_2)
\le 100\,C_{\mathcal{K}}\,\|f_1\|_\infty\|f_2\|_\infty .
\end{gathered}
\end{equation}
\item[(f)] For fixed $s$ and every witness pair $(f_1,f_2)$, the bounds \eqref{8.1:eq-kernel-bounds-b}
hold with $\mathcal{K}_{K,s}(f_1,f_2)$ replaced by its lower limit and by its upper limit as
$K\to\infty$; the lower bounds there hold when $\rho\ge C_0^{\mathrm{ker}}(L)$.
\end{itemize}
\end{lemma}

\begin{proof}
\emph{(a).} The first bound in \eqref{8.1:eq-kernel-bounds-a} concerns the covariance
$C^{\mathrm{bg}}(p,q)=\mathbb{E}_\gamma[F_B(p)F_B(q)]$ of the background curvature
$F_B=dA_B$. Here $A_B=\mathbb{E}_\gamma[\mathcal{A}\,|\,B]$ is the linear minimiser of
Definition~\ref{8.1:def-background-kernel}. The second bound concerns the steps $C^{\mathrm{step}}_i$
of the orthogonal tower of Lemma~\ref{8.1:lem-background-tower}. These steps are the unit-coupling
tower of Ba{\l}aban above the level $k_0$.

The two bounds of \eqref{8.1:eq-kernel-bounds-a} are the bounds that hypothesis (h1) supplies for
these kernels. Since $A_B$ is the linear function of $B$ given by the minimiser of the linearised
problem, $F_B=dA_B$ is the image of $B$ under the composite of that linear map with $d$. Hence
$C^{\mathrm{bg}}(p,q)$ is expressed through the kernel of this composite and the covariance of $B$. The
bounds \cite[Theorem~1, (9)--(10)]{B-CMP102b} on the kernel of the linear minimiser, together with the
operator theorem at free current, give the two displayed bounds for the background covariance and for
its tower steps, with constants $C_u$ and $\delta_{\square}>0$.

\emph{(b).} The field $F_B=dA_B$ is a linear function of the centred Gaussian field $\mathcal{A}$, so
it is centred. So is $(dB)(P)$. By the flux identity \eqref{8.1:eq-kernel-bounds-1} and the definition
of $C^{\mathrm{bg}}$,
\begin{equation*}
\mathfrak{g}_{k_0}=\mathbb{E}_\gamma\bigl[(dB)(P)^2\bigr]
=\sum_{q,q'}\varphi_P(q)\,\varphi_P(q')\,C^{\mathrm{bg}}(q,q').
\end{equation*}
The Cauchy--Schwarz inequality for $\mathbb{E}_\gamma$ gives
$|C^{\mathrm{bg}}(q,q')|\le C^{\mathrm{bg}}(q,q)^{1/2}C^{\mathrm{bg}}(q',q')^{1/2}$. Each factor on the
right is at most the square root of $\max_{q''\in\operatorname{supp}\varphi_P}C^{\mathrm{bg}}(q'',q'')$.
Hence
\begin{align*}
\mathfrak{g}_{k_0}
&\le\Bigl(\sum_q|\varphi_P(q)|\Bigr)^2\max_{q\in\operatorname{supp}\varphi_P}C^{\mathrm{bg}}(q,q)\\
&=\eta^{-4}\max_{q\in\operatorname{supp}\varphi_P}C^{\mathrm{bg}}(q,q),
\end{align*}
where the equality uses $\|\varphi_P\|_1=\eta^{-2}$. The lower bound $\mathfrak{g}_{k_0}\ge c_d$ of
(h2) now gives (b).

\emph{(c).} Let $f\in\mathfrak{F}^{+}(\mathsf{d},\vartheta)$. Then $f\ge0$ and $f\not\equiv0$, so
$\|f\|_\infty>0$. Let $x_0$ be a point with $f(x_0)=\|f\|_\infty$. The class condition
$\operatorname{Lip}_{k_0}f\le\|f\|_\infty/(\vartheta\rho)$ gives $f(x)\ge\frac12\|f\|_\infty$ whenever
$|x-x_0|\le\vartheta\rho/2$. The hypothesis $\rho\ge2(C_d+2)/\vartheta$ says $\vartheta\rho/2\ge C_d+2$.

The points of the block lattice of level $k_0$ are one cell apart, so every point of the torus lies
within distance $1$ of one of them. Choose a plaquette $P$ of the block lattice of level $k_0$ attached
at such a point, so that $|x_P-x_0|\le2$. By the support condition of (h2), every
$q\in\operatorname{supp}\varphi_P$ satisfies $|x_q-x_P|\le C_d$, hence $|x_q-x_0|\le C_d+2$. Then
$f(x_q)\ge\frac12\|f\|_\infty$ for all $q\in\operatorname{supp}\varphi_P$. By (b) there is a
$q\in\operatorname{supp}\varphi_P$ with $C^{\mathrm{bg}}(q,q)\ge c_d\eta^4$.

Every term of
\begin{equation*}
\mathcal{K}_{K,s}(f,f)=\sum_{p}\sum_{p'}f(x_p)f(x_{p'})\,C^{\mathrm{bg}}(p,p')^2
\end{equation*}
is non-negative. Hence $\mathcal{K}_{K,s}(f,f)$ is at least its diagonal term at this $q$:
\begin{equation*}
\mathcal{K}_{K,s}(f,f)\ge f(x_q)^2\,C^{\mathrm{bg}}(q,q)^2\ge\tfrac14\|f\|_\infty^2\,c_d^2\,\eta^8>0 .
\end{equation*}
This bound carries the factor $\eta^8$ and so depends on $K$. The form of the lower bound that does not
depend on the scale is the one in (e), which rests on (d).

\emph{(d).} We first compute the summed square \eqref{8.1:eq-kernel-bounds-3}. Write
$T_{\mu\nu\rho\sigma}:=I_{\mu\rho}I_{\nu\sigma}-I_{\mu\sigma}I_{\nu\rho}$.

The matrix $I$ is symmetric, and $I^2=\delta$, since $\hat{x}\cdot\hat{x}=1$ gives
\begin{equation*}
(I^2)_{\mu\nu}=\delta_{\mu\nu}-4\hat{x}_\mu\hat{x}_\nu+4\hat{x}_\mu(\hat{x}\cdot\hat{x})\hat{x}_\nu
=\delta_{\mu\nu}.
\end{equation*}
Hence, with all sums
running over the four space indices,
$\sum_{\mu,\rho}I_{\mu\rho}^2=\sum_\mu(I^2)_{\mu\mu}=4$ and $\sum_\mu(I^4)_{\mu\mu}=4$. Expanding the
square,
\begin{align*}
\sum_{\mu,\nu,\rho,\sigma}T_{\mu\nu\rho\sigma}^2
&=2\Bigl(\sum_{\mu,\rho}I_{\mu\rho}^2\Bigr)^2
-2\sum_{\mu,\nu,\rho,\sigma}I_{\mu\rho}I_{\rho\nu}I_{\nu\sigma}I_{\sigma\mu}\\
&=2\cdot4^2-2\cdot4=24 .
\end{align*}
The tensor $T$ is antisymmetric in $(\mu,\nu)$ and in $(\rho,\sigma)$, so it vanishes when $\mu=\nu$ or
$\rho=\sigma$. The unrestricted sum is therefore four times the sum over $\mu<\nu$, $\rho<\sigma$, which
equals $6$. Multiplying by $(\pi^2|x|^4)^{-2}$ gives \eqref{8.1:eq-kernel-bounds-3}.

Next, the vanishing of $\langle F_{12}(x)F_{34}(0)\rangle$. We have
\begin{equation*}
I_{13}I_{24}-I_{14}I_{23}=4\hat{x}_1\hat{x}_3\hat{x}_2\hat{x}_4-4\hat{x}_1\hat{x}_4\hat{x}_2\hat{x}_3=0,
\end{equation*}
so $\langle F_{12}(x)F_{34}(0)\rangle=0$ for every $x\ne0$.

Next, the change of sign. We have
\begin{equation*}
\pi^2|x|^4\langle F_{12}(x)F_{12}(0)\rangle=(1-2\hat{x}_1^2)(1-2\hat{x}_2^2)-4\hat{x}_1^2\hat{x}_2^2
=1-2\hat{x}_1^2-2\hat{x}_2^2 .
\end{equation*}
This equals $1$ when $\hat{x}$ is the unit vector in the third coordinate direction and $-1$ when it is
the unit vector in the first. So the components change sign, while the summed square
\eqref{8.1:eq-kernel-bounds-3} does not vanish.

Finally, the lower bound on $k$. Let $K-s\ge k_W(L)$ and $|x-y|\ge C_0^{\mathrm{ker}}(L)$. The premises
of (h3) hold: the bounds \eqref{8.1:eq-kernel-bounds-a} were proved in (a), and the operator theorem at
free current is part of (h1). The conclusion of (h3), together with \eqref{8.1:eq-kernel-bounds-3} at the
point $x-y$, gives
\begin{equation*}
\eta^{-8}\,k(x,y)\ge\frac13\cdot\frac{6}{\pi^4|x-y|^8}=\frac{2}{\pi^4}\,|x-y|^{-8}
=c_*\,|x-y|^{-8},
\end{equation*}
which is (d).

\emph{(e).} Let $(f_1,f_2)$ be a witness pair. We regroup the double sum defining $\mathcal{K}_{K,s}$ by
the points $x_p$, $x_q$ to which the plaquettes $p$, $q$ are attached. Each plaquette is attached to
exactly one point, and $f_1$, $f_2$ are evaluated at these points. This gives
\begin{equation*}
\mathcal{K}_{K,s}(f_1,f_2)=\sum_{x}\sum_{y}f_1(x)\,f_2(y)\,k(x,y),
\end{equation*}
with $x\in\operatorname{supp}f_1$ and $y\in\operatorname{supp}f_2$. By
Definition~\ref{8.1:def-witness-class}, such point pairs satisfy $\rho\le|x-y|\le10\rho$, measured in
cells.

\emph{Upper bound.} Let $p$ be attached at $x\in\operatorname{supp}f_1$ and $q$ at
$y\in\operatorname{supp}f_2$. Then $|x_p-x_q|=|x-y|\ge\rho$, so $1+|x_p-x_q|\ge\rho$. By
\eqref{8.1:eq-kernel-bounds-a}, $C^{\mathrm{bg}}(p,q)^2\le C_u^2\eta^8\rho^{-8}$. At each point there are
six attached plaquettes, one in each coordinate plane, so
\begin{equation*}
k(x,y)\le36\,C_u^2\,\eta^8\rho^{-8}.
\end{equation*}
Using $f_1,f_2\ge0$ and $\mathfrak{n}(f)=\eta^4\sum_x f(x)$,
\begin{equation*}
\mathcal{K}_{K,s}(f_1,f_2)\le 36\,C_u^2\,\rho^{-8}\,\mathfrak{n}(f_1)\,\mathfrak{n}(f_2)
=C_{\mathcal{K}}\,\rho^{-8}\,\mathfrak{n}(f_1)\,\mathfrak{n}(f_2).
\end{equation*}

\emph{Lower bound.} Assume $K-s\ge k_W(L)$ and $\rho\ge C_0^{\mathrm{ker}}(L)$. Then
$|x-y|\ge\rho\ge C_0^{\mathrm{ker}}(L)$, so (d) applies. Together with $|x-y|\le10\rho$ it gives
\begin{equation*}
k(x,y)\ge c_*\,\eta^8(10\rho)^{-8}=c_*\,10^{-8}\,\eta^8\rho^{-8}.
\end{equation*}
Summing against $f_1(x)f_2(y)\ge0$ gives the lower bound in the first line of
\eqref{8.1:eq-kernel-bounds-b}, since $c_*10^{-8}=c_{\mathcal{K}}$.

\emph{Second line.} Definition~\ref{8.1:def-witness-class} gives
$\vartheta\|f\|_\infty\rho^4\le\mathfrak{n}(f)\le10\rho^4\|f\|_\infty$. Hence
$\rho^{-8}\mathfrak{n}(f_1)\mathfrak{n}(f_2)$ lies between $\vartheta^2\|f_1\|_\infty\|f_2\|_\infty$ and
$100\|f_1\|_\infty\|f_2\|_\infty$. Inserting this into the first line yields the second line of
\eqref{8.1:eq-kernel-bounds-b}.

The bounds in the second line depend neither on $\rho$ nor on $\mathsf{d}$; this is the marginal scaling
of the kernel. The constants $c_{\mathcal{K}}$ and $C_{\mathcal{K}}$ of (e) are those of clause (iv) of
Theorem~\ref{3.1:thm-main}.

\emph{(f).} The constants in \eqref{8.1:eq-kernel-bounds-b} do not depend on $K$, and the condition
$K-s\ge k_W(L)$ holds for all large $K$ at fixed $s$. Hence each inequality of
\eqref{8.1:eq-kernel-bounds-b} holds for all large $K$, and passing to the lower and the upper limit as
$K\to\infty$ preserves it.
\end{proof}

\begin{lemma}[Decay of a moment-free quadratic form against a distant square]
\label{8.1:lem-moment-free-decay}
Work in the Gaussian model, that is, for the free curvature field $F$ on the plaquettes
$\Lambda^2(\mathbb{F})$ of the bare torus lattice $\mathbb{F}$, distributed by the Gaussian law
$\gamma_{\mathrm{ex}}$ with covariance $P_{\mathrm{ex}}$. Write $F^2(p)$ for the square of $F$ at
the plaquette $p$, and $P_{\mathrm{ex}}(q,p)$ for the covariance of the values of $F$ at $q$ and
$p$. The components of $F$ are suppressed, and products are contracted as in $F^2$.

Let $\varrho>0$ (a radius, used in this sense only in this lemma), let $x$ be a point, and let
$D\ge 4\varrho$. Assume the following size and variation bounds on $P_{\mathrm{ex}}$; they are the
input that this lemma takes from the continuum kernel, and below they are used only through
\eqref{8.1:eq-moment-free-decay-1}.
There is a constant $C_\infty$ such that, for all $r>0$ and all plaquettes $q,q',p$ with
$|x(q)-x(p)|\ge r$ and $|x(q')-x(p)|\ge r$,
\begin{equation}\label{8.1:eq-moment-free-decay-1}
|P_{\mathrm{ex}}(q,p)|\le C_\infty\,r^{-4},\qquad
|P_{\mathrm{ex}}(q,p)-P_{\mathrm{ex}}(q',p)|\le C_\infty\,|x(q)-x(q')|\,r^{-5}.
\end{equation}

Let $f_1$ be a function supported within $\varrho$ of $x$. Let
\begin{equation*}
\mathcal{Q}=\sum_{q,q'}a(q,q')\,F(q)F(q')
\end{equation*}
be a quadratic form in $F$ obtained from $\sum_p f_1(x(p))F^2(p)$ by the subtraction of its
$F^2$-moment, with the following three properties:
\begin{itemize}
\item it is localised within $\varrho$ of $x$, that is, $a(q,q')=0$ unless $x(q)$ and $x(q')$
both lie within $\varrho$ of $x$;
\item it has zero $F^2$-moment, that is, $\sum_{q,q'}a(q,q')=0$;
\item $\sum_{q,q'}|a(q,q')|\le\mathfrak{n}(f_1)$.
\end{itemize}
Let $f_2$ be a function whose support lies at distance at least $D$ from $x$, and put
$\|f_2\|_1=\sum_p|f_2(x(p))|$. Then
\begin{equation}\label{8.1:eq-moment-free-decay-a}
\Bigl|\operatorname{Cov}\Bigl(\mathcal{Q},\ \sum_{p} f_2(x(p))\,F^2(p)\Bigr)\Bigr|
\le C\,\mathfrak{n}(f_1)\,\|f_2\|_1\,\varrho\,D^{-9},
\end{equation}
where $C=8(4/3)^9C_\infty^2$ and the covariance is taken under $\gamma_{\mathrm{ex}}$.
\end{lemma}

This lemma supplies the relative factor $\varrho/D$ in part~(b) of the Gaussian extraction.

\begin{proof}
The covariance on the left of \eqref{8.1:eq-moment-free-decay-a} is linear in $f_2$. It equals
$\sum_p f_2(x(p))\operatorname{Cov}(\mathcal{Q},F^2(p))$, the sum running over the plaquettes $p$
with $f_2(x(p))\neq0$, for which $|x(p)-x|\ge D$. It therefore suffices to show that for every
plaquette $p$ with $|x(p)-x|\ge D$
\begin{equation}\label{8.1:eq-moment-free-decay-2}
|\operatorname{Cov}(\mathcal{Q},F^2(p))|\le C\,\mathfrak{n}(f_1)\,\varrho\,D^{-9}.
\end{equation}
Summing \eqref{8.1:eq-moment-free-decay-2} against $|f_2(x(p))|$ then produces the factor
$\|f_2\|_1$.

Fix such a $p$. For plaquettes $q,q'$ we have
\begin{equation*}
F(q)F(q')=\tfrac14\bigl[(F(q)+F(q'))^2-(F(q)-F(q'))^2\bigr].
\end{equation*}
Lemma~\ref{8.1:lem-gaussian-squares}, applied to the square $(F(q)\pm F(q'))^2$ and to $F^2(p)$,
gives
\begin{equation*}
\operatorname{Cov}\bigl((F(q)\pm F(q'))^2,F^2(p)\bigr)
=2\bigl(P_{\mathrm{ex}}(q,p)\pm P_{\mathrm{ex}}(q',p)\bigr)^2.
\end{equation*}
Subtracting the two cases and dividing by four gives
$\operatorname{Cov}(F(q)F(q'),F^2(p))=2P_{\mathrm{ex}}(q,p)P_{\mathrm{ex}}(q',p)$. Hence
\begin{equation}\label{8.1:eq-moment-free-decay-3}
\operatorname{Cov}(\mathcal{Q},F^2(p))=2\sum_{q,q'}a(q,q')\,P_{\mathrm{ex}}(q,p)\,P_{\mathrm{ex}}(q',p).
\end{equation}

If all coefficients $a(q,q')$ vanish there is nothing to prove. Otherwise fix a plaquette $q_0$
with $a(q_0,q_1)\neq0$ for some $q_1$; by localisation $x(q_0)$ lies within $\varrho$ of $x$.
Since $\sum_{q,q'}a(q,q')=0$, we may subtract $2\sum_{q,q'}a(q,q')P_{\mathrm{ex}}(q_0,p)^2=0$ from
\eqref{8.1:eq-moment-free-decay-3}. This gives
\begin{equation*}
\operatorname{Cov}(\mathcal{Q},F^2(p))
=2\sum_{q,q'}a(q,q')\bigl[P_{\mathrm{ex}}(q,p)P_{\mathrm{ex}}(q',p)-P_{\mathrm{ex}}(q_0,p)^2\bigr].
\end{equation*}
The bracket can be rewritten as
\begin{equation*}
\bigl(P_{\mathrm{ex}}(q,p)-P_{\mathrm{ex}}(q_0,p)\bigr)P_{\mathrm{ex}}(q',p)
+P_{\mathrm{ex}}(q_0,p)\bigl(P_{\mathrm{ex}}(q',p)-P_{\mathrm{ex}}(q_0,p)\bigr).
\end{equation*}

Only pairs with $x(q),x(q')$ within $\varrho$ of $x$ contribute. Since $D\ge4\varrho$, each of
$q,q',q_0$ satisfies $|x(\cdot)-x(p)|\ge D-\varrho\ge 3D/4$. Moreover
$|x(q)-x(q_0)|\le2\varrho$ and $|x(q')-x(q_0)|\le2\varrho$. Now apply
\eqref{8.1:eq-moment-free-decay-1} with $r=3D/4$. Each of the two terms of the bracket is then at
most $C_\infty\cdot2\varrho\,(3D/4)^{-5}\cdot C_\infty(3D/4)^{-4}$, so the bracket is at most
$4(4/3)^9C_\infty^2\,\varrho\,D^{-9}$. Multiplying by $2\sum_{q,q'}|a(q,q')|\le2\mathfrak{n}(f_1)$
gives \eqref{8.1:eq-moment-free-decay-2} with $C=8(4/3)^9C_\infty^2$, and hence
\eqref{8.1:eq-moment-free-decay-a}.

Without the zero-moment subtraction the bracket would be
$P_{\mathrm{ex}}(q,p)P_{\mathrm{ex}}(q',p)$, of size $C_\infty^2(4/3)^8D^{-8}$; the gain is the
relative factor $\varrho/D$. Under the hypotheses of the lemma only $\varrho/D\le1/4$ holds. In
part~(b) of the Gaussian extraction the lemma is applied with $\varrho/D\le C'/C_W(s)$ for a
constant $C'$, so the gain there is supplied by the room function $C_W(s)$, not by the grading.
\end{proof}

\begin{remark}[The leading term in one line]\label{8.1:rem-leading-term}
Let $k_0=K-s$ be the witness level. Denote by $\mathcal{O}^{(k_0)}(f)$ the
block-field version at level $k_0$ of the smeared plaquette action density $\mathcal{O}(f)$, and by
$F^a_B$, $a=1,\dots,N^2-1$, the components of the background curvature $F_B$ in an orthonormal
basis of $\mathfrak{su}(N)$. The sums over $p$ below run over the plaquettes as in
Lemma~\ref{8.1:lem-kernel-positivity}, and $f(p)$ stands for $f(x(p))$. At the level of the
leading term,
\begin{equation}\label{8.1:eq-leading-term-1}
\mathcal{O}^{(k_0)}(f)\;\longmapsto\;\tfrac12\,g_{K-s}^2\sum_p f(p)\sum_a\bigl(F^a_B(p)\bigr)^2 ,
\end{equation}
and for two test functions $f_1,f_2$ the covariance of the corresponding leading terms, under the
Gaussian law of the background tower of Lemma~\ref{8.1:lem-background-tower} above $k_0$, is
\begin{equation}\label{8.1:eq-leading-term-2}
\begin{split}
&\operatorname{Cov}\Bigl(\tfrac12\,g_{K-s}^2\sum_p f_1(p)\sum_a\bigl(F^a_B(p)\bigr)^2,\\
&\qquad\qquad\tfrac12\,g_{K-s}^2\sum_p f_2(p)\sum_a\bigl(F^a_B(p)\bigr)^2\Bigr)\\
&\qquad=b_0\,g_{K-s}^4\,\mathcal{K}_{K,s}(f_1,f_2),\qquad b_0=\tfrac12\,(N^2-1).
\end{split}
\end{equation}
Call a \emph{line} one factor $C^{\mathrm{bg}}(p,p')$ of the background covariance; the
single-line exchange is the part of a covariance with one such factor, and the \emph{bubble} is
the part with two. Invariance under $\mathrm{Ad}$ removes the term of first order in the background
curvature and the single-line exchange; it does not remove the bubble. The coupling $g_{K-s}$ in
\eqref{8.1:eq-leading-term-2} is the one read off from the exponent at level $k_0$ of the
correlation theorem, and the marginal response of the finer levels is carried by the normalisation
$N_{k_0}$.
\end{remark}

\begin{proof}
The right side of \eqref{8.1:eq-leading-term-1} equals $g_{K-s}^2X^{(2)}(f)$, where $X^{(2)}(f)$
is the quadratic curvature functional of Lemma~\ref{8.1:lem-kernel-positivity}; this is the
definition of $X^{(2)}(f)$ written out. Hence the left side of \eqref{8.1:eq-leading-term-2} is
$g_{K-s}^4$ times the covariance of $X^{(2)}(f_1)$ and $X^{(2)}(f_2)$. The latter is given by
\eqref{8.1:eq-kernel-positivity-a} with $C$ the background covariance $C^{\mathrm{bg}}$, the
covariance of the tower of Lemma~\ref{8.1:lem-background-tower} above $k_0$, that is,
\begin{equation*}
\operatorname{Cov}\bigl(X^{(2)}(f_1),X^{(2)}(f_2)\bigr)
=b_0\sum_{p,p'}f_1(p)\,f_2(p')\,C^{\mathrm{bg}}(p,p')^2 ;
\end{equation*}
the double sum on the right is, by the definition of $\mathcal{K}_{K,s}$ in the glossary
(Definition~\ref{0.2:not-glossary}), the kernel $\mathcal{K}_{K,s}(f_1,f_2)$, and this yields
\eqref{8.1:eq-leading-term-2}. The vanishing of the term of first order follows from parts (i)
and (ii) of Lemma~\ref{8.1:lem-invariant-facts}, and that of the single-line exchange from its
part (iii), by which the expectation of the product of an $\mathrm{Ad}$-invariant function with an
$\mathrm{Ad}$-covariant one vanishes. The bubble, the part with two lines, is the right side of
\eqref{8.1:eq-kernel-positivity-a} itself, and it is not of either of these two kinds. That the
coupling is $g_{K-s}$ is the statement that the exponent at level $k_0$ of the correlation theorem
fixes it, as in the hypothesis of the two-point witness theorem that takes the level-$k_0$
exponent from the correlation theorem; the response of the levels finer than $k_0$ enters only
through $N_{k_0}$.
\end{proof}

\begin{lemma}[A linear positive history decomposition]\label{8.1:lem-history-decomposition}
Let $s$ be the witness depth and $k_0=K-s$. Let the bare torus lattice $\mathbb{F}$, the observables
$\mathcal{O}(f_i)$, the complex source $w=\sum_i z_if_i$ and the sourced bare density $\rho_0^w$ be as
in Notation~\ref{8.1:not-curvature-observable}, and let $\rho^w_{k_0}$ be the sourced effective density at
level $k_0$ of the correlation theorem at the witness depth (Proposition~\ref{8.1:prop-correlation-witness-depth}).
Write $\rho_0$ for $\rho_0^w$ at $z=0$, $\mathcal{O}_i:=\mathcal{O}(f_i)$ and $\partial_i:=\partial/\partial z_i$.
Assume the following.
\begin{itemize}
\item[(a)] The statement of Proposition~\ref{8.1:prop-correlation-witness-depth} at level $0$: the sourced
effective density at level $0$ is $\rho_0^w$, and the operations performed at the $k_0$ steps from level $0$
to level $k_0$ do not depend on $z$.
\item[(b)] The conservation of the sourced integrals stated in
Proposition~\ref{8.1:prop-correlation-witness-depth}: for every $z$ for which $w$ lies in the domain on which
that proposition defines $\rho^w_{k_0}$, one has $\int\rho^w_{k_0}=\int\rho_0^w$.
\item[(c)] The history representation of the effective density at level $k_0$, in the un-resummed form of
the expansion of Proposition~\ref{8.1:prop-correlation-witness-depth}, which uses \cite[(0.4)]{B-CMP122a}:
there is a finite set of histories $h$, and a linear map $\mathsf{P}_{k_0}$, composed of the operations
of \textup{(a)}, which assigns to every bounded measurable complex function $\varphi$ of the bare configuration
a family $(\mathsf{P}_{k_0}(\varphi)_h)_h$ of measurable functions of the level-$k_0$ configuration, such
that $\rho^w_{k_0}=\sum_h\mathsf{P}_{k_0}(\rho_0^w)_h$ for every $z$ as in \textup{(b)}.
\item[(d)] Positivity of the history terms: $\mathsf{P}_{k_0}(\varphi)_h\ge0$ for every $h$ whenever
$\varphi\ge0$.
\end{itemize}
Then $\mathsf{P}_{k_0}$ is linear, positive and independent of $z$, and the history terms
$(\tau^w_h)_h:=\mathsf{P}_{k_0}(\rho_0^w)$, defined for every $z$, satisfy, for every $z$ as in \textup{(b)},
\begin{equation}\label{8.1:eq-history-decomposition-1}
\rho^w_{k_0}=\sum_h\tau^w_h,\qquad \sum_h\int\tau^w_h=\int\rho_0^w,
\end{equation}
and
\begin{equation}\label{8.1:eq-history-decomposition-2}
\partial_i\tau^w_h\big|_{z=0}=\mathsf{P}_{k_0}(\mathcal{O}_i\rho_0)_h,\qquad
\partial_1\partial_2\tau^w_h\big|_{z=0}=\mathsf{P}_{k_0}(\mathcal{O}_1\mathcal{O}_2\rho_0)_h .
\end{equation}
If the terms of the expansion of the effective density at level $k_0$ in
Proposition~\ref{8.1:prop-correlation-witness-depth} are the sums of the terms $\tau^w_h$ over the classes
of a partition of the set of histories, the same conclusions hold when the histories are replaced by the
index of that expansion.
\end{lemma}

\begin{proof}
Linearity of $\mathsf{P}_{k_0}$ is part of (c), positivity is (d), and independence of $z$ holds because
$\mathsf{P}_{k_0}$ is composed of the operations of (a), which do not depend on $z$. The first identity
in \eqref{8.1:eq-history-decomposition-1} is the identity of (c) with $\tau^w_h=\mathsf{P}_{k_0}(\rho_0^w)_h$.

A domination bound. Let $\varphi$ be real with $|\varphi|\le c\,\rho_0$ for a constant $c$. Then
$c\rho_0+\varphi\ge0$ and $c\rho_0-\varphi\ge0$, so by (d) and linearity $\mathsf{P}_{k_0}(\varphi)_h$
is real and $|\mathsf{P}_{k_0}(\varphi)_h|\le c\,\tau^0_h$, where $\tau^0_h=\mathsf{P}_{k_0}(\rho_0)_h\ge0$.
Splitting a complex $\varphi$ with $|\varphi|\le c\rho_0$ into real and imaginary parts gives
\begin{equation}\label{8.1:eq-history-decomposition-3}
|\mathsf{P}_{k_0}(\varphi)_h|\le 2c\,\tau^0_h .
\end{equation}

The form of $\rho_0^w$. The weighted action of Definition~\ref{1.1:def-wilson-action} is linear in its
weight, and its part linear in $w$ is $-\mathcal{O}(w)$ with $\mathcal{O}(w)=\sum_iz_i\mathcal{O}_i$
(Notation~\ref{8.1:not-curvature-observable}); hence $\rho_0^w=e^{\mathcal{O}(w)}\rho_0$. Since
$0\le a_p(U)\le2$ and $\mathbb{F}$ is finite,
$|\mathcal{O}_i|\le \Gamma_i:=2|\Lambda^2(\mathbb{F})|\,\|f_i\|_\infty$,
so $|\rho_0^w|\le e^{\sum_i|z_i|\Gamma_i}\rho_0$.

Integrals. By \eqref{8.1:eq-history-decomposition-3} applied to $\varphi=\rho_0^w$,
$|\tau^w_h|\le2e^{\sum_i|z_i|\Gamma_i}\tau^0_h$. At $z=0$ the identity of (c) and (b) give
$\sum_h\int\tau^0_h=\int\rho_0<\infty$, and each $\tau^0_h\ge0$; so every $\tau^w_h$ is integrable.
The set of histories being finite, integration commutes with the sum, and by (b), for every $z$ as in (b),
\begin{equation*}
\sum_h\int\tau^w_h=\int\rho^w_{k_0}=\int\rho_0^w ,
\end{equation*}
which is the second identity in \eqref{8.1:eq-history-decomposition-1}.

Derivatives. Fix $z$, let $\xi\in\mathbb{C}\setminus\{0\}$, and let $z^\xi$ differ from $z$ by
$z_i\mapsto z_i+\xi$, with source $w^\xi$. Then
\begin{equation*}
\psi_\xi:=\frac{\rho_0^{w^\xi}-\rho_0^w}{\xi}-\mathcal{O}_i\rho_0^w
=\rho_0^w\Bigl(\frac{e^{\xi\mathcal{O}_i}-1}{\xi}-\mathcal{O}_i\Bigr).
\end{equation*}
For $|a|\le\Gamma_i$ the power series of the exponential gives
\begin{equation*}
\Bigl|\frac{e^{\xi a}-1}{\xi}-a\Bigr|\le|\xi|\,\Gamma_i^2e^{|\xi|\Gamma_i}.
\end{equation*}
Hence $|\psi_\xi|\le c(\xi)\rho_0$, where
\begin{equation*}
c(\xi):=|\xi|\,\Gamma_i^2e^{|\xi|\Gamma_i}e^{\sum_j|z_j|\Gamma_j}\longrightarrow0\qquad(\xi\to0).
\end{equation*}
By linearity and independence of $z$,
\begin{equation*}
\frac{\tau^{w^\xi}_h-\tau^w_h}{\xi}-\mathsf{P}_{k_0}(\mathcal{O}_i\rho_0^w)_h
=\mathsf{P}_{k_0}(\psi_\xi)_h ,
\end{equation*}
whose modulus is at most $2c(\xi)\tau^0_h$ by \eqref{8.1:eq-history-decomposition-3}. Thus, pointwise
in the level-$k_0$ configuration, $\partial_i\tau^w_h=\mathsf{P}_{k_0}(\mathcal{O}_i\rho_0^w)_h$ for all $z$;
no restriction on $z$ enters here, since $\tau^w_h=\mathsf{P}_{k_0}(\rho_0^w)_h$ is defined for all $z$.
The same argument with $\rho_0^w$ replaced by $\mathcal{O}_1\rho_0^w$, which satisfies
$|\mathcal{O}_1\rho_0^w|\le \Gamma_1e^{\sum_j|z_j|\Gamma_j}\rho_0$, and with $i=2$, gives
$\partial_2\partial_1\tau^w_h=\mathsf{P}_{k_0}(\mathcal{O}_1\mathcal{O}_2\rho_0^w)_h$. Exchanging the roles
of $1$ and $2$ gives the same right-hand side for $\partial_1\partial_2\tau^w_h$, since
$\mathcal{O}_1\mathcal{O}_2=\mathcal{O}_2\mathcal{O}_1$. Setting $z=0$, so that
$\rho_0^w=\rho_0$, yields \eqref{8.1:eq-history-decomposition-2}.

The two indexings. Suppose that the terms of the expansion in
Proposition~\ref{8.1:prop-correlation-witness-depth} are the sums of the $\tau^w_h$ over the classes
$\bar h$ of a partition of the set of histories, and put
$\mathsf{P}'_{k_0}(\varphi)_{\bar h}:=\sum_{h\in\bar h}\mathsf{P}_{k_0}(\varphi)_h$. A finite sum of linear,
positive, $z$-independent maps is linear, positive and $z$-independent, and
$\tau'^w_{\bar h}:=\mathsf{P}'_{k_0}(\rho_0^w)_{\bar h}=\sum_{h\in\bar h}\tau^w_h$. Since the classes
partition the histories, summing \eqref{8.1:eq-history-decomposition-1} and
\eqref{8.1:eq-history-decomposition-2} over the classes gives the same identities for
$\tau'^w_{\bar h}$ and $\mathsf{P}'_{k_0}$.
\end{proof}

\begin{lemma}[The total covariance formula over histories]\label{8.1:lem-total-covariance}
Assume the hypotheses of Lemma~\ref{8.1:lem-history-decomposition} and use its notation: the witness
level $k_0=K-s$, the test functions $f_1,f_2$ with $\mathcal{O}_i:=\mathcal{O}(f_i)$, the complex
source $w=z_1f_1+z_2f_2$, the sourced bare density $\rho_0^w$, the linear positive $z$-free map
$\mathsf{P}_{k_0}$ and the history terms $(\tau^w_h)_h=\mathsf{P}_{k_0}(\rho_0^w)$, functions of the
block field $V$ at level $k_0$, integrated against the reference measure $dV$ of that lemma. Put
$\tau_h:=\tau^w_h|_{z=0}$, $Z(z):=\int\rho_0^w$ and $Z:=Z(0)$. Let $\tilde{\nu}$ be the probability
law on pairs $(h,V)$ with density $\tau_h(V)/Z$; its marginal in $V$ is the block-field law
$\nu^B_{K,s}$. On $\{\tau_h>0\}$ put
\begin{equation*}
\Phi^h_i:=\partial_{z_i}\log\tau^w_h\big|_{z=0}\quad(i=1,2),\qquad
\Psi^h:=\partial_{z_1}\partial_{z_2}\log\tau^w_h\big|_{z=0}.
\end{equation*}
Then $\Phi^h_1,\Phi^h_2,\Psi^h$ are $\tilde{\nu}$-integrable, and the connected two-point function
$S^T_{2;K}=\partial_{z_1}\partial_{z_2}\log Z(z)|_{z=0}$ of $\mathcal{O}_1,\mathcal{O}_2$ satisfies
\begin{equation}\label{8.1:eq-total-covariance-a}
S^T_{2;K}=\mathbb{E}_{\tilde{\nu}}\big[\Psi^h\big]
+\operatorname{Cov}_{\tilde{\nu}}\big(\Phi^h_1,\Phi^h_2\big),
\end{equation}
where $\operatorname{Cov}_{\tilde{\nu}}(A,B):=\mathbb{E}_{\tilde{\nu}}[AB]
-\mathbb{E}_{\tilde{\nu}}[A]\,\mathbb{E}_{\tilde{\nu}}[B]$.
\end{lemma}

\begin{proof}
\emph{The law $\tilde{\nu}$.} By Lemma~\ref{8.1:lem-history-decomposition} at $z=0$,
$\sum_h\int\tau_h\,dV=\int\rho_0=Z$, and $\tau_h\ge 0$ because $\mathsf{P}_{k_0}$ is positive and
$\rho_0\ge0$; so $\tau_h/Z$ is a probability density on pairs $(h,V)$. Since
$\sum_h\tau_h=\rho_{k_0}$, its marginal in $V$ has density $\rho_{k_0}/Z$, which is $\nu^B_{K,s}$.

\emph{Derivatives of the history terms are dominated by $\tau_h$.} By
Lemma~\ref{8.1:lem-history-decomposition},
\begin{equation}\label{8.1:eq-total-covariance-1}
\partial_{z_i}\tau^w_h\big|_{z=0}=\mathsf{P}_{k_0}(\mathcal{O}_i\rho_0)_h,\qquad
\partial_{z_1}\partial_{z_2}\tau^w_h\big|_{z=0}=\mathsf{P}_{k_0}(\mathcal{O}_1\mathcal{O}_2\rho_0)_h.
\end{equation}
Since $\operatorname{tr}\mathbf{1}=1$, $0\le1-\operatorname{Re}\operatorname{tr}U(\partial p)\le2$
for every plaquette $p$, so $|\mathcal{O}(f)|\le c(f):=2\|f\|_\infty\,|\Lambda^2(\mathbb{F})|$. For
real-valued $F$ with $|F|\le c$ we have $-c\rho_0\le F\rho_0\le c\rho_0$, and linearity and
positivity of $\mathsf{P}_{k_0}$ give $|\mathsf{P}_{k_0}(F\rho_0)_h|\le c\,\tau_h$. Applying this to
the real and imaginary parts of $\mathcal{O}_i$ and of $\mathcal{O}_1\mathcal{O}_2$, the right-hand
sides of \eqref{8.1:eq-total-covariance-1} are bounded in modulus by $2c(f_i)\tau_h$ and by
$4c(f_1)c(f_2)\tau_h$ respectively. In particular both derivatives vanish wherever $\tau_h=0$.

\emph{The logarithmic derivatives.} Where $\tau_h>0$, $\tau^w_h$ is twice differentiable in $z$ and
non-zero for $z$ near $0$, and the chain rule gives
\begin{equation}\label{8.1:eq-total-covariance-2}
\partial_{z_i}\tau^w_h\big|_{z=0}=\tau_h\,\Phi^h_i,\qquad
\partial_{z_1}\partial_{z_2}\tau^w_h\big|_{z=0}=\tau_h\big(\Psi^h+\Phi^h_1\Phi^h_2\big).
\end{equation}
By the previous paragraph $|\Phi^h_i|\le2c(f_i)$ and
$|\Psi^h|\le4c(f_1)c(f_2)+4c(f_1)c(f_2)$, so all three are bounded and $\tilde{\nu}$-integrable.

\emph{The derivatives of $Z(z)$.} Differentiating the identity
$\sum_h\int\tau^w_h\,dV=\int\rho_0^w=Z(z)$ of Lemma~\ref{8.1:lem-history-decomposition} at $z=0$,
the integrands vanish on $\{\tau_h=0\}$, and on $\{\tau_h>0\}$
\eqref{8.1:eq-total-covariance-2} applies; hence
\begin{equation*}
\partial_{z_i}Z\big|_{z=0}=\sum_h\int_{\{\tau_h>0\}}\tau_h\,\Phi^h_i\,dV
=Z\,\mathbb{E}_{\tilde{\nu}}\big[\Phi^h_i\big],
\end{equation*}
\begin{equation*}
\partial_{z_1}\partial_{z_2}Z\big|_{z=0}=\sum_h\int_{\{\tau_h>0\}}\tau_h\big(\Psi^h
+\Phi^h_1\Phi^h_2\big)\,dV
=Z\Big(\mathbb{E}_{\tilde{\nu}}\big[\Psi^h\big]+\mathbb{E}_{\tilde{\nu}}
\big[\Phi^h_1\Phi^h_2\big]\Big).
\end{equation*}

\emph{Conclusion.} Since $Z>0$,
\begin{equation*}
\partial_{z_1}\partial_{z_2}\log Z(z)\big|_{z=0}
=\frac{\partial_{z_1}\partial_{z_2}Z|_{z=0}}{Z}
-\frac{\partial_{z_1}Z|_{z=0}\;\partial_{z_2}Z|_{z=0}}{Z^2},
\end{equation*}
and inserting the two previous displays gives
\begin{equation*}
S^T_{2;K}=\mathbb{E}_{\tilde{\nu}}[\Psi^h]+\mathbb{E}_{\tilde{\nu}}[\Phi^h_1\Phi^h_2]
-\mathbb{E}_{\tilde{\nu}}[\Phi^h_1]\,\mathbb{E}_{\tilde{\nu}}[\Phi^h_2],
\end{equation*}
which is \eqref{8.1:eq-total-covariance-a}.
\end{proof}

\begin{remark}
The argument is algebraic and uses no smallness condition. It compares no two laws: the only law
is $\tilde{\nu}$, and $\Phi^h_i$, $\Psi^h$ are logarithmic derivatives of the history terms
$\tau^w_h$ themselves, not of the conditional probability $\tau_h/\rho_{k_0}$ of the history $h$
given $V$, which does not enter $\Phi^h_i$.
\end{remark}

\begin{lemma}[The running normalisation at the witness level]\label{8.1:lem-running-normalisation}
Let $k_0\le K$. Let $N_{k_0}=\zeta'_{k_0}$ be the normalisation at level $k_0$, that is, the
derivative at zero of the running shift. Let $b_j(\zeta)$, $y_j$, $d_j(\varrho)$ be the running
quantities and $r$, $\lambda$, $c_2$ the constants of the correlation theorem at the witness depth
(Proposition~\ref{8.1:prop-correlation-witness-depth}), with $b_j'(0)$ the derivative of $b_j(\zeta)$
at $\zeta=0$. Let $\kappa_0\ge 7$ be the irrelevance exponent of the correlation theorem at the
witness depth. Assume $x_{k_0}>4c_2$, and put $\hat{g}:=(x_{k_0}-4c_2)^{-1/2}$.
Assume that the correlation theorem at the witness depth
(Proposition~\ref{8.1:prop-correlation-witness-depth}) provides (i)--(iv) below, and assume (v):
\begin{enumerate}
\item[(i)] its normalisation identity: the normalisation is the sum of the increments'
derivatives,
\begin{equation*}
N_{k_0}=1+\sum_{j\le k_0}b_j'(0);
\end{equation*}
\item[(ii)] its bound on the increments: for every $j\le k_0$ one has $|b_j'(0)|\le d_{j-1}(2r)$;
\item[(iii)] its reference bound: the bounds in (ii) have total at most one third,
\begin{equation*}
\sum_{j\le k_0}d_{j-1}(2r)\le\tfrac13;
\end{equation*}
\item[(iv)] its lower bound on $y_i$ below the witness level: for every level $i<k_0$, with
$\nu=k_0-i$, one has $y_i\ge x_{k_0}+\lambda\nu-4c_2$;
\item[(v)] the bound obtained by summing the running-shift statement of the correlation theorem
at the witness depth over the levels strictly below $k_0$: if the lower bound (iv) holds at every
level $i<k_0$, then there is a constant $C_{\mathfrak{D}}'$, depending on none of $K$, $k_0$, $s$,
such that
\begin{equation}\label{8.1:eq-running-normalisation-1}
\begin{split}
\sum_{i<k_0}d_i(2r)
&\le C_{\mathfrak{D}}'\Bigl[\lambda^{-1}\hat{g}\,(\log\hat{g}^{-2})^{p_0}
+(1+\lambda^{-1})\,\frac{\hat{g}^{\kappa_0-6}}{r}\Bigr]\\
&=:\mathfrak{d}_W(x_{k_0}).
\end{split}
\end{equation}
\end{enumerate}
Then, with the constant $C_{\mathfrak{D}}'$ of (v) and $\mathfrak{d}_W$ as in
\eqref{8.1:eq-running-normalisation-1},
\begin{equation}\label{8.1:eq-running-normalisation-a}
N_{k_0}\in\bigl[\tfrac23,\tfrac43\bigr],\qquad N_{k_0}^2\in\bigl[\tfrac49,\tfrac{16}9\bigr],\qquad
|N_{k_0}-1|\le\mathfrak{d}_W(x_{k_0}).
\end{equation}
Moreover, every set $J$ of levels $j\le k_0$ satisfies
\begin{equation}\label{8.1:eq-running-normalisation-2}
\Bigl|\sum_{j\in J}b_j'(0)\Bigr|\le\mathfrak{d}_W(x_{k_0}).
\end{equation}
\end{lemma}

\begin{proof}
By (ii), for every $j\le k_0$ we have $d_{j-1}(2r)\ge|b_j'(0)|\ge0$, so all the numbers
$d_{j-1}(2r)$ with $j\le k_0$ are non-negative.

We first prove \eqref{8.1:eq-running-normalisation-2}. Let $J$ be a set of levels $j\le k_0$.
We apply the triangle inequality, then (ii) term by term, and then enlarge the sum from $J$ to
all $j\le k_0$, which is allowed because the added terms are non-negative. This gives
\begin{equation}\label{8.1:eq-running-normalisation-3}
\Bigl|\sum_{j\in J}b_j'(0)\Bigr|\le\sum_{j\in J}|b_j'(0)|\le\sum_{j\in J}d_{j-1}(2r)
\le\sum_{j\le k_0}d_{j-1}(2r).
\end{equation}
Now substitute $i=j-1$. As $j$ runs over the levels $j\le k_0$, the index $i$ runs over the levels
$i\le k_0-1$, which are exactly the levels strictly below $k_0$. So the last sum in
\eqref{8.1:eq-running-normalisation-3} equals $\sum_{i<k_0}d_i(2r)$. At each of these levels the
lower bound $y_i\ge x_{k_0}+\lambda(k_0-i)-4c_2$ of (iv) holds. Hence the hypothesis of (v) is
satisfied, and (v), with its constant $C_{\mathfrak{D}}'$, gives
\begin{equation*}
\sum_{j\le k_0}d_{j-1}(2r)=\sum_{i<k_0}d_i(2r)\le\mathfrak{d}_W(x_{k_0}).
\end{equation*}
Combining this with \eqref{8.1:eq-running-normalisation-3} proves
\eqref{8.1:eq-running-normalisation-2}.

Next we prove the rate bound, the third assertion of \eqref{8.1:eq-running-normalisation-a}. By (i),
\begin{equation*}
N_{k_0}-1=\sum_{j\le k_0}b_j'(0).
\end{equation*}
This is \eqref{8.1:eq-running-normalisation-2} for the set $J$ of all levels $j\le k_0$, so
$|N_{k_0}-1|\le\mathfrak{d}_W(x_{k_0})$.

Finally we prove the first two assertions of \eqref{8.1:eq-running-normalisation-a}. By
\eqref{8.1:eq-running-normalisation-3} for the set of all levels $j\le k_0$, together with (iii),
\begin{equation*}
|N_{k_0}-1|\le\sum_{j\le k_0}d_{j-1}(2r)\le\tfrac13.
\end{equation*}
Hence $N_{k_0}\in[\tfrac23,\tfrac43]$. Since $N_{k_0}\ge\tfrac23>0$ and squaring is increasing on
$[0,\infty)$, it follows that $N_{k_0}^2\in[\tfrac49,\tfrac{16}9]$.
\end{proof}

\begin{remark}
In the identity of hypothesis (i) of Lemma~\ref{8.1:lem-running-normalisation}, the marginal
response to the source of the large-field regions healed at the levels $j\le k_0$ is contained in
the derivatives $b_j'(0)$, by the correlation theorem at the witness depth
(Proposition~\ref{8.1:prop-correlation-witness-depth}) and its running-shift statement; see also
\cite[(6.1)]{B-CMP119}. The bounds \eqref{8.1:eq-running-normalisation-a} and
\eqref{8.1:eq-running-normalisation-2} therefore include that response; it requires no separate
estimate. In the two-point witness theorem the normalisation enters through the interval for
$N_{k_0}^2$ in \eqref{8.1:eq-running-normalisation-a}; the rate $\mathfrak{d}_W$ enters only its
remainder $R_K$ and its statement of what it does not assert.
\end{remark}

\begin{proposition}[Structure of the first logarithmic derivative of a history term]
\label{8.1:prop-logarithmic-derivative}
Let $(f_1,f_2)$ be a witness pair in the sense of Definition~\ref{8.1:def-witness-class}, with
the neighbourhoods $N_i$ fixed there, and let $k_0=K-s$ be the witness level. Let $h$ be a history
of the linear positive history decomposition of Lemma~\ref{8.1:lem-history-decomposition}, with
history term $\tau^w_h$, and let $\Phi^h_i$ and $\Psi^h$ be the first and the mixed logarithmic
derivatives of $\tau^w_h$ at $z=0$ defined in Lemma~\ref{8.1:lem-total-covariance}. The large-field
regions indexing $h$ are called the \emph{labels} of $h$; a label is \emph{live} if it still carries
its large-field factor at level $k_0$, and \emph{healed} if that factor has been absorbed back into
the small-field format by a later step; $j(Y)$ is the level at which the label $Y$ was created.
Lemma~\ref{8.1:lem-history-decomposition} admits two history indices: the \emph{summed index},
which is the history index of the correlation theorem, and the \emph{refined index}, its
un-resummed refinement, of which the summed index is a partial sum.
For a function $f$ on the bare lattice and a plaquette $p$ we write $f(p)$ for $f(x(p))$, and
$\|f\|_\infty$ for the sup norm. Assume:
\begin{itemize}
\item[(P1)] the correlation theorem at the witness depth, Proposition~\ref{8.1:prop-correlation-witness-depth},
that is, the correlation theorem with $m:=s$: its representation of the sourced effective density at
level $k_0$ as a sum over histories, with all its parts, in which the history term of index $h$
carries the sourced families $\mathbf{T}^w_h$, $\mathbf{E}^w$, $\mathbf{D}^w$, $\mathbf{R}^w$,
$\mathbf{B}^w$, the sourced term $E^w$ subtracted from them, and the running shift
$\zeta_{k_0}(p)$ multiplying the plaquette
term $a_p(U_{k_0}(V))$ of the block field; its running-shift statement; its statement at every
level $j\le k_0$; together with the representation of the running shift by weights
$\varpi^{(j)}_y(p)$ and plaquettes $p_y$ in \cite[(3.40), (3.65)]{B-CMP119}, and the
localisation functions $\varphi_j$ of the history at level $j$ in \cite[(2.24)]{B-CMP119};
\item[(P2)] the profile absorption for the running shift at level $k_0$: for every plaquette $p$
of level $k_0$ the identity \eqref{8.1:eq-logarithmic-derivative-a} of item (b) below holds, with
$\delta^{\mathrm{reg}}_i$, $\theta^h$ and $\delta^{h,\mathrm{lab}}_i$ as defined there; the weights
$\varpi^{(j)}_y(p)$ vanish unless the distance from $p$ to $p_y$ is at most $C_{\mathrm{pr}}M$; and
the profile term $\delta^{\mathrm{reg}}_i$ is Lipschitz, so that, divided by $N_{k_0}$, it is absorbed
into the test function $f_i$ as a Lipschitz modification of it;
\item[(P3)] for the last two assertions of \textup{(e)} only: the source lemma.
\end{itemize}
Let $\mathfrak{d}_W(x_{k_0})$ be the rate of \eqref{8.1:eq-running-normalisation-a}, and let $r$ be as
in Lemma~\ref{8.1:lem-running-normalisation}; it is the radius in the sources of the analyticity
domains $\mathcal{W}_j$ of the correlation theorem, in the sense that a Cauchy estimate on these
domains costs a factor $\|f_i\|_\infty/r$ per derivative in $z_i$. Then on the event
$\{\tau_h>0\}$, for $i=1,2$, the following hold.
\begin{itemize}
\item[(a)] With $a_p$ the plaquette term of (P1),
\begin{equation}\label{8.1:eq-logarithmic-derivative-1}
\begin{split}
\Phi^h_i={}&\partial_{z_i}\log\mathbf{T}^w_h\big|_0
+\sum_p\partial_{z_i}\zeta_{k_0}(p)\big|_0\,a_p\bigl(U_{k_0}(V)\bigr)\\
&+\partial_{z_i}\bigl[\mathbf{E}^w+\mathbf{D}^w+\mathbf{R}^w+\mathbf{B}^w-E^w\bigr]\big|_0 .
\end{split}
\end{equation}
Here $\log\tau^w_h$ is, up to terms independent of $z$, the sum of $\log\mathbf{T}^w_h$ and of the
\emph{exponent} $A^w_{k_0,h}$ of the history term, and $\partial_{z_i}A^w_{k_0,h}|_0$ is the sum of
the last two terms of \eqref{8.1:eq-logarithmic-derivative-1}.
\item[(b)] For every plaquette $p$ of level $k_0$,
\begin{equation}\label{8.1:eq-logarithmic-derivative-a}
\partial_{z_i}\zeta_{k_0}(p)\big|_0
=N_{k_0}f_i(p)+\delta^{\mathrm{reg}}_i(p)+\theta^h(p)f_i(p)+\delta^{h,\mathrm{lab}}_i(p),
\end{equation}
where
\begin{align*}
\delta^{\mathrm{reg}}_i(p)&:=\sum_{j\le k_0}b_j'(0)\sum_y\varpi^{(j)}_y(p)\bigl[f_i(p_y)-f_i(p)\bigr],\\
\theta^h(p)&:=-\sum_{j\le k_0}\bigl(1-\varphi_j(p)\bigr)b_j'(0),\\
\delta^{h,\mathrm{lab}}_i(p)&:=-\sum_{j\le k_0}\bigl(1-\varphi_j(p)\bigr)b_j'(0)
\sum_y\varpi^{(j)}_y(p)\bigl[f_i(p_y)-f_i(p)\bigr].
\end{align*}
The term $\delta^{\mathrm{reg}}_i$ does not depend on $h$, is supported in
$(\operatorname{supp}f_i)^{+C_{\mathrm{pr}}M}$, and
\begin{equation}\label{8.1:eq-logarithmic-derivative-2}
|\delta^{\mathrm{reg}}_i|\le\mathfrak{d}_W(x_{k_0})\,\frac{C_{\mathrm{pr}}M\|f_i\|_\infty}{\vartheta\rho}.
\end{equation}
The terms $\theta^h$ and $\delta^{h,\mathrm{lab}}_i$ vanish outside the labels of $h$, and
\begin{equation}\label{8.1:eq-logarithmic-derivative-3}
|\theta^h|\le\mathfrak{d}_W(x_{k_0}),\qquad
|\delta^{h,\mathrm{lab}}_i|\le\mathfrak{d}_W(x_{k_0})\,\frac{C_{\mathrm{pr}}M\|f_i\|_\infty}{\vartheta\rho}.
\end{equation}
For plaquettes $p$ of level $k_0$ put $\tilde f_i(p):=f_i(p)+\delta^{\mathrm{reg}}_i(p)/N_{k_0}$ and
$\varepsilon^h_i(p):=\theta^h(p)f_i(p)+\delta^{h,\mathrm{lab}}_i(p)$, and for a function $f$ on the
plaquettes of level $k_0$, or on the bare lattice with the convention above,
\begin{equation*}
\mathcal{O}^{(k_0)}(f)(V):=\sum_pf(p)\,a_p\bigl(U_{k_0}(V)\bigr).
\end{equation*}
Then $\tilde f_i$ is Lipschitz, $\varepsilon^h_i$ vanishes outside the labels of $h$,
\begin{equation}\label{8.1:eq-logarithmic-derivative-4}
\begin{gathered}
|\tilde f_i-f_i|\le\tfrac32\,\mathfrak{d}_W(x_{k_0})\,\frac{C_{\mathrm{pr}}M\|f_i\|_\infty}{\vartheta\rho},\\
|\varepsilon^h_i|\le\mathfrak{d}_W(x_{k_0})\Bigl(1+\frac{C_{\mathrm{pr}}M}{\vartheta\rho}\Bigr)\|f_i\|_\infty,
\end{gathered}
\end{equation}
and the plaquette sum in \eqref{8.1:eq-logarithmic-derivative-1} equals
$N_{k_0}\mathcal{O}^{(k_0)}(\tilde f_i)(V)+\mathcal{O}^{(k_0)}(\varepsilon^h_i)(V)$.
\item[(c)] The function
\begin{equation}\label{8.1:eq-logarithmic-derivative-b}
\begin{split}
\mathrm{Irr}_i(V):={}&\partial_{z_i}\bigl[\mathbf{E}^w+\mathbf{D}^w+\mathbf{R}^w-E^w\bigr]\big|_0(V)\\
&-\partial_{z_i}\bigl[\mathbf{E}^w+\mathbf{D}^w+\mathbf{R}^w-E^w\bigr]\big|_0(\mathbf{1})
\end{split}
\end{equation}
does not depend on $h$. Its contribution from each level $j\le k_0$ is the $z_i$-derivative at $0$
of whole-lattice families from which the coupling increments have been extracted, and whose
level-$j$ remainders $R^{(j)}_\zeta$ satisfy $\|R^{(j)}_\zeta\|^{\flat}\le g_j^{\kappa_0}$ in the norm
$\|\cdot\|^{\flat}$ in which the correlation theorem estimates them; by the irrelevance exponent of
\cite[(0.29)]{B-CMP109}, equivalently by the per-cell scale sums of the correlation theorem, these
contributions are summable over the levels $j\le k_0$. We write $c^h_i$ for the value at
$V=\mathbf{1}$ of the right-hand side of \eqref{8.1:eq-logarithmic-derivative-1}, and
$c^h_i=\Phi^h_i(\mathbf{1})$.
\item[(d)] On the summed index the set $\mathcal{D}_i$ and the shadow $\mathrm{Sh}^h_i$ defined below
satisfy $\mathcal{D}_i=\emptyset$ and $\mathrm{Sh}^h_i\equiv0$; the rest of this item concerns the
refined index. Put $\mathfrak{l}:=L\check{R}_{k_0}+c_L$, with $c_L$ a constant, and let
$\mathcal{D}_i$ be the set of healed labels
$Y$ of $h$ meeting $N_i$ with $j(Y)<k_0-\mathfrak{l}$ (the \emph{deep} labels); the threshold
$k_0-\mathfrak{l}$ is not a fixed number of levels below $k_0$ but depends on $k_0$ through
$\check{R}_{k_0}$. For $Y\in\mathcal{D}_i$ the terms of the exponent
$A^w_{k_0,h}$ and of $\log\mathbf{T}^w_h$ \emph{attached} to $Y$ are: the terms $V^{\sharp}$ of the
history representation of the correlation theorem and their descendants; the terms of the primed
form $\mathbf{B}'$ of Ba{\l}aban's large-field operation $\mathbf{B}$, in the case in which the family
$\mathcal{P}_F$ entering its definition is empty; the flat deficits
$-\sum_{X\not\subset\Lambda_j(h)}e_j$ \cite{B-CMP119} of the flat-deficit item of that
representation and of the torus bound of the correlation theorem
(Proposition~\ref{8.1:prop-correlation-witness-depth}), $\Lambda_j(h)$ being the level-$j$ region
of $h$ in that
item; and the terms of $\mathcal{O}^{(k_0)}(\varepsilon^h_i)$ with $p$ in $Y$. These terms are indexed
by domains meeting the reading closure $Y^{\boxplus}$. For such a domain $D$ let
$\psi^h_{i,Y}(D)(V)$ be the $z_i$-derivative at $0$ of the sum of the terms attached to $Y$ and
indexed by $D$, and let $\psi^h_{i,Y}(V):=\sum_D\psi^h_{i,Y}(D)(V)$ be the $z_i$-derivative at $0$
of the sum of all the terms attached to $Y$. Let the
\emph{shadow} $\mathrm{Sh}^h_i$ be $\sum_{Y\in\mathcal{D}_i}\psi^h_{i,Y}$ together with the
$z_i$-derivatives at $0$ of the terms attached to healed labels $Y$ of $h$ with
$j(Y)<k_0-\mathfrak{l}$ not meeting $N_i$ whose supports reach $\operatorname{supp}f_i$. Then for
every $Y\in\mathcal{D}_i$ and every domain $D$ meeting
$Y^{\boxplus}$ that indexes terms attached to $Y$,
\begin{equation}\label{8.1:eq-logarithmic-derivative-c}
\begin{split}
&\sup_V|\psi^h_{i,Y}(D)|\le\sigma_Y,\\
&\sigma_Y:=C_\sigma\bigl(M\check{R}_{j(Y)}\bigr)^{c_\sigma}
\Bigl(c_1^{\flat}(g_{j(Y)})\vee B_0\vee C_e+\mathfrak{d}_W(x_{k_0})+g_{j(Y)}^{\kappa_0}\Bigr)
\frac{\|f_i\|_\infty}{r},
\end{split}
\end{equation}
where $\vee$ denotes the maximum, $c_1^{\flat}(g)\in\{e^{-p_0},\alpha^{1/3}\}$ as in
\cite{B-CMP122b}, $B_0$ and $C_e$ are the constants bounding the sizes of the attached terms in the
correlation theorem, and $C_\sigma$, $c_\sigma$ are constants independent of $h$ and of $Y$;
$\sigma_Y$ depends on $Y$ only through $j(Y)$. Consequently
$\sup_V|\psi^h_{i,Y}|\le n(Y)\,\sigma_Y$, where $n(Y)$ is the number of the domains $D$ meeting
$Y^{\boxplus}$ that index terms attached to $Y$. The
terms attached to a healed label $Y$ with
$j(Y)<k_0-\mathfrak{l}$ not meeting $N_i$ that
reach $\operatorname{supp}f_i$ have length at least $\rho L^\nu/(4M)$, $\nu=k_0-j(Y)$, and size at
most $(B_0\vee C_e)e^{-\kappa\rho L^\nu/(4M)}$.
\item[(e)] Let $\mathrm{LF}^h_i$ be the rest of $\Phi^h_i$:
\begin{equation}\label{8.1:eq-logarithmic-derivative-d}
\mathrm{LF}^h_i:=\Phi^h_i-c^h_i-N_{k_0}\mathcal{O}^{(k_0)}(\tilde f_i)
-\mathcal{O}^{(k_0)}(\varepsilon^h_i)-\mathrm{Irr}_i-\mathrm{Sh}^h_i .
\end{equation}
In what follows each term of $\mathbf{T}^w_h$ and of $\mathbf{B}^w$ enters $\mathrm{LF}^h_i$ taken
minus its value at $\mathbf{1}$. The function $\mathrm{LF}^h_i$ is formed from the following terms
of \eqref{8.1:eq-logarithmic-derivative-1}: the terms of the live closures of $\mathbf{T}^w_h$, the
live $\mathbf{B}^w$-terms and the terms of $\mathbf{T}^w_h$ and $\mathbf{B}^w$ coming from
characteristic functions, for live labels meeting $N_i$, of any level of creation, whose sum, so
taken, we denote by $\mathrm{LF}^{h,\mathrm{near}}_i$; the reach terms; and, for the refined index,
the terms attached to healed labels $Y$ of $h$ with $j(Y)\ge k_0-\mathfrak{l}$ (the
\emph{coarse-healed} labels). The characteristic functions just mentioned are those internal to
$\mathbf{T}^w_h$ and $\mathbf{B}^w$ in (P1); the characteristic functions of the history do not
depend on $z$ by Lemma~\ref{8.1:lem-total-covariance} and contribute nothing to $\Phi^h_i$. A
\emph{reach term} is a term anchored at a live label of $h$ not meeting $N_i$ whose support extends
to $\operatorname{supp}f_i$; let $\mathrm{LF}^{h,\mathrm{reach}}_i$ be the sum of the reach terms.
Accordingly
\begin{equation*}
\mathrm{LF}^h_i=\mathrm{LF}^{h,\mathrm{near}}_i+\mathrm{LF}^{h,\mathrm{reach}}_i
-\mathrm{LF}^{h,\mathrm{reach}}_i(\mathbf{1})+\mathrm{LF}^{h,\mathrm{heal}}_i,
\end{equation*}
where $\mathrm{LF}^{h,\mathrm{heal}}_i:=0$ on the summed index, and on the refined index
$\mathrm{LF}^{h,\mathrm{heal}}_i$ is determined by this identity. Then
$\mathrm{LF}^{h,\mathrm{near}}_i=0$ unless some live
label of $h$ meets $N_i$. Under \textup{(P3)},
\begin{equation}\label{8.1:eq-logarithmic-derivative-5}
\sup_V\bigl|\mathrm{LF}^{h,\mathrm{reach}}_i\bigr|\le C\|f_i\|_\infty\rho^4e^{-\varkappa\rho/8},
\end{equation}
with a constant $C$, the bound coming from the anchored counting by which the source lemma bounds
its far piece $\mathfrak{F}^h_i$; and every term
of $\mathrm{LF}^h_i$ is a term of the live piece $\mathfrak{L}^h_i$ of the source lemma.
\item[(f)] With these definitions,
\begin{equation}\label{8.1:eq-logarithmic-derivative-e}
\begin{aligned}
\Phi^h_i&=c^h_i+N_{k_0}\mathcal{O}^{(k_0)}(\tilde f_i)+\mathcal{O}^{(k_0)}(\varepsilon^h_i)
+\mathrm{Irr}_i+\mathrm{Sh}^h_i+\mathrm{LF}^h_i,\\
\Psi^h&=\partial_{z_1}\partial_{z_2}A^w_{k_0,h}\big|_0
+\partial_{z_1}\partial_{z_2}\log\mathbf{T}^w_h\big|_0,
\end{aligned}
\end{equation}
and in $\Psi^h$ the part coming from the first item of the history representation of
Proposition~\ref{8.1:prop-correlation-witness-depth} is zero.
\end{itemize}
\end{proposition}

\begin{proof}
Throughout, $h$ is fixed and we work on $\{\tau_h>0\}$, where $\Phi^h_i$ and $\Psi^h$ are defined
(Lemma~\ref{8.1:lem-total-covariance}).

\emph{Step 1: the first decomposition.} By (P1), $\tau^w_h$ is the term of index $h$ of the
history representation of the sourced effective density at level $k_0$ furnished by
Proposition~\ref{8.1:prop-correlation-witness-depth} with $m:=s$. On $\{\tau_h>0\}$ this term is
the product of characteristic functions of the history, which do not depend on $z$
(Lemma~\ref{8.1:lem-total-covariance}), of $\mathbf{T}^w_h$, and of the exponential of the exponent
$A^w_{k_0,h}$, whose $z$-dependent part is the sum of the plaquette terms
$\zeta_{k_0}(p)\,a_p(U_{k_0}(V))$ and of
$\mathbf{E}^w+\mathbf{D}^w+\mathbf{R}^w+\mathbf{B}^w-E^w$. Taking the logarithm, the characteristic
functions contribute terms independent of $z$, and differentiating in $z_i$ at $z=0$ gives
\eqref{8.1:eq-logarithmic-derivative-1}; the plaquette term $a_p(U_{k_0}(V))$ does not depend on $z$,
so only the running shift is differentiated in it. This is (a).

\emph{Step 2: the running shift.} (P2) is the identity \eqref{8.1:eq-logarithmic-derivative-a}, with
the four terms defined in (b), the weights $\varpi^{(j)}_y(p)$ and plaquettes $p_y$ being those of
\cite[(3.40), (3.65)]{B-CMP119} and the localisation functions $\varphi_j$ those of
\cite[(2.24)]{B-CMP119}, supplied by (P1). The term
$\delta^{\mathrm{reg}}_i$ is built from $b_j'(0)$, the weights and $f_i$ only; the history enters the
running shift only through the factors $1-\varphi_j(p)$, which occur in $\theta^h$ and
$\delta^{h,\mathrm{lab}}_i$ and not in $\delta^{\mathrm{reg}}_i$. Hence $\delta^{\mathrm{reg}}_i$ does
not depend on $h$. If $\delta^{\mathrm{reg}}_i(p)\ne0$, there are $j$ and $y$ with
$\varpi^{(j)}_y(p)\ne0$ and $f_i(p_y)\ne f_i(p)$; then $p$ or $p_y$ lies in $\operatorname{supp}f_i$,
and the distance from $p_y$ to $p$ is at most $C_{\mathrm{pr}}M$ by (P2), so
$p\in(\operatorname{supp}f_i)^{+C_{\mathrm{pr}}M}$. For every $y$ with $\varpi^{(j)}_y(p)\ne0$,
Definition~\ref{8.1:def-witness-class} gives $\operatorname{Lip}_{k_0}f_i\le\|f_i\|_\infty/(\vartheta\rho)$,
and therefore
\begin{equation*}
|f_i(p_y)-f_i(p)|\le\operatorname{Lip}_{k_0}f_i\cdot C_{\mathrm{pr}}M
\le\frac{C_{\mathrm{pr}}M\|f_i\|_\infty}{\vartheta\rho}.
\end{equation*}
Summing over $y$ with the bounds on the weights in \cite[(3.40), (3.65)]{B-CMP119}, and over the
levels $j\le k_0$ with the bound of Lemma~\ref{8.1:lem-running-normalisation} on the coefficients
$b_j'(0)$ summed over the levels below $k_0$ only, gives \eqref{8.1:eq-logarithmic-derivative-2}. It is
this partial sum below the witness level, and not the sum over all levels of the bounds $d_k$ of the
correlation theorem, that is used: the latter would carry an implicit restriction on the depth.

By \cite[(2.24)]{B-CMP119}, $1-\varphi_j(p)$ is non-zero only inside the large-field regions of the
history, that is, inside the labels of $h$; hence $\theta^h$ and $\delta^{h,\mathrm{lab}}_i$ vanish
outside the labels of $h$. The bound $|\theta^h|\le\mathfrak{d}_W(x_{k_0})$ is the bound of
Lemma~\ref{8.1:lem-running-normalisation} on partial sums of the $b_j'(0)$ over levels below $k_0$,
and the bound on $\delta^{h,\mathrm{lab}}_i$ in \eqref{8.1:eq-logarithmic-derivative-3} follows from
the same partial sums and the bound on $|f_i(p_y)-f_i(p)|$ just displayed, exactly as for
$\delta^{\mathrm{reg}}_i$.

\emph{Step 3: absorption of the profile.} By (P2) the $h$-independent profile term is Lipschitz and
is absorbed, divided by the number $N_{k_0}$, into the test function:
$\tilde f_i=f_i+\delta^{\mathrm{reg}}_i/N_{k_0}$ is Lipschitz.
Lemma~\ref{8.1:lem-running-normalisation} gives $N_{k_0}\ge\tfrac23$, so
$|\tilde f_i-f_i|=|\delta^{\mathrm{reg}}_i|/N_{k_0}\le\tfrac32|\delta^{\mathrm{reg}}_i|$, and
\eqref{8.1:eq-logarithmic-derivative-2} gives the first bound of
\eqref{8.1:eq-logarithmic-derivative-4}. Since $\theta^h$ and $\delta^{h,\mathrm{lab}}_i$ vanish
outside the labels of $h$, so does $\varepsilon^h_i$, and by \eqref{8.1:eq-logarithmic-derivative-3}
\begin{equation*}
|\varepsilon^h_i|\le|\theta^h|\,\|f_i\|_\infty+|\delta^{h,\mathrm{lab}}_i|
\le\mathfrak{d}_W(x_{k_0})\Bigl(1+\frac{C_{\mathrm{pr}}M}{\vartheta\rho}\Bigr)\|f_i\|_\infty .
\end{equation*}
Finally, inserting \eqref{8.1:eq-logarithmic-derivative-a} into the plaquette sum of
\eqref{8.1:eq-logarithmic-derivative-1},
\begin{equation*}
\begin{split}
\sum_p\partial_{z_i}\zeta_{k_0}(p)\big|_0\,a_p\bigl(U_{k_0}(V)\bigr)
&=N_{k_0}\sum_p\Bigl(f_i(p)+\frac{\delta^{\mathrm{reg}}_i(p)}{N_{k_0}}\Bigr)a_p\bigl(U_{k_0}(V)\bigr)\\
&\quad+\sum_p\bigl(\theta^h(p)f_i(p)+\delta^{h,\mathrm{lab}}_i(p)\bigr)a_p\bigl(U_{k_0}(V)\bigr),
\end{split}
\end{equation*}
which is $N_{k_0}\mathcal{O}^{(k_0)}(\tilde f_i)(V)+\mathcal{O}^{(k_0)}(\varepsilon^h_i)(V)$. This
proves (b).

\emph{Step 4: the irrelevant part.} The families $\mathbf{E}^w$, $\mathbf{D}^w$, $\mathbf{R}^w$ and the
term $E^w$ of the correlation theorem are defined on the whole lattice, with the coupling increments
extracted, and carry no history; hence $\mathrm{Irr}_i$ does not depend on $h$. By the correlation
theorem at every level $j\le k_0$ (P1), their level-$j$ remainders satisfy
$\|R^{(j)}_\zeta\|^{\flat}\le g_j^{\kappa_0}$, and the irrelevance exponent of
\cite[(0.29)]{B-CMP109}, which is the same input as the per-cell scale sums of the correlation
theorem, makes these contributions summable over $j\le k_0$. By construction
$\mathrm{Irr}_i(\mathbf{1})=0$. This proves (c).

\emph{Step 5: the shadow.} On the summed index, the regions of \cite{B-CMP122b} that carry the first
group of attached terms in the list of (d), the terms $V^{\sharp}$ and their descendants, are
integrated within the step at which they arise and so are live at that step, and these terms attach
to the large-field region live at that step, never to a healed one; a healed region is summed out and
does not index a history of the summed index. Hence no
healed label of $h$ exists on the summed index, so $\mathcal{D}_i=\emptyset$, the reach terms of
healed labels are absent, and $\mathrm{Sh}^h_i\equiv0$.

On the refined index, fix $Y\in\mathcal{D}_i$ and a domain $D$ meeting $Y^{\boxplus}$ that indexes
terms attached to $Y$. The attached terms indexed by $D$ are analytic in the sources on the domains
$\mathcal{W}_j$ of the correlation theorem taken with $m:=s$, and the Cauchy estimate on these
domains costs one factor
$\|f_i\|_\infty/r$ per derivative in $z_i$. Their size at level $j(Y)$ is bounded by the bracket in
$\sigma_Y$, with $c_1^{\flat}(g)\in\{e^{-p_0},\alpha^{1/3}\}$ as in \cite{B-CMP122b}. The oscillation
of the terms $V^{\sharp}$ of the history representation is bounded by a constant times a
power of $\check{R}_{j(Y)}$, and this factor is absorbed into $(M\check{R}_{j(Y)})^{c_\sigma}$. This
gives $\sup_V|\psi^h_{i,Y}(D)|\le\sigma_Y$, which is the first line of
\eqref{8.1:eq-logarithmic-derivative-c}. Summing over the $n(Y)$ domains $D$ meeting
$Y^{\boxplus}$ that index terms attached to $Y$,
\begin{equation*}
\sup_V|\psi^h_{i,Y}|\le\sum_D\sup_V|\psi^h_{i,Y}(D)|\le n(Y)\,\sigma_Y .
\end{equation*}

Let now $Y$ be a healed label of $h$ with $j(Y)<k_0-\mathfrak{l}$ not meeting $N_i$, so that $Y$ lies
at distance at least $\rho/4$ from $\operatorname{supp}f_i$, and consider an attached term of $Y$
reaching $\operatorname{supp}f_i$. By the decay of the level-$j(Y)$ terms in the correlation theorem,
such a term has length at least $\rho L^\nu/(4M)$, $\nu=k_0-j(Y)$, and size at most
$(B_0\vee C_e)e^{-\kappa\rho L^\nu/(4M)}$. These bounds are super-exponentially small in $\nu$, and
folding these terms into $\mathrm{Sh}^h_i$ uses no input beyond the correlation theorem. This
proves (d).

\emph{Step 6: the live part.} A term whose support does not meet $\operatorname{supp}f_i$ does not
depend on $z_i$, the source $w=\sum_iz_if_i$ entering only through $\operatorname{supp}f_i$, and
contributes nothing to $\Phi^h_i$. Hence, by (P1), the terms of
\eqref{8.1:eq-logarithmic-derivative-1} not
accounted for in Steps~2--5 are those listed in (e): the terms of live labels meeting $N_i$, the
reach terms and, on the refined index, the terms attached to coarse-healed labels; and by Step~7
below $c^h_i$ collects their
values at $\mathbf{1}$, so that each of them enters $\mathrm{LF}^h_i$ taken minus its value at
$\mathbf{1}$; this gives the displayed decomposition of $\mathrm{LF}^h_i$ in (e). On the summed index
there are no healed labels and $\mathrm{Sh}^h_i\equiv0$ (Step~5), so the terms listed in (e) are the
terms of live labels meeting $N_i$ and the reach terms only, and
$\mathrm{LF}^{h,\mathrm{heal}}_i=0$ there. The terms forming $\mathrm{LF}^{h,\mathrm{near}}_i$ are
terms of live labels meeting $N_i$; they are absent when no live label of $h$ meets $N_i$, that is,
off the first clause of the contact event $\mathcal{B}_i$, the set of histories in which some live
label meets $N_i$ or in which the clause on coarse-healed labels holds; its first clause is that
some live label of $h$ meets $N_i$. Under (P3),
the anchored counting by which the source lemma bounds its far piece $\mathfrak{F}^h_i$, applied to
the reach terms, gives
\eqref{8.1:eq-logarithmic-derivative-5}. On the summed index the reach terms enter the later
estimates paired with bounded factors, and their contribution is controlled there by the
second-derivative room lemma. By part (d) of the source lemma, every term of
$\mathrm{LF}^h_i$ is a term of its live piece $\mathfrak{L}^h_i$. This proves (e).

\emph{Step 7: the final formula.} By (a) and (b),
\begin{equation*}
\begin{split}
\Phi^h_i={}&\partial_{z_i}\log\mathbf{T}^w_h\big|_0+N_{k_0}\mathcal{O}^{(k_0)}(\tilde f_i)
+\mathcal{O}^{(k_0)}(\varepsilon^h_i)\\
&+\partial_{z_i}\bigl[\mathbf{E}^w+\mathbf{D}^w+\mathbf{R}^w-E^w\bigr]\big|_0
+\partial_{z_i}\mathbf{B}^w\big|_0 ,
\end{split}
\end{equation*}
and by \eqref{8.1:eq-logarithmic-derivative-b} the fourth term is $\mathrm{Irr}_i$ plus its value at
$\mathbf{1}$. The plaquette term $a_p$, the block-field analogue of $1-\operatorname{Re}\operatorname{tr}
U(\partial p)$, vanishes at the identity configuration, and $U_{k_0}(\mathbf{1})=\mathbf{1}$; hence the
two $\mathcal{O}^{(k_0)}$-terms vanish at $\mathbf{1}$, and
\begin{equation*}
c^h_i=\partial_{z_i}\log\mathbf{T}^w_h\big|_0(\mathbf{1})
+\partial_{z_i}\bigl[\mathbf{E}^w+\mathbf{D}^w+\mathbf{R}^w+\mathbf{B}^w-E^w\bigr]\big|_0(\mathbf{1}).
\end{equation*}
Subtracting,
\begin{equation*}
\Phi^h_i-c^h_i-N_{k_0}\mathcal{O}^{(k_0)}(\tilde f_i)-\mathcal{O}^{(k_0)}(\varepsilon^h_i)
-\mathrm{Irr}_i
\end{equation*}
is the sum of the $\mathbf{T}^w_h$- and $\mathbf{B}^w$-terms, each taken minus its
value at $\mathbf{1}$. Separating the shadow of (d), and defining $\mathrm{LF}^h_i$ as the rest, as
in (e), gives the first line of \eqref{8.1:eq-logarithmic-derivative-e}. Since the characteristic
functions do not depend on $z$, the mixed derivative of $\log\tau^w_h$ at $z=0$ is the mixed
derivative of $\log\mathbf{T}^w_h+A^w_{k_0,h}$, which is the second line of
\eqref{8.1:eq-logarithmic-derivative-e}; by the correlation theorem, the first item of its history
representation contributes zero to it. This proves (f).
\end{proof}

\begin{remark}[The shadow on the refined index]
On the refined index the bound $\sigma_Y$ of \eqref{8.1:eq-logarithmic-derivative-c} depends on $Y$
only through $j(Y)$, while the terms attached to a healed label $Y$ are indexed by the domains $D$
meeting $Y^{\boxplus}$, whose number $n(Y)$ enters the bound
$\sup_V|\psi^h_{i,Y}|\le n(Y)\,\sigma_Y$ of
Proposition~\ref{8.1:prop-logarithmic-derivative}\textup{(d)}. The number $n(Y)$ grows with the
number of cubes of level $j(Y)$ in $Y\cap N_i^{+}$, where $N_i^{+}$ is an enlargement of $N_i$. For a
domain to be renormalised, \cite{B-CMP122a} requires that it be contained in a cube of side
$100M\check{R}_k$ of its level $k$, $\check{R}_k$ being Ba{\l}aban's radius function at level $k$;
this requirement does not bound the number of domains meeting
$Y^{\boxplus}$. The growth of $n(Y)$ is harmless for the sums
\begin{equation*}
\sum_{Y\ni\mathfrak{c}}\tilde\nu(\{Y\in h\})
\end{equation*}
over healed labels $Y$ containing a fixed cell $\mathfrak{c}$, since these events are disjoint for
distinct $Y$, but not for sums of the form
\begin{equation*}
\sum_Y n(Y)\,\sigma_Y\,\tilde\nu(\{Y\in h\})\,(\cdots)
\end{equation*}
weighted by further factors. For these sums the shadow may be indexed by pairs consisting of a level
and a cube of that level, each such index carrying $\sigma_Y$; or $n(Y)$ may be carried in the bound,
as in (d), with its moments taken from the size clause of the large-field block bound, healed part.
In either form the resulting sums over $Y$ converge without any additional power of $L$, by the
birth-level rarity weight $\tilde q_j$. On the summed index every component of the class of domains
renormalised within a single cube of their level is a single cube, and $\mathcal{D}_i=\emptyset$ in
any case.
\end{remark}

\begin{lemma}[Locality of the mixed second derivative]\label{8.1:lem-mixed-derivative-locality}
Let $(f_1,f_2)$ be a pair of test functions in the setting of
Proposition~\ref{8.1:prop-correlation-witness-depth}, in particular with disjoint supports, of
radius $\rho$ in cells and separated as there; let $w=z_1f_1+z_2f_2$ be the complex source,
and write $\partial_i=\partial/\partial z_i$ and $|_0$ for evaluation at $z_1=z_2=0$. Work at the
witness level $k_0$, let $0\le j\le k_0$ and $\nu=k_0-j$, and let $\ell_\nu$ be the length defined
in Lemma~\ref{8.1:lem-spanning-split}. Assume:
\begin{itemize}
\item[(a)] the correlation theorem at the witness depth,
Proposition~\ref{8.1:prop-correlation-witness-depth}, with its analyticity domains
$\mathcal{W}_j(X)$, each containing the polydisc $|z_i|<r/(2\|f_i\|_\sharp)$, $i=1,2$, where
$r>0$ is the radius of that proposition;
\item[(b)] the source lemma (live marked slots and annealed shadow).
\end{itemize}
Then the following hold.
\begin{itemize}
\item[(i)] For every plaquette $p$, the running shift $\zeta_{k_0}(p)$ is, by
Proposition~\ref{8.1:prop-correlation-witness-depth}, a sum over indices $(j,y)$, $0\le j\le k_0$,
of level-$j$ terms $\mathfrak{i}^{(j)}_y$, each of which depends on the source only through the
single value
\begin{equation*}
w(x(p_y))=z_1f_1(x(p_y))+z_2f_2(x(p_y))
\end{equation*}
at one plaquette $p_y$. For every $j\le k_0$ and every index $y$,
$\partial_1\partial_2\mathfrak{i}^{(j)}_y|_0$ is a multiple of $f_1(x(p_y))f_2(x(p_y))$, hence
$\partial_1\partial_2\mathfrak{i}^{(j)}_y|_0=0$; consequently
$\partial_1\partial_2\zeta_{k_0}(p)|_0=0$.
\item[(ii)] Let $T$ be a local term at level $j$ localised in $X$, and let $X^{+}$ be the
enlargement of $X$, as in Lemma~\ref{8.1:lem-spanning-split}. If $T$ depends on both $z_1$ and
$z_2$, then $X^{+}$ meets both supports, and therefore $d_j(X)\ge\ell_\nu$. Moreover
\begin{equation}\label{8.1:eq-mixed-derivative-locality-a}
\bigl|(\partial_1\partial_2T)(0,0)\bigr|
\le\sup_{\mathcal{W}_j(X)}|T|\cdot\frac{4\|f_1\|_\sharp\|f_2\|_\sharp}{r^2}.
\end{equation}
\item[(iii)] For a history $h$, with live exponent $\mathfrak{A}^w_h$ and its normalised law
$\mathbb{E}_h$,
\begin{equation}\label{8.1:eq-mixed-derivative-locality-b}
\begin{split}
&\operatorname{Cov}_{\mathbb{E}_h}\bigl(\partial_{z_1}\mathfrak{A}^w_h|_0,\,
\partial_{z_2}\mathfrak{A}^w_h|_0\bigr)=0\\
&\qquad\text{unless the live labels of $h$, one label or two, reach both $S_1^{+}$ and
$S_2^{+}$,}
\end{split}
\end{equation}
where $S_1^{+}$ and $S_2^{+}$ are the sets that the source lemma attaches to the supports of
$f_1$ and $f_2$.
\end{itemize}
\end{lemma}

\begin{proof}
(i) By hypothesis (a), each level-$j$ term $\mathfrak{i}^{(j)}_y$ depends on the source $w$ only
through the single value $w(x(p_y))=z_1f_1(x(p_y))+z_2f_2(x(p_y))$. By the chain rule, its mixed
derivative $\partial_1\partial_2\mathfrak{i}^{(j)}_y|_0$ is the second derivative of
$\mathfrak{i}^{(j)}_y$ with respect to that value, at $0$, multiplied by
$f_1(x(p_y))f_2(x(p_y))$. The supports of $f_1$ and $f_2$ are disjoint, so at the point
$x(p_y)$ at least one factor vanishes, and $\partial_1\partial_2\mathfrak{i}^{(j)}_y|_0=0$.
Summing over the indices $(j,y)$ of the decomposition of $\zeta_{k_0}(p)$ gives
$\partial_1\partial_2\zeta_{k_0}(p)|_0=0$.

(ii) By hypothesis (a), a local term $T$ at level $j$ localised in $X$ depends on the sources only
through the restrictions of $f_1$ and $f_2$ to $X^{+}$. If $X^{+}$ misses the support of $f_i$,
then $T$ does not depend on $z_i$, and so $\partial_1\partial_2T|_0=0$. Consequently a term that
depends on both sources has $X^{+}$ meeting both supports. By Lemma~\ref{8.1:lem-spanning-split},
a level-$j$ set $X$ whose enlargement $X^{+}$ meets both supports has $d_j(X)\ge\ell_\nu$, where
$\nu=k_0-j$; it is the separation of the supports, assumed in the setting of
Proposition~\ref{8.1:prop-correlation-witness-depth}, that this lemma converts into the lower
bound $d_j(X)\ge\ell_\nu$.

For the bound, recall that by hypothesis (a), $T$ is analytic in $(z_1,z_2)$ on
$\mathcal{W}_j(X)$, and this domain contains the polydisc $|z_i|<r/(2\|f_i\|_\sharp)$, $i=1,2$.
Fix a number $\xi$ with $0<\xi<1$. The two-variable Cauchy formula on the product of the circles
$|z_i|=\xi r/(2\|f_i\|_\sharp)$, $i=1,2$, gives
\begin{equation*}
\partial_1\partial_2T|_0
=\frac{1}{(2\pi\mathrm{i})^2}\oint_{|z_1|=\xi r/(2\|f_1\|_\sharp)}
\oint_{|z_2|=\xi r/(2\|f_2\|_\sharp)}\frac{T(z_1,z_2)}{z_1^2z_2^2}\,dz_2\,dz_1 .
\end{equation*}
Estimating the integrand by its supremum gives
\begin{equation*}
\bigl|(\partial_1\partial_2T)(0,0)\bigr|
\le\frac{4\|f_1\|_\sharp\|f_2\|_\sharp}{\xi^2r^2}\sup_{\mathcal{W}_j(X)}|T|.
\end{equation*}
Letting $\xi\uparrow1$ yields \eqref{8.1:eq-mixed-derivative-locality-a}.

(iii) Here we use hypothesis (b). The source lemma describes how the live exponent
$\mathfrak{A}^w_h$ depends on $z_1$ and $z_2$ through the live labels of $h$; this is the content of
the first and fourth clauses of its description of the form of $\mathfrak{A}^w_h$, according to
which $\partial_{z_i}\mathfrak{A}^w_h|_0$ receives contributions only from live labels of $h$ that
reach $S_i^{+}$, $i=1,2$. Hence, if no live label of $h$ reaches $S_1^{+}$, then
$\partial_{z_1}\mathfrak{A}^w_h|_0=0$ and the covariance under $\mathbb{E}_h$ vanishes; likewise,
if no live label of $h$ reaches $S_2^{+}$, then $\partial_{z_2}\mathfrak{A}^w_h|_0=0$ and the
covariance vanishes. So the covariance is non-zero only if the live labels, one label or two, reach
both $S_1^{+}$ and $S_2^{+}$, which is \eqref{8.1:eq-mixed-derivative-locality-b}.
\end{proof}

\begin{definition}[The position average and perturbed profiles]\label{8.1:not-position-average}
Fix the depth $s$ and the witness level $k_0=K-s$. Put
\begin{equation*}
S_0:=L^2M\check{R}_{K-s},
\end{equation*}
where $M=L^{m'}$ is the auxiliary parameter of the glossary (one of the parameters $\mathfrak{p}$)
and $\check{R}_j$ is Ba{\l}aban's radius function. Let $\mathfrak{P}_s$ be the sublattice
$S_0\mathbb{Z}^4$ taken modulo the periods of the torus. For a quantity $G$ and
$v\in\mathfrak{P}_s$, $G^v$ denotes the translate of $G$ by $v$. The \emph{position average} of
$G$ is
\begin{equation}\label{8.1:eq-position-average-1}
\langle\!\langle G\rangle\!\rangle:=|\mathfrak{P}_s|^{-1}\sum_{v\in\mathfrak{P}_s}G^v .
\end{equation}
Let $C_h$ denote the constant that bounds the perturbation of the profiles $\tilde{f}_i$ entering
the first logarithmic derivative \eqref{8.1:eq-logarithmic-derivative-a}; it enters the floor
$C_0^{\mathrm{prof}}$ among the constants of \eqref{8.1:eq-constants-ceiling-c}. The class of
\emph{perturbed profiles} is
\begin{equation}\label{8.1:eq-position-average-2}
\begin{split}
\tilde{\mathfrak{F}}(\mathsf{d},\vartheta):=\Bigl\{f+\delta:\ & f\in\mathfrak{F}^{+}(\mathsf{d},\vartheta),\ \delta\ \text{Lipschitz},\
\operatorname{supp}\delta\subset(\operatorname{supp}f)^{+C_hM},\\
&|\delta|\le C_hM\,\|f\|_\infty/(\vartheta\rho)\Bigr\}.
\end{split}
\end{equation}
Here $\mathfrak{F}^{+}(\mathsf{d},\vartheta)$ is the class of
Definition~\ref{8.1:def-witness-class}. No bound is imposed on the Lipschitz constant of $\delta$;
that constant is kept symbolic. In this chapter the Gaussian extraction is applied to profiles of
the class \eqref{8.1:eq-position-average-2}; its hypotheses involve no bound on the Lipschitz
constant of the profile, which is why that constant of $\delta$ is left free.

By Proposition~\ref{8.1:prop-logarithmic-derivative}, the profiles $\tilde{f}_i$ that enter the
first logarithmic derivative \eqref{8.1:eq-logarithmic-derivative-a} have the form
$f_i+\delta_i$, with $(f_1,f_2)$ the witness pair and $\delta_i$ subject to the conditions in
\eqref{8.1:eq-position-average-2}; they therefore lie in $\tilde{\mathfrak{F}}(\mathsf{d},\vartheta)$.
Assume the floor
\begin{equation}\label{8.1:eq-position-average-3}
\rho\ge C_0^{\mathrm{prof}}\ge 160\,C_hM/\vartheta^2
\end{equation}
of \eqref{8.1:eq-constants-ceiling-c}. Then, for $\tilde{f}=f+\delta\in\tilde{\mathfrak{F}}(\mathsf{d},\vartheta)$, one has $\mathfrak{n}(\tilde{f})\in[\tfrac12,2]\,\mathfrak{n}(f)$. For a witness pair perturbed in this way, the supports of $\tilde{f}_1$ and $\tilde{f}_2$ are at distance at least $\rho/2$.
\end{definition}

\begin{proof}
\emph{A bound on $\vartheta$.} Let $f\in\mathfrak{F}^{+}(\mathsf{d},\vartheta)$. By
Definition~\ref{8.1:def-witness-class}, $f$ is Lipschitz, $f\not\equiv0$, $f$ vanishes outside a
ball of radius $\rho$ cells, and $\operatorname{Lip}_{k_0}f\le\|f\|_\infty/(\vartheta\rho)$. Being
Lipschitz on the compact torus, $f$ attains $\|f\|_\infty$ at a point $\xi$; since
$f(\xi)=\|f\|_\infty>0$, the point $\xi$ lies in that ball. Let $\xi'$ be the point at distance
$3\rho/2$ from the centre of the ball on a ray from the centre through $\xi$ (any ray if $\xi$ is
the centre). Then $\xi'$ lies outside the ball, so $f(\xi')=0$, and $|\xi-\xi'|\le2\rho$. Hence
\begin{equation*}
\|f\|_\infty=f(\xi)-f(\xi')\le\operatorname{Lip}_{k_0}f\cdot2\rho\le2\|f\|_\infty/\vartheta .
\end{equation*}
Since $\|f\|_\infty>0$, we obtain $\vartheta\le2$. With \eqref{8.1:eq-position-average-3} this gives
\begin{equation*}
C_hM\le\vartheta^2\rho/160\le\rho/40,\qquad C_hM/\vartheta\le\vartheta\rho/160 .
\end{equation*}

\emph{The mass.} The mass $\mathfrak{n}$ of Definition~\ref{8.1:def-witness-class} is linear in its
argument, so
\begin{equation*}
\mathfrak{n}(\tilde{f})=\mathfrak{n}(f)+\mathfrak{n}(\delta).
\end{equation*}
By Definition~\ref{8.1:def-witness-class}, $\operatorname{supp}f$ lies in a ball of radius $\rho$
cells. Hence $\operatorname{supp}\delta$ lies in the ball of radius $\rho+C_hM$ concentric with
that support ball.

We use the count of bare sites on which the bound $\mathfrak{n}(f)\le 10\rho^4\|f\|_\infty$ of
Definition~\ref{8.1:def-witness-class} rests: a ball of radius $r\ge\rho$ cells contains at most
$10r^4\eta^{-4}$ bare sites. At $r=\rho+C_hM$ it bounds $\eta^4$ times the number of bare sites in
the ball of radius $\rho+C_hM$ by $10(\rho+C_hM)^4$.
With the pointwise bound on $\delta$ this yields
\begin{equation*}
|\mathfrak{n}(\delta)|\le 10(\rho+C_hM)^4\,\frac{C_hM\,\|f\|_\infty}{\vartheta\rho}.
\end{equation*}
Since $C_hM\le\rho/40$,
\begin{equation*}
(\rho+C_hM)^4\le(41/40)^4\rho^4\le\tfrac{9}{8}\,\rho^4 .
\end{equation*}
Using $C_hM/\vartheta\le\vartheta\rho/160$ and the lower bound
$\vartheta\|f\|_\infty\rho^4\le\mathfrak{n}(f)$ of Definition~\ref{8.1:def-witness-class}, we obtain
\begin{equation*}
|\mathfrak{n}(\delta)|\le \frac{45}{4}\,\rho^3\|f\|_\infty\,\frac{\vartheta\rho}{160}
=\frac{9}{128}\,\vartheta\rho^4\|f\|_\infty\le\frac{9}{128}\,\mathfrak{n}(f).
\end{equation*}
Hence $\mathfrak{n}(\tilde{f})\in[\tfrac{119}{128},\tfrac{137}{128}]\,\mathfrak{n}(f)\subset[\tfrac12,2]\,\mathfrak{n}(f)$.

\emph{The distance.} Write $\tilde{f}_i=f_i+\delta_i$ with $f_i$ and $\delta_i$ as in
\eqref{8.1:eq-position-average-2}. Since $\operatorname{supp}\tilde{f}_i\subset(\operatorname{supp}f_i)^{+C_hM}$
by \eqref{8.1:eq-position-average-2}, the triangle inequality gives
\begin{equation*}
\operatorname{dist}(\operatorname{supp}\tilde{f}_1,\operatorname{supp}\tilde{f}_2)
\ge\operatorname{dist}(\operatorname{supp}f_1,\operatorname{supp}f_2)-2C_hM .
\end{equation*}
For a witness pair of Definition~\ref{8.1:def-witness-class}, points of the two supports are at
distance at least $\rho$ cells. By the bound on $\vartheta$ above, applied to $f_1$,
\begin{equation*}
2C_hM\le\vartheta^2\rho/80\le\rho/20\le\rho/2 .
\end{equation*}
The right-hand side of the distance inequality is therefore at least $\rho/2$.
\end{proof}

\begin{remark}[The error terms of the Gaussian extraction at the witness scale]
\label{8.1:rem-extraction-errors}
We work under the hypotheses of clause (a) of the Gaussian extraction: $\rho\in[C_W(s),\,LC_W(s)[$,
$s-r\ge s_*$, $2B^{\sharp}(r+1)\le x_{k_0}$ and $r=r(s)$, with the two profiles taken in
the class $\tilde{\mathfrak{F}}(\mathsf{d},\vartheta)$ of
Definition~\ref{8.1:not-position-average} as perturbations of a pair of the witness class
$\mathfrak{T}(\mathsf{d};\vartheta)$. We recall from clause (a) of the Gaussian extraction
the bound on its error $\mathcal{E}_a$:
\begin{equation}\label{8.1:eq-extraction-errors-1}
\begin{split}
|\mathcal{E}_a|\le C_a\,g_{k_0}^4\,\mathfrak{n}(\tilde f_1)\,\mathfrak{n}(\tilde f_2)\Bigl\{
&\Bigl[g_{k_0}^{\sigma_1}\Pi_r+g_{k_0}^2(1+r)\\
&\quad+D_{\square}^4(1+r)\rho^4p_0(g_{k_0})^4e^{-c_ap_0(g_{k_0})}
+p_0^4e^{-c_aD_{\square}}\Bigr]\\
&+\rho^{-8}\Bigl[L^{-4(r-r_1)}+L^{-8(r-r_1)}p_0(g_{k_0+r})^4\Bigr]\Bigr\},
\end{split}
\end{equation}
with $\Pi_r:=rD_{\square}^4L^4p_0^2$ and with $r_1$ as in the Gaussian extraction. Each
member of
\eqref{8.1:eq-extraction-errors-1} is produced by one or two of the inputs of the proof of
clause (a) listed below, and at each input that produces a member we say which member it
produces. In particular, at the witness scale the product of the two
observables is bounded by the regularity lemma that accompanies the source lemma through
the bound of
order $\eta^2$ alone, $\eta$ being the bare spacing $L^{-k_0}$ in witness units, and no
large-field operation enters that bound.
\end{remark}

\begin{proof}
The proof of clause (a) of the Gaussian extraction uses the following inputs.
\begin{enumerate}
\item By the position-average identity, a consequence of the translation invariance of the
lattice measure (Lemma~\ref{1.1:lem-invariance}), the connected two-point function equals
its position average in the sense of Definition~\ref{8.1:not-position-average}:
\begin{equation*}
S^T_{2;K,m}(f_1,f_2)=\langle\!\langle S^T_{2;K,m}(f_1^v,f_2^v)\rangle\!\rangle .
\end{equation*}
Every bound of the assembly of the two-point witness theorem, in particular every bound
entering below, holds uniformly in the position $v\in\mathfrak{P}_s$, and
therefore passes to the average.
\item Clause (ii-a) of Theorem~\ref{3.1:thm-main} is used, in the segment form in which
the proof of the Gaussian extraction applies it.
\item The complement of the event singled out as clean in the proof of the Gaussian
extraction has weight at most
$C D_{\square}^4(1+r)\rho^4e^{-cp_0(g_{k_0})}$, where $C$ and $c$ denote positive
constants. This bound follows from the
one-point form of the large-field block bound (whose live part is discharged according to
the age of the live regions: for young regions by the corollary on the rarity of young
large-field blocks, and for old healed regions by the healed old-region bound, while old
never-healed regions are accounted for by the finite-volume closure),
together with an auxiliary estimate of the two-point witness construction,
and clause (0) of Theorem~\ref{3.1:thm-main}. For the never-healed regions the bound is
therefore used under the per-torus depth floor of the finite-volume closure.
Multiplied by the suprema of item 8, it yields the
third member of \eqref{8.1:eq-extraction-errors-1}.
\item The running of the coupling over the $r$ levels, under the hypothesis
$2B^{\sharp}(r+1)\le x_{k_0}$, yields the member $g_{k_0}^2(1+r)$.
\item The tail of the tower of covariance steps yields the member with the factor
$L^{-4(r-r_1)}$.
\item The background member is estimated through the window condition
of~\cite[(2.17)]{B-CMP119} at level $k_0+r$. It yields the member with the factor
$L^{-8(r-r_1)}p_0(g_{k_0+r})^4$.
\item The comparison of the covariances under the measure restricted to the small-field
window with the covariances of
its Gaussian part yields the member $g_{k_0}^{\sigma_1}\Pi_r$ with $\sigma_1=1$.
\item On the complement of the clean event, the weight of item 3 is multiplied by the
suprema of the
two observables. Write $\Xi_i:=\mathcal{O}^{(k_0)}(\tilde f_i)^v$ for $i=1,2$, where
$\tilde f_i=f_i+\delta_i$ as in Definition~\ref{8.1:not-position-average}, so that, under
the radius condition of that definition, $\mathfrak{n}(\tilde f_i)\in[\tfrac12,2]\,
\mathfrak{n}(f_i)$. The product of the two obeys
\begin{equation}\label{8.1:eq-extraction-errors-2}
|\Xi_1\Xi_2|\le \frac{C_X^2\,\mathfrak{n}_1\mathfrak{n}_2}{\vartheta^2},
\qquad \mathfrak{n}_i:=\mathfrak{n}(f_i),
\end{equation}
by the regularity lemma at the witness scale, in the form of the corollary of the source
lemma that rests on the bound of order $\eta^2$ alone. Indeed,
\cite[Theorem~1]{B-CMP102b} applied at the top
scale gives
\begin{equation*}
0\le \mathcal{O}^{(k_0)}(f_i)(V)\le \frac{80\,C_{\square}B_3^2a_1^2}{\vartheta}\,
\mathfrak{n}(f_i)
\end{equation*}
wherever the background configuration $U_{k_0}(V)$ of
Definition~\ref{1.5:def-domains} is defined near the support $S_i$ of the $i$-th profile,
with $B_3$ and $a_1$ constants. Off that
domain the observable is fixed by the convention of the two-point witness theorem.
Applied with the translated perturbed profile $\tilde f_i^{\,v}$ in place of $f_i$, and
combined with $\mathfrak{n}(\tilde f_i)\le 2\,\mathfrak{n}(f_i)$, the mass being unchanged
under translation, this bound gives $|\Xi_i|$ at most a numerical multiple of
$C_{\square}B_3^2a_1^2\,\mathfrak{n}_i/\vartheta$; multiplying the two bounds yields
\eqref{8.1:eq-extraction-errors-2}, the constant $C_X$ absorbing the numerical factor. No
large-field operation is used in \eqref{8.1:eq-extraction-errors-2}.
\item The locality input of the proof yields the member
$p_0^4e^{-c_aD_{\square}}$.
\end{enumerate}
Items 3--9 account for every member of \eqref{8.1:eq-extraction-errors-1}.
By item 1, the resulting bound holds for the position average.
\end{proof}

\begin{theorem}[The two-point witness theorem]\label{8.1:thm-two-point-witness}
The gauge group is $\mathrm{SU}(2)$; the formulas
below are written for $\mathrm{SU}(N)$. Work in the quantifier prefix of Theorem~\ref{3.1:thm-main}
(Notation~\ref{2.1:not-prefix}), inside the regime of the correlation theorem, that is, for parameter
values at which that theorem holds, with the coupling threshold $\gamma_J$ that it provides for
those values; the reference coupling is restricted below to $g\le\gamma_W\le\gamma_J$. Fix a
shape constant $\vartheta\in(0,1]$.

Assume that the following statements hold:
\begin{enumerate}
\item[(1)] the correlation theorem, applied below at depth $s=s(\mathsf{d})$ below the cutoff;
\item[(2)] the Gaussian extraction, in each of its three parts;
\item[(3)] the large-field block bound, healed part, for healed regions (those whose large-field
factor has been absorbed back into the small-field format by a later step), and, for regions still
carrying their large-field factor at the level read, the following statements: the corollary on
rarity of young large-field blocks, for young regions; the healed old-region bound, for old regions
healed at a later step; the proposition on never-healed old regions, which bounds at each age the
weight of the old regions not healed by the last lattice of the run, for those regions; and the
additive theorem, which bounds the weight of the event that such regions meet the neighbourhoods of
both supports of Definition~\ref{8.1:def-witness-class} by the smaller of the cell counts of the
two neighbourhoods times the one-cell weight of such regions, so that this weight enters linearly,
adds to this the series over their ages, and so combines the preceding bounds into the bound on the
contribution of such regions used in the proof. Here the \emph{age} at level $k$ of a large-field
region created at step $j$ is $k-j$, the number of steps since the step that created it, and a
region is \emph{young} if its age is at most the level count $\mathfrak{l}$ of the source lemma and
\emph{old} if its age exceeds it;
\item[(4)] the source lemma;
\item[(5)] the continuum kernel;
\item[(6)] the further inputs:
\begin{itemize}
\item clauses (0) and (ii-a) of Theorem~\ref{3.1:thm-main};
\item the envelope \eqref{8.1:eq-coupling-envelope-f} of the running inverse coupling;
\item the operator theorem at free current;
\item the estimate, together with its linearised form, on which, with the operator theorem at
free current, the upper bound in Lemma~\ref{8.1:lem-kernel-bounds} rests;
\item the statement from which the lower bound in Lemma~\ref{8.1:lem-kernel-bounds} is derived,
in the form used in this chapter;
\item the positivity dichotomy of the large-field operations, on which the union bounds over
single cells in the proof rest;
\item the remaining statements of the setting of Theorem~\ref{3.1:thm-main} that the proof uses;
\item the hypotheses of Lemma~\ref{8.1:lem-kernel-bounds} and of
Lemma~\ref{8.1:lem-running-normalisation}.
\end{itemize}
\end{enumerate}

Then there exist constants $C_0$, $r_0$, $C_D$, $k_W$ and $\gamma_W\le\gamma_J$ depending only on
$L$, $M$, $M_1$, the auxiliary parameters $\mathfrak{p}$ (among them $\kappa$), $N$, $\vartheta$,
and the constants of the statements assumed in (1)--(6), among them $\kappa_0$, $E_0$, $B_0$,
$B_3$, $\alpha_0$ and $r$, with the following property; $C_0$, $r_0$, $C_D$ and $k_W$ are the
constants of Definition~\ref{8.1:def-coupling-envelope}, where $C_0$ enters the room function
$C_W(s)$, $r_0$ is the constant depth offset of the witness construction and $k_W$ is the level
floor in the conditions below, and their values are fixed in the proof. Let:
\begin{itemize}
\item $g\in(0,\gamma_W]$;
\item $m$ be any torus exponent admitted by the torus clause of Theorem~\ref{3.1:thm-main},
with no further
condition;
\item $g_0\in(0,g_-(g)]$, and $K=K(g_0,g)$ the first passage;
\item $\mathsf{d}$ satisfy the one-sided conditions $s(\mathsf{d})\ge s_W(g)$ and
$K-s(\mathsf{d})\ge k_W$ and, when $m>m_V(g)$, the one-sided per-torus depth floor
$s(\mathsf{d})\ge s_{\mathrm{vol}}(g,m)$ of Definition~\ref{8.18:def-volume-horizon};
\item $(f_1,f_2)$ be any witness pair in
$\mathfrak{F}^{+}(\mathsf{d},\vartheta)\times\mathfrak{F}^{+}(\mathsf{d},\vartheta)$, in the sense of
Definition~\ref{8.1:def-witness-class}.
\end{itemize}
Here $s_W(g)$ is the $g$-dependent depth floor of Definition~\ref{8.1:def-coupling-envelope}, and
the first two conditions on $\mathsf{d}$ are those displayed in \eqref{8.1:eq-coupling-envelope-d};
$m_V(g)$ denotes the threshold on the torus exponent up to which the per-torus depth floor is not
needed. Write $s:=s(\mathsf{d})$. Then
\begin{equation}\label{8.1:eq-two-point-witness-1}
S^T_{2;K,m}(f_1,f_2)=N_{K-s}^{2}\, b_0\, g_{K-s}^{4}\,
\mathcal{K}_{K,s}(f_1,f_2)\,\bigl(1+R_K\bigr),
\end{equation}
with
\begin{equation}\label{8.1:eq-two-point-witness-2}
|R_K|\le\tfrac12,\qquad\text{indeed}\qquad |R_K|\le\tfrac14 .
\end{equation}
Here $b_0=(N^2-1)/2$, and $\mathcal{K}_{K,s}$ is the two-point kernel of
Definition~\ref{8.1:def-background-kernel}. The normalisation
$N_{K-s}$ of Lemma~\ref{8.1:lem-running-normalisation} (see
\eqref{8.1:eq-running-normalisation-a}) satisfies
\begin{equation}\label{8.1:eq-two-point-witness-3}
N_{K-s}\in\bigl[\,1-\mathfrak{d}_W(\underline{x}_s),\,1+\mathfrak{d}_W(\underline{x}_s)\,\bigr]
\subset\bigl[\tfrac23,\tfrac43\bigr],
\end{equation}
where $\underline{x}_s$ is the lower envelope of Definition~\ref{8.1:def-coupling-envelope}.
The constants $C_0$, $r_0$, $C_D$, $k_W$, $\gamma_W$ are independent of $K$, $m$, the parameter
$N_T$ of the correlation theorem, $g_0$ and of the position of the supports of $f_1$ and $f_2$.
\end{theorem}

Part (b) of the Gaussian extraction is stated; its proof is incomplete at the step marked
$(\ast)$ in that proposition. For healed
regions, hypothesis (3) is the large-field block bound, healed part, which is proved as an
implication from the statements listed in its own statement. For regions still carrying their
large-field factor at the level read, the statements named in hypothesis (3) stand as follows.
For young regions the bound is supplied by the corollary on rarity of young large-field blocks,
which is proved as an implication from the statements named in it. For old regions healed at a
later step it is supplied by the healed old-region bound, likewise proved as an implication from
the statements named in it. For old regions that are never healed it is supplied by the
proposition on never-healed old regions, proved as an implication from the statements named in
it, one of which, an input shared with clause (v) of Theorem~\ref{3.1:thm-main}, has its proof
outstanding; that
proposition bounds the weight of such regions at each age, and the resulting series in the age is
summed under the conditions on $\mathsf{d}$ listed in Theorem~\ref{8.1:thm-two-point-witness},
among them the per-torus depth floor $s\ge s_{\mathrm{vol}}(g,m)$ when $m>m_V(g)$. The annealed
live-tail statement, which is not proved, is not used.

The proof of Theorem~\ref{8.1:thm-two-point-witness} occupies the remainder of this section.

\begin{remark}[Healed components are summed out of the history index]\label{8.1:rem-healed-summed-out}
Let $k_0=K-s$ be the witness level. By the \emph{index of the correlation theorem} we mean the
index set of the decomposition of the sourced effective density $\rho^w_{k_0}$ furnished by the
correlation theorem at level $k_0$. Its elements $h$ are the \emph{histories}, and the large-field
regions recorded by a history $h$ are called its \emph{labels}. The
decomposition $(\tau^w_h)_h=\mathsf{P}_{k_0}(\rho^w_0)$ of
Lemma~\ref{8.1:lem-history-decomposition} is also available on a finer index, obtained by
refining the one just described without resumming; on that index no label, once created, is ever
discarded. We call it the \emph{finer index}.
\begin{enumerate}
\item[(a)] The index of the correlation theorem is a partial sum of the finer index: for each
history $h$ of it, $\tau^w_h$ is the sum of those
terms of the finer index that the operation $\mathbf{R}$ resums into it. Both decompositions are
exact. The following statements hold on the index of the correlation theorem exactly as stated:
the linear positive history
decomposition (Lemma~\ref{8.1:lem-history-decomposition}), the total covariance formula
(Lemma~\ref{8.1:lem-total-covariance}), the structure of the first logarithmic derivative
(Proposition~\ref{8.1:prop-logarithmic-derivative}, in particular
\eqref{8.1:eq-logarithmic-derivative-c}) and the locality of the mixed second derivative
(Lemma~\ref{8.1:lem-mixed-derivative-locality}).
\item[(b)] Every label healed by $\mathbf{R}$, whatever its age, is summed out at the step of
$\mathbf{R}$ at which it is healed, and is therefore absent from every history $h$. The labels of
the histories are the regions live at level $k_0$, of any birth level.
\item[(c)] We use the notation of the decomposition of the main term and of its remainders set
out below. Let $\mathfrak{l}=L\check{R}_{k_0}+c_L$ be the level count used in this chapter, and
let $\mathcal{D}_i$ be the
family of healed labels of $h$ that lie under $N_i$ and have levels below $k_0-\mathfrak{l}$ (the
\emph{deep} labels). Then:
\begin{itemize}
\item $\mathcal{D}_i=\emptyset$;
\item the deep pieces $\mathrm{Sh}^h_i$, $\mathring{\mathfrak{S}}_i$ and
$\mathfrak{S}^{\flat,h}_i$ vanish identically;
\item $m_i(h)$ is identically equal to $c^0_i$, its value when no deep label is present
under $N_i$;
\item $\theta^h$ and $\delta^{h,\mathrm{lab}}_i$ are supported in the live regions of $h$, and
they vanish unless $h$ lies in the contact event $\mathcal{B}_i$;
\item the remainders $R^{m\mathfrak{L}}$ and $(\mathrm{sh})'$ vanish.
\end{itemize}
\item[(d)] The residue of the healed components consists of:
\begin{itemize}
\item the whole-lattice covariant family $R^{(j)}$, which enters $\mathrm{Irr}_i$ and the mixed
second derivative of Lemma~\ref{8.1:lem-mixed-derivative-locality}; it is of order
$g_j^{\kappa_0}$ uniformly in the history, so that the supremum over histories may be taken in
its bound;
\item contributions to the coupling-constant terms $b_j$ and to the normalisation $N_{k_0}$.
\end{itemize}
\item[(e)] On the index of the correlation theorem the flat shell bound (stated; not proved) is
not used. Nor are the following used: the spanning
event $B_{\times}$; the quantities $\mu^{\sharp}_Y$ and $\chi_Y$; the conditional and healed
clauses of the large-field block bound, healed part; and the condition on $\gamma$ that
accompanies them. The conditions (M1)--(M3) of the source lemma are not required; they may be
imposed or omitted.
\end{enumerate}
On the finer index the situation is different. Deep healed labels are present under $N_i$ in
histories of the no-contact event $\mathcal{B}_i^{\mathrm{c}}$, and their number is not bounded
(by part (a) of the counting
lemma of this chapter). The flat shell bound, parts (a) and (b), is then an input, and the
following enter as they are displayed in this chapter: the remainders $R^{m\mathfrak{L}}$ and
$(\mathrm{sh})'$, the accompanying condition on $\gamma$, the conditional and healed clauses of
the large-field block bound, healed part, the conditions (M1)--(M3) of the source lemma, and the
coarse-healed clause of $\mathcal{B}_i$. The
statement that a history of $\mathcal{B}_i^{\mathrm{c}}$ carries no deep healed label under $N_i$
is a
property of the index of the correlation theorem, the index produced by $\mathbf{R}$. On the finer
index it fails, and the flat
shell bound takes its place. Here $\mathcal{C}$ is the $\sigma$-algebra generated by $V_{k_0}$
and the labels of $h$ of levels at least $k_0-\mathfrak{l}$. A proof on the finer index would
need the following four points, displayed here without proof.
\begin{itemize}
\item[(1)] \emph{Functional form.} Let $F_1,F_2$ be bounded, with $F_i$ measurable with respect
to the configuration of deep labels under $N_i^{++}$ and to the block field there; for the member
carrying $\mu^{\sharp}$ the block field is taken on the whole lattice. Then
$|\operatorname{Cov}_{\tilde{\nu}}(F_1,F_2)|$ is bounded by the forms of parts (a) and (b) of the
flat shell bound summed
over the atoms, with the conditional supremum-expectation $E^{\natural}$. The function $m_i(h)$ is
such an $F_i$.
\item[(2)] \emph{Growth of $\sigma_Y$ in $|Y|$.} This growth is controlled in one of two ways.
The first is to index by cubes. For a cube $\mathfrak{c}$ of level $j$, let
$\psi_{i,(j,\mathfrak{c})}$ be the attached terms anchored in $\mathfrak{c}$, let
$\sigma_{(j,\mathfrak{c})}$ be the quantity $\sigma_Y$ formed for the cube $\mathfrak{c}$ in
place of $Y$, and take as presence indicator the
indicator that $\mathfrak{c}$ lies in a healed label of level $j$ of $h$. The second is to carry
the number of cubes of $Y$ in $\sigma_Y$ and to use the size clause of the large-field block
bound, healed part, at the resolution of level $j$. Either way only the constants change.
\item[(3)] \emph{Measurability.} In part (a) of the flat shell bound the functionals $X$ are
measurable with respect to
the block field $V_{k_0}$ on the whole lattice together with the labels meeting $N_2^{++}$. For
the functional $X=\mathfrak{L}^h_2$ entering $R^{m\mathfrak{L}}$, the wider $\sigma$-algebra
generated by $\mathcal{C}$ and $h\cap\mathcal{D}_2$ is
used. Consider the part of $\mathfrak{L}^h_2$ due to deep labels in its interior: it is a
covariance under $\mathbb{E}_h$ against the field parts attached to deep labels, of second order
and of size $L^{-4\nu'}$ per deep label $Y'$ of level $k_0-\nu'$. This part is moved into the
piece of $\mathfrak{S}^{\flat,h}_2$ attached to $Y'$, so that $\mathfrak{L}^h_2$ is
$\mathcal{C}$-measurable.
\item[(4)] \emph{Part (b) of the flat shell bound on $B_{\times}$.} Here $\mathcal{C}$ does not
contain the internal
fields of a live region spanning both neighbourhoods, integrated in the plain class. The deep
labels on the
two sides are coupled through these fields at some order, and part (b) of the flat shell bound on
$B_{\times}$ must control this coupling.
\end{itemize}
\end{remark}

\begin{proof}
\emph{Item} (a). The finer index is the one on which
Lemma~\ref{8.1:lem-history-decomposition} constructs the history terms. The index of the
correlation theorem is obtained from it by summing, for each of its histories $h$, the finer terms
that $\mathbf{R}$ resums into $\tau^w_h$.
Exactness is preserved, because the decomposition is obtained by a regrouping of the finer
decomposition $\rho^w_{k_0}=\sum_{h'}\tau^w_{h'}$, the sum running over the finer index. The
operator $\mathsf{P}_{k_0}$ is linear, positive and independent of
$z$, and each of the following properties passes from the finer terms to their partial sums:
\begin{itemize}
\item positivity;
\item the identity $\rho^w_{k_0}=\sum_h\tau^w_h$;
\item the identity $\sum_h\int\tau^w_h=\int\rho^w_0$;
\item the identities
\begin{align*}
\partial_{z_i}\tau^w_h|_0&=\mathsf{P}_{k_0}(\mathcal{O}(f_i)\rho_0)_h,\\
\partial_{z_1}\partial_{z_2}\tau^w_h|_0&=\mathsf{P}_{k_0}(\mathcal{O}(f_1)\mathcal{O}(f_2)\rho_0)_h,
\end{align*}
where $\rho_0$ is the bare density $\rho^w_0$ at $z=0$.
\end{itemize}
Each of these is preserved when terms are added together, so it survives the summation. Thus
Lemma~\ref{8.1:lem-history-decomposition} holds on the index of the correlation theorem.

The total covariance formula of Lemma~\ref{8.1:lem-total-covariance} is derived algebraically from
these identities. It needs no smallness of the history terms that is uniform in $h$, so it holds
on this index as well. Proposition~\ref{8.1:prop-logarithmic-derivative} and
Lemma~\ref{8.1:lem-mixed-derivative-locality} are proved from the correlation theorem on that
theorem's own index.

\emph{Item} (b). In \cite{B-CMP122a} the summation over the large-field regions is applied to the
quotients of the densities. By \cite[(1.33)--(1.35), (1.71)--(1.72), (1.98)--(1.102)]{B-CMP122b},
at the step of $\mathbf{R}$ at which a
component is healed, the terms attached to that component are summed there, and what remains of
them is absorbed into the whole-lattice family $R^{(j)}$ of the effective action and into the
coupling-constant terms $b_j$. The corresponding constants enter the normalisation $N_{k_0}$.
Hence no term of the expansion at level $k_0$ is anchored at a
component healed at an earlier step, and the index at level $k_0$ records only the regions still
live there.

This argument uses only the fact that the component has been healed, not its age. In particular,
labels are not removed from the histories according to their age. A live
region of any birth level remains a label.

\emph{Item} (c). By definition $\mathcal{D}_i$ consists of healed labels of $h$. By item (b), no
history has healed labels, so $\mathcal{D}_i=\emptyset$. The pieces $\mathrm{Sh}^h_i$,
$\mathring{\mathfrak{S}}_i$ and $\mathfrak{S}^{\flat,h}_i$ are carried by the labels of
$\mathcal{D}_i$, and $m_i(h)$ depends on $h$ only through the configuration of deep labels;
hence the pieces vanish and $m_i(h)$ takes the single value $c^0_i$. The terms $\theta^h$ and
$\delta^{h,\mathrm{lab}}_i$ can be attached only to
labels of $h$, that is, to live regions; by the correlation theorem's description of the current
regions, they vanish unless $h\in\mathcal{B}_i$. The remainders
$R^{m\mathfrak{L}}$ and $(\mathrm{sh})'$ belong to the bookkeeping of the deep healed labels on
the finer index, described after item (e); since $\mathcal{D}_i=\emptyset$, they vanish as well.

\emph{Item} (d). By item (b), the residue of the healed components consists of the whole-lattice
covariant family $R^{(j)}$ and contributions to the coupling-constant terms $b_j$ and to the
normalisation $N_{k_0}$. Inside $\mathrm{Irr}_i$ and in the
mixed second derivative of Lemma~\ref{8.1:lem-mixed-derivative-locality}, the family $R^{(j)}$
enters with size $g_j^{\kappa_0}$, as in \cite{B-CMP119}. This size bound holds uniformly in the
history; hence, in the bounds of these terms, taking the supremum over histories is legitimate.

\emph{Item} (e). The atoms of the flat shell bound are deep healed labels
$Y\in\mathcal{D}_i$, and $\mu^{\sharp}_Y$ and $\chi_Y$ are defined for such $Y$ only. By item (c)
there are no atoms on the index of the correlation theorem. The spanning event $B_{\times}$ enters
only through part (b) of the flat shell bound, which is asserted also on $B_{\times}$, and the
conditional and healed clauses of the large-field block bound, healed part,
concern healed labels. Hence neither is needed on this index, and neither is the accompanying
condition on $\gamma$.

The argument for the two-point witness theorem in this chapter is carried out on the index of the
correlation theorem and does not use the flat shell bound. On the finer index the same conclusion
would follow from the flat shell bound, which is stated and not proved, together with the points
(1)--(4) above. The choice between the two indices affects only the length of the argument, not
its conclusion.
\end{proof}

\begin{lemma}[Covariance against a deep healed label]\label{8.1:lem-healed-label-covariance}
Assume the following two hypotheses.
\begin{itemize}
\item[(a)] Mechanism~(1) and mechanism~(2) of the flat shell bound hold.
\item[(b)] The first clause of the large-field block bound holds, that is, its one-point annealed
bound at every level.
\end{itemize}
Here and below, the \emph{regime without chain terms} means that mechanisms~(1) and~(2) are
applied to the flat shell bound with every one of its terms that carries a chain-decay factor
$\chi_Y$ omitted.
Let $\tilde{\nu}$ be the joint law of the history and the block field, and write $E$ for expectation
under $\tilde{\nu}$. Let $\mathcal{C}$ be the coarse $\sigma$-algebra. It contains
$\sigma(V_{k_0})$, where the block field $V_{k_0}$ is taken on the whole torus, and the coarse
labels of the history, that is, its large-field labels at the coarse levels, are
$\mathcal{C}$-measurable wherever on the torus they lie. Let $Y$ be a deep healed label
created at level $j$, write $\mathbf{1}_Y$ for the indicator that $Y$ is present in the history,
and write $\tilde{\nu}(Y):=E[\mathbf{1}_Y]$ for the probability that $Y$ is present. Let $p_Y$ be
the conditional probability, given $\mathcal{C}$, that $Y$ is present:
\begin{equation*}
p_Y:=E[\mathbf{1}_Y\mid\mathcal{C}],
\end{equation*}
so that $E[p_Y]=\tilde{\nu}(Y)$. Use the same letters with a prime for a second deep healed label
$Y'\neq Y$, created at level $j'$.

Then, in the regime without chain terms, the following hold.
\begin{enumerate}
\item[(i)] For every bounded $\mathcal{C}$-measurable random variable $\Xi$,
\begin{equation}\label{8.1:eq-healed-label-covariance-1}
|\operatorname{Cov}(\mathbf{1}_Y,\Xi)|\le \tilde{q}_j\,\mu^{\sharp}_Y\,E|\Xi|.
\end{equation}
\item[(ii)] The covariance of two such indicators satisfies
\begin{equation}\label{8.1:eq-healed-label-covariance-2}
\begin{aligned}
\operatorname{Cov}(\mathbf{1}_Y,\mathbf{1}_{Y'})
&=E\bigl[(p_Y-\tilde{\nu}(Y))(p_{Y'}-\tilde{\nu}(Y'))\bigr]\\
&\le \tilde{q}_j\,\tilde{q}_{j'}\,\mu^{\sharp}_Y\,\mu^{\sharp}_{Y'}.
\end{aligned}
\end{equation}
\end{enumerate}
\end{lemma}

The lemma is used only in the alternative bookkeeping in which healed labels are kept in the
history index rather than summed out as in Remark~\ref{8.1:rem-healed-summed-out}, namely in the
bound on the remainder $R^{m\mathfrak{L}}$ of $R_K$. The terms carrying a chain-decay factor remain
part of the flat shell bound as it is stated; only the regime without them is used here.

\begin{proof}
We first prove (i). Let $\Xi$ be bounded and $\mathcal{C}$-measurable. Then
\begin{equation*}
\operatorname{Cov}(\mathbf{1}_Y,\Xi)=E\bigl[(\mathbf{1}_Y-\tilde{\nu}(Y))\Xi\bigr].
\end{equation*}
Condition on $\mathcal{C}$. Since $\Xi$ is $\mathcal{C}$-measurable, the tower property gives
\begin{equation*}
E[\mathbf{1}_Y \Xi]=E\bigl[E[\mathbf{1}_Y\mid\mathcal{C}]\,\Xi\bigr]=E[p_Y \Xi].
\end{equation*}
Combining the two displays,
\begin{equation}\label{8.1:eq-healed-label-covariance-3}
\operatorname{Cov}(\mathbf{1}_Y,\Xi)=E\bigl[(p_Y-\tilde{\nu}(Y))\,\Xi\bigr].
\end{equation}
Now apply hypotheses~(a) and~(b) in the regime without chain terms. Mechanisms~(1) and~(2) of the
flat shell bound estimate the absolute value of the right-hand side of
\eqref{8.1:eq-healed-label-covariance-3} by the coupling size $\mu^{\sharp}_Y$ of the flat shell
bound times $E|\Xi|$ times the presence weight of the birth level $j$ of $Y$, and the one-point
annealed bound of hypothesis~(b) at the level $j$ bounds that weight by the birth-level weight
$\tilde{q}_j$. This is the bound \eqref{8.1:eq-healed-label-covariance-1}.

We now prove (ii). In the regime without chain terms, hypothesis~(a) gives that the conditional
covariance of $\mathbf{1}_Y$ and $\mathbf{1}_{Y'}$ given $\mathcal{C}$ has expectation zero, that is,
\begin{equation*}
E[\mathbf{1}_Y\mathbf{1}_{Y'}]=E[p_Y\,p_{Y'}].
\end{equation*}
By the definition of the covariance and this identity,
\begin{equation*}
\operatorname{Cov}(\mathbf{1}_Y,\mathbf{1}_{Y'})
=E[\mathbf{1}_Y\mathbf{1}_{Y'}]-\tilde{\nu}(Y)\,\tilde{\nu}(Y')
=E[p_Y\,p_{Y'}]-\tilde{\nu}(Y)\,\tilde{\nu}(Y').
\end{equation*}
Since $E[p_Y]=\tilde{\nu}(Y)$ and $E[p_{Y'}]=\tilde{\nu}(Y')$, expanding the product shows that the
right-hand side equals $E[(p_Y-\tilde{\nu}(Y))(p_{Y'}-\tilde{\nu}(Y'))]$.
This is the equality in \eqref{8.1:eq-healed-label-covariance-2}.

For the inequality, let $\Xi'$ be the sign of $p_{Y'}-\tilde{\nu}(Y')$, taken to be $0$ where this
difference vanishes. Since $p_{Y'}$ is $\mathcal{C}$-measurable, $\Xi'$ is bounded and
$\mathcal{C}$-measurable, and $E|\Xi'|\le1$. Identity \eqref{8.1:eq-healed-label-covariance-3} and
part~(i), both for $Y'$ in place of $Y$, give
\begin{equation*}
E|p_{Y'}-\tilde{\nu}(Y')|
=\operatorname{Cov}(\mathbf{1}_{Y'},\Xi')
\le \tilde{q}_{j'}\,\mu^{\sharp}_{Y'}\,E|\Xi'|
\le \tilde{q}_{j'}\,\mu^{\sharp}_{Y'}.
\end{equation*}
Now $\Xi:=p_{Y'}-\tilde{\nu}(Y')$ is bounded and
$\mathcal{C}$-measurable, so the equality already proved, identity
\eqref{8.1:eq-healed-label-covariance-3} and part~(i) for $Y$ give
\begin{align*}
\operatorname{Cov}(\mathbf{1}_Y,\mathbf{1}_{Y'})
&=E\bigl[(p_Y-\tilde{\nu}(Y))(p_{Y'}-\tilde{\nu}(Y'))\bigr]
=\operatorname{Cov}\bigl(\mathbf{1}_Y,\,p_{Y'}-\tilde{\nu}(Y')\bigr)\\
&\le \tilde{q}_j\,\mu^{\sharp}_Y\,E|p_{Y'}-\tilde{\nu}(Y')|
\le \tilde{q}_j\,\tilde{q}_{j'}\,\mu^{\sharp}_Y\,\mu^{\sharp}_{Y'}.
\end{align*}
This is the inequality in \eqref{8.1:eq-healed-label-covariance-2}.
\end{proof}

\begin{lemma}[The bad-history event and its annealed mass]\label{8.1:lem-bad-event}
We use the following data.
\begin{itemize}
\item $(f_1,f_2)$ is a witness pair, and $N_1,N_2$ are the neighbourhoods of $\operatorname{supp}f_1,
\operatorname{supp}f_2$ fixed in Definition~\ref{8.1:def-witness-class}.
\item $k_0=K-s$ is the witness level and $x=x_{k_0}$.
\item $\tilde{\nu}$ is the joint law of the history $h$ and the block field $V$ at level $k_0$.
\item $\check{R}_j=\check{R}(x_j)$ is Ba{\l}aban's radius function at level $j$.
\item For $k\le k_0$, a \emph{level-$k$ cell} is a cube of side $\check{R}_k$ spacings of the level-$k$
lattice, that is, of side $\check{R}_kL^{k-k_0}$ level-$k_0$ spacings; the radius $\rho$ of
Definition~\ref{8.1:def-witness-class} is measured in level-$k_0$ spacings.
\end{itemize}
We put
\begin{equation}\label{8.1:eq-bad-event-b}
\mathfrak{l}:=L\,\check{R}_{k_0}+c_L,\qquad \check{R}_{k_0}:=\check{R}(x_{k_0}),\qquad
\mathfrak{l}(y):=L\,\check{R}(y)+c_L ,
\end{equation}
where $c_L$ is the constant offset of the level count, not to be confused with the constant $C_L$ of
hypothesis~(a) below. The function $\mathfrak{l}(y)$ is read for $y$ in the
envelope of Definition~\ref{8.1:def-coupling-envelope}. For $k\le k_0$, $n_k(N_i)$ denotes the number of
level-$k$ cells meeting $N_i$. We say that $V$ is \emph{off the window} at a point if it violates the
level-$k_0$ small-field condition there. We define the \emph{contact event}
\begin{equation*}
\mathcal{B}_i:=\{\text{some live label of }h\text{ meets }N_i\}\cup
\{V\text{ is off the level-}k_0\text{ window on }N_i\},
\end{equation*}
with no restriction on the birth level of the label. The \emph{conditional variant} of the argument is the one
that uses the conditional rarity clause of the large-field block bound; in it, $\mathcal{B}_i$ is enlarged by
the event that a label of $h$ healed at step $k$, for some $k$ with $k_0-\mathfrak{l}\le k\le k_0$, meets $N_i$.
We put $\mathcal{B}:=\mathcal{B}_1\cup\mathcal{B}_2$, the \emph{bad-history event}; it contains the spanning
event $B_{\times}$ of the flat
shell bound, by the definition of that event in the flat shell bound. We let $\mathcal{B}_i^{\mathrm{c}}$, the
\emph{no-contact event}, be the complement of
$\mathcal{B}_i$. By convention $\mathcal{O}^{(k_0)}$ and $\mathrm{Irr}_i$ are set equal to $0$ where the field
is off the window; this set lies inside $\mathcal{B}$.

Assume the following.
\begin{enumerate}
\item[(a)] The live-at-the-level clause of the large-field block bound holds in the following unconditional
form. For every level-$k_0$ cell $\mathfrak{c}$,
\begin{equation*}
\tilde{\nu}\bigl(\text{a component live at level }k_0\text{ contains }\mathfrak{c}\bigr)
\le C_L\,e^{-\theta p_0(g_{k_0})},
\end{equation*}
whatever the birth level of the component, uniformly in $K$, in the volume and in the position of
$\mathfrak{c}$. (The large-field block bound states this bound also conditionally on events that do not depend
on the ancestry of the component; only the unconditional form is used here.) In its birth-level form, for every
$j\le k_0$ and every level $k'$ with $j\le k'\le k_0$, a label of birth level $j$ that is live (in the
conditional variant: or healed) at level $k'$ meets a given level-$j$ cell with $\tilde{\nu}$-probability at most
\begin{equation*}
\tilde{q}_j:=C(M\check{R}_j)^4e^{-c_q p_0(g_j)} .
\end{equation*}
Here $C$ is a constant independent of $j$, $k'$, $K$, the volume and the position of the cell, and
$c_q=\theta$ for live labels and $c_q=3/2$ for healed ones.
\item[(b)] The large-field block bound, healed part, holds. For every $k\le k_0$ and every level-$k$ cell
$\mathfrak{c}$,
\begin{equation*}
\tilde{\nu}\bigl(\text{a label of }h\text{ healed at step }k\text{ contains }\mathfrak{c}\bigr)
\le K_H\,e^{-(3/2)p_0(g_k)} .
\end{equation*}
\item[(c)] Along the levels read the degree function drops by at most one: $p_0(g_k)\ge p_0(g_{k_0})-1$ for
every $k\le k_0$.
\item[(d)] Ba{\l}aban's condition $p_0\ge 5r$ of \cite{B-CMP119} holds on the exponent of the degree
function $p_0(g')=A_0(\log g'^{-2})^{p_0}$.
\item[(e)] The coupling satisfies $g\le\gamma_W$, with the member of \eqref{8.1:eq-constants-ceiling-a}
that gives the inequality $L^{4\mathfrak{l}(x)}\cdot\mathrm{polylog}\le e^{\frac12c_qp_0(g_{k_0})}$ along the
envelope.
\item[(f)] The live size-and-age clause of the large-field block bound holds: the event that a label of $h$
born before $k_0-\mathfrak{l}$ and still live at $k_0$ meets $N_1\cup N_2$, which we call the
\emph{old-live-component event}, has $\tilde{\nu}$-probability at most a constant times
$\rho^4e^{-c\,p_0(g_{k_0})}$, for some $c>0$, this bound being a sum over ages of contributions that converge
super-exponentially in the age.
\item[(g)] The comparison of rarity weights along the window, stated together with the joint-presence clause
of the large-field block bound, holds: for $k_0-\mathfrak{l}\le j\le k_0$,
\begin{equation*}
e^{-c_qp_0(g_j)}\le e^{A_0}e^{-c_qp_0(g_{k_0})} .
\end{equation*}
\end{enumerate}
Then the following hold.
\begin{enumerate}
\item[(i)] For $k\le k_0$, $n_k(N_i)\le(8\rho L^{k_0-k}/\check{R}_k+2)^4$.
\item[(ii)] Up to a $\tilde{\nu}$-null set, the off-window clause of $\mathcal{B}_i$ is contained in its label
clause.
\item[(iii)] For $i=1,2$,
\begin{equation}\label{8.1:eq-bad-event-1}
\tilde{\nu}(\mathcal{B}_i)\le n_{k_0}(N_i)\,C_L\,e^{-\theta p_0(g_{k_0})}
+e^{3/2}K_H\sum_{k=k_0-\mathfrak{l}}^{k_0}n_k(N_i)\,e^{-(3/2)p_0(g_{k_0})} .
\end{equation}
The second term is present only in the conditional variant.
\item[(iv)] There are constants $c'>0$ and $C''$, depending only on $\theta$, on $C$, $M$, $A_0$, $C_L$,
$K_H$ and on the constants of hypothesis~(f), and in particular not on $K$, $g_0$, $m$, $s$, the history or
the position, such that
\begin{equation}\label{8.1:eq-bad-event-2}
\tilde{\nu}(\mathcal{B})\le C''\rho^4e^{-c'p_0(g_{k_0})}
\end{equation}
and
\begin{equation}\label{8.1:eq-bad-event-a}
\mathfrak{q}_W:=2\sum_{j\ge k_0-\mathfrak{l}}(8\rho L^{k_0-j})^4\tilde{q}_j
\le C''\rho^4e^{-c'p_0(g_{k_0})} ,
\end{equation}
where $\tilde{q}_j$ is taken with $c_q=\theta$. Read by birth level, $\tilde{\nu}(\mathcal{B})$ is at most
$\mathfrak{q}_W$, which accounts for the live labels born at steps $j\ge k_0-\mathfrak{l}$, plus the
$\tilde{\nu}$-probability of the old-live-component event of hypothesis~(f), plus, in the conditional variant,
the sum over $i=1,2$ of the second term of \eqref{8.1:eq-bad-event-1}. For the
two terms of \eqref{8.1:eq-bad-event-1} one may take $c'=\tfrac12\min\{\theta,3/2\}$.
\item[(v)] On $\mathcal{B}_i^{\mathrm{c}}$, in the decomposition \eqref{8.1:eq-logarithmic-derivative-e},
\begin{equation*}
\Phi^h_i=c^h_i+N_{k_0}\mathcal{O}^{(k_0)}(\tilde{f}_i)+\mathcal{O}^{(k_0)}(\varepsilon^h_i)
+\mathrm{Irr}_i+\mathrm{Sh}^h_i .
\end{equation*}
Off $\mathcal{B}$, $\Psi^h$ is the sum of $\partial_{z_1}\partial_{z_2}T$ over the spanning local terms $T$ of
the flat shell bound.
\end{enumerate}
\end{lemma}

\begin{proof}
\emph{Why the event contains live labels of every birth.}
Ba{\l}aban's operation $\mathbf{R}$ heals a large-field component only under two conditions,
\cite[\S1]{B-CMP122a}. The component must be contained in a cube of side $100MR_k$; here $R_k=R(x_k)$, and
$R(y)$, a function of the inverse coupling $y$, is the radius of the correlation theorem, the least power of
$L$ not smaller than $(\log y)^{r}$, to be
distinguished from Ba{\l}aban's radius function $\check{R}$. In addition, no new
large-field regions may have been created inside it in the preceding $R_j$ renormalization steps, where $j$ is
the step at which it was created (Ba{\l}aban writes $N$ for this number). All other components
remain live. By \cite[(1.80)--(1.81)]{B-CMP122b}, the lifetime of a component $Z$ created at step $j$ is at
most $2\log_L d'_j(Z)+\check{R}_j+2$, which depends on the size $d'_j(Z)$ of $Z$ in the sense of
\cite[(1.80)]{B-CMP122b}. By
\cite[(1.85)--(1.86)]{B-CMP122b}, merged or refreshed components inherit the oldest birth index and remain live.

It follows that a component live at level $k_0$ need not have been created within the last $\mathfrak{l}$
levels. Suppose the label clause of $\mathcal{B}_i$ were restricted to labels of level at least
$k_0-\mathfrak{l}$. Then it would omit the histories in which an old live label meets $N_i$. On these histories
the decomposition \eqref{8.1:eq-logarithmic-derivative-e} carries a large-field term $\mathrm{LF}^h_i$, and the
representation in (v) fails. For this reason the event names live labels of any birth, and the lifetime enters
only the estimate.

\emph{The level count.}
The number $R_j$ of preceding steps without new large-field regions that healing requires is read at the birth
of a component, \cite[\S1]{B-CMP122b}. The
function $\check{R}$ is
non-increasing along the flow, while $R$ can change by one power of $L$ across the levels at which a component
is live. The count $\check{R}(x_{k_0})+c_L$ can therefore fall short of the birth-level radius by up to a
factor $L$. The count $\mathfrak{l}$ of \eqref{8.1:eq-bad-event-b} contains the birth-level radius with one power
of $L$ of room.

The count $\mathfrak{l}$ introduces no cap in the sense of Definition~\ref{2.3:def-cap}. The quantity
$\mathfrak{l}$ depends only on $x_{k_0}$,
$L$ and $c_L$; in particular it does not depend on the history $h$ or on the block field. The coupling
thresholds that use it depend on $L$, and $L$ is chosen before
$\gamma_W$.

\emph{Proof of (i).}
By Definition~\ref{8.1:def-witness-class}, $\operatorname{supp}f_i$ lies in a ball of radius $\rho$, measured
in spacings of the level-$k_0$ lattice (the cells of Definition~\ref{8.1:def-witness-class} are these
spacings; a level-$k$ cell is a cube of side $\check{R}_k$ level-$k$ spacings). Hence
$N_i=(\operatorname{supp}f_i)^{+\rho/4}$ lies in a ball of radius
$5\rho/4$. That ball is contained in a cube of side $5\rho/2\le 8\rho$ level-$k_0$ spacings.

A cube of side $\ell$ meets at most $(\ell/a+2)^4$ cubes of side $a$ of a lattice cover. A level-$k$ cell is a
cube of side $\check{R}_k$ level-$k$ spacings, that is, of side $\check{R}_kL^{k-k_0}$ level-$k_0$ spacings.
With $\ell=8\rho$ and $a=\check{R}_kL^{k-k_0}$ this gives (i).

\emph{Proof of (ii).}
By the construction of the history decomposition $\mathsf{P}_{k_0}$, through its characteristic functions
$\chi_h$, a history term $\tau_h$ vanishes where the
block field is off the level-$k_0$ window at a point of $N_i$ not covered by a live label of $h$. The
off-window clause of $\mathcal{B}_i$ therefore adds only a $\tilde{\nu}$-null set to the label clause.

\emph{Proof of (iii).}
By (ii),
\begin{equation*}
\tilde{\nu}(\mathcal{B}_i)\le
\tilde{\nu}(\text{some live label meets }N_i)+\tilde{\nu}(\text{the healed clause holds}),
\end{equation*}
where the second term is present only in the conditional variant.

Every live label meeting $N_i$ is a component live at level $k_0$ containing one of the $n_{k_0}(N_i)$
level-$k_0$ cells that meet $N_i$. A union bound over these cells and hypothesis~(a) give the first term of
\eqref{8.1:eq-bad-event-1}. No sum over birth levels is needed, because the live component at level $k_0$
is itself the object estimated.

Each label healed at step $k$, $k_0-\mathfrak{l}\le k\le k_0$, that meets $N_i$ contains one of the
$n_k(N_i)$ level-$k$ cells meeting $N_i$. A union bound over $k$ and over these cells and hypothesis~(b) bound
the healed clause by
\begin{equation*}
K_H\sum_{k=k_0-\mathfrak{l}}^{k_0}n_k(N_i)\,e^{-(3/2)p_0(g_k)} .
\end{equation*}
Inserting the lower bound (c) on $p_0(g_k)$, which gives
\begin{equation*}
e^{-(3/2)p_0(g_k)}\le e^{3/2}e^{-(3/2)p_0(g_{k_0})}
\end{equation*}
for every $k\le k_0$, yields the second term of \eqref{8.1:eq-bad-event-1}.

\emph{The source of hypotheses~(a) and~(f).}
Bounds of the forms required in hypotheses~(a) and~(f), with constants independent of $K$, of the position
and of the history, are supplied by sorting the live labels at level $k_0$ by their age and by their fate at
the last lattice of the run. Live labels of
age at most $\mathfrak{l}$ are covered by the rarity of young large-field blocks and its corollary. Live labels
older than $\mathfrak{l}$ that are healed at a later step of the run are covered by the healed old-region bound.
Live labels older than $\mathfrak{l}$ that are never healed are covered by the age bound at the last lattice of
the run, the run being prolonged on the same lattice. When the torus exponent $m$ exceeds a volume threshold
$m_V(g)$, this requires $s\ge s_{\mathrm{vol}}(g,m)$, where $s_{\mathrm{vol}}(g,m)$ is the per-torus depth
floor, a lower bound on $s$ depending on $g$ and on $m$; when $m\le m_V(g)$, it requires no condition on $s$
beyond $s\ge s_W(g)$. The uniformity in the volume stated in~(a) is read under that floor; removing the
floor's dependence on the torus exponent is not part of this lemma. The healed part, hypothesis~(b), enters
only the second term.

\emph{Proof of the inequality in \eqref{8.1:eq-bad-event-a}.}
The sum defining $\mathfrak{q}_W$ runs over $j$ with $k_0-\mathfrak{l}\le j\le k_0$, that is, over at most
$\mathfrak{l}+1$ terms. For each such term, $L^{4(k_0-j)}\le L^{4\mathfrak{l}}$ and, by~(g),
\begin{equation*}
e^{-c_qp_0(g_j)}\le e^{A_0}e^{-c_qp_0(g_{k_0})} .
\end{equation*}
Hence
\begin{equation*}
\mathfrak{q}_W\le 2\cdot8^4e^{A_0}CM^4\rho^4\,L^{4\mathfrak{l}}\,e^{-c_qp_0(g_{k_0})}
\sum_{j=k_0-\mathfrak{l}}^{k_0}\check{R}_j^4 .
\end{equation*}
The factor $\sum_j\check{R}_j^4$, a sum of at most $\mathfrak{l}+1$ terms, is the polylogarithmic factor
$\mathrm{polylog}$ of the absorption below. Indeed, for these $j$ we have $s\le K-j\le s+\mathfrak{l}$, so
Definition~\ref{8.1:def-coupling-envelope} gives
\begin{equation*}
x_j\le\bar{x}_{K-j}\le\bar{x}_s+B^{\sharp}\mathfrak{l},\qquad x_{k_0}\le\bar{x}_s ;
\end{equation*}
together with $\check{R}_j\le R(x_j)<L(\log x_j)^r$ and $\mathfrak{l}\le L\,R(x_{k_0})+c_L$ (see below), this
bounds $\mathfrak{l}+1$ and $\max_j\check{R}_j^4$ by a polynomial in $\log\bar{x}_s$ whose coefficients
depend on $L$, $c_L$, $B^{\sharp}$ and $r$.

We now bound the power of $L$. Recall that $R(y)$ is the least power of $L$
not smaller than $(\log y)^{r}$. Then $R(y)<L(\log y)^r$, so that
\begin{equation*}
L^{4R(y)}\le\exp\bigl(4L\log L\,(\log y)^r\bigr).
\end{equation*}
In $L^{4\mathfrak{l}(y)}$ we use $\check{R}(y)\le R(y)$; then
\begin{equation*}
L^{4\mathfrak{l}(y)}\le L^{4c_L}\exp\bigl(4L^2\log L\,(\log y)^r\bigr),
\end{equation*}
so the factor $L$ and the additive constant $c_L$ of \eqref{8.1:eq-bad-event-b} change
only two things: a constant depending on $L$ in the coefficient of $(\log y)^r$ in the exponent, and a factor
$L^{4c_L}$.

On the other side, $p_0(g_{k_0})=A_0(\log x)^{p_0}$ with $p_0\ge 5r$ by~(d). Consequently
\begin{equation*}
L^{4\mathfrak{l}(x)}\cdot\mathrm{polylog}\le e^{\frac12 c_q p_0(g_{k_0})}
\end{equation*}
as soon as $x$ exceeds a threshold depending on $L$. Along the envelope this is the ceiling condition
in~(e). Therefore $\mathfrak{q}_W\le C''\rho^4e^{-\frac12c_qp_0(g_{k_0})}$, which gives the bound with $c'$.

\emph{The birth-level reading in (iv).}
First, $\tilde{\nu}(\mathcal{B})\le\tilde{\nu}(\mathcal{B}_1)+\tilde{\nu}(\mathcal{B}_2)$, which is the
factor $2$. Each term is bounded by \eqref{8.1:eq-bad-event-1}. We organise this bound by the birth level of
the labels.

The labels born at some $j\ge k_0-\mathfrak{l}$ and live at $k_0$ meet $N_i$ through one of the level-$j$
cells meeting $N_i$. With the cube of side $5\rho/2$ from the proof of (i), at most
\begin{equation*}
\bigl(5\rho L^{k_0-j}/(2\check{R}_j)+2\bigr)^4\le(8\rho L^{k_0-j})^4
\end{equation*}
level-$j$ cells meet $N_i$, since $\rho\ge1$ (Definition~\ref{8.1:def-witness-class}) and $\check{R}_j\ge1$.
By the birth-level form of~(a), each such cell carries weight at most
$\tilde{q}_j$. These contributions make up the sum in \eqref{8.1:eq-bad-event-a}.

The labels born before $k_0-\mathfrak{l}$ and still live at $k_0$ form the old-live-component event of
hypothesis~(f). Its annealed share is bounded by hypothesis~(f), the live size-and-age clause of the large-field block
bound. That bound
converges super-exponentially in the age, and its total is of the same form
$C''\rho^4e^{-c'p_0(g_{k_0})}$. The off-window clause contributes nothing, by (ii), and in the conditional
variant the healed clause of $\mathcal{B}_i$ is bounded by the second term of \eqref{8.1:eq-bad-event-1}, as in
the proof of (iii). Thus, in the birth-level form, $\tilde{\nu}(\mathcal{B})$ is at most
$\mathfrak{q}_W$ plus the annealed mass of the old-live-component event, plus, in the conditional variant, the
sum over $i=1,2$ of the second term of \eqref{8.1:eq-bad-event-1}; each of the first two is at most
$C''\rho^4e^{-c'p_0(g_{k_0})}$, and the third is bounded below together with the first term of
\eqref{8.1:eq-bad-event-1}.

Finally, we absorb the first form \eqref{8.1:eq-bad-event-1} directly; this bounds $\tilde{\nu}(\mathcal{B})$
by the right-hand side of \eqref{8.1:eq-bad-event-2}. Since $\rho\ge1$
(Definition~\ref{8.1:def-witness-class}) and $\check{R}_k\ge1$, part (i) gives
$n_k(N_i)\le(10\rho)^4L^{4(k_0-k)}$. The sum over $k$ is therefore at most
$(\mathfrak{l}+1)(10\rho)^4L^{4\mathfrak{l}}$. The absorption above, with $c_q=\theta$ for the first term and
$c_q=3/2$ for the second, bounds twice the right-hand side of \eqref{8.1:eq-bad-event-1} by
\begin{equation*}
2\cdot10^4(C_L+e^{3/2}K_H)\rho^4e^{-c'p_0(g_{k_0})}.
\end{equation*}

\emph{Coarser labels.}
Labels of level above $k_0$ do not occur in the histories at step $k_0$. Their effect on
$\nu^B_{K,s}$ is treated in the Gaussian extraction.

\emph{Proof of (v).}
By Proposition~\ref{8.1:prop-logarithmic-derivative}, on $\{\tau_h>0\}$ the derivative $\Phi^h_i$ decomposes as
in \eqref{8.1:eq-logarithmic-derivative-e}. On $\mathcal{B}_i^{\mathrm{c}}$, no live label of $h$ meets $N_i$, so the term
$\mathrm{LF}^h_i$ of that decomposition is absent. On $\mathcal{B}_i^{\mathrm{c}}$ the field is in the window on $N_i$,
so the conventions on $\mathcal{O}^{(k_0)}$ and $\mathrm{Irr}_i$ do not intervene. What remains is the displayed
formula.

Off $\mathcal{B}$, neither $N_1$ nor $N_2$ meets a live label. The mixed derivative $\Psi^h$ then reduces to the
sum of $\partial_{z_1}\partial_{z_2}T$ over the spanning local terms $T$ of the flat shell bound.
\end{proof}

\begin{proof}[Proof of Theorem~\ref{8.1:thm-two-point-witness}: the main term and its decomposition into remainders]
\label{8.1:prf-main-term}
We set up the main term and split the rest of the connected two-point function into named
remainders. By the position-average identity we may replace every quantity below by its
position average $\langle\!\langle\cdot\rangle\!\rangle$ over the sublattice $\mathfrak{P}_s$ of
Notation~\ref{8.1:not-position-average}. To each translate we apply the total covariance formula
of Lemma~\ref{8.1:lem-total-covariance}, that is, \eqref{8.1:eq-total-covariance-a}:
\begin{equation*}
S^T_{2;K}=\mathbb{E}_{\tilde{\nu}}\bigl[\Psi^h\bigr]
+\operatorname{Cov}_{\tilde{\nu}}\bigl(\Phi^h_1,\Phi^h_2\bigr).
\end{equation*}
Here $\tilde{\nu}$ is the joint law on pairs $(h,V)$ of a history and a block field, and
$\Phi^h_i$, $\Psi^h$ are the first and mixed logarithmic derivatives at zero source of that lemma.

\emph{Units.} Throughout the rest of the proof we write $N_{k_0}$ for the running normalisation at
the witness level $k_0=K-s$, and $g_{k_0}=x_{k_0}^{-1/2}$. We put
\begin{equation}\label{8.1:eq-main-term-a}
\begin{aligned}
\mathfrak{S}&:=g_{k_0}^4\,\mathfrak{n}(f_1)\,\mathfrak{n}(f_2)\,\rho^{-8},
\qquad
\varepsilon:=g_{k_0}\,p_0(g_{k_0}),\\
\mathfrak{M}&:=N_{k_0}^2\,b_0\,g_{k_0}^4\,\mathcal{K}_{K,s}(f_1,f_2)
\ \ge\ \tfrac{4}{9}\,b_0\,c_{\mathcal{K}}\,\mathfrak{S}
=:c_M\,\mathfrak{S}
\ \ge\ c_M\,\vartheta^2\,g_{k_0}^4\,\|f_1\|_\infty\,\|f_2\|_\infty .
\end{aligned}
\end{equation}
The factor $\tfrac49$ in the first inequality comes from $N_{k_0}\in[\tfrac23,\tfrac43]$, which is
part of the statement of Lemma~\ref{8.1:lem-running-normalisation}, so that $N_{k_0}^2\ge\tfrac49$.
The constant $c_{\mathcal{K}}$ is the one in the lower bound \eqref{8.1:eq-kernel-bounds-b} of
Lemma~\ref{8.1:lem-kernel-bounds}. We abbreviate $x:=x_{k_0}$. By the envelope of the running
coupling \eqref{8.1:eq-coupling-envelope-f}, taken at $k_0=K-s$,
\begin{equation*}
x\in[\underline{x}_s,\bar{x}_s].
\end{equation*}
We write $\mathrm{polylog}(y)$ for a polynomial in $\log y$, and a quantity is called
$\mathrm{polylog}$ if it is a polynomial in $\log\bar{x}_s$, that is, of the form
$\mathrm{polylog}(\bar{x}_s)$. The
quantities $\rho$, $r$, $D_{\square}$, $\Pi_r$, $\mathfrak{p}_{\mathrm{src}}$ and
$\mathfrak{p}_{\mathrm{tot}}$ are of this kind.

\emph{Sizes.} The source lemma and the bad-event lemma (Lemma~\ref{8.1:lem-bad-event}) supply the
following bounds for the slot sizes, the sizes of the far and deep pieces, and the masses of the
exceptional events:
\begin{equation}\label{8.1:eq-main-term-1}
\begin{aligned}
\mathfrak{s}^{\sharp}_i&\le 10\,C_{\mathrm{live}}\,\bar{\Theta}\,
\mathfrak{p}_{\mathrm{src}}(\bar{x}_s)\,\rho^8\,\|f_i\|_\infty/\vartheta,\\
\mathfrak{f}_i&\le C_{\mathrm{src}}\,C_{\mathrm{far}}\,\|f_i\|_\infty\,\rho^8\,
e^{-\varkappa\rho/16},\\
\mathring{\mathfrak{S}}_i&\le C'\rho^4\,\|f_i\|_\infty\,r^{-1}\varepsilon^2 e^{-c\,p_0},\\
\mathfrak{q}_W&\le C''\rho^4\,e^{-c'p_0(g_{k_0})},\\
\mathfrak{q}_{\times}&\le C''''\rho^4\,e^{-c'p_0(g_{k_0})}.
\end{aligned}
\end{equation}
Here $C_{\mathrm{live}}$, $C_{\mathrm{src}}$, $C_{\mathrm{far}}$, $C'$, $C''$, $C''''$, $c$ and
$c'$ are constants supplied by the source lemma and by Lemma~\ref{8.1:lem-bad-event}.
In the bound on $\mathfrak{s}^{\sharp}_i$ we have used $\mathfrak{n}(f_i)\le 10\,\rho^4\|f_i\|_\infty$.
The bound on $\mathfrak{q}_W$ is the bound on the annealed mass \eqref{8.1:eq-bad-event-a} of the
bad-history event.

\emph{Decomposition of the first logarithmic derivative.} We apply clause~(d) of the source lemma.
It gives a representation that holds on the whole of the support of $\tau_h$, that is,
$\tilde{\nu}$-almost surely:
\begin{equation}\label{8.1:eq-main-term-b}
\Phi^h_i=m_i(h)+G_i(V)+\mathfrak{S}^{\flat,h}_i(V)+\mathfrak{L}^h_i(V).
\end{equation}
Its four terms have the following properties.
\begin{itemize}
\item The term $m_i(h)$ depends on the history only and not on the block field $V$.
\item The term
\begin{equation*}
G_i:=N_{k_0}\,\mathcal{O}^{(k_0)}(\tilde{f}_i)+\mathrm{Irr}_i
\end{equation*}
does not depend on the history. Here $\tilde{f}_i$ is the perturbed profile of $f_i$ in the class
$\tilde{\mathfrak{F}}(\mathsf{d},\vartheta)$ of Notation~\ref{8.1:not-position-average}, and
$\mathrm{Irr}_i$ is the irrelevant contribution to $\Phi^h_i$ in
Proposition~\ref{8.1:prop-logarithmic-derivative}. The term $G_i$ vanishes at the identity
configuration, $G_i(\mathbf{1})=0$, and
\begin{equation*}
|G_i|\le C_G\,\mathfrak{n}(f_i)\,\varepsilon^2 ,
\end{equation*}
with a constant $C_G$ supplied by the source lemma.
\item The deep term is a sum over the family $\mathcal{Y}_i$ of potential deep labels,
\begin{equation*}
\mathfrak{S}^{\flat,h}_i=\sum_{Y\in\mathcal{Y}_i}\mathbf{1}_{Y\in h}\,\xi^h_{i,Y},
\qquad
|\xi^h_{i,Y}|\le C_{\Theta}\,\bar{\Theta}\,\mathring{\sigma}_Y .
\end{equation*}
Each $\xi_{i,Y}$ is local. It is the summand belonging to $Y$ of the deep piece
$\mathfrak{S}^h_i$ of the source lemma, with the interior deep labels evaluated inside their live
closure, and this evaluation costs the factor $C_{\Theta}\bar{\Theta}$ in the bound; the number
$\mathring{\mathfrak{S}}_i$ of \eqref{8.1:eq-main-term-1} bounds the size of that deep piece.
For a label $Y$ created at step $j$, with $\nu=k_0-j$,
\begin{equation}\label{8.1:eq-main-term-2}
\mathring{\sigma}_Y
:=C\,(M\check{R}_j)^{c}\,x_j\,\mathrm{polylog}(x_j)\,L^{-4\nu}\,\varepsilon^2
\bigl(c_1^{\flat}\vee B_0\vee C_e+\mathfrak{d}_W+g_j^{\kappa_0}+\mathfrak{d}_W(x)\bigr)\,
\|f_i\|_\infty ,
\end{equation}
where $C$, $c$ and $C_e$ are constants supplied by the source lemma.
The factor $L^{-4\nu}$ is present for every label $Y$. It comes from the inversion estimate,
applied at the radii of \cite[(2.28)]{B-CMP119}, together with the volume bound
$|Y|\le(100M\check{R}_j)^4L^{-4\nu}$ per cube of the class of $Y$, which holds under the
condition on $\kappa$ in the order of choice (Definition~\ref{2.3:not-stages-middle}).
\item The last term splits as
\begin{equation*}
\mathfrak{L}^h_i=\mathfrak{L}^{\mathrm{near},h}_i+\mathfrak{F}^h_i .
\end{equation*}
Off the old-live event $\mathfrak{O}_i$ of the source lemma, the near part satisfies
$\sup|\mathfrak{L}^{\mathrm{near},h}_i|\le\mathfrak{s}^{\sharp}_i$. The near part vanishes off the
part of the contact event $\mathcal{B}_i$ (the bad-history event of
Lemma~\ref{8.1:lem-bad-event}) defined by its condition on the labels of $h$. The far part satisfies
$\sup|\mathfrak{F}^h_i|\le\mathfrak{f}_i$.
\end{itemize}
No indicator of the no-contact event $\mathcal{B}_i^{\mathrm{c}}$ occurs in
\eqref{8.1:eq-main-term-b}, because the
representation holds on the whole of the support of $\tau_h$. Consequently no correction term
arises from using the representation on $\mathcal{B}_i$. In the same way,
$\mathbb{E}_{\tilde{\nu}}\Psi$ needs no splitting into good and bad histories.

\emph{The remainders.} Suppressing the superscript $h$ and the argument $V$, we have for $i=1,2$
\begin{equation*}
\Phi_i=(m_i+G_i+\mathfrak{S}^{\flat}_i)+\mathfrak{L}_i
\end{equation*}
by \eqref{8.1:eq-main-term-b}. Covariance is bilinear; hence
$\operatorname{Cov}_{\tilde{\nu}}(\Phi_1,\Phi_2)$ equals
\begin{equation*}
\operatorname{Cov}(m_1+G_1+\mathfrak{S}^{\flat}_1,\,m_2+G_2+\mathfrak{S}^{\flat}_2)
+\operatorname{Cov}(\mathfrak{L}_1,\,m_2+G_2+\mathfrak{S}^{\flat}_2+\mathfrak{L}_2)
+\operatorname{Cov}(m_1+G_1+\mathfrak{S}^{\flat}_1,\,\mathfrak{L}_2).
\end{equation*}
In the second and third covariances we split off the parts containing $m_2$ and $m_1$. This gives
\begin{equation}\label{8.1:eq-main-term-3}
\begin{aligned}
\operatorname{Cov}_{\tilde{\nu}}(\Phi_1,\Phi_2)
&=\operatorname{Cov}(m_1+G_1+\mathfrak{S}^{\flat}_1,\,m_2+G_2+\mathfrak{S}^{\flat}_2)\\
&\quad+\bigl[\operatorname{Cov}(\mathfrak{L}_1,\,G_2+\mathfrak{S}^{\flat}_2+\mathfrak{L}_2)
+\operatorname{Cov}(G_1+\mathfrak{S}^{\flat}_1,\,\mathfrak{L}_2)\bigr]\\
&\quad+\bigl[\operatorname{Cov}(m_1,\mathfrak{L}_2)+\operatorname{Cov}(\mathfrak{L}_1,m_2)\bigr].
\end{aligned}
\end{equation}
All covariances here are taken under $\tilde{\nu}$. We name the three groups as follows.
\begin{itemize}
\item The first line of \eqref{8.1:eq-main-term-3} is $\operatorname{Cov}(G_1,G_2)+(\mathrm{sh})'$.
This defines $(\mathrm{sh})'$ as the difference between
$\operatorname{Cov}(m_1+G_1+\mathfrak{S}^{\flat}_1,\,m_2+G_2+\mathfrak{S}^{\flat}_2)$ and
$\operatorname{Cov}(G_1,G_2)$.
\item The bracket on the second line is the $\mathfrak{L}$-group, denoted $R^{\mathfrak{L}}$.
\item The bracket on the third line is denoted $R^{m\mathfrak{L}}$.
\end{itemize}
For the expectation term, the source lemma gives
\begin{equation}\label{8.1:eq-main-term-4}
\mathbb{E}_{\tilde{\nu}}\Psi
=\mathbb{E}_{\tilde{\nu}}\Psi^{\mathrm{span}}+\mathbb{E}_{\tilde{\nu}}\operatorname{Cov}_{\mathbb{E}_h}
=:R^{\mathrm{loc}}+R^{\Psi,\mathrm{live}} ,
\end{equation}
where $\Psi^{\mathrm{span}}$ is the spanning part of the mixed derivative and
$\operatorname{Cov}_{\mathbb{E}_h}$ is the covariance term taken under the normalised law
$\mathbb{E}_h$ of the live exponent.

\emph{The history index of Remark~\ref{8.1:rem-healed-summed-out}.} The rest of the proof is
carried out in that history index, in which every healed component is summed out; there:
\begin{itemize}
\item $m_i$ reduces to a constant independent of the history, which we denote $c^0_i$:
$m_i\equiv c^0_i$;
\item the family of deep healed labels is empty, $\mathcal{D}_i=\emptyset$;
\item the deep piece $\mathfrak{S}^h_i$ of the source lemma, the deep term
$\mathfrak{S}^{\flat,h}_i$ of \eqref{8.1:eq-main-term-b} and the term $\bar{D}$ of the source
lemma all vanish identically: $\mathfrak{S}^h_i\equiv\mathfrak{S}^{\flat,h}_i\equiv0$ and
$\bar{D}\equiv0$.
\end{itemize}
Therefore
\begin{equation*}
\Phi^h_i=c^0_i+G_i+\mathfrak{L}^h_i,
\qquad
\Psi^h=\Psi^{\mathrm{span},h}+\operatorname{Cov}_{\mathbb{E}_h}.
\end{equation*}
A constant has zero covariance with every random variable, so $R^{m\mathfrak{L}}=0$. Since moreover
$\mathfrak{S}^{\flat}_i\equiv0$, the first line of \eqref{8.1:eq-main-term-3} reduces to
$\operatorname{Cov}(G_1,G_2)$, and $(\mathrm{sh})'=0$. Then \eqref{8.1:eq-main-term-3} becomes
\begin{equation}\label{8.1:eq-main-term-5}
\operatorname{Cov}_{\tilde{\nu}}(\Phi_1,\Phi_2)
=\operatorname{Cov}(G_1,G_2)
+\bigl[\operatorname{Cov}(\mathfrak{L}_1,G_2)+\operatorname{Cov}(G_1,\mathfrak{L}_2)
+\operatorname{Cov}(\mathfrak{L}_1,\mathfrak{L}_2)\bigr].
\end{equation}
The bracket in \eqref{8.1:eq-main-term-5} is $R^{\mathfrak{L}}$ with
$\mathfrak{S}^{\flat}_1=\mathfrak{S}^{\flat}_2=0$. The Gaussian extraction, applied at depth $s$ to
the perturbed profiles $\tilde{f}_1,\tilde{f}_2$, writes
\begin{equation*}
\operatorname{Cov}(G_1,G_2)=\mathfrak{M}+R^{\mathrm{prof}}+R^{\mathrm{GX}}+R^{\mathrm{irr}} ,
\end{equation*}
with three remainders $R^{\mathrm{prof}}$, $R^{\mathrm{GX}}$, $R^{\mathrm{irr}}$ that are specified
by the Gaussian extraction and bounded in the parts of the proof that follow.
Combining this with \eqref{8.1:eq-main-term-4}, \eqref{8.1:eq-main-term-5} and the total covariance
formula, we obtain, for each translate,
\begin{equation}\label{8.1:eq-main-term-6}
S^T_{2;K}=\mathfrak{M}+R^{\mathrm{prof}}+R^{\mathrm{GX}}+R^{\mathrm{irr}}+R^{\mathfrak{L}}
+R^{\mathrm{loc}}+R^{\Psi,\mathrm{live}} .
\end{equation}
Averaging over $\mathfrak{P}_s$ by the position-average identity, the same identity holds with
$\langle\!\langle\cdot\rangle\!\rangle$ applied to every term.
The six remainders
\begin{equation*}
R^{\mathrm{prof}},\ R^{\mathrm{GX}},\ R^{\mathrm{irr}},\ R^{\mathfrak{L}},\ R^{\mathrm{loc}},\
R^{\Psi,\mathrm{live}}
\end{equation*}
are each bounded in absolute value by $\mathfrak{M}/32$; these six bounds are proved in what
follows.
\end{proof}

\begin{proof}[The profile remainder]\label{8.1:prf-profile-remainder}
Let $\tilde{f}_i=f_i+\delta_i$, $i=1,2$, be the perturbed profiles, with $f_i\in\mathfrak{F}^{+}(\mathsf{d},\vartheta)$
and $\tilde{f}_i\in\tilde{\mathfrak{F}}(\mathsf{d},\vartheta)$ as in Notation~\ref{8.1:not-position-average}; the
perturbations $\delta_1,\delta_2$ are those produced by the profile absorption.
We write $\mathcal{K}=\mathcal{K}_{K,s}$. The main term $\mathfrak{M}$ is defined by
\begin{equation*}
\mathfrak{M}=N_{k_0}^2\,b_0\,g_{k_0}^4\,\mathcal{K}(f_1,f_2),
\end{equation*}
as in \eqref{8.1:eq-main-term-a}. We work under the standing condition $\rho\ge C_0\ge C_0^{\mathrm{prof}}$, where
$C_0^{\mathrm{prof}}=C_0^{\mathrm{prof}}(\vartheta,C_{\mathcal{K}}/c_{\mathcal{K}},C_hM)$ is the member of $C_0$ fixed below
and listed in
\eqref{8.1:eq-constants-ceiling-c}, and $C_0^{\mathrm{prof}}\ge 160C_hM/\vartheta^2$ as in
Notation~\ref{8.1:not-position-average}. We put
\begin{equation*}
R^{\mathrm{prof}}:=N_{k_0}^2\,b_0\,g_{k_0}^4\,\mathcal{K}(\tilde{f}_1,\tilde{f}_2)-\mathfrak{M},
\end{equation*}
the profile remainder, and $r^{\mathrm{prof}}:=|R^{\mathrm{prof}}|/\mathfrak{M}$ when $\mathfrak{M}>0$,
$r^{\mathrm{prof}}:=0$ when $\mathfrak{M}=0$. We claim that
\begin{equation}\label{8.1:eq-profile-remainder-1}
\bigl|N_{k_0}^2b_0g_{k_0}^4\mathcal{K}(\tilde{f}_1,\tilde{f}_2)-\mathfrak{M}\bigr|\le \mathfrak{M}\,r^{\mathrm{prof}},
\qquad
r^{\mathrm{prof}}\le C_{\mathrm{prof}}\,\frac{C_{\mathcal{K}}}{c_{\mathcal{K}}}\,
\frac{\mathfrak{d}_W(x_{k_0})\,C_hM}{\vartheta^2\rho},
\end{equation}
with an absolute constant $C_{\mathrm{prof}}$, and that $r^{\mathrm{prof}}\le 1/32$.

\emph{Inputs.} From the kernel bounds \eqref{8.1:eq-kernel-bounds-a} and \eqref{8.1:eq-kernel-bounds-b} of
Lemma~\ref{8.1:lem-kernel-bounds} we use two consequences.

First, an upper bound. Let $(u,v)$ be one of the pairs $(\delta_1,f_2)$, $(f_1,\delta_2)$, $(\delta_1,\delta_2)$,
whose supports are at inner distance at least $\rho/2$ (see below). Then
\begin{equation*}
|\mathcal{K}(u,v)|\le 2^8C_{\mathcal{K}}\rho^{-8}\,\mathfrak{n}(|u|)\,\mathfrak{n}(|v|).
\end{equation*}

Second, a lower bound:
\begin{equation*}
\mathcal{K}(f_1,f_2)\ge c_{\mathcal{K}}\rho^{-8}\,\mathfrak{n}(f_1)\,\mathfrak{n}(f_2).
\end{equation*}

By Notation~\ref{8.1:not-position-average}, the condition $\rho\ge C_0^{\mathrm{prof}}$ ensures that the pieces attached to
$f_1$ (namely $f_1$ and $\delta_1$) and those attached to $f_2$ are at inner distance at least $\rho/2$. Hence the
upper bound applies to each of the pairs $(\delta_1,f_2)$, $(f_1,\delta_2)$ and $(\delta_1,\delta_2)$.

Two further bounds concern the profiles. By the definition of the witness test class
$\mathfrak{F}^{+}(\mathsf{d},\vartheta)$, the mass of a witness test function is bounded below by
$\mathfrak{n}(f_i)\ge\vartheta\rho^4\|f_i\|_\infty$. For the perturbations produced by the profile absorption,
\begin{equation*}
\mathfrak{n}(|\delta_i|)\le \frac{160\,\rho^4\,\mathfrak{d}_W\,C_hM\,\|f_i\|_\infty}{\vartheta\rho},
\end{equation*}
where $\mathfrak{d}_W=\mathfrak{d}_W(x_{k_0})$ is the rate of Lemma~\ref{8.1:lem-running-normalisation}.

If $\|f_1\|_\infty\|f_2\|_\infty=0$, then by the bound on $|\delta_i|$ in the definition of
$\tilde{\mathfrak{F}}(\mathsf{d},\vartheta)$ one of the pairs $f_i,\delta_i$ vanishes identically. Both sides of
\eqref{8.1:eq-profile-remainder-1} are then zero. We assume from now on that both norms are positive. The lower bound
then gives $\mathcal{K}(f_1,f_2)>0$, and since $N_{k_0}\ge 2/3$ by Lemma~\ref{8.1:lem-running-normalisation}, also
$\mathfrak{M}>0$.

\emph{Bilinearity.} Since $\mathcal{K}$ is bilinear,
\begin{equation}\label{8.1:eq-profile-remainder-2}
\mathcal{K}(\tilde{f}_1,\tilde{f}_2)-\mathcal{K}(f_1,f_2)
=\mathcal{K}(\delta_1,f_2)+\mathcal{K}(f_1,\delta_2)+\mathcal{K}(\delta_1,\delta_2).
\end{equation}
We multiply \eqref{8.1:eq-profile-remainder-2} by $N_{k_0}^2b_0g_{k_0}^4$ and divide by $\mathfrak{M}$. Applying the
upper bound to each term on the right and the lower bound to $\mathcal{K}(f_1,f_2)$ gives
\begin{equation}\label{8.1:eq-profile-remainder-3}
\frac{\bigl|N_{k_0}^2b_0g_{k_0}^4\mathcal{K}(\tilde{f}_1,\tilde{f}_2)-\mathfrak{M}\bigr|}{\mathfrak{M}}
\le 2^8\frac{C_{\mathcal{K}}}{c_{\mathcal{K}}}\Bigl(\xi_1+\xi_2+\xi_1\xi_2\Bigr),
\qquad \xi_i:=\frac{\mathfrak{n}(|\delta_i|)}{\mathfrak{n}(f_i)}.
\end{equation}
Here we used $\mathfrak{n}(|f_i|)=\mathfrak{n}(f_i)$ for the non-negative profiles of $\mathfrak{F}^{+}$. Dividing the
bound on $\mathfrak{n}(|\delta_i|)$ by the lower bound on $\mathfrak{n}(f_i)$, the factors $\rho^4\|f_i\|_\infty$ cancel,
and
\begin{equation*}
\xi_i\le \xi:=\frac{160\,\mathfrak{d}_W\,C_hM}{\vartheta^2\rho}.
\end{equation*}
By the coupling ceiling \eqref{8.1:eq-constants-ceiling-a}, $\mathfrak{d}_W(x_{k_0})\le 1/3$, a bound free of the
couplings and uniform in the depth $s$. Together with
$\rho\ge 160C_hM/\vartheta^2$ this gives $\xi\le 1/3$, so that $\xi_1+\xi_2+\xi_1\xi_2\le 2\xi+\xi^2\le 3\xi$. Inserted
into \eqref{8.1:eq-profile-remainder-3}, this is \eqref{8.1:eq-profile-remainder-1}, with an absolute constant
$C_{\mathrm{prof}}\ge 3\cdot 2^8\cdot 160$, which we fix.

\emph{Smallness.} Using $\mathfrak{d}_W(x_{k_0})\le 1/3$ once more in \eqref{8.1:eq-profile-remainder-1},
\begin{equation*}
r^{\mathrm{prof}}\le \frac{C_{\mathrm{prof}}}{3}\,\frac{C_{\mathcal{K}}}{c_{\mathcal{K}}}\,\frac{C_hM}{\vartheta^2\rho}.
\end{equation*}
The right side is at most $1/32$ as soon as
\begin{equation*}
\rho\ge \frac{32}{3}\,C_{\mathrm{prof}}\,\frac{C_{\mathcal{K}}}{c_{\mathcal{K}}}\,\frac{C_hM}{\vartheta^2}.
\end{equation*}
We set
\begin{equation*}
C_0^{\mathrm{prof}}(\vartheta,C_{\mathcal{K}}/c_{\mathcal{K}},C_hM)
:=\max\Bigl\{\frac{32}{3}\,C_{\mathrm{prof}}\,\frac{C_{\mathcal{K}}}{c_{\mathcal{K}}}\,\frac{C_hM}{\vartheta^2},\;
\frac{160\,C_hM}{\vartheta^2}\Bigr\};
\end{equation*}
this is the member of $C_0$ entered in \eqref{8.1:eq-constants-ceiling-c}. Under $\rho\ge C_0\ge C_0^{\mathrm{prof}}$ we
therefore have $r^{\mathrm{prof}}\le 1/32$, uniformly in the couplings and in the depth.

\emph{Pieces not in $R^{\mathrm{prof}}$.} The label-supported pieces $\varepsilon^h_i$ of the first logarithmic
derivative (Proposition~\ref{8.1:prop-logarithmic-derivative}) do not enter $R^{\mathrm{prof}}$. Those carried by deep
labels are kept in the term $\mathfrak{S}^{\flat,h}_i$ of the decomposition \eqref{8.1:eq-main-term-b}. Those carried by
near labels that are live or coarse go into $\mathfrak{L}^h_i$, which is supported on the contact event $\mathcal{B}_i$.
There are no other label-supported pieces.
\end{proof}

\begin{proof}[Proof of the two-point witness theorem, continued: the large-field remainders]
\label{8.1:prf-large-field-remainders}
In this part of the proof we bound the sum
\begin{equation}\label{8.1:eq-large-field-remainders-1}
R^{\mathcal{B}\prime}:=R^{\mathfrak{L}}+R^{m\mathfrak{L}}+R^{\Psi,\mathrm{live}}
\end{equation}
of the three members of the remainder $R_K$ that arise, through the decomposition
\eqref{8.1:eq-main-term-b} of $\Phi^h_i$, from the large-field parts $\mathfrak{L}^h_i$ or from the
covariance of the live exponent.

\emph{Conventions.} Throughout, $E$ and $\operatorname{Cov}$ denote expectation and covariance under the
joint law $\tilde{\nu}$ of the history $h$ and the block field $V$. We write
$\mathfrak{n}_i:=\mathfrak{n}(f_i)$ for $i=1,2$. The letters $\bar{\Theta}$, $C_{\Theta}$,
$\mathfrak{s}^{\sharp}_i$, $\mathfrak{f}_i$, $\mathfrak{S}^h_i$, $\mathring{\mathfrak{S}}_i$,
$\mathfrak{S}^{\flat,h}_i$, $\mathfrak{L}^h_i=\mathfrak{L}^{\mathrm{near},h}_i+\mathfrak{F}^h_i$,
$m_i(h)$, $G_i$, $\mathfrak{q}_{\times}$, $\mathsf{X}$, $\mathbb{E}_h$, $\mathfrak{A}^w_h$ and
$\mathfrak{r}^{\mathrm{loc}}$ are those of the source lemma, and $C_G$ is the constant with
$|G_i|\le C_G\mathfrak{n}_i\varepsilon^2$ in its clause~(d). We omit the superscript $h$ where no
confusion arises. We write $S_i:=\operatorname{supp}f_i$,
and $S_i^{+}$ for the enlargement of $S_i$ fixed in the source lemma. We write
$\mathcal{B}_i$ for the bad-history event of Lemma~\ref{8.1:lem-bad-event},
$\mathcal{B}_i^{\mathrm{lab}}$ for the event defined by its label clause, and
$\mathcal{B}_i^{\mathrm{lab},\mathfrak{l}}$ for the part of $\mathcal{B}_i^{\mathrm{lab}}$ carried by
labels of level at least $k_0-\mathfrak{l}$ (for a live label: of age at most the level count
$\mathfrak{l}$ of \eqref{8.1:eq-bad-event-b}); by Lemma~\ref{8.1:lem-bad-event},
$\tilde{\nu}(\mathcal{B}_i^{\mathrm{lab},\mathfrak{l}})\le\mathfrak{q}_W$, with $\mathfrak{q}_W$ as in
\eqref{8.1:eq-bad-event-a}. Live labels of larger age are treated separately, through the events
$\mathfrak{O}_i$ below. We put
\begin{equation*}
B_{12}:=\mathcal{B}_1^{\mathrm{lab},\mathfrak{l}}\cap\mathcal{B}_2^{\mathrm{lab},\mathfrak{l}}.
\end{equation*}
A letter written without index ($\mathfrak{s}^{\sharp}$, $\mathfrak{f}$, $\mathfrak{n}$,
$\mathring{\mathfrak{S}}$, $\|f\|_\infty$) stands for the larger of its two indexed values. We speak
of the
\emph{kept bookkeeping} when the deep healed labels under $N_i$ are kept in $m_i(h)$ and in
$\mathfrak{S}^{\flat,h}_i$, and of the \emph{summed bookkeeping} when they are summed out; statements
marked ``(kept bookkeeping)'' below are used only in the former. The proof of the theorem follows
the summed bookkeeping; the kept bookkeeping is an alternative.

\emph{Inputs.} The bounds below are derived from the following statements, and from no others
besides Lemmas~\ref{8.1:lem-mixed-derivative-locality}, \ref{8.1:lem-healed-label-covariance}
and~\ref{8.1:lem-bad-event}:
\begin{enumerate}
\item clauses (b)--(e) of the source lemma, including the geometric statement and the bounds
(M1)--(M4) of its clause~(e), and its statement on the event $\mathfrak{O}_i$ described below;
\item the large-field block bound: its one-point annealed form, which gives the masses
$\mathfrak{q}_W$ and $\mathfrak{q}_{\times}$; and (kept bookkeeping) its joint-presence clause, used
for pairs of labels that are not both live, and its conditional clause, used for $E^{\natural}$ below
and in (M1)--(M3);
\item the theorem of this chapter that bounds the large-field remainders additively in terms of the
one-cell live weight and the one-cell age series (below, the additive theorem), through its one-point
organisation for live labels: the joint presence of two live labels is bounded by the presence of one
of them, so that one live weight enters linearly; the one-cell live weight is supplied by the
corollary on rarity of young large-field blocks; and the ages above $\mathfrak{l}$ enter through the
one-cell age series, whose part coming from regions healed after their creation is bounded by the
healed old-region bound and whose part coming from regions never healed is paid in advance at the
last lattice of the run, under the per-torus depth floor that is one of the displayed one-sided
conditions of clause~(iv) of Theorem~\ref{3.1:thm-main} (below, the finite-volume closure);
\item clause (a) of the flat shell bound, used for $R^{m\mathfrak{L}}$ only.
\end{enumerate}
The covariance $\operatorname{Cov}(m_1,m_2)$ does not enter $R^{\mathcal{B}\prime}$. It is bounded
once (kept bookkeeping), in the remainder $(\mathrm{sh})'$, by clause (b) of the flat shell bound
with the spanning
event $B_{\times}$ included; a bound in two pieces would set two divergent sums against a single
small factor.

\emph{An elementary inequality.} For bounded random variables $X,Y$ on the same probability space,
$\operatorname{Cov}(X,Y)=E[X(Y-EY)]$ and $|Y-EY|\le2\sup|Y|$, hence
\begin{equation}\label{8.1:eq-large-field-remainders-2}
|\operatorname{Cov}(X,Y)|\le2\sup|Y|\cdot E|X|,\qquad
|\operatorname{Cov}(X,Y)|\le2\sup|X|\cdot\sup|Y|.
\end{equation}

\emph{The events $\mathfrak{O}_i$.} Let $\mathfrak{O}_i$ be the event that $h$ has a live label
near $S_i$ (within distance $\rho/8$, in the sense in which clause (d) of the source lemma uses the
word) whose age exceeds the level count $\mathfrak{l}$ of
\eqref{8.1:eq-bad-event-b}. The
supremum bounds on $\mathfrak{L}^{\mathrm{near},h}_i$ and on $\mathfrak{L}^h_i$ used below hold for
every history outside $\mathfrak{O}_i$. On $\mathfrak{O}_i$ the source lemma bounds each
$\mathfrak{L}$-factor by the same supremum bound multiplied by $L^{4(a-\mathfrak{l})}$, where $a$ is
the largest age of a live label within distance $\rho/8$ of $S_i$. The shares of these events are paid
by the one-cell age series of the additive theorem. By that theorem, with $n_i:=n_{k_0}(N_i)$ the
number of $k_0$-cells of $N_i$ and $q^{=a}_{k_0}$ its one-cell age weight, a bound on the probability
that a given $k_0$-cell lies in a live large-field region of age $a$, the share of one
$\mathfrak{L}$-factor on $\mathfrak{O}_i$ is at most
\begin{equation*}
\mathfrak{s}^{\sharp}_i\,n_i\sum_{a>\mathfrak{l}}L^{4(a-\mathfrak{l})}q^{=a}_{k_0},
\end{equation*}
and that of a product of two $\mathfrak{L}$-factors at most
\begin{equation*}
\mathfrak{s}^{\sharp}_1\mathfrak{s}^{\sharp}_2\,(n_1+n_2)\sum_{a>\mathfrak{l}}L^{8(a-\mathfrak{l})}
q^{=a}_{k_0}.
\end{equation*}
In the age series, the regions healed after their creation are bounded by the healed old-region
bound, and the regions not healed by the last lattice of the run are paid in advance, on the
reference torus, under the per-torus depth floor of the finite-volume closure; the resulting shares
lie within the budget of the ceiling condition on $\gamma_{\mathrm{src}}$ in
\eqref{8.1:eq-constants-ceiling-b}. In what follows every
$\mathfrak{L}$-factor is taken outside $\mathfrak{O}_i$, that is, $\mathfrak{L}_i$ stands for
$\mathfrak{L}_i\mathbf{1}_{\mathfrak{O}_i^{\mathrm{c}}}$, and every event whose mass is used is taken
outside $\mathfrak{O}_1\cup\mathfrak{O}_2$; the shares of $\mathfrak{O}_1$ and $\mathfrak{O}_2$ are
the ones just bounded.

\emph{Tools.} (T1) By clause (d) of the source lemma,
$\sup_V|\mathfrak{L}^{\mathrm{near},h}_i|\le\mathfrak{s}^{\sharp}_i$, and
$\mathfrak{L}^{\mathrm{near},h}_i$ vanishes unless $h$ has a live label near $S_i$ or (kept
bookkeeping) a coarse healed label of level at least $k_0-\mathfrak{l}$ meeting $N_i$. Outside
$\mathfrak{O}_i$ a live label near $S_i$ has age at most $\mathfrak{l}$, that is, level at least
$k_0-\mathfrak{l}$, and it meets $N_i$; so, outside $\mathfrak{O}_i$, this event is contained in
$\mathcal{B}_i^{\mathrm{lab},\mathfrak{l}}$. Hence
\begin{equation}\label{8.1:eq-large-field-remainders-3}
E|\mathfrak{L}^{\mathrm{near}}_i|\le\mathfrak{s}^{\sharp}_i\,
\tilde{\nu}(\mathcal{B}_i^{\mathrm{lab},\mathfrak{l}})
\le\mathfrak{s}^{\sharp}_i\mathfrak{q}_W .
\end{equation}
(T2) By the same clause, $|\mathfrak{F}^h_i|\le\mathfrak{f}_i$ for every $h$ and $V$. With (T1),
$E|\mathfrak{L}_i|\le\mathfrak{s}^{\sharp}_i\mathfrak{q}_W+\mathfrak{f}_i$ and
$\sup|\mathfrak{L}_i|\le\mathfrak{s}^{\sharp}_i+\mathfrak{f}_i$.

(T3) Let $q^{\mathrm{lv}}_{k_0}$ be the one-cell live weight at level $k_0$, a bound on the
probability that a given $k_0$-cell lies in a live large-field region of any age, and put
\begin{equation*}
\mathfrak{q}_{12}:=\min\{n_1,n_2\}\,q^{\mathrm{lv}}_{k_0}+C_{LF}\mathfrak{q}_W^2,
\end{equation*}
where the second term is present only in the kept bookkeeping. We claim
\begin{equation}\label{8.1:eq-large-field-remainders-4}
\tilde{\nu}(B_{12})\le\mathfrak{q}_{12}+\mathfrak{q}_{\times}.
\end{equation}
On the complement of $\mathsf{X}$, a history in $B_{12}$ carries two labels, one near each $S_i$, and
by the geometric statement of clause (e) of the source lemma their collars are disjoint. If both
labels are live, then, since a live label near $S_i$ meets $N_i$, the history lies in
$B^{\mathrm{lv}}_1\cap B^{\mathrm{lv}}_2$, where $B^{\mathrm{lv}}_i$ is the event that a live label of
$h$ meets $N_i$. The joint presence is bounded by the presence of one label, and the union bound over
the $k_0$-cells of $N_i$ gives
\begin{equation*}
\tilde{\nu}\bigl(B^{\mathrm{lv}}_1\cap B^{\mathrm{lv}}_2\bigr)
\le\min\bigl\{\tilde{\nu}(B^{\mathrm{lv}}_1),\tilde{\nu}(B^{\mathrm{lv}}_2)\bigr\}
\le\min\{n_1,n_2\}\,q^{\mathrm{lv}}_{k_0};
\end{equation*}
this is the one-point organisation of the additive theorem, in which one live weight enters
linearly, and $q^{\mathrm{lv}}_{k_0}$ is supplied by the corollary on rarity of young large-field
blocks. If one of the two labels is a coarse healed label (kept bookkeeping), summing the
joint-presence clause of the large-field block bound over the list of labels in the definition of
$\mathcal{B}_i$ in Lemma~\ref{8.1:lem-bad-event} gives the mass $C_{LF}\mathfrak{q}_W^2$. On
$\mathsf{X}$ we use (M4): $\tilde{\nu}(\mathsf{X})=\mathfrak{q}_{\times}$.

(T4) (Kept bookkeeping.) For every $\mathfrak{S}$-factor that meets an $\mathfrak{L}$-factor we use
(M1) in its crude
form, $E[\mathfrak{S}^h_i\mathbf{1}_A]\le\mathring{\mathfrak{S}}_i$ for every event $A$; for the
product of two $\mathfrak{S}$-factors we use (M2),
$E[\mathfrak{S}^h_1\mathfrak{S}^h_2]\le C_{LF}\mathring{\mathfrak{S}}_1\mathring{\mathfrak{S}}_2$.

\emph{The member $R^{\mathfrak{L}}$.} We bound the covariances of $\mathfrak{L}^h_1$ with $G_2$,
$\mathfrak{S}^{\flat,h}_2$ and $\mathfrak{L}^h_2$; those with the indices exchanged are bounded in the
same way. First, by \eqref{8.1:eq-large-field-remainders-2} with $X=\mathfrak{L}_1$, $Y=G_2$, and by
(T1)--(T2),
\begin{equation*}
|\operatorname{Cov}(\mathfrak{L}_1,G_2)|\le2\sup|G_2|\cdot E|\mathfrak{L}_1|
\le2C_G\mathfrak{n}_2\varepsilon^2\bigl(\mathfrak{s}^{\sharp}_1\mathfrak{q}_W+\mathfrak{f}_1\bigr).
\end{equation*}
Second (kept bookkeeping), clause (d) of the source lemma gives
$|\mathfrak{S}^{\flat,h}_2|\le C_{\Theta}\bar{\Theta}\,\mathfrak{S}^h_2$, by summing its termwise bound
over $Y\in h\cap\mathcal{D}_2$. By \eqref{8.1:eq-large-field-remainders-2} with the roles of $X$ and $Y$
exchanged, by (T2) and by (T4),
\begin{equation*}
|\operatorname{Cov}(\mathfrak{L}_1,\mathfrak{S}^{\flat,h}_2)|
\le2\sup|\mathfrak{L}_1|\cdot E|\mathfrak{S}^{\flat,h}_2|
\le2C_{\Theta}\bar{\Theta}\bigl(\mathfrak{s}^{\sharp}_1+\mathfrak{f}_1\bigr)\mathring{\mathfrak{S}}_2 .
\end{equation*}
Third, by bilinearity $\operatorname{Cov}(\mathfrak{L}_1,\mathfrak{L}_2)$ is the sum of the four
covariances between $\mathfrak{L}^{\mathrm{near}}_i$ or $\mathfrak{F}_i$ on each side. The product
$\mathfrak{L}^{\mathrm{near}}_1\mathfrak{L}^{\mathrm{near}}_2$ vanishes off $B_{12}$. Hence, by (T1) and
\eqref{8.1:eq-large-field-remainders-4} for the first line below, by
\eqref{8.1:eq-large-field-remainders-3} for the second, by the first inequality of
\eqref{8.1:eq-large-field-remainders-2} with (T1)--(T2) for the third (which holds also with the
indices exchanged), and by the second inequality of \eqref{8.1:eq-large-field-remainders-2} for the
fourth,
\begin{align*}
|E[\mathfrak{L}^{\mathrm{near}}_1\mathfrak{L}^{\mathrm{near}}_2]|
&\le\mathfrak{s}^{\sharp}_1\mathfrak{s}^{\sharp}_2\bigl(\mathfrak{q}_{12}+\mathfrak{q}_{\times}\bigr),\\
|E\mathfrak{L}^{\mathrm{near}}_1\,E\mathfrak{L}^{\mathrm{near}}_2|
&\le\mathfrak{s}^{\sharp}_1\mathfrak{s}^{\sharp}_2\mathfrak{q}_W^2,\\
|\operatorname{Cov}(\mathfrak{L}^{\mathrm{near}}_1,\mathfrak{F}_2)|
&\le2\mathfrak{f}_2\mathfrak{s}^{\sharp}_1\mathfrak{q}_W,\\
|\operatorname{Cov}(\mathfrak{F}_1,\mathfrak{F}_2)|&\le2\mathfrak{f}_1\mathfrak{f}_2 .
\end{align*}
Altogether
\begin{equation*}
|\operatorname{Cov}(\mathfrak{L}_1,\mathfrak{L}_2)|
\le\mathfrak{s}^{\sharp}_1\mathfrak{s}^{\sharp}_2
\bigl(\mathfrak{q}_{12}+\mathfrak{q}_{\times}+\mathfrak{q}_W^2\bigr)
+\bigl(\mathfrak{s}^{\sharp}_1\mathfrak{f}_2+\mathfrak{f}_1\mathfrak{s}^{\sharp}_2\bigr)\,2\mathfrak{q}_W
+2\mathfrak{f}_1\mathfrak{f}_2 .
\end{equation*}
Adding the first two bounds and their counterparts with the indices exchanged to the third, and
replacing indexed
letters by the larger of their two values where they occur symmetrically, we obtain
\begin{equation}\label{8.1:eq-large-field-remainders-5}
\begin{split}
R^{\mathfrak{L}}\le{}&4C_G\mathfrak{n}\varepsilon^2\bigl(\mathfrak{s}^{\sharp}\mathfrak{q}_W+\mathfrak{f}\bigr)
+4C_{\Theta}\bar{\Theta}\bigl(\mathfrak{s}^{\sharp}+\mathfrak{f}\bigr)\mathring{\mathfrak{S}}\\
&+\mathfrak{s}^{\sharp}_1\mathfrak{s}^{\sharp}_2\bigl(\mathfrak{q}_{12}+\mathfrak{q}_W^2+\mathfrak{q}_{\times}\bigr)
+4\mathfrak{f}\mathfrak{s}^{\sharp}\mathfrak{q}_W+2\mathfrak{f}_1\mathfrak{f}_2 .
\end{split}
\end{equation}

\emph{The member $R^{m\mathfrak{L}}$ (kept bookkeeping).} Since
$m_1(h)=c^0_1+\sum_{Y\in h\cap\mathcal{D}_1}\psi_{1,Y}(\mathbf{1})$ and the history-free constant
$c^0_1$ of clause (d) of the source lemma has
covariance zero with every bounded random variable,
\begin{equation*}
\operatorname{Cov}(m_1,\mathfrak{L}_2)
=\sum_{Y\in\mathcal{Y}_1}\psi_{1,Y}(\mathbf{1})\,
\operatorname{Cov}\bigl(\mathbf{1}_{Y\in h},\mathfrak{L}^h_2\bigr).
\end{equation*}
Put $X:=\mathfrak{L}^h_2$. It is bounded, by $\mathfrak{s}^{\sharp}_2+\mathfrak{f}_2$ off
$\mathfrak{O}_2$ by (T2), and it is measurable with respect to
$\sigma(\mathcal{C}\vee h\cap\mathcal{D}_2)$, the class to which
Lemma~\ref{8.1:lem-healed-label-covariance} applies: the block field $V$ is global and the coarse
labels may lie anywhere, and the dependence of $\mathfrak{L}^h_2$ on the deep healed labels in the
interior of the live closures has been placed in $\mathfrak{S}^{\flat,h}_2$ by the decomposition of
clause (d) of the source lemma. Clause (a) of the flat shell bound therefore gives
\begin{equation*}
|\operatorname{Cov}(\mathbf{1}_Y,X)|\le\tilde{q}_{j(Y)}\,(\mu^{\sharp}_Y+\chi_Y)\,E^{\natural}|X| ,
\end{equation*}
with $j(Y)$ the level of the label $Y$.
Off $B_{\times}$, the conditional clause of the large-field block bound applies, because the labels
of $\mathcal{B}_2$ have collars disjoint from the collars of the labels under $N_1^{++}$, and gives,
the supremum being over the conditioning events $A$ admitted by that clause,
\begin{equation*}
E^{\natural}|X|\le\mathfrak{s}^{\sharp}_2\,
\sup_A\tilde{\nu}\bigl(\mathcal{B}_2^{\mathrm{lab},\mathfrak{l}}\,\big|\,A\bigr)+\mathfrak{f}_2
\le\mathfrak{s}^{\sharp}_2C_{LF}\mathfrak{q}_W+\mathfrak{f}_2 ;
\end{equation*}
the event $B_{\times}$ is covered by the clause of the flat shell bound that treats $B_{\times}$
itself. Summing over $Y\in\mathcal{Y}_1$ exactly as in the bound (M0) of the source lemma for
$\mathring{\mathfrak{S}}_1$ (the number of positions of a label at level difference $\nu$ is of order
$L^{4\nu}$, against the factor $L^{-4\nu}$ carried by $\mu^{\sharp}_Y$, while $\chi_Y$ decays
super-exponentially), with the constants $C'$ and $c$ of that bound (M0), we obtain
\begin{equation}\label{8.1:eq-large-field-remainders-6}
R^{m\mathfrak{L}}\le2C'\rho^4\|f\|_\infty\,r^{-1}
\bigl(\mathfrak{s}^{\sharp}C_{LF}\mathfrak{q}_W+\mathfrak{f}\bigr)\,e^{-c\,p_0(g_{k_0})} .
\end{equation}
This member uses clause (a) of the flat shell bound and the one-point form of the large-field block
bound together with its conditional clause.

\emph{The member $R^{\Psi,\mathrm{live}}$.} By clauses (c) and (d) of the source lemma,
\begin{equation*}
|E\,\Psi^h|\le\mathfrak{r}^{\mathrm{loc}}
+E\bigl|\operatorname{Cov}_{\mathbb{E}_h}\bigl(\partial_{z_1}\mathfrak{A}^w_h|_0,
\partial_{z_2}\mathfrak{A}^w_h|_0\bigr)\bigr| ;
\end{equation*}
the first term is the deterministic spanning sum of the main-term decomposition and is not part of
$R^{\Psi,\mathrm{live}}$. For every $h$ the covariance is at most
\begin{equation}\label{8.1:eq-large-field-remainders-7}
2\bigl(\mathfrak{s}^{\sharp}_1+\mathfrak{f}_1+C_{\Theta}\bar{\Theta}\mathfrak{S}^h_1\bigr)
\bigl(\mathfrak{s}^{\sharp}_2+\mathfrak{f}_2+C_{\Theta}\bar{\Theta}\mathfrak{S}^h_2\bigr),
\end{equation}
and it vanishes unless live labels reach both $S_1^+$ and $S_2^+$. This is part~(iii) of
Lemma~\ref{8.1:lem-mixed-derivative-locality} in the form
\eqref{8.1:eq-mixed-derivative-locality-b}: the covariance is taken under the joint measure
$\mathbb{E}_h$ of all live components. Live labels reach both supports in one of two
ways: either, outside $\mathfrak{O}_1\cup\mathfrak{O}_2$, $h\in B_{12}$, an event of mass at most
$\mathfrak{q}_{12}+\mathfrak{q}_{\times}$ by
\eqref{8.1:eq-large-field-remainders-4}; or one side is reached only through the long terms of far
components, whose whole share is contained in $\mathfrak{f}_i$ as a supremum over every $h$. We take
the expectation of \eqref{8.1:eq-large-field-remainders-7} term by term. The product
$\mathfrak{s}^{\sharp}_1\mathfrak{s}^{\sharp}_2$ is paid by the mass of $B_{12}$. The products of
$\mathfrak{s}^{\sharp}_i$ with $C_{\Theta}\bar{\Theta}\mathfrak{S}^h_{i'}$, $i\ne i'$, are bounded by
(M1) in its crude form, and the product of the two $\mathfrak{S}$-factors by (M2). The terms
containing a factor $\mathfrak{f}$ are bounded with $\mathfrak{f}$ as a supremum over every $h$, the
factor $\mathfrak{s}^{\sharp}$ of the other side entering with the mass $\mathfrak{q}_W$ and the factor
$C_{\Theta}\bar{\Theta}\mathfrak{S}^h$ through (M1) in its crude form. This gives
\begin{equation}\label{8.1:eq-large-field-remainders-8}
\begin{split}
R^{\Psi,\mathrm{live}}:={}&2\mathfrak{s}^{\sharp}_1\mathfrak{s}^{\sharp}_2
\bigl(\mathfrak{q}_{12}+\mathfrak{q}_{\times}\bigr)
+2C_{\Theta}\bar{\Theta}\bigl(\mathfrak{s}^{\sharp}_1\mathring{\mathfrak{S}}_2
+\mathfrak{s}^{\sharp}_2\mathring{\mathfrak{S}}_1\bigr)\\
&+2C_{\Theta}^2\bar{\Theta}^2C_{LF}\,\mathring{\mathfrak{S}}_1\mathring{\mathfrak{S}}_2
+4\mathfrak{f}\bigl(\mathfrak{s}^{\sharp}\mathfrak{q}_W
+C_{\Theta}\bar{\Theta}\mathring{\mathfrak{S}}+\mathfrak{f}\bigr).
\end{split}
\end{equation}

\emph{Relative sizes.} We compare each term of \eqref{8.1:eq-large-field-remainders-5},
\eqref{8.1:eq-large-field-remainders-6} and \eqref{8.1:eq-large-field-remainders-8} with
$c_M\vartheta^2g_{k_0}^4\|f_1\|_\infty\|f_2\|_\infty$, which is at most $\mathfrak{M}$ by
\eqref{8.1:eq-main-term-a}. Here $x=g_{k_0}^{-2}$, which is at most $\bar{x}_s$ by
\eqref{8.1:eq-coupling-envelope-f}, $\varepsilon^2=g_{k_0}^2p_0(g_{k_0})^2$,
$\rho\le LC_W(s)=\mathrm{polylog}(\bar{x}_s)$ and
$\mathfrak{p}_{\mathrm{src}}=\mathrm{polylog}(\bar{x}_s)$, where $\mathrm{polylog}$ means a polynomial in
$\log\bar{x}_s$. Let $c$ be the exponent of $p_0(g_{k_0})$ in
\eqref{8.1:eq-large-field-remainders-6} and in the bound for $\mathring{\mathfrak{S}}_i$, let $c'$ be
that in the bounds for $\mathfrak{q}_W$ and
$\mathfrak{q}_{\times}$, and put $c_{\mathrm{src}}:=\min\{c,c'\}$. Every term without a factor
$\mathfrak{f}$, divided by $c_M\vartheta^2g_{k_0}^4\|f_1\|_\infty\|f_2\|_\infty$, is at most
\begin{equation}\label{8.1:eq-large-field-remainders-9}
C_{\mathrm{tot}}\,\bar{\Theta}^2\,\mathfrak{p}_{\mathrm{tot}}(\bar{x}_s)\,\bar{x}_s^2\,
e^{-c_{\mathrm{src}}p_0(g_{k_0})}\,\vartheta^{-4},
\end{equation}
with $\mathfrak{p}_{\mathrm{tot}}$ a polylog. For instance, up to constant factors, the relative
sizes are
\begin{align*}
\mathfrak{s}^{\sharp}_1\mathfrak{s}^{\sharp}_2\mathfrak{q}_W^2:&\quad
\bar{\Theta}^2\mathfrak{p}_{\mathrm{src}}^2\rho^{24}x^2e^{-2c'p_0(g_{k_0})}\vartheta^{-4},\\
\mathfrak{s}^{\sharp}_1\mathfrak{s}^{\sharp}_2\mathfrak{q}_{\times}:&\quad
\bar{\Theta}^2\mathfrak{p}_{\mathrm{src}}^2\rho^{20}x^2e^{-c'p_0(g_{k_0})}\vartheta^{-4},\\
\mathfrak{n}\varepsilon^2\mathfrak{s}^{\sharp}\mathfrak{q}_W:&\quad
\bar{\Theta}\mathfrak{p}_{\mathrm{src}}\rho^{16}p_0(g_{k_0})^2\,x\,e^{-c'p_0(g_{k_0})}\vartheta^{-1},\\
\text{first member of }R^{m\mathfrak{L}}:&\quad
\bar{\Theta}\mathfrak{p}_{\mathrm{src}}\rho^{16}x^2e^{-(c+c')p_0(g_{k_0})}(r\vartheta)^{-1},\\
\bar{\Theta}^2\mathring{\mathfrak{S}}_1\mathring{\mathfrak{S}}_2:&\quad
\bar{\Theta}^2\rho^8p_0(g_{k_0})^4e^{-2c\,p_0(g_{k_0})}r^{-2};
\end{align*}
in the first two, $\mathfrak{s}^{\sharp}_1\mathfrak{s}^{\sharp}_2$ carries $\vartheta^{-2}$ and the
normalisation $c_M\vartheta^2$ another $\vartheta^{-2}$. All of them are dominated by
\eqref{8.1:eq-large-field-remainders-9}. The terms in which the one-cell live weight
$q^{\mathrm{lv}}_{k_0}$ enters through $\mathfrak{q}_{12}$, and the shares of the events
$\mathfrak{O}_i$, are bounded in the additive theorem in the same form, one exponential in
$p_0(g_{k_0})$ against $\bar{\Theta}^2$, a polylog and $\bar{x}_s^2$, with $c_{\mathrm{src}}$
replaced by $\min\{c_{\mathrm{src}},c_{\mathrm{lv}}\}$, where $c_{\mathrm{lv}}$ is the exponent of
$p_0(g_{k_0})$ in the bound for $q^{\mathrm{lv}}_{k_0}$ supplied by the corollary on rarity of young
large-field blocks. Every term with a factor $\mathfrak{f}$, divided by the same
quantity, is at most
\begin{equation}\label{8.1:eq-large-field-remainders-10}
C_{\mathrm{tot}}'\,\bar{\Theta}\,\mathfrak{p}_{\mathrm{tot}}'(\bar{x}_s)\,\rho^{24}\,\bar{x}_s^2\,
e^{-\varkappa\rho/16}\,\vartheta^{-2},
\end{equation}
with $\mathfrak{p}_{\mathrm{tot}}'$ a polylog, the term $\mathfrak{f}_1\mathfrak{f}_2$, which carries
$e^{-\varkappa\rho/8}$, included.

\emph{Conclusion.} Under the conditions \eqref{8.1:eq-constants-ceiling-b} (the ceiling condition on
$\gamma_{\mathrm{src}}$, with $c_{\mathrm{src}}$ replaced by $\min\{c_{\mathrm{src}},c_{\mathrm{lv}}\}$
in its exponential member, the condition on $\varepsilon$, and the floor $C_0^{\mathrm{src}}$), and,
for the part of the age series coming from regions never healed, on the reference torus under the
per-torus depth floor of the finite-volume closure, the bounds
\eqref{8.1:eq-large-field-remainders-9} and \eqref{8.1:eq-large-field-remainders-10}, together with
the terms in $q^{\mathrm{lv}}_{k_0}$ and the shares of the events $\mathfrak{O}_i$, which lie within
the budget of the ceiling condition on $\gamma_{\mathrm{src}}$, give
\begin{equation}\label{8.1:eq-large-field-remainders-11}
R^{\mathcal{B}\prime}\le\mathfrak{M}/32 .
\end{equation}
In the summed bookkeeping every $\mathring{\mathfrak{S}}$-term is $0$, and so is the term
$C_{LF}\mathfrak{q}_W^2$ of $\mathfrak{q}_{12}$. Of the source lemma only
clauses (b)--(d), its corollary and the bound (M4) are used there, together with the one-point form
of the large-field block bound, the one-point organisation of the additive theorem with the one-cell
live weight and the one-cell age series described above, and the ceiling condition on
$\gamma_{\mathrm{src}}$ and the floor
$C_0^{\mathrm{src}}$ of \eqref{8.1:eq-constants-ceiling-b}; the bounds (M1)--(M3) of clause (e), the
joint-presence and conditional clauses of the large-field block bound and clause (a) of the flat
shell bound are not needed.
\end{proof}

\begin{lemma}[The spanning local remainder]\label{8.1:prf-local-remainder}
Let $k_0=K-s$ be the witness level, let $j\le k_0$ denote the level of a local term and
$\nu:=k_0-j$ the level difference, let $f_1,f_2$ be the two test functions with
$\rho/4$-neighbourhoods $N_1,N_2$ of their supports, and let $r$ be the constant of
Lemma~\ref{8.1:lem-spanning-split}, so that the sources range over the polydisc
$|z_i|<r/(2\|f_i\|_\sharp)$, $i=1,2$. Assume the following two statements.
\begin{itemize}
\item[(a)] Clauses (c) and (d) of the correlation theorem at level $k_0$ hold, together with the
torus bound of that theorem for the local exponent terms $e_j(X;z)$. Let $E_0$, $B_0$ be the
constants of those clauses and $C_e$ the constant of the torus bound for the local exponent terms,
let $\delta\ge0$ be the fraction by which the decay rate $\kappa$ is reduced there, for which
$\kappa'':=\min\{(1-\tfrac14\delta)\kappa,\,(1-\delta)\kappa\}$ satisfies $\kappa''\ge\kappa/2$.
Through the supremum bounds of the
correlation theorem, measured in its flat norm $\|\cdot\|^{\flat}$, every local term $T$ at level
$j$ localised in $X$ of the families $\mathbf{E}^w$, $\mathbf{D}^w$, $\mathbf{B}^w$ and every local
exponent term $e_j(X;z)$ satisfies
\begin{equation*}
\sup_{\mathcal{W}_j(X)}|T|\le(E_0+B_0+C_e+1)\,e^{-\kappa'' d_j(X)},
\end{equation*}
and every local term $T$ at level $j$ localised in $X$ of the family $\mathbf{R}^w$ satisfies
\begin{equation*}
\sup_{\mathcal{W}_j(X)}|T|\le g_j^{\kappa_0}\,e^{-\kappa d_j(X)}.
\end{equation*}
Assume moreover that $g_j\le1$ for every level $j\le k_0$.
\item[(b)] Clauses (c) and (d) of the source lemma: for every history $h$,
\begin{equation}\label{8.1:eq-local-remainder-1}
|\Psi^{\mathrm{span},h}|\le \mathfrak{r}^{\mathrm{loc}},
\end{equation}
where $\Psi^{\mathrm{span},h}$ is the part of the mixed logarithmic derivative $\Psi^h$ made of
local terms spanning both supports, exterior and interior spanning terms together, and
\begin{equation}\label{8.1:eq-local-remainder-2}
\mathfrak{r}^{\mathrm{loc}}
:=\sum_{j\le k_0}\ \sum_{X:\ X^{+}\ \text{meets both supports}}\ \sum_{T}
\sup_{\mathcal{W}_j(X)}|T|\cdot\frac{4\|f_1\|_\sharp\,\|f_2\|_\sharp}{r^2},
\end{equation}
the inner sum running over all local terms $T$ at level $j$ localised in $X$ that can occur, each
counted once, and $X^{+}$ being the enlargement of $X$ of
Lemma~\ref{8.1:lem-mixed-derivative-locality}; the sum includes the local terms attached to one
support but localised far from it, and the far
live terms of \eqref{8.1:eq-logarithmic-derivative-d} that reach one support, all of these being
local terms of the kinds listed in hypothesis (a).
\end{itemize}
Put $R^{\mathrm{loc}}:=\mathfrak{r}^{\mathrm{loc}}$ for every history $h$, and note that
$\kappa''=(1-\delta)\kappa\le\kappa$ since $\delta\ge0$. Then:
\begin{itemize}
\item[(i)] there is a constant $C_{\mathrm{loc}}$, depending only on $E_0$, $B_0$,
$C_e$, $\kappa$ and the constant of the tree estimate \cite[(1.26)]{B-CMP116}, such that for every
$\rho\ge 24M$
\begin{equation}\label{8.1:eq-local-remainder-3}
R^{\mathrm{loc}}\le C_{\mathrm{loc}}\,r^{-2}\,\|f_1\|_\sharp\,\|f_2\|_\sharp
\Big(\frac{\rho}{M}\Big)^{4}L^{4}\,e^{-\kappa''\rho/(32M)};
\end{equation}
\item[(ii)] with $\mathfrak{M}$, $c_M$, $\vartheta$ and $x=x_{k_0}$ as in
\eqref{8.1:eq-main-term-a}, there is a constant $C'$, depending only on $C_{\mathrm{loc}}$, such
that $R^{\mathrm{loc}}\le\mathfrak{M}/32$
as soon as $\rho\ge 24M$ and
\begin{equation}\label{8.1:eq-local-remainder-4}
\rho\ge C_W(s)\ge\frac{32M}{\kappa''}\Big[2\log\bar{x}_s+4\log\frac{\rho}{M}+4\log L
+\log\frac{32C'}{\vartheta^2r^2c_M}\Big];
\end{equation}
under \eqref{8.1:eq-coupling-envelope-a} with $C_0\ge C_0^{\mathrm{loc}}$ of
\eqref{8.1:eq-constants-ceiling-c}, for every $s$ one has $C_W(s)\ge128M/\kappa''$, and the second
inequality of \eqref{8.1:eq-local-remainder-4} holds at $\rho=C_W(s)$; consequently
$R^{\mathrm{loc}}\le\mathfrak{M}/32$ for every $\rho\ge C_W(s)$ with $\rho\ge24M$.
\end{itemize}
\end{lemma}

\begin{proof}
\emph{The remainder and its summands.} By hypothesis (b), the spanning part of $\Psi^h$ is bounded
by $\mathfrak{r}^{\mathrm{loc}}$ for every $h$, and $\mathfrak{r}^{\mathrm{loc}}$ does not depend on
$h$; this is why $R^{\mathrm{loc}}$ may be defined as $\mathfrak{r}^{\mathrm{loc}}$ for every
history at once. The shape of the summands in \eqref{8.1:eq-local-remainder-2} is that of
\eqref{8.1:eq-mixed-derivative-locality-a}. A local term $T$ at level $j$ is analytic in the
sources on $\mathcal{W}_j(X)$, and by Cauchy's inequality on the polydisc
$|z_i|<r/(2\|f_i\|_\sharp)$, $i=1,2$, its mixed derivative at zero obeys
\begin{equation*}
|\partial_1\partial_2 T|_0|\le\sup_{\mathcal{W}_j(X)}|T|\cdot
\frac{4\|f_1\|_\sharp\,\|f_2\|_\sharp}{r^2}.
\end{equation*}
Moreover, by \eqref{8.1:eq-mixed-derivative-locality-a}, $T$ depends on both sources only if
$X^{+}$ meets both supports, and then $d_j(X)\ge\ell_\nu$, with
$\ell_\nu=[\rho L^\nu/(4M)-3]_{+}$ the lengths of Lemma~\ref{8.1:lem-spanning-split}.

\emph{Supremum bounds.} By hypothesis (a), every local term of the families $\mathbf{E}^w$,
$\mathbf{D}^w$, $\mathbf{B}^w$ and every local exponent term $e_j(X;z)$ has supremum on
$\mathcal{W}_j(X)$ at most $(E_0+B_0+C_e+1)\,e^{-\kappa'' d_j(X)}$, while every local term of the
family $\mathbf{R}^w$ has supremum at most $g_j^{\kappa_0}\,e^{-\kappa d_j(X)}$.
Since $\kappa''\le(1-\delta)\kappa\le\kappa$ for $\delta\ge0$, the factor $e^{-\kappa d_j(X)}$ is
at most $e^{-\kappa'' d_j(X)}$, and $g_j^{\kappa_0}\le1$ since $g_j\le1$ for $j\le k_0$ by
hypothesis (a) and $\kappa_0>0$, so that both kinds of terms are bounded by a constant times
$e^{-\kappa''d_j(X)}$. Each of the families $\mathbf{E}^w$, $\mathbf{D}^w$, $\mathbf{B}^w$,
$\mathbf{R}^w$ is a sum, over localisation regions, of one term localised in each region, and the
local exponent terms $e_j(X;z)$ are indexed by the regions in the same way. For fixed $(j,X)$ the
inner sum in \eqref{8.1:eq-local-remainder-2} therefore contains
at most one term of each of the families $\mathbf{E}^w$, $\mathbf{D}^w$, $\mathbf{B}^w$,
$\mathbf{R}^w$ and the one local exponent term $e_j(X;z)$, each potential term being counted once;
hence
\begin{equation*}
\sum_{T}\sup_{\mathcal{W}_j(X)}|T|\le\big(4(E_0+B_0+C_e+1)+1\big)\,e^{-\kappa''d_j(X)}.
\end{equation*}

\emph{Anchors and trees.} Fix $\nu\ge0$, that is the level $j=k_0-\nu$. Every region $X$ whose
enlargement $X^{+}$ meets both supports is anchored at a level-$j$ cube meeting $N_1$; there are
at most $(4L\rho L^\nu/M+3)^4$ such cubes. For a fixed anchor, the sum of
$e^{-\kappa''d_j(X)}$ over the regions $X$ containing it is carried out by the tree estimate
\cite[(1.26)]{B-CMP116}; because $d_j(X)\ge\ell_\nu$ on every region that contributes, this sum is
at most a constant times $e^{-\frac14\kappa''\ell_\nu}$. Collecting the three preceding paragraphs,
\begin{equation*}
R^{\mathrm{loc}}\le C\,r^{-2}\,\|f_1\|_\sharp\,\|f_2\|_\sharp
\sum_{\nu\ge0}\Big(\frac{4L\rho L^\nu}{M}+3\Big)^{4}e^{-\frac14\kappa''\ell_\nu},
\end{equation*}
with a constant $C$ formed from the constants of the supremum bounds, the factor $4$ and the tree
constant.

\emph{The term $\nu=0$.} Let $\rho\ge24M$. Then $\rho/(8M)\ge3$, hence
\begin{equation*}
\ell_0\ge\frac{\rho}{4M}-3\ge\frac{\rho}{8M},
\qquad
e^{-\frac14\kappa''\ell_0}\le e^{-\kappa''\rho/(32M)},
\end{equation*}
and $3\le L\rho/M$, hence
\begin{equation*}
\frac{4L\rho}{M}+3\le\frac{5L\rho}{M},
\qquad
\Big(\frac{4L\rho}{M}+3\Big)^{4}\le 5^4\Big(\frac{\rho}{M}\Big)^{4}L^{4}.
\end{equation*}
The sum over $\nu$ has the shape of the bound in Lemma~\ref{8.1:lem-spanning-split}, and its term
$\nu=0$ dominates. Indeed, for $\nu\ge1$ also $\rho L^\nu\ge24M$, so the same two estimates give
$\ell_\nu\ge\rho L^\nu/(8M)$ and $(4L\rho L^\nu/M+3)^4\le5^4(\rho/M)^4L^{4+4\nu}$; hence the
$\nu$-th term is at most $5^4(\rho/M)^4L^4e^{-\kappa''\rho/(32M)}$ times
\begin{equation*}
L^{4\nu}e^{-\kappa''\rho(L^\nu-1)/(32M)}\le L^{4\nu}e^{-\frac34\kappa''(L^\nu-1)}
\le \sup_{y\ge3^\nu}y^4e^{-\frac38\kappa(y-1)},
\end{equation*}
where we used $\rho/M\ge24$, $\kappa''\ge\kappa/2$ and $y=L^\nu\ge3^\nu$, $L$ being an odd
blocking factor, hence $L\ge3$. The
right-hand side is summable over $\nu\ge1$, with a sum depending only on $\kappa$. Absorbing
$5^4$, $C$ and one plus this sum into $C_{\mathrm{loc}}$ gives \eqref{8.1:eq-local-remainder-3}.
This proves (i).

\emph{The relative bound.} By \eqref{8.1:eq-main-term-a},
$\mathfrak{M}\ge c_M\vartheta^2g_{k_0}^4\|f_1\|_\sharp\,\|f_2\|_\sharp$, and $g_{k_0}^4=x^{-2}$ since
$x=x_{k_0}=g_{k_0}^{-2}$. Dividing \eqref{8.1:eq-local-remainder-3} by this lower bound, the norms
cancel, and with $C':=C_{\mathrm{loc}}$
\begin{equation}\label{8.1:eq-local-remainder-5}
\frac{R^{\mathrm{loc}}}{\mathfrak{M}}
\le\frac{C'x^2(\rho/M)^4L^4\,e^{-\kappa''\rho/(32M)}}{\vartheta^2r^2c_M}.
\end{equation}
Taking logarithms, the right-hand side of \eqref{8.1:eq-local-remainder-5} is at most $1/32$ if and
only if
\begin{equation}\label{8.1:eq-local-remainder-6}
\frac{\kappa''\rho}{32M}\ge 2\log x+4\log\frac{\rho}{M}+4\log L
+\log\frac{32C'}{\vartheta^2r^2c_M}.
\end{equation}
Since $x\in[\underline{x}_s,\bar{x}_s]$ by \eqref{8.1:eq-main-term-a}, we have
$\log x\le\log\bar{x}_s$, so this inequality follows from the second inequality of
\eqref{8.1:eq-local-remainder-4} together with $\rho\ge C_W(s)$, on multiplying by
$\kappa''/(32M)>0$. Hence $R^{\mathrm{loc}}\le\mathfrak{M}/32$.

\emph{The threshold is met.} In \eqref{8.1:eq-coupling-envelope-a} the room function $C_W(s)$
contains the term $C_0(\log\bar{x}_s)^{2p_0}$, with $p_0$ the fixed exponent of the degree function
$p_0(\cdot)$. The member $C_0^{\mathrm{loc}}$ of \eqref{8.1:eq-constants-ceiling-c} reserved for
this bound is fixed so that, for every $s$ and every $C_0\ge C_0^{\mathrm{loc}}$, the term
$C_0(\log\bar{x}_s)^{2p_0}$ is at least $128M/\kappa''$ and at least the right-hand side of
\eqref{8.1:eq-local-remainder-4} taken at $\rho=C_W(s)$; since $C_W(s)$ contains this term,
$C_W(s)\ge128M/\kappa''$ and the second inequality of \eqref{8.1:eq-local-remainder-4} holds at
$\rho=C_W(s)$. Here $\rho$ occurs on both sides of \eqref{8.1:eq-local-remainder-6};
since the derivative of
$\rho\mapsto\kappa''\rho/(32M)-4\log(\rho/M)$ is $\kappa''/(32M)-4/\rho$, which is positive for
$\rho>128M/\kappa''$, \eqref{8.1:eq-local-remainder-6} holds for every $\rho\ge C_W(s)$ as soon as
it holds at $\rho=C_W(s)$ and $C_W(s)\ge128M/\kappa''$. At $\rho=C_W(s)$,
\eqref{8.1:eq-local-remainder-6} follows from the second
inequality of \eqref{8.1:eq-local-remainder-4} at $\rho=C_W(s)$, as in the preceding paragraph.
Hence, for every $\rho\ge C_W(s)$ with $\rho\ge24M$, the right-hand side of
\eqref{8.1:eq-local-remainder-5} is at most $1/32$ and $R^{\mathrm{loc}}\le\mathfrak{M}/32$. Since
$C_0^{\mathrm{loc}}$ depends neither on $g$ nor on $s$, one choice $C_0\ge C_0^{\mathrm{loc}}$ gives
this for every $s$. This proves (ii).
\end{proof}

\begin{lemma}[The Gaussian-extraction remainder]\label{8.1:prf-extraction-remainder}
Assume the following two statements:
\begin{itemize}
\item[(A)] part~(a) of the Gaussian extraction;
\item[(B)] the corollary of the source lemma through which item~(8) of
Remark~\ref{8.1:rem-extraction-errors} is obtained.
\end{itemize}
Let $k_0=K-s$ be the witness level and $r=r(s)$ the extraction depth, and put
\begin{equation}\label{8.1:eq-extraction-remainder-1}
D_{\square}(s):=\max\Bigl\{\Bigl\lceil c_a^{-1}\bigl(4p_0\log\log\bar{x}_s
+8\log\bigl(L\,C_W(s)\bigr)+\log C_D\bigr)\Bigr\rceil,\ \ell+\ell_0\,r(s)+D_0\Bigr\},
\end{equation}
where $C_W(s)$ and the constant $C_D$ are those of
Definition~\ref{8.1:def-coupling-envelope}, and $\ell,\ell_0,D_0$ are the constants, independent
of $K$, of the second lower bound required of $D_{\square}$ in the Gaussian extraction. Write
$D_{\square}$ for $D_{\square}(s)$, let $r_1$ be as in the depth floor
\eqref{8.1:eq-coupling-envelope-c}, and put
\begin{align}
\mathcal{T}_1&:=\rho^8 g_{k_0}^{\sigma_1}\Pi_r,
&\mathcal{T}_2&:=\rho^8 g_{k_0}^2(1+r),\notag\\
\mathcal{T}_3&:=\rho^{12}D_{\square}^4(1+r)\,p_0(g_{k_0})^4e^{-c_ap_0(g_{k_0})},
&\mathcal{T}_4&:=\rho^8p_0(g_{k_0})^4e^{-c_aD_{\square}},\label{8.1:eq-extraction-remainder-2}\\
\mathcal{T}_5&:=L^{-4(r-r_1)},
&\mathcal{T}_6&:=L^{-8(r-r_1)}\,p_0(g_{k_0+r})^4.\notag
\end{align}
Then the Gaussian-extraction member $R^{\mathrm{GX}}$ of the decomposition
\eqref{8.1:eq-main-term-a} satisfies
\begin{equation}\label{8.1:eq-extraction-remainder-3}
\bigl|R^{\mathrm{GX}}\bigr|\le c_M\mathfrak{S}\cdot\frac{16}{9}\cdot\frac{C_a}{c_M}\cdot4
\cdot\sum_{j=1}^{6}\mathcal{T}_j .
\end{equation}
Suppose moreover that $r_0\ge r_0(L,A_0,C_a,c_M)$, a threshold depending only on
$L,A_0,C_a,c_M$; that $C_D\ge C_D(A_0,C_a,c_M)$, a threshold depending only on
$A_0,C_a,c_M$; and that the two ceiling conditions of \eqref{8.1:eq-constants-ceiling-a} on the
suprema
\begin{equation*}
\sup_y\bigl[y^{-\sigma_1/2}\,\mathrm{polylog}_2(y)\bigr],\qquad
\sup_y\bigl[y^{2}\,\mathrm{polylog}_3(y)\,e^{-c_aA_0(\log y)^{p_0}}\bigr]
\end{equation*}
hold, with the polynomials $\mathrm{polylog}_2$, $\mathrm{polylog}_3$ and the bounding constants
fixed there. Then for each $j=1,\dots,6$
\begin{equation}\label{8.1:eq-extraction-remainder-4}
c_M\mathfrak{S}\cdot\frac{16}{9}\cdot\frac{C_a}{c_M}\cdot4\cdot\mathcal{T}_j
\le\frac{\mathfrak{M}}{6\cdot32},
\qquad\text{hence}\qquad
\bigl|R^{\mathrm{GX}}\bigr|\le\frac{\mathfrak{M}}{32}.
\end{equation}
The number $D_{\square}(s)$ is independent of $K$ and is bounded above by a polynomial in
$\log\bar{x}_s$ whose coefficients do not depend on $K$.
\end{lemma}

\begin{proof}
Throughout, $p_0(g')=A_0(\log g'^{-2})^{p_0}$ is the degree function, and $g_k^{-2}=x_k$.
The coupling $g_{k_0+r}$ enters the error of part~(a), and the couplings are defined at the
levels $0,\dots,K$; hence $k_0+r\le K$, that is, $k_0+r=K-(s-r)$ with $0\le s-r\le s$. By the
envelope of Definition~\ref{8.1:def-coupling-envelope},
\begin{equation*}
\begin{aligned}
&\underline{x}_s\le x_{k_0}\le\bar{x}_s,\\
&x_{k_0+r}\le\bar{x}_{s-r}=X+B^{\sharp}(s-r+1)\le\bar{x}_s .
\end{aligned}
\end{equation*}
The coupling ceiling of Definition~\ref{8.1:def-coupling-envelope} is used in the form
$\underline{x}_{s'}\ge e$ for $0\le s'\le s$. Since $x_{k_0}\ge\underline{x}_s$ and, by the
envelope at depth $s-r$, $x_{k_0+r}\ge\underline{x}_{s-r}$, this gives $\log x_{k_0}\ge1$,
$\log x_{k_0+r}\ge1$ and $\log\bar{x}_s\ge\log\underline{x}_s\ge1$; on $[e,\infty)$ the map
$t\mapsto A_0(\log t)^{p_0}$ is increasing.
Consequently
\begin{equation}\label{8.1:eq-extraction-remainder-5}
\begin{gathered}
p_0(g_{k_0})^4\le A_0^4(\log\bar{x}_s)^{4p_0},\qquad
p_0(g_{k_0+r})^4\le A_0^4(\log\bar{x}_s)^{4p_0},\\
p_0(g_{k_0})\ge A_0(\log\underline{x}_s)^{p_0}.
\end{gathered}
\end{equation}

\emph{The bound \eqref{8.1:eq-extraction-remainder-3}.} Part~(a) of the Gaussian extraction
is applied at depth $r=r(s)$ with $D_{\square}=D_{\square}(s)$. The quantity it estimates is
\begin{equation*}
\operatorname{Cov}_{\nu^B_{K,s}}\bigl(N_{k_0}\mathcal{O}^{(k_0)}(\tilde{f}_1),\,
N_{k_0}\mathcal{O}^{(k_0)}(\tilde{f}_2)\bigr),
\end{equation*}
the covariance, under the block-field marginal $\nu^B_{K,s}$, of the block-field
observables of the perturbed profiles $\tilde{f}_i\in\tilde{\mathfrak{F}}(\mathsf{d},\vartheta)$
of Notation~\ref{8.1:not-position-average}, each multiplied by the normalisation
$N_{k_0}=\zeta'_{k_0}$ at the witness level, as in
part~(a); this is the object of part~(a) itself, with the
conventions under which part~(a) is stated, and item~(8) of
Remark~\ref{8.1:rem-extraction-errors} is supplied by hypothesis~(B). The error of part~(a) for
this covariance is $R^{\mathrm{GX}}$, and part~(a) bounds it, relative to $c_M\mathfrak{S}$, by the
sum of the six terms in \eqref{8.1:eq-extraction-remainder-3}. The first entry of the maximum in
\eqref{8.1:eq-extraction-remainder-1} is built from $\log\bar{x}_s$ and from
$C_W(s)=C_0p_0(\bar{x}_s^{-1/2})^2$, and the second from $r(s)$; none of these depends on
$K$. Moreover $p_0(\bar{x}_s^{-1/2})=A_0(\log\bar{x}_s)^{p_0}$, so that
\begin{equation*}
\log\bigl(L\,C_W(s)\bigr)=\log\bigl(L\,C_0A_0^2\bigr)+2p_0\log\log\bar{x}_s ,
\end{equation*}
and $\log\log\bar{x}_s\le\log\bar{x}_s$ because $\log\bar{x}_s\ge1>0$, as noted above,
and $\log t<t$ for $t>0$. With $\lceil a\rceil\le a+1$, the
first entry of the maximum is therefore at most
\begin{equation*}
1+c_a^{-1}\bigl(20p_0\log\bar{x}_s+8\log\bigl(L\,C_0A_0^2\bigr)+\log C_D\bigr),
\end{equation*}
a polynomial in $\log\bar{x}_s$ with coefficients independent of $K$. The second entry
$\ell+\ell_0r(s)+D_0$ is bounded by such a polynomial because $r(s)$ is, as stated with
\eqref{8.1:eq-main-term-a}, and $\ell,\ell_0,D_0$ are constants. This gives the last assertion.

\emph{The tower terms $\mathcal{T}_5,\mathcal{T}_6$.} The depth floor
\eqref{8.1:eq-coupling-envelope-c} gives
\begin{equation*}
r-r_1\ge\Bigl\lceil\Bigl(\frac{p_0}{2}+1\Bigr)\log_L\log\bar{x}_s\Bigr\rceil+r_0-1 .
\end{equation*}
Using $\lceil a\rceil\ge a$ and $L^{\log_L\log\bar{x}_s}=\log\bar{x}_s$, this yields
$L^{-8(r-r_1)}\le L^{8-8r_0}(\log\bar{x}_s)^{-4p_0-8}$, and with the second inequality of
\eqref{8.1:eq-extraction-remainder-5} and $(\log\bar{x}_s)^{-8}\le1$,
\begin{equation*}
\mathcal{T}_6\le A_0^4L^{8-8r_0}(\log\bar{x}_s)^{-8}\le A_0^4L^{8-8r_0};
\end{equation*}
the factor $(\log\bar{x}_s)^{4p_0}$ coming from $p_0(g_{k_0+r})^4$ is cancelled. In
particular $r-r_1\ge r_0-1$, so that $\mathcal{T}_5\le L^{4-4r_0}$.

\emph{The term $\mathcal{T}_4$.} By the first entry of the maximum in
\eqref{8.1:eq-extraction-remainder-1},
\begin{equation*}
e^{-c_aD_{\square}}\le C_D^{-1}(\log\bar{x}_s)^{-4p_0}\bigl(L\,C_W(s)\bigr)^{-8},
\end{equation*}
and the factor $(\log\bar{x}_s)^{-4p_0}$ compensates the bound on $p_0(g_{k_0})^4$ in
\eqref{8.1:eq-extraction-remainder-5}; with these,
\begin{equation*}
\mathcal{T}_4\le A_0^4C_D^{-1}\,\rho^8\bigl(L\,C_W(s)\bigr)^{-8}.
\end{equation*}
With the comparison $\rho\le L\,C_W(s)$ between the support radius $\rho$ in cells and the room
function, $\rho^8(L\,C_W(s))^{-8}\le1$, and hence $\mathcal{T}_4\le A_0^4C_D^{-1}$.

\emph{The terms $\mathcal{T}_1,\mathcal{T}_2,\mathcal{T}_3$.} The quantities $\rho$, $r$,
$D_{\square}$ and $\Pi_r$ are bounded by polynomials in $\log\bar{x}_s$, as stated with
\eqref{8.1:eq-main-term-a} (for $D_{\square}$ see also the last assertion proved above), and so
is the bound on $p_0(g_{k_0})^4$ in \eqref{8.1:eq-extraction-remainder-5}. Since
$g_{k_0}^{\sigma_1}=x_{k_0}^{-\sigma_1/2}\le\underline{x}_s^{-\sigma_1/2}$,
$g_{k_0}^2=x_{k_0}^{-1}\le\underline{x}_s^{-1}$, and, by the third inequality of
\eqref{8.1:eq-extraction-remainder-5},
$e^{-c_ap_0(g_{k_0})}\le e^{-c_aA_0(\log\underline{x}_s)^{p_0}}$, each of $\mathcal{T}_1$,
$\mathcal{T}_2$, $\mathcal{T}_3$ is at most a polynomial in $\log\bar{x}_s$ times, respectively,
\begin{equation*}
\underline{x}_s^{-\sigma_1/2},\qquad\underline{x}_s^{-1},\qquad
e^{-c_aA_0(\log\underline{x}_s)^{p_0}} .
\end{equation*}
These bounds are read at $y=\underline{x}_s$, the polynomials in $\log\bar{x}_s$ being written
through the affine envelope $\bar{x}_s\le\bar{x}(\underline{x}_s)$ of
Definition~\ref{8.1:def-coupling-envelope}. Here $\sigma_1=1$ (item~(7) of
Remark~\ref{8.1:rem-extraction-errors}) and $\underline{x}_s\ge e>1$, as noted after the
envelope, so that $\underline{x}_s^{-1}\le\underline{x}_s^{-\sigma_1/2}$
and $\underline{x}_s^{-2}\le1$. Thus $\mathcal{T}_1$ and $\mathcal{T}_2$ are at most
$\underline{x}_s^{-\sigma_1/2}$ times a polynomial in $\log\bar{x}(\underline{x}_s)$, which is the
quantity controlled by the first ceiling condition of \eqref{8.1:eq-constants-ceiling-a}; and,
writing $P$ for the polynomial that bounds the factors of $\mathcal{T}_3$ other than
$e^{-c_ap_0(g_{k_0})}$, expressed in $\log\bar{x}(\underline{x}_s)$,
\begin{equation*}
\mathcal{T}_3\le\underline{x}_s^{-2}\cdot
\Bigl[y^{2}\,P\bigl(\log\bar{x}(y)\bigr)\,e^{-c_aA_0(\log y)^{p_0}}\Bigr]_{y=\underline{x}_s},
\end{equation*}
where the bracket is the quantity controlled by the second ceiling condition and
$\underline{x}_s^{-2}\le1$. The two ceiling conditions of
\eqref{8.1:eq-constants-ceiling-a} are used in the form in which
$\mathrm{polylog}_2$ and $\mathrm{polylog}_3$ are the polynomial factors just described,
multiplied by $6\cdot32\cdot\frac{16}{9}\cdot\frac{C_a}{c_M}\cdot4$, and the bounding constants
are equal to one.

\emph{Conclusion.} Under $r_0\ge r_0(L,A_0,C_a,c_M)$ the bounds for $\mathcal{T}_5$ and
$\mathcal{T}_6$, under $C_D\ge C_D(A_0,C_a,c_M)$ the bound for $\mathcal{T}_4$, and under
the two ceiling conditions of \eqref{8.1:eq-constants-ceiling-a} the bounds for
$\mathcal{T}_1,\mathcal{T}_2,\mathcal{T}_3$ give
\begin{equation*}
\frac{16}{9}\cdot\frac{C_a}{c_M}\cdot4\cdot\mathcal{T}_j\le\frac{1}{6\cdot32},
\qquad j=1,\dots,6 .
\end{equation*}
Multiplying by $c_M\mathfrak{S}$ and using $c_M\mathfrak{S}\le\mathfrak{M}$ from
\eqref{8.1:eq-main-term-a} gives the first assertion of
\eqref{8.1:eq-extraction-remainder-4}. Summing the six inequalities and inserting the result
in \eqref{8.1:eq-extraction-remainder-3} gives
$|R^{\mathrm{GX}}|\le6\cdot\mathfrak{M}/(6\cdot32)=\mathfrak{M}/32$.
\end{proof}

\begin{proof}[Proof of the two-point witness theorem, continued: the irrelevant remainder and the room]\label{8.1:prf-irrelevant-remainder}
We bound the irrelevant remainder $R^{\mathrm{irr}}$, one of the named members of the total remainder
$R_K$, against the main term $\mathfrak{M}$ defined in \eqref{8.1:eq-main-term-a}. The only input is
clause~(b) of the Gaussian extraction, whose proof is incomplete at one step; we take that clause as
the hypothesis of this step.

Clause~(b) of the Gaussian extraction gives
\begin{equation}\label{8.1:eq-irrelevant-remainder-1}
  R^{\mathrm{irr}} \le 7C_b\,\rho^{-2}\,p_0(g_{k_0})^4\,\mathfrak{M}
  + \bigl(\text{terms of the type of } \mathcal{E}_a\bigr),
\end{equation}
where $C_b$ and $\mathcal{E}_a$ are the constant and the error term of the Gaussian extraction, and
$\rho$ is the support radius in cells.

In the first term of \eqref{8.1:eq-irrelevant-remainder-1} we use the room function
$C_W(s)=C_0\,p_0(\bar{x}_s^{-1/2})^2$ of \eqref{8.1:eq-coupling-envelope-a}, read at the envelope
$\bar{x}_s$ rather than at $x_{k_0}$: at the witness depth the support radius satisfies
$\rho\ge C_W(s)$, so that
\begin{equation*}
  \rho^{-2}\le C_W(s)^{-2}=C_0^{-2}\,p_0(\bar{x}_s^{-1/2})^{-4}.
\end{equation*}
The first term is therefore at most
\begin{equation}\label{8.1:eq-irrelevant-remainder-2}
  \mathfrak{M}\cdot 7C_b\,C_0^{-2}
  \left[\frac{p_0(g_{k_0})}{p_0(\bar{x}_s^{-1/2})}\right]^4 .
\end{equation}
The terms of the type of $\mathcal{E}_a$ are bounded exactly as the corresponding terms in the bound
of the Gaussian-extraction remainder $R^{\mathrm{GX}}$ proved above.

The ratio in \eqref{8.1:eq-irrelevant-remainder-2} is at most $1$, because the monotonicity runs in
the right direction. By the upper envelope bound of \eqref{8.1:eq-coupling-envelope-a},
$x_{k_0}=x_{K-s}\le\bar{x}_s$, which is the same as $g_{k_0}\ge\bar{x}_s^{-1/2}$. Hence
\begin{equation*}
  \log g_{k_0}^{-2}=\log x_{k_0}\le\log\bar{x}_s=\log\bigl(\bar{x}_s^{-1/2}\bigr)^{-2}.
\end{equation*}
The degree function $p_0(g')=A_0(\log g'^{-2})^{p_0}$ is increasing in $\log g'^{-2}$. It follows that
$p_0(g_{k_0})\le p_0(\bar{x}_s^{-1/2})$.

So the first term of \eqref{8.1:eq-irrelevant-remainder-1} is at most $7C_bC_0^{-2}\mathfrak{M}$.
Since $7/448=1/64$, this is at most $\mathfrak{M}/64$ once $C_0^2\ge 448C_b$. The constants are
ordered as in \eqref{8.1:eq-constants-ceiling-c}. The terms of the type of $\mathcal{E}_a$, estimated
as for $R^{\mathrm{GX}}$, are at most $\mathfrak{M}/64$ under the same order of constants; adding the
two bounds gives
\begin{equation}\label{8.1:eq-irrelevant-remainder-3}
  R^{\mathrm{irr}}\le\frac{\mathfrak{M}}{32}\qquad\text{once } C_0^2\ge 448\,C_b .
\end{equation}
The constant $1/32$ is independent of $g$: the condition $C_0^2\ge448C_b$ belongs to the order of
constants \eqref{8.1:eq-constants-ceiling-c}, and the ratio in \eqref{8.1:eq-irrelevant-remainder-2}
is at most $1$ for every $g$ because the room function is read at the envelope $\bar{x}_s$.
\end{proof}

\begin{remark}[The shadow terms when healed labels are kept]\label{8.1:rem-shadow-terms}
In Remark~\ref{8.1:rem-healed-summed-out} every healed label is summed out of the history index,
so that a history $h$ contains no deep healed label, and each of the covariances bounded
below is zero. In this remark the history index $h$ is instead the refinement of the
history decomposition that is not resummed, so that a healed label, once present, stays in
$h$. Deep healed labels under $N_i$ then occur in $h$, and the covariances below do not vanish.

\emph{Notation.} Let $i\in\{1,2\}$. We write
$\Phi_i=m_i+G_i+\mathfrak{S}^{\flat}_i+\mathfrak{L}_i$ for the decomposition
\eqref{8.1:eq-main-term-b}, and $m_i=m_i(h)$ for its first part. Here $\Phi_i$,
$\mathfrak{S}^{\flat}_i$ and $\mathfrak{L}_i$ stand for $\Phi^h_i$, $\mathfrak{S}^{\flat,h}_i$
and $\mathfrak{L}^h_i$, the superscript $h$ being suppressed where the history is not
displayed. A letter $Y$ denotes a
potential deep label under $N_1$, that is, an element of $\mathcal{Y}_1$, and $Y'$ one under
$N_2$. We write $j=j(Y)$ for the birth level of $Y$, $j'=j(Y')$, $\nu=k_0-j$ and $\nu'=k_0-j'$.
Further, $\sigma_Y$ is the amplitude of the piece of $\Phi^h_1$ attached to $Y$ in
\eqref{8.1:eq-logarithmic-derivative-c}, and $\mathring{\sigma}_Y$ is the amplitude of the
interior piece. We write $\mathfrak{S}^{\flat,h}_{2,Y'}$ for the part of
$\mathfrak{S}^{\flat,h}_2$ carried by $Y'$, and
$\tilde{\nu}(Y')=\tilde{\nu}(Y'\in h)$. The letters $\mu^{\sharp}_Y$, $\chi_Y$, $B_{\times}$
and $E^{\natural}$ are those of the flat shell bound, and $\chi$ denotes the chain-decay
member of its clause~(b). The constants $\mathfrak{M}$, $c_M$, $\varepsilon$, $x$ and
$\mathrm{polylog}$ are those of \eqref{8.1:eq-main-term-a}, and $C_G$ is the constant in the
size $|G_i|\le C_G\,\mathfrak{n}(f_i)\,\varepsilon^2$ of the part $G_i$ of the decomposition
\eqref{8.1:eq-main-term-b}, recorded with the main term.
The properties (M1), (M2) and (M3) are those stated, together with the source lemma, for the
conditional sup-expectation $E^{\natural}$.

\emph{Hypotheses.}
\begin{itemize}
\item[(A)] Clauses (a) and (b) of the flat shell bound hold.
\item[(B)] Clause (1) of the large-field block bound holds in its conditional form.
\item[(C)] The lemma supplying the estimate of order $\eta^2$, where $\eta$ is the bare
spacing $L^{-k_0}$ in units of the witness level, holds.
\item[(D)] The sizes \eqref{8.1:eq-logarithmic-derivative-c} hold.
\item[(E)] Clause (S3) of the source lemma holds, in the form of that lemma on which the main
term \eqref{8.1:eq-main-term-a} rests.
\end{itemize}
Hypothesis (E) is used only to bound the interior labels' pieces with the factor
$C_{\Theta}\bar{\Theta}$. No other clause or form of the source lemma is used.

\emph{Conclusion.} Under (A)--(E) the following bounds hold:
\begin{align}
|\operatorname{Cov}_{\tilde\nu}(m_1,m_2)|
&\le\Bigl[\sum_Y\sigma_Y\tilde q_j\mu^{\sharp}_Y\Bigr]^2+\chi
\label{8.1:eq-shadow-terms-1}\\
&\le C'^2\rho^8\|f_1\|\|f_2\|r^{-2}e^{-2cp_0(g_{k_0})}\notag\\
&\qquad+C\rho^8\|f_1\|\|f_2\|r^{-2}\tilde q^{\,2}e^{-c'\rho/(400M\check{R})},\notag\\
|\operatorname{Cov}_{\tilde\nu}(m_1,G_2)|
&\le\sum_Y\sigma_Y\tilde q_j(\mu^{\sharp}_Y+\chi_Y)\,C_G\,\mathfrak{n}(f_2)\,\varepsilon^2,
\label{8.1:eq-shadow-terms-2}\\
|\operatorname{Cov}_{\tilde\nu}(m_1,\mathfrak{S}^{\flat}_2)|
&\le C_{LF}C_{\Theta}\bar{\Theta}
\Bigl[\sum_Y\sigma_Y\tilde q_j(\mu^{\sharp}_Y+\chi_Y)\Bigr]
\Bigl[\sum_{Y'}\mathring{\sigma}_{Y'}\tilde q_{j'}\Bigr],
\label{8.1:eq-shadow-terms-3}\\
|\operatorname{Cov}_{\tilde\nu}(G_1,\mathfrak{S}^{\flat}_2)|
&\le\sum_{Y'}2\sup|G_1|\,C_{\Theta}\bar{\Theta}\,\mathring{\sigma}_{Y'}\,\tilde{\nu}(Y'),
\label{8.1:eq-shadow-terms-4}\\
|\operatorname{Cov}_{\tilde\nu}(\mathfrak{S}^{\flat}_1,\mathfrak{S}^{\flat}_2)|
&\le(C_{LF}+1)C_{\Theta}^2\bar{\Theta}^2\sum_Y\sum_{Y'}
\mathring{\sigma}_Y\mathring{\sigma}_{Y'}\tilde q_j\tilde q_{j'}.
\label{8.1:eq-shadow-terms-5}
\end{align}
The effects transmitted through the background are contained in the quantity
$\mu^{\sharp}$ of clause~(a) of the flat shell bound and are counted once there.

Each right-hand side in \eqref{8.1:eq-shadow-terms-1}--\eqref{8.1:eq-shadow-terms-5},
divided by $\mathfrak{M}$, is at most
\begin{equation*}
\frac{C_{\Theta}^2\bar{\Theta}^2\cdot\mathrm{polylog}\cdot x^2e^{-cp_0(g_{k_0})}}
{\vartheta^2r^2c_M}.
\end{equation*}
Under the coupling-ceiling condition of \eqref{8.1:eq-constants-ceiling-b}, the sum of these
covariances is at most $\mathfrak{M}/32$.
\end{remark}

\begin{proof}
\emph{The common count.} Fix a level $j=k_0-\nu$. The deep labels of level $j$ under $N_1$
number at most $C(\rho L^{\nu})^4(M\check{R}_j)^c$. Each term of the sums above carries a
factor $\tilde q_jL^{-4\nu}x_j\,\mathrm{polylog}$. The product of the count and this factor
is therefore at most
\begin{equation*}
C\rho^4(M\check{R}_j)^c\,\tilde q_j\,x_j\,\mathrm{polylog},
\end{equation*}
since $(\rho L^{\nu})^4L^{-4\nu}=\rho^4$. Summed over the birth levels $j$, these bounds give
at most $C'\rho^4e^{-cp_0(g_{k_0})}$, by the level-sum estimate for the birth-level rarity
weights $\tilde q_j$.

The power $L^{4\nu}$ coming from the count is cancelled exactly by $L^{-4\nu}$. No spare
power of $L$ is left over and none is needed. Without the factor $L^{-4\nu}$ the sum
diverges, by clause (a) of the counting lemma for deep healed labels. By symmetry the same
count holds for the labels $Y'$ under $N_2$.

\emph{The functional reading of clause (b).} Clause (b) of the flat shell bound is used for
functionals. Let $\Xi_1$ and $\Xi_2$ be bounded functionals. Suppose $\Xi_i$ is measurable
with respect to the configuration of deep labels under $N_i^{++}$ and of $V$ there, with
$V$ global for the $\mu^{\sharp}$-member. Then $|\operatorname{Cov}_{\tilde\nu}(\Xi_1,\Xi_2)|$
is bounded by the right-hand sides of clauses (a) and (b) summed over the presence atoms,
with the conditional sup-expectation $E^{\natural}$ in place of the expectation of the
functional. The functional $m_i(h)$ is of this kind.

\emph{Bound \eqref{8.1:eq-shadow-terms-1}.} Apply clause (b) in this functional form to the
pair $(m_1,m_2)$. This gives the first inequality. In the second, the first term is obtained
by applying the common count to each factor of the square, while the $\chi$-member of
clause (b) contributes the second term, with its chain-decay factor
$e^{-c'\rho/(400M\check{R})}$.

The contribution of the spanning event $B_{\times}$ is included in this estimate and is
bounded once, by clause (b) applied to the pair $(m_1,m_2)$. Bounding it separately for each
side would set two sums, each divergent by the counting lemma, against a single small factor.

\emph{Bound \eqref{8.1:eq-shadow-terms-2}.} Apply clause (a) of the flat shell bound with
$G_2$ as the bounded functional, with $V$ global. This is admissible because $G_2$ belongs
to the class of bounded functionals of Lemma~\ref{8.1:lem-healed-label-covariance}, namely
those measurable with respect to the coarse $\sigma$-algebra $\mathcal{C}$, which contains
$\sigma(V_{k_0})$. The size of $G_2$,
namely $|G_2|\le C_G\mathfrak{n}(f_2)\varepsilon^2$, is the one recorded with the main term
for the part $G_2$ of the decomposition \eqref{8.1:eq-main-term-b}.

\emph{Bound \eqref{8.1:eq-shadow-terms-3}.} Apply clause (a) of the flat shell bound with
the functional $\mathbf{1}_{Y'\in h}\,\mathfrak{S}^{\flat,h}_{2,Y'}$. By hypothesis (E) this
functional is bounded by $C_{\Theta}\bar{\Theta}\mathring{\sigma}_{Y'}$. It is measurable
with respect to $\mathcal{C}\vee\sigma(Y')$. The conditional sup-expectation
satisfies
\begin{equation*}
E^{\natural}\le C_{LF}C_{\Theta}\bar{\Theta}\mathring{\sigma}_{Y'}\tilde q_{j'}
\end{equation*}
by property (M3). Summing over $Y$ and $Y'$ gives
\eqref{8.1:eq-shadow-terms-3}.

\emph{Bounds \eqref{8.1:eq-shadow-terms-4} and \eqref{8.1:eq-shadow-terms-5}.} The first
follows label by label from property (M1), with the bound
$C_{\Theta}\bar{\Theta}\mathring{\sigma}_{Y'}$ on the piece carried by $Y'$.

The second follows from property (M2). No decoupling between the two sides is needed,
because the factor $L^{-4\nu'}$ is already contained in $\mathring{\sigma}$.

\emph{Effects through the background.} These effects are contained in the quantity
$\mu^{\sharp}$ of clause (a). They were therefore counted, once, in each of the bounds where
clause (a) was applied.

\emph{Relative size.} By \eqref{8.1:eq-main-term-a},
\begin{equation*}
\mathfrak{M}\ge c_M\vartheta^2g_{k_0}^4\|f_1\|\|f_2\|,
\end{equation*}
and $g_{k_0}^{-4}=x^2$ since $x=x_{k_0}=g_{k_0}^{-2}$. The quantities $\rho$ and $r$ are
$\mathrm{polylog}$. Hence each of the right-hand sides above, divided by $\mathfrak{M}$, is
at most
\begin{equation*}
\frac{C_{\Theta}^2\bar{\Theta}^2\cdot\mathrm{polylog}\cdot x^2e^{-cp_0(g_{k_0})}}
{\vartheta^2r^2c_M}.
\end{equation*}
The coupling-ceiling condition of \eqref{8.1:eq-constants-ceiling-b} implies condition
$(\gamma\text{-}1)$ of \eqref{8.1:eq-constants-ceiling-a} with the additional factor
$C_{\Theta}^2\bar{\Theta}^2$. Under it, the sum of these quotients is at most $1/32$.
Multiplying by $\mathfrak{M}$ shows that the total is at most $\mathfrak{M}/32$.
\end{proof}

\begin{proof}[Proof of the two-point witness theorem, concluded: order of the constants and the
coupling ceiling]\label{8.1:prf-constants-ceiling}
We first collect the terms of the total covariance formula \eqref{8.1:eq-total-covariance-a}.
Its two terms are the expectation of $\Psi^h$ and the covariance of $\Phi^h_1$ and $\Phi^h_2$.
With $\Phi^h_i$ split as in \eqref{8.1:eq-main-term-b}, these two terms are distributed among the
remainders of the preceding parts of the proof as follows:
\begin{equation}\label{8.1:eq-constants-ceiling-1}
\begin{aligned}
\mathbb{E}_{\tilde{\nu}}\Psi^h
&=R^{\mathrm{loc}}+R^{\Psi,\mathrm{live}},\\
\operatorname{Cov}_{\tilde{\nu}}\bigl(\Phi^h_1,\Phi^h_2\bigr)
&=\bigl[\mathfrak{M}+R^{\mathrm{prof}}+R^{\mathrm{GX}}+R^{\mathrm{irr}}\bigr]
+R^{\mathfrak{L}}+R^{m\mathfrak{L}}+(\mathrm{sh})'.
\end{aligned}
\end{equation}
The terms in which a live large-field factor is paired with any other factor are contained in
$R^{\Psi,\mathrm{live}}$. The covariance $\operatorname{Cov}(m_1,m_2)$ of the first components in
the split \eqref{8.1:eq-main-term-b} is the first of the shadow terms of
Remark~\ref{8.1:rem-shadow-terms}, and is therefore a part of $(\mathrm{sh})'$. The sum
$R^{\mathfrak{L}}+R^{m\mathfrak{L}}+R^{\Psi,\mathrm{live}}$ is the remainder
$R^{\mathcal{B}\prime}$ bounded in the part of this proof treating the large-field remainders.
Hence,
writing $S^T_{2;K}:=S^T_{2;K,m}(f_1,f_2)$ for the connected two-point function of the witness
pair, the torus index $m$ being fixed, the difference $S^T_{2;K}-\mathfrak{M}$ is bounded in
absolute value by the sum of the absolute values of the six remainders
\begin{equation*}
R^{\mathrm{prof}},\quad R^{\mathcal{B}\prime},\quad R^{\mathrm{loc}},\quad R^{\mathrm{GX}},\quad
R^{\mathrm{irr}},\quad (\mathrm{sh})'.
\end{equation*}

It remains to fix the constants on which the bounds of these remainders depend, and to check that
each remainder is at most $\mathfrak{M}/32$ for every depth $s$. The constants are fixed in the
following order; each choice reads only $L$, $M$, $\mathfrak{p}$, $N$, $\vartheta$ and the constants
of the estimates established above, and none reads $g$, $K$, $m$ or $s$.

First, $C_0$ is chosen subject to
\begin{equation}\label{8.1:eq-constants-ceiling-c}
\begin{aligned}
C_0\ge\max\Bigl\{&\frac{40M}{\vartheta},\ 4C_hM,\ 16C_hM,\ \frac{C_hM}{\vartheta},\
\frac{2(C_d+2)}{\vartheta},\
400M\sup_y\frac{\check{R}(y)^4}{p_0(y^{-1/2})^2},\ 3200M,\\
&C_0^{\mathrm{ker}},\ C_0^{\mathrm{loc}},\ C_0^{\mathrm{prof}},\ C_0^{\mathrm{src}},\
(448C_b)^{1/2},\ 24M,\ 2\Bigr\}.
\end{aligned}
\end{equation}
Here $C_0^{\mathrm{prof}}$ is the lower bound on $C_0$ above which the profile remainder is at most
$\mathfrak{M}/32$ (in the part of this proof treating the profile remainder), $(448C_b)^{1/2}$ is
the lower bound on $C_0$ needed for the irrelevant remainder, that is, $C_0^2\ge 448C_b$
(in the part of this proof treating the irrelevant remainder),
$C_0^{\mathrm{loc}}$ is the floor of the spanning local remainder
(in the part of this proof treating that remainder), $C_0^{\mathrm{ker}}$ is the floor required by
the kernel bounds, and $C_0^{\mathrm{src}}$ is the floor $(C_0\text{-src})$ of
\eqref{8.1:eq-constants-ceiling-b} below. The members $16C_hM$, $C_hM/\vartheta$, $3200M$ and
$C_0^{\mathrm{src}}$ are lower bounds, independent of $g$, on separations measured in cells; they
are of the same kind as the room function \eqref{8.1:eq-coupling-envelope-a} and the floor
$C_0^{\mathrm{loc}}$.

Next $r_0$ is fixed, then $C_D$, then $k_W(L)$.

Finally, with $X=g^{-2}$, $\gamma_W$ is defined as the largest $g\le\gamma_J$ for which the
following conditions hold: the conditions
\begin{equation}\label{8.1:eq-constants-ceiling-a}
\begin{aligned}
&(\gamma\text{-}0)\quad X\ge X_0;\\
&(\gamma\text{-}1)\text{--}(\gamma\text{-}3)\quad\text{three ceilings on $g$ of the form
\eqref{8.1:eq-constants-ceiling-3} below,}\\
&\qquad\text{required by the remainder estimates above;}\\
&(\gamma\text{-}4)\quad r(s)\le C\,\frac{p_0(\underline{x}_s^{-1/2})}{\log L}
\quad\text{for every } s;\\
&(\gamma\text{-}5)\quad\text{void};\\
&(\gamma\text{-}6)\quad L^{4\mathfrak{l}(y)}\bigl(M\check{R}(y)\bigr)^4e^{A_0}
\le e^{\frac12 c_q\,p_0(y^{-1/2})}\quad\text{for } y\ge X-c_5;\\
&(\gamma\text{-}7)\quad \mathfrak{d}_W(X-4c_2)\le\tfrac13;\\
&(\gamma\text{-}8)\quad \mu^{\sharp}_Y\le\tfrac12;
\end{aligned}
\end{equation}
and the two conditions
$(\gamma\text{-src})$,
$(\gamma\text{-}\varepsilon)$ of
\begin{equation}\label{8.1:eq-constants-ceiling-b}
\begin{aligned}
&(\gamma\text{-src})\quad
\sup_{y\ge X-c_5}\mathfrak{F}_{\mathrm{src}}(y)\le\frac{\vartheta^4c_M}{64\,C_{\mathrm{tot}}};\\
&(\gamma\text{-}\varepsilon)\quad
\sup_{y\ge X-c_5}\frac{3\lambda\bigl(\mathfrak{l}(\bar{x}(y))+1\bigr)+2c_2}{y\log 2}
\le 2^{1/p_0}-1;\\
&(C_0\text{-src})\quad C_0\ge C_0^{\mathrm{src}}.
\end{aligned}
\end{equation}
Condition $(\gamma\text{-}5)$ is void because the requirement on the depth is the depth floor
\eqref{8.1:eq-coupling-envelope-d} on $s$, which is assumed throughout; it is not a
condition on $g$.
In $(C_0\text{-src})$, $C_0^{\mathrm{src}}$ is a constant, independent of $g$ and $s$, such that
for every $C_0\ge C_0^{\mathrm{src}}$ and every $s$, with $\rho$ the support radius in cells
(which satisfies $\rho\ge C_W(s)$ by \eqref{8.1:eq-coupling-envelope-a}),
\begin{equation*}
C_{\mathrm{tot}}'\,\bar{\Theta}\,P(\bar{x}_s)\,\rho^{24}\,\bar{x}_s^2\,
e^{-\varkappa\rho/16}\le\frac{\vartheta^2c_M}{64}.
\end{equation*}
In every member the variable $y$ ranges over $y\ge X-c_5$, the range of $\underline{x}_s$ by
\eqref{8.1:eq-coupling-envelope-f}.
In these displays $\mathfrak{l}(y):=L\,\check{R}(y)+c_L$ is the level count
\eqref{8.1:eq-bad-event-b} read as a function of the inverse coupling $y$;
$\bar{x}(y)=(B^{\sharp}/\lambda^{\sharp})(y+c_5)+B^{\sharp}$ is the affine envelope function of
\eqref{8.1:eq-coupling-envelope-f}; $\bar{\Theta}:=L^{4(\mathfrak{l}(\bar{x}_s)+1)}$ is the envelope
age factor of the source lemma; $P(\bar{x}_s)$ denotes a fixed polynomial in $\log\bar{x}_s$, one
of the polylogarithmic factors listed where the main term is decomposed into remainders;
$\lambda$ is the
constant of the level window in Lemma~\ref{8.1:lem-running-normalisation}; $c_q$ is the constant of
the rarity weight $\tilde{q}_j$ in Lemma~\ref{8.1:lem-bad-event}; $C_h$ is the constant entering
the profile remainder through the factor $C_hM$ (see the part of this proof treating the profile
remainder); $C_D$ is the constant entering the cube size $D_{\square}(s)$ of the
Gaussian-extraction remainder (in the part of this proof treating that remainder);
$C_{\mathrm{tot}}'$ is a
summary constant of the remainder estimates, of the same kind as $C_{\mathrm{tot}}$; and
\begin{equation}\label{8.1:eq-constants-ceiling-2}
\mathfrak{F}_{\mathrm{src}}(y):=L^{8(\mathfrak{l}(\bar{x}(y))+1)}\,
\mathfrak{p}_{\mathrm{tot}}\bigl(\bar{x}(y)\bigr)\,\bar{x}(y)^2\,
\exp\bigl(-c_{\mathrm{src}}A_0(\log y)^{p_0}\bigr).
\end{equation}
The floor $(C_0\text{-src})$ is a member of \eqref{8.1:eq-constants-ceiling-c} and is met when
$C_0$ is fixed; the existence of $C_0^{\mathrm{src}}$ is shown below. We now go through the members
whose form is not evident from the display.

In $(\gamma\text{-}0)$, $X_0$ denotes the largest of the lower thresholds on the inverse coupling
required by the running normalisation lemma
(Lemma~\ref{8.1:lem-running-normalisation}). Condition $(\gamma\text{-}4)$ is the side condition
of the Gaussian extraction, with the constant $C$ of that side condition; its left side grows like
$\log\log\bar{x}_s$, while its right side is $C A_0(\log\underline{x}_s)^{p_0}/\log L$, since
$p_0(y^{-1/2})=A_0(\log y)^{p_0}$ by the definition of the degree function $p_0(\cdot)$.

In $(\gamma\text{-}6)$, Ba{\l}aban's radius function satisfies
$\check{R}(x)\le C_R(\log x)^{r_{\check{R}}}+1$, with the constant $C_R$ and the exponent
$r_{\check{R}}$ of the growth bound for $\check{R}$ in the correlation theorem; $r_{\check{R}}$ is
the exponent written $r$ in the condition $p_0\ge 5r$ stated in the glossary
(Notation~\ref{0.2:not-glossary}). Hence
\begin{equation*}
L^{4\mathfrak{l}(y)}
\le\exp\Bigl(4\log L\,\bigl(L C_R(\log y)^{r_{\check{R}}}+L+c_L\bigr)\Bigr),
\end{equation*}
whereas the exponent on the right of $(\gamma\text{-}6)$ is $\frac12 c_qA_0(\log y)^{p_0}$. The
remaining factor $(M\check{R}(y))^4e^{A_0}$ on the left of $(\gamma\text{-}6)$ is at most a
constant times a polynomial in $\log y$. Since the auxiliary exponent satisfies
$p_0\ge 5r_{\check{R}}$
\cite{B-CMP119}, the power
$(\log y)^{p_0}$ wins. The threshold in $y$ beyond which $(\gamma\text{-}6)$ holds depends on $L$;
this is admissible, because $\gamma_W$ is chosen after $L$.

In $(\gamma\text{-}7)$, $\mathfrak{d}_W$ is the rate function \eqref{8.1:eq-running-normalisation-a}
of the running normalisation lemma. Condition $(\gamma\text{-}8)$ is used only when healed labels
are kept, for the shadow terms of Remark~\ref{8.1:rem-shadow-terms}. In the bound for
$\mu^{\sharp}_Y$ used there, the factor $L^{4\mathfrak{l}(y)}$ dominates $y$ times a polynomial in
$\log y$, so that $\mu^{\sharp}_Y\le\frac12$ for every $Y$ once $y$ is large; level differences
$\nu\ge\mathfrak{l}$ only improve the bound.

Condition $(\gamma\text{-src})$ is a ceiling of the same kind as $(\gamma\text{-}1)$, and it
implies $(\gamma\text{-}1)$. It is written in the single envelope function $\bar{x}(\cdot)$ of
\eqref{8.1:eq-coupling-envelope-f}. Three facts reduce the corresponding remainder bounds to the
function \eqref{8.1:eq-constants-ceiling-2}.
\begin{itemize}
\item Every polynomial in $\log\bar{x}_s$ is at most a polynomial in $\log y$, since
$\bar{x}_s\le\bar{x}(y)$ and $\bar{x}$ is affine.
\item The bound $\bar{\Theta}^2\le L^{8(\mathfrak{l}(\bar{x}(y))+1)}$ holds by the monotonicity of
$\check{R}$.
\item The bound
$e^{-c_{\mathrm{src}}p_0(g_{k_0})}\le e^{-c_{\mathrm{src}}A_0(\log\underline{x}_s)^{p_0}}$ holds
because $p_0(g_{k_0})=A_0(\log x_{k_0})^{p_0}$ with $x_{k_0}=g_{k_0}^{-2}\ge\underline{x}_s$
by \eqref{8.1:eq-coupling-envelope-f}.
\end{itemize}
Taking logarithms in \eqref{8.1:eq-constants-ceiling-2},
\begin{equation*}
\log\mathfrak{F}_{\mathrm{src}}(y)
=8\log L\,\bigl(\mathfrak{l}(\bar{x}(y))+1\bigr)
+\log\bigl(\mathfrak{p}_{\mathrm{tot}}(\bar{x}(y))\,\bar{x}(y)^2\bigr)
-c_{\mathrm{src}}A_0(\log y)^{p_0}.
\end{equation*}
The first term is at most
$8\log L\,\bigl(LC_R(\log\bar{x}(y))^{r_{\check{R}}}+L+c_L+1\bigr)$, and the second is
$O(\log y)$. Since $p_0\ge 5r_{\check{R}}$ and $p_0>1$, the last term dominates. Hence
$\mathfrak{F}_{\mathrm{src}}$ tends to zero, and it is decreasing for $y\ge y_0$, where $y_0$
depends on $L$, $\mathfrak{p}$ and the constants fixed before. Thus $(\gamma\text{-src})$ is a
condition $g\le\gamma_{\mathrm{src}}$, with $\gamma_{\mathrm{src}}$ depending on $L$, $M$,
$\mathfrak{p}$, $N$, $\vartheta$, $C_0$ and the constants fixed before it. It is one ceiling for
every $s$.

Condition $(\gamma\text{-}\varepsilon)$ serves the constant $C_\Theta$ of the source lemma. Put
$\varepsilon_j:=g_j\,p_0(g_j)$, so that $\varepsilon_{k_0}$ is the number $\varepsilon$ of
\eqref{8.1:eq-main-term-a}. Since $p_0(g_j)=A_0(\log x_j)^{p_0}$,
\begin{equation*}
\Bigl(\frac{\varepsilon_j}{\varepsilon_{k_0}}\Bigr)^2
=\frac{x_{k_0}}{x_j}\Bigl(\frac{\log x_j}{\log x_{k_0}}\Bigr)^{2p_0}.
\end{equation*}
Let $j$ be one of the levels at which the live large-field regions entering the source lemma were
created. For $X\ge 4c_2$ one has $x_{k_0}/x_j\le 2$ on these levels. Together with this,
$(\gamma\text{-}\varepsilon)$ yields $(\varepsilon_j/\varepsilon_{k_0})^2\le 8$ on these levels.
The left side of $(\gamma\text{-}\varepsilon)$ is
the supremum of a function decreasing to zero. The factor $\log 2$ in its denominator enters
through the lower bound $\log\underline{x}_s\ge\log 2$.

The floor $C_0^{\mathrm{src}}$ of $(C_0\text{-src})$ exists, is independent of $g$, and is uniform
in $s$. Indeed, by \eqref{8.1:eq-coupling-envelope-a},
\begin{equation*}
\rho\ge C_W(s)=C_0\,p_0(\bar{x}_s^{-1/2})^2=C_0A_0^2(\log\bar{x}_s)^{2p_0},
\end{equation*}
while
\begin{equation*}
\log\bigl(\bar{\Theta}\,P(\bar{x}_s)\,\rho^{24}\,\bar{x}_s^2\bigr)
=O\bigl(L(\log\bar{x}_s)^{r_{\check{R}}}\bigr)+O(\log\bar{x}_s)
+O\bigl(\log C_0+\log\log\bar{x}_s\bigr).
\end{equation*}
Since $2p_0>r_{\check{R}}$ and the left side depends on $C_0$ only logarithmically, this is at
most $\varkappa C_0A_0^2(\log\bar{x}_s)^{2p_0}/32$ for $C_0$ large and all
$\bar{x}_s\ge X\ge\gamma_J^{-2}$. That bound equals $\varkappa C_W(s)/32\le\varkappa\rho/32$. Hence
the left side of the inequality in $(C_0\text{-src})$ is at most
\begin{equation*}
C_{\mathrm{tot}}'\,e^{\varkappa\rho/32-\varkappa\rho/16}=C_{\mathrm{tot}}'\,e^{-\varkappa\rho/32}
\le C_{\mathrm{tot}}'\,e^{-\varkappa C_0A_0^2(\log\bar{x}_s)^{2p_0}/32}.
\end{equation*}
For $C_0$ large this is at most $\vartheta^2c_M/64$, uniformly in $s$, because
$\bar{x}_s\ge\gamma_J^{-2}$.

We now conclude. Each member of the list defining $\gamma_W$ is a condition of the form
\begin{equation}\label{8.1:eq-constants-ceiling-3}
\sup_{y\ge X-c_5}F(y)\le\mathrm{const},\qquad F\ \text{decreasing to zero}.
\end{equation}
The supremum in \eqref{8.1:eq-constants-ceiling-3} is taken over a set that shrinks as $X$ grows.
Hence a condition of this form that holds at some $X$ holds at every larger $X$, and all the
conditions hold for every $g\le\gamma_W$.

Each bound obtained for the remainders in the preceding parts of the proof is the product of two
factors. The first is decreasing in $x_{k_0}$, and is therefore at most its value at
$\underline{x}_s$, since $x_{k_0}\ge\underline{x}_s$ by \eqref{8.1:eq-coupling-envelope-f}. The
second is a polynomial in $\log\bar{x}_s$, and $\bar{x}_s$ is bounded by the affine function
$\bar{x}(\underline{x}_s)$. Each such bound is therefore at most a function of
$y=\underline{x}_s\ge X-c_5$
controlled by one of the members of the list defining $\gamma_W$, each of the form
\eqref{8.1:eq-constants-ceiling-3}. Consequently each remainder is at
most $\mathfrak{M}/32$ for all $s$ once $g\le\gamma_W$, that is, once $X\ge\gamma_W^{-2}$.

There is no circularity. The quantities $C_W(s)$, $\rho$, $r(s)$ and $D_{\square}(s)$ are
determined by $(g,s)$ once the constants have been fixed. The order of choice is: the members of
\eqref{8.1:eq-constants-ceiling-c} for $C_0$; then $r_0$ and $C_D$; then the members of the list
defining $\gamma_W$. No choice returns to $L$, $K$, $m$ or $s$.

Therefore, by the bounds obtained in the preceding parts of this proof for the profile remainder,
the large-field remainders, the spanning local remainder, the Gaussian-extraction remainder and
the irrelevant remainder, and by Remark~\ref{8.1:rem-shadow-terms}, under the conditions just
fixed,
\begin{equation}\label{8.1:eq-constants-ceiling-4}
\begin{aligned}
\bigl|S^T_{2;K}-\mathfrak{M}\bigr|
&\le|R^{\mathrm{prof}}|+|R^{\mathcal{B}\prime}|+|R^{\mathrm{loc}}|+|R^{\mathrm{GX}}|
+|R^{\mathrm{irr}}|+|(\mathrm{sh})'|\\
&\le\frac{6\,\mathfrak{M}}{32}<\frac{\mathfrak{M}}{4}.
\end{aligned}
\end{equation}
In \eqref{8.1:eq-constants-ceiling-4} the profile remainder contributes at most $\mathfrak{M}/32$,
and the five other remainders contribute at most $\mathfrak{M}/32$ each. Hence the remainder $R_K$
of the two-point witness theorem, which measures $S^T_{2;K}-\mathfrak{M}$ relative to
$\mathfrak{M}$, satisfies $|R_K|\le\frac14\le\frac12$.

All conditions imposed above are one-sided. The members $(\gamma\text{-src})$ and
$(\gamma\text{-}\varepsilon)$, like the other members of the list defining $\gamma_W$, are suprema
over $y\ge X-c_5$ of functions decreasing to zero; they are ceilings on $g$, uniform in $s$, $K$
and $m$. The members $(C_0\text{-src})$, $16C_hM$, $C_hM/\vartheta$ and $3200M$ are floors,
independent of $g$, on separations measured in cells. The factor $\bar{\Theta}$ is not a condition
on anything and in particular not a cap in the sense of Definition~\ref{2.3:def-cap}: it is a factor
inside a remainder, paid for by the factor $e^{-c_{\mathrm{src}}p_0(g_{k_0})}$ through
$(\gamma\text{-src})$. The
blocking factor $L$ enters the constant $C_\Theta$ and the other constants of the source lemma, and
the factor $\bar{\Theta}$, only as a parameter, and no lower bound on $L$ reads any of them. The
ages of live
large-field regions enter only through the level count
$\mathfrak{l}$, and the depth of deep labels is unbounded and summed in the annealed measure.
Nothing bounds $L$, $K$, $m$, $s$, the ages or the depth from above, or a coupling from below.
\end{proof}

\begin{remark}[What the two-point witness theorem does not assert]\label{8.1:rem-witness-not-asserted}
The two-point witness theorem (Theorem~\ref{8.1:thm-two-point-witness}) asserts its displayed
conclusion and nothing more. In particular, it does not assert any of the following.
\begin{itemize}
\item The existence of a state on the observable algebra of clause (iii)(a) of
Theorem~\ref{3.1:thm-main}. That sub-clause is proved as an implication from two statements
named as its antecedents in Theorem~\ref{3.1:thm-main}, and
Theorem~\ref{8.1:thm-two-point-witness} is not one of them.
\item That the law of the smeared curvature field at finite cutoff, or any limit of it, differs
from a Gaussian law.
\item Anything about a single observable, that is, about the one-point function ($n=1$ in
$S^T_{n;K,m}$).
\item Anything about test functions with coinciding supports.
\item Anything about pairs whose supports lie at distance greater than $6\mathsf{d}$.
\item Anything about test functions that change sign.
\item Anything about the limit $\vartheta\to0$. The parameter $\vartheta\in\,]0,1]$ is fixed in
the theorem, and its constants depend on $\vartheta$.
\item The case $N\ge3$. The theorem is stated for $\mathrm{SU}(2)$, and its formulas are written
for $\mathrm{SU}(N)$; the case $N\ge3$ requires a further statement beyond those on which the
theorem rests.
\item The sharp form of the estimate, that is, the Gaussian extraction in relative precision.
The theorem does give the bound \eqref{8.1:eq-running-normalisation-a} on the running
normalisation at the witness level, by which $|N_{K-s}-1|$ is at most the rate $\mathfrak{d}_W$
displayed there, uniformly in $K$; convergence of $N_{K-s}$ as $K\to\infty$ at fixed $s$ is not
asserted.
\item Any conclusion below the depth floor $s(\mathsf{d})\ge s_W(g)$ of
Theorem~\ref{8.1:thm-two-point-witness}; the level condition $K-s(\mathsf{d})\ge k_W$ is likewise
not removed.
\item Invariance under $O(4)$.
\item A mass gap.
\item A bound on the slope per single renormalisation step.
\item A bound on $\Phi^h_i-m_i(h)$ that is uniform over all histories $h$ and all block fields
$V$. Here $\Phi^h_i$ is the first logarithmic derivative of the history term in the $i$-th source
variable, and $m_i(h)$ is the history-dependent term of the source lemma that is subtracted from
$\Phi^h_i$ there.
\end{itemize}
\end{remark}

\begin{corollary}[Two-sided two-point bounds and non-triviality]\label{8.1:cor-two-sided-bounds}
Work in the setting of the two-point witness theorem,
Theorem~\ref{8.1:thm-two-point-witness}: the quantifier prefix of Theorem~\ref{3.1:thm-main},
referred to below as Theorem~0 (Notation~\ref{2.1:not-prefix}), and the regime of the
correlation theorem, that is, a number $\varrho_J>0$, the threshold $\gamma_J$ which the
correlation theorem attaches to $\varrho_J$, and $g\le\gamma_J$.
Assume the following.
\begin{itemize}
\item The conclusion of the two-point witness theorem,
Theorem~\ref{8.1:thm-two-point-witness}.
\item Clause (0) of Theorem~0, in the form of the envelope
\eqref{8.1:eq-coupling-envelope-f} of the running inverse coupling
$x_{K-s}=g_{K-s}^{-2}$ (Definition~\ref{8.1:def-coupling-envelope}).
\item Clause (ii) of Theorem~0: the limits of its item (b) along its sequence
$K_i\to\infty$, with the sub-items (b1)--(b2), and its statement that the curvature Hilbert
space $\mathcal{H}^{\mathrm{curv}}$ has dimension at least two.
\item The kernel bounds \eqref{8.1:eq-kernel-bounds-b} of
Lemma~\ref{8.1:lem-kernel-bounds}.
\end{itemize}
Fix $\vartheta\in\,]0,1]$. There are a threshold $\gamma_W\le\gamma_J$ and a level $k_W$,
namely those of Theorem~\ref{8.1:thm-two-point-witness},
and constants $c_{\mathcal{K}},C_{\mathcal{K}}>0$, namely those of
Lemma~\ref{8.1:lem-kernel-bounds}, all independent of $K$, $g_0$,
$m$, $s$, $\mathsf{d}$ and of the position of the supports, together with the room function
$C_W(s)$ of Notation~\ref{8.1:not-test-pairs-kernel}, given by
\eqref{8.1:eq-coupling-envelope-a}, and a depth floor $s_W(g)$, with the following property.
The floor $s_W(g)$ is a one-sided lower bound
on the depth, depending on $g$ at most through $\log_L g^{-2}$; it is the floor
\eqref{8.1:eq-coupling-envelope-d}, of the form
$\max\{O(\log_L\log g^{-2}),\,m_{\mathbb{T}}(g)\}$, and no assertion is made that this floor,
or its dependence on $g$, can be removed. For every $g\le\gamma_W$, every
$g_0\in\,]0,g_{-}(g)]$, $K=K(g_0,g)$, every torus
exponent $m$ admitted by the torus clause of the prefix, every
$\mathsf{d}$ whose depth $s(\mathsf{d})$ (Notation~\ref{8.1:not-test-pairs-kernel}) satisfies
\begin{equation*}
s_W(g)\le s(\mathsf{d})\le K-k_W
\end{equation*}
and, when Theorem~\ref{8.1:thm-two-point-witness} attaches its per-torus depth floor to the
exponent $m$, is not smaller than that floor,
and every $(f_1,f_2)\in\mathfrak{T}(\mathsf{d};\vartheta)$, one has, with $s:=s(\mathsf{d})$ and
with the torus exponent suppressed from the notation $S^T_{2;K}$:
\begin{enumerate}
\item[(a)] the representation
\begin{equation}\label{8.1:eq-two-sided-bounds-1}
S^T_{2;K}(f_1,f_2)=(\zeta'_{K-s})^2\,b_0\,g_{K-s}^4\,\mathcal{K}_s(f_1,f_2)\,(1+R_K),
\qquad |R_K|\le\tfrac12,\qquad \zeta'_{K-s}\in\bigl[\tfrac23,\tfrac43\bigr],
\end{equation}
together with the scale-free kernel bounds
\begin{equation}\label{8.1:eq-two-sided-bounds-2}
c_{\mathcal{K}}\,\vartheta^2\,\|f_1\|_\infty\|f_2\|_\infty\le\mathcal{K}_s(f_1,f_2)
\le 100\,C_{\mathcal{K}}\|f_1\|_\infty\|f_2\|_\infty ;
\end{equation}
\item[(b)] uniformly in $K$, $g_0$, $m$ and the position of the supports,
\begin{equation}\label{8.1:eq-two-sided-bounds-a}
\tfrac29\,b_0\,g_{K-s}^4\,\mathcal{K}_s(f_1,f_2)\le S^T_{2;K}(f_1,f_2)
\le\tfrac83\,b_0\,g_{K-s}^4\,\mathcal{K}_s(f_1,f_2),
\end{equation}
and $S^T_{2;K}(f_1,f_2)>0$;
\item[(c)] the bounds
\begin{equation}\label{8.1:eq-two-sided-bounds-3}
\begin{split}
\tfrac29\,b_0\,c_{\mathcal{K}}\,\vartheta^2\,\|f_1\|_\infty\|f_2\|_\infty\,
\bigl(g^{-2}+B^{\sharp}(s+1)\bigr)^{-2}
&\le S^T_{2;K}(f_1,f_2)\\
&\le\tfrac83\cdot100\,b_0\,C_{\mathcal{K}}\,\|f_1\|_\infty\|f_2\|_\infty\,
\bigl(g^{-2}+\lambda^{\sharp}s-c_5\bigr)^{-2};
\end{split}
\end{equation}
\item[(d)] every limit of $S^T_{2;K}$ provided by item (b) of clause (ii) of
Theorem~0, along its sequence $K_i\to\infty$ and in the sense of its sub-items (b1)--(b2),
which we denote by $S^T_2$ (the torus exponent being suppressed as above),
satisfies, for every $\mathsf{d}$ with $s(\mathsf{d})\ge s_W(g)$ and, when
Theorem~\ref{8.1:thm-two-point-witness} attaches its per-torus depth floor to $m$, with
$s(\mathsf{d})$ not smaller than that floor, and every
$(f_1,f_2)\in\mathfrak{T}(\mathsf{d};\vartheta)$,
\begin{equation}\label{8.1:eq-two-sided-bounds-4}
\begin{split}
\tfrac29\,b_0\,c_{\mathcal{K}}\,\vartheta^2\,\|f_1\|_\infty\|f_2\|_\infty\,
\bigl(g^{-2}+B^{\sharp}(s+1)\bigr)^{-2}
&\le S^T_2(f_1,f_2)\\
&\le\tfrac83\cdot100\,b_0\,C_{\mathcal{K}}\,\|f_1\|_\infty\|f_2\|_\infty\,
\bigl(g^{-2}+\lambda^{\sharp}s-c_5\bigr)^{-2};
\end{split}
\end{equation}
\item[(e)] $\dim\mathcal{H}^{\mathrm{curv}}\ge2$ (this is the third hypothesis above, recorded
here among the properties of the limits);
\item[(f)] the sup-normalised quantity $S^T_2(f_1,f_2)/(\|f_1\|_\infty\|f_2\|_\infty)$ is
strictly positive and tends to zero logarithmically in $1/\mathsf{d}$ as $\mathsf{d}\to0$,
uniformly over $(f_1,f_2)\in\mathfrak{T}(\mathsf{d};\vartheta)$, and hence along
sup-normalised dilates of a fixed pair; consequently no such limit is scale invariant, none is
the free Maxwell field at any charge, and
each is non-trivial in the sense of Definition~\ref{3.1:def-non-trivial}: not identically zero,
and the sup-normalised connected two-point function of the action density tends to zero with
the scale more slowly than every power of it;
\item[(g)] the bounds of (d), divided by $\|f_1\|_\infty\|f_2\|_\infty$, are those of the fourth
power of the running coupling at the depth
$s=s(\mathsf{d})$, with the slopes of clause (0):
\begin{equation*}
\tfrac29\,b_0\,c_{\mathcal{K}}\,\vartheta^2\,\bar{x}_s^{-2}
\le\frac{S^T_2(f_1,f_2)}{\|f_1\|_\infty\|f_2\|_\infty}
\le\tfrac{800}3\,b_0\,C_{\mathcal{K}}\,\underline{x}_s^{-2},
\end{equation*}
where $\underline{x}_s=g^{-2}+\lambda^{\sharp}s-c_5$ and $\bar{x}_s=g^{-2}+B^{\sharp}(s+1)$ are
affine in $s$ with the slopes $\lambda^{\sharp}$ and $B^{\sharp}$; when these envelopes are
written in terms of $\log_L(1/\mathsf{d})$ instead of $s(\mathsf{d})$, the additive constant of
$\bar{x}_s$ acquires the further term $B^{\sharp}\log_L C_W(s(\mathsf{d}))$, which is of lower
order than $\log_L(1/\mathsf{d})$ as $\mathsf{d}\to0$.
\end{enumerate}
The corollary does not assert that $\zeta'_{K-s}\to1$: concerning $\zeta'_{K-s}$ it uses only
the running normalisation of Lemma~\ref{8.1:lem-running-normalisation}, in the form
\eqref{8.1:eq-running-normalisation-a} read at the witness level,
\begin{equation*}
\limsup_{K\to\infty}|\zeta'_{K-s}-1|\le\mathfrak{d}_W\bigl(g^{-2}+\lambda^{\sharp}s-c_5\bigr),
\end{equation*}
a quantity that tends to zero as $g^{-2}+\lambda^{\sharp}s\to\infty$, and no bound of the form
$g^\sigma$ with $\sigma\ge2$ is asserted for it. Nor does the corollary cover the limit
$\vartheta\to0$, test functions of variable sign, or supports at mutual distance greater than
$6\mathsf{d}$.
\end{corollary}

\begin{proof}
\emph{The representation.} The quantifiers of the corollary are those of the two-point witness
theorem, Theorem~\ref{8.1:thm-two-point-witness}, and we take $\gamma_W$, $k_W$, $s_W(g)$ and
$C_W(s)$ from it; the floor is \eqref{8.1:eq-coupling-envelope-d} and the room function is
\eqref{8.1:eq-coupling-envelope-a}. Under these quantifiers that theorem gives
\eqref{8.1:eq-two-sided-bounds-1}, with $\zeta'_{K-s}=N_{K-s}\in[\frac23,\frac43]$ as in
\eqref{8.1:eq-running-normalisation-a}. The kernel bounds \eqref{8.1:eq-two-sided-bounds-2},
with constants $c_{\mathcal{K}}\vartheta^2$ and $100\,C_{\mathcal{K}}$ independent of the
scale, are
\eqref{8.1:eq-kernel-bounds-b} of Lemma~\ref{8.1:lem-kernel-bounds}. This proves (a).

\emph{Positivity of the kernel.} Each $f_i$ is non-negative and not identically zero, because
$(f_1,f_2)\in\mathfrak{T}(\mathsf{d};\vartheta)$; hence $\|f_i\|_\infty>0$. Since also
$\vartheta>0$, the lower bound in
\eqref{8.1:eq-two-sided-bounds-2} therefore gives $\mathcal{K}_s(f_1,f_2)>0$. Moreover
$b_0=(N^2-1)/2>0$ and $g_{K-s}>0$, so the factor $b_0\,g_{K-s}^4\,\mathcal{K}_s(f_1,f_2)$ in
\eqref{8.1:eq-two-sided-bounds-1} is strictly positive.

\emph{The constants $\frac29$ and $\frac83$.} From $\zeta'_{K-s}\in[\frac23,\frac43]$ and
$|R_K|\le\frac12$ we obtain $(\zeta'_{K-s})^2\in[\frac49,\frac{16}9]$ and
$1+R_K\in[\frac12,\frac32]$, and therefore
\begin{equation*}
\tfrac29=\tfrac49\cdot\tfrac12\le(\zeta'_{K-s})^2(1+R_K)\le\tfrac{16}9\cdot\tfrac32=\tfrac83 .
\end{equation*}
Multiplying by the strictly positive factor $b_0\,g_{K-s}^4\,\mathcal{K}_s(f_1,f_2)$ and using
\eqref{8.1:eq-two-sided-bounds-1} gives \eqref{8.1:eq-two-sided-bounds-a}.
The numbers $\frac29$ and $\frac83$ do not depend on $K$, $g_0$, $m$ or the position of the
supports, and \eqref{8.1:eq-two-sided-bounds-1} holds for all values of these admitted by the
quantifiers; so \eqref{8.1:eq-two-sided-bounds-a} is uniform in them. Its left-hand side is
strictly positive, whence $S^T_{2;K}(f_1,f_2)>0$. This proves (b).

\emph{The envelope.} By clause (0) of Theorem~0 in the form \eqref{8.1:eq-coupling-envelope-f},
\begin{equation*}
\underline{x}_s=g^{-2}+\lambda^{\sharp}s-c_5\le x_{K-s}\le g^{-2}+B^{\sharp}(s+1)=\bar{x}_s ,
\end{equation*}
with $x_{K-s}=g_{K-s}^{-2}>0$. Since the map $u\mapsto u^{-2}$ is decreasing on $]0,\infty[$, the
right-hand inequality gives $g_{K-s}^4\ge\bar{x}_s^{-2}$. The lower envelope is positive:
$\lambda^{\sharp}s\ge0$, and $g^{-2}\ge\gamma_W^{-2}>c_5$ by the choice of the coupling ceiling
in Definition~\ref{8.1:def-coupling-envelope}; so the left-hand inequality gives
$g_{K-s}^4\le\underline{x}_s^{-2}$. Inserting
these bounds and \eqref{8.1:eq-two-sided-bounds-2} into the two sides of
\eqref{8.1:eq-two-sided-bounds-a} gives \eqref{8.1:eq-two-sided-bounds-3}. This proves (c).

\emph{Limits.} The two envelopes $\bar{x}_s$ and $\underline{x}_s$ do not depend on $K$, and the
constants in \eqref{8.1:eq-two-sided-bounds-3} do not depend on $K$, $g_0$ or $m$. Fix
$\mathsf{d}$ with $s(\mathsf{d})\ge s_W(g)$ and, when Theorem~\ref{8.1:thm-two-point-witness}
attaches its per-torus depth floor to $m$, with $s(\mathsf{d})$ not smaller than that floor, and
fix $(f_1,f_2)\in\mathfrak{T}(\mathsf{d};\vartheta)$. These lower bounds are conditions on
$s(\mathsf{d})$ for the given $g$ and $m$ and do not involve $K$.
Since $K_i\to\infty$, the condition $s(\mathsf{d})\le K_i-k_W$ holds for all sufficiently large
$i$, so \eqref{8.1:eq-two-sided-bounds-3} holds along the tail of the sequence of item (b) of
clause (ii) of Theorem~0, with bounds independent of $i$. Passing to the limit along that
sequence, item (b) of clause (ii) with its sub-items (b1)--(b2) carries these bounds over to
the limit $S^T_2(f_1,f_2)$ with the same constants;
this is \eqref{8.1:eq-two-sided-bounds-4}.
This proves (d); in particular $S^T_2(f_1,f_2)>0$.

\emph{Dimension.} By the statement of clause (ii) of Theorem~0 that the curvature Hilbert space
has dimension at least two, $\dim\mathcal{H}^{\mathrm{curv}}\ge2$; this is the third hypothesis,
recorded in (e) without further argument. This proves (e).

\emph{Decay, scale invariance, free field.} Divide \eqref{8.1:eq-two-sided-bounds-4} by
$\|f_1\|_\infty\|f_2\|_\infty>0$. The resulting quotient lies between
$\frac29 b_0c_{\mathcal{K}}\vartheta^2\bar{x}_s^{-2}>0$ and
$\frac{800}3 b_0C_{\mathcal{K}}\underline{x}_s^{-2}$, with $s=s(\mathsf{d})$. The depth
$s(\mathsf{d})$ of Notation~\ref{8.1:not-test-pairs-kernel} increases without bound as
$\mathsf{d}\to0$, logarithmically in $1/\mathsf{d}$, and no upper bound on it is imposed in (d).
Since $\lambda^{\sharp}\ge\underline{\lambda}>0$, the upper bound is the inverse square of an
affine function of $s(\mathsf{d})$ with positive slope, and so is the lower bound, with slope
$B^{\sharp}$. Hence the quotient is strictly positive and tends to zero as $\mathsf{d}\to0$,
logarithmically in $1/\mathsf{d}$ and in particular more slowly than every power of
$\mathsf{d}$; the convergence is uniform over $(f_1,f_2)\in\mathfrak{T}(\mathsf{d};\vartheta)$,
since both bounds depend on the pair only through $s(\mathsf{d})$. Now let
$(\varphi_1,\varphi_2)\in\mathfrak{T}(\mathsf{d}_1;\vartheta)$ and $0<\mathsf{d}\le\mathsf{d}_1$,
and let $(f_1,f_2)$ be obtained by shrinking $(\varphi_1,\varphi_2)$ about a point by the factor
$\mathsf{d}/\mathsf{d}_1$. Then:
\begin{itemize}
\item the supports are contained in balls of radius $\mathsf{d}$;
\item the sup norms are unchanged;
\item the integrals are multiplied by $(\mathsf{d}/\mathsf{d}_1)^4$;
\item the Lipschitz constants are multiplied by $\mathsf{d}_1/\mathsf{d}$;
\item the distance of the supports is multiplied by $\mathsf{d}/\mathsf{d}_1$.
\end{itemize}
Hence the conditions defining $\mathfrak{T}(\mathsf{d}_1;\vartheta)$ for the pair
$(\varphi_1,\varphi_2)$ become exactly those defining $\mathfrak{T}(\mathsf{d};\vartheta)$ for
$(f_1,f_2)$, that is, $(f_1,f_2)\in\mathfrak{T}(\mathsf{d};\vartheta)$. Along such
sup-normalised dilates the quotient is constant in $\mathsf{d}$ for a scale-invariant law of
the action density and for the free Maxwell field at any charge.
By what was just shown, the quotient is strictly positive and tends to zero as $\mathsf{d}\to0$.
It is therefore not constant, so the limit is neither scale invariant nor a free Maxwell field
at any charge. Being not identically zero and decaying more slowly than every power of
$\mathsf{d}$, it is non-trivial in the sense of Definition~\ref{3.1:def-non-trivial}. This
proves (f).

\emph{Asymptotic freedom.} The display in (g) is \eqref{8.1:eq-two-sided-bounds-4} divided by
$\|f_1\|_\infty\|f_2\|_\infty$, with $\frac83\cdot100=\frac{800}3$, and with $\bar{x}_s$,
$\underline{x}_s$ the envelopes of \eqref{8.1:eq-coupling-envelope-f}, which are affine in $s$
with the slopes $\lambda^{\sharp}$ and $B^{\sharp}$ of clause (0). The decay is therefore that of
the fourth power of the running coupling at depth $s(\mathsf{d})$. The depth $s(\mathsf{d})$ is
defined in Notation~\ref{8.1:not-test-pairs-kernel} through the room function $C_W$, and
writing $\bar{x}_{s(\mathsf{d})}$ in terms of $\log_L(1/\mathsf{d})$ adds the term
$B^{\sharp}\log_L C_W(s(\mathsf{d}))$ to its additive constant. By
\eqref{8.1:eq-coupling-envelope-a} and the degree function
$p_0(g')=A_0(\log g'^{-2})^{p_0}$, evaluated at $g'=\bar{x}_s^{-1/2}$,
\begin{equation*}
C_W(s)=C_0\,p_0\bigl(\bar{x}_s^{-1/2}\bigr)^2=C_0A_0^2\,(\log\bar{x}_s)^{2p_0},
\qquad
\log_L C_W(s)=\log_L\bigl(C_0A_0^2\bigr)+2p_0\log_L\log\bar{x}_s .
\end{equation*}
Since $\bar{x}_s$ is affine in $s$, $\log\bar{x}_s=O(\log s)$, and hence
$\log_L C_W(s)=O(\log\log s)$; since $s(\mathsf{d})=O(\log_L(1/\mathsf{d}))$, the term
$B^{\sharp}\log_L C_W(s(\mathsf{d}))$ is $O(\log\log\log(1/\mathsf{d}))$, hence of lower
order than $\log_L(1/\mathsf{d})$. This proves (g).
\end{proof}

\begin{definition}[The block field at depth $s$, thresholds and the full average]
\label{8.2:not-depth-block-field}
Throughout this chapter $K$ is the cutoff, so that the bare spacing is $L^{-K}$, and $s$ is the
depth. We put
\begin{equation*}
  k_0 := K - s, \qquad x_k := g_k^{-2}, \qquad n := N^2 - 1 = \dim \mathfrak{su}(N).
\end{equation*}
The trace is normalised by $\operatorname{tr} 1 = 1$, where $1$ is the identity matrix. On
$\mathfrak{g} := \mathfrak{su}(N)$ we use
the inner product $\langle X, Y\rangle := -\operatorname{tr}(XY)$ and the norm
$|X|^2 := \langle X, X\rangle$.

\emph{The block field at depth $s$.} We write $V := V^{(s)} = V_{k_0}$ for the \emph{depth-$s$ block
field}. It is the block field $V_{k_0}$ of Ba{\l}aban's construction at level $k_0$, that is, the
variable integrated in the renormalisation step at level $k_0$. It lives on the block lattice
$\mathbb{T}^{(k_0)}$ of Definition~\ref{1.5:def-block-average}, which we call the depth-$s$ unit
lattice; in its own lattice units it is a torus of side $2L^{m+s}$. The block field one level
further is $W := V_{k_0+1}$, and later levels are written in the same way. Let $\rho_k$ be the
effective density at level $k$ (Definition~\ref{1.5:def-densities}). We put
\begin{equation}\label{8.2:eq-depth-block-field-1}
  \nu^{B}_{K,s} := \frac{\rho_{k_0}\, dV}{\int \rho_0},
\end{equation}
where $\rho_0$ is the density at level $0$.

\emph{The plaquette observable.} For a plaquette $p$ of $\mathbb{T}^{(k_0)}$ we put
$\mathcal{O}_{s,p} := 1 - \operatorname{Re}\operatorname{tr} V(\partial p)$. If
$V(\partial p) = e^{Y}$ with $Y \in \mathfrak{g}$, then
\begin{equation}\label{8.2:eq-depth-block-field-2}
  \mathcal{O}_{s,p} = \tfrac12 |Y|^2 - \tfrac1{24}\operatorname{tr} Y^4 + O(|Y|^6).
\end{equation}
To see this, expand the exponential series and write $Y = iH$ with $H$ Hermitian. The traces
$\operatorname{tr} Y^{2j+1} = \pm i \operatorname{tr} H^{2j+1}$ are purely imaginary, so the odd
powers do not contribute to the real part. The even terms give
\begin{equation*}
  \operatorname{Re}\operatorname{tr} e^{Y}
  = 1 + \tfrac12 \operatorname{tr} Y^2 + \tfrac1{24}\operatorname{tr} Y^4 + O(|Y|^6).
\end{equation*}
Finally $-\operatorname{tr} Y^2 = |Y|^2$.

\emph{The degree function and the thresholds.} We write $r_B$ for the lifetime exponent of
Ba{\l}aban's construction in \cite{B-CMP119}, where it is denoted $r$; we rename it so that the
letter $r$ remains free for numbers of renormalisation steps. It satisfies $r_B \ge 2$. The degree
function is
\begin{equation*}
  p_0(g) := A_0 \bigl(\log g^{-2}\bigr)^{p_0},
\end{equation*}
as in \cite{B-CMP119}, and its exponent satisfies $p_0 \ge 5 r_B$. The thresholds at level $k$ are
\begin{equation}\label{8.2:eq-depth-block-field-3}
  \varepsilon_k := g_k\, p_0(g_k), \qquad
  \delta_k := \frac{A_1}{A_0}\, \varepsilon_k, \qquad
  \mathfrak{b}_k := \frac{\delta_k}{g_k} = A_1 (\log x_k)^{p_0}.
\end{equation}
The last equality holds because $g_k^{-2} = x_k$ and
\begin{equation*}
  \frac{\delta_k}{g_k} = \frac{A_1}{A_0}\,p_0(g_k) = A_1\bigl(\log g_k^{-2}\bigr)^{p_0}.
\end{equation*}

\emph{Averages and further symbols of the witness construction.} We use the average
$\mathrm{Av}_p$ over plaquettes $p$, taken in the witness construction, and the sublattice
\begin{equation*}
  \mathfrak{P} := S_0 \mathbb{Z}^4 / 2L^{m+s}\mathbb{Z}^4, \qquad
  S_0 := L^2 M \check{R}_{K-s}.
\end{equation*}
Here $\check{R}_{K-s}$ is a number of the witness construction attached to level $K-s$, so that
$S_0$ is the mesh of the sublattice $\mathfrak{P}$. We also use the depth $s_*$, the lower end of
the range of depths in item~(b) below; the term $\delta^B(s) \asymp L^{-4s}$, the additive term of
the bound in item~(b); the level $k_*$ of the witness construction; and the depth
\begin{equation*}
  s_{\mathrm{up}}(g) := \tfrac14 \log_L g^{-2} + C_{\mathrm{up}},
\end{equation*}
in which $C_{\mathrm{up}}$ is a constant of the witness construction that does not depend on $g$.
We write $\langle\!\langle \cdot \rangle\!\rangle_s$ for the \emph{full average} over all depth-$s$
plaquettes of the torus. It is distinct from the average $\mathrm{Av}_p$.

\emph{Inputs.} The following statements are used with this notation, or, in item~(d), referred to.
\begin{itemize}
  \item[(a)] Clause~(0) of Theorem~\ref{3.1:thm-main} gives, for $i \le j$,
  \begin{equation*}
    \lambda_\sharp (j - i) - c_5 \;\le\; x_i - x_j \;\le\; B_\sharp (j - i),
  \end{equation*}
  and $g_k \in (0, g)$.
  \item[(b)] For $s_* \le s \le K-1$ we use the upper bound of clause~(ii-d) of
  Theorem~\ref{3.1:thm-main},
  \begin{equation*}
    \mathrm{Av}_p[\mathcal{O}] \le C(N)\, g_k^2\, p_0(g_k)^2 + 2\delta^B(s).
  \end{equation*}
  \item[(c)] By Lemma~\ref{1.1:lem-invariance}, the expectation
  $\langle\cdot\rangle_{K,m}$ is exactly invariant under translations and under $W_4$.
  \item[(d)] The first sentence of clause~(iii)(a) of Theorem~\ref{3.1:thm-main} is the statement
  towards which the arguments of this chapter are directed; it is identified here by its clause
  letter and is not among the inputs (a)--(c).
\end{itemize}
\end{definition}

\begin{definition}[Ba{\l}aban's small-field step and the recovered field]
\label{8.2:not-small-field-step}
We recall the small-field renormalisation transformation of \cite[\S2]{B-CMP109}, in the
notation used there, and fix the map that expresses the integration variable of one step
through the Gaussian fluctuation field.

\emph{The step.} At step $k$ the small-field transformation acts on the density by
\cite[(2.1)]{B-CMP109}:
\begin{equation}\label{8.2:eq-small-field-step-2}
\begin{split}
W\longmapsto\int dV\,\delta\bigl(\mathcal{M}(V)W^{-1}\bigr)\,\chi_k
\exp\bigl[&-g_k^{-2}\mathbf{G}(V)-g_k^{-2}A\bigl(U_k(V)\bigr)\\
&+E_k\bigl(U_k(V)\bigr)\bigr].
\end{split}
\end{equation}
Here $\mathbf{G}$ is the exponential gauge-fixing function of
\cite[(0.14)--(0.16)]{B-CMP109}, inserted by the Faddeev--Popov procedure; $U_k(V)$ is the
$k$-step minimiser of \cite[(0.21)]{B-CMP109}; and $V^{(k)}=V^{(k)}(W)$ is the one-step
critical configuration of \cite[(2.3)]{B-CMP109}. The variable $W$ is the bond configuration
of level $k+1$ on which the transformed density depends; $\mathcal{M}$ is the block-averaging
map (Definition~\ref{1.5:def-block-average}), so that the delta function imposes the
averaging constraint on $V$; $E_k$ is the correction to the Wilson term in the exponent of
the density at level $k$; and $U_{k+1}$, which enters below, is the configuration of level $k+1$
at which, after the step, the action is evaluated. The transformation is carried out by the
following operations of \cite[\S2]{B-CMP109}.
\begin{enumerate}
\item[(a)] The integration variable is written $V=V'V^{(k)}$, and $V'=:e^{iB'}$.
\item[(b)] The configuration $U_k(V'V^{(k)})$ is re-described by gauge transformations (from
the axial to the Landau gauge, with the transformation $u_k$). These re-descriptions rewrite
the background inside $A$ and $E_k$; they do not change the integration variable.
\item[(c)] The terms are expanded by Taylor's formula, \cite[(2.4)--(2.8)]{B-CMP109}, and the
window $\chi_k$ of \cite[(2.9)]{B-CMP109} is the product over bonds $b$ of the indicator
functions of $|B'(b)|$ lying below the threshold fixed there.
\item[(d)] The analytic change of variables $B'=B-h\tilde D(B)$, with the operator $h$ and the
map $\tilde D$ of \cite[p.~267]{B-CMP109} (this $h$ is Ba{\l}aban's operator and is unrelated
to the histories $h$), linearises the constraint; $\tilde D$ begins with quadratic terms, and
the Jacobian is
$\exp\operatorname{Tr}\log\bigl(I-h\,\delta\tilde D/\delta B\bigr)$
\cite[p.~267]{B-CMP109}.
\item[(e)] The scaling $B=g_kB''$ is performed, after which, in the words of
\cite[pp.~267--268]{B-CMP109}, ``the only term with a negative power of $g_k$ is the action
evaluated at the configuration $U_{k+1}$'', and the terms of order zero are
\cite[(2.11)]{B-CMP109}, ``the
quadratic form in $B'$ \dots\ equal to $-\frac12\langle B',\Delta^{(k)}B'\rangle$, see
(3.156) [13]'', where [13] is \cite[(3.156)]{B-CMP99b}.
\item[(f)] The variables $B''(b_0(c))$ (with $b_0(c)$ the bonds so labelled in
\cite[p.~268]{B-CMP109}; this letter is Ba{\l}aban's and is unrelated to the constant
$b_0=(N^2-1)/2$) are eliminated by the delta function $\delta(\tilde QB'')$, in which
$\tilde Q$ is the operator of \cite[p.~268]{B-CMP109} expressing the linearised averaging
constraint: ``Denoting the remaining variables by $B$ we have $B'=CB$, $C$ is the operator
determined by the configuration $V^{(k)}$, and the measure becomes a Gaussian measure in
variables $B$, with the covariance $C^{(k)}=C^{(k)}(U_{k+1})=(C^*\Delta^{(k)}C)^{-1}$''
\cite[p.~268]{B-CMP109}.
\end{enumerate}
A comparison of the sentences quoted in (e) and (f) with \eqref{8.2:eq-small-field-step-1}
below shows that the variable denoted $B'$ there is the scaled field of (e), so that $B''=CB$.

\emph{The result.} The result of the step is \cite[(2.12)]{B-CMP109}, which gives the new
action $A_{k+1}$ as a function of $U_{k+1}$ (the letter $A_{k+1}$ is Ba{\l}aban's and is
distinct from the Wilson action $A$ and from the parameters $A_0$, $A_1$):
\begin{equation}\label{8.2:eq-small-field-step-3}
\begin{split}
A_{k+1}(U_{k+1})={}&-g_k^{-2}A(U_{k+1})+E_k(U_{k+1})
+\bigl[\log Z^{(k)}(U_{k+1})-\log Z^{(k)}(1)\bigr]\\
&+\log\Bigl[(N''_k)^{-1}\int d\mu_{C^{(k)}}(B)\,\chi_k\exp\Bigl[
\operatorname{Tr}\log\Bigl(I-h\frac{\delta\tilde D}{\delta B}\Bigr)(g_kCB)\\
&\quad+\log\sigma\bigl(g_kCB-h\tilde D(g_kCB)\bigr)
+g_k^{-2}\bigl\langle H_1h\tilde D_3(g_kCB),J\bigr\rangle
-g_k^{-2}\mathbf{G}_3(g_kB)\\
&\quad+g_k^{-2}\bigl\langle H_1g_kCB,\Delta_1H_1h\tilde D(g_kCB)\bigr\rangle\\
&\quad-\tfrac12g_k^{-2}\bigl\langle H_1h\tilde D(g_kCB),
\Delta_1H_1h\tilde D(g_kCB)\bigr\rangle\\
&\quad-g_k^{-2}\mathcal{V}\bigl(H_1(g_kCB-h\tilde D(g_kCB))\bigr)\\
&\quad+\bigl\{E_k\bigl(U_k(\exp i[g_kCB-h\tilde D(g_kCB)]\,V^{(k)})\bigr)
-E_k\bigl(U_k(V^{(k)})\bigr)\bigr\}\Bigr]\Bigr],
\end{split}
\end{equation}
where $Z^{(k)}$, $\sigma$, $H_1$, $\Delta_1$, $\tilde D_3$, $J$, $\mathbf{G}_3$ and the
function written $\mathcal{V}$ here are those of \cite[\S2]{B-CMP109}, and $N''_k$ is the
integral in \eqref{8.2:eq-small-field-step-3} evaluated at $U_{k+1}=1$. We write
$P^{(k)}(g_k,U_{k+1},B)$ for everything under the exponential in
\eqref{8.2:eq-small-field-step-3} except the curly bracket. Then \cite[(2.13)]{B-CMP109}
reads
\begin{equation}\label{8.2:eq-small-field-step-4}
E^{(k+1)}(g_k,U_{k+1})=\log\int d\mu_{C^{(k)}}(B)\,\chi_k
\exp\bigl[P^{(k)}(g_k,U_{k+1},B)+\{\dots\}\bigr],
\end{equation}
with $\{\dots\}$ the curly bracket of \eqref{8.2:eq-small-field-step-3}, and
\cite{B-CMP109} states that ``the expression under the exponential above vanishes at
$g_k=0$''. Under a gauge transformation $u$, acting on the configuration by
$U_{k+1}\mapsto U^u_{k+1}$ and on the fluctuation field by $B'\mapsto R(u)B'$, with $R(u)$ the
map so written in \cite[(2.16)]{B-CMP109} (this $R$ is Ba{\l}aban's and is unrelated to the
radius $R$ and to the operation $\mathbf{R}$), the measure and all the expressions are
invariant \cite[(2.16)]{B-CMP109}. According to \cite{B-CMP109}, $P^{(k)}$ ``is an analytic
function of the
fluctuation field $B$ defined on the unit lattice $T_1^{(k)}$ \dots\ bounded by a second
order polynomial in $B$, with coefficients of the order $O(g_k)$, hence it can be treated as a
small perturbation of the basic quadratic form in the Gaussian measure
$d\mu_{C^{(k)}}$''; its analytic extension $P^{(k)}(g_k,U,J,B)$ is
\cite[(3.1)]{B-CMP109}. The construction of a cluster expansion of the integral in
\eqref{8.2:eq-small-field-step-4} is the proof of \cite[Theorem~3]{B-CMP109}.

\emph{The recovered field.} For the Gaussian variable $B$ we put
\begin{equation}\label{8.2:eq-small-field-step-1}
\begin{gathered}
\mathfrak{v}_k(B;W):=\exp\bigl(i[g_kCB-h\tilde D(g_kCB)]\bigr)\,V^{(k)}(W),\\
Y_k(B):=i\bigl[g_kCB-h\tilde D(g_kCB)\bigr]=i\,g_kCB+O\bigl(g_k^2|CB|^2\bigr).
\end{gathered}
\end{equation}
Thus $\mathfrak{v}_k(B;W)=e^{Y_k(B)}V^{(k)}(W)$, and $\mathfrak{v}_k(B;W)$ is the integration
variable $V$ of \eqref{8.2:eq-small-field-step-2} expressed through $B$.
\end{definition}

\begin{proof}[Proof of the assertions in Definition~\ref{8.2:not-small-field-step}]
We show that $\mathfrak{v}_k(B;W)$ is the integration variable $V$ of
\eqref{8.2:eq-small-field-step-2}, and we prove the expansion of $Y_k(B)$ in
\eqref{8.2:eq-small-field-step-1}.

In \eqref{8.2:eq-small-field-step-2} the term $E_k(U_k(V))$ is evaluated at the integration
variable $V$. By (b), the gauge re-descriptions act on the background inside $A$ and $E_k$
and leave the integration variable unchanged; the operations (d)--(f) are changes of the
integration variable and of the measure, and (c) only reorganises the integrand. Hence the
first term in the curly bracket of
\eqref{8.2:eq-small-field-step-3} is $E_k(U_k(V))$ with $V$ written in the final variable
$B$, and its argument displays this variable as
\[
V=\exp i\bigl[g_kCB-h\tilde D(g_kCB)\bigr]\,V^{(k)}=\mathfrak{v}_k(B;W).
\]
The second term of the curly bracket is the same expression at $B=0$, since
$\tilde D(0)=0$ ($\tilde D$ begins with quadratic terms).

The same expression results from composing the changes of variables. By (a),
$V=e^{iB'}V^{(k)}$. By (d), $B'=B-h\tilde D(B)$ with $B$ the variable introduced there; by
(e) this variable equals $g_kB''$; and by (f) the scaled field is $B''=CB$ in the remaining
variables $B$. Substituting successively,
\[
B'=g_kB''-h\tilde D(g_kB'')=g_kCB-h\tilde D(g_kCB),
\]
so that, by the display above, $V=\mathfrak{v}_k(B;W)$ and $iB'=Y_k(B)$. This
is consistent with the covariance in (f): by (e) the quadratic form of order zero in the
scaled field is $-\frac12\langle B'',\Delta^{(k)}B''\rangle$, and with $B''=CB$ it equals
$-\frac12\langle B,C^*\Delta^{(k)}CB\rangle$, the quadratic form of the Gaussian measure with
covariance $(C^*\Delta^{(k)}C)^{-1}=C^{(k)}$.

It remains to prove the expansion. By (d), $\tilde D$ begins with quadratic terms, so
$h\tilde D(g_kCB)$ is of second order in $g_kCB$, that is,
$h\tilde D(g_kCB)=O(g_k^2|CB|^2)$. Therefore
\[
Y_k(B)=i\,g_kCB-i\,h\tilde D(g_kCB)=i\,g_kCB+O\bigl(g_k^2|CB|^2\bigr),
\]
which is the second line of \eqref{8.2:eq-small-field-step-1}.
\end{proof}

\begin{definition}[The effective action of the full large-field expansion]
\label{8.2:not-full-effective-action}
The renormalisation step whose small-field form is fixed in Definition~\ref{8.2:not-small-field-step}
is used in this chapter also in the form it takes in the complete expansion of \cite{B-CMP119},
where the large-field regions are included. We fix the notation for that form. All objects are
those of \cite{B-CMP119}, denoted by its letters, with one exception. In
\cite[(2.21), (2.23)]{B-CMP119} the letter $A$ is used both for the weighted Wilson action
$A(\cdot,\cdot)$ of Definition~\ref{1.1:def-wilson-action} and for the fluctuation fields $A_j$.
Ba{\l}aban's effective action $A_k$ is therefore written $A^{\mathrm{eff}}_k$ here, so that the
letter $A$ does not receive a third use.

\emph{The density.} At level $k$ the effective density $\rho_k$ of Definition~\ref{1.5:def-densities}
is
\begin{equation}\label{8.2:eq-full-effective-action-1}
\rho_k=\sum\mathbf{T}_k\bigl(\exp A^{\mathrm{eff}}_k\bigr),
\end{equation}
the sum running over the admissible sequences of domains of \cite[(2.18)]{B-CMP119}. Here
$\mathbf{T}_k$ is the operation of \cite[(2.18)]{B-CMP119}, and $A^{\mathrm{eff}}_k$
is the effective action defined in \eqref{8.2:eq-full-effective-action-2} below. The small-field
characteristic functions that enter are those of \cite[(2.17)]{B-CMP119}. We write $\Lambda_j$
for the domain with index $j$ of an admissible sequence of domains in
\eqref{8.2:eq-full-effective-action-1}.

\emph{One step.} A step $j\to j+1$, ``if it is not changed by an R-operation'', has the form
\cite[(2.21)]{B-CMP119}. It integrates the block field $V_j$ on the domain $\Omega_{j+1}$ of
\cite[(2.21)]{B-CMP119} (Ba{\l}aban's letter for a lattice domain, not the vacuum vector
$\Omega$ of Definition~\ref{0.2:not-glossary}), and it integrates the fluctuation fields $A_j$. The
integrand is a small-field characteristic function $\chi$, of the kind defined in
\cite[(2.17)]{B-CMP119}, times an exponential whose exponent is minus one half of a quadratic
form in $A_j$ plus further terms. About this quadratic form, \cite[after (2.21)]{B-CMP119}
states: ``the
quadratic form in the exponential couples only the fields $A_j$ in the same component \dots{} The
other terms of this quadratic form are included into the effective action into $B$-terms''. The
second argument $\mathbf{A}$ of $B_k$ in \eqref{8.2:eq-full-effective-action-2} is the letter of
\cite[(2.23)]{B-CMP119} for the fluctuation fields $A_j$, on which, by the sentence just quoted,
the $B$-terms depend.

\emph{The effective action.} By \cite[(2.23)]{B-CMP119},
\begin{equation}\label{8.2:eq-full-effective-action-2}
A^{\mathrm{eff}}_k=-A\bigl(g_k^{-2}(\cdot),U_k\bigr)+E_k(U_k)+R_k(U_k)+B_k(U_k,\mathbf{A})-E_k .
\end{equation}
The first term is the weighted Wilson action of Definition~\ref{1.1:def-wilson-action}. It is
evaluated at the block field $U_k$ and at the position-dependent inverse coupling
$g_k^{-2}(\cdot)$. The last term is written $E_k$ without an argument, as in
\cite[(2.23)]{B-CMP119}; the letter $E_k$ thus has two uses in the display, distinguished by the
presence of the argument.
Of the three remaining terms, \cite[after (2.23)]{B-CMP119} states the following.
\begin{itemize}
\item $E_k$ is ``fully renormalized''.
\item $R_k$ is described as ``the effect of the R-operations \dots{} vacuum energy
renormalization only''.
\item $B_k$ is ``localized closely to large field regions. It does not require any
renormalization''.
\end{itemize}

\emph{The coupling.} Position-dependent couplings are introduced in
Definition~\ref{1.5:def-domains}. The one used here is that of \cite[(2.24)]{B-CMP119}. It is
built from functions $\varphi_j$ and domains $\Lambda_j^{\sim}$, one for each domain $\Lambda_j$,
that enter \cite[(2.24)]{B-CMP119}, and \cite[(2.24)]{B-CMP119} states:
``$\varphi_j=1$ on $\Lambda_j^{\sim}$. Thus $g_j(x)=g_j$ on
the last domain''.

\emph{The localised representation of $E_k$.} The term $E_k(U_k)$ has the localised
representation of \cite[(2.25)--(2.27)]{B-CMP119}, whose terms are indexed by a level $j$ and a
domain $X$. Each term has the following four properties.
\begin{itemize}
\item[(i)] It depends on $U_k$ only through the restriction $U_k\restriction X$.
\item[(ii)] It extends analytically to the set $\mathbf{U}_j(X,\alpha_{0j},\alpha_{1j})$ of
\cite[(2.25)--(2.27)]{B-CMP119}, which is determined by the domain $X$ and by the radii
$\alpha_{0j}$, $\alpha_{1j}$ given below.
\item[(iii)] It is gauge invariant.
\item[(iv)] It satisfies the bound \cite[(1.1.18)]{B-CMP119}.
\end{itemize}
The radii in (ii) are those of \cite[(2.28)]{B-CMP119}:
\begin{equation}\label{8.2:eq-full-effective-action-3}
\alpha_{0j}=g_jC_0\bigl(\log g_j^{-2}\bigr)^{q_0},\qquad
\alpha_{1j}=g_jC_1\bigl(\log g_j^{-2}\bigr)^{q_1}.
\end{equation}

\emph{The $\mathbf{R}$-operation terms.} By \cite[(2.30)]{B-CMP119},
\begin{equation}\label{8.2:eq-full-effective-action-4}
R_k(U_k)=\sum_j\Bigl[R^{(j)}(\Lambda_j,U_k)-R^{(j)}(\Lambda_j,1)\Bigr],
\end{equation}
where $1$ denotes the configuration equal to the identity on every bond. The function
$R^{(j)}(\Lambda_j,\cdot)$ is a sum of terms $R^{(j)}(X,\cdot)$ over the domains $X$ of the
family $\mathcal{D}_j$ of \cite[(2.30)]{B-CMP119}. These terms have properties (i)--(iii) above.
In addition they satisfy the following bound, stronger than (iv), for the arguments $(U,J)$ of
\cite[(2.31)]{B-CMP119}:
\begin{equation}\label{8.2:eq-full-effective-action-5}
\bigl|R^{(j)}\bigl(X,(U,J)\bigr)\bigr|\le g_j^{\kappa_0}\exp\bigl(-\kappa\,d_j(X)\bigr).
\end{equation}
Here $d_j(X)$ and $\kappa$ are as in Definition~\ref{0.2:not-glossary}, and $\kappa_0$ is the
exponent of \cite[(2.31)]{B-CMP119}, of which \cite[after (2.31)]{B-CMP119} states:
``$\kappa_0$ can be chosen arbitrarily large''.

\emph{Marginal terms.} Let $\eta$ denote the lattice spacing of the lattice on which the
expansion of \cite{B-CMP119} starts. About the terms $R^{(j)}(X,\cdot)$ entering
\eqref{8.2:eq-full-effective-action-4}, bounded in \eqref{8.2:eq-full-effective-action-5},
\cite[after (2.31)]{B-CMP119} states:
\begin{quote}
``After the vacuum energy renormalization we obtain a sum of marginal terms, i.e., terms with
bounds $O(1)(L^j\eta)^4g_j^{\kappa_0}\exp(-\kappa d_j(X))$. The sum over $X$ is controlled by the
exponential factor, and by the factor $(L^j\eta)^4$. The sum over $j$ is controlled by
$g_j^{\kappa_0}$.''
\end{quote}
\end{definition}

\begin{definition}[Local gauge-invariant holomorphic observables]\label{8.2:def-local-observables}
Let $0\le k\le K$ and let $\mathbb{T}^{(k)}$ be the block lattice at level $k$ of
Definition~\ref{1.5:def-block-average}. We write $\mathfrak{su}(N)$ for the Lie algebra of
$\mathrm{SU}(N)$, $\mathfrak{su}(N)_{\mathbb{C}}=\mathfrak{su}(N)+i\,\mathfrak{su}(N)$ for its
complexification, $|\cdot|$ for the operator norm of complex $N\times N$ matrices, and
$\mathbf{1}$ for the unit matrix.
\begin{itemize}
\item[(1)] For a plaquette or a bond $p$ of $\mathbb{T}^{(k)}$ and an integer $\ell\ge1$,
$\mathfrak{B}_\ell(p)$ is the set of bonds of $\mathbb{T}^{(k)}$ within graph distance $\ell$
of $p$.
\item[(2)] For $\varrho>0$, the complex $\varrho$-tube about $\mathrm{SU}(N)^{\mathfrak{B}_\ell(p)}$ is
\begin{equation*}
\mathbb{U}_\varrho=\bigl\{V:\ V(b)=e^{Z_b}V_0(b),\ V_0(b)\in\mathrm{SU}(N),\
Z_b\in\mathfrak{su}(N)_{\mathbb{C}},\ |Z_b|<\varrho,\ b\in\mathfrak{B}_\ell(p)\bigr\}.
\end{equation*}
\item[(3)] For $\mathsf{M}>0$, the class $\mathfrak{G}_k(p;\ell,\mathsf{M},\varrho)$ consists of
the real functions $G$ of the restriction of the configuration $V$ to $\mathfrak{B}_\ell(p)$ that
are (a) invariant under all $\mathrm{SU}(N)$-valued lattice gauge transformations, and (b)
holomorphic on $\mathbb{U}_\varrho$, with $\sup_{\mathbb{U}_\varrho}|G|\le\mathsf{M}$.
\item[(4)] For $U\in\mathrm{SU}(N)^{\mathfrak{B}_\ell(p)}$ and $j\ge0$, $D^jG(U)$ is the $j$-th
differential at $Y=0$ of $Y\mapsto G(Ue^{Y})$ on $\mathfrak{su}(N)^{\mathfrak{B}_\ell(p)}$, where
$(Ue^Y)(b)=U(b)e^{Y_b}$ (left-invariant derivatives); thus
\begin{equation*}
D^jG(U)[Y,\dots,Y]=\frac{d^j}{dt^j}G(Ue^{tY})\Big|_{t=0}.
\end{equation*}
We put
\begin{equation*}
\|D^jG\|_\infty=\sup\Bigl\{\,\bigl|D^jG(U)[Y,\dots,Y]\bigr|:\
U\in\mathrm{SU}(N)^{\mathfrak{B}_\ell(p)},\ \max_{b}|Y_b|\le1\Bigr\}.
\end{equation*}
\item[(5)] For $U\in\mathrm{SU}(N)^{\mathfrak{B}_\ell(p)}$ the \emph{local curvature} is
$\operatorname{curv}_{\ell,p}(U)=\max_{p'\subset\mathfrak{B}_\ell(p)}|U(\partial p')-\mathbf{1}|$,
the maximum over the plaquettes $p'$ whose bonds lie in $\mathfrak{B}_\ell(p)$. We write
$\operatorname{Hess}_UG=D^2G(U)$, a quadratic form on
$\mathfrak{su}(N)^{\mathfrak{B}_\ell(p)}$, and, for a centred Gaussian law $\Gamma$ on
$\mathfrak{su}(N)^{\mathfrak{B}_\ell(p)}$,
\begin{equation*}
\operatorname{Tr}_\Gamma[\operatorname{Hess}_UG]=\mathbb{E}_{Z\sim\Gamma}\,D^2G(U)[Z,Z].
\end{equation*}
\end{itemize}
\end{definition}

Every $G\in\mathfrak{G}_k(p;\ell,\mathsf{M},\varrho)$ satisfies, for all $j\ge0$,
\begin{equation}\label{8.2:eq-local-observables-1}
\|D^jG\|_\infty\le\mu_j,\qquad \mu_j=j!\,\mathsf{M}\,\varrho^{-j}.
\end{equation}
The observable $O_{s,p}$ of Notation~\ref{8.2:not-depth-block-field}, for a plaquette $p$ of
$\mathbb{T}^{(k_0)}$, where $k_0$ is the level at which $O_{s,p}$ is defined there, satisfies,
for every $\varrho>0$,
\begin{equation}\label{8.2:eq-local-observables-2}
O_{s,p}\in\mathfrak{G}_{k_0}\bigl(p;1,\mathsf{M}_O(\varrho),\varrho\bigr),\qquad
\mathsf{M}_O(\varrho)=1+e^{4\varrho},
\end{equation}
and $O_{s,p}^2$ belongs likewise to $\mathfrak{G}_{k_0}(p;1,\mathsf{M}_O(\varrho)^2,\varrho)$.

\begin{proof}
We prove \eqref{8.2:eq-local-observables-1} by the Cauchy estimate. Fix
$U\in\mathrm{SU}(N)^{\mathfrak{B}_\ell(p)}$ and $Y$ with $\max_b|Y_b|\le1$. For
$z\in\mathbb{C}$ with $|z|<\varrho$ and each $b$,
\begin{equation*}
U(b)e^{zY_b}=e^{zU(b)Y_bU(b)^{-1}}\,U(b),
\end{equation*}
where $zU(b)Y_bU(b)^{-1}\in\mathfrak{su}(N)_{\mathbb{C}}$ and, conjugation by the unitary matrix
$U(b)$ preserving the operator norm, $|zU(b)Y_bU(b)^{-1}|=|z|\,|Y_b|<\varrho$. Hence
$Ue^{zY}\in\mathbb{U}_\varrho$, with $V_0=U$. The function $f(z)=G(Ue^{zY})$ is holomorphic on
$\{|z|<\varrho\}$, as the composition of the entire maps $z\mapsto e^{zY_b}$ with $G$, and
$|f|\le\mathsf{M}$ there. On the real axis $f(t)=G(Ue^{tY})$, so
$f^{(j)}(0)=D^jG(U)[Y,\dots,Y]$. For $0<r<\varrho$ the Cauchy formula on the circle $|z|=r$
gives $|f^{(j)}(0)|\le j!\,\mathsf{M}\,r^{-j}$; letting $r\uparrow\varrho$ and taking the
supremum over $U$ and $Y$ gives \eqref{8.2:eq-local-observables-1}.

For \eqref{8.2:eq-local-observables-2}: $O_{s,p}$, a function of the bond variables in
$\mathfrak{B}_1(p)$, is holomorphic on $\mathbb{U}_\varrho$ for every $\varrho>0$ (it is entire
in the coordinates $Z_b$), so condition (b) holds for every $\varrho>0$, and its modulus
on $\mathbb{U}_\varrho$ is at most $\mathsf{M}_O(\varrho)=1+e^{4\varrho}$. For $O_{s,p}^2$: the
square of a real, gauge-invariant function is real and gauge invariant; the square of a
holomorphic function on $\mathbb{U}_\varrho$ is holomorphic there; and
\begin{equation*}
\sup_{\mathbb{U}_\varrho}|O_{s,p}^2|\le\Bigl(\sup_{\mathbb{U}_\varrho}|O_{s,p}|\Bigr)^2
\le\mathsf{M}_O(\varrho)^2.
\end{equation*}
\end{proof}

\begin{lemma}[Gauge-invariant observables near a flat background]\label{8.2:lem-flat-background}
Let $G\in\mathfrak{G}_k(p;\ell,M,\varrho)$ be a local gauge-invariant holomorphic observable in the sense of
Definition~\ref{8.2:def-local-observables}, and let $\mu_j=j!\,M\varrho^{-j}$ be the bounds on its
left-invariant derivatives $D^jG$ on $\mathrm{SU}(N)^{\mathfrak{B}_\ell(p)}$ stated there. Let $U_0$ be
\emph{flat} on $\mathfrak{B}_\ell(p)$, that is, pure gauge: gauge equivalent to the configuration
$\mathbf{1}$. For $U\in\mathrm{SU}(N)^{\mathfrak{B}_\ell(p)}$ let $\eta(U)$ be the plaquette-size functional of $U$ on
$\mathfrak{B}_\ell(p)$ in the lemma on the complete axial gauge on a box. Put
\begin{equation}\label{8.2:eq-flat-background-1}
\ell' := 3(2\ell+2),\qquad
\eta_{\mathrm{ax}} := \frac{\sin(\pi/N)}{3\ell'},\qquad
C_{\mathrm{ax}} := \frac{\pi}{3}\cdot 3\ell' .
\end{equation}
Then:
\begin{enumerate}
\item[(a)] $DG(U_0)=0$.
\item[(b)] For every $U\in\mathrm{SU}(N)^{\mathfrak{B}_\ell(p)}$ with $\eta(U)\le\eta_{\mathrm{ax}}$,
\begin{equation}\label{8.2:eq-flat-background-2}
\|DG(U)\|\le C_{\mathrm{ax}}\,\mu_2\,\eta(U),\qquad
|G(U)-G(U_0)|\le \tfrac12\,C_{\mathrm{ax}}^2\,\mu_2\,\eta(U)^2 .
\end{equation}
\item[(c)] $D^2G(U_0)$, read in a gauge in which $U_0=\mathbf{1}$, is invariant under the diagonal adjoint
action $Z\mapsto(\operatorname{Ad}_uZ_b)_b$ on $Z=(Z_b)_b\in\mathfrak{su}(N)^{\mathfrak{B}_\ell(p)}$,
$u\in\mathrm{SU}(N)$. For the single-plaquette action density at $p$, written
$O_{s,p}(U)=1-\operatorname{Re}\operatorname{tr}U(\partial p)$,
\begin{equation}\label{8.2:eq-flat-background-3}
D^2O_{s,p}(\mathbf{1})[Z,Z]=\Bigl|\sum_{b\in\partial p}\pm Z_b\Bigr|^2=|\operatorname{curl}Z(p)|^2 ,
\qquad \operatorname{curl}Z(p):=\sum_{b\in\partial p}\pm Z_b ,
\end{equation}
the signs being the orientations of the bonds $b$ in $\partial p$.
\end{enumerate}
\end{lemma}

\begin{proof}
Throughout, $\mathfrak{B}=\mathfrak{B}_\ell(p)$, $|\xi|=(-\operatorname{tr}\xi^2)^{1/2}$ for
$\xi\in\mathfrak{su}(N)$, the norm in which the tube of Definition~\ref{8.2:def-local-observables}, and
hence the bounds $\mu_j$, are measured, and $\|Z\|_\infty=\max_{b\in\mathfrak{B}}|Z_b|$ for
$Z\in\mathfrak{su}(N)^{\mathfrak{B}}$. A gauge transformation $\lambda$ carries $Ue^{t'Z}$ to the transformed
configuration times $e^{t'Z'}$, where $Z'_b=\operatorname{Ad}_{\lambda(y_b)}Z_b$, $\lambda(y_b)$ being the value
of $\lambda$ at the final endpoint $y_b$ of $b$.
Since $G$ is gauge invariant and $\operatorname{Ad}$ preserves the norm, $DG(U)=0$ if and only if the
derivative vanishes at the transformed configuration, and $\|DG\|$ and $|G|$ take the same values at the two
configurations.

(a) By gauge invariance of $G$ we may transform $U_0$ to $\mathbf{1}$. A constant gauge transformation
$\lambda\equiv u\in\mathrm{SU}(N)$ fixes $\mathbf{1}$ and acts on every bond variable by conjugation,
$V(b)\mapsto uV(b)u^{-1}$. Since $ue^{t'Z_b}u^{-1}=e^{t'\operatorname{Ad}_uZ_b}$, its differential at
$\mathbf{1}$ is $\operatorname{Ad}_u$ on each summand of $\mathfrak{su}(N)^{\mathfrak{B}}$. Invariance of $G$
therefore gives
\[
DG(\mathbf{1})[\operatorname{Ad}_uZ]=DG(\mathbf{1})[Z]\qquad(u\in\mathrm{SU}(N)).
\]
Here $\operatorname{Ad}_uZ$ means $(\operatorname{Ad}_uZ_b)_b$. For each bond $b$, let $\delta_b$ denote the
indicator of $b$, so that $\xi\delta_b$ is the element of $\mathfrak{su}(N)^{\mathfrak{B}}$ equal to $\xi$ at $b$
and to $0$ at every other bond. The linear functional $\varphi_b(\xi)=DG(\mathbf{1})[\xi\delta_b]$ on
$\mathfrak{su}(N)$ is then $\operatorname{Ad}$-invariant. Differentiating
$\varphi_b(\operatorname{Ad}_{e^{t\xi'}}\xi)=\varphi_b(\xi)$ at $t=0$ gives $\varphi_b([\xi',\xi])=0$ for all
$\xi,\xi'\in\mathfrak{su}(N)$. Since $[\mathfrak{su}(N),\mathfrak{su}(N)]=\mathfrak{su}(N)$, it follows that
$\varphi_b=0$ for every $b$, and hence $DG(\mathbf{1})=0$.

(b) The set $\mathfrak{B}$ lies in a box of side $\ell'$. By the lemma on the complete axial gauge on a box,
the complete axial gauge on that box writes every bond variable as an ordered product of at most
$(d-1)\ell'=3\ell'$ plaquette variables. Consequently $U$ is gauge equivalent to a configuration $e^{A}$,
$A\in\mathfrak{su}(N)^{\mathfrak{B}}$, with
\[
\|A\|_\infty\le C_{\mathrm{ax}}\,\eta(U);
\]
this chart is valid for $3\ell'\eta(U)<\sin(\pi/N)$. By gauge invariance, $G(U)=G(e^A)$ and
$\|DG(U)\|=\|DG(e^A)\|$. Moreover $G(U_0)=G(\mathbf{1})$, because $U_0$ is gauge equivalent to
$\mathbf{1}$.

On each bond $e^{(t+t')A_b}=e^{tA_b}e^{t'A_b}$. Hence $t\mapsto G(e^{tA})$ is $C^2$, with derivative
$DG(e^{tA})[A]$ and second derivative $D^2G(e^{tA})[A,A]$. The latter is bounded by
$\mu_2\|A\|_\infty^2$. In the same way, for fixed $Z$ the derivative in $t$ of $DG(e^{tA})[Z]$ is
$D^2G(e^{tA})[A,Z]$, which is bounded by $\mu_2\|A\|_\infty\|Z\|_\infty$.

By (a), applied to the flat configuration $\mathbf{1}$, we have $DG(\mathbf{1})=0$. Integrating from $t=0$
to $t=1$ gives
\[
\|DG(e^A)\|\le\mu_2\|A\|_\infty .
\]
Integrating twice from $t=0$ gives
\[
|G(e^A)-G(\mathbf{1})|\le\tfrac12\mu_2\|A\|_\infty^2 .
\]
Inserting $\|A\|_\infty\le C_{\mathrm{ax}}\eta(U)$ yields \eqref{8.2:eq-flat-background-2}.

(c) Take $U_0=\mathbf{1}$ and let $A,Z\in\mathfrak{su}(N)^{\mathfrak{B}}$ be arbitrary. For a constant gauge
transformation $u$, we have
$ue^{tA}e^{t'Z}u^{-1}=e^{t\operatorname{Ad}_uA}e^{t'\operatorname{Ad}_uZ}$ bondwise. Invariance of $G$ gives
$G(e^{t\operatorname{Ad}_uA}e^{t'\operatorname{Ad}_uZ})=G(e^{tA}e^{t'Z})$. Comparing Taylor polynomials at
$t=t'=0$, this is
\[
D^2G(\mathbf{1})[\operatorname{Ad}_uA,\operatorname{Ad}_uZ]=D^2G(\mathbf{1})[A,Z]\qquad(u\in\mathrm{SU}(N)).
\]

For $O_{s,p}$, the plaquette variable of $e^{tZ}$ is the ordered product of the factors $e^{\pm tZ_b}$,
$b\in\partial p$. By the Baker--Campbell--Hausdorff formula it equals $e^{W(t)}$ with
\[
W(t)=t\operatorname{curl}Z(p)+t^2Y_p+O(t^3)\qquad(t\to0),
\]
where $Y_p$ is a sum of commutators of the $Z_b$.
Expanding the exponential,
\[
\operatorname{Re}\operatorname{tr}e^{W}=1+\operatorname{Re}\operatorname{tr}W
+\tfrac12\operatorname{Re}\operatorname{tr}W^2+O(|W|^3).
\]
Since $\operatorname{tr}W=0$ on $\mathfrak{su}(N)$, the first-order term
vanishes. Furthermore,
\[
\operatorname{tr}W^2=t^2\operatorname{tr}(\operatorname{curl}Z(p))^2
+2t^3\operatorname{tr}\bigl(\operatorname{curl}Z(p)\,Y_p\bigr)+O(t^4).
\]
Thus the commutator terms enter $O_{s,p}$ only at cubic order in $Z$. It follows that
\[
O_{s,p}(e^{tZ})=-\tfrac12 t^2\operatorname{tr}(\operatorname{curl}Z(p))^2+O(t^3),
\]
and therefore
$D^2O_{s,p}(\mathbf{1})[Z,Z]=-\operatorname{tr}(\operatorname{curl}Z(p))^2=|\operatorname{curl}Z(p)|^2$;
this is \eqref{8.2:eq-flat-background-3}.
\end{proof}

\begin{lemma}[The small-field step with an inserted observable]\label{8.2:lem-conditional-step}
Let $F$ be a bounded gauge-invariant Borel function of the configurations
$V\in\mathrm{SU}(N)^{\mathrm{bonds}(\mathbb{T}^{(k)})}$ on the level-$k$ lattice
$\mathbb{T}^{(k)}$ of Definition~\ref{1.5:def-block-average}. Let $W$ be a configuration on
$\mathbb{T}^{(k+1)}$, regular in the sense of \cite{B-CMP109}. Apply the six operations
\textup{(op1)}--\textup{(op6)} of Ba{\l}aban's small-field step, listed in
Notation~\ref{8.2:not-small-field-step}, to the starting integral of Ba{\l}aban's $k$-th step
\cite[\S2]{B-CMP109}, which is the left side of~\eqref{8.2:eq-conditional-step-1} at
$F\equiv1$, with its integrand multiplied by $F(V)$. Then
\begin{equation}\label{8.2:eq-conditional-step-1}
\begin{split}
&\int dV\,\delta\bigl(\mathcal{M}(V)W^{-1}\bigr)\,\chi_k\,
 e^{-g_k^{-2}\mathbf{G}(V)-g_k^{-2}A(U_k(V))+E_k(U_k(V))}\,F(V)\\
&\qquad=\mathfrak{N}_k(W)\int d\mu_{C^{(k)}(W)}(B)\,\chi_k(B)\,
 e^{P^{(k)}(g_k,U_{k+1}(W),B)+\Phi_k(B;W)}\,F\bigl(\mathfrak{v}_k(B;W)\bigr).
\end{split}
\end{equation}
Here $\mathcal{M}$ is the block-averaging map of Definition~\ref{1.5:def-block-average}, from
level $k$ to level $k+1$, $\mathbf{G}$ is the gauge-fixing functional, and $E_k$ is the
effective-action term of \cite{B-CMP109}. The function $P^{(k)}(g_k,U_{k+1}(W),B)$ is the term
of the exponent of \cite[(2.12)]{B-CMP109} so denoted there. The function $\Phi_k(B;W)$ denotes
the terms of the
exponent of \cite[(2.12)]{B-CMP109} other than $P^{(k)}$ and the $B$-independent ones collected
in $\mathfrak{N}_k(W)$ below. The map $\mathfrak{v}_k$ is the recovered field of
Notation~\ref{8.2:not-small-field-step}. The operator
$C^{(k)}(W)=(C^{*}\Delta^{(k)}C)^{-1}$ is Ba{\l}aban's fluctuation covariance at the background
determined by $W$; here $C$ is the linear map that expresses the variables eliminated by the
$\delta$-function through the remaining ones, $B''=CB$, as in the proof below and in
Notation~\ref{8.2:not-small-field-step}, and $\Delta^{(k)}$ is the operator so denoted in
\cite[\S2]{B-CMP109}.

The prefactor $\mathfrak{N}_k(W)$ does not depend on $B$, and it is the same as at $F\equiv1$:
\begin{equation*}
\mathfrak{N}_k(W)=\exp\bigl[-g_k^{-2}A(U_{k+1})+E_k(U_{k+1})\bigr]\,
Z^{(k)}(U_{k+1})\,\mathfrak{c}_k ,
\end{equation*}
where $U_{k+1}=U_{k+1}(W)$, $Z^{(k)}(U_{k+1})$ is the factor so denoted in
\cite[(2.12)]{B-CMP109}, and
$\mathfrak{c}_k$ stands for the constant factors of Ba{\l}aban's step (they involve the constants
$N'_k$ and $\sigma_0$ of \cite{B-CMP109}).

Let $\mathbb{E}_k[F\mid W]$ denote the one-step conditional expectation of $F$
given the next field $W$, that is, the quotient of the left side
of~\eqref{8.2:eq-conditional-step-1} by its value at $F\equiv1$ (the blackboard letter
distinguishes the expectation from the effective-action term $E_k$). Then, in the pure
small-field format, it is the ratio
\begin{equation}\label{8.2:eq-conditional-step-2}
\mathbb{E}_k[F\mid W]
=\frac{\displaystyle\int d\mu_{C^{(k)}(W)}(B)\,\chi_k\,e^{\mathcal{P}_k(B;W)}\,
F\bigl(\mathfrak{v}_k(B;W)\bigr)}
{\displaystyle\int d\mu_{C^{(k)}(W)}(B)\,\chi_k\,e^{\mathcal{P}_k(B;W)}},
\end{equation}
where
\begin{equation*}
\mathcal{P}_k(B;W):=P^{(k)}\bigl(g_k,U_{k+1}(W),B\bigr)+\Phi_k(B;W).
\end{equation*}
\end{lemma}

\begin{proof}
Each of Ba{\l}aban's operations, namely the gauge-fixing insertion
\cite[(0.15)--(0.16)]{B-CMP109} and \textup{(op1)}--\textup{(op6)}, commutes with
multiplication by $F$, once $F$ is expressed in the new variables. We go through them in order.

\emph{The gauge-fixing insertion.} The passage from \cite[(0.15)]{B-CMP109} to
\cite[(0.16)]{B-CMP109} is valid for every gauge-invariant integrand. By gauge invariance with
respect to the transformations $u$ used there, the integrand does not depend on $u$, and the
integral over $u$ equals $1$.

Now consider the product of $F$ with the integrand of the starting integral. It is gauge
invariant apart from the gauge-fixing factor $e^{-g_k^{-2}\mathbf{G}(V)}$, because each of its
remaining factors is gauge invariant, namely:
\begin{itemize}
\item the $\delta$-function of the covariant block averages;
\item the factor $\chi_k$;
\item the part of the exponent that does not contain $\mathbf{G}$;
\item and $F$ itself, by hypothesis.
\end{itemize}
Hence the insertion applies to the integral of $F$ times these gauge-invariant factors, and
carries it to the left side of~\eqref{8.2:eq-conditional-step-1}, in which the gauge-fixing
factor appears.

\emph{\textup{(op1)}.} This operation is a translation of Haar measure, $V=V'V^{(k)}$, where
$V^{(k)}=V^{(k)}(W)$ is the configuration on $\mathbb{T}^{(k)}$ determined by $W$
(Notation~\ref{8.2:not-small-field-step}). By invariance of Haar measure the inserted factor
becomes
\begin{equation*}
F(V)\longmapsto F(V'V^{(k)}).
\end{equation*}

\emph{\textup{(op2)}.} This operation rewrites $U_k(V'V^{(k)})$ inside $A$ and inside $E_k$ by
gauge transformations of the fine lattice that equal $1$ on $\mathbb{T}^{(k)}$. It does not act
on the variable $V$. Hence it does not act on $F$.

\emph{\textup{(op3)}.} This operation is bookkeeping of the exponent together with the
definition of $\chi_k$. It leaves the factor $F$ unchanged.

\emph{\textup{(op4)}.} This operation is an analytic change of the integration variable,
\begin{equation*}
V'=e^{i(B-h\widetilde{D}(B))},
\end{equation*}
with $h$ and $\widetilde{D}$ as in Notation~\ref{8.2:not-small-field-step}.
Its Jacobian is written into the exponent: it is the $\operatorname{Tr}\log$ term of
\cite[(2.12)]{B-CMP109}. The inserted factor becomes
\begin{equation*}
F(V'V^{(k)})\longmapsto F\bigl(e^{i(B-h\widetilde{D}(B))}V^{(k)}\bigr).
\end{equation*}

\emph{\textup{(op5)}.} This operation is the linear scaling $B=g_kB''$. The inserted factor
becomes
\begin{equation*}
F\bigl(e^{i(B-h\widetilde{D}(B))}V^{(k)}\bigr)\longmapsto
F\bigl(e^{i(g_kB''-h\widetilde{D}(g_kB''))}V^{(k)}\bigr).
\end{equation*}

\emph{\textup{(op6)}.} This operation integrates the $\delta$-function over the variables
$B''(b^{\circ}(c))$, where $b^{\circ}(c)$ denotes the bonds written $b_0(c)$ in Ba{\l}aban's
construction \cite[\S2]{B-CMP109}, $c$ ranging over the index set of that construction. The
integration expresses these
variables through the remaining ones, denoted again by $B$, so that $B''=CB$. Inserting
$B''=CB$ into the factor obtained
after \textup{(op5)} gives
\begin{equation*}
F\bigl(e^{i(g_kCB-h\widetilde{D}(g_kCB))}V^{(k)}\bigr)=F\bigl(\mathfrak{v}_k(B;W)\bigr),
\end{equation*}
by the definition of $\mathfrak{v}_k$ in Notation~\ref{8.2:not-small-field-step}.

No operation uses any property of the integrand beyond gauge invariance and integrability.
Gauge invariance of $F$ is a hypothesis. Integrability holds because $F$ is bounded and Borel.
Moreover, none of the prefactors that Ba{\l}aban's step pulls out of the integral involves the
integration variable. These prefactors are
\begin{equation*}
e^{-g_k^{-2}A(U_{k+1})},\qquad e^{E_k(U_{k+1})},\qquad Z^{(k)}(U_{k+1}),\qquad N'_k,\qquad
\sigma_0 .
\end{equation*}
The operations therefore carry the integral with $F$ inserted to the right side
of~\eqref{8.2:eq-conditional-step-1}. The prefactor $\mathfrak{N}_k(W)$ is the one obtained at
$F\equiv1$, and the integrand acquires the single additional factor
$F(\mathfrak{v}_k(B;W))$. This proves~\eqref{8.2:eq-conditional-step-1}.

At $F\equiv1$, identity~\eqref{8.2:eq-conditional-step-1} is the identity
\cite[(2.12)]{B-CMP109} before the logarithm is taken and before normalisation by the value at
$U_{k+1}=1$.

Finally, divide~\eqref{8.2:eq-conditional-step-1} by its case $F\equiv1$. The common prefactor
$\mathfrak{N}_k(W)$ cancels. Writing the exponent as $\mathcal{P}_k(B;W)$, as defined in the
statement, in both integrals gives~\eqref{8.2:eq-conditional-step-2}.
\end{proof}

\begin{remark}[Gauge covariance of the one-step conditional expectation]\label{8.2:rem-step-covariance}
Fix a level $k$ and a regular block configuration $W$. Write $\mathcal{S}_k(B;W)$ for the full
exponent of the Gaussian integral on the right of the identity of
Lemma~\ref{8.2:lem-conditional-step}. For a bounded Borel function $F$ of the fine
configuration define the one-step conditional expectation of $F$ given $W$ by
\begin{equation}\label{8.2:eq-step-covariance-1}
  \mathbb{E}_k[F\,|\,W]
  =\frac{\int d\mu_{C^{(k)}(W)}(B)\,\chi_k(B)\,e^{\mathcal{S}_k(B;W)}\,
         F(\mathfrak{v}_k(B;W))}
        {\int d\mu_{C^{(k)}(W)}(B)\,\chi_k(B)\,e^{\mathcal{S}_k(B;W)}} .
\end{equation}
For gauge-invariant $F$, Lemma~\ref{8.2:lem-conditional-step}, whose prefactor
$\mathfrak{N}_k(W)$ is the same with and without $F$, identifies
\eqref{8.2:eq-step-covariance-1} with the ratio of the left side of the identity of
Lemma~\ref{8.2:lem-conditional-step} with $F$ inserted to the same with $F=1$. For general $F$, and
for the Lie-algebra-valued functions $\Phi$ of item (a), we take
\eqref{8.2:eq-step-covariance-1} as the definition.
\begin{enumerate}
\item[(a)] Let $u$ be a gauge transformation of the block lattice, $W^{u}$ the transformed block
  configuration, assumed regular, so that the left sides below are defined, and $\tilde u$ the
  gauge transformation of the fine lattice by which the gauge-fixed configuration over $W^{u}$
  differs from that over $W$. If $F$ is gauge invariant, then
  \begin{equation*}
    \mathbb{E}_k[F\,|\,W^{u}]=\mathbb{E}_k[F\,|\,W].
  \end{equation*}
  Let $\mathfrak{g}$ be the Lie algebra of $\mathrm{SU}(N)$, and let $\Phi$ be a bounded Borel
  function of the fine configuration whose values are $\mathfrak{g}$-valued fields on the sites
  of the fine lattice, covariant under gauge transformations:
  \begin{equation*}
    \Phi(V^{\lambda})(x)=\operatorname{Ad}(\lambda(x))\,\Phi(V)(x)
  \end{equation*}
  for every gauge transformation $\lambda$ of the fine lattice and every site $x$, where
  $\operatorname{Ad}(\lambda(x))$ denotes the adjoint action of $\lambda(x)$ on $\mathfrak{g}$.
  Then, for every site $x$,
  \begin{equation*}
    \mathbb{E}_k[\Phi\,|\,W^{u}](x)=\operatorname{Ad}(\tilde u(x))\,\mathbb{E}_k[\Phi\,|\,W](x),
  \end{equation*}
  that is, $W\mapsto \mathbb{E}_k[\Phi\,|\,W]$ is equivariant.
\item[(b)] In the complete expansion of \cite{B-CMP119}, the step-$k$ operation on a term of
  that expansion has the form \cite[(2.21)]{B-CMP119} wherever that term is not
  changed by an $\mathbf{R}$-operation. On the clean event at a plaquette $p$ (defined further
  on in this chapter), this is the case within the box about $p$ attached to that event. There:
  \begin{itemize}
  \item the thresholds of $\chi_k$ are those of \cite[(2.17)]{B-CMP119}, namely
    $\varepsilon_k,\delta_k$, together with the window of $\chi_k$ on the bond values $B_b$ of
    the Gaussian variable, $|B_b|\lesssim\mathfrak{b}_k$ for every bond $b$, where
    $\mathfrak{b}_k$ is the window width;
  \item the exponent, written $\mathcal{P}_k$ in the notation of \cite{B-CMP119}, is augmented
    by the increments of the quantities written $R_k$ and $B_k$ in \cite[(2.23)]{B-CMP119}
    (that paper's letters);
  \item the exponent is further augmented by the coupling-lag term of
    \cite[(2.24)]{B-CMP119}.
  \end{itemize}
  These three augmentations are displayed further on in this chapter.
  The identity of Lemma~\ref{8.2:lem-conditional-step} holds unchanged, obtained by the same
  operations. In particular, for gauge-invariant $F$ and for the step carrying $\mathbf{T}$,
  the numerator at the first step that touches the support of $F$ is Ba{\l}aban's gauge-fixed
  representation with $F$ inserted. This is the form in which it enters clause (ii-a) of
  Theorem~\ref{3.1:thm-main}.
\item[(c)] The Gaussian variable $B$ has covariance
  \begin{equation*}
    C^{(k)}(W)=\bigl(C^{*}\Delta^{(k)}C\bigr)^{-1},
  \end{equation*}
  where $\Delta^{(k)}$ is the quadratic form of \cite[(3.156)]{B-CMP99b} with exponential
  gauge fixing, the unadorned $C$ on the right is the elimination operator determined by the
  background configuration $V^{(k)}(W)$, and $C^{*}$ is its adjoint. The field increment is
  $Y_k(B)$ of the operations recalled in
  Lemma~\ref{8.2:lem-conditional-step}. Every Gaussian quantity in this chapter is taken with
  this $C^{(k)}$.
\end{enumerate}
\end{remark}

\begin{proof}
Items (b) and (c) restate, from the papers cited there, the form of the step and of its
covariance; only item (a) requires proof.

\emph{The invariance used.} By \cite[(2.16)]{B-CMP109}, the measure and the exponent are
invariant under the simultaneous substitution
$U_{k+1}\mapsto U_{k+1}^{u}$, $B\mapsto R(u)B$, in the notation $R(u)$ of
\cite[(2.16)]{B-CMP109} for the action of $u$ on the fluctuation field $B$. Passing from $W$ to
$W^{u}$ carries $U_{k+1}$ to $U_{k+1}^{u}$. We use this invariance, together with the invariance
of $\chi_k$ under the same substitution, in the form that the substitution leaves
$\chi_k(B)\,e^{\mathcal{S}_k(B;W)}\,d\mu_{C^{(k)}(W)}(B)$ unchanged. In the same substitution
the argument of the inserted factor is transformed by the gauge transformation $\tilde u$
of item (a): at $(R(u)B',W^{u})$ it is $\mathfrak{v}_k(B';W)^{\tilde u}$.

\emph{Invariant $F$.} In the numerator and the denominator of
\eqref{8.2:eq-step-covariance-1} with $W^{u}$ in place of $W$, substitute $B=R(u)B'$. By the
invariance just stated, the measure, $\chi_k$ and the exponent become those at $(B',W)$, and
the inserted factor becomes
$F\bigl(\mathfrak{v}_k(B';W)^{\tilde u}\bigr)$, which equals $F(\mathfrak{v}_k(B';W))$ since
$F$ is gauge invariant. The denominator is the case $F=1$. Therefore
\begin{equation*}
  \mathbb{E}_k[F\,|\,W^{u}]=\mathbb{E}_k[F\,|\,W].
\end{equation*}

\emph{Covariant $\Phi$.} Fix a site $x$. By covariance,
\begin{equation*}
  \Phi\bigl(\mathfrak{v}_k(B';W)^{\tilde u}\bigr)(x)
  =\operatorname{Ad}(\tilde u(x))\,\Phi(\mathfrak{v}_k(B';W))(x).
\end{equation*}
After the same substitution, the numerator for $\Phi$ at $W^{u}$, evaluated at $x$,
becomes
\begin{equation*}
  \int d\mu_{C^{(k)}(W)}(B')\,\chi_k(B')\,e^{\mathcal{S}_k(B';W)}\,
  \operatorname{Ad}(\tilde u(x))\,\Phi(\mathfrak{v}_k(B';W))(x).
\end{equation*}
The map $\operatorname{Ad}(\tilde u(x))$ is linear and, $\tilde u$ being fixed by $u$ and $W$,
does not depend on $B'$, so it commutes with the integral. The denominator is unchanged, as in
the invariant case. Hence
\begin{equation*}
  \mathbb{E}_k[\Phi\,|\,W^{u}](x)=\operatorname{Ad}(\tilde u(x))\,\mathbb{E}_k[\Phi\,|\,W](x),
\end{equation*}
which, $x$ being arbitrary, is the asserted equivariance.
\end{proof}

\begin{definition}[Boxes about the successive images of a plaquette]\label{8.2:not-plaquette-boxes}
Throughout, $k_0$, $r$ and $D$ are as fixed earlier in this section, and $k$ denotes a step with
$k\in[k_0,k_0+r)$. Let $p$ be a plaquette of the block lattice $\mathbb{T}^{(k_0)}$ of
Definition~\ref{1.5:def-block-average}.

The \emph{image} $p_k$ of $p$ at step $k$ is defined recursively:
\begin{itemize}
\item[(a)] $p_{k_0}:=p$;
\item[(b)] $p_{k+1}$ is the plaquette, or cell, of $\mathbb{T}^{(k+1)}$ that contains the block
image of $p_k$.
\end{itemize}

At step $k$ we write $\mathfrak{B}:=\mathfrak{B}_D(p_k)$ for the $D$-box about $p_k$. This is the set
of bonds of $\mathbb{T}^{(k)}$ within graph distance $D$ of $p_k$, that is, the set
$\mathfrak{B}_\ell(p_k)$ of Definition~\ref{8.2:def-local-observables} with $\ell=D$.

The number of real Gaussian variables in the box is
\begin{equation*}
n_{\mathfrak{B}}:=n\cdot\#\mathfrak{B}\le C_d\,n\,D^4 ,
\end{equation*}
where $n$ is as fixed earlier in this section and $C_d$ is a constant depending only on the
dimension $d$.
\end{definition}

\begin{definition}[Insertion of an observable along histories]\label{8.2:def-segment-insertion}
Let $s$ be a depth below the cutoff and put $k_0:=K-s$; let $r$ be the number of steps of the segment
from level $k_0$ to level $k_0+r$, and consider the expansion at level $k_0+r$: the terms
of \cite[(2.18)]{B-CMP119} at $k=k_0+r$, grouped by their histories $h$ as in clause~(ii-a) of the
two-point witness theorem. Let
$G$ be a bounded gauge-invariant Borel function of the block field $V_{k_0}$, and write $\sup|G|$
for its supremum norm. Run Ba{\l}aban's procedure from level $k_0$, with initial density
$\rho_{k_0}$ and with the factor $G(V_{k_0})$ carried along through the operations of each history
in the following way:
\begin{itemize}
\item[(a)] at each application of the operation $\mathbf{T}$, the carried function is pushed
forward;
\item[(b)] at each decomposition of unity, the carried function goes with the shares of the
decomposition;
\item[(c)] at a pre-integration, the carried function enters only the numerator of the quotient
formed in Ba{\l}aban's pre-integration \cite{B-CMP119};
\item[(d)] at each removal of a $\delta$-function, the carried function is multiplied by the
corresponding Radon--Nikodym factor.
\end{itemize}
For each history $h$ this produces a function $G^{\langle h\rangle}(W)$ of the block field
$W:=V_{k_0+r}$, called the \emph{function carried along $h$}. Let $\tau_h$ denote the measure of
the history $h$, and let $\nu^{B}_{K,s-r}$ denote the measure on configurations of the block field
at depth $s-r$, that is at level $k_0+r$, which enters the two-point witness theorem. We say
that the \emph{insertion of an observable along histories} holds for $G$ on the segment from level
$k_0$ to level $k_0+r$ if the following three requirements are met.
\begin{itemize}
\item[(i)] For every history $h$ of the expansion at level $k_0+r$,
\begin{equation*}
\bigl|G^{\langle h\rangle}\bigr|\le\sup|G|.
\end{equation*}
\item[(ii)] The carried functions and the measures of the histories satisfy
\begin{equation}\label{8.2:eq-segment-insertion-1}
\begin{gathered}
\sum_h\int G^{\langle h\rangle}(W)\,d\tau_h(W)=\int G(V)\,\rho_{k_0}(V)\,dV,\\
\sum_h\int d\tau_h=\int\rho_{k_0}=\int\rho_0,\\
\Bigl(\int\rho_0\Bigr)^{-1}\sum_h\tau_h(dW)=\nu^{B}_{K,s-r}(dW).
\end{gathered}
\end{equation}
\item[(iii)] Let $p$ be a plaquette and let $h$ be a history that is clean at $p$ through the
segment, in the sense of the definition of histories clean at a plaquette through a segment in
this chapter. Then, near $p$, the carried function is the composition of the one-step conditional
expectations $\mathbb{E}_{k_0},\mathbb{E}_{k_0+1},\dots,\mathbb{E}_{k_0+r-1}$ of the small-field
step with an inserted observable (Lemma~\ref{8.2:lem-conditional-step}),
\begin{equation*}
G^{\langle h\rangle}
=\mathbb{E}_{k_0+r-1}\Bigl[\cdots \mathbb{E}_{k_0+1}\bigl[\mathbb{E}_{k_0}[G]\bigr]\cdots\Bigr]
\quad\text{near }p,
\end{equation*}
and the operations of $h$ outside the successive localisation boxes around $p$ act on
$G^{\langle h\rangle}$ only through the tails of part~(i) of the localisation of the sourced step.
\end{itemize}
\end{definition}

This property is the segment reading of clause~(ii-a) of the two-point witness theorem, which is
an identity: the tower property of the one-step conditional expectations, disintegrated along
histories.

\begin{lemma}[Analytic form and parity splitting of the box exponent]\label{8.2:lem-box-exponent}
Let $k\in[k_0,k_0+r)$ be a step. Let $p_k$, the box $\mathfrak{B}=\mathfrak{B}_D(p_k)$ and the number
$n_{\mathfrak{B}}$ of real Gaussian variables in $\mathfrak{B}$ be as in
Notation~\ref{8.2:not-plaquette-boxes}. Write $B\in\mathbb{R}^{n_{\mathfrak{B}}}$ for the box
variables, and let $C$, $U_k$, $\mathfrak{v}_k(B;W)$ and $V^{(k)}(W)$ be as in
Notation~\ref{8.2:not-small-field-step}. Let $W$ lie in the level-$(k+1)$ small-field window near
$p_{k+1}$.
For a level $j<k$ let $\mathfrak{l}_j$ be a number such that every large-field region created at
level $j$ is live for at most $\mathfrak{l}_j$ steps. For a history
$h$ put
\begin{equation}\label{8.2:eq-box-exponent-1}
\mathfrak{S}'_k(h):=\sum_{j<k}\ \sum_{\square}(\mathfrak{l}_j+1)\,|\square|_k,
\end{equation}
where the inner sum runs over the level-$j$ large-field cubes $\square$ of $h$ under $\mathfrak{B}$,
and $|\square|_k$ is the volume of $\square$ in step-$k$ units. Let $\tau_*>0$, and let $h$ be a
$p$-clean history, that is, a history on the clean event at $p$ (on such a history the nearest live
or current large-field object is at distance at least $D$ from $p_k$), on which
$\mathfrak{S}'_k(h)\le\tau_*\#\mathfrak{B}$.
Let $\mathcal{P}_k(B;W)$, the \emph{exponent of the step}, be the sum of five summands:
the main action $P^{(k)}$; the increments
\begin{gather*}
E_k\bigl(U_k(\mathfrak{v}_k(B))\bigr)-E_k\bigl(U_k(V^{(k)})\bigr),\qquad
R_k\bigl(U_k(\mathfrak{v}_k(B))\bigr)-R_k\bigl(U_k(V^{(k)})\bigr),\\
B_k\bigl(U_k(\mathfrak{v}_k(B))\bigr)-B_k\bigl(U_k(V^{(k)})\bigr)
\end{gather*}
of the densities $E_k$, $R_k$, $B_k$ of Notation~\ref{8.2:not-full-effective-action}; and the lag
of the running coupling, the term of \cite[(2.24)]{B-CMP119}. We call these summands the main
action, the $E_k$-, $R_k$- and $B_k$-increments, and the lag.
Let $\mathfrak{b}_k$ be the window width and $c_w,C_w>0$ the constants of the window, all fixed by
the small-field window convention of hypothesis~(f) below, let $\varrho_{\mathcal{P}}>0$ be the
imaginary width of the tube of hypothesis~(b) below, and put
\begin{equation*}
\mathcal{T}_k:=\bigl\{B\in\mathbb{C}^{n_{\mathfrak{B}}}:\ |\operatorname{Re}(CB)_b|<C_w\mathfrak{b}_k+1,\
|\operatorname{Im}(CB)_b|<\varrho_{\mathcal{P}}\ \text{for all }b\bigr\}.
\end{equation*}
Assume the following.
\begin{itemize}
\item[(a)] (Analyticity of the effective densities.) Along $h$ the following holds. The localised
terms of the effective action at
every level $j\le k$ are analytic and gauge invariant on Ba{\l}aban's complex domains. On those
domains each term attached to a localisation domain $X$ of level $j$ is bounded by a constant
times $e^{-\kappa d_j(X)}$.
\item[(b)] (Uniform tube bound.) The localised terms of the exponent $\mathcal{P}_k$ of the step
are
holomorphic on a complex tube about the small-field window, of imaginary width
$\varrho_{\mathcal{P}}$, which contains the tube $\mathcal{T}_k$, and each term
attached to a localisation domain $X$ of level $j$ is bounded there, uniformly in $K$, by a
constant times $e^{-\kappa d_j(X)}$; tube and window do not depend on $K$, and the width
$\varrho_{\mathcal{P}}>0$ and the constant depend on $(L,\mathfrak{p},N)$ only.
\item[(c)] The effective action at level $k$ has the format of \cite[Theorem~3]{B-CMP109}.
\item[(d)] The inductive hypotheses \cite[(2.23)--(2.31)]{B-CMP119} of the complete expansion hold
along $h$.
\item[(e)] The coupling-flow bounds of clause~(0) of Theorem~\ref{3.1:thm-main} hold: for
$\nu\ge1$,
\begin{equation*}
x_{k-\nu}\ \ge\ x_k+\lambda^{\sharp}\nu-c_5,
\end{equation*}
and $\beta_j\le B^{\sharp}$ for every level $j$.
\item[(f)] (Small-field window convention.) The small-field characteristic functions of the step
are those of \cite[(2.9)]{B-CMP109} for the small-field part of the expansion, and those of
\cite[(2.17), (2.21)]{B-CMP119} for the complete expansion. The parameters $\varepsilon_k$,
$\delta_k$, the window width $\mathfrak{b}_k$ and the constants $c_w$, $C_w$ of the window are
fixed by this convention.
\item[(g)] (Covariance bound.) For every box variable $i\in\mathfrak{B}$ and every bond $b$,
\begin{equation*}
(C^{(k)})_{ii}\le 2(1+\|C\|^2)/\gamma_0,\qquad (CC^{(k)}C^*)_{bb}\le 2(1+\|C\|^2)/\gamma_0 ,
\end{equation*}
where $C^{(k)}$ is the covariance of the Gaussian fluctuation field of step $k$ and $\gamma_0>0$ is
a constant fixed independently of $k$, of $K$ and of the history $h$.
\item[(h)] (Curvature bound.) For every $B\in\mathcal{T}_k$, every level $j<k$ and every
localisation domain $X$ of level $j$ meeting $\mathfrak{B}$, the level-$j$ curvature on $X$ of the
configuration $U_k(\mathfrak{v}_k(B;W))$, built from the step-$k$ background $V^{(k)}$ and the
fluctuation $B$, is at most
\begin{equation*}
C_{\mathrm{curv}}\,L^{-2(k-j)}\bigl(g_k|dB|+\kappa(V^{(k)})\bigr)
\end{equation*}
with a constant $C_{\mathrm{curv}}$ depending on $(L,\mathfrak{p},N)$ only. Here $dB$ is the
lattice derivative of $B$, and
$\kappa(V^{(k)})$ is the value of the curvature functional $\kappa(\cdot)$ of
Lemma~\ref{8.2:lem-flat-background} (not the decay rate $\kappa$) at the background
$V^{(k)}=V^{(k)}(W)$.
\end{itemize}
There is a constant $C_{\mathcal{P}}<\infty$, depending on $(L,\mathfrak{p},N)$ only, such that the
following hold.
\begin{itemize}
\item[(w)] (Window.) The characteristic functions of the step near $p$ define a measurable set
$W_k\subset\mathbb{R}^{n_{\mathfrak{B}}}$ of box variables such that
\begin{equation}\label{8.2:eq-box-exponent-2}
\{B\in\mathbb{R}^{n_{\mathfrak{B}}}:\ \|B\|_\infty<c_w\mathfrak{b}_k\}\ \subset\ W_k\ \subset\
\{B:\ |(CB)_b|<C_w\mathfrak{b}_k\ \text{for all bonds }b\}.
\end{equation}
Let $\gamma_k$ be the marginal on the box variables of the centred Gaussian measure with
covariance $C^{(k)}$ of hypothesis~(g). Put
\begin{align*}
v_k&:=\max\Bigl\{\max_{i\in\mathfrak{B}}(C^{(k)})_{ii},\ \max_b(CC^{(k)}C^*)_{bb}\Bigr\},\\
t_k&:=2n_{\mathfrak{B}}\,e^{-c_w^2\mathfrak{b}_k^2/(2v_k)}.
\end{align*}
Then $v_k\le 2(1+\|C\|^2)/\gamma_0$, and
\begin{equation}\label{8.2:eq-box-exponent-3}
\gamma_k(W_k^{c})\le t_k,\qquad \gamma_k\bigl(W_k\setminus(-W_k)\bigr)\le t_k .
\end{equation}
We call this property of the window \emph{$t_k$-symmetry}.
\item[(an)] (Analytic form.) Let $\mathcal{P}_{k,\mathfrak{B}}(B;W)$, the \emph{box part of the
exponent}, be the sum of those localised terms of $\mathcal{P}_k$ whose localisation domains lie
inside $\mathfrak{B}$. Then $\mathcal{P}_{k,\mathfrak{B}}$ is holomorphic on the complex tube
$\mathcal{T}_k$.
Let $\mathcal{P}^{\mathrm{far}}_{k,\mathfrak{B}}$ be the sum of the localised terms of the
$B_k$-increment whose localisation domains lie inside $\mathfrak{B}$. Under $B\mapsto-B$ the
difference $\mathcal{P}_{k,\mathfrak{B}}-\mathcal{P}^{\mathrm{far}}_{k,\mathfrak{B}}$ splits by
parity as $P_k^{-}+P_k^{+}$, so that
\begin{equation*}
\mathcal{P}_{k,\mathfrak{B}}=P_k^{-}+P_k^{+}+\mathcal{P}^{\mathrm{far}}_{k,\mathfrak{B}},
\end{equation*}
with
\begin{equation}\label{8.2:eq-box-exponent-4}
\begin{aligned}
\sup_{\mathcal{T}_k}|P_k^{-}|&\le\vartheta_k:=C_{\mathcal{P}}\,g_k\,n_{\mathfrak{B}}\,
\mathfrak{b}_k^{3}+C_{\mathcal{P}}\,g_k^{2}\,B^{\sharp}\tau_*\,n_{\mathfrak{B}}\,\mathfrak{b}_k^{2}\\
&\qquad+C_{\mathcal{P}}\,n_{\mathfrak{B}}\,g_k^{\kappa_0-4},\\
\sup_{\mathcal{T}_k}|P_k^{+}|&\le\vartheta_k^{+}:=C_{\mathcal{P}}\,g_k^{2}\,n_{\mathfrak{B}}\,
\mathfrak{b}_k^{4}\,(1+B^{\sharp}\tau_*)+C_{\mathcal{P}}\,n_{\mathfrak{B}}\,g_k^{\kappa_0-4}.
\end{aligned}
\end{equation}
Here $\kappa_0$ is the exponent of the bound \cite[(2.31)]{B-CMP119} on the terms $R^{(j)}(X)$.
The terms $\mathcal{P}^{\mathrm{far}}_{k,\mathfrak{B}}$ are not counted in
\eqref{8.2:eq-box-exponent-4}: on a $p$-clean history their box part is at most a constant times
$e^{-\kappa D}\#\mathfrak{B}$, and they are absorbed into the tail of part~(i) of the localisation
of the sourced step.
\end{itemize}
\end{lemma}

\begin{proof}
\emph{The window.} The set $W_k$ is the set of box variables $B$ at which every small-field
characteristic function of the step whose argument lies in $\mathfrak{B}$ equals one; it is
measurable, being the set where a finite product of characteristic functions of measurable sets
equals one. Its inclusions \eqref{8.2:eq-box-exponent-2} are read off the small-field
characteristic functions of hypothesis~(f).

For the small-field part of the expansion these are the factors
$\prod_b\chi(|B'(b)|<\varepsilon_1)$ of \cite[(2.9)]{B-CMP109}. For the variables denoted $b_0(c)$
in \cite[p.~267]{B-CMP109} (Ba{\l}aban's notation, distinct from the constant $b_0=(N^2-1)/2$ of
Notation~\ref{0.2:not-glossary}) they are
the same conditions with $\varepsilon_1$ replaced by $O(\varepsilon_1)$, as stated in
\cite[p.~267]{B-CMP109}. For the complete expansion they are the characteristic functions of
\cite[(2.17), (2.21)]{B-CMP119}, with $\varepsilon_k$, $\delta_k$ and $\mathfrak{b}_k$ as fixed by
the convention of hypothesis~(f).

The bound $v_k\le 2(1+\|C\|^2)/\gamma_0$ follows from hypothesis~(g), since $v_k$ is the maximum of
the diagonal entries bounded there.

Now let $B\in W_k^{c}$. By the first inclusion in \eqref{8.2:eq-box-exponent-2}, $B$ does not lie in
the cube $\{\|B\|_\infty<c_w\mathfrak{b}_k\}$. Hence
\begin{equation*}
W_k^{c}\subset\bigcup_{i}\{|B_i|\ge c_w\mathfrak{b}_k\},
\end{equation*}
where the union runs over the $n_{\mathfrak{B}}$ box variables. Under $\gamma_k$ the variable $B_i$
is a centred one-dimensional Gaussian of variance $(C^{(k)})_{ii}\le v_k$. The one-dimensional
Gaussian tail bound therefore gives
\begin{equation*}
\gamma_k(|B_i|\ge c_w\mathfrak{b}_k)\le 2e^{-c_w^2\mathfrak{b}_k^2/(2v_k)}.
\end{equation*}
Summing over the $n_{\mathfrak{B}}$ indices gives the first bound in \eqref{8.2:eq-box-exponent-3}.

For the second bound, note that $W_k\setminus(-W_k)\subset(-W_k)^{c}$. The cube
$\{\|B\|_\infty<c_w\mathfrak{b}_k\}$ is invariant under $B\mapsto-B$, so $-W_k$ also contains it.
The argument just given therefore applies to $-W_k$ in place of $W_k$. It yields
$\gamma_k((-W_k)^{c})\le t_k$, and hence the second bound in \eqref{8.2:eq-box-exponent-3}.

\emph{The tube and the parity splitting.} The map $B\mapsto CB$ is linear, and the conditions
defining $\mathcal{T}_k$ involve only $|\operatorname{Re}(CB)_b|$ and
$|\operatorname{Im}(CB)_b|$. Hence $-\mathcal{T}_k=\mathcal{T}_k$. For $B\in\mathcal{T}_k$ we put
\begin{equation*}
\begin{split}
P_k^{\mp}(B):=\tfrac12\Bigl(&\bigl(\mathcal{P}_{k,\mathfrak{B}}
-\mathcal{P}^{\mathrm{far}}_{k,\mathfrak{B}}\bigr)(B;W)\\
&\mp\bigl(\mathcal{P}_{k,\mathfrak{B}}-\mathcal{P}^{\mathrm{far}}_{k,\mathfrak{B}}\bigr)(-B;W)\Bigr).
\end{split}
\end{equation*}
Then $\mathcal{P}_{k,\mathfrak{B}}=P_k^{-}+P_k^{+}+\mathcal{P}^{\mathrm{far}}_{k,\mathfrak{B}}$. The
function $P_k^{-}$ is odd and $P_k^{+}$ is even, and, since $\mathcal{T}_k=-\mathcal{T}_k$, both are
holomorphic on $\mathcal{T}_k$ as soon as
$\mathcal{P}_{k,\mathfrak{B}}-\mathcal{P}^{\mathrm{far}}_{k,\mathfrak{B}}$ is.

Four of the five summands of $\mathcal{P}_k$ are treated below as functions of the field $g_kCB$,
with explicit prefactors; the $B_k$-increment is bounded separately. Terms of odd order in $g_kCB$
are odd in $B$, and terms of even order are
even. The splitting is thus read off this structure, summand by summand, and not off any bound. The
width $\varrho_{\mathcal{P}}$ is that of hypothesis~(b). The constant $C_{\mathcal{P}}$ collects
the constant of hypothesis~(b) and the constants absorbed below; it depends on $(L,\mathfrak{p},N)$
only.

Hypothesis~(b) is used in the form stated there: the bound holds for each localisation domain,
with the weights $e^{-\kappa d_j(X)}$. Estimate \eqref{8.2:eq-box-exponent-4} uses it in this form.
Holomorphy of $\mathcal{P}_{k,\mathfrak{B}}$ on $\mathcal{T}_k$ follows from hypotheses~(a)
and~(b) and from the analyticity of each summand recorded below. We take the five summands in turn.

\emph{The main action and the measure.} By hypothesis~(c), $P^{(k)}$ has the format of
\cite[Theorem~3]{B-CMP109}. In particular $P^{(k)}$ vanishes at $g_k=0$ \cite[\S2]{B-CMP109}. It is
the sum of two kinds of term:
\begin{itemize}
\item $g_k^{-2}$ times an analytic function of the field $g_kCB$ beginning at third order;
\item the Jacobian terms and the terms $\sigma$, $\mathbf{G}_3$ and $V$ of
\cite[(2.12)]{B-CMP109}, each an analytic function of $g_kCB$
beginning at first order, with a prefactor that makes it $O(g_k)$.
\end{itemize}
The map $\tilde D$ entering $\mathfrak{v}_k$ has an expansion beginning with quadratic terms,
as stated in \cite[p.~267]{B-CMP109}.

The linear term of the pure action cancels by criticality. Indeed, \cite[\S2]{B-CMP109} states that,
in its notation, $\langle\delta A',J\rangle=0$ (this is the criticality just invoked), and hence
that the only term with a negative power of $g_k$ is the
action evaluated at $U_{k+1}$. In the first kind of term, the third-order and fourth-order
contributions carry the factors
\begin{equation*}
g_k^{-2}(g_k|CB|)^3=g_k|CB|^3,\qquad g_k^{-2}(g_k|CB|)^4=g_k^2|CB|^4 .
\end{equation*}
The second kind of term contributes $O(g_k)$ at first order and $O(g_k^2)$ at second order.
Consequently, per localisation domain on the tube, the part of this summand odd in $B$ is
$O(g_k)(|B|+|B|^3)$. Its even part is $O(g_k^2)(|B|^2+|B|^4)$.

\emph{The $E_k$-increment.} This summand is
\begin{equation*}
E_k\bigl(U_k(\mathfrak{v}_k(B))\bigr)-E_k\bigl(U_k(V^{(k)})\bigr),\qquad
E_k=\sum_j\sum_X E^{(j)}(X).
\end{equation*}
By hypothesis~(a), in the form of \cite[(1.18)]{B-CMP109} and \cite[(2.25)--(2.27)]{B-CMP119}, each
$E^{(j)}(X)$ is analytic and gauge invariant. Moreover
$|E^{(j)}(X)|\le E_0e^{-\kappa d_j(X)}$, with $E_0$ the constant of \cite[(1.18)]{B-CMP109}, on the
domains with the absolute radii $\alpha_0$, $\alpha_1$ of \cite[(1.18)]{B-CMP109}. By
hypothesis~(b) the increment is holomorphic and bounded uniformly in $K$ on the $K$-independent
tube.

By the Cauchy estimate in the field, the first-order part of the increment is $O(g_k|B|)$ times a
covariant coefficient. Its second-order part is $O(g_k^2|B|^2)$.

The ultraviolet sum over the levels $j$ converges by the irrelevance bound
\cite[(0.29)--(0.30)]{B-CMP109}. With $\eta=L^{-k}$, so that $L^j\eta=L^{-(k-j)}$, the level-$j$
contributions carry the factor $(L^j\eta)^{4+\alpha}$, where $\alpha>0$ is the exponent of the
irrelevance bound (not to be confused with the radii $\alpha_0$, $\alpha_1$), and the level sum is
at most
\begin{equation*}
O(1)\,(1-L^{-\alpha})^{-1}.
\end{equation*}
This is where the running coupling $g_k$, and not $g_0$, enters. The marginal $F^2$-content of
every finer level was extracted into the coefficients $\beta_j$ in Ba{\l}aban's construction
\cite[\S0]{B-CMP109}.

\emph{The $R_k$-increment.} This summand is
\begin{equation*}
R_k\bigl(U_k(\mathfrak{v}_k(B))\bigr)-R_k\bigl(U_k(V^{(k)})\bigr).
\end{equation*}
By hypothesis~(d), in the form of \cite[(2.30)--(2.31)]{B-CMP119}, it is a sum over the terms
$R^{(j)}(X)$ present in the history. Each $R^{(j)}(X)$ is analytic and gauge invariant on
Ba{\l}aban's complex domain $U_j(X,\alpha_{0j},\alpha_{1j})$ of level-$j$ configurations on $X$
with radii $\alpha_{0j}$, $\alpha_{1j}$, and by \cite[(2.31)]{B-CMP119}
\begin{equation*}
|R^{(j)}(X)|\le g_j^{\kappa_0}e^{-\kappa d_j(X)},
\end{equation*}
with $\kappa_0$ arbitrarily large \cite[\S2]{B-CMP119}. On backgrounds smooth at scale $j$ the
vacuum-subtracted value of $R^{(j)}(X)$ is
\begin{equation*}
O(1)\,(L^j\eta)^4\,g_j^{\kappa_0}\,e^{-\kappa d_j(X)}.
\end{equation*}
This is the treatment of marginal terms in \cite[\S2]{B-CMP119}, and it is used here as stated
there.

The mechanism producing the factor $(L^j\eta)^4=L^{-4(k-j)}$ is
Lemma~\ref{8.2:lem-flat-background}\,(a)--(b), applied at level $j$. A gauge-invariant analytic
function of $U$ restricted to $X$ that vanishes at $U=1$ is of second order in the
level-$j$ curvature of $U$ on $X$. For the configuration $U_k(\mathfrak{v}_k(B;W))$ of the step-$k$
background and fluctuation, with $B\in\mathcal{T}_k$, this curvature is at
most
\begin{equation*}
C_{\mathrm{curv}}(L^j\eta)^2\bigl(g_k|dB|+\kappa(V^{(k)})\bigr)
\end{equation*}
by hypothesis~(h), since $(L^j\eta)^2=L^{-2(k-j)}$.

We now sum over the box. The count below holds for every history, and no event is needed. Let
$\mathcal{D}_j$ be the set of localisation domains of level $j$. The number of
$X\in\mathcal{D}_j$ with $X\cap\mathfrak{B}\ne\emptyset$ and $d\le d_j(X)<d+1$ is at most
\begin{equation*}
C_{\mathrm{count}}\,\#\mathfrak{B}\,L^{4(k-j)}e^{c'd},
\end{equation*}
where $C_{\mathrm{count}}$ and $c'$ are constants, $c'$ being the exponential rate at which this
number grows with $d$.
The vacuum subtraction is a second-order Cauchy estimate in the field. It costs the factor
$\alpha_{\cdot j}^{-2}$, where $\alpha_{\cdot j}$ denotes the radii $\alpha_{0j}$, $\alpha_{1j}$.
The radii of \cite[(2.28)]{B-CMP119} are $g_j$ times a factor polynomial in $\log x_j$. Hence
$\alpha_{\cdot j}^{-2}\le x_j\operatorname{polylog}(x_j)$, where $\operatorname{polylog}(x_j)$
denotes a factor polynomial in $\log x_j$.

The factor $(L^j\eta)^4=L^{-4(k-j)}$ compensates the factor $L^{4(k-j)}$ of the count. We also have
$g_j^{\kappa_0}=x_j^{-\kappa_0/2}$. With $\nu=k-j$ we therefore obtain
\begin{align*}
&\sum_{j<k}\sum_{X}(L^j\eta)^4g_j^{\kappa_0}e^{-\kappa d_j(X)}\,\alpha_{\cdot j}^{-2}\\
&\qquad\le C_{\mathrm{sum}}\,\#\mathfrak{B}\sum_{\nu\ge1}x_{k-\nu}^{1-\kappa_0/2}
\operatorname{polylog}(x_{k-\nu})\\
&\qquad\le C'_{\mathrm{sum}}\,\#\mathfrak{B}\,x_k^{2-\kappa_0/2},
\end{align*}
with constants $C_{\mathrm{sum}}$ and $C'_{\mathrm{sum}}$.
The last inequality is hypothesis~(e), $x_{k-\nu}\ge x_k+\lambda^{\sharp}\nu-c_5$, applied to each
term of the sum over $\nu$. Since $x_k^{2-\kappa_0/2}=g_k^{\kappa_0-4}$ and
$\#\mathfrak{B}\le n_{\mathfrak{B}}$, this bounds the box part of the $R_k$-increment on
$\mathcal{T}_k$ by a constant times $n_{\mathfrak{B}}g_k^{\kappa_0-4}$. Since
$\mathcal{T}_k=-\mathcal{T}_k$, the odd and the even part of a function $f$ on $\mathcal{T}_k$ are
each bounded there by its supremum, because
$\tfrac12|f(B)\mp f(-B)|\le\sup_{\mathcal{T}_k}|f|$ for $B\in\mathcal{T}_k$. The odd and the even
part of the $R_k$-increment therefore give the member
$C_{\mathcal{P}}n_{\mathfrak{B}}g_k^{\kappa_0-4}$ of $\vartheta_k$ and of $\vartheta_k^{+}$ in
\eqref{8.2:eq-box-exponent-4}, respectively. For
$\kappa_0\ge8$ this member is at most $g_k^4$ times $C_{\mathcal{P}}n_{\mathfrak{B}}$.

\emph{The $B_k$-increment.} Its terms are localised close to large-field regions
\cite[\S2]{B-CMP119}, and they are bounded and exponentially localised. On a $p$-clean history the
nearest live or current large-field object is at distance at least $D$ from $p_k$. Their box part
is therefore at most
\begin{equation*}
C_{\mathrm{far}}e^{-\kappa D}\,\#\mathfrak{B}
\end{equation*}
for a constant $C_{\mathrm{far}}$.
This summand depends on $B$ only through terms whose supremum is at most
$C_{\mathrm{far}}e^{-\kappa D}$. It is
absorbed into the tail of part~(i) of the localisation of the sourced step and is not counted in
\eqref{8.2:eq-box-exponent-4}.

\emph{The lag of the running coupling.} By \cite[(2.24)]{B-CMP119}, on the fine volume under
$\mathfrak{B}$ the main action is
\begin{equation*}
\sum_x\eta^4g_k(x)^{-2}|F_x|^2,\qquad g_k(x)^{-2}=g_k^{-2}\bigl(1+\Delta(x)\bigr),
\end{equation*}
where $x$ runs over the fine points under $\mathfrak{B}$, $F_x$ is the curvature of the fine field
at $x$, and
\begin{equation*}
\Delta(x):=g_k^2\sum_{j\text{ missed at }x}\beta_j .
\end{equation*}
Here a level $j$ is missed at $x$ if $x$ lay outside the region $\Lambda_j^{\sim}$ of
\cite[(2.24)]{B-CMP119}. By hypothesis~(e), $\beta_j\le B^{\sharp}$, and therefore
\begin{equation*}
\Delta(x)\in\bigl[0,\ g_k^2B^{\sharp}\cdot\#\{j:\ x\notin\Lambda_j^{\sim}\}\bigr].
\end{equation*}

A fine point lies outside $\Lambda_j^{\sim}$ only while it lies inside the enlargement of a
large-field region that is live at step $j$. By the choice of $\mathfrak{l}_j$ in the statement, a
region created at level $j$ is live for at most $\mathfrak{l}_j$ steps. Hence, with
$\mathfrak{S}'_k(h)$ as in \eqref{8.2:eq-box-exponent-1},
\begin{equation*}
\sum_{x\text{ under }\mathfrak{B}}\eta^4\Delta(x)\le g_k^2B^{\sharp}\,\mathfrak{S}'_k(h).
\end{equation*}

The contribution of the lag to the exponent is, per fine point,
\begin{equation*}
\Delta(x)\Bigl[2g_k^{-1}\bigl\langle F_{\mathrm{bg}},\,g_k\,\delta F_1(B)\bigr\rangle
+|\delta F_1(B)|^2+\cdots\Bigr],
\end{equation*}
where $F_{\mathrm{bg}}$ is the curvature of the background and $\delta F_1(B)$ the first variation
of the curvature in the box variables. Its odd part is $O(\Delta\,\kappa(V^{(k)})|B|)$ and its even
part is $O(\Delta|B|^2)$. By hypothesis, $\mathfrak{S}'_k(h)\le\tau_*\#\mathfrak{B}$.
Under this hypothesis the box supremum of the lag is at most
\begin{equation*}
C_{\mathrm{lag}}\,g_k^2\,B^{\sharp}\,\tau_*\,n_{\mathfrak{B}}\,\mathfrak{b}_k^2 ,
\end{equation*}
for a constant $C_{\mathrm{lag}}$. By the bound on parity parts stated for the $R_k$-increment,
this bound holds for the odd part and for the even part of the lag separately. The odd part of the
lag thus gives the member $C_{\mathcal{P}}g_k^2B^{\sharp}\tau_*n_{\mathfrak{B}}\mathfrak{b}_k^2$ of
$\vartheta_k$, and its even part is contained in $\vartheta_k^{+}$, as the member
$B^{\sharp}\tau_*$ of
\eqref{8.2:eq-box-exponent-4}.

Among the contributions to the box part of the exponent, the lag is the only one that is not small
for every history. In the worst case $\mathfrak{S}'_k(h)$ is of order
$\#\mathfrak{B}\,\mathfrak{l}_0$, and the lifetime $\mathfrak{l}_0$ of regions created at level $0$
cannot in general be chosen independently of $K$.

\emph{Collection.} The main action and measure, the $E_k$-increment, the $R_k$-increment and the
lag are each holomorphic on $\mathcal{T}_k$, with the odd and even parts displayed above. On
$\mathcal{T}_k$ every bond $b$ satisfies
\begin{equation*}
|(CB)_b|\le|\operatorname{Re}(CB)_b|+|\operatorname{Im}(CB)_b|<C_w\mathfrak{b}_k+1+\varrho_{\mathcal{P}},
\end{equation*}
and in the members below this quantity is replaced by a constant multiple of $\mathfrak{b}_k$, the
constant being absorbed into $C_{\mathcal{P}}$. The per-domain bounds, with their weights
$e^{-\kappa d_j(X)}$ of hypothesis~(b), are summed over the localisation domains inside
$\mathfrak{B}$ on $\mathcal{T}_k$; this gives the following members of
\eqref{8.2:eq-box-exponent-4}:
\begin{itemize}
\item the odd parts of the main action and measure and of the $E_k$-increment, of orders $g_k|B|$
and $g_k|B|^3$, give the member $C_{\mathcal{P}}g_kn_{\mathfrak{B}}\mathfrak{b}_k^3$ of
$\vartheta_k$;
\item the even parts of these two summands, of orders $g_k^2|B|^2$ and $g_k^2|B|^4$, give the
member $C_{\mathcal{P}}g_k^2n_{\mathfrak{B}}\mathfrak{b}_k^4$ of $\vartheta_k^{+}$;
\item the odd part of the lag gives the member
$C_{\mathcal{P}}g_k^2B^{\sharp}\tau_*n_{\mathfrak{B}}\mathfrak{b}_k^2$ of $\vartheta_k$, and its even
part the member $C_{\mathcal{P}}g_k^2n_{\mathfrak{B}}\mathfrak{b}_k^4B^{\sharp}\tau_*$ of
$\vartheta_k^{+}$;
\item the odd and the even part of the $R_k$-increment give the member
$C_{\mathcal{P}}n_{\mathfrak{B}}g_k^{\kappa_0-4}$ of $\vartheta_k$ and of $\vartheta_k^{+}$,
respectively.
\end{itemize}
The $B_k$-increment, whose box terms form $\mathcal{P}^{\mathrm{far}}_{k,\mathfrak{B}}$, is absorbed
into the tail of part~(i) of the localisation of the sourced step, as described
above. This proves \eqref{8.2:eq-box-exponent-4}.
\end{proof}

\begin{definition}[The localisation hypothesis for one clean sourced step]
\label{8.2:def-clean-step-localisation}
We use the step $k$ in the form of the small-field step with an inserted observable
(Lemma~\ref{8.2:lem-conditional-step}). There $W$ is the block configuration at level $k+1$,
$C^{(k)}(W)$ is the covariance of the Gaussian integration over the fluctuation field $B$, and
$\mathfrak{v}_k(B;W)$ is the configuration into which the observable is inserted;
$U_{k+1}(W)$ is the configuration entering the prefactor of that lemma, and $V^{(k)}(W)$ is the
background configuration of the step. For an observable $G$ we write $\mathbb{E}_k[G\,|\,W]$ for
the conditional expectation at step $k$ of $G$ given $W$: the integral of
Lemma~\ref{8.2:lem-conditional-step} with $F=G$, divided by the same integral with $F=1$.

Fix a plaquette $p$. Let $p_k$ and $p_{k+1}$ denote the data at levels $k$ and $k+1$ attached
to $p$; observables are admitted near them. Let $\mathfrak{G}_k(p_k;\ell,M,\varrho)$ denote the
class of bounded, local, analytic, gauge-invariant observables admitted at level $k$ near $p_k$,
with parameters $(\ell,M,\varrho)$. Here $M$ and $\varrho$ are
parameters of the class; they are not the cube size $M=L^{m'}$ among the parameters
$\mathfrak{p}$, nor the constant $\varrho$ of the glossary (Notation~\ref{0.2:not-glossary}).
Steps that are $p$-clean with box
radius $D$ are understood in the sense of Lemma~\ref{8.2:lem-box-exponent}; for such a step we
write $\mathfrak{B}$ for its box, of radius $D$, and $b,b',b''$ for bonds; we further use the
following objects.
\begin{itemize}
\item[--] The measure $\gamma_k$ is the marginal on the box variables of the centred Gaussian
measure $d\mu_{C^{(k)}(W)}$ of Lemma~\ref{8.2:lem-conditional-step}.
\item[--] The set $\mathcal{W}_k$ is the window of the box variables of
Lemma~\ref{8.2:lem-box-exponent}: the set of box variables cut out by the step's characteristic
functions, restricted to the box; $\mathbf{1}_{\mathcal{W}_k}$ is its indicator.
\item[--] The function $\mathcal{P}_{k,\mathfrak{B}}(B;W)$ is the box exponent of
Lemma~\ref{8.2:lem-box-exponent}: the exponent of the step of
Lemma~\ref{8.2:lem-conditional-step} restricted to the box.
\item[--] The number $\varrho_{\mathcal{P}}$ is the radius of the analytic format of the step's
exponent, Lemma~\ref{8.2:lem-box-exponent}.
\item[--] The \emph{window of configurations} is the small-field window at level $k+1$ near
$p_{k+1}$ of Lemma~\ref{8.2:lem-box-exponent}, a set of configurations $W$; the
\emph{complex window} is the complex extension of that set of configurations.
\end{itemize}
The \emph{box expectation} of a function $F$ of the box variables is the finite-dimensional
integral
\begin{equation}\label{8.2:eq-clean-step-localisation-1}
\langle F\rangle_{k,\mathfrak{B}}(W)
:=\frac{\displaystyle\int d\gamma_k(B)\,\mathbf{1}_{\mathcal{W}_k}(B)\,
e^{\mathcal{P}_{k,\mathfrak{B}}(B;W)}\,F(B)}
{\displaystyle\int d\gamma_k(B)\,\mathbf{1}_{\mathcal{W}_k}(B)\,
e^{\mathcal{P}_{k,\mathfrak{B}}(B;W)}} .
\end{equation}
For a bounded function $G$ we put
$\operatorname{osc}(G):=\sup G-\inf G$.

The \emph{localisation hypothesis for one clean sourced step} is the following assertion.
There are constants $C_{\mathrm{loc}}$, $c_{\mathrm{loc}}$, $\ell_0$, $D_0$, $C_{(1)}$,
$C_{\mathrm{OP}}$, $c_{\mathrm{OP}}$, depending only on $L$, $\mathfrak{p}$ and $N$, with the
following property. Let $G\in\mathfrak{G}_k(p_k;\ell,M,\varrho)$, and let the step $k$ be
$p$-clean with box radius $D\ge \ell+D_0$. Then:
\begin{enumerate}
\item[(i)] \emph{(box decoupling)} for every configuration $W$ in the window of configurations,
one has
\begin{equation}\label{8.2:eq-clean-step-localisation-2}
\bigl|\,\mathbb{E}_k[G\,|\,W]-\langle G\circ\mathfrak{v}_k\rangle_{k,\mathfrak{B}}(W)\bigr|
\le C_{\mathrm{loc}}\,\operatorname{osc}(G)\,e^{-c_{\mathrm{loc}}(D-\ell)} ;
\end{equation}
\item[(ii)] \emph{(re-localisation)} put $\ell':=\lceil \ell/L\rceil+\ell_0$. There are a number
$\varrho'\ge\min\{\varrho_{\mathcal{P}},\varrho\}$ and an observable
$G'\in\mathfrak{G}_{k+1}(p_{k+1};\ell',C_{(1)}M,\varrho')$
such that, with the supremum taken over configurations $W$ in the window of configurations,
\begin{equation}\label{8.2:eq-clean-step-localisation-3}
\sup\,\bigl|\,\mathbb{E}_k[G\,|\,\cdot\,]-G'\bigr|
\le C_{\mathrm{loc}}\,\operatorname{osc}(G)\,e^{-c_{\mathrm{loc}}D} .
\end{equation}
\item[(iii)] \emph{(covariance and background locality)} for $b,b'\in\mathfrak{B}$,
\begin{equation}\label{8.2:eq-clean-step-localisation-4}
\bigl|\bigl(C^{(k)}(W)\bigr)_{bb'}\bigr|\le C_{\mathrm{OP}}\,e^{-c_{\mathrm{OP}}|b-b'|} .
\end{equation}
Moreover, the following quantities are analytic functions of $W$ on the complex window,
and their variation in the single variable $W(b'')$ obeys the Cauchy bound
$C_{\mathrm{OP}}\,e^{-c_{\mathrm{OP}}\operatorname{dist}(b'',\mathfrak{B})}$:
\begin{itemize}
\item[--] the restriction of $C^{(k)}(W)$ to $\mathfrak{B}$;
\item[--] the restriction of $V^{(k)}(W)$ to $\mathfrak{B}$;
\item[--] the part of $U_{k+1}(W)$ lying under $\mathfrak{B}$.
\end{itemize}
\end{enumerate}
\end{definition}

Items (i) and (ii) are the content of the lemma on localisation of the sourced step; item
(iii) is that of the operator theorem at free current. Items (i) and (ii) express the localised
analytic representation of the step integral with one bounded, local, analytic, gauge-invariant
source term; without the source term this representation is \cite[Theorem~3]{B-CMP109}, the
representation \cite[(1.7)]{B-CMP109} with the bounds \cite[(1.18)]{B-CMP109}, obtained by the
cluster expansion of \cite[\S\S3--5]{B-CMP109}. Item (iii) uses \cite[Proposition~1.2]{B-CMP95}
and \cite{B-CMP99b}. Item (iii) alone does not imply item (i): item (iii) controls only the decay
of the Gaussian covariance, whereas the characteristic functions and the exponential of the
step's exponent in the density of Lemma~\ref{8.2:lem-conditional-step} also couple the box to
its complement, and decoupling them is the content of the cluster expansion.

\begin{lemma}[Positivity of the fluctuation form]\label{8.2:lem-fluctuation-positivity}
Let $C$ be the linear map of Notation~\ref{8.2:not-small-field-step}, through which the fluctuation
field $B$ enters the step as $CB$, and let $\Delta^{(k)}$ be the symmetric operator of the level-$k$
fluctuation quadratic form. Write
\begin{equation*}
\mathcal{Q}=\mathcal{Q}(U,0):=C^{*}\Delta^{(k)}C
\end{equation*}
for the fluctuation form at the background $U$, with its second argument equal to $0$. It is
regarded as a symmetric quadratic form on the constraint surface, the finite-dimensional real space
of fluctuation fields $B$ over which the fluctuation integral of the step is taken. This space
carries the Euclidean scalar product $\langle\cdot,\cdot\rangle$.

Assume, as the antecedent of this lemma, the lower bound for the fluctuation form at the background
$U$:
\begin{equation}\label{8.2:eq-fluctuation-positivity-1}
\mathcal{Q}(U,0)\ \ge\ \frac{\gamma_0}{2}\qquad\text{on the constraint surface,}
\end{equation}
where $\gamma_0>0$ is the positivity constant of that bound.
Then $\mathcal{Q}$ is invertible on the constraint surface, and $C^{(k)}:=\mathcal{Q}^{-1}$ satisfies
\begin{equation}\label{8.2:eq-fluctuation-positivity-2}
C^{(k)}\ \le\ \frac{2}{\gamma_0}
\end{equation}
as an operator. Let $\|C\|$ be the operator norm of $C$ from the constraint surface to the space of
the box values $(CB)_b$, both spaces carrying their Euclidean norms. Let
$v_k:=\sup_b\operatorname{Var}\bigl((CB)_b\bigr)$ be the box variance. The variance is taken under
the centred Gaussian measure on the constraint surface with covariance $C^{(k)}$, and the supremum
is taken over the boxes $b$. Then
\begin{equation}\label{8.2:eq-fluctuation-positivity-3}
v_k\ \le\ \frac{2\|C\|^{2}}{\gamma_0}.
\end{equation}
\end{lemma}

\begin{proof}
By \eqref{8.2:eq-fluctuation-positivity-1}, $\langle x,\mathcal{Q}x\rangle\ge(\gamma_0/2)|x|^{2}$ for
every $x$ in the constraint surface. Since $\mathcal{Q}$ is symmetric on this finite-dimensional space,
it has an orthonormal eigenbasis. Applying the inequality to each eigenvector shows that every
eigenvalue is at least $\gamma_0/2>0$. Hence $\mathcal{Q}$ is invertible. Moreover
$C^{(k)}=\mathcal{Q}^{-1}$ is diagonal in the same basis, with eigenvalues in $(0,2/\gamma_0]$. This is
\eqref{8.2:eq-fluctuation-positivity-2}, and it gives $\|C^{(k)}\|\le 2/\gamma_0$.

The covariance of $CB$ is $C\,C^{(k)}C^{*}$. Let $e_b$ be the unit vector of the box $b$. Then,
using $\|C^{*}e_b\|\le\|C^{*}\|=\|C\|$,
\begin{equation*}
\operatorname{Var}\bigl((CB)_b\bigr)=\langle C^{*}e_b,\,C^{(k)}C^{*}e_b\rangle
\le\|C^{(k)}\|\,\|C^{*}e_b\|^{2}
\le\frac{2}{\gamma_0}\,\|C\|^{2}.
\end{equation*}
Taking the supremum over $b$ gives \eqref{8.2:eq-fluctuation-positivity-3}.
\end{proof}

\begin{lemma}[Contraction of curvature by the background fields]\label{8.2:lem-curvature-contraction}
For a level $j$ and a plaquette $\bar p$ of $\mathbb{T}^{(j)}$, write $\bar p_{j+1}$ for a plaquette of
$\mathbb{T}^{(j+1)}$, fixed once for all for each $\bar p$, at minimal Euclidean distance from the
barycentre $x(\bar p)$ of $\bar p$, and
for $j'>j+1$ put $\bar p_{j'}:=(\bar p_{j'-1})_{j'}$; thus $(\bar p_{j'})_{j''}=\bar p_{j''}$ whenever
$j<j'<j''$. Fix a
level $k_0$ and a plaquette $p_{k_0}$ of $\mathbb{T}^{(k_0)}$, and for $j>k_0$ put $p_j:=(p_{k_0})_j$;
throughout, $k$ denotes a level with $k\ge k_0$, so that $p_k$ and $p_{k+1}=(p_k)_{k+1}$
are defined. Let $\ell\ge 1$, and for a plaquette $\bar p$ of any level let
$\mathfrak{B}_\ell(\bar p)$ be as in Definition~\ref{8.2:def-local-observables}. For a plaquette $\bar p$
of any level, $\mathcal{N}(\bar p)$ denotes a fixed finite neighbourhood of $\bar p$ in lattice units of
its level. For a
configuration $W$ on $\mathbb{T}^{(k+1)}$ in the window at level $k+1$ of Ba{\l}aban's small-field step,
$V^{(k)}(W)$ is the background configuration on $\mathbb{T}^{(k)}$ of
Notation~\ref{8.2:not-small-field-step}. For an integer $i\ge 1$ we put
\begin{equation}\label{8.2:eq-curvature-contraction-1}
\mathcal{V}^{(i)}_{k_0}:=V^{(k_0)}\circ V^{(k_0+1)}\circ\cdots\circ V^{(k_0+i-1)},
\end{equation}
a map from configurations on $\mathbb{T}^{(k_0+i)}$ to configurations on $\mathbb{T}^{(k_0)}$, defined on
the window at level $k_0+i$ by hypothesis (III) below.
Let $\mathfrak{g}$ be the Lie algebra of $\mathrm{SU}(N)$. For a configuration $U$ on a lattice
$\mathbb{T}^{(j)}$, a $\mathfrak{g}$-valued $1$-form $Z$ on its bonds and a plaquette $\bar p$ of
$\mathbb{T}^{(j)}$ we put
\begin{equation*}
(\operatorname{curl}_UZ)(\bar p):=\frac{\mathrm{d}}{\mathrm{d}t}\,(e^{tZ}U)(\partial \bar p)\Big|_{t=0},
\end{equation*}
where $e^{tZ}U$ is the configuration $b\mapsto e^{tZ_b}U(b)$; the same formula defines
$\operatorname{curl}_UZ$ for complex configurations $U$ and $1$-forms $Z$ with values in the
complexification $\mathfrak{g}_{\mathbb{C}}$ of $\mathfrak{g}$. Assume the following four statements.
\begin{itemize}
\item[(I)] The following consequence of the bounds \cite[Theorem~1, (8)--(10)]{B-CMP102b}, which hold
for every $m$ and with constants depending on $d$ and $L$ only (in particular they impose no cap in the
sense of Definition~\ref{2.3:def-cap}): there is a constant $B_3$,
depending on $d$ and $L$ only, such that for every level $j$, every plaquette $\bar p$ of
$\mathbb{T}^{(j)}$, every integer $n\ge1$, every $V$ in the window at level $j+n$ and every plaquette
$p'\subset\mathfrak{B}_\ell(\bar p)$ the following holds. The Wilson action $A$ of
Definition~\ref{1.1:def-wilson-action} attains its minimum over the configurations on the bare lattice
$\mathbb{T}^{(0)}$ whose $(j+n)$-fold block average under the block-averaging map of
Definition~\ref{1.5:def-block-average} is $V$ at a configuration, the minimiser of the variational problem
posed by $V$, which is unique up to a gauge transformation and whose block averages at every level are
determined up to a gauge transformation of that level; let $V_{j,n}$, the level-$j$ image of that
minimiser, be its $j$-fold block average. Then
\begin{equation}\label{8.2:eq-curvature-contraction-5}
\bigl|V_{j,n}(\partial p')-1\bigr|
\le B_3L^{-2n}\max_{p''\in\mathcal{N}(\bar p_{j+n})}\bigl|V(\partial p'')-1\bigr| .
\end{equation}
\item[(II)] The operator theorem at free current, with its constants $C_{\mathrm{op}}$ and
$c_{\mathrm{op}}$, which are the constants entering clause (iii) of the localisation hypothesis of
Definition~\ref{8.2:def-clean-step-localisation}.
\item[(III)] The complex-space form of the background contraction, with one constant for all numbers of
steps: the maps $W\mapsto V^{(k)}(W)$ and
$V\mapsto\mathcal{V}^{(i)}_{k_0}(V)$, with their values restricted to the bonds of the box of radius $D$
about $p_k$ in $\mathbb{T}^{(k)}$ of Definition~\ref{8.2:def-clean-step-localisation}, resp.\ of the box
of radius $D$ about $p_{k_0}$ in $\mathbb{T}^{(k_0)}$ defined as there with $k_0$ in place of $k$,
are holomorphic on the complex extension of the window
(in what follows, the complex window); for every $i\ge1$ the successive images of a configuration in
the complex window at level $k_0+i$ under the factors of \eqref{8.2:eq-curvature-contraction-1} lie in
the complex windows at the levels $k_0+i-1,\dots,k_0+1$, and those of a configuration in the window at
level $k_0+i$ lie in the windows at the same levels, so that \eqref{8.2:eq-curvature-contraction-1}
is defined there; the bound \eqref{8.2:eq-curvature-contraction-5} holds for $V$ in the complex window
with the constant $B_3$ of \textup{(I)}, its left-hand side being read as $V^{(k)}(W)(\partial p')-1$
when $j=k$, $\bar p=p_k$, $n=1$, $V=W$, and as $\mathcal{V}^{(i)}_{k_0}(V)(\partial p')-1$ when $j=k_0$,
$\bar p=p_{k_0}$, $n=i$.
\item[(IV)] The first-order form of the background contraction, with one constant for all numbers of
steps: there is one constant $B_{\mathrm{lin}}$ such that for every $k\ge k_0$, every
$i\ge1$, every $W$ in the complex window at level $k+1$, every $V$ in the complex window at level
$k_0+i$, every $\mathfrak{g}_{\mathbb{C}}$-valued $1$-form $Z$ on bonds near $p_{k+1}$ (resp.\
$p_{k_0+i}$) and every $p'\subset\mathfrak{B}_\ell(p_k)$ (resp.\ $p'\subset\mathfrak{B}_\ell(p_{k_0})$)
\begin{align}
\bigl|\operatorname{curl}_{V^{(k)}(W)}\bigl(\mathrm{d}V^{(k)}(W)Z\bigr)(p')\bigr|
&\le B_{\mathrm{lin}}L^{-2}\max_{p''}\bigl|(\operatorname{curl}_WZ)(p'')\bigr| ,
\label{8.2:eq-curvature-contraction-7}\\
\bigl|\operatorname{curl}_{\mathcal{V}^{(i)}_{k_0}(V)}\bigl(\mathrm{d}\mathcal{V}^{(i)}_{k_0}(V)Z\bigr)(p')\bigr|
&\le B_{\mathrm{lin}}L^{-2i}\max_{p''}\bigl|(\operatorname{curl}_VZ)(p'')\bigr| ,
\label{8.2:eq-curvature-contraction-8}
\end{align}
the maxima being taken over the plaquettes $p''$ of $\mathbb{T}^{(k+1)}$, resp.\
$\mathbb{T}^{(k_0+i)}$, and $\mathrm{d}V^{(k)}(W)$, $\mathrm{d}\mathcal{V}^{(i)}_{k_0}(V)$ denoting the
differentials of the two maps at $W$ and at $V$.
\end{itemize}
Then the following hold.
\begin{itemize}
\item[(a)] For every $W$ in the window at level $k+1$ and every plaquette $p'\subset\mathfrak{B}_\ell(p_k)$
\begin{equation}\label{8.2:eq-curvature-contraction-2}
\bigl|V^{(k)}(W)(\partial p')-1\bigr|
\le B_3L^{-2}\max_{p''\in\mathcal{N}(p_{k+1})}\bigl|W(\partial p'')-1\bigr| .
\end{equation}
For every $i\ge1$ and every $V$ in the window at level $k_0+i$,
$\mathcal{V}^{(i)}_{k_0}(V)$ coincides, up to a gauge transformation on $\mathbb{T}^{(k_0)}$, with the
level-$k_0$ image $V_{k_0,i}$ of the minimiser of the variational
problem posed by $V$, in the sense of \textup{(I)}, and for every plaquette
$p'\subset\mathfrak{B}_\ell(p_{k_0})$
\begin{equation}\label{8.2:eq-curvature-contraction-3}
\bigl|\mathcal{V}^{(i)}_{k_0}(V)(\partial p')-1\bigr|
\le B_3L^{-2i}\max_{p''\in\mathcal{N}(p_{k_0+i})}\bigl|V(\partial p'')-1\bigr| ,
\end{equation}
with the constant $B_3$ of \textup{(I)}, for every $i$ and every $m$.
\item[(b)] For every $W$ in the window at level $k+1$, every plaquette $p'\subset\mathfrak{B}_\ell(p_k)$
and every multi-index $\eta$, the derivative of order $\eta$ at $W$ of the function
$W'\mapsto V^{(k)}(W')(\partial p')$ with respect to the plaquette variables near
$p_{k+1}$ is bounded by $c_\eta B_3L^{-2}$, where $c_\eta$ depends only on $\eta$, on the
radius $\delta$ of a polydisc about $W$ contained in the complex window and on the supremum over $W'$ in
that polydisc of $\max_{p''\in\mathcal{N}(p_{k+1})}|W'(\partial p'')-1|$. For every $i\ge1$, every $V$
in the window at level $k_0+i$ and every $p'\subset\mathfrak{B}_\ell(p_{k_0})$
the same holds at $V$ for $V'\mapsto\mathcal{V}^{(i)}_{k_0}(V')(\partial p')$, with $c_\eta B_3L^{-2i}$,
$\mathcal{N}(p_{k_0+i})$ and a polydisc about $V$ in place of $c_\eta B_3L^{-2}$, $\mathcal{N}(p_{k+1})$
and a polydisc about $W$.
\item[(c)] For every level $k\ge k_0$ and every $W$ in the window at level $k+1$ let
$\mathcal{J}_k:=\mathrm{d}V^{(k)}(W)$ be the differential at the point $W$ of the map
$W\mapsto V^{(k)}(W)$; it carries $1$-forms on $\mathbb{T}^{(k+1)}$ (tangent vectors at $W$) to
$1$-forms on $\mathbb{T}^{(k)}$ (tangent vectors at the background $V^{(k)}(W)$). For $i\ge1$ and $V$ in
the window at level $k_0+i$ let $\mathcal{J}^{(i)}_{k_0}:=\mathrm{d}\mathcal{V}^{(i)}_{k_0}(V)$. The map
$\mathcal{J}_k$ takes a $1$-form $Z$ on $\mathbb{T}^{(k+1)}$ near $p_{k+1}$ to one on $\mathbb{T}^{(k)}$
near $p_k$, and there is one constant $B_{\mathrm{lin}}$ such that for every $W$ in the window at level
$k+1$, every $1$-form $Z$ on $\mathbb{T}^{(k+1)}$ near $p_{k+1}$ and every $p'\subset\mathfrak{B}_\ell(p_k)$
\begin{equation}\label{8.2:eq-curvature-contraction-4}
\bigl|\operatorname{curl}_{V^{(k)}(W)}(\mathcal{J}_kZ)(p')\bigr|
\le B_{\mathrm{lin}}L^{-2}\max_{p''}\bigl|(\operatorname{curl}_WZ)(p'')\bigr| ,
\end{equation}
the maximum over the plaquettes $p''$ of $\mathbb{T}^{(k+1)}$,
and, with the same $B_{\mathrm{lin}}$, for every $1$-form $Z$ on $\mathbb{T}^{(k_0+i)}$ near
$p_{k_0+i}$ and every $p'\subset\mathfrak{B}_\ell(p_{k_0})$
\begin{equation}\label{8.2:eq-curvature-contraction-6}
\bigl|\operatorname{curl}_{\mathcal{V}^{(i)}_{k_0}(V)}(\mathcal{J}^{(i)}_{k_0}Z)(p')\bigr|
\le B_{\mathrm{lin}}L^{-2i}\max_{p''}\bigl|(\operatorname{curl}_VZ)(p'')\bigr| ,
\end{equation}
the maximum over the plaquettes $p''$ of $\mathbb{T}^{(k_0+i)}$.
\item[(d)] If the localisation hypothesis of Definition~\ref{8.2:def-clean-step-localisation} holds for
the step $k$ with box radius $D\ge\ell+D_0$, then on the event that this step is $p_k$-clean the
backgrounds of
Ba{\l}aban's complete construction \cite[(2.10)--(2.16)]{B-CMP119}, whose determining sets avoid the
large-field regions, coincide near $p_k$ with the pure small-field background $V^{(k)}(W)$, $W$ the
level-$(k+1)$ configuration of the step, up to an error at most
$C_{\mathrm{op}}e^{-c_{\mathrm{op}}D}$.
\end{itemize}
\end{lemma}

\begin{proof}
\emph{(a), one step.} The configuration $V^{(k)}(W)$ is $W_{k,1}$ in the notation of (I), the level-$k$
image of the minimiser of the variational problem posed by $W$, and $(p_k)_{k+1}=p_{k+1}$, so
\eqref{8.2:eq-curvature-contraction-2} is the case $j=k$, $\bar p=p_k$, $n=1$
of \eqref{8.2:eq-curvature-contraction-5} in hypothesis (I), the consequence of
\cite[Theorem~1, (8)--(10)]{B-CMP102b} assumed there, whose constants depend on $d$ and $L$ only and
which holds for every $m$; hence $B_3$ does not depend on $k$, on $m$ or on $W$ (see also \cite{B-CMP119}).

\emph{(a), composition.} By \cite[(0.21), (2.3)]{B-CMP109}, minimal orbits are nested: the minimiser of
the variational problem at level $k+1$ is admissible for, and minimal in, the problem at level $k$
posed by its own $k$-fold average. We show by induction on $i\ge1$ that for every $V$ in the window at
level $k_0+i$
\begin{equation*}
\mathcal{V}^{(i)}_{k_0}(V)=V_{k_0,i}\quad\text{up to a gauge transformation on }\mathbb{T}^{(k_0)} .
\end{equation*}
For $i=1$ this is the identity $V^{(k_0)}(V)=V_{k_0,1}$ of the one-step case above. Let $i\ge2$ and put
$V':=V^{(k_0+i-1)}(V)$; it is the $(k_0+i-1)$-fold block average $V_{k_0+i-1,1}$ of a minimiser $U$ of
the problem posed by $V$, and it lies in the window at level $k_0+i-1$ by hypothesis (III). By the
nesting property, applied with $k+1=k_0+i$, $U$ is admissible for, and minimal in, the problem posed by
$V'$; by the uniqueness in (I) the minimiser of that problem is $U$ up to a gauge transformation, and
therefore, by the last clause of (I), its $k_0$-fold block average $V'_{k_0,i-1}$ coincides with the
$k_0$-fold block average $V_{k_0,i}$ of $U$ up to a gauge transformation on $\mathbb{T}^{(k_0)}$. Since
\eqref{8.2:eq-curvature-contraction-1} gives
$\mathcal{V}^{(i)}_{k_0}(V)=\mathcal{V}^{(i-1)}_{k_0}(V')$, the induction hypothesis applied to $V'$
gives $\mathcal{V}^{(i)}_{k_0}(V)=V'_{k_0,i-1}$ up to a gauge transformation, and the induction is
complete. The nesting property is thus used once for each of the $i-1$ passages between consecutive
factors of \eqref{8.2:eq-curvature-contraction-1}; this is the second assertion of (a). A gauge
transformation conjugates each plaquette value $\mathcal{V}^{(i)}_{k_0}(V)(\partial p')$ by an element of
$\mathrm{SU}(N)$, which leaves the left-hand side of \eqref{8.2:eq-curvature-contraction-3} unchanged.
Hypothesis (I) with
$j=k_0$, $\bar p=p_{k_0}$ and $n=i$ applies to this single minimiser, and $(p_{k_0})_{k_0+i}=p_{k_0+i}$;
since \eqref{8.2:eq-curvature-contraction-5} holds
for every $n$ and every $m$ with one constant depending on $d$ and $L$ only, it gives
\eqref{8.2:eq-curvature-contraction-3} with the same constant $B_3$ for every $i$.

\emph{(b).} By hypothesis (III) the maps in question are holomorphic on the complex
window and satisfy there \eqref{8.2:eq-curvature-contraction-5} in the readings stated in (III), that
is, \eqref{8.2:eq-curvature-contraction-2}, respectively \eqref{8.2:eq-curvature-contraction-3}, hold at
the points of the complex window. Recall Cauchy's inequality: if $f$ is holomorphic on a closed
polydisc of radius $\delta$ in each variable about a point $u$ and $|f|\le S$ there, then
\begin{equation*}
|\partial^\eta f(u)|\le \eta!\,S\,\delta^{-|\eta|}
\end{equation*}
for every multi-index $\eta$. We apply it to the matrix entries of
$V^{(k)}(W')(\partial p')-1$, as functions of the plaquette variables of $W'$ near $p_{k+1}$, on a
polydisc of radius $\delta$ about $W$ contained in the complex window; by
\eqref{8.2:eq-curvature-contraction-2} on that window,
$S$ is $B_3L^{-2}$ times the supremum over $W'$ in the polydisc of
$\max_{p''\in\mathcal{N}(p_{k+1})}|W'(\partial p'')-1|$. Every derivative at $W$ is therefore
$B_3L^{-2}$ times a factor depending only on its order, on $\delta$ and on that supremum, which is the
assertion. The same argument with \eqref{8.2:eq-curvature-contraction-3}, on a polydisc about $V$, gives
the bound $c_\eta B_3L^{-2i}$ for $\mathcal{V}^{(i)}_{k_0}$.

\emph{(c).} For a $1$-form $Z$ supported on bonds near $p_{k+1}$ we consider the restriction of
$\mathcal{J}_kZ$ to bonds near $p_k$; in this sense $\mathcal{J}_k$ takes $1$-forms near $p_{k+1}$ to
$1$-forms near $p_k$. Fix $W$ in the window at level $k+1$ and $Z$, and for complex $t$ of small modulus
put $W_t:=e^{tZ}W$, which lies in the complex window. At $t=0$ the derivative of
$V^{(k)}(W_t)(\partial p')$ is, by the definitions of the curl and of the differential $\mathcal{J}_k$
at $W$, $\operatorname{curl}_{V^{(k)}(W)}(\mathcal{J}_kZ)(p')$, and that of $W_t(\partial p'')$ is
$(\operatorname{curl}_WZ)(p'')$. Thus \eqref{8.2:eq-curvature-contraction-4} bounds the first-order
variation of the plaquette values of the background by the first-order variation of the plaquette
values of $W$. It is \eqref{8.2:eq-curvature-contraction-7} of hypothesis (IV), read at the real point
$W$ and for the $\mathfrak{g}$-valued form $Z$, and $B_{\mathrm{lin}}$ is the
constant of (IV). For the $i$-fold transport, the chain rule applied to
\eqref{8.2:eq-curvature-contraction-1} gives
\begin{equation*}
\mathcal{J}^{(i)}_{k_0}=\mathcal{J}_{k_0}\circ\mathcal{J}_{k_0+1}\circ\cdots\circ\mathcal{J}_{k_0+i-1},
\end{equation*}
each factor taken at the configuration obtained from $V$ by the factors of
\eqref{8.2:eq-curvature-contraction-1} to its right. Iterating \eqref{8.2:eq-curvature-contraction-4}
along this factorisation would multiply the constants of the factors, so the single constant in
\eqref{8.2:eq-curvature-contraction-6} is not obtained that way but from hypothesis (IV). With
$V_t:=e^{tZ}V$, the derivative at $t=0$ of
$\mathcal{V}^{(i)}_{k_0}(V_t)(\partial p')$ is
$\operatorname{curl}_{\mathcal{V}^{(i)}_{k_0}(V)}(\mathcal{J}^{(i)}_{k_0}Z)(p')$ and that of
$V_t(\partial p'')$ is $(\operatorname{curl}_VZ)(p'')$, as above; the bound
\eqref{8.2:eq-curvature-contraction-6}, with $L^{-2i}$ and the single constant $B_{\mathrm{lin}}$, is
\eqref{8.2:eq-curvature-contraction-8} of hypothesis (IV), which holds with one constant for all $i$,
read at the real point $V$.

\emph{(d).} This is clause (iii) of the localisation hypothesis of
Definition~\ref{8.2:def-clean-step-localisation}, applied to the step $k$ with box radius
$D\ge\ell+D_0$ on the event that it is $p_k$-clean;
its constants are those of the operator theorem at free current, hypothesis (II).
\end{proof}

\begin{definition}[{The flat quadratic form is the free field's;
\cite[(1.19), (1.23)--(1.24)]{B-CMP95}}]\label{8.2:def-flat-quadratic-form}
Let $0\le k\le K$. Here $d$, $\|\cdot\|$ and $\langle\cdot,\cdot\rangle$ denote, as in
\cite[\S1]{B-CMP95} and in the normalisation printed there, the lattice exterior derivative, the
norm on real-valued functions on plaquettes and the inner product on real-valued functions on bonds;
no confusion with the dimension $d=4$ arises.
We use the rescaled units of \cite[\S1]{B-CMP95}: lengths are multiplied by $L^{K-k}$, so that the
lattice of level $k$ of Definition~\ref{1.5:def-block-average} has unit spacing and the bare
lattice has spacing $L^{-k}$.
Let $Q_k$ be the linearised block-averaging map over $k$ steps (the map $Q$ of the glossary,
Notation~\ref{0.2:not-glossary}, taken over $k$ steps); it acts on real-valued fields
$\mathcal{A}$ on the bonds of the bare lattice (the configurations of the free lattice field of
Notation~\ref{0.2:not-glossary}) and takes values in real-valued fields $B$ on the
bonds of the lattice of level $k$, the unit lattice in these units. The operator $\Delta_k$ is
fixed by
\begin{equation}\label{8.2:eq-flat-quadratic-form-1}
\tfrac12\,\langle B,\Delta_k B\rangle
=\inf\Bigl\{\tfrac12\,\|d\mathcal{A}\|^2 \;:\; Q_k\mathcal{A}=B\Bigr\},
\end{equation}
and the quadratic form $B\mapsto\tfrac12\langle B,\Delta_kB\rangle$ so defined is the
\emph{flat quadratic form} of level~$k$.
The \emph{flat background} is the configuration in which every bond variable equals the
identity, and the \emph{flat orbit} is its gauge orbit.

Ba{\l}aban states in \cite{B-CMP95} and in \cite[(0.21)--(0.22)]{B-CMP109} the following, which is
used here as stated there. For a block field $V$ on the lattice of level $k$, let $U_k(V)$ be the
configuration on the bare lattice that Ba{\l}aban's step at level $k$ associates with $V$ (a
bare-lattice configuration, not the block field $U_k$ of the glossary), and let $A$ be the Wilson
action of Definition~\ref{1.1:def-wilson-action}. At flat background, the Hessian, at the flat
orbit and on gauge-invariant functionals, of the map $V\mapsto A(U_k(V))$ is the direct sum of
$N^2-1$ copies of the flat quadratic form of level $k$, one for each real colour component of the
Lie algebra of $\mathrm{SU}(N)$, the number $N^2-1$ being the real dimension of that Lie algebra.
This Hessian is the step-$k$ quadratic form for the unit-lattice fluctuation of the gauge field;
per colour component it is the form \eqref{8.2:eq-flat-quadratic-form-1}.
The paper \cite{B-CMP95} expresses this identification in the passage
\begin{quotation}
Lattice gauge theories may be looked at as perturbations of the theory of a vector field with a
Gaussian action
\end{quotation}
and, further on in the same paper, in the passage
\begin{quotation}
Main terms \ldots\ equal to the simplest quadratic action for Abelian vector fields
\end{quotation}
\end{definition}

The operator $\Delta_k$ and the variational formula \eqref{8.2:eq-flat-quadratic-form-1} are
those of \cite[(1.19), (1.23)--(1.24)]{B-CMP95}, used here as stated there.

In this chapter the identification is used only at the level, $k_0$ say, at which the lemma of
this chapter on the free lattice field under block averaging is applied. There the
flat-background quadratic coefficients of the step, summed over the colour components, are equated
with the constant that the free lattice field $\mathcal{A}$ carries at level $k_0$ in that lemma.
The Gaussian extraction does not use it.

\begin{definition}[The admissible range and the clean event]\label{8.2:def-clean-event}
We use the plaquette $p$, the step levels $k_0\le k<k_0+r$, the images $p_k$ of $p$ (with
$p_{k_0}=p$) and the boxes $\mathfrak{B}_D(p_k)$ of Notation~\ref{8.2:not-plaquette-boxes}.
Fix integers $s$, $r$ and a box radius $D$ in the \emph{admissible range}
\begin{equation}\label{8.2:eq-clean-event-1}
\begin{gathered}
s_*\le s-r,\qquad 1\le r,\qquad 2B^\sharp(r+1)\le x_{k_0},\qquad s\le K-1,\\
D\ge D_0+\ell_r ,
\end{gathered}
\end{equation}
where $\ell_r$ is the largest support radius of the functionals met in the analysis of the
level-$(k_0+r)$ expansion carried out in the remainder of this chapter; under the admissible
choices of the parameters it satisfies
\begin{equation}\label{8.2:eq-clean-event-2}
\ell_r\le \ell+r\ell_0\le \ell+C\log x .
\end{equation}
Here $D_0$ is the offset of the box radius in the localisation hypothesis for one clean sourced
step (Definition~\ref{8.2:def-clean-step-localisation}); $\ell$ is the support-radius parameter
of the class of functionals in that definition and $\ell_0$ the increase of the support radius in
one step; $C$ is a constant
fixed by the admissible choices; and $x$ is the inverse coupling in terms of which these choices
are made.
The \emph{clean event} $\mathcal{C}_p(D)$ is the set of histories $h$ of the level-$(k_0+r)$
expansion that satisfy the following three conditions.
\begin{itemize}
\item[(c1)] \emph{Current levels.} For every $j$ with $k_0\le j\le k_0+r$, no large-field object
of level $j$ lies within distance $D$, measured in units of the level-$j$ lattice, of $p_j$.
Here a large-field object of level $j$ is, in the large-field vocabulary of this chapter, any of
the following: a large-field region belonging
to one of the three classes $P'$, $Q'$, $R'$ of regions created at step $j$; an
$R'$-pre-integrated component, of either class; a term of the family $\mathsf{K}(I)$ of
large-field terms of level $j$, with $I$ ranging over the index set of that family; a large-field
term $\mathsf{R}^{\mathrm{lo}}$ of level $j$.
\item[(c2)] \emph{Coupling lag under $p$.} For every $j$ with $k_0\le j<k_0+r$, the
volume-weighted coupling lag $\mathfrak{S}'_j(h)$ of the history $h$ at step $j$ satisfies
\begin{equation*}
\mathfrak{S}'_j(h)\le \tau_*\,\#\mathfrak{B}_D(p_j),\qquad \tau_*:=1 .
\end{equation*}
The lag is thus required to be small in this volume-weighted sense, not to be absent. Any
other fixed value of $\tau_*$ may be used; it enters only the constant $C_{\mathcal{P}}$ of the
bound on the box part of the exponent of the step in Lemma~\ref{8.2:lem-box-exponent}.
\item[(c3)] \emph{No live region under $p$.} No large-field region, of any level, that is live
at step $k_0$ (that is, still carrying its large-field factor and not yet integrated by
Ba{\l}aban's operation $\mathbf{R}$) meets the fine volume under $\mathfrak{B}_D(p)$.
\end{itemize}
\end{definition}

Let $\mathbf{A}$ denote a large-field fluctuation variable, that is, a fluctuation variable of a
large-field region that has not yet been integrated by $\mathbf{R}$. The clean event has the
following consequences, proved below. On $\mathcal{C}_p(D)$, at
every step $j\in[k_0,k_0+r)$ and within $\mathfrak{B}_D(p_j)$, the following hold.
\begin{enumerate}
\item[(a)] The step has the un-modified form of \cite[(2.21)]{B-CMP119}.
\item[(b)] No un-integrated large-field fluctuation variable $\mathbf{A}$ is coupled into the
box, except through terms in the fluctuation variables $B_j$ of step $j$ coming from distance
at least $D$.
\item[(c)] The box part of the exponent of the step obeys the bound of
Lemma~\ref{8.2:lem-box-exponent} on the box part of the exponent, with the constant
$C_{\mathcal{P}}$ of that bound.
\end{enumerate}

\begin{proof}
(a) By (c1), no large-field object of level $j$ lies within distance $D$ of $p_j$; in
particular none of the large-field modifications of level $j$ acts within
$\mathfrak{B}_D(p_j)$. Within the box the step is therefore of the un-modified form
\cite[(2.21)]{B-CMP119}.

(b) Let $X$ be a large-field region that is live under the box $\mathfrak{B}_D(p_j)$ at a step
$j\in[k_0,k_0+r)$. If $j=k_0$, the box is $\mathfrak{B}_D(p)$, and such an $X$ is excluded by
(c3). If $j>k_0$, then either $X$ was created at a level in $[k_0,j]$ within the box, which is
excluded by (c1), or $X$ was already live at step $k_0$, which is excluded by (c3).
Hence no live region lies under the box at any
step of the range, and an un-integrated variable $\mathbf{A}$ can enter the box only through
terms in $B_j$ originating at distance at least $D$.

(c) The contribution of the coupling lag to the box part of the exponent is controlled by
(c2). The remaining contributions satisfy the bound of Lemma~\ref{8.2:lem-box-exponent}
deterministically, that is, on every history in $\mathcal{C}_p(D)$, without further condition.
\end{proof}

The definition does not exclude large-field regions under $p$ that have already been integrated
by $\mathbf{R}$ and are therefore no longer live. At deep levels such regions
are present under $p$ with probability close to one, so that no event excluding them at every
step of an initial segment of the expansion has probability bounded away from zero uniformly in
the depth. The terms $R$ of \cite[(2.31)]{B-CMP119} attached to them are deterministic by that
display, and their contribution to the coupling lag is the one controlled by (c2).

\begin{definition}[Bounded lifetime of large-field regions]\label{8.2:def-large-field-lifetime}
In this definition $r_B>0$ is the exponent of $\log\check{x}_j$ in the definition
\eqref{8.2:eq-large-field-lifetime-2} of $\check{R}_j$ below, $C_{\mathfrak{l}}$ is a constant
not depending on the level $j$, and $\check{x}_j>1$ is the inverse coupling at level $j$
in Ba{\l}aban's large-field construction
\cite{B-CMP122a}. The exponent $p_0$ of the degree
function $p_0(\cdot)$ satisfies
\begin{equation}\label{8.2:eq-large-field-lifetime-1}
p_0\ \ge\ 5r_B .
\end{equation}
For a level $j$ let $\check{R}_j$ be the least power of $L$ that is at least
$(\log\check{x}_j)^{r_B}$:
\begin{equation}\label{8.2:eq-large-field-lifetime-2}
\check{R}_j=\min\bigl\{L^n:\ n\ \text{a non-negative integer},\
L^n\ge(\log\check{x}_j)^{r_B}\bigr\},
\end{equation}
and put
\begin{equation}\label{8.2:eq-large-field-lifetime-3}
\mathfrak{l}_j=\check{R}_j+C_{\mathfrak{l}} .
\end{equation}
A large-field region created at level $j$ has \emph{age} $k-j$ at level $k$. It is
\emph{live} at level $k$ if it still carries its large-field factor at that level.

We say that large-field regions have \emph{bounded lifetime} when the following holds: a
large-field region created at level $j$ is integrated by Ba{\l}aban's large-field operation
$\mathbf{R}$ (Definition~\ref{1.5:def-densities}) within $\mathfrak{l}_j$ further steps; a
large-field region created at level $j$ that is still live at an age exceeding
$\mathfrak{l}_j$ contains a large-field region created at a later level.

The numbers $\check{R}_j$ and the lifetime $\mathfrak{l}_j$ follow Ba{\l}aban's definitions
in \cite{B-CMP122a}. Of this definition only the
inequality $\mathfrak{l}_j\le\check{R}_j+C_{\mathfrak{l}}$ (here an equality) and the
following arithmetic fact are used later: for every fixed $L>1$,
\begin{equation}\label{8.2:eq-large-field-lifetime-4}
\frac{(\log x)^{p_0}}{(\log L)\,(\log x)^{r_B}}\ \longrightarrow\ +\infty
\qquad (x\to\infty).
\end{equation}
\end{definition}

\begin{proof}
We prove \eqref{8.2:eq-large-field-lifetime-4}.

Let $L>1$ be fixed, so that $\log L>0$, and let $x>e$, so that $\log x>1$. Then
\begin{equation*}
\frac{(\log x)^{p_0}}{(\log L)\,(\log x)^{r_B}}
=\frac{(\log x)^{p_0-r_B}}{\log L}.
\end{equation*}
By \eqref{8.2:eq-large-field-lifetime-1} and $r_B>0$, $p_0-r_B\ge 4r_B>0$; in particular
$p_0>r_B$. Since $\log x>1$, raising
$\log x$ to a larger exponent does not decrease it, so
\begin{equation*}
\frac{(\log x)^{p_0-r_B}}{\log L}\ \ge\ \frac{(\log x)^{4r_B}}{\log L}.
\end{equation*}
Because $r_B>0$, the right-hand side tends to $+\infty$ as $x\to\infty$, with $L$ held
fixed. This proves \eqref{8.2:eq-large-field-lifetime-4}.
\end{proof}

\begin{lemma}[Comparison of couplings along a segment]\label{8.2:lem-coupling-comparison}
Let $k_0\ge 0$ and $r\ge 1$ be integers with $k_0+r\le K$ (in the application, the level $k_0$ and
the integer $r$ of Definition~\ref{8.2:def-clean-event}), and let $x_k=g_k^{-2}$ be the running
inverse couplings of Definition~\ref{1.2:def-couplings}. Let $B^\sharp$ and $c_5=2c_2$ be the flow
constants of Definition~\ref{1.2:def-flow-constants}. Assume the part
\begin{equation}\label{8.2:eq-coupling-comparison-1}
2B^\sharp(r+1)\le x_{k_0}
\end{equation}
of the range condition of Definition~\ref{8.2:def-clean-event}. Assume that the lower and upper
flow bounds for the running inverse couplings, and the continuity statement accompanying them,
hold at the levels $k_0,\dots,k_0+r$; in particular, for every integer $i$ with $0\le i\le r$:
\begin{itemize}
\item[(a)] the lower flow bound $x_{k_0+i}\ge x_{k_0}-B^\sharp i$;
\item[(b)] the upper flow bound $x_{k_0+i}\le x_{k_0}+c_5$.
\end{itemize}
Then for every integer $i$ with $0\le i\le r$,
\begin{gather}
\tfrac12\,x_{k_0}\le x_{k_0}-B^\sharp i\le x_{k_0+i}\le x_{k_0}+c_5,
\label{8.2:eq-coupling-comparison-2}\\
g_{k_0+i}^2\le 2g_{k_0}^2,\qquad
\Bigl|\frac{g_{k_0+i}^2}{g_{k_0}^2}-1\Bigr|\le (2B^\sharp i+c_5)\,g_{k_0}^2.
\label{8.2:eq-coupling-comparison-3}
\end{gather}
Moreover there is a constant $C\ge 2$ such that $x_{k_0}\ge C$ implies, for $0\le i\le r$,
\begin{equation*}
\mathfrak{b}_{k_0+i}\le 2\,\mathfrak{b}_{k_0},\qquad p_0(g_{k_0+i})\le 2\,p_0(g_{k_0}),
\end{equation*}
where $\mathfrak{b}_k$ is the amplitude attached to level $k$ in this chapter and
$p_0(\cdot)$ is the degree function of Definition~\ref{1.5:def-domains}.
\end{lemma}

\begin{proof}
Fix $0\le i\le r$. Since $i\le r<r+1$, \eqref{8.2:eq-coupling-comparison-1} gives
$B^\sharp i\le B^\sharp(r+1)\le x_{k_0}/2$, which is the first inequality of
\eqref{8.2:eq-coupling-comparison-2}; the second and third inequalities are hypotheses (a)
and (b).

Since $g_k^2=1/x_k$ and $x_{k_0+i}\ge x_{k_0}/2>0$, we get
\begin{equation*}
g_{k_0+i}^2=\frac{1}{x_{k_0+i}}\le\frac{2}{x_{k_0}}=2g_{k_0}^2 .
\end{equation*}
For the relative bound write
\begin{equation*}
\frac{g_{k_0+i}^2}{g_{k_0}^2}-1=\frac{x_{k_0}-x_{k_0+i}}{x_{k_0+i}} .
\end{equation*}
If $x_{k_0+i}\ge x_{k_0}$, then by (b) the modulus is at most
$c_5/x_{k_0+i}\le c_5/x_{k_0}=c_5\,g_{k_0}^2$. If $x_{k_0+i}<x_{k_0}$, then by (a) the numerator is
at most $B^\sharp i$ and by \eqref{8.2:eq-coupling-comparison-2} the denominator is at least
$x_{k_0}/2$, so the modulus is at most $2B^\sharp i\,g_{k_0}^2$. In both cases the modulus is at
most $(2B^\sharp i+c_5)g_{k_0}^2$, which is \eqref{8.2:eq-coupling-comparison-3}.

For the degree function, $p_0(g_k)=A_0(\log x_k)^{p_0}$ with $A_0>0$. Choose $C\ge 2$ with
$C\ge c_5$ and $(1+\log 2/\log C)^{p_0}\le 2$; such a $C$ exists since the left side tends to $1$
as $C\to\infty$, and it depends only on $c_5$ and $p_0$. If $x_{k_0}\ge C$, then
$x_{k_0+i}\ge x_{k_0}/2\ge 1$ by \eqref{8.2:eq-coupling-comparison-2}, so $\log x_{k_0+i}\ge0$, and
by (b) and $c_5\le x_{k_0}$,
\begin{equation*}
0\le\log x_{k_0+i}\le \log(2x_{k_0})
=(\log x_{k_0})\Bigl(1+\frac{\log 2}{\log x_{k_0}}\Bigr).
\end{equation*}
Raising to the power $p_0$ and using $\log x_{k_0}\ge\log C$ gives
$(\log x_{k_0+i})^{p_0}\le 2(\log x_{k_0})^{p_0}$; multiplying by $A_0$ gives
$p_0(g_{k_0+i})\le 2p_0(g_{k_0})$.

Under the same threshold $x_{k_0}\ge C$, the bound $\mathfrak{b}_{k_0+i}\le 2\,\mathfrak{b}_{k_0}$
is a consequence of \eqref{8.2:eq-coupling-comparison-2} and
\eqref{8.2:eq-coupling-comparison-3}.
\end{proof}

In the sums over ultraviolet levels that occur in this chapter, only the lower flow bound,
hypothesis~(a) of Lemma~\ref{8.2:lem-coupling-comparison}, is used.

\begin{lemma}[Laplace extraction with parity on a nearly symmetric window]\label{8.2:lem-laplace-parity}
Let $D\ge 1$ be an integer and $\Sigma$ a positive-definite covariance matrix on $\mathbb{R}^{D}$.
Let $\mu_{\Sigma}$ be the centred Gaussian measure on $\mathbb{R}^{D}$ with covariance $\Sigma$, and
put $v:=\max_{1\le i\le D}\Sigma_{ii}$. Let $0\le\eta\le\frac14$. A measurable set
$\mathcal{W}\subset\mathbb{R}^{D}$ is an \emph{$\eta$-symmetric window} if
\begin{equation*}
\mu_{\Sigma}(\mathcal{W}^{c})\le\eta
\quad\text{and}\quad
\mu_{\Sigma}\bigl(\mathcal{W}\setminus(-\mathcal{W})\bigr)\le\eta ,
\qquad \mathcal{W}^{c}:=\mathbb{R}^{D}\setminus\mathcal{W}.
\end{equation*}
Fix such a window and write
$\mathcal{W}_{\mathrm{s}}:=\mathcal{W}\cap(-\mathcal{W})$ and
$\mathcal{W}_{\mathrm{a}}:=\mathcal{W}\setminus(-\mathcal{W})$.

For measurable functions $\Phi,\Psi$ on $\mathcal{W}$ put
$\mathbb{E}_0[\Phi]:=\mu_{\Sigma}(\mathcal{W})^{-1}\int_{\mathcal{W}}\Phi\,d\mu_{\Sigma}$ whenever the
integral converges absolutely,
$\operatorname{Cov}_0(\Phi,\Psi):=\mathbb{E}_0[\Phi\Psi]-\mathbb{E}_0[\Phi]\,\mathbb{E}_0[\Psi]$
and $\widehat{\Phi}:=\Phi-\mathbb{E}_0[\Phi]$. For a measurable $\Phi$ on $\mathbb{R}^{D}$, or on
$\mathcal{W}$ extended by zero to $\mathbb{R}^{D}$, put
$\|\Phi\|_{2}:=(\int\Phi^{2}\,d\mu_{\Sigma})^{1/2}$. The \emph{even and odd parts} $\Phi^{\pm}$ of a
function $\Phi$ on $\mathcal{W}$ are the functions on $\mathcal{W}$ defined by
$\Phi^{\pm}(y):=\frac12\bigl(\Phi(y)\pm\Phi(-y)\bigr)$ for $y\in\mathcal{W}_{\mathrm{s}}$, and by
$\Phi^{+}:=\Phi$, $\Phi^{-}:=0$ on $\mathcal{W}_{\mathrm{a}}$. Thus $\Phi=\Phi^{+}+\Phi^{-}$ on
$\mathcal{W}$, $\Phi^{+}$ is even and $\Phi^{-}$ is odd on $\mathcal{W}_{\mathrm{s}}$, and
$\mu_{\Sigma}(\mathcal{W}_{\mathrm{a}})\le\eta$.
For a function on $\mathbb{R}^{D}$ the notions $\mathbb{E}_0$, $\operatorname{Cov}_0$,
$\widehat{\Phi}$ and the even and odd parts refer to its restriction to $\mathcal{W}$.

Let $P:\mathcal{W}\to\mathbb{R}$ be measurable with
$\sup_{\mathcal{W}}|P|\le\delta\le\frac14$, and let $\delta^{+}$ be a number with
$\sup_{\mathcal{W}}|P^{+}|\le\delta^{+}\le\delta$.
The \emph{tilted expectation} is
$\langle\Phi\rangle_{P}:=\mathbb{E}_0[\Phi e^{P}]/\mathbb{E}_0[e^{P}]$.

Then for every measurable $G:\mathbb{R}^{D}\to\mathbb{R}$ with $\|G\|_{2}<\infty$, with even and odd
parts $G^{\pm}$, the following hold.
\begin{enumerate}
\item[(a)]
\begin{equation*}
\langle G\rangle_{P}=\mathbb{E}_0G+\operatorname{Cov}_0(G,P)+\epsilon_{2}(G),
\qquad
|\epsilon_{2}(G)|\le 3\delta^{2}\,\mathbb{E}_0|\widehat{G}|.
\end{equation*}
\item[(b)] $|\mathbb{E}_0G^{-}|\le 4\|G^{-}\|_{2}\,\eta^{1/2}$;
\begin{equation*}
\operatorname{Cov}_0(G,P)=\operatorname{Cov}_0(G^{+},P^{+})+\mathbb{E}_0[G^{-}P^{-}]
+\epsilon_{\mathrm{par}},
\qquad
|\epsilon_{\mathrm{par}}|\le 16\,\delta\bigl(\|G^{+}\|_{2}+\|G^{-}\|_{2}\bigr)\eta^{1/2};
\end{equation*}
and
\begin{equation*}
|\operatorname{Cov}_0(G^{+},P^{+})|\le 2\delta^{+}\,\mathbb{E}_0|G^{+}-\mathbb{E}_0G^{+}|,
\qquad
|\mathbb{E}_0[G^{-}P^{-}]|\le\delta\,\mathbb{E}_0|G^{-}|.
\end{equation*}
\item[(c)] For every measurable $\Phi:\mathbb{R}^{D}\to\mathbb{R}$ with $\|\Phi\|_{2}<\infty$,
$|\mathbb{E}_0\Phi-\int\Phi\,d\mu_{\Sigma}|\le 4\|\Phi\|_{2}\,\eta^{1/2}$ and
$\mathbb{E}_0|\Phi|\le\frac43\int|\Phi|\,d\mu_{\Sigma}$.
\item[(d)] (second order)
\begin{equation*}
\begin{split}
\langle G\rangle_{P}&=\mathbb{E}_0G+\operatorname{Cov}_0(G,P)
+\tfrac12\operatorname{Cov}_0(G,P^{2})-\operatorname{Cov}_0(G,P)\,\mathbb{E}_0P+\epsilon_{3}(G),\\
|\epsilon_{3}(G)|&\le 4\delta^{3}\,\mathbb{E}_0|\widehat{G}|.
\end{split}
\end{equation*}
\end{enumerate}
\end{lemma}

\begin{proof}
\emph{Preliminaries.} Since $\mu_{\Sigma}(\mathcal{W}^{c})\le\eta\le\frac14$, we have
$\mu_{\Sigma}(\mathcal{W})\ge\frac34$. For $\Phi$ with $\|\Phi\|_{2}<\infty$ the Cauchy--Schwarz
inequality gives $\mathbb{E}_0|\Phi|\le(\mathbb{E}_0[\Phi^{2}])^{1/2}\le(4/3)^{1/2}\|\Phi\|_{2}$, so all
expectations below are finite; the products with $P$, $P^{\pm}$, $e^{P}$ are covered as well, these
functions being bounded on $\mathcal{W}$. The measure $\mu_{\Sigma}$ is invariant under $y\mapsto-y$,
and $\mathcal{W}_{\mathrm{s}}$ is invariant under the same map. Hence
$\int_{\mathcal{W}_{\mathrm{s}}}G(-y)^{2}\,d\mu_{\Sigma}(y)=\int_{\mathcal{W}_{\mathrm{s}}}G^{2}\,d\mu_{\Sigma}$.
Pointwise on $\mathcal{W}_{\mathrm{s}}$ we have $(G^{\pm})^{2}\le\frac12(G(y)^{2}+G(-y)^{2})$, while
on $\mathcal{W}_{\mathrm{a}}$ we have $(G^{+})^{2}=G^{2}$ and $G^{-}=0$. Therefore
$\|G^{\pm}\|_{2}\le\|G\|_{2}<\infty$. For any two functions we have
$\operatorname{Cov}_0(\Phi,\Psi)=\mathbb{E}_0[\widehat{\Phi}\,\Psi]=\mathbb{E}_0[\widehat{\Phi}\,\widehat{\Psi}]$,
because $\mathbb{E}_0\widehat{\Phi}=0$. In particular covariances with a constant vanish.

\emph{Proof of} (c). Since $\int_{\mathcal{W}}\Phi=\int\Phi-\int_{\mathcal{W}^{c}}\Phi$ and
$\mu_{\Sigma}(\mathcal{W})=1-\mu_{\Sigma}(\mathcal{W}^{c})$, we have
\begin{equation*}
\mathbb{E}_0\Phi-\int\Phi\,d\mu_{\Sigma}
=\frac{-\int_{\mathcal{W}^{c}}\Phi\,d\mu_{\Sigma}
+\mu_{\Sigma}(\mathcal{W}^{c})\int\Phi\,d\mu_{\Sigma}}{\mu_{\Sigma}(\mathcal{W})}.
\end{equation*}
By the Cauchy--Schwarz inequality,
\begin{equation*}
\Bigl|\int_{\mathcal{W}^{c}}\Phi\,d\mu_{\Sigma}\Bigr|
\le\|\Phi\|_{2}\,\mu_{\Sigma}(\mathcal{W}^{c})^{1/2}\le\|\Phi\|_{2}\eta^{1/2}
\end{equation*}
and $|\int\Phi\,d\mu_{\Sigma}|\le\|\Phi\|_{2}$. Since $\eta\le\frac14$, we also have
$\eta\le\frac12\eta^{1/2}$. Hence
\begin{equation*}
\Bigl|\mathbb{E}_0\Phi-\int\Phi\,d\mu_{\Sigma}\Bigr|
\le\frac{\|\Phi\|_{2}\,(\eta^{1/2}+\eta)}{3/4}
\le 2\|\Phi\|_{2}\,\eta^{1/2}\le 4\|\Phi\|_{2}\,\eta^{1/2}.
\end{equation*}
Finally, $\mathbb{E}_0|\Phi|=\mu_{\Sigma}(\mathcal{W})^{-1}\int_{\mathcal{W}}|\Phi|\,d\mu_{\Sigma}
\le\frac43\int|\Phi|\,d\mu_{\Sigma}$.

\emph{Proof of} (a). For $y\in\mathcal{W}$ write $e^{P}=1+P+\Delta_{2}$. By Taylor's formula with
Lagrange remainder, $\Delta_{2}=\frac12P^{2}e^{\xi}$ with $|\xi|\le|P|$, so that
$|\Delta_{2}|\le\frac12P^{2}e^{|P|}\le\frac12\delta^{2}e^{1/4}$. Since $(1.3)^{4}=2.8561>e$, we have
$e^{1/4}<1.3$, and therefore
\begin{equation}\label{8.2:eq-laplace-parity-1}
|\Delta_{2}|\le 0.65\,\delta^{2}\qquad\text{on }\mathcal{W}.
\end{equation}
Put $\Xi:=\mathbb{E}_0[e^{P}]=1+\mathbb{E}_0P+\mathbb{E}_0\Delta_{2}$. Then
$|\Xi-1|\le\delta+0.65\,\delta^{2}\le\delta+\delta^{2}$. Moreover
$e^{P}\ge e^{-1/4}>1/1.3>\frac34$ on $\mathcal{W}$, so $\Xi\ge\frac34$. Next,
\begin{equation*}
\langle G\rangle_{P}-\mathbb{E}_0G
=\frac{\mathbb{E}_0[Ge^{P}]-\mathbb{E}_0G\,\mathbb{E}_0[e^{P}]}{\Xi}
=\frac{\operatorname{Cov}_0(G,e^{P})}{\Xi}
=\frac{\operatorname{Cov}_0(G,P)+\operatorname{Cov}_0(G,\Delta_{2})}{\Xi},
\end{equation*}
because $\operatorname{Cov}_0(G,1)=0$. Using $1/\Xi=1-(\Xi-1)/\Xi$, this becomes
$\langle G\rangle_{P}=\mathbb{E}_0G+\operatorname{Cov}_0(G,P)+\epsilon_{2}(G)$ with
\begin{equation*}
\epsilon_{2}(G)=\frac{\operatorname{Cov}_0(G,\Delta_{2})}{\Xi}
-\operatorname{Cov}_0(G,P)\,\frac{\Xi-1}{\Xi}.
\end{equation*}
By \eqref{8.2:eq-laplace-parity-1},
$|\operatorname{Cov}_0(G,\Delta_{2})|=|\mathbb{E}_0[\widehat{G}\Delta_{2}]|\le 0.65\,\delta^{2}\,\mathbb{E}_0|\widehat{G}|$.
Likewise $|\operatorname{Cov}_0(G,P)|=|\mathbb{E}_0[\widehat{G}P]|\le\delta\,\mathbb{E}_0|\widehat{G}|$.
With $\Xi\ge\frac34$ and $\delta\le\frac14$ we obtain
\begin{equation*}
|\epsilon_{2}(G)|\le\frac43\bigl(0.65\,\delta^{2}+\delta(\delta+\delta^{2})\bigr)\mathbb{E}_0|\widehat{G}|
\le\frac43\,(1.9)\,\delta^{2}\,\mathbb{E}_0|\widehat{G}|\le 3\delta^{2}\,\mathbb{E}_0|\widehat{G}|.
\end{equation*}

\emph{Proof of} (d). We argue as in (a), now with $e^{P}=1+P+\frac12P^{2}+\Delta_{3}$ on $\mathcal{W}$.
By Taylor's formula, $|\Delta_{3}|\le\frac16|P|^{3}e^{|P|}\le\frac16\delta^{3}e^{\delta}$. Put
\begin{equation*}
\Gamma_{1}:=\operatorname{Cov}_0(G,P),\qquad
\Gamma_{2}:=\operatorname{Cov}_0(G,P^{2}),\qquad
\Gamma_{3}:=\operatorname{Cov}_0(G,\Delta_{3}).
\end{equation*}
As in (a), $\langle G\rangle_{P}-\mathbb{E}_0G=(\Gamma_{1}+\frac12\Gamma_{2}+\Gamma_{3})/\Xi$. The remainder
$\epsilon_{3}(G)$ is defined by the display in (d), that is,
$\epsilon_{3}(G)=\langle G\rangle_{P}-\mathbb{E}_0G-\Gamma_{1}-\frac12\Gamma_{2}+\Gamma_{1}\mathbb{E}_0P$.
Collecting terms, it equals
\begin{equation*}
\epsilon_{3}(G)=\Gamma_{1}\Bigl(\frac1\Xi-1+\mathbb{E}_0P\Bigr)+\frac12\Gamma_{2}\Bigl(\frac1\Xi-1\Bigr)
+\frac{\Gamma_{3}}{\Xi}.
\end{equation*}
We write $\Xi=1+\mathbb{E}_0P+\mathbb{E}_0\Delta_{2}$ as in (a). Expanding the product
$\Xi\,\mathbb{E}_0P$ gives
\begin{equation*}
\frac1\Xi-1+\mathbb{E}_0P=\frac{1-\Xi+\Xi\,\mathbb{E}_0P}{\Xi}
=\frac{(\mathbb{E}_0P)^{2}-\mathbb{E}_0\Delta_{2}+\mathbb{E}_0P\,\mathbb{E}_0\Delta_{2}}{\Xi}.
\end{equation*}
By \eqref{8.2:eq-laplace-parity-1} and $\Xi\ge\frac34$, this gives
$|\frac1\Xi-1+\mathbb{E}_0P|\le\frac43(1.65\,\delta^{2}+0.65\,\delta^{3})$. Also
$|\frac1\Xi-1|=|\Xi-1|/\Xi\le\frac43(\delta+0.65\,\delta^{2})$. The three covariances are bounded
through $\operatorname{Cov}_0(G,\Psi)=\mathbb{E}_0[\widehat{G}\Psi]$:
\begin{equation*}
|\Gamma_{1}|\le\delta\,\mathbb{E}_0|\widehat{G}|,\qquad
|\Gamma_{2}|\le\delta^{2}\,\mathbb{E}_0|\widehat{G}|,\qquad
|\Gamma_{3}|\le\tfrac16\delta^{3}e^{\delta}\,\mathbb{E}_0|\widehat{G}|.
\end{equation*}
Inserting these bounds, and using $\delta\le\frac14$ and $e^{\delta}\le e^{1/4}<1.3$, we obtain
\begin{align*}
|\epsilon_{3}(G)|
&\le\frac43\,\delta^{3}\,\mathbb{E}_0|\widehat{G}|
\Bigl(1.65+0.65\,\delta+\tfrac12(1+0.65\,\delta)+\tfrac16e^{\delta}\Bigr)\\
&\le\frac43\,\bigl(1.65+0.1625+0.5+0.08125+\tfrac{1.3}{6}\bigr)\,\delta^{3}\,\mathbb{E}_0|\widehat{G}|
\le 4\delta^{3}\,\mathbb{E}_0|\widehat{G}|,
\end{align*}
since the bracket in the second line is less than $2.62$. Thus the terms of order at most two in $P$
are exactly those displayed in (d), and the remainder obeys the stated bound.

\emph{Proof of} (b). If $\mathcal{W}$ were symmetric, $\mathbb{E}_0$ would be invariant under
$y\mapsto-y$. Then $\mathbb{E}_0$ of an odd function would vanish, and so would the covariance of an
even function with an odd one. In general we proceed as follows. Let $\Phi$ be a function on
$\mathcal{W}$ with $\|\Phi\|_{2}<\infty$ (extension by zero off $\mathcal{W}$) which is odd on
$\mathcal{W}_{\mathrm{s}}$. The change of variables $y\mapsto-y$ leaves $\mu_{\Sigma}$ and
$\mathcal{W}_{\mathrm{s}}$ invariant, so
$\int_{\mathcal{W}_{\mathrm{s}}}\Phi\,d\mu_{\Sigma}=-\int_{\mathcal{W}_{\mathrm{s}}}\Phi\,d\mu_{\Sigma}=0$.
Hence $\int_{\mathcal{W}}\Phi\,d\mu_{\Sigma}=\int_{\mathcal{W}_{\mathrm{a}}}\Phi\,d\mu_{\Sigma}$. The
Cauchy--Schwarz inequality and $\mu_{\Sigma}(\mathcal{W})\ge\frac34$ then give
\begin{equation}\label{8.2:eq-laplace-parity-2}
|\mathbb{E}_0\Phi|\le\frac{\|\Phi\|_{2}\,\mu_{\Sigma}(\mathcal{W}_{\mathrm{a}})^{1/2}}{\mu_{\Sigma}(\mathcal{W})}
\le\frac43\,\|\Phi\|_{2}\,\eta^{1/2}.
\end{equation}
Applied to $\Phi=G^{-}$, \eqref{8.2:eq-laplace-parity-2} gives
$|\mathbb{E}_0G^{-}|\le\frac43\|G^{-}\|_{2}\eta^{1/2}\le 4\|G^{-}\|_{2}\eta^{1/2}$, the first
assertion of (b).

Next we note the bounds on the parts of $P$ and of $\widehat{G}$ used below. On
$\mathcal{W}_{\mathrm{s}}$ we have $|P^{-}(y)|\le\frac12(|P(y)|+|P(-y)|)\le\delta$, and $P^{-}=0$ on
$\mathcal{W}_{\mathrm{a}}$. Thus $|P^{-}|\le\delta$ and $|P^{+}|\le\delta^{+}\le\delta$ on
$\mathcal{W}$. Put $\widehat{G}^{\pm}:=G^{\pm}-\mathbb{E}_0G^{\pm}$ on $\mathcal{W}$, extended by zero
off $\mathcal{W}$. Then
$\|\widehat{G}^{\pm}\|_{2}\le\|G^{\pm}\|_{2}+|\mathbb{E}_0G^{\pm}|\,\mu_{\Sigma}(\mathcal{W})^{1/2}$.
Since
$|\mathbb{E}_0G^{\pm}|\le(\mathbb{E}_0[(G^{\pm})^{2}])^{1/2}\le\|G^{\pm}\|_{2}\,\mu_{\Sigma}(\mathcal{W})^{-1/2}$,
it follows that $\|\widehat{G}^{\pm}\|_{2}\le 2\|G^{\pm}\|_{2}$.

Since $G=G^{+}+G^{-}$ and $P=P^{+}+P^{-}$ on $\mathcal{W}$ and $\mathbb{E}_0$ is linear, we have
$\widehat{G}=\widehat{G}^{+}+\widehat{G}^{-}$ on $\mathcal{W}$. Therefore
\begin{equation*}
\operatorname{Cov}_0(G,P)=\mathbb{E}_0[\widehat{G}P]
=\mathbb{E}_0[\widehat{G}^{+}P^{+}]+\mathbb{E}_0[\widehat{G}^{-}P^{-}]
+\mathbb{E}_0[\widehat{G}^{+}P^{-}]+\mathbb{E}_0[\widehat{G}^{-}P^{+}].
\end{equation*}
We treat the four terms in turn.

The integrand $\widehat{G}^{+}P^{-}$ is odd on $\mathcal{W}_{\mathrm{s}}$, being the product of the even
function $G^{+}-\mathbb{E}_0G^{+}$ and the odd function $P^{-}$. Its norm satisfies
$\|\widehat{G}^{+}P^{-}\|_{2}\le\delta\|\widehat{G}^{+}\|_{2}\le 2\delta\|G^{+}\|_{2}$. By
\eqref{8.2:eq-laplace-parity-2},
$|\mathbb{E}_0[\widehat{G}^{+}P^{-}]|\le\frac83\,\delta\,\|G^{+}\|_{2}\,\eta^{1/2}$.

For the last term we write
$\mathbb{E}_0[\widehat{G}^{-}P^{+}]=\mathbb{E}_0[G^{-}P^{+}]-\mathbb{E}_0G^{-}\,\mathbb{E}_0P^{+}$.
The integrand $G^{-}P^{+}$ is odd on $\mathcal{W}_{\mathrm{s}}$ and has
$\|G^{-}P^{+}\|_{2}\le\delta\|G^{-}\|_{2}$, so \eqref{8.2:eq-laplace-parity-2} bounds its expectation
by $\frac43\delta\|G^{-}\|_{2}\eta^{1/2}$. The product term is controlled by the smallness of
$\mathbb{E}_0G^{-}$ just shown, together with $|\mathbb{E}_0P^{+}|\le\delta$; it is at most
$\frac43\delta\|G^{-}\|_{2}\eta^{1/2}$. Hence
$|\mathbb{E}_0[\widehat{G}^{-}P^{+}]|\le\frac83\,\delta\,\|G^{-}\|_{2}\,\eta^{1/2}$.

The first term is
$\mathbb{E}_0[\widehat{G}^{+}P^{+}]=\operatorname{Cov}_0(G^{+},P^{+})$. The second is
$\mathbb{E}_0[\widehat{G}^{-}P^{-}]=\mathbb{E}_0[G^{-}P^{-}]-\mathbb{E}_0G^{-}\,\mathbb{E}_0P^{-}$. Its
correction satisfies
$|\mathbb{E}_0G^{-}\,\mathbb{E}_0P^{-}|\le\frac43\,\delta\,\|G^{-}\|_{2}\,\eta^{1/2}$, again by
\eqref{8.2:eq-laplace-parity-2} and $|P^{-}|\le\delta$. We put
\begin{equation*}
\epsilon_{\mathrm{par}}:=\mathbb{E}_0[\widehat{G}^{+}P^{-}]+\mathbb{E}_0[\widehat{G}^{-}P^{+}]
-\mathbb{E}_0G^{-}\,\mathbb{E}_0P^{-}.
\end{equation*}
Then
\begin{equation*}
\operatorname{Cov}_0(G,P)=\operatorname{Cov}_0(G^{+},P^{+})+\mathbb{E}_0[G^{-}P^{-}]+\epsilon_{\mathrm{par}},
\end{equation*}
and the three bounds just obtained give
\begin{equation*}
|\epsilon_{\mathrm{par}}|
\le\Bigl(\tfrac83\|G^{+}\|_{2}+\bigl(\tfrac83+\tfrac43\bigr)\|G^{-}\|_{2}\Bigr)\delta\,\eta^{1/2}
\le 16\,\delta\bigl(\|G^{+}\|_{2}+\|G^{-}\|_{2}\bigr)\eta^{1/2}.
\end{equation*}

Finally,
$\operatorname{Cov}_0(G^{+},P^{+})=\mathbb{E}_0[\widehat{G}^{+}(P^{+}-\mathbb{E}_0P^{+})]$. Since
$|P^{+}|\le\delta^{+}$ on $\mathcal{W}$, we have $|P^{+}-\mathbb{E}_0P^{+}|\le 2\delta^{+}$, and hence
$|\operatorname{Cov}_0(G^{+},P^{+})|\le 2\delta^{+}\,\mathbb{E}_0|G^{+}-\mathbb{E}_0G^{+}|$. Since
$|P^{-}|\le\delta$ on $\mathcal{W}$, we also have
$|\mathbb{E}_0[G^{-}P^{-}]|\le\delta\,\mathbb{E}_0|G^{-}|$. This completes the proof of (b).
\end{proof}

\begin{lemma}[Laplace extraction for a tilted covariance]\label{8.2:lem-laplace-covariance}
Work in the setting and notation of Lemma~\ref{8.2:lem-laplace-parity}: the centred Gaussian
measure $\gamma$, the window $W$ with $\gamma(W^c)\le t$ and $\gamma(W\setminus(-W))\le t$, where
$t\le\frac14$, the normalised average $E_0$, the norm $M_2$, the centring $\hat H=H-E_0H$, the
function $P$ with $\sup_W|P|\le\vartheta\le\frac14$ and its even and odd parts $P^{\pm}$, and the
tilted expectation $\langle\cdot\rangle$. We write
$\operatorname{Cov}_0(H_1,H_2):=E_0[H_1H_2]-E_0[H_1]\,E_0[H_2]$ for the covariance under $E_0$,
$W_s:=W\cap(-W)$ for the symmetric part of the window, and $H^{\wedge}:=\hat H$. For a function $H$
on $W$, $H^{\pm}$ denote its even and odd parts in the sense of
Lemma~\ref{8.2:lem-laplace-parity}, and $\hat G_i^{\pm}$ are the even and odd parts of $\hat G_i$.

Let $G_1,G_2$ satisfy $M_2(G_i^2)<\infty$ for $i=1,2$. Then the tilted covariance
$\langle G_1;G_2\rangle:=\langle G_1G_2\rangle-\langle G_1\rangle\langle G_2\rangle$ satisfies
\begin{equation}\label{8.2:eq-laplace-covariance-1}
\langle G_1;G_2\rangle=\operatorname{Cov}_0(G_1,G_2)+\operatorname{Cov}_0(\hat G_1\hat G_2,P)
-\operatorname{Cov}_0(G_1,P)\operatorname{Cov}_0(G_2,P)+\rho^{\mathrm{cov}},
\end{equation}
where the remainder $\rho^{\mathrm{cov}}$ obeys
\begin{equation}\label{8.2:eq-laplace-covariance-2}
\begin{split}
|\rho^{\mathrm{cov}}|\le{}&6\vartheta^2E_0|\hat G_1\hat G_2|
+10\vartheta^3E_0|\hat G_1|\,E_0|\hat G_2|\\
&+3\vartheta^2\big(E_0|\hat G_1|\cdot|\operatorname{Cov}_0(G_2,P)|
+E_0|\hat G_2|\cdot|\operatorname{Cov}_0(G_1,P)|\big).
\end{split}
\end{equation}
Sorted by parity on $W_s$, the two leading terms satisfy
\begin{equation}\label{8.2:eq-laplace-covariance-3}
\begin{split}
\operatorname{Cov}_0(G_1,G_2)={}&\operatorname{Cov}_0(G_1^{+},G_2^{+})+E_0[G_1^{-}G_2^{-}]\\
&+O\big((M_2(G_1^{+})M_2(G_2^{-})+M_2(G_1^{-})M_2(G_2^{+}))\,t^{1/2}\big)
\end{split}
\end{equation}
and
\begin{equation}\label{8.2:eq-laplace-covariance-4}
\begin{split}
\operatorname{Cov}_0(\hat G_1\hat G_2,P)
={}&E_0\big[(\hat G_1^{+}\hat G_2^{+}+\hat G_1^{-}\hat G_2^{-})^{\wedge}P^{+}\big]\\
&+E_0\big[(\hat G_1^{+}\hat G_2^{-}+\hat G_1^{-}\hat G_2^{+})P^{-}\big]+\mathcal{E},
\end{split}
\end{equation}
where $\mathcal{E}=O(\vartheta\,t^{1/2})$ times products of $M_2$-norms of the functions built from
$G_1$ and $G_2$.
\end{lemma}

\begin{proof}
\emph{Integrability.} Since $\gamma$ is a probability measure, the Cauchy--Schwarz inequality gives
$M_2(G_i)\le M_2(G_i^2)^{1/2}<\infty$. Moreover $\gamma(W)\ge 1-t\ge\frac34$, so $E_0G_i$ is finite.
From $(G_i-c)^2\le 2G_i^2+2c^2$ we get $M_2(\hat G_i^{\,2})<\infty$. The Cauchy--Schwarz inequality
then gives
\[
M_2(\hat G_1\hat G_2)\le M_2(\hat G_1^{\,2})^{1/2}M_2(\hat G_2^{\,2})^{1/2}<\infty .
\]
Because $|P|\le\vartheta$ on $W$, all the tilted expectations written below are finite.

\emph{Reduction to centred functions.} The functional $\langle\cdot\rangle$ is linear and satisfies
$\langle 1\rangle=1$. Hence, for constants $a,b$,
\begin{align*}
&\langle(G_1+a)(G_2+b)\rangle-\langle G_1+a\rangle\langle G_2+b\rangle\\
&\quad=\langle G_1G_2\rangle+a\langle G_2\rangle+b\langle G_1\rangle+ab
-(\langle G_1\rangle+a)(\langle G_2\rangle+b)
=\langle G_1;G_2\rangle .
\end{align*}
Taking $a=-E_0G_1$ and $b=-E_0G_2$ gives
\begin{equation}\label{8.2:eq-laplace-covariance-5}
\langle G_1;G_2\rangle=\langle\hat G_1\hat G_2\rangle-\langle\hat G_1\rangle\langle\hat G_2\rangle .
\end{equation}

\emph{Part (a) of Lemma~\ref{8.2:lem-laplace-parity}.} We use part (a) in the following form. For $H$
with $M_2(H)<\infty$,
\[
\langle H\rangle=E_0[H]+\operatorname{Cov}_0(H,P)+\rho_2(H),\qquad
|\rho_2(H)|\le3\vartheta^2E_0|H-E_0H| .
\]
Here $\rho_2(H)$ denotes the remainder of that part.

\emph{The product of the means.} Apply this to $H=\hat G_i$. We have $E_0\hat G_i=0$, so
$\hat G_i-E_0\hat G_i=\hat G_i$. The covariance $\operatorname{Cov}_0(\cdot,P)$ does not change when a
constant is added to its first argument, so $\operatorname{Cov}_0(\hat G_i,P)=\operatorname{Cov}_0(G_i,P)$.
Therefore
\[
\langle\hat G_i\rangle=\operatorname{Cov}_0(G_i,P)+\rho_2(\hat G_i),\qquad
|\rho_2(\hat G_i)|\le3\vartheta^2E_0|\hat G_i| .
\]
Multiplying the two identities gives
\begin{equation}\label{8.2:eq-laplace-covariance-6}
\langle\hat G_1\rangle\langle\hat G_2\rangle=\operatorname{Cov}_0(G_1,P)\operatorname{Cov}_0(G_2,P)+r,
\end{equation}
where
\[
r=\rho_2(\hat G_1)\operatorname{Cov}_0(G_2,P)+\rho_2(\hat G_2)\operatorname{Cov}_0(G_1,P)
+\rho_2(\hat G_1)\rho_2(\hat G_2).
\]
The remainder is bounded by
\[
|r|\le3\vartheta^2E_0|\hat G_1|\,|\operatorname{Cov}_0(G_2,P)|
+3\vartheta^2E_0|\hat G_2|\,|\operatorname{Cov}_0(G_1,P)|+9\vartheta^4E_0|\hat G_1|\,E_0|\hat G_2| .
\]

\emph{The mean of the product.} Apply part (a) to $H=\hat G_1\hat G_2$, which is admissible by the first
paragraph. By definition of $\operatorname{Cov}_0$,
\[
E_0[\hat G_1\hat G_2]=E_0[G_1G_2]-E_0G_1\,E_0G_2=\operatorname{Cov}_0(G_1,G_2).
\]
Hence
\begin{equation}\label{8.2:eq-laplace-covariance-7}
\langle\hat G_1\hat G_2\rangle=\operatorname{Cov}_0(G_1,G_2)+\operatorname{Cov}_0(\hat G_1\hat G_2,P)
+\rho_2(\hat G_1\hat G_2).
\end{equation}
Since $|E_0[\hat G_1\hat G_2]|\le E_0|\hat G_1\hat G_2|$, the triangle inequality gives
\[
|\rho_2(\hat G_1\hat G_2)|\le3\vartheta^2E_0\big|\hat G_1\hat G_2-E_0[\hat G_1\hat G_2]\big|
\le6\vartheta^2E_0|\hat G_1\hat G_2| .
\]

\emph{Subtraction.} Subtract \eqref{8.2:eq-laplace-covariance-6} from
\eqref{8.2:eq-laplace-covariance-7} and use \eqref{8.2:eq-laplace-covariance-5}. This gives
\eqref{8.2:eq-laplace-covariance-1} with $\rho^{\mathrm{cov}}=\rho_2(\hat G_1\hat G_2)-r$. The bounds on
$\rho_2(\hat G_1\hat G_2)$ and $r$ give \eqref{8.2:eq-laplace-covariance-2}, once we note that
\[
9\vartheta^4\le\tfrac94\vartheta^3\le10\vartheta^3
\]
because $\vartheta\le\frac14$.

\emph{Parity sorting.} Write $G_i=G_i^{+}+G_i^{-}$ and expand $\operatorname{Cov}_0(G_1,G_2)$ by
bilinearity:
\begin{align*}
\operatorname{Cov}_0(G_1,G_2)={}&\operatorname{Cov}_0(G_1^{+},G_2^{+})+E_0[G_1^{-}G_2^{-}]
-E_0G_1^{-}\,E_0G_2^{-}\\
&+\operatorname{Cov}_0(G_1^{+},G_2^{-})+\operatorname{Cov}_0(G_1^{-},G_2^{+}).
\end{align*}
Each term on the second line, and the product $E_0G_1^{-}\,E_0G_2^{-}$, is built from $E_0$-averages of
functions that are odd on $W_s$. These are estimated by the parity sorting of part (b) of
Lemma~\ref{8.2:lem-laplace-parity}, which gives \eqref{8.2:eq-laplace-covariance-3}.

For the second leading term, the definition of $\operatorname{Cov}_0$ gives
$\operatorname{Cov}_0(\hat G_1\hat G_2,P)=E_0[(\hat G_1\hat G_2)^{\wedge}P]$. Write
$\hat G_i=\hat G_i^{+}+\hat G_i^{-}$. Then
\[
\hat G_1\hat G_2=(\hat G_1^{+}\hat G_2^{+}+\hat G_1^{-}\hat G_2^{-})
+(\hat G_1^{+}\hat G_2^{-}+\hat G_1^{-}\hat G_2^{+}),
\]
the sum of a function even on $W_s$ and a function odd on $W_s$. Insert $P=P^{+}+P^{-}$. This gives the
two terms displayed in \eqref{8.2:eq-laplace-covariance-4}, together with the following remaining
contributions:
\begin{itemize}
\item the even part paired with $P^{-}$;
\item the odd part paired with $P^{+}$;
\item the difference between $(\cdot)^{\wedge}$ and $(\cdot)$ on the odd part paired with $P^{-}$, namely
\[
-E_0[\hat G_1^{+}\hat G_2^{-}+\hat G_1^{-}\hat G_2^{+}]\,E_0[P^{-}].
\]
\end{itemize}
All of these involve $E_0$-averages of functions odd on $W_s$, each carrying a factor bounded by
$\vartheta$. They are estimated by the parity sorting of part (b) of Lemma~\ref{8.2:lem-laplace-parity},
and together they form the error $\mathcal{E}$ of \eqref{8.2:eq-laplace-covariance-4}.
\end{proof}

\begin{remark}\label{8.2:lem-one-step-extraction}
The statement of this lemma is not written out here; parts of its proof are printed below.
\end{remark}

\begin{proof}[Proof of Lemma~\ref{8.2:lem-one-step-extraction}, continued: collection of the terms]
\label{8.2:prf-step-collection}
We keep the notation fixed in the first two steps of this proof. The conditional expectation
$E_k[G\,|\,W]$ has been reduced, up to the localisation error $C_{\mathrm{loc}}\operatorname{osc}(G)
e^{-c_{\mathrm{loc}}(D-\ell)}$, to the box expectation $\langle\mathbf{G}\rangle$. Here
$\mathbf{G}(B)=G(\mathfrak{v}_k(B;W))$, the weight is $t$-symmetric with $t=t_k$, and the tilt is
$P=\mathcal{P}_{k,\mathfrak{B}}$ with parity parts $P^{\pm}$ and smallness $\vartheta=\vartheta_k
+\vartheta^{+}_k$. The field on $\mathfrak{B}_\ell(p_k)$ is expanded as
$Y=g_kY_1+g_k^2Y_2+g_k^3Y_3+\dots$, with $Y_1=iCB$ and $Y_2=-i\,h\tilde{D}_2(CB,CB)$. The parity
parts of $\mathbf{G}$ are $\mathbf{G}^{-}$ and $\mathbf{G}^{+}-G(\mathrm{bg})=g_k^2\mathbf{G}_2
+\mathbf{G}^{+}_{\ge 4}$, with the even second-order jet $\mathbf{G}_2$ and the even remainder
$\mathbf{G}^{+}_{\ge4}$ as in the previous step.
We write $\lambda=\|DG(\mathrm{bg})\|$ and $m_j=E_\gamma\|CB\|^j_{\mathfrak{B}_\ell'}$. The letter
$\gamma$ denotes the box Gaussian law $N(0,C^{(k)}(W))$, and $E_\gamma$ is its expectation.

\emph{The expansion.} We apply items (a)--(c) of the expansion lemma for expectations under a
tilted, $t$-symmetric box weight, in the notation fixed there. In that notation, $E_0$ and
$\operatorname{Cov}_0$ are the expectation and the covariance under the untilted $t$-symmetric box
weight; $M_2(\cdot)$ is the functional of that lemma through which the difference between $E_0$ and
the Gaussian expectation $E_\gamma$ is bounded by a multiple of $t_k^{1/2}$; $\rho_{\mathrm{par}}$
is the remainder of that lemma which is bounded in the same way by a multiple of $t_k^{1/2}$; and
$\rho_2(\cdot)$ is its remainder of second order in the tilt. The lemma gives
\begin{equation}\label{8.2:eq-step-collection-1}
\begin{split}
\langle\mathbf{G}\rangle
={}&G(\mathrm{bg})+E_0\bigl[g_k^2\mathbf{G}_2\bigr]+E_0\bigl[\mathbf{G}^{+}_{\ge4}\bigr]
+E_0\mathbf{G}^{-}+\operatorname{Cov}_0(\mathbf{G}^{+},P^{+})\\
&+E_0\bigl[\mathbf{G}^{-}P^{-}\bigr]+\rho_{\mathrm{par}}+\rho_2(\mathbf{G}).
\end{split}
\end{equation}
We treat the terms on the right one at a time.

\emph{The second-order term.} By the same lemma, $E_0$ may be replaced by the Gaussian expectation
on $g_k^2\mathbf{G}_2$ at the cost of a $t$-term:
\begin{equation}\label{8.2:eq-step-collection-2}
E_0\bigl[g_k^2\mathbf{G}_2\bigr]
=g_k^2\int\mathbf{G}_2\,d\gamma+O\bigl(g_k^2M_2(\mathbf{G}_2)\,t_k^{1/2}\bigr).
\end{equation}
Insert the definition of $\mathbf{G}_2$ from the previous step, namely $DG(\mathrm{bg})[Y_2]$ plus
one half of $D^2G(\mathrm{bg})[Y_1,Y_1]$, and use the linearity of $DG(\mathrm{bg})$ to move the
expectation inside the first differential. This gives
\begin{equation}\label{8.2:eq-step-collection-3}
\begin{split}
\int\mathbf{G}_2\,d\gamma
&=\tfrac12E_\gamma D^2G(\mathrm{bg})[Y_1,Y_1]+DG(\mathrm{bg})\bigl[E_\gamma Y_2\bigr]\\
&=\tfrac12\operatorname{Tr}_{C^{(k)}}\bigl[\operatorname{Hess}_{\mathrm{bg}}G\bigr]
+\bigl\langle DG(\mathrm{bg}),\mathfrak{d}_k\bigr\rangle .
\end{split}
\end{equation}
The first equality in the second line holds because $Y_1=iCB$ restricted to
$\mathfrak{B}_\ell(p_k)$. Hence $E_\gamma D^2G(\mathrm{bg})[Y_1,Y_1]$ is, by definition, the
Gaussian trace $\operatorname{Tr}_{C^{(k)}}[\operatorname{Hess}_{\mathrm{bg}}G]$ of part~(i) of
Lemma~\ref{8.2:lem-one-step-extraction}. In the second term, $\mathfrak{d}_k$ is the mean-field
shift of part~(iii) of Lemma~\ref{8.2:lem-one-step-extraction}. Since $h\tilde{D}_2$ is quadratic,
$h\tilde{D}_2(g_kCB)=g_k^2\,h\tilde{D}_2(CB,CB)$, and therefore
\[
\mathfrak{d}_k=-\frac{E_\gamma\bigl[i\,h\tilde{D}_2(g_kCB)\bigr]}{g_k^2}
=-E_\gamma\bigl[i\,h\tilde{D}_2(CB,CB)\bigr]=E_\gamma Y_2 .
\]
We write the action of $DG(\mathrm{bg})$ on a vector as the pairing
$\langle DG(\mathrm{bg}),\cdot\rangle$. Thus $DG(\mathrm{bg})[E_\gamma Y_2]=\langle
DG(\mathrm{bg}),\mathfrak{d}_k\rangle$. By the jet bound $\|Y_2\|\le C_2\|CB\|^2_{\mathfrak{B}_\ell'}$
of the previous step,
\[
\|\mathfrak{d}_k\|\le E_\gamma\|Y_2\|\le C_2m_2 .
\]

\emph{The even term of order four.} Since $|E_0X|\le E_0|X|$, the moment bound
$E_0|\mathbf{G}^{+}_{\ge4}|\le Cg_k^4\mu m_4$ of the previous step gives
\[
\bigl|E_0\mathbf{G}^{+}_{\ge4}\bigr|\le Cg_k^4\mu m_4 .
\]

\emph{The odd term.} In the notation of the expansion lemma, the odd part satisfies
\[
\bigl|E_0\mathbf{G}^{-}\bigr|\le 4M_2(\mathbf{G}^{-})\,t_k^{1/2}.
\]

\emph{The even covariance.} Subtracting the constant $G(\mathrm{bg})$ from $\mathbf{G}^{+}$ does
not change the covariance with $P^{+}$, and the even part of the tilt has smallness
$\vartheta^{+}_k$. The bound on $E_0|\mathbf{G}^{+}-G(\mathrm{bg})|$ from the previous step then
gives
\[
\bigl|\operatorname{Cov}_0(\mathbf{G}^{+},P^{+})\bigr|
\le2\vartheta^{+}_kE_0\bigl|\mathbf{G}^{+}-G(\mathrm{bg})\bigr|
\le C\vartheta^{+}_k\bigl(g_k^2(\lambda+\mu_2)m_2+g_k^4\mu m_4\bigr).
\]
Now insert the size $C g_k^2n_{\mathfrak{B}}\mathfrak{b}_k^4(1+B^{\sharp}\tau_*)$ of
$\vartheta^{+}_k$, the smallness of the even part of the tilt identified in the first step. The
right-hand side becomes
\[
C g_k^2n_{\mathfrak{B}}\mathfrak{b}_k^4(1+B^{\sharp}\tau_*)
\bigl(g_k^2(\lambda+\mu_2)m_2+g_k^4\mu m_4\bigr).
\]
This term is therefore of order $g_k^4$ times a polynomial prefactor, plus a term of order
$g_k^4\lambda$. Throughout this proof, a polynomial prefactor means a factor of the kind entering
$\Pi_1$ and $\Pi_1'$, that is, polynomial in $D$ and $\log x_k$; we write $\Pi$ for a prefactor of
this kind, whose value may change from one occurrence to the next.

\emph{The odd covariance.} By the parity decomposition of the previous step, $\mathbf{G}^{-}$
equals $g_kDG(\mathrm{bg})[Y_1]$ plus odd terms of order at least three in $g_k$. The odd part of
the tilt has smallness $\vartheta_k$. The terms of order at least three in $g_k$ therefore
contribute $O(g_k^3\mu m_3\vartheta_k)$, and
\[
E_0\bigl[\mathbf{G}^{-}P^{-}\bigr]=g_kE_0\bigl[DG(\mathrm{bg})[Y_1]P^{-}\bigr]
+O\bigl(g_k^3\mu m_3\vartheta_k\bigr).
\]
By linearity of $DG(\mathrm{bg})$, the first term equals $g_k\langle
DG(\mathrm{bg}),E_0[Y_1P^{-}]\rangle$. We define the two vectors $\mathfrak{t}_k$ and
$\mathfrak{l}_k$ by splitting $P^{-}$ according to its two sources:
\begin{itemize}
\item the part of $P^{-}$ coming from the action and measure vertices gives the tadpole
$\mathfrak{t}_k$;
\item the part linear in $B$ coming from the increments of the correction terms of the level-$k$
effective density $\rho_k$ (its localised part, its remainder and the lag terms) gives $C^{(k)}$
applied to the coefficient form of that linear part, which is $\mathfrak{l}_k$.
\end{itemize}
This splitting gives
\begin{equation}\label{8.2:eq-step-collection-4}
E_0\bigl[Y_1P^{-}\bigr]=g_k(\mathfrak{t}_k+\mathfrak{l}_k)\bigl(1+O(\vartheta)\bigr).
\end{equation}
Both are $\mathfrak{g}$-valued vectors on $\mathfrak{B}_\ell'$, and, by the size of the odd tilt
smallness $\vartheta_k$ identified in the first step,
\[
\|\mathfrak{t}_k\|,\ \|\mathfrak{l}_k\|\le\frac{\vartheta_km_1}{g_k}
\le Cn_{\mathfrak{B}}\mathfrak{b}_k^3m_1 .
\]
This is the polylogarithmic bound on the vectors in part~(iii) of
Lemma~\ref{8.2:lem-one-step-extraction}. Since $|\langle DG(\mathrm{bg}),v\rangle|\le\lambda\|v\|$,
we conclude that
\[
\bigl|E_0[\mathbf{G}^{-}P^{-}]\bigr|\le g_k^2\lambda\cdot Cn_{\mathfrak{B}}\mathfrak{b}_k^3m_1
+Cg_k^3\mu m_3\cdot g_kn_{\mathfrak{B}}\mathfrak{b}_k^3 .
\]

\emph{The second-order remainder of the tilt.} The bounds on $E_0|\mathbf{G}^{-}|$ and on
$E_0|\mathbf{G}^{+}-G(\mathrm{bg})|$ from the previous step control $E_0|\mathbf{G}-E_0\mathbf{G}|$.
This gives
\[
|\rho_2(\mathbf{G})|\le3(\vartheta_k+\vartheta^{+}_k)^2E_0\bigl|\mathbf{G}-E_0\mathbf{G}\bigr|
\le Cg_k^2n_{\mathfrak{B}}^2\mathfrak{b}_k^6\bigl(g_k\lambda m_1+g_k^2\mu_2m_2+\dots\bigr).
\]
Here the dots stand for the further terms of that bound, namely $g_k^2\lambda m_2$,
$g_k^3\mu m_3$ and $g_k^4\mu m_4$, each with a constant. This remainder is therefore of order
$g_k^3\lambda$ times a polynomial prefactor, plus $g_k^4$ times a polynomial prefactor.

\emph{The $t$-terms.} The remainder $\rho_{\mathrm{par}}$ and the $t$-term of
\eqref{8.2:eq-step-collection-2} are bounded by $C\operatorname{osc}(G)\,t_k^{1/2}$, which is the
form of the corresponding tail in part~(ii) of Lemma~\ref{8.2:lem-one-step-extraction}. The same
holds for the odd term bounded above.

\emph{Collection.} Combine \eqref{8.2:eq-step-collection-1} with the localisation of the first
step. The term exact at order $g_k^2$ is
\[
\tfrac12g_k^2\operatorname{Tr}_{C^{(k)}}\bigl[\operatorname{Hess}_{\mathrm{bg}}G\bigr],
\]
which comes from \eqref{8.2:eq-step-collection-2} and \eqref{8.2:eq-step-collection-3}. The terms
carrying $\lambda$ at order $g_k^2$ have two origins. The first is the shift term of
\eqref{8.2:eq-step-collection-3}. The second is the first term of the odd covariance, evaluated
through \eqref{8.2:eq-step-collection-4}. Together with the $O(\vartheta)$ correction of
\eqref{8.2:eq-step-collection-4}, they form the $\kappa$-member
\[
g_k^2\Bigl[\bigl\langle DG(\mathrm{bg}),\mathfrak{d}_k\bigr\rangle
+\bigl\langle DG(\mathrm{bg}),\mathfrak{t}_k+\mathfrak{l}_k\bigr\rangle
+O(\vartheta)\,\lambda\,\Pi\Bigr].
\]
The two remaining terms carrying $\lambda$ are also of the form $g_k^2\,O(\vartheta)\,\lambda\,\Pi$
and are counted in the last member of this display. The first is the term of order $g_k^4\lambda$
of the even covariance, which is $g_k^2$ times $\vartheta^{+}_k$ times $C\lambda m_2$. The second is
the term of order $g_k^3\lambda\Pi$ in $\rho_2(\mathbf{G})$; it equals
$3(\vartheta_k+\vartheta^{+}_k)$ times $(\vartheta_k+\vartheta^{+}_k)g_k\lambda m_1$ with a
constant, and $(\vartheta_k+\vartheta^{+}_k)g_k\le Cg_k^2\Pi$ by the sizes of $\vartheta_k$ and
$\vartheta^{+}_k$ used above, so that it is at most $g_k^2\lambda\Pi$. We use part~(b) of the
gradient bound for gauge-invariant local observables at the background:
\[
\lambda=\|DG(\mathrm{bg})\|\le C_{\mathrm{ax}}\mu_2\kappa_{\mathrm{bg}}.
\]
Together with the bounds on $\mathfrak{d}_k$, $\mathfrak{t}_k$ and $\mathfrak{l}_k$ above, this
shows that the $\kappa$-member is at most $g_k^2\Pi_1'\mu_2\kappa_{\mathrm{bg}}$. It consists of
pairings of $DG(\mathrm{bg})$ with the three vectors $\mathfrak{t}_k$, $\mathfrak{d}_k$ and
$\mathfrak{l}_k$, up to the $O(\vartheta)$ correction. Every other term is either at most
$Cg_k^4\Pi_1\mu$ or a tail:
\begin{itemize}
\item the even term of order four;
\item the parts of the two covariances that do not carry $\lambda$;
\item the part of $\rho_2(\mathbf{G})$ that does not carry $\lambda$;
\item the localisation error and the $t$-terms.
\end{itemize}
Finally we define $\mathfrak{L}_k(G;W)+\mathcal{E}^{(1)}_k(G;W)$ to be everything beyond the two
terms $G(\mathrm{bg})$ and $\tfrac12g_k^2\operatorname{Tr}_{C^{(k)}}[\operatorname{Hess}_{\mathrm{bg}}
G]$. With this definition the collection above gives the identity of part~(i) of
Lemma~\ref{8.2:lem-one-step-extraction}; the bounds collected are part~(ii), with $C_{\mathrm{SX}}$
the largest of the constants $C$ above; and the description of the $\kappa$-member is part~(iii).
\end{proof}

Clause~(iv) of Lemma~\ref{8.2:lem-one-step-extraction} concerns each of the $\mathfrak{g}$-valued one-line vectors
$\mathfrak{w}\in\{\mathfrak{t}_k,\mathfrak{d}_k,\mathfrak{l}_k\}$: the tadpole, the mean-field shift and the linear-increment vector. The proof takes the following four statements as its hypotheses.
\begin{itemize}
\item[(a)] The hypothesis of analyticity of the step on the complex window.
\item[(b)] Clause (iii) of localisation of the sourced step.
\item[(c)] The theorem on gauge-equivariant Lie-algebra-valued functions of a lattice gauge field. It carries standing restrictions concerning the centre of the structure group, the topology of the domain and analyticity; they are checked at the end of the proof below. Under these restrictions, let $\mathfrak{w}$ be a $\mathfrak{g}$-valued function of a gauge field which
  \begin{itemize}
  \item is gauge-equivariant,
  \item is holomorphic on a complex neighbourhood of the real configurations, with supremum at most $M$ there, and
  \item depends only on the links of a finite box.
  \end{itemize}
  Then $\mathfrak{w}$ vanishes at flat configurations. Moreover, its norm is at most $M$ times the maximum of finitely many geometric constants of the box, times the curvature size of the box. Each of these geometric constants is at most linear in the side of the box.
\item[(d)] The lemma stating that averaging an equivariant function against a gauge-covariant family of measures yields an equivariant function.
\end{itemize}

\begin{proof}[Proof of clause (iv) of Lemma~\ref{8.2:lem-one-step-extraction}]
\label{8.2:prf-one-line-equivariance}
Fix $W$ in the complex window at level $k+1$. Each of $\mathfrak{t}_k(W)$, $\mathfrak{d}_k(W)$, $\mathfrak{l}_k(W)$ is an expectation of a $\mathfrak{g}$-valued function of the pair formed by the fluctuation field $B$ and the background. For the main terms the expectation is taken under the box Gaussian law $N(0,C^{(k)}(W))$. For the corrections of order $\vartheta$ it is taken under the box-tilted law.

These functions are built from the analytic covariant objects of Ba{\l}aban's construction:
\begin{itemize}
\item the field $CB$;
\item the analytic covariant operator $\tilde{D}_2$ of Ba{\l}aban's construction;
\item the first-order vertices;
\item the coefficient forms of the coupling increment.
\end{itemize}

\emph{Equivariance.} Let $u$ be a gauge transformation of $W$. The background transforms as a gauge field, and the fluctuation field $B$ transforms by the linear action induced by $u$ on fluctuation fields.

Write $d\mu_{C^{(k)}}$ for the Gaussian measure with covariance $C^{(k)}(W)$, and $\mathcal{P}_k$ for the tilt of the step measure. The measure
\begin{equation*}
d\mu_{C^{(k)}}\,\chi_k\,e^{\mathcal{P}_k}
\end{equation*}
is invariant under this transformation; this is \cite[(2.16)]{B-CMP109}.

The integrand is built from the covariant objects listed above, so it transforms covariantly under $u$; in addition, the measure is invariant. Together with the transformation laws of the background and of $B$, these facts are exactly the covariance hypothesis of the lemma (d). The lemma therefore shows that $W\mapsto\mathfrak{w}(W)$ is gauge-equivariant. This is the equivariance hypothesis of the theorem (c).

\emph{Holomorphy and the bound.} Under the tilted law, $\mathfrak{w}(W)$ is a ratio of two box integrals, the denominator being the normalising integral. By (a) and (b), the integrands are holomorphic in $W$ on the complex window. By the same two statements, the normalising integral stays bounded away from $0$ for $\vartheta\le\frac14$. The ratio $\mathfrak{w}(W)$ is therefore holomorphic on the complex window. For a constant $C_1$ it satisfies
\begin{equation*}
\sup|\mathfrak{w}|\le C_1\,n_{\mathfrak{B}}\,\mathfrak{b}_k^{3}\,m_1=:M_{\mathfrak{w}} .
\end{equation*}
Here $n_{\mathfrak{B}}$ and $m_1$ are the quantities entering the bounds of (a) and (b). Holomorphy holds on a domain whose radii are of the type of $\varrho_{\mathcal{P}}$ and depend on $(L,\mathfrak{p},N)$ only. This is the holomorphy hypothesis of the theorem (c).

\emph{Finite dependence.} The value $\mathfrak{w}(W)$ depends on $W$ only through the links of a box of side of order $D$. This is the finite-dependence hypothesis of the theorem (c).

\emph{Conclusion from the theorem (c).} The three hypotheses of (c) hold, so two consequences follow. First, $\mathfrak{w}(W)=0$ at flat $W$. Second, for $W$ in the complex window,
\begin{equation}\label{8.2:eq-one-line-equivariance-1}
\|\mathfrak{w}(W)\|\le K_{\mathfrak{w}}\,\kappa_{\mathfrak{B}}(W).
\end{equation}
Here $\kappa_{\mathfrak{B}}(W)$ is the curvature size of the box from Lemma~\ref{8.2:lem-one-step-extraction}. The constant $K_{\mathfrak{w}}$ equals $M_{\mathfrak{w}}$ times the maximum of the geometric constants of (c) for this box. By (c), each of these constants is at most linear in the side of the box, which is of order $D$. One of these constants is $(d-1)D/\sin(\pi/N)$. Hence, for a constant $C_2$,
\begin{equation*}
K_{\mathfrak{w}}\le C_2\,D\,n_{\mathfrak{B}}\,\mathfrak{b}_k^{3}\,m_1 .
\end{equation*}

\emph{The hypotheses of (c) on the group, the domain and analyticity.} It remains to check the restrictions of the theorem (c) concerning the centre of the group, the topology of the domain and analyticity.
\begin{itemize}
\item The structure group is $\mathrm{SU}(N)$ (centre-free).
\item The box is simply connected, and the exterior configuration is extended over it.
\item The analyticity used is that of Ba{\l}aban's construction. The statements on the scope of the theorem (c) that concern analyticity show the following. The factors $\chi$ that are not analytic are constant on the real window. They therefore factor out of the ratio defining $\mathfrak{w}$, up to the tail terms of Lemma~\ref{8.2:lem-one-step-extraction}.
\end{itemize}
Hence all hypotheses of the theorem (c) hold, and \eqref{8.2:eq-one-line-equivariance-1} holds for each $\mathfrak{w}\in\{\mathfrak{t}_k,\mathfrak{d}_k,\mathfrak{l}_k\}$. Consider the term carrying the factor $\kappa_{\mathrm{bg}}$ in the collection of the terms, part~\ref{8.2:prf-step-collection} of the proof of Lemma~\ref{8.2:lem-one-step-extraction}. Inserting \eqref{8.2:eq-one-line-equivariance-1} into that term, for each such $\mathfrak{w}$, gives clause~(iv).
\end{proof}

\begin{proof}[Proof of Lemma~\ref{8.2:lem-one-step-extraction}\textup{(v)}: the one-step conditional covariance]
\label{8.2:prf-one-step-covariance}
This part of the proof derives part~(v) from the following inputs, together with the
hypotheses of Lemma~\ref{8.2:lem-one-step-extraction}:
\begin{itemize}
\item[(a)] the covariance form of the expansion of the tilted box measure. For two
functions of the fluctuation $B$ on the box it expresses the conditional covariance through
four groups of terms, each of which is bounded below: the untilted covariance
$\operatorname{Cov}_0(\mathbf{G}_1,\mathbf{G}_2)$;
a term of first order in the box tilt $P=P^{+}+P^{-}$, given by
$\operatorname{Cov}_0(\hat{\mathbf{G}}_1\hat{\mathbf{G}}_2,P)$; a product term given by
$\operatorname{Cov}_0(\mathbf{G}_1,P)\operatorname{Cov}_0(\mathbf{G}_2,P)$; and a remainder
$\rho^{\mathrm{cov}}$. The remainder is bounded by
$6\vartheta_k^2E_0|\hat{\mathbf{G}}_1\hat{\mathbf{G}}_2|$ plus terms of third order in
$\vartheta_k$;
\item[(b)] the bound $\|DG(\mathrm{bg})\|\le C_{\mathrm{ax}}\mu_2\kappa_{\mathrm{bg}}$ on the
first differential of a gauge-invariant local observable at a background of curvature size
$\kappa_{\mathrm{bg}}$;
\item[(c)] for the sharpened form of the $\kappa$-linear member only, the gauge equivariance
of the one-line vectors.
\end{itemize}
Here $E_0$ and $\operatorname{Cov}_0$ denote expectation and covariance under the untilted step
measure. For a function $\mathbf{F}$ of $B$ we write $\hat{\mathbf{F}}:=\mathbf{F}-E_0\mathbf{F}$
for its centring.
Throughout, $\mathrm{poly}$ denotes a polynomial prefactor in $D$ and $\log x_k$ whose value may
change from one occurrence to the next.

\emph{Set-up.} For $a=1,2$ let $\mathrm{bg}_a:=V^{(k)}(W)\restriction\mathfrak{B}_\ell(p_{a,k})$
and $\kappa_{\mathrm{bg}_a}:=\kappa(\mathrm{bg}_a)$. Both supports lie in the one box
$\mathfrak{B}$. We put $\mathbf{G}_a:=G_a\circ\mathfrak{v}_k$, regarded as a function of $B$, and
apply input~(a) to the pair $(\mathbf{G}_1,\mathbf{G}_2)$.

For each $a$ we recall the jet expansion of $\mathbf{G}_a$ obtained earlier in the proof of
Lemma~\ref{8.2:lem-one-step-extraction}. The field on the box is expanded as
$Y=g_kY_1+g_k^2Y_2+g_k^3Y_3+\dots$, where $Y_1=iCB$ is odd in $B$, $Y_2$ is even, and
$Y_j$ has parity $(-1)^j$. The odd part $\mathbf{G}_a^{-}$ is $g_kDG_a(\mathrm{bg}_a)[Y_1]$ plus
terms of order $g_k^3$ and higher, and
\begin{align*}
\mathbf{G}_a^{+}-G_a(\mathrm{bg}_a)&=g_k^2\mathbf{G}_{a,2}+\mathbf{G}^{+}_{a,\ge4},\qquad
\mathbf{G}_{a,2}:=DG_a(\mathrm{bg}_a)[Y_2]+\tfrac12D^2G_a(\mathrm{bg}_a)[Y_1,Y_1],\\
|\mathbf{G}^{+}_{a,\ge4}|&\le Cg_k^4\mu\bigl(\|CB\|^4+\|CB\|^2\bigr).
\end{align*}
We write $\lambda_a:=\|DG_a(\mathrm{bg}_a)\|$. By input~(b),
$\lambda_a\le C_{\mathrm{ax}}\mu_2\kappa_{\mathrm{bg}_a}$. The Gaussian moments
$m_j:=E_\gamma\|CB\|^j_{\mathfrak{B}_{\ell}'}$ are taken on the bounded enlargement
$\mathfrak{B}_\ell'$ of the box on which $Y$ reads $CB$; they are bounded as in the earlier
step of the proof, and every moment $m_j$ is absorbed into $\mathrm{poly}$ below.

\emph{The untilted covariance.} We sort $\operatorname{Cov}_0(\mathbf{G}_1,\mathbf{G}_2)$ by
parity in $B$. The untilted measure is even in $B$. Hence $E_0\mathbf{G}_a^{-}=0$, the mixed
products $\mathbf{G}_1^{\pm}\mathbf{G}_2^{\mp}$ have vanishing expectation, and
\begin{equation*}
\operatorname{Cov}_0(\mathbf{G}_1,\mathbf{G}_2)=E_0[\mathbf{G}_1^{-}\mathbf{G}_2^{-}]
+\operatorname{Cov}_0(\mathbf{G}_1^{+},\mathbf{G}_2^{+}).
\end{equation*}

For the odd part, we insert the expansions of $\mathbf{G}_a^{-}$. The product of the two
leading terms gives $g_k^2E_\gamma[DG_1(\mathrm{bg}_1)[Y_1]\,DG_2(\mathrm{bg}_2)[Y_1]]$ after
$E_0$ is replaced by $E_\gamma$, up to tail contributions. Here and below, the tail
contributions are those collected at the end into the last member
$M^2(e^{-c_{\mathrm{loc}}(D-2\ell)}+t_k^{1/2})$ of the bound on $\mathcal{E}^{(2)}_k$.
The cross terms between a leading
term and a term of order $g_k^3$ are bounded by $Cg_k^4\lambda\mu m_4$, with $\lambda$ the
corresponding $\lambda_a$. Since $Y_1=iCB$, the leading expectation is by definition the
one-line pairing. Input~(b) converts $\lambda_a$ into
$C_{\mathrm{ax}}\mu_2\kappa_{\mathrm{bg}_a}$, and so
\begin{equation}\label{8.2:eq-one-step-covariance-1}
E_0[\mathbf{G}_1^{-}\mathbf{G}_2^{-}]
=g_k^2\bigl\langle DG_1(\mathrm{bg}_1),C^{(k)}(W)DG_2(\mathrm{bg}_2)\bigr\rangle
+O\bigl(g_k^4\mu_2\mu(\kappa_{\mathrm{bg}_1}+\kappa_{\mathrm{bg}_2})\,\mathrm{poly}\bigr),
\end{equation}
again up to tail contributions.

For the even part, the constants $G_a(\mathrm{bg}_a)$ do not contribute to the covariance. The
leading terms $g_k^2\mathbf{G}_{a,2}$ give, after the same replacement of $E_0$ by $E_\gamma$,
\begin{equation*}
\operatorname{Cov}_0(\mathbf{G}_1^{+},\mathbf{G}_2^{+})
=g_k^4\operatorname{Cov}_\gamma(\mathbf{G}_{1,2},\mathbf{G}_{2,2})
+O(g_k^6\mu^2m_6)
\end{equation*}
up to tail contributions. Expanding $\mathbf{G}_{a,2}$,
\begin{equation*}
\operatorname{Cov}_\gamma(\mathbf{G}_{1,2},\mathbf{G}_{2,2})
=\tfrac14\operatorname{Cov}_\gamma\bigl(D^2G_1(\mathrm{bg}_1)[Y_1,Y_1],
D^2G_2(\mathrm{bg}_2)[Y_1,Y_1]\bigr)+r_{12}.
\end{equation*}
Here $r_{12}$ consists of the terms with one factor $DG_a(\mathrm{bg}_a)[Y_2]$, each
bounded by $C\lambda_a\mu_2m_4$. By input~(b), their total contribution
$g_k^4r_{12}$ is of the size of the $\kappa$-linear member.

\emph{The two-line term.} For a real centred Gaussian vector $\xi$ with covariance $\Sigma$ and
symmetric bilinear forms $q_1,q_2$, the Isserlis formula gives
\begin{equation*}
\operatorname{Cov}\bigl(q_1(\xi,\xi),q_2(\xi,\xi)\bigr)=2\operatorname{Tr}(q_1\Sigma q_2\Sigma),
\end{equation*}
that is, twice the sum over the pairs of contractions. We apply it with $\xi=CB$,
$B\sim N(0,C^{(k)}(W))$, and $q_a=D^2G_a(\mathrm{bg}_a)$. The identity
$D^2G_a[iCB,iCB]=-D^2G_a[CB,CB]$ introduces one sign in each argument, and the two signs
cancel in the covariance. Therefore
\begin{equation*}
\tfrac14g_k^4\operatorname{Cov}_\gamma\bigl(D^2G_1[Y_1,Y_1],D^2G_2[Y_1,Y_1]\bigr)
=\tfrac12g_k^4\sum_{\text{pairs of $C^{(k)}$-contractions}}
\operatorname{Hess}G_1\otimes\operatorname{Hess}G_2,
\end{equation*}
which is the two-line term of part~(v).

\emph{First order in $P$.} We sort $\operatorname{Cov}_0(\hat{\mathbf{G}}_1\hat{\mathbf{G}}_2,P)$
by parity. Only products of total even parity survive, and we treat them in three groups.

\emph{Even times even against $P^{+}$.} The centred product
$\hat{\mathbf{G}}_1^{+}\hat{\mathbf{G}}_2^{+}$ against $P^{+}$ is bounded by
$2\vartheta^{+}_kE_0|\hat{\mathbf{G}}_1^{+}\hat{\mathbf{G}}_2^{+}|$, and
\begin{equation*}
2\vartheta^{+}_kE_0|\hat{\mathbf{G}}_1^{+}\hat{\mathbf{G}}_2^{+}|
\le Cg_k^2\cdot g_k^4\,\mathrm{poly}=g_k^6\,\mathrm{poly}.
\end{equation*}

\emph{Odd times odd against $P^{+}$.} The term
$\hat{\mathbf{G}}_1^{-}\hat{\mathbf{G}}_2^{-}P^{+}$ is bounded by
$Cg_k^2\cdot g_k^2\lambda_1\lambda_2\,\mathrm{poly}$. By input~(b) this is
$g_k^4\mu_2^2\kappa_{\mathrm{bg}_1}\kappa_{\mathrm{bg}_2}\,\mathrm{poly}$.

\emph{Mixed terms against $P^{-}$.} The term
$(\hat{\mathbf{G}}_1^{+}\hat{\mathbf{G}}_2^{-}+\hat{\mathbf{G}}_1^{-}\hat{\mathbf{G}}_2^{+})P^{-}$
is bounded by
\begin{equation*}
C\cdot g_k^2\mu_2\cdot g_k\lambda\cdot g_kn_{\mathfrak{B}}\mathfrak{b}_k^3\,\mathrm{poly}
=g_k^4\kappa\,\mathrm{poly},
\end{equation*}
where $\lambda$ and $\kappa$ stand for the $\lambda_a$ and $\kappa_{\mathrm{bg}_a}$ of the odd
factor, $n_{\mathfrak{B}}$ is the box factor of the bound on the odd part $P^{-}$ of the tilt,
and $n_{\mathfrak{B}}\mathfrak{b}_k^3$ is absorbed into $\mathrm{poly}$. This crude bound is of
the shape of the $\kappa$-linear member.

\emph{The sharpened form.} The leading part of $\hat{\mathbf{G}}_2^{-}$ is
$g_kDG_2(\mathrm{bg}_2)[Y_1]$. Hence the first mixed summand, at leading order, is the pairing
of the vector $E_0[\,Y_1\,\hat{\mathbf{G}}_1^{+}P^{-}]$ with $DG_2(\mathrm{bg}_2)$. By input~(c)
this vector is equivariant in $W$, and is therefore $O(\kappa_{\mathfrak{B}})$. Combined with
$\lambda_2\le C_{\mathrm{ax}}\mu_2\kappa_{\mathrm{bg}_2}$, this gives a bound of the form
$\kappa_{\mathrm{bg}}\kappa_{\mathfrak{B}}$ in place of $\kappa$. The second mixed summand is
the first with the indices $1,2$ exchanged. The remaining contributions to the
$\kappa$-linear member, listed in the collection below, are brought to the same form by
input~(c) applied as in the proof of part~(iv) of Lemma~\ref{8.2:lem-one-step-extraction}.

\emph{The product term.} Each factor satisfies
\begin{equation*}
|\operatorname{Cov}_0(\mathbf{G}_a,P)|\le C(g_k^2\lambda_a+g_k^4)\,\mathrm{poly}.
\end{equation*}
Multiplying out the two factors,
\begin{equation*}
|\operatorname{Cov}_0(\mathbf{G}_1,P)\operatorname{Cov}_0(\mathbf{G}_2,P)|
\le C\bigl(g_k^4\lambda_1\lambda_2+g_k^6(\lambda_1+\lambda_2)+g_k^8\bigr)\mathrm{poly}.
\end{equation*}
By input~(b) this is $g_k^4\kappa^2\,\mathrm{poly}$ plus terms of order $g_k^6$.

\emph{The remainder $\rho^{\mathrm{cov}}$.} Inserting the jet expansions of
$\hat{\mathbf{G}}_1$ and $\hat{\mathbf{G}}_2$ into the product and absorbing the Gaussian
moments into $\mathrm{poly}$ as before,
\begin{align*}
6\vartheta_k^2E_0|\hat{\mathbf{G}}_1\hat{\mathbf{G}}_2|
&\le Cg_k^2\,\mathrm{poly}\cdot(g_k\lambda_1+g_k^2\mu_2)(g_k\lambda_2+g_k^2\mu_2)\,\mathrm{poly}\\
&=g_k^4\kappa^2\,\mathrm{poly}+g_k^5\kappa\,\mathrm{poly}+g_k^6\,\mathrm{poly},
\end{align*}
where the last line again uses input~(b). The terms of third order in $\vartheta_k$ are
bounded in the same way.

\emph{Collection.} Adding the contributions gives the assertion of part~(v):
\begin{itemize}
\item The one-line term is the principal part of \eqref{8.2:eq-one-step-covariance-1}.
\item The two-line term is the Isserlis contraction sum.
\item The errors of order $g_k^6$ form the member $g_k^6\Pi_2\mu^2$ of the bound on
$\mathcal{E}^{(2)}_k$. These are $O(g_k^6\mu^2m_6)$, the even-even first-order term, the
$g_k^6$ parts of the product term and of $\rho^{\mathrm{cov}}$, and the terms of third order in
$\vartheta_k$.
\item The errors linear in $\kappa_{\mathrm{bg}_1}+\kappa_{\mathrm{bg}_2}$ form the member
$g_k^4\Pi_2'\mu_2\mu(\kappa_{\mathrm{bg}_1}+\kappa_{\mathrm{bg}_2})$. These are the cross terms
in \eqref{8.2:eq-one-step-covariance-1}, the terms $g_k^4r_{12}$, the mixed
first-order terms, and the $g_k^5\kappa$ part of $\rho^{\mathrm{cov}}$ (which carries one
further power of $g_k$). In the form using input~(c), applied as in the proof of part~(iv),
this member becomes $\kappa_{\mathrm{bg}}\kappa_{\mathfrak{B}}$.
\item The errors quadratic in the curvature form the member
$g_k^4\Pi_2''\mu_2^2\kappa_{\mathrm{bg}_1}\kappa_{\mathrm{bg}_2}$. These are the odd-odd
first-order term, the $g_k^4\kappa^2$ part of the product term, and the $g_k^4\kappa^2$ part of
$\rho^{\mathrm{cov}}$.
\item The tail contributions form the last member
$M^2(e^{-c_{\mathrm{loc}}(D-2\ell)}+t_k^{1/2})$.
\end{itemize}
\end{proof}

\begin{remark}[The one-loop term at flat background]\label{8.2:rem-one-loop-flat}
Let $k_0$, $r$, $k\in[k_0,k_0+r)$, $G\in\mathfrak{G}_k(p_k;\ell,M,\varrho)$, the box $\mathfrak{B}$
and $W$ be as in
the one-step extraction lemma (Lemma~\ref{8.2:lem-one-step-extraction}). As there, $\mathbf{G}$
denotes $G$ composed with the one-step field map $\mathfrak{v}_k$, regarded as a function of the
fluctuation $B$ on the box, and $\langle\mathbf{G}\rangle$ its expectation under the step measure.
We write $\mathbf{G}_j$ for the jet of $\mathbf{G}$ of order $g_k^j$, with parts even and odd in the
fluctuation $B$ denoted
$\mathbf{G}_j^{\mathrm{even}}$, $\mathbf{G}_j^{\mathrm{odd}}$. We write $P=P^{+}+P^{-}$ for the
parity decomposition of the box tilt $\mathcal{P}_{k,\mathfrak{B}}$, and put $P_1^{-}$ for the
part of $P^{-}$ of order $g_k$ and $P_2^{+}$ for the part of $P^{+}$ of order $g_k^2$.
Finally, $E_\gamma$ and $\operatorname{Cov}_\gamma$ denote expectation and covariance under the
centred Gaussian law $N(0,C^{(k)}(W))$ on the box.
We write $n_{\mathfrak{B}}$ for the constant, depending on the box $\mathfrak{B}$, through which the
smallness parameters $\vartheta_k$, $\vartheta_k^{+}$ of the odd and even parts of the tilt are
bounded in the proof of Lemma~\ref{8.2:lem-one-step-extraction}; for the even part this bound reads
\begin{equation*}
\vartheta_k^{+}\le C\,g_k^2\,n_{\mathfrak{B}}\,\mathfrak{b}_k^4\,\bigl(1+B^{\sharp}\tau_*\bigr),
\end{equation*}
with $C$ a constant and $\mathfrak{b}_k$ the window radius of that lemma.

Suppose the background is flat, $\kappa_{\mathrm{bg}}=0$. In the expansion of
$\langle\mathbf{G}\rangle$ given by part (d) of the expansion lemma for the tilted step measure,
with $\mathbf{G}$ expanded to sixth order and $P$ to second order, the terms of total order $g_k^4$
are
\begin{equation}\label{8.2:eq-one-loop-flat-a}
\begin{split}
\mathfrak{L}_k^{\mathrm{flat}}(G)
:= g_k^4\cdot\biggl(\Bigl\{
&E_\gamma\bigl[\mathbf{G}_4^{\mathrm{even}}\bigr]
+\operatorname{Cov}_\gamma\bigl(\mathbf{G}_2,P_2^{+}\bigr)\\
&+E_\gamma\bigl[\mathbf{G}_3^{\mathrm{odd}}P_1^{-}\bigr]
+\tfrac12\operatorname{Cov}_\gamma\bigl(\mathbf{G}_2,(P_1^{-})^2\bigr)\Bigr\}\Big/g_k^4\biggr).
\end{split}
\end{equation}
The factor in parentheses is the coefficient of $g_k^4$.
The quantity $\mathfrak{L}_k^{\mathrm{flat}}(G)$ is a finite sum of Wick integrals of the vertices
of Ba{\l}aban's step against the jets of $G$. They comprise:
\begin{itemize}
\item the quartic Taylor term of $G$, and the square of the second-order shift of the field map;
\item $\operatorname{Hess}G$ against, in turn, the quartic vertex of the action, the Hessian
(second variation) vertex of the step, and the quadratic parts of the Jacobian and of the
accompanying functional of Ba{\l}aban's step;
\item the cubic jet of $G$ against the cubic vertex;
\item $\operatorname{Hess}G$ against two contracted cubic vertices.
\end{itemize}
The terms of orders $g_k$, $g_k^3$ and $g_k^5$ in $\langle\mathbf{G}\rangle$ vanish identically at
flat background. With $P$ expanded to second order, the remainder after the terms of order at most
four, of which $\mathfrak{L}_k^{\mathrm{flat}}(G)$ is the order-four part, is
\begin{equation}\label{8.2:eq-one-loop-flat-1}
O\bigl(g_k^6\,n_{\mathfrak{B}}^3\,\mathfrak{b}_k^{12}\,\mu\bigr)
\end{equation}
at flat background, where $\mu$ is the constant of the Cauchy bounds of the jets of $G$, as in
Lemma~\ref{8.2:lem-one-step-extraction}. At a curved background the same display holds, with the
$\kappa$-members of parts (iii) and (iv) of Lemma~\ref{8.2:lem-one-step-extraction} added.

The display \eqref{8.2:eq-one-loop-flat-a} is extracted as a formula only. No value of it is
computed or claimed, and it is not used in what follows. The one-loop coefficient of the $r$-step
expansion contains, besides contributions of the form
\eqref{8.2:eq-one-loop-flat-a}, terms from the running of the coupling across the steps and from
the linearisation of the backgrounds. No symmetry fixes its sign, and it is not asserted here to be
non-zero.
\end{remark}

\begin{proof}
We apply part (d) of the expansion lemma for the tilted step measure to $\langle\mathbf{G}\rangle$.
In this application $\mathbf{G}$ is expanded to sixth order in its jets and the tilt $P$ to second
order.

The expansion of $\mathbf{G}$ to sixth order is legitimate because $G$ is holomorphic in the tube
of radius $\varrho$. The jets of $G$ therefore obey the Cauchy bounds $\mu_j=j!M\varrho^{-j}$.

With these expansions, part (d) gives $\langle\mathbf{G}\rangle$ as a finite sum of terms of
definite total order in $g_k$ and a remainder. We sort the finite sum by total order in $g_k$.

Consider first the terms that depend on the curvature of the background. At flat background
$\kappa_{\mathrm{bg}}=0$, so the curvature parameter of part (d), which is bounded by a constant
times $\mu_2\kappa_{\mathrm{bg}}$, vanishes, and so do the
$\mathfrak{g}$-valued one-line vectors:
\begin{equation*}
\mathfrak{t}_k=\mathfrak{d}_k=\mathfrak{l}_k=0 .
\end{equation*}
Consequently none of the terms that carry these quantities contributes at any order.

Consider next the terms of odd total order $g_k$, $g_k^3$ or $g_k^5$. By part (b) of the same lemma,
at flat background each such term is a Gaussian integral of a polynomial of odd total degree in
$B$. Since $N(0,C^{(k)}(W))$ is centred, every such integral is zero. The terms of these orders
therefore vanish identically.

Consider finally the terms of total order four. Part (d), with $\mathbf{G}$ expanded to sixth order
and $P$ to second order, gives these at flat background as exactly the four terms of
\eqref{8.2:eq-one-loop-flat-a}. Writing each jet and each part of $P$
through the vertices of the step, every one of them is a Wick integral of the kind listed in the
remark.

All terms of total order six and higher, together with the remainder of part (d) when $P$ is taken
to second order, form the remainder. Part (d) bounds it at flat background by
\eqref{8.2:eq-one-loop-flat-1}.

At a curved background we carry out the same sorting without setting $\kappa_{\mathrm{bg}}=0$. The
terms carrying the curvature parameter and $\mathfrak{t}_k$, $\mathfrak{d}_k$, $\mathfrak{l}_k$ then
remain and are the $\kappa$-members of parts (iii) and (iv) of
Lemma~\ref{8.2:lem-one-step-extraction}.

The display thus follows from part (d) of the expansion lemma together with the parity statement
of part (b); no evaluation of the Wick integrals enters.
\end{proof}

\begin{definition}[Backgrounds and transported step laws]\label{8.2:def-transported-laws}
Fix the depth $s$ and the corresponding level $k_0=K-s$, the level $k_0+r$, the plaquette $p$ at
depth $s$, and the side lengths $\ell$ and $D$ of the
boxes $\mathfrak{B}_\ell(\cdot)$ and $\mathfrak{B}_D(\cdot)$. Write $\mathfrak{g}=\mathfrak{su}(N)$ and
$n=\dim\mathfrak{g}$. For each level $j$, $V^{(j)}$ denotes the one-step background map, which takes a
level-$(j+1)$ configuration to the level-$j$ background determined by it; it enters the one-step
extraction lemma (Lemma~\ref{8.2:lem-one-step-extraction}) as $V^{(k)}$. All backgrounds and laws below are
determined by a level-$(k_0+r)$ field
\begin{equation*}
W:=V_{k_0+r}
\end{equation*}
in the level-$(k_0+r)$ window, as in Lemma~\ref{8.2:lem-one-step-extraction} with $k=k_0+r-1$.

\emph{The background at depth $s$.} Let $\mathfrak{B}^{(r)}$ be the $r$-step minimiser. Here and below
the letter $\mathfrak{B}$ with a superscript $(i)$ or with subscripts $i,r$ denotes a background map, and
with a subscript $\ell$ or $D$ and a plaquette argument it denotes a box. We put
\begin{equation}\label{8.2:eq-transported-laws-1}
\mathrm{bg}_r(W):=\mathfrak{B}^{(r)}(W)\restriction\mathfrak{B}_\ell(p),
\end{equation}
the depth-$s$ image of the $r$-step minimiser of $W$, restricted to the box $\mathfrak{B}_\ell(p)$ of side $\ell$
about $p$, in the sense of the first part of the statement on the depth-$s$ image of the minimiser.

\emph{Intermediate backgrounds and linear transport.} For $0\le i<r$ we put
\begin{equation}\label{8.2:eq-transported-laws-2}
W_i:=\mathfrak{B}_{i+1,r}(W):=V^{(k_0+i+1)}\circ\cdots\circ V^{(k_0+r-1)}(W),
\end{equation}
the level-$(k_0+i+1)$ background determined by $W$ (for $i=r-1$ the composition is empty and
$W_{r-1}=W$). With the same convention
\begin{equation*}
\mathfrak{B}_{i,r}(W)=V^{(k_0+i)}\circ\cdots\circ V^{(k_0+r-1)}(W)
\end{equation*}
is the level-$(k_0+i)$ background determined by $W$. The \emph{$i$-fold linear transport to depth $s$} is the
differential of the $i$-step minimiser at this background:
\begin{equation}\label{8.2:eq-transported-laws-3}
\mathfrak{L}^{(i)}_W:=D\mathfrak{B}^{(i)}\bigl(\mathfrak{B}_{i,r}(W)\bigr),
\end{equation}
in the sense of the third part of the statement on the depth-$s$ image of the minimiser.

\emph{The transported step laws.} For $0\le i<r$ let $B_i$ be distributed according to the centred Gaussian
law $N\bigl(0,C^{(k_0+i)}(W_i)\bigr)$, where $C^{(k_0+i)}(W_i)$ is the covariance of the step at level $k_0+i$ in
the background $W_i$, and let $iCB_i$ be the field formed from $B_i$ as in
Lemma~\ref{8.2:lem-one-step-extraction}; the factor $i$ multiplying $CB_i$ is the factor $i$ of the field
$iCB$ in that lemma and is unrelated to the index $i$. The \emph{transported step law} $\Gamma_i(W)$ is the
centred Gaussian law on $\mathfrak{g}^{\mathfrak{B}_\ell(p)}$ of the field
\begin{equation}\label{8.2:eq-transported-laws-4}
\mathfrak{L}^{(i)}_W\bigl(iCB_i\bigr)\restriction\mathfrak{B}_\ell(p).
\end{equation}

\emph{Step coefficients.} Extending the notation of Lemma~\ref{8.2:lem-one-step-extraction}, where the
subscript of the trace is a covariance, for a centred Gaussian law $\Gamma$ on
$\mathfrak{g}^{\mathfrak{B}_\ell(p)}$, a background $b$ and an
observable $G$ we write $\operatorname{Tr}_{\Gamma}[\operatorname{Hess}_b G]$ for
the expectation of the second differential $D^2G(b)[X,X]$ over a
field $X$ distributed according to $\Gamma$. Let $O_{s,p}$ be the curvature observable at depth $s$ at the
plaquette $p$. For $0\le i<r$ we put
\begin{equation}\label{8.2:eq-transported-laws-5}
\mathfrak{g}^{\mathrm{step}}_{k_0,i}(W)
:=n^{-1}\operatorname{Tr}_{\Gamma_i(W)}\bigl[\operatorname{Hess}_{\mathrm{bg}_r(W)}O_{s,p}\bigr],
\end{equation}
so that, for every $g'>0$, multiplying \eqref{8.2:eq-transported-laws-5} by $\tfrac{n}{2}g'^2$ gives
\begin{equation}\label{8.2:eq-transported-laws-6}
\tfrac12 g'^2\operatorname{Tr}_{\Gamma_i(W)}\bigl[\operatorname{Hess}_{\mathrm{bg}_r(W)}O_{s,p}\bigr]
=\tfrac{n}{2}\,g'^2\,\mathfrak{g}^{\mathrm{step}}_{k_0,i}(W).
\end{equation}
We write $\mathfrak{g}^{\mathrm{step}}_{k_0,i}$, without argument, for the value of
\eqref{8.2:eq-transported-laws-5} at a flat field $W$, that is, one with $W(\partial p'')=\mathbf{1}$ for
every plaquette $p''$.

\emph{The background value of the observable.} We put
\begin{equation}\label{8.2:eq-transported-laws-7}
\Theta_{\mathrm{bg}}(W):=O_{s,p}\bigl(\mathrm{bg}_r(W)\bigr)\ \ge\ 0 .
\end{equation}

\emph{Curvature sizes.} Let $p_{k_0+r}$ be the image of $p$ at level $k_0+r$, let
$\mathcal{N}(p_{k_0+r})$ be its neighbourhood and $B_3$ the constant, both as in the first part of the
statement on the depth-$s$ image of the minimiser. We put
\begin{equation}\label{8.2:eq-transported-laws-8}
\kappa_W:=B_3L^{-2r}\max_{p''\in\mathcal{N}(p_{k_0+r})}\bigl|W(\partial p'')-\mathbf{1}\bigr|,
\end{equation}
and we let $\kappa^{\mathfrak{B}}_W$ be the maximum of $|W(\partial p'')-\mathbf{1}|$ over the plaquettes $p''$
of the level-$(k_0+r)$ images of the boxes $\mathfrak{B}_D(p_k)$, $k_0\le k<k_0+r$, of side $D$ about the
images $p_k$ of $p$. For every such $W$,
\begin{equation}\label{8.2:eq-transported-laws-9}
\kappa^{\mathfrak{B}}_W\le 2\varepsilon_{k_0+r},
\end{equation}
where $\varepsilon_{k_0+r}$ is the small-field radius at level $k_0+r$.
\end{definition}

\begin{remark}\label{8.2:lem-extraction-one-point}
The statement of this lemma is not written out here; parts of its proof are printed below.
\end{remark}

\begin{proof}[Proof of Lemma~\ref{8.2:lem-extraction-one-point}, continued: one step applied to the
carried main part]\label{8.2:prf-carried-step}
We continue the step from stage $i$ to stage $i+1$ of the induction, for some $0\le i<r$. By the
inductive claim at stage $i$, established in the preceding steps of the proof, the main part of the
carried function on the clean event $\mathcal{C}_p(D)$ is the function of the level-$(k_0+i)$ field $V$
\begin{equation*}
\mathfrak{M}_i(V)=G\bigl(\mathfrak{B}^{(i)}(V)\bigr)
+\frac12\sum_{i'<i}g^2_{k_0+i'}\,
\operatorname{Tr}_{\Gamma^{(i)}_{i'}(V)}\bigl[\operatorname{Hess}_{\mathfrak{B}^{(i)}(V)}G\bigr],
\end{equation*}
with $\Gamma^{(i)}_{i'}(V)$ the law $\Gamma_{i'}$ of Definition~\ref{8.2:def-transported-laws} with $V$
in place of $W$ as the conditioning field at level $k_0+i$; the Cauchy bounds $\mu_j$ and $\mu$, the
gradient norm $\lambda^{(i)}:=\|DG(\mathfrak{B}^{(i)}V)\|$ and the support size $\ell_i$ are those of
the preceding step of the proof. Besides the one-step
extraction lemma, Lemma~\ref{8.2:lem-one-step-extraction}, this part of the proof uses part (iii) of
localisation of the sourced step, part (b) of the bounds on the composed background maps and their
linear transports, part (b) of the gradient bound for gauge-invariant observables at a background, and
the equivariance theorem for gauge-equivariant functions on a finite link set.

Let $V'$ be a field at level $k_0+i+1$ in its window. We apply parts (i)--(iii) of
Lemma~\ref{8.2:lem-one-step-extraction}, the corresponding display, to the
$D$-local truncation of $\mathfrak{M}_i$ at step $k_0+i$, with box $\mathfrak{B}_D(p_{k_0+i})$.
Part (i) gives
\begin{equation}\label{8.2:eq-carried-step-1}
\begin{split}
E\bigl[\mathfrak{M}_i\,\big|\,V'\bigr]
={}&\mathfrak{M}_i\bigl(V^{(k_0+i)}(V')\bigr)
+\frac12\,g^2_{k_0+i}\,\operatorname{Tr}_{C^{(k_0+i)}(V')}
\bigl[\operatorname{Hess}_{V^{(k_0+i)}(V')}\mathfrak{M}_i\bigr]\\
&+\bigl(\mathfrak{L}+\mathcal{E}^{(1)}\bigr)_{k_0+i}\bigl(\mathfrak{M}_i;V'\bigr).
\end{split}
\end{equation}
We treat the three terms on the right in turn.

\emph{First term.} The $(i+1)$-step background is the $i$-step background of the one-step background,
$\mathfrak{B}^{(i)}\circ V^{(k_0+i)}=\mathfrak{B}^{(i+1)}$. Moreover, by the definition of
$\Gamma^{(i)}_{i'}$, inserting $V=V^{(k_0+i)}(V')$ as conditioning field at level $k_0+i$ is the same
as inserting $V'$ as conditioning field at level $k_0+i+1$, that is,
$\Gamma^{(i)}_{i'}(V^{(k_0+i)}(V'))=\Gamma^{(i+1)}_{i'}(V')$. Hence
\begin{equation}\label{8.2:eq-carried-step-2}
\mathfrak{M}_i\circ V^{(k_0+i)}
=G\bigl(\mathfrak{B}^{(i+1)}(\cdot)\bigr)
+\frac12\sum_{i'<i}g^2_{k_0+i'}\,
\operatorname{Tr}_{\Gamma^{(i+1)}_{i'}}\bigl[\operatorname{Hess}_{\mathfrak{B}^{(i+1)}}G\bigr],
\end{equation}
and the right side of \eqref{8.2:eq-carried-step-2} consists of the first $i+1$ summands of
$\mathfrak{M}_{i+1}$.

\emph{Second term.} The Hessian of $\mathfrak{M}_i$ is
\begin{equation*}
\operatorname{Hess}(\mathfrak{M}_i)=\operatorname{Hess}\bigl(G\circ\mathfrak{B}^{(i)}\bigr)
+\sum_{i'<i}\bigl(\text{Hessian of the $i'$-th trace summand}\bigr),
\end{equation*}
where the $i'$-th trace summand is $O(g^2_{k_0+i'}\mu_2L^{-4i'})$ times a function whose Hessian is
$O(L^{-4i})$ by the holonomy argument that bounds the second jet of $G\circ\mathfrak{B}^{(i)}$ (the
invariant $G$ depends on the background only through at most $C_\ell$ plaquette holonomies, each
holomorphic with all derivatives $O(B_3L^{-2i})$).

For the first part we use the chain rule
\begin{equation*}
D^2\bigl(G\circ\mathfrak{B}^{(i)}\bigr)[h,h]
=D^2G\bigl[D\mathfrak{B}^{(i)}h,\,D\mathfrak{B}^{(i)}h\bigr]
+DG\bigl[D^2\mathfrak{B}^{(i)}[h,h]\bigr]
\end{equation*}
at the point $V^{(k_0+i)}(V')$, with $h=iCB$,
$C=C^{(k_0+i)}(V')$ and $B\sim N(0,C^{(k_0+i)}(V'))$; here and below the letter $i$ in $iCB$ is the
imaginary unit, not the induction index. Writing $\mathfrak{L}^{(i)}$ for
$D\mathfrak{B}^{(i)}$ at $V^{(k_0+i)}(V')$ and using
$\mathfrak{B}^{(i)}(V^{(k_0+i)}(V'))=\mathfrak{B}^{(i+1)}(V')$, we obtain
\begin{equation}\label{8.2:eq-carried-step-3}
\begin{split}
\frac12\,g^2_{k_0+i}\operatorname{Tr}_{C}\bigl[\operatorname{Hess}(G\circ\mathfrak{B}^{(i)})\bigr]
={}&\frac12\,g^2_{k_0+i}\,E_B\,D^2G\bigl(\mathfrak{B}^{(i+1)}V'\bigr)
\bigl[\mathfrak{L}^{(i)}iCB,\mathfrak{L}^{(i)}iCB\bigr]\\
&+\frac12\,g^2_{k_0+i}\,DG\bigl(\mathfrak{B}^{(i+1)}V'\bigr)
\bigl[E_B\,D^2\mathfrak{B}^{(i)}[iCB,iCB]\bigr].
\end{split}
\end{equation}
By the definition of $\Gamma_i$ in Definition~\ref{8.2:def-transported-laws} (the law of the
$i$-fold linear transport of $iCB_i$, restricted to $\mathfrak{B}_\ell(p)$, with the level-$(k_0+i+1)$
background as conditioning field, here $V'$), the first line of the right side of
\eqref{8.2:eq-carried-step-3} equals
\begin{equation*}
\frac12\,g^2_{k_0+i}\,\operatorname{Tr}_{\Gamma^{(i+1)}_i(V')}
\bigl[\operatorname{Hess}_{\mathfrak{B}^{(i+1)}(V')}G\bigr],
\end{equation*}
which is the last summand of $\mathfrak{M}_{i+1}$. The second line of \eqref{8.2:eq-carried-step-3}
is the term linear in the background curvature. Write $m_2$ for a constant bounding the second moment
of the box fluctuation $iCB$. By the bound $\lambda^{(i+1)}\le C_{\mathrm{ax}}\mu_2B_3L^{-2i-2}\kappa(V')$
on the gradient norm (part (b) of the gradient bound for gauge-invariant observables at a background,
applied at the background
$\mathfrak{B}^{(i+1)}V'$) and by $\|D^2\mathfrak{B}^{(i)}\|=O(L^{-2i})$, this term is bounded by
\begin{equation}\label{8.2:eq-carried-step-4}
g^2_{k_0+i}\,\lambda^{(i+1)}\,C\,L^{-2i}m_2
\le C\,g^2_{k_0+i}\,\mu_2\,C_{\mathrm{ax}}B_3\,L^{-4i-2}\,\kappa(V')\,m_2 .
\end{equation}
For the improved form of the remainder bound of Lemma~\ref{8.2:lem-extraction-one-point} (the case
$\sigma_1=2$), note in addition that the vector $E_B[D^2\mathfrak{B}^{(i)}[iCB,iCB]]$ is
gauge-equivariant (covariantly) as a function of $V'$, so by the equivariance theorem it is
$O(\kappa_{\mathfrak{B}}(V'))$, with $\kappa_{\mathfrak{B}}$ the curvature size over the box as in
Lemma~\ref{8.2:lem-one-step-extraction}; the term is therefore of second order in the curvature, of
order $\kappa(V')\,\kappa_{\mathfrak{B}}(V')$.

The Hessians of the trace summands contribute at most
\begin{equation}\label{8.2:eq-carried-step-5}
C\,g^2_{k_0+i}\sum_{i'<i}g^2_{k_0+i'}\,\mu_2\,L^{-4i'}m_2=O\bigl(g^4_{k_0}\mu_2\bigr)
\end{equation}
per step, the sum over $i'$ carrying the geometric factor $L^{-4i'}$. Indeed, the dependence of the
$i'$-th trace summand on $V$ enters through the covariance $C^{(k_0+i')}(\cdot)$ and the transport
$\mathfrak{L}^{(i')}$, which have $O(1)$ derivatives by part (iii) of localisation of the sourced step
and part (b) of the bounds on the composed background maps, and through $\operatorname{Hess}G$, which
has $O(L^{-2i})$ derivatives. Summed over the $r$ steps of the segment,
\eqref{8.2:eq-carried-step-5} gives $O(g^4_{k_0}\mu_2\,r)$, which is contained in the term
$g^{2+\sigma_1}_{k_0}\Pi_{r,D}$ of the remainder bound of Lemma~\ref{8.2:lem-extraction-one-point},
whose prefactor $\Pi_{r,D}$ is polynomial in $r$.

\emph{Third term.} We apply part (ii) of Lemma~\ref{8.2:lem-one-step-extraction},
the corresponding display, with $(M,\mu,\mu_2)$ replaced by the jets of
$\mathfrak{M}_i$,
\begin{equation*}
\mu^{(i)}\le C\mu,\qquad
\mu_2^{(i)}\le C\mu_2L^{-4i}+C\lambda^{(i)}L^{-2i}+Cg^2_{k_0}\mu_2,
\end{equation*}
with $\ell$ replaced by the support size $\ell_i$, and with the background curvature
\begin{equation*}
\kappa_{\mathrm{bg}}=\kappa\bigl(V^{(k_0+i)}(V')\restriction\mathfrak{B}_{\ell_i}(p_{k_0+i})\bigr)
\le B_3L^{-2}\kappa(V').
\end{equation*}
With $\Pi_1(\ell_i)$, $\Pi_1'(\ell_i)$ the prefactors of that display computed with $\ell_i$ in place
of $\ell$, and with $C'$ the constant of part (ii) of Lemma~\ref{8.2:lem-one-step-extraction}, this
gives
\begin{equation}\label{8.2:eq-carried-step-6}
\begin{split}
\bigl|\bigl(\mathfrak{L}+\mathcal{E}^{(1)}\bigr)_{k_0+i}(\mathfrak{M}_i;V')\bigr|
\le C'\Bigl[&\,g^4_{k_0+i}\Pi_1(\ell_i)\,\mu
+g^2_{k_0+i}\Pi_1'(\ell_i)\,\mu_2L^{-4i}B_3L^{-2}\kappa(V')\\
&+\operatorname{osc}(\mathfrak{M}_i)\bigl(e^{-c_{\mathrm{loc}}(D-\ell_i)}
+t_{k_0+i}^{1/2}\bigr)\Bigr],
\end{split}
\end{equation}
where the oscillation of $\mathfrak{M}_i$ over the window satisfies
\begin{equation*}
\operatorname{osc}(\mathfrak{M}_i)\le C\mu_2L^{-4i}\,\mathfrak{b}_{k_0+i}^{\,2}+Cg^2_{k_0}\mu_2 ,
\end{equation*}
with $\mathfrak{b}_{k_0+i}$ the radius of the level-$(k_0+i)$ window.

It remains to collect \eqref{8.2:eq-carried-step-2}--\eqref{8.2:eq-carried-step-6}, together with
the carried remainder, into the inductive claim at stage $i+1$; this is done in the next step of the
proof.
\end{proof}

\begin{proof}[Proof of Lemma~\ref{8.2:lem-extraction-one-point}, continued: recursion for the
extraction remainder]
\label{8.2:prf-remainder-recursion}
We keep the notation of the inductive claim. On the clean event $\mathcal{C}_p(D)$ the carried
function $G_i$ after $i$ clean steps is a function of $V=V_{k_0+i}$ in the level-$(k_0+i)$
window. It splits as $G_i=\mathfrak{M}_i+\mathfrak{J}_i$, where $\mathfrak{M}_i$ is the main part
and $\sup|\mathfrak{J}_i|\le J_i$, the supremum being taken over that window.

Constants are as follows.
\begin{itemize}
\item $C$ denotes constants depending only on $(L,\mathfrak{p},N,\ell)$ and the step constants,
as $C_{\mathrm{GX}}$ does.
\item $\Pi_1$, $\Pi_1'$, $c_{\mathrm{loc}}$ and $t_k$ are the polynomial prefactors, the
localisation rate and the tail term of the one-step extraction lemma,
the corresponding display of Lemma~\ref{8.2:lem-one-step-extraction}.
\item $m_2$ is the second-moment constant of the Gaussian law of the fluctuation of one
renormalisation step.
\item $B_3$ is the constant in the derivative bounds $O(B_3L^{-2i})$ of the background maps given
by parts~(a) and~(b) of the transport estimates for the background map.
\item $c_{\mathrm{OP}}$ is the rate of the localisation tails supplied by the localisation of the
sourced step.
\item $\ell_i\le \ell/L^i+i\ell_0+C$ is the support size of the first summand of
$\mathfrak{M}_i$.
\end{itemize}
All of these enter through the step of this proof in which one renormalisation step is applied to
the carried main part. The hypotheses used below are those of
Lemma~\ref{8.2:lem-extraction-one-point}.

\emph{The recursion.} The inductive claim at $i+1$ follows by collecting four contributions:
\begin{itemize}
\item the remainder $E_{k_0+i}[\mathfrak{J}_i\,|\,\cdot\,]$, carried in supremum norm;
\item the first, second and third terms of the one-step expansion of
$E_{k_0+i}[\mathfrak{M}_i\,|\,V']$ obtained in that step;
\item the contributions of the Hessians of the trace summands;
\item the localisation tails moved into the remainder.
\end{itemize}
The claim then holds with
\begin{equation}\label{8.2:eq-remainder-recursion-1}
\begin{split}
J_{i+1}:=J_i+C\Big[\,&g_{k_0+i}^4\,\Pi_1(\ell_i)\,\mu+g_{k_0}^4\,\mu_2m_2\\
&+g_{k_0+i}^2\,\mu_2\big(L^{-4i}+g_{k_0}^2\big)\big(\Pi_1'(\ell_i)+C\big)B_3L^{-2}
\sup_{V_{k_0+i+1}}\kappa(V_{k_0+i+1})\\
&+\mu\big(e^{-c_{\mathrm{loc}}(D-\ell_i)}+e^{-c_{\mathrm{OP}}D}+t_{k_0+i}^{1/2}\big)\Big],
\end{split}
\end{equation}
where the supremum runs over the level-$(k_0+i+1)$ window. We call the third member of the
bracket the \emph{$\kappa$-member} of step $i$.

\emph{The curvature in the window.} Let $\varepsilon_k$ denote the small-field parameter of
level $k$; the curvature in the level-$k$ window is at most $2\varepsilon_k$. We have
\begin{equation}\label{8.2:eq-remainder-recursion-2}
\sup_{V_{k_0+i+1}}\kappa(V_{k_0+i+1})\le 2\varepsilon_{k_0+i+1}\le 4g_{k_0}\,p_0(g_{k_0}),
\end{equation}
where $p_0(\cdot)$ is the degree function. The first inequality is the bound just stated, at
level $k_0+i+1$. The second uses the small-field threshold $\varepsilon_k\le 2g_k\,p_0(g_k)$ at
$k=k_0+i+1$, together with clause~(0), listed among the hypotheses of
Lemma~\ref{8.2:lem-extraction-one-point}, for the comparison of the running couplings at the
levels $k_0+i+1$ and $k_0$.

\emph{The $\kappa$-members for $i<r-1$.} Insert \eqref{8.2:eq-remainder-recursion-2} into the
$\kappa$-member of \eqref{8.2:eq-remainder-recursion-1}. The couplings $g_{k_0+i}$ are compared
with $g_{k_0}$ in the same way. Summing over $i<r-1$ gives
\begin{equation}\label{8.2:eq-remainder-recursion-3}
\sum_{i<r-1}(\text{$\kappa$-member of step }i)
\le C g_{k_0}^3\,p_0(g_{k_0})\,\mu_2\,\Pi_1'\sum_{i}L^{-4i-2}
\le C g_{k_0}^3\,\Pi'_{r,D}\,\mu_2 .
\end{equation}
In the second inequality the geometric sum is bounded by a constant, and $p_0(g_{k_0})\Pi_1'$ is
absorbed into $\Pi'_{r,D}$. The result is of order $g^3_{k_0}$. This is the place where the
crude form of the bound gives order $g_{k_0}^3$ rather than the order $g_{k_0}^4$ reached in
assertion~(iii), and it is the reason for the value $\sigma_1=1$ in
the corresponding display.

\emph{The $\kappa$-member of the last step.} At $i=r-1$ the conditioning field at level
$k_0+r$ is $W$ itself. Here we do not bound the curvature through
\eqref{8.2:eq-remainder-recursion-2}; it is kept as
$\max_{p''\in\mathcal{N}(p_{k_0+r})}|W(\partial p'')-\mathbf{1}|$. The factor $L^{-4i}$ at
$i=r-1$ is $L^{-4(r-1)}$. By the definition of $\kappa_W$, which carries the factor
$B_3L^{-2r}$, $\kappa_W$ is $B_3L^{-2r}$ times this maximum.
Hence the $\kappa$-member of the last step is
\begin{equation}\label{8.2:eq-remainder-recursion-4}
\begin{split}
g_{k_0+r-1}^2\,\mu_2\,\Pi_1'\,B_3L^{-4(r-1)-2}
\max_{p''\in\mathcal{N}(p_{k_0+r})}|W(\partial p'')-\mathbf{1}|
&=g_{k_0+r-1}^2\,\mu_2\,\Pi_1'\,L^{2-2r}\,\kappa_W\\
&\le 2g_{k_0}^2\,\Pi'_{r,D}\,\mu_2\,\kappa_W .
\end{split}
\end{equation}
The equality is the identity $B_3L^{-4r+2}=L^{2-2r}\cdot B_3L^{-2r}$. The inequality uses
$L^{2-2r}\le 1$, which holds because $r\ge1$, together with the comparison of $g_{k_0+r-1}$ with
$g_{k_0}$. Since $\mu_2=2M\varrho^{-2}$, the right side of
\eqref{8.2:eq-remainder-recursion-4} is the displayed $\kappa_W$-member
$g^2_{k_0}\Pi'_{r,D}\varrho^{-2}\kappa_W$ of the corresponding display, up to the
factor $C_{\mathrm{GX}}M$.

\emph{The members of order $g^4$.} Summing the first two members of the bracket in
\eqref{8.2:eq-remainder-recursion-1} over $i<r$ gives at most $C g_{k_0}^4\,\Pi_1\,\mu\cdot r$.
This is the only place where $r$ enters $\Pi_{r,D}$, through its last factor $r$. Every member
except the $\Pi_1$-term of the one-step extraction lemma itself carries a factor $L^{-4i}$, which
could be used to remove this factor $r$; we keep it.

\emph{The tails.} Summing the tail members of \eqref{8.2:eq-remainder-recursion-1} over $i<r$
gives
\begin{equation}\label{8.2:eq-remainder-recursion-5}
\sum_{i<r}\big(e^{-c_{\mathrm{loc}}(D-\ell_i)}+e^{-c_{\mathrm{OP}}D}+t_{k_0+i}^{1/2}\big)
\le r\big(e^{-c_{\mathrm{loc}}(D-\ell_r)}+e^{-c_w\mathfrak{b}_{k_0}^2/C}\big).
\end{equation}
We now compare the last exponential with the degree function. By definition
$\mathfrak{b}_{k_0}=A_1(\log x_{k_0})^{p_0}$ and $p_0(g_{k_0})=A_0(\log g_{k_0}^{-2})^{p_0}$, and
$x_{k_0}=g_{k_0}^{-2}$. Hence $(\log x_{k_0})^{2p_0}=A_0^{-2}p_0(g_{k_0})^2$, and with
$c:=c_w/C$ the constant in the exponent of \eqref{8.2:eq-remainder-recursion-5},
\begin{equation}\label{8.2:eq-remainder-recursion-6}
e^{-c\,\mathfrak{b}_{k_0}^2}=e^{-cA_1^2(\log x_{k_0})^{2p_0}}\le e^{-c_w\,p_0(g_{k_0})^2}.
\end{equation}

\emph{Conclusion at $i=r$.} By the definitions of the backgrounds and transported step laws
(Definition~\ref{8.2:def-transported-laws}), $\mathfrak{B}^{(r)}(W)\restriction\mathfrak{B}_\ell(p)
=\operatorname{bg}_r(W)$, and $\Gamma^{(r)}_{i'}=\Gamma_{i'}$ by the definition of
$\Gamma^{(i)}_{i'}$ in the inductive claim. At $i=r$ we therefore have
\begin{equation*}
G^{\langle h\rangle}=G_r=\mathfrak{M}_r+\mathfrak{J}_r .
\end{equation*}
Here $\mathfrak{M}_r(W)$ is the main term
\begin{equation*}
G(\operatorname{bg}_r(W))+\tfrac12\sum_{i=0}^{r-1}g_{k_0+i}^2
\operatorname{Tr}_{\Gamma_i(W)}[\operatorname{Hess}_{\operatorname{bg}_r(W)}G]
\end{equation*}
of the corresponding display. The remainder obeys $|\mathfrak{J}_r|\le J_r$. By
\eqref{8.2:eq-remainder-recursion-3}--\eqref{8.2:eq-remainder-recursion-6} and the bound on the
members of order $g^4$, $J_r$ is bounded by the right side of
the corresponding display with $\sigma_1=1$. With
$\mathcal{E}_r(G;W):=\mathfrak{J}_r(W)$ this proves the corresponding displays.

\emph{Assertion \textup{(iii)}.} We use the corresponding refinement of
the one-step extraction lemma. It rests on the gauge equivariance of the one-line vectors. With
it, every $\kappa$-member carries the product $\kappa(V')\,\kappa^{\mathfrak{B}}(V')$ in place of
$\kappa(V')$. For $i<r-1$ both factors are bounded as in \eqref{8.2:eq-remainder-recursion-2}, so
\begin{equation*}
\kappa(V')\,\kappa^{\mathfrak{B}}(V')\le\big(4g_{k_0}\,p_0(g_{k_0})\big)^2 .
\end{equation*}
The corresponding members are therefore of order $g_{k_0}^4\,p_0(g_{k_0})^2$ and are contained
in $g^4_{k_0}\Pi_{r,D}$. At $i=r-1$ the computation \eqref{8.2:eq-remainder-recursion-4} is
repeated with $\kappa_W\,\kappa^{\mathfrak{B}}_W$ in place of $\kappa_W$. Hence
the corresponding display holds with $\sigma_1=2$ and with its $\kappa$-member replaced
by $g_{k_0}^2\Pi''_{r,D}\varrho^{-2}\kappa_W\,\kappa^{\mathfrak{B}}_W$.

\emph{Assertion \textup{(i)}.} Let $n=\dim\mathfrak{g}$. For $G=O_{s,p}$ we use part~(c) of the
algebraic lemma on gauge-invariant holomorphic functions of a box configuration. It shows that at
the flat configuration the step coefficient is the per-colour second moment of the linearised
plaquette variable of the transported fluctuation. Part~(c) of the transport estimates for the
background map, a hypothesis of Lemma~\ref{8.2:lem-extraction-one-point}, bounds the linear
transport $\mathfrak{L}^{(i)}$, composed with the plaquette curl, by $B_{\mathrm{lin}}L^{-2i}$.
Let $E$ denote expectation in $B_i\sim N(0,C^{(k_0+i)}(W_i))$. The positivity bound on the step
covariance, a hypothesis of Lemma~\ref{8.2:lem-extraction-one-point}, bounds
$n^{-1}E\|C^{(k_0+i)}(W_i)B_i\|^2$ by a constant. Write $\mathfrak{g}^{\mathrm{step}}_{k_0,i}$
for the value of $\mathfrak{g}^{\mathrm{step}}_{k_0,i}(W)$ at the flat configuration
$W=\mathbf{1}$. Together,
\begin{equation}\label{8.2:eq-remainder-recursion-7}
\mathfrak{g}^{\mathrm{step}}_{k_0,i}
=n^{-1}E\big|\operatorname{curl}\big(\mathfrak{L}^{(i)}(\mathrm{i}\,C^{(k_0+i)}(W_i)B_i)\big)(p)
\big|^2\ \ge 0,
\end{equation}
\begin{equation}\label{8.2:eq-remainder-recursion-8}
\mathfrak{g}^{\mathrm{step}}_{k_0,i}
\le B_{\mathrm{lin}}^2L^{-4i}\,n^{-1}E\|C^{(k_0+i)}(W_i)B_i\|^2\le C_\eta L^{-4i}.
\end{equation}

Now consider the dependence on $W$. The map $W\mapsto\mathfrak{g}^{\mathrm{step}}_{k_0,i}(W)$ is
the trace of the gauge-invariant Hessian of $O_{s,p}$ at $\operatorname{bg}_r(W)$ against the
gauge-covariant Gaussian law $\Gamma_i(W)$. It is therefore a gauge-invariant holomorphic function
of $W$ on the box, and its supremum on the tube is at most $C_\eta L^{-4i}$, the order of the
flat bound \eqref{8.2:eq-remainder-recursion-8}. This gives
\begin{equation}\label{8.2:eq-remainder-recursion-9}
\big|\mathfrak{g}^{\mathrm{step}}_{k_0,i}(W)-\mathfrak{g}^{\mathrm{step}}_{k_0,i}\big|
\le C\cdot C_\eta L^{-4i}\cdot\big(\kappa^{\mathfrak{B}}_W\,C_{\mathrm{ax}}(D)\big)^{a},
\end{equation}
where $C_{\mathrm{ax}}(D)$ is the constant of part~(b) of the algebraic lemma at box side $D$,
at most $D$ times a polylogarithmic factor. The mean value theorem on the tube gives
\eqref{8.2:eq-remainder-recursion-9} with $a=1$, the linear bound; part~(b) of the algebraic
lemma at the scale of the box improves it to $a=2$, the quadratic bound.

\emph{Assertion \textup{(ii)}.} We use only two facts. The first is the exact identity
$|U-\mathbf{1}|_{\mathrm{HS}}^2=2N\big(1-\operatorname{Re}\operatorname{tr}U\big)$ for
$U\in\mathrm{SU}(N)$, with $\operatorname{tr}\mathbf{1}=1$. Indeed,
\begin{equation*}
|U-\mathbf{1}|_{\mathrm{HS}}^2=N\operatorname{tr}\big((U-\mathbf{1})^*(U-\mathbf{1})\big)
=2N-2N\operatorname{Re}\operatorname{tr}U .
\end{equation*}
The second is part~(a) of the transport estimates for the background map. They give
\begin{equation}\label{8.2:eq-remainder-recursion-10}
\begin{split}
\Theta_{\mathrm{bg}}(W)=O_{s,p}(\operatorname{bg}_r(W))
&=(2N)^{-1}\big|\operatorname{bg}_r(W)(\partial p)-\mathbf{1}\big|_{\mathrm{HS}}^2\\
&\le C_\eta L^{-4r}\max_{p''\in\mathcal{N}(p_{k_0+r})}
\big|W(\partial p'')-\mathbf{1}\big|_{\mathrm{HS}}^2\\
&=C_\eta L^{-4r}\cdot 2N\max_{p''\in\mathcal{N}(p_{k_0+r})}O_{s-r,p''}(W).
\end{split}
\end{equation}
The first equality is the definition of $\Theta_{\mathrm{bg}}$. The second uses
$O_{s,p}(U)=1-\operatorname{Re}\operatorname{tr}U(\partial p)$ and the identity above.
The inequality is part~(a) of the transport estimates, with the factor $(2N)^{-1}$ absorbed into
$C_\eta$. The last line is the identity again.
\end{proof}

\begin{remark}[Exactness of the Gaussian term]\label{8.2:rem-gaussian-exact}
Let $G$, $h$, $r$ and $W$ be as in the Gaussian extraction lemma for one local observable
(Lemma~\ref{8.2:lem-extraction-one-point}).
\begin{enumerate}
\item[(a)] The Gaussian term of the corresponding display, namely
\begin{equation*}
\frac12\sum_{i=0}^{r-1} g_{k_0+i}^{2}\,
\operatorname{Tr}_{\Gamma_i(W)}\bigl[\operatorname{Hess}_{\mathrm{bg}_r(W)}G\bigr],
\end{equation*}
is exact: it is the second-order Wick term of $G$ at the true $r$-step background
$\mathrm{bg}_r(W)$, taken against the sum of Ba{\l}aban's own step covariances
$C^{(k_0+i)}$, $0\le i\le r-1$, each transported by Ba{\l}aban's own linearised minimisers.
No model Gaussian is substituted for the step laws.
\item[(b)] Let $G=O_{s,p}$ be the plaquette observable of
Lemma~\ref{8.2:lem-extraction-one-point}, with its sup bound $M=M_O$ and with $\ell=1$. We
call $W$ flat if $\kappa_W=0$, that is, if $W(\partial p')=\mathbf{1}$ for every plaquette
$p'$ entering $\kappa_W$, and we write $(\mathrm{flat})$ for this argument. The value of the
Gaussian term at
flat $W$ is identified with the value $(n/2)\,\mathfrak{c}_{k_0}$, where
$n=\dim\mathfrak{su}(N)=N^2-1$, of the lemma on Ba{\l}aban's level-$k_0$ Gaussian
coefficient, up to the tail of the tower, only through two inputs. The first input is the
lemma summing the Gaussian tower of flat step coefficients, which gives
\begin{equation}\label{8.2:eq-gaussian-exact-1}
\sum_{i\ge 0}\mathfrak{g}^{\mathrm{step}}_{k_0,i}(\mathrm{flat})
=\mathfrak{g}_{k_0}=\mathfrak{c}_{k_0},
\qquad
\sum_{i\ge r}\mathfrak{g}^{\mathrm{step}}_{k_0,i}(\mathrm{flat})\le C_\eta L^{-4r},
\end{equation}
with the constant $C_\eta$ of that lemma. The second input is the
flat-background identification, which is printed without proof in Ba{\l}aban's papers and
is cited as printed.
\end{enumerate}
\end{remark}

\begin{proof}
(a) We unwind the definitions entering the corresponding display of Lemma~\ref{8.2:lem-extraction-one-point}. For each
$i$ with $0\le i\le r-1$, $\Gamma_i(W)$ is, by definition, the centred Gaussian law of
Ba{\l}aban's step at level $k_0+i$, of covariance $C^{(k_0+i)}$, transported by the
linearised minimisers. The symbol
$\operatorname{Tr}_{\Gamma_i(W)}[\operatorname{Hess}_{\mathrm{bg}_r(W)}G]$ denotes the
expectation, under $\Gamma_i(W)$, of the second differential of $G$ at $\mathrm{bg}_r(W)$.
The point $\mathrm{bg}_r(W)$ is the background produced by the $r$ steps themselves. Hence
each summand, weighted by $\tfrac12 g_{k_0+i}^2$, is the second-order Wick term of $G$ at
$\mathrm{bg}_r(W)$ for the transported covariance of step $k_0+i$. The sum over $i$ is the
second-order Wick term for the covariance $\sum_{i<r} g^2_{k_0+i}\,C^{(k_0+i)}$, each
summand transported by its minimiser. Every Gaussian datum entering the term (the laws
$\Gamma_i(W)$ and the point $\mathrm{bg}_r(W)$) is built from the covariances and
minimisers of Ba{\l}aban's step, so no model Gaussian is substituted.

(b) Let $G=O_{s,p}$, $M=M_O$, $\ell=1$. By the case $G=O_{s,p}$ of
Lemma~\ref{8.2:lem-extraction-one-point}, the main
term of the corresponding display is
\begin{equation*}
\Theta_{\mathrm{bg}}(W)+\frac n2\sum_{i<r} g_{k_0+i}^{2}\,
\mathfrak{g}^{\mathrm{step}}_{k_0,i}(W).
\end{equation*}
The Gaussian part of this expression is a finite weighted sum of the per-colour step
fluctuation coefficients $\mathfrak{g}^{\mathrm{step}}_{k_0,i}$. At flat $W$ it becomes a
weighted sum of the coefficients $\mathfrak{g}^{\mathrm{step}}_{k_0,i}(\mathrm{flat})$.

Neither the corresponding display of Lemma~\ref{8.2:lem-extraction-one-point} nor part~(a) relates the coefficients
$\mathfrak{g}^{\mathrm{step}}_{k_0,i}(\mathrm{flat})$ to $\mathfrak{c}_{k_0}$. Passing to
$(n/2)\,\mathfrak{c}_{k_0}$ requires exactly the two inputs named in part~(b):
\eqref{8.2:eq-gaussian-exact-1} supplies the summation of the full tower and the bound on
its tail beyond $r$ steps, and the flat-background identification supplies the remaining
step. Thus the value at
flat $W$ agrees with $(n/2)\,\mathfrak{c}_{k_0}$, up to the tail bounded in
\eqref{8.2:eq-gaussian-exact-1}, through these two inputs and through no other.
\end{proof}

\begin{definition}[Two smeared action densities in one box]\label{8.2:def-two-densities}
Let $\mathfrak{g}$ be the Lie algebra of $\mathrm{SU}(N)$, so that the constant $b_0$ of the
two-point witness theorem (see the glossary, Definition~\ref{0.2:not-glossary}) is
$b_0=(N^2-1)/2=\tfrac12\dim\mathfrak{g}$. Fix a depth $s$ and a level $k_0$ as in
Definition~\ref{8.2:def-transported-laws}, the constant $C_W=C_W(s)$, and two test functions
$f_1,f_2$.

\emph{The two densities.} For $a=1,2$, $X^{(k_0)}(f_a)$ denotes the smeared background action
density at level $k_0$. It is a gauge-invariant holomorphic function of the level-$k_0$
block field $V_{k_0}$, supported on the union of the depth-$s$ cells meeting
$\operatorname{supp}f_a$. The test function $f_a$ is taken so that $\operatorname{supp}f_a$
meets a number of depth-$s$ cells comparable to $C_W^4$.

\emph{Separation.} The supports of $X^{(k_0)}(f_1)$ and $X^{(k_0)}(f_2)$ are disjoint. Their
distance, measured in depth-$s$ cells, lies between $C_W$ and $6C_W$. This two-sided window is a
condition on the test functions $f_1,f_2$.

\emph{The kernel.} The letters $G^{(s)}_K$ and $\mathcal{K}_s$ are those of the two-point witness
theorem (see the glossary, Definition~\ref{0.2:not-glossary}). Here $G^{(s)}_K(p,q)$ is the
covariance of the plaquette fluxes through $p$ and through $q$ of the linearised level-$k_0$
background field, at unit coupling and flat background, when the fluctuation field from which
that background field is built has covariance equal to the sum of Ba{\l}aban's fluctuation
covariances of the levels $K-s,\dots,K$ at flat background, transported to level $K-s$; and
\begin{equation}\label{8.2:eq-two-densities-1}
\mathcal{K}_s(f_1,f_2)=\sum_{p,q}f_1(x(p))\,f_2(x(q))\,\bigl[G^{(s)}_K(p,q)\bigr]^2 .
\end{equation}

\emph{The target display.} The pair is the object for which the covariance identity of clause
(iii)(c) of Theorem~\ref{3.1:thm-main},
\begin{equation}\label{8.2:eq-two-densities-2}
\begin{gathered}
\operatorname{Cov}_{\nu^{B}_{K,s}}\bigl(X^{(k_0)}(f_1),X^{(k_0)}(f_2)\bigr)
=b_0\,g_{k_0}^4\,\mathcal{K}_s(f_1,f_2)\,(1+R_0),\\
|R_0|\le \tfrac1{16}+C\,g_{k_0}^2,
\end{gathered}
\end{equation}
in which $R_0$ is a relative error and $C$ a constant, is to be established under
Ba{\l}aban's block law $\nu^{B}_{K,s}$ at depth $s$. The identity
\eqref{8.2:eq-two-densities-2} is asserted for $C_W$ not smaller than a threshold independent
of $g$, and under an a priori bound on fourth moments under $\nu^{B}_{K,s}$, stated below with
the Gaussian extraction for the pair.

\emph{The pair as local observables.} In the notation of
Lemma~\ref{8.2:lem-extraction-one-point} we write
\begin{equation*}
G_a:=X^{(k_0)}(f_a)\in\mathfrak{G}_{k_0}(p_a;\ell_W,M_W,\varrho),\qquad a=1,2,
\end{equation*}
where $p_a$ is a plaquette chosen so that the box $\mathfrak{B}_{\ell_W}(p_a)$ contains
$\operatorname{supp}G_a$. Here $\ell_W$ is the support parameter of the class, comparable to
$C_W$, and $M_W$ is the bound for $|G_a|$ on the holomorphy tube of radius $\varrho$; $M_W$ is
left symbolic, being fixed by the normalisation in which the observables are used. The
separation condition reads
\begin{equation*}
\operatorname{dist}(\operatorname{supp}G_1,\operatorname{supp}G_2)\in[C_W,6C_W].
\end{equation*}

\emph{The joint clean event.} Let $\mathcal{C}_p(D)$, the step count $r$ and the constants
$\ell_0$, $D_0$ be as in the hypotheses of Lemma~\ref{8.2:lem-extraction-one-point}. We put
\begin{equation}\label{8.2:eq-two-densities-3}
\mathcal{C}_{p_1,p_2}(D):=\mathcal{C}_{p_1}(D)\cap\mathcal{C}_{p_2}(D),
\qquad D\ge 8C_W+r\ell_0+D_0 .
\end{equation}
Here both clean events are taken with respect to one common box $\mathfrak{B}_D$ of side $D$,
chosen once with
\begin{equation*}
\mathfrak{B}_D\supseteq\operatorname{supp}G_1\cup\operatorname{supp}G_2 .
\end{equation*}
\end{definition}

Display \eqref{8.2:eq-two-densities-2} is derived in two steps, both below: the Gaussian
extraction for the pair, over $r$ clean steps on $\mathcal{C}_{p_1,p_2}(D)$, and its corollary
for the pair $X^{(k_0)}(f_1),X^{(k_0)}(f_2)$.

In this derivation the covariance in \eqref{8.2:eq-two-densities-2} has a one-line and a
two-line contribution: the term in which the two observables are joined by one covariance line
and the term in which they are joined by two. Gauge invariance of $X^{(k_0)}(f_a)$ removes the
one-line contribution, not the two-line one.

\begin{proposition}[The Gaussian extraction]\label{8.2:prop-gaussian-extraction}
Assume the following.
\begin{itemize}
\item The hypotheses of Lemma~\ref{8.2:lem-extraction-one-point} (Gaussian extraction for one local
observable). Among them are the localisation of the sourced step and the condition on the range
of~$r$.
\item For the second form of the bound on the remainder in~(b): the equivariance hypothesis, under
which the corresponding display of Lemma~\ref{8.2:lem-extraction-one-point} holds with $\sigma_1=2$ and the members linear in the
curvature size of the remainder bounds of the one-step
extraction lemma, Lemma~\ref{8.2:lem-one-step-extraction}, acquire the factor
$\kappa_{\mathfrak{B}}$.
\item For the value in~(c): the flat-background hypothesis.
\end{itemize}

Let $G_a\in\mathfrak{G}_{k_0}(p_a;\ell,M,\varrho)$, $a=1,2$, have both supports in one box
$\mathfrak{B}_D$ as in Definition~\ref{8.2:def-two-densities}. Let
$h\in\mathcal{C}_{p_1,p_2}(D)$ be a history that is clean for $p_1$ and for $p_2$ with this one box,
and let $W$ lie in the level-$(k_0+r)$ window. Write $\mathrm{bg}_r(W)$ for the depth-$s$ image of
the $r$-step minimiser, restricted to $\mathfrak{B}_\ell(p_1)\cup\mathfrak{B}_\ell(p_2)$, as in
Definition~\ref{8.2:def-transported-laws} for $p=p_1$ and for $p=p_2$; the function $G_a$ and its
differentials read its restriction to $\mathfrak{B}_\ell(p_a)$. Write $\kappa_{W,a}$ for $\kappa_W$
of that definition taken at $p=p_a$.

For $0\le i<r$ let $\Gamma^{12}_i(W)$ be the joint centred Gaussian law on
$\mathfrak{g}^{\mathfrak{B}_\ell(p_1)\cup\mathfrak{B}_\ell(p_2)}$ of
$\mathfrak{L}^{(i)}_W(iCB_i)$ restricted to both supports, with $B_i$, $C$ and
$\mathfrak{L}^{(i)}_W$ as in Definition~\ref{8.2:def-transported-laws}. Its two marginals are the
laws $\Gamma_i(W)$ of that definition for $p=p_1$ and for $p=p_2$. Its cross-covariance between the
two supports carries the covariance line joining $p_1$ and $p_2$.

For a centred Gaussian law $\Gamma$ on this space we put
\begin{equation*}
\langle DG_1(\mathrm{bg}_r),\Gamma\,DG_2(\mathrm{bg}_r)\rangle
:=E_{\xi\sim\Gamma}\bigl[DG_1(\mathrm{bg}_r)[\xi]\cdot DG_2(\mathrm{bg}_r)[\xi]\bigr].
\end{equation*}
For two such laws $\Gamma,\Gamma'$ and the Hessians $H_a=\operatorname{Hess}_{\mathrm{bg}_r}G_a$
we put
\begin{equation*}
(H_1\otimes H_2)[\Gamma,\Gamma']:=E\bigl[H_1[\xi,\xi']\,H_2[\xi,\xi']\bigr],
\end{equation*}
where $\xi\sim\Gamma$ and $\xi'\sim\Gamma'$ are independent. This is the full contraction in which
one argument of $H_1$ is paired with one argument of $H_2$ through the cross-covariance of $\Gamma$,
and the remaining two arguments are paired through the cross-covariance of $\Gamma'$. Since $H_1$ and
$H_2$ are symmetric, $(H_1\otimes H_2)[\Gamma,\Gamma']=(H_1\otimes H_2)[\Gamma',\Gamma]$. Finally,
$n=\dim\mathfrak{g}$. The letters $C_1$, $C_\eta$ and $C_{\mathrm{ext},2}$ below denote positive
constants not depending on $W$, $h$, $k_0$, $r$ or $D$. Then the following hold.
\begin{enumerate}
\item[(a)] (Joint moment.) The product $G_1G_2$ lies in
$\mathfrak{G}_{k_0}(\cdot\,;\ell+6C_W,M^2,\varrho)$. The quantity $(G_1G_2)^{\langle h\rangle}(W)$
is given by Lemma~\ref{8.2:lem-extraction-one-point} applied to $G_1G_2$, that is, by
the corresponding display with the corresponding bound, with
jets up to fourth order. Its main term is the sum of two pieces: the Gaussian functional with the
structure displayed in~(b), and the product of the main terms of
the corresponding display of Lemma~\ref{8.2:lem-extraction-one-point} for $G_1$ and for $G_2$.
\item[(b)] (Truncated moment.) The carried conditional covariance
\begin{equation*}
\mathcal{T}^{(r)}(W):=(G_1G_2)^{\langle h\rangle}(W)-G_1^{\langle h\rangle}(W)\,
G_2^{\langle h\rangle}(W)
\end{equation*}
satisfies
\begin{equation}\label{8.2:eq-gaussian-extraction-a}
\begin{split}
\mathcal{T}^{(r)}(W)={}&\sum_{i<r}g_{k_0+i}^2\,
\bigl\langle DG_1(\mathrm{bg}_r),\Gamma^{12}_i(W)\,DG_2(\mathrm{bg}_r)\bigr\rangle\\
&+\frac12\sum_{i,i'<r}g_{k_0+i}^2g_{k_0+i'}^2
\bigl(\operatorname{Hess}_{\mathrm{bg}_r}G_1\otimes\operatorname{Hess}_{\mathrm{bg}_r}G_2\bigr)
\bigl[\Gamma^{12}_i,\Gamma^{12}_{i'}\bigr]
+\mathcal{E}_{2,r}(W).
\end{split}
\end{equation}
The first sum, the one-line term, consists of a single transported covariance line between the two
gradients. It vanishes at flat $W$, and it always satisfies
\begin{equation}\label{8.2:eq-gaussian-extraction-1}
\Bigl|\sum_{i<r}g_{k_0+i}^2\,
\bigl\langle DG_1(\mathrm{bg}_r),\Gamma^{12}_i(W)\,DG_2(\mathrm{bg}_r)\bigr\rangle\Bigr|
\le C_1\sum_{i<r}g_{k_0+i}^2L^{-4i}\,C_{\mathrm{ax}}^2\mu_2^2\,\kappa_{W,1}\kappa_{W,2},
\end{equation}
where $C_{\mathrm{ax}}$ and $\mu_2=2M\varrho^{-2}$ are as in
Lemma~\ref{8.2:lem-one-step-extraction}.

The second sum, the two-line term, consists of two transported covariance lines. Take
$G_a=O_{s,p_a}$, where $O_{s,p}(U):=1-\operatorname{Re}\operatorname{tr}U(\partial p)$ is the action
density of the plaquette $p$ at depth $s$, and take $W$ flat.
Then the two-line term equals
\begin{equation}\label{8.2:eq-gaussian-extraction-2}
\frac n2\Bigl(\sum_{i<r}g_{k_0+i}^2\,\mathfrak{g}^{\mathrm{step}}_{12,i}\Bigr)^2,
\end{equation}
where, for $i<r$, $\mathfrak{g}^{\mathrm{step}}_{12,i}$ is the flat-background covariance per colour
of the fluxes through $p_1$ and through $p_2$ of the $i$-th transported step fluctuation:
\begin{equation*}
\mathfrak{g}^{\mathrm{step}}_{12,i}:=n^{-1}E_{\xi\sim\Gamma^{12}_i(W)}
\bigl[\langle\ell_{p_1}(\xi),\ell_{p_2}(\xi)\rangle\bigr]\quad\text{at flat }W,
\end{equation*}
with $\ell_p(\xi)\in\mathfrak{g}$ the linearised flux of $\xi$ through the plaquette $p$.

The remainder satisfies
\begin{equation}\label{8.2:eq-gaussian-extraction-3}
\begin{split}
|\mathcal{E}_{2,r}(W)|\le{}&C_{\mathrm{ext},2}M^2\Bigl[g_{k_0}^6\Pi_2
+g_{k_0}^4\Pi_2'\bigl(\kappa_{W,1}+\kappa_{W,2}+g_{k_0}p_0(g_{k_0})\bigr)
+g_{k_0}^4\Pi_2''\kappa_{W,1}\kappa_{W,2}\Bigr]\\
&+M^2\bigl(e^{-c_{\mathrm{loc}}D}+e^{-c_{\mathrm{w}}p_0(g_{k_0})^2}\bigr).
\end{split}
\end{equation}
Relative to $g_{k_0}^4$, this is of order $g_{k_0}p_0(g_{k_0})$ and $\kappa_{W,a}$ times the
prefactors. Under the equivariance hypothesis the bound improves to
\begin{equation}\label{8.2:eq-gaussian-extraction-4}
\begin{split}
|\mathcal{E}_{2,r}(W)|\le{}&C_{\mathrm{ext},2}M^2\Bigl[g_{k_0}^6\Pi_2
+g_{k_0}^4\Pi_2'\bigl((\kappa_{W,1}+\kappa_{W,2})\kappa^{\mathfrak{B}}_W
+g_{k_0}^2p_0(g_{k_0})^2\bigr)
+g_{k_0}^4\Pi_2''\kappa_{W,1}\kappa_{W,2}\Bigr]\\
&+M^2\bigl(e^{-c_{\mathrm{loc}}D}+e^{-c_{\mathrm{w}}p_0(g_{k_0})^2}\bigr).
\end{split}
\end{equation}
Relative to $g_{k_0}^4$, this is of order $g_{k_0}^2$ and
$\kappa_{W,a}\kappa^{\mathfrak{B}}_W$ times the prefactors, plus the exponential terms.

In both bounds $c_{\mathrm{loc}}$ and $c_{\mathrm{w}}$ are the constants of
the corresponding display of Lemma~\ref{8.2:lem-extraction-one-point}. The prefactors $\Pi_2,\Pi_2',\Pi_2''$ are polynomials in
$D$, $r$, $\log x_{k_0}$ and $C_W$, and none of them contains a factor $L^{4r}$.
\item[(c)] (Two-point tower.) Let $G_a=O_{s,p_a}$. Let $\mathcal{A}$ be the free lattice field and
$\mathcal{G}_i:=\sigma(Q^{(i)}\mathcal{A})$ its nested conditionings. Extend
$\mathfrak{g}^{\mathrm{step}}_{12,i}$ to all $i\ge0$ as the per-colour covariance of the fluxes
through $p_1$ and through $p_2$ of the $i$-th fluctuation of $\mathcal{A}$,
\begin{equation*}
\begin{split}
\mathfrak{g}^{\mathrm{step}}_{12,i}&:=n^{-1}E\bigl[\langle\Delta_i\ell_{p_1},
\Delta_i\ell_{p_2}\rangle\bigr],\\
\Delta_i\ell_p&:=E\bigl[\ell_p(\mathcal{A})\,\big|\,\mathcal{G}_i\bigr]
-E\bigl[\ell_p(\mathcal{A})\,\big|\,\mathcal{G}_{i+1}\bigr].
\end{split}
\end{equation*}
By the flat-background hypothesis, for $i<r$ this agrees with the definition of
$\mathfrak{g}^{\mathrm{step}}_{12,i}$ in~(b). Then
\begin{equation}\label{8.2:eq-gaussian-extraction-5}
\begin{split}
&\sum_{i\ge0}\mathfrak{g}^{\mathrm{step}}_{12,i}=\mathfrak{g}_{12}
:=n^{-1}E\bigl[\langle\ell_{p_1}(\mathcal{A}),\ell_{p_2}(\mathcal{A})\rangle\bigr],\\
&\Bigl|\sum_{i\ge r}\mathfrak{g}^{\mathrm{step}}_{12,i}\Bigr|\le C_\eta L^{-4r}.
\end{split}
\end{equation}
Here $\mathfrak{g}_{12}$ is the per-colour covariance of the fluxes through $p_1$ and $p_2$ of the
free lattice field $\mathcal{A}$. The identity is the law of total covariance for $\mathcal{A}$ along
its nested conditionings $\sigma(Q^{(i)}\mathcal{A})$.

Now take the smeared densities $X^{(k_0)}(f_a)$ of Definition~\ref{8.2:def-two-densities} in place
of $G_a$, with $\mathcal{K}_s$ and $G^{(s)}_K$ as in that definition. The same identity, applied term
by term, gives the structure $b_0\mathcal{K}_s(f_1,f_2)$ of clause~(iii)(c) of
Theorem~\ref{3.1:thm-main}, with $G^{(s)}_K$ replaced by its truncation to the $r$ levels
$k_0,\dots,k_0+r-1$. The levels $k_0+i$ with $i\ge r$, coarser than those retained, contribute for
each pair of plaquettes the tail $\sum_{i\ge r}\mathfrak{g}^{\mathrm{step}}_{12,i}$ bounded in
\eqref{8.2:eq-gaussian-extraction-5}.
\end{enumerate}
\end{proposition}

\begin{proof}
\emph{Part~(a).} This is Lemma~\ref{8.2:lem-extraction-one-point} for the
observable $G_1G_2$ of the class displayed in~(a). The product $G_1G_2$ is gauge invariant and
holomorphic on the tube of radius $\varrho$, as a product of two such functions, and its supremum
there is at most $M^2$. Its support is the union of the supports of $G_1$ and $G_2$, which lie in
the one box of Definition~\ref{8.2:def-two-densities} at mutual distance at most $6C_W$; the class
displayed in~(a) records this support by the size $\ell+6C_W$, and the condition on $D$ in that
definition is the one under which Lemma~\ref{8.2:lem-extraction-one-point} applies to an
observable with this support. By the definition of $\mathcal{T}^{(r)}$,
\begin{equation*}
(G_1G_2)^{\langle h\rangle}(W)=\mathcal{T}^{(r)}(W)
+G_1^{\langle h\rangle}(W)\,G_2^{\langle h\rangle}(W).
\end{equation*}
So the main term of $(G_1G_2)^{\langle h\rangle}(W)$ is the sum of two terms. The first is the
main term of \eqref{8.2:eq-gaussian-extraction-a}. The second is the product of the main terms of
the corresponding display for $G_1$ and for $G_2$. The first is computed in the proof of
part~(b), which follows.

\emph{Part~(b): the law of total covariance.} For $0\le i\le r$ let $G_a^{(i)}$ and
$(G_1G_2)^{(i)}$ be the carried functions after $i$ of the $r$ steps on the history $h$. Then
\begin{equation*}
G_a^{(0)}=G_a,\quad (G_1G_2)^{(0)}=G_1G_2,\quad
G_a^{(r)}=G_a^{\langle h\rangle}(W),\quad
(G_1G_2)^{(r)}=(G_1G_2)^{\langle h\rangle}(W).
\end{equation*}
On the joint clean event each step acts by its conditional expectation; this is exact, by the
algebra of conditional expectations and the segregation hypothesis, one of the hypotheses of
Lemma~\ref{8.2:lem-extraction-one-point}:
\begin{equation*}
G_a^{(i+1)}=E_{k_0+i}\bigl[G_a^{(i)}\,\big|\,\cdot\,\bigr],\qquad
(G_1G_2)^{(i+1)}=E_{k_0+i}\bigl[(G_1G_2)^{(i)}\,\big|\,\cdot\,\bigr],
\end{equation*}
where the conditioning is on $V_{k_0+i+1}$. Put
$\mathcal{T}^{(i)}:=(G_1G_2)^{(i)}-G_1^{(i)}G_2^{(i)}$, so that $\mathcal{T}^{(0)}=0$. Insert
$(G_1G_2)^{(i)}=\mathcal{T}^{(i)}+G_1^{(i)}G_2^{(i)}$ and use the definition of the conditional
covariance:
\begin{equation*}
\begin{split}
(G_1G_2)^{(i+1)}
&=E_{k_0+i}\bigl[\mathcal{T}^{(i)}\,\big|\,\cdot\,\bigr]
+E_{k_0+i}\bigl[G_1^{(i)}G_2^{(i)}\,\big|\,\cdot\,\bigr]\\
&=E_{k_0+i}\bigl[\mathcal{T}^{(i)}\,\big|\,\cdot\,\bigr]
+\operatorname{Cov}_{k_0+i}\bigl(G_1^{(i)},G_2^{(i)}\,\big|\,\cdot\,\bigr)
+G_1^{(i+1)}G_2^{(i+1)}.
\end{split}
\end{equation*}
This gives, exactly,
\begin{equation*}
\mathcal{T}^{(i+1)}=E_{k_0+i}\bigl[\mathcal{T}^{(i)}\,\big|\,\cdot\,\bigr]
+\operatorname{Cov}_{k_0+i}\bigl(G_1^{(i)},G_2^{(i)}\,\big|\,\cdot\,\bigr).
\end{equation*}
For $0\le i<r$ write
$\mathbf{E}_i:=E_{k_0+r-1}\circ\dots\circ E_{k_0+i+1}$ for the composed conditional expectation
given $W=V_{k_0+r}$; for $i=r-1$ this is the identity. Starting from $\mathcal{T}^{(0)}=0$, induction
on $j$ and the linearity of conditional expectations give
\begin{equation*}
\mathcal{T}^{(j)}=\sum_{i<j}E_{k_0+j-1}\circ\dots\circ E_{k_0+i+1}
\bigl[\operatorname{Cov}_{k_0+i}\bigl(G_1^{(i)},G_2^{(i)}\,\big|\,\cdot\,\bigr)\bigr].
\end{equation*}
At $j=r$ this reads
\begin{equation}\label{8.2:eq-gaussian-extraction-6}
\mathcal{T}^{(r)}=\sum_{i<r}\mathbf{E}_i\Bigl[\operatorname{Cov}_{k_0+i}
\bigl(G_1^{(i)},G_2^{(i)}\,\big|\,\cdot\,\bigr)\Bigr].
\end{equation}

\emph{Part~(b): splitting off the remainders.} Here and below $\Pi$, $\Pi'$, $\Pi''$ denote
polynomial prefactors in $D$, $r$, $\log x_{k_0}$ and $C_W$ of the type of $\Pi_{r,D}$,
$\Pi'_{r,D}$ in the corresponding display; products of such prefactors, and their
products with powers of $p_0(g_{k_0})=A_0(\log x_{k_0})^{p_0}$, are counted among prefactors of the
same type. The letter $C$ denotes a positive constant not depending on $W$, $h$, $k_0$, $r$ or $D$,
whose value may change from line to line. For each $i$, split
$G_a^{(i)}=\mathfrak{M}_{a,i}+\mathfrak{J}_{a,i}$ into the carried main part $\mathfrak{M}_{a,i}$
and the remainder $\mathfrak{J}_{a,i}$, as in the steps ``one step applied to the carried main
part'' (Step~\ref{8.2:prf-carried-step}) and ``recursion for the extraction remainder''
(Step~\ref{8.2:prf-remainder-recursion}) of the proof of
Lemma~\ref{8.2:lem-extraction-one-point}. Let $J_{a,i}$ bound
$|\mathfrak{J}_{a,i}|$, and put
$\hat{\mathfrak{M}}_{b,i}:=\mathfrak{M}_{b,i}-E_{k_0+i}[\mathfrak{M}_{b,i}|\cdot]$. By bilinearity,
\begin{equation*}
\operatorname{Cov}_{k_0+i}\bigl(G_1^{(i)},G_2^{(i)}\,\big|\,\cdot\,\bigr)
=\operatorname{Cov}_{k_0+i}\bigl(\mathfrak{M}_{1,i},\mathfrak{M}_{2,i}\,\big|\,\cdot\,\bigr)
+R_i,
\end{equation*}
where $R_i$ collects the covariances containing at least one $\mathfrak{J}$. We have
\begin{equation*}
\operatorname{Cov}_{k_0+i}\bigl(\mathfrak{M}_{b,i},\mathfrak{J}_{a,i}\,\big|\,\cdot\,\bigr)
=E_{k_0+i}\bigl[\hat{\mathfrak{M}}_{b,i}\mathfrak{J}_{a,i}\,\big|\,\cdot\,\bigr],
\end{equation*}
and this is bounded by $J_{a,i}E_{k_0+i}\bigl[|\hat{\mathfrak{M}}_{b,i}|\,\big|\,\cdot\,\bigr]$,
which we abbreviate to $J_{a,i}E|\hat{\mathfrak{M}}_{b,i}|$. The covariance of the two remainders is
bounded by a numerical multiple of $J_{1,i}J_{2,i}$. Hence $|R_i|$ is bounded, up to a numerical
factor, by
\begin{equation*}
\sum_{a\ne b}J_{a,i}\,E|\hat{\mathfrak{M}}_{b,i}|+J_{1,i}J_{2,i}.
\end{equation*}
By the step ``recursion for the extraction remainder'' of the proof of
Lemma~\ref{8.2:lem-extraction-one-point}, $J_{a,i}\le Cg^3\Pi M$, and by the step ``one step
applied to the carried main part'' of that proof, $E|\hat{\mathfrak{M}}_{b,i}|\le
C(g\kappa+g^2)M$. Therefore this is at most
\begin{equation*}
C\,(g^3\Pi)(g\kappa+g^2)M^2+C(g^3\Pi)^2M^2,
\end{equation*}
with $g=g_{k_0}$ and $\kappa$ the curvature size of the background at that step. These
contributions are absorbed in \eqref{8.2:eq-gaussian-extraction-3}: the product
$(g^3\Pi)(g\kappa)M^2$ in the members with $\Pi_2'$, using the bounds on $\kappa$ recorded below for
the one-step errors ($\kappa\le4g_{k_0}p_0(g_{k_0})$ for $i<r-1$, and $\kappa$ of the type
$\kappa_{W,a}$ at $i=r-1$); the product $(g^3\Pi)g^2M^2$ in the member
$g_{k_0}^4\Pi_2'\,g_{k_0}p_0(g_{k_0})$; and $(g^3\Pi)^2$ in the member $g_{k_0}^6\Pi_2$. Applying
$\mathbf{E}_i$ does not increase their supremum.

For the bound \eqref{8.2:eq-gaussian-extraction-4} the remainder recursion is used instead in its
form under the equivariance hypothesis, in which the remainder bounds are taken in their
equivariant form; it gives
\begin{equation*}
J_{a,i}\le C\bigl(g^4\Pi+g^2\bar{\kappa}\bar{\kappa}^{\mathfrak{B}}\Pi'\bigr)M,
\end{equation*}
where $\bar{\kappa}$ and $\bar{\kappa}^{\mathfrak{B}}$ are the largest curvature sizes of the
background and of the box at the steps $k_0,\dots,k_0+i-1$ at which the remainder is produced.
These steps use the backgrounds $V_{k_0+j+1}$ with $j\le i-1<r-1$, so
$\bar{\kappa}\le4g_{k_0}p_0(g_{k_0})$ by the bound recorded below for the one-step errors. With
$E|\hat{\mathfrak{M}}_{b,i}|\le C(g\kappa+g^2)M$, where $\kappa$ is the curvature size at the step
$k_0+i$, the products $J_{a,i}E|\hat{\mathfrak{M}}_{b,i}|$ consist of the four members
\begin{equation*}
g^6\Pi M^2,\qquad g^5\kappa\Pi M^2,\qquad
g^4\bar{\kappa}\bar{\kappa}^{\mathfrak{B}}\Pi'M^2,\qquad
g^3\kappa\bar{\kappa}\bar{\kappa}^{\mathfrak{B}}\Pi'M^2.
\end{equation*}
Since $\bar{\kappa}\le4g_{k_0}p_0(g_{k_0})$, the fourth is at most
$4g_{k_0}^4p_0(g_{k_0})\,\kappa\bar{\kappa}^{\mathfrak{B}}\Pi'M^2$; with the factor
$p_0(g_{k_0})$ taken into the prefactor, it has the form of the third. The product
$J_{1,i}J_{2,i}$ consists of the members
\begin{equation*}
g^8\Pi^2M^2,\qquad g^6\bar{\kappa}\bar{\kappa}^{\mathfrak{B}}\Pi\Pi'M^2,\qquad
g^4\bigl(\bar{\kappa}\bar{\kappa}^{\mathfrak{B}}\Pi'\bigr)^2M^2;
\end{equation*}
since $g_{k_0}\le1$ and the curvature sizes are bounded by a constant, each of them is at most a
constant times a member of the form $g^6\Pi M^2$ or of the form
$g^4\bar{\kappa}\bar{\kappa}^{\mathfrak{B}}\Pi'M^2$, the prefactors being multiplied. The members
$g^6\Pi M^2$ and those of the form $g^4\kappa\kappa^{\mathfrak{B}}\Pi'M^2$ are placed in
\eqref{8.2:eq-gaussian-extraction-4} as the corresponding members of the one-step errors below.
For the member $g^5\kappa\Pi M^2$, if $i<r-1$ then $\kappa\le4g_{k_0}p_0(g_{k_0})$, and since
$p_0(g_{k_0})\le\tfrac12(1+p_0(g_{k_0})^2)$,
\begin{equation*}
g_{k_0}^5\kappa\Pi\le4g_{k_0}^6p_0(g_{k_0})\Pi
\le2g_{k_0}^6\Pi+2g_{k_0}^4\,g_{k_0}^2p_0(g_{k_0})^2\,\Pi,
\end{equation*}
which lies in the members $g_{k_0}^6\Pi_2$ and $g_{k_0}^4\Pi_2'g_{k_0}^2p_0(g_{k_0})^2$ of
\eqref{8.2:eq-gaussian-extraction-4}. If $i=r-1$, then $\kappa$ is of the type $\kappa_{W,a}$. By
Definition~\ref{8.2:def-transported-laws}, $\kappa_{W,a}$ is at most a constant times
$\kappa^{\mathfrak{B}}_W$, the neighbourhood of the image plaquette over which $\kappa_{W,a}$ is
taken lying in the level-$(k_0+r)$ image of the box. With
$g_{k_0}\kappa_{W,a}\le\tfrac12(g_{k_0}^2+\kappa_{W,a}^2)$ this gives
\begin{equation*}
g_{k_0}^5\kappa_{W,a}\Pi\le\tfrac12g_{k_0}^6\Pi+\tfrac12g_{k_0}^4\kappa_{W,a}^2\Pi
\le\tfrac12g_{k_0}^6\Pi+Cg_{k_0}^4\kappa_{W,a}\kappa^{\mathfrak{B}}_W\Pi,
\end{equation*}
which lies in the members $g_{k_0}^6\Pi_2$ and
$g_{k_0}^4\Pi_2'(\kappa_{W,1}+\kappa_{W,2})\kappa^{\mathfrak{B}}_W$ of
\eqref{8.2:eq-gaussian-extraction-4}.

\emph{Part~(b): one step.} The main parts $\mathfrak{M}_{1,i},\mathfrak{M}_{2,i}$ are local
functions of the level-$(k_0+i)$ field with both supports in the one box of the joint clean event.
Apply part~(v) of Lemma~\ref{8.2:lem-one-step-extraction}, the corresponding display, to this pair at the step $k_0+i$. It gives
\begin{equation*}
\operatorname{Cov}_{k_0+i}\bigl(\mathfrak{M}_{1,i},\mathfrak{M}_{2,i}\,\big|\,V'\bigr)
=\Phi_i(V')+\Psi_i(V')+\mathcal{E}^{(2)}_{k_0+i},\qquad V':=V_{k_0+i+1}.
\end{equation*}
The one-line term is
\begin{equation*}
\Phi_i(V'):=g_{k_0+i}^2\bigl\langle D\mathfrak{M}_{1,i}(\mathrm{bg}_1),C^{(k_0+i)}(V')\,
D\mathfrak{M}_{2,i}(\mathrm{bg}_2)\bigr\rangle,
\end{equation*}
where $\mathrm{bg}_a$ is the background of the step near $p_a$, as in part~(v) of
Lemma~\ref{8.2:lem-one-step-extraction}.
In terms of $G_1$ and $G_2$ it reads
\begin{equation*}
\Phi_i(V')=g_{k_0+i}^2\bigl\langle DG_1(\mathfrak{B}^{(i+1)}(V'))[\mathfrak{L}^{(i)}\,\cdot\,],
\,C^{(k_0+i)}\,DG_2(\mathfrak{B}^{(i+1)}(V'))[\mathfrak{L}^{(i)}\,\cdot\,]\bigr\rangle
+\mathcal{E}^{\Phi}_i(V').
\end{equation*}
Here $\mathcal{E}^{\Phi}_i(V')$ is the correction made by replacing the differentials of the main
parts $\mathfrak{M}_{a,i}$ by those of $G_a$ at the transported background; it is of order
$g_{k_0+i}^4$ multiplied by further factors, and it is collected with the errors below.
The two-line term is
\begin{equation*}
\Psi_i(V'):=\frac14g_{k_0+i}^4\operatorname{Cov}_\gamma\bigl(
\operatorname{Hess}\mathfrak{M}_{1,i}[iCB,iCB],
\operatorname{Hess}\mathfrak{M}_{2,i}[iCB,iCB]\bigr),
\end{equation*}
where $\gamma=N(0,C^{(k_0+i)}(V'))$. The error $\mathcal{E}^{(2)}_{k_0+i}$ obeys the bound of
the corresponding display of Lemma~\ref{8.2:lem-one-step-extraction}. Both $\Phi_i$ and $\Psi_i$ are gauge-invariant holomorphic
local functions of $V'$.

\emph{Part~(b): carrying $\Psi_i$.} We carry $\Phi_i$ and $\Psi_i$ through the remaining $r-i-1$
steps. We use Lemma~\ref{8.2:lem-extraction-one-point} with one insertion, initial level
$k_0+i+1$ and $r-i-1$ steps. For the bound \eqref{8.2:eq-gaussian-extraction-3} the form of
Lemma~\ref{8.2:lem-extraction-one-point} without the equivariance hypothesis suffices, because
$\Phi_i$ and $\Psi_i$ already carry $g^2$ and $g^4$. For the bound
\eqref{8.2:eq-gaussian-extraction-4} we use instead the form of
Lemma~\ref{8.2:lem-extraction-one-point} under the equivariance hypothesis, in which the remainder
bounds of
Lemma~\ref{8.2:lem-one-step-extraction} are taken in their equivariant form: the exponent
$\sigma_1$ of the corresponding display equals~$2$, and the member linear in the
curvature size carries the further factor $\kappa^{\mathfrak{B}}_W$. For $\Psi_i$ it gives
\begin{equation*}
\mathbf{E}_i[\Psi_i]=\Psi_i\bigl(\mathfrak{B}_{i+1,r}(W)\bigr)+O(g^2\cdot g^4).
\end{equation*}
The value on the right is the diagonal two-line term
\begin{equation*}
\frac14g_{k_0+i}^4\operatorname{Cov}_{\xi\sim\Gamma^{12}_i(W)}\bigl(
\operatorname{Hess}_{\mathrm{bg}_r(W)}G_1[\xi,\xi],\operatorname{Hess}_{\mathrm{bg}_r(W)}G_2[\xi,\xi]
\bigr).
\end{equation*}
By the Isserlis formula, as in the corresponding display of Lemma~\ref{8.2:lem-one-step-extraction}, with
$H_a=\operatorname{Hess}_{\mathrm{bg}_r(W)}G_a$ and $\xi,\xi'$ independent with law
$\Gamma^{12}_i(W)$,
\begin{equation*}
\operatorname{Cov}\bigl(H_1[\xi,\xi],H_2[\xi,\xi]\bigr)=2E\bigl[H_1[\xi,\xi']\,H_2[\xi,\xi']\bigr]
=2(H_1\otimes H_2)[\Gamma^{12}_i,\Gamma^{12}_i].
\end{equation*}
Hence this value equals $\tfrac12g_{k_0+i}^4(H_1\otimes H_2)[\Gamma^{12}_i,\Gamma^{12}_i]$, which is
the term $i'=i$ of the two-line sum in \eqref{8.2:eq-gaussian-extraction-a}.

\emph{Part~(b): carrying $\Phi_i$.} For $\Phi_i$, Lemma~\ref{8.2:lem-extraction-one-point} without
the equivariance hypothesis gives
\begin{equation*}
\mathbf{E}_i[\Phi_i]=\Phi_i\bigl(\mathfrak{B}_{i+1,r}(W)\bigr)
+\frac12\sum_{i<i'<r}g_{k_0+i'}^2\operatorname{Tr}_{\Gamma^{12}_{i'}}[\operatorname{Hess}\Phi_i]
+O(g^2\cdot g^3\Pi).
\end{equation*}
The value $\Phi_i(\mathfrak{B}_{i+1,r}(W))$ has its gradients at $\mathrm{bg}_r(W)$ and its
transported fluctuation distributed according to $\Gamma^{12}_i(W)$. It is therefore the $i$-th
term of the one-line sum in \eqref{8.2:eq-gaussian-extraction-a} plus
$\mathcal{E}^{\Phi}_i(\mathfrak{B}_{i+1,r}(W))$, which is collected with the errors below. That
term is a product of two background
gradients. It vanishes at flat $W$ by part~(a) of the algebraic lemma on gauge-invariant local
functions, and it is always at most a constant times $\kappa_{W,1}\kappa_{W,2}$; this is the factor
through which \eqref{8.2:eq-gaussian-extraction-1} depends on $W$. The factor
$C_{\mathrm{ax}}^2\mu_2^2$ multiplying the product of the curvature sizes of the two backgrounds is
that of the one-line bound in the corresponding display of Lemma~\ref{8.2:lem-one-step-extraction}; in
\eqref{8.2:eq-gaussian-extraction-1} the $i$-th term carries in addition the weight
$g_{k_0+i}^2L^{-4i}$.

For the Hessian of $\Phi_i$ at the flat configuration we have
\begin{equation*}
\begin{split}
\Phi_i(e^{Z}\cdot\mathrm{flat})={}&g_{k_0+i}^2\bigl\langle
D^2G_1(\mathrm{flat})[\mathfrak{L}^{(i+1)}Z,\mathfrak{L}^{(i)}\,\cdot\,],\\
&\qquad C^{(k_0+i)}\,D^2G_2(\mathrm{flat})[\mathfrak{L}^{(i+1)}Z,\mathfrak{L}^{(i)}\,\cdot\,]
\bigr\rangle+O(Z^3).
\end{split}
\end{equation*}
Hence $\tfrac12g_{k_0+i'}^2\operatorname{Tr}_{\Gamma^{12}_{i'}}[\operatorname{Hess}\Phi_i]$ equals
\begin{equation*}
g_{k_0+i}^2g_{k_0+i'}^2\,
\bigl(\operatorname{Hess}G_1\otimes\operatorname{Hess}G_2\bigr)\bigl[\Gamma^{12}_i,\Gamma^{12}_{i'}\bigr],
\end{equation*}
the contraction in which one pair is contracted through $\Gamma^{12}_i$ and the other through
$\Gamma^{12}_{i'}$. Since this contraction is symmetric in its two laws, the two ordered pairs
$(i,i')$ and $(i',i)$ of the two-line sum in \eqref{8.2:eq-gaussian-extraction-a}, each with the
factor $\tfrac12$, together give it exactly once. These are the cross-level terms
$i\ne i'$ of the two-line sum in \eqref{8.2:eq-gaussian-extraction-a}.

\emph{Part~(b): the combinatorial factor for $G_a=O_{s,p_a}$.} Let $G_a=O_{s,p_a}$. Then
\begin{equation*}
D^2O_{s,p_a}(\mathbf{1})[Z,Z']=\langle\ell_{p_a}(Z),\ell_{p_a}(Z')\rangle,
\end{equation*}
and the colour structure is diagonal. Consider each unordered pair $\{i,i'\}$ with $i\ne i'$. It
arises exactly once, namely from carrying $\Phi_{\min(i,i')}$ through the step $\max(i,i')$, and it
contributes
\begin{equation*}
n\,g_{k_0+i}^2g_{k_0+i'}^2\,
\mathfrak{g}^{\mathrm{step}}_{12,i}\mathfrak{g}^{\mathrm{step}}_{12,i'}.
\end{equation*}
The term $\Psi_i$ contributes $\tfrac n2g_{k_0+i}^4(\mathfrak{g}^{\mathrm{step}}_{12,i})^2$. This
follows from the Isserlis formula: let $\xi,\eta$ be jointly centred Gaussian vectors in
$\mathfrak{g}$ with $\operatorname{Cov}(\xi_a,\eta_b)=\delta_{ab}c$. Then
$\operatorname{Cov}(\xi_a^2,\eta_b^2)=2\operatorname{Cov}(\xi_a,\eta_b)^2$, and therefore
\begin{equation*}
\operatorname{Cov}\bigl(\tfrac12|\xi|^2,\tfrac12|\eta|^2\bigr)
=\frac12\sum_{a,b}\operatorname{Cov}(\xi_a,\eta_b)^2=\frac n2c^2.
\end{equation*}
We apply this with $c=\mathfrak{g}^{\mathrm{step}}_{12,i}$. Adding the diagonal and the cross-level
contributions gives
\begin{equation*}
\sum_{i<r}\frac n2g_{k_0+i}^4(\mathfrak{g}^{\mathrm{step}}_{12,i})^2
+\sum_{i<i'<r}n\,g_{k_0+i}^2g_{k_0+i'}^2
\mathfrak{g}^{\mathrm{step}}_{12,i}\mathfrak{g}^{\mathrm{step}}_{12,i'}
=\frac n2\Bigl(\sum_{i<r}g_{k_0+i}^2\mathfrak{g}^{\mathrm{step}}_{12,i}\Bigr)^2,
\end{equation*}
which is \eqref{8.2:eq-gaussian-extraction-2}.

\emph{Part~(b): the errors.} There are four kinds.
\begin{itemize}
\item The one-step errors $\mathcal{E}^{(2)}_{k_0+i}$ of the corresponding display of Lemma~\ref{8.2:lem-one-step-extraction}.
Each has a member of the form $g^6\Pi$, a member of the form $g^4\kappa\Pi'$ and a member of the
form $g^4\kappa_{\mathrm{bg}_1}\kappa_{\mathrm{bg}_2}\Pi''$; under the equivariance hypothesis the
second becomes $g^4\kappa\kappa^{\mathfrak{B}}\Pi'$. Here $\kappa=\kappa(V_{k_0+i+1})$, and
$\kappa_{\mathrm{bg}_a}$ are the curvature sizes of the step's backgrounds near $p_1$ and $p_2$.
\begin{itemize}
\item For $i<r-1$ this satisfies, in the supremum, $\kappa\le4g_{k_0}p_0(g_{k_0})$. Without the
equivariance hypothesis this produces the member $g_{k_0}^4\cdot g_{k_0}p_0(g_{k_0})\Pi_2'$ of
\eqref{8.2:eq-gaussian-extraction-3}. It is of relative order $g_{k_0}p_0(g_{k_0})$, the two-point
analogue of the exponent $\sigma_1=1$ of the corresponding display. Under the
equivariance hypothesis it produces $g_{k_0}^4(g_{k_0}p_0(g_{k_0}))^2$, which is the member
$g_{k_0}^4\Pi_2'g_{k_0}^2p_0(g_{k_0})^2$ of \eqref{8.2:eq-gaussian-extraction-4}. With the same
bound for the curvature sizes $\kappa_{\mathrm{bg}_a}$ of the step, the member
$g^4\kappa_{\mathrm{bg}_1}\kappa_{\mathrm{bg}_2}\Pi''$ is at most
$16g_{k_0}^6p_0(g_{k_0})^2\Pi''$; with $p_0(g_{k_0})^2$ taken into the prefactor, it lies in the
member $g_{k_0}^6\Pi_2$ of both bounds.
\item At $i=r-1$ the curvature size is of the type $\kappa_{W,a}$. This gives the members in
$\kappa_{W,1}+\kappa_{W,2}$, respectively $(\kappa_{W,1}+\kappa_{W,2})\kappa^{\mathfrak{B}}_W$,
together with the member in $\kappa_{W,1}\kappa_{W,2}$.
\end{itemize}
\item The corrections $\mathcal{E}^{\Phi}_i$ from rewriting the one-line term in terms of $G_1$ and
$G_2$, each of order $g_{k_0+i}^4$ multiplied by further factors; they are collected into
$\mathcal{E}_{2,r}(W)$.
\item The errors from carrying $\Psi_i$ and $\Phi_i$. Without the equivariance hypothesis these are
of relative order $g^2$ and $g^3\Pi$ against the factors $g^4$ and $g^2$ that $\Psi_i$ and $\Phi_i$
carry. Under the equivariance hypothesis, the form of Lemma~\ref{8.2:lem-extraction-one-point}
used above bounds the error of carrying $\Phi_i$ by a constant times
\begin{equation*}
g_{k_0}^2M^2\bigl(g_{k_0}^4\Pi+g_{k_0}^2\Pi'\,(\kappa_{W,1}+\kappa_{W,2})\,
\kappa^{\mathfrak{B}}_W\bigr),
\end{equation*}
and that of carrying $\Psi_i$ by $g_{k_0}^2$ times this; they lie in the members
$g_{k_0}^6\Pi_2$ and $g_{k_0}^4\Pi_2'(\kappa_{W,1}+\kappa_{W,2})\kappa^{\mathfrak{B}}_W$ of
\eqref{8.2:eq-gaussian-extraction-4}.
\item The cross terms $R_i$ containing a remainder $\mathfrak{J}$, estimated above.
\end{itemize}
Collecting these errors, together with the exponential terms of the applied estimates, gives
$\mathcal{E}_{2,r}(W)$. It satisfies \eqref{8.2:eq-gaussian-extraction-3}, and, under the
equivariance hypothesis, \eqref{8.2:eq-gaussian-extraction-4}. Substituting the contributions
identified above into \eqref{8.2:eq-gaussian-extraction-6} gives
\eqref{8.2:eq-gaussian-extraction-a}.

Part~(c) is proved below, from the flat-background hypothesis and the law of total covariance for
$\mathcal{A}$ along its nested conditionings $\sigma(Q^{(i)}\mathcal{A})$.
\end{proof}

\begin{proof}[Proof of clause (c) of Proposition~\ref{8.2:prop-gaussian-extraction}: the two-point tower identity]
\label{8.2:prf-two-point-tower}
Clause (c) concerns $G_a=O_{s,p_a}$, $a=1,2$. It asserts that
\begin{equation*}
\sum_{i\ge 0}\mathfrak{g}^{\mathrm{step}}_{12,i}(\mathrm{flat})=\mathfrak{g}_{12}
:=\operatorname{Cov}(\ell_{p_1},\ell_{p_2}),
\end{equation*}
the covariance of the plaquette fluxes at $p_1$ and $p_2$ of the free lattice field $\mathcal{A}$, and
that $\bigl|\sum_{i\ge r}\mathfrak{g}^{\mathrm{step}}_{12,i}(\mathrm{flat})\bigr|\le C_\eta L^{-4r}$.
For $i\ge0$ let $\mathcal{G}_i$ be the $\sigma$-algebra generated by the field $Q^{(i)}\mathcal{A}$,
with $Q^{(0)}$ the identity, so that $\mathcal{G}_0$ is the $\sigma$-algebra generated by $\mathcal{A}$.
These conditionings are nested, $\mathcal{G}_{i+1}\subset\mathcal{G}_i$.
The clause is derived from three inputs, taken as hypotheses:
\begin{itemize}
\item[(1)] the lemma on the free lattice field $\mathcal{A}$ and its plaquette fluxes $\ell_p$ under
the linearised block-averaging maps $Q^{(i)}$;
\item[(2)] the flat-background hypothesis of this chapter, which identifies the step covariances: at
flat background, the covariance of the $i$-th conditional increment of $\mathcal{A}$ defined below is
Ba{\l}aban's $i$-th step covariance transported by the linearised minimiser;
\item[(3)] the one-point tower identity for a single plaquette flux, through its variance tail: for
every $I\ge0$ and $p=p_1,p_2$ the variance of $E[\ell_p\mid\mathcal{G}_I]$ is at most
$C_\eta L^{-4I}$.
\end{itemize}

For $i\ge0$ put $X_i:=E[\ell_{p_1}\mid\mathcal{G}_i]$ and $Y_i:=E[\ell_{p_2}\mid\mathcal{G}_i]$. Since
$\ell_{p_a}$ is $\mathcal{G}_0$-measurable, $X_0=\ell_{p_1}$ and $Y_0=\ell_{p_2}$. Since
$\mathcal{G}_{i+1}\subset\mathcal{G}_i$, the tower property gives $E[X_i\mid\mathcal{G}_{i+1}]=X_{i+1}$
and $E[Y_i\mid\mathcal{G}_{i+1}]=Y_{i+1}$.

Now apply the law of total covariance to the pair $(X_i,Y_i)$, conditioning on $\mathcal{G}_{i+1}$:
\begin{equation}\label{8.2:eq-two-point-tower-1}
\operatorname{Cov}(X_i,Y_i)
=E\bigl[\operatorname{Cov}(X_i,Y_i\mid\mathcal{G}_{i+1})\bigr]+\operatorname{Cov}(X_{i+1},Y_{i+1}).
\end{equation}
Sum \eqref{8.2:eq-two-point-tower-1} over $0\le i<I$; the terms $\operatorname{Cov}(X_i,Y_i)$ with
$1\le i<I$ cancel, and we obtain
\begin{equation}\label{8.2:eq-two-point-tower-2}
\operatorname{Cov}(\ell_{p_1},\ell_{p_2})
=\sum_{i<I}E\bigl[\operatorname{Cov}(X_i,Y_i\mid\mathcal{G}_{i+1})\bigr]
+\operatorname{Cov}(X_I,Y_I).
\end{equation}

The field $\mathcal{A}$ is Gaussian, so every $X_i$ and $Y_i$ is a linear functional of $\mathcal{A}$,
and the conditional covariances of jointly Gaussian variables are deterministic. Hence
\begin{equation*}
\operatorname{Cov}(X_i,Y_i\mid\mathcal{G}_{i+1})
=E\bigl[(X_i-X_{i+1})(Y_i-Y_{i+1})\mid\mathcal{G}_{i+1}\bigr]
\end{equation*}
is a constant. It therefore equals its own expectation,
\begin{equation*}
E\bigl[(X_i-X_{i+1})(Y_i-Y_{i+1})\bigr]=\operatorname{Cov}(X_i-X_{i+1},Y_i-Y_{i+1}),
\end{equation*}
the increments having mean zero because, by the tower property, $E[X_i]=E[\ell_{p_1}]$ and
$E[Y_i]=E[\ell_{p_2}]$ for every $i$. So the $i$-th summand in
\eqref{8.2:eq-two-point-tower-2} is the $p_1$--$p_2$ covariance of the $i$-th conditional increment
$(X_i-X_{i+1},\,Y_i-Y_{i+1})$.

Input (2) identifies this covariance with Ba{\l}aban's $i$-th step covariance at flat background,
transported by the linearised minimiser. The linearised minimiser is the Gaussian conditional mean,
\cite[(1.23)--(1.24)]{B-CMP95}. By the definition of $\mathfrak{g}^{\mathrm{step}}_{12,i}$, the $i$-th
summand is therefore $\mathfrak{g}^{\mathrm{step}}_{12,i}(\mathrm{flat})$, and
\eqref{8.2:eq-two-point-tower-2} becomes
\begin{equation}\label{8.2:eq-two-point-tower-3}
\mathfrak{g}_{12}
=\sum_{i<I}\mathfrak{g}^{\mathrm{step}}_{12,i}(\mathrm{flat})+\operatorname{Cov}(X_I,Y_I)
\qquad(I\ge0).
\end{equation}

Write $\operatorname{Var}(X):=\operatorname{Cov}(X,X)$ for the variance under the law of $\mathcal{A}$.
For the remainder we use the Cauchy--Schwarz inequality and input (3), applied first at a general
level $I$ and, below, at $I=r$:
\begin{equation}\label{8.2:eq-two-point-tower-4}
\bigl|\operatorname{Cov}(X_I,Y_I)\bigr|
\le\operatorname{Var}(X_I)^{1/2}\operatorname{Var}(Y_I)^{1/2}
\le C_\eta L^{-4I}.
\end{equation}
Since $L>1$, the right side tends to $0$ as $I\to\infty$. By \eqref{8.2:eq-two-point-tower-3}, the
partial sums converge to $\mathfrak{g}_{12}$, which is the identity
$\sum_{i\ge0}\mathfrak{g}^{\mathrm{step}}_{12,i}(\mathrm{flat})=\mathfrak{g}_{12}$.

Subtracting \eqref{8.2:eq-two-point-tower-3} with $I=r$ from this identity gives
$\sum_{i\ge r}\mathfrak{g}^{\mathrm{step}}_{12,i}(\mathrm{flat})=\operatorname{Cov}(X_r,Y_r)$. By
\eqref{8.2:eq-two-point-tower-4} with $I=r$, its modulus is at most $C_\eta L^{-4r}$. This proves
clause (c).
\end{proof}

\begin{corollary}[Conditional second moment of the plaquette observable]\label{8.2:cor-second-moment}
Assume the hypotheses of Lemma~\ref{8.2:lem-extraction-one-point} (the Gaussian extraction for one local
observable), in particular the localisation of the sourced step, with $G=O_{s,p}$, so that $M=M_O$ and
$\ell=1$. Let $r$, $D$, $h\in\mathcal{C}_p(D)$ and $W$ in the level-$(k_0+r)$ window be as there.
Write $n=\dim\mathfrak{g}$,
$\Theta_{\mathrm{bg}}(W)=O_{s,p}(\mathrm{bg}_r(W))$, and let $\mathcal{T}^{(r)}(O,O)(W)$ be the carried
conditional covariance of Proposition~\ref{8.2:prop-gaussian-extraction} (the Gaussian extraction) with
$G_1=G_2=O_{s,p}$. Then the following hold.
\begin{enumerate}
\item[(a)] Exactly,
\begin{equation}\label{8.2:eq-second-moment-1}
(O_{s,p}^2)^{\langle h\rangle}(W)=\bigl(O_{s,p}^{\langle h\rangle}(W)\bigr)^2
+\mathcal{T}^{(r)}(O,O)(W).
\end{equation}
\item[(b)] To fourth order in the couplings,
\begin{equation}\label{8.2:eq-second-moment-2}
\begin{split}
(O_{s,p}^2)^{\langle h\rangle}(W)
&=\Bigl[\Theta_{\mathrm{bg}}(W)+\frac{n}{2}\sum_{i<r}g_{k_0+i}^2\,
\mathfrak{g}^{\mathrm{step}}_{k_0,i}(W)\Bigr]^2\\
&\quad+\sum_{i<r}g_{k_0+i}^2\,\bigl\langle DO(\mathrm{bg}_r),\Gamma_i(W)\,DO(\mathrm{bg}_r)\bigr\rangle\\
&\quad+\frac12\sum_{i,i'<r}g_{k_0+i}^2\,g_{k_0+i'}^2\sum_{\mathrm{contractions}}
\bigl(\operatorname{Hess}_{\mathrm{bg}_r}O_{s,p}\otimes\operatorname{Hess}_{\mathrm{bg}_r}O_{s,p}\bigr)
\bigl[\Gamma_i(W),\Gamma_{i'}(W)\bigr]\\
&\quad+\mathcal{E}'_r(W),
\end{split}
\end{equation}
where the sum over contractions is the one in the two-line term of
Proposition~\ref{8.2:prop-gaussian-extraction}(b).
\item[(c)] At flat $W$ the main term of \eqref{8.2:eq-second-moment-2} equals
\begin{equation}\label{8.2:eq-second-moment-3}
\Bigl(1+\frac{2}{n}\Bigr)\Bigl(\frac{n}{2}\sum_{i<r}g_{k_0+i}^2\,
\mathfrak{g}^{\mathrm{step}}_{k_0,i}\Bigr)^2 .
\end{equation}
\item[(d)] The remainder satisfies
\begin{equation}\label{8.2:eq-second-moment-4}
\begin{split}
|\mathcal{E}'_r(W)|\le C_{\mathrm{ex}}M_O^2\Bigl[\,&g_{k_0}^{4+\sigma_1}\Pi
+g_{k_0}^4\Pi'\bigl(\kappa_W+g_{k_0}\,p_0(g_{k_0})\bigr)
+g_{k_0}^2\,\Theta_{\mathrm{bg}}(W)\,\Pi\\
&+e^{-c'_{\mathrm{loc}}D}+e^{-c'_w p_0(g_{k_0})^2}\Bigr],
\end{split}
\end{equation}
where $C_{\mathrm{ex}}$ is a constant depending only on $(L,\mathfrak{p},N)$ and the step constants, as
the constant of the corresponding display of Lemma~\ref{8.2:lem-extraction-one-point} does, $\sigma_1$ is the
exponent of the corresponding display of Lemma~\ref{8.2:lem-extraction-one-point} ($\sigma_1=1$ in this form), $\Pi$ and $\Pi'$ are
polynomials in $D$, $r$, $\log x_{k_0}$ and $C_W$ containing no power of $L^{4r}$, and
$c'_{\mathrm{loc}}$, $c'_w$ are positive constants depending only on the rates $c_{\mathrm{loc}}$, $c_w$
of the corresponding display of Lemma~\ref{8.2:lem-extraction-one-point} and of the remainder bound in
Proposition~\ref{8.2:prop-gaussian-extraction}(b).
\item[(e)] If in addition the equivariance hypothesis under which the sharpened form of
Lemma~\ref{8.2:lem-extraction-one-point} holds is assumed, then \eqref{8.2:eq-second-moment-4} holds with
$\sigma_1=2$, with $\kappa_W\kappa^{\mathfrak{B}}_W$ in place of $\kappa_W$, and with
$g_{k_0}^2\,p_0(g_{k_0})^2$ in place of $g_{k_0}\,p_0(g_{k_0})$.
\end{enumerate}
\end{corollary}

\begin{proof}
\emph{Applicability of the Gaussian extraction.} Proposition~\ref{8.2:prop-gaussian-extraction} is
applied with $G_1=G_2=O_{s,p}$ and $p_1=p_2=p$. Coincident supports are allowed there: its hypotheses
require both supports to lie in one box, and no separation between them is used. With $p_1=p_2=p$ the two
curvature sizes coincide, $\kappa_{W,1}=\kappa_{W,2}=\kappa_W$, and the joint law $\Gamma^{12}_i(W)$ lives
on $(\mathfrak{g})^{\mathfrak{B}_1(p)}$; its marginal is $\Gamma_i(W)$, so $\Gamma^{12}_i(W)=\Gamma_i(W)$.

\emph{Step 1: the exact identity.} By the definition of the carried conditional covariance in
Proposition~\ref{8.2:prop-gaussian-extraction}(b), with $G_1=G_2=O_{s,p}$,
\begin{equation*}
\mathcal{T}^{(r)}(O,O)(W)=(O_{s,p}^2)^{\langle h\rangle}(W)-O_{s,p}^{\langle h\rangle}(W)\,
O_{s,p}^{\langle h\rangle}(W).
\end{equation*}
Rearranging gives \eqref{8.2:eq-second-moment-1}.

\emph{Step 2: the carried mean.} By the corresponding display of Lemma~\ref{8.2:lem-extraction-one-point} for $G=O_{s,p}$, whose main
term Lemma~\ref{8.2:lem-extraction-one-point} identifies,
\begin{equation*}
O_{s,p}^{\langle h\rangle}(W)=\Theta_{\mathrm{bg}}(W)+\frac{n}{2}\sum_{i<r}g_{k_0+i}^2\,
\mathfrak{g}^{\mathrm{step}}_{k_0,i}(W)+\mathcal{E}_r(O_{s,p};W),
\end{equation*}
with $\mathcal{E}_r(O_{s,p};W)$ bounded by the corresponding display of Lemma~\ref{8.2:lem-extraction-one-point}. Squaring gives the
first line of \eqref{8.2:eq-second-moment-2} plus the cross term and the square
\begin{equation*}
2\Bigl[\Theta_{\mathrm{bg}}(W)+\frac n2\sum_{i<r}g_{k_0+i}^2\,
\mathfrak{g}^{\mathrm{step}}_{k_0,i}(W)\Bigr]\mathcal{E}_r(O_{s,p};W)
+\mathcal{E}_r(O_{s,p};W)^2 .
\end{equation*}

\emph{Step 3: the carried variance.} By Proposition~\ref{8.2:prop-gaussian-extraction}(b) with
$G_1=G_2=O_{s,p}$, $\mathcal{T}^{(r)}(O,O)(W)$ is the sum of a one-line term, a two-line term and a
remainder $\mathcal{E}_{2,r}(W)$. Since $\Gamma^{12}_i(W)=\Gamma_i(W)$ and $G_1=G_2=O_{s,p}$, the
one-line term is the second line of \eqref{8.2:eq-second-moment-2} and the two-line term is the third
line of \eqref{8.2:eq-second-moment-2}. At flat $W$ that proposition evaluates the two-line term as
$\tfrac n2\bigl(\sum_{i<r}g_{k_0+i}^2\mathfrak{g}^{\mathrm{step}}_{12,i}\bigr)^2$, and for $p_1=p_2=p$ the
$p_1$--$p_2$ flux covariance per colour $\mathfrak{g}^{\mathrm{step}}_{12,i}$ is the one-point
coefficient $\mathfrak{g}^{\mathrm{step}}_{k_0,i}$; hence at flat $W$ the third line of
\eqref{8.2:eq-second-moment-2} equals
\begin{equation*}
\frac n2\Bigl(\sum_{i<r}g_{k_0+i}^2\,\mathfrak{g}^{\mathrm{step}}_{k_0,i}\Bigr)^2 .
\end{equation*}

\emph{Step 4: the remainder.} Define $\mathcal{E}'_r(W)$ as the sum of the cross term and the square of
Step~2 and of $\mathcal{E}_{2,r}(W)$. With this definition,
\eqref{8.2:eq-second-moment-1} and Steps~2--3 give \eqref{8.2:eq-second-moment-2}. The bound
\eqref{8.2:eq-second-moment-4} follows by inserting the corresponding display of Lemma~\ref{8.2:lem-extraction-one-point} into the
cross term and the square, and the remainder bound of Proposition~\ref{8.2:prop-gaussian-extraction}(b)
with $\kappa_{W,1}=\kappa_{W,2}=\kappa_W$, both taken in the form without the equivariance hypothesis
($\sigma_1=1$). The cross term
$2\Theta_{\mathrm{bg}}(W)\,\mathcal{E}_r(O_{s,p};W)$ is the origin of the member
$g_{k_0}^2\Theta_{\mathrm{bg}}(W)\Pi$. Under the equivariance hypothesis the same insertion, with the
sharpened forms of both remainder bounds, gives (e).

\emph{Step 5: flat $W$.} At flat $W$ every plaquette observable of $W$ vanishes. Since
$\Theta_{\mathrm{bg}}(W)=O_{s,p}(\mathrm{bg}_r(W))\ge0$ (each plaquette term
$1-\operatorname{Re}\operatorname{tr}U$ is non-negative) and the bound on $\Theta_{\mathrm{bg}}(W)$ in
Lemma~\ref{8.2:lem-extraction-one-point}(ii) vanishes at flat $W$, we get
$\Theta_{\mathrm{bg}}(W)=0$. The one-line term vanishes at flat $W$ by
Proposition~\ref{8.2:prop-gaussian-extraction}(b). Moreover $\kappa^{\mathfrak{B}}_W=0$, so
Lemma~\ref{8.2:lem-extraction-one-point}(i) gives $\mathfrak{g}^{\mathrm{step}}_{k_0,i}(W)=
\mathfrak{g}^{\mathrm{step}}_{k_0,i}$, and the third line of \eqref{8.2:eq-second-moment-2} has the value
found in Step~3. Put
$b=\tfrac n2\sum_{i<r}g_{k_0+i}^2\mathfrak{g}^{\mathrm{step}}_{k_0,i}$.
The main term is then
\begin{equation*}
b^2+\frac n2\Bigl(\frac 2n b\Bigr)^2=b^2+\frac 2n b^2
=\Bigl(1+\frac2n\Bigr)b^2,
\end{equation*}
which is \eqref{8.2:eq-second-moment-3}.
\end{proof}

\begin{corollary}[Two-point covariance averaged over joint translates]\label{8.2:cor-pair-covariance}
Let $s$ and $D$, an integer $r\ge 1$, the level $k_0=K-s$, the plaquettes $p_1,p_2$ and the box
$\mathfrak{B}_D$ containing both supports be as in the Gaussian extraction
(Proposition~\ref{8.2:prop-gaussian-extraction}), and write $n=\dim\mathfrak{g}=N^2-1$, so that
$b_0=n/2$. For $a=1,2$ let $G_a$ be either the smeared background action density
$X^{(k_0)}(f_a)$ of Definition~\ref{8.2:def-two-densities} or the plaquette observable
$O_{s,p_a}$ at depth $s$. In either case $G_a\in\mathfrak{G}_{k_0}(p_a;\ell,M,\varrho)$, as in
Proposition~\ref{8.2:prop-gaussian-extraction}.

Let $\nu^B_{K,s}$ be Ba{\l}aban's block law at depth $s$, and let $W$ be the field at level
$k_0+r$, distributed according to $\nu^B_{K,s-r}$. Let $\mathcal{C}=\mathcal{C}_{p_1,p_2}(D)$ be
the joint clean event. Let $G_a^{\langle h\rangle}(W)$ be the carried functions and
$\mathcal{T}^{(r)}(W)$ the carried conditional covariance of
\eqref{8.2:eq-gaussian-extraction-a}. Let $\Theta_a(W)$, the plaquette value at the $r$-step
background, and the curvature size $\kappa_{W,a}$ be those of that proposition for $p_a$.
Let $\mathcal{N}_a$ be the neighbourhood of the image of $p_a$ at level $k_0+r$.

For a translation $v$ of the depth-$s$ torus, write $(\tau_vf)(x)=f(x-v)$, and let $G_a^v$ be
$G_a$ with $p_a$ replaced by $p_a+v$, resp.\ with $f_a$ replaced by $\tau_vf_a$. Write
$\langle\!\langle\cdot\rangle\!\rangle_v$ for the average over $v$ in the depth-$s$ torus (or
over the class $\mathfrak{P}$ of translates). For a function $\varphi$ of a depth-$(s-r)$
plaquette $p''$, let $\operatorname{Av}_{p''}\varphi$ be its average over the translates of
$p''$ in $\mathfrak{P}$.

Assume the following.
\begin{itemize}
\item[(a)] The conclusions of Proposition~\ref{8.2:prop-gaussian-extraction} hold for $G_1,G_2$
on $\mathcal{C}$.
\item[(b)] The annealed one-point large-field block bound holds at every level $j\le K$ and for
every cube of the level-$j$ partition; this bound is not proved in this book. The lifetime bound
for live large-field regions holds.
\item[(c)] Clause~(0) of Theorem~0 holds.
\item[(d)] The transported curvature bounds hold for $a=1,2$:
\begin{equation}\label{8.2:eq-pair-covariance-1}
\begin{gathered}
\Theta_a(W)\le C_\eta L^{-4r}\,2N\max_{p''\in\mathcal{N}_a}O_{s-r,p''}(W),\\
\kappa_{W,a}\le B_3L^{-2r}\Bigl(2N\max_{p''\in\mathcal{N}_a}O_{s-r,p''}(W)\Bigr)^{1/2}.
\end{gathered}
\end{equation}
Here $C_\eta,B_3$ are the constants of those bounds. In addition, for $a=1,2$, the sum
$\mathfrak{g}_{k_0,a}(W)=\sum_{i<r}\mathfrak{g}^{\mathrm{step}}_{k_0,i}(W)$ of the one-point step
coefficients for $p_a$ at the background $W$, whose value at flat background is
$\mathfrak{g}_{k_0}$, satisfies, for a fixed exponent $C'$,
\begin{equation*}
\begin{split}
E_{\nu^B_{K,s-r}}\bigl|\mathfrak{g}_{k_0,a}(W)-\mathfrak{g}_{k_0}\bigr|^2
\le C\,L^{-4r}(1+\log x_{k_0})^{C'}\,
E_{\nu^B_{K,s-r}}\Bigl[\max_{p''\in\mathcal{N}_1\cup\mathcal{N}_2}O_{s-r,p''}(W)\Bigr].
\end{split}
\end{equation*}
\item[(e)] The flat-background value statement holds. It gives, for the smeared pair, the value
$b_0g^4_{k_0}\mathcal{K}_s^{(<r)}(f_1,f_2)$ of the two-line term of
\eqref{8.2:eq-gaussian-extraction-a}. Here $\mathcal{K}_s^{(<r)}$ denotes the two-line kernel
with the sum over steps restricted to $i<r$, that is, to the $r$ fluctuation steps at levels
$k_0,\dots,k_0+r-1$.
\item[(f)] The two-line kernel by which the members of the remainder below are divided is
bounded below with a constant $c_{\mathcal{K}}>0$, at the pairs considered (in the average over
$v$ for part~$(\alpha)$, at the pair itself for part~$(\beta)$). For the plaquette pair this
means $(\sum_{i<r}\mathfrak{g}^{\mathrm{step}}_{12,i})^2\ge c_{\mathcal{K}}$. For the smeared
pair it means
\begin{equation}\label{8.2:eq-pair-covariance-2}
c_{\mathcal{K}}\,\mathfrak{n}_s(f_1)\,\mathfrak{n}_s(f_2)\le C_W^8\,\mathcal{K}_s^{(<r)}(f_1,f_2),
\end{equation}
in which $\mathfrak{n}_s(f_a)$ denotes the normalisation of $f_a$ at depth $s$ fixed in
clause~(iii)(c) of Theorem~0, and $C_W$ is the constant of that clause for which the supports
of $f_1,f_2$ are separated by at least $C_W$ and at most $6C_W$ depth-$s$ cells; the kernel
lower bound used there is the same inequality with the
untruncated kernel $\mathcal{K}_s$ in place of $\mathcal{K}_s^{(<r)}$.
\end{itemize}
For part~$(\alpha)$ below we assume in addition $s_*\le s-r\le K-1$ and:
\begin{itemize}
\item[(g)] The upper bound of clause~(ii-d) of Theorem~0 holds, in class-average form, at every
depth $s'$ with $s_*\le s'\le K-1$; in particular, at $s'=s-r$,
\begin{equation}\label{8.2:eq-pair-covariance-3}
\operatorname{Av}_{p''}\,\nu^B_{K,s-r}(O_{s-r,p''})
\le C_{\mathrm{AP}}\,g^2_{k_0+r}\,p_0(g_{k_0+r})^2+C_{\mathrm{AP}}\,\delta^B(s-r).
\end{equation}
Here $p_0(\cdot)$ is the degree function, $p_0(g')=A_0(\log g'^{-2})^{p_0}$;
$C_{\mathrm{AP}}$ is the constant and $\delta^B(\cdot)$ the depth-dependent term of that
upper bound, and the depth term satisfies $\delta^B(s')\le CL^{-4s'}$.
\end{itemize}
For part~$(\beta)$ below we assume, in place of (g):
\begin{itemize}
\item[(g$'$)] The position-uniform a-priori bound holds:
\begin{equation}\label{8.2:eq-pair-covariance-4}
\sup_{p''}\,\nu^B_{K,s-r}(O_{s-r,p''})
\le C_{\mathrm{AP}}\,g^2_{k_0+r}\,p_0(g_{k_0+r})^2+C_{\mathrm{AP}}\,\delta^B(s-r),
\end{equation}
with $\delta^B(s-r)\le CL^{-4(s-r)}$. This is the upper bound of clause~(ii-d) with the class
average removed; it is not proved in this book. It would follow from
\eqref{8.2:eq-pair-covariance-3} if $\nu^B_{K,s-r}(O_{s-r,p''})$ were independent of $p''$
within its class of $\mathfrak{P}$-translates, an invariance of $\nu^B_{K,s-r}$ that is not
proved here.
\end{itemize}

Then the following hold.

\emph{The three-member decomposition.} We have
\begin{equation}\label{8.2:eq-pair-covariance-5}
\operatorname{Cov}_{\nu^B_{K,s}}(G_1,G_2)
=E_{\nu^B_{K,s-r}}\bigl[\mathcal{T}^{(r)}\,\mathbf{1}_{\mathcal{C}}\bigr]
+\mathcal{E}^{\mathrm{unc}}
+\operatorname{Cov}_{\nu^B_{K,s-r}}\bigl(G_1^{\langle h\rangle},G_2^{\langle h\rangle}\bigr),
\end{equation}
whose three members are as follows.
\begin{itemize}
\item[(1)] The main term of the first member is the two-line term of
\eqref{8.2:eq-gaussian-extraction-a}; the one-line term and the $W$-dependent $\kappa$-members
of $\mathcal{T}^{(r)}$ are carried into the remainder $R_0$ below. For
the plaquette pair the main term is
\begin{equation*}
\frac{n}{2}\,g^4_{k_0}\Bigl(\sum_{i<r}\mathfrak{g}^{\mathrm{step}}_{12,i}\Bigr)^2
\bigl(1+O(g^2_{k_0}rB^{\sharp})+O(\text{averaged $\kappa$-members})\bigr).
\end{equation*}
For the smeared pair it is $b_0g^4_{k_0}\mathcal{K}_s^{(<r)}(f_1,f_2)$ times a factor
$1+\cdots$ of the same kind.
\item[(2)] $\mathcal{E}^{\mathrm{unc}}$ is the unclean part, the contribution of histories
outside $\mathcal{C}$. It satisfies
\begin{equation*}
|\mathcal{E}^{\mathrm{unc}}|\le CM^2D^4(1+r)\,e^{-c\,p_0(g_{k_0})}.
\end{equation*}
\item[(3)] The last member is the covariance of the carried functions.
\end{itemize}

\emph{Part $(\alpha)$.} Under (a)--(g), averaging over simultaneous translates
$(p_1+v,p_2+v)$, resp.\ $(\tau_vf_1,\tau_vf_2)$, gives
\begin{equation}\label{8.2:eq-pair-covariance-6}
\bigl\langle\!\bigl\langle\operatorname{Cov}_{\nu^B_{K,s}}(G_1^v,G_2^v)\bigr\rangle\!\bigr\rangle_v
=b_0\,g^4_{k_0}\,\bigl\langle\!\bigl\langle\mathcal{K}^{v}\bigr\rangle\!\bigr\rangle_v\,(1+R_0).
\end{equation}
Here $\mathcal{K}^v$ is the two-line kernel at the translated pair: for the smeared pair
$\mathcal{K}_s^{(<r)}(\tau_vf_1,\tau_vf_2)$, for the plaquette pair
$(\sum_{i<r}\mathfrak{g}^{\mathrm{step}}_{12,i})^2$ at $(p_1+v,p_2+v)$. The remainder satisfies
\begin{equation}\label{8.2:eq-pair-covariance-7}
\begin{split}
|R_0|\le C\Bigl[\,&g^2_{k_0}\Pi_2+g_{k_0}p_0(g_{k_0})\Pi_2'
+L^{-4r}(1+\log x_{k_0})^{C'}\\
&+L^{-4r}g^2_{k_0}p_0(g_{k_0})^2(1+\log x_{k_0})^{C'}
+L^{-8r}x_{k_0}p_0(g_{k_0})^2\\
&+L^{-4r-4s}x^2_{k_0}+L^{-4s}(1+\log x_{k_0})^{C'}\\
&+D^4x^2_{k_0}(1+r)e^{-c\,p_0(g_{k_0})}+x^2_{k_0}e^{-c_{\mathrm{loc}}D}\Bigr].
\end{split}
\end{equation}
The constant $C$ depends on $M$, $C_W$ and the lower bound $c_{\mathcal{K}}$ of (f), and $C'$ is
a fixed exponent. The members of the
remainder are normalised by division by the two-line kernel. When the remainder of the Gaussian
extraction carries the exponent $\sigma_1=2$ (with equivariance), the member
$g_{k_0}p_0(g_{k_0})\Pi_2'$ is replaced by $g^2_{k_0}p_0(g_{k_0})^2\Pi_2'$.

\emph{Part $(\beta)$.} Under (a)--(f) and (g$'$), the identity \eqref{8.2:eq-pair-covariance-6}
holds pointwise in $(p_1,p_2)$, resp.\ $(f_1,f_2)$, without the average over $v$. The remainder
obeys the same bound \eqref{8.2:eq-pair-covariance-7}.

\emph{Floors.} Two members of \eqref{8.2:eq-pair-covariance-7} are controlled only under
one-sided floors. The member $L^{-8r}x_{k_0}p_0(g_{k_0})^2$ requires the extraction depth
\begin{equation}\label{8.2:eq-pair-covariance-8}
r=\lceil\theta\log_Lx_{k_0}\rceil\quad\text{with }\theta>\tfrac18 .
\end{equation}
With $s_{\mathrm{up}}(g)$ the depth parameter entering the upper bound on the running inverse
couplings $x_k$, the member $L^{-4r-4s}x^2_{k_0}$ satisfies
\begin{equation}\label{8.2:eq-pair-covariance-9}
L^{-4r-4s}x^2_{k_0}\le C\,L^{-4(s+r-2s_{\mathrm{up}}(g))}\bigl(1+B^{\sharp}(s+1)g^2\bigr)^2 ,
\end{equation}
and this requires the depth floor
\begin{equation}\label{8.2:eq-pair-covariance-10}
s\ge 2s_{\mathrm{up}}(g)-r+C .
\end{equation}

\emph{Use.} For the two-point witness theorem, the connected two-point function
$S^T_{2;K,m}(f_1,f_2)$ of the
lattice measure is invariant under joint translations of $(f_1,f_2)$. Consequently the averaged
form $(\alpha)$ applies to it with
no loss, provided the covariances above are averaged over joint translates before they are
used.

The relative bound $|R_0|\le 1/16+Cg^2_{k_0}$ used in clause~(iii)(c) of Theorem~0 is the form
$(\alpha)$ under the following choices: $r$ and $D$ as fixed for the two-point witness theorem,
$g\le\gamma_W$, and $s\ge s_W(g)$. The constant $1/16$ absorbs the truncation of the kernel at
$r$ steps and the loss in the constant $C_W$; for the smeared pair, the floor used there is the
kernel lower bound of clause~(iii)(c) for the untruncated kernel $\mathcal{K}_s$, and the
passage from $\mathcal{K}_s$ to $\mathcal{K}_s^{(<r)}$ in \eqref{8.2:eq-pair-covariance-2} is
part of the truncation absorbed in $1/16$.
\end{corollary}

\begin{proof}
\emph{The decomposition.} We apply the two-point tower identity, and then the law of total
covariance over the
randomness of $W$ that remains after the $r$ steps. This splits
$\operatorname{Cov}_{\nu^B_{K,s}}(G_1,G_2)$ into three pieces:
\begin{itemize}
\item the expectation under $\nu^B_{K,s-r}$ of the conditional covariance on clean histories,
which by hypothesis~(a) is $\mathcal{T}^{(r)}(W)$ of \eqref{8.2:eq-gaussian-extraction-a};
\item the contribution $\mathcal{E}^{\mathrm{unc}}$ of histories outside $\mathcal{C}$;
\item the covariance under $\nu^B_{K,s-r}$ of the conditional means, which are the carried
functions $G_a^{\langle h\rangle}$.
\end{itemize}
This is \eqref{8.2:eq-pair-covariance-5}.

\emph{The clean main term.} By hypothesis~(a), the main term of the first member of
\eqref{8.2:eq-pair-covariance-5} is the two-line term of
\eqref{8.2:eq-gaussian-extraction-a}; its one-line term and its $W$-dependent $\kappa$-members
are carried into $R_0$ and bounded below. For the plaquette pair that display gives the main
term in the
form (1), with relative corrections $O(g^2_{k_0}rB^{\sharp})$ and the averaged
$\kappa$-members. For the smeared pair, hypothesis~(e) identifies its value as
$b_0g^4_{k_0}\mathcal{K}_s^{(<r)}(f_1,f_2)$. Since $b_0=(N^2-1)/2=n/2$, the two forms carry
the same coefficient.

\emph{The unclean part.} Both $G_a$ are bounded in modulus by $M$. The contribution of the
histories outside $\mathcal{C}$ to the covariance is therefore at most $4M^2$ times the weight
of $\mathcal{C}^c$. By definition
$\mathcal{C}_{p_1,p_2}(D)=\mathcal{C}_{p_1}(D)\cap\mathcal{C}_{p_2}(D)$, so
\begin{equation*}
\mathcal{C}^c=\mathcal{C}_{p_1}(D)^c\cup\mathcal{C}_{p_2}(D)^c,
\end{equation*}
and the weight of $\mathcal{C}^c$ is at most the sum of the weights of the two complements.
Each of these is bounded by the weight bound \eqref{8.2:eq-clean-complement-a} for the
complement of the clean event. That bound is derived from the annealed one-point large-field
block bound at every level, the lifetime bound (both in hypothesis~(b)) and clause~(0) of
Theorem~0 (hypothesis~(c)). It gives $CD^4(1+r)e^{-c\,p_0(g_{k_0})}$ for each plaquette.
Hence
\begin{equation*}
|\mathcal{E}^{\mathrm{unc}}|\le 4M^2\cdot 2\,CD^4(1+r)e^{-c\,p_0(g_{k_0})},
\end{equation*}
which is (2) after renaming $C$.

\emph{The second moments.} The last member of \eqref{8.2:eq-pair-covariance-5} and the
$W$-averages of the $\kappa$-members and of the one-line term are second moments under
$\nu^B_{K,s-r}$. They are of four kinds:
\begin{itemize}
\item $E[\kappa_{W,1}\kappa_{W,2}]$;
\item $E[\Theta_1\Theta_2]$;
\item the expectation of $\Theta_a$ times the $\mathfrak{g}_{k_0,b}(W)$-term of
$G_b^{\langle h\rangle}$;
\item $\operatorname{Cov}_W(G^{\mathrm{main}}_1,G^{\mathrm{main}}_2)$.
\end{itemize}
Here $G^{\mathrm{main}}_a(W)$ is the main part of the carried function for $G_a$, and
$\mathfrak{g}_{k_0,b}(W)$ is the sum of hypothesis~(d).

To bound these moments we use hypothesis~(d), \eqref{8.2:eq-pair-covariance-1}, with $\max O$
standing for $\max_{p''\in\mathcal{N}_1\cup\mathcal{N}_2}O_{s-r,p''}(W)$. Products are
linearised as follows. One has $0\le O_{s-r,p''}\le 2$. Hence every power of $\max O$ of
degree at least one satisfies
\begin{equation*}
(\max O)^{\mathrm{deg}}=\max O\,(\max O)^{\mathrm{deg}-1}\le 2^{\mathrm{deg}-1}\max O .
\end{equation*}
For instance,
\begin{equation*}
\Theta_1\Theta_2\le C_\eta^2(2N)^2L^{-8r}(\max O)^2\le 2\,C_\eta^2(2N)^2L^{-8r}\max O,
\end{equation*}
and
\begin{equation*}
\kappa_{W,1}\kappa_{W,2}\le B_3^2\,2N\,L^{-4r}\max O .
\end{equation*}
In this way every one of the four moments is bounded by
\begin{equation}\label{8.2:eq-pair-covariance-11}
C\bigl(L^{-8r}+g^4_{k_0}L^{-4r}+g^2_{k_0}L^{-6r}\bigr)\,
E_{\nu^B_{K,s-r}}\Bigl[\max_{p''\in\mathcal{N}_1\cup\mathcal{N}_2}O_{s-r,p''}(W)\Bigr].
\end{equation}
What remains is to bound the expectation in \eqref{8.2:eq-pair-covariance-11}. Under
hypothesis~(g) this is done only after averaging over $v$, as follows.

\emph{Averaging over joint translates; part $(\alpha)$.} Since $O_{s-r,p''}\ge0$, the maximum
over $\mathcal{N}_1\cup\mathcal{N}_2$ is at most the sum over $p''\in\mathcal{N}_1\cup\mathcal{N}_2$.
Replace $(p_1,p_2)$ by $(p_1+v,p_2+v)$, resp.\ $(f_1,f_2)$ by $(\tau_vf_1,\tau_vf_2)$, and
average over $v$. The neighbourhoods are then translated jointly. The sum over the
neighbourhood becomes $|\mathcal{N}_1\cup\mathcal{N}_2|$ times an average over the depth-$(s-r)$
torus of the class averages $\operatorname{Av}_{p''}\nu^B_{K,s-r}(O_{s-r,p''})$. Each of these
is bounded by the upper bound \eqref{8.2:eq-pair-covariance-3} of clause~(ii-d), which applies
at depth $s-r$ by the assumption $s_*\le s-r\le K-1$. Hypothesis~(g) bounds only these class
averages. For this reason
the averages over simultaneous translates are taken on both sides of
\eqref{8.2:eq-pair-covariance-6}.

\emph{The infrared members.} We now measure each contribution against the main term
$b_0g^4_{k_0}$ times the kernel.
\begin{itemize}
\item \emph{Covariance of the main parts.} The covariance of the main parts is at most
$g^4_{k_0}$ times the larger of the second moments of $\mathfrak{g}_{k_0,a}(W)-\mathfrak{g}_{k_0}$,
$a=1,2$, times a power
of $\log x_{k_0}$, and by hypothesis~(d) each of these second moments is at most $L^{-4r}$ times
the expectation of $\max O$, up to a further power of $\log x_{k_0}$:
\begin{equation*}
\begin{split}
\operatorname{Cov}_W(G^{\mathrm{main}}_1,G^{\mathrm{main}}_2)
&\le g^4_{k_0}\,\max_{a=1,2}E\bigl|\mathfrak{g}_{k_0,a}(W)-\mathfrak{g}_{k_0}\bigr|^2
(1+\log x_{k_0})^{C'}\\
&\le g^4_{k_0}L^{-4r}\,E[\max O]\,(1+\log x_{k_0})^{C'}.
\end{split}
\end{equation*}
After the average over $v$ and hypothesis~(g), this gives relative to the main term the
members
\begin{equation*}
L^{-4r}\bigl(g^2_{k_0}p_0(g_{k_0})^2+L^{-4(s-r)}\bigr)(1+\log x_{k_0})^{C'}.
\end{equation*}
Here the right side of \eqref{8.2:eq-pair-covariance-3} is used in the form
$C(g^2_{k_0}p_0(g_{k_0})^2+L^{-4(s-r)})$, that is, with $\delta^B(s-r)\le CL^{-4(s-r)}$ from
hypothesis~(g) and, by
hypothesis~(c), $g^2_{k_0+r}p_0(g_{k_0+r})^2\le Cg^2_{k_0}p_0(g_{k_0})^2$.
These are the fourth and seventh members of \eqref{8.2:eq-pair-covariance-7}, since
$L^{-4r}L^{-4(s-r)}=L^{-4s}$.
\item \emph{The product of the plaquette values.} By the linearisation above,
\begin{equation*}
E[\Theta_1\Theta_2]\le CL^{-8r}\cdot 2\cdot E[\max O].
\end{equation*}
This member carries no factor $g^4_{k_0}$. Dividing by $g^4_{k_0}=x_{k_0}^{-2}$, it gives
relative to the main term, after the average over $v$ and hypothesis~(g), used in the same form,
\begin{equation*}
L^{-8r}\bigl(x_{k_0}p_0(g_{k_0})^2+x^2_{k_0}L^{-4(s-r)}\bigr)
=L^{-8r}x_{k_0}p_0(g_{k_0})^2+L^{-4r-4s}x^2_{k_0}.
\end{equation*}
These are the fifth and sixth members of \eqref{8.2:eq-pair-covariance-7}.
\item \emph{The cross terms.} The expectations of $\Theta_a$ times the
$\mathfrak{g}_{k_0,b}(W)$-term of $G_b^{\langle h\rangle}$ are bounded in the same way by
\eqref{8.2:eq-pair-covariance-11} and are smaller than the two members just derived.
\end{itemize}
The members carrying $x_{k_0}$ and $x^2_{k_0}$ arise because an absolute error is divided by
the main term, which carries the factor $g^4_{k_0}$.

\emph{Collection of the members.} The unclean part (2), divided by $b_0g^4_{k_0}$ times the
kernel, gives $D^4x^2_{k_0}(1+r)e^{-c\,p_0(g_{k_0})}$; the factors $M^2$ and $C_W^8$ are
absorbed in $C$, which therefore depends on $M$, $C_W$ and $c_{\mathcal{K}}$.
The remaining members of \eqref{8.2:eq-pair-covariance-7} are, as listed there, those of the
clean main term and of the remainder of the Gaussian extraction, measured against the same
quantity:
\begin{itemize}
\item the relative remainder $g^2_{k_0}\Pi_2+g_{k_0}p_0(g_{k_0})\Pi_2'$, or its replacement
when $\sigma_1=2$;
\item the member $L^{-4r}(1+\log x_{k_0})^{C'}$;
\item the member $x^2_{k_0}e^{-c_{\mathrm{loc}}D}$.
\end{itemize}
Every member is obtained by dividing by the two-line kernel. This division is controlled by
hypothesis~(f), which for the smeared pair is \eqref{8.2:eq-pair-covariance-2}; this is why
$C$ depends on $c_{\mathcal{K}}$. This proves $(\alpha)$.

\emph{Part $(\beta)$.} Hypothesis~(g$'$) bounds $\nu^B_{K,s-r}(O_{s-r,p''})$ uniformly in
$p''$. The expectation in \eqref{8.2:eq-pair-covariance-11} is therefore at most
$|\mathcal{N}_1\cup\mathcal{N}_2|$ times the right side of \eqref{8.2:eq-pair-covariance-4},
without any average over $v$. The derivation of the infrared members then goes through for
each fixed $(p_1,p_2)$, resp.\ $(f_1,f_2)$. The decomposition, the clean main term, the
unclean part and the division by the kernel do not involve the average. This proves
$(\beta)$.

\emph{Floors.} Suppose $r\ge\theta\log_Lx_{k_0}$, as in \eqref{8.2:eq-pair-covariance-8}. Then
$L^{-8r}\le x_{k_0}^{-8\theta}$, so
\begin{equation*}
L^{-8r}x_{k_0}p_0(g_{k_0})^2\le x_{k_0}^{1-8\theta}p_0(g_{k_0})^2 .
\end{equation*}
Since $p_0(g_{k_0})=A_0(\log x_{k_0})^{p_0}$, the right side tends to zero as
$x_{k_0}\to\infty$ when $\theta>1/8$. For the corresponding one-point statement, in the form
averaged over translates of $p$, any $\theta>0$ suffices. For the member
$L^{-4r-4s}x^2_{k_0}$, the bound \eqref{8.2:eq-pair-covariance-9} follows from the upper
bound on the running inverse couplings $x_k$ in which $s_{\mathrm{up}}(g)$ enters. It is small
under the depth floor \eqref{8.2:eq-pair-covariance-10}. Both floors are one-sided.

\emph{Use.} The invariance of $S^T_{2;K,m}$ under joint translations of $(f_1,f_2)$ follows
from the invariance of the lattice measure under lattice isometries
(Lemma~\ref{1.1:lem-invariance}).
\end{proof}

\begin{remark}[Pointwise statements and averages over translates]\label{8.2:rem-pointwise-averages}
The following four points concern the form in which the statements of this section hold.
\begin{enumerate}
\item[(a)] Lemma~\ref{8.2:lem-extraction-one-point} and Proposition~\ref{8.2:prop-gaussian-extraction}
are pointwise in $p$ (respectively in the pair $(p_1,p_2)$), and they hold history by history.
They are statements about the carried functions $G^{\langle h\rangle}(W)$ on the clean event
$\mathcal{C}_p(D)$ (respectively $\mathcal{C}_{p_1,p_2}(D)$), and no average over positions
enters them.
\item[(b)] Their error members that depend on $W$, and the restriction to the clean event, are
controlled in expectation over $W$ distributed according to $\nu^{B}_{K,s-r}$ only in averaged
form. The error members in question are the curvature sizes $\kappa_W$ and $\kappa^{\mathfrak{B}}_W$,
the value $\Theta_{\mathrm{bg}}(W)$ and the one-line term. Accordingly, the statements that use
them as inputs are stated in averaged form:
\begin{itemize}
\item the one-point witness bound for Ba{\l}aban's block law is stated for the full average
$\langle\!\langle\cdot\rangle\!\rangle_s$ over the depth-$s$ torus;
\item Corollary~\ref{8.2:cor-pair-covariance} is stated for the average
$\langle\!\langle\cdot\rangle\!\rangle$ over joint translates of the pair.
\end{itemize}
\item[(c)] For the lattice measure itself these averages cost nothing. The comparison of the block
law with the lattice measure (for the one-point statement) and the route of the two-point witness
theorem (for clause (iii-c) of Theorem~\ref{3.1:thm-main}) meet the block-law statements in the
averaged form without loss.
\item[(d)] Nothing in this section asserts or uses invariance of $\nu^{B}_{K,s}$ under the
translates in $\mathfrak{P}$.
\end{enumerate}
\end{remark}

\begin{proof}
Point (a) is the form of the two statements. Lemma~\ref{8.2:lem-extraction-one-point} and
Proposition~\ref{8.2:prop-gaussian-extraction} are formulated for a fixed plaquette $p$
(respectively a fixed pair $(p_1,p_2)$), a fixed history $h$ in the clean event and a fixed $W$
in the window. Their conclusions are identities and bounds for $G^{\langle h\rangle}(W)$
(respectively for the joint moment), and none of them involves an average.

For (b), write $O_{s-r,p''}$ for the plaquette observable $O_{s,p}$ that enters
Corollary~\ref{8.2:cor-pair-covariance}, taken at depth $s-r$ and plaquette $p''$.
Write $\mathcal{N}(p_{k_0+r})$ for the neighbourhood of the image plaquette $p_{k_0+r}$ of $p$ at
level $k_0+r$. In expectation over $W$ under $\nu^{B}_{K,s-r}$, the $W$-dependent error members
listed in (b) and the restriction to $\mathcal{C}_p(D)$ are controlled only through the sum
\begin{equation*}
\sum_{p''\in\mathcal{N}(p_{k_0+r})}E_{\nu^{B}_{K,s-r}}\bigl[O_{s-r,p''}(W)\bigr].
\end{equation*}
The only bound on these expectations available in this book is the upper bound of clause (ii-d)
of Theorem~0 (Theorem~\ref{3.1:thm-main}). That bound is stated for the class averages
$\operatorname{Av}_{p''}$ over the translates in $\mathfrak{P}$, at depths in $[s_*,K-1]$.

No statement of this book bounds a single-plaquette value $\nu^{B}_{K,s-r}(O_{s-r,p''})$. Such a
bound would follow from the class-average bound if $\nu^{B}_{K,s-r}$ were invariant under the
translates in $\mathfrak{P}$, and no such invariance is asserted here. The full average over $p$
removes this difficulty. It turns the sum displayed above into
\begin{equation*}
|\mathcal{N}(p_{k_0+r})|\cdot
\langle\!\langle E_{\nu^{B}_{K,s-r}}[O_{s-r,\cdot}(W)]\rangle\!\rangle_{s-r}.
\end{equation*}
This is an average of class averages $\operatorname{Av}_{p''}$. Each of them is bounded by the
upper bound of clause (ii-d) at depth $s-r$, which applies since the range hypothesis of
Lemma~\ref{8.2:lem-extraction-one-point} gives
\begin{equation*}
s_*\le s-r\le K-1 .
\end{equation*}

Equivalently, the one-point witness bound can be read for class averages $\operatorname{Av}_p$ in
place of the full average. For the pair, the same reasoning applies to both supports at once. It
therefore requires averaging over simultaneous translates $(p_1+v,p_2+v)$, with $v$ ranging over
the depth-$s$ torus, respectively over the pair of test functions $f_1,f_2$ both translated by the
same vector $v$, and this is the form of Corollary~\ref{8.2:cor-pair-covariance}.

For (c), we use the exact invariance of the lattice measure under translations of the torus. By
this invariance, the expectation $\langle O_{s,p}\rangle_K$ under the lattice measure is
independent of $p$, so
\begin{equation*}
\langle O_{s,p}\rangle_K=\langle\!\langle\langle O_{s,\cdot}\rangle_K\rangle\!\rangle_s .
\end{equation*}
By the same invariance, the connected two-point function $S^T_{2;K}$ is invariant under joint
translations of its two arguments. Hence an averaged block-law statement, combined with the
comparison of the block law with the lattice measure, yields the pointwise statement for the
lattice measure. Likewise, since $S^T_{2;K}(f_1,f_2)$ is unchanged by averaging over joint
translates, the two-point witness theorem can read Corollary~\ref{8.2:cor-pair-covariance} in its
averaged form without loss.

Point (d) follows from the three preceding points. Every bound used above is read either in
class-average form or for the lattice measure, whose translation invariance is exact. No step uses
invariance of $\nu^{B}_{K,s}$ under $\mathfrak{P}$.
\end{proof}

\begin{lemma}[Weight of the complement of the clean event]\label{8.2:lem-clean-complement}
Let $k_0$, $r$ and the plaquette $p$ be as in Definition~\ref{8.2:def-transported-laws}. Here and
below $C$ denotes a constant depending only on $(L,\mathfrak{p},N)$ and on the constants of
hypotheses (a), (b) and (d) below, not on $K$, $k_0$, $r$, $D$; it may change from line to line. Let
$\mathfrak{B}=\mathfrak{B}_D(p)$ be the box of side $D$ about $p$, write $\#\mathfrak{B}$ for the number
of level-$k_0$ cubes under $\mathfrak{B}$, and let $\mathcal{C}_p(D)$ be the clean event: the set of
$p$-clean histories with box of side $D$. It is defined by three conditions, which we use in the
following form. Condition (c1): at each level $j$ with $k_0\le j\le k_0+r$, no level-$j$ large-field
object of $h$ meets any of the at most $CD^4$ level-$j$ cubes under $\mathfrak{B}$. Condition (c2$'$):
at each of the $r$ levels $k$ of the clean steps, the count $\mathfrak{S}'_k$ of old large-field cubes
under $\mathfrak{B}$, created at levels $j<k_0$, satisfies $\mathfrak{S}'_k<\tau_*\,\#\mathfrak{B}$.
Condition (c3): no large-field region is live at level $k_0$ under $\mathfrak{B}$. For two plaquettes
$p_1,p_2$ the joint clean event is $\mathcal{C}_{p_1,p_2}(D)=\mathcal{C}_{p_1}(D)\cap\mathcal{C}_{p_2}(D)$.
The \emph{weight} of a set of histories $h$ is the sum of $\int d\tau_h$ over the
histories $h$ of the set, divided by the total $\sum_h\int d\tau_h$. Write
$q_j:=C_q\,e^{-c_q\,p_0(g_j)}$, with constants $C_q,c_q>0$ depending only on $(L,\mathfrak{p},N)$.
Assume the following.
\begin{enumerate}
\item[(a)] \emph{The annealed one-point large-field block bound.} For every $K$, every level $j\le K$
and every cube $\square$ of the level-$j$ partition of the depth-$(K-j)$ lattice,
\begin{equation}\label{8.2:eq-clean-complement-1}
\sum_{\substack{h:\ h \text{ carries a level-}j\text{ large-field}\\ \text{object meeting } \square}}
\int d\tau_h\;\le\; q_j\sum_h\int d\tau_h ,
\end{equation}
where the large-field objects are those of the list of large-field objects fixed together with the
clean event $\mathcal{C}_p(D)$. For finitely many disjoint cubes $\square_1,\dots,\square_m$, of
levels $j_1,\dots,j_m$, the product bound holds:
\begin{equation*}
\sum_{\substack{h:\ \text{for each } i,\ h \text{ carries a level-}j_i\\
\text{large-field object meeting } \square_i}}
\int d\tau_h\;\le\;\Bigl(\prod_{i=1}^m q_{j_i}\Bigr)\sum_h\int d\tau_h .
\end{equation*}
Moreover, for every set $S$, the expected number, under the normalised weights
$\int d\tau_h/\sum_{h'}\int d\tau_{h'}$, of level-$j$ large-field cubes of $h$ in $S$ is at
most $q_j$ times the number of level-$j$ cubes in $S$. The bound is annealed (it concerns the mass
$\sum_h\int d\tau_h$, not a conditional mass), concerns one cube, and holds at every level $j$.
\item[(b)] \emph{The lifetime bound}, with the lifetime $\mathfrak{l}_j$ and the lifetime radius
$\check{R}_j$ at level $j$, an offset constant $C_{\mathfrak{l}}$, constants $C'_{\mathfrak{l}}$ and
$C_R$, and the exponent $r_B>0$, all fixed with the lifetime bound. A
large-field region live at level $k_0$ under $\mathfrak{B}$ was
created at a level $\ge k_0-\mathfrak{l}_{k_0}-C_{\mathfrak{l}}$; this includes the case of a newer
creation inside such a region. At every level $j$ the lifetimes and lifetime radii satisfy
$\mathfrak{l}_j+1\le \check{R}_j+C'_{\mathfrak{l}}$ and $\check{R}_j\le C_R(\log x_j)^{r_B}$.
\item[(c)] \emph{Clause \textup{(0)} of Theorem~0} (Theorem~\ref{3.1:thm-main}), in the form: for the
levels $j$ with $k_0\le j\le k_0+r$ one has $x_j\ge x_{k_0}/2$.
\item[(d)] \emph{The ultraviolet-ward flow bound}, in the two forms used below:
$p_0(g_{k_0-\nu})\ge p_0(g_{k_0})(1-o(1))$ as $x_{k_0}\to\infty$, uniformly in $\nu\ge1$; and for
each $c'>0$ there are constants $C,c''>0$ with
\begin{equation}\label{8.2:eq-clean-complement-4}
\sum_{\nu\ge1}e^{-c'p_0(g_{k_0-\nu})}\;\le\;C\,e^{-c''p_0(g_{k_0})}.
\end{equation}
\item[(e)] The exponent $p_0$ of the degree function $p_0(\cdot)$ satisfies $p_0\ge 5r_B$; in
particular $p_0>r_B$.
\end{enumerate}
Then there are constants $C,c>0$, depending only on $(L,\mathfrak{p},N)$ and on the constants of
hypotheses (a), (b) and (d), and not on $K$, $k_0$, $r$, $D$, such that
\begin{equation}\label{8.2:eq-clean-complement-a}
\mathrm{weight}\bigl(\mathcal{C}_p(D)^c\bigr)\;\le\;C\,D^4(1+r)\,e^{-c\,p_0(g_{k_0})},
\end{equation}
and for a pair of plaquettes $p_1,p_2$ the joint clean event satisfies
\begin{equation}\label{8.2:eq-clean-complement-5}
\mathrm{weight}\bigl(\mathcal{C}_{p_1,p_2}(D)^c\bigr)\;\le\;2\,C\,D^4(1+r)\,e^{-c\,p_0(g_{k_0})}.
\end{equation}
\end{lemma}

\begin{proof}
Let $\mathcal{F}_1,\mathcal{F}_2,\mathcal{F}_3$ be the sets of histories at which
condition (c1), (c2$'$), (c3), respectively, fails. Since $\mathcal{C}_p(D)^c$ is contained in
$\mathcal{F}_1\cup\mathcal{F}_2\cup\mathcal{F}_3$ and the weight is subadditive,
\begin{equation}\label{8.2:eq-clean-complement-6}
\mathrm{weight}\bigl(\mathcal{C}_p(D)^c\bigr)\le
\mathrm{weight}(\mathcal{F}_1)+\mathrm{weight}(\mathcal{F}_2)+\mathrm{weight}(\mathcal{F}_3).
\end{equation}
We bound the three terms in turn.

\emph{The failure of} (c1). Condition (c1) concerns the levels $j$ with $k_0\le j\le k_0+r$, and at
each such level at most $CD^4$ cubes. A history in $\mathcal{F}_1$ carries, at one of these $r+1$
levels $j$, a level-$j$ large-field object meeting one of these cubes. Applying
\eqref{8.2:eq-clean-complement-1} to each of these cubes at each of these levels and summing,
\begin{equation*}
\mathrm{weight}(\mathcal{F}_1)\le C D^4 (r+1)\max_{k_0\le j\le k_0+r} q_j .
\end{equation*}
By hypothesis (c), every level $j$ of this range has $x_j\ge x_{k_0}/2$. Since
$p_0(g)=A_0(\log g^{-2})^{p_0}$, that is $p_0(g_j)=A_0(\log x_j)^{p_0}$, this gives
\begin{equation}\label{8.2:eq-clean-complement-3}
p_0(g_j)\;\ge\;A_0\bigl(\log x_{k_0}-\log 2\bigr)^{p_0}
\;=\;p_0(g_{k_0})\Bigl(1-\frac{\log 2}{\log x_{k_0}}\Bigr)^{p_0},
\end{equation}
so that $p_0(g_j)\ge p_0(g_{k_0})(1-o(1))$ as $x_{k_0}\to\infty$; once the term $o(1)$ is at
most $\tfrac12$, this gives $q_j\le C_q e^{-\frac12 c_q p_0(g_{k_0})}$. In the remaining range of
$x_{k_0}$, which is bounded, $p_0(g_{k_0})$ is bounded, and $q_j\le C_q\le C e^{-\frac12 c_q
p_0(g_{k_0})}$ after enlarging $C$. Absorbing $C_q$ into $C$,
\begin{equation}\label{8.2:eq-clean-complement-7}
\mathrm{weight}(\mathcal{F}_1)\le C D^4(r+1)\,e^{-\frac12 c_q\,p_0(g_{k_0})}.
\end{equation}

\emph{The failure of} (c2$'$). Condition (c2$'$) concerns the levels $j<k_0$, through the counts
$\mathfrak{S}'_k$ of old large-field cubes under $\mathfrak{B}$, one for each of the $r$ levels $k$
of the clean steps. Write $E$ for the annealed expectation of hypothesis (a), that is, the average
against the weights $\int d\tau_h$ normalised by
$\sum_h\int d\tau_h$. A large-field cube created at level $k_0-\nu$, $\nu\ge1$, enters
$\mathfrak{S}'_k$ with the volume $|\square|_k$ and with the multiplicity factor
$\mathfrak{l}_{k_0-\nu}+1$ of the lifetime bound (b). The expected-number form of hypothesis (a) at
the level $k_0-\nu$ therefore gives
\begin{equation*}
E\bigl[\mathfrak{S}'_k\bigr]\le\sum_{\nu\ge1}\bigl(\mathfrak{l}_{k_0-\nu}+1\bigr)\,q_{k_0-\nu}\cdot
\#\{\text{level-}(k_0-\nu)\text{ cubes under }\mathfrak{B}\}\cdot|\square|_k .
\end{equation*}
The number of level-$(k_0-\nu)$ cubes under $\mathfrak{B}$ is $\#\mathfrak{B}$ times the footprint
factor $L^{4\nu}$, and, measuring volumes in units in which a level-$k_0$ cube has side one, the
volume is $|\square|_k=L^{-4\nu}$; the
footprint factor $L^{4\nu}$ is cancelled exactly by the factor $L^{-4\nu}$ of the volume. With
$\mathfrak{l}_j+1\le\check{R}_j+C'_{\mathfrak{l}}$ from (b) and $C'_{\mathfrak{l}}\le C$, this yields
\begin{equation*}
E\bigl[\mathfrak{S}'_k\bigr]\le\#\mathfrak{B}\cdot
\sum_{\nu\ge1}\bigl(\check{R}_{k_0-\nu}+C\bigr)\,q_{k_0-\nu}.
\end{equation*}
By the second inequality of (b), $\check{R}_j+C\le C_R(\log x_j)^{r_B}+C$. Since $p_0>r_B$ by
hypothesis (e), the function $t\mapsto (C_Rt^{r_B}+C)\,e^{-\frac12 c_qA_0t^{p_0}}$ is bounded on
$t\ge0$, and therefore, with $c':=\tfrac12 c_q$,
\begin{equation*}
\bigl(\check{R}_j+C\bigr)\,q_j\le C\,e^{-c'p_0(g_j)} .
\end{equation*}
The sum over $\nu\ge1$ runs in the ultraviolet-ward direction, towards the levels below $k_0$, which
is the direction in which hypothesis (d) applies; by the second form
\eqref{8.2:eq-clean-complement-4} of hypothesis (d),
\begin{equation*}
\sum_{\nu\ge1}e^{-c'p_0(g_{k_0-\nu})}\le C\,e^{-c''p_0(g_{k_0})}.
\end{equation*}
Hence $E[\mathfrak{S}'_k]\le C\,\#\mathfrak{B}\,e^{-c''p_0(g_{k_0})}$. Markov's inequality applied to
$\mathfrak{S}'_k$ at the threshold $\tau_*\,\#\mathfrak{B}$ of condition (c2$'$) gives, at the level $k$,
\begin{equation*}
\mathrm{weight}\bigl(\mathfrak{S}'_k\ge\tau_*\,\#\mathfrak{B}\bigr)
\le\frac{E\bigl[\mathfrak{S}'_k\bigr]}{\tau_*\,\#\mathfrak{B}}
\le\frac{C\,e^{-c''p_0(g_{k_0})}}{\tau_*}.
\end{equation*}
Summing over the $r$ levels $k$ and absorbing the factor $1/\tau_*$ into the constant,
\begin{equation}\label{8.2:eq-clean-complement-8}
\mathrm{weight}(\mathcal{F}_2)\le C\,r\,e^{-c''p_0(g_{k_0})}.
\end{equation}

\emph{The failure of} (c3). A history in $\mathcal{F}_3$ carries a large-field region live at level
$k_0$ under $\mathfrak{B}$. By the lifetime bound (b), including its clause on a newer creation inside
the region, such a region was created at a level $k_0-\nu$ with
$0\le\nu\le\mathfrak{l}_{k_0}+C_{\mathfrak{l}}$. At
relative level $\nu$ the region can occupy at most $CD^4L^{4\nu}$ positions under $\mathfrak{B}$.
By the first form of hypothesis (d), which is uniform in $\nu\ge1$, there is a threshold for
$x_{k_0}$ above which the term $o(1)$ is at most $\tfrac12$ for all $\nu\ge1$ simultaneously; there
$p_0(g_{k_0-\nu})\ge\tfrac12\,p_0(g_{k_0})$, hence $q_{k_0-\nu}\le C_q e^{-\frac12 c_q p_0(g_{k_0})}$
for every $\nu\ge1$, and for $\nu=0$ this holds trivially. In the remaining range of $x_{k_0}$, which
is bounded, $p_0(g_{k_0})$ is bounded, and $q_{k_0-\nu}\le C_q\le Ce^{-\frac12 c_q p_0(g_{k_0})}$
after enlarging $C$. Applying \eqref{8.2:eq-clean-complement-1} to each position, using
this bound on $q_{k_0-\nu}$, and bounding each $L^{4\nu}$ by its largest value,
\begin{align*}
\mathrm{weight}(\mathcal{F}_3)
&\le C D^4\sum_{\nu\le \mathfrak{l}_{k_0}+C_{\mathfrak{l}}}L^{4\nu}\,q_{k_0-\nu}\\
&\le C D^4\,\bigl(\mathfrak{l}_{k_0}+C_{\mathfrak{l}}+1\bigr)\,
L^{4(\mathfrak{l}_{k_0}+C_{\mathfrak{l}})}\,e^{-\frac12 c_q\,p_0(g_{k_0})}.
\end{align*}
By the first inequality of (b),
$\mathfrak{l}_{k_0}+C_{\mathfrak{l}}\le\check{R}_{k_0}+C'_{\mathfrak{l}}+C_{\mathfrak{l}}$, and by the
second, $\check{R}_{k_0}\le C_R(\log x_{k_0})^{r_B}$. Using
$\log(\mathfrak{l}_{k_0}+C_{\mathfrak{l}}+1)\le\mathfrak{l}_{k_0}+C_{\mathfrak{l}}$,
\begin{align*}
\log\Bigl[\bigl(\mathfrak{l}_{k_0}+C_{\mathfrak{l}}+1\bigr)\,
L^{4(\mathfrak{l}_{k_0}+C_{\mathfrak{l}})}\Bigr]
&\le(1+4\log L)\bigl(\mathfrak{l}_{k_0}+C_{\mathfrak{l}}\bigr)\\
&\le C(1+4\log L)\bigl((\log x_{k_0})^{r_B}+1\bigr).
\end{align*}
Since $p_0(g_{k_0})=A_0(\log x_{k_0})^{p_0}$ and $p_0\ge 5r_B>r_B$ by hypothesis (e), there is a
constant $C'$, depending only on the constants above, such that for all $x_{k_0}$ the right-hand side
is at most $\tfrac14 c_qA_0(\log x_{k_0})^{p_0}+C'=\tfrac14 c_q\,p_0(g_{k_0})+C'$. Hence the factor
$(\mathfrak{l}_{k_0}+C_{\mathfrak{l}}+1)\,L^{4(\mathfrak{l}_{k_0}+C_{\mathfrak{l}})}$ is absorbed by
part of the exponential, and with $c''':=\tfrac14 c_q$,
\begin{equation}\label{8.2:eq-clean-complement-9}
\mathrm{weight}(\mathcal{F}_3)\le C D^4\,e^{-c'''p_0(g_{k_0})}.
\end{equation}

\emph{Total.} Insert \eqref{8.2:eq-clean-complement-7}, \eqref{8.2:eq-clean-complement-8} and
\eqref{8.2:eq-clean-complement-9} into \eqref{8.2:eq-clean-complement-6}. Put
$c:=\min\{\tfrac12 c_q,\,c'',\,c'''\}$. Each of the three bounds is then at most a constant times
$D^4(1+r)e^{-c\,p_0(g_{k_0})}$, the side $D$ of the box being at least one. This is
\eqref{8.2:eq-clean-complement-a}. For the pair $p_1,p_2$, since
$\mathcal{C}_{p_1,p_2}(D)=\mathcal{C}_{p_1}(D)\cap\mathcal{C}_{p_2}(D)$, the complement
$\mathcal{C}_{p_1,p_2}(D)^c$ is contained in $\mathcal{C}_{p_1}(D)^c\cup\mathcal{C}_{p_2}(D)^c$;
subadditivity of the weight and \eqref{8.2:eq-clean-complement-a} at $p_1$ and at $p_2$ bound its
weight by twice the right-hand side of \eqref{8.2:eq-clean-complement-a}; this is
\eqref{8.2:eq-clean-complement-5}.
\end{proof}

\begin{corollary}[The one-point witness under the block law]\label{8.2:cor-witness-block}
Work under the quantifier prefix of Theorem~0 in the form it takes inside clause (iv)
(Notation~\ref{2.1:not-prefix}), with
$g_K\ge g_-(g)$, and assume:
\begin{itemize}
\item[(a)] the statements from which Lemma~\ref{8.2:lem-extraction-one-point} (Gaussian extraction
for one local observable) is proved, so that the corresponding displays and the clauses of that lemma, among them
the corresponding display of Lemma~\ref{8.2:lem-extraction-one-point}, are available;
\item[(b)] the large-field block bound for one plaquette (whose proof is outstanding);
\item[(c)] the lifetime statement for large-field regions;
\item[(d)] the upper bound of clause (ii-d) of Theorem~0;
\item[(e)] clause (0) of Theorem~0;
\item[(f)] the flat-background identification of the step coefficients, which fixes the value of
the coefficient $b^{(k)}$ below.
\end{itemize}
Let $O_{s,p}$, the level $k_0$, $n=\dim\mathfrak{g}$, $\ell_0$, $D_0$, $\Pi_{r,D}$, the range
condition and the constant $M_O$ (the value of $M$ for $G=O_{s,p}$) be as in
Lemma~\ref{8.2:lem-extraction-one-point}, and put
\begin{equation*}
b^{(k)}:=\frac{n}{2}\,\mathfrak{c}_k ,
\end{equation*}
where the $\mathfrak{c}_k$ are the level-$k$ Gaussian coefficients of the lemma that introduces them;
they satisfy $b^{(0)}=n/4$ and $c_b\,n/2\le b^{(k)}\le n/2$ for every $k$, with a constant $c_b>0$ of
that lemma, and $\mathfrak{c}_k\to\mathfrak{c}_\infty$ as $k\to\infty$ with
$\mathfrak{c}_\infty\in[8/81,4/9]$.
Then there are constants $C_W,c_W>0$, depending only on $(L,\mathfrak{p},N)$, such that for every
$K$, $s$, $r$, $D$ satisfying the range condition and $D\ge 1+r\ell_0+D_0$,
\begin{equation}\label{8.2:eq-witness-block-1}
\langle\!\langle\nu^{B}_{K,s}(O_{s,p})\rangle\!\rangle_s
=b^{(k_0)}\,g^2_{k_0}\,\bigl(1+R^B_{K,s}\bigr),
\end{equation}
where
\begin{equation}\label{8.2:eq-witness-block-a}
\begin{split}
|R^B_{K,s}|\le C_W\Bigl[\,&g_{k_0}^{\sigma_1}\Pi_{r,D}+g^2_{k_0}
+L^{-2r}g_{k_0}\,p_0(g_{k_0})\,D\,\operatorname{polylog}(x_{k_0})
+L^{-4r}p_0(g_{k_0})^2+L^{-2s}\\
&+L^{-4(s-s_{\mathrm{up}}(g))}\bigl(1+B^{\sharp}(s+1)g^2\bigr)
+D^4x_{k_0}(1+r)\,e^{-c_Wp_0(g_{k_0})}+x_{k_0}\,e^{-c_WD}\Bigr].
\end{split}
\end{equation}
Here $\sigma_1=1$, or $\sigma_1=2$ in the situation of clause (iii) of
Lemma~\ref{8.2:lem-extraction-one-point}; $p_0(\cdot)$ is the degree function;
$\operatorname{polylog}(x_{k_0})$ denotes a factor bounded by a fixed polynomial in $\log x_{k_0}$;
and $s_{\mathrm{up}}(g)$ is the depth attached to the reference coupling $g$ in the upper flow
bound on the running couplings.
\end{corollary}

The average
$\langle\!\langle\cdot\rangle\!\rangle_s$ is the full average over the plaquettes $p$ of the depth-$s$
torus; the right-hand side of \eqref{8.2:eq-witness-block-1} does not depend on $p$.

\begin{proof}
The proof has four steps. The constants $C$ and $c$ below depend only on $(L,\mathfrak{p},N)$ and
may change from line to line.

\emph{Step 1: the starting representation.} We write $\nu^B_{K,s}(O_{s,p})$ in the form in which
Lemma~\ref{8.2:lem-extraction-one-point} is stated, by the representation of the block law over
$r$ of Ba{\l}aban's steps on which that lemma rests: as an expectation, over the level-$(k_0+r)$
field $W\sim\nu^B_{K,s-r}$ and over the histories $h$ of the $r$ steps between depth $s-r$ and
depth $s$, of the carried function $O_{s,p}^{\langle h\rangle}(W)$. We split this expectation
according to whether $h$ lies in the clean event $\mathcal{C}_p(D)$ or in its complement.

\emph{Step 2: histories outside the clean event.} By
Lemma~\ref{8.2:lem-clean-complement} the weight of the complement of the clean event satisfies
\eqref{8.2:eq-clean-complement-a}:
\begin{equation*}
\operatorname{weight}\bigl(\mathcal{C}_p(D)^c\bigr)\le C\,D^4(1+r)\,e^{-c\,p_0(g_{k_0})}.
\end{equation*}
Hypotheses (b), the large-field block bound for one plaquette, and (c), the lifetime statement
for large-field regions, are, respectively, a hypothesis of the representation used in Step~1 and
a hypothesis of Lemma~\ref{8.2:lem-clean-complement}, as stated there. With this
weight bound and hypothesis (b), the histories outside
$\mathcal{C}_p(D)$ contribute to $\langle\!\langle\nu^{B}_{K,s}(O_{s,p})\rangle\!\rangle_s$ at most
\begin{equation*}
C\,M_O\,D^4(1+r)\,e^{-c\,p_0(g_{k_0})}.
\end{equation*}

\emph{Step 3: histories in the clean event.} For $h\in\mathcal{C}_p(D)$ and $W$ in the
level-$(k_0+r)$ window we apply
Lemma~\ref{8.2:lem-extraction-one-point} with $G=O_{s,p}$, $M=M_O$ and $\ell=1$; its condition
$D\ge\ell+r\ell_0+D_0$ is then the hypothesis $D\ge1+r\ell_0+D_0$ of the corollary, and the range
condition is assumed. Its tube condition $\varrho\ge\varrho_{\mathcal{P}}$ is met with
$\varrho=\varrho_{\mathcal{P}}$: with $\operatorname{Re}\operatorname{tr}U(\partial p)$ continued as
$\tfrac12\bigl(\operatorname{tr}U(\partial p)+\operatorname{tr}U(\partial p)^{-1}\bigr)$, and with the
inverse of a matrix of determinant one equal to its adjugate, $O_{s,p}$ is a polynomial in the
entries of the bond variables, hence holomorphic on every tube. By
the corresponding display of Lemma~\ref{8.2:lem-extraction-one-point} and the
identification of the main
term for $G=O_{s,p}$ in that lemma,
\begin{equation}\label{8.2:eq-witness-block-2}
O_{s,p}^{\langle h\rangle}(W)=\Theta_{\mathrm{bg}}(W)
+\frac{n}{2}\sum_{i=0}^{r-1}g^2_{k_0+i}\,\mathfrak{g}^{\mathrm{step}}_{k_0,i}(W)
+\mathcal{E}_r(O_{s,p};W),
\end{equation}
with $\mathcal{E}_r$ bounded by the corresponding display of Lemma~\ref{8.2:lem-extraction-one-point} (with $\sigma_1=2$ and the
$\kappa$-member of the corresponding display of Lemma~\ref{8.2:lem-extraction-one-point} modified as in clause (iii) of the lemma,
in the situation of that clause). We treat the members of
\eqref{8.2:eq-witness-block-2} in four parts.

(3a) \emph{The Gaussian term at flat $W$.} Write $\mathfrak{g}^{\mathrm{step}}_{k_0,i}$ for the value
of $\mathfrak{g}^{\mathrm{step}}_{k_0,i}(W)$ at flat $W$. Clause (0) of Theorem~0 controls the running
of the coupling across the $r$ levels:
\begin{equation*}
\Bigl|\frac{g^2_{k_0+i}}{g^2_{k_0}}-1\Bigr|\le\bigl(2B^{\sharp}i+c_5\bigr)\,g^2_{k_0}
\qquad(0\le i<r).
\end{equation*}
By the lemma on the sum of the flat step coefficients, recalled in
Remark~\ref{8.2:rem-gaussian-exact}, together with the flat-background identification (f), the
sum of the flat step coefficients over the first $r$ steps is $\mathfrak{c}_{k_0}$ up to the tail
of that sum:
\begin{equation*}
\sum_{i=0}^{r-1}\mathfrak{g}^{\mathrm{step}}_{k_0,i}=\mathfrak{c}_{k_0}-\delta_r,
\qquad|\delta_r|\le C_\eta L^{-4r},
\end{equation*}
where $\delta_r$ denotes this tail and $C_\eta$ is the constant of that lemma; the same constant
bounds the step coefficients in clauses (i) and (ii) of
Lemma~\ref{8.2:lem-extraction-one-point}. Hence, with $b^{(k_0)}=(n/2)\mathfrak{c}_{k_0}$,
\begin{equation}\label{8.2:eq-witness-block-3}
\frac{n}{2}\sum_{i=0}^{r-1}g^2_{k_0+i}\,\mathfrak{g}^{\mathrm{step}}_{k_0,i}
=b^{(k_0)}g^2_{k_0}\Bigl[1+\frac{1}{\mathfrak{c}_{k_0}}
\Bigl(-\delta_r+\sum_{i=0}^{r-1}\Bigl(\frac{g^2_{k_0+i}}{g^2_{k_0}}-1\Bigr)
\mathfrak{g}^{\mathrm{step}}_{k_0,i}\Bigr)\Bigr].
\end{equation}
By clause (i) of Lemma~\ref{8.2:lem-extraction-one-point},
$0\le\mathfrak{g}^{\mathrm{step}}_{k_0,i}\le C_\eta L^{-4i}$, so that
\begin{equation*}
\Bigl|\sum_{i=0}^{r-1}\Bigl(\frac{g^2_{k_0+i}}{g^2_{k_0}}-1\Bigr)
\mathfrak{g}^{\mathrm{step}}_{k_0,i}\Bigr|
\le g^2_{k_0}\,C_\eta\sum_{i\ge0}\bigl(2B^{\sharp}i+c_5\bigr)L^{-4i},
\end{equation*}
and the series on the right converges. Finally $\mathfrak{c}_{k_0}=2b^{(k_0)}/n\ge c_b$ by the lemma on
the Gaussian coefficients, so the bracket in \eqref{8.2:eq-witness-block-3} differs from $1$ by at
most $C(L^{-4r}+g^2_{k_0})$. The passage from $\mathfrak{g}^{\mathrm{step}}_{k_0,i}(W)$ to its flat
value is controlled by clause (i) of the lemma:
\begin{equation*}
\bigl|\mathfrak{g}^{\mathrm{step}}_{k_0,i}(W)-\mathfrak{g}^{\mathrm{step}}_{k_0,i}\bigr|
\le C_\eta L^{-4i}\,\kappa^{\mathfrak{B}}_W\,D\,\operatorname{polylog}(x_{k_0}).
\end{equation*}

(3b)--(3c) \emph{The members depending on $W$.} The members of
\eqref{8.2:eq-witness-block-2} and of the corresponding display of Lemma~\ref{8.2:lem-extraction-one-point} that depend on $W$ are
those carrying the curvature sizes $\kappa_W$, $\kappa^{\mathfrak{B}}_W$ (among them the difference
just displayed, and the member of the corresponding display of Lemma~\ref{8.2:lem-extraction-one-point} carrying the factor
$g^2_{k_0}\Pi'_{r,D}\varrho^{-2}\kappa_W$) and the background value $\Theta_{\mathrm{bg}}(W)$. By
clause (ii) of Lemma~\ref{8.2:lem-extraction-one-point},
\begin{equation*}
\Theta_{\mathrm{bg}}(W)\le C_\eta L^{-4r}\cdot 2N
\max_{p''\in\mathcal{N}(p_{k_0+r})}O_{s-r,p''}(W),
\end{equation*}
and, since $|U-\mathbf{1}|^2_{\mathrm{HS}}=2N(1-\operatorname{Re}\operatorname{tr}U)$ for
$U\in\mathrm{SU}(N)$ with $\operatorname{tr}\mathbf{1}=1$, the Cauchy--Schwarz inequality bounds the
expectations of the members linear in $|W(\partial p'')-\mathbf{1}|$ by square roots of expectations
of $O_{s-r,p''}(W)$. Both kinds of $W$-dependent member are thereby controlled in expectation over
$W\sim\nu^B_{K,s-r}$ through the expectations of $O_{s-r,p''}(W)$, for $p''$ in
$\mathcal{N}(p_{k_0+r})$ (for $\kappa_W$ and $\Theta_{\mathrm{bg}}$) and for $p''$ in the
level-$(k_0+r)$ image of the $D$-boxes (for $\kappa^{\mathfrak{B}}_W$).

(3d) \emph{The full average.} Lemma~\ref{8.2:lem-extraction-one-point} holds for each plaquette and
each clean history separately, and a single-plaquette value $\nu^B_{K,s-r}(O_{s-r,p''})$ is not
controlled by the hypotheses. We therefore take the full average $\langle\!\langle\cdot\rangle\!\rangle_s$
over $p$. We bound the maximum over $p''\in\mathcal{N}(p_{k_0+r})$ in clause (ii) by the sum over
$p''\in\mathcal{N}(p_{k_0+r})$, each $O_{s-r,p''}(W)$ being non-negative. As set out in
Remark~\ref{8.2:rem-pointwise-averages}, the full average turns this sum into
$|\mathcal{N}|$ times a depth-$(s-r)$ average, which is an
average of class averages $\operatorname{Av}_{p''}$, and each class average is bounded by the
upper bound of clause (ii-d) of Theorem~0 at depth $s-r\in[s_*,K-1]$. This gives, for the
average of the contribution of the background value $\Theta_{\mathrm{bg}}$, divided by the
normaliser $\mathfrak{m}$ of that contribution, the bound
\begin{equation*}
C\bigl[L^{-4r}p_0(g_{k_0})^2+L^{-4s}x_{k_0}\bigr],
\end{equation*}
and, for the contribution of the curvature members, the bound
\begin{equation*}
C\bigl[L^{-2r}g_{k_0}\,p_0(g_{k_0})+L^{-2s}\bigr].
\end{equation*}
The term $L^{-4s}x_{k_0}$ is estimated by the upper flow bound together with the lower bound
$g_K\ge g_-(g)$:
\begin{equation*}
L^{-4s}x_{k_0}\le L^{-4(s-s_{\mathrm{up}}(g))}\,L^{-4C_{\mathrm{up}}}
\bigl(1+B^{\sharp}(s+1)g^2\bigr),
\end{equation*}
where $C_{\mathrm{up}}$ is the constant of the upper flow bound.

\emph{Step 4: assembly.} We add the contribution of Step~2 and the averaged contributions of
Step~3: the main term \eqref{8.2:eq-witness-block-3} with the estimates of (3a), the
$W$-dependent members estimated in (3b)--(3d), and the average of $\mathcal{E}_r(O_{s,p};W)$
bounded by the corresponding display of Lemma~\ref{8.2:lem-extraction-one-point}. The term $CL^{-4r}$ of (3a) is absorbed into the
member $C_WL^{-4r}p_0(g_{k_0})^2$ of \eqref{8.2:eq-witness-block-a}, since $p_0(g_{k_0})\ge1$ for the
couplings of the prefix. Dividing by $b^{(k_0)}g^2_{k_0}$, as the
definition of $R^B_{K,s}$ through \eqref{8.2:eq-witness-block-1} requires, and using
$b^{(k_0)}\ge c_b\,n/2$ and $g^{-2}_{k_0}=x_{k_0}$, the sum of these bounds is bounded by the
right-hand side of \eqref{8.2:eq-witness-block-a}, with constants $C_W,c_W$ depending only on
$(L,\mathfrak{p},N)$. This proves \eqref{8.2:eq-witness-block-1} and
\eqref{8.2:eq-witness-block-a}.
\end{proof}

\begin{corollary}[The one-point witness in the true measure]\label{8.2:cor-witness-true}
Work under the quantifier prefix and with the constants $C_W,c_W>0$ of
Corollary~\ref{8.2:cor-witness-block}, and let $O_{s,p}$, $k_0$, $b^{(k)}$, $r$, $D$, $s_*$, $k_*$ and the
full average
$\langle\!\langle\cdot\rangle\!\rangle_s$ over the depth-$s$ torus be as there. Write
$\langle\cdot\rangle_K$ for the expectation, in the measure of Definition~\ref{1.1:def-wilson-action}
on the bare lattice of spacing $L^{-K}$ (the true measure), of functions of the depth-$s$ block
field obtained from the bare configuration through the block-averaging map of
Definition~\ref{1.5:def-block-average}, as opposed to Ba{\l}aban's block law
$\nu^{B}_{K,s}$. Assume that, for every $K,s,r,D$ satisfying the hypotheses of
Corollary~\ref{8.2:cor-witness-block} with $s_*+r\le s\le K-k_*$:
\begin{itemize}
\item[(a)] the one-point witness under the block law, Corollary~\ref{8.2:cor-witness-block}, holds;
\item[(b)] translation invariance of the true measure: for every plaquette $p$ of the depth-$s$
lattice,
\begin{equation}\label{8.2:eq-witness-true-1}
\langle O_{s,p}\rangle_K=\langle\!\langle \langle O_{s,\cdot}\rangle_K\rangle\!\rangle_s ;
\end{equation}
\item[(c)] the comparison of the true measure with the block law applies to the observable
$O_{s,\cdot}$, its only requirement on an observable $F$ being a bound on the quantity
$x^{-1-\sigma}\|F\|_{C^2}$, with $x$, $\sigma$ and the norm $\|\cdot\|_{C^2}$ as in the comparison
of the true measure with the block law; applied to $O_{s,\cdot}$, it yields for the full averages
of $O_{s,\cdot}$ in the two measures
\begin{equation}\label{8.2:eq-witness-true-2}
\bigl|\langle\!\langle \langle O_{s,\cdot}\rangle_K\rangle\!\rangle_s
-\langle\!\langle \nu^{B}_{K,s}(O_{s,\cdot})\rangle\!\rangle_s\bigr|
\le b^{(k_0)}g_{k_0}^2\, C_W\, x_{k_0}\,e^{-c_W p_0(g_{k_0})} ,
\end{equation}
where $p_0(\cdot)$ is the degree function (recall $g_{k_0}^2x_{k_0}=1$).
\end{itemize}
Then for every such $K,s,r,D$ and for every plaquette $p$ of the depth-$s$ lattice,
\begin{equation}\label{8.2:eq-witness-true-3}
\langle O_{s,p}\rangle_K=b^{(k_0)}g_{k_0}^2\,(1+R_{K,s}),
\end{equation}
where $R_{K,s}$ does not depend on $p$ and $|R_{K,s}|$ is at most the right-hand side of
\eqref{8.2:eq-witness-block-a} plus $C_W\,x_{k_0}\,e^{-c_W p_0(g_{k_0})}$.
\end{corollary}

\begin{proof}
Fix $K,s,r,D$ as in the statement and a plaquette $p$ of the depth-$s$ lattice. By hypothesis~(b),
the value $\langle O_{s,p}\rangle_K$ equals the full average
$\langle\!\langle \langle O_{s,\cdot}\rangle_K\rangle\!\rangle_s$, which does not depend on $p$.

By hypothesis~(a), that is, Corollary~\ref{8.2:cor-witness-block}, whose lower bound on $b^{(k)}$
gives $b^{(k_0)}>0$, we have
\begin{equation*}
\langle\!\langle \nu^{B}_{K,s}(O_{s,\cdot})\rangle\!\rangle_s
=b^{(k_0)}g_{k_0}^2\,(1+R^{B}_{K,s}),
\end{equation*}
with $|R^{B}_{K,s}|$ at most the right-hand side of \eqref{8.2:eq-witness-block-a}.
Define $R_{K,s}$ by \eqref{8.2:eq-witness-true-3}; this is possible because
$b^{(k_0)}g_{k_0}^2>0$, and by the first paragraph $R_{K,s}$ is independent of $p$. Subtracting the last
display from \eqref{8.2:eq-witness-true-3} and using \eqref{8.2:eq-witness-true-1},
\begin{equation*}
b^{(k_0)}g_{k_0}^2\,\bigl(R_{K,s}-R^{B}_{K,s}\bigr)
=\langle\!\langle \langle O_{s,\cdot}\rangle_K\rangle\!\rangle_s
-\langle\!\langle \nu^{B}_{K,s}(O_{s,\cdot})\rangle\!\rangle_s .
\end{equation*}
Dividing by $b^{(k_0)}g_{k_0}^2$ and inserting \eqref{8.2:eq-witness-true-2} gives
$|R_{K,s}-R^{B}_{K,s}|\le C_W x_{k_0}e^{-c_W p_0(g_{k_0})}$. The triangle inequality then yields
\begin{equation*}
|R_{K,s}|\le |R^{B}_{K,s}|+|R_{K,s}-R^{B}_{K,s}|
\le |R^{B}_{K,s}|+C_W\,x_{k_0}\,e^{-c_W p_0(g_{k_0})},
\end{equation*}
and $|R^{B}_{K,s}|$ is at most the right-hand side of \eqref{8.2:eq-witness-block-a}; this is
the asserted bound on $|R_{K,s}|$.
\end{proof}

\begin{corollary}[Admissible box sizes and extraction depths]\label{8.2:cor-admissible-depths}
Let $\theta\in\left]0,\tfrac14\right]$, and let
$\sigma(\theta)=(1-\epsilon)\min\{\sigma_1,8\theta\}$ with $\epsilon\in\left]0,1\right[$ fixed;
the constants written $C_\theta$ below depend on $\epsilon$ as well as on $\theta$.
Throughout, $k_0$, $x_{k_0}=g_{k_0}^{-2}$ and $R_{K,s}$ are the level, the inverse squared running
coupling at that level and the remainder of Corollary~\ref{8.2:cor-witness-true}, $k_0$ being the
level at which that corollary reads its main term. We assume $\log x_{k_0}\ge1$, so that the degree
function $p_0(g_{k_0})=A_0(\log x_{k_0})^{p_0}$ used in the choices below is positive.
Choose the box size
\begin{equation}\label{8.2:eq-admissible-depths-1}
D:=\Bigl\lceil \tfrac{2}{c_W}\,p_0(g_{k_0})\Bigr\rceil+r\ell_0+D_0 ,
\end{equation}
and the extraction depth $r$ by one of
\begin{equation}\label{8.2:eq-admissible-depths-2}
\text{(r1)}\quad r=\Bigl\lceil \tfrac34\log_L p_0(g_{k_0})\Bigr\rceil ,
\qquad\qquad
\text{(r2)}\quad r=\bigl\lceil \theta\log_L x_{k_0}\bigr\rceil ,
\end{equation}
where $c_W$, $\ell_0$, $D_0$ are the constants of Corollary~\ref{8.2:cor-witness-true} named in
(a) below. For the depth $s$ we consider the range
\begin{equation}\label{8.2:eq-admissible-depths-3}
\bigl(\tfrac14+\theta\bigr)\log_L g^{-2}+C_{\mathrm{up}}+s_*\;\le\; s\;\le\; K-k_* ,
\end{equation}
with $C_{\mathrm{up}}$ of (b) and $s_*$, $k_*$ of (a) below.
Assume the following.
\begin{itemize}
\item[(a)] The one-point witness in the true measure, Corollary~\ref{8.2:cor-witness-true}, holds.
Its constants are $C_W$, $c_W$, $\ell_0$, $D_0$, $s_*$, $k_*$, its exponent is
$\sigma_1\in\{1,2\}$, and $R_{K,s}$ is its remainder. Its hypothesis on the box size is
$D\ge 1+r\ell_0+D_0$.
\item[(b)] The upper flow bound holds, in the following form. It fixes the depth offset
$s_{\mathrm{up}}(g)$, which enters the remainder bound \eqref{8.2:eq-witness-block-a}, and the
constant $C_{\mathrm{up}}$, which enters the range \eqref{8.2:eq-admissible-depths-3}. Moreover
there is $\gamma_W(\theta)>0$ such that, for $g\le\gamma_W(\theta)$ and $r$ chosen by (r2), every
$s$ in the range \eqref{8.2:eq-admissible-depths-3} satisfies $s_*+r\le s\le K-k_*$ and
$g_{k_0}^{\sigma(\theta)}\le C_\theta\, g_{K-s}^{\sigma(\theta)}$.
\item[(c)] The conversion of the infrared members holds in the following form. For
$g\le\gamma_W(\theta)$, $r$ chosen by (r2) and every $s$ in the range
\eqref{8.2:eq-admissible-depths-3}, each of the two infrared members of
\eqref{8.2:eq-witness-block-a}, those carrying $L^{-2s}$ and $L^{-4(s-s_{\mathrm{up}}(g))}$, is at
most $C_\theta\, g_{K-s}^{\sigma(\theta)}$.
\end{itemize}
Then:
\begin{itemize}
\item[(i)] For either choice of $r$ one has $D\ge 1+r\ell_0+D_0$. For (r1) one has
$L^{-4r}\le p_0(g_{k_0})^{-3}$; consequently the member $L^{-4r}p_0(g_{k_0})^2$ of
\eqref{8.2:eq-witness-block-a} is at most $p_0(g_{k_0})^{-1}$, so that under (r1) it decays at an
inverse-polylogarithmic rate in $x_{k_0}$.
\item[(ii)] For (r2) one has $L^{-4r}\le g_{k_0}^{8\theta}$ and
\begin{equation*}
g_{k_0}^{\sigma_1}\Pi_{r,D}\le C_\theta\, g_{k_0}^{\sigma_1}\operatorname{polylog}(x_{k_0});
\end{equation*}
the right-hand side carries no power of $x_{k_0}$, since $\Pi_{r,D}$ contains no factor $L^{4r}$.
\item[(iii)] For (r2) and every $s$ with $s_*+r\le s\le K-k_*$ one has
\begin{equation}\label{8.2:eq-admissible-depths-4}
|R_{K,s}|\le C_\theta\, g_{k_0}^{\sigma(\theta)}
+C_W\Bigl[L^{-2s}+L^{-4(s-s_{\mathrm{up}}(g))}\bigl(1+B^{\sharp}(s+1)g^2\bigr)\Bigr].
\end{equation}
If moreover $g\le\gamma_W(\theta)$ and $s$ lies in the range \eqref{8.2:eq-admissible-depths-3},
one has
\begin{equation}\label{8.2:eq-admissible-depths-5}
|R_{K,s}|\le C_\theta\, g_{K-s}^{\sigma(\theta)} .
\end{equation}
\end{itemize}
Here $\operatorname{polylog}(x_{k_0})$ denotes a quantity bounded by a constant times a fixed power
of $\log x_{k_0}$.
The conditions of this corollary are one-sided. The depth is bounded below only, by a floor of the
form
\begin{equation*}
s\ge s_W(g):=s_*+s_{\mathrm{up}}(g)+\theta\log_L g^{-2}+C
\end{equation*}
with a constant $C$. The lower end of the range \eqref{8.2:eq-admissible-depths-3} is likewise a
lower bound on $s$ only. The number $\gamma_W(\theta)$ is a ceiling on the coupling (in the sense
of Notation~\ref{2.3:not-stages-first}). Neither condition is a cap in the sense of
Definition~\ref{2.3:def-cap}. They are of the admissible shapes of
Lemma~\ref{2.3:lem-choice-shapes}.
\end{corollary}

\begin{proof}
By the standing condition $\log x_{k_0}\ge1$ we have $x_{k_0}>1$, $g_{k_0}<1$ and
$p_0(g_{k_0})>0$, and every constant is bounded by a constant times a power of $\log x_{k_0}$.

\emph{(i)} The first summand in \eqref{8.2:eq-admissible-depths-1} is the ceiling of a positive
number, hence at least $1$. This gives $D\ge 1+r\ell_0+D_0$, which is the hypothesis of
Corollary~\ref{8.2:cor-witness-true} on the box size. For (r1) we have
$r\ge\frac34\log_L p_0(g_{k_0})$, and therefore
\begin{equation*}
L^{-4r}\le L^{-3\log_L p_0(g_{k_0})}=p_0(g_{k_0})^{-3} .
\end{equation*}
Hence
\begin{equation*}
L^{-4r}p_0(g_{k_0})^2\le p_0(g_{k_0})^{-1}=A_0^{-1}(\log x_{k_0})^{-p_0} .
\end{equation*}

\emph{(ii)} For (r2) we have $r\ge\theta\log_L x_{k_0}$, and therefore
\begin{equation*}
L^{-4r}\le x_{k_0}^{-4\theta}=g_{k_0}^{8\theta},
\qquad
L^{-2r}\le g_{k_0}^{4\theta} .
\end{equation*}
Moreover $r\le\theta\log_L x_{k_0}+1$, and $p_0(g_{k_0})$ is a fixed power of $\log x_{k_0}$. By
\eqref{8.2:eq-admissible-depths-1}, both $r$ and $D$ are therefore bounded by
$C_\theta\operatorname{polylog}(x_{k_0})$.
By Lemma~\ref{8.2:lem-extraction-one-point} (the corresponding display and the
definitions following it), $\Pi_{r,D}$ is a polynomial in $D$, $r$ and $\log x_{k_0}$, and it
contains no factor $L^{4r}$. Hence the second bound of (ii) holds.
No power of $x_{k_0}$ arises, since the only $r$-dependence of $\Pi_{r,D}$ is polynomial.

\emph{(iii)} Let $r$ be chosen by (r2), and let $s_*+r\le s\le K-k_*$. By (i),
$D\ge1+r\ell_0+D_0$. So Corollary~\ref{8.2:cor-witness-true} applies. It bounds $|R_{K,s}|$ by
the right-hand side of \eqref{8.2:eq-witness-block-a} plus
$C_W x_{k_0}e^{-c_Wp_0(g_{k_0})}$. We bound the members in turn.

We use one elementary fact repeatedly. For $\delta>0$ and $x_{k_0}>1$, the quantity
$\operatorname{polylog}(x_{k_0})\,x_{k_0}^{-\delta/2}$ is bounded. Hence
$g_{k_0}^{a}\operatorname{polylog}(x_{k_0})\le C_\theta\,g_{k_0}^{\sigma(\theta)}$ whenever
$a-\sigma(\theta)>0$, with $C_\theta$ depending on $a-\sigma(\theta)$.
\begin{enumerate}
\item[(1)] By (ii),
\begin{equation*}
C_Wg_{k_0}^{\sigma_1}\Pi_{r,D}\le C_\theta g_{k_0}^{\sigma_1}\operatorname{polylog}(x_{k_0}) .
\end{equation*}
Moreover $\sigma_1-\sigma(\theta)\ge\epsilon\sigma_1>0$. So this member is at most
$C_\theta g_{k_0}^{\sigma(\theta)}$.
\item[(2)] Since $\sigma(\theta)<\sigma_1\le2$ and $g_{k_0}<1$, we have
$g_{k_0}^2\le g_{k_0}^{\sigma(\theta)}$.
\item[(3)] By (ii), and since $p_0(g_{k_0})$ and $D$ are at most
$C_\theta\operatorname{polylog}(x_{k_0})$,
\begin{equation*}
L^{-2r}g_{k_0}\,p_0(g_{k_0})\,D\operatorname{polylog}(x_{k_0})
\le C_\theta\, g_{k_0}^{1+4\theta}\operatorname{polylog}(x_{k_0}) .
\end{equation*}
We claim $1+4\theta\ge\min\{\sigma_1,8\theta\}$. If $\sigma_1=1$ this is clear. If $\sigma_1=2$,
then $8\theta\le1+4\theta$ because $\theta\le\frac14$. Hence
$1+4\theta-\sigma(\theta)\ge\epsilon\min\{\sigma_1,8\theta\}>0$, and this member is at most
$C_\theta g_{k_0}^{\sigma(\theta)}$.
\item[(4)] By (ii), $L^{-4r}p_0(g_{k_0})^2\le g_{k_0}^{8\theta}\operatorname{polylog}(x_{k_0})$. Also
$8\theta-\sigma(\theta)\ge\epsilon\min\{\sigma_1,8\theta\}>0$. So this member is at most
$C_\theta g_{k_0}^{\sigma(\theta)}$.
\item[(5)] Consider the members carrying an exponential. By \eqref{8.2:eq-admissible-depths-1},
$c_WD\ge2p_0(g_{k_0})$, so
$x_{k_0}e^{-c_WD}\le x_{k_0}e^{-2p_0(g_{k_0})}$. The members
\begin{equation*}
D^4x_{k_0}(1+r)e^{-c_Wp_0(g_{k_0})},
\qquad
x_{k_0}e^{-c_WD},
\qquad
C_Wx_{k_0}e^{-c_Wp_0(g_{k_0})}
\end{equation*}
are therefore at most
\begin{equation*}
C_\theta\operatorname{polylog}(x_{k_0})\,x_{k_0}\,e^{-cA_0(\log x_{k_0})^{p_0}},
\qquad c=\min\{c_W,2\}.
\end{equation*}
The exponent $p_0$ of the degree function is larger than one, by the one-sided condition under
which it is chosen (see Notation~\ref{0.2:not-glossary}). Hence this quantity, multiplied by
any power of $x_{k_0}$, is bounded. In particular it is at most
$C_\theta g_{k_0}^{\sigma(\theta)}$.
\item[(6)] The remaining members of \eqref{8.2:eq-witness-block-a} are the infrared members
\begin{equation*}
L^{-2s}
\quad\text{and}\quad
L^{-4(s-s_{\mathrm{up}}(g))}\bigl(1+B^{\sharp}(s+1)g^2\bigr).
\end{equation*}
They contain the depth offset $s_{\mathrm{up}}(g)$ of the upper flow bound.
\end{enumerate}
Adding items (1)--(6) gives \eqref{8.2:eq-admissible-depths-4}.

Now let moreover $g\le\gamma_W(\theta)$ and let $s$ lie in the range
\eqref{8.2:eq-admissible-depths-3}. By hypothesis (b), $s_*+r\le s\le K-k_*$, so
\eqref{8.2:eq-admissible-depths-4} holds. By hypothesis (b) again,
$g_{k_0}^{\sigma(\theta)}\le C_\theta g_{K-s}^{\sigma(\theta)}$, and by hypothesis (c) the two
infrared members on the right-hand side of \eqref{8.2:eq-admissible-depths-4} are each at most
$C_\theta g_{K-s}^{\sigma(\theta)}$. This gives \eqref{8.2:eq-admissible-depths-5}.
\end{proof}

\begin{remark}[What the one-point witness does not assert]\label{8.2:rem-not-asserted}
The one-point witness under the block law,
Corollary~\ref{8.2:cor-witness-block}, together with its bound
\eqref{8.2:eq-witness-block-a} on the remainder $R^{B}_{K,s}$, does not assert
any of the following three statements.
\begin{enumerate}
\item The removal of the infrared member of \eqref{8.2:eq-witness-block-a},
namely the term
\begin{equation*}
L^{-4(s-s_{\mathrm{up}}(g))}\bigl(1+B^{\sharp}(s+1)g^{2}\bigr),
\end{equation*}
and with it the removal of the floor $s_{\mathrm{up}}(g)$ from the range of
depths $s$. Such a removal would require a Gaussian evaluation of the torus
integral at depth $0$, in which the moduli of flat connections and the constant
modes are treated as collective coordinates. No such evaluation is carried out
in this book, and the bound \eqref{8.2:eq-witness-block-a} is used only in the
form in which it is stated, with this member present.
\item Any value of the one-loop slope. In particular neither its sign nor its
non-vanishing is asserted; the one-loop term at flat background is displayed in
Remark~\ref{8.2:rem-one-loop-flat} as a formula only.
\item That the theory fails to be a Gaussian field. No such conclusion is
drawn from Corollary~\ref{8.2:cor-witness-block}.
\end{enumerate}
\end{remark}

\begin{definition}[Step, box, window and exterior datum]\label{8.3:not-step-box-window}
The following letters are fixed for the rest of this chapter. Properties (a)--(d) at the end
are proved below.

\emph{Ambient objects.} The dimension is $d=4$ and $N\ge 2$. We write
$\mathfrak{g}:=\mathfrak{su}(N)$ and $n:=\dim\mathfrak{g}=N^2-1$. For $X,Y\in\mathfrak{g}$ we
put $\langle X,Y\rangle:=-N\operatorname{tr}(XY)$, that is, minus the matrix trace of $XY$
($\operatorname{tr}$ being the normalised trace of the glossary), and
$|X|_2^2:=\langle X,X\rangle$, the Euclidean (Hilbert--Schmidt) norm. $|X|$ denotes the
operator norm of \cite[(19)]{B-CMP98}, and $c_N\ge 1$ the norm-equivalence number of
\cite[(20)]{B-CMP98}, which depends on $N$ only; for $N=2$ the two norms are proportional. Of
$|\cdot|$ and $c_N$ only the following is used: $|\cdot|$ is an $\operatorname{Ad}$-invariant
norm on $\mathfrak{g}$, and for every $\mathfrak{b}>0$
\begin{equation}\label{8.3:eq-step-box-window-1}
\{X\in\mathfrak{g}:|X|_2<\mathfrak{b}\}\subset\{X\in\mathfrak{g}:|X|<\mathfrak{b}\}
\subset\{X\in\mathfrak{g}:|X|_2<c_N\mathfrak{b}\}.
\end{equation}
On the complexification $\mathfrak{g}_{\mathbb{C}}=\mathfrak{g}+\mathrm{i}\mathfrak{g}$ we write
$\operatorname{Re}(X+\mathrm{i}Y):=X$ and $\operatorname{Im}(X+\mathrm{i}Y):=Y$, and the same
bond by bond for fields.

\emph{Lattices.} $\mathbb{T}^{(k)}$ is the lattice of the $k$-th step
(Definition~\ref{1.5:def-block-average}), measured in level-$k$ units, so that it has unit
spacing; $b$ denotes a bond and $p$ a plaquette. For a block $c$ of $\mathbb{T}^{(k+1)}$,
$\hat{b}(c)$ is the distinguished bond of $\mathbb{T}^{(k)}$ that is eliminated by the
$\delta$-function constraint of Ba{\l}aban's step \cite[\S2]{B-CMP109}, where it is written
$b_0(c)$; we write $\hat{b}(c)$ because $b_0$ denotes the number $(N^2-1)/2$ of the glossary.
The \emph{free bonds} are the bonds of $\mathbb{T}^{(k)}$ other than the bonds $\hat{b}(c)$. The
letter $d(\cdot,\cdot)$ denotes the lattice distance between bonds of $\mathbb{T}^{(k)}$, in
level-$k$ units.

\emph{Couplings and the box half-width.} The couplings are $g_k$ and $x_k:=g_k^{-2}$
(Definition~\ref{1.2:def-couplings}). With the degree function $p_0(\cdot)$ of the glossary
(Notation~\ref{0.2:not-glossary}), $p_0(g)=A_0(\log g^{-2})^{p_0}$, whose exponent $p_0$
satisfies the one-sided condition recorded there, we put
\begin{equation}\label{8.3:eq-step-box-window-2}
\varepsilon_k:=g_k\,p_0(g_k),\qquad
\delta_k:=\frac{A_1}{A_0}\,\varepsilon_k,\qquad
\mathfrak{b}_k:=\frac{\delta_k}{g_k}.
\end{equation}
We call $\mathfrak{b}_k$ the \emph{box half-width}; it is the half-width of the fluctuation
window in the scaled variable.

\emph{Levels.} $K$ is the number of steps and $s$ the depth; $k_0:=K-s$; $r$ is the number of
extracted steps. The extracted steps form the segment $k_0\le j<k_0+r$, that is, the $r$ steps
$j=k_0,\dots,k_0+r-1$, and $i:=j-k_0$, so that $0\le i<r$. This range of $j$ and of $i$ is the
one used throughout this chapter.

\emph{The step.} Ba{\l}aban's step starts from \cite[(2.1)]{B-CMP109} and arrives at
\cite[(2.12)--(2.13)]{B-CMP109}. After the six preparatory operations by which the Gaussian
extraction inserts the observable into the step and brings the step into this form,
the $(k+1)$-st transformation restricted to the small-field region is a Gaussian integral
$d\mu_{C^{(k)}}(B)$ \cite[\S2]{B-CMP109}. Its covariance is
\begin{equation}\label{8.3:eq-step-box-window-3}
C^{(k)}(U_{k+1})=\bigl(\mathbf{C}^*\Delta^{(k)}\mathbf{C}\bigr)^{-1},
\end{equation}
where $\mathbf{C}$ is Ba{\l}aban's linear elimination operator ($B'=\mathbf{C}B$) and
$\Delta^{(k)}$ is the operator of Ba{\l}aban's quadratic form in
\cite[(2.12)--(2.13)]{B-CMP109}. The integral carries the
small-field window $\chi_k$ and the total exponent $P^{\mathrm{tot}}_k$, which consists of the
term $P^{(k)}$ of the exponent of the integrand in \cite[\S2]{B-CMP109} together with the
further terms of that exponent there. When the large-field operations are
present, $P^{\mathrm{tot}}_k$ also contains the increments of the terms of the operation
written $\mathbf{E}$ in \cite[(2.23)--(2.24)]{B-CMP119}, of the large-field operations
$\mathbf{R}$ and $\mathbf{B}$, and of the frozen-coupling terms of
\cite[(2.23)--(2.24)]{B-CMP119}. An observable $G$ of $V_k$ is carried as the factor
$G(\mathcal{V}_W(B))$, where
\begin{equation}\label{8.3:eq-step-box-window-4}
\mathcal{V}_W(B):=\exp\bigl(\mathrm{i}\bigl[g_k\mathbf{C}B-\eta\tilde{D}(g_k\mathbf{C}B)\bigr]\bigr)\,
V^{(k)}(W).
\end{equation}
Here $V^{(k)}(W)$ is the configuration of \cite[\S2]{B-CMP109} on which the exponential factor
acts, and $\eta\tilde{D}(g_k\mathbf{C}B)$ is the correction term of Ba{\l}aban's formula there,
with his operator $\tilde{D}$; the factor written $\eta$ in \eqref{8.3:eq-step-box-window-4} is
written $h$ in that formula, and we write $\eta$ so that $h$ denotes only the history that
indexes the conditional expectations below.

\emph{Conditional expectations, observables and the step operator.}
$\mathbb{E}_{i,h}[\cdot](V_{j+1})$, where $h$ is a history (Notation~\ref{0.2:not-glossary}),
is the branch conditional expectation of the segment decomposition of the extracted steps. It is
a positive normalised linear functional on bounded functions of $V_j$, and
$\operatorname{Cov}_{\mathbb{E}_{i,h}}$ is its covariance. $\mathfrak{F}_j(S;w,\varrho)$ is the
class of local, gauge-invariant, analytic observables indexed by a set $S$, the complex source
$w$ and the parameter $\varrho$; it is defined in the proof of the Gaussian extraction, where its
invariance and membership properties are also proved. By the operator theorem at free current,
in the form in which the Gaussian extraction uses it, the operator
$A^{\Gamma}_j:=\Gamma_j^{-1}=\mathbf{C}^*\Delta^{(j)}\mathbf{C}$, where $\Gamma_j=C^{(j)}$, is
real symmetric on $\mathfrak{g}^{\text{free bonds}}$, and
\begin{equation}\label{8.3:eq-step-box-window-5}
a_-\le A^{\Gamma}_j\le a_+\ \ \text{as forms},\qquad
\|A^{\Gamma}_j(b,b')\|\le a_+\,e^{-c_A d(b,b')}
\end{equation}
(norms of $n\times n$ blocks), with constants $0<a_-\le a_+$ and $c_A>0$; $a_-$ is
Ba{\l}aban's $\gamma_0$. These constants do not depend on $j$, $K$, $m$, $U$ or the volume.
The same bounds hold for the operator restricted to a component of $\Omega_{j+1}$.

The Gaussian extraction further uses a decay estimate, proved in its proof, with a decay rate
and constants $C_{\mathrm{dec}}$, $\mathfrak{v}:=n/a_-$ and $\mathfrak{v}_{\mathfrak{c}}$. It
also uses
\begin{equation}\label{8.3:eq-step-box-window-6}
\vartheta_j:=n\cdot n_D\cdot C_P\,g_j\,\mathfrak{b}_j^2\cdot C_w\,\mathfrak{b}_j,
\end{equation}
where $n_D$ is the number of bonds of the box, $C_P$ is the constant of the localised
exponent bounds below, and $C_w$ is a constant fixed in the proof of the Gaussian extraction.
The number $\vartheta_j$ is not the parameter $\vartheta$ of $\mathfrak{T}(\mathsf{d};\vartheta)$.
The letters fixed here are those used in the lemmas of the proof of the Gaussian extraction.

\emph{Letters of the step on \cite[(3.25)]{B-CMP119}.} Put $j=k$ and $t:=k+1$.
$\Omega_{k+1}$ is the new small-field domain \cite[(3.5)]{B-CMP119}. $\Lambda_{k+1}\subset
\Omega_{k+1}$ is the domain of \cite[(3.20)]{B-CMP119}, on which only the small-field
characteristic functions of \cite[(3.16)]{B-CMP119} occur in the integral. $\mathbb{B}^+$ is the
set of free bonds of $\Omega_{k+1}$, that is $\Omega^{(k)*}_{k+1}$, and $\Lambda^{(k)*}_{k+1}$
is the set of free bonds of $\Lambda_{k+1}$. The integration variable of
\cite[(3.25)]{B-CMP119} is the scaled fluctuation field
$A=(A(b))_{b\in\Lambda^{(k)*}_{k+1}}$. The field $A_k$ on
$\Omega_{k+1}\cap\Lambda^c_{k+1}$ is not an integration variable of \cite[(3.25)]{B-CMP119};
it remains an argument of the operation $\mathbf{T}^{(k+1)}$ \cite[\S3]{B-CMP119}.

Let $\mathsf{Q}:=\mathbf{C}^*\Delta^{(k)}\mathbf{C}$ on $\mathfrak{g}^{\mathbb{B}^+}$; this is
the operator $A^{\Gamma}_k$ above. For a box $I\subset\Lambda^{(k)*}_{k+1}$ of free bonds, put
$O:=\mathbb{B}^+\setminus I$. We write $\mathsf{Q}_{II},\mathsf{Q}_{IO},\dots$ for the blocks of
$\mathsf{Q}$, and set
\begin{equation}\label{8.3:eq-step-box-window-7}
\begin{gathered}
\mathsf{A}:=\mathsf{Q}_{II},\qquad \mathsf{J}:=\mathsf{Q}_{IO},\qquad
\Gamma:=\mathsf{Q}^{-1},\qquad
\mathcal{S}:=\mathsf{J}(\mathsf{Q}_{OO})^{-1}\mathsf{J}^{\mathsf{T}},\\
\mathcal{K}:=\mathsf{Q}_{II}-\mathsf{Q}_{IO}(\mathsf{Q}_{OO})^{-1}\mathsf{Q}_{OI}.
\end{gathered}
\end{equation}
$\mathsf{J}$ is the boundary (source) operator. $\mathcal{K}$ is the precision of the box
marginal, the marginal on $I$ of $\mathcal{N}(0,\Gamma)$, which is $\mathcal{N}(0,\Gamma_{II})$;
here $\mathcal{N}(0,\Sigma)$ denotes the centred Gaussian measure with covariance $\Sigma$, so
that $\mathcal{N}(0,\Gamma)$ is a measure on $\mathfrak{g}^{\mathbb{B}^+}$; see (b).
The precision $\mathcal{K}$ carries no subscript and is not the witness kernel $\mathcal{K}_s$ of
the glossary; windows are denoted by $\mathsf{K}$ and $\mathcal{W}$, never by $\mathcal{K}$.

\emph{The window.} Put
\begin{equation}\label{8.3:eq-step-box-window-8}
\mathsf{K}:=\{X\in\mathfrak{g}:|X|<\mathfrak{b}_k\},\qquad
\mathcal{W}_I:=\prod_{b\in I}\mathsf{K},\qquad
\hat{\mathfrak{b}}_k:=\sup_{X\in\mathsf{K}}|X|_2.
\end{equation}

\emph{The exterior datum.} Let $y:=(A(b))_{b\in I}$. The \emph{exterior datum} is
$\mathsf{e}\in\mathsf{E}$, where $\mathsf{E}$ is the set of all integration variables of the
term of \cite[(3.25)]{B-CMP119} under consideration other than $y$. Fix
$(h,V_{k+1},\mathsf{e})$. Then $\mathsf{e}$ determines the background
$\mathfrak{u}(\mathsf{e})=(U_{k+1},J_{k+1})$, a pair whose first component is the block
configuration $U_{k+1}$ and whose second component $J_{k+1}$ is the further background datum
on which, together with $U_{k+1}$, the operator $\mathsf{Q}$ depends; it determines the operator
$\mathsf{Q}^{\mathsf{e}}:=\mathsf{Q}[\mathfrak{u}(\mathsf{e})]$ with
$\mathsf{J}^{\mathsf{e}}:=\mathsf{Q}^{\mathsf{e}}_{IO}$, the exterior field
$A_O(\mathsf{e})$, and the map $\mathcal{V}_{\mathsf{e}}$, the counterpart at the datum
$\mathsf{e}$ of the map $\mathcal{V}_W$ of \eqref{8.3:eq-step-box-window-4}. The subscript $O$
in $A_O(\mathsf{e})$, and in $C_O$ below, is the exterior bond set $O$, not the digit zero. With
$P_{\mathsf{e}}(y):=P^{\mathrm{tot}}_k(y;\mathsf{e})$ we put
\begin{equation}\label{8.3:eq-step-box-window-9}
\begin{aligned}
H_{\mathsf{e}}(y)&:=\tfrac12\langle y,\mathsf{Q}^{\mathsf{e}}_{II}\,y\rangle
+\langle y,\mathsf{J}^{\mathsf{e}}A_O(\mathsf{e})\rangle-P_{\mathsf{e}}(y),\\
E_{\mathsf{e}}(dy)&:=Z_{\mathsf{e}}^{-1}\,\mathbf{1}_{\mathcal{W}_I}(y)\,
e^{-H_{\mathsf{e}}(y)}\,dy.
\end{aligned}
\end{equation}
Here $Z_{\mathsf{e}}$ is the normalising constant, and
$\langle\cdot,\cdot\rangle$ on $\mathfrak{g}^I$ is the sum over bonds.

\emph{The reference datum.} $\mathsf{e}_0$ is any admissible exterior datum with all
fluctuation variables on $O$ set to $0$, and $\mathcal{V}_0:=\mathcal{V}_{\mathsf{e}_0}$. Put
\begin{equation}\label{8.3:eq-step-box-window-10}
\begin{gathered}
P_0(y):=P^{\mathrm{tot}}_k(y;\mathsf{e}_0)-P^{\mathrm{tot}}_k(0;\mathsf{e}_0),\qquad
\gamma^{\mathrm{box}}:=\mathcal{N}\bigl(0,((\mathsf{Q}^{\mathsf{e}_0})^{-1})_{II}\bigr),\\
\langle\Phi\rangle^{\mathrm{box}}:=
\frac{\int_{\mathcal{W}_I}\Phi\,e^{P_0}\,d\gamma^{\mathrm{box}}}
{\int_{\mathcal{W}_I}e^{P_0}\,d\gamma^{\mathrm{box}}}.
\end{gathered}
\end{equation}
$\operatorname{Cov}^{\mathrm{box}}$ is the covariance of $\langle\cdot\rangle^{\mathrm{box}}$.

\emph{The set $S^+$ on which an observable is read.} Let the observable $G$ depend on $V_k$
only through its values on a bond set
$\mathfrak{S}_G$. The set $S^+$ of bonds on which the step reads it, fixed with the
localisation of the step, satisfies $S^+\subset\{b:d(b,\mathfrak{S}_G)\le 2L\}$. We
write $n_{S^+}:=\#S^+$ and $D_\star:=d(S^+,O)$.

\emph{The tube.} For $b'\in O$ let $C_O(b')\ge1$ be the exterior-field constant of the step at
the bond $b'$; these constants depend on $A_0/A_1$ and $N$ only. $C_O(b')$ equals $1$ if $b'$ is
a bond of $\Lambda_{k+1}$, and is of order $A_0/A_1$ if $b'$ is a bond of
$\Omega_{k+1}\setminus\Lambda_{k+1}$, by \cite[(3.14)]{B-CMP119}. We put
$C_O:=\max_{b'\in O}C_O(b')\ge1$. The
\emph{exterior range} is the set of exterior values with
$|B(b')|_2\le C_O(b')\hat{\mathfrak{b}}_k$
for $b'\in O$. Let $\|\mathbf{C}\|$ be the $\ell^\infty\to\ell^\infty$ norm of the block-local
real operator $\mathbf{C}$ with respect to $|\cdot|_2$, that is, the least number with
$|(\mathbf{C}B)(b)|_2\le\|\mathbf{C}\|\max_{b'}|B(b')|_2$ for real $B$ and all $b$. With
$\varrho_{\mathcal{P}}$ a radius of Ba{\l}aban's construction, we let $\mathbf{C}$ act
complex-linearly on $(\mathfrak{g}_{\mathbb{C}})^{\mathbb{B}^+}$ and put
\begin{equation}\label{8.3:eq-step-box-window-11}
\mathcal{T}_k:=\Bigl\{B\in(\mathfrak{g}_{\mathbb{C}})^{\mathbb{B}^+}:\ 
\begin{aligned}
&|\operatorname{Re}(\mathbf{C}B)(b)|_2<(1\vee C_O)\,c_N\|\mathbf{C}\|\,\mathfrak{b}_k+1,\\
&|\operatorname{Im}(\mathbf{C}B)(b)|_2<\varrho_{\mathcal{P}}\ \text{ for all } b
\end{aligned}\Bigr\}.
\end{equation}
The bounds on $P^{\mathrm{tot}}_k$ on which the Gaussian extraction rests, read on
$\mathcal{T}_k$, are called the \emph{localised exponent bounds}. Their constants $C_P$, $c_P$ and
$\varrho_{\mathcal{P}}$ depend on $(L,\mathfrak{p},N)$ only.

With these letters the following hold.
\begin{itemize}
\item[(a)] $\mathfrak{b}_k=A_1(\log x_k)^{p_0}$.
\item[(b)] $\mathsf{Q}_{OO}$ is invertible, $\mathcal{S}\ge0$, $\mathcal{K}=\mathsf{A}-\mathcal{S}$,
and $\mathcal{K}=((\mathsf{Q}^{-1})_{II})^{-1}=(\Gamma_{II})^{-1}$.
\item[(c)] $\mathsf{K}$ is open, bounded, convex, symmetric and $\operatorname{Ad}$-invariant,
and $\hat{\mathfrak{b}}_k\le c_N\mathfrak{b}_k$.
\item[(d)] $\mathcal{T}_k$ contains, with a real margin, the closed real window
$\prod_{b\in I}\{|y_b|\le\mathfrak{b}_k\}$ times the exterior range. Precisely: every real $B$
whose restriction to $I$ lies in this closed window and whose values on $O$ lie in the exterior
range satisfies, for all $b$,
\begin{equation}\label{8.3:eq-step-box-window-12}
|\operatorname{Re}(\mathbf{C}B)(b)|_2\le(1\vee C_O)\,c_N\|\mathbf{C}\|\,\mathfrak{b}_k,
\qquad \operatorname{Im}(\mathbf{C}B)(b)=0,
\end{equation}
and hence lies in $\mathcal{T}_k$.
\end{itemize}
\end{definition}

\begin{proof}
(a) By \eqref{8.3:eq-step-box-window-2} and the form of the degree function,
\begin{equation*}
\mathfrak{b}_k=\frac{\delta_k}{g_k}=\frac{A_1}{A_0}\,p_0(g_k)
=\frac{A_1}{A_0}\,A_0\,(\log g_k^{-2})^{p_0}=A_1(\log x_k)^{p_0},
\end{equation*}
since $x_k=g_k^{-2}$.

(b) The operator $\mathsf{Q}$ is the operator $A^{\Gamma}_k$ of \eqref{8.3:eq-step-box-window-5},
read on the free bonds of $\Omega_{k+1}$. It is therefore real symmetric with
$\mathsf{Q}\ge a_-$ as a form, where $a_->0$. Three consequences follow. Throughout,
$\langle\cdot,\cdot\rangle$ on $\mathfrak{g}^O$ and on $\mathfrak{g}^{\mathbb{B}^+}$ is the sum
over bonds, as on $\mathfrak{g}^I$.

First, $\mathsf{Q}_{OI}=\mathsf{J}^{\mathsf{T}}$, by symmetry.

Second, $\mathsf{Q}_{OO}$ is invertible. For $z\in\mathfrak{g}^O$, extend $z$ by zero to
$\mathfrak{g}^{\mathbb{B}^+}$; then $\langle z,\mathsf{Q}_{OO}z\rangle\ge a_-\langle z,z\rangle$.
So $\mathsf{Q}_{OO}$ is positive definite, and so is $(\mathsf{Q}_{OO})^{-1}$.

Third, $\mathcal{S}\ge0$. Indeed, for every $y\in\mathfrak{g}^I$,
\begin{equation*}
\langle y,\mathcal{S}y\rangle=\langle\mathsf{J}^{\mathsf{T}}y,(\mathsf{Q}_{OO})^{-1}
\mathsf{J}^{\mathsf{T}}y\rangle\ge0.
\end{equation*}

Since $\mathsf{Q}_{OI}=\mathsf{J}^{\mathsf{T}}$, the definition in
\eqref{8.3:eq-step-box-window-7} reads $\mathcal{K}=\mathsf{A}-\mathcal{S}$.

Next, $\mathcal{K}$ is invertible. Given $y\in\mathfrak{g}^I$, put
$z:=-(\mathsf{Q}_{OO})^{-1}\mathsf{J}^{\mathsf{T}}y$ and $v:=(y,z)$. Expanding the form by
blocks gives
\begin{equation*}
\begin{aligned}
\langle v,\mathsf{Q}v\rangle
&=\langle y,\mathsf{A}y\rangle+2\langle y,\mathsf{J}z\rangle+\langle z,\mathsf{Q}_{OO}z\rangle\\
&=\langle y,\mathsf{A}y\rangle-2\langle y,\mathcal{S}y\rangle+\langle y,\mathcal{S}y\rangle
=\langle y,\mathcal{K}y\rangle.
\end{aligned}
\end{equation*}
Hence
\begin{equation*}
\langle y,\mathcal{K}y\rangle=\langle v,\mathsf{Q}v\rangle\ge a_-\langle v,v\rangle
\ge a_-\langle y,y\rangle,
\end{equation*}
and $\mathcal{K}$ is positive definite.

Finally, we identify $\mathcal{K}^{-1}$ with $(\mathsf{Q}^{-1})_{II}$. Put
$\mathsf{X}:=\mathcal{K}^{-1}$ and
$\mathsf{Y}:=-(\mathsf{Q}_{OO})^{-1}\mathsf{J}^{\mathsf{T}}\mathcal{K}^{-1}$.
Then
\begin{equation*}
\begin{aligned}
\mathsf{A}\mathsf{X}+\mathsf{J}\mathsf{Y}
&=\bigl(\mathsf{A}-\mathsf{J}(\mathsf{Q}_{OO})^{-1}\mathsf{J}^{\mathsf{T}}\bigr)
\mathcal{K}^{-1}=\mathrm{id}_{\mathfrak{g}^I},\\
\mathsf{J}^{\mathsf{T}}\mathsf{X}+\mathsf{Q}_{OO}\mathsf{Y}
&=\mathsf{J}^{\mathsf{T}}\mathcal{K}^{-1}-\mathsf{J}^{\mathsf{T}}\mathcal{K}^{-1}=0.
\end{aligned}
\end{equation*}
Thus the column $(\mathsf{X},\mathsf{Y})$ is the $I$-column of $\mathsf{Q}^{-1}$, and
$(\mathsf{Q}^{-1})_{II}=\mathsf{X}=\mathcal{K}^{-1}$. As $\Gamma=\mathsf{Q}^{-1}$, this gives
$\mathcal{K}=(\Gamma_{II})^{-1}$. Since the marginal of $\mathcal{N}(0,\Gamma)$ on $I$ is
$\mathcal{N}(0,\Gamma_{II})$, $\mathcal{K}$ is its precision.

(c) $\mathsf{K}$ is the open ball of radius $\mathfrak{b}_k$ of the norm $|\cdot|$ on the
finite-dimensional space $\mathfrak{g}$.
\begin{itemize}
\item It is open, because a norm on a finite-dimensional space is continuous.
\item It is convex: by homogeneity and the triangle inequality the norm $|\cdot|$ is a convex
function on $\mathfrak{g}$, and a strict sublevel set of a convex function is convex.
\item It is symmetric, since $|-X|=|X|$.
\item It is $\operatorname{Ad}$-invariant, because $|\cdot|$ is. Indeed, the norm $|\cdot|$ is
the operator norm, and for
$u\in\mathrm{SU}(N)$ multiplication by $u$ and by $u^{-1}$ preserves the Euclidean length of
vectors of $\mathbb{C}^N$. Hence $|uXu^{-1}|=|X|$.
\end{itemize}

For the bound on $\hat{\mathfrak{b}}_k$, let $X\in\mathfrak{g}$ and $\mathfrak{b}>|X|$. The second
inclusion of \eqref{8.3:eq-step-box-window-1} gives $|X|_2<c_N\mathfrak{b}$. Letting
$\mathfrak{b}\downarrow|X|$ yields
\begin{equation}\label{8.3:eq-step-box-window-13}
|X|_2\le c_N|X|\qquad(X\in\mathfrak{g}).
\end{equation}
In particular $\mathsf{K}$ is bounded, and $\hat{\mathfrak{b}}_k\le c_N\mathfrak{b}_k$.

(d) Let $B$ be real, with $|B(b')|\le\mathfrak{b}_k$ for $b'\in I$ and
$|B(b')|_2\le C_O(b')\hat{\mathfrak{b}}_k$ for $b'\in O$. By
\eqref{8.3:eq-step-box-window-13}, $|B(b')|_2\le c_N\mathfrak{b}_k$ for $b'\in I$. By the
definition of $C_O$ and by (c),
\begin{equation*}
|B(b')|_2\le C_O(b')\hat{\mathfrak{b}}_k\le C_O\hat{\mathfrak{b}}_k\le C_Oc_N\mathfrak{b}_k
\qquad(b'\in O).
\end{equation*}
Hence
$\max_{b'}|B(b')|_2\le(1\vee C_O)c_N\mathfrak{b}_k$. Since $B$ is real and $\mathbf{C}$ is real
linear, $\mathbf{C}B$ is real. By the definition of $\|\mathbf{C}\|$, for every $b$,
\begin{equation*}
|\operatorname{Re}(\mathbf{C}B)(b)|_2=|(\mathbf{C}B)(b)|_2
\le\|\mathbf{C}\|\max_{b'}|B(b')|_2
\le(1\vee C_O)\,c_N\|\mathbf{C}\|\,\mathfrak{b}_k,
\end{equation*}
and $\operatorname{Im}(\mathbf{C}B)(b)=0$. This is \eqref{8.3:eq-step-box-window-12}. The first
quantity is at most the half-width of \eqref{8.3:eq-step-box-window-11} minus $1$, and the
second is $0<\varrho_{\mathcal{P}}$. Hence $B\in\mathcal{T}_k$, with margin at least $1$ in the
real direction and $\varrho_{\mathcal{P}}$ in the imaginary direction.
\end{proof}

\begin{definition}[The clean event with a per-bond small-piece clause]\label{8.3:def-clean-event}
Let $k_0$ and $r$ be as in Notation~\ref{8.3:not-step-box-window}, and let $D$ be the parameter of
the boxes $\mathfrak{B}^{(i)}_D$ of the clean event of the preceding section. For $i\in[0,r]$, the box
$\mathfrak{B}^{(i)}_D$, its clean collar $\tilde{\mathfrak{B}}^{(i)}_D$ and its free bonds, the
conditions (c1) and (c3), and the level offset $\mathfrak{l}:=\check R_{k_0}+\ell_0$, together with
the quantities $\check R_k$ ($k$ a level) and $\ell_0$ entering it, are those of the
clean event of the preceding section, taken here without change.

The clean collar $\tilde{\mathfrak{B}}^{(i)}_D$ is the enlargement of $\mathfrak{B}^{(i)}_D$ by
$100M\check R_{k_0+i}$ level-$(k_0+i)$ cubes, each of these cubes being a cube of side $M$ in
level-$(k_0+i)$ units.

For $i\in[0,r]$ we write $j:=k_0+i$. For a level $i'<k_0-\mathfrak{l}$ we write
$\nu:=j-i'$, the number of steps from level $i'$ to level $j$.

Let $h$ be a history. A \emph{pocket} of level $i'$ is a region whose fluctuation variables were
integrated out (pre-integrated) at level $i'$. The terms of the level-$j$ effective action of $h$
that are attached to a pocket of level $i'$, as fixed with the clean event of the preceding section,
are the \emph{pocket-attached terms of level $i'$}. Each such term is indexed by a region
$X$ with a box $\tilde X$, and we use the following quantities.
\begin{itemize}
\item[(a)] $d(X)$ is the side of $\tilde X$ in level-$i'$ units, and
$\kappa_{\mathrm{ax}}(X):=\tfrac{\pi}{3}(d-1)\,d(X)$.
\item[(b)] $M_X$ is the supremum of the modulus of the term on its domain of analyticity. For the
member of the family with frozen coupling, $M_X$ includes the lag factor $(1+B^\sharp\nu_X)$ of that
member, $\nu_X$ being its lag, in steps, as fixed for that member with the clean event of the
preceding section. This factor satisfies
\[
(1+B^\sharp\nu_X)\le x_{i'} .
\]
\item[(c)] $C$ is a constant depending only on $(L,M,d)$.
\item[(d)] $c_P$ is a constant with $c_P\le\kappa_{\mathrm{loc}}/4$, where $\kappa_{\mathrm{loc}}$ is
the decay rate of the localisation factor $e^{-\kappa_{\mathrm{loc}}d(X)}$ carried by some members of
the family.
\item[(e)] $c'$ is the constant of the lemma of the preceding section that bounds the weight of its
clean event.
\end{itemize}

Let $b$ be a free bond of $\mathfrak{B}^{(i)}_D$. Its \emph{dependence range} is the set of points
within $R_{\mathbf{C}}+1$ level-$(j+1)$ blocks of $b$, enlarged by the reach of the localisation of
the quadratic fluctuation term; here $R_{\mathbf{C}}$ is the range constant, in level-$(j+1)$ blocks,
fixed with the clean event of the preceding section. This is the dependence shell used in the bound
on the pocket-attached terms that accompanies that clean event. A box $\tilde X$ lies
\emph{under the range of $b$} if it lies within this dependence range. Since $b$ lies in
$\mathfrak{B}^{(i)}_D$, every such box lies within
$R_{\mathbf{C}}+2$ level-$(j+1)$ blocks of $\mathfrak{B}^{(i)}_D$.

The \emph{per-bond small-piece clause} $(\mathrm{c}2)^\sharp$ requires the following: for every
$i\in[0,r]$, with $j=k_0+i$, and every free bond $b$ of $\mathfrak{B}^{(i)}_D$,
\begin{equation}\label{8.3:eq-clean-event-a}
\begin{split}
\sum_{i'<k_0-\mathfrak{l}}\ \sum_{X:\ \tilde X\ \text{under the range of}\ b}
 &L^{-4\nu}(1+B^\sharp\nu)\,\kappa_{\mathrm{ax}}(X)^2\\
 &\times\bigl(d(X)L^{-\nu}+C\bigr)^4 e^{c_P(d(X)L^{-\nu}+C)}\,M_X
 \;\le\; e^{-\frac12 c'p_0(g_{k_0})} .
\end{split}
\end{equation}
Here $\nu=j-i'$, and the inner sum runs over the pocket-attached terms of level $i'$ in the level-$j$
effective action of $h$ whose box $\tilde X$ lies under the range of $b$.

The \emph{clean event} $\mathcal{C}_S(D)$ is the set of histories satisfying (c1),
$(\mathrm{c}2)^\sharp$ and (c3):
\[
\mathcal{C}_S(D):=(\mathrm{c}1)\wedge(\mathrm{c}2)^\sharp\wedge(\mathrm{c}3).
\]
Thus $(\mathrm{c}2)^\sharp$ takes the place of condition (c2) of the clean event of the preceding
section. It
differs from (c2) in three ways: it is imposed bond by bond; it concerns boxes in the dependence shell
of $R_{\mathbf{C}}+2$ blocks of $\mathfrak{B}^{(i)}_D$ rather than boxes meeting
$\mathfrak{B}^{(i)}_D$; and it carries the Schur weight $e^{c_P\,d(b,b')}$, an increasing function of
the level-$j$ distance $d(b,b')$ between bonds $b$ and $b'$.

The factor $(d(X)L^{-\nu}+C)^4e^{c_P(d(X)L^{-\nu}+C)}$ in \eqref{8.3:eq-clean-event-a} is the
per-term count of bonds together with the Schur weight. Both are measured in level-$j$ units, in the
following sense. The diameter of $\tilde X$ in level-$j$ units is $d(X)L^{-\nu}$. Moreover, every bond
$b'$ whose fluctuation variable reaches a term with box $\tilde X$ under the range of $b$ lies within
$d(X)L^{-\nu}+C$ level-$j$ units of $b$.

On $\mathcal{C}_S(D)$ the following hold within every $\mathfrak{B}^{(i)}_D$ and its clean collar:
\begin{itemize}
\item[(1)] the characteristic functions of the steps $k_0,\dots,k_0+r$ are the small-field ones;
\item[(2)] the coupling is $g_{k_0+i}$;
\item[(3)] the small-field format of the effective action applies;
\item[(4)] the boxes $\mathfrak{B}^{(i)}_D$ lie in $\Lambda_{j+1}$, the region of level $j+1$ whose
complement is the large-field region of that level.
\end{itemize}
Items (1)--(3) hold as on the clean event of the preceding section; item (4) is hypothesis (a) of
the lemma below that lists the hypotheses of the localisation of the sourced step.
\end{definition}

\begin{proof}[Proof of the geometric assertion in Definition~\ref{8.3:def-clean-event}]
Fix $i$, $j=k_0+i$, a free bond $b$ of $\mathfrak{B}^{(i)}_D$, a level $i'<k_0-\mathfrak{l}$, and a
pocket-attached term of level $i'$ whose box $\tilde X$ lies under the range of $b$.

\emph{The diameter of $\tilde X$ in level-$j$ units.} By definition, $d(X)$ is the side of $\tilde X$
measured in level-$i'$ units. One level-$j$ unit consists of $L^{j-i'}=L^{\nu}$ level-$i'$ units.
Hence the diameter of $\tilde X$, measured in level-$j$ units, is $d(X)L^{-\nu}$.

\emph{The distance of $b'$ from $b$.} The Schur weight is a function of level-$j$ distances $d(b,b')$
between bonds. Let $b'$ be a bond whose fluctuation variable reaches the term with box $\tilde X$. Then
that term depends on the variable at $b'$, so $b'$ lies within the dependence range, on its own side,
of a bond of $\tilde X$; and $\tilde X$ lies under the range of $b$. A
path from $b$ to $b'$ therefore consists of three pieces:
\begin{itemize}
\item[(i)] from $b$ into $\tilde X$, within the dependence range of $b$;
\item[(ii)] across $\tilde X$, over at most its level-$j$ diameter $d(X)L^{-\nu}$;
\item[(iii)] from $\tilde X$ to $b'$, within the dependence range on the side of $b'$.
\end{itemize}
Each dependence range has a length in level-$j$ units bounded by a constant depending only on
$(L,M,d)$. Collecting the two ranges in the constant $C$ gives
\[
d(b,b')\le d(X)L^{-\nu}+C .
\]

\emph{The count and the weight.} Every such $b'$ lies in the level-$j$ cube about $b$ of side
$2(d(X)L^{-\nu}+C)$. In dimension $d=4$ this cube contains at most
$\bigl(2(d(X)L^{-\nu}+C)+1\bigr)^4$ lattice points, and each bond in it has its initial point among
them, with $4$ bonds per initial point. Hence the number of bonds in the cube is at most
\[
4\bigl(2(d(X)L^{-\nu}+C)+1\bigr)^4 ,
\]
a numerical multiple of $(d(X)L^{-\nu}+C)^4$ when $C\ge1$; the numerical factor is not carried in
\eqref{8.3:eq-clean-event-a}, which defines the event. For each such $b'$ the Schur weight satisfies
$e^{c_P d(b,b')}\le e^{c_P(d(X)L^{-\nu}+C)}$, because $d(b,b')\le d(X)L^{-\nu}+C$. These are the two
factors displayed in \eqref{8.3:eq-clean-event-a}.
\end{proof}

The clause $(\mathrm{c}2)^\sharp$ is an event, not a deterministic bound, for the following reason.
Under a level-$j$ bond the level-$i'$ terms have a footprint of $L^{4\nu}$ positions. This footprint is
cancelled, but not beaten, by the curvature factor $L^{-4\nu}$ in \eqref{8.3:eq-clean-event-a}.

The clause enters the weight of the clean event through the following bound on the weight of its
complement:
\begin{equation*}
\operatorname{weight}\bigl(((\mathrm{c}2)^\sharp)^{\mathrm{c}}\bigr)
\le C'D^4(1+r)\,e^{-c_{\mathrm{wt}}p_0(g_{k_0})},
\end{equation*}
with constants $C'$ and $c_{\mathrm{wt}}>0$ fixed in that bound. This bound is the weight bound for
the per-bond
small-piece clause, which rests on the large-field block bound, healed part, in its pointwise form at
every finer level, and on the lower bound in the running-shift statement of the correlation theorem;
it is used here as stated.
Its proof runs the Markov inequality of the lemma of the preceding section on the weight of its clean
event bond by bond over the at most $C_2(DL+100M\check R)^4$ bonds, $C_2$ a constant,
and over the $r+1$ steps, and it is a bound of the same form as the weight bound of that
lemma. Neither \eqref{8.3:eq-clean-event-a} nor the definition of $\mathcal{C}_S(D)$ involves
this bound.

The bound on the complement holds under the one-sided floor
$x_{k_0}\ge x_{\mathrm{geo}}(L,M,\mathfrak{p})$ of \eqref{8.3:eq-depth-coupling-conditions-b}, which
makes $(100M\check R_{i'}+C)L^{-\nu}\le1$ for all $i'<k_0-\mathfrak{l}$ and all $j\in[k_0,k_0+r]$.
Together with $c_P\le\kappa_{\mathrm{loc}}/4$ of item (d), the floor enters its Markov inequality as
follows. There
the sum over $X$, at each position, of the weights displayed in \eqref{8.3:eq-clean-event-a} is
absorbed against the convergent factor $q_{i'}$ of the level-$i'$ terms, in the same way as the power
$x_{i'}^{a_P}$ of the inverse coupling is absorbed in the lemma of the preceding section on the weight
of its clean event, with the exponent
$a_P$ enlarged and with constants that depend on $M$. The members of the family fall into two kinds.

\emph{Members carrying the localisation factor.} These are the members of the two kinds of
pocket-attached terms labelled $B$ and $\mathbf{B}'$ with the clean event of the preceding section,
which carry $e^{-\kappa_{\mathrm{loc}}d(X)}$. Since $\nu\ge1$,
we have $L^{-\nu}\le1$. Together with $c_P\le\kappa_{\mathrm{loc}}/4$ this gives
\[
e^{c_Pd(X)L^{-\nu}}\le e^{c_Pd(X)}\le e^{\kappa_{\mathrm{loc}}d(X)/4} .
\]

\emph{The frozen-coupling member.} This member carries no localisation factor. By the containment of
a pre-integrated region in \cite{B-CMP122a}, a region integrated out at level $k$ is contained in a
cube of side $100MR_k$. In the form in which this containment is used with the clean event of the
preceding section, applied at level $k=i'$ to the pocket of this member, it gives
\[
d(X)\le100M\check R_{i'}+C .
\]
Under the floor $x_{k_0}\ge x_{\mathrm{geo}}(L,M,\mathfrak{p})$ therefore $d(X)L^{-\nu}\le1$, and the
weight of this member is at most $e^{c_P(1+C)}$. Its remaining factors are polynomial: for a constant
$C_1=C_1(L,M,d)$,
\[
\bigl(d(X)L^{-\nu}+C\bigr)^4\kappa_{\mathrm{ax}}(X)^2\,|X|\le C_1(M\check R_{i'})^{C_1},
\qquad
(1+B^\sharp\nu_X)\le x_{i'} .
\]

Consequently no factor exponential in $M\check R_{i'}$, such as $e^{100c_PM\check R_{i'}}$, arises
for the frozen-coupling member. No condition on $x$
beyond $x_{\mathrm{geo}}$ and the conditions of the lemma of the preceding section on the weight of its
clean event enters the
bound on the complement, and the floor $x_{\mathrm{geo}}$ enters only through this bound.

\begin{definition}[One-sided conditions on the depth coupling]\label{8.3:def-depth-coupling-conditions}
Let $k_0$ and $r$ be as in Definition~\ref{8.3:def-clean-event}, and let $j=k_0+i$, $0\le i\le r$,
be the levels of the segment considered there. We write $x_j=g_j^{-2}$ for the running inverse
couplings of Definition~\ref{1.2:def-couplings}, and $p_0(\cdot)$ for the degree function of the
coupling. The following notation is used; the symbols $c'$ and $c_P$ are those of
\eqref{8.3:eq-clean-event-a}:
\begin{itemize}
\item $\mathfrak{b}_j$ and $\hat{\mathfrak{b}}_j$ are the inner and outer per-bond radii at level $j$
of the domain of the free bond variables. This domain contains every configuration whose coordinate
on each bond has norm less than $\mathfrak{b}_j$. It is contained in the set of configurations whose
coordinate on each bond has norm less than $\hat{\mathfrak{b}}_j$;
\item $c'$ is the constant in the exponent $e^{-\frac12 c'p_0(g_{k_0})}$ on the right of
\eqref{8.3:eq-clean-event-a};
\item $C_P$ is the constant that bounds, on the event \eqref{8.3:eq-clean-event-a}, the members
attached to large-field regions together with the members produced by the large-field operation
$\mathbf{R}$. The number $c_P$ is the constant in the weight $e^{c_P(d(X)L^{-\nu}+C)}$ of
\eqref{8.3:eq-clean-event-a};
\item $\mathbf{C}$ is the operator whose operator norm $\|\mathbf{C}\|$ enters the bounds on the
members attached to large-field regions, that is, the first two inequalities of
\eqref{8.3:eq-depth-coupling-conditions-1} below;
\item $C_\delta$ and $C_{\mathrm{grad}}$ are the constants of these two bounds, and $C$ is the
constant of the third inequality of \eqref{8.3:eq-depth-coupling-conditions-1}; elsewhere in this
definition the letter $C$ denotes constants whose values may differ from one occurrence to the
next;
\item $\kappa_0\ge 4$ is the exponent in the factor $x_j^{(3-\kappa_0)/2}$ of the third
inequality of \eqref{8.3:eq-depth-coupling-conditions-1};
\item $\mathcal{P}$ is a fixed polynomial.
\end{itemize}
The one-sided conditions used in this section are the following.

\smallskip
\noindent(a) \emph{The floor $x_{\mathrm{pk}}$.} We require
\begin{equation}\label{8.3:eq-depth-coupling-conditions-a}
x_{k_0}\ \ge\ x_{\mathrm{pk}}(L,\mathfrak{p},N,M).
\end{equation}
The threshold $x_{\mathrm{pk}}$ depends on $(L,\mathfrak{p},N,M)$ only.
Here $x_{\mathrm{pk}}(L,\mathfrak{p},N,M)$ is the least number $x$ such that, whenever
$x_{k_0}\ge x$, the three inequalities \eqref{8.3:eq-depth-coupling-conditions-1} below hold for
every $j$ of the segment. If $x_{k_0}\ge x$, then $x_j\ge x/2$ for every $j$ of the segment, by
the lower flow bound for the running inverse couplings together with the range condition
$2B^\sharp(r+1)\le x_{k_0}$. The three inequalities are
\begin{equation}\label{8.3:eq-depth-coupling-conditions-1}
\begin{aligned}
C_\delta\,\|\mathbf{C}\|^2\,\mathfrak{b}_j^2\,e^{-\frac12 c'p_0(g_{k_0})}
&\le C_P\,g_j\,\mathfrak{b}_j,\\
C_{\mathrm{grad}}\,\|\mathbf{C}\|\,\hat{\mathfrak{b}}_j\,
\bigl(\mathfrak{b}_j^2+p_0(g_j)\bigr)\,e^{-\frac12 c'p_0(g_{k_0})}
&\le C_P\,g_j\,\mathfrak{b}_j^2,\\
C\,g_j\,x_j^{(3-\kappa_0)/2}\,\mathcal{P}(\log x_j)\,\hat{\mathfrak{b}}_j^{-1}
&\le C_P\,g_j\,\mathfrak{b}_j .
\end{aligned}
\end{equation}
The second inequality of \eqref{8.3:eq-depth-coupling-conditions-1} controls the per-bond
gradient of the members attached to large-field regions on the event
\eqref{8.3:eq-clean-event-a}.

\smallskip
\noindent(b) \emph{The floor $x_{\mathrm{geo}}$.} For a level $k$, let $R_k$ be the
large-field radius at level $k$ of Ba{\l}aban's construction: a large-field region created at
step $k$ is contained in a cube of side $100MR_k$ in level-$k$ units \cite{B-CMP122a}.
We write $\check R_k$ for the radius at level $k$ that enters the lag of
Definition~\ref{8.3:def-clean-event}. Let $\mathfrak{l}=\check R_{k_0}+\ell_0$ be that lag, which
appears in the summation range $i'<k_0-\mathfrak{l}$ of \eqref{8.3:eq-clean-event-a}; here
$\ell_0$ is the fixed offset of the lag in the clean event. We require
\begin{equation}\label{8.3:eq-depth-coupling-conditions-b}
x_{k_0}\ \ge\ x_{\mathrm{geo}}(L,M,\mathfrak{p}).
\end{equation}
Here $x_{\mathrm{geo}}(L,M,\mathfrak{p})$ is the least number $x$ such that, whenever
$x_{k_0}\ge x$, every $i'<k_0-\mathfrak{l}$ and every $j\in[k_0,k_0+r)$ satisfy
\eqref{8.3:eq-depth-coupling-conditions-2} below, where $C$ is the additive constant of the bound
$d(X)\le 100M\check R_{i'}+C$, in level-$i'$ units, on the diameter of the box of a member $X$ of
\eqref{8.3:eq-clean-event-a} with frozen coupling created at step $i'$:
\begin{equation}\label{8.3:eq-depth-coupling-conditions-2}
\bigl(100M\check R_{i'}+C\bigr)\,L^{-\nu}\ \le\ 1,
\qquad \nu:=j-i'>\mathfrak{l}.
\end{equation}
The floor \eqref{8.3:eq-depth-coupling-conditions-b} is used only to bound the weight
$e^{c_P(d(X)L^{-\nu}+C)}$ of the members $X$ of \eqref{8.3:eq-clean-event-a} by $e^{c_P(1+C)}$;
no display of the localisation of the sourced step or of the Gaussian extraction depends on it.

\smallskip
\noindent(c) \emph{The floor $D_0$.} Let $D$ be the free width parameter of the class
$\mathcal{C}_S(D)$ of histories considered in this section, let $\ell_S$ be the length attached
to the set $S$ of that class, and let $\ell_{\mathrm{geo}}$ be the geometric constant of the
corollary of this section in which the floor below is used. We require
\begin{equation}\label{8.3:eq-depth-coupling-conditions-c}
D\ \ge\ D_0:=\ell_{\mathrm{geo}}+1 .
\end{equation}
This makes the length $\ell_S$ at least $1$. The floor
\eqref{8.3:eq-depth-coupling-conditions-c} lies inside the lower bound on $D$ under which the
Gaussian extraction is stated.

Put $\ell_\star:=M+2L+1$. Let $O$ be the complement, in the set $\mathbb{B}^+$ of bonds of the
localisation of the sourced step, of the set of free bonds of $\mathfrak{B}^{(i)}_D$ (the free
bonds of \eqref{8.3:eq-clean-event-a}), and let $S^+$ be the set attached to $S$ in the class
$\mathcal{C}_S(D)$. The lower bound
\begin{equation*}
D_\star:=\operatorname{dist}(S^+,O)\ \ge\ DL-\ell_\star
\end{equation*}
requires nothing beyond $D\ge 1+\lceil \ell_\star/L\rceil$ to be non-vacuous, and this condition
likewise lies inside the Gaussian extraction's floor on $D$.

\smallskip
\noindent(d) \emph{Further one-sided conditions, displayed with the statements that use them.}
Three further conditions are used in this section. Each is displayed with the statement proved
under it and is only collected here.
\begin{itemize}
\item[(d1)] The first is the ceiling on the running coupling $g_j$ under which the theorem is
stated that treats, for every history of $\mathcal{C}_S(D)$, the integral over the free bonds of
$\mathfrak{B}^{(i)}_D$ against a centred Gaussian measure on those bonds. It bounds
$g_j\mathfrak{b}_j$ times a
number depending on $C_P$ and $c_P$. It is used on the event \eqref{8.3:eq-clean-event-a}
under the floor \eqref{8.3:eq-depth-coupling-conditions-a}. Since
$g_j\mathfrak{b}_j=A_1g_j(\log x_j)^{p_0}$ and the function
$g\mapsto g(\log g^{-2})^{p_0}$ is increasing for $g\le e^{-p_0}$, this ceiling is implied
by $g_j\le\gamma_H$, for a number $\gamma_H$ depending only on the dimension, on $L$, $N$,
$A_1$, $p_0$, $C_P$, $c_P$, and on the constant on the right of the ceiling. The Gaussian
extraction uses the statement proved under this ceiling only in the regime where the ceiling
holds.
\item[(d2)] The second is a ceiling of the same kind. It bounds a combination of
$g_j\mathfrak{b}_j$ and of the bounds on the members attached to large-field regions and on the
members produced by $\mathbf{R}$. The corollary referred to in part (c) is stated under it. It
holds on the whole segment once $g_{k_0}$ is below a threshold of the same kind, since
$g_j\le\sqrt2\,g_{k_0}$ and $\mathfrak{b}_j\le 2\mathfrak{b}_{k_0}$ by the two-sided flow bounds
for the running inverse couplings and the range condition.
\item[(d3)] The third is an optional floor on $D$, of the same kind as the lower bound on $D$
under which the Gaussian extraction is stated. Each of the lower bounds $D\ge C\log x_{k_0}$
imposed by the statements that use the Gaussian extraction implies it.
It is used only to fix the values $C_H=6$ and $c_H=\mu_\star/4$ of two constants of
the statement proved under (d1). Without it, the prefactor $n_{S^+}\mathfrak{b}_j^2$ of that
statement is carried along in the bounds. The optional floor involves the block parameter $M$
through $\ell_\star$, through the constant $C_\star$ of that statement, and through the number
$\mu_\star$ of part (e); through $C_\star$ it also involves the further constants of that
statement. It is a lower bound on $D$ only. It places no upper bound on $M$, which
remains a fixed parameter.
\end{itemize}
Each of (d1)--(d3) is one-sided. None of them bounds any coupling from below, and none of them
imposes a condition on the parameters of the localisation of the sourced step.

\smallskip
\noindent(e) \emph{Definitions that are not conditions.} The numbers $\mu_1$, $\mu_\star$,
$\mu_A$ and two further decay exponents of the proofs of this section are defined as the largest
numbers satisfying displayed inequalities; they impose no condition on any parameter. Auxiliary
parameters internal to those proofs are sent to their limits. The shell and the collar of width a
bounded number of blocks are part of the definitions; their cost is polynomial and is absorbed in
the clean-event lemma (distinct from
Definition~\ref{8.3:def-clean-event}, and stated under the floors on $x_{k_0}$ listed in part (f)).

\smallskip
\noindent(f) \emph{Ambient conditions.} The Gaussian extraction carries the following one-sided
conditions:
\begin{itemize}
\item the range condition $s_*\le s-r$, $2B^\sharp(r+1)\le x_{k_0}$;
\item the floors $x_{k_0}\ge x_{\mathrm{cl}}$ and $x_{k_0}\ge x_{\mathrm{life}}$;
\item the ceiling on $g_{k_0}$ of the Gaussian extraction;
\item the bounds $r\le x_{k_0}$ and $D\le x_{k_0}$. These two enter only the last clause
of the per-bond gradient estimate governed by the second inequality of
\eqref{8.3:eq-depth-coupling-conditions-1}.
\end{itemize}
In this section the defining list of the ceiling on $g_{k_0}$ of the Gaussian extraction also
contains the two ceilings (d1), (d2) and the floor \eqref{8.3:eq-depth-coupling-conditions-a}.
The floors on $x_{k_0}$ under which the clean-event lemma is used are $x_{\mathrm{cl}}$,
$x_{\mathrm{life}}$ and $x_{\mathrm{geo}}$.
\end{definition}

\begin{proof}
We show that \eqref{8.3:eq-depth-coupling-conditions-a}--\eqref{8.3:eq-depth-coupling-conditions-c}
are one-sided, and that the number $x_{\mathrm{geo}}$ exists.

\emph{One-sidedness.} Since $x_{k_0}=g_{k_0}^{-2}$, each of
\eqref{8.3:eq-depth-coupling-conditions-a} and \eqref{8.3:eq-depth-coupling-conditions-b} is a
lower bound on $x_{k_0}$. Equivalently, each is an upper bound on $g_{k_0}$. The condition
\eqref{8.3:eq-depth-coupling-conditions-a} is of the same arithmetic kind as the floor
$x_{\mathrm{life}}$ of the clean-event lemma. Under that floor, a factor $L^{4\mathfrak{l}}$ times
a polynomial factor is absorbed by $e^{-c_qp_0(g_{k_0})}$ at the cost of half of its exponent:
for $x_{k_0}\ge x_{\mathrm{life}}$ their product is at most $e^{-\frac12 c_qp_0(g_{k_0})}$.
Here $c_q$ is the constant in the exponent of that lemma's bound. The condition
\eqref{8.3:eq-depth-coupling-conditions-c} is a lower bound on $D$.

\emph{Meaning of \eqref{8.3:eq-depth-coupling-conditions-2}.} Let $X$ be a member of
\eqref{8.3:eq-clean-event-a} with frozen coupling, created at step $i'$, and let $\tilde X$ be its
box. By Ba{\l}aban's containment of such a region in a cube of side $100MR_{i'}$
\cite{B-CMP122a}, in the form in which the clean-event lemma uses it, we have
\begin{equation*}
d(X)\le 100M\check R_{i'}+C
\end{equation*}
in level-$i'$ units. One level-$i'$ unit is $L^{-\nu}$ level-$j$ units. Therefore
\eqref{8.3:eq-depth-coupling-conditions-2} says that the level-$j$ diameter of $\tilde X$ is at
most one level-$j$ unit. Consequently $d(X)L^{-\nu}\le1$, and the factor
$e^{c_P(d(X)L^{-\nu}+C)}$ of $X$ in \eqref{8.3:eq-clean-event-a} is at most $e^{c_P(1+C)}$.

\emph{Existence of $x_{\mathrm{geo}}$.} Let $i'<k_0-\mathfrak{l}$ and $j\in[k_0,k_0+r)$, and put
$\nu=j-i'>\mathfrak{l}$. By the two-sided flow bounds for the running inverse couplings,
\begin{equation*}
x_{i'}\le x_j+B^\sharp\nu
\qquad\text{and}\qquad
x_j\le x_{k_0}+c_5 ,
\end{equation*}
hence
\begin{equation*}
x_{i'}\ \le\ x_{k_0}+c_5+B^\sharp\nu .
\end{equation*}
The radii satisfy $\check R_{i'}\asymp(\log x_{i'})^{r_{\mathrm{B}}}$, where $r_{\mathrm{B}}$ is
Ba{\l}aban's exponent in the size of these radii (denoted $r$ in Ba{\l}aban's construction; written
$r_{\mathrm{B}}$ here to distinguish it from the segment length $r$). Consequently the left-hand
factor $100M\check R_{i'}+C$ of \eqref{8.3:eq-depth-coupling-conditions-2} grows at most like a
fixed power of $\log(x_{k_0}+c_5+B^\sharp\nu)$.

Against this factor stands
\begin{equation*}
L^{\nu}=L^{\mathfrak{l}}\cdot L^{\nu-\mathfrak{l}},
\qquad
\mathfrak{l}\ \ge\ \check R_{k_0}\asymp(\log x_{k_0})^{r_{\mathrm{B}}} .
\end{equation*}
The required inequality thus compares $100M$ times a fixed power of
$\log(x_{k_0}+c_5+B^\sharp\nu)$, plus a constant, with
$L^{\mathfrak{l}}L^{\nu-\mathfrak{l}}$. This is a comparison of the same arithmetic kind as the
one under which the floor $x_{\mathrm{life}}$ is met. The factor $L^{\nu-\mathfrak{l}}$ absorbs
the growth of that power in
$\nu$, uniformly in $\nu>\mathfrak{l}$. The factor $L^{\mathfrak{l}}$, with
$\mathfrak{l}\ge\check R_{k_0}$, absorbs its growth in $x_{k_0}$. Hence
\eqref{8.3:eq-depth-coupling-conditions-2} holds for every $i'<k_0-\mathfrak{l}$ and every
$j\in[k_0,k_0+r)$ once $x_{k_0}$ is large enough. The least such threshold is
$x_{\mathrm{geo}}(L,M,\mathfrak{p})$.

Larger $L$ only helps in \eqref{8.3:eq-depth-coupling-conditions-2}. The floor
$x_{\mathrm{geo}}$ is therefore a lower bound on $x_{k_0}$, equivalently an upper bound on
$g_{k_0}$, of the same kind as $x_{\mathrm{cl}}$ and $x_{\mathrm{life}}$.
\end{proof}

\begin{definition}[Ba{\l}aban's step with the conditional fluctuation integral;
{\cite[\S3, (3.1)--(3.25)]{B-CMP119}}]\label{8.3:not-fluctuation-integral}
Throughout, $k$ is the level of the step, which passes from level $k$ to level $k+1$. Lengths
are measured in level-$k$ units, so that the lattice is the block lattice $\mathbb{T}^{(k)}$ of
Definition~\ref{1.5:def-block-average}. We write $\delta_k$ for Ba{\l}aban's small-field
threshold for the fluctuation field at level $k$, and we put
\begin{equation*}
\mathfrak{b}_k:=g_k^{-1}\delta_k=A_1(\log x_k)^{p_0},\qquad x_k=g_k^{-2}.
\end{equation*}
Here $A_1$ is one of Ba{\l}aban's auxiliary parameters and $p_0$ is the fixed exponent
of this threshold.
The letter $A_k$ has two meanings in \cite[\S3]{B-CMP119}, and both are kept here. In the
exponent of the first display below, and whenever it is written with arguments (a coupling and a
configuration), $A_k$ is the effective action at level $k$. Written without arguments, from
item~(d) on, $A_k$ is the fluctuation field. The same convention applies to the letter $A$:
with two arguments it is the weighted Wilson action of Definition~\ref{1.1:def-wilson-action};
without arguments, in \eqref{8.3:eq-fluctuation-integral-a}, it is the fluctuation field on
$\Lambda_{k+1}$.

In \cite[\S3]{B-CMP119} the renormalisation transformation $T$ applied to the effective density
$\rho_k$ is computed by the following operations, in this order.
In the displays and quotations of this definition the letters $h$, $C$, $M$, $P$, $Q$, $R$, $S$,
$T$ and $\Omega$, with or without indices, denote Ba{\l}aban's objects named there, not the
glossary entries of the same letters.
\begin{enumerate}
\item[(a)] The starting point is \cite[(3.1)]{B-CMP119}:
\begin{equation*}
(T\rho_k)(V_{k+1})=\sum_{\{\Omega_j\},\{\Lambda_j\}}\int dV_k\,
\delta\bigl(\bar V_kV_{k+1}^{-1}\bigr)\,\chi_k\,\mathbf{T}_k\exp A_k .
\end{equation*}
The sum has one term for each history, that is, for each pair of sequences of domains
$\{\Omega_j\},\{\Lambda_j\}$, and $\bar V_k$ is the block average of $V_k$.
\item[(b)] Next, in \cite[(3.2)--(3.5)]{B-CMP119}, the following are introduced:
\begin{itemize}
\item decompositions of unity in $V_{k+1}$, indexed by sets $P_{k+1}$;
\item the axial gauge fixing, as in \cite[(1.5)]{B-CMP119};
\item the decomposition \cite[(3.3)]{B-CMP119}, indexed by sets $Q_{k+1}$, which restricts the
approximate fluctuation field on the bond sets $(\square'^{\sim2})^{(k)*}$.
\end{itemize}
The domain $\Omega_{k+1}$ is then defined by \cite[(3.5)]{B-CMP119}.
\item[(c)] In \cite[(3.6)--(3.9)]{B-CMP119}, with $\chi_{\mathrm{Ax}}$ the characteristic function
of the axial gauge fixing:
\begin{quotation}
Now we prove that on $\Omega^{\sim}_{k+1}$ the functions $\chi_{\mathrm{Ax}}$ and $\chi_k$ are
equal to $1$ \dots{} Thus $\chi_k(\square)=1$ for $\square\subset\Omega^{\sim}_{k+1}$ \dots{} Thus
$\chi_{\mathrm{Ax}}((\Omega^{\sim}_{k+1})^{(k)})=1$
\end{quotation}
Consequently no characteristic function of an earlier level, and no gauge-fixing characteristic
function, remains on $\Omega_{k+1}$.
\item[(d)] In \cite[(3.10)--(3.14)]{B-CMP119} the background is $V^{(k)}=M^k(U_{k+1})$, the
image of the minimal configuration $U_{k+1}$ under the map $M^k$ of
\cite[(3.10)--(3.14)]{B-CMP119}, and
\begin{quotation}
the minimal configuration $U_{k+1}$ is determined by the sequence of domains
$\{\Omega_{k+1},\Omega_k,\dots\}$, and the corresponding sequence of fields
$\{V_{k+1}\lceil\Gamma_{k+1},V_k\lceil\Gamma_k,\dots\}$
\end{quotation}
The fluctuation field is $V'_k=V_k(V^{(k)})^{-1}$ on the bond set $B(\Gamma_{k+1})$ attached to
the domain $\Gamma_{k+1}$ of \cite[(3.10)--(3.14)]{B-CMP119}, and
\begin{quotation}
We define $A_k=(1/i)\log V'_k$, hence $|A_k|<O(1)\varepsilon_k$ also
\end{quotation}
The field $A_k$ so defined is the field with respect to which the expansions and the
integrations of the following items are made.
\item[(e)] Then, following \cite[\S3]{B-CMP119}:
\begin{quotation}
The next steps are exactly the same as in Sect.~C [16], and in fact the same as in Sect.~2 [I].
We expand all the expressions with respect to $A_k$, and we linearize the expressions in the
$\delta$-functions for bonds in $\Gamma_{k+1}$. Finally we remove the $\delta$-functions using
the operator $C$, and we perform the scaling transformation $A_k=g_kA_k$
\end{quotation}
In \cite{B-CMP119} the reference [16] is \cite{B-CMP116}, whose Sect.~C is
\cite[App.~C]{B-CMP116}, and the reference [I] is \cite{B-CMP109}. After these operations the
exponent of the integrand is that of \cite[(3.15)]{B-CMP119}:
\begin{equation*}
\begin{split}
&-\tfrac12\langle A_k,C^*\Delta^{(k)}CA_k\rangle+P^{(k)}\\
&\quad+\Bigl\{(\mathbf{E}_k+\mathbf{R}_k+\mathbf{B}_k)
\bigl(U_k(\exp i[g_kCA_k-h\tilde D(g_kCA_k)]V^{(k)})\bigr)
-(\mathbf{E}_k+\mathbf{R}_k+\mathbf{B}_k)(U_{k+1})\\
&\qquad+A\bigl(\tfrac{1}{g_k^2}(\cdot)-\tfrac{1}{g_k^2},U_k(\dots)\bigr)
-A\bigl(\tfrac{1}{g_k^2}(\cdot)-\tfrac{1}{g_k^2},U_{k+1}\bigr)\Bigr\}.
\end{split}
\end{equation*}
Here $A(\cdot,\cdot)$ is the weighted Wilson action of
Definition~\ref{1.1:def-wilson-action}, and $\tfrac{1}{g_k^2}(\cdot)$ is the position-dependent
inverse coupling of Definition~\ref{1.5:def-domains}. The operator $C$ is the one with which the
$\delta$-functions are removed in the quotation above; the symbols $\Delta^{(k)}$, $P^{(k)}$,
$\mathbf{E}_k$, $\mathbf{R}_k$, $\mathbf{B}_k$, $h$, $\tilde D$ and $U_k(\cdot)$ are
Ba{\l}aban's, as in \cite[(3.15)]{B-CMP119}. The six summands of this exponent (the quadratic
form, $P^{(k)}$, and the four terms inside the brace) are the six terms with which the
Gaussian extraction deals.
\item[(f)] In \cite[(3.16)--(3.19)]{B-CMP119}:
\begin{quotation}
We want to have characteristic functions depending on the fluctuation field only, hence we
introduce another decomposition of unity \dots{}
\begin{align*}
1&=\sum_{R_{k+1}}\prod_{\square'\subset R^c_{k+1}}
\chi\Bigl(\Bigl\{\sup_{b\in(\square'^{\sim2})^{(k)*}}|A_k(b)|<\delta_k\Bigr\}\Bigr)\\
&\qquad\prod_{\square'\subset R_{k+1}}
\chi\Bigl(\Bigl\{\sup_{b\in(\square'^{\sim2})^{(k)*}}|A_k(b)|\ge\delta_k\Bigr\}\Bigr)\\
&=\sum_{R_{k+1}}\chi^{(k)}(R^c_{k+1})\,\chi^{(k)c}(R_{k+1})
\end{align*}
\end{quotation}
The text continues, in \cite[(3.17)--(3.19)]{B-CMP119}:
\begin{quotation}
For a cube $\square'\subset R^c_{k+1}$ there are now two small field characteristic functions in
the integral, one from the decomposition (3.3), and another from the decomposition (3.16). The
first is equal to $1$
\end{quotation}
The two decompositions referred to in this quotation are those of \cite[(3.3)]{B-CMP119}
(item~(b)) and \cite[(3.16)]{B-CMP119} (the display just quoted).
\cite[(3.16)]{B-CMP119} states the threshold $\delta_k$. The first-step counterpart
\cite[(1.18)--(1.19)]{B-CMP119}, written for the scaled variable, states the threshold
$g_0^{-1}\delta_0$, which is $\mathfrak{b}_0$. Whether the threshold is attached to the scaled or
to the unscaled fluctuation field affects only the radius of the set in the characteristic
function; the set of bonds over which the supremum is taken is the same in both cases.
\item[(g)] The domain $\Lambda_{k+1}$ is defined in \cite[(3.20)]{B-CMP119}, and
\begin{quotation}
On the domain $\Lambda_{k+1}$, in fact on $\Lambda^{\sim}_{k+1}$ also, there are only the small
field characteristic functions from (3.16) in the integral, their product is
$\chi^{(k)}(\Lambda_{k+1})$.
\end{quotation}
\item[(h)] Finally:
\begin{quotation}
For each term of the obtained sum we consider the conditional integral with respect to $A_k$
restricted to $\Lambda_{k+1}$. The remaining integration with respect to $A_k$ restricted to
$\Omega_{k+1}\cap\Lambda^c_{k+1}$ is included into the operation $\mathbf{T}^{(k+1)}$
\end{quotation}
See \cite[(3.23)--(3.24)]{B-CMP119}, where $\mathbf{T}_{k+1}=\mathbf{T}^{(k+1)}\mathbf{T}_k$.
\end{enumerate}

The result of these operations is \cite[(3.25)]{B-CMP119}:
\begin{equation}\label{8.3:eq-fluctuation-integral-a}
\begin{split}
(T\rho_k)(V_{k+1})
&=\sum_{\{\Omega_j\},\{V_j\}}\chi_{k+1}(\Omega_{k+1})\,\mathbf{T}_{k+1}\,
\exp\Bigl[A_k\bigl(\tfrac{1}{g_k^2},U_{k+1}\bigr)+E_0^{(k)}\Bigr]\\
&\qquad\cdot z^{(k)}\int dA\lceil_{\Lambda_{k+1}}\,\chi^{(k)}(\Lambda_{k+1})\\
&\qquad\cdot\exp\Bigl[-\tfrac12\langle A,C^*\Delta^{(k)}CA\rangle
-\langle A,C^*\Delta^{(k)}CA_k\rangle\\
&\qquad\qquad-\tfrac12\bigl\langle A_k,C^*\Delta^{(k)}CC^{(k)}(\Lambda_{k+1})
C^*\Delta^{(k)}CA_k\bigr\rangle\\
&\qquad\qquad+\mathbf{P}^{(k)}\bigl(g_k,(A_k,A)\bigr)+\{\dots\}
+V^{(k)}\bigl(S_{k+1},(A_k,A)\bigr)\Bigr].
\end{split}
\end{equation}
In this display $\{\dots\}$ stands for the brace term as printed in \cite[(3.25)]{B-CMP119}, and
\begin{quotation}
Here $A_k$ denotes the fluctuation field on $\Omega_{k+1}\cap\Lambda^c_{k+1}$, and $A$ denotes
this field on $\Lambda_{k+1}$.
\end{quotation}
The operator $C$ and the symbol $\Delta^{(k)}$ are as in item~(e); the symbols $z^{(k)}$,
$E_0^{(k)}$, $C^{(k)}(\Lambda_{k+1})$, $\mathbf{P}^{(k)}$ and $S_{k+1}$ are Ba{\l}aban's, as in
\cite[(3.25)]{B-CMP119}. In \eqref{8.3:eq-fluctuation-integral-a} the term
$V^{(k)}(S_{k+1},(A_k,A))$ is the functional so written in \cite[(3.25)]{B-CMP119}, not the
background configuration $V^{(k)}$ of item~(d).

In the pure small-field case of \cite[\S2]{B-CMP109} the same display is
\cite[(2.12)--(2.13)]{B-CMP109}. There $\Lambda_{k+1}=\Omega_{k+1}=\mathbb{T}^{(k)}$ is the whole
lattice, the field $A_k$ on $\Omega_{k+1}\cap\Lambda^c_{k+1}$ is absent, and the operation
$\mathbf{T}_{k+1}$ is absent.

This is \cite[\S3, (3.1)--(3.25)]{B-CMP119}; the derivation of each display is Ba{\l}aban's and
is not reproduced.
\end{definition}

\begin{definition}[{The fluctuation window as a product of norm balls;
\cite[(2.9)]{B-CMP109}, \cite[(2.3)]{B-CMP116}, \cite[(3.16)--(3.20)]{B-CMP119}}]
\label{8.3:def-window-product}
Let $0\le k<K$, and let $\Lambda_{k+1}$ be the region on which the fluctuation field is
integrated in Ba{\l}aban's step with the conditional fluctuation integral,
Notation~\ref{8.3:not-fluctuation-integral}, display~\eqref{8.3:eq-fluctuation-integral-a}.
Write $\Lambda^{(k)}_{k+1}$ for the set of bonds of the lattice of level $k$ contained in
$\Lambda_{k+1}$. For each element $c$ of the lattice of level $k+1$ contained in
$\Lambda_{k+1}$, \cite[(2.9)]{B-CMP109} and \cite[(1.9)]{B-CMP119} single out a bond $\hat b(c)$
of level $k$ whose variable is expressed in terms of the variables on the remaining bonds;
write $\Lambda^{(k+1)}_{k+1}$ for the set of these $c$. Following \cite[(1.9)]{B-CMP119}, the
\emph{free bonds} of $\Lambda_{k+1}$ are
\begin{equation*}
\Lambda^{(k)*}_{k+1}
=\Lambda^{(k)}_{k+1}\setminus\bigl\{\hat b(c): c\in\Lambda^{(k+1)}_{k+1}\bigr\}.
\end{equation*}
Let $\mathfrak{g}$ be the Lie algebra of $\mathrm{SU}(N)$ and $|\cdot|$ an
$\operatorname{Ad}$-invariant norm on $\mathfrak{g}$, the one in which the cited papers bound the
fields. Let $\mathfrak{b}_k$ be the box radius of the preceding section,
$\mathfrak{b}_k=A_1(\log x_k)^{p_0}$, and put
\begin{equation*}
\mathsf{K}_k=\{\xi\in\mathfrak{g}:\ |\xi|<\mathfrak{b}_k\},
\end{equation*}
with indicator function $\mathbf{1}_{\mathsf{K}_k}$. The \emph{fluctuation window} of the step
from level $k$ to level $k+1$ is the function of the integration variable
$A=(A(b))_{b\in\Lambda^{(k)*}_{k+1}}$ of \eqref{8.3:eq-fluctuation-integral-a}, taken in the
rescaled normalisation in which the unscaled fluctuation field is $g_kA$, given by
\begin{equation}\label{8.3:eq-window-product-a}
\chi^{(k)}(\Lambda_{k+1})
=\prod_{b\in\Lambda^{(k)*}_{k+1}}\mathbf{1}_{\mathsf{K}_k}\bigl(A(b)\bigr).
\end{equation}
The set $\mathsf{K}_k$ is open, bounded, convex, symmetric and $\operatorname{Ad}$-invariant.
\end{definition}

The product form \eqref{8.3:eq-window-product-a}, one norm ball for each free bond, is the
definition of \cite[(2.9)]{B-CMP109} and of \cite[(3.16), (3.17)--(3.20)]{B-CMP119}, and is the
form used in the expansion \cite[(2.3)]{B-CMP116}. In \cite{B-CMP119} the window is written for
the unscaled fluctuation field, with a radius $\delta_k$; the radius $\mathfrak{b}_k$ in
\eqref{8.3:eq-window-product-a} is the box radius of the rescaled variable, the unscaled field
being $g_k$ times the rescaled one. In \cite[(2.9)]{B-CMP109}, in whose notation $T^{(k)}$ and
$T^{(k+1)}$ are the lattices of levels $k$ and $k+1$, the characteristic function is defined in
the variables $B'(b)$ and with a radius $\varepsilon_1$, as
\begin{equation*}
\chi_k=\prod_{b\in T^{(k)}\setminus\{\hat b(c):\,c\in T^{(k+1)}\}}
\chi\bigl(\{|B'(b)|<\varepsilon_1\}\bigr),
\end{equation*}
where $\chi(\cdot)$ is the characteristic function of the set indicated. The bonds $\hat b(c)$
carry no factor: \cite[pp.~267--268]{B-CMP109} observes that the restrictions above imply
restrictions on the variables $B'(\hat b(c))$, with $\varepsilon_1$ replaced by
$O(\varepsilon_1)$, because these variables can be expressed in terms of the remaining ones; the
restrictions on the bonds $\hat b(c)$ are thus consequences of the factors in the product and not
factors of it. Moreover the change of variables by which \cite[pp.~267--268]{B-CMP109} passes from
the variables $B$ to the variables $B'$ alters only the variables on the set
$\{\hat b(c): c\in T^{(k+1)}\}$, since the function by which the two differ vanishes everywhere
outside that set; it therefore leaves the variable on every free bond unchanged. The cluster
expansion of \cite[(2.3)]{B-CMP116} uses the bondwise product: for a set $Y_0$ it takes as free
bonds inside $Y_0$ the bonds $b\in T^{(k)}$ with $b\subset Y_0$ other than the bonds $\hat b(c)$,
and writes
\begin{equation*}
\chi_k=\chi_{k,Y_0}\,\chi_{k,Y_0^c}
=\sum_P(-1)^{|P|}\chi_{k,Y_0}\,\chi^c_{k,P},
\end{equation*}
where $\chi_{k,Y_0}$ and $\chi_{k,Y_0^c}$ are the partial products over the free bonds inside
and outside $Y_0$, the sum runs over sets $P$ of bonds, $|P|$ is the number of bonds in $P$, and
$\chi^c_{k,P}$ is the factor attached to the bond set $P$ in \cite[(2.3)]{B-CMP116}.
The corresponding fact at the first step is stated in \cite{B-CMP119}: writing $\Omega_1^{(1)}$
for the set of elements $c$ of the lattice of level one in the region $\Omega_1$ of that step,
the linearizing transformation there differs from the identity only on the bonds $\hat b(c)$,
$c\in\Omega_1^{(1)}$, hence the characteristic function $\chi'_0(\Omega_1)$ defined in
\cite[(1.8)]{B-CMP119} is unchanged under it; the free bonds of a set $X$ are the set $X^*$ of
\cite[(1.9)]{B-CMP119}, obtained from the bonds of $X$ by removing the bonds $\hat b(c)$ for $c$ in
the set $X^{(1)}$ of elements of the lattice of level one in $X$. The same reasoning at level $k$
gives that the function $\chi^{(k)}(\Lambda_{k+1})$ of \cite[(3.16)--(3.20)]{B-CMP119} has the
product form \eqref{8.3:eq-window-product-a} over $\Lambda^{(k)*}_{k+1}$, with radius $\delta_k$
in the unscaled fluctuation field and radius $\mathfrak{b}_k$ in the rescaled variable $A$.

In particular, given any fluctuation field on the bonds outside $\Lambda_{k+1}$, the window on
$\Lambda_{k+1}$ is \eqref{8.3:eq-window-product-a} and does not depend on that field.

\begin{remark}[The product form and a two-sided inclusion of windows]
The product form is stronger than the two-sided inclusion \eqref{8.3:eq-window-product-1}
between norm conditions. To state it, let $\mathcal{E}$ be the set of bonds carrying the
fluctuation field of the step, inside and outside $\Lambda_{k+1}$: the free bonds
$\Lambda^{(k)*}_{k+1}$, carrying the variable $A$, together with the bonds of
$\Omega_{k+1}\cap\Lambda^c_{k+1}$ carrying the field $A_k$ of
\eqref{8.3:eq-fluctuation-integral-a}. We consider configurations $y=(y(b))_{b\in\mathcal{E}}$
of the field over $\mathcal{E}$, in the rescaled normalisation of
Definition~\ref{8.3:def-window-product}; every maximum $\max_b$ below runs over $b\in\mathcal{E}$.
Let $\Gamma$ be a linear map acting on such configurations, let $\|\Gamma\|$ be its operator norm
for $\max_b|\cdot|$, and suppose $\|\Gamma\|\ge1$. Let $\mathcal{W}$ be the set of
configurations $y$ with $|y(b)|<\mathfrak{b}_k$ for every $b\in\mathcal{E}$; its factor over
$\Lambda^{(k)*}_{k+1}$ is the product \eqref{8.3:eq-window-product-a}. Then
\begin{equation}\label{8.3:eq-window-product-1}
\bigl\{\max_b|y(b)|<\mathfrak{b}_k/\|\Gamma\|\bigr\}
\subset\mathcal{W}\subset
\bigl\{\max_b|(\Gamma y)(b)|<\|\Gamma\|\,\mathfrak{b}_k\bigr\}.
\end{equation}
The first inclusion holds because $\mathfrak{b}_k/\|\Gamma\|\le\mathfrak{b}_k$; the second
because, on $\mathcal{W}$,
\begin{equation*}
\max_b|(\Gamma y)(b)|\le\|\Gamma\|\max_b|y(b)|<\|\Gamma\|\,\mathfrak{b}_k .
\end{equation*}
A window condition, that is a set of configurations over $\mathcal{E}$, satisfying
\eqref{8.3:eq-window-product-1} alone need not have the product form, and its restriction to
$\Lambda_{k+1}$ given the exterior field --- the set of configurations over the free bonds of
$\Lambda_{k+1}$ which, together with the given values on the bonds of $\mathcal{E}$ outside
$\Lambda_{k+1}$, lie in it --- may depend on the exterior field. Fix $\eta>1$ and
take for $\Gamma$ the map multiplying the field on every bond of $\mathcal{E}$ by $\eta$, so that
$\|\Gamma\|=\eta$; write $\mathfrak{b}'=\mathfrak{b}_k/\eta$ and $\mathfrak{b}=\mathfrak{b}_k$, so
that $\mathfrak{b}'<\mathfrak{b}$. Since $|\eta\,y(b)|<\eta\,\mathfrak{b}_k$ exactly when
$|y(b)|<\mathfrak{b}$, the inclusions \eqref{8.3:eq-window-product-1} for a set $\mathcal{W}'$ in
place of $\mathcal{W}$ read
\begin{equation*}
\bigl\{\max_b|y(b)|<\mathfrak{b}'\bigr\}\subset\mathcal{W}'\subset
\bigl\{\max_b|y(b)|<\mathfrak{b}\bigr\}.
\end{equation*}
For a nearest-neighbour pair $\{b,b'\}$ of bonds of $\mathcal{E}$ let $\mathcal{W}_{bb'}$ be the
set of configurations $y$ with either $|y(b)|<\mathfrak{b}'$ and $|y(b')|<\mathfrak{b}'$, or both
$|y(b)|$ and $|y(b')|$ in $[\mathfrak{b}',\mathfrak{b})$, and let $\mathcal{W}'$ be the
intersection of the sets $\mathcal{W}_{bb'}$ over all nearest-neighbour pairs in $\mathcal{E}$. If
$|y(b)|<\mathfrak{b}'$ for every $b$, every pair satisfies the first alternative, so
$y\in\mathcal{W}'$; if $y\in\mathcal{W}'$, then, since every bond of $\mathcal{E}$ has a nearest
neighbour in $\mathcal{E}$, every bond satisfies $|y(b)|<\mathfrak{b}$. Thus $\mathcal{W}'$
satisfies \eqref{8.3:eq-window-product-1} for this $\Gamma$. Now freeze the bonds of
$\mathcal{E}$ outside $\Lambda_{k+1}$ and adjacent to it at values with
$|y|\in[\mathfrak{b}',\mathfrak{b})$. For a free bond $b$ of $\Lambda_{k+1}$ neighbouring
such an exterior bond $b'$, the condition $\mathcal{W}_{bb'}$ excludes the first alternative and
forces $|y(b)|\in[\mathfrak{b}',\mathfrak{b})$; repeating this along chains of nearest-neighbour
pairs, every free bond of $\Lambda_{k+1}$ joined to the exterior by such a chain is forced into
$[\mathfrak{b}',\mathfrak{b})$. For this exterior field the restriction of $\mathcal{W}'$ to
$\Lambda_{k+1}$ given the exterior field therefore excludes the ball of radius $\mathfrak{b}'$
about the origin, and localisation of the sourced step
fails for it. Localisation of the sourced step uses the product form
\eqref{8.3:eq-window-product-a}, whose restriction to $\Lambda_{k+1}$ does not involve the
exterior field; the inclusions \eqref{8.3:eq-window-product-1} alone do not suffice for it, as
the example shows.
\end{remark}

\begin{definition}[Box precision, boundary operator and marginal precision]\label{8.3:def-box-precision}
Let $\mathbb{B}^{+}$ (the free bonds of the new small-field domain), the box $I\subset\mathbb{B}^{+}$,
its exterior $O:=\mathbb{B}^{+}\setminus I$ and the real symmetric form $\mathsf{Q}$ on
$\mathfrak{g}^{\mathbb{B}^{+}}$ be as in Notation~\ref{8.3:not-step-box-window}, and write
$\operatorname{dist}(b,b')$ for the lattice distance between bonds $b,b'$ used there. All inner products,
adjoints and order relations between
operators refer to the Euclidean inner product $\langle\cdot,\cdot\rangle$ of
$\mathfrak{g}^{\mathbb{B}^{+}}$. For an operator $\mathsf{M}$ on $\mathfrak{g}^{\mathbb{B}^{+}}$ and
$X,Y\subset\mathbb{B}^{+}$, $\mathsf{M}_{XY}$ is its block from $\mathfrak{g}^{Y}$ to
$\mathfrak{g}^{X}$, and $\mathsf{M}(b,b')$ is its entry, a linear map of $\mathfrak{g}$, with operator
norm $\|\mathsf{M}(b,b')\|$.

From the operator theorem at free current we take the following hypothesis. At each real background
configuration of the step, $\mathsf{Q}$ is real symmetric and
\begin{equation}\label{8.3:eq-box-precision-1}
a_-\le\mathsf{Q}\le a_+\ \text{as forms},\qquad
\|\mathsf{Q}(b,b')\|\le a_+e^{-c_A \operatorname{dist}(b,b')}\quad(b,b'\in\mathbb{B}^{+}),
\end{equation}
where $0<a_-\le a_+$ and $c_A>0$ do not depend on the level, on $K$, $m$, the background
configuration or the volume. Here $a_-$ is the lower form bound of that theorem, $a_+$
is the larger of its upper form bound and the constant of its entry bound, and $c_A$ is the decay
rate of its entry bound.

We call
\begin{itemize}
\item[(a)] $\mathsf{A}:=\mathsf{Q}_{II}$ the \emph{box precision};
\item[(b)] $\mathsf{J}:=\mathsf{Q}_{IO}$ the \emph{boundary operator} (or source operator);
\item[(c)] $\mathcal{K}:=\mathsf{Q}_{II}-\mathsf{Q}_{IO}(\mathsf{Q}_{OO})^{-1}\mathsf{Q}_{OI}$ the
\emph{marginal precision};
\item[(d)] $\mathcal{S}:=\mathsf{J}(\mathsf{Q}_{OO})^{-1}\mathsf{J}^{\top}$ the \emph{boundary
correction}.
\end{itemize}
These objects are well defined, and the following hold.
\begin{itemize}
\item[(i)] For $y\in\mathfrak{g}^{I}$ and $A_O\in\mathfrak{g}^{O}$,
\begin{equation}\label{8.3:eq-box-precision-2}
\tfrac12\langle y\oplus A_O,\mathsf{Q}(y\oplus A_O)\rangle
=\tfrac12\langle y,\mathsf{A}y\rangle+\langle y,\mathsf{J}A_O\rangle
+\tfrac12\langle A_O,\mathsf{Q}_{OO}A_O\rangle .
\end{equation}
Hence $\mathsf{A}$ is the precision of the law of the box field conditioned on the exterior, and a
frozen exterior field $A_O$ acts on the box through the linear term $\langle y,\mathsf{J}A_O\rangle$.
With the integrated bond set taken to be $I$, these are the inverse covariance and the source
operator of the conditional integral in the decomposition of Ba{\l}aban's form,
\cite[(3.23), (3.25)]{B-CMP119}.
\item[(ii)] $\mathsf{J}$ is local:
$\|\mathsf{J}(b,o)\|\le a_+e^{-c_A \operatorname{dist}(b,o)}$ for $b\in I$ and $o\in O$.
\item[(iii)] $\mathcal{K}=\bigl((\mathsf{Q}^{-1})_{II}\bigr)^{-1}=(\Gamma_{II})^{-1}$, where
$\Gamma:=\mathsf{Q}^{-1}$. Thus
$\mathcal{K}$ is the precision of the marginal $\mathcal{N}(0,\Gamma_{II})$ on $\mathfrak{g}^{I}$ of
the Gaussian law $\mathcal{N}(0,\Gamma)$ on $\mathfrak{g}^{\mathbb{B}^{+}}$, and
$a_-\le\mathcal{K}\le a_+$.
\item[(iv)] $\mathsf{A}-\mathcal{K}=\mathcal{S}$ and $\mathcal{S}\ge0$.
\end{itemize}
When $\mathsf{Q}$ is evaluated at the background of the reference exterior datum of
Notation~\ref{8.3:not-step-box-window}, the marginal $\mathcal{N}(0,\Gamma_{II})$ is the box measure
$\gamma^{\mathrm{box}}$ used in the localisation of the sourced step. The boundary correction
$\mathcal{S}$ is localised at the edge of the box from both sides; this is proved in the lemma on
$\mathcal{S}$ below.
\end{definition}

\begin{proof}
Restrict the first bound in \eqref{8.3:eq-box-precision-1} to vectors supported in $O$, that is,
to $0\oplus A_O$. This gives $\langle A_O,\mathsf{Q}_{OO}A_O\rangle\ge a_-\langle A_O,A_O\rangle$.
Hence $\mathsf{Q}_{OO}$
is positive definite and invertible, and $(\mathsf{Q}_{OO})^{-1}\ge0$. In the same way
$a_-\le\mathsf{A}\le a_+$. Since $\mathsf{Q}\ge a_->0$, $\mathsf{Q}$ is invertible. So
$\Gamma$, $\mathcal{K}$ and $\mathcal{S}$ are well defined. Because $\mathsf{Q}$ is symmetric,
$\mathsf{Q}_{OI}=(\mathsf{Q}_{IO})^{\top}=\mathsf{J}^{\top}$.

(i) Expand the form according to $\mathfrak{g}^{\mathbb{B}^{+}}=\mathfrak{g}^{I}\oplus\mathfrak{g}^{O}$.
This gives $\frac12\langle y,\mathsf{Q}_{II}y\rangle$, the cross terms
$\frac12\langle y,\mathsf{Q}_{IO}A_O\rangle$ and $\frac12\langle A_O,\mathsf{Q}_{OI}y\rangle$, and
$\frac12\langle A_O,\mathsf{Q}_{OO}A_O\rangle$. By $\mathsf{Q}_{OI}=\mathsf{J}^{\top}$ the two cross
terms are equal, and their sum is $\langle y,\mathsf{J}A_O\rangle$. This is
\eqref{8.3:eq-box-precision-2}.

The Gaussian law with covariance $\Gamma$ has density proportional to
$\exp(-\frac12\langle y\oplus A_O,\mathsf{Q}(y\oplus A_O)\rangle)$. Fix $A_O$. By
\eqref{8.3:eq-box-precision-2}, the conditional density in $y$ is proportional to
$\exp(-\frac12\langle y,\mathsf{A}y\rangle-\langle y,\mathsf{J}A_O\rangle)$, because the last term of
\eqref{8.3:eq-box-precision-2} does not depend on $y$. This is a Gaussian density with precision
$\mathsf{A}$, which is positive by the first paragraph, and with linear term
$\langle y,\mathsf{J}A_O\rangle$. This is the decomposition of \cite[(3.23), (3.25)]{B-CMP119} with
the integrated set taken to be $I$.

(ii) For $b\in I$ and $o\in O$, $\mathsf{J}(b,o)=\mathsf{Q}(b,o)$ by definition of the block. The
second bound in \eqref{8.3:eq-box-precision-1} with $b'=o$ then gives the claim.

(iii) Put $\tilde A_O:=A_O+(\mathsf{Q}_{OO})^{-1}\mathsf{J}^{\top}y$. Expanding
$\langle\tilde A_O,\mathsf{Q}_{OO}\tilde A_O\rangle$ and using the symmetry of $\mathsf{Q}_{OO}$ gives
\begin{equation*}
\langle\tilde A_O,\mathsf{Q}_{OO}\tilde A_O\rangle
=\langle A_O,\mathsf{Q}_{OO}A_O\rangle+2\langle y,\mathsf{J}A_O\rangle
+\langle y,\mathsf{J}(\mathsf{Q}_{OO})^{-1}\mathsf{J}^{\top}y\rangle .
\end{equation*}
Together with \eqref{8.3:eq-box-precision-2} and $\mathsf{J}^{\top}=\mathsf{Q}_{OI}$, this yields
\begin{equation}\label{8.3:eq-box-precision-3}
\langle y\oplus A_O,\mathsf{Q}(y\oplus A_O)\rangle
=\langle\tilde A_O,\mathsf{Q}_{OO}\tilde A_O\rangle+\langle y,\mathcal{K}y\rangle .
\end{equation}
Integrate the density $\exp(-\frac12\langle y\oplus A_O,\mathsf{Q}(y\oplus A_O)\rangle)$ over
$A_O\in\mathfrak{g}^{O}$. For fixed $y$, the change of variables $A_O\mapsto\tilde A_O$ is a
translation. By
\eqref{8.3:eq-box-precision-3}, the integral therefore equals
$\exp(-\frac12\langle y,\mathcal{K}y\rangle)$ times the $y$-independent number
$\int\exp(-\frac12\langle\tilde A_O,\mathsf{Q}_{OO}\tilde A_O\rangle)\,d\tilde A_O$, which is finite
since $\mathsf{Q}_{OO}\ge a_-$. So the marginal of $\mathcal{N}(0,\Gamma)$ on $\mathfrak{g}^{I}$ has
density proportional to $\exp(-\frac12\langle y,\mathcal{K}y\rangle)$.

The marginal of a centred Gaussian law on a block of coordinates is the centred Gaussian law whose
covariance is the corresponding block of the covariance. The marginal is therefore
$\mathcal{N}(0,\Gamma_{II})$, and comparing densities gives $\mathcal{K}=(\Gamma_{II})^{-1}$. This is
the Schur complement identity.

For the bounds, the spectrum of the symmetric operator $\mathsf{Q}$ lies in $[a_-,a_+]$ by
\eqref{8.3:eq-box-precision-1}. So the spectrum of $\Gamma=\mathsf{Q}^{-1}$ lies in
$[1/a_+,1/a_-]$, that is, $1/a_+\le\Gamma\le1/a_-$. Restrict this to vectors $y\oplus0$ supported in
$I$: since $\langle y\oplus0,\Gamma(y\oplus0)\rangle=\langle y,\Gamma_{II}y\rangle$, we get
$1/a_+\le\Gamma_{II}\le1/a_-$. The symmetric operator $\Gamma_{II}$ thus has spectrum in
$[1/a_+,1/a_-]$. Its inverse $\mathcal{K}$ has spectrum in $[a_-,a_+]$, that is,
$a_-\le\mathcal{K}\le a_+$.

(iv) By the definitions of $\mathsf{A}$ and $\mathcal{K}$ and by $\mathsf{Q}_{IO}=\mathsf{J}$,
$\mathsf{Q}_{OI}=\mathsf{J}^{\top}$,
\begin{equation*}
\mathsf{A}-\mathcal{K}=\mathsf{Q}_{IO}(\mathsf{Q}_{OO})^{-1}\mathsf{Q}_{OI}
=\mathsf{J}(\mathsf{Q}_{OO})^{-1}\mathsf{J}^{\top}=\mathcal{S}.
\end{equation*}
For $y\in\mathfrak{g}^{I}$,
$\langle y,\mathcal{S}y\rangle=\langle\mathsf{J}^{\top}y,(\mathsf{Q}_{OO})^{-1}\mathsf{J}^{\top}y\rangle$.
This is $\ge0$ because $(\mathsf{Q}_{OO})^{-1}\ge0$ by the first paragraph. Hence $\mathcal{S}\ge0$.
\end{proof}

\begin{lemma}[The read set of the inserted observable]\label{8.3:lem-read-set}
Fix a level $k$ and the fluctuation integral of Ba{\l}aban's step in the form of
Notation~\ref{8.3:not-fluctuation-integral}. There $A$ is the fluctuation field on
$\Lambda_{k+1}$, $V^{(k)}$ is the background configuration and $U_{k+1}$ is the configuration
on the next block lattice. We use the following objects of \cite{B-CMP119}.
\begin{itemize}
\item[--] The distinguished bonds of \cite{B-CMP119} form a family indexed by $c$; the
\emph{distinguished bond} with index $c$ is written $b_0(c)$ there. We write it $\hat b(c)$, so
that $b_0$ keeps its meaning from the glossary.
\item[--] Attached to $c$ are two bond sets $B(c_-)$ and $B(c_+)$, which contain the contours
of Ba{\l}aban's construction, the \emph{corridor} $B(c)$, and a coefficient $\eta(c)$, written
$h(c)$ in \cite{B-CMP119} (the letter $h$ denotes histories in this book).
\item[--] A bond is \emph{free} if it is not of the form $\hat b(c)$ for any $c$.
\item[--] $\operatorname{dist}(b,S)$ denotes the lattice distance from a bond $b$ to a set
$S$ of bonds.
\end{itemize}
Let $\mathcal{V}$ denote the map $A\mapsto V_k$. Let $G$ be a function of the block field
$V_k$ that depends on $V_k$ only through the values $V_k(b)$, $b\in\mathfrak{S}_G$, for a set
of bonds $\mathfrak{S}_G$. Assume the following.
\begin{itemize}
\item[(i)] (Parametrisation of the block field.) The fluctuation field $A$ is indexed by the
free bonds. At a free bond $b$ one has
$V_k(b)=\exp\bigl(ig_kA(b)\bigr)V^{(k)}(b)$. At a distinguished bond one has
\begin{equation}\label{8.3:eq-read-set-1}
V_k(\hat b(c))=\exp\Bigl(i\bigl[g_k(CA)(\hat b(c))
-\eta(c)\,\tilde D(g_kCA)(c)\bigr]\Bigr)V^{(k)}(\hat b(c)).
\end{equation}
Here $C$ is the operator of \cite[p.~268]{B-CMP109}, which is determined by the
configuration $V^{(k)}$; in the notation of that paper the fluctuation variables are written
$B$ and $B'=CB$, so that $CA$ here corresponds to $B'$ there. The background $V^{(k)}$,
$U_{k+1}$ and the coefficients $\eta(c)$ do not depend on $A$.
\item[(ii)] (Locality of $C$.) $(CA)(\hat b(c))$ is a linear function of the remaining
variables of \cite{B-CMP119}, that is, of the values $A(b)$ at the free bonds $b$ of
$B(c_-)\cup B(c_+)$.
\item[(iii)] (Finite locality of $\tilde D$.) $\tilde D(g_kCA)(c)$ is the solution at $c$ of
the equation defining $\tilde D$ in \cite[p.~268]{B-CMP109}. It depends on $A$ only through
the values $A(b)$ at the free bonds $b$ of $B(c_-)\cup B(c_+)\cup B(c)$.
\item[(iv)] (Geometry.) Every bond of $B(c_-)\cup B(c_+)\cup B(c)$ lies at distance at
most $2L$ from $\hat b(c)$.
\end{itemize}
Hypotheses~(ii) and~(iii) are locality statements of Ba{\l}aban's papers, used here as stated
there. In \cite[p.~268]{B-CMP109}, $B'=CB$, where $C$ is the operator determined by the
configuration $V^{(k)}$; \cite{B-CMP116} states that the function $B'$ defined by
\cite[(3.2)]{B-CMP109} is a local function of $B$; and \cite{B-CMP119} states that
$(CA)(\hat b(c))$ is a linear function of the remaining variables on the bonds of
$B(c_-)\cup B(c_+)$ other than $\hat b(c)$. The data of the equation defining
$\tilde D(\cdot)(c)$ in \cite[p.~268]{B-CMP109} are the bonds entering $(\tilde QB')(c)$,
where $\tilde Q$ is the operator of that equation; these are bonds of the contours in
$B(c_-)$, $B(c_+)$ and of the corridor $B(c)$, and the corridor contains no distinguished
bond $\hat b(c')$ with $c'\neq c$.

Put
\begin{equation}\label{8.3:eq-read-set-2}
S^+:=\bigl(\mathfrak{S}_G\setminus\{\hat b(c)\}_c\bigr)\;\cup
\bigcup_{c\,:\,\hat b(c)\in\mathfrak{S}_G}
\bigl\{\text{free bonds of }B(c_-)\cup B(c_+)\cup B(c)\bigr\},
\end{equation}
and call $S^+$ the \emph{read set} of $G$.
Then $G\circ\mathcal{V}$ is a function of the values $A(b)$, $b\in S^+$, alone, and
\begin{equation}\label{8.3:eq-read-set-3}
S^+\subset\{b:\operatorname{dist}(b,\mathfrak{S}_G)\le 2L\}.
\end{equation}
\end{lemma}

\begin{proof}
Let $A$ and $A'$ be two fluctuation fields with $A(b)=A'(b)$ for every $b\in S^+$. We show
that $V_k(b)$ takes the same value at $A$ and at $A'$ for every $b\in\mathfrak{S}_G$. Since $G$
reads $V_k$ only on $\mathfrak{S}_G$, this gives $G(\mathcal{V}(A))=G(\mathcal{V}(A'))$.

Let first $b\in\mathfrak{S}_G$ be free. Then $b\in\mathfrak{S}_G\setminus\{\hat b(c)\}_c\subset
S^+$, so $A(b)=A'(b)$. By hypothesis~(i), $V_k(b)=\exp(ig_kA(b))V^{(k)}(b)$, and $V^{(k)}$ does
not depend on the fluctuation field. Hence $V_k(b)$ is the same at $A$ and at $A'$.

Let next $b=\hat b(c)\in\mathfrak{S}_G$. By~\eqref{8.3:eq-read-set-2}, every free bond of
$B(c_-)\cup B(c_+)\cup B(c)$ lies in $S^+$, so $A$ and $A'$ agree on all these bonds. By
hypothesis~(ii), $(CA)(\hat b(c))=(CA')(\hat b(c))$, since both are the same linear function of
the values at the free bonds of $B(c_-)\cup B(c_+)$. By hypothesis~(iii),
$\tilde D(g_kCA)(c)=\tilde D(g_kCA')(c)$, since both depend on the field only through its
values at the free bonds of $B(c_-)\cup B(c_+)\cup B(c)$. By hypothesis~(i), the coefficient
$\eta(c)$ and the factor $V^{(k)}(\hat b(c))$ do not depend on the fluctuation field. Hence the
right side of~\eqref{8.3:eq-read-set-1} is the same at $A$ and at $A'$.

These two cases exhaust $\mathfrak{S}_G$, which proves the first assertion.

For~\eqref{8.3:eq-read-set-3}, take $b\in S^+$. If
$b\in\mathfrak{S}_G\setminus\{\hat b(c)\}_c$, then
$\operatorname{dist}(b,\mathfrak{S}_G)=0$. Otherwise $b$ is a bond of
$B(c_-)\cup B(c_+)\cup B(c)$ for some $c$ with $\hat b(c)\in\mathfrak{S}_G$. By
hypothesis~(iv), $\operatorname{dist}(b,\hat b(c))\le 2L$, and therefore
$\operatorname{dist}(b,\mathfrak{S}_G)\le 2L$.
\end{proof}

In particular $G\circ\mathcal{V}$ has exactly finite range in $A$: no estimate that uses
Lemma~\ref{8.3:lem-read-set} needs decay of $\tilde D$ in distant variables, nor a truncation
of the variables outside a box.

Fix a level $k$, a history branch $h$ together with the corresponding term of Ba{\l}aban's conditional
fluctuation integral \eqref{8.3:eq-fluctuation-integral-a} (Notation~\ref{8.3:not-fluctuation-integral}),
a field $V_{k+1}$, and a set $I\subset\Lambda^{(k)*}_{k+1}$ of free bonds. The box variable is
$y:=A\restriction I$, where $A$ is the fluctuation field on $\Lambda_{k+1}$ of
\eqref{8.3:eq-fluctuation-integral-a}. The \emph{exterior datum} $\mathsf{e}$ is the collection of all
integration variables of that term other than $y$:
\begin{itemize}
\item the restriction $A\restriction(\Lambda^{(k)*}_{k+1}\setminus I)$;
\item the field $A_k$ on $\Omega_{k+1}\cap\Lambda_{k+1}^{c}$;
\item the restriction $V_k\restriction\Omega_{k+1}^{c}$;
\item every variable of $\mathbf{T}_k$;
\item if later operations of the step integrate part of $V_{k+1}$ far from $I$ (the operation
$\mathbf{R}'$ of \cite[(1.76), (1.100)]{B-CMP122a}), those variables as well.
\end{itemize}
Let $\mathsf{E}$ be the set of exterior data. For fixed $(h,V_{k+1},\mathsf{e})$ the following are fixed:
\begin{itemize}
\item the background $\mathfrak{u}(\mathsf{e}):=(U_{k+1},J_{k+1})$, the background pair on which the
form of Definition~\ref{8.3:def-box-precision} depends;
\item the form $\mathsf{Q}^{\mathsf{e}}:=\mathsf{Q}[\mathfrak{u}(\mathsf{e})]$ of
Definition~\ref{8.3:def-box-precision}, with boundary operator
$\mathsf{J}^{\mathsf{e}}=\mathsf{Q}^{\mathsf{e}}_{IO}$;
\item the exterior field $A_O(\mathsf{e})$ on $O:=\mathbb{B}^{+}\setminus I$, where $\mathbb{B}^{+}$ is
the set of free bonds of $\Omega_{k+1}$;
\item the map $\mathcal{V}_{\mathsf{e}}$ of Notation~\ref{8.3:not-step-box-window}.
\end{itemize}
The integrand of the term, as a function of $y$, is
$\operatorname{const}(\mathsf{e})\,\mathbf{1}_{\mathcal{W}_I}(y)\,e^{-H_{\mathsf{e}}(y)}$, where
\begin{equation}\label{8.3:eq-box-law-mixture-1}
\begin{gathered}
\mathcal{W}_I:=\prod_{b\in I}\mathsf{K},\qquad
H_{\mathsf{e}}(y):=\tfrac12\langle y,\mathsf{Q}^{\mathsf{e}}_{II}\,y\rangle
+\langle y,\mathsf{J}^{\mathsf{e}}A_O(\mathsf{e})\rangle-P_{\mathsf{e}}(y),\\
P_{\mathsf{e}}(y):=P^{\mathrm{tot}}_k(y;\mathsf{e}).
\end{gathered}
\end{equation}
Here $\mathsf{K}$ is the per-bond window of Notation~\ref{8.3:not-step-box-window}. The function
$P^{\mathrm{tot}}_k(y;\mathsf{e})$ is the sum of $\mathbf{P}^{(k)}$, the terms in braces in
\eqref{8.3:eq-fluctuation-integral-a}, and $V^{(k)}$, all evaluated at $(y;\mathsf{e})$. Let
$E_{\mathsf{e}}$ be the probability measure
$Z_{\mathsf{e}}^{-1}\mathbf{1}_{\mathcal{W}_I}(y)\,e^{-H_{\mathsf{e}}(y)}\,dy$, where $Z_{\mathsf{e}}$ is
the normalising integral.

\begin{lemma}[The branch expectation as a mixture of box laws]\label{8.3:lem-box-law-mixture}
With the notation just fixed, assume the following.
\begin{enumerate}
\item[(a)] \emph{Form of the step.} On the branch $h$, the step near $I$ is a term of
\eqref{8.3:eq-fluctuation-integral-a}, with $I\subset\Lambda^{(k)*}_{k+1}$.
\item[(b)] \emph{Positivity of the operations.} The large-field operations of the step are of the four
kinds $(\alpha)$--$(\delta)$ admitted in the segmentation lemma. In particular, every weight in
$\mathbf{T}_{k+1}$ is non-negative.
\item[(c)] \emph{Insertion of the observable.} In the passage from Ba{\l}aban's integral at level $k$ to
the term of \eqref{8.3:eq-fluctuation-integral-a}, the Faddeev--Popov insertion leaves a gauge-invariant
function of the field unchanged, and every other stage is a change of the integration variable or a
rewriting of factors that do not involve the inserted function.
\end{enumerate}
Let $G$ be a bounded Borel function of $V_k\restriction\mathfrak{S}_G$, the restriction of $V_k$ to a
set of bonds $\mathfrak{S}_G$, and let $S^{+}$ be the read set built from $\mathfrak{S}_G$ in
Lemma~\ref{8.3:lem-read-set}. Assume that $S^{+}\subset I$ and that $G$ is gauge invariant; in the pure
setting, in which the expectation is that of the step law of Notation~\ref{8.3:not-step-box-window}
with its field $B$, $G$ may instead be any bounded Borel function. Then, for
$\tau_h$-almost every $V_{k+1}$,
\begin{equation}\label{8.3:eq-box-law-mixture-2}
\mathbb{E}_{i,h}[G](V_{k+1})
=\int_{\mathsf{E}}\pi_{h,V_{k+1}}(d\mathsf{e})\;E_{\mathsf{e}}\bigl[G\circ\mathcal{V}_{\mathsf{e}}\bigr].
\end{equation}
Here $\pi_{h,V_{k+1}}$ is a probability measure on the set $\mathsf{E}$ of exterior data, and
$\mathbb{E}_{i,h}[G]$ is the branch expectation of the segmentation lemma at the extracted step of the
branch $h$ indexed by $i$, whose level is the level $k$ fixed above.
\end{lemma}

\begin{proof}
Let $\tau^{G}_h$ be the measure obtained from $\tau_h$ by inserting the factor $G$ into the branch's
term. Parts (b) and (c) of the segmentation lemma identify the branch expectation as a Radon--Nikodym
derivative:
\begin{equation*}
\mathbb{E}_{i,h}[G](V_{k+1})=\frac{d\tau^{G}_h}{d\tau_h}(V_{k+1}).
\end{equation*}
This is the ratio, formed at $\tau_h$-almost every $V_{k+1}$, of the branch's term computed with the
factor $G$ to the same term computed without it.

We first bring the numerator into the form of \eqref{8.3:eq-fluctuation-integral-a}, using
hypothesis~(c). The Faddeev--Popov insertion leaves $G$ unchanged, because $G$ is gauge invariant. Each
of the remaining stages is a change of the integration variable or a rewriting of factors that do not
involve $G$. Such a stage carries $G$ along as a function of the new integration variables. At fixed
exterior datum, $V_k$ is then expressed through $y$ by the map $\mathcal{V}_{\mathsf{e}}$. The
numerator is therefore the term of \eqref{8.3:eq-fluctuation-integral-a} with the factor
$G(\mathcal{V}_{\mathsf{e}}(y))$, $y=A\restriction I$, placed under the last integral. The denominator
is the same term without this factor.

Next we check that Fubini's theorem applies to both terms. All integrands in question fall into one of
two classes. Either they are non-negative, or they are bounded by $\sup|G|$ times a non-negative
integrable integrand; the latter holds because $|G\circ\mathcal{V}_{\mathsf{e}}|\le\sup|G|$. Hence
Fubini's theorem applies to the numerator and to the denominator, and we may integrate over $y$ first
at fixed $\mathsf{e}$.

We now carry out the $y$-integration. By the form of the integrand stated before
\eqref{8.3:eq-box-law-mixture-1}, the $y$-integral at fixed $\mathsf{e}$ of the numerator's integrand
is
\begin{equation*}
\operatorname{const}(\mathsf{e})\int\mathbf{1}_{\mathcal{W}_I}(y)\,e^{-H_{\mathsf{e}}(y)}\,
G(\mathcal{V}_{\mathsf{e}}(y))\,dy
=\operatorname{const}(\mathsf{e})\,Z_{\mathsf{e}}\,E_{\mathsf{e}}\bigl[G\circ\mathcal{V}_{\mathsf{e}}\bigr].
\end{equation*}
The same computation without $G$ gives $\operatorname{const}(\mathsf{e})\,Z_{\mathsf{e}}$. Consequently
\begin{equation}\label{8.3:eq-box-law-mixture-3}
\text{numerator}
=\int_{\mathsf{E}}m(d\mathsf{e})\,Z_{\mathsf{e}}\,E_{\mathsf{e}}\bigl[G\circ\mathcal{V}_{\mathsf{e}}\bigr],
\qquad
\text{denominator}=\int_{\mathsf{E}}m(d\mathsf{e})\,Z_{\mathsf{e}}.
\end{equation}
Here $m(d\mathsf{e})$ is the product of three factors:
\begin{itemize}
\item the weights of $\mathbf{T}_{k+1}$;
\item the measure $dA\restriction(\Lambda_{k+1}\setminus I)\,\chi^{(k)}(\Lambda_{k+1}\setminus I)$, in
which $\chi^{(k)}(\Lambda_{k+1}\setminus I)$ is the characteristic function $\chi^{(k)}$ of
\eqref{8.3:eq-fluctuation-integral-a} taken on $\Lambda_{k+1}\setminus I$ in place of $\Lambda_{k+1}$;
\item the number $\operatorname{const}(\mathsf{e})>0$.
\end{itemize}

We show that $m(d\mathsf{e})\ge0$. By hypothesis~(b), the operations of the step are of the kinds
$(\alpha)$--$(\delta)$. Hence every weight of $\mathbf{T}_{k+1}$ is a characteristic function, a
$\delta$-function, a factor of a partition of unity, or the exponential of a real function, and each of
these is non-negative. The second factor of $m$ is non-negative, being Lebesgue measure times a
characteristic function. The third factor is positive. Hence $m(d\mathsf{e})\ge0$.

Define
\begin{equation*}
\pi_{h,V_{k+1}}(d\mathsf{e})
:=\frac{m(d\mathsf{e})\,Z_{\mathsf{e}}}{\int_{\mathsf{E}}m(d\mathsf{e}')\,Z_{\mathsf{e}'}}.
\end{equation*}
This is a non-negative measure of total mass one on $\mathsf{E}$, that is, a probability measure.
Dividing the first identity of \eqref{8.3:eq-box-law-mixture-3} by the second gives
\eqref{8.3:eq-box-law-mixture-2}.

It remains to treat later operations of the same step. Suppose that positive operations of the same
step act on variables outside $I$, namely shares in $V_{k+1}$ and the $\mathbf{R}'$-integrations far
from $I$. These operations replace $\pi_{h,V_{k+1}}$ by a further mixture and replace $\mathsf{E}$ by a
larger set. Nothing else in the argument changes.

In the pure setting, $G$ may be any bounded Borel function and need not be gauge invariant. There the
exterior datum is $\mathsf{e}=B\restriction O$, the restriction to $O$ of the field $B$ of the step
law, and $\pi_{h,V_{k+1}}$ is the $O$-marginal of the step law. The identity
\eqref{8.3:eq-box-law-mixture-2} then uses nothing beyond hypothesis~(c).
\end{proof}

\begin{definition}[Block matrices, exponential weights and the row-sum test]
\label{8.3:not-block-matrices}
Let $\mathfrak{B}$ be a finite set carrying a metric $d$, and let $n\ge 1$ be an integer. We put
\begin{equation*}
E_{\mathfrak{B}}:=(\mathbb{R}^n)^{\mathfrak{B}},
\end{equation*}
whose elements are written $y=(y_b)_{b\in\mathfrak{B}}$ with $y_b\in\mathbb{R}^n$. We equip
$E_{\mathfrak{B}}$ with the Euclidean inner product $\langle\cdot,\cdot\rangle$ and norm
$|y|^2:=\sum_{b\in\mathfrak{B}}|y_b|^2$, where $|y_b|$ is the Euclidean norm on $\mathbb{R}^n$.
A matrix $M$ on $E_{\mathfrak{B}}$ is given by its blocks $M(b,b')\in\mathbb{R}^{n\times n}$,
$b,b'\in\mathfrak{B}$, so that $(My)_b=\sum_{b'}M(b,b')y_{b'}$. We call the $n$ coordinates of
$y_b$ the colour block at $b$. We write $\|\cdot\|$ for operator norms: for $M$ with respect to
$|\cdot|$ on $E_{\mathfrak{B}}$, and for a block $M(b,b')$ with respect to the Euclidean norm on
$\mathbb{R}^n$.

For $c>0$ we put
\begin{equation*}
S(c):=\sup_{b\in\mathfrak{B}}\sum_{b'\in\mathfrak{B}}e^{-c\,d(b,b')},\qquad
S^1(c):=\sup_{b\in\mathfrak{B}}\sum_{b'\in\mathfrak{B}}d(b,b')\,e^{-c\,d(b,b')}.
\end{equation*}
In the application, $\mathfrak{B}$ is a subset of the bonds of a lattice of unit spacing in
dimension four, and $S(c)$, $S^1(c)$ are bounded by the corresponding numbers $S_4(c)$,
$S^1_4(c)$ of that lattice, which do not depend on the volume.

For a function $\varphi\colon\mathfrak{B}\to\mathbb{R}$ we put
\begin{equation*}
\mathcal{T}_\varphi:=\bigoplus_{b\in\mathfrak{B}}e^{\varphi(b)}\,\mathbf{1}_n,
\end{equation*}
where $\mathbf{1}_n$ is the identity matrix of size $n$; thus $\mathcal{T}_\varphi$ is
block-diagonal, with blocks indexed by $\mathfrak{B}$, and acts on the colour block at $b$ as
multiplication by the scalar $e^{\varphi(b)}$. For $\mu\ge0$ we write $\operatorname{Lip}(\varphi)\le\mu$
to mean that $|\varphi(b)-\varphi(b')|\le\mu\,d(b,b')$ for all $b,b'\in\mathfrak{B}$.

Finally, every matrix $M$ on $E_{\mathfrak{B}}$ satisfies the row-sum test (Schur test)
\begin{equation}\label{8.3:eq-block-matrices-1}
\|M\|\le\max\Bigl\{\sup_{b}\sum_{b'}\|M(b,b')\|,\ \sup_{b'}\sum_{b}\|M(b,b')\|\Bigr\}.
\end{equation}
\end{definition}

\begin{proof}[Proof of \eqref{8.3:eq-block-matrices-1}]
Write $R:=\sup_b\sum_{b'}\|M(b,b')\|$ and $C:=\sup_{b'}\sum_b\|M(b,b')\|$, and let
$y\in E_{\mathfrak{B}}$. For each $b$, the definition of the block norms gives
$|(My)_b|\le\sum_{b'}\|M(b,b')\|\,|y_{b'}|$. Writing
$\|M(b,b')\|\,|y_{b'}|=\|M(b,b')\|^{1/2}\cdot\|M(b,b')\|^{1/2}|y_{b'}|$ and applying the
Cauchy--Schwarz inequality to the sum over $b'$, we obtain
\begin{equation*}
|(My)_b|^2\le\Bigl(\sum_{b'}\|M(b,b')\|\Bigr)\Bigl(\sum_{b'}\|M(b,b')\|\,|y_{b'}|^2\Bigr)
\le R\sum_{b'}\|M(b,b')\|\,|y_{b'}|^2 .
\end{equation*}
Summing over $b$ and exchanging the two finite sums,
\begin{equation*}
|My|^2\le R\sum_{b'}|y_{b'}|^2\sum_{b}\|M(b,b')\|\le R\,C\,|y|^2 .
\end{equation*}
Hence $\|M\|\le\sqrt{RC}$, and $\sqrt{RC}\le\max\{R,C\}$, which is
\eqref{8.3:eq-block-matrices-1}.
\end{proof}

\begin{lemma}[Exponential weights, decay of inverses and Schur complements]
\label{8.3:lem-weighted-inverse-decay}
Let $\mathfrak{B}$, its metric $d$, $n$, $E_{\mathfrak{B}}$, the blocks $M(b,b')$ of a matrix $M$ on
$E_{\mathfrak{B}}$, the operator norms $\|\cdot\|$, the numbers $S(\cdot)$ and $S^{1}(\cdot)$, the
weights $\mathcal{T}_\varphi$ and the condition $\operatorname{Lip}(\varphi)\le\mu$ be as in
Notation~\ref{8.3:not-block-matrices}. For $J\subset\mathfrak{B}$ we put
$E_J:=(\mathbb{R}^n)^J\subset E_{\mathfrak{B}}$. For $J,J'\subset\mathfrak{B}$ and a matrix $A$ on
$E_{\mathfrak{B}}$ we write $A_{JJ'}$ for the matrix from $E_{J'}$ to $E_J$ whose blocks are $A(b,b')$,
$b\in J$, $b'\in J'$. Let $A$ be real symmetric on $E_{\mathfrak{B}}$ with $A\ge a_->0$ in the sense
of quadratic forms and
\begin{equation*}
\|A(b,b')\|\le a_+e^{-c\,d(b,b')}\qquad (b,b'\in\mathfrak{B}),
\end{equation*}
and let $0<\mu<c$ satisfy
\begin{equation}\label{8.3:eq-weighted-inverse-decay-a}
a_+\,\mu\,S^{1}(c-\mu)\le \tfrac14\,a_- .
\end{equation}
Then for every $\varphi\colon\mathfrak{B}\to\mathbb{R}$ with $\operatorname{Lip}(\varphi)\le\mu$ the
following hold.
\begin{enumerate}
\item[(i)] $\|\mathcal{T}_\varphi A\mathcal{T}_\varphi^{-1}-A\|\le \tfrac14 a_-$; hence
$\langle\xi,\mathcal{T}_\varphi A\mathcal{T}_\varphi^{-1}\xi\rangle\ge\tfrac34 a_-|\xi|^2$ for all real
$\xi\in E_{\mathfrak{B}}$.
\item[(ii)] For all $b,b'\in\mathfrak{B}$,
\begin{equation*}
\|A^{-1}(b,b')\|\le \frac{4}{3a_-}\,e^{-\mu d(b,b')} .
\end{equation*}
\item[(iii)] For every non-empty $J\subset\mathfrak{B}$, (i) and (ii) hold verbatim for the principal
submatrix $A_{JJ}$ on $E_J$, with $d$ and $\varphi$ restricted to $J$ and with the same constants
$a_-$, $\mu$.
\item[(iv)] If a matrix $M$ on $E_{\mathfrak{B}}$ satisfies $\|M(b,b')\|\le\varepsilon e^{-c\,d(b,b')}$
for all $b,b'$, then $\|\mathcal{T}_\varphi M\mathcal{T}_\varphi^{-1}\|\le\varepsilon\,S(c-\mu)$.
\item[(v)] Let $\mathfrak{B}=I\sqcup O$ with $I$ and $O$ non-empty (if $O=\emptyset$, then
$\mathcal{S}=0$ and $\mathcal{K}=A$, and the assertions reduce to $a_-\le A\le\|A\|$), and put
\begin{equation*}
\mathcal{S}:=A_{IO}(A_{OO})^{-1}A_{OI},\qquad \mathcal{K}:=A_{II}-\mathcal{S}.
\end{equation*}
Then $\mathcal{S}\ge0$, $\mathcal{K}=((A^{-1})_{II})^{-1}$, and $a_-\le\mathcal{K}\le\|A\|$ in the
sense of quadratic forms on $E_I$. Moreover, with $C_{\mathcal{S}}:=4a_+^2/(3a_-)$ and
$d(b,O):=\min_{o\in O}d(b,o)$, for all $b,b'\in I$,
\begin{align}
\|\mathcal{S}(b,b')\|&\le C_{\mathcal{S}}\,S(c-\mu)^2\,e^{-\mu d(b,b')},
\label{8.3:eq-weighted-inverse-decay-1}\\
\sum_{b'\in I}\|\mathcal{S}(b,b')\|&\le C_{\mathcal{S}}\,S(\mu)\,S(c)\,S(c/2)\,
e^{-(c/2)d(b,O)}.\label{8.3:eq-weighted-inverse-decay-2}
\end{align}
\end{enumerate}
\end{lemma}

\begin{proof}
The argument for (i) and (ii) is the weighted-norm argument of Combes and Thomas \cite{CT73}.
Throughout, ``the Schur test'' means the row- and column-sum bound for $\|M\|$ of
Notation~\ref{8.3:not-block-matrices}.

(i) Since $\mathcal{T}_\varphi$ is block diagonal and equal to the scalar $e^{\varphi(b)}$ on the block
of $b$, the $(b,b')$-block of $\mathcal{T}_\varphi A\mathcal{T}_\varphi^{-1}-A$ is
$A(b,b')\bigl(e^{\varphi(b)-\varphi(b')}-1\bigr)$. For real $x$ one has
$e^x-1=\int_0^1xe^{tx}\,dt$, hence $|e^x-1|\le|x|e^{|x|}$. With $x=\varphi(b)-\varphi(b')$ we have
$|x|\le\mu d(b,b')$, and $t\mapsto te^t$ is increasing on $[0,\infty)$; therefore
\begin{equation*}
\bigl\|A(b,b')\bigl(e^{\varphi(b)-\varphi(b')}-1\bigr)\bigr\|
\le a_+e^{-c\,d(b,b')}\,\mu d(b,b')\,e^{\mu d(b,b')}
= a_+\mu\, d(b,b')\,e^{-(c-\mu)d(b,b')} .
\end{equation*}
For fixed $b$ the sum of these bounds over $b'$ is at most $a_+\mu S^{1}(c-\mu)$ by the definition of
$S^{1}$; since $d$ is symmetric, the same bound holds for the sum over $b$ at fixed $b'$. By the Schur
test and \eqref{8.3:eq-weighted-inverse-decay-a},
$\|\mathcal{T}_\varphi A\mathcal{T}_\varphi^{-1}-A\|\le a_+\mu S^{1}(c-\mu)\le\tfrac14a_-$. Consequently,
for real $\xi$,
\begin{equation*}
\langle\xi,\mathcal{T}_\varphi A\mathcal{T}_\varphi^{-1}\xi\rangle
\ge\langle\xi,A\xi\rangle-\tfrac14a_-|\xi|^2\ge\tfrac34a_-|\xi|^2 .
\end{equation*}

(ii) Put $N:=\mathcal{T}_\varphi A\mathcal{T}_\varphi^{-1}$. By the Cauchy--Schwarz inequality and (i),
$|N\xi|\,|\xi|\ge\langle\xi,N\xi\rangle\ge\tfrac34a_-|\xi|^2$, so $|N\xi|\ge\tfrac34a_-|\xi|$. Hence
$N$ is injective, thus invertible on the finite-dimensional space $E_{\mathfrak{B}}$, and
$\|N^{-1}\|\le 4/(3a_-)$. Since $A\ge a_->0$, $A$ is invertible and
$N^{-1}=\mathcal{T}_\varphi A^{-1}\mathcal{T}_\varphi^{-1}$, whose $(b,b')$-block is
$e^{\varphi(b)-\varphi(b')}A^{-1}(b,b')$. A block of a matrix is the matrix composed with the isometric
embedding of $\mathbb{R}^n$ into the slot $b'$ and the coordinate projection onto the slot $b$, so
its norm is at most the norm of the matrix; thus
$e^{\varphi(b)-\varphi(b')}\|A^{-1}(b,b')\|\le 4/(3a_-)$ for every admissible $\varphi$. Fix $b'$ and
take $\varphi:=\mu d(\cdot,b')$; by the triangle inequality
$|\varphi(b_1)-\varphi(b_2)|\le\mu d(b_1,b_2)$, so $\operatorname{Lip}(\varphi)\le\mu$, and
$\varphi(b)-\varphi(b')=\mu d(b,b')$. This gives
$\|A^{-1}(b,b')\|\le\frac{4}{3a_-}e^{-\mu d(b,b')}$.

(iii) The matrix $A_{JJ}$ is real symmetric. For $\xi\in E_J$ let $\tilde\xi\in E_{\mathfrak{B}}$ be its
extension by zero; then
\begin{equation*}
\langle\xi,A_{JJ}\xi\rangle=\langle\tilde\xi,A\tilde\xi\rangle\ge a_-|\tilde\xi|^2=a_-|\xi|^2,
\end{equation*}
so $A_{JJ}\ge a_-$ by restriction of the form. The blocks of $A_{JJ}$ are the blocks
$A(b,b')$, $b,b'\in J$, so they obey the same bound $a_+e^{-c\,d(b,b')}$. The numbers $S(\cdot)$ and
$S^{1}(\cdot)$ formed on $J$ with the restricted metric are sums of non-negative terms over fewer
indices, hence at most those formed on $\mathfrak{B}$; therefore
\eqref{8.3:eq-weighted-inverse-decay-a} holds on $J$, and the row and column sums in the proof of (i)
are bounded by the same quantity. The restriction of $\varphi$ to $J$ satisfies
$\operatorname{Lip}\le\mu$, and $\mu d(\cdot,b')$ restricted to $J$ is the weight used in (ii). The
proofs of (i) and (ii) therefore apply to $A_{JJ}$ word for word.

(iv) The $(b,b')$-block of $\mathcal{T}_\varphi M\mathcal{T}_\varphi^{-1}$ is
$e^{\varphi(b)-\varphi(b')}M(b,b')$, of norm at most
$\varepsilon e^{-c\,d(b,b')}e^{\mu d(b,b')}=\varepsilon e^{-(c-\mu)d(b,b')}$. Its row sums and, $d$
being symmetric, its column sums are at most $\varepsilon S(c-\mu)$; the Schur test gives the claim.

(v) By (iii), $A_{OO}\ge a_->0$, so $A_{OO}$ is invertible. Since $A$ is symmetric,
$A_{OI}=A_{IO}^{\mathsf{T}}$, and multiplying out the block matrices shows that $A$ factorises as
\begin{equation*}
A=\begin{pmatrix}1&A_{IO}A_{OO}^{-1}\\0&1\end{pmatrix}
\begin{pmatrix}\mathcal{K}&0\\0&A_{OO}\end{pmatrix}
\begin{pmatrix}1&0\\A_{OO}^{-1}A_{OI}&1\end{pmatrix}
\end{equation*}
with respect to $E_{\mathfrak{B}}=E_I\oplus E_O$, where $1$ denotes the identity on $E_I$, resp.\
$E_O$. The outer factors are invertible (unipotent), and
$A$ is invertible, so $\mathcal{K}$ is invertible; inverting the three factors and reading off the
$II$-block gives the block-inversion formula
\begin{equation*}
(A^{-1})_{II}=(A_{II}-A_{IO}A_{OO}^{-1}A_{OI})^{-1}=\mathcal{K}^{-1},
\end{equation*}
that is, $\mathcal{K}=((A^{-1})_{II})^{-1}$. The spectrum of the symmetric matrix $A$ lies in
$[a_-,\|A\|]$, so
$1/\|A\|\le A^{-1}\le 1/a_-$; the compression $(A^{-1})_{II}$ of $A^{-1}$ to $E_I$ obeys the same
form bounds, $1/\|A\|\le(A^{-1})_{II}\le1/a_-$. Thus the eigenvalues of the symmetric matrix
$(A^{-1})_{II}$ lie in $[1/\|A\|,1/a_-]$, those of its inverse $\mathcal{K}$ in $[a_-,\|A\|]$, and
$a_-\le\mathcal{K}\le\|A\|$. Further $\mathcal{S}=A_{IO}A_{OO}^{-1}A_{IO}^{\mathsf{T}}$ and
$A_{OO}^{-1}>0$, so
$\langle\xi,\mathcal{S}\xi\rangle=\langle A_{IO}^{\mathsf{T}}\xi,A_{OO}^{-1}A_{IO}^{\mathsf{T}}\xi\rangle
\ge0$.

For $b,b'\in I$,
\begin{equation*}
\mathcal{S}(b,b')=\sum_{o,o'\in O}A(b,o)\,A_{OO}^{-1}(o,o')\,A(o',b').
\end{equation*}
By (ii) and (iii) applied to
$A_{OO}$, $\|A_{OO}^{-1}(o,o')\|\le\frac{4}{3a_-}e^{-\mu d(o,o')}$; together with the bound on the
blocks of $A$ this gives
\begin{equation}\label{8.3:eq-weighted-inverse-decay-3}
\|\mathcal{S}(b,b')\|\le C_{\mathcal{S}}\sum_{o,o'\in O}
e^{-c\,d(b,o)-\mu d(o,o')-c\,d(o',b')} .
\end{equation}
By the triangle inequality $d(b,b')\le d(b,o)+d(o,o')+d(o',b')$, so
\begin{align*}
-c\,d(b,o)-\mu d(o,o')-c\,d(o',b')
&=-\mu\bigl(d(b,o)+d(o,o')+d(o',b')\bigr)-(c-\mu)\bigl(d(b,o)+d(o',b')\bigr)\\
&\le-\mu d(b,b')-(c-\mu)\bigl(d(b,o)+d(o',b')\bigr).
\end{align*}
Inserting this into \eqref{8.3:eq-weighted-inverse-decay-3} and using
$\sum_{o\in O}e^{-(c-\mu)d(b,o)}\le S(c-\mu)$ and, by symmetry of $d$,
$\sum_{o'\in O}e^{-(c-\mu)d(o',b')}\le S(c-\mu)$, we obtain \eqref{8.3:eq-weighted-inverse-decay-1}.

For \eqref{8.3:eq-weighted-inverse-decay-2}, sum \eqref{8.3:eq-weighted-inverse-decay-3} over
$b'\in I$, then over $o'$, then over $o$. First,
$\sum_{b'\in I}e^{-c\,d(o',b')}\le S(c)$; next, $\sum_{o'\in O}e^{-\mu d(o,o')}\le S(\mu)$; finally,
since $d(b,o)\ge d(b,O)$ for $o\in O$,
\begin{equation*}
\sum_{o\in O}e^{-c\,d(b,o)}
=\sum_{o\in O}e^{-(c/2)d(b,o)}e^{-(c/2)d(b,o)}
\le e^{-(c/2)d(b,O)}\,S(c/2).
\end{equation*}
Multiplying the three bounds gives
\begin{equation*}
\sum_{b'\in I}\|\mathcal{S}(b,b')\|\le C_{\mathcal{S}}\,S(c)\,S(\mu)\,S(c/2)\,e^{-(c/2)d(b,O)},
\end{equation*}
which is \eqref{8.3:eq-weighted-inverse-decay-2}.
\end{proof}

\begin{lemma}[Smooth retraction and convex wall for a convex body]\label{8.3:lem-retraction-wall}
Let $n\ge1$, let $K\subset\mathbb{R}^n$ be open, bounded and convex, let $m>0$, and put
$K':=K+B(0,m)$, where $B(0,m)$ is the open ball of radius $m$ about the origin and the sum is the
Minkowski sum; in this lemma and its proof all letters are local and unrelated to the quantities
that carry the same names elsewhere in the book. Write $\mathbf{1}$ for the identity matrix and order symmetric matrices as quadratic
forms. There are maps $\mathsf{r}\in C^\infty(\mathbb{R}^n;\mathbb{R}^n)$,
$\psi\in C^\infty(\mathbb{R}^n;[0,\infty))$ and numbers $C_\psi>0$, $c_\psi>0$ and $R_\psi$ such
that
\begin{equation}\label{8.3:eq-retraction-wall-a}
\begin{aligned}
&\text{(e1)}\quad \mathsf{r}(y)=y\ \text{for } y\in\overline{K};\qquad
  \mathsf{r}(\mathbb{R}^n)\subset\overline{K'};\\
&\text{(e2)}\quad D\mathsf{r}(y)\ \text{is symmetric and}\ 0\le D\mathsf{r}(y)\le\mathbf{1},\
  \text{so}\ |\mathsf{r}(y)-\mathsf{r}(y')|\le|y-y'|;\\
&\text{(e3)}\quad \psi\ \text{is convex},\ 0\le D^2\psi(y)\le C_\psi\mathbf{1},\
  \psi=0\ \text{on}\ \overline{K},\ \psi>0\ \text{off}\ \overline{K},\\
&\phantom{\text{(e3)}\quad}\ \psi(y)\ge c_\psi\,(|y|-R_\psi)_+^2;\\
&\text{(e4)}\quad |D^2\mathsf{r}(y)[\xi,\xi]|\le D^2\psi(y)[\xi,\xi]\quad
  \text{for all } y,\xi\in\mathbb{R}^n,
\end{aligned}
\end{equation}
where (e2) and (e3) hold for all $y,y'\in\mathbb{R}^n$.
\end{lemma}

\begin{proof}
The set $K'$ is open, bounded and convex, being the Minkowski sum of two open bounded convex
sets; hence $\overline{K'}$ is a non-empty, closed, bounded, convex set. Let $\pi$ be the metric
projection onto $\overline{K'}$, that is, $\pi(y)$ is the unique point of $\overline{K'}$ nearest
to $y$, and put
\begin{equation*}
f(y):=\tfrac12\operatorname{dist}\bigl(y,\overline{K'}\bigr)^2=\tfrac12\,|y-\pi(y)|^2,
\qquad y\in\mathbb{R}^n.
\end{equation*}
We record four properties of $f$.

(f1) \emph{$f$ is convex.} The function $\operatorname{dist}(\cdot,\overline{K'})$ is non-negative
and convex, because $\overline{K'}$ is convex; and $f$ is its composition with
$t\mapsto\frac12t^2$, which is convex and non-decreasing on $[0,\infty)$. Hence $f$ is convex.

(f2) \emph{For all $y,\xi\in\mathbb{R}^n$,}
\begin{equation*}
f(y+\xi)+f(y-\xi)-2f(y)\le|\xi|^2 .
\end{equation*}
Indeed, put $a:=y-\pi(y)$. Since $\pi(y)\in\overline{K'}$, we have
$\operatorname{dist}(y\pm\xi,\overline{K'})\le|y\pm\xi-\pi(y)|=|a\pm\xi|$, and therefore
\begin{equation*}
f(y+\xi)+f(y-\xi)\le\tfrac12|a+\xi|^2+\tfrac12|a-\xi|^2=|a|^2+|\xi|^2=2f(y)+|\xi|^2 .
\end{equation*}

(f3) \emph{$f$ is differentiable, and}
\begin{equation}\label{8.3:eq-retraction-wall-b}
\nabla f=\operatorname{id}-\pi .
\end{equation}
Let $y,\xi\in\mathbb{R}^n$ and $a:=y-\pi(y)$. The comparison used in (f2) gives
\begin{equation*}
f(y+\xi)\le\tfrac12|a+\xi|^2=f(y)+\langle y-\pi(y),\xi\rangle+\tfrac12|\xi|^2 .
\end{equation*}
Exchanging the roles of $y$ and $y+\xi$, that is, comparing $\operatorname{dist}(y,\overline{K'})$
with $|y-\pi(y+\xi)|=|(y+\xi-\pi(y+\xi))-\xi|$, we obtain
\begin{equation*}
f(y+\xi)\ge f(y)+\langle y+\xi-\pi(y+\xi),\xi\rangle-\tfrac12|\xi|^2 .
\end{equation*}
Projections onto closed convex sets are non-expansive, so $|\pi(y+\xi)-\pi(y)|\le|\xi|$, and hence
\begin{equation*}
\langle y+\xi-\pi(y+\xi),\xi\rangle
=\langle y-\pi(y),\xi\rangle+|\xi|^2-\langle\pi(y+\xi)-\pi(y),\xi\rangle
\ge\langle y-\pi(y),\xi\rangle .
\end{equation*}
The upper and the lower bound for $f(y+\xi)$ are therefore both
$f(y)+\langle y-\pi(y),\xi\rangle+O(|\xi|^2)$; more
precisely, $|f(y+\xi)-f(y)-\langle y-\pi(y),\xi\rangle|\le\frac12|\xi|^2$. Thus $f$ is
differentiable at $y$ with gradient $y-\pi(y)$, which is \eqref{8.3:eq-retraction-wall-b}. Since
$\pi$ is non-expansive, $\nabla f$ is continuous, so $f\in C^1(\mathbb{R}^n)$.

(f4) \emph{If $\overline{K'}\subset\overline{B}(0,R')$, then $f(y)\ge\frac12(|y|-R')_+^2$.} For
every $x\in\overline{K'}$ we have $|y-x|\ge|y|-|x|\ge|y|-R'$, so
$\operatorname{dist}(y,\overline{K'})\ge(|y|-R')_+$.

\emph{Second derivatives of mollifications of $f$.} For $\eta\in\mathbb{R}^n$ write
\begin{equation*}
\Delta^2_\eta f(x):=f(x+\eta)+f(x-\eta)-2f(x).
\end{equation*}
Let $0\le\phi\in C_c^\infty(\mathbb{R}^n)$. Since $f$ is continuous and $\phi$ is smooth with
compact support, $f*\phi\in C^\infty$ (all derivatives fall on $\phi$). For a $C^2$ function $g$,
Taylor's formula gives $t^{-2}\bigl(g(y+t\xi)+g(y-t\xi)-2g(y)\bigr)\to D^2g(y)[\xi,\xi]$ as
$t\to0$; applied to $g=f*\phi$, and using $(f*\phi)(y\pm t\xi)=\int f(y\pm t\xi-z)\phi(z)\,dz$,
this yields
\begin{equation}\label{8.3:eq-retraction-wall-c}
D^2(f*\phi)(y)[\xi,\xi]
=\lim_{t\to0}t^{-2}\int\Delta^2_{t\xi}f(y-z)\,\phi(z)\,dz
\ \in\ \Bigl[0,\ |\xi|^2\!\int\phi\Bigr].
\end{equation}
The inclusion holds by (f1) and (f2): by convexity of $f$ one has $\Delta^2_{t\xi}f\ge0$, and by
(f2) one has $\Delta^2_{t\xi}f\le t^2|\xi|^2$; so, since $\phi\ge0$, each quotient
$t^{-2}\int\Delta^2_{t\xi}f(y-z)\phi(z)\,dz$ lies in the closed interval
$[0,|\xi|^2\int\phi]$, and so does the limit. The derivation of the limit formula in
\eqref{8.3:eq-retraction-wall-c} did not use the sign of $\phi$; the same limit formula therefore
holds for a signed kernel $\phi\in C_c^\infty(\mathbb{R}^n)$.

\emph{The kernels.} Choose $\omega\in C_c^\infty(B(0,m/2))$ with $\omega\ge0$, $\omega$ even and
$\int\omega=1$. Let $\beta(z):=\exp\bigl(-1/(m^2-|z|^2)\bigr)$ for $|z|<m$ and $\beta(z):=0$
otherwise, and put
\begin{equation*}
\tilde\omega:=\frac{\sup|\nabla\omega|}{\inf_{|z|\le m/2}\beta(z)}\;\beta .
\end{equation*}
The infimum in the denominator equals $\exp(-4/(3m^2))>0$, and $\sup|\nabla\omega|>0$ because
$\omega$ is a non-zero function of compact support. Hence $\tilde\omega\in C_c^\infty$,
$\{\tilde\omega>0\}=B(0,m)$, and $\tilde\omega\ge|\nabla\omega|$ pointwise: on $|z|\le m/2$ the
prefactor times $\beta(z)$ is at least $\sup|\nabla\omega|$, and for $|z|>m/2$ we have
$\nabla\omega(z)=0\le\tilde\omega(z)$. Put
\begin{equation*}
\sigma:=f*\omega,\qquad\psi:=f*\tilde\omega,\qquad\mathsf{r}:=\operatorname{id}-\nabla\sigma .
\end{equation*}
Then $\sigma,\psi\in C^\infty$, so $\mathsf{r}\in C^\infty(\mathbb{R}^n;\mathbb{R}^n)$; and
$\psi\ge0$ because $f\ge0$ and $\tilde\omega\ge0$.

\emph{Proof of (e2).} We have $D\mathsf{r}=\mathbf{1}-D^2\sigma$, which is symmetric since
$D^2\sigma$ is the Hessian of a smooth function. By \eqref{8.3:eq-retraction-wall-c} with
$\phi=\omega$ and $\int\omega=1$, $D^2\sigma(y)[\xi,\xi]\in[0,|\xi|^2]$, that is,
$0\le D^2\sigma(y)\le\mathbf{1}$, and so $0\le D\mathsf{r}(y)\le\mathbf{1}$. A symmetric matrix
with spectrum in $[0,1]$ has operator norm at most $1$; writing
\begin{equation*}
\mathsf{r}(y)-\mathsf{r}(y')=\int_0^1D\mathsf{r}(y'+\tau(y-y'))(y-y')\,d\tau ,
\end{equation*}
we obtain $|\mathsf{r}(y)-\mathsf{r}(y')|\le|y-y'|$.

\emph{Proof of (e1).} Let $y\in K+B(0,m/2)$ and $z\in\operatorname{supp}\omega\subset B(0,m/2)$.
Then
\begin{equation*}
y-z\in K+B(0,m/2)+B(0,m/2)\subset K+B(0,m)=K',
\end{equation*}
so $f(y-z)=0$. Hence $\sigma$ vanishes on
the open set $K+B(0,m/2)$, which contains $\overline{K}$ (every point of $\overline{K}$ lies at
distance less than $m/2$ from a point of $K$). Therefore $\nabla\sigma=0$ there, and
$\mathsf{r}=\operatorname{id}$ there, in particular on $\overline{K}$. For the second assertion,
since $f\in C^1$ by (f3), differentiation under the integral sign gives
$\nabla\sigma=(\nabla f)*\omega$; by \eqref{8.3:eq-retraction-wall-b}, and since
$\int z\,\omega(z)\,dz=0$ ($\omega$ is even) and $\int\omega=1$,
\begin{equation*}
\mathsf{r}(y)=y-\int\bigl(y-z-\pi(y-z)\bigr)\omega(z)\,dz=\int\pi(y-z)\,\omega(z)\,dz .
\end{equation*}
This is an average, with respect to the probability density $\omega$, of the continuous function
$z\mapsto\pi(y-z)$ with values in the closed convex set $\overline{K'}$, and a closed convex set
contains every such average. Hence $\mathsf{r}(\mathbb{R}^n)\subset\overline{K'}$.

\emph{Proof of (e3).} By \eqref{8.3:eq-retraction-wall-c} with $\phi=\tilde\omega$,
$0\le D^2\psi(y)\le C_\psi\mathbf{1}$ with $C_\psi:=\int\tilde\omega$, which is positive because
$\tilde\omega>0$ on $B(0,m)$; the lower bound $D^2\psi\ge0$ gives the convexity of $\psi$.

Let $y\in\overline{K}$ and $|z|<m$. Choose $y_j\in K$ with $y_j\to y$; then $y_j-z\in K'$, hence
$y-z\in\overline{K'}$ and $f(y-z)=0$. Since $\tilde\omega$ vanishes outside $B(0,m)$, we get
$\psi(y)=\int_{B(0,m)}f(y-z)\tilde\omega(z)\,dz=0$.

Let $y\notin\overline{K}$. Separating $y$ from the closed convex set $\overline{K}$, we find a unit
vector $u$ with $\langle u,y\rangle>\mu_u:=\sup_{x\in K}\langle u,x\rangle$. Since
$K'=K+B(0,m)$, we have $\sup_{x\in\overline{K'}}\langle u,x\rangle=\mu_u+m$. The interval
$\bigl(\max\{0,\,m-(\langle u,y\rangle-\mu_u)\},\,m\bigr)$ is non-empty because
$\langle u,y\rangle-\mu_u>0$; for $t$ in it,
\begin{equation*}
\langle u,y+tu\rangle=\langle u,y\rangle+t>\mu_u+m=\sup_{x\in\overline{K'}}\langle u,x\rangle ,
\end{equation*}
so $y+tu\notin\overline{K'}$ and $f(y+tu)>0$, and by continuity $f>0$ near $y+tu$. Now
$y+tu=y-z$ with $z:=-tu\in B(0,m)$, where $\tilde\omega(z)>0$. The integrand
$z'\mapsto f(y-z')\tilde\omega(z')$ is continuous, non-negative, and positive at $z'=z$; hence
$\psi(y)>0$.

For the growth, $\overline{K'}$ is bounded, so $\overline{K'}\subset\overline{B}(0,R')$ for some
$R'$. For $|z|\le m$ we have $|y-z|\ge|y|-m$, and $\tau\mapsto(\tau-R')_+^2$ is non-decreasing, so
(f4) gives $f(y-z)\ge\frac12(|y|-m-R')_+^2$. Integrating against $\tilde\omega$, which is supported
in $\overline{B}(0,m)$,
\begin{equation*}
\psi(y)\ge\tfrac12C_\psi\,(|y|-R_\psi)_+^2,\qquad R_\psi:=R'+m,
\end{equation*}
which is the growth bound of (e3) with $c_\psi:=\frac12C_\psi>0$.

\emph{Proof of (e4).} Let $e$ be a unit vector. Since $\langle e,\mathsf{r}(y)\rangle
=\langle e,y\rangle-\partial_e\sigma(y)$ and derivatives of the smooth function $\sigma$ commute,
and since $\partial_e\sigma=f*\partial_e\omega$,
\begin{equation*}
\begin{aligned}
\langle e,D^2\mathsf{r}(y)[\xi,\xi]\rangle
&=-\partial_e\bigl(D^2\sigma[\xi,\xi]\bigr)(y)
=-D^2(f*\partial_e\omega)(y)[\xi,\xi]\\
&=-\lim_{t\to0}t^{-2}\int\Delta^2_{t\xi}f(y-z)\,\partial_e\omega(z)\,dz ,
\end{aligned}
\end{equation*}
the last equality being the limit formula of \eqref{8.3:eq-retraction-wall-c} for the signed
kernel $\partial_e\omega\in C_c^\infty$. Since $\Delta^2_{t\xi}f\ge0$ by (f1), and since
$|\partial_e\omega(z)|=|\langle e,\nabla\omega(z)\rangle|$ is at most $|\nabla\omega(z)|$, hence at
most $\tilde\omega(z)$, each quotient
satisfies
\begin{equation*}
\Bigl|\,t^{-2}\int\Delta^2_{t\xi}f(y-z)\,\partial_e\omega(z)\,dz\Bigr|
\le t^{-2}\int\Delta^2_{t\xi}f(y-z)\,\tilde\omega(z)\,dz .
\end{equation*}
Letting $t\to0$ and using \eqref{8.3:eq-retraction-wall-c} with $\phi=\tilde\omega$, we obtain
\begin{equation*}
\begin{split}
\bigl|\langle e,D^2\mathsf{r}(y)[\xi,\xi]\rangle\bigr|
&\le\lim_{t\to0}t^{-2}\int\Delta^2_{t\xi}f(y-z)\,\tilde\omega(z)\,dz\\
&=D^2\psi(y)[\xi,\xi].
\end{split}
\end{equation*}
Taking the
supremum over unit vectors $e$ gives $|D^2\mathsf{r}(y)[\xi,\xi]|\le D^2\psi(y)[\xi,\xi]$.
\end{proof}

\begin{lemma}[Hamiltonian flow preserves the Gibbs measure]\label{8.3:lem-flow-invariance}
Let $E$ be a real Euclidean space of finite dimension $\dim E$, with scalar product
$\langle\cdot,\cdot\rangle$, norm $|\cdot|$ and Lebesgue measure $dq$. Let $H\in C^2(E)$ satisfy
$\sup_{q\in E}\|H''(q)\|<\infty$ and $\int_E e^{-H(q)}\,dq<\infty$. Then the system
\begin{equation}\label{8.3:eq-flow-invariance-1}
\dot q=p,\qquad \dot p=-\nabla H(q)
\end{equation}
has a global flow $\Phi_t$ on $E\times E$, and $\Phi_t$ preserves the probability measure
$\nu_H\otimes\gamma$, where $\nu_H$ is the probability measure on $E$ with density proportional
to $e^{-H(q)}$ with respect to $dq$ and $\gamma:=\mathcal{N}(0,I_E)$ is the standard
Gaussian measure on $E$, that is, the centred Gaussian measure whose covariance is the identity
map $I_E$ of $E$. In particular, if $(q,p)$ is distributed according to
$\nu_H\otimes\gamma$ and $(q(t),p(t))=\Phi_t(q,p)$, then $q(t)$ is distributed according to
$\nu_H$ for every $t$.
\end{lemma}

\begin{proof}
Put $X(q,p):=(p,-\nabla H(q))$, the vector field of \eqref{8.3:eq-flow-invariance-1} on
$E\times E$. Since $H\in C^2(E)$, $X$ is $C^1$. Since $\sup\|H''\|<\infty$, the map $\nabla H$ is
Lipschitz with constant $\sup\|H''\|$, and $(q,p)\mapsto p$ is Lipschitz with constant $1$;
hence $X$ is globally Lipschitz. A globally Lipschitz $C^1$ vector field has solutions for all
times and all initial data, depending $C^1$ on the initial data. Thus \eqref{8.3:eq-flow-invariance-1}
has a global $C^1$ flow $\Phi_t$, with $\Phi_t\circ\Phi_s=\Phi_{t+s}$; in particular each $\Phi_t$ is
a $C^1$ diffeomorphism of $E\times E$ with inverse $\Phi_{-t}$.

The divergence of $X$ vanishes: in orthonormal coordinates $q_i,p_i$ of $E$,
\begin{equation*}
\operatorname{div}X=\sum_i\partial_{q_i}p_i-\sum_i\partial_{p_i}\partial_{q_i}H(q)=0,
\end{equation*}
because the first component does not depend on $q$ and the second does not depend on $p$. The
derivative $D\Phi_t(z)$ solves the variational equation
$\partial_t D\Phi_t(z)=DX(\Phi_t z)\,D\Phi_t(z)$, so by Jacobi's formula for the derivative of a
determinant, Liouville's formula holds:
\begin{equation*}
\partial_t\det D\Phi_t(z)=(\operatorname{div}X)(\Phi_t z)\,\det D\Phi_t(z).
\end{equation*}
As $\operatorname{div}X=0$ and $\det D\Phi_0=1$, we obtain $\det D\Phi_t\equiv1$ for every $t$.

Put $\mathcal{H}(q,p):=H(q)+\tfrac12|p|^2$. Along a solution of \eqref{8.3:eq-flow-invariance-1},
\begin{equation*}
\frac{d}{dt}\,\mathcal{H}(q(t),p(t))=\langle\nabla H(q),\dot q\rangle+\langle p,\dot p\rangle
=\langle\nabla H(q),p\rangle-\langle p,\nabla H(q)\rangle=0,
\end{equation*}
so $\mathcal{H}\circ\Phi_t=\mathcal{H}$ for every $t$.

The measure $\nu_H\otimes\gamma$ has density proportional to $e^{-\mathcal{H}(z)}$ with respect to
Lebesgue measure $dz$ on $E\times E$, the normalising constant being
$(2\pi)^{\dim E/2}\int_E e^{-H}\,dq<\infty$. Let $F$ be a
bounded measurable function on $E\times E$ and fix $t$. The change of variables $z'=\Phi_t(z)$,
that is $z=\Phi_{-t}(z')$, has Jacobian $|\det D\Phi_{-t}(z')|=1$; together with the conservation
of $\mathcal{H}$ it gives
\begin{equation}\label{8.3:eq-flow-invariance-2}
\int F(\Phi_t z)\,e^{-\mathcal{H}(z)}\,dz=\int F(z')\,e^{-\mathcal{H}(\Phi_{-t}z')}\,dz'
=\int F(z')\,e^{-\mathcal{H}(z')}\,dz'.
\end{equation}
Dividing by the normalising constant, \eqref{8.3:eq-flow-invariance-2} states that the image of
$\nu_H\otimes\gamma$ under $\Phi_t$ is $\nu_H\otimes\gamma$.

Finally, if $(q,p)$ is distributed according to $\nu_H\otimes\gamma$, then $(q(t),p(t))=\Phi_t(q,p)$
is distributed according to $\nu_H\otimes\gamma$ by what was just shown; taking
$F(q',p')=f(q')$ in \eqref{8.3:eq-flow-invariance-2} for bounded measurable $f$ on $E$, the first
component $q(t)$ is distributed according to the first marginal of $\nu_H\otimes\gamma$, which is
$\nu_H$.
\end{proof}

\begin{lemma}[Synchronous coupling of two Hamiltonian flows]\label{8.3:lem-synchronous-coupling}
Let $E$ be a finite-dimensional real Euclidean space with inner product
$\langle\cdot,\cdot\rangle$ and norm $|\cdot|$. Linear maps of $E$ are written as real matrices,
and $\|\cdot\|$ is the operator norm. Let $H_0,H_1\in C^2(E)$ have bounded Hessians $H_0''$,
$H_1''$. Put $\mathcal{D}:=H_1-H_0$, and let $\mathcal{T}$ be an invertible real matrix. Suppose
that there are $\lambda>0$ and $\tilde\Lambda$ such that for all $x\in E$ and all $\xi\in E$
\begin{equation}\label{8.3:eq-synchronous-coupling-1}
\bigl\langle \xi,\mathcal{T}H_0''(x)\mathcal{T}^{-1}\xi\bigr\rangle\ \ge\ \lambda|\xi|^2,
\qquad
\bigl\|\mathcal{T}H_0''(x)\mathcal{T}^{-1}\bigr\|\ \le\ \tilde\Lambda .
\end{equation}
In this lemma $q$ and $p$ denote position and momentum in $E$.
For $a=0,1$ let $(q^a(t),p^a(t))$ solve the Hamiltonian system in which $H_a$ is the potential
and the kinetic energy is $\tfrac12|p|^2$, that is,
$\dot q^a=p^a$ and $\dot p^a=-\nabla H_a(q^a)$, with $q^a(0)=q^a$ and $p^a(0)=p$. The initial
momentum $p\in E$ is the same for $a=0$ and for $a=1$. Put $v(t):=\mathcal{T}(q^1(t)-q^0(t))$.
Then for every $T>0$ with $T^2\tilde\Lambda\le 1$,
\begin{equation}\label{8.3:eq-synchronous-coupling-2}
|v(T)|\ \le\ \Bigl(1-\tfrac12\lambda T^2+\tilde\Lambda^2T^4\Bigr)|v(0)|
+\bigl(1+T^2\tilde\Lambda\bigr)\int_0^T (T-s)\,\bigl|\mathcal{T}\nabla\mathcal{D}(q^1(s))\bigr|\,ds .
\end{equation}
\end{lemma}

\begin{proof}
Since $H_a\in C^2(E)$, the trajectories $q^a$ are twice continuously differentiable. Moreover,
$\ddot q^a=\dot p^a=-\nabla H_a(q^a)$. Put $\delta:=q^1-q^0$. Because $\nabla H_1=\nabla H_0+\nabla\mathcal{D}$,
\begin{equation*}
\ddot\delta=-\nabla H_1(q^1)+\nabla H_0(q^0)
=-\bigl[\nabla H_0(q^1)-\nabla H_0(q^0)\bigr]-\nabla\mathcal{D}(q^1).
\end{equation*}
Apply the fundamental theorem of calculus to $u\mapsto\nabla H_0(q^0+u\delta)$ on $[0,1]$. This gives
$\nabla H_0(q^1)-\nabla H_0(q^0)=\mathcal{M}(t)\delta(t)$, where
\begin{equation*}
\mathcal{M}(t):=\int_0^1 H_0''\bigl(q^0(t)+u\delta(t)\bigr)\,du .
\end{equation*}
Hence $\ddot\delta=-\mathcal{M}(t)\delta-\nabla\mathcal{D}(q^1(t))$. Also
$\dot\delta(0)=p^1(0)-p^0(0)=p-p=0$, because the two initial momenta coincide.

Multiply by $\mathcal{T}$ and write $v=\mathcal{T}\delta$. Put
$\widetilde{\mathcal{M}}(t):=\mathcal{T}\mathcal{M}(t)\mathcal{T}^{-1}$ and
$\tilde e(t):=\mathcal{T}\nabla\mathcal{D}(q^1(t))$. Then
$\ddot v=-\widetilde{\mathcal{M}}v-\tilde e$ and $\dot v(0)=0$. Write $v_0:=v(0)$. Taylor's formula
with integral remainder, applied to the $C^2$ function $v$, gives for $t\ge0$
\begin{equation}\label{8.3:eq-synchronous-coupling-3}
v(t)=v_0-\int_0^t (t-s)\bigl[\widetilde{\mathcal{M}}(s)v(s)+\tilde e(s)\bigr]\,ds .
\end{equation}
For each $s$, the matrix
\begin{equation*}
\widetilde{\mathcal{M}}(s)=\int_0^1\mathcal{T}H_0''(q^0+u\delta)\mathcal{T}^{-1}\,du
\end{equation*}
is an average over $u\in[0,1]$ of matrices satisfying \eqref{8.3:eq-synchronous-coupling-1}. By
linearity of the integral, $\langle\xi,\widetilde{\mathcal{M}}(s)\xi\rangle\ge\lambda|\xi|^2$ for all
$\xi\in E$. By the triangle inequality for the integral, $\|\widetilde{\mathcal{M}}(s)\|\le\tilde\Lambda$.

Put
\begin{equation*}
\varepsilon_T:=\int_0^T(T-s)|\tilde e(s)|\,ds,
\qquad
\bar v:=\sup_{[0,T]}|v| .
\end{equation*}
Here $\bar v$ is finite because $v$ is continuous on $[0,T]$. Let $t\in[0,T]$. For $0\le s\le t$ we
have $0\le t-s\le T-s$, and $|\tilde e|\ge0$. Therefore
$\int_0^t(t-s)|\tilde e(s)|\,ds\le\varepsilon_T$. Using $\|\widetilde{\mathcal{M}}(s)\|\le\tilde\Lambda$,
$|v(s)|\le\bar v$ and $\int_0^t(t-s)\,ds=\tfrac12t^2\le\tfrac12T^2$,
\eqref{8.3:eq-synchronous-coupling-3} gives
\begin{equation}\label{8.3:eq-synchronous-coupling-4}
|v(t)-v_0|\ \le\ \tfrac12T^2\tilde\Lambda\,\bar v+\varepsilon_T
\qquad(0\le t\le T).
\end{equation}
In particular $|v(t)|\le|v_0|+\tfrac12T^2\tilde\Lambda\bar v+\varepsilon_T$ for all $t\in[0,T]$.
Taking the supremum over $t$ gives
$\bar v\le|v_0|+\tfrac12T^2\tilde\Lambda\bar v+\varepsilon_T$. Since $T^2\tilde\Lambda\le1$, the
coefficient of $\bar v$ on the right is at most $\tfrac12$. Hence $\tfrac12\bar v\le|v_0|+\varepsilon_T$,
that is, $\bar v\le 2(|v_0|+\varepsilon_T)$. Inserting this into \eqref{8.3:eq-synchronous-coupling-4}
gives
\begin{equation}\label{8.3:eq-synchronous-coupling-5}
|v(s)-v_0|\ \le\ T^2\tilde\Lambda\bigl(|v_0|+\varepsilon_T\bigr)+\varepsilon_T
\qquad(0\le s\le T).
\end{equation}

Put $\mathcal{B}:=\int_0^T(T-s)\widetilde{\mathcal{M}}(s)\,ds$. Write $v(s)=v_0+(v(s)-v_0)$ in the
integrand of \eqref{8.3:eq-synchronous-coupling-3} at $t=T$. This gives
\begin{equation}\label{8.3:eq-synchronous-coupling-6}
v(T)=(\mathbf{1}-\mathcal{B})v_0-\int_0^T(T-s)\widetilde{\mathcal{M}}(s)\bigl(v(s)-v_0\bigr)\,ds
-\int_0^T(T-s)\tilde e(s)\,ds .
\end{equation}
We bound the three terms in turn.

\emph{First term.} Expanding the square of the Euclidean norm,
\begin{equation*}
|(\mathbf{1}-\mathcal{B})v_0|^2=|v_0|^2-2\langle v_0,\mathcal{B}v_0\rangle+|\mathcal{B}v_0|^2 .
\end{equation*}
By the lower bound on $\widetilde{\mathcal{M}}(s)$,
\begin{equation*}
\langle v_0,\mathcal{B}v_0\rangle=\int_0^T(T-s)\langle v_0,\widetilde{\mathcal{M}}(s)v_0\rangle\,ds
\ge\tfrac12\lambda T^2|v_0|^2 .
\end{equation*}
By the upper bound, $\|\mathcal{B}\|\le\tfrac12T^2\tilde\Lambda$, so
$|\mathcal{B}v_0|^2\le\tfrac14\tilde\Lambda^2T^4|v_0|^2$. Hence
\begin{equation*}
|(\mathbf{1}-\mathcal{B})v_0|^2\le|v_0|^2\bigl(1-\eta\bigr),
\qquad \eta:=\lambda T^2-\tfrac14\tilde\Lambda^2T^4 .
\end{equation*}
For a unit vector $\xi$, \eqref{8.3:eq-synchronous-coupling-1} gives
$\lambda\le\langle\xi,\mathcal{T}H_0''(x)\mathcal{T}^{-1}\xi\rangle\le\tilde\Lambda$. Hence
\begin{equation*}
\eta\ \le\ \lambda T^2\ \le\ \tilde\Lambda T^2\ \le\ 1 .
\end{equation*}
For $\eta\le1$ one has $\sqrt{1-\eta}\le1-\tfrac12\eta$, because
$1-\tfrac12\eta>0$ and $(1-\tfrac12\eta)^2=1-\eta+\tfrac14\eta^2\ge1-\eta$. Therefore
\begin{equation*}
|(\mathbf{1}-\mathcal{B})v_0|\le|v_0|\bigl(1-\tfrac12\lambda T^2+\tfrac18\tilde\Lambda^2T^4\bigr).
\end{equation*}

\emph{Second term.} By $\|\widetilde{\mathcal{M}}(s)\|\le\tilde\Lambda$,
\eqref{8.3:eq-synchronous-coupling-5} and $\int_0^T(T-s)\,ds=\tfrac12T^2$, its norm is at most
\begin{equation*}
\tfrac12T^2\tilde\Lambda\Bigl[T^2\tilde\Lambda\bigl(|v_0|+\varepsilon_T\bigr)+\varepsilon_T\Bigr].
\end{equation*}

\emph{Third term.} Its norm is at most $\int_0^T(T-s)|\tilde e(s)|\,ds=\varepsilon_T$.

\emph{Collecting.} Apply the triangle inequality to \eqref{8.3:eq-synchronous-coupling-6} and insert
the three bounds. The coefficient of $|v_0|$ is
\begin{equation*}
1-\tfrac12\lambda T^2+\bigl(\tfrac18+\tfrac12\bigr)\tilde\Lambda^2T^4
\ \le\ 1-\tfrac12\lambda T^2+\tilde\Lambda^2T^4 .
\end{equation*}
The coefficient of $\varepsilon_T$ is
\begin{equation*}
1+\tfrac12T^2\tilde\Lambda\bigl(T^2\tilde\Lambda+1\bigr)\ \le\ 1+T^2\tilde\Lambda,
\end{equation*}
where we used $T^2\tilde\Lambda\le1$. Since $v_0=v(0)$ and, by the definition of $\tilde e$,
$\varepsilon_T$ is the integral on the right of \eqref{8.3:eq-synchronous-coupling-2}, this is
\eqref{8.3:eq-synchronous-coupling-2}.
\end{proof}

\begin{proposition}[Exterior-uniform Lipschitz stability of windowed Gibbs measures]
\label{8.3:prop-lipschitz-stability}
Let $I$ be a finite set with a metric $d$, let $n\ge1$, and let $E:=E_I=(\mathbb{R}^n)^I$ with the
Euclidean scalar product $\langle\cdot,\cdot\rangle$ and norm $|\cdot|$ (the space $E_{\mathfrak{B}}$
of Notation~\ref{8.3:not-block-matrices} with $\mathfrak{B}=I$). For $\varphi\colon I\to\mathbb{R}$ let
$\mathcal{T}_\varphi=\bigoplus_{b}e^{\varphi(b)}\mathbf{1}_n$ be the weight of
Notation~\ref{8.3:not-block-matrices}. Fix numbers $r_{\mathrm{in}}>0$ and $\hat{\mathfrak{b}}$, the
same for all $b\in I$. For each $b\in I$ let $K_b\subset\mathbb{R}^n$ be an open convex
set with $B(0,r_{\mathrm{in}})\subset K_b\subset B(0,\hat{\mathfrak{b}})$, and put
$\mathcal{W}:=\prod_{b\in I}K_b\subset E$.

A \emph{Hamiltonian on $\mathcal{W}$} is a function
$H(y)=\tfrac12\langle y,Ay\rangle+\langle\mathsf{h},y\rangle-P(y)$ with $A$ a real symmetric matrix on $E$,
$\mathsf{h}\in E$, and $P\in C^2(\mathcal{W})$ such that $\sup_{\mathcal{W}}|\partial_bP|$ and
$\sup_{\mathcal{W}}\|\partial_b\partial_{b'}P\|$ are finite for all $b,b'\in I$; here $\partial_bP$ is
the $b$-block of the gradient, a vector of $\mathbb{R}^n$, and $\partial_b\partial_{b'}P$ the
$(b,b')$-block of the Hessian $P''$, with its operator norm. Its Gibbs measure is
$\nu_H:=Z_H^{-1}\mathbf{1}_{\mathcal{W}}e^{-H}\,dy$, $Z_H$ the normalisation. $H$ is \emph{coercive for
$\varphi$}, with constants $\lambda_A,\varepsilon_0$, if
\begin{equation}\label{8.3:eq-lipschitz-stability-a}
\begin{gathered}
\langle\xi,\mathcal{T}_\varphi A\mathcal{T}_\varphi^{-1}\xi\rangle\ge\lambda_A|\xi|^2
\qquad(\xi\in E),\\
\sup_{y\in\mathcal{W}}\|\mathcal{T}_\varphi P''(y)\mathcal{T}_\varphi^{-1}\|\le\varepsilon_0<\lambda_A ,
\end{gathered}
\end{equation}
and then $\lambda:=\lambda_A-\varepsilon_0$.

Let $H_0,H_1$ be Hamiltonians on $\mathcal{W}$, with data $A_a,\mathsf{h}_a,P_a$ ($a=0,1$), and let $H_0$ be
coercive for $\varphi$ with $\lambda=\lambda_A-\varepsilon_0>0$. Nothing further is assumed on $H_1$: in
particular neither convexity nor smallness. Put
\begin{equation}\label{8.3:eq-lipschitz-stability-b}
\delta_b:=\sup_{y\in\mathcal{W}}\bigl|\bigl((A_1-A_0)y+\mathsf{h}_1-\mathsf{h}_0\bigr)_b\bigr|
+\sup_{y\in\mathcal{W}}\bigl|\partial_b(P_1-P_0)(y)\bigr|\qquad(b\in I).
\end{equation}
Then for every $G\colon\mathcal{W}\to\mathbb{R}$ with
$|G(y)-G(y')|\le\ell_G\,|\mathcal{T}_\varphi(y-y')|$ for all $y,y'\in\mathcal{W}$,
\begin{equation}\label{8.3:eq-lipschitz-stability-1}
\bigl|\nu_{H_1}(G)-\nu_{H_0}(G)\bigr|
\le\frac{\ell_G}{\lambda}\Bigl(\sum_{b\in I}e^{2\varphi(b)}\delta_b^2\Bigr)^{1/2}.
\end{equation}
\end{proposition}

\begin{proof}
Write $\mathcal{T}:=\mathcal{T}_\varphi$. All quantities in \eqref{8.3:eq-lipschitz-stability-1} are
finite: $\mathcal{W}$ is bounded and convex, so $|P_a(y)-P_a(0)|\le\sup_{\mathcal{W}}|\nabla P_a|\,|y|$
and $H_a$ is bounded on $\mathcal{W}$, whence $0<Z_{H_a}<\infty$; and $G$, being Lipschitz on the
bounded set $\mathcal{W}$, is bounded there.

\medskip\noindent Step~0: a margin around the window. Fix $\theta\in(0,1)$, put
$P_a^\theta(y):=P_a(\theta y)$ for $y\in\theta^{-1}\mathcal{W}$, and
$m:=(\theta^{-1/2}-1)r_{\mathrm{in}}>0$. We have $K_b+B(0,m)\subset(1+m/r_{\mathrm{in}})K_b$: for
$y\in K_b$ and $|z|<m$,
\begin{equation*}
y+z=\Bigl(1+\frac{m}{r_{\mathrm{in}}}\Bigr)\Bigl[\frac{1}{1+m/r_{\mathrm{in}}}\,y
+\frac{m/r_{\mathrm{in}}}{1+m/r_{\mathrm{in}}}\,\frac{r_{\mathrm{in}}z}{m}\Bigr],
\end{equation*}
and the bracket is a convex combination of $y\in K_b$ and $r_{\mathrm{in}}z/m\in B(0,r_{\mathrm{in}})\subset
K_b$. Since $1+m/r_{\mathrm{in}}=\theta^{-1/2}$, this gives $K_b+B(0,m)\subset\theta^{-1/2}K_b$. The
closure of the latter is $\theta^{-1/2}\overline{K}_b$, and it lies in $\theta^{-1}K_b$: for
$y\in\overline{K}_b$ one has $\theta^{-1/2}y=\theta^{-1}(\theta^{1/2}y)$, and $\theta^{1/2}y\in K_b$
because $0$ is an interior point of the convex set $K_b$, so that every point of the segment from $0$
to a point of $\overline{K}_b$, other than that endpoint, is interior. Put
$K_b':=K_b+B(0,m)$. Then $\prod_b\overline{K_b'}\subset\theta^{-1}\mathcal{W}$, an open set on which
$P_a^\theta$ is $C^2$ with
\begin{equation*}
\partial_bP_a^\theta(y)=\theta\,\partial_bP_a(\theta y),\qquad
\partial_b\partial_{b'}P_a^\theta(y)=\theta^2\,\partial_b\partial_{b'}P_a(\theta y),\qquad
\theta y\in\mathcal{W},
\end{equation*}
so its first and second derivatives are bounded there. Moreover, by
\eqref{8.3:eq-lipschitz-stability-a}, for $y\in\theta^{-1}\mathcal{W}$,
\begin{equation}\label{8.3:eq-lipschitz-stability-2}
\|\mathcal{T}(P_0^\theta)''(y)\mathcal{T}^{-1}\|\le\theta^2\varepsilon_0\le\varepsilon_0,
\qquad
|\partial_bP_0^\theta(y)|\le\kappa_P:=\max_{b\in I}\sup_{\mathcal{W}}|\partial_bP_0| .
\end{equation}

\medskip\noindent Step~1: soft walls. For each $b$ apply
Lemma~\ref{8.3:lem-retraction-wall} to the open, bounded, convex set $K_b$ and the number $m>0$; it
gives $\mathsf{r}_b\in C^\infty(\mathbb{R}^n;\mathbb{R}^n)$, $\psi_b\in C^\infty(\mathbb{R}^n;[0,\infty))$
and numbers $C_{\psi,b},c_{\psi,b}>0$, $R_{\psi,b}$ with the properties (e1)--(e4) of
\eqref{8.3:eq-retraction-wall-a}. Let $\mathsf{r}:=\prod_b\mathsf{r}_b$, i.e.
$\mathsf{r}(y)_b=\mathsf{r}_b(y_b)$; by (e1), $\mathsf{r}(E)\subset\prod_b\overline{K_b'}$. For
$\varepsilon>0$ and $a=0,1$ define on all of $E$
\begin{equation}\label{8.3:eq-lipschitz-stability-3}
H_a^\varepsilon(y):=\tfrac12\langle y,A_ay\rangle+\langle\mathsf{h}_a,y\rangle-P_a^\theta(\mathsf{r}(y))
+(\kappa_P+\varepsilon^{-1})\sum_{b\in I}\psi_b(y_b).
\end{equation}
This is well defined and $C^2$ because $\mathsf{r}$ takes values in $\prod_b\overline{K_b'}$, where
$P_a^\theta$ is $C^2$ (Step~0). Its Hessian is bounded: by the chain rule (displayed below for $a=0$)
it is built from the bounded derivatives of $P_a^\theta$ on $\prod_b\overline{K_b'}$, from
$D\mathsf{r}$, bounded by (e2), from $D^2\mathsf{r}_b$, bounded because
$|D^2\mathsf{r}_b(y)[\xi,\xi]|\le D^2\psi_b(y)[\xi,\xi]\le C_{\psi,b}|\xi|^2$ by (e4) and (e3) and
polarisation, and from $D^2\psi_b\le C_{\psi,b}$.

Integrability. $P_a^\theta$ is continuous on the compact set $\prod_b\overline{K_b'}$, so
$c_{\mathrm{bd}}:=\max_a\sup_E|P_a^\theta\circ\mathsf{r}|<\infty$. For $u\ge0$ and $v\ge0$ one has
$(u-v)_+^2\ge\tfrac12u^2-v^2$ (for $u\ge v$ the difference is $\tfrac12(u-2v)^2$; for $u<v$ the right
side is negative). With (e3), $\psi_b\ge0$ and $\kappa_P\ge0$, putting
$c_{\min}:=\min_bc_{\psi,b}$,
\begin{equation*}
H_a^\varepsilon(y)\ge\tfrac12\bigl(\varepsilon^{-1}c_{\min}-\|A_a\|\bigr)|y|^2-|\mathsf{h}_a|\,|y|
-c_{\mathrm{bd}}-\varepsilon^{-1}\sum_bc_{\psi,b}R_{\psi,b}^2 .
\end{equation*}
Choose $\varepsilon_1>0$ with $\varepsilon_1^{-1}c_{\min}\ge\max_a\|A_a\|+2$. Since
$\varepsilon\mapsto H_a^\varepsilon(y)$ is non-increasing, $H_a^\varepsilon\ge H_a^{\varepsilon_1}$ for
$\varepsilon\le\varepsilon_1$, and $H_a^{\varepsilon_1}(y)\ge|y|^2-|\mathsf{h}_a||y|-C'$ for a constant $C'$;
hence $\int_E(1+|y|)e^{-H_a^\varepsilon}\,dy<\infty$ for $\varepsilon\le\varepsilon_1$.

The walls cancel in the difference:
\begin{equation}\label{8.3:eq-lipschitz-stability-4}
\mathcal{D}^\theta:=H_1^\varepsilon-H_0^\varepsilon
=\tfrac12\langle y,(A_1-A_0)y\rangle+\langle\mathsf{h}_1-\mathsf{h}_0,y\rangle
-(P_1^\theta-P_0^\theta)(\mathsf{r}(y)),
\end{equation}
which does not depend on $\varepsilon$. Its gradient
\begin{equation*}
\nabla\mathcal{D}^\theta(y)=(A_1-A_0)y+\mathsf{h}_1-\mathsf{h}_0
-D\mathsf{r}(y)\,\nabla(P_1^\theta-P_0^\theta)(\mathsf{r}(y))
\end{equation*}
is continuous and of at most linear growth, the last term being bounded.

Hessian of $H_0^\varepsilon$. Since $\mathsf{r}_b$ depends on $y_b$ only, the chain rule gives, for
$\xi\in E$,
\begin{equation}\label{8.3:eq-lipschitz-stability-5}
\begin{split}
D^2H_0^\varepsilon(y)[\xi,\xi]&=\langle\xi,A_0\xi\rangle
-(P_0^\theta)''(\mathsf{r}(y))\bigl[D\mathsf{r}(y)\xi,D\mathsf{r}(y)\xi\bigr]\\
&\quad+\sum_{b\in I}\Bigl\{(\kappa_P+\varepsilon^{-1})D^2\psi_b(y_b)[\xi_b,\xi_b]
-\bigl\langle\partial_bP_0^\theta(\mathsf{r}(y)),D^2\mathsf{r}_b(y_b)[\xi_b,\xi_b]\bigr\rangle\Bigr\}.
\end{split}
\end{equation}
Let $\Gamma_b(y)$ be the symmetric $n\times n$ matrix whose quadratic form is the $b$-th summand in
braces, and $\Gamma(y):=\bigoplus_b\Gamma_b(y)$; let $D\mathsf{r}(y)=\bigoplus_bD\mathsf{r}_b(y_b)$. Then
\begin{equation*}
(H_0^\varepsilon)''=A_0-D\mathsf{r}\,(P_0^\theta)''(\mathsf{r})\,D\mathsf{r}+\Gamma,
\end{equation*}
where $D\mathsf{r}$ is symmetric by (e2), and $D\mathsf{r}$ and $\Gamma$ are bond-diagonal. Since
$\mathsf{r}(y)\in\prod_b\overline{K_b'}$, \eqref{8.3:eq-lipschitz-stability-2} gives
$|\partial_bP_0^\theta(\mathsf{r}(y))|\le\kappa_P$, so by (e4)
\begin{equation*}
\langle\xi_b,\Gamma_b(y)\xi_b\rangle\ge(\kappa_P+\varepsilon^{-1})D^2\psi_b(y_b)[\xi_b,\xi_b]
-\kappa_PD^2\psi_b(y_b)[\xi_b,\xi_b]=\varepsilon^{-1}D^2\psi_b(y_b)[\xi_b,\xi_b]\ge0,
\end{equation*}
the last inequality by (e3). A bond-diagonal matrix commutes with $\mathcal{T}$, which acts on each block
as the scalar $e^{\varphi(b)}$. Hence
\begin{equation*}
\mathcal{T}(H_0^\varepsilon)''\mathcal{T}^{-1}=\mathcal{T}A_0\mathcal{T}^{-1}
-D\mathsf{r}\,\bigl(\mathcal{T}(P_0^\theta)''(\mathsf{r})\mathcal{T}^{-1}\bigr)\,D\mathsf{r}+\Gamma .
\end{equation*}
For the middle term, put $\Pi(y):=\mathcal{T}(P_0^\theta)''(\mathsf{r}(y))\mathcal{T}^{-1}$; then
$\|\Pi(y)\|\le\varepsilon_0$ by \eqref{8.3:eq-lipschitz-stability-2}, and
$\langle\xi,D\mathsf{r}\,\Pi\,D\mathsf{r}\,\xi\rangle=\langle D\mathsf{r}\xi,\Pi D\mathsf{r}\xi\rangle$
with $|D\mathsf{r}\xi|\le|\xi|$ since
$0\le D\mathsf{r}_b\le\mathbf{1}$ by (e2). With \eqref{8.3:eq-lipschitz-stability-a} we obtain, for all
$y\in E$, $\xi\in E$ and $\varepsilon>0$,
\begin{equation}\label{8.3:eq-lipschitz-stability-c}
\langle\xi,\mathcal{T}(H_0^\varepsilon)''(y)\mathcal{T}^{-1}\xi\rangle
\ge\lambda_A|\xi|^2-\varepsilon_0|D\mathsf{r}(y)\xi|^2+\langle\xi,\Gamma(y)\xi\rangle
\ge\lambda|\xi|^2 .
\end{equation}
Since the Hessian is bounded,
\begin{equation*}
\tilde{\Lambda}_\varepsilon:=\sup_y\|\mathcal{T}(H_0^\varepsilon)''(y)\mathcal{T}^{-1}\|<\infty .
\end{equation*}

\medskip\noindent Step~2: coupling. Fix $\varepsilon\le\varepsilon_1$, let
$\nu_a^\varepsilon:=\bigl(\int_Ee^{-H_a^\varepsilon}\bigr)^{-1}e^{-H_a^\varepsilon}dy$ on $E$, and let
$\gamma$ be the standard Gaussian measure $\mathcal{N}(0,\mathbf{1}_E)$ of
Lemma~\ref{8.3:lem-flow-invariance}. On one probability space, with expectation $\mathbb{E}$, take
$(q_0^0,q_0^1)$ with law $\nu_0^\varepsilon\otimes\nu_1^\varepsilon$ and $p_0,p_1,\dots$ independent,
each of law $\gamma$, and independent of $(q_0^0,q_0^1)$. Fix $T>0$. For $l\ge0$ and $a=0,1$ let
$s\mapsto(q_l^a(s),p_l^a(s))$ be the solution of $\dot q=p$, $\dot p=-\nabla H_a^\varepsilon(q)$ with
initial data $(q_l^a,p_l)$ --- the same $p_l$ for $a=0$ and $a=1$ --- and put $q_{l+1}^a:=q_l^a(T)$.
The flows exist globally by Lemma~\ref{8.3:lem-flow-invariance}, whose hypotheses hold for
$H_a^\varepsilon$ by Step~1.

By induction on $l$: $q_l^a$ is a function of $(q_0^0,q_0^1,p_0,\dots,p_{l-1})$, hence independent of
$p_l$; if $q_l^a$ has law $\nu_a^\varepsilon$, then $(q_l^a,p_l)$ has law $\nu_a^\varepsilon\otimes\gamma$,
and Lemma~\ref{8.3:lem-flow-invariance} shows that $q_l^a(s)$ has law $\nu_a^\varepsilon$ for every
$s\ge0$; in particular $q_{l+1}^a$ does. Thus $q_l^a\sim\nu_a^\varepsilon$ for all $l$, and
$q_l^1(s)\sim\nu_1^\varepsilon$ for every intermediate time $s\in[0,T]$.

Let $\eta_l:=\mathcal{T}(q_l^1-q_l^0)$. Apply Lemma~\ref{8.3:lem-synchronous-coupling} to the pair
$H_0^\varepsilon,H_1^\varepsilon$, the matrix $\mathcal{T}$ and $\mathcal{D}=\mathcal{D}^\theta$: both
functions are $C^2$ with bounded Hessians (Step~1), and its coercivity and norm hypotheses hold with
$\lambda$ and $\tilde{\Lambda}_\varepsilon$ by \eqref{8.3:eq-lipschitz-stability-c} and the definition
of $\tilde{\Lambda}_\varepsilon$. For $T^2\tilde{\Lambda}_\varepsilon\le1$ it gives, pathwise,
\begin{equation*}
|\eta_{l+1}|\le\bigl(1-\tfrac12\lambda T^2+\tilde{\Lambda}_\varepsilon^2T^4\bigr)|\eta_l|
+(1+T^2\tilde{\Lambda}_\varepsilon)\int_0^T(T-s)\,|\mathcal{T}\nabla\mathcal{D}^\theta(q_l^1(s))|\,ds .
\end{equation*}
The integrand is a non-negative measurable function of $s$ and of the randomness (the flow is
continuous), so by Fubini and $q_l^1(s)\sim\nu_1^\varepsilon$,
\begin{equation*}
\mathbb{E}\int_0^T(T-s)|\mathcal{T}\nabla\mathcal{D}^\theta(q_l^1(s))|\,ds
=\int_0^T(T-s)\,\nu_1^\varepsilon(|\mathcal{T}\nabla\mathcal{D}^\theta|)\,ds
=\tfrac12T^2\,\nu_1^\varepsilon(|\mathcal{T}\nabla\mathcal{D}^\theta|),
\end{equation*}
which is finite by the linear growth of $\nabla\mathcal{D}^\theta$ and the integrability of Step~1.
Hence, with $\sigma_T:=1-\tfrac12\lambda T^2+\tilde{\Lambda}_\varepsilon^2T^4$,
\begin{equation}\label{8.3:eq-lipschitz-stability-6}
\mathbb{E}|\eta_{l+1}|\le\sigma_T\,\mathbb{E}|\eta_l|
+(1+T^2\tilde{\Lambda}_\varepsilon)\,\tfrac12T^2\,\nu_1^\varepsilon(|\mathcal{T}\nabla\mathcal{D}^\theta|).
\end{equation}
Here $\mathbb{E}|\eta_0|<\infty$, since $\nu_a^\varepsilon$ has a finite first moment, and by
\eqref{8.3:eq-lipschitz-stability-6} all $\mathbb{E}|\eta_l|$ are finite. Now $\sigma_T<1$ if and only if
$2\tilde{\Lambda}_\varepsilon^2T^2<\lambda$; and $\sigma_T\ge0$, because
$\lambda\le\tilde{\Lambda}_\varepsilon$ (by \eqref{8.3:eq-lipschitz-stability-c}) and
$T^2\tilde{\Lambda}_\varepsilon\le1$ give $\tfrac12\lambda T^2\le\tfrac12$. For such $T$, iterating
\eqref{8.3:eq-lipschitz-stability-6} yields
$\mathbb{E}|\eta_l|\le\sigma_T^l\,\mathbb{E}|\eta_0|+c_T/(1-\sigma_T)$, with $c_T$ the last term of
\eqref{8.3:eq-lipschitz-stability-6}, and therefore
\begin{equation}\label{8.3:eq-lipschitz-stability-7}
\limsup_{l\to\infty}\mathbb{E}|\eta_l|
\le\frac{(1+T^2\tilde{\Lambda}_\varepsilon)\,\nu_1^\varepsilon(|\mathcal{T}\nabla\mathcal{D}^\theta|)}
{\lambda-2\tilde{\Lambda}_\varepsilon^2T^2},
\end{equation}
using
\begin{equation*}
\frac{\tfrac12T^2}{\tfrac12\lambda T^2-\tilde{\Lambda}_\varepsilon^2T^4}
=\frac{1}{\lambda-2\tilde{\Lambda}_\varepsilon^2T^2}.
\end{equation*}

Extend $G$ to $E$ by $\tilde G(y):=\inf_{y'\in\mathcal{W}}\bigl[G(y')+\ell_G|\mathcal{T}(y-y')|\bigr]$.
For $y\in\mathcal{W}$, the choice $y'=y$ gives $\tilde G(y)\le G(y)$, while
$G(y')+\ell_G|\mathcal{T}(y-y')|\ge G(y)$ by the Lipschitz condition; so $\tilde G=G$ on $\mathcal{W}$.
For any $y_0\in\mathcal{W}$ the same condition and the triangle inequality give
$\tilde G(y)\ge G(y_0)-\ell_G|\mathcal{T}(y-y_0)|>-\infty$. For $y,z\in E$ and $y'\in\mathcal{W}$, the
triangle inequality $|\mathcal{T}(y-y')|\le|\mathcal{T}(z-y')|+|\mathcal{T}(y-z)|$, multiplied by
$\ell_G$ and added to $G(y')$, bounds the bracket at $y$ by the bracket at $z$ plus
$\ell_G|\mathcal{T}(y-z)|$; taking the
infimum over $y'$ and exchanging $y,z$ gives $|\tilde G(y)-\tilde G(z)|\le\ell_G|\mathcal{T}(y-z)|$. In
particular $\tilde G$ is continuous and of at most linear growth, hence $\nu_a^\varepsilon$-integrable.
For every $l$, since $q_l^a\sim\nu_a^\varepsilon$,
\begin{equation*}
|\nu_1^\varepsilon(\tilde G)-\nu_0^\varepsilon(\tilde G)|
=\bigl|\mathbb{E}[\tilde G(q_l^1)-\tilde G(q_l^0)]\bigr|\le\ell_G\,\mathbb{E}|\eta_l| .
\end{equation*}
Letting $l\to\infty$ and using \eqref{8.3:eq-lipschitz-stability-7}, then letting $T\to0$ (the
conditions $T^2\tilde{\Lambda}_\varepsilon\le1$ and $2\tilde{\Lambda}_\varepsilon^2T^2<\lambda$ hold for
all small $T$), we obtain
\begin{equation}\label{8.3:eq-lipschitz-stability-d}
|\nu_1^\varepsilon(\tilde G)-\nu_0^\varepsilon(\tilde G)|
\le\frac{\ell_G}{\lambda}\,\nu_1^\varepsilon(|\mathcal{T}\nabla\mathcal{D}^\theta|)
\qquad\text{for every }\varepsilon\le\varepsilon_1 ;
\end{equation}
the constant $\tilde{\Lambda}_\varepsilon$ no longer appears.

\medskip\noindent Step~3: hard walls. Let
$H_a^\theta(y):=\tfrac12\langle y,A_ay\rangle+\langle\mathsf{h}_a,y\rangle-P_a^\theta(y)$. On
$\overline{\mathcal{W}}$ one has $\psi_b(y_b)=0$ for all $b$ by (e3) and $\mathsf{r}(y)=y$ by (e1), so
$H_a^\varepsilon=H_a^\theta$ there for every $\varepsilon$. Off $\overline{\mathcal{W}}$ some
$y_b\notin\overline{K}_b$, so $\psi_b(y_b)>0$ by (e3) and $H_a^\varepsilon(y)\to+\infty$ as
$\varepsilon\downarrow0$. Since $H_a^\varepsilon$ is non-increasing in $\varepsilon$, $e^{-H_a^\varepsilon}$
decreases pointwise, as $\varepsilon\downarrow0$, to
$\mathbf{1}_{\overline{\mathcal{W}}}e^{-H_a^\theta}$, and is dominated by $e^{-H_a^{\varepsilon_1}}$ for
$\varepsilon\le\varepsilon_1$. The set $\overline{\mathcal{W}}\setminus\mathcal{W}$ is contained in
$\bigcup_b\{y:y_b\in\partial K_b\}$, and the boundary of a bounded convex set of $\mathbb{R}^n$ is
Lebesgue-null; so $\mathbf{1}_{\overline{\mathcal{W}}}=\mathbf{1}_{\mathcal{W}}$ almost everywhere. By
dominated convergence, applied to numerators and to the (positive) denominators, for every continuous $F$
with $|F(y)|\le C(1+|y|)$,
\begin{equation*}
\nu_a^\varepsilon(F)\longrightarrow\nu_{H_a^\theta}(F)
:=Z_{H_a^\theta}^{-1}\int_{\mathcal{W}}F\,e^{-H_a^\theta}\,dy\qquad(\varepsilon\downarrow0).
\end{equation*}
Both $\tilde G$ and $|\mathcal{T}\nabla\mathcal{D}^\theta|$ are such functions (Steps~1 and~2). Passing
to the limit in \eqref{8.3:eq-lipschitz-stability-d}, and using that $\nu_{H_a^\theta}$ lives on
$\mathcal{W}$ where $\tilde G=G$,
\begin{equation*}
|\nu_{H_1^\theta}(G)-\nu_{H_0^\theta}(G)|
\le\frac{\ell_G}{\lambda}\,\nu_{H_1^\theta}(|\mathcal{T}\nabla\mathcal{D}^\theta|).
\end{equation*}
On $\overline{\mathcal{W}}$, $\mathsf{r}(y)=y$, and $D\mathsf{r}_b=\mathbf{1}$ on the open set $K_b$,
hence on $\overline{K}_b$ by continuity; so by \eqref{8.3:eq-lipschitz-stability-4}
\begin{equation*}
\partial_b\mathcal{D}^\theta(y)=\bigl((A_1-A_0)y+\mathsf{h}_1-\mathsf{h}_0\bigr)_b
-\theta\,\partial_b(P_1-P_0)(\theta y).
\end{equation*}
For $y\in\mathcal{W}$ one has $\theta y\in\mathcal{W}$ and $\theta<1$, so by
\eqref{8.3:eq-lipschitz-stability-b} $|\partial_b\mathcal{D}^\theta(y)|\le\delta_b$. Since
$(\mathcal{T}v)_b=e^{\varphi(b)}v_b$, for $y\in\mathcal{W}$
\begin{equation*}
|\mathcal{T}\nabla\mathcal{D}^\theta(y)|
=\Bigl(\sum_be^{2\varphi(b)}|\partial_b\mathcal{D}^\theta(y)|^2\Bigr)^{1/2}
\le\Bigl(\sum_be^{2\varphi(b)}\delta_b^2\Bigr)^{1/2}.
\end{equation*}
Therefore
\begin{equation}\label{8.3:eq-lipschitz-stability-8}
|\nu_{H_1^\theta}(G)-\nu_{H_0^\theta}(G)|
\le\frac{\ell_G}{\lambda}\Bigl(\sum_{b\in I}e^{2\varphi(b)}\delta_b^2\Bigr)^{1/2}.
\end{equation}

\medskip\noindent Step~4: $\theta\uparrow1$. For $y\in\mathcal{W}$, $\theta y\in\mathcal{W}$ and
$P_a^\theta(y)=P_a(\theta y)\to P_a(y)$ by continuity of $P_a$ on $\mathcal{W}$; and
$|P_a^\theta|\le\sup_{\mathcal{W}}|P_a|<\infty$ (shown at the beginning of the proof). Hence
$e^{-H_a^\theta}\to e^{-H_a}$ pointwise on $\mathcal{W}$, boundedly in $\theta$, and $G$ is bounded on
the bounded set $\mathcal{W}$. Dominated convergence in $\int_{\mathcal{W}}Ge^{-H_a^\theta}$ and in
$\int_{\mathcal{W}}e^{-H_a^\theta}$, whose limit $Z_{H_a}$ is positive, gives
$\nu_{H_a^\theta}(G)\to\nu_{H_a}(G)$. The right side of \eqref{8.3:eq-lipschitz-stability-8} does not
depend on $\theta$, and \eqref{8.3:eq-lipschitz-stability-1} follows.
\end{proof}

\begin{remark}[Where the structure of the window enters]
(1) The product structure and the convexity of $\mathcal{W}$ enter the proof of
Proposition~\ref{8.3:prop-lipschitz-stability} only through the fact that the wall matrix $\Gamma$ is
bond-diagonal and non-negative, i.e.\ $\mathcal{T}\Gamma\mathcal{T}^{-1}=\Gamma\ge0$. For a convex
constraint that couples bonds, a wall matrix $\Gamma=vv^{\mathsf{T}}$ with $v$ spread over several bonds
gives $\mathcal{T}\Gamma\mathcal{T}^{-1}=(\mathcal{T}v)(\mathcal{T}^{-1}v)^{\mathsf{T}}$, whose symmetric
part has the eigenvalue
\begin{equation*}
\tfrac12\bigl(\langle\mathcal{T}v,\mathcal{T}^{-1}v\rangle-|\mathcal{T}v|\,|\mathcal{T}^{-1}v|\bigr),
\end{equation*}
negative when $\mathcal{T}v$ and $\mathcal{T}^{-1}v$ are not parallel, and unbounded as the wall hardens;
the weighted estimate \eqref{8.3:eq-lipschitz-stability-c} would fail. For a non-convex window
$\Gamma\not\ge0$, as the example given with \eqref{8.3:eq-window-product-a} shows.

(2) Gaussian case. Take, formally, $\mathcal{W}=E$, $P_a=0$, $A_1=A_0$, $\mathsf{h}_0=0$. Then
$\varepsilon_0=0$, $\lambda=\lambda_A$, $\delta_b=|(\mathsf{h}_1)_b|$, and
$\nu_1(y_b)-\nu_0(y_b)=-(A_0^{-1}\mathsf{h}_1)_b$. Testing with $G(y)=\langle u,y_b\rangle$ for a unit
vector $u\in\mathbb{R}^n$, for which $\ell_G=e^{-\varphi(b)}$, the bound
\eqref{8.3:eq-lipschitz-stability-1} reads
\begin{equation*}
|(A_0^{-1}\mathsf{h}_1)_b|\le\lambda^{-1}e^{-\varphi(b)}\,|\mathcal{T}_\varphi\mathsf{h}_1| ,
\end{equation*}
which is a weighted bound on the blocks of $A_0^{-1}$.

(3) No property of $H_1$ enters beyond the numbers $\delta_b$. This is why an exterior frozen at the
window's edge, or a non-small boundary term, costs nothing but its size times the weight.
\end{remark}

\begin{theorem}[Two-sided density bound on the read set]\label{8.3:thm-density-bound}
Let the finite index set $I$ with its metric $d$, the open convex sets $K_b\subset\mathbb{R}^n$
\textup{(}$b\in I$\textup{)}, the product $\mathcal{W}=\prod_{b\in I}K_b$, the Hamiltonians
$H_0,H_1$ on $\mathcal{W}$ with their data $A_a,h_a,P_a$ \textup{(}$a=0,1$\textup{)}, the
probability measures $\nu_{H_a}$, and the numbers $\delta_b$ \textup{(}$b\in I$\textup{)} of
\eqref{8.3:eq-lipschitz-stability-b} be as in Proposition~\ref{8.3:prop-lipschitz-stability}.
Let $S^+\subset I$ \textup{(}the \emph{read set}\textup{)}, put $R:=I\setminus S^+$
\textup{(}a letter local to this theorem\textup{)},
and let $\mu>0$. Suppose that $H_0$ is coercive, in the sense of
\eqref{8.3:eq-lipschitz-stability-a}, with one and the same pair $(\lambda_A,\varepsilon_0)$ for
every function $\varphi$ on $I$ with $\operatorname{Lip}(\varphi)\le\mu$, and let
$\lambda:=\lambda_A-\varepsilon_0$. For $s\in S^+$ put
\begin{align}
\mathfrak{L}_s&:=\sup_{y\in\mathcal{W}}\Bigl(\sum_{b\in R}e^{2\mu d(b,s)}
\bigl\|\bigl(A_0-P_0''(y)\bigr)(b,s)\bigr\|^2\Bigr)^{1/2},\label{8.3:eq-density-bound-1}\\
\mathfrak{D}_s&:=\Bigl(\sum_{b\in R}e^{-2\mu d(b,s)}\delta_b^2\Bigr)^{1/2},
\label{8.3:eq-density-bound-2}
\end{align}
and
\begin{equation}\label{8.3:eq-density-bound-3}
\Xi:=\sum_{s\in S^+}\operatorname{diam}(K_s)\Bigl[\delta_s+\frac{\mathfrak{L}_s\mathfrak{D}_s}{\lambda}\Bigr].
\end{equation}
Let $m_a$ be the density of the marginal of $\nu_{H_a}$ on
$\mathcal{W}_{S^+}:=\prod_{s\in S^+}K_s$. Then
\begin{equation}\label{8.3:eq-density-bound-4}
e^{-\Xi}\le\frac{m_1(u)}{m_0(u)}\le e^{\Xi}\qquad\text{for every }u\in\mathcal{W}_{S^+}.
\end{equation}
\end{theorem}

Here, for a real symmetric matrix $B$ on $E_I$ and $b,b'\in I$, $B(b,b')$ denotes the
$n\times n$ block of $B$ coupling the components at $b'$ to those at $b$, and $\|\cdot\|$ is the
operator norm of an $n\times n$ matrix. For $J,J'\subset I$, $B_{JJ'}$ is the block of $B$ with
rows in $J$ and columns in $J'$, and $h_J$ is the restriction of $h\in E_I$ to $J$; $P''$ is the
Hessian of $P$ and $\partial_s$ is the gradient with respect to the component $y_s\in\mathbb{R}^n$.

\begin{proof}
Write $\mathcal{W}_R:=\prod_{b\in R}K_b$ and $y=(u,v)$ with $u\in\mathcal{W}_{S^+}$ and
$v\in\mathcal{W}_R$. Since $\mathcal{W}$ is a product, $\mathcal{W}=\mathcal{W}_{S^+}\times\mathcal{W}_R$, and
\begin{equation*}
m_a(u)=Z_a^{-1}\int_{\mathcal{W}_R}e^{-H_a(u,v)}\,dv ,
\end{equation*}
where $Z_a$ is the normalisation of $\nu_{H_a}$. The domain of integration $\mathcal{W}_R$ does not
depend on $u$; this is where the product structure of $\mathcal{W}$ enters. The function $H_a$ is
of class $C^1$ on $\mathcal{W}$ with bounded gradient, because $A_a$ and $h_a$ are fixed, $\mathcal{W}$
is bounded, and $\sup_{\mathcal{W}}|\partial_bP_a|$ is finite for every $b$. Hence one may
differentiate under the integral sign, and
\begin{equation*}
\partial_s\log m_a(u)=-\nu_a^u\bigl(\partial_sH_a(u,\cdot)\bigr),\qquad s\in S^+ ,
\end{equation*}
where $\nu_a^u:=\nu_{H_a(u,\cdot)}$ is the probability measure on $\mathcal{W}_R$ with density
proportional to $e^{-H_a(u,v)}$. With $\mathcal{D}:=H_1-H_0$ we have
$\partial_sH_1=\partial_s\mathcal{D}+\partial_sH_0$. Subtracting the two identities therefore gives
\begin{equation}\label{8.3:eq-density-bound-5}
\partial_s\log\frac{m_1}{m_0}(u)=-\nu_1^u\bigl(\partial_s\mathcal{D}(u,\cdot)\bigr)
-\Bigl[\nu_1^u\bigl(\partial_sH_0(u,\cdot)\bigr)-\nu_0^u\bigl(\partial_sH_0(u,\cdot)\bigr)\Bigr].
\end{equation}

\emph{The first term.} By \eqref{8.3:eq-lipschitz-stability-b}, $|\partial_s\mathcal{D}|\le\delta_s$
on $\mathcal{W}$. Since $\nu_1^u$ is a probability measure, the first term of
\eqref{8.3:eq-density-bound-5} has modulus at most $\delta_s$.

\emph{The restricted Hamiltonians.} For the bracket we apply
Proposition~\ref{8.3:prop-lipschitz-stability} on $\mathcal{W}_R$, with index set $R$ and the
restriction of $d$, to the Hamiltonians
\begin{equation*}
v\mapsto H_a(u,v)=\tfrac12\langle v,(A_a)_{RR}v\rangle
+\langle (h_a)_R+(A_a)_{RS^+}u,\,v\rangle-P_a(u,v)+\mathrm{const},
\end{equation*}
$a=0,1$. The constant depends on $u$ only and does not change $\nu_a^u$. The set $\mathcal{W}_R$ is
the product of the $K_b$, $b\in R$, so it is of the form required there, and $v\mapsto P_a(u,v)$
has the required regularity. The numbers that \eqref{8.3:eq-lipschitz-stability-b} attaches to these
two Hamiltonians are at most $\delta_b$, $b\in R$: they are the same expressions as those defining
$\delta_b$, evaluated only at the points $(u,v)\in\mathcal{W}$.

Let $\varphi:=-\mu d(\cdot,s)$ on $I$. By the triangle inequality $\operatorname{Lip}(\varphi)\le\mu$,
so $H_0$ is coercive for $\varphi$ with the pair $(\lambda_A,\varepsilon_0)$. Write
$\mathcal{T}_R$ for the weight of Proposition~\ref{8.3:prop-lipschitz-stability} attached to the
restriction $\varphi|_R$, and $\mathcal{T}=\mathcal{T}_\varphi$. Then $H_0(u,\cdot)$ is coercive for
$\varphi|_R$ with the same constants. Indeed, for $\xi$ supported on $R$,
\begin{equation*}
\langle\xi,\mathcal{T}_R(A_0)_{RR}\mathcal{T}_R^{-1}\xi\rangle
=\langle\xi,\mathcal{T}A_0\mathcal{T}^{-1}\xi\rangle\ge\lambda_A|\xi|^2 .
\end{equation*}
Moreover, the Hessian of $v\mapsto P_0(u,v)$ is $P_0''(u,v)_{RR}$, and
$\mathcal{T}_RP_0''(u,v)_{RR}\mathcal{T}_R^{-1}$ is the compression to $R$ of
$\mathcal{T}P_0''(u,v)\mathcal{T}^{-1}$. A compression does not increase the norm, so its norm is at
most $\varepsilon_0$. The constant $\lambda=\lambda_A-\varepsilon_0$ is the same.

\emph{The test functions.} For a unit vector $e\in\mathbb{R}^n$ put
$G_e(v):=\langle e,\partial_sH_0(u,v)\rangle$, a function in $C^1(\mathcal{W}_R)$. Differentiating
the quadratic part and $P_0$ gives
\begin{equation*}
\nabla_bG_e(v)=\bigl(A_0-P_0''(u,v)\bigr)(b,s)\,e,\qquad b\in R .
\end{equation*}
Hence $|\nabla_bG_e(v)|\le\|(A_0-P_0''(u,v))(b,s)\|$. For $\varphi=-\mu d(\cdot,s)$ the weight
$\mathcal{T}_R^{-1}$ carries the factor $e^{\mu d(b,s)}$ at the site $b$, so
\eqref{8.3:eq-density-bound-1} gives
\begin{equation*}
\sup_{v\in\mathcal{W}_R}\bigl|\mathcal{T}_R^{-1}\nabla G_e(v)\bigr|\le\mathfrak{L}_s .
\end{equation*}
The set $\mathcal{W}_R$ is convex. For $v,v'\in\mathcal{W}_R$ put $v_t:=v'+t(v-v')$. Since
$\mathcal{T}_R$ is the diagonal weight,
\begin{equation*}
G_e(v)-G_e(v')=\int_0^1\bigl\langle\mathcal{T}_R^{-1}\nabla G_e(v_t),\,
\mathcal{T}_R(v-v')\bigr\rangle\,dt .
\end{equation*}
Applying the Cauchy--Schwarz inequality inside the integral,
\begin{equation*}
|G_e(v)-G_e(v')|\le\sup_{\mathcal{W}_R}\bigl|\mathcal{T}_R^{-1}\nabla G_e\bigr|\,
\bigl|\mathcal{T}_R(v-v')\bigr|\le\mathfrak{L}_s\bigl|\mathcal{T}_R(v-v')\bigr| .
\end{equation*}
Proposition~\ref{8.3:prop-lipschitz-stability}, applied with the weight $\varphi|_R$, the Lipschitz
constant $\mathfrak{L}_s$ and the numbers $\delta_b$, $b\in R$, whose weighted norm is
$\mathfrak{D}_s$ by \eqref{8.3:eq-density-bound-2}, gives
\begin{equation*}
\bigl|\nu_1^u(G_e)-\nu_0^u(G_e)\bigr|\le\frac{\mathfrak{L}_s}{\lambda}\,\mathfrak{D}_s .
\end{equation*}
The left-hand side is the modulus of the scalar product of the bracket in
\eqref{8.3:eq-density-bound-5} with $e$. Since $e$ is an arbitrary unit vector, the bracket has
modulus at most $\mathfrak{L}_s\mathfrak{D}_s/\lambda$.

\emph{Integration.} Combining the two bounds in \eqref{8.3:eq-density-bound-5},
\begin{equation*}
\Bigl|\partial_s\log\frac{m_1}{m_0}(u)\Bigr|\le\delta_s+\frac{\mathfrak{L}_s\mathfrak{D}_s}{\lambda}
\qquad\text{for all }s\in S^+,\ u\in\mathcal{W}_{S^+}.
\end{equation*}
The set $\mathcal{W}_{S^+}$ is convex. For $u,u'\in\mathcal{W}_{S^+}$ put $u_t:=u'+t(u-u')$. Since
$|u_s-u'_s|\le\operatorname{diam}(K_s)$,
\begin{align*}
\Bigl|\log\frac{m_1}{m_0}(u)-\log\frac{m_1}{m_0}(u')\Bigr|
&=\Bigl|\int_0^1\sum_{s\in S^+}\Bigl\langle\partial_s\log\frac{m_1}{m_0}(u_t),\,u_s-u'_s\Bigr\rangle
\,dt\Bigr|\\
&\le\sum_{s\in S^+}\operatorname{diam}(K_s)\Bigl[\delta_s+\frac{\mathfrak{L}_s\mathfrak{D}_s}{\lambda}\Bigr]
=\Xi .
\end{align*}
Thus the oscillation of $\log(m_1/m_0)$ on $\mathcal{W}_{S^+}$ is at most $\Xi$.

\emph{Normalisation.} The densities $m_0,m_1$ are continuous and strictly positive on
$\mathcal{W}_{S^+}$. They are continuous because they are differentiable, and strictly positive
because $e^{-H_a}>0$ and $\mathcal{W}_R$ is a non-empty open set. Both have integral $1$ over
$\mathcal{W}_{S^+}$. If $m_1>m_0$ held everywhere, then $\int m_1>\int m_0$; if $m_1<m_0$ held
everywhere, then $\int m_1<\int m_0$. Neither is possible, so
\begin{equation*}
\inf_{\mathcal{W}_{S^+}}\frac{m_1}{m_0}\le1\le\sup_{\mathcal{W}_{S^+}}\frac{m_1}{m_0}.
\end{equation*}
Fix $u\in\mathcal{W}_{S^+}$. For every $u'$ the oscillation bound gives
$\log\frac{m_1}{m_0}(u)\ge\log\frac{m_1}{m_0}(u')-\Xi$. Taking the supremum over $u'$ yields
$\log\frac{m_1}{m_0}(u)\ge-\Xi$. Likewise $\log\frac{m_1}{m_0}(u)\le\log\frac{m_1}{m_0}(u')+\Xi$
for every $u'$, and taking the infimum over $u'$ yields $\log\frac{m_1}{m_0}(u)\le\Xi$. This is
\eqref{8.3:eq-density-bound-4}.
\end{proof}

\begin{corollary}[Expectations and covariances under mixtures]\label{8.3:cor-mixture-covariances}
Let $(E,\mathcal{E})$ be a measurable space and let $\pi$ be a probability measure on it. For
each $\mathsf{e}\in E$ let $H_{\mathsf{e}}$ be such that the law $\nu_{H_{\mathsf{e}}}$ satisfies
Theorem~\ref{8.3:thm-density-bound} against the same $H_0$, with an exponent
$\Xi_{\mathsf{e}}\le\Xi$. This means that the pair $(H_0,H_{\mathsf{e}})$ in the roles of
$(H_0,H_1)$ satisfies the conclusion of that theorem with $\Xi_{\mathsf{e}}$ in place of $\Xi$.
Let
\begin{equation*}
\nu=\int\pi(d\mathsf{e})\,\nu_{H_{\mathsf{e}}}
\end{equation*}
be the corresponding mixture, and write $\nu_0:=\nu_{H_0}$. Let $m_0$ be the density of the
$S^+$-marginal of $\nu_0$ on $\mathcal{W}_{S^+}$, as in Theorem~\ref{8.3:thm-density-bound}. Then
the $S^+$-marginal of $\nu$ has a density $m$ with
\begin{equation}\label{8.3:eq-mixture-covariances-1}
e^{-\Xi}m_0\le m\le e^{\Xi}m_0 .
\end{equation}
For a bounded function $F$ on $\mathcal{W}$ put
\begin{equation*}
\|F\|_{\mathcal{W}}:=\sup_{y\in\mathcal{W}}|F(y)|,
\end{equation*}
and for a probability measure $\mu$ put
\begin{equation*}
\operatorname{Cov}_{\mu}(F,F'):=\mu(FF')-\mu(F)\mu(F').
\end{equation*}
Let $G,G'$ be bounded Borel functions of the restriction $y|_{S^+}$, and let $c,c'$ be real
numbers. Then
\begin{enumerate}
\item[(i)] the expectations satisfy
\begin{equation}\label{8.3:eq-mixture-covariances-2}
|\nu(G)-\nu_0(G)|\le(e^{\Xi}-1)\,\nu_0(|G-c|)\le(e^{\Xi}-1)\,\|G-c\|_{\mathcal{W}} ;
\end{equation}
\item[(ii)] the covariances satisfy
\begin{equation}\label{8.3:eq-mixture-covariances-3}
\bigl|\operatorname{Cov}_{\nu}(G,G')-\operatorname{Cov}_{\nu_0}(G,G')\bigr|
\le 3(e^{\Xi}-1)\,\|G-c\|_{\mathcal{W}}\,\|G'-c'\|_{\mathcal{W}} .
\end{equation}
\end{enumerate}
\end{corollary}

\begin{proof}
\emph{The density bound \eqref{8.3:eq-mixture-covariances-1}.} This holds because convex
combinations preserve the two-sided bound. In detail, let $m_{\mathsf{e}}$ be the density of the
$S^+$-marginal of $\nu_{H_{\mathsf{e}}}$. By hypothesis,
\begin{equation*}
e^{-\Xi_{\mathsf{e}}}m_0\le m_{\mathsf{e}}\le e^{\Xi_{\mathsf{e}}}m_0
\end{equation*}
on $\mathcal{W}_{S^+}$. Since $\Xi_{\mathsf{e}}\le\Xi$, this gives
$e^{-\Xi}m_0\le m_{\mathsf{e}}\le e^{\Xi}m_0$. Let $B\subset\mathcal{W}_{S^+}$ be a Borel set.
Integrating this inequality over $B$ against the reference measure of the densities gives
\begin{equation*}
e^{-\Xi}\nu_0^{+}(B)\le\nu_{H_{\mathsf{e}}}^{+}(B)\le e^{\Xi}\nu_0^{+}(B),
\end{equation*}
where the superscript $+$ denotes the $S^+$-marginal. The $S^+$-marginal of the mixture is
$\nu^{+}(B)=\int\pi(d\mathsf{e})\,\nu_{H_{\mathsf{e}}}^{+}(B)$. Integrating the last display
against the probability measure $\pi$ therefore gives
\begin{equation*}
e^{-\Xi}\nu_0^{+}(B)\le\nu^{+}(B)\le e^{\Xi}\nu_0^{+}(B)
\end{equation*}
for every Borel $B$. Hence $\nu^{+}$ is absolutely continuous with respect to $\nu_0^{+}$, and
so with respect to the reference measure. Its density $m$ satisfies
\eqref{8.3:eq-mixture-covariances-1} almost everywhere, and a version of $m$ can be chosen for
which \eqref{8.3:eq-mixture-covariances-1} holds everywhere.

\emph{Proof of \textup{(i)}.} Both $m$ and $m_0$ are probability densities, so
$\int c\,(m-m_0)=0$. Since $G$ depends only on $y|_{S^+}$,
\begin{equation*}
\nu(G)-\nu_0(G)=\int (G-c)\,(m-m_0)=\int (G-c)\Bigl(\frac{m}{m_0}-1\Bigr)m_0 ,
\end{equation*}
where the ratio is taken where $m_0>0$; by \eqref{8.3:eq-mixture-covariances-1}, $m=0$ where
$m_0=0$. By \eqref{8.3:eq-mixture-covariances-1},
\begin{equation*}
e^{-\Xi}-1\le\frac{m}{m_0}-1\le e^{\Xi}-1 .
\end{equation*}
Moreover $1-e^{-\Xi}\le e^{\Xi}-1$, because $e^{\Xi}+e^{-\Xi}\ge2$. Hence
$|m/m_0-1|\le e^{\Xi}-1$. Consequently
\begin{equation*}
|\nu(G)-\nu_0(G)|\le(e^{\Xi}-1)\int|G-c|\,m_0=(e^{\Xi}-1)\,\nu_0(|G-c|).
\end{equation*}
The second inequality in \eqref{8.3:eq-mixture-covariances-2} holds because $\nu_0$ is a
probability measure and $|G-c|\le\|G-c\|_{\mathcal{W}}$ pointwise.

\emph{Proof of \textup{(ii)}.} Put $X:=G-c$ and $X':=G'-c'$. A covariance is unchanged when
constants are subtracted from its arguments. Hence, for $\mu\in\{\nu,\nu_0\}$,
\begin{equation*}
\operatorname{Cov}_{\mu}(G,G')=\operatorname{Cov}_{\mu}(X,X')=\mu(X\,X')-\mu(X)\mu(X').
\end{equation*}
We estimate the two terms separately, writing $\|\cdot\|=\|\cdot\|_{\mathcal{W}}$.

\emph{The first term.} The product $X\,X'$ is a bounded Borel function of $y|_{S^+}$. Applying
\textup{(i)} to it with constant $0$ gives
\begin{equation*}
|\nu(X\,X')-\nu_0(X\,X')|\le(e^{\Xi}-1)\,\|X\,X'\|\le(e^{\Xi}-1)\,\|X\|\,\|X'\| .
\end{equation*}

\emph{The second term.} We have
\begin{align*}
|\nu(X)\nu(X')-\nu_0(X)\nu_0(X')|
&\le|\nu(X)-\nu_0(X)|\,|\nu(X')|+|\nu_0(X)|\,|\nu(X')-\nu_0(X')|\\
&\le 2(e^{\Xi}-1)\,\|X\|\,\|X'\| .
\end{align*}
Here \textup{(i)}, applied to $G$ with constant $c$ and to $G'$ with constant $c'$, gives
$|\nu(X)-\nu_0(X)|\le(e^{\Xi}-1)\|X\|$ and $|\nu(X')-\nu_0(X')|\le(e^{\Xi}-1)\|X'\|$. The
remaining factors are bounded by $|\nu(X')|\le\|X'\|$ and $|\nu_0(X)|\le\|X\|$, because $\nu$ and
$\nu_0$ are probability measures.

Adding the two estimates gives \eqref{8.3:eq-mixture-covariances-3}.
\end{proof}

\begin{corollary}[Choice of box exponent and removal of perturbations]
\label{8.3:cor-box-exponent-choice}
We use the data and the notation of Theorem~\ref{8.3:thm-density-bound}: the read set
$S^{+}\subset I$, its complement $R=I\setminus S^{+}$, the rate $\mu>0$, the window
$\mathcal{W}$, the numbers $\delta_b$ of \eqref{8.3:eq-lipschitz-stability-b}, and the
quantities $\mathfrak{L}_s$, $\mathfrak{D}_s$ and $\Xi$ formed there.
\begin{enumerate}
\item[(i)] Let $H_0$ and $H_0'$ both satisfy the hypotheses that
Theorem~\ref{8.3:thm-density-bound} places on $H_0$, with the same $A_0$, and suppose
that their perturbations $P_0$ and $P_0'$ satisfy
\begin{equation*}
|\partial_b(P_0-P_0')|\le\delta_b
\end{equation*}
in the sense of \eqref{8.3:eq-lipschitz-stability-b}. Let $m_0$ and $m_0'$ be the
densities of the $S^{+}$-marginals of $\nu_{H_0}$ and $\nu_{H_0'}$ on
$\mathcal{W}_{S^{+}}$. Then $e^{-\Xi}\le m_0'(u)/m_0(u)\le e^{\Xi}$ for every
$u\in\mathcal{W}_{S^{+}}$, with $\Xi$ formed from the data of $H_0$.
\item[(ii)] Theorem~\ref{8.3:thm-density-bound} holds with $P_0:=0$, in which case
$\varepsilon_0=0$, and with $\delta_b$ replaced, for every $b$, by
$\delta_b+\sup_{\mathcal{W}}|\partial_bP_1|$.
\end{enumerate}
\end{corollary}

\begin{proof}
Both parts are Theorem~\ref{8.3:thm-density-bound} applied to other data.

(i) Apply Theorem~\ref{8.3:thm-density-bound} with $H_0$ as the reference and
$H_1:=H_0'$. The hypotheses on the reference, coercivity with the constants
$(\lambda_A,\varepsilon_0)$ for every $\varphi$ with $\operatorname{Lip}(\varphi)\le\mu$,
hold because $H_0$ is assumed to satisfy them. Since $H_0$ and $H_0'$ have the same $A_0$,
they differ only through their perturbations. Hence the bound
\eqref{8.3:eq-lipschitz-stability-b}, which the theorem requires of the pair, is the
assumed bound $|\partial_b(P_0-P_0')|\le\delta_b$. The conclusion of the theorem, with
$m_1=m_0'$, is the asserted two-sided bound. Since $H_0'$ satisfies the same hypotheses,
the roles of the two may be exchanged. This gives the same bound with $\Xi$ formed from
the data of $H_0'$.

(ii) Apply Theorem~\ref{8.3:thm-density-bound} to the data in which the reference
perturbation is $P_0:=0$. For these data the coercivity hypothesis is read with
$\varepsilon_0=0$. The bound \eqref{8.3:eq-lipschitz-stability-b} is read with
$\delta_b+\sup_{\mathcal{W}}|\partial_bP_1|$ in place of $\delta_b$. The theorem then
gives its conclusion with these data. In the quantities entering $\Xi$, $\delta_s$ becomes
$\delta_s+\sup_{\mathcal{W}}|\partial_sP_1|$, and $\mathfrak{D}_s$ is formed with the
replaced $\delta_b$. Since $P_0''=0$, the supremum over $y\in\mathcal{W}$ in
$\mathfrak{L}_s$ is taken over a quantity independent of $y$, so that
\begin{equation*}
\mathfrak{L}_s=\Bigl(\sum_{b\in R}e^{2\mu d(b,s)}\,\|A_0(b,s)\|^2\Bigr)^{1/2},
\qquad s\in S^{+}.
\end{equation*}
\end{proof}

\begin{corollary}[Sub-Gaussian tails of linear functionals]\label{8.3:cor-subgaussian-tails}
Let $H$ be a Hamiltonian on $\mathcal{W}$ in the sense of
Proposition~\ref{8.3:prop-lipschitz-stability}. Suppose $H$ is coercive for $\varphi\equiv0$ in the
sense of \eqref{8.3:eq-lipschitz-stability-a}, with $\lambda>0$. Let $\ell\in E$ and put
$F:=\langle \ell,\cdot\rangle$. Then for every real $s$
\begin{equation}\label{8.3:eq-subgaussian-tails-1}
\nu_H\bigl(e^{s(F-\nu_H(F))}\bigr)\le e^{s^2|\ell|^2/2\lambda},
\end{equation}
and hence, for every $t>0$,
\begin{equation}\label{8.3:eq-subgaussian-tails-2}
\nu_H\bigl(|F-\nu_H(F)|\ge t\bigr)\le 2\,e^{-\lambda t^2/2|\ell|^2}.
\end{equation}
\end{corollary}

\begin{proof}
The set $\mathcal{W}$ is bounded, so $|F|\le|\ell|\sup_{y\in\mathcal{W}}|y|<\infty$ on
$\mathcal{W}$. For real $s$ put $\Lambda(s):=\log\nu_H(e^{sF})$.

We first compute $\Lambda'$. Since $e^{sF}e^{-H}=e^{-(H-sF)}$, we have
$\nu_H(e^{sF})=Z_{H-sF}/Z_H$. Because $F$ is bounded on $\mathcal{W}$ and $\nu_H$ is a
probability measure, we may differentiate under the integral. This gives
\begin{equation*}
\Lambda'(s)=\frac{\int_{\mathcal{W}}F\,e^{-(H-sF)}\,dy}{Z_{H-sF}}=\nu_{H-sF}(F).
\end{equation*}

Next we check that $H-sF$ is itself a Hamiltonian. It is obtained from $H$ by the replacement
$h\mapsto h-s\ell$ in the linear term, with the same $A$ and the same $P$. It is therefore a
Hamiltonian on $\mathcal{W}$. Since coercivity for $\varphi\equiv0$ involves only $A$ and $P$,
$H-sF$ satisfies \eqref{8.3:eq-lipschitz-stability-a} with the same constants, and hence with the
same $\lambda$.

We now apply Proposition~\ref{8.3:prop-lipschitz-stability} with $H_0:=H$, $H_1:=H-sF$,
$\varphi\equiv0$ and observable $G:=F$. The two Hamiltonians differ only in their linear terms, by
$s\ell$. Consequently the quantities of that proposition take the values $\delta_b=|s|\,|\ell_b|$
for $b\in I$ and $\ell_G=|\ell|$. The proposition then yields
\begin{equation*}
|\Lambda'(s)-\Lambda'(0)|=|\nu_{H-sF}(F)-\nu_H(F)|\le\frac{|s|\,|\ell|^2}{\lambda}.
\end{equation*}

Integrating this bound from $0$ to $s$ uses $\Lambda(0)=0$ and $\Lambda'(0)=\nu_H(F)$. We obtain
\begin{equation*}
\Lambda(s)-s\,\nu_H(F)=\int_0^s\bigl(\Lambda'(u)-\Lambda'(0)\bigr)\,du
\le\frac{|\ell|^2}{\lambda}\int_0^{|s|}u\,du=\frac{s^2|\ell|^2}{2\lambda},
\end{equation*}
for $s$ of either sign. The left-hand side equals $\log\nu_H\bigl(e^{s(F-\nu_H(F))}\bigr)$, which
proves \eqref{8.3:eq-subgaussian-tails-1}.

For the tail bound, let $t>0$ and $s>0$. Chebyshev's inequality applied to the positive function
$e^{s(F-\nu_H(F))}$, together with \eqref{8.3:eq-subgaussian-tails-1}, gives
\begin{equation*}
\nu_H\bigl(F-\nu_H(F)\ge t\bigr)\le e^{-st}\,\nu_H\bigl(e^{s(F-\nu_H(F))}\bigr)
\le e^{-st+s^2|\ell|^2/2\lambda}.
\end{equation*}
The choice $s=\lambda t/|\ell|^2$ optimises the exponent and yields the bound
$e^{-\lambda t^2/2|\ell|^2}$. The same argument with $-s$ in place of $s$ in
\eqref{8.3:eq-subgaussian-tails-1} bounds $\nu_H\bigl(F-\nu_H(F)\le -t\bigr)$ by the same
quantity. Adding the two bounds gives \eqref{8.3:eq-subgaussian-tails-2}.
\end{proof}

\begin{lemma}[Polymer format implies locality of the exterior force]\label{8.3:lem-exterior-force-locality}
Let $\mathfrak{B}_{\square}$ be a set of bonds of the lattice, called the box region, let
$\mathfrak{B}_{\square}^{c}$ be its complement, and let $I\subset\mathfrak{B}_{\square}$ be a set of bonds,
called the free bonds of $\mathfrak{B}_{\square}$. Assume that every member of an exterior datum
$\mathsf{e}$ other than its background
$\mathfrak{u}(\mathsf{e})$ is located in $\mathfrak{B}_{\square}^{c}$, and that for every $b\in I$ the
distances $d(b,O)$ and $d(b,\mathfrak{B}_{\square}^{c})$ differ by at most $1$, where $O$ is as in
Notation~\ref{8.3:not-step-box-window}.

Let $\mathfrak{Y}$ be a family of lattice domains with a size function $|Y|_{\star}\ge 0$, and let
$\ell_{\mathfrak{Y}}>0$, $\kappa_{\mathfrak{Y}}\ge 0$ and $C_{\mathfrak{Y}}$ be constants such that
\begin{itemize}
\item[\textup{(g1)}] $\operatorname{diam} Y\le \ell_{\mathfrak{Y}}\,(1+|Y|_{\star})$ for every
$Y\in\mathfrak{Y}$;
\item[\textup{(g2)}] $\sum_{Y\in\mathfrak{Y},\,Y\ni b}e^{-\kappa_{\mathfrak{Y}}|Y|_{\star}}\le C_{\mathfrak{Y}}$
for every bond $b$.
\end{itemize}
For every exterior datum $\mathsf{e}$ and every $y$ in the window of
Notation~\ref{8.3:not-step-box-window} put
\begin{equation*}
F_{\mathsf{e}}(y)=\sum_{Y\in\mathfrak{Y}}F(Y;y,\mathsf{e}).
\end{equation*}
Here $F(Y;y,\mathsf{e})$ depends on $y$ only through its values on $Y\cap I$. It depends on $\mathsf{e}$
only through the members of $\mathsf{e}$ located in $Y$ and the values of $\mathfrak{u}(\mathsf{e})$ on $Y$.
We assume that there are constants $v_1,v_2\ge 0$ such that, for every $Y\in\mathfrak{Y}$, all
exterior data $\mathsf{e},\mathsf{e}'$, every such $y$ and every $b\in Y\cap I$,
\begin{itemize}
\item[\textup{(v1)}] $|\partial_b F(Y;y,\mathsf{e})|\le v_1\,e^{-2\kappa_{\mathfrak{Y}}|Y|_{\star}}$;
\item[\textup{(v2)}] if $Y\subset\mathfrak{B}_{\square}$, then
\begin{equation*}
|\partial_b F(Y;y,\mathsf{e})-\partial_b F(Y;y,\mathsf{e}')|
\le v_2\,e^{-2\kappa_{\mathfrak{Y}}|Y|_{\star}}
\sup_{z\in Y}|\mathfrak{u}(\mathsf{e})-\mathfrak{u}(\mathsf{e}')|(z).
\end{equation*}
\end{itemize}
Suppose that the exterior data $\mathsf{e},\mathsf{e}'$ satisfy, for constants $\epsilon_0\ge 0$ and
$c_{\eta}\ge 0$,
\begin{equation}\label{8.3:eq-exterior-force-locality-1}
|\mathfrak{u}(\mathsf{e})-\mathfrak{u}(\mathsf{e}')|(z)\le \epsilon_0\,e^{-c_{\eta}d(z,\mathfrak{B}_{\square}^{c})}
\qquad\text{for } z\in\mathfrak{B}_{\square}.
\end{equation}
Put $c_{\partial}:=\min\{c_{\eta},\kappa_{\mathfrak{Y}}/\ell_{\mathfrak{Y}}\}$. Then for every $b\in I$ and
every $y$ in the window,
\begin{equation}\label{8.3:eq-exterior-force-locality-2}
|\partial_b(F_{\mathsf{e}}-F_{\mathsf{e}'})(y)|
\le (2v_1+v_2\epsilon_0)\,e^{\kappa_{\mathfrak{Y}}}C_{\mathfrak{Y}}\,
e^{-c_{\partial}d(b,\mathfrak{B}_{\square}^{c})}.
\end{equation}
\end{lemma}

\begin{proof}
Fix $b\in I$ and $y$ in the window, and write $D_b:=d(b,\mathfrak{B}_{\square}^{c})$.

If $b\notin Y$, then $F(Y;\cdot,\mathsf{e})$ and $F(Y;\cdot,\mathsf{e}')$ do not depend on the value of $y$ at
$b$, so $\partial_b F(Y;y,\cdot)=0$. Differentiating the sum defining $F_{\mathsf{e}}$ term by term therefore
gives
\begin{equation*}
\partial_b(F_{\mathsf{e}}-F_{\mathsf{e}'})(y)
=\sum_{Y\in\mathfrak{Y},\,Y\ni b}\bigl(\partial_b F(Y;y,\mathsf{e})-\partial_b F(Y;y,\mathsf{e}')\bigr).
\end{equation*}
We split the sum into the domains $Y\ni b$ with $Y\not\subset\mathfrak{B}_{\square}$ and the domains
$Y\ni b$ with $Y\subset\mathfrak{B}_{\square}$, and bound the two parts separately.

\emph{Domains not contained in the box.} Let $Y\ni b$ with $Y\not\subset\mathfrak{B}_{\square}$. Then $Y$
contains a point $z'$ of $\mathfrak{B}_{\square}^{c}$, and
\begin{equation*}
D_b\le d(b,z')\le \operatorname{diam}Y.
\end{equation*}
Combining this with (g1) gives $D_b\le\ell_{\mathfrak{Y}}(1+|Y|_{\star})$, that is,
$|Y|_{\star}\ge D_b/\ell_{\mathfrak{Y}}-1$. Hence
\begin{equation*}
e^{-2\kappa_{\mathfrak{Y}}|Y|_{\star}}=e^{-\kappa_{\mathfrak{Y}}|Y|_{\star}}e^{-\kappa_{\mathfrak{Y}}|Y|_{\star}}
\le e^{-\kappa_{\mathfrak{Y}}|Y|_{\star}}\,e^{\kappa_{\mathfrak{Y}}}\,
e^{-(\kappa_{\mathfrak{Y}}/\ell_{\mathfrak{Y}})D_b}.
\end{equation*}
Apply (v1) once to $\mathsf{e}$ and once to $\mathsf{e}'$, and use the triangle inequality. Each such $Y$
then contributes at most
\begin{equation*}
2v_1e^{-2\kappa_{\mathfrak{Y}}|Y|_{\star}}
\le 2v_1\,e^{\kappa_{\mathfrak{Y}}}\,e^{-\kappa_{\mathfrak{Y}}|Y|_{\star}}\,
e^{-(\kappa_{\mathfrak{Y}}/\ell_{\mathfrak{Y}})D_b}.
\end{equation*}
Summing over these $Y\ni b$ and using (g2), they contribute together at most
$2v_1e^{\kappa_{\mathfrak{Y}}}C_{\mathfrak{Y}}e^{-(\kappa_{\mathfrak{Y}}/\ell_{\mathfrak{Y}})D_b}$. Since
$c_{\partial}\le\kappa_{\mathfrak{Y}}/\ell_{\mathfrak{Y}}$, this is at most
\begin{equation*}
2v_1\,e^{\kappa_{\mathfrak{Y}}}C_{\mathfrak{Y}}\,e^{-c_{\partial}D_b}.
\end{equation*}

\emph{Domains contained in the box.} Let $Y\ni b$ with $Y\subset\mathfrak{B}_{\square}$, and let $z\in Y$.
Since $d(b,z)\le\operatorname{diam}Y$, the triangle inequality gives
$d(z,\mathfrak{B}_{\square}^{c})\ge D_b-\operatorname{diam}Y$. We bound the background difference in
three steps:
\begin{itemize}
\item[(a)] by \eqref{8.3:eq-exterior-force-locality-1}, which applies because $z\in\mathfrak{B}_{\square}$,
together with $c_{\partial}\le c_{\eta}$, we have
\begin{equation*}
|\mathfrak{u}(\mathsf{e})-\mathfrak{u}(\mathsf{e}')|(z)\le\epsilon_0e^{-c_{\partial}d(z,\mathfrak{B}_{\square}^{c})};
\end{equation*}
\item[(b)] since $t\mapsto e^{-c_{\partial}t}$ is non-increasing, the lower bound on
$d(z,\mathfrak{B}_{\square}^{c})$ and (g1) give
\begin{equation*}
e^{-c_{\partial}d(z,\mathfrak{B}_{\square}^{c})}
\le e^{-c_{\partial}D_b}\,e^{c_{\partial}\ell_{\mathfrak{Y}}(1+|Y|_{\star})};
\end{equation*}
\item[(c)] since $c_{\partial}\ell_{\mathfrak{Y}}\le\kappa_{\mathfrak{Y}}$, we have
$e^{c_{\partial}\ell_{\mathfrak{Y}}(1+|Y|_{\star})}\le e^{\kappa_{\mathfrak{Y}}(1+|Y|_{\star})}$.
\end{itemize}
Taking the supremum over $z\in Y$,
\begin{equation*}
\sup_{z\in Y}|\mathfrak{u}(\mathsf{e})-\mathfrak{u}(\mathsf{e}')|(z)
\le\epsilon_0\,e^{-c_{\partial}D_b}\,e^{\kappa_{\mathfrak{Y}}(1+|Y|_{\star})}.
\end{equation*}
Inserting this into (v2), each such $Y$ contributes at most
\begin{equation*}
v_2\,e^{-2\kappa_{\mathfrak{Y}}|Y|_{\star}}\,\epsilon_0\,e^{-c_{\partial}D_b}\,
e^{\kappa_{\mathfrak{Y}}(1+|Y|_{\star})}
=v_2\epsilon_0\,e^{\kappa_{\mathfrak{Y}}}\,e^{-\kappa_{\mathfrak{Y}}|Y|_{\star}}\,e^{-c_{\partial}D_b}.
\end{equation*}
Summing over these $Y\ni b$ and using (g2), they contribute together at most
\begin{equation*}
v_2\epsilon_0\,e^{\kappa_{\mathfrak{Y}}}C_{\mathfrak{Y}}\,e^{-c_{\partial}D_b}.
\end{equation*}

Adding the two contributions gives \eqref{8.3:eq-exterior-force-locality-2}.
\end{proof}

The lemma is applied with two choices of $F$ and of the family of domains. In each of them,
(g1), (g2), (v1), (v2) and \eqref{8.3:eq-exterior-force-locality-1}, with their constants, enter as
hypotheses of the application.

In the first choice, $F_{\mathsf{e}}$ is the sum $P^{\mathrm{tot}}_j$, over domains $Y$, of the localised
terms. When the exterior datum consists of a single background, these are the terms of the
representation in \cite[Lemma~2, (1.36), (1.41)--(1.43)]{B-CMP116}; in the general case they are the
terms of the $\mathbf{E}+\mathbf{R}+\mathbf{B}$ representation of \cite[Theorem~2]{B-CMP119} with the
bounds \cite[(2.27)--(2.35)]{B-CMP119}. The family $\mathfrak{Y}$
is the family $\mathbf{D}_j$ of domains carrying these terms, $|Y|_{\star}=d_j(Y)$, and $\ell_{\mathfrak{Y}}$
is comparable to $M$. Condition (g2) is then \cite[(1.26)]{B-CMP116}. Conditions (v1) and (v2) are the
Cauchy bounds for these terms, in the gauge-field variable on radii exceeding the window,
respectively in the variables of the representation of \cite[Theorem~2]{B-CMP119} on radii
exceeding the variation of the backgrounds.

In the second choice,
\begin{equation*}
F_{\mathsf{e}}(y)=\tfrac12\langle y,\mathsf{Q}^{\mathsf{e}}_{II}y\rangle.
\end{equation*}
Here $\mathsf{Q}=\sum_{\nu}T_{\nu}$ is the decomposition provided by clauses (D), (E) and (H) of the
operator theorem at free current, and the family $\mathfrak{Y}$ consists of the domains
$Y=\operatorname{dom}\nu$ of the terms $T_{\nu}$. In
this choice $v_1$ and $v_2$ are proportional to the window parameter $\hat{\mathfrak{b}}_j$ of
Notation~\ref{8.3:not-step-box-window}.

In both choices, hypothesis \eqref{8.3:eq-exterior-force-locality-1} on the backgrounds is the
localisation of the background difference. In Ba{\l}aban's construction, the bounds
\cite[(3.17)--(3.19)]{B-CMP119} are an instance of it. Consequently, in the general case, the
exterior-locality bounds for $b\in I$ (the bounds on the differences, between an exterior datum and a
reference exterior datum, of the force and of the $I$-block of the quadratic form, each by a constant
times $e^{-c_{\partial}d(b,O)}$) follow from the lemma, provided the two formats cited above and
clauses (D), (E), (H) and (A) of the operator theorem at free current supply (g1), (g2), (v1) and
(v2), and provided the localisation of the background difference supplies
\eqref{8.3:eq-exterior-force-locality-1}, with the background flow statement as a further
hypothesis; all these provisions are taken here as hypotheses. The distance
$d(b,\mathfrak{B}_{\square}^{c})$ in
\eqref{8.3:eq-exterior-force-locality-2} is replaced by $d(b,O)$ at the cost of a factor
$e^{c_{\partial}}$, since for $b\in I$ the two distances differ by at most $1$ and $c_{\partial}\ge 0$.

\begin{lemma}[Hypotheses of box localisation and their verification]
\label{8.3:lem-localisation-hypotheses}
Let the step be the one of Notation~\ref{8.3:not-step-box-window} at level $j=k$, where
$j=k_0+i$ with $0\le i\le r$, and $k_0$, $r$ and the box $\mathfrak{B}^{(i)}_D$ are as in
Definition~\ref{8.3:def-clean-event}; let $h$ be a
history, and let $V_{j+1}$ be fixed. Let $I\subset\mathbb{B}^{+}$ be a box of free bonds and
$O:=\mathbb{B}^{+}\setminus I$. Let the inserted observable $G$ read $V_j$ on a bond set
$\mathfrak{S}_G$, and let $S^{+}\subset I$ be its read set
(Lemma~\ref{8.3:lem-read-set}); put $D_\star:=d(S^{+},O)$ and $n_{S^{+}}:=\#S^{+}$. Write $|\cdot|$
for the norm on $\mathfrak{g}$ defining the window, $|\cdot|_2$ for the Euclidean norm in the
coordinates $X^a$, $a=1,\dots,\dim\mathfrak{g}$, and $\|\cdot\|$ for the operator norm with respect
to $|\cdot|_2$; write $\|\mathbf{C}\|$ for the $\ell^\infty\to\ell^\infty$ norm, with respect to
$|\cdot|_2$, of Ba{\l}aban's block-local operator $\mathbf{C}$ of
Notation~\ref{8.3:not-step-box-window}. Put
\begin{equation*}
\hat{\mathfrak{b}}_j:=\sup\{|X|_2:\ X\in\mathfrak{g},\ |X|<\mathfrak{b}_j\}\le c_N\mathfrak{b}_j ,
\end{equation*}
where $c_N$ is the norm-equivalence constant of \cite[(20)]{B-CMP98}, a number depending on $N$
only. For a function $F$ of $y=(y_b)_{b\in I}$, $\partial_bF$ is its gradient in $y_b$ and
$\partial_b\partial_{b'}F$ is the $\mathfrak{g}$-block of its Hessian at $(b,b')$.

The \emph{hypotheses of box localisation} at the step are the following six conditions.
\begin{itemize}
\item[(h1)] (Placement.) For the observable $G$ and for $\tau_h$-almost every $V_{j+1}$,
\begin{equation*}
\mathbb{E}_{i,h}[G](V_{j+1})
=\int_{\mathsf{E}}\pi_{h,V_{j+1}}(d\mathsf{e})\,E_{\mathsf{e}}[G\circ\mathcal{V}_{\mathsf{e}}],
\end{equation*}
where $\pi_{h,V_{j+1}}$ is a probability measure on the set $\mathsf{E}$ of exterior data and
$E_{\mathsf{e}}$ is the probability measure on $\mathcal{W}_I$ with density with respect to $dy$
proportional to $\mathbf{1}_{\mathcal{W}_I}(y)e^{-H_{\mathsf{e}}(y)}$,
\begin{equation*}
H_{\mathsf{e}}(y)=\tfrac12\langle y,\mathsf{Q}^{\mathsf{e}}_{II}y\rangle
+\langle y,\mathsf{J}^{\mathsf{e}}A_O(\mathsf{e})\rangle-P_{\mathsf{e}}(y),
\qquad \mathsf{J}^{\mathsf{e}}:=\mathsf{Q}^{\mathsf{e}}_{IO}.
\end{equation*}
\item[(h2)] (Window.) $\mathcal{W}_I=\prod_{b\in I}\{X\in\mathfrak{g}:|X|<\mathfrak{b}_j\}$, a
product of norm balls of $\mathfrak{g}$.
\item[(h3)] (Form.) For every admissible exterior datum $\mathsf{e}$, $\mathsf{Q}^{\mathsf{e}}$ is
real symmetric on $\mathfrak{g}^{\mathbb{B}^{+}}$ and
$\|\mathsf{Q}^{\mathsf{e}}(b,b')\|\le a_+e^{-c_A d(b,b')}$; for the reference datum
$\mathsf{e}_0$ of \textup{(h5)} moreover $\mathsf{Q}^{\mathsf{e}_0}\ge a_->0$; the constants
$a_-,a_+,c_A$ do not depend on the level.
\item[(h4)] (Exterior field.) $|A_O(\mathsf{e})(o)|_2\le C_O\hat{\mathfrak{b}}_j$ for every
$o\in O$ and every admissible $\mathsf{e}$.
\item[(h5)] (Reference exponent.) There are a reference datum $\mathsf{e}_0$ and
$P_0\in C^2(\mathcal{W}_I)$ with
\begin{equation*}
\sup_{y\in\mathcal{W}_I}|\partial_bP_0(y)|_2<\infty,\qquad
\|\partial_b\partial_{b'}P_0(y)\|\le\varepsilon_P\,e^{-c_Pd(b,b')}
\quad(b,b'\in I,\ y\in\mathcal{W}_I),
\end{equation*}
under the one-sided condition
\begin{equation}\label{8.3:eq-localisation-hypotheses-1}
\varepsilon_P\,S(c_P/2)\le a_-/4 ,
\end{equation}
$S(\cdot)$ being the lattice sum of Lemma~\ref{8.3:lem-weighted-inverse-decay}.
\item[(h6)] (Exterior locality of the force.) There are $\eta_{\mathsf{Q}},\varepsilon_\partial\ge0$
and $c_\partial>0$ such that for every admissible $\mathsf{e}$, every $b\in I$ and every
$y\in\mathcal{W}_I$
\begin{equation}\label{8.3:eq-localisation-hypotheses-2}
\begin{aligned}
|((\mathsf{Q}^{\mathsf{e}}_{II}-\mathsf{Q}^{\mathsf{e}_0}_{II})y)_b|_2
&\le\eta_{\mathsf{Q}}\,\hat{\mathfrak{b}}_j\,e^{-c_\partial d(b,O)},\\
|\partial_b(P_{\mathsf{e}}-P_0)(y)|_2&\le\varepsilon_\partial\,e^{-c_\partial d(b,O)}.
\end{aligned}
\end{equation}
\end{itemize}

Assume the following.
\begin{itemize}
\item[(a)] The form hypothesis and the positivity hypothesis of
Lemma~\ref{8.3:lem-box-law-mixture} hold at the step for the box $I$: near $I$ the history's step
is a term of the fluctuation integral \eqref{8.3:eq-fluctuation-integral-a} with
$I\subset\Lambda^{(j)*}_{j+1}$, and the operations are of the kinds admitted there, every weight in
$\mathbf{T}_{j+1}$ being nonnegative. The observable $G$ is a bounded Borel function of $V_j$
restricted to $\mathfrak{S}_G$, with $S^{+}\subset I$, gauge invariant (any bounded Borel $G$ in
the small-field setting of \cite{B-CMP109}).
\item[(b)] The window of the step is the product of norm balls of
Definition~\ref{8.3:def-window-product}, display \eqref{8.3:eq-window-product-a}, at $k=j$.
\item[(c)] Clauses (B), (D) and (I) of the operator theorem at free current hold in the following
form: for every background $\mathfrak{u}$ in the real domain $\mathcal{U}_{\mathbb{R}}$ of that
theorem, $\mathsf{Q}[\mathfrak{u}]$ is real symmetric on $\mathfrak{g}^{\mathbb{B}^{+}}$ with
$\|\mathsf{Q}[\mathfrak{u}](b,b')\|\le a_+e^{-c_Ad(b,b')}$, with $a_+,c_A$ independent of the level.
For the background $\mathfrak{u}(\mathsf{e}_0)$ of the reference datum, in addition,
$\mathsf{Q}[\mathfrak{u}(\mathsf{e}_0)]\ge a_->0$, with $a_-$ independent of the level.
Moreover the background $\mathfrak{u}(\mathsf{e})$ of every admissible exterior datum
$\mathsf{e}$, and the background $\mathfrak{u}(\mathsf{e}_0)$ of the reference datum, lie in
$\mathcal{U}_{\mathbb{R}}$.
\item[(d)] The exterior field on $\Omega_{j+1}\setminus\Lambda_{j+1}$ obeys
\cite[(3.13)--(3.14)]{B-CMP119}: $|A_k|<O(1)\varepsilon_k$, in Ba{\l}aban's unscaled variable $A_k$
and with his small-field parameter $\varepsilon_k$ of \cite[(3.14)]{B-CMP119}; that is, in the
scaled variable,
$|A_O(\mathsf{e})(o)|\le C(A_0/A_1)\mathfrak{b}_j$ with $C$ a numerical constant. Let $C_O\ge1$ be
the constant, depending on $A_0/A_1$ and $N$ only, for which this bound gives
$|A_O(\mathsf{e})(o)|_2\le C_O\hat{\mathfrak{b}}_j$.
\item[(e)] (Analytic locality of the total potential on the tube.) Write
$P^{\mathrm{tot}}_j=P^{\mathrm{sf}}_j+P^{\mathbf{E}}_j+P^{\mathbf{R}}_j+P^{\mathrm{pk}}_j$ for the
decomposition of the total potential into its small-field member, its $\mathbf{E}$-member, its
$\mathbf{R}$-term member and its member attached to large-field regions of earlier levels
(its terms, including the frozen-coupling ones, as in Definition~\ref{8.3:def-clean-event}). Write
$B=(B(b))_{b\in\mathbb{B}^{+}}$ for the field whose restriction to $I$ is $y$ and whose restriction
to $O$ is $A_O(\mathsf{e})$, and $B^a(b)$ for its components. For every
exterior datum $\mathsf{e}$ whose background lies in $\mathcal{U}_{\mathbb{R}}$ (in particular, by
(c), for every admissible $\mathsf{e}$ and for $\mathsf{e}_0$),
$P^{\mathrm{tot}}_j(\,\cdot\,;\mathsf{e})$ is analytic on a complex tube about the closed
real window, defined by bounds on $(\mathbf{C}B)(b)$ with real half-width
$(1\vee C_O)c_N\|\mathbf{C}\|\mathfrak{b}_j+1$ in $|\cdot|_2$ and imaginary half-width a fixed
radius depending on $(L,\mathfrak{p},N)$ only, so that its real part contains, with a
real margin, the closed real window $\prod_{b\in I}\{|y_b|\le\mathfrak{b}_j\}$ together with the
exterior range of (d); and on it the sum $P^{\mathrm{sf}}_j+P^{\mathbf{E}}_j$ obeys, for
$a,a'=1,\dots,\dim\mathfrak{g}$ and bonds $b,b'$ of $\mathbb{B}^{+}$ at distance $\ell$ (box and
exterior bonds alike, so that the mixed derivatives in box and exterior variables are included),
\begin{equation*}
\Bigl|\frac{\partial(P^{\mathrm{sf}}_j+P^{\mathbf{E}}_j)}{\partial B^a(b)}\Bigr|
\le C_Pg_j\mathfrak{b}_j^2,
\qquad
\Bigl|\frac{\partial^2(P^{\mathrm{sf}}_j+P^{\mathbf{E}}_j)}{\partial B^a(b)\,\partial B^{a'}(b')}\Bigr|
\le C_Pg_j\mathfrak{b}_j\,e^{-c_P\ell},
\end{equation*}
with $C_P,c_P$ depending on $(L,\mathfrak{p},N)$ only.
\item[(f)] This item concerns the members $P^{\mathrm{pk}}_j$ and $P^{\mathbf{R}}_j$ of the
decomposition in (e). The box $I$ consists of free bonds of $\mathfrak{B}^{(i)}_D$, so that the
clause of \eqref{8.3:eq-clean-event-a} applies at every $b\in I$; the history $h$ lies in the
event of \eqref{8.3:eq-clean-event-a}; and $x_{k_0}$ satisfies the conditions
\eqref{8.3:eq-depth-coupling-conditions-a}. For every term $X$ of $P^{\mathrm{pk}}_j$ of level
$i'$, with $\nu:=j-i'$, and with $d(X)$, $\kappa_{\mathrm{ax}}(X)$, $M_X$, $\mathfrak{l}$ and $c'$
as in \eqref{8.3:eq-clean-event-a}, the second-order size $\varepsilon_X$ of $X$ in the
fluctuation variables at the bonds of $\mathbb{B}^{+}$, on the
tube of (e), in the sense of the second-order clause of the equivalence lemmas, satisfies
\begin{equation*}
\varepsilon_X\le C\,M_X\,\kappa_{\mathrm{ax}}(X)^2L^{-4\nu}(1+B^\sharp\nu)\,\mathfrak{b}_j^2 .
\end{equation*}
This bound holds at the ratio bound $(g_j/R^{\mathrm{an}}_{i'})^2\le C(1+B^\sharp\nu)$,
$R^{\mathrm{an}}_{i'}$ denoting the analyticity radius of the level-$i'$ term, which follows from
the upper bound of the coupling flow and \cite[(2.28)]{B-CMP119}. On the event of
\eqref{8.3:eq-clean-event-a}, the clause of the decomposition lemma of the Gaussian extraction on
the contribution of these terms, run with the clause of \eqref{8.3:eq-clean-event-a} summed over
the bonds of $I$ and with the same ratio bound, gives for every $b\in I$
\begin{equation*}
|\partial_bP^{\mathrm{pk}}_j|_2
\le C_{\mathrm{grad}}\|\mathbf{C}\|\,\hat{\mathfrak{b}}_j\bigl(\mathfrak{b}_j^2+p_0(g_j)\bigr)
e^{-\frac12c'p_0(g_{k_0})},
\end{equation*}
$C_{\mathrm{grad}}$ being the constant of the second inequality defining $x_{\mathrm{pk}}$ in
Definition~\ref{8.3:def-depth-coupling-conditions}. For a
function $F$ of the field $B$ of (e) we call
\begin{equation*}
\sup_{b\in I}\ \sup_{B}\ \sum_{b'\in\mathbb{B}^{+}}
\|\partial_b\partial_{b'}F(B)\|\,e^{c_Pd(b,b')}
\end{equation*}
the \emph{weighted Hessian norm} of $F$; here $\partial_{b'}$ denotes for $b'\in O$ the gradient
in $B(b')$, and the second supremum runs over the real fields $B$ whose restriction to $I$ lies
in the closure $\overline{\mathcal{W}_I}$ and which satisfy $|B(o)|_2\le C_O\hat{\mathfrak{b}}_j$
for every $o\in O$, a set contained, with a real margin, in the tube of (e). On every history
the weighted Hessian norm of $P^{\mathbf{R}}_j$ is at most
$Cg_jx_j^{(3-\kappa_0)/2}\,\mathrm{polylog}(x_j)/\mathfrak{b}_j$, $\mathrm{polylog}(x_j)$ denoting a
polynomial in $\log x_j$; this is the bound on the $\mathbf{R}$-term member in the decomposition
lemma of the Gaussian extraction, which rests on \cite[(2.31)]{B-CMP119} with Ba{\l}aban's
exponent $\kappa_0\ge4$ of that display, combined with Cauchy's estimate.
\item[(g)] For the members of the exterior datum that move the background (fields on large-field
zones), the dependence on the exterior variable at a bond $b''$ of each of the following three
quantities is at most $C_{\mathrm{op}}e^{-c_{\mathrm{op}}d(b'',I)}$: Ba{\l}aban's operator
$C^{(j)}$ of the step restricted to the box, the effective potential restricted to the box, and
$U_{j+1}$ under the box; this bound holds with the constants of the operator
theorem at free current, analytically and with these
Cauchy bounds on the complex window (clauses (D), (E), (H) and (A) and the geometric clause of the
operator theorem at free current); the box-scale localisation bound among the inputs of the
Gaussian extraction holds;
the terms generated by $\mathbf{B}$ and the effective potential on the large-field region of the
step are localised close to the large-field regions in the form
\cite[(2.34)--(2.35)]{B-CMP119}, with the factor $e^{-\kappa d(Y)}$; and the resulting families
of terms, with their localisation domains and size functions, satisfy the two geometric
conditions (g1), (g2) and the amplitude conditions of
Lemma~\ref{8.3:lem-exterior-force-locality}, with the amplitude parameters $v_1,v_2$ of that lemma
at most a constant times $\hat{\mathfrak{b}}_j$ and its parameter $\epsilon_0$ at most a constant
times Ba{\l}aban's small-field parameter $\varepsilon_j$.
\end{itemize}
Then \textup{(h1)}--\textup{(h4)} hold. \textup{(h5)} holds with the reference datum
$\mathsf{e}_0$ obtained from any admissible exterior datum by setting all fluctuation variables on
$O$ to $0$, with
\begin{equation*}
P_0(y):=P^{\mathrm{tot}}_j(y;\mathsf{e}_0)-P^{\mathrm{tot}}_j(0;\mathsf{e}_0),\qquad
\varepsilon_P:=(\dim\mathfrak{g})\,C_Pg_j\mathfrak{b}_j ,
\end{equation*}
$C_P$ being replaced by a larger constant, still written $C_P$, so chosen that the Hessian bound of
\textup{(e)} covers the sum of the three members $P^{\mathrm{sf}}_j+P^{\mathbf{E}}_j$,
$P^{\mathrm{pk}}_j$ and $P^{\mathbf{R}}_j$,
whenever \eqref{8.3:eq-localisation-hypotheses-1} holds; with this
enlarged $C_P$, the Hessian bound of \textup{(e)} holds for the whole of $P^{\mathrm{tot}}_j$, for
$b\in I$ and $b'\in\mathbb{B}^{+}$ and at every real field $B$ of the set over which the weighted
Hessian norm of \textup{(f)} is taken, on the
event of \eqref{8.3:eq-clean-event-a}. On that event, moreover,
$|\partial_bP^{\mathrm{pk}}_j|_2\le C_Pg_j\mathfrak{b}_j^2$ for every $b\in I$.
\textup{(h6)} holds: in the small-field setting of \cite{B-CMP109}, by the Hessian bound of
\textup{(e)} for the whole of $P^{\mathrm{tot}}_j$, with
\begin{equation*}
\eta_{\mathsf{Q}}=0,\qquad c_\partial=c_P/2,\qquad
\varepsilon_\partial=(\dim\mathfrak{g})\,C_Pg_j\mathfrak{b}_j\,\hat{\mathfrak{b}}_j\,S(c_P/2);
\end{equation*}
and in general with the constants furnished by Lemma~\ref{8.3:lem-exterior-force-locality} under
\textup{(g)}, which need not be small.
\end{lemma}

\begin{proof}
We verify the six hypotheses in turn.

\emph{Placement \textup{(h1)}.} By (a), the hypotheses of Lemma~\ref{8.3:lem-box-law-mixture} hold
for the box $I$ and the observable $G$. Its conclusion is the mixture formula of (h1) for
$\tau_h$-almost every $V_{j+1}$, with $\pi_{h,V_{j+1}}$ a probability measure on the exterior data
and $E_{\mathsf{e}}$ the box law of Notation~\ref{8.3:not-step-box-window}, whose density is
proportional to $\mathbf{1}_{\mathcal{W}_I}e^{-H_{\mathsf{e}}}$ with $H_{\mathsf{e}}$ as displayed;
the term $\langle y,\mathsf{J}^{\mathsf{e}}A_O(\mathsf{e})\rangle$ is the action of the frozen
exterior field on the box through the boundary operator $\mathsf{J}^{\mathsf{e}}=\mathsf{Q}^{\mathsf{e}}_{IO}$
(Definition~\ref{8.3:def-box-precision}). The form hypothesis refers to the fluctuation integral
\eqref{8.3:eq-fluctuation-integral-a}. For histories $h$ of
the clean event $\mathcal{C}_S(D)$ of Definition~\ref{8.3:def-clean-event}, the box
$\mathfrak{B}^{(i)}_D$ lies in $\Lambda_{j+1}$; this is the statement that the form hypothesis
makes for boxes $I$ of free bonds of $\mathfrak{B}^{(i)}_D$.

\emph{Window \textup{(h2)}.} By (b), the characteristic function of the small-field window on the
free bonds of $\Lambda_{j+1}$ is the product over these bonds of the indicator of the ball
$\{X\in\mathfrak{g}:|X|<\mathfrak{b}_j\}$, display \eqref{8.3:eq-window-product-a}. The product shape
is Ba{\l}aban's definition \cite[(2.9)]{B-CMP109}, \cite[(3.16)]{B-CMP119}, and, on
$\Lambda_{j+1}$, only these characteristic functions enter \cite[(3.20)]{B-CMP119}; the radius
$\mathfrak{b}_j$ in the scaled variable is the one fixed in
Definition~\ref{8.3:def-window-product}. Restricting the product to the bonds of $I$ gives
$\mathcal{W}_I$ as in (h2).

\emph{Form \textup{(h3)}.} Let $\mathsf{e}$ be an admissible exterior datum. By the second part of
(c), its background $\mathfrak{u}(\mathsf{e})$ lies in $\mathcal{U}_{\mathbb{R}}$, so the first
part of (c) applies to $\mathsf{Q}^{\mathsf{e}}=\mathsf{Q}[\mathfrak{u}(\mathsf{e})]$: it is real
symmetric on $\mathfrak{g}^{\mathbb{B}^{+}}$ with $\|\mathsf{Q}^{\mathsf{e}}(b,b')\|\le
a_+e^{-c_Ad(b,b')}$. The same applies to the reference datum, whose background also lies in
$\mathcal{U}_{\mathbb{R}}$, and (c) gives in addition $\mathsf{Q}^{\mathsf{e}_0}\ge a_->0$. The
constants are those of (c) and do not depend on the level.

\emph{Exterior field \textup{(h4)}.} Let $o\in O$. If $o$ is a free bond of
$\Lambda_{j+1}\setminus I$, then $A_O(\mathsf{e})(o)$ is an integration variable of
\eqref{8.3:eq-fluctuation-integral-a} restricted by the window of (h2) extended to all free bonds
of $\Lambda_{j+1}$; hence $|A_O(\mathsf{e})(o)|<\mathfrak{b}_j$ and, by the definition of
$\hat{\mathfrak{b}}_j$, $|A_O(\mathsf{e})(o)|_2\le\hat{\mathfrak{b}}_j$, which is (h4) with constant
$1$. If $o$ is a bond of $\Omega_{j+1}\setminus\Lambda_{j+1}$, then (d) gives
$|A_O(\mathsf{e})(o)|\le C(A_0/A_1)\mathfrak{b}_j$; since $X\mapsto|X|_2$ and $X\mapsto|X|$ are
homogeneous, a
vector $X$ with $|X|<c\,\mathfrak{b}_j$ satisfies $|X|_2\le c\,\hat{\mathfrak{b}}_j$, and the
constant $C_O\ge1$ of (d) gives $|A_O(\mathsf{e})(o)|_2\le C_O\hat{\mathfrak{b}}_j$. As $C_O\ge1$,
(h4) holds on all of $O$. By the choice of the half-width in (e), the exterior range so obtained
lies in the real part of the tube of (e); this is used for (h6).

\emph{Reference exponent \textup{(h5)}.} Take $\mathsf{e}_0$ and $P_0$ as in the statement. Since
the closure of $\mathcal{W}_I$ lies, with a real margin, in the tube of (e), $P_0$ is analytic on
a neighbourhood of that closure, in particular $P_0\in C^2(\mathcal{W}_I)$, and its derivatives in
$y$ are those of
$P^{\mathrm{tot}}_j(\,\cdot\,;\mathsf{e}_0)$, the subtracted constant not depending on $y$. The
closure of $\mathcal{W}_I$, a finite product of closed balls of $\mathfrak{g}$, is compact, and
the gradient of $P_0$ is continuous on it; hence
$\sup_{\mathcal{W}_I}|\partial_bP_0|_2<\infty$. For the Hessian block of the members
$P^{\mathrm{sf}}_j+P^{\mathbf{E}}_j$, each entry is at most
$C_Pg_j\mathfrak{b}_je^{-c_Pd(b,b')}$ by (e), and a $\dim\mathfrak{g}\times\dim\mathfrak{g}$ matrix
with entries bounded by $\beta$ has operator norm at most $(\dim\mathfrak{g})\beta$; hence
\begin{equation*}
\|\partial_b\partial_{b'}(P^{\mathrm{sf}}_j+P^{\mathbf{E}}_j)(y;\mathsf{e}_0)\|
\le(\dim\mathfrak{g})\,C_Pg_j\mathfrak{b}_j\,e^{-c_Pd(b,b')}
=\varepsilon_P\,e^{-c_Pd(b,b')}\qquad(y\in\mathcal{W}_I).
\end{equation*}

The member $P^{\mathrm{pk}}_j$ attached to large-field regions of earlier levels and the
$\mathbf{R}$-term
member $P^{\mathbf{R}}_j$ are placed inside $\varepsilon_P$ as follows. For $P^{\mathrm{pk}}_j$
only the weighted Hessian norm of (f) is used, and it is estimated on the event
\eqref{8.3:eq-clean-event-a}, which is contained in $\mathcal{C}_S(D)$. Fix $b\in I$. A bond $b'$
contributes to $\partial_b\partial_{b'}P^{\mathrm{pk}}_j$ only through terms $X$ of some level
$i'<k_0-\mathfrak{l}$ whose box $\tilde X$ lies under the dependence range of $b$ and whose
dependence also reaches the fluctuation variable at $b'$; these are the terms over which the sum
in \eqref{8.3:eq-clean-event-a} runs, the terms of $P^{\mathrm{pk}}_j$ being those of
Definition~\ref{8.3:def-clean-event}. The side of $\tilde X$ is $d(X)$ in
level-$i'$ units, so its diameter in level-$j$ units is $d(X)L^{-\nu}$, $\nu=j-i'$. Every such $b'$
therefore lies within $d(X)L^{-\nu}+C$ level-$j$ units of $b$: the dependence range on both sides
contributes the constant $C$, the box its level-$j$ diameter. Consequently, for each term $X$, the
number of bonds $b'$ it couples to $b$ is at most a constant depending on the dimension only times
$(d(X)L^{-\nu}+C)^4$, that constant being absorbed into the constant $C_\delta$ below, and each
carries the weight
$e^{c_Pd(b,b')}\le e^{c_P(d(X)L^{-\nu}+C)}$. The second-order bound of (f) for each term, together
with Cauchy's estimate on the margin of the tube of (e), bounds the corresponding Hessian blocks.
Summing over $b'$ with the weight, then over $X$ and over $i'<k_0-\mathfrak{l}$, produces
$C_\delta\|\mathbf{C}\|^2\mathfrak{b}_j^2$ times the left-hand side of \eqref{8.3:eq-clean-event-a},
with the constant $C_\delta$ of the first inequality defining $x_{\mathrm{pk}}$ in
Definition~\ref{8.3:def-depth-coupling-conditions}, which depends on the
margin of the tube of (e).
On that event this left-hand side is at most $e^{-\frac12c'p_0(g_{k_0})}$, and therefore
\begin{equation*}
\sum_{b'}\|\partial_b\partial_{b'}P^{\mathrm{pk}}_j\|\,e^{c_Pd(b,b')}
\le C_\delta\|\mathbf{C}\|^2\mathfrak{b}_j^2\,e^{-\frac12c'p_0(g_{k_0})}\le C_Pg_j\mathfrak{b}_j ,
\end{equation*}
the last inequality being the first of the inequalities defining $x_{\mathrm{pk}}$ in
Definition~\ref{8.3:def-depth-coupling-conditions},
valid because $x_{k_0}\ge x_{\mathrm{pk}}$. The factor $\mathfrak{b}_j^2$ of the second-order bound
is kept in this summed bound. If instead Cauchy's estimate is taken with a margin proportional to
$\hat{\mathfrak{b}}_j$, the factor $\mathfrak{b}_j^2$ cancels against the square of that margin,
$C_\delta$ becomes proportional to the inverse square of the constant of proportionality, and one
obtains the smaller bound $C_\delta\|\mathbf{C}\|^2e^{-\frac12c'p_0(g_{k_0})}$; either bound is
absorbed by the condition $x_{k_0}\ge x_{\mathrm{pk}}$. Each summand of a weighted row sum is at
most the sum,
so $\|\partial_b\partial_{b'}P^{\mathrm{pk}}_j\|\le C_Pg_j\mathfrak{b}_je^{-c_Pd(b,b')}$ for all
bonds $b'$ of $\mathbb{B}^{+}$ and all real fields $B$ of the set of (f). The first-order bound of
(f) gives, for every
bond $b\in I$, on the event \eqref{8.3:eq-clean-event-a},
\begin{equation*}
|\partial_bP^{\mathrm{pk}}_j|_2
\le C_{\mathrm{grad}}\|\mathbf{C}\|\,\hat{\mathfrak{b}}_j\bigl(\mathfrak{b}_j^2+p_0(g_j)\bigr)
e^{-\frac12c'p_0(g_{k_0})}\le C_Pg_j\mathfrak{b}_j^2 ,
\end{equation*}
the last inequality being the second of the inequalities defining $x_{\mathrm{pk}}$ in
Definition~\ref{8.3:def-depth-coupling-conditions}; this is the gradient bound for these members. The
per-bond Hessian bound is not obtained on every history from the other clauses of
Definition~\ref{8.3:def-clean-event} alone: the number $L^{4(j-i')}$ of level-$i'$ terms under a
level-$j$ bond is exactly compensated, and not beaten, by the factor $L^{-4(j-i')}$ of the
second-order bound, which is why it is established only on the event
\eqref{8.3:eq-clean-event-a}. The weight
of the complement of that event enters no display of this lemma.

The $\mathbf{R}$-term member is bounded on every history. By the bound of (f) on this member,
resting on \cite[(2.31)]{B-CMP119} with Ba{\l}aban's exponent $\kappa_0\ge4$ of that display, and
Cauchy's estimate, its weighted Hessian
norm is at most
\begin{equation*}
Cg_jx_j^{(3-\kappa_0)/2}\,\mathrm{polylog}(x_j)/\mathfrak{b}_j
\le c_N\,Cg_jx_j^{(3-\kappa_0)/2}\,\mathrm{polylog}(x_j)/\hat{\mathfrak{b}}_j ,
\end{equation*}
since
$\mathfrak{b}_j\ge\hat{\mathfrak{b}}_j/c_N$; with $c_N$ absorbed into $C$, by the last of the
inequalities defining $x_{\mathrm{pk}}$ in Definition~\ref{8.3:def-depth-coupling-conditions} this
is at most $C_Pg_j\mathfrak{b}_j$, and,
as for $P^{\mathrm{pk}}_j$, each of its Hessian blocks satisfies the corresponding per-pair bound.
Both members are thus inside $\varepsilon_P$ once $C_P$ is
enlarged: adding the three bounds, on the event \eqref{8.3:eq-clean-event-a} the Hessian bound of
(e) holds for the whole of $P^{\mathrm{tot}}_j$, for $b\in I$ and $b'\in\mathbb{B}^{+}$ and at every
real field of the set of (f), with
$C_P$ replaced by a larger constant, still
written $C_P$, and (h5) holds with $\varepsilon_P=(\dim\mathfrak{g})C_Pg_j\mathfrak{b}_j$ whenever
\eqref{8.3:eq-localisation-hypotheses-1} holds. With this $\varepsilon_P$,
\eqref{8.3:eq-localisation-hypotheses-1} reads
$(\dim\mathfrak{g})C_PS(c_P/2)\,g_j\mathfrak{b}_j\le a_-/4$, a one-sided upper bound on the running
product $g_j\mathfrak{b}_j$; the lemma on localisation of the sourced step is stated under it, and,
once the constant $a_-$ is absorbed, it is weaker than the coupling ceiling under which the
Gaussian extraction is stated, to whose defining conditions it is adjoined.

Only the smallness of the weighted Hessian norm is used, not its order in $g_j$. If the
$\mathbf{E}$-member is estimated with the radii of \cite[(2.28)]{B-CMP119}, of order $g_j$ times a
power of $\log x_j$, namely $\alpha_{0,j}=C_0g_j(\log x_j)^{q_0}$, then the Cauchy radius in the
scaled variable is
$\alpha_{0,j}/(Cg_j)=(C_0/C)(\log x_j)^{q_0}$, and Cauchy's estimate gives a Hessian bound, hence
an admissible value of $\varepsilon_P$, of order $O((\log x_j)^{-2q_0})$, which tends to $0$; the
condition
\eqref{8.3:eq-localisation-hypotheses-1} is then again a one-sided smallness condition. This
alternative presupposes that $q_0$ exceeds, with a margin, the exponent of $\log x_j$ in the
half-width of the window.

\emph{Exterior locality of the force \textup{(h6)}, small-field setting.} In the setting of
\cite{B-CMP109} the exterior datum is the exterior field $B_O:=A_O(\mathsf{e})$ alone and all
exterior data share one background. Hence $\mathsf{Q}^{\mathsf{e}}=\mathsf{Q}^{\mathsf{e}_0}$ and
the first inequality of \eqref{8.3:eq-localisation-hypotheses-2} holds with $\eta_{\mathsf{Q}}=0$.
For the second, $P_{\mathsf{e}}(y)=P^{\mathrm{tot}}_j(y;B_O)$, and $P_0(y)$ differs from
$P^{\mathrm{tot}}_j(y;0)$ by a constant, since $\mathsf{e}_0$ has all fluctuation variables on $O$
equal to $0$. For $t\in[0,1]$ the field $tB_O$ stays in the window, a product of balls being
star-shaped about $0$, and the segments $tA_O(\mathsf{e})$, $t\in[0,1]$, for exterior data in the
range of (h4) lie in the tube of (e). By the fundamental theorem of calculus along
$t\mapsto(y,tB_O)$,
\begin{equation*}
\partial_b\bigl(P^{\mathrm{tot}}_j(y;B_O)-P^{\mathrm{tot}}_j(y;0)\bigr)
=\int_0^1\sum_{o\in O}\partial_o\partial_bP^{\mathrm{tot}}_j(y;tB_O)\,[B_O(o)]\,dt .
\end{equation*}
By the Hessian bound of (e) for the whole of $P^{\mathrm{tot}}_j$, established above at every real
field of the set of (f), with
$b\in I$ and $o\in O$, as in the bound on
$\|\partial_b\partial_{b'}P_0\|$ above, and by $|B_O(o)|_2\le\hat{\mathfrak{b}}_j$, which holds
because in this setting every exterior variable is restricted by the small-field window, so that
(h4) holds with constant $1$ as on the free bonds of $\Lambda_{j+1}$ (and so, $C_O$ being at least
$1$, the field $(y,tB_O)$ belongs to the set of (f)), the modulus is
at most
\begin{equation*}
(\dim\mathfrak{g})\,C_Pg_j\mathfrak{b}_j\,\hat{\mathfrak{b}}_j\sum_{o\in O}e^{-c_Pd(b,o)} .
\end{equation*}
Since $d(b,o)\ge d(b,O)$ for $o\in O$,
\begin{equation*}
\sum_{o\in O}e^{-c_Pd(b,o)}
\le e^{-\frac{c_P}{2}d(b,O)}\sum_{o\in O}e^{-\frac{c_P}{2}d(b,o)}
\le e^{-\frac{c_P}{2}d(b,O)}\,S(c_P/2),
\end{equation*}
$S(c)$ bounding $\sum_{b'}e^{-c\,d(b,b')}$ uniformly in $b$. This is the second inequality of
\eqref{8.3:eq-localisation-hypotheses-2} with
$\varepsilon_\partial=(\dim\mathfrak{g})C_Pg_j\mathfrak{b}_j\hat{\mathfrak{b}}_jS(c_P/2)$ and
$c_\partial=c_P/2$.

\emph{Exterior locality of the force \textup{(h6)}, general case.} For the fluctuation members of
$\mathsf{e}$ the argument is the one just given, the segments from $0$ to the exterior fields
lying in the tube of (e) by (h4). For the members of $\mathsf{e}$ that move the background, the
exterior dependence of $\mathsf{Q}^{\mathsf{e}}_{II}$ and of $P_{\mathsf{e}}$ is given by the
bounds collected in (g): the exponential decay in the distance of the exterior variable from the
box, analytic with Cauchy bounds on the complex window, from clauses (D), (E), (H) and (A) and the
geometric clause of the operator theorem at free current;
the box-scale localisation bound among the inputs of the Gaussian extraction; and, for the terms
generated by $\mathbf{B}$ and the
effective potential on the large-field region, the localisation close to the large-field regions
of \cite[(2.34)--(2.35)]{B-CMP119} with the factor $e^{-\kappa d(Y)}$, combined with Cauchy's
estimate. By (c) every admissible exterior datum has its background in $\mathcal{U}_{\mathbb{R}}$,
so these bounds apply to every $\mathsf{e}$. By (g) the resulting families satisfy the hypotheses
of Lemma~\ref{8.3:lem-exterior-force-locality}, which turns them into the two inequalities of
\eqref{8.3:eq-localisation-hypotheses-2}. The constants $\varepsilon_\partial$ and
$\eta_{\mathsf{Q}}$ need not be small: in the lemma on localisation of the sourced step they are
multiplied by a
factor decaying exponentially in $D_\star=d(S^{+},O)$, so any bound polynomial in $\mathfrak{b}_j$
serves, and, the amplitude parameters $v_1,v_2$ supplied by (g) being at most a constant times
$\hat{\mathfrak{b}}_j$ and $\epsilon_0$ at most a constant times $\varepsilon_j$,
Lemma~\ref{8.3:lem-exterior-force-locality} gives
$\varepsilon_\partial\le C\mathfrak{b}_j$ and $\eta_{\mathsf{Q}}\hat{\mathfrak{b}}_j\le C\mathfrak{b}_j$.
\end{proof}

We use the notation of Lemma~\ref{8.3:lem-localisation-hypotheses} and hypotheses (h1)--(h6) listed there. This includes the level $j$ of the step, the box $I\subset\mathbb{B}^{+}$, $O:=\mathbb{B}^{+}\setminus I$, the read set $S^{+}\subset I$, $D_\star:=d(S^{+},O)$, $n_{S^{+}}:=\#S^{+}$, the window $\mathcal{W}_I=\prod_{b\in I}\mathsf{K}$ and $\hat{\mathfrak{b}}_j:=\sup_{X\in \mathsf{K}}|X|_2$. It also includes the exterior data $\mathsf{e}$ with reference datum $\mathsf{e}_0$, the matrices $\mathsf{Q}^{\mathsf{e}}$ and their box--exterior blocks $\mathsf{J}^{\mathsf{e}}=\mathsf{Q}^{\mathsf{e}}_{IO}$, the exterior field $A_O(\mathsf{e})$, the exponents $P_{\mathsf{e}}$ and $P_0$, and the constants $a_\pm$, $c_A$, $c_P$, $\varepsilon_P$, $C_O$, $\eta_{\mathsf{Q}}$, $\varepsilon_\partial$, $c_\partial$. Together with the read set $S^{+}$, Lemma~\ref{8.3:lem-localisation-hypotheses} gives a fixed set $\mathfrak{S}_G$ of level-$j$ variables and maps $\mathcal{V}_{\mathsf{e}}$ expressing $V_j|_{\mathfrak{S}_G}$ as a function of the box variables $y|_{S^{+}}$; this is what makes $S^{+}$ the read set. The observables $G,G'$ below are bounded Borel functions of $V_j|_{\mathfrak{S}_G}$.

The lattice sums $S(\cdot)$, $S^{1}(\cdot)$ and the weighting operators $\mathcal{T}_\varphi$ are those of Lemma~\ref{8.3:lem-weighted-inverse-decay}. For a bounded function $F$ on $\mathcal{W}_I$ we write $\|F\|_{\mathcal{W}_I}:=\sup_{y\in\mathcal{W}_I}|F(y)|$, and $y|_{S^{+}}$ denotes the restriction of $y$ to $S^{+}$.

We use the following constants. Let $\mu_1$ be the largest $\mu\le c_A/2$ with $a_+\mu S^{1}(c_A/2)\le a_-/4$. Put
\[
C_{\mathcal{S}}:=\frac{4a_+^2}{3a_-},\qquad
a_+^{\mathcal{K}}:=a_++C_{\mathcal{S}}\,S(c_A/2)^2 .
\]
Let $\mu_\star$ be the largest
\[
\mu\le\min\{\mu_1/2,\ c_P/2,\ c_\partial,\ c_A/4\}\quad\text{with}\quad a_+^{\mathcal{K}}\,\mu\,S^{1}(\mu_1/2)\le a_-/4 .
\]
Put
\begin{equation}\label{8.3:eq-box-localisation-1}
\begin{gathered}
\lambda:=\frac{a_-}{2},\qquad
\mathfrak{L}:=(a_+^{\mathcal{K}}+\varepsilon_P)\,S(2\mu_\star)^{1/2},\\
\bar\delta:=\bigl[a_+C_O\,S(c_A/2)+C_{\mathcal{S}}\,S(\mu_1)S(c_A)S(c_A/2)+\eta_{\mathsf{Q}}\bigr]
\,\hat{\mathfrak{b}}_j+\varepsilon_\partial ,
\end{gathered}
\end{equation}
and
\begin{equation}\label{8.3:eq-box-localisation-2}
\Xi:=2\hat{\mathfrak{b}}_j\,n_{S^{+}}\,\bar\delta\,
\Bigl[1+\frac{2\mathfrak{L}}{a_-}\,S(\mu_\star)^{1/2}\Bigr]\,e^{-\mu_\star D_\star/2}.
\end{equation}

For $0<\mu\le\mu_1$ we have $c_A-\mu\ge c_A/2$, hence $S^{1}(c_A-\mu)\le S^{1}(c_A/2)$. The defining inequality of $\mu_1$ therefore implies condition \eqref{8.3:eq-weighted-inverse-decay-a} for $A=\mathsf{Q}^{\mathsf{e}_0}$ at $c=c_A$, $\mu=\mu_1$. Likewise $\mu_\star\le\mu_1/2$ gives $S^{1}(\mu_1-\mu_\star)\le S^{1}(\mu_1/2)$. The defining inequality of $\mu_\star$ therefore implies condition \eqref{8.3:eq-weighted-inverse-decay-a} for the matrix $\mathcal{K}$ of the proof below, at $c=\mu_1$, $\mu=\mu_\star$.

The constants $\mu_1$, $\mu_\star$, $C_{\mathcal{S}}$ and $a_+^{\mathcal{K}}$ depend only on $(a_-,a_+,c_A,c_P,c_\partial)$ and on the metric $d$ through the sums $S(\cdot)$, $S^{1}(\cdot)$. The same holds for $\mathfrak{L}$ when $\varepsilon_P\le a_-$, since then $\mathfrak{L}\le(a_+^{\mathcal{K}}+a_-)S(2\mu_\star)^{1/2}$. None of them depends on the level $j$, on the number of renormalisation steps, on the torus exponent, on the bare or the reference coupling, on the depth below the cutoff, on any position, on the size of the box $I$ (in particular not on $D_\star$), on the extraction depth, on the volume, or on the history outside $I$.

With the values of $\eta_{\mathsf{Q}}$ and $\varepsilon_\partial$ supplied in Lemma~\ref{8.3:lem-localisation-hypotheses}, and since $\hat{\mathfrak{b}}_j\le c_N\mathfrak{b}_j$, with $c_N$ the norm-equivalence constant of Lemma~\ref{8.3:lem-localisation-hypotheses}, which depends on $N$ only, we have $\bar\delta\le C\mathfrak{b}_j$ with
\[
C:=c_N\bigl[a_+C_O\,S(c_A/2)+C_{\mathcal{S}}\,S(\mu_1)S(c_A)S(c_A/2)+\eta_{\mathsf{Q}}\bigr]+\varepsilon_\partial/\mathfrak{b}_j .
\]
Consequently, when $\varepsilon_P\le a_-$, so that $\mathfrak{L}\le(a_+^{\mathcal{K}}+a_-)S(2\mu_\star)^{1/2}$,
\[
\Xi\le C_\star\,n_{S^{+}}\,\mathfrak{b}_j^2\,e^{-\mu_\star D_\star/2},
\]
where
\[
C_\star:=2c_N C\Bigl[1+\frac{2}{a_-}\,(a_+^{\mathcal{K}}+a_-)\,S(2\mu_\star)^{1/2}S(\mu_\star)^{1/2}\Bigr].
\]
The condition $\varepsilon_P\le a_-$ follows from the one-sided condition $\varepsilon_P\,S(c_P/2)\le a_-/4$ of (h5) whenever $S(c_P/2)\ge1$. Besides the constants listed above, $C_\star$ depends on $C_O$ (hence on the ratio $A_0/A_1$ of Ba{\l}aban's auxiliary parameters of Notation~\ref{2.1:not-prefix} and on $N$), on $\eta_{\mathsf{Q}}$ and on $\varepsilon_\partial/\mathfrak{b}_j$. For general exterior data it also depends on the constants of the estimates from which $\eta_{\mathsf{Q}}$ and $\varepsilon_\partial$ are obtained. None of the quantities on which $C$ and $C_\star$ thus depend is a level, the volume, a position or a datum of the history.

Since $\mu_\star\le c_\partial$, and for general exterior data $c_\partial$ is at most a level-free constant divided by a length comparable with the block parameter $M$, the constant $\mu_\star$ depends on $M$ in that case, and so does every constant defined from it, in particular the constant $\mu_\star/4$ that serves as the decay rate in Theorem~\ref{8.3:thm-sourced-step-localisation}, and the depth floor defined from $\mu_\star$. None of these is a level, the volume, a position or a datum of the history. The parameter $M$ is a fixed power of $L$ and is not bounded from above here.

\begin{theorem}[Box localisation under the six hypotheses]\label{8.3:thm-box-localisation}
Assume \textup{(h1)}--\textup{(h6)} of Lemma~\ref{8.3:lem-localisation-hypotheses}, and let $\Xi$ be given by \eqref{8.3:eq-box-localisation-2}. Let $\Gamma:=(\mathsf{Q}^{\mathsf{e}_0})^{-1}$ on $\mathfrak{g}^{\mathbb{B}^{+}}$, let $\gamma^{\mathrm{box}}$ be the centred Gaussian measure on $\mathfrak{g}^{I}$ with covariance $\Gamma_{II}$, and for bounded Borel $\Phi$ on $\mathcal{W}_I$ put
\[
\langle\Phi\rangle^{\mathrm{box}}:=
\frac{\int_{\mathcal{W}_I}\Phi\,e^{P_0}\,d\gamma^{\mathrm{box}}}
{\int_{\mathcal{W}_I}e^{P_0}\,d\gamma^{\mathrm{box}}} .
\]
Write $\operatorname{Cov}^{\mathrm{box}}$ for the covariance under $\langle\cdot\rangle^{\mathrm{box}}$, and
\[
\operatorname{Cov}_{\mathbb{E}_{i,h}}(G,G'):=\mathbb{E}_{i,h}[GG']-\mathbb{E}_{i,h}[G]\,\mathbb{E}_{i,h}[G'],
\]
all evaluated at $V_{j+1}$. The configuration $V_{j+1}$ is fixed, and in \textup{(C)} we write $\mathbb{E}_{i,h}[G]$ for $\mathbb{E}_{i,h}[G](V_{j+1})$. Then the following hold.
\begin{enumerate}
\item[\textup{(A)}] \emph{(Density form.)} For every exterior datum $\mathsf{e}$, the law of $y|_{S^{+}}$ under $E_{\mathsf{e}}$ has a density with respect to its law under $\langle\cdot\rangle^{\mathrm{box}}$, with values in $[e^{-\Xi},e^{\Xi}]$. The same holds for every mixture of the $E_{\mathsf{e}}$.
\item[\textup{(B)}] \emph{(One background.)} Suppose $\mathcal{V}_{\mathsf{e}}=\mathcal{V}$ does not depend on $\mathsf{e}$, for instance when $\mathsf{e}$ ranges over exterior data with one background. Then for all bounded Borel functions $G,G'$ of $V_j|_{\mathfrak{S}_G}$ and all real $c,c'$,
\begin{align*}
\bigl|\mathbb{E}_{i,h}[G](V_{j+1})-\langle G\circ\mathcal{V}\rangle^{\mathrm{box}}\bigr|
&\le(e^{\Xi}-1)\,\|G\circ\mathcal{V}-c\|_{\mathcal{W}_I},\\
\bigl|\operatorname{Cov}_{\mathbb{E}_{i,h}}(G,G')
-\operatorname{Cov}^{\mathrm{box}}(G\circ\mathcal{V},G'\circ\mathcal{V})\bigr|
&\le 3(e^{\Xi}-1)\,\|G\circ\mathcal{V}-c\|_{\mathcal{W}_I}\,
\|G'\circ\mathcal{V}-c'\|_{\mathcal{W}_I}.
\end{align*}
\item[\textup{(C)}] \emph{(General exterior data.)} In general, put $\mathcal{V}_0:=\mathcal{V}_{\mathsf{e}_0}$ and $\eta_G:=\sup_{\mathsf{e}}\|G\circ\mathcal{V}_{\mathsf{e}}-G\circ\mathcal{V}_0\|_{\mathcal{W}_I}$. Then for $G,G',c,c'$ as in \textup{(B)},
\[
\bigl|\mathbb{E}_{i,h}[G]-\langle G\circ\mathcal{V}_0\rangle^{\mathrm{box}}\bigr|
\le(e^{\Xi}-1)\,\|G\circ\mathcal{V}_0-c\|_{\mathcal{W}_I}+\eta_G ,
\]
and
\begin{align*}
\bigl|\operatorname{Cov}_{\mathbb{E}_{i,h}}(G,G')
-\operatorname{Cov}^{\mathrm{box}}&(G\circ\mathcal{V}_0,G'\circ\mathcal{V}_0)\bigr|\\
&\le 3(e^{\Xi}-1)\,\|G\circ\mathcal{V}_0-c\|_{\mathcal{W}_I}\,
\|G'\circ\mathcal{V}_0-c'\|_{\mathcal{W}_I}\\
&\quad+2\eta_G\bigl(\|G'\circ\mathcal{V}_0-c'\|_{\mathcal{W}_I}+\eta_{G'}\bigr)
+2\eta_{G'}\,\|G\circ\mathcal{V}_0-c\|_{\mathcal{W}_I}.
\end{align*}
\end{enumerate}
\end{theorem}

\begin{proof}
Throughout, $\mathcal{W}:=\mathcal{W}_I$. By (h2), $\mathsf{K}$ is an open norm ball of $\mathfrak{g}$, so each factor of $\mathcal{W}$ is open, convex and contains a ball about $0$. By the definition of $\hat{\mathfrak{b}}_j$, every $y\in\mathcal{W}$ satisfies $|y_{b'}|_2\le\hat{\mathfrak{b}}_j$ for all $b'\in I$.

\emph{The reference Hamiltonian.} By (h3), $\mathsf{Q}^{\mathsf{e}_0}$ is real symmetric on $\mathfrak{g}^{\mathbb{B}^{+}}$ with $\mathsf{Q}^{\mathsf{e}_0}\ge a_->0$ and $\|\mathsf{Q}^{\mathsf{e}_0}(b,b')\|\le a_+e^{-c_Ad(b,b')}$. As noted after \eqref{8.3:eq-box-localisation-2}, condition \eqref{8.3:eq-weighted-inverse-decay-a} holds for $A=\mathsf{Q}^{\mathsf{e}_0}$ at $c=c_A$ and $\mu=\mu_1$, with $0<\mu_1<c_A$. Hence Lemma~\ref{8.3:lem-weighted-inverse-decay} applies to $\mathsf{Q}^{\mathsf{e}_0}$ with the decomposition $\mathbb{B}^{+}=I\sqcup O$. Put
\[
\mathcal{S}:=\mathsf{J}^{\mathsf{e}_0}(\mathsf{Q}^{\mathsf{e}_0}_{OO})^{-1}\mathsf{J}^{\mathsf{e}_0\,T},
\qquad
\mathcal{K}:=\bigl(\Gamma_{II}\bigr)^{-1}=\mathsf{Q}^{\mathsf{e}_0}_{II}-\mathcal{S},
\]
where the second expression for $\mathcal{K}$ is Lemma~\ref{8.3:lem-weighted-inverse-decay}(v). On $\mathcal{W}$ set
\[
H_0(y):=\tfrac12\langle y,\mathcal{K}y\rangle-P_0(y).
\]
By (h5), $P_0\in C^2(\mathcal{W})$ with bounded first and second derivatives, so $H_0$ is a Hamiltonian on $\mathcal{W}$ in the sense of Proposition~\ref{8.3:prop-lipschitz-stability}. The measure $\gamma^{\mathrm{box}}$ has density proportional to $e^{-\frac12\langle y,\mathcal{K}y\rangle}$. Hence
\[
\langle\Phi\rangle^{\mathrm{box}}
=\frac{\int_{\mathcal{W}}\Phi\,e^{-\frac12\langle y,\mathcal{K}y\rangle+P_0(y)}\,dy}
{\int_{\mathcal{W}}e^{-\frac12\langle y,\mathcal{K}y\rangle+P_0(y)}\,dy}
=\nu_{H_0}(\Phi),
\]
that is, $\nu_{H_0}=\langle\cdot\rangle^{\mathrm{box}}$. This measure has three features: it is a ratio over the box variables alone; its Gaussian part is the box marginal of the centred Gaussian measure with covariance $\Gamma$, so $H_0$ has no linear term; and the window $\mathcal{W}_I$ appears in both numerator and denominator. By (h1), $E_{\mathsf{e}}=\nu_{H_{\mathsf{e}}}$ for the Hamiltonian $H_{\mathsf{e}}$ displayed there.

\emph{Coercivity of $H_0$.} Since $\mathsf{Q}^{\mathsf{e}_0}\ge a_-$, we have $\Gamma\le a_-^{-1}$. Hence the principal block satisfies $\Gamma_{II}\le a_-^{-1}$ on $\mathfrak{g}^{I}$, and $\Gamma_{II}$ is positive, so its inverse satisfies $\mathcal{K}\ge a_-$ (this is also Lemma~\ref{8.3:lem-weighted-inverse-decay}(v)).

For the entries, $\|\mathsf{Q}^{\mathsf{e}_0}_{II}(b,b')\|\le a_+e^{-c_Ad(b,b')}$ by (h3). Lemma~\ref{8.3:lem-weighted-inverse-decay}(v), applied at $\mu=\mu_1$, gives $\|\mathcal{S}(b,b')\|\le C_{\mathcal{S}}S(c_A-\mu_1)^2e^{-\mu_1d(b,b')}$. We have $c_A-\mu_1\ge c_A/2$, the function $S$ is decreasing, and $e^{-c_Ad}\le e^{-\mu_1d}$. Therefore
\[
\|\mathcal{K}(b,b')\|
\le a_+e^{-c_Ad(b,b')}+C_{\mathcal{S}}S(c_A-\mu_1)^2e^{-\mu_1d(b,b')}
\le a_+^{\mathcal{K}}\,e^{-\mu_1d(b,b')}.
\]
So $\mathcal{K}$ is real symmetric on $\mathfrak{g}^{I}$ and satisfies the hypotheses of Lemma~\ref{8.3:lem-weighted-inverse-decay} with constants $(a_-,a_+^{\mathcal{K}})$ and decay rate $c=\mu_1$. Condition \eqref{8.3:eq-weighted-inverse-decay-a} holds at $\mu=\mu_\star<\mu_1$, as observed after \eqref{8.3:eq-box-localisation-2}, because $S^{1}(\mu_1-\mu_\star)\le S^{1}(\mu_1/2)$. Lemma~\ref{8.3:lem-weighted-inverse-decay}(i) then gives
\[
\langle\xi,\mathcal{T}_\varphi\mathcal{K}\mathcal{T}_\varphi^{-1}\xi\rangle\ge\tfrac34a_-|\xi|^2
\qquad\text{for all real }\xi\text{ and every }\varphi\text{ with }\operatorname{Lip}(\varphi)\le\mu_\star .
\]

By (h5), $\|\partial_b\partial_{b'}P_0(y)\|\le\varepsilon_Pe^{-c_Pd(b,b')}$ on $\mathcal{W}$. Since $\mu_\star\le c_P/2<c_P$, Lemma~\ref{8.3:lem-weighted-inverse-decay}(iv) applies to the matrix $P_0''(y)$ with $\varepsilon=\varepsilon_P$, $c=c_P$ and $\mu=\mu_\star$. Using that $S$ is decreasing and the one-sided condition $\varepsilon_P\,S(c_P/2)\le a_-/4$ of (h5), we get, for every $y\in\mathcal{W}$ and $\operatorname{Lip}(\varphi)\le\mu_\star$,
\[
\|\mathcal{T}_\varphi P_0''(y)\mathcal{T}_\varphi^{-1}\|
\le\varepsilon_P\,S(c_P-\mu_\star)\le\varepsilon_P\,S(c_P/2)\le\tfrac14a_- .
\]
Hence $H_0$ satisfies the coercivity condition \eqref{8.3:eq-lipschitz-stability-a} with $\lambda_A=\tfrac34a_-$ and $\varepsilon_0=\tfrac14a_-<\lambda_A$, for every $\varphi$ with $\operatorname{Lip}(\varphi)\le\mu_\star$. The corresponding constant is $\lambda=\lambda_A-\varepsilon_0=a_-/2$, as in \eqref{8.3:eq-box-localisation-1}.

\emph{The discrepancies $\delta_b$.} Fix an exterior datum $\mathsf{e}$ and compare $H_0$ with $H_{\mathsf{e}}$, written in the form of Proposition~\ref{8.3:prop-lipschitz-stability}. In that notation, the quadratic forms of $H_{\mathsf{e}}$ and $H_0$ differ by $(\mathsf{Q}^{\mathsf{e}}_{II}-\mathsf{Q}^{\mathsf{e}_0}_{II})+\mathcal{S}$, the linear term of $H_{\mathsf{e}}$ is $\mathsf{J}^{\mathsf{e}}A_O(\mathsf{e})$ while that of $H_0$ vanishes, and the perturbations differ by $P_{\mathsf{e}}-P_0$.
By its definition and the triangle inequality, the quantity $\delta_b$ of \eqref{8.3:eq-lipschitz-stability-b} for this pair is at most the supremum over $y\in\mathcal{W}$ of the sum of four terms. These are the norms of the $b$-components of $(\mathsf{Q}^{\mathsf{e}}_{II}-\mathsf{Q}^{\mathsf{e}_0}_{II})y$, of $\mathcal{S}y$, of $\mathsf{J}^{\mathsf{e}}A_O(\mathsf{e})$, and of $\partial_b(P_{\mathsf{e}}-P_0)(y)$. We bound them in turn for $y\in\mathcal{W}$ and $b\in I$.

\emph{The term $\mathcal{S}y$.} Since $|y_{b'}|_2\le\hat{\mathfrak{b}}_j$, we have $|(\mathcal{S}y)_b|\le\hat{\mathfrak{b}}_j\sum_{b'}\|\mathcal{S}(b,b')\|$. By Lemma~\ref{8.3:lem-weighted-inverse-decay}(v),
\[
|(\mathcal{S}y)_b|
\le C_{\mathcal{S}}\,S(\mu_1)S(c_A)S(c_A/2)\,\hat{\mathfrak{b}}_j\,e^{-(c_A/2)d(b,O)} .
\]

\emph{The term $\mathsf{J}^{\mathsf{e}}A_O(\mathsf{e})$.} The entries of $\mathsf{J}^{\mathsf{e}}$ are bounded by (h3), and $|A_O(\mathsf{e})(o)|_2\le C_O\hat{\mathfrak{b}}_j$ for $o\in O$ by (h4). Hence
\[
|(\mathsf{J}^{\mathsf{e}}A_O(\mathsf{e}))_b|
\le a_+C_O\,\hat{\mathfrak{b}}_j\sum_{o\in O}e^{-c_Ad(b,o)} .
\]
Write $e^{-c_Ad(b,o)}=e^{-(c_A/2)d(b,o)}\,e^{-(c_A/2)d(b,o)}$. Bound the first factor by $e^{-(c_A/2)d(b,O)}$, and sum the second over $o$, which gives at most $S(c_A/2)$. Thus
\[
|(\mathsf{J}^{\mathsf{e}}A_O(\mathsf{e}))_b|
\le a_+C_O\,S(c_A/2)\,\hat{\mathfrak{b}}_j\,e^{-(c_A/2)d(b,O)} .
\]

\emph{The remaining two terms.} By (h6),
\[
|((\mathsf{Q}^{\mathsf{e}}_{II}-\mathsf{Q}^{\mathsf{e}_0}_{II})y)_b|_2\le\eta_{\mathsf{Q}}\hat{\mathfrak{b}}_j\,e^{-c_\partial d(b,O)},
\qquad
|\partial_b(P_{\mathsf{e}}-P_0)(y)|_2\le\varepsilon_\partial\,e^{-c_\partial d(b,O)} .
\]

Since $\mu_\star\le\min\{c_A/4,c_\partial\}$, each of the four exponentials is at most $e^{-\mu_\star d(b,O)}$. Adding the four bounds and comparing with \eqref{8.3:eq-box-localisation-1} gives
\begin{equation}\label{8.3:eq-box-localisation-3}
\delta_b\le\bar\delta\,e^{-\mu_\star d(b,O)}\qquad(b\in I).
\end{equation}
Thus every discrepancy between the step at a frozen exterior and the box ratio enters as a force localised at the edge of the box.

\emph{The sums of the density bound.} We apply Theorem~\ref{8.3:thm-density-bound} to the pair $(H_0,H_{\mathsf{e}})$, with $\mu=\mu_\star$, read set $S^{+}$ and $R:=I\setminus S^{+}$. The coercivity hypothesis was established above, with the same $(\lambda_A,\varepsilon_0)=(\tfrac34a_-,\tfrac14a_-)$ for every $\varphi$ with $\operatorname{Lip}(\varphi)\le\mu_\star$. We now bound the quantities $\delta_s$, $\mathfrak{D}_s$, $\mathfrak{L}_s$ and $\operatorname{diam}(\mathsf{K})$ that enter its exponent.

\emph{The bound on $\delta_s$.} Let $s\in S^{+}$. Then $d(s,O)\ge D_\star$, so \eqref{8.3:eq-box-localisation-3} gives $\delta_s\le\bar\delta\,e^{-\mu_\star D_\star}$.

\emph{The bound on $\mathfrak{D}_s$.} For $b\in R$, inequality \eqref{8.3:eq-box-localisation-3} gives
\begin{align*}
e^{-2\mu_\star d(b,s)}\delta_b^2
&\le\bar\delta^2\,e^{-\mu_\star d(b,s)}\,e^{-\mu_\star(d(b,s)+2d(b,O))}\\
&\le\bar\delta^2\,e^{-\mu_\star d(b,s)}\,e^{-\mu_\star(d(b,s)+d(b,O))}\\
&\le\bar\delta^2\,e^{-\mu_\star d(b,s)}\,e^{-\mu_\star D_\star}.
\end{align*}
In the last step we used $d(b,s)+d(b,O)\ge d(s,O)\ge D_\star$. This holds because $d(s,o)\le d(s,b)+d(b,o)$ for every $o\in O$; take the infimum over $o$. Summing over $b\in R$ and using $\sum_b e^{-\mu_\star d(b,s)}\le S(\mu_\star)$, we get
\[
\mathfrak{D}_s\le\bar\delta\,S(\mu_\star)^{1/2}\,e^{-\mu_\star D_\star/2}.
\]

\emph{The bound on $\mathfrak{L}_s$.} The matrix of Theorem~\ref{8.3:thm-density-bound} for the reference Hamiltonian $H_0$ is $\mathcal{K}-P_0''(y)$. By the entry bound on $\mathcal{K}$ and by (h5),
\[
\|(\mathcal{K}-P_0''(y))(b,s)\|
\le a_+^{\mathcal{K}}e^{-\mu_1d(b,s)}+\varepsilon_Pe^{-c_Pd(b,s)}
\le(a_+^{\mathcal{K}}+\varepsilon_P)\,e^{-2\mu_\star d(b,s)},
\]
since $\mu_1\ge2\mu_\star$ and $c_P\ge2\mu_\star$. Hence
\[
\mathfrak{L}_s
\le(a_+^{\mathcal{K}}+\varepsilon_P)\Bigl(\sum_{b}e^{2\mu_\star d(b,s)}e^{-4\mu_\star d(b,s)}\Bigr)^{1/2}
=(a_+^{\mathcal{K}}+\varepsilon_P)\Bigl(\sum_{b}e^{-2\mu_\star d(b,s)}\Bigr)^{1/2}
\le\mathfrak{L}.
\]

\emph{The diameter.} Since $\mathsf{K}$ lies in the ball of radius $\hat{\mathfrak{b}}_j$ for $|\cdot|_2$, we have $\operatorname{diam}(\mathsf{K})\le2\hat{\mathfrak{b}}_j$.

\emph{Conclusion for the exponent.} Using $\lambda=a_-/2$ and $e^{-\mu_\star D_\star}\le e^{-\mu_\star D_\star/2}$, the exponent of Theorem~\ref{8.3:thm-density-bound} is bounded as follows:
\begin{align*}
\sum_{s\in S^{+}}\operatorname{diam}(\mathsf{K})\Bigl[\delta_s+\frac{\mathfrak{L}_s\mathfrak{D}_s}{\lambda}\Bigr]
&\le n_{S^{+}}\cdot2\hat{\mathfrak{b}}_j\Bigl[\bar\delta\,e^{-\mu_\star D_\star}
+\frac{2\mathfrak{L}}{a_-}\,\bar\delta\,S(\mu_\star)^{1/2}e^{-\mu_\star D_\star/2}\Bigr]\\
&\le2\hat{\mathfrak{b}}_j\,n_{S^{+}}\,\bar\delta\Bigl[1+\frac{2\mathfrak{L}}{a_-}S(\mu_\star)^{1/2}\Bigr]
e^{-\mu_\star D_\star/2}
=\Xi .
\end{align*}

\emph{Proof of \textup{(A)}.} Let $m_0$ and $m_{\mathsf{e}}$ be the densities of the $S^{+}$-marginals of $\nu_{H_0}=\langle\cdot\rangle^{\mathrm{box}}$ and $\nu_{H_{\mathsf{e}}}=E_{\mathsf{e}}$ on $\prod_{s\in S^{+}}\mathsf{K}$. Since its exponent is at most $\Xi$ by the preceding paragraph, Theorem~\ref{8.3:thm-density-bound} gives $e^{-\Xi}\le m_{\mathsf{e}}/m_0\le e^{\Xi}$ pointwise. Thus $m_{\mathsf{e}}/m_0$ is a density of the law of $y|_{S^{+}}$ under $E_{\mathsf{e}}$ with respect to its law under $\langle\cdot\rangle^{\mathrm{box}}$, with values in $[e^{-\Xi},e^{\Xi}]$. This holds for every exterior datum $\mathsf{e}$. Every $E_{\mathsf{e}}$ therefore satisfies Theorem~\ref{8.3:thm-density-bound} against the same $H_0$ with exponent at most $\Xi$. Corollary~\ref{8.3:cor-mixture-covariances} then gives the same two-sided density bound for every mixture of the $E_{\mathsf{e}}$.

\emph{Proof of \textup{(B)}.} Let $\nu:=\int\pi(d\mathsf{e})\,E_{\mathsf{e}}$, with $\pi$ the mixing measure of (h1). Since $\mathcal{V}_{\mathsf{e}}=\mathcal{V}$ and $S^{+}$ is the read set, $G\circ\mathcal{V}$ and $G'\circ\mathcal{V}$ are bounded Borel functions of $y|_{S^{+}}$. By (h1), applied to $G$, to $G'$ and to $GG'$,
\[
\mathbb{E}_{i,h}[G](V_{j+1})=\nu(G\circ\mathcal{V}),
\qquad
\operatorname{Cov}_{\mathbb{E}_{i,h}}(G,G')=\operatorname{Cov}_\nu(G\circ\mathcal{V},G'\circ\mathcal{V}).
\]
Since $\nu_{H_0}=\langle\cdot\rangle^{\mathrm{box}}$, Corollary~\ref{8.3:cor-mixture-covariances}(i), applied to $\nu$ and the function $G\circ\mathcal{V}$ with the constant $c$, gives the first inequality of (B). Corollary~\ref{8.3:cor-mixture-covariances}(ii), applied to $\nu$ and the pair $G\circ\mathcal{V}$, $G'\circ\mathcal{V}$ with the constants $c,c'$, gives the second.

\emph{Proof of \textup{(C)}, the expectation.} By (h1),
\[
\mathbb{E}_{i,h}[G]
=\int\pi(d\mathsf{e})\,E_{\mathsf{e}}[G\circ\mathcal{V}_0]
+\int\pi(d\mathsf{e})\,E_{\mathsf{e}}\bigl[G\circ\mathcal{V}_{\mathsf{e}}-G\circ\mathcal{V}_0\bigr].
\]
The first term is $\nu(G\circ\mathcal{V}_0)$, with $\nu$ as above. The argument of (B), with $\mathcal{V}_0$ in place of $\mathcal{V}$, bounds its distance from $\langle G\circ\mathcal{V}_0\rangle^{\mathrm{box}}$ by $(e^{\Xi}-1)\|G\circ\mathcal{V}_0-c\|_{\mathcal{W}_I}$. In the second term, each $E_{\mathsf{e}}$ is a probability measure on $\mathcal{W}_I$, so
\[
\bigl|E_{\mathsf{e}}[G\circ\mathcal{V}_{\mathsf{e}}-G\circ\mathcal{V}_0]\bigr|
\le\|G\circ\mathcal{V}_{\mathsf{e}}-G\circ\mathcal{V}_0\|_{\mathcal{W}_I}\le\eta_G .
\]
Since $\pi$ is a probability measure, the second term is at most $\eta_G$ in absolute value. This proves the first inequality of (C).

\emph{Proof of \textup{(C)}, the covariance.} Let $\hat\nu(d\mathsf{e},dy):=\pi(d\mathsf{e})\,E_{\mathsf{e}}(dy)$ be the joint law of $(\mathsf{e},y)$. Put
\begin{align*}
X&:=G\circ\mathcal{V}_{\mathsf{e}}(y)-c, & X_0&:=G\circ\mathcal{V}_0(y)-c,\\
X'&:=G'\circ\mathcal{V}_{\mathsf{e}}(y)-c', & X'_0&:=G'\circ\mathcal{V}_0(y)-c'.
\end{align*}
By (h1), applied to $G$, $G'$ and $GG'$, and because covariances are unchanged by adding constants, $\operatorname{Cov}_{\mathbb{E}_{i,h}}(G,G')=\operatorname{Cov}_{\hat\nu}(X,X')$. Bilinearity of the covariance gives
\[
\operatorname{Cov}_{\hat\nu}(X,X')-\operatorname{Cov}_{\hat\nu}(X_0,X'_0)
=\operatorname{Cov}_{\hat\nu}(X-X_0,X')+\operatorname{Cov}_{\hat\nu}(X_0,X'-X'_0).
\]
For bounded $U,V$ we have
\[
|\operatorname{Cov}(U,V)|\le|\mathbb{E}[UV]|+|\mathbb{E}U|\,|\mathbb{E}V|\le2\sup|U|\,\sup|V|,
\]
the suprema taken over $\mathcal{W}_I$ and the range of $\mathsf{e}$. Moreover,
\begin{align*}
\sup|X-X_0|&\le\eta_G,\qquad \sup|X'-X'_0|\le\eta_{G'},\\
\sup|X_0|&=\|G\circ\mathcal{V}_0-c\|_{\mathcal{W}_I},\\
\sup|X'|&\le\|X'_0\|_{\mathcal{W}_I}+\eta_{G'}
=\|G'\circ\mathcal{V}_0-c'\|_{\mathcal{W}_I}+\eta_{G'}.
\end{align*}
Therefore
\[
\bigl|\operatorname{Cov}_{\hat\nu}(X,X')-\operatorname{Cov}_{\hat\nu}(X_0,X'_0)\bigr|
\le2\eta_G\bigl(\|G'\circ\mathcal{V}_0-c'\|_{\mathcal{W}_I}+\eta_{G'}\bigr)
+2\eta_{G'}\,\|G\circ\mathcal{V}_0-c\|_{\mathcal{W}_I}.
\]
The variables $X_0$ and $X'_0$ do not depend on $\mathsf{e}$, and the $y$-marginal of $\hat\nu$ is $\nu=\int\pi(d\mathsf{e})E_{\mathsf{e}}$. Hence $\operatorname{Cov}_{\hat\nu}(X_0,X'_0)=\operatorname{Cov}_\nu(G\circ\mathcal{V}_0,G'\circ\mathcal{V}_0)$. By Corollary~\ref{8.3:cor-mixture-covariances}(ii), applied as in (B) with $\mathcal{V}_0$ in place of $\mathcal{V}$, this differs from $\operatorname{Cov}^{\mathrm{box}}(G\circ\mathcal{V}_0,G'\circ\mathcal{V}_0)$ by at most
\[
3(e^{\Xi}-1)\|G\circ\mathcal{V}_0-c\|_{\mathcal{W}_I}\|G'\circ\mathcal{V}_0-c'\|_{\mathcal{W}_I}.
\]
The triangle inequality now gives the second inequality of (C).
\end{proof}

\begin{theorem}[Localisation of the sourced step]\label{8.3:thm-sourced-step-localisation}
Let $K$, $s$, $k_0:=K-s$, $r\ge 1$, $D\ge 1$, the support $S\subset Q_S$ (a cube of side $R_S$)
and the block parameter $M$ be as in the data of the Gaussian extraction, and let
$\mathcal{C}_S(D)$ be the clean event of Definition~\ref{8.3:def-clean-event}. Let
$j\in[k_0,k_0+r)$ and $i:=j-k_0$.

The following objects are those of Lemmas~\ref{8.3:lem-box-law-mixture}
and~\ref{8.3:lem-localisation-hypotheses}: the bond set $\mathbb{B}^{+}$ on which the quadratic
form of the step acts, the read set $S^{+}$ with
$n_{S^+}:=\#S^{+}$, the radius $\hat{\mathfrak{b}}_j$, the
fluctuation variable $y$, the exterior data $\mathsf{e}$ and their admissibility, the
reference exterior datum $\mathsf{e}_0$, the branch expectation $\mathbb{E}_{i,h}$, the laws
$E_{\mathsf{e}}$ and the maps $\mathcal{V}_{\mathsf{e}}$. The following are those of the
Gaussian extraction: the radius $\mathfrak{b}_j$, the number $\vartheta_j$, the box
$\mathfrak{B}^{(i)}_D$ with its centre $S^{(i)}$, and the number $n$. The number $n$ enters
the Gaussian extraction through its variance bound $\mathfrak{v}=n/a_-$, with $a_-$ the lower
constant of the operator theorem at free current (see (c) below), and through
$\vartheta_j$; it is unrelated to the order of a Schwinger function. The domain
$\Lambda^{(j)*}_{j+1}$ is that of Ba{\l}aban's representation \cite[(3.25)]{B-CMP119}.

Two cases are distinguished. The \emph{pure small-field case} is that of \cite{B-CMP109} and
\cite{B-CMP116}. The \emph{complete case} is that of \cite{B-CMP119}, \cite{B-CMP122a} and
\cite{B-CMP122b} on a clean branch.

Assume the following.
\begin{itemize}
\item[(a)] \emph{Placement and positivity.} Near the box, the step of the branch at level $j$
is a term of Ba{\l}aban's representation \cite[(3.25)]{B-CMP119}, and the free bonds of the
box lie in $\Lambda^{(j)*}_{j+1}$. The operations of the branch are of the four kinds of
the segment lemma of the Gaussian extraction, and every weight in $\mathbf{T}_{j+1}$ is
non-negative. Clauses (c) and (d) of that segment lemma hold, together with the
identification of the type of the step which enters
Lemma~\ref{8.3:lem-box-law-mixture}.
\item[(b)] \emph{The window.} On the free bonds of the small-field domain the fluctuation
window is the product over bonds of norm balls of radius $\mathfrak{b}_j$,
as in \eqref{8.3:eq-window-product-a}.
\item[(c)] \emph{The quadratic form.} For every admissible exterior datum $\mathsf{e}$ the
quadratic form $\mathsf{Q}^{\mathsf{e}}$ of the step obeys the bounds of the operator
theorem at free current, with constants $a_-$, $a_+$, $c_A$. The background of every
admissible exterior datum, and that of the reference datum $\mathsf{e}_0$, lies in the real
domain of that theorem.
\item[(d)] \emph{The exterior fields.} The exterior fields of every admissible exterior datum
obey Ba{\l}aban's bounds \cite[(3.13)--(3.14)]{B-CMP119}, with the exterior-field constant
$C_O\ge 1$.
\item[(e)] \emph{The interaction.} For every admissible exterior datum $\mathsf{e}$ and for
$\mathsf{e}=\mathsf{e}_0$, on a complex tube about the window
$\prod_b\{|y_b|<\mathfrak{b}_j\}$ that also contains the range of the exterior fields
allowed by (d), the interaction $P^{\mathrm{tot}}_j(\cdot\,;\mathsf{e})$ of the step has the
following properties:
\begin{itemize}
\item it is holomorphic;
\item its bond gradient is at most $C_Pg_j\mathfrak{b}_j^2$;
\item its bond Hessians and mixed Hessians, in bonds $b,b'$, are at most
$C_Pg_j\mathfrak{b}_je^{-c_Pd(b,b')}$;
\item it vanishes identically at $g_j=0$.
\end{itemize}
The terms of $P^{\mathrm{tot}}_j$ attached to large-field regions are estimated on the
event of the per-bond small-piece clause of Definition~\ref{8.3:def-clean-event}.
\item[(f)] \emph{Locality of the exterior force.} The hypothesis of locality of the exterior
force among those of Lemma~\ref{8.3:lem-localisation-hypotheses} holds, with constants
$\eta_{\mathsf{Q}}$, $\varepsilon_\partial$ and decay rate $c_\partial$: for every admissible
exterior datum $\mathsf{e}$ and every free bond $b$ of the box, the quadratic form and the
force of the step at $b$ differ from those of $\mathsf{e}_0$ by at most
$\eta_{\mathsf{Q}}e^{-c_\partial d(b,O)}$ and $\varepsilon_\partial e^{-c_\partial d(b,O)}$
respectively, where $O$ is the set of bonds of $\mathbb{B}^{+}$ outside the box (see below),
as displayed in that lemma. In the
complete case (f) is assumed; it is the conclusion of
Lemma~\ref{8.3:lem-exterior-force-locality} once the hypotheses of that lemma are verified
from the localised formats of \cite[Lemma~2]{B-CMP116} and \cite[Theorem~2]{B-CMP119}, the
operator theorem at free current, (g) and the real-domain clause of (c). In the pure
small-field case (f) follows from the mixed-Hessian clause of (e)
(Lemma~\ref{8.3:lem-localisation-hypotheses}), and nothing further is assumed.
\item[(g)] \emph{Locality of the background.} Let $\mathcal{X}$ be the set of bonds carrying
the large-field data of the exterior data. For the bond set $\mathfrak{S}_G$ of every
observable $G$ considered in (C) below,
\begin{equation*}
\sup_{\mathsf{e}}\max_{b\in\mathfrak{S}_G}\bigl|V^{(j)}_{\mathsf{e}}(b)V^{(j)}_{\mathsf{e}_0}(b)^{-1}
-\mathbf{1}\bigr|\le C_\eta\varepsilon_j e^{-c_\eta\operatorname{dist}(\mathfrak{S}_G,\mathcal{X})},
\end{equation*}
modulo gauge; here $V^{(j)}_{\mathsf{e}}$ denotes the background configuration of level $j$
determined by the exterior datum $\mathsf{e}$, $\varepsilon_j$ is the small-field parameter of
level $j$ (Ba{\l}aban's $\varepsilon_k$ at $k=j$), and Ba{\l}aban's instance of this bound is
\cite[(3.19)]{B-CMP119}.
\item[(h)] \emph{Finite range of the observable map.} The field $V_j$ read by the
observables is a local function of the fluctuation variable $y$.
\item[(i)] \emph{Inputs of the small-piece clause.} The form of the large-field terms
\cite[(2.31)]{B-CMP119} holds, together with the two equivalence lemmas and the quadratic
background-deviation bounds of the Gaussian extraction, through which the clause of
Definition~\ref{8.3:def-clean-event} controls the terms named at the end of (e).
\item[(j)] \emph{The flow along the segment.} The lifetime and coupling-comparison statements
of the Gaussian extraction hold; by them, $x_j\ge x_{k_0}/2$ for $k_0\le j<k_0+r$.
\end{itemize}
Suppose that the running coupling at level $j$ obeys the ceiling
\begin{equation}\label{8.3:eq-sourced-step-localisation-1}
\varepsilon_P\,S_d(c_P/2)\le \frac{a_-}{4},\qquad
\varepsilon_P:=nC_Pg_j\mathfrak{b}_j,\qquad
S_d(c):=\sup_b\sum_{b'}e^{-c\,d(b,b')}.
\end{equation}
This ceiling holds for all $g_j\le\gamma_H$, where $\gamma_H$ is the
largest number such that \eqref{8.3:eq-sourced-step-localisation-1} holds for every
$g_j\le\gamma_H$; the number $\gamma_H$ depends on
$d,L,N,A_1,\mathfrak{p}_0,C_P,c_P,a_-$ only, where $\mathfrak{p}_0$ is the parameter of that
name among the data of the Gaussian extraction. Suppose also that
$x_{k_0}\ge x_{\mathrm{pk}}$, as in \eqref{8.3:eq-depth-coupling-conditions-a}.

For $h\in\mathcal{C}_S(D)$ and a configuration $V_{j+1}$ let $I$ be the set of
free bonds of the box $\mathfrak{B}^{(i)}_D$; by (a), $I\subset\Lambda^{(j)*}_{j+1}$. Put
$O:=\mathbb{B}^{+}\setminus I$, $D_\star:=d(S^+,O)$ and
$\Gamma:=(\mathsf{Q}^{\mathsf{e}_0})^{-1}$, and let
$\gamma^{\mathrm{box}}:=\mathcal{N}(0,\Gamma\restriction I)$. Let
$\mathcal{W}_I:=\prod_{b\in I}\{\xi:|\xi|<\mathfrak{b}_j\}$, and put
\begin{equation*}
P_0:=P^{\mathrm{tot}}_j(\cdot\,;\mathsf{e}_0)-P^{\mathrm{tot}}_j(0;\mathsf{e}_0),\qquad
\langle\Phi\rangle^{\mathrm{box}}
:=\frac{\int_{\mathcal{W}_I}\Phi\,e^{P_0}\,d\gamma^{\mathrm{box}}}
{\int_{\mathcal{W}_I}e^{P_0}\,d\gamma^{\mathrm{box}}}.
\end{equation*}
Let $\mathcal{V}_0:=\mathcal{V}_{\mathsf{e}_0}$, let $\|F\|_{\mathcal{W}_I}:=\sup_{\mathcal{W}_I}|F|$,
and let $E_{\mathsf{e}}$, $\mathcal{V}_{\mathsf{e}}$ be as in
Lemmas~\ref{8.3:lem-box-law-mixture} and~\ref{8.3:lem-localisation-hypotheses}. Let $\Xi$ be
the number of Theorem~\ref{8.3:thm-box-localisation}. By its definition there and
$\hat{\mathfrak{b}}_j\le c_N\mathfrak{b}_j$, with $c_N$ the norm-equivalence constant of
Lemma~\ref{8.3:lem-localisation-hypotheses}, it obeys
\begin{equation}\label{8.3:eq-sourced-step-localisation-2}
\Xi\le C_\star\,n_{S^+}\,\mathfrak{b}_j^2\,e^{-\mu_\star D_\star/2},\qquad
D_\star\ge DL-\ell_\star,\qquad \ell_\star:=M+2L+1,
\end{equation}
where the constants $C_\star$ and $\mu_\star$ depend on the following quantities only:
\begin{itemize}
\item $a_-,a_+,c_A,c_P,c_\partial,d,N$;
\item for $C_\star$, in addition $C_O$ (hence $A_0$ and $A_1$) and the ratios
$\eta_{\mathsf{Q}}/\hat{\mathfrak{b}}_j$, $\varepsilon_\partial/\mathfrak{b}_j$;
\item for $\mu_\star$ in the complete case, in addition the block parameter $M$, through
$c_\partial$: there $c_\partial$ is the minimum of $c_\eta$ and the quotient of the polymer
decay rate of Lemma~\ref{8.3:lem-exterior-force-locality} by the diameter scale of that lemma,
which is comparable to $M$.
\end{itemize}
They do not depend on $j$, $K$, $m$, $g_0$, $g$, $s$, $r$, $D$, the position, the volume, or
the history outside $I$. Then, for every such $h$ and $\tau_h$-almost every
$V_{j+1}$, the following hold.
\begin{itemize}
\item[(A)] \emph{Density form.} For every exterior datum $\mathsf{e}$, the law of
$y\restriction S^+$ under $E_{\mathsf{e}}$ has a density with respect to its law under
$\langle\cdot\rangle^{\mathrm{box}}$, with values in $[e^{-\Xi},e^{\Xi}]$. The same holds for
every mixture of the $E_{\mathsf{e}}$ with non-negative weights.
\item[(B)] \emph{Pure small-field case, or exterior data sharing one background.} In this
case $\mathcal{V}$ denotes the common
value of the maps $\mathcal{V}_{\mathsf{e}}$. Let $G$, $G'$ be bounded Borel functions of
$V_j\restriction\mathfrak{S}_G$ and $V_j\restriction\mathfrak{S}_{G'}$, where
$\mathfrak{S}_G$, $\mathfrak{S}_{G'}$ lie within $R_S+M$ of the centre $S^{(i)}$ (gauge
invariant in the complete case), and let $c,c'$ be real. Then
\begin{equation}\label{8.3:eq-sourced-step-localisation-3}
\begin{aligned}
&\bigl|\mathbb{E}_{i,h}[G](V_{j+1})-\langle G\circ\mathcal{V}\rangle^{\mathrm{box}}\bigr|
\le(e^{\Xi}-1)\,\|G\circ\mathcal{V}-c\|_{\mathcal{W}_I},\\
&\bigl|\operatorname{Cov}_{\mathbb{E}_{i,h}}(G,G')
-\operatorname{Cov}^{\mathrm{box}}(G\circ\mathcal{V},G'\circ\mathcal{V})\bigr|\\
&\qquad\le 3(e^{\Xi}-1)\,\|G\circ\mathcal{V}-c\|_{\mathcal{W}_I}\,
\|G'\circ\mathcal{V}-c'\|_{\mathcal{W}_I}.
\end{aligned}
\end{equation}
\item[(C)] \emph{Complete case.} The bounds of (B) hold with $\mathcal{V}_0$ in place of
$\mathcal{V}$, together with the additional terms of part~(C) of
Theorem~\ref{8.3:thm-box-localisation}. Those terms are built from
\begin{equation*}
\eta_G:=\sup_{\mathsf{e}}\|G\circ\mathcal{V}_{\mathsf{e}}-G\circ\mathcal{V}_0\|_{\mathcal{W}_I}
\end{equation*}
and from the corresponding $\eta_{G'}$. For $G$ in the family
$\mathfrak{F}_j(S;\mathfrak{a},\varrho)$ of observables of the Gaussian extraction,
where $\mathfrak{a}$ is the amplitude parameter of the family, one has
$\eta_G\le C_{\mathfrak{F}}\mathfrak{a}(\kappa+\epsilon)\epsilon$, where $\epsilon$ denotes
the left side of the display in (g), so that by (g)
\begin{equation*}
\epsilon\le C_\eta\varepsilon_je^{-c_\eta\operatorname{dist}(\mathfrak{S}_G,\mathcal{X})}.
\end{equation*}
\end{itemize}
Suppose in addition the optional one-sided floor on the free parameter $D$,
\begin{equation}\label{8.3:eq-sourced-step-localisation-4}
\tfrac14\mu_\star DL\ge\log\bigl(C_\star n_{S^+}\mathfrak{b}_j^2\bigr)+\tfrac12\mu_\star\ell_\star .
\end{equation}
Through $\mu_\star$, $\ell_\star$ and $C_\star$ this floor depends on $M$ (in the complete case
also through $c_\partial$), on $C_O$ and on the constants of (f). Under it,
\begin{equation}\label{8.3:eq-sourced-step-localisation-5}
\Xi\le e^{-\mu_\star DL/4}\le 1,\qquad e^{\Xi}-1\le 2\Xi .
\end{equation}
Consequently the localisation hypothesis of the Gaussian extraction holds with $C_H:=6$ and
$c_H:=\mu_\star/4$: the two left sides of \eqref{8.3:eq-sourced-step-localisation-3} are at
most
\begin{align*}
&6e^{-\mu_\star DL/4}\|G\circ\mathcal{V}-c\|_{\mathcal{W}_I},\\
&6e^{-\mu_\star DL/4}\|G\circ\mathcal{V}-c\|_{\mathcal{W}_I}\|G'\circ\mathcal{V}-c'\|_{\mathcal{W}_I},
\end{align*}
respectively. Without \eqref{8.3:eq-sourced-step-localisation-4}, the prefactor
$n_{S^+}\mathfrak{b}_j^2$ multiplies the two terms of the Gaussian extraction that carry the
factor $e^{-c''D}$.

Finally, the box datum has the properties that the Gaussian extraction uses:
\begin{itemize}
\item $\gamma^{\mathrm{box}}$ is centred;
\item its coordinate variances are at most $\mathfrak{v}=n/a_-$;
\item its plaquette covariances inside the box are exactly those of
$\Gamma=(\mathsf{Q}^{\mathsf{e}_0})^{-1}$;
\item $\sup_{\mathcal{W}_I}|P_0|\le\vartheta_j$;
\item the sandwich
\begin{equation}\label{8.3:eq-sourced-step-localisation-6}
\{|y_b|_2<\mathfrak{b}_j\ \forall b\}\subset\mathcal{W}_I\subset\{|y_b|_2<\hat{\mathfrak{b}}_j\ \forall b\}
\end{equation}
holds, so the sandwich hypothesis of the Gaussian extraction holds as a corollary.
\end{itemize}
\end{theorem}

\begin{proof}
\emph{Reduction to box localisation.} Fix $h\in\mathcal{C}_S(D)$ and let $V_{j+1}$ range
outside a $\tau_h$-null set. By Lemma~\ref{8.3:lem-localisation-hypotheses}, the following
inputs give the six hypotheses of Theorem~\ref{8.3:thm-box-localisation} for the set $I$ of
free bonds of the box, the
exterior $O$, the reference datum $\mathsf{e}_0$ and the exponent $P_0$:
\begin{itemize}
\item hypotheses (a)--(j);
\item the ceiling \eqref{8.3:eq-sourced-step-localisation-1};
\item the floor $x_{k_0}\ge x_{\mathrm{pk}}$ of \eqref{8.3:eq-depth-coupling-conditions-a}.
\end{itemize}
In particular, by Lemma~\ref{8.3:lem-box-law-mixture}, which applies under (a), we have
\begin{equation*}
\mathbb{E}_{i,h}[G](V_{j+1})=\int\pi_{h,V_{j+1}}(d\mathsf{e})\,
E_{\mathsf{e}}[G\circ\mathcal{V}_{\mathsf{e}}],
\end{equation*}
where $\pi_{h,V_{j+1}}$ is a probability measure on the exterior data. Parts (A), (B) and (C)
are parts (A), (B) and (C) of Theorem~\ref{8.3:thm-box-localisation}, read with the present
data, and, by the definition of $\Xi$ in that theorem and
$\hat{\mathfrak{b}}_j\le c_N\mathfrak{b}_j$, $\Xi$ obeys the bound
\eqref{8.3:eq-sourced-step-localisation-2} in terms of
$D_\star=d(S^+,O)$. It remains to supply five things:
\begin{itemize}
\item the geometry behind $D_\star$ and the form under the floor;
\item the oscillation form;
\item the bound on the box exponent;
\item the role of invariance;
\item the terms $\eta_G$.
\end{itemize}
The properties of the box datum are treated at the end.

\emph{Geometry.} The set $I$ consists of the bonds of $\mathfrak{B}^{(i)}_D$, which is the
centre $S^{(i)}$ enlarged by $D$ blocks of level $j+1$. Such a block has side $L$ in units of
level $j$, so the distance from the centre set to $O$ is at least $DL$. The sets
$\mathfrak{S}_G$ lie within $R_S+M$ of the centre, and $S^+$ lies within a further $2L$. The
cube $Q_S$, of side $R_S$, lies inside the centre $S^{(i)}$, so $R_S$ does not enter the lower
bound on $D_\star$. Hence
$D_\star\ge DL-\ell_\star$ with $\ell_\star:=M+2L+1$, and $n_{S^+}\le C_d(R_S+2M+4L)^4$ with
$C_d$ depending on $d$ only.

Take $c:=G(\mathcal{V}(0))$, the value of $G$ at the value of the observable map at zero
fluctuation. Write $\rho_G$ and $\rho_{GG'}$ for the two differences on the left of
\eqref{8.3:eq-sourced-step-localisation-3}. Part (B) gives
\begin{equation*}
|\rho_G|\le(e^{\Xi}-1)\|G\circ\mathcal{V}-G(\mathcal{V}(0))\|_{\mathcal{W}_I},
\end{equation*}
and, for arbitrary real $c,c'$,
\begin{equation*}
|\rho_{GG'}|\le3(e^{\Xi}-1)\|G\circ\mathcal{V}-c\|_{\mathcal{W}_I}
\|G'\circ\mathcal{V}-c'\|_{\mathcal{W}_I}.
\end{equation*}
Inserting $D_\star\ge DL-\ell_\star$ into \eqref{8.3:eq-sourced-step-localisation-2} gives
\begin{equation}\label{8.3:eq-sourced-step-localisation-7}
\Xi\le C_\star n_{S^+}\mathfrak{b}_j^2\,e^{\mu_\star\ell_\star/2}\,e^{-\mu_\star DL/2}.
\end{equation}

\emph{The form under the floor.} Exponentiating \eqref{8.3:eq-sourced-step-localisation-4}
gives $C_\star n_{S^+}\mathfrak{b}_j^2e^{\mu_\star\ell_\star/2}\le e^{\mu_\star DL/4}$. With
\eqref{8.3:eq-sourced-step-localisation-7} this yields
$\Xi\le e^{\mu_\star DL/4}e^{-\mu_\star DL/2}=e^{-\mu_\star DL/4}$, and the right side is at
most $1$. For $0\le x\le1$ convexity gives $e^x-1\le(e-1)x\le 2x$, hence $e^{\Xi}-1\le2\Xi$.
This proves \eqref{8.3:eq-sourced-step-localisation-5}. It follows that
\begin{align*}
|\rho_G|&\le2e^{-\mu_\star DL/4}\|G\circ\mathcal{V}-G(\mathcal{V}(0))\|_{\mathcal{W}_I},\\
|\rho_{GG'}|&\le6e^{-\mu_\star DL/4}\|G\circ\mathcal{V}-c\|_{\mathcal{W}_I}
\|G'\circ\mathcal{V}-c'\|_{\mathcal{W}_I}.
\end{align*}
Since $2\le 6$, both bounds have the form $C_He^{-c_HDL}$ with $C_H:=6$ and
$c_H:=\mu_\star/4$.

\emph{Without the floor.} We keep the prefactor. The two terms of the Gaussian extraction that
carry the factor $e^{-c''D}$ then gain the factor $n_{S^+}\mathfrak{b}_j^2$, with
$c''\le\mu_\star L/2$.

In the complete case, $\mu_\star$ depends on $M$ through $c_\partial$, and so do the floor
\eqref{8.3:eq-sourced-step-localisation-4} and $c_H$. The floor remains one-sided.

\emph{Oscillation form.} For a real bounded function $F$ on $\mathcal{W}_I$ write
$\operatorname{osc}_{\mathcal{W}_I}F:=\sup_{\mathcal{W}_I}F-\inf_{\mathcal{W}_I}F$, and write
$\operatorname{osc}(G):=\sup G-\inf G$ over the domain of $G$, so that
$\operatorname{osc}_{\mathcal{W}_I}(G\circ\mathcal{V})\le\operatorname{osc}(G)$. Then
$\inf_c\|F-c\|_{\mathcal{W}_I}=\tfrac12\operatorname{osc}_{\mathcal{W}_I}F$, the infimum being
attained at the midpoint of the range of $F$. Applied to $F=G\circ\mathcal{V}$ this gives
\begin{equation*}
\inf_c\|G\circ\mathcal{V}-c\|_{\mathcal{W}_I}
=\tfrac12\operatorname{osc}_{\mathcal{W}_I}(G\circ\mathcal{V})\le\tfrac12\operatorname{osc}(G).
\end{equation*}
The first bound of \eqref{8.3:eq-sourced-step-localisation-3} therefore gives
\begin{equation*}
\bigl|\mathbb{E}_{i,h}[G](V_{j+1})-\langle G\circ\mathcal{V}\rangle^{\mathrm{box}}\bigr|
\le\tfrac12(e^{\Xi}-1)\operatorname{osc}(G).
\end{equation*}
Under \eqref{8.3:eq-sourced-step-localisation-4}, by \eqref{8.3:eq-sourced-step-localisation-5},
the right side is at most $\operatorname{osc}(G)\,e^{-\mu_\star DL/4}$. For every $\ell\ge0$,
\begin{equation*}
e^{-\mu_\star DL/4}\le e^{-(\mu_\star L/4)(D-\ell)}.
\end{equation*}
Hence the oscillation form of the localisation hypothesis of the Gaussian extraction, a bound
$C_{\mathrm{loc}}\operatorname{osc}(G)e^{-c_{\mathrm{loc}}(D-\ell)}$ with $\ell\ge0$ the
offset of that hypothesis, holds with $C_{\mathrm{loc}}:=1$ and
$c_{\mathrm{loc}}:=\mu_\star L/4$.

\emph{The box exponent.} By definition $P_0(0)=0$. The window $\mathcal{W}_I$ is a product of
balls, so it is convex and contains $0$. The mean-value inequality along the segment from
$0$ to $y\in\mathcal{W}_I$ therefore gives
\begin{equation*}
|P_0(y)|\le\sum_{b\in I}\sup_{\mathcal{W}_I}|\partial_bP_0|\cdot|y_b|_2,
\end{equation*}
and $|y_b|_2<\hat{\mathfrak{b}}_j$ for every $b$ on $\mathcal{W}_I$, by the upper inclusion in
\eqref{8.3:eq-sourced-step-localisation-6}, proved below independently of $P_0$. The gradient
$\partial_bP_0=\partial_bP^{\mathrm{tot}}_j(\cdot\,;\mathsf{e}_0)$ is bounded on
$\mathcal{W}_I$ by (e), applied with $\mathsf{e}=\mathsf{e}_0$: the bond gradient is at most
$C_Pg_j\mathfrak{b}_j^2$, and since $n\ge1$ we use this in the form
$\sup_{\mathcal{W}_I}|\partial_bP_0|\le nC_Pg_j\mathfrak{b}_j^2$ for every $b\in I$. With
$n_D:=\#I$ we obtain
\begin{equation*}
\sup_{\mathcal{W}_I}|P_0|\le\sum_{b\in I}\sup_{\mathcal{W}_I}|\partial_bP_0|\cdot
\hat{\mathfrak{b}}_j\le n_D\cdot nC_Pg_j\mathfrak{b}_j^2\cdot\hat{\mathfrak{b}}_j .
\end{equation*}
The quantity $\vartheta_j$ of the Gaussian extraction is $n\cdot n_D\cdot C_Pg_j
\mathfrak{b}_j^2\cdot C_w\mathfrak{b}_j$, with $C_w$ the window constant of the Gaussian
extraction; since $\hat{\mathfrak{b}}_j\le C_w\mathfrak{b}_j$, the right side above is at most
$\vartheta_j$. The Gaussian
extraction also uses a splitting of the exponent by parity in $g_j$ and by size. This
splitting is inherited by $P_0$: it is $P^{\mathrm{tot}}_j$ at zero exterior fluctuation, a
function of $(g_j,y)$ through the same monomials.

One may prefer as box exponent the sum of Ba{\l}aban's localised terms whose localisation lies
inside $I$ (\cite[Lemma~2, (1.36), (1.41)--(1.43)]{B-CMP116}). By
Corollary~\ref{8.3:cor-box-exponent-choice}(i), that sum defines the same box law near $S^+$
up to $e^{\pm\Xi'}$, with $\Xi'$ of the same form as $\Xi$, provided two conditions hold:
\begin{itemize}
\item the omitted localised terms $\mathsf{P}_j(Y)$ of \cite[Lemma~2]{B-CMP116}, those
attached to domains $Y\ni b$ with $Y\not\subset I$, satisfy
\begin{equation*}
\Bigl|\sum_{Y\ni b,\,Y\not\subset I}\partial_b\mathsf{P}_j(Y)\Bigr|
\le\varepsilon_\partial e^{-c_\partial d(b,O)};
\end{equation*}
\item the kept sum obeys hypothesis (h5) of Lemma~\ref{8.3:lem-localisation-hypotheses}.
\end{itemize}
Both conditions follow from the sum \cite[(1.26)]{B-CMP116}, which carries the factor
$e^{-\kappa d_j(Y)}$, together with $d_j(Y)\ge d(b,O)/M$ and the Cauchy inequality. That form
of the exponent thus follows from \cite[Lemma~2, (1.26), (1.36), (1.41)--(1.43)]{B-CMP116}
read in this way; it is not used below. The form $P_0$ needs no polymer
representation at all, and it is the one adopted. The choice does not affect the conclusions.

\emph{Invariance.} Neither analyticity nor gauge invariance of $G$ is used in
Theorem~\ref{8.3:thm-box-localisation}. Invariance enters only through the Faddeev--Popov
step of Lemma~\ref{8.3:lem-box-law-mixture}, in the complete case.

\emph{The terms $\eta_G$.} Let $G\in\mathfrak{F}_j(S;\mathfrak{a},\varrho)$ and put
\begin{equation*}
\epsilon:=\sup_{\mathsf{e}}\max_{b\in\mathfrak{S}_G}\bigl|V^{(j)}_{\mathsf{e}}(b)
V^{(j)}_{\mathsf{e}_0}(b)^{-1}-\mathbf{1}\bigr|\quad\text{modulo gauge.}
\end{equation*}
By the two clauses of the invariance statement for $\mathfrak{F}_j$,
$\eta_G\le C_{\mathfrak{F}}\mathfrak{a}(\kappa+\epsilon)\epsilon$. By (g),
\begin{equation*}
\epsilon\le C_\eta\varepsilon_je^{-c_\eta\operatorname{dist}(\mathfrak{S}_G,\mathcal{X})};
\end{equation*}
Ba{\l}aban's own instance of this bound is \cite[(3.19)]{B-CMP119}. These are exactly the
tails that the Gaussian extraction already contains in its term with $e^{-c''D}$, so no new
term arises.

\emph{The box datum.} The law $\gamma^{\mathrm{box}}=\mathcal{N}(0,\Gamma\restriction I)$ is
centred by definition. It is the marginal on $I$ of $\mathcal{N}(0,\Gamma)$, so its
covariances inside the box are those of $\Gamma=(\mathsf{Q}^{\mathsf{e}_0})^{-1}$ exactly, and
there is no mismatch. The coordinate variances are at most $\mathfrak{v}=n/a_-$, by the lower
bound $a_-$ on the quadratic form in (c). The bound $\sup_{\mathcal{W}_I}|P_0|\le\vartheta_j$
was proved above. The inclusions \eqref{8.3:eq-sourced-step-localisation-6} compare the window
with Euclidean balls. The upper inclusion is the definition of $\hat{\mathfrak{b}}_j$ in
Lemma~\ref{8.3:lem-localisation-hypotheses} as the supremum of $|\xi|_2$ over the single-bond
window set, a number at most $c_N\mathfrak{b}_j$ with $c_N$ depending on $N$ only. The lower
inclusion states that on each bond the Euclidean ball $\{\xi:|\xi|_2<\mathfrak{b}_j\}$ lies in
the norm ball $\{\xi:|\xi|<\mathfrak{b}_j\}$ of the window \eqref{8.3:eq-window-product-a}.
The sandwich hypothesis of the Gaussian extraction is a consequence of these inclusions.
\end{proof}

\begin{remark}[The form obtained and the localisation hypothesis of the Gaussian extraction]
We compare the conclusions of Theorem~\ref{8.3:thm-sourced-step-localisation} with the
localisation hypothesis of the Gaussian extraction.
\begin{itemize}
\item[($\alpha$)] \emph{The error.} The error bound takes two forms, both uniform in the
exterior datum:
\begin{itemize}
\item the density bound (A);
\item the bounds (B)--(C), with the factor $e^{\Xi}-1$ and
$\Xi\le C_\star n_{S^+}\mathfrak{b}_j^2e^{-\mu_\star D_\star/2}$.
\end{itemize}
The form $C_He^{-c_HDL}$, with $C_H=6$ and $c_H=\mu_\star/4$, holds under
\eqref{8.3:eq-sourced-step-localisation-4}. Without that floor it holds with the prefactor
$n_{S^+}\mathfrak{b}_j^2$, which is of polylogarithmic size in $x_j$. The floor
\eqref{8.3:eq-sourced-step-localisation-4} is a lower bound on $D$ of the same kind
as the lower bound on $D$ among the data of the Gaussian extraction. In the applications of
the Gaussian extraction, $D\ge C\log x_{k_0}$ and $D=O(p_0(g_{k_0}))$, with the degree
function $p_0(\cdot)$; those values satisfy the floor, since
$n_{S^+}\le C_d(R_S+2M+4L)^4$ and $\log\mathfrak{b}_j=O(\log\log x_j)$. Without the floor the
prefactor is immaterial there as well: $D$ is a free parameter that enters the Gaussian
extraction only through the product of $x_{k_0}$, a polynomial in the window radius and
$e^{-c''D}$, and this product is made small by $D\ge C\log x_{k_0}$. The gradient-form
corollary of the second proof of box localisation gives a constant $C_H$ that carries the
factor $n_G\hat{\mathfrak{b}}_j^2$, with $n_G$ the number of bonds in $\mathfrak{S}_G$, when the
error is measured in the supremum norm, and is free of the level when it is measured in the
gradient norm.
\item[($\beta$)] \emph{The Gaussian.} The Gaussian is
$\gamma^{\mathrm{box}}=\mathcal{N}(0,\Gamma\restriction I)$. Here
$\Gamma=(\mathsf{Q}^{\mathsf{e}_0})^{-1}$ is the inverse of Ba{\l}aban's quadratic form
$\mathbf{C}^*\Delta^{(j)}\mathbf{C}$ of the fluctuation action (\cite{B-CMP109}), with
$\mathbf{C}$ the linear map taking the fluctuation variable to the bond field, taken
on the free bonds of $\Omega_{j+1}$, the
small-field domain of level $j+1$, built on the
branch's nested domains and on the reference background. So $\gamma^{\mathrm{box}}$ is the box
marginal; it is centred, and its precision is the Schur complement of
$\mathsf{Q}^{\mathsf{e}_0}$ with respect to the exterior bonds. The identification corollary
for the box covariance identifies it, near the support and within $e^{-\mu_1(DL-\ell_\star)}$,
where $\mu_1$ is the decay rate of that corollary, with the whole-lattice
covariance $C^{(j)}(U)$ of \cite[(2.13)]{B-CMP109}. That corollary is derived from the
covariance comparison lemma, from the operator theorem at free current, from (g), and from a
further statement of the Gaussian extraction which matches its box objects with Ba{\l}aban's.
Ba{\l}aban's Dirichlet box object $(\mathsf{Q}_{II})^{-1}=C^{(j)}(I)$ of \cite[(2.5)]{B-CMP116}
is also obtained, at the price $\mathfrak{v}_{\mathfrak{c}}e^{-2\mu\operatorname{dist}(\cdot,\partial I)}$
given by the third clause of the decoupling lemma, where $\mathfrak{v}_{\mathfrak{c}}$ and
$\mu$ are the variance constant and the decay rate of that lemma.
\item[($\gamma$)] \emph{The exponent.} The exponent is $P_0$, in place of the sum of
Ba{\l}aban's localised terms inside $I$. What the Gaussian extraction uses is the bound
$\sup_{\mathcal{W}_I}|P_0|\le\vartheta_j$ and the parity structure in $g_j$. By
Corollary~\ref{8.3:cor-box-exponent-choice}(i), the other form defines the same box law near
$S^+$ up to $e^{\pm\Xi'}$. That other form is proved as an implication from
\cite[Lemma~2, (1.26), (1.36), (1.41)--(1.43)]{B-CMP116}, as shown in the proof above. The
choice is immaterial.
\item[($\delta$)] \emph{The window.} The window is
$\mathcal{W}_I=\prod_{b\in I}\{\xi:|\xi|<\mathfrak{b}_j\}$, the product over the box's free bonds
of Ba{\l}aban's norm balls of radius $\mathfrak{b}_j$; see
\eqref{8.3:eq-window-product-a}. The sandwich \eqref{8.3:eq-sourced-step-localisation-6} is a
corollary of this and is strictly weaker. With the sandwich alone, the statement uniform in
the exterior datum is false, as the example accompanying \eqref{8.3:eq-window-product-a} shows.
\item[($\varepsilon$)] \emph{The observable.} By (h), the map $\mathcal{V}$ has no tails to
cut: $G\circ\mathcal{V}$ reads $y$ exactly on $S^+\subset I$. In the pure case, $G$ and $G'$
are arbitrary bounded Borel functions. In the complete case, the background read by
$\mathcal{V}_{\mathsf{e}}$ depends on the exterior data through the fields on large-field
zones (\cite[(3.10)--(3.11)]{B-CMP119}). This gives rise to the terms $\eta_G$ of (C), which
are controlled for $G\in\mathfrak{F}_j$.
\end{itemize}
\end{remark}

Under the law $\gamma^{\mathrm{box}}$ of Theorem~\ref{8.3:thm-box-localisation}, the covariance of $y_b$ and $y_{b'}$, for $b,b'\in I$, is exactly the entry $\Gamma(b,b')$ of $\Gamma=(\mathsf{Q}^{\mathsf{e}_0})^{-1}$. The reason is that $\gamma^{\mathrm{box}}$ is the centred Gaussian law with covariance the restriction of $\Gamma$ to $I$, which is the marginal on $I$ of the centred Gaussian law with covariance $\Gamma$.

The two-point quantities produced by the Gaussian extraction are therefore the normalised traces over the Lie algebra of the covariances, under the centred Gaussian law with covariance $\Gamma=(\mathsf{Q}^{\mathsf{e}_0})^{-1}$, of the differentiated observables of that extraction at two positions. They are built on the inverse of the form $\mathsf{Q}^{\mathsf{e}_0}$ of Theorem~\ref{8.3:thm-box-localisation}. This form is Ba{\l}aban's form $\mathbf{C}^*\Delta^{(j)}\mathbf{C}$, in the notation of \cite[(2.13)]{B-CMP109}, taken on the free bonds of the region of level $j+1$, and it is formed with the nested sequence of domains of the branch of the history under consideration and with the background of the exterior datum $\mathsf{e}_0$.

The statements used for clause~(iii) of Theorem~\ref{3.1:thm-main}, and the statements that express the kernel of the Gaussian extraction through Ba{\l}aban's propagator, are written with a different covariance. They use the propagator $C^{(j)}(U)=(\mathbf{C}^*\Delta^{(j)}\mathbf{C})^{-1}$ of \cite[(2.13)]{B-CMP109}, in the notation of that paper, taken on the whole lattice of level $j$, without a nest, at $U\equiv\mathbf{1}$. If instead one takes for $\Gamma$ the covariance of \cite[(3.23)]{B-CMP119}, namely $C^{(j)}(\Lambda_{j+1})$ in the notation of that paper, which equals $(\mathsf{Q}_{\mathbb{B}^*\mathbb{B}^*})^{-1}$, the inverse of the principal block of the form $\mathsf{Q}$ on the bond set $\mathbb{B}^*$, then Theorem~\ref{8.3:thm-box-localisation} holds verbatim with $\mathbb{B}^*\setminus I$ as the exterior index set in the Schur complement of Lemma~\ref{8.3:lem-weighted-inverse-decay}\,(v) and with the bonds of $\mathbb{B}^+\setminus\mathbb{B}^*$ counted in the exterior part of the form. The two box marginals so obtained differ as in Lemma~\ref{8.3:lem-inverse-comparison} below.

The following lemma compares, entry by entry on a common subset of their index sets, the inverses of two exponentially local positive forms.

\begin{lemma}[Comparison of inverses of exponentially local forms]\label{8.3:lem-inverse-comparison}
Let $\mathfrak{B}^1$ and $\mathfrak{B}^2$ be index sets carrying one metric $d$. For $a=1,2$ let $A^a$ be a real symmetric operator on $E_{\mathfrak{B}^a}$, with blocks $A^a(b,b')$, in the sense of Lemma~\ref{8.3:lem-weighted-inverse-decay}. Assume that for $a=1,2$
\begin{equation*}
A^a\ge a_->0,\qquad \|A^a(b,b')\|\le a_+e^{-c\,d(b,b')}\quad(b,b'\in\mathfrak{B}^a),
\end{equation*}
and that $0<\mu<c$ satisfies \eqref{8.3:eq-weighted-inverse-decay-a} for $A^1$ and for $A^2$. Put $\Gamma^a:=(A^a)^{-1}$.

Let $J\subset\mathfrak{B}^1\cap\mathfrak{B}^2$, and put
\begin{equation*}
d_\partial(b):=\min_{a=1,2}d(b,\mathfrak{B}^a\setminus J),
\end{equation*}
with the convention $d(b,\emptyset):=+\infty$. Suppose that, for some $\eta_0\ge0$,
\begin{equation}\label{8.3:eq-inverse-comparison-1}
\|(A^1-A^2)(x,x')\|\le\eta_0\,e^{-c\,d(x,x')}\,e^{-\mu d_\partial(x)}\qquad(x,x'\in J).
\end{equation}
Then for all $b,b'\in J$
\begin{equation}\label{8.3:eq-inverse-comparison-2}
\|\Gamma^1(b,b')-\Gamma^2(b,b')\|
\le\Bigl(\frac{4}{3a_-}\Bigr)^{2}S(c)S(\mu)
\Bigl[2a_+e^{-\frac{\mu}{2}(d_\partial(b)+d_\partial(b'))}+\eta_0e^{-\mu d_\partial(b)}\Bigr].
\end{equation}
Here, for $\nu>0$, $S(\nu)$ is the larger over $a=1,2$ of the lattice sums
\begin{equation*}
\sup_{b\in\mathfrak{B}^a}\sum_{b'\in\mathfrak{B}^a}e^{-\nu d(b,b')},
\end{equation*}
the sums of Lemma~\ref{8.3:lem-weighted-inverse-decay} formed on $\mathfrak{B}^1$ and on $\mathfrak{B}^2$.
\end{lemma}

\begin{proof}
Fix $a\in\{1,2\}$ and put $O:=O^a:=\mathfrak{B}^a\setminus J$. We write operators on $E_{\mathfrak{B}^a}$ in blocks according to $\mathfrak{B}^a=J\sqcup O$. For $b,b'\in J$ we have $\Gamma^a(b,b')=\Gamma^a_{JJ}(b,b')$.

\emph{Block inversion.} Since $A^a\ge a_-$, the operator $\Gamma^a$ is positive definite. Its principal block $\Gamma^a_{OO}$ is therefore positive definite, and $A^a_{JJ}\ge a_-$ is invertible. The identity $\Gamma^aA^a=\mathbf{1}$, read in the blocks $(J,J)$ and $(O,J)$, gives
\begin{equation*}
\Gamma^a_{JJ}A^a_{JJ}+\Gamma^a_{JO}A^a_{OJ}=\mathbf{1}_J,\qquad \Gamma^a_{OJ}A^a_{JJ}+\Gamma^a_{OO}A^a_{OJ}=0.
\end{equation*}
The second equation gives $A^a_{OJ}=-(\Gamma^a_{OO})^{-1}\Gamma^a_{OJ}A^a_{JJ}$. Inserting this into the first equation and multiplying on the right by $(A^a_{JJ})^{-1}$ yields
\begin{equation}\label{8.3:eq-inverse-comparison-3}
\Gamma^a_{JJ}-(A^a_{JJ})^{-1}=\Pi^a,\qquad \Pi^a:=\Gamma^a_{JO}(\Gamma^a_{OO})^{-1}\Gamma^a_{OJ}.
\end{equation}
If $O=\emptyset$, we read $\Pi^a=0$. By \eqref{8.3:eq-inverse-comparison-3}, for $b,b'\in J$,
\begin{equation}\label{8.3:eq-inverse-comparison-4}
\Gamma^1(b,b')-\Gamma^2(b,b')=\Pi^1(b,b')-\Pi^2(b,b')
+\bigl[(A^1_{JJ})^{-1}-(A^2_{JJ})^{-1}\bigr](b,b').
\end{equation}
We bound the three terms in turn.

\emph{The norm of $(\Gamma^a_{OO})^{-1}$.} The spectrum of $A^a$ lies in $[a_-,\|A^a\|]$, so $\Gamma^a\ge\|A^a\|^{-1}$. For $\xi\in E_O$, extended by zero to $E_{\mathfrak{B}^a}$, this gives
\begin{equation*}
\langle\xi,\Gamma^a_{OO}\xi\rangle=\langle\xi,\Gamma^a\xi\rangle\ge\|A^a\|^{-1}|\xi|^2.
\end{equation*}
Hence $\|(\Gamma^a_{OO})^{-1}\|\le\|A^a\|$.

Next we bound $\|A^a\|$. Since $d$ is symmetric, so is $e^{-c\,d(b,b')}$. For vectors $\xi,\xi'$ with components $\xi_b,\xi'_{b'}$,
\begin{align*}
|\langle\xi,A^a\xi'\rangle|
&\le\sum_{b,b'}a_+e^{-c\,d(b,b')}|\xi_b||\xi'_{b'}|\\
&\le\frac12\sum_{b,b'}a_+e^{-c\,d(b,b')}\bigl(|\xi_b|^2+|\xi'_{b'}|^2\bigr)\\
&\le\frac12\,a_+S(c)\bigl(|\xi|^2+|\xi'|^2\bigr).
\end{align*}
Taking $|\xi|=|\xi'|=1$ gives $\|A^a\|\le a_+S(c)$. Altogether,
\begin{equation}\label{8.3:eq-inverse-comparison-5}
\|(\Gamma^a_{OO})^{-1}\|\le\|A^a\|\le a_+S(c).
\end{equation}

\emph{Square sums of $\Gamma^a$ against $O$.} By Lemma~\ref{8.3:lem-weighted-inverse-decay}\,(ii), applied to $A^a$ (its hypotheses hold because $A^a\ge a_-$, the blocks of $A^a$ obey the assumed bound, and $\mu$ satisfies \eqref{8.3:eq-weighted-inverse-decay-a}), we have $\|\Gamma^a(b,o)\|\le\tfrac{4}{3a_-}e^{-\mu d(b,o)}$. By symmetry of $\Gamma^a$, $\|\Gamma^a(o,b)\|=\|\Gamma^a(b,o)\|$. Since $e^{-\mu d(b,o)}\le e^{-\mu d(b,O)}$ for $o\in O$, we obtain for $b\in J$
\begin{equation}\label{8.3:eq-inverse-comparison-6}
\begin{split}
\sum_{o\in O}\|\Gamma^a(o,b)\|^2
&\le\Bigl(\frac{4}{3a_-}\Bigr)^2\sum_{o\in O}e^{-2\mu d(b,o)}\\
&\le\Bigl(\frac{4}{3a_-}\Bigr)^2S(\mu)\,e^{-\mu d(b,O)}.
\end{split}
\end{equation}

\emph{The terms $\Pi^a$.} Let $u$ and $v$ be unit vectors in the component spaces of $E_{\mathfrak{B}^a}$ at $b$ and at $b'$. Let $\Gamma^a_{Ob}u\in E_O$ be the vector with components $\Gamma^a(o,b)u$, $o\in O$, and define $\Gamma^a_{Ob'}v$ in the same way. Since $\Gamma^a(b,o)$ is the transpose of $\Gamma^a(o,b)$,
\begin{equation*}
\langle u,\Pi^a(b,b')v\rangle=\langle\Gamma^a_{Ob}u,(\Gamma^a_{OO})^{-1}\Gamma^a_{Ob'}v\rangle.
\end{equation*}
Moreover $|\Gamma^a_{Ob}u|^2\le\sum_{o\in O}\|\Gamma^a(o,b)\|^2$, and likewise for $\Gamma^a_{Ob'}v$. By the Cauchy--Schwarz inequality, \eqref{8.3:eq-inverse-comparison-5} and \eqref{8.3:eq-inverse-comparison-6},
\begin{equation*}
\|\Pi^a(b,b')\|\le a_+S(c)\Bigl(\frac{4}{3a_-}\Bigr)^2S(\mu)\,e^{-\frac{\mu}{2}(d(b,O^a)+d(b',O^a))}.
\end{equation*}
By definition of $d_\partial$ we have $d(b,O^a)\ge d_\partial(b)$ and $d(b',O^a)\ge d_\partial(b')$. Applying this for $a=1$ and for $a=2$ gives
\begin{equation}\label{8.3:eq-inverse-comparison-7}
\|\Pi^1(b,b')\|+\|\Pi^2(b,b')\|
\le\Bigl(\frac{4}{3a_-}\Bigr)^2S(c)S(\mu)\,2a_+e^{-\frac{\mu}{2}(d_\partial(b)+d_\partial(b'))}.
\end{equation}

\emph{The difference of the inverses of the principal blocks.} Multiplying out the right-hand side shows
\begin{equation*}
(A^1_{JJ})^{-1}-(A^2_{JJ})^{-1}=(A^1_{JJ})^{-1}(A^2-A^1)_{JJ}(A^2_{JJ})^{-1}.
\end{equation*}
By Lemma~\ref{8.3:lem-weighted-inverse-decay}\,(iii), the bound of part~(ii) of that lemma holds for the principal submatrices $A^a_{JJ}$. Thus $\|(A^a_{JJ})^{-1}(x,x')\|\le\tfrac{4}{3a_-}e^{-\mu d(x,x')}$ for $x,x'\in J$. With \eqref{8.3:eq-inverse-comparison-1}, the $(b,b')$-block is therefore bounded by
\begin{equation*}
\Bigl(\frac{4}{3a_-}\Bigr)^2\eta_0\sum_{x,x'\in J}
e^{-\mu d(b,x)-\mu d_\partial(x)}\,e^{-c\,d(x,x')}\,e^{-\mu d(x',b')}.
\end{equation*}

Next, $d(b,x)+d_\partial(x)\ge d_\partial(b)$. Indeed, let $a'$ attain the minimum defining $d_\partial(x)$. Then $d(b,o)\le d(b,x)+d(x,o)$ for every $o\in O^{a'}$, and so
\begin{equation*}
d_\partial(b)\le d(b,O^{a'})\le d(b,x)+d_\partial(x).
\end{equation*}
Hence the first exponential is at most $e^{-\mu d_\partial(b)}$. Summing first over $x$ and then over $x'$ gives
\begin{equation*}
\sum_{x'\in J}e^{-\mu d(x',b')}\sum_{x\in J}e^{-c\,d(x,x')}\le S(c)S(\mu).
\end{equation*}
Therefore
\begin{equation}\label{8.3:eq-inverse-comparison-8}
\bigl\|\bigl[(A^1_{JJ})^{-1}-(A^2_{JJ})^{-1}\bigr](b,b')\bigr\|
\le\Bigl(\frac{4}{3a_-}\Bigr)^2\eta_0\,e^{-\mu d_\partial(b)}S(c)S(\mu).
\end{equation}

\emph{Conclusion.} Take norms in \eqref{8.3:eq-inverse-comparison-4}. The triangle inequality together with \eqref{8.3:eq-inverse-comparison-7} and \eqref{8.3:eq-inverse-comparison-8} gives \eqref{8.3:eq-inverse-comparison-2}.
\end{proof}

\begin{corollary}[The extracted covariance as Ba{\l}aban's fluctuation covariance]
\label{8.3:cor-extracted-covariance}
Fix a level $j$ with $K-s\le j\le K$ and a clean branch of the Gaussian extraction, with label
$\mathsf{e}_0$. Let
$\mathfrak{B}_D$ be
the clean enlarged box of the branch at step $j$, reaching $D$ blocks of side $L$, that
is, to distance $DL$, beyond its core set $\mathfrak{Y}$, let $I$ be the set of free bonds of
$\mathfrak{B}_D$, and put $J:=I$.
A \emph{datum} is a nest of domains
together with a background. The datum of the branch consists of the nest of the branch and the
background $\mathfrak{u}(\mathsf{e}_0)$ of the branch; it is a datum of the complete model in the
sense of the operator theorem at free current, the kind of datum a branch of the extraction
carries. A
\emph{pure} datum has the trivial nest, all of whose domains are the whole lattice. A clean
enlarged box has, in particular, the following three properties, and its cleanness is used here
only through them: the nest of the branch coincides with the trivial nest on a neighbourhood of
$J$; $J$ is contained in the set $\mathfrak{B}^1$ of free bonds of the $(j+1)$-th domain of that
nest; and every bond of the lattice of level $j$ at distance less than $DL$ from $\mathfrak{Y}$
lies in $J$. The forms compared are:
\begin{itemize}
\item $A^1$, the form $\mathbf{C}^*\Delta^{(j)}\mathbf{C}$ on $\mathfrak{B}^1$, built on the nest
of the branch and on the background $\mathfrak{u}(\mathsf{e}_0)$, and $\Gamma^1:=(A^1)^{-1}$;
\item $A^2$, the form $\mathbf{C}^*\Delta^{(j)}\mathbf{C}$ of the pure datum with
background $U$, on the set $\mathfrak{B}^2$ of all bonds of the whole lattice of level $j$, and
$\Gamma^2:=(A^2)^{-1}$, which is
Ba{\l}aban's fluctuation covariance $C^{(j)}(U)$ of \cite[(2.13)]{B-CMP109}.
\end{itemize}
Here $\mathbf{C}$ and $\Delta^{(j)}$ are the operators of \cite[(2.13)]{B-CMP109} from which
Ba{\l}aban's fluctuation covariance is built. The sets $\mathfrak{B}^1,\mathfrak{B}^2$ carry one
metric $d$, and $J\subset\mathfrak{B}^1\cap\mathfrak{B}^2$.
For a covariance $\Gamma$ on bond variables $y$, write $\operatorname{Cov}_{\Gamma}$ for the
covariance under the centred Gaussian measure with covariance $\Gamma$. The colour-averaged
plaquette covariances are
\begin{align*}
\mathfrak{c}^{(j)}_{\mathrm{br}}(q,q')
 &:= n_{\mathfrak{g}}^{-1}\operatorname{tr}_{\mathfrak{g}}
 \operatorname{Cov}_{\Gamma^1}\bigl((\mathrm{d}_{\mathfrak{u}(\mathsf{e}_0)}\mathbf{C}y)(q),
 (\mathrm{d}_{\mathfrak{u}(\mathsf{e}_0)}\mathbf{C}y)(q')\bigr),\\
\mathfrak{c}^{(j)}(q,q';U)
 &:= n_{\mathfrak{g}}^{-1}\operatorname{tr}_{\mathfrak{g}}
 \operatorname{Cov}_{\Gamma^2}\bigl((\mathrm{d}_U\mathbf{C}y)(q),(\mathrm{d}_U\mathbf{C}y)(q')\bigr).
\end{align*}
Here $\mathbf{C}$ and the covariant derivative $\mathrm{d}$ are taken at the background
$\mathfrak{u}(\mathsf{e}_0)$ in the first line and at $U$ in the second, as the subscripts of
$\mathrm{d}$ indicate;
$\operatorname{tr}_{\mathfrak{g}}$ is the trace over the Lie algebra $\mathfrak{g}$ of
$\mathrm{SU}(N)$, and $n_{\mathfrak{g}}$ is the dimension of $\mathfrak{g}$, so that
$n_{\mathfrak{g}}^{-1}\operatorname{tr}_{\mathfrak{g}}$ is the colour average.

Both forms are real symmetric; let $c_A>0$ be the rate of hypothesis (a) below, and let
$a_-,a_+>0$ be such that $A^a\ge a_-$ and
$\|A^a(b,b')\|\le a_+e^{-c_A d(b,b')}$ for $a=1,2$, the bounds on $A^2$ holding with these same
constants at every background $U$ of the class of \cite[(2.5)--(2.11)]{B-CMP109}; the forms are
thus assumed to decay at the
same rate $c_A$ as in (a). Let $\mu_1\le c_A/2$ be a rate
admissible for
$\mu$ in Lemma~\ref{8.3:lem-inverse-comparison}. Let $d_\partial$ and $S(\cdot)$ be as in that
lemma, for the domains $\mathfrak{B}^1,\mathfrak{B}^2$ and for $J$. Let $\ell_{\mathrm{pl}}>0$, with
$\ell_{\mathrm{pl}}+2L<DL$, be the
distance from the core set $\mathfrak{Y}$ within which the compared plaquettes lie,
$\epsilon_0>0$ the amplitude of the allowed discrepancy between backgrounds, and
$\|\mathbf{C}\|$ a number with $\sup_b\sum_{b'}\|\mathbf{C}(b,b')\|\le\|\mathbf{C}\|$ for
$\mathbf{C}$ taken at $\mathfrak{u}(\mathsf{e}_0)$ and at every background of that class.
Finally, $C'$ is a constant with the
following property, which enters here as a hypothesis. Let $\Gamma=A^{-1}$, where $A$ is a real
symmetric form on $\mathfrak{B}^2$ with $A\ge a_-$ and $\|A(b,b')\|\le a_+e^{-c_Ad(b,b')}$, and
let $\mathfrak{u},\mathfrak{u}'$ be two backgrounds, each either the background
$\mathfrak{u}(\mathsf{e}_0)$ of the branch or a background of the class of
\cite[(2.5)--(2.11)]{B-CMP109}, with
$|\mathfrak{u}-\mathfrak{u}'|(z)\le\epsilon_0e^{-\mu_1d_\partial(z)}$ at $z\in J$. Then, for
plaquettes $q,q'$ within $\ell_{\mathrm{pl}}$ of $\mathfrak{Y}$,
\begin{equation*}
\begin{split}
&\Bigl|n_{\mathfrak{g}}^{-1}\operatorname{tr}_{\mathfrak{g}}
\operatorname{Cov}_{\Gamma}\bigl((\mathrm{d}_{\mathfrak{u}}\mathbf{C}y)(q),
(\mathrm{d}_{\mathfrak{u}}\mathbf{C}y)(q')\bigr)
-n_{\mathfrak{g}}^{-1}\operatorname{tr}_{\mathfrak{g}}
\operatorname{Cov}_{\Gamma}\bigl((\mathrm{d}_{\mathfrak{u}'}\mathbf{C}y)(q),
(\mathrm{d}_{\mathfrak{u}'}\mathbf{C}y)(q')\bigr)\Bigr|\\
&\qquad\le C'\epsilon_0e^{-\mu_1(DL-\ell_{\mathrm{pl}})},
\end{split}
\end{equation*}
where $\mathbf{C}$ and $\mathrm{d}$ are taken at $\mathfrak{u}$ in the first covariance and at
$\mathfrak{u}'$ in the second. The datum of the branch, and the pure datum with any background of
the class of \cite[(2.5)--(2.11)]{B-CMP109}, are admissible in the sense of the operator theorem
at free current; this also enters here as a hypothesis. The numbers $a_\pm$, $c_A$, $\mu_1$,
$S(c_A)$, $S(\mu_1)$,
$\|\mathbf{C}\|$, $\ell_{\mathrm{pl}}$ and $C'$, and the constants $C_{\mathrm{geo}}$, $C_\eta$,
$c_\eta$
of (a) and (b) below, are taken independent of $D$ and $M$.

Besides Lemma~\ref{8.3:lem-inverse-comparison}, the corollary assumes the following.
\begin{itemize}
\item[(a)] \emph{Geometric comparison of the fluctuation forms} (read from the operator theorem at
free current; it enters here as a hypothesis). Consider two data, admissible in the sense of that
theorem, each of the
complete model or pure, whose nests coincide on a neighbourhood of $J$ and whose backgrounds
satisfy $|\mathfrak{u}^1-\mathfrak{u}^2|\le\epsilon(x)$ at $x\in J$, modulo a common gauge
transformation. Then their fluctuation forms, again written $A^1,A^2$ and defined on the free
bonds $\mathfrak{B}^1,\mathfrak{B}^2$ of the $(j+1)$-th domains of the two nests, satisfy, for
$x,x'\in J$ and with a constant
$C_{\mathrm{geo}}$ supplied by that theorem,
\begin{equation}\label{8.3:eq-extracted-covariance-1}
\|(A^1-A^2)(x,x')\|
\le C_{\mathrm{geo}}e^{-c_Ad(x,x')}\Bigl[e^{-c_Ad_\partial(x)}
 +\sup_{z}\epsilon(z)e^{-c_Ad(z,x)}\Bigr].
\end{equation}
\item[(b)] \emph{Localisation of the backgrounds near the box.} Let $C_\eta>0$ and $c_\eta>0$ be
the constants supplied by the localisation statement for the backgrounds near the box. Say that
the backgrounds of a configuration $\Phi$ of the extraction are \emph{flat on $J$ up to
the localisation tails} if the background $\mathfrak{u}(\mathsf{e}_0)$ of the branch at
$\Phi$ agrees on $J$ with the flat background $U\equiv\mathbf{1}$ within
$\epsilon(z)\le C_\eta e^{-c_\eta M}e^{-\mu_1d_\partial(z)}$; that is, if the condition on the
background in (1) below holds with $U\equiv\mathbf{1}$ and with
$\epsilon_0:=C_\eta e^{-c_\eta M}$.
\item[(c)] \emph{The dictionary between Ba{\l}aban's flat covariances and the free tower} (taken
from the dictionary identifying Ba{\l}aban's flat step covariances with the increments of the free
tower; it enters here as a hypothesis).
Let $k_0$ be the starting level of the free tower (in the normalisation of clause~(iii)(c) of
Theorem~\ref{3.1:thm-main}, $k_0=K-s$), its increments $\mathfrak{c}_i$ being
indexed by $i=0,1,\dots,s$ from that level. The covariance
$\mathfrak{c}^{(j)}(\cdot,\cdot;\mathbf{1})$,
composed with the parallel transports to level $K-s$, equals $\mathfrak{c}_{j-k_0}(\cdot,\cdot)$,
the increment of the free tower at step $j$.
\end{itemize}
Then the following hold.
\begin{itemize}
\item[(1)] Let $U$ be any global small-field background of the class
\cite[(2.5)--(2.11)]{B-CMP109} that
agrees with $\mathfrak{u}(\mathsf{e}_0)$ on $J$ within $\epsilon(z)$, where
$\epsilon(z)\le\epsilon_0e^{-\mu_1d_\partial(z)}$. Then for plaquettes $q,q'$ within
$\ell_{\mathrm{pl}}$
of $\mathfrak{Y}$,
\begin{equation}\label{8.3:eq-extracted-covariance-2}
\begin{split}
\bigl|\mathfrak{c}^{(j)}_{\mathrm{br}}(q,q')-\mathfrak{c}^{(j)}(q,q';U)\bigr|
&\le 16\|\mathbf{C}\|^2\Bigl(\frac{4}{3a_-}\Bigr)^{2}S(c_A)S(\mu_1)
 \bigl[2a_++C_{\mathrm{geo}}(1+\epsilon_0)\bigr]e^{-\mu_1(DL-\ell_{\mathrm{pl}}-2L)}\\
&\quad+C'\epsilon_0e^{-\mu_1(DL-\ell_{\mathrm{pl}})}.
\end{split}
\end{equation}
\item[(2)] Suppose that the branch is taken at a configuration $\Phi$ as in
\textup{(b)}, and take $\epsilon_0:=C_\eta e^{-c_\eta M}$ in the standing assumptions. Then, for
plaquettes $q,q'$ within $\ell_{\mathrm{pl}}$ of $\mathfrak{Y}$,
\[
\mathfrak{c}^{(j)}_{\mathrm{br}}(q,q')
=\mathfrak{c}^{(j)}(q,q';\mathbf{1})+O\bigl(e^{-\mu_1DL}+e^{-c_\eta M}\bigr);
\]
more precisely, the difference is bounded by the right side of
\eqref{8.3:eq-extracted-covariance-2} evaluated at this $\epsilon_0$.
\item[(3)] The increments $\mathfrak{c}_i$ sum over the levels to the covariance $G^{(s)}_K$ of
clause~(iii)(c) of Theorem~\ref{3.1:thm-main}. Hence the main term of the Gaussian extraction,
\begin{equation}\label{8.3:eq-extracted-covariance-3}
2n_{\mathfrak{g}}\sum_{i,i'<r}g^2g^2\operatorname{Tr}\bigl(H^{F_1}\mathfrak{c}_iH^{F_2}
\mathfrak{c}_{i'}\bigr),
\end{equation}
is the $r$-truncation of $b_0g^4\mathcal{K}_s$ with no change of covariance: it is the sum over
$i,i'<r$ of the terms of an expression, formed with the same increments $\mathfrak{c}_i$, whose
sum over all $0\le i,i'\le s$ equals $b_0g^4\mathcal{K}_s(f_1,f_2)$. Here $r=r(s)$ is the
extraction depth, $F_a=\mathcal{O}(f_a)$ \textup{(}$a=1,2$\textup{)} are the two observables
whose correlation the Gaussian extraction computes, $H^{F_a}$ is the kernel of the quadratic
term of $F_a$ in the Gaussian extraction, $\operatorname{Tr}$ denotes the trace over
plaquette indices of the product of kernels, the factor $g^2g^2=g^4$ carries one factor $g^2$ of
the reference coupling for each of the two covariances, and $\mathcal{K}_s=\mathcal{K}_s(f_1,f_2)$.
\end{itemize}
\end{corollary}

\begin{proof}
\emph{The hypothesis of Lemma~\ref{8.3:lem-inverse-comparison}.} Apply hypothesis (a) to two
data: the datum of the branch, with background $\mathfrak{u}(\mathsf{e}_0)$, and the pure datum
with background $U$. The first is a datum of the complete model and the second is pure, and both
are admissible in the sense of (a) by the standing assumption on admissibility stated before (a).
Since the enlarged box is clean, the nest of the branch
coincides with the
trivial nest of the pure datum on a neighbourhood of $J$, so the two nests coincide there as (a)
requires. By assumption their backgrounds differ on $J$ by at most $\epsilon(z)$. We now bound the
bracket of \eqref{8.3:eq-extracted-covariance-1}.

Since $\mu_1\le c_A/2\le c_A$, we have $e^{-c_Ad_\partial(x)}\le e^{-\mu_1d_\partial(x)}$. Each
function $b\mapsto d(b,\mathfrak{B}^a\setminus J)$ is $1$-Lipschitz, and so is their minimum
$d_\partial$. Hence $d_\partial(x)\le d_\partial(z)+d(z,x)$, and for $z\in J$
\[
\epsilon(z)e^{-c_Ad(z,x)}\le\epsilon_0e^{-\mu_1d_\partial(z)}e^{-\mu_1d(z,x)}
\le\epsilon_0e^{-\mu_1d_\partial(x)}.
\]
So the bracket is at most $(1+\epsilon_0)e^{-\mu_1d_\partial(x)}$. By
\eqref{8.3:eq-extracted-covariance-1}, for $x,x'\in J$,
\[
\|(A^1-A^2)(x,x')\|
\le C_{\mathrm{geo}}(1+\epsilon_0)e^{-c_Ad(x,x')}e^{-\mu_1d_\partial(x)}.
\]
This is the hypothesis of Lemma~\ref{8.3:lem-inverse-comparison} with the lemma's decay rate equal
to $c_A$, with $\mu=\mu_1$ and
$\eta_0=C_{\mathrm{geo}}(1+\epsilon_0)$. The remaining hypotheses of that lemma, namely symmetry,
the lower bound $a_-$ and the decay $a_+e^{-c_Ad}$, are the standing conditions on $a_\pm$, and the
inclusion $J\subset\mathfrak{B}^1\cap\mathfrak{B}^2$, with one metric $d$ on both domains, holds
because $J\subset\mathfrak{B}^1$ by the cleanness of the enlarged box and $\mathfrak{B}^2$ consists
of all bonds of level $j$. The
lemma gives, for $b,b'\in J$,
\begin{equation}\label{8.3:eq-extracted-covariance-4}
\begin{split}
\|\Gamma^1(b,b')-\Gamma^2(b,b')\|&\le\Bigl(\frac{4}{3a_-}\Bigr)^{2}S(c_A)S(\mu_1)\\
&\quad\times\Bigl[2a_+e^{-\frac{\mu_1}{2}(d_\partial(b)+d_\partial(b'))}
+C_{\mathrm{geo}}(1+\epsilon_0)e^{-\mu_1d_\partial(b)}\Bigr].
\end{split}
\end{equation}

\emph{From bonds to plaquettes.} The variable $(\mathrm{d}\mathbf{C}y)(q)$ is a sum of four terms,
one for each bond of $\partial q$. Its covariance with $(\mathrm{d}\mathbf{C}y)(q')$ is therefore
a sum of at most $4\cdot4=16$ blocks of $\mathbf{C}\Gamma\mathbf{C}^*$. Each of the four terms is
$\pm$ a parallel transport, of norm one, of $(\mathbf{C}y)(b_1)$ with $b_1\in\partial q$, and
$|n_{\mathfrak{g}}^{-1}\operatorname{tr}_{\mathfrak{g}}X|\le\|X\|$ for every linear map $X$ of
$\mathfrak{g}$; so each block, with $b_1\in\partial q$ and $b_2\in\partial q'$, contributes to the
difference of the colour-averaged covariances at
most $\|(\mathbf{C}(\Gamma^1-\Gamma^2)\mathbf{C}^*)(b_1,b_2)\|$. The operator $\mathbf{C}$
reads only bonds within distance $2L$ of the bond on which it is evaluated. Hence a plaquette $q$
within $\ell_{\mathrm{pl}}$ of $\mathfrak{Y}$ reads only bonds $b$ within $\ell_{\mathrm{pl}}+2L$ of
$\mathfrak{Y}$.
By the third property of the clean enlarged box, every such bond $b$, being at distance at most
$\ell_{\mathrm{pl}}+2L<DL$ from $\mathfrak{Y}$, lies in $J$, and every bond of
$\mathfrak{B}^1\setminus J$ or of $\mathfrak{B}^2\setminus J$ lies at distance at least $DL$ from
$\mathfrak{Y}$; by the triangle inequality, $d_\partial(b)\ge DL-\ell_{\mathrm{pl}}-2L$. Thus, for
the bonds read by
$q$ and $q'$, both terms of the bracket in \eqref{8.3:eq-extracted-covariance-4} are bounded as
follows:
\begin{gather*}
2a_+e^{-\frac{\mu_1}{2}(d_\partial(b)+d_\partial(b'))}\le2a_+e^{-\mu_1(DL-\ell_{\mathrm{pl}}-2L)},\\
C_{\mathrm{geo}}(1+\epsilon_0)e^{-\mu_1d_\partial(b)}
\le C_{\mathrm{geo}}(1+\epsilon_0)e^{-\mu_1(DL-\ell_{\mathrm{pl}}-2L)}.
\end{gather*}
For a kernel $\Xi$ on bonds, the definition of $\|\mathbf{C}\|$ gives
\[
\|(\mathbf{C}\Xi\mathbf{C}^*)(b_1,b_2)\|
\le\sum_{b,b'}\|\mathbf{C}(b_1,b)\|\,\|\Xi(b,b')\|\,\|\mathbf{C}(b_2,b')\|
\le\|\mathbf{C}\|^2\sup_{b,b'}\|\Xi(b,b')\|,
\]
the supremum being over the bonds read. Summing the
sixteen blocks with $\Xi=\Gamma^1-\Gamma^2$, while keeping $\mathbf{C}$ and $\mathrm{d}$ at the
background $\mathfrak{u}(\mathsf{e}_0)$, gives the first member of the right side of
\eqref{8.3:eq-extracted-covariance-2}.

\emph{Variation of $\mathbf{C}$ and of the covariant derivative.} The covariance
$\mathfrak{c}^{(j)}_{\mathrm{br}}$ is formed with $\mathbf{C}$ and $\mathrm{d}$ at
$\mathfrak{u}(\mathsf{e}_0)$, while $\mathfrak{c}^{(j)}(\cdot;U)$ is formed with them at $U$. On
$J$ the two backgrounds differ by at most
$\epsilon(z)\le\epsilon_0e^{-\mu_1d_\partial(z)}$, so the defining property of $C'$, applied with
$\Gamma=\Gamma^2$, which is of the kind admitted there because $A^2$ satisfies the standing bounds
at the background $U$ of the class of \cite[(2.5)--(2.11)]{B-CMP109},
$\mathfrak{u}=\mathfrak{u}(\mathsf{e}_0)$ and $\mathfrak{u}'=U$, shows that the
variation of $\mathbf{C}$ and $\mathrm{d}_U$ themselves with the background contributes the second
member, $C'\epsilon_0e^{-\mu_1(DL-\ell_{\mathrm{pl}})}$. The
triangle inequality, applied to the two differences, gives \eqref{8.3:eq-extracted-covariance-2}.
This proves (1).

\emph{Flat backgrounds.} The flat background $U\equiv\mathbf{1}$ is a global small-field
background of the class of \cite[(2.5)--(2.11)]{B-CMP109}, its plaquette variables being all equal
to $\mathbf{1}$. If the branch is taken at a configuration $\Phi$ as in hypothesis (b),
then by (b) this background agrees with $\mathfrak{u}(\mathsf{e}_0)$ on $J$ within
$\epsilon(z)\le\epsilon_0e^{-\mu_1d_\partial(z)}$, with $\epsilon_0=C_\eta e^{-c_\eta M}$.
Hence (1) applies with $U\equiv\mathbf{1}$ and gives \eqref{8.3:eq-extracted-covariance-2} for
$\mathfrak{c}^{(j)}(q,q';\mathbf{1})$, which is the second assertion of (2). Its first member equals
$C''e^{-\mu_1DL}$, where
\[
\begin{split}
C''&:=16\|\mathbf{C}\|^2\Bigl(\frac{4}{3a_-}\Bigr)^{2}S(c_A)S(\mu_1)\\
&\quad\times\bigl[2a_++C_{\mathrm{geo}}(1+\epsilon_0)\bigr]e^{\mu_1(\ell_{\mathrm{pl}}+2L)},
\end{split}
\]
and $C''$ is bounded independently of $D$ and $M$, because the quantities entering it other than
$\epsilon_0$ are taken independent of $D$ and $M$, and $\epsilon_0=C_\eta e^{-c_\eta M}\le C_\eta$.
Its second member is at most $C'C_\eta e^{-c_\eta M}$, because
$e^{-\mu_1(DL-\ell_{\mathrm{pl}})}\le1$ by
$\ell_{\mathrm{pl}}<DL$. The difference is therefore $O(e^{-\mu_1DL}+e^{-c_\eta M})$. This
proves (2).

\emph{Identification with the free tower.} By hypothesis (c),
$\mathfrak{c}^{(j)}(\cdot,\cdot;\mathbf{1})$, transported to level $K-s$, is the increment
$\mathfrak{c}_{j-k_0}$ of the free tower. By its definition in clause~(iii)(c) of
Theorem~\ref{3.1:thm-main}, the covariance $G^{(s)}_K$ is the sum of Ba{\l}aban's fluctuation
covariances of the levels $K-s,\dots,K$ at flat background, transported to level $K-s$; these are
the covariances $C^{(k)}$ of \cite[(2.13)]{B-CMP109} at $U\equiv\mathbf{1}$. Hence the increments
$\mathfrak{c}_i$, $i=j-k_0$ with $j=K-s,\dots,K$, sum to $G^{(s)}_K$.

The kernel $\mathcal{K}_s$ of the two-point witness theorem is formed with the covariance
$G^{(s)}_K$. The main term \eqref{8.3:eq-extracted-covariance-3} of the Gaussian extraction is
formed with the increments $\mathfrak{c}_i$, $\mathfrak{c}_{i'}$ for $i,i'<r$. By the Gaussian
extraction, the same expression summed over all $0\le i,i'\le s$ is
$b_0g^4\mathcal{K}_s(f_1,f_2)$; restricted to $i,i'<r$, it is therefore its $r$-truncation,
formed with the same covariance: what separates the two is
only the levels $i\ge r$, not any change of covariance. This proves (3).

The levels $i\ge r$ of the kernel are not part of \eqref{8.3:eq-extracted-covariance-3}. They are
estimated within the two-point witness theorem, by the second-derivative room lemma and by the
witness theorem's tail bound on the kernel. The dependence on $\Phi$ through non-flat
backgrounds is the Lipschitz member of the Gaussian extraction, and the corollary does not alter
it.
\end{proof}

\begin{definition}[The decoupled family and the coupling observable]\label{8.3:def-decoupled-family}
\emph{The setting $\mathfrak{S}$: weakly perturbed Gaussian measures with convex windows at each site.}
The data are the following.
\begin{itemize}
\item A finite set $\mathfrak{X}$ of \emph{sites} with a metric $\operatorname{dist}$ and constants $C_d$
and $d$ such that, for every $x\in\mathfrak{X}$ and every $a\ge 0$,
\begin{equation*}
\#\{x'\in\mathfrak{X}:\operatorname{dist}(x,x')\le a\}\le C_d(1+a)^d .
\end{equation*}
In the application the sites are the free bonds of the box $I$ of the decoupled family below, with the
lattice distance of the level $j$ under consideration, and $d=4$.
\item A finite-dimensional real vector space $\mathfrak{g}$ with a Euclidean norm $|\cdot|$, one copy at
each site; configurations are $y=(y_x)_{x\in\mathfrak{X}}\in\mathfrak{g}^{\mathfrak{X}}$.
\item Radii $r_x\in[r_-,r_+]$ and a number $\delta\in(0,1]$; for each site an open convex set
$K_x\subset\mathfrak{g}$ with $0\in K_x\subset\{|y_x|<r_x\}$ (in the application, the norm ball of
Ba{\l}aban's construction at level $j$, and $r_x$ its Euclidean out-radius), and
$K_x^+:=K_x+\{|\cdot|<\delta r_-\}$. We put $\mathcal{W}:=\prod_x K_x$ and
$\mathcal{W}^+:=\prod_x K_x^+$; both are open and convex.
\item A real symmetric operator $\mathsf{Q}=(\mathsf{Q}_{xx'})$ on $\mathfrak{g}^{\mathfrak{X}}$, with blocks
$\mathsf{Q}_{xx'}$ acting on $\mathfrak{g}$, such that $q_-\mathbf{1}\le\mathsf{Q}\le q_+\mathbf{1}$ and
$\|\mathsf{Q}_{xx'}\|\le q_+e^{-c_{\mathsf{Q}}\operatorname{dist}(x,x')}$.
\item Optionally, a \emph{frozen exterior}: a finite set $F$ of further sites, disjoint from
$\mathfrak{X}$, with $\operatorname{dist}$ extended to $\mathfrak{X}\sqcup F$ obeying the same growth bound,
and frozen values $\eta=(\eta_f)_{f\in F}$ ranging in
$\mathcal{W}_F:=\prod_{f\in F}\{|\eta_f|<r_f\}$, $r_f\le r_+$. In the application these are the far
fluctuation variables $y_F$, held fixed.
\item A vector $v=\mathring{v}+\mathsf{R}\eta\in\mathfrak{g}^{\mathfrak{X}}$ with $\mathring{v}$ arbitrary and
$\|\mathsf{R}_{xf}\|\le q_+e^{-c_{\mathsf{Q}}\operatorname{dist}(x,f)}$.
\item A function $P=P(y,\eta)\in C^2(\mathcal{W}^+\times\mathcal{W}_F;\mathbb{R})$ (of class $C^\infty$ in the
application) such that, uniformly in $\eta\in\mathcal{W}_F$, on $\mathcal{W}^+$,
\begin{equation*}
\|\partial_xP\|\le p_1,\qquad
\|\partial_x\partial_{x'}P\|\le p_2e^{-c_P\operatorname{dist}(x,x')}\quad(x=x'\text{ included}),
\end{equation*}
and $\|\partial_x\partial_{\eta_f}P\|\le p_2e^{-c_P\operatorname{dist}(x,f)}$ for $x\in\mathfrak{X}$, $f\in F$.
\end{itemize}
With $F=\emptyset$ this is the plain setting. Every statement made in this setting is uniform in $\eta$,
which is a parameter and is never integrated. We put
$S_d(c):=\sup_x\sum_{x'}e^{-c\operatorname{dist}(x,x')}$, so that
$S_d(c)\le C_d\sum_{a\ge0}(1+a)^de^{-ca}$, and we assume the smallness condition
\begin{equation*}
p_2\,S_d(c_P)\le q_-/4 .\tag{sm}
\end{equation*}
Condition \textup{(sm)} constrains $p_2$ only; in the estimates of this section the second derivative of
the retraction of Lemma~\ref{8.3:lem-retraction-wall} is absorbed by the convex wall through property (e4)
of that lemma. For each frozen $\eta$ the measure is
\begin{equation}\label{8.3:eq-decoupled-family-1}
\mu(dy):=\mathcal{Z}^{-1}\mathbf{1}_{\mathcal{W}}(y)\,e^{-\Phi(y)}\,dy,\qquad
\Phi(y):=\tfrac12\langle y,\mathsf{Q}y\rangle+\langle v,y\rangle-P(y,\eta),
\end{equation}
with $\mathcal{Z}$ the normalisation; $\mathfrak{S}$ denotes the class of these measures. The
\emph{constants of the class} are $m_{\mathfrak{S}}:=q_-/2$,
$\kappa_{\mathfrak{S}}(x,x'):=q_+e^{-c_{\mathsf{Q}}\operatorname{dist}(x,x')}+p_2e^{-c_P\operatorname{dist}(x,x')}$
for $x\neq x'$, $\nu_0:=\tfrac12\min(c_{\mathsf{Q}},c_P)$,
\begin{equation*}
\mathfrak{K}:=\sup_x\sum_{x'\neq x}\kappa_{\mathfrak{S}}(x,x')e^{\nu_0\operatorname{dist}(x,x')}
\le(q_++p_2)S_d(\nu_0),\qquad
\nu:=\min\Bigl\{\frac{\nu_0}{2},\ \frac{e\,\nu_0\,m_{\mathfrak{S}}}{4\mathfrak{K}}\Bigr\},
\end{equation*}
and $\bar{\mathfrak{K}}:=(q_++p_2)S_d(\nu_0)$, so that $\bar{\mathfrak{K}}\ge\mathfrak{K}$. All of them are
free of $|\mathfrak{X}|$, $|F|$, $v$, $\eta$ and the radii (given \textup{(sm)}).

\emph{Closure of the class.} Let $\mu\in\mathfrak{S}$ and $T\subset\mathfrak{X}$; write
$T^c:=\mathfrak{X}\setminus T$ and $\bar{\mathcal{W}}_T:=\prod_{x\in T}\bar K_x$.
\begin{enumerate}
\item[(a)] For every $y_T\in\bar{\mathcal{W}}_T$ the conditional law of $y':=y|_{T^c}$ given $y_T$,
\begin{equation*}
\mu^{y_T}(dy')\propto\mathbf{1}_{\mathcal{W}_{T^c}}(y')\,e^{-\Phi(y_T,y')}\,dy',
\end{equation*}
is a measure of $\mathfrak{S}$ on the sites $T^c$ with the same $q_\pm$, $c_{\mathsf{Q}}$, $p_1$, $p_2$, $c_P$,
$\delta$, $r_\pm$: its data are $\mathsf{Q}_{T^cT^c}$, $v_{T^c}+\mathsf{Q}_{T^cT}y_T$ and $P(y_T,\cdot,\eta)$,
that is, $T$ joins the frozen exterior. Hence the constants $m_{\mathfrak{S}}$, $\mathfrak{K}$, $\nu$ of $\mu$
serve for $\mu^{y_T}$ as well: its own $m_{\mathfrak{S}}$ is $q_-/2$, its own $\mathfrak{K}$ is at most
$\mathfrak{K}$, and its own $\nu$ is at least $\nu$.
Likewise any $\mathsf{Q}'$ with the same bounds may stand in place of $\mathsf{Q}$, and any
$P'=(1-\sigma)P_1+\sigma P_2$ in place of $P$, where $\sigma\in[0,1]$ and $P_1,P_2$ obey the bounds of $P$.
\end{enumerate}

\emph{The decoupled family.} Fix a partition $\mathfrak{X}=I\sqcup O$ into a \emph{box} $I$ and an
\emph{outside} $O$ (in the application, $O=\emptyset$, the frozen exterior $F$ is the set of all other
free bonds, and $\eta=y_F$), and assume the
separation condition
\begin{equation}\label{8.3:eq-decoupled-family-2}
\operatorname{dist}(x,f)\ge d_O(x):=\operatorname{dist}(x,O)\qquad
\text{for all }x\in I,\ f\in F,
\end{equation}
that is, the frozen sites lie beyond $O$ as seen from $I$. For $\sigma\in[0,1]$ put
\begin{equation}\label{8.3:eq-decoupled-family-3}
\begin{aligned}
\mathsf{Q}^\sigma&:=\mathsf{Q}_{II}\oplus\mathsf{Q}_{OO}+\sigma(\mathsf{Q}_{IO}+\mathsf{Q}_{OI}),\\
v^\sigma&:=\bigl(\mathring{v}_I+\sigma(\mathsf{R}\eta)_I\bigr)\oplus\bigl(\mathring{v}_O+(\mathsf{R}\eta)_O\bigr),\\
P^\sigma(y,\eta)&:=(1-\sigma)\bigl[P(y_I,0_O,0_F)+P(0_I,y_O,\eta)\bigr]+\sigma P(y_I,y_O,\eta),
\end{aligned}
\end{equation}
let $\Phi^\sigma(y):=\tfrac12\langle y,\mathsf{Q}^\sigma y\rangle+\langle v^\sigma,y\rangle-P^\sigma(y,\eta)$, and
let $\mu_\sigma$ be the measure with the data $(\mathsf{Q}^\sigma,v^\sigma,P^\sigma)$ on the same window
$\mathcal{W}$. Then:
\begin{enumerate}
\item[(b)] $\mu_\sigma\in\mathfrak{S}$ with the same constants, for every $\sigma\in[0,1]$;
\item[(c)] $\mu_{\sigma=1}=\mu$;
\item[(d)] $\mu_{\sigma=0}=\mu^{\mathrm{box}}\otimes\mu^{\mathrm{out}}$ factorises, with the \emph{box datum}
\begin{equation}\label{8.3:eq-decoupled-family-4}
\mu^{\mathrm{box}}(dy_I):=(\mathcal{Z}^{\mathrm{box}})^{-1}\mathbf{1}_{\mathcal{W}_I}(y_I)
\exp\Bigl[-\tfrac12\langle y_I,\mathsf{Q}_{II}y_I\rangle-\langle \mathring{v}_I,y_I\rangle
+P(y_I,0_O,0_F)\Bigr]dy_I ,
\end{equation}
which is independent of $\eta$, of $\mathring{v}_O$, of $y_O$ and of the values of $P$ off the subspace
$\{y_O=0,\ \eta=0\}$ (in the application $\mathring{v}=0$, and the box datum is centred and free of the far
field), and where $\mu^{\mathrm{out}}$ is the probability measure on $\mathcal{W}_O$ with density
proportional to
\begin{equation*}
\exp\Bigl[-\tfrac12\langle y_O,\mathsf{Q}_{OO}y_O\rangle
-\langle \mathring{v}_O+(\mathsf{R}\eta)_O,y_O\rangle+P(0_I,y_O,\eta)\Bigr].
\end{equation*}
\end{enumerate}

\emph{The coupling observable} is
\begin{equation}\label{8.3:eq-decoupled-family-5}
\begin{aligned}
\mathfrak{W}:=-\partial_\sigma\Phi^\sigma
={}&-\langle y_I,\mathsf{Q}_{IO}y_O\rangle-\langle(\mathsf{R}\eta)_I,y_I\rangle\\
&+\bigl[P(y_I,y_O,\eta)-P(y_I,0_O,0_F)-P(0_I,y_O,\eta)\bigr],
\end{aligned}
\end{equation}
and it does not depend on $\sigma$.
\end{definition}

\begin{proof}
\emph{Normalisations.} The closure $\bar{\mathcal{W}}=\prod_x\bar K_x$ is compact, being a product of
closed subsets of the balls $\{|y_x|\le r_x\}$, and it is contained in the open set $\mathcal{W}^+$; the
function $P(\cdot,\eta)$ is continuous on $\mathcal{W}^+$, hence bounded on $\bar{\mathcal{W}}$, and the
quadratic and linear terms are bounded there as well. Since $\mathcal{W}$ is open and non-empty, the
normalisation $\mathcal{Z}$ in \eqref{8.3:eq-decoupled-family-1} is finite and positive; the same argument
applies to every normalisation below.

\emph{Proof of (a).} Write $y=(y_T,y')$. Since $\mathsf{Q}$ is symmetric,
$\mathsf{Q}_{TT^c}=\mathsf{Q}_{T^cT}^{\mathsf{T}}$, and
\begin{equation*}
\begin{aligned}
\Phi(y_T,y')={}&\tfrac12\langle y',\mathsf{Q}_{T^cT^c}y'\rangle
+\langle v_{T^c}+\mathsf{Q}_{T^cT}y_T,\ y'\rangle-P(y_T,y',\eta)\\
&+\tfrac12\langle y_T,\mathsf{Q}_{TT}y_T\rangle+\langle v_T,y_T\rangle .
\end{aligned}
\end{equation*}
The last line does not depend on $y'$ and is absorbed in the normalisation. Thus $\mu^{y_T}$ has the form
\eqref{8.3:eq-decoupled-family-1} on the sites $T^c$ with window $\mathcal{W}_{T^c}$, operator
$\mathsf{Q}_{T^cT^c}$, vector $v_{T^c}+\mathsf{Q}_{T^cT}y_T$ and perturbation $P(y_T,\cdot,\eta)$, the frozen
exterior being $F\sqcup T$ with values $(\eta,y_T)$. We check the bounds. For $u\in\mathfrak{g}^{T^c}$,
$\langle u,\mathsf{Q}_{T^cT^c}u\rangle=\langle(0_T,u),\mathsf{Q}(0_T,u)\rangle$ lies between $q_-|u|^2$ and
$q_+|u|^2$, and the blocks of $\mathsf{Q}_{T^cT^c}$ are blocks of $\mathsf{Q}$, so its form bounds and entry
bounds are inherited. The new vector is $\mathring{v}_{T^c}+\mathsf{R}'(\eta,y_T)$ with
$\mathsf{R}'_{xf}=\mathsf{R}_{xf}$ for $f\in F$ and $\mathsf{R}'_{xs}=\mathsf{Q}_{xs}$ for $s\in T$, and
$\|\mathsf{Q}_{xs}\|\le q_+e^{-c_{\mathsf{Q}}\operatorname{dist}(x,s)}$. Since $y_T\in\bar{\mathcal{W}}_T$ and
$\bar K_s\subset K_s^+$, the point $(y_T,y')$ lies in $\mathcal{W}^+$ whenever $y'\in\mathcal{W}^+_{T^c}$, so
$P(y_T,\cdot,\eta)$ is of class $C^2$ there; its derivatives $\partial_x$ and $\partial_x\partial_{x'}$
($x,x'\in T^c$) and $\partial_x\partial_{\eta_f}$ are those of $P$ at $(y_T,y',\eta)$, and its mixed
derivatives $\partial_x\partial_{y_s}$ ($s\in T$) are the derivatives $\partial_x\partial_sP$, bounded by
$p_2e^{-c_P\operatorname{dist}(x,s)}$. The metric on $T^c\sqcup F\sqcup T$ is the given one, so the growth
bound holds, and the radii and $\delta$ are unchanged. Hence $q_\pm$, $c_{\mathsf{Q}}$, $p_1$, $p_2$, $c_P$,
$\delta$, $r_\pm$ are the same. Consequently $m_{\mathfrak{S}}=q_-/2$ and $\nu_0$ are unchanged,
$\bar{\mathfrak{K}}$ does not increase, since $S_d(\nu_0)$ for the sites $T^c$ is a supremum over fewer
sites of sums over fewer sites, and the kernel $\kappa_{\mathfrak{S}}$ on $T^c$ is the restriction of that of
$\mu$; the supremum and the sum defining $\mathfrak{K}$ for $\mu^{y_T}$ run over subsets of those for $\mu$
with non-negative terms, so that value does not exceed $\mathfrak{K}$. Since $\nu$ is non-increasing in
$\mathfrak{K}$ at fixed $\nu_0$ and $m_{\mathfrak{S}}$, the rate $\nu$ of $\mu^{y_T}$ is at least that of $\mu$.

For the last sentence of (a): replacing $\mathsf{Q}$ by $\mathsf{Q}'$ with $q_-\mathbf{1}\le\mathsf{Q}'\le
q_+\mathbf{1}$ and $\|\mathsf{Q}'_{xx'}\|\le q_+e^{-c_{\mathsf{Q}}\operatorname{dist}(x,x')}$ leaves all
requirements of the setting satisfied with the same numbers. If $P_1,P_2$ obey the bounds of $P$ and
$\sigma\in[0,1]$, every derivative $D$ occurring in those bounds satisfies
$\|D((1-\sigma)P_1+\sigma P_2)\|\le(1-\sigma)\|DP_1\|+\sigma\|DP_2\|$, and the right side is at most the
common bound because $(1-\sigma)+\sigma=1$; and $(1-\sigma)P_1+\sigma P_2$ is of class $C^2$ on
$\mathcal{W}^+\times\mathcal{W}_F$ because $P_1$ and $P_2$ are.

\emph{Proof of (b).} First, $\mathsf{Q}^\sigma=(1-\sigma)(\mathsf{Q}_{II}\oplus\mathsf{Q}_{OO})+\sigma\mathsf{Q}$,
since $\mathsf{Q}=\mathsf{Q}_{II}\oplus\mathsf{Q}_{OO}+\mathsf{Q}_{IO}+\mathsf{Q}_{OI}$. The principal blocks
inherit the form bounds, as in the proof of (a):
\begin{equation*}
\langle y,(\mathsf{Q}_{II}\oplus\mathsf{Q}_{OO})y\rangle
=\langle(y_I,0_O),\mathsf{Q}(y_I,0_O)\rangle+\langle(0_I,y_O),\mathsf{Q}(0_I,y_O)\rangle ,
\end{equation*}
which lies between $q_-(|y_I|^2+|y_O|^2)=q_-|y|^2$ and $q_+|y|^2$. So $\mathsf{Q}^\sigma$ is a convex
combination of two symmetric operators between $q_-\mathbf{1}$ and $q_+\mathbf{1}$, hence lies between them;
its blocks are $\mathsf{Q}_{xx'}$, $\sigma\mathsf{Q}_{xx'}$ or $0$, so its entries are bounded by those of
$\mathsf{Q}$. Next, $v^\sigma=\mathring{v}+\mathsf{R}^\sigma\eta$ with $\mathsf{R}^\sigma_{xf}:=\sigma\mathsf{R}_{xf}$
for $x\in I$ and $\mathsf{R}^\sigma_{xf}:=\mathsf{R}_{xf}$ for $x\in O$, and
$\|\mathsf{R}^\sigma_{xf}\|\le\|\mathsf{R}_{xf}\|$.

For $P^\sigma$, put $B(y,\eta):=P(y_I,0_O,0_F)+P(0_I,y_O,\eta)$. The two pieces are $P$ restricted to
coordinate subspaces through $0\in\mathcal{W}\times\mathcal{W}_F$: since $0\in K_x\subset K_x^+$, the points
$(y_I,0_O)$ and $(0_I,y_O)$ lie in $\mathcal{W}^+$ for $y\in\mathcal{W}^+$, and $0_F\in\mathcal{W}_F$, so $B$
is of class $C^2$ on $\mathcal{W}^+\times\mathcal{W}_F$, and the bounds of $P$ hold at the points where the
pieces are evaluated. For $x\in I$, $\partial_xB=(\partial_xP)(y_I,0_O,0_F)$; for $x\in O$,
$\partial_xB=(\partial_xP)(0_I,y_O,\eta)$; in both cases $\|\partial_xB\|\le p_1$. For $x,x'\in I$,
$\partial_x\partial_{x'}B=(\partial_x\partial_{x'}P)(y_I,0_O,0_F)$; for $x,x'\in O$ it equals
$(\partial_x\partial_{x'}P)(0_I,y_O,\eta)$; for $x\in I$, $x'\in O$ or conversely it is $0$; in every case
its norm is at most $p_2e^{-c_P\operatorname{dist}(x,x')}$. Finally $\partial_x\partial_{\eta_f}B=0$ for
$x\in I$, and $\partial_x\partial_{\eta_f}B=(\partial_x\partial_{\eta_f}P)(0_I,y_O,\eta)$ for $x\in O$, of
norm at most $p_2e^{-c_P\operatorname{dist}(x,f)}$. Thus $B$ obeys the bounds of $P$, and
$P^\sigma=(1-\sigma)B+\sigma P$ obeys them by the last sentence of (a). The window $\mathcal{W}$, the radii,
$\delta$ and the metric are unchanged. Hence $\mu_\sigma\in\mathfrak{S}$ with the same $q_\pm$,
$c_{\mathsf{Q}}$, $p_1$, $p_2$, $c_P$, $\delta$, $r_\pm$ on the same set of sites, so
$m_{\mathfrak{S}}$, $\kappa_{\mathfrak{S}}$, $\nu_0$, $\mathfrak{K}$, $\bar{\mathfrak{K}}$ and $\nu$ are the
same; condition \textup{(sm)} involves only $p_2$, $c_P$, $q_-$ and the metric, and therefore continues to hold.

\emph{Proof of (c).} At $\sigma=1$, $\mathsf{Q}^1=\mathsf{Q}_{II}\oplus\mathsf{Q}_{OO}+\mathsf{Q}_{IO}
+\mathsf{Q}_{OI}=\mathsf{Q}$, $v^1=\mathring{v}+\mathsf{R}\eta=v$ and $P^1=P$; hence $\Phi^1=\Phi$ and
$\mu_{\sigma=1}=\mu$.

\emph{Proof of (d).} At $\sigma=0$,
\begin{equation*}
\begin{aligned}
\Phi^0(y)={}&\tfrac12\langle y_I,\mathsf{Q}_{II}y_I\rangle+\langle \mathring{v}_I,y_I\rangle-P(y_I,0_O,0_F)\\
&+\tfrac12\langle y_O,\mathsf{Q}_{OO}y_O\rangle+\langle \mathring{v}_O+(\mathsf{R}\eta)_O,y_O\rangle
-P(0_I,y_O,\eta),
\end{aligned}
\end{equation*}
and $\mathbf{1}_{\mathcal{W}}(y)=\mathbf{1}_{\mathcal{W}_I}(y_I)\mathbf{1}_{\mathcal{W}_O}(y_O)$ because
$\mathcal{W}$ is a product. The density of $\mu_{\sigma=0}$ is therefore the product of a function of $y_I$
and a function of $y_O$, and $\mu_{\sigma=0}=\mu^{\mathrm{box}}\otimes\mu^{\mathrm{out}}$ with
$\mu^{\mathrm{box}}$ as in
\eqref{8.3:eq-decoupled-family-4} and
\begin{equation*}
\mu^{\mathrm{out}}(dy_O)\propto\mathbf{1}_{\mathcal{W}_O}(y_O)
\exp\Bigl[-\tfrac12\langle y_O,\mathsf{Q}_{OO}y_O\rangle-\langle \mathring{v}_O+(\mathsf{R}\eta)_O,y_O\rangle
+P(0_I,y_O,\eta)\Bigr]dy_O ,
\end{equation*}
each factor normalised as explained at the beginning of the proof. The right side of
\eqref{8.3:eq-decoupled-family-4} contains only $\mathsf{Q}_{II}$, $\mathring{v}_I$, the window $\mathcal{W}_I$ and the
values of $P$ at points $(y_I,0_O,0_F)$; it therefore does not involve $\eta$, $\mathring{v}_O$, $y_O$, or the values of
$P$ off the subspace $\{y_O=0,\ \eta=0\}$. When $\mathring{v}=0$, the linear term in
\eqref{8.3:eq-decoupled-family-4} vanishes.

\emph{The coupling observable.} By \eqref{8.3:eq-decoupled-family-3}, $\mathsf{Q}^\sigma$, $v^\sigma$ and
$P^\sigma$ are affine in $\sigma$, hence so is $\Phi^\sigma$, and $\partial_\sigma\Phi^\sigma$ does not depend
on $\sigma$. Term by term,
\begin{equation*}
\partial_\sigma\tfrac12\langle y,\mathsf{Q}^\sigma y\rangle
=\tfrac12\langle y,(\mathsf{Q}_{IO}+\mathsf{Q}_{OI})y\rangle=\langle y_I,\mathsf{Q}_{IO}y_O\rangle ,
\end{equation*}
using $\mathsf{Q}_{OI}=\mathsf{Q}_{IO}^{\mathsf{T}}$; next
$\partial_\sigma\langle v^\sigma,y\rangle=\langle(\mathsf{R}\eta)_I,y_I\rangle$; and
\begin{equation*}
\partial_\sigma P^\sigma(y,\eta)=P(y_I,y_O,\eta)-P(y_I,0_O,0_F)-P(0_I,y_O,\eta).
\end{equation*}
Since $\Phi^\sigma=\tfrac12\langle y,\mathsf{Q}^\sigma y\rangle+\langle v^\sigma,y\rangle-P^\sigma$, changing the
sign of the sum of the first two derivatives and adding the third gives
\eqref{8.3:eq-decoupled-family-5}.
\end{proof}

In what follows, $\mathfrak{S}$ denotes the class of windowed Gibbs measures defined in this section. We list the data of a measure $\mu\in\mathfrak{S}$ and the properties of the class that Lemma~\ref{8.3:lem-correlation-decay} and its proof use, with the notation for them.
\begin{itemize}
\item[(i)] Sites. $\mathfrak{X}$ is a finite set with a metric $\operatorname{dist}$. A configuration is $y=(y_x)_{x\in\mathfrak{X}}\in\mathbb{R}^{n|\mathfrak{X}|}$ with $y_x\in\mathbb{R}^n$, the single-site space. For $S\subset\mathfrak{X}$ we write $y_S:=(y_x)_{x\in S}$ and $S^c:=\mathfrak{X}\setminus S$. The symbol $\partial_x$ denotes the gradient in the variable $y_x$.
\item[(ii)] Window. For each $x\in\mathfrak{X}$, $K_x\subset\mathbb{R}^n$ is an open bounded convex set with $\operatorname{diam}K_x\le 2(1+\delta)r_x$, and $K_x^+\subset\mathbb{R}^n$ is an open bounded set with $\overline{K_x+B(0,\delta r_-/2)}\subset K_x^+$. Here $\delta>0$, $r_x>0$ and $r_->0$ are the numbers of the class that bound the diameter of $K_x$ and fix the width of the enlargement $K_x^+$. We write
\begin{equation*}
\mathcal{K}:=\prod_{x\in\mathfrak{X}}K_x,\qquad
\mathcal{K}^+:=\prod_{x\in\mathfrak{X}}K_x^+,\qquad
\mathcal{K}_S:=\prod_{x\in S}K_x .
\end{equation*}
\item[(iii)] Density. The measure is $\mu(dy)=\mathbf{1}_{\overline{\mathcal{K}}}(y)e^{-\Phi(y)}\,dy\big/\!\int_{\overline{\mathcal{K}}}e^{-\Phi}$, where
\begin{equation*}
\Phi(y)=\tfrac12\langle y,\mathsf{Q}y\rangle+\langle\mathsf{h},y\rangle-P(y).
\end{equation*}
Here $\mathsf{Q}$ is real symmetric with $\mathsf{q}_-\mathbf{1}\le\mathsf{Q}\le\mathsf{q}_+\mathbf{1}$ and $\mathsf{q}_->0$. The vector $\mathsf{h}$ lies in $\mathbb{R}^{n|\mathfrak{X}|}$, and $P$ is a bounded function of class $C^2$ on $\mathcal{K}^+$ satisfying the bounds of item~(iv).

Membership in $\mathfrak{S}$ does not include the smoothness of $P$, or the boundedness of every derivative of $P$, on a neighbourhood of the compact set $\prod_x\overline{K_x+B(0,\delta r_-/2)}$. In the abstract setting of the class this is a hypothesis made in addition to the regularity of $P$ required there, and the lemma below assumes it explicitly.

In the application, $P$ is holomorphic on a complex tube containing $\mathcal{K}^+$, by the localisation and nesting properties established for it. Cauchy's estimates on the remaining margin then bound every derivative of $P$ on a neighbourhood of the compact set $\prod_x\overline{K_x+B(0,\delta r_-/2)}$, which contains the image of the retraction $\mathsf{r}$ of Step~0 below. The intersection of this neighbourhood with $\mathcal{K}^+$ is then an open set $\mathcal{N}$ as in the lemma.
\item[(iv)] Bounds on $P$. On $\mathcal{K}^+$ one has $\|\partial_xP\|\le\mathsf{p}_1$ for every $x$. Also, $P''$ has operator norm at most $\varpi$, where $\varpi$ is the bound of the class on the Schur norm of $P''$ over site blocks on $\mathcal{K}^+$. The smallness condition of the class gives $\varpi\le\mathsf{q}_-/4$, and we put $\mathsf{m}:=\mathsf{q}_-/2$.
\item[(v)] Coupling. For $x\neq x'$, $\mathfrak{c}(x,x')\ge0$ is the coupling bound of the class between the sites $x$ and $x'$. The property of it used below is
\begin{equation*}
\|\mathsf{Q}_{xx'}\|+\sup_{\mathcal{K}^+}\|\partial_x\partial_{x'}P\|\le\mathfrak{c}(x,x'),
\end{equation*}
which gives in particular $\|\mathsf{Q}_{xx'}-\partial_x\partial_{x'}P(y)\|\le\mathfrak{c}(x,x')$ for $y\in\mathcal{K}^+$. The class carries constants $\nu_0>0$ and $\mathfrak{K}<\infty$ satisfying the first of the following inequalities, and a decay rate $\nu$ chosen to satisfy the other two:
\begin{equation*}
\sup_{x}\sum_{x'\neq x}\mathfrak{c}(x,x')e^{\nu_0\operatorname{dist}(x,x')}\le\mathfrak{K},\qquad
0<\nu\le\frac{\nu_0}{2},\qquad
\nu\le\frac{e\nu_0\mathsf{m}}{4\mathfrak{K}} .
\end{equation*}
\item[(vi)] Conditioning. Let $S\subset\mathfrak{X}$ and $y_S\in\mathcal{K}_S$. The conditional law $\mu^{y_S}$ of $y_{S^c}$ given $y_S$ is carried by the fixed set $\mathcal{K}_{S^c}$ and has density proportional to $\mathbf{1}_{\mathcal{K}_{S^c}}(y')e^{-\Phi(y_S,y')}$. It is again a member of $\mathfrak{S}$, with the same constants $\mathsf{m}$, $\mathfrak{K}$ and $\nu$.
\end{itemize}

\begin{lemma}[Exponential decay of correlations in the class]\label{8.3:lem-correlation-decay}
Assume the following statement, the Helffer--Sj\"ostrand covariance representation. Let $\mathsf{N}\in\mathbb{N}$ and $V\in C^\infty(\mathbb{R}^{\mathsf{N}})$. Suppose that $\lambda_-\mathbf{1}\le\operatorname{Hess}V(y)\le\lambda_+\mathbf{1}$ for all $y$, for some $0<\lambda_-\le\lambda_+<\infty$, and that all derivatives of $V$ of order $\ge2$ are bounded on $\mathbb{R}^{\mathsf{N}}$. Then the measure $\mu_V(dy):=e^{-V(y)}dy/\int e^{-V}$ and the operator $\mathcal{L}_V:=\Delta-\nabla V\cdot\nabla$ satisfy the following:
\begin{itemize}
\item[(1)] $-\mathcal{L}_V$ on $C_c^\infty(\mathbb{R}^{\mathsf{N}})\subset L^2(\mu_V)$ is essentially self-adjoint and $\ge0$;
\item[(2)] the operator $\mathfrak{L}_V$ on vector fields $v=(v_k)_{k\le\mathsf{N}}$, defined by $(\mathfrak{L}_Vv)_k:=-\mathcal{L}_Vv_k+\sum_l(\partial_k\partial_lV)v_l$, is essentially self-adjoint on $C_c^\infty(\mathbb{R}^{\mathsf{N}};\mathbb{R}^{\mathsf{N}})\subset L^2(\mu_V;\mathbb{R}^{\mathsf{N}})$, and $\mathfrak{L}_V\ge\lambda_-$;
\item[(3)] for $F,G\in C^\infty(\mathbb{R}^{\mathsf{N}})$ with bounded derivatives of all orders $\ge1$,
\begin{equation}\label{8.3:eq-correlation-decay-1}
\operatorname{Cov}_{\mu_V}(F,G)=\langle\nabla F,\mathfrak{L}_V^{-1}\nabla G\rangle_{L^2(\mu_V;\mathbb{R}^{\mathsf{N}})} .
\end{equation}
\end{itemize}
Let $\mu\in\mathfrak{S}$. Assume in addition, beyond the data listed in item~(iii) above, that there is an open set $\mathcal{N}$ with
\begin{equation*}
\prod_{x\in\mathfrak{X}}\overline{K_x+B(0,\delta r_-/2)}\subset\mathcal{N}\subset\mathcal{K}^+
\end{equation*}
on which $P$ is $C^\infty$ and has bounded derivatives of every order. Then the following hold.
\begin{itemize}
\item[(a)] Let $F,G$ be $C^\infty$ on $\mathcal{K}^+$ with bounded derivatives of all orders on $\mathcal{K}^+$. Then
\begin{equation}\label{8.3:eq-correlation-decay-2}
|\operatorname{Cov}_\mu(F,G)|\le\frac{2}{\mathsf{m}}\sum_{x,x'\in\mathfrak{X}}
e^{-\nu\operatorname{dist}(x,x')}\,\|\partial_xF\|_{\infty,\mathcal{K}^+}\,
\|\partial_{x'}G\|_{\infty,\mathcal{K}^+}.
\end{equation}
\item[(b)] Let $S\subset\mathfrak{X}$, and let $F$ be bounded measurable and depend on $y_S$ only. Let $c$ be any constant, and let $G$ be $C^\infty$ on $\mathcal{K}^+$ with bounded derivatives of all orders on $\mathcal{K}^+$. Then
\begin{equation}\label{8.3:eq-correlation-decay-3}
\begin{split}
|\operatorname{Cov}_\mu(F,G)|\le\|F-c\|_\infty\sum_{x\in S}2(1+\delta)r_x
\Bigl[\,&\|\partial_xG\|_\infty\\
&+\frac{2}{\mathsf{m}}\sum_{x'',x'''\in\mathfrak{X}\setminus S}
e^{-\nu\operatorname{dist}(x'',x''')}\|\partial_{x''}G\|_\infty\,
\mathfrak{c}(x,x''')\Bigr].
\end{split}
\end{equation}
Here $\|\partial_xG\|_\infty$ and $\|\partial_{x''}G\|_\infty$ denote suprema over $\mathcal{K}^+$, as in (a), and $\|F-c\|_\infty$ is the supremum norm of $F-c$.
\end{itemize}
\end{lemma}

\begin{proof}
Write $\mathsf{N}:=n|\mathfrak{X}|$.

\medskip\noindent Step~0: reduction to a smooth potential on the whole space.

\emph{Smoothing devices.} For each $x\in\mathfrak{X}$ we apply Lemma~\ref{8.3:lem-retraction-wall} to the convex body $K_x$, with the enlargement radius of that lemma equal to $\delta r_-/2$. We write $K_x':=K_x+B(0,\delta r_-/2)$. The lemma provides $\mathsf{r}_x\in C^\infty(\mathbb{R}^n;\mathbb{R}^n)$, a function $\psi_x\in C^\infty(\mathbb{R}^n;[0,\infty))$ and a number $C_{\psi_x}$. By \eqref{8.3:eq-retraction-wall-a} they have the following properties:
\begin{itemize}
\item $\mathsf{r}_x=\mathrm{id}$ on $\overline{K_x}$;
\item $\mathsf{r}_x(\mathbb{R}^n)\subset\overline{K_x'}\subset K_x^+$;
\item $D\mathsf{r}_x$ is symmetric with $0\le D\mathsf{r}_x\le\mathbf{1}$;
\item $\psi_x$ is convex with $0\le D^2\psi_x\le C_{\psi_x}$, and $\psi_x=0$ exactly on $\overline{K_x}$, so that $\psi_x>0$ off $\overline{K_x}$;
\item $|D^2\mathsf{r}_x(u)[\xi,\xi]|\le D^2\psi_x(u)[\xi,\xi]$ for all $u,\xi\in\mathbb{R}^n$.
\end{itemize}
Put $\mathsf{r}(y):=(\mathsf{r}_x(y_x))_{x\in\mathfrak{X}}$, $\Psi(y):=\sum_x\psi_x(y_x)$ and $C_\psi:=\max_xC_{\psi_x}$. For $\mathsf{J}\in\mathbb{N}$ put
\begin{equation}\label{8.3:eq-correlation-decay-4}
\Phi_{\mathsf{J}}(y):=\tfrac12\langle y,\mathsf{Q}y\rangle+\langle\mathsf{h},y\rangle
-P(\mathsf{r}(y))+(\mathsf{p}_1+\mathsf{J})\Psi(y),\qquad y\in\mathbb{R}^{\mathsf{N}},
\end{equation}
and $\mu_{\mathsf{J}}:=e^{-\Phi_{\mathsf{J}}}dy/\int e^{-\Phi_{\mathsf{J}}}$.

(0-a) $\Phi_{\mathsf{J}}\in C^\infty(\mathbb{R}^{\mathsf{N}})$. The $\mathsf{Q}$-part and the linear part are polynomials. The composite $P\circ\mathsf{r}$ is smooth because $\mathsf{r}$ is smooth with values in the compact set $\prod_x\overline{K_x'}\subset\mathcal{N}$, where $P$ is $C^\infty$ by the additional hypothesis of the lemma, and the function $\Psi$ is smooth.

(0-a$'$) All derivatives of $\Phi_{\mathsf{J}}$ of order $\ge2$ are bounded on $\mathbb{R}^{\mathsf{N}}$. Moreover, suppose $F$ is $C^\infty$ with bounded derivatives of every order on a neighbourhood of the compact set $\prod_x\overline{K_x'}$, where
\begin{equation*}
\mathsf{r}(\mathbb{R}^{\mathsf{N}})\subset\prod_x\overline{K_x'}\subset\mathcal{K}^+ .
\end{equation*}
Then $\tilde F:=F\circ\mathsf{r}$ has bounded derivatives of every order $\ge1$.

To see this, recall the construction in the proof of Lemma~\ref{8.3:lem-retraction-wall}. There $\mathsf{r}_x=\mathrm{id}-\nabla(f_x*\omega)$ and $\psi_x=f_x*\tilde\omega$. Here $f_x(u):=\frac12\operatorname{dist}(u,\overline{K_x'})^2$ for $u\in\mathbb{R}^n$, and $\pi_x$ is the metric projection onto $\overline{K_x'}$. The functions $\omega,\tilde\omega\in C_c^\infty(\mathbb{R}^n)$ are the two kernels of that proof: the mollifier and the kernel dominating its gradient.

The $\mathsf{Q}$-part of $\Phi_{\mathsf{J}}$ is quadratic, so its derivatives of order $\ge2$ are constant or zero; those of the linear part vanish.

By \eqref{8.3:eq-retraction-wall-b}, $\nabla f_x=\mathrm{id}-\pi_x$. By that identity and the non-expansiveness of $\pi_x$, this map is $1$-Lipschitz. For $j\ge2$ we have $\int D^{j-1}\omega=0$, and therefore, for $u\in\mathbb{R}^n$,
\begin{equation*}
D^j(f_x*\omega)(u)=\int\bigl(\nabla f_x(u-z)-\nabla f_x(u)\bigr)\cdot D^{j-1}\omega(z)\,dz .
\end{equation*}
This has norm at most $\int|z|\,|D^{j-1}\omega(z)|\,dz$, and likewise with $\tilde\omega$ in place of $\omega$. Hence $\mathsf{r}_x=\mathrm{id}-\nabla(f_x*\omega)$ has all derivatives of order $\ge1$ bounded, and $\psi_x=f_x*\tilde\omega$ has all derivatives of order $\ge2$ bounded. The same therefore holds for $\mathsf{r}$ and for $\Psi$.

Next, $P$ has bounded derivatives of every order on the open set $\mathcal{N}$, by the additional hypothesis of the lemma. The set $\mathcal{N}$ is a neighbourhood of the compact set $\prod_x\overline{K_x'}$, which contains the image of $\mathsf{r}$.

By the chain rule, every derivative of $P\circ\mathsf{r}$, and of $F\circ\mathsf{r}$, of order $j\ge1$ is a finite sum of products. Each product consists of a derivative of $P$ (respectively $F$) at $\mathsf{r}(y)$ and derivatives of $\mathsf{r}$ of orders between $1$ and $j$. All these factors are bounded, so every derivative of $P\circ\mathsf{r}$ and of $F\circ\mathsf{r}$ of order $\ge1$ is bounded. This proves (0-a$'$).

Consequently, once (0-b) below is established, the Helffer--Sj\"ostrand representation applies to $V:=\Phi_{\mathsf{J}}$, and its identity \eqref{8.3:eq-correlation-decay-1} applies to the pairs $(\tilde F,\tilde G)$ met below.

(0-b) $\mathsf{m}\mathbf{1}\le\operatorname{Hess}\Phi_{\mathsf{J}}\le\Lambda_{\mathsf{J}}\mathbf{1}$ with $\mathsf{m}=\mathsf{q}_-/2$ and a finite $\Lambda_{\mathsf{J}}$. Since $\mathsf{r}$ and $\Psi$ act site by site, for $\xi=(\xi_x)_x$,
\begin{equation}\label{8.3:eq-correlation-decay-5}
\begin{split}
\operatorname{Hess}\Phi_{\mathsf{J}}(y)[\xi,\xi]={}&\langle\xi,\mathsf{Q}\xi\rangle
-P''(\mathsf{r}(y))[D\mathsf{r}\,\xi,D\mathsf{r}\,\xi]\\
&+\sum_x\Bigl\{(\mathsf{p}_1+\mathsf{J})D^2\psi_x(y_x)[\xi_x,\xi_x]
-\bigl\langle\partial_xP(\mathsf{r}(y)),D^2\mathsf{r}_x(y_x)[\xi_x,\xi_x]\bigr\rangle\Bigr\}.
\end{split}
\end{equation}
We bound the terms in turn.
\begin{itemize}
\item The second term. Since $\mathsf{r}(y)\in\mathcal{K}^+$, the Schur bound of item~(iv) over site blocks applies. Together with $\|D\mathsf{r}_x\|\le1$, it shows that the modulus of this term is at most $\varpi|\xi|^2$.
\item The braces. On $\mathcal{K}^+\ni\mathsf{r}(y)$ we have $\|\partial_xP\|\le\mathsf{p}_1$, and $|D^2\mathsf{r}_x[\xi_x,\xi_x]|\le D^2\psi_x[\xi_x,\xi_x]$. Hence each brace satisfies
\begin{equation*}
\{\cdots\}\ge(\mathsf{p}_1+\mathsf{J}-\mathsf{p}_1)D^2\psi_x[\xi_x,\xi_x]
=\mathsf{J}D^2\psi_x[\xi_x,\xi_x]\ge0 .
\end{equation*}
Thus the wall absorbs the curvature of the retraction. This is the same mechanism as in \eqref{8.3:eq-lipschitz-stability-c} in the proof of Proposition~\ref{8.3:prop-lipschitz-stability}.
\item Lower bound. We have $\mathsf{Q}\ge\mathsf{q}_-$. The second term is at least $-\varpi|\xi|^2\ge-\tfrac14\mathsf{q}_-|\xi|^2$, by the smallness condition of item~(iv). Each brace is non-negative. Hence
\begin{equation*}
\operatorname{Hess}\Phi_{\mathsf{J}}\ge\mathsf{q}_--\mathsf{q}_-/4\ge\mathsf{m}.
\end{equation*}
\item Upper bound. Each brace is at most $(\mathsf{p}_1+\mathsf{J})C_\psi|\xi_x|^2+\mathsf{p}_1C_\psi|\xi_x|^2$, so
\begin{equation*}
\operatorname{Hess}\Phi_{\mathsf{J}}\le\mathsf{q}_++\mathsf{q}_-/4+(2\mathsf{p}_1+\mathsf{J})C_\psi=:\Lambda_{\mathsf{J}}.
\end{equation*}
\end{itemize}

(0-c) For $x\neq x'$ the off-diagonal block of the Hessian is
\begin{equation*}
(\operatorname{Hess}\Phi_{\mathsf{J}})_{xx'}(y)=\mathsf{Q}_{xx'}
-(D\mathsf{r}_x)^{\mathsf{T}}(\partial_x\partial_{x'}P)(\mathsf{r}(y))\,D\mathsf{r}_{x'},
\end{equation*}
because $\Psi$ and the terms containing $D^2\mathsf{r}$ are site-diagonal. We have $\|D\mathsf{r}_x\|\le1$, $\|D\mathsf{r}_{x'}\|\le1$ and $\mathsf{r}(y)\in\mathcal{K}^+$. Hence item~(v) gives $\|(\operatorname{Hess}\Phi_{\mathsf{J}})_{xx'}(y)\|\le\mathfrak{c}(x,x')$ for all $y$ and all $\mathsf{J}$.

(0-d) As $\mathsf{J}\to\infty$, $\mu_{\mathsf{J}}\to\mu$ weakly. Moreover, $\operatorname{Cov}_{\mu_{\mathsf{J}}}(\tilde F,\tilde G)\to\operatorname{Cov}_\mu(\tilde F,\tilde G)$ for bounded continuous $\tilde F,\tilde G$ on $\mathbb{R}^{\mathsf{N}}$.

Domination. Since $\Psi\ge0$ and $\mathsf{r}(y)\in\mathcal{K}^+$, the densities $e^{-\Phi_{\mathsf{J}}}$ are dominated by
\begin{equation*}
\exp\Bigl[-\tfrac12\langle y,\mathsf{Q}y\rangle-\langle\mathsf{h},y\rangle+\sup_{\mathcal{K}^+}P\Bigr].
\end{equation*}
This function is in $L^1$: the dimension is finite, $\mathsf{Q}\ge\mathsf{q}_->0$, and $P$ is bounded on the bounded set $\mathcal{K}^+$. The bound depends on $|\mathfrak{X}|$. This is harmless here, since no constant below depends on it.

Pointwise limit. On $\overline{\mathcal{K}}$ we have $\mathsf{r}=\mathrm{id}$ and $\Psi=0$, so $\Phi_{\mathsf{J}}=\Phi$. Off $\overline{\mathcal{K}}$ we have $\Psi>0$, so $e^{-(\mathsf{p}_1+\mathsf{J})\Psi}\to0$. Hence the densities converge pointwise to $\mathbf{1}_{\overline{\mathcal{K}}}(y)e^{-\Phi(y)}$.

Since $\mathcal{K}$ is convex, its boundary $\partial\mathcal{K}$ is Lebesgue-null. We apply dominated convergence in the numerators $\int\tilde F\tilde Ge^{-\Phi_{\mathsf{J}}}$, $\int\tilde Fe^{-\Phi_{\mathsf{J}}}$, $\int\tilde Ge^{-\Phi_{\mathsf{J}}}$ and in the denominators $\int e^{-\Phi_{\mathsf{J}}}$. This gives the claim, since the denominators converge to $\int_{\overline{\mathcal{K}}}e^{-\Phi}>0$.

(0-e) For $F$ as in (a) put $\tilde F:=F\circ\mathsf{r}$. By (0-a$'$), $\tilde F\in C^\infty(\mathbb{R}^{\mathsf{N}})$ has bounded derivatives of every order $\ge1$. It is bounded and continuous, since it takes the values of the continuous function $F$ on the set $\mathsf{r}(\mathbb{R}^{\mathsf{N}})$, which lies in the compact subset $\prod_x\overline{K_x'}$ of $\mathcal{K}^+$. Since $\mathsf{r}=\mathrm{id}$ on $\overline{\mathcal{K}}$, we have $\tilde F=F$ there. Since $\mathsf{r}$ acts site by site,
\begin{equation*}
\partial_x\tilde F(y)=(D\mathsf{r}_x(y_x))^{\mathsf{T}}(\partial_xF)(\mathsf{r}(y))
\end{equation*}
involves only $\partial_xF$, and $\|D\mathsf{r}_x\|\le1$ gives $\|\partial_x\tilde F\|_\infty\le\|\partial_xF\|_{\infty,\mathcal{K}^+}$. The same holds for $\tilde G:=G\circ\mathsf{r}$.

Since $\mu$ is carried by $\overline{\mathcal{K}}$, (0-d) gives
\begin{equation*}
\operatorname{Cov}_\mu(F,G)=\operatorname{Cov}_\mu(\tilde F,\tilde G)
=\lim_{\mathsf{J}\to\infty}\operatorname{Cov}_{\mu_{\mathsf{J}}}(\tilde F,\tilde G).
\end{equation*}
It therefore suffices to prove \eqref{8.3:eq-correlation-decay-2} for $\mu_{\mathsf{J}}$ and the pair $(\tilde F,\tilde G)$, with constants free of $\mathsf{J}$, and then to let $\mathsf{J}\to\infty$.

\medskip\noindent Step~1: the covariance representation.

By (0-a), (0-a$'$) and (0-b), the Helffer--Sj\"ostrand representation applies to $V=\Phi_{\mathsf{J}}$ with $\lambda_-=\mathsf{m}$ and $\lambda_+=\Lambda_{\mathsf{J}}$. Let $\mathcal{L}_{\mathsf{J}}:=\mathcal{L}_{\Phi_{\mathsf{J}}}$, and let $\mathfrak{L}_{\mathsf{J}}:=(-\mathcal{L}_{\mathsf{J}})\otimes\mathbf{1}+\operatorname{Hess}\Phi_{\mathsf{J}}$ be the (closure of the) operator of item~(2). It is self-adjoint with $\mathfrak{L}_{\mathsf{J}}\ge\mathsf{m}$, and
\begin{equation*}
\operatorname{Cov}_{\mu_{\mathsf{J}}}(\tilde F,\tilde G)
=\langle\nabla\tilde F,\mathfrak{L}_{\mathsf{J}}^{-1}\nabla\tilde G\rangle_{L^2(\mu_{\mathsf{J}};\mathbb{R}^{\mathsf{N}})} .
\end{equation*}

\medskip\noindent Step~2: the site-resolved resolvent bound, by the argument of Combes and Thomas \cite{CT73}.

Index the components of a vector field $v\in L^2(\mu_{\mathsf{J}};\mathbb{R}^{\mathsf{N}})$ by sites: $v=(v_x)_{x\in\mathfrak{X}}$ with $v_x\in L^2(\mu_{\mathsf{J}};\mathbb{R}^n)$. Let $\Pi_x$ be the orthogonal projection onto the $x$-components. We claim that for all $x,x'\in\mathfrak{X}$
\begin{equation}\label{8.3:eq-correlation-decay-6}
\|\Pi_x\mathfrak{L}_{\mathsf{J}}^{-1}\Pi_{x'}\|\le\frac{2}{\mathsf{m}}e^{-\nu\operatorname{dist}(x,x')}.
\end{equation}

\emph{The weighted operator.} Fix $x_0\in\mathfrak{X}$ and put $\varphi(x):=\nu\operatorname{dist}(x,x_0)$. Let $\mathcal{U}$ be the operator $(\mathcal{U}v)_x:=e^{\varphi(x)}v_x$. It multiplies all components at site $x$ by the same constant. Since $\mathfrak{X}$ is finite, it is bounded and boundedly invertible.

The operator $\mathfrak{M}:=(-\mathcal{L}_{\mathsf{J}})\otimes\mathbf{1}$ acts on each scalar component separately. Hence $\mathcal{U}\mathfrak{M}\mathcal{U}^{-1}=\mathfrak{M}$ on the common core $C_c^\infty$. The operator $\mathcal{U}$ preserves this core and preserves the graph norm up to constants, so the identity extends to $\operatorname{Dom}(\mathfrak{L}_{\mathsf{J}})$.

Let $\mathcal{B}$ denote multiplication by $\operatorname{Hess}\Phi_{\mathsf{J}}(y)$. Then $\mathcal{U}\mathcal{B}\mathcal{U}^{-1}=\mathcal{B}+\mathcal{D}$, where $\mathcal{D}$ is multiplication by the matrix with blocks
\begin{equation*}
\mathcal{D}(y)_{xx'}=\bigl(e^{\varphi(x)-\varphi(x')}-1\bigr)(\operatorname{Hess}\Phi_{\mathsf{J}}(y))_{xx'},
\qquad \mathcal{D}(y)_{xx}=0 .
\end{equation*}

\emph{Bound on the perturbation.} We have $|\varphi(x)-\varphi(x')|\le\nu\operatorname{dist}(x,x')$, and $|e^\tau-1|\le e^{|\tau|}-1\le|\tau|e^{|\tau|}$. Hence, by (0-c),
\begin{equation*}
\begin{split}
\|\mathcal{D}(y)_{xx'}\|
&\le\|(\operatorname{Hess}\Phi_{\mathsf{J}}(y))_{xx'}\|\,\nu\operatorname{dist}(x,x')
e^{\nu\operatorname{dist}(x,x')}\\
&\le\mathfrak{c}(x,x')\,\nu\operatorname{dist}(x,x')e^{\nu\operatorname{dist}(x,x')}.
\end{split}
\end{equation*}
The middle expression of this display is symmetric in $(x,x')$, because the Hessian is symmetric. So both the row sums and the column sums of the block norms $\|\mathcal{D}(y)_{xx'}\|$ are bounded by the row sums of the middle expression, and hence by the supremum over $x$ of the sums of the right-hand side. We use $\sup_{\tau\ge0}\tau e^{-(\nu_0-\nu)\tau}=1/(e(\nu_0-\nu))$. Schur's test with block operator norms then gives, for every $y$,
\begin{equation*}
\begin{split}
\|\mathcal{D}(y)\|_{\ell^2\to\ell^2}
&\le\sup_x\sum_{x'\neq x}\mathfrak{c}(x,x')\,\nu\operatorname{dist}(x,x')e^{\nu\operatorname{dist}(x,x')}\\
&\le\nu\,\sup_x\sum_{x'\neq x}\mathfrak{c}(x,x')e^{\nu_0\operatorname{dist}(x,x')}\,
\sup_{\tau\ge0}\tau e^{-(\nu_0-\nu)\tau}\\
&\le\frac{\nu\mathfrak{K}}{e(\nu_0-\nu)}\le\frac{2\nu\mathfrak{K}}{e\nu_0}\le\frac{\mathsf{m}}{2}.
\end{split}
\end{equation*}
The last two inequalities use the choice of $\nu$ in item~(v): $\nu\le\nu_0/2$ gives $\nu_0-\nu\ge\nu_0/2$, and $\nu\le e\nu_0\mathsf{m}/(4\mathfrak{K})$. Since the bound holds fibrewise, $\|\mathcal{D}\|_{L^2(\mu_{\mathsf{J}};\mathbb{R}^{\mathsf{N}})\to L^2}\le\mathsf{m}/2$.

\emph{Inversion.} We now have $\mathcal{U}\mathfrak{L}_{\mathsf{J}}\mathcal{U}^{-1}=\mathfrak{L}_{\mathsf{J}}+\mathcal{D}$ on $\operatorname{Dom}(\mathfrak{L}_{\mathsf{J}})$. Since $\mathfrak{L}_{\mathsf{J}}\ge\mathsf{m}$, it is invertible with $\|\mathfrak{L}_{\mathsf{J}}^{-1}\|\le1/\mathsf{m}$, and so $\|\mathfrak{L}_{\mathsf{J}}^{-1}\mathcal{D}\|\le\frac12$. By the Neumann series, $\mathfrak{L}_{\mathsf{J}}+\mathcal{D}=\mathfrak{L}_{\mathsf{J}}(1+\mathfrak{L}_{\mathsf{J}}^{-1}\mathcal{D})$ is invertible with $\|(\mathfrak{L}_{\mathsf{J}}+\mathcal{D})^{-1}\|\le2/\mathsf{m}$.

Therefore $\mathfrak{L}_{\mathsf{J}}^{-1}=\mathcal{U}^{-1}(\mathfrak{L}_{\mathsf{J}}+\mathcal{D})^{-1}\mathcal{U}$. Since $\varphi(x_0)=0$,
\begin{equation*}
\begin{split}
\Pi_x\mathfrak{L}_{\mathsf{J}}^{-1}\Pi_{x_0}
&=e^{-\varphi(x)}\Pi_x(\mathfrak{L}_{\mathsf{J}}+\mathcal{D})^{-1}\Pi_{x_0}e^{\varphi(x_0)}\\
&=e^{-\nu\operatorname{dist}(x,x_0)}\Pi_x(\mathfrak{L}_{\mathsf{J}}+\mathcal{D})^{-1}\Pi_{x_0},
\end{split}
\end{equation*}
which has norm at most $(2/\mathsf{m})e^{-\nu\operatorname{dist}(x,x_0)}$. Taking $x_0=x'$ gives \eqref{8.3:eq-correlation-decay-6}.

\medskip\noindent Step~3: conclusion of (a).

We decompose $\nabla\tilde F$ and $\nabla\tilde G$ into their site components in the identity of Step~1:
\begin{equation*}
\operatorname{Cov}_{\mu_{\mathsf{J}}}(\tilde F,\tilde G)=\sum_{x,x'}
\bigl\langle\Pi_x\nabla\tilde F,\;\Pi_x\mathfrak{L}_{\mathsf{J}}^{-1}\Pi_{x'}\,\Pi_{x'}\nabla\tilde G\bigr\rangle .
\end{equation*}
We use the Cauchy--Schwarz inequality, \eqref{8.3:eq-correlation-decay-6}, the fact that $\mu_{\mathsf{J}}$ is a probability measure, and (0-e). This gives
\begin{equation*}
\begin{split}
|\operatorname{Cov}_{\mu_{\mathsf{J}}}(\tilde F,\tilde G)|
&\le\frac{2}{\mathsf{m}}\sum_{x,x'}e^{-\nu\operatorname{dist}(x,x')}
\|\partial_x\tilde F\|_{L^2(\mu_{\mathsf{J}})}\|\partial_{x'}\tilde G\|_{L^2(\mu_{\mathsf{J}})}\\
&\le\frac{2}{\mathsf{m}}\sum_{x,x'}e^{-\nu\operatorname{dist}(x,x')}
\|\partial_xF\|_{\infty,\mathcal{K}^+}\|\partial_{x'}G\|_{\infty,\mathcal{K}^+}.
\end{split}
\end{equation*}
The right-hand side does not depend on $\mathsf{J}$. Letting $\mathsf{J}\to\infty$ as in Step~0 gives \eqref{8.3:eq-correlation-decay-2}. This proves (a).

The argument of Steps~0--3 uses, of the functions $F$ and $G$, only the following two properties. They are bounded and continuously differentiable on $\mathcal{K}^+$ with bounded first derivatives there; this enters through (0-d) and (0-e). They are $C^\infty$ with bounded derivatives of every order on $\mathcal{N}$; this enters through (0-a$'$), since $\mathcal{N}$ is a neighbourhood of $\prod_x\overline{K_x'}$. Hence \eqref{8.3:eq-correlation-decay-2} holds for every pair $F,G$ with these two properties.

\medskip\noindent Step~4: proof of (b).

Let $F=F(y_S)$ be bounded measurable, $c$ a constant, and $G$ as in (b). For $y_S\in\mathcal{K}_S$ put
\begin{equation*}
\bar G(y_S):=\langle G(y_S,\cdot)\rangle_{\mu^{y_S}},
\end{equation*}
with $\mu^{y_S}$ the conditional law of item~(vi); $\bar G$ is the conditional expectation of $G$ given $y_S$. The window is a product, $\mathcal{K}=\mathcal{K}_S\times\mathcal{K}_{S^c}$, so $\mu^{y_S}$ is carried by the fixed set $\mathcal{K}_{S^c}$ for every $y_S\in\mathcal{K}_S$ (item~(vi)).

\emph{Reduction to an oscillation.} The function $F-c$ is $y_S$-measurable. Hence
\begin{equation*}
\operatorname{Cov}_\mu(F,G)=\int(F-c)\bigl(G-\mu(G)\bigr)d\mu=\int(F-c)\bigl(\bar G-\mu(G)\bigr)d\mu .
\end{equation*}
The $y_S$-marginal of $\mu$ is absolutely continuous and carried by $\overline{\mathcal{K}_S}$. The boundary of $\overline{\mathcal{K}_S}$ is Lebesgue-null, so the marginal is carried by $\mathcal{K}_S$. The value $\mu(G)=\int\bar G\,d\mu$ lies in the closed convex hull of the range of $\bar G$. Hence $|\bar G(y_S)-\mu(G)|\le\operatorname{osc}_{\mathcal{K}_S}\bar G$, where
\begin{equation*}
\operatorname{osc}_{\mathcal{K}_S}\bar G:=\sup_{y_S^{(0)},y_S^{(1)}\in\mathcal{K}_S}
\bigl|\bar G(y_S^{(1)})-\bar G(y_S^{(0)})\bigr| .
\end{equation*}
Hence
\begin{equation*}
|\operatorname{Cov}_\mu(F,G)|\le\|F-c\|_\infty\int|\bar G-\mu(G)|\,d\mu\le\|F-c\|_\infty\operatorname{osc}_{\mathcal{K}_S}\bar G .
\end{equation*}

\emph{Differentiability of $\bar G$.} We have
\begin{equation*}
\bar G(y_S)=\frac{\int G(y_S,y')e^{-\Phi(y_S,y')}\mathbf{1}_{\mathcal{K}_{S^c}}(y')\,dy'}
{\int e^{-\Phi(y_S,y')}\mathbf{1}_{\mathcal{K}_{S^c}}(y')\,dy'} .
\end{equation*}
The integrands are $C^1$ in $y_S$ and bounded with bounded derivatives on $\mathcal{K}^+$, and the domain of $y'$ does not depend on $y_S$. So we may differentiate under the integral. The derivative of the numerator is $\int(\partial_xG-G\,\partial_x\Phi)e^{-\Phi}$, and that of the denominator is $-\int\partial_x\Phi\,e^{-\Phi}$. Hence $\bar G\in C^1(\mathcal{K}_S)$ and, for $x\in S$,
\begin{equation}\label{8.3:eq-correlation-decay-7}
\partial_{y_x}\bar G(y_S)=\langle\partial_xG\rangle_{\mu^{y_S}}
-\operatorname{Cov}_{\mu^{y_S}}\bigl(G(y_S,\cdot),\partial_x\Phi(y_S,\cdot)\bigr).
\end{equation}

\emph{Bound on the derivative.} By item~(vi), $\mu^{y_S}\in\mathfrak{S}$ on the sites $S^c$, with the same $\mathsf{m}$, $\mathfrak{K}$ and $\nu$. Its density is proportional to $\mathbf{1}_{\mathcal{K}_{S^c}}e^{-\Phi(y_S,\cdot)}$. The part of $\Phi(y_S,\cdot)$ that is not a polynomial of degree at most two in $y'$ is $-P(y_S,\cdot)$. Put $\mathcal{N}^{y_S}:=\{y':(y_S,y')\in\mathcal{N}\}$. This set is open and contained in $\prod_{x'''\in S^c}K^+_{x'''}$. It contains $\prod_{x'''\in S^c}\overline{K'_{x'''}}$, because $y_S\in\mathcal{K}_S\subset\prod_{x\in S}\overline{K_x'}$. On $\mathcal{N}^{y_S}$ the function $P(y_S,\cdot)$ is $C^\infty$ and has bounded derivatives of every order. So the additional hypothesis of the lemma holds for $\mu^{y_S}$ as well, with $\mathcal{N}^{y_S}$ in place of $\mathcal{N}$. We apply (a) to $\mu^{y_S}$ and, for each unit vector $\theta\in\mathbb{R}^n$, to the pair of functions of $y'$
\begin{equation*}
G(y_S,\cdot)\quad\text{and}\quad\langle \theta,\partial_x\Phi(y_S,\cdot)\rangle .
\end{equation*}
The function $G(y_S,\cdot)$ is $C^\infty$ in $y'$ with all derivatives bounded on $\prod_{x'''\in S^c}K^+_{x'''}$, by hypothesis.

The function $\langle\theta,\partial_x\Phi(y_S,\cdot)\rangle$, where $\partial_x\Phi=(\mathsf{Q}y+\mathsf{h})_x-\partial_xP$, is bounded and continuously differentiable on $\prod_{x'''\in S^c}K^+_{x'''}$. Its first derivatives are bounded there by item~(v), as computed below, and it is $C^\infty$ with bounded derivatives of every order on $\mathcal{N}^{y_S}$.

Both functions therefore have the two properties listed at the end of Step~3, read for $\mu^{y_S}$ with $\mathcal{N}^{y_S}$ in place of $\mathcal{N}$. Hence the conclusion of (a) holds for this pair.

The $x'''$-gradient of the second function is $(\partial_{x'''}\partial_x\Phi)^{\mathsf{T}}\theta$, of norm at most $\|\partial_{x'''}\partial_x\Phi\|$. Applying (a) in this form to the pair and taking the supremum over unit vectors $\theta$ gives
\begin{equation*}
\begin{split}
&\bigl|\operatorname{Cov}_{\mu^{y_S}}\bigl(G(y_S,\cdot),\partial_x\Phi(y_S,\cdot)\bigr)\bigr|\\
&\qquad\le\frac{2}{\mathsf{m}}\sum_{x'',x'''\in S^c}e^{-\nu\operatorname{dist}(x'',x''')}
\|\partial_{x''}G\|_\infty\,\|\partial_{x'''}\partial_x\Phi\|_\infty .
\end{split}
\end{equation*}
For $x\in S$ and $x'''\notin S$ we have $x'''\neq x$. So $\partial_{x'''}\partial_x\Phi=\mathsf{Q}_{x'''x}-\partial_{x'''}\partial_xP$, which has norm at most $\mathfrak{c}(x,x''')$ by item~(v). Together with \eqref{8.3:eq-correlation-decay-7},
\begin{equation*}
\|\partial_{y_x}\bar G\|_\infty\le\|\partial_xG\|_\infty+\frac{2}{\mathsf{m}}\sum_{x'',x'''\in S^c}
e^{-\nu\operatorname{dist}(x'',x''')}\|\partial_{x''}G\|_\infty\,\mathfrak{c}(x,x''').
\end{equation*}

\emph{Bound on the oscillation.} The set $\mathcal{K}_S$ is convex. For $y_S^{(0)},y_S^{(1)}\in\mathcal{K}_S$ the segment from $y_S^{(0)}$ to $y_S^{(1)}$ lies in $\mathcal{K}_S$, and the chain rule over site blocks gives
\begin{equation*}
\begin{split}
&\bar G(y_S^{(1)})-\bar G(y_S^{(0)})\\
&\qquad=\int_0^1\sum_{x\in S}\bigl\langle\partial_{y_x}\bar G\bigl(y_S^{(0)}+t(y_S^{(1)}-y_S^{(0)})\bigr),
y_x^{(1)}-y_x^{(0)}\bigr\rangle dt .
\end{split}
\end{equation*}
Since $|y_x^{(1)}-y_x^{(0)}|\le\operatorname{diam}K_x\le2(1+\delta)r_x$,
\begin{equation*}
\operatorname{osc}_{\mathcal{K}_S}\bar G\le\sum_{x\in S}\operatorname{diam}(K_x)\|\partial_{y_x}\bar G\|_\infty
\le\sum_{x\in S}2(1+\delta)r_x\|\partial_{y_x}\bar G\|_\infty .
\end{equation*}
We combine this with the bound on $\|\partial_{y_x}\bar G\|_\infty$ and with the reduction to the oscillation. This gives \eqref{8.3:eq-correlation-decay-3} and proves (b).
\end{proof}

\begin{remark}[Scope of the covariance representation and the window]
\label{8.3:rem-representation-scope}
In this remark $\Phi$ is the potential of a measure in the class of Lemma~\ref{8.3:lem-correlation-decay}
(exponential decay of correlations in the class). We write $\mathsf{m}:=\inf\operatorname{Hess}\Phi$ for
its uniform convexity constant, as in part~(a) of that lemma, where the prefactor of the bound is
$2/\mathsf{m}$. The
letters $\mathfrak{X}$ (the set of sites),
$\operatorname{dist}$ and $\nu$ (the decay
rate) are used as in part~(a) of that lemma. The operator $-\mathcal{L}$ is the operator of the
covariance representation on which part~(a) rests, acting on scalar functions. The operator
$\mathfrak{L}$ is the operator of that
representation acting on vector-valued functions, with one component for each coordinate at each site.
Both are defined in the proof of that lemma.

\smallskip\noindent
\emph{(i) What the covariance representation is used for.} From the covariance representation the proof
of Lemma~\ref{8.3:lem-correlation-decay}(a) uses only three things:
\begin{itemize}
\item the self-adjointness of $\mathfrak{L}$, together with $\mathfrak{L}\ge\mathsf{m}$;
\item the commutation of $(-\mathcal{L})\otimes\mathbf{1}$ with rescalings of the components by a
constant at each site, which is immediate;
\item the covariance identity of the representation, displayed in the proof of that lemma.
\end{itemize}
It uses no spectral gap of $-\mathcal{L}$, no variance inequality for log-concave measures in terms of
the inverse Hessian, and no semiclassical analysis. It also uses no maximum principle, that is, no sign
property of the kernel, for Ba{\l}aban's operator $\mathbf{C}^{*}\Delta^{(j)}\mathbf{C}$; what such a
principle would supply is supplied instead by uniform convexity, namely the lower bound $\gamma_0$ on the
symmetric form of that operator which Ba{\l}aban's construction provides.

\smallskip\noindent
\emph{(ii) Why the window costs nothing here.} Here a window is the convex set in which the field at a
site is constrained to lie; $I$ is a box of sites with boundary $\partial I$, $y_{O}$ is the field outside
it, and $\sigma$, $\tau$ are the interpolation parameters of the decoupling and localisation theorem of
this section.
Consider the alternative argument that proceeds by
conditioning alone. It fixes the exterior field $y_{O}$ and must then remove the box window under the
conditional Gaussian measure. For exterior data at the edge of the window, the mean of that measure at
the sites of the box $I$ adjacent to $\partial I$ is, by the second part of the decay estimate for the
conditional Gaussian mean, of the size of the window. The complement of the
window is therefore not small there.

The proof of the decoupling and localisation theorem, by contrast, never removes a window. The convex
window at each site is kept inside
every measure of the interpolations in the parameters $\sigma$ and $\tau$, and inside every conditioning.
Only two properties are used: log-concavity, and locality of the Hessian. Neither depends on where in its
window the exterior field lies.

The box datum obtained at $\sigma=0$ still carries its own window $\mathcal{W}_I$. That window is removed
separately, inside the finite-dimensional box, where the tail bound for its complement does not depend
on the volume.

\smallskip\noindent
\emph{(iii) A second route to part~(a) of Lemma~\ref{8.3:lem-correlation-decay}.} Let
$\Phi_{\mathsf{K}}$ be the finite-volume potential of the proof of that lemma. It is $C^\infty$, it
satisfies $\mathsf{m}\le\operatorname{Hess}\Phi_{\mathsf{K}}\le\Lambda_{\mathsf{K}}$ for a finite constant
$\Lambda_{\mathsf{K}}$, the upper Hessian bound of that proof, and all its derivatives
of order at least two are bounded. Let $\mu_{\mathsf{K}}$ be the probability measure with density
proportional to $e^{-\Phi_{\mathsf{K}}}$, and let $\mathcal{L}_{\mathsf{K}}$ be the operator of that
proof,
\begin{equation*}
\mathcal{L}_{\mathsf{K}}u=\Delta u-\langle\nabla\Phi_{\mathsf{K}},\nabla u\rangle .
\end{equation*}
Let $\mathcal{B}_t$ be a standard Brownian motion. The following are standard results of the theory of
stochastic flows; in part~(iii) they are used by statement and are not proved in this book. The equation
\begin{equation}\label{8.3:eq-representation-scope-1}
dY_t=-\nabla\Phi_{\mathsf{K}}(Y_t)\,dt+\sqrt{2}\,d\mathcal{B}_t
\end{equation}
has a $C^\infty$ stochastic flow $y\mapsto Y^y_t$. Its Jacobian $J_t:=\partial_yY^y_t$ solves
\begin{equation*}
\partial_tJ_t=-\operatorname{Hess}\Phi_{\mathsf{K}}(Y^y_t)\,J_t,\qquad J_0=\mathbf{1}.
\end{equation*}
For $\tilde G\in C^2_b$ the function $u(t,y):=\mathbb{E}[\tilde G(Y^y_t)]$ lies in $C^{1,2}_b$ and
satisfies
\begin{equation}\label{8.3:eq-representation-scope-2}
\partial_tu=\mathcal{L}_{\mathsf{K}}u,\qquad
\nabla u(t,y)=\mathbb{E}\bigl[J_t^{\mathsf{T}}\nabla\tilde G(Y^y_t)\bigr].
\end{equation}
From these facts, in place of the covariance representation, one obtains the bound of
Lemma~\ref{8.3:lem-correlation-decay}(a) for $\mu_{\mathsf{K}}$, with the same constant $2/\mathsf{m}$: for
$\tilde F\in C^1_b$ and $\tilde G\in C^2_b$,
\begin{equation}\label{8.3:eq-representation-scope-3}
\bigl|\operatorname{Cov}_{\mu_{\mathsf{K}}}(\tilde F,\tilde G)\bigr|
\le\frac{2}{\mathsf{m}}\sum_{x,x'\in\mathfrak{X}}e^{-\nu\operatorname{dist}(x,x')}
\|\partial_x\tilde F\|_\infty\|\partial_{x'}\tilde G\|_\infty .
\end{equation}
The passage from $\mu_{\mathsf{K}}$ and $\tilde F,\tilde G$ to the measures of the class and the
observables of part~(a) is that of the proof of Lemma~\ref{8.3:lem-correlation-decay}, unchanged.
Whichever of the two routes is taken, part~(a) of that lemma rests on a result of analysis that is used
by statement and not proved in this book.
\end{remark}

\begin{proof}[Proof of \eqref{8.3:eq-representation-scope-3}]
The argument has four steps.

\emph{Step $(\alpha)$: pathwise decay of the Jacobian.} Conjugate by $e^{\varphi}$, with the weight
$\varphi$ used for the conjugation in the proof of Lemma~\ref{8.3:lem-correlation-decay}. Put
\begin{equation*}
\mathsf{H}(t):=\operatorname{Hess}\Phi_{\mathsf{K}}(Y^y_t),\qquad
\mathsf{E}(t):=e^{\varphi}\mathsf{H}(t)e^{-\varphi}-\mathsf{H}(t).
\end{equation*}
Then $\langle v,\mathsf{H}(t)v\rangle\ge\mathsf{m}|v|^2$, since $\operatorname{Hess}\Phi_{\mathsf{K}}\ge\mathsf{m}$,
and for this weight $\|\mathsf{E}(t)\|\le\mathsf{m}/2$; both bounds hold along every path. By the equation
for $J_t$, the matrix $\tilde J_t:=e^{\varphi}J_te^{-\varphi}$ satisfies
$\partial_t\tilde J_t=-e^{\varphi}\mathsf{H}(t)e^{-\varphi}\,\tilde J_t$, that is,
\begin{equation*}
\partial_t\tilde J_t=-(\mathsf{H}+\mathsf{E})(t)\,\tilde J_t .
\end{equation*}
Hence, for
every vector $v$,
\begin{equation*}
\partial_t|\tilde J_tv|^2=-2\langle\tilde J_tv,(\mathsf{H}+\mathsf{E})(t)\tilde J_tv\rangle
\le-2\Bigl(\mathsf{m}-\frac{\mathsf{m}}{2}\Bigr)|\tilde J_tv|^2=-\mathsf{m}|\tilde J_tv|^2 .
\end{equation*}
Since $\tilde J_0=\mathbf{1}$, this gives $|\tilde J_tv|^2\le e^{-\mathsf{m}t}|v|^2$, that is,
$\|\tilde J_t\|\le e^{-\mathsf{m}t/2}$. The passage from this bound to the blocks $(J_t)_{x'x}$ of $J_t$ is
the one made in that proof, with the weights $\varphi$ used there; it gives, for every path,
\begin{equation}\label{8.3:eq-representation-scope-4}
\|(J_t)_{x'x}\|\le e^{-\mathsf{m}t/2-\nu\operatorname{dist}(x,x')}\qquad(x,x'\in\mathfrak{X}).
\end{equation}
The equation for $J_t$ and $\operatorname{Hess}\Phi_{\mathsf{K}}\ge\mathsf{m}$ also give, for every
vector $v$,
\begin{equation*}
\partial_t|J_tv|^2=-2\langle J_tv,\operatorname{Hess}\Phi_{\mathsf{K}}(Y^y_t)J_tv\rangle
\le-2\mathsf{m}|J_tv|^2 ,
\end{equation*}
hence, since $J_0=\mathbf{1}$,
\begin{equation}\label{8.3:eq-representation-scope-5}
\|J_t\|\le e^{-\mathsf{m}t}.
\end{equation}

\emph{Step $(\beta)$: integration by parts and invariance.} Let $\tilde F\in C^1_b$ and $u\in C^2_b$.
Integrate by parts against the density $e^{-\Phi_{\mathsf{K}}}$ on large balls. Since
$\operatorname{Hess}\Phi_{\mathsf{K}}\ge\mathsf{m}$, this density decays like a Gaussian, and $\tilde F$, $u$
and $\nabla u$ are bounded; so the boundary terms vanish. The result is
\begin{equation}\label{8.3:eq-representation-scope-6}
\int\tilde F\,\mathcal{L}_{\mathsf{K}}u\,d\mu_{\mathsf{K}}
=-\int\langle\nabla\tilde F,\nabla u\rangle\,d\mu_{\mathsf{K}} .
\end{equation}
Take $\tilde F\equiv1$ and $u=u(t,\cdot)$. By \eqref{8.3:eq-representation-scope-2} and
\eqref{8.3:eq-representation-scope-6},
\begin{equation*}
\frac{d}{dt}\int u(t,\cdot)\,d\mu_{\mathsf{K}}=\int\mathcal{L}_{\mathsf{K}}u(t,\cdot)\,d\mu_{\mathsf{K}}=0 .
\end{equation*}
Thus $\mu_{\mathsf{K}}$ is invariant, and $\int u(T,\cdot)\,d\mu_{\mathsf{K}}=\mu_{\mathsf{K}}\tilde G$ for
every $T\ge0$.

\emph{Step $(\gamma)$: splitting the covariance.} Fix $T>0$. The following identity holds: expanding the
second integral and using $\int u(T,\cdot)\,d\mu_{\mathsf{K}}=\mu_{\mathsf{K}}\tilde G$ reduces its right
side to $\int\tilde F\tilde G\,d\mu_{\mathsf{K}}-\mu_{\mathsf{K}}\tilde F\,\mu_{\mathsf{K}}\tilde G$.
\begin{equation*}
\operatorname{Cov}_{\mu_{\mathsf{K}}}(\tilde F,\tilde G)
=\int\tilde F\bigl(\tilde G-u(T,\cdot)\bigr)d\mu_{\mathsf{K}}
+\int\bigl(\tilde F-\mu_{\mathsf{K}}\tilde F\bigr)\bigl(u(T,\cdot)-u(T,0)\bigr)d\mu_{\mathsf{K}} .
\end{equation*}

For the second term, \eqref{8.3:eq-representation-scope-2} and \eqref{8.3:eq-representation-scope-5} give
\begin{equation*}
|u(T,y)-u(T,0)|\le|y|\,\sup_z|\nabla u(T,z)|\le|y|\,\|\nabla\tilde G\|_\infty e^{-\mathsf{m}T}.
\end{equation*}
Together with $|\tilde F-\mu_{\mathsf{K}}\tilde F|\le2\|\tilde F\|_\infty$, the second term is bounded by
\begin{equation*}
2\|\tilde F\|_\infty\|\nabla\tilde G\|_\infty e^{-\mathsf{m}T}\int|y|\,d\mu_{\mathsf{K}}(y).
\end{equation*}
This tends to $0$ as $T\to\infty$, because the first moment is finite by the Gaussian decay of the density.

For the first term, $u(0,\cdot)=\tilde G$, so \eqref{8.3:eq-representation-scope-2} gives
\begin{equation*}
\tilde G-u(T,\cdot)=-\int_0^T\mathcal{L}_{\mathsf{K}}u(t,\cdot)\,dt .
\end{equation*}
All integrands are bounded, so we may exchange the order of integration. Applying
\eqref{8.3:eq-representation-scope-6} with $u=u(t,\cdot)$ for each $t$ then gives
\begin{equation*}
\int\tilde F\bigl(\tilde G-u(T,\cdot)\bigr)d\mu_{\mathsf{K}}
=\int_0^T\!\!\int\langle\nabla\tilde F,\nabla u(t,\cdot)\rangle\,d\mu_{\mathsf{K}}\,dt .
\end{equation*}

\emph{Step $(\delta)$: the bound.} By \eqref{8.3:eq-representation-scope-2}, the $x$-component of
$\nabla u(t,y)$ is
\begin{equation*}
\mathbb{E}\Bigl[\sum_{x'}(J_t)_{x'x}^{\mathsf{T}}\,\partial_{x'}\tilde G(Y^y_t)\Bigr].
\end{equation*}
Inserting \eqref{8.3:eq-representation-scope-4}, we obtain
\begin{align*}
\Bigl|\int\langle\nabla\tilde F,\nabla u(t,\cdot)\rangle\,d\mu_{\mathsf{K}}\Bigr|
&\le\sum_{x,x'\in\mathfrak{X}}\|\partial_x\tilde F\|_\infty\|\partial_{x'}\tilde G\|_\infty
\sup_{\text{paths}}\|(J_t)_{x'x}\|\\
&\le\sum_{x,x'\in\mathfrak{X}}e^{-\mathsf{m}t/2-\nu\operatorname{dist}(x,x')}
\|\partial_x\tilde F\|_\infty\|\partial_{x'}\tilde G\|_\infty .
\end{align*}
This bound does not involve $T$. Integrating it in $t$ over $[0,\infty)$, and using
$\int_0^\infty e^{-\mathsf{m}t/2}\,dt=2/\mathsf{m}$, bounds the first term of step $(\gamma)$ by the right side
of \eqref{8.3:eq-representation-scope-3} for every $T$. Letting $T\to\infty$ in the identity of
step $(\gamma)$, and using the bound on the second term obtained there, gives
\eqref{8.3:eq-representation-scope-3}.
\end{proof}

\begin{lemma}[Principal blocks, Schur complements and interpolations]\label{8.3:lem-principal-blocks}
Let $\mathfrak{Y}$ be a finite subset of the lattice on which the decay lemma for inverses of
positive operators with exponentially decaying blocks, stated in the section preceding this one
(in what follows, the decay lemma for inverses), is stated, with that lemma's distance
$\operatorname{dist}$, and let $n\ge1$. For $c>0$
let $S_d(c)$ be the lattice-sum constant of that lemma. It is used here only through the bound
\begin{equation*}
\sup_{x\in\mathfrak{Y}}\sum_{y\in\mathfrak{Y}}e^{-c\operatorname{dist}(x,y)}\le S_d(c),
\end{equation*}
which $\mathfrak{Y}$ inherits from that lattice.
Let $a_-$, $a_+$ and $c_A$ be positive constants. Let $A$ be a real symmetric operator on
$\bigoplus_{x\in\mathfrak{Y}}\mathbb{R}^n$, with $n\times n$ blocks $A(x,x')$, such that
\begin{equation*}
a_-\mathbf{1}\le A\le a_+\mathbf{1},\qquad
\|A(x,x')\|\le a_+e^{-c_A\operatorname{dist}(x,x')}\quad(x,x'\in\mathfrak{Y}).
\end{equation*}
Put $\Gamma:=A^{-1}$. Let $I\subset\mathfrak{Y}$ and $\mathcal{O}:=\mathfrak{Y}\setminus I$. For an
operator $T$ on $\bigoplus_{x\in\mathfrak{Y}}\mathbb{R}^n$ and $J,J'\subset\mathfrak{Y}$, write
$T_{JJ'}$ for the block of $T$ with rows in $J$ and columns in $J'$. Thus $\Gamma_{II}$ is the
restriction of $\Gamma$ to $\bigoplus_{x\in I}\mathbb{R}^n$.
Put $d_{\mathcal{O}}(x):=\operatorname{dist}(x,\mathcal{O})$. Let $\mu_A$ be the decay rate that
the decay lemma for inverses assigns to the data $(a_-,a_+,c_A)$, namely the first quantity in
\eqref{8.3:eq-principal-blocks-1}, and define $C_{sc}$ as there:
\begin{equation}\label{8.3:eq-principal-blocks-1}
\mu_A:=\min\Bigl\{\frac{c_A}{2},\ \frac{a_-c_A}{2a_+S_d(c_A)}\Bigr\},\qquad
C_{sc}:=\frac{2a_+^2}{a_-}\,S_d(c_A-\mu_A)^2 .
\end{equation}
Then:
\begin{itemize}
\item[(a)] as forms on $\bigoplus_{x\in I}\mathbb{R}^n$,
\begin{equation*}
a_-\le A_{II}\le a_+,\qquad a_-\le(\Gamma_{II})^{-1}\le A_{II}\le a_+ ;
\end{equation*}
\item[(b)] $(\Gamma_{II})^{-1}=A_{II}-S_c$, where
$S_c:=A_{I\mathcal{O}}(A_{\mathcal{O}\mathcal{O}})^{-1}A_{\mathcal{O}I}\ge0$, and for $x,x'\in I$
\begin{equation}\label{8.3:eq-principal-blocks-2}
\|S_c(x,x')\|\le C_{sc}\,e^{-(\mu_A/2)\operatorname{dist}(x,x')}\,
e^{-(\mu_A/2)(d_{\mathcal{O}}(x)+d_{\mathcal{O}}(x'))};
\end{equation}
\item[(c)] for $\sigma,\tau\in[0,1]$ the operators
\begin{equation*}
(1-\sigma)\bigl(A_{II}\oplus A_{\mathcal{O}\mathcal{O}}\bigr)+\sigma A,\qquad
(1-\tau)A_{II}+\tau(\Gamma_{II})^{-1}
\end{equation*}
lie between $a_-$ and $a_+$. Their off-diagonal blocks at $(x,x')$ are bounded in norm by
$(a_++C_{sc})e^{-\min(c_A,\mu_A/2)\operatorname{dist}(x,x')}$.
\end{itemize}
\end{lemma}

\begin{proof}
In (b) we use the decay lemma for inverses, which fixes the lattice-sum constant $S_d(\cdot)$ and
the rate function whose value at $(a_-,a_+,c_A)$ is $\mu_A$. We use it in the following form. Let
$J\subset\mathfrak{Y}$, and let $B$ be a real
symmetric operator on $\bigoplus_{y\in J}\mathbb{R}^n$ with $a_-\le B\le a_+$ and
$\|B(y,y')\|\le a_+e^{-c_A\operatorname{dist}(y,y')}$. Then
\begin{equation}\label{8.3:eq-principal-blocks-3}
\|B^{-1}(y,y')\|\le\frac{2}{a_-}\,e^{-\mu_A\operatorname{dist}(y,y')}\qquad(y,y'\in J),
\end{equation}
with $\mu_A$ as in \eqref{8.3:eq-principal-blocks-1}. The constant $2/a_-$ and the rate $\mu_A$ in
\eqref{8.3:eq-principal-blocks-3} are those of the decay lemma for inverses; this bound is not taken
from Lemma~\ref{8.3:lem-weighted-inverse-decay}, whose bound for inverses carries the constant
$4/(3a_-)$ and a rate subject to that lemma's own smallness condition.

\emph{(a)} Let $\xi\in\bigoplus_{x\in I}\mathbb{R}^n$, extended by zero to $\mathfrak{Y}$. Then
$\langle\xi,A_{II}\xi\rangle=\langle\xi,A\xi\rangle$, so the form bounds of $A$ restrict to
$a_-\le A_{II}\le a_+$; in the same way $a_-\le A_{\mathcal{O}\mathcal{O}}\le a_+$. In particular
$A_{\mathcal{O}\mathcal{O}}$ is positive definite, hence invertible. The block-inversion formula
for the decomposition $\mathfrak{Y}=I\sqcup\mathcal{O}$ therefore gives
\begin{equation*}
\Gamma_{II}=\bigl(A_{II}-A_{I\mathcal{O}}A_{\mathcal{O}\mathcal{O}}^{-1}A_{\mathcal{O}I}\bigr)^{-1},
\qquad\text{that is,}\qquad (\Gamma_{II})^{-1}=A_{II}-S_c .
\end{equation*}
Since $A$ is symmetric, $A_{\mathcal{O}I}=(A_{I\mathcal{O}})^{\mathsf T}$. Hence
$\langle\xi,S_c\xi\rangle=\langle\eta,A_{\mathcal{O}\mathcal{O}}^{-1}\eta\rangle$ with
$\eta=A_{\mathcal{O}I}\xi$. As $A_{\mathcal{O}\mathcal{O}}>0$, also $A_{\mathcal{O}\mathcal{O}}^{-1}>0$,
and therefore $S_c\ge0$. Hence $(\Gamma_{II})^{-1}\le A_{II}\le a_+$. For the lower bound, note that
$\Gamma_{II}$ is the restriction of the positive operator $\Gamma$, and
\begin{equation*}
\Gamma_{II}\le\|\Gamma\|\,\mathbf{1}=\lambda_{\min}(A)^{-1}\mathbf{1}\le a_-^{-1}\mathbf{1}.
\end{equation*}
Since inversion reverses the order of positive definite operators, $(\Gamma_{II})^{-1}\ge a_-$.

\emph{(b)} The identity $(\Gamma_{II})^{-1}=A_{II}-S_c$ and the bound $S_c\ge0$ were shown in (a).
The principal block $B=A_{\mathcal{O}\mathcal{O}}$ satisfies the hypotheses of
\eqref{8.3:eq-principal-blocks-3} with the same constants: $a_-\le A_{\mathcal{O}\mathcal{O}}\le a_+$
by the proof of (a), and its blocks are blocks of $A$. Hence
$\|(A_{\mathcal{O}\mathcal{O}})^{-1}(o,o')\|\le(2/a_-)e^{-\mu_A\operatorname{dist}(o,o')}$ for
$o,o'\in\mathcal{O}$. Writing out the block product and using the entry bound of $A$, we obtain for
$x,x'\in I$
\begin{equation*}
\|S_c(x,x')\|\le\sum_{o,o'\in\mathcal{O}}a_+e^{-c_A\operatorname{dist}(x,o)}\,
\frac{2}{a_-}e^{-\mu_A\operatorname{dist}(o,o')}\,a_+e^{-c_A\operatorname{dist}(o',x')} .
\end{equation*}
Since $\mu_A\le c_A/2$, we may split the first factor as
\begin{equation*}
e^{-c_A\operatorname{dist}(x,o)}
=e^{-\mu_A\operatorname{dist}(x,o)}\,e^{-(c_A-\mu_A)\operatorname{dist}(x,o)},
\end{equation*}
and likewise the last factor. The factors carrying $\mu_A$ combine to $e^{-\mu_A\ell}$ with
$\ell:=\operatorname{dist}(x,o)+\operatorname{dist}(o,o')+\operatorname{dist}(o',x')$. By the
triangle inequality, $\ell\ge\operatorname{dist}(x,x')$. Since $o,o'\in\mathcal{O}$, also
\begin{equation*}
\ell\ge\operatorname{dist}(x,o)+\operatorname{dist}(o',x')\ge d_{\mathcal{O}}(x)+d_{\mathcal{O}}(x').
\end{equation*}
Taking half of each of these two lower bounds,
\begin{equation*}
e^{-\mu_A\ell}\le e^{-(\mu_A/2)\operatorname{dist}(x,x')}\,
e^{-(\mu_A/2)(d_{\mathcal{O}}(x)+d_{\mathcal{O}}(x'))},
\end{equation*}
uniformly in $o,o'$. Since $c_A-\mu_A>0$, the property of $S_d$ recalled in the statement bounds the
remaining sum:
\begin{equation*}
\sum_{o\in\mathcal{O}}e^{-(c_A-\mu_A)\operatorname{dist}(x,o)}
\sum_{o'\in\mathcal{O}}e^{-(c_A-\mu_A)\operatorname{dist}(o',x')}\le S_d(c_A-\mu_A)^2 .
\end{equation*}
Multiplying by the prefactor $2a_+^2/a_-$ gives \eqref{8.3:eq-principal-blocks-2} with $C_{sc}$ as in
\eqref{8.3:eq-principal-blocks-1}.

\emph{(c)} By (a) and the restriction argument in its proof, $A_{II}$ and $A_{\mathcal{O}\mathcal{O}}$
lie between $a_-$ and $a_+$; hence so does $A_{II}\oplus A_{\mathcal{O}\mathcal{O}}$. The operator
$A$ does so by hypothesis, and $A_{II}$ and $(\Gamma_{II})^{-1}$ do so by (a). A convex combination
of operators each lying between $a_-$ and $a_+$ lies between $a_-$ and $a_+$.

For the blocks, consider first $(1-\sigma)(A_{II}\oplus A_{\mathcal{O}\mathcal{O}})+\sigma A$. Its
block at $(x,x')$ equals $A(x,x')$ if $x,x'$ both lie in $I$ or both in $\mathcal{O}$, and
$\sigma A(x,x')$ otherwise. In either case its norm is at most
$a_+e^{-c_A\operatorname{dist}(x,x')}$. Next, by (b),
$(1-\tau)A_{II}+\tau(\Gamma_{II})^{-1}=A_{II}-\tau S_c$. By the entry bound of $A$ and
\eqref{8.3:eq-principal-blocks-2}, dropping the factor involving $d_{\mathcal{O}}$, which is at most
one, the block of this operator at $(x,x')$ has norm at most
\begin{equation*}
a_+e^{-c_A\operatorname{dist}(x,x')}+C_{sc}e^{-(\mu_A/2)\operatorname{dist}(x,x')}.
\end{equation*}
Both bounds are at most $(a_++C_{sc})e^{-\min(c_A,\mu_A/2)\operatorname{dist}(x,x')}$.
\end{proof}

\begin{lemma}[Boundary profile of the coupling observable]\label{8.3:lem-boundary-profile}
Let the partition $\mathfrak{X}=I\sqcup O$, the frozen sites $F$ with their variables
$z=(z_f)_{f\in F}$, the measure $\mu\in\mathfrak{S}$ with data $(Q,\ell^0,R,P)$ and the coupling
observable $\mathfrak{W}$ be as in Definition~\ref{8.3:def-decoupled-family}. Assume the
separation condition of that definition, $\operatorname{dist}(x,f)\ge\operatorname{dist}(x,O)$
for all $x\in I$ and $f\in F$. Let $q_+$, $c_Q$, $p_1$, $p_2$ and $c_P$ be the
constants, and $S_d(\cdot)$ the lattice-sum bound, of the class $\mathfrak{S}$. Let
$\mathcal{K}^+$ be the enlarged window of the variables $y$ and $\mathcal{K}_F$ the window of the
frozen variables $z$, and let $\bar r>0$ be such that every coordinate $y_u$ of a point of
$\mathcal{K}^+$ and every coordinate $z_f$ of a point of $\mathcal{K}_F$ has norm at most
$\bar r$. Let $c_1>0$ and $\varpi_0>0$ satisfy
\begin{equation*}
c_1\le\min\{c_Q,c_P\},\qquad
\varpi_0\ge 2\bar r\bigl(q_+S_d(c_Q/2)+p_2S_d(c_P/2)\bigr).
\end{equation*}
For $x\in I$ put
$d_O(x):=\operatorname{dist}(x,O)$ and $d_F(x):=\operatorname{dist}(x,F)$, and for $o\in O$
put $d_I(o):=\operatorname{dist}(o,I)$. Then on
$\mathcal{K}^+\times\mathcal{K}_F$ the following hold.
\begin{enumerate}
\item[(a)] For $x\in I$,
\begin{equation}\label{8.3:eq-boundary-profile-1}
\begin{aligned}
\|\partial_x\mathfrak{W}\|
&\le \sum_{f\in F}\|R_{xf}\|\,\bar r+\sum_{o\in O}\|Q_{xo}\|\,\bar r
 +\bigl\|\partial_xP(y_I,y_O,z)-\partial_xP(y_I,0_O,0_F)\bigr\|\\
&\le \bar r q_+S_d(c_Q/2)\bigl[e^{-(c_Q/2)d_F(x)}+e^{-(c_Q/2)d_O(x)}\bigr]\\
&\quad+\bar r p_2S_d(c_P/2)\bigl[e^{-(c_P/2)d_O(x)}+e^{-(c_P/2)d_F(x)}\bigr]
 \le \varpi_0e^{-(c_1/2)d_O(x)}.
\end{aligned}
\end{equation}
\item[(b)] For $o\in O$,
\begin{equation}\label{8.3:eq-boundary-profile-2}
\|\partial_o\mathfrak{W}\|
\le \bar r\bigl(q_+S_d(c_Q/2)+p_2S_d(c_P/2)\bigr)e^{-(c_1/2)d_I(o)}
\le \varpi_0e^{-(c_1/2)d_I(o)}.
\end{equation}
\item[(c)] On the $I$-coordinate part $\mathcal{K}_I^+$ of $\mathcal{K}^+$,
\begin{equation}\label{8.3:eq-boundary-profile-3}
\sup_{\mathcal{K}_I^+}\bigl|P(y_I,0_O,0_F)-P(0,0,0)\bigr|
\le p_1\sum_{x\in I}\bar r=p_1\bar r\,|I|.
\end{equation}
\end{enumerate}
Consequently, with $d_\partial:=d_O$ on $I$, $d_\partial:=d_I$ on $O$, and
$\omega(x):=\varpi_0e^{-(c_1/2)d_\partial(x)}$, one has
$\|\partial_x\mathfrak{W}\|_\infty\le\omega(x)$ for every $x\in\mathfrak{X}$, uniformly in $z$.
In the case $O=\emptyset$, \textup{(a)} reads
\begin{equation}\label{8.3:eq-boundary-profile-4}
\begin{aligned}
\|\partial_x\mathfrak{W}\|
&\le \bar r q_+S_d(c_Q/2)e^{-(c_Q/2)d_F(x)}+\bar r p_2S_d(c_P/2)e^{-(c_P/2)d_F(x)}\\
&\le \varpi_0e^{-(c_1/2)d_F(x)},
\end{aligned}
\end{equation}
\textup{(b)} is void, and $\omega(x):=\varpi_0e^{-(c_1/2)d_F(x)}$.
\end{lemma}

\begin{proof}
We recall the properties that enter. By the choice of $\bar r$, every
coordinate $y_u$ of a point of $\mathcal{K}^+$ and every coordinate $z_f$ of a point of
$\mathcal{K}_F$ has norm at most $\bar r$. From the definition of the class $\mathfrak{S}$, the
entries satisfy
$\|Q_{uv}\|\le q_+e^{-c_Q\operatorname{dist}(u,v)}$ and
$\|R_{uf}\|\le q_+e^{-c_Q\operatorname{dist}(u,f)}$. The first derivatives of $P$ are bounded
by $p_1$, and the mixed second derivatives satisfy
$\|\partial_u\partial_vP\|\le p_2e^{-c_P\operatorname{dist}(u,v)}$, the variables $z_f$ counting
as sites. For $c>0$, the lattice sums $\sum_v e^{-c\operatorname{dist}(u,v)}$ are bounded by
$S_d(c)$. From the last property we obtain, for a site $u$, a set $Y$ of sites and $\delta\ge0$
with $\operatorname{dist}(u,y)\ge\delta$ for all $y\in Y$, and for $c>0$,
\begin{equation}\label{8.3:eq-boundary-profile-5}
\begin{aligned}
\sum_{y\in Y}e^{-c\operatorname{dist}(u,y)}
&=\sum_{y\in Y}e^{-(c/2)\operatorname{dist}(u,y)}e^{-(c/2)\operatorname{dist}(u,y)}\\
&\le e^{-(c/2)\delta}\sum_{y\in Y}e^{-(c/2)\operatorname{dist}(u,y)}
\le S_d(c/2)\,e^{-(c/2)\delta}.
\end{aligned}
\end{equation}
Recall from Definition~\ref{8.3:def-decoupled-family} that
\begin{equation*}
\mathfrak{W}=-\langle y_I,Q_{IO}y_O\rangle-\langle(Rz)_I,y_I\rangle
+\bigl[P(y_I,y_O,z)-P(y_I,0_O,0_F)-P(0_I,y_O,z)\bigr].
\end{equation*}

\emph{(a)} Let $x\in I$. The term $P(0_I,y_O,z)$ does not depend on $y_x$. Hence
\begin{equation*}
\partial_x\mathfrak{W}=-\sum_{o\in O}Q_{xo}y_o-\sum_{f\in F}R_{xf}z_f
+\partial_xP(y_I,y_O,z)-\partial_xP(y_I,0_O,0_F).
\end{equation*}
The triangle inequality and the bound $\bar r$ on the coordinates $y_o$ and $z_f$ give the
first inequality of \eqref{8.3:eq-boundary-profile-1}. For the $P$-term we write
\begin{equation*}
\begin{aligned}
\partial_xP(y_I,y_O,z)-\partial_xP(y_I,0_O,0_F)
&=\int_0^1\frac{d}{ds}\,\partial_xP(y_I,sy_O,sz)\,ds\\
&=\int_0^1\Bigl[\sum_{o\in O}(\partial_o\partial_xP)(y_I,sy_O,sz)[y_o]
+\sum_{f\in F}(\partial_{z_f}\partial_xP)(y_I,sy_O,sz)[z_f]\Bigr]ds.
\end{aligned}
\end{equation*}
The segment $s\mapsto(y_I,sy_O,sz)$ joins $(y_I,0_O,0_F)$ to $(y_I,y_O,z)$. It lies in
$\mathcal{K}^+\times\mathcal{K}_F$, which is convex and contains $0$, so the bounds of
$\mathfrak{S}$ apply along it. Each term of the integrand is therefore at most
$p_2e^{-c_P\operatorname{dist}(x,o)}\bar r$, respectively
$p_2e^{-c_P\operatorname{dist}(x,f)}\bar r$.

We apply \eqref{8.3:eq-boundary-profile-5} with $u=x$ four times. For the $R$-sum and the
$z$-part of the $P$-term we take $Y=F$, $\delta=d_F(x)$ and $c=c_Q$, respectively $c=c_P$. For
the $Q$-sum and the $y_O$-part of the $P$-term we take $Y=O$, $\delta=d_O(x)$ and $c=c_Q$,
respectively $c=c_P$. This gives the second inequality of \eqref{8.3:eq-boundary-profile-1}.

For the last inequality, the separation condition gives $\operatorname{dist}(x,f)\ge d_O(x)$
for every $f\in F$, hence $d_F(x)\ge d_O(x)$. Each exponential in the bracket with $d_F(x)$ is
therefore at most the corresponding exponential with $d_O(x)$. The middle member of
\eqref{8.3:eq-boundary-profile-1} is thus at most
\begin{equation*}
2\bar r q_+S_d(c_Q/2)e^{-(c_Q/2)d_O(x)}+2\bar r p_2S_d(c_P/2)e^{-(c_P/2)d_O(x)}.
\end{equation*}
Since $d_O(x)\ge0$ and $c_1\le\min\{c_Q,c_P\}$, both $e^{-(c_Q/2)d_O(x)}$ and
$e^{-(c_P/2)d_O(x)}$ are at most $e^{-(c_1/2)d_O(x)}$; since moreover
$2\bar r(q_+S_d(c_Q/2)+p_2S_d(c_P/2))\le\varpi_0$, the last display is at most
$\varpi_0e^{-(c_1/2)d_O(x)}$.

\emph{(b)} Let $o\in O$. The term $\langle(Rz)_I,y_I\rangle$ does not depend on $y_o$, so its
$\partial_o$-derivative is $0$. The piece $P(y_I,0_O,0_F)$ is also free of $y_o$. Hence
\begin{equation*}
\partial_o\mathfrak{W}=-\sum_{x\in I}Q_{ox}y_x+\partial_oP(y_I,y_O,z)-\partial_oP(0_I,y_O,z).
\end{equation*}
The first sum has norm at most $\sum_{x\in I}\|Q_{ox}\|\,\bar r$. By
\eqref{8.3:eq-boundary-profile-5} with $u=o$, $Y=I$, $\delta=d_I(o)$ and $c=c_Q$, this is at most
$\bar r q_+S_d(c_Q/2)e^{-(c_Q/2)d_I(o)}$.

The $P$-difference is treated as in (a), by the segment in $y_I$ only:
\begin{equation*}
\partial_oP(y_I,y_O,z)-\partial_oP(0_I,y_O,z)
=\int_0^1\sum_{x\in I}(\partial_x\partial_oP)(sy_I,y_O,z)[y_x]\,ds.
\end{equation*}
Each term of the integrand is at most $p_2e^{-c_P\operatorname{dist}(o,x)}\bar r$. By
\eqref{8.3:eq-boundary-profile-5} with $u=o$, $Y=I$, $\delta=d_I(o)$ and $c=c_P$, the integral
is at most $\bar r p_2S_d(c_P/2)e^{-(c_P/2)d_I(o)}$.

Adding the two bounds, and using that $c_1\le\min\{c_Q,c_P\}$ and $d_I(o)\ge0$ give
$e^{-(c_Q/2)d_I(o)}\le e^{-(c_1/2)d_I(o)}$ and $e^{-(c_P/2)d_I(o)}\le e^{-(c_1/2)d_I(o)}$, gives
the first inequality of \eqref{8.3:eq-boundary-profile-2}. The second follows from
$\bar r(q_+S_d(c_Q/2)+p_2S_d(c_P/2))\le\varpi_0$.

\emph{(c)} For $y_I\in\mathcal{K}_I^+$ we use the segment from $0$:
\begin{equation*}
P(y_I,0_O,0_F)-P(0,0,0)
=\int_0^1\sum_{x\in I}(\partial_xP)(sy_I,0_O,0_F)[y_x]\,ds.
\end{equation*}
Each term is at most $p_1\bar r$, which gives \eqref{8.3:eq-boundary-profile-3}.

\emph{Conclusion.} Bounds (a) and (b) hold at every point of $\mathcal{K}^+\times\mathcal{K}_F$,
and their right-hand sides do not depend on the point. Taking the supremum gives
$\|\partial_x\mathfrak{W}\|_\infty\le\varpi_0e^{-(c_1/2)d_O(x)}=\omega(x)$ for $x\in I$, and
$\|\partial_o\mathfrak{W}\|_\infty\le\varpi_0e^{-(c_1/2)d_I(o)}=\omega(o)$ for $o\in O$. The
bound is uniform in $z\in\mathcal{K}_F$.

If $O=\emptyset$, the sums over $O$ in (a) are empty. The computation of (a) then leaves only the
$R$-sum and the $z$-part of the $P$-term, which gives the first inequality of
\eqref{8.3:eq-boundary-profile-4}. The second inequality follows since $c_1\le\min\{c_Q,c_P\}$
and $d_F(x)\ge0$ bound both exponentials by $e^{-(c_1/2)d_F(x)}$, and
$\bar r(q_+S_d(c_Q/2)+p_2S_d(c_P/2))\le\varpi_0$. There is no site $o\in O$, so (b) is void,
and the profile becomes $\omega(x)=\varpi_0e^{-(c_1/2)d_F(x)}$.
\end{proof}

Throughout Theorem~\ref{8.3:thm-decoupling-localisation} and its proof, $\mathfrak{S}$ is the class of weakly perturbed Gaussian measures with
convex per-site windows in which Lemma~\ref{8.3:lem-correlation-decay} is stated, and we use its
notation:
\begin{itemize}
\item the finite site set $\mathfrak{X}$, with its metric $\operatorname{dist}$ and dimension $d$;
\item each site $x$ carries a copy of $\mathbb{R}^n$ with its Euclidean norm, and
$y=(y_x)_{x\in\mathfrak{X}}$;
\item the windows $K_x$ and their enlargements $K_x^+$ by $\delta r_-$, the products $\mathcal{K}$
and $\mathcal{K}^+$, and, for a set $T$ of sites, the products $\mathcal{K}_T$ and $\mathcal{K}_T^+$
over $T$;
\item the radii $r_x\in[r_-,r_+]$ and $\delta\in(0,1]$;
\item the frozen exterior $F$, with frozen values $z\in\mathcal{K}_F$;
\item the data $Q$, $\ell^0$, $R$, $P$ of the energy $\Phi$, with the constants $q_\pm$, $c_Q$,
$p_1$, $p_2$, $c_P$ (the function $P$ is of class $C^2$, and $C^\infty$ in the application made
later in this chapter);
\item the lattice sum $S_d(c)$ and the smallness condition $p_2S_d(c_P)\le q_-/4$;
\item the class constants $m_{\mathfrak{S}}$, $\kappa(x,x')$, $\nu_0$, $\mathfrak{K}$, $\nu$ and
$\bar{\mathfrak{K}}=(q_++p_2)S_d(\nu_0)\ge\mathfrak{K}$.
\end{itemize}

Given a partition $\mathfrak{X}=I\sqcup O$ satisfying the separation condition $\operatorname{dist}(x,f)\ge d_O(x):=\operatorname{dist}(x,O)$ for all $x\in I$ and $f\in F$, Definition~\ref{8.3:def-decoupled-family} provides the following objects:
\begin{itemize}
\item the decoupled family $(\mu_\sigma)_{\sigma\in[0,1]}$, with energies $\Phi^\sigma=\tfrac12\langle y,Q^\sigma y\rangle+\langle\ell^\sigma,y\rangle-P^\sigma(y,z)$ affine in $\sigma$;
\item the end points $\mu_1=\mu$ and $\mu_0=\mu^{\mathrm{bx}}\otimes\mu^{\mathrm{out}}$, where $\mu^{\mathrm{bx}}$ is the box datum on the sites $I$ and $\mu^{\mathrm{out}}$ is the factor on $O$;
\item the coupling observable $\mathfrak{W}:=-\partial_\sigma\Phi^\sigma$, which does not depend on $\sigma$.
\end{itemize}

For a probability measure $\mu'$ we write $\langle X\rangle_{\mu'}$ for the expectation of $X$ and $\operatorname{Cov}_{\mu'}(X,Y):=\langle XY\rangle_{\mu'}-\langle X\rangle_{\mu'}\langle Y\rangle_{\mu'}$. We abbreviate $\langle\cdot\rangle_\sigma:=\langle\cdot\rangle_{\mu_\sigma}$; with $\bar r$ as in Theorem~\ref{8.3:thm-decoupling-localisation}, note that $|y_x|\le\bar r$ on $\mathcal{K}^+$.

\begin{theorem}[Box localisation by decoupling interpolation]\label{8.3:thm-decoupling-localisation}
Let $\mu\in\mathfrak{S}$. Let $\mathfrak{X}=I\sqcup O$ be a partition satisfying the separation condition, with the decoupled family of Definition~\ref{8.3:def-decoupled-family}. Let $S\subset I$ and $\ell_S:=\operatorname{dist}(S,O)$. Put
\begin{gather*}
c_1:=\min(c_Q,c_P),\qquad \bar r:=(1+\delta)r_+,\qquad
\omega_0:=4\bar r(q_++p_2)S_d(c_1/2),\qquad \nu_1:=\tfrac12\min(\nu,c_1/2),\\
C_{\mathrm{lat}}=C_{\mathrm{lat}}(d,\nu):=2\max\{1,\,2S_d(\nu/2)\}.
\end{gather*}
Then the following hold, uniformly in $\sigma\in[0,1]$ and in $z\in\mathcal{K}_F$.
\begin{enumerate}
\item[(i)] For every bounded measurable $G=G(y_S)$ and every constant $c$,
\begin{equation}\label{8.3:eq-decoupling-localisation-1}
\begin{aligned}
\bigl|\langle G\rangle_\mu-\langle G\rangle_{\mu^{\mathrm{bx}}}\bigr|
&\le\sup_{\sigma\in[0,1]}\bigl|\operatorname{Cov}_{\mu_\sigma}(G,\mathfrak{W})\bigr|\\
&\le C_{\mathrm{lat}}\,\|G-c\|_\infty\,|S|\,\bar r\,\omega_0
\bigl(1+\bar{\mathfrak{K}}/m_{\mathfrak{S}}\bigr)e^{-\nu_1\ell_S}.
\end{aligned}
\end{equation}
\item[(ii)] For bounded measurable functions $G_1,G_2$ of $y_S$ and constants $c',c''$,
\begin{equation}\label{8.3:eq-decoupling-localisation-2}
\begin{aligned}
&\bigl|\operatorname{Cov}_\mu(G_1,G_2)-\operatorname{Cov}_{\mu^{\mathrm{bx}}}(G_1,G_2)\bigr|\\
&\qquad\le 4C_{\mathrm{lat}}\,\|G_1-c'\|_\infty\|G_2-c''\|_\infty\,|S|\,\bar r\,\omega_0
\bigl(1+\bar{\mathfrak{K}}/m_{\mathfrak{S}}\bigr)e^{-\nu_1\ell_S}.
\end{aligned}
\end{equation}
\item[(iii)] Suppose the following.
\begin{itemize}
\item $Q=A_{\mathfrak{X}\mathfrak{X}}$ is the principal block of a real symmetric $A$ on $\bigoplus_{x\in\mathfrak{Y}}\mathbb{R}^n$, for a larger finite site set $\mathfrak{Y}\supset\mathfrak{X}$.
\item $a_-\le A\le a_+$ and $\|A(x,x')\|\le a_+e^{-c_A\operatorname{dist}(x,x')}$, so that $q_-=a_-$, $q_+=a_+$ and $c_Q=c_A$.
\end{itemize}
Let $\Gamma:=A^{-1}$ and $\mathcal{O}:=\mathfrak{Y}\setminus I$, and let $\Gamma_{II}$ be the principal block of $\Gamma$ on $I$. Define $\mu^{\mathrm{mg}}$ as $\mu^{\mathrm{bx}}$ with $Q_{II}=A_{II}$ replaced by
\[(\Gamma_{II})^{-1}=A_{II}-A_{I\mathcal{O}}(A_{\mathcal{O}\mathcal{O}})^{-1}A_{\mathcal{O}I}\]
(Lemma~\ref{8.3:lem-principal-blocks}), that is,
\begin{equation*}
\mu^{\mathrm{mg}}(dy_I):=Z_{\mathrm{mg}}^{-1}\mathbf{1}_{\mathcal{K}_I}(y_I)
\exp\Bigl[-\tfrac12\langle y_I,(\Gamma_{II})^{-1}y_I\rangle-\langle\ell^0_I,y_I\rangle
+P(y_I,0_O,0_F)\Bigr]dy_I.
\end{equation*}
Then (i) and (ii) hold with $\mu^{\mathrm{mg}}$ in place of $\mu^{\mathrm{bx}}$, and also for the pair $(\mu^{\mathrm{bx}},\mu^{\mathrm{mg}})$, with the following changes:
\begin{itemize}
\item $\ell_S:=\operatorname{dist}(S,\mathfrak{Y}\setminus I)$; in the application made later in
this chapter, with $O=\emptyset$ as in the remark following the proof,
$\mathfrak{Y}\setminus I$ is the complement of the box $I$, the distance to which equals $d_F$ on
$I$, so that this $\ell_S$ is $\operatorname{dist}(S,F)$, the same $\ell_S$ as in that remark;
\item $(\omega_0,c_1)$ is replaced by $\bigl(\omega_0+\bar rC_{sc}S_d(\mu_A/2),\ \min(c_1,\mu_A/2)\bigr)$;
\item the class constants $(\mathfrak{K},\bar{\mathfrak{K}},\nu)$ are recomputed for $q_+\mapsto a_++C_{sc}$ and $c_Q\mapsto\min(c_A,\mu_A/2)$.
\end{itemize}
Here $C_{sc}$ and $\mu_A$ are the constants of Lemma~\ref{8.3:lem-principal-blocks}, functions of $a_+/a_-$, $c_A$ and $d$ only.
\end{enumerate}
All constants are free of $|\mathfrak{X}|$, $|\mathfrak{Y}|$, $|O|$, $|F|$, $\ell^0$, $z$ and $\sigma$. They are free of the radii except through $\bar r$ and the smallness condition $p_2S_d(c_P)\le q_-/4$.
\end{theorem}

\begin{proof}
\emph{Part (i), the interpolation formula.} Fix $z\in\mathcal{K}_F$ and a bounded measurable function $\Psi$ of $y$. For every $\sigma$,
\[\langle \Psi\rangle_\sigma=\frac{\int_{\mathcal{K}}\Psi e^{-\Phi^\sigma}\,dy}{\int_{\mathcal{K}}e^{-\Phi^\sigma}\,dy}.\]
Both integrals are over the fixed bounded set $\mathcal{K}$. The energy $\Phi^\sigma$ is affine in $\sigma$ with $\partial_\sigma\Phi^\sigma=-\mathfrak{W}$, so $\Phi^\sigma=\Phi^0-\sigma\mathfrak{W}$.

The observable $\mathfrak{W}$ is continuous on the compact closure of $\mathcal{K}$, which lies in $\mathcal{K}^+$; hence it is bounded on $\mathcal{K}$. The integrands $\Psi e^{-\Phi^\sigma}$ and $e^{-\Phi^\sigma}$ are therefore bounded and smooth in $\sigma$, and both integrals may be differentiated under the integral sign. Differentiating the ratio gives
\begin{equation}\label{8.3:eq-decoupling-localisation-3}
\begin{split}
\frac{d}{d\sigma}\langle \Psi\rangle_\sigma
&=\langle \Psi\mathfrak{W}\rangle_\sigma-\langle \Psi\rangle_\sigma\langle\mathfrak{W}\rangle_\sigma\\
&=-\operatorname{Cov}_{\mu_\sigma}(\Psi,\partial_\sigma\Phi^\sigma)
=\operatorname{Cov}_{\mu_\sigma}(\Psi,\mathfrak{W}),
\end{split}
\end{equation}
and the right side is continuous in $\sigma$. Take $\Psi=G$.

At $\sigma=1$ we have $\mu_1=\mu$. At $\sigma=0$ we have $\mu_0=\mu^{\mathrm{bx}}\otimes\mu^{\mathrm{out}}$, and $G$ depends on $y_S$ only, with $S\subset I$. Integrating out $y_O$ against the probability measure $\mu^{\mathrm{out}}$ gives $\langle G\rangle_{\mu_0}=\langle G\rangle_{\mu^{\mathrm{bx}}}$. Hence
\begin{equation}\label{8.3:eq-decoupling-localisation-4}
\langle G\rangle_\mu-\langle G\rangle_{\mu^{\mathrm{bx}}}=\int_0^1\operatorname{Cov}_{\mu_\sigma}(G,\mathfrak{W})\,d\sigma,
\end{equation}
which gives the first inequality of \eqref{8.3:eq-decoupling-localisation-1}.

\emph{Part (i), the covariance bound.} For each $\sigma$, Definition~\ref{8.3:def-decoupled-family}
gives $\mu_\sigma\in\mathfrak{S}$ with the same class constants as $\mu$. Put
$\omega(x):=\sup\|\partial_x\mathfrak{W}\|$, the supremum being taken over
$\mathcal{K}^+\times\mathcal{K}_F$. Lemma~\ref{8.3:lem-boundary-profile} gives
\[\omega(x)\le\omega_0e^{-(c_1/2)d_O(x)}\quad(x\in I),\qquad\omega(x)\le\omega_0\quad(x\in O);\]
the first of these bounds is part (a) of that lemma.
We apply part (b) of Lemma~\ref{8.3:lem-correlation-decay}. We use it for $\mu_\sigma$, with the
bounded measurable function $G$ of $y_S$, the constant $c$, and the second function
$\mathfrak{W}$, which is $C^\infty$ on $\mathcal{K}^+$ when $P$ is, as in the application made later
in this chapter. Together with the profile bound this gives
\begin{equation}\label{8.3:eq-decoupling-localisation-5}
\begin{aligned}
\bigl|\operatorname{Cov}_{\mu_\sigma}(G,\mathfrak{W})\bigr|
&\le\|G-c\|_\infty\sum_{x\in S}2\bar r\Bigl[\omega(x)\\
&\qquad+\frac{2}{m_{\mathfrak{S}}}\sum_{x'',x'''\in\mathfrak{X}\setminus S}e^{-\nu\operatorname{dist}(x'',x''')}\,
\omega(x'')\,\kappa(x,x''')\Bigr].
\end{aligned}
\end{equation}

We estimate the two terms in the bracket. First, let $x\in S$. Then $d_O(x)\ge\operatorname{dist}(S,O)=\ell_S$, so $\omega(x)\le\omega_0e^{-(c_1/2)\ell_S}$.

For the double sum, put $\nu_2:=\min(c_1/2,\nu/2,\nu_0)$, and let $x\in S$ and $x'',x'''\in\mathfrak{X}$. We claim
\begin{equation}\label{8.3:eq-decoupling-localisation-6}
\omega(x'')\le\omega_0e^{-\nu_2\ell_S}e^{\nu_2(\operatorname{dist}(x,x''')+\operatorname{dist}(x''',x''))}.
\end{equation}
\begin{itemize}
\item If $x''\in I$, the triangle inequality gives
\[d_O(x'')\ge d_O(x)-\operatorname{dist}(x,x'')\ge\ell_S-\operatorname{dist}(x,x''')-\operatorname{dist}(x''',x'').\]
Since $\nu_2\le c_1/2$ and $d_O(x'')\ge0$, we get $\omega(x'')\le\omega_0e^{-(c_1/2)d_O(x'')}\le\omega_0e^{-\nu_2d_O(x'')}$. Inserting the lower bound for $d_O(x'')$ gives \eqref{8.3:eq-decoupling-localisation-6}.
\item If $x''\in O$, then $\operatorname{dist}(x,x'')\ge\ell_S$. Combined with $\operatorname{dist}(x,x'')\le\operatorname{dist}(x,x''')+\operatorname{dist}(x''',x'')$, this shows that the right side of \eqref{8.3:eq-decoupling-localisation-6} is at least $\omega_0$. It therefore dominates $\omega(x'')$, by the bound $\omega(x'')\le\omega_0$ on $O$ recalled above.
\end{itemize}

Next, $c_Q\ge c_1=2\nu_0$ and $c_P\ge 2\nu_0$ by the definitions of $c_1$ and $\nu_0$. Hence, for
$x'''\ne x$,
\[\kappa(x,x''')\le(q_++p_2)e^{-2\nu_0\operatorname{dist}(x,x''')}.\]
In the double sum we have $x'''\in\mathfrak{X}\setminus S$ and $x\in S$, so $x'''\ne x$. We insert
this bound and \eqref{8.3:eq-decoupling-localisation-6} into the double sum, whose terms are
non-negative, and then enlarge the ranges of $x''$ and $x'''$ to all of $\mathfrak{X}$. This gives
\begin{align*}
&\sum_{x'',x'''\in\mathfrak{X}\setminus S}e^{-\nu\operatorname{dist}(x'',x''')}\omega(x'')\kappa(x,x''')\\
&\quad\le\omega_0e^{-\nu_2\ell_S}\sum_{x''',x''\in\mathfrak{X}}
e^{-\nu\operatorname{dist}(x'',x''')+\nu_2\operatorname{dist}(x''',x'')}\\
&\qquad\qquad\times(q_++p_2)e^{-2\nu_0\operatorname{dist}(x,x''')}\,e^{\nu_2\operatorname{dist}(x,x''')}\\
&\quad\le\omega_0e^{-\nu_2\ell_S}\sum_{x'''}(q_++p_2)e^{-(2\nu_0-\nu_2)\operatorname{dist}(x,x''')}
\sum_{x''}e^{-(\nu-\nu_2)\operatorname{dist}(x'',x''')}\\
&\quad\le\omega_0e^{-\nu_2\ell_S}(q_++p_2)S_d(\nu_0)S_d(\nu/2)
=\omega_0e^{-\nu_2\ell_S}\,\bar{\mathfrak{K}}\,S_d(\nu/2).
\end{align*}
In the last line we used three facts: $2\nu_0-\nu_2\ge\nu_0$ and $\nu-\nu_2\ge\nu/2$, both because $\nu_2\le\min(\nu_0,\nu/2)$; and $c\mapsto S_d(c)$ is non-increasing.

Inserting both estimates into \eqref{8.3:eq-decoupling-localisation-5}, and using $\nu_2\le c_1/2$ to replace $e^{-(c_1/2)\ell_S}$ by $e^{-\nu_2\ell_S}$, we find
\begin{equation*}
\bigl|\operatorname{Cov}_{\mu_\sigma}(G,\mathfrak{W})\bigr|
\le\|G-c\|_\infty|S|\,2\bar r\omega_0\Bigl[1+\frac{2}{m_{\mathfrak{S}}}(q_++p_2)S_d(\nu_0)S_d(\nu/2)\Bigr]e^{-\nu_2\ell_S}.
\end{equation*}
Two elementary remarks finish the argument.
\begin{itemize}
\item With $a:=\bar{\mathfrak{K}}/m_{\mathfrak{S}}$ and $b:=2S_d(\nu/2)$ we have $1+ab\le\max\{1,b\}(1+a)$. Hence
\[2\Bigl[1+\tfrac{2}{m_{\mathfrak{S}}}\bar{\mathfrak{K}}S_d(\nu/2)\Bigr]\le C_{\mathrm{lat}}\bigl(1+\bar{\mathfrak{K}}/m_{\mathfrak{S}}\bigr).\]
\item By definition of $\nu$ we have $\nu\le\nu_0/2\le\nu_0$, so $\nu/2\le\nu_0$ and $\nu_2=\min(c_1/2,\nu/2)$. Therefore
\[\nu_2\ge\tfrac12\min(c_1/2,\nu)=\nu_1,\qquad e^{-\nu_2\ell_S}\le e^{-\nu_1\ell_S}.\]
\end{itemize}

This is the second inequality of \eqref{8.3:eq-decoupling-localisation-1}. It holds for every $\sigma\in[0,1]$ and every $z\in\mathcal{K}_F$, since the class constants of $\mu_\sigma$ do not depend on $\sigma$ or $z$. This proves (i).

\emph{Part (ii).} Apply \eqref{8.3:eq-decoupling-localisation-3} to $\Psi=G_1G_2$, $\Psi=G_1$ and $\Psi=G_2$, which are bounded measurable functions of $y_S$. This shows that $\sigma\mapsto\operatorname{Cov}_{\mu_\sigma}(G_1,G_2)$ is $C^1$, being a polynomial in the three expectations $\langle G_1G_2\rangle_\sigma$, $\langle G_1\rangle_\sigma$ and $\langle G_2\rangle_\sigma$, with
\begin{align*}
\frac{d}{d\sigma}\operatorname{Cov}_{\mu_\sigma}(G_1,G_2)
&=\operatorname{Cov}_{\mu_\sigma}(G_1G_2,\mathfrak{W})-\langle G_2\rangle_\sigma\operatorname{Cov}_{\mu_\sigma}(G_1,\mathfrak{W})
-\langle G_1\rangle_\sigma\operatorname{Cov}_{\mu_\sigma}(G_2,\mathfrak{W})\\
&=\bigl\langle(G_1-\langle G_1\rangle_\sigma)(G_2-\langle G_2\rangle_\sigma)
(\mathfrak{W}-\langle\mathfrak{W}\rangle_\sigma)\bigr\rangle_\sigma,
\end{align*}
the third joint cumulant of $G_1$, $G_2$ and $\mathfrak{W}$ under $\mu_\sigma$. The second equality follows by expanding both sides into moments of $\mu_\sigma$. Put
\[H_\sigma:=(G_1-\langle G_1\rangle_\sigma)(G_2-\langle G_2\rangle_\sigma).\]
Since $\langle\mathfrak{W}-\langle\mathfrak{W}\rangle_\sigma\rangle_\sigma=0$, subtracting the constant $\langle H_\sigma\rangle_\sigma$ from $H_\sigma$ does not change $\langle H_\sigma(\mathfrak{W}-\langle\mathfrak{W}\rangle_\sigma)\rangle_\sigma$; hence the derivative above equals $\operatorname{Cov}_{\mu_\sigma}(H_\sigma,\mathfrak{W})$.

For fixed $\sigma$, $H_\sigma$ is a bounded measurable function of $y_S$. Since $|\langle G_1\rangle_\sigma-c'|\le\|G_1-c'\|_\infty$, we have $\|G_1-\langle G_1\rangle_\sigma\|_\infty\le2\|G_1-c'\|_\infty$, and likewise for $G_2$ with $c''$. Hence
\[\|H_\sigma\|_\infty\le\|G_1-\langle G_1\rangle_\sigma\|_\infty\,\|G_2-\langle G_2\rangle_\sigma\|_\infty\le2\|G_1-c'\|_\infty\cdot2\|G_2-c''\|_\infty.\]
The covariance bound of (i), applied to $H_\sigma$ with the constant $c=0$, gives
\begin{equation*}
\Bigl|\frac{d}{d\sigma}\operatorname{Cov}_{\mu_\sigma}(G_1,G_2)\Bigr|\le4\|G_1-c'\|_\infty\|G_2-c''\|_\infty\,
C_{\mathrm{lat}}|S|\bar r\omega_0\bigl(1+\bar{\mathfrak{K}}/m_{\mathfrak{S}}\bigr)e^{-\nu_1\ell_S}.
\end{equation*}
We integrate over $\sigma\in[0,1]$. At $\sigma=1$ the covariance is $\operatorname{Cov}_\mu(G_1,G_2)$. At $\sigma=0$ the measure $\mu_0=\mu^{\mathrm{bx}}\otimes\mu^{\mathrm{out}}$ factorises and $G_1,G_2$ depend on $y_S$ with $S\subset I$, so the covariance is $\operatorname{Cov}_{\mu^{\mathrm{bx}}}(G_1,G_2)$. This gives \eqref{8.3:eq-decoupling-localisation-2} and proves (ii).

\emph{Part (iii).} For $\tau\in[0,1]$, let $\mu^\tau$ be $\mu^{\mathrm{bx}}$ with $Q_{II}$ replaced by
\[Q^\tau_{II}:=(1-\tau)A_{II}+\tau(\Gamma_{II})^{-1}=A_{II}-\tau S_c,\]
where $S_c=A_{I\mathcal{O}}(A_{\mathcal{O}\mathcal{O}})^{-1}A_{\mathcal{O}I}$. The last equality is part (b) of Lemma~\ref{8.3:lem-principal-blocks}.

By parts (a) and (c) of Lemma~\ref{8.3:lem-principal-blocks}, $\mu^\tau\in\mathfrak{S}$ on the sites $I$, with:
\begin{itemize}
\item $q_-=a_-$, since $Q^\tau_{II}\ge a_-$ by part (a), both $A_{II}$ and $(\Gamma_{II})^{-1}$ being bounded below by $a_-$;
\item $q_+\mapsto a_++C_{sc}$ and $c_Q\mapsto\min(c_A,\mu_A/2)$;
\item the same perturbation $P(\cdot,0_O,0_F)$, the same window and the same smallness condition.
\end{itemize}

Moreover $\mu^0=\mu^{\mathrm{bx}}$ and $\mu^1=\mu^{\mathrm{mg}}$. The energy of $\mu^\tau$ is affine in $\tau$ with
\[\partial_\tau\Phi^\tau=-\tfrac12\langle y_I,S_cy_I\rangle=:-\mathfrak{V}.\]
Hence the computation \eqref{8.3:eq-decoupling-localisation-3}, with $\mathfrak{V}$ in place of $\mathfrak{W}$, gives
\[\frac{d}{d\tau}\langle G\rangle_{\mu^\tau}=\operatorname{Cov}_{\mu^\tau}(G,\mathfrak{V}).\]

Since $S_c$ is symmetric, $\partial_x\mathfrak{V}=(S_cy_I)_x$. By the entry bound of part (b) of Lemma~\ref{8.3:lem-principal-blocks}, dropping the factor $e^{-(\mu_A/2)d_{\mathcal{O}}(x')}\le1$ and using $|y_{x'}|\le\bar r$ on $\mathcal{K}^+$, we get
\begin{align*}
\|\partial_x\mathfrak{V}\|=\|(S_cy_I)_x\|
&\le\sum_{x'}C_{sc}e^{-(\mu_A/2)d_{\mathcal{O}}(x)}e^{-(\mu_A/2)\operatorname{dist}(x,x')}\,\bar r\\
&\le\bar rC_{sc}S_d(\mu_A/2)\,e^{-(\mu_A/2)d_{\mathcal{O}}(x)}.
\end{align*}
This is a boundary profile of the shape in Lemma~\ref{8.3:lem-boundary-profile}, with
$(\omega_0,c_1)$ replaced by $(\bar rC_{sc}S_d(\mu_A/2),\mu_A)$. For $x\in S$ we have
$d_{\mathcal{O}}(x)\ge\ell_S=\operatorname{dist}(S,\mathfrak{Y}\setminus I)$, by the definition of
$\ell_S$ in (iii). In the application with $O=\emptyset$ made later in this chapter,
$\mathcal{O}$ is the complement of the box $I$ and $d_{\mathcal{O}}=d_F$ on $I$, so that
$\ell_S=\operatorname{dist}(S,F)$.

The proofs of (i) and (ii) now repeat, step for step, with $\mathfrak{V}$ in place of $\mathfrak{W}$, with $\tau$ in place of $\sigma$, with $d_{\mathcal{O}}$ in place of $d_O$, and with the class constants of $\mu^\tau$. This gives (i) and (ii) for the pair $(\mu^{\mathrm{bx}},\mu^{\mathrm{mg}})$. The statements for the pair $(\mu,\mu^{\mathrm{mg}})$ follow from those for $(\mu,\mu^{\mathrm{bx}})$ and $(\mu^{\mathrm{bx}},\mu^{\mathrm{mg}})$ by the triangle inequality.
\end{proof}

\begin{remark}[The case of an empty outside]
In the application of Theorem~\ref{8.3:thm-decoupling-localisation} made later in this chapter, $O:=\emptyset$: the decoupling is between the box $\mathfrak{X}=I$ and the frozen exterior $F$ only. Then $Q^\sigma=Q$, $\ell^\sigma=\ell^0+\sigma Rz$ and $P^\sigma(y,z)=(1-\sigma)P(y,0_F)+\sigma P(y,z)$. Moreover
\begin{gather*}
\mu^{\mathrm{bx}}(dy)=Z_{\mathrm{bx}}^{-1}\mathbf{1}_{\mathcal{K}}(y)
\exp\bigl[-\tfrac12\langle y,Qy\rangle-\langle\ell^0,y\rangle+P(y,0_F)\bigr]dy,\\
\mathfrak{W}=-\langle Rz,y\rangle+\bigl[P(y,z)-P(y,0_F)\bigr].
\end{gather*}
In this case:
\begin{itemize}
\item the separation condition is void;
\item throughout Definition~\ref{8.3:def-decoupled-family}, Lemma~\ref{8.3:lem-boundary-profile} and the proof of (i) and (ii), $d_O:=\operatorname{dist}(\cdot,O)$ is replaced by $d_F:=\operatorname{dist}(\cdot,F)$, and $\ell_S:=\operatorname{dist}(S,F)$;
\item the terms indexed by $o\in O$ are absent, and every step of the proof goes through unchanged.
\end{itemize}

Part (iii) is unchanged, with $\mathcal{O}:=\mathfrak{Y}\setminus I$.
\end{remark}

\begin{corollary}[Gradient form of box localisation]
Under the hypotheses of Theorem~\ref{8.3:thm-decoupling-localisation}, there is a lattice-sum constant $C'_{\mathrm{lat}}=C'_{\mathrm{lat}}(d,\nu)$ with the following property. For every $G\in C^1(\mathcal{K}_I^+)$ depending on $y_S$,
\begin{equation}\label{8.3:eq-decoupling-localisation-7}
\bigl|\langle G\rangle_\mu-\langle G\rangle_{\mu^{\mathrm{bx}}}\bigr|
\le\frac{2}{m_{\mathfrak{S}}}\,C'_{\mathrm{lat}}\,\omega_0\sum_{x\in S}\|\partial_xG\|_{\infty,\mathcal{K}_S^+}\,e^{-\nu_1\ell_S}.
\end{equation}
Suppose moreover that $G$ is holomorphic on the complex $\delta r$-neighbourhood of the window, and
let $\operatorname{osc}^{\mathbb{C}}(G)$ be its oscillation there. Then
\[\sum_{x\in S}\|\partial_xG\|_\infty\le|S|\cdot2n\cdot\operatorname{osc}^{\mathbb{C}}(G)/(\delta r_-).\]
Consequently $\bigl|\langle G\rangle_\mu-\langle G\rangle_{\mu^{\mathrm{bx}}}\bigr|$ is at most $C\,|S|\,\operatorname{osc}^{\mathbb{C}}(G)\,(\omega_0/r_-)\,e^{-\nu_1\ell_S}$, with $C:=8nC'_{\mathrm{lat}}/(m_{\mathfrak{S}}\delta)$, and
\[\frac{\omega_0}{r_-}=4(1+\delta)\frac{r_+}{r_-}(q_++p_2)S_d(c_1/2)\]
contains the radii only through the ratio $r_+/r_-$.
\end{corollary}

\begin{proof}
In the proof of part (i) of Theorem~\ref{8.3:thm-decoupling-localisation}, use part (a) of Lemma~\ref{8.3:lem-correlation-decay} in place of part (b). Since $\partial_xG=0$ for $x\notin S$, this gives, for every $\sigma$,
\begin{equation*}
\bigl|\operatorname{Cov}_{\mu_\sigma}(G,\mathfrak{W})\bigr|\le\frac{2}{m_{\mathfrak{S}}}\sum_{x\in S}\sum_{x''\in\mathfrak{X}}
e^{-\nu\operatorname{dist}(x,x'')}\,\|\partial_xG\|_\infty\,\omega(x'').
\end{equation*}

We perform the remaining lattice convolution. Let $x\in S$ and $x''\in\mathfrak{X}$.
\begin{itemize}
\item If $x''\in I$, then $d_O(x'')\ge\ell_S-\operatorname{dist}(x,x'')$.
\item If $x''\in O$, then $\operatorname{dist}(x,x'')\ge\ell_S$.
\end{itemize}

As in the proof of \eqref{8.3:eq-decoupling-localisation-6}, with $x'''=x$, both cases give $\omega(x'')\le\omega_0e^{-\nu_2\ell_S}e^{\nu_2\operatorname{dist}(x,x'')}$. Therefore
\[\sum_{x''}e^{-\nu\operatorname{dist}(x,x'')}\omega(x'')\le\omega_0e^{-\nu_2\ell_S}S_d(\nu-\nu_2)\le S_d(\nu/2)\,\omega_0e^{-\nu_1\ell_S},\]
so one may take $C'_{\mathrm{lat}}=S_d(\nu/2)$. Integrating over $\sigma$ by \eqref{8.3:eq-decoupling-localisation-4} gives \eqref{8.3:eq-decoupling-localisation-7}.

For the holomorphic case, fix $x\in S$ and a point $y$ of the real $\delta r/3$-interior. The
oscillation of $G$ on the complex $\delta r$-neighbourhood bounds the modulus of $G-G(y)$ there.
The Cauchy estimate, applied to $G-G(y)$ in each of the $n$ complex coordinates of the component at
$x$, gives on the real $\delta r/3$-interior
\[\|\partial_xG\|_\infty\le2n\operatorname{osc}^{\mathbb{C}}(G)/(\delta r_x).\]
Since $r_x\ge r_-$, summing over $x\in S$ gives the stated bound on
$\sum_{x\in S}\|\partial_xG\|_\infty$.

The function $G$ is complex-valued, while \eqref{8.3:eq-decoupling-localisation-7} concerns
real-valued functions. At real points $\partial_x\operatorname{Re}G=\operatorname{Re}\partial_xG$ and
$\partial_x\operatorname{Im}G=\operatorname{Im}\partial_xG$, so
$\|\partial_x\operatorname{Re}G\|_\infty$ and $\|\partial_x\operatorname{Im}G\|_\infty$ are at most
$\|\partial_xG\|_\infty$. Applying \eqref{8.3:eq-decoupling-localisation-7} to the real and
imaginary parts of $G$ separately and adding costs a factor $2$. Inserting the bound on
$\sum_{x\in S}\|\partial_xG\|_\infty$ therefore gives the factor
$2\cdot\tfrac{2}{m_{\mathfrak{S}}} C'_{\mathrm{lat}}\cdot2n/\delta=C$ in front of $|S|\operatorname{osc}^{\mathbb{C}}(G)(\omega_0/r_-)e^{-\nu_1\ell_S}$. Finally,
\[\frac{\omega_0}{r_-}=\frac{4\bar r(q_++p_2)S_d(c_1/2)}{r_-}=4(1+\delta)\frac{r_+}{r_-}(q_++p_2)S_d(c_1/2),\]
by $\bar r=(1+\delta)r_+$.
\end{proof}

\begin{proposition}[The step on a clean history lies in the class]\label{8.3:prop-class-membership}
Let $k_0$, $r$, $D$, $\mathfrak{l}$, the sets
$\mathfrak{B}^{(i)}_D\subset\widetilde{\mathfrak{B}}^{(i)}_D$ and the clean
event $\mathcal{C}_S(D)$ be as in Definition~\ref{8.3:def-clean-event}. For $0\le i<r$ put $j:=k_0+i$,
so that $k_0\le j<k_0+r$, and let $I$ be the set of free bonds of $\mathfrak{B}^{(i)}_D$. Let $\mathbb{B}^+$ be
as in Lemma~\ref{8.3:lem-localisation-hypotheses}, here the set of all free bonds of $\Omega_{j+1}$, where
$\Omega_{j+1}\supset\Lambda_{j+1}$ are the domains of the step at level $j$ in \cite[\S3]{B-CMP119} (in the
notation there), and let
$F:=\mathbb{B}^+\setminus I$, the frozen exterior, consisting of all remaining free bonds of $\Omega_{j+1}$.
We write $y=(y_b)$, $y_b\in\mathfrak{g}$, for the fluctuation variables of the step at level $j$, $y_I$ and
$y_F$ for their restrictions to $I$ and to $F$, $B^a(b)$, $a=1,\dots,N^2-1$, for the coordinates of $y_b$ in a
fixed basis of $\mathfrak{g}$ orthonormal for the Euclidean norm $|\cdot|_2$, and $P^{\mathrm{tot}}_j$ for the
exponent of the step-$j$ integrand other than its Gaussian part, as in
Lemma~\ref{8.3:lem-localisation-hypotheses}. The distance $\operatorname{dist}(b,b')$ between bonds is the
level-$j$ lattice distance, and for $c>0$ we write $S_d(c)$ for the Schur constant of the class,
\[
S_d(c)=\sup_b\sum_{b'}e^{-c\operatorname{dist}(b,b')},
\]
the sum running over the bonds of the level-$j$ lattice in dimension $d=4$; thus $S_d(c)\ge1$. Write
$\mathsf{K}_j:=\{\xi\in\mathfrak{g}:|\xi|<\mathfrak{b}_j\}$ for the norm ball of
\eqref{8.3:eq-window-product-a} at level $j$, $\widehat{\mathfrak{b}}_j:=\sup_{\xi\in\mathsf{K}_j}|\xi|_2\le
c_N\mathfrak{b}_j$ for its Euclidean out-radius, with $c_N$ the norm-equivalence constant of
Lemma~\ref{8.3:lem-localisation-hypotheses}, depending on $N$ only, $\mathcal{W}_I:=\prod_{b\in I}\mathsf{K}_j$,
and $\mathcal{W}_F$ for the product of the window balls over the free bonds of $\Lambda_{j+1}\setminus I$ and of
the far-field balls of radius $(1\vee C_O)\widehat{\mathfrak{b}}_j$ over the free bonds of
$\Omega_{j+1}\setminus\Lambda_{j+1}$, where $C_O\ge1$ is the exterior-field constant of
Lemma~\ref{8.3:lem-localisation-hypotheses}. Let $\mathbf{C}$ be the linear operator of \cite{B-CMP109}
acting on the fluctuation variables, $\|\mathbf{C}\|$ its norm and $R_{\mathbf{C}}$ its range measured in
level-$(j+1)$ blocks, and put
\begin{equation}\label{8.3:eq-class-membership-7}
\mathfrak{h}_{\mathrm{pk}}:=C_\delta\|\mathbf{C}\|^2\mathfrak{b}_j^2\,e^{-\frac12c'p_0(g_{k_0})},\qquad
\mathfrak{h}_{\mathrm{R}}:=C'g_jx_j^{(3-\kappa_0)/2}\operatorname{polylog}(x_j)/\widehat{\mathfrak{b}}_j,
\end{equation}
where $C_\delta$ is a constant depending on the margin $\delta_w$ of hypothesis (b) below, $c'$ is the constant
of \eqref{8.3:eq-clean-event-a}, $p_0(\cdot)$ is the degree function, $\kappa_0\ge4$ is the exponent of
\cite[(2.31)]{B-CMP119}, $C'$ is a constant and $\operatorname{polylog}(x_j)$
denotes a polynomial in $\log x_j$. Assume the following.
\begin{itemize}
\item[(a)] The operator theorem at free current, in the following form: the background $\mathfrak{u}$ of every
admissible exterior datum lies in the real domain of that theorem, and the quadratic form $A^{\Gamma}_j$ of
the step's fluctuation integral, with blocks $A_{II}:=A^{\Gamma}_j\restriction_{I\times I}$ and
$A_{IF}:=A^{\Gamma}_j\restriction_{I\times F}$, is
real symmetric, satisfies
$\|A^{\Gamma}_j(b,b')\|\le a_+e^{-c_A\operatorname{dist}(b,b')}$ and is bounded below by $a_->0$; the constants
$a_-$, $a_+$, $c_A$ are level-free.
\item[(b)] The locality and analyticity bound on the tube: there is a complex tube $\mathcal{T}_j$ about the
window containing $\mathcal{W}_I^+\times\mathcal{W}_F$, far-field radii included, with real margin $\delta_w$
over the window, $\delta_w\widehat{\mathfrak{b}}_j\le1$, where $\mathcal{W}_I^+$ denotes the set of real $y_I$
each of whose components lies within Euclidean distance $\delta_w\widehat{\mathfrak{b}}_j$ of $\mathsf{K}_j$;
and on $\mathcal{T}_j$, for all bonds $b,b'$ (pairs with one bond
in $I$ and one in $F$ included) and $a,a'=1,\dots,N^2-1$,
\begin{equation}\label{8.3:eq-class-membership-5}
\Bigl|\frac{\partial P^{\mathrm{tot}}_j}{\partial B^a(b)}\Bigr|\le C_Pg_j\mathfrak{b}_j^2,\qquad
\Bigl|\frac{\partial^2P^{\mathrm{tot}}_j}{\partial B^a(b)\,\partial B^{a'}(b')}\Bigr|
\le C_Pg_j\mathfrak{b}_je^{-c_P\operatorname{dist}(b,b')},
\end{equation}
where $C_P,c_P>0$ are level-free constants, independent of $j$, $K$, $m$, $s$, the position and the
volume, as in Lemma~\ref{8.3:lem-localisation-hypotheses}.
\item[(c)] The window of Definition~\ref{8.3:def-window-product}: the fluctuation window is the product over the
free bonds of the open, bounded, convex, symmetric, $\mathrm{Ad}$-invariant balls $\mathsf{K}_j$.
\item[(d)] The two hypotheses of Lemma~\ref{8.3:lem-box-law-mixture}: the step of the history near $I$ is a term
of the integral \cite[(3.25)]{B-CMP119}, as in Lemma~\ref{8.3:lem-box-law-mixture}, with $I$ contained in its
set of integration bonds $\Lambda^{(j)*}_{j+1}$, and every weight of the large-field operations is
non-negative.
\item[(e)] The collar condition: let $\mathfrak{r}_{\mathrm{dec}}$ be the diameter, in level-$j$ units, of the
neighbourhood $\widetilde{\square}$ of \cite[(2.16)]{B-CMP119} of one of the level-$j$ decision cubes $\square$
of \cite{B-CMP119} (the cubes of the size $LM_2R_j$ in the notation there; $\widetilde{\square}$ adds one
layer), plus $R_{\mathbf{C}}+2$ blocks, and let $I^+$ be the set of bonds within distance
$\mathfrak{r}_{\mathrm{dec}}$ of $I$; then $I^+\subset\widetilde{\mathfrak{B}}^{(i)}_D$, and every level-$j$
decision cube $\square$ whose neighbourhood $\widetilde{\square}$ meets $I^+$ meets
$\widetilde{\mathfrak{B}}^{(i)}_D$.
\item[(f)] The finiteness hypothesis on the clean collar $\widetilde{\mathfrak{B}}^{(i)}_D$ that accompanies
Definition~\ref{8.3:def-clean-event}.
\item[(g)] The per-bond clause \eqref{8.3:eq-clean-event-a}, with $x_{k_0}\ge x_{\mathrm{pk}}(L,\mathfrak{p},N,M)$
as in \eqref{8.3:eq-depth-coupling-conditions-a}.
\item[(h)] The bound \cite[(2.31)]{B-CMP119} for the $\mathbf{R}$-terms and the far-field bound
\cite[(3.13)--(3.14)]{B-CMP119}.
\item[(i)] The comparison of the running couplings along the segment, and the analyticity radii of the
level-$i'$ terms summed in \eqref{8.3:eq-clean-event-a}: for $k_0\le j<k_0+r$ one has
$g_j\le\sqrt2\,g_{k_0}$, $\mathfrak{b}_j\le2\mathfrak{b}_{k_0}$, $x_j\ge x_{k_0}/2$, and for
$i'<k_0-\mathfrak{l}$, $\nu:=j-i'$, the ratio bound
$(g_j/R^{\mathrm{an}}_{i'})^2\le(c_{\mathrm{rad}}C_0)^{-2}(1+B^\sharp\nu)$, where $R^{\mathrm{an}}_{i'}$
denotes the analyticity radius of the level-$i'$ localised terms and $c_{\mathrm{rad}}$ is the constant in the
lower bound $R^{\mathrm{an}}_{i'}\ge c_{\mathrm{rad}}C_0g_{i'}$.
\item[(j)] The ceiling: $g_{k_0}\le\gamma_{\mathrm{GX}}$, where $\gamma_{\mathrm{GX}}$ is the ceiling of the
Gaussian extraction and the defining list of one-sided conditions of
$\gamma_{\mathrm{GX}}$ is understood to contain the floor $x_{k_0}\ge x_{\mathrm{pk}}$ and, for
$k_0\le j<k_0+r$, the condition
\begin{equation}\label{8.3:eq-class-membership-1}
g_j\mathfrak{b}_j\,C_PS_d(c_P)+(\mathfrak{h}_{\mathrm{pk}}+\mathfrak{h}_{\mathrm{R}})S_d(c_P)\le\tfrac14a_-,
\end{equation}
with $\mathfrak{h}_{\mathrm{pk}}$ and $\mathfrak{h}_{\mathrm{R}}$ as in \eqref{8.3:eq-class-membership-7}.
\end{itemize}
Then for every history $h\in\mathcal{C}_S(D)$, every $j$ with $k_0\le j<k_0+r$, every block field $V_{j+1}$ in
its window and every $y_F\in\mathcal{W}_F$, the conditional law of $y_I$ given $h$, $V_{j+1}$, $y_F$ and the
background $\mathfrak{u}$ (that is, given also the members of the exterior datum that determine the
background), which has density proportional to
\begin{equation}\label{8.3:eq-class-membership-6}
\mathbf{1}_{\mathcal{W}_I}(y_I)\exp\Bigl[-\tfrac12\langle y_I,A_{II}y_I\rangle-\langle y_I,A_{IF}y_F\rangle
+P^{\mathrm{tot}}_j(y_I,y_F)\Bigr],
\end{equation}
is a measure of the class $\mathfrak{S}$ of box laws on the sites $I$ with the level-free constants
\begin{equation}\label{8.3:eq-class-membership-2}
\begin{gathered}
m_{\mathfrak{S}}=\tfrac12a_-,\qquad\nu_0=\tfrac12\min(c_A,c_P),\\
\mathfrak{K}\le(a_++a_-)S_d(\nu_0),\qquad\bar{\mathfrak{K}}\le(a_++a_-)S_d(\nu_0),\qquad
\nu_{\mathfrak{S}}=\nu_{\mathfrak{S}}(a_-,a_+,c_A,c_P,d),
\end{gathered}
\end{equation}
uniformly in $h$, $V_{j+1}$, $y_F$, the position of $I$, the volume, $K$, $m$ and $s$. The kernel
$\kappa_{\mathfrak{S}}$ of the class is the one formed with the rate $\nu_0$; the subscript $\mathfrak{S}$ on
$m_{\mathfrak{S}}$ and $\nu_{\mathfrak{S}}$ separates these constants of the class from the torus exponent $m$
and from $\nu=j-i'$. The data of the
class are: the sites $I$, with fibre $\mathfrak{g}$ and the level-$j$ lattice distance; the quadratic form
$A_{II}$; the frozen variable $y_F$ with linear term $A_{IF}y_F$; the windows $\mathsf{K}_j$ at the sites of
$I$, of radius $\widehat{\mathfrak{b}}_j$, and the exterior radii at most $(1\vee C_O)\widehat{\mathfrak{b}}_j$;
the margin $\delta_w$; and the exponent $P^{\mathrm{tot}}_j(\cdot,y_F)$ with gradient bound
$p_1=C_Pg_j\mathfrak{b}_j^2$, Hessian bound
$p_2=C_Pg_j\mathfrak{b}_j+\mathfrak{h}_{\mathrm{pk}}+\mathfrak{h}_{\mathrm{R}}$ and decay rate $c_P$; the
separation condition of the class is void and its smallness condition is \eqref{8.3:eq-class-membership-1}.
\end{proposition}

\begin{proof}
The proof has four steps: the identification of the data of the class together with the ceiling
\eqref{8.3:eq-class-membership-1}; the geometry of the near set; the bounds on the attached member and the
$\mathbf{R}$-term member of $P^{\mathrm{tot}}_j$ that enter $p_2$ beside \eqref{8.3:eq-class-membership-5}; and
the conclusion.

\emph{Step 1: the data of the class, the convexity margin and the ceiling.} The sites of the class are the set
$I$ of free bonds of $\mathfrak{B}^{(i)}_D$; by hypothesis (d), $I\subset\Lambda^{(j)*}_{j+1}$. The fibre at
each site is $\mathfrak{g}$, of dimension $N^2-1$, and the distance is the level-$j$ lattice distance. The
quadratic form of the class is $A_{II}$, with lower bound $a_-$, upper bound $a_+$ and decay rate $c_A$, by
hypothesis (a). The frozen exterior is
the set $F$ of the statement, consisting of all remaining free bonds of $\Omega_{j+1}$: those of
$\Lambda_{j+1}\setminus I$ carry their window balls, and those of $\Omega_{j+1}\setminus\Lambda_{j+1}$ carry the
far field of \cite[(3.13)--(3.14)]{B-CMP119}, of size $O(1)\varepsilon_j$ (in the notation there, with $k=j$);
in the scaled variable this gives $|y_f|_2\le C_O\widehat{\mathfrak{b}}_j$ on $\Omega_{j+1}\setminus\Lambda_{j+1}$,
with $C_O\ge1$ the exterior-field constant, and $|y_f|_2<\widehat{\mathfrak{b}}_j$ on
$\Lambda_{j+1}\setminus I$, so that $|y_f|_2\le(1\vee C_O)\widehat{\mathfrak{b}}_j$ for every $f\in F$. The
frozen variable of the
class is $\mathsf{z}:=y_F$, its linear term is $\mathsf{l}^0:=0$ with coupling operator $\mathsf{R}:=A_{IF}$, so
that the linear term of the class is $\mathsf{l}=A_{IF}y_F$. The windows are $\mathsf{K}_x:=\mathsf{K}_j$ for
$x\in I$, the norm ball of Definition~\ref{8.3:def-window-product}; its Euclidean out-radius is
$\widehat{\mathfrak{b}}_j\le c_N\mathfrak{b}_j$. The radii of the class are $\mathsf{r}_x:=\widehat{\mathfrak{b}}_j$
for $x\in I$ and $\mathsf{r}_f\le(1\vee C_O)\widehat{\mathfrak{b}}_j$ for $f\in F$, so that
$\mathsf{r}_-=\widehat{\mathfrak{b}}_j$ and $\mathsf{r}_+:=(1\vee C_O)\widehat{\mathfrak{b}}_j$. The margin of the
class is $\delta_w$, the real margin of the tube $\mathcal{T}_j$ of hypothesis (b), with
$\delta_w\widehat{\mathfrak{b}}_j\le1$; by hypothesis (b) the tube contains
$\mathcal{W}_I^+\times\mathcal{W}_F$ with the far-field radii $\mathsf{r}_f$ included. The exponent of the class
is $P(y_I,\mathsf{z}):=P^{\mathrm{tot}}_j(y_I,y_F)$ for $y_I\in\mathcal{W}_I^+$, with gradient bound
$p_1:=C_Pg_j\mathfrak{b}_j^2$, Hessian bound
$p_2:=C_Pg_j\mathfrak{b}_j+\mathfrak{h}_{\mathrm{pk}}+\mathfrak{h}_{\mathrm{R}}$ and decay rate $c_P$: the terms
$C_Pg_j\mathfrak{b}_j^2$, $C_Pg_j\mathfrak{b}_j$ and the rate $c_P$ come from \eqref{8.3:eq-class-membership-5}
on $\mathcal{W}_I^+\subset\mathcal{T}_j$, the mixed Hessians between $I$ and $F$ being included because
\eqref{8.3:eq-class-membership-5} holds for all pairs of bonds, and $\mathfrak{h}_{\mathrm{pk}}$,
$\mathfrak{h}_{\mathrm{R}}$ bound the two members treated in Step~3. The separation condition of
Definition~\ref{8.3:def-decoupled-family} is void, since the outside set $O$ is empty here. The smallness
condition of the class reads
\[
(C_Pg_j\mathfrak{b}_j+\mathfrak{h}_{\mathrm{pk}}+\mathfrak{h}_{\mathrm{R}})S_d(c_P)\le\tfrac14a_-,
\]
which is \eqref{8.3:eq-class-membership-1}.

We show that \eqref{8.3:eq-class-membership-1} is a one-sided ceiling on the running coupling. By the definition
of the window radius $\mathfrak{b}_j$,
\[
g_j\mathfrak{b}_j=\delta_j=(A_1/A_0)\varepsilon_j=A_1g_j(\log x_j)^{p_0},
\]
where $\delta_j$ and $\varepsilon_j$ are Ba{\l}aban's small-field thresholds in \cite{B-CMP119}; this tends
to $0$ as $x_j\to\infty$; and by Step~3 below,
$\mathfrak{h}_{\mathrm{pk}}+\mathfrak{h}_{\mathrm{R}}\le2C_Pg_j\mathfrak{b}_j$ on $\mathcal{C}_S(D)$ once
$x_{k_0}\ge x_{\mathrm{pk}}$. Hence \eqref{8.3:eq-class-membership-1} is a one-sided ceiling on $g_j$,
equivalently a floor $x_j\ge x_H(L,\mathfrak{p},N)$. By hypothesis (i), $g_j\le\sqrt2\,g_{k_0}$ and
$\mathfrak{b}_j\le2\mathfrak{b}_{k_0}$ on the whole segment; therefore there is $\gamma'>0$, depending on $L$,
$\mathfrak{p}$, $N$ and $M$ only, such that
\eqref{8.3:eq-class-membership-1} holds for all $j$ with $k_0\le j<k_0+r$ whenever $g_{k_0}\le\gamma'$. This
condition is of the same kind as the one-sided conditions defining the ceiling $\gamma_{\mathrm{GX}}$ fixed with
the Gaussian extraction, and by hypothesis (j) it is adjoined to their list. Above the ceiling the displays of
the Gaussian extraction hold trivially with enlarged constants; the present proposition is stated under the
ceiling. No further condition is imposed: no lower bound on any coupling enters, and no upper bound on $L$, $K$,
$m$, $s$, $r$ or $D$.

The constants of the class are then as follows. The mass is $m_{\mathfrak{S}}=a_-/2$, the rate is
$\nu_0=\tfrac12\min(c_A,c_P)$, and the kernel $\kappa_{\mathfrak{S}}$ of the class is formed with the rate
$\nu_0$. For $\mathfrak{K}$ we have
$\mathfrak{K}\le(a_++p_2)S_d(\nu_0)$ with
$p_2=C_Pg_j\mathfrak{b}_j+\mathfrak{h}_{\mathrm{pk}}+\mathfrak{h}_{\mathrm{R}}$, and $p_2\le\tfrac14a_-$ by
\eqref{8.3:eq-class-membership-1}, since $p_2S_d(c_P)\le\tfrac14a_-$ and $S_d\ge1$. Thus
\[
\mathfrak{K}\le(a_++C_Pg_j\mathfrak{b}_j+\mathfrak{h}_{\mathrm{pk}}+\mathfrak{h}_{\mathrm{R}})S_d(\nu_0)
\le(a_++a_-)S_d(\nu_0).
\]
For $\bar{\mathfrak{K}}$, the same inequality $p_2\le\tfrac14a_-$ gives
\[
\bar{\mathfrak{K}}=(a_++p_2)S_d(\nu_0)\le(a_++\tfrac14a_-)S_d(\nu_0)\le(a_++a_-)S_d(\nu_0).
\]
Finally $\nu_{\mathfrak{S}}=\nu_{\mathfrak{S}}(a_-,a_+,c_A,c_P,d)$. All of these are level-free: they do not depend on
$K$, $m$, $s$, $r$, $D$, the position, $U_{j+1}$, $h$ or the volume. This is \eqref{8.3:eq-class-membership-2}.

\emph{Step 2: the near set, its dependence collar and the clean collar.} By
Definition~\ref{8.3:def-clean-event}, $\mathfrak{B}^{(i)}_D$ is the set $S^{(i)}$ enlarged by $D$ level-$(j+1)$
blocks on every side, and $\widetilde{\mathfrak{B}}^{(i)}_D$ is $\mathfrak{B}^{(i)}_D$ enlarged by
$100M\check{R}_{k_0+i}$ level-$(k_0+i)$ cubes. The clauses (c1) and (c3) of the clean event forbid every
large-field creation of level in $[k_0-\mathfrak{l},k_0+i]$ at cubes meeting $\widetilde{\mathfrak{B}}^{(i)}_D$,
and on $\mathcal{C}_S(D)$, within every $\mathfrak{B}^{(i)}_D$, the characteristic functions of the steps
$k_0,\dots,k_0+r$ are the small-field ones. Consequently the window product of hypothesis (c) and the quadratic
bound on the fluctuations apply there.

The set $I$ of free bonds of $\mathfrak{B}^{(i)}_D$ is a union of level-$(j+1)$ blocks. Let
$\mathfrak{r}_{\mathrm{dec}}$ and $I^+$ be as in hypothesis (e). Hypothesis (e) requires
$I^+\subset\widetilde{\mathfrak{B}}^{(i)}_D$ and that every decision cube $\square$ whose $\widetilde{\square}$
meets $I^+$ meets $\widetilde{\mathfrak{B}}^{(i)}_D$; in other
words, the collar of $\widetilde{\mathfrak{B}}^{(i)}_D$ over $\mathfrak{B}^{(i)}_D$ is at least
$\mathfrak{r}_{\mathrm{dec}}$ wide.

From hypothesis (e), the clauses (c1) and (c3), the window product and its convexity (hypothesis (c)), and
hypothesis (f), we draw four consequences.
\begin{itemize}
\item[(g1)] Every window factor of the step-$j$ integrand that depends on $y_I$ is a per-site indicator
$\mathbf{1}_{\mathsf{K}_j}(y_b)$, $b\in I$, of the norm ball. Indeed, by \cite[\S3]{B-CMP119}, on the domain
$\Lambda_{j+1}$, and also on the enlarged domain $\widetilde{\Lambda}_{j+1}$ of \cite[\S3]{B-CMP119} (in the
notation there), only the small-field characteristic functions of
\cite[(3.16)]{B-CMP119} occur in the integral, and their product is $\chi^{(j)}(\Lambda_{j+1})$, each cube's factor
being a product over its bonds of balls, by Definition~\ref{8.3:def-window-product}. The characteristic functions
of the type \cite[(3.3)]{B-CMP119} are identically $1$ there, by \cite[(3.17)--(3.19)]{B-CMP119}. In the pure
small-field setting the product is the display \cite[(2.9)]{B-CMP109} verbatim. This is the
convexity of the window, stated for the norm balls of Definition~\ref{8.3:def-window-product}.
\item[(g2)] Every factor of the integrand other than these indicators is either the Gaussian factor, or
$e^{P^{\mathrm{tot}}_j}$, or independent of $y_I$.
\item[(g3)] Hence the placement form of Lemma~\ref{8.3:lem-box-law-mixture} holds on $I$: the joint density is
\[
\Psi(y_F)\,\mathbf{1}_{\mathcal{W}_I}(y_I)\exp\Bigl[-\tfrac12\langle y_I,A_{II}y_I\rangle
-\langle y_I,A_{IF}y_F\rangle+P^{\mathrm{tot}}_j(y_I,y_F)\Bigr],\qquad\Psi\ge0.
\]
This is the law $E_{\mathsf{e}}$ of Lemma~\ref{8.3:lem-localisation-hypotheses}, with an exterior datum
$\mathsf{e}$ containing $y_F$; the inequality $\Psi\ge0$ follows from the positivity hypothesis in (d).
\item[(g4)] For an observable $G$ depending on $V_j$ only within $R_S+M$ level-$j$ units of the centre of
$S^{(i)}$, with $R_S$ as in the theorem on localisation of the sourced step, its
read set $S^+$ (Lemma~\ref{8.3:lem-read-set}) is contained in
$S^{(i)}$ enlarged by $M+R_{\mathbf{C}}+1$ blocks, while the boundary $\partial I$ lies $D$ level-$(j+1)$
blocks, that is $DL$ level-$j$ units, further out. Therefore, in level-$j$ units,
\[
\operatorname{dist}(S^+,F)\ge DL-(M+R_{\mathbf{C}}+2)L=:\varpi_S .
\]
\end{itemize}

\emph{Step 3: the attached member and the $\mathbf{R}$-term member of $P^{\mathrm{tot}}_j$ on the box.} We write
$P^{\mathrm{tot}}_j=P^{\mathrm{sf}}_j+P^{\mathrm{E}}_j+P^{\mathrm{R}}_j+P^{\mathrm{pk}}_j$, where
$P^{\mathrm{sf}}_j$ is the small-field member, $P^{\mathrm{E}}_j$ the member coming from the
$\mathbf{E}$-terms of \cite[\S2]{B-CMP119}, whose analyticity radii are those of \cite[(2.28)]{B-CMP119},
$P^{\mathrm{pk}}_j$ collects the terms of
levels $i'<k_0-\mathfrak{l}$ summed in
\eqref{8.3:eq-clean-event-a} and $P^{\mathrm{R}}_j$ the $\mathbf{R}$-terms.

(3a) On $\mathcal{C}_S(D)$, for every bond $b\in I$ and every $y\in\mathcal{W}_I^+\times\mathcal{W}_F$,
\begin{equation}\label{8.3:eq-class-membership-3}
\sum_{b'}\|\partial_b\partial_{b'}P^{\mathrm{pk}}_j(y)\|\,e^{c_P\operatorname{dist}(b,b')}\le
\mathfrak{h}_{\mathrm{pk}}=C_\delta\|\mathbf{C}\|^2\mathfrak{b}_j^2\,e^{-\frac12c'p_0(g_{k_0})}.
\end{equation}
The bound is obtained term by term, and the terms are then summed. For each term we use its axial-gauge
representative, given by the second equivalence lemma, and the vanishing of its gradient at the flat
configuration, given by part (i) of the first equivalence lemma. By the Schwarz lemma, as in part (c) of the lemma
bounding the attached members, this gives the second-order bound for the term. In that bound we insert three
estimates: the quadratic bound on the fluctuations and its linearised form, which give
$f^{\max}_X\le C''L^{-2\nu}g_j\mathfrak{b}_j$ for the maximal size $f^{\max}_X$ of the fluctuation field over
$\widetilde{X}$; the lower bound $R^{\mathrm{an}}_{i'}\ge c_{\mathrm{rad}}C_0g_{i'}$
on the analyticity radii of the level-$i'$ terms; and the ratio bound
$(g_j/R^{\mathrm{an}}_{i'})^2\le(c_{\mathrm{rad}}C_0)^{-2}(1+B^\sharp\nu)$ of hypothesis (i). The result is that the
term with localisation domain $\widetilde{X}$ has second-order size at most a constant times
$M_X\kappa_{\mathrm{ax}}(X)^2L^{-4\nu}(1+B^\sharp\nu)\mathfrak{b}_j^2$. Cauchy's estimate on the margin of the tube
then bounds its Hessian.

It remains to sum over the terms and over $b'$. The side $d(X)$ of $\widetilde{X}$ is measured in level-$i'$ units,
so the diameter of $\widetilde{X}$ in level-$j$ units is $d(X)L^{-\nu}$. Every bond $b'$ whose fluctuation variable
reaches a level-$i'$ term whose localisation domain lies under the dependence range of $b$ lies within
$d(X)L^{-\nu}+C$ level-$j$
units of $b$, where $C$ is the constant of \eqref{8.3:eq-clean-event-a}: this distance consists of the dependence
range on both sides plus the level-$j$ diameter. Hence, for
one term, the number of such $b'$ is at most $(d(X)L^{-\nu}+C)^4$ and each carries weight at most
$e^{c_P(d(X)L^{-\nu}+C)}$. The sum over $b'$ with the weight $e^{c_P\operatorname{dist}(b,b')}$, summed over the
terms, is therefore bounded by a constant times the left-hand side of \eqref{8.3:eq-clean-event-a}, which on
$\mathcal{C}_S(D)$ is at most $e^{-\frac12c'p_0(g_{k_0})}$. This gives \eqref{8.3:eq-class-membership-3}. No
condition on $x$ beyond $x_{\mathrm{pk}}$ enters.

(3b) On every history, the $\mathbf{R}$-term member has bond-gradient at most
$C'g_jx_j^{(3-\kappa_0)/2}\operatorname{polylog}(x_j)$, where $\operatorname{polylog}(x_j)$ denotes a polynomial in
$\log x_j$. This follows from part (b) of the lemma bounding the attached members and from
\cite[(2.31)]{B-CMP119}: the marginal volume factor $(L^{i'}\eta)^4$ in \cite[(2.31)]{B-CMP119} (in the
notation there) cancels the count of level-$i'$ localisation domains, as stated in \cite[\S2]{B-CMP119}. We
combine this gradient bound with Cauchy's estimate on the margin $\delta_w\widehat{\mathfrak{b}}_j/3$ and with
the locality factor $e^{-\kappa_{\mathrm{loc}}d(X)}$ of the terms, $\kappa_{\mathrm{loc}}$ being their
localisation rate. This gives a weighted Hessian at most
\begin{equation}\label{8.3:eq-class-membership-4}
\mathfrak{h}_{\mathrm{R}}=C'g_jx_j^{(3-\kappa_0)/2}\operatorname{polylog}(x_j)/\widehat{\mathfrak{b}}_j .
\end{equation}

Consequently $\mathfrak{h}_{\mathrm{pk}}+\mathfrak{h}_{\mathrm{R}}\le2C_Pg_j\mathfrak{b}_j$ for
$x_{k_0}\ge x_{\mathrm{pk}}(L,\mathfrak{p},N,M)$. Here $\kappa_0\ge4$; by \cite[\S2]{B-CMP119}, the exponent
$\kappa_0$ of \cite[(2.31)]{B-CMP119} can be chosen
arbitrarily large. By hypothesis (i), $x_j\ge x_{k_0}/2$ on the segment. The first and the last of the conditions in
\eqref{8.3:eq-depth-coupling-conditions-a} are the inequalities $\mathfrak{h}_{\mathrm{pk}}\le C_Pg_j\mathfrak{b}_j$
and $\mathfrak{h}_{\mathrm{R}}\le C_Pg_j\mathfrak{b}_j$. This is a one-sided floor on $x_{k_0}$, and by hypothesis (j)
it is adjoined to the list defining $\gamma_{\mathrm{GX}}$, together with \eqref{8.3:eq-class-membership-1}.

Two further kinds of contribution remain. Let $\mathfrak{d}_{\mathrm{col}}$ denote the width, in
level-$(j+1)$ blocks, of the collar of $\widetilde{\mathfrak{B}}^{(i)}_D$ beyond $I$. The first kind consists of the
large-field objects of levels in
$[k_0-\mathfrak{l},j]$ that are live outside $\widetilde{\mathfrak{B}}^{(i)}_D$, which (c1) and (c3) allow off
$\widetilde{\mathfrak{B}}^{(i)}_D$. The second consists of the localisation domains $X$ that span far enough to reach
$I$. Both reach $y_I$ only through the tails of the localised quadratic bound and through the bound
$M_X\le O(1)e^{-\kappa d(X)}$ with $d(X)\ge\mathfrak{d}_{\mathrm{col}}$, $\kappa$ being Ba{\l}aban's parameter
$\kappa$ of the glossary. By the majorants of the localised terms and Cauchy's estimate on their radii of analyticity, their
contributions to $p_1$ and $p_2$ are polynomial in $p_0(g_{k_0})$ times $e^{-c\,\mathfrak{d}_{\mathrm{col}}L}$,
smaller than
every bound displayed above. In the same way the exponent $P^{\mathrm{tot}}_j(\cdot,0_F)$ of the box depends on $h$
through such far tails, at most a polynomial in $p_0(g_{k_0})$ times $e^{-\kappa\,\mathfrak{d}_{\mathrm{col}}}$ in
supremum over
$\mathcal{W}_I$. This is immaterial for the lemmas of this chapter on the single step and on the Gaussian part,
which use of that exponent only its supremum bound and the Gaussian factor.

\emph{Step 4: conclusion.} Let $h\in\mathcal{C}_S(D)$, let $k_0\le j<k_0+r$, and let $g_{k_0}\le\gamma_{\mathrm{GX}}$ with
\eqref{8.3:eq-class-membership-1} and the floor $x_{k_0}\ge x_{\mathrm{pk}}$ in its defining list. Fix $V_{j+1}$ in its
window and $y_F\in\mathcal{W}_F$. By (g3), the joint density of $(y_I,y_F)$ given $h$, $V_{j+1}$ and the background
$\mathfrak{u}$ is $\Psi(y_F)$ times the expression \eqref{8.3:eq-class-membership-6}. Conditioning further on $y_F$, the
factor $\Psi(y_F)$ is a constant and disappears in the normalisation, so the conditional law of $y_I$ has density
proportional to \eqref{8.3:eq-class-membership-6}.

We now check that its data are those of Step~1. The window is the product of the convex balls $\mathsf{K}_j$ over
$I$, by (g1). The quadratic form $A_{II}$ satisfies the bounds $a_-$, $a_+$, $c_A$ by hypothesis (a), which applies
because the background $\mathfrak{u}$ lies in the real domain of the operator theorem at free current. The linear term
is $A_{IF}y_F$, with $y_F$ in the product of balls of radii at most $(1\vee C_O)\widehat{\mathfrak{b}}_j$. The exponent
$P^{\mathrm{tot}}_j(\cdot,y_F)$ has gradient bound $p_1$, Hessian bound $p_2$ and decay rate $c_P$ on
$\mathcal{W}_I^+$, mixed Hessians included, by \eqref{8.3:eq-class-membership-5} and Step~3. The margin is
$\delta_w$, by hypothesis (b). The separation condition is void, and the smallness condition is
\eqref{8.3:eq-class-membership-1}, which holds under the ceiling. Hence the conditional law is a measure of the class
$\mathfrak{S}$ on the sites $I$ with the constants \eqref{8.3:eq-class-membership-2}. These depend only on $a_-$,
$a_+$, $c_A$, $c_P$ and $d$; the membership therefore holds uniformly in $h$, $V_{j+1}$, $y_F$, the position of
$I$, the volume,
$K$, $m$ and $s$.
\end{proof}

\begin{remark}[Width of the clean collar]
Hypothesis (e) of Proposition~\ref{8.3:prop-class-membership} holds as soon as the collar of
$\widetilde{\mathfrak{B}}^{(i)}_D$ over $\mathfrak{B}^{(i)}_D$ in Definition~\ref{8.3:def-clean-event} is at least
$\mathfrak{r}_{\mathrm{dec}}$ wide. If the clean event of Definition~\ref{8.3:def-clean-event} is defined with the
collar $100M\check{R}+4\mathfrak{r}_{\mathrm{dec}}$ in place of $100M\check{R}$, hypothesis (e) holds in every
case, since that collar is at least $4\mathfrak{r}_{\mathrm{dec}}$ wide. This enlargement of the clean event
changes only the cube counts in the lemma bounding the weight of
the complement of the clean event, by the further polynomial factor
\[
\Bigl(1+\frac{4\mathfrak{r}_{\mathrm{dec}}}{100M\check{R}}\Bigr)^4\le C'L^8,
\]
for a constant $C'$, which is absorbed into the factor $e^{-\frac14c_qp_0}$ of that lemma, $c_q$ being its
constant, in the same way
as the factor $(r+1)^2L^{4r}(M\check{R})^4$ is absorbed there, because $p_0(g_{k_0})=A_0(\log x_{k_0})^{p_0}$
exceeds every fixed multiple of $\log L$ once $x_{k_0}\ge x_{\mathrm{cl}}(L,\mathfrak{p})$, the floor of that lemma.
The weights of the complement, a constant times $D^4(1+r)e^{-cp_0}$, are unchanged in form. Likewise, if the
collar of $\widetilde{\mathfrak{B}}^{(i)}_D$ over $\mathfrak{B}^{(i)}_D$ is taken to be $111\,LM\check{R}_{j+1}$
level-$j$ units wide, the inclusion $I\subset\Lambda_{j+1}$ of hypothesis (d) holds; this enlargement changes the
same weight by a factor at most $(2L)^4$ and nothing else.
\end{remark}

Throughout what follows we use the data and notation of localisation of the sourced step, Theorem~\ref{8.3:thm-sourced-step-localisation}. These are the integers $K$, $s$, $k_0=K-s$, $r\ge1$ and $D\ge1$; the support $S$ with the side $R_S$ of its cube, and $M$; the event $\mathcal{C}_S(D)$ of clean histories; the branch expectation $\mathbb{E}_{i,h}$; the box $\mathfrak{B}^{(i)}_D$, the enlarged box $\tilde{\mathfrak{B}}^{(i)}_D$ and the centre $S^{(i)}$; and the constants of the class $\mathfrak{S}$ fixed in Proposition~\ref{8.3:prop-class-membership}, namely $a_-,a_+,c_A,C_P,c_P,p_2,\nu,\nu_0,C_O,\delta_w$. We write $S_d(c):=\sup_b\sum_{b'}e^{-c\,d(b,b')}$ for the lattice sum.

Further notation is as follows. $n_{\mathfrak{g}}$ is the number of real components of the scaled fluctuation variable at one bond, that is, the fibre dimension in Lemma~\ref{8.3:lem-principal-blocks}. $\mathfrak{b}_j$ is the radius of the per-bond window ball at level $j$ in Ba{\l}aban's norm, and $\hat{\mathfrak{b}}_j\le c_N\mathfrak{b}_j$ is its Euclidean envelope, $c_N$ being the norm-equivalence number. $\mathbf{C}$ is the operator through which the scaled fluctuation field enters the bond variables, and $R_{\mathbf{C}}$ is its range. For a set $J$ of bonds, $\|\Phi\|_{\mathcal{W}_J}$ is the supremum of $|\Phi|$ over the window $\mathcal{W}_J$.

\emph{The conditional law of the fluctuation field.} Fix $h\in\mathcal{C}_S(D)$, $0\le i<r$, $j:=k_0+i$, and $V_{j+1}$ in the support of $\tau_h$; in particular $V_{j+1}$ lies in its window on $\tilde{\mathfrak{B}}^{(i)}_D$ \cite{B-CMP119}. By the lemma on the segments of a history, clause~(c), the function carried along $h$ passes from level $j$ to level $j+1$ through $\mathbb{E}_{i,h}$. Let $\mathbb{B}^+$ be the set of free bonds of the region $\Omega_{j+1}$ of the branch, so that $\mathbb{B}^+$ contains the bonds of $\tilde{\mathfrak{B}}^{(i)}_D$. Let $\mathcal{V}=\mathcal{V}_{V_{j+1}}$ be the map expressing the level-$j$ field through $V_{j+1}$ and the scaled fluctuation field $y$ on $\mathbb{B}^+$.

By Lemma~\ref{8.3:lem-box-law-mixture}, under its first hypothesis (the step of the branch near the box $\mathfrak{B}^{(i)}_D$ is a term of \eqref{8.3:eq-fluctuation-integral-a}), the following holds for every $G$ supported in the box: $\mathbb{E}_{i,h}[G](V_{j+1})$ is the normalised expectation of $G\circ\mathcal{V}$ under the conditional law $\mathbf{P}:=\mathbf{P}(\,\cdot\mid h,V_{j+1})$ of $y$ (see \cite[(3.25)]{B-CMP119}). This law has density
\begin{equation}\label{8.3:eq-interpolation-localisation-1}
\mathbf{P}(dy)\;\propto\;\chi_j(y)\,
\exp\Bigl[-\tfrac12\langle y,A^{\Gamma}_jy\rangle+P^{\mathrm{tot}}_j(y)\Bigr]\,dy ,
\end{equation}
where $A^{\Gamma}_j$ is the quadratic form of the step on $\mathbb{B}^+$ and $P^{\mathrm{tot}}_j$ is its perturbation. Put $\Gamma_j:=(A^{\Gamma}_j)^{-1}$.

In the complete model the background near the box $\mathfrak{B}^{(i)}_D$ also depends on the fields on the large-field zones \cite[(3.10)--(3.11)]{B-CMP119}. From here to Corollary~\ref{8.3:cor-interpolation-localisation}, every object is taken, in the complete model, at a fixed exterior background $\mathfrak{u}$, that is, after conditioning also on the background-determining members of the exterior datum. The objects $\mathbf{P}$, $A^{\Gamma}_j=\mathsf{Q}[\mathfrak{u}]$, $P^{\mathrm{tot}}_j(\cdot\,;\mathfrak{u})$ and $\mathcal{V}=\mathcal{V}_{\mathfrak{u}}$ are then those of that background.

Let $I$ be the set of free bonds of $\mathfrak{B}^{(i)}_D$ and $I^{\mathrm{c}}:=\mathbb{B}^+\setminus I$. Let $\mathcal{W}_I$ be the product of the windows on $I$, and write $\mathcal{W}_{I^{\mathrm{c}}}$ for the far window product on $I^{\mathrm{c}}$ of Proposition~\ref{8.3:prop-class-membership}. Disintegrating $\mathbf{P}$ over $y_{I^{\mathrm{c}}}$ gives
\begin{equation}\label{8.3:eq-interpolation-localisation-2}
\mathbb{E}_{i,h}[G]=\int\varpi(dy_{I^{\mathrm{c}}})\,
\langle G\circ\mathcal{V}\rangle_{I;y_{I^{\mathrm{c}}}},
\qquad
\varpi(dy_{I^{\mathrm{c}}})\propto\Psi(y_{I^{\mathrm{c}}})\,
\mathcal{Z}_I(y_{I^{\mathrm{c}}})\,dy_{I^{\mathrm{c}}} .
\end{equation}
The ingredients are as follows.
\begin{itemize}
\item $\mathcal{Z}_I(y_{I^{\mathrm{c}}})>0$ is the integral over $\mathcal{W}_I$ of the density \eqref{8.3:eq-interpolation-localisation-1} in $y_I$ at frozen $y_{I^{\mathrm{c}}}$.
\item $\Psi(y_{I^{\mathrm{c}}})$ is the remaining factor, which is non-negative by the second hypothesis of Lemma~\ref{8.3:lem-box-law-mixture}.
\item Hence $\varpi$ is a probability measure on $\mathcal{W}_{I^{\mathrm{c}}}$.
\item $\langle\cdot\rangle_{I;y_{I^{\mathrm{c}}}}$ is the conditional law of $y_I$, which lies in the class $\mathfrak{S}$ by Proposition~\ref{8.3:prop-class-membership}.
\end{itemize}

\emph{The box data.} The one-insertion and two-insertion estimates of the Gaussian extraction are applied to a box datum. It is obtained from $\langle\cdot\rangle_{I;y_{I^{\mathrm{c}}}}$ by switching off the couplings between $I$ and $I^{\mathrm{c}}$. In the setting of Theorem~\ref{8.3:thm-decoupling-localisation} this means: interpolation parameter $\sigma=0$, $O=\emptyset$, frozen exterior variable $y_{I^{\mathrm{c}}}$ on $I^{\mathrm{c}}$, linear term $0$, coupling block $A^{\Gamma}_{II^{\mathrm{c}}}$, and perturbation $P^{\mathrm{tot}}_j(y_I,y_{I^{\mathrm{c}}})$. Explicitly,
\begin{equation}\label{8.3:eq-interpolation-localisation-3}
\mu^{\mathrm{bx}}_j(dy_I)\;\propto\;\mathbf{1}_{\mathcal{W}_I}(y_I)\,
\exp\Bigl[-\tfrac12\langle y_I,A^{\Gamma}_{II}y_I\rangle
+P^{\mathrm{tot}}_j(y_I,0_{I^{\mathrm{c}}})\Bigr]\,dy_I .
\end{equation}
Its Gaussian part is
\[
\gamma^{\mathrm{bx}}_j=\mathcal{N}\bigl(0,(A^{\Gamma}_{II})^{-1}\bigr)
=\mathcal{N}\bigl(0,C^{(j)}(I)\bigr),
\]
where, as in \cite[(2.5)--(2.6)]{B-CMP116}, we write $C^{(j)}(I)$ for the step covariance $C^{(j)}(U_{j+1})$ conditioned on the complement of $I$; the background is suppressed in this notation. It is centred, and its window is the box's own per-bond balls.

Its perturbation $P^{\mathrm{bx}}_j:=P^{\mathrm{tot}}_j(\cdot,0_{I^{\mathrm{c}}})$ does not depend on $y_{I^{\mathrm{c}}}$. It depends on $h$ only through the definition of the clean event and the far tails of Proposition~\ref{8.3:prop-class-membership}. Near the box, $P^{\mathrm{tot}}_j$ is the perturbation of Ba{\l}aban's representation together with the increments of the far spent pieces of the small-piece statement.

Equivalently, up to an error of the form \eqref{8.3:eq-interpolation-localisation-6} (Theorem~\ref{8.3:thm-decoupling-localisation}\,(iii)), one may use the datum $\mu^{\mathrm{mg}}_j$. It is formed as in \eqref{8.3:eq-interpolation-localisation-3} with Gaussian part $\gamma^{\mathrm{mg}}_j=\mathcal{N}(0,\Gamma_j\lceil I)$, where $\Gamma_j=C^{(j)}(U_{j+1})$ is the step covariance on the free bonds of $\Omega_{j+1}$. This is exactly the box datum of the one-insertion form of the Gaussian extraction, and its Gaussian part $\gamma^{\mathrm{mg}}_j$ is the box Gaussian $\gamma^{\mathrm{box}}$ of Theorem~\ref{8.3:thm-box-localisation} (in the pure model; in the complete model, at the reference background).

In the complete model the background $\mathfrak{u}$ determines not only $\mathcal{V}$ but also $A^{\Gamma}_j=\mathsf{Q}[\mathfrak{u}]$ and $P^{\mathrm{tot}}_j(\cdot,0_{I^{\mathrm{c}}};\mathfrak{u})$, hence $\mu^{\mathrm{bx}}_j$ and $\mu^{\mathrm{mg}}_j$ themselves. Conditioning on $(h,V_{j+1},y_{I^{\mathrm{c}}})$ does not fix them. For this reason every object is taken at a fixed background. The passage to the reference background, with the mixture over backgrounds, is the content of Theorem~\ref{8.3:thm-box-localisation}.

\emph{The box data meet the hypotheses of the step estimate.} We check the three hypotheses: box Gaussian, box window and box perturbation.

Box Gaussian. Both Gaussian parts are centred. Their coordinate variances satisfy
\[
(A^{\Gamma}_{II})^{-1}(b,b)\le1/a_-,
\qquad
(\Gamma_j\lceil I)(b,b)\le1/a_-,
\]
by the decay lemma for exponentially local forms, clause~(iv). So the variance parameter is $\mathfrak{v}=n_{\mathfrak{g}}/a_-$, as in the one-insertion extraction.

Box window. We have $\mathcal{W}_I=\prod_{b\in I}\mathcal{B}^{\mathrm{w}}_j$, where $\mathcal{B}^{\mathrm{w}}_j$ is Ba{\l}aban's norm ball of radius $\mathfrak{b}_j$, and
\[
\{|y_b|_2<\mathfrak{b}_j\}\subset\mathcal{B}^{\mathrm{w}}_j\subset\{|y_b|_2<\hat{\mathfrak{b}}_j\}.
\]
The window sandwich required by the Gaussian extraction therefore holds, as a corollary, with the constant $C_w:=\|\mathbf{C}\|$, multiplied by $c_N$ when the Euclidean norm is used on the fibre. Let $n_D:=|I|$. The window-exit term of the step estimate is then unchanged:
\[
t_j\le2n_{\mathfrak{g}}n_D\exp\bigl(-\mathfrak{b}_j^2/(8n_{\mathfrak{g}}\|\mathbf{C}\|^2\mathfrak{v})\bigr).
\]

Box perturbation. $P^{\mathrm{bx}}_j$ is real. The window $\mathcal{W}_I$ is star-shaped, so integrating along the segment from $0$, with first-derivative bound $C_Pg_j\mathfrak{b}_j^2$ per bond, and $|I|=n_D$ bonds, gives
\begin{equation}\label{8.3:eq-interpolation-localisation-4}
\sup_{\mathcal{W}_I}\bigl|P^{\mathrm{bx}}_j-P^{\mathrm{bx}}_j(0)\bigr|
\le C_Pg_j\mathfrak{b}_j^2\,n_D\,\hat{\mathfrak{b}}_j\le\Pi_j .
\end{equation}
Here $\hat{\mathfrak{b}}_j\le\sqrt{N}\,\mathfrak{b}_j\le n_{\mathfrak{g}}C_w\mathfrak{b}_j$, and $\Pi_j$ is the perturbation bound of the one-insertion extraction. The constant $P^{\mathrm{bx}}_j(0)$ cancels in every normalised expectation and covariance.

\emph{Plaquette covariances inside the box.} We write $\mathfrak{c}^{(j)}_{\mathrm{br}}$ for the plaquette covariance of the branch, the form determined by the nest of the branch (the nested sequence of its regions), and keep $\mathfrak{c}^{(j)}$ for Ba{\l}aban's whole-lattice plaquette covariance. For $\mu^{\mathrm{mg}}_j$ the plaquette covariance is exactly the branch's plaquette covariance, because $\gamma^{\mathrm{mg}}_j$ is the marginal:
\[
\mathfrak{c}^{\mathrm{box},(j)}(q,q')=\mathfrak{c}^{(j)}_{\mathrm{br}}(q,q';U_{j+1}).
\]
Its identification near the support with Ba{\l}aban's whole-lattice covariance $\mathfrak{c}^{(j)}$ is Corollary~\ref{8.3:cor-extracted-covariance}.

For $\mu^{\mathrm{bx}}_j$ we have $\mathfrak{c}^{\mathrm{box},(j)}=\mathfrak{c}^{(j)}_{\mathrm{br}}+\delta\mathfrak{c}$. For plaquettes $q,q'$ near the support, the decay lemma, clauses~(iii) and~(iv), gives
\[
|\delta\mathfrak{c}(q,q')|
\le16\|\mathbf{C}\|^2\frac{a_+}{a_-^2}\,C'_{\mathrm{dec}}\,
e^{-\mu_A(d_{I^{\mathrm{c}}}(q)+d_{I^{\mathrm{c}}}(q'))}
\le C\,\mathfrak{v}_{\mathfrak{c}}\,e^{-2\mu_A(DL-M-R_{\mathbf{C}}-2)} .
\]
Here $\mu_A$ and $C'_{\mathrm{dec}}$ are the rate and the constant of that lemma for $(a_-,a_+,c_A)$, $d_{I^{\mathrm{c}}}(q)$ is the distance from $q$ to $I^{\mathrm{c}}$, and $\mathfrak{v}_{\mathfrak{c}}$ is the plaquette-covariance scale of the one-insertion extraction. This is a term of order $e^{-c''D}$, of the same kind as the terms $g_j^2\mathfrak{b}_j^2e^{-c_{\mathrm{ip}}DL}$ of that extraction, with $c_{\mathrm{ip}}$ the rate of Corollary~\ref{8.3:cor-interpolation-localisation}.

In either case the Gaussian moment lemmas for one and two insertions, the Wick lemma and the step estimate apply to the box datum as stated. With $\mu^{\mathrm{mg}}_j$ the one- and two-insertion estimates of the Gaussian extraction run unchanged. With $\mu^{\mathrm{bx}}_j$ they run with $\mathfrak{c}^{\mathrm{box},(j)}=\mathfrak{c}^{(j)}_{\mathrm{br}}$ replaced by $\mathfrak{c}^{\mathrm{box},(j)}=\mathfrak{c}^{(j)}_{\mathrm{br}}+\delta\mathfrak{c}$, and with one more term of order $e^{-c''D}$, subject to the corresponding adjustment of the cut in the two-insertion estimate of the Gaussian extraction.

\begin{corollary}[Localisation of the sourced step by interpolation]
\label{8.3:cor-interpolation-localisation}
Assume the following.
\begin{itemize}
\item The antecedents of localisation of the sourced step, Theorem~\ref{8.3:thm-sourced-step-localisation}, as listed in its statement.
\item The hypothesis from which exponential decay of correlations in the class, Lemma~\ref{8.3:lem-correlation-decay}, is proved, as stated in that lemma.
\item The finite-range hypothesis on the map through which the box-exterior variables enter $\mathcal{V}$.
\end{itemize}

Let $h\in\mathcal{C}_S(D)$, let $k_0\le j<k_0+r$ and $i=j-k_0$, and let $V_{j+1}$ lie in the support of $\tau_h$. Let $g_{k_0}\le\gamma_{\mathrm{ex}}$, where $\gamma_{\mathrm{ex}}$ is the coupling ceiling of the Gaussian extraction; it incorporates the condition defining $\gamma_H$ and the condition $x_{k_0}\ge x_{\mathrm{pk}}$ of \eqref{8.3:eq-depth-coupling-conditions-a}. Let $I$ be the set of free bonds of $\mathfrak{B}^{(i)}_D$. Let $G,G'$ be bounded measurable functions that depend on $V_j$ only through bonds within distance $R_S+M$ of $S^{(i)}$. Put
\[
\ell_{\mathrm{geo}}:=M+R_{\mathbf{C}}+2
\]
in units of the level-$j$ lattice, the units in which $DL$ and $\ell_S$ below are also measured (no further factor of $L$ enters); this number depends on $L,M,d$ only. Let $D\ge D_0:=\ell_{\mathrm{geo}}+1$, as in \eqref{8.3:eq-depth-coupling-conditions-c}.
\begin{enumerate}
\item[(a)] (Pure model, $\mathcal{V}:=\mathcal{V}_{V_{j+1}}$.) Under this finite-range hypothesis, no truncation of $\mathcal{V}$ is needed, and $G\circ\mathcal{V}$ depends on $y$ only through $y\lceil\tilde{S}_G$ for a set of bonds $\tilde{S}_G\subset I$ at distance $\ell_S\ge DL-\ell_{\mathrm{geo}}$ from $I^{\mathrm{c}}$. For the box data $\mu^{\mathrm{bx}}_j$ of \eqref{8.3:eq-interpolation-localisation-3} and $\mu^{\mathrm{mg}}_j$ defined after it, with the window product $\mathcal{W}_I$ defined above, and for $\natural\in\{\mathrm{bx},\mathrm{mg}\}$,
\begin{equation}\label{8.3:eq-interpolation-localisation-5}
\begin{aligned}
\mathbb{E}_{i,h}[G](V_{j+1})&=\langle G\circ\mathcal{V}\rangle_{\mu^{\natural}_j}+\Delta_G,\\
\operatorname{Cov}_{\mathbb{E}_{i,h}}(G,G')
&=\operatorname{Cov}_{\mu^{\natural}_j}(G\circ\mathcal{V},G'\circ\mathcal{V})+\Delta_{GG'} .
\end{aligned}
\end{equation}
For all constants $c,c'$,
\begin{equation}\label{8.3:eq-interpolation-localisation-6}
\begin{aligned}
|\Delta_G|&\le C_{\mathrm{ip}}\|G\circ\mathcal{V}-c\|_{\mathcal{W}_I}\,e^{-c_{\mathrm{ip}}DL},\\
|\Delta_{GG'}|&\le C_{\mathrm{ip}}\|G\circ\mathcal{V}-c\|_{\mathcal{W}_I}
\|G'\circ\mathcal{V}-c'\|_{\mathcal{W}_I}\,e^{-c_{\mathrm{ip}}DL}\\
&\quad+4C_{\mathrm{ip}}^2\|G\circ\mathcal{V}-c\|_{\mathcal{W}_I}
\|G'\circ\mathcal{V}-c'\|_{\mathcal{W}_I}\,e^{-2c_{\mathrm{ip}}DL} .
\end{aligned}
\end{equation}
The constants are as follows.
\begin{itemize}
\item $c_{\mathrm{ip}}:=\nu_1>0$ is the rate of Theorem~\ref{8.3:thm-decoupling-localisation}. It depends only on $a_-,a_+,c_A,C_P,c_P,n_{\mathfrak{g}},d$ and not on the level.
\item We put
\[
C_{\mathrm{ip}}:=8\,C\,n_G\,\hat{\mathfrak{b}}_j^{\,2}\,(a_++a_-)
\bigl(1+\bar{\mathfrak{K}}/m_\star\bigr)\,e^{\nu_1\ell_{\mathrm{geo}}} ,
\]
where $n_G:=\#\tilde{S}_G\le C(R_S+M+R_{\mathbf{C}})^4$.
\item $C=C(d,\nu;1\vee C_O,\delta_w,c_N)$ absorbs $(1\vee C_O)(1+\delta_w)c_N$.
\item $m_\star$ is the convexity constant of the class in Lemma~\ref{8.3:lem-correlation-decay}.
\item $\bar{\mathfrak{K}}\le(a_++a_-)S_d(\nu_0)$ is the constant of Theorem~\ref{8.3:thm-decoupling-localisation}.
\end{itemize}

If $G\circ\mathcal{V}$ is holomorphic on the complex box window of the holomorphic corollary of Theorem~\ref{8.3:thm-decoupling-localisation}, with oscillation $\operatorname{osc}^{\mathbb{C}}(G\circ\mathcal{V})$ there, then
\[
|\Delta_G|\le C_{\mathrm{ip}}^{\nabla}\,n_G\,\operatorname{osc}^{\mathbb{C}}(G\circ\mathcal{V})\,e^{-c_{\mathrm{ip}}DL}.
\]
Here $C_{\mathrm{ip}}^{\nabla}$ does not depend on the level and contains no power of $\mathfrak{b}_j$. By the membership statement for the carried class, this applies to the observables carried through the extraction and to their main parts.
\item[(b)] (Complete model, given the exterior background $\mathfrak{u}$.) Replace $\mathbb{E}_{i,h}[\cdot]$ by the branch expectation further conditioned on the background-determining members of the exterior datum, written $\mathbb{E}_{i,h}[\,\cdot\mid\mathfrak{u}](V_{j+1})$. Let $\mu^{\natural}_j[\mathfrak{u}]$ be the box datum formed, for $\natural=\mathrm{bx}$, with the form $\mathsf{Q}_{II}[\mathfrak{u}]$ and, for $\natural=\mathrm{mg}$, with its Schur complement $\bigl(\mathsf{Q}[\mathfrak{u}]^{-1}\lceil I\bigr)^{-1}$, in both cases with exponent $P^{\mathrm{tot}}_j(\cdot,0_{I^{\mathrm{c}}};\mathfrak{u})$. Let $\mathcal{V}=\mathcal{V}_{\mathfrak{u}}$. Then \eqref{8.3:eq-interpolation-localisation-5} and \eqref{8.3:eq-interpolation-localisation-6} hold with the same $C_{\mathrm{ip}},c_{\mathrm{ip}}$, uniformly in $\mathfrak{u}$. Every admissible background lies, by the fluctuation bound on admissible backgrounds, in the real domain of the operator theorem at free current, so the class constants of Proposition~\ref{8.3:prop-class-membership} apply.
\end{enumerate}
In the complete model the unconditional statement, with the mixture over backgrounds, is the content of Theorem~\ref{8.3:thm-box-localisation}\,(A),~(C) together with the comparison of mixtures of laws used in its proof.
\end{corollary}

\begin{proof}
We argue in the pure model. In the complete model every object is read at a fixed background $\mathfrak{u}$. Under that background the box datum $\mu^{\natural}_j[\mathfrak{u}]$ does not depend on $y_{I^{\mathrm{c}}}$, which is the only property used in the covariance step below. Uniformity in $\mathfrak{u}$ holds because the constants used are those of the class, which apply to every admissible background.

\emph{One point, $\natural=\mathrm{bx}$.} Since $\varpi$ in \eqref{8.3:eq-interpolation-localisation-2} is a probability measure and $\mu^{\mathrm{bx}}_j$ does not depend on $y_{I^{\mathrm{c}}}$,
\[
\mathbb{E}_{i,h}[G]-\langle G\circ\mathcal{V}\rangle_{\mu^{\mathrm{bx}}_j}
=\int\varpi(dy_{I^{\mathrm{c}}})\Bigl[\langle G\circ\mathcal{V}\rangle_{I;y_{I^{\mathrm{c}}}}
-\langle G\circ\mathcal{V}\rangle_{\mu^{\mathrm{bx}}_j}\Bigr].
\]
Theorem~\ref{8.3:thm-decoupling-localisation}\,(i), specialised to $O=\emptyset$, bounds the bracket uniformly in the frozen $y_{I^{\mathrm{c}}}$. The hypotheses of that theorem are supplied as follows.
\begin{itemize}
\item The class constants are those of Proposition~\ref{8.3:prop-class-membership}.
\item The separation condition is void, because $O=\emptyset$.
\item $c_1=\min(c_A,c_P)$.
\item $\bar{r}\le2(1\vee C_O)\hat{\mathfrak{b}}_j(1+\delta_w)$.
\end{itemize}
The constant of that theorem, written $\bar\omega$ here, satisfies
\[
\bar\omega=4\bar{r}(a_++p_2)S_d(c_1/2)\le C'\mathfrak{b}_j(a_++a_-).
\]
Indeed, $\bar{r}a_+$ is of order $\mathfrak{b}_ja_+$. For the $p_2$-term, the bound $p_2S_d(c_P)\le a_-/4$ on the class constants, displayed in Proposition~\ref{8.3:prop-class-membership}, gives
\[
\bar{r}p_2S_d(c_1/2)\le\frac{\bar{r}a_-S_d(c_1/2)}{4S_d(c_P)}=O(\mathfrak{b}_ja_-).
\]
(In fact the $p_2$-term is smaller than the Gaussian one by a factor of order $g_j\mathfrak{b}_j$; the display shows that $(a_++a_-)$ is in any case an envelope for both.)

Hence the bracket is at most
\[
C_1\|G\circ\mathcal{V}-c\|_{\mathcal{W}_I}\,n_G\,\bar{r}\,\bar\omega\,
\bigl(1+\bar{\mathfrak{K}}/m_\star\bigr)e^{-\nu_1\ell_S},
\]
with $C_1=C_1(d,\nu)$ the lattice-sum constant of Theorem~\ref{8.3:thm-decoupling-localisation}. We now collect the factors.
\begin{itemize}
\item $\ell_S\ge DL-\ell_{\mathrm{geo}}$ gives $e^{-\nu_1\ell_S}\le e^{\nu_1\ell_{\mathrm{geo}}}e^{-\nu_1DL}$.
\item $\bar{r}\,\bar\omega\le2(1\vee C_O)(1+\delta_w)C'\,\hat{\mathfrak{b}}_j\mathfrak{b}_j(a_++a_-)$ and $\mathfrak{b}_j\le\hat{\mathfrak{b}}_j$.
\item $C_1$, $C'$, $(1\vee C_O)$, $(1+\delta_w)$ and $c_N$ are absorbed into $C$.
\end{itemize}
With these, the bound is at most $C_{\mathrm{ip}}\|G\circ\mathcal{V}-c\|_{\mathcal{W}_I}e^{-c_{\mathrm{ip}}DL}$. Integrating against $\varpi$ gives the bound on $\Delta_G$ in \eqref{8.3:eq-interpolation-localisation-6}.

\emph{Two points, $\natural=\mathrm{bx}$.} The law of total covariance for the disintegration \eqref{8.3:eq-interpolation-localisation-2} gives
\[
\begin{aligned}
\operatorname{Cov}_{\mathbb{E}_{i,h}}(G,G')
&=\int\varpi(dy_{I^{\mathrm{c}}})\operatorname{Cov}_{I;y_{I^{\mathrm{c}}}}(G\circ\mathcal{V},G'\circ\mathcal{V})\\
&\quad+\operatorname{Cov}_{\varpi}\bigl(\langle G\circ\mathcal{V}\rangle_{I;y_{I^{\mathrm{c}}}},
\langle G'\circ\mathcal{V}\rangle_{I;y_{I^{\mathrm{c}}}}\bigr).
\end{aligned}
\]
For every $y_{I^{\mathrm{c}}}$, Theorem~\ref{8.3:thm-decoupling-localisation}\,(ii) bounds
\[
\bigl|\operatorname{Cov}_{I;y_{I^{\mathrm{c}}}}(G\circ\mathcal{V},G'\circ\mathcal{V})
-\operatorname{Cov}_{\mu^{\mathrm{bx}}_j}(G\circ\mathcal{V},G'\circ\mathcal{V})\bigr| ;
\]
with the constants evaluated as in the one-point case, this bound is the first term of the bound on $\Delta_{GG'}$. Since $\varpi$ is a probability measure, the same bound holds after integration.

For the second term we use $|\operatorname{Cov}_{\varpi}(X,Y)|\le\operatorname{osc}(X)\operatorname{osc}(Y)$. Because the box datum does not depend on $y_{I^{\mathrm{c}}}$,
\[
\operatorname{osc}_{y_{I^{\mathrm{c}}}}\langle G\circ\mathcal{V}\rangle_{I;y_{I^{\mathrm{c}}}}
\le2\sup_{y_{I^{\mathrm{c}}}}\bigl|\langle G\circ\mathcal{V}\rangle_{I;y_{I^{\mathrm{c}}}}
-\langle G\circ\mathcal{V}\rangle_{\mu^{\mathrm{bx}}_j}\bigr|
\le2C_{\mathrm{ip}}\|G\circ\mathcal{V}-c\|_{\mathcal{W}_I}e^{-c_{\mathrm{ip}}DL},
\]
and likewise for $G'$. So the second term is at most
\[
4C_{\mathrm{ip}}^2\|G\circ\mathcal{V}-c\|_{\mathcal{W}_I}\|G'\circ\mathcal{V}-c'\|_{\mathcal{W}_I}\,
e^{-2c_{\mathrm{ip}}DL},
\]
which is the second term of \eqref{8.3:eq-interpolation-localisation-6}.

\emph{The case $\natural=\mathrm{mg}$.} We apply Theorem~\ref{8.3:thm-decoupling-localisation}\,(iii) with $A:=A^{\Gamma}_j$ on $\mathbb{B}^+$. Taking $\mathfrak{Y}:=\mathbb{B}^+$ in Lemma~\ref{8.3:lem-principal-blocks}, the complement $\mathcal{O}=I^{\mathrm{c}}$ consists of the frozen far bonds; the blocks $A^{\Gamma}_{II^{\mathrm{c}}}$ obey that lemma's bounds, which depend only on $(a_-,a_+,c_A)$. This passes from $\mu^{\mathrm{bx}}_j$ to $\mu^{\mathrm{mg}}_j$ with an error of the same form.

\emph{The holomorphic form.} The bound in terms of $\operatorname{osc}^{\mathbb{C}}$ is the holomorphic corollary of Theorem~\ref{8.3:thm-decoupling-localisation}, applied with the same profile.
\end{proof}

Along this route the localisation used by the Gaussian extraction holds with prefactor $C_{\mathrm{ip}}$ and rate $c_{\mathrm{ip}}$; these two letters are reserved for the constants of Corollary~\ref{8.3:cor-interpolation-localisation}. The prefactor $C_{\mathrm{ip}}$ carries $\mathfrak{b}_j^2$, which is polylogarithmic in $x_j$. It is needed where the carried two-point functions are only bounded, and it is absorbed by the free choice of $D$ in each application. In the complete model the unconditional statement has the following ingredients. For an exterior datum $\mathsf{e}$ with background $\mathfrak{u}(\mathsf{e})$ and the reference datum $\mathsf{e}_0$ of Theorem~\ref{8.3:thm-box-localisation}, they are: the passage from $\mu^{\natural}_j[\mathfrak{u}(\mathsf{e})]$ to the reference datum at $\mathsf{e}_0$; the mixture over backgrounds; the error terms $\eta_G$ of Theorem~\ref{8.3:thm-box-localisation}\,(C); the variation of the form with the background; and the background part of the boundary term of the exterior localisation statement, all of which belong to Theorem~\ref{8.3:thm-box-localisation}\,(A),~(C) and the comparison of mixtures of laws used in its proof.

\emph{The extracted covariance.} At step $j$, the one-insertion extraction extracts the coefficient
\[
n_{\mathfrak{g}}\,g_j^2\langle H^{\mathcal{G}}(V_0),\mathfrak{c}^{\mathrm{box},(j)}(V_0)\rangle
\]
for an observable $\mathcal{G}$. Here $H^{\mathcal{G}}$ is the kernel attached to $\mathcal{G}$ in the Gaussian extraction and $V_0:=V^{(j)}(V_{j+1})$ is the background of the extraction. The two-line term of the two-insertion extraction is
\[
2n_{\mathfrak{g}}\,g_j^4\operatorname{tr}\bigl(H^{\mathcal{G}_1}\mathfrak{c}^{\mathrm{box},(j)}
H^{\mathcal{G}_2}\mathfrak{c}^{\mathrm{box},(j)}\bigr).
\]
By the comparison of plaquette covariances above, $\mathfrak{c}^{\mathrm{box},(j)}=\mathfrak{c}^{(j)}_{\mathrm{br}}(\cdot\,;U_{j+1})$, where
\[
\mathfrak{c}^{(j)}_{\mathrm{br}}(q,q';U):=n_{\mathfrak{g}}^{-1}\operatorname{tr}_{\mathfrak{g}}
\operatorname{Cov}_{\mu_{\Gamma_j(U)}}\bigl((\partial_U\mathbf{C}y)(q),(\partial_U\mathbf{C}y)(q')\bigr)
\]
is the colour-averaged plaquette covariance of Ba{\l}aban's step-$j$ Gaussian \cite[(2.13)]{B-CMP109} on the free bonds of the branch; with Ba{\l}aban's whole-lattice step covariance $C^{(j)}(U)$ of \cite[(2.13)]{B-CMP109} in place of $\Gamma_j(U)$, the same expression is the whole-lattice covariance $\mathfrak{c}^{(j)}(q,q';U)$ of Corollary~\ref{8.3:cor-extracted-covariance}. Here $(\partial_U\mathbf{C}y)(q)$ is the plaquette variable of the bond field $\mathbf{C}y$ at $q$ in the background $U$, as in the statement of the Gaussian extraction, $\operatorname{tr}_{\mathfrak{g}}$ is the trace over the colour components, and $\mu_{\Gamma}$ is the centred Gaussian measure of covariance $\Gamma$. The identity is exact for $\mu^{\mathrm{mg}}_j$ and holds up to $\delta\mathfrak{c}$ for $\mu^{\mathrm{bx}}_j$. Between the nest form $\mathfrak{c}^{(j)}_{\mathrm{br}}$ and the whole-lattice form $\mathfrak{c}^{(j)}$ it holds up to the term of order $e^{-\mu_1(DL-\ell_\star)}$ of Corollary~\ref{8.3:cor-extracted-covariance}.

The dictionary identification on which Corollary~\ref{8.3:cor-extracted-covariance} rests identifies $\mathfrak{c}^{(j)}(\cdot\,;\mathbf{1})$, Ba{\l}aban's whole-lattice plaquette covariance at the flat background $U\equiv\mathbf{1}$, with the increment $\mathfrak{c}_{j-k_0}$ of the free tower, the sequence of free step covariances transported to level $k_0$. The tower covariance at level $k_0$ is
\[
\mathfrak{c}^{\mathrm{tw}}_{k_0}=\sum_{i\ge0}\mathfrak{c}_i,
\]
whose tail $\sum_{i\ge r}\mathfrak{c}_i$ is bounded by $C_\eta L^{-4r}$. The witness kernel of the two-point witness theorem is
\[
\mathcal{K}_s(f_1,f_2)=\sum_{p,q}f_1(x(p))f_2(x(q))\bigl[C^{\mathrm{bg}}_{K,s}(p,q)\bigr]^2 ,
\]
the sum running over pairs of plaquettes, with $C^{\mathrm{bg}}_{K,s}$ the kernel $G^{(s)}_K$ of clause~(iii)(c) of Theorem~\ref{3.1:thm-main}. This kernel is the full tower.

The main term of the two-insertion extraction,
\[
\mathcal{T}_r(\mathbf{1})=2n_{\mathfrak{g}}\sum_{i,i'<r}g^2g^2
\operatorname{tr}\bigl(H^{\mathcal{G}_1}\mathfrak{c}_iH^{\mathcal{G}_2}\mathfrak{c}_{i'}\bigr),
\]
with the couplings as they enter the two-insertion estimate of the Gaussian extraction, is its truncation at $r$. The two differ by the tail levels $\ge r$, which remain with the two-point witness theorem.

We therefore have the following chain. The covariance used by the one- and two-insertion estimates is Ba{\l}aban's $C^{(j)}(U_{j+1})$, restricted to the box. At flat background it is the free step increment. Summed over the extracted levels, it is the truncated tower, that is, $C^{\mathrm{bg}}_{K,s}$ minus the tail of the witness theorem. No link of this chain depends on the couplings at the boundary of the box except through the factor $e^{-c_{\mathrm{ip}}DL}$ of Theorem~\ref{8.3:thm-decoupling-localisation} and the factor $e^{-2\mu_ADL}$ of the decay lemma, clause~(iii).

\begin{definition}[The step, its zones and the determining set with its data]\label{8.4:not-step-zones}
\emph{Matrices.} The dimension is $d=4$. Let $G=\mathrm{SU}(N)$ with $N\ge 2$. Let $\mathfrak{g}$ be the real
vector space of Hermitian traceless $N\times N$ matrices, so that $\exp(iX)\in G$ for $X\in\mathfrak{g}$. We
write $G^{c}=\mathrm{SL}(N,\mathbb{C})$ and $\mathfrak{g}^{c}=\mathfrak{sl}(N,\mathbb{C})=\mathfrak{g}\oplus
i\mathfrak{g}$. For $z\in\mathfrak{g}^{c}$ we write $z=z'+iz''$, where
\begin{equation*}
z':=\tfrac12\,(z+z^{*}),\qquad z'':=\tfrac{1}{2i}\,(z-z^{*}),
\end{equation*}
and $z',z''\in\mathfrak{g}$. The norm $|\cdot|$ is the operator norm on complex $N\times N$ matrices, and
$\operatorname{tr}:=\operatorname{Tr}/N$, so that $|\operatorname{tr}A|\le|A|$ and
$|\operatorname{tr}(AB)|\le|A|\,|B|$. For $g\in\mathrm{U}(N)$ we put
$\operatorname{Ad}_{g}X:=gXg^{-1}$; then $|\operatorname{Ad}_{g}X|=|X|$ and
\begin{equation*}
|\operatorname{Ad}_{g}X-X|\le 2\,|g-1|\,|X|.
\end{equation*}

\emph{The step.} A step index $j\ge1$ with $k_{0}\le j\le k_{0}+r$ is fixed, where $k_{0}$ is the first and
$k_{0}+r$ the last of the consecutive steps considered. The blocking factor $L$ is an
odd integer with $L\ge L_{0}$, where $L_{0}=L_{0}(N)$ is the constant of hypothesis (H2)
(Definition~\ref{2.2:not-hypotheses}); no further lower bound on $L$ is imposed in what follows. The fine
lattice $\mathbb{T}$ is a finite periodic lattice of spacing $\eta$, and lengths are normalised so that the
level-$j$ lattice has unit spacing, $L^{j}\eta=1$. For $0\le n\le j$ the level-$n$ lattice
$\mathbb{T}^{(n)}$ consists of the centres $z$ of the level-$n$ blocks $B^{n}(z)$; these are cubes of
$L^{nd}$ fine sites, of side $L^{n}\eta=L^{n-j}$. The bonds of $\mathbb{T}^{(n)}$ join nearest level-$n$
neighbours. A configuration on a set of oriented bonds is a map $U$ into $G$ (or into $G^{c}$) with
$U(b^{-1})=U(b)^{-1}$, where $b^{-1}$ is $b$ with the opposite orientation. Let $e_{\mu}$,
$\mu=1,\dots,d$, be the unit coordinate vectors. For a fine plaquette $p$ with base vertex $x_{p}$ in the
coordinate directions $\mu<\nu$ we take the bonds
\begin{align*}
b_{1}&=\langle x_{p},\,x_{p}+\eta e_{\mu}\rangle, &
b_{2}&=\langle x_{p}+\eta e_{\mu},\,x_{p}+\eta e_{\mu}+\eta e_{\nu}\rangle,\\
b_{3}&=\langle x_{p}+\eta e_{\nu},\,x_{p}+\eta e_{\mu}+\eta e_{\nu}\rangle, &
b_{4}&=\langle x_{p},\,x_{p}+\eta e_{\nu}\rangle,
\end{align*}
and put $U(\partial p):=U(b_{1})U(b_{2})U(b_{3})^{-1}U(b_{4})^{-1}$.

\emph{The domains of the term.} The history $h$ selects, in the density of step $j$
(\cite[(2.18)]{B-CMP119}, or its form modified by the large-field operation $\mathbf{R}$,
\cite[(1.12)]{B-CMP122a}), a term with an admissible sequence of domains
\begin{equation*}
\mathbb{T}\supseteq\Omega_{1}\supseteq\Omega_{2}\supseteq\cdots\supseteq\Omega_{j}\supseteq\Omega_{j+1},
\end{equation*}
each $\Omega_{n}$ being a union of level-$n$ blocks, regarded as a set of fine sites
(\cite[(2.1), pp.~254--256]{B-CMP119}); $\Omega_{j+1}$ is the domain on which the fluctuation field of
step $j$ lives (\cite[(3.10)--(3.11), (3.15)]{B-CMP119}). The \emph{zones} are
\begin{equation*}
Z_{0}:=\Omega_{1}^{c},\qquad Z_{n}:=\Omega_{n}\setminus\Omega_{n+1}\quad(1\le n\le j-1),\qquad
Z_{j}:=\Omega_{j}.
\end{equation*}
For a fine site $x$, $n(x)$ is the index of the zone containing $x$; for a fine bond
$b=\langle x,x+\eta e_{\mu}\rangle$ we put $n(b):=n(x)$, and for a fine plaquette $p$ we put
$n(p):=n(x_{p})$. The \emph{determining set} is
\begin{equation*}
\mathbb{B}_{j}:=\bigcup_{n=0}^{j}\Gamma_{n},\qquad
\Gamma_{n}:=\{z\in\mathbb{T}^{(n)}:\ B^{n}(z)\subset Z_{n}\}
\end{equation*}
(\cite[(2.2)]{B-CMP119}); by admissibility of the sequence of domains, each $Z_{n}$ is a union of
level-$n$ blocks.

\emph{Data on the determining set.} The data on $\mathbb{B}_{j}$ are $V=(V_{n})_{0\le n\le j}$, where
$V_{n}$ is a $G$-valued configuration on the bonds of $\mathbb{T}^{(n)}$ between points of $\Gamma_{n}$
(\cite[(2.10)]{B-CMP119}). Here $V_{0}$ on $\Gamma_{0}$ is the bare field on $\Omega_{1}^{c}$; for
$1\le n\le j-1$, $V_{n}$ are the level-$n$ variables on $Z_{n}$ that have not yet been integrated out,
as in \cite[p.~256]{B-CMP119}: ``In $j$-th step the region is $\Omega_{j}^{c}$, and gauge field
variables $V_{j-1}$ are defined on bonds of $\Omega_{j}^{(j-1)c}$''; and $V_{j}$ on $\Gamma_{j}$ is the top
datum. With $M^{n}$ the $n$-fold block average of \cite[(2.11)]{B-CMP119} and $M^{0}$ the identity, we put
\begin{equation*}
M_{\mathbb{B}_{j}}(U):=\bigl(M^{n}(U)\big|_{\Gamma_{n}}\bigr)_{0\le n\le j}.
\end{equation*}
We write $U(\mathbb{B}_{j},V)$ for the minimal configuration of the Wilson functional $A$
over the set $\{U:M_{\mathbb{B}_{j}}(U)=V\}$, in the axial gauge
(\cite[(2.12)]{B-CMP119}; \cite[Theorem~1]{B-CMP102b}), whenever this minimal configuration exists and is
unique.

\emph{Bond sets.} $E_{1}$ is the set of fine bonds with at least one endpoint in $\Omega_{1}$; $E_{0}$ is the
set of fine bonds with both endpoints in $\Omega_{1}^{c}$; $C_{0}\subset E_{0}$, the
\emph{$1$-collar}, is
the set of bonds of $E_{0}$ that lie in some fine plaquette containing a bond of $E_{1}$. A zone-$0$
plaquette of $\mathbb{B}_{j}$ is called \emph{deep} if all its bonds lie in
$E_{0}\setminus C_{0}$.

\emph{The regularity hypothesis and its restricted form.} For a plaquette $p'$ of the determining set we
write $(\partial V)(p')$ for the plaquette variable of $V$ at $p'$, formed from the bond variables of $V$
around $p'$ as $U(\partial p)$ is formed above; $p'$ is a level-$n$ plaquette when it lies in zone $n$,
with the convention stated in the text accompanying \cite[(7)]{B-CMP102b} at interfaces between zones.
Hypothesis
\cite[(7)]{B-CMP102b} requires, for the constant of that hypothesis, which we write $\hat\varepsilon_{1}$,
and for every
plaquette $p'$ of the determining set,
\begin{equation}\label{8.4:eq-step-zones-1}
|(\partial V)(p')-1|<\hat\varepsilon_{1},
\end{equation}
where $1$ is the identity matrix. For a number $\bar\varepsilon>0$, the \emph{restricted regularity
hypothesis} $(7^{\flat})(\bar\varepsilon)$ is the same inequality with $\hat\varepsilon_{1}=\bar\varepsilon$,
required at every plaquette of $\mathbb{B}_{j}$ except the deep zone-$0$ plaquettes:
\begin{equation}\label{8.4:eq-step-zones-a}
\begin{gathered}
|(\partial V)(p')-1|<\bar\varepsilon\\
\text{for every plaquette $p'$ of $\mathbb{B}_{j}$ other than the deep zone-$0$ plaquettes.}
\end{gathered}
\end{equation}
\end{definition}

\begin{proof}[Proof of the elementary assertions about matrices]
Let $X\in\mathfrak{g}$. Since $X$ is Hermitian, $iX$ is anti-Hermitian, so
$\exp(iX)^{*}=\exp(-iX)=\exp(iX)^{-1}$ and $\exp(iX)$ is unitary; moreover
$\det\exp(iX)=\exp(i\operatorname{Tr}X)=1$ because $X$ is traceless. Hence $\exp(iX)\in G$.

Let $z\in\mathfrak{g}^{c}$, so $\operatorname{Tr}z=0$ and therefore
$\operatorname{Tr}z^{*}=\overline{\operatorname{Tr}z}=0$. The matrix $z'=\tfrac12(z+z^{*})$ satisfies
$(z')^{*}=z'$, and
\begin{equation*}
(z'')^{*}=\frac{1}{-2i}\,(z^{*}-z)=\frac{1}{2i}\,(z-z^{*})=z'',
\end{equation*}
so both are Hermitian; both are traceless because $\operatorname{Tr}z=\operatorname{Tr}z^{*}=0$. Thus
$z',z''\in\mathfrak{g}$, and
$z'+iz''=\tfrac12(z+z^{*})+\tfrac12(z-z^{*})=z$.

Let $A$ be a complex $N\times N$ matrix and $(v_{1},\dots,v_{N})$ an orthonormal basis of $\mathbb{C}^{N}$.
Since $|\langle v_{a},Av_{a}\rangle|\le|A|$ for each $a$,
\begin{equation*}
|\operatorname{tr}A|=\frac1N\Bigl|\sum_{a=1}^{N}\langle v_{a},Av_{a}\rangle\Bigr|\le|A|.
\end{equation*}
Applying this to $AB$ and using $|AB|\le|A|\,|B|$ gives $|\operatorname{tr}(AB)|\le|A|\,|B|$.

Let $g\in\mathrm{U}(N)$. Multiplication by a unitary matrix on either side preserves the operator norm,
so $|\operatorname{Ad}_{g}X|=|gXg^{-1}|=|X|$; the same fact gives $|g^{-1}|=1$ and
\begin{equation*}
|g^{-1}-1|=|g^{*}-1|=|(g-1)^{*}|=|g-1|.
\end{equation*}
From the identity
$gXg^{-1}-X=(g-1)Xg^{-1}+X(g^{-1}-1)$ we obtain
\begin{equation*}
|\operatorname{Ad}_{g}X-X|\le|g-1|\,|X|\,|g^{-1}|+|X|\,|g^{-1}-1|=2\,|g-1|\,|X|.
\end{equation*}
\end{proof}

\begin{definition}[The box variation, tubes, covariant curl and distances]\label{8.4:not-box-variation}
We use the step $j$, the lattices $\mathbb{T}$ and $\mathbb{T}^{(n)}$, the fine spacing $\eta$
with $L^{j}\eta=1$, the domains $\Omega_{n}$, the zones and the zone index $n(\cdot)$, the
determining set $\mathbf{B}_{j}=\bigcup_{n=0}^{j}\Gamma_{n}$, the data
$V=(V_{n})_{0\le n\le j}$ on it and the minimiser $U(\mathbf{B}_{j},V)$, all as fixed in
Definition~\ref{8.4:not-step-zones}, together with the decomposition $z=z'+iz''$ of
$z\in\mathfrak{g}^{c}$ into $z',z''\in\mathfrak{g}$ and the operator norm $|\cdot|$ fixed there.

\emph{The box and the variation.} Let $I$ be a finite set of bonds of $\mathbb{T}^{(j)}$, the
bond set of the box window. The variables $y=(y_{b})_{b\in I}\in(\mathfrak{g}^{c})^{I}$ are the
scaled fluctuation variables of the step $j$ on the box, and $\tilde{y}:=y\oplus 0$ denotes $y$
extended by $0$ to all bonds of $\mathbb{T}^{(j)}$ that carry the fluctuation field of the step
$j$. Let $\mathbf{C}$ be Ba{\l}aban's linear operator solving the linearised constraint
\cite{B-CMP119}, let $\mathbf{h}$ be Ba{\l}aban's corridor operator of \cite{B-CMP109}, which
Ba{\l}aban writes $h$, and let
$\tilde{D}$ be Ba{\l}aban's analytic map of \cite{B-CMP109}, written $\tilde{D}$ in
\cite[(3.15)]{B-CMP119}.
We put
\begin{equation}\label{8.4:eq-box-variation-a}
B(y):=g_{j}\mathbf{C}\tilde{y}-\mathbf{h}\tilde{D}\bigl(g_{j}\mathbf{C}\tilde{y}\bigr),\qquad
V'(y):=\exp\bigl(iB(y)\bigr)
\end{equation}
on the bonds of $\mathbb{T}^{(j)}$, with $V'(y):=1$ where $B(y)=0$.

The \emph{base data} are $V^{\mathrm{base}}:=(V_{0},\dots,V_{j-1},V_{j}^{\mathrm{base}})$, where
$V_{j}^{\mathrm{base}}:=V^{(j)}(V_{j+1})$ on the level-$j$ bonds with both endpoints in
$\Omega_{j+1}$, $V_{j+1}$ being Ba{\l}aban's level-$(j+1)$ variables and $V^{(j)}(V_{j+1})$
Ba{\l}aban's next-level background averaged to level $j$
\cite[(3.10)--(3.13)]{B-CMP119}, and $V_{j}^{\mathrm{base}}:=V_{j}$, the unintegrated level-$j$
variable, on the remaining bonds of $\Gamma_{j}$. The \emph{varied data} $V^{(y)}$ agree with
$V^{\mathrm{base}}$ in the zones $n<j$, and on the top zone they are the bondwise product
\begin{equation*}
V^{(y)}_{j}(b):=V'(y)(b)\,V_{j}^{\mathrm{base}}(b)\qquad\text{for the bonds $b$ between points
of $\Gamma_{j}$.}
\end{equation*}
Then $V^{(0)}=V^{\mathrm{base}}$. We write $\mathcal{V}(y\oplus 0)$ for the top datum of
$V^{(y)}$. For real $y$
(that is, $y_{b}\in\mathfrak{g}$ for all $b\in I$) we put $U^{y}:=U(\mathbf{B}_{j},V^{(y)})$,
and $U_{0}:=U^{0}=U(\mathbf{B}_{j},V^{\mathrm{base}})$.

\emph{Tubes.} For $Y_{R},Y_{I}>0$ put
\begin{equation}\label{8.4:eq-box-variation-b}
\mathcal{T}(Y_{R},Y_{I}):=\bigl\{y\in(\mathfrak{g}^{c})^{I}:\ |y_{b}'|<Y_{R},\
|y_{b}''|<Y_{I}\ \text{for all } b\in I\bigr\},\qquad Y:=Y_{R}+Y_{I}.
\end{equation}
The set $\mathcal{T}(Y_{R},Y_{I})$ is convex and contains $0$, and its real slice
$\mathcal{T}(Y_{R},Y_{I})\cap\mathfrak{g}^{I}$ is nonempty and open in $\mathfrak{g}^{I}$.

\emph{Directional derivatives.} For $b\in I$ and $\zeta\in\mathfrak{g}^{c}$ with
$|\zeta|\le 1$, let $\zeta e_{b}\in(\mathfrak{g}^{c})^{I}$ be the element equal to $\zeta$ at $b$
and to $0$ at the other bonds of $I$. For a holomorphic function $\varphi$ we put
\begin{equation*}
\partial_{b,\zeta}\varphi(y):=\frac{d}{ds}\,\varphi(y+s\zeta e_{b})\Big|_{s=0},
\qquad s\in\mathbb{C}.
\end{equation*}

\emph{Covariant curl.} Let $Z$ be a $\mathfrak{g}^{c}$-valued function on the fine bonds,
extended to reversed bonds by $Z(b^{-1}):=-\operatorname{Ad}_{U_{0}(b)^{-1}}Z(b)$. For a fine
plaquette $p$ with bonds $b_{1},b_{2},b_{3},b_{4}$ as in Definition~\ref{8.4:not-step-zones}, and
with $U_{i}:=U_{0}(b_{i})$, put
\begin{equation}\label{8.4:eq-box-variation-c}
(D_{U_{0}}Z)(p):=\eta^{-1}\bigl[Z(b_{1})+\operatorname{Ad}_{U_{1}}Z(b_{2})
-\operatorname{Ad}_{U_{4}}Z(b_{3})-Z(b_{4})\bigr].
\end{equation}

\emph{Distances.} For fine sites $x,x'$, $\operatorname{dist}_{j}(x,x')$ is $\eta$ times the
graph distance of $x$ and $x'$ in the fine lattice, that is, their distance in level-$j$ units.
The \emph{scaled path distance} is
\begin{equation*}
\rho_{\mathrm{sc}}(x,x'):=\min_{\gamma}\ \sum_{\langle v,v'\rangle\in\gamma}L^{-n(v)},
\end{equation*}
the minimum being taken over the nearest-neighbour fine paths $\gamma$ from $x$ to $x'$ and the
sum over the bonds $\langle v,v'\rangle$ of $\gamma$; each bond is thus counted in units of the
block size of the zone in which it starts. One has
$\rho_{\mathrm{sc}}\ge\operatorname{dist}_{j}$. For a bond $b=\langle v,v+e_{\mu}\rangle$ of
$\mathbb{T}^{(j)}$ and $\rho\in\{\operatorname{dist}_{j},\rho_{\mathrm{sc}}\}$ we put
$\rho(x,b):=\min\bigl(\rho(x,v),\rho(x,v+e_{\mu})\bigr)$, and for sets we take the minimum over
pairs. Finally $d(\cdot,\cdot)$ denotes Ba{\l}aban's distance function of
\cite[(190)]{B-CMP102b}; we do not reproduce its definition.

\emph{Thresholds and constants.} The thresholds of Ba{\l}aban's characteristic functions are
\begin{equation*}
\varepsilon_{n}:=g_{n}A_{0}\,p_{0}(g_{n}),\qquad
\delta_{n}:=g_{n}A_{1}\,p_{0}(g_{n})/A_{0}
\end{equation*}
\cite[after~(3.3)]{B-CMP119}, with $p_{0}(\cdot)$ the degree function; $x_{n}:=g_{n}^{-2}$. We use
Ba{\l}aban's constants $a_{0},a_{1},B_{3}$ of \cite[Theorem~1]{B-CMP102b}, $B_{5}$ of
\cite[Proposition~9]{B-CMP102b}, $\delta_{0}$, and the constant of the $O(1)$ in
\cite[(190)]{B-CMP102b}, which we denote $c_{*}$. The symbol $\delta_{0}$ always denotes this
constant of Ba{\l}aban's, never the threshold $\delta_{n}$ at $n=0$; the threshold is written
$\delta_{n}$ only for $n\ge 1$, and at $n=0$ it is written out as
$g_{0}A_{1}\,p_{0}(g_{0})/A_{0}$.
\end{definition}

\begin{proof}
We prove the three assertions made in Definition~\ref{8.4:not-box-variation}.

\emph{The identity $V^{(0)}=V^{\mathrm{base}}$.} At $y=0$ the extension $\tilde{y}=0\oplus 0$ is
the zero field, so $g_{j}\mathbf{C}\tilde{y}=0$ because $\mathbf{C}$ is linear. Hence, by
\eqref{8.4:eq-box-variation-a},
\begin{equation*}
B(0)=0-\mathbf{h}\tilde{D}(0)=-\mathbf{h}(0)=0,
\end{equation*}
where we used $\tilde{D}(0)=0$ \cite{B-CMP109} and that the operator $\mathbf{h}$ maps $0$ to $0$.
Therefore $V'(0)=1$ on every bond of $\mathbb{T}^{(j)}$, and the bondwise product on the top zone
gives $V^{(0)}_{j}=V_{j}^{\mathrm{base}}$; in the zones $n<j$ the data are unchanged by
definition. Thus $V^{(0)}=V^{\mathrm{base}}$, and in particular
$U^{0}=U(\mathbf{B}_{j},V^{\mathrm{base}})$, so that the two expressions given for $U_{0}$ agree.

\emph{The tube.} For each $b\in I$ the maps $y\mapsto y_{b}'=(y_{b}+y_{b}^{*})/2$ and
$y\mapsto y_{b}''=(y_{b}-y_{b}^{*})/(2i)$ are real-linear, and $|\cdot|$ is a norm. If
$y^{(1)},y^{(2)}\in\mathcal{T}(Y_{R},Y_{I})$ and $0\le\theta\le 1$, then
\begin{equation*}
\bigl|\bigl(\theta y^{(1)}+(1-\theta)y^{(2)}\bigr)_{b}'\bigr|
\le\theta\bigl|(y^{(1)}_{b})'\bigr|+(1-\theta)\bigl|(y^{(2)}_{b})'\bigr|<Y_{R},
\end{equation*}
and in the same way $|(\theta y^{(1)}+(1-\theta)y^{(2)})_{b}''|<Y_{I}$, for every $b\in I$; so
$\mathcal{T}(Y_{R},Y_{I})$ is convex. Since $0'=0''=0$ and $Y_{R},Y_{I}>0$, it contains $0$. For
$y\in\mathfrak{g}^{I}$ each $y_{b}$ is Hermitian, so $y_{b}'=y_{b}$ and $y_{b}''=0$; the
condition $|y_{b}''|<Y_{I}$ then holds because $Y_{I}>0$, and
\begin{equation*}
\mathcal{T}(Y_{R},Y_{I})\cap\mathfrak{g}^{I}
=\bigl\{y\in\mathfrak{g}^{I}:\ |y_{b}|<Y_{R}\ \text{for all } b\in I\bigr\}.
\end{equation*}
This is the product over $b\in I$ of the open balls $\{y_{b}\in\mathfrak{g}:|y_{b}|<Y_{R}\}$ of
$\mathfrak{g}$, hence open in $\mathfrak{g}^{I}$, and it contains $0$, hence is nonempty.

\emph{The comparison $\rho_{\mathrm{sc}}\ge\operatorname{dist}_{j}$.} Every fine site $v$ has
$n(v)\le j$, and $L>1$, so $L^{-n(v)}\ge L^{-j}=\eta$. For a nearest-neighbour fine path
$\gamma$ from $x$ to $x'$ consisting of $\#\gamma$ bonds this gives
\begin{equation*}
\sum_{\langle v,v'\rangle\in\gamma}L^{-n(v)}\ \ge\ \eta\,\#\gamma\ \ge\
\eta\cdot(\text{fine graph distance of $x$ and $x'$})=\operatorname{dist}_{j}(x,x').
\end{equation*}
Taking the minimum over $\gamma$ gives $\rho_{\mathrm{sc}}(x,x')\ge\operatorname{dist}_{j}(x,x')$.
\end{proof}

\begin{theorem}[{Ba{\l}aban's analyticity of the minimiser in the data, with derivative decay;
\cite[Theorem~1, Proposition~9, (190)]{B-CMP102b}}]\label{8.4:thm-minimiser-analyticity}
Let $\mathfrak{B}_{k}=\bigcup_{n=0}^{k}\Lambda_{n}$ be a determining set in the sense of
\cite{B-CMP102b}, with zones $\Lambda_{n}$, on the fine lattice $\mathbb{T}$ of spacing $\eta$. For
data $V$ on $\mathfrak{B}_{k}$, write $U_{k}(V)$ for the representative, in the axial gauge of
\cite{B-CMP102b}, of the minimal gauge orbit of the Wilson action under the constraint that its
block averages on $\mathfrak{B}_{k}$ equal $V$, as in
\cite{B-CMP102b}.
\begin{itemize}
\item[(a)] There are constants $a_{0}$, $a_{1}$, $B_{3}$, depending only on $d$ and $L$, with the
following property. Let $V$ satisfy the regularity condition \cite[(7)]{B-CMP102b} with a
constant $\hat{\varepsilon}_{1}\le a_{1}$, and let $\hat{\varepsilon}_{0}$ satisfy
$B_{3}\hat{\varepsilon}_{1}\le\hat{\varepsilon}_{0}\le a_{0}$. Then the minimal gauge orbit of this
constrained variational problem exists, it is unique in the space \cite[(6)]{B-CMP102b} defined
with the parameter $\hat{\varepsilon}_{0}$, and it has the regularity \cite[(9)--(10)]{B-CMP102b}
on the special cubes of \cite[Theorem~1]{B-CMP102b}. Here $\hat{\varepsilon}_{1}$ and
$\hat{\varepsilon}_{0}$ are the constants written $\varepsilon_{1}$ and $\varepsilon_{0}$ in
\cite[Theorem~1]{B-CMP102b}; the hats keep them apart from the small-field thresholds
$\varepsilon_{n}$ used elsewhere in this section.
\item[(b)] Let $V_{0}$ be data on $\mathfrak{B}_{k}$ to which \textup{(a)} applies. Then the
minimal configuration $U_{k}(V'V_{0})$ extends analytically to $G^{c}$-valued $V'$ on
$\mathfrak{B}_{k}$ sufficiently close to the identity, in the sense of
\cite[Proposition~9]{B-CMP102b}; here $V'V_{0}$ is the bondwise product. Put
$B=\frac{1}{i}\log V'$. The function $U_{k}(V'V_{0})\,U_{k}(V_{0})^{-1}$, transformed to the
Landau gauge, equals $\exp\bigl(i\eta\,\mathcal{H}(B)\bigr)$, where the $1$-form
$\mathcal{H}(B)=(\mathcal{H}_{\mu}(B,x))$ is analytic in $B$. Its functional derivative obeys the
following bounds: for $x\in\Delta(y)$, $y\in\Lambda_{n}$ and $y'\in\Lambda_{n'}$,
\begin{equation}\label{8.4:eq-minimiser-analyticity-a}
\begin{gathered}
\Bigl|\frac{\delta}{\delta B_{\nu}(y')}\mathcal{H}_{\mu}(B,x)\Bigr|
\le c_{*}\,(L^{n}\eta)^{-1}(L^{n'}\eta)^{-d}\,e^{-\delta_{0}d(y,y')/8},\\
\Bigl|\nabla_{x}\frac{\delta}{\delta B_{\nu}(y')}\mathcal{H}_{\mu}(B,x)\Bigr|
\le c_{*}\,(L^{n}\eta)^{-2}(L^{n'}\eta)^{-d}\,e^{-\delta_{0}d(y,y')/8}.
\end{gathered}
\end{equation}
Here $\Delta(y)$ are the localisation sets, $d(y,y')$ is the distance and $\delta_{0}$ the decay
constant of \cite[(190)]{B-CMP102b}, and $c_{*}$ is the constant written $O(1)$ there.
\end{itemize}
\end{theorem}

This is \cite[Theorem~1]{B-CMP102b} together with \cite[Proposition~9 and (190),
pp.~308--309]{B-CMP102b}.

The theorem enters this book through the following identification. In
\cite[p.~256, after (2.12)]{B-CMP119}, the variational problem \cite[(2.12)]{B-CMP119} for the
determining set $\mathbf{B}_{j}$ is referred
to \cite{B-CMP102b}. Accordingly, the minimiser $U(\mathbf{B}_{j},V)$ of
Notation~\ref{8.4:not-step-zones} is the configuration $U_{k}(V)$ above with $k=j$,
$\mathfrak{B}_{k}=\mathbf{B}_{j}$ and $\Lambda_{n}=\Gamma_{n}$. In this identification
\cite{B-CMP119} applies part~\textup{(a)} to the problem \cite[(2.12)]{B-CMP119}, and the bounds
\eqref{8.4:eq-minimiser-analyticity-a} to the response of its minimiser. Both parts are used here
in this form.

\begin{definition}[Hypotheses on the response estimate, the distance and data regularity]
\label{8.4:def-response-hypotheses}
Let the step index $j\ge 1$, the fine lattice $\mathbb{T}$ of spacing $\eta$ with $L^{j}\eta=1$, the
admissible sequence $\mathbb{T}\supseteq\Omega_{1}\supseteq\cdots\supseteq\Omega_{j+1}$, the zones
$Z_{n}$ and the zone index $n(x)$, the determining set $\mathbf{B}_{j}=\bigcup_{n}\Gamma_{n}$ with its
data $V=(V_{n})_{0\le n\le j}$, the averages $M^{n}$, the minimiser $U(\mathbf{B}_{j},V)$ and the bond
sets $E_{1},E_{0},C_{0}$ be as in Definition~\ref{8.4:not-step-zones}. Let the box bond set $I$, the
box variables $y\in(\mathfrak{g}^{c})^{I}$ and their extension $\tilde{y}$ by zero, the multiplier
$B(y)$ with $V'(y)=\exp(iB(y))$, the base data $V^{\mathrm{base}}$, the base minimiser
$U_{0}=U(\mathbf{B}_{j},V^{\mathrm{base}})$, the tubes $\mathcal{T}(Y_{R},Y_{I})$ with
$Y=Y_{R}+Y_{I}$, the directional derivatives $\partial_{b,\zeta}$, the covariant curl $D_{U_{0}}$ of
\eqref{8.4:eq-box-variation-c} and the distances $\operatorname{dist}_{j}$ and
$\rho_{\mathrm{sc}}$ be as in Definition~\ref{8.4:not-box-variation}. Let $\mathcal{H}$, the constants
$c_{*}$, $\delta_{0}$, $B_{5}$, the determining set $\mathfrak{B}_{k}$ with its zones $\Lambda_{n}$,
the sets $\Delta(y)$ and Ba{\l}aban's distance $d(y,y')$ be as in
Theorem~\ref{8.4:thm-minimiser-analyticity}. Throughout, the level $k$ of \cite{B-CMP102b} is $j$,
its determining set $\mathfrak{B}_{k}$ is $\mathbf{B}_{j}$ and its zones $\Lambda_{n}$ are the sets
$\Gamma_{n}$, the variational problem \cite[(2.12)]{B-CMP119} for $\mathbf{B}_{j}$ being the problem
of \cite{B-CMP102b} \cite[p.~256]{B-CMP119}. In \textup{(a1)}--\textup{(a4)} the letters $V_{0}$,
$V'$, $k$ are those of Theorem~\ref{8.4:thm-minimiser-analyticity}: $V_{0}$ denotes data on
$\mathfrak{B}_{k}$ to which \cite[Theorem~1]{B-CMP102b} applies (not the zone-zero component of
$V$), $V'$ a small $G^{c}$-valued multiplier and $k$ the level; in the applications
$V_{0}=V^{\mathrm{base}}$ and $V'=V'(y)$. In \textup{(a1)}--\textup{(b2)} the letters $y,y'$
denote points of the zones $\Lambda_{n}$; in \textup{(c8)}--\textup{(c9)}, in
\eqref{8.4:eq-response-hypotheses-a} and in the proof below, the letter $y$ denotes the
box variables. The following conditions are the hypotheses \textup{(a1)}--\textup{(a7)},
\textup{(b1)}--\textup{(b2)} and \textup{(c1)}--\textup{(c9)} of this definition; a statement that
uses one of them lists it among its hypotheses.

\smallskip
\noindent\emph{Hypotheses on the response estimate.}
\begin{itemize}
\item[(a1)] In the assertion of Theorem~\ref{8.4:thm-minimiser-analyticity} that
$U_{k}(V'V_{0})U_{k}(V_{0})^{-1}$, transformed to the Landau gauge, is $\exp(i\eta\mathcal{H}(B))$,
the base $U_{k}(V_{0})$ is a fixed configuration in a fixed gauge, independent of $B$; and for
$G$-valued $V'$ the configurations $U_{k}(V'V_{0})$ and $\exp(i\eta\mathcal{H}(B))\cdot U_{k}(V_{0})$
(bondwise product) are related by a $G$-valued gauge transformation, which may depend on $B$.
\item[(a2)] $\mathcal{H}(0)=0$: at $V'=1$ the represented object is the identity, whose Landau-gauge
form is the identity.
\item[(a3)] For $y'$ in the top zone $\Lambda_{k}$, where $L^{k}\eta=1$, the functional derivative
$\delta/\delta B_{\nu}(y')$ in \cite[(182)]{B-CMP102b} and in
\eqref{8.4:eq-minimiser-analyticity-a} is the derivative with respect to the single
$\mathfrak{g}^{c}$-valued top-bond variable $B_{\nu}(y')$, and the bound is on the operator norm of
the resulting linear map $\mathfrak{g}^{c}\to\mathfrak{g}^{c}$.
\item[(a4)] The difference $\nabla_{x}$ in the second inequality of
\eqref{8.4:eq-minimiser-analyticity-a} is the covariant difference with respect to the base
background: for each top bond $(y',\nu)$ and each $\zeta\in\mathfrak{g}^{c}$ the $1$-form
$Z:=(\delta\mathcal{H}(B,\cdot)/\delta B_{\nu}(y'))[\zeta]$ (a letter local to this item) satisfies,
for every plaquette $p$ at a site $x\in\Delta(y)$ with $y\in\Lambda_{n}$ (and $L^{k}\eta=1$),
\begin{equation*}
\eta^{2}\,\bigl|(D_{U_{0}}Z)(p)\bigr|
\le 2c_{*}\,\eta^{2}(L^{n}\eta)^{-2}e^{-\delta_{0}d(y,y')/8}\,|\zeta| .
\end{equation*}
\item[(a5)] For every fine site $x$ there is $y_{x}\in\Gamma_{n(x)}$ with $x\in\Delta(y_{x})$.
\item[(a6)] \eqref{8.4:eq-minimiser-analyticity-a} holds at every $B$ of the analyticity domain
\cite[(172)]{B-CMP102b}, and this domain contains $\{B:\sup|B|<a_{9}\}$ for some $a_{9}>0$.
\item[(a7)] The constants $c_{*}$, $\delta_{0}$, $a_{9}$, $B_{5}$ depend neither on $j$, nor on the
admissible sequence $\{\Omega_{n}\}$, nor on the size of $\mathbb{T}$; they may depend on $d$, $L$,
$N$ and the fixed parameters of \cite{B-CMP102b}.
\end{itemize}

\smallskip
\noindent\emph{Hypotheses on the distance.} For a top bond $(y',\nu)$ write $b(y',\nu)$ for the
level-$j$ bond from $y'$ in the direction $\nu$ (the top bond carrying the variable $B_{\nu}(y')$).
\begin{itemize}
\item[(b1)] There are $c_{d}>0$ and $c_{d}'\ge 0$, depending on $d$ and $L$ only, such that
\begin{equation*}
d(y_{x},y')\ge c_{d}\operatorname{dist}_{j}\bigl(x,b(y',\nu)\bigr)-c_{d}'
\end{equation*}
whenever $x\in\Delta(y_{x})$ and $(y',\nu)$ is a top bond.
\item[(b2)] The same holds with $\rho_{\mathrm{sc}}$ in place of $\operatorname{dist}_{j}$.
\end{itemize}

\smallskip
\noindent\emph{Hypotheses on the regularity and locality of the data.}
\begin{itemize}
\item[(c1)] The conclusions of \cite[Theorem~1]{B-CMP102b} and of
Theorem~\ref{8.4:thm-minimiser-analyticity} hold for $(\mathbf{B}_{j},V^{\mathrm{base}})$ when
$(7^{\flat})(\bar\varepsilon)$ of \eqref{8.4:eq-step-zones-a} holds with $\bar\varepsilon\le a_{1}$,
the base configuration $U(\mathbf{B}_{j},V^{\mathrm{base}})$ being equal to $V_{0}$ on $E_{0}$ and,
on $E_{1}$, to the minimiser of \cite[(2.12)]{B-CMP119} with $U=V_{0}$ on $E_{0}$ supplied by
\cite[Theorem~1]{B-CMP102b}, the data at the deep zone-zero plaquettes not entering this minimiser
by the lemma of this chapter on the deep zone-zero data. This is used here as in
\cite{B-CMP119}, which invokes \cite[Theorem~1]{B-CMP102b} for the problem
\cite[(1.12)]{B-CMP119}, whose zone-zero data contain the large-field regions, and which also uses
the estimate \eqref{8.4:eq-minimiser-analyticity-a} \cite{B-CMP119}.
\item[(c2)] For $1\le n\le j-1$ the unintegrated data $V_{n}$ satisfy
$|\partial V_{n}-1|<c_{\Gamma}L^{2}\varepsilon_{n}$
at the level-$n$ plaquettes of $\Gamma_{n}$, and the data are regular on $\mathbf{B}_{j}$ in the sense
of \cite[(7)]{B-CMP102b} (see \eqref{8.4:eq-step-zones-a}) at the zone interfaces. In
\cite[p.~256, sentence before (2.10)]{B-CMP119} it is assumed that $V_{j-1}$ is regular
on $\Gamma_{j-1}$ in the sense $|\partial V_{j-1}-1|<O(L^{2})\varepsilon_{j-1}$; this is part of the
inductive description of the density there, maintained by the constructions of
\cite[\S3]{B-CMP119} and \cite[(1.18)--(1.20)]{B-CMP122a}, and \cite[p.~256]{B-CMP119} states that a
regular configuration on the determining set determines the minimal orbit.
\item[(c3)] Let $V^{(j)}=V^{(j)}(V_{j+1})$ denote the next-level background averaged to level $j$
\cite[(3.10)--(3.13)]{B-CMP119}, as in Definition~\ref{8.4:not-box-variation}. On the support of
$\chi_{j}(\Omega_{j})$ and of the small-field functions of step $j+1$,
$|\partial V_{j}-1|<2\varepsilon_{j}$ at the level-$j$ plaquettes of $\Omega_{j}$, and
\begin{equation*}
\bigl|V_{j}(b)\bigl(V^{(j)}(b)\bigr)^{-1}-1\bigr|<c'\varepsilon_{j}
\end{equation*}
on the level-$j$ bonds $b$ of $\Omega_{j+1}$. The characteristic function $\chi_{k}$ of
\cite{B-CMP109} is defined so that $U_{k}(V)$ lies in the set of regular configurations of
\cite{B-CMP109} with parameter $\varepsilon$, which, by a proposition cited there, gives
$|V(\partial p')-1|<2\varepsilon$ at the plaquettes $p'$ on its support; and, with $V^{(k)}$ the
background of level $k$ defined in the same way, \cite[(3.13)--(3.14)]{B-CMP119} give
$|V_{k}(V^{(k)})^{-1}-1|<O(1)\varepsilon_{k}$ on $\Omega_{k+1}$.
\item[(c4)] $|\partial V_{0}-1|<\varepsilon_{0}$ at every fine plaquette of
$\tilde{\Omega}_{1}\setminus\Omega_{1}$, in particular on $C_{0}$ (since $LMR_{1}\ge 2$). By
\cite{B-CMP119}, the distance between $\Omega_{1}$ and the union of the large-field regions is at
least $2LMR_{1}$, and on $\tilde{\Omega}_{1}$ only small-field characteristic functions occur.
\item[(c5)] The constraint $M^{0}(U)\restriction\Gamma_{0}=V_{0}$ is $U(b)=V_{0}(b)$ for all
$b\in E_{0}$, and the unknowns of \cite[(2.12)]{B-CMP119} are the restriction $U\restriction E_{1}$;
here \cite[(1.12)]{B-CMP119} sets $U=V_{0}$ on $\Omega_{1}^{c}$. If the bonds crossing
$\partial\Omega_{1}$ are assigned to $\Omega_{1}^{c}$ as well, every statement using this hypothesis
holds with $C_{0}$ enlarged by one layer.
\item[(c6)] For $n\ge 1$ and a level-$n$ bond $c=\langle z,z'\rangle$ of $\Gamma_{n}$, $M^{n}(U)(c)$
depends only on $U$ on the fine bonds with both endpoints in $B^{n}(z)\cup B^{n}(z')$, a set contained
in $\Omega_{1}$; the axial gauge conditions concern bonds of $E_{1}$ (\cite[(2.11)]{B-CMP119},
\cite{B-CMP98}).
\item[(c7)] The functional $A^{\eta}$ of \cite[(2.12)]{B-CMP119} is, with a constant $a_{\eta}>0$ and
the sum over all fine plaquettes,
\begin{equation*}
A^{\eta}(U)=a_{\eta}\sum_{p\subset\mathbb{T}}\bigl(1-\operatorname{Re}\operatorname{tr}U(\partial p)\bigr).
\end{equation*}
\item[(c8)] Let $c$ range over the index set of the constraint of \cite{B-CMP109}, with the attached
bond $b_{0}(c)$, the points $c_{-}$, $c_{+}$, the blocks of $c_{-}$ and of $c_{+}$, and the corridor
of $c$ as there. For a $\mathfrak{g}^{c}$-valued function $\Phi$ on bonds, $\mathbf{C}$ is linear with
$(\mathbf{C}\Phi)(b')=\Phi(b')$ for $b'$ not of the form $b_{0}(c)$, and $(\mathbf{C}\Phi)(b_{0}(c))$
is a linear function of $\Phi$ on the bonds of the blocks of $c_{-}$ and $c_{+}$ of norm at most
$c_{C}$ (replacing $c_{C}$ by $\max\{1,c_{C}\}$, we may and do assume $c_{C}\ge 1$). Further, for a
function $\Psi$ indexed by $c$, $(\mathbf{h}\Psi)(b_{0}(c))=\mathbf{h}(c)\Psi(c)$ with
$|\mathbf{h}(c)|\le c_{h}$, and
$\mathbf{h}\Psi=0$ elsewhere. $\tilde{D}$ is holomorphic on $\{X:\sup|X|<b_{D}\}$, with $\tilde{D}(0)=0$,
$d\tilde{D}(0)=0$, and for $\sup|X|\le b_{D}/2$
\begin{equation*}
|\tilde{D}(X)(c)|\le c_{D}\sup|X|^{2},\qquad
|d\tilde{D}_{X}[W](c)|\le c_{D}\sup|X|\,\sup|W| ;
\end{equation*}
$\tilde{D}(X)(c)$ depends on $X$ only on the bonds of the blocks of $c_{-}$ and $c_{+}$ and of the
corridor of $c$. The constants depend on $d$ and $L$ only (\cite{B-CMP119}, \cite{B-CMP109}). The
following consequence of these properties is also part of \textup{(c8)}, used in the form stated:
for $b\in I$, $\zeta\in\mathfrak{g}^{c}$ and $y\in\mathcal{T}(Y_{R},Y_{I})$ with
$g_{j}c_{C}Y\le\min\{b_{D}/2,(2c_{h}c_{D})^{-1}\}$, the function $\partial_{b,\zeta}B(y)$ is
supported on a set $S(b)$ of at most $N_{B}:=(4d+2)^{2}$ top bonds, all within
$\operatorname{dist}_{j}$-distance $4L+2$ of $b$.
\item[(c9)] Every top bond within $\operatorname{dist}_{j}$-distance $4L+2$ of $I$ has both endpoints
in $\Omega_{j+1}$. This holds on the clean event of the Gaussian extraction, on which the
characteristic functions within every box are the small-field ones.
\end{itemize}
Under \textup{(c8)}, for $y\in\mathcal{T}(Y_{R},Y_{I})$, $b\in I$ and $\zeta\in\mathfrak{g}^{c}$,
provided $g_{j}c_{C}Y\le\min\{b_{D}/2,(2c_{h}c_{D})^{-1}\}$,
\begin{equation}\label{8.4:eq-response-hypotheses-a}
\bigl|\partial_{b,\zeta}B(y)(b')\bigr|\le 2g_{j}c_{C}|\zeta|,\qquad
\sup_{b'}|B(y)(b')|\le 2g_{j}c_{C}Y ;
\end{equation}
and under \textup{(c8)} and \textup{(c9)},
$\operatorname{supp}B(y)\subset S(I):=\bigcup_{b\in I}S(b)$, which consists of bonds of the top zone
$\Gamma_{j}$.
\end{definition}

\begin{proof}[Proof of \eqref{8.4:eq-response-hypotheses-a} and of the support inclusion]
Fix $Y_{R},Y_{I}>0$ with $g_{j}c_{C}Y\le\min\{b_{D}/2,(2c_{h}c_{D})^{-1}\}$, a point
$y\in\mathcal{T}(Y_{R},Y_{I})$, a bond $b\in I$ and $\zeta\in\mathfrak{g}^{c}$. Write $\zeta e_{b}$
for the element of $(\mathfrak{g}^{c})^{I}$ equal to $\zeta$ at $b$ and $0$ elsewhere, and put
(local notation for this proof)
\begin{equation*}
X:=g_{j}\mathbf{C}\tilde{y},\qquad W:=g_{j}\mathbf{C}\,\widetilde{\zeta e_{b}},
\end{equation*}
where $\widetilde{\zeta e_{b}}$ is $\zeta e_{b}$ extended by zero.

\emph{Sup bounds for $\mathbf{C}$.} For a $\mathfrak{g}^{c}$-valued $\Phi$ on bonds, (c8) gives
$|(\mathbf{C}\Phi)(b')|=|\Phi(b')|\le\sup|\Phi|$ when $b'$ is not of the form $b_{0}(c)$, and
$|(\mathbf{C}\Phi)(b_{0}(c))|\le c_{C}\sup|\Phi|$; since $c_{C}\ge 1$, we get
$\sup|\mathbf{C}\Phi|\le c_{C}\sup|\Phi|$. For $b''\in I$ we have $y_{b''}=y_{b''}'+iy_{b''}''$, so by
the triangle inequality for the operator norm and the definition of the tube,
\begin{equation*}
|y_{b''}|\le|y_{b''}'|+|y_{b''}''|<Y_{R}+Y_{I}=Y .
\end{equation*}
Hence $\sup|\tilde{y}|<Y$, so $\sup|X|\le g_{j}c_{C}Y\le b_{D}/2$, and likewise
$\sup|W|\le g_{j}c_{C}|\zeta|$.

\emph{The derivative.} By the definition of $B(y)$ in Definition~\ref{8.4:not-box-variation} and the
linearity of $\mathbf{C}$, for complex $s$
\begin{equation*}
B(y+s\zeta e_{b})=X+sW-\mathbf{h}\tilde{D}(X+sW).
\end{equation*}
Since $\sup|X|\le b_{D}/2<b_{D}$ and $\tilde{D}$ is holomorphic on $\{\sup|X|<b_{D}\}$, the right side
is defined and holomorphic for $|s|$ small, and by the linearity of $\mathbf{h}$ and the chain rule
\begin{equation*}
\partial_{b,\zeta}B(y)=W-\mathbf{h}\,d\tilde{D}_{X}[W].
\end{equation*}

\emph{First bound in \eqref{8.4:eq-response-hypotheses-a}.} Let $b'$ be a top bond. If $b'$ is not of
the form $b_{0}(c)$, then $(\mathbf{h}\,d\tilde{D}_{X}[W])(b')=0$ by (c8), so
$|\partial_{b,\zeta}B(y)(b')|=|W(b')|\le g_{j}c_{C}|\zeta|$. If $b'=b_{0}(c)$, then by (c8), applicable
since $\sup|X|\le b_{D}/2$,
\begin{align*}
\bigl|(\mathbf{h}\,d\tilde{D}_{X}[W])(b_{0}(c))\bigr|
&=|\mathbf{h}(c)|\,\bigl|d\tilde{D}_{X}[W](c)\bigr|
\le c_{h}c_{D}\sup|X|\,\sup|W|\\
&\le c_{h}c_{D}\,(g_{j}c_{C}Y)\,(g_{j}c_{C}|\zeta|)
\le\tfrac12\,g_{j}c_{C}|\zeta|,
\end{align*}
the last step by $g_{j}c_{C}Y\le(2c_{h}c_{D})^{-1}$. Together with $|W(b')|\le g_{j}c_{C}|\zeta|$ this
gives $|\partial_{b,\zeta}B(y)(b')|\le\frac32 g_{j}c_{C}|\zeta|\le 2g_{j}c_{C}|\zeta|$.

\emph{Second bound in \eqref{8.4:eq-response-hypotheses-a}.} Here $B(y)=X-\mathbf{h}\tilde{D}(X)$. At
a bond
$b'$ not of the form $b_{0}(c)$, $|B(y)(b')|=|X(b')|\le g_{j}c_{C}Y$. At $b'=b_{0}(c)$, by (c8),
\begin{align*}
|B(y)(b_{0}(c))|
&\le|X(b_{0}(c))|+|\mathbf{h}(c)|\,|\tilde{D}(X)(c)|
\le g_{j}c_{C}Y+c_{h}c_{D}\,(g_{j}c_{C}Y)^{2}\\
&\le g_{j}c_{C}Y\bigl(1+\tfrac12\bigr)\le 2g_{j}c_{C}Y ,
\end{align*}
again by $g_{j}c_{C}Y\le(2c_{h}c_{D})^{-1}$.

\emph{Support.} By linearity of $\mathbf{C}$, $\mathbf{C}0=0$, and $\tilde{D}(0)=0$, so $B(0)=0$. The
tube is convex and contains $0$, so $ty\in\mathcal{T}(Y_{R},Y_{I})$ for $t\in[0,1]$. On the tube the
map $y\mapsto g_{j}\mathbf{C}\tilde{y}$ is linear with values in $\{\sup|X|\le b_{D}/2\}$ by the sup
bounds above, $\tilde{D}$ is holomorphic on $\{\sup|X|<b_{D}\}$ and $\mathbf{h}$ is linear, so $B$ is
holomorphic on $\mathcal{T}(Y_{R},Y_{I})$. The directional derivative is linear in the direction, so
$\partial_{b,y_{b}}$ is defined by the same formula, and by the chain rule
$\frac{d}{dt}B(ty)=\sum_{b\in I}\partial_{b,y_{b}}B(ty)$. Hence
\begin{equation*}
B(y)=B(y)-B(0)=\int_{0}^{1}\frac{d}{dt}B(ty)\,dt
=\int_{0}^{1}\sum_{b\in I}\partial_{b,y_{b}}B(ty)\,dt .
\end{equation*}
By (c8), applied at $ty$ with $\zeta=y_{b}$, each integrand $\partial_{b,y_{b}}B(ty)$ vanishes off
$S(b)$, so $B(y)(b')=0$ for $b'\notin S(I)$. Every bond of $S(I)$ is a top bond within
$\operatorname{dist}_{j}$-distance $4L+2$ of some $b\in I$, hence by (c9) has both endpoints in
$\Omega_{j+1}\subset\Omega_{j}=Z_{j}$.
Since $Z_{j}$ is a union of level-$j$ blocks and a level-$j$ point lies in its own block, the
level-$j$ block of each endpoint lies in $Z_{j}$, so both endpoints belong to $\Gamma_{j}$.
\end{proof}

\begin{definition}[Hypotheses on background regularity, zone thickness and far members]
\label{8.4:def-far-member-hypotheses}
The following hypotheses are used only in the corollaries of this section. Throughout, the step
index $j\ge1$, the fine lattice $\mathbb{T}$ of spacing $\eta$, the level-$n$ lattices
$\mathbb{T}^{(n)}$, the admissible domains $\Omega_{n}$, the zones $Z_{n}=\Omega_{n}\setminus
\Omega_{n+1}$, the zone index $n(\cdot)$ and the bond set $E_{1}$ are those of
Notation~\ref{8.4:not-step-zones}. The box bond set $I$, the tube
$\mathcal{T}(Y_{R},Y_{I})$ with $Y=Y_{R}+Y_{I}$, the base minimiser $U_{0}$, the varied
minimiser $U^{y}$, the covariant curl $D_{U_{0}}$ of~\eqref{8.4:eq-box-variation-c} and the
level-$j$ distance $\operatorname{dist}_{j}$ are those of Notation~\ref{8.4:not-box-variation}.
\begin{enumerate}
\item[\textup{(a)}] Regularity of the base background. For the base data $V^{\mathrm{base}}$ of a
term of the expansion, that is, the data whose minimiser is $U_{0}$,
\begin{equation}\label{8.4:eq-far-member-hypotheses-1}
|U_{0}(\partial p)-1|\le C_{\mathrm{reg}}\,\varepsilon_{n(p)}\,L^{-2n(p)}
\end{equation}
at every fine plaquette $p$ having a bond in $E_{1}$. Here $C_{\mathrm{reg}}$ is a constant of
Ba{\l}aban's construction depending only on $d$, $L$ and $M$. The corresponding statement in
\cite[(3.12)]{B-CMP119} is that the field $U_{k+1}$ satisfies the regularity condition
\begin{equation*}
|\partial U_{k+1}-1|<2B_{3}\,\varepsilon_{k+1}\,(L^{-1}\eta)^{2}
\quad\text{on }\Omega_{k+1},
\end{equation*}
where $\partial U_{k+1}$ stands for the plaquette variables $U_{k+1}(\partial p)$,
$p\subset\Omega_{k+1}$; it is derived there from an earlier theorem cited in that paper and from
the restrictions on the data. For the lower zones compare \cite[(9)--(10)]{B-CMP102b}.
\item[\textup{(b)}] Thickness of the zones. For $1\le n\le j-1$, every fine path from
$\Omega_{n+1}$ to $\Omega_{n}^{c}$ has at least
\begin{equation*}
M\cdot L^{n-j}/\eta=M\cdot L^{n}
\end{equation*}
fine bonds in $Z_{n}$; that is, the zone $Z_{n}$ has thickness at least $M$ in units of
level-$n$ blocks. Compare the lower bound $2MR_{j}$ on the distance between boundaries in
\cite[p.~256]{B-CMP119}, the distance at least $2LMR_{1}$ between $\Omega_{1}$ and the union of
the large-field regions in \cite[\S1]{B-CMP119}, and the separation by one layer of $M$-cubes in
\cite[p.~179]{B-CMP122a}. If for some determining set of Ba{\l}aban's construction only the
thickness $M_{1}$ is available, then the floor condition $M\ge M_{0}(L)$ of the corollaries of
this section, where $M_{0}(L)$ denotes the floor on $M$ fixed in those corollaries, is imposed on
$M_{1}$ in place of $M$.
\item[\textup{(c)}] The frozen-coupling member. Put
\begin{equation*}
\xi(x):=g_{j}(x)^{-2}-g_{j}^{-2},
\end{equation*}
where $g_{j}(x)$ is the position-dependent coupling (Definition~\ref{1.5:def-domains}). The
letter $\xi$ denotes this coupling difference throughout; $\Delta(y)$ denotes the localisation
sets of Definition~\ref{8.4:def-response-hypotheses}. The
\emph{frozen-coupling member} is
\begin{equation}\label{8.4:eq-far-member-hypotheses-2}
\mathcal{L}(y):=-\sum_{p}\xi(x_{p})\bigl(1-\operatorname{Re}\operatorname{tr}
U^{y}(\partial p)\bigr),
\end{equation}
the sum running over the fine plaquettes $p$, where $x_{p}$ denotes the site of $p$ at which the
position-dependent coupling is read. The member $\mathcal{L}$ is the position-dependent term of
\cite[(2.23)--(2.24)]{B-CMP119}. As in \cite[\S2]{B-CMP119}, $\xi(x)=0$ for $x$ in the set
$\Lambda_{j}^{\sim -1}$ of Ba{\l}aban's construction, on which $g_{j}(x)=g_{j}$. The hypothesis
is that
\begin{equation*}
|\xi(x_{p})|\le C_{\Delta}\,\bigl(1+j-n(p)\bigr)
\end{equation*}
for every fine plaquette $p$. This form rests on the one-sided bound on the running inverse
couplings,
\begin{equation*}
x_{n}\le x_{j}+\lambda'(j-n)+c_{\lambda}\qquad\text{for } n\le j,
\end{equation*}
with constants $\lambda'$ and $c_{\lambda}$.
\item[\textup{(d)}] The clean collar. Let $W\ge 4L+3$ denote the width of the clean collar in
level-$j$ units; it is kept as a symbolic constant, equal to $100M\check{R}_{j}$ up to an
additive correction, where $\check{R}_{j}$ is a radius with $\check{R}_{j}\ge R_{j}$. The
hypothesis is that
\begin{equation*}
\operatorname{dist}_{j}(p,I)\ge W
\quad\text{and}\quad
\operatorname{dist}_{j}(\tilde{X},I)\ge W
\end{equation*}
for every fine plaquette $p$ with $\xi(x_{p})\neq0$ and for the reading domain $\tilde{X}$
of every far member in the sense of~\textup{(e)}. It expresses the geometry of the event on
which, near the box, all characteristic functions are the small-field ones; it is part of the
clean-collar geometry of the Gaussian extraction.
\item[\textup{(e)}] Far members. A \emph{far member} is a tuple
$(f,\tilde{X},\rho^{(1)},\rho^{(2)},M_{f})$ with the following data:
\begin{itemize}
\item $\tilde{X}$, its \emph{reading domain}, is a finite set of fine bonds with
$\operatorname{dist}_{j}(\tilde{X},I)\ge W$;
\item $\rho^{(1)}$ and $\rho^{(2)}$ satisfy $0<\rho^{(1)}\le1$ and $0<\rho^{(2)}\le1$;
\item $M_{f}\ge0$.
\end{itemize}
On $\mathfrak{g}^{c}$-valued functions $Z$ on $\tilde{X}$, define the norm
\begin{equation}\label{8.4:eq-far-member-hypotheses-a}
\begin{split}
\|Z\|_{\tilde X}:=\max\Bigl\{\;
&\sup_{b\in\tilde{X}}\frac{\eta\,|Z(b)|\,L^{n(b)}}{\rho^{(1)}}\,,\\
&\sup_{p:\ \text{bonds of }p\subset\tilde{X}}
\frac{\eta^{2}\,|(D_{U_{0}}Z)(p)|\,L^{2n(p)}}{\rho^{(2)}}\Bigr\}.
\end{split}
\end{equation}
The tuple is required to satisfy the following conditions.
\begin{itemize}
\item[\textup{(m1)}] The map $Z\mapsto f\bigl(\exp(i\eta Z)\cdot U_{0}|_{\tilde X}\bigr)$, with
the bondwise product, is defined and holomorphic on $\{\|Z\|_{\tilde X}<1\}$ and is bounded
there by $M_{f}$.
\item[\textup{(m2)}] $f$ is invariant under $G$-valued gauge transformations.
\item[\textup{(m3)}] For real $y\in\mathcal{T}(Y_{R},Y_{I})$, that is, with $y_{b}\in\mathfrak{g}$
for all $b\in I$, the value of the member in the exponent of the step is
$f(U^{y}|_{\tilde X})$.
\end{itemize}
This schema is meant for the terms of the step exponent in
\cite[(2.27)(i)--(iii), (2.31), (2.34)--(2.37)]{B-CMP119} lying at level-$j$ distance at least
$W$ from the box; that these terms have the form \textup{(m1)}--\textup{(m3)} is assumed here,
not proved.
\item[\textup{(f)}] Counting of far members. For a far member $f$ write $\tilde{X}_{f}$,
$\rho^{(1)}_{f}$, $\rho^{(2)}_{f}$, $M_{f}$ for its data. For each far member put
\begin{equation*}
w_{f}:=g_{j}\,M_{f}\max\bigl\{1/\rho^{(1)}_{f},\,1/\rho^{(2)}_{f}\bigr\},
\end{equation*}
and call it the \emph{weight} of the member. For each far member let $x_{f}\in
V(\tilde{X}_{f})$ be its \emph{anchor site}, the site fixed for it in the corollaries of this
section.
Here $V(\tilde{X}_{f})$ is the set of endpoints
of the bonds of $\tilde{X}_{f}$. The hypothesis is that there are constants
$K_{\mathrm{cnt}}\ge0$ and $q\ge0$ such that, for every level-$j$ cell, that is, every level-$j$
block $B^{j}(z)$, and every zone index $n$,
\begin{equation}\label{8.4:eq-far-member-hypotheses-3}
\sum\bigl\{w_{f}:\ x_{f}\in B^{j}(z)\cap Z_{n}\bigr\}
\le K_{\mathrm{cnt}}\,(1+j-n)^{q}\,L^{4(j-n)}.
\end{equation}
Compare the counting in \cite[(2.31)]{B-CMP119}.
\item[\textup{(g)}] Growth of the radii and of the window. Ba{\l}aban's $R_{k}$
(\cite[\S3, the sentences following (3.12)]{B-CMP119}) is a function of $g_{k}$ such that,
for every fixed exponent $a$, there is a coupling ceiling $\gamma(a)$ with
\begin{equation*}
e^{-R_{k}}\,p_{0}(g_{k})^{a}\le1
\quad\text{for } g_{k}\le\gamma(a),
\end{equation*}
where $p_{0}(\cdot)$ is the degree function.
Since $\check{R}_{j}\ge R_{j}$ and $p_{0}(g_{j})^{a}\ge0$, it follows that
\begin{equation*}
e^{-\check{R}_{j}}\,p_{0}(g_{j})^{a}\le e^{-R_{j}}\,p_{0}(g_{j})^{a}\le1
\quad\text{for } g_{j}\le\gamma(a).
\end{equation*}
Moreover, the tube size $Y=Y_{R}+Y_{I}$ satisfies
\begin{equation*}
Y\le c_{Y}\,p_{0}(g_{j})^{m_{Y}}.
\end{equation*}
\end{enumerate}
\end{definition}

\begin{definition}[The decay weight and the response size]\label{8.4:not-decay-weight}
Let $I$, $Y$ and the fine lattice $\mathbb{T}$ be as in Notation~\ref{8.4:not-box-variation}. Let
$c_{d}>0$, $c_{d}'\ge 0$, $c_{*}$ and $\delta_{0}>0$ be the constants of
Definition~\ref{8.4:def-response-hypotheses}. Let $c_{C}$ be the bound, fixed in that definition,
on the norm of Ba{\l}aban's constraint-solving operator $\mathbf{C}$ of
Notation~\ref{8.4:not-box-variation}, and let
$N_{B}=(4d+2)^{2}$, with $d=4$ the dimension. Put $c_{1}:=c_{d}(4L+2)+c_{d}'$.

Definition~\ref{8.4:def-response-hypotheses} compares Ba{\l}aban's distance $d(\cdot,\cdot)$
of \cite[(190)]{B-CMP102b}
with a distance $\rho$ between fine sites and bonds, and the choice of $\rho$ depends on which
comparison is assumed:
\begin{itemize}
\item if the comparison is assumed with the scaled path distance, put
  $\rho:=\rho_{\mathrm{sc}}$;
\item if it is assumed with the level-$j$ distance only, put $\rho:=\operatorname{dist}_{j}$.
\end{itemize}

For a fine site $x$ and a bond $b\in I$, the \emph{decay weight} is
\begin{equation}\label{8.4:eq-decay-weight-a}
E(x,b):=\min\Bigl\{1,\ \exp\Bigl(-\frac{\delta_{0}}{8}\bigl(c_{d}\,\rho(x,b)-c_{1}\bigr)\Bigr)\Bigr\}.
\end{equation}
In this section, the \emph{response constant} $K$ and the \emph{response size}
$\Theta(x)$ at a fine site $x$ are
\begin{equation}\label{8.4:eq-decay-weight-b}
K:=4c_{*}c_{C}N_{B},\qquad \Theta(x):=g_{j}\,Y\sum_{b\in I}E(x,b).
\end{equation}
The factor $4$ in $K$ is chosen so that the response bound below holds with this
single constant.

Since $\rho_{\mathrm{sc}}\ge\operatorname{dist}_{j}$ for the two distances of
Notation~\ref{8.4:not-box-variation}, the weight~\eqref{8.4:eq-decay-weight-a}
formed with $\rho=\rho_{\mathrm{sc}}$ is pointwise at most the weight formed with
$\rho=\operatorname{dist}_{j}$. Consequently, every bound stated below in terms of $E$ holds
a fortiori when $\operatorname{dist}_{j}$ is used. Here the bounds meant are upper bounds whose
right-hand sides are nondecreasing functions of the values of $E$, in particular $\Theta$.
\end{definition}

\begin{proof}
Fix a fine site $x$ and a bond $b\in I$. Write $E^{\mathrm{sc}}$ for the weight formed with
$\rho_{\mathrm{sc}}$ and $E^{\mathrm{dist}}$ for the weight formed with $\operatorname{dist}_{j}$.

Since $c_{d}>0$ and $\rho_{\mathrm{sc}}(x,b)\ge\operatorname{dist}_{j}(x,b)$, we have
\begin{equation*}
c_{d}\,\rho_{\mathrm{sc}}(x,b)-c_{1}\ \ge\ c_{d}\operatorname{dist}_{j}(x,b)-c_{1}.
\end{equation*}
Multiplying by $-\delta_{0}/8<0$ reverses this inequality. Since the exponential is
increasing, it follows that
\begin{equation*}
\exp\Bigl(-\frac{\delta_{0}}{8}\bigl(c_{d}\rho_{\mathrm{sc}}(x,b)-c_{1}\bigr)\Bigr)
\le\exp\Bigl(-\frac{\delta_{0}}{8}\bigl(c_{d}\operatorname{dist}_{j}(x,b)-c_{1}\bigr)\Bigr).
\end{equation*}
The map $t\mapsto\min\{1,t\}$ is nondecreasing, so $E^{\mathrm{sc}}(x,b)\le E^{\mathrm{dist}}(x,b)$.

Finally, let $F$ be a right-hand side that is a nondecreasing function of the values of $E$.
By the pointwise inequality just proved, replacing $E^{\mathrm{sc}}$ by $E^{\mathrm{dist}}$ does
not decrease $F$. Hence an upper bound by $F$ formed with $E^{\mathrm{sc}}$ implies the same
upper bound formed with $E^{\mathrm{dist}}$. In particular this applies to $\Theta(x)$, which is
nondecreasing in the values of $E$ since $g_{j}Y\ge 0$.
\end{proof}

\begin{lemma}[Three elementary facts on holomorphic functions and matrix exponentials]
\label{8.4:lem-elementary-facts}
Let $I$, the tubes $\mathcal{T}(Y_{R},Y_{I})\subset(\mathfrak{g}^{c})^{I}$ and the directional
derivatives $\partial_{b,\zeta}$ be as in Notation~\ref{8.4:not-box-variation}
and~\eqref{8.4:eq-box-variation-b}. In items (a) and (b), for $\xi\in\mathfrak{g}^{c}$ we write
$\xi=\xi'+i\xi''$ with $\xi',\xi''\in\mathfrak{g}$, namely $\xi'=\frac12(\xi+\xi^{*})$ and
$\xi''=\frac1{2i}(\xi-\xi^{*})$; since $|\xi^{*}|=|\xi|$, this gives $|\xi'|\le|\xi|$ and
$|\xi''|\le|\xi|$. In item (c) a prime denotes the derivative $d/dt$. We write $|\cdot|$ for the
operator norm of complex $N\times N$ matrices,
$1$ for the identity matrix, and $\operatorname{Tr}$ for the unnormalised trace. The letters
$s$, $g$, $Y$, $A$, $r$, $t$, $\rho$, $\tau$, $u$, $\sigma$, $\xi$, $\psi$, $\Phi$,
the multi-index $\alpha$, the basis
$T_{a}$, the coordinates $z_{b,a}$ and all summation indices are local to this lemma and its proof.
\begin{itemize}
\item[(a)] \emph{Identity principle on the real slice.} Let
$\mathcal{T}=\mathcal{T}(Y_{R},Y_{I})$ be a tube as in~\eqref{8.4:eq-box-variation-b}. Let
$\varphi$ be holomorphic on $\mathcal{T}$ with $\varphi=0$ on
$\mathcal{T}\cap\mathfrak{g}^{I}$. Then
\begin{equation}\label{8.4:eq-elementary-facts-a}
\varphi(y)=0\qquad\text{for all } y\in\mathcal{T}.
\end{equation}
\item[(b)] \emph{Cauchy estimate.} Let $Y_{R}'<Y_{R}$ and $Y_{I}'<Y_{I}$, and put
$r:=\min(Y_{R}-Y_{R}',\,Y_{I}-Y_{I}')$. Let $\varphi$ be holomorphic on
$\mathcal{T}(Y_{R},Y_{I})$ with $|\varphi|\le A$ there. Then for every
$y\in\mathcal{T}(Y_{R}',Y_{I}')$, every $b\in I$ and every $\zeta\in\mathfrak{g}^{c}$ with
$|\zeta|\le1$,
\begin{equation}\label{8.4:eq-elementary-facts-b}
|\partial_{b,\zeta}\varphi(y)|\le \frac{A}{r}.
\end{equation}
\item[(c)] \emph{Exponential estimates.} Let $X_{1},\dots,X_{4}$ be complex $N\times N$ matrices
with $s:=\sum_{i}|X_{i}|\le1$, and put $g:=e^{X_{1}}e^{X_{2}}e^{X_{3}}e^{X_{4}}$. If
$|g-1|\le\frac14$, let $Y$ be the principal logarithm of $g$,
\begin{equation}\label{8.4:eq-elementary-facts-1}
Y=\log g:=\sum_{m\ge1}\frac{(-1)^{m+1}}{m}(g-1)^{m}.
\end{equation}
Then
\begin{equation}\label{8.4:eq-elementary-facts-c}
\begin{gathered}
\Bigl|g-1-\sum_{i=1}^{4}X_{i}\Bigr|\le 2s^{2},
\qquad
\Bigl|g'-\sum_{i=1}^{4}X_{i}'\Bigr|\le 3s\sum_{i=1}^{4}|X_{i}'|,\\
|Y|\le\tfrac43\,|g-1|,\qquad |Y'|\le\tfrac43\,|g'|,\qquad \operatorname{Tr}Y=0,
\end{gathered}
\end{equation}
under the following conditions: the first bound as stated; the second when
$X_{i}=X_{i}(t)$ are $C^{1}$ functions of a real variable $t$ with
derivatives $X_{i}'$, at every $t$ at which $s\le1$; the third when $|g-1|\le\frac14$; the
fourth when $g$ is $C^{1}$ in $t$, at every $t$ at which $|g(t)-1|\le\frac14$; the fifth when
$|g-1|\le\frac14$, $\det g=1$ and moreover $\tfrac43N|g-1|<2\pi$. The last condition holds
automatically when $N\le18$, since then $\tfrac43N|g-1|\le\frac N3<2\pi$.
\end{itemize}
\end{lemma}

\begin{proof}
\emph{Item (a).} Choose a basis $T_{1},\dots,T_{N^{2}-1}$ of the real vector space
$\mathfrak{g}$, which has real dimension $N^{2}-1$. It is also a basis of
$\mathfrak{g}^{c}=\mathfrak{g}+i\mathfrak{g}$ over $\mathbb{C}$. Writing
$y_{b}=\sum_{a}z_{b,a}T_{a}$ identifies $(\mathfrak{g}^{c})^{I}$ with
$\mathbb{C}^{|I|(N^{2}-1)}$, with coordinates $z_{b,a}$, $b\in I$, $1\le a\le N^{2}-1$. Under this
identification $y_{b}'$ and $y_{b}''$ have coordinates
$\operatorname{Re}z_{b,a}$ and $\operatorname{Im}z_{b,a}$. Hence the real slice
$\mathfrak{g}^{I}$ is the set of real points, and holomorphy in $y$ is holomorphy in the
coordinates $z$. For $\zeta\in\mathfrak{g}^{c}$ and $b\in I$ we write $\zeta e_{b}$ for the
element of $(\mathfrak{g}^{c})^{I}$ with entry $\zeta$ at $b$ and entry $0$ at every other bond;
thus $T_{a}e_{b}$ is the unit vector of the coordinate $z_{b,a}$.

The set $\mathcal{T}$ is open, because it is defined by strict inequalities between continuous
functions. It is convex, because $y\mapsto y_{b}'$ and $y\mapsto y_{b}''$ are real-linear and the
norm is convex. By Notation~\ref{8.4:not-box-variation}, its real slice
$\mathcal{T}\cap\mathfrak{g}^{I}$ is nonempty and open in $\mathfrak{g}^{I}$.

Let $y_{0}\in\mathcal{T}\cap\mathfrak{g}^{I}$, and let $\psi$ be holomorphic on $\mathcal{T}$. For
real $\sigma$ close to $0$, the point $y_{0}+\sigma T_{a}e_{b}$ stays in the real slice, since
$T_{a}\in\mathfrak{g}$ and the slice is open in $\mathfrak{g}^{I}$. Since the
complex derivative may be computed along real increments,
\begin{equation*}
\frac{\partial\psi}{\partial z_{b,a}}(y_{0})
=\lim_{\sigma\to0,\,\sigma\in\mathbb{R}}\frac{\psi(y_{0}+\sigma T_{a}e_{b})-\psi(y_{0})}{\sigma}
=\frac{\partial}{\partial(\operatorname{Re}z_{b,a})}
\bigl(\psi|_{\mathcal{T}\cap\mathfrak{g}^{I}}\bigr)(y_{0}).
\end{equation*}
Apply this successively to $\psi=\varphi$ and to its holomorphic derivatives. We obtain, at
every point $y_{0}$ of the real slice and for every multi-index, that the holomorphic derivative
$\partial^{\alpha}\varphi(y_{0})$ equals the corresponding real partial derivative of the
restriction of $\varphi$ to the real slice. That restriction vanishes identically on an open
subset of $\mathfrak{g}^{I}$ containing $y_{0}$, so $\partial^{\alpha}\varphi(y_{0})=0$ for
all $\alpha$.

Since $\mathcal{T}$ is open, $\varphi$ equals its Taylor series at $y_{0}$ on a polydisc centred at
$y_{0}$ and contained in $\mathcal{T}$. All coefficients of that series vanish, so $\varphi=0$ near
$y_{0}$. Finally, $\mathcal{T}$ is convex, hence connected. By the identity theorem for
holomorphic functions of several variables, a holomorphic function on a connected open set that
vanishes on a nonempty open subset vanishes identically. This gives~\eqref{8.4:eq-elementary-facts-a}.

\emph{Item (b).} As in item (a), $\zeta e_{b}$ denotes the element of $(\mathfrak{g}^{c})^{I}$
with entry $\zeta$ at $b$ and entry $0$ at every other bond.

Let $s\in\mathbb{C}$ with $|s|<r$. By the inequalities $|\xi'|\le|\xi|$ and $|\xi''|\le|\xi|$
fixed in the statement, applied to $\xi=s\zeta$,
\begin{equation*}
|(s\zeta)'|\le|s||\zeta|\le|s|,\qquad |(s\zeta)''|\le|s||\zeta|\le|s|.
\end{equation*}
The point $y+s\zeta e_{b}$ differs from $y$ only at $b$. There
\begin{equation*}
|(y_{b}+s\zeta)'|\le|y_{b}'|+|s|<Y_{R}'+(Y_{R}-Y_{R}')=Y_{R},
\end{equation*}
and in the same way $|(y_{b}+s\zeta)''|<Y_{I}$. Hence $y+s\zeta e_{b}\in\mathcal{T}(Y_{R},Y_{I})$.

The function $\Phi(s):=\varphi(y+s\zeta e_{b})$ is therefore holomorphic on the disc $|s|<r$ and
bounded there by $A$. By definition $\Phi'(0)=\partial_{b,\zeta}\varphi(y)$. For $0<\rho<r$, the
one-variable Cauchy formula gives
\begin{equation*}
\Phi'(0)=\frac1{2\pi i}\oint_{|s|=\rho}\frac{\Phi(s)}{s^{2}}\,ds,
\qquad\text{hence}\qquad
|\Phi'(0)|\le\frac{A}{\rho}.
\end{equation*}
Letting $\rho\to r$ gives~\eqref{8.4:eq-elementary-facts-b}.

\emph{Item (c), the product.} Expand every exponential, $e^{X_{i}}=\sum_{n\ge0}X_{i}^{n}/n!$.
This gives the absolutely convergent expansion
\begin{equation*}
g=\sum_{n_{1},\dots,n_{4}\ge0}
\frac{X_{1}^{n_{1}}X_{2}^{n_{2}}X_{3}^{n_{3}}X_{4}^{n_{4}}}{n_{1}!\,n_{2}!\,n_{3}!\,n_{4}!}.
\end{equation*}
The term of total degree $0$ is $1$, and the terms of total degree $1$ are $X_{1},\dots,X_{4}$.
Hence $g-1-\sum_{i}X_{i}$ is the sum of the terms of total degree at least $2$. By
submultiplicativity, each such term is bounded in norm by the corresponding term of the same
expansion with every $X_{i}$ replaced by $|X_{i}|$. That expansion sums to
$e^{|X_{1}|}\cdots e^{|X_{4}|}=e^{s}$, with degree-$0$ part $1$ and degree-$1$ part $s$.

Since $k!\ge2\,(k-2)!$ for $k\ge2$, and $s\le1$, we obtain
\begin{align*}
\Bigl|g-1-\sum_{i}X_{i}\Bigr|&\le e^{s}-1-s=\sum_{k\ge2}\frac{s^{k}}{k!}\\
&\le\frac{s^{2}}{2}\sum_{j\ge0}\frac{s^{j}}{j!}=\frac{s^{2}e^{s}}{2}\le\frac{e}{2}\,s^{2}\le2s^{2}.
\end{align*}

\emph{Item (c), the derivative of the product.} The derivative of the exponential of a $C^{1}$
matrix function is
\begin{equation*}
\frac{d}{dt}e^{X}=\int_{0}^{1}e^{\tau X}X'e^{(1-\tau)X}\,d\tau .
\end{equation*}
By the product rule,
\begin{equation*}
g'=\sum_{i=1}^{4}e^{X_{1}}\cdots e^{X_{i-1}}\Bigl(\frac{d}{dt}e^{X_{i}}\Bigr)
e^{X_{i+1}}\cdots e^{X_{4}}.
\end{equation*}
Expand
$e^{X_{j}}$ for $j\neq i$, and $e^{\tau X_{i}}$, $e^{(1-\tau)X_{i}}$, in power series. Then the
$i$-th summand becomes a sum of terms, each a product of powers of the $X$'s multiplying the one
factor $X_{i}'$. Its term of degree $0$ in the $X$'s is $\int_{0}^{1}X_{i}'\,d\tau=X_{i}'$.
Therefore $g'-\sum_{i}X_{i}'$ is the sum of the terms of degree at least $1$ in the $X$'s.

Replace every matrix by its norm. The scalar majorants of all terms of the $i$-th summand then sum
to
\begin{equation*}
|X_{i}'|\prod_{j\neq i}e^{|X_{j}|}\int_{0}^{1}e^{\tau|X_{i}|}e^{(1-\tau)|X_{i}|}\,d\tau
=|X_{i}'|\,e^{s}.
\end{equation*}
The degree-$0$ term contributes exactly $|X_{i}'|$. Hence the terms of degree at least $1$
contribute at most $|X_{i}'|(e^{s}-1)$. Summing over $i$, and using, for $s\le1$,
\begin{equation*}
e^{s}-1=\sum_{k\ge1}\frac{s^{k}}{k!}\le s\,e^{s}\le e\,s\le3s,
\end{equation*}
we obtain
\begin{equation*}
\Bigl|g'-\sum_{i}X_{i}'\Bigr|\le\Bigl(\sum_{i}|X_{i}'|\Bigr)(e^{s}-1)\le3s\sum_{i}|X_{i}'|.
\end{equation*}
This proves the first two bounds of~\eqref{8.4:eq-elementary-facts-c}.

\emph{Item (c), the logarithm.} For $|g-1|<1$ the series~\eqref{8.4:eq-elementary-facts-1}
converges absolutely. If $|g-1|\le\frac14$, then
$1-|g-1|\ge\frac34$, and therefore
\begin{equation*}
|Y|\le\sum_{m\ge1}|g-1|^{m}=\frac{|g-1|}{1-|g-1|}\le\frac43\,|g-1|,
\end{equation*}
which is the third bound of~\eqref{8.4:eq-elementary-facts-c}.

Now let $g$ be $C^{1}$ in $t$, and let $t_{0}$ satisfy $|g(t_{0})-1|\le\frac14$. By continuity
there is a compact neighbourhood of $t_{0}$ on which $|g-1|\le\frac12$. On it,
\begin{equation*}
\frac{d}{dt}(g-1)^{m}=\sum_{j=0}^{m-1}(g-1)^{j}g'(g-1)^{m-1-j},
\end{equation*}
which has norm at most $m|g-1|^{m-1}|g'|$. The termwise differentiated series is therefore
dominated by $\sum_{m}2^{1-m}\sup|g'|$, so it converges uniformly and may be differentiated
termwise. At $t_{0}$ this gives
\begin{equation*}
|Y'|\le\sum_{m\ge1}|g-1|^{m-1}|g'|=\frac{|g'|}{1-|g-1|}\le\frac43\,|g'|,
\end{equation*}
which is the fourth bound of~\eqref{8.4:eq-elementary-facts-c}.

Finally, let $\det g=1$. The series~\eqref{8.4:eq-elementary-facts-1} is a power series in the
single matrix $g-1$. The scalar identity $\exp\bigl(\sum_{m\ge1}(-1)^{m+1}u^{m}/m\bigr)=1+u$ for
$u\in\mathbb{C}$, $|u|<1$, therefore gives $e^{Y}=g$. Hence
\begin{equation*}
\det g=\det e^{Y}=e^{\operatorname{Tr}Y}=1 ,
\end{equation*}
so $\operatorname{Tr}Y\in2\pi i\mathbb{Z}$. The trace of an $N\times N$ matrix is the sum of its
$N$ eigenvalues, each of modulus at most $|Y|$; with the third bound and the hypothesis
$\tfrac43N|g-1|<2\pi$ this gives
\begin{equation*}
|\operatorname{Tr}Y|\le N|Y|\le\tfrac43N|g-1|<2\pi .
\end{equation*}
The only element of $2\pi i\mathbb{Z}$ of modulus less than $2\pi$ is $0$, so
$\operatorname{Tr}Y=0$. This is the last assertion
of~\eqref{8.4:eq-elementary-facts-c}.
\end{proof}

\begin{lemma}[Deep zone-0 data are invisible to the minimiser]\label{8.4:lem-zone-zero-invisible}
We use the step index $j$, the fine lattice $\mathbb{T}$ of spacing $\eta$, the domains
$\Omega_{n}$ with zones $Z_{n}$, the determining set $\mathbf{B}_{j}=\bigcup_{n}\Gamma_{n}$,
the block averages $M^{n}$ with their restriction $M_{\mathbf{B}_{j}}$ to the determining set,
and the bond sets $E_{1}$, $E_{0}$, $C_{0}$ of Notation~\ref{8.4:not-step-zones}. Thus
$E_{1}$ is the set of fine bonds with at least one endpoint in $\Omega_{1}$, $E_{0}$ is the set of
fine bonds with both endpoints in $\Omega_{1}^{c}$, and $C_{0}\subset E_{0}$ is the set of bonds
of $E_{0}$ that lie in some fine plaquette containing a bond of $E_{1}$. Assume the following.
\begin{enumerate}
\item[(i)] (The zone-0 constraint.) The constraint $M^{0}(U)\restriction\Gamma_{0}=V_{0}$ is the
condition $U(b)=V_{0}(b)$ for all $b\in E_{0}$, and the unknowns of the variational problem of
\cite[(2.12)]{B-CMP119} are the values $U\restriction E_{1}$; this is the form of the constraint of
\cite[(1.12)]{B-CMP119} used here.
\item[(ii)] (Locality of the averaging constraints.) For $n\ge1$ and a level-$n$ bond
$c=\langle z,z'\rangle$ of $\Gamma_{n}$, the value $M^{n}(U)(c)$ depends only on the values of
$U$ on fine bonds with both endpoints in $B^{n}(z)\cup B^{n}(z')$, and
$B^{n}(z)\cup B^{n}(z')\subset\Omega_{1}$; the axial gauge conditions concern bonds of $E_{1}$
only. This is the averaging operation of \cite[(2.11)]{B-CMP119} and \cite{B-CMP98}.
\item[(iii)] (The Wilson form of the functional.) The functional of \cite[(2.12)]{B-CMP119} is
\begin{equation*}
A(U)=a_{\eta}\sum_{p\subset\mathbb{T}}\bigl(1-\operatorname{Re}\operatorname{tr}U(\partial p)\bigr),
\end{equation*}
with a constant $a_{\eta}>0$, the sum running over all fine plaquettes of $\mathbb{T}$; in this
lemma and its proof the letter $A$ denotes this functional.
\end{enumerate}
Let $V=(V_{0},\dots,V_{j})$ and $V^{\sharp}=(V^{\sharp}_{0},\dots,V^{\sharp}_{j})$ be data on
$\mathbf{B}_{j}$ with $V_{n}=V^{\sharp}_{n}$ for $1\le n\le j$ and $V_{0}=V^{\sharp}_{0}$ on
$C_{0}$. For a configuration $u$ on $E_{1}$ let $U_{V}(u)$ be the configuration equal to $u$ on
$E_{1}$ and to $V_{0}$ on $E_{0}$, and $U_{V^{\sharp}}(u)$ the configuration equal to $u$ on
$E_{1}$ and to $V^{\sharp}_{0}$ on $E_{0}$. Then:
\begin{enumerate}
\item[(a)] $A(U_{V}(u))-A(U_{V^{\sharp}}(u))=\kappa(V,V^{\sharp})$ does not depend on $u$;
\item[(b)] $M_{\mathbf{B}_{j}}(U_{V}(u))=V$ if and only if
$M_{\mathbf{B}_{j}}(U_{V^{\sharp}}(u))=V^{\sharp}$;
\item[(c)] hence $u$ is a constrained critical point (respectively a constrained minimiser,
respectively the axially gauged minimiser) for the data $V$ if and only if it is one for the data
$V^{\sharp}$, and the second differentials of the two constrained functionals along variations of
$u$ coincide.
\end{enumerate}
\end{lemma}

\begin{proof}
Every fine bond either has at least one endpoint in $\Omega_{1}$ or has both endpoints in
$\Omega_{1}^{c}$, and not both; so $E_{1}$ and $E_{0}$ partition the fine bonds, and $U_{V}(u)$,
$U_{V^{\sharp}}(u)$ are well-defined configurations on $\mathbb{T}$.

(a) By hypothesis (iii), $A(U)=a_{\eta}\sum_{p}(1-\operatorname{Re}\operatorname{tr}U(\partial p))$.
Split the fine plaquettes into the two classes
\begin{equation*}
\mathcal{P}_{0}:=\{p:\ \text{all bonds of }p\text{ lie in }E_{0}\},\qquad
\mathcal{P}_{1}:=\{p:\ \text{some bond of }p\text{ lies in }E_{1}\};
\end{equation*}
by the partition just noted, every plaquette lies in exactly one of them. Let $p\in\mathcal{P}_{1}$.
Every bond of $p$ that is not in $E_{1}$ lies in $E_{0}$, and since $p$ contains a bond of $E_{1}$,
such a bond lies in $C_{0}$ by the definition of $C_{0}$ in Notation~\ref{8.4:not-step-zones}.
The holonomy $U(\partial p)$ is the ordered product of the bond variables around $\partial p$;
hence $U_{V}(u)(\partial p)$ and $U_{V^{\sharp}}(u)(\partial p)$ are the same ordered product of
values of $u$ and of $V_{0}\restriction C_{0}=V^{\sharp}_{0}\restriction C_{0}$, and the terms of
the two actions indexed by $p$ are equal. Let $p\in\mathcal{P}_{0}$. All bonds of $p$ lie in
$E_{0}$, so
the holonomy $U_{V}(u)(\partial p)=V_{0}(\partial p)$ involves only $V_{0}$, and likewise
$U_{V^{\sharp}}(u)(\partial p)=V^{\sharp}_{0}(\partial p)$ involves only $V^{\sharp}_{0}$;
neither involves $u$. Subtracting, the terms with $p\in\mathcal{P}_{1}$ cancel, and for
$p\in\mathcal{P}_{0}$ one has
\begin{equation*}
\bigl(1-\operatorname{Re}\operatorname{tr}V_{0}(\partial p)\bigr)
-\bigl(1-\operatorname{Re}\operatorname{tr}V^{\sharp}_{0}(\partial p)\bigr)
=\operatorname{Re}\operatorname{tr}V^{\sharp}_{0}(\partial p)
-\operatorname{Re}\operatorname{tr}V_{0}(\partial p).
\end{equation*}
Therefore
\begin{equation}\label{8.4:eq-zone-zero-invisible-1}
\begin{split}
A(U_{V}(u))-A(U_{V^{\sharp}}(u))
&=a_{\eta}\sum_{p\in\mathcal{P}_{0}}\bigl[\operatorname{Re}\operatorname{tr}V^{\sharp}_{0}(\partial p)
-\operatorname{Re}\operatorname{tr}V_{0}(\partial p)\bigr]\\
&=:\kappa(V,V^{\sharp}),
\end{split}
\end{equation}
and the right-hand side does not depend on $u$. This proves (a).

(b) The condition $M_{\mathbf{B}_{j}}(U)=V$ consists of the constraints
$M^{n}(U)\restriction\Gamma_{n}=V_{n}$ for $0\le n\le j$. For $n=0$: by hypothesis (i) the zone-0
constraint for the data $V$ is $U(b)=V_{0}(b)$ for all $b\in E_{0}$, which $U_{V}(u)$ satisfies by
construction; in the same way $U_{V^{\sharp}}(u)$ satisfies the zone-0 constraint for the data
$V^{\sharp}$ by construction. For $1\le n\le j$: by hypothesis (ii), the constraint
$M^{n}(U)\restriction\Gamma_{n}=V_{n}$ involves only the values of $U$ on fine bonds with both
endpoints
in blocks of $Z_{n}$, and these blocks lie in $\Omega_{1}$; such bonds have their endpoints in
$\Omega_{1}$ and therefore belong to $E_{1}$. On $E_{1}$ one has $U_{V}(u)=U_{V^{\sharp}}(u)=u$,
and $V_{n}=V^{\sharp}_{n}$ by assumption. Hence for each $n\ge1$ the constraint
$M^{n}(U_{V}(u))\restriction\Gamma_{n}=V_{n}$ is the same condition on $u$ as the constraint
$M^{n}(U_{V^{\sharp}}(u))\restriction\Gamma_{n}=V^{\sharp}_{n}$. Together with the case $n=0$ this
gives the equivalence in (b). Moreover, by hypothesis (ii) the axial gauge conditions concern bonds
of $E_{1}$ only, on which the two configurations both equal $u$; so they are the same conditions on
$u$ for the data $V$ and for the data $V^{\sharp}$.

(c) By (b), the set of configurations $u$ on $E_{1}$ satisfying the constraints for the data $V$
coincides with the set of those satisfying the constraints for the data $V^{\sharp}$, and by (a), on
this common set the two constrained functionals $u\mapsto A(U_{V}(u))$ and
$u\mapsto A(U_{V^{\sharp}}(u))$ differ by the constant $\kappa(V,V^{\sharp})$. Two functions on the
same constraint set that differ by a constant have the same constrained critical points and the same
constrained minimisers, and their second differentials along variations of $u$ coincide. Since the
axial gauge conditions are, by the last sentence of the proof of (b), the same conditions on $u$ for
both data, the axially gauged minimiser for $V$ is the axially gauged minimiser for $V^{\sharp}$.
This proves (c).
\end{proof}

\begin{lemma}[The restricted regularity hypothesis holds for the base data]\label{8.4:lem-base-regularity}
Let the step index $j\ge1$, the admissible domains $\Omega_{n}$, the zones $Z_{n}=\Omega_{n}\setminus\Omega_{n+1}$,
the determining set $\mathbf{B}_{j}=\bigcup_{n}\Gamma_{n}$ and the collar bond set $C_{0}$ be as in
Notation~\ref{8.4:not-step-zones}, let the base data $V^{\mathrm{base}}$ be as in
Notation~\ref{8.4:not-box-variation}, and let $\varepsilon_{n}$ be
Ba{\l}aban's small-field thresholds. Let $c_{\Gamma}$ and $c'$ be constants for which the following
hypotheses hold for the configuration of every term in the support of the step-$j$ density.
\begin{itemize}
\item[(a)] \emph{Regularity below the top level}, in the form stated in \cite[p.~256]{B-CMP119}: for
$1\le n\le j-1$ the live data $V_{n}$ satisfy $|\partial V_{n}-1|<c_{\Gamma}L^{2}\varepsilon_{n}$ at
the level-$n$ plaquettes of $\Gamma_{n}$, and at the zone interfaces the data are regular on
$\mathbf{B}_{j}$ in the sense of \eqref{8.4:eq-step-zones-a}, read with the interface convention of
\cite{B-CMP102b}.
\item[(b)] \emph{Regularity at the top level}, as stated in \cite{B-CMP109} and
\cite[(3.13)--(3.14)]{B-CMP119}: $|\partial V_{j}-1|<2\varepsilon_{j}$
at the level-$j$ plaquettes of $\Omega_{j}$, and
$|V_{j}(b)(V^{(j)}(b))^{-1}-1|<c'\varepsilon_{j}$ on the level-$j$ bonds $b$ of $\Omega_{j+1}$; here
$V^{(j)}$ is the datum carried by $V^{\mathrm{base}}$ on the level-$j$ bonds of $\Omega_{j+1}$
(Notation~\ref{8.4:not-box-variation}). The support in question lies in that of the small-field
characteristic function $\chi_{j}(\Omega_{j})$ and of the step-$(j+1)$ small-field functions, on which
these bounds are stated in \cite[(3.13)--(3.14)]{B-CMP119}.
\item[(c)] \emph{Regularity on the zone-$0$ collar}, as stated in \cite{B-CMP119}: the distance
between $\Omega_{1}$ and the union of the large-field regions is at least $2LMR_{1}$, only small-field
characteristic functions occur on Ba{\l}aban's enlarged first domain $\tilde{\Omega}_{1}$, and
$|\partial V_{0}-1|<\varepsilon_{0}$ at every fine plaquette of $\tilde{\Omega}_{1}\setminus\Omega_{1}$.
\item[(d)] $c'\varepsilon_{j}\le\tfrac14$.
\end{itemize}
Put
\begin{equation}\label{8.4:eq-base-regularity-a}
\bar\varepsilon:=c_{7}\max_{0\le n\le j}\varepsilon_{n},\qquad
c_{7}:=\max\{c_{\Gamma}L^{2},\,2+8c',\,1\}.
\end{equation}
Then for every term selected by a history $h$ in the set $\mathcal{C}_{S}(D)$ of clean histories
of the localisation of the sourced step, with $D$ as there
(indeed by any history in the support of the step-$j$ density) the base data $V^{\mathrm{base}}$
satisfy $(7^{\flat})(\bar\varepsilon)$ of \eqref{8.4:eq-step-zones-a} at all plaquettes of the zones
$1\le n\le j$, at the interface plaquettes, and at the zone-$0$ plaquettes meeting $C_{0}$. At deep
zone-$0$ plaquettes lying in level-$0$ large-field regions the inequality fails in general.
\end{lemma}

\begin{proof}
Fix a term selected by a history in the support of the step-$j$ density; every history in
$\mathcal{C}_{S}(D)$ is of this kind. Hypotheses (a)--(c) therefore apply to the configuration of the
term. Throughout, for a
configuration $V$ and a plaquette $q$ we write $(\partial V)(q)=V(\partial q)$ for the ordered product
of $V$ around $\partial q$.

\emph{Zones $1\le n\le j-1$.} On zone $n$ the base data are the live data $V_{n}$. By (a), at every
level-$n$ plaquette of $\Gamma_{n}$,
\begin{equation*}
|\partial V_{n}-1|<c_{\Gamma}L^{2}\varepsilon_{n}\le c_{7}\varepsilon_{n}\le\bar\varepsilon,
\end{equation*}
because $c_{7}\ge c_{\Gamma}L^{2}$ and $\varepsilon_{n}\le\max_{0\le n'\le j}\varepsilon_{n'}$.
At the interface plaquettes, (a) supplies the regularity in the sense of
\eqref{8.4:eq-step-zones-a}, with the convention of \cite{B-CMP102b} at zone interfaces.

\emph{Top zone.} On the level-$j$ bonds outside $\Omega_{j+1}$ the base datum is $V_{j}$, and (b) gives
\begin{equation*}
|\partial V_{j}-1|<2\varepsilon_{j}\le c_{7}\varepsilon_{j}\le\bar\varepsilon
\end{equation*}
at every level-$j$
plaquette all of whose bonds lie there, since $c_{7}\ge2+8c'\ge2$. On the level-$j$ bonds inside
$\Omega_{j+1}$ the base datum is $V^{(j)}$. Define, for such bonds $b$,
$V_{j}'(b):=V_{j}(b)(V^{(j)}(b))^{-1}$, so that
\begin{equation*}
V_{j}(b)=V_{j}'(b)\,V^{(j)}(b),\qquad |V_{j}'(b)-1|<c'\varepsilon_{j}
\end{equation*}
by (b), while $|\partial V_{j}-1|<2\varepsilon_{j}$. Since $V_{j}(b)$ and $V^{(j)}(b)$ lie in
$\mathrm{SU}(N)$, so does $V_{j}'(b)$, and $|V_{j}'(b)^{-1}-1|=|V_{j}'(b)^{*}-1|=|V_{j}'(b)-1|$.

Let $q$ be a level-$j$ plaquette with bonds $c_{1},\dots,c_{4}$ inside $\Omega_{j+1}$. Write
$V^{(j)}(c_{i})=V_{j}'(c_{i})^{-1}V_{j}(c_{i})$; for a bond traversed against its orientation the
factor is $V^{(j)}(c_{i})^{-1}=V_{j}(c_{i})^{-1}V_{j}'(c_{i})$. Hence $V^{(j)}(\partial q)$ is obtained
from the product $V_{j}(\partial q)$ of four unitary factors by inserting four unitary factors
$B_{1},\dots,B_{4}$, each with $|B_{i}-1|<\delta:=c'\varepsilon_{j}$. Removing the inserted factors one
at a time, i.e.\ writing the difference as a telescoping sum of four terms, the $i$-th term is a product
of unitary factors, of the not yet removed factors $B_{l}$, $l<i$ (each of norm at most $1+\delta$), and
of one factor $B_{i}-1$. Hence
\begin{align*}
|V^{(j)}(\partial q)-V_{j}(\partial q)|&\le\sum_{i=1}^{4}\delta(1+\delta)^{i-1}
\le4c'\varepsilon_{j}(1+c'\varepsilon_{j})^{3}\le8c'\varepsilon_{j},\\
4(1+c'\varepsilon_{j})^{3}&\le4\cdot(5/4)^{3}=125/16<8,
\end{align*}
where the second line uses (d) and gives the last inequality of the first.
Therefore
\begin{equation*}
|\partial V^{(j)}(q)-1|\le|V^{(j)}(\partial q)-V_{j}(\partial q)|+|V_{j}(\partial q)-1|
<(2+8c')\varepsilon_{j}\le c_{7}\varepsilon_{j}\le\bar\varepsilon .
\end{equation*}
Mixed plaquettes, with some bonds inside $\Omega_{j+1}$ and some outside, are treated identically with
fewer insertions, and the telescoping sum then has fewer terms, so the same bound holds.

We note that hypothesis (d) holds whenever $\bar\varepsilon\le1$: indeed, by
\eqref{8.4:eq-base-regularity-a},
\begin{equation*}
c'\varepsilon_{j}\le\frac{c'\bar\varepsilon}{c_{7}}\le\frac{c'}{2+8c'}\le\frac18 .
\end{equation*}
In particular it follows from the ceiling \eqref{8.4:eq-far-member-a} when the smallness constant
$a_{1}$ there satisfies $a_{1}\le1$, and, if $a_{1}>1$, once that ceiling is read with
$\min(a_{1},1)$ in place of $a_{1}$.

\emph{Zone-$0$ collar.} By (c), $|\partial V_{0}-1|<\varepsilon_{0}$ at every fine plaquette of
$\tilde{\Omega}_{1}\setminus\Omega_{1}$. Since $LMR_{1}\ge2$, the zone-$0$ plaquettes meeting $C_{0}$ are
among these, and there
\begin{equation*}
|\partial V_{0}-1|<\varepsilon_{0}\le c_{7}\varepsilon_{0}\le\bar\varepsilon,
\end{equation*}
because $c_{7}\ge1$.

\emph{Deep zone $0$.} Inside a level-$0$ large-field region of the term the bare plaquettes violate the
step-$1$ small-field condition, by the definition of that region in \cite{B-CMP119}. Hence the
inequality of \eqref{8.4:eq-step-zones-a} fails there in general.
\end{proof}

\begin{lemma}[The far-member lemma]\label{8.4:lem-far-member}
Let the step index $j$, the fine lattice $\mathbb{T}$ of spacing $\eta$, the zone index $n(x)$ of a
fine site $x$, the sets $\Omega_{n}$ and $\Gamma_{n}$, the determining set $\mathbf{B}_{j}$, the
bond sets $E_{1}$, $E_{0}$, $C_{0}$ and the data be as in Notation~\ref{8.4:not-step-zones}. Let the box
$I$, the variables $y=(y_{b})_{b\in I}$ and their extension $\tilde{y}$ by zero, the multiplier
$V'(y)=\exp(iB(y))$ of \eqref{8.4:eq-box-variation-a}, the base data $V^{\mathrm{base}}$, the varied
data $V^{(y)}=V'(y)V^{\mathrm{base}}$, the varied minimiser $U^{y}$, the base minimiser $U_{0}$,
the tubes $\mathcal{T}(Y_{R},Y_{I})$ with $Y=Y_{R}+Y_{I}$, the derivatives $\partial_{b,\zeta}$ and
the covariant curl $D_{U_{0}}$ be as in Notation~\ref{8.4:not-box-variation}; a point $y$ of a
tube is called real if $y_{b}\in\mathfrak{g}$ for every $b\in I$. Let $E(x,b)$, $K$ and
$\Theta(x)$ be as in Notation~\ref{8.4:not-decay-weight}.

Assume the following statements.
\begin{itemize}
\item Ba{\l}aban's analyticity of the minimiser in the data,
Theorem~\ref{8.4:thm-minimiser-analyticity}, that is \cite[Proposition~9]{B-CMP102b} with the
derivative bound \cite[(190)]{B-CMP102b}.
\item The identification stated in \cite[p.~256]{B-CMP119}: the variational problem
\cite[(2.12)]{B-CMP119} for $\mathbf{B}_{j}$ is the variational problem of \cite{B-CMP102b}, so that
$U(\mathbf{B}_{j},V)$ is the minimal configuration $U_{k}(V)$ of \cite{B-CMP102b} with $k=j$ and
determining set $\mathfrak{B}_{k}=\mathbf{B}_{j}$, whose zones $\Lambda_{n}$ are the sets $\Gamma_{n}$.
\item Applicability to the base data: if $V^{\mathrm{base}}$ satisfies the restricted regularity
condition \eqref{8.4:eq-step-zones-a} with a threshold $\bar\varepsilon\le a_{1}$, then the conclusions
of \cite[Theorem~1]{B-CMP102b} and of \cite[Proposition~9]{B-CMP102b} hold for
$(\mathbf{B}_{j},V^{\mathrm{base}})$, the base configuration $U(\mathbf{B}_{j},V^{\mathrm{base}})$
being the minimiser on $E_{1}$ given by Lemma~\ref{8.4:lem-zone-zero-invisible} together with $V_{0}$
on $E_{0}$. In this form \cite[Theorem~1]{B-CMP102b} is applied in \cite{B-CMP119}.
\item The hypotheses on the response estimate of Definition~\ref{8.4:def-response-hypotheses}.
\item Ba{\l}aban's distance $d(\cdot,\cdot)$ of \cite[(190)]{B-CMP102b} satisfies one of the two
lower bounds
\begin{equation*}
d(y_{x},y')\ge c_{d}\operatorname{dist}_{j}\bigl(x,b(y',\nu)\bigr)-c_{d}',\qquad
d(y_{x},y')\ge c_{d}\rho_{\mathrm{sc}}\bigl(x,b(y',\nu)\bigr)-c_{d}',
\end{equation*}
each required for every fine site $x$, every site $y_{x}$ with $x\in\Delta(y_{x})$, where the sets
$\Delta(\cdot)$ are the localisation sets of \cite[(190)]{B-CMP102b}, and every bond $b(y',\nu)$ of
the top zone $\Gamma_{j}$ with initial site $y'$ and direction $\nu$; here $c_{d}>0$ and
$c_{d}'\ge0$ depend on $d$ and $L$ only. Notation~\ref{8.4:not-decay-weight} refers to these two
lower bounds in fixing $\rho$.
\item The properties of Ba{\l}aban's operators $\mathbf{C}$, $h$, $\tilde{D}$
(\cite{B-CMP119}, \cite{B-CMP109}) stated in Definition~\ref{8.4:def-response-hypotheses}: the
linearity and locality of $\mathbf{C}$ with the norm bound $c_{C}$, the bound $c_{h}$ on $h$, and the
holomorphy of $\tilde{D}$ on $\{X:\sup|X|<b_{D}\}$ with $\tilde{D}(0)=0$, its bounds with constant
$c_{D}$ and its locality; the constants $c_{C}$, $c_{h}$, $c_{D}$, $b_{D}$ depend on $d$ and $L$ only.
For a coarse bond $c$ (the bonds indexing the values $h(c)$ and $\tilde{D}(X)(c)$) we write $c_{-}$,
$c_{+}$ for its endpoints, $\Box(c_{\pm})$ for their blocks, $\operatorname{cor}(c)$ for its
corridor, and $b_{0}(c)$ for the bond of $\mathbb{T}^{(j)}$ at which $\mathbf{C}$ and $h$ place
their values.
\end{itemize}
Assume moreover:
\begin{itemize}
\item[(a)] the restricted regularity condition \eqref{8.4:eq-step-zones-a} holds for
$V^{\mathrm{base}}$ with a threshold satisfying the ceiling
\begin{equation}\label{8.4:eq-far-member-a}
\bar\varepsilon\le a_{1};
\end{equation}
\item[(b)] the box is clean: every bond of $\mathbb{T}^{(j)}$ at $\operatorname{dist}_{j}$-distance
at most $4L+2$ from $I$ has both endpoints in $\Omega_{j+1}$;
\item[(c)] the tube is small:
\begin{equation}\label{8.4:eq-far-member-b}
g_{j}c_{C}Y\le y_{*}:=\min\Bigl\{\frac{b_{D}}{2},\ \frac{1}{2c_{h}c_{D}},\ \frac{a_{9}}{2}\Bigr\},
\end{equation}
where $a_{9}>0$ is the radius of a ball $\{B:\sup|B|<a_{9}\}$ contained in the analyticity domain
\cite[(172)]{B-CMP102b}, as in Definition~\ref{8.4:def-response-hypotheses}.
\end{itemize}
Put $\mathcal{T}:=\mathcal{T}(Y_{R},Y_{I})$. Then the following hold.
\begin{itemize}
\item[(i)] (Representation, holomorphy.) There are a map $y\mapsto\mathcal{H}^{(y)}$ from
$\mathcal{T}$ to $\mathfrak{g}^{c}$-valued functions on the fine bonds of $\mathbb{T}$, holomorphic in
$y$ (each value $\mathcal{H}^{(y)}(b)$ is a holomorphic function on $\mathcal{T}$), with
$\mathcal{H}^{(0)}=0$, and for each real $y\in\mathcal{T}$ a $G$-valued gauge transformation
$u_{y}$, such that for real $y\in\mathcal{T}$
\begin{equation}\label{8.4:eq-far-member-c}
U^{y}=\bigl(\mathcal{W}^{(y)}\bigr)^{u_{y}},\qquad
\mathcal{W}^{(y)}:=\exp\bigl(i\eta\mathcal{H}^{(y)}\bigr)\cdot U_{0},
\end{equation}
the product being bondwise: $\mathcal{W}^{(y)}(b)=\exp(i\eta\mathcal{H}^{(y)}(b))\,U_{0}(b)$.
\item[(ii)] (Two-point decay of the first derivative, uniformly on the tube.) For all
$y\in\mathcal{T}$, $b\in I$, $\zeta\in\mathfrak{g}^{c}$ with $|\zeta|\le1$, every fine bond $b_{x}$
with initial site $x$ and every fine plaquette $p$ with base site $x$, with $n:=n(x)$,
\begin{align}
\eta\,\bigl|\partial_{b,\zeta}\mathcal{H}^{(y)}(b_{x})\bigr|
&\le K g_{j}L^{-n}E(x,b),\label{8.4:eq-far-member-d}\\
\eta^{2}\bigl|\bigl(D_{U_{0}}\partial_{b,\zeta}\mathcal{H}^{(y)}\bigr)(p)\bigr|
&\le K g_{j}L^{-2n}E(x,b).\label{8.4:eq-far-member-e}
\end{align}
Here $E$ is the decay weight \eqref{8.4:eq-decay-weight-a}, with $\rho=\operatorname{dist}_{j}$
under the lower bound by $c_{d}\operatorname{dist}_{j}-c_{d}'$ and $\rho=\rho_{\mathrm{sc}}$ under the
lower bound by $c_{d}\rho_{\mathrm{sc}}-c_{d}'$; the constant $K=4c_{*}c_{C}N_{B}$ of
\eqref{8.4:eq-decay-weight-b}, with $N_{B}:=(4d+2)^{2}$, depends on $d$, $L$, $N$ and Ba{\l}aban's
constants only, not on $j$, $n$, the history $h$ in the sense of Notation~\ref{0.2:not-glossary}
(not the operator $h$ above), $I$, the clean collar width $W$, or the volume. The bound
\eqref{8.4:eq-far-member-e} is the one that uses the hypothesis of
Definition~\ref{8.4:def-response-hypotheses} identifying the difference $\nabla_{x}$ of
\cite[(190)]{B-CMP102b} with the covariant difference with respect to the base background.
\item[(iii)] (Size on the tube.) For all $y\in\mathcal{T}$, $b_{x}$, $p$ and $n$ as in (ii),
\begin{equation*}
\eta\,|\mathcal{H}^{(y)}(b_{x})|\le K L^{-n}\Theta(x),\qquad
\eta^{2}\bigl|(D_{U_{0}}\mathcal{H}^{(y)})(p)\bigr|\le K L^{-2n}\Theta(x).
\end{equation*}
\item[(iv)] (Zone $0$ does not respond.) For every fine plaquette $p$ all of whose bonds lie in
$E_{0}$, $\operatorname{tr}\mathcal{W}^{(y)}(\partial p)=\operatorname{tr}V_{0}(\partial p)$ for all
$y\in\mathcal{T}$.
\item[(v)] (Mixed second derivatives.) Let $Y_{R}'<Y_{R}$, $Y_{I}'<Y_{I}$, and write
$r_{\mathcal{T}}:=\min(Y_{R}-Y_{R}',\,Y_{I}-Y_{I}')$ for the margin between the two tubes.
For all $y\in\mathcal{T}(Y_{R}',Y_{I}')$, $b,b'\in I$,
$\zeta,\zeta'\in\mathfrak{g}^{c}$ with $|\zeta|,|\zeta'|\le1$, and $b_{x}$, $n$ as in (ii),
\begin{equation}\label{8.4:eq-far-member-f}
\eta\,\bigl|\partial_{b',\zeta'}\partial_{b,\zeta}\mathcal{H}^{(y)}(b_{x})\bigr|
\le\frac{Kg_{j}}{r_{\mathcal{T}}}\,L^{-n}\min\{E(x,b),E(x,b')\},
\end{equation}
and the same bound holds with
$\eta^{2}|(D_{U_{0}}\partial_{b',\zeta'}\partial_{b,\zeta}\mathcal{H}^{(y)})(p)|$ on the left and
$L^{-2n}$ in place of $L^{-n}$ on the right, for every fine plaquette $p$ with base site $x$.
\item[(vi)] (Holomorphic extension of gauge-invariant functionals.) Let $F$ be a function of
configurations on a set $\tilde{X}$ of fine bonds, invariant under $G$-valued gauge
transformations, such that $Z\mapsto F\bigl(\exp(i\eta Z)\cdot(U_{0}\restriction\tilde{X})\bigr)$ is
holomorphic on an open set containing
$\{\mathcal{H}^{(y)}\restriction\tilde{X}:y\in\mathcal{T}\}$. Then
$y\mapsto F(\mathcal{W}^{(y)}\restriction\tilde{X})$ is holomorphic on $\mathcal{T}$ and is the unique
holomorphic extension of the function $y\mapsto F(U^{y}\restriction\tilde{X})$ of real
$y\in\mathcal{T}$.
\end{itemize}
\end{lemma}

\begin{proof}[Proof of clause \textup{(i)} of Lemma~\ref{8.4:lem-far-member}]
The proof of clause (i) proceeds in three steps: the variational problem and its data; the
variation as a small multiplier of the top-zone datum, holomorphic in $y$; the definition of
$\mathcal{H}^{(y)}$ and the representation \eqref{8.4:eq-far-member-c}.

\emph{First step: the variational problem and its data.} By the identification stated in
\cite[p.~256]{B-CMP119}, the map $V\mapsto U(\mathbf{B}_{j},V)$ is the map $V\mapsto U_{k}(V)$ of
\cite{B-CMP102b} for $k=j$ and determining set $\mathfrak{B}_{k}=\mathbf{B}_{j}$. By hypothesis (a)
the base data $V^{\mathrm{base}}$ satisfy the restricted regularity condition
\eqref{8.4:eq-step-zones-a}, with threshold $\bar\varepsilon\le a_{1}$ by
\eqref{8.4:eq-far-member-a}. By Lemma~\ref{8.4:lem-zone-zero-invisible}, the constrained problem for the
unknowns $U\restriction E_{1}$ does not change if $V_{0}$ is modified on $E_{0}\setminus C_{0}$. By the
applicability statement assumed above, \cite[Theorem~1]{B-CMP102b} and
\cite[Proposition~9]{B-CMP102b} therefore apply to $(\mathbf{B}_{j},V^{\mathrm{base}})$: the base
minimiser $U_{0}=U(\mathbf{B}_{j},V^{\mathrm{base}})$ exists and is unique in the axial gauge, and
\cite[Proposition~9]{B-CMP102b}, in the form of Theorem~\ref{8.4:thm-minimiser-analyticity}, is
available with $V^{\mathrm{base}}$ in the role of the configuration called $V_{0}$ there (not to be
confused with the zone-$0$ dat\-um $V_{0}$ of $V^{\mathrm{base}}$).

\emph{Second step: the variation is a small top-zone multiplier, holomorphic in $y$.} By
\eqref{8.4:eq-box-variation-a},
\begin{equation*}
B(y)=g_{j}\mathbf{C}\tilde{y}-h\tilde{D}(g_{j}\mathbf{C}\tilde{y}).
\end{equation*}
The operator $\mathbf{C}$ is linear, $h$ is linear, and $\tilde{D}$ is holomorphic on
$\{X:\sup|X|<b_{D}\}$; \cite{B-CMP109} describes it as an analytic function of its argument. For
$y\in\mathcal{T}$ we have $\max_{b}|y_{b}|\le Y$ (Notation~\ref{8.4:not-box-variation}), hence, by the
norm bound $c_{C}$ on $\mathbf{C}$ and by the first entry of the minimum in
\eqref{8.4:eq-far-member-b},
\begin{equation*}
\sup|g_{j}\mathbf{C}\tilde{y}|\le g_{j}c_{C}\max_{b}|y_{b}|\le g_{j}c_{C}Y\le\frac{b_{D}}{2}.
\end{equation*}
Thus $y\mapsto g_{j}\mathbf{C}\tilde{y}$ is a linear map of $\mathcal{T}$ into the domain of
holomorphy of $\tilde{D}$, and $y\mapsto B(y)(b')$ is holomorphic on $\mathcal{T}$ for every bond $b'$
of $\mathbb{T}^{(j)}$, as the difference of a linear function and the composition of the linear map
$h$ with the holomorphic function $y\mapsto\tilde{D}(g_{j}\mathbf{C}\tilde{y})$.

For $b\in I$ and $\zeta\in\mathfrak{g}^{c}$ let $\zeta e_{b}$ denote the element of
$(\mathfrak{g}^{c})^{I}$ equal to $\zeta$ at $b$ and to $0$ elsewhere, extended by zero to all
bonds of $\mathbb{T}^{(j)}$. Since $\tilde{y}$ is linear in $y$, the chain rule gives
\begin{equation*}
\partial_{b,\zeta}B(y)=g_{j}\mathbf{C}(\zeta e_{b})
-h\,d\tilde{D}_{g_{j}\mathbf{C}\tilde{y}}\bigl[g_{j}\mathbf{C}(\zeta e_{b})\bigr].
\end{equation*}
We now derive the bound \eqref{8.4:eq-response-hypotheses-a}. The argument
$X=g_{j}\mathbf{C}\tilde{y}$ satisfies $\sup|X|\le g_{j}c_{C}Y\le b_{D}/2$, so the bound on
$d\tilde{D}_{X}$ applies; the increment $X'=g_{j}\mathbf{C}(\zeta e_{b})$ satisfies
$\sup|X'|\le g_{j}c_{C}|\zeta|$; and $|h(c)|\le c_{h}$. Hence, for every bond $b'$ of
$\mathbb{T}^{(j)}$,
\begin{align*}
|\partial_{b,\zeta}B(y)(b')|
&\le g_{j}\bigl|(\mathbf{C}(\zeta e_{b}))(b')\bigr|
+c_{h}\bigl|d\tilde{D}_{g_{j}\mathbf{C}\tilde{y}}[g_{j}\mathbf{C}(\zeta e_{b})]\bigr|\\
&\le g_{j}c_{C}|\zeta|\bigl(1+c_{h}c_{D}g_{j}c_{C}Y\bigr)
\le 2g_{j}c_{C}|\zeta|,
\end{align*}
the last step because $c_{h}c_{D}g_{j}c_{C}Y\le\tfrac12$ by the second entry of the minimum in
\eqref{8.4:eq-far-member-b}. In the same way, from $|\tilde{D}(X)(c)|\le c_{D}\sup|X|^{2}$ and
$|h(c)|\le c_{h}$,
\begin{equation*}
\sup|B(y)|\le g_{j}c_{C}Y+c_{h}c_{D}(g_{j}c_{C}Y)^{2}\le 2g_{j}c_{C}Y\le a_{9},
\end{equation*}
where the middle inequality uses $c_{h}c_{D}g_{j}c_{C}Y\le\tfrac12$ once more and the last one uses the
third entry of the minimum in \eqref{8.4:eq-far-member-b}.

Next, the support. By the description of $\mathbf{C}$, the configuration $\mathbf{C}(\zeta e_{b})$ is
supported on $b$ and on the bonds $b_{0}(c)$ for the at most $4d$ coarse bonds $c$ one of whose blocks
contains an endpoint of $b$. By the locality of $\tilde{D}$, $d\tilde{D}_{X}[X'](c)\neq0$ only if
$X'\neq0$ on a bond of $\Box(c_{-})\cup\Box(c_{+})\cup\operatorname{cor}(c)$, and $h$ places the value
at $b_{0}(c)$. Let
$S(b)$ be the support of $\partial_{b,\zeta}B(y)$. Counting the coarse bonds whose block pair meets the
support of $\mathbf{C}(\zeta e_{b})$ gives
\begin{equation*}
|S(b)|\le(4d+1)+(4d+1)^{2}=(4d+1)(4d+2)\le(4d+2)^{2}=N_{B},
\end{equation*}
and all bonds of $S(b)$ lie at $\operatorname{dist}_{j}$-distance at most $4L+2$ from $b$. Put
$S(I):=\bigcup_{b\in I}S(b)$. The same description applies to $B(y)$ itself: $\mathbf{C}\tilde{y}$ is
supported on $I$ and on the bonds $b_{0}(c)$ above, and $\tilde{D}(X)(c)$, which vanishes at $X=0$ and
depends on $X$ only on $\Box(c_{-})\cup\Box(c_{+})\cup\operatorname{cor}(c)$, vanishes when $X$
vanishes there; hence
$B(y)$ is supported in $S(I)$.

By hypothesis (b), every bond of $S(b)$, lying at $\operatorname{dist}_{j}$-distance at most $4L+2$
from $I$, has both endpoints in
$\Omega_{j+1}\subset\Omega_{j}$; thus $S(I)$ consists of bonds of the top zone $\Gamma_{j}$. Hence
$V'(y)=\exp(iB(y))$ multiplies the top-zone datum only: $V^{(y)}=V'(y)V^{\mathrm{base}}$ is exactly
of the form $V'V_{0}$ of \cite[Proposition~9]{B-CMP102b}, with $V'$ supported in $S(I)$, a set of
bonds of $\Gamma_{j}$, and with
\begin{equation*}
B:=\frac{1}{i}\log V'(y)=B(y),\qquad \sup|B|<a_{9}.
\end{equation*}
By the hypothesis of Definition~\ref{8.4:def-response-hypotheses} that the analyticity domain
\cite[(172)]{B-CMP102b} contains $\{B:\sup|B|<a_{9}\}$ and that \cite[(190)]{B-CMP102b} holds at
every point of it, $B(y)$ lies in the domain \cite[(172)]{B-CMP102b} for every $y\in\mathcal{T}$.

\emph{Third step: definition of $\mathcal{H}^{(y)}$, holomorphy, and the representation.} Define
\begin{equation*}
\mathcal{H}^{(y)}:=\mathcal{H}(B(y)),
\end{equation*}
where $\mathcal{H}$ is the function of Theorem~\ref{8.4:thm-minimiser-analyticity} for the base data
$V^{\mathrm{base}}$, which is available by the first step. By
Theorem~\ref{8.4:thm-minimiser-analyticity}, $\mathcal{H}$ is analytic in $B$ on the domain
\cite[(172)]{B-CMP102b}. By the second step, $y\mapsto B(y)$ is holomorphic on $\mathcal{T}$ with
values in that domain; the composition $y\mapsto\mathcal{H}^{(y)}(b)$ is therefore holomorphic on
$\mathcal{T}$ for every fine bond $b$. Next, $B(0)=0$, because $\tilde{0}=0$, $\mathbf{C}$ is linear
and $\tilde{D}(0)=0$; hence
\begin{equation*}
\mathcal{H}^{(0)}=\mathcal{H}(B(0))=\mathcal{H}(0)=0,
\end{equation*}
the last equality being the hypothesis of Definition~\ref{8.4:def-response-hypotheses} that
$\mathcal{H}(0)=0$ (at $V'=1$ the represented configuration is the identity, whose Landau-gauge form is
the identity).

Let $y\in\mathcal{T}$ be real. Then $B(y)$ is $\mathfrak{g}$-valued, because $\tilde{y}$ is
$\mathfrak{g}$-valued and $\mathbf{C}$, $h$ and $\tilde{D}$ preserve $\mathfrak{g}$-valuedness
($\tilde{D}$ is introduced in \cite{B-CMP109} as a $\mathfrak{g}$-valued function of its argument).
Hence $V'(y)=\exp(iB(y))$ is $G$-valued, since $\exp(iX)\in G$ for $X\in\mathfrak{g}$
(Notation~\ref{8.4:not-step-zones}), and $U^{y}=U_{k}(V'(y)V^{\mathrm{base}})$ is the $G$-valued
axial-gauge minimiser. By Theorem~\ref{8.4:thm-minimiser-analyticity}, read with the hypothesis of
Definition~\ref{8.4:def-response-hypotheses} on the Landau gauge (the base $U_{k}(V^{\mathrm{base}})$
is a fixed configuration in a fixed gauge independent of $B$, and for $G$-valued $V'$ the
configurations $U_{k}(V'V^{\mathrm{base}})$ and
$\exp(i\eta\mathcal{H}(B))\cdot U_{k}(V^{\mathrm{base}})$ are related by a $G$-valued gauge
transformation, which may depend on $B$), applied with $V'=V'(y)$, $B=B(y)$ and
$V^{\mathrm{base}}$ in the role of the base, the configuration $U^{y}$ is related by a $G$-valued gauge
transformation to
\begin{equation*}
\exp\bigl(i\eta\mathcal{H}(B(y))\bigr)\cdot U_{k}(V^{\mathrm{base}})
=\exp\bigl(i\eta\mathcal{H}^{(y)}\bigr)\cdot U_{0}=\mathcal{W}^{(y)},
\end{equation*}
where $U_{k}(V^{\mathrm{base}})=U_{0}$ by the first step. Writing $u_{y}$ for the $G$-valued gauge
transformation with $U^{y}=(\mathcal{W}^{(y)})^{u_{y}}$, this is \eqref{8.4:eq-far-member-c}. For real
$y$ the configuration $\mathcal{W}^{(y)}$ is $G$-valued, so $\eta\mathcal{H}^{(y)}$ is
$\mathfrak{g}$-valued there.

This proves clause (i). Clauses (ii)--(vi) are proved next.
\end{proof}

\begin{proof}[Proof of the far-member lemma: decay, size and extension]
\label{8.4:prf-far-member-proof}
We prove clauses (ii)--(vi) of Lemma~\ref{8.4:lem-far-member}. From clause (i) we use the
following:
\begin{itemize}
\item each value $\mathcal{H}^{(y)}(b)$ is a holomorphic function of $y$ on the tube
$\mathcal{T}=\mathcal{T}(Y_{R},Y_{I})$;
\item $\mathcal{H}^{(0)}=0$;
\item the representation \eqref{8.4:eq-far-member-c} of $U^{y}$, for real $y$, through the
gauge transformation $u_{y}$ and the fixed-gauge configuration $\mathcal{W}^{(y)}$.
\end{itemize}
Throughout, $\rho=\operatorname{dist}_{j}$ under the distance comparison with
$\operatorname{dist}_{j}$ of Definition~\ref{8.4:def-response-hypotheses}, and
$\rho=\rho_{\mathrm{sc}}$ under the comparison with $\rho_{\mathrm{sc}}$; $E(x,b)$ is the decay
weight \eqref{8.4:eq-decay-weight-a} and $c_{1}=c_{d}(4L+2)+c_{d}'$.

\emph{Clause (ii), the bound \eqref{8.4:eq-far-member-d}.}
Fix $y\in\mathcal{T}$, $b\in I$ and $\zeta\in\mathfrak{g}^{c}$ with $|\zeta|\le1$. Fix a fine bond
$b_{x}=\langle x,x+\eta e_{\mu}\rangle$ and put $n:=n(x)$. By the covering hypothesis of
Definition~\ref{8.4:def-response-hypotheses}, there is a point $y_{x}\in\Gamma_{n}$ with
$x\in\Delta(y_{x})$; we fix one.

The response $\mathcal{H}^{(y)}$ is the composition of $B\mapsto\mathcal{H}(B)$ with
$y\mapsto B(y)$. The first map is analytic in the finitely many variables $B_{\nu}(y')$, and the
second is holomorphic in $y$. By the consequence of the locality and finiteness properties of the
constraint-solving operator $\mathbf{C}$, the corridor operator and the analytic map $\tilde D$
stated in Definition~\ref{8.4:def-response-hypotheses}, the
derivative $\partial_{b,\zeta}B(y)$ is supported on the set $S(b)$ of at most $N_{B}$ top bonds.
The chain rule therefore gives
\begin{equation}\label{8.4:eq-far-member-proof-1}
\partial_{b,\zeta}\mathcal{H}^{(y)}_{\mu}(x)
=\sum_{(y',\nu)\in S(b)}
\frac{\delta\mathcal{H}_{\mu}(B,x)}{\delta B_{\nu}(y')}\Big|_{B=B(y)}
\bigl[\partial_{b,\zeta}B_{\nu}(y')(y)\bigr],
\end{equation}
where $\mathcal{H}^{(y)}_{\mu}(x)=\mathcal{H}^{(y)}(b_{x})$.

We now bound each summand, using three hypotheses of
Definition~\ref{8.4:def-response-hypotheses}:
\begin{itemize}
\item for a top bond $(y',\nu)$ the functional derivative is the derivative with respect to the
single $\mathfrak{g}^{c}$-valued variable $B_{\nu}(y')$, and the bound concerns the operator norm
of the resulting linear map $\mathfrak{g}^{c}\to\mathfrak{g}^{c}$;
\item the first quantity of \eqref{8.4:eq-minimiser-analyticity-a}, that is of
\cite[(190)]{B-CMP102b}, is bounded at every point of the analyticity domain, which contains all
$B$ with $\sup|B|<a_{9}$; this is applied at $B=B(y)$;
\item in the top zone $L^{j}\eta=1$.
\end{itemize}
The bounds \eqref{8.4:eq-response-hypotheses-a} hold under the proviso
$g_{j}c_{C}Y\le\min\{b_{D}/2,(2c_{h}c_{D})^{-1}\}$; this proviso is met, because the smallness
hypothesis of Lemma~\ref{8.4:lem-far-member} on the tube gives $g_{j}c_{C}Y\le y_{*}$, and
$y_{*}\le\min\{b_{D}/2,(2c_{h}c_{D})^{-1}\}$ by the definition of $y_{*}$.
Moreover, by the second bound of \eqref{8.4:eq-response-hypotheses-a}, the same smallness
hypothesis and $y_{*}\le a_{9}/2$,
\begin{equation*}
\sup_{b'}\bigl|B(y)(b')\bigr|\le2g_{j}c_{C}Y\le2y_{*}\le a_{9};
\end{equation*}
the bound on the first quantity is used at $B=B(y)$ by the hypothesis that
\eqref{8.4:eq-minimiser-analyticity-a} holds on the analyticity domain.
In \eqref{8.4:eq-minimiser-analyticity-a} the level of the bond variable $B_{\nu}(y')$ is the top
level $k=j$, so that the volume factor $(L^{j}\eta)^{-d}$ equals $1$, and we use
$\eta(L^{n}\eta)^{-1}=L^{-n}$.
Together with the first bound of \eqref{8.4:eq-response-hypotheses-a} and $|\zeta|\le1$, each
summand of \eqref{8.4:eq-far-member-proof-1} satisfies
\begin{align*}
\eta\,\Bigl|\frac{\delta\mathcal{H}_{\mu}(B,x)}{\delta B_{\nu}(y')}\Big|_{B=B(y)}
\bigl[\partial_{b,\zeta}B_{\nu}(y')(y)\bigr]\Bigr|
&\le c_{*}L^{-n}e^{-\delta_{0}d(y_{x},y')/8}\,
\bigl|\partial_{b,\zeta}B_{\nu}(y')(y)\bigr|\\
&\le c_{*}L^{-n}e^{-\delta_{0}d(y_{x},y')/8}\cdot2g_{j}c_{C}.
\end{align*}

Next we replace $d(y_{x},y')$ by $\rho(x,b)$. Write $b(y',\nu)$ for the top bond with initial
site $y'$ and direction $\nu$. Since $x\in\Delta(y_{x})$, the distance comparison of
Definition~\ref{8.4:def-response-hypotheses} (with $\operatorname{dist}_{j}$, respectively with
$\rho_{\mathrm{sc}}$) applies. Every element of $S(b)$ lies within top distance $4L+2$ of $b$.
Hence
\begin{equation}\label{8.4:eq-far-member-proof-2}
\begin{aligned}
d(y_{x},y')&\ge c_{d}\,\rho\bigl(x,b(y',\nu)\bigr)-c_{d}'\\
&\ge c_{d}\bigl(\rho(x,b)-(4L+2)\bigr)-c_{d}'\\
&=c_{d}\,\rho(x,b)-c_{1}.
\end{aligned}
\end{equation}
For $\rho=\rho_{\mathrm{sc}}$ the second inequality rests on two facts:
\begin{itemize}
\item the triangle inequality
\begin{equation*}
\rho_{\mathrm{sc}}(x,b)\le\rho_{\mathrm{sc}}\bigl(x,b(y',\nu)\bigr)
+\rho_{\mathrm{sc}}\bigl(b(y',\nu),b\bigr);
\end{equation*}
\item $\rho_{\mathrm{sc}}(b(y',\nu),b)\le4L+2$, because $b(y',\nu)$ and $b$ are joined by a
top-zone path of at most $4L+2$ top bonds, each consisting of $L^{j}$ fine bonds weighted $L^{-j}$.
\end{itemize}
The same two facts hold for $\operatorname{dist}_{j}$. Since also $d(y_{x},y')\ge0$, we have
$e^{-\delta_{0}d(y_{x},y')/8}\le1$, and by \eqref{8.4:eq-far-member-proof-2}
\begin{equation*}
e^{-\delta_{0}d(y_{x},y')/8}\le\exp\bigl(-(\delta_{0}/8)(c_{d}\rho(x,b)-c_{1})\bigr).
\end{equation*}
By \eqref{8.4:eq-decay-weight-a}, therefore, $e^{-\delta_{0}d(y_{x},y')/8}\le E(x,b)$.

Summing over the at most $N_{B}$ elements of $S(b)$ in \eqref{8.4:eq-far-member-proof-1}, we obtain
\begin{equation}\label{8.4:eq-far-member-proof-3}
\eta\,\bigl|\partial_{b,\zeta}\mathcal{H}^{(y)}(b_{x})\bigr|
\le N_{B}\cdot c_{*}\cdot2g_{j}c_{C}\cdot L^{-n}E(x,b)
=2c_{*}c_{C}N_{B}\,g_{j}L^{-n}E(x,b).
\end{equation}

\emph{Clause (ii), the bound \eqref{8.4:eq-far-member-e}.}
Let $p$ be a fine plaquette with base site $x$. For $(y',\nu)\in S(b)$ put
\begin{equation*}
Z_{y',\nu}:=\frac{\delta\mathcal{H}(B,\cdot)}{\delta B_{\nu}(y')}\Big|_{B=B(y)}
\bigl[\partial_{b,\zeta}B_{\nu}(y')(y)\bigr].
\end{equation*}
The covariant curl $D_{U_{0}}$ of \eqref{8.4:eq-box-variation-c} is linear. By
\eqref{8.4:eq-far-member-proof-1} we therefore have
\begin{equation*}
\bigl(D_{U_{0}}\partial_{b,\zeta}\mathcal{H}^{(y)}\bigr)(p)
=\sum_{(y',\nu)\in S(b)}\bigl(D_{U_{0}}Z_{y',\nu}\bigr)(p).
\end{equation*}
Definition~\ref{8.4:def-response-hypotheses} contains the hypothesis that the difference operator
in the second quantity of \eqref{8.4:eq-minimiser-analyticity-a} is the covariant difference with
respect to the base background. We apply it at the point $y_{x}$ (recall $x\in\Delta(y_{x})$),
with $\zeta$ replaced by $\partial_{b,\zeta}B_{\nu}(y')(y)$. Using $\eta^{2}(L^{n}\eta)^{-2}=L^{-2n}$
and the first bound of \eqref{8.4:eq-response-hypotheses-a}, each term satisfies
\begin{equation*}
\eta^{2}\bigl|(D_{U_{0}}Z_{y',\nu})(p)\bigr|
\le2c_{*}L^{-2n}e^{-\delta_{0}d(y_{x},y')/8}\cdot2g_{j}c_{C}.
\end{equation*}
The exponential is bounded by $E(x,b)$ exactly as above. Summing over the at most $N_{B}$ elements
of $S(b)$ gives
\begin{equation}\label{8.4:eq-far-member-proof-4}
\eta^{2}\bigl|\bigl(D_{U_{0}}\partial_{b,\zeta}\mathcal{H}^{(y)}\bigr)(p)\bigr|
\le N_{B}\cdot2c_{*}\cdot2g_{j}c_{C}\cdot L^{-2n}E(x,b)
=4c_{*}c_{C}N_{B}\,g_{j}L^{-2n}E(x,b).
\end{equation}

\emph{The value of $K$.}
We take the constant $K$ of \eqref{8.4:eq-decay-weight-b} to be $K=4c_{*}c_{C}N_{B}$; compared
with $2c_{*}c_{C}N_{B}$ this only enlarges $K$ in \eqref{8.4:eq-far-member-d}. With this
value:
\begin{itemize}
\item \eqref{8.4:eq-far-member-proof-4} is \eqref{8.4:eq-far-member-e};
\item \eqref{8.4:eq-far-member-proof-3} gives \eqref{8.4:eq-far-member-d}, since
$2c_{*}c_{C}N_{B}\le K$.
\end{itemize}
All constants built from $K$ below are built from this value.

\emph{Independence of $K$.}
By the uniformity hypothesis of Definition~\ref{8.4:def-response-hypotheses}, $c_{*}$ and
$\delta_{0}$ depend neither on $j$, nor on the admissible sequence $\{\Omega_{n}\}$, nor on the size
of $\mathbb{T}$. Also $c_{C}$ and $N_{B}=(4d+2)^{2}$ depend on $d$ and $L$ only. Hence $K$ depends
on $d$, $L$, $N$ and the constants of \cite{B-CMP102b} only, and not on $j$, $n$, the history $h$,
$I$, $W$ or the volume. This proves clause (ii).

\emph{Clause (iii).}
The tube $\mathcal{T}$ is convex and contains $0$. Hence $ty\in\mathcal{T}$ for $y\in\mathcal{T}$
and $t\in[0,1]$, and
\begin{equation}\label{8.4:eq-far-member-proof-5}
\mathcal{H}^{(y)}-\mathcal{H}^{(0)}
=\int_{0}^{1}\frac{d}{dt}\mathcal{H}^{(ty)}\,dt
=\int_{0}^{1}\sum_{b\in I}\partial_{b,y_{b}}\mathcal{H}^{(ty)}\,dt .
\end{equation}
Here $\partial_{b,y_{b}}$ is the directional derivative in the direction $y_{b}$. Terms with
$y_{b}=0$ vanish. Otherwise
$\partial_{b,y_{b}}=|y_{b}|\,\partial_{b,y_{b}/|y_{b}|}$, and the unit direction
$y_{b}/|y_{b}|$ is admissible in \eqref{8.4:eq-far-member-d}. Write $y_{b}=y_{b}'+iy_{b}''$ with
$y_{b}',y_{b}''\in\mathfrak{g}$, the decomposition entering the definition
\eqref{8.4:eq-box-variation-b} of the tube. Then
\begin{equation*}
|y_{b}|\le|y_{b}'|+|y_{b}''|<Y_{R}+Y_{I}=Y .
\end{equation*}
We insert $\mathcal{H}^{(0)}=0$ into \eqref{8.4:eq-far-member-proof-5}, evaluate at $b_{x}$, and
apply \eqref{8.4:eq-far-member-d} at each point $ty\in\mathcal{T}$. By the definition of
$\Theta(x)$ in \eqref{8.4:eq-decay-weight-b},
\begin{equation*}
\eta\,\bigl|\mathcal{H}^{(y)}(b_{x})\bigr|
\le\sum_{b\in I}Kg_{j}L^{-n}E(x,b)\cdot Y=KL^{-n}\Theta(x).
\end{equation*}
The operator $D_{U_{0}}$ is linear, so it may be applied under the integral in
\eqref{8.4:eq-far-member-proof-5}. The same argument, with \eqref{8.4:eq-far-member-e} in place
of \eqref{8.4:eq-far-member-d}, gives
$\eta^{2}|(D_{U_{0}}\mathcal{H}^{(y)})(p)|\le KL^{-2n}\Theta(x)$.

\emph{Clause (iv).}
Let $p$ be a fine plaquette all of whose bonds lie in $E_{0}$, and let $x_{p}$ be its base site.

First let $y\in\mathcal{T}$ be real. By the hypothesis of
Definition~\ref{8.4:def-response-hypotheses} on the zone-$0$ constraint, the constraint on
$\Gamma_{0}$ reads $U(b)=V_{0}(b)$ for all $b\in E_{0}$, and the data on $\Gamma_{0}$ are not
varied. Hence $U^{y}|_{E_{0}}=V_{0}$, and
$\operatorname{tr}U^{y}(\partial p)=\operatorname{tr}V_{0}(\partial p)$. By
\eqref{8.4:eq-far-member-c} and the gauge covariance of holonomies,
\begin{equation*}
\mathcal{W}^{(y)}(\partial p)=u_{y}(x_{p})^{-1}\,U^{y}(\partial p)\,u_{y}(x_{p}),
\end{equation*}
or the inverse conjugation, according to the convention for $U\mapsto U^{u}$. In either case
$\operatorname{tr}\mathcal{W}^{(y)}(\partial p)=\operatorname{tr}V_{0}(\partial p)$ for real $y$.

Now consider the function
$y\mapsto\operatorname{tr}\mathcal{W}^{(y)}(\partial p)-\operatorname{tr}V_{0}(\partial p)$ on
$\mathcal{T}$. It is holomorphic, for two reasons:
\begin{itemize}
\item each $\mathcal{H}^{(y)}(b)$ is holomorphic by clause (i);
\item the exponential, products and inverses in $\mathrm{GL}(N,\mathbb{C})$, and the trace, are
holomorphic.
\end{itemize}
It vanishes on the real slice. By the identity principle on the real slice,
\eqref{8.4:eq-elementary-facts-a} of Lemma~\ref{8.4:lem-elementary-facts}, it vanishes on all of
$\mathcal{T}$.

\emph{Clause (v).}
Let $Y_{R}'<Y_{R}$ and $Y_{I}'<Y_{I}$, and put $r=\min(Y_{R}-Y_{R}',Y_{I}-Y_{I}')$. The function
$y\mapsto\eta\,\partial_{b,\zeta}\mathcal{H}^{(y)}(b_{x})$ is holomorphic on $\mathcal{T}$, being
a derivative of the holomorphic function $y\mapsto\mathcal{H}^{(y)}(b_{x})$. By
\eqref{8.4:eq-far-member-d} it is bounded on $\mathcal{T}$ by $Kg_{j}L^{-n}E(x,b)$. The Cauchy
estimate \eqref{8.4:eq-elementary-facts-b} of Lemma~\ref{8.4:lem-elementary-facts} gives, for
$y\in\mathcal{T}(Y_{R}',Y_{I}')$, $b'\in I$ and $|\zeta'|\le1$,
\begin{equation*}
\bigl|\partial_{b',\zeta'}\partial_{b,\zeta}\,\eta\mathcal{H}^{(y)}(b_{x})\bigr|
\le\frac{Kg_{j}L^{-n}E(x,b)}{r}.
\end{equation*}
Mixed derivatives of a holomorphic function commute, so
$\partial_{b',\zeta'}\partial_{b,\zeta}=\partial_{b,\zeta}\partial_{b',\zeta'}$ on
$\mathcal{H}^{(y)}(b_{x})$. The same argument applied to
$y\mapsto\eta\,\partial_{b',\zeta'}\mathcal{H}^{(y)}(b_{x})$ therefore gives the bound with
$E(x,b')$ in place of $E(x,b)$. Together the two bounds give the minimum in
\eqref{8.4:eq-far-member-f}. The version with $\eta^{2}D_{U_{0}}$ and $L^{-2n}$ is proved
identically, using \eqref{8.4:eq-far-member-e} in place of \eqref{8.4:eq-far-member-d}.

\emph{Clause (vi).}
Let $F$ be as in clause (vi). For real $y$, by \eqref{8.4:eq-far-member-c} and the invariance of
$F$ under $G$-valued gauge transformations,
\begin{equation*}
F\bigl(U^{y}|_{\tilde{X}}\bigr)
=F\bigl((\mathcal{W}^{(y)})^{u_{y}}|_{\tilde{X}}\bigr)
=F\bigl(\mathcal{W}^{(y)}|_{\tilde{X}}\bigr).
\end{equation*}
On $\mathcal{T}$ we have
\begin{equation*}
F\bigl(\mathcal{W}^{(y)}|_{\tilde{X}}\bigr)
=F\bigl((\exp(i\eta\mathcal{H}^{(y)})U_{0})|_{\tilde{X}}\bigr).
\end{equation*}
The function $y\mapsto F(\mathcal{W}^{(y)}|_{\tilde{X}})$ is therefore the composition of two
holomorphic maps:
\begin{itemize}
\item $y\mapsto\mathcal{H}^{(y)}|_{\tilde{X}}$, holomorphic by clause (i);
\item $Z\mapsto F((\exp(i\eta Z)U_{0})|_{\tilde{X}})$, holomorphic, by the hypothesis of clause
(vi), on an open set containing the values of the first map.
\end{itemize}
Hence it is holomorphic on $\mathcal{T}$. It extends the real function
$y\mapsto F(U^{y}|_{\tilde{X}})$. If two holomorphic functions on $\mathcal{T}$ both extend
this function, their difference vanishes on the real slice, hence on $\mathcal{T}$ by
\eqref{8.4:eq-elementary-facts-a}. This proves uniqueness.

\emph{Where the clean event enters.}
The clean event enters the far-member lemma in two places only.
\begin{itemize}
\item It enters through the clean-event hypothesis of Lemma~\ref{8.4:lem-far-member}: every top
bond within top distance $4L+2$ of $I$ has both endpoints in $\Omega_{j+1}$. The box bonds and
their $(4L+2)$-neighbourhood therefore carry the step-$j$ fluctuation field, so that the
variation is a top-zone multiplier, and $\operatorname{supp}B(y)\subset S(I)$ consists of bonds of
the top zone $\Gamma_{j}$.
\item It enters through the regularity hypothesis $(7^{\flat})(\bar\varepsilon)$ of the same
lemma, which holds for every term of the expansion by the lemma asserting that every such term
satisfies the restricted regularity hypothesis $(7^{\flat})(\bar\varepsilon)$.
\end{itemize}
The constants obtained above do not depend on the history $h$. The conclusions of
Lemma~\ref{8.4:lem-far-member} therefore hold uniformly over clean histories. This completes the
proof of Lemma~\ref{8.4:lem-far-member}.
\end{proof}

\begin{remark}[The response in top-relative units and the age factor]\label{8.4:rem-age-factor}
Let $I$ be the set of box bonds and $\mathcal{T}=\mathcal{T}(Y_{R},Y_{I})$ the tube of Lemma~\ref{8.4:lem-far-member}, and let $y=(y_{b})_{b\in I}\in\mathcal{T}$. Let $x$ be a site at or near a large-field region live at step $j$. Let $U(\mathcal{V}(y\oplus 0))$ be the background determined by the top datum $\mathcal{V}(y\oplus 0)$, which, when the varied data of Lemma~\ref{8.4:lem-far-member} are the top datum $\mathcal{V}(y\oplus 0)$, is the varied minimiser $U^{y}$ of that lemma. Consider the following requirement on how $U(\mathcal{V}(y\oplus 0))$ at $x$ depends on the box variable $y_{b}$, $b\in I$:
\begin{itemize}
\item[(a)] the dependence is bounded by $C e^{-c\operatorname{dist}_{j}(b,x)}$, with $C$ not exponentially large in $\operatorname{dist}_{j}(b,x)$;
\item[(b)] it is holomorphic in $y$ on the tube $\mathcal{T}$;
\item[(c)] it holds uniformly over clean histories.
\end{itemize}
Clauses (i)--(iii), (v) and (vi) of the far-member lemma (Lemma~\ref{8.4:lem-far-member}) yield this requirement with
\[
c=\delta_{0}c_{d}/8 .
\]
They yield it in the following form:
\begin{itemize}
\item in the fixed-gauge representation of clause~(i) of that lemma;
\item for the gauge-invariant readings of the background, which are all that gauge-invariant members and plaquette traces see, through clause~(vi) and the corollary on plaquette traces;
\item in the natural unit of the zone of $x$, namely $L^{-n}$ per bond and $L^{-2n}$ per plaquette, where $n=n(x)$;
\item uniformly in the history $h$.
\end{itemize}
The lemma differs from the requirement in three respects.
\begin{itemize}
\item[$(\alpha)$] It is a statement about $\mathcal{H}^{(y)}$, the Landau-gauge response $1$-form. It is not a statement about the $G^{c}$-valued bond variable $U(\mathcal{V}(y\oplus 0))(x)$ in an unspecified gauge.
\item[$(\beta)$] Measured in top-relative units, the zone-$n$ response $L^{-2n}=L^{2(j-n)}\cdot L^{-2j}$ is larger than the top-zone unit $L^{-2j}$, the plaquette unit $L^{-2n}$ at zone index $n=j$, by the factor $L^{2(j-n)}$. This factor is not exponentially large in the distance, but it is exponentially large in the age $j-n$ of the live region. It is compensated only through the lower bound
\begin{equation}\label{8.4:eq-age-factor-1}
\rho_{\mathrm{sc}}\;\ge\;\tfrac12\bigl(\operatorname{dist}_{j}+M(j-1-n)_{+}\bigr),
\end{equation}
where $(\,\cdot\,)_{+}$ denotes the positive part.
This bound is available under the scaled distance hypothesis of Lemma~\ref{8.4:lem-far-member}, in which the distance argument $\rho$ of the weight $E(x,b)$ of that lemma is $\rho_{\mathrm{sc}}$, together with the admissibility hypothesis assumed in the corollaries on frozen-coupling members and on generic far members; the compensation is carried out in those corollaries.
\item[$(\gamma)$] Under the scaled distance hypothesis the decay is in $\rho_{\mathrm{sc}}$, and $\rho_{\mathrm{sc}}\ge\operatorname{dist}_{j}$. This is strictly stronger than the requirement.
\end{itemize}
\end{remark}

\begin{proof}
We check each assertion against the clauses of Lemma~\ref{8.4:lem-far-member}, in the order stated.

\emph{Representation and holomorphy.} By clause~(i) of Lemma~\ref{8.4:lem-far-member}, the map $y\mapsto\mathcal{H}^{(y)}$ is holomorphic on $\mathcal{T}$ with $\mathcal{H}^{(0)}=0$, and for real $y\in\mathcal{T}$ there is an $\mathrm{SU}(N)$-valued gauge transformation $u_{y}$ with $U^{y}=(W^{(y)})^{u_{y}}$, where $W^{(y)}=\exp(i\eta\mathcal{H}^{(y)})\,U_{0}$ bondwise, the exponential factor standing to the left of $U_{0}$ on each fine bond, and $\eta$ is the fine lattice spacing. In the fixed-gauge representative $W^{(y)}$ of clause~(i) the dependence on $y$ is through $\mathcal{H}^{(y)}$ alone. Hence item~(b), read in that representation, is the holomorphy of $\mathcal{H}^{(y)}$.

\emph{Decay.} For every fine bond $b_{x}$ with initial site $x$ and every fine plaquette $p$ with base site $x$, clause~(ii) bounds $\eta|\partial_{b,\zeta}\mathcal{H}^{(y)}(b_{x})|$ by $K g_{j}L^{-n}E(x,b)$, where $K$ is the response constant of Lemma~\ref{8.4:lem-far-member}, for every $y\in\mathcal{T}$ and every $\zeta\in\mathfrak{g}^{c}$ with $|\zeta|\le1$. It bounds $\eta^{2}|(D_{U_{0}}\partial_{b,\zeta}\mathcal{H}^{(y)})(p)|$ by $Kg_{j}L^{-2n}E(x,b)$. Under the unscaled distance hypothesis of Lemma~\ref{8.4:lem-far-member}, $\rho=\operatorname{dist}_{j}$, and the weight $E(x,b)$ carries the decay rate $\delta_{0}c_{d}/8$ in $\operatorname{dist}_{j}(b,x)$; under the scaled distance hypothesis $\rho=\rho_{\mathrm{sc}}$, see point~$(\gamma)$ below. The prefactor $Kg_{j}$ does not involve $\operatorname{dist}_{j}(b,x)$. The constant $K$ depends on $d$, $L$, $N$ and Ba{\l}aban's constants only, and not on $j$, $n$, $h$, $I$, the clean collar width $W$, or the volume. This gives item~(a) in the zone units $L^{-n}$ and $L^{-2n}$, and item~(c), since the constants do not see $h$. Clause~(iii) bounds the size of $\mathcal{H}^{(y)}$ and of $D_{U_{0}}\mathcal{H}^{(y)}$ on $\mathcal{T}$ in the same zone units,
\begin{align*}
\eta|\mathcal{H}^{(y)}(b_{x})|&\le K L^{-n}\Theta(x),\\
\eta^{2}|(D_{U_{0}}\mathcal{H}^{(y)})(p)|&\le K L^{-2n}\Theta(x),
\end{align*}
where $\Theta(x)$ is the response size of Lemma~\ref{8.4:lem-far-member}; these bounds carry neither the factor $g_{j}$ nor the weight $E(x,b)$. Clause~(v) bounds the mixed second derivatives $\partial_{b',\zeta'}\partial_{b,\zeta}\mathcal{H}^{(y)}$, for $b,b'\in I$ and $|\zeta|,|\zeta'|\le1$, on every smaller tube $\mathcal{T}(Y_{R}',Y_{I}')$ with $Y_{R}'<Y_{R}$ and $Y_{I}'<Y_{I}$, in the same zone units, with prefactor $Kg_{j}/r$, where $r=\min(Y_{R}-Y_{R}',Y_{I}-Y_{I}')$, and with the weight $\min\{E(x,b),E(x,b')\}$.

\emph{Gauge-invariant readings.} Let $F$ be a function of configurations on a bond set $\tilde X$, invariant under $\mathrm{SU}(N)$-valued gauge transformations and satisfying the holomorphy hypothesis of clause~(vi); writing $\cdot\,|_{\tilde X}$ for the restriction of a configuration to $\tilde X$, that clause shows that $y\mapsto F(W^{(y)}|_{\tilde X})$ is holomorphic on $\mathcal{T}$ and is the unique holomorphic extension of $y\mapsto F(U^{y}|_{\tilde X})$. Together with the corollary on plaquette traces, this covers what gauge-invariant members and plaquette traces see.

\emph{Point $(\alpha)$.} The bounds of clause~(ii) concern $\mathcal{H}^{(y)}$. By clause~(i) the varied minimiser $U^{y}$ coincides with $W^{(y)}$ only after the gauge transformation $u_{y}$, and only for real $y$.

\emph{Point $(\beta)$.} By the identity in $(\beta)$, the zone-$n$ response exceeds the top-zone unit by the factor $L^{2(j-n)}=e^{2(j-n)\log L}$, which depends on $j$ and $n$ only, not on $\operatorname{dist}_{j}(b,x)$. It is therefore exponentially large in the age $j-n$ and not in the distance. The only compensation is \eqref{8.4:eq-age-factor-1}, under the hypotheses named in $(\beta)$.

\emph{Point $(\gamma)$.} Under the scaled distance hypothesis, $\rho=\rho_{\mathrm{sc}}$ and $\rho_{\mathrm{sc}}\ge\operatorname{dist}_{j}$. A bound decaying in $\rho_{\mathrm{sc}}$ therefore implies the one decaying in $\operatorname{dist}_{j}$.
\end{proof}

\begin{remark}[The far-member lemma for the large-field determining sets]\label{8.4:rem-large-field-sets}
Ba{\l}aban's large-field operation $\mathbf{R}$ replaces the sequence of admissible domains
$\{\Omega_{n}\}$ by a sequence of domains $\{\Omega''_{n}\}$, and the determining set
$\mathbf{B}_{j}$ by the determining sets $\mathbb{B}''_{k}$, $\mathbb{B}^{(n)}_{k}$ and
$\mathbb{B}^{(n)}_{k}(Z)$ of \cite[(1.12)--(1.20)]{B-CMP122a}, for the level $k$, the index $n$ and
large-field regions $Z$ as there. Let $\mathbb{B}$ denote any one of these determining sets whose
top zone contains $S(I)$. Then Lemma~\ref{8.4:lem-far-member}, the far-member lemma, holds with
$\mathbf{B}_{j}$ replaced by $\mathbb{B}$, that is, for the minimiser $U(\mathbb{B},\cdot)$ in
place of $U(\mathbf{B}_{j},\cdot)$, with the same constants.
\end{remark}

\begin{proof}
That the domains underlying these determining sets are admissible is the following sentence of
\cite[p.~179]{B-CMP122a}:
\begin{quote}
the complements of these sets form an admissible sequence of domains
\end{quote}
This sentence is printed without proof in \cite[p.~179]{B-CMP122a}; cited as printed.

The data carried by these determining sets are of two kinds. The first kind consists of
restrictions of the configuration $V$. The second consists of block averages
$M^{\cdot}(Q_{k}^{s*}V_{k})$ built from the block field $V_{k}$, of the orders prescribed in
\cite[(1.12)--(1.20)]{B-CMP122a}. Restrictions of a regular configuration are regular. For the
block averages, \cite{B-CMP102b} states that averages of regular configurations are again
regular: they satisfy the regularity condition whose restricted form is the regularity hypothesis
$(7^{\flat})(\bar\varepsilon)$ of Lemma~\ref{8.4:lem-far-member}, with a threshold
$\varepsilon_{1}=O(\varepsilon_{0})$.
Consequently the data of $\mathbb{B}$ are regular in this sense whenever $V_{k}$ is regular.

Thus $\mathbb{B}$ is a determining set over an admissible sequence of domains, with regular data,
whose top zone contains $S(I)$. These are the properties of
$\mathbf{B}_{j}$ on which the statement and the proof of Lemma~\ref{8.4:lem-far-member} draw, and
the lemma with its proof applies verbatim to $U(\mathbb{B},\cdot)$. The constants are the same,
by the uniformity statement for the constants of Lemma~\ref{8.4:lem-far-member}; in particular,
as stated in clause \textup{(ii)} of Lemma~\ref{8.4:lem-far-member}, the constant $K$ of the
lemma depends only on $d$, $L$, $N$ and Ba{\l}aban's constants, and not on $j$, $n$, $h$, $I$,
$W$ or the volume.
\end{proof}

\begin{corollary}[Plaquette traces of the responding background]\label{8.4:cor-plaquette-traces}
Let $\mathcal{T}:=\mathcal{T}(Y_{R},Y_{I})$ be the tube of complex box variables of
Lemma~\ref{8.4:lem-far-member}. Assume the three hypotheses of that lemma (the restricted
regularity of the base data with threshold $\bar\varepsilon$, the box hypothesis, and the smallness
of the tube), and assume that the following two statements hold:
\begin{itemize}
\item the far-member lemma, Lemma~\ref{8.4:lem-far-member}, that is, its conclusions
\textup{(i)}--\textup{(vi)} for the tube $\mathcal{T}$;
\item the regularity bound on the base background of
Definition~\ref{8.4:def-far-member-hypotheses}.
\end{itemize}
Let $p$ be a fine plaquette with a bond in $E_{1}$, let $x$ be its base site and $n:=n(x)$, and
assume
\begin{equation}\label{8.4:eq-plaquette-traces-a}
\Theta(x)\le\theta_{*}:=\bigl(64K(1+C_{\mathrm{reg}}\bar\varepsilon)\bigr)^{-1},
\end{equation}
with $\Theta(x)$ as in \eqref{8.4:eq-decay-weight-b} and $K$ the constant there, which for
plaquettes meeting a zone interface is multiplied by $L$ (see the proof). Put
\begin{equation*}
t_{p}(y):=\operatorname{tr}W^{(y)}(\partial p),\qquad
\bar t_{p}(y):=\operatorname{tr}W^{(y)}(\partial p)^{-1},
\end{equation*}
with $W^{(y)}$ as in \eqref{8.4:eq-far-member-c}. Then $t_{p}$ and $\bar t_{p}$ are holomorphic
on $\mathcal{T}$ and extend, from real $y$, the function $\operatorname{tr}U^{y}(\partial p)$
and its complex conjugate. For $y\in\mathcal{T}$, $b\in I$ and $\zeta\in\mathfrak{g}^{c}$ with
$|\zeta|\le1$,
\begin{gather}
\bigl|W^{(y)}(\partial p)U_{0}(\partial p)^{-1}-1\bigr|\le K'L^{-2n}\Theta(x),
\qquad K':=2K(1+4C_{\mathrm{reg}}\bar\varepsilon),\label{8.4:eq-plaquette-traces-b}\\
\begin{split}
|\partial_{b,\zeta}t_{p}(y)|+|\partial_{b,\zeta}\bar t_{p}(y)|
&\le K_{T}\,g_{j}L^{-4n}E(x,b)\bigl(\bar\varepsilon+\Theta(x)\bigr),\\
K_{T}&:=16K(1+4C_{\mathrm{reg}}\bar\varepsilon)(C_{\mathrm{reg}}+2K'),
\end{split}\label{8.4:eq-plaquette-traces-c}
\end{gather}
with $E(x,b)$ as in \eqref{8.4:eq-decay-weight-a}. Moreover, let $Y_{R}'<Y_{R}$,
$Y_{I}'<Y_{I}$ and $r_{\mathcal{T}}:=\min(Y_{R}-Y_{R}',\,Y_{I}-Y_{I}')$. Then for all
$y\in\mathcal{T}(Y_{R}',Y_{I}')$, all $b,b'\in I$ and all $\zeta,\zeta'\in\mathfrak{g}^{c}$ with
$|\zeta|,|\zeta'|\le1$, each of the mixed second derivatives
$\partial_{b',\zeta'}\partial_{b,\zeta}t_{p}(y)$ and
$\partial_{b',\zeta'}\partial_{b,\zeta}\bar t_{p}(y)$ is bounded in modulus by
\begin{equation}\label{8.4:eq-plaquette-traces-1}
\frac{K_{T}\,g_{j}}{r_{\mathcal{T}}}\,L^{-4n}\min\{E(x,b),E(x,b')\}\bigl(\bar\varepsilon+\Theta(x)\bigr).
\end{equation}
\end{corollary}

\begin{proof}
Throughout, $|\cdot|$ is the operator norm on $N\times N$ matrices, in which the facts of
Lemma~\ref{8.4:lem-elementary-facts} are applied; we use that conjugation by a unitary matrix is
isometric for it and that $|\operatorname{tr}(AB)|\le|A|\,|B|$, where
$\operatorname{tr}=\operatorname{Tr}/N$. All references to parts (i)--(vi) and to
\eqref{8.4:eq-far-member-c}--\eqref{8.4:eq-far-member-e} are to the far-member lemma,
Lemma~\ref{8.4:lem-far-member}, which is among the hypotheses.

We fix $p$, $x$, $n$ as in the statement and $y\in\mathcal{T}$. We label the bonds of $p$ as
$b_{1},\dots,b_{4}$ so that, for every configuration $V$ on the fine bonds,
\begin{equation*}
V(\partial p)=V(b_{1})V(b_{2})V(b_{3})^{-1}V(b_{4})^{-1};
\end{equation*}
thus $b_{1}$ and $b_{4}$ have initial site $x$, and the initial sites $x_{2}$, $x_{3}$ of
$b_{2}$, $b_{3}$ differ from $x$ by one fine step. We put $x_{1}=x_{4}:=x$,
\begin{equation*}
U_{0,i}:=U_{0}(b_{i}),\qquad \mathcal{H}_{i}:=\mathcal{H}^{(y)}(b_{i})\qquad(i=1,\dots,4),
\qquad \xi:=KL^{-n}\Theta(x).
\end{equation*}
The matrices $U_{0,i}$ are unitary and do not depend on $y$. Part (iii) gives
$\eta|\mathcal{H}_{i}|\le\xi$ for each $i$.

The bounds of part (iii) and \eqref{8.4:eq-far-member-d} for a bond are stated at the initial
site of that bond; the plaquette bounds of part (iii) and \eqref{8.4:eq-far-member-e} are read at
the base site $x$ of $p$. The sites $x_{2}$, $x_{3}$ lie in the zone of $x$ unless $p$ meets a zone
interface, in which case they may lie in a neighbouring zone. To avoid case distinctions we apply
the bond bounds at the initial site $x_{i}$ of each $b_{i}$, and when $n(x_{i})=n-1$ we use
$L^{-n(x_{i})}=L\cdot L^{-n}$ (when $n(x_{i})=n+1$ one has $L^{-n(x_{i})}\le L^{-n}$). This costs a
factor $L$, which for plaquettes at a zone interface is absorbed into $K$; we keep writing $K$.

The regularity bound on the base background of Definition~\ref{8.4:def-far-member-hypotheses},
applied at the plaquette $p$, which has a bond in $E_{1}$, bounds $|U_{0}(\partial p)-1|$ by
$C_{\mathrm{reg}}\varepsilon_{n(p)}L^{-2n(p)}$; we use it in the form
\begin{equation}\label{8.4:eq-plaquette-traces-2}
|U_{0}(\partial p)-1|\le C_{\mathrm{reg}}\bar\varepsilon L^{-2n},
\end{equation}
with the threshold $\varepsilon_{n(p)}$ replaced by the uniform threshold $\bar\varepsilon$ and the
zone index $n(p)$ by $n$.
Since $U_{0}(\partial p)$ is unitary, $|U_{0}(\partial p)^{-1}-1|=|U_{0}(\partial p)-1|$. We also
use, for a unitary matrix $\mathsf{u}$ and a matrix $\mathsf{v}$,
\begin{equation}\label{8.4:eq-plaquette-traces-3}
|\operatorname{Ad}_{\mathsf{u}}\mathsf{v}-\mathsf{v}|\le2|\mathsf{u}-1|\,|\mathsf{v}|,
\end{equation}
which follows from
$\operatorname{Ad}_{\mathsf{u}}\mathsf{v}-\mathsf{v}=(\mathsf{u}-1)\mathsf{v}\mathsf{u}^{-1}
+\mathsf{v}(\mathsf{u}^{-1}-1)$ and
$|\mathsf{u}^{-1}-1|=|\mathsf{u}^{-1}(1-\mathsf{u})|=|\mathsf{u}-1|$. Finally $n\ge0$ and $L\ge1$,
so $L^{-n}\le1$.

\smallskip
\noindent\textit{Holomorphy and extension.} Let $F$ be either
$V\mapsto\operatorname{tr}V(\partial p)$ or $V\mapsto\operatorname{tr}V(\partial p)^{-1}$,
functions of configurations on the set of the four bonds of $p$. Both are
invariant under $G$-valued gauge transformations, since $V(\partial p)$ is the product of the bond
variables around the closed loop $\partial p$ and a gauge transformation changes it by a conjugation.
For $\psi$ a $\mathfrak{g}^{c}$-valued function on the four bonds of $p$, the map
$\psi\mapsto F(\exp(i\eta\psi)U_{0})$, with $U_{0}$ restricted to these bonds, is holomorphic on
the whole space of such functions: the entries of $\exp(i\eta\psi(b_{i}))$ are entire functions of
$\psi$, and
$(\exp(i\eta\psi(b_{i}))U_{0,i})^{-1}=U_{0,i}^{-1}\exp(-i\eta\psi(b_{i}))$. Part (vi), applied with
the set of the four bonds of $p$ as bond set, therefore shows that
$t_{p}$ and $\bar t_{p}$ are holomorphic on $\mathcal{T}$ and are the unique holomorphic extensions
of $y\mapsto\operatorname{tr}U^{y}(\partial p)$ and $y\mapsto\operatorname{tr}U^{y}(\partial p)^{-1}$
from the real points of $\mathcal{T}$. At real $y$ the matrix $U^{y}(\partial p)$ is unitary, so
$\operatorname{tr}U^{y}(\partial p)^{-1}$ is the complex conjugate of
$\operatorname{tr}U^{y}(\partial p)$.

\smallskip
\noindent\textit{The conjugation identity.} By \eqref{8.4:eq-far-member-c},
$W^{(y)}(b_{i})=e^{i\eta\mathcal{H}_{i}}U_{0,i}$, hence
\begin{equation*}
W^{(y)}(\partial p)=e^{i\eta\mathcal{H}_{1}}U_{0,1}\,e^{i\eta\mathcal{H}_{2}}U_{0,2}\,
U_{0,3}^{-1}e^{-i\eta\mathcal{H}_{3}}\,U_{0,4}^{-1}e^{-i\eta\mathcal{H}_{4}}.
\end{equation*}
We move the exponential factors to the left through the unitary products preceding them, each
step being $AE=\operatorname{Ad}_{A}(E)\,A$ with $\operatorname{Ad}_{A}(e^{Z})=e^{\operatorname{Ad}_{A}Z}$;
successively,
\begin{align*}
U_{0,1}e^{i\eta\mathcal{H}_{2}}
&=\operatorname{Ad}_{U_{0,1}}(e^{i\eta\mathcal{H}_{2}})\,U_{0,1},\\
U_{0,1}U_{0,2}U_{0,3}^{-1}e^{-i\eta\mathcal{H}_{3}}
&=\operatorname{Ad}_{U_{0,1}U_{0,2}U_{0,3}^{-1}}(e^{-i\eta\mathcal{H}_{3}})\,
U_{0,1}U_{0,2}U_{0,3}^{-1},\\
U_{0}(\partial p)e^{-i\eta\mathcal{H}_{4}}
&=\operatorname{Ad}_{U_{0}(\partial p)}(e^{-i\eta\mathcal{H}_{4}})\,U_{0}(\partial p),
\end{align*}
where in the last step we used $U_{0,1}U_{0,2}U_{0,3}^{-1}U_{0,4}^{-1}=U_{0}(\partial p)$. This gives
\begin{equation}\label{8.4:eq-plaquette-traces-4}
W^{(y)}(\partial p)=e^{X_{1}}e^{X_{2}}e^{X_{3}}e^{X_{4}}\,U_{0}(\partial p),
\end{equation}
where
\begin{align*}
X_{1}&:=i\eta\mathcal{H}_{1}, &
X_{2}&:=i\eta\operatorname{Ad}_{U_{0,1}}\mathcal{H}_{2},\\
X_{3}&:=-i\eta\operatorname{Ad}_{U_{0,1}U_{0,2}U_{0,3}^{-1}}\mathcal{H}_{3}, &
X_{4}&:=-i\eta\operatorname{Ad}_{U_{0}(\partial p)}\mathcal{H}_{4}.
\end{align*}
Since $U_{0,1}U_{0,2}U_{0,3}^{-1}=U_{0}(\partial p)U_{0,4}$, we have
$\operatorname{Ad}_{U_{0,1}U_{0,2}U_{0,3}^{-1}}=\operatorname{Ad}_{U_{0}(\partial p)}\operatorname{Ad}_{U_{0,4}}$,
and therefore
\begin{equation*}
\sum_{i=1}^{4}X_{i}
=i\eta\bigl(\mathcal{H}_{1}+\operatorname{Ad}_{U_{0,1}}\mathcal{H}_{2}
-\operatorname{Ad}_{U_{0,4}}\mathcal{H}_{3}-\mathcal{H}_{4}\bigr)
-i\eta\bigl(\operatorname{Ad}_{U_{0}(\partial p)}-1\bigr)
\bigl(\operatorname{Ad}_{U_{0,4}}\mathcal{H}_{3}+\mathcal{H}_{4}\bigr).
\end{equation*}
By the definition of the covariant curl in \eqref{8.4:eq-box-variation-c}, the first term is
$i\eta^{2}(D_{U_{0}}\mathcal{H}^{(y)})(p)$, so that
\begin{equation}\label{8.4:eq-plaquette-traces-5}
\sum_{i=1}^{4}X_{i}=i\eta^{2}(D_{U_{0}}\mathcal{H}^{(y)})(p)
-i\eta\bigl(\operatorname{Ad}_{U_{0}(\partial p)}-1\bigr)
\bigl(\operatorname{Ad}_{U_{0,4}}\mathcal{H}_{3}+\mathcal{H}_{4}\bigr).
\end{equation}

\smallskip
\noindent\textit{Proof of \eqref{8.4:eq-plaquette-traces-b}.} Put
\begin{equation*}
\Pi:=W^{(y)}(\partial p)U_{0}(\partial p)^{-1}=e^{X_{1}}e^{X_{2}}e^{X_{3}}e^{X_{4}},
\qquad \varsigma:=\sum_{i=1}^{4}|X_{i}|,
\end{equation*}
the second expression for $\Pi$ by \eqref{8.4:eq-plaquette-traces-4}. Conjugation by unitaries is
isometric, so $|X_{i}|=\eta|\mathcal{H}_{i}|\le\xi$, and $\varsigma\le4\xi\le4K\Theta(x)$. By
\eqref{8.4:eq-plaquette-traces-a},
\begin{equation*}
4K\Theta(x)\le4K\theta_{*}=\bigl(16(1+C_{\mathrm{reg}}\bar\varepsilon)\bigr)^{-1}\le\tfrac{1}{16};
\end{equation*}
hence $\varsigma\le1/16$. By \eqref{8.4:eq-plaquette-traces-5}, the plaquette bound of part (iii),
\eqref{8.4:eq-plaquette-traces-3} and \eqref{8.4:eq-plaquette-traces-2},
\begin{align*}
\Bigl|\sum_{i}X_{i}\Bigr|
&\le\eta^{2}|(D_{U_{0}}\mathcal{H}^{(y)})(p)|
+2|U_{0}(\partial p)-1|\,\eta\bigl(|\mathcal{H}_{3}|+|\mathcal{H}_{4}|\bigr)\\
&\le KL^{-2n}\Theta(x)+2C_{\mathrm{reg}}\bar\varepsilon L^{-2n}\cdot2\xi
\le K(1+4C_{\mathrm{reg}}\bar\varepsilon)L^{-2n}\Theta(x),
\end{align*}
the last step because $2\xi=2KL^{-n}\Theta(x)\le2K\Theta(x)$. The exponential estimates
\eqref{8.4:eq-elementary-facts-c} of Lemma~\ref{8.4:lem-elementary-facts} apply since
$\varsigma\le1$, and $2\varsigma^{2}\le32\xi^{2}=32K^{2}L^{-2n}\Theta(x)^{2}$; thus
\begin{align*}
|\Pi-1|&\le\Bigl|\sum_{i}X_{i}\Bigr|+2\varsigma^{2}
\le K(1+4C_{\mathrm{reg}}\bar\varepsilon)L^{-2n}\Theta(x)+32K^{2}L^{-2n}\Theta(x)^{2}\\
&\le 2K(1+4C_{\mathrm{reg}}\bar\varepsilon)L^{-2n}\Theta(x)=K'L^{-2n}\Theta(x),
\end{align*}
using $32K\Theta(x)\le32K\theta_{*}\le1$ from \eqref{8.4:eq-plaquette-traces-a}. This is
\eqref{8.4:eq-plaquette-traces-b}. Moreover, since
$1+4C_{\mathrm{reg}}\bar\varepsilon\le4(1+C_{\mathrm{reg}}\bar\varepsilon)$,
\begin{equation*}
K'L^{-2n}\Theta(x)\le K'\theta_{*}
=\frac{2(1+4C_{\mathrm{reg}}\bar\varepsilon)}{64(1+C_{\mathrm{reg}}\bar\varepsilon)}\le\frac18,
\end{equation*}
so $|\Pi-1|\le1/4$.

\smallskip
\noindent\textit{The first derivative of $\Pi$.} Fix $b\in I$ and $\zeta\in\mathfrak{g}^{c}$
with $|\zeta|\le1$, and write $\zeta e_{b}$ for the element of $(\mathfrak{g}^{c})^{I}$ equal to
$\zeta$ at $b$ and to $0$ at the other elements of $I$. All functions of $y$ involved are holomorphic
on $\mathcal{T}$, so $\partial_{b,\zeta}$ at $y$ is the derivative at $\sigma=0$ along the real
segment $\sigma\mapsto y+\sigma\zeta e_{b}$, $\sigma$ real and small; we denote it by a prime. Since
the $U_{0,i}$ do not depend on $y$,
\begin{align*}
X_{1}'&=i\eta\,\partial_{b,\zeta}\mathcal{H}^{(y)}(b_{1}), &
X_{2}'&=i\eta\operatorname{Ad}_{U_{0,1}}\partial_{b,\zeta}\mathcal{H}^{(y)}(b_{2}),\\
X_{3}'&=-i\eta\operatorname{Ad}_{U_{0,1}U_{0,2}U_{0,3}^{-1}}\partial_{b,\zeta}\mathcal{H}^{(y)}(b_{3}),
&
X_{4}'&=-i\eta\operatorname{Ad}_{U_{0}(\partial p)}\partial_{b,\zeta}\mathcal{H}^{(y)}(b_{4}),
\end{align*}
each a fixed unitary conjugation of $\pm i\eta\,\partial_{b,\zeta}\mathcal{H}^{(y)}(b_{i})$. By
\eqref{8.4:eq-far-member-d},
\begin{equation*}
\sum_{i}|X_{i}'|\le4Kg_{j}L^{-n}E(x,b).
\end{equation*}
Differentiating \eqref{8.4:eq-plaquette-traces-5}, which is linear in $\mathcal{H}^{(y)}$,
\begin{equation*}
\sum_{i}X_{i}'=i\eta^{2}(D_{U_{0}}\partial_{b,\zeta}\mathcal{H}^{(y)})(p)
-i\eta\bigl(\operatorname{Ad}_{U_{0}(\partial p)}-1\bigr)
\bigl(\operatorname{Ad}_{U_{0,4}}\partial_{b,\zeta}\mathcal{H}^{(y)}(b_{3})
+\partial_{b,\zeta}\mathcal{H}^{(y)}(b_{4})\bigr),
\end{equation*}
and by \eqref{8.4:eq-far-member-e}, \eqref{8.4:eq-plaquette-traces-3},
\eqref{8.4:eq-plaquette-traces-2} and \eqref{8.4:eq-far-member-d},
\begin{align*}
\Bigl|\sum_{i}X_{i}'\Bigr|
&\le Kg_{j}L^{-2n}E(x,b)+4C_{\mathrm{reg}}\bar\varepsilon L^{-2n}\cdot Kg_{j}L^{-n}E(x,b)\\
&\le K(1+4C_{\mathrm{reg}}\bar\varepsilon)g_{j}L^{-2n}E(x,b).
\end{align*}
By \eqref{8.4:eq-elementary-facts-c}, $|\Pi'-\sum_{i}X_{i}'|\le3\varsigma\sum_{i}|X_{i}'|$, and
$3\varsigma\le12\xi$; hence
\begin{align*}
|\Pi'|&\le\Bigl|\sum_{i}X_{i}'\Bigr|+3\varsigma\sum_{i}|X_{i}'|
\le K(1+4C_{\mathrm{reg}}\bar\varepsilon)g_{j}L^{-2n}E(x,b)+12\xi\cdot4Kg_{j}L^{-n}E(x,b)\\
&\le K(1+4C_{\mathrm{reg}}\bar\varepsilon)g_{j}L^{-2n}E(x,b)
+48K^{2}\Theta(x)g_{j}L^{-2n}E(x,b)\\
&\le2K(1+4C_{\mathrm{reg}}\bar\varepsilon)g_{j}L^{-2n}E(x,b)=:K''g_{j}L^{-2n}E(x,b),
\end{align*}
using
\begin{equation*}
48K\Theta(x)\le48K\theta_{*}\le1\le1+4C_{\mathrm{reg}}\bar\varepsilon,
\end{equation*}
which follows from \eqref{8.4:eq-plaquette-traces-a}.

\smallskip
\noindent\textit{Tracelessness and the logarithm.} Each $X_{i}$ is traceless: $\mathcal{H}^{(y)}$
takes values in $\mathfrak{g}^{c}=\mathfrak{sl}(N,\mathbb{C})$ by part (i), and conjugation
preserves the trace. The tracelessness of $\mathcal{H}^{(y)}$ also follows from
\eqref{8.4:eq-far-member-c}: at real $y$ both sides of it have determinant one, since $U^{y}$ and
$U_{0}$ are $G$-valued and a gauge transformation does not change determinants, so
$e^{i\eta\operatorname{Tr}\mathcal{H}^{(y)}}=1$ bondwise; on the real points of $\mathcal{T}$, which
form a convex set, the continuous function $y\mapsto\operatorname{Tr}\mathcal{H}^{(y)}$ therefore
takes values in $(2\pi/\eta)\mathbb{Z}$ at each bond and vanishes at $y=0$ because
$\mathcal{H}^{(0)}=0$, so it vanishes at every real point of $\mathcal{T}$; the identity principle
\eqref{8.4:eq-elementary-facts-a} of Lemma~\ref{8.4:lem-elementary-facts}, applied to the
holomorphic function $y\mapsto\operatorname{Tr}\mathcal{H}^{(y)}$ at each bond, extends this to all
of $\mathcal{T}$. Hence
$\det\Pi=\prod_{i}e^{\operatorname{Tr}X_{i}}=1$. Since $|\Pi-1|\le1/4$ at every point of
$\mathcal{T}$, the principal logarithm
$\Xi:=\log\Pi=\sum_{m\ge1}(-1)^{m+1}(\Pi-1)^{m}/m$ is defined on all of $\mathcal{T}$, the
series converging uniformly there, so that $\Xi$ is a holomorphic, in particular continuous,
function of $y$ on $\mathcal{T}$; and \eqref{8.4:eq-elementary-facts-c} gives
$|\Xi'|\le\frac43|\Pi'|$. Moreover $e^{\operatorname{Tr}\Xi}=\det e^{\Xi}=\det\Pi=1$, so the
continuous function $y\mapsto\operatorname{Tr}\Xi$ takes values in $2\pi i\mathbb{Z}$ on
$\mathcal{T}$, which is convex and hence connected. At $y=0$ we have $\mathcal{H}^{(0)}=0$, hence
$X_{i}=0$ for each $i$, $\Pi=1$ and $\Xi=0$. Therefore $\operatorname{Tr}\Xi=0$ at every point of
$\mathcal{T}$, and, as $\operatorname{Tr}\Xi$ vanishes identically along the segment, also
$\operatorname{Tr}\Xi'=0$.

\smallskip
\noindent\textit{The derivative of $t_{p}$.} By \eqref{8.4:eq-plaquette-traces-4},
$t_{p}=\operatorname{tr}(\Pi\,U_{0}(\partial p))$, and $U_{0}(\partial p)$ does not depend on
$y$, so
\begin{equation*}
\partial_{b,\zeta}t_{p}=\operatorname{tr}(\Pi'U_{0}(\partial p))
=\operatorname{tr}(\Pi')+\operatorname{tr}\bigl(\Pi'(U_{0}(\partial p)-1)\bigr).
\end{equation*}
By cyclicity of the trace, $d\operatorname{tr}e^{\Xi}=\operatorname{tr}(e^{\Xi}d\Xi)$; with
$e^{\Xi}=\Pi$ and $\operatorname{tr}\Xi'=0$,
\begin{equation*}
\operatorname{tr}(\Pi')=\operatorname{tr}(e^{\Xi}\Xi')
=\operatorname{tr}(\Xi')+\operatorname{tr}\bigl((e^{\Xi}-1)\Xi'\bigr)
=\operatorname{tr}\bigl((\Pi-1)\Xi'\bigr).
\end{equation*}
Therefore, by \eqref{8.4:eq-plaquette-traces-b}, \eqref{8.4:eq-plaquette-traces-2} and the bound
on $|\Pi'|$,
\begin{align*}
|\partial_{b,\zeta}t_{p}|&\le|\Pi-1|\cdot\tfrac43|\Pi'|+|\Pi'|\,C_{\mathrm{reg}}\bar\varepsilon L^{-2n}
\le K''g_{j}L^{-2n}E(x,b)\Bigl[\tfrac43K'L^{-2n}\Theta(x)
+C_{\mathrm{reg}}\bar\varepsilon L^{-2n}\Bigr]\\
&\le K''(C_{\mathrm{reg}}+2K')\,g_{j}L^{-4n}E(x,b)\bigl(\bar\varepsilon+\Theta(x)\bigr).
\end{align*}

\smallskip
\noindent\textit{The derivative of $\bar t_{p}$.} We have
$\bar t_{p}=\operatorname{tr}(U_{0}(\partial p)^{-1}\Pi^{-1})$ and
$(\Pi^{-1})'=-\Pi^{-1}\Pi'\Pi^{-1}$, so by cyclicity
\begin{align*}
\partial_{b,\zeta}\bar t_{p}
&=-\operatorname{tr}\bigl(\Pi'\,\Pi^{-1}U_{0}(\partial p)^{-1}\Pi^{-1}\bigr)\\
&=-\operatorname{tr}(\Pi')-\operatorname{tr}\bigl(\Pi'(\Pi^{-1}U_{0}(\partial p)^{-1}\Pi^{-1}-1)\bigr).
\end{align*}
Put $\delta_{\Pi}:=|\Pi^{-1}-1|$. From $\Pi^{-1}-1=-\Pi^{-1}(\Pi-1)$ and
$|\Pi^{-1}|\le(1-|\Pi-1|)^{-1}$ we get $\delta_{\Pi}\le\frac43|\Pi-1|\le\frac13$. For the unitary
matrix $\mathsf{u}:=U_{0}(\partial p)^{-1}$,
\begin{gather*}
\Pi^{-1}\mathsf{u}\Pi^{-1}-\mathsf{u}
=(\Pi^{-1}-1)\mathsf{u}\Pi^{-1}+\mathsf{u}(\Pi^{-1}-1),\\
|\Pi^{-1}\mathsf{u}\Pi^{-1}-\mathsf{u}|\le\delta_{\Pi}(1+\delta_{\Pi})+\delta_{\Pi}
\le3\delta_{\Pi},
\end{gather*}
and so
\begin{equation*}
|\Pi^{-1}U_{0}(\partial p)^{-1}\Pi^{-1}-1|\le|U_{0}(\partial p)^{-1}-1|+3|\Pi^{-1}-1|
\le C_{\mathrm{reg}}\bar\varepsilon L^{-2n}+4K'L^{-2n}\Theta(x).
\end{equation*}
Using $|\operatorname{tr}(\Pi')|\le\frac43|\Pi-1|\,|\Pi'|$ from the preceding paragraph,
\begin{align*}
|\partial_{b,\zeta}\bar t_{p}|
&\le|\Pi'|\Bigl[\tfrac43K'L^{-2n}\Theta(x)+C_{\mathrm{reg}}\bar\varepsilon L^{-2n}
+4K'L^{-2n}\Theta(x)\Bigr]\\
&\le K''(C_{\mathrm{reg}}+6K')\,g_{j}L^{-4n}E(x,b)\bigl(\bar\varepsilon+\Theta(x)\bigr).
\end{align*}
Adding the two bounds,
\begin{equation*}
|\partial_{b,\zeta}t_{p}|+|\partial_{b,\zeta}\bar t_{p}|
\le K''(2C_{\mathrm{reg}}+8K')\,g_{j}L^{-4n}E(x,b)\bigl(\bar\varepsilon+\Theta(x)\bigr),
\end{equation*}
and
\begin{equation*}
K''(2C_{\mathrm{reg}}+8K')=4K(1+4C_{\mathrm{reg}}\bar\varepsilon)(C_{\mathrm{reg}}+4K')\le K_{T},
\end{equation*}
because $C_{\mathrm{reg}}+4K'\le4(C_{\mathrm{reg}}+2K')$. This is
\eqref{8.4:eq-plaquette-traces-c}.

\smallskip
\noindent\textit{Mixed second derivatives.} The function $\partial_{b,\zeta}t_{p}$ is holomorphic
on $\mathcal{T}=\mathcal{T}(Y_{R},Y_{I})$, and by \eqref{8.4:eq-plaquette-traces-c} it is bounded
there by $\mathsf{a}:=K_{T}g_{j}L^{-4n}E(x,b)(\bar\varepsilon+\Theta(x))$. The Cauchy estimate
\eqref{8.4:eq-elementary-facts-b} of Lemma~\ref{8.4:lem-elementary-facts} gives
$|\partial_{b',\zeta'}\partial_{b,\zeta}t_{p}(y)|\le\mathsf{a}/r_{\mathcal{T}}$ for
$y\in\mathcal{T}(Y_{R}',Y_{I}')$
and $|\zeta'|\le1$. Since mixed derivatives of a holomorphic function commute,
$\partial_{b',\zeta'}\partial_{b,\zeta}t_{p}=\partial_{b,\zeta}\partial_{b',\zeta'}t_{p}$, and the same
argument applied to $\partial_{b',\zeta'}t_{p}$ bounds it by
$K_{T}g_{j}r_{\mathcal{T}}^{-1}L^{-4n}E(x,b')(\bar\varepsilon+\Theta(x))$. Taking the smaller of the
two bounds
gives \eqref{8.4:eq-plaquette-traces-1} for $t_{p}$. The same argument with $\bar t_{p}$ in place
of $t_{p}$ gives \eqref{8.4:eq-plaquette-traces-1} for $\bar t_{p}$.
\end{proof}

\begin{corollary}[The frozen-coupling member at live regions]\label{8.4:cor-frozen-coupling}
We work in the setting of this section: $j$ is the level, $I$ the set of level-$j$ bonds carrying
the box variables (these bonds lie in $\Omega_{j+1}$), $\mathcal{T}=\mathcal{T}(Y_{R},Y_{I})$ the
tube of complex box variables with
$Y=Y_{R}+Y_{I}$, and $\mathcal{T}(Y_{R}',Y_{I}')\subset\mathcal{T}$ and $r$ are as in
Corollary~\ref{8.4:cor-plaquette-traces}. For a fine plaquette $p$ we write $x_{p}$ for its base
site and $n(p):=n(x_{p})$ for its zone index. $E_{1}$ is the set of fine bonds meeting
$\Omega_{1}$. Assume hypotheses \textup{(H1)}--\textup{(H3)} of Notation~\ref{2.2:not-hypotheses},
and assume the following.
\begin{enumerate}
\item[(a)] The conclusions of Corollary~\ref{8.4:cor-plaquette-traces} (plaquette traces of the
responding background) hold at every fine plaquette $p$ with a bond in $E_{1}$ whose base site
satisfies \eqref{8.4:eq-plaquette-traces-a}.
\item[(b)] The frozen-coupling member $\mathcal{L}$ of the effective action at level $j$, in the
form of \cite[(2.23)--(2.24)]{B-CMP119}, the coupling difference $\Delta$ and its bound
$C_{\Delta}(1+j-n(p))$ are those fixed in Definition~\ref{8.4:def-far-member-hypotheses}, with
$x_{p}$ the base site of $p$ as above.
\item[(c)] The clean-collar hypothesis of Definition~\ref{8.4:def-far-member-hypotheses} holds,
with the clean collar width $W$ in level-$j$ units fixed there.
\item[(d)] The zone-thickness hypothesis of Definition~\ref{8.4:def-far-member-hypotheses} holds.
\item[(e)] The decay weight $E(x,b)$ of \eqref{8.4:eq-decay-weight-a}, and with it the response
size $\Theta(x)$ of \eqref{8.4:eq-decay-weight-b}, is taken with $\rho=\rho_{\mathrm{sc}}$, the
scaled path distance.
\end{enumerate}
Let $c_{4}$ be a lattice constant such that at most $c_{4}L^{4j}$ fine plaquettes have their base
site in a given level-$j$ cell (unit cube of the level-$j$ lattice), and let $c_{4}'$ be a lattice
constant
such that, for every $m\ge0$, at most $c_{4}'(m+1)^{d-1}$ level-$j$ bonds lie at level-$j$ distance
in $[m,m+1)$ from a given site, and at most $c_{4}'(m+1)^{d-1}$ level-$j$ cells lie at level-$j$
distance in $[m,m+1)$ from a given level-$j$ bond. Put $\mu:=\nu:=\delta_{0}c_{d}/16$, and let
$C_{\nu}$ and $C_{2\nu}$ be the constants
\begin{equation*}
C_{\nu}:=e^{\nu}\sum_{m'\ge0}(2+m')^{d-1}e^{-\nu m'},\qquad
C_{2\nu}:=e^{2\nu}\sum_{m'\ge0}(2+m')^{d-1}e^{-2\nu m'} .
\end{equation*}
Assume finally the floor on $M$ and the collar condition
\begin{equation}\label{8.4:eq-frozen-coupling-a}
\begin{gathered}
L^{4}e^{-\mu M}\le\tfrac12;\\
g_{j}Y\,c_{4}'C_{2\nu}\,e^{\delta_{0}c_{1}/8}\,W^{d-1}e^{-2\nu W}\le\theta_{*},
\end{gathered}
\end{equation}
with $\theta_{*}$ the threshold of \eqref{8.4:eq-plaquette-traces-a}. Then, with
$t_{p}(y):=\operatorname{tr}W^{(y)}(\partial p)$ and
$\bar t_{p}(y):=\operatorname{tr}W^{(y)}(\partial p)^{-1}$, defined by this formula for every fine
plaquette $p$ (for plaquettes with a bond in $E_{1}$ these are the plaquette traces of
Corollary~\ref{8.4:cor-plaquette-traces}), the function
\begin{equation*}
\mathcal{L}^{c}(y):=-\sum_{p}\Delta(x_{p})\Bigl(1-\tfrac12\bigl(t_{p}(y)+\bar t_{p}(y)\bigr)\Bigr)
\end{equation*}
is a holomorphic extension of $\mathcal{L}$ to $\mathcal{T}$, and for all $y\in\mathcal{T}$,
$b\in I$ and $\zeta\in\mathfrak{g}^{c}$ with $|\zeta|\le1$,
\begin{equation}\label{8.4:eq-frozen-coupling-b}
\begin{gathered}
|\partial_{b,\zeta}\mathcal{L}^{c}(y)|\le\Lambda_{\mathrm{lag}}
:=K_{L}\,g_{j}(\bar\varepsilon+\theta_{*})\,W^{d-1}e^{-\nu W},\\
K_{L}:=8K_{T}C_{\Delta}c_{4}L^{4}c_{4}'C_{\nu}\,e^{\delta_{0}c_{1}/8+2\nu}.
\end{gathered}
\end{equation}
Let $S_{\Delta}$ be the set of base sites $x_{p}$ of the fine plaquettes $p$ with
$\Delta(x_{p})\neq0$. On $\mathcal{T}(Y_{R}',Y_{I}')$ the mixed box Hessian of $\mathcal{L}^{c}$ is
bounded entrywise by $\Lambda_{\mathrm{lag}}/r$, and indeed its entry at $(b,b')$ is bounded by
$\Lambda_{\mathrm{lag}}/r$ with the factor $e^{-\nu W}$ replaced by
\begin{equation*}
\exp\bigl(-\nu\max\{\operatorname{dist}_{j}(b,S_{\Delta}),\,\operatorname{dist}_{j}(b',S_{\Delta})\}\bigr).
\end{equation*}
The deep plaquettes of zone $0$, that is, the plaquettes with no bond in
$E_{1}$, where $1-\operatorname{Re}\operatorname{tr}V_{0}(\partial p)$ may be of order one and
$|\Delta|$ of order $1+j$, contribute exactly zero to $\partial_{b,\zeta}\mathcal{L}^{c}$ and to its
Hessian.
\end{corollary}

\begin{proof}
\emph{Holomorphy and extension.} For $U\in\mathrm{SU}(N)$ we have $U^{-1}=U^{*}$, hence
$\operatorname{tr}U^{-1}=\overline{\operatorname{tr}U}$ and
\begin{equation*}
1-\operatorname{Re}\operatorname{tr}U=1-\tfrac12\bigl(\operatorname{tr}U+\operatorname{tr}U^{-1}\bigr).
\end{equation*}
Terms with $\Delta(x_{p})=0$ are absent from $\mathcal{L}$ and from $\mathcal{L}^{c}$. For a
plaquette $p$ with $\Delta(x_{p})\neq0$ and a bond in $E_{1}$, Corollary~\ref{8.4:cor-plaquette-traces}
(hypothesis (a)) states that $t_{p}$ and $\bar t_{p}$ are holomorphic on $\mathcal{T}$ and extend
$\operatorname{tr}U^{y}(\partial p)$ and $\operatorname{tr}U^{y}(\partial p)^{-1}$ from real $y$;
its hypothesis \eqref{8.4:eq-plaquette-traces-a} at the base site $x_{p}$ is verified in the next
step. For a deep plaquette of zone $0$ (no bond in $E_{1}$), clause~(iv) of the far-member lemma
states that $t_{p}$ and $\bar t_{p}$ are constant. Hence each term of $\mathcal{L}^{c}$ is
holomorphic on $\mathcal{T}$ and coincides on real $y$ with the corresponding term of
$\mathcal{L}$, so $\mathcal{L}^{c}$ is a holomorphic extension of $\mathcal{L}$. Differentiating
term by term,
\begin{equation*}
\partial_{b,\zeta}\mathcal{L}^{c}
=\tfrac12\sum_{p}\Delta(x_{p})\bigl(\partial_{b,\zeta}t_{p}+\partial_{b,\zeta}\bar t_{p}\bigr),
\end{equation*}
where the sum runs over the plaquettes with $\Delta(x_{p})\neq0$ and a bond in $E_{1}$: the deep
plaquettes of zone $0$ contribute zero, since their $t_{p}$, $\bar t_{p}$ are constant. This
proves the assertion on the deep plaquettes of zone $0$, since a term with constant $t_{p}$,
$\bar t_{p}$ contributes nothing to first or to second derivatives.

\emph{The threshold \eqref{8.4:eq-plaquette-traces-a} at the plaquettes of the sum.} By
hypothesis (c), $\operatorname{dist}_{j}(x_{p},I)\ge W$ for the plaquettes of the sum. Let $x$ be
a fine site with $\operatorname{dist}_{j}(x,I)\ge W$. By \eqref{8.4:eq-decay-weight-a},
$E(x,b)\le\exp\bigl(-(\delta_{0}/8)(c_{d}\rho(x,b)-c_{1})\bigr)$, and
$\rho=\rho_{\mathrm{sc}}\ge\operatorname{dist}_{j}$; since $\delta_{0}c_{d}/8=2\nu$,
\eqref{8.4:eq-decay-weight-b} gives
\begin{align*}
\Theta(x)=g_{j}Y\sum_{b\in I}E(x,b)
&\le g_{j}Y e^{\delta_{0}c_{1}/8}\sum_{b\in I}e^{-2\nu\operatorname{dist}_{j}(x,b)}\\
&\le g_{j}Y e^{\delta_{0}c_{1}/8}\sum_{m\ge W}
\#\{b\in I:\operatorname{dist}_{j}(x,b)\in[m,m+1)\}\,e^{-2\nu m}\\
&\le g_{j}Y e^{\delta_{0}c_{1}/8}\,c_{4}'\sum_{m\ge W}(m+1)^{d-1}e^{-2\nu m},
\end{align*}
where in the second line we used $\operatorname{dist}_{j}(x,b)\ge\operatorname{dist}_{j}(x,I)\ge W$
for $b\in I$, and in the third the definition of $c_{4}'$. Write $m=W+m'$ with $m'\ge0$. Since
$W\ge1$ (the collar width of hypothesis (c) is at least $4L+3$),
\begin{equation*}
W+m'+1\le W(2+m'),
\end{equation*}
because the difference of the two sides is $(W-1)(1+m')\ge0$; hence
$(m+1)^{d-1}\le W^{d-1}(2+m')^{d-1}$ and
\begin{equation*}
\sum_{m\ge W}(m+1)^{d-1}e^{-2\nu m}
\le W^{d-1}e^{-2\nu W}\sum_{m'\ge0}(2+m')^{d-1}e^{-2\nu m'}
\le C_{2\nu}W^{d-1}e^{-2\nu W},
\end{equation*}
using $e^{2\nu}\ge1$ in the definition of $C_{2\nu}$. Therefore
\begin{equation*}
\Theta(x)\le g_{j}Y\,c_{4}'C_{2\nu}\,e^{\delta_{0}c_{1}/8}\,W^{d-1}e^{-2\nu W}\le\theta_{*}
\end{equation*}
by the collar condition in \eqref{8.4:eq-frozen-coupling-a}. Thus
\eqref{8.4:eq-plaquette-traces-a} holds at $x_{p}$ for every plaquette of the sum,
Corollary~\ref{8.4:cor-plaquette-traces} applies there, and \eqref{8.4:eq-plaquette-traces-c}
holds with $\bar\varepsilon+\Theta(x_{p})\le\bar\varepsilon+\theta_{*}$.

\emph{The bound per plaquette.} For a plaquette $p$ of the sum, hypothesis (b) and
\eqref{8.4:eq-plaquette-traces-c}, applied with $n=n(p)$, give
\begin{equation*}
|\Delta(x_{p})|\cdot\tfrac12\bigl(|\partial_{b,\zeta}t_{p}|+|\partial_{b,\zeta}\bar t_{p}|\bigr)
\le C_{\Delta}(1+j-n(p))\cdot K_{T}g_{j}L^{-4n(p)}E(x_{p},b)(\bar\varepsilon+\theta_{*}).
\end{equation*}
We group the plaquettes according to the level-$j$ cell $c$ containing $x_{p}$ and to the zone
$n$. By the definition of $c_{4}$,
\begin{equation*}
\sum_{p:\,x_{p}\in c\cap Z_{n}}L^{-4n}\le c_{4}L^{4(j-n)}.
\end{equation*}

\emph{A lower bound for $\rho_{\mathrm{sc}}$.} Let $x$ be a site of $c\cap Z_{n}$ and $b\in I$.
The bond $b$ lies in $\Omega_{j+1}\subset\Omega_{j}$, hence in $\Omega_{m+1}$ for every
$m\le j-1$, while $x\in Z_{n}$ lies in $\Omega_{m}^{c}$ for every $m>n$. Consequently a fine path
from $x$ to $b$ must, for each $m$ with $n<m<j$, contain a piece running from $\Omega_{m+1}$ to
$\Omega_{m}^{c}$ inside $Z_{m}$. By hypothesis (d) this piece has at least $ML^{m}$ fine bonds in
$Z_{m}$; each carries weight $L^{-m}$ in $\rho_{\mathrm{sc}}$, so the piece contributes at least
$M$ to $\rho_{\mathrm{sc}}$. There are $(j-1-n)_{+}$ such values of $m$, hence
$\rho_{\mathrm{sc}}(x,b)\ge M(j-1-n)_{+}$. Also $\rho_{\mathrm{sc}}(x,b)\ge\operatorname{dist}_{j}(x,b)$,
and so
\begin{equation*}
\rho_{\mathrm{sc}}(x_{p},b)\ge\tfrac12\bigl(\operatorname{dist}_{j}(x_{p},b)+M(j-1-n)_{+}\bigr).
\end{equation*}
Inserting this in \eqref{8.4:eq-decay-weight-a}, using $\delta_{0}c_{d}/16=\nu=\mu$ and, for
$x_{p}\in c$, $\operatorname{dist}_{j}(x_{p},b)\ge\operatorname{dist}_{j}(c,b)-1$, we obtain
\begin{equation}\label{8.4:eq-frozen-coupling-1}
E(x_{p},b)\le e^{\delta_{0}c_{1}/8}\,e^{-\nu\operatorname{dist}_{j}(c,b)+\nu}\,
e^{-\mu M(j-1-n)_{+}} .
\end{equation}

\emph{Summation.} The cells $c$ containing a base site $x_{p}$ of the sum satisfy
$\operatorname{dist}_{j}(c,I)\ge W-1$: since
$\operatorname{dist}_{j}(c,I)\ge\operatorname{dist}_{j}(x_{p},I)-1$ for $x_{p}\in c$, this follows
from hypothesis (c). Summing the bound per plaquette, with \eqref{8.4:eq-frozen-coupling-1} and the count
of plaquettes per cell and zone, and writing $a:=j-n$ (so $0\le a\le j$, $1+j-n=1+a$,
$L^{4(j-n)}=L^{4a}$ and $(j-1-n)_{+}=(a-1)_{+}$), we get
\begin{equation}\label{8.4:eq-frozen-coupling-2}
\begin{split}
|\partial_{b,\zeta}\mathcal{L}^{c}|
\le{}&K_{T}C_{\Delta}c_{4}\,g_{j}(\bar\varepsilon+\theta_{*})\,e^{\delta_{0}c_{1}/8+\nu}\\
&\times\Bigl[\sum_{a=0}^{j}(1+a)L^{4a}e^{-\mu M(a-1)_{+}}\Bigr]
\cdot\Bigl[\sum_{c:\,\operatorname{dist}_{j}(c,I)\ge W-1}e^{-\nu\operatorname{dist}_{j}(c,b)}\Bigr].
\end{split}
\end{equation}

\emph{The first bracket.} The term $a=0$ is $1$ and the term $a=1$ is $2L^{4}$. For $a\ge2$,
\begin{equation*}
(1+a)L^{4a}e^{-\mu M(a-1)}=(1+a)L^{4}\bigl(L^{4}e^{-\mu M}\bigr)^{a-1}\le(1+a)L^{4}2^{-(a-1)}
\end{equation*}
by the floor on $M$ in \eqref{8.4:eq-frozen-coupling-a}. Extending the sum to all $a\ge2$ and
writing $m=a-1$,
\begin{equation*}
\sum_{a=0}^{j}(1+a)L^{4a}e^{-\mu M(a-1)_{+}}
\le1+2L^{4}+L^{4}\sum_{m\ge1}(m+2)2^{-m}=1+2L^{4}+4L^{4}\le8L^{4},
\end{equation*}
since $\sum_{m\ge1}m2^{-m}=2$ and $\sum_{m\ge1}2^{-m}=1$, and $1\le L^{4}$.

\emph{The second bracket.} Since $b\in I$, every cell $c$ with $\operatorname{dist}_{j}(c,I)\ge W-1$
has $\operatorname{dist}_{j}(c,b)\ge W-1$, and by the definition of $c_{4}'$ at most
$c_{4}'(m+1)^{d-1}$ cells lie at level-$j$ distance in $[m,m+1)$ from $b$. Hence the second
bracket is at most $c_{4}'\sum_{m\ge W-1}(m+1)^{d-1}e^{-\nu m}$. Writing $m=W-1+m'$ with $m'\ge0$,
we have $m+1=W+m'\le W(2+m')$ and $e^{-\nu m}=e^{-\nu(W-1)}e^{-\nu m'}$, so, as in the
verification of the threshold,
\begin{equation*}
c_{4}'\sum_{m\ge W-1}(m+1)^{d-1}e^{-\nu m}\le c_{4}'C_{\nu}\,W^{d-1}e^{-\nu(W-1)} .
\end{equation*}

\emph{Collecting.} Inserting the two bracket bounds in \eqref{8.4:eq-frozen-coupling-2} and
writing $e^{\nu}e^{-\nu(W-1)}=e^{2\nu}e^{-\nu W}$,
\begin{equation*}
|\partial_{b,\zeta}\mathcal{L}^{c}|
\le8K_{T}C_{\Delta}c_{4}L^{4}c_{4}'C_{\nu}\,e^{\delta_{0}c_{1}/8+2\nu}\,
g_{j}(\bar\varepsilon+\theta_{*})\,W^{d-1}e^{-\nu W}=\Lambda_{\mathrm{lag}},
\end{equation*}
which is \eqref{8.4:eq-frozen-coupling-b}.

\emph{The Hessian.} The entrywise bound $\Lambda_{\mathrm{lag}}/r$ on
$\mathcal{T}(Y_{R}',Y_{I}')$ is \eqref{8.4:eq-elementary-facts-b} of
Lemma~\ref{8.4:lem-elementary-facts}, applied to the function
$y\mapsto\partial_{b,\zeta}\mathcal{L}^{c}(y)$, which is holomorphic on $\mathcal{T}$ and bounded
there by $\Lambda_{\mathrm{lag}}$. The improved form is obtained by the same computation, with the
bound of Corollary~\ref{8.4:cor-plaquette-traces} on mixed second derivatives, which carries
$\min\{E(x,b),E(x,b')\}$, and with the cell sum kept with both $b$ and $b'$.

\emph{With $\rho=\operatorname{dist}_{j}$ alone.} If the decay weight is taken with
$\rho=\operatorname{dist}_{j}$, the lower bound $\rho_{\mathrm{sc}}\ge M(j-1-n)_{+}$ is not
available, the factor $e^{-\mu M(j-1-n)_{+}}$ is absent from \eqref{8.4:eq-frozen-coupling-1},
and the first bracket becomes $\sum_{a\le j}(1+a)L^{4a}$, which is unbounded in $j$; the argument
gives no bound uniform in $j$, and the corresponding statement is not proved.
\end{proof}

\begin{corollary}[Generic far members: extension and summed gradient bound]\label{8.4:cor-generic-far-members}
Assume the following.
\begin{itemize}
\item[(a)] The hypotheses of the far-member lemma, Lemma~\ref{8.4:lem-far-member}, hold. We write
$\mathcal{T}:=\mathcal{T}(Y_{R},Y_{I})$ for the tube of that lemma, $\mathcal{H}^{(y)}$ for its response
and $W^{(y)}$ for the configuration of \eqref{8.4:eq-far-member-c}.
\item[(b)] Hypotheses \textup{(b)} and \textup{(d)} of Definition~\ref{8.4:def-far-member-hypotheses} hold,
and the decay weight $E(x,b)$ of \eqref{8.4:eq-decay-weight-a} is formed with
$\rho=\rho_{\mathrm{sc}}$ (Notation~\ref{8.4:not-decay-weight}). The constant $K$ and the response size
$\Theta(x)$ are those of \eqref{8.4:eq-decay-weight-b}.
\item[(c)] $(f,\tilde{X},\rho^{(1)},\rho^{(2)},M_{f})$ is a far member in the sense of
Definition~\ref{8.4:def-far-member-hypotheses}, with the norm $\|\cdot\|_{\tilde{X}}$ of
\eqref{8.4:eq-far-member-hypotheses-a} and the conditions \textup{(m1)}--\textup{(m3)} stated there, and
\begin{equation}\label{8.4:eq-generic-far-members-a}
K\max\Bigl\{\frac{1}{\rho^{(1)}},\frac{1}{\rho^{(2)}}\Bigr\}\,\Theta_{\tilde{X}}\le\frac12,
\qquad
\Theta_{\tilde{X}}:=\max_{x\in V(\tilde{X})}\Theta(x),
\end{equation}
where $V(\tilde{X})$ is the set of endpoints of bonds of $\tilde{X}$.
\end{itemize}
Write $W^{(y)}|_{\tilde{X}}$ for the restriction of $W^{(y)}$ to the bonds of $\tilde{X}$. Then
$\varphi_{f}(y):=f(W^{(y)}|_{\tilde{X}})$ is the holomorphic extension to $\mathcal{T}$ of the member's value,
$|\varphi_{f}|\le M_{f}$ on $\mathcal{T}$, and for $y\in\mathcal{T}$, $b\in I$ and $\zeta\in\mathfrak{g}^{c}$
with $|\zeta|\le1$,
\begin{equation}\label{8.4:eq-generic-far-members-b}
|\partial_{b,\zeta}\varphi_{f}(y)|
\le 2M_{f}Kg_{j}\max\Bigl\{\frac{1}{\rho^{(1)}},\frac{1}{\rho^{(2)}}\Bigr\}E_{\tilde{X}}(b),
\qquad
E_{\tilde{X}}(b):=\max_{x\in V(\tilde{X})}E(x,b).
\end{equation}
Let $Y_{R}'<Y_{R}$, $Y_{I}'<Y_{I}$ and $r:=\min(Y_{R}-Y_{R}',Y_{I}-Y_{I}')$. For
$y\in\mathcal{T}(Y_{R}',Y_{I}')$, $b,b'\in I$ and $|\zeta|,|\zeta'|\le1$, the mixed second derivatives satisfy
\begin{equation}\label{8.4:eq-generic-far-members-1}
|\partial_{b',\zeta'}\partial_{b,\zeta}\varphi_{f}(y)|
\le\frac{2M_{f}Kg_{j}}{r}\max\Bigl\{\frac{1}{\rho^{(1)}},\frac{1}{\rho^{(2)}}\Bigr\}
\min\{E_{\tilde{X}}(b),E_{\tilde{X}}(b')\},
\end{equation}
that is, the bound \eqref{8.4:eq-generic-far-members-b} divided by $r$, with $E_{\tilde{X}}(b)$ replaced by
$\min\{E_{\tilde{X}}(b),E_{\tilde{X}}(b')\}$.

Assume in addition that every far member $(f,\tilde{X}_{f},\rho^{(1)}_{f},\rho^{(2)}_{f},M_{f})$ satisfies
\textup{(c)}, that \eqref{8.4:eq-frozen-coupling-a} holds, with $\mu=\nu=\delta_{0}c_{d}/16$ as there, and that
the following counting condition holds: with
$w_{f}:=g_{j}M_{f}\max\{1/\rho^{(1)}_{f},1/\rho^{(2)}_{f}\}$, there are constants $K_{\mathrm{cnt}}\ge0$ and
$q\ge0$ such that, for every choice of sites $x_{f}\in V(\tilde{X}_{f})$, one for each far member, every
level-$j$ cell $c$ and every zone index $n$,
\begin{equation}\label{8.4:eq-generic-far-members-2}
\sum_{f:\,x_{f}\in c\cap Z_{n}}w_{f}\le K_{\mathrm{cnt}}(1+j-n)^{q}L^{4(j-n)}.
\end{equation}
Then, with $c_{4}'$ and $C_{\nu}$ the lattice constants of the cell sum in the proof of
Corollary~\ref{8.4:cor-frozen-coupling}, for all $b\in I$ and $|\zeta|\le1$,
\begin{equation}\label{8.4:eq-generic-far-members-c}
\begin{gathered}
\sum_{f}\sup_{\mathcal{T}}|\partial_{b,\zeta}\varphi_{f}|\le\Lambda_{\mathrm{mem}}:=K_{M}W^{d-1}e^{-\nu W},\\
K_{M}:=32c_{q}L^{4}KK_{\mathrm{cnt}}c_{4}'C_{\nu}e^{\delta_{0}c_{1}/8+2\nu},
\qquad
c_{q}:=\sum_{m\ge0}(m+3)^{q}2^{-m},
\end{gathered}
\end{equation}
and, for all $b,b'\in I$ and $|\zeta|,|\zeta'|\le1$,
\begin{equation}\label{8.4:eq-generic-far-members-3}
\sum_{f}\sup_{\mathcal{T}(Y_{R}',Y_{I}')}|\partial_{b',\zeta'}\partial_{b,\zeta}\varphi_{f}|
\le\frac{\Lambda_{\mathrm{mem}}}{r}.
\end{equation}
\end{corollary}

\begin{proof}
\emph{Size of the response on $\tilde{X}$.} Let $y\in\mathcal{T}$. By the definition
\eqref{8.4:eq-far-member-hypotheses-a} of the norm, $\|\mathcal{H}^{(y)}|_{\tilde{X}}\|_{\tilde{X}}$ is the larger
of the supremum over bonds $b_{x}\in\tilde{X}$ of $\eta|\mathcal{H}^{(y)}(b_{x})|L^{n}/\rho^{(1)}$ and the
supremum over plaquettes $p$ with all bonds in $\tilde{X}$ of
$\eta^{2}|(D_{U_{0}}\mathcal{H}^{(y)})(p)|L^{2n}/\rho^{(2)}$, where $x$ is the initial site of the bond,
respectively the base site of the plaquette, and $n=n(x)$. In both cases $x$ is an endpoint of a bond of
$\tilde{X}$, so $x\in V(\tilde{X})$ and $\Theta(x)\le\Theta_{\tilde{X}}$. Clause~(iii) of
Lemma~\ref{8.4:lem-far-member} bounds the first quantity by $K\Theta(x)/\rho^{(1)}$ and the second by
$K\Theta(x)/\rho^{(2)}$. Hence, by \eqref{8.4:eq-generic-far-members-a},
\begin{equation*}
\|\mathcal{H}^{(y)}|_{\tilde{X}}\|_{\tilde{X}}
\le K\max\Bigl\{\frac{1}{\rho^{(1)}},\frac{1}{\rho^{(2)}}\Bigr\}\Theta_{\tilde{X}}\le\frac12
\qquad(y\in\mathcal{T}).
\end{equation*}

\emph{Holomorphic extension and its size.} By the preceding display, $\mathcal{H}^{(y)}|_{\tilde{X}}$ lies in
the open set $\{Z:\|Z\|_{\tilde{X}}<1\}$ for every $y\in\mathcal{T}$. By condition \textup{(m1)},
$Z\mapsto f(\exp(i\eta Z)U_{0}|_{\tilde{X}})$ is defined, holomorphic and bounded by $M_{f}$ on this set, and by
condition \textup{(m2)} the function $f$ is invariant under $G$-valued gauge transformations. Clause~(vi) of
Lemma~\ref{8.4:lem-far-member}, applied with $F=f$, therefore shows that
\begin{equation*}
\varphi_{f}(y)=f(W^{(y)}|_{\tilde{X}})=f(e^{i\eta\mathcal{H}^{(y)}}U_{0}|_{\tilde{X}})
\end{equation*}
(the second equality is the bondwise definition of $W^{(y)}$ in \eqref{8.4:eq-far-member-c}) is holomorphic on
$\mathcal{T}$ and is the unique holomorphic extension of the real function $y\mapsto f(U^{y}|_{\tilde{X}})$,
which by condition \textup{(m3)} is the member's value. Since $\mathcal{H}^{(y)}|_{\tilde{X}}$ lies in the set on
which \textup{(m1)} gives the bound $M_{f}$, we have $|\varphi_{f}|\le M_{f}$ on $\mathcal{T}$.

\emph{The first derivative.} Fix $y\in\mathcal{T}$, $b\in I$ and $\zeta$ with $|\zeta|\le1$, and put
$Z':=\partial_{b,\zeta}\mathcal{H}^{(y)}|_{\tilde{X}}$. For complex $s$ write $y+s\zeta e_{b}$ for the point
obtained from $y$ by adding $s\zeta$ to its box variable at $b$. By clause~(i) of
Lemma~\ref{8.4:lem-far-member} the map $y\mapsto\mathcal{H}^{(y)}$ is holomorphic, so
$\mathcal{H}^{(y+s\zeta e_{b})}|_{\tilde{X}}=\mathcal{H}^{(y)}|_{\tilde{X}}+sZ'+O(s^{2})$ as $s\to0$. The
$s^{2}$-remainder does not contribute to a first derivative of a holomorphic composition; by the chain rule,
\begin{equation}\label{8.4:eq-generic-far-members-4}
\partial_{b,\zeta}\varphi_{f}(y)
=\frac{d}{ds}\Big|_{s=0}f\bigl(e^{i\eta\mathcal{H}^{(y+s\zeta e_{b})}}U_{0}|_{\tilde{X}}\bigr)
=\frac{d}{ds}\Big|_{s=0}f\bigl(e^{i\eta(\mathcal{H}^{(y)}+sZ')}U_{0}|_{\tilde{X}}\bigr).
\end{equation}
If $Z'=0$, the right-hand side of \eqref{8.4:eq-generic-far-members-4} is the derivative of a constant, so
$\partial_{b,\zeta}\varphi_{f}(y)=0$ and \eqref{8.4:eq-generic-far-members-b} holds. Otherwise define, for
complex $t$ with $|t|<\tfrac12$,
\begin{equation*}
\psi(t):=f\Bigl(\exp\Bigl(i\eta\Bigl(\mathcal{H}^{(y)}+t\,\frac{Z'}{\|Z'\|_{\tilde{X}}}\Bigr)\Bigr)
U_{0}|_{\tilde{X}}\Bigr).
\end{equation*}
By the triangle inequality and the first step, the argument $\mathcal{H}^{(y)}|_{\tilde{X}}+tZ'/\|Z'\|_{\tilde{X}}$
has norm at most $\tfrac12+|t|<1$; it depends affinely on $t$, so by \textup{(m1)} the function $\psi$ is
holomorphic on the disc $|t|<\tfrac12$ and $|\psi|\le M_{f}$ there. Cauchy's estimate on the circle of radius
$t'$, for every $t'<\tfrac12$, gives $|\psi'(0)|\le M_{f}/t'$, hence $|\psi'(0)|\le2M_{f}$. Substituting
$s=t/\|Z'\|_{\tilde{X}}$ in \eqref{8.4:eq-generic-far-members-4} shows
$\partial_{b,\zeta}\varphi_{f}(y)=\|Z'\|_{\tilde{X}}\,\psi'(0)$. Therefore
\begin{equation*}
|\partial_{b,\zeta}\varphi_{f}(y)|\le2M_{f}\,\|\partial_{b,\zeta}\mathcal{H}^{(y)}|_{\tilde{X}}\|_{\tilde{X}}.
\end{equation*}
We bound the norm on the right through its definition \eqref{8.4:eq-far-member-hypotheses-a}. For a bond
$b_{x}\in\tilde{X}$ with initial site $x$, and for a plaquette $p$ with all bonds in $\tilde{X}$ and base site
$x$, with $n=n(x)$ in each case, \eqref{8.4:eq-far-member-d} and \eqref{8.4:eq-far-member-e} give
respectively
\begin{align*}
\frac{\eta|\partial_{b,\zeta}\mathcal{H}^{(y)}(b_{x})|L^{n}}{\rho^{(1)}}
&\le\frac{Kg_{j}E(x,b)}{\rho^{(1)}},\\
\frac{\eta^{2}|(D_{U_{0}}\partial_{b,\zeta}\mathcal{H}^{(y)})(p)|L^{2n}}{\rho^{(2)}}
&\le\frac{Kg_{j}E(x,b)}{\rho^{(2)}}.
\end{align*}
In both cases $x\in V(\tilde{X})$, so the supremum over bonds and plaquettes of $\tilde{X}$ of $E$ at their sites
is at most $E_{\tilde{X}}(b)$. Hence
\begin{equation*}
\|\partial_{b,\zeta}\mathcal{H}^{(y)}|_{\tilde{X}}\|_{\tilde{X}}
\le Kg_{j}\max\Bigl\{\frac{1}{\rho^{(1)}},\frac{1}{\rho^{(2)}}\Bigr\}E_{\tilde{X}}(b),
\end{equation*}
which together with the preceding display is \eqref{8.4:eq-generic-far-members-b}.

\emph{Mixed second derivatives.} Fix $b,b'\in I$ and $\zeta,\zeta'$ with $|\zeta|,|\zeta'|\le1$. The function
$\partial_{b,\zeta}\varphi_{f}$ is holomorphic on $\mathcal{T}$, as a derivative of the holomorphic function
$\varphi_{f}$, and by \eqref{8.4:eq-generic-far-members-b} it is bounded on $\mathcal{T}$ by the right-hand side
of \eqref{8.4:eq-generic-far-members-b}. The elementary fact \eqref{8.4:eq-elementary-facts-b} of
Lemma~\ref{8.4:lem-elementary-facts}, applied to this function, bounds its derivative in the direction
$(b',\zeta')$ at points of $\mathcal{T}(Y_{R}',Y_{I}')$ by that bound divided by $r$; this is
\eqref{8.4:eq-generic-far-members-1} with $E_{\tilde{X}}(b)$ in place of the minimum. For the symmetrisation:
mixed derivatives of the holomorphic function $\varphi_{f}$ commute, so
$\partial_{b',\zeta'}\partial_{b,\zeta}\varphi_{f}=\partial_{b,\zeta}\partial_{b',\zeta'}\varphi_{f}$, and the same
argument with the roles of $(b,\zeta)$ and $(b',\zeta')$ exchanged gives the bound with $E_{\tilde{X}}(b')$. The
smaller of the two bounds is \eqref{8.4:eq-generic-far-members-1}.

\emph{Summation.} Fix $b\in I$ and $\zeta$ with $|\zeta|\le1$. For each far member choose a site
$x_{f}\in V(\tilde{X}_{f})$ with $E(x_{f},b)=E_{\tilde{X}_{f}}(b)$; this is possible because $V(\tilde{X}_{f})$
is finite. The bound below is uniform in this choice, since it uses only $x_{f}\in V(\tilde{X}_{f})$ and
$\operatorname{dist}_{j}(x_{f},I)\ge W$, the latter from $\operatorname{dist}_{j}(\tilde{X}_{f},I)\ge W$ in the
definition of a far member; and the counting condition \eqref{8.4:eq-generic-far-members-2} is assumed for every
such choice. By \eqref{8.4:eq-generic-far-members-b}, which holds uniformly in $y\in\mathcal{T}$, and the
definition of $w_{f}$,
\begin{equation*}
\sum_{f}\sup_{\mathcal{T}}|\partial_{b,\zeta}\varphi_{f}|\le2K\sum_{f}w_{f}E(x_{f},b).
\end{equation*}
Grouping the far members according to the level-$j$ cell $c$ and the zone $Z_{n}$ containing $x_{f}$, and
bounding $E(x_{f},b)$ by its maximum over $c\cap Z_{n}$, the right-hand side is at most
\begin{equation*}
2K\sum_{c,n}\Bigl[\sum_{f:\,x_{f}\in c\cap Z_{n}}w_{f}\Bigr]\max_{x\in c\cap Z_{n}}E(x,b).
\end{equation*}
The cells that contain a site $x_{f}$ satisfy $\operatorname{dist}_{j}(c,I)\ge W-1$. For a site $x$ of a level-$j$
cell $c$ lying in $Z_{n}$, the bound for the decay weight established in the proof of
Corollary~\ref{8.4:cor-frozen-coupling} gives, with $(t)_{+}:=\max\{t,0\}$,
\begin{equation}\label{8.4:eq-generic-far-members-5}
E(x,b)\le e^{\delta_{0}c_{1}/8+\nu}\,e^{-\nu\operatorname{dist}_{j}(c,b)}\,e^{-\mu M(j-1-n)_{+}}.
\end{equation}
Inserting \eqref{8.4:eq-generic-far-members-2} and \eqref{8.4:eq-generic-far-members-5}, we obtain
\begin{equation*}
\sum_{f}\sup_{\mathcal{T}}|\partial_{b,\zeta}\varphi_{f}|
\le2KK_{\mathrm{cnt}}\,e^{\delta_{0}c_{1}/8+\nu}
\sum_{c:\,\operatorname{dist}_{j}(c,I)\ge W-1}e^{-\nu\operatorname{dist}_{j}(c,b)}
\sum_{n}(1+j-n)^{q}L^{4(j-n)}e^{-\mu M(j-1-n)_{+}}.
\end{equation*}
Writing $a:=j-n$, the $n$-sum is at most
$\sum_{a\ge0}(1+a)^{q}L^{4a}e^{-\mu M(a-1)_{+}}$. The term $a=0$ equals $1$ and the term $a=1$ equals
$2^{q}L^{4}$; for $a\ge2$ put $m:=a-1\ge1$, so that the term equals $L^{4}(m+2)^{q}(L^{4}e^{-\mu M})^{m}$, which
is at most $L^{4}(m+2)^{q}2^{-m}$ by the first inequality of \eqref{8.4:eq-frozen-coupling-a},
$L^{4}e^{-\mu M}\le\tfrac12$ (this is the only part of \eqref{8.4:eq-frozen-coupling-a} used). Hence
\begin{equation*}
\sum_{a\ge0}(1+a)^{q}L^{4a}e^{-\mu M(a-1)_{+}}
\le1+2^{q}L^{4}+L^{4}\sum_{m\ge1}(m+2)^{q}2^{-m}\le16c_{q}L^{4}.
\end{equation*}
For the last inequality: $c_{q}\ge3^{q}\ge\max(1,2^{q})$, so $1\le c_{q}L^{4}$ and $2^{q}L^{4}\le c_{q}L^{4}$;
and, shifting the summation index by one,
\begin{equation*}
\sum_{m\ge1}(m+2)^{q}2^{-m}=\frac12\sum_{m\ge0}(m+3)^{q}2^{-m}=\frac12c_{q};
\end{equation*}
the sum of the three terms is at
most $\tfrac52c_{q}L^{4}\le16c_{q}L^{4}$. The $c$-sum is bounded as in the proof of
Corollary~\ref{8.4:cor-frozen-coupling}:
\begin{equation*}
\sum_{c:\,\operatorname{dist}_{j}(c,I)\ge W-1}e^{-\nu\operatorname{dist}_{j}(c,b)}
\le c_{4}'C_{\nu}W^{d-1}e^{-\nu(W-1)}.
\end{equation*}
Multiplying the three factors,
\begin{equation*}
\sum_{f}\sup_{\mathcal{T}}|\partial_{b,\zeta}\varphi_{f}|
\le2KK_{\mathrm{cnt}}\cdot16c_{q}L^{4}\cdot e^{\delta_{0}c_{1}/8+\nu}\cdot c_{4}'C_{\nu}W^{d-1}e^{-\nu W+\nu}
=K_{M}W^{d-1}e^{-\nu W},
\end{equation*}
which is \eqref{8.4:eq-generic-far-members-c}.

For \eqref{8.4:eq-generic-far-members-3}: since $\min\{E_{\tilde{X}_{f}}(b),E_{\tilde{X}_{f}}(b')\}
\le E_{\tilde{X}_{f}}(b)$, the bound \eqref{8.4:eq-generic-far-members-1} for each far member is at most the
right-hand side of \eqref{8.4:eq-generic-far-members-b} divided by $r$, uniformly in
$y\in\mathcal{T}(Y_{R}',Y_{I}')$. The summation just carried out, applied to these bounds, therefore gives
$\Lambda_{\mathrm{mem}}/r$.
\end{proof}

\begin{proposition}[The small-piece statement. Stated; proof incomplete at the step marked $(\ast)$]
\label{8.4:prop-small-piece}
Assume the hypotheses of Corollary~\ref{8.4:cor-frozen-coupling} (the frozen-coupling member at
live regions), among them the floor on $M$ and the collar-width condition of
\eqref{8.4:eq-frozen-coupling-a}, and, for every
generic far member, the hypotheses of Corollary~\ref{8.4:cor-generic-far-members} (generic far
members: extension and summed gradient bound). Put $\nu=\delta_{0}c_{d}/16$. Let $W$ be the width
of the clean collar in level-$j$ units, so that $W=100M\check R_{j}$ up to an additive correction
coming from the construction of the collar, and let $\Lambda_{\mathrm{lag}}$, $K_{L}$ and
$\Lambda_{\mathrm{mem}}$, $K_{M}$ be
the quantities of \eqref{8.4:eq-frozen-coupling-b} and \eqref{8.4:eq-generic-far-members-c}.
Let $\mathcal{T}(Y_{R}',Y_{I}')\subset\mathcal{T}$ be the smaller tube and $r>0$ the margin with
which Corollaries~\ref{8.4:cor-frozen-coupling} and~\ref{8.4:cor-generic-far-members} bound the
mixed box Hessian, through \eqref{8.4:eq-elementary-facts-b} of
Lemma~\ref{8.4:lem-elementary-facts}.
\begin{enumerate}
\item[(a)] The far members, namely the frozen-coupling member together with the generic far
members, contribute in total to the supremum over the tube of the box gradient of the step
exponent, and to the supremum of its mixed box Hessian, at most
\begin{equation}\label{8.4:eq-small-piece-a}
\Lambda_{\mathrm{lag}}+\Lambda_{\mathrm{mem}}\ \text{ on }\mathcal{T},
\qquad
\frac{\Lambda_{\mathrm{lag}}+\Lambda_{\mathrm{mem}}}{r}\ \text{ on }\mathcal{T}(Y_{R}',Y_{I}').
\end{equation}
Each of these carries the factor $W^{d-1}e^{-\nu W}$, multiplied by factors that are polynomial
in $L$ and in the age-independent constants. No factor $g_{j}^{-2}$ survives.
\item[(b)] Let $\mathfrak{b}$ be the budget of the smallness hypothesis on the step exponent in
the localisation of the sourced step, and $\mathfrak{b}'$ its companion for the mixed Hessian.
The comparison of
\eqref{8.4:eq-small-piece-a} with this budget is the one-sided condition
\begin{equation}\label{8.4:eq-small-piece-b}
\Lambda_{\mathrm{lag}}+\Lambda_{\mathrm{mem}}\le\mathfrak{b},
\qquad
\frac{\Lambda_{\mathrm{lag}}+\Lambda_{\mathrm{mem}}}{r}\le\mathfrak{b}'.
\end{equation}
This is a lower bound on $W=100M\check R_{j}$ of the form
\begin{equation}\label{8.4:eq-small-piece-1}
\nu W-(d-1)\log W\ \ge\
\log\frac{K_{L}g_{j}(\bar\varepsilon+\theta_{*})+K_{M}}{\mathfrak{b}}.
\end{equation}
It holds for all $g_{j}\le\gamma$ and all $M\ge M_{0}(L)$ if and only if $\check R_{j}$ grows at
least like $\log(1/\mathfrak{b}(g_{j}))/(100\nu M)$. When $\mathfrak{b}$ is a power of $g_{j}$,
or a power of $p_{0}(g_{j})$, this requirement is the growth condition on Ba{\l}aban's radii.
Here the growth condition states that $\check R_{j}\ge R_{j}$, where $\check R_{j}$ is the radius
that fixes the clean collar width $W$ above and $R_{j}$ is Ba{\l}aban's radius $R_{k}$ of
\cite{B-CMP119} at level $k=j$, a function of $g_{j}$ such that, for every fixed $m$,
$e^{-R_{j}}p_{0}(g_{j})^{m}\le1$ for all $g_{j}\le\gamma(m)$, and $R_{j}\to\infty$ as
$g_{j}\to0$.
When $\mathfrak{b}=e^{-c\,p_{0}(g_{k_{0}})}$, with the degree function $p_{0}(\cdot)$, where
$k_{0}$ is a level fixed by the localisation of the sourced step and $c>0$ is a constant, the
budget is a function of the coupling $g_{k_{0}}$ at level $k_{0}$, not
of $g_{j}$. The far members that carry this budget are called, in the terminology of the
localisation of the sourced step, the \emph{far spent-pocket pieces}. For them the growth
condition on Ba{\l}aban's radii,
being a condition in $g_{j}$, does not by itself yield \eqref{8.4:eq-small-piece-b}, and
condition \eqref{8.4:eq-small-piece-b}, as stated, is then a cap in the sense of
Definition~\ref{2.3:def-cap}.
\item[(c)] Apart from the far spent-pocket pieces of item (b), condition
\eqref{8.4:eq-small-piece-b} imposes no lower bound on a coupling. It
imposes no upper bound on any of $L$, $M$, the number of renormalisation steps, $N$, $D$
(the parameter of the localisation of the sourced step that enters its set of clean histories
$\mathcal{C}_{S}(D)$), $r$, $j$. If $\check R_{j}$ were
bounded, \eqref{8.4:eq-small-piece-b} would force a lower bound on $g_{j}$, that is, a cap in
the sense of Definition~\ref{2.3:def-cap}. Under the growth condition on Ba{\l}aban's radii
$\check R_{j}$ is unbounded, so this does not occur; the only cap that remains is that of
item (b) for the far spent-pocket pieces.
\end{enumerate}
\end{proposition}

\begin{proof}
\emph{The bound \eqref{8.4:eq-small-piece-a}.} The frozen-coupling member is treated by
Corollary~\ref{8.4:cor-frozen-coupling}. By \eqref{8.4:eq-frozen-coupling-b}, the gradient
$\partial_{b,\zeta}$ of its holomorphic extension $\mathcal{L}^{c}$ is bounded on $\mathcal{T}$
by $\Lambda_{\mathrm{lag}}$, for every $b\in I$, the index set of the box variables of the tube
$\mathcal{T}$, and every direction $|\zeta|\le1$. The generic far members are
treated by Corollary~\ref{8.4:cor-generic-far-members}. By
\eqref{8.4:eq-generic-far-members-c}, the sum over all of them of the suprema on $\mathcal{T}$
of $|\partial_{b,\zeta}\varphi_{f}|$ is at most $\Lambda_{\mathrm{mem}}$. Adding the two bounds
gives the first half of \eqref{8.4:eq-small-piece-a}.

The mixed box Hessian on $\mathcal{T}(Y_{R}',Y_{I}')$ is bounded by this gradient bound divided
by $r$, by \eqref{8.4:eq-elementary-facts-b} of Lemma~\ref{8.4:lem-elementary-facts} (three
elementary facts on holomorphic functions and matrix exponentials). This is the second half of
\eqref{8.4:eq-small-piece-a}.

By the definitions in \eqref{8.4:eq-frozen-coupling-b} and
\eqref{8.4:eq-generic-far-members-c},
\begin{equation}\label{8.4:eq-small-piece-2}
\Lambda_{\mathrm{lag}}+\Lambda_{\mathrm{mem}}
=\bigl(K_{L}g_{j}(\bar\varepsilon+\theta_{*})+K_{M}\bigr)\,W^{d-1}e^{-\nu W}.
\end{equation}
The constants $K_{L}$ and $K_{M}$, as displayed there, are products of $L^{4}$ and of
age-independent constants. Hence \eqref{8.4:eq-small-piece-a} carries $W^{d-1}e^{-\nu W}$
against factors of the stated kind.

\emph{The condition \eqref{8.4:eq-small-piece-b}.} Consider the function
$W\mapsto W^{d-1}e^{-\nu W}$. Its derivative is
\[
W^{d-2}e^{-\nu W}\bigl((d-1)-\nu W\bigr),
\]
so the function is decreasing for $W\ge(d-1)/\nu$, and it tends to $0$ as $W\to\infty$. By
\eqref{8.4:eq-small-piece-2}, the condition \eqref{8.4:eq-small-piece-b} is therefore
equivalent to $W\ge\max\{W_{*}(\mathfrak{b}),W_{*}(r\mathfrak{b}')\}$, for an explicit threshold
$W_{*}(\cdot)$ that is finite at every positive argument.

Both sides of the first inequality of \eqref{8.4:eq-small-piece-b} are positive, so we may take
logarithms in \eqref{8.4:eq-small-piece-2}. The first inequality is then equivalent to
\eqref{8.4:eq-small-piece-1}. The second inequality is equivalent to the same display with
$\mathfrak{b}$ replaced by $r\mathfrak{b}'$.

Now put $W=100M\check R_{j}$, up to the additive correction coming from the construction of the
collar, and
$\mathfrak{b}=\mathfrak{b}(g_{j})$. Then \eqref{8.4:eq-small-piece-b} holds for all
$g_{j}\le\gamma$ provided that, for all $g_{j}\le\gamma$,
\begin{equation}\label{8.4:eq-small-piece-3}
100\nu M\check R_{j}-(d-1)\log\bigl(100M\check R_{j}\bigr)
\ \ge\
\log\frac{K_{L}g_{j}(\bar\varepsilon+\theta_{*})+K_{M}}{\mathfrak{b}(g_{j})},
\end{equation}
together with the same inequality with $\mathfrak{b}(g_{j})$ replaced by $r\,\mathfrak{b}'(g_{j})$.

\emph{The case of a budget bounded below in terms of $g_{j}$ and $p_{0}(g_{j})$.} Suppose that
there are constants
$c_{\mathfrak{b}}>0$ and
$m_{1},m_{2},m_{3}\ge0$ such that
\[
\mathfrak{b}(g)\ \ge\ c_{\mathfrak{b}}\,g^{m_{1}}p_{0}(g)^{-m_{2}}e^{-m_{3}p_{0}(g)}.
\]
Suppose also that the degree function $p_{0}(\cdot)$ has at most logarithmic growth in
$x=g^{-2}$. Taking the logarithm of the lower bound gives
\[
\log\frac{1}{\mathfrak{b}(g_{j})}
\ \le\ \log\frac{1}{c_{\mathfrak{b}}}+\frac{m_{1}}{2}\log x_{j}+m_{2}\log p_{0}(g_{j})
+m_{3}\,p_{0}(g_{j}).
\]
Hence the right side of \eqref{8.4:eq-small-piece-3} is at most
$C\bigl(1+\log x_{j}+p_{0}(g_{j})\bigr)$, with a constant $C$ that depends only on $K_{L}$,
$K_{M}$, $\gamma$, $\theta_{*}$, an upper bound for $\bar\varepsilon$, and $c_{\mathfrak{b}}$,
$m_{1}$, $m_{2}$, $m_{3}$; in particular $C$ does not depend on $M$, on $j$, or on
$g_{j}\le\gamma$.

Under the growth condition of item (b), $\check R_{j}\ge R_{j}$, with $R_{j}$ as described
there.
This condition supplies the inequality \eqref{8.4:eq-small-piece-3} after a ceiling on $g_{j}$
and the floor $M\ge1$. $(\ast)$

\emph{The far spent-pocket pieces.} For the far spent-pocket pieces of item (b) the budget is
$\mathfrak{b}=e^{-c\,p_{0}(g_{k_{0}})}$, a function of the coupling
$g_{k_{0}}$ at level $k_{0}$ and not of $g_{j}$. The growth condition of item (b) is a
condition in $g_{j}$, so for these pieces it does not by itself yield
\eqref{8.4:eq-small-piece-b}, and the condition \eqref{8.4:eq-small-piece-b}, as stated, is then
a cap in the sense of Definition~\ref{2.3:def-cap}. This cap is not removed by the present
proposition. A removal of it goes through a variant of the present statement adapted to these
pieces together with a bound for healed regions of this kind, neither of which is part of the
present proposition.

\emph{The case of a faster-decaying budget.} If the budget decays faster than every
$e^{-aR_{j}}$, $a>0$, then \eqref{8.4:eq-small-piece-b} is an additional lower bound on
$\check R_{j}$. The proposition displays this bound as a condition and does not assert it.

\emph{The direction of the condition.} In no case is an upper bound on $L$, $M$, the number of
renormalisation steps, $N$, $D$, $r$, $j$ used, nor a lower bound on a coupling. For members
whose budget is a function of $g_{j}$, condition \eqref{8.4:eq-small-piece-b} could
become a cap, that is, a lower bound on $g_{j}$, only if $\check R_{j}$ were bounded, and this is
excluded by the growth condition on Ba{\l}aban's radii. The cap for the far spent-pocket pieces
treated above has a different origin, namely the dependence of their budget on $g_{k_{0}}$.

\emph{The absence of $g_{j}^{-2}$.} Consider the per-plaquette size
$|\Delta|\bigl(1-\operatorname{Re}\operatorname{tr}U(\partial p)\bigr)$. It can be of order
$(1+j)\cdot1$ only at the deep plaquettes of zone $0$, where the fields are bare large fields.
There the derivatives in the box variables $y$ vanish identically, and these plaquettes
contribute exactly zero; this is the
last sentence of Corollary~\ref{8.4:cor-frozen-coupling}.

At every plaquette that responds to $y$, the base background is regular in the sense fixed for
the base data in this section, and at such a plaquette the bound on the plaquette traces at a
regular background, also fixed in this section, shows that the first
derivative of the Wilson density carries the factor
\[
|U_{0}(\partial p)-1|\ \le\ C_{\mathrm{reg}}\,\bar\varepsilon\,L^{-2n},
\]
where $n$ is the zone index of the plaquette. The coupling difference satisfies
$|\Delta|\le C_{\Delta}(1+\mathrm{age})$, where age denotes the age at level $j$ of the
large-field region to which the coupling of the plaquette belongs, and this factor is absorbed by
the factor
$e^{-\mu M\cdot\mathrm{age}}$ of the proof of Corollary~\ref{8.4:cor-frozen-coupling}.
\end{proof}

\begin{remark}[Own-scale smallness: where needed and where supplied]\label{8.4:rem-own-scale-smallness}
Throughout this remark $V^{\mathrm{base}}$ denotes the base data and $V'$ the variation, in the
notation of \cite{B-CMP102b}, where the base data are written $V_{0}$.

\emph{Where smallness is needed.} \cite[Proposition~9]{B-CMP102b} assumes two things: the
regularity hypothesis \cite[(7)]{B-CMP102b} of
\cite[Theorem~1]{B-CMP102b} for the base data $V^{\mathrm{base}}$ on the whole determining set
$\mathfrak{B}_{k}$, that is, on every zone, and hence also on the zone that contains the
observation point; and smallness of the variation $V'$ on $\mathfrak{B}_{k}$. Under these
assumptions its conclusion \eqref{8.4:eq-minimiser-analyticity-a} holds for all
$y\in\Lambda_{n}$ and $y'\in\Lambda_{n'}$, with no restriction on the positions of $y$ and $y'$.
Consequently, at an end of the response that lies in a live large-field region, what is needed is
smallness of the own-scale plaquettes of the data there, namely of the live variables $V_{n}$ on
$\Gamma_{n}$ inside the region. Smallness of the field relative to the box is not needed, and
neither is a condition of fluctuation type.

\emph{Where it is supplied, and where it is not needed.} At the zones $n\ge 1$, including the
interiors of live large-field regions, this smallness is supplied by the level-$n$ small-field
plaquette functions that the region retains. This is the inductive regularity of the live data,
\cite[p.~256]{B-CMP119}. A large-field region created at step $n+1$ inside $\Omega_{n}$ is cut out
either by conditions of fluctuation type on the level-$n$ data, which read (in \cite{B-CMP119}
they are written with $k$ in place of $n$)
\begin{equation*}
  \bigl|V_{n}(b)\bigl(V^{(n)}_{\square}(b)\bigr)^{-1}-1\bigr|\ge 2\delta_{n}
  \quad\text{\cite[(3.3)]{B-CMP119}},
  \qquad
  |A_{n}(b)|\ge\delta_{n}\quad\text{\cite[(3.16)]{B-CMP119}},
\end{equation*}
where $V^{(n)}_{\square}(b)$ and $A_{n}(b)$ are the fields entering the large-field conditions
\cite[(3.3)]{B-CMP119} and \cite[(3.16)]{B-CMP119} respectively ($A_{n}$ is not the Wilson action
$A$), or by a plaquette condition on the new level-$(n+1)$ field. None of these conditions
revokes the factor $\chi_{n}(\Omega_{n})$. At the top zone the smallness is supplied by the
regularity of the top-level data, \cite[(3.13)--(3.14)]{B-CMP119}, and on the zone-$0$ collar by
the regularity on that collar, \cite[\S1]{B-CMP119}. Lemma~\ref{8.4:lem-base-regularity} assembles
these three regularity statements into the restricted regularity
hypothesis $(7^{\flat})(\bar\varepsilon)$ of \eqref{8.4:eq-step-zones-a} for $V^{\mathrm{base}}$.

The smallness is not supplied inside large-field regions of level~$0$, which are defined by bare
plaquettes and lie in deep zone~$0$; there the regularity hypothesis \cite[(7)]{B-CMP102b} is not
satisfied. It is not needed there,
for three reasons. First, the zero-fold average $M^{0}$ is the identity, so the base configuration
$U_{0}$ equals the zone-$0$ data on the bonds of $E_{0}$. Second, the response vanishes identically
there:
for every fine plaquette $p$ all of whose bonds lie in $E_{0}$, and every $y$ in the tube
$\mathcal{T}(Y_{R},Y_{I})$ of Lemma~\ref{8.4:lem-far-member}, the
plaquette trace of the varied configuration equals that of the zone-$0$ data, by clause~(iv) of
Lemma~\ref{8.4:lem-far-member}. Third, the variational problem on $E_{1}$ does not see these data,
by Lemma~\ref{8.4:lem-zone-zero-invisible}. For base data containing level-$0$ large-field
regions, \cite[Theorem~1]{B-CMP102b} and \cite[Proposition~9]{B-CMP102b} enter through their
transfer to such data, which is one of the hypotheses from which Lemma~\ref{8.4:lem-far-member} is
proved: the conclusions of these two results are taken to hold for
$(\mathbf{B}_{j},V^{\mathrm{base}})$ whenever
$(7^{\flat})(\bar\varepsilon)$ holds with $\bar\varepsilon\le a_{1}$, the base configuration
$U(\mathbf{B}_{j},V^{\mathrm{base}})$ being the minimiser on $E_{1}$ given by
Lemma~\ref{8.4:lem-zone-zero-invisible}, completed by the zone-$0$ data on $E_{0}$. This is the form
in which \cite[\S\S1,~3]{B-CMP119} uses these two results: there \cite[Theorem~1]{B-CMP102b} is
applied to obtain exactly one critical point of the functional \cite[(1.12)]{B-CMP119} in the
domain of integration of \cite[(1.11)]{B-CMP119}, whose zone-$0$ data contain the large-field
regions. The transfer is used here in that form.

\emph{The domains of the localised terms.} We call \emph{localised terms} the terms of
Ba{\l}aban's expansion whose analyticity domains are prescribed in \cite[(2.34)--(2.37)]{B-CMP119}.
Their domains are not kept at a distance from the region where smallness holds, and they need not be.
For such a term, write $\mathcal{D}$ for its localisation domain (denoted $X$ in \cite{B-CMP119}).
By \cite[(2.34)--(2.37)]{B-CMP119}, its analyticity radii are prescribed zone by zone. With
$\alpha_{0,n}$, $\beta$ and $\xi$ as in \cite[(2.34)--(2.37)]{B-CMP119} (this $\beta$ is not the
coupling $\beta_{k+1}$), the radius
on $\mathcal{D}\cap(\Omega_{n}\setminus\Omega_{n+1})$ for $n=1,\dots,j-1$, and on
$\mathcal{D}\cap\Omega_{j}$ for $n=j$, is
\begin{equation}\label{8.4:eq-own-scale-smallness-1}
  \bigl(1-\beta\bigl(1-2^{-(j-n)}\bigr)\bigr)\,\alpha_{0,n}\,\xi^{2}\,(L^{n}\xi)^{-2}.
\end{equation}
No condition is imposed on $\mathcal{D}\cap\Omega_{1}^{c}$. A localised term may therefore read the
background down to and including zone~$0$.

At a zone $n\ge 1$ the factor $\xi^{2}(L^{n}\xi)^{-2}$ in \eqref{8.4:eq-own-scale-smallness-1}
equals $L^{-2n}$. Since $1-2^{-(j-n)}\le 1$, the radius \eqref{8.4:eq-own-scale-smallness-1} is at
least $(1-\beta)\alpha_{0,n}L^{-2n}$. The base background takes up at most
$C_{\mathrm{reg}}\varepsilon_{n}L^{-2n}$ of it, where $C_{\mathrm{reg}}$ is the regularity constant
of the base background. The Cauchy margin of the term at zone~$n$ is therefore at least
\begin{equation}\label{8.4:eq-own-scale-smallness-2}
  \bigl((1-\beta)\alpha_{0,n}-C_{\mathrm{reg}}\varepsilon_{n}\bigr)L^{-2n}.
\end{equation}
This is of the same scale $L^{-2n}$ as the response: by the second bound of clause~(ii) of
Lemma~\ref{8.4:lem-far-member},
\begin{equation*}
  \eta^{2}\bigl|(D_{U_{0}}\partial_{b,\zeta}\mathcal{H}^{(y)})(p)\bigr|
  \le K g_{j}L^{-2n}E(x,b)
\end{equation*}
for every $y$ in the tube, every bond $b$ carrying a box variable, every direction
$\zeta\in\mathfrak{g}^{c}$ with $|\zeta|\le 1$ and every fine plaquette $p$ whose base site $x$
lies in zone $n$. Division of this bound by the margin
\eqref{8.4:eq-own-scale-smallness-2} in the Cauchy estimate shows that the gradient of the
localised term in the box variables carries the factor
\begin{equation*}
  \frac{g_{j}E}{(1-\beta)\alpha_{0,n}-C_{\mathrm{reg}}\varepsilon_{n}},
\end{equation*}
with $E=E(x,b)$. Since $\alpha_{0,n}$ is proportional to $g_{n}$, this factor is polynomial in the
age. No distance from the live region is required.
\end{remark}

\begin{remark}[The one-sided conditions and their order of choice]
\label{8.4:rem-one-sided-conditions}
The constants entering the conditions of this section are chosen in the
following order. First $d$ and $N$. Then $L\ge L_{0}$, with $L_{0}=L_{0}(N)$ the constant of
hypothesis (H2) of Theorem~\ref{3.1:thm-main}; no property of $L_{0}$ is used in this section.
Then Ba{\l}aban's constants
$a_{1}$, $a_{9}$, $b_{D}$, $c_{C}$, $c_{h}$, $c_{D}$, $c_{*}$, $\delta_{0}$, $c_{d}$, $c_{d}'$,
$C_{\mathrm{reg}}$, $C_{\Delta}$. Then $M$. Then the ceilings on the running couplings. Last,
the collar width $W$ and the budget $\mathfrak{b}$. Each condition is one-sided in the sense of
Definition~\ref{2.3:def-order-table} and Lemma~\ref{2.4:lem-no-cycle}, as follows; for condition (f)
this holds under the growth condition on the radius $\check R_{j}$.
\begin{itemize}
\item[(a)] \emph{Ceiling on the running couplings.} With $c_{7}$ as in
\eqref{8.4:eq-base-regularity-a},
\begin{equation}\label{8.4:eq-one-sided-conditions-1}
\bar\varepsilon=c_{7}\max_{0\le n\le j}\varepsilon_{n}\le\min(a_{1},1),
\end{equation}
that is, $g_{n}A_{0}p_{0}(g_{n})\le\min(a_{1},1)/c_{7}$ for all $0\le n\le j$. This condition
secures hypothesis \eqref{8.4:eq-far-member-a} of the far-member lemma
(Lemma~\ref{8.4:lem-far-member}).
\item[(b)] \emph{Ceiling on $g_{j}Y$.} With $Y=Y_{R}+Y_{I}$,
\begin{equation}\label{8.4:eq-one-sided-conditions-2}
g_{j}c_{C}(Y_{R}+Y_{I})\le y_{*}=\min\{b_{D}/2,\,(2c_{h}c_{D})^{-1},\,a_{9}/2\},
\end{equation}
which is the tube condition \eqref{8.4:eq-far-member-b} of Lemma~\ref{8.4:lem-far-member}.
Within the window $Y\le c_{Y}p_{0}(g_{j})^{m_{Y}}$ on the tube radii supplied by the second part
of the growth condition on the radius $\check R_{j}$, it is a ceiling on $g_{j}$.
\item[(c)] \emph{Floor on $M$, chosen after $L$.} With $\mu=\delta_{0}c_{d}/16$,
\begin{equation}\label{8.4:eq-one-sided-conditions-3}
L^{4}e^{-\mu M}\le\tfrac12,
\quad\text{that is,}\quad
M\ge M_{0}(L):=(4\log L+\log 2)\cdot\frac{16}{\delta_{0}c_{d}}.
\end{equation}
This is the first half of \eqref{8.4:eq-frozen-coupling-a}. It is a hypothesis of
Corollaries~\ref{8.4:cor-frozen-coupling} and~\ref{8.4:cor-generic-far-members}. If some zone of
Ba{\l}aban's determining set has thickness only $M_{1}$, the same floor is imposed on $M_{1}$.
\item[(d)] \emph{Floor on the collar width.} With $\nu=\delta_{0}c_{d}/16$,
\begin{equation}\label{8.4:eq-one-sided-conditions-4}
g_{j}Y\,c_{4}'C_{2\nu}e^{\delta_{0}c_{1}/8}\,W^{d-1}e^{-2\nu W}\le\theta_{*},
\end{equation}
the second half of \eqref{8.4:eq-frozen-coupling-a}. Given (b), it is implied by a lower bound
on $W$ by a constant depending on $d$, $L$, $N$ and Ba{\l}aban's constants.
\item[(e)] \emph{Per-member domain condition.} Condition \eqref{8.4:eq-generic-far-members-a} of
Corollary~\ref{8.4:cor-generic-far-members}, for the member $f$ with reading domain
$\tilde X_{f}$ and radii $\rho^{(1)}_{f},\rho^{(2)}_{f}$, reads
\begin{equation}\label{8.4:eq-one-sided-conditions-5}
K\max\{1/\rho^{(1)}_{f},\,1/\rho^{(2)}_{f}\}\,\Theta_{\tilde X_{f}}\le\tfrac12 .
\end{equation}
For members beyond the collar it is implied by a lower bound on $W$ of the same kind as in (d),
in which $\max\{1/\rho^{(1)}_{f},1/\rho^{(2)}_{f}\}$ enters polynomially.
\item[(f)] \emph{The small-piece comparison.} The condition \eqref{8.4:eq-small-piece-b} of the
small-piece statement (Proposition~\ref{8.4:prop-small-piece}, whose proof is incomplete at one
step), that the bounds in \eqref{8.4:eq-small-piece-a} do not exceed the budgets
$\mathfrak{b}$ and $\mathfrak{b}'$, is a lower
bound on $W=100M\check R_{j}$ relative to $\log(1/\mathfrak{b})$. Given the growth
condition on the radius $\check R_{j}$, it is a floor on $W$. It is kept as the exact condition
and is used only as a
hypothesis; it is not asserted unconditionally.
Were $\check R_{j}$ bounded, it would force a lower bound on $g_{j}$, that is, a cap in the
sense of Definition~\ref{2.3:def-cap}; wherever Proposition~\ref{8.4:prop-small-piece} records
\eqref{8.4:eq-small-piece-b} as a cap, that cap is not removed by this remark.
\end{itemize}
No floor on $L$ beyond $L\ge L_{0}$ of Theorem~\ref{3.1:thm-main} is introduced, and no floor
on $D$, the parameter of the set
$\mathcal{C}_{S}(D)$ of clean histories in Lemma~\ref{8.4:lem-base-regularity}. No condition is
imposed on the number of renormalisation steps, on the parameter $r$ of \eqref{8.4:eq-small-piece-a},
on $N$,
on the volume, or on the age of any region.
\end{remark}

\begin{proof}
We take the conditions in turn. For each we check three things: which letters
(Definition~\ref{2.3:def-order-table}) it reads, that
these letters are chosen no later than the letter it bounds, and the direction of the bound.

(a) The constant $c_{7}=\max\{c_{\Gamma}L^{2},\,2+8c',\,1\}$ of
\eqref{8.4:eq-base-regularity-a} depends only on $L$ and on constants fixed before the
couplings. Since $c_{7}\ge1>0$, \eqref{8.4:eq-one-sided-conditions-1} holds if and only if
$c_{7}\varepsilon_{n}\le\min(a_{1},1)$ for each $0\le n\le j$. With
$\varepsilon_{n}=g_{n}A_{0}p_{0}(g_{n})$ this is the displayed inequality for the couplings.
Since the degree function $p_{0}(\cdot)$ grows only like a power of $\log g'^{-2}$, we have
$g'p_{0}(g')^{m}\to0$ as $g'\to0$ for every fixed $m$; with $m=1$, the inequality holds for all
sufficiently small $g_{n}$. It is
therefore a ceiling, of the form in which Ba{\l}aban requires $g_{k}$ to be sufficiently small
\cite{B-CMP119}. Finally, $\min(a_{1},1)\le a_{1}$, so \eqref{8.4:eq-far-member-a} follows.

(b) The threshold $y_{*}$ reads only $b_{D}$, $c_{h}$, $c_{D}$, $a_{9}$, all chosen before the
couplings. Inside the window $Y\le c_{Y}p_{0}(g_{j})^{m_{Y}}$ the left side of
\eqref{8.4:eq-one-sided-conditions-2} is at most $c_{C}c_{Y}\,g_{j}p_{0}(g_{j})^{m_{Y}}$.
By the property of the degree function used in (a), with $m=m_{Y}$,
this tends to $0$ as $g_{j}\to0$. Hence \eqref{8.4:eq-one-sided-conditions-2} holds for all
sufficiently small $g_{j}$, which is a ceiling on $g_{j}$.

(c) Taking logarithms, and using $\mu=\delta_{0}c_{d}/16>0$,
\begin{equation*}
L^{4}e^{-\mu M}\le\tfrac12
\iff 4\log L-\mu M\le-\log2
\iff M\ge\frac{4\log L+\log2}{\mu}=M_{0}(L).
\end{equation*}
The threshold $M_{0}(L)$ reads only $L$, $\delta_{0}$ and $c_{d}$, all chosen before $M$. The
condition is a lower bound on $M$ alone.

(d) By (b), $g_{j}Y\le y_{*}/c_{C}$. Hence \eqref{8.4:eq-one-sided-conditions-4} is implied by
\begin{equation}\label{8.4:eq-one-sided-conditions-6}
W^{d-1}e^{-2\nu W}\le\frac{\theta_{*}\,c_{C}}{y_{*}\,c_{4}'C_{2\nu}}\,e^{-\delta_{0}c_{1}/8}.
\end{equation}
The derivative of $W\mapsto W^{d-1}e^{-2\nu W}$ is $W^{d-2}e^{-2\nu W}(d-1-2\nu W)$. It is
negative for $W>(d-1)/(2\nu)$, and the function tends to $0$ as $W\to\infty$. So
\eqref{8.4:eq-one-sided-conditions-6} holds for every $W$ not smaller than a threshold. That
threshold depends only on $d$, $\nu$, $\theta_{*}$, $y_{*}$, $c_{C}$, $c_{4}'$, $C_{2\nu}$,
$\delta_{0}$ and
$c_{1}=c_{d}(4L+2)+c_{d}'$, hence on $d$, $L$, $N$ and Ba{\l}aban's constants. All of these are
chosen before $W$. The condition is a lower bound on $W$.

(e) Condition \eqref{8.4:eq-one-sided-conditions-5} bounds the response size
$\Theta_{\tilde X_{f}}$ at the member, through the response constant $K$ and the radii of the
member. It imposes no lower bound on a coupling and no upper bound on $L$, $M$
or $W$. For members beyond the collar it follows, as stated in (e), from a lower bound on $W$
of the kind in (d).

(f) By the small-piece statement (Proposition~\ref{8.4:prop-small-piece}, whose proof is
incomplete at one step), the condition \eqref{8.4:eq-small-piece-b} takes the form
\begin{equation*}
\nu W-(d-1)\log W\ge\log\frac{K_{L}g_{j}(\bar\varepsilon+\theta_{*})+K_{M}}{\mathfrak{b}} .
\end{equation*}
The left side has derivative $\nu-(d-1)/W$. It is therefore increasing for $W>(d-1)/\nu$ and
tends to $\infty$. The condition is thus a lower bound on $W=100M\check R_{j}$ by a quantity
that grows with $\log(1/\mathfrak{b})$. The letters $W$ and $\mathfrak{b}$ are chosen last.
The budget $\mathfrak{b}=\mathfrak{b}(g_{j})$ depends on $g_{j}$. By the same proposition, the
condition holds for all $g_{j}\le\gamma$ and all $M\ge M_{0}(L)$ if and only if $\check R_{j}$
grows at least like $\log(1/\mathfrak{b}(g_{j}))/(100\nu M)$; in the cases listed in
Proposition~\ref{8.4:prop-small-piece}, among them $\mathfrak{b}$ a power of
$g_{j}$ or of $p_{0}(g_{j})$, this is the growth condition on the radius $\check R_{j}$. Under
that condition the lower bound is a floor, and no lower bound on $g_{j}$ is imposed. Were
$\check R_{j}$ bounded, the condition would force a lower bound on $g_{j}$, that is, a cap in the
sense of Definition~\ref{2.3:def-cap}; under the growth condition on the radius $\check R_{j}$
this does not occur.
Where Proposition~\ref{8.4:prop-small-piece} records the condition as a cap, that cap is not
removed here, and the conclusion of this paragraph is drawn only under the growth condition.

Conditions (a) and (b) are upper bounds on the couplings. Conditions (c), (d) and, under the
growth condition on the radius $\check R_{j}$, (f) are lower
bounds on $M$ and $W$, and (e) is a condition at each member which, beyond the collar, follows
from a lower bound on $W$. In (a)--(f) the letter $L$ is never itself bounded: it enters only
through thresholds fixed before the letter each condition bounds, such as $c_{7}$, $c_{1}$,
$M_{0}(L)$ and, in (f), the constants $K_{L}$, $K_{M}$, which carry factors polynomial in $L$.
None of (a)--(f) bounds $D$, the number of renormalisation steps, $r$, $N$, the volume or
the age of a region.
\end{proof}

\begin{definition}[Scales, levels and decay rates]\label{8.5:not-scales-levels}
Throughout this section $d=4$, and the following letters are used.
\begin{itemize}
\item[(a)] \emph{Blocking and localisation.} $L$ is the blocking factor. $M$ is the side of the
cubes of Ba{\l}aban's localisation partitions. It is a power of $L$ with
\begin{equation*}
M\ge L^{3}M_1 .
\end{equation*}
\item[(b)] \emph{Decay rates.} $\kappa$ is the decay rate in the distances $d_j(X)$ of
\cite{B-CMP109}.

$\kappa_R$ is the decay rate printed in \cite[(2.31)]{B-CMP119}, and $\kappa_0$ is the
exponent of the running coupling $g_j$ in the bound \cite[(2.31)]{B-CMP119}. No relation
between $\kappa_R$, $\kappa_0$ and $\kappa$ is assumed.

We put
\begin{equation}\label{8.5:eq-scales-levels-1}
\hat\kappa:=\kappa/2,\qquad \hat\kappa_R:=\kappa_R/2 .
\end{equation}
These are the decay rates of the majorant masses used below, after half of each decay rate
has been used to absorb the radius of the Cauchy estimates in the complex source $w$.
\item[(c)] \emph{Order of choice.} The letters in (a) and (b) are fixed before every other
constant of this section. They are subject to the lower bounds that the correlation theorem
places on them.
\item[(d)] \emph{Steps and levels.} $K$ is the number of renormalisation steps. $s\ge1$ is the
depth at which the two-point witness theorem is read. The corresponding level is
\begin{equation*}
k_0:=K-s ,
\end{equation*}
which we call the witness level. $m$ is the torus exponent.
\item[(e)] \emph{Units.} Positions are measured in the units of
Definition~\ref{1.1:def-lattice}: those of the final unit lattice, the block lattice reached
after all $K$ steps, on the torus $\mathbb{T}_m$ of that definition. The bare spacing is
$\varepsilon:=L^{-K}$. We write
\begin{equation*}
\ell:=L^{-s}
\end{equation*}
for the side of a cell of the level-$k_0$ lattice.
\item[(f)] \emph{Spacings after $k$ steps.} After $k$ steps the current block lattice has
spacing $L^{-(K-k)}$. Ba{\l}aban's parameter $\eta$ of \cite{B-CMP109} at step $k$ is the bare
spacing measured in the units of step $k$:
\begin{equation*}
\eta_k=L^{-k}.
\end{equation*}
A term born at level $j\le k$ lives on the lattice whose spacing in the units of step $k$ is
\begin{equation}\label{8.5:eq-scales-levels-2}
L^{j}\eta_k=L^{-(k-j)} .
\end{equation}
\item[(g)] \emph{The value at the witness level.} At $k=k_0$ we write $\eta$ for $\eta_{k_0}$.
Since $\varepsilon/\ell=L^{-K}L^{s}$, we have
\begin{equation}\label{8.5:eq-scales-levels-3}
\eta:=\eta_{k_0}=L^{-k_0}=L^{-(K-s)}=\frac{\varepsilon}{\ell} .
\end{equation}
\end{itemize}
\end{definition}

\begin{definition}[Cells, cube partitions and the partition lattice]\label{8.5:not-partition-lattice}
We use $L$, $M$, $K$, $s$, the witness level $k_0=K-s$ and the bare spacing
$\varepsilon=L^{-K}$ as fixed in Definition~\ref{8.5:not-scales-levels}.

\emph{Cells.} A \emph{cell} $c$ is a unit cube of the level-$k_0$ lattice. Its side is
\begin{equation*}
\ell:=L^{-s}=\varepsilon L^{k_0},
\end{equation*}
and its centre is denoted $x_c$. We write $\mathsf{C}\cong\mathbb{Z}^4$ for the cell lattice of
the whole-lattice family. Elements of $\mathsf{C}$ also serve as translations by whole cells.

\emph{Partitions.} For each level $j$, $\pi_j$ is Ba{\l}aban's partition of the $j$-th unit
lattice into cubes $\square\in\pi_j$ of side $M$ \cite{B-CMP109}. We fix the partitions so that
the origin is a corner of a cube of every $\pi_j$, $j\le k_0$; in level-$j$ units, $\pi_j$ is
then the set of translates, by the vectors of $M\mathbb{Z}^4$, of the cube of
side $M$ that has a corner at the origin. Measured in cells of the level-$k_0$ lattice, the
cubes of $\pi_j$ have side $M\cdot L^{-(k_0-j)}$. The localisation domains
$X\in\mathcal{D}_j$ are unions of cubes of $\pi_j$ \cite{B-CMP109}.

\emph{The partition lattice.} The \emph{partition lattice at level $k_0$} is
\begin{equation}\label{8.5:eq-partition-lattice-1}
\Gamma:=\{a\in\mathsf{C}:\ \pi_j+a=\pi_j\ \text{for every } j\le k_0\}.
\end{equation}

\emph{Period blocks.} A \emph{period block} $\mathfrak{B}$ is a cube of $\pi_{k_0}$. It
consists of $M^4$ cells.

The following two assertions are proved below:
\begin{itemize}
\item[(a)] the partitions $\pi_j$, $j\le k_0$, are nested;
\item[(b)] $\Gamma=M\mathbb{Z}^4$ in cells, and the coarsest member $\pi_{k_0}$ alone
determines $\Gamma$.
\end{itemize}

\emph{Assumption on the localisation of \cite{B-CMP119}.} We assume that the further blocks
that the localisation of \cite{B-CMP119} uses beside the cubes of the partitions $\pi_j$ have
side $M_1$ in cells, with corners on $M_1\mathbb{Z}^4$, and that $M_1$ divides $M$; and that
the domain of each of the $\beta$-extracted terms of \cite{B-CMP119} (the family
$\operatorname{Irr}_1^E$ of this chapter) is a union of cubes of one of the partitions $\pi_j$,
$j\le k_0$. Under this assumption $\Gamma$ is the same lattice for the $\beta$-extracted terms
of \cite{B-CMP119}. A family of terms indexed instead by cubes of $R_{k_0}M$ cells on a side,
where $R_{k_0}$ denotes the ratio of the side of these cubes to the side $M$ of the cubes of
$\pi_{k_0}$, would have period $R_{k_0}M$.
\end{definition}

\begin{proof}
Let $j\le k_0$. The spacing of the level-$j$ lattice is $L^{-(K-j)}$, and the side of a
cell is $L^{-(K-k_0)}$. Dividing the first by the second, a length of $M$ in level-$j$ units
equals
\begin{equation*}
\frac{M\,L^{-(K-j)}}{L^{-(K-k_0)}}=M\cdot L^{-(k_0-j)}
\end{equation*}
cells. This is the side of the cubes of $\pi_j$ measured in cells.

For the same reason, a translation by $a\in\mathsf{C}$, written in level-$j$ units, is the
translation by $L^{k_0-j}a$.

\emph{Nestedness.} Let $j<k_0$. The level-$(j+1)$ spacing is $L$ times the level-$j$ spacing.
Hence, in level-$j$ units, the cubes of $\pi_{j+1}$ have side $ML$, and their corners form
$ML\mathbb{Z}^4$. The corner set $ML\mathbb{Z}^4$ of $\pi_{j+1}$ is contained in the corner set
$M\mathbb{Z}^4$ of $\pi_j$ (recall that $M$ is a power of $L$, so that $L\mid M$).
Therefore each cube of $\pi_{j+1}$ is the union of the $L^4$ cubes of $\pi_j$ that it
contains. By induction on $j'-j$, every cube of $\pi_j$ lies in a cube of $\pi_{j'}$ whenever
$j\le j'\le k_0$. This proves (a).

\emph{The partition lattice.} Let $a\in\Gamma$. The cubes of $\pi_{k_0}$ have side $M$
cells, and their corners form $M\mathbb{Z}^4$ in cells. Hence $\pi_{k_0}+a=\pi_{k_0}$ holds
exactly when $a\in M\mathbb{Z}^4$. This shows $\Gamma\subseteq M\mathbb{Z}^4$.

Conversely, let $a\in M\mathbb{Z}^4$ and $j\le k_0$. In level-$j$ units the translation is
by $L^{k_0-j}a$, and
\begin{equation*}
L^{k_0-j}a\in ML^{k_0-j}\mathbb{Z}^4\subseteq M\mathbb{Z}^4 .
\end{equation*}
Every vector of $M\mathbb{Z}^4$ maps the set of translates defining $\pi_j$ onto itself, so
$\pi_j+a=\pi_j$. Hence $a\in\Gamma$. Thus $\Gamma=M\mathbb{Z}^4$, and membership in $\Gamma$
is decided by $\pi_{k_0}$ alone. This proves (b).

\emph{The $\beta$-extracted terms of \cite{B-CMP119}.} Under the assumption just stated, the
domain of each of these terms is a union of cubes of some $\pi_j$, $j\le k_0$, and by (b) a
translation by a vector $a$ of $M\mathbb{Z}^4$ in cells satisfies $\pi_j+a=\pi_j$ for every
$j\le k_0$, hence maps unions of cubes of $\pi_j$ to unions of cubes of $\pi_j$. The further
blocks have side $M_1$ and corners on $M_1\mathbb{Z}^4$, and $M_1$ divides $M$. The corner set
$M_1\mathbb{Z}^4$ of these blocks therefore contains $M\mathbb{Z}^4$, and the same translations
map the family of these blocks onto itself. So $\Gamma$ is unchanged for these terms.

Finally, for the partition of the level-$k_0$ lattice into cubes of side $R_{k_0}M$ cells with
a corner at the origin, the argument for (b) with $M$ replaced by $R_{k_0}M$ shows that the cell
translations preserving it form $R_{k_0}M\mathbb{Z}^4$.
\end{proof}

\begin{definition}[Separation scale, support geometry and plane-independent weights]
\label{8.5:def-separation-scale}
Let $k_0=K-s$ be the witness level, $\ell=L^{-s}$ the side of a cell and $x_c$ the centre of a
cell $c$, as in Notation~\ref{8.5:not-partition-lattice}. For a plaquette $p$ we write $x(p)$
for its barycentre and $\pi(p)$ for the plane of $\mathbb{R}^4$ (spanned by two coordinate axes)
to which $p$ is parallel; $\mathfrak{O}$ denotes the set of the six such planes.
\begin{itemize}
\item[(a)] \emph{Separation scale.} We put
\begin{equation}\label{8.5:eq-separation-scale-1}
\mathsf{d}:=C_W\,\ell ,\qquad C_W\ge C_0 .
\end{equation}
The constant $C_W$ is a floor in the sense of Notation~\ref{2.3:not-stages-first}: it is chosen
after $M$, $\kappa$, $\kappa_R$, $\rho_\sharp$, $\Theta_0$, $C_\nabla$, $C_\sigma$, $A_c$, $r_w$ and
the remaining constants on which the estimates of this section depend, and no condition bounds it
from above. Here $A_c$ and $\Theta_0$ are the constants of items~(b) and~(e) below, $C_\nabla$ is the
constant of the gradient bound of this section, $C_\sigma$ the constant of its growth bound, and
$r_w$ the parameter whose inverse is the level-free route weight. The choice of $C_W$ is also
subject to the accompanying condition on the period of the localisation domains stated with the
partition lattice $\Gamma$ in Notation~\ref{8.5:not-partition-lattice}.
\item[(b)] \emph{Geometry.} The two centres are points $x_1,x_2$ of the torus with
\begin{equation}\label{8.5:eq-separation-scale-2}
3\mathsf{d}\;\le\;|x_1-x_2|\;\le\;A_c\,\mathsf{d},
\end{equation}
where $A_c\ge4$ is a pure number, the \emph{centre-separation constant}. In the geometry of the
two-point witness theorem, $A_c=8$.
\item[(c)] \emph{Balls and supports.} We put
\begin{equation*}
B_1:=B\bigl(x_1,\tfrac54\mathsf{d}\bigr),\qquad B_2:=B\bigl(x_2,\tfrac54\mathsf{d}\bigr).
\end{equation*}
The \emph{marked weight} $f_1$ is a real, possibly signed, Lipschitz function of position. It
enters only through the plaquette weight $p\mapsto f_1(x(p))$, and $\operatorname{supp}f_1\subset B_1$.
The weight $h_2$, which occupies the Gaussian slot of the second-derivative room lemma, is
Lipschitz, and $\operatorname{supp}h_2\subset B_2$ up to an exponentially small
profile tail that decays with the distance. This tail arises where $h_2$ is obtained from the
sourced format of the correlation theorem, and it is estimated separately in clause~(vi) of the
second-derivative room lemma.
\item[(d)] \emph{Plane independence.} The marked weight $f_1$ is \emph{plane-independent}: the
value it assigns to $p$ depends on $x(p)$ only and not on $\pi(p)$. This is the class of the
smeared curvature field $\mathcal{O}(f)$ of Definition~\ref{1.3:def-curvature-field}. It is also
the class of the source weights of the correlation theorem. The second test function $f_2$, from
which $h_2$ is derived, is plane-independent as well.
\item[(e)] \emph{Shape functional and class.} For a Lipschitz function $f$, a point $x$ and a
radius $R>0$ we put
\begin{equation}\label{8.5:eq-separation-scale-3}
\Theta(f;x,R):=\frac{R\cdot\operatorname{Lip}f\cdot|B(x,R)|}{\|f\|_1},\qquad
\mathfrak{F}(x,R;\Theta_0):=\bigl\{f:\ \Theta(f;x,R)\le\Theta_0\bigr\},
\end{equation}
where $|B(x,R)|$ is the volume of the ball and $\|f\|_1$ the integral of $|f|$. The class
$\mathfrak{F}(x,R;\Theta_0)$ increases with $\Theta_0$. We may therefore assume $\Theta_0\ge1$.
\end{itemize}
\end{definition}

\emph{Use of the upper bound in \eqref{8.5:eq-separation-scale-2}.} The upper bound
$|x_1-x_2|\le A_c\mathsf{d}$ enters in two places only, and in neither is it used as a cap on a
constant chosen as a floor.

The first place is the maximal support distance
\begin{equation}\label{8.5:eq-separation-scale-4}
D_{\max}:=\bigl(A_c+\tfrac52\bigr)\mathsf{d}.
\end{equation}
Let $y_1\in B_1$ and $y_2\in B_2$. By the triangle inequality and \eqref{8.5:eq-separation-scale-2},
\begin{equation*}
|y_1-y_2|\le|y_1-x_1|+|x_1-x_2|+|x_2-y_2|
\le\tfrac54\mathsf{d}+A_c\mathsf{d}+\tfrac54\mathsf{d}=D_{\max}.
\end{equation*}
The quantity $D_{\max}$ is used in the conversion step of this section. There it enters only
through a pure number that depends on $(A_c,\Theta_0)$ alone.

The second place is in the entropy bound of this section, where the spacing ratio $\eta$, the
quantity $\mathfrak{r}$ and the torus clause of the two-point witness theorem are read at plaquette
pairs whose distance is at most $(A_c+3)\mathsf{d}$.

\emph{The centre-separation constant and the pair factor.} The constant $A_c$ bounds the distance
of the two centres. It is distinct from the plaquette-pair factor $A_*$ of the two-point witness
theorem, which bounds the distance of the plaquettes of a pair: $\operatorname{dist}(p,p')\le A_*\mathsf{d}$.
For plaquettes whose barycentres lie in balls of radii $a_1\mathsf{d}$ and $a_2\mathsf{d}$ about
$x_1$ and $x_2$, the triangle inequality gives
$\operatorname{dist}(p,p')\le(A_c+a_1+a_2)\mathsf{d}$. On the geometry of that theorem $A_*=10$, and
on its variant geometry $A_*=21/2$.

The upper separation bound used in the decay estimate of that theorem is this plaquette-pair bound.
It is not the window $C_0\le C_W<C_0L$. That window expresses the choice of $s$ from $\mathsf{d}$
made in the construction of the witness, and it enters no constant. In this section it is read once,
to place the two members of the clauses on $m$ in the entropy bound and in the gradient bound of this
section, and nowhere else.

\emph{Plane-dependent weights.} We explain why plane independence is imposed on $f_1$ in
item~(d).

Suppose instead that $f_1=(f_1^P)_{P\in\mathfrak{O}}$ is a plane-dependent plaquette weight. The
two-point witness theorem admits such weights in the additivity clause of its weight class
$\mathfrak{F}$, whose shape condition is the bound $\Theta\le\Theta_0$ of item~(e),
although its relative clause is scalar. For such $f_1$ the locally constant approximant of $f_1$ on
the reading neighbourhood $X^+$ is the anisotropic constant profile
$p\mapsto f_1^{\pi(p)}(x_c)$, where $c=c(X)$ is the anchor cell.

Let $\mathbf{1}_P$ denote the profile equal to $1$ on the plaquettes parallel to $P$ and to $0$ on
the others. The sum over the domains anchored at $c$ of the constant-flux forms then reads
\begin{equation}\label{8.5:eq-separation-scale-5}
\sum_{c(X)=c}q_X[f_1]=\sum_{P\in\mathfrak{O}}f_1^P(x_c)\,q_P(c)[\mathbf{1}_P]+e(c).
\end{equation}
Here $q_P(c)[\mathbf{1}_P]$ is the corresponding sum of the forms marked by $\mathbf{1}_P$, and
$e(c)$ is the contribution of the difference between $f_1$ and its approximant. The relevant
residual family is the $\mathbf{1}_P$-marked family
\begin{equation*}
\partial_P\mathcal{R}^E:=\{D[\mathbf{1}_P]V_X\}.
\end{equation*}

The mean-zero identity of the residual family is proved at the $W_4$-scalar profile
$\mathbf{1}=\sum_P\mathbf{1}_P$ only. What it gives is therefore
\begin{equation}\label{8.5:eq-separation-scale-6}
\sum_{P\in\mathfrak{O}}\mathfrak{q}_M\bigl(\partial_P\mathcal{R}^E\bigr)=0,
\end{equation}
where $\mathfrak{q}_M$ is the block mean. It does not give
$\mathfrak{q}_M(\partial_P\mathcal{R}^E)=0$ for each $P$.

Each block mean $\mathfrak{q}_M(\partial_P\mathcal{R}^E)$ is an element of
$\operatorname{Sym}^2(\mathbb{R}^{\mathfrak{O}})$. The sign changes of the coordinate axes map every
plane to itself, so each $\mathfrak{q}_M(\partial_P\mathcal{R}^E)$ is diagonal, of the form
\begin{equation*}
\lambda_0\,e_P\otimes e_P+\lambda_1\sum_{|Q\cap P|=1}e_Q\otimes e_Q
+\lambda_2\,e_{P^c}\otimes e_{P^c}.
\end{equation*}
Here $|Q\cap P|$ is the number of axes shared by $Q$ and $P$, and $P^c$ is the plane spanned by the
two axes not in $P$.

Fix a plane $Q$ and count the planes $P$ in each case. The plane $Q$ equals $P$ for exactly one $P$.
It shares exactly one axis with $P$ for exactly four $P$. It equals $P^c$ for exactly one $P$.
Hence the coefficient of $e_Q\otimes e_Q$ in the sum over $P$ is $\lambda_0+4\lambda_1+\lambda_2$,
and \eqref{8.5:eq-separation-scale-6} says exactly that
\begin{equation*}
\lambda_0+4\lambda_1+\lambda_2=0 .
\end{equation*}
This does not force $\lambda_0$, $\lambda_1$ and $\lambda_2$ to vanish.

An isotropic coupling correction therefore leaves the anisotropic dimension-four content of an
anisotropic coupling shift unextracted. For that part the zero-mean step in the proof of the
second-derivative room lemma has no zero mean, and no room is obtained from it here.

Accordingly, the irrelevance hypothesis of the two-point witness theorem, in the form supplied by
this section, serves its additivity clause only over plane-independent weights. There are two
possible cases.
\begin{itemize}
\item[(1)] The weight class $\mathfrak{F}$ of that theorem is taken plane-independent. This matches
the scalar observable $\mathcal{O}(f)$ of Definition~\ref{1.3:def-curvature-field}, the only source
functional read by the corresponding clause of the correlation theorem. The kernel lemmas of the
two-point witness theorem may keep their greater generality.
\item[(2)] One has in addition an anisotropic extraction clause: for every plane $P$ the block mean
$\mathfrak{q}_M(\partial_P\mathcal{R}^E)$ of the $\mathbf{1}_P$-marked residual family vanishes,
established under the hypotheses of the correlation theorem.
\end{itemize}
Ba{\l}aban's coupling correction is a single $W_4$-invariant scalar per level
(\cite[(1.22)]{B-CMP109}, \cite[(3.61)]{B-CMP119}). It is not a clause of the kind in~(2).

For the weight $h_2$, the hypotheses of the second-derivative room lemma are a sup bound and a
distance tail on $h_2$, and any plaquette weight satisfying them is admissible there. The derivation
of these two bounds from the sourced format of the correlation theorem is carried out for the scalar
class of Definition~\ref{1.3:def-curvature-field}. This is why $f_2$ is taken plane-independent in
item~(d).

\begin{definition}[Planes, symmetric forms and the hyperoctahedral action]
\label{8.5:not-planes-forms}
\emph{Planes and constant fluxes.} We write
\begin{equation*}
\mathfrak{O}:=\bigl\{P=\{\mu,\nu\}:\ 1\le\mu<\nu\le4\bigr\}
\end{equation*}
for the set of the six coordinate planes. The space $\mathbb{R}^{\mathfrak{O}}$ has the basis $(e_P)_{P\in\mathfrak{O}}$. We identify it with $\Lambda^2\mathbb{R}^4$ by sending $e_P$, for $P=\{\mu,\nu\}$ with $\mu<\nu$, to $e_\mu\wedge e_\nu$. Here $e_1,\dots,e_4$ is the standard basis of $\mathbb{R}^4$.

For $c',c'',c\in\mathbb{R}^{\mathfrak{O}}$ we put
\begin{equation*}
\langle c',c''\rangle:=\sum_{P\in\mathfrak{O}}c'_P\,c''_P,
\qquad
|c|_\infty:=\max_{P\in\mathfrak{O}}|c_P|.
\end{equation*}
For a plaquette $p$ we write $\pi(p)\in\mathfrak{O}$ for the plane of $p$. A \emph{constant flux} is an element $c\in\mathbb{R}^{\mathfrak{O}}$. It is read as the $2$-cochain $p\mapsto c_{\pi(p)}$ on whichever box it is restricted to.

\emph{Symmetric forms.} We write $\operatorname{Sym}^2(\mathbb{R}^{\mathfrak{O}})$ for the space of symmetric bilinear forms $q$ on $\mathbb{R}^{\mathfrak{O}}$. On it we use the norm
\begin{equation*}
|q|:=\sup\bigl\{|q[c',c'']|:\ |c'|_\infty\le1,\ |c''|_\infty\le1\bigr\}.
\end{equation*}
For $u,v\in\mathbb{R}^{\mathfrak{O}}$ we put $\langle q,u\otimes v\rangle:=q[u,v]$.

\emph{The hyperoctahedral group.} The hyperoctahedral group $W_4$ is the group of signed permutations of the four axes. Each $r\in W_4$ determines a permutation $\mu\mapsto r(\mu)$ of $\{1,2,3,4\}$ and signs $\sigma_\mu\in\{\pm1\}$, and acts on $\mathbb{R}^4$ by
\begin{equation*}
r\,e_\mu=\sigma_\mu\,e_{r(\mu)}.
\end{equation*}

It acts on $\mathbb{R}^{\mathfrak{O}}$ by the induced signed permutation of planes:
\begin{equation*}
r\,(e_\mu\wedge e_\nu)=\sigma_\mu\sigma_\nu\,e_{r(\mu)}\wedge e_{r(\nu)}.
\end{equation*}
Thus, with $P'=\{r(\mu),r(\nu)\}$, the right-hand side equals $\sigma_\mu\sigma_\nu\,e_{P'}$ if $r(\mu)<r(\nu)$ and $-\sigma_\mu\sigma_\nu\,e_{P'}$ if $r(\mu)>r(\nu)$.

On $\operatorname{Sym}^2(\mathbb{R}^{\mathfrak{O}})$ the group acts by
\begin{equation}\label{8.5:eq-planes-forms-1}
(r\cdot q)[c',c'']:=q\bigl[r^{-1}c',\,r^{-1}c''\bigr].
\end{equation}

\emph{The action on the lattice.} On the lattice, coordinate reflections are taken about cube centres, with the convention of \cite{B-CMP119} for bonds. That is, the reflection in the $\mu$-th coordinate is
\begin{equation*}
r x=(x_1,\dots,-x_\mu,\dots,x_4)-e_\mu .
\end{equation*}
Permutations of the axes are taken about a cube corner.

These maps are used generator by generator. They are not used as a single point-group action with a common fixed point. Let $\pi_j$ be Ba{\l}aban's cube partitions and $\Gamma$ the partition lattice, as in Notation~\ref{8.5:not-partition-lattice}. Each generator permutes the cubes of every $\pi_j$, normalises $\Gamma$, and maps $\Gamma$-classes to $\Gamma$-classes.
The lemma of this chapter on the action of $W_4$ uses the sign changes and the permutations separately.
\end{definition}

\begin{definition}[Bond fields and affine test fields of constant flux]\label{8.5:not-affine-test-field}
In this section $\mathfrak{g}=\mathfrak{su}(2)$; item~(c) of the lemma on constant fluxes holds for any
simple compact Lie algebra $\mathfrak{g}$.

\emph{Bond fields.} The letters $B$, $B'$ denote $\mathfrak{g}$-valued $1$-cochains on the bonds of the
bare lattice: the value on a reversed bond is the negative of the value on the bond. The
corresponding bond variables are $U=\exp iB$, that is, $U(b)=\exp\bigl(iB(b)\bigr)$ for every bond $b$.
The letter $\lambda$ denotes a $\mathfrak{g}$-valued $0$-cochain on the sites. Write
$\varepsilon=L^{-K}$ for the bare spacing, $e_\nu$ for the unit vector of axis $\nu$ (not to be
confused with the plane basis vectors $e_P$ of Definition~\ref{8.5:not-planes-forms}), and
$B_\nu(x)$ for the value of $B$ on the bond from $x$ to $x+\varepsilon e_\nu$. The coboundaries are
the usual ones: $d\lambda$ assigns $\lambda(x+\varepsilon e_\nu)-\lambda(x)$ to that bond, and for the
plaquette $p$ with lower corner $x$, spanned by $e_\mu,e_\nu$ ($\mu<\nu$) and positively oriented,
\begin{equation}\label{8.5:eq-affine-test-field-1}
dB(p)=B_\mu(x)+B_\nu(x+\varepsilon e_\mu)-B_\mu(x+\varepsilon e_\nu)-B_\nu(x).
\end{equation}

\emph{Affine test fields.} For $c\in\mathbb{R}^{\mathfrak{O}}$ (Definition~\ref{8.5:not-planes-forms})
put $c_{\mu\nu}:=c_P$ for $P=\{\mu,\nu\}$ with $\mu<\nu$, $c_{\nu\mu}:=-c_{\mu\nu}$ and
$c_{\mu\mu}:=0$. For a site $y$ the \emph{affine test field} is the $1$-cochain
\begin{equation}\label{8.5:eq-affine-test-field-2}
B^{c,y}_\nu(x):=\frac12\sum_{\mu=1}^{4}c_{\mu\nu}\,(x-y)_\mu ,
\end{equation}
defined for $x$ in any box of the bare lattice, the coordinates of $x$ and $y$ being read in one
chart of $\mathbb{T}_m$ containing the box, where $(x-y)_\mu$ is the $\mu$-th coordinate
of $x-y$ counted in bare lattice steps. Then
\begin{equation}\label{8.5:eq-affine-test-field-3}
dB^{c,y}\equiv c ,
\end{equation}
that is, $dB^{c,y}(p)=c_{\pi(p)}$ for every positively oriented plaquette $p$ of the box: $B^{c,y}$
has the constant flux $c$ in the sense of Definition~\ref{8.5:not-planes-forms}.
\end{definition}

\begin{proof}
Let $p$ have lower corner $x$ and span $e_\mu,e_\nu$, $\mu<\nu$, so that $\pi(p)=\{\mu,\nu\}$.
Passing from $x$ to $x+\varepsilon e_\mu$ raises $(x-y)_\mu$ by one and leaves the other coordinates
unchanged. By \eqref{8.5:eq-affine-test-field-2},
\[
B^{c,y}_\nu(x+\varepsilon e_\mu)-B^{c,y}_\nu(x)=\tfrac12 c_{\mu\nu}.
\]
In the same way,
\[
B^{c,y}_\mu(x+\varepsilon e_\nu)-B^{c,y}_\mu(x)=\tfrac12 c_{\nu\mu}=-\tfrac12 c_{\mu\nu}.
\]
Inserting both differences into \eqref{8.5:eq-affine-test-field-1} gives
\[
dB^{c,y}(p)=\tfrac12c_{\mu\nu}+\tfrac12c_{\mu\nu}=c_{\mu\nu}=c_{\pi(p)} .
\]
\end{proof}

\begin{definition}[Localised terms: hull, anchor, radius and fine-field Hessian]
\label{8.5:def-localised-terms}
We use the cells $c$ of side $\ell$, their centres $x_c$, the cell lattice $\mathsf{C}$ and the
partition lattice $\Gamma$ of Notation~\ref{8.5:not-partition-lattice}. We also use the bond fields
$B$, $B'$ and the relation $U=\exp iB$ of Definition~\ref{8.5:not-affine-test-field}. A
\emph{partition cube} is a cube of the partition $\pi_{k_0}$, that is, a cube of side $M$ cells
with corners in $\Gamma$.

\emph{Terms.} A \emph{term} is a function $T$ of the bond variables $U$ with the following properties.
It is gauge invariant. It is analytic in a neighbourhood of the identity configuration $\mathbf{1}$.
It depends on $U$ only through the restriction $U|_X$, where $X$ is a bounded union of partition
cubes; we say that $T$ is localised in $X$.

\emph{Filled hull.} Let $X^+$ be the reading neighbourhood of $X$, defined later in this chapter. The
\emph{filled hull} $\hat{X}$ is the smallest box of cells containing $X^+$. A box of cells is a set of
cells whose coordinates in $\mathsf{C}\cong\mathbb{Z}^4$ range over a product of four integer
intervals. The intersection of two such boxes containing $X^+$ is again such a box, so the smallest
one exists: its coordinate intervals run, in each direction, from the least to the greatest
coordinate of the cells meeting $X^+$.

\emph{Anchor.} Order the bare lattice sites lexicographically by their coordinates in the whole-lattice
family. Since $\hat{X}$ is a box, its lexicographically least site is its corner with least
coordinates in every direction. A cell of $\hat{X}$ containing this corner has it as its own corner
with least coordinates, so exactly one cell of $\hat{X}$ contains it. The \emph{anchor} $c(X)$ is the
cell containing the lexicographically least site of $\hat{X}$; in particular $c(X)\subset\hat{X}$.

For every $a\in\mathsf{C}$ we have $c(X+a)=c(X)+a$. To see this, note that the reading
neighbourhood, defined later in this chapter, commutes with translations by cell vectors,
$(X+a)^+=X^++a$, and note three further facts.
\begin{itemize}
\item[(a)] The smallest box of cells containing a translate of a set is the translate of the smallest
box containing the set, so $\widehat{X+a}=\hat{X}+a$.
\item[(b)] A translation by a cell vector $a$ maps bare sites to bare sites, because $\ell/\varepsilon$
is an integer. It preserves the lexicographic order, so it maps the least site of $\hat{X}$ to the
least site of $\hat{X}+a$.
\item[(c)] The same translation maps the cell of $\hat{X}$ containing that site to the cell of
$\hat{X}+a$ containing its image.
\end{itemize}

\emph{Radius.} Centres of cells differ by integer multiples of $\ell$ in each coordinate. The
\emph{radius} of $X$ is the integer
\begin{equation}\label{8.5:eq-localised-terms-1}
r(X):=\max\Bigl\{1,\ \max_{c'\subset\hat{X}}\ \ell^{-1}\,\bigl|x_{c'}-x_{c(X)}\bigr|_\infty\Bigr\},
\end{equation}
the radius of $\hat{X}$ in cells about $x_{c(X)}$ in the maximum norm $|\cdot|_\infty$. By definition
$r(X)\ge1$.

Euclidean distances within $\hat{X}$ from $x_{c(X)}$ are bounded in terms of $r(X)$. Every point $y$
of $\hat{X}$ lies in a cell $c'\subset\hat{X}$, and $|y-x_{c'}|_\infty\le\ell/2$. By
\eqref{8.5:eq-localised-terms-1} we therefore have
$|y-x_{c(X)}|_\infty\le(r(X)+\tfrac12)\ell$. In $\mathbb{R}^4$ the Euclidean norm is at most
$\sqrt4=2$ times the maximum norm $|\cdot|_\infty$, and $2(r(X)+1)\le4r(X)$ because $r(X)\ge1$. Hence
\begin{equation}\label{8.5:eq-localised-terms-2}
|y-x_{c(X)}|\le 2\bigl(r(X)+1\bigr)\ell\le 4\,r(X)\,\ell
\qquad\text{for all } y\in\hat{X}.
\end{equation}

\emph{Fine-field Hessian.} For a term $T$ and bond fields $B$, $B'$ we put
\begin{equation}\label{8.5:eq-localised-terms-3}
\operatorname{Hess}T[B,B']
:=\partial_\sigma\partial_\tau\,T\bigl(\exp i(\sigma B+\tau B')\bigr)\Big|_{\sigma=\tau=0}.
\end{equation}
This is well defined because $T$ is analytic near $\mathbf{1}$. It is the Hessian in the fine field,
that is, in the bare bond variables $U=\exp iB$, and not in the block fields. The background
configuration of Definition~\ref{1.5:def-domains}, as a function $U_{k_0}(\cdot)$ of the block field
of level $k_0$, enters later in this chapter through its differential at the identity configuration,
the linearised minimiser $M_{k_0}$:
\begin{equation*}
DU_{k_0}(\mathbf{1})=M_{k_0}.
\end{equation*}
The quadratic forms $Q_T$, $q_T$ and the mass $\mathfrak{m}_T$ of a term are defined later in this
chapter.
\end{definition}

\begin{definition}[The frozen sourced family at the witness level]\label{8.5:def-frozen-sourced-family}
Let $k_0=K-s$ be the witness level, $\ell=L^{-s}$ the side of its cells, $\varepsilon=L^{-K}$ the
bare spacing and $\eta=\varepsilon/\ell$. Terms, filled hulls $\hat{X}$, reading neighbourhoods
$X^+$, anchors $c(X)$, radii $r(X)$ and the fine-field Hessian $\operatorname{Hess}$ are those of
Definition~\ref{8.5:def-localised-terms}. For a term $T$, $Q_T$ is its flux form, $q_T$ its
constant-flux form, $\mathfrak{m}_T$ its mass in the all-direction operator norm and
$\mathfrak{m}^{\mathrm{reg}}_T$ its mass on regular directions, all as in the subsection on flux
forms of this chapter. We write $\operatorname{Hess}^{\mathfrak{n}}$ for the fine-field Hessian in
the normalisation of the block law of the background fluxes.

\begin{itemize}
\item[(a)] \emph{The frozen sourced family.} Let $w$ be a complex source profile. The
\emph{frozen sourced level-$k_0$ family with profile $w$} is the level-$k_0$ family of the
sourced-format lemma of the correlation theorem, with profile $w$; \emph{frozen} means that every
threshold, cube size, decomposition, background and reference law is that of the run at $w=0$,
and that the family is the whole-lattice one. On the small-field region
near $B(x_1,2\mathsf{d})$ its density is
\begin{equation}\label{8.5:eq-frozen-sourced-family-1}
\exp\Bigl[-\sum_p\bigl(x_{k_0}-\zeta_{k_0}(p)[w]\bigr)\,O_p
+\sum_{j\le k_0}\sum_X V^{(j),w}_X(U)
+\text{(terms not meeting the region)}\Bigr],
\end{equation}
where $U=U_{k_0}(V)$, $V^{(j),w}_X$ are the localised terms of level $j$ of that family, and
$O_p$ is the density of the main term, \cite[(2.23)]{B-CMP119}. The following is the regularity
hypothesis of the two-point witness theorem on this region: the $z$-marked family present there
(with $w=zf$ as in (b)) coincides with the whole-lattice frozen family up to terms of length at
least $c_\ast C_W/M$, for a constant $c_\ast>0$, whose mass is absorbed into the constant $C_1$. It
is among the hypotheses of the second-derivative room lemma and is not proved here.

\item[(b)] \emph{Marked forms.} For a profile $f$ put $w=zf$ and
$D[f]:=\partial_z\big|_{z=0}$. The \emph{marked forms} are
\begin{equation}\label{8.5:eq-frozen-sourced-family-2}
H_X[f]:=D[f]\operatorname{Hess}V^{(j),zf}_X,
\end{equation}
where the level $j$ is read off $X$ and is suppressed from the notation. The forms $q_X[f]$ are
their constant-flux forms, and $Q_X=Q_X[f]$ their flux forms, the profile being suppressed from
the notation $Q_X$.

\item[(c)] \emph{The two classes.} The index set $\operatorname{Irr}_1$ of the localised families
of the sourced format is split as
$\operatorname{Irr}_1=\operatorname{Irr}_1^E\sqcup\operatorname{Irr}_1^R$.
The \emph{E-class} $\operatorname{Irr}_1^E$ consists of the families whose point-attributed
translation sum enters the coupling increment extracted at some step $\le k_0$. These are the
terms $E^{(j)}$ of \cite[(1.7), (1.20)]{B-CMP109} and $E$ of \cite[(3.50), (3.57)]{B-CMP119}. In
the sourced format of the correlation theorem they are the $\beta^0$-, $(\mathsf{P}-L)$- and
$Q$-families that enter the decomposition of the coupling increment in the running-shift
statement (families so named in the running-shift statement; this $L$ is not the blocking factor,
and this $Q$ is not a quadratic form of this definition).
The \emph{R-class} $\operatorname{Irr}_1^R$ consists of the remaining families. These are the
R-terms of \cite[(2.30)--(2.31)]{B-CMP119}, the aged large-field remainders not passed
through the extraction that produces the E-class, and the frozen remainder family $R_\zeta$ of the
sourced format.
The dichotomy is required to be exhaustive in the following sense: every non-extracted $z$-marked
family of the level-$k_0$ sourced format belongs to the class of \cite[(2.31)]{B-CMP119}; the
frozen families of that format are $E_\zeta$, $P^\flat_\zeta$, $Q_\zeta$, $R_\zeta$. This
exhaustiveness is read from the running-shift statement; it is taken as a hypothesis wherever the
dichotomy is used and is not proved here.

\item[(d)] \emph{The first quadratic form.} Let $f_1$ be the first test function. Put
\begin{equation}\label{8.5:eq-frozen-sourced-family-3}
Q_1:=\sum_\iota\sum_X H_X[f_1][\cdot,\cdot],
\end{equation}
evaluated on the background fluxes of the block law, where $\iota$ runs over the index set of the
outer sum in the block law of the background fluxes; $\iota$ is a summation index only, distinct
from the variables of the Hessian of Definition~\ref{8.5:def-localised-terms}.
Splitting the sum over $X$ according to (c) gives $Q_1=Q_1^E+Q_1^R$.
We record three conventions explicitly.
\begin{itemize}
\item[--] The quadratic part of a term in the fine field is $\tfrac12\operatorname{Hess}$ (the
factor-$2$ convention of the Hessian of Definition~\ref{8.5:def-localised-terms}).
\item[--] The block law carries the factor $g_{k_0}^2(2N)$ times the trace normalisation.
\item[--] $M_{k_0}=DU_{k_0}(\mathbf{1})$ (see (j)) is $\operatorname{Ad}$-equivariant. By Schur's
lemma for the simple Lie algebra
$\mathfrak{g}$, this is used for the reduction to a single colour exactly as in the block law.
\end{itemize}
Accordingly the same factor $g_{k_0}^2/2$ stands in front of $Q_1$ and in front of the second
quadratic form $Q_2$, defined in the block-law subsection. By the flux-form statements of this
chapter, by the statement that $U_{k_0}$ is $C^2$ at the flat configuration (recalled in (j)),
and by the colour reduction of the block law, every gauge-invariant analytic $T\circ U_{k_0}$ has
quadratic part $(g_{k_0}^2/2)\operatorname{Hess}^{\mathfrak{n}}T[\psi,\psi]$ in $b$, the Gaussian
letters $b$ and $\psi$ being those of (j). Hence the Wick computation of the block-law subsection
produces the main term's own prefactor $b_0g_{k_0}^4$, and the conversion to the normalisation of
the two-point witness theorem is free of $g$.

\item[(e)] \emph{Constant shift.} For a constant $\zeta$ put
$V^\zeta_X:=V^{(j),w\equiv\zeta}_X$; at $\zeta=0$ this is the unmarked term, written
$V^{(j)}_X:=V^{(j),w\equiv0}_X$. We write $\mathcal{R}^{E,\zeta}$ for the residual E-family at
shift $\zeta$.

\item[(f)] \emph{The main-term mass per cell.} Let $\mathfrak{n}$ be $\operatorname{Hess}^{\mathfrak{n}}A$
per level-$k_0$ cell, on unit constant flux per bare plaquette:
\begin{equation}\label{8.5:eq-frozen-sourced-family-4}
\mathfrak{n}=6\,(\ell/\varepsilon)^4 .
\end{equation}
The value \eqref{8.5:eq-frozen-sourced-family-4} refers to the normalisation
$\operatorname{Hess}^{\mathfrak{n}}$ of the block-law subsection; $6(\ell/\varepsilon)^4$ is also
the number of bare plaquettes in a level-$k_0$ cell. Thus $\mathfrak{n}$ is the
mass of the main term per cell, and $\mathfrak{n}\asymp(\ell/\varepsilon)^4=\eta^{-4}$.
The coupling grading of (h) bears on the ratio $\bar{\mathfrak{a}}/\mathfrak{n}$ of the
cell-aggregate masses of (i) to $\mathfrak{n}$.

\item[(g)] \emph{Majorant masses.} Let $\bar{\mathfrak{w}}$ be a level-free weight: a number fixed
before $C_0$ and read at the witness coupling $g_{K-s}$. It is fixed in
\eqref{8.5:eq-frozen-sourced-family-11}; two further readings of it are named in (h). Put
\[
\hat{\mathfrak{m}}_X:=\bar{\mathfrak{w}}\,\hat{\mathfrak{m}}^0_X ,
\]
where $\hat{\mathfrak{m}}^0_X$ is Ba{\l}aban's majorant of the flux-form mass of the unmarked
term. For an E-term of level $j$,
\begin{equation}\label{8.5:eq-frozen-sourced-family-5}
\hat{\mathfrak{m}}^0_X:=C_H\,\mathfrak{n}\,L^{-(4+\alpha/2)(k_0-j)}\,e^{-\hat\kappa d_j(X)} .
\end{equation}
Here $(L^j\eta)^{4+\alpha}e^{-\kappa d_j(X)}$ is the shape displayed in \cite[(0.29)]{B-CMP109}
for one term $V^{(j)}_X$, with $\alpha>0$ the exponent printed there (this $\alpha$ is the
exponent of \cite[(0.29)]{B-CMP109}, not one of the radii of Notation~\ref{2.1:not-prefix}). This
shape bounds a term, not a density; the count of terms over the
$\pi_j$-cubes of a cell is made in the tail count of the localised terms. Since $\eta=L^{-k_0}$,
\begin{equation}\label{8.5:eq-frozen-sourced-family-6}
(L^j\eta)^{4+\alpha}e^{-\kappa d_j(X)}=L^{-(4+\alpha)(k_0-j)}e^{-\kappa d_j(X)} .
\end{equation}
The spare half $L^{-\alpha(k_0-j)/2}$ is a margin. Because $\bar{\mathfrak{w}}$ is level-free, no
growth in the level is absorbed anywhere.
The constant $C_H=C_H(\alpha_0,M,\kappa)$ is the Hessian conversion by Cauchy's inequality in the
fine bond field, carried out in the proof below. The rate $\hat\kappa$ is half of $\kappa$.

\emph{The fine-field form of Ba{\l}aban's bound.} The majorant \eqref{8.5:eq-frozen-sourced-family-5}
is used on the fine-field domain of the following statement. It is taken as a hypothesis and is
not proved here. The residual terms $V^{(j)}_X$, composed with
\cite[(1.15)]{B-CMP109} at step $k$, are analytic on the complex fine-field neighbourhood
\[
\bigl\{\,|\partial U-1|<\alpha_0\eta_k^2,\ U\text{ regular at scale }L^j\eta_k\,\bigr\}
\subset\mathfrak{U}^c_k ,
\]
regularity being meant in the sense of \cite[(3.32), (4.4)]{B-CMP109}. On this neighbourhood they
satisfy the bound $(L^j\eta)^{4+\alpha}e^{-\kappa d_j(X)}$ of \cite[(0.29)]{B-CMP109}.
This hypothesis is the fine-field form of the passage of \cite{B-CMP109} giving
\cite[(0.29)]{B-CMP109} in place of \cite[(0.25)]{B-CMP109} for regular gauge field
configurations, used together with \cite[(1.15)--(1.16)]{B-CMP109}. In \cite[Theorem~1]{B-CMP109} the
bound \cite[(0.29)]{B-CMP109} is stated on the
coarse small-field space $\{V:\ |\partial U_k(V)-1|<\varepsilon_0\eta^2\}$, $\varepsilon_0$ being
the constant of \cite[Theorem~1]{B-CMP109}. The spaces
$\mathfrak{U}^c_j(X,\alpha_0,\alpha_1)$ of \cite[(1.18)]{B-CMP109} carry the bound
$E_0e^{-\kappa d_j(X)}$, with the constant $E_0$ of \cite[(1.18)]{B-CMP109} and with no power of
$L^j\eta$.
The power $\alpha$ is what makes the masses $\bar{\mathfrak{a}}$ (the cell-aggregate masses of
(i)) and the sum over $j$ uniform in
$K$ and $s$. Every statement below that is uniform in $K$ and $s$ uses this fine-field form.

For an R-term of level $j$,
\begin{equation}\label{8.5:eq-frozen-sourced-family-7}
\hat{\mathfrak{m}}_X:=\bar{\mathfrak{w}}^R\,C_H\,\mathfrak{n}\,g_j^{\kappa_0}\,
L^{-4(k_0-j)}\,e^{-\hat\kappa_R d_j(X)},
\end{equation}
which is the shape of \cite[(2.31)]{B-CMP119}, with a fixed exponent $\kappa_0>0$ of that shape;
$\hat\kappa_R$ is the halved form of the rate
$\kappa_R$ printed there. The weight $\bar{\mathfrak{w}}^R$ is fixed in (h).

\emph{Domination.} Let $\|\cdot\|_X$ be the norm on regular directions of the flux forms. We
assume the following five statements, none of which is proved here:
\begin{itemize}
\item[--] the fine-field form of Ba{\l}aban's bound stated above;
\item[--] for R-terms, the bound of \cite[(2.30)--(2.31)]{B-CMP119}, read on the same
fine-field domain in the same way as \cite[(0.29)]{B-CMP109};
\item[--] the analyticity of the terms in $w$ taken as a hypothesis in (h);
\item[--] the regularity of the paired directions given by the H\"older bound
\eqref{8.5:eq-frozen-sourced-family-9}--\eqref{8.5:eq-frozen-sourced-family-10} below;
\item[--] the lower bound of the tail count of the localised terms,
\[
d_j(X)\ge (r(X)-r_{\mathrm{read}})/M-1,
\]
where $r_{\mathrm{read}}$ is the width of the reading collar.
\end{itemize}
Under these hypotheses, and with $\bar{\mathfrak{w}}$ the choice
\eqref{8.5:eq-frozen-sourced-family-11} of (h), the majorants dominate the regular-direction mass:
\begin{equation}\label{8.5:eq-frozen-sourced-family-8}
|q_X[f]|\le\mathfrak{m}^{\mathrm{reg}}_X[f]\le\hat{\mathfrak{m}}_X\sup_{X^+}|f| .
\end{equation}
Equivalently,
\[
|Q_X[u,u']|\le\hat{\mathfrak{m}}_X\,\sup|f|\,\|u\|_X\|u'\|_X
\]
for regular directions $u,u'$. The domination is not asserted for the all-direction operator norm
$\mathfrak{m}_X$. A rough unit flux direction leaves the regularity clause of the fine-field form
at radius of order $\eta^3$, not $\eta^2$. This already happens for the spaces of
\cite[(1.18)]{B-CMP109}, through the bound on $d^*F/\xi^2$, in the notation of \cite{B-CMP109},
that enters their definition.

\emph{Regular directions.} The directions paired in the two following subsections are regular.
Their regularity is the fine-oscillation H\"older bound for the linearised minimiser. This bound
is taken here by statement, in the form proved as an implication from the following four
statements:
\begin{itemize}
\item[--] the kernel bound for $\bar{M}_{k_0}$ on the torus, read from
\cite[Propositions~2.3, 2.6, 2.7]{B-CMP96}, \cite{B-CMP98} and \cite{B-CMP99b};
\item[--] the straight-segment profile statement (a statement of this book, stated with its
proof incomplete) for Ba{\l}aban's averaging \cite[(0.12)]{B-CMP109}, \cite[(1.18)]{B-CMP95};
\item[--] two statements of classical lattice potential theory.
\end{itemize}
The bound reads as follows. Let $\bar u$ be a level-$k_0$ $2$-cochain and let
$\varphi=\bar{M}_{k_0}\bar u$ be an exact fine $2$-cochain in the range
of the linearised minimiser, on a level-$k_0$ torus of side at least $8$ cells. Let $\varphi_0$ be
any fine $2$-cochain whose value at $p$ depends only on the plane $\pi(p)$ of $p$. Then for every
fine bond $\mathfrak{b}$,
\begin{equation}\label{8.5:eq-frozen-sourced-family-9}
|d^*\varphi(\mathfrak{b})|\le C_{\mathrm{reg}}\,\eta\sum_{c'}
e^{-a\operatorname{dist}(c',\mathfrak{b})}\sup_{c'}|\varphi-\varphi_0| .
\end{equation}
Moreover, let $p,p'$ be translates of each other along an axis with
$|x(p)-x(p')|\le\ell$, and let $\beta\in[0,1)$. Then
\begin{equation}\label{8.5:eq-frozen-sourced-family-10}
|\varphi(p)-\varphi(p')|\le\frac{C_{\mathrm{reg}}}{1-\beta}
\Bigl(\frac{|x(p)-x(p')|}{\ell}\Bigr)^{\beta}
\sum_{c'}e^{-a\operatorname{dist}(c',p)}\sup_{c'}|\varphi-\varphi_0| .
\end{equation}
Here $c'$ runs over cells, $\sup_{c'}$ denotes the supremum over the fine plaquettes contained in
$c'$, distances are measured in cells, and $a$, $C_{\mathrm{reg}}$ depend on
$(d,L,M)$ only. Thus the H\"older seminorm of $\varphi$ of any exponent $\beta_0\in[0,1)$ on
$\hat{X}$, and
$\eta^{-1}|d^*\varphi|$, are bounded by the cell-scale oscillation of $\varphi$ about any constant.
No power of $\eta$ is lost or gained. Put
\[
C_\Gamma':=\sup_p\sum_{c'}e^{-a\operatorname{dist}(c',p)/2},\qquad
C_\Gamma'':=\sup_p\sum_{c'}\bigl(1+\operatorname{dist}(c',p)\bigr)
e^{-a\operatorname{dist}(c',p)/2};
\]
both are numbers depending on $(d,L,M)$.
The H\"older bound is used wherever a non-constant direction is paired: in the pairing constant
$\varrho_X$ of the first line of the subsection in which the constant row-sum identity is stated,
in the pairing constant of the second-derivative room lemma, and in the covariance constants of
the defect bounds and of the remainder-column bound. The Cauchy bound in $\zeta$ is taken on
constant directions, which are
regular trivially. When the coarse-weight form is wanted, two facts are used with it: the constant
row-sum identity and the torus-versus-whole-lattice clause for the kernel of $\bar{M}_{k_0}$
described in (j). Also used is the fact that Ba{\l}aban's block average has a plane-preserving
non-negative profile, which is the content of the straight-segment profile statement.

\item[(h)] \emph{The weight.} We take as a hypothesis, not proved here, the localised form of the
sourced-format lemma of the correlation theorem: there is a radius $r_w$ such that on the sup-ball
of radius $r_w$ in $w$ the terms $V^{(j),w}_X$ are analytic in $w$, read $w$ on $X^+$ only, and
obey Ba{\l}aban's bounds at field radius $\alpha_0$. With both block-law factors displayed as in
(d), and all
Hessians in the normalisation $\operatorname{Hess}^{\mathfrak{n}}$, the domination
\eqref{8.5:eq-frozen-sourced-family-8} of the regular-direction mass holds with the unweighted
choice
\begin{equation}\label{8.5:eq-frozen-sourced-family-11}
\bar{\mathfrak{w}}:=r_w^{-1}.
\end{equation}
It follows from that hypothesis, without the coupling grading, by the Cauchy inequality in $w$
carried out in the proof below, together with the remaining hypotheses listed before
\eqref{8.5:eq-frozen-sourced-family-8}. Under this choice every statement below holds. This
includes the absolute ones, namely
the defect bounds, the second-derivative room lemma and the remainder-column bound, as well as the
relative comparison of the E-class with the main term, in terms of the cell-aggregate mass
$\bar{\mathfrak{a}}^E$ of (i):
\begin{equation}\label{8.5:eq-frozen-sourced-family-12}
\bar{\mathfrak{a}}^E/\mathfrak{n}\le C_{\mathfrak{a}}(\kappa,L,M,\alpha)\,
C_H(\alpha_0,M,\kappa)\,r_w^{-1}=:C_{\mathrm{gr}}\,r_w^{-1}.
\end{equation}
This bound is free of $g$. It is the content of the tail count of the localised terms: the number
of terms per cell is controlled by the marginal factor $(L^j\eta)^4$ of \cite[(0.29)]{B-CMP109},
and the sum over $j$ by the spare factor.

\emph{The coupling grading.} The two-point witness theorem carries a coupling-grading clause. It
is a named optional sharpening; it is not proved here, and nothing below uses it.
It has two forms.
\begin{itemize}
\item[--] Under the first grading, $\mathfrak{w}_H:=1$, and the clause asserts that the $z$-content
of a fluctuation-generated term carries $g_j^2$, the source being a shift of $x_j$ (the device of
\cite[Lemma~3]{B-CMP116}), and hence that $\bar{\mathfrak{a}}^E/\mathfrak{n}\le Cg_{k_0}^2$ (in the
normalisation of the grading). This is an additional factor gained in
the $z$-derivative, and would multiply the relative comparison
\eqref{8.5:eq-frozen-sourced-family-12} by a further $g_{k_0}^2$. Whether that bound also counts
the block-law factor written in (d) on both quadratic parts is immaterial here.
\item[--] Under the second grading, $\mathfrak{w}_W:=p_0(g_{k_0})^4\varrho_J^{-2}$, with $p_0(\cdot)$
the degree function of the coupling. Here
$w$-analyticity holds only on the region of (a), the field fluctuations are of size
$g_{k_0}p_0(g_{k_0})$, and the radius is the radius $\varrho_J$ of the correlation theorem. The
product $r_w^{-1}C_H$ is then read as
$\mathfrak{w}_W$, with its own one-sided floor
\[
C_W\ge C_0\,p_0(g_{K-s})^4\varrho_J^{-2}.
\]
Which of the two supports of the sup-ball applies is a clause of the sourced-format lemma of the
correlation theorem and is not decided in this definition. This is the only place at which a
dependence on $g_{K-s}$ can enter $C_0$, and it is one-sided.
\end{itemize}

For the R-class, likewise, $\bar{\mathfrak{w}}^R:=r_w^{-1}$ on the unweighted choice. Under the
second grading $\bar{\mathfrak{w}}^R:=1$. Under the first it is the weight assigned by the
R-analogue of the coupling-grading clause, which puts a further $g_j^2$ on the R-class terms, is
optional and is not proved here.

\item[(i)] \emph{Cell aggregates.} We write $X\in E$ and $X\in R$ for
$X\in\operatorname{Irr}_1^E$ and $X\in\operatorname{Irr}_1^R$. For a cell $c$ put
\begin{equation}\label{8.5:eq-frozen-sourced-family-13}
q^E(c)[f]:=\sum_{X\in E:\ c(X)=c}q_X[f],\qquad
\mathfrak{a}^E_c:=\sum_{X\in E:\ c(X)=c}\hat{\mathfrak{m}}_X,\qquad
\bar{\mathfrak{a}}^E:=\sup_c\mathfrak{a}^E_c .
\end{equation}
Define $q^R(c)[f]$, $\mathfrak{a}^R_c$ and $\bar{\mathfrak{a}}^R$ in the same way with $R$ in
place of $E$. The \emph{mass-weighted radius} is
\begin{equation}\label{8.5:eq-frozen-sourced-family-14}
\rho_\sharp:=\max_{\bullet\in\{E,R\}}\ \sup_{c:\ \mathfrak{a}^\bullet_c>0}\
\frac{1}{\mathfrak{a}^\bullet_c}\sum_{X\in\bullet:\ c(X)=c}r(X)\,\hat{\mathfrak{m}}_X .
\end{equation}
It is taken class by class, E-sums over the E-class and R-sums over the R-class, as the
second-derivative room lemma and the remainder-column bound use it. Cells with
$\mathfrak{a}^\bullet_c=0$ are excluded from the supremum. The weight cancels in the ratio, so
$\rho_\sharp$ depends on $\hat{\mathfrak{m}}^0$ only. It is a mean radius and is used for
first-order defect terms only. A hard dependence radius $R_{\mathrm{dep}}$ carries a tail
$e^{-\hat\kappa R_{\mathrm{dep}}/(2M)}$, treated in the subsection on localisation.

Thus $\mathfrak{a}_c$, $\bar{\mathfrak{a}}$ and $\rho_\sharp$ are majorant masses and radii, and
every use of them is as an upper bound. The second-derivative room lemma, the defect bounds and the
tail count are of degree one in the masses. The radius $\rho_\sharp$ and the tail ratio are of
degree zero. Hence no result changes when $\bar{\mathfrak{w}}$ is read through the coupling
grading or through $r_w^{-1}$.

\item[(j)] \emph{Gaussian letters and the linearised maps.} The letters $\bar{Q}$, $\bar{M}$, $v$,
$b$, $\psi$, $\mathcal{C}:=\mathcal{C}^{\mathrm{bg}}$ and $C_\nabla$ are those of the block law.
The quantity $\mathsf{S}_q$ of the block law, for which the block-law subsection gives an explicit
scale $\sigma_{\mathsf{d}}$ and an upper comparison bound, is taken there over the ball
$B(x_1,3\mathsf{d}/2)$ rather than over the ball $B(x_1,2\mathsf{d})$ taken by the covariance
decomposition of the block law.
Put $M_{k_0}:=DU_{k_0}(\mathbf{1})$. By the statements of the block-law subsection, $U_{k_0}$ is
$C^2$ at the flat configuration and
\[
d\circ M_{k_0}=\bar{M}_{k_0}\circ d .
\]
$K(p,\mathfrak{P})$ denotes the whole-lattice kernel of $\bar{M}_{k_0}$. It has three properties:
exponential decay; a clause comparing it on the torus with the whole lattice; and a constant
row-sum identity.
\end{itemize}
\end{definition}

\begin{proof}[Proof of the domination \eqref{8.5:eq-frozen-sourced-family-8} from the hypotheses
listed before it]
We derive \eqref{8.5:eq-frozen-sourced-family-8} from the five hypotheses listed before it in (g).
We first treat the unmarked term, then the marking in $w$, then the remaining assertions.

\emph{The unmarked term.} Let $X$ index an E-term of level $j$ and let $T=V^{(j)}_X$. Let $u$ be a
regular exact flux on $\hat{X}$, that is, one of the form $\varphi|_{\hat{X}}$ with $\varphi$ as in
\eqref{8.5:eq-frozen-sourced-family-9}--\eqref{8.5:eq-frozen-sourced-family-10}, bounded by one per
bare plaquette. Its regularity is supplied by
\eqref{8.5:eq-frozen-sourced-family-9}--\eqref{8.5:eq-frozen-sourced-family-10}, as assumed.

Choose an axial gauge on $\hat{X}$. In it, $u$ is realised as $u=dB$ by a bond field $B$ with
\[
|B|\lesssim r(X)/\eta .
\]
The reason is that an axial path in $\hat{X}$ from the anchor has at most of order
$r(X)\ell/\varepsilon=r(X)/\eta$ bare bonds. For complex $\sigma$,
\begin{equation}\label{8.5:eq-frozen-sourced-family-15}
|\partial(e^{i\sigma B})-1|\le|\sigma|\,|dB|+C|\sigma|^2|B|^2 .
\end{equation}
Hence the configuration $e^{i\sigma B}$ stays inside the fine-field domain of (g) for
\[
|\sigma|<\eta^2\min\{\alpha_0,\,c_\sigma/r(X)\},
\]
with a constant $c_\sigma>0$. Only on this domain is the fine-field form of Ba{\l}aban's bound
used. We denote the radius in $\sigma$ by $\sigma_X:=\eta^2\min\{\alpha_0,c_\sigma/r(X)\}$.

The function $\sigma\mapsto T(e^{i\sigma B})$ is analytic on $|\sigma|<\sigma_X$. By the
fine-field form of \cite[(0.29)]{B-CMP109} and by \eqref{8.5:eq-frozen-sourced-family-6}, it is
bounded there by $L^{-(4+\alpha)(k_0-j)}e^{-\kappa d_j(X)}$. Cauchy's inequality for the second
derivative at $\sigma=0$ gives
\[
|\operatorname{Hess}T[B,B]|\le 2\,\sigma_X^{-2}\,L^{-(4+\alpha)(k_0-j)}e^{-\kappa d_j(X)} .
\]
The mixed form $\operatorname{Hess}T[B,B']$ is recovered from the diagonal one by polarisation,
applied to $B\pm B'$:
\[
\operatorname{Hess}T[B,B']=\tfrac14\bigl(\operatorname{Hess}T[B+B',B+B']
-\operatorname{Hess}T[B-B',B-B']\bigr).
\]
Since $B\pm B'$ have flux at most two per bare plaquette, the radius in $\sigma$ for them is
$\sigma_X/2$; the resulting factor $4$ is absorbed in $C_H$.

Next we bound $\sigma_X^{-2}$:
\[
\sigma_X^{-2}=\eta^{-4}\max\{\alpha_0^{-2},\,r(X)^2c_\sigma^{-2}\}
\le\eta^{-4}\,(\alpha_0^{-2}+c_\sigma^{-2})\,(1+r(X))^2 .
\]
Here $\eta^{-4}=(\ell/\varepsilon)^4=\mathfrak{n}/6$ by \eqref{8.5:eq-frozen-sourced-family-4}.
This factor $\eta^{-4}$ is the factor $\mathfrak{n}$ of \eqref{8.5:eq-frozen-sourced-family-5}. It
comes from the two $\sigma$-derivatives, and the factor $\alpha_0^{-2}$ goes into $C_H$; nothing
else of $\eta$ enters.

The remaining factor $(1+r(X))^2$ is converted into a power of $d_j(X)$, using the lower bound
$d_j(X)\ge (r(X)-r_{\mathrm{read}})/M-1$ of the tail count assumed before
\eqref{8.5:eq-frozen-sourced-family-8}. From it,
\[
(1+r(X))^2\le C\bigl(M(d_j(X)+2)\bigr)^2 .
\]
For $\kappa\ge1$ we have $\sup_{n\ge0}(n+2)^2e^{-\kappa n/2}\le C$. Therefore
\begin{align*}
(1+r(X))^2e^{-\kappa d_j(X)}
&\le CM^2\,(d_j(X)+2)^2e^{-\kappa d_j(X)/2}\,e^{-\kappa d_j(X)/2}\\
&\le CM^2\,e^{-\kappa d_j(X)/2} .
\end{align*}
This is the origin of the halved rate $\hat\kappa=\kappa/2$ and of the dependence of $C_H$ on $M$.
Since $M$ is fixed before $C_0$, no uniformity in $K$, $s$ or $m$ is affected.

Finally, $L^{-(4+\alpha)(k_0-j)}\le L^{-(4+\alpha/2)(k_0-j)}$. Collecting the factors into
$C_H=C_H(\alpha_0,M,\kappa)$, the flux form of the unmarked term is bounded on regular unit
directions by $\hat{\mathfrak{m}}^0_X$ of \eqref{8.5:eq-frozen-sourced-family-5}.

For an R-term the two steps are the same, with the bound of \cite[(2.30)--(2.31)]{B-CMP119}
on the fine-field domain, assumed before \eqref{8.5:eq-frozen-sourced-family-8}, in place of
\cite[(0.29)]{B-CMP109}. They give the factor $C_H\mathfrak{n}$ and the rate
$\hat\kappa_R$ of \eqref{8.5:eq-frozen-sourced-family-7}.

\emph{The marking in $w$.} By the hypothesis of (h) on the localised form of the sourced-format
lemma, the terms
$V^{(j),w}_X$ are analytic in $w$ on the sup-ball of radius $r_w$, they read $w$ on $X^+$ only, and
they obey the bounds used above there. Hence for $w=zf$ the function
$z\mapsto\operatorname{Hess}V^{(j),zf}_X[B,B]$ is analytic for $|z|\sup_{X^+}|f|<r_w$. On that disc
it obeys the bound established for the unmarked term. Cauchy's inequality for the first derivative
at $z=0$ gives
\[
|H_X[f][B,B]|\le r_w^{-1}\sup_{X^+}|f|\cdot\hat{\mathfrak{m}}^0_X
\]
on regular unit directions. Passing to constant-flux forms gives
\[
|q_X[f]|\le\mathfrak{m}^{\mathrm{reg}}_X[f]\le r_w^{-1}\hat{\mathfrak{m}}^0_X\sup_{X^+}|f| .
\]
With $\bar{\mathfrak{w}}=r_w^{-1}$ as in \eqref{8.5:eq-frozen-sourced-family-11}, this is
\eqref{8.5:eq-frozen-sourced-family-8}.

\emph{The weight cancels in $\rho_\sharp$.} Within each class the weight, $\bar{\mathfrak{w}}$ or
$\bar{\mathfrak{w}}^R$, is a level-free constant. It therefore factors out of numerator and
denominator in \eqref{8.5:eq-frozen-sourced-family-14}, and $\rho_\sharp$ is the same when every
$\hat{\mathfrak{m}}_X$ is replaced by $\hat{\mathfrak{m}}^0_X$.

\emph{Degrees.} A bound of degree one in the masses scales by the same constant as the weight. A
quantity of degree zero does not change. Consequently, reading $\bar{\mathfrak{w}}$ through the
coupling grading of (h) instead of through $r_w^{-1}$ leaves the form of every bound stated in
terms of these majorants unchanged; that \eqref{8.5:eq-frozen-sourced-family-8} holds under either
grading of (h) is the coupling-grading hypothesis, which is not proved here.
\end{proof}

\begin{lemma}[Constant-flux forms of gauge-invariant localised terms]
\label{8.5:lem-constant-flux-forms}
Let $\hat{X}$ be a box of the bare lattice which does not wrap around the torus $\mathbb{T}_m$, let
$X\subset\hat{X}$, and let $T$ be a function of the restriction $U|_X$ of the bond variables to $X$,
analytic on an open neighbourhood of the identity configuration $\mathbf{1}$. Fix an
$\operatorname{Ad}$-invariant inner product $(\cdot,\cdot)_{\mathfrak{g}}$ on the Lie algebra
$\mathfrak{g}$ of the gauge group, and measure $\mathfrak{g}$-valued cochains in the norm it induces.
Let $\operatorname{Hess}T$ be the fine-field Hessian of
Definition~\ref{8.5:def-localised-terms}, and for a $\mathfrak{g}$-valued bond field $B$ put
\[
DT(\mathbf{1})[B]:=\partial_\tau T\bigl(\exp(i\tau B)\bigr)\big|_{\tau=0}.
\]
Assume that $T$ is gauge invariant: $T(U^u)=T(U)$ for every gauge transformation $u$, where
$U^u$ is the gauge transform of $U$ by $u$ (Lemma~\ref{1.1:lem-invariance}). Then the following
hold.
\begin{itemize}
\item[(a)] $\operatorname{Hess}T[d\lambda,B']=0$ for every $\mathfrak{g}$-valued function
$\lambda$ on sites and every $\mathfrak{g}$-valued bond field $B'$, and $DT(\mathbf{1})=0$.
\item[(b)] $\operatorname{Hess}T[B,B']$ depends on $B,B'$ only through
$(dB,dB')|_{\hat{X}}$. Consequently
\[
Q_T[\varphi,\varphi']:=\operatorname{Hess}T[B,B'],\qquad dB=\varphi,\ dB'=\varphi'\ \text{on }\hat{X},
\]
is a well-defined symmetric bilinear form on pairs of $\mathfrak{g}$-valued $2$-cochains
exact on $\hat{X}$.
\item[(c)] Let $\mathfrak{g}$ be compact and simple, and let $(T^a)$ be a basis of $\mathfrak{g}$
orthonormal for $(\cdot,\cdot)_{\mathfrak{g}}$. There is a symmetric bilinear form
$Q^{\mathbb{R}}_T$ on real bond fields, depending only on $(dB,dB')|_{\hat{X}}$, such that
\[
\operatorname{Hess}T[T^a\otimes B,\,T^b\otimes B']=\delta^{ab}\,Q^{\mathbb{R}}_T[B,B']
\]
for all real bond fields $B,B'$ and all $a,b$. For semisimple $\mathfrak{g}$ the same holds on
each simple ideal separately.
\item[(d)] For $c\in\mathbb{R}^{\mathfrak{O}}$ and a point $y$ let $B^{c,y}$ be an affine bond
field of constant flux $c$ about $y$: a real $1$-cochain whose value on each bond is a linear
function of the displacement of the bond from $y$, and whose coboundary is
$dB^{c,y}(p)=c_{\pi(p)}$ on every plaquette $p\subset\hat{X}$. Then
\[
\begin{aligned}
q_T[c',c'']&:=Q^{\mathbb{R}}_T\bigl[B^{c',y},B^{c'',y}\bigr]\\
&\phantom{:}=\operatorname{Hess}T\bigl[T^a\otimes B^{c',y},\,T^a\otimes B^{c'',y}\bigr]
\quad\text{for any }a,
\end{aligned}
\]
the equality by (c), is independent of $y$; it defines the \emph{constant-flux form}
$q_T\in\operatorname{Sym}^2(\mathbb{R}^{\mathfrak{O}})$. For every lattice motion $\mathsf{g}$,
\[
q_{T\circ\mathsf{g}^{-1}}=\mathsf{g}\cdot q_T .
\]
Here $\mathsf{g}$ transports bond fields, $(T\circ\mathsf{g}^{-1})(U):=T(\mathsf{g}^{-1}\cdot U)$,
$\mathsf{g}\cdot c\in\mathbb{R}^{\mathfrak{O}}$ is the constant flux
$\mathsf{g}\cdot(dB^{c,y})$, and
$(\mathsf{g}\cdot q)[c',c'']:=q[\mathsf{g}^{-1}\cdot c',\mathsf{g}^{-1}\cdot c'']$.
\item[(e)] The quantity
\[
\mathfrak{m}_T:=\sup\bigl\{|Q_T[\varphi,\varphi']|:\ \varphi,\varphi'\ \text{exact on }\hat{X},\
|\varphi|_\infty,|\varphi'|_\infty\le1\ \text{on }\hat{X}\bigr\}
\]
is finite, and $|q_T|\le\mathfrak{m}_T$, where
$|q_T|:=\sup\{|q_T[c',c'']|:|c'|_\infty,|c''|_\infty\le1\}$.
For $\varphi$ exact on $\hat{X}$ put
\begin{gather*}
{}[\varphi]_{\beta_0;\hat{X}}:=\sup\frac{|\varphi(p')-\varphi(p)|}
{\bigl(|x(p')-x(p)|/\ell\bigr)^{\beta_0}},\\
\mathfrak{j}_{\hat{X}}(\varphi):=\eta^{-1}\max_{b\subset\hat{X}}|d^*\varphi(b)|,\qquad
\mathfrak{reg}_X(\varphi):=\max\bigl\{[\varphi]_{\beta_0;\hat{X}},\,
\mathfrak{j}_{\hat{X}}(\varphi)\bigr\},
\end{gather*}
where the supremum in the first line is over pairs $p,p'\subset\hat{X}$ that are translates of
one another along a coordinate axis with $0<|x(p')-x(p)|\le\ell$,
$\beta_0\in(0,1)$ is the H\"older exponent of \cite[(3.32), (4.18)]{B-CMP109} and $d^*$
is the plain bare co-differential, the signed sum over the six plaquettes containing the
bond $b$. With $r(X)$ the radius of $\hat{X}$ in cells
(Definition~\ref{8.5:def-localised-terms}), assume in addition the admissibility condition:
there is $c_{\mathrm{adm}}=c_{\mathrm{adm}}(\alpha_0,\alpha_1,\beta_0)>0$ such that, for every
box $\hat{X}$ not wrapping the torus, every $\varphi$ exact on $\hat{X}$ and every single-colour
direction $T^a$, with $B_\varphi$ a bond primitive of $\varphi$ on $\hat{X}$, the second
line of the following display holds, the first and third lines being definitions:
\begin{equation}\label{8.5:eq-constant-flux-forms-a}
\begin{gathered}
\|\varphi\|_X:=\max\bigl\{|\varphi|_\infty,\ \mathfrak{reg}_X(\varphi)/r(X)\bigr\},\\
\|\varphi\|_X\le1,\ \ |\sigma|\le c_{\mathrm{adm}}\,\eta^2/r(X)\ \Longrightarrow\
\exp(i\sigma T^aB_\varphi)\in\mathfrak{U}^{\mathrm{f}}_j(\hat{X})\ \text{ for every } j\le k_0,\\
\mathfrak{m}^{\mathrm{reg}}_T:=\sup\bigl\{|Q_T[\varphi,\varphi']|:\
\|\varphi\|_X,\|\varphi'\|_X\le1\bigr\}.
\end{gathered}
\end{equation}
Here $\mathfrak{U}^{\mathrm{f}}_j(\hat{X})$ is the complex fine-field domain of the fine-field
reading of \cite[(0.29)]{B-CMP109} for the level-$j$ terms supported on $\hat{X}$, a reading
built from \cite[(1.11)--(1.16), (3.32), (4.4)]{B-CMP109}, and $\alpha_0,\alpha_1$ are the radii of
Ba{\l}aban's complex configuration spaces $\mathfrak{U}^c_j$. For the configurations in the
second line the single-colour bond variables commute, so the plaquette holonomy is
$\exp(i\sigma T^a\varphi(p))$ exactly and the covariant co-differential is the plain $d^*$;
since $r(X)\ge1$,
\[
\begin{gathered}
|\sigma|\,|\varphi(p)|\le c_{\mathrm{adm}}\eta^2,\\
|\sigma|\,|\varphi(p')-\varphi(p)|\le c_{\mathrm{adm}}\eta^2\bigl(|x(p')-x(p)|/\ell\bigr)^{\beta_0},\\
|\sigma|\,|d^*\varphi(b)|\le c_{\mathrm{adm}}\eta^3,
\end{gathered}
\]
for all plaquettes $p\subset\hat{X}$, bonds $b\subset\hat{X}$ and pairs $p,p'$ as in the
definition of $[\varphi]_{\beta_0;\hat{X}}$: the curvature is at most $c_{\mathrm{adm}}\eta^2$,
the H\"older quotient has the shape of \cite[(4.18)]{B-CMP109}, and the current has the shape
of \cite[(4.17)]{B-CMP109}. Independently of this condition,
\[\mathfrak{m}^{\mathrm{reg}}_T\le\mathfrak{m}_T\qquad\text{and}\qquad
|q_T|\le\mathfrak{m}^{\mathrm{reg}}_T,
\]
constant fluxes having $\mathfrak{reg}_X=0$. Every direction paired with these forms in what
follows is measured in $\|\cdot\|_X$.
\end{itemize}
\end{lemma}

Gauge invariance of each term is the hypothesis under which the lemma is applied. For
Ba{\l}aban's effective terms it is stated in \cite[(1.19)]{B-CMP109}, and
\cite[(4.9)]{B-CMP109} is stated for every gauge-invariant analytic function. For the terms of
the frozen sourced family of Definition~\ref{8.5:def-frozen-sourced-family} it is part of the
format in which that family is given.

\begin{proof}
We write $\Phi(\mathcal{Y}):=T(\exp(i\mathcal{Y}))$ for $\mathfrak{g}$-valued bond fields
$\mathcal{Y}$ near $0$. By the definitions,
$\operatorname{Hess}T[B,B']=D^2\Phi(0)[B,B']$ and $DT(\mathbf{1})[C]=D\Phi(0)[C]$.
Since $\Phi$ is analytic near $0$, the form $D^2\Phi(0)$ is symmetric and bilinear.

(a) Fix $\lambda$ and $B'$. For real $\sigma,\tau$ near $0$ put $u:=\exp(i\sigma\lambda)$ and
$U:=\exp(i\tau B')$. By gauge invariance $T(U^u)=T(\exp(i\tau B'))$, so the function
$(\sigma,\tau)\mapsto T(U^u)$ does not depend on $\sigma$. By the expansion
\cite[(4.8)]{B-CMP109} of the gauge transform, whose correction terms are commutators,
$U^u=\exp(i\mathcal{Y}(\sigma,\tau))$ with
\[
\mathcal{Y}(\sigma,\tau)=\tau B'-\sigma\,d\lambda+\sigma\tau\,\Lambda(\lambda,B')+\Xi(\sigma,\tau).
\]
Here $\Lambda(\lambda,B')$ is the bilinear commutator term of that expansion, and every term of
$\Xi$ is divisible by $\sigma^2$ or by $\tau^2$. Hence $\partial_\sigma\mathcal{Y}(0,0)=-d\lambda$,
$\partial_\tau\mathcal{Y}(0,0)=B'$ and
$\partial_\sigma\partial_\tau\mathcal{Y}(0,0)=\Lambda(\lambda,B')$, and $\Xi$
contributes nothing to these derivatives. The chain rule gives
\[
\partial_\sigma\partial_\tau\Phi(\mathcal{Y}(\sigma,\tau))\big|_{0}
=D^2\Phi(0)\bigl[\partial_\sigma\mathcal{Y}(0),\partial_\tau\mathcal{Y}(0)\bigr]
+D\Phi(0)\bigl[\partial_\sigma\partial_\tau\mathcal{Y}(0)\bigr].
\]
Since $\Phi(\mathcal{Y}(\sigma,\tau))=T(U^u)$ is independent of $\sigma$, its mixed derivative
vanishes, and therefore
\begin{equation}\label{8.5:eq-constant-flux-forms-1}
0=-\operatorname{Hess}T[d\lambda,B']+DT(\mathbf{1})\bigl[\Lambda(\lambda,B')\bigr].
\end{equation}

We show that $DT(\mathbf{1})=0$. Take $\lambda\equiv\xi$ constant, with $\xi\in\mathfrak{g}$.
Then $u$ is constant, and $U^u=uUu^{-1}$ on every bond. Since
$u\exp(i\tau B')u^{-1}=\exp(i\tau\operatorname{Ad}_uB')$, gauge invariance gives
$T(\exp(i\tau\operatorname{Ad}_uB'))=T(\exp(i\tau B'))$. Differentiating in $\tau$ at $0$
gives
\[
DT(\mathbf{1})\bigl[\operatorname{Ad}_{\exp(i\sigma\xi)}B'\bigr]=DT(\mathbf{1})[B']
\qquad\text{for all }\sigma.
\]
Differentiating in $\sigma$ at $0$ then gives
$DT(\mathbf{1})[\operatorname{ad}_\xi B']=0$ for all $B'$. Here
$\operatorname{ad}_\xi B'$ is the derivative at $\sigma=0$ of
$\operatorname{Ad}_{\exp(i\sigma\xi)}B'$, a non-zero multiple of the commutator of $\xi$
with $B'(b)$ on each bond $b$.

Represent the linear functional $DT(\mathbf{1})$ as
$DT(\mathbf{1})[C]=\sum_b(\Delta(b),C(b))_{\mathfrak{g}}$, with $\Delta(b)$ in the
complexification of $\mathfrak{g}$ and $(\cdot,\cdot)_{\mathfrak{g}}$ extended bilinearly. By
$\operatorname{Ad}$-invariance,
\[
\bigl(\Delta(b),[\xi,C(b)]\bigr)_{\mathfrak{g}}=-\bigl([\xi,\Delta(b)],C(b)\bigr)_{\mathfrak{g}}.
\]
Since $B'$ is arbitrary, $[\xi,\Delta(b)]=0$ for all $\xi\in\mathfrak{g}$ and all $b$. Thus each
$\Delta(b)$ lies in the centre, which is trivial because $\mathfrak{g}$ has trivial centre
\cite[(4.14)]{B-CMP109}. Hence $DT(\mathbf{1})=0$, and
\eqref{8.5:eq-constant-flux-forms-1} gives $\operatorname{Hess}T[d\lambda,B']=0$.

(b) Let $dB=d\tilde{B}$ on $\hat{X}$. Then $B-\tilde{B}$ has vanishing coboundary on the box
$\hat{X}$. This box does not wrap the torus and is therefore contractible, so
$B-\tilde{B}=d\lambda$ on $\hat{X}$ for a $\mathfrak{g}$-valued function $\lambda$ on the sites
of $\hat{X}$, which we extend arbitrarily to all sites. Since $T$ reads only $U|_X$ and
$X\subset\hat{X}$, the value $\operatorname{Hess}T[C,B']$ depends only on $C|_{\hat{X}}$.
Hence, by bilinearity and (a),
\[
\operatorname{Hess}T[B,B']-\operatorname{Hess}T[\tilde{B},B']
=\operatorname{Hess}T[d\lambda,B']=0 .
\]
By symmetry of the Hessian the same holds in the second argument. Every $2$-cochain exact on
$\hat{X}$ has a primitive on $\hat{X}$, which extends to a bond field on the whole lattice.
Therefore $Q_T$ is well defined, and it is symmetric and bilinear because $\operatorname{Hess}T$
is.

(c) Fix real bond fields $B,B'$ and put
$\varpi_{B,B'}(\xi,\xi'):=\operatorname{Hess}T[\xi\otimes B,\xi'\otimes B']$ for
$\xi,\xi'\in\mathfrak{g}$.
For a constant gauge transformation $u$ we have
$u\exp(i\mathcal{Y})u^{-1}=\exp(i\operatorname{Ad}_u\mathcal{Y})$,
so gauge invariance gives $\Phi(\operatorname{Ad}_u\mathcal{Y})=\Phi(\mathcal{Y})$.
Differentiating twice yields
$\varpi_{B,B'}(\operatorname{Ad}_u\xi,\operatorname{Ad}_u\xi')=\varpi_{B,B'}(\xi,\xi')$. Thus the
bilinear form $(a,b)\mapsto\operatorname{Hess}T[T^a\otimes B,T^b\otimes B']$ is
$\operatorname{Ad}$-invariant.

Write $\varpi_{B,B'}(\xi,\xi')=(\xi,\Psi_{B,B'}\xi')_{\mathfrak{g}}$ with $\Psi_{B,B'}$ a
complex-linear map of the complexification of $\mathfrak{g}$ into itself, and
$(\cdot,\cdot)_{\mathfrak{g}}$ extended bilinearly. Since
$(\cdot,\cdot)_{\mathfrak{g}}$ is $\operatorname{Ad}$-invariant, $\Psi_{B,B'}$ commutes with the
adjoint action. For compact simple $\mathfrak{g}$ whose complexification is simple, the linear
maps of the complexification commuting with the adjoint action are the multiples of the
identity; this is the case for $\mathfrak{su}(N)$ (compare \cite{B-CMP109}, where the
corresponding form is stated to be proportional to $\delta_{ab}$). Hence $\Psi_{B,B'}$ is a
multiple of the identity; we let $Q^{\mathbb{R}}_T[B,B']$ be that multiple, and orthonormality of
$(T^a)$ gives the identity of (c). This form is symmetric and bilinear in $(B,B')$
because $\operatorname{Hess}T$ is. By (b), and because $d(T^a\otimes B)=T^a\otimes dB$, it
depends only on $(dB,dB')|_{\hat{X}}$. For semisimple $\mathfrak{g}$ the argument applies to each
simple ideal.

(d) For two points $y,y'$ the difference $B^{c,y}-B^{c,y'}$ is a constant $1$-cochain, equal
on $\hat{X}$ to $d\lambda$ for a linear function $\lambda$. Hence $dB^{c,y}=dB^{c,y'}$ on
$\hat{X}$, and by (c), which
rests on (b) and hence on (a), $q_T[c',c'']$ is independent of $y$.

For the covariance, transporting fields commutes with the exponential, so
\begin{align*}
\operatorname{Hess}(T\circ\mathsf{g}^{-1})[B,B']
&=\operatorname{Hess}T[\mathsf{g}^{-1}\cdot B,\mathsf{g}^{-1}\cdot B'],\\
Q^{\mathbb{R}}_{T\circ\mathsf{g}^{-1}}[B,B']
&=Q^{\mathbb{R}}_T[\mathsf{g}^{-1}\cdot B,\mathsf{g}^{-1}\cdot B'].
\end{align*}
The coboundary $d$ commutes with $\mathsf{g}$. So $\mathsf{g}^{-1}\cdot B^{c,y}$ has constant
flux $\mathsf{g}^{-1}\cdot c$ and is linear in the displacement from $\mathsf{g}^{-1}y$. It
therefore differs from $B^{\mathsf{g}^{-1}\cdot c,\mathsf{g}^{-1}y}$ by a field with vanishing
coboundary on the box, which is exact there as in (b):
$\mathsf{g}^{-1}\cdot B^{c,y}=B^{\mathsf{g}^{-1}\cdot c,\mathsf{g}^{-1}y}+d\lambda$.
By (b) and (c) the $d\lambda$ does not contribute, and by the independence of the base point
\[
q_{T\circ\mathsf{g}^{-1}}[c',c'']
=q_T\bigl[\mathsf{g}^{-1}\cdot c',\mathsf{g}^{-1}\cdot c''\bigr]
=(\mathsf{g}\cdot q_T)[c',c''].
\]

(e) The $\mathfrak{g}$-valued $2$-cochains exact on $\hat{X}$ form a finite-dimensional space.
The bilinear form $Q_T$ is continuous on it, and the set over which $\mathfrak{m}_T$ is a
supremum is bounded, so $\mathfrak{m}_T<\infty$.

Let $|c'|_\infty,|c''|_\infty\le1$ and fix $a$. By (c) and (b),
\[
q_T[c',c'']=\operatorname{Hess}T\bigl[T^a\otimes B^{c',y},T^a\otimes B^{c'',y}\bigr]
=Q_T\bigl[T^a\otimes dB^{c',y},\,T^a\otimes dB^{c'',y}\bigr].
\]
The cochain $\varphi:=T^a\otimes dB^{c,y}$ is exact on $\hat{X}$, with primitive
$T^a\otimes B^{c,y}$. Its value on a plaquette $p$ is $T^ac_{\pi(p)}$, so
$|\varphi|_\infty\le|c|_\infty\le1$, since $|T^a|=1$. Hence $|q_T[c',c'']|\le\mathfrak{m}_T$.

Since $|\varphi|_\infty\le\|\varphi\|_X$, the set over which $\mathfrak{m}^{\mathrm{reg}}_T$ is
a supremum is contained in the set over which $\mathfrak{m}_T$ is one, so
$\mathfrak{m}^{\mathrm{reg}}_T\le\mathfrak{m}_T$.

For the constant flux $\varphi=T^a\otimes dB^{c,y}$, a translate $p'$ of $p$ along a coordinate
axis lies in the same plane as $p$, so $\varphi(p')=\varphi(p)$ and
$[\varphi]_{\beta_0;\hat{X}}=0$. The six plaquettes containing a bond $b$ come in pairs lying in
a common plane, the two members of a pair entering $d^*$ with opposite signs and carrying the
same value of $\varphi$; hence $d^*\varphi(b)=0$ and $\mathfrak{j}_{\hat{X}}(\varphi)=0$. Thus
$\mathfrak{reg}_X(\varphi)=0$, so $\|\varphi\|_X=|\varphi|_\infty\le1$, and the displayed
identity gives $|q_T[c',c'']|\le\mathfrak{m}^{\mathrm{reg}}_T$, that is,
$|q_T|\le\mathfrak{m}^{\mathrm{reg}}_T$.

It remains to verify the consequences of the second line of
\eqref{8.5:eq-constant-flux-forms-a} listed after it. Let $\|\varphi\|_X\le1$ and
$|\sigma|\le c_{\mathrm{adm}}\eta^2/r(X)$. The bond variables $\exp(i\sigma T^aB_\varphi(b))$ all
lie in the one-parameter subgroup generated by $T^a$, so they commute; the ordered product round
a plaquette $p\subset\hat{X}$ is therefore $\exp(i\sigma T^a\,dB_\varphi(p))$, which equals
$\exp(i\sigma T^a\varphi(p))$, and their adjoint action fixes $T^a$, so the covariant
co-differential of $T^a\otimes\varphi$ reduces to $T^a\otimes d^*\varphi$. By the definition of
$\|\varphi\|_X$ we have $|\varphi|_\infty\le1$ on $\hat{X}$ and
$\mathfrak{reg}_X(\varphi)\le r(X)$, hence $[\varphi]_{\beta_0;\hat{X}}\le r(X)$ and
$\mathfrak{j}_{\hat{X}}(\varphi)\le r(X)$. Using also $r(X)\ge1$, for plaquettes and bonds of
$\hat{X}$ and pairs $p,p'$ as in the definition of $[\varphi]_{\beta_0;\hat{X}}$,
\begin{align*}
|\sigma|\,|\varphi(p)|&\le c_{\mathrm{adm}}\eta^2/r(X)\le c_{\mathrm{adm}}\eta^2,\\
|\sigma|\,|\varphi(p')-\varphi(p)|
&\le|\sigma|\,[\varphi]_{\beta_0;\hat{X}}\bigl(|x(p')-x(p)|/\ell\bigr)^{\beta_0}\\
&\le c_{\mathrm{adm}}\eta^2\bigl(|x(p')-x(p)|/\ell\bigr)^{\beta_0},\\
|\sigma|\,|d^*\varphi(b)|&\le|\sigma|\,\eta\,\mathfrak{j}_{\hat{X}}(\varphi)
\le c_{\mathrm{adm}}\eta^3.
\end{align*}
\end{proof}

All the forms of the lemma are forms in the fine field, supported in $\hat{X}$. On the torus a
constant flux $c$ means the restriction to the box $\hat{X}$ of the curl of the affine field
$B^{c,y}$, which is always exact there. The background map does not enter the lemma; in the
estimates that follow it enters only through the directions of the linearised minimiser
$M_{k_0}$.

\begin{lemma}[Hyperoctahedral-invariant forms on two-forms are scalar]\label{8.5:lem-invariant-forms-scalar}
Let $\mathfrak{O}$, $\mathbb{R}^{\mathfrak{O}}$, its basis $e_P$, the pairing
$\langle\cdot,\cdot\rangle$, $\operatorname{Sym}^2(\mathbb{R}^{\mathfrak{O}})$ and the action of $W_4$
be as in Notation~\ref{8.5:not-planes-forms}.
\begin{enumerate}
\item[(i)] Every $W_4$-invariant $q\in\operatorname{Sym}^2(\mathbb{R}^{\mathfrak{O}})$ is a multiple of
$\langle\cdot,\cdot\rangle$.
\item[(ii)] Every $W_4$-equivariant linear map $\Lambda\colon\mathbb{R}^{\mathfrak{O}}\to\mathbb{R}^{\mathfrak{O}}$
is a scalar multiple of the identity.
\end{enumerate}
\end{lemma}

\begin{proof}
For $1\le\mu\le4$ let $r_\mu\in W_4$ be the sign change of the axis $\mu$. Under
$\mathbb{R}^{\mathfrak{O}}\cong\Lambda^2\mathbb{R}^4$ it sends $e_P\mapsto-e_P$ if $\mu\in P$ and
$e_P\mapsto e_P$ if $\mu\notin P$, since exactly one factor of the wedge product changes sign when
$\mu\in P$. Put $\sigma_P:=-1$ if $\mu\in P$ and $\sigma_P:=1$ otherwise; thus $r_\mu$ is diagonal in the
basis $(e_P)$, and $r_\mu^{-1}=r_\mu$.

(i) Write $q_{PQ}:=q[e_P,e_Q]$. Invariance under $r_\mu$ gives
\begin{equation*}
q_{PQ}=(r_\mu\cdot q)[e_P,e_Q]=q[r_\mu e_P,r_\mu e_Q]=\sigma_P\sigma_Q\,q_{PQ}.
\end{equation*}
If $\mu$ lies in exactly one of $P$, $Q$, then $\sigma_P\sigma_Q=-1$, hence $q_{PQ}=0$. Let
$P\neq Q$. Both are two-element subsets of $\{1,2,3,4\}$, so $|P\cap Q|\le1$. If $|P\cap Q|=1$, take
$\mu\in P\setminus Q$, which is non-empty. If $P\cap Q=\emptyset$, take any $\mu\in P$. In either case
$\mu$ lies in exactly one of $P$, $Q$, so $q_{PQ}=0$ for all $P\neq Q$.

The reflections $r_\mu$ cannot be dispensed with here. Let $P^c:=\{1,2,3,4\}\setminus P$, and let
$\operatorname{sgn}(P)\in\{\pm1\}$ be the sign of the permutation that lists the elements of $P$ and
then those of $P^c$, each in increasing order. Consider the Hodge pairing
\begin{equation*}
h[c',c'']:=\sum_{P\in\mathfrak{O}}\operatorname{sgn}(P)\,c'_P\,c''_{P^c}.
\end{equation*}
It satisfies $c'\wedge c''=h[c',c'']\,e_1\wedge e_2\wedge e_3\wedge e_4$, since $e_P\wedge e_Q=0$ unless
$Q=P^c$. Moreover $h$ is symmetric, because $\operatorname{sgn}(P^c)=\operatorname{sgn}(P)$. It is not
a multiple of $\langle\cdot,\cdot\rangle$. An orthogonal map $g$ multiplies $c'\wedge c''$ by
$\det g$, so $h$ is invariant under the elements of $W_4$ of determinant one. A single sign change
maps $h$ to $-h$: in the display above, $\mu$ lies in exactly one of $P$ and $P^c$, so
$\sigma_P\sigma_{P^c}=-1$. Hence statement~(i) holds for $W_4$ but fails for its subgroup
$W_4\cap\mathrm{SO}(4)$. The lattice symmetries considered in \cite{B-CMP119} and \cite{B-CMP109}
include the coordinate reflections.

Now let $P,Q\in\mathfrak{O}$. There is a permutation $\pi$ of the axes with $\pi(P)=Q$, and its induced
action satisfies $\pi e_P=\theta\,e_Q$ for some $\theta\in\{\pm1\}$. Invariance under $\pi^{-1}$ gives
\begin{equation*}
q_{QQ}=q[\pi e_P,\pi e_P]=\theta^2\,q_{PP}=q_{PP}.
\end{equation*}
Thus the permutations act transitively on $\mathfrak{O}$ and $q_{PP}$ is a constant $\lambda$. Since the
off-diagonal entries vanish, $q[c',c'']=\lambda\sum_P c'_Pc''_P=\lambda\langle c',c''\rangle$.

(ii) The same argument applies entrywise. Put $\Lambda_{PQ}:=\langle e_P,\Lambda e_Q\rangle$.
Equivariance and $r_\mu^2=\mathbf{1}$ give $\Lambda=r_\mu\Lambda r_\mu$. Since $r_\mu$ is diagonal with
entries $\pm1$, it is self-adjoint for $\langle\cdot,\cdot\rangle$, and therefore
\begin{equation*}
\Lambda_{PQ}=\langle r_\mu e_P,\Lambda r_\mu e_Q\rangle=\sigma_P\sigma_Q\,\Lambda_{PQ}.
\end{equation*}
The case analysis of~(i) then gives $\Lambda_{PQ}=0$ for $P\neq Q$, so $\Lambda e_P=\Lambda_{PP}e_P$.
Now take $\pi$ with $\pi(P)=Q$ and $\pi e_P=\theta e_Q$. Equivariance $\Lambda\pi=\pi\Lambda$, applied
to $e_P$, yields $\theta\Lambda_{QQ}e_Q=\theta\Lambda_{PP}e_Q$. Hence $\Lambda_{QQ}=\Lambda_{PP}$, and
$\Lambda$ is a scalar.

The proof has used the sign changes $r_\mu$ and the permutations of the axes separately, as
generators of $W_4$, in accordance with the generator-wise action fixed in
Notation~\ref{8.5:not-planes-forms}.
\end{proof}

\begin{definition}[Partition-lattice and hyperoctahedral covariance of localised terms. Stated; proof
incomplete at the step marked $(\ast)$]\label{8.5:def-term-covariance}
We use the cells, the cube partitions $\pi_j$, the localisation domains $\mathcal{D}_j$ and the partition
lattice $\Gamma=M\mathbb{Z}^4$ (in cells) of Notation~\ref{8.5:not-partition-lattice}. We also use the
hyperoctahedral group $W_4$ with its action of Notation~\ref{8.5:not-planes-forms}, and the frozen sourced
level-$k_0$ family with its terms $V^{(j),w}_X$ and the source shift $\zeta_{k_0}(p)[w]$ of the inverse
coupling $x_{k_0}$ in its main term, as in Definition~\ref{8.5:def-frozen-sourced-family}. For $j\le k_0$ we
write $\{E^{(j)}(X,\cdot)\}_{X\in\mathcal{D}_j}$ for the family of localised terms of the $j$-th term of
Ba{\l}aban's expansion \cite[(1.3)]{B-CMP109} in its localised representation \cite[(1.7)]{B-CMP109};
we suppress the arguments of $E^{(j)}(X,\cdot)$ other than the configuration $U$.
We write $\{V^{(j)}_X\}_{X\in\mathcal{D}_j}$ for the residual family of \cite[(0.28)]{B-CMP109}, which
remains after the coefficient $\beta_j$ has been extracted.
For $a\in\Gamma$ we write $X+a$ and $p+a$ for the translates of a domain $X$ and of a plaquette $p$ by
$a$ cells, that is, by the physical vector $a\ell$, and $x-a$, $b-a$ for the translates of a point $x$
and of a bond $b$ by the vector $-a\ell$; we put
\begin{equation*}
  (t_a w)(x):=w(x-a),\qquad (t_aU)(b):=U(b-a)
\end{equation*}
for complex source profiles $w$, points $x$, configurations $U$ and bonds $b$. For $r\in W_4$ we write $rX$
and $rp$ for the images of $X$ and $p$ under the action of Notation~\ref{8.5:not-planes-forms}, we put
$(r\cdot w)(x):=w(r^{-1}x)$, and we write $r\cdot U$ for the transformed configuration, the action of
reflections on bond variables being the one written out in \cite[(3.58)--(3.59)]{B-CMP119}.

\begin{itemize}
\item[(a)] \emph{Covariance under the partition lattice.} For every
$a\in\Gamma$, every $j\le k_0$, every $X\in\mathcal{D}_j$, every profile $w$, every configuration $U$ and every
plaquette $p$,
\begin{align}
  \mathcal{D}_j+a&=\mathcal{D}_j,\label{8.5:eq-term-covariance-1}\\
  V^{(j),t_aw}_{X+a}(t_aU)&=V^{(j),w}_X(U),\label{8.5:eq-term-covariance-2}\\
  \zeta_{k_0}(p+a)[t_aw]&=\zeta_{k_0}(p)[w],\label{8.5:eq-term-covariance-3}
\end{align}
and the family $\{E^{(j)}(X,\cdot)\}_{X\in\mathcal{D}_j}$ satisfies
\begin{equation}\label{8.5:eq-term-covariance-4}
  E^{(j)}(X+a,t_aU)=E^{(j)}(X,U).
\end{equation}
Moreover, let $\mathcal{N}$, a power of $L$, denote the side of an auxiliary torus of the correlation
theorem. Then for
every $X$ whose filled hull $\hat{X}$ does not wrap around that torus, the localised term of the
whole-lattice family indexed by $X$ coincides with the corresponding term of the torus family up to an
error $O(e^{-c\mathcal{N}/M})$ per term, with a constant $c>0$, uniformly summable over the terms. This
error arises from the torus propagators and from the
difference between the coefficient $\beta^{(\mathcal{N})}_j$ extracted on the torus and the coefficient
$\beta_j$ defined in infinite volume by \cite[(1.22)]{B-CMP109}.
\item[(b)] \emph{Hyperoctahedral covariance.} For every $r\in W_4$ (with the action fixed above),
every $j\le k_0$, every $X\in\mathcal{D}_j$, every profile $w$,
every configuration $U$ and every plaquette $p$,
\begin{align}
  r\mathcal{D}_j&=\mathcal{D}_j,\label{8.5:eq-term-covariance-5}\\
  V^{(j),r\cdot w}_{rX}(r\cdot U)&=V^{(j),w}_X(U),\label{8.5:eq-term-covariance-6}\\
  V^{(j)}_{rX}(r\cdot U)&=V^{(j)}_X(U),\label{8.5:eq-term-covariance-7}\\
  \zeta_{k_0}(rp)[r\cdot w]&=\zeta_{k_0}(p)[w].\label{8.5:eq-term-covariance-8}
\end{align}
Identities \eqref{8.5:eq-term-covariance-6} and \eqref{8.5:eq-term-covariance-8} concern the sourced
terms and the source shift of the frozen sourced family of
Definition~\ref{8.5:def-frozen-sourced-family}; identity \eqref{8.5:eq-term-covariance-7} concerns the
residual family $\{V^{(j)}_X\}_X$; and the family
$\{E^{(j)}(X,\cdot)\}_{X\in\mathcal{D}_j}$ satisfies
\begin{equation}\label{8.5:eq-term-covariance-9}
  E^{(j)}(rX,r\cdot U)=E^{(j)}(X,U).
\end{equation}
\end{itemize}
\end{definition}

\begin{proof}
\emph{Step 1: the summed functional in \cite{B-CMP109}.} \cite[\S1]{B-CMP109} states that the explicitly defined
expressions in the $j$-th term of \cite[(1.3)]{B-CMP109} are invariant under the Euclidean transformations
of the lattice $\mathbb{T}^{(j+1)}$ (Definition~\ref{1.5:def-block-average}), and it assumes that the same
holds for all expressions in this term. The terms $V^{(j)}_X$ of the residual family belong to the
expressions of the $j$-th term to which that assumption applies. The invariance is stated for the $j$-th
term, that is, for the sum over $X$ of the
localised terms. The unit translations of $\mathbb{T}^{(j+1)}$ do not preserve the partition $\pi_j$, so
the statement concerns the summed functional and not the individual terms indexed by $X$.
\cite[(5.2)--(5.7)]{B-CMP109} states the invariance of the second variation of the whole $j$-th term
$\sum_{X\in\mathcal{D}_j}E^{(j)}(X,\cdot)$ of \cite[(1.3)]{B-CMP109} under all Euclidean symmetries of the
lattice, reflections included; this too concerns the summed functional.

\emph{Step 2: the point-attributed family in \cite{B-CMP119}.} \cite[(3.58)]{B-CMP119} states Euclidean
covariance, with reflections included, following from \cite[(2.29)]{B-CMP119}, term by term for the family
of localised terms attributed to the points $z$ of the lattice.
The translation invariance that \cite[\S3]{B-CMP119} states for the
whole-lattice function $E^{(j)}(U_j,z)$ is likewise a property of a function attributed to the points $z$.
Both are indexed by points, not by domains $X\in\mathcal{D}_j$.

\emph{Step 3 $(\ast)$: the termwise identities for the families indexed by $X$.} The scalar reduction of the
second-order terms is applied to the family $\{E^{(j)}(X,\cdot)\}_X$ and to the residual family
$\{V^{(j)}_X\}_X$ (respectively $\{V^{(j),w}_X\}_X$ for the frozen sourced family) separately, and the
reflections contained in $W_4$ are essential for it; this is why the hyperoctahedral
identities are required for both families. Every $a\in\Gamma$ permutes the cubes of $\pi_j$ for
$j\le k_0$ (Notation~\ref{8.5:not-partition-lattice}), and the domains of $\mathcal{D}_j$ are unions of
such cubes; identity \eqref{8.5:eq-term-covariance-1} states that this permutation maps the class
$\mathcal{D}_j$ onto itself, and it is obtained together with \eqref{8.5:eq-term-covariance-4} at the
present step. In the same way every generator of $W_4$ permutes the cubes of every $\pi_j$, and identity
\eqref{8.5:eq-term-covariance-5}, which states that the class $\mathcal{D}_j$ is mapped onto itself,
belongs to the same step. The termwise identities at issue are
\eqref{8.5:eq-term-covariance-4} and \eqref{8.5:eq-term-covariance-9} for $\{E^{(j)}(X,\cdot)\}_X$, and
\eqref{8.5:eq-term-covariance-7} for the residual family $\{V^{(j)}_X\}_X$ of \cite[(0.28)]{B-CMP109},
term by term. These identities are what the indexing of the cluster expansion of
\cite[\S\S2--3]{B-CMP109} by the cubes of the $\Gamma$-invariant partitions $\pi_j$, which are permuted by
every $a\in\Gamma$ and by every generator of $W_4$, is taken to give term by term. The termwise
identities are not derived from that construction; the proof is incomplete at this step. The same holds
for the frozen sourced family of
Definition~\ref{8.5:def-frozen-sourced-family}, whose construction is part of the correlation theorem: the
identities \eqref{8.5:eq-term-covariance-2}, \eqref{8.5:eq-term-covariance-3},
\eqref{8.5:eq-term-covariance-6} and \eqref{8.5:eq-term-covariance-8}, which concern the sourced terms
$V^{(j),w}_X$ and the source shift $\zeta_{k_0}$ of that definition, are the frozen-family part of the same
incomplete step. In
that family every threshold, cube size, decomposition, background and reference law is frozen, and the
position occurs in no constant. The thresholds are numerical constants and are therefore unchanged under
the transformations of $W_4$ and of $\Gamma$. The source
\begin{equation*}
  \mathcal{O}(f)=\sum_p f(x(p))\bigl[1-\operatorname{Re}\operatorname{tr}U(\partial p)\bigr]
\end{equation*}
is covariant under $W_4$ when $f$ and $U$ are transformed together (Lemma~\ref{1.1:lem-invariance}).
These facts concern the constants and
the source. The termwise covariance clause of the inductive format of the frozen family is not proved
here; it is part of this incomplete step.

The torus clause of part (a) belongs to the same step. In the correlation theorem
the constants and families that do not depend on the volume are obtained on auxiliary tori whose side is
a power of $L$. The
infinite-volume limit defining $\beta_j$ \cite[(1.22)]{B-CMP109} exists by the localised representation
\cite[(1.7)]{B-CMP109}. The coincidence up to $O(e^{-c\mathcal{N}/M})$ per term, uniformly
summable over the terms, is not proved; it belongs to the step marked $(\ast)$. This form, weaker than
exact coincidence of the terms, suffices for the
comparison with the torus made later in this
section, because
\begin{equation*}
  \Bigl(\frac{\mathcal{N}}{M}\Bigr)^4e^{-c\mathcal{N}/M}\longrightarrow0
\end{equation*}
and
\begin{equation*}
  \bigl(\beta^{(\mathcal{N})}_j-\beta_j\bigr)\operatorname{Hess}A[dB,dB]=o(\mathcal{N}^4)\,\eta^4 .
\end{equation*}
Here $\operatorname{Hess}A$ is the Hessian of the Wilson action $A$ and $dB$ is the test increment of that
comparison.

\emph{Why $\Gamma$ and not $\mathbb{Z}^4$.} The lattice $\Gamma$ in part (a) cannot be replaced by the cell
lattice $\mathbb{Z}^4$. Let $a\in\mathbb{Z}^4\setminus M\mathbb{Z}^4$, and choose an axis $\mu$ with
$a_\mu\notin M\mathbb{Z}$. Let $X\in\mathcal{D}_{k_0}$ be non-empty and bounded. The cubes of $\pi_{k_0}$
have side $M$ cells, $M>1$, and are permuted by the translations of $\Gamma$
(Notation~\ref{8.5:not-partition-lattice}); hence in the direction $\mu$ they cut the cell coordinate into
consecutive blocks of $M$ integers, and the first elements of these blocks form one coset of
$M\mathbb{Z}$. Since $X$ is a union of cubes of $\pi_{k_0}$, the $\mu$-coordinates of its cells form a
union of such blocks, so the least of them, $n_\mu$, is the first element of a block. Because
$a_\mu\notin M\mathbb{Z}$, the integer $n_\mu+a_\mu$ is not the first element of a block.

Let $y$ be a cell of $X+a$ with $\mu$-coordinate $n_\mu+a_\mu$. The cube of $\pi_{k_0}$ containing $y$ has
lower $\mu$-coordinate strictly smaller than $n_\mu+a_\mu$, while $n_\mu+a_\mu$ is the least
$\mu$-coordinate among the cells of $X+a$. Thus that cube contains cells outside $X+a$ while meeting
$X+a$.

Therefore $X+a$ is not a union of cubes of $\pi_{k_0}$, and $\mathcal{D}_{k_0}+a\ne\mathcal{D}_{k_0}$. For such
$a$ and $j=k_0$, the left side of \eqref{8.5:eq-term-covariance-2} is indexed by
$X+a\notin\mathcal{D}_{k_0}$ and is therefore not a member of the family. Hence covariance of the
family indexed by $X$ under unit translations fails.

The translation invariance of the point-attributed function $E^{(j)}(U_j,z)$ of \cite[\S3]{B-CMP119},
recalled in Step~2, refers to a different indexing and is not in conflict with this.
\end{proof}

\begin{lemma}[Covariance makes block-summed constant-flux forms scalar]
\label{8.5:lem-block-sum-scalar}
Let $\mathsf{C}\cong\mathbb{Z}^4$ be the cell lattice and $\Gamma=M\mathbb{Z}^4\subset\mathsf{C}$ the partition
lattice. Let $\mathcal{V}=\{V_X\}$ be a family of terms each of which satisfies the hypotheses of
Lemma~\ref{8.5:lem-constant-flux-forms}; we index the family by its set of localisation domains and
write $X\in\mathcal{V}$ for $V_X\in\mathcal{V}$. For $V_X\in\mathcal{V}$ write $q_X\in
\operatorname{Sym}^2(\mathbb{R}^{\mathfrak{O}})$ for its constant-flux form, as in that lemma, and
$c(X)\in\mathsf{C}$ for its anchor cell. Assume:
\begin{itemize}
\item[(i)] the family $\mathcal{V}$ is covariant under $\Gamma$ in the sense of
Definition~\ref{8.5:def-term-covariance};
\item[(ii)] the family $\mathcal{V}$ is covariant under each generator of $W_4$ (the
sign changes of single axes and the permutations of the axes), in the sense of
Definition~\ref{8.5:def-term-covariance}.
\end{itemize}
Both hypotheses are instances of the covariance property of
Definition~\ref{8.5:def-term-covariance}, which is stated there with its proof incomplete at one
step; in the present lemma they are assumptions.
Define the cell aggregate and the $\Gamma$-class sum by
\begin{equation}\label{8.5:eq-block-sum-scalar-1}
q_{\mathcal{V}}(c):=\sum_{X\in\mathcal{V}:\,c(X)=c}q_X ,\qquad
\mathfrak{q}_M(\mathcal{V}):=\sum_{X\in\mathcal{V}/\Gamma}q_X ,
\end{equation}
where the second sum takes one representative in each $\Gamma$-class (it does not depend on the
choice of representatives, since $q_{X+a}=q_X$ for $a\in\Gamma$, as shown in the proof of (a)). Then:
\begin{itemize}
\item[(a)] $q_{\mathcal{V}}(c)$ is $\Gamma$-periodic in $c$;
\item[(b)] for a period block $\mathfrak{B}$ (a cube of $M^4$ cells of $\pi_{k_0}$) the sum
$\Phi_{\mathfrak{B}}:=\sum_{c\in\mathfrak{B}}q_{\mathcal{V}}(c)$ is independent of $\mathfrak{B}$,
is equal to $\mathfrak{q}_M(\mathcal{V})$, and is a multiple
$\delta(\mathcal{V})\langle\cdot,\cdot\rangle$ of the pairing on $\mathbb{R}^{\mathfrak{O}}$.
\end{itemize}
\end{lemma}

\begin{proof}
The properties of constant-flux forms used below are clauses~(a) and~(d) of
Lemma~\ref{8.5:lem-constant-flux-forms}.

(a) Fix $a\in\Gamma$ and a cell $c$. By the $\Gamma$-covariance hypothesis (i) the translate $X+a$ of
the localisation domain of a term of $\mathcal{V}$ is again the localisation domain of a term of
$\mathcal{V}$; and the anchor cell is translation-equivariant, $c(X+a)=c(X)+a$. Hence $X\mapsto X+a$
maps $\{X\in\mathcal{V}: c(X)=c\}$ bijectively onto $\{X\in\mathcal{V}: c(X)=c+a\}$, the inverse
being $X\mapsto X-a$. It remains to show $q_{X+a}=q_X$. By hypothesis (i), $V_{X+a}(t_aU)=V_X(U)$
(for terms carrying a source, the source is translated together with the configuration, as in
Definition~\ref{8.5:def-term-covariance}); the
identity configuration $\mathbf{1}$ is translation invariant, so differentiating twice at
$\mathbf{1}$ shows that the Hessian of the translated term on translated test fields is the
Hessian of the original term on the original ones:
\begin{equation*}
\operatorname{Hess}V_{X+a}[t_aB,t_aB']=\operatorname{Hess}V_X[B,B'] .
\end{equation*}
Apply this to the affine test fields $B^{c',y}$ of constant flux $c'$, on which the constant-flux
form is read. By clause~(d) of Lemma~\ref{8.5:lem-constant-flux-forms},
\begin{equation*}
t_aB^{c',y}=B^{c',y+a}=B^{c',y}+d\lambda
\end{equation*}
for a linearised gauge direction $d\lambda$. By clause~(a) of the same lemma the Hessian of a
gauge-invariant term vanishes whenever one argument is of the form $d\lambda$, so the gauge
directions drop out of $\operatorname{Hess}V_{X+a}[t_aB^{c',y},t_aB^{c'',y}]$. Therefore the form
of $V_{X+a}$ on the pair $(B^{c',y},B^{c'',y})$ coincides with the form of $V_X$ on the same pair,
that is, $q_{X+a}=q_X$. Summing over the bijection above gives $q_{\mathcal{V}}(c+a)=q_{\mathcal{V}}(c)$.

(b) A period block $\mathfrak{B}$ contains exactly one cell of each class of $\mathsf{C}/\Gamma$, and
any two period blocks are therefore complete systems of representatives of $\mathsf{C}/\Gamma$. By
(a), $q_{\mathcal{V}}$ is a function on $\mathsf{C}/\Gamma$, so $\Phi_{\mathfrak{B}}$ does not depend
on $\mathfrak{B}$.

Next, each $\Gamma$-class $\{X+a: a\in\Gamma\}$ of $\mathcal{V}$ has exactly one member anchored in
$\mathfrak{B}$: the anchors $c(X+a)=c(X)+a$ run through the $\Gamma$-class of $c(X)$, which meets
$\mathfrak{B}$ in exactly one cell. By the identity $q_{X+a}=q_X$ established in the proof of (a),
\begin{equation*}
\Phi_{\mathfrak{B}}=\sum_{c\in\mathfrak{B}}\ \sum_{X\in\mathcal{V}:\,c(X)=c}q_X
=\sum_{X\in\mathcal{V}:\,c(X)\in\mathfrak{B}}q_X=\mathfrak{q}_M(\mathcal{V}) .
\end{equation*}

Let $r$ be one of the generators of $W_4$ in hypothesis (ii). By the $W_4$-covariance hypothesis
(ii), $X\mapsto rX$
permutes $\mathcal{V}$. Since $r$ is a signed permutation of the axes, $r\Gamma=\Gamma$, and
$r(X+a)=rX+ra$ with $ra\in\Gamma$; so $r$ maps $\Gamma$-classes of $\mathcal{V}$ to
$\Gamma$-classes, and $[X]\mapsto[rX]$ is a bijection of $\mathcal{V}/\Gamma$. Let $r\cdot q$ denote
the action of $r$ on $\operatorname{Sym}^2(\mathbb{R}^{\mathfrak{O}})$ induced by its action on the
planes $\mathfrak{O}$, as in Lemma~\ref{8.5:lem-invariant-forms-scalar}. Reindexing the class
sum along this bijection, and then using the covariance of constant-flux forms under signed
permutations, $q_{rX}=r\cdot q_X$, which is clause~(d) of Lemma~\ref{8.5:lem-constant-flux-forms}
together with hypothesis (ii), we obtain
\begin{equation}\label{8.5:eq-block-sum-scalar-2}
\mathfrak{q}_M(\mathcal{V})=\sum_{X\in\mathcal{V}/\Gamma}q_{rX}
=\sum_{X\in\mathcal{V}/\Gamma}r\cdot q_X=r\cdot\mathfrak{q}_M(\mathcal{V}) .
\end{equation}
Thus $\mathfrak{q}_M(\mathcal{V})$ is invariant under every sign change of a single axis and every
permutation of the axes, hence under $W_4$. Identity \eqref{8.5:eq-block-sum-scalar-2} gives
invariance under the sign changes and under the permutations separately, which is what part~(i) of
the lemma on hyperoctahedral-invariant forms (Lemma~\ref{8.5:lem-invariant-forms-scalar}) requires.
Hence $\mathfrak{q}_M(\mathcal{V})=\delta(\mathcal{V})\langle\cdot,\cdot\rangle$ for some
$\delta(\mathcal{V})\in\mathbb{R}$. Since $\Phi_{\mathfrak{B}}=
\mathfrak{q}_M(\mathcal{V})$, this proves (b).
\end{proof}

\begin{remark}[Per-cell annihilation of constant flux is unavailable]\label{8.5:rem-per-cell-annihilation}
Let $\mathcal{V}$ be a family of terms as in Lemma~\ref{8.5:lem-block-sum-scalar}, with cell
aggregates $q_{\mathcal{V}}(c)=\sum_{X\in\mathcal{V}:\,c(X)=c}q_X$ and with
$\mathfrak{q}_M(\mathcal{V})$ the block sum of that lemma.
\begin{itemize}
\item[(i)] Suppose that $\mathcal{V}$ were covariant under all translations of the cell lattice
$\mathsf{C}$ by unit cells, and not only under the partition lattice $\Gamma=M\mathbb{Z}^4$. Then
part (a) of Lemma~\ref{8.5:lem-block-sum-scalar} would hold for every $a\in\mathbb{Z}^4$, so
$q_{\mathcal{V}}(c)$ would be independent of $c$; summing over the $M^4$ cells of a period block
and using part (b) of that lemma would give
\begin{equation}\label{8.5:eq-per-cell-annihilation-1}
q_{\mathcal{V}}(c)=M^{-4}\,\mathfrak{q}_M(\mathcal{V})\qquad\text{for every }c\in\mathsf{C}.
\end{equation}
The mean identity for constant flux would then give
$q_{\mathcal{V}}(c)[\mathbf{1}]=0$ for every cell $c$. This is \emph{per-cell annihilation of
constant flux}.
\item[(ii)] The family obtained from the partition into cubes of $M^4$ cells is covariant under
$\Gamma$, but not under all unit-cell translations. The argument in (i) therefore does not
apply to it. Per-cell annihilation of constant flux is not proved in this book. It is not used
anywhere in what follows.
A re-indexing of the family averaged over the offsets of the partition would restore
covariance under all unit-cell translations, because the total functional, the sum of all the
terms of the family, is covariant under unit-cell translations. That route is not used.
\end{itemize}
\end{remark}

\begin{definition}[Identity of representations for the shifted frozen run; proof outstanding]
\label{8.5:def-representation-identity}
Let $\varrho_J>0$ be the constant of the correlation theorem for which the running shift of that
theorem,
\begin{equation*}
\zeta_{k+1}:=\zeta_k+b_{k+1}(\zeta)-b_{k+1}(0),\qquad \zeta_0=\zeta,
\end{equation*}
is analytic in $\zeta$ on the disc $D(0,2\varrho_J)$ and real for real $\zeta$. Here $b_{k+1}(\zeta)$ is
the $\beta$-number, in the sense fixed below, of the level-$(k+1)$ terms of the whole-lattice family
at bare weight $x_0-\zeta$
with every threshold, cube size, decomposition, background and reference law frozen at those of the
$\zeta=0$ run. Let $r_w>0$ be the constant of the localised form of the sourced-format statement of
the correlation theorem whose inverse $r_w^{-1}=\bar{\mathfrak{w}}$ is the level-free route weight. We put
\begin{equation*}
\zeta_*:=\min\{2\varrho_J,\,r_w\}.
\end{equation*}

\emph{Analyticity in the constant shift.} The correlation theorem contains an analyticity clause
(stated; proof incomplete). It reads as follows. For complex $\zeta\in D(0,\zeta_*)$,
take the frozen sourced whole-lattice family at constant profile $w\equiv\zeta$. At level $k_0$ this
family has the representation of Definition~\ref{8.5:def-frozen-sourced-family}, with index sets and
classes of terms that do not depend on $\zeta$. Every term $V^{\zeta}_X$ of this representation (the
term $V^{(j),w}_X$ of Definition~\ref{8.5:def-frozen-sourced-family} at $w\equiv\zeta$) is gauge
invariant, is $X$-local and is analytic in $\zeta$. The amount moved into the weight of the main term is
\begin{equation*}
\sum_{k<k_0}b_{k+1}(\zeta).
\end{equation*}
This is the clause. It is used below in the following way. When the terms $V^{\zeta}_X$ of extracted
type in this representation (the residual family of extracted type at constant shift $\zeta$) are
differentiated termwise in $\zeta$, the weighted majorants
$\hat{\mathfrak{m}}_X=r_w^{-1}\hat{\mathfrak{m}}^0_X$ serve as Cauchy majorants. Per anchor cell $c$
they give
\begin{equation}\label{8.5:eq-representation-identity-1}
\sum_{c(X)=c}\ \sup_{|\zeta|\le\zeta_*/2}\bigl|\partial_\zeta q_{V^{\zeta}_X}[\mathbf{1}]\bigr|
\le\frac{2r_w}{\zeta_*}\,\mathfrak{a}^E_c
\le 2\max\Bigl\{1,\frac{r_w}{2\varrho_J}\Bigr\}\,\mathfrak{a}^E_c ,
\end{equation}
where $q_{V^{\zeta}_X}[\mathbf{1}]$ is the constant-flux form of the single term $V^{\zeta}_X$ at the
constant profile $\mathbf{1}$, and $\mathfrak{a}^E_c$ is the cell-aggregate majorant mass of the
residual family of extracted type at the anchor cell $c$; termwise differentiation is then justified by the
Weierstrass majorant test. The second inequality
in \eqref{8.5:eq-representation-identity-1} follows from the definition of $\zeta_*$. If $\zeta_*=r_w$,
then $2r_w/\zeta_*=2$. If $\zeta_*=2\varrho_J$, then $2r_w/\zeta_*=2\cdot r_w/(2\varrho_J)$.

\emph{Families and their $\beta$-numbers.} Consider a family of localised terms of the expansion.
The extraction by Ward--Takahashi identities that leads to \cite[(0.28)]{B-CMP109} groups together,
out of the family, those terms whose sum is $\beta$ times an expansion of the Wilson action, for some
coefficient $\beta$.
This coefficient is the family's \emph{$\beta$-number}; for the families extracted at step $k+1$ it is
their contribution to the coupling increment $\beta_{k+1}$. It is the number
defined in \cite[(3.61)]{B-CMP109}. For a family
$\mathcal{F}$ we write $\beta[\mathcal{F}]$ for its $\beta$-number. We call the families from which the
coupling term is extracted the \emph{families of extracted type}. The running-shift statement of the
correlation theorem (stated; proof incomplete) decomposes the number $b_{k+1}=b_{k+1}(\zeta)$ of the
frozen family as
\begin{equation}\label{8.5:eq-representation-identity-2}
b_{k+1}=\beta^0_{k+1}+\beta\bigl[\mathsf{P}^{(k+1)}-\mathcal{L}^{(k+1)}\bigr]
+\beta\bigl[\mathcal{Q}^{(k+1)}\bigr],
\end{equation}
where the difference $\mathsf{P}^{(k+1)}-\mathcal{L}^{(k+1)}$ is the family of level-$(k+1)$ terms
carried by the regions described in the small-piece statement of this chapter,
$\mathsf{P}^{(k+1)}$ and $\mathcal{L}^{(k+1)}$ being the two level-$(k+1)$ families of the
running-shift statement of which it is the difference; and $\mathcal{Q}^{(k+1)}$ is the
quadratic-difference family.
In \eqref{8.5:eq-representation-identity-2} the three terms are read as follows.
\begin{itemize}
\item $\beta^0_{k+1}$ is the $\beta$-number of the frozen family of extracted type
$\{V^{(k+1),0}_X\}_X$ carrying no source mark (the superscript $0$). These are the $\zeta=0$ run's own
irrelevant terms born at level $k+1$.
\item The second term is the $\beta$-number of the difference family
$\mathsf{P}^{(k+1)}-\mathcal{L}^{(k+1)}$ described above.
\item The third term is the $\beta$-number of the quadratic-difference family $\mathcal{Q}^{(k+1)}$.
\end{itemize}
This reading of $\beta^0_{k+1}$ names the family from which it comes. It therefore places the residual
of the extraction among the families of extracted type.

\emph{The identity.} For every real $\zeta$ with $|\zeta|<\zeta_*$ and every torus of the construction
whose side, counted in cells, is a multiple of $M$, let the \emph{frozen run} at bare weight
$x_0-\zeta$ be the run with every threshold, cube size, decomposition, background and reference law
frozen at those of the $\zeta=0$ run, and let $A^{\zeta}_{k_0}$ be its effective action at level $k_0$.
The \emph{identity of representations for the shifted frozen run} is the following assertion.

This run carries an identity of analytic functions of the coarse variable $V$ of the type
\cite[(0.28)]{B-CMP109}, as Ba{\l}aban's own run does. It holds on the level-$k_0$ small-field space
$\{V:|\partial U_{k_0}(V)-\mathbf{1}|<\varepsilon_0\eta^2\}$ of $T_\eta$. Here $T_\eta$ is the bare
lattice, of spacing $\eta$ in the units of level $k_0$; $\partial U_{k_0}(V)$ stands for the plaquette
variables $U_{k_0}(V)(\partial p)$ of the minimiser; and $\varepsilon_0$ is Ba{\l}aban's small-field
constant, the threshold of the small-field renormalisation transformations
\cite[(0.17)--(0.20)]{B-CMP109}. The identity reads
\begin{equation}\label{8.5:eq-representation-identity-a}
\begin{split}
A^{\zeta}_{k_0}\bigl(U_{k_0}(V)\bigr)
&=\bigl(\text{the localised representation of the $\zeta$-run at frozen}\\
&\qquad\text{thresholds, before the extraction}\bigr)\\
&=-\bigl(x_{k_0}-\zeta_{k_0}\bigr)A\bigl(U_{k_0}(V)\bigr)
+\sum_{j\le k_0}\sum_{X}V^{\zeta}_X\bigl(U_{k_0}(V)\bigr).
\end{split}
\end{equation}
In \eqref{8.5:eq-representation-identity-a}, $\zeta_{k_0}$ is the running shift at level $k_0$ defined
above, and the weight $x_{k_0}-\zeta_{k_0}$ is the weight $x_{k_0}-\zeta_{k_0}(p)[w]$ of
Definition~\ref{8.5:def-frozen-sourced-family} at constant profile $w\equiv\zeta$. For each
$j\le k_0$, $X$ ranges over the index set of level $j$ of
Definition~\ref{8.5:def-frozen-sourced-family}, and $V^{\zeta}_X$ is the term of that level with index
$X$, that is, $V^{(j),w}_X$ of that definition at $w\equiv\zeta$. Moreover, the following is asserted
together with \eqref{8.5:eq-representation-identity-a}.
\begin{itemize}
\item For each $k<k_0$, the families of extracted type of level $k+1$ of the $\zeta$-run are exactly
the families whose $\beta$-number is $b_{k+1}(\zeta)$, read as in
\eqref{8.5:eq-representation-identity-2}.
\item $U_{k_0}(\cdot)$ is the same minimiser as for the $\zeta=0$ run, since the background is frozen.
\end{itemize}
Only the form of \eqref{8.5:eq-representation-identity-a} is used in what follows.
\end{definition}

Proof outstanding.

\medskip
\noindent\textit{Route.} For Ba{\l}aban's own run, \cite[Theorem~1]{B-CMP109} gives the two
representations of the level-$k$ effective action, whose identity is the one expressed by
\cite[(0.28)]{B-CMP109}. This theorem and \cite[Theorem~3]{B-CMP109} concern the actions defined
inductively by the small-field renormalisation transformations \cite[(0.17)--(0.20)]{B-CMP109}, whose
thresholds and couplings are those of the $\zeta=0$ run.

The definition of the frozen whole-lattice family and of $b_{k+1}(\zeta)$ recalled above keeps those
thresholds and couplings unchanged in the run started at $x_0-\zeta$. It defines numbers and
whole-lattice families only, and asserts nothing about the level-$k_0$ density of that run on the torus.

The analyticity clause recalled above gives three things: analyticity in $\zeta$, index sets that do
not depend on $\zeta$, and the amount $\sum_{k<k_0}b_{k+1}(\zeta)$. It does not give that the
level-$k_0$ density of the $\zeta$-run admits both representations in
\eqref{8.5:eq-representation-identity-a} as the same analytic function of $V$. That is the property
used in the proof of the two-point witness theorem, where a quantity on the torus is differentiated
in $\zeta$.

A proof would establish \eqref{8.5:eq-representation-identity-a} inside the inductive format of the
correlation theorem at constant profile $w\equiv\zeta$. At level $k$ that format is the inductive
hypothesis of the correlation theorem, carrying Ba{\l}aban's bounds under that theorem's hypotheses.
The proof would combine this with
the running-shift statement that $b_{k+1}(\zeta)$ is the $\beta$-number of the families of extracted
type of the $\zeta$-run. Both are clauses of the correlation theorem (stated; proof incomplete).

Throughout, cochains on a lattice take values in a fixed finite-dimensional real inner-product space, with pointwise norm $|\cdot|$.

Let $\varepsilon$ be the bare spacing and $\varepsilon_j$ the spacing of the level-$j$ lattice, and put $\eta_j=\varepsilon/\varepsilon_j$, the bare spacing in level-$j$ units. We work on the infinite lattices $\varepsilon\mathbb{Z}^4$ (bare) and $\varepsilon_j\mathbb{Z}^4$ (level $j$), and write $e_\mu$ for the unit vector of direction $\mu\in\{1,2,3,4\}$.

On $\varepsilon\mathbb{Z}^4$ a $1$-cochain $a$ assigns a value $a(x,\mu)$ to the bond from $x$ to $x+\varepsilon e_\mu$. Its coboundary on the plaquette $(x;\mu,\nu)$, $\mu<\nu$, is
\begin{equation*}
da(x;\mu,\nu)=a(x,\mu)+a(x+\varepsilon e_\mu,\nu)-a(x+\varepsilon e_\nu,\mu)-a(x,\nu),
\end{equation*}
and the same conventions are used on $\varepsilon_j\mathbb{Z}^4$. A cochain is \emph{constant} if its value on a bond (plaquette) depends only on the direction $\mu$ (the plane $\mu<\nu$).

Let $\mathcal{T}_j$ be the group of translations by vectors of $\varepsilon_j\mathbb{Z}^4$, acting on cochains of either lattice by $(ta)(x,\mu)=a(x-y,\mu)$ for the translation $t$ by $y$. A $1$-cochain $a$ is \emph{affine} if $ta-a$ is a constant $1$-cochain for every $t\in\mathcal{T}_j$.

We write $\langle\cdot,\cdot\rangle$ for the sum over plaquettes of the pointwise inner product of two $2$-cochains, and the same symbol for the analogous sum over bonds of two $1$-cochains, whenever the sum is finite, and $\|\cdot\|$ for the corresponding norm.

\begin{lemma}[Affine solutions of the constrained first-variation equations]
\label{8.5:lem-affine-first-variation}
Let $Q_j$ be the linearisation at the identity configuration of the $j$-fold block-averaging map of Definition~\ref{1.5:def-block-average}, acting from bare $1$-cochains to level-$j$ $1$-cochains. Assume the closedness property of the linearised block average in the following form:
\begin{itemize}
\item[(a)] $Q_j$ is linear and maps closed bare $1$-cochains onto closed level-$j$ $1$-cochains;
\item[(b)] $Q_j$ is translation covariant, that is, $Q_j(ta)=t(Q_ja)$ for every $t\in\mathcal{T}_j$.
\end{itemize}
Then $Q_j$ maps constant $1$-cochains to constant ones and affine $1$-cochains to affine ones.

Moreover, for every constant level-$j$ $2$-cochain $u$ there is an affine bare $1$-cochain $a_0$ with the following properties: $Q_ja_0$ is affine; $dQ_ja_0=u$; and $da_0=:c$ is constant with $c=\eta_j^2u$.

Finally, $a_0$ satisfies the first- and second-variation equations of the minimisation of $\tfrac12\|da\|^2$ under the constraint $Q_ja=Q_ja_0$, in the following sense: for every finitely supported bare $1$-cochain $\delta$ with $Q_j\delta=0$ the sum
\begin{equation*}
\sum_p\bigl(|d(a_0+\delta)(p)|^2-|da_0(p)|^2\bigr)
\end{equation*}
has finitely many non-zero terms, its linear part in $\delta$ is $2\langle da_0,d\delta\rangle$, its quadratic part is $\|d\delta\|^2$, and
\begin{equation}\label{8.5:eq-affine-first-variation-1}
\langle da_0,d\delta\rangle=0,\qquad \|d\delta\|^2\ge0 .
\end{equation}
\end{lemma}

\begin{proof}
\emph{Constants.} Let $\xi$ be a constant bare $1$-cochain. Then $t\xi=\xi$ for every $t\in\mathcal{T}_j$, so by (b)
\begin{equation*}
t(Q_j\xi)=Q_j(t\xi)=Q_j\xi .
\end{equation*}
The group $\mathcal{T}_j$ acts transitively on the level-$j$ bonds of each direction. Hence $Q_j\xi$ is constant.

\emph{Affine cochains.} Let $a$ be affine and $t\in\mathcal{T}_j$, and write $ta=a+\xi_t$ with $\xi_t$ constant. By (b) and the linearity in (a),
\begin{equation*}
t(Q_ja)=Q_j(ta)=Q_ja+Q_j\xi_t ,
\end{equation*}
and $Q_j\xi_t$ is constant by the first step. So $Q_ja$ is affine.

\emph{Construction of $a_0$.} Let $u$ be given and put $c_{\mu\nu}=\eta_j^2u_{\mu\nu}$ for $\mu<\nu$. Extend by $c_{\nu\mu}=-c_{\mu\nu}$ and $c_{\mu\mu}=0$, and define
\begin{equation*}
a_0(x,\mu)=\frac{1}{2\varepsilon}\sum_{\rho=1}^4 c_{\rho\mu}\,x_\rho .
\end{equation*}
For every $y\in\varepsilon\mathbb{Z}^4$, $a_0(x+y,\mu)-a_0(x,\mu)$ does not depend on $x$, so $a_0$ is affine.

Moreover,
\begin{align*}
a_0(x+\varepsilon e_\mu,\nu)-a_0(x,\nu)&=\tfrac12c_{\mu\nu},\\
a_0(x+\varepsilon e_\nu,\mu)-a_0(x,\mu)&=-\tfrac12c_{\mu\nu}.
\end{align*}
Hence $da_0(x;\mu,\nu)=c_{\mu\nu}$, and $da_0=c$ is constant.

By the previous step, $Q_ja_0$ is affine. For $t\in\mathcal{T}_j$ write $t(Q_ja_0)=Q_ja_0+\xi$ with $\xi$ constant. Since the coboundary of a constant $1$-cochain vanishes and $d$ commutes with translations,
\begin{equation*}
t(dQ_ja_0)=d\,t(Q_ja_0)=dQ_ja_0+d\xi=dQ_ja_0 .
\end{equation*}
By transitivity of $\mathcal{T}_j$ on the level-$j$ plaquettes of each plane, $dQ_ja_0$ is constant.

Its value is read from the explicit form of the block-averaging map in Definition~\ref{1.5:def-block-average}, by counting plaquettes: a level-$j$ plaquette $\mathfrak{P}$ of the plane $(\mu,\nu)$ is the union of $\eta_j^{-2}$ bare plaquettes $p$ of that plane, and the level-$j$ circulation of $Q_ja_0$ around $\partial\mathfrak{P}$ is the sum of $da_0$ over them, so that
\begin{equation*}
dQ_ja_0(\mathfrak{P})=\sum_{p\subset\mathfrak{P}}da_0(p)=\eta_j^{-2}c_{\mu\nu}=u_{\mu\nu} .
\end{equation*}
Hence $dQ_ja_0=u$.

\emph{The variations.} Let $\delta$ be finitely supported with $Q_j\delta=0$. Then $a_0+\delta$ satisfies the constraint, and $d\delta$ is finitely supported. Since $da_0=c$, the following sum has finitely many non-zero terms, and expanding the square gives
\begin{equation}\label{8.5:eq-affine-first-variation-2}
\sum_p\bigl(|d(a_0+\delta)(p)|^2-|da_0(p)|^2\bigr)=2\langle c,d\delta\rangle+\|d\delta\|^2 .
\end{equation}
The first variation is $\langle da_0,d\delta\rangle=\langle c,d\delta\rangle$.

Summation by parts, which is legitimate because $\delta$ is finitely supported, gives $\langle c,d\delta\rangle=\langle d^*c,\delta\rangle$. Here $(d^*c)(b)$ is the sum, over the plaquettes $p$ containing the bond $b$, of $\pm c(p)$, the sign being the orientation of $b$ in $\partial p$.

In each plane containing the direction of $b=(x,\mu)$, say the plane of the directions $\mu$ and $\nu$, exactly two plaquettes contain $b$. These are the plaquettes based at $x$ and at $x-\varepsilon e_\nu$, and $b$ enters their boundaries with opposite signs. Since $c$ is constant the two contributions cancel, so $d^*c=0$.

Hence $\langle da_0,d\delta\rangle=0$. The second variation $\|d\delta\|^2$ is non-negative.

This proves \eqref{8.5:eq-affine-first-variation-1}. By \eqref{8.5:eq-affine-first-variation-2}, $a_0$ therefore satisfies the equations of the constrained minimisation problem.
\end{proof}

\begin{proposition}[The constant row-sum identity of the minimiser kernel. Stated; proof incomplete
at the step marked $(\ast)$]\label{8.5:prop-row-sum-identity}
For every level $j$, let $\bar{M}_j$ be the linearised minimiser on fluxes at level $j$, let $Q_j$ be
the linearised block average at level $j$ and $\bar{Q}_j$ the linearised block average on fluxes at
level $j$, let $\eta_j$ be the ratio of the bare spacing to the
spacing of the level-$j$ lattice, and let $K_j(p,\mathfrak{P})$, for a bare plaquette $p$ and a
level-$j$ plaquette $\mathfrak{P}$, be the whole-lattice kernel of $\bar{M}_j$ on bounded
$2$-cochains, that is, the kernel to which the decay bound for the kernel of the linearised
minimiser applies. Then for every constant $u\in\mathbb{R}^{\mathfrak{O}}$
\begin{equation}\label{8.5:eq-row-sum-identity-a}
  \sum_{\mathfrak{P}} K_j(p,\mathfrak{P})\,u(\pi(\mathfrak{P}))
  \;=\;\eta_j^{2}\,u(\pi(p))
  \qquad\text{for every bare plaquette } p .
\end{equation}
In words, $\bar{M}_j$ maps a constant coarse flux to the constant fine flux multiplied by
$\eta_j^2$. Equivalently, the affine solution $a_0$ of the constrained first-variation equations
of Lemma~\ref{8.5:lem-affine-first-variation},
whose flux $d a_0$ is bounded, is identified with the output of the kernel on data of unbounded
support:
\begin{equation}\label{8.5:eq-row-sum-identity-1}
  d a_0 \;=\; \bar{M}_j(u).
\end{equation}
\end{proposition}

Stated; the proof below is incomplete at the step marked $(\ast)$.

On the torus $\mathbb{T}_m$, the identity \eqref{8.5:eq-row-sum-identity-a} is read on the box
$\hat{X}$ on which the kernel enters the two-point witness theorem, where it holds up to the
difference between the torus kernel and the whole-lattice kernel of $\bar{M}_j$, as bounded in the
comparison of the torus kernel of $\bar{M}_j$ with its whole-lattice kernel.

\begin{proof}
Fix a level $j$ and a constant $u\in\mathbb{R}^{\mathfrak{O}}$.

\emph{The right-hand side.} By Lemma~\ref{8.5:lem-affine-first-variation}, there is an affine bare
$1$-cochain $a_0$ with $Q_j a_0$ affine, $d Q_j a_0 = u$, and flux $d a_0$ constant, hence
bounded, and equal to $\eta_j^2 u$, that is, $d a_0(p)=\eta_j^2 u(\pi(p))$ for every bare
plaquette $p$. By the same lemma, $a_0$ satisfies the equations of the constrained minimisation
problem against every finitely supported competitor $\delta$ with $Q_j\delta=0$. Thus the
right-hand side of \eqref{8.5:eq-row-sum-identity-a} is $d a_0(p)$.

\emph{The left-hand side.} The left-hand side of \eqref{8.5:eq-row-sum-identity-a} is the value at
$p$ of $\bar{M}_j(u)$, the output of the kernel $K_j$ on the datum $u$; the datum is bounded but not
finitely supported. The sum converges absolutely by the decay bound for the kernel of the
linearised minimiser, which is the inequality
\begin{equation}\label{8.5:eq-row-sum-identity-2}
  |K_j(p,\mathfrak{P})|
  \;\le\; C_\Gamma\,\eta_j^{2}\,
  e^{-\delta_\Gamma \operatorname{dist}(p,\mathfrak{P})/\ell_\Gamma},
\end{equation}
with the distance measured in cells.

$(\ast)$ \emph{Identification.} It remains to show that $\bar{M}_j(u)=d a_0$, which is
\eqref{8.5:eq-row-sum-identity-1}. Both are bounded and both are defined on all of the lattice.
The identification requires a uniqueness statement in $\ell^\infty$, of Liouville type, for bounded
solutions of the constrained problem, or, alternatively, the decay of the kernel of
$(\bar{Q}_j\bar{Q}_j^{*})^{-1}$. Neither is contained in the statements used above: the input of
Lemma~\ref{8.5:lem-affine-first-variation} is that $Q_j$ maps closed cochains onto closed cochains,
and the decay bound is the inequality \eqref{8.5:eq-row-sum-identity-2}. The identity
\eqref{8.5:eq-row-sum-identity-a} is compatible with the
inequality expressing the $\eta_j^2$-scaling of the minimiser stated in \cite{B-CMP109}, but it is
not derived from that inequality.
\end{proof}

\emph{Route.} Two routes are available; neither is carried out here. Route ($\alpha$) would supply
the step $(\ast)$ under further hypotheses on the kernel; route ($\beta$) would make
\eqref{8.5:eq-row-sum-identity-a} unnecessary at the place where the two-point witness theorem
uses it.

(\emph{$\alpha$}) A periodicity argument on the one-cell torus. Granting the four properties of the
kernel listed at the end of this paragraph, the function $\varphi_u:=K_j u$ would be cell-periodic
and exact. Then
\begin{equation*}
  \psi:=\varphi_u-d a_0 ,
\end{equation*}
where $d a_0=\eta_j^2\,u(\pi(\cdot))$ by Lemma~\ref{8.5:lem-affine-first-variation}, would be
cell-periodic and exact with $\bar{Q}_j\psi=0$, so that, after subtraction of a constant,
$\psi=d\delta'$ with $Q_j\delta'=0$. The Euler--Lagrange relation at the level of the kernel,
$d^{*}(dM_j\xi)=Q_j^{*}\Delta_j\xi$ for coarse data $\xi$, where $M_j$ is the linearised
minimiser on bond variables at level $j$ and $\Delta_j$ is the operator on coarse data furnished,
together with this relation, by the Gaussian algebra of the finite-volume problem,
with $\Delta_j\xi$ a constant, denoted $\mu_0$, when $\xi$ is affine, would give on the one-cell
torus
\begin{equation*}
  \|d\delta'\|^{2}=\langle Q_j\delta',\mu_0\rangle=0 .
\end{equation*}
This route uses four properties of the kernel beyond \eqref{8.5:eq-row-sum-identity-2}: its columns
are exact; it is invariant under coarse translations; $\bar{Q}_jK_j$ is the identity on coarse
fluxes; and the Euler--Lagrange
relation above, obtained from the Gaussian algebra of the finite-volume problem passed to the limit
$\mathbb{Z}^4$ under the decay bound. Route ($\alpha$) takes these four properties as additional
hypotheses on the kernel.

(\emph{$\beta$}) A route internal to the torus, at level $k_0$. Let $v$ be the coarse datum of the
two-point witness theorem, a triangle wave; the finitely many layers of cells at which its slope
changes sign are the kink layers; off them the flux of $v$ is $\sigma c$, with $c$ a constant and
$\sigma$ a sign. Take a
bare $\tilde{a}$ with $Q_{k_0}\tilde{a}=v$ exactly and $d\tilde{a}=\eta_{k_0}^2\sigma c$ off the kink
layers, by the local right inverse of $Q_{k_0}$ on closed cochains furnished by the statement that
$Q_{k_0}$ maps closed cochains onto closed cochains, the input of
Lemma~\ref{8.5:lem-affine-first-variation}. Then
\begin{equation*}
  M_{k_0}v=\tilde{a}-G_Q\bigl(d^{*}d\tilde{a}\bigr),
\end{equation*}
with $d^{*}d\tilde{a}$ supported on the kink layers, where $G_Q$ is the kernel of the constrained
problem on $\ker Q_{k_0}$, which is the fluctuation kernel $C^{\mathrm{fl}}$; one uses its decay in
the flux variables, which is the decay lemma of the two-point witness theorem. This route
bypasses both \eqref{8.5:eq-row-sum-identity-a} and the comparison of the torus kernel of
$\bar{M}_j$ with its whole-lattice kernel.

\emph{Use of the identity.} By \eqref{8.5:eq-row-sum-identity-a}, the constant part of a coarse flux
is carried by $\bar{M}_j$ to the constant part of the fine flux, at every level of the sequence of
levels in \cite[(0.28)]{B-CMP109}, up to the factor $\eta_j^2$ (in dimension
$d=4$ the integral over a cell of a density of dimension four does not change under a change of
scale). This needs the derivative at the identity configuration $\mathbf{1}$ of Ba{\l}aban's
minimising configuration
$U_{k_0}(V)$, regarded as a function of the block field $V$, to be the linearised minimiser:
$U_{k_0}(\cdot)$ is twice
continuously differentiable at $\mathbf{1}$, with $DU_{k_0}(\mathbf{1})=M_{k_0}$. The
identification of the differential is used in the form stated in \cite{B-CMP95}, and
$U_{k_0}(\cdot)$ is analytic at $\mathbf{1}$ in Ba{\l}aban's construction of the
minimising configuration; the twice-differentiability, and its validity at every level, are
clauses of this book's formulation of the statement, and this book's proof of these two clauses
is incomplete.
This statement is an input of the two-point witness theorem: the proof of that theorem applies the
chain rule to
second order to $V\mapsto T(U_{k_0}(V))$, for a function $T$ of the bare configuration, and tests the
result on the directions $M_{k_0}v$.

\begin{definition}[Kernel decay and torus versus whole-lattice agreement. Stated; proof incomplete at the step marked $(\ast)$]
\label{8.5:def-kernel-decay-torus}
Let $k_0$ be the witness level, let $\eta$ be the bare spacing measured in units of the
level-$k_0$ spacing, and let $\bar{M}_{k_0}$ be the linearised minimiser at level $k_0$.
For a bare plaquette $p$ and a level-$k_0$ plaquette $\mathfrak{P}$ write $K(p,\mathfrak{P})$
for the kernel of $\bar{M}_{k_0}$ on the whole lattice, and $\operatorname{dist}(p,\mathfrak{P})$
for the distance between $p$ and $\mathfrak{P}$ measured in cells. The constants
$C_\Gamma$, $\delta_\Gamma$, $\ell_\Gamma$ below depend on $d$, $L$ and $M$ only.
\begin{enumerate}
\item[(i)] \emph{Kernel decay.} For all $p$ and $\mathfrak{P}$,
\begin{equation}\label{8.5:eq-kernel-decay-torus-1}
  |K(p,\mathfrak{P})|
  \le C_\Gamma\,\eta^2\,e^{-\delta_\Gamma \operatorname{dist}(p,\mathfrak{P})/\ell_\Gamma}.
\end{equation}
\item[(ii)] \emph{Torus versus whole-lattice agreement.} Let $\mathcal{N}$ be a multiple of $M$,
and let $K_{\mathrm{tor}}(p,\mathfrak{P})$ be the kernel of the linearised minimiser
$\bar{M}_{k_0}$ on the torus of side $\mathcal{N}$ cells. Then
\begin{equation}\label{8.5:eq-kernel-decay-torus-a}
  \sup_{p}\sum_{\mathfrak{P}}\bigl|K_{\mathrm{tor}}(p,\mathfrak{P})-K(p,\mathfrak{P})\bigr|
  \le C_\Gamma\,\eta^2\,e^{-\delta_\Gamma\mathcal{N}/(2\ell_\Gamma)} ,
\end{equation}
where, for $p$ and $\mathfrak{P}$ on the torus, the whole-lattice kernel $K(p,\mathfrak{P})$
is read through lifts of $p$ and $\mathfrak{P}$ to the whole lattice, the torus images being
summed.
\end{enumerate}
Part~(i) is taken by statement from \cite[Props.~2.3, 2.6, 2.7]{B-CMP96}, together with
\cite{B-CMP98} and \cite[\S\S C--E]{B-CMP99b}; it is used here as stated there, and its
proof is Ba{\l}aban's. Part~(ii) is stated; the proof below is incomplete at the step
marked $(\ast)$.

Part~(ii) is used at one place only, in the proof of the two-point witness theorem. On the torus a
constant $2$-cochain $c$ is closed but not exact, so that $\bar{M}_{k_0}$ applied to $c$ on the
torus is not defined; the identification, away from the kinks, of the coboundary $dB$ of the
field $B$ constructed in that proof with
$\pm\eta^2 c$ therefore uses both the whole-lattice row-sum identity for
$\bar{M}_{k_0}$ and the agreement \eqref{8.5:eq-kernel-decay-torus-a}. Part~(i) is used for
the kink layers of the coarse triangle wave and, together with the whole-lattice row-sum
statement for $\bar{M}_{k_0}$, in the same identification away from the kinks.
\end{definition}

\begin{proof}[Proof of part~\textup{(ii)}]
The argument for the bound \eqref{8.5:eq-kernel-decay-torus-a} rests on the finite-volume
construction of Ba{\l}aban's minimiser and propagator kernels in
\cite[Props.~2.3, 2.6, 2.7]{B-CMP96}, \cite{B-CMP98} and \cite[\S\S C--E]{B-CMP99b}; see also
the auxiliary tori of \cite{B-CMP109}. The route is to express both $K_{\mathrm{tor}}$ and $K$
through the random-walk and finite-range representations of Ba{\l}aban's
propagators on which the decay \eqref{8.5:eq-kernel-decay-torus-1} of part~(i) rests.

$(\ast)$ The comparison of the two expressions in the sup row-sum norm, the torus images
being summed, and the bound $C_\Gamma\eta^2e^{-\delta_\Gamma\mathcal{N}/(2\ell_\Gamma)}$
of \eqref{8.5:eq-kernel-decay-torus-a} obtained from it, are not carried out here.
\end{proof}

\begin{definition}[Profile locality and analyticity of sourced terms. Stated; proof incomplete at the
step marked $(\ast)$]\label{8.5:def-profile-locality}
Let $k_0=K-s$ be the witness level, $\ell=L^{-s}$ the side of a cell of the level-$k_0$ lattice,
$\eta=L^{-K}/\ell$ the bare spacing in units of $\ell$, and $x_1$ the centre of the first witness
ball $B_1=B(x_1,5\mathsf{d}/4)$. Consider the frozen sourced level-$k_0$ family with profile $w$ of
Definition~\ref{8.5:def-frozen-sourced-family}; its terms are $V^{(j),w}_X$, $j\le k_0$, and its
running weight is $\zeta_{k_0}(p)[w]$. For a functional $F$ of the profile and a profile $g$ put
\begin{equation*}
D[g]F:=\partial_z F[zg]\big|_{z=0},
\end{equation*}
the derivative at $z=0$ along the profile $zg$, and write $\mathbf{1}$ for the constant profile $1$.
For a term $V^{(j),w}_X$ let $\hat{X}$ be the filled hull of $X$, $c(X)$ its anchor cell with centre
$x_{c(X)}$, and $X^+\subset\hat{X}$ its reading neighbourhood, of radius $r(X)$ cells about
$x_{c(X)}$, where $r(X)$ is obtained by summing, over the levels $k\le k_0$, the localisation
radius of the step at level $k$, expressed in cells of the level-$k_0$ lattice (a geometric series).

The frozen sourced family is said to have \emph{profile locality and analyticity} if the following
hold: for every term $V^{(j),w}_X$ of the family meeting $B(x_1,2\mathsf{d})$, clauses (a) and (b)
below hold; and, for every plaquette $p$ of the small-field region near $B(x_1,2\mathsf{d})$ on
which the running weight is defined and every $R\ge 1$, the running weight splits into two
parts $\zeta^{\le R}_{k_0}(p)[w]$ and $\zeta^{>R}_{k_0}(p)[w]$ as in the first line of (a$'$), in
such a way that, for every profile $g$, the remaining lines of (a$'$) hold, and the two parts
satisfy (b):
\begin{equation}\label{8.5:eq-profile-locality-a}
\begin{aligned}
&\text{(a)}\quad && V^{(j),w}_X=V^{(j),w'}_X\quad\text{whenever } w|_{X^+}=w'|_{X^+};\\
&\text{(a}'\text{)}\quad && \zeta_{k_0}(p)[w]=\zeta^{\le R}_{k_0}(p)[w]+\zeta^{>R}_{k_0}(p)[w],\\
& && \zeta^{\le R}_{k_0}(p)[w]=\zeta^{\le R}_{k_0}(p)[w']\quad
  \text{whenever } w|_{B(p,R\ell)}=w'|_{B(p,R\ell)},\\
& && \bigl|D[g]\zeta^{\le R}_{k_0}(p)-g(x(p))\,D[\mathbf{1}]\zeta^{\le R}_{k_0}(p)\bigr|
  \le \tfrac{4}{3}\cdot 2\rho_\sharp\ell\,\operatorname{Lip} g,\\
& && \bigl|D[g]\zeta^{>R}_{k_0}(p)\bigr|\le C_w\,e^{-\hat\kappa R/(2M)}\sup|g|;\\
&\text{(b)}\quad && w\mapsto V^{(j),w}_X,\ \ w\mapsto\zeta^{\le R}_{k_0}(p)[w],\\
& && w\mapsto\zeta^{>R}_{k_0}(p)[w]\ \text{ are analytic on }\{\sup|w|<r_w\}\\
& && \text{and satisfy there the bounds of Ba{\l}aban's construction}\\
& && \text{in \cite{B-CMP109}, with one } r_w>0 \text{ independent of } K,\ s,\ m\\
& && \text{and the position.}
\end{aligned}
\end{equation}
In (a) the term depends on $w$ through $w|_{X^+}$ only. In (a$'$) the part
$\zeta^{\le R}_{k_0}(p)$ reads $w$ on the ball $B(p,R\ell)$ only, and the third line says that its
first-order, mass-weighted reading radius is at most $2\rho_\sharp\ell$, where $\rho_\sharp$ is the
mass-weighted mean radius; $\hat\kappa$ is the halved decay rate of the majorant masses and $C_w$ a
constant. In (b) the small-field conditions under which these bounds are stated are read, for this
use, as a regularity space of fine fields of the kind of the complex configuration spaces
$\mathfrak{U}^c_k$ of \cite[(0.29)]{B-CMP109}. As a consequence of (a) and (b), the map
$f\mapsto H_X[f]$ to the marked Hessian is linear, and
\begin{equation}\label{8.5:eq-profile-locality-1}
\begin{aligned}
\mathfrak{m}^{\mathrm{reg}}_X[f]
&\le r_w^{-1}\sup\bigl\{\mathfrak{m}^{\mathrm{reg}}_{V^{(j),w}_X}:\ \sup|w|<r_w\bigr\}
   \cdot\sup_{X^+}|f|\\
&\le \hat{\mathfrak{m}}_X\sup_{X^+}|f|,
\end{aligned}
\end{equation}
where $\mathfrak{m}^{\mathrm{reg}}$ is the regular-direction mass of item~(e) of
Lemma~\ref{8.5:lem-constant-flux-forms} (see \eqref{8.5:eq-constant-flux-forms-a}); the supremum in
the middle member is at most the majorant $\hat{\mathfrak{m}}^0_X$ given by these bounds, and
$\hat{\mathfrak{m}}_X$ is the majorant formed from $\hat{\mathfrak{m}}^0_X$ in which the route
weight is taken to be $\bar{\mathfrak{w}}:=r_w^{-1}$ (the unweighted choice). The weighted
domination, with $\bar{\mathfrak{w}}:=C\mathfrak{w}_H$ or $\bar{\mathfrak{w}}:=C\mathfrak{w}_W$ for a
suitable constant $C$, is a separate statement of this chapter and is not part of (b).

\emph{On the norm in} (b). Clause (b) is written in the supremum norm of $w$, following the bound
\begin{equation*}
|\zeta_k(p)|\le\tfrac{4}{3}\sup|w|
\end{equation*}
of the sourced-format lemma of the correlation theorem, whereas the torus bound of the
correlation theorem weighs the test functions by
$\|f\|_\sharp=\max\{\|f\|_\infty,40M\operatorname{Lip}f\}$ of
Definition~\ref{1.3:def-curvature-field}, in units of the bare lattice. Whether the analyticity
ball in $w$ of the correlation theorem at level $k_0$ is a ball for the supremum norm or a ball for
$\|\cdot\|_\sharp$ is not decided here. If it is a
$\|\cdot\|_\sharp$-ball, with weight $40M\operatorname{Lip}$ in bare units, the extraction constant
$C_{\mathrm{ext}}$ remains absolute and uniform in $s$.

The frozen sourced family of Definition~\ref{8.5:def-frozen-sourced-family} has profile locality
and analyticity.
\end{definition}

\begin{proof}
$(\ast)$ The frozen sourced family of Definition~\ref{8.5:def-frozen-sourced-family} is the family
of the sourced-format lemma of the correlation theorem (stated; proof incomplete). That statement
asserts, among other things, that the terms are analytic in $w$, that they read $w$ on
$X^+$ only, that they satisfy the bounds of (b) under the small-field conditions, and that
$|\zeta_k(p)|\le\frac{4}{3}\sup|w|$. Clauses (a), (a$'$) --- including the splitting of the running
weight and its two bounds --- and the analyticity assertion of (b) are taken from that statement as
stated, in the form (a), (a$'$), (b) above, with the analyticity
domain in $w$ taken to be the sup-ball $\{\sup|w|<r_w\}$, with $r_w>0$ independent of $K$, $s$,
$m$ and position, and with the small-field conditions read as a regularity space of fine fields of
the kind of \cite[(0.29)]{B-CMP109}. The value of the radius $r_w$, and whether the ball is a
sup-ball or a $\|\cdot\|_\sharp$-ball, are not derived here: the proof is incomplete at this step.

\emph{The consequences of} (b). Since the terms are analytic in $w$ on the sup-ball, which is a
neighbourhood of $w=0$, the map $f\mapsto H_X[f]$ is linear. For the first inequality in
\eqref{8.5:eq-profile-locality-1}, fix a profile $f$, and put $f^+:=f$ on $X^+$ and $f^+:=0$
elsewhere. For $|z|$ small enough that $zf$ and $zf^+$ both lie in $\{\sup|w|<r_w\}$, the profiles
$zf$ and $zf^+$ agree on $X^+$, so by (a) $V^{(j),zf}_X=V^{(j),zf^+}_X$; hence
$H_X[f]=H_X[f^+]$. If $\sup_{X^+}|f|=0$, then $f^+=0$, so $V^{(j),zf^+}_X=V^{(j),0}_X$ for every
$z\in\mathbb{C}$, $H_X[f]=0$, and both sides of the first inequality vanish. Otherwise the profile
$zf^+$, $z\in\mathbb{C}$, has supremum $|z|\sup_{X^+}|f|$, and it lies in $\{\sup|w|<r_w\}$ for
\begin{equation*}
|z|<r_w\Bigl(\sup_{X^+}|f|\Bigr)^{-1}.
\end{equation*}
By (b), $z\mapsto V^{(j),zf^+}_X$ is therefore analytic on this disc. The Cauchy estimate in $z$ on
this disc, whose reciprocal radius is $r_w^{-1}\sup_{X^+}|f|$, together with the Cauchy estimate in
the field, gives the first inequality. The mass on the left is the regular-direction mass of
item~(e) of Lemma~\ref{8.5:lem-constant-flux-forms}: the Cauchy disc in the field realises the
domain of \cite[(0.29)]{B-CMP109} on regular directions only. The second inequality in
\eqref{8.5:eq-profile-locality-1} holds because the supremum in the middle member is at most the
majorant $\hat{\mathfrak{m}}^0_X$ given by the bounds of (b), with the route weight
$\bar{\mathfrak{w}}:=r_w^{-1}$ (unweighted) in $\hat{\mathfrak{m}}_X$.
\end{proof}

This is the only place in the chapter where the size of the domain of \cite[(0.29)]{B-CMP109}
enters. It makes the ratio $\bar{\mathfrak{a}}^E/\mathfrak{n}$ and the sum over $j$ finite,
uniformly in $K$ and $s$: the factor $(L^j\eta)^4$ controls the count, and the factor
$L^{-\alpha(k_0-j)/2}$ controls the sum.
Of the sourced-format lemma of the correlation theorem, exactly the content of (a), (a$'$) and (b)
is used in this chapter; a constant profile enters only through the case $f\equiv\mathbf{1}$, in
the mean statement of this chapter.

\begin{lemma}[Radius tails of the majorant masses]\label{8.5:lem-radius-tails}
Let $k_0=K-s$ be the witness level, and consider the terms $V^{(j),w}_X$, $j\le k_0$,
$X\in\mathcal{D}_j$, of the frozen sourced family of
Definition~\ref{8.5:def-frozen-sourced-family}. For each such $X$ let
$\hat{X}\supset X^+\supset X$ be its filled hull and its reading neighbourhood, let $c(X)$ be
its anchor cell, that is, the level-$k_0$ cell containing the anchoring corner of $\hat{X}$,
let $r(X)$ be the radius of $\hat{X}$ in level-$k_0$ cells, and let $r_{\mathrm{read}}$ be the
reading collar, all as in \eqref{8.5:eq-profile-locality-a}. Put $\hat\kappa=\kappa/2$.
We recall the majorant masses of this chapter. Here and below a term is identified with its
pair $(j,X)$, the level $j$ being suppressed in the notation, and sums over
$X\in\operatorname{Irr}_1^E$ (or $X\in\operatorname{Irr}_1^R$) run over the terms of that class
of all levels $j\le k_0$. The majorant mass of a term of
level $j$ of the extracted class $\operatorname{Irr}_1^E$ is
\begin{equation}\label{8.5:eq-radius-tails-1}
  \hat{\mathfrak{m}}_X=\bar{\mathfrak{w}}\,C_H\,\mathfrak{n}\,
  L^{-(4+\alpha/2)(k_0-j)}\,e^{-\hat\kappa d_j(X)},
\end{equation}
where $\alpha>0$ is the exponent of hypothesis (i) below and $\alpha_0$ its field radius,
$\mathfrak{n}=6(\ell/\varepsilon)^4$ is the main-term mass per cell,
$C_H=C_H(\alpha_0,M,\kappa)$ is the constant of the Cauchy estimate in the fine bond field by
which the bound of hypothesis (i) below is converted into a bound on Hessians, and
$\bar{\mathfrak{w}}=r_w^{-1}$ is the inverse of the radius $r_w$ of the sup-ball of
hypothesis (iii), a number independent of $j$ and of $X$. The majorant mass of a term of level
$j$ of the remainder class $\operatorname{Irr}_1^R$ is
\begin{equation}\label{8.5:eq-radius-tails-9}
  \hat{\mathfrak{m}}_X=\bar{\mathfrak{w}}^R\,C_H\,\mathfrak{n}\,g_j^{\kappa_0}\,
  L^{-4(k_0-j)}\,e^{-\hat\kappa_R d_j(X)},
\end{equation}
of the shape of \cite[(2.31)]{B-CMP119}, where $\kappa_0$ is the exponent of the coupling in
that bound, $\hat\kappa_R$ is the decay rate of the remainder majorants and
$\bar{\mathfrak{w}}^R$ is a weight independent of the level $j$. For a cell $c$ put
\begin{equation*}
  \mathfrak{a}^E_c=\sum_{X\in\operatorname{Irr}_1^E,\ c(X)=c}\hat{\mathfrak{m}}_X,
  \qquad \bar{\mathfrak{a}}^E=\sup_c\mathfrak{a}^E_c,
\end{equation*}
and define $\mathfrak{a}^R_c$, $\bar{\mathfrak{a}}^R$ in the same way for the remainder class
$\operatorname{Irr}_1^R$. Assume:
\begin{itemize}
\item[(i)] the bound of \cite[(0.29)]{B-CMP109} on the fine-field domain: with
  $\eta=\eta_{k_0}$, the terms $V^{(j)}(X)$ of the effective action of \cite{B-CMP109} at
  step $k_0$ are analytic on the complex fine-field domain of configurations
  $U$ with $|\partial U-1|<\alpha_0\eta^2$ and $U$ regular at scale $L^j\eta$ in the sense of
  \cite{B-CMP109}, a subset of the complex configuration spaces
  $\mathfrak{U}^c_{k_0}$ of \cite[(1.18)]{B-CMP109}, and satisfy there
  $|V^{(j)}(X)|\le(L^j\eta)^{4+\alpha}e^{-\kappa d_j(X)}$ with an exponent $\alpha>0$,
  where $L^j\eta=L^{-(k_0-j)}$;
\item[(ii)] the animal sum \cite[(0.26)]{B-CMP109}: for every cube $\square$ of $\pi_j$ and
  every integer $n\ge0$, the number of $X\in\mathcal{D}_j$ with $\square\subset X$ and
  $d_j(X)=n$ is at most $C_{\mathsf{a}}^{\,n}$, where $C_{\mathsf{a}}$ is the constant of that
  bound;
\item[(iii)] the sup-ball clause of profile locality and analyticity of sourced terms,
  Definition~\ref{8.5:def-profile-locality}, display \eqref{8.5:eq-profile-locality-a},
  which supplies the radius $r_w$ of the sup-ball on which the sourced terms are analytic in
  the profile $w$; by it the weight $\bar{\mathfrak{w}}=r_w^{-1}$ in
  \eqref{8.5:eq-radius-tails-1} is one number, independent of the level $j$;
\item[(iv)] (for (b) only) the rate $\hat\kappa_R$ satisfies
  $\hat\kappa_R\ge2\log(2C_{\mathsf{a}})+2$, and there is a number $C_R\ge1$, independent of
  $K$, $s$, $m$ and of position, such that
  \begin{equation}\label{8.5:eq-radius-tails-10}
    \bar{\mathfrak{w}}^R\,C_H\,\mathfrak{n}\sum_{j\le k_0}g_j^{\kappa_0}
    \le C_R\,\bar{\mathfrak{a}}^R ;
  \end{equation}
\end{itemize}
and assume
\begin{equation}\label{8.5:eq-radius-tails-2}
  \kappa\ge 4\log(2C_{\mathsf{a}})+4,\qquad\text{equivalently}\qquad
  \hat\kappa\ge 2\log(2C_{\mathsf{a}})+2 .
\end{equation}
Condition \eqref{8.5:eq-radius-tails-2} is a one-sided lower bound on $\kappa$, one of the
lower bounds imposed on $\kappa$ at the stage of the order of choice at which $\kappa$ is
fixed (Notation~\ref{2.3:not-stages-middle}); the constant $C_w$ below increases with
$\kappa$ and enters only constants chosen after $\kappa$.
Then the following hold.
\begin{itemize}
\item[(a)] There is $C_w=C_w(d,L,M,\kappa,\alpha)\ge1$, $d$ the dimension, depending in
  addition on the constant $C_{\mathsf{a}}$ of (ii) and on $r_{\mathrm{read}}$, such that for
  every $R\ge1$ and every cell $c$, with $\bar{\mathfrak{a}}=\bar{\mathfrak{a}}^E$,
  \begin{equation}\label{8.5:eq-radius-tails-a}
    \sum_{\substack{X\in\operatorname{Irr}_1^E\\ c(X)=c,\ r(X)>R}}
    \bigl(1+r(X)\bigr)\,\hat{\mathfrak{m}}_X
    \;\le\; C_w\,e^{-\hat\kappa R/(2M)}\,\bar{\mathfrak{a}} .
  \end{equation}
\item[(b)] If in addition (iv) holds, the bound \eqref{8.5:eq-radius-tails-a} holds for the
  remainder class, with $\operatorname{Irr}_1^R$ in place of $\operatorname{Irr}_1^E$, the rate
  $\hat\kappa_R$ in place of $\hat\kappa$, $\bar{\mathfrak{a}}^R$ in place of
  $\bar{\mathfrak{a}}^E$, and a constant $C^R_w\ge1$ in place of $C_w$, depending on $d$,
  $M$, $\hat\kappa_R$, $C_{\mathsf{a}}$, $r_{\mathrm{read}}$ and $C_R$ only.
\item[(c)] For the extracted class, that is, for the terms indexed by the localisation domains
  of \cite{B-CMP109} with majorants of the shape \eqref{8.5:eq-radius-tails-1}, the
  mass-weighted mean radius
  \begin{equation*}
    \rho^E_\sharp=\sup_{c:\ \mathfrak{a}^E_c>0}\
    \frac{1}{\mathfrak{a}^E_c}\sum_{\substack{X\in\operatorname{Irr}_1^E\\ c(X)=c}}
    r(X)\,\hat{\mathfrak{m}}_X ,
  \end{equation*}
  cells with $\mathfrak{a}^E_c=0$ being excluded from the supremum, satisfies
  $\rho^E_\sharp\le C_\rho M$ with $C_\rho=C_\rho(\kappa,L,M,\alpha)$, depending in addition
  on $C_{\mathsf{a}}$ and $r_{\mathrm{read}}$ through $C_w$. Of the mass-weighted mean radius
  $\rho_\sharp=\max\{\rho^E_\sharp,\rho^R_\sharp\}$, with $\rho^R_\sharp$ defined in the same
  way over the remainder class, only the first member is bounded here; the members of the
  remainder class of the shape of \cite[(2.31)]{B-CMP119} are not covered by (c).
\end{itemize}
The constants $C_w$, $C^R_w$ and $C_\rho$ depend on none of $K$, $s$, $m$ and on no position.
\end{lemma}

Hypothesis (i) uses the shape $(L^j\eta)^{4+\alpha}e^{-\kappa d_j(X)}$ of the bound
\cite[(0.29)]{B-CMP109} on the fine-field domain described in (i), which is a domain other
than the one on which \cite[Theorem~1]{B-CMP109} states that bound. This fine-field form of the
bound is stated; proof incomplete, the step not supplied being the passage from the domain of
configurations on which \cite[Theorem~1]{B-CMP109} states \cite[(0.29)]{B-CMP109} to the
fine-field domain of (i). The clauses of Definition~\ref{8.5:def-profile-locality} used in (iii)
and for the reading collar $r_{\mathrm{read}}$ are likewise stated; proof incomplete at one
step. The lemma is proved as an implication from
(i), (ii) and (iii), and, for (b), (iv). Hypothesis (iv), on which part (b) alone depends,
concerns the remainder class and its rate $\hat\kappa_R$ only.

\begin{proof}
Throughout, $\hat{\mathfrak{m}}_X$ denotes the majorant \eqref{8.5:eq-radius-tails-1}, whose
shape in $j$ and $d_j(X)$ is the shape of the bound in hypothesis (i): the factor
$(L^j\eta)^{4+\alpha}=L^{-(4+\alpha)(k_0-j)}$ of (i) splits into the marginal factor
$(L^j\eta)^4=L^{-4(k_0-j)}$ and the spare factor $(L^j\eta)^\alpha=L^{-\alpha(k_0-j)}$, of
which \eqref{8.5:eq-radius-tails-1} retains $L^{-\alpha(k_0-j)/2}$. The marginal factor will
pay the number of domains anchored in a cell, and the retained part of the spare factor will
pay the sum over levels $j$.

\emph{Radius and tree length.} Let $X\in\mathcal{D}_j$ have filled hull
$\hat{X}\supset X^+\supset X$ of radius $r=r(X)$ cells. By the definition of $X^+$ and of
$r_{\mathrm{read}}$ in \eqref{8.5:eq-profile-locality-a}
(Definition~\ref{8.5:def-profile-locality}), $X^+$ lies in $\hat{X}$ and differs from $X$ by
at most the reading collar $r_{\mathrm{read}}$, so $X$ spans a distance at least
$r-r_{\mathrm{read}}$ cells. Since $L^j\eta=L^{-(k_0-j)}\le1$, the cubes of $\pi_j$ have side
at most $M$ cells; a tree through a union of such cubes spanning distance $r-r_{\mathrm{read}}$
has length at least $(r-r_{\mathrm{read}})/M$ in units of $M$-cubes. Hence
\begin{equation}\label{8.5:eq-radius-tails-3}
  d_j(X)\ge \frac{r(X)-r_{\mathrm{read}}}{M}-1,
  \qquad\text{that is,}\qquad
  r(X)\le M\bigl(d_j(X)+1\bigr)+r_{\mathrm{read}} .
\end{equation}
Multiplying the first inequality by $-\hat\kappa/2$ and exponentiating,
\begin{equation}\label{8.5:eq-radius-tails-4}
  e^{-\hat\kappa d_j(X)/2}\le e^{\hat\kappa/2}\,e^{\hat\kappa r_{\mathrm{read}}/(2M)}\,
  e^{-\hat\kappa r(X)/(2M)} .
\end{equation}
The factor $e^{\hat\kappa/2}$ comes from the term $-1$ in \eqref{8.5:eq-radius-tails-3}, and
the factor $e^{\hat\kappa r_{\mathrm{read}}/(2M)}$ from the collar; both will be put into
$C_w$.

\emph{Counting anchored domains.} Fix a cell $c$, a level $j\le k_0$ and $n\ge0$. A domain
$X\in\mathcal{D}_j$ with $c(X)=c$ is anchored by the corner of its hull $\hat{X}\supset X^+$,
not by a cube of $X$ lying in $c$. If $d_j(X)=n$, then by \eqref{8.5:eq-radius-tails-3} the
hull, and hence every cube of $\pi_j$ contained in $X$, lies within distance
$(n+1)M+r_{\mathrm{read}}$ of $c$. Let $\mathfrak{p}_{\mathsf{a}}(n)$ be the number of
level-$k_0$ cells at distance at most $(n+1)M+r_{\mathrm{read}}$ from $c$. These cells lie in
the cube of side $2(n+1)M+2r_{\mathrm{read}}+3$ cells centred at $c$, so that
\begin{equation*}
  \mathfrak{p}_{\mathsf{a}}(n)\le\bigl(2(n+1)M+2r_{\mathrm{read}}+3\bigr)^4\le C_d\,2^n,
  \qquad
  C_d:=\sup_{n\ge0}2^{-n}\bigl(2(n+1)M+2r_{\mathrm{read}}+3\bigr)^4 ;
\end{equation*}
the number $C_d\ge1$ is finite and depends on $M$ and $r_{\mathrm{read}}$. A level-$k_0$ cell
contains the lower corners of at most $1+L^{4(k_0-j)}M^{-4}$ cubes of $\pi_j$. Every $X$
counted contains a cube of $\pi_j$ whose lower corner lies in one of the
$\mathfrak{p}_{\mathsf{a}}(n)$ cells, so at most
$(1+L^{4(k_0-j)}M^{-4})\,\mathfrak{p}_{\mathsf{a}}(n)$ cubes are to be considered, and each of
them lies, by hypothesis (ii), in at most $C_{\mathsf{a}}^{\,n}$ domains $X\in\mathcal{D}_j$
with $d_j(X)=n$. Therefore
\begin{equation}\label{8.5:eq-radius-tails-5}
\begin{split}
  \#\{X\in\mathcal{D}_j:\ c(X)=c,\ d_j(X)=n\}
  &\le\bigl(1+L^{4(k_0-j)}M^{-4}\bigr)\,\mathfrak{p}_{\mathsf{a}}(n)\,C_{\mathsf{a}}^{\,n}\\
  &\le C_d\bigl(1+L^{4(k_0-j)}M^{-4}\bigr)(2C_{\mathsf{a}})^n .
\end{split}
\end{equation}

\emph{The marginal factor pays the count.} Since $k_0\ge j$ and $M\ge1$,
\begin{equation}\label{8.5:eq-radius-tails-6}
  L^{-4(k_0-j)}\bigl(1+L^{4(k_0-j)}M^{-4}\bigr)=L^{-4(k_0-j)}+M^{-4}\le 2 .
\end{equation}

\emph{Summation over shells.} Fix $j$ and $n$, and consider the domains counted in
\eqref{8.5:eq-radius-tails-5} with $r(X)>R$. For each of them write
$e^{-\hat\kappa d_j(X)}=e^{-\hat\kappa n/2}\,e^{-\hat\kappa d_j(X)/2}$, bound the second
factor by \eqref{8.5:eq-radius-tails-4} and then by
$e^{\hat\kappa/2}e^{\hat\kappa r_{\mathrm{read}}/(2M)}e^{-\hat\kappa R/(2M)}$ (since
$r(X)>R$), and bound $1+r(X)$ by $1+r_{\mathrm{read}}+M(n+1)$ using
\eqref{8.5:eq-radius-tails-3}. Inserting \eqref{8.5:eq-radius-tails-1},
\eqref{8.5:eq-radius-tails-5} and \eqref{8.5:eq-radius-tails-6}, and using hypothesis (iii),
by which $\bar{\mathfrak{w}}$ does not depend on $j$, we get
\begin{equation}\label{8.5:eq-radius-tails-7}
\begin{split}
  \sum_{\substack{X\in\mathcal{D}_j,\ c(X)=c\\ d_j(X)=n,\ r(X)>R}}
  \bigl(1+r(X)\bigr)\hat{\mathfrak{m}}_X
  &\le 2C_d\,\bar{\mathfrak{w}}C_H\mathfrak{n}\,L^{-\alpha(k_0-j)/2}\,
  e^{\hat\kappa/2}e^{\hat\kappa r_{\mathrm{read}}/(2M)}e^{-\hat\kappa R/(2M)}\\
  &\quad\times\bigl(1+r_{\mathrm{read}}+M(n+1)\bigr)
  (2C_{\mathsf{a}})^n e^{-\hat\kappa n/2} .
\end{split}
\end{equation}
By \eqref{8.5:eq-radius-tails-2}, $2C_{\mathsf{a}}e^{-\hat\kappa/2}\le e^{-1}$, so that
$(2C_{\mathsf{a}})^ne^{-\hat\kappa n/2}\le e^{-n}$. With
$\sum_{n\ge0}e^{-n}=(1-e^{-1})^{-1}$ and $\sum_{n\ge0}(n+1)e^{-n}=(1-e^{-1})^{-2}$,
\begin{equation*}
  \sum_{n\ge0}\bigl(1+r_{\mathrm{read}}+M(n+1)\bigr)e^{-n}
  \le\bigl(1+r_{\mathrm{read}}+M\bigr)(1-e^{-1})^{-2} .
\end{equation*}

\emph{The spare factor pays the sum over levels.} Finally
\begin{equation}\label{8.5:eq-radius-tails-8}
  \sum_{j\le k_0}L^{-\alpha(k_0-j)/2}\le\bigl(1-L^{-\alpha/2}\bigr)^{-1},
\end{equation}
a bound on a geometric series, valid since $\alpha>0$ and $L>1$. Summing
\eqref{8.5:eq-radius-tails-7} over $n\ge0$ and $j\le k_0$, the left member of
\eqref{8.5:eq-radius-tails-a} is at most
\begin{equation*}
  2C_d\,e^{\hat\kappa/2}e^{\hat\kappa r_{\mathrm{read}}/(2M)}
  \bigl(1+r_{\mathrm{read}}+M\bigr)(1-e^{-1})^{-2}\bigl(1-L^{-\alpha/2}\bigr)^{-1}\,
  \bar{\mathfrak{w}}C_H\mathfrak{n}\,e^{-\hat\kappa R/(2M)} .
\end{equation*}
The quantity $\bar{\mathfrak{a}}=\bar{\mathfrak{a}}^E$ is a supremum over cells of sums of the
same majorants \eqref{8.5:eq-radius-tails-1}. We use the following property of the
localisation domains, the same on which the bound on the ratio in (c) below rests: every cell
at a corner of a cube of $\pi_{k_0}$ anchors the one-cube domain of level $k_0$ formed by that
cube, a member of the extracted class with $d_{k_0}(X)=0$ and hence, by
\eqref{8.5:eq-radius-tails-1}, of majorant $\bar{\mathfrak{w}}C_H\mathfrak{n}$. Hence
$\bar{\mathfrak{w}}C_H\mathfrak{n}\le\bar{\mathfrak{a}}$, and the shells sum to
$C_we^{-\hat\kappa R/(2M)}\bar{\mathfrak{a}}$, where $C_w\ge1$ collects the displayed
factors; through $C_d$, the collar and the exponent $\alpha$ it depends on $d$, $L$, $M$,
$\kappa$, $\alpha$, $C_{\mathsf{a}}$ and $r_{\mathrm{read}}$. This proves (a). No quantity in
$C_w$ depends on $K$, $s$, $m$ or $c$:
the dependence on $k_0-j$ has been summed in \eqref{8.5:eq-radius-tails-8} by means of the
exponent $\alpha$ supplied by hypothesis (i), and $\mathfrak{n}$ and $\bar{\mathfrak{w}}$
occur only through the comparison with $\bar{\mathfrak{a}}$.

For (b), let $X$ range over the remainder class, with the majorants
\eqref{8.5:eq-radius-tails-9}. The bounds \eqref{8.5:eq-radius-tails-3},
\eqref{8.5:eq-radius-tails-5} and \eqref{8.5:eq-radius-tails-6} do not involve the majorants
and hold unchanged, \eqref{8.5:eq-radius-tails-4} holds with $\hat\kappa_R$ in place of
$\hat\kappa$, and the factor $L^{-4(k_0-j)}$ of \eqref{8.5:eq-radius-tails-9} is the marginal
factor of \eqref{8.5:eq-radius-tails-6}. Writing
$e^{-\hat\kappa_R d_j(X)}=e^{-\hat\kappa_R n/2}e^{-\hat\kappa_R d_j(X)/2}$ and proceeding as
for \eqref{8.5:eq-radius-tails-7}, with $\bar{\mathfrak{w}}^R$ independent of $j$, we get
\begin{equation*}
\begin{split}
  \sum_{\substack{X\in\mathcal{D}_j,\ c(X)=c\\ d_j(X)=n,\ r(X)>R}}
  \bigl(1+r(X)\bigr)\hat{\mathfrak{m}}_X
  &\le 2C_d\,\bar{\mathfrak{w}}^RC_H\mathfrak{n}\,g_j^{\kappa_0}\,
  e^{\hat\kappa_R/2}e^{\hat\kappa_R r_{\mathrm{read}}/(2M)}e^{-\hat\kappa_R R/(2M)}\\
  &\quad\times\bigl(1+r_{\mathrm{read}}+M(n+1)\bigr)
  (2C_{\mathsf{a}})^n e^{-\hat\kappa_R n/2} .
\end{split}
\end{equation*}
By the floor on $\hat\kappa_R$ in (iv), $(2C_{\mathsf{a}})^ne^{-\hat\kappa_R n/2}\le e^{-n}$,
and the sum over $n\ge0$ is bounded as in (a). The sum over $j\le k_0$ of
$\bar{\mathfrak{w}}^RC_H\mathfrak{n}\,g_j^{\kappa_0}$ is at most $C_R\bar{\mathfrak{a}}^R$ by
\eqref{8.5:eq-radius-tails-10}; this replaces \eqref{8.5:eq-radius-tails-8} and the comparison
with the one-cube domain. Hence (b) holds with
\begin{equation*}
  C^R_w=2C_d\,e^{\hat\kappa_R/2}e^{\hat\kappa_R r_{\mathrm{read}}/(2M)}
  \bigl(1+r_{\mathrm{read}}+M\bigr)(1-e^{-1})^{-2}\,C_R\ \ge1 .
\end{equation*}

For (c), fix a cell $c$ with $\mathfrak{a}^E_c>0$; all sums below run over
$X\in\operatorname{Irr}_1^E$. Splitting the sum at $r(X)=R$, using $r(X)\le R$ in the first
part and $r(X)\le1+r(X)$ together with (a) in the second,
\begin{equation*}
  \frac{1}{\mathfrak{a}^E_c}\sum_{c(X)=c}r(X)\hat{\mathfrak{m}}_X
  \le R\,\frac{1}{\mathfrak{a}^E_c}\sum_{\substack{c(X)=c\\ r(X)\le R}}\hat{\mathfrak{m}}_X
  +\frac{1}{\mathfrak{a}^E_c}\sum_{\substack{c(X)=c\\ r(X)>R}}
  \bigl(1+r(X)\bigr)\hat{\mathfrak{m}}_X ,
\end{equation*}
and therefore
\begin{equation*}
  \rho^E_\sharp\le R+C_we^{-\hat\kappa R/(2M)}\,
  \frac{\bar{\mathfrak{a}}}{\min\{\mathfrak{a}^E_c:\ \mathfrak{a}^E_c>0\}} .
\end{equation*}
Bound (a) is relative to $\bar{\mathfrak{a}}$, while the denominator in $\rho^E_\sharp$ is
$\mathfrak{a}^E_c$; this is why the ratio appears. For the terms of the extracted class the
ratio $\bar{\mathfrak{a}}/\min\{\mathfrak{a}^E_c:\mathfrak{a}^E_c>0\}$ is a number depending
on $L$, $M$, $\alpha$ and $r_{\mathrm{read}}$ only. For the numerator, multiplying the count
\eqref{8.5:eq-radius-tails-5} by the majorant \eqref{8.5:eq-radius-tails-1}, without the
restriction $r(X)>R$ and without the factor $1+r(X)$, using \eqref{8.5:eq-radius-tails-6},
the bound $(2C_{\mathsf{a}})^ne^{-\hat\kappa n}\le(2C_{\mathsf{a}})^ne^{-\hat\kappa n/2}
\le e^{-n}$ from \eqref{8.5:eq-radius-tails-2}, and \eqref{8.5:eq-radius-tails-8}, and
summing over $n\ge0$ and $j\le k_0$, gives for every cell $c$
\begin{equation*}
  \mathfrak{a}^E_c\le 2C_d\,(1-e^{-1})^{-1}\bigl(1-L^{-\alpha/2}\bigr)^{-1}\,
  \bar{\mathfrak{w}}C_H\mathfrak{n},
\end{equation*}
and hence the same bound for $\bar{\mathfrak{a}}$. For the denominator, cells not at corners
of cubes of $\pi_{k_0}$ anchor only levels
$j$ with $L^{k_0-j}>M$, of mass of order $(ML)^{-\alpha}$ relative to the one-cube majorant
$\bar{\mathfrak{w}}C_H\mathfrak{n}$,
while cells with $\mathfrak{a}^E_c=0$ are excluded from the supremum defining $\rho^E_\sharp$.
The weight $\bar{\mathfrak{w}}$ cancels in the ratio. Taking $R=M\ge1$ gives
\begin{equation*}
  \rho^E_\sharp\le M+C_we^{-\hat\kappa/2}\,
  \frac{\bar{\mathfrak{a}}}{\min\{\mathfrak{a}^E_c:\ \mathfrak{a}^E_c>0\}},
\end{equation*}
which is at most $C_\rho M$ with $C_\rho$ depending on $\kappa$, $L$, $M$, $\alpha$, and on
$C_{\mathsf{a}}$ and $r_{\mathrm{read}}$ through $C_w$. Since $M$ and $L$ are fixed before
the floor $C_0$ of the constant $C_W$ is chosen, and none of the constants depends on $K$,
$s$, $m$ or position, (c) follows.
\end{proof}

\begin{remark}\label{8.5:prop-extraction-complete}
The statement of this proposition is not written out here; parts of its proof are printed below.
\end{remark}

\begin{proof}[Proof of Proposition~\ref{8.5:prop-extraction-complete}, continued: the torus test on
linearised-minimiser directions]
\label{8.5:prf-torus-test}
In this part of the proof the steps at which the proof is incomplete are marked $(\ast)$. Here we
derive the vanishing of the $\Gamma$-class sum of the constant-flux forms of the residual family
from the corresponding identity. We first do so for Ba{\l}aban's own
$\beta$-extracted family, and then for the frozen sourced family of the correlation theorem at
constant shift $\zeta$ $(\ast)$, which is the $\beta$-extracted part $\operatorname{Irr}_1^E$ of the
two-point witness.

\emph{The identity and the directions on which it is used.} The corresponding identity compares two representations of the same effective action
$A_{k_0}$. The first is the counterterm representation of \cite[Theorem~3, (1.3)--(1.22)]{B-CMP109},
resp.\ its inductive form \cite[(0.22)--(0.25)]{B-CMP109}, whose terms $E^{(j)}$ satisfy the bounds
\cite[(1.18)]{B-CMP109}. The second is the renormalised representation of
\cite[Theorem~1]{B-CMP109}, with the bounds \cite[(0.29)]{B-CMP109}. The identity is summed over
$j\le k_0$, and only the sum over $j$ is used below; the form for each $j$ separately would require
the regrouping of \cite{B-CMP109} level by level and is not used.

It is used only along linearised-minimiser directions $B=M_{k_0}v$, with $v$ a coarse $1$-cochain,
and not along every bare direction $B$. In \cite[Theorem~1]{B-CMP109}, \cite[(0.28)]{B-CMP109} is an
identity in the coarse variable $V$, and we use it only in that form. Read instead as an identity of
analytic functions of the fine field $U$ near $\mathbf{1}$, in local coordinates $U=\exp i\eta H$
(the extraction being performed on $E^{(j)}(X,\exp i\eta H)$ by Taylor expansion in $H$ and the
Ward--Takahashi identities \cite[(4.9)]{B-CMP109}), it would admit bare test fields, and
neither the torus-versus-whole-lattice comparison of the kernels nor the row-sum identity
\eqref{8.5:eq-row-sum-identity-a} would be needed; that reading is not used here.

Restricting to minimiser directions has a price. It requires four inputs: the kernel decay bound,
the comparison of torus and whole-lattice kernels (both in \eqref{8.5:eq-kernel-decay-torus-a}), the
row-sum identity \eqref{8.5:eq-row-sum-identity-a}, and the regularity of the background map at the
flat configuration.

\emph{Scalar reduction.} Write $\{E^{(j)}_X\}_{j\le k_0,X}$ for the $E$-family and
$\{V^{(j)}_X\}_{j\le k_0,X}$ for the residual family. Each family is covariant under translations
by $\Gamma$ and under $W_4$. The first property is the $\Gamma$-covariance of the localised terms. The
second is their hypercubic covariance, reflections included; termwise for the $X$-indexed families
it is used as a hypothesis, and the step is marked $(\ast)$. For the frozen sourced family
at constant shift $\zeta$ both covariance properties are marked $(\ast)$.

We apply the lemma that the $\Gamma$-class sum of the constant-flux forms of a $\Gamma$- and
$W_4$-covariant family is a $W_4$-scalar, that is, a multiple of $\langle\cdot,\cdot\rangle$. We
apply it to the two families separately. We also use the lemma that the constant-flux form of $A$ on
a period block is a $W_4$-scalar. Together these show that the three per-block class sums
$\mathfrak{q}_M(\sum_j\sum_XE^{(j)}_X)$, $q_{\sum_j\beta_jA|_{\mathfrak{B}}}$ and
$\mathfrak{q}_M(\sum_j\sum_XV^{(j)}_X)$ are $W_4$-scalars. Hence
\begin{equation}\label{8.5:eq-torus-test-1}
\Delta:=\mathfrak{q}_M\Bigl(\sum_j\sum_XE^{(j)}_X\Bigr)-q_{\sum_j\beta_jA|_{\mathfrak{B}}}
-\mathfrak{q}_M\Bigl(\sum_j\sum_XV^{(j)}_X\Bigr)=\delta\,\langle\cdot,\cdot\rangle
\end{equation}
for one real number $\delta$. To show $\delta=0$ it suffices to show $\Delta[c,c]=0$ for one $c$ with
$\langle c,c\rangle\ne0$. We test one $c$ that is a multiple of a single basis plane $e_P$; we call
such $c$ decomposable.

\emph{The test field: a coarse wave pushed through $M_{k_0}$.} Let $\mathcal{N}$ be a multiple of $M$.
Let $T_1^{(k_0)}$ be the level-$k_0$ unit lattice on a torus of side $\mathcal{N}$ cells, among the
tori along which the corresponding display holds. Put
\begin{equation*}
v:=\tau^1\otimes e_2\,h(x_1),
\end{equation*}
where $\tau^1$ is one fixed colour direction, $x_1$ is the first coordinate and $e_\mu$ is the unit
$1$-cochain in direction $\mu$, so that $e_1\wedge e_2=e_{12}$. One colour direction suffices, by the
statement on constant-flux forms of gauge-invariant localised terms,
\eqref{8.5:eq-constant-flux-forms-a}. Here $h$ is the period-$\mathcal{N}$ triangle wave of slope
$\pm1$ per cell.

The blocking factor $L$ is odd, so $M$ is odd and the multiple $\mathcal{N}$ of $M$ may be odd. If
$\mathcal{N}$ is odd we insert one flat layer, on which $h'=0$, so that the heights run
$0\to\lfloor\mathcal{N}/2\rfloor\to\lfloor\mathcal{N}/2\rfloor\to0$; this creates a third defect
hyperplane. Let
\begin{equation*}
\mathfrak{K}:=\{x_1=0\}\cup\{x_1=\lfloor\mathcal{N}/2\rfloor\},
\end{equation*}
together with the flat layer when $\mathcal{N}$ is odd.
Then $dv=h'(x_1)\,e_1\wedge e_2$, and on every coarse plaquette off $\mathfrak{K}$ the function $h'$
is constant, equal to $+1$ or $-1$. Hence $dv=\pm c$ there, with $c:=e_{12}$. Put $B:=M_{k_0}v$.
By the commutation $d\,M_{k_0}=\bar{M}_{k_0}\,d$ of the linearised minimiser with the coboundary, and
the kernel representation of its action $\bar{M}_{k_0}$ on fluxes, $dB=\bar{M}_{k_0}(dv)$, that is,
\begin{equation*}
(dB)_p=\sum_{\mathfrak{P}}K_{\mathrm{tor}}(p,\mathfrak{P})\,dv(\mathfrak{P})
\qquad\text{for every bare plaquette }p,
\end{equation*}
where $K_{\mathrm{tor}}$ is the kernel of $\bar{M}_{k_0}$ on the torus and $K$ its whole-lattice
counterpart.

\emph{Off the defect set, $dB$ is locally constant up to exponentially small errors.} Let $p$ be a
bare plaquette at $\ell^\infty$-distance $t\ge1$ cells from $\mathfrak{K}$, in a region where
$h'\equiv\sigma\in\{\pm1\}$. Regard $dv$ as a periodic flux on the whole level-$k_0$ lattice. The
row-sum identity \eqref{8.5:eq-row-sum-identity-a} gives
\begin{equation*}
\sum_{\mathfrak{P}}K(p,\mathfrak{P})\,\sigma c_{\pi(\mathfrak{P})}=\eta^2\sigma c_{\pi(p)}.
\end{equation*}
Therefore
\begin{equation}\label{8.5:eq-torus-test-2}
(dB)_p-\eta^2\sigma c_{\pi(p)}
=\sum_{\mathfrak{P}}K(p,\mathfrak{P})\bigl(dv(\mathfrak{P})-\sigma c_{\pi(\mathfrak{P})}\bigr)
+\sum_{\mathfrak{P}}(K_{\mathrm{tor}}-K)(p,\mathfrak{P})\,dv(\mathfrak{P}).
\end{equation}
The use of \eqref{8.5:eq-row-sum-identity-a} is marked $(\ast)$.

In the first sum, the summand vanishes when $\operatorname{dist}(p,\mathfrak{P})<t$, because $dv=\sigma c$
on such plaquettes. Elsewhere $|dv(\mathfrak{P})-\sigma c_{\pi(\mathfrak{P})}|\le2$. The second sum is
controlled by the torus-versus-whole-lattice comparison in \eqref{8.5:eq-kernel-decay-torus-a};
that use is marked $(\ast)$. With the kernel decay bound of
\eqref{8.5:eq-kernel-decay-torus-a} this gives
\begin{align}
\bigl|(dB)_p-\eta^2\sigma c_{\pi(p)}\bigr|
&\le2C_\Gamma\eta^2\sum_{\operatorname{dist}(p,\mathfrak{P})\ge t}
e^{-\delta_\Gamma\operatorname{dist}(p,\mathfrak{P})/\ell_\Gamma}
+C_\Gamma\eta^2e^{-\delta_\Gamma\mathcal{N}/(2\ell_\Gamma)}\notag\\
&\le C_\Gamma'\eta^2e^{-\delta_\Gamma t/(2\ell_\Gamma)}.\label{8.5:eq-torus-test-3}
\end{align}
In the last step the shell counts, at most $C(1+n)^3$ plaquettes at distance $n$, are absorbed by
halving the rate, and $t\le\mathcal{N}/2$ is used. On each region where $h'\equiv\sigma$ put
$\xi_p:=\eta^{-2}(dB)_p-\sigma c_{\pi(p)}$; by \eqref{8.5:eq-torus-test-3},
$|\xi_p|\le C_\Gamma'e^{-\delta_\Gamma t/(2\ell_\Gamma)}$ at plaquettes $p$ at distance $t\ge1$
from $\mathfrak{K}$.

Everywhere we also have $|dB|_\infty\le C_\Gamma''\eta^2$. Indeed, the row sum of $|K_{\mathrm{tor}}|$
is at most the row sum of $|K|$ plus the torus-versus-whole-lattice difference. Thus, off
$\mathfrak{K}$, $dB=\pm\eta^2c$ up to the error \eqref{8.5:eq-torus-test-3}; this rests on the two
kernel statements of \eqref{8.5:eq-kernel-decay-torus-a} and on \eqref{8.5:eq-row-sum-identity-a}.

\emph{Terms with non-wrapping $\hat{X}$, $r(X)<\mathcal{N}/2$.} Put $t(X):=\operatorname{dist}(\hat{X},\mathfrak{K})$.

\emph{The case $t(X)\ge1$.} Then $\hat{X}$ lies in one region, and on $\hat{X}$
\begin{align*}
&dB=\eta^2(\sigma_Xc+\xi_X),\qquad \sigma_X=\pm1,\\
&|\xi_X|_\infty\le\min\bigl\{C_\Gamma'e^{-\delta_\Gamma t(X)/(2\ell_\Gamma)},\,C_\Gamma''+1\bigr\},
\end{align*}
where $\sigma_X$ is the value of $h'$ on that region and $\xi_X:=\xi|_{\hat{X}}$,
by \eqref{8.5:eq-torus-test-3} and the bound on $|dB|_\infty$.

On $\hat{X}$, $B$ differs from $\eta^2B^{\sigma_Xc,y}$, the affine field of constant flux
$\sigma_Xc$ anchored at a point $y$ of $\hat{X}$, by $d$ of a quadratic polynomial plus the part
carrying $\xi_X$. By the statement on constant-flux forms of gauge-invariant localised terms,
\eqref{8.5:eq-constant-flux-forms-a}, the Hessian of a localised term along the
gradient part vanishes. The same statement defines $q_X$ on constant fluxes and the flux form
$Q_X$, and bounds $Q_X$ by $\|\cdot\|_X$ and $\mathfrak{m}^{\mathrm{reg}}_X$. Using bilinearity, and
$\sigma_X^2=1$ to remove the sign, we get
\begin{align}
\bigl|Q_X[dB,dB]-\eta^4q_X[c,c]\bigr|
&=\eta^4\bigl|2Q_X[\sigma_Xc,\xi_X]+Q_X[\xi_X,\xi_X]\bigr|\notag\\
&\le(C_\Gamma''+3)\min\bigl\{C_\Gamma'e^{-\delta_\Gamma t(X)/(2\ell_\Gamma)},\,C_\Gamma''+1\bigr\}
\eta^4\mathfrak{m}^{\mathrm{reg}}_X.\label{8.5:eq-torus-test-4}
\end{align}
In the shell sum below, for $X$ with $t(X)\ge1$ anchored at distance $u<8$ from $\mathfrak{K}$, we
use the crude member $(C_\Gamma''+3)(C_\Gamma''+1)$.

\emph{Admissibility of the directions.} The inequality \eqref{8.5:eq-torus-test-4} uses that $c$
and $\xi_X$ are admissible directions for the majorant in the norm $\|\cdot\|_X$ of
\eqref{8.5:eq-constant-flux-forms-a}. This admissibility statement is marked $(\ast)$.

The constant $c$ is admissible because it is constant. For $\xi_X$, note that
$\xi_X=\eta^{-2}dB|_{\hat{X}}-\sigma_Xc$ with $dB=\bar{M}_{k_0}(dv)\in\operatorname{ran}\bar{M}_{k_0}$.
Apply the derivative form of the kernel decay bound, with constants $C_{\mathrm{reg}}$, $\beta_0$ and
$a$ ($\beta_0$ the regularity exponent of that bound, not a value of $\beta_j$), to
$\varphi:=\eta^{-2}dB$ with the fine constant $\sigma_Xc$. Write $\operatorname{reg}_X$ for the
regularity seminorm on $\hat{X}$ that enters $\|\cdot\|_X$; it vanishes on constants. This gives
\begin{align*}
\operatorname{reg}_X(\xi_X)&=\operatorname{reg}_X(\varphi|_{\hat{X}})\\
&\le C_{\mathrm{reg}}(1-\beta_0)^{-1}\max_{p\subset\hat{X}}\sum_{\square'}
e^{-a\operatorname{dist}(\square',p)}
\sup_{p'\subset\square'}\bigl|\eta^{-2}(dB)_{p'}-\sigma_Xc_{\pi(p')}\bigr|,
\end{align*}
the sum running over cells $\square'$ and the supremum over the bare plaquettes $p'$ in $\square'$.
We bound the summand in two zones.
\begin{itemize}
\item On the cells $\square'$ of the region containing $\hat{X}$ with
$\operatorname{dist}(\square',\hat{X})<t(X)$, the summand is at most
$\sup_{p'\subset\square'}|\xi_{p'}|$. The cell $\square'$ lies at distance at least
$t(X)-\operatorname{dist}(\square',\hat{X})$ from
$\mathfrak{K}$, so by \eqref{8.5:eq-torus-test-3} at the cell $\square'$ this is at most
$C_\Gamma'e^{-\delta_\Gamma(t(X)-\operatorname{dist}(\square',\hat{X}))/(2\ell_\Gamma)}$.
\item Beyond that distance the summand is at most $C_\Gamma''+2$.
\end{itemize}
Hence the weighted sum is at most $C_\Gamma'''e^{-\min\{a,\delta_\Gamma/(2\ell_\Gamma)\}t(X)}$ for a
constant $C_\Gamma'''$ depending on $C_\Gamma'$, $C_\Gamma''$, $a$, $\delta_\Gamma$, $\ell_\Gamma$ only.
This is
the same exponential form as the bound on $|\xi_X|_\infty$. Combining the two bounds, and using
$|c|=1$,
\begin{equation*}
\|\xi_X\|_X\le\bigl(C_\Gamma'+C_{\mathrm{reg}}(1-\beta_0)^{-1}C_\Gamma'''\bigr)
e^{-\min\{a,\delta_\Gamma/(2\ell_\Gamma)\}t(X)}.
\end{equation*}
Therefore \eqref{8.5:eq-torus-test-4} and the displays below hold after three replacements:
$C_\Gamma'$ by $C_\Gamma'+C_{\mathrm{reg}}(1-\beta_0)^{-1}C_\Gamma'''$, $C_\Gamma''$ likewise, and
$\delta_\Gamma/(2\ell_\Gamma)$ by $\min\{a,\delta_\Gamma/(2\ell_\Gamma)\}$. These change only the
constants and the rate of the kink sum. The conclusion $\Delta[c,c]=O(M^4/\mathcal{N})+o(1)$ below is
unchanged. Without the regularity input, the bound by Cauchy's estimate on two radii would leave a
member of order $\eta\mathcal{N}^4$ in the kink sum.

\emph{The case $t(X)=0$.} Then
\begin{equation*}
\bigl|Q_X[dB,dB]-\eta^4q_X[c,c]\bigr|\le\bigl((C_\Gamma'')^2+1\bigr)\eta^4\mathfrak{m}^{\mathrm{reg}}_X,
\end{equation*}
and $\mathfrak{m}^{\mathrm{reg}}_X\le\hat{\mathfrak{m}}_X$ by the statement on constant-flux forms of
gauge-invariant localised terms, \eqref{8.5:eq-constant-flux-forms-a}.

\emph{The shell sum.} There are at most $6\mathcal{N}^3$ anchors at $\ell^\infty$-distance $u$ cells
from $\mathfrak{K}$: three hyperplanes, each with two sides. Let $X$ be anchored at such an anchor.
Either $r(X)\ge u/4$, or $r(X)<u/4$; in the latter case
\begin{equation*}
t(X)\ge u-2(r(X)+1)\ge u/2-2\ge u/4\qquad(u\ge8).
\end{equation*}

Summing shell by shell with the radius tails \eqref{8.5:eq-radius-tails-a}, and bounding the mass of
the members anchored at one cell by $\bar{\mathfrak{a}}^0$, we obtain
\begin{align}
&\sum_{X\ \text{non-wrapping}}\bigl|Q_X[dB,dB]-\eta^4q_X[c,c]\bigr|\notag\\
&\quad\le\eta^4\cdot6\mathcal{N}^3\sum_{u\ge0}\Bigl[\bigl((C_\Gamma'')^2+1\bigr)
\sup_{z_0}\sum_{c(X)=z_0,\ r(X)\ge u/4}\hat{\mathfrak{m}}^0_X\notag\\
&\qquad+(C_\Gamma''+3)C_\Gamma'e^{-\delta_\Gamma u/(8\ell_\Gamma)}\bar{\mathfrak{a}}^0
+[u<8]\,(C_\Gamma''+3)(C_\Gamma''+1)\bar{\mathfrak{a}}^0\Bigr].
\label{8.5:eq-torus-test-5}
\end{align}
Here $z_0$ runs over anchor cells, $[u<8]$ equals $1$ if $u<8$ and $0$ otherwise, and in the second
member $e^{-\delta_\Gamma t(X)/(2\ell_\Gamma)}\le
e^{-\delta_\Gamma u/(8\ell_\Gamma)}$ because $t(X)\ge u/4$. By \eqref{8.5:eq-radius-tails-a},
\begin{equation*}
\sup_{z_0}\sum_{c(X)=z_0,\ r(X)\ge u/4}\hat{\mathfrak{m}}^0_X
\le[u<4]\,\bar{\mathfrak{a}}^0+C_we^{-\hat\kappa(u/4-1)/(2M)}\bar{\mathfrak{a}}^0,
\end{equation*}
where $C_w$ is the constant of \eqref{8.5:eq-radius-tails-a}; this is summable in $u$. In all, the
right side of \eqref{8.5:eq-torus-test-5} is at most
\begin{equation*}
\eta^4\mathcal{N}^3\,C_{\mathrm{kink}}(1+\rho_\sharp)
\bar{\mathfrak{a}}^0=o(\mathcal{N}^4)\eta^4,
\end{equation*}
where $C_{\mathrm{kink}}$ depends on $\kappa$, $M$ and the kernel and regularity constants only.

This holds for two families.
\begin{itemize}
\item For the residual family the majorants are the unweighted ones $\hat{\mathfrak{m}}^0_X$, at
$\zeta=0$, resp.\ at constant shift $\zeta$ for the frozen family. They come from the size bounds
\cite[(0.29)]{B-CMP109}, resp.\ from their analogue at constant shift, marked $(\ast)$ for the
frozen family, used without weights; no route weight is involved.
\item For the pre-extraction $E$-family $\{E^{(j)}_X\}$ the size is $E_0e^{-\kappa d_j(X)}$, with the
constant $E_0$ of \cite[(1.18)]{B-CMP109}, with no power of $L^j\eta$. Here the argument runs with
this family's own
per-anchor mass, which is independent of $\mathcal{N}$, and its own shell tails of the form
$e^{-\hat\kappa d_j}$. These follow from \cite[(1.18)]{B-CMP109}, the count
\cite[(0.26)]{B-CMP109} and the same Cauchy conversion, at fixed $(j,k_0)$. The result is summed over
the finitely many $j\le k_0$.
\end{itemize}

\emph{The term $A$.} Write $A_\eta$ for the Wilson action on the bare lattice of spacing $\eta$
(not to be confused with the effective action $A_{k_0}$). Its Hessian satisfies
\begin{equation}\label{8.5:eq-torus-test-6}
\operatorname{Hess}A_\eta[dB,dB]=\eta^4\Bigl[\mathcal{N}^4\cdot\frac{\mathfrak{n}}{6}\langle c,c\rangle
+O(\mathcal{N}^3C_{\mathrm{kink}}\mathfrak{n})\Bigr].
\end{equation}
The Hessian of $A$ per cell on the constant flux $c$ is $(\ell/\varepsilon)^4\langle c,c\rangle
=(\mathfrak{n}/6)\langle c,c\rangle$, where $\mathfrak{n}=6(\ell/\varepsilon)^4$ is its value on
$c\equiv(1,\dots,1)$. This normalisation is immaterial for the vanishing statement. The defect term
comes from the flat layer and from the neighbourhoods of the kinks, where
$|c+\xi|^2\neq|c|^2$. It is bounded by the same shell sum and is of order $\mathcal{N}^3$.

\emph{Wrapping domains.} Consider $X$ with $r(X)\ge\mathcal{N}/2$, so that
$d_j(X)\ge\mathcal{N}/(2M)-1$. These lie outside the statement on constant-flux forms,
\eqref{8.5:eq-constant-flux-forms-a}: their Hessian reads the holonomies of $B$ around the torus, and
$\sup|B|$, for which only the bound $\sup|B|\le C\eta\mathcal{N}$ with $C$ independent of
$\mathcal{N}$ is available, is not small. We bound
them by Cauchy's estimate in the bond
field at the absolute radius $\mathfrak{r}_T$ of the analyticity domain of \cite[(1.18)]{B-CMP109},
resp.\ \cite[(0.29)]{B-CMP109}:
\begin{equation*}
|\operatorname{Hess}T[B,B]|\le2\operatorname{size}(T)\bigl(\sup|B|/\mathfrak{r}_T\bigr)^2,
\end{equation*}
where $\operatorname{size}(T)$ is the bound on $|T|$ on that domain given there.

The number of $X\in\mathcal{D}_j$ with $d_j(X)=n$ meeting the torus is at most
$\mathcal{N}^4L^{4(k_0-j)}C_{\mathsf{a}}^n$, where $C_{\mathsf{a}}$ is the constant of the count
\cite[(0.26)]{B-CMP109}. The wrapping total is therefore at most
\begin{equation*}
C(j,k_0,M)\,\mathcal{N}^6\sum_{n\ge\mathcal{N}/(2M)-1}\bigl(C_{\mathsf{a}}e^{-\kappa}\bigr)^n=o(1)
\qquad(\mathcal{N}\to\infty)
\end{equation*}
at fixed $(j,k_0,M,\kappa)$. Here $\kappa>\log C_{\mathsf{a}}$ holds inside the lower bound on $\kappa$
in the order of choice (Notation~\ref{2.3:not-stages-middle}). This applies to the $E$-family and the
residual family alike. For $A$ there is no wrapping issue.

\emph{Torus versus whole-lattice class sums.} By the torus-coincidence clause of the $\Gamma$-covariance
of the localised terms, the torus family's terms with non-wrapping $\hat{X}$ coincide with the
whole-lattice family's terms up to $O(e^{-c_{\mathrm{tor}}\mathcal{N}/M})$ per term, uniformly
summably, with $c_{\mathrm{tor}}>0$ the constant of that clause. The total effect on the class sum is
at most $(\mathcal{N}/M)^4\cdot Ce^{-c_{\mathrm{tor}}\mathcal{N}/M}\bar{\mathfrak{a}}^0\to0$.
The torus-extracted coefficient $\beta^{(\mathcal{N})}_j$ replaces $\beta_j$ at cost
\begin{equation*}
(\beta^{(\mathcal{N})}_j-\beta_j)\operatorname{Hess}A_\eta[dB,dB]=o(\mathcal{N}^4)\eta^4.
\end{equation*}

Every $\Gamma$-class of radius $<\mathcal{N}/2$ has exactly $(\mathcal{N}/M)^4$ members anchored in
the torus. By \eqref{8.5:eq-radius-tails-a}, the classes of radius $\ge\mathcal{N}/2$ have total mass
\begin{equation*}
\le M^4\sup_{z_0}\sum_{c(X)=z_0,\ r(X)\ge\mathcal{N}/2}\hat{\mathfrak{m}}^0_X
\le M^4C_we^{-\hat\kappa(\mathcal{N}/2-1)/(2M)}\bar{\mathfrak{a}}^0;
\end{equation*}
for the pre-extraction $E$-family its own shell tail is used instead. Hence, for each of the two
$X$-indexed families, with $\mathfrak{q}_M$ standing for the class sum of the family in question,
\begin{align}
\sum_{X\ \text{non-wrapping, torus}}q_X
&=\Bigl(\frac{\mathcal{N}}{M}\Bigr)^4\Bigl[\mathfrak{q}_M
+O\bigl(M^4C_we^{\hat\kappa/(2M)}e^{-\hat\kappa\mathcal{N}/(4M)}\bar{\mathfrak{a}}^0\bigr)\notag\\
&\qquad+O\bigl(e^{-c_{\mathrm{tor}}\mathcal{N}/M}\bar{\mathfrak{a}}^0\bigr)\Bigr].
\label{8.5:eq-torus-test-7}
\end{align}

\emph{Conclusion of the torus test.} Insert \eqref{8.5:eq-torus-test-4}--\eqref{8.5:eq-torus-test-7}
and the wrapping bound into the corresponding display with $B=M_{k_0}v$, and divide by
$\eta^4$; here $\eta=L^{-k_0}$ is fixed at fixed $k_0$. The identity reads
\begin{equation*}
\Bigl(\frac{\mathcal{N}}{M}\Bigr)^4\Delta[c,c]=O(\mathcal{N}^3)+o(\mathcal{N}^4),
\end{equation*}
with constants depending on $k_0$, $M$, $\kappa$ and the kernel and regularity constants only. That is,
$\Delta[c,c]=O(M^4/\mathcal{N})+o(1)\to0$ as $\mathcal{N}\to\infty$. Since $\Delta[c,c]=\delta\langle c,c\rangle$
and $\langle c,c\rangle=1$, we get $\delta=0$, so
\begin{equation}\label{8.5:eq-torus-test-8}
\mathfrak{q}_M\Bigl(\sum_j\sum_XV^{(j)}_X\Bigr)
=\sum_j\Bigl[\mathfrak{q}_M\Bigl(\sum_XE^{(j)}_X\Bigr)-q_{\beta_jA|_{\mathfrak{B}}}\Bigr].
\end{equation}
Only this sum over $j$ is used; the per-$j$ form is not used.

Only decomposable $c$ can be tested in this way. A coarse field with $dv$ piecewise constant, equal
to $\pm c$ on slabs, exists only when $c$ is a multiple of a single $e_P$. This is exactly the gap
exhibited by the quadratic part of the lattice density $\sum F\wedge F$ of the curvature, and the
scalar reduction, in which the reflections in $W_4$ are used, closes it.

\emph{Identification of the levels.} Several statements identify the constant flux at the coarse
level of \cite[(0.28)]{B-CMP109} with a constant bare flux. They are: the constant-flux property of
the linearised minimiser, the row-sum identity \eqref{8.5:eq-row-sum-identity-a}, the two kernel
statements of \eqref{8.5:eq-kernel-decay-torus-a}, the regularity of the background map at the flat
configuration with differential the linearised minimiser, and the lemma factorising the linearised
minimiser through an intermediate level. The factor is $\eta^2$ (at level $j$, $\eta_j^2$), exactly,
by \eqref{8.5:eq-row-sum-identity-a}.

There is a further point. For $j<k_0$ the terms $E^{(j)}_X$ in the corresponding display
are the carried terms composed with \cite[(1.15)]{B-CMP109}, functions of $U_{k_0}(V)$. The kernel
$\Pi^{(j)}$ used below is built from the terms $E^{(j)}(\cdot,U_j(\cdot))$ at the level $j$ at which
they are created. Their constant-flux forms agree by the linear nesting of the background maps
alone. Let $\mathcal{M}^j$ be the block-averaging map to level $j$ entering the composition of
\cite[(1.15)]{B-CMP109}, with differential $Q_j$ at $\mathbf{1}$. Then
\begin{equation}\label{8.5:eq-torus-test-9}
D(U_j\circ\mathcal{M}^j)(\mathbf{1})\,M_{k_0}=M_jQ_jM_{k_0}=M_{k_0}.
\end{equation}
This uses three facts. The first is the lemma that the linearised block average inverts the
linearised minimiser, $Q\mathsf{M}=\mathrm{id}$. The second is the factorisation
$\mathsf{M}=\mathsf{M}_1\mathsf{M}_2$ of the lemma factorising the linearised minimiser through an
intermediate level. The third is the statement that the
background map is $C^2$ at the flat configuration with $DU_j(\mathbf{1})=M_j$, read at every
$j\le k_0$ and not only at $k_0$. This last statement is given in \cite{B-CMP95} at $k_0$; its $C^2$
clause and its form at every level are marked $(\ast)$.

Since $DT(\mathbf{1})=0$ for every gauge-invariant localised term $T$, by the statement on
constant-flux forms of gauge-invariant localised terms, \eqref{8.5:eq-constant-flux-forms-a}, the
term of the chain rule containing $D^2(U_j\circ\mathcal{M}^j)(\mathbf{1})$ vanishes. So the Hessian
of the carried composite along
$M_{k_0}v$ equals the Hessian at the level of creation along $M_j(Q_jM_{k_0}v)$. The row-sum
identity \eqref{8.5:eq-row-sum-identity-a} at the intermediate level $j$, marked $(\ast)$, makes
$Q_jM_{k_0}v$ affine off the kinks. The function $J$ of \cite[(1.15)]{B-CMP109} is a gauge-covariant
function of $U$. The composite is therefore gauge invariant, the constant-flux statement applies to
it, and no separate statement on $J$ is needed.

\emph{The class sum of the $E$-family is that of $\beta_jA$, for each $j$.} We use the point
attribution $E^{(j)}(X,U)=\sum_zE^{(j)}(X,U,z)$ of \cite[(3.50)]{B-CMP119}, of which only the gauge
invariance, analyticity and $\Gamma$-covariant attribution of the pieces $E_z$ are used. The
kernel $\Pi^{(j)}_{\mu\nu}(x,y,z)$ of \cite[(3.57)]{B-CMP119} is the second derivative of
$E_z\circ U_j(\exp i\,\cdot)$ in the coarse field.

Let $\tilde B^{\tilde c,y_0}$ be the coarse affine field of constant flux $\tilde c$ anchored at
$y_0$. By the row-sum identity \eqref{8.5:eq-row-sum-identity-a} at level $j$ $(\ast)$ and
$DU_j(\mathbf{1})=M_j$ $(\ast)$, the field $M_j\tilde B^{\tilde c,y_0}$ is affine on the fine
lattice, with flux $\eta_j^2\tilde c$, up to an exact form $d\lambda$, with $\lambda$ a function on
the fine sites. Hence
\begin{align}
\eta_j^4\,q^{\mathrm{fine}}_{E_z}[\tilde c,\tilde c]
&=\operatorname{Hess}\bigl(E_z\circ U_j\circ\exp i\,\cdot\bigr)[\tilde B^{\tilde c,y_0},\tilde B^{\tilde c,y_0}]
\notag\\
&=\operatorname{Hess}\bigl(E_z\circ U_j\circ\exp i\,\cdot\bigr)[\tilde B^{\tilde c,z},\tilde B^{\tilde c,z}]
\notag\\
&=\frac14\sum_{\mu,\nu,\mu',\nu'}\tilde c_{\mu'\mu}\tilde c_{\nu'\nu}\sum_{x,y}
\Pi^{(j)}_{\mu\nu}(x,y,z)(x-z)_{\mu'}(y-z)_{\nu'}\notag\\
&=\frac14\sum_{\mu,\nu,\mu',\nu'}\tilde c_{\mu'\mu}\tilde c_{\nu'\nu}\Pi^{(j)}_{\mu\nu,\mu'\nu'}.
\label{8.5:eq-torus-test-10}
\end{align}
Here $q^{\mathrm{fine}}_{E_z}=q_{E_z}$ denotes the constant-flux form of $E_z$ on constant fine
fluxes, and the indices are those of \cite[(3.57)]{B-CMP119}, with the two indices that it pairs
with the coordinates $x-z$ and $y-z$ written $\mu'$ and $\nu'$.
The second equality holds because $\tilde B^{\tilde c,y_0}-\tilde B^{\tilde c,z}$ is $d$ of a linear
function, so gauge invariance applies. This is the step from the first Ward--Takahashi identity
\cite[(4.15)]{B-CMP109} used in \cite[(3.60)]{B-CMP119}. The last member is the number of
\cite[(3.57)]{B-CMP119}, which is independent of $z$.

On conventions: \cite[(3.57)]{B-CMP119} carries the factor $\frac12$ and the restricted indices
$\mu'<\mu$, $\nu'<\nu$, while the factor $\frac14$ in \eqref{8.5:eq-torus-test-10} runs over
unrestricted indices. The factor $\eta_j^4$ appears on both sides and is immaterial.

The normalisation is not immaterial. The vanishing is an exact equality of two numbers, and an
unreconciled factor would leave a $W_4$-scalar residue of the size of the main term. Such a factor
could come from the $\frac12$ on restricted indices against the $\frac14$ on unrestricted ones, from
the sign in \cite[(1.22)]{B-CMP109}, or from the normalisation of $A$. The normalisation is carried
by Ba{\l}aban's definition of $\beta$ as the coefficient of the subtracted $A$-term. That
\cite[(3.56), (3.61)]{B-CMP119} make the class-summed second moment equal to $\beta_j$ times the
quadratic form of $A$, in the normalisation of \cite{B-CMP109}, is used here as stated in
\cite{B-CMP119}.

As a consistency check, which is not used in the argument: gauge invariance under the quadratic
gauge functions $\lambda(x)=\frac12Sx\cdot x$, $S$ a symmetric $4\times4$ matrix, gives
$\Pi_{\mu\nu,\mu'\nu'}=-\Pi_{\mu'\nu,\mu\nu'}$.
Hence the quarter-sum over unrestricted indices equals the sum over restricted indices, which equals
the Hessian, in agreement with the factor $\frac12$ of \cite[(3.57)]{B-CMP119}.

By \cite[(3.58)--(3.61)]{B-CMP119}, reflections leave only $\mu=\nu$, $\mu'=\nu'$, and
permutations make the surviving entries equal; this is the content of the lemma on $W_4$-invariance
used above. So \eqref{8.5:eq-torus-test-10} equals
\begin{equation*}
\beta_j\mathfrak{n}^{\mathrm{coarse}}_j\langle\tilde c,\tilde c\rangle
=\eta_j^4\beta_j\mathfrak{n}_j\langle\tilde c,\tilde c\rangle.
\end{equation*}
Here $\mathfrak{n}_j\langle c,c\rangle$ is the Hessian of $A$ per
level-$j$ site along a constant fine flux $c$; it is a $W_4$-scalar by the same lemma.
The number $\mathfrak{n}^{\mathrm{coarse}}_j$ is its coarse-normalised counterpart, with the same
factor $\eta_j^4$ on both sides.

The text following \cite[(3.61)]{B-CMP119} states that this identity may serve as a definition of
the $\beta$-function and that it agrees with \cite[(1.22)]{B-CMP109}. The identification of the
coefficient $\beta_j$ of
\cite[(0.28)]{B-CMP109} with that of \cite[(5.16)]{B-CMP109} and of \cite[(1.22)]{B-CMP109} is used
as stated in \cite{B-CMP109}.

A period block $\mathfrak{B}$ contains $(ML^{k_0-j})^4$ level-$j$ sites. Summing $z$ over
$\mathfrak{B}$,
\begin{equation*}
\sum_{z\in\mathfrak{B}}q_{E_z}=(ML^{k_0-j})^4\beta_j\mathfrak{n}_j\langle\cdot,\cdot\rangle
=q_{\beta_jA|_{\mathfrak{B}}}.
\end{equation*}
The block sum of the point-attributed family equals the $\Gamma$-class sum of the $X$-indexed family,
for three reasons. Each class of $X$ meets the attribution exactly through the pieces
$E^{(j)}(X,U,z)$, $\Gamma$-periodically: termwise $\Gamma$-covariance makes $q_X$ unchanged under
translation by $\Gamma$. The block $\mathfrak{B}$ is a fundamental domain of $\Gamma$. And Hessians
add.

With \eqref{8.5:eq-torus-test-8}, this gives
\begin{equation*}
\mathfrak{q}_M\Bigl(\sum_j\sum_XV^{(j)}_X\Bigr)
=\sum_j\Bigl[\mathfrak{q}_M\Bigl(\sum_XE^{(j)}_X\Bigr)-q_{\beta_jA|_{\mathfrak{B}}}\Bigr]=0
\end{equation*}
for the terms $E$ of \cite{B-CMP109} and of \cite{B-CMP119}.

For the frozen sourced family of the correlation theorem at constant shift $\zeta$, the running-shift
statement makes $b_{k+1}(\zeta)$ the number of \cite[(3.61)]{B-CMP119} for the families from which it
is built at shift $\zeta$. The identity differentiated is then the identity of representations
\eqref{8.5:eq-representation-identity-a} $(\ast)$.

\emph{What the argument rests on.} For Ba{\l}aban's own $\beta$-extracted family, the argument above
derives the conclusion from the following statements.
\begin{itemize}
\item \cite[(0.28)]{B-CMP109}, an identity in the coarse variable on the small-field space of the
lattice $T_\eta$ of \cite[Theorem~1]{B-CMP109}, tested on minimiser directions and summed over $j$;
together with \cite[(1.20)--(1.22), (5.1), (5.16)]{B-CMP109}.
\item \cite[(3.50), (3.56)--(3.61)]{B-CMP119}, used as stated there, with the identification of the
coefficients in \cite{B-CMP109}.
\item The $\Gamma$-covariance of the localised terms, including its torus-coincidence clause
(agreement up to uniformly summable $O(e^{-c_{\mathrm{tor}}\mathcal{N}/M})$ per term).
\item The hypercubic covariance, reflections included, of the $E$-family and the residual family.
It is stated in \cite[(5.2)--(5.7)]{B-CMP109} for the summed functional and in
\cite[(3.58)]{B-CMP119} for the point-attributed family; termwise for the $X$-indexed families it is
marked $(\ast)$.
\item The kernel decay bound of \eqref{8.5:eq-kernel-decay-torus-a}, as a decay statement, and its
derivative form, with the constants $C_{\mathrm{reg}}$, $\beta_0$ and $a$ of that statement, used
for the admissibility of $\xi_X$.
\item The torus-versus-whole-lattice comparison of \eqref{8.5:eq-kernel-decay-torus-a}, marked
$(\ast)$. It rests on the finite-volume construction of the kernel in \cite{B-CMP96},
\cite{B-CMP98}, \cite{B-CMP99b} and on the infinite-volume limit in \cite{B-CMP109}.
\item The row-sum identity \eqref{8.5:eq-row-sum-identity-a} at levels $k_0$ and $j$, marked $(\ast)$.
It would follow from the constant-flux property of the linearised minimiser together with a
uniqueness statement of Liouville type for bounded solutions, or from the decay of the kernel of
$(\bar{Q}\bar{Q}^*)^{-1}$.
\item The regularity of the background map at the flat configuration, with
$DU_{k_0}(\mathbf{1})=M_{k_0}$, as stated in \cite{B-CMP95}; its $C^2$ clause and its form at every
level $j\le k_0$ are marked $(\ast)$.
\item The constant-flux property of the linearised minimiser, which is part of the description of
the linearised minimiser; its first-variation step is proved in the first part of this proof.
\item The radius tails \eqref{8.5:eq-radius-tails-a} for the unweighted masses
$\hat{\mathfrak{m}}^0_X$; for the pre-extraction $E$-family, \cite[(1.18), (0.26)]{B-CMP109} and
Cauchy's estimate at fixed $(j,k_0)$.
\item The admissibility statement of \eqref{8.5:eq-constant-flux-forms-a}, marked $(\ast)$.
\end{itemize}

For the frozen sourced family at constant shift $\zeta$, the argument is incomplete in three further
places, marked $(\ast)$ above: the constant-shift form of the correlation theorem, the identity of
representations \eqref{8.5:eq-representation-identity-a}, and the $\Gamma$- and hypercubic covariance
of that family. With these granted, the torus test and the computation of the class sum of the
$E$-family apply verbatim, with \eqref{8.5:eq-representation-identity-a} in place of
\cite[(0.28)]{B-CMP109}.

No statement of this kind is made, for either family, for the remainder part
$\operatorname{Irr}_1^R$: its terms are outside $b_{k+1}$. The second-order structure is treated
separately for $\operatorname{Irr}_1^E$ and $\operatorname{Irr}_1^R$, and that treatment does not
depend on which of the two carries the marginal, $z$-linear content of the old large-field
remainders.
\end{proof}

\begin{proposition}[Zero block mean of the marked residual form. Stated; proof incomplete at the step marked $(\ast)$]\label{8.5:prop-zero-block-mean}
Let $k_0=K-s$ be the witness level, let $\Gamma=M\mathbb{Z}^4$ be the partition lattice in cells, and
let a \emph{period block} $\mathfrak{B}$ of $\pi_{k_0}$ be a cube of $M^4$ cells of the partition
$\pi_{k_0}$. For real $\zeta$ with $|\zeta|<\zeta_*$, where $\zeta_*>0$ is the radius of
Proposition~\ref{8.5:prop-extraction-complete}, let
$\mathcal{R}^{E,\zeta}=\{V^{\zeta}_X\}_{X\in E}$ be the residual $E$-family at constant shift $\zeta$,
the union of the step-$j$ families over $j\le k_0$, and let $\mathfrak{q}_M(\mathcal{R}^{E,\zeta})$ be its
$\Gamma$-class sum of constant-flux forms, both as in Proposition~\ref{8.5:prop-extraction-complete};
here $E$ denotes the index set of the family. The residual $E$-family marked by the first source
$z_1$ at constant profile $\mathbf{1}$ is
\begin{equation}\label{8.5:eq-zero-block-mean-1}
\partial\mathcal{R}^E:=\{H_X[\mathbf{1}]\}_{X\in E}
=\bigl\{\partial_\zeta V^{\zeta}_X\big|_{\zeta=0}\bigr\}_{X\in E}.
\end{equation}
Write $q_X[\mathbf{1}]$ for the constant-flux form of $H_X[\mathbf{1}]$, and put
\begin{equation}\label{8.5:eq-zero-block-mean-2}
q^E(c)[\mathbf{1}]:=\sum_{X\in E:\,c(X)=c}q_X[\mathbf{1}],
\qquad
\mathfrak{q}_M(\partial\mathcal{R}^E):=\sum_{X\in E/\Gamma}q_X[\mathbf{1}],
\end{equation}
the second sum taken with one representative per $\Gamma$-class. The statement concerns two
settings.

\emph{First setting.} $\mathcal{R}^{E,0}$ is Ba{\l}aban's own $\beta$-extracted residual family of
\cite[(0.28)]{B-CMP109}, and $\mathcal{R}^{E,\zeta}$ is the corresponding family of Ba{\l}aban's
sequence of renormalisation steps (the run)
restarted at bare weight $x_0-\zeta$, where $x_0=g_0^{-2}$; the mark is thus a constant shift of the
bare weight of that run. The hypotheses are:
\begin{itemize}
\item[(i)] Proposition~\ref{8.5:prop-extraction-complete} for these families, whose proof is
incomplete at one step;
\item[(ii)] the first of the paired conditions fixed in the preliminaries of this section;
\item[(iii)] the second of these paired conditions;
\item[(iv)] the decomposition clause fixed in the preliminaries of this section, in the primed form
of its second part;
\item[(v)] Ba{\l}aban's minimiser $V\mapsto U_{k_0}(V)$ is of class $C^2$ at the flat configuration,
with $DU_{k_0}(\mathbf{1})v=M_{k_0}v$ for coarse directions $v$;
\item[(vi)] Lemma~\ref{8.5:lem-block-sum-scalar} (covariance makes block-summed constant-flux forms
scalar);
\item[(vii)] Lemma~\ref{8.5:lem-radius-tails} (radius tails of the majorant masses), together with the
form of the bound \cite[(0.29)]{B-CMP109} in the fine field on the complex configuration spaces
$\mathfrak{U}^c_k$, on which that lemma rests;
\item[(viii)] for $\zeta$ in a complex disc about $0$, the terms $V^{\zeta}_X$ of Ba{\l}aban's run
restarted at bare weight $x_0-\zeta$ are analytic in $\zeta$, and their index sets do not depend
on $\zeta$.
\end{itemize}
The first setting is asserted for the families of \cite{B-CMP109} only.

\emph{Second setting.} $\mathcal{R}^{E,\zeta}$ is the frozen sourced family of the correlation
theorem at constant shift $\zeta$, the family that enters the two-point witness theorem as
$\operatorname{Irr}_1^E$. The statements used are: Proposition~\ref{8.5:prop-extraction-complete} for
this family, whose proof is incomplete at one step, read through the identity of representations for
the shifted frozen run \eqref{8.5:eq-representation-identity-a}, whose proof is outstanding;
the analyticity in $\zeta$ of the terms of the frozen sourced run of the correlation theorem, with
$\zeta$-independent index sets; and the statements listed in (ii)--(vii), the form of the bound
\cite[(0.29)]{B-CMP109} entering for this family through Ba{\l}aban's bounds as carried by the
correlation theorem.

In either setting:
\begin{itemize}
\item[(1)] $q^E(c)[\mathbf{1}]$ is $\Gamma$-periodic in $c$;
\item[(2)] $|q^E(c)[\mathbf{1}]|\le\mathfrak{a}^E_c$ for every cell $c$, where $\mathfrak{a}^E_c$ is
the cell-aggregate majorant mass of the family indexed by $E$;
\item[(3)] for every period block $\mathfrak{B}$ of $\pi_{k_0}$,
\begin{equation}\label{8.5:eq-zero-block-mean-a}
\sum_{c\in\mathfrak{B}}q^E(c)[\mathbf{1}]=\mathfrak{q}_M(\partial\mathcal{R}^E)=0
\qquad\text{in }\operatorname{Sym}^2(\mathbb{R}^{\mathfrak{O}}).
\end{equation}
\end{itemize}
\end{proposition}

In the first setting the proposition is proved as an implication from (i)--(viii); in the second
setting it is stated, and the proof below is incomplete at the step marked $(\ast)$.

\begin{proof}
\emph{Step 1 $(\ast)$: differentiation in $\zeta$.} We first record that the index set $E$ does not
depend on $\zeta$ near $0$ and that each $q_{V^{\zeta}_X}$ is real-analytic in $\zeta$. In the first
setting both facts are hypothesis (viii). In the second setting they are the statement that the terms
of the frozen sourced run of the correlation theorem are analytic in the constant shift $\zeta$ with
$\zeta$-independent index sets; this statement is one of the two at which the proof is incomplete.

The $\Gamma$-class sum $\mathfrak{q}_M(\mathcal{R}^{E,\zeta})$ converges uniformly on compact subsets
of the interval $(-\zeta_*,\zeta_*)$; only a neighbourhood of
$0$ is used. The uniform convergence comes from the majorant masses of
Lemma~\ref{8.5:lem-radius-tails}, which dominate the terms of the class sum. Hence
$\zeta\mapsto\mathfrak{q}_M(\mathcal{R}^{E,\zeta})$ is differentiable, with derivative at $\zeta=0$
\[
\frac{d}{d\zeta}\Big|_{\zeta=0}\mathfrak{q}_M(\mathcal{R}^{E,\zeta})
=\sum_{X\in E/\Gamma}\partial_\zeta q_{V^{\zeta}_X}\Big|_{\zeta=0}
=\sum_{X\in E/\Gamma}q_X[\mathbf{1}]
=\mathfrak{q}_M(\partial\mathcal{R}^E).
\]
The middle equality holds because the constant-flux form of a term depends linearly on the term, so
that the $\zeta$-derivative of $q_{V^{\zeta}_X}$ at $0$ is the constant-flux form of
$\partial_\zeta V^{\zeta}_X|_{\zeta=0}=H_X[\mathbf{1}]$ (see \eqref{8.5:eq-zero-block-mean-1}). The
last equality is the definition \eqref{8.5:eq-zero-block-mean-2}. The index set over which the sums
run is the same for all $\zeta$ near $0$ by the first paragraph.

Proposition~\ref{8.5:prop-extraction-complete}, whose proof is incomplete at one step, states that
$\mathfrak{q}_M(\mathcal{R}^{E,\zeta})=0$ for every real $|\zeta|<\zeta_*$. In the first setting this
is that proposition for Ba{\l}aban's run at bare weight $x_0-\zeta$, whose representation is
\cite[Theorem~1]{B-CMP109} for that run; this is hypothesis (i). In the second setting the vanishing
is read through the identity of representations for the shifted frozen run
(the identity \eqref{8.5:eq-representation-identity-a}, whose proof is outstanding), which supplies the
identity that is differentiated. This identity is the second statement at which the proof is
incomplete. A function that vanishes identically on an interval about $0$ has derivative $0$ there, so
\begin{equation}\label{8.5:eq-zero-block-mean-3}
\mathfrak{q}_M(\partial\mathcal{R}^E)=0.
\end{equation}
Granted the two statements named in this step, the remaining steps are complete.

\emph{Step 2: covariance.} At constant profile the source is invariant under translations and under
$W_4$, so for each $\zeta$ the family $\mathcal{R}^{E,\zeta}$ is $\Gamma$- and $W_4$-covariant. The
index set does not depend on $\zeta$, so differentiating the covariance relations at $\zeta=0$ shows
that the family $\partial\mathcal{R}^E$ is $\Gamma$- and $W_4$-covariant as well. Hence
Lemma~\ref{8.5:lem-block-sum-scalar} applies to $\mathcal{R}^{E,\zeta}$ and to $\partial\mathcal{R}^E$.
Its part (a), applied to $\partial\mathcal{R}^E$, gives the $\Gamma$-periodicity of
$q^E(c)[\mathbf{1}]$ in $c$, which is assertion (1). Its part (b) shows that
$\sum_{c\in\mathfrak{B}}q^E(c)[\mathbf{1}]$ does not depend on the period block $\mathfrak{B}$ and
equals $\mathfrak{q}_M(\partial\mathcal{R}^E)$. This is the first equality in
\eqref{8.5:eq-zero-block-mean-a}.

Part (b), applied to $\mathcal{R}^{E,\zeta}$, gives
$\mathfrak{q}_M(\mathcal{R}^{E,\zeta})=\delta(\zeta)\langle\cdot,\cdot\rangle$ with
$\delta(\zeta):=\delta(\mathcal{R}^{E,\zeta})$. By Step~1 the left side is differentiable at
$0$, so $\delta$ is, and
\[
\mathfrak{q}_M(\partial\mathcal{R}^E)=\delta'(0)\langle\cdot,\cdot\rangle .
\]
Thus the anisotropic part of the block sum vanishes by part (b) of
Lemma~\ref{8.5:lem-block-sum-scalar} alone. Proposition~\ref{8.5:prop-extraction-complete}, whose
proof is incomplete at one step, is needed, through \eqref{8.5:eq-zero-block-mean-3}, exactly for the
scalar $\delta'(0)=0$. With
\eqref{8.5:eq-zero-block-mean-3} this gives the second equality in \eqref{8.5:eq-zero-block-mean-a}.

\emph{Step 3: size.} The constant-flux form of each marked term $H_X[\mathbf{1}]$ is bounded by the
mass $\mathfrak{m}_T$ of the term $T=H_X[\mathbf{1}]$, written $\mathfrak{m}_X[\mathbf{1}]$
(clause (e) of the statement on flux forms of terms).
That mass is bounded by the majorant mass, $\mathfrak{m}_X[\mathbf{1}]\le\hat{\mathfrak{m}}_X$
(the mass comparison of the section on majorant masses). Summing over the $X\in E$ with $c(X)=c$, as
in \eqref{8.5:eq-zero-block-mean-2}, these two bounds give
$|q^E(c)[\mathbf{1}]|\le\mathfrak{a}^E_c$, which is assertion (2).
\end{proof}

\begin{remark}
\cite{B-CMP109} states that the effective actions are analytic in the coupling
$g_{j-1}\in[0,\gamma]$. That the terms of Ba{\l}aban's run are analytic in the bare weight, with
$\zeta$-independent index sets, is read from these statements through
\cite[(0.18), (0.20)]{B-CMP109}; this reading, which is hypothesis (viii), is stated, and its proof
is incomplete. For the families
of \cite{B-CMP119}, whose thresholds depend
polylogarithmically on $g_k^{-2}$ and so move with the bare weight, no statement is made here in the
first setting. Proposition~\ref{8.5:prop-extraction-complete} at $\zeta=0$, whose proof is incomplete
at one step, is stated for the families of both \cite{B-CMP109} and \cite{B-CMP119}; for the latter
its argument uses the representation theorem
\cite[(2.23)--(2.24)]{B-CMP119} in place of \cite[Theorem~1]{B-CMP109}.

The proof separates two contributions to \eqref{8.5:eq-zero-block-mean-a}, which concerns column
sums. Every component of the block-summed residual form other than its $W_4$-scalar part vanishes by
covariance alone. The scalar part vanishes if and only if the complete $W_4$-scalar coefficient of
the family is extracted, which is the content of Proposition~\ref{8.5:prop-extraction-complete},
whose proof is incomplete at one step: in
the first setting for Ba{\l}aban's family, in the second setting, at the step marked $(\ast)$, for the
witness's family. The row sums, that is the individual cell aggregates $q^E(c)[\mathbf{1}]$, do not
vanish. They form a $\Gamma$-periodic coefficient of zero block mean and of size at most
$\mathfrak{a}^E_c$.

The first setting does not contain the witness's family. The family $\operatorname{Irr}_1^E$ of the
two-point witness theorem is the frozen family of the second setting: its mark is a derivative in $z$
at frozen thresholds. The $\zeta$-derivative of the first setting differentiates Ba{\l}aban's run
itself. It therefore also differentiates that run's $\zeta$-dependent thresholds, cube sizes and
decompositions. The two derivatives differ exactly by the quantities that the run of the correlation
theorem holds frozen. Only the identity of representations for the shifted frozen run,
\eqref{8.5:eq-representation-identity-a}, whose proof is outstanding, asserts that the frozen family
still satisfies the identity \cite[(0.28)]{B-CMP109}.
\end{remark}

\begin{remark}[The remainder species: a scalar block mean small in the coupling]
\label{8.5:rem-remainder-species}
Let $\operatorname{Irr}_1^R$ be the remainder species. No part of the coupling increment
$\beta_{k+1}$ is extracted from $\operatorname{Irr}_1^R$ in Ba{\l}aban's construction
\cite{B-CMP119}; the fact is recorded in the running-shift statement. For a cell $c$ write
$q^R(c)$ for the cell aggregate $q_{\mathcal{V}}(c)$ of
Lemma~\ref{8.5:lem-block-sum-scalar} with $\mathcal{V}=\operatorname{Irr}_1^R$, and $q^R(c)[\mathbf{1}]$
for its value along the constant profile $\mathbf{1}$. Let $\mathfrak{B}$ be a period block, that
is, a cube of $M^4$ cells of $\pi_{k_0}$ forming a period of $\Gamma$. Then there is a number
$\delta_R$ such that
\begin{equation}\label{8.5:eq-remainder-species-1}
\sum_{c\in\mathfrak{B}} q^R(c)[\mathbf{1}] = \delta_R\,\langle\cdot,\cdot\rangle,
\qquad
|\delta_R|\,M^{-4}\le \bar{\mathfrak{a}}^R .
\end{equation}
The left side of \eqref{8.5:eq-remainder-species-1} is a scalar multiple of the pairing, hence
$W_4$-invariant, of the same form as the main term, and small only because the majorant
$\bar{\mathfrak{a}}^R$ is small in the coupling: no structural fact is available for the remainder
species, its block mean need not vanish, and no room is gained. The remainder species
is therefore paid for by its size alone, through the bound in
\eqref{8.5:eq-remainder-species-1}, in the estimate of the remainder contribution later in this
chapter.
\end{remark}

\begin{proof}
The remainder species $\operatorname{Irr}_1^R$ is a family of terms of the kind considered in
Lemma~\ref{8.5:lem-block-sum-scalar}, covariant generator-wise under the partition lattice $\Gamma$
and under $W_4$ by the translation-covariance statement and the $W_4$-covariance statement, the
two antecedents of Lemma~\ref{8.5:lem-block-sum-scalar}. The lemma does not require that a coupling
increment has been extracted from the
family, so it applies with $\mathcal{V}=\operatorname{Irr}_1^R$.

Since the constant profile $\mathbf{1}$ is invariant under translations by $\Gamma$ and under
$W_4$, the marked forms $q_X[\mathbf{1}]$, $X\in\operatorname{Irr}_1^R$, again form a family
covariant under $\Gamma$ and $W_4$, and part~(b) of the lemma applies to them. It shows that the
sum of the aggregates $q^R(c)[\mathbf{1}]$ over a period block
$\mathfrak{B}$ does not depend on the block, and that this sum is a scalar multiple of the pairing
$\langle\cdot,\cdot\rangle$; we write $\delta_R$ for this scalar (the scalar $\delta(\mathcal{V})$ of
the lemma for $\mathcal{V}=\operatorname{Irr}_1^R$, read along $\mathbf{1}$). This gives the identity in
\eqref{8.5:eq-remainder-species-1}.

By definition of the cell-aggregate majorant mass $\mathfrak{a}^R_c$, for
every cell $c$ and every $v\in\mathbb{R}^{\mathfrak{O}}$,
\begin{equation}\label{8.5:eq-remainder-species-2}
\bigl|q^R(c)[\mathbf{1}](v,v)\bigr|\le \mathfrak{a}^R_c\,\langle v,v\rangle,
\qquad
\mathfrak{a}^R_c\le\bar{\mathfrak{a}}^R:=\sup_{c'}\mathfrak{a}^R_{c'} .
\end{equation}
Evaluate the identity in \eqref{8.5:eq-remainder-species-1} at a vector $v$ with
$\langle v,v\rangle>0$ and sum \eqref{8.5:eq-remainder-species-2} over the $M^4$ cells of
$\mathfrak{B}$. This gives $|\delta_R|\,\langle v,v\rangle\le M^4\,\bar{\mathfrak{a}}^R\,\langle v,v\rangle$,
hence $|\delta_R|\,M^{-4}\le\bar{\mathfrak{a}}^R$.
\end{proof}

\begin{remark}[Per-term annihilation of constant flux and its obstacles]\label{8.5:rem-per-term-annihilation}
For a step $j$ and a region $X$, let $E^{(j)}_X$ be the localised term attached to $(j,X)$. Let $q_X$ be the constant-flux form of its marked Hessian, in the sense of Lemma~\ref{8.5:lem-constant-flux-forms}. The per-term form of the annihilation of constant flux is the assertion
\begin{equation}\label{8.5:eq-per-term-annihilation-1}
q_X=0 \qquad\text{for every localised term } E^{(j)}_X .
\end{equation}
This assertion is not proved here and is not used anywhere in this book; two routes towards it, and the obstacle on each, are as follows.

\emph{First route: scaling.} Consider the following \emph{fixed-function condition}. For fixed $j$ and $X$, the images of $E^{(j)}_X$ at all steps $k\ge j$ are given by one and the same analytic function, which satisfies the bound \cite[(0.29)]{B-CMP109} with one and the same constant $O(1)$ for all $k\ge j$.

Under this condition one would get $q_X=0$ for the fixed pair $(j,X)$ by letting $k-j\to\infty$, using the spare factor governed by $\alpha$ in \cite[(0.29)]{B-CMP109}. The reason is that the marginal content of a term carried from step $j$ to step $k$ scales exactly as $L^{-4(k-j)}$, by the scaling law in $\eta^2$. Here $\eta$ is the bare spacing in units of the step-$k$ lattice.

The fixed-function condition is not proved here, and no statement in Ba{\l}aban's papers supports it. Two facts speak against the consistency in $k$ of the quantity subtracted from a single term at each step:
\begin{itemize}
\item the passage of \cite{B-CMP109} stating that certain terms ``do not behave in a regular way under the renormalization transformations'';
\item the part of Ba{\l}aban's action $A_\eta(U_k)$ of the block field $U_k$ (notation of \cite{B-CMP109}, with $\eta$ as above) that depends on $k$.
\end{itemize}
Finally, for a term created at the witness level, $j=k_0$, the spare factor of the first route equals $1$.

\emph{Second route: extraction term by term.} This route requires a \emph{centring condition}: it replaces $E^{(j)}_X$ by $E^{(j)}_X$ minus its own point-attributed part of dimension four, that part being centred so that it is the full scalar of the term.

The centring condition is consistent with \cite{B-CMP109} and with \cite{B-CMP119}, but it is stated in neither; accordingly \eqref{8.5:eq-per-term-annihilation-1} is not obtained by this route here. At second order the object extracted by this route is the class tensor.

Neither route is pursued.
\end{remark}

\begin{definition}[The Gaussian frame at the witness level]\label{8.5:def-gaussian-frame}
Let $k_0=K-s$ be the witness level, $\ell=L^{-s}$ the side of a level-$k_0$ cell,
$\varepsilon=L^{-K}$ the bare spacing and $\mathsf{d}=C_W\ell$ the separation scale.
Let $f_1$ and $h_2$ be the two test functions of the witness, read on bare plaquettes, attached
to the balls $B_1=B(x_1,5\mathsf{d}/4)$ and $B_2=B(x_2,5\mathsf{d}/4)$ of the geometry of the
two-point witness theorem, with $\operatorname{supp}h_2\subset B_2$ and
$|x_1-x_2|\ge 3\mathsf{d}$. Let $n$ be the number of colours, so that
$b_0=n/2=(N^2-1)/2$, and let $T^1,\dots,T^n$ be the generators of the Lie algebra
$\mathfrak{g}$ of $\mathrm{SU}(N)$ used below, with trace normalisation $\tau$.
For a bare plaquette $p$, $O_p$ is the density of the main term, as in
Definition~\ref{8.5:def-frozen-sourced-family}, and the action is
$A=\sum_p O_p$.
\begin{itemize}
\item[(a)] \emph{The block law.} Under the tree block law at level $k_0$ (the facts stated in
(a)--(c), and the normalisation of $O_p$ in (d), are those of the block-law description in the
two-point witness construction), the fields
$b^1,\dots,b^n$ are independent centred Gaussian $1$-cochains on the level-$k_0$ lattice
with quadratic form $\tfrac12\langle b,\Delta_{k_0}b\rangle$, where $\Delta_{k_0}$ is the
operator of that law, and the block field is
\begin{equation}\label{8.5:eq-gaussian-frame-1}
V=\exp\Bigl(i\,g_{k_0}(2N)^{1/2}\sum_{a=1}^{n}b^aT^a\Bigr).
\end{equation}
\item[(b)] \emph{The background flux.} For each colour $a$ put
$\psi^a:=\bar{M}_{k_0}(db^a)$, a $2$-cochain on bare plaquettes. Then
\begin{equation*}
\operatorname{Cov}\bigl(\psi^a_p,\psi^{a'}_q\bigr)=\delta^{aa'}\,\mathcal{C}_{pq},
\qquad \mathcal{C}:=\mathcal{C}^{\mathrm{bg}},
\end{equation*}
and $\mathcal{C}$ is exact-valued: every column $\mathcal{C}_{\cdot q}$ lies in the space
$\mathcal{E}$ of exact $2$-cochains on the bare lattice, the range of $\mathcal{C}^{\mathrm{bg}}$
being contained in $\mathcal{E}$.
\item[(c)] \emph{The derivative at the identity.} The map $U_{k_0}(\cdot)$ is twice
differentiable at $\mathbf{1}$, and $M_{k_0}$ denotes its derivative there; $M_{k_0}$ is
$\operatorname{Ad}$-equivariant and intertwines the coboundaries, so that
$d(M_{k_0}b^a)=\bar{M}_{k_0}(db^a)=\psi^a$ for each colour $a$.
\item[(d)] \emph{The normalised Hessian and the mass per cell.} For $T$ as in
Lemma~\ref{8.5:lem-constant-flux-forms} (a functional; no confusion with the generators $T^a$
arises), $\operatorname{Hess}_UT$ denotes its Hessian at
$\mathbf{1}$ on one colour component, read as a form in the fluxes, and
$\operatorname{Hess}^{\mathfrak{n}}:=(2N)\tau\cdot\operatorname{Hess}_U$, as recalled in
\eqref{8.5:eq-gaussian-frame-b}; in this
normalisation $\operatorname{Hess}^{\mathfrak{n}}O_p[\psi,\psi]=\sum_a(\psi^a_p)^2$, so that a
unit flux has unit mass on each bare plaquette. We let $\mathfrak{n}$ be the value of
$\operatorname{Hess}^{\mathfrak{n}}A$ per level-$k_0$ cell on the flux equal to $1$ on every
bare plaquette (unit constant flux per bare plaquette); then
\begin{equation}\label{8.5:eq-gaussian-frame-2}
\mathfrak{n}=6(\ell/\varepsilon)^4 .
\end{equation}
\item[(e)] \emph{The two quadratic parts.} For a profile $f$ let $D[f]$ and the localised
terms $V^{zf}_X$ be those of Definition~\ref{8.5:def-frozen-sourced-family}, put
$H_X[f_1]:=D[f_1]\operatorname{Hess}^{\mathfrak{n}}V^{zf_1}_X$ and let $Q_X$ be the flux form of
$H_X[f_1]$. The constant-flux forms $q_X$, the masses $\mathfrak{m}_X$ and the majorants
$\hat{\mathfrak{m}}_X$ of these terms are taken in the same normalisation; the unweighted
majorant $\hat{\mathfrak{m}}^0_X$, which is $\mathfrak{n}$ times a constant times a shape
factor, is already stated relative to $\mathfrak{n}$. Let $Q_1$ be the $b$-quadratic part
at $\mathbf{1}$ of $\sum_XD[f_1]\bigl(V^{zf_1}_X\circ U_{k_0}\bigr)$, and let
$Q_2:=\mathcal{O}^{(k_0)}_{\mathrm{quad}}(h_2)$ be the $b$-quadratic part at $\mathbf{1}$ of
the main term $\sum_qh_2(q)\,O_q\circ U_{k_0}$ of the second observable at level $k_0$. Then
\begin{equation}\label{8.5:eq-gaussian-frame-b}
\begin{gathered}
\operatorname{Hess}^{\mathfrak{n}}=(2N)\tau\cdot\operatorname{Hess}_U,\\
Q_1=\frac{g^2_{k_0}}{2}\sum_a\sum_XQ_X[\psi^a,\psi^a],\qquad
Q_2=\frac{g^2_{k_0}}{2}\sum_qh_2(q)\sum_a(\psi^a_q)^2 ,
\end{gathered}
\end{equation}
with the same factor $g^2_{k_0}/2$ in front of $Q_1$ and of $Q_2$.
\item[(f)] \emph{The covariance size on the near ball.} Put
$\mathfrak{A}:=B(x_1,3\mathsf{d}/2)$ and
\begin{equation}\label{8.5:eq-gaussian-frame-3}
\mathsf{S}_q:=\max\bigl\{|\mathcal{C}_{pq}|:\ p\text{ bare in }\mathfrak{A}\bigr\}.
\end{equation}
For $q\in\operatorname{supp}h_2$ every bare $p$ in $\mathfrak{A}$ satisfies
$\operatorname{dist}(p,q)\ge\mathsf{d}/4$. The larger ball $B(x_1,2\mathsf{d})$ would not have
this property: it meets $B(x_2,5\mathsf{d}/4)$ whenever $|x_1-x_2|<13\mathsf{d}/4$, so that a
maximum over it can be a near-coincident value. The quantity $\mathsf{S}_q$ enters only the
estimates on $\mathfrak{A}$; the contributions of terms reaching outside $\mathfrak{A}$ are
controlled through $\sup_{\hat{X}}|\mathcal{C}_{\cdot q}|$ by a growth bound, and never through
$\mathsf{S}_q$.
\item[(g)] \emph{Scale and unit.} For $r>0$ put $\sigma(r):=\varepsilon^4(\pi^2r^4)^{-1}$ and,
for bare plaquettes $p,p'$, $\sigma_{pp'}:=\sigma(|x(p)-x(p')|)$; this is the positive radial
majorant of the kernel entries used in the two-point witness construction, a function of the
distance only and not a kernel entry. Let $N_1$ be the number of level-$k_0$ cells meeting
$B(x_1,3\mathsf{d}/2)$. We put
\begin{equation}\label{8.5:eq-gaussian-frame-a}
\begin{gathered}
\sigma_{\mathsf{d}}:=2\,\sigma(\mathsf{d}/4)
=\frac{2\varepsilon^4}{\pi^2(\mathsf{d}/4)^4}
=2\cdot4^4\pi^{-2}\varepsilon^4C_W^{-4}\ell^{-4},\\
\mathfrak{U}_q:=\|f_1\|_\infty N_1\sigma_{\mathsf{d}}^2,\qquad
\mathfrak{U}:=b_0\,g^4_{k_0}\sum_q|h_2(q)|\,\mathfrak{U}_q .
\end{gathered}
\end{equation}
The unit $\mathfrak{U}$ contains no mass, no moment and no kernel maximum; it is the unit of
the second-derivative room lemma and carries the two block-law factors
$(g^2_{k_0}/2)^2$ and the factor $2n$ of Wick's theorem. The two-point witness theorem
measures every remainder term in its own unit, $b_0g^4_{k_0}$ times the kernel constant of
that theorem; this unit is a different quantity, is always written $\mathfrak{U}_W$, or
$\mathfrak{U}_W(h)$, and enters where estimates in the unit $\mathfrak{U}$ are converted
into it.
\end{itemize}
\end{definition}

\begin{proof}[Proof of the identities in Definition~\ref{8.5:def-gaussian-frame}]
\emph{Step 1: the $b$-quadratic part of a gauge-invariant analytic term.} Let $T$ be gauge
invariant and analytic on an open neighbourhood of $\mathbf{1}$, as in
Lemma~\ref{8.5:lem-constant-flux-forms}, which is proved as an implication from per-term
gauge invariance. Put $Y(b):=g_{k_0}(2N)^{1/2}\sum_ab^aT^a$, so that $V=\exp(iY(b))$ by
\eqref{8.5:eq-gaussian-frame-1}, and $Y$ is linear in $b$. By the twice-differentiability of
$U_{k_0}$ at $\mathbf{1}$ (Definition~\ref{8.5:def-gaussian-frame}\,(c)), we expand
$T\circ U_{k_0}(\exp(iY(b)))$ at $b=0$ to second order. The first-order term is
$DT(\mathbf{1})\,M_{k_0}[iY(b)]$. The second-order term has two parts: $DT(\mathbf{1})$
applied to the second-order part of $U_{k_0}(\exp(iY(b)))$, which collects the term of
$D^2U_{k_0}$ and the image under $M_{k_0}$ of the quadratic term of the exponential; and
$\tfrac12\operatorname{Hess}T$ at $\mathbf{1}$ evaluated on the linear part $M_{k_0}[iY(b)]$ in
both slots. By Lemma~\ref{8.5:lem-constant-flux-forms}\,(a), $DT(\mathbf{1})=0$.
Hence there is no first-order term, the first part of the second-order term vanishes, and
the $b$-quadratic part is
$\tfrac12\operatorname{Hess}T\bigl[M_{k_0}[iY(b)],M_{k_0}[iY(b)]\bigr]$. By
Definition~\ref{8.5:def-gaussian-frame}\,(c), $M_{k_0}$ carries $b^a$ to a bond field whose
coboundary is $\psi^a$, colour by colour; by Lemma~\ref{8.5:lem-constant-flux-forms}\,(b) the
Hessian depends on its arguments only through their coboundaries, so, read as a form in the
fluxes, it is evaluated on the flux whose colour-$a$ component is $g_{k_0}(2N)^{1/2}\psi^a$.
Expanding
by bilinearity gives a double sum over colours $a,a'$. Since $M_{k_0}$ is
$\operatorname{Ad}$-equivariant and $\mathfrak{g}$ is simple, Schur's lemma reduces this
double sum to one colour at a time: the terms with $a\ne a'$ vanish and each diagonal term
carries the trace normalisation $\tau$ of the generators $T^a$. The $b$-quadratic part of
$T\circ U_{k_0}$ is therefore
\begin{equation}\label{8.5:eq-gaussian-frame-4}
\frac12\,g^2_{k_0}(2N)\tau\sum_a\operatorname{Hess}_UT[\psi^a,\psi^a]
=\frac{g^2_{k_0}}{2}\sum_a\operatorname{Hess}^{\mathfrak{n}}T[\psi^a,\psi^a],
\end{equation}
where the equality is the definition of $\operatorname{Hess}^{\mathfrak{n}}$ in
\eqref{8.5:eq-gaussian-frame-b}.

\emph{Step 2: the observable and $Q_2$.} For $T=O_p$ the block-law description recalled in
(a) gives that the left side of \eqref{8.5:eq-gaussian-frame-4} equals
$(g^2_{k_0}/2)\sum_a(\psi^a_p)^2$; equivalently
$\operatorname{Hess}^{\mathfrak{n}}O_p[\psi,\psi]=\sum_a(\psi^a_p)^2$, which is the
normalisation stated in Definition~\ref{8.5:def-gaussian-frame}\,(d). Multiplying by
$h_2(q)$ with $p=q$ and summing over bare plaquettes $q$, the $b$-quadratic part of
$\sum_qh_2(q)\,O_q\circ U_{k_0}$ is $(g^2_{k_0}/2)\sum_qh_2(q)\sum_a(\psi^a_q)^2$, which is
the formula for $Q_2$ in \eqref{8.5:eq-gaussian-frame-b}.

\emph{Step 3: the marked term and $Q_1$.} Each localised term $V^{zf_1}_X$ is, for every $z$,
gauge invariant and analytic near $\mathbf{1}$, this being the format of
Definition~\ref{8.5:def-frozen-sourced-family}; so the antecedent of
Lemma~\ref{8.5:lem-constant-flux-forms} holds for it, and part (a) of that lemma gives
$DV^{zf_1}_X(\mathbf{1})=0$ for every $z$. Since $D[f_1]$ is the derivative in $z$ at
$z=0$, the differential at $\mathbf{1}$ of $D[f_1]V^{zf_1}_X$ vanishes as well. Step~1
therefore applies to $V^{zf_1}_X$ for each $z$: its $b$-quadratic part is
$(g^2_{k_0}/2)\sum_a\operatorname{Hess}^{\mathfrak{n}}V^{zf_1}_X[\psi^a,\psi^a]$. Applying the
linear operation $D[f_1]$, the $b$-quadratic part of $D[f_1](V^{zf_1}_X\circ U_{k_0})$ is
\begin{equation*}
\frac{g^2_{k_0}}{2}\sum_aH_X[f_1][\psi^a,\psi^a]
=\frac{g^2_{k_0}}{2}\sum_aQ_X[\psi^a,\psi^a],
\end{equation*}
the reduction to one colour at a time being the same use of the $\operatorname{Ad}$-equivariance
of $M_{k_0}$ and of Schur's lemma as in Step~1 for $O_p$. Summing over $X$ gives the
formula for $Q_1$ in \eqref{8.5:eq-gaussian-frame-b}, with the same factor $g^2_{k_0}/2$ as
in front of $Q_2$.

\emph{Step 4: the value of $\mathfrak{n}$.} By Step~2, a flux equal to $1$ on every bare
plaquette, in one colour, contributes $\operatorname{Hess}^{\mathfrak{n}}O_p[\psi,\psi]=1$
for each bare plaquette $p$. Since $A=\sum_pO_p$, the value of
$\operatorname{Hess}^{\mathfrak{n}}A$ per level-$k_0$ cell on this flux is the number of bare
plaquettes per cell. There are six bare plaquettes per bare site, one for each of the
$\binom42=6$ coordinate planes, and $(\ell/\varepsilon)^4$ bare sites per cell of side
$\ell$. Hence $\mathfrak{n}=6(\ell/\varepsilon)^4$ exactly, which is
\eqref{8.5:eq-gaussian-frame-2}.

\emph{Step 5: the distance in (f).} Let $q\in\operatorname{supp}h_2\subset B(x_2,5\mathsf{d}/4)$
and let $p$ be bare in $\mathfrak{A}=B(x_1,3\mathsf{d}/2)$. By the triangle inequality and
$|x_1-x_2|\ge3\mathsf{d}$,
\begin{equation*}
\operatorname{dist}(p,q)\ \ge\ |x_1-x_2|-\frac{3\mathsf{d}}{2}-\frac{5\mathsf{d}}{4}
\ \ge\ 3\mathsf{d}-\frac{3\mathsf{d}}{2}-\frac{5\mathsf{d}}{4}=\frac{\mathsf{d}}{4}.
\end{equation*}
For the ball $B(x_1,2\mathsf{d})$, the distance between $x_1$ and $x_2$ is less than the sum
$2\mathsf{d}+5\mathsf{d}/4=13\mathsf{d}/4$ of the two radii whenever
$|x_1-x_2|<13\mathsf{d}/4$, and then the two balls meet.

\emph{Step 6: the constants in (g).} With $\mathsf{d}=C_W\ell$ one has
$(\mathsf{d}/4)^4=C_W^4\ell^4/4^4$, so
\begin{equation*}
\frac{2\varepsilon^4}{\pi^2(\mathsf{d}/4)^4}
=\frac{2\cdot4^4\varepsilon^4}{\pi^2C_W^4\ell^4}
=2\cdot4^4\pi^{-2}\varepsilon^4C_W^{-4}\ell^{-4},
\end{equation*}
which is the chain in \eqref{8.5:eq-gaussian-frame-a}. Finally, the two block-law factors
and the factor of Wick's theorem combine to the prefactor of $\mathfrak{U}$:
\begin{equation*}
\Bigl(\frac{g^2_{k_0}}{2}\Bigr)^2\cdot2n=\frac{n}{2}\,g^4_{k_0}=b_0\,g^4_{k_0}.
\end{equation*}
\end{proof}

\begin{lemma}[Upper bound of the background covariance on the separated range]
\label{8.5:lem-covariance-upper}
Let $\mathsf{d}=C_W\ell$ be the separation scale and let $x_1,x_2$ be the centres of
Definition~\ref{8.5:def-separation-scale}, so that $3\mathsf{d}\le|x_1-x_2|\le A_c\mathsf{d}$,
where $A_c\ge4$ is the centre-separation number fixed there. Let $\mathcal{C}=\mathcal{C}^{\mathrm{bg}}$
be the background flux covariance of Definition~\ref{8.5:def-gaussian-frame}. Let $G$ be the free
kernel and $C^{\mathrm{fl}}$ the fluctuation kernel of the two-point witness theorem, so that
$\mathcal{C}=G-C^{\mathrm{fl}}$. For bare
plaquettes $p$ we write $x_p=x(p)$ for the barycentre. For $r>0$ put
\[
\sigma(r):=\frac{\varepsilon^4}{\pi^2r^4},
\]
so that the radial majorant of Definition~\ref{8.5:def-gaussian-frame} is
$\sigma_{pq}=\sigma(|x_p-x_q|)$, and $\sigma_{\mathsf{d}}=2\sigma(\mathsf{d}/4)$ as in
\eqref{8.5:eq-gaussian-frame-a}. Recall from Definition~\ref{8.5:def-gaussian-frame} that, for a
bare plaquette $q$,
\[
\mathsf{S}_q=\max\bigl\{|\mathcal{C}_{pq}|:\ p\text{ bare},\ x_p\in B(x_1,3\mathsf{d}/2)\bigr\}.
\]
The pairs considered are bare plaquettes $p,q$ with
\[
\tfrac14\mathsf{d}\le|x_p-x_q|\le(A_c+3)\mathsf{d};
\]
for every pair of the conclusion below, $|x_p-x_q|$ lies in
$[\tfrac14\mathsf{d},(A_c+\tfrac{11}4)\mathsf{d}]$. The
lower pair factor $a_*$ of the free-kernel lemma and the fluctuation decay lemma, and the
plaquette-pair factor $A_*$ of Definition~\ref{8.5:def-separation-scale}, are here replaced by
$\tfrac14$ and $A_c+3$.

We use the following constants. $C_{\mathrm{Green}}$ is the constant of the lattice Green's function
asymptotics, and $z_0$ is the threshold in the distance clause of the free-kernel lemma. $C_T$ and
$\mathfrak{N}$ are the constant and the size
parameter in which the torus comparison lemma expresses its correction. $C''$, $\delta'$, $\ell'$ and
$C_0'$ are the constants in which the fluctuation decay lemma states its bound. With these constants
define
\begin{align}
\eta_G&:=C_{\mathrm{Green}}\Bigl(\frac{4\varepsilon}{\mathsf{d}}\Bigr)^2
+C_T\pi^2\Bigl(\frac{(A_c+3)\mathsf{d}}{\varepsilon}\Bigr)^4\mathfrak{N}^{-4},
\label{8.5:eq-covariance-upper-1}\\
\mathfrak{r}&:=C''L^4\bigl((A_c+3)C_W\bigr)^4e^{-\delta'C_W/(4\ell')}.
\label{8.5:eq-covariance-upper-2}
\end{align}
Here $\eta_G$ is the closeness parameter of the free kernel; it is not the ratio $\varepsilon/\ell$.

Hypotheses (a) and (b) below are, respectively, the conclusion of the lattice Green's function
asymptotics combined with the torus comparison lemma, in the form in which the proof of the
free-kernel lemma uses it, and the conclusion of the fluctuation decay lemma, whose proof derives it
from the decay bound for the fluctuation propagator; each is stated for the pairs considered here.
Assume the following.
\begin{enumerate}
\item[(a)] \emph{Free kernel.} Let $D_{PQ}(x)$,
$x\in\mathbb{R}^4\setminus\{0\}$, indexed by pairs of planes $P,Q$, be the continuum flux
two-point tensor of the flux-tensor lemma. Let $p,q$ be bare plaquettes in the range displayed
above. Suppose that
$|x_p-x_q|\le L^m/2$ (the hypothesis of the torus comparison lemma). Suppose also that the distance
clause of the free-kernel lemma holds at the minimal pair distance $\mathsf{d}/4$; there this clause
reads $\mathsf{d}/(8\varepsilon)\ge z_0$. Then
\[
\bigl|G_{pq}-\varepsilon^4D_{\pi(p)\pi(q)}(x_p-x_q)\bigr|\le\eta_G\,\sigma_{pq}.
\]
\item[(b)] \emph{Fluctuation kernel.} Let $p,q$ be bare plaquettes in the range displayed above, and
suppose $C_W\ge C_0^{\mathrm{dec}}(\tfrac14)$, with $C_0^{\mathrm{dec}}(\tfrac14)$ as in
\eqref{8.5:eq-covariance-upper-3} below. Then
\[
|C^{\mathrm{fl}}_{pq}|\le\mathfrak{r}\,\sigma_{pq}.
\]
\item[(c)] \emph{Floors.} The following one-sided conditions hold:
\begin{align}
C_W&\ge C_0^{\mathrm{dec}}(\tfrac14)
:=\max\Bigl\{C_0',\ \frac{64\,\ell'\log L}{\delta'(L-1)}\Bigr\},
\label{8.5:eq-covariance-upper-3}\\
L^{m+s}&\ge2(A_c+3)C_W,\label{8.5:eq-covariance-upper-4}\\
C_WL^{K-s}&\ge8z_0.\label{8.5:eq-covariance-upper-5}
\end{align}
In addition,
\[
\eta_G+\mathfrak{r}\le1.
\]
\end{enumerate}
Condition \eqref{8.5:eq-covariance-upper-4} is one of the lower bounds on $m$ making up the torus
clause $m\ge m_{\mathbb{T}}$ of the quantifier prefix (Notation~\ref{2.1:not-prefix}). Condition
\eqref{8.5:eq-covariance-upper-5} is one of the lower bounds on the witness level $K-s$ making up
the floor $K-s\ge k^W_*$, where $k^W_*$ is the level threshold of the two-point witness theorem.
None of these conditions is a cap in the sense of Definition~\ref{2.3:def-cap}.

Under these assumptions, let $p$ be a bare plaquette with $x_p\in B(x_1,3\mathsf{d}/2)$, and let $q$
be a bare plaquette with $x_q\in B_2=B(x_2,5\mathsf{d}/4)$. Then
\begin{equation}\label{8.5:eq-covariance-upper-a}
|\mathcal{C}_{pq}|=|G_{pq}-C^{\mathrm{fl}}_{pq}|
\le(1+\eta_G+\mathfrak{r})\,\sigma_{pq}
\le2\sigma(\mathsf{d}/4)=\sigma_{\mathsf{d}}.
\end{equation}
Hence $\mathsf{S}_q\le\sigma_{\mathsf{d}}$ for every bare $q$ with $x_q\in B_2$.
\end{lemma}

\begin{proof}
\emph{The pair-distance range.} Let $p$ and $q$ be as in the conclusion. By the triangle
inequality and the lower separation bound $|x_1-x_2|\ge3\mathsf{d}$ of
Definition~\ref{8.5:def-separation-scale},
\[
|x_p-x_q|\ \ge\ |x_1-x_2|-|x_p-x_1|-|x_q-x_2|
\ \ge\ 3\mathsf{d}-\tfrac32\mathsf{d}-\tfrac54\mathsf{d}=\tfrac14\mathsf{d}.
\]
By the upper separation bound $|x_1-x_2|\le A_c\mathsf{d}$ of the same definition,
\[
|x_p-x_q|\ \le\ |x_1-x_2|+|x_p-x_1|+|x_q-x_2|
\ \le\ \bigl(A_c+\tfrac32+\tfrac54\bigr)\mathsf{d}=\bigl(A_c+\tfrac{11}4\bigr)\mathsf{d}.
\]
Since $A_c+\tfrac{11}4\le A_c+3$, every such pair lies in the range
\[
\tfrac14\mathsf{d}\le|x_p-x_q|\le(A_c+3)\mathsf{d}.
\]
This is the range for which hypotheses (a) and (b) are stated. In the free-kernel lemma and the
fluctuation decay lemma the pair factors are those of the witness's own geometry: lower factor $1$
(or $\tfrac12$ on the second geometry for which these lemmas are stated) and upper factor the
witness's plaquette-pair factor. Here they are
replaced by $\tfrac14$ and $A_c+3$. This replacement enters at four places: the two members
of $\eta_G$ in \eqref{8.5:eq-covariance-upper-1}, the torus hypothesis, the distance clause at the
minimal pair distance, and $\mathfrak{r}$ in \eqref{8.5:eq-covariance-upper-2}.

\emph{The provisos of (a) and (b).} First, the torus hypothesis. We have $\mathsf{d}=C_W\ell$ with
$\ell=L^{-s}$. Hence the maximal pair distance satisfies
\[
(A_c+3)\mathsf{d}=(A_c+3)C_WL^{-s}\le\tfrac12L^m,
\]
which is equivalent to \eqref{8.5:eq-covariance-upper-4}. Therefore
$|x_p-x_q|\le(A_c+3)\mathsf{d}\le L^m/2$, and the hypothesis of the torus comparison lemma holds.

Second, the distance clause. Since $\varepsilon=L^{-K}$,
\[
\frac{\mathsf{d}}{8\varepsilon}=\frac{C_WL^{K-s}}{8}.
\]
By \eqref{8.5:eq-covariance-upper-5} this is at least $z_0$, so the distance clause of the
free-kernel lemma at the minimal pair distance $\mathsf{d}/4$ holds.

Third, the condition $C_W\ge C_0^{\mathrm{dec}}(\tfrac14)$ of (b) is
\eqref{8.5:eq-covariance-upper-3}. All three provisos are therefore met for the pair $(p,q)$.

\emph{The free kernel.} By hypothesis (a),
\[
\bigl|G_{pq}-\varepsilon^4D_{\pi(p)\pi(q)}(x_p-x_q)\bigr|\le\eta_G\sigma_{pq}.
\]
The continuum tensor of the flux-tensor lemma satisfies the entrywise bound
$|D_{PQ}(x)|^2\le\pi^{-4}|x|^{-8}$, on which the proof of the free-kernel lemma rests; taking
square roots gives $|D_{PQ}(x)|\le\pi^{-2}|x|^{-4}$. Hence, for each entry,
\[
\bigl|\varepsilon^4D_{\pi(p)\pi(q)}(x_p-x_q)\bigr|
\le\frac{\varepsilon^4}{\pi^2|x_p-x_q|^4}=\sigma_{pq}.
\]
Combining the two estimates,
\begin{align*}
|G_{pq}|&\le\bigl|\varepsilon^4D_{\pi(p)\pi(q)}(x_p-x_q)\bigr|
+\bigl|G_{pq}-\varepsilon^4D_{\pi(p)\pi(q)}(x_p-x_q)\bigr|\\
&\le(1+\eta_G)\,\sigma_{pq}.
\end{align*}

\emph{The fluctuation kernel.} Hypothesis (b) is the conclusion of the fluctuation decay lemma on the
present range: the derivation of that lemma from the decay bound for the fluctuation propagator,
with the witness's pair factors replaced by $\tfrac14$ and $A_c+3$, gives the constant
$\mathfrak{r}$ of \eqref{8.5:eq-covariance-upper-2} and the floor $C_0^{\mathrm{dec}}(\tfrac14)$ of
\eqref{8.5:eq-covariance-upper-3}. By hypothesis (b), whose proviso was checked above,
$|C^{\mathrm{fl}}_{pq}|\le\mathfrak{r}\,\sigma_{pq}$.

\emph{The covariance.} Since $\mathcal{C}_{pq}=G_{pq}-C^{\mathrm{fl}}_{pq}$, the two bounds give
\[
|\mathcal{C}_{pq}|\le|G_{pq}|+|C^{\mathrm{fl}}_{pq}|\le(1+\eta_G+\mathfrak{r})\,\sigma_{pq}.
\]
This is the first inequality in \eqref{8.5:eq-covariance-upper-a}. By (c), $\eta_G+\mathfrak{r}\le1$,
so the right-hand side is at most $2\sigma_{pq}$. The function $r\mapsto\sigma(r)$ is decreasing on
$(0,\infty)$, and $|x_p-x_q|\ge\mathsf{d}/4$. Hence
$\sigma_{pq}=\sigma(|x_p-x_q|)\le\sigma(\mathsf{d}/4)$, and therefore
\[
(1+\eta_G+\mathfrak{r})\,\sigma_{pq}\le2\sigma(\mathsf{d}/4).
\]
By \eqref{8.5:eq-gaussian-frame-a}, $2\sigma(\mathsf{d}/4)=\sigma_{\mathsf{d}}$. Explicitly, using
$\mathsf{d}=C_W\ell$,
\[
2\sigma(\mathsf{d}/4)=\frac{2\cdot4^4\varepsilon^4}{\pi^2\mathsf{d}^4}
=2\cdot4^4\pi^{-2}\varepsilon^4C_W^{-4}\ell^{-4}.
\]
This proves \eqref{8.5:eq-covariance-upper-a}.

\emph{The maximum.} Fix a bare $q$ with $x_q\in B_2$. The quantity $\mathsf{S}_q$ is the maximum of
$|\mathcal{C}_{pq}|$ over the finitely many bare $p$ with $x_p\in B(x_1,3\mathsf{d}/2)$. Each of
these entries is at most $\sigma_{\mathsf{d}}$ by \eqref{8.5:eq-covariance-upper-a}. Hence
$\mathsf{S}_q\le\sigma_{\mathsf{d}}$.
\end{proof}

Along the floors $K-s\ge k^W_*$, $m\ge m_{\mathbb{T}}$ and $C_W\ge C_0^{\mathrm{dec}}(\tfrac14)$,
the numbers $\eta_G$ and $\mathfrak{r}$ tend to zero: they are the closeness parameters of the
free-kernel lemma and of the fluctuation decay lemma, with the pair factors replaced as above, which
enter the rate $r_{\mathrm{ker}}$ of the kernel comparison lemma, and the derivations of their decay
there are unchanged by the replacement. Lemma~\ref{8.5:lem-covariance-upper} uses only
$\eta_G+\mathfrak{r}\le1$.

The floor \eqref{8.5:eq-covariance-difference-a} of the cell-scale difference bound for the
background covariance consists of the two conditions
\[
C_W\ge C_\nabla^{\mathrm{dec}},
\qquad
L^{4(m+s)}\ge\frac{C_T\pi^2(A_c+3)^4C_W^5}{C_\nabla^0},
\]
where $C_\nabla^{\mathrm{dec}}$ and $C_\nabla^0$ are the constants of that bound. These conditions are
not hypotheses of Lemma~\ref{8.5:lem-covariance-upper}. The second of them is a lower bound on $m$
belonging to the torus clause; neither is a cap in the sense of Definition~\ref{2.3:def-cap}.

The ball $B_2$ contains the support of the weight $h_2$ of
Definition~\ref{8.5:def-separation-scale}, up to the distance-decaying profile tail of $h_2$. That
tail is estimated in clause (vi) of the second-derivative room lemma, and the present lemma is not
used there.

The bounds of Lemma~\ref{8.5:lem-covariance-upper} are upper bounds on moduli only. No lower bound
on an individual entry $\mathcal{C}_{pq}$, or on $\mathsf{S}_q$, is used anywhere in this section.
For an individual entry no such bound holds. For $\mathsf{S}_q$, whose maximum runs over plaquettes
$p$ of all six planes, a lower bound does hold, up to a term of order $\varepsilon/\mathsf{d}$ times
the radial majorant $\sigma$. That term comes from
the offsets, of order $\varepsilon$, of the barycentres of the six plaquettes at a site, which make
up the bracket of the cell-scale difference bound; this lower bound is not needed.

The reason no entrywise lower bound holds is that the flux two-point function is a sign-indefinite
tensor in the pair of planes $(\pi(p),\pi(q))$ and in the direction $x_p-x_q$. For the continuum
tensor $D$ of the flux-tensor lemma, writing $(ij)$ for the coordinate plane spanned by the $i$-th
and $j$-th axes, that lemma gives $D_{(12),(34)}\equiv0$, and
\[
D_{(12),(12)}(x)\ \propto\ \bigl(x_3^2+x_4^2-x_1^2-x_2^2\bigr)|x|^{-6}
\]
vanishes on a cone. Thus the entries vanish identically between orthogonal planes and on
cones within a plane. Moreover, by clause (v) of the kernel-form lemma, on the cone of non-negative
plaquette weights the ratio of the continuum form $\mathcal{K}^{\mathrm{cont}}$ to the majorant form
$\mathcal{K}^{+}$ can be arbitrarily small for weights that are not scalar.

The positive objects that the two-point witness theorem bounds from below are the following:
\begin{itemize}
\item the majorant form $\mathcal{K}^{+}$ (clause (iii) of the kernel-form lemma);
\item on the non-negative cone of scalar weights, the continuum form $\mathcal{K}^{\mathrm{cont}}$,
which there equals $\mathcal{K}^{+}$ (clauses (i) and (iv) of the kernel-form lemma), and with it
the background form $\mathcal{K}^{\mathrm{bg}}$, since
$\mathcal{K}^{\mathrm{bg}}=\mathcal{K}^{\mathrm{cont}}\bigl(1+O(r_{\mathrm{ker}})\bigr)$ by the
kernel comparison lemma.
\end{itemize}
The conversion step of this section uses exactly these.

\begin{lemma}[Cell-scale difference bound for the background covariance]
\label{8.5:lem-covariance-difference}
Let $\mathcal{C}=\mathcal{C}^{\mathrm{bg}}$ be the covariance of the background flux in the Gaussian
frame at the witness level $k_0=K-s$ (Definition~\ref{8.5:def-gaussian-frame}). Let
$\varepsilon=L^{-K}$ be the bare spacing, $\ell=L^{-s}$ the side of a level-$k_0$ cell and
$\mathsf{d}=C_W\ell$ the separation scale. Let $x_1,x_2$ be the centres of the two supports, and assume
$3\mathsf{d}\le|x_1-x_2|\le A_c\mathsf{d}$, the separation window of the geometry used in
Lemma~\ref{8.5:lem-covariance-upper} (a condition on the supports of the two test functions). Put
\begin{equation*}
\mathfrak{A}:=B(x_1,3\mathsf{d}/2),\qquad \mathfrak{A}':=B(x_1,13\mathsf{d}/8),\qquad
B_2:=B(x_2,5\mathsf{d}/4).
\end{equation*}
For a bare plaquette $p$ write $x_p:=x(p)$ for its barycentre and $\pi(p)$ for its plane, and write
$p\in E$ for a set $E$ when $x_p\in E$. For $t>0$
put $\sigma(t):=\varepsilon^4(\pi^2t^4)^{-1}$ and $\sigma_{pq}:=\sigma(|x_p-x_q|)$, and let
$\sigma_{\mathsf{d}}=2\sigma(\mathsf{d}/4)=2\cdot4^4\pi^{-2}\varepsilon^4\mathsf{d}^{-4}$ be the scale
fixed in \eqref{8.5:eq-gaussian-frame-a}. Assume the following.
\begin{enumerate}
\item[(a)] The identity $\mathcal{C}=G-C^{\mathrm{fl}}$ holds entrywise. Here $G$ is the free kernel and
$C^{\mathrm{fl}}$ the fluctuation kernel of the construction of the background covariance.
\item[(b)] The lattice Green's function asymptotics, with constant $C_{\mathrm{Green}}$, and the torus
lemma, with constant $C_T$, hold.
\item[(c)] The decay bounds for Ba{\l}aban's propagators hold.
\item[(d)] The explicit-tensor lemma for the continuum flux two-point tensor $D=(D_{PQ})$ holds.
Here $D$ is indexed by pairs $(P,Q)$ of planes. Its entries are rational functions of degree $-4$,
with quadratic numerators over $|z|^6$. They are therefore smooth on
$\mathbb{R}^4\setminus\{0\}$ and satisfy $|\nabla D_{PQ}(z)|\le C_D|z|^{-5}$ with an absolute constant
$C_D$, and $C_D\le 8\pi^{-2}$ in the normalisation of that lemma.
\end{enumerate}
Assume also the floors of Lemma~\ref{8.5:lem-covariance-upper}.

Let $\eta$ and $\mathfrak{r}$ be the relative errors of Lemma~\ref{8.5:lem-covariance-upper}: for bare
pairs $p,q$ with $\mathsf{d}/4\le|x_p-x_q|\le(A_c+3)\mathsf{d}$, the range on which the two proofs are
re-run in Lemma~\ref{8.5:lem-covariance-upper} (with minimal pair distance $\mathsf{d}/4$ and upper
pair factor $A_c+3$),
\begin{equation*}
G_{pq}=\varepsilon^4D_{\pi(p)\pi(q)}(x_p-x_q)+\rho_{pq},\qquad |\rho_{pq}|\le\eta\,\sigma_{pq},
\end{equation*}
where $\eta$ is at most the sum of the two terms
\begin{equation*}
16C_{\mathrm{Green}}C_W^{-2}L^{-2(K-s)}
\qquad\text{and}\qquad
C_T\pi^2(A_c+3)^4C_W^4L^{-4(m+s)};
\end{equation*}
here $A_c$ is the centre-separation constant of the geometry of the two supports. Further
$|C^{\mathrm{fl}}_{pq}|\le\mathfrak{r}\,\sigma_{pq}$ on the same range, with
\begin{equation*}
\mathfrak{r}=\mathfrak{r}(C_W):=C''L^4\bigl((A_c+3)C_W\bigr)^4e^{-\delta'C_W/(4\ell')},
\end{equation*}
where $C''$, $\delta'$, $\ell'$ are the constants of the decay lemma for the fluctuation kernel
($\ell'$ is a length constant of that lemma, unrelated to the cell side $\ell$). Put
\begin{equation}\label{8.5:eq-covariance-difference-1}
C_\nabla^0:=\max\{2\pi^2C_D,\,1\},\qquad
C_\nabla:=C_\nabla^0\cdot\max\bigl\{1,\;30\,C_{\mathrm{reg}}(1-\beta_0)^{-1}C_\Gamma''\bigr\},
\end{equation}
with $C_{\mathrm{reg}}$, $\beta_0$, $C_\Gamma''$ the kernel and regularity constants. The second
factor depends on the dimension, $L$ and $M$ only, and $C_\nabla^0\le C_\nabla$. Suppose finally the
one-sided condition
\begin{equation}\label{8.5:eq-covariance-difference-a}
\eta+\mathfrak{r}\;\le\;\frac{4C_\nabla^0\,\ell}{\mathsf{d}}\;=\;\frac{4C_\nabla^0}{C_W}.
\end{equation}

Let $p,p'$ be bare plaquettes with $p'$ a translate of $p$ (same plane and orientation), and let $q$
be a bare plaquette. Then the first line of the display below holds whenever $p,p'\in\mathfrak{A}$,
$\operatorname{dist}(x_q,\mathfrak{A})\ge\mathsf{d}/4$ and
\begin{equation*}
|x_p-x_q|\le(A_c+3)\mathsf{d},\qquad |x_{p'}-x_q|\le(A_c+3)\mathsf{d}.
\end{equation*}
This is the case in particular for every $q$ with $x_q\in B_2$: the distance condition holds since
$|x_1-x_2|\ge3\mathsf{d}$ and $3-\tfrac{3}{2}-\tfrac{5}{4}=\tfrac{1}{4}$, and the upper bounds hold since
$|x_1-x_2|\le A_c\mathsf{d}$ gives $|x_p-x_q|\le(A_c+11/4)\mathsf{d}\le(A_c+3)\mathsf{d}$, and the same
for $p'$. The second line
holds whenever $p,p'\in\mathfrak{A}'$ and $x_q\in B_2$:
\begin{equation}\label{8.5:eq-covariance-difference-b}
\begin{aligned}
|\mathcal{C}_{pq}-\mathcal{C}_{p'q}|
&\le C_\nabla\,\sigma_{\mathsf{d}}\,\frac{|x_p-x_{p'}|+4\ell}{\mathsf{d}},\\
|\mathcal{C}_{pq}-\mathcal{C}_{p'q}|
&\le 32\,C_\nabla^0\,\sigma_{\mathsf{d}}\,\frac{|x_p-x_{p'}|+4\ell}{\mathsf{d}}.
\end{aligned}
\end{equation}
For the second line, let $\eta'$ and $\mathfrak{r}'$ be the relative errors obtained from the same two
proofs read at minimal pair distance $\mathsf{d}/8$ (that is, with $a_*=1/8$, where $a_*$ is the
minimal pair distance in units of $\mathsf{d}$) and with upper pair factor $A_c+25/8$: they bound
$|\rho_{pq}|$ by $\eta'\sigma_{pq}$ and $|C^{\mathrm{fl}}_{pq}|$ by $\mathfrak{r}'\sigma_{pq}$ for
$\mathsf{d}/8\le|x_p-x_q|\le(A_c+25/8)\mathsf{d}$. The corresponding three floor members at
$a_*=1/8$ are assumed, and so is
\begin{equation}\label{8.5:eq-covariance-difference-5}
\eta'+\mathfrak{r}'\;\le\;\frac{4C_\nabla^0}{C_W}.
\end{equation}

Finally, for the quantities $\eta$ and $\mathfrak{r}$ of the first line, let
$C_\nabla^{\mathrm{dec}}$ be the least number such that
\begin{equation*}
\mathfrak{r}(C_W)\,C_W\le 2C_\nabla^0\qquad\text{for all } C_W\ge C_\nabla^{\mathrm{dec}}.
\end{equation*}
Then, when $C_W\ge1$, the floor condition \eqref{8.5:eq-covariance-difference-a} holds as soon as
\begin{equation}\label{8.5:eq-covariance-difference-2}
C_W\ge C_\nabla^{\mathrm{dec}},
\end{equation}
\begin{equation}\label{8.5:eq-covariance-difference-3}
L^{2(K-s)}\ge 16\,C_{\mathrm{Green}}/C_\nabla^0,
\end{equation}
\begin{equation}\label{8.5:eq-covariance-difference-4}
L^{4(m+s)}\ge C_T\pi^2(A_c+3)^4C_W^5/C_\nabla^0.
\end{equation}
\end{lemma}

\begin{proof}
Fix $p,p'$ and $q$ as in the first line of \eqref{8.5:eq-covariance-difference-b} and write
$\Delta:=|x_p-x_{p'}|$. The argument has four steps.

\emph{Step 1.} By hypothesis~(a),
\begin{equation*}
\mathcal{C}_{pq}-\mathcal{C}_{p'q}
=\bigl(G_{pq}-G_{p'q}\bigr)-\bigl(C^{\mathrm{fl}}_{pq}-C^{\mathrm{fl}}_{p'q}\bigr).
\end{equation*}

\emph{Step 2.} The points $x_p$ and $x_{p'}$ lie in $\mathfrak{A}$ and
$\operatorname{dist}(x_q,\mathfrak{A})\ge\mathsf{d}/4$. Hence $|x_p-x_q|\ge\mathsf{d}/4$ and
$|x_{p'}-x_q|\ge\mathsf{d}/4$; together with the hypothesis $|x_p-x_q|,|x_{p'}-x_q|\le(A_c+3)\mathsf{d}$,
both pairs lie in the pair-distance range on which
Lemma~\ref{8.5:lem-covariance-upper} re-runs the two proofs named in the statement. As in that
lemma, the proof of the comparison of the free kernel with the continuum kernel gives
\begin{equation*}
G_{pq}=\varepsilon^4D_{\pi(p)\pi(q)}(x_p-x_q)+\rho_{pq},\qquad |\rho_{pq}|\le\eta\,\sigma_{pq}.
\end{equation*}
Since $t\mapsto\sigma(t)$ is decreasing and $|x_p-x_q|\ge\mathsf{d}/4$, we have
$\sigma_{pq}\le\sigma(\mathsf{d}/4)=\sigma_{\mathsf{d}}/2$, hence
$|\rho_{pq}|\le\eta\,\sigma_{\mathsf{d}}/2$. The proof of the decay lemma for the fluctuation kernel
gives, in the same way,
\begin{equation*}
|C^{\mathrm{fl}}_{pq}|\le\mathfrak{r}\,\sigma_{pq}\le\mathfrak{r}\,\sigma_{\mathsf{d}}/2.
\end{equation*}
The same statements hold with $p'$ in place of $p$.

\emph{Step 3.} Since $p'$ is a translate of $p$, the two plaquettes have a common plane. Put
$P:=\pi(p)=\pi(p')$ and $Q:=\pi(q)$. Then the continuum terms of $G_{pq}$ and $G_{p'q}$ are values of
the same entry $D_{PQ}$. For $t\in[0,1]$ put
\begin{equation*}
z(t):=(1-t)(x_p-x_q)+t(x_{p'}-x_q)=y(t)-x_q,\qquad y(t):=(1-t)x_p+t\,x_{p'}.
\end{equation*}
The ball $\mathfrak{A}$ is convex, so $y(t)\in\mathfrak{A}$, and therefore
$|z(t)|\ge\operatorname{dist}(x_q,\mathfrak{A})\ge\mathsf{d}/4$. The segment
$[x_p-x_q,\,x_{p'}-x_q]$ thus stays at distance at least $\mathsf{d}/4$ from $0$, where $D_{PQ}$ is
smooth by hypothesis~(d). The mean value inequality along this segment, together with
$|\nabla D_{PQ}(z)|\le C_D|z|^{-5}$, gives
\begin{equation*}
\bigl|\varepsilon^4D_{PQ}(x_p-x_q)-\varepsilon^4D_{PQ}(x_{p'}-x_q)\bigr|
\le\varepsilon^4C_D\Bigl(\frac{\mathsf{d}}{4}\Bigr)^{-5}\Delta.
\end{equation*}
To express the right-hand side in terms of $\sigma_{\mathsf{d}}$, use
$\sigma_{\mathsf{d}}=2\cdot4^4\varepsilon^4\pi^{-2}\mathsf{d}^{-4}$, that is
$4^4\varepsilon^4\mathsf{d}^{-4}=\pi^2\sigma_{\mathsf{d}}/2$. This gives
\begin{equation*}
\varepsilon^4C_D\Bigl(\frac{\mathsf{d}}{4}\Bigr)^{-5}\Delta
=4^5C_D\,\varepsilon^4\mathsf{d}^{-5}\Delta
=2\pi^2C_D\,\frac{\sigma_{\mathsf{d}}\,\Delta}{\mathsf{d}}.
\end{equation*}

\emph{Step 4.} Insert Step 2 into Step 1 and estimate the difference of the continuum terms by
Step 3. The four value errors $\rho_{pq}$, $\rho_{p'q}$, $C^{\mathrm{fl}}_{pq}$, $C^{\mathrm{fl}}_{p'q}$
contribute at most $2(\eta+\mathfrak{r})\sigma_{\mathsf{d}}/2$. Hence
\begin{equation*}
|\mathcal{C}_{pq}-\mathcal{C}_{p'q}|
\le 2\pi^2C_D\,\frac{\sigma_{\mathsf{d}}\Delta}{\mathsf{d}}+(\eta+\mathfrak{r})\,\sigma_{\mathsf{d}}.
\end{equation*}
By \eqref{8.5:eq-covariance-difference-1} we have $2\pi^2C_D\le C_\nabla^0$. By
\eqref{8.5:eq-covariance-difference-a} we have
$(\eta+\mathfrak{r})\sigma_{\mathsf{d}}\le 4C_\nabla^0\ell\,\sigma_{\mathsf{d}}/\mathsf{d}$. Therefore
\begin{equation*}
|\mathcal{C}_{pq}-\mathcal{C}_{p'q}|
\le C_\nabla^0\,\sigma_{\mathsf{d}}\,\frac{\Delta+4\ell}{\mathsf{d}}
\le C_\nabla\,\sigma_{\mathsf{d}}\,\frac{\Delta+4\ell}{\mathsf{d}},
\end{equation*}
the last inequality because $C_\nabla^0\le C_\nabla$. This is the first line of
\eqref{8.5:eq-covariance-difference-b}.

\emph{The enlarged ball.} Now let $p,p'\in\mathfrak{A}'$ be translates of each other and let
$x_q\in B_2$. Since $|x_1-x_2|\ge3\mathsf{d}$ and $3-\tfrac{13}{8}-\tfrac{5}{4}=\tfrac{1}{8}$, every
point of $\mathfrak{A}'$ is at distance at least $\mathsf{d}/8$ from $B_2$. So
$\operatorname{dist}(x_q,\mathfrak{A}')\ge\mathsf{d}/8$, and $\mathfrak{A}'$ is convex. We run Steps
1--4 again with $\mathfrak{A}'$ in place of $\mathfrak{A}$ and $\mathsf{d}/8$ in place of
$\mathsf{d}/4$. Since $|x_1-x_2|\le A_c\mathsf{d}$, we also have
$|x_p-x_q|\le(A_c+23/8)\mathsf{d}\le(A_c+25/8)\mathsf{d}$, and the same for $p'$. In Step 2 the
comparison and the decay bound are the re-readings at $a_*=1/8$ with
pair factor $A_c+25/8$ described in the statement, under their floor members at $a_*=1/8$, with
relative errors $\eta'$ and $\mathfrak{r}'$.

In Step 3 the segment now stays at distance at least $\mathsf{d}/8$ from $0$. Since
$(\mathsf{d}/8)^{-5}=32\,(\mathsf{d}/4)^{-5}$, the continuum difference is at most
\begin{equation*}
32\cdot2\pi^2C_D\,\frac{\sigma_{\mathsf{d}}\Delta}{\mathsf{d}}
\le 32\,C_\nabla^0\,\frac{\sigma_{\mathsf{d}}\Delta}{\mathsf{d}}.
\end{equation*}
In Step 4 we now have $\sigma_{pq},\sigma_{p'q}\le\sigma(\mathsf{d}/8)$ in place of
$\sigma(\mathsf{d}/4)$, so each of the four value errors is at most $\eta'\,\sigma(\mathsf{d}/8)$ or
$\mathfrak{r}'\,\sigma(\mathsf{d}/8)$. Since
$\sigma(\mathsf{d}/8)=16\,\sigma(\mathsf{d}/4)=8\,\sigma_{\mathsf{d}}$, their sum is at most
\begin{equation*}
2(\eta'+\mathfrak{r}')\,\sigma(\mathsf{d}/8)=16(\eta'+\mathfrak{r}')\,\sigma_{\mathsf{d}}
\le 64\,C_\nabla^0\,\frac{\sigma_{\mathsf{d}}\,\ell}{\mathsf{d}},
\end{equation*}
the inequality by \eqref{8.5:eq-covariance-difference-5}. Adding the two contributions,
\begin{equation*}
|\mathcal{C}_{pq}-\mathcal{C}_{p'q}|
\le 32\,C_\nabla^0\,\sigma_{\mathsf{d}}\,\frac{\Delta+2\ell}{\mathsf{d}},
\end{equation*}
which is at most the second line of \eqref{8.5:eq-covariance-difference-b}.

\emph{The floor condition \eqref{8.5:eq-covariance-difference-a} follows from
\eqref{8.5:eq-covariance-difference-2}--\eqref{8.5:eq-covariance-difference-4}.} The condition is
one-sided. Of the three conditions, \eqref{8.5:eq-covariance-difference-2} and
\eqref{8.5:eq-covariance-difference-4} are among the floors displayed in
Lemma~\ref{8.5:lem-covariance-upper}; \eqref{8.5:eq-covariance-difference-3} is a further one-sided
lower bound on $K-s$.
\begin{enumerate}
\item[(1)] The quantity $\mathfrak{r}(C_W)C_W$ is exponentially small in $C_W$, so the number
$C_\nabla^{\mathrm{dec}}$ of the statement exists, and by \eqref{8.5:eq-covariance-difference-2} and
the definition of $C_\nabla^{\mathrm{dec}}$ we have $\mathfrak{r}(C_W)\,C_W\le 2C_\nabla^0$.
The condition \eqref{8.5:eq-covariance-difference-2} is one more member of the floor on $C_W$, of the
same kind as the floor $C_W\ge C_0^{\mathrm{dec}}(1/4)$ of Lemma~\ref{8.5:lem-covariance-upper}.
\item[(2)] For the first term of $\eta$ we require
\begin{equation*}
16C_{\mathrm{Green}}C_W^{-2}L^{-2(K-s)}\cdot C_W\le C_\nabla^0.
\end{equation*}
Since $C_W\ge1$, this follows from
\eqref{8.5:eq-covariance-difference-3}, which is a one-sided member of the floor on $K-s$.
\item[(3)] For the torus term of $\eta$ we require
\begin{equation*}
C_T\pi^2(A_c+3)^4C_W^4L^{-4(m+s)}\cdot C_W\le C_\nabla^0,
\end{equation*}
which is equivalent to \eqref{8.5:eq-covariance-difference-4}.
This is a one-sided member of the torus clause of the quantifier prefix
(Notation~\ref{2.1:not-prefix}), read after $C_0$. At fixed $(m,s)$ it bounds $C_W$ from above.
Once $C_W$ is taken in the window $[C_0,C_0L)$ in which the two-point witness theorem reads it, it
is implied by the same inequality with $C_W$ replaced by $C_0L$, which is a floor on $m+s$ alone. No
upper bound on any parameter is imposed.
\end{enumerate}
Under (1)--(3), and since $\eta$ is at most the sum of its two terms,
$(\eta+\mathfrak{r})C_W\le C_\nabla^0+C_\nabla^0+2C_\nabla^0=4C_\nabla^0$. This is
\eqref{8.5:eq-covariance-difference-a}.
\end{proof}

\begin{remark}
We record four comments on the scope of Lemma~\ref{8.5:lem-covariance-difference}.

\emph{Only translates.} For plaquettes of different planes the bound
\eqref{8.5:eq-covariance-difference-b} would be false. The entry $\mathcal{C}_{pq}$ changes with the
plane of $p$ by an amount of order $\sigma$, not of order $\sigma|x_p-x_{p'}|/\mathsf{d}$. The
second-derivative room lemma compares $\mathcal{C}_{pq}$ only with $\mathcal{C}_{p'q}$ for a reference
plaquette $p'$ of the plane $P=\pi(p)$ in a cell. It also compares its cell quantities plane by
plane. Only translates therefore occur.

\emph{Consequences in the form used by the second-derivative room lemma.} Two instances of the first
line of \eqref{8.5:eq-covariance-difference-b} are used there.
\begin{itemize}
\item A displacement $|x_p-x_{p'}|\le 2(r+1)\ell$ with $r\ge1$ gives the relative bound
\begin{equation*}
C_\nabla\,\frac{2r+6}{C_W}\le \frac{8C_\nabla r}{C_W}\le\theta_r:=\frac{16C_\nabla r}{C_W}.
\end{equation*}
\item A displacement at most $2M\ell$ gives
$C_\nabla(2M+4)/C_W\le 8MC_\nabla/C_W$.
\end{itemize}
The second factor in \eqref{8.5:eq-covariance-difference-1} enlarges the constant so that the same
letter $C_\nabla$ also dominates the regularity part of the directions paired in the regularity
bookkeeping of that lemma. The second line of \eqref{8.5:eq-covariance-difference-b} is used only
for that bookkeeping.

\emph{Why $\sigma_{\mathsf{d}}$ and not the kernel maximum.} Let
$\mathsf{S}_q:=\max\{|\mathcal{C}_{pq}|:p \text{ bare in } \mathfrak{A}\}$. Writing $\mathsf{S}_q$ in
place of $\sigma_{\mathsf{d}}$ on the right-hand side would require a lower bound on $\mathsf{S}_q$.
For the maximum over the six planes such a bound holds: the per-plane identity in the proof of the
first clause of the witness's kernel-form lemma, together with the value-closeness bounds
$|\rho_{pq}|\le\eta\,\sigma_{pq}$ and $|C^{\mathrm{fl}}_{pq}|\le\mathfrak{r}\,\sigma_{pq}$ recalled in
the statement of Lemma~\ref{8.5:lem-covariance-difference}, gives
$\mathsf{S}_q\ge(6^{-1/2}-\eta-\mathfrak{r})\,\sigma(|y-x_q|)$, where $y$ is the point of
$\mathfrak{A}$ nearest to $x_q$, up to a term of order $(\varepsilon/\mathsf{d})\,\sigma(|y-x_q|)$.
For an individual entry no such bound holds. No lower bound is used.

\emph{Below the cell scale.} For gradients at the bare resolution, $\varepsilon^{-1}\nabla\mathcal{C}$,
the additive term $(\eta+\mathfrak{r})\sigma_{\mathsf{d}}$ in the proof of the first line of
\eqref{8.5:eq-covariance-difference-b} cannot be absorbed. Such a bound
would need the derivative half of the propagator decay bounds for $\bar{M}_{k_0}$ and for the
propagators of the earlier steps, of the kind stated in \cite[(3.9)]{B-CMP109} (compare the H\"older
form \cite[(4.17)--(4.18)]{B-CMP109}). It would also need a corresponding form of the lattice Green's
function asymptotics for differences. The second-derivative room lemma uses nothing at that
resolution: it needs
$\theta_r=16C_\nabla r/C_W$ with $r\ge1$, and displacements $2M\ell$, both covered above. In
particular, a gradient bound
$|\varepsilon^{-1}\nabla_\mu\mathcal{C}_{pq}|\le C_\nabla\mathsf{S}_q/\operatorname{dist}(p,q)$ is not
used. Such a bound cannot come from a comparison of the gradient of $\varepsilon^4D$ with its modulus:
the entries of $D$ vanish on cones where their gradients do not. The comparison does hold with
$\sigma$ in place of the modulus.

The regularity of the paired directions is nevertheless used elsewhere. The directions that the
second-derivative room lemma pairs with the majorant must lie in the regularity class of the complex
configuration spaces
$\mathfrak{U}^c_k$ of \cite{B-CMP109}. That membership is supplied by the H\"older form of the fine
oscillations of the propagator decay bounds, together with the admissibility reading. It holds
under the floors of the regularity bookkeeping: $C_W\ge16C_\nabla$,
$C_W\ge C_{\mathrm{reg}}^{\mathrm{dec}}$, where $C_{\mathrm{reg}}^{\mathrm{dec}}$ is the regularity
floor constant of the bookkeeping of the second-derivative room lemma, and the members at $a_*=1/8$.
\end{remark}

\begin{lemma}[Growth of the background covariance in the explicit scale]
\label{8.5:lem-covariance-growth}
Let $\mathcal{C}=\mathcal{C}^{\mathrm{bg}}$ be the background flux covariance of the Gaussian frame at the witness level
(Definition~\ref{8.5:def-gaussian-frame}). Let $\sigma_{\mathsf{d}}$ be the explicit scale fixed in~\eqref{8.5:eq-gaussian-frame-a}.
Assume the following two hypotheses.
\begin{itemize}
\item[(a)] The Gaussian decay statement for the block-averaged free field holds at the levels $j\le K$.
\item[(b)] The same Gaussian statement holds for the iterated block average of the unit-lattice free field on the torus
$\mathbb{T}_m$, at the levels $K<j\le K+m$.
\end{itemize}
Then there is a number $C_\sigma$, depending only on the dimension $d=4$, on $L$, $M$ and on the kernel constants $C_\Gamma$,
$\delta_\Gamma$, $\ell_\Gamma$, such that for every bare plaquette $q$
\begin{equation}\label{8.5:eq-covariance-growth-a}
\sup_{p}\,\lvert\mathcal{C}_{pq}\rvert\;\le\;C_\sigma\,C_W^4\,\sigma_{\mathsf{d}} .
\end{equation}
Here the supremum runs over the bare plaquettes $p$ of $\mathbb{T}_m$. The constant $C_\sigma$ is given by~\eqref{8.5:eq-covariance-growth-3} below.
\end{lemma}

\begin{proof}
Throughout, $\varepsilon=L^{-K}$ is the bare spacing and $\ell=L^{-s}$ is the cell side of the level-$k_0$ lattice, with $k_0=K-s$.
The separation scale is $\mathsf{d}=C_W\ell$.

\emph{Reduction to the coincident values.}
By Definition~\ref{8.5:def-gaussian-frame}, for every colour $a$ we have $\mathcal{C}_{pq}=\operatorname{Cov}(\psi^a_p,\psi^a_q)$.
Here $(\psi^a_p)_p$ is a centred Gaussian family indexed by the bare plaquettes. Hence $\mathcal{C}$ is a covariance matrix.
The Cauchy--Schwarz inequality for covariances gives, for all bare plaquettes $p,q$,
\begin{equation*}
\lvert\mathcal{C}_{pq}\rvert\le(\mathcal{C}_{pp}\,\mathcal{C}_{qq})^{1/2}\le\sup_{p'}\mathcal{C}_{p'p'} .
\end{equation*}
It therefore suffices to bound the coincident values of $\mathcal{C}$. We show that there is a constant $C_{\mathrm{bg}}$,
depending only on $d$, $L$, $M$ and the kernel constants $C_\Gamma$, $\delta_\Gamma$, $\ell_\Gamma$, such that
\begin{equation}\label{8.5:eq-covariance-growth-1}
\sup_{p}\,\mathcal{C}_{pp}\;\le\;C_{\mathrm{bg}}\,\varepsilon^4\ell^{-4}.
\end{equation}
Only this diagonal bound is used; neither off-diagonal decay of $\mathcal{C}$ nor any estimate of the free lattice
Green's function $G$ is needed.

\emph{The coincident bound from the background tower.}
Consider the coincident-cell curvature variance $\mathcal{C}_{pp}$. It is expanded by the tower decomposition of the
background covariance into summands indexed by the levels $j>k_0$. The last summand is the torus-top summand, obtained by
conditioning down to the top level at which $\mathbb{T}_m$ is a single block.

The decay of the background tower summands rests on a per-step implication. From the Gaussian decay statement
at level $j$, it bounds the coincident value of the level-$j$ summand by $C_{\mathrm{step}}(\varepsilon L^{K-j+1})^4$, where
$C_{\mathrm{step}}$ depends only on $d$, $L$, $M$ and the kernel constants $C_\Gamma$, $\delta_\Gamma$, $\ell_\Gamma$. In the decay
bound for the tower summands this implication is applied at the levels $1\le j\le k_0$. We apply the same implication to the
summands $j>k_0$. Its input at the
levels $j\le K$ is hypothesis~(a), and its input at the levels $K<j\le K+m$ is hypothesis~(b).
Since $\varepsilon=L^{-K}$, the prefactor is
\begin{equation*}
(\varepsilon L^{K-j+1})^4=L^{-4(j-1)} .
\end{equation*}

For $j\ge k_0+1$ the prefactor satisfies $(\varepsilon L^{K-j+1})^4\le(\varepsilon/\ell)^4$. Indeed, since $L>1$,
\begin{equation*}
\varepsilon L^{K-j+1}=L^{-(j-1)}\le L^{-k_0}=L^{-K}L^{s}=\varepsilon/\ell .
\end{equation*}
Summing over the levels $j>k_0$ of the summands other than the torus-top summand, their coincident values add up to at most
\begin{equation*}
C_{\mathrm{step}}\sum_{j\ge k_0+1}L^{-4(j-1)}
=C_{\mathrm{step}}\,L^{-4k_0}\sum_{i\ge0}L^{-4i}
=C_{\mathrm{step}}\,(1-L^{-4})^{-1}\Bigl(\frac{\varepsilon}{\ell}\Bigr)^4 ,
\end{equation*}
where we used $L^{-k_0}=\varepsilon/\ell$. The constant on the right does not depend on $K$, $s$ or $m$.

The torus-top summand is treated in the same way. At the top level let $\bar{M}_{\mathrm{top}}$ be the linearised minimiser,
$\bar{Q}_{\mathrm{top}}$ the linearised block average and $v_{\mathrm{top}}$ the Gaussian variable conditioned there.
The summand is
$\bar{M}_{\mathrm{top}}\operatorname{Cov}(v_{\mathrm{top}})\bar{M}_{\mathrm{top}}^{*}$, and it lives on a torus whose number of
plaquettes is bounded by a constant multiple of $L^4$. Its middle factor is
$\operatorname{Cov}(v_{\mathrm{top}})=\bar{Q}_{\mathrm{top}}\bar{Q}_{\mathrm{top}}^{*}$, whose entries (equivalently, whose
operator norm) are bounded by a constant depending on $L$ only.
Hence the torus-top summand obeys a bound of the same shape as the per-step bound that hypotheses~(a) and~(b) furnish for
the other summands: its coincident value is at most $C_{\mathrm{top}}(\varepsilon L^{K-j+1})^4$, where $j$ now denotes the
index of the top level of the tower and $C_{\mathrm{top}}$ depends only on $L$ and on the quantities on which
$C_{\mathrm{step}}$ depends. This index satisfies $j\ge k_0+1$, so that by the computation above the prefactor is at most
$(\varepsilon/\ell)^4$.

The background tower thus bounds the coincident-cell curvature variance by $C_{\mathrm{bg}}\varepsilon^4\ell^{-4}$ with
\begin{equation*}
C_{\mathrm{bg}}:=C_{\mathrm{step}}\,(1-L^{-4})^{-1}+C_{\mathrm{top}} ,
\end{equation*}
which is~\eqref{8.5:eq-covariance-growth-1}.

\emph{Conversion to the explicit scale.}
By its definition in~\eqref{8.5:eq-gaussian-frame-a},
$\sigma_{\mathsf{d}}=2\cdot4^4\pi^{-2}\varepsilon^4(C_W\ell)^{-4}$. Hence
$C_W^4\sigma_{\mathsf{d}}=2\cdot4^4\pi^{-2}(\varepsilon/\ell)^4$. We put
\begin{equation}\label{8.5:eq-covariance-growth-2}
C^0_\sigma:=\frac{C_{\mathrm{bg}}\,\pi^2}{2\cdot4^4}.
\end{equation}
Combining the reduction to coincident values with~\eqref{8.5:eq-covariance-growth-1}, we obtain for every bare plaquette $q$
\begin{equation*}
\sup_{p}\lvert\mathcal{C}_{pq}\rvert
\le C_{\mathrm{bg}}\Bigl(\frac{\varepsilon}{\ell}\Bigr)^4
=\frac{C_{\mathrm{bg}}\,\pi^2}{2\cdot4^4}\,C_W^4\,\sigma_{\mathsf{d}}
=C^0_\sigma\,C_W^4\,\sigma_{\mathsf{d}} .
\end{equation*}

\emph{The constant.}
We define
\begin{equation}\label{8.5:eq-covariance-growth-3}
C_\sigma:=C^0_\sigma\cdot\max\bigl\{1,\;C_{\mathrm{reg}}(1-\beta_0)^{-1}C_\Gamma'\bigr\}.
\end{equation}
Here $C_{\mathrm{reg}}$ and $\beta_0<1$ are the regularity constants and $C_\Gamma'$ is a kernel constant of the background
kernel, each determined by $d$, $L$, $M$ and $C_\Gamma$, $\delta_\Gamma$, $\ell_\Gamma$.
The maximum is at least $1$, so $C^0_\sigma\le C_\sigma$. The last display therefore gives~\eqref{8.5:eq-covariance-growth-a}.

The enlargement by the factor $\max\{1,C_{\mathrm{reg}}(1-\beta_0)^{-1}C_\Gamma'\}$ is not used in the present bound.
It is made here so that the same letter $C_\sigma$ also serves where the regular-direction norm of the column
$\mathcal{C}_{\cdot q}$ restricted to a filled hull $\hat{X}$ is estimated; that estimate is not part of the present lemma.

The factors in~\eqref{8.5:eq-covariance-growth-2} and~\eqref{8.5:eq-covariance-growth-3} are fixed by $d$, $L$, $M$ and the
kernel constants $C_\Gamma$, $\delta_\Gamma$, $\ell_\Gamma$.

The bound~\eqref{8.5:eq-covariance-growth-a} is used only in the exponentially suppressed tail terms of the witness estimate.
\end{proof}

\begin{lemma}[The Wick identity for the marked covariance]\label{8.5:lem-wick-identity}
Let $Q_1$, $Q_2$, the background fluxes $\psi^a$, their covariance
$\mathcal{C}=\mathcal{C}^{\mathrm{bg}}$, the flux forms $Q_X$ and the weight $h_2$ be those of
the Gaussian frame at the witness level (Definition~\ref{8.5:def-gaussian-frame}), with the
normalisation~\eqref{8.5:eq-gaussian-frame-b}. Let $n:=N^2-1$ be the number of colours, and write
$\mathcal{C}_{\cdot q}$ for the column of $\mathcal{C}$ at the bare plaquette $q$. Then
\begin{equation}\label{8.5:eq-wick-identity-1}
\operatorname{Cov}(Q_1,Q_2)
=b_0\,g_{k_0}^4\sum_q h_2(q)\sum_X Q_X[\mathcal{C}_{\cdot q},\mathcal{C}_{\cdot q}],
\qquad b_0=\frac{n}{2}.
\end{equation}
\end{lemma}

\begin{proof}
All expectations are taken under the tree block law at level $k_0$ of
Definition~\ref{8.5:def-gaussian-frame}. We use the normalisation~\eqref{8.5:eq-gaussian-frame-b},
in which
\begin{equation*}
Q_1=\frac{g_{k_0}^2}{2}\sum_a\sum_X Q_X[\psi^a,\psi^a],\qquad
Q_2=\frac{g_{k_0}^2}{2}\sum_q h_2(q)\sum_a(\psi^a_q)^2 .
\end{equation*}

\emph{One colour.} Let $\psi$ be a centred Gaussian vector indexed by the bare plaquettes, with
covariance $\mathcal{C}$, and let $H$ be the matrix of the bilinear form $\sum_X Q_X$, so that
$\psi^{\top}H\psi=\sum_X Q_X[\psi,\psi]$. By Isserlis' theorem,
\begin{equation*}
\mathbb{E}[\psi_p\psi_{p'}\psi_q\psi_q]
=\mathcal{C}_{pp'}\mathcal{C}_{qq}+2\,\mathcal{C}_{pq}\mathcal{C}_{p'q}.
\end{equation*}
The first term equals $\mathbb{E}[\psi_p\psi_{p'}]\,\mathbb{E}[\psi_q^2]$. Subtracting it and
summing against $H_{pp'}$ therefore gives
\begin{equation}\label{8.5:eq-wick-identity-2}
\operatorname{Cov}\bigl(\psi^{\top}H\psi,(\psi_q)^2\bigr)
=2\sum_{p,p'}H_{pp'}\mathcal{C}_{pq}\mathcal{C}_{p'q}
=2\,(\mathcal{C}e_q)^{\top}H(\mathcal{C}e_q),
\end{equation}
where $e_q$ is the unit vector at $q$. Since $\mathcal{C}$ is a covariance, it is symmetric, so
$\mathcal{C}e_q=\mathcal{C}_{\cdot q}$. Hence the right side of~\eqref{8.5:eq-wick-identity-2}
equals $2\sum_X Q_X[\mathcal{C}_{\cdot q},\mathcal{C}_{\cdot q}]$.

\emph{Colours.} By Definition~\ref{8.5:def-gaussian-frame}, the $n$ colours $b^a$ are independent
centred Gaussian cochains. Each $\psi^a$ is a linear image of $b^a$, and
$\operatorname{Cov}(\psi^a_p,\psi^{a'}_q)=\delta^{aa'}\mathcal{C}_{pq}$. Hence the families
$(\psi^a)_a$ are independent, and each has covariance $\mathcal{C}$. In the double sum over the
colours $a$ of $Q_1$ and $a'$ of $Q_2$, the terms with $a\neq a'$ therefore have zero covariance.
Each of the $n$ diagonal terms is given by~\eqref{8.5:eq-wick-identity-2}.

\emph{Prefactors.} By bilinearity of the covariance, the two factors $g_{k_0}^2/2$, the factor $2$
of~\eqref{8.5:eq-wick-identity-2} and the factor $n$ from the colours combine to
\begin{equation*}
\Bigl(\frac{g_{k_0}^2}{2}\Bigr)^2\cdot 2n=\frac{n}{2}\,g_{k_0}^4=b_0\,g_{k_0}^4 .
\end{equation*}
Summing over $q$ with the weights $h_2(q)$ gives~\eqref{8.5:eq-wick-identity-1}.

This prefactor is the prefactor of the Gaussian main term of the two-point witness theorem. In
that theorem the covariance of the quadratic parts at level $k_0$ of the curvature observables
with weights $h_1$ and $h_2$ is $b_0$ times the fourth power of the coupling $g_{k_0}$ of the
witness level times $\mathcal{K}^{\mathrm{bg}}(h_1,h_2)$, where $\mathcal{K}^{\mathrm{bg}}$
denotes the background two-point kernel in which that theorem writes its main term.

Finally, Definition~\ref{8.5:def-gaussian-frame} gives $\mathcal{C}_{\cdot q}\in\mathcal{E}$, the
space of exact fluxes, since the range of $\mathcal{C}^{\mathrm{bg}}$ lies in $\mathcal{E}$. So
each $Q_X[\mathcal{C}_{\cdot q},\mathcal{C}_{\cdot q}]$ is the flux form evaluated on exact fluxes,
as provided by part~(b) of Lemma~\ref{8.5:lem-constant-flux-forms}.
\end{proof}

\begin{proposition}[The source-profile defect]\label{8.5:prop-profile-defect}
Assume the following: clauses (a) and (b) below of the statement of profile locality and analyticity
of sourced terms (Definition~\ref{8.5:def-profile-locality}, whose proof is incomplete at the step
marked $(\ast)$), for the terms $V^{(j),w}_X$ of the frozen sourced family at the witness level
(Definition~\ref{8.5:def-frozen-sourced-family}) that meet $B(x_1,2\mathsf{d})$, and the weighted
domination (c) below:
\begin{itemize}
\item[(a)] clause (a) of \eqref{8.5:eq-profile-locality-a}: $w\mapsto V^{(j),w}_X$ depends on
$w|_{X^+}$ only, where $X^+\subset\hat{X}$ has radius $r(X)$ cells about the reference point
$x_{c(X)}$ of the anchor cell $c(X)$;
\item[(b)] clause (b) of \eqref{8.5:eq-profile-locality-a}: these terms are analytic in $w$ on
$\{\sup|w|<r_w\}$, with $r_w>0$ independent of $K$, $s$, $m$ and of position, $f\mapsto H_X[f]$
is linear, and the regular-direction mass satisfies
$\mathfrak{m}^{\mathrm{reg}}_X[f]\le\hat{\mathfrak{m}}_X\sup_{X^+}|f|$ with the unweighted route
weight $\bar{\mathfrak{w}}=r_w^{-1}$, $\hat{\mathfrak{m}}_X$ being the majorant mass of the term
given by Ba{\l}aban's bounds, read on the complex configuration spaces $\mathfrak{U}^c_k$ of
\cite{B-CMP109};
\item[(c)] the weighted domination of the majorant masses: for the weighted route weight
$\bar{\mathfrak{w}}$, a fixed constant multiple of $\mathfrak{w}_H$ or of $\mathfrak{w}_W$,
$\mathfrak{m}^{\mathrm{reg}}_X[f]\le\hat{\mathfrak{m}}_X\sup_{X^+}|f|$.
\end{itemize}
Let $f$ be a real, possibly signed, Lipschitz function of position on the torus, the plaquette
weight being $p\mapsto f(x(p))$, independent of the plane of $p$. (For plaquette weights $(f^P)_P$
that depend on the plane, the corresponding reduction is to the anisotropic constant profile
$\sum_P f^P(x_{c(X)})\mathbf{1}_P$, where $\mathbf{1}_P$ is the profile equal to $1$ on plaquettes of
the plane $P$ and to $0$ on the others; this reduction is not used in this book.) Then for every
domain $X$ of such a term, anchored at the cell $c=c(X)$,
\begin{equation}\label{8.5:eq-profile-defect-1}
q_X[f]=f(x_c)\,q_X[\mathbf{1}]+e_X,\qquad
|e_X|\le 4\,r(X)\,\ell\,\operatorname{Lip}f\cdot\hat{\mathfrak{m}}_X .
\end{equation}
Summing over the terms anchored at $c$ within one of the two families $\operatorname{Irr}_1^E$,
$\operatorname{Irr}_1^R$, and putting $q(c)[\mathbf{1}]:=\sum_{c(X)=c}q_X[\mathbf{1}]$ and
$e(c):=\sum_{c(X)=c}e_X$,
\begin{equation}\label{8.5:eq-profile-defect-a}
\sum_{c(X)=c}q_X[f]=f(x_c)\,q(c)[\mathbf{1}]+e(c),\qquad
|e(c)|\le 4\,\rho_\sharp\,\ell\,\operatorname{Lip}f\cdot\mathfrak{a}_c ,
\end{equation}
where $\mathfrak{a}_c$ is $\mathfrak{a}^E_c$ or $\mathfrak{a}^R_c$ according to the family, and the
same route weight $\bar{\mathfrak{w}}$ enters $\hat{\mathfrak{m}}_X$ and $\mathfrak{a}_c$. Since
$\ell=L^{-s}$, the defect is of first order per cell: relative to the cell mass $\mathfrak{a}_c$
it is at most $4\rho_\sharp L^{-s}\operatorname{Lip}f$. This is the defect bound, with constant
$C_{\mathrm{ext}}=4$, that enters the two-point witness theorem.
\end{proposition}

\begin{proof}
Fix a domain $X$ of such a term, anchored at $c$. By hypothesis (b), the map
$f\mapsto H_X[f]$ is linear; the constant-flux form $q_X[\,\cdot\,]$ of the marked Hessian is
obtained from $H_X[\,\cdot\,]$ by a linear operation, so $f\mapsto q_X[f]$ is linear, and
\begin{equation*}
q_X[f]-f(x_c)\,q_X[\mathbf{1}]=q_X\bigl[f-f(x_c)\mathbf{1}\bigr].
\end{equation*}
Thus $e_X=q_X[f-f(x_c)\mathbf{1}]$. By hypothesis (a), the term
$V^{(j),w}_X$ depends on $w|_{X^+}$ only;
hence $q_X[\varphi]$ depends only on the restriction $\varphi|_{X^+}$ of the profile $\varphi$.

Put $\psi:=(f-f(x_c))|_{X^+}$ and $S:=\sup_{X^+}|f-f(x_c)|$. Let $\tilde\psi$ be the McShane
extension of $\psi$ to the torus,
\begin{equation*}
\tilde\psi(x):=\inf_{y\in X^+}\bigl(\psi(y)+\operatorname{Lip}f\cdot|x-y|\bigr),
\end{equation*}
with $|x-y|$ the distance on the torus. Since $\operatorname{Lip}\psi\le\operatorname{Lip}f$, the
function $\tilde\psi$ coincides with $\psi$ on $X^+$ and
$\operatorname{Lip}\tilde\psi\le\operatorname{Lip}f$.
Let $\tilde\varphi:=\max\{-S,\min\{S,\tilde\psi\}\}$ be $\tilde\psi$ clipped at $\pm S$. Clipping
does not increase the Lipschitz constant, and it does not change $\tilde\psi$ on $X^+$, where
$|\tilde\psi|=|\psi|\le S$. Therefore
\begin{equation*}
\tilde\varphi=f-f(x_c)\ \text{on }X^+,\qquad
\operatorname{Lip}\tilde\varphi\le\operatorname{Lip}f,\qquad
\sup|\tilde\varphi|\le S .
\end{equation*}
Since $X^+$ lies in the region of radius $r(X)$ cells, of side $\ell$, about $x_c$, every point of
$X^+$ lies within sup-norm distance $(r(X)+1)\ell$ of $x_c$, the additional cell accounting for the
position of $x_c$ within the anchor cell, and hence within Euclidean distance
$2(r(X)+1)\ell$, the Euclidean norm on $\mathbb{R}^4$ being at most twice the sup norm. As $f$ is
Lipschitz,
\begin{equation*}
\sup|\tilde\varphi|\le S\le 2\bigl(r(X)+1\bigr)\ell\operatorname{Lip}f
\le 4\,r(X)\,\ell\operatorname{Lip}f,
\end{equation*}
the last inequality holding because $r(X)\ge1$. The function $\tilde\varphi$ is a real Lipschitz
function of position on the torus, hence a profile of the class to which hypotheses (b) and (c)
apply. Because $\tilde\varphi$ and $f-f(x_c)\mathbf{1}$ agree on $X^+$, the dependence statement of
the first paragraph gives $e_X=q_X[\tilde\varphi]$.

We now bound $q_X[\tilde\varphi]$. With the unweighted route weight $\bar{\mathfrak{w}}=r_w^{-1}$,
hypothesis (b)
gives $\mathfrak{m}^{\mathrm{reg}}_X[\tilde\varphi]\le\hat{\mathfrak{m}}_X\sup_{X^+}|\tilde\varphi|$.
With the weighted route weight, a fixed constant multiple of $\mathfrak{w}_H$ or of
$\mathfrak{w}_W$, hypothesis (c) gives the same domination. Part (e) of
Lemma~\ref{8.5:lem-constant-flux-forms} bounds the constant-flux form by the regular-direction
mass, $|q_X[\tilde\varphi]|\le\mathfrak{m}^{\mathrm{reg}}_X[\tilde\varphi]$. Combining these two
steps with the bound on $\sup|\tilde\varphi|$ above,
\begin{align*}
|e_X|=|q_X[\tilde\varphi]|
&\le\mathfrak{m}^{\mathrm{reg}}_X[\tilde\varphi]
\le\hat{\mathfrak{m}}_X\sup_{X^+}|\tilde\varphi|\\
&\le 4\,r(X)\,\ell\operatorname{Lip}f\cdot\hat{\mathfrak{m}}_X ,
\end{align*}
which is \eqref{8.5:eq-profile-defect-1}.

Finally, sum \eqref{8.5:eq-profile-defect-1} over the domains $X$ of the terms of one family with
$c(X)=c$. The identity in \eqref{8.5:eq-profile-defect-a} follows from the definitions of
$q(c)[\mathbf{1}]$ and $e(c)$. For the bound,
\begin{equation*}
|e(c)|\le 4\,\ell\operatorname{Lip}f\sum_{c(X)=c}r(X)\,\hat{\mathfrak{m}}_X
\le 4\,\rho_\sharp\,\ell\operatorname{Lip}f\cdot\mathfrak{a}_c ,
\end{equation*}
since $\mathfrak{a}_c$ is the aggregate of the majorant masses $\hat{\mathfrak{m}}_X$ of the terms
of that family anchored at $c$, computed with the same route weight, and $\rho_\sharp$ bounds their
mass-weighted mean radius.
\end{proof}

\begin{lemma}[The profile-corrected weight and its tail]\label{8.5:lem-corrected-weight}
Let $k_0=K-s$ be the level of the two-point witness theorem, let $\ell=L^{-s}$ be the side of a
cell of the level-$k_0$ lattice, and let $\mathsf{d}=C_W\ell$. For a plaquette $p$, with
barycentre $x(p)$, and a bounded function $g$, put
\begin{equation}\label{8.5:eq-corrected-weight-1}
\mathcal{Z}_p[g]:=D[g]\,\zeta_{k_0}(p)[zg]
=\partial_z\big|_{z=0}\,\zeta_{k_0}(p)[zg],
\qquad
\zeta'_{k_0}(p):=\mathcal{Z}_p[\mathbf{1}] .
\end{equation}
For $x\in\mathbb{T}_m$ and $A\subset\mathbb{T}_m$ write $\operatorname{dist}(x,A)$ for the
distance from $x$ to $A$. Let $r_w>0$ be the radius of the analyticity domain in part (b) of
Definition~\ref{8.5:def-profile-locality}; it does not depend on $(K,s,m)$ or on $p$. Assume the
following.
\begin{itemize}
\item[(L)] The linear bound on the running shift stated with the correlation theorem. Here
$\zeta_{k_0}(p)[w]$ denotes the running shift of the correlation theorem at level $k_0$ and
plaquette $p$, as a function of the complex source $w$; for every plaquette $p$ and every $w$
with $\sup|w|<r_w$, $|\zeta_{k_0}(p)[w]|\le\tfrac43\sup|w|$.
\item[(P)] Part (a$'$) of the profile-locality statement, Definition~\ref{8.5:def-profile-locality},
stated there with proof incomplete at the step marked $(\ast)$, in the form stated there. For
every $R\ge1$,
$\zeta_{k_0}(p)=\zeta^{\le R}_{k_0}(p)+\zeta^{>R}_{k_0}(p)$, where
$\zeta^{\le R}_{k_0}(p)[w]$ depends on $w|_{B(x(p),R\ell)}$ only, and for every bounded $g$
\begin{equation}\label{8.5:eq-corrected-weight-2}
\begin{aligned}
\big|D[g]\zeta^{\le R}_{k_0}(p)-g(x(p))\,D[\mathbf{1}]\zeta^{\le R}_{k_0}(p)\big|
&\le\tfrac43\cdot2\rho_\sharp\ell\operatorname{Lip}g,\\
\big|D[g]\zeta^{>R}_{k_0}(p)\big|&\le C_w e^{-\hat\kappa R/(2M)}\sup|g| ,
\end{aligned}
\end{equation}
where $C_w\ge1$ is the constant of \eqref{8.5:eq-radius-tails-a}.
\item[(S)] Part (b) of the profile-locality statement, Definition~\ref{8.5:def-profile-locality},
stated there with proof incomplete at the step marked $(\ast)$. The functions $\zeta_{k_0}(p)$,
$\zeta^{\le R}_{k_0}(p)$ and $\zeta^{>R}_{k_0}(p)$ are analytic in $w$ on $\{\sup|w|<r_w\}$.
\item[(J)] As stated with the correlation theorem, $\zeta'_{k_0}(p)$ does not depend on $p$ and
coincides with the running shift $\zeta'_{k_0}$ of the correlation theorem, and
$\zeta'_{k_0}\ge\tfrac23$.
\item[(F)] $C_W\ge16$ and $C_W\ge C_1$. Here $C_1=C_1(\min\{\hat\kappa,\hat\kappa_R\},M,C_w',C_\sigma)$
is the least number such that
\begin{equation}\label{8.5:eq-corrected-weight-3}
C_w'\,e^{-\min\{\hat\kappa,\hat\kappa_R\}C/(64M)}\,(2+C_\sigma^2C^8)\,C\le1
\qquad\text{for all }C\ge C_1 ,
\end{equation}
with $C_w'\ge C_w\ge1$ and $C_\sigma$ the constant of \eqref{8.5:eq-covariance-growth-a}.
\end{itemize}
Let $x_1,x_2\in\mathbb{T}_m$ and, for $i=1,2$, let $f_i$ be a real-valued bounded Lipschitz
function with $\operatorname{supp}f_i\subset B(x_i,\mathsf{d})$. Define the \emph{profile-corrected
weight} $h_i(p):=\mathcal{Z}_p[f_i]/\zeta'_{k_0}$ and put $\delta f_i(p):=h_i(p)-f_i(x(p))$. Let
$B_i:=B(x_i,5\mathsf{d}/4)$, and split $h_i=h_i^{\mathrm{in}}+h_i^{\mathrm{out}}$ by setting
$h_i^{\mathrm{in}}(p):=h_i(p)$ if $x(p)\in B_i$ and $h_i^{\mathrm{in}}(p):=0$ otherwise. Then the
following hold.
\begin{itemize}
\item[(i)] For every $p$ and every bounded $g$, $|\mathcal{Z}_p[g]|\le\tfrac43\sup|g|$.
\item[(ii)] For every $p$,
\begin{equation}\label{8.5:eq-corrected-weight-4}
|\delta f_i(p)|\le 4\frac{\rho_\sharp}{C_W}\,\mathsf{d}\operatorname{Lip}f_i
+3C_we^{-\hat\kappa C_W/(32M)}\|f_i\|_\infty ,
\end{equation}
and the second summand is at most $C_W^{-2}\|f_i\|_\infty$.
\item[(iii)] If $t:=\operatorname{dist}(x(p),\operatorname{supp}f_i)>2\ell$, then
\begin{equation}\label{8.5:eq-corrected-weight-5}
|\delta f_i(p)|\le\tfrac32\,C_we^{-\hat\kappa(t/\ell-2)/(2M)}\|f_i\|_\infty .
\end{equation}
\item[(iv)] $\operatorname{supp}h_i^{\mathrm{in}}\subset B_i$, and for every $p$
\begin{equation}\label{8.5:eq-corrected-weight-a}
\begin{aligned}
|h_i(p)|&\le2\|f_i\|_\infty,\\
|h_i^{\mathrm{out}}(p)|&\le
3C_we^{-\hat\kappa(\operatorname{dist}(x(p),B(x_i,\mathsf{d}))/\ell-2)/(2M)}\|f_i\|_\infty .
\end{aligned}
\end{equation}
\end{itemize}
\end{lemma}

\begin{proof}
\emph{(i).} Fix $p$ and a bounded $g\ne0$. By (S), $z\mapsto\zeta_{k_0}(p)[zg]$ is analytic for
$|z|<r_w/\sup|g|$. By (L) with $w=0$, $\zeta_{k_0}(p)[0]=0$. Hence
\[
\mathcal{Z}_p[g]=\lim_{z\to0}\frac{\zeta_{k_0}(p)[zg]}{z}.
\]
By (L) with $w=zg$, each quotient has modulus at most $\tfrac43\sup|g|$, so the limit does too.
The bound uses only (L) and (S). Combining (i) for $g=f_i$ with (J) gives
\[
|h_i(p)|=\frac{|\mathcal{Z}_p[f_i]|}{\zeta'_{k_0}}
\le\frac43\cdot\frac32\,\|f_i\|_\infty=2\|f_i\|_\infty ,
\]
which is the first line of \eqref{8.5:eq-corrected-weight-a}.

\emph{(ii).} Let $f$ be a real-valued bounded function with $\operatorname{Lip}f<\infty$.
Take $R:=C_W/16$ in (P); $R\ge1$ by (F). The decomposition of (P) holds as an
identity of functions of $w$, and $D[\cdot]$ is a derivative. Hence
$\mathcal{Z}_p[f]=D[f]\zeta^{\le R}_{k_0}(p)+D[f]\zeta^{>R}_{k_0}(p)$, and in particular
$\zeta'_{k_0}=D[\mathbf{1}]\zeta^{\le R}_{k_0}(p)+D[\mathbf{1}]\zeta^{>R}_{k_0}(p)$. Therefore
\begin{multline*}
\mathcal{Z}_p[f]-f(x(p))\zeta'_{k_0}
=\big(D[f]\zeta^{\le R}_{k_0}(p)-f(x(p))D[\mathbf{1}]\zeta^{\le R}_{k_0}(p)\big)\\
+D[f]\zeta^{>R}_{k_0}(p)-f(x(p))\,D[\mathbf{1}]\zeta^{>R}_{k_0}(p).
\end{multline*}
The function $\zeta^{\le R}_{k_0}(p)$ depends on $w$ only through $w|_{B(x(p),R\ell)}$, and $x(p)$
lies in this ball. Write $B:=B(x(p),R\ell)$ and let $\tilde f$ be the clipped extension of $f|_B$:
the function
\[
y\mapsto\inf_{x\in B}\big(f(x)+\operatorname{Lip}f\cdot\operatorname{dist}(x,y)\big),
\]
truncated to the interval $[\inf_Bf,\sup_Bf]$. The infimum is finite because $f$ is bounded; it
equals $f(y)$ for $y\in B$, since $f(x)+\operatorname{Lip}f\cdot\operatorname{dist}(x,y)\ge f(y)$
for $x,y\in B$ with equality at $x=y$; as an infimum of functions with Lipschitz constant
$\operatorname{Lip}f$ it has Lipschitz constant at most $\operatorname{Lip}f$, and truncation does not
increase the Lipschitz constant. Thus $\tilde f=f$ on $B$, $\operatorname{Lip}\tilde f\le
\operatorname{Lip}f$ and $\sup|\tilde f|\le\sup_B|f|$. Consequently
$D[\tilde f]\zeta^{\le R}_{k_0}(p)=D[f]\zeta^{\le R}_{k_0}(p)$ and $\tilde f(x(p))=f(x(p))$, so the
first bracket is unchanged if $f$ is replaced by $\tilde f$. By the first line of
\eqref{8.5:eq-corrected-weight-2} applied to $g=\tilde f$, the bracket is at most
$\tfrac43\cdot2\rho_\sharp\ell\operatorname{Lip}\tilde f\le\tfrac43\cdot2\rho_\sharp\ell\operatorname{Lip}f$.

By the second line of \eqref{8.5:eq-corrected-weight-2}, applied to $g=f$ and to
$g=\mathbf{1}$, and using $|f(x(p))|\le\|f\|_\infty$ and $\sup|\mathbf{1}|=1$, the last two terms
are each at most $C_we^{-\hat\kappa R/(2M)}\|f\|_\infty$. Since $\hat\kappa R/(2M)=\hat\kappa
C_W/(32M)$, we obtain
\[
\big|\mathcal{Z}_p[f]-f(x(p))\zeta'_{k_0}\big|
\le\tfrac43\cdot2\rho_\sharp\ell\operatorname{Lip}f
+2C_we^{-\hat\kappa C_W/(32M)}\|f\|_\infty .
\]
With $f=f_i$, we have $\delta f_i(p)=\big(\mathcal{Z}_p[f_i]-f_i(x(p))\zeta'_{k_0}\big)/\zeta'_{k_0}$.
Dividing by $\zeta'_{k_0}\ge\tfrac23$ (J) multiplies the right side by at most $\tfrac32$. Since
$\rho_\sharp\ell=(\rho_\sharp/C_W)\mathsf{d}$, this gives \eqref{8.5:eq-corrected-weight-4}.

For the second summand, $\hat\kappa\ge\min\{\hat\kappa,\hat\kappa_R\}$ gives
\[
e^{-\hat\kappa C_W/(32M)}\le\big(e^{-\min\{\hat\kappa,\hat\kappa_R\}C_W/(64M)}\big)^2 .
\]
Since $C_W\ge C_1$, \eqref{8.5:eq-corrected-weight-3} at $C=C_W$ bounds the inner exponential by
$\big(C_w'(2+C_\sigma^2C_W^8)C_W\big)^{-1}$. With $C_w'\ge C_w\ge1$ this yields
\[
C_we^{-\hat\kappa C_W/(32M)}
\le\frac{C_w}{C_w'^{\,2}}\,(2+C_\sigma^2C_W^8)^{-2}C_W^{-2}
\le C_w^{-1}(2+C_\sigma^2C_W^8)^{-2}C_W^{-2}.
\]
Since $(2+C_\sigma^2C_W^8)^{-2}\le\tfrac14$ and $C_w\ge1$, three times the right side is at most
$\tfrac34C_W^{-2}\le C_W^{-2}$. In the second-derivative room lemma (the lemma with floors
\eqref{8.5:eq-room-lemma-a}), this summand is absorbed into the constant of the first summand in
the same way as the corresponding exponentially small terms are absorbed there.

\emph{(iii).} Property (a$'$) in (P) holds for every $R\ge1$, so we may choose $R$ depending on
$p$. Let $t=\operatorname{dist}(x(p),\operatorname{supp}f_i)>2\ell$ and $R:=\lfloor t/\ell\rfloor-1$.
Then $\lfloor t/\ell\rfloor\ge2$, so $R\ge1$. Moreover $R\ell\le t-\ell<t$, hence $f_i\equiv0$ on
$B(x(p),R\ell)$.

Because $\zeta^{\le R}_{k_0}(p)[w]$ depends on $w|_{B(x(p),R\ell)}$ only, the map
$z\mapsto\zeta^{\le R}_{k_0}(p)[zf_i]$ is constant, equal to $\zeta^{\le R}_{k_0}(p)[0]$. Thus
$D[f_i]\zeta^{\le R}_{k_0}(p)=0$, and
$\mathcal{Z}_p[f_i]=D[f_i]\zeta^{>R}_{k_0}(p)$. By the second line of
\eqref{8.5:eq-corrected-weight-2},
\[
|\mathcal{Z}_p[f_i]|\le C_we^{-\hat\kappa R/(2M)}\|f_i\|_\infty
\le C_we^{-\hat\kappa(t/\ell-2)/(2M)}\|f_i\|_\infty ,
\]
using $R\ge t/\ell-2$. Since $t>0$, $f_i(x(p))=0$, so $\delta f_i(p)=h_i(p)=\mathcal{Z}_p[f_i]/\zeta'_{k_0}$.
Dividing by $\zeta'_{k_0}\ge\tfrac23$ gives \eqref{8.5:eq-corrected-weight-5}.

\emph{(iv).} By definition $h_i^{\mathrm{in}}$ vanishes at every $p$ with $x(p)\notin B_i$, and
$h_i^{\mathrm{out}}(p)=0$ whenever $x(p)\in B_i$. Let $x(p)\notin B_i$. Then, by (F),
\[
\operatorname{dist}(x(p),B(x_i,\mathsf{d}))\ge\tfrac14\mathsf{d}=\tfrac14C_W\ell\ge4\ell .
\]
Since
$\operatorname{supp}f_i\subset B(x_i,\mathsf{d})$, the number $t$ of (iii) satisfies
$t\ge\operatorname{dist}(x(p),B(x_i,\mathsf{d}))>2\ell$. Moreover $f_i(x(p))=0$, so
$h_i^{\mathrm{out}}(p)=h_i(p)=\delta f_i(p)$. By \eqref{8.5:eq-corrected-weight-5}, together with the
monotonicity of the exponential in $t$,
\[
|h_i^{\mathrm{out}}(p)|
\le\tfrac32\,C_we^{-\hat\kappa(\operatorname{dist}(x(p),B(x_i,\mathsf{d}))/\ell-2)/(2M)}
\|f_i\|_\infty ,
\]
which implies the second line of \eqref{8.5:eq-corrected-weight-a}. The first line was proved
in (i).
\end{proof}

\begin{remark}
The bound \eqref{8.5:eq-corrected-weight-a} exhibits $h_i^{\mathrm{out}}$ as a tail that is small in
$L^1$ and decays with the distance from $B(x_i,\mathsf{d})$. Its $L^1$ mass, that is, the integral
over $x\in\mathbb{T}_m\setminus B_i$ of the majorant in the second line of
\eqref{8.5:eq-corrected-weight-a} with $x(p)$ replaced by $x$, is at most
\[
C(M,\hat\kappa)\,C_w\,\ell\,|\partial B_i|\,\|f_i\|_\infty\,e^{-\hat\kappa C_W/(8M)},
\]
where $|\partial B_i|$ denotes the three-dimensional surface measure of the sphere bounding $B_i$
and the constant $C(M,\hat\kappa)$ depends on $M$ and $\hat\kappa$ only; in particular the bound
does not depend on $m$ or $K$. In the second-derivative room lemma (the lemma with floors
\eqref{8.5:eq-room-lemma-a}) this tail is bounded using \eqref{8.5:eq-covariance-growth-a} under the
floor \eqref{8.5:eq-room-tail-a}. A bound of the same
size but independent of $p$ would, summed over $\mathbb{T}_m$, grow with $m$ and $K$; the decay in
the distance makes the sum independent of both.

Only $h_2$ enters the two-point witness theorem; the profile attached there to $f_1$ is $f_1$ itself,
which has no tail. Since $\operatorname{supp}f_i\subset B(x_i,\mathsf{d})\subset B_i$, one has
$f_i(x(p))=0$ whenever $x(p)\notin B_i$; hence
\[
\delta f_i(p)=\big(h_i^{\mathrm{in}}(p)-f_i(x(p))\big)+h_i^{\mathrm{out}}(p),
\]
where the first summand, as a function of $p$, vanishes unless $x(p)\in B_i$, and the second is
bounded by the second line of \eqref{8.5:eq-corrected-weight-a}. In particular the inclusion
$\operatorname{supp}\delta f_i\subset B_i$ holds only up to $h_i^{\mathrm{out}}$.

Hypothesis (S) is analyticity on a ball in the sup norm of $w$, matching the sup norm in (L).
Write $C_{\mathrm{ext}}$ for the constant controlling the extension $\tilde f$ of step (ii) of the
proof relative to $f$ in the norm of the analyticity domain. Suppose instead the analyticity
domain were a ball in the norm $\|\cdot\|_\sharp$ of Definition~\ref{1.3:def-curvature-field},
with the weight $40M\cdot\operatorname{Lip}$ in bare units. Then $C_{\mathrm{ext}}$
would stay absolute. With a weight $M\cdot\operatorname{Lip}$ in cell units, $C_{\mathrm{ext}}$ would
be of order $M$, and the factor $C_W^{-1}$ of \eqref{8.5:eq-corrected-weight-4} would be kept. Only a
weight in units coarser than the cell would destroy uniformity in $s$, and the sup-norm ball of
(S) excludes this.
\end{remark}

\begin{remark}[The period of members on finer threshold cubes. Stated; proof incomplete at the step
marked $(\ast)$]\label{8.5:rem-threshold-period}
Let $k_0=K-s$ be the witness level, and measure lengths in cells of the level-$k_0$ lattice, of
side $\ell=L^{-s}$, as in Notation~\ref{8.5:not-partition-lattice}. Consider those members of the
frozen sourced family of Definition~\ref{8.5:def-frozen-sourced-family} that come from the
expansion of \cite{B-CMP119} and are localised on cubes of side $R_jM$ in level-$j$ units. Here
$R_j$ is the radius of \cite{B-CMP119} at level $j$, a polylogarithmic function of
$x_j=g_j^{-2}$. Put
\begin{equation}\label{8.5:eq-threshold-period-1}
\Lambda:=\operatorname{lcm}\bigl\{\,R_jML^{-(k_0-j)}\;:\;j\le k_0,\ L^{k_0-j}\le R_jM\,\bigr\}.
\end{equation}
\begin{enumerate}
\item[(a)] For these members, the mass-weighted mean radius $\rho_\sharp$ becomes
$\rho_\sharp(g_{K-s})$. The partition lattice $\Gamma=M\mathbb{Z}^4$ of
Notation~\ref{8.5:not-partition-lattice} becomes $\Lambda\mathbb{Z}^4$ (in cells), with $\Lambda$
given by \eqref{8.5:eq-threshold-period-1} under the hypothesis of (b) below, and as described in
(c) otherwise. Accordingly:
  \begin{itemize}
  \item the entries $8\rho_\sharp+8$ and $16M$ of the floor $C_0$ of $C_W$ are read as
  $8\rho_\sharp(g_{K-s})+8$ and $16\Lambda$;
  \item the factor $M\rho_\sharp$ of the room radius of the second-derivative room lemma is read
  as $\Lambda\rho_\sharp(g_{K-s})$.
  \end{itemize}
\item[(b)] Assume the \emph{power-of-$L$ hypothesis}: every $R_j$ is a power of $L$ (only the
levels $j\le k_0$ are used below). It is one of the hypotheses of the second-derivative room lemma
for these members. Then $\Lambda=R_{k_0}M$. The two floor entries listed in (a) are then one-sided
floors on $C_W$, and all three quantities of (a) read the witness coupling $g_{K-s}$ alone. Hence
they depend on $s$ only through $x_{K-s}$, and they depend neither on $K$ nor on $m$.
\item[(c)] Without the power-of-$L$ hypothesis, write $R_jM=L^{\nu_j}\mu_j$ with $\nu_j\ge0$
maximal, so that $\mu_j$ is the $L$-free part of $R_jM$. Then every level $j\le k_0$, including
the levels excluded from \eqref{8.5:eq-threshold-period-1}, imposes a cell period divisible by
$\mu_j$, once $L^{k_0-j}$ absorbs $L^{\nu_j}$. On this branch we write $\Lambda$ for the common
period of the cube partitions of all levels $j\le k_0$; it is a common multiple of every $\mu_j$,
$j\le k_0$, and so is at least $\operatorname{lcm}_{j\le k_0}\mu_j$. This $\Lambda$ can depend on
$K$ through the range of $x_j$, which extends up to $x_0$. For these members, uniformity in $K$ of
the bound of the second-derivative room lemma is not proved in this book.
\end{enumerate}
Stated; proof incomplete at the step marked $(\ast)$.
\end{remark}

\begin{proof}
\emph{The cube sides at level $k_0$.} The unit of the level-$j$ lattice is $\varepsilon L^{j}$,
and a cell has side $\ell=\varepsilon L^{k_0}$. A cube of side $R_jM$ in level-$j$ units
therefore has side
\begin{equation*}
R_jM\,\varepsilon L^{j}=R_jML^{-(k_0-j)}\,\ell ,
\end{equation*}
that is, $R_jML^{-(k_0-j)}$ cells. The condition $L^{k_0-j}\le R_jM$ in
\eqref{8.5:eq-threshold-period-1} says exactly that this side is at least one cell.

$(\ast)$ Three statements about the constructions of \cite{B-CMP119} are taken here as hypotheses
and are not proved in this book:
\begin{itemize}
\item the blocks that \cite{B-CMP119} introduces in addition to the cubes of the partitions
$\pi_j$ of Notation~\ref{8.5:not-partition-lattice} have side $M_1$, with $M_1$ dividing $M$;
\item the localisation domains of the $\beta$-extracted terms of \cite{B-CMP119} (the class
$\operatorname{Irr}_1^E$) are unions of cubes of side $M$;
\item the power-of-$L$ hypothesis of (b), together with the identification $\Lambda=R_{k_0}M$,
that is, that the largest of the sides in \eqref{8.5:eq-threshold-period-1} is the one at
$j=k_0$.
\end{itemize}
The first two are recorded for the period claim of (a): they ensure that the localisation domains
of these members are unions of cubes of side $M$, tiled by the blocks of side $M_1$, so that a
translation preserving the cube partitions considered below permutes the domains; the third is
used in (b).

\emph{The period and the radius.} Under the hypothesis of (b) the following holds; without it the
period is the one described in (c), treated below. Since $M=L^{m'}$, every $R_jM$ is then a power
of $L$. Each side $R_jML^{-(k_0-j)}$ with $L^{k_0-j}\le R_jM$ is therefore an integral power of
$L$ with non-negative exponent; in particular it is an integer number of cells. Each of the cube
partitions in question, seen in cells, is a grid partition: its cubes are the translates of one
cube by the integer multiples of its side in cells along the four coordinate axes, and every such
translation maps the partition onto itself.
Each side $R_jML^{-(k_0-j)}$ occurring in \eqref{8.5:eq-threshold-period-1} divides $\Lambda$.
Hence translations by $\Lambda\mathbb{Z}^4$ preserve every one of these cube partitions at once,
and $\Lambda\mathbb{Z}^4$ takes the place of the partition lattice $\Gamma=M\mathbb{Z}^4$ for these
members. The mass-weighted mean radius is read at the witness coupling, which gives
$\rho_\sharp(g_{K-s})$.

The entries $8\rho_\sharp+8$ and $16M$ of the floor $C_0$, and the factor $M\rho_\sharp$ of the
room radius, enter through $\rho_\sharp$ and through the period $M$ of $\Gamma$. Making the two
substitutions $\rho_\sharp\mapsto\rho_\sharp(g_{K-s})$ and $M\mapsto\Lambda$ gives the quantities
listed in (a).

\emph{Under the power-of-$L$ hypothesis.} As shown above, each side $R_jML^{-(k_0-j)}$ occurring
in \eqref{8.5:eq-threshold-period-1} is an integral power of $L$ with non-negative exponent. The
least common multiple of finitely many powers of $L$ is the largest of
them. That this largest side is the one at $j=k_0$, so that $\Lambda=R_{k_0}M$, is taken as part
of the step marked $(\ast)$.

The quantities of (a) are then
\begin{equation*}
8\rho_\sharp(g_{K-s})+8,\qquad 16R_{k_0}M,\qquad R_{k_0}M\,\rho_\sharp(g_{K-s}).
\end{equation*}
Here $R_{k_0}$ is a polylogarithmic function of $x_{k_0}=x_{K-s}$, and $M=L^{m'}$ is one of the
auxiliary parameters collected in $\mathfrak{p}$. Each of the three therefore reads no coupling
other than $g_{K-s}$;
besides it only $L$ and the auxiliary parameters enter. The first
two enter as lower bounds on $C_W$, so they are one-sided floors; the third is the factor of the
room radius. Each depends on $s$ only through $x_{K-s}$, and on neither $K$ nor $m$. This is (b).

\emph{Without the power-of-$L$ hypothesis.} Write $R_jM=L^{\nu_j}\mu_j$ with $\nu_j$ maximal.
Once $L^{k_0-j}$ absorbs the factor $L^{\nu_j}$, the cubes of level $j$, seen in cells, repeat
only under translations by multiples of $\mu_j$. This holds for every $j\le k_0$, including the
levels excluded from \eqref{8.5:eq-threshold-period-1}. The common period $\Lambda$ of all these
partitions, in the sense of (c), is therefore a common multiple of every $\mu_j$, $j\le k_0$, and
so is at least $\operatorname{lcm}_{j\le k_0}\mu_j$.

For $j\le k_0$ the inverse couplings $x_j$ range from $x_{k_0}$ up to the bare value $x_0$, which
depends on $K$. The functions $R_j$, and with them the $\mu_j$, therefore need not be bounded
in terms of $g_{K-s}$ alone, so that $\Lambda$ can depend on $K$. On this branch the floor entry
$16\Lambda$ and the room-radius factor $\Lambda\rho_\sharp(g_{K-s})$ are not shown to be
independent of $K$. Uniformity
in $K$ of the bound of the second-derivative room lemma is therefore not proved for these
members. This is (c).
\end{proof}

\begin{remark}[The sourced format on the physical torus]\label{8.5:rem-torus-format}
Let $k_0=K-s$ be the witness level, $\ell=L^{-s}$ the side of its cells, and $\mathsf{d}=C_W\ell$
the separation scale. The construction of the two-point witness theorem takes the following
three facts as its standing setting.
\begin{itemize}
\item[(a)] The sourced-format statement of the correlation theorem applies at the witness level
$k_0$ on the physical torus $\mathbb{T}_m$. The resulting family is the frozen sourced family of
Definition~\ref{8.5:def-frozen-sourced-family}.
\item[(b)] The torus $\mathbb{T}_m$ holds the two supports $B(x_1,5\mathsf{d}/4)$ and
$B(x_2,5\mathsf{d}/4)$ of the test functions. A term of the family at scale $j$ whose
localisation domain $X\in\mathcal{D}_j$ winds around $\mathbb{T}_m$ satisfies, for a constant $C$,
\begin{equation}\label{8.5:eq-torus-format-1}
d_j(X)\;\ge\;\frac{L^m}{\ell M}-C ,
\end{equation}
and the summed majorant mass of these terms is bounded within the constant $C_1$ once
\begin{equation}\label{8.5:eq-torus-format-2}
L^m\ell^{-1}\;\ge\;3C_W .
\end{equation}
\item[(c)] On the window, the family carrying the source coincides with the whole-lattice
family up to terms of large length. This fact is one of the hypotheses of the second-derivative
room lemma.
\end{itemize}
The integer $m$ enters none of the constants.
\end{remark}

\begin{theorem}[The second-derivative room lemma. Stated; proof incomplete at the step marked $(\ast)$]\label{8.5:thm-room-lemma}
Let $\mathsf{d}=C_W\ell$ be the separation scale of Definition~\ref{8.5:def-separation-scale}. Let
$x_1,x_2$ satisfy the centre separation of Definition~\ref{8.5:def-separation-scale},
$3\mathsf{d}\le|x_1-x_2|\le A_c\mathsf{d}$. This is a condition on the positions of the supports of
the two weights. Its upper member enters no constant of \eqref{8.5:eq-room-lemma-b} below. It is
read only in the range of Lemma~\ref{8.5:lem-covariance-upper}, with the floor members displayed
there, and in the conversion to the unit \eqref{8.5:eq-room-lemma-4} of the two-point witness
theorem, after the proof. Let
$\mathcal{C}=\mathcal{C}^{\mathrm{bg}}$, $\sigma_{\mathsf{d}}$, $N_1$ and the unit $\mathfrak{U}$ be as in
the Gaussian frame at the witness level (Definition~\ref{8.5:def-gaussian-frame},
\eqref{8.5:eq-gaussian-frame-a}). Let $Q_2$ be the Gaussian quadratic form of that definition with
a weight $h_2$ specified below. Let $Q_1^E$ be the part of the marked quadratic form $Q_1$ of that
definition that is
carried by the $\beta$-extracted family $\operatorname{Irr}_1^E$. Let $\bar{\mathfrak{a}}^E$ be the
supremum of the cell-aggregate majorant masses $\mathfrak{a}^E_c$ of that family. Assume the
following.
\begin{itemize}
\item[(1)] $(\ast)$ The zero block mean of the marked residual form,
Proposition~\ref{8.5:prop-zero-block-mean}, for the family to which it is applied here. That family is
the frozen sourced family at the witness level of Definition~\ref{8.5:def-frozen-sourced-family}.
For the frozen sourced family that proposition is stated; proof incomplete, and this appeal is
the step marked $(\ast)$.
\item[(2)] The source-profile defect, Proposition~\ref{8.5:prop-profile-defect}.
\item[(3)] The cell-scale difference bound for the background covariance,
Lemma~\ref{8.5:lem-covariance-difference}. It holds for translates $p,p'$ (same plane and
orientation), at cell resolution, under its floor \eqref{8.5:eq-covariance-difference-a}.
\item[(4)] The upper bound of the background covariance on the separated range,
Lemma~\ref{8.5:lem-covariance-upper}.
\item[(5)] The growth of the background covariance in the explicit scale,
Lemma~\ref{8.5:lem-covariance-growth}. This is a bound from above only.
\item[(6)] The radius tails of the majorant masses, Lemma~\ref{8.5:lem-radius-tails}. They are stated
for the majorant masses $\hat{\mathfrak{m}}_X$ of that lemma, with the constant $C_w\ge1$ and the rate
$\hat\kappa$ fixed there, and $\hat\kappa_R$ the rate of the same tails for the remainder species.
\item[(7)] The distance-decaying tail of the profile-corrected weight, in
Lemma~\ref{8.5:lem-corrected-weight}, \eqref{8.5:eq-corrected-weight-a}. It is used in
step~(vi) of the proof.
\item[(8)] The pairing of the domains that wrap around the torus $\mathbb{T}_m$ (compare
Remark~\ref{8.5:rem-torus-format}), in the form used in step~(iv) of the proof (stated; proof
incomplete).
\item[(9)] Coincidence of the family present near $B(x_1,2\mathsf{d})$ with the whole-lattice family
(stated; proof incomplete). Its content is as follows. The $z$-marked family present on the
small-field region near $B(x_1,2\mathsf{d})$ of Definition~\ref{8.5:def-frozen-sourced-family}
coincides with the whole-lattice frozen family of that definition.
Hypothesis~(1) asserts the zero block means of that whole-lattice family. The coincidence holds up to
terms $X$ of tree length $d_j(X)\ge c'\,C_W/M$, with $c'>0$ a number supplied by the
coincidence statement itself; such terms are called long. The coincidence holds in both directions:
the family present on that region is
contained in the union of the whole-lattice family and the long terms, and every whole-lattice term
anchored near that region but absent from it is long. By the tree-length form of
Lemma~\ref{8.5:lem-radius-tails} (the shell sum in $d_j(X)$ of its proof), the total majorant mass
of these long terms is absorbed by the floor
member $C_W\ge C_1$ of \eqref{8.5:eq-room-lemma-a} below. They are paid explicitly in the split into
long terms of step~(iii) of the proof (compare Remark~\ref{8.5:rem-torus-format}).
\item[(10)] For members of the frozen sourced family of
Definition~\ref{8.5:def-frozen-sourced-family} localised on the cubes of side $R_jM$ of
\cite{B-CMP119}: the hypothesis of Remark~\ref{8.5:rem-threshold-period} (stated; proof
incomplete), that every $R_j$ is a power of $L$.
\item[(11)] The fine-oscillation H\"older form of the decay bound for Ba{\l}aban's propagators.
Together with hypothesis~(12) it places the directions that the proof pairs with the
majorant masses, among them the columns $\mathcal{C}_{\cdot q}$, in the regularity class carried by
those masses.
\item[(12)] Admissibility: the directions that the proof pairs with the majorant masses, among them
the columns $\mathcal{C}_{\cdot q}$, are admissible arguments of the bound carried by the majorant
masses of Lemma~\ref{8.5:lem-radius-tails}.
\end{itemize}
Hypotheses~(11) and~(12) are used under the floor members $C_W\ge16C_\nabla$,
$C_W\ge C_{\mathrm{reg}}^{\mathrm{dec}}$ and the members at the minimal pair distance $\mathsf{d}/8$
of \eqref{8.5:eq-room-lemma-a} below.

Suppose the following one-sided floors hold. Each is displayed, and none is read from above:
\begin{equation}\label{8.5:eq-room-lemma-a}
\begin{gathered}
C_W\ge\max\bigl\{8\rho_\sharp+8,\ 16M,\ 16C_\nabla,\ C_1,\ C_1',\
C_0^{\mathrm{dec}}(\tfrac14),\ C_0^{\mathrm{dec}}(\tfrac18),\
C_\nabla^{\mathrm{dec}},\ C_{\mathrm{reg}}^{\mathrm{dec}}\bigr\},\\
C_W L^{K-s}\ge 16\,z_0,\qquad L^{2(K-s)}\ge 64\,C_{\mathrm{Green}}/C_\nabla^0 ,\\
C_w'\,e^{-\min\{\hat\kappa,\hat\kappa_R\}C/(64M)}\,\bigl(2+C_\sigma^2C^8\bigr)\,C\le 1
\qquad\text{for all } C\ge C_1,\\
C_1=C_1(\min\{\hat\kappa,\hat\kappa_R\},M,C_w',C_\sigma).
\end{gathered}
\end{equation}
The floor also includes the members bounding $K-s$ and $m+s$ from below that are displayed with
Lemma~\ref{8.5:lem-covariance-upper} and at \eqref{8.5:eq-covariance-difference-a}. The two members
in the second line of \eqref{8.5:eq-room-lemma-a} belong to the minimal pair distance $\mathsf{d}/8$
of step~(i) of the proof. Here $C_1'$ is the floor member fixed in step~(vi) of the proof,
$C_w'\ge C_w$ the constant of step~(iv) of the proof, and $C_1$ the least number with the property
in the third line of \eqref{8.5:eq-room-lemma-a}; one number $C_1$ serves both the
$\beta$-extracted and the remainder families. Further, $\eta$, $\mathfrak{r}$, $z_0$,
$C_{\mathrm{Green}}$ (the Green's-function constant) and $C_0^{\mathrm{dec}}(\tfrac14)$ are the
quantities of
Lemma~\ref{8.5:lem-covariance-upper}; $C_\nabla^0$ and $C_\nabla^{\mathrm{dec}}$ are those of
Lemma~\ref{8.5:lem-covariance-difference}; and $C_0^{\mathrm{dec}}(\tfrac18)$,
$C_{\mathrm{reg}}^{\mathrm{dec}}$ are fixed in step~(i) of the proof. Finally, $C_\nabla$ is the
constant of Lemma~\ref{8.5:lem-covariance-difference} and $C_\sigma$ that of
Lemma~\ref{8.5:lem-covariance-growth}.

The member $C_W\ge16C_\nabla$ contains the member $C_W\ge 4C_\nabla$. With it, the floor
\eqref{8.5:eq-covariance-difference-a} gives $\eta+\mathfrak{r}\le1$, the smallness used in
Lemma~\ref{8.5:lem-covariance-upper}. This member is implied by $C_W\ge16M$ once $C_\nabla\le 4M$;
in any case it is one-sided.

Let $f_1$ be a plane-independent Lipschitz weight in the sense of
Definition~\ref{8.5:def-separation-scale}, with $\operatorname{supp}f_1\subset B_1=B(x_1,5\mathsf{d}/4)$
and $f_1\in\mathfrak{F}(x_1,5\mathsf{d}/4;\Theta_0)$, the class of that definition. Let $f_2$ be
plane-independent with $f_2\in\mathfrak{F}(x_2,5\mathsf{d}/4;\Theta_0)$ and
$\operatorname{supp}f_2\subset B(x_2,\mathsf{d})$. Let
$h_2=h_2^{\mathrm{in}}+h_2^{\mathrm{out}}$, where
$\operatorname{supp}h_2^{\mathrm{in}}\subset B_2=B(x_2,5\mathsf{d}/4)$, with
$|h_2(q)|\le2\|f_2\|_\infty$ for all $q$, and, off $B_2$,
\begin{equation}\label{8.5:eq-room-lemma-1}
|h_2^{\mathrm{out}}(q)|\le 3C_w\,
e^{-\hat\kappa\left(\operatorname{dist}(x_q,B(x_2,\mathsf{d}))/\ell-2\right)/(2M)}\,
\|f_2\|_\infty .
\end{equation}
Let $\mathfrak{U}$ be read with $\sum_q|h_2(q)|$ replaced by
\begin{equation}\label{8.5:eq-room-lemma-2}
\sum_q|h_2^{\mathrm{in}}(q)|+2\varepsilon^{-4}|B_2|\,\|f_2\|_\infty .
\end{equation}
Then
\begin{equation}\label{8.5:eq-room-lemma-b}
\begin{split}
\bigl|\operatorname{Cov}(Q_1^E,Q_2)\bigr|
&\le r^{M}_{\mathrm{room}}\,\bar{\mathfrak{a}}^E\,\mathfrak{U},\\
r^{M}_{\mathrm{room}}
&:=C_W^{-1}\bigl[2M(\Theta_0+8C_\nabla)+48C_\nabla\rho_\sharp+4\rho_\sharp\Theta_0+2\bigr]\\
&\le 64\,M\rho_\sharp(1+\Theta_0+C_\nabla)/C_W .
\end{split}
\end{equation}
In \eqref{8.5:eq-room-lemma-b} the summand $2$ in the bracket is the one unit taken by the profile tail
of $h_2$, in step~(vi) of the proof.

This bound is absolute: no grading is consumed. The majorant masses carry only the level-free route
weight $\bar{\mathfrak{w}}=r_w^{-1}$, where $r_w>0$ is the radius of the sup-ball of
clause~(b) of \eqref{8.5:eq-profile-locality-a}, and the grading letters $\mathfrak{w}_H$,
$\mathfrak{w}_W$ do not enter. No lower bound on an entry $\mathcal{C}_{pq}$, or on
the kernel maximum over $B(x_1,3\mathsf{d}/2)$, is used: the kernel enters only through the upper
bounds of Lemma~\ref{8.5:lem-covariance-difference}, Lemma~\ref{8.5:lem-covariance-upper} and
Lemma~\ref{8.5:lem-covariance-growth}. No constant depends on $K$, $s$, $m$ or the position. For
members of the frozen sourced
family of Definition~\ref{8.5:def-frozen-sourced-family}
localised on the cubes of side $R_jM$ of \cite{B-CMP119}, this is asserted under hypothesis~(10), with
$M$ replaced there by $\Lambda=R_{k_0}M$. The coupling enters only through the factor
$b_0g^4_{k_0}$ of the block law, which is the main term's own prefactor, taken at $g_{K-s}$.
\end{theorem}

\begin{remark}[The weight $h_2$, the shape constant and the unit in
Theorem~\ref{8.5:thm-room-lemma}]
For the weight $h_2=f_2+\delta f_2$ of the two-point witness theorem, with
$\operatorname{supp}f_2\subset B(x_2,\mathsf{d})$, the decomposition and the bound
\eqref{8.5:eq-room-lemma-1} are the two bounds \eqref{8.5:eq-corrected-weight-a}. These also give
$|h_2|\le2\|f_2\|_\infty$ everywhere. If $h_2^{\mathrm{out}}\equiv0$, the hypotheses on $h_2$ reduce to
$\operatorname{supp}h_2\subset B_2$ and $|h_2|\le2\|f_2\|_\infty$. In \eqref{8.5:eq-room-lemma-2} the
second summand is the budget for the tail, paid in step~(vi) of the proof.

The test-function class of the two-point witness theorem satisfies
\begin{equation}\label{8.5:eq-room-lemma-3}
\mathfrak{F}(x_i,\mathsf{d};\Theta_0^W)\subset
\mathfrak{F}\bigl(x_i,5\mathsf{d}/4;(5/4)^5\Theta_0^W\bigr),
\end{equation}
where $\Theta_0^W$ is its shape constant. For that theorem, $\Theta_0:=(5/4)^5\Theta_0^W$ is
therefore to be read in every constant above.

The kernel maximum
\begin{equation*}
\mathsf{S}_q=\max\bigl\{|\mathcal{C}_{pq}|:\ p\ \text{bare in}\ B(x_1,3\mathsf{d}/2)\bigr\}
\end{equation*}
of Definition~\ref{8.5:def-gaussian-frame} appears only inside the proof. It is
bounded there by $\sigma_{\mathsf{d}}$ through Lemma~\ref{8.5:lem-covariance-upper}. The unit
$\mathfrak{U}$ carries the explicit scale $\sigma_{\mathsf{d}}$.

Suppose the unit is built on $\ell^{-4}\|f_1\|_1$ in place of $\|f_1\|_\infty N_1$. Then the right
side of \eqref{8.5:eq-room-lemma-b} is to be multiplied by $3\Theta_0$.

The conversion of \eqref{8.5:eq-room-lemma-b} to the unit of the two-point witness theorem,
\begin{equation}\label{8.5:eq-room-lemma-4}
\mathfrak{U}_W=b_0g^4_{k_0}\,\mathcal{K}^+(f_1,f_2),
\end{equation}
where $\mathcal{K}^+(f_1,f_2)$ is the positive majorant form of the kernel in the two-point witness
theorem, is displayed after the proof of Theorem~\ref{8.5:thm-room-lemma}. It is free of $g$.
It uses only the following:
Lemma~\ref{8.5:lem-covariance-upper}; the lower bound on the majorant form $\mathcal{K}^+$;
clause~(b) of \eqref{8.5:eq-profile-locality-a}; and the bounds \eqref{8.5:eq-corrected-weight-a}
on $h_2$.
\end{remark}

The part of the proof of the second-derivative room lemma given below treats the quadratic form of the family $E$ against the background covariance. It covers the bulk, the grouping of the constant-flux part by anchor cells, the summation by parts at the block scale, and the tail. The lemma is Theorem~\ref{8.5:thm-room-lemma}, which is stated with its proof incomplete. The conclusion of this part is the bound \eqref{8.5:eq-room-bulk-a}. Stated; proof incomplete at the step marked $(\ast)$.

\begin{proof}[Proof of Theorem~\ref{8.5:thm-room-lemma}: bulk, anchor grouping and summation by
parts]
\label{8.5:prf-room-bulk}
The hypotheses of Theorem~\ref{8.5:thm-room-lemma} are in force throughout; in particular:
\begin{itemize}
\item $f_1$ is plane-independent and belongs to $\mathfrak{F}(x_1,5\mathsf{d}/4;\Theta_0)$;
\item $C_W$ satisfies the one-sided floors \eqref{8.5:eq-room-lemma-a}.
\end{itemize}
We write $\mathsf{d}=C_W\ell$, $B_1=B(x_1,5\mathsf{d}/4)$ and $B_2=B(x_2,5\mathsf{d}/4)$, and we put
\begin{equation}\label{8.5:eq-room-bulk-1}
\mathfrak{A}:=B(x_1,3\mathsf{d}/2),\qquad
\mathfrak{A}':=B(x_1,13\mathsf{d}/8),\qquad
\mathsf{S}_q:=\max\{|\mathcal{C}_{pq}|:\ p\ \text{bare in}\ \mathfrak{A}\}.
\end{equation}
Here $\mathcal{C}=\mathcal{C}^{\mathrm{bg}}$ is the background flux covariance. The forms $Q_X$, $q_X[f_1]$ and the majorants $\mathfrak{m}^{\mathrm{reg}}_X$, $\hat{\mathfrak{m}}_X$ are in the normalisation \eqref{8.5:eq-gaussian-frame-b}, and $\sigma_{\mathsf{d}}$, $N_1$ and $\mathfrak{U}_q$ are as in \eqref{8.5:eq-gaussian-frame-a}.

Fix $q\in\operatorname{supp}h_2$ and one colour. All sums below are over $X\in E$; here $E$ denotes the family of marked terms of Theorem~\ref{8.5:thm-room-lemma} present on the window, the windowed family. A term is \emph{long} if its tree length $d_j$ is at least the threshold $n_{\mathrm{GR}}$ of step (iii). We use the statement that the windowed family coincides with the whole-lattice frozen family up to long terms; that statement is stated with its proof incomplete. By it, $E$ is the whole-lattice frozen family of the marked terms, up to long terms, which are bounded separately in the long-term split of step (iii) and in the tail of step (iv).

For a term $X$ we use the following notation:
\begin{itemize}
\item $\hat{X}$ is its filled hull and $r(X)$ the radius of $\hat{X}$ in cells;
\item $X^+$ is its reading neighbourhood and $\hat{X}^+$ the enlargement of $\hat{X}$ by one cell;
\item $c(X)$ is its anchor cell.
\end{itemize}
$\mathfrak{O}$ is the set of planes, $\pi(p)$ the plane of $p$, and $p_P(c)$ the bare plaquette of plane $P\in\mathfrak{O}$ at the least corner of the cell $c$.

\emph{(i) Bulk: $r(X)\le R_0:=C_W/8-2$.}
Since $M\ge L\ge2$, the floor $C_W\ge16M$ gives $C_W\ge32$, and hence $R_0\ge C_W/16$.

A term $X$ that contributes satisfies $X^+\cap\operatorname{supp}f_1\neq\emptyset$, by part (a) of the localisation statement for the marked family. Moreover $\operatorname{supp}f_1\subset B_1$. The Euclidean reach of a filled hull of radius at most $R_0$ cells is $2(R_0+1)\ell=\mathsf{d}/4-2\ell$. Therefore
\begin{equation}\label{8.5:eq-room-bulk-2}
\begin{gathered}
\hat{X}\subset B\bigl(x_1,5\mathsf{d}/4+2(R_0+1)\ell\bigr)=B(x_1,3\mathsf{d}/2-2\ell),\\
\hat{X}^+\subset B(x_1,3\mathsf{d}/2-\ell)\subset\mathfrak{A}.
\end{gathered}
\end{equation}
The anchor cell of $X$ lies in $\hat{X}$, hence in $\mathfrak{A}$. This is the ball on which the cell-scale difference bound \eqref{8.5:eq-covariance-difference-b} holds, under the floor \eqref{8.5:eq-covariance-difference-a}.

For each cell $c\in\mathfrak{A}$ define $\mu_c\in\mathbb{R}^{\mathfrak{O}}$ by $\mu_c(P):=\mathcal{C}_{p_P(c),q}$. By the upper bound \eqref{8.5:eq-covariance-upper-a},
\begin{equation}\label{8.5:eq-room-bulk-3}
|\mu_c|_\infty\le\mathsf{S}_q\le\sigma_{\mathsf{d}}.
\end{equation}

For $X$ anchored at $c$ we write, on $\hat{X}$,
\begin{equation*}
\mathcal{C}_{pq}=\mu_c(\pi(p))+\xi_p .
\end{equation*}
Both summands are exact on $\hat{X}$. The first is a constant flux restricted to a box. The second is a difference of two exact fields.

For $p\in\hat{X}$ the plaquettes $p$ and $p_{\pi(p)}(c)$ are translates of each other, and $|x_p-x_{p_{\pi(p)}(c)}|\le2(r(X)+1)\ell$. The bound \eqref{8.5:eq-covariance-difference-b} is radial and at cell resolution. Applying it to this pair, and using $r(X)\ge1$, we get
\begin{equation}\label{8.5:eq-room-bulk-4}
\sup_{\hat{X}}|\xi|_\infty\le C_\nabla\bigl(2r(X)+6\bigr)\frac{\ell}{\mathsf{d}}\,\sigma_{\mathsf{d}}
\le16C_\nabla\frac{r(X)}{C_W}\,\sigma_{\mathsf{d}}=:\theta_r\sigma_{\mathsf{d}} .
\end{equation}

By bilinearity,
\begin{equation}\label{8.5:eq-room-bulk-5}
Q_X[\mathcal{C}_{\cdot q},\mathcal{C}_{\cdot q}]
=q_X[f_1][\mu_c,\mu_c]+\bigl(2Q_X[\mu_c,\xi]+Q_X[\xi,\xi]\bigr).
\end{equation}
The directions are measured in the norm of \eqref{8.5:eq-constant-flux-forms-a},
\begin{equation*}
\|\cdot\|_X=\max\{|\cdot|_\infty,\ \mathfrak{reg}_X(\cdot)/r(X)\}.
\end{equation*}
Here the regularity functional $\mathfrak{reg}_X$ is built from the H\"older-$\beta_0$ increments and from $\eta^{-1}|d^*\cdot|$. In this norm the majorant applies to the directions, by the admissibility statement of \eqref{8.5:eq-constant-flux-forms-a}, which is stated with its proof incomplete. That is, $|Q_X[u,v]|\le\mathfrak{m}^{\mathrm{reg}}_X[f_1]\,\|u\|_X\|v\|_X$ for the directions used below.

The direction $\mu_c$ is constant. Hence $\mathfrak{reg}_X(\mu_c)=0$ and $\|\mu_c\|_X=|\mu_c|_\infty\le\sigma_{\mathsf{d}}$.

We now bound the regularity of $\xi=(\mathcal{C}_{\cdot q}-\mu_c)|_{\hat{X}}$. We use the fine-oscillation H\"older form of Ba{\l}aban's propagator decay. Its hypotheses are met as follows:
\begin{itemize}
\item the profile is $\mathcal{C}_{\cdot q}$, which lies in the range of $\bar{M}_{k_0}$ because $\mathcal{C}^{\mathrm{bg}}=\bar{M}_{k_0}\bar{Q}P_{\mathcal{E}}$ in the formula of the witness construction for the background covariance, $\mathcal{E}$ being the space of exact fields;
\item the fine constant is $\mu_c$;
\item the H\"older exponent is $\beta=\beta_0$.
\end{itemize}
Since $\mu_c$ is constant, its increments and its $d^*$ vanish. The fine-oscillation H\"older form then gives
\begin{equation}\label{8.5:eq-room-bulk-6}
\mathfrak{reg}_X(\xi)=\mathfrak{reg}_X(\mathcal{C}_{\cdot q}|_{\hat{X}})
\le\frac{C_{\mathrm{reg}}}{1-\beta_0}\max_{p\subset\hat{X}}\sum_{c'}
e^{-a\operatorname{dist}(c',p)}\sup_{c'}|\mathcal{C}_{\cdot q}-\mu_c| .
\end{equation}

Fix $p\subset\hat{X}$, write $r=r(X)$, and split the cells $c'$ in \eqref{8.5:eq-room-bulk-6} into near cells and far cells.

\emph{Near cells,} $\operatorname{dist}(p,c')<C_W/16$. Such cells satisfy $c'\subset\mathfrak{A}'$, at distance at least $\mathsf{d}/8$ from $B_2$. We apply the bound \eqref{8.5:eq-covariance-difference-b} on $\mathfrak{A}'$; it carries the constant $32C_\nabla^0$ there. Here $C_\nabla^0\le C_\nabla$ is the constant of \eqref{8.5:eq-covariance-difference-b} before its enlargement, with which that bound already holds. We compare with the anchor's corner $p_P(c)\in\hat{X}$, plane by plane. This gives
\begin{equation}\label{8.5:eq-room-bulk-7}
\begin{split}
\sup_{c'}|\mathcal{C}_{\cdot q}-\mu_c|
&\le32C_\nabla^0\sigma_{\mathsf{d}}\bigl(2\operatorname{dist}(c',p)+2r+10\bigr)\ell/\mathsf{d}\\
&=32C_\nabla^0\sigma_{\mathsf{d}}\bigl(2\operatorname{dist}(c',p)+2r+10\bigr)/C_W .
\end{split}
\end{equation}

\emph{Far cells,} $\operatorname{dist}(p,c')\ge C_W/16$. The growth bound \eqref{8.5:eq-covariance-growth-a} applied to $\mathcal{C}_{\cdot q}$, together with \eqref{8.5:eq-room-bulk-3} applied to $\mu_c$, gives
\begin{equation*}
\sup_{c'}|\mathcal{C}_{\cdot q}-\mu_c|\le(C_\sigma^0C_W^4+1)\sigma_{\mathsf{d}} .
\end{equation*}
Here $C_\sigma^0\le C_\sigma$ is the constant of \eqref{8.5:eq-covariance-growth-a} before its enlargement, with which the sup bound of \eqref{8.5:eq-covariance-growth-a} already holds.
On the other hand,
\begin{equation*}
\sum_{\operatorname{dist}\ge C_W/16}e^{-a\operatorname{dist}}\le e^{-aC_W/32}C_\Gamma' .
\end{equation*}
The product of the two is at most $\sigma_{\mathsf{d}}/C_W$ under the one-sided floor member $C_W\ge C^{\mathrm{dec}}_{\mathrm{reg}}$ of \eqref{8.5:eq-room-lemma-a}, a floor member of the same kind as the decay floor $C_0^{\mathrm{dec}}(\tfrac18)$ described below. Here $C^{\mathrm{dec}}_{\mathrm{reg}}$ is the least number such that
\begin{equation}\label{8.5:eq-room-bulk-8}
(C_\sigma^0C^4+1)\,C_\Gamma'\cdot C\cdot e^{-aC/32}\le1
\qquad\text{for all }C\ge C^{\mathrm{dec}}_{\mathrm{reg}} .
\end{equation}

\emph{Sum of the two parts.} Summing the near cells with the constant $C_\Gamma''$ and adding the far cells gives
\begin{equation}\label{8.5:eq-room-bulk-9}
\sum_{c'}e^{-a\operatorname{dist}(c',p)}\sup_{c'}|\mathcal{C}_{\cdot q}-\mu_c|
\le\bigl[32C_\nabla^0C_\Gamma''(2r+12)+1\bigr]\frac{\sigma_{\mathsf{d}}}{C_W}
\le30\cdot16\,C_\nabla^0C_\Gamma''\,\frac{r\,\sigma_{\mathsf{d}}}{C_W}.
\end{equation}
The second inequality holds because $r\ge1$ and $C_\nabla^0,C_\Gamma''\ge1$, so that $32(2r+12)+1\le449r\le480r$. Inserting \eqref{8.5:eq-room-bulk-9} into \eqref{8.5:eq-room-bulk-6} gives
\begin{equation*}
\mathfrak{reg}_X(\xi)\le30\,\frac{C_{\mathrm{reg}}}{1-\beta_0}\,C_\Gamma''\cdot
\frac{16C_\nabla^0r}{C_W}\,\sigma_{\mathsf{d}} .
\end{equation*}

\emph{The norm of $\xi$.} Recall from \eqref{8.5:eq-covariance-difference-b} that
\begin{equation*}
C_\nabla=C_\nabla^0\max\{1,\,30C_{\mathrm{reg}}(1-\beta_0)^{-1}C_\Gamma''\}.
\end{equation*}
The second factor of the maximum gives
\begin{equation}\label{8.5:eq-room-bulk-10}
\frac{\mathfrak{reg}_X(\xi)}{r}
\le480\,\frac{C_{\mathrm{reg}}}{1-\beta_0}\,C_\Gamma''C_\nabla^0\,\frac{\sigma_{\mathsf{d}}}{C_W}
\le\frac{16C_\nabla\,\sigma_{\mathsf{d}}}{C_W}\le\theta_r\sigma_{\mathsf{d}},
\end{equation}
with $\theta_r$ as in \eqref{8.5:eq-room-bulk-4}.
For the sup part, \eqref{8.5:eq-covariance-difference-b} gives, as in \eqref{8.5:eq-room-bulk-4},
\begin{equation*}
|\xi|_\infty\le C_\nabla^0(2r+6)\sigma_{\mathsf{d}}/C_W\le\theta_r\sigma_{\mathsf{d}} .
\end{equation*}
Hence $\|\xi\|_X\le\theta_r\sigma_{\mathsf{d}}$.

\emph{The norm of $\mathcal{C}_{\cdot q}|_{\hat{X}}$.} The middle member of \eqref{8.5:eq-room-bulk-10} does not depend on $r$. Under the displayed one-sided floor member $C_W\ge16C_\nabla$ it therefore also gives
\begin{equation*}
\mathfrak{reg}_X(\mathcal{C}_{\cdot q}|_{\hat{X}})/r=\mathfrak{reg}_X(\xi)/r\le\sigma_{\mathsf{d}} .
\end{equation*}
Together with $|\mathcal{C}_{pq}|\le\mathsf{S}_q\le\sigma_{\mathsf{d}}$ for $p\in\hat{X}\subset\mathfrak{A}$, this gives $\|\mathcal{C}_{\cdot q}|_{\hat{X}}\|_X\le\sigma_{\mathsf{d}}$. No power of $\eta$ enters these two bounds.

The bounds \eqref{8.5:eq-covariance-upper-a} and \eqref{8.5:eq-covariance-difference-b} are used here for pair distances down to $\mathsf{d}/8$ on $\mathfrak{A}'$. They hold there by the same implications, under three one-sided floor members listed with \eqref{8.5:eq-room-lemma-a}:
\begin{equation*}
C_W\ge C_0^{\mathrm{dec}}(\tfrac18),\qquad C_WL^{K-s}\ge16z_0,\qquad
L^{2(K-s)}\ge64C_{\mathrm{Green}}/C_\nabla^0 .
\end{equation*}
Here $C_0^{\mathrm{dec}}(\tfrac18)$ is the floor of the decay bound for the fluctuation kernel read at the pair-distance fraction $\tfrac18$, $z_0$ is the threshold of the $z_0$-clause of the comparison of the free kernel with its continuum counterpart, and $C_{\mathrm{Green}}$ is the Green's-function constant, the constant in the asymptotics of the lattice Green's function; all three are fixed with \eqref{8.5:eq-covariance-upper-a} and \eqref{8.5:eq-covariance-difference-a}.
Every bound below holds with the constants as displayed. There are two cases.

\emph{Case $\theta_r\le1$.} The bracket in \eqref{8.5:eq-room-bulk-5} is at most
\begin{equation*}
(2\theta_r+\theta_r^2)\,\mathfrak{m}^{\mathrm{reg}}_X[f_1]\sigma_{\mathsf{d}}^2
\le3\theta_r\,\mathfrak{m}^{\mathrm{reg}}_X[f_1]\sigma_{\mathsf{d}}^2 .
\end{equation*}
This uses $\|\mu_c\|_X\le\sigma_{\mathsf{d}}$ and $\|\xi\|_X\le\theta_r\sigma_{\mathsf{d}}$. Here $\mathfrak{m}^{\mathrm{reg}}_X[f_1]$ is the regular-direction mass of \eqref{8.5:eq-constant-flux-forms-a}.

\emph{Case $\theta_r>1$.} The constant-flux form satisfies $q_X[f_1][\mu_c,\mu_c]=Q_X[\mu_c,\mu_c]$. Using $\|\mathcal{C}_{\cdot q}|_{\hat{X}}\|_X\le\sigma_{\mathsf{d}}$ and $\|\mu_c\|_X\le\sigma_{\mathsf{d}}$, we get
\begin{equation*}
|Q_X[\mathcal{C}_{\cdot q},\mathcal{C}_{\cdot q}]-q_X[f_1][\mu_c,\mu_c]|
\le2\mathfrak{m}^{\mathrm{reg}}_X[f_1]\sigma_{\mathsf{d}}^2
\le2\theta_r\mathfrak{m}^{\mathrm{reg}}_X[f_1]\sigma_{\mathsf{d}}^2 .
\end{equation*}

In both cases we use the majorant bound $\mathfrak{m}^{\mathrm{reg}}_X[f_1]\le\hat{\mathfrak{m}}_X\sup_{X^+}|f_1|$ for the regular-direction mass, which is the sup-ball bound of part (b) of the localisation statement for the marked family, read with the majorant masses of \eqref{8.5:eq-constant-flux-forms-a}; and $3\theta_r=48C_\nabla r(X)/C_W$. This gives
\begin{equation}\label{8.5:eq-room-bulk-11}
|Q_X[\mathcal{C}_{\cdot q},\mathcal{C}_{\cdot q}]-q_X[f_1][\mu_c,\mu_c]|
\le48C_\nabla\frac{r(X)}{C_W}\,\hat{\mathfrak{m}}_X\|f_1\|_\infty\sigma_{\mathsf{d}}^2 .
\end{equation}
Now sum over the bulk terms $X$ anchored at $c$. Recall that $\rho_\sharp$ is the mass-weighted mean radius, so that $\sum_{c(X)=c}r(X)\hat{\mathfrak{m}}_X\le\rho_\sharp\mathfrak{a}^E_c$. The sum of the right sides of \eqref{8.5:eq-room-bulk-11} is therefore at most
\begin{equation*}
48C_\nabla C_W^{-1}\|f_1\|_\infty\rho_\sharp\mathfrak{a}^E_c\sigma_{\mathsf{d}}^2 .
\end{equation*}

\emph{(ii) The constant-flux part, grouped by anchor.}
For a cell $c$ let $q^E(c)[\mathbf{1}]$ be the constant-flux form at the constant profile $\mathbf{1}$, aggregated over the terms anchored at $c$. The source-profile defect \eqref{8.5:eq-profile-defect-a} gives, for every $c\in\mathfrak{A}$, with $x_c$ the reference point of the cell $c$ in \eqref{8.5:eq-profile-defect-a},
\begin{equation}\label{8.5:eq-room-bulk-12}
\sum_{c(X)=c,\ r(X)\le R_0}q_X[f_1]
=f_1(x_c)\bigl[q^E(c)[\mathbf{1}]-q^{\mathrm{tail}}(c)\bigr]+e(c),
\qquad
q^{\mathrm{tail}}(c):=\sum_{c(X)=c,\ r(X)>R_0}q_X[\mathbf{1}].
\end{equation}
The defect satisfies $|e(c)|\le4\rho_\sharp\,\ell\operatorname{Lip}f_1\cdot\mathfrak{a}^E_c$.

For the tail form we use the radius tails \eqref{8.5:eq-radius-tails-a} of the majorant masses. These are proved as an implication from the fine-field form of Ba{\l}aban's bound on localised terms, which is stated with its proof incomplete. They give
\begin{equation}\label{8.5:eq-room-bulk-13}
|q^{\mathrm{tail}}(c)|\le C_w e^{-\hat\kappa R_0/(2M)}\bar{\mathfrak{a}}^E .
\end{equation}
Evaluated at $(\mu_c,\mu_c)$, with $|\mu_c|_\infty\le\sigma_{\mathsf{d}}$, the $e$-piece and the tail piece together contribute, per anchor, at most
\begin{equation*}
\bigl[4\rho_\sharp\,\ell\operatorname{Lip}f_1\,\mathfrak{a}^E_c
+\|f_1\|_\infty C_w e^{-\hat\kappa R_0/(2M)}\bar{\mathfrak{a}}^E\bigr]\sigma_{\mathsf{d}}^2 .
\end{equation*}

\emph{(iii) Summation by parts at the block scale.}
The remaining piece is
\begin{equation*}
\sum_{c\in\mathfrak{A}}\langle q^E(c)[\mathbf{1}],\Phi(c)\rangle,
\qquad
\Phi(c):=f_1(x_c)\,\mu_c\otimes\mu_c .
\end{equation*}
We have $\Phi(c)=0$ unless $x_c\in B_1$. Since $\Phi$ vanishes off $B_1$, the sum may be taken over all cells of the period blocks $\mathfrak{B}$ that meet $B_1$; these are the cubes of $M^4$ cells of $\pi_{k_0}$. Each such block is contained in
\begin{equation*}
B(x_1,5\mathsf{d}/4+2M\ell)\subset B(x_1,11\mathsf{d}/8)\subset\mathfrak{A},
\end{equation*}
since $2M\ell\le\mathsf{d}/8$ by the floor $C_W\ge16M$.

\emph{Windowed family and whole-lattice family.} The zero block mean \eqref{8.5:eq-zero-block-mean-a} is asserted for the whole-lattice frozen family. For the frozen family of the two-point witness theorem it is stated with its proof incomplete. The aggregates here, however, are those of the windowed family.

We use the coincidence statement for the two families in both directions:
\begin{itemize}
\item the windowed family is contained in the whole-lattice family together with long terms;
\item whole-lattice terms anchored in these blocks but absent from the window are long.
\end{itemize}
This coincidence is stated with its proof incomplete. By it, the symmetric difference of the two families consists of terms of tree length $d_j\ge n_{\mathrm{GR}}$, where
\begin{equation*}
n_{\mathrm{GR}}:=c_{\mathrm{GR}}C_W/M
\end{equation*}
is the length threshold of the coincidence, $c_{\mathrm{GR}}>0$ being the constant of the coincidence statement. Such terms need not have large radius, and so they may lie in the bulk of step (i).

We therefore write the whole-lattice aggregate as $q^{\mathrm{win}}(c)+q^{\mathrm{long}}(c)$. Here $q^{\mathrm{win}}(c)$ is the windowed aggregate $q^E(c)[\mathbf{1}]$ of (ii), and $q^{\mathrm{long}}(c)$ is the contribution of the terms of the symmetric difference, counted with sign. We use the form of \eqref{8.5:eq-radius-tails-a} indexed by the length $d_j$ instead of the radius. Its proof sums the shells of fixed $d_j$ at ratio at most $e^{-\hat\kappa/2-1}$, and gives the rate $\hat\kappa/2$:
\begin{equation*}
|q^{\mathrm{long}}(c)|\le C_we^{-\hat\kappa n_{\mathrm{GR}}/2}\bar{\mathfrak{a}}^E .
\end{equation*}
We pair $q^{\mathrm{long}}$ with $\Phi(c)$ exactly as $q^{\mathrm{tail}}$ is paired in (ii). The total is at most
\begin{equation}\label{8.5:eq-room-bulk-14}
N_1\|f_1\|_\infty\sigma_{\mathsf{d}}^2\cdot C_we^{-\hat\kappa n_{\mathrm{GR}}/2}\bar{\mathfrak{a}}^E .
\end{equation}
This is absorbed by the floor $C_W\ge C_1$ of \eqref{8.5:eq-room-lemma-a} once
\begin{equation*}
\frac{\hat\kappa n_{\mathrm{GR}}}{2C_W}=\frac{\hat\kappa c_{\mathrm{GR}}}{2M}
\ge\frac{\min\{\hat\kappa,\hat\kappa_R\}}{64M},
\end{equation*}
in particular once $c_{\mathrm{GR}}\ge1/32$. Failing that, it is absorbed after one more one-sided member is added to $C_1$, listed with the rate condition on the torus-coincidence bound in step (iv).

\emph{Use of the zero block mean.} Pick $c_{\mathfrak{B}}\in\mathfrak{B}$. The zero block mean \eqref{8.5:eq-zero-block-mean-a}, applied to the whole-lattice aggregates, gives
\begin{equation*}
\sum_{c\in\mathfrak{B}}\langle q^E(c)[\mathbf{1}],\Phi(c)\rangle
=\sum_{c\in\mathfrak{B}}\langle q^E(c)[\mathbf{1}],\Phi(c)-\Phi(c_{\mathfrak{B}})\rangle,
\end{equation*}
up to the long-term piece \eqref{8.5:eq-room-bulk-14}. Moreover
\begin{equation}\label{8.5:eq-room-bulk-15}
\begin{split}
\Phi(c)-\Phi(c_{\mathfrak{B}})
&=\bigl(f_1(x_c)-f_1(x_{c_{\mathfrak{B}}})\bigr)\mu_c\otimes\mu_c\\
&\quad+f_1(x_{c_{\mathfrak{B}}})\bigl[(\mu_c-\mu_{c_{\mathfrak{B}}})\otimes\mu_c
+\mu_{c_{\mathfrak{B}}}\otimes(\mu_c-\mu_{c_{\mathfrak{B}}})\bigr].
\end{split}
\end{equation}
The factors in \eqref{8.5:eq-room-bulk-15} are bounded as follows:
\begin{itemize}
\item $|x_c-x_{c_{\mathfrak{B}}}|\le2M\ell$;
\item $|\mu_c|_\infty,\ |\mu_{c_{\mathfrak{B}}}|_\infty\le\mathsf{S}_q\le\sigma_{\mathsf{d}}$;
\item the difference of the constant fluxes satisfies
\begin{equation*}
|\mu_c-\mu_{c_{\mathfrak{B}}}|_\infty
\le C_\nabla\sigma_{\mathsf{d}}(2M\ell+4\ell)/\mathsf{d}
\le8MC_\nabla\sigma_{\mathsf{d}}\ell/\mathsf{d}
=4C_\nabla\sigma_{\mathsf{d}}\cdot2M\ell/\mathsf{d}.
\end{equation*}
\end{itemize}
The last bound uses $2M+4\le8M$ for $M\ge1$. It comes from \eqref{8.5:eq-covariance-difference-b}, applied plane by plane to translates, with both corners in $B(x_1,11\mathsf{d}/8)\subset\mathfrak{A}$. Hence, with $|q^E(c)[\mathbf{1}]|\le\mathfrak{a}^E_c$,
\begin{equation}\label{8.5:eq-room-bulk-16}
|\langle q^E(c)[\mathbf{1}],\Phi(c)-\Phi(c_{\mathfrak{B}})\rangle|
\le\mathfrak{a}^E_c\cdot2M\ell\bigl[\operatorname{Lip}f_1+8C_\nabla\|f_1\|_\infty/\mathsf{d}\bigr]
\sigma_{\mathsf{d}}^2 .
\end{equation}
This step uses the smoothness of the profile ($\operatorname{Lip}f_1$, at scale $\mathsf{d}$) and that of the kernel (the bound \eqref{8.5:eq-covariance-difference-b}, at scale $\mathsf{d}$), both at the block scale $2M\ell\le\mathsf{d}/8$. It uses nothing finer than the class $\mathfrak{F}(x_1,5\mathsf{d}/4;\Theta_0)$ and the cell-scale difference bound.

\emph{(iv) Tail: $r(X)>R_0\ge C_W/16$.}
We first consider a non-wrapping term $X$, one whose hull $\hat{X}$ is a box not wrapping the torus. The wrapping domains are treated at the step marked $(\ast)$ below.

By the admissibility statement of \eqref{8.5:eq-constant-flux-forms-a},
\begin{equation}\label{8.5:eq-room-bulk-17}
|Q_X[\mathcal{C}_{\cdot q},\mathcal{C}_{\cdot q}]|
\le\mathfrak{m}^{\mathrm{reg}}_X[f_1]\,\|\mathcal{C}_{\cdot q}|_{\hat{X}}\|_X^2
\le\|f_1\|_\infty\hat{\mathfrak{m}}_X\,C_\sigma^2C_W^8\sigma_{\mathsf{d}}^2 .
\end{equation}
The second inequality is justified in three steps.
\begin{itemize}
\item The sup part is at most $C_\sigma^0C_W^4\sigma_{\mathsf{d}}$ by \eqref{8.5:eq-covariance-growth-a}.
\item For the regularity part, we apply the fine-oscillation H\"older form of the propagator decay with fine constant $0$, together with \eqref{8.5:eq-covariance-growth-a} in every cell. This gives
\begin{equation*}
\mathfrak{reg}_X(\mathcal{C}_{\cdot q}|_{\hat{X}})
\le C_{\mathrm{reg}}(1-\beta_0)^{-1}C_\Gamma'\,C_\sigma^0C_W^4\sigma_{\mathsf{d}} .
\end{equation*}
\item Since $\|\cdot\|_X$ is the maximum of the sup part and of the regularity part divided by $r(X)\ge1$, the two preceding items combine to
\begin{equation*}
\|\mathcal{C}_{\cdot q}|_{\hat{X}}\|_X
\le\max\{1,\,C_{\mathrm{reg}}(1-\beta_0)^{-1}C_\Gamma'\}\,C_\sigma^0C_W^4\sigma_{\mathsf{d}}
=C_\sigma C_W^4\sigma_{\mathsf{d}},
\end{equation*}
with $C_\sigma=C_\sigma^0\max\{1,\,C_{\mathrm{reg}}(1-\beta_0)^{-1}C_\Gamma'\}$ the constant of \eqref{8.5:eq-covariance-growth-a}. Taking $C_\sigma$ this large is harmless because \eqref{8.5:eq-covariance-growth-a} is used from above only.
\end{itemize}
The last factor in \eqref{8.5:eq-room-bulk-17} also uses $\mathfrak{m}^{\mathrm{reg}}_X[f_1]\le\hat{\mathfrak{m}}_X\|f_1\|_\infty$.

\emph{Number of anchor cells.} The anchors of tail terms of radius $r$ lie within Euclidean distance $2(r+1)\ell$ of $B_1$. Their number is at most
\begin{equation*}
\bigl(1+4(r+1)\ell/\mathsf{d}\bigr)^4\cdot2^4N_1\le16(1+r)^4N_1,
\end{equation*}
because $4(r+1)/C_W\le r$ for $r\ge1$ and $C_W\ge8$.

\emph{Summation over $r$.} We use \eqref{8.5:eq-radius-tails-a} at radius $r$ and sum over $r>R_0$. The factor $16(1+r)^4e^{-\hat\kappa r/(4M)}$ is folded into the constant
\begin{equation}\label{8.5:eq-room-bulk-18}
C_w':=C_w\,e^{\hat\kappa/(2M)}\bigl(1-e^{-\hat\kappa/(4M)}\bigr)^{-1}
\sup_{r\ge1}16(1+r)^4e^{-\hat\kappa r/(4M)} .
\end{equation}
This constant contains the shell shift $e^{\hat\kappa/(2M)}$ and the geometric series in $e^{-\hat\kappa/(4M)}$. The factor multiplying $C_w$ depends on $(M,\hat\kappa)$ only; $C_w'$ is one-sided and, like $C_w$, independent of $K$, $s$, $m$ and the position. The floors $C_1$ and $C_1'$ are defined with $C_w'$ in place of $C_w$.

Combining this with \eqref{8.5:eq-room-bulk-17}, the contribution of the non-wrapping tail terms is at most
\begin{equation}\label{8.5:eq-room-bulk-19}
\|f_1\|_\infty\,\bar{\mathfrak{a}}^E N_1\,C_w'e^{-\hat\kappa C_W/(64M)}C_\sigma^2C_W^8
\sigma_{\mathsf{d}}^2 .
\end{equation}

\emph{$(\ast)$ The wrapping domains.} These are the terms whose hull wraps the torus; they have length at least $L^m/(2\ell)$. Their mass is absorbed by the floor $C_W\ge C_1$ under the bound on wrapping domains of the two-point witness theorem, which that theorem takes as stated, with its proof incomplete. For their pairing with the background flux $\psi_q$ at $q$ of the Gaussian frame \eqref{8.5:eq-gaussian-frame-b} the proof uses the following input: this pairing is absorbed by the same floor $C_W\ge C_1$. This is provided the exponential rate of the torus-coincidence bound for the marked family is at least $\min\{\hat\kappa,\hat\kappa_R\}/64$; otherwise one more one-sided member is added to $C_1$. The parameter $m$ enters no constant. This pairing bound is stated, and the proof given here is incomplete at this step.

\emph{(v) Collection.}
Every kernel factor in steps (i)--(iii) is $\sigma_{\mathsf{d}}$. This holds because $|\mu_c|_\infty\le\mathsf{S}_q\le\sigma_{\mathsf{d}}$ by \eqref{8.5:eq-covariance-upper-a}, and because \eqref{8.5:eq-covariance-difference-b} is radial. We also use three bounds.
\begin{itemize}
\item Summation over anchors: since $\mathfrak{A}=B(x_1,3\mathsf{d}/2)$ is the ball defining $N_1$,
\begin{equation*}
\sum_{c\in\mathfrak{A}}\mathfrak{a}^E_c\le\bar{\mathfrak{a}}^EN_1 .
\end{equation*}
\item Lipschitz ratio: since $f_1\in\mathfrak{F}(x_1,5\mathsf{d}/4;\Theta_0)$ and $\|f_1\|_1\le\|f_1\|_\infty|B_1|$,
\begin{equation*}
\frac{\ell\operatorname{Lip}f_1}{\|f_1\|_\infty}
=\frac{\mathsf{d}\operatorname{Lip}f_1}{C_W\|f_1\|_\infty}
\le\frac{1}{C_W}\cdot\frac45\cdot\frac{5\mathsf{d}}{4}\,\frac{\operatorname{Lip}f_1\,|B_1|}{\|f_1\|_1}
\le\frac{\Theta_0}{C_W}.
\end{equation*}
\item Block-scale factor: $2M\ell\cdot8C_\nabla/\mathsf{d}=16MC_\nabla/C_W$.
\end{itemize}

We now insert these bounds into each contribution.
\begin{itemize}
\item Step (i) contributes at most $48C_\nabla\rho_\sharp C_W^{-1}\bar{\mathfrak{a}}^E\mathfrak{U}_q$.
\item The $e$-piece of step (ii) contributes at most $4\rho_\sharp\Theta_0C_W^{-1}\bar{\mathfrak{a}}^E\mathfrak{U}_q$.
\item The tail piece of step (ii), together with $R_0\ge C_W/16$ and $C_w\le C_w'$, contributes at most one unit $C_w'e^{-\hat\kappa C_W/(64M)}\bar{\mathfrak{a}}^E\mathfrak{U}_q$.
\item Step (iii), by \eqref{8.5:eq-room-bulk-16}, contributes at most $(2M\Theta_0+16MC_\nabla)C_W^{-1}\bar{\mathfrak{a}}^E\mathfrak{U}_q$.
\item Step (iv), by \eqref{8.5:eq-room-bulk-19}, contributes the term $C_\sigma^2C_W^8$ in the curly bracket below.
\item The long-term piece \eqref{8.5:eq-room-bulk-14} of step (iii) and the wrapping-domain pieces of $(\ast)$ are not among the contributions collected in the first inequality of \eqref{8.5:eq-room-bulk-a}; under the floor $C_W\ge C_1$ and the rate conditions stated in step (iii) and at $(\ast)$ they are absorbed, together with the exponential member, in its last inequality.
\end{itemize}
Here $\mathfrak{U}_q=\|f_1\|_\infty N_1\sigma_{\mathsf{d}}^2$. Hence
\begin{equation}\label{8.5:eq-room-bulk-a}
\begin{split}
\Bigl|\sum_{X\in E}Q_X[\mathcal{C}_{\cdot q},\mathcal{C}_{\cdot q}]\Bigr|
&\le\bar{\mathfrak{a}}^E\,\mathfrak{U}_q\Bigl\{C_W^{-1}\bigl[48C_\nabla\rho_\sharp
+4\rho_\sharp\Theta_0+2M\Theta_0+16MC_\nabla\bigr]\\
&\qquad\qquad+C_w'e^{-\hat\kappa C_W/(64M)}\bigl(2+C_\sigma^2C_W^8\bigr)\Bigr\}\\
&\le r^M_{\mathrm{room}}\,\bar{\mathfrak{a}}^E\,\mathfrak{U}_q .
\end{split}
\end{equation}
In the factor $2+C_\sigma^2C_W^8$, one unit is for the tail form $q^{\mathrm{tail}}$ of step (ii), and $C_\sigma^2C_W^8$ is for step (iv). The second unit is not used by steps (i)--(iv).

The last inequality in \eqref{8.5:eq-room-bulk-a} comes from two facts. First, under the floor $C_W\ge C_1$ of \eqref{8.5:eq-room-lemma-a}, the exponential member of the curly bracket is at most $C_W^{-1}$, because $\min\{\hat\kappa,\hat\kappa_R\}\le\hat\kappa$. Second, the remaining members, together with the one unit $C_W^{-1}$ just obtained, are among the members of $r^M_{\mathrm{room}}$ in \eqref{8.5:eq-room-lemma-b}, since $2M(\Theta_0+8C_\nabla)=2M\Theta_0+16MC_\nabla$; the second unit of its summand $2C_W^{-1}$ is reserved for the profile tail of $h_2$.
\end{proof}

\begin{proof}[Proof of the second-derivative room lemma, continued: profile tail and envelope constant]
\label{8.5:prf-room-tail}
We continue with the notation of the preceding part of this proof, the part on the bulk, the anchor grouping and the summation by parts, which ends in the collected bound \eqref{8.5:eq-room-bulk-a}. That part is incomplete at the step marked $(\ast)$ there, and every appeal below to its steps \textup{(i)}--\textup{(v)} is subject to this. The second inequality of the collected bound \eqref{8.5:eq-room-bulk-a} is where the floor member $C_1$ of \eqref{8.5:eq-room-lemma-a} enters. For every $C\ge0$ one has
\begin{equation*}
e^{-\hat\kappa C/(64M)}\le e^{-\min\{\hat\kappa,\hat\kappa_R\}C/(64M)},
\end{equation*}
since $\min\{\hat\kappa,\hat\kappa_R\}\le\hat\kappa$; so the member $C_1$ of \eqref{8.5:eq-room-lemma-a}, written at the rate $\min\{\hat\kappa,\hat\kappa_R\}$, implies the corresponding inequality at the rate $\hat\kappa$, and under $C_1$ the exponential term of \eqref{8.5:eq-room-bulk-a} is absorbed into one unit of the constant summand of $r^M_{\mathrm{room}}$ in \eqref{8.5:eq-room-lemma-b}. This bounds the $q$-column
\begin{equation*}
\Bigl|\sum_{X\in E}Q_X[\mathcal{C}_{\cdot q},\mathcal{C}_{\cdot q}]\Bigr|
\end{equation*}
for every $q$ with $\operatorname{dist}(x_q,\mathfrak{A})\ge\mathsf{d}/4$, and in particular for every $q$ with $x_q\in B_2$. We multiply this bound by $b_0g_{k_0}^4|h_2^{\mathrm{in}}(q)|$ and sum over $q$ in the Wick pairing.

\emph{Step \textup{(vi)}: the profile tail of $h_2$.}
Let now $x_q\notin B_2$. No separation of $x_q$ from $\mathfrak{A}=B(x_1,3\mathsf{d}/2)$ is available, so we bound the whole $q$-column crudely by the growth of the background covariance (Lemma~\ref{8.5:lem-covariance-growth}, display \eqref{8.5:eq-covariance-growth-a}).

Every contributing $X$ falls into one of two classes:
\begin{itemize}
\item[(a)] it is a bulk term, and is then anchored in $\mathfrak{A}$ (step \textup{(i)}, in the preceding part, whose proof is incomplete at $(\ast)$);
\item[(b)] it is a tail term, counted as in step \textup{(iv)}.
\end{itemize}
For each of them we use, as in step \textup{(iv)} (in the preceding part, whose proof is incomplete at $(\ast)$), for every direction $\varphi$ admissible in the sense of \eqref{8.5:eq-constant-flux-forms-a}, the bound $|Q_X[\varphi,\varphi]|\le\mathfrak{m}^{\mathrm{reg}}_X[f_1]\,\|\varphi\|_X^2$ together with $\mathfrak{m}^{\mathrm{reg}}_X[f_1]\le\hat{\mathfrak{m}}_X\|f_1\|_\infty$. The bulk majorant masses are then summed anchor by anchor into $\sum_{c\in\mathfrak{A}}\mathfrak{a}^E_c$, and the tail majorant masses, by step \textup{(iv)}, into $\bar{\mathfrak{a}}^EN_1C_w'e^{-\hat\kappa C_W/(64M)}$. The direction $\mathcal{C}_{\cdot q}|_{\hat X}$ is measured in the norm $\|\cdot\|_X$ of \eqref{8.5:eq-constant-flux-forms-a} as in step \textup{(iv)}, with $C_\sigma$ the enlarged constant fixed together with the Gaussian frame \eqref{8.5:eq-gaussian-frame-a}. This gives $\|\mathcal{C}_{\cdot q}|_{\hat X}\|_X\le C_\sigma C_W^4\sigma_{\mathsf{d}}$ for every contributing $X$.

Three further facts are used. By step \textup{(v)} (in the preceding part, whose proof is incomplete at $(\ast)$), $\sum_{c\in\mathfrak{A}}\mathfrak{a}^E_c\le\bar{\mathfrak{a}}^EN_1$. Under $C_1$ the factor $C_w'e^{-\hat\kappa C_W/(64M)}$ does not exceed $1$, by the absorption recalled at the beginning of this part. Finally, $\mathfrak{U}_q=N_1\|f_1\|_\infty\sigma_{\mathsf{d}}^2$ by \eqref{8.5:eq-gaussian-frame-a}. Hence
\begin{align*}
\Bigl|\sum_{X\in E}Q_X[\mathcal{C}_{\cdot q},\mathcal{C}_{\cdot q}]\Bigr|
&\le\Bigl(\sum_{c\in\mathfrak{A}}\mathfrak{a}^E_c
+\bar{\mathfrak{a}}^EN_1C_w'e^{-\hat\kappa C_W/(64M)}\Bigr)
\|f_1\|_\infty\sup_X\|\mathcal{C}_{\cdot q}|_{\hat X}\|_X^2\\
&\le 2\bar{\mathfrak{a}}^EN_1\|f_1\|_\infty C_\sigma^2C_W^8\sigma_{\mathsf{d}}^2
=2C_\sigma^2C_W^8\,\bar{\mathfrak{a}}^E\mathfrak{U}_q .
\end{align*}

On the other hand, the weight $h_2^{\mathrm{out}}$ is exponentially small and decays with the distance. Put
\begin{equation*}
t_q:=\operatorname{dist}\bigl(x_q,B(x_2,\mathsf{d})\bigr).
\end{equation*}
Since $x_q\notin B_2=B(x_2,5\mathsf{d}/4)$, we have $t_q\ge\mathsf{d}/4$. As $\operatorname{supp}h_2^{\mathrm{in}}\subset B_2$, the weight at such $q$ is $h_2^{\mathrm{out}}(q)$. By Lemma~\ref{8.5:lem-corrected-weight}, in the form \eqref{8.5:eq-corrected-weight-a},
\begin{equation}\label{8.5:eq-room-tail-1}
|h_2^{\mathrm{out}}(q)|\le 3C_we^{-\hat\kappa(t_q/\ell-2)/(2M)}\|f_2\|_\infty .
\end{equation}

We count the bare plaquettes $q$ with $x_q\notin B_2$ in the shells
\begin{equation*}
\Sigma_n:=\{q:\ t\le t_q<t+\ell\},\qquad t=\mathsf{d}/4+n\ell,\quad n\ge0.
\end{equation*}
These shells cover all $q$ with $x_q\notin B_2$, because $t_q\ge\mathsf{d}/4$ for all of them. The barycentres $x_q$ of the plaquettes $q\in\Sigma_n$ lie in the spherical shell
\begin{equation*}
\frac{5\mathsf{d}}{4}+n\ell\le|x-x_2|<\frac{5\mathsf{d}}{4}+(n+1)\ell
\end{equation*}
of thickness $\ell$. The Euclidean ball of radius $R$ in four dimensions has volume $|B(R)|=\pi^2R^4/2$, so $d|B(R)|/dR=2\pi^2R^3$. This derivative is increasing in $R$, so the shell has volume at most $2\pi^2(5\mathsf{d}/4+(n+1)\ell)^3\ell$. The bare sites in a shell of thickness $\ell$ number at most $2\varepsilon^{-4}$ times its volume. There are six plaquettes per site. Writing
\begin{equation*}
\frac{5\mathsf{d}}{4}+(n+1)\ell=\frac{5\mathsf{d}}{4}\Bigl(1+\frac45\cdot\frac{n+1}{C_W}\Bigr)
\end{equation*}
by $\ell/\mathsf{d}=C_W^{-1}$, and then using $4/5\le1$ and $C_W\ge1$ (so that $1+(n+1)/C_W\le2+n$), we obtain
\begin{align}
\#\Sigma_n&\le 6\cdot2\varepsilon^{-4}\cdot2\pi^2\Bigl(\frac{5\mathsf{d}}{4}+(n+1)\ell\Bigr)^3\ell
\le 24\pi^2\varepsilon^{-4}\Bigl(\frac{5\mathsf{d}}{4}\Bigr)^3\ell\Bigl(1+\frac{n+1}{C_W}\Bigr)^3
\notag\\
&\le 24\pi^2\varepsilon^{-4}\Bigl(\frac{5\mathsf{d}}{4}\Bigr)^3\ell\,(2+n)^3 .
\label{8.5:eq-room-tail-2}
\end{align}
Since $|B_2|=\pi^2(5\mathsf{d}/4)^4/2$, we have $(5\mathsf{d}/4)^3=8|B_2|/(5\pi^2\mathsf{d})$, and therefore
\begin{equation*}
\Bigl(\frac{5\mathsf{d}}{4}\Bigr)^3\ell
=|B_2|\cdot\frac{8}{5\pi^2}\cdot\frac{\ell}{\mathsf{d}}
=|B_2|\cdot\frac{8}{5\pi^2C_W}.
\end{equation*}
Thus $\#\Sigma_n\le(192/5)\varepsilon^{-4}|B_2|C_W^{-1}(2+n)^3$.

On $\Sigma_n$ one has $t_q/\ell\ge C_W/4+n$. Hence, by \eqref{8.5:eq-room-tail-1},
\begin{equation*}
|h_2^{\mathrm{out}}(q)|\le 3C_we^{-\hat\kappa(C_W/4-2)/(2M)}e^{-\hat\kappa n/(2M)}\|f_2\|_\infty .
\end{equation*}
Put
\begin{equation*}
C_t:=\sum_{n\ge0}(2+n)^3e^{-\hat\kappa n/(2M)},
\end{equation*}
a number depending on $(M,\hat\kappa)$ only. Summing over $n$, and using $3\cdot192/5\le120$ together with
\begin{equation*}
e^{-\hat\kappa(C_W/4-2)/(2M)}=e^{\hat\kappa/M}e^{-\hat\kappa C_W/(8M)},
\end{equation*}
we get
\begin{align}
\sum_{x_q\notin B_2}|h_2^{\mathrm{out}}(q)|
&\le 3C_w\|f_2\|_\infty\cdot\frac{192}{5}\,\varepsilon^{-4}|B_2|C_W^{-1}
\cdot e^{-\hat\kappa(C_W/4-2)/(2M)}\sum_{n\ge0}(2+n)^3e^{-\hat\kappa n/(2M)}
\notag\\
&\le 120\,C_wC_te^{\hat\kappa/M}e^{-\hat\kappa C_W/(8M)}C_W^{-1}
\cdot\varepsilon^{-4}|B_2|\,\|f_2\|_\infty .
\label{8.5:eq-room-tail-3}
\end{align}

Multiply the column bound above by $b_0g_{k_0}^4|h_2^{\mathrm{out}}(q)|$ and sum over $x_q\notin B_2$ by \eqref{8.5:eq-room-tail-3}. The total of the tail is then at most
\begin{equation}\label{8.5:eq-room-tail-4}
\Bigl[120\,C_wC_te^{\hat\kappa/M}e^{-\hat\kappa C_W/(8M)}\cdot2C_\sigma^2C_W^8\Bigr]
\cdot C_W^{-1}\cdot\bar{\mathfrak{a}}^E\mathfrak{U}_q\cdot b_0g_{k_0}^4\,
\varepsilon^{-4}|B_2|\,\|f_2\|_\infty .
\end{equation}
We impose one more displayed member of the floor, in which $C_1'$ is the least number with the property stated in the second line:
\begin{equation}\label{8.5:eq-room-tail-a}
\begin{split}
&C_W\ge C_1',\qquad C_1'=C_1'(\hat\kappa,M,C_w',C_\sigma),\\
&240\,C_w'C_te^{\hat\kappa/M}e^{-\hat\kappa C/(8M)}C_\sigma^2C^8\le1
\quad\text{for all }C\ge C_1'.
\end{split}
\end{equation}
This member is one-sided: it is a lower bound on $C_W$ only. The number $C_1'$ depends neither on $K$, $s$, $m$ nor on the positions of the supports.

The bracket in \eqref{8.5:eq-room-tail-4} equals $240\,C_wC_te^{\hat\kappa/M}e^{-\hat\kappa C_W/(8M)}C_\sigma^2C_W^8$. We have $C_w\le C_w'$: in the definition of $C_w'$ in step \textup{(iv)}, one has $(1-e^{-\hat\kappa/(4M)})^{-1}\ge1$, and, evaluating the supremum at $r=1$,
\begin{equation*}
e^{\hat\kappa/(2M)}\sup_{r\ge1}16(1+r)^4e^{-\hat\kappa r/(4M)}\ge256\,e^{\hat\kappa/(4M)}\ge1 .
\end{equation*}
Hence, under \eqref{8.5:eq-room-tail-a}, the bracket is at most $1$, and the total of the tail is at most
\begin{equation*}
C_W^{-1}\cdot\bar{\mathfrak{a}}^E\mathfrak{U}_q\cdot b_0g_{k_0}^4\cdot
2\varepsilon^{-4}|B_2|\,\|f_2\|_\infty .
\end{equation*}
The tail therefore costs a second unit $C_W^{-1}\cdot1$ against the factor $2\varepsilon^{-4}|B_2|\|f_2\|_\infty$ provided for it in $\mathfrak{U}$ (see \eqref{8.5:eq-gaussian-frame-a}). Together with the unit absorbed at the beginning of this part, this accounts for the constant summand $2$ of $r^M_{\mathrm{room}}$ in \eqref{8.5:eq-room-lemma-b}.

\emph{Unit conversion.}
Under the floor member $C_W\ge16M$ of \eqref{8.5:eq-room-lemma-a}, and since $M\ge1$, we have $\mathsf{d}=C_W\ell\ge16\ell$, that is, $2\ell\le\mathsf{d}/8$. A cell of side $\ell$ has Euclidean diameter $2\ell$. Hence the cells meeting $B(x_1,3\mathsf{d}/2)$ lie in
\begin{equation*}
B(x_1,3\mathsf{d}/2+2\ell)\subset B(x_1,13\mathsf{d}/8).
\end{equation*}
The number $N_1$ counts the cells of the ball $\mathfrak{A}=B(x_1,3\mathsf{d}/2)$ used in step \textup{(v)}; these cells meet $B(x_1,3\mathsf{d}/2)$, are disjoint, and have volume $\ell^4$ each. Since $(13\mathsf{d}/8)/(5\mathsf{d}/4)=13/10$ and $(13/10)^4\le3$,
\begin{equation}\label{8.5:eq-room-tail-5}
\|f_1\|_\infty N_1\ell^4\le\|f_1\|_\infty\bigl|B(x_1,13\mathsf{d}/8)\bigr|
=\Bigl(\frac{13}{10}\Bigr)^4\|f_1\|_\infty|B_1|\le3\|f_1\|_\infty|B_1| .
\end{equation}
A Lipschitz function $f$ vanishing off $B_1=B(x_1,5\mathsf{d}/4)$ satisfies $\|f\|_\infty\le(5\mathsf{d}/4)\operatorname{Lip}f$. By the definition of the shape functional this gives
\begin{equation*}
\|f\|_\infty|B_1|\le\Theta(f;x_1,5\mathsf{d}/4)\,\|f\|_1 .
\end{equation*}
For $f_1$, whose shape functional is at most $\Theta_0$, we obtain $\|f_1\|_\infty|B_1|\le\Theta_0\|f_1\|_1$. Combined with \eqref{8.5:eq-room-tail-5} and $\mathfrak{U}_q=N_1\|f_1\|_\infty\sigma_{\mathsf{d}}^2$, this yields
\begin{equation*}
\mathfrak{U}_q\le3\Theta_0\,\ell^{-4}\|f_1\|_1\sigma_{\mathsf{d}}^2 .
\end{equation*}

\emph{The envelope constant.}
Consider the five summands $2M\Theta_0$, $16MC_\nabla$, $48C_\nabla\rho_\sharp$, $4\rho_\sharp\Theta_0$ and $2$. The first four are those of the bracket in \eqref{8.5:eq-room-bulk-a}; the last is the summand $2$ just obtained. Since $M,\rho_\sharp\ge1$, each of them is at most $M\rho_\sharp$ times, respectively, $2\Theta_0$, $16C_\nabla$, $48C_\nabla$, $4\Theta_0$ and $2$. Hence
\begin{equation}\label{8.5:eq-room-tail-6}
2M(\Theta_0+8C_\nabla)+48C_\nabla\rho_\sharp+4\rho_\sharp\Theta_0+2
\le M\rho_\sharp(64C_\nabla+6\Theta_0+2)
\le 64M\rho_\sharp(1+\Theta_0+C_\nabla).
\end{equation}
The left side of \eqref{8.5:eq-room-tail-6} is the bracket of the collected bound \eqref{8.5:eq-room-bulk-a} together with the constant summand $2$ obtained above, so \eqref{8.5:eq-room-tail-6} bounds these summands of $r^M_{\mathrm{room}}$ in \eqref{8.5:eq-room-lemma-b} by the single envelope $64M\rho_\sharp(1+\Theta_0+C_\nabla)$.
This completes the proof, subject to the step marked $(\ast)$ in the preceding part.
\end{proof}

\begin{corollary}[The room in the two-point witness's unit, free of the coupling; stated; proof incomplete at the step marked $(\ast)$]\label{8.5:cor-room-witness-unit}
Let the hypotheses of the second-derivative room lemma, Theorem~\ref{8.5:thm-room-lemma}, be in force,
together with its one-sided floor \eqref{8.5:eq-room-lemma-a} on $C_W$; in particular the data are:
\begin{itemize}
\item $f_1\in\mathfrak{F}(x_1,5\mathsf{d}/4;\Theta_0)$ is independent of the plane;
\item $f_2\in\mathfrak{F}(x_2,5\mathsf{d}/4;\Theta_0)$ and $\operatorname{supp}f_2\subset B(x_2,\mathsf{d})$;
\item the centres satisfy the separation condition under which that theorem is stated, whose upper
member is $|x_1-x_2|\le A_c\mathsf{d}$;
\item $h_2=h_2^{\mathrm{in}}+h_2^{\mathrm{out}}$ is the profile-corrected weight of $f_2$ of
Lemma~\ref{8.5:lem-corrected-weight}, with $h_2^{\mathrm{in}}$ its restriction to the plaquettes
centred in $B_2$ and $h_2^{\mathrm{out}}:=h_2-h_2^{\mathrm{in}}$.
\end{itemize}
For a weight $f=(f^P)_P$, one test function on the torus for each of the six planes $P$ (a
plaquette $q$ in the plane $P$ receiving the value $f^P(x(q))$), write
$|f|(x):=\max_P|f^P(x)|$ for its envelope and $\|f\|_1:=\int|f|(x)\,d^4x$ for the $L^1$-norm of
the envelope. Put
\begin{equation*}
\mathcal{K}^+(f_1,f_2):=\frac{6}{\pi^4}\int\!\!\int|f_1|(x)\,|f_2|(y)\,|x-y|^{-8}\,d^4x\,d^4y,
\qquad
\mathfrak{U}_W:=b_0\,g_{k_0}^4\,\mathcal{K}^+(f_1,f_2).
\end{equation*}
Here $k_0=K-s$ is the witness level.
Let $r_w$ be the radius of the ball $\{\sup|w|<r_w\}$ of complex sources on which clause~(b) of
Definition~\ref{8.5:def-profile-locality} asserts analyticity. Let
$C_{\mathrm{gr}}:=C_{\mathfrak{a}}(\kappa,L,M,\alpha)\,C_H(\alpha_0,M,\kappa)$ be the constant of
Definition~\ref{8.5:def-frozen-sourced-family}, the product of the
counting constant $C_{\mathfrak{a}}$ of the cell aggregates of the majorant masses and the Hessian
conversion constant $C_H$ of the majorants $\hat{\mathfrak{m}}^0_X$; here $\alpha$ is the exponent
in the bound \cite[(0.29)]{B-CMP109} on the residual terms, and $\alpha_0$ the field radius of the
complex fine-field domain, contained in the configuration spaces $\mathfrak{U}^c_k$ of
\cite[(1.18)]{B-CMP109}, on which that bound is used.
Define
\begin{equation}\label{8.5:eq-room-witness-unit-a}
\begin{gathered}
C_{\mathrm{conv}}(A_c,\Theta_0):=312\,\Theta_0^2\bigl(4(A_c+\tfrac52)\bigr)^8,\qquad
C_{\mathrm{irr}}\,(1+\Theta_0):=64\,C_{\mathrm{conv}}(A_c,\Theta_0)\,C_{\mathrm{gr}}\,r_w^{-1},\\
r^E_{\mathbb{X}}:=C_{\mathrm{irr}}\,
\frac{M\rho_\sharp\,(1+\Theta_0+C_\nabla)(1+\Theta_0)}{C_W}.
\end{gathered}
\end{equation}
Here $C_\nabla$ is the constant of the bound on the variation of the background covariance used in
Theorem~\ref{8.5:thm-room-lemma}, and the subscript $\mathbb{X}$ refers to the irreducible
remainder member of the two-point witness theorem, into whose bound $r^E_{\mathbb{X}}$ enters.
Let $Q_1,Q_2$ be the quadratic forms of \eqref{8.5:eq-gaussian-frame-b}, and let $Q_1^E$ be the
part of $Q_1$ carried by the extracted families $\operatorname{Irr}_1^E$ of
Definition~\ref{8.5:def-frozen-sourced-family}. Then
\begin{equation}\label{8.5:eq-room-witness-unit-1}
\bigl|\operatorname{Cov}(Q_1^E,Q_2)\bigr|\le r^E_{\mathbb{X}}\,\mathfrak{U}_W .
\end{equation}

For the weights of the two-point witness theorem, which lie in the classes
$\mathfrak{F}(x_i,\mathsf{d};\Theta_0^W)$ with that theorem's shape constant $\Theta_0^W$, $f_1$
being independent of the plane, the statement applies with
$\Theta_0:=(5/4)^5\Theta_0^W$. At this value
\begin{equation*}
C_{\mathrm{conv}}=312\,(5/4)^{10}\,(\Theta_0^W)^2\bigl(4(A_c+\tfrac52)\bigr)^8 .
\end{equation*}
\end{corollary}

For the profile-corrected weights $h_1,h_2$ of $f_1,f_2$ (Lemma~\ref{8.5:lem-corrected-weight}) the
two-point witness theorem uses the unit
$\mathfrak{U}_W(h):=b_0g_{k_0}^4\mathcal{K}^+(h_1,h_2)$ and shows
$\mathfrak{U}_W(h)\le2\,\mathfrak{U}_W$; this is not used below.
The number $C_{\mathrm{conv}}(A_c,\Theta_0)$ depends on $A_c$ and $\Theta_0$ only. With $r_w$ as in
the statement, no coupling constant enters
$r^E_{\mathbb{X}}$, and the coupling enters the right-hand side of
\eqref{8.5:eq-room-witness-unit-1} only through the factor $b_0g_{k_0}^4$ of $\mathfrak{U}_W$. If
clause~(b) holds with analyticity in the complex source $w$ only on the region near
$B(x_1,2\mathsf{d})$ of Definition~\ref{8.5:def-frozen-sourced-family}, the factor $C_Hr_w^{-1}$
contained in $r^E_{\mathbb{X}}$ is replaced by the grading letter $\mathfrak{w}_W$ of the two-point
witness theorem, which depends on $g_{k_0}$; this form is described after the proof.

\begin{proof}
Throughout, $\|\cdot\|_1$ is the $L^1$-norm of the envelope. For the plane-independent $f_1$ it is
simply $\int|f_1|$. For the plaquette weight $f_2$, the six planes are carried explicitly by the
count in the second step below and by the envelope in $\mathcal{K}^+$. Hence no further factor
coming from the planes arises.

\emph{The absolute bound.}
We apply Theorem~\ref{8.5:thm-room-lemma} (stated; proof incomplete). Under the
hypotheses above it gives \eqref{8.5:eq-room-lemma-b}, namely
\begin{equation}\label{8.5:eq-room-witness-unit-2}
\bigl|\operatorname{Cov}(Q_1^E,Q_2)\bigr|\le r^M_{\mathrm{room}}\,\bar{\mathfrak{a}}^E\,\mathfrak{U},
\qquad
r^M_{\mathrm{room}}\le 64\,\frac{M\rho_\sharp(1+\Theta_0+C_\nabla)}{C_W}.
\end{equation}
Here the sum $\sum_q|h_2(q)|$ in the unit $\mathfrak{U}$ of \eqref{8.5:eq-gaussian-frame-a} is read
as $\sum_q|h_2^{\mathrm{in}}(q)|+2\varepsilon^{-4}|B_2|\,\|f_2\|_\infty$, so that
\begin{equation}\label{8.5:eq-room-witness-unit-3}
\mathfrak{U}=b_0\,g_{k_0}^4\Bigl[\sum_q|h_2^{\mathrm{in}}(q)|
+2\varepsilon^{-4}|B_2|\,\|f_2\|_\infty\Bigr]\,\|f_1\|_\infty N_1\,\sigma_{\mathsf{d}}^2 .
\end{equation}
Here $N_1$ is the number of level-$k_0$ cells meeting $B(x_1,3\mathsf{d}/2)$, as in
\eqref{8.5:eq-gaussian-frame-a}.
In this bound the kernel of the background covariance enters only through the explicit scale
$\sigma_{\mathsf{d}}$ of \eqref{8.5:eq-gaussian-frame-a}. By Lemma~\ref{8.5:lem-covariance-upper},
this scale majorises the kernel on the separated range.

$(\ast)$ We convert $\bar{\mathfrak{a}}^E\mathfrak{U}$ into $\mathfrak{U}_W$ entirely in terms of
the weights $f_1,f_2$. Of the steps of this conversion, the fifth rests on the analyticity asserted
in clause~(b) of Definition~\ref{8.5:def-profile-locality}, whose proof is incomplete.
The corrected weight $h_2$ enters only through the two bounds
\eqref{8.5:eq-corrected-weight-a} of Lemma~\ref{8.5:lem-corrected-weight}:
\begin{itemize}
\item $|h_2(q)|\le 2\|f_2\|_\infty$ for every bare plaquette $q$. By
Lemma~\ref{8.5:lem-corrected-weight}, $h_2(q)$ is the derivative at $z=0$ of the running shift at
$q$ along the profile $zf_2$, divided by $\zeta'_{k_0}$. That derivative is at most
$(4/3)\|f_2\|_\infty$ in modulus by the linear bound on the running shift on which that lemma
rests, read there together with the lower bound $\zeta'_{k_0}\ge 2/3$ on the slope.
\item The tail of $h_2$ off $B_2$ decays with the distance. This tail is used in the profile-tail
step of the proof of Theorem~\ref{8.5:thm-room-lemma}, where it is paid for by the summand
$2\varepsilon^{-4}|B_2|\,\|f_2\|_\infty$ in \eqref{8.5:eq-room-witness-unit-3}.
\end{itemize}

\emph{The class of the weights of the two-point witness theorem.}
Let $f_i\in\mathfrak{F}(x_i,\mathsf{d};\Theta_0^W)$. Then
$\operatorname{supp}f_i\subset B(x_i,\mathsf{d})\subset B(x_i,5\mathsf{d}/4)$. The functional
$\Theta(\cdot\,;x_i,R)$ depends monotonically on $R$ through the factor
$R\,|B(x_i,R)|=\pi^2R^5/2$. This factor is multiplied by $(5/4)^5$ when $R=\mathsf{d}$ is replaced
by $5\mathsf{d}/4$. Hence $f_i\in\mathfrak{F}(x_i,5\mathsf{d}/4;(5/4)^5\Theta_0^W)$, which is the
last assertion of the statement. Substituting $\Theta_0=(5/4)^5\Theta_0^W$ in
\eqref{8.5:eq-room-witness-unit-a} gives the displayed value of $C_{\mathrm{conv}}$.

\emph{First step: the shape class.} We claim
\begin{equation}\label{8.5:eq-room-witness-unit-4}
\|f_1\|_\infty N_1\le 3\,\Theta_0\,\ell^{-4}\,\|f_1\|_1 .
\end{equation}
A cell of side $\ell$ has diameter $2\ell$. So every cell meeting $B(x_1,3\mathsf{d}/2)$ lies in
$B(x_1,3\mathsf{d}/2+2\ell)$. We have $2\ell\le\mathsf{d}/8$, since $\mathsf{d}=C_W\ell$ and the
floor \eqref{8.5:eq-room-lemma-a} contains $C_W\ge16M$ with $M\ge1$. Hence these cells lie in
$B(x_1,13\mathsf{d}/8)$. They are disjoint up to boundaries, each of volume $\ell^4$, and therefore
\begin{equation*}
N_1\ell^4\le|B(x_1,13\mathsf{d}/8)|=\bigl(\tfrac{13}{10}\bigr)^4|B_1|\le 3\,|B_1| .
\end{equation*}
Let $f$ be Lipschitz and vanish off $B_1=B(x_1,5\mathsf{d}/4)$. Every point of $B_1$ lies within
$5\mathsf{d}/4$ of the boundary sphere of $B_1$, on which $f$ vanishes by continuity, so
$\|f\|_\infty\le(5\mathsf{d}/4)\operatorname{Lip}f$. Hence
\begin{equation*}
\|f\|_\infty|B_1|\le\tfrac54\mathsf{d}\,\operatorname{Lip}f\,|B_1|
\le\Theta(f;x_1,5\mathsf{d}/4)\,\|f\|_1,
\end{equation*}
the last inequality by the definition of $\Theta(\cdot\,;x_1,R)$ at $R=5\mathsf{d}/4$. Since
$\Theta(f_1;x_1,5\mathsf{d}/4)\le\Theta_0$, this gives $\|f_1\|_\infty|B_1|\le\Theta_0\|f_1\|_1$.
Multiplying the two inequalities gives \eqref{8.5:eq-room-witness-unit-4}.

\emph{Second step: counting the sum over $q$.} The index $q$ runs over bare plaquettes. Assign each
bare plaquette to its least corner, so that each bare site carries six plaquettes, and consider the
bare sites at least one of whose plaquettes has its centre in $B_2$. Each such
site lies within $\varepsilon$ of $B_2$. The cube of side $\varepsilon$ at such a site has diameter
$2\varepsilon$, so it lies in $B(x_2,5\mathsf{d}/4+3\varepsilon)$. These cubes are disjoint up to
boundaries. Hence there are at most
\begin{equation*}
\varepsilon^{-4}\,|B(x_2,5\mathsf{d}/4+3\varepsilon)|
=\varepsilon^{-4}\Bigl(1+\frac{12\varepsilon}{5\mathsf{d}}\Bigr)^4|B_2|
\le 2\,\varepsilon^{-4}|B_2|
\end{equation*}
such sites. The last inequality holds because $\mathsf{d}=C_W\ell\ge16M\ell$ under the floor
$C_W\ge16M$, and $16M\ell\ge16\varepsilon$ as $\ell\ge\varepsilon$ and $M\ge1$. Therefore
$12\varepsilon/(5\mathsf{d})\le 3/20$ and $(23/20)^4\le2$. No new floor is introduced.

Now $h_2^{\mathrm{in}}$ vanishes off $B_2$ and $|h_2^{\mathrm{in}}(q)|\le|h_2(q)|\le2\|f_2\|_\infty$.
Therefore
\begin{equation*}
\sum_q|h_2^{\mathrm{in}}(q)|+2\varepsilon^{-4}|B_2|\,\|f_2\|_\infty
\le(6\cdot2\cdot2+2)\,\varepsilon^{-4}|B_2|\,\|f_2\|_\infty
=26\,\varepsilon^{-4}|B_2|\,\|f_2\|_\infty .
\end{equation*}
The weight $f_2$ is Lipschitz and vanishes off $B_2$, so the inequality of the first step gives
$\|f_2\|_\infty|B_2|\le\Theta(f_2;x_2,5\mathsf{d}/4)\|f_2\|_1$, and
$\Theta(f_2;x_2,5\mathsf{d}/4)\le\Theta_0$. Hence
\begin{equation}\label{8.5:eq-room-witness-unit-5}
\sum_q|h_2^{\mathrm{in}}(q)|+2\varepsilon^{-4}|B_2|\,\|f_2\|_\infty
\le 26\,\Theta_0\,\varepsilon^{-4}\,\|f_2\|_1 .
\end{equation}

\emph{Third step: the scale.} By \eqref{8.5:eq-gaussian-frame-a},
$\sigma_{\mathsf{d}}=2\varepsilon^4\pi^{-2}(\mathsf{d}/4)^{-4}$. Hence
\begin{equation}\label{8.5:eq-room-witness-unit-6}
\sigma_{\mathsf{d}}^2=4\,\varepsilon^8\pi^{-4}(\mathsf{d}/4)^{-8}
=4\cdot4^8\,\varepsilon^8\,\pi^{-4}\,\mathsf{d}^{-8}.
\end{equation}

\emph{Fourth step: a lower bound for the majorant kernel.} This is of the same kind as the lower
bound in the kernel lemma of the two-point witness theorem. There the balls have radius $\mathsf{d}$;
here they have radius $5\mathsf{d}/4$. The bound follows directly from the definition of
$\mathcal{K}^+$, read for $f_1,f_2$ themselves.

Take $x\in\operatorname{supp}f_1\subset B_1$ and $y\in\operatorname{supp}f_2\subset B_2$. By the
upper member of the separation condition,
\begin{equation*}
|x-y|\le|x_1-x_2|+\tfrac52\mathsf{d}\le D_{\max}:=(A_c+\tfrac52)\,\mathsf{d}.
\end{equation*}
So $|x-y|^{-8}\ge D_{\max}^{-8}$ on the product of the supports, and integration gives
\begin{equation}\label{8.5:eq-room-witness-unit-7}
\mathcal{K}^+(f_1,f_2)\ge\frac{6}{\pi^4}\,D_{\max}^{-8}\,\|f_1\|_1\|f_2\|_1 .
\end{equation}
No tail of a profile enters here, because $f_1,f_2$ are the compactly supported weights of the
statement.

\emph{Fifth step: the mass.} By \eqref{8.5:eq-gaussian-frame-b}, the main term's mass per
level-$k_0$ cell on unit constant flux per bare plaquette is $\mathfrak{n}=6(\ell/\varepsilon)^4$.
The maximal cell aggregate of the majorant masses of the extracted families satisfies
\begin{equation}\label{8.5:eq-room-witness-unit-8}
\bar{\mathfrak{a}}^E/\mathfrak{n}\le C_{\mathrm{gr}}\,r_w^{-1}.
\end{equation}
This is the unweighted domination stated in Definition~\ref{8.5:def-frozen-sourced-family},
obtained there from the analyticity asserted in clause~(b) of
Definition~\ref{8.5:def-profile-locality}, whose proof is incomplete,
together with the shape of the bound \cite[(0.29)]{B-CMP109} on the residual terms on
which the majorants $\hat{\mathfrak{m}}^0_X$ are modelled. The bound
\eqref{8.5:eq-room-witness-unit-8} contains no
coupling constant. The coupling-grading clause of the two-point witness theorem is not used.

\emph{Multiplication.} Combining \eqref{8.5:eq-room-witness-unit-3}--\eqref{8.5:eq-room-witness-unit-6}
we obtain
\begin{align*}
\bar{\mathfrak{a}}^E\,\mathfrak{U}
&\le b_0\,g_{k_0}^4\,\bar{\mathfrak{a}}^E\cdot26\,\Theta_0\,\varepsilon^{-4}\|f_2\|_1
\cdot3\,\Theta_0\,\ell^{-4}\|f_1\|_1\cdot4\cdot4^8\varepsilon^8\pi^{-4}\mathsf{d}^{-8}\\
&=\frac{\bar{\mathfrak{a}}^E}{\mathfrak{n}}\,b_0\,g_{k_0}^4\cdot6\Bigl(\frac{\ell}{\varepsilon}\Bigr)^4
\cdot312\cdot4^8\,\Theta_0^2\cdot\varepsilon^4\ell^{-4}\,\pi^{-4}\mathsf{d}^{-8}\,
\|f_1\|_1\|f_2\|_1\\
&=\frac{\bar{\mathfrak{a}}^E}{\mathfrak{n}}\,b_0\,g_{k_0}^4\cdot1872\cdot4^8\,\Theta_0^2\,
\pi^{-4}\mathsf{d}^{-8}\,\|f_1\|_1\|f_2\|_1 .
\end{align*}
In the second line we used $26\cdot3\cdot4=312$ and $\varepsilon^{-4}\varepsilon^8=\varepsilon^4$.
In the third line every power of $\varepsilon$ and of $\ell$ cancels, since
$(\ell/\varepsilon)^4\,\varepsilon^{4}\,\ell^{-4}=1$. The only coupling factor is
$b_0g_{k_0}^4$. It arises from the block law's factor $g_{k_0}^2/2$, which
\eqref{8.5:eq-gaussian-frame-b} places on $Q_1$ and on $Q_2$ alike, and it is the same prefactor
as that of the main term.

By \eqref{8.5:eq-room-witness-unit-7}, with $D_{\max}^8=(A_c+\frac52)^8\mathsf{d}^8$,
\begin{equation*}
\pi^{-4}\mathsf{d}^{-8}\,\|f_1\|_1\|f_2\|_1\le\tfrac16\,(A_c+\tfrac52)^8\,\mathcal{K}^+(f_1,f_2).
\end{equation*}
Therefore
\begin{equation}\label{8.5:eq-room-witness-unit-9}
\bar{\mathfrak{a}}^E\,\mathfrak{U}
\le\frac{1872}{6}\,4^8\,\Theta_0^2\,(A_c+\tfrac52)^8\cdot
\frac{\bar{\mathfrak{a}}^E}{\mathfrak{n}}\,\mathfrak{U}_W
=C_{\mathrm{conv}}(A_c,\Theta_0)\,\frac{\bar{\mathfrak{a}}^E}{\mathfrak{n}}\,\mathfrak{U}_W ,
\end{equation}
because $1872/6=312$ and $4^8(A_c+\frac52)^8=(4(A_c+\frac52))^8$. Thus $C_{\mathrm{conv}}$ is a
number depending on $(A_c,\Theta_0)$ only.

\emph{Conclusion.} Inserting \eqref{8.5:eq-room-witness-unit-9} and then
\eqref{8.5:eq-room-witness-unit-8} into \eqref{8.5:eq-room-witness-unit-2} gives
\begin{equation*}
\bigl|\operatorname{Cov}(Q_1^E,Q_2)\bigr|
\le64\,\frac{M\rho_\sharp(1+\Theta_0+C_\nabla)}{C_W}\,C_{\mathrm{conv}}(A_c,\Theta_0)\,
C_{\mathrm{gr}}\,r_w^{-1}\,\mathfrak{U}_W .
\end{equation*}
By the definition of $C_{\mathrm{irr}}$ in \eqref{8.5:eq-room-witness-unit-a}, the right-hand side
equals $r^E_{\mathbb{X}}\,\mathfrak{U}_W$. This proves \eqref{8.5:eq-room-witness-unit-1}.
\end{proof}

\emph{The form in which the two-point witness theorem uses this bound.} That theorem employs, for
its irreducible remainder member $\mathbb{X}$, a bound of the form
$|\mathbb{X}|\le2r_{\mathbb{X}}\mathfrak{U}_W$, where
$r_{\mathbb{X}}=r^E_{\mathbb{X}}+r^R_{\mathbb{X}}$, with $r^R_{\mathbb{X}}$ the share of the
remaining families; the present corollary supplies $r^E_{\mathbb{X}}$. The extracted part
$r^E_{\mathbb{X}}$ has there the form
\begin{equation*}
C_{\mathrm{irr}}\,\mathfrak{w}\,M\rho_\sharp(1+\Theta_0+C_\nabla)(1+\Theta_0)/C_W ,
\end{equation*}
with $\mathfrak{w}$ one of that theorem's two grading letters $\mathfrak{w}_H$, $\mathfrak{w}_W$.
The bound \eqref{8.5:eq-room-witness-unit-1} has exactly
this form with $\mathfrak{w}=\mathfrak{w}_H:=1$, and no coupling enters anywhere.

The dependence on $\Theta_0$ and $A_c$ matches the factor $C_{\mathrm{irr}}(1+\Theta_0)$ of that form
only in notation. The constant $C_{\mathrm{irr}}$ defined in \eqref{8.5:eq-room-witness-unit-a}
carries the factors $\Theta_0^2/(1+\Theta_0)$ and $(A_c+\frac52)^8$. This is harmless, because
$C_0$ is chosen after $\Theta_0$ and $A_c$.

\emph{The form of clause~(b) with analyticity only near $B(x_1,2\mathsf{d})$.} If clause~(b) of
Definition~\ref{8.5:def-profile-locality} holds with analyticity in the complex source $w$ only on
the region near $B(x_1,2\mathsf{d})$ of Definition~\ref{8.5:def-frozen-sourced-family},
then the factor $r_w^{-1}C_H$ is replaced by the grading letter $\mathfrak{w}_W$ of the two-point
witness theorem,
\begin{equation*}
\mathfrak{w}_W=p_0(g_{k_0})^4\varrho_J^{-2},
\end{equation*}
where $p_0(\cdot)$ is the degree function and $\varrho_J$ is the analyticity radius, in the field
fluctuations, of the sourced terms on that region. It then comes with the one-sided floor
\begin{equation*}
C_W\ge C_0\,p_0(g_{K-s})^4\varrho_J^{-2}
\end{equation*}
that the two-point witness theorem attaches to this reading. This is the only way in which a
dependence on $g_{K-s}$ can enter $C_0$.

\emph{Constants.} The floor $C_W\ge C_0$ is one-sided; no cap in the sense of
Definition~\ref{2.3:def-cap} is imposed. Here $C_0$ is chosen after
$M,\kappa,\kappa_R,\rho_\sharp,\Theta_0,C_\nabla,C_\sigma,A_c,r_w$, where $C_\sigma$ is the
constant of the growth bound of the background covariance in units of $\sigma_{\mathsf{d}}$ used in
the tail steps of the proof of Theorem~\ref{8.5:thm-room-lemma}. For members of the sourced family
localised on the larger cubes of \cite{B-CMP119}, the independence of these constants from the
position rests, as in that theorem, on the hypothesis named there that the side of every such cube
is a power of $L$. The constant $C_0$ is chosen so that
$r^E_{\mathbb{X}}$ satisfies the bound $\le1/192$ in which the two-point witness theorem uses it.

\emph{The coupling-grading clause.} This optional clause of the two-point witness theorem is not
used here. Granted as an additional factor $g_{k_0}^2$ in the dependence of the
fluctuation-generated terms on the source, it would multiply $r^E_{\mathbb{X}}$ by a further
factor $g_{k_0}^2$.

\emph{The meaning of ``relative to the main term''.} Suppose $f_1,f_2$ are scalar and non-negative,
which is the class for which the two-point witness theorem states its bound relative to the main
term. Then the identification of $\mathcal{K}^+$ with the continuum kernel form for such weights,
in the kernel lemma of that theorem, gives $\mathcal{K}^+(f_1,f_2)=\mathcal{K}_s(f_1,f_2)$, the
continuum kernel form of the main term. In that case \eqref{8.5:eq-room-witness-unit-1} reads
\begin{equation*}
\bigl|\operatorname{Cov}(Q_1^E,Q_2)\bigr|\le r^E_{\mathbb{X}}\,b_0\,g_{k_0}^4\,
\mathcal{K}_s(f_1,f_2).
\end{equation*}
This is relative to the main term $\zeta'^2_{K-s}\,b_0\,g_{k_0}^4\,\mathcal{K}_s(f_1,f_2)$ of the
two-point witness theorem, up to the slope factor $\zeta'^2_{K-s}\in[4/9,16/9]$, exactly as for
every other member of its error.

Likewise, suppose the weights lie in the cone $\mathcal{K}_s\ge\theta\,\mathcal{K}^+>0$ of that
relative bound, with $\theta>0$ the cone parameter. Then the same holds with an additional factor
$\theta^{-1}$.

This bound does not address the comparison of $\mathcal{K}_s(f_1,f_2)$ with the lattice form
$\mathcal{K}^{\mathrm{bg}}(h_1,h_2)$ actually produced at level $k_0$. That comparison is the
two-point witness theorem's own lattice-kernel comparison and is not proved again here. The
profile-corrected weights $h_i$ need be neither scalar nor non-negative.

For general real or non-scalar weights, the main term may vanish while $\mathcal{K}^+>0$. The
examples of a vanishing main term in the kernel lemma exhibit this for orthogonal planes and for
directions outside the cone. For such weights only the statement relative to $\mathcal{K}^+$ is
meaningful, exactly as in the additive error bound of the two-point witness theorem. No step above
uses a lower bound on an individual kernel entry $|\mathcal{C}_{pq}|$ or on the kernel maximum
$\mathsf{S}_q$ of \eqref{8.5:eq-gaussian-frame-a}.

\begin{proposition}[Covariance of the remainder column, small in the running coupling]
\label{8.5:prop-remainder-covariance}
Work at the witness level $k_0=K-s$, in the Gaussian frame
\eqref{8.5:eq-gaussian-frame-a}--\eqref{8.5:eq-gaussian-frame-b}: test functions $f_1,f_2$,
balls $B_1,B_2$, background covariance $\mathcal{C}=\mathcal{C}^{\mathrm{bg}}$, scale
$\sigma_{\mathsf{d}}$, units $\mathfrak{U}_q$ and $\mathfrak{U}$, and main-term mass per cell
$\mathfrak{n}=6(\ell/\varepsilon)^4$. The profile-corrected weight $h_2$ is that of
\eqref{8.5:eq-corrected-weight-a}, split as $h_2=h_2^{\mathrm{in}}+h_2^{\mathrm{out}}$, where
$h_2^{\mathrm{in}}(q)=h_2(q)$ if the centre $x_q$ of the bare plaquette $q$ lies in $B_2$ and
$h_2^{\mathrm{in}}(q)=0$ otherwise. The masses $\mathfrak{m}^{\mathrm{reg}}_X$, $\hat{\mathfrak{m}}_X$ and
the norms $\|\cdot\|_X$ are those of \eqref{8.5:eq-constant-flux-forms-a}.

Write $Q_1^R$ for the part of $Q_1$ formed by the terms of the remainder class
$\operatorname{Irr}_1^R$. These are the terms carrying the source that are not extracted into
the class $\operatorname{Irr}_1^E$. Write $\mathfrak{a}^R_c$ for their cell-aggregate majorant
mass at the cell $c$, and put $\bar{\mathfrak{a}}^R:=\sup_c\mathfrak{a}^R_c$.

Assume the following.
\begin{enumerate}
\item[(a)] The bound of \cite[(2.31)]{B-CMP119} on the remainder terms $R^{(j)}(X)$ of step $j$
localised on $X$:
\begin{equation}\label{8.5:eq-remainder-covariance-1}
|R^{(j)}(X)|\le g_j^{\kappa_0}\,(L^j\eta)^4\,e^{-\kappa_R d_j(X)},\qquad j\le k_0 .
\end{equation}
Here $\kappa_R$ is the decay rate of that display. We assume that its exponent $\kappa_0$ is the
exponent $\kappa_0\ge7$ of the statement of the correlation theorem. No relation between
$\kappa_0$ and $\kappa$ is assumed.
\item[(b)] Every family of the level-$k_0$ format of the correlation theorem that carries the
source and is not extracted consists of terms of the kind bounded in
\eqref{8.5:eq-remainder-covariance-1} (stated; proof incomplete; subject to the condition on the
running shift $\zeta'_k$ stated with it).
\item[(c)] The comparison of running couplings: for $j\le k_0$,
$x_j\ge x_{k_0}+\lambda(k_0-j)-2c_2$, where $\lambda$ is the positive slope constant of that
comparison (stated; proof incomplete).
\item[(d)] Profile locality and analyticity of sourced terms, \eqref{8.5:eq-profile-locality-a}
of Definition~\ref{8.5:def-profile-locality} (stated; proof incomplete at one step). The terms
are analytic in $w$ on the disc of radius $r_w$ given there.
\item[(e)] The growth bound of Lemma~\ref{8.5:lem-covariance-growth},
\eqref{8.5:eq-covariance-growth-a}, with $C_\sigma$ as enlarged in the Gaussian frame.
\item[(f)] The radius-tail bound of Lemma~\ref{8.5:lem-radius-tails},
\eqref{8.5:eq-radius-tails-a}, for the remainder class. It is taken at the rate
$\hat\kappa_R:=\kappa_R/2$, with \eqref{8.5:eq-remainder-covariance-1} read on the fine-field
configuration spaces exactly as Lemma~\ref{8.5:lem-radius-tails} reads the majorant from which it
starts, and with the majorant masses replaced by those of the remainder class. It is obtained by
the same derivation, the factor $(L^j\eta)^4$ paying the count. We assume the one-sided
condition
\begin{equation}\label{8.5:eq-remainder-covariance-2}
\hat\kappa_R\ge 2\log(2C_{\mathsf{a}})+2 ,
\end{equation}
where $C_{\mathsf{a}}$ is the counting constant of that derivation; under it the derivation of
Lemma~\ref{8.5:lem-radius-tails} applies to the remainder class.
\end{enumerate}

Assume further:
\begin{itemize}
\item the floor \eqref{8.5:eq-room-lemma-a} of the second-derivative room lemma with all its
members, among them $C_W\ge16C_\nabla$, where $C_\nabla$ is the constant of
\eqref{8.5:eq-covariance-difference-b}, and the decay member $C_W\ge C^{\mathrm{dec}}_{\mathrm{reg}}$
displayed there; its member $C_1$ is read at the rate $\min\{\hat\kappa,\hat\kappa_R\}$;
\item the floor $C_1'$ of \eqref{8.5:eq-room-tail-a};
\item the geometric conditions on the centres $x_1,x_2$, whose upper member is
$|x_1-x_2|\le A_c\mathsf{d}$, with centre-separation constant $A_c$;
\item the budget for the profile tail of $h_2$ contained in $\mathfrak{U}$ as in the
second-derivative room lemma;
\item the witness level satisfies $x_{k_0}\ge4c_2$, which holds for members of $\gamma_W$, and
$x_{k_0}\ge1$.
\end{itemize}
Then the following hold.
\begin{equation}\label{8.5:eq-remainder-covariance-3}
|\operatorname{Cov}(Q_1^R,Q_2)|\le 3\,\bar{\mathfrak{a}}^R\,\mathfrak{U},
\qquad
\bar{\mathfrak{a}}^R/\mathfrak{n}\le C_R\,g_{K-s}^{\kappa_0-2}.
\end{equation}

The second bound holds for the unweighted majorant. The route weight
$\bar{\mathfrak{w}}^R:=r_w^{-1}$ is absorbed into
$C_R=C_R(d,L,M,\kappa_0,\kappa_R,\lambda,c_2,r_w)$, which is independent of the position and of
$K$, $s$ and $m$. If the optional grading gain for the remainder class holds, it multiplies the
right side by a further factor $2Cg_{K-s}^2$, where $C$ is the constant of that gain.

Relative to the unit $\mathfrak{U}_W=b_0g_{k_0}^4\mathcal{K}^+(f_1,f_2)$ of the two-point witness
theorem, where $\mathcal{K}^+$ is the envelope kernel form fixed with the conversion of
Corollary~\ref{8.5:cor-room-witness-unit} (whose proof is incomplete at one step), and with
$C_{\mathrm{conv}}(A_c,\Theta_0)$ as in \eqref{8.5:eq-room-witness-unit-a},
\begin{equation}\label{8.5:eq-remainder-covariance-4}
|\operatorname{Cov}(Q_1^R,Q_2)|\le 3\,C_{\mathrm{conv}}(A_c,\Theta_0)\,C_R\,g_{K-s}^{\kappa_0-2}\,
\mathfrak{U}_W .
\end{equation}
If the grading gain holds, the coefficient is
$6\,C_{\mathrm{conv}}(A_c,\Theta_0)\,C_R\,C\,g_{K-s}^{\kappa_0}$.

If $f_1,f_2$ are scalar and non-negative, the same bounds hold with $\mathfrak{U}_W$ replaced by
the main term $b_0g_{k_0}^4\mathcal{K}_s(f_1,f_2)$. In any unit $\mathfrak{U}'\ge\mathfrak{U}_W$
the coefficient is at most the same. In particular this applies to the coincident-scale unit of
the column of $\gamma_W$ used in the Gaussian extraction.
\end{proposition}

\begin{proof}
No moment condition on the aggregated forms is used: every term is bounded in absolute value.
Fix a bare plaquette $q$ and one component in the Lie algebra $\mathfrak{g}$. All sums over $X$
below run over the terms of $Q_1^R$.

\emph{Step 1: bulk.} Let $R_0$ be the bulk radius of the second-derivative room lemma, and
consider the terms with $r(X)\le R_0$. As in the bulk part of the proof of that lemma
(part~\ref{8.5:prf-room-bulk} of that proof, which is incomplete at one step), such a term has
$\hat{X}\subset B(x_1,3\mathsf{d}/2)$. For $x_q\in B_2$ the majorant of
\eqref{8.5:eq-constant-flux-forms-a} gives
\begin{equation}\label{8.5:eq-remainder-covariance-5}
|Q_X[\mathcal{C}_{\cdot q},\mathcal{C}_{\cdot q}]|
\le\mathfrak{m}^{\mathrm{reg}}_X[f_1]\,\|\mathcal{C}_{\cdot q}|_{\hat X}\|_X^2
\le\hat{\mathfrak{m}}_X\,\|f_1\|_\infty\,\sigma_{\mathsf{d}}^2 .
\end{equation}
The second inequality rests on two bounds:
\begin{itemize}
\item $\mathfrak{m}^{\mathrm{reg}}_X[f_1]\le\hat{\mathfrak{m}}_X\sup_{X^+}|f_1|$;
\item $\|\mathcal{C}_{\cdot q}|_{\hat X}\|_X\le\sigma_{\mathsf{d}}$.
\end{itemize}
The first holds by the majorant \eqref{8.5:eq-constant-flux-forms-a}, the profile $f_1$ being
read on $X^+\subset\hat X$ only, by item~(a) of Definition~\ref{8.5:def-profile-locality}.
The second of these is obtained exactly as in that bulk part. Its supremum part is
Lemma~\ref{8.5:lem-covariance-upper}, \eqref{8.5:eq-covariance-upper-a}. Its regularity part comes
from the H\"older-regularity bound for the range of the linearised minimiser $\bar{M}_{k_0}$ used
there, together with the difference bound \eqref{8.5:eq-covariance-difference-b} of
Lemma~\ref{8.5:lem-covariance-difference}, under the members $C_W\ge16C_\nabla$ and
$C_W\ge C^{\mathrm{dec}}_{\mathrm{reg}}$ of the floor assumed in the statement.

The anchor cells of the bulk terms lie in $B(x_1,3\mathsf{d}/2)$. At most $N_1$ cells meet this
ball, $N_1$ being the cell count that enters $\mathfrak{U}_q$ in \eqref{8.5:eq-gaussian-frame-a}.
Summing \eqref{8.5:eq-remainder-covariance-5} over the bulk terms therefore gives at most
\[
\sum_c\mathfrak{a}^R_c\,\|f_1\|_\infty\sigma_{\mathsf{d}}^2
\le N_1\bar{\mathfrak{a}}^R\|f_1\|_\infty\sigma_{\mathsf{d}}^2
=\bar{\mathfrak{a}}^R\mathfrak{U}_q,
\]
by the definition of $\mathfrak{U}_q$ in \eqref{8.5:eq-gaussian-frame-a}.

\emph{Step 2: tails.} The terms with $r(X)>R_0$ are treated as in the tail step of the proof of
the second-derivative room lemma (part~\ref{8.5:prf-room-bulk} of that proof, which is incomplete
at one step), with two changes:
\begin{itemize}
\item the radius tails are taken from hypothesis~(f) at the rate $\hat\kappa_R$;
\item the kernel is bounded on every cell through hypothesis~(e), with $C_\sigma$ as enlarged.
\end{itemize}
Adding Step~1, the $q$-column satisfies, for $x_q\in B_2$,
\begin{equation}\label{8.5:eq-remainder-covariance-6}
\Bigl|\sum_X Q_X[\mathcal{C}_{\cdot q},\mathcal{C}_{\cdot q}]\Bigr|
\le\bigl(1+C_we^{-\hat\kappa_RC_W/(64M)}C_\sigma^2C_W^8\bigr)\bar{\mathfrak{a}}^R\mathfrak{U}_q
\le2\,\bar{\mathfrak{a}}^R\mathfrak{U}_q .
\end{equation}
Here $C_w$ is the constant of \eqref{8.5:eq-radius-tails-a} at the rate $\hat\kappa_R$.
The last inequality holds because the floor $C_1$, read at $\min\{\hat\kappa,\hat\kappa_R\}$,
makes the second summand at most $1$.

\emph{Step 3: the profile tail of $h_2$ and the summation over $q$.} As in the summation step of
the proof of the second-derivative room lemma (part~\ref{8.5:prf-room-tail} of that proof, which is
incomplete at one step), the covariance is bounded by
\begin{equation}\label{8.5:eq-remainder-covariance-7}
|\operatorname{Cov}(Q_1^R,Q_2)|\le b_0g_{k_0}^4\sum_q|h_2(q)|\,
\Bigl|\sum_X Q_X[\mathcal{C}_{\cdot q},\mathcal{C}_{\cdot q}]\Bigr| .
\end{equation}
The sum over $q$ in \eqref{8.5:eq-remainder-covariance-7} splits into two parts.

For $x_q\in B_2$ we have $|h_2(q)|=|h_2^{\mathrm{in}}(q)|$. By
\eqref{8.5:eq-remainder-covariance-6} these plaquettes contribute to the right side of
\eqref{8.5:eq-remainder-covariance-7} at most
$2\bar{\mathfrak{a}}^R\mathfrak{U}_q\,b_0g_{k_0}^4\sum_q|h_2^{\mathrm{in}}(q)|$.

For $x_q\notin B_2$ we proceed as in the profile-tail step of the same proof
(part~\ref{8.5:prf-room-tail}, which is incomplete at one step). There the whole $q$-column
is bounded crudely through hypothesis~(e). The tail terms in $X$ are counted with the rate
$\hat\kappa_R$ of hypothesis~(f); the floor $C_1$, read at $\min\{\hat\kappa,\hat\kappa_R\}$,
serves this tail in $X$ for the remainder class. The column bound is weighted with the decaying
tail of $h_2^{\mathrm{out}}$ from \eqref{8.5:eq-corrected-weight-a}. Under the floor $C_1'$ these
plaquettes contribute at most
$\bar{\mathfrak{a}}^R\mathfrak{U}_q\,b_0g_{k_0}^4\,2\varepsilon^{-4}|B_2|\|f_2\|_\infty$.

By \eqref{8.5:eq-gaussian-frame-a},
\[
\mathfrak{U}=b_0g_{k_0}^4\Bigl[\sum_q|h_2^{\mathrm{in}}(q)|
+2\varepsilon^{-4}|B_2|\|f_2\|_\infty\Bigr]\mathfrak{U}_q .
\]
The first part is therefore at most $2\bar{\mathfrak{a}}^R\mathfrak{U}$ and the second at most
$\bar{\mathfrak{a}}^R\mathfrak{U}$. This proves the first inequality of
\eqref{8.5:eq-remainder-covariance-3}.

\emph{Step 4: the mass.} By hypothesis~(b), every term entering $\mathfrak{a}^R_c$ is a term
$R^{(j)}(X)$, $j\le k_0$, bounded by \eqref{8.5:eq-remainder-covariance-1}.

Fix a level-$j$ cell and sum the factor $e^{-\kappa_R d_j(X)}$ over the domains $X$ containing
it. This is the count carried out for hypothesis~(f) under
\eqref{8.5:eq-remainder-covariance-2}. It is bounded by a constant
$C_{\mathrm{cell}}=C_{\mathrm{cell}}(d,L,M,\kappa_R)$, independent of the position and of $K$,
$s$, $m$. Hence the terms of step $j$ contribute, by \eqref{8.5:eq-remainder-covariance-1}, at
most $C_{\mathrm{cell}}\,g_j^{\kappa_0}(L^j\eta)^4$ per level-$j$ unit cell.

Since $\varepsilon=L^{-K}$ and $\ell=L^{-s}$,
\[
\eta=\varepsilon/\ell=L^{-K}/L^{-s}=L^{-k_0},\qquad (L^j\eta)^4=L^{-4(k_0-j)} .
\]
This factor cancels exactly the number $L^{4(k_0-j)}$ of level-$j$ cells contained in one
level-$k_0$ cell. This power is marginal: it leaves no geometric gain in $k_0-j$ and no room. The
mass of the terms of step $j$ per level-$k_0$ cell, measured relative to $\mathfrak{n}$, is
therefore at most $C_{\mathrm{cell}}\,g_j^{\kappa_0}$.

The mass entering $\mathfrak{a}^R_c$ is that of the marked part of these terms. By the
analyticity in $w$ of item~(b) of Definition~\ref{8.5:def-profile-locality} (hypothesis~(d),
whose proof is incomplete at one step), Cauchy's estimate on the disc of radius $r_w$ bounds it by
$r_w^{-1}$ times the mass of the terms themselves.

For the fresh terms, $j=k_0$, this gives, for every cell $c$, a contribution to
$\mathfrak{a}^R_c/\mathfrak{n}$ of at most $C_{\mathrm{cell}}r_w^{-1}g_{k_0}^{\kappa_0}$.

For the aged terms that are not extracted, if present at level $k_0$, we must bound
$\sum_{j\le k_0}g_j^{\kappa_0}$. Put $t:=k_0-j$. By hypothesis~(c),
$x_j\ge x_{k_0}+\lambda t-2c_2$. Since $x_{k_0}\ge4c_2$, we have $2c_2\le x_{k_0}/2$, so
$x_j\ge x_{k_0}/2+\lambda t$. With $g_j^{\kappa_0}=x_j^{-\kappa_0/2}$ it follows that
$g_j^{\kappa_0}\le(x_{k_0}/2+\lambda t)^{-\kappa_0/2}$.

This function of $t$ is decreasing. Its sum over $t\ge0$ is therefore at most its value at $t=0$
plus its integral over $[0,\infty)$, which converges because $\kappa_0>2$:
\begin{align}
\sum_{t\ge0}\Bigl(\frac{x_{k_0}}{2}+\lambda t\Bigr)^{-\kappa_0/2}
&\le\Bigl(\frac{x_{k_0}}{2}\Bigr)^{-\kappa_0/2}
+\int_0^\infty\Bigl(\frac{x_{k_0}}{2}+\lambda t\Bigr)^{-\kappa_0/2}dt
\notag\\
&=2^{\kappa_0/2}x_{k_0}^{-\kappa_0/2}
+\frac{2^{\kappa_0/2}\,x_{k_0}^{1-\kappa_0/2}}{\lambda(\kappa_0-2)}
\notag\\
&\le2^{\kappa_0/2}\Bigl(1+\frac{1}{\lambda(\kappa_0-2)}\Bigr)x_{k_0}^{1-\kappa_0/2}
=:C_\lambda\,g_{k_0}^{\kappa_0-2}.
\label{8.5:eq-remainder-covariance-8}
\end{align}
In the second line the integral equals
$(x_{k_0}/2)^{1-\kappa_0/2}/(\lambda(\kappa_0/2-1))$. The last inequality uses
$x_{k_0}^{-\kappa_0/2}\le x_{k_0}^{1-\kappa_0/2}$, valid for $x_{k_0}\ge1$, and
$x_{k_0}^{1-\kappa_0/2}=g_{k_0}^{\kappa_0-2}$. The estimate is used for $\kappa_0\ge4$, which
holds since $\kappa_0\ge7$.

Every coupling $g_j$, $j\le k_0$, has thus been bounded through the monotone comparison~(c) by a
power of $g_{k_0}=g_{K-s}$. The sum \eqref{8.5:eq-remainder-covariance-8} contains the fresh
terms as its term $t=0$. Hence
\[
\bar{\mathfrak{a}}^R/\mathfrak{n}\le C_{\mathrm{cell}}\,r_w^{-1}\,C_\lambda\,g_{K-s}^{\kappa_0-2},
\]
which is the second inequality of \eqref{8.5:eq-remainder-covariance-3} with
$C_R:=C_{\mathrm{cell}}\,r_w^{-1}\,C_\lambda$. This constant depends only on
$(d,L,M,\kappa_0,\kappa_R,\lambda,c_2,r_w)$.

If the grading gain holds, it multiplies the marked mass by a further $2Cg_{K-s}^2$.

\emph{Step 5: the unit of the two-point witness theorem.} The conversion steps of
Corollary~\ref{8.5:cor-room-witness-unit} (whose proof is incomplete at one step) run entirely in
$f_1,f_2$ and are linear in the mass. Applied with $\bar{\mathfrak{a}}^R/\mathfrak{n}$ in place of
$\bar{\mathfrak{a}}^E/\mathfrak{n}$, they give
\[
\bar{\mathfrak{a}}^R\mathfrak{U}\le C_{\mathrm{conv}}(A_c,\Theta_0)\,
(\bar{\mathfrak{a}}^R/\mathfrak{n})\,\mathfrak{U}_W ,
\]
with no coupling factor. Inserting the two inequalities of
\eqref{8.5:eq-remainder-covariance-3} gives \eqref{8.5:eq-remainder-covariance-4}. With the
grading gain, the coefficient becomes
\[
3C_{\mathrm{conv}}(A_c,\Theta_0)\,C_R\cdot2Cg_{K-s}^2\cdot g_{K-s}^{\kappa_0-2}
=6\,C_{\mathrm{conv}}(A_c,\Theta_0)\,C_R\,C\,g_{K-s}^{\kappa_0}.
\]

For scalar non-negative $f_1,f_2$ we have $\mathcal{K}^+(f_1,f_2)=\mathcal{K}_s(f_1,f_2)$, by the
comparison of kernel forms that belongs to the two-point witness theorem. There
$\mathfrak{U}_W$ is therefore the main term $b_0g_{k_0}^4\mathcal{K}_s(f_1,f_2)$.

Finally, if $\mathfrak{U}'\ge\mathfrak{U}_W$, then a bound of $|\operatorname{Cov}(Q_1^R,Q_2)|$ by
a non-negative multiple of $\mathfrak{U}_W$ is also a bound by the same multiple of
$\mathfrak{U}'$.
\end{proof}

\begin{remark}
The bound of Proposition~\ref{8.5:prop-remainder-covariance} concerns the remainder terms present
at the witness level in the format of the correlation theorem. It uses no history of healed
large-field regions and no Schwarz inequality in two slots. The exponent it produces is
$\kappa_0-2$, or $\kappa_0$ if the grading gain holds, where $\kappa_0\ge7$ is the exponent of
\cite[(2.31)]{B-CMP119} as assumed in hypothesis~(a). The exponent $\kappa_0-2$ is the exponent of
the remainder coefficient of the two-point witness theorem. The bound
\eqref{8.5:eq-remainder-covariance-4} carries no factor $(C_0L)^8$ and uses no condition
$C_W<C_0L$.

Suppose the aged remainder terms, $j<k_0$, are included in $b_{k+1}$
(Notation~\ref{0.2:not-glossary}). They then belong to the extracted class
$\operatorname{Irr}_1^E$ and leave this column. The column consists of
the fresh terms $R^{(k_0)}$ alone, and the bound holds unchanged: it is the term $t=0$ of
\eqref{8.5:eq-remainder-covariance-8}. Hypothesis~(b) carries a condition on the running shift
$\zeta'_k$; that condition is part of (b) as assumed here and is not discharged in this
proposition.
\end{remark}

\begin{remark}[Irrelevance at second order as vanishing column sums]\label{8.5:rem-second-order-irrelevance}
The content of irrelevance at second order is the following. The column sums of the marked residual
form over a period block vanish: the anisotropic block mean and, after the extraction, the scalar
block mean. The vanishing of these block means is the subject of
Proposition~\ref{8.5:prop-zero-block-mean} (stated; proof incomplete at the step marked $(\ast)$
there), whose identity~\eqref{8.5:eq-zero-block-mean-a} expresses the zero block mean of the marked
residual form. The row sums of the form do not vanish.

The room of order $C_W^{-1}$ in the second-derivative room lemma
(Theorem~\ref{8.5:thm-room-lemma}, stated; proof incomplete at the step marked $(\ast)$ there) is
the sum of two contributions.
\begin{itemize}
\item[(a)] A coefficient that is $M$-periodic and has zero mean is averaged against the profile and
against the kernel. Both of these are smooth at the scale $\mathsf{d}=C_W\ell$, that is, $C_W$ cells.
\item[(b)] The decay estimate of the covariance contributes the term $\rho_\sharp/C_W$ for the
locality radius. Here $\rho_\sharp$ is the mass-weighted mean radius.
\end{itemize}
\end{remark}

\begin{remark}[Completeness fails only by a scalar of main-term type]\label{8.5:rem-completeness-failure}
Write $k_0=K-s$, and let $\mathcal{R}^{E,\zeta}$ be the residual family of the $\beta$-extracted species
$\operatorname{Irr}_1^E$ at constant shift $\zeta$. The statement of Proposition~\ref{8.5:prop-extraction-complete},
whose proof is incomplete at one step, can fail only by a $W_4$-scalar, block-constant form, that is, by a
block-summed form of $\mathcal{R}^{E,\zeta}$ equal to a non-zero multiple of $\langle\cdot,\cdot\rangle$. Such a
residue is of the type of the main term. It can be accounted for as a change of the running shift
$\zeta'_{k_0}$ by at most $\bar{\mathfrak{a}}^E/\mathfrak{n}$. Along the route on which the majorant mass
$\bar{\mathfrak{a}}^E$ is controlled in powers of the running coupling, this bound is
$O(g_{K-s}^2)$. The price of this accounting is a re-identification of the prefactor in the statement of the
two-point witness theorem. The re-identification is confined to the statement of that theorem.
The scalar $\delta_R$ of the remainder species (Remark~\ref{8.5:rem-remainder-species}) is a shift of
exactly this kind. Its size is at most $C_R\,g_{K-s}^{\kappa_0-2}$, where $C_R$ is a constant of the
remainder-column bound of Proposition~\ref{8.5:prop-remainder-covariance} and $\kappa_0$ the exponent of that
bound, and it is
contained in the estimate of that proposition. A family escaping the membership hypothesis of that proposition
would give a residue of the same type.
\end{remark}

\begin{proof}
Lemma~\ref{8.5:lem-block-sum-scalar} uses only the $\Gamma$- and $W_4$-covariance of the family to which it is
applied. It does not use the conclusion of Proposition~\ref{8.5:prop-extraction-complete}. The family
$\mathcal{R}^{E,\zeta}$ is $\Gamma$- and $W_4$-covariant by the $\Gamma$- and $W_4$-covariance statements from
which Lemma~\ref{8.5:lem-block-sum-scalar} is proved as an implication, so the lemma applies to it. For every
period block
$\mathfrak{B}$, its part (b) gives that the sum over $c\in\mathfrak{B}$ of the cell aggregates
$q_{\mathcal{R}^{E,\zeta}}(c)$ is independent of $\mathfrak{B}$, equals the $\Gamma$-class sum
$\mathfrak{q}_M(\mathcal{R}^{E,\zeta})$, and is a multiple $\delta(\mathcal{R}^{E,\zeta})\langle\cdot,\cdot\rangle$.
Every component other than the $W_4$-scalar therefore vanishes by covariance alone.
Proposition~\ref{8.5:prop-extraction-complete}, whose proof is incomplete at one step, asserts that
$\mathfrak{q}_M(\mathcal{R}^{E,\zeta})=0$, that is, $\delta(\mathcal{R}^{E,\zeta})=0$. If it fails, what remains
is the block-constant $W_4$-scalar form $\delta(\mathcal{R}^{E,\zeta})\langle\cdot,\cdot\rangle$.
By the bound of the cell aggregates of the $\beta$-extracted species by their majorant masses, each
$q_{\mathcal{R}^{E,\zeta}}(c)$ has size at most $\mathfrak{a}^E_c\le\bar{\mathfrak{a}}^E$. The block
$\mathfrak{B}$ contains $M^4$ cells, so
\begin{equation}\label{8.5:eq-completeness-failure-1}
  |\delta(\mathcal{R}^{E,\zeta})|\,M^{-4}\le\bar{\mathfrak{a}}^E .
\end{equation}
A multiple of $\langle\cdot,\cdot\rangle$ that is constant over blocks has the form of the main term, whose mass
per cell is $\mathfrak{n}$. Dividing the per-cell size \eqref{8.5:eq-completeness-failure-1} by $\mathfrak{n}$
shows that the residue is a change of $\zeta'_{k_0}$ by at most $\bar{\mathfrak{a}}^E/\mathfrak{n}$; absorbing
it into $\zeta'_{k_0}$ is the re-identification of the prefactor described above.

For the remainder species $\operatorname{Irr}_1^R$, from which no $\beta$ is extracted,
Remark~\ref{8.5:rem-remainder-species} obtains from the same lemma the block sum $\delta_R\langle\cdot,\cdot\rangle$
with $|\delta_R|M^{-4}\le\bar{\mathfrak{a}}^R$, a shift of the kind just described, which
Proposition~\ref{8.5:prop-remainder-covariance} bounds by size as part of the remainder column.

Finally, suppose a family of the level-$k_0$ format were neither extracted nor of the species covered by the
membership hypothesis of Proposition~\ref{8.5:prop-remainder-covariance}. The argument of the first paragraph
uses only the $\Gamma$- and $W_4$-covariance of the family, and such a family has it by the same covariance
statements; the argument therefore applies unchanged, and the block-summed form of the family is again a
block-constant $W_4$-scalar of main-term type.
\end{proof}

\begin{remark}[What the bounds of Ba{\l}aban's papers give for a term created at an earlier step]
Let a localised term with localisation domain $X$ be created at step $j$ and read at the witness
level $k_0$, so that its age
at level $k_0$ is $k_0-j$. Assume that, for fixed $j$ and $X$, the term is one analytic function,
independent of $k$, obeying \cite[(0.29)]{B-CMP109} with one constant for all $k\ge j$; this
hypothesis is not proved (see Remark~\ref{8.5:rem-per-term-annihilation}). Under it,
\cite[(0.29)]{B-CMP109} gives for the marginal residue of the term, the part of it seen by
constant fluxes, only the relative smallness
\begin{equation*}
  L^{-\alpha(k_0-j)},
\end{equation*}
where $\alpha>0$ is the exponent by which the power $4+\alpha$ of the small factor in
\cite[(0.29)]{B-CMP109} exceeds the marginal power $4$; for $j=k_0$ this factor equals $1$. The
bound
\cite[(1.18)]{B-CMP109} imposes no smallness on this residue. The bound \cite[(2.31)]{B-CMP119}
contains no exponent $\alpha$ at all: the terms it controls are marginal by their shape and small
only through the coupling.
\end{remark}

\begin{remark}[No contribution escapes through large domains]
The coefficient extracted from the localised terms into the coupling is the quantity defined by
\cite[(3.61)]{B-CMP119}, with the sum there extended to all localisation domains as in
\cite[\S3]{B-CMP119}; hence no part of it is lost through large domains. The spatial tails of
these sums are those controlled by Lemma~\ref{8.5:lem-radius-tails}.
\end{remark}

\begin{remark}[Independence of the constant-flux form from kernel and level]
\label{8.5:rem-flux-form-independence}
Let $T$ be a localised term as in Lemma~\ref{8.5:lem-constant-flux-forms}: gauge invariant,
analytic on an open neighbourhood of $\mathbf{1}$, and depending on $U|_X$ with $X\subset\hat{X}$,
where the filled hull $\hat{X}$ is a box of the bare lattice not wrapping the torus. Let $q_T$ be
its constant-flux form, the restriction of the flux form $Q_T$ of that lemma to fluxes that are
constant on $\hat{X}$, that is, take one value $u(\pi)$ on all plaquettes of each plane
$\pi\in\mathfrak{O}$; so $q_T$ is a quadratic form on $\mathbb{R}^{\mathfrak{O}}$. Then $q_T$
does not depend on the choice of the kernel through which a
constant flux is realised by a field. Assume moreover the constant-flux statement for the
linearised minimiser, the constant row-sum identity of
Proposition~\ref{8.5:prop-row-sum-identity} (stated; proof incomplete), and the kernel
decay bound fixed in Definition~\ref{8.5:def-kernel-decay-torus}. Then, for
every level $j$, the value of $q_T$ on the
fine flux produced from a constant coarse flux at level $j$ equals $\eta_j^4$ times its value on
that constant flux. Here $\eta_j$ denotes the spacing ratio appearing in
\eqref{8.5:eq-row-sum-identity-a}, the bare spacing in units of the spacing of the level-$j$
lattice.
\end{remark}

\begin{proof}
We first show that $q_T$ does not depend on the kernel. By clause~(b) of
Lemma~\ref{8.5:lem-constant-flux-forms}, $\operatorname{Hess} T[B,B']$ depends on the pair
$(B,B')$ only through the restrictions $(dB,dB')|_{\hat{X}}$ of their fluxes, and clause~(d) of
the same lemma gives the restriction of the resulting flux form $Q_T$ to constant fluxes; hence
two fields whose fluxes on $\hat{X}$ are the same constant flux give the same value of $q_T$,
and since a kernel enters only as a means of producing a field with a prescribed flux, $q_T$
carries no dependence on which kernel is chosen.

We next compare coarse and fine fluxes. Fix a level $j$ and a constant
$u\in\mathbb{R}^{\mathfrak{O}}$, that is, one value for each plane, regarded as the bounded
$2$-cochain $\mathfrak{P}\mapsto u(\pi(\mathfrak{P}))$ on the level-$j$ lattice, again denoted
$u$. Let $K_j(p,\mathfrak{P})$ be the whole-lattice kernel
of $\bar{M}_j$. The kernel decay bound fixed in Definition~\ref{8.5:def-kernel-decay-torus},
assumed in the statement, is stated at the level $K-s$ for every depth $s$, hence at every
level $j$, with constants $C_\Gamma,\delta_\Gamma,\ell_\Gamma$ depending on $(d,L,M)$ only; it
gives
\begin{equation*}
  |K_j(p,\mathfrak{P})| \le C_\Gamma\,\eta_j^2
  e^{-\delta_\Gamma \operatorname{dist}(p,\mathfrak{P})/\ell_\Gamma},
\end{equation*}
the distance being measured in cells of the level-$j$ lattice. Since $u$ is bounded, this bound
makes the sum
\begin{equation*}
  \sum_{\mathfrak{P}} K_j(p,\mathfrak{P})\,u(\pi(\mathfrak{P}))
\end{equation*}
absolutely convergent for every bare plaquette $p$, so $\bar{M}_j(u)$ is defined for this datum,
which is not of bounded support.

The constant-flux statement for the linearised minimiser that accompanies the constant row-sum
identity, together with Proposition~\ref{8.5:prop-row-sum-identity} (stated; proof incomplete
at the step marked $(\ast)$ there), identifies the field produced from $u$ by the linearised
minimiser with the output of this kernel. The row-sum identity
\eqref{8.5:eq-row-sum-identity-a} then gives, for every bare plaquette $p$,
\begin{equation}\label{8.5:eq-flux-form-independence-1}
  \bigl(\bar{M}_j(u)\bigr)(p)
  = \sum_{\mathfrak{P}} K_j(p,\mathfrak{P})\,u(\pi(\mathfrak{P}))
  = \eta_j^2\,u(\pi(p)).
\end{equation}
So the fine flux produced from the constant coarse flux $u$ is itself constant on $\hat{X}$
and equals $\eta_j^2 u$, exactly and with no error term.

By the first part of the proof, the value of the form on this fine flux is $q_T$ evaluated at
$\eta_j^2 u$. Since $q_T$ is the restriction of the quadratic form $Q_T$, it is homogeneous
of degree two, so
\begin{equation}\label{8.5:eq-flux-form-independence-2}
  q_T\bigl[\eta_j^2 u\bigr] = \eta_j^4\, q_T[u].
\end{equation}
Thus the value of $q_T$ on the fine flux produced from $u$ at level $j$ and its value on $u$
differ exactly by the factor $\eta_j^4$.
\end{proof}

\begin{remark}[No entrywise lower bound on the flux two-point function]
\label{8.5:rem-no-entrywise-bound}
Write $\mathfrak{O}$ for the set of the six coordinate planes, $\pi(p)\in\mathfrak{O}$ for the plane of a
plaquette $p$, and $\mathcal{C}=\mathcal{C}^{\mathrm{bg}}$ for the background flux two-point function.
At separated arguments,
\begin{equation}\label{8.5:eq-no-entrywise-bound-1}
\mathcal{C}_{pq}=\varepsilon^4 D_{\pi(p)\pi(q)}\bigl(x(p)-x(q)\bigr)\,\bigl(1+o(1)\bigr).
\end{equation}
Here $o(1)$ refers to the limit along the floors of the two-point witness, on which the error terms
of Lemma~\ref{8.5:lem-covariance-upper} tend to zero. The tensor $D=(D_{PQ})_{P,Q\in\mathfrak{O}}$ is
the sign-indefinite continuum curvature tensor of the curvature-tensor lemma of the two-point
witness, which gives in particular
\begin{equation}\label{8.5:eq-no-entrywise-bound-2}
D_{12,34}\equiv 0,\qquad
D_{12,12}(x)\propto(x_3^2+x_4^2-x_1^2-x_2^2)\,|x|^{-6}.
\end{equation}
Consequently no lower bound on an individual entry $\mathcal{C}_{pq}$ can hold. We use the radial
majorant of Lemma~\ref{8.5:lem-covariance-upper},
\[
\sigma(r)=\varepsilon^4\pi^{-2}r^{-4},\qquad \sigma_{pq}=\sigma\bigl(|x(p)-x(q)|\bigr),\qquad
\sigma_{\mathsf{d}}=2\sigma(\mathsf{d}/4).
\]
The objects that are positive and radial are the majorant $\sigma$, whose square $\sigma_{pq}^2$ is
$\varepsilon^8$ times the plane sum \eqref{8.5:eq-no-entrywise-bound-3} at $x=x(p)-x(q)$, and the
majorant form $\mathcal{K}^+$ of the two-point witness theorem. A weight $f$ may depend on the plane,
$f=(f^P)_{P\in\mathfrak{O}}$; it is called scalar when $f^P$ does not depend on $P$. The form
$\mathcal{K}^+(f_1,f_2)$ is $6\pi^{-4}$ times the integral of
$|f_1|(x)\,|f_2|(y)\,|x-y|^{-8}$ over $x$ and $y$, with the envelope $|f|(x)=\max_P|f^P(x)|$. It
coincides with the continuum bilinear form $\mathcal{K}^{\mathrm{cont}}$ of the kernel $D$ exactly on
the cone of scalar non-negative weights, where both equal the form $\mathcal{K}_s$ of the glossary.
Every comparison in this section therefore bounds $\mathcal{C}$ from above
through $\sigma_{\mathsf{d}}$ and bounds positive quantities from below through $\mathcal{K}^+$.
\end{remark}

\begin{proof}
By \eqref{8.5:eq-no-entrywise-bound-2}, the entries of $D$ between the two orthogonal planes
$\{1,2\}$ and $\{3,4\}$ vanish identically. Within the plane $\{1,2\}$, the diagonal entry vanishes on
the cone $x_1^2+x_2^2=x_3^2+x_4^2$. Its sign on the two sides of that cone is the sign of
$x_3^2+x_4^2-x_1^2-x_2^2$, so it takes both signs. Through \eqref{8.5:eq-no-entrywise-bound-1}
the leading term of $\mathcal{C}_{pq}$ is zero whenever $\pi(p)$ and $\pi(q)$ are orthogonal, or
$x(p)-x(q)$ lies on such a cone. For this reason no inequality of the form
$|\mathcal{C}_{pq}|\ge c\,\sigma_{pq}$ with $c>0$ is available.

The same obstruction appears at the level of bilinear forms. By part (v) of the majorant-form lemma
of the two-point witness, on the non-negative cone the quotient
$\mathcal{K}^{\mathrm{cont}}/\mathcal{K}^+$ can be made arbitrarily small for non-scalar weights.
Hence no inequality $\mathcal{K}^{\mathrm{cont}}\ge\theta\,\mathcal{K}^+$ with a fixed $\theta>0$
holds uniformly over non-negative non-scalar weights.

The majorant rests on the per-plane identity established in the proof of part (i) of that lemma:
\begin{equation}\label{8.5:eq-no-entrywise-bound-3}
\sum_{Q\in\mathfrak{O}}D_{PQ}(x)^2=\pi^{-4}|x|^{-8}\qquad(P\in\mathfrak{O}).
\end{equation}
Each term of the sum is non-negative, so $|D_{PQ}(x)|\le\pi^{-2}|x|^{-4}$. This gives the entrywise
upper bound
\[
|\varepsilon^4D_{\pi(p)\pi(q)}(x(p)-x(q))|\le\sigma_{pq}.
\]
On the separated range this is turned into $|\mathcal{C}_{pq}|\le\sigma_{\mathsf{d}}$ by
Lemma~\ref{8.5:lem-covariance-upper}.

Since $\mathfrak{O}$ has six elements, \eqref{8.5:eq-no-entrywise-bound-3} also yields the lower
bound on the maximum over the six planes at a point:
\[
\max_{Q\in\mathfrak{O}}D_{PQ}(x)^2\ge\tfrac16\,\pi^{-4}|x|^{-8}.
\]
That maximum is therefore bounded below. The lower bound is used nowhere in this section, and for
this reason the cell-scale difference bound, Lemma~\ref{8.5:lem-covariance-difference}, is stated
with the radial majorant and not with the maximum over planes.

By part (iii) of the majorant-form lemma, $\mathcal{K}^+$ is positive. By parts (i) and (iv) of that
lemma, $\mathcal{K}^+=\mathcal{K}^{\mathrm{cont}}$ on scalar non-negative weights. The conversion to
the two-point witness's unit (Corollary~\ref{8.5:cor-room-witness-unit}, whose proof is incomplete
at one step) accordingly bounds $|\mathcal{C}_{pq}|$ from above by $\sigma_{\mathsf{d}}$ and bounds
$\mathcal{K}^+(f_1,f_2)$ from below by a multiple of $\|f_1\|_1\|f_2\|_1$, where $\|f\|_1$ is the
$L^1$ norm of the envelope $|f|$; every lower bound thus runs through $\mathcal{K}^+$.

For the same reason the two-point witness theorem states its remainder bound in two forms. The
first is relative to $\mathcal{K}^+$ for all weights of the Lipschitz shape class of its statement.
The second is relative to its main
term only on the cone of scalar non-negative weights, where the main term's kernel form coincides
with $\mathcal{K}^+$.
\end{proof}

\begin{definition}[Notation for the averaged extraction estimate]\label{8.5:not-extraction-notation}
The following letters are fixed throughout this section.

\emph{Standing data.} The correlation theorem is used as an input, with its quantifier prefix and
its order of choice of constants (Definitions~\ref{2.3:not-stages-first},
\ref{2.3:not-stages-middle} and~\ref{2.3:not-stages-last}). We write $r>0$ for its source radius,
which is also the source radius of the two-point witness theorem and is fixed by hypothesis (H7) of
Definition~\ref{2.2:not-hypotheses} before every ceiling of the order of choice. We take
$g\le\gamma_J$, $x_0:=g_0^{-2}>X:=g^{-2}$, and $K=K(g_0,g)$, the first passage of
Definition~\ref{1.2:def-flow-constants}; the letter $X$ for regions and polymers below is unrelated
to $X=g^{-2}$. The correlation theorem's parameter $m$ is set equal to $s$, so that its lower
condition reads $s\ge m_{\mathbb{T}}(g)$. The integer $N_T\ge1$, for which the unit lattice
$T^{(n)}$ has $N_TL^s$ sites per side, is free, and
\begin{equation}\label{8.5:eq-extraction-notation-1}
n:=k_0:=K'=K-s .
\end{equation}
Here $K'$ is the correlation theorem's level $K-m$, which for $m=s$ is the witness level $n$.
The inductive statements of the correlation theorem at all levels $k\le K$, together with its
standing conditions, are in force. Lengths are measured in cells of the level-$n$ unit lattice
$T^{(n)}$ unless another level is named. On $\mathbb{R}^4$, $|\cdot|$ is the sup norm, and for a set
$A$ and $a\ge0$ we put $A^{+a}:=\{x:\operatorname{dist}(x,A)\le a\}$.

\emph{The history $h_0$.} We write $h_0$ for the history in which the small-field domains
$\Omega_i=\Lambda_i$ are the whole torus for every index, the cutoff functions are
$\varphi_j\equiv1$, the region where groups are placed at level $j$ is
$\Lambda_j^{\circledR}=T^{(j)}$, and there is neither a large-field region $Z$ nor a boundary collar
of one. For $h_0$ the
summation of boundary terms at finer zones and the boundary terms of the correlation theorem do not
occur. By the sentence of \cite{B-CMP119} quoted in the standing conditions of the correlation
theorem, the terms $E^{(j)}(X,\cdot,y;w)$ and $E^{(j)}_{\hat{w}}(X,\cdot,y)$ of the sourced effective
action, in any history and for $X\subset\Lambda_j$, are the terms of the history $h_0$. They are
defined by the formula that produces them.

\emph{Levels.} In the averaged extraction estimate and the level sums that bound it we take $j\le n$
and put $t:=L^{j-n}$ and $\nu:=1+n-j$. In the bounds for a single step below the witness level,
where the current step is $k<n$, we put instead $t:=L^{j-k}$ and $\nu:=1+k-j$.

\emph{Functions of the coupling.} With the exponent $\mathfrak{p}_0$ of the degree function
$p_0(\cdot)$ and $T(g'):=A_1(\log g'^{-2})^{\mathfrak{p}_0}$ we put
\begin{align*}
\delta_k&:=g_kT(g_k)=g_kA_1(\log x_k)^{\mathfrak{p}_0}, &
\varepsilon_k&:=g_kp_0(g_k)=g_kA_0(\log x_k)^{\mathfrak{p}_0},\\
\alpha_{*,j}&=g_jC_*(\log x_j)^{q}, & \mathfrak{a}_j&:=C_{\mathfrak{a}}\alpha_{*,j},\\
\alpha_{3,j}&:=c_3\alpha_{*,j}, & \mathfrak{r}_j&:=c_T\alpha_{*,j}.
\end{align*}
Here $\alpha_{*,j}=\min\{\alpha_{0,j},\alpha_{1,j}\}$ is the minimum of the two radii
$\alpha_{i,j}=g_jC_i(\log x_j)^{q}$ of the tube of \cite[(2.28)]{B-CMP119}, taken with
$q_0=q_1=q$, so that $C_*=\min\{C_0,C_1\}$. The function $\theta$ of the correlation theorem is
\begin{equation}\label{8.5:eq-extraction-notation-2}
\theta_k:=\frac{K_R\delta_k}{\alpha_{*,k}}=C_\theta(\log x_k)^{\mathfrak{p}_0-q},\qquad
C_\theta:=\frac{K_RA_1}{C_*},
\end{equation}
and we put $\hat{\theta}_k:=\max_{i\le k}\theta_i$; the correlation theorem's ratio $\rho_{k,j}$ of
the tube radii at levels $k$ and $j$ equals $\mathfrak{a}_k/\mathfrak{a}_j$, because $q_0=q_1=q$.
With the enlargement parameter $\beta_s=\tfrac14$ of the correlation theorem (a fixed number; its
subscript is not the depth), the enlarged letters are
\begin{equation}\label{8.5:eq-extraction-notation-3}
\hat{\mathfrak{a}}_k:=(1+2\beta_s)\mathfrak{a}_k=\tfrac32\mathfrak{a}_k,\qquad
\hat{\rho}_{k,j}:=\tfrac32\rho_{k,j}.
\end{equation}

\emph{Level series.} Let $\hat{\alpha}$ be the H\"older exponent of the correlation theorem, the
exponent in its factor $t^{4+\hat{\alpha}}$; it is distinct from the enlarged radii
$\hat{\alpha}_k$. We put
\begin{align*}
S'_I&:=\sup_k\sum_{\nu\ge1}\nu^3\bigl(1+\hat{\rho}_{k,k+1-\nu}\bigr)L^{-\hat{\alpha}(\nu-1)},\\
\hat{S}_I&:=\sup_k\sum_{\nu\ge1}\nu^3\hat{\rho}^2_{k,k+1-\nu}
\bigl(1+\hat{\rho}_{k,k+1-\nu}\bigr)L^{-\hat{\alpha}(\nu-1)},\\
S'_\omega&:=\sum_{a\ge0}L^{-a/2},\qquad
C_{t\rho}:=\sup_{k,\nu}L^{-(\nu-1)}\rho_{k,k+1-\nu}.
\end{align*}

\emph{Counting constants.} Let $\pi_j$ be the partition into $M$-cubes at level $j$, $D_j$ the set
of polymers at level $j$, $d_j(X)$ the size of a polymer and $|X|_{\pi_j}$ its number of
$\pi_j$-cubes. Let $\delta>0$ be the constant of the correlation theorem that enters its floor
$\kappa\ge C_\delta(d)/\delta$; it is a fixed number, not to be confused with $\delta_k$.
For $a\in\{0,2\}$ we put
\begin{equation}\label{8.5:eq-extraction-notation-4}
\mathsf{S}_a:=\sup_{j,\;c\in\pi_j}\;\sum_{X\in D_j,\;X\ni c}
|X|_{\pi_j}\bigl(d_j(X)+5\bigr)^a e^{-\frac14\delta\kappa d_j(X)} .
\end{equation}

\emph{Constants of the correlation theorem.} The following constants are used as the correlation
theorem fixes them. All of them are chosen before $\gamma$ in its order of choice, and none of them
depends on $r$:
\begin{gather*}
E_0,\; C_I,\; C_M=2dC_{\mathfrak{a}}/c_3,\; C_B\le5,\; C_S\ge4,\; C_{\mathfrak{a}},\; c_3,\;
c_T,\; c'',\; K_R,\\
C_A,\; \nu_0,\; \alpha_4,\; b^{\mathrm{adm}},\;
K_{CE}=8b^{\mathrm{adm}}C_A/(\nu_0\alpha_4),\; C_1,\; C_6,\; \bar{\varepsilon},\;
\varepsilon_1,\; C_{\mathfrak{m}},\; \lambda,\; c_2,\; b_+,\; \kappa_0\ge7 .
\end{gather*}
Here $C_1$ is the constant with $|B'|\le C_1\delta_k$ on $\operatorname{supp}\chi_k$
\cite[(1.20)]{B-CMP116}; it is distinct from the radius constant $C_1$ of
\cite[(2.28)]{B-CMP119} above. The constant $C_6$ is the one in
$|G_y(B')|\le C_6L^4\sup_{B(y)}|B'|^2$ for the gauge-fixing term of the correlation theorem.
Finally, $\bar{\varepsilon}$ and $\varepsilon_1$ are constants of the analysis carried out in the
correlation theorem.
\end{definition}

\begin{proof}[Proof of the identities and finiteness statements above]
The enlarged letters satisfy
\begin{equation}\label{8.5:eq-extraction-notation-5}
\hat{\rho}^2_{k,j}=w_{k,j},\qquad \hat{\rho}^2_{k,j}\mathfrak{a}_j^2=\hat{\mathfrak{a}}_k^2 .
\end{equation}
The first identity holds because the weight of the correlation theorem is
$w_{k,j}=(1+2\beta_s)^2\rho^2_{k,j}$, and $(1+2\beta_s)^2=\tfrac94$ for $\beta_s=\tfrac14$. For the
second, \eqref{8.5:eq-extraction-notation-3} and the definition of $\rho_{k,j}$ give
\begin{equation*}
\hat{\rho}^2_{k,j}\mathfrak{a}_j^2=\tfrac94\,\frac{\mathfrak{a}_k^2}{\mathfrak{a}_j^2}\,
\mathfrak{a}_j^2=\bigl(\tfrac32\mathfrak{a}_k\bigr)^2=\hat{\mathfrak{a}}_k^2 .
\end{equation*}
The second member of \eqref{8.5:eq-extraction-notation-2} follows by inserting the definitions of
$\delta_k$ and $\alpha_{*,k}$, since the factors $g_k$ cancel:
\begin{equation*}
\frac{K_R g_kA_1(\log x_k)^{\mathfrak{p}_0}}{g_kC_*(\log x_k)^{q}}
=C_\theta(\log x_k)^{\mathfrak{p}_0-q}.
\end{equation*}

The series $S'_I$, $\hat{S}_I$ and $S'_\omega$ and the constant $C_{t\rho}$ do not depend on $K$, and
they are finite and bounded uniformly in $g\le\gamma_J$. This follows by the argument of the closing
sentence of the correlation theorem's table of per-cell bounds, which rests on
\begin{equation*}
\rho^2_{k,j}\le\bigl(1+o(1)\bigr)\bigl(1+(3\lambda(k-j)+2c_2)g^2\bigr),
\end{equation*}
where $\lambda$ is the correlation theorem's flow-rate constant.
For $j=k+1-\nu$ we have $k-j=\nu-1$. Hence $\rho_{k,k+1-\nu}$, and with it
$\hat{\rho}_{k,k+1-\nu}$, is bounded by the square root of a polynomial of degree one in $\nu$, with
coefficients independent of $k$ and $K$. Each summand of $S'_I$ and $\hat{S}_I$ is therefore bounded
by a polynomial in $\nu$ times the geometric weight $L^{-\hat{\alpha}(\nu-1)}$. The factor
$L^{-(\nu-1)}$ in $C_{t\rho}$ dominates the same square-root growth. The series $S'_\omega$ is
geometric with ratio $L^{-1/2}$.

The constants $\mathsf{S}_a$ of \eqref{8.5:eq-extraction-notation-4} depend only on
$(d,L,M,\kappa,\delta)$. They are finite by \cite[(1.26)]{B-CMP116}, applied at the decay rate
$\tfrac14\delta\kappa$, which the floor $\kappa\ge C_\delta(d)/\delta$ of the correlation theorem
permits. Their finiteness is that of the constant $C(\kappa)$ in the correlation theorem's lemma on
boundary terms.
Finally, every $\pi_j$-cube contains $M^4$ points of $T^{(j)}$, so that
$\#\bigl(X\cap T^{(j)}\bigr)=M^4|X|_{\pi_j}$.
\end{proof}

\begin{lemma}[Decomposition of the first-order irrelevant response]\label{8.5:lem-response-decomposition}
Let $n:=K-s$ be the witness level, and assume:
\begin{itemize}
\item the representation of the effective action with complex sources in the correlation theorem, parts (a)--(f), together with the extraction functional $\beta[\cdot]$ of the correlation theorem and that theorem's definition $b_j(\zeta):=\beta[E^{(j)}_\zeta]$ of the coupling increments;
\item parts (b) and (c) of the inductive statement of the correlation theorem;
\item the representation of the response of a branch in the two-point witness theorem.
\end{itemize}
Let $f$ be Lipschitz and $z\in\mathbb{C}$ with $|z|\,\|f\|_\sharp<r$, where $r$ is the source radius of the classes $\mathcal{W}_j$ of the correlation theorem. By the class definitions of the correlation theorem, $zf\in\mathcal{W}_{K'}\subset\mathcal{W}_j$ for every $j\le K'=n$. Put
\begin{align*}
\dot{S}^{(j)}(X,\cdot,y)&:=\partial_\zeta E^{(j)}_\zeta(X,\cdot,y)\big|_{\zeta=0},\qquad
\dot{D}^{(j)}(X,\cdot,y)[f]:=\partial_z D^{(j)}(X,\cdot,y;zf)\big|_{z=0},\\
\dot{R}^{(j)}(X,\cdot)[f]&:=\partial_z R^{(j)}(X,\cdot;zf)\big|_{z=0},\qquad
F^{\mathrm{vac}}:=F-F(\mathbf{1}).
\end{align*}
Then $\beta[\dot{S}^{(j)}]=b_j'(0)$. Put
\begin{equation*}
\mathfrak{i}_y[\dot{S}^{(j)}]:=\sum_{X\ni y}\bigl(\dot{S}^{(j)}(X,\cdot,y)\bigr)^{\mathrm{vac}}-b_j'(0)\,A(h^{(j)}_y,\cdot),
\end{equation*}
where $h^{(j)}_y$ is the weight of part (a) of the representation and $A(h^{(j)}_y,\cdot)$ the corresponding weighted Wilson action. Then, on the whole-torus history $h_0$ of the two-point witness theorem, for $w=zf$ and every fine configuration $\mathfrak{b}$ at which the terms are evaluated,
\begin{equation}\label{8.5:eq-response-decomposition-a}
\begin{split}
&\partial_z\bigl[\mathbf{E}^w_n+\mathbf{D}^w_n+\mathbf{R}^w_n-E^w_n\bigr]\big|_{z=0}(\mathfrak{b})
-\partial_z\bigl[\mathbf{E}^w_n+\mathbf{D}^w_n+\mathbf{R}^w_n-E^w_n\bigr]\big|_{z=0}(\mathbf{1})\\
&\quad=\sum_{j\le n}\Bigl\{\sum_{y\in T^{(j)}}f(p_y)\,\mathfrak{i}_y[\dot{S}^{(j)}](\mathfrak{b})
+\sum_{X\in D_j}\ \sum_{y\in X\cap T^{(j)}}\bigl(\dot{D}^{(j)}(X,\cdot,y)[f]\bigr)^{\mathrm{vac}}(\mathfrak{b})\\
&\qquad\qquad+\sum_{X\in D_j}\bigl(\dot{R}^{(j)}(X,\cdot)[f]\bigr)^{\mathrm{vac}}(\mathfrak{b})\Bigr\}.
\end{split}
\end{equation}
Here $p_y$ is the plaquette at which part (b) of the representation evaluates the source for the point $y$.
\end{lemma}

\begin{proof}
We first note that the source values entering part (b) lie in the disc on which the inductive statement gives analyticity. Since $\|f\|_\sharp\ge\|f\|_\infty$, every plaquette $p$ satisfies
\begin{equation*}
|zf(p)|\le|z|\,\|f\|_\infty\le|z|\,\|f\|_\sharp<r<2r.
\end{equation*}

\emph{The sum $\mathbf{E}^w_n$.} By part (b) of the representation, since on the whole-torus history $h_0$ the region $\Lambda_j^{\circledR}$ in which that part places the groups at level $j$ is all of $T^{(j)}$,
\begin{equation*}
\mathbf{E}^w_n=\sum_{j\le n}\sum_{y\in T^{(j)}}\mathfrak{i}^{(j)}_y\Big|_{\hat{w}=zf(p_y)},
\end{equation*}
where
\begin{equation*}
\mathfrak{i}^{(j)}_y=\sum_{X\ni y}\bigl[E^{(j)}_{\hat{w}}(X,\cdot,y)-E^{(j)}_{\hat{w}}(X,\mathbf{1},y)\bigr]
-b_j(\hat{w})\,A(h^{(j)}_y,\cdot).
\end{equation*}
By parts (b) and (c) of the inductive statement, the maps $\hat{w}\mapsto E^{(j)}_{\hat{w}}$ and $\hat{w}\mapsto b_j(\hat{w})$ are analytic on $D(0,2r)$. The inner map $z\mapsto zf(p_y)$ is linear with derivative $f(p_y)$ and, by the first paragraph, takes values in $D(0,2r)$. The chain rule therefore gives
\begin{align*}
\partial_z E^{(j)}_{zf(p_y)}(X,\cdot,y)\big|_{z=0}&=f(p_y)\,\dot{S}^{(j)}(X,\cdot,y),\\
\partial_z b_j(zf(p_y))\big|_{z=0}&=f(p_y)\,b_j'(0).
\end{align*}
Hence the $z$-derivative at $0$ of $\mathfrak{i}^{(j)}_y$ at $\hat{w}=zf(p_y)$ equals
\begin{equation*}
f(p_y)\Bigl[\sum_{X\ni y}\bigl(\dot{S}^{(j)}(X,\cdot,y)\bigr)^{\mathrm{vac}}-b_j'(0)\,A(h^{(j)}_y,\cdot)\Bigr]
=f(p_y)\,\mathfrak{i}_y[\dot{S}^{(j)}].
\end{equation*}
This gives the first sum on the right of \eqref{8.5:eq-response-decomposition-a}.

\emph{The sums $\mathbf{D}^w_n$ and $\mathbf{R}^w_n$.} By parts (c) and (d) of the representation, $\mathbf{D}^w_n$ and $\mathbf{R}^w_n$ are sums over $j\le n$ and $X\in D_j$ (and, for $\mathbf{D}^w_n$, over $y\in X\cap T^{(j)}$) of terms that are already vacuum-subtracted, that is, of the form $F-F(\mathbf{1})$. Their local terms are analytic in the source on $\mathcal{W}_j(X)$, and $zf$ lies in that class. The sums are finite, since they run over polymers of a torus. We may therefore differentiate term by term. By the definitions of $\dot{D}^{(j)}$ and $\dot{R}^{(j)}$, this yields the second and third sums on the right of \eqref{8.5:eq-response-decomposition-a}.

\emph{The field-free term and the subtraction at $\mathbf{1}$.} By part (f) of the representation, $E^w_n$ does not depend on the field. Its $z$-derivative is therefore the same at $\mathfrak{b}$ and at $\mathbf{1}$, and it cancels against the corresponding term of the expression evaluated at $\mathbf{1}$.

Every term of the right member vanishes at $\mathbf{1}$. Each is either of the form $F^{\mathrm{vac}}$, which vanishes at $\mathbf{1}$ by definition, or a multiple of the weighted Wilson action $A(h^{(j)}_y,\cdot)$. The latter vanishes at $\mathbf{1}$ because each plaquette term $1-\operatorname{Re}\operatorname{tr}U(\partial p)$ does, since $\operatorname{tr}\mathbf{1}=1$ (Definition~\ref{1.1:def-wilson-action}). Thus the derivative of $\mathbf{E}^w_n+\mathbf{D}^w_n+\mathbf{R}^w_n$ evaluated at $\mathbf{1}$ is zero. Subtracting it leaves exactly the right member evaluated at $\mathfrak{b}$, which proves \eqref{8.5:eq-response-decomposition-a}.

\emph{The family $\dot{S}^{(j)}$.} The family $\dot{S}^{(j)}$ is a pointed $E$-family at scale $j$ on the radii $\alpha_{\cdot,j}$. The three defining conditions of a pointed $E$-family are linear, and they hold for every $E^{(j)}_\zeta$ by the correlation theorem's proposition on the frozen families. They therefore hold for the difference quotients $\zeta^{-1}(E^{(j)}_\zeta-E^{(j)}_0)$. Since these quotients converge to $\dot{S}^{(j)}$ in the norm $\|\cdot\|^{\flat}$, that is, uniformly on the tubes $\mathcal{U}^{\flat}_j(X)$, the conditions hold for the derivative $\dot{S}^{(j)}$ as well.

\emph{The identity $\beta[\dot{S}^{(j)}]=b_j'(0)$.} By the correlation theorem's definition of the coupling increments, $b_j(\zeta)=\beta[E^{(j)}_\zeta]$. The functional $\beta$ is linear and bounded on the space of families normed by $\|\cdot\|^{\flat}$, by the lemma on the boundedness of the extraction functional in the $\flat$-norm. Applying $\beta$ to the difference quotients of $\zeta\mapsto E^{(j)}_\zeta$ at $0$ gives the difference quotients of $b_j$ at $0$. Passing to the limit, $\beta$ commutes with $\partial_\zeta$, and $b_j'(0)=\beta[\dot{S}^{(j)}]$. This justifies writing $b_j'(0)$ in place of $\beta[\dot{S}^{(j)}]$ in $\mathfrak{i}_y[\dot{S}^{(j)}]$.

Finally, in the representation of the response of a branch in the two-point witness theorem, the Wilson member
\begin{equation*}
\sum_p\partial_z\zeta_n(p)\big|_{z=0}\,a_p
\end{equation*}
is the main curvature insertion $N\mathcal{O}(\tilde{f})$. It is not part of the left member of \eqref{8.5:eq-response-decomposition-a}, which, for the single source $w=zf$, is the irrelevant part $\operatorname{Irr}$ of that response.
\end{proof}

\begin{lemma}[The unit-sum bound at one renormalisation step]\label{8.5:lem-unit-sum}
Notation is that of Notation~\ref{8.5:not-extraction-notation}; we assume the following.
\begin{itemize}
\item[(H)] The correlation theorem, and in particular these parts of it:
\begin{itemize}
\item its activity lemma;
\item its core estimate for renormalised groups, with the note that follows it;
\item its single-term estimate, parts (i) and (ii), with the remark that follows it;
\item part (b) of its Wilson-piece bound;
\item its tube lemma;
\item the first step of the proof of its complex-background lemma;
\item the holomorphy statement for its response map $\mathsf{R}_p$;
\item the properties (c1)--(c3) of its admissible backgrounds;
\item its table of units;
\item its comparison inequalities for polymer sizes, for the ratios $\rho_{k,j}$ and for $\kappa$;
\item the count of units in its bound on the increments of the effective action.
\end{itemize}
\item[(H$'$)] The cluster-expansion lemma, with its notions of a unit $v=(j_v,S_v,f_v)$, of the
domain condition, of the oscillation $\omega(v)$ and of the marked activity sum $W^m$.
\end{itemize}
Fix a step $k<n$ along the history $h_0$, and assume
\begin{equation*}
\theta_k\le \tfrac13 .
\end{equation*}
Let $\mathfrak{M}$ be a finite family of units of the following three kinds.
\begin{itemize}
\item[(a)] \emph{Groups.} For each pair $(j,y)$ with $j\le k$ and $y\in T^{(j)}$, $\mathfrak{M}$
contains at most one group $\mathfrak{i}_y[F]$. Here $F=F^{(j,y)}$ is a pointed $E$-family at
scale $j$ on the radii $\alpha_{\cdot,j}$ with $\|F\|^{\flat}\le\mathfrak{n}^G_j$. The group is
presented in the way the first row of the table of units presents a frozen group at the site of
the operation $\mathbf{T}$. The
\emph{core} consists of the terms with polymers $X\ni y$ lying inside $2^d$ cubes of $\pi_k$;
the core, taken together with the counterterm piece $\beta[F]A(h_y,\cdot)$, is one unit.
Every other constituent of the group, called a \emph{tail}, is a unit of
its own.
\item[(b)] \emph{Singles.} For each triple $(j,X,y)$ (or pair $(j,X)$) with $j\le k$,
$\mathfrak{M}$ contains at most one single $H(X,\cdot)-H(X,1)$. Here $H(X,\cdot)$ has the
properties (f1) and (f2) of the single-term estimate on $\mathcal{U}^{\flat}_j(X)$, and
\begin{equation*}
|H(X,\cdot)|\le\mathfrak{n}^M_j\,e^{-(1-\frac14\delta)\kappa d_j(X)} .
\end{equation*}
\item[(c)] \emph{Wilson pieces.} For each shift cube $\square$, $\mathfrak{M}$ contains at most
one Wilson piece $A(c\hat{\chi}_\square,\cdot)$, where $\sup|c|\le\bar{c}$.
\end{itemize}
Then the domain condition of the cluster-expansion lemma holds for every unit of $\mathfrak{M}$
on the whole output tube, whose radii are $\hat{\alpha}_k$. Moreover
\begin{equation}\label{8.5:eq-unit-sum-a}
\begin{split}
W^m(\mathfrak{M})\le 8\theta_k\Bigl\{C_{\boxplus}\Bigl[&\sum_{j\le k}\nu^3\hat{\rho}_{k,j}^2
(1+\hat{\rho}_{k,j})L^{-\hat{\alpha}(\nu-1)}\mathfrak{n}^G_j\\
&+\sum_{j\le k}\hat{\rho}_{k,j}^2\mathfrak{n}^M_j\Bigr]+C_5\,\bar{c}\,\hat{\mathfrak{a}}_k^2\Bigr\},
\end{split}
\end{equation}
where $\nu:=1+k-j$ in the summand of index $j$, where $\hat{\alpha}\le 1$ is the exponent of
the core estimate (not to be confused with the radii $\hat{\alpha}_k$ of the output tube), where
$C_I$, $C_M$ and $C_S$ are, respectively, the constant of the core estimate, the constant of the
single-term estimate and the constant of the Wilson-piece row of the table of units, and where
\begin{equation}\label{8.5:eq-unit-sum-1}
\begin{aligned}
C_{\boxplus}&:=\max\bigl\{(5M)^4e^{2^d\kappa}C_I+2\cdot 3^dM^4\mathsf{S}_0,\;
3^dM^4(32C_M^2M^2+2)\mathsf{S}_2\bigr\},\\
C_5&:=2^d\,e^{(2^d-1)\kappa}\,C_S\,(2M)^4 .
\end{aligned}
\end{equation}
\end{lemma}

\begin{proof}
For a level $j\le k$ we put $t:=L^{j-k}$ and $\nu:=1+k-j$. Thus $t=L^{-(\nu-1)}$ and
$t^{\hat{\alpha}}=L^{-\hat{\alpha}(\nu-1)}$. We abbreviate $\hat{\rho}:=\hat{\rho}_{k,j}$.

\emph{The domain condition and the oscillation.} On the support of $\chi_k$ the fluctuation
field $B$ is real and satisfies $|B|<T_k$. Hence $K_Rg_k\sup|B|<\theta_k\alpha_{*,k}$. Put
\begin{equation*}
\mathsf{w}_k:=\frac{1}{\theta_k}-1 .
\end{equation*}
Since $\theta_k\le\frac13$, we have $\mathsf{w}_k\ge 2$. The response map
$\mathsf{R}_p(\mathsf{t},\sigma)$ of the holomorphy statement is indexed by a pair
$p=(\square,Y_0)$ of a shift cube $\square$ and a region $Y_0$, carries a factor $\eta_p$
attached to $p$, and depends on a complex parameter $\mathsf{t}$, on a function $\sigma$ of the
unit cubes $\Delta$ and on a fluctuation datum $\mathsf{B}$. For complex $\mathsf{t}$ with
$\eta_p|\mathsf{t}|\le\mathsf{w}_k$,
\begin{equation*}
(1+\eta_p|\mathsf{t}|)\,K_Rg_k\sup|B|<(1+\mathsf{w}_k)\,\theta_k\alpha_{*,k}=\alpha_{*,k} .
\end{equation*}
So the condition of the holomorphy statement for the response map holds with
$\mathsf{B}=B$.

That statement therefore applies. For $p=(\square,Y_0)$, for all $(\mathbf{U},\mathbf{J})$ in the
output tube over $Y_0$, and for all $\sigma$ with $|\sigma(\Delta)|\le e^{\kappa_1}$, the response
map $\mathsf{R}_p(\mathsf{t},\sigma)$, restricted to $S_v$, is an admissible background. It is
jointly holomorphic in its arguments. The domain condition of the cluster-expansion lemma
concerns the values $|w|\le 2$ of its complex parameter $w$, which multiplies the shift.
Measured in units of the shift, $w$ corresponds to $\eta_p\mathsf{t}$, so that the disc
$\eta_p|\mathsf{t}|\le\mathsf{w}_k$ is the disc $|w|\le\mathsf{w}_k$; because $\mathsf{w}_k\ge2$,
it contains the sub-disc $|w|\le 2$. The domain condition
therefore holds for every unit on the whole output tube, which is the first assertion.

For a unit $v$, let $\mathfrak{N}(v)$ be a number such that
$|f_v\circ\mathsf{R}_p-f_v(1)|\le\mathfrak{N}(v)$ on the disc
$\eta_p|\mathsf{t}|\le\mathsf{w}_k$. Measured in units of the shift, this disc has radius
$\mathsf{w}_k$. The difference between the values of $f_v\circ\mathsf{R}_p$ at the shift $w$ and
at the shift $0$ is a holomorphic function of $w$ on this disc; it vanishes at $w=0$ and is
bounded there by $2\mathfrak{N}(v)$. By the Schwarz lemma it is bounded by
$2\mathfrak{N}(v)\,|w|/\mathsf{w}_k$, hence by $2\mathfrak{N}(v)\cdot2/\mathsf{w}_k$ on the
sub-disc $|w|\le2$. Therefore
\begin{equation*}
\omega(v)\le 2\mathfrak{N}(v)\cdot\frac{2}{\mathsf{w}_k}=\frac{4\mathfrak{N}(v)}{\mathsf{w}_k} .
\end{equation*}
This is the
oscillation bound stated in the activity lemma. Since $1-\theta_k\ge\frac23$,
\begin{equation*}
\frac{4}{\mathsf{w}_k}=\frac{4\theta_k}{1-\theta_k}\le 6\theta_k\le 8\theta_k .
\end{equation*}
So each unit $v$ with $S_v\cap\Delta\neq\emptyset$ contributes at most
$8\theta_k\,\mathfrak{N}(v)\,e^{(1-\frac12\delta)\kappa d_k(\bar{S}_v)}$ to the sum whose
supremum over the unit cubes $\Delta$ is $W^m$.

\emph{Sizes of the units on the admissible backgrounds of this step.} The output tube has the
radii $\hat{\alpha}_k$. The backgrounds it generates satisfy (c1)--(c3) at the amplitude
$C_{\mathfrak{a}}\hat{\alpha}_k=\hat{\mathfrak{a}}_k$, because (c1)--(c3) involve the radii only
through ratios or ceilings, as does the tube lemma. The remark following the single-term
estimate reads the ratio $(\hat{\alpha}_k/\alpha_j)^2=w_{k,j}$ on exactly this tube, and the
count of units in the bound on the increments uses it for the units bounded by the single-term
estimate.

The core estimate, the single-term estimate and the Wilson-piece bound read a background only
through (c1)--(c3), as the first step of the proof of the complex-background lemma shows. Let
$S$ be the radius of analyticity in the proof of the single-term estimate. Then
\begin{equation*}
S^{-1}\le\frac{2d\,D_X\,\hat{\mathfrak{a}}_k\,t^2}{c_3\alpha_{*,j}}=C_MD_Xt^2\hat{\rho} ,
\end{equation*}
where $d=4$ is the dimension and $C_M=2dC_{\mathfrak{a}}/c_3$.
Hence the following values of $\mathfrak{N}$ may be taken.
\begin{itemize}
\item[(1)] For a core with its counterterm piece, the core estimate gives
\begin{equation*}
\mathfrak{N}=C_I\,\mathfrak{n}^G_j\,\nu^3\hat{\rho}^2(1+\hat{\rho})\,t^{4+\hat{\alpha}} .
\end{equation*}
The note following the core estimate covers the absence of the tails. The two cores differ
by constituents of the tail class: one consists of the $X\ni y$ inside $2^d$ cubes of $\pi_k$,
the other of the $X\subset\square^{\sim\nu}$.
\item[(2)] For a tail with polymer $X$,
\begin{equation*}
\mathfrak{N}=2\mathfrak{n}^G_j\,e^{-\kappa d_j(X)} .
\end{equation*}
\item[(3)] For a single with $X\subset\tilde{\square}^2$, part (i) of the single-term estimate
gives
\begin{equation*}
\mathfrak{N}=2\mathfrak{n}^M_j\,e^{-(1-\frac14\delta)\kappa d_j(X)}(C_MD_Xt^2\hat{\rho})^2 .
\end{equation*}
\item[(4)] For any other single, $d_j(X)\ge t^{-1}$, and part (ii) gives
\begin{equation*}
\mathfrak{N}=2\mathfrak{n}^M_j\,e^{-(1-\frac14\delta)\kappa d_j(X)} .
\end{equation*}
\item[(5)] For a Wilson piece, part (b) of the Wilson-piece bound gives the bound
$4\sup|c|\,\hat{\mathfrak{a}}_k^2$ on each unit cell. The cube $\square$ consists of $(2M)^4$
cells, and $A(c,1)=0$. Hence
\begin{equation*}
\mathfrak{N}=C_S(2M)^4\,\bar{c}\,\hat{\mathfrak{a}}_k^2 .
\end{equation*}
\end{itemize}

\emph{The sum over the units meeting one unit cube.} Fix a unit cube $\Delta$ of the lattice
of level $k$. We bound the sum of $\mathfrak{N}(v)\,e^{(1-\frac12\delta)\kappa d_k(\bar{S}_v)}$
over the units $v\in\mathfrak{M}$ with $S_v\cap\Delta\neq\emptyset$. Throughout we use
$d_k(\bar{X})\le t\,d_j(X)$, from the comparison inequalities.

\emph{Cores.} For the unit formed by a core with its counterterm, $\bar{S}_v$ consists of at
most $2^d$ cubes of $\pi_k$. Hence $d_k(\bar{S}_v)\le 2^d-1$, and the weight is at most
$e^{(2^d-1)\kappa}\le e^{2^d\kappa}$. The condition $S_v\cap\Delta\neq\emptyset$ places $y$ in
the cube of side $5M$ about $\Delta$, which contains at most $(5M)^4t^{-4}$ points of
$T^{(j)}$. Since $t^{-4}\cdot t^{4+\hat{\alpha}}=L^{-\hat{\alpha}(\nu-1)}$, and there is at most
one group for each $(j,y)$, the cores of level $j$ contribute at most
\begin{equation*}
(5M)^4e^{2^d\kappa}C_I\,\nu^3\hat{\rho}^2(1+\hat{\rho})L^{-\hat{\alpha}(\nu-1)}\mathfrak{n}^G_j .
\end{equation*}

\emph{Tails.} First let $\nu\ge 2$, so that $t\le\frac13$. Because $d_j(X)\ge\nu-1$ and
$\kappa\ge 10\log L$ (comparison inequalities), we have
\begin{equation*}
e^{-\kappa d_j(X)/2}\le L^{-5d_j(X)}\le L^{-5(\nu-1)}=t^5 .
\end{equation*}
Therefore
\begin{equation*}
\begin{split}
e^{-\kappa d_j(X)}\,e^{(1-\frac12\delta)\kappa t\,d_j(X)}
&\le e^{-\frac23\kappa d_j(X)}\\
&\le t^5e^{-\kappa d_j(X)/6} .
\end{split}
\end{equation*}
For $\nu=1$ instead,
\begin{equation*}
e^{-\kappa d_j(X)}\,e^{(1-\frac12\delta)\kappa d_j(X)}=e^{-\frac12\delta\kappa d_j(X)} .
\end{equation*}

We anchor each pair $(X,y)$ with $X\cap\Delta\neq\emptyset$ at a cube $c\in\pi_j$ with
$c\subset X$ and $c\cap\Delta\neq\emptyset$. A cube of $\pi_j$ has side $Mt$ in units of level
$k$. It meets the unit cube $\Delta$ in at most
\begin{equation*}
\Bigl(\frac{1}{Mt}+2\Bigr)^4\le 3^d\max\{(Mt)^{-4},1\}
\end{equation*}
positions. For each position there are $M^4|X|_{\pi_j}$ points $y$, and
$\max\{(Mt)^{-4},1\}M^4\le M^4t^{-4}$. Since both remaining exponentials are at most
$e^{-\frac14\delta\kappa d_j(X)}$, the sum of $|X|_{\pi_j}$ times the remaining exponential over
$X\ni c$ is at most $\mathsf{S}_0$. The tails of level $j$ therefore contribute at
most
\begin{equation*}
3^dM^4t^{-4}\,\mathsf{S}_0\cdot 2\mathfrak{n}^G_j\cdot t^5 ,
\end{equation*}
with the factor $t^5$ replaced by $1$ when $\nu=1$. Moreover,
\begin{equation*}
t\le\nu^3\hat{\rho}^2(1+\hat{\rho})L^{-\hat{\alpha}(\nu-1)} ,
\end{equation*}
because $\hat{\alpha}\le 1$ and, by the comparison inequalities,
$\hat{\rho}^2\ge\frac94(1-o(1))\ge 1$. Together with the cores, the groups of level $j$
contribute at most
\begin{equation*}
\bigl((5M)^4e^{2^d\kappa}C_I+2\cdot 3^dM^4\mathsf{S}_0\bigr)\,
\nu^3\hat{\rho}^2(1+\hat{\rho})L^{-\hat{\alpha}(\nu-1)}\mathfrak{n}^G_j .
\end{equation*}

\emph{Singles with $X\subset\tilde{\square}^2$.} Here $D_X\le 4M(d_j(X)+5)$. Since $t\le1$,
the exponent left after the weight is
\begin{equation*}
-\kappa d_j(X)\Bigl[\bigl(1-\tfrac14\delta\bigr)-\bigl(1-\tfrac12\delta\bigr)t\Bigr]
\le-\tfrac14\delta\kappa d_j(X)
\end{equation*}
for every $\nu\ge1$. The same anchoring gives the bound
\begin{equation*}
3^dM^4t^{-4}\cdot 2\mathfrak{n}^M_j\,(4C_MM)^2\,t^4\,\hat{\rho}^2\,\mathsf{S}_2
=3^dM^4\cdot 32C_M^2M^2\,\mathsf{S}_2\,\hat{\rho}^2\mathfrak{n}^M_j .
\end{equation*}

\emph{Singles with $X\not\subset\tilde{\square}^2$.} Then $\nu\ge2$, so $t\le\frac13$ and the
weight is at most $e^{\kappa d_j(X)/3}$. By the comparison inequalities,
$e^{-\frac{2}{15}\kappa d_j(X)}\le t^4$. Moreover $(1-\frac14\delta)-\frac13-\frac{2}{15}\ge\frac14$,
so a factor $e^{-\frac14\kappa d_j(X)}$ remains. The same anchoring gives
\begin{equation*}
3^dM^4t^{-4}\cdot 2\mathfrak{n}^M_j\,t^4\,\mathsf{S}_0
\le 2\cdot 3^dM^4\,\mathsf{S}_2\,\hat{\rho}^2\mathfrak{n}^M_j .
\end{equation*}
Here we used $\mathsf{S}_0\le\mathsf{S}_2$ and $\hat{\rho}^2\ge1$. Hence the singles of level
$j$ contribute at most
\begin{equation*}
3^dM^4(32C_M^2M^2+2)\,\mathsf{S}_2\,\hat{\rho}^2\mathfrak{n}^M_j .
\end{equation*}

\emph{Wilson pieces.} The cube $\Delta$ meets at most $2^d$ shift cubes, and
$d_k(\square)\le 2^d-1$. The Wilson pieces therefore contribute at most
\begin{equation*}
2^de^{(2^d-1)\kappa}C_S(2M)^4\,\bar{c}\,\hat{\mathfrak{a}}_k^2=C_5\,\bar{c}\,\hat{\mathfrak{a}}_k^2 .
\end{equation*}

\emph{Conclusion.} By the definition \eqref{8.5:eq-unit-sum-1} of $C_{\boxplus}$, the
contribution of the groups of level $j$ is at most
$C_{\boxplus}\nu^3\hat{\rho}^2(1+\hat{\rho})L^{-\hat{\alpha}(\nu-1)}\mathfrak{n}^G_j$, and that of
the singles of level $j$ is at most $C_{\boxplus}\hat{\rho}^2\mathfrak{n}^M_j$. We sum over
$j\le k$ and add the Wilson pieces. Multiplying by the factor $8\theta_k$ of the first part of
the proof and taking the supremum over $\Delta$ gives \eqref{8.5:eq-unit-sum-a}.

The count in the last part of this proof is the count by which the correlation theorem fixes
the constant of its bound on the increments of the effective action, carried out level by
level, with $(5M)^4$ in place of $(3M)^4$ and with the factor $3^d$ of the anchoring count. It
uses no property of the radii and none of the individual terms.
\end{proof}

\begin{lemma}[The marked-sum bound with directly given potentials]\label{8.5:lem-direct-potentials}
Fix a level $k$, and assume the following three statements:
\begin{itemize}
\item[(a)] the cluster-expansion lemma for marked units at one renormalisation step, together
with its corollary bounding the marked sub-sum, with the constants $C_A$, $\nu_0$, $\alpha_4$,
$b^{\mathrm{adm}}$, the number $\delta$, and
$K_{CE}:=8b^{\mathrm{adm}}C_A/(\nu_0\alpha_4)$ fixed there; the proof of that lemma yields
analyticity in $\lambda$ and the uniform majorant of the output kernels, and it uses the marked
potential only through its localisation in $Y$, its holomorphy, its invariance under the
transformations of \cite[(3.29)]{B-CMP109}, and the bound
\begin{equation}\label{8.5:eq-direct-potentials-1}
|V^m_k(Y)|\le C_AW^me^{-(1-2\delta)\kappa d_k(Y)} ;
\end{equation}
\item[(b)] the lemma on the locality and the activity of the marked potential in the proof of
the correlation theorem;
\item[(c)] the third step of the proof of the complex-source lemma for the correlation
theorem, which files the gauge-fixing pieces $\zeta_k(p_y)G_y(B')$ among the potentials, $p_y$
being the plaquette at which the running shift $\zeta_k$ enters the gauge-fixing coefficient
$x_k-\zeta_k(p_y)$ at the site $y$, with the bound, in modulus, on one block
\begin{equation}\label{8.5:eq-direct-potentials-2}
|\zeta_k(p_y)G_y(B')|\le C L^4 r\,(C_1g_kT^{\mathbb{C}})^2 ,
\end{equation}
in which the constant $C$, the number $r$ and the complex-tube parameter $T^{\mathbb{C}}$ are
those fixed in that step.
\end{itemize}
Let $\mathfrak{M}$ be a finite family of marked units satisfying the domain condition of the
cluster-expansion lemma, with marked activity sum $W^m$. Let $V^m_k$ be the marked potential
generated by $\mathfrak{M}$, so that \eqref{8.5:eq-direct-potentials-1} holds for every
polymer $Y$.
Let $V^{\mathrm{dir}}=\{V^{\mathrm{dir}}(Y)\}_Y$ be a marked potential given directly: each
$V^{\mathrm{dir}}(Y)$ is localised in $Y$, holomorphic in the fields, invariant under the
transformations of \cite[(3.29)]{B-CMP109}, and
\begin{equation}\label{8.5:eq-direct-potentials-3}
|V^{\mathrm{dir}}(Y)|\le \tilde{a}^{\mathrm{dir}}\,e^{-(1-2\delta)\kappa d_k(Y)} .
\end{equation}
Put $\tilde{a}:=C_AW^m+\tilde{a}^{\mathrm{dir}}$, assume $\tilde{a}\le \nu_0\alpha_4/8$, and put
$\lambda_0:=\nu_0\alpha_4/(4\tilde{a})$.
For complex $\lambda$ with $|\lambda|<\lambda_0$ let $t_\lambda(X)$ be the output kernel at
level $k+1$ with domain
$X$ of the cluster expansion in which the marked potential $V^m_k+V^{\mathrm{dir}}$ is
multiplied by $\lambda$ and the unmarked potential is left unchanged. Then $\lambda_0\ge2$,
$t_\lambda(X)$ is analytic on $|\lambda|<\lambda_0$ with
$|t_\lambda(X)|\le b^{\mathrm{adm}}e^{-\kappa d_{k+1}(X)}$ there, and
\begin{equation}\label{8.5:eq-direct-potentials-a}
|t_1(X)-t_0(X)|\le \frac{2b^{\mathrm{adm}}}{\lambda_0}\,e^{-\kappa d_{k+1}(X)}
=\frac{K_{CE}}{C_A}\,\tilde{a}\,e^{-\kappa d_{k+1}(X)} .
\end{equation}
Moreover, under (c), the gauge-fixing pieces $c(y)G_y(B')$ with $|c(y)|\le\bar{c}$ are marked
potentials given directly, and for them
\begin{equation}\label{8.5:eq-direct-potentials-4}
\tilde{a}^{\mathrm{dir}}\le C_6M^4\bar{c}\,C_1^2\delta_k^2 ,
\end{equation}
where $C_6$ is a number such that, on one block,
$|c(y)G_y(B')|\le C_6L^4\bar{c}\,C_1^2\delta_k^2$ whenever $|c(y)|\le\bar{c}$.
\end{lemma}

\begin{proof}
By (a), the part of the proof of the cluster-expansion lemma that yields analyticity in
$\lambda$ and the uniform majorant of the output kernels uses the marked potential only through
four of its properties: localisation in $Y$, holomorphy, invariance under the transformations of
\cite[(3.29)]{B-CMP109}, and the bound
\eqref{8.5:eq-direct-potentials-1}. The lemma in (b) is formulated in
exactly these terms, the activity of the marked potential being its only quantitative input.

The sum $V^m_k+V^{\mathrm{dir}}$ has the first three properties, since each is preserved
under addition of two potentials having it. By the triangle inequality, the bound on
$V^m_k$ and \eqref{8.5:eq-direct-potentials-3},
\begin{equation*}
|V^m_k(Y)+V^{\mathrm{dir}}(Y)|\le (C_AW^m+\tilde{a}^{\mathrm{dir}})
e^{-(1-2\delta)\kappa d_k(Y)}=\tilde{a}\,e^{-(1-2\delta)\kappa d_k(Y)} .
\end{equation*}
Hence the conclusions of the cluster-expansion lemma hold with $C_AW^m$ replaced by
$\tilde{a}$. Its smallness hypothesis $W^m\le\nu_0\alpha_4/(8C_A)$ becomes
$\tilde{a}\le\nu_0\alpha_4/8$, which is assumed, and its radius $\nu_0\alpha_4/(4C_AW^m)$
becomes $\lambda_0$; from $\tilde{a}\le\nu_0\alpha_4/8$,
\begin{equation*}
\lambda_0=\frac{\nu_0\alpha_4}{4\tilde{a}}\ge\frac{\nu_0\alpha_4}{4\nu_0\alpha_4/8}=2 .
\end{equation*}
The lemma thus gives: $t_\lambda(X)$ is analytic on $|\lambda|<\lambda_0$, and
$|t_\lambda(X)|\le b^{\mathrm{adm}}e^{-\kappa d_{k+1}(X)}$ there.

Put $\varphi(\lambda):=t_\lambda(X)-t_0(X)$. It is analytic on $|\lambda|<\lambda_0$, vanishes
at $\lambda=0$, and $|\varphi(\lambda)|\le 2b^{\mathrm{adm}}e^{-\kappa d_{k+1}(X)}$ there.
The Schwarz lemma, applied on the unit disc to the function
\begin{equation*}
\mu\longmapsto\frac{\varphi(\lambda_0\mu)}{2b^{\mathrm{adm}}e^{-\kappa d_{k+1}(X)}} ,
\end{equation*}
gives
\begin{equation*}
|\varphi(\lambda)|\le 2b^{\mathrm{adm}}e^{-\kappa d_{k+1}(X)}\,\frac{|\lambda|}{\lambda_0}
\qquad\text{for } |\lambda|<\lambda_0 .
\end{equation*}
Since $\lambda_0\ge2>1$, the point $\lambda=1$ lies in the disc, and this is the inequality in
\eqref{8.5:eq-direct-potentials-a}. The equality there is algebra:
\begin{equation*}
\frac{2b^{\mathrm{adm}}}{\lambda_0}=\frac{8b^{\mathrm{adm}}\tilde{a}}{\nu_0\alpha_4}
=\frac{K_{CE}}{C_A}\,\tilde{a} .
\end{equation*}
For $\tilde{a}^{\mathrm{dir}}=0$ the right side is $K_{CE}W^me^{-\kappa d_{k+1}(X)}$, the
bound of the corollary in (a).

By (c), the gauge-fixing pieces are among the potentials; they are therefore localised,
holomorphic and invariant, that is, marked potentials given directly. Their activity is
obtained by summing a bound on one block over the blocks of one $\pi_k$-cube. A $\pi_k$-cube
contains $(M/L)^4$ blocks, and
\begin{equation*}
L^4\Bigl(\frac{M}{L}\Bigr)^4=M^4 ;
\end{equation*}
thus summing \eqref{8.5:eq-direct-potentials-2} over the blocks of one cube gives, for the
pieces with coefficient $\zeta_k(p_y)$, the bound $C\,r\,M^4C_1^2(g_kT^{\mathbb{C}})^2$. For
coefficients with $|c(y)|\le\bar{c}$, the bound on one block defining $C_6$, summed over the
same $(M/L)^4$ blocks, gives at most $C_6M^4\bar{c}\,C_1^2\delta_k^2$, the factor $\bar{c}$
entering once through that bound; this is \eqref{8.5:eq-direct-potentials-4}.
\end{proof}

\begin{lemma}[Quadratic smallness of the source dependence]\label{8.5:lem-source-quadratic}
We use the notation of Notation~\ref{8.5:not-extraction-notation}. Here $n$ is the witness level, $r$
is the radius of the complex sources, with $r\le r_j<2r$, and $E^{(j)}(X,\cdot,y;w)$ are the local
families of the correlation theorem. Assume the following.
\begin{enumerate}
\item[(i)] The correlation theorem, in the following parts:
\begin{itemize}
\item its inductive hypotheses at every level $k\le n$, including the bounds
$\|E^{(j)}_\zeta\|^{\flat}\le E_0$ and, for the memory terms $R^{(j)}_\zeta$,
$\|R^{(j)}_\zeta\|^{\flat}\le g_j^{\kappa_0}$ for $|\zeta|<2r$,
and the bound $|\zeta_k|\le\frac43\sup|w|$ on the running shift;
\item the locality of the outputs of a step, their independence of the region, and their freezing
at constant sources;
\item the list of the places at which the source enters a renormalisation step;
\item the bounds for renormalised groups, for memory terms, for Wilson pieces and for the variation
of a local family when the source is frozen at a point;
\item the anchoring of units per cell, the summation of marked activities and the homogeneity of the
response map, on which the unit-sum bound (Lemma~\ref{8.5:lem-unit-sum}) rests;
\item the table of units;
\item the marking of a renormalised group by the difference of two families, carrying its own
counterterm coefficient;
\item the monotonicity of $\theta(\cdot)$ and $\alpha_*(\cdot)$ on $\mathopen]0,\gamma_J]$ and the bound
$g_j<g$ for $j<K$;
\item the elementary bounds on the flow of the couplings.
\end{itemize}
\item[(ii)] The cluster-expansion lemma for one renormalisation step, with its constants $C_A$,
$\nu_0$, $\alpha_4$, $b^{\mathrm{adm}}$, $\bar{\varepsilon}$ and
$K_{CE}:=8b^{\mathrm{adm}}C_A/(\nu_0\alpha_4)$. We write $\varepsilon_1$ for the analyticity radius
of the source-independent potentials $\mathbf{P}^{(k)}$ of a step.
\item[(iii)] The following conditions, in which
\begin{equation*}
C_s:=10K_{CE}c_S,\qquad
c_S:=\frac{64}{3}C_5+1+\frac{8}{3}\,\frac{C_6M^4C_1^2}{C_AK_R^2C_{\mathfrak{a}}^2}.
\end{equation*}
Here $C_6$ is the constant of the activity bound for the gauge-fixing terms
$\zeta^w_k(p_y)G_y(B')$ of a renormalisation step. Further, $\delta$ is the parameter by which the
bounds of the cluster-expansion lemma give up decay, as in its bound $e^{-(1-2\delta)\kappa d_k(Y)}$
for the marked potentials, and
\begin{equation*}
C_B:=\sup_X 2\Bigl(1+\frac{\mathrm{diam}_j(X^+)}{10M}\Bigr)e^{-\frac14\delta\kappa d_j(X)}
\end{equation*}
is the constant of the correlation theorem's bound for freezing the source at a point.
In these conditions $T(g'):=A_1(\log g'^{-2})^{\mathfrak{p}_0}$, so that $\delta_k=g_kT(g_k)$, and
$\theta(\cdot)$, $\mathfrak{m}(\cdot)$, $\alpha_*(\cdot)$ are the functions of the coupling of the
correlation theorem, with $\theta_k=\theta(g_k)$ and $\alpha_{*,k}=\alpha_*(g_k)$:
\begin{equation}\label{8.5:eq-source-quadratic-a}
\sup_{g'\le g}\Bigl[2\,\frac{g'T(g')}{\varepsilon_1}\,\bar{\varepsilon}
+8C_AC_{\boxplus}\,\theta(g')\bigl(\hat{S}_IE_0+\mathfrak{m}(g')\bigr)\Bigr]
\le\frac14\,\nu_0\alpha_4;
\end{equation}
\begin{equation}\label{8.5:eq-source-quadratic-b}
\sup_{g'\le g}40K_{CE}C_{\boxplus}\,\theta(g')\bigl(S'_I+C_BS'_\omega\bigr)\le\frac12
\quad\text{and}\quad \theta(g)\le\frac13;
\end{equation}
\begin{equation}\label{8.5:eq-source-quadratic-c}
\sup_{y\ge g^{-2}-4c_2}16C_{\boxplus}C_{\mathfrak{m}}\,y^{-(\kappa_0-2)/2}\cdot
\frac{2y}{C_{\mathfrak{a}}^2C_*^2(\log y)^{2q}}\le r;
\end{equation}
\begin{equation}\label{8.5:eq-source-quadratic-d}
\mathfrak{a}_k^2\le 2\,\mathfrak{a}_{k+1}^2\qquad\text{for } k<K;
\end{equation}
\begin{equation}\label{8.5:eq-source-quadratic-e}
C_s\,r\,\theta(g)\,C_{\mathfrak{a}}^2\,\alpha_*(g)^2\le b^{\mathrm{adm}}.
\end{equation}
\end{enumerate}

Then for $1\le j\le n$, $X\in D_j$, $y\in X\cap T^{(j)}$ and $w\in\mathcal{W}_j(X)$,
\begin{equation}\label{8.5:eq-source-quadratic-f}
\bigl|E^{(j)}(X,\cdot,y;w)-E^{(j)}(X,\cdot,y;0)\bigr|
\le C_s\,r\,\hat{\theta}_{j-1}\,\mathfrak{a}_j^2\,e^{-\kappa d_j(X)}
\qquad\text{on }\mathcal{U}^{\flat}_j(X).
\end{equation}
\end{lemma}

The left-hand sides of \eqref{8.5:eq-source-quadratic-a}--\eqref{8.5:eq-source-quadratic-c} are
suprema, over $g'\le g$ or over $y\ge g^{-2}-4c_2$, of functions that tend to zero as $g'\to0$,
respectively as $y\to\infty$. Each is
compared with a number fixed before $g$. Condition \eqref{8.5:eq-source-quadratic-d} follows, for $g$
small once $q$ is fixed, from the bound $x_{k+1}/x_k\in[1-b_+g^2,1+2c_2g^2]$ on the flow of the
couplings in the correlation theorem, with its constant $b_+$, together with
$\mathfrak{a}_k^2=C_{\mathfrak{a}}^2C_*^2(\log x_k)^{2q}/x_k$. All five are conditions that $g$ be small.

\begin{proof}
\emph{Set-up.} For $1\le j\le n$, $X\in D_j$, $y\in X\cap T^{(j)}$ and $w\in\mathcal{W}_j(X)$ put
\begin{equation*}
S^{(j)}(X,\cdot,y;w):=E^{(j)}(X,\cdot,y;w)-E^{(j)}(X,\cdot,y;0).
\end{equation*}
Put also
\begin{equation*}
\mathsf{s}_j:=\sup_{X,y}\ \sup_{w\in\mathcal{W}_j(X)}\ \sup_{\mathcal{U}^{\flat}_j(X)}
e^{\kappa d_j(X)}\bigl|S^{(j)}(X,\cdot,y;w)\bigr|.
\end{equation*}
We prove
\begin{equation}\label{8.5:eq-source-quadratic-1}
\mathsf{s}_j\le\sigma_j\,\mathfrak{a}_j^2,\qquad \sigma_j:=C_s\,r\,\hat{\theta}_{j-1},
\end{equation}
for $1\le j\le n$. This is \eqref{8.5:eq-source-quadratic-f}.

The proof is by induction on $j=k+1$, for $0\le k<n$. Put $\bar{\sigma}_k:=\max_{j\le k}\sigma_j$
and $\bar{\sigma}_0:=0$. At step $k$ we assume \eqref{8.5:eq-source-quadratic-1} for all
$1\le j\le k$. For $k=0$ this assumption is empty, and every sum over older levels below is empty.

\emph{Three facts used throughout.}
\begin{itemize}
\item The functions $\theta(\cdot)$ and $\alpha_*(\cdot)$ are nondecreasing on $\mathopen]0,\gamma_J]$, and
$g_j<g$ for $j\le k<K$. Hence
$\hat{\theta}_k\le\theta(g)$ and $\mathfrak{a}_k=C_{\mathfrak{a}}\alpha_*(g_k)\le C_{\mathfrak{a}}\alpha_*(g)$.
In particular $\theta_k\le\hat{\theta}_k\le\frac13$, by \eqref{8.5:eq-source-quadratic-b}.
\item Since $\hat{\theta}$ is a running maximum, $\sigma_j\le C_sr\hat{\theta}_{k-1}$ for $j\le k$.
Hence, for $k\ge1$, $\bar{\sigma}_k\le C_sr\hat{\theta}_{k-1}\le C_sr\hat{\theta}_k$; for $k=0$,
$\bar{\sigma}_0=0$.
\item The step from level $k$ produces its output on the enlarged tube with radii
$\hat{\alpha}_k=(1+2\beta_s)\alpha_\cdot(g_k)$, where $\beta_s=\frac14$. On this tube admissible
backgrounds have amplitude $\hat{\mathfrak{a}}_k=\frac32\mathfrak{a}_k$, and the ratios of amplitudes
are $\hat{\rho}_{k,j}=\frac32\rho_{k,j}$ with $\rho_{k,j}=\mathfrak{a}_k/\mathfrak{a}_j$. Hence
\begin{equation}\label{8.5:eq-source-quadratic-2}
\hat{\rho}_{k,j}^2\,\mathfrak{a}_j^2=\hat{\mathfrak{a}}_k^2=\frac94\,\mathfrak{a}_k^2,
\qquad \hat{\rho}_{k,j}^2=w_{k,j}.
\end{equation}
\end{itemize}

\emph{Step 1: the two kernels.} Fix $X\in D_{k+1}$, $y\in X\cap T^{(k+1)}$ and
$w\in\mathcal{W}_{k+1}(X)$. The classes are nested: $\mathcal{W}_{k+1}\subset\mathcal{W}_j$ for
$j\le k$. Hence for every $j\le k$ and every polymer $X'$ of level $j$ with $X'^+\subset X^+$, the
restriction of $w$ to $\mathrm{plaq}(X'^+)$ belongs to $\mathcal{W}_j(X')$.

By the locality statement of the correlation theorem for a step at level $k+1\le n$, the outputs
are analytic on $\mathcal{W}_{k+1}(X)$, local, and frozen at constant sources. By the same statement,
$E^{(k+1)}(X,\cdot,y;w)$ is the pinned output kernel with domain $X$, in the sense of
\cite[(3.47)]{B-CMP119}, of the final logarithm of \cite[(3.28)]{B-CMP119}. This logarithm is
computed on the potentials $V^w(Y)$, $Y\subset X$. The potential map of the cluster-expansion lemma
produces these potentials from three kinds of units:
\begin{enumerate}
\item[(a)] the source-independent potentials $\mathbf{P}^{(k)}$ of the step;
\item[(b)] the Wilson vertex $A(\zeta^w_k,\cdot)$, cut into the pieces
$A(\zeta^w_k\hat{\chi}_\square,\cdot)$ by the shift-cube cutoffs $\hat{\chi}_\square$, and the
gauge-fixing terms $\zeta^w_k(p_y)G_y(B')$;
\item[(c)] the renormalised inputs
\begin{equation*}
\sum_{j\le k}\Bigl\{\sum_y\mathfrak{i}_y\bigl[E^{(j)}_{w(p_y)}\bigr]
+\sum\bigl(D^{(j)}(w)\bigr)^{\mathrm{vac}}\Bigr\}
+\sum_{j\le k_1(k)}\bigl(R^{(j)}(w)\bigr)^{\mathrm{vac}},
\end{equation*}
where $D^{(j)}(X',\cdot,y;w):=E^{(j)}(X',\cdot,y;w)-E^{(j)}_{w(p_y)}$ is the local source
dependence of the correlation theorem, $R^{(j)}$ are the terms of the correlation theorem that we
call memory terms, $F^{\mathrm{vac}}:=F-F(1)$ is the subtraction of the value at the unit
configuration, and
$k_1(k)\le k$ is the level up to which the terms $R^{(j)}$ enter the step, as in Ba{\l}aban's
construction \cite{B-CMP119}.
\end{enumerate}
The kernel $E^{(k+1)}(X,\cdot,y;0)$ is the same kernel, computed on the potentials $V^0(Y)$.

The source enters the step at no other place. By the correlation theorem it enters a
renormalisation step at exactly four places:
\begin{itemize}
\item the Wilson weight $x_k(\cdot)-\zeta_k$;
\item the coefficient $a_k(y):=x_k-\zeta_k(p_y)$ of the exponential gauge-fixing factor of the
block $y$, together with the constant $\log\mathfrak{z}(a_k(y))$;
\item the terms;
\item constants.
\end{itemize}
The first place gives the Wilson pieces in (b). For the second,
$\exp[-a_k(y)G_y]=\exp[-x_kG_y]\exp[\zeta_k(p_y)G_y]$, and the second factor gives the
gauge-fixing terms in (b). The third place gives the units (c).

The following objects do not depend on the source: the real Gaussian measure, cut off by $\chi_k$
at $T_k$; the backgrounds; the operators; the partitions; the pointing
\cite[(3.30)--(3.42)]{B-CMP119}; and the expansion \cite[(3.44)--(3.47)]{B-CMP119}. The expression
\cite[(3.37)]{B-CMP119} has three members. The one-loop term and the term with the derivative of
the cut-off function do not depend on the source, so they cancel in $S^{(k+1)}$. The constants of
the fourth place and $\log\mathfrak{z}(a_k(y))$ do not depend on the field. They belong to the
constant $E^w$ of the sourced effective action, not to the family $E^{(k+1)}$.

\emph{Step 2: unmarked and marked units.} The potential map of the cluster-expansion lemma is
linear in the functional on which it acts. We therefore write $V^w=V^0+(V^w-V^0)$.

\emph{The unmarked part $V^0$.} At $w=0$ the running shift $\zeta^0_k$ vanishes, so the units (b)
vanish. The unmarked part therefore comes from three kinds of units: the potentials $\mathbf{P}^{(k)}$, the
groups $\mathfrak{i}_y[E^{(j)}_0]$, and the memory terms at $w=0$. It satisfies the format hypothesis
of the cluster-expansion lemma by the inductive hypothesis of the correlation theorem at source
zero. Its activity hypothesis is condition \eqref{8.5:eq-source-quadratic-a}, as the following two
estimates show.
\begin{itemize}
\item As functions of their field variable $\mathsf{A}$, the potentials $\mathbf{P}^{(k)}$ vanish at
$\mathsf{A}=0$ and are analytic on
$|\mathsf{A}|<\varepsilon_1$. Recall the Schwarz bound: a function analytic on a disc about a point
$c$ and bounded there satisfies $|f(z)-f(c)|\le 2\,(\sup|f|)\,|z-c|$ divided by the radius of the
disc. We apply it with centre $0$, radius $\varepsilon_1$, the bound $\bar{\varepsilon}$ for
$|\mathbf{P}^{(k)}|$ on that disc, and $|\mathsf{A}|\le\delta_k$ on the support of $\chi_k$. It gives
that the activity of the $\mathbf{P}^{(k)}$ on the support of $\chi_k$ is at most
$2(\delta_k/\varepsilon_1)\bar{\varepsilon}$.
\item For the groups at source zero and the memory terms at $w=0$ we apply the unit-sum bound,
Lemma~\ref{8.5:lem-unit-sum}, with the norm bounds $\mathfrak{n}^G_j=E_0$ and
$\mathfrak{n}^M_j=g_j^{\kappa_0}$. Its hypothesis $\theta_k\le\frac13$ holds by the first fact
above. It gives, for the activity sum $W$ of these units and with
$\mu_k:=\sum_{j\le k}w_{k,j}x_j^{-\kappa_0/2}$,
$C_AW\le 8C_AC_{\boxplus}\theta_k(\hat{S}_IE_0+\mu_k)$.
\end{itemize}
Taking $g'=g_k\le g$ in \eqref{8.5:eq-source-quadratic-a}, with $\delta_k=g_kT(g_k)$ and the bound
$\mu_k\le\mathfrak{m}_k$ of the correlation theorem, the total activity of the unmarked part is at
most $\frac14\nu_0\alpha_4$.

\emph{The marked part $V^w-V^0$.} This is the image, under the potential map, of the units listed
below. For each we give the norm bound with which it enters Lemma~\ref{8.5:lem-unit-sum}. Let
$\omega_{k,j}:=L^{-(k-j)/2}$, and recall the parameter $\delta$ and the constant $C_B$ from
hypothesis (iii) of the lemma.
\begin{enumerate}
\item[(1)] \emph{Groups.} Put $\hat{w}:=w(p_y)$ and $S^{(j)}_{\hat{w}}:=E^{(j)}_{\hat{w}}-E^{(j)}_0$.
Since the counterterm coefficient of a group is linear in the family, the difference of the two
groups is the renormalised group of $S^{(j)}_{\hat{w}}$, with its own counterterm coefficient
$b_j(\hat{w})-b_j(0)$:
\begin{equation*}
\mathfrak{i}_y\bigl[E^{(j)}_{\hat{w}}\bigr]-\mathfrak{i}_y\bigl[E^{(j)}_0\bigr]
=\mathfrak{i}_y\bigl[S^{(j)}_{\hat{w}}\bigr].
\end{equation*}
Here $S^{(j)}_{\hat{w}}$ is a pointed $E$-family, the difference of two. This is the device by
which the correlation theorem marks groups.

We show that $\|S^{(j)}_{\hat{w}}\|^{\flat}\le\mathsf{s}_j$. The constant function $\hat{w}$ on
$\mathrm{plaq}(X'^+)$ satisfies $|\hat{w}|<r_{k+1}\le r_j$ and has vanishing Lipschitz seminorm, so
it belongs to $\mathcal{W}_j(X')$. By freezing at constants, $E^{(j)}(X',\cdot,y;w)=E^{(j)}_{\hat{w}}$
when $w\equiv\hat{w}$ on $X'^+$. Therefore $S^{(j)}_{\hat{w}}$ is $S^{(j)}$ evaluated at the constant
source $\hat{w}$, and the bound follows. We take
\begin{equation*}
\mathfrak{n}^G_j=\mathsf{s}_j.
\end{equation*}

\item[(2)] \emph{Local source dependence.} The units are $(D^{(j)}(X',\cdot,y;w))^{\mathrm{vac}}$,
where $D^{(j)}\equiv0$ at $w=0$. By the definition of $D^{(j)}$, and since
$E^{(j)}(\cdot;0)=E^{(j)}_0$,
\begin{equation*}
D^{(j)}(w)=E^{(j)}(w)-E^{(j)}_{w(p_y)}=S^{(j)}(w)-S^{(j)}_{w(p_y)}.
\end{equation*}
The source-independent parts cancel.

Fix $(\mathbf{U},\mathbf{J})\in\mathcal{U}^{\flat}_j(X')$. The function
$w\mapsto S^{(j)}(X',(\mathbf{U},\mathbf{J}),y;w)$ is analytic on $\mathcal{W}_j(X')$ and bounded
by $\mathsf{s}_je^{-\kappa d_j(X')}$.

For $j<k$ we use the correlation theorem's bound for freezing the source at a point. It states:
if $F(X';\cdot)$ is analytic on $\mathcal{W}_j(X')$ with $|F|\le\mathfrak{n}$, $j<k$,
$w\in\mathcal{W}_k(X')$ and $c\in\mathrm{plaq}(X'^+)$, then
\begin{equation*}
|F(X';w)-F(X';w(c))|\le 2\mathfrak{n}\,\omega_{k,j}\,\bigl(1+\mathrm{diam}_j(X'^+)/10M\bigr).
\end{equation*}
We apply it with $c:=p_y$ and obtain
\begin{equation*}
|D^{(j)}|\le 2\mathsf{s}_j\,\omega_{k,j}\Bigl(1+\frac{\mathrm{diam}_j(X'^+)}{10M}\Bigr)
e^{-\kappa d_j(X')}
\le C_B\,\omega_{k,j}\,\mathsf{s}_j\,e^{-(1-\frac14\delta)\kappa d_j(X')}.
\end{equation*}

For $j=k$ that bound does not apply, since it requires $j<k$. We use instead
\begin{equation*}
|D^{(k)}|\le|S^{(k)}(w)|+|S^{(k)}_{w(p_y)}|\le 2\mathsf{s}_ke^{-\kappa d_k(X')}
\le C_B\,\omega_{k,k}\,\mathsf{s}_k\,e^{-(1-\frac14\delta)\kappa d_k(X')}.
\end{equation*}
This holds because $\omega_{k,k}=1$ and $C_B\ge2$: the quantity in the supremum defining $C_B$ is at
least $2$ at a polymer with $d_j=0$. We take
\begin{equation*}
\mathfrak{n}^M_j=C_B\,\omega_{k,j}\,\mathsf{s}_j.
\end{equation*}

\item[(3)] \emph{Memory terms.} The units are
$(R^{(j)}(X',\cdot;w)-R^{(j)}(X',\cdot;0))^{\mathrm{vac}}$ for $j\le k_1(k)$. By the inductive
hypothesis, $\|R^{(j)}\|^{\flat}\le g_j^{\kappa_0}$ uniformly on $\mathcal{W}_j$, so we take
\begin{equation*}
\mathfrak{n}^M_j=2g_j^{\kappa_0}.
\end{equation*}

\item[(4)] \emph{Wilson pieces.} The units are $A(\zeta^w_k\hat{\chi}_\square,\cdot)$. The inductive
hypothesis gives $|\zeta_k|\le\frac43\sup|w|<\frac83r$, so we take
\begin{equation*}
\bar{c}=\frac83\,r.
\end{equation*}

\item[(5)] \emph{The gauge-fixing potential.} The term $\zeta^w_k(p_y)G_y(B')$ is a directly given
marked potential. Its activity satisfies
\begin{equation*}
\tilde{a}^{\mathrm{dir}}\le\frac83\,r\,C_6M^4C_1^2\,\delta_k^2,
\end{equation*}
where the factor $\frac83r$ is the bound on $|\zeta^w_k(p_y)|$.
\end{enumerate}

\emph{Step 3: the sum.} By the induction hypothesis and \eqref{8.5:eq-source-quadratic-2},
\begin{equation*}
\hat{\rho}_{k,j}^2\,\mathsf{s}_j\le\sigma_j\,\hat{\rho}_{k,j}^2\,\mathfrak{a}_j^2
=\sigma_j\,\hat{\mathfrak{a}}_k^2\le\bar{\sigma}_k\,\hat{\mathfrak{a}}_k^2\qquad(j\le k).
\end{equation*}
The unit-sum bound reads the norm bounds of Step~2 through sums over the levels, with
$\nu=1+k-j$ and with $\hat{\alpha}$ the H\"older exponent of the correlation theorem's bound for
renormalised groups (not to be confused with the radii $\hat{\alpha}_k$). By the definition of the
$K$-free level series $S'_I$ and $S'_\omega$ of the
correlation theorem this gives
\begin{align*}
\sum_{j\le k}\nu^3\hat{\rho}_{k,j}^2(1+\hat{\rho}_{k,j})L^{-\hat{\alpha}(\nu-1)}\mathsf{s}_j
&\le S'_I\,\bar{\sigma}_k\,\hat{\mathfrak{a}}_k^2,\\
\sum_{j\le k}\hat{\rho}_{k,j}^2\,C_B\,\omega_{k,j}\,\mathsf{s}_j
&\le C_BS'_\omega\,\bar{\sigma}_k\,\hat{\mathfrak{a}}_k^2.
\end{align*}
For the memory terms, $\hat{\rho}_{k,j}^2=w_{k,j}$ and $g_j^{\kappa_0}=x_j^{-\kappa_0/2}$. Recalling
$\mu_k$ from Step~2 and the bound $\mu_k\le\mathfrak{m}_k$ of the correlation theorem,
\begin{equation*}
\sum_{j\le k_1(k)}\hat{\rho}_{k,j}^2\cdot 2g_j^{\kappa_0}\le 2\mu_k\le 2\mathfrak{m}_k.
\end{equation*}

For the direct potential, $\theta_k=K_R\delta_k/\alpha_{*,k}$, $\mathfrak{a}_k=C_{\mathfrak{a}}\alpha_{*,k}$,
$\mathfrak{a}_k\le\hat{\mathfrak{a}}_k$ and $\theta_k\le1$. Hence
\begin{equation*}
\delta_k^2=\frac{\theta_k^2\alpha_{*,k}^2}{K_R^2}
\le\frac{\theta_k\,\hat{\mathfrak{a}}_k^2}{(K_RC_{\mathfrak{a}})^2}.
\end{equation*}

The total marked activity is $\tilde{a}:=C_AW^m+\tilde{a}^{\mathrm{dir}}$, where $W^m$ is the
activity sum of the marked units. The marked part
is bounded by the unit-sum bound, Lemma~\ref{8.5:lem-unit-sum} in the form
\eqref{8.5:eq-unit-sum-a}, with the norm bounds $\mathfrak{n}^G_j$, $\mathfrak{n}^M_j$, $\bar{c}$ of
Step~2 and the sums above. The direct potential is added by the marked-sum bound with directly
given potentials, Lemma~\ref{8.5:lem-direct-potentials}. Together they give
\begin{equation}\label{8.5:eq-source-quadratic-3}
\begin{split}
\tilde{a}&\le C_A\theta_k\hat{\mathfrak{a}}_k^2\Bigl\{8C_{\boxplus}(S'_I+C_BS'_\omega)\bar{\sigma}_k
+\frac{64}{3}C_5r+\frac83\,r\,\frac{C_6M^4C_1^2}{C_AK_R^2C_{\mathfrak{a}}^2}\Bigr\}\\
&\quad+16C_AC_{\boxplus}\theta_k\mathfrak{m}_k .
\end{split}
\end{equation}

We bound the last term. By definition $\mathfrak{m}_k=C_{\mathfrak{m}}y_k^{-(\kappa_0-2)/2}$ and
\begin{equation*}
\mathfrak{a}_k^2=\frac{C_{\mathfrak{a}}^2C_*^2(\log x_k)^{2q}}{x_k}.
\end{equation*}
The flow of the couplings gives $x_k\in[y_k,2y_k]$, so
\begin{equation*}
\mathfrak{a}_k^2\ge\frac{C_{\mathfrak{a}}^2C_*^2(\log y_k)^{2q}}{2y_k}.
\end{equation*}
Condition \eqref{8.5:eq-source-quadratic-c} at $y=y_k$ therefore gives
\begin{equation*}
16C_{\boxplus}\mathfrak{m}_k\le r\,\mathfrak{a}_k^2\le r\,\hat{\mathfrak{a}}_k^2.
\end{equation*}
So the last term of \eqref{8.5:eq-source-quadratic-3} is at most $C_A\theta_k\hat{\mathfrak{a}}_k^2\cdot r$,
which accounts for the summand $1$ in $c_S$. Hence
\begin{equation}\label{8.5:eq-source-quadratic-4}
\tilde{a}\le C_A\theta_k\hat{\mathfrak{a}}_k^2\bigl\{8C_{\boxplus}(S'_I+C_BS'_\omega)\bar{\sigma}_k
+c_Sr\bigr\}.
\end{equation}

\emph{Step 4: closing.} We first verify the hypothesis $\tilde{a}\le\nu_0\alpha_4/8$ of
Lemma~\ref{8.5:lem-direct-potentials}.

The definition of $K_{CE}$ means
\begin{equation*}
\frac{\nu_0\alpha_4}{8}=\frac{C_Ab^{\mathrm{adm}}}{K_{CE}}.
\end{equation*}
Condition \eqref{8.5:eq-source-quadratic-b} at $g'=g_k\le g$ gives
\begin{equation*}
8C_{\boxplus}(S'_I+C_BS'_\omega)\theta_k\le\frac{1}{10K_{CE}},
\end{equation*}
and $c_S=C_s/(10K_{CE})$ by definition. Inserting both into \eqref{8.5:eq-source-quadratic-4},
using $\hat{\mathfrak{a}}_k^2=\frac94\mathfrak{a}_k^2$ and the bounds
$\bar{\sigma}_k\le C_sr\hat{\theta}_k$ (second fact) and $\theta_k\le\hat{\theta}_k$, gives
\begin{align*}
\tilde{a}&\le\frac{C_A}{K_{CE}}\,\hat{\mathfrak{a}}_k^2\Bigl\{\frac{\bar{\sigma}_k}{10}
+\frac{C_sr\theta_k}{10}\Bigr\}
\le\frac{C_A}{K_{CE}}\cdot\frac94\,\mathfrak{a}_k^2\cdot\frac{2}{10}\,C_sr\hat{\theta}_k\\
&=\frac{9}{20}\,\frac{C_A}{K_{CE}}\,C_sr\hat{\theta}_k\mathfrak{a}_k^2.
\end{align*}
By the first fact, $\hat{\theta}_k\le\theta(g)$ and $\mathfrak{a}_k\le C_{\mathfrak{a}}\alpha_*(g)$.
Condition \eqref{8.5:eq-source-quadratic-e} then gives
\begin{equation*}
\tilde{a}\le\frac{9}{20}\,\frac{C_Ab^{\mathrm{adm}}}{K_{CE}}<\frac{\nu_0\alpha_4}{8}.
\end{equation*}

We now apply \eqref{8.5:eq-direct-potentials-a} to the interpolation $V^0+\lambda(V^w-V^0)$. By
Step~1, its kernel at $\lambda=1$ is $E^{(k+1)}(X,\cdot,y;w)$ and at $\lambda=0$ is
$E^{(k+1)}(X,\cdot,y;0)$, so the difference is $S^{(k+1)}(X,\cdot,y;w)$. The bound holds on the
output tube with radii $\hat{\alpha}_k$. This tube contains $\mathcal{U}^{\flat}_{k+1}(X)$, by the
comparison $(1+\beta_s)\alpha_{\cdot,k+1}\le\hat{\alpha}_k$ of the correlation theorem. Since $X$,
$y$ and $w$ were arbitrary, \eqref{8.5:eq-source-quadratic-4} gives
\begin{equation*}
\mathsf{s}_{k+1}\le\frac{K_{CE}}{C_A}\,\tilde{a}
\le K_{CE}\theta_k\hat{\mathfrak{a}}_k^2\bigl\{8C_{\boxplus}(S'_I+C_BS'_\omega)\bar{\sigma}_k
+c_Sr\bigr\}.
\end{equation*}

By \eqref{8.5:eq-source-quadratic-d},
\begin{equation*}
\hat{\mathfrak{a}}_k^2=\frac94\,\mathfrak{a}_k^2\le\frac92\,\mathfrak{a}_{k+1}^2
\le 5\,\mathfrak{a}_{k+1}^2.
\end{equation*}
Therefore
\begin{align*}
\mathsf{s}_{k+1}&\le\mathfrak{a}_{k+1}^2\bigl\{40K_{CE}C_{\boxplus}\theta_k(S'_I+C_BS'_\omega)
\bar{\sigma}_k+5K_{CE}c_Sr\theta_k\bigr\}\\
&\le\mathfrak{a}_{k+1}^2\Bigl\{\frac12\,\bar{\sigma}_k+\frac12\,C_sr\theta_k\Bigr\}\\
&\le\mathfrak{a}_{k+1}^2\,C_sr\hat{\theta}_k=\sigma_{k+1}\,\mathfrak{a}_{k+1}^2 .
\end{align*}
The second inequality uses \eqref{8.5:eq-source-quadratic-b} at $g'=g_k$ and the identity
$5K_{CE}c_S=\frac12C_s$. The third uses $\bar{\sigma}_k\le C_sr\hat{\theta}_{k-1}\le C_sr\hat{\theta}_k$
for $k\ge1$ (and $\bar{\sigma}_0=0$ for $k=0$) and $\theta_k\le\hat{\theta}_k$.

This is \eqref{8.5:eq-source-quadratic-1} for $j=k+1$. It completes the induction, and with it the
proof of \eqref{8.5:eq-source-quadratic-f}.
\end{proof}

\begin{corollary}[Derivative bounds free of the source radius]\label{8.5:cor-source-derivative}
Assume the hypotheses of Lemma~\ref{8.5:lem-source-quadratic}, and let $n=K-s$ be the witness level.
For $1\le j\le n$ put
\begin{equation}\label{8.5:eq-source-derivative-1}
\mathfrak{n}_j:=C_s\,\hat{\theta}_{j-1}\,\mathfrak{a}_j^2 ,
\end{equation}
with $C_s$ the constant and $r$ the source radius of \eqref{8.5:eq-source-quadratic-f}. Let
$E^{(j)}_{\hat{w}}$, $D^{(j)}$ and
$R^{(j)}$ be the frozen families and the local difference and remainder terms of the sourced format
of the effective action in the correlation theorem. For $|\zeta|<r$ write
$S^{(j)}_\zeta:=E^{(j)}_\zeta-E^{(j)}_0$ for the source-dependent part of the frozen family, and
$\dot{S}^{(j)}:=\partial_\zeta S^{(j)}_\zeta|_{\zeta=0}$. For a test function $f$
(Definition~\ref{1.3:def-curvature-field}), read on plaquettes as $p\mapsto f(x(p))$, put
$\dot{D}^{(j)}(X,\cdot,y)[f]:=\partial_z D^{(j)}(X,\cdot,y;zf)|_{z=0}$ and
$\dot{R}^{(j)}(X,\cdot)[f]:=\partial_z R^{(j)}(X,\cdot;zf)|_{z=0}$. Then for $1\le j\le n$,
$X\in D_j$ and $y\in X\cap T^{(j)}$:
\begin{enumerate}
\item[(i)] $\|\dot{S}^{(j)}\|^{\flat}\le\mathfrak{n}_j$;
\item[(ii)] on $\mathcal{U}^{\flat}_j(X)$,
\begin{equation}\label{8.5:eq-source-derivative-2}
\bigl|\dot{D}^{(j)}(X,\cdot,y)[f]\bigr|
\le C_B\,\mathfrak{n}_j\,\omega_{n,j}\,\|f\|_\sharp\,e^{-(1-\frac14\delta)\kappa d_j(X)},
\end{equation}
and $\dot{D}^{(j)}(X,\cdot,y)[f]=0$ if $f\equiv0$ on $X^+$;
\item[(iii)] for every such $f$,
\begin{equation}\label{8.5:eq-source-derivative-3}
\bigl|\dot{R}^{(j)}(X,\cdot)[f]\bigr|\le g_j^{\kappa_0}\,\|f\|_\sharp\,r^{-1}\,e^{-\kappa d_j(X)},
\end{equation}
and $\dot{R}^{(j)}(X,\cdot)[f]=0$ if $f\equiv0$ on $X^+$.
\end{enumerate}
Here $\omega_{n,j}=L^{-(n-j)/2}$; $C_B$ is the constant, and $\delta$ the fraction of the decay
rate $\kappa$ given up, in the bound at level $k$ on the difference terms $D^{(j)}$, read at $k=n$.
\end{corollary}

\begin{proof}
Throughout we use Cauchy's inequality in the following form: if $\phi$ is analytic on an open disc
centred at $0$ and bounded in modulus there, then $|\phi'(0)|$ is at most the supremum of $|\phi|$ on
the disc divided by the radius of the disc.

\emph{Part (i).} Fix $j$, $X\in D_j$, $y\in X\cap T^{(j)}$ and a point $(\mathbf{U},\mathbf{J})$ of
$\mathcal{U}^{\flat}_j(X)$. Let $\zeta\in\mathbb{C}$ with $|\zeta|<r$. The constant source $w\equiv\zeta$
satisfies $|w|<r\le r_j$ and has vanishing Lipschitz seminorm, which is smaller than $\mathfrak{l}_j>0$;
by the definition of the classes of complex sources it therefore belongs to $\mathcal{W}_j(X)$, and so
does $w\equiv0$. Thus $D(0,r)$, read as the set of constant sources, lies in $\mathcal{W}_j(X)$. By the
clause of the sourced format of the effective action in the correlation theorem on the local terms
$E^{(j)}$ and $D^{(j)}$, the local term
$E^{(j)}(X,\cdot,y;w)$ is analytic in $w\in\mathcal{W}_j(X)$ and equals $E^{(j)}_{\hat{w}}$ when
$w\equiv\hat{w}$ on $X^+$. Restricting to constant sources, $\zeta\mapsto S^{(j)}_\zeta(X,(\mathbf{U},
\mathbf{J}),y)$ is analytic on $D(0,r)$, and it vanishes at $\zeta=0$ by its definition. Applying
\eqref{8.5:eq-source-quadratic-f} with $w\equiv\zeta$ gives, for $|\zeta|<r$,
\begin{equation}\label{8.5:eq-source-derivative-4}
\bigl|S^{(j)}_\zeta(X,(\mathbf{U},\mathbf{J}),y)\bigr|
\le C_s\,r\,\hat{\theta}_{j-1}\,\mathfrak{a}_j^2\,e^{-\kappa d_j(X)} .
\end{equation}
Cauchy's inequality at $0$ on $D(0,r)$ divides this bound by $r$:
\begin{equation*}
\bigl|\dot{S}^{(j)}(X,(\mathbf{U},\mathbf{J}),y)\bigr|
\le C_s\,\hat{\theta}_{j-1}\,\mathfrak{a}_j^2\,e^{-\kappa d_j(X)}=\mathfrak{n}_j\,e^{-\kappa d_j(X)} .
\end{equation*}
The radius $r$ has cancelled. Since $X$, $y$ and the point of $\mathcal{U}^{\flat}_j(X)$ were
arbitrary, this is the bound $\|\dot{S}^{(j)}\|^{\flat}\le\mathfrak{n}_j$.

\emph{Part (ii).} By the same clause, $D^{(j)}(X,\cdot,y;w)$ is analytic in $w$ and local:
it reads $w$ only on $X^+$. If $f\equiv0$ on $X^+$, then $zf\equiv0$ on $X^+$ for every $z$. Hence
$z\mapsto D^{(j)}(X,\cdot,y;zf)$ is constant and its derivative at $0$ vanishes. This proves the
second assertion of part (ii). It also gives \eqref{8.5:eq-source-derivative-2} when $\|f\|_\sharp=0$,
because $\|f\|_\sharp\ge\|f\|_\infty$ forces $f\equiv0$.

Let now $\|f\|_\sharp>0$. For $|z|<r/\|f\|_\sharp$ the restriction of $zf$ to $\mathrm{plaq}(X^+)$
belongs to $\mathcal{W}_n(X)$, by the inclusion property of the classes of complex sources for sources
of $\|\cdot\|_\sharp$-norm less than $r$: the modulus is controlled through
$\|zf\|_\infty\le|z|\,\|f\|_\sharp<r\le r_n$, and the Lipschitz seminorm through
$40M\,\mathrm{Lip}(zf)\le|z|\,\|f\|_\sharp$. Since
$\mathcal{W}_n(X)\subset\mathcal{W}_j(X)$ for $j\le n$, the map
$z\mapsto D^{(j)}(X,\cdot,y;zf)$ is analytic on this disc. The bound at level $k$ on the difference
terms $D^{(j)}$, read at $k:=n=K-s$, gives
on this disc, on $\mathcal{U}^{\flat}_j(X)$,
\begin{equation*}
\bigl|D^{(j)}(X,\cdot,y;zf)\bigr|\le C_B\,\omega_{n,j}\,\mathsf{s}_j\,e^{-(1-\frac14\delta)\kappa d_j(X)} ,
\end{equation*}
where $\mathsf{s}_j$ is the supremum of the source dependence
$|E^{(j)}(X,\cdot,y;w)-E^{(j)}(X,\cdot,y;0)|$ of the local terms over $X\in D_j$,
$y\in X\cap T^{(j)}$,
$w\in\mathcal{W}_j(X)$ and fields in $\mathcal{U}^{\flat}_j(X)$. By
\eqref{8.5:eq-source-quadratic-f}, valid for every $w\in\mathcal{W}_j(X)$ on
$\mathcal{U}^{\flat}_j(X)$, and $e^{-\kappa d_j(X)}\le1$, we have
$\mathsf{s}_j\le C_s r\hat{\theta}_{j-1}\mathfrak{a}_j^2=r\,\mathfrak{n}_j$. Cauchy's inequality on the
disc of radius $r/\|f\|_\sharp$ divides the bound by this radius:
\begin{equation*}
\bigl|\dot{D}^{(j)}(X,\cdot,y)[f]\bigr|
\le\frac{\|f\|_\sharp}{r}\,C_B\,\omega_{n,j}\,r\,\mathfrak{n}_j\,e^{-(1-\frac14\delta)\kappa d_j(X)} .
\end{equation*}
The right-hand side is that of \eqref{8.5:eq-source-derivative-2}; again $r$ cancels.

\emph{Part (iii).} By the clause of the sourced format in the correlation theorem on the remainder
terms $R^{(j)}$, $R^{(j)}(X,\cdot;w)$ is analytic and local in
$w\in\mathcal{W}_j(X)$. It is also bounded there, uniformly in $w$, by $g_j^{\kappa_0}e^{-\kappa d_j(X)}$.
If $f\equiv0$ on $X^+$, then $z\mapsto R^{(j)}(X,\cdot;zf)$ is constant by locality, so its derivative
at $0$ vanishes. If $f\not\equiv0$, then $\|f\|_\sharp>0$. As in part (ii), the restriction of $zf$ to
$\mathrm{plaq}(X^+)$ lies in $\mathcal{W}_j(X)$ for $|z|<r/\|f\|_\sharp$, so
$z\mapsto R^{(j)}(X,\cdot;zf)$ is analytic on this disc and
bounded there by $g_j^{\kappa_0}e^{-\kappa d_j(X)}$. Cauchy's inequality on the disc of radius
$r/\|f\|_\sharp$ gives \eqref{8.5:eq-source-derivative-3}.
\end{proof}

\begin{remark}[The radius-free cost of the field Hessian]\label{8.5:rem-hessian-cost}
Fix a level $j$ and recall $\alpha_{*,j}=g_jC_*(\log x_j)^q$, $\alpha_{3,j}=c_3\alpha_{*,j}$ and
$\mathfrak{a}_j=C_{\mathfrak{a}}\alpha_{*,j}$. On the tube $\mathcal{U}^{\flat}_j(X)$, the field
Hessian of a term of $\|\cdot\|^{\flat}$-norm $\mathfrak{n}$ costs the dimensionless number
$4\mathfrak{n}/\alpha_{3,j}^2$.
\begin{enumerate}
\item[(a)] Suppose the $\zeta$-derivative of the source-dependent part is bounded by the Cauchy
estimate from an a-priori supremum $E_0$ on a disc of source radius $\mathsf{r}$, so that
$\mathfrak{n}=E_0/(2\mathsf{r})$. Then the number is
\begin{equation}\label{8.5:eq-hessian-cost-1}
\mathfrak{h}_j(\mathsf{r})=\frac{4E_0}{2\mathsf{r}\,\alpha_{3,j}^2}
=\frac{2E_0}{c_3^2C_*^2}\cdot\frac{x_j}{\mathsf{r}(\log x_j)^{2q}}
\asymp\frac{x_j}{\mathsf{r}(\log x_j)^{2q}} .
\end{equation}
\item[(b)] With $\mathfrak{n}=\mathfrak{n}_j=C_s\hat{\theta}_{j-1}\mathfrak{a}_j^2$, the bound of
Corollary~\ref{8.5:cor-source-derivative}\textup{(i)}, the number is
\begin{equation}\label{8.5:eq-hessian-cost-2}
\frac{4\mathfrak{n}_j}{\alpha_{3,j}^2}=\frac{4C_sC_{\mathfrak{a}}^2}{c_3^2}\,\hat{\theta}_{j-1}.
\end{equation}
It does not depend on the source radius, its constant does not depend on the coupling, and it
decreases to zero with $\hat{\theta}_{j-1}$.
\end{enumerate}
\end{remark}

\begin{proof}
(a) Since $x_j=g_j^{-2}$, we have
$\alpha_{3,j}^2=c_3^2C_*^2g_j^2(\log x_j)^{2q}=c_3^2C_*^2(\log x_j)^{2q}/x_j$.
Substituting this into $4E_0/(2\mathsf{r}\alpha_{3,j}^2)$ gives the middle expression in
\eqref{8.5:eq-hessian-cost-1}, and the comparability holds because the prefactor
$2E_0/(c_3^2C_*^2)$ is fixed.

(b) By definition $\mathfrak{a}_j^2=C_{\mathfrak{a}}^2\alpha_{*,j}^2$ and
$\alpha_{3,j}^2=c_3^2\alpha_{*,j}^2$. Hence
\begin{equation*}
\frac{4\mathfrak{n}_j}{\alpha_{3,j}^2}
=\frac{4C_s\hat{\theta}_{j-1}C_{\mathfrak{a}}^2\alpha_{*,j}^2}{c_3^2\alpha_{*,j}^2}
=\frac{4C_sC_{\mathfrak{a}}^2}{c_3^2}\,\hat{\theta}_{j-1},
\end{equation*}
which is \eqref{8.5:eq-hessian-cost-2}.

The bound $\mathfrak{n}_j$ of Corollary~\ref{8.5:cor-source-derivative}\textup{(i)} contains no
source radius. In that corollary the factor $r$ in the supremum
$C_sr\hat{\theta}_{j-1}\mathfrak{a}_j^2$ on $D(0,r)$ cancels against the $1/r$ of the Cauchy
estimate at $0$. The right side of \eqref{8.5:eq-hessian-cost-2} is therefore free of the source
radius. Its prefactor $4C_sC_{\mathfrak{a}}^2/c_3^2$ is built from fixed constants only, so it does
not depend on the coupling. Since this prefactor is fixed, the number decreases to zero exactly as
$\hat{\theta}_{j-1}$ does, where $\hat{\theta}_{j-1}=\max_{i\le j-1}\theta_i$ and
$\theta_i=C_\theta(\log x_i)^{\mathfrak{p}_0-q}$.
\end{proof}

\begin{definition}[Regular critical backgrounds over an enlarged support]\label{8.5:def-critical-backgrounds}
Let $f$ be a test function whose support $S:=\operatorname{supp}f$ lies in a ball $B(x,\rho)$. For
$r>0$ write $S^{+r}$ for the $r$-neighbourhood of $S$, and put
\begin{equation*}
  \mathsf{C}:=S^{+\rho/8},\qquad \mathsf{C}^{\sharp}:=S^{+3\rho/16}.
\end{equation*}
Thus $\mathsf{C}\subset\mathsf{C}^{\sharp}\subset S^{+\rho/4}$. When $S=S_i$, the set $S^{+\rho/4}$ is the set
$N_i$ of the two-point witness theorem. Let $R_1M_1$ be the size of the big blocks in the class of
cubes of \cite[Theorem~1]{B-CMP102b}. Let $M_Q$ be the least multiple of $R_1M_1$ such that
\begin{equation*}
  2M_Q\ \ge\ 13M+2R_1M_1 .
\end{equation*}
The number $M_Q$ depends only on $d$, $L$, $M$ and $M_1$.

Let $n=K-s$ be the level at depth $s$ below the cutoff. Let $a_0$, $a_1$, $B_3$ be the constants of
\cite[Theorem~1]{B-CMP102b}. For $0<\varepsilon'\le a_1$, the class
$\mathfrak{R}(S;\varepsilon')$ is the set of real, that is $\mathrm{SU}(N)$-valued,
configurations $\mathfrak{b}$ on the bonds of the bare lattice lying over $\mathsf{C}^{\sharp}$ with
the following property. The configuration $\mathfrak{b}$ is the restriction to those bonds of a
minimal configuration $U_{\mathfrak{B}'}(V')$ of \cite[Theorem~1]{B-CMP102b} (the minimal
configuration of \cite[(2.12)]{B-CMP119}), for some determining set
\begin{equation*}
  \mathfrak{B}'=\bigcup_{j\le n}\Lambda'_j
\end{equation*}
with $\Lambda'_j$ the part of $\mathfrak{B}'$ on the lattice of level $j$, built over some
admissible sequence of domains $\{\Omega'_j\}_{j\le n}$ of \cite[Theorem~1]{B-CMP102b}, and for
some configuration $V'$, such that:
\begin{enumerate}
  \item[$(\mathfrak{R}1)$] $\mathsf{C}^{\sharp}\subset\Omega'_n$, so that over $\mathsf{C}^{\sharp}$ the
  level of $\mathfrak{B}'$ is $n$;
  \item[$(\mathfrak{R}2)$] $V'$ satisfies the condition \cite[(7)]{B-CMP102b} with
  $\varepsilon_1\le\varepsilon'$.
\end{enumerate}
Away from $\mathsf{C}^{\sharp}$, nothing is required of $\mathfrak{B}'$ and $V'$ beyond the condition
\cite[(7)]{B-CMP102b} of $(\mathfrak{R}2)$.

The definition does not assert the existence of a minimal configuration for given data. The class
$\mathfrak{R}(S;\varepsilon')$ consists of configurations that are minimal configurations of
\cite[Theorem~1]{B-CMP102b} and for which the conclusions \cite[(8)--(10)]{B-CMP102b} hold.
Wherever the class is used with particular data, existence for those data is supplied at that place.

\emph{The current.} For $\mathfrak{b}\in\mathfrak{R}(S;\varepsilon')$ let $J(\mathfrak{b})$ be the
current of \cite[(1.8)]{B-CMP109} evaluated at the background $\mathfrak{b}$. It is a local,
gauge-covariant function of $\mathfrak{b}$.

\emph{Functionals of a background.} Let $F$ be a functional of a polymer $X$ and of a pair
consisting of a configuration and a current, possibly pointed at a point $y$. For such $F$ we write
\begin{equation*}
  F(X,\mathfrak{b}):=F\bigl(X,(\mathfrak{b},J(\mathfrak{b}))|_{X}\bigr),\qquad
  F(X,\mathfrak{b},y):=F\bigl(X,(\mathfrak{b},J(\mathfrak{b}))|_{X},y\bigr),
\end{equation*}
where $|_{X}$ denotes restriction to $X$.

\emph{The ceiling on $\varepsilon'$.} In what follows $\varepsilon'$ is assumed to satisfy
\begin{equation}\label{8.5:eq-critical-backgrounds-a}
  B_3\varepsilon'\ \le\ \min\Bigl\{\bar{c}_2,\ \tfrac12 a_1,\ \frac{a_0}{2B_3}\Bigr\}
  \qquad\text{and}\qquad
  M_Q\ \le\ \frac{R_1M_1\,a_1}{\varepsilon'} ,
\end{equation}
where $\bar{c}_2$ is the threshold of \cite[Proposition~2]{B-CMP98}. The numbers $a_0$, $a_1$,
$B_3$, $\bar{c}_2$ do not depend on the coupling. Condition \eqref{8.5:eq-critical-backgrounds-a}
is therefore a smallness condition on $\varepsilon'$. In its uses, $\varepsilon'$ is of the order
of the coupling times a power of its logarithm.
\end{definition}

\begin{lemma}[Regularity, regular gauge and nesting of critical backgrounds]
\label{8.5:lem-background-regularity}
Let $0<\varepsilon'\le a_1$ satisfy \eqref{8.5:eq-critical-backgrounds-a}, let
$\mathfrak{b}\in\mathfrak{R}(S;\varepsilon')$ be a configuration of
Definition~\ref{8.5:def-critical-backgrounds}, and let $\mathfrak{B}'$, $\{\Omega'_m\}_{m\le n}$,
$V'$ be data for which $\mathfrak{b}$ is the restriction to the bare bonds over
$\mathsf{C}^{\sharp}$ of the minimal configuration $U_{\mathfrak{B}'}(V')$, as in that definition.
Here $a_0,a_1,B_3$ are the constants of \cite[Theorem~1]{B-CMP102b}, which depend on $d$ and
$L$ only, and $\eta:=L^{-n}$ is the spacing of the bare lattice in units of the level-$n$
lattice, so that $L^n\eta=1$. For $j\le n$ put $t:=L^{j-n}$, and for $m<j$ put
$t_m:=L^{m-n}$. Assume the following two statements.
\begin{enumerate}
\item[(a)] Ba{\l}aban's regularity theorem \cite[Theorem~1]{B-CMP102b}, used here as stated
there through its displays \cite[(2), (8), (9)]{B-CMP102b}.
\item[(b)] The nesting statement of the correlation theorem: a minimiser restricted to a box
with margin is the box minimiser of its own determining data. In the form used here: let
$j\le n$, let $\square_0$ be a box that is a union of cubes of $\pi_j$, with determining set
$\mathfrak{B}(\square_0)$, which consists of a part of level $j$ and of layers of levels
$m<j$ (the \emph{graded rim} of $\square_0$; a layer of level $m$ is said to be of order $m$);
let $\square_0^{+t}$ denote the enlargement of $\square_0$ by the margin with which the nesting
statement is stated; and let $\alpha>0$ satisfy
$\alpha\le\bar{c}_2$, $2\alpha\le a_1$ and $2B_3\alpha\le a_0$, where $\bar{c}_2$ is the
smallness constant of the nesting statement. If a configuration $U$ is regular at $\alpha$ in
level-$j$ units, that is, satisfies the conditions \cite[(2)]{B-CMP102b} defining
$\mathfrak{U}_j$ with $\alpha$ in place of $\varepsilon_0$, and is critical for a determining
set of level $\ge j$ over $\square_0^{+t}$, then the family $W:=M_{\mathfrak{B}(\square_0)}(U)$
of block averages of $U$ over $\mathfrak{B}(\square_0)$ obeys \cite[(7)]{B-CMP102b} with
$\varepsilon_1=2\alpha$ on the part of level $j$ and with $\varepsilon_1=2\alpha L^{-2(j-m)}$
on the layer of order $m<j$ (the latter by the averaging bound
\cite[Proposition~2, (54)]{B-CMP98}), and $U\restriction\square_0$ belongs to
the minimal orbit of $(\mathfrak{B}(\square_0),W)$.
\end{enumerate}
Then the following hold.
\begin{enumerate}
\item[(i)] For every bare plaquette $q$ over $\mathsf{C}^{\sharp}$,
\begin{equation}\label{8.5:eq-background-regularity-a}
|\mathfrak{b}(\partial q)-1|<B_3\varepsilon'\eta^2 .
\end{equation}
\item[(ii)] Let $\square$ be a cube of the class of \cite[Theorem~1]{B-CMP102b} of size
$2M_Q$ whose concentric cube of size $2M_Q+4R_1M_1$ lies in $\mathsf{C}^{\sharp}$. There is
an $\mathrm{SU}(N)$-valued gauge transformation $u$ such that
\begin{equation}\label{8.5:eq-background-regularity-b}
\mathfrak{b}^{u^{-1}}=e^{i\eta A}\ \text{on }\square,\qquad
|A|<B_3M_Q\varepsilon',\qquad |\nabla^{\eta}A|<B_3M_Q\varepsilon' .
\end{equation}
\item[(iii)] Let $j\le n$ and let $\square_0$ be a box that is a union of cubes of $\pi_j$ with
$\square_0^{+t}\subset\mathsf{C}^{\sharp}$. Put $\mathsf{V}:=M^{\cdot}(\mathfrak{b})$ on the
determining set of $\square_0$. Then, modulo gauge,
\begin{equation}\label{8.5:eq-background-regularity-c}
\begin{gathered}
(\mathfrak{b},J(\mathfrak{b}))\restriction\square_0
=\bigl(U_j(\square_0,\mathsf{V}),J_j(\square_0,\mathsf{V})\bigr),\\
|\partial\mathsf{V}-1|\le\mathfrak{e}\,t^2\ \text{per level-}j\text{ plaquette},\qquad
\mathfrak{e}:=2B_3\varepsilon',
\end{gathered}
\end{equation}
with $t_m^2\le t^2$ in place of $t^2$ on the plaquettes of level $m<j$ of the determining set
of $\square_0$ (its graded rim). Here $J(\mathfrak{b})$
is the current attached to $\mathfrak{b}$ and $J_j(\square_0,\mathsf{V})$ that of the box
minimiser $U_j(\square_0,\mathsf{V})$, both as in the correlation theorem.
\end{enumerate}
\end{lemma}

\begin{proof}
Write $U:=U_{\mathfrak{B}'}(V')$. By the second condition of
Definition~\ref{8.5:def-critical-backgrounds}, $V'$ obeys \cite[(7)]{B-CMP102b} with some
$\varepsilon_1\le\varepsilon'$; since that condition is an upper bound, $V'$ obeys it with
$\varepsilon_1=\varepsilon'$, and $\varepsilon'\le a_1$. Hence \cite[Theorem~1]{B-CMP102b}
applies with $\varepsilon_1=\varepsilon'$.

\emph{Proof of (i).} By \cite[(8)]{B-CMP102b}, the minimal configuration lies in the regular
space $\mathfrak{U}_n(\{\Omega'_m\},B_3\varepsilon')$. By the first condition of
Definition~\ref{8.5:def-critical-backgrounds}, $\mathsf{C}^{\sharp}\subset\Omega'_n$, so over
$\mathsf{C}^{\sharp}$ the level of $\mathfrak{B}'$ is $n$, and the plaquette condition of
\cite[(2)]{B-CMP102b} is read at the level $n$: it is
$|U(\partial q)-1|<B_3\varepsilon'\eta^2(L^n\eta)^{-2}$ for every bare plaquette $q$ over
$\mathsf{C}^{\sharp}$. Since $L^n\eta=1$ and $\mathfrak{b}$ coincides with $U$ there, this is
\eqref{8.5:eq-background-regularity-a}.

\emph{Proof of (ii).} We apply \cite[(9)]{B-CMP102b} at $j=n$, with $\varepsilon_1=\varepsilon'$
and $M=M_Q$. The number $M_Q$ is a multiple of $R_1M_1$, and $M_Q\le M(\varepsilon')$, the
bound on cube sizes in \cite[Theorem~1]{B-CMP102b}, by
\eqref{8.5:eq-critical-backgrounds-a}. At level $n$ a cube of the class has size
$2M_QL^n\eta=2M_Q$, and the class requires the concentric cube of size $2M_Q+4R_1M_1$ to lie
in $B^n(\Lambda'_n)$, the union of the level-$n$ blocks over the layer $\Lambda'_n$ of
$\mathfrak{B}'=\bigcup_{m\le n}\Lambda'_m$, in the notation of \cite[Theorem~1]{B-CMP102b};
this holds because that cube lies in
$\mathsf{C}^{\sharp}\subset\Omega'_n$ by the first condition of
Definition~\ref{8.5:def-critical-backgrounds}. Then \cite[(9)]{B-CMP102b} gives a gauge
transformation $u$, defined on a neighbourhood of $\square$, with $U^{u^{-1}}=e^{i\eta A}$ on
$\square$ and
\begin{equation*}
|A|<B_3M_Q\varepsilon'(L^n\eta)^{-1},\qquad
|\nabla^{\eta}A|<B_3M_Q\varepsilon'(L^n\eta)^{-2}.
\end{equation*}
With $L^n\eta=1$ and $U=\mathfrak{b}$ on $\square\subset\mathsf{C}^{\sharp}$ this is
\eqref{8.5:eq-background-regularity-b}.

\emph{Proof of (iii).} We apply statement (b) to $U$ with $\alpha:=B_3\varepsilon't^2$.
First, $U$ is regular at $\alpha$ in level-$j$ units over $\mathsf{C}^{\sharp}$. By
\cite[(8)]{B-CMP102b}, $U$ belongs to $\mathfrak{U}_n(\{\Omega'_m\},B_3\varepsilon')$, so over
$\mathsf{C}^{\sharp}\subset\Omega'_n$ it satisfies the conditions of \cite[(2)]{B-CMP102b} at
the level $n$ with $B_3\varepsilon'$ in place of $\varepsilon_0$; these are read in level-$j$
units through $L^j\eta=L^{j-n}=t$. For the plaquette member the reading is
\begin{equation*}
B_3\varepsilon'\eta^2=B_3\varepsilon't^2\,\eta^2\,t^{-2}=\alpha\,\eta^2(L^j\eta)^{-2},
\end{equation*}
so that by \eqref{8.5:eq-background-regularity-a} the plaquette bound of
\cite[(2)]{B-CMP102b} at level $j$ holds over $\mathsf{C}^{\sharp}$ with $\alpha$ in place of
$\varepsilon_0$. Second, $U$, being a minimal configuration for $\mathfrak{B}'$, is critical
for $\mathfrak{B}'$, whose level over $\mathsf{C}^{\sharp}$ is $n\ge j$; since
$\square_0^{+t}\subset\mathsf{C}^{\sharp}$, $U$ is critical for a determining set of level
$\ge j$ over $\square_0^{+t}$. Third, since $j\le n$ we have $t\le1$, hence
$\alpha\le B_3\varepsilon'$; the conditions $\alpha\le\bar{c}_2$, $2\alpha\le a_1$,
$2B_3\alpha\le a_0$ are upper bounds on $\alpha$, and at $\alpha=B_3\varepsilon'$ (the case
$t=1$) they are the conditions \eqref{8.5:eq-critical-backgrounds-a}.

Statement (b) therefore gives: $W:=M_{\mathfrak{B}(\square_0)}(U)$ obeys
\cite[(7)]{B-CMP102b} on the part of level $j$ of $\mathfrak{B}(\square_0)$ with
\begin{equation*}
\varepsilon_1=2\alpha=2B_3\varepsilon't^2=\mathfrak{e}\,t^2,
\end{equation*}
and on the layer of order $m<j$ with $\varepsilon_1=2\alpha L^{-2(j-m)}$, and
$U\restriction\square_0$ belongs to the minimal orbit of $(\mathfrak{B}(\square_0),W)$.
By definition $W$ is the family of averages of $U=\mathfrak{b}$ on the determining set of
$\square_0$, that is $W=\mathsf{V}=M^{\cdot}(\mathfrak{b})$; the bound of
\cite[(7)]{B-CMP102b} with $\varepsilon_1=\mathfrak{e}t^2$ is the bound
$|\partial\mathsf{V}-1|\le\mathfrak{e}t^2$ per level-$j$ plaquette in
\eqref{8.5:eq-background-regularity-c}. On the layer of the graded rim of order $m<j$,
since $t=L^{j-n}$,
\begin{equation*}
2\alpha L^{-2(j-m)}=\mathfrak{e}\,L^{2(j-n)}L^{-2(j-m)}=\mathfrak{e}\,L^{2(m-n)}
=\mathfrak{e}\,t_m^2 ,
\end{equation*}
which is the bound with $t_m^2$ in place of $t^2$; and $t_m^2\le t^2$ because $m<j$. The minimal
orbit of $(\mathfrak{B}(\square_0),\mathsf{V})$ is the orbit of the box minimiser
$U_j(\square_0,\mathsf{V})$, so modulo gauge
$\mathfrak{b}\restriction\square_0=U_j(\square_0,\mathsf{V})$. Finally, the current is local
and gauge covariant, so the gauge transformation carrying $\mathfrak{b}\restriction\square_0$
to $U_j(\square_0,\mathsf{V})$ carries $J(\mathfrak{b})\restriction\square_0$ to
$J_j(\square_0,\mathsf{V})$; the two pairs are therefore equivalent in the sense of the
correlation theorem's equivalence of configuration--current pairs, which is the meaning of
``modulo gauge'' in the first line of \eqref{8.5:eq-background-regularity-c}.
\end{proof}

\begin{lemma}[Vacuum subtraction at the field's own amplitude]\label{8.5:lem-vacuum-subtraction}
Let $n$ be the witness level, let $j\le n$, and put $t:=L^{j-n}$. Let $X\in D_j$ be a polymer at
level $j$, and let $F(X,\cdot)$ be a family on the tube $\mathcal{U}^{\flat}_j(X)$ that satisfies there
the two conditions imposed on families in the correlation theorem, with $|F|\le\mathfrak{n}$ on
$\mathcal{U}^{\flat}_j(X)$. Let $\square_X$, the box around $X$, be $X$ enlarged by two layers of
$\pi_j$-cubes; its side is at most
\begin{equation*}
D_X:=\operatorname{diam}_jX+5M\le 4M\bigl(d_j(X)+5\bigr),
\end{equation*}
and we assume $\square_X^{+t}\subset\mathsf{C}^{\sharp}$. Let
$\mathfrak{b}\in\mathfrak{R}(S;\varepsilon')$ (Definition~\ref{8.5:def-critical-backgrounds}), let
$\mathfrak{e}:=2B_3\varepsilon'$ be the field's own amplitude, and put
$\varrho_j:=\mathfrak{e}/\mathfrak{a}_j$. Let $C_M:=2dC_{\mathfrak{a}}/c_3$ be the constant of the
amplitude bound of the correlation theorem. Assume the following statements.
\begin{itemize}
\item[(a)] The analyticity statement for box minimisers of the correlation theorem: for a box
$\square_0$ around $X$ and complex data $B'$ on the determining set of $\square_0$ with
$|B'|<\alpha_{3,j}$, the pair $(U_j(\square_0,e^{iB'}),J_j(\square_0,e^{iB'}))$, restricted to $X$,
lies in $\mathcal{U}^{\flat}_j(X)$ and depends analytically on $B'$.
\item[(b)] The comb-gauge representation from the proof of the amplitude bound of the correlation
theorem: in the comb gauge centred at a point $y$, the box data are $\mathsf{V}=e^{iB^{\mathrm{ax}}}$,
and the variable of $\mathsf{V}$ at a bond $b$ is an ordered product of at most $d(1+|b-y|)$ plaquette
variables of $\mathsf{V}$, each conjugated by a transport along the comb.
\item[(c)] The Schwarz-type bound of the correlation theorem: if $f$ is analytic on the disc of centre
$c$ and radius $R$, $|f|\le M'$ there, and $f'(c)=0$, then $|f(z)-f(c)|\le 2M'|z-c|^2/R^2$.
\item[(d)] The identity \cite[(4.14)]{B-CMP109}, which is used below to see that the first
$\sigma$-derivative at $\sigma=0$ of the function $\varphi$ of the proof vanishes.
\end{itemize}
If $C_MD_Xt^2\varrho_j\le1$, then
\begin{equation}\label{8.5:eq-vacuum-subtraction-1}
\bigl|F(X,\mathfrak{b})-F(X,1)\bigr|\le 2\mathfrak{n}\bigl(C_MD_Xt^2\varrho_j\bigr)^2 .
\end{equation}
\end{lemma}

\begin{proof}
Let $\mathsf{V}$ denote the level-$j$ data of $\mathfrak{b}$ on the determining set of $\square_X$.
Since $\square_X^{+t}\subset\mathsf{C}^{\sharp}$, the nesting statement
\eqref{8.5:eq-background-regularity-c} of Lemma~\ref{8.5:lem-background-regularity} identifies the
restriction of $\mathfrak{b}$ to the box with the
box minimiser of its own determining data. Together with the second of the two conditions on $F$ this
gives
\begin{equation}\label{8.5:eq-vacuum-subtraction-2}
F(X,\mathfrak{b})=F\bigl(X,U_j(\square_X,\mathsf{V}),J_j(\square_X,\mathsf{V})\bigr).
\end{equation}

We now use the comb gauge of (b), centred at a point $y$ of $X$. In it
$\mathsf{V}=e^{iB^{\mathrm{ax}}}$, where at each bond $B^{\mathrm{ax}}(b)$ is Hermitian with spectrum
in $(-\pi,\pi]$.
The variable of $\mathsf{V}$ at a bond $b$ is an ordered product of at most
$d(1+|b-y|)\le dD_X$ plaquette variables of $\mathsf{V}$, each conjugated by a transport. Since
$\mathfrak{b}$ is real, that is $G$-valued, these transports are unitary. Two facts about unitary
matrices are used. First, for unitaries $P,Q$ one has
\begin{equation*}
|PQ-1|\le|P(Q-1)|+|P-1|=|Q-1|+|P-1|,
\end{equation*}
so by induction $|\prod_sP_s-1|\le\sum_s|P_s-1|$. Second, conjugation by a unitary does not change
$|\,\cdot-1|$. Each plaquette variable $\partial\mathsf{V}$ at level $j$ satisfies
$|\partial\mathsf{V}-1|\le\mathfrak{e}t^2$; this is the level-$j$ reading of the regularity clause of
Lemma~\ref{8.5:lem-background-regularity}, by which $\mathfrak{b}$ is the restriction of a
configuration in Ba{\l}aban's regular space with parameter $B_3\varepsilon'$. Therefore
\begin{equation}\label{8.5:eq-vacuum-subtraction-3}
\bigl|e^{iB^{\mathrm{ax}}(b)}-1\bigr|\le dD_X\,\mathfrak{e}\,t^2 .
\end{equation}

The right-hand side of \eqref{8.5:eq-vacuum-subtraction-3} can be rewritten using
$C_M=2dC_{\mathfrak{a}}/c_3$, $\mathfrak{a}_j=C_{\mathfrak{a}}\alpha_{*,j}$ and
$\alpha_{3,j}=c_3\alpha_{*,j}$. With the hypothesis $C_MD_Xt^2\varrho_j\le1$ this gives
\begin{equation*}
dD_X\mathfrak{e}t^2=\tfrac12c_3\alpha_{*,j}\cdot C_MD_Xt^2\varrho_j\le\tfrac12\alpha_{3,j}\le1 .
\end{equation*}
Now let $B$ be Hermitian with spectrum in $(-\pi,\pi]$ and $|e^{iB}-1|\le1$. For every eigenvalue
$\theta$ of $B$ we have $|e^{i\theta}-1|=2|\sin(\theta/2)|\le1$, hence $|\theta|\le\pi/3$. On
$[0,\pi/3]$ concavity of the sine gives $2\sin(\theta/2)\ge(3/\pi)\theta$. Hence
\begin{equation*}
|B|\le\tfrac{\pi}{3}\,|e^{iB}-1|\le2\,|e^{iB}-1| .
\end{equation*}
Applied at every bond, together with \eqref{8.5:eq-vacuum-subtraction-3}, this yields
\begin{equation}\label{8.5:eq-vacuum-subtraction-4}
\sup|B^{\mathrm{ax}}|\le 2dD_X\mathfrak{e}t^2=\alpha_{3,j}\cdot C_MD_Xt^2\varrho_j .
\end{equation}

For complex $\sigma$ define
\begin{equation*}
\varphi(\sigma):=F\bigl(X,U_j(\square_X,e^{i\sigma B^{\mathrm{ax}}}),
J_j(\square_X,e^{i\sigma B^{\mathrm{ax}}})\bigr),
\qquad
\Sigma:=\frac{\alpha_{3,j}}{\sup|B^{\mathrm{ax}}|}.
\end{equation*}
By \eqref{8.5:eq-vacuum-subtraction-4} and the hypothesis,
$\Sigma\ge(C_MD_Xt^2\varrho_j)^{-1}\ge1$. For $|\sigma|<\Sigma$ the complex data
$\sigma B^{\mathrm{ax}}$ satisfy $|\sigma B^{\mathrm{ax}}|<\alpha_{3,j}$. By the analyticity
statement (a), the pair entering $\varphi$
therefore lies, restricted to $X$, in $\mathcal{U}^{\flat}_j(X)$ and depends analytically on $\sigma$.
Hence $\varphi$ is analytic on $|\sigma|<\Sigma$, and $|\varphi|\le\mathfrak{n}$ there.

At $\sigma=1$ the data are $\mathsf{V}$, so $\varphi(1)=F(X,\mathfrak{b})$ by
\eqref{8.5:eq-vacuum-subtraction-2}. At $\sigma=0$ the data are trivial and $\varphi(0)=F(X,1)$. By
\cite[(4.14)]{B-CMP109}, $\varphi'(0)=0$. The Schwarz-type bound (c) applies with $c=0$, $z=1$,
$R=\Sigma$ and $M'=\mathfrak{n}$, and gives
\begin{equation*}
|F(X,\mathfrak{b})-F(X,1)|=|\varphi(1)-\varphi(0)|\le2\mathfrak{n}\,\Sigma^{-2}
\le2\mathfrak{n}\bigl(C_MD_Xt^2\varrho_j\bigr)^2,
\end{equation*}
which is \eqref{8.5:eq-vacuum-subtraction-1}.

The comb-gauge bound of the correlation theorem carries a restriction on the side $D$ of the box,
namely $D\le5M/t$. That
restriction controls the distortion produced by conjugation with complex transports. Here all
transports are unitary, so it plays no role.
\end{proof}

\begin{lemma}[The core estimate on a critical background]\label{8.5:lem-core-estimate}
Let $j\le n$, put $\nu:=1+n-j$ and $t:=L^{j-n}$, let $y\in T^{(j)}$, and let $\square\in\pi_j$ be the
cube containing $y$. Let $F$ be a pointed $E$-family at scale $j$ on the radii $\alpha_{\cdot,j}$ with
$\|F\|^{\flat}\le\mathfrak{n}$. Write $(F(X,\cdot,y))^{\mathrm{vac}}:=F(X,\cdot,y)-F(X,1,y)$ for the
vacuum subtraction, and put
\begin{equation}\label{8.5:eq-core-estimate-1}
\mathfrak{i}^c_y[F]:=\sum_{X\ni y,\ X\subset\square^{\sim\nu}}\bigl(F(X,\cdot,y)\bigr)^{\mathrm{vac}}
-\beta[F]\,A(h_y,\cdot),
\end{equation}
and
\begin{equation}\label{8.5:eq-core-estimate-2}
\mathfrak{N}_I:=C_I\,\mathfrak{n}\,\nu^3\rho_{n,j}^2(1+\rho_{n,j})\,t^{4+\hat{\alpha}}.
\end{equation}
Here $h_y$, the counterterm $\beta[F]A(h_y,\cdot)$, the constant $C_I$, the notion of a pointed
$E$-family and that of a background admissible at scale $n$ are those of the core estimate in the
proof of the correlation theorem.
Let $Q_y$ be the smallest box of $\pi_n$-cubes with $Q_y^{\sim-2}\supset\square^{\sim\nu+1}$ (it has at
most $11$ cubes on a side), and let $Q_y^+$ be $Q_y$ with one layer of $\pi_n$-cubes added. For a
configuration $\mathfrak{u}$ on $Q_y^+$ we write
$\operatorname{ext}_n(\mathfrak{u}):=(U_n,J_n)(Q_y,M^n(\mathfrak{u}))$ for the box minimiser on $Q_y$ of
the level-$n$ data of $\mathfrak{u}$, together with its current. Assume:
\begin{enumerate}
\item[(a)] (the core estimate of the proof of the correlation theorem, for the core and the
counterterm alone) for every background $\mathfrak{b}'$ admissible at scale $n$ on a set containing
$\square^{\sim\nu}\cup\operatorname{supp}h_y$ one has
$|\mathfrak{i}^c_y[F](\mathfrak{b}')|\le\mathfrak{N}_I$;
of the family, this estimate uses only the bound $\|F\|\le\mathfrak{n}$, its defining properties as a
pointed family, and the coefficient $\beta[F]$;
\item[(b)] (admissibility of box extensions) with $c''\le1$ a fixed number and $\mathfrak{r}_n$ the
radius for which the pairs $(e^{i\eta A},J)$ with $\max\{|A|,|\nabla A|,|J|\}\le\mathfrak{r}_n$ on a
set $X$ lie in the $\flat$-tube $\mathcal{U}^{\flat}_n(X)$: if, in one complex gauge,
$\mathfrak{u}=e^{i\eta A}$ on the bonds
of $Q_y^+$ with $\max\{|A|,|\nabla A|\}<c''\mathfrak{r}_n$, then $\operatorname{ext}_n(\mathfrak{u})$
is an admissible background at scale $n$ on $Q_y^{\sim-2}$; $\operatorname{ext}_n$ is analytic on this
ball, and $\operatorname{ext}_n(1)=(1,0)$; moreover, for every configuration $\mathfrak{u}$ on
$Q_y^+$ and every gauge transformation $v$ on the bonds of $Q_y^+$, the pair
$\operatorname{ext}_n(\mathfrak{u}^{v})$ is gauge equivalent on $Q_y$ to
$\operatorname{ext}_n(\mathfrak{u})$;
\item[(c)] the nesting property \eqref{8.5:eq-background-regularity-c} of
Lemma~\ref{8.5:lem-background-regularity};
\item[(d)] for every polymer $X\ni y$, the vanishing $DF(X;1,0)=0$ of the first derivative of
$F(X,\cdot,y)$ at $(1,0)$, in the form of \cite[(4.14)]{B-CMP109}.
\end{enumerate}
Suppose that $Q_y^+$ is contained in a cube to which the regular gauge
\eqref{8.5:eq-background-regularity-b} applies. Let $\mathfrak{b}\in\mathfrak{R}(S;\varepsilon')$
(Definition~\ref{8.5:def-critical-backgrounds}) and suppose
\begin{equation}\label{8.5:eq-core-estimate-3}
\epsilon:=\frac{B_3M_Q\varepsilon'}{c''\mathfrak{r}_n}\le1 .
\end{equation}
Then
\begin{equation}\label{8.5:eq-core-estimate-4}
\bigl|\mathfrak{i}^c_y[F](\mathfrak{b})\bigr|
\le2\epsilon^2\cdot C_I\,\mathfrak{n}\,\nu^3\rho_{n,j}^2(1+\rho_{n,j})\,t^{4+\hat{\alpha}}
=2\epsilon^2\,\mathfrak{N}_I.
\end{equation}
\end{lemma}

\begin{proof}
\emph{The gauge.} Since a cube of the class in \eqref{8.5:eq-background-regularity-b} contains
$Q_y^+$, that display gives a real gauge in which $\mathfrak{b}=e^{i\eta A}$ on $Q_y^+$, with
\begin{equation*}
\max\{|A|,|\nabla A|\}<B_3M_Q\varepsilon'=\epsilon\,c''\mathfrak{r}_n ,
\end{equation*}
the equality being the definition \eqref{8.5:eq-core-estimate-3}. The display
\eqref{8.5:eq-background-regularity-b} bounds $|\nabla^{\eta}A|$; in the units of level $n$, in which
the lemma is read, this is $|\nabla A|$. A real gauge is in
particular a complex gauge.

\emph{The function $\Phi$.} For complex $\sigma$ put
\begin{equation*}
\Phi(\sigma):=\mathfrak{i}^c_y[F]\bigl(\operatorname{ext}_n(e^{i\sigma\eta A})\bigr).
\end{equation*}
For $|\sigma|<1/\epsilon$ the potential $\sigma A$ satisfies
$\max\{|\sigma A|,|\nabla(\sigma A)|\}<c''\mathfrak{r}_n$ on the bonds of $Q_y^+$, so
$e^{i\sigma\eta A}$ lies in the ball of hypothesis (b) on the bonds of $Q_y^+$, that is, in $c''$
times the ball of radius $\mathfrak{r}_n$ of the $\flat$-tube inclusion. By hypothesis (b),
$\operatorname{ext}_n(e^{i\sigma\eta A})$ is then an admissible background at scale $n$ on
\begin{equation*}
Q_y^{\sim-2}\supset\square^{\sim\nu+1}\supset\square^{\sim\nu}\cup\operatorname{supp}h_y ,
\end{equation*}
and it depends analytically on $\sigma$, since $\operatorname{ext}_n$ is analytic on that ball and
$\sigma\mapsto e^{i\sigma\eta A}$ is entire.

\emph{The bound on the disc.} Hypothesis (a) applies on every admissible background, to the core and
the counterterm alone. Of the family it uses only the norm bound by $\mathfrak{n}$, the defining
properties of a pointed family and $\beta[F]$, all of which $F$ has. Applied at
$\mathfrak{b}'=\operatorname{ext}_n(e^{i\sigma\eta A})$ it gives
\begin{equation*}
|\Phi(\sigma)|\le\mathfrak{N}_I\qquad(|\sigma|<1/\epsilon).
\end{equation*}

\emph{Vanishing at $\sigma=0$.} At $\sigma=0$ the argument is
$\operatorname{ext}_n(1)=(1,0)$. There each vacuum-subtracted term of
\eqref{8.5:eq-core-estimate-1} vanishes. The counterterm $A(h_y,\cdot)$ vanishes as well, since it
is built from the plaquette functions $a_p$, which vanish at $1$. Hence $\Phi(0)=0$.

\emph{Vanishing of the derivative at $\sigma=0$.} By hypothesis (d), $DF(X;1,0)=0$. The map
$\sigma\mapsto\operatorname{ext}_n(e^{i\sigma\eta A})$ is analytic and takes $\sigma=0$ to
$(1,0)$. By the chain rule, each term $F(X,\operatorname{ext}_n(e^{i\sigma\eta A}),y)$ has
zero derivative at $\sigma=0$. The functions $a_p$ vanish to second order at $1$, so
$A(h_y,\operatorname{ext}_n(e^{i\sigma\eta A}))$ also has zero derivative at $\sigma=0$. Hence
$\Phi'(0)=0$.

\emph{The Schwarz-lemma bound.} Let $f$ be analytic on the disc $|z-c|<R$ with $|f|\le M'$ there and
$f'(c)=0$. The function $w\mapsto(f(c+Rw)-f(c))/(2M')$ is analytic on the unit disc, bounded by $1$,
and vanishes to second order at $w=0$. By the Schwarz lemma it is therefore bounded by $|w|^2$, so
that
\begin{equation*}
|f(z)-f(c)|\le2M'|z-c|^2/R^2 .
\end{equation*}
We apply this to $f=\Phi$ and $c=0$. Put $a:=\max\{|A|,|\nabla A|\}$, the maximum over the bonds of
$Q_y^+$; by the strict bound of the first step, $a<\epsilon c''\mathfrak{r}_n$. Let
$R':=c''\mathfrak{r}_n/a$ if $a>0$ and $R':=\infty$ if $a=0$; then $R'>1/\epsilon$. For
$|\sigma|<R'$ the potential $\sigma A$ still satisfies
$\max\{|\sigma A|,|\nabla(\sigma A)|\}<c''\mathfrak{r}_n$ on the bonds of $Q_y^+$, so the arguments
of the two steps defining and bounding $\Phi$ apply on this larger disc: $\Phi$ is analytic and
$|\Phi|\le\mathfrak{N}_I$ there. Take $R:=R'$ if $R'$ is finite and $R:=1/\epsilon+1$ otherwise. Then
$1/\epsilon<R\le R'$, and $R>1$ because $\epsilon\le1$, so $z=1$ lies in the disc $|z|<R$. With
$M'=\mathfrak{N}_I$, and since $\Phi(0)=0$ and $R>1/\epsilon$, we obtain
\begin{equation*}
|\Phi(1)|\le\frac{2\mathfrak{N}_I}{R^2}\le2\mathfrak{N}_I\epsilon^2 .
\end{equation*}

\emph{Identification of $\Phi(1)$.} Write $\mathfrak{b}_{Q_y^+}$ for the restriction of $\mathfrak{b}$
to the bonds of $Q_y^+$. The gauge of the first step is given by a real gauge transformation $u$ with
$\mathfrak{b}^{u^{-1}}=e^{i\eta A}$ on $Q_y^+$, as in \eqref{8.5:eq-background-regularity-b}; hence
$\Phi(1)=\mathfrak{i}^c_y[F](\operatorname{ext}_n(\mathfrak{b}^{u^{-1}}_{Q_y^+}))$. By the gauge
covariance of $\operatorname{ext}_n$ in hypothesis (b), the pair
$\operatorname{ext}_n(\mathfrak{b}^{u^{-1}}_{Q_y^+})$ is gauge equivalent on $Q_y$ to
$\operatorname{ext}_n(\mathfrak{b}_{Q_y^+})$. By the nesting property
\eqref{8.5:eq-background-regularity-c}, read with $(j,\square_0)$ replaced by $(n,Q_y)$, the pair
$\operatorname{ext}_n(\mathfrak{b}_{Q_y^+})$ is equivalent on $Q_y$, in the sense of the equivalence of
background pairs (a field together with its current) fixed with the correlation theorem, to the pair
$(\mathfrak{b},J(\mathfrak{b}))$ formed by $\mathfrak{b}$ and its current. The functional
$\mathfrak{i}^c_y[F]$ is invariant under these equivalences for two reasons. The terms
$F(X,\cdot,y)$ are invariant by the gauge property among the defining properties of a pointed
family. The counterterm is invariant because each $a_p$ is a class function. Therefore
$\Phi(1)=\mathfrak{i}^c_y[F](\mathfrak{b})$, and the bound on $|\Phi(1)|$ is
\eqref{8.5:eq-core-estimate-4}.
\end{proof}

\begin{remark}[Cancellation of the radii at the field amplitude]\label{8.5:rem-radii-cancel}
Let $j\le n$. Let
\begin{equation*}
\mathfrak{n}_j:=C_s\hat{\theta}_{j-1}\mathfrak{a}_j^2
\end{equation*}
be the bound on $\|\dot{S}^{(j)}\|^{\flat}$ given by
Corollary~\ref{8.5:cor-source-derivative}\,(i). Let $\varrho_j:=\mathfrak{e}/\mathfrak{a}_j$ be the ratio of
Lemma~\ref{8.5:lem-vacuum-subtraction}, where $\mathfrak{e}$ is the field amplitude of that lemma.
Let $\epsilon:=B_3M_Q\varepsilon'/(c''\mathfrak{r}_n)$ be the fraction of
Lemma~\ref{8.5:lem-core-estimate}.
Let $\rho_{n,j}:=\mathfrak{a}_n/\mathfrak{a}_j$. With $\mathfrak{n}=\mathfrak{n}_j$ in those two lemmas
the following identities hold:
\begin{align}
\mathfrak{n}_j\varrho_j^2&=C_s\hat{\theta}_{j-1}\mathfrak{e}^2,\label{8.5:eq-radii-cancel-1}\\
\epsilon^2\mathfrak{n}_j\rho_{n,j}^2&=C_s\hat{\theta}_{j-1}\epsilon^2\mathfrak{a}_n^2
=C_s\hat{\theta}_{j-1}\mathfrak{e}_Q^2,\label{8.5:eq-radii-cancel-2}
\end{align}
where
\begin{equation}\label{8.5:eq-radii-cancel-3}
\mathfrak{e}_Q:=\epsilon\,\mathfrak{a}_n=\frac{C_{\mathfrak{a}}B_3M_Q}{c''c_T}\,\varepsilon'.
\end{equation}
We put $\mathfrak{e}_\star:=\max\{\mathfrak{e},\mathfrak{e}_Q\}$.

The field amplitude, rather than the tube amplitude, has to be carried through for the following
reason. Let $y\in T^{(j)}$ and let $\mathfrak{b}\in\mathfrak{R}(S;\varepsilon')$; such a real background
belongs to the list of admissible backgrounds of the correlation theorem. If one applies to
$F=\dot{S}^{(j)}$ the general bound for the renormalised group $\mathfrak{i}_y[F]$ of a pointed family on
admissible backgrounds as it stands, with the tube amplitude, one obtains
\begin{equation*}
\bigl|\mathfrak{i}_y[\dot{S}^{(j)}](\mathfrak{b})\bigr|
\le C_I\mathfrak{n}_j\nu^3\rho_{n,j}^2(1+\rho_{n,j})\,t^{4+\hat{\alpha}}
=C_IC_s\hat{\theta}_{j-1}\mathfrak{a}_n^2\,\nu^3(1+\rho_{n,j})\,t^{4+\hat{\alpha}} .
\end{equation*}
Here $t=L^{j-n}$, $\nu=1+n-j$, and $C_I$ and $\hat{\alpha}$ are the constant and the H\"older exponent
of that bound. The amplitude in
this bound is that of the tube, $\mathfrak{a}_n$, a fixed multiple of $g_n(\log x_n)^q$. It is not the
amplitude of the field, $\varepsilon_n=g_nA_0(\log x_n)^{\mathfrak{p}_0}$. The supremum of the
irrelevant functional $\operatorname{Irr}$ would then be of order $\hat{\theta}_{n-1}\mathfrak{a}_n^2$
(since $\hat{\theta}_{j-1}\le\hat{\theta}_{n-1}$ for $j\le n$), that is, of order
$g_n^2(\log x_n)^{\mathfrak{p}_0+q}$. This grows with $q$, and no lower bound on $q$ could compensate
the factor $\rho^8$, where $\rho$ is the letter of the two-point witness theorem and not the ratio
$\rho_{n,j}$.

The factor $\epsilon^2=(\mathfrak{e}_Q/\mathfrak{a}_n)^2$ in Lemma~\ref{8.5:lem-core-estimate}
restores the field's own amplitude, and so does the factor $\varrho_j^2$ of
Lemma~\ref{8.5:lem-vacuum-subtraction}, which, set against the factor $\rho_{n,j}^2$ of the bound
above, amounts to $\varrho_j^2\rho_{n,j}^{-2}=(\mathfrak{e}/\mathfrak{a}_n)^2$. In each case this costs
one application of the Schwarz lemma,
because both functionals vanish to second order at the unit configuration~$1$.
\end{remark}

\begin{proof}
For \eqref{8.5:eq-radii-cancel-1} we insert $\varrho_j^2=\mathfrak{e}^2/\mathfrak{a}_j^2$ into the
definition of $\mathfrak{n}_j$:
\begin{equation*}
\mathfrak{n}_j\varrho_j^2
=C_s\hat{\theta}_{j-1}\mathfrak{a}_j^2\cdot\frac{\mathfrak{e}^2}{\mathfrak{a}_j^2}
=C_s\hat{\theta}_{j-1}\mathfrak{e}^2 .
\end{equation*}

For \eqref{8.5:eq-radii-cancel-2} we use $\rho_{n,j}^2=\mathfrak{a}_n^2/\mathfrak{a}_j^2$. This gives
$\mathfrak{n}_j\rho_{n,j}^2=C_s\hat{\theta}_{j-1}\mathfrak{a}_n^2$. Multiplying by $\epsilon^2$ gives the
first equality. The second equality is the definition $\mathfrak{e}_Q:=\epsilon\mathfrak{a}_n$.

It remains to evaluate $\epsilon\mathfrak{a}_n$ as in \eqref{8.5:eq-radii-cancel-3}. By definition,
$\mathfrak{a}_n=C_{\mathfrak{a}}\alpha_{*,n}$ and $\mathfrak{r}_n=c_T\alpha_{*,n}$. Hence
\begin{equation*}
\epsilon\,\mathfrak{a}_n
=\frac{B_3M_Q\varepsilon'}{c''c_T\alpha_{*,n}}\cdot C_{\mathfrak{a}}\alpha_{*,n}
=\frac{C_{\mathfrak{a}}B_3M_Q}{c''c_T}\,\varepsilon' .
\end{equation*}

Thus the tube amplitude $\mathfrak{a}_j$ cancels in \eqref{8.5:eq-radii-cancel-1}. The radius
$\alpha_{*,n}$ cancels in \eqref{8.5:eq-radii-cancel-3}, through $\mathfrak{a}_n$ and $\mathfrak{r}_n$.
Apart from the constant $C_s$, the right-hand sides contain only $\hat{\theta}_{j-1}$ and the field
amplitudes $\mathfrak{e}$ and $\mathfrak{e}_Q$, both at most $\mathfrak{e}_\star$; no tube radius remains.
\end{proof}

\begin{definition}[The kept set and the shortened response]\label{8.5:def-kept-set}
We keep the conventions of Notation~\ref{8.5:not-extraction-notation}: the witness level $n$, the
test function $f$ with support $S$, the set $S^+=S^{+C_hM}$ (where $C_hM$ is the width by which
the two-point witness theorem enlarges the support $S$), the sets
$\mathsf{C}=S^{+\rho/8}\subset\mathsf{C}^{\sharp}=S^{+3\rho/16}$,
the scale $\rho$, and the big-block sizes $R_1M_1$ and $M_Q$. For a set $A$ and $r\ge0$, $A^{+r}$ is
the set of points at $n$-distance at most $r$ from $A$. Subscripts or prefixes $n$ on distances,
diameters and sides mean that they are measured in units of the lattice $T^{(n)}$. For $j\le n$ we
write $t:=L^{j-n}$ and $\nu:=1+n-j$.

We take the following objects from the lemmas and the corollary listed here.
\begin{itemize}
\item From Lemma~\ref{8.5:lem-vacuum-subtraction}, for $X\in D_j$: the box $\square_X$ and its side
bound $D_X=\operatorname{diam}_jX+5M$, the constant $C_M$, and $\varrho_j=\mathfrak{e}/\mathfrak{a}_j$.
\item From Lemma~\ref{8.5:lem-core-estimate}, for $y\in T^{(j)}$ with $\square\in\pi_j$ the cube
containing $y$: the enlarged cube $\square^{\sim\nu}$, the boxes $Q_y$ and $Q_y^+$, the core
$\mathfrak{i}^c_y[\cdot]$ and the number $\epsilon$.
\item From Corollary~\ref{8.5:cor-source-derivative}: the derivatives $\dot{S}^{(j)}$, $\dot{D}^{(j)}$,
$\dot{R}^{(j)}$.
\end{itemize}

\emph{Evaluation plaquettes.} For $y\in T^{(j)}$ let $\square_y\in\pi_j$ be the cube containing $y$, and
let $\square_y^+$ be $\square_y$ enlarged by one $\pi_j$-layer. The evaluation plaquette $p_y$ is a
plaquette with $p_y\in\operatorname{plaq}(\square_y^+)$, and nothing further about its choice is used.
In the proof of the functional bound for $\operatorname{Irr}^{\mathrm{sh}}[f]$, the comparison lemma
of the correlation theorem, which compares a source-dependent analytic family $F(X;\cdot)$ at $w$
with its value at the constant source equal to $w(c)$, is applied with $c:=p_y$ to every pair
$(X,y)$ with $y\in X$, and it requires $c\in\operatorname{plaq}(X^+)$. For $X:=\square_y$ this
requirement is exactly $p_y\in\operatorname{plaq}(\square_y^+)$; since every $X\in D_j$ with $y\in X$
is a union of cubes of $\pi_j$ and so contains $\square_y$, this choice meets the requirement for all
such $X$. We also take $C_h\ge2$. This costs nothing: enlarging $C_h$ only raises lower bounds on the
constants of the separation scale, namely $16C_hM$, $C_hM/\vartheta$ and the first member
of~\eqref{8.5:eq-kept-set-a}, and the bound
\begin{equation*}
|S^+|\le(2\rho+2C_hM+2)^4\le80\rho^4
\end{equation*}
remains true under $16C_hM\le\rho$.

\emph{The kept set, the kept tails and the shortened response.} For $j\le n$ and $y\in T^{(j)}$ put
\begin{equation}\label{8.5:eq-kept-set-c}
\begin{split}
\mathcal{K}_j&:=\{X\in D_j:\ D_Xt\le\rho/16\ \text{and}\ C_MD_Xt^2\varrho_j\le1\},\\
\mathcal{T}_y&:=\{X\in\mathcal{K}_j:\ X\ni y,\ X\not\subset\square^{\sim\nu}\},\\
\operatorname{Irr}^{\mathrm{sh}}[f](\mathfrak{b})
&:=\sum_{j\le n}\Bigl\{\sum_{y\in T^{(j)}}f(p_y)\Bigl[\mathfrak{i}^c_y[\dot{S}^{(j)}](\mathfrak{b})
+\sum_{X\in\mathcal{T}_y}\bigl(\dot{S}^{(j)}(X,\cdot,y)\bigr)^{\mathrm{vac}}(\mathfrak{b})\Bigr]\\
&\qquad+\sum_{X\in\mathcal{K}_j}\sum_{y\in X\cap T^{(j)}}
\bigl(\dot{D}^{(j)}(X,\cdot,y)[f]\bigr)^{\mathrm{vac}}(\mathfrak{b})
+\sum_{X\in\mathcal{K}_j}\bigl(\dot{R}^{(j)}(X,\cdot)[f]\bigr)^{\mathrm{vac}}(\mathfrak{b})\Bigr\},
\end{split}
\end{equation}
for every configuration $\mathfrak{b}$ on the bare bonds over $\mathsf{C}^{\sharp}$. We call
$\mathcal{K}_j$ the \emph{kept set} at level $j$, $\mathcal{T}_y$ the \emph{kept tails} at $y$, and
$\operatorname{Irr}^{\mathrm{sh}}[f]$ the \emph{shortened response}.

\emph{Standing conditions.} We impose
\begin{equation}\label{8.5:eq-kept-set-a}
C_0A_0^2\ \ge\ \max\{8(C_hM+2M_Q+4R_1M_1+2),\ 16C_hM,\ 3200M\}.
\end{equation}
This is a lower bound on the constant $C_0A_0^2$ of the separation scale of the two-point witness
theorem, and it is free of $g$. Its last two members are lower bounds which that theorem already
imposes. Since $\rho\ge C_W(s)\ge C_0A_0^2$, the bound~\eqref{8.5:eq-kept-set-a} holds with
$\rho$ in place of $C_0A_0^2$. With $\varepsilon'=C_w\varepsilon_n$ we also impose
\begin{equation}\label{8.5:eq-kept-set-b}
\epsilon\le1\qquad\text{and}\qquad\mathfrak{e}\le\mathfrak{a}_n .
\end{equation}
Each of $\epsilon$ and $\mathfrak{e}/\mathfrak{a}_n$ is a constant times
$(\log x_n)^{\mathfrak{p}_0-q}$, and the degree $\mathfrak{p}_0-q$ is negative because the constants
are chosen with $q>\mathfrak{p}_0$.
So~\eqref{8.5:eq-kept-set-b} is one-sided: it asks $x_n$ to be large.

\emph{Geometry.} Assume~\eqref{8.5:eq-kept-set-a}, $C_h\ge2$, and the relations
$\rho\ge C_W(s)\ge C_0A_0^2$ and $L^s>20\rho$, which hold under the prefix of the two-point witness
theorem. Then the following hold.
\begin{enumerate}
\item[(a)] If $f(p_y)\neq0$, then $y\in S^{+2M}\subset S^+$.
\item[(b)] If $f(p_y)\neq0$, then some cube of the big-block grid of side $2M_Q$ (the cubes
of~\eqref{8.5:eq-background-regularity-b}) contains $Q_y^+$. Its enlargement in the sense
of~\eqref{8.5:eq-background-regularity-b} lies in $S^{+C_hM+2M_Q+4R_1M_1}\subset\mathsf{C}$.
\item[(c)] If $f(p_y)\neq0$ and $X\in\mathcal{T}_y$, then $\square_X^{+t}\subset\mathsf{C}$.
\item[(d)] Let $X\in\mathcal{K}_j$. Then $\dot{D}^{(j)}(X,\cdot,y)[f]$ and
$\dot{R}^{(j)}(X,\cdot)[f]$ vanish unless $X^+\cap S\neq\emptyset$, and in that case
$\square_X^{+t}\subset\mathsf{C}$.
\item[(e)] Each polymer $X$ of a core or of $\mathcal{K}_j$, each box $\square_X$, each $Q_y^+$, and
$S^{+\rho/4}$ has $n$-diameter at most $3\rho<\mathsf{N}_s/6$, where $\mathsf{N}_s:=N_TL^s$ is the
number of sites per side of $T^{(n)}$, namely the torus side of the correlation theorem read at
$m:=s$.
\end{enumerate}
Consequently every term of~\eqref{8.5:eq-kept-set-c} that is not identically zero is a function of
$\mathfrak{b}$ over $\mathsf{C}$, with the enlargements required by the hypotheses of
Lemma~\ref{8.5:lem-vacuum-subtraction} and Lemma~\ref{8.5:lem-core-estimate}.
\end{definition}

\begin{proof}
(a) The cube $\square_y$ has $n$-side $Mt$, and $x(p_y)\in\square_y^+$. Hence
\begin{equation*}
\operatorname{dist}_n(x(p_y),y)\le2Mt\le2M,
\end{equation*}
because $t=L^{j-n}\le1$ for $j\le n$. If $f(p_y)\neq0$, then $x(p_y)\in S$, so $y\in S^{+2M}$. Since
$C_h\ge2$, we have $S^{+2M}\subset S^{+C_hM}=S^+$.

(b) The cube $\square^{\sim\nu+1}$ has $n$-side $M(2\nu+3)t=M(2\nu+3)L^{1-\nu}$. For $\nu\ge1$ and
$L\ge3$ this is at most $5M$. The box $Q_y$ has at most eleven $\pi_n$-cubes of $n$-side $M$ a side,
so $Q_y^+$ has $n$-side at most $13M$. Since $2M_Q\ge13M+2R_1M_1$, some cube of the big-block grid
of side $2M_Q$ contains $Q_y^+$. By (a) we have $y\in S^+$, so the enlargement of that cube in the
sense of~\eqref{8.5:eq-background-regularity-b} lies in
$S^{+C_hM+2M_Q+4R_1M_1}$. By~\eqref{8.5:eq-kept-set-a} and $\rho\ge C_0A_0^2$,
\begin{equation*}
C_hM+2M_Q+4R_1M_1\le\rho/8-2 .
\end{equation*}
Hence this set lies in $S^{+\rho/8}=\mathsf{C}$.

(c) Let $X\in\mathcal{T}_y$. Then $y\in X\subset\square_X$. The box $\square_X$ has side at most
$D_X$ in units of $T^{(j)}$, and therefore $n$-diameter at most $D_Xt\le\rho/16$, because
$X\in\mathcal{K}_j$. By (a) we have $y\in S^+$, and $t\le1$, so
$\square_X^{+t}\subset S^{+C_hM+1+\rho/16}$. This set lies in $\mathsf{C}$ because
$C_hM+1\le\rho/16$.

(d) Suppose $X^+\cap S=\emptyset$. Then $f\equiv0$ on $X^+$, and both $\dot{D}^{(j)}(X,\cdot,y)[f]$
and $\dot{R}^{(j)}(X,\cdot)[f]$ vanish by Corollary~\ref{8.5:cor-source-derivative}. Suppose instead
$X^+\cap S\neq\emptyset$. Since $\square_X\supset X^+$, the box $\square_X$ meets $S$. As in (c),
$\square_X$ has $n$-diameter at most $\rho/16$, and $t\le1$, so
$\square_X^{+t}\subset S^{+1+\rho/16}$. This set lies in $\mathsf{C}$ because
$1+\rho/16\le\rho/8$, which holds since $\rho\ge3200M$ by~\eqref{8.5:eq-kept-set-a}.

(e) These sets have $n$-diameter at most $3\rho$. Moreover $\mathsf{N}_s\ge L^s$, and
\begin{equation*}
L^s\ge L^{r(s)}\ge20L\,C_W(s)>20\rho,
\end{equation*}
by the depth conditions $s\ge r(s)$ and $L^{r(s)}\ge20L\,C_W(s)$ of the two-point witness theorem
and $\rho<LC_W(s)$. Hence $3\rho<\mathsf{N}_s/6$. So all these sets embed in $\mathbb{Z}^4$, and no
further condition excluding wrapping around the torus is needed.

Finally, (c) and (d) give $\square_X^{+t}\subset\mathsf{C}\subset\mathsf{C}^{\sharp}$ for every
polymer whose term in~\eqref{8.5:eq-kept-set-c} is not identically zero. This is the
inclusion required by Lemma~\ref{8.5:lem-vacuum-subtraction}. Part (b) supplies the cube containing
$Q_y^+$ that Lemma~\ref{8.5:lem-core-estimate} requires, with its enlargement inside $\mathsf{C}$.
\end{proof}

\begin{theorem}[Bound on the shortened irrelevant response]\label{8.5:thm-shortened-bound}
Work in the setting of this section. Assume the ceilings under which
Lemma~\ref{8.5:lem-source-quadratic} is stated, the regularity ceiling of
Lemma~\ref{8.5:lem-background-regularity}, the ceiling~\eqref{8.5:eq-kept-set-b}, and the geometric
floor of this section. Let $f$ be a Lipschitz function with $\operatorname{supp} f=S\subset B(x,\rho)$,
where $\rho\ge C_W(s)$, and let $\mathfrak{b}\in\mathfrak{R}(S;\varepsilon')$ be arbitrary. Write
$S^+=S^{+C_hM}$, and for a set $A$ let $|A|$ be the number of $n$-cells meeting $A$. Then the shortened
irrelevant response $\operatorname{Irr}^{\mathrm{sh}}[f]$ of~\eqref{8.5:eq-kept-set-c} satisfies
\begin{equation}\label{8.5:eq-shortened-bound-1}
\bigl|\operatorname{Irr}^{\mathrm{sh}}[f](\mathfrak{b})\bigr|
\le\bigl[C^{\flat}\,\hat{\theta}_{n-1}\,\mathfrak{e}_\star^2+C_D\,r^{-1}\mathfrak{m}_n\bigr]
\,|S^+|\,\|f\|_\sharp ,
\end{equation}
where $\mathfrak{e}_\star$ is the amplitude of Remark~\ref{8.5:rem-radii-cancel} and
\begin{equation}\label{8.5:eq-shortened-bound-2}
C^{\flat}:=C_s\bigl(2C_IS'_I+64C_M^2M^2\mathsf{S}_2+C_BC_DS'_\omega\bigr),
\qquad C_D:=32\cdot 3^dC_M^2M^2\mathsf{S}_2 .
\end{equation}
Every constant in~\eqref{8.5:eq-shortened-bound-1} and~\eqref{8.5:eq-shortened-bound-2} depends only
on letters chosen before $r$ in the order of choice of the correlation theorem; none of $K$, $g_0$,
$m$, $N_T$, $s$, $\rho$, the position of $S$, or the choice of $(\mathfrak{B}',V')$ enters.
\end{theorem}

\begin{proof}
Throughout, $j\le n$ is a level, $t=L^{j-n}$ and $\nu=1+n-j$, so that $t\le1$ and
$t=L^{-(\nu-1)}$. We have $S^+=S^{+C_hM}\supset S^{+M}$. By the definition
\eqref{8.5:eq-kept-set-c} of $\operatorname{Irr}^{\mathrm{sh}}[f]$ and the triangle inequality,
$|\operatorname{Irr}^{\mathrm{sh}}[f](\mathfrak{b})|$ is at most the sum over $j\le n$ of four
contributions: the cores $\sum_{y}|f(p_y)|\,|\mathfrak{i}^c_y[\dot{S}^{(j)}](\mathfrak{b})|$, the kept
tails $\sum_y|f(p_y)|\sum_{X\in\mathcal{T}_y}|(\dot{S}^{(j)}(X,\cdot,y))^{\mathrm{vac}}(\mathfrak{b})|$,
the terms of $\dot{D}^{(j)}$ and the terms of $\dot{R}^{(j)}$, the last two summed over
$X\in\mathcal{K}_j$. We bound the four contributions in turn, after isolating a counting estimate
common to the last two.

\smallskip\noindent\textit{(i) The cores.} Every $y\in T^{(j)}$ with $f(p_y)\neq0$ lies in $S^+$,
because $S^+\supset S^{+M}$. A cell of level $n$ contains $t^{-4}$ points of $T^{(j)}$, hence
\begin{equation}\label{8.5:eq-shortened-bound-3}
\#\{y\in T^{(j)}:f(p_y)\neq0\}\le\#\bigl(T^{(j)}\cap S^+\bigr)\le|S^+|\,t^{-4}.
\end{equation}
For each such $y$ we apply the core estimate, Lemma~\ref{8.5:lem-core-estimate}, to the family
$F=\dot{S}^{(j)}$. By Corollary~\ref{8.5:cor-source-derivative}\,(i) its norm obeys
$\|\dot{S}^{(j)}\|^{\flat}\le\mathfrak{n}_j=C_s\hat{\theta}_{j-1}\mathfrak{a}_j^2$, so we may take
$\mathfrak{n}=\mathfrak{n}_j$; and by Lemma~\ref{8.5:lem-response-decomposition} the counterterm
coefficient is $\beta[\dot{S}^{(j)}]=b_j'(0)$, so that the quantity
$\mathfrak{i}^c_y[\dot{S}^{(j)}]$ of~\eqref{8.5:eq-kept-set-c} is the one the core estimate bounds.
The core estimate gives
\begin{equation*}
|\mathfrak{i}^c_y[\dot{S}^{(j)}](\mathfrak{b})|
\le2\epsilon^2C_I\,\mathfrak{n}_j\,\nu^3\rho_{n,j}^2(1+\rho_{n,j})\,t^{4+\hat{\alpha}} .
\end{equation*}
By Remark~\ref{8.5:rem-radii-cancel} the radii cancel:
$\epsilon^2\mathfrak{n}_j\rho_{n,j}^2=C_s\hat{\theta}_{j-1}\mathfrak{e}_Q^2$. Multiplying by
$|f(p_y)|\le\|f\|_\infty$ and by the count~\eqref{8.5:eq-shortened-bound-3}, the cores of level $j$
contribute at most
\begin{equation*}
\|f\|_\infty|S^+|\,t^{-4}\cdot2C_IC_s\hat{\theta}_{j-1}\mathfrak{e}_Q^2\,\nu^3(1+\rho_{n,j})\,
t^{4+\hat{\alpha}}
=2C_IC_s\hat{\theta}_{j-1}\mathfrak{e}_Q^2|S^+|\|f\|_\infty\,
\nu^3(1+\rho_{n,j})L^{-\hat{\alpha}(\nu-1)} .
\end{equation*}
Since $\hat{\theta}_k=\max_{i\le k}\theta_i$ is non-decreasing in $k$, we have
$\hat{\theta}_{j-1}\le\hat{\theta}_{n-1}$. Writing $j=n+1-\nu$, the sum over $j\le n$ becomes a sum
over $\nu\ge1$, and
\begin{equation*}
\sum_{\nu\ge1}\nu^3(1+\rho_{n,n+1-\nu})L^{-\hat{\alpha}(\nu-1)}\le S'_I .
\end{equation*}
Hence the cores contribute in total at most
\begin{equation}\label{8.5:eq-shortened-bound-4}
2C_IC_sS'_I\,\hat{\theta}_{n-1}\,\mathfrak{e}_Q^2\,|S^+|\,\|f\|_\infty .
\end{equation}

\smallskip\noindent\textit{(ii) The kept tails.} Let $y\in T^{(j)}$ and $X\in\mathcal{T}_y$, so
$X\ni y$, $X\not\subset\square^{\sim\nu}$ and $X\in\mathcal{K}_j$. By the correlation theorem, a
polymer $X\ni y$ with $X\not\subset\square^{\sim\nu}$ has $d_j(X)\ge\nu-1$. The family
$\dot{S}^{(j)}(X,\cdot,y)$ is bounded on $\mathcal{U}^{\flat}_j(X)$ by
$\mathfrak{n}_je^{-\kappa d_j(X)}$, by Corollary~\ref{8.5:cor-source-derivative}\,(i); and
$C_MD_Xt^2\varrho_j\le1$ because $X\in\mathcal{K}_j$. Lemma~\ref{8.5:lem-vacuum-subtraction} with
$\mathfrak{n}=\mathfrak{n}_je^{-\kappa d_j(X)}$ and the bound $D_X\le4M(d_j(X)+5)$ give
\begin{equation*}
\bigl|(\dot{S}^{(j)}(X,\cdot,y))^{\mathrm{vac}}(\mathfrak{b})\bigr|
\le2\mathfrak{n}_je^{-\kappa d_j(X)}\bigl(C_MD_Xt^2\varrho_j\bigr)^2
\le2\mathfrak{n}_j\varrho_j^2(4C_MM)^2(d_j(X)+5)^2t^4e^{-\kappa d_j(X)} .
\end{equation*}
We sum over the polymers. Writing
$e^{-\kappa d_j(X)}=e^{-\frac12\kappa d_j(X)}e^{-\frac12\kappa d_j(X)}$, using $d_j(X)\ge\nu-1$ in the
first factor and $\frac12\kappa\ge\frac14\delta\kappa$ in the second, and then the definition of
$\mathsf{S}_2$ (whose summands carry the further factor $|X|_{\pi_j}\ge1$), we get
\begin{equation*}
\sum_{X\ni y,\ d_j(X)\ge\nu-1}(d_j(X)+5)^2e^{-\kappa d_j(X)}
\le e^{-\frac12\kappa(\nu-1)}\mathsf{S}_2\le L^{-5(\nu-1)}\mathsf{S}_2 ,
\end{equation*}
the last step because $\kappa\ge10\log L$, the floor on $\kappa$ under which the correlation theorem
is proved. By Remark~\ref{8.5:rem-radii-cancel},
$\mathfrak{n}_j\varrho_j^2=C_s\hat{\theta}_{j-1}\mathfrak{e}^2$. Multiplying by $|f(p_y)|\le\|f\|_\infty$
and by~\eqref{8.5:eq-shortened-bound-3}, the kept tails of level $j$ contribute at most
\begin{equation*}
\|f\|_\infty|S^+|\,t^{-4}\cdot32C_M^2M^2\mathsf{S}_2\,C_s\hat{\theta}_{j-1}\mathfrak{e}^2\,t^4
L^{-5(\nu-1)} .
\end{equation*}
Since $\hat{\theta}_{j-1}\le\hat{\theta}_{n-1}$ and $\sum_{\nu\ge1}L^{-5(\nu-1)}\le2$, the kept tails
contribute in total at most
\begin{equation}\label{8.5:eq-shortened-bound-5}
64C_M^2M^2\mathsf{S}_2\,C_s\,\hat{\theta}_{n-1}\,\mathfrak{e}^2\,|S^+|\,\|f\|_\infty .
\end{equation}

\smallskip\noindent\textit{(iii) A counting estimate.} Let $T(X,\cdot,y)$, or $T(X,\cdot)$ in the
absence of a marked point, be a family defined for $X\in D_j$ (and $y\in X\cap T^{(j)}$) which vanishes
unless $X^+\cap S\neq\emptyset$ and which obeys, for every $X\in\mathcal{K}_j$,
\begin{equation}\label{8.5:eq-shortened-bound-6}
|T^{\mathrm{vac}}(X,\mathfrak{b},y)|\le2\mathfrak{n}_T\,e^{-(1-\frac14\delta)\kappa d_j(X)}
\bigl(C_MD_Xt^2\varrho_j\bigr)^2
\end{equation}
(respectively the same bound for $T^{\mathrm{vac}}(X,\mathfrak{b})$). We claim
\begin{equation}\label{8.5:eq-shortened-bound-7}
\sum_{X\in\mathcal{K}_j}\ \sum_{y\in X\cap T^{(j)}}|T^{\mathrm{vac}}(X,\mathfrak{b},y)|
\le C_D\,|S^+|\,\mathfrak{n}_T\,\varrho_j^2 ,
\end{equation}
and the same bound for $\sum_{X\in\mathcal{K}_j}|T^{\mathrm{vac}}(X,\mathfrak{b})|$. If
$X^+\cap S\neq\emptyset$, then $X$ contains a cube $c'\in\pi_j$ equal or adjacent to a cube
$c\in\pi_j$ meeting $S$. The cubes $c\in\pi_j$ meeting $S$ are disjoint, have volume $(Mt)^4$ in units
of $n$-cells, and lie in $S^{+Mt}\subset S^{+M}\subset S^+$ (recall $t\le1$); as $S^+$ is covered by
the $|S^+|$ cells of level $n$ meeting it, their number is at most
$\operatorname{vol}(S^+)(Mt)^{-4}\le|S^+|(Mt)^{-4}$, whatever the value of $Mt$. Each such $c$ has at
most $3^d$ cubes $c'$ equal or adjacent to it. Given $c'$, a polymer $X\ni c'$ carries at most
$M^4|X|_{\pi_j}$ points $y\in X\cap T^{(j)}$. Inserting $D_X\le4M(d_j(X)+5)$
into~\eqref{8.5:eq-shortened-bound-6} and using $(1-\frac14\delta)\kappa\ge\frac14\delta\kappa$, the
sum over $X\ni c'$ and the points $y$ is at most
\begin{equation*}
M^4\cdot2\mathfrak{n}_T(4C_MM)^2t^4\varrho_j^2
\sum_{X\ni c'}|X|_{\pi_j}(d_j(X)+5)^2e^{-\frac14\delta\kappa d_j(X)}
\le M^4\cdot2\mathfrak{n}_T(4C_MM)^2t^4\varrho_j^2\,\mathsf{S}_2 .
\end{equation*}
Multiplying by the number of pairs $(c,c')$ gives
\begin{equation*}
\sum_{X,y}|T^{\mathrm{vac}}|
\le|S^+|(Mt)^{-4}\cdot3^d\cdot M^4\cdot2\mathfrak{n}_T(4C_MM)^2t^4\varrho_j^2\cdot\mathsf{S}_2
=C_D\,|S^+|\,\mathfrak{n}_T\,\varrho_j^2 ,
\end{equation*}
which is~\eqref{8.5:eq-shortened-bound-7}. Without a marked point the factor $M^4|X|_{\pi_j}\ge1$ is
simply not used, and the same bound results.

\smallskip\noindent\textit{(iv) The terms of $\dot{D}^{(j)}$.} By
Corollary~\ref{8.5:cor-source-derivative}\,(ii), on $\mathcal{U}^{\flat}_j(X)$,
\begin{equation*}
|\dot{D}^{(j)}(X,\cdot,y)[f]|\le C_B\mathfrak{n}_j\omega_{n,j}\|f\|_\sharp\,
e^{-(1-\frac14\delta)\kappa d_j(X)} ,
\end{equation*}
and $\dot{D}^{(j)}(X,\cdot,y)[f]=0$ if $f\equiv0$ on $X^+$; the latter holds when $X^+\cap S=\emptyset$,
since $f$ vanishes off $S=\operatorname{supp}f$. For $X\in\mathcal{K}_j$ we have
$C_MD_Xt^2\varrho_j\le1$, and Lemma~\ref{8.5:lem-vacuum-subtraction} with $\mathfrak{n}$ equal to the
right-hand side of the last display yields
\eqref{8.5:eq-shortened-bound-6} with $\mathfrak{n}_T=C_B\mathfrak{n}_j\omega_{n,j}\|f\|_\sharp$. By
\eqref{8.5:eq-shortened-bound-7} and $\mathfrak{n}_j\varrho_j^2=C_s\hat{\theta}_{j-1}\mathfrak{e}^2$
(Remark~\ref{8.5:rem-radii-cancel}), level $j$ contributes at most
$C_BC_DC_s\hat{\theta}_{j-1}\mathfrak{e}^2\omega_{n,j}|S^+|\|f\|_\sharp$. With
$\hat{\theta}_{j-1}\le\hat{\theta}_{n-1}$ and $\sum_{j\le n}\omega_{n,j}\le S'_\omega$, the terms of
$\dot{D}^{(j)}$ contribute in total at most
\begin{equation}\label{8.5:eq-shortened-bound-8}
C_BC_DC_sS'_\omega\,\hat{\theta}_{n-1}\,\mathfrak{e}^2\,|S^+|\,\|f\|_\sharp .
\end{equation}

\smallskip\noindent\textit{(v) The terms of $\dot{R}^{(j)}$.} By
Corollary~\ref{8.5:cor-source-derivative}\,(iii) together with
Lemma~\ref{8.5:lem-vacuum-subtraction}, applied as in step~(iv), the family
$\dot{R}^{(j)}(X,\cdot)[f]$ satisfies the hypotheses of step~(iii) with
$\mathfrak{n}_T=g_j^{\kappa_0}\|f\|_\sharp/r$. Since $\rho_{n,j}=\mathfrak{a}_n/\mathfrak{a}_j$,
\begin{equation*}
\varrho_j=\frac{\mathfrak{e}}{\mathfrak{a}_j}=\frac{\mathfrak{e}}{\mathfrak{a}_n}\,\rho_{n,j}
\le\rho_{n,j}
\end{equation*}
by the ceiling~\eqref{8.5:eq-kept-set-b}. By~\eqref{8.5:eq-shortened-bound-7}, level $j$ contributes
at most $C_Dr^{-1}\rho_{n,j}^2g_j^{\kappa_0}|S^+|\|f\|_\sharp$. Now $\rho_{n,j}^2\le w_{n,j}$, because
$w_{n,j}=\hat{\rho}_{n,j}^2=\frac94\rho_{n,j}^2$, and $g_j^{\kappa_0}=x_j^{-\kappa_0/2}$, because
$x_j=g_j^{-2}$; hence, by the level-sum bound $\mu_n\le\mathfrak{m}_n$ established with the correlation
theorem,
\begin{equation*}
\sum_{j\le n}\rho_{n,j}^2g_j^{\kappa_0}\le\sum_{j\le n}w_{n,j}x_j^{-\kappa_0/2}=\mu_n
\le\mathfrak{m}_n .
\end{equation*}
The terms of $\dot{R}^{(j)}$ therefore contribute in total at most
\begin{equation}\label{8.5:eq-shortened-bound-9}
C_D\,r^{-1}\mathfrak{m}_n\,|S^+|\,\|f\|_\sharp .
\end{equation}

\smallskip\noindent\textit{Conclusion.} We add~\eqref{8.5:eq-shortened-bound-4},
\eqref{8.5:eq-shortened-bound-5}, \eqref{8.5:eq-shortened-bound-8} and~\eqref{8.5:eq-shortened-bound-9}.
Since $\mathfrak{e}\le\mathfrak{e}_\star$ and $\mathfrak{e}_Q\le\mathfrak{e}_\star$ by the definition
$\mathfrak{e}_\star=\max\{\mathfrak{e},\mathfrak{e}_Q\}$, and
$\|f\|_\infty\le\|f\|_\sharp=\max\{\|f\|_\infty,40M\operatorname{Lip}f\}$, the first three bounds
add up to at most
\begin{equation*}
C_s\bigl(2C_IS'_I+64C_M^2M^2\mathsf{S}_2+C_BC_DS'_\omega\bigr)\hat{\theta}_{n-1}\mathfrak{e}_\star^2
|S^+|\,\|f\|_\sharp
=C^{\flat}\hat{\theta}_{n-1}\mathfrak{e}_\star^2|S^+|\,\|f\|_\sharp ,
\end{equation*}
and the fourth is~\eqref{8.5:eq-shortened-bound-9}. This is~\eqref{8.5:eq-shortened-bound-1}.

\smallskip\noindent\textit{The level sums.} For each level $j$, with $\nu=n-j+1$, the number of terms
times their weight in the four contributions is as follows. For the cores,
$t^{-4}\cdot\nu^3(1+\rho_{n,j})t^{4+\hat{\alpha}}$ reduces to
$\nu^3(1+\rho_{n,j})L^{-\hat{\alpha}(\nu-1)}$, summed by $S'_I$. For the kept tails,
$t^{-4}\cdot t^4\cdot L^{-5(\nu-1)}$ is a geometric series in $\nu$. For the terms of
$\dot{D}^{(j)}$, $t^{-4}\cdot t^4\cdot\omega_{n,j}$ reduces to $\omega_{n,j}$, decaying in $\nu$ as
$L^{-(\nu-1)/2}$, summed by $S'_\omega$. For the terms of $\dot{R}^{(j)}$,
$t^{-4}\cdot t^4\cdot\rho_{n,j}^2g_j^{\kappa_0}$ is summed by the level-sum bound
$\mu_n\le\mathfrak{m}_n$, with
$\kappa_0>4$. No weight of a constituent without vacuum subtraction is summed over a level, and no
sum $\sum_{j\le n}1$ occurs. Hence every level sum is bounded by one of $S'_I$, $2$, $S'_\omega$,
$\mathfrak{m}_n$, and the constants $C^{\flat}$, $C_D$ of~\eqref{8.5:eq-shortened-bound-2} are formed
from $C_s$, $C_I$, $C_M$, $C_B$, $M$, $\mathsf{S}_2$, $S'_I$, $S'_\omega$ only, each of which is chosen
before $r$ in the order of choice of the correlation theorem; none of $K$, $g_0$, $m$, $N_T$, $s$,
$\rho$, the position of $S$ or the choice of $(\mathfrak{B}',V')$ enters $C^{\flat}$ or $C_D$.
\end{proof}

\begin{lemma}[The long constituents outside the kept set]\label{8.5:lem-long-constituents}
Work in the setting and under the hypotheses of Theorem~\ref{8.5:thm-shortened-bound}, with its
constants $C^{\flat}$ and $C_D$, with the kept sets $\mathcal{K}_j$ of
Definition~\ref{8.5:def-kept-set}, and with $\varepsilon'=C_w\varepsilon_n$, where the constant $C_w$
is fixed in the branch form of the response bound. Let $f$ be Lipschitz with
$\operatorname{supp}f=S\subset B(x,\rho)$, where $\rho\ge C_W(s)$. For $j\le n$ put $t:=L^{j-n}$ and
$\nu:=1+n-j$. With the notation of Definition~\ref{8.5:def-kept-set} and
\eqref{8.5:eq-kept-set-c}, for $y\in T^{(j)}$ put
\begin{equation*}
\mathcal{T}^{lg}_y:=\{X\in D_j\setminus\mathcal{K}_j:\ y\in X,\ X\not\subset\square^{\sim\nu}\}.
\end{equation*}
The \emph{long part} is
\begin{equation}\label{8.5:eq-long-constituents-1}
\begin{split}
\operatorname{Irr}^{lg}[f](\mathfrak{b}):=\sum_{j\le n}\Bigl\{
&\sum_{y\in T^{(j)}}f(p_y)\sum_{X\in\mathcal{T}^{lg}_y}
\bigl(\dot{S}^{(j)}(X,\cdot,y)\bigr)^{\mathrm{vac}}(\mathfrak{b})\\
&+\sum_{X\in D_j\setminus\mathcal{K}_j}\ \sum_{y\in X\cap T^{(j)}}
\bigl(\dot{D}^{(j)}(X,\cdot,y)[f]\bigr)^{\mathrm{vac}}(\mathfrak{b})
+\sum_{X\in D_j\setminus\mathcal{K}_j}\bigl(\dot{R}^{(j)}(X,\cdot)[f]\bigr)^{\mathrm{vac}}(\mathfrak{b})\Bigr\}.
\end{split}
\end{equation}
It is evaluated at every configuration $\mathfrak{b}$ such that
$(\mathfrak{b},J)\restriction X\in\mathcal{U}^{\flat}_j(X)$ for every polymer $X$ occurring in
\eqref{8.5:eq-long-constituents-1}, where $J$ is the current of the box minimisers attached to
$\mathfrak{b}$; this is the format of the correlation theorem. Assume moreover that
$\kappa\ge10\log L$, $L\ge3$ and $q>\mathfrak{p}_0$, and that the witness level satisfies
$x_n\ge g^{-2}-c_5$, where $g$ is the reference coupling and $c_5=2c_2$. Let
$\mathfrak{e}=2B_3\varepsilon'$,
$\mathfrak{e}_\star=\max\{\mathfrak{e},\mathfrak{e}_Q\}$, $\varrho_j=\mathfrak{e}/\mathfrak{a}_j$, and put
\begin{gather*}
\mathsf{a}_\star:=\min\Bigl\{\frac{\rho}{64M},\ \frac{\mathfrak{a}_n}{4C_MMC_{t\rho}\mathfrak{e}}\Bigr\},\qquad
c_{\mathsf{a}}:=\frac{C_{\mathfrak{a}}C_*}{8C_MMC_{t\rho}B_3C_wA_0},\\
\mathsf{a}_\downarrow(x):=\min\Bigl\{\frac{C_0A_0^2(\log x)^{2\mathfrak{p}_0}}{64M},\
c_{\mathsf{a}}(\log x)^{q-\mathfrak{p}_0}\Bigr\},\\
C_{lg}:=8\cdot3^d\mathsf{S}_0e^{5\kappa/2}(1+C_B),\qquad C^{\flat}_1:=C^{\flat}/C_s .
\end{gather*}
Define $\mathfrak{f}^{lg}[f]$ and impose the ceiling on the constants as follows:
\begin{equation}\label{8.5:eq-long-constituents-a}
\begin{gathered}
\mathfrak{f}^{lg}[f]:=C_{lg}e^{-\frac12\kappa\mathsf{a}_\star}
\bigl[C_s\hat{\theta}_{n-1}\mathfrak{a}_n^2+r^{-1}\mathfrak{m}_n\bigr]|S^+|\,\|f\|_\sharp,\\
\sup_{x\ge g^{-2}-c_5}C_{lg}e^{-\frac12\kappa\mathsf{a}_\downarrow(x)}
\max\Bigl\{\frac{(4C_MMC_{t\rho}c_{\mathsf{a}})^2(\log x)^{2(q-\mathfrak{p}_0)}}{C^{\flat}_1},\
\frac{1}{C_D}\Bigr\}\le1,\\
c_{\mathsf{a}}\bigl(\log(g^{-2}-c_5)\bigr)^{q-\mathfrak{p}_0}\ge1 .
\end{gathered}
\end{equation}
Under the floor \eqref{8.5:eq-kept-set-a} and the ceiling, which consists of the two conditions in
the second and third rows of \eqref{8.5:eq-long-constituents-a} (called below the first and the
second of the two conditions in \eqref{8.5:eq-long-constituents-a}),
\begin{equation}\label{8.5:eq-long-constituents-2}
\bigl|\operatorname{Irr}^{lg}[f](\mathfrak{b})\bigr|\le\mathfrak{f}^{lg}[f]
\le\bigl[C^{\flat}\hat{\theta}_{n-1}\mathfrak{e}_\star^2+C_Dr^{-1}\mathfrak{m}_n\bigr]|S^+|\,\|f\|_\sharp .
\end{equation}
\end{lemma}

\begin{proof}
\emph{The long polymers are large.} Let $X\in D_j\setminus\mathcal{K}_j$. By
Definition~\ref{8.5:def-kept-set}, either $D_Xt>\rho/16$ or $C_MD_Xt^2\varrho_j>1$. Since
$D_X\le4M(d_j(X)+5)$, the first alternative gives $d_j(X)+5>\rho/(64Mt)$. The second gives
$d_j(X)+5>1/(4C_MMt^2\varrho_j)$. Hence
\begin{equation}\label{8.5:eq-long-constituents-3}
d_j(X)>a_j:=\min\Bigl\{\frac{\rho}{64Mt},\ \frac{1}{4C_MMt^2\varrho_j}\Bigr\}-5 .
\end{equation}

\emph{Each long constituent.} Write $d=d_j(X)$. For a family $F$, the vacuum-subtracted value
$F^{\mathrm{vac}}=F-F(1)$ is a difference of two values of $F$. At a configuration with
$(\mathfrak{b},J)\restriction X\in\mathcal{U}^{\flat}_j(X)$, each of these two values is bounded by
the bound of $F$ on $\mathcal{U}^{\flat}_j(X)$: the first by the hypothesis on $\mathfrak{b}$, the
second because the unit configuration, restricted to $X$, lies in $\mathcal{U}^{\flat}_j(X)$. So
$F^{\mathrm{vac}}$ is at most twice that bound.
By Corollary~\ref{8.5:cor-source-derivative}:
\begin{itemize}
\item part (i) bounds a long tail, weighted by $f(p_y)$, by $2\mathfrak{n}_je^{-\kappa d}\|f\|_\infty$;
\item part (ii) bounds a long $\dot{D}^{(j)}$-term by
$2C_B\mathfrak{n}_j\omega_{n,j}\|f\|_\sharp e^{-(1-\frac14\delta)\kappa d}$;
\item part (iii) bounds a long $\dot{R}^{(j)}$-term by $2g_j^{\kappa_0}\|f\|_\sharp r^{-1}e^{-\kappa d}$.
\end{itemize}
Here $\mathfrak{n}_j=C_s\hat{\theta}_{j-1}\mathfrak{a}_j^2$.

\emph{One level.} We have $e^{-\kappa d}\le e^{-(1-\frac14\delta)\kappa d}$. By
\eqref{8.5:eq-long-constituents-3}, $d>a_j$, hence
\begin{equation*}
e^{-(1-\frac14\delta)\kappa d}\le e^{-\frac12\kappa a_j}e^{-\frac14\delta\kappa d}.
\end{equation*}
By part (ii) of Corollary~\ref{8.5:cor-source-derivative}, a $\dot{D}^{(j)}$-term vanishes when
$f\equiv0$ on $X^+$. We now use the anchoring count of the proof of
Theorem~\ref{8.5:thm-shortened-bound}. It bounds the sum of $e^{-\frac14\delta\kappa d_j(X)}$ over
the constituents of one kind at level $j$, each counted with its anchor points, by
$3^d\mathsf{S}_0|S^+|t^{-4}$. We use also $\|f\|_\infty\le\|f\|_\sharp$ and $\omega_{n,j}\le1$.
Level $j$ then contributes at most
\begin{equation}\label{8.5:eq-long-constituents-4}
3^d\mathsf{S}_0|S^+|\,t^{-4}e^{-\frac12\kappa a_j}\cdot
2\bigl[(1+C_B)\mathfrak{n}_j+g_j^{\kappa_0}/r\bigr]\|f\|_\sharp .
\end{equation}

\emph{The threshold in units of $\mathsf{a}_\star$.} Since $t=L^{-(\nu-1)}$ and
$\varrho_j=(\mathfrak{e}/\mathfrak{a}_n)\rho_{n,j}$, the bound
$L^{-(\nu-1)}\rho_{n,j}\le C_{t\rho}$ on the level ratios gives
\begin{equation*}
t^2\varrho_j=L^{-2(\nu-1)}\frac{\mathfrak{e}}{\mathfrak{a}_n}\rho_{n,j}
\le C_{t\rho}\frac{\mathfrak{e}}{\mathfrak{a}_n}L^{-(\nu-1)} .
\end{equation*}
Together with $1/t=L^{\nu-1}$, the definition \eqref{8.5:eq-long-constituents-3} yields
$a_j+5\ge L^{\nu-1}\mathsf{a}_\star$.

\emph{The geometric series.} Put
$u_\nu:=L^{4(\nu-1)}e^{-\frac12\kappa(L^{\nu-1}\mathsf{a}_\star-5)}$. By the preceding step,
$t^{-4}e^{-\frac12\kappa a_j}\le u_\nu$. Moreover,
\begin{equation*}
\frac{u_{\nu+1}}{u_\nu}=L^4e^{-\frac12\kappa\mathsf{a}_\star L^{\nu-1}(L-1)}
\le L^4e^{-\frac12\kappa\mathsf{a}_\star(L-1)} .
\end{equation*}
The right side is at most $e^{-1}\le\frac12$ as soon as
$\frac12\kappa\mathsf{a}_\star(L-1)\ge4\log L+1$. This holds when $\mathsf{a}_\star\ge1$,
$\kappa\ge10\log L$ and $L\ge3$, the last two being hypotheses of the lemma: in that case
$\frac12\kappa\mathsf{a}_\star(L-1)\ge10\log L$, and $6\log L\ge1$. That $\mathsf{a}_\star\ge1$
follows from the floor \eqref{8.5:eq-kept-set-a} and the ceiling
\eqref{8.5:eq-long-constituents-a}; this is shown in the last step. Hence
\begin{equation*}
\sup_\nu u_\nu=u_1=e^{5\kappa/2}e^{-\frac12\kappa\mathsf{a}_\star},
\qquad \sum_\nu u_\nu\le2u_1 .
\end{equation*}

\emph{Sum over levels.} For $j\le n$ we have $\hat{\theta}_{j-1}\le\hat{\theta}_{n-1}$, since
$\hat{\theta}$ is a running maximum. Also
$\mathfrak{a}_j^2=\mathfrak{a}_n^2/\rho_{n,j}^2\le2\mathfrak{a}_n^2$, by the lower bound
$\rho_{n,j}^2\ge\frac12$, valid for all $j\le n$, on the level ratios. Therefore
$\mathfrak{n}_j\le2C_s\hat{\theta}_{n-1}\mathfrak{a}_n^2$, and
\begin{equation*}
\sum_jt^{-4}e^{-\frac12\kappa a_j}\mathfrak{n}_j\le4u_1C_s\hat{\theta}_{n-1}\mathfrak{a}_n^2 .
\end{equation*}
By the same lower bound, $g_j^{\kappa_0}\le2\rho_{n,j}^2g_j^{\kappa_0}$. Since
$w_{n,j}=\hat{\rho}_{n,j}^2\ge\rho_{n,j}^2$ and $g_j^{\kappa_0}=x_j^{-\kappa_0/2}$, the level-sum
bound $\mu_n\le\mathfrak{m}_n$ gives
\begin{equation*}
\sum_jt^{-4}e^{-\frac12\kappa a_j}g_j^{\kappa_0}\le u_1\cdot2\mu_n\le2u_1\mathfrak{m}_n .
\end{equation*}
We sum \eqref{8.5:eq-long-constituents-4} over $j\le n$, then use $2\cdot4=8$ and
$2\cdot2\le8(1+C_B)$. This gives
\begin{equation*}
\bigl|\operatorname{Irr}^{lg}[f](\mathfrak{b})\bigr|
\le8\cdot3^d\mathsf{S}_0(1+C_B)u_1
\bigl[C_s\hat{\theta}_{n-1}\mathfrak{a}_n^2+r^{-1}\mathfrak{m}_n\bigr]|S^+|\,\|f\|_\sharp .
\end{equation*}
The right side is $\mathfrak{f}^{lg}[f]$, which proves the first inequality of
\eqref{8.5:eq-long-constituents-2}.

\emph{The ceiling.} We have $\varepsilon'=C_w\varepsilon_n$, $\varepsilon_n=g_nA_0(\log x_n)^{\mathfrak{p}_0}$
and $\mathfrak{a}_n=C_{\mathfrak{a}}g_nC_*(\log x_n)^q$. Hence
\begin{equation*}
\frac{\mathfrak{a}_n}{\mathfrak{e}}=\frac{C_{\mathfrak{a}}C_*}{2B_3C_wA_0}(\log x_n)^{q-\mathfrak{p}_0}
=4C_MMC_{t\rho}c_{\mathsf{a}}(\log x_n)^{q-\mathfrak{p}_0} .
\end{equation*}
Also $\rho\ge C_W(s)\ge C_0A_0^2(\log x_n)^{2\mathfrak{p}_0}$, the second inequality being the lower
bound on $C_W(s)$ from the two-point witness theorem. Inserting both into the definition of
$\mathsf{a}_\star$ gives $\mathsf{a}_\star\ge\mathsf{a}_\downarrow(x_n)$. Since
$\mathfrak{e}_\star\ge\mathfrak{e}$,
\begin{equation*}
(\mathfrak{a}_n/\mathfrak{e}_\star)^2\le(4C_MMC_{t\rho}c_{\mathsf{a}})^2(\log x_n)^{2(q-\mathfrak{p}_0)} .
\end{equation*}
The second member of $\mathsf{a}_\downarrow(x_n)$ is at least $1$ by the second of the two
conditions in \eqref{8.5:eq-long-constituents-a}: by hypothesis $x_n\ge g^{-2}-c_5$ and
$q-\mathfrak{p}_0>0$, so that $(\log x)^{q-\mathfrak{p}_0}$ is nondecreasing in $x$ on this range.
Together with
the floor \eqref{8.5:eq-kept-set-a}, this gives the inequality $\mathsf{a}_\star\ge1$ used in the
series step.

Finally, write $C_s\hat{\theta}_{n-1}\mathfrak{a}_n^2$ as
$C_s(\mathfrak{a}_n/\mathfrak{e}_\star)^2\hat{\theta}_{n-1}\mathfrak{e}_\star^2$, and use
$e^{-\frac12\kappa\mathsf{a}_\star}\le e^{-\frac12\kappa\mathsf{a}_\downarrow(x_n)}$. The first of the
two conditions in \eqref{8.5:eq-long-constituents-a} is then applied at $x=x_n$, which lies in the
range of the supremum because $x_n\ge g^{-2}-c_5$, to each of the two
summands of $\mathfrak{f}^{lg}[f]$. For the first summand,
\begin{equation*}
C_{lg}e^{-\frac12\kappa\mathsf{a}_\star}C_s\hat{\theta}_{n-1}\mathfrak{a}_n^2
\le C^{\flat}_1C_s\hat{\theta}_{n-1}\mathfrak{e}_\star^2=C^{\flat}\hat{\theta}_{n-1}\mathfrak{e}_\star^2 .
\end{equation*}
For the second summand,
$C_{lg}e^{-\frac12\kappa\mathsf{a}_\star}r^{-1}\mathfrak{m}_n\le C_Dr^{-1}\mathfrak{m}_n$. This is
the second inequality of \eqref{8.5:eq-long-constituents-2}.
\end{proof}

In the sense of Definition~\ref{2.3:def-order-table}, the first of the two conditions in
\eqref{8.5:eq-long-constituents-a} has degree $-\infty$: an exponential of minus a positive power
of $\log x$ stands against a power of $\log x$. In the order of choice it is imposed after $C_0$.

The long part $\operatorname{Irr}^{lg}[f]$ and the bound \eqref{8.5:eq-long-constituents-2} are used
only in the branch form of the response bound and in the concluding estimate of this section; they
do not enter Theorem~\ref{8.5:thm-shortened-bound}.

\begin{proposition}[Every reading evaluates the shortened functional]
\label{8.5:prop-readings-shortened}
Let $n=K-s$ be the level of the two-point witness theorem, and let $f_1,f_2$ be test functions
of that theorem. For $i\in\{1,2\}$ let $S_i=\operatorname{supp}f_i$, let $S_i^+$ be the set so
denoted in that theorem, and let $\mathsf{C}_i^{\sharp}=S_i^{+3\rho/16}$ and
$N_i=S_i^{+\rho/4}$. Assume the floor \eqref{8.5:eq-kept-set-a}. Assume, moreover, the
following statements.
\begin{itemize}
\item[(A)] For a union $\Omega$ of level-$n$ $M_1$-cubes, $\mathbf{B}_n(\Omega)$ is the
determining set of \cite[(2.13)]{B-CMP119}, and its level-$n$ domain contains
$\Omega^{\sim-2}$; the domain-local minimiser $U_{n,\Omega}(Q_nV)$ is the minimal configuration
$U_{\mathbf{B}_n(\Omega)}(Q_nV)$ for the determining set $\mathbf{B}_n(\Omega)$ and the data
$Q_nV$ used in \cite[(2.16)--(2.17)]{B-CMP119}, and it is a function of the restriction of $V$
to the bonds of $\Omega$.
\item[(B)] The prolongation lemma for domain-local minimisers, in the following three parts.
Existence: for $V$ in the level-$n$ window on $N_i$ (defined below), the prolonged data of the
domain-local minimiser exist and lie, at every level $k$ of the tower, in Ba{\l}aban's regular
space $\mathfrak{U}_k(\{\Omega_m\},2B_3\alpha)$. Plaquettes: on every layer, and across layers
with Ba{\l}aban's barred bonds (the bonds of the determining set that join its layers),
the plaquette variables of $Q_nV$ are $\mathbf{1}$ or conjugates of some $V(\partial P)$ with
$P$ a level-$n$ plaquette of the domain. Nesting: if no live label of a history $h$ meets
$N_i$, then $N_i$ lies in the level-$n$ small-field region $\Omega_n(h)$ of $h$, and
$N_i\subset\Lambda_n(h)$.
\item[(C)] From the two-point witness theorem: its regularity statement, by which, for every
value of the variables $\mathsf{v}$ that $\mathbf{T}_h$ still integrates, the data of the
background $U^h_n$, the minimal configuration for the determining set $\mathfrak{B}_h$ of $h$
with the data given by $V$ and $\mathsf{v}$, satisfy the smallness condition of
Definition~\ref{8.5:def-critical-backgrounds} with
$\varepsilon_1\le C_rL^2\max_{j\le n}\varepsilon_j$; its positivity statement, by which the
measure $\mathbb{E}_h(d\mathsf{v})$ of $\mathsf{v}$, read with the whole exponent at $z=0$
under the integral, is a probability measure on $\operatorname{supp}m_h$, where $m_h$ is the
measure against which $\mathbf{T}_h$ integrates; the inequality
$M\check{R}_jL^{j-n}\le\rho/8$ for $j\le n$, which that theorem obtains from its lower bound on
$\rho$, where $M\check{R}_j$ is, in level-$j$ units, the side of the cubes of which one layer
lies between the correlation theorem's region $\Lambda_j^{\circledR}$ (where it places its
level-$j$ groups) and the complement of $\Lambda_j$, so that the left member is the width of
one such layer in level-$n$ units; and its exceptional set, denoted here $\mathcal{E}_i$: the
set of pairs $(h,V)$ such that some live label of $h$ meets $N_i$ or $V$ is not in the
level-$n$ window on $N_i$.
\item[(D)] From the correlation theorem: terms whose polymer $X$ lies in $\Lambda_j$ do not
depend on $\Lambda_j$ and coincide with the corresponding terms arising from the expressions
defined on the whole lattice, a statement it takes from \cite{B-CMP119}; the comparison
$x_j\ge x_n-2c_2$ for $j\le n$ of the running inverse couplings; for $p\ge0$ and $0<g<1$, the
equivalence of
$\frac{d}{dg}\bigl[g(\log g^{-2})^{p}\bigr]>0$ with $\log g^{-2}>2p$; the nesting
$\Lambda_j^{c}\subset\Lambda_k^{c}$ for $j\le k$; on $N_i$ the localising functions
$\varphi_j$ of the correlation theorem are identically $1$; and its sourced-format statement,
by which the terms at level $n$ of the summand indexed by the history $h$ (the branch $h$) are
evaluated on $U^h_n$.
\end{itemize}
A real configuration $V$ is in the \emph{level-$n$ window on $N_i$} if
$|V(\partial P)-1|\le\varepsilon^{\mathrm{win}}:=C_rL^2\varepsilon_n$ for every level-$n$
plaquette $P\subset N_i$ \cite{B-CMP119}. Put
\begin{equation*}
\varepsilon':=C_w\varepsilon_n,\qquad C_w:=\sqrt{2}\,C_rL^2,\qquad
C_{\mathfrak{e}}:=C_wB_3\max\Bigl\{2,\frac{C_{\mathfrak{a}}M_Q}{c''c_T}\Bigr\}.
\end{equation*}
Here $B_3$ and $C_r$ are constants of the regularity statement of the two-point witness
theorem, and $C_{\mathfrak{a}}$, $c''$, $c_T$, $M_Q$ are the constants of this section entering
$\mathfrak{a}_n=C_{\mathfrak{a}}\alpha_{*,n}$, $\mathfrak{r}_n=c_T\alpha_{*,n}$ and $\epsilon$.
Let $\Omega_i$ be the union of the level-$n$ $M_1$-cubes that meet
$S_i^{+3\rho/16+2M_1}$, and define
\begin{equation}\label{8.5:eq-readings-shortened-1}
\operatorname{Irr}_i(V):=
\begin{cases}
\operatorname{Irr}^{\mathrm{sh}}[f_i]\bigl(U_{n,\Omega_i}(Q_nV)\bigr),
& V\text{ in the level-$n$ window on }N_i,\\
0, & \text{otherwise.}
\end{cases}
\end{equation}
By its definition, $\operatorname{Irr}_i$ is a function of the restriction of $V$ to the bonds
of $\Omega_i$, gauge invariant, equal to $0$ at $V=\mathbf{1}$, and no history enters it.
Then the following hold.
\begin{itemize}
\item[(a)] $\mathfrak{e}_\star=C_{\mathfrak{e}}\varepsilon_n$.
\item[(b)] $\Omega_i^{\sim-2}\supset\mathsf{C}_i^{\sharp}$ and
$\Omega_i\subset S_i^{+3\rho/16+3M_1}\subset N_i$. For $V$ in the level-$n$ window on $N_i$,
the restriction of $U_{n,\Omega_i}(Q_nV)$ to the bare bonds over $\mathsf{C}_i^{\sharp}$
belongs to $\mathfrak{R}(S_i;\varepsilon')$; thus, for $V$ in the level-$n$ window on $N_i$,
$\operatorname{Irr}_i(V)$ is a value of $\operatorname{Irr}^{\mathrm{sh}}[f_i]$ on
$\mathfrak{R}(S_i;\varepsilon')$, and off the window $\operatorname{Irr}_i(V)=0$.
\item[(c)] Let $(h,V)\notin\mathcal{E}_i$, and let $x_n\ge4c_2$ and
$\log(x_n/2)>2\mathfrak{p}_0$. For every value of the variables that $\mathbf{T}_h$ still
integrates, the restriction of $U^h_n$ to the bare bonds over $\mathsf{C}_i^{\sharp}$ belongs
to $\mathfrak{R}(S_i;\varepsilon')$, and
\begin{equation}\label{8.5:eq-readings-shortened-2}
\begin{split}
&\partial_{z_i}\bigl[\mathbf{E}^w+\mathbf{D}^w+\mathbf{R}^w\bigr]_h\big|_0(U^h_n)
-\partial_{z_i}\bigl[\mathbf{E}^w+\mathbf{D}^w+\mathbf{R}^w\bigr]_h\big|_0(\mathbf{1})\\
&\qquad=\operatorname{Irr}^{\mathrm{sh}}[f_i](U^h_n)+\mathcal{L}^h_i,
\end{split}
\end{equation}
where $\mathcal{L}^h_i$ is the sum of those long constituents of
Lemma~\ref{8.5:lem-long-constituents} (the constituents of
\eqref{8.5:eq-response-decomposition-a} outside the core groups whose polymer is not in
$\mathcal{K}_j$) whose polymer $X$ satisfies $X\subset\Lambda_j(h)$; and, under the hypotheses
of Lemma~\ref{8.5:lem-long-constituents} on $f_i$ and $\rho$,
$|\mathcal{L}^h_i|\le\mathfrak{f}^{lg}[f_i]$. We
write $\Delta^h_i$ for the left member of \eqref{8.5:eq-readings-shortened-2}.
\item[(d)] Wherever the global minimiser $U_n(V)$ exists and its data satisfy the smallness
condition of Definition~\ref{8.5:def-critical-backgrounds} with $\varepsilon_1\le\varepsilon'$,
its restriction to the bare bonds over $\mathsf{C}_i^{\sharp}$ belongs to
$\mathfrak{R}(S_i;\varepsilon')$, and the left member of
\eqref{8.5:eq-response-decomposition-a} for $f=f_i$, evaluated at $U_n(V)$, equals
$\operatorname{Irr}^{\mathrm{sh}}[f_i](U_n(V))$ plus the sum of all long constituents.
If $C_{\mathrm{ext}}$ is a constant and $V^{\natural}$ a real configuration with
$\sup_p|V^{\natural}(\partial p)-1|\le C_{\mathrm{ext}}\varepsilon^{\mathrm{win}}$, the same
holds for $U_n(V^{\natural})$ in place of $U_n(V)$ when $C_w$ is replaced by
$C_{\mathrm{ext}}C_rL^2$, with $\varepsilon'$ and the kept set read at this value of $C_w$.
\end{itemize}
On the triples $(h,V,\mathsf{v})$ we write $\tilde{\nu}$ for the law on $(h,V)$ of the
two-point witness theorem and, with $\mathbb{E}_h$ the probability measure of (C), we put
\begin{equation*}
\tilde{\mu}(dh,dV,d\mathsf{v}):=\tilde{\nu}(dh,dV)\,\mathbb{E}_h(d\mathsf{v}),
\end{equation*}
the normalised integrand at $z=0$ itself, and, with $\mathbf{1}_{\mathcal{E}_i^{c}}$ the
indicator of the complement of $\mathcal{E}_i$,
\begin{align*}
\mathsf{i}^h_i&:=\Delta^h_i\,\mathbf{1}_{\mathcal{E}_i^{c}}(h,V),\\
\mathsf{o}^h_i&:=N\sum_p\tilde{f}_i(x_p)\,a_p(U^h_n)\,\mathbf{1}_{\mathcal{E}_i^{c}}(h,V),
\qquad
\mathbb{I}^h_i(V):=\mathbb{E}_h\bigl[\mathsf{i}^h_i\bigr],
\end{align*}
where $N\in[\tfrac23,\tfrac43]$ is the normalisation and $\tilde{f}_i$ are the corrected test
functions of the two-point witness theorem, and
\begin{equation*}
a_p(U)=\tfrac12\operatorname{tr}\bigl[(\mathbf{1}-U(\partial p))
(\mathbf{1}-U(\partial p))^{*}\bigr]=1-\operatorname{Re}\operatorname{tr}U(\partial p).
\end{equation*}
\end{proposition}

\begin{proof}
\emph{The properties listed after} \eqref{8.5:eq-readings-shortened-1}. By (A), the
configuration $U_{n,\Omega_i}(Q_nV)$ entering \eqref{8.5:eq-readings-shortened-1} is a function
of the restriction of $V$ to the bonds of $\Omega_i$, and no history enters
\eqref{8.5:eq-readings-shortened-1}. A gauge transformation of $V$ replaces each
$V(\partial P)$ by a conjugate, which leaves $|V(\partial P)-1|$, and hence the window
condition, unchanged; the minimal configuration of the transformed data is the transform of the
minimal configuration, and the constituents of $\operatorname{Irr}^{\mathrm{sh}}[f_i]$ are gauge
invariant (Definition~\ref{8.5:def-kept-set}); so $\operatorname{Irr}_i$ is gauge invariant. At
$V=\mathbf{1}$ the window condition holds, $Q_n\mathbf{1}=\mathbf{1}$, the minimal configuration
for the data $\mathbf{1}$ is $\mathbf{1}$, and every constituent of \eqref{8.5:eq-kept-set-c}
vanishes at $\mathbf{1}$ (Definition~\ref{8.5:def-kept-set}); so
$\operatorname{Irr}_i(\mathbf{1})=0$.

\emph{Part} (a). By the definitions of $\mathfrak{e}$, $\mathfrak{e}_Q$ and $\epsilon$ in this
section, $\mathfrak{e}=2B_3\varepsilon'$,
$\mathfrak{e}_Q=\epsilon\,\mathfrak{a}_n$ with $\epsilon=B_3M_Q\varepsilon'/(c''\mathfrak{r}_n)$,
$\mathfrak{r}_n=c_T\alpha_{*,n}$ and $\mathfrak{a}_n=C_{\mathfrak{a}}\alpha_{*,n}$, so that
$\mathfrak{e}_Q=B_3\varepsilon'C_{\mathfrak{a}}M_Q/(c''c_T)$. Hence
\begin{equation*}
\mathfrak{e}_\star=\max\{\mathfrak{e},\mathfrak{e}_Q\}
=B_3\varepsilon'\max\Bigl\{2,\frac{C_{\mathfrak{a}}M_Q}{c''c_T}\Bigr\}
=C_{\mathfrak{e}}\varepsilon_n .
\end{equation*}

\emph{Part} (b), \emph{geometry}. By the definition of $\Omega_i$ as the union of the level-$n$
$M_1$-cubes that meet $S_i^{+3\rho/16+2M_1}$, one has
$\Omega_i^{\sim-2}\supset\mathsf{C}_i^{\sharp}=S_i^{+3\rho/16}$ and
$\Omega_i\subset S_i^{+3\rho/16+3M_1}$. Finally $3M_1\le\rho/16$, by the floor
\eqref{8.5:eq-kept-set-a} together with $M_1<M$ \cite{B-CMP119}, whence
$S_i^{+3\rho/16+3M_1}\subset S_i^{+\rho/4}=N_i$.

\emph{Part} (b), \emph{membership}. Put $\mathfrak{B}':=\mathbf{B}_n(\Omega_i)$
\cite[(2.13)]{B-CMP119} and $V':=Q_nV$; by (A), $U_{n,\Omega_i}(Q_nV)$ is the minimal
configuration $U_{\mathfrak{B}'}(V')$. The level-$n$ domain $\Omega_n'$ of $\mathfrak{B}'$
satisfies, by (A), $\Omega_n'\supset\Omega_i^{\sim-2}\supset\mathsf{C}_i^{\sharp}$, which is
condition $(\mathfrak{R}1)$ of Definition~\ref{8.5:def-critical-backgrounds}. By the plaquette
part of (B), the plaquette variables of $V'$ are $\mathbf{1}$ or conjugates of some
$V(\partial P)$ with $P\subset\Omega_i$; since conjugation does not change $|\cdot-1|$, and
since $\Omega_i\subset N_i$ and $V$ is in the window,
\begin{equation*}
\varepsilon_1\le\sup_{P\subset\Omega_i}|V(\partial P)-1|\le\varepsilon^{\mathrm{win}}
=C_rL^2\varepsilon_n\le\sqrt{2}\,C_rL^2\varepsilon_n=\varepsilon'.
\end{equation*}
This is condition $(\mathfrak{R}2)$. Existence of $U_{n,\Omega_i}(Q_nV)$, with its tower in
$\mathfrak{U}_k(\{\Omega_m\},2B_3\alpha)$, is the existence part of (B); these are the
hypotheses under which \cite[(2.16)--(2.17)]{B-CMP119} are applied. Hence the restriction lies in
$\mathfrak{R}(S_i;\varepsilon')$.

\emph{Part} (c), \emph{condition $(\mathfrak{R}1)$}. Since $(h,V)\notin\mathcal{E}_i$, no live
label of $h$ meets $N_i\supset\mathsf{C}_i^{\sharp}$; by the nesting statement of (B), $N_i$
lies in $\Omega_n(h)$, so over $\mathsf{C}_i^{\sharp}$ the level of $\mathfrak{B}_h$ is $n$.

\emph{Part} (c), \emph{condition $(\mathfrak{R}2)$}. By the regularity statement of (C), the
data of $U^h_n$ satisfy the smallness condition with
$\varepsilon_1\le C_rL^2\max_{j\le n}\varepsilon_j$.
We show $\varepsilon_j\le\sqrt{2}\,\varepsilon_n$ for $j\le n$. By the comparison of couplings
in (D), $x_j\ge x_n-2c_2$, that is,
$g_j^{-2}\ge g_n^{-2}(1-2c_2g_n^2)$. As $x_n\ge4c_2$ gives $2c_2g_n^2\le\frac12$,
\begin{equation*}
g_j\le(1-2c_2g_n^2)^{-1/2}g_n\le\sqrt{2}\,g_n .
\end{equation*}
Let $\varepsilon(g):=gA_0(\log g^{-2})^{\mathfrak{p}_0}$, so that
$\varepsilon_k=\varepsilon(g_k)$. By the equivalence in (D) with $p=\mathfrak{p}_0$,
$\varepsilon$ is increasing on the interval where
$\log g^{-2}>2\mathfrak{p}_0$. That interval contains $\sqrt{2}\,g_n$, because
$\log\bigl((\sqrt{2}g_n)^{-2}\bigr)=\log(x_n/2)>2\mathfrak{p}_0$, and with it $g_j$. Hence
\begin{align*}
\varepsilon_j=\varepsilon(g_j)\le\varepsilon(\sqrt{2}\,g_n)
&=\sqrt{2}\,g_nA_0\bigl(\log(x_n/2)\bigr)^{\mathfrak{p}_0}\\
&\le\sqrt{2}\,g_nA_0(\log x_n)^{\mathfrak{p}_0}=\sqrt{2}\,\varepsilon_n .
\end{align*}
So $\varepsilon_1\le\sqrt{2}\,C_rL^2\varepsilon_n=\varepsilon'$, which is
$(\mathfrak{R}2)$, and the restriction of $U^h_n$ lies in $\mathfrak{R}(S_i;\varepsilon')$.

\emph{Part} (c), \emph{the identity}. On $N_i$ the functions $\varphi_j$ equal $1$ by (D). By
the nesting statement of (B), $N_i\subset\Lambda_n$; by $\Lambda_j^c\subset\Lambda_n^c$ from
(D), $\Lambda_n\subset\Lambda_j$ for $j\le n$. Moreover
$S_i^+\subset\Lambda_j^{\circledR}$: with $t=L^{j-n}$, one layer of
$M\check{R}_j$-cubes has width $M\check{R}_jt\le\rho/8$ in level-$n$ units, by (C). By the
whole-lattice coincidence of (D), the constituents of
$\partial_{z_i}[\mathbf{E}^w+\mathbf{D}^w+\mathbf{R}^w]_h|_0$ that belong to the core groups
$\mathfrak{i}^c_y[\dot{S}^{(j)}]$ of Definition~\ref{8.5:def-kept-set} or whose polymer $X$
lies in $\mathcal{K}_j$ are those of \eqref{8.5:eq-response-decomposition-a}
(Lemma~\ref{8.5:lem-response-decomposition}). By \eqref{8.5:eq-kept-set-c}
(Definition~\ref{8.5:def-kept-set}) these constituents add up to
$\operatorname{Irr}^{\mathrm{sh}}[f_i]$; the remaining ones are the long constituents with
$X\subset\Lambda_j(h)$. The sourced-format statement of (D) evaluates all these terms on
$U^h_n$, which gives \eqref{8.5:eq-readings-shortened-2}. Under the hypotheses of
Lemma~\ref{8.5:lem-long-constituents} on $f_i$ and $\rho$, the bound
$\mathfrak{f}^{lg}[f_i]$ of \eqref{8.5:eq-long-constituents-a}
(Lemma~\ref{8.5:lem-long-constituents}) is a sum of termwise bounds on the long constituents,
so it bounds each of their sub-sums, in particular $|\mathcal{L}^h_i|$.

\emph{Part} (d). For the global minimiser the level-$n$ domain is the whole torus, so
$(\mathfrak{R}1)$ holds; $(\mathfrak{R}2)$ is the hypothesis. By
Lemma~\ref{8.5:lem-response-decomposition}, the left member of
\eqref{8.5:eq-response-decomposition-a} for $f=f_i$ at $U_n(V)$ is the sum of the constituents
displayed there; by \eqref{8.5:eq-kept-set-c} (Definition~\ref{8.5:def-kept-set}) those that
belong to the core groups or whose polymer lies in $\mathcal{K}_j$ add up to
$\operatorname{Irr}^{\mathrm{sh}}[f_i](U_n(V))$, and the remaining ones are all the long
constituents. The argument uses no property of $V$ or of the value of $C_w$; with $C_w$
replaced by $C_{\mathrm{ext}}C_rL^2$, for which
$\varepsilon'=C_{\mathrm{ext}}C_rL^2\varepsilon_n=C_{\mathrm{ext}}\varepsilon^{\mathrm{win}}$, it
applies verbatim to $U_n(V^{\natural})$, which gives the last sentence of (d).
\end{proof}

\begin{proposition}[One sup bound for every reading]\label{8.5:prop-reading-sup}
Work in the setting, and under the ceilings and the floor, of
Theorem~\ref{8.5:thm-shortened-bound}, and under the ceiling \eqref{8.5:eq-long-constituents-a}
of Lemma~\ref{8.5:lem-long-constituents}. Let $f_i$, $S_i$, $S_i^+$ and the function
$\operatorname{Irr}_i$ be those of Proposition~\ref{8.5:prop-readings-shortened}. Let
$\operatorname{Irr}_i^{\natural}$ and $\operatorname{Irr}_i^{\mathrm{gl}}$ be the functions of the
second and third readings of the irrelevant response of $f_i$ in
Proposition~\ref{8.5:prop-readings-shortened}, and $\mathsf{i}^h_i$ the response of the
branch $h$ defined there. Let $\mathbb{I}^h_i$ denote the $\mathbb{E}_h$-mean of $\mathsf{i}^h_i$. Put
\begin{equation}\label{8.5:eq-reading-sup-a}
\mathsf{I}_i:=2\bigl[C^{\flat}\,\hat{\theta}_{n-1}\,\mathfrak{e}_\star^{2}
+C_D\,r^{-1}\mathfrak{m}_n\bigr]\,|S_i^+|\,\|f_i\|_\sharp .
\end{equation}
Then
\begin{equation}\label{8.5:eq-reading-sup-1}
\sup|\operatorname{Irr}_i|\le\tfrac12\,\mathsf{I}_i ,
\end{equation}
and each of the four quantities
$\sup|\mathsf{i}^h_i|$, $\sup|\mathbb{I}^h_i|$, $\sup|\operatorname{Irr}_i^{\natural}|$ and
$\sup|\operatorname{Irr}_i^{\mathrm{gl}}|$ is at most $\mathsf{I}_i$.
\end{proposition}

\begin{proof}
\emph{The reading through $\Omega_i$.} By Proposition~\ref{8.5:prop-readings-shortened},
$\operatorname{Irr}_i(V)$ vanishes for $V$ outside the window on $N_i$. For $V$ in that window,
\[
\operatorname{Irr}_i(V)=\operatorname{Irr}^{\mathrm{sh}}[f_i]\bigl(U_{n,\Omega_i}(Q_nV)\bigr),
\]
and by the same proposition the configuration $U_{n,\Omega_i}(Q_nV)$ belongs to
$\mathfrak{R}(S_i;\varepsilon')$. Theorem~\ref{8.5:thm-shortened-bound} holds for every element of
$\mathfrak{R}(S_i;\varepsilon')$. Applied with $f=f_i$ and $S=S_i$, it gives
\[
|\operatorname{Irr}_i(V)|\le\bigl[C^{\flat}\hat{\theta}_{n-1}\mathfrak{e}_\star^2
+C_Dr^{-1}\mathfrak{m}_n\bigr]|S_i^+|\,\|f_i\|_\sharp=\tfrac12\,\mathsf{I}_i .
\]
The right-hand side does not depend on $V$, and the value $0$ taken off the window obeys the same
bound. This proves \eqref{8.5:eq-reading-sup-1}.

\emph{The branch responses and the other two readings.} By
Proposition~\ref{8.5:prop-readings-shortened}, each of $\mathsf{i}^h_i$,
$\operatorname{Irr}_i^{\natural}$ and $\operatorname{Irr}_i^{\mathrm{gl}}$ is the sum of two parts:
a kept part, which is the shortened functional $\operatorname{Irr}^{\mathrm{sh}}[f_i]$ evaluated at a
configuration in $\mathfrak{R}(S_i;\varepsilon')$, and the long constituents outside the kept set of
Lemma~\ref{8.5:lem-long-constituents}.

For the kept part, Theorem~\ref{8.5:thm-shortened-bound} gives the bound $\tfrac12\mathsf{I}_i$,
exactly as in the first paragraph.

For the long constituents, Lemma~\ref{8.5:lem-long-constituents} bounds them in absolute value by
$\mathfrak{f}^{lg}[f_i]$. Under the ceiling \eqref{8.5:eq-long-constituents-a},
\[
\mathfrak{f}^{lg}[f_i]\le
\bigl[C^{\flat}\hat{\theta}_{n-1}\mathfrak{e}_\star^2+C_Dr^{-1}\mathfrak{m}_n\bigr]|S_i^+|\,\|f_i\|_\sharp
=\tfrac12\,\mathsf{I}_i .
\]

Adding the two contributions gives
$\sup|\mathsf{i}^h_i|\le\mathsf{I}_i$, $\sup|\operatorname{Irr}_i^{\natural}|\le\mathsf{I}_i$ and
$\sup|\operatorname{Irr}_i^{\mathrm{gl}}|\le\mathsf{I}_i$.

\emph{The mean.} In the construction of the two-point witness theorem, $\mathbb{E}_h$ is the
expectation of a probability measure. Positivity
alone therefore gives $|\mathbb{E}_h[X]|\le\sup|X|$ for every bounded $X$. Taking
$X=\mathsf{i}^h_i$ yields
\[
|\mathbb{I}^h_i|=|\mathbb{E}_h[\mathsf{i}^h_i]|\le\sup|\mathsf{i}^h_i|\le\mathsf{I}_i ,
\]
and so $\sup|\mathbb{I}^h_i|\le\mathsf{I}_i$.

No comparison between the different readings enters the proof: each reading is bounded directly by
\eqref{8.5:eq-reading-sup-a}.
\end{proof}

\begin{lemma}[Sup bound on the main curvature insertion]\label{8.5:lem-insertion-sup}
Let $i\in\{1,2\}$. Let $n=k_0$ be the witness level, let $S_i$, $S_i^+$, $\rho$ and $f_i$ be as in the
two-point witness theorem, let $\tilde{f}_i$ be the modified test function of that theorem, of which
only the properties (b) below are used, and let $\mathsf{C}_i=S_i^{+\rho/8}$ and
$\mathsf{C}_i^{\sharp}=S_i^{+3\rho/16}$.
Let $\varepsilon'=C_w\varepsilon_n$ and assume $\varepsilon'\le a_1$, where $a_1$ is the constant of
Definition~\ref{8.5:def-critical-backgrounds}, so that the class $\mathfrak{R}(S_i;\varepsilon')$ of
that definition is defined.
For an $\mathrm{SU}(N)$-valued configuration $U$ given on the bare bonds of every bare plaquette $p$
with $x(p)\in\operatorname{supp}\tilde{f}_i$ put
\begin{equation*}
\mathcal{O}[\tilde{f}_i](U):=\sum_{p:\,x(p)\in\operatorname{supp}\tilde{f}_i}\tilde{f}_i(x(p))\,a_p(U),
\end{equation*}
the sum over bare plaquettes, with $a_p(U)=1-\operatorname{Re}\operatorname{tr}U(\partial p)$ as in (a)
below. Assume the following statements of the construction of the two-point
witness theorem.
\begin{itemize}
\item[(a)] (Plaquette sums on regular configurations.) For every level $k$, every admissible sequence
$\{\Omega_j\}_{j\le k}$, every $\varepsilon$ for which Ba{\l}aban's regular space
$\mathfrak{U}_k(\{\Omega_j\},\varepsilon)$ is defined, every $U$ in that space and every cell
$\Delta'\subset\Omega_k$ of the level-$k$ lattice (an own-scale cell),
\begin{equation*}
\sum_{p\subset\Delta'}a_p(U)\le C_{\square}\varepsilon^2,\qquad C_{\square}:=6C_a,
\end{equation*}
where
\begin{equation*}
a_p(U)=\tfrac12\operatorname{tr}\bigl[(1-U(\partial p))(1-U(\partial p))^*\bigr]
=1-\operatorname{Re}\operatorname{tr}U(\partial p)
\end{equation*}
(since $\operatorname{tr}1=1$), the plaquette density of the smeared curvature field of
Definition~\ref{1.3:def-curvature-field}, and $a_p(U)\le C_a|U(\partial p)-1|^2$.
\item[(b)] (The modified test function.) $\operatorname{supp}\tilde{f}_i\subset S_i^+$, and
$|\tilde{f}_i-f_i|\le \tfrac32\,\mathfrak{d}_W C_hM\|f_i\|_\infty/(\vartheta\rho)$, where
$\mathfrak{d}_W$ and $C_h$ are constants of the two-point witness theorem; moreover
$\mathfrak{d}_W\le\frac13$ and $\rho\ge C_hM/\vartheta$.
\item[(c)] (The normalisation.) The normalisation constant $N$ of the witness construction (a
positive number; not the $N$ of the gauge group) satisfies $N\in[\frac23,\frac43]$.
\item[(d)] (The local insertion.)
$X^{(k_0)}(\tilde{f}_i)=\mathcal{O}[\tilde{f}_i](U_{n,\Omega_{\tilde{f}_i}}(Q_nV))$
for $V$ in the window on $N_i$ of
Proposition~\ref{8.5:prop-readings-shortened} and $X^{(k_0)}(\tilde{f}_i)=0$
otherwise, where $U_{n,\Omega}(Q_nV)$ is the minimiser of \cite[(2.12)]{B-CMP119} taken with the
determining set $\mathbf{B}_n(\Omega)$ of \cite[(2.13)]{B-CMP119}, and $\Omega_{\tilde{f}_i}$ is a
union of level-$n$ $M_1$-cubes containing a neighbourhood of $\operatorname{supp}\tilde{f}_i$, given
by the two-radius device used in Proposition~\ref{8.5:prop-readings-shortened}; this minimiser lies in
$\mathfrak{U}_n(\{\Omega'_m\}_{m\le n},B_3\varepsilon^{\mathrm{win}})$, where $B_3$ is the constant of
Lemma~\ref{8.5:lem-background-regularity}, for an admissible sequence
$\{\Omega'_m\}_{m\le n}$ and has level $n$ over $\operatorname{supp}\tilde{f}_i$.
\end{itemize}
Assume moreover $\varepsilon^{\mathrm{win}}\le\varepsilon'=C_w\varepsilon_n$.
For $\mathfrak{b}\in\mathfrak{R}(S_i;\varepsilon')$, a configuration on the bare bonds over
$\mathsf{C}_i^{\sharp}\supset S_i^+\supset\operatorname{supp}\tilde{f}_i$, the value
$\mathcal{O}[\tilde{f}_i](\mathfrak{b})$ is defined, and
\begin{equation}\label{8.5:eq-insertion-sup-1}
|\mathcal{O}[\tilde{f}_i](\mathfrak{b})|\le \sup|\tilde{f}_i|\cdot|S_i^+|\cdot
C_{\square}(B_3\varepsilon')^2,
\end{equation}
where $|S_i^+|$ denotes the number of own-scale cells of level $n$ meeting $S_i^+$; and
\begin{equation}\label{8.5:eq-insertion-sup-a}
\sup_{\mathfrak{R}(S_i;\varepsilon')}|N\mathcal{O}[\tilde{f}_i]|\le\mathsf{O}_i
:=C_{\mathcal{O}}\,\varepsilon_n^2\,|S_i^+|\,\|f_i\|_\infty,
\qquad C_{\mathcal{O}}:=\tfrac83\,C_{\square}B_3^2C_w^2 .
\end{equation}
In particular $\sup|\mathsf{o}^h_i|\le\mathsf{O}_i$ for the insertion $\mathsf{o}^h_i$ of every
branch $h$ of the two-point witness theorem, this insertion being $N\mathcal{O}[\tilde{f}_i]$
evaluated on a configuration of $\mathfrak{R}(S_i;\varepsilon')$. The local insertion obeys the same
bound: $\sup_V|NX^{(k_0)}(\tilde{f}_i)|\le\mathsf{O}_i$.
\end{lemma}

\begin{proof}
Fix $\mathfrak{b}\in\mathfrak{R}(S_i;\varepsilon')$. Each $a_p(\mathfrak{b})$ is one half of the
trace of the positive semidefinite matrix
$(1-\mathfrak{b}(\partial p))(1-\mathfrak{b}(\partial p))^*$, hence non-negative.
Only plaquettes with $x(p)\in\operatorname{supp}\tilde{f}_i\subset S_i^+$ contribute, by (b), so
\begin{equation*}
|\mathcal{O}[\tilde{f}_i](\mathfrak{b})|\le\sup|\tilde{f}_i|\sum_{p:\,x(p)\in S_i^+}a_p(\mathfrak{b}).
\end{equation*}
Since $S_i^+\subset\mathsf{C}_i\subset\mathsf{C}_i^{\sharp}$, these plaquettes lie over
$\mathsf{C}_i^{\sharp}$. By Definition~\ref{8.5:def-critical-backgrounds}, over $\mathsf{C}_i^{\sharp}$
the configuration $\mathfrak{b}$ is the restriction of a minimal configuration $U_{\mathfrak{B}'}(V')$
of \cite[Theorem~1]{B-CMP102b} over an admissible sequence $\{\Omega'_m\}_{m\le n}$ with
$\mathsf{C}_i^{\sharp}\subset\Omega'_n$, so that over $\mathsf{C}_i^{\sharp}$ the level is $n$, and
with $V'$ subject to the smallness condition of that definition; as in the proof of
Lemma~\ref{8.5:lem-background-regularity}, this gives
$U_{\mathfrak{B}'}(V')\in\mathfrak{U}_n(\{\Omega'_m\}_{m\le n},B_3\varepsilon')$.
Hypothesis (a), applied with $k=n$, the
sequence $\{\Omega'_m\}$ and $\varepsilon=B_3\varepsilon'$, bounds the sum over each own-scale cell of
level $n$ in $\mathsf{C}_i^{\sharp}$ by $C_{\square}(B_3\varepsilon')^2$; summing over the
$|S_i^+|$ own-scale cells of level $n$ meeting $S_i^+$ gives \eqref{8.5:eq-insertion-sup-1}.

By (b), using $\mathfrak{d}_W\le\frac13$ and $C_hM/(\vartheta\rho)\le1$,
\begin{equation*}
|\tilde{f}_i|\le|f_i|+|\tilde{f}_i-f_i|
\le\|f_i\|_\infty+\tfrac32\cdot\tfrac13\,\|f_i\|_\infty
\le2\|f_i\|_\infty .
\end{equation*}
Inserting this, $N\le\frac43$ from (c) and $\varepsilon'=C_w\varepsilon_n$ into
\eqref{8.5:eq-insertion-sup-1},
\begin{equation*}
|N\mathcal{O}[\tilde{f}_i](\mathfrak{b})|
\le\tfrac43\cdot2\|f_i\|_\infty\,|S_i^+|\,C_{\square}B_3^2C_w^2\varepsilon_n^2
=\mathsf{O}_i ,
\end{equation*}
and taking the supremum over $\mathfrak{b}$ gives \eqref{8.5:eq-insertion-sup-a}. The assertion
on $\mathsf{o}^h_i$ is \eqref{8.5:eq-insertion-sup-a} read for the insertion of the branch $h$,
which is $N\mathcal{O}[\tilde{f}_i]$ evaluated on a configuration of $\mathfrak{R}(S_i;\varepsilon')$.

For the local insertion, off the window $X^{(k_0)}(\tilde{f}_i)=0$ by (d). In the window, by (d) the
minimiser $U_{n,\Omega_{\tilde{f}_i}}(Q_nV)$ lies in
$\mathfrak{U}_n(\{\Omega'_m\}_{m\le n},B_3\varepsilon^{\mathrm{win}})$
and has level $n$ over $\operatorname{supp}\tilde{f}_i$. Hence hypothesis (a) applies to it on the
own-scale cells covering $\operatorname{supp}\tilde{f}_i$, and the two estimates above, with
$B_3\varepsilon^{\mathrm{win}}$ in the place of $B_3\varepsilon'$ and
$(\varepsilon^{\mathrm{win}})^2\le(\varepsilon')^2=C_w^2\varepsilon_n^2$ from the assumption
$\varepsilon^{\mathrm{win}}\le\varepsilon'$, give the same bound
$|NX^{(k_0)}(\tilde{f}_i)|\le\mathsf{O}_i$ for every $V$.
\end{proof}

Throughout this subsection Notation~\ref{8.5:not-extraction-notation} is in force. We write
$n:=k_0:=K-s$ for the witness level and $x:=x_n=g_n^{-2}$.

\emph{Two elementary facts.} Let $(\mathcal{X},\mu)$ be a probability space, and let $a,b$ be
bounded real random variables on it. Then
$\operatorname{Cov}_\mu(a,b)=E_\mu[a\,(b-E_\mu b)]$. Hence
$|\operatorname{Cov}_\mu(a,b)|\le\sup|a|\,E_\mu|b-E_\mu b|$. Moreover
$E_\mu|b-E_\mu b|\le E_\mu|b|+|E_\mu b|\le 2\sup|b|$.

For the second fact, let $W$ be a parameter. Suppose that for each value of $W$ the map
$F\mapsto F^{\langle\rangle}(W)$ is a composition of positive normalised linear maps. Then it is
integration against a probability measure $\mu_W$. The expression
\begin{equation*}
T^{\langle\rangle}[F_1,F_2](W):=(F_1F_2)^{\langle\rangle}(W)
-F_1^{\langle\rangle}(W)\,F_2^{\langle\rangle}(W)
\end{equation*}
is therefore the covariance $\operatorname{Cov}_{\mu_W}(F_1,F_2)$, and the first fact applies to
it. We obtain
\begin{equation}\label{8.5:eq-averaged-extraction-b}
\begin{gathered}
|\operatorname{Cov}_\mu(a,b)|\le 2\sup|a|\,\sup|b|,\\
|T^{\langle\rangle}[F_1,F_2](W)|\le 2\sup|F_1|\,\sup|F_2|.
\end{gathered}
\end{equation}

\emph{Sizes along the envelope.} We impose the following ceiling on $g$:
\begin{equation}\label{8.5:eq-averaged-extraction-c}
\sup_{y\ge X-c_5}\Bigl(\frac{\log y}{\log(y-2c_2)}\Bigr)^{q-\mathfrak{p}_0}\le 2
\qquad\text{and}\qquad X-c_5\ge 8c_2 .
\end{equation}
By the envelope of the two-point witness theorem, $x_n\ge\underline{x}_s\ge X-c_5$. For $i\le n$,
the comparison of the running inverse couplings at different levels gives $x_i\ge x_n-2c_2$.

The function $\theta(\cdot)$ is nondecreasing in the coupling. Hence for $i\le n-1$ we have
$\theta_i=C_\theta(\log x_i)^{\mathfrak{p}_0-q}\le C_\theta(\log(x-2c_2))^{\mathfrak{p}_0-q}$.
Applying \eqref{8.5:eq-averaged-extraction-c} at $y=x$ gives the first bound below.

Recall $\mathfrak{m}_n=C_{\mathfrak{m}}y_n^{-(\kappa_0-2)/2}$ with $y_n=\min_{i\le n}x_i-2c_2$.
Then $y_n\ge x_n-4c_2$, and $x_n-4c_2\ge x/2$ because $x\ge X-c_5\ge 8c_2$. This gives the second
bound below:
\begin{equation}\label{8.5:eq-averaged-extraction-1}
\hat{\theta}_{n-1}\le C_\theta(\log(x-2c_2))^{\mathfrak{p}_0-q}\le 2\theta_n,
\qquad
\mathfrak{m}_n\le C_{\mathfrak{m}}2^{(\kappa_0-2)/2}x^{-(\kappa_0-2)/2}.
\end{equation}
Moreover $\varepsilon_n^2=A_0^2(\log x)^{2\mathfrak{p}_0}/x$ and
$\mathfrak{e}_\star^2=C_{\mathfrak{e}}^2\varepsilon_n^2$.

By the class of the test functions, $|S_i^+|\le 80\rho^4$ and
$\mathfrak{n}(f_i)\ge\vartheta\|f_i\|_\infty\rho^4$. Hence
$|S_i^+|\,\|f_i\|_\infty\le 80\rho^4\|f_i\|_\infty\le 80\,\mathfrak{n}(f_i)/\vartheta$. By the
floor $\rho\ge 40M/\vartheta$ of the two-point witness theorem, also
$\|f_i\|_\sharp=\|f_i\|_\infty$.

We put
\begin{equation}\label{8.5:eq-averaged-extraction-d}
\begin{aligned}
\mathsf{i}(x)&:=4C^{\flat}C_\theta C_{\mathfrak{e}}^2A_0^2(\log x)^{3\mathfrak{p}_0-q}
+2C_DC_{\mathfrak{m}}2^{(\kappa_0-2)/2}r^{-1}x^{-(\kappa_0-4)/2},\\
\mathsf{o}(x)&:=C_{\mathcal{O}}A_0^2(\log x)^{2\mathfrak{p}_0},\\
\mathsf{E}^{\flat}(x)&:=12800\,\vartheta^{-2}\bigl[2\,\mathsf{i}(x)\,\mathsf{o}(x)
+\mathsf{i}(x)^2\bigr].
\end{aligned}
\end{equation}
Insert \eqref{8.5:eq-averaged-extraction-1} and the value of $\varepsilon_n^2$ into
\eqref{8.5:eq-reading-sup-a} and \eqref{8.5:eq-insertion-sup-a}. This gives
\begin{equation}\label{8.5:eq-averaged-extraction-2}
\mathsf{I}_i\le x^{-1}\mathsf{i}(x)\,|S_i^+|\,\|f_i\|_\sharp,
\qquad
\mathsf{O}_i=x^{-1}\mathsf{o}(x)\,|S_i^+|\,\|f_i\|_\infty .
\end{equation}
Indeed,
\begin{align*}
&2C^{\flat}\cdot 2C_\theta(\log x)^{\mathfrak{p}_0-q}\cdot C_{\mathfrak{e}}^2A_0^2
(\log x)^{2\mathfrak{p}_0}x^{-1}
=x^{-1}\cdot 4C^{\flat}C_\theta C_{\mathfrak{e}}^2A_0^2(\log x)^{3\mathfrak{p}_0-q},\\
&2C_Dr^{-1}C_{\mathfrak{m}}2^{(\kappa_0-2)/2}x^{-(\kappa_0-2)/2}
=x^{-1}\cdot 2C_DC_{\mathfrak{m}}2^{(\kappa_0-2)/2}r^{-1}x^{-(\kappa_0-4)/2},
\end{align*}
and the right-hand sides are $x^{-1}$ times the two terms of $\mathsf{i}(x)$.

\emph{The floor on $q$.} We impose
\begin{equation}\label{8.5:eq-averaged-extraction-e}
q\ge 21\,\mathfrak{p}_0+1 .
\end{equation}
This is a lower bound on Ba{\l}aban's integers $q_0=q_1=q$ of \cite[(2.28)]{B-CMP119}, which are
stated there to be integers greater than $1$.

The exponent $\mathfrak{p}_0$ is the fixed exponent $p_0$ of the degree function $p_0(\cdot)$. In
\cite{B-CMP119}, $\delta_k=g_k(A_1/A_0)p_0(g_k)$. Thus \eqref{8.5:eq-averaged-extraction-e} is a
relation between two of Ba{\l}aban's exponents, $q$ against $p_0$.

Let $r_{\mathrm{pr}}$ denote the integer written $r$ in the condition $p_0\ge 5r$ of
\cite{B-CMP119}. In the order of choice of Notation~\ref{2.3:not-stages-first},
\eqref{8.5:eq-averaged-extraction-e} is a Stage-0 exponent. It is chosen after $r_{\mathrm{pr}}$
and $\mathfrak{p}_0$ and before $\kappa,\dots,\gamma$. The floor reads $\mathfrak{p}_0$ only.

The correlation theorem requires a lower bound on $q$. This bound is implied by
\eqref{8.5:eq-averaged-extraction-e} as soon as $20\,\mathfrak{p}_0\ge 7r_{\mathrm{pr}}$. That
inequality follows from $\mathfrak{p}_0\ge 5r_{\mathrm{pr}}$.

No ceiling bounds $q$ from above. Exactly one floor depends on $q$: a floor on $M$, chosen after
$q$. Raising $q$ therefore raises that floor on $M$, hence the constants $M_Q, C_{\boxplus},\dots$
chosen after $M$ and the floors on $C_0$. It lowers the ceilings $\gamma_J,\gamma_W$. There is no
cycle, since \eqref{8.5:eq-averaged-extraction-e} reads $\mathfrak{p}_0$ alone.

\emph{The ceiling.} Put
$\mathsf{E}^{\flat}_\downarrow(y):=\sup_{x\ge y}\mathsf{E}^{\flat}(x)$, and impose
\begin{equation}\label{8.5:eq-averaged-extraction-f}
\sup_{y\ge X-c_5}\bigl(LC_0A_0^2(\log\bar{x}(y))^{2\mathfrak{p}_0}\bigr)^8\,
\mathsf{E}^{\flat}_\downarrow(y)\le\frac{c_M}{32}.
\end{equation}
The quantity under the supremum is a constant times
$(\log y)^{16\mathfrak{p}_0}(\log y)^{5\mathfrak{p}_0-q}(1+o(1))$. Its degree is
$21\mathfrak{p}_0-q\le -1$ by \eqref{8.5:eq-averaged-extraction-e}, so it tends to $0$. Thus
\eqref{8.5:eq-averaged-extraction-f} is one ceiling on $g$, the same for every $s,K,m$. It reads
$C_0,\vartheta,c_M,r$ and is chosen after them.

\begin{theorem}[The averaged second-order extraction estimate]\label{8.5:thm-averaged-extraction}
Assume the following statements.
\begin{itemize}
\item[(a)] The correlation theorem, read at the value $s$ of its parameter, so that it is read at
the level $K-s$, together with the cluster-expansion lemma on which its proof depends.
\item[(b)] Ba{\l}aban's regularity theorem \cite[Theorem~1]{B-CMP102b}, and the statements
\cite[(4.14)]{B-CMP109}, \cite[(1.26)]{B-CMP116}, \cite[(2.13)]{B-CMP119},
\cite[(2.28)]{B-CMP119} and \cite[(3.28)--(3.47)]{B-CMP119}, used as stated there.
\item[(c)] From the two-point witness theorem:
\begin{itemize}
\item its bound on plaquette sums over Ba{\l}aban's regular spaces
$\mathfrak{U}_k(\{\Omega_j\},\varepsilon)$;
\item its proposition on the first logarithmic derivatives $\Phi^h_i$;
\item the lower bound $\mathfrak{M}\ge c_M\mathfrak{S}$ on its main term, where
$\mathfrak{S}:=g_{k_0}^4\mathfrak{n}(f_1)\mathfrak{n}(f_2)\rho^{-8}$.
\end{itemize}
\end{itemize}
Assume the floors \eqref{8.5:eq-averaged-extraction-e} and \eqref{8.5:eq-kept-set-a}. Assume the
following ceilings, all read at $\varepsilon'=C_w\varepsilon_n$:
\eqref{8.5:eq-source-quadratic-a-bis}--\eqref{8.5:eq-source-quadratic-e},
\eqref{8.5:eq-critical-backgrounds-a}, \eqref{8.5:eq-kept-set-b},
\eqref{8.5:eq-long-constituents-a} and \eqref{8.5:eq-averaged-extraction-c}.

Let $K$ and $s$ lie in the quantifier range of the Gaussian extraction, so that
$\rho:=L^s\mathsf{d}\in[C_W(s),LC_W(s))$, and let $n=K-s$. We assume the two-point witness
theorem's floors $\rho\ge 40M/\vartheta$ and $\rho\ge C_hM/\vartheta$. Let
$(f_1,f_2)\in\mathfrak{F}^+(\mathsf{d},\vartheta)^2$ be a pair as in the two-point witness
theorem.
\begin{itemize}
\item[(A)] For every $\mathfrak{b}_i\in\mathfrak{R}(S_i;\varepsilon')$,
\begin{equation*}
|\operatorname{Irr}^{\mathrm{sh}}[f_i](\mathfrak{b}_i)|\le\tfrac12\mathsf{I}_i
\qquad\text{and}\qquad
|N\mathcal{O}[\tilde{f}_i](\mathfrak{b}_i)|\le\mathsf{O}_i .
\end{equation*}
\item[(B)] Let $(\mathcal{X},\mu)$ be a probability space. Let $I_i,O_i\colon\mathcal{X}\to
\mathbb{R}$ be measurable with $|I_i|\le\mathsf{I}_i$ and $|O_i|\le\mathsf{O}_i$, $i=1,2$. Then
\begin{equation}\label{8.5:eq-averaged-extraction-3}
\begin{split}
|\operatorname{Cov}_\mu(I_1,O_2)|+|\operatorname{Cov}_\mu(O_1,I_2)|
&+|\operatorname{Cov}_\mu(I_1,I_2)|\\
&\le g_n^4\,\mathfrak{n}(f_1)\mathfrak{n}(f_2)\,\mathsf{E}^{\flat}(x_n).
\end{split}
\end{equation}
\item[(C)] Let $\operatorname{Irr}_i$ be the function of
Proposition~\ref{8.5:prop-readings-shortened}. Take either of the following choices.
\begin{itemize}
\item $G_1=\operatorname{Irr}_1$ and
$G_2\in\{NX^{(k_0)}(\tilde{f}_2),\,N\mathcal{O}^{(k_0)}(\tilde{f}_2),\,\operatorname{Irr}_2\}$.
Here $\mathcal{O}^{(k_0)}(\tilde{f}_2)$ is taken in any reading that evaluates
$\mathcal{O}[\tilde{f}_2]$ on $\mathfrak{R}(S_2;\varepsilon')$.
\item Symmetrically, $G_2=\operatorname{Irr}_2$ and $G_1$ in the corresponding set for $f_1$.
\end{itemize}
Then, uniformly in $K,g_0,m,N_T$, in $s$ in the range above and in the translate
$v\in\mathfrak{P}_s$ (no average over $v$ is needed),
\begin{equation}\label{8.5:eq-averaged-extraction-4}
|\operatorname{Cov}_{\nu^B_{K,s}}(G_1^v,G_2^v)|
\le g_{k_0}^4\,\mathfrak{n}(f_1)\mathfrak{n}(f_2)\,\mathsf{E}^{\flat}(x_{k_0}).
\end{equation}
The same bound holds for
$|\langle\!\langle\operatorname{Cov}_{\nu^B_{K,s}}(G_1^v,G_2^v)\rangle\!\rangle|$. In other words,
the covariance bound in part~(b) of the Gaussian extraction holds with its constant $C_b$ equal
to $0$ and with the single additional term
$g_{k_0}^4\mathfrak{n}(f_1)\mathfrak{n}(f_2)\mathsf{E}^{\flat}(x_{k_0})$.
\item[(D)] Assume the ceiling \eqref{8.5:eq-averaged-extraction-f}. For $i=1,2$, let $X_i$ be
$X^{(k_0)}(\tilde{f}_i)$ or $\mathcal{O}^{(k_0)}(\tilde{f}_i)$, and let $\mathfrak{M}$ be the main
term of the two-point witness theorem. Denote by $R^{\mathrm{irr}}$ the following sum, through
which the irrelevant remainder of the two-point witness theorem is bounded:
\begin{equation*}
\begin{split}
R^{\mathrm{irr}}:={}&|\langle\!\langle\operatorname{Cov}_{\nu^B_{K,s}}
(\operatorname{Irr}_1^v,(NX_2)^v)\rangle\!\rangle|
+|\langle\!\langle\operatorname{Cov}_{\nu^B_{K,s}}((NX_1)^v,\operatorname{Irr}_2^v)
\rangle\!\rangle|\\
&+|\langle\!\langle\operatorname{Cov}_{\nu^B_{K,s}}(\operatorname{Irr}_1^v,
\operatorname{Irr}_2^v)\rangle\!\rangle| .
\end{split}
\end{equation*}
Then
\begin{equation}\label{8.5:eq-averaged-extraction-a}
R^{\mathrm{irr}}\le g_{k_0}^4\,\mathfrak{n}(f_1)\mathfrak{n}(f_2)\,\mathsf{E}^{\flat}(x_{k_0})
\le\frac{\mathfrak{M}}{32}.
\end{equation}
The two-point witness theorem lists the condition $C_0\ge(448C_b)^{1/2}$ among its conditions on
$C_0$. With $C_b=0$ this condition holds for every $C_0$. The bound $|R_K|\le 6/32$ of that
theorem is unchanged.
\item[(E)] Consider the joint space of $(h,V,\mathsf{v})$ with the joint law $\tilde{\mu}$. Its
$(h,V)$-marginal is $\tilde{\nu}$, and the conditional law of $\mathsf{v}$ given $(h,V)$ is
$\mathbb{E}_h$. Let $\Phi^h_i$ and $\Psi^h$ be the first logarithmic derivatives in the sources
$z_i$ and the mixed second logarithmic derivative, at zero source, of the sourced density
$\tau^w_h$ of the history $h$, as in the decomposition
\begin{equation*}
S^T_{2;K,m}=E_{\tilde{\nu}}[\Psi^h]+\operatorname{Cov}_{\tilde{\nu}}(\Phi^h_1,\Phi^h_2)
\end{equation*}
of the two-point witness theorem. Then
\begin{equation}\label{8.5:eq-averaged-extraction-5}
\begin{split}
|\operatorname{Cov}_{\tilde{\mu}}(\mathsf{i}^h_1,\mathsf{o}^h_2)|
+|\operatorname{Cov}_{\tilde{\mu}}(\mathsf{o}^h_1,\mathsf{i}^h_2)|
&+|\operatorname{Cov}_{\tilde{\mu}}(\mathsf{i}^h_1,\mathsf{i}^h_2)|\\
&\le g_{k_0}^4\,\mathfrak{n}(f_1)\mathfrak{n}(f_2)\,\mathsf{E}^{\flat}(x_{k_0}).
\end{split}
\end{equation}
Each of these joint covariances is the sum of two shares: the share of that pair in
$E_{\tilde{\nu}}[\Psi^h]$, and its share in $\operatorname{Cov}_{\tilde{\nu}}(\Phi^h_1,\Phi^h_2)$.
This follows from
\begin{equation}\label{8.5:eq-averaged-extraction-6}
\operatorname{Cov}_{\tilde{\mu}}(a,b)
=E_{\tilde{\nu}}\bigl[\operatorname{Cov}_{\mathbb{E}_h}(a,b)\bigr]
+\operatorname{Cov}_{\tilde{\nu}}\bigl(\mathbb{E}_h[a],\mathbb{E}_h[b]\bigr).
\end{equation}
Consequently, the two-point witness theorem may carry the irrelevant response $\mathsf{i}^h_i$ of
each history $h$ itself, with no requirement that it be independent of $h$. The share of each
pair in $E_{\tilde{\nu}}[\Psi^h]$, namely the $\tilde{\nu}$-mean of its conditional covariance
under $\mathbb{E}_h$, is accounted for by the same bound on the joint covariance, and no separate
estimate of it is needed.
\end{itemize}
\end{theorem}

\begin{proof}
The estimate rests on four steps.

\emph{Step (I).} The irrelevant response is, as an identity, a sum over the levels $j\le n$ whose
terms are the renormalised groups of $\dot{S}^{(j)}:=\partial_\zeta E^{(j)}_\zeta|_0$ and the
vacuum-subtracted single terms $\dot{D}^{(j)}$ and $\dot{R}^{(j)}$. The shortened functional
$\operatorname{Irr}^{\mathrm{sh}}[f]$ of Definition~\ref{8.5:def-kept-set} is a sum of this form.

\emph{Step (II).} Lemma~\ref{8.5:lem-source-quadratic-bis} bounds, on the correlation theorem's own
radii and at its own $r$, the source-dependent part of every family $E^{(j)}$ by
$C_s\,r\,\hat{\theta}_{j-1}\,\mathfrak{a}_j^2e^{-\kappa d_j(X)}$.

\emph{Step (III).} Evaluate at a real regular critical fine configuration of curvature
$\varepsilon'$ at level $n$. There a group is at most $2\epsilon^2$ times the weight of the group
lemma of the correlation theorem, in the form used in the proof of
Theorem~\ref{8.5:thm-shortened-bound}. A single term is at most twice the norm bound of its
family times $(C_MD_Xt^2\mathfrak{e}/\mathfrak{a}_j)^2$. The amplitudes cancel, since
$\mathfrak{a}_j^2(\mathfrak{e}/\mathfrak{a}_j)^2=\mathfrak{e}^2$. This is the content of
Theorem~\ref{8.5:thm-shortened-bound}.

\emph{Step (IV).} A covariance is bounded by twice the product of the two sups; this is the
first line of \eqref{8.5:eq-averaged-extraction-b}.

Hypotheses (a) and (b) supply the setting of Theorem~\ref{8.5:thm-shortened-bound} and of
Lemma~\ref{8.5:lem-source-quadratic-bis}; the first item of (c) is an input of
Lemma~\ref{8.5:lem-insertion-sup}; the second and third items of (c) enter in Parts (E) and (D).

We now prove the parts in order.

\emph{Part (A).} We apply Theorem~\ref{8.5:thm-shortened-bound} to $f=f_i$ with $S=S_i$. Its
hypotheses hold for the following reasons:
\begin{itemize}
\item $f_i$ is Lipschitz, with support $S_i$ contained in a ball of radius $\rho$, by the class
$\mathfrak{F}^+(\mathsf{d},\vartheta)$;
\item $\rho\ge C_W(s)$, by the range of $s$;
\item the ceilings \eqref{8.5:eq-source-quadratic-a-bis}--\eqref{8.5:eq-source-quadratic-e},
\eqref{8.5:eq-critical-backgrounds-a} and \eqref{8.5:eq-kept-set-b} and the floor
\eqref{8.5:eq-kept-set-a} are among our assumptions;
\item $\mathfrak{b}_i\in\mathfrak{R}(S_i;\varepsilon')$.
\end{itemize}
Its conclusion is
\begin{equation*}
|\operatorname{Irr}^{\mathrm{sh}}[f_i](\mathfrak{b}_i)|\le
\bigl[C^{\flat}\hat{\theta}_{n-1}\mathfrak{e}_\star^2+C_Dr^{-1}\mathfrak{m}_n\bigr]
|S_i^+|\,\|f_i\|_\sharp .
\end{equation*}
By \eqref{8.5:eq-reading-sup-a} the right-hand side is $\tfrac12\mathsf{I}_i$. The second
inequality of (A) is the bound of Lemma~\ref{8.5:lem-insertion-sup} at
$\mathfrak{b}_i\in\mathfrak{R}(S_i;\varepsilon')$, with $\mathsf{O}_i$ as in
\eqref{8.5:eq-insertion-sup-a}.

\emph{Part (B).} By the first line of \eqref{8.5:eq-averaged-extraction-b}, each covariance is at
most twice the product of the sups. So the left-hand side of \eqref{8.5:eq-averaged-extraction-3}
is at most $2\mathsf{I}_1\mathsf{O}_2+2\mathsf{O}_1\mathsf{I}_2+2\mathsf{I}_1\mathsf{I}_2$.

We use \eqref{8.5:eq-averaged-extraction-2} with $\|f_i\|_\sharp=\|f_i\|_\infty$, where
$\mathsf{i},\mathsf{o}$ are evaluated at $x=x_n$, and then
$|S_i^+|\,\|f_i\|_\infty\le 80\,\mathfrak{n}(f_i)/\vartheta$. We get
\begin{equation*}
\begin{split}
2\mathsf{I}_1\mathsf{O}_2+2\mathsf{O}_1\mathsf{I}_2+2\mathsf{I}_1\mathsf{I}_2
&\le x^{-2}|S_1^+|\,|S_2^+|\,\|f_1\|_\infty\|f_2\|_\infty
\bigl[4\mathsf{i}\mathsf{o}+2\mathsf{i}^2\bigr]\\
&\le x^{-2}\bigl(80/\vartheta\bigr)^2\mathfrak{n}(f_1)\mathfrak{n}(f_2)
\bigl[4\mathsf{i}\mathsf{o}+2\mathsf{i}^2\bigr]\\
&=x^{-2}\mathfrak{n}(f_1)\mathfrak{n}(f_2)\cdot 6400\,\vartheta^{-2}\cdot 2
\bigl[2\mathsf{i}\mathsf{o}+\mathsf{i}^2\bigr]\\
&=x^{-2}\mathfrak{n}(f_1)\mathfrak{n}(f_2)\,\mathsf{E}^{\flat}(x),
\end{split}
\end{equation*}
by \eqref{8.5:eq-averaged-extraction-d}. Since $g_n^4=x^{-2}$, this is
\eqref{8.5:eq-averaged-extraction-3}.

\emph{Part (C).} Put $\mu:=\nu^B_{K,s}$, and let $\mathcal{X}$ be the space of level-$n$ fields
$V$. Fix $v\in\mathfrak{P}_s$. Put $I_i:=\operatorname{Irr}_i^v$ and
$O_i:=(NX^{(k_0)}(\tilde{f}_i))^v$, respectively the chosen reading of
$(N\mathcal{O}^{(k_0)}(\tilde{f}_i))^v$; when $G_2=\operatorname{Irr}_2$ the pair is $(I_1,I_2)$,
and symmetrically for $G_1$.

The translate by $v$ is a translate by a level-$n$ lattice vector. It maps the class of
$(f_1,f_2)$ to itself, with the same $\vartheta$, $\rho$, norms, $\mathfrak{n}$ and $|S^+|$.
Every constant in (A) is independent of position. Hence $\mathsf{I}_i$ and $\mathsf{O}_i$ are the
same for the translated pair.

Outside the window on $N_i$, both functions vanish, by the convention of the two-point witness
theorem. Inside the window, Proposition~\ref{8.5:prop-readings-shortened} shows that the
configuration at which $\operatorname{Irr}_i$ is evaluated lies in
$\mathfrak{R}(S_i;\varepsilon')$. So (A) gives
$|I_i|\le\tfrac12\mathsf{I}_i\le\mathsf{I}_i$. Lemma~\ref{8.5:lem-insertion-sup} gives
$|O_i|\le\mathsf{O}_i$, for $X^{(k_0)}(\tilde{f}_i)$ as well as for
$\mathcal{O}^{(k_0)}(\tilde{f}_i)$.

Thus the hypotheses of (B) hold. Each single covariance in (C) is one of the three terms on the
left-hand side of \eqref{8.5:eq-averaged-extraction-3}. It is therefore at most the bound there,
with $n=k_0$. This is \eqref{8.5:eq-averaged-extraction-4}.

Finally, $\langle\!\langle\cdot\rangle\!\rangle$ is the mean over the finite set
$\mathfrak{P}_s$. The absolute value of a mean is at most the mean of the absolute values, and
the mean of numbers bounded by a number $a$ is bounded by $a$. This gives the averaged form.

\emph{Part (D).} For each $v$, the three covariances in $R^{\mathrm{irr}}$ are the three terms of
\eqref{8.5:eq-averaged-extraction-3}. The inputs are the same as in Part~(C), with
$I_i=\operatorname{Irr}_i^v$ and $O_i=(NX_i)^v$. Their sum is therefore at most
$g_{k_0}^4\mathfrak{n}(f_1)\mathfrak{n}(f_2)\mathsf{E}^{\flat}(x_{k_0})$. Taking absolute values
inside the mean over $\mathfrak{P}_s$ and then the mean of this sum gives the first inequality of
\eqref{8.5:eq-averaged-extraction-a}.

For the second inequality, the lower bound of hypothesis~(c) reads
\begin{equation*}
\mathfrak{M}\ge c_M\,g_{k_0}^4\mathfrak{n}(f_1)\mathfrak{n}(f_2)\rho^{-8}.
\end{equation*}
It therefore suffices that $\rho^8\mathsf{E}^{\flat}(x_{k_0})\le c_M/32$. We collect:
\begin{itemize}
\item $\rho<LC_W(s)=LC_0A_0^2(\log\bar{x}_s)^{2\mathfrak{p}_0}$, by the definition of $C_W(s)$;
\item $\bar{x}_s\le\bar{x}(\underline{x}_s)$, by the envelope of the two-point witness theorem;
\item $x_{k_0}\ge\underline{x}_s$, so
$\mathsf{E}^{\flat}(x_{k_0})\le\mathsf{E}^{\flat}_\downarrow(\underline{x}_s)$;
\item $\underline{x}_s\ge X-c_5$.
\end{itemize}
Since the logarithm is increasing, these give
\begin{equation*}
\rho^8\mathsf{E}^{\flat}(x_{k_0})\le
\bigl(LC_0A_0^2(\log\bar{x}(\underline{x}_s))^{2\mathfrak{p}_0}\bigr)^8
\mathsf{E}^{\flat}_\downarrow(\underline{x}_s)\le\frac{c_M}{32},
\end{equation*}
by \eqref{8.5:eq-averaged-extraction-f} at $y=\underline{x}_s$.

By (C), $C_b=0$, so the condition $C_0\ge(448C_b)^{1/2}$ holds for every $C_0$. The bound
$R^{\mathrm{irr}}\le\mathfrak{M}/32$ is \eqref{8.5:eq-averaged-extraction-a}, so the bound
$|R_K|\le 6/32$ is unchanged.

\emph{Part (E), the bound.} Put $\mu:=\tilde{\mu}$. Proposition~\ref{8.5:prop-reading-sup} gives
$\sup|\mathsf{i}^h_i|\le\mathsf{I}_i$, and Lemma~\ref{8.5:lem-insertion-sup} gives
$\sup|\mathsf{o}^h_i|\le\mathsf{O}_i$. Part~(B), which holds under any probability measure,
applies with $I_i:=\mathsf{i}^h_i$ and $O_i:=\mathsf{o}^h_i$. With $n=k_0$ it gives
\eqref{8.5:eq-averaged-extraction-5}.

\emph{Part (E), the decomposition.} Let $\mathsf{A}^w$ denote the full exponent of the integrand
in the sourced density, so that $\tau^w_h=\chi_h\int dm_h\,e^{\mathsf{A}^w}$, where $dm_h$ is the
measure over the integrated variables $\mathsf{v}$ of $h$. Write $\partial_i:=\partial_{z_i}$ and
$\mathsf{A}^0:=\mathsf{A}^w|_{z=0}$. The factor $\chi_h$ carries no source, and
$\mathbb{E}_h$ is the normalised measure $e^{\mathsf{A}^0}dm_h/\int dm_h\,e^{\mathsf{A}^0}$ on the
support of $m_h$. Differentiating the logarithm of an integral of an exponential once and twice
at $z=0$ gives
\begin{equation*}
\begin{split}
\Phi^h_i&=\mathbb{E}_h\bigl[\partial_i\mathsf{A}^w|_0\bigr],\\
\Psi^h&=\mathbb{E}_h\bigl[\partial_1\partial_2\mathsf{A}^w|_0\bigr]
+\operatorname{Cov}_{\mathbb{E}_h}\bigl(\partial_1\mathsf{A}^w|_0,\partial_2\mathsf{A}^w|_0\bigr).
\end{split}
\end{equation*}
Indeed, $\partial_1\partial_2\log\int e^{\mathsf{A}^w}$ equals
\begin{equation*}
\frac{\int e^{\mathsf{A}^w}\bigl(\partial_1\partial_2\mathsf{A}^w
+\partial_1\mathsf{A}^w\,\partial_2\mathsf{A}^w\bigr)}{\int e^{\mathsf{A}^w}}
\end{equation*}
minus the product of the two first logarithmic derivatives.

Insert these expressions into the decomposition of the two-point witness theorem,
\begin{equation*}
S^T_{2;K,m}=E_{\tilde{\nu}}[\Psi^h]
+\operatorname{Cov}_{\tilde{\nu}}(\Phi^h_1,\Phi^h_2).
\end{equation*}
Then apply \eqref{8.5:eq-averaged-extraction-6}. This is the law of total covariance, an identity
for square-integrable random variables on the joint space with marginal $\tilde{\nu}$ and
conditional law $\mathbb{E}_h$. It gives
\begin{equation*}
S^T_{2;K,m}=E_{\tilde{\nu}}\bigl[\mathbb{E}_h[\partial_1\partial_2\mathsf{A}^w|_0]\bigr]
+\operatorname{Cov}_{\tilde{\mu}}\bigl(\partial_1\mathsf{A}^w|_0,\partial_2\mathsf{A}^w|_0\bigr).
\end{equation*}

Now split $\partial_i\mathsf{A}^w|_0$ into its members, among which are $\mathsf{i}^h_i$ and
$\mathsf{o}^h_i$. By bilinearity of both sides of \eqref{8.5:eq-averaged-extraction-6}, the joint
covariance of each pair of members equals the sum of two terms. The first is its share in
$E_{\tilde{\nu}}[\operatorname{Cov}_{\mathbb{E}_h}(\cdot,\cdot)]$, which is part of
$E_{\tilde{\nu}}[\Psi^h]$. The second is its share in
$\operatorname{Cov}_{\tilde{\nu}}(\Phi^h_1,\Phi^h_2)$.

For the three pairs of \eqref{8.5:eq-averaged-extraction-5}, the joint covariance is thus the sum
of the two shares, and it is the joint covariance that \eqref{8.5:eq-averaged-extraction-5}
bounds. The share in $E_{\tilde{\nu}}[\Psi^h]$, the $\tilde{\nu}$-mean of the conditional
covariance under $\mathbb{E}_h$, therefore enters only through that sum and needs no separate
estimate; no property of $\mathsf{i}^h_i$ beyond its sup bound has been used.
\end{proof}

\begin{corollary}[Joint sup bound on insertion and irrelevant response]\label{8.5:cor-joint-sup}
In the setting of Theorem~\ref{8.5:thm-averaged-extraction}, let $i\in\{1,2\}$, let
$n=k_0=K-s$ be the witness level, let $\varepsilon$ and $x$ denote the small-field threshold
and the inverse squared coupling at which part~(d) of the source lemma is stated, let
$\mathsf{O}_i$ and $\mathsf{I}_i$ be the sup bounds \eqref{8.5:eq-insertion-sup-a} and
\eqref{8.5:eq-reading-sup-a}, and put
\begin{equation*}
C_G:=80\,\vartheta^{-1}\bigl(C_{\mathcal{O}}+4C^{\flat}C_\theta C_{\mathfrak{e}}^2+1\bigr).
\end{equation*}
Besides the ceiling \eqref{8.5:eq-averaged-extraction-c}, which is among the hypotheses of
Theorem~\ref{8.5:thm-averaged-extraction}, assume
\begin{equation}\label{8.5:eq-joint-sup-1}
2C_DC_{\mathfrak{m}}\,2^{(\kappa_0-2)/2}\,r^{-1}\,x^{-(\kappa_0-4)/2}\le A_0^2 ,
\end{equation}
where $r$ is the source radius entering \eqref{8.5:eq-reading-sup-a} (the radius of the disc
$D(0,r)$ of source constants in Corollary~\ref{8.5:cor-source-derivative-bis}) and $\kappa_0$ is the
exponent in $\mathfrak{m}_n=C_{\mathfrak{m}}y_n^{-(\kappa_0-2)/2}$.
Then
\begin{equation}\label{8.5:eq-joint-sup-2}
\bigl|N X^{(k_0)}(\tilde{f}_i)+\operatorname{Irr}_i\bigr|
\le \mathsf{O}_i+\mathsf{I}_i
\le C_G\,\mathfrak{n}(f_i)\,\varepsilon^2 .
\end{equation}
In particular, the bound $|G_i|\le C_G\mathfrak{n}(f_i)\varepsilon^2$ of part~(d) of the
source lemma, whose $G_i$ is the quantity bounded in \eqref{8.5:eq-joint-sup-2}, holds at the
radii of the correlation theorem and at fixed $r$.
\end{corollary}

\begin{proof}
\emph{The first inequality.} By Lemma~\ref{8.5:lem-insertion-sup}, the local insertion
$X^{(k_0)}(\tilde{f}_i)$ obeys the same bound as the main curvature insertion, so that
$|NX^{(k_0)}(\tilde{f}_i)|\le\mathsf{O}_i$ with
$\mathsf{O}_i=C_{\mathcal{O}}\varepsilon_n^2|S_i^+|\,\|f_i\|_\infty$.
By Proposition~\ref{8.5:prop-reading-sup}, every reading of the irrelevant response obeys
$\sup|\operatorname{Irr}_i|\le\mathsf{I}_i$ (and the reading on $\Omega_i$ even obeys
$\sup|\operatorname{Irr}_i|\le\frac12\mathsf{I}_i$). The triangle inequality gives the first
inequality of \eqref{8.5:eq-joint-sup-2}, whichever reading of $\operatorname{Irr}_i$ is used.

\emph{The second inequality.} By \eqref{8.5:eq-reading-sup-a},
\begin{equation*}
\mathsf{I}_i=2\bigl[C^{\flat}\hat{\theta}_{n-1}\mathfrak{e}_\star^2
+C_Dr^{-1}\mathfrak{m}_n\bigr]|S_i^+|\,\|f_i\|_\sharp .
\end{equation*}
The hypothesis on $\hat{\theta}$ read here is the ceiling \eqref{8.5:eq-averaged-extraction-c},
and the only property of $\hat{\theta}_{n-1}$ used is the bound $\hat{\theta}_{n-1}\le C_\theta$;
together with $\mathfrak{e}_\star=C_{\mathfrak{e}}\varepsilon_n$ this gives
$C^{\flat}\hat{\theta}_{n-1}\mathfrak{e}_\star^2\le C^{\flat}C_\theta C_{\mathfrak{e}}^2\varepsilon_n^2$.
In particular the floor \eqref{8.5:eq-averaged-extraction-e}, although among the hypotheses of
Theorem~\ref{8.5:thm-averaged-extraction}, is not used in this argument. Since
$\|f_i\|_\sharp=\max\{\|f_i\|_\infty,40M\operatorname{Lip}f_i\}\ge\|f_i\|_\infty$, we obtain
\begin{equation}\label{8.5:eq-joint-sup-3}
\begin{aligned}
\mathsf{O}_i+\mathsf{I}_i
\le{}& \Bigl[\bigl(C_{\mathcal{O}}+2C^{\flat}C_\theta C_{\mathfrak{e}}^2\bigr)\varepsilon_n^2\\
&\quad+2C_Dr^{-1}\mathfrak{m}_n\Bigr]|S_i^+|\,\|f_i\|_\sharp .
\end{aligned}
\end{equation}
In \eqref{8.5:eq-joint-sup-3} the first term carries the factor $\varepsilon_n^2$, and the second
term carries $\mathfrak{m}_n=C_{\mathfrak{m}}y_n^{-(\kappa_0-2)/2}$; it is to this second term that
the hypothesis \eqref{8.5:eq-joint-sup-1} applies. Under the ceiling
\eqref{8.5:eq-averaged-extraction-c} and the hypothesis \eqref{8.5:eq-joint-sup-1}, the right-hand
side of \eqref{8.5:eq-joint-sup-3} is at most $C_G\mathfrak{n}(f_i)\varepsilon^2$, with $C_G$ as in
the statement. This is the second inequality of \eqref{8.5:eq-joint-sup-2}.
\end{proof}

\begin{corollary}[The pointwise form for every top field]\label{8.5:cor-pointwise-form}
Let the hypotheses of the averaged second-order extraction estimate
(Theorem~\ref{8.5:thm-averaged-extraction}) hold, and let $T^{\langle h_0\rangle}$ be the
truncated expectation of the every-volume route, formed with any tower of positive normalised
maps. In it, $\operatorname{Irr}_1=\operatorname{Irr}^{\mathrm{sh}}[f_1]$ is the shortened
irrelevant functional of the first test function, and $G_2$ is the curvature insertion of the
second. Consider the lattice theory with $K'=n=K-s$ renormalisation steps (bare spacing
$L^{-n}$), the one to which the correlation theorem is applied. Then
\begin{equation}\label{8.5:eq-pointwise-form-1}
\bigl|T^{\langle h_0\rangle}[\operatorname{Irr}_1,G_2](W)\bigr|
\le g^4\,\mathfrak{n}\,\mathfrak{n}\,\mathsf{E}^{\flat}(x)
\end{equation}
for every top field $W$, every pair of shapes $S_1,S_2$ satisfying the condition
\eqref{8.5:eq-kept-set-a} and every separation $\mathsf{d}$ of the supports. Here
$\mathsf{E}^{\flat}(x)$ is the envelope function of \eqref{8.5:eq-averaged-extraction-d}, and
$g$, $x$ and the two factors $\mathfrak{n}$ are those of
\eqref{8.5:eq-averaged-extraction-b}. The bound \eqref{8.5:eq-pointwise-form-1} carries no
decay in the distance. It also carries only a negative power of $\log x$ through
$\mathsf{E}^{\flat}(x)$.
\end{corollary}

\begin{proof}
A tower of positive normalised maps, evaluated at a fixed top field $W$, is a positive
normalised linear functional, that is, the expectation of a probability measure, and
$T^{\langle h_0\rangle}[\operatorname{Irr}_1,G_2](W)$ is the covariance of
$\operatorname{Irr}_1$ and $G_2$ under it. By Part~(A) of
Theorem~\ref{8.5:thm-averaged-extraction} the two insertions are bounded by $\mathsf{I}_1$ and
$\mathsf{O}_2$, so the second estimate of \eqref{8.5:eq-averaged-extraction-b}, which holds for
any such measure, applies; read at each top field $W$ separately, it is the bound
\eqref{8.5:eq-pointwise-form-1}. Its right-hand side does not depend on $W$ or on the distance
between the supports, and the estimate holds for every pair of shapes satisfying
\eqref{8.5:eq-kept-set-a}. Hence the same bound holds for every top field $W$, every separation
$\mathsf{d}$ and every pair of shapes satisfying \eqref{8.5:eq-kept-set-a}. This is
\eqref{8.5:eq-pointwise-form-1}.
\end{proof}

\begin{corollary}[The original form with unit constant]\label{8.5:cor-unit-constant}
Assume the hypotheses of the averaged second-order extraction estimate
(Theorem~\ref{8.5:thm-averaged-extraction}), with $K$, $s$, the witness level $k_0=K-s$, the
ratio $\rho$, the inverse couplings $x_k=g_k^{-2}$ and the pair $(f_1,f_2)$ as there; let
$\mathsf{E}^{\flat}$ be the function of \eqref{8.5:eq-averaged-extraction-d},
$\mathsf{E}^{\flat}_\downarrow$ its envelope and $\bar{x}(\cdot)$ the function fixed there, and
let $\mathcal{K}_{K,s}(f_1,f_2)$ be the kernel of the continuum-kernel statement of the two-point
witness theorem. Assume in addition the ceiling
\begin{equation}\label{8.5:eq-unit-constant-a}
\sup_{y\ge X-c_5}
\bigl(LC_0A_0^2(\log\bar{x}(y))^{2\mathfrak{p}_0}\bigr)^{10}
\bigl(A_0(\log y)^{\mathfrak{p}_0}\bigr)^{-4}\,
\mathsf{E}^{\flat}_\downarrow(y)\;\le\; b_0c_{\mathcal{K}} ,
\end{equation}
where $X=g^{-2}$. Then
\begin{equation}\label{8.5:eq-unit-constant-1}
g^4\,\mathfrak{n}(f_1)\,\mathfrak{n}(f_2)\,\mathsf{E}^{\flat}(x_{k_0})
\;\le\;\rho^{-2}\,p_0(g_{k_0})^4\,b_0\,g_{k_0}^4\,\mathcal{K}_{K,s}(f_1,f_2).
\end{equation}
Consequently the covariance bound of part (b) of the Gaussian extraction holds in its original
form with the constant $C_b$ of that bound equal to $1$ and without any error term, and the
condition $C_0^2\ge 448\,C_b$ used in the proof of the two-point witness theorem becomes the
condition $C_0^2\ge 448$; the ceiling \eqref{8.5:eq-unit-constant-a} reads $C_0$ and is imposed
after it, so no cycle enters the order of choice.
\end{corollary}

\begin{proof}
We first record the size of the quantity under the supremum in
\eqref{8.5:eq-unit-constant-a}. Counted in powers of the logarithm, its three factors have
degrees $20\mathfrak{p}_0$, $-4\mathfrak{p}_0$ and $5\mathfrak{p}_0-q$, so that its degree is
\begin{equation*}
20\mathfrak{p}_0-4\mathfrak{p}_0+5\mathfrak{p}_0-q\;=\;21\mathfrak{p}_0-q\;\le\;-1 ,
\end{equation*}
by the floor \eqref{8.5:eq-averaged-extraction-e} on $q$, one of the hypotheses of
Theorem~\ref{8.5:thm-averaged-extraction}. The ceiling
\eqref{8.5:eq-unit-constant-a} is thus a bound by $b_0c_{\mathcal{K}}$ on a quantity of negative
degree in the logarithm.

By the lower bound in part (e) of the continuum-kernel statement of the two-point witness
theorem,
\begin{equation}\label{8.5:eq-unit-constant-2}
\mathcal{K}_{K,s}(f_1,f_2)\;\ge\;c_{\mathcal{K}}\,\rho^{-8}\,\mathfrak{n}(f_1)\,\mathfrak{n}(f_2).
\end{equation}
Multiplying \eqref{8.5:eq-unit-constant-2} by the positive number
$\rho^{-2}p_0(g_{k_0})^4b_0g_{k_0}^4$, we obtain
\begin{equation*}
\rho^{-2}p_0(g_{k_0})^4\,b_0\,g_{k_0}^4\,\mathcal{K}_{K,s}(f_1,f_2)
\;\ge\;\rho^{-10}\,b_0c_{\mathcal{K}}\,p_0(g_{k_0})^4\,g_{k_0}^4\,
\mathfrak{n}(f_1)\,\mathfrak{n}(f_2).
\end{equation*}
Hence \eqref{8.5:eq-unit-constant-1} follows once
\begin{equation}\label{8.5:eq-unit-constant-3}
g^4\,\mathsf{E}^{\flat}(x_{k_0})\;\le\;\rho^{-10}\,b_0c_{\mathcal{K}}\,p_0(g_{k_0})^4\,g_{k_0}^4 .
\end{equation}
By the definition of the degree function, and since $x_{k_0}=g_{k_0}^{-2}$, we have
$p_0(g_{k_0})=A_0(\log x_{k_0})^{\mathfrak{p}_0}$; thus the second factor under the supremum in
\eqref{8.5:eq-unit-constant-a}, read at $y=x_{k_0}$, is $p_0(g_{k_0})^{-4}$.

We read \eqref{8.5:eq-unit-constant-a} at $y=x_{k_0}$. Three comparisons enter, each concerning
objects fixed in the setting of Theorem~\ref{8.5:thm-averaged-extraction}. First,
$x_{k_0}\ge X-c_5$, so that $y=x_{k_0}$ lies in the range of the supremum; this is the range
condition on the inverse couplings at the levels $k\le K$ that comes with the first-passage level
$K$ of Definition~\ref{1.2:def-flow-constants}. Second, the hypotheses of
Theorem~\ref{8.5:thm-averaged-extraction} place $\rho$ in the window
$C_W(s)\le\rho<LC_W(s)$, and the definition of $C_W(s)$ through the function $\bar{x}(\cdot)$
bounds $C_W(s)$ by $C_0A_0^2(\log\bar{x}(x_{k_0}))^{2\mathfrak{p}_0}$; together they give
\begin{equation*}
\rho^{10}\;\le\;\bigl(LC_0A_0^2(\log\bar{x}(x_{k_0}))^{2\mathfrak{p}_0}\bigr)^{10}.
\end{equation*}
Third, by the definition of the envelope $\mathsf{E}^{\flat}_\downarrow$ in
Theorem~\ref{8.5:thm-averaged-extraction},
\begin{equation*}
g^4\,\mathsf{E}^{\flat}(x_{k_0})\;\le\;g_{k_0}^4\,\mathsf{E}^{\flat}_\downarrow(x_{k_0}).
\end{equation*}
With these, the ceiling at $y=x_{k_0}$ and the second comparison give
\begin{align*}
\mathsf{E}^{\flat}_\downarrow(x_{k_0})
&\le b_0c_{\mathcal{K}}\,p_0(g_{k_0})^4
\bigl(LC_0A_0^2(\log\bar{x}(x_{k_0}))^{2\mathfrak{p}_0}\bigr)^{-10}\\
&\le \rho^{-10}\,b_0c_{\mathcal{K}}\,p_0(g_{k_0})^4 ,
\end{align*}
and multiplying by $g_{k_0}^4$ and using the third comparison yields
\eqref{8.5:eq-unit-constant-3}. This proves \eqref{8.5:eq-unit-constant-1}.

By parts (A) and (B) of Theorem~\ref{8.5:thm-averaged-extraction}, with the functions of
\eqref{8.5:eq-averaged-extraction-d} read at $x=x_{k_0}$, the translation-averaged covariances
of the pairs considered in part (b) of the Gaussian extraction are at most the left-hand side of
\eqref{8.5:eq-unit-constant-1}. By
\eqref{8.5:eq-unit-constant-1} these averaged covariances are therefore bounded by
$\rho^{-2}p_0(g_{k_0})^4b_0g_{k_0}^4\mathcal{K}_{K,s}(f_1,f_2)$, which is the right-hand side of
part (b) with $C_b=1$ and with no error term.

The constant $C_b=1$ depends on no other constant. The condition $C_0^2\ge 448\,C_b$ is therefore
the lower bound $C_0^2\ge 448$ on $C_0$ alone. The ceiling \eqref{8.5:eq-unit-constant-a} reads
$C_0$ and is imposed after $C_0$ has been fixed, so the order of choice contains no cycle.
\end{proof}

Either this corollary or Theorem~\ref{8.5:thm-averaged-extraction} itself suffices for the bound
on the irrelevant remainder in the proof of the two-point witness theorem.

\begin{remark}[Logarithmic gain as decay in the radius]\label{8.5:rem-logarithmic-gain}
The envelope $\mathsf{E}^{\flat}(x)$ in \eqref{8.5:eq-averaged-extraction-d} of the averaged
second-order extraction estimate (Theorem~\ref{8.5:thm-averaged-extraction}) carries the factor
$(\log x)^{5\mathfrak{p}_0-q}$
and no power of $g$. That estimate therefore gives no bound that decays as a power of the coupling
for pairs of test functions at large separation.

It does give decay in the radius. In the range
\begin{equation*}
C_W(s)\le\rho<LC_W(s)
\end{equation*}
of the two-point witness theorem, a negative power of $\log x_{k_0}$, the logarithm of the inverse
squared coupling at the witness level $k_0=K-s$, is bounded by a negative power of $\rho$.
Let $\mathfrak{M}$ be the main term of the two-point witness theorem and $R^{\mathrm{irr}}$ the
contribution of the shortened irrelevant functionals $\operatorname{Irr}^{\mathrm{sh}}[f_i]$ of
Theorem~\ref{8.5:thm-averaged-extraction} to the averaged extraction. Then
\begin{equation}\label{8.5:eq-logarithmic-gain-1}
R^{\mathrm{irr}}\le C\,\rho^{-a}\,\mathfrak{M},\qquad
a:=\frac{q-21\mathfrak{p}_0}{2\mathfrak{p}_0}>0,
\end{equation}
with a constant $C$ that does not depend on $\rho$.
\end{remark}

\begin{proof}
Let $\rho$ lie in the range of the two-point witness theorem, and let $k_0=K-s$ be the witness
level. In this range the envelope quantities $\bar{x}_s$ and $\underline{x}_s$ of that theorem satisfy
\begin{equation}\label{8.5:eq-logarithmic-gain-2}
\log\bar{x}_s>\Bigl(\frac{\rho}{LC_0A_0^2}\Bigr)^{1/(2\mathfrak{p}_0)}
\end{equation}
and
\begin{equation*}
\log\underline{x}_s\ge\log\bar{x}_s-C' ,
\end{equation*}
with a constant $C'$ of the two-point witness theorem.
Moreover $x_{k_0}=g_{k_0}^{-2}$ satisfies $\log x_{k_0}\ge\log\underline{x}_s$. Combining these
three inequalities gives
\begin{equation}\label{8.5:eq-logarithmic-gain-3}
\log x_{k_0}\ge\log\underline{x}_s\ge\log\bar{x}_s-C'
>\Bigl(\frac{\rho}{LC_0A_0^2}\Bigr)^{1/(2\mathfrak{p}_0)}-C' .
\end{equation}
Hence, for $b>0$, the negative power $(\log x_{k_0})^{-b}$ is bounded by a constant times
$\rho^{-b/(2\mathfrak{p}_0)}$. Here $L$, $C_0$, $A_0$ and $\mathfrak{p}_0$ do not depend on
$\rho$, so this constant does not either.

In the bound on the irrelevant remainder $R^{\mathrm{irr}}$ by the main term $\mathfrak{M}$
obtained from Theorem~\ref{8.5:thm-averaged-extraction}, the quantity $\log x_{k_0}$ enters with
total exponent $-(q-21\mathfrak{p}_0)$. Applying the preceding paragraph with
$b=q-21\mathfrak{p}_0$ converts this power into $\rho^{-b/(2\mathfrak{p}_0)}=\rho^{-a}$, with $a$ as
in \eqref{8.5:eq-logarithmic-gain-1}, and this gives \eqref{8.5:eq-logarithmic-gain-1}.
Finally, the positivity $a>0$ is the inequality $q>21\mathfrak{p}_0$ between the exponents, which
holds under the floor on $q$, \eqref{8.5:eq-averaged-extraction-e}, assumed in
Theorem~\ref{8.5:thm-averaged-extraction}.
\end{proof}

\begin{lemma}[The irrelevant remainder as a sum over levels]\label{8.5:lem-remainder-level-sum}
Let $k_0=K-s$. Work with the sourced format of the effective action at level $k_0$ given by the
correlation theorem, read on the whole lattice. That is, the localisation functions $\varphi_j$
are identically $1$, and every level-$j$ localisation domain is the whole level-$j$ block lattice
$\mathbb{T}^{(j)}$ of Definition~\ref{1.5:def-block-average}. The field-free constant $E^w$ of
that format does not depend on the configuration.

For $j\le k_0$ and a source $w$ we suppress the remaining arguments of the members of the level-$j$
family of the format, and write $E^{(j)}(X,U,y;w)$ for the member attached to the region $X$ and
the point $y$. Put
\begin{equation*}
E^{(j)}(U;w):=\sum_{y\in\mathbb{T}^{(j)}}\sum_{X\ni y}E^{(j)}(X,U,y;w)
\end{equation*}
for its total. Write $R^{(j)}(U;w)$ for the level-$j$ total of the remainder family, so that
\begin{equation*}
\mathbf{R}^w(U)-\mathbf{R}^w(1)=\sum_{j\le k_0}\big[R^{(j)}(U;w)-R^{(j)}(1;w)\big].
\end{equation*}

Let $f$ be a test function and $z$ a complex variable, and take $w=zf$. Put
\begin{equation*}
\dot{E}^{(j)}[f]:=\partial_zE^{(j)}(\cdot\,;zf)\big|_{z=0},\qquad
\dot{R}^{(j)}[f]:=\partial_zR^{(j)}(\cdot\,;zf)\big|_{z=0},
\end{equation*}
\begin{equation*}
I_jf:=\sum_{y\in\mathbb{T}^{(j)}}f(p_y)\,h^{(j)}_y,\qquad F^{\mathrm{vac}}:=F-F(1).
\end{equation*}
Here $p_y$ is the evaluation plaquette and $h^{(j)}_y$ the profile attached to $y$ in the
correlation theorem's format, and $1$ denotes the configuration with every bond variable equal to
the identity.
Let $\operatorname{Irr}[f]$ be the irrelevant remainder of the carried source in the two-point
witness theorem, namely
\begin{equation*}
\begin{split}
\operatorname{Irr}[f]={}&\partial_z\big[\mathbf{E}^w+\mathbf{D}^w+\mathbf{R}^w-E^w\big]_{w=zf}
\Big|_{z=0}(U_{k_0})\\
&-\partial_z\big[\mathbf{E}^w+\mathbf{D}^w+\mathbf{R}^w-E^w\big]_{w=zf}\Big|_{z=0}(1).
\end{split}
\end{equation*}
Then the following hold.
\begin{itemize}
\item[(a)] The remainder is a sum over levels:
\begin{equation}\label{8.5:eq-remainder-level-sum-a}
\begin{split}
\operatorname{Irr}[f]={}&\sum_{j\le k_0}\Big\{\big(\dot{E}^{(j)}[f]\big)^{\mathrm{vac}}(U_{k_0})
-b_j'(0)\,A\big(I_jf,U_{k_0}\big)\Big\}\\
&+\sum_{j\le k_0}\big(\dot{R}^{(j)}[f]\big)^{\mathrm{vac}}(U_{k_0}).
\end{split}
\end{equation}
\item[(b)] The families $E^{(j)}$ enter $\operatorname{Irr}[f]$ only through their totals
$E^{(j)}(U;w)$. How the members are localised, that is, the regions $X$ to which they are
attached, and how they are pointed, that is, the points $y$, may therefore be chosen freely.
\end{itemize}
\end{lemma}

\begin{proof}
All sums below are finite, because $\mathbb{T}_m$ and its block lattices are finite. Derivatives
in $z$ may therefore be taken term by term.

\emph{The $\mathbf{E}$- and $\mathbf{D}$-families.} By the definition of $\mathbf{E}^w$ in the
sourced format, $\mathbf{E}^w=\sum_{j\le k_0}\sum_{y}\mathfrak{i}^{(j)}_y$. Here
$\mathfrak{i}^{(j)}_y$ is built on the family $E^{(j)}_{w(p_y)}$ frozen at the source value
$w(p_y)$, and its $\beta$-coefficient is $b_j(w(p_y))$. Explicitly,
\begin{equation*}
\mathfrak{i}^{(j)}_y(U)=\sum_{X\ni y}\big[E^{(j)}_{w(p_y)}(X,U,y)-E^{(j)}_{w(p_y)}(X,1,y)\big]
-b_j(w(p_y))\,A\big(h^{(j)}_y,U\big).
\end{equation*}
The format defines $D^{(j)}(X,U,y;w):=E^{(j)}(X,U,y;w)-E^{(j)}_{w(p_y)}(X,U,y)$, and
$\mathbf{D}^w$ is the sum over $j\le k_0$, $y$, $X\ni y$ of
$D^{(j)}(X,U,y;w)-D^{(j)}(X,1,y;w)$.

Add the two sums. For each triple $(j,y,X)$, the frozen member $E^{(j)}_{w(p_y)}(X,U,y)$ of
$\mathfrak{i}^{(j)}_y$ cancels the term $-E^{(j)}_{w(p_y)}(X,U,y)$ of $D^{(j)}(X,U,y;w)$.
In the same way, $-E^{(j)}_{w(p_y)}(X,1,y)$ cancels $+E^{(j)}_{w(p_y)}(X,1,y)$ coming from
$-D^{(j)}(X,1,y;w)$. What remains is
\begin{equation}\label{8.5:eq-remainder-level-sum-1}
(\mathbf{E}^w+\mathbf{D}^w)(U)=\sum_{j\le k_0}\Big\{E^{(j)}(U;w)-E^{(j)}(1;w)
-\sum_{y\in\mathbb{T}^{(j)}}b_j(w(p_y))\,A\big(h^{(j)}_y,U\big)\Big\}.
\end{equation}
Now set $U=1$. The difference $E^{(j)}(1;w)-E^{(j)}(1;w)$ vanishes. Moreover
\begin{equation*}
A(h,1)=\sum_ph(p)\,\big[1-\operatorname{Re}\operatorname{tr}U(\partial p)\big]\Big|_{U\equiv1}=0,
\end{equation*}
because the trace of the identity is $1$ in the normalisation of
Definition~\ref{1.1:def-wilson-action}; here $A(h,U)$ is the weighted form of the Wilson action of
that definition. Hence the right side of
\eqref{8.5:eq-remainder-level-sum-1} vanishes at $U=1$.

\emph{The constant.} $E^w$ does not depend on the configuration. Its values at $U_{k_0}$ and
at $1$ are therefore equal, and $\partial_zE^{zf}|_{z=0}$ cancels in the difference that
defines $\operatorname{Irr}[f]$.

\emph{Differentiation.} Put $w=zf$ in \eqref{8.5:eq-remainder-level-sum-1}. Then
\begin{equation*}
\partial_z\big[E^{(j)}(U;zf)-E^{(j)}(1;zf)\big]\big|_{z=0}
=\dot{E}^{(j)}[f](U)-\dot{E}^{(j)}[f](1)=\big(\dot{E}^{(j)}[f]\big)^{\mathrm{vac}}(U).
\end{equation*}
Also $\partial_zb_j(zf(p_y))|_{z=0}=f(p_y)\,b_j'(0)$. The weighted action $A(h,U)$ is linear in
the weight $h$, so
\begin{equation*}
\sum_{y\in\mathbb{T}^{(j)}}b_j'(0)f(p_y)\,A\big(h^{(j)}_y,U\big)=b_j'(0)\,A(I_jf,U).
\end{equation*}
The same computation gives
\begin{equation*}
\partial_z\big[\mathbf{R}^{zf}(U)-\mathbf{R}^{zf}(1)\big]\big|_{z=0}
=\sum_{j\le k_0}\big(\dot{R}^{(j)}[f]\big)^{\mathrm{vac}}(U).
\end{equation*}

Combine the three contributions, subtract the value at $1$ (which is zero for the $\mathbf{E}$-
and $\mathbf{D}$-part), and evaluate at $U=U_{k_0}$. This gives
\eqref{8.5:eq-remainder-level-sum-a}, which is (a).

\emph{Part (b).} In \eqref{8.5:eq-remainder-level-sum-a} the families $E^{(j)}$ appear only
through $\dot{E}^{(j)}[f]$. This is the $z$-derivative of the total $E^{(j)}(\cdot\,;zf)$, which
does not record the regions $X$ or the points $y$ carrying the individual members.
\end{proof}

\begin{definition}[The two systems of field radii]\label{8.5:not-field-radii}
\emph{Standing setting.} In what follows the constants are those of the quantifier prefix of
the correlation theorem, chosen in the order of choice fixed in
Definitions~\ref{2.3:not-stages-first}, \ref{2.3:not-stages-middle}
and~\ref{2.3:not-stages-last}. Further:
\begin{itemize}
\item $r>0$ is the radius of the correlation theorem, in which the source parameter $\zeta$
ranges over the disc $D(0,2r)$, and $g\le\gamma_J$;
\item $K=K(g_0,g)$ is the first-passage number;
\item $m$ satisfies the lower bound on the torus exponent under which the correlation theorem
is stated, and $K':=K-m$.
\end{itemize}
In the application made in this chapter, $m:=s$, so that $K'=K-s=k_0$.

Every renormalisation step considered is carried out on the whole lattice, on the physical
torus, at levels $\le K'$. Concretely:
\begin{itemize}
\item the nested domains $\Omega_i\supset\Lambda_i$ on which Ba{\l}aban's level-$i$ density is
defined are the whole torus;
\item the partition functions $\varphi_j$ that localise the fluctuation integrals are
identically $1$;
\item there is no large-field region $Z$ and no rim;
\item the boundary terms, and the summation of boundary terms at finer zones, do not occur.
\end{itemize}
By \cite{B-CMP119}, the terms of items (b) and (c) of the decomposition to which these
estimates are applied are, in this setting, of this whole-lattice form. At step $k$, $k<K'$,
the units that enter are of level $j\le k$. For such $j$ we put
\begin{equation*}
t:=L^{j-k},\qquad \nu:=1+k-j,
\end{equation*}
and there is one zone, $n=k$.

\emph{Fluctuation cut.} $\delta_k:=g_kT_k$, where $T_k$ is the cut on the real fluctuation
fields in the setting of the correlation theorem.

\emph{The two systems of radii.}
\begin{itemize}
\item[($\flat$)] Shrinking radii. At level $j$ the radii $a_j=(\alpha_{i,j})_i$ are the
level-dependent radii $\alpha_{i,j}=\alpha_i(g_j)$ of the correlation theorem's setting,
and the minimal radius is
\begin{equation*}
\alpha_{*,j}:=g_jC_*(\log x_j)^q ,
\end{equation*}
where $C_*$ is the constant of the minimal radius in that setting.
\item[($\circ$)] Absolute radii. The radii $a_j=(\alpha_i)_i$ do not depend on $j$, and
$\alpha_*:=\min\{\alpha_0,\alpha_1\}$.
\end{itemize}
The regularity sizes are $\mathfrak{a}_j:=C_{\mathfrak{a}}\alpha_{*,j}$ and
$\mathfrak{a}_{\circ}:=C_{\mathfrak{a}}\alpha_*$, where $C_{\mathfrak{a}}$ is the constant of
condition (c1) in the proof of the correlation theorem.

\emph{Cut-to-radius ratios.} We put
\begin{equation*}
\theta_k:=\frac{K_R\,\delta_k}{\alpha_{*,k}},\qquad
\theta^{\circ}_k:=\frac{K_R\,\delta_k}{\alpha_*},\qquad
\hat\theta_k:=\max_{i\le k}\theta_i,\qquad
\hat\theta^{\circ}_k:=\max_{i\le k}\theta^{\circ}_i .
\end{equation*}
Here $K_R$ is the constant of the homogeneous containment lemma in the proof of the correlation
theorem, described below.
The ratio $\theta_k$ is the ratio $\theta$ of the correlation theorem's setting, where it is
written $\theta_k=C_{\theta}(\log x_k)^{p_0-q}$ with $C_{\theta}:=K_RA_1/C_*$.

\emph{Room of the shift disc.} Consider real fluctuation fields $B$ with $|B|<T_k$. The room of
the disc of complex shifts in the local fluctuation step of the correlation theorem's induction
is the radius supplied by the homogeneous containment
lemma, which bounds Ba{\l}aban's analyticity containment by one number, homogeneous
of degree one in the field. That radius is
\begin{equation}\label{8.5:eq-field-radii-1}
\mathsf{w}_k:=\frac{1}{\theta_k}-1\quad\text{for the system }(\flat),\qquad
\mathsf{w}^{\circ}_k:=\frac{1}{\theta^{\circ}_k}-1\quad\text{for the system }(\circ).
\end{equation}
We write $\mathsf{w}$ for either of them, and $\theta$ for the corresponding ratio.

\emph{Ratios between levels.} For the system $(\flat)$ we put
$\rho_{k,j}:=\alpha_i(g_k)/\alpha_i(g_j)$ for $i=0,1$. The radii $\alpha_{0,j}$ and
$\alpha_{1,j}$ are each $g_j$ times a constant times $(\log x_j)^{q_i}$, with the same exponent
$q_0=q_1=q$, so this ratio is the same for both radii and equals
$\mathfrak{a}_k/\mathfrak{a}_j$. For the system $(\circ)$ we put $\rho_{k,j}:=1$.

\emph{Level sums.} $S_I$ and $S_{\omega}$ are the level sums attached to the table of level
estimates in the proof of the correlation theorem. We also put
\begin{gather*}
S'_I:=\sup_k\sum_{\nu}\nu^3\bigl(1+\rho_{k,k+1-\nu}\bigr)L^{-\hat\alpha(\nu-1)},\qquad
S'_{\omega}:=1+\sum_{a\ge1}L^{-a/2},\\
S_I^{\circ}:=2\sum_{\nu\ge1}\nu^3L^{-\hat\alpha(\nu-1)} .
\end{gather*}
In $S'_I$ the sum runs over the values $\nu=1+k-j$ for the levels $j\le k$. Here $\hat\alpha$ is
the exponent of the factor $t^{4+\hat\alpha}$ in the bound on pointed groups in the proof of the
correlation theorem. These numbers do not
depend on $K$, as noted after the definition of the level sums in the proof of the correlation
theorem.

\emph{Constants.} The following constants are used in what follows.
\begin{itemize}
\item $C_5:=2^de^{(2^d-1)\kappa}C_S(2M)^4$, where $C_S$ is a constant of the proof of the
correlation theorem. The counting factors come from the fifth line of the table of level
estimates: a cell meets at most $2^d$ cubes
$\square$, and each such cube has $d_k(\square)\le2^d-1$ (by part~(b) of the lemma in the proof
of the correlation theorem that supplies these counts).
\item $C_6$ is the constant $C$ of the bound $CL^4r|B'|^2$ in the fifth line of that table. It
is the constant for which $|G_y(B')|\le C_6L^4\sup_{B(y)}|B'|^2$, where $B'$ is the fluctuation
field entering the vertex of the correlation theorem's expansion and $G_y(B')$ is the term
attached to $y$ there. On the support of
$\chi_k$ one has $|B'|\le C_1\delta_k$, by \cite[(1.20)]{B-CMP116}.
\item $\mathsf{K}:=K_{CE}/C_A$, which equals $8b^{\mathrm{adm}}/(\nu_0\alpha_4)$, in terms of
the constants of the marked-expansion lemma in the proof of the correlation theorem.
\item $C_U^{+}:=\max\{C_U,4C^{(2)},4C^{(3)}\}$. Here $C_U$ is the constant of the proposition
bounding the marked weight $\mathsf{W}^m$ in that proof. The constants $C^{(2)}$, $C^{(3)}$ are
the constants $C$ in the last column of the second and third lines of the table of level
estimates. Thus a single constant serves groups, singles indexed by pairs $(X,y)$, and singles
indexed by $X$.
\item $C^{(5)}$ is the constant $C$ of the bound $\le Crg_n^2(\log x_n)^{2q}$ in the last
column of the fifth line.
\end{itemize}

\emph{Absolute-radius forms.} The absolute-radius form of the homogeneous containment lemma is
the same statement with $\alpha_{*,k}$ replaced by $\alpha_*$. On the whole lattice it is the
statement of \cite[(3.14)--(3.17)]{B-CMP109} and of \cite[(1.17)--(1.22)]{B-CMP116}; the latter
states that the expression \cite[(1.10)]{B-CMP116} extends analytically in the variables
$t_{\square}$ subject to
\begin{equation*}
|t_{\square}|\,4B_0C_1e^{16\kappa_1}g_k|B|\le\tfrac12\alpha_2 ,
\end{equation*}
with the constant $B_0$ and the radius $\alpha_2$ as in \cite{B-CMP116}.
The three respects in which the correlation theorem's setting differs from that of these
papers do not arise on the whole lattice at absolute radii. These three respects are the
level-dependent radii, the layers of level $n<k$, and the rim. The smallness condition
$e^{32\kappa_1}C_1\delta_k\le c_{\mathrm{abs}}$ of the absolute-radius response theorem, with
its constant $c_{\mathrm{abs}}$, remains in force.

The absolute-radius form of the conditions (c1)--(c3) and $(\ell0)$ of the proof of the
correlation theorem is obtained by replacing $\mathfrak{a}_n$ by $\mathfrak{a}_{\circ}$, and the
radius $\alpha_{3,j}$ of those conditions by $\alpha_3:=c_3\alpha_*$, where $c_3$ is the constant
of those conditions. In this form they are
\cite[(3.18), (3.25), (3.32)]{B-CMP109}, \cite[Lemma~4]{B-CMP109}, \cite[(4.18)]{B-CMP109} and
the clause labelled (iv) in the same paper, all used as stated there.

\emph{Room for small ratios.} If $\theta\in(0,\tfrac13]$, then
\begin{equation}\label{8.5:eq-field-radii-2}
\mathsf{w}\ge2\qquad\text{and}\qquad \frac{4}{\mathsf{w}}\le6\theta\le8\theta .
\end{equation}
\end{definition}

\begin{proof}[Proof of \eqref{8.5:eq-field-radii-2}]
Let $\theta\in(0,\tfrac13]$. By \eqref{8.5:eq-field-radii-1},
\begin{equation*}
\mathsf{w}=\frac{1}{\theta}-1=\frac{1-\theta}{\theta} .
\end{equation*}
Since $1/\theta\ge3$, we get $\mathsf{w}\ge3-1=2$, which is the first inequality in
\eqref{8.5:eq-field-radii-2}.

Since $\theta\le\tfrac13$, we also have $1-\theta\ge\tfrac23$. Therefore
$\mathsf{w}\ge 2/(3\theta)$, and
\begin{equation*}
\frac{4}{\mathsf{w}}\le 4\cdot\frac{3\theta}{2}=6\theta\le8\theta ,
\end{equation*}
the last step because $\theta>0$. This is the second inequality in
\eqref{8.5:eq-field-radii-2}.

The argument applies verbatim to $\theta=\theta_k$ with $\mathsf{w}=\mathsf{w}_k$, and to
$\theta=\theta^{\circ}_k$ with $\mathsf{w}=\mathsf{w}^{\circ}_k$.
\end{proof}

\begin{lemma}[The unit-sum bound for a marked family]\label{8.5:lem-unit-sum-bound}
Fix a step $k$ and one of the two radius systems $a\in\{\flat,\circ\}$ of
Notation~\ref{8.5:not-field-radii}, with the field radii $a_j$, the ratios $\rho_{k,j}$, the
regularity sizes $\mathfrak{a}_k$, $\mathfrak{a}_{\circ}$, the room $\mathsf{w}_k$,
$\mathsf{w}^{\circ}_k$ and the constants $C_U^{+}$, $C_5$ fixed there. Write $\theta_a:=\theta_k$
and $\mathsf{w}_a:=\mathsf{w}_k$ if $a=\flat$, and $\theta_a:=\theta^{\circ}_k$ and
$\mathsf{w}_a:=\mathsf{w}^{\circ}_k$ if $a=\circ$, and assume $\theta_a\le 1/3$. For a level
$j\le k$ put $\nu:=1+k-j$ and $t:=L^{j-k}$, and write $\hat\alpha$ for the exponent in the group
bound of the correlation theorem. Let $\mathcal{M}$ be a finite family of marked
units of the following three kinds.
\begin{itemize}
\item[(G)] \emph{Groups.} For each pair $(j,y)$ with $j\le k$ and $y\in\mathbb{T}^{(j)}$,
at most one group $c\,\mathfrak{i}_y[F^{(j)}]$ with $|c|\le 1$, where $F^{(j)}$ is a pointed
$E$-family at scale $j$ on the radii $a_j$ with norm at most $\mathfrak{n}^G_j$, presented as
the table of unit bounds in the proof of the correlation theorem presents groups: the core
and its counterterm piece on one background, the tails singly, in accordance with the
hypotheses (Hn), (Hf) of the marked-expansion lemma.
\item[(M)] \emph{Singles.} For each triple $(j,X,y)$ (or for each pair $(j,X)$) at most one
single $c\,[H(X,\cdot)-H(X,1)]$ with $|c|\le 1$, where $H(X,\cdot)$ satisfies the two conditions
(f1), (f2) imposed on singles in the single bound of the correlation theorem, on the complex
background space $\mathcal{U}^{c}_j(X,a_j)$ at the radii $a_j$, and
$|H(X,\cdot)|\le\mathfrak{n}^M_j\,e^{-(1-\frac14\delta)\kappa d_j(X)}$.
\item[(W)] \emph{Wilson pieces.} For each shift cube $\square$ at most one Wilson piece
$A(c\,\hat\chi_{\square},\cdot)$ with $\sup|c|\le\bar c$.
\end{itemize}
Then every unit of $\mathcal{M}$ satisfies the domain condition (D) of the marked-expansion
lemma on the whole output space, and
\begin{equation}\label{8.5:eq-unit-sum-bound-a}
\begin{split}
W^m(\mathcal{M})\le{}& 2C_U^{+}\theta_a\Big[\sum_{j\le k}\nu^3\rho_{k,j}^2(1+\rho_{k,j})
L^{-\hat\alpha(\nu-1)}\mathfrak{n}^G_j+\sum_{j\le k}\rho_{k,j}^2\mathfrak{n}^M_j\Big]\\
&+8\theta_aC_5\bar c\,\mathfrak{a}^2,
\end{split}
\end{equation}
where $\mathfrak{a}=\mathfrak{a}_k$ if $a=\flat$ and $\mathfrak{a}=\mathfrak{a}_{\circ}$ if
$a=\circ$.
\end{lemma}

\begin{proof}
Recall that a marked unit $v$ carries a level $j_v$, a support $S_v$ and a functional $f_v$,
that $\omega(v)$ denotes its oscillation and $\mathfrak{N}(v)$ its size, and that the marked
weight of $\mathcal{M}$ is
\begin{equation}\label{8.5:eq-unit-sum-bound-1}
W^m(\mathcal{M})=\sup_{\Delta}\sum_{v\in\mathcal{M}:\,S_v\cap\Delta\neq\emptyset}
\omega(v)\,e^{(1-\frac12\delta)\kappa d_k(\bar S_v)},
\end{equation}
the supremum running over the current cells $\Delta$. The proof has four steps.

\emph{Step one: the domain condition.} Condition (D) of the marked-expansion lemma requires
\begin{equation}\label{8.5:eq-unit-sum-bound-2}
W\exp(i\xi w\,\delta\mathfrak{v})\in\mathfrak{D}_v\qquad
\text{for all }(W,\delta\mathfrak{v})\in\mathfrak{E}_v\text{ and all }|w|\le 2,
\end{equation}
where $\mathfrak{D}_v$ is the domain and $\mathfrak{E}_v$ the set of pairs $(W,\delta\mathfrak{v})$
attached to the unit $v$ in the marked-expansion lemma, and $\xi$ is as there.
By the homogeneity lemma for the containment condition (for $a=\circ$, by its reading at
absolute radii, in which $\alpha_{*,k}$ is replaced by $\alpha_*$), the response maps
$\mathsf{R}_p(\mathsf{t},\sigma)$ of that lemma, in the shift parameter $\mathsf{t}$ with weight
$\eta_p$ and the cell parameters $\sigma(\Delta)$, take values in admissible backgrounds whenever
$\eta_p|\mathsf{t}|\le\mathsf{w}_a$, $|\sigma(\Delta)|\le e^{\kappa_1}$, the fluctuation field
$B$ is real with $|B|<T_k$, and the output variables $(\mathbf{U},\mathbf{J})$ lie in the output
tube of the step. Since
$\theta_a\le 1/3$,
\begin{equation*}
\mathsf{w}_a=\theta_a^{-1}-1\ge 3-1=2 .
\end{equation*}
Hence the region on which the response maps land in admissible backgrounds contains the
region required in \eqref{8.5:eq-unit-sum-bound-2}, which asks only $|w|\le 2$. This proves
that (D) holds for every unit of $\mathcal{M}$ on the whole output space.

\emph{Step two: the unit sizes.} On all admissible backgrounds
$|f_v-f_v(1)|\le\mathfrak{N}(v)$, where $\mathfrak{N}(v)$ is as follows. For a core with its
counterterm piece, by the group bound, with $C_I$ the constant of that bound,
\begin{equation*}
\mathfrak{N}(v)=C_I\,\mathfrak{n}^G_j\,\nu^3\rho_{k,j}^2(1+\rho_{k,j})\,t^{4+\hat\alpha}.
\end{equation*}
For a tail, by the tail bound, $\mathfrak{N}(v)=2\mathfrak{n}^G_j\,e^{-\kappa d_j(X)}$. For a
single, by part (i) of the single bound, and by its part (ii) beyond the region
$\tilde\square^{2}$ of that bound,
\begin{equation*}
\mathfrak{N}(v)=2\mathfrak{n}^M_j\big(C_MD_Xt^2\rho_{k,j}\big)^2
e^{-(1-\frac14\delta)\kappa d_j(X)}.
\end{equation*}
For a Wilson piece, by part (b) of the Wilson-piece bound, with $C_S$ the constant there,
$\mathfrak{N}(v)=C_S(2M)^4\bar c\,\mathfrak{a}^2$. The group bound, the tail bound, the single
bound and the
Wilson-piece bound use the background only through the regularity conditions (c1)--(c3),
which at $a=\circ$ are read with $\mathfrak{a}_n$ replaced by $\mathfrak{a}_{\circ}$; they
therefore hold in either radius system.

\emph{Step three: the oscillations.} The Schwarz lemma on the shift disc, in the form of the
closing remark of the oscillation lemma for marked units, gives
$\omega(v)\le 4\mathfrak{N}(v)/\mathsf{w}_a$. Since $\mathsf{w}_a=(1-\theta_a)/\theta_a$ and
$1-\theta_a\ge 2/3$,
\begin{equation}\label{8.5:eq-unit-sum-bound-3}
\omega(v)\le\frac{4\mathfrak{N}(v)}{\mathsf{w}_a}
=\frac{4\theta_a}{1-\theta_a}\,\mathfrak{N}(v)\le 6\theta_a\mathfrak{N}(v)
\le 8\theta_a\mathfrak{N}(v).
\end{equation}

\emph{Step four: the sum over one cell.} Inserting \eqref{8.5:eq-unit-sum-bound-3} into
\eqref{8.5:eq-unit-sum-bound-1} gives
\begin{equation}\label{8.5:eq-unit-sum-bound-4}
W^m(\mathcal{M})\le 8\theta_a\sup_{\Delta}\sum_{v\in\mathcal{M}:\,S_v\cap\Delta\neq\emptyset}
\mathfrak{N}(v)\,e^{(1-\frac12\delta)\kappa d_k(\bar S_v)} .
\end{equation}
The sum on the right is counted as follows. A current cell contains $t^{-4}$ anchors $y$ of
level $j$, and this factor is set against $t^4$. For a core the weight is at most
$e^{(2^d-1)\kappa}$, and $t^{-4}\,t^{4+\hat\alpha}=t^{\hat\alpha}=L^{-\hat\alpha(\nu-1)}$; this
is the source of the factor $\nu^3\rho_{k,j}^2(1+\rho_{k,j})L^{-\hat\alpha(\nu-1)}$ in
\eqref{8.5:eq-unit-sum-bound-a}. For singles, $t^{-4}(t^2)^2=1$, and the count is
\begin{equation*}
\sum_{X\ni y}\big(d_j(X)+5\big)^2e^{-\frac14\delta\kappa d_j(X)},
\end{equation*}
the diameter factor $D_X^2$ and the decay $e^{-(1-\frac14\delta)\kappa d_j(X)}$ being set
against the weight as in the unit-sum proposition; the remaining factor is $\rho_{k,j}^2$.
For the tails, the $t^{-4}$ anchors per cell are set against a factor $t^{5}$ in place of
$t^{4}$. This count
is the one that defines the constant $C_U$ in the unit-sum proposition of the correlation
theorem, where groups are bounded by the group bound and singles by the single bound, and it
is the one in the last column of the entries for $(X,y)$-indexed and for $X$-indexed singles
in the table of unit bounds, with constants $C^{(2)}$ and $C^{(3)}$ there; it does not depend
on the radius system. In the unit-sum proposition the oscillation enters through the factor
$4\theta_k$ and the result carries the constant $C_U\theta_k$. Here
\eqref{8.5:eq-unit-sum-bound-3} gives $8\theta_a$ in place of $4\theta_k$, so the groups
contribute at most
\begin{equation*}
2C_U\theta_a\sum_{j\le k}\nu^3\rho_{k,j}^2(1+\rho_{k,j})L^{-\hat\alpha(\nu-1)}\mathfrak{n}^G_j,
\end{equation*}
and the singles are bounded by the same count, with the constants $C^{(2)}$, $C^{(3)}$ of
those entries according to the indexing and with the oscillation factor $8\theta_a$, hence by
\begin{equation*}
2C_U^{+}\theta_a\sum_{j\le k}\rho_{k,j}^2\mathfrak{n}^M_j .
\end{equation*}
Since
$C_U^{+}=\max\{C_U,4C^{(2)},4C^{(3)}\}$, one constant serves groups, $(X,y)$-indexed singles
and $X$-indexed singles, and together these contributions are bounded by the first line of
\eqref{8.5:eq-unit-sum-bound-a}.

For the Wilson pieces, a current cell meets at most $2^d$ shift cubes $\square$, and
$d_k(\square)\le 2^d-1$, so their weights are at most $e^{(2^d-1)\kappa}$. By
\eqref{8.5:eq-unit-sum-bound-4} and step two their contribution is at most
\begin{equation*}
8\theta_a\,2^de^{(2^d-1)\kappa}C_S(2M)^4\bar c\,\mathfrak{a}^2=8\theta_aC_5\bar c\,\mathfrak{a}^2,
\end{equation*}
because $C_5=2^de^{(2^d-1)\kappa}C_S(2M)^4$. Adding the three contributions gives
\eqref{8.5:eq-unit-sum-bound-a}.
\end{proof}

\begin{lemma}[Marked potentials given directly]\label{8.5:lem-direct-potentials-bis}
Let the unmarked potential $V^u$ and a finite marked family $\mathcal{M}$ of weight $W^m$ satisfy
the hypotheses of the marked-expansion lemma at step $k$, and let $V^{\mathrm{dir}}$ be a marked
potential that is given directly rather than built from $\mathcal{M}$. Assume that
$V^{\mathrm{dir}}$ satisfies the holomorphy hypotheses which that lemma imposes on its marked input,
and that
\begin{equation}\label{8.5:eq-direct-potentials-1-bis}
|V^{\mathrm{dir}}(Y)|\le \tilde a^{\mathrm{dir}}\,e^{-(1-2\delta)\kappa d_k(Y)}
\end{equation}
for every $Y$ over which the marked-expansion lemma states its activity hypothesis, with the
fraction $\delta$ of $\kappa$ used in that lemma (not to be confused with the fluctuation cut
$\delta_k$ of~\eqref{8.5:eq-direct-potentials-2-bis}). Let $V^m$ be the sum of
the marked potential built from $\mathcal{M}$ and
$V^{\mathrm{dir}}$, and let $t_\lambda(X)$, $t^{[1]}(X)$ and $t_1-t_0$ be formed as in the
marked-expansion lemma from $V=V^u+\lambda V^m$. Write $\mathsf{K}:=K_{CE}/C_A$, and set
$\tilde a$
and $\lambda_0$ by the first line of~\eqref{8.5:eq-direct-potentials-a-bis}. Then the conclusions of
the marked-expansion lemma hold with $\tilde a$ in place of $C_AW^m$; in particular $t_\lambda(X)$
is analytic on $|\lambda|<\lambda_0$. If moreover $\tilde a\le \nu_0\alpha_4/8$, the second line
of the following display holds:
\begin{equation}\label{8.5:eq-direct-potentials-a-bis}
\begin{aligned}
\tilde a&:=C_A W^m+\tilde a^{\mathrm{dir}},\qquad
\lambda_0:=\frac{\nu_0\alpha_4}{4\tilde a},\\
|t^{[1]}(X)|&\le \tfrac12\,\mathsf{K}\,\tilde a\,e^{-\kappa d_{k+1}(X)},\qquad
|t_1-t_0|(X)\le \mathsf{K}\,\tilde a\,e^{-\kappa d_{k+1}(X)} .
\end{aligned}
\end{equation}

The gauge-fixing pieces $c(y)G_y(B')$ of step $k$, indexed by blocks $y$, with $|c|\le\bar c$,
where $G_y$ is a polynomial in the configuration $B'=C\mathsf{A}-h\tilde D(C\mathsf{A})$ of
Ba{\l}aban's construction, are marked potentials of this kind: each
piece is localised in the cube of Ba{\l}aban's partition $\pi_k$ that contains the block $y$ and is
invariant under the transformations of
\cite[(3.29)]{B-CMP109}. For each cube $Y$ of $\pi_k$, the constant $\tilde a^{\mathrm{dir}}$
of~\eqref{8.5:eq-direct-potentials-1-bis} for the pieces localised in $Y$ satisfies
\begin{equation}\label{8.5:eq-direct-potentials-2-bis}
\tilde a^{\mathrm{dir}}\le C_6 M^4\,\bar c\,C_1^2\,\delta_k^2 .
\end{equation}
\end{lemma}

\begin{proof}
The marked-expansion lemma uses only two properties of its marked input: a bound of $|V^m(Y)|$
by a constant, the activity of $V^m$, times $e^{-(1-2\delta)\kappa d_k(Y)}$, and the holomorphy
hypotheses. The marked
sub-sum lemma that accompanies the correlation theorem is stated in exactly this form: if
$V=V^u+\lambda V^m$, the marked sub-sum
$t_1-t_0$ is bounded by $K_{CE}$ times the activity of $V^m$, divided by $C_A$, times
$e^{-\kappa d_{k+1}(X)}$. Its proof uses analyticity in $\lambda$: the integrand is an entire
function of $\lambda$, and the Schwarz estimate is applied on the disc of radius
$\nu_0\alpha_4/4$ divided by the activity of $V^m$. In that lemma the marked potential is
built from
$\mathcal{M}$ and its activity is $C_AW^m$, whence the radius
$\nu_0\alpha_4/(4C_AW^m)$.

Now let $V^m$ contain in addition $V^{\mathrm{dir}}$. The sum of two potentials satisfying the
holomorphy hypotheses satisfies them, and by the triangle inequality together
with~\eqref{8.5:eq-direct-potentials-1-bis} the activity of the sum is at most
$C_AW^m+\tilde a^{\mathrm{dir}}=\tilde a$. Since nothing else about the marked input enters, the
marked-expansion lemma applies with the activity $\tilde a$ in place of $C_AW^m$.
Thus $t_\lambda(X)$ is analytic
in $|\lambda|<\lambda_0$ with $\lambda_0$ as in~\eqref{8.5:eq-direct-potentials-a-bis}, and
$|t^{[m]}(X)|\le b^{\mathrm{adm}}\lambda_0^{-m}e^{-\kappa d_{k+1}(X)}$. Here $t^{[m]}$ is the
sub-sum carrying exactly $m$ marked pieces. With $m=1$ and $\lambda_0^{-1}=4\tilde a/(\nu_0\alpha_4)$
this gives
\begin{equation*}
|t^{[1]}(X)|\le \frac{4b^{\mathrm{adm}}}{\nu_0\alpha_4}\,\tilde a\,e^{-\kappa d_{k+1}(X)}
=\tfrac12\,\mathsf{K}\,\tilde a\,e^{-\kappa d_{k+1}(X)},
\end{equation*}
because, by the definition of $K_{CE}$ in the marked-expansion lemma,
$K_{CE}=8b^{\mathrm{adm}}C_A/(\nu_0\alpha_4)$, so that
$\mathsf{K}=8b^{\mathrm{adm}}/(\nu_0\alpha_4)$.

For the difference, each marked piece carries one factor $\lambda$, so
$t_\lambda=\sum_{m\ge0}\lambda^m t^{[m]}$ on the disc. If $\tilde a\le\nu_0\alpha_4/8$, then
$\lambda_0\ge2$, so $\lambda=1$ lies in the disc and $\sum_{m\ge1}\lambda_0^{-m}\le 2\lambda_0^{-1}$.
Hence
\begin{equation*}
|t_1-t_0|(X)\le\sum_{m\ge1}|t^{[m]}(X)|
\le 2b^{\mathrm{adm}}\lambda_0^{-1}e^{-\kappa d_{k+1}(X)}
=\mathsf{K}\,\tilde a\,e^{-\kappa d_{k+1}(X)} .
\end{equation*}
This is also the bound of the marked sub-sum lemma with the activity $\tilde a$ in place of
$C_AW^m$.

Each gauge-fixing piece is a polynomial in $B'$, localised in one cube of $\pi_k$ and invariant
under the transformations of \cite[(3.29)]{B-CMP109}; these properties make the pieces marked
potentials given directly in the sense of the first part of the lemma, and
\eqref{8.5:eq-direct-potentials-2-bis} is the bound on their activity over one cube $Y$, which enters
as $\tilde a^{\mathrm{dir}}$.
\end{proof}

In the following proposition the symbols $K'$, $D_j$, $X^{+}$, $E_0$, $\kappa_0$, $C_B$,
$C_{\mathfrak{m}}$, $T_k$, $C^{(k)}$, $V^{(k)}$, $B'$, $G_y$, $\Psi_i$, $\Phi_i$, $J(U)$, $M^j$,
$b_j'(0)$ and the renormalisation coefficient $\beta[\cdot]$ have the meaning they have in the proof
of the correlation theorem. So do the symbols $k_1(k)$, $\omega_{k,j}$, $w_{k,j}$ and
$\hat\alpha_k$ used in the proof.

\begin{proposition}[Splitting the sourced terms on shrinking radii]\label{8.5:prop-local-source-split}
Work with the shrinking radii $\flat$ of Notation~\ref{8.5:not-field-radii}. Assume that the
inductive statement in the proof of the correlation theorem holds at every level $k\le K'$.
Assume also the two ceilings below. In them, $\theta(g')$ and $\mathfrak{m}(g')$ denote the ratio
$\theta_k$ of Notation~\ref{8.5:not-field-radii} and the bound $\mathfrak{m}_k$, read as functions
of the coupling $g'=g_k$:
\begin{equation}\label{8.5:eq-local-source-split-a}
\begin{gathered}
\sup_{g'\le g}2K_{CE}\,C_U^{+}\,\theta(g')\bigl(1+S_I+C_BS_{\omega}\bigr)\le\tfrac12,
\qquad \theta(g)\le\tfrac13,\\
C_{\mathfrak{m}}\,(g^{-2}-4c_2)^{-(\kappa_0-2)/2}\le b^{\mathrm{adm}},
\end{gathered}
\end{equation}
\begin{equation}\label{8.5:eq-local-source-split-b}
\sup_{g'\le g}8C_A\,\theta(g')\Bigl[2\bar C_{\mathsf{W}}+C^{(3)}\mathfrak{m}(g')
+C^{(5)}r\,g'^2\bigl(\log g'^{-2}\bigr)^{2q}\Bigr]\le\tfrac14\,\nu_0\alpha_4 .
\end{equation}
Then for every $1\le j\le K'$, every $X\in D_j$ and every $y\in X$ there are functions
$P^{(j)}(X,\cdot,y;w)$ and $Q^{(j)}(X,\cdot,y;w)$ on
$\mathcal{U}^{\flat}_j(X)\times\mathcal{W}_j(X)$, the dot standing for the pair
$(\mathbf{U},\mathbf{J})$, with the following properties.
\begin{enumerate}
\item[(i)] $E^{(j)}(X,\cdot,y;w)=P^{(j)}(X,\cdot,y;w)+Q^{(j)}(X,\cdot,y;w)$.
\item[(ii)] $P^{(j)}$ and $Q^{(j)}$ are analytic, invariant under the symmetries of
\cite[(1.19)]{B-CMP109}, and depend on $(\mathbf{U},\mathbf{J})$ only through its restriction
to $X$ and on $w$ only through its restriction to $X^{+}$.
\item[(iii)] If $w\equiv\hat w$ on $X^{+}$ for a constant $\hat w$, then
$P^{(j)}(X,\cdot,y;w)=P^{\flat(j)}_{\hat w}(X,\cdot,y)$ and
$Q^{(j)}(X,\cdot,y;w)=Q^{(j)}_{\hat w}(X,\cdot,y)$.
\item[(iv)] The bounds
\begin{equation}\label{8.5:eq-local-source-split-d}
\bigl|Q^{(j)}(X,\cdot,y;w)\bigr|\le\mathfrak{m}_{j-1}\,e^{-\kappa d_j(X)},
\qquad
\bigl|P^{(j)}(X,\cdot,y;w)\bigr|\le 2E_0\,e^{-\kappa d_j(X)}
\end{equation}
hold.
\item[(v)] Let $U=U_j(W)$ be a real minimal configuration of level $j$, with $W$ real and
regular on $\mathbb{T}^{(j)}$. This class contains $U_n(V)$ for every $n\ge j$, since
$U_n=U_j(M^j(U_n))$ modulo gauge \cite[(3.3)]{B-CMP109}; this is also stated in the correlation
theorem's description of minimal configurations. Let $g$ be a Lipschitz function from
the plaquettes to $\mathbb{C}$ (here and in the recursion below $g$ denotes this function on
plaquettes, not the coupling). Then
\begin{equation}\label{8.5:eq-local-source-split-1}
\partial_z\Big|_{z=0}\sum_{X,y}P^{(j)}\bigl(X,U,J(U),y;zg\bigr)
=\dot{\mathcal{T}}^{(j)}[g](U)+\mathrm{const},
\end{equation}
where the constant does not depend on $U$ and $\dot{\mathcal{T}}^{(j)}[g]$ is defined by the
recursion below.
\end{enumerate}
The recursion is as follows. Put $\dot{\mathcal{T}}^{(0)}[g]:=0$ and, for $k\ge0$,
\begin{gather*}
\dot\zeta_k[g]:=g+\sum_{j\le k}b_j'(0)\,I_jg,\qquad
\mathsf{b}^P_j:=\partial_{\zeta}\beta\bigl[P^{\flat(j)}_{\zeta}\bigr]\Big|_{\zeta=0},\\
\bigl(\dot{\mathcal{T}}^{(j)}[g]\bigr)^{\mathrm{ren}}
:=\bigl(\dot{\mathcal{T}}^{(j)}[g]\bigr)^{\mathrm{vac}}-\mathsf{b}^P_j\,A(I_jg,\cdot),
\end{gather*}
with $A(f,\cdot)$ the weighted Wilson action of Definition~\ref{1.1:def-wilson-action}, and
\begin{equation}\label{8.5:eq-local-source-split-c}
\begin{split}
\dot{\mathcal{T}}^{(k+1)}[g](U)&:=\bigl\langle\dot{\mathfrak{B}}_k[g]\bigr\rangle^U_k,\\
\dot{\mathfrak{B}}_k[g]&:=\Bigl\{A(\dot\zeta_k[g],\cdot)
+\sum_{j\le k}\bigl(\dot{\mathcal{T}}^{(j)}[g]\bigr)^{\mathrm{ren}}\Bigr\}
\bigl(U_k(e^{iB'}V^{(k)})\bigr)\\
&\qquad-\Bigl\{A(\dot\zeta_k[g],\cdot)
+\sum_{j\le k}\bigl(\dot{\mathcal{T}}^{(j)}[g]\bigr)^{\mathrm{ren}}\Bigr\}(U_{k+1})
+\sum_{y\in\mathbb{T}^{(k+1)}}\dot\zeta_k[g](p_y)\,G_y(B').
\end{split}
\end{equation}
Here $\langle\cdot\rangle^U_k$ is the normalised measure proportional to
\begin{equation}\label{8.5:eq-local-source-split-2}
d\mu_{C^{(k)}}(B)\prod_b\chi\bigl(|B(b)|<T_k\bigr)
\exp\Bigl[g_k^{-2}\sum\Psi_i+\sum\Phi_i
+\bigl\{\mathsf{P}_0^{\mathrm{ren}}\bigl(U_k(e^{iB'}V^{(k)})\bigr)
-\mathsf{P}_0^{\mathrm{ren}}(U_{k+1})\bigr\}\Bigr]
\end{equation}
at $U_{k+1}=U$, with $\mathsf{P}_0:=(\mathsf{P}^{(j)}_{\xi=0})_{j\le k}$. This is the measure of
part~(d) of the proposition on the total of the potentials of the fluctuation step, taken at the
parameter value $s=1$.
\end{proposition}

Here $\theta(g')$ is the ratio $\theta_k$ of Notation~\ref{8.5:not-field-radii} read as a
function of the coupling: $\theta_k=\theta(g_k)=C_{\theta}(\log x_k)^{\mathfrak{p}_0-q}$. By the
condition relating the exponents $\mathfrak{p}_0$ and $q$ among the constants of the correlation
theorem, $\theta(g')$ decreases to zero as $g'\downarrow0$. In the standing setting of the
correlation theorem, within which the proposition is stated, the couplings $g_k$, $k\le K'$,
satisfy $g_k\le g$. Hence $\theta_k\le\theta(g)$ for every $k\le K'$.

Since $C_U^{+}\ge C_U$ and $C_BS_{\omega}\ge0$, the first condition in
\eqref{8.5:eq-local-source-split-a} implies $K_{CE}C_U\theta_k(1+S_I)\le\frac12$. This is the
smallness condition of the proposition on the $Q$-families in the proof of the correlation
theorem (below: the $Q$-family proposition).

The third condition in \eqref{8.5:eq-local-source-split-a} ensures that
$\tilde a\le\nu_0\alpha_4/8$, that is $\lambda_0\ge2$, in the corollary of the marked-expansion
lemma. Here $\tilde a$ and $\lambda_0$ are as in Lemma~\ref{8.5:lem-direct-potentials-bis}.

Condition \eqref{8.5:eq-local-source-split-b} is the smallness condition that the marked-expansion
lemma imposes on the unmarked input, here the local unmarked input of the proof below. The constant
$\bar C_{\mathsf{W}}$ is fixed among the constants of the correlation theorem: it bounds the
unit-sum weight $\mathsf{W}$ of the unmarked potentials of the local step, whose families there
have norm at most $E_0$. It enters doubled because $|P^{(j)}|\le 2E_0$.

\begin{proof}
We argue by induction on $j=k+1$. We assume (i)--(v) at all levels $\le k$ and prove them at
level $k+1$. For $k=0$ the inductive statement of the correlation theorem at level $0$ is empty,
and there are no levels $j\le0$ to which (i)--(v) refer.

\emph{The output kernel and its inputs.} We use the local step statement in the proof of the
correlation theorem: the step at level $k$ performs the same operations as \cite{B-CMP119}. By it,
$E^{(k+1)}(X,\cdot,y;w)$ is the output kernel, fixed as in \cite[(3.47)]{B-CMP119}, of the last
logarithm of \cite[(3.28)]{B-CMP119}. That logarithm is taken of an integral whose exponent is
built from three inputs:
\begin{itemize}
\item[(a)] the potentials of the step at level $k$ that do not depend on the source variable $z$;
we write them $\mathcal{P}^{(k)}$, to distinguish them from the families $P^{(j)}(w)$
constructed here;
\item[(b)] the vertex $\mathfrak{S}_k$ of the correlation theorem's three-family split, with
$\zeta_k=\zeta^w_k$ the sourced shift of the correlation theorem, together with the gauge-fixing
pieces $\zeta_k(p_y)G_y$;
\item[(c)] the differences between the values at $U_k(e^{iB'}V^{(k)})$ and at $U_{k+1}$ of
\begin{equation*}
\sum_{j\le k}\bigl(E^{(j)}(w)\bigr)^{\mathrm{ren},\mathrm{loc}}
+\sum_{j\le k_1(k)}\bigl(R^{(j)}(w)\bigr)^{\mathrm{vac}},
\end{equation*}
where
\begin{equation*}
\bigl(E^{(j)}(w)\bigr)^{\mathrm{ren},\mathrm{loc}}
=\sum_y\mathfrak{i}_y\bigl[E^{(j)}_{w(p_y)}\bigr]+\sum\bigl(D^{(j)}\bigr)^{\mathrm{vac}};
\end{equation*}
the first sum consists of the renormalised pointed groups and the second of the difference
family of the sourced format.
\end{itemize}

\emph{Splitting the inputs.} At the levels $j\le k$ we use (i)--(iii), which hold by the
induction hypothesis. For a constant source value $\hat w$, (i) and (iii) give
$E^{(j)}_{\hat w}=P^{\flat(j)}_{\hat w}+Q^{(j)}_{\hat w}$. The renormalisation coefficient $\beta$
is linear, so $(E^{(j)})^{\mathrm{ren}}=(P^{\flat(j)})^{\mathrm{ren}}+(Q^{(j)})^{\mathrm{ren}}$.
Hence
\begin{equation*}
\mathfrak{i}_y[E_{\hat w}]=\mathfrak{i}_y[P^{\flat}_{\hat w}]+\mathfrak{i}_y[Q_{\hat w}],
\end{equation*}
and each of the two groups on the right is renormalised with its own coefficient $\beta$. This
is the device of the $Q$-family proposition. In the same way the difference family splits as
$D^{(j)}=D^P+D^Q$, where
\begin{equation*}
D^P:=P^{(j)}(w)-P^{\flat}_{w(p_y)},\qquad D^Q:=Q^{(j)}(w)-Q_{w(p_y)}.
\end{equation*}

\emph{The marking.} We mark the following pieces as units in the sense of
Lemma~\ref{8.5:lem-unit-sum-bound}. Both families of singles below are units of the kind listed
as (M) in that lemma.
\begin{itemize}
\item[(G)] The groups $\mathfrak{i}_y[Q^{(j)}_{w(p_y)}]$, with norm
$\mathfrak{n}^G_j=\mathfrak{m}_{j-1}$. This is the bound on the $Q$-families contained in the
inductive statement of the correlation theorem.
\item[(M$_1$)] The singles $(D^Q)^{\mathrm{vac}}$. The function $Q^{(j)}(X,\cdot,y;\cdot)$ is
analytic on $\mathcal{W}_j(X)$ and bounded by $\mathfrak{m}_{j-1}e^{-\kappa d_j(X)}$. By the lemma
on the variation of a family in the source, therefore,
\begin{equation*}
|D^Q|\le2\mathfrak{m}_{j-1}\,\omega_{k,j}\Bigl(1+\frac{\operatorname{diam}_j(X^{+})}{10M}\Bigr)
e^{-\kappa d_j(X)}
\le C_B\,\mathfrak{m}_{j-1}\,\omega_{k,j}\,e^{-(1-\frac14\delta)\kappa d_j(X)} .
\end{equation*}
So these singles have norm $\mathfrak{n}^M_j=C_B\,\omega_{k,j}\,\mathfrak{m}_{j-1}$.
\item[(M$_2$)] The singles $(R^{(j)}(X,\cdot;w))^{\mathrm{vac}}$ with $j\le k_1(k)$. Their norm
is $\mathfrak{n}^M_j=g_j^{\kappa_0}$, by the remainder bound of the sourced format. They are read
on the radius $\hat\alpha_k$, with the weight $w_{k,j}=(\hat\alpha_k/\alpha_j)^2$ of the
$Q$-family proposition, as in the closing remark of the background-comparison lemma of the
correlation theorem.
\end{itemize}
All remaining pieces are unmarked. These are the $z$-free potentials $\mathcal{P}^{(k)}$, the
vertex with its gauge-fixing pieces, the groups $\mathfrak{i}_y[P^{\flat}_{w(p_y)}]$ (for which
$\|P^{\flat}\|^{\flat}\le E_0$), and the singles $(D^P)^{\mathrm{vac}}$. The bound on the latter
follows from the lemma on the variation of a family in the source together with $|P|\le2E_0$.

The unmarked input is in the format required by the marked-expansion lemma. Indeed, it consists
of the potentials of the local step statement with $E$ replaced by $P$ in the groups and in the
singles. The lemma that estimates the potentials generated by a unit is linear in the unit
$f_v$ and depends on a unit only through its size $\mathfrak{N}(v)$. For real $B$, the bound on
the complexified potentials gives
\begin{equation*}
\text{activity of the unmarked input}
\le8C_A\theta_k\bigl[2\bar C_{\mathsf{W}}+C^{(5)}r\,g_k^2(\log g_k^{-2})^{2q}\bigr].
\end{equation*}
The supremum in \eqref{8.5:eq-local-source-split-b} runs over all couplings $g'\le g$, and
$g_k\le g$ (see the paragraphs following the proposition). So the right-hand side above is at most
the left-hand side of \eqref{8.5:eq-local-source-split-b}. Hence the smallness condition of the
marked-expansion lemma on the unmarked input is \eqref{8.5:eq-local-source-split-b}.

\emph{The marked weight.} Since $\theta_k\le\theta(g)$ (see the paragraphs following the
proposition), the second condition in \eqref{8.5:eq-local-source-split-a} gives
$\theta_k\le\frac13$. We apply
Lemma~\ref{8.5:lem-unit-sum-bound} at the radius system $\flat$ to the marked family just
described. It contains no Wilson pieces, so the last term of
\eqref{8.5:eq-unit-sum-bound-a} is absent.

We insert $\mathfrak{n}^G_j=\mathfrak{m}_{j-1}$ and
$\mathfrak{n}^M_j=C_B\omega_{k,j}\mathfrak{m}_{j-1}$ into the two remaining sums. The level sums
$S_I$ and $S_{\omega}$ of the correlation theorem bound the corresponding sums over $j\le k$ of the
ratio and decay factors. This turns the two sums into $S_I\max_{j\le k}\mathfrak{m}_{j-1}$ and
$C_BS_{\omega}\max_{j\le k}\mathfrak{m}_{j-1}$. For the singles (M$_2$), read at the radius
$\hat\alpha_k$, the ratio $\rho^2_{k,j}$ is the weight $w_{k,j}$. We obtain
\begin{equation}\label{8.5:eq-local-source-split-3}
\begin{split}
W^m&\le2C_U^{+}\theta_k\Bigl[(S_I+C_BS_{\omega})\max_{j\le k}\mathfrak{m}_{j-1}
+\sum_{j\le k_1(k)}w_{k,j}\,g_j^{\kappa_0}\Bigr]\\
&\le2C_U^{+}\theta_k\bigl[(S_I+C_BS_{\omega})\,\mathfrak{m}_{k-1}+\mu_k\bigr].
\end{split}
\end{equation}
In the second line we used that $\mathfrak{m}$ is nondecreasing in its index, so that the
maximum is $\mathfrak{m}_{k-1}$, and that the sum over $j\le k_1(k)$ is bounded by the
remainder-memory sum $\mu_k$.

\emph{Definition of the two families at level $k+1$.} For complex $\lambda$ let $t_{\lambda}$ be
the output kernel of the marked expansion when the marked input is multiplied by $\lambda$ and
the unmarked input is left unchanged. Let $t_0$ and $t_1$ be its values at $\lambda=0$ and
$\lambda=1$; both are kernels fixed as in \cite[(3.47)]{B-CMP119}. The marked-expansion lemma
covers every output kernel of the fluctuation step of \cite{B-CMP119}, at the level-dependent
radii used there. We put
\begin{equation*}
P^{(k+1)}:=t_0,\qquad Q^{(k+1)}:=t_1-t_0 .
\end{equation*}
At $\lambda=1$ the input is the full input listed above. Hence $t_1=E^{(k+1)}$, which is (i) at
level $k+1$.

\emph{The bound on $Q^{(k+1)}$.} The third condition in \eqref{8.5:eq-local-source-split-a}
gives $\lambda_0\ge2$, with $\lambda_0$ as in Lemma~\ref{8.5:lem-direct-potentials-bis}. The
corollary of the marked-expansion lemma then gives
\begin{equation*}
|Q^{(k+1)}(X,\cdot,y;w)|\le K_{CE}\,W^m\,e^{-\kappa d_{k+1}(X)} .
\end{equation*}
We insert \eqref{8.5:eq-local-source-split-3}. The first condition in
\eqref{8.5:eq-local-source-split-a}, evaluated at $g'=g_k$, bounds both
$2K_{CE}C_U^{+}\theta_k(S_I+C_BS_{\omega})$ and $2K_{CE}C_U^{+}\theta_k$ by $\frac12$. Hence
\begin{equation*}
|Q^{(k+1)}(X,\cdot,y;w)|
\le\bigl(\tfrac12\mathfrak{m}_{k-1}+\tfrac12\mu_k\bigr)e^{-\kappa d_{k+1}(X)}
\le\mathfrak{m}_k\,e^{-\kappa d_{k+1}(X)} .
\end{equation*}
The last inequality holds because $\mu_k\le\mathfrak{m}_k$, by the inequality between these two
bounds stated in the correlation theorem, and because $\mathfrak{m}_{k-1}\le\mathfrak{m}_k$
($\mathfrak{m}$ is nondecreasing in $k$). This is the first bound in
\eqref{8.5:eq-local-source-split-d} at level $k+1$.

\emph{Proof of (ii).} Analyticity and invariance under the symmetries of
\cite[(1.19)]{B-CMP109} are the analyticity and invariance statements of the marked-expansion
lemma, applied to $t_0$ and $t_1$. For locality we use two facts. In the marked expansion an
output with domain $X$ depends only on the potentials $V(Y)$ with $Y\subset X$. By the local step
statement, each potential in $Y$ depends only on $\zeta_k$ restricted to $Y$ and on the old terms
in $Y$, that is, on $w$ restricted to $X^{+}$. Together these give the locality in (ii).

\emph{Proof of (iii).} Suppose $w\equiv\hat w$ on $X^{+}$. Then every potential in $X$, marked or
not, is the frozen one. The $Q$-family proposition defines $Q_{\hat w}$ through the same marking.
Hence $t_1-t_0=Q^{(k+1)}_{\hat w}$ and $t_0=P^{\flat(k+1)}_{\hat w}$ on $X$.

\emph{The bound on $P^{(k+1)}$.} We use three facts: (i); the bound
$|E^{(k+1)}|\le E_0e^{-\kappa d_{k+1}(X)}$ of the inductive statement of the correlation theorem
at level $k+1\le K'$, which is assumed; and the bound on $Q^{(k+1)}$ just proved. They give
\begin{equation*}
|P^{(k+1)}|\le|E^{(k+1)}|+|Q^{(k+1)}|
\le(E_0+\mathfrak{m}_k)\,e^{-\kappa d_{k+1}(X)}\le2E_0\,e^{-\kappa d_{k+1}(X)} .
\end{equation*}
The last inequality is $\mathfrak{m}_k\le E_0$. This completes (iv) at level $k+1$.

\emph{Proof of (v).} At $\lambda=0$ we use the identity \cite[(3.42)]{B-CMP119} together with
\cite[(2.11)--(2.13)]{B-CMP116}. We also use the lemma of the correlation theorem re-summing the
polymer expansion on a finite torus: it states that on a finite torus the exponential of the sum
of the kernels over all polymers equals the integral from which they are extracted. This gives
\begin{multline*}
\sum_{X,y}P^{(k+1)}(X,U_{k+1},\cdot,y;w)+\mathrm{const}(w)
=\log\int d\mu_{C^{(k)}}\,\chi\,\exp\Bigl[\mathcal{P}^{(k)}+\mathfrak{S}^w_k\\
+\sum_{j\le k}\Bigl\{\bigl(P^{(j)}(w)\bigr)^{\mathrm{ren},\mathrm{loc}}
\bigl(U_k(e^{iB'}V^{(k)})\bigr)
-\bigl(P^{(j)}(w)\bigr)^{\mathrm{ren},\mathrm{loc}}(U_{k+1})\Bigr\}\Bigr].
\end{multline*}
In this identity the dot on the left now stands for the current $\mathbf{J}$ alone, and
\begin{equation*}
\bigl(P^{(j)}(w)\bigr)^{\mathrm{ren},\mathrm{loc}}
:=\sum_y\mathfrak{i}_y\bigl[P^{\flat(j)}_{w(p_y)}\bigr]+\sum\bigl(D^P\bigr)^{\mathrm{vac}}
\end{equation*}
is the unmarked part of $(E^{(j)}(w))^{\mathrm{ren},\mathrm{loc}}$; $\mathfrak{S}^w_k$ is the
vertex with $\zeta_k=\zeta^w_k$. The inputs enter here only through their totals. Indeed, by the
statement of the correlation theorem on the reading of older terms, the integrand of every step
depends on the older terms only through their totals.

At $w=0$ the integrand is that of part~(d) of the proposition on the total of the potentials of
the fluctuation step, so the normalised measure at $w=0$ is the measure
\eqref{8.5:eq-local-source-split-2}. Indeed, $P^{\flat}_0$ agrees with $\mathsf{P}_0$ by the
three-family split of the correlation theorem, and $\beta[P^{\flat}_0]=\beta[\mathsf{P}_0]$ by the
lemma identifying the two renormalisation coefficients.

We now set $w=zg$ and differentiate at $z=0$. The derivative in $z$ at $z=0$ of the logarithm of
the integral is the expectation, in the normalised measure at $z=0$, of the $z$-derivative at
$z=0$ of the exponent. The potentials $\mathcal{P}^{(k)}$ do not depend on $z$ and contribute
nothing. We treat the other terms in turn.
\begin{itemize}
\item[(1)] Differentiating the formula for the sourced shift $\zeta^w_k$ in the correlation
theorem at $w=zg$ gives $\partial_z\zeta^{zg}_k|_{z=0}=\dot\zeta_k[g]$. The vertex depends
linearly on $\zeta_k$ and
$A(f,\cdot)$ is linear in $f$. Hence $\partial_z\mathfrak{S}^{zg}_k|_{z=0}$ is the sum of the
Wilson and gauge-fixing members of $\dot{\mathfrak{B}}_k[g]$, namely
\begin{equation*}
A(\dot\zeta_k[g],U_k(e^{iB'}V^{(k)}))-A(\dot\zeta_k[g],U_{k+1})
+\sum_{y}\dot\zeta_k[g](p_y)\,G_y(B').
\end{equation*}
\item[(2)] For $j\le k$ we have
\begin{equation*}
\partial_z\bigl(P^{(j)}(zg)\bigr)^{\mathrm{ren},\mathrm{loc}}\Big|_{z=0}
=\Bigl(\partial_z\sum P^{(j)}\Bigr)^{\mathrm{vac}}
-\sum_y\partial_{\zeta}\beta\bigl[P^{\flat}_{\zeta}\bigr]\Big|_{\zeta=0}\,g(p_y)\,A(h_y,\cdot).
\end{equation*}
Since $I_jg=\sum_yg(p_y)h_y$ and $A(f,\cdot)$ is linear in $f$, the second term equals
$\mathsf{b}^P_jA(I_jg,\cdot)$. On the support of the cut-off $\chi$ the configuration
$U_k(e^{iB'}V^{(k)})$ is real and regular, so it lies in the class of (v) at level $j\le k$ by the
inclusion noted there. So does $U_{k+1}=U$, being $U_{k+1}(W)$ with $k+1\ge j$. Hence (v) at
level $j$ applies at both configurations. It identifies $\partial_z\sum P^{(j)}$ with
$\dot{\mathcal{T}}^{(j)}[g]$ up to a constant, and the constant cancels in the vacuum
subtraction. Hence
\begin{equation*}
\partial_z\bigl(P^{(j)}(zg)\bigr)^{\mathrm{ren},\mathrm{loc}}\Big|_{z=0}
=\bigl(\dot{\mathcal{T}}^{(j)}[g]\bigr)^{\mathrm{ren}},
\end{equation*}
evaluated at both configurations.
\end{itemize}
Adding (1) and the differences from (2), the $z$-derivative of the exponent at $z=0$ is
$\dot{\mathfrak{B}}_k[g]$. Taking the expectation in \eqref{8.5:eq-local-source-split-2} at
$U_{k+1}=U$ and using \eqref{8.5:eq-local-source-split-c}, we obtain
\begin{equation*}
\partial_z\Big|_{z=0}\sum_{X,y}P^{(k+1)}\bigl(X,U,J(U),y;zg\bigr)
=\bigl\langle\dot{\mathfrak{B}}_k[g]\bigr\rangle^U_k+\mathrm{const}
=\dot{\mathcal{T}}^{(k+1)}[g](U)+\mathrm{const}.
\end{equation*}
This is \eqref{8.5:eq-local-source-split-1} at level $k+1$ and completes the induction.
\end{proof}

\begin{remark}
\leavevmode
\begin{enumerate}
\item[(1)] None of the statements that make up the induction of the correlation theorem uses
Proposition~\ref{8.5:prop-local-source-split}. The proposition uses that inductive statement and
itself at lower levels only, so the argument is not circular.
\item[(2)] The reason the bound on $Q^{(k+1)}$ closes is the one in the $Q$-family proposition.
Each $Q$-family is renormalised by its own coefficient; the differences $D^Q$ decay like
$\omega_{k,j}$; and the remainders of earlier levels are paid for by the factors
$g_j^{\kappa_0}$.
\item[(3)] Put $\mathsf{b}^Q_j:=\partial_{\zeta}\beta[Q^{(j)}_{\zeta}]|_{\zeta=0}$. By the
running-shift statement, $\mathsf{b}^Q_j=b_j'(0)-\mathsf{b}^P_j$. The running-shift statement also
bounds $|\beta[Q^{(j)}_{\zeta}]|$ by $v_{j-1}$ on the disc $D(0,2r)$, on which
$\zeta\mapsto\beta[Q^{(j)}_{\zeta}]$ is analytic. The Cauchy estimate on $D(0,2r)$, together with
the statement of the correlation theorem on the real values of the flow coefficients, then gives
\begin{equation*}
|\mathsf{b}^Q_j|\le\frac{v_{j-1}}{2r}.
\end{equation*}
\end{enumerate}
\end{remark}

\begin{lemma}[Quadratic smallness of the source-dependent part]\label{8.5:lem-source-quadratic-bis}
We work in the setting and notation of Notation~\ref{8.5:not-field-radii}. Put
\begin{equation*}
c_S:=\frac{64}{3}C_5+1+\frac{8}{3}\,\frac{C_6M^4C_1^2}{C_AK_R^2C_{\mathfrak{a}}^2},
\qquad C_s:=4K_{CE}\,c_S .
\end{equation*}
Assume that the conclusions of the correlation theorem hold at every level $k\le K'$, and
assume the ceilings
\begin{equation}\label{8.5:eq-source-quadratic-a-bis}
\begin{gathered}
\sup_{g'\le g}4K_{CE}C_U^{+}\theta(g')\bigl(S'_I+C_BS'_{\omega}\bigr)\le\tfrac12,
\qquad \theta(g)\le\tfrac13;\\
\sup_{x\ge g^{-2}-4c_2}\frac{4C_U^{+}C_{\mathfrak{m}}\,x^{-(\kappa_0-2)/2}\cdot 2x}
{C_{\mathfrak{a}}^2C_*^2(\log x)^{2q}}\le r;\\
\mathfrak{a}_k^2\le 2\mathfrak{a}_{k+1}^2\quad(0\le k<K');\\
\sup_{g'\le g}C_s\,r\,\theta(g')\,C_{\mathfrak{a}}^2\alpha_*(g')^2\le b^{\mathrm{adm}}.
\end{gathered}
\end{equation}
In the second line $\kappa_0$ is the exponent of the remainder bound of the correlation theorem,
and $C_{\mathfrak{m}}$ is the constant of its bound on the remainder memory sums
$\mathfrak{m}_k$ recalled in the proof below.
Here the third line follows from the bound
$x_{k+1}/x_k\in[1-b_{+}g^2,\,1+2c_2g^2]$ on consecutive inverse couplings, and the fourth line
ensures $\lambda_0\ge2$ in Lemma~\ref{8.5:lem-direct-potentials-bis}.
Assume finally the ceiling \eqref{8.5:eq-local-source-split-b} at its $z$-free members, that is,
the activity hypothesis of the marked cluster expansion for the unmarked $z$-free run. Then for
$1\le j\le K'$, for every pointed term $E^{(j)}(X,\cdot,y;w)$ of the sourced effective action at
level $j$, in the format of the correlation theorem, and for every
$w\in\mathcal{W}_j(X)$,
\begin{equation}\label{8.5:eq-source-quadratic-b-bis}
\bigl|E^{(j)}(X,\cdot,y;w)-E^{(j)}(X,\cdot,y;0)\bigr|
\le C_s\,r\,\hat\theta_{j-1}\,\mathfrak{a}_j^2\,e^{-\kappa d_j(X)}
\qquad\text{on }\mathcal{U}^{\flat}_j(X).
\end{equation}
\end{lemma}

\begin{proof}
For $1\le j\le K'$ let
\begin{equation*}
S^{(j)}(w):=E^{(j)}(w)-E^{(j)}(0)
\end{equation*}
be the source-dependent part, let $\mathsf{s}_j$ be its $\|\cdot\|^{\flat}$-norm as a pointed
$E$-family at scale $j$ on the shrinking radii, the supremum being taken over the pointed pairs
$(X,y)$, over backgrounds in $\mathcal{U}^{\flat}_j(X)$ and over
$w\in\mathcal{W}_j(X)$, and put $\hat\sigma_j:=C_s\,r\,\hat\theta_{j-1}$. We prove
\begin{equation}\label{8.5:eq-source-quadratic-1-bis}
\mathsf{s}_j\le\hat\sigma_j\,\mathfrak{a}_j^2,\qquad 1\le j\le K',
\end{equation}
which gives \eqref{8.5:eq-source-quadratic-b-bis}, the $\|\cdot\|^{\flat}$-norm of a pointed
$E$-family at scale $j$ being the supremum of $e^{\kappa d_j(X)}|\cdot|$ over backgrounds in
$\mathcal{U}^{\flat}_j(X)$ (Notation~\ref{8.5:not-field-radii}). The proof is an induction over
the steps
$k\to k+1$, $0\le k<K'$: we assume \eqref{8.5:eq-source-quadratic-1-bis} for $1\le j\le k$, and we
prove it for $j=k+1$. Write $\bar\sigma_k:=\max_{1\le j\le k}\hat\sigma_j$, with
$\bar\sigma_0:=0$ (empty maximum); the groups and the singles from the difference family below
are indexed by the levels $1\le j\le k$ at which $S^{(j)}$ is defined, so that at step $0\to1$
there are none.

\emph{Unmarked and marked potentials.} At step $k\to k+1$ we take as unmarked potentials those
of the $z$-free run, that is, of the run at $w=0$. The hypotheses of the marked cluster
expansion on the unmarked input hold for them by the inductive bounds of the $z$-free run
supplied by the correlation theorem at $w=0$, together with \eqref{8.5:eq-local-source-split-b}
at its $z$-free members. The marked potentials are the differences between the sourced run and
the $z$-free run. They are units of the three kinds of the unit-sum bound
(Lemma~\ref{8.5:lem-unit-sum-bound}), and one family given directly:
\begin{itemize}
\item[(G)] Groups. For a constant source $\hat w$ (constants belong to $\mathcal{W}_j(X)$),
\begin{equation*}
\mathfrak{i}^{(j)}_y\bigl[E^{(j)}_{\hat w}\bigr]-\mathfrak{i}^{(j)}_y\bigl[E^{(j)}_{0}\bigr]
=\mathfrak{i}^{(j)}_y\bigl[S^{(j)}_{\hat w}\bigr],
\qquad S^{(j)}_{\hat w}:=E^{(j)}_{\hat w}-E^{(j)}_0 .
\end{equation*}
As the difference of two pointed $E$-families, $S^{(j)}_{\hat w}$ is a pointed $E$-family at
scale $j$, with its own $\beta$-coefficient $b_j(\hat w)-b_j(0)$; its norm is at most
$\mathsf{s}_j$. Thus $\mathfrak{n}^G_j=\mathsf{s}_j$.
\item[(M)] Singles from the difference family. The difference family is
$D^{(j)}=E^{(j)}(w)-E^{(j)}_{w(p_y)}$. Subtracting and adding $E^{(j)}(0)=E^{(j)}_0$, the
$z$-free parts cancel:
\begin{equation*}
D^{(j)}=\bigl(E^{(j)}(w)-E^{(j)}(0)\bigr)-\bigl(E^{(j)}_{w(p_y)}-E^{(j)}_0\bigr)
=S^{(j)}(w)-S^{(j)}_{w(p_y)} .
\end{equation*}
The marked single is $(D^{(j)})^{\mathrm{vac}}$. The bound for pointed families under freezing
of the source at $p_y$, applied to $F:=S^{(j)}(X,\cdot,y;\cdot)$, gives
$\mathfrak{n}^M_j=C_B\,\omega_{k,j}\,\mathsf{s}_j$, with the constant $C_B$ and the weights
$\omega_{k,j}$ of that bound.
\item[(M)] Singles from the remainder. For $j\le k_1(k)$, where $k_1(k)$ is the memory cutoff of
the remainder families $R^{(j)}$ at step $k$ fixed in the correlation theorem, the marked
single is
$(R^{(j)}(w)-R^{(j)}(0))^{\mathrm{vac}}$, with $\mathfrak{n}^M_j=2g_j^{\kappa_0}$.
\item[(W)] Wilson pieces. For each shift cube the marked piece is
$A(\zeta^w_k\hat\chi_{\square},\cdot)$, with $\hat\chi_{\square}$ the localising function of the
shift cube $\square$ in the correlation theorem. The bound $|\zeta_k|\le\tfrac43\sup|w|$ on the
running source shift, which is part of the correlation theorem, gives $\bar c=\tfrac83\,r$ for
$w\in\mathcal{W}_j(X)$, on which $\sup|w|\le2r$ by the definition of $\mathcal{W}_j(X)$ in the
correlation theorem.
\item[] Direct. The gauge-fixing pieces $\zeta_k(p_y)G_y$, with $G_y$ the gauge-fixing term at
the block of $y$, are marked potentials given directly
in the sense of Lemma~\ref{8.5:lem-direct-potentials-bis}, with $|\zeta_k(p_y)|\le\bar c$. Per cube
their activity obeys
\begin{equation*}
\tilde a^{\mathrm{dir}}\le\frac83\,r\,C_6M^4C_1^2\,\delta_k^2,
\qquad
\delta_k^2=\frac{\theta_k^2\,\mathfrak{a}_k^2}{(K_RC_{\mathfrak{a}})^2},
\end{equation*}
the identity for $\delta_k^2$ being the definition $\theta_k=K_R\delta_k/\alpha_{*,k}$
combined with $\mathfrak{a}_k=C_{\mathfrak{a}}\alpha_{*,k}$.
\end{itemize}

\emph{The unit sums.} We apply Lemma~\ref{8.5:lem-unit-sum-bound} with the shrinking radii,
$a=\flat$, $\theta_{\flat}=\theta_k$ and $\mathfrak{a}=\mathfrak{a}_k$; its hypothesis
$\theta_k\le\tfrac13$ follows from $\theta(g)\le\tfrac13$ in the first line of
\eqref{8.5:eq-source-quadratic-a-bis}, since $\theta_k=\theta(g_k)$, the function
$\theta(g')=C_{\theta}(\log g'^{-2})^{p_0-q}$ is nondecreasing in $g'$ for $q\ge p_0$
(Notation~\ref{8.5:not-field-radii}), and $g_k\le g$ for $k\le K'$. By the induction hypothesis
$\mathsf{s}_j\le\hat\sigma_j
\mathfrak{a}_j^2\le\bar\sigma_k\mathfrak{a}_j^2$ for $1\le j\le k$, and
$\rho_{k,j}^2\mathfrak{a}_j^2=\mathfrak{a}_k^2$ by the definition of $\rho_{k,j}$. The group
term of \eqref{8.5:eq-unit-sum-bound-a} is therefore at most
\begin{equation*}
2C_U^{+}\theta_k\,\mathfrak{a}_k^2\,\bar\sigma_k
\sum_{1\le j\le k}\nu^3(1+\rho_{k,j})L^{-\hat\alpha(\nu-1)}
\le 2C_U^{+}\theta_k\,\mathfrak{a}_k^2\,S'_I\,\bar\sigma_k,
\end{equation*}
because with $\nu=1+k-j$ the sum runs over $1\le\nu\le k$ with $j=k+1-\nu$ and is bounded by
$S'_I$, by the definition of $S'_I$ (with the same decay exponent $\hat\alpha$). The singles
from the difference family contribute at most
\begin{equation*}
2C_U^{+}\theta_k\,\mathfrak{a}_k^2\,C_B\,\bar\sigma_k\sum_{1\le j\le k}\omega_{k,j}
\le 2C_U^{+}\theta_k\,\mathfrak{a}_k^2\,C_BS'_{\omega}\,\bar\sigma_k ,
\end{equation*}
since the weights $\omega_{k,j}$ of the freezing bound satisfy, by that bound,
$\sum_{1\le j\le k}\omega_{k,j}\le S'_{\omega}=1+\sum_{a\ge1}L^{-a/2}$. The singles from the
remainder contribute
$2C_U^{+}\theta_k\sum_{j\le k_1(k)}\rho_{k,j}^2\cdot2g_j^{\kappa_0}$, which is at most
$4C_U^{+}\theta_k\mathfrak{m}_k$, where $\mathfrak{m}_k$ is the bound fixed in the correlation
theorem at step $k$ for the sum $\sum_{j\le k_1(k)}\rho_{k,j}^2g_j^{\kappa_0}$. The
Wilson term is $8\theta_kC_5\bar c\,\mathfrak{a}_k^2=\tfrac{64}{3}C_5\,r\,\theta_k
\mathfrak{a}_k^2$. For the direct activity, $\theta_k\le1$ gives
\begin{equation*}
\tilde a^{\mathrm{dir}}
\le C_A\theta_k\,\mathfrak{a}_k^2\cdot
\frac83\,\frac{C_6M^4C_1^2}{C_AK_R^2C_{\mathfrak{a}}^2}\,r .
\end{equation*}
Adding, and recalling
$c_S-1=\tfrac{64}{3}C_5+\tfrac83C_6M^4C_1^2/(C_AK_R^2C_{\mathfrak{a}}^2)$, we obtain for the
activity $\tilde a=C_AW^m+\tilde a^{\mathrm{dir}}$ of Lemma~\ref{8.5:lem-direct-potentials-bis}
\begin{equation*}
\tilde a\le C_A\theta_k\,\mathfrak{a}_k^2
\Bigl\{2C_U^{+}\bigl(S'_I+C_BS'_{\omega}\bigr)\bar\sigma_k+(c_S-1)\,r\Bigr\}
+4C_AC_U^{+}\theta_k\,\mathfrak{m}_k .
\end{equation*}
By the correlation theorem,
\begin{equation*}
\mathfrak{m}_k\le C_{\mathfrak{m}}y_k^{-(\kappa_0-2)/2},
\qquad g^{-2}-4c_2\le y_k\le x_k\le2y_k,
\end{equation*}
where $y_k$ is a number (not a lattice point). The second line of
\eqref{8.5:eq-source-quadratic-a-bis} at $x=y_k$ therefore gives
\begin{equation*}
4C_U^{+}\mathfrak{m}_k
\le r\,\frac{C_{\mathfrak{a}}^2C_*^2(\log y_k)^{2q}}{2y_k}
\le r\,\frac{C_{\mathfrak{a}}^2C_*^2(\log x_k)^{2q}}{x_k}
=r\,\mathfrak{a}_k^2,
\end{equation*}
the middle inequality because $y_k\le x_k\le2y_k$, and the last identity because
$\alpha_{*,k}=g_kC_*(\log x_k)^q$ and $x_k=g_k^{-2}$. Hence
\begin{equation}\label{8.5:eq-source-quadratic-2-bis}
\tilde a\le C_A\theta_k\,\mathfrak{a}_k^2
\Bigl\{2C_U^{+}\bigl(S'_I+C_BS'_{\omega}\bigr)\bar\sigma_k+c_S\,r\Bigr\}.
\end{equation}

\emph{The marked bound.} The fourth line of \eqref{8.5:eq-source-quadratic-a-bis} ensures the
hypothesis $\tilde a\le\nu_0\alpha_4/8$, that is $\lambda_0\ge2$ (recall
$\lambda_0=\nu_0\alpha_4/(4\tilde a)$), under which
Lemma~\ref{8.5:lem-direct-potentials-bis} bounds the marked sub-sum by
\eqref{8.5:eq-direct-potentials-a-bis},
$|t_1-t_0|(X)\le\mathsf{K}\tilde a\,e^{-\kappa d_{k+1}(X)}$ with
$\mathsf{K}=K_{CE}/C_A$. The source-dependent part at level $k+1$ is this marked sub-sum, so
\eqref{8.5:eq-source-quadratic-2-bis} gives
\begin{equation*}
\mathsf{s}_{k+1}\le K_{CE}\theta_k\,\mathfrak{a}_k^2
\Bigl\{2C_U^{+}\bigl(S'_I+C_BS'_{\omega}\bigr)\bar\sigma_k+c_S\,r\Bigr\}.
\end{equation*}
By the third line of \eqref{8.5:eq-source-quadratic-a-bis}, $\mathfrak{a}_k^2\le2\mathfrak{a}_{k+1}^2$;
by the first line, $4K_{CE}C_U^{+}\theta_k(S'_I+C_BS'_{\omega})\le\tfrac12$. Therefore
\begin{equation*}
\mathsf{s}_{k+1}
\le\mathfrak{a}_{k+1}^2\Bigl\{\tfrac12\bar\sigma_k+2K_{CE}c_S\,r\,\theta_k\Bigr\}
\le\hat\sigma_{k+1}\,\mathfrak{a}_{k+1}^2 .
\end{equation*}
The last inequality holds because $\bar\sigma_k=\max_{1\le j\le k}C_sr\hat\theta_{j-1}\le
C_sr\hat\theta_k$, $2K_{CE}c_S=\tfrac12C_s$ and $\theta_k\le\hat\theta_k$, so that the brace is
at most $\tfrac12C_sr\hat\theta_k+\tfrac12C_sr\hat\theta_k=\hat\sigma_{k+1}$.

Finally, the fresh outputs of step $k$ are defined on the output domain of that step, at the
enlarged radii of the correlation theorem, and this domain contains
$\mathcal{U}^{\flat}_{k+1}(X)$ by the correlation theorem (its statement on the
output domains of the fresh terms); so the bound on
$\mathsf{s}_{k+1}$ holds on $\mathcal{U}^{\flat}_{k+1}(X)$. This proves
\eqref{8.5:eq-source-quadratic-1-bis} for $j=k+1$ and completes the induction.
\end{proof}

\begin{corollary}[The first source derivative on shrinking radii]\label{8.5:cor-source-derivative-bis}
Assume the hypotheses of Lemma~\ref{8.5:lem-source-quadratic-bis}, and let $1\le j\le K'$.
\begin{itemize}
\item[(a)] If $f$ is Lipschitz and $z\in\mathbb{C}$ satisfies $|z|\,\|f\|_{\sharp}<r$, then
$zf\in\mathcal{W}_j$.
\item[(b)] Write $E^{(j)}_{\zeta}:=E^{(j)}(X,\cdot,y;\zeta)$ for the family at the constant
source $w\equiv\zeta$. The derivative
\begin{equation*}
\dot{S}^{(j)}:=\partial_{\zeta}E^{(j)}_{\zeta}\big|_{\zeta=0}
\end{equation*}
is a pointed $E$-family. It satisfies
\begin{equation}\label{8.5:eq-source-derivative-1-bis}
\|\dot{S}^{(j)}\|^{\flat}\le C_s\,\hat\theta_{j-1}\,\mathfrak{a}_j^2,
\qquad
\beta[\dot{S}^{(j)}]=b_j'(0).
\end{equation}
\item[(c)] For Lipschitz $f$ put
\begin{equation*}
\dot{D}^{(j)}(X,\cdot,y)[f]:=\partial_z D^{(j)}(X,\cdot,y;zf)\big|_{z=0}.
\end{equation*}
Here $C_B$, $\omega_{K',j}$ and $\delta$ are the constant, the level weight and the parameter of
the decay exponent in the lemma on the Lipschitz remainder of a pointed family.
Then, on $\mathcal{U}^{\flat}_j(X)$,
\begin{equation}\label{8.5:eq-source-derivative-2-bis}
\big|\dot{D}^{(j)}(X,\cdot,y)[f]\big|
\le C_B\,C_s\,\hat\theta_{j-1}\,\mathfrak{a}_j^2\,\omega_{K',j}\,\|f\|_{\sharp}\,
e^{-(1-\frac14\delta)\kappa d_j(X)}.
\end{equation}
\end{itemize}
The radius $r$ does not occur in \eqref{8.5:eq-source-derivative-1-bis} or
\eqref{8.5:eq-source-derivative-2-bis}.
\end{corollary}

\begin{proof}
(a) By the display of the correlation theorem that describes the admissible sources, a source
$zf$ with $|z|\,\|f\|_{\sharp}<r$ lies in $\mathcal{W}_j$. Hence $zf\in\mathcal{W}_j$.

(b) Take $f\equiv1$. Then $\|f\|_{\sharp}=\max\{1,0\}=1$, and (a) gives $\zeta\in\mathcal{W}_j$ whenever
$|\zeta|<r$. For such $\zeta$ put $S^{(j)}_{\zeta}:=E^{(j)}_{\zeta}-E^{(j)}_0$. By the proof of
Lemma~\ref{8.5:lem-source-quadratic-bis}, $S^{(j)}_{\zeta}$ is a pointed $E$-family, being the difference of two
such families. Its own coefficient is $\beta[S^{(j)}_{\zeta}]=b_j(\zeta)-b_j(0)$. Its
$\|\cdot\|^{\flat}$-norm is at most $\mathsf{s}_j$, the supremum over $\mathcal{W}_j$. The same proof
establishes
\begin{equation*}
\mathsf{s}_j\le\hat\sigma_j\,\mathfrak{a}_j^2=C_s\,r\,\hat\theta_{j-1}\,\mathfrak{a}_j^2,
\end{equation*}
which is the norm form of \eqref{8.5:eq-source-quadratic-b-bis}.

The family $E^{(j)}_0$ does not depend on $\zeta$, so $\dot{S}^{(j)}=\partial_{\zeta}S^{(j)}_{\zeta}|_{\zeta=0}$.
The family $E^{(j)}(X,\cdot,y;w)$ is analytic in the source $w\in\mathcal{W}_j$ (the correlation
theorem), so $\zeta\mapsto S^{(j)}_{\zeta}$ is analytic on $|\zeta|<r$, and for $0<r'<r$
Cauchy's formula gives
\begin{equation*}
\dot{S}^{(j)}=\frac{1}{2\pi i}\oint_{|\zeta|=r'}\frac{S^{(j)}_{\zeta}}{\zeta^2}\,d\zeta .
\end{equation*}
Thus $\dot{S}^{(j)}$ is a pointed $E$-family. Taking norms under the integral,
\begin{equation*}
\|\dot{S}^{(j)}\|^{\flat}
\le\frac{1}{r'}\sup_{|\zeta|=r'}\|S^{(j)}_{\zeta}\|^{\flat}
\le\frac{C_s\,r\,\hat\theta_{j-1}\,\mathfrak{a}_j^2}{r'} .
\end{equation*}
Letting $r'\uparrow r$ gives the first bound in \eqref{8.5:eq-source-derivative-1-bis}. The same
representation, applied to $\beta[S^{(j)}_{\zeta}]=b_j(\zeta)-b_j(0)$, gives
$\beta[\dot{S}^{(j)}]=b_j'(0)$.

(c) If $f=0$, both sides of \eqref{8.5:eq-source-derivative-2-bis} vanish. Otherwise $\|f\|_{\sharp}>0$, and by
(a) the source $zf$ lies in $\mathcal{W}_j$ for $|z|<r/\|f\|_{\sharp}$. As in the proof of
Lemma~\ref{8.5:lem-source-quadratic-bis}, for such sources
\begin{equation*}
D^{(j)}(X,\cdot,y;w)=S^{(j)}(w)-S^{(j)}_{w(p_y)} .
\end{equation*}
The $z$-free parts cancel in this difference. We now apply the lemma on the Lipschitz remainder
$F(w)-F_{w(p_y)}$ of a pointed family, at level $K'$, to $F:=S^{(j)}(X,\cdot,y;\cdot)$, whose
$\|\cdot\|^{\flat}$-norm is at most $\mathsf{s}_j$. On $\mathcal{U}^{\flat}_j(X)$ it gives
\begin{equation*}
|D^{(j)}(X,\cdot,y;zf)|
\le C_B\,\omega_{K',j}\,C_s\,r\,\hat\theta_{j-1}\,\mathfrak{a}_j^2\,e^{-(1-\frac14\delta)\kappa d_j(X)}.
\end{equation*}
For the same reason as in (b), $z\mapsto D^{(j)}(X,\cdot,y;zf)$ is analytic on
$|z|<r/\|f\|_{\sharp}$, and Cauchy's estimate on the circle $|z|=r''$, for $r''<r/\|f\|_{\sharp}$, bounds
$|\dot{D}^{(j)}(X,\cdot,y)[f]|$ by this right-hand side divided by $r''$. Letting
$r''\uparrow r/\|f\|_{\sharp}$ replaces the factor $r$ by $\|f\|_{\sharp}$, which is
\eqref{8.5:eq-source-derivative-2-bis}. In (b) and in (c) the factor $r$ of
\eqref{8.5:eq-source-quadratic-b-bis} is divided by the radius of the Cauchy circle, so $r$ cancels.
\end{proof}

\begin{definition}[Block indicators, Lipschitz moduli and constants]\label{8.5:not-block-indicators}
We use the block lattices $\mathbb{T}^{(j)}$ of Definition~\ref{1.5:def-block-average}, the barycentre
$x(p)$ of a plaquette $p$, the symbols of the two systems of field radii fixed in
Notation~\ref{8.5:not-field-radii}, and the constants $C_5$, $C_6$, $\mathsf{K}$, $C_U^{+}$,
$S_I^{\circ}$ and $K_R$ (the constant in the ratio $\theta$) fixed there; $C_A$, $\alpha_4$, $\nu_0$
are the constants of the marked-expansion lemma. Write $\operatorname{dist}_j$ for the sup-distance
measured in level-$j$ units, and $\pi_j$ for Ba{\l}aban's partition of the level-$j$ lattice into
cubes of side $M$ \cite{B-CMP109}.

\emph{Lipschitz moduli.} For a bounded function $g$ from the plaquettes of the torus to
$\mathbb{C}$ (for such functions we keep the letter $g$ of the correlation theorem's first-order
responses; the reference coupling, also written $g$, is always called the coupling below) and a set
$A$ of plaquettes put
\begin{equation*}
\operatorname{Lip}_j(g;A)
:=\sup\Bigl\{\frac{|g(p)-g(p')|}{\operatorname{dist}_j(x(p),x(p'))}:\ p,p'\in A,\ p\neq p'\Bigr\}.
\end{equation*}
For $i\le j$ one has
\begin{equation*}
\operatorname{Lip}_i(g;A)=L^{-(j-i)}\operatorname{Lip}_j(g;A).
\end{equation*}

\emph{Evaluation plaquettes, profiles, cubes.} For $y\in\mathbb{T}^{(j)}$, $p_y$ is the evaluation
plaquette of $y$ used in the correlation theorem; the only property of $p_y$ used below is
$\operatorname{dist}_j(x(p_y),y)\le 1$. We write $h_y:=h^{(j)}_y$ for the functions of
\cite[(3.40), (3.65)]{B-CMP119}: $h_y\ge 0$, $\sum_y h_y=1$, $\operatorname{supp}h_y\subset
y+[-1,1]^4$ in level-$j$ units, and $h_{ry}=h_y\circ r^{-1}$ for the lattice symmetries $r$. For a
set $X$, $X^{\oplus}$ is $X$ enlarged by two layers of cubes of $\pi_j$. For $y\in\mathbb{T}^{(j)}$,
$\square_y\in\pi_j$ is the cube containing $y$, and $\square_y^1$ is the cube of side $3M$ about
$\square_y$ \cite{B-CMP109}; it is a union of $3^d$ cubes of $\pi_j$.

\emph{Block indicators.} For $y\in\mathbb{T}^{(j)}$ let $B(y)$, the level-$j$ unit block of $y$, be
the closed cube of side $1$, in level-$j$ units, centred at $y$. For a site $y'\in\mathbb{T}^{(i)}$
with $i<j$ put
\begin{equation*}
\mathsf{h}_y(y'):=
\begin{cases}
1, & y'\in B(y),\\
0, & \text{otherwise}.
\end{cases}
\end{equation*}
Since $L$ is odd, no site of $\mathbb{T}^{(i)}$ lies on the boundary of a block; the blocks $B(y)$,
$y\in\mathbb{T}^{(j)}$, partition $\mathbb{T}^{(i)}$; $rB(y)=B(ry)$ for every symmetry $r$
preserving $\mathbb{T}^{(j)}$; each cube of $\pi_j$ is a union of unit blocks; and
$B(y)\subset\square_y$.
In \cite[(3.40)]{B-CMP119} the function $h_z$ is given as a formula in the position,
$h_z(x)=\prod_{\mu}h(x_\mu-z_\mu)$ with $h(t)=\max\{1-L^{-1}|t|,0\}$ in level-$(j-1)$ units; read
at $x=y'$ it is covariant under the lattice symmetries and would serve equally as site weight.
As site weight we use the block indicator $\mathsf{h}_y$, and $h_y$ remains the profile.

\emph{Normalisation of the running profile.} With $\dot\zeta_k[g]$ the running source profile of
the first-order recursion and $b_j$ the flow coefficients of the correlation theorem, put
\begin{equation*}
N_k:=\dot\zeta_k[1]=1+\sum_{j\le k}b_j'(0),
\end{equation*}
where $b_j'(0)$ is the derivative at $0$ of the flow coefficient $b_j$ of the correlation theorem.
Then $N_k\in[\tfrac23,\tfrac43]$ and, for every bounded function $g$ on the plaquettes,
$|\dot\zeta_k[g]|\le\tfrac43\sup|g|$.

\emph{Constants.} Put
\begin{align*}
c_V&:=\tfrac{32}{3}\,C_A C_5\,\mathfrak{a}_{\circ}^2+\tfrac{4}{3}\,C_6M^4C_1^2\,\alpha_*/K_R,\\
\mathsf{c}_{\square}&:=2^{3^d-1}e^{(2^d+4)\kappa},\\
C_{LR}&:=6\,\mathsf{c}_{\square}\mathsf{K}\,c_V .
\end{align*}
They do not depend on the coupling $g$, on $K$, nor on the parameter $r>0$ of the correlation
theorem's quantifier prefix, and are fixed before $\gamma$.

\emph{Ceilings.} Write $T(\cdot)$ for the function of the coupling with $T_k=T(g_k)$, so that the
fluctuation cut is $\delta_k=g_kT(g_k)$, and, for the reference coupling $g$, put
$\theta^{\circ}(g):=K_RgT(g)/\alpha_*$. Let $c_{\mathrm{abs}}$ be
the constant of the analyticity condition at absolute radii under which the expansion of
\cite[(1.17)--(1.22)]{B-CMP116} is used, and let $\varepsilon_1$, $\bar\varepsilon$ be the smallness
parameters of the marked-expansion lemma. The ceilings are the inequalities
\eqref{8.5:eq-block-indicators-a} and \eqref{8.5:eq-block-indicators-b} below; the first three
inequalities in \eqref{8.5:eq-block-indicators-a} and the inequality
\eqref{8.5:eq-block-indicators-b} are conditions on the coupling $g$, and the fourth inequality in
\eqref{8.5:eq-block-indicators-a} is a consequence recorded with them:
\begin{equation}\label{8.5:eq-block-indicators-a}
\begin{gathered}
\theta^{\circ}(g)\le\tfrac13,\qquad e^{32\kappa_1}C_1\,gT(g)\le c_{\mathrm{abs}},\qquad
8\,\frac{gT(g)}{\varepsilon_1}\,\frac{\bar\varepsilon}{\alpha_4}\le\nu_0,\\
\delta_i\le 2\delta_k\quad(i\le k),
\end{gathered}
\end{equation}
and
\begin{equation}\label{8.5:eq-block-indicators-b}
2\,\mathsf{c}_{\square}\mathsf{K}\,C_AC_U^{+}\bigl(3S_I^{\circ}+1\bigr)\,\theta^{\circ}(g)\le\tfrac12 .
\end{equation}
The third inequality in \eqref{8.5:eq-block-indicators-a} is the smallness hypothesis of the
marked-expansion lemma, read for the potentials of the fluctuation step at the cut $gT(g)$.
\end{definition}

\begin{proof}[Proof of the assertions in Notation~\ref{8.5:not-block-indicators}]
\emph{Lipschitz moduli at two levels.} Let $i\le j$. One level-$j$ unit is $L^{j-i}$ level-$i$
units, so $\operatorname{dist}_i=L^{j-i}\operatorname{dist}_j$. Dividing $|g(p)-g(p')|$ by the two
distances and taking the supremum over $p\neq p'$ in $A$ gives
\begin{equation*}
\operatorname{Lip}_i(g;A)=L^{-(j-i)}\operatorname{Lip}_j(g;A).
\end{equation*}

\emph{The blocks.} Measure in level-$j$ units, with $\mathbb{T}^{(j)}$ at the integer points. Then
the faces of the blocks $B(y)$ lie on hyperplanes $\{x_\mu=n+\tfrac12\}$, $n\in\mathbb{Z}$, while
the sites of $\mathbb{T}^{(i)}$, $i<j$, have coordinates in $L^{-(j-i)}\mathbb{Z}$. If
$L^{-(j-i)}a=n+\tfrac12$ with $a\in\mathbb{Z}$, then $2a=L^{j-i}(2n+1)$, which is impossible because
$L$ is odd. Hence no site of a finer lattice lies on a block boundary. The closed blocks cover
space and two distinct blocks meet only on their boundaries, so each site of $\mathbb{T}^{(i)}$
lies in exactly one block: the blocks partition $\mathbb{T}^{(i)}$. A symmetry $r$ preserving
$\mathbb{T}^{(j)}$ is an isometry for the sup-distance, so it maps the unit cube centred at $y$ onto
the unit cube centred at $ry$, that is $rB(y)=B(ry)$, and hence
$\mathsf{h}_{ry}(ry')=\mathsf{h}_y(y')$. The cubes of $\pi_j$ have side $M=L^{m'}$ and centres in
$\mathbb{T}^{(j+m')}$ \cite{B-CMP109}; since $M$ is odd, each of them is a union of level-$j$ unit
blocks. The cube $\square_y$ contains $y$, and the only unit block containing $y$ is $B(y)$, the
block centred at $y$; hence $B(y)\subset\square_y$.

\emph{The running profile.} The correlation theorem, in its bound on the running source profile
at level $k$, gives $N_k\in[\tfrac23,\tfrac43]$ and, for every bounded function $g$ on the
plaquettes,
\begin{equation*}
|\dot\zeta_k[g]|\le\tfrac43\,\sup|g| .
\end{equation*}

\emph{The ceilings.} The last condition in \eqref{8.5:eq-block-indicators-a} follows from the
inequality $x_i\ge x_k-2c_2$ ($i\le k$) of the correlation theorem together with the correlation
theorem's statement on the cut $g\mapsto gT(g)$ on $]0,\gamma_J]$; the same statement gives that
$gT(g)$ is nondecreasing on $]0,\gamma_J]$. Since $g_k<g$, with $g$ the reference coupling, and
$gT(g)$ is nondecreasing, each of the
conditions on the coupling $g$ in \eqref{8.5:eq-block-indicators-a} and
\eqref{8.5:eq-block-indicators-b} that
holds at $g$ holds with $g$ replaced by $g_k$, for every level $k$; in particular
$\theta^{\circ}_k=\theta^{\circ}(g_k)\le\theta^{\circ}(g)$. Each of them is therefore one ceiling
on the coupling $g$.
\end{proof}

\begin{remark}\label{8.5:thm-linear-response}
The statement of this theorem is not written out here; parts of its proof are printed below.
\end{remark}

\begin{proof}[Proof of Theorem~\ref{8.5:thm-linear-response}, continued]
\label{8.5:prf-linear-response-proof}
We continue the induction on $j=k+1$. The unmarked potentials $V^u$, the marked families
$\mathcal{M}_y[g]$ with their units of the kinds (m1), (m2), (m5), (m6), the first coefficients
$t^{[1]}_y(X_0)[g]$ with their domains $X=X_0\cup\square_y^1$, and the functionals
$\dot{P}^{(k+1)}(X,\cdot,y)[g]$ are those introduced above in this proof, and clause (a) of
Theorem~\ref{8.5:thm-linear-response} holds at level $k+1$. We prove clauses (b), (d) and (c)
at level $k+1$, in this order.

\smallskip
\emph{Clause} (b)\emph{: the bound on $F^{(k+1)}$.} Take $g\equiv 1$ and read off the units of
$\mathcal{M}_y[1]$.
\begin{itemize}
\item[(m2)] These units vanish. By definition
\[
\mathsf{D}^{(i)}(X',\cdot,y')[1]=\dot{P}^{(i)}(X',\cdot,y')[1]-1\cdot F^{(i)}(X',\cdot,y')=0 ,
\]
so every single of kind (m2) is zero.
\item[(m1)] The coefficient of the group $\mathfrak{i}_{y'}[F^{(i)}]$ is
$\mathsf{h}_y(y')g(p_{y'})=\mathsf{h}_y(y')$. Its modulus $|c|$ is at most $1$, and the group
norm is $\mathfrak{n}^G_i=\mathfrak{n}_i$.
\item[(m5)] The Wilson coefficient bound is $\bar c=4/3$.
\item[(m6)] The directly given activity satisfies
\begin{equation*}
\tilde a^{\mathrm{dir}}\le \tfrac43 C_6L^4C_1^2\delta_k^2\le \tfrac43 C_6M^4C_1^2\delta_k^2 ,
\end{equation*}
the second inequality because $L\le M$.
\end{itemize}
We apply the unit-sum bound for a marked family (Lemma~\ref{8.5:lem-unit-sum-bound},
display~\eqref{8.5:eq-unit-sum-bound-a}) at the absolute radius system $a=\circ$. There all
ratios $\rho_{k,j}$ of field radii between levels equal $1$. Together with the definition
$\tilde a:=C_AW^m+\tilde a^{\mathrm{dir}}$ of Lemma~\ref{8.5:lem-direct-potentials-bis}, this gives
\begin{equation}\label{8.5:eq-linear-response-proof-1}
\tilde a\le C_A\Big[2C_U^{+}\theta^{\circ}_kS_I^{\circ}\bar{\mathfrak{n}}_k
+\tfrac{32}{3}\theta^{\circ}_kC_5\mathfrak{a}_{\circ}^2\Big]
+\tfrac43 C_6M^4C_1^2\delta_k^2,
\qquad \bar{\mathfrak{n}}_k:=\max_{i\le k}\mathfrak{n}_i .
\end{equation}
Now $\delta_k^2\le\delta_k$ and $\delta_k=\theta^{\circ}_k\alpha_*/K_R$, with the constant
$K_R$ entering the definition of $\theta^{\circ}_k$. Every term
of~\eqref{8.5:eq-linear-response-proof-1} therefore carries the factor $\theta^{\circ}_k$, and
\begin{align*}
\tilde a&\le\theta^{\circ}_k\Big[2C_AC_U^{+}S_I^{\circ}\bar{\mathfrak{n}}_k
+\tfrac{32}{3}C_AC_5\mathfrak{a}_{\circ}^2+\tfrac43 C_6M^4C_1^2\,\alpha_*/K_R\Big]\\
&\le\theta^{\circ}_k\big[2C_AC_U^{+}S_I^{\circ}\bar{\mathfrak{n}}_k+c_V\big],
\end{align*}
with the constant $c_V$ of Notation~\ref{8.5:not-block-indicators}.

By Lemma~\ref{8.5:lem-direct-potentials-bis}, display~\eqref{8.5:eq-direct-potentials-a-bis}, each first
coefficient $t^{[1]}_y(X_0)[1]$ is bounded by
$\tfrac12\mathsf{K}\tilde a\,e^{-\kappa d_{k+1}(X_0)}$. The functional
$F^{(k+1)}(X,\cdot,y)=\dot{P}^{(k+1)}(X,\cdot,y)[1]$ is the sum of the coefficients
$t^{[1]}_y(X_0)[1]$ over the $X_0\in D_{k+1}$ with $X_0\cup\square_y^1=X$. There are at most
$2^{3^d}$ of them, because each such $X_0$ contains $X\setminus\square_y^1$. For each of them
$d_{k+1}(X)\le d_{k+1}(X_0)+2^d+4$, hence
$e^{-\kappa d_{k+1}(X_0)}\le e^{(2^d+4)\kappa}e^{-\kappa d_{k+1}(X)}$. The factor
$\tfrac12\cdot 2^{3^d}e^{(2^d+4)\kappa}$ is collected in the constant
$\mathsf{c}_{\square}$ of Notation~\ref{8.5:not-block-indicators}. This gives
$|F^{(k+1)}(X,\cdot,y)|\le\mathfrak{n}_{k+1}e^{-\kappa d_{k+1}(X)}$ with
\begin{equation}\label{8.5:eq-linear-response-proof-2}
\mathfrak{n}_{k+1}\le\mathsf{c}_{\square}\mathsf{K}\theta^{\circ}_k
\big[2C_AC_U^{+}S_I^{\circ}\bar{\mathfrak{n}}_k+c_V\big].
\end{equation}

\smallskip
\emph{Clause} (b)\emph{: the bound on $\mathsf{D}^{(k+1)}$.} Both $\dot{P}^{(k+1)}(X,\cdot,y)[g]$
and $F^{(k+1)}(X,\cdot,y)$ are built by the same construction. The first coefficients are linear
in the marked family. Hence
\[
\mathsf{D}^{(k+1)}(X,\cdot,y)[g]
=\dot{P}^{(k+1)}(X,\cdot,y)[g]-g(p_y)F^{(k+1)}(X,\cdot,y)
\]
is the same construction applied to the marked family that consists of the units of
$\mathcal{M}_y[g]$ together with the units of $\mathcal{M}_y[1]$ multiplied by $-g(p_y)$.
Combining equal units, its units are as follows.
\begin{itemize}
\item[(m1)] The group $\mathfrak{i}_{y'}[F^{(i)}]$ carries the coefficient
$\mathsf{h}_y(y')[g(p_{y'})-g(p_y)]$. Its modulus is at most
$3\operatorname{Lip}_{k+1}(g)$, since $|\mathsf{h}_y(y')|\le 1$ as before and
$\operatorname{dist}_{k+1}(x(p_{y'}),x(p_y))\le 1+1+1$.
\item[(m2)] These units are unchanged, since the corresponding units of $\mathcal{M}_y[1]$ vanish.
By clause (b) of the induction hypothesis at level $i\le k$,
\begin{equation*}
|\mathsf{D}^{(i)}(X',\cdot,y')[g]|
\le\mathfrak{d}_iL^{-(k+1-i)}\operatorname{Lip}_{k+1}(g)\,e^{-\kappa d_i(X')} .
\end{equation*}
\item[(m5)] The Wilson weight is $[\dot\zeta_k[g]-g(p_y)N_k]h_y$. On $\operatorname{supp}h_y$,
\begin{align*}
|\dot\zeta_k[g](p)-g(p_y)\dot\zeta_k[1](p)|
&\le|g(p)-g(p_y)|
+\sum_i|b_i'(0)|\sum_{y''}h^{(i)}_{y''}(p)\,|g(p_{y''})-g(p_y)|\\
&\le\tfrac43\cdot 3\operatorname{Lip}_{k+1}(g).
\end{align*}
\item[(m6)] The coefficient of the direct potential $G_y(B')$ is
\begin{equation*}
\dot\zeta_k[g](p_y)-g(p_y)N_k
=\sum_ib_i'(0)\sum_{y''}h^{(i)}_{y''}(p_y)\,[g(p_{y''})-g(p_y)] .
\end{equation*}
Its modulus is at most $\tfrac13\cdot 3\operatorname{Lip}_{k+1}(g)$.
\end{itemize}

It remains to see that only the values of $g$ on $X^{\oplus}$ enter. Recall that $t^{[1]}_y(X_0)$
is the sub-sum of
the cluster expansion carrying exactly one marked piece, and that this piece lies in a polymer
$Y\subset X_0$. Hence $t^{[1]}_y(X_0)$ depends only on the marked units $v$ whose support
$S_v$ lies in $X_0$. We therefore apply Lemma~\ref{8.5:lem-unit-sum-bound} and
Lemma~\ref{8.5:lem-direct-potentials-bis} to this sub-family, whose coefficients read $g$ on
$X^{\oplus}$ only. Every Lipschitz modulus above may thus be taken to be
$\operatorname{Lip}_{k+1}(g;X^{\oplus})$.

We divide by $\operatorname{Lip}_{k+1}(g;X^{\oplus})$ and use
$\sum_{i\le k}L^{-(k+1-i)}\le1$. The four kinds then enter the bound of the case $g\equiv1$ as
follows:
\begin{itemize}
\item[(m1)] $3\bar{\mathfrak{n}}_k$ takes the place of $\bar{\mathfrak{n}}_k$;
\item[(m2)] the singles contribute
$\bar{\mathfrak{d}}_k:=\max_{i\le k}\mathfrak{d}_i$;
\item[(m5)] the Wilson coefficient bound becomes $\bar c=4=3\cdot\tfrac43$;
\item[(m6)] the coefficient bound $1$ satisfies $1\le3\cdot\tfrac43$.
\end{itemize}
The bound
\[
|\mathsf{D}^{(k+1)}(X,\cdot,y)[g]|\le\mathfrak{d}_{k+1}
\operatorname{Lip}_{k+1}(g;X^{\oplus})\,e^{-\kappa d_{k+1}(X)}
\]
therefore holds with
\begin{equation}\label{8.5:eq-linear-response-proof-3}
\mathfrak{d}_{k+1}\le\mathsf{c}_{\square}\mathsf{K}\theta^{\circ}_k
\big[2C_AC_U^{+}\big(3S_I^{\circ}\bar{\mathfrak{n}}_k+\bar{\mathfrak{d}}_k\big)+3c_V\big].
\end{equation}

\smallskip
\emph{Closing the induction.} Put
$\mathfrak{x}_k:=\max\{\bar{\mathfrak{n}}_k,\bar{\mathfrak{d}}_k\}$. The right-hand side
of~\eqref{8.5:eq-linear-response-proof-2} is at most that
of~\eqref{8.5:eq-linear-response-proof-3}. By~\eqref{8.5:eq-block-indicators-b}, and with
$C_{LR}=6\mathsf{c}_{\square}\mathsf{K}c_V$,
\begin{align*}
\max\{\mathfrak{n}_{k+1},\mathfrak{d}_{k+1}\}
&\le\mathsf{c}_{\square}\mathsf{K}\theta^{\circ}_k
\big[2C_AC_U^{+}(3S_I^{\circ}+1)\mathfrak{x}_k+3c_V\big]\\
&\le\tfrac12\mathfrak{x}_k+\tfrac12C_{LR}\theta^{\circ}_k .
\end{align*}
By the induction hypothesis $\mathfrak{x}_k\le C_{LR}\hat\theta^{\circ}_{k-1}$. Moreover
$\hat\theta^{\circ}$ is the running maximum of $\theta^{\circ}$. The right-hand side is therefore
at most $C_{LR}\hat\theta^{\circ}_k$, and so is $\mathfrak{x}_k$. Hence
\[
\mathfrak{x}_{k+1}=\max\{\mathfrak{x}_k,\mathfrak{n}_{k+1},\mathfrak{d}_{k+1}\}
\le C_{LR}\hat\theta^{\circ}_k .
\]
This is clause (b) at level $k+1$.

The smallness condition $W^m\le\nu_0\alpha_4/(8C_A)$ on the marked weight, under which the
cluster expansion with a marked potential is set up, imposes no condition here. The coefficient
$t^{[1]}$ is homogeneous of degree one in the marked family. One may therefore apply the bounds
above to the marked family multiplied by a small positive number and divide the result by that
number.

\smallskip
\emph{Clause} (d)\emph{.} On the finite torus the output kernels $t_\lambda(X_0)$ of the step run
on $\{V^u\}\cup\{\lambda\mathcal{M}_y[g]\}$ satisfy
\begin{equation}\label{8.5:eq-linear-response-proof-4}
\sum_{X_0}t_\lambda(X_0)
=\log\int d\mu_{C^{(k)}}\,\chi\,\exp\sum_Y\big(V^u+\lambda V^m\big)(Y).
\end{equation}
By the lemma identifying the output kernels of the step with the kernels of the cluster
expansion, and by the statement of the cluster expansion with a marked potential that its
identity holds for every value of the marking parameter, this identity, which is
\cite[(2.11)--(2.13)]{B-CMP116}, holds on the finite torus for every $\lambda$. The derivative of
the left-hand side with respect to $\lambda$ at $\lambda=0$ is the sum of the first coefficients,
$\sum_{X_0}t^{[1]}_y(X_0)[g]$. Each $X_0$ determines exactly one $X=X_0\cup\square_y^1$, so this
sum equals $\sum_{X\ni y}\dot{P}^{(k+1)}(X,\cdot,y)[g]$. The sum $\sum_YV^u(Y)$ is the exponent of
the fluctuation expectation $\langle\cdot\rangle^U_k$. The $\lambda$-derivative at $0$ of the
right-hand side of~\eqref{8.5:eq-linear-response-proof-4} is therefore
\begin{equation*}
\frac{\int d\mu_{C^{(k)}}\,\chi\,\big(\sum_YV^m_{\mathcal{M}_y[g]}(Y)\big)
\exp\sum_YV^u(Y)}{\int d\mu_{C^{(k)}}\,\chi\,\exp\sum_YV^u(Y)}
=\Big\langle\sum_YV^m_{\mathcal{M}_y[g]}(Y)\Big\rangle^U_k .
\end{equation*}
We sum over $y$ and use the corresponding display. For every real level-$(k+1)$ minimal
configuration $U$ this gives
\begin{equation*}
\sum_{X,y}\dot{P}^{(k+1)}(X,U,J(U),y)[g]=\big\langle\dot{\mathfrak{B}}_k[g]\big\rangle^U_k
=\dot{\mathcal{T}}^{(k+1)}[g](U),
\end{equation*}
the last equality being the recursion~\eqref{8.5:eq-local-source-split-c}. This is clause (d) at
level $k+1$.

\smallskip
\emph{Clause} (c)\emph{.} Conditions (f1) and (f2) of the definition of a pointed $E$-family,
whose conditions are (f1)--(f3), follow from clause (a).

For (f3), let $U$ be a real regular configuration. Carry out the computation of clause (d) with
$g\equiv1$, without the sum over $y$. It shows that the pointed total is
\begin{equation*}
F^{(k+1)}(U,y):=\sum_{X\ni y}F^{(k+1)}(X,U,J(U),y)
=\big\langle\dot{\mathfrak{B}}^{(y)}_k\big\rangle^U_k ,
\end{equation*}
where
\begin{equation}\label{8.5:eq-linear-response-proof-5}
\dot{\mathfrak{B}}^{(y)}_k:=\Big\{A(N_kh_y,\cdot)
+\sum_{i\le k}\sum_{y'\in\mathbb{T}^{(i)}}\mathsf{h}_y(y')\,\mathfrak{i}_{y'}[F^{(i)}]\Big\}
\Big|^{U_k(e^{\mathrm{i}B'}V^{(k)})}_{U_{k+1}}+N_kG_y(B').
\end{equation}
No evaluation plaquette $p_y$ occurs in~\eqref{8.5:eq-linear-response-proof-5}. Let $\sigma$ be a
Euclidean symmetry of the torus preserving $\mathbb{T}^{(k+1)}$. We check four invariances.
\begin{itemize}
\item[($\alpha$)] The fluctuation expectation $\langle\cdot\rangle^U_k$ is covariant in $U$. In
\cite[(2.1), (2.18)]{B-CMP109} it is stated: ``By their definitions the expressions in (2.1) are
invariant with respect to these transformations \dots\ The transformation (2.18) generates an
orthogonal transformation of the fluctuation field $B'$, hence all the remaining operations
preserve the invariance.'' The inputs $\mathsf{P}_0$ are invariant
in total, because members of $\mathbb{B}_k$ satisfy (f3).
\item[($\beta$)] We have $A(N_kh_{\sigma y},\sigma U)=A(N_kh_y,U)$. This follows because the
plaquette term $1-\operatorname{Re}\operatorname{tr}U(\partial p)$ of the weighted Wilson action
(Definition~\ref{1.1:def-wilson-action}) takes at $\sigma p$ and $\sigma U$ the value it takes
at $p$ and $U$, and because $h_{\sigma y}=h_y\circ\sigma^{-1}$.
\item[($\gamma$)] We have
$\mathfrak{i}_{\sigma y'}[F^{(i)}](\sigma U)=\mathfrak{i}_{y'}[F^{(i)}](U)$. Indeed,
$\mathfrak{i}_{y'}[F^{(i)}]$ is the pointed total of $F^{(i)}$, minus its value at the
configuration $1$, minus $\beta[F^{(i)}]A(h_{y'},\cdot)$. The pointed total is invariant by (f3)
at level $i$, since $\sigma$ preserves $\mathbb{T}^{(i)}$, which contains $\mathbb{T}^{(k+1)}$.
The value at $1$ is invariant because $\sigma1=1$. The last term is invariant by ($\beta$).
Finally, $\mathsf{h}_{\sigma y}(\sigma y')=\mathsf{h}_y(y')$.
\item[($\delta$)] We have $G_{\sigma y}(\sigma B')=G_y(B')$. \cite{B-CMP109} states the invariance
of $G=\sum_yG_y$. Each $G_y$ reads $B'$ on the block $B(y)$ only, the blocks are disjoint,
$G_y(0)=0$, and $\sigma$ maps $B(y)$ onto $B(\sigma y)$. Let $B'$ be supported in $B(y)$. Then
$G(B')=G_y(B')$ and $G(\sigma B')=G_{\sigma y}(\sigma B')$, so the invariance of the sum gives
the claim for such $B'$. For general $B'$ it follows by restricting $B'$ to $B(y)$, since
$\sigma B'$ restricted to $B(\sigma y)$ is $\sigma$ applied to that restriction.
\end{itemize}
By ($\alpha$)--($\delta$), $F^{(k+1)}(\sigma U,\sigma y)=F^{(k+1)}(U,y)$. This is the invariance
(f3), in the form \cite[(2.29)]{B-CMP119}.

It remains to identify $\beta[F^{(k+1)}]$. By clause (d) with $g\equiv1$ and by clauses (iii)
and (v) of Proposition~\ref{8.5:prop-local-source-split} (see~\eqref{8.5:eq-local-source-split-d}),
\begin{equation*}
\sum_yF^{(k+1)}(U,y)=\dot{\mathcal{T}}^{(k+1)}[1](U)
=\partial_\zeta\big|_{\zeta=0}\sum_{X}P^{\flat(k+1)}_{\zeta}(X,U)+c ,
\end{equation*}
with a constant $c$ independent of $U$. We use the lemma identifying the coefficient $\beta$ of a
family through its total, which has three parts: $\beta$ reads the total only; it is linear;
and for a pointed family it is the number which \cite{B-CMP119} states ``coincides with the
definition (I.1.22)'', that is, with the definition \cite[(1.22)]{B-CMP109}. This lemma gives
\begin{equation*}
\beta[F^{(k+1)}]=\partial_\zeta\beta[P^{\flat(k+1)}_{\zeta}]\big|_0=\mathsf{b}^P_{k+1}.
\end{equation*}
That lemma asks the totals to agree on real regular configurations of every torus. This holds
here: the construction of Proposition~\ref{8.5:prop-local-source-split} and that of this theorem
are the same on every torus whose side is divisible as required for the block lattices of
Definition~\ref{1.5:def-block-average},
and their non-wrapping terms are torus-independent by clause (a). This proves clause (c) at level
$k+1$ and completes the induction.
\end{proof}

\begin{remark}
Three comments on the proof.
\begin{enumerate}
\item The source radius occurs nowhere in Theorem~\ref{8.5:thm-linear-response}, because a
first-order statement needs no source radius. Of the correlation theorem the proof uses only
three things: the numbers $b_j'(0)$, through $\sum_j|b_j'(0)|\le\tfrac13$; the membership
$\mathsf{P}_0\in\mathbb{B}$; and the couplings.
\item The gain is the room $\mathsf{w}^{\circ}$, of the order of $\alpha_*/(K_R\delta_k)$. Every
output is the expectation of a bracket that vanishes with the fluctuation field. The one-loop
term, which is the reason why the constant $E_0$ is of order one, is independent of the source,
by the lemma stating that the one-loop term does not depend on the source. It therefore never
enters.
\item No analytic continuation of the correlation theorem's own pointed and local pure families
to the spaces $\mathcal{U}_j(X)$ is used. By~\eqref{8.5:eq-remainder-level-sum-a} only the first
derivative of the total enters. Clause (d), together with clause (v) of
Proposition~\ref{8.5:prop-local-source-split}, identifies it.
\end{enumerate}
\end{remark}

\begin{definition}[Configurations regular at a given scale]\label{8.5:def-regular-configurations}
Let $j$ and $n$ be levels with $n\ge j$, and put
\begin{equation*}
t:=L^{j-n},\qquad \nu:=1+n-j .
\end{equation*}
Let $\mathfrak{e}>0$. A real configuration $\mathfrak{b}$ on the fine lattice of the torus,
that is, one with values in $\mathrm{SU}(N)$, is called
\emph{$\mathfrak{e}$-regular at scale $n$} if the conditions (c1)--(c3) under which the
correlation theorem admits backgrounds at scale $n$ hold for every box $\square_0$, with the
regularity size $\mathfrak{a}_n$ of Notation~\ref{8.5:not-field-radii} replaced by $\mathfrak{e}$.
The box $\square_0$ is taken with a margin of two layers of level-$j$ blocks of the partition
$\pi_j$, or, in the form used in Ba{\l}aban's construction, with a margin of one layer about the
enlarged box $\square^{\sim\nu}$.

Written out, the conditions require the following for every such box $\square_0$.
\begin{itemize}
\item Modulo a gauge transformation, the restriction of $\mathfrak{b}$ to the domain attached to
$\square_0$ in those conditions coincides with the configuration $U_j(\square_0,\mathsf{V})$
built there from box data $\mathsf{V}$, and these satisfy
\begin{equation}\label{8.5:eq-regular-configurations-1}
|\partial\mathsf{V}-1|\le \mathfrak{e}\,t^{2}.
\end{equation}
\item In the comb gauge centred at $y$, the potential $B^{\mathrm{ax}}$ satisfies, for every bond
$b$,
\begin{equation}\label{8.5:eq-regular-configurations-2}
|B^{\mathrm{ax}}(b)|\le 2d\,(1+|b-y|)\,\mathfrak{e}\,t^{2}.
\end{equation}
\item In the same gauge, $B^{\mathrm{ax}}$ has, with a remainder $\eta_{\mu}(x)$, the form
\begin{equation}\label{8.5:eq-regular-configurations-3}
B^{\mathrm{ax}}_{\mu}(x)=\sum_{\sigma<\mu}(x_{\sigma}-y_{\sigma})\,F_{\sigma\mu}(y)+\eta_{\mu}(x),
\end{equation}
where
\begin{equation}\label{8.5:eq-regular-configurations-4}
|F|\le \mathfrak{e}\,t^{2},\qquad
|\eta_{\mu}(x)|\le C\,(1+|x-y|)^{2}\,\mathfrak{e}\,t^{2+\hat\alpha},
\end{equation}
with the constant $C$ and the exponent $\hat\alpha$ of those conditions.
\end{itemize}

Note that for real configurations the restriction $D\le 5M/t$ on the diameter $D$ of the domain,
contained in the third of those conditions, (c3), is void. Its only use is to guarantee that
conjugations distort by a factor at most $e^{O(M\alpha_1)}\le 2$, and conjugation by unitary
matrices does not distort at all.
\end{definition}

\begin{lemma}[Minimisers of small-curvature data are regular]\label{8.5:lem-minimiser-regular}
Let $V$ be a real configuration on the block lattice $\mathbb{T}^{(n)}$ of
Definition~\ref{1.5:def-block-average}, and let $\varepsilon'>0$ satisfy
\begin{equation*}
\sup_p\,\lvert V(\partial p)-1\rvert\le\varepsilon'\le a_1 ,
\end{equation*}
the supremum running over the plaquettes $p$ of $\mathbb{T}^{(n)}$. Here $a_1=a_1(d,L)$ is the
smallness constant of \cite[Theorem~1]{B-CMP102b}. Then the background configuration $U_n(V)$
(Definition~\ref{1.5:def-domains}) is $\mathfrak{e}$-regular at scale $n$ in the sense of
Definition~\ref{8.5:def-regular-configurations}, with
\begin{equation}\label{8.5:eq-minimiser-regular-a}
\mathfrak{e}:=C_{\mathrm{reg}}\,\varepsilon',
\end{equation}
where $C_{\mathrm{reg}}=C_{\mathrm{reg}}(d,L,M)$ is the constant fixed at the end of the proof.
\end{lemma}

\begin{proof}
Fix $j\le n$, put $t:=L^{j-n}$, and fix a box $\square_0$ and the region $X$ attached to it in
Definition~\ref{8.5:def-regular-configurations}. We verify conditions (c1)--(c3) of
Definition~\ref{8.5:def-regular-configurations} for $\square_0$ with the
size $\mathfrak{e}$ of \eqref{8.5:eq-minimiser-regular-a}.

First, the hypothesis $\varepsilon'\le a_1$ is the smallness condition of
\cite[Theorem~1]{B-CMP102b}. That theorem provides, for the data $V$, a minimal gauge
orbit in Ba{\l}aban's space of configurations that
are regular of size $B_3\varepsilon'$ with respect to the sequence of domains of
that theorem; $U_n(V)$ is a representative of this orbit. The constants $a_1$ and $B_3$ depend on
$d$ and $L$ only. The regularity bounds \cite[(9)--(10)]{B-CMP102b} of that theorem are linear in
the size parameter $\varepsilon_1$ of the data, here $\varepsilon'$, and their constant $B_4$
depends on $d$, $L$ and the H\"older exponent $\hat\alpha$.

Second, we use the restriction lemma for minimisers stated with the correlation theorem; it
rests on property~(iv) of the minimal configurations listed in \cite{B-CMP109}: the restriction
of a minimiser to a box with margin is the box minimiser of its
own data. Applied to $U_n(V)$ and $\square_0$, it gives, modulo gauge, that the restriction of
$U_n(V)$ to $X$ equals $U_j(\square_0,\mathsf{V})$, where $\mathsf{V}$ is the data of $U_n(V)$ on
$\square_0$, namely its level-$j$ averages.

Third, \cite[p.~263]{B-CMP109} states that if a configuration $U$ on a lattice of spacing $\xi$
satisfies $\lvert\partial U-1\rvert<\alpha_0'\xi^2$, then its averages $\bar U^{j}$ over $j$
blocking steps obey
\begin{equation*}
\lvert\partial\bar U^{j}-1\rvert<2\alpha_0'(L^{j}\xi)^2 .
\end{equation*}
By the first step, $U_n(V)$ is regular of size $B_3\varepsilon'$; membership in that space
includes the plaquette bound $\lvert\partial U_n(V)-1\rvert\le B_3\varepsilon'L^{-2n}$ on the
fine lattice, of spacing $L^{-n}$, which is the hypothesis of the averaging bound.
We apply this to $U_n(V)$ with $\alpha_0'=B_3\varepsilon'$ and $\xi=L^{-n}$, so that
$L^{j}\xi=t$. This gives
$\lvert\partial\mathsf{V}-1\rvert\le 2B_3\varepsilon't^2$. Together with the second step, this is
(c1) with the constant $2B_3$.

The comb-gauge bounds in (c2) and (c3), centred at $y$, are read from the regularity bounds
\cite[(9)--(10)]{B-CMP102b}. These bounds are linear in $\varepsilon_1$, so each of them holds with
a constant depending on $d$, $L$ and $M$ times $\varepsilon'$. In particular, the H\"older
remainder $\mathfrak{e}_{\mu}$ is bounded by \cite[(9)]{B-CMP102b}, and the constant produced for
(c3) is $B_4Mc_d$, where $c_d$ is a dimensional constant.

Define $C_{\mathrm{reg}}$ as the largest of the constants produced in this way, among them $2B_3$
for (c1) and $B_4Mc_d$ for (c3). Then (c1)--(c3) hold with the size
$\mathfrak{e}=C_{\mathrm{reg}}\varepsilon'$ for the box $\square_0$. Since $\square_0$ was
arbitrary, $U_n(V)$ is $\mathfrak{e}$-regular at scale $n$.
\end{proof}

\begin{lemma}[Single terms on small real backgrounds]\label{8.5:lem-single-terms}
Let $d=4$, let $0\le j\le n$ and $t=L^{j-n}$. Let $X$ be a set of level $j$, and let
$F(X,\cdot)$ have properties \textup{(f1)} and \textup{(f2)} of the format of the correlation
theorem on the complex background space $\mathcal{U}^{c}_j(X,\alpha_j)$, with
$|F|\le\mathfrak{n}$ there. Let $\mathfrak{b}$ be a real configuration that is
$\mathfrak{e}$-regular at scale $n$ in the sense of
Definition~\ref{8.5:def-regular-configurations}, and put
$\varrho_j:=\mathfrak{e}/(C_{\mathfrak{a}}\alpha_{*,j})$. Put
\begin{equation}\label{8.5:eq-single-terms-1}
D_X:=\operatorname{diam}_jX+5M,\qquad
C_M:=\frac{2d\,C_{\mathfrak{a}}}{c_3},
\end{equation}
where $\operatorname{diam}_jX$ is the diameter of $X$ in level-$j$ lattice units and $c_3$ is
the constant of condition \textup{(c2)} of Definition~\ref{8.5:def-regular-configurations};
note that $D_X\le 4M\bigl(d_j(X)+5\bigr)$. Let $\Sigma_n$ denote the side of the torus
$\mathbb{T}^{(n)}$. Then for every $X$
whose own box fits in the torus, that is, with $D_X\le\tfrac12\Sigma_n t^{-1}$,
\begin{equation}\label{8.5:eq-single-terms-a}
\bigl|F(X,\mathfrak{b})-F(X,1)\bigr|
\le 2\mathfrak{n}\,\min\Bigl\{1,\bigl(C_M D_X t^2\varrho_j\bigr)^2\Bigr\}.
\end{equation}
For every other $X$ (those whose box wraps around the torus) the lemma asserts only the
bound $|F(X,\mathfrak{b})-F(X,1)|\le 2\mathfrak{n}$.
\end{lemma}

\begin{proof}
The bound $2\mathfrak{n}$, which is the only assertion for wrapping $X$ and the first member
of the minimum in \eqref{8.5:eq-single-terms-a}, follows from the triangle inequality and
$|F|\le\mathfrak{n}$:
\begin{equation*}
|F(X,\mathfrak{b})-F(X,1)|\le|F(X,\mathfrak{b})|+|F(X,1)|\le 2\mathfrak{n}.
\end{equation*}

Let now $D_X\le\tfrac12\Sigma_n t^{-1}$. We repeat the argument of the lemma in the proof of
the correlation theorem that compares single terms at an admissible background with the
trivial background. Here it is carried out on the box $\square_0$ of $X$ itself, which fits
in the torus by the assumption on $D_X$. No containment of $X$ in a fixed enlarged cube is
needed, because only one zone occurs and the field is real. By condition \textup{(c1)} of
Definition~\ref{8.5:def-regular-configurations}, modulo a gauge transformation,
$\mathfrak{b}$ restricted to $X$ is $U_j(\square_0,\mathsf{V})$. Writing the box data in the
comb gauge centred at a point $y$ with the potential $B^{\mathrm{ax}}$, we put, for complex
$\sigma$,
\begin{equation*}
\varphi(\sigma):=F\Bigl(X,\,U_j\bigl(\square_0,e^{i\sigma B^{\mathrm{ax}}}\bigr),\,
J_j\bigl(\square_0,e^{i\sigma B^{\mathrm{ax}}}\bigr)\Bigr),
\qquad
S:=\frac{c_3\,\alpha_{*,j}}{\sup|B^{\mathrm{ax}}|}.
\end{equation*}
Here $J_j(\square_0,\cdot)$ is the companion of $U_j(\square_0,\cdot)$ that the format of the
correlation theorem attaches to box data on $\square_0$: the pair
$\bigl(U_j(\square_0,\cdot),J_j(\square_0,\cdot)\bigr)$ fills the background argument of
$F(X,\cdot)$.
The function $\varphi$ has the following three properties.
\begin{itemize}
\item[(a)] By condition \textup{(c2)}, for $|\sigma|<S$ the arguments of $F$ in
$\varphi(\sigma)$ lie in $\mathcal{U}^{c}_j(X,\alpha_j)$, on which $F(X,\cdot)$ has
properties \textup{(f1)}, \textup{(f2)} and satisfies $|F|\le\mathfrak{n}$; accordingly
$\varphi$ is analytic on the disc $|\sigma|<S$ and satisfies $|\varphi|\le\mathfrak{n}$ there.
\item[(b)] By \textup{(c1)} and the gauge invariance \textup{(f2)}, $\varphi(1)=F(X,\mathfrak{b})$.
At $\sigma=0$ the configuration is the trivial one, so $\varphi(0)=F(X,1)$.
\item[(c)] By \cite[(4.14)]{B-CMP109}, $\varphi'(0)=0$.
\end{itemize}

We apply the Schwarz lemma to
$\psi(\tau):=\bigl(\varphi(S\tau)-\varphi(0)\bigr)/(2\mathfrak{n})$. This function is
analytic on the unit disc and satisfies $|\psi|\le1$ there, by (a). It also satisfies
$\psi(0)=0$, and $\psi'(0)=0$ by (c). Hence $\psi(\tau)/\tau^2$ is analytic on the unit
disc, and by the maximum principle on the circles $|\tau|=r$ it is bounded there by
$r^{-2}$ for every $r<1$, hence, letting $r\uparrow1$, by $1$. Thus
$|\psi(\tau)|\le|\tau|^2$ for $|\tau|<1$. If
$S>1$, taking $\tau=S^{-1}$ and using (b) gives
\begin{equation}\label{8.5:eq-single-terms-2}
|F(X,\mathfrak{b})-F(X,1)|=|\varphi(1)-\varphi(0)|\le 2\mathfrak{n}\,S^{-2}.
\end{equation}
If $S\le1$, then $S^{-2}\ge1$, and the first paragraph already gives
$2\mathfrak{n}\le 2\mathfrak{n}\min\{1,S^{-2}\}$. In either case
$|F(X,\mathfrak{b})-F(X,1)|\le 2\mathfrak{n}\min\{1,S^{-2}\}$.

It remains to bound $S^{-1}$. Condition \textup{(c3)} of
Definition~\ref{8.5:def-regular-configurations} bounds $\sup|B^{\mathrm{ax}}|$ over the box
of $X$ by $2dD_X\mathfrak{e}t^2$, whence
\begin{equation*}
S^{-1}\le\frac{2dD_X\mathfrak{e}t^2}{c_3\alpha_{*,j}}
=\frac{2dC_{\mathfrak{a}}}{c_3}\,D_Xt^2\,\frac{\mathfrak{e}}{C_{\mathfrak{a}}\alpha_{*,j}}
=C_MD_Xt^2\varrho_j,
\end{equation*}
by the definitions of $C_M$ in \eqref{8.5:eq-single-terms-1} and of $\varrho_j$. Since
$\min\{1,\cdot\}$ is non-decreasing, inserting this into the bound
$2\mathfrak{n}\min\{1,S^{-2}\}$ gives \eqref{8.5:eq-single-terms-a}.
\end{proof}

\begin{lemma}[Renormalised groups on small real backgrounds]\label{8.5:lem-renormalised-groups}
Let $F$ be a pointed $E$-family at scale $j$ on the radii $a_j$, with $\|F\|\le\mathfrak{n}$, and
let $\mathfrak{i}_y$ be its renormalised pointed group at the point $y$. Let $n\ge j$, $t=L^{j-n}$
and $\nu=1+n-j$, and let $\mathfrak{b}$ be a real configuration which is $\mathfrak{e}$-regular at
scale $n$ in the sense of Definition~\ref{8.5:def-regular-configurations}. Let
$\varrho_j=\mathfrak{e}/(C_{\mathfrak{a}}a_{*,j})$ and $C_M=2dC_{\mathfrak{a}}/c_3$ be as in
Lemma~\ref{8.5:lem-single-terms}. Assume that $\varrho_j\le\rho_{n,j}$, that is
$\mathfrak{e}\le\mathfrak{a}_n$, in the system $\flat$ of shrinking radii, and that
$\varrho_j\le1$ in the system $\circ$ of absolute radii. Assume moreover that
\begin{equation}\label{8.5:eq-renormalised-groups-b}
\sup_{x\ge g^{-2}-c_5}\,2C(\kappa)e^{5\kappa}\sup_{\tau\ge1}\tau^{9}
e^{-\kappa L^{11}(\log x)^{r_{\mathrm{pr}}}\tau/16}
\Bigl(\frac{\mathfrak{a}_{\circ}}{C_{\mathfrak{e}}\,\varepsilon(x^{-1/2})}\Bigr)^{2}\le1,
\end{equation}
where $g$ is the reference coupling and $c_5=2c_2$; the function
$\varepsilon(g')=g'A_0(\log g'^{-2})^{\mathfrak{p}_0}$ is the small-field size at coupling $g'$;
$r_{\mathrm{pr}}$ is the exponent of the logarithm in the function $R(\cdot)$ that enters the
torus exponent, written here as a function of the inverse coupling, for which
$r_{\mathrm{pr}}\ge2$ \cite{B-CMP119}; $\mathfrak{a}_{\circ}$ is the regularity size
$C_{\mathfrak{a}}\alpha$ of the system $\circ$ of absolute radii and $C_{\mathfrak{e}}$ the
constant relating $\mathfrak{e}$ to the window size, as in Notation~\ref{0.2:not-glossary}; and
$C(\kappa)$ is the constant of the bound
\begin{equation*}
\sum_{X\ni y,\ d_j(X)\ge d'}(d_j(X)+5)^{2}e^{-\kappa d_j(X)}\le C(\kappa)e^{-\kappa d'/2},
\end{equation*}
valid for every $d'\ge0$ by \cite[(1.26)]{B-CMP116}. Assume finally that
\begin{equation}\label{8.5:eq-renormalised-groups-2}
\varrho_j\ge\frac{C_{\mathfrak{e}}\,\varepsilon(x_n^{-1/2})}{\mathfrak{a}_{\circ}},
\end{equation}
with $x_n=g_n^{-2}$ the running inverse coupling at level $n$.
Then
\begin{equation}\label{8.5:eq-renormalised-groups-a}
|\mathfrak{i}_y(\mathfrak{b})|\le C_I'\,\mathfrak{n}\,\nu^{3}\varrho_j^{2}(1+\varrho_j)\,
t^{4+\hat\alpha},
\end{equation}
with a constant $C_I'=C_I'(d,L,M,\kappa)$ which does not depend on the radii.
\end{lemma}

In the system $\flat$ one has $\varrho_j=(\mathfrak{e}/\mathfrak{a}_n)\rho_{n,j}$, which may exceed
$1$ at very deep levels, exactly as $\rho_{n,j}$ may in the bound for renormalised pointed groups
at general radii in the proof of the correlation theorem; the factor $1+\varrho_j$
in~\eqref{8.5:eq-renormalised-groups-a} is the factor $1+\rho_{n,j}$ of that bound. The
expression under the outer supremum in~\eqref{8.5:eq-renormalised-groups-b} tends to zero, as
$x\to\infty$, faster than every power of $x^{-1}$, because $r_{\mathrm{pr}}\ge2$.

\begin{proof}
We follow the proof of the bound for renormalised pointed groups at general radii in the proof of
the correlation theorem, which gives
\begin{equation*}
|\mathfrak{i}_y(\mathfrak{b})|\le C_I\mathfrak{n}\nu^{3}\rho_{n,j}^{2}(1+\rho_{n,j})t^{4+\hat\alpha}
\end{equation*}
for backgrounds admissible at scale $n$. We make in it the substitutions
$\mathfrak{a}_n\mapsto\mathfrak{e}$ and $\rho_{n,j}\mapsto\varrho_j$. The estimates of the near
levels and of the core remain valid after these substitutions and are written out below; the
estimate of the tail of the sum over domains is replaced by a different one. Throughout,
$D=M(2\nu+3)$.

\emph{Near levels.} Suppose $C_MDt^{2}\varrho_j>\tfrac12$. Since $t=L^{-(\nu-1)}$, this reads
$C_MM(2\nu+3)L^{-2(\nu-1)}\varrho_j>\tfrac12$. The count used in the proof named above, which
rests on
\begin{equation*}
\varrho_j^{2}\le\rho_{n,j}^{2}\le1+(3\lambda\nu+2c_2)g^{2}+o(1)
\end{equation*}
in the system $\flat$, with $\lambda$ the linear flow constant of the upper bound on the running
inverse couplings used in the wrapping estimate below, then
restricts $\nu$ to $\nu\le\nu_*(L,M)$; in the system $\circ$ the hypothesis $\varrho_j\le1$ serves
in place of this count. As $t\le1$ and $\nu\le\nu_*$, we get
$C_MM(2\nu_*+3)\varrho_j\ge C_MDt^{2}\varrho_j>\tfrac12$, that is,
$\varrho_j>\varrho_*:=(2C_MM(2\nu_*+3))^{-1}$. The definition of the renormalised pointed group,
in which $\beta[F]$ is the coefficient of $F$ and $A(h_y,\cdot)$ the action weighted by the
profile $h_y$, gives the trivial bound $|\mathfrak{i}_y(\mathfrak{b})|\le2C(\kappa)\mathfrak{n}
+|\beta[F]|\,|A(h_y,\mathfrak{b})|$, and the right-hand side of this bound is at most
$C\mathfrak{n}(1+\varrho_j^{2})$.
Since $1<\varrho_*^{-2}\varrho_j^{2}$, we have $1+\varrho_j^{2}\le(1+\varrho_*^{-2})\varrho_j^{2}$;
since $4+\hat\alpha\le5$ and $\nu-1<\nu_*$, we have
$t^{-(4+\hat\alpha)}=L^{(4+\hat\alpha)(\nu-1)}\le L^{5\nu_*}$. Therefore
\begin{equation*}
|\mathfrak{i}_y(\mathfrak{b})|\le C\mathfrak{n}(1+\varrho_j^{2})
\le C(1+\varrho_*^{-2})L^{5\nu_*}\,\mathfrak{n}\varrho_j^{2}t^{4+\hat\alpha},
\end{equation*}
which is~\eqref{8.5:eq-renormalised-groups-a}, since $\nu^{3}(1+\varrho_j)\ge1$ and $\nu_*$ depends
on $(L,M)$ only. From now on $C_MDt^{2}\varrho_j\le\tfrac12$.

\emph{Tail.} In the proof named above the domains $X\ni y$ with $X\not\subset\square^{\sim\nu}$ were
bounded by
\begin{equation*}
2\mathfrak{n}e^{-\kappa d_j(X)}\le2\mathfrak{n}e^{-\kappa d_j(X)/2}t^{5},
\end{equation*}
a bound without a factor $\rho_{n,j}^{2}$; it sufficed there because $\rho_{n,j}\ge\tfrac12$ in
that range. Here $\varrho_j$ may
be small, and these domains are estimated as follows. Let $X\ni y$, $X\not\subset\square^{\sim\nu}$;
then $d_j(X)\ge\nu-1$. If $X$ is not wrapping, that is $D_X\le\tfrac12Nt^{-1}$ with $N$ the side of
$\mathbb{T}^{(n)}$ (the letter of Lemma~\ref{8.5:lem-single-terms}; here $N$ is not the rank of the
gauge group), Lemma~\ref{8.5:lem-single-terms}, in the form~\eqref{8.5:eq-single-terms-a}
together with $D_X\le4M(d_j(X)+5)$, gives
\begin{equation}\label{8.5:eq-renormalised-groups-1}
|F(X,\mathfrak{b},y)-F(X,1,y)|\le2\mathfrak{n}e^{-\kappa d_j(X)}
\bigl(4C_MM(d_j(X)+5)\bigr)^{2}t^{4}\varrho_j^{2}.
\end{equation}
By the definition of $C(\kappa)$, applied with $d'=\nu-1$, the sum over domains satisfies
\begin{equation*}
\sum_{X\ni y,\ d_j(X)\ge\nu-1}(d_j(X)+5)^{2}e^{-\kappa d_j(X)}\le C(\kappa)e^{-\kappa(\nu-1)/2}
\le C(\kappa)t^{5},
\end{equation*}
the last inequality because the lower bound $\kappa\ge10\log L$, imposed where $\kappa$ is chosen
in the order of choice (Notation~\ref{2.3:not-stages-middle}), gives
$e^{-\kappa(\nu-1)/2}\le L^{-5(\nu-1)}=t^{5}$. Summing~\eqref{8.5:eq-renormalised-groups-1}, the
non-wrapping tail domains contribute at most $32C_M^{2}M^{2}C(\kappa)\,\mathfrak{n}\varrho_j^{2}t^{9}$.

If $X$ is wrapping, $D_X>\tfrac12Nt^{-1}$, then $4M(d_j(X)+5)\ge D_X>\tfrac12Nt^{-1}$ forces
$d_j(X)\ge Nt^{-1}/(8M)-5$. Moreover
\begin{equation*}
\begin{split}
N\ge L^{m}\ge N_{\mathbb{T}}(g,m)&\ge C_{\mathbb{T}}R(g^{-2}+b_{+}+4c_2+3\lambda m)\\
&\ge L^{11}M\,R(x_{K'}).
\end{split}
\end{equation*}
Here $N_{\mathbb{T}}(g,m)$ is the torus side prescribed in the definition of the torus exponent,
$C_{\mathbb{T}}$ is the constant in that definition, $b_{+}$ is the additive constant in the upper
bound on the running inverse couplings, and $K'$ is the level at which that bound is read. The
first inequality holds because the side $N$ of $\mathbb{T}^{(n)}$ is at least $L^{m}$, the second
is the lower bound on the torus exponent $m$, the third is the definition of $N_{\mathbb{T}}$ with
the source parameter $u$ of that definition set to $0$, and the fourth uses the bound
$x_{K'}\le g^{-2}+b_{+}+4c_2+3\lambda m$ on the inverse coupling at level $K'$; here $n=K'$, so
that $R(x_{K'})=R(x_n)$. Hence $d_j(X)\ge L^{11}R(x_n)t^{-1}/8-5$. For wrapping domains
Lemma~\ref{8.5:lem-single-terms} claims only the trivial bound $2\mathfrak{n}e^{-\kappa d_j(X)}$ on
each term, and summing it over the wrapping $X\ni y$ gives
\begin{equation*}
\sum_{X\ni y\ \mathrm{wrapping}}2\mathfrak{n}e^{-\kappa d_j(X)}
\le2C(\kappa)\mathfrak{n}e^{5\kappa}e^{-\kappa L^{11}R(x_n)t^{-1}/16}
\le\mathfrak{n}\varrho_j^{2}t^{9},
\end{equation*}
the last inequality following from the condition~\eqref{8.5:eq-renormalised-groups-b} read at
$x=x_n$ and at $\tau=t^{-1}\ge1$. This reading uses three inequalities: $x_n\ge g^{-2}-c_5$, so
that $x_n$ lies in the range of the outer supremum; $(\log x_n)^{r_{\mathrm{pr}}}\le R(x_n)$,
through which $R(x_n)$ enters; and
\begin{equation*}
\varrho_j^{-2}\le\Bigl(\frac{\mathfrak{a}_{\circ}}{C_{\mathfrak{e}}\,\varepsilon(x_n^{-1/2})}\Bigr)^{2},
\end{equation*}
which is the hypothesis~\eqref{8.5:eq-renormalised-groups-2}, and through which
$\varrho_j^{-2}$ enters. Since $t\le1$,
$4+\hat\alpha\le9$ and $\nu^{3}(1+\varrho_j)\ge1$, both tail contributions are bounded by the
right-hand side of~\eqref{8.5:eq-renormalised-groups-a} with a constant depending on
$(d,L,M,\kappa)$.

\emph{Core.} Let $X\ni y$ with $X\subset\square^{\sim\nu}$, and let $\varphi_X$ be the function of
$\sigma$ defined in the proof of Lemma~\ref{8.5:lem-single-terms} for the term $F(X,\cdot,y)$, with
the comb-gauge potential $B^{\mathrm{ax}}$ of Definition~\ref{8.5:def-regular-configurations};
its contribution to $\mathfrak{i}_y(\mathfrak{b})$ is $\varphi_X(1)-\varphi_X(0)$. That proof gives
$\varphi_X'(0)=0$ and, for the radius $S$ of the disc of analyticity,
\begin{equation*}
S^{-1}\le\frac{2dD\mathfrak{e}t^{2}}{c_3a_{*,j}}=C_MDt^{2}\varrho_j\le\tfrac12,
\end{equation*}
so Taylor's formula at $\sigma=0$ gives
\begin{equation*}
\varphi_X(1)-\varphi_X(0)=\tfrac12\varphi_X''(0)
+O\bigl(\mathfrak{n}e^{-\kappa d_j(X)}(C_MD\varrho_j)^{3}t^{6}\bigr).
\end{equation*}
The second derivative $\varphi_X''(0)$ is a quadratic form in $B^{\mathrm{ax}}$. In what follows
$\mu'$ and $\mu$ range over directions (the direction index written $\kappa$ in
Definition~\ref{8.5:def-regular-configurations} is written $\mu'$ here, since $\kappa$ denotes the
decay constant). Replacing $B^{\mathrm{ax}}$ by its main term
\begin{equation*}
B^{\mathrm{main}}_{\mu}(x)=\sum_{\mu'<\mu}(x_{\mu'}-y_{\mu'})F_{\mu'\mu}(y),
\end{equation*}
that is, dropping
the remainder $\mathfrak{e}_\mu$ with
$|\mathfrak{e}_\mu(x)|\le C(1+|x-y|)^{2}\mathfrak{e}t^{2+\hat\alpha}$, costs at most
$C\mathfrak{n}e^{-\kappa d_j(X)}D^{3}\varrho_j^{2}t^{4+\hat\alpha}$. Next, the box average
$U_j(\square_0,\cdot)$ is replaced by the free $U_j$, and the sum over $X$ is extended from the
domains contained in $\square^{\sim\nu}$ to all domains $X\ni y$; these two operations change the
kernels of the quadratic forms by at most $e^{-\delta_0M\nu}$, with $\delta_0>0$ the rate of
exponential decay of the difference between the box average and the free average, and
$e^{-\kappa\nu/2}$ respectively,
each at most $t^{5}$, and the kernels act on $B^{\mathrm{main}}$, for which
$|B^{\mathrm{main}}|^{2}\le(D\mathfrak{e}t^{2})^{2}$. By the property (f3) of pointed $E$-families,
the quadratic part so obtained equals
\begin{equation*}
\beta[F]\cdot\tfrac12\sum_{\mu'<\mu}\operatorname{tr}F_{\mu'\mu}(y)^{2}.
\end{equation*}
On the other hand
\begin{align*}
A(h_y,\mathfrak{b})&=\tfrac12\sum_{\mu'<\mu}\operatorname{tr}F_{\mu'\mu}(y)^{2}
+O(\nu^{3}\mathfrak{e}^{2}t^{4+\hat\alpha}),\\
|\beta[F]|&\le C\mathfrak{n}(c_3a_{*,j})^{-2};
\end{align*}
since $\mathfrak{e}/a_{*,j}=C_{\mathfrak{a}}\varrho_j$, the difference between
$\beta[F]A(h_y,\mathfrak{b})$ and the quadratic part is at most
$C(C_{\mathfrak{a}}/c_3)^{2}\mathfrak{n}\nu^{3}\varrho_j^{2}t^{4+\hat\alpha}$.

Every error term of the core is proportional to $\varrho_j^{2}(1+\varrho_j)$: after summation over
$X$ with the weights $e^{-\kappa d_j(X)}$, and using $D\le5M\nu$, $t\le1$ and
$\varrho_j^{3}\le\varrho_j^{2}(1+\varrho_j)$, each of them is at most a constant depending on
$(d,L,M,\kappa)$ times $\mathfrak{n}\nu^{3}\varrho_j^{2}(1+\varrho_j)t^{4+\hat\alpha}$. Adding the
core, the non-wrapping tail and the wrapping tail gives~\eqref{8.5:eq-renormalised-groups-a} when
$C_MDt^{2}\varrho_j\le\tfrac12$; the near levels were treated above. None of the constants
depends on the radii $a_j$, $a_{*,j}$.
\end{proof}

\begin{lemma}[The regular extension lemma]\label{8.5:lem-regular-extension}
Let $Q'\subset Q\subset\mathbb{T}^{(n)}$ be concentric boxes in the lattice $\mathbb{T}^{(n)}$ of
Definition~\ref{1.5:def-block-average}. Let $R$ be the side of $Q$, and put
$\mathsf{m}:=\operatorname{dist}(Q',\partial Q)$. Assume $\mathsf{m}\ge 2$. Let $\varepsilon'>0$,
and let $V$ be a real configuration on the bonds of $Q$ such that
$|V(\partial p)-\mathbf{1}|\le\varepsilon'$ for every plaquette $p\subset Q$, and assume
\begin{equation}\label{8.5:eq-regular-extension-a}
d^2R^2\varepsilon'\le 1 .
\end{equation}
Then there is a real configuration $V^{\natural}$ on all bonds of $\mathbb{T}^{(n)}$ with the
following properties. It is gauge equivalent to $V$ on $Q'$. It equals $\mathbf{1}$ on every bond
not in $Q$. It satisfies
\begin{equation}\label{8.5:eq-regular-extension-1}
\sup_p|V^{\natural}(\partial p)-\mathbf{1}|\le C_{\mathrm{ext}}\,\varepsilon',\qquad
C_{\mathrm{ext}}:=137+\frac{8\pi dR}{\mathsf{m}} .
\end{equation}
Finally, $V^{\natural}$ is determined by the restriction of $V$ to the bonds of $Q$ up to one
constant rotation.
\end{lemma}

Throughout, $|\cdot|$ is the operator norm of an $N\times N$ matrix, $\mathbf{1}$ is the unit
matrix, and distances are measured in units of the spacing of $\mathbb{T}^{(n)}$. For a bond $b$ we
write $b_-$ for its initial site.

\begin{proof}
We first record two elementary facts. If $U_1,\dots,U_r$ and $U_1',\dots,U_r'$ are unitary, then
\begin{equation*}
U_1\cdots U_r-U_1'\cdots U_r'=\sum_{j=1}^r U_1\cdots U_{j-1}(U_j-U_j')U_{j+1}'\cdots U_r' .
\end{equation*}
Hence $|U_1\cdots U_r-U_1'\cdots U_r'|\le\sum_j|U_j-U_j'|$. Second, for real $\theta$ with
$|\theta|\le\pi$ we have $|e^{i\theta}-1|=2|\sin(\theta/2)|\ge 2|\theta|/\pi$. By the spectral
theorem it follows that $|A|\le(\pi/2)|e^{iA}-\mathbf{1}|$ for every Hermitian $A$ with spectrum in
$[-\pi,\pi]$. Likewise $|e^{itA}-e^{isA}|\le|t-s|\,|A|$ for Hermitian $A$ and real $s,t$.

\emph{The comb gauge.} Let $c$ be the centre of $Q$. We pass to the comb gauge on $Q$ centred at
$c$, and write $u$ for the gauge transformation on the sites of $Q$ that achieves it. The comb gauge
centred at $c$ is the gauge in which $V^u(b)=\mathbf{1}$ on every bond $b$ of the spanning tree of
$Q$ formed by the coordinate-ordered lattice paths from $c$ (first along direction $1$, then along
direction $2$, and so on up to direction $d$); such a $u$ exists and is unique up to its value
$u(c)$. For a bond $b$ we write $|b-c|$ for the sup-distance from $c$ to the farther endpoint of $b$.
In this gauge
each bond variable $V^u(b)$, $b$ a bond of $Q$, is a product of at most
$d(1+|b-c|)\le dR$ plaquette variables $V^u(\partial p)$, $p\subset Q$, each conjugated by a unitary
matrix.

A gauge transformation conjugates $V(\partial p)$ by the value of $u$ at the base site of $p$, and
the operator norm is unitarily invariant. Hence each of these conjugated factors lies within
$\varepsilon'$ of $\mathbf{1}$. The first fact above then gives
$|V^u(b)-\mathbf{1}|\le dR\varepsilon'$.

Let $B(b)$ be the Hermitian matrix with spectrum in $(-\pi,\pi]$ such that $V^u(b)=e^{iB(b)}$. By
the second fact,
\begin{equation}\label{8.5:eq-regular-extension-2}
|B(b)|\le\bar\beta:=\frac{\pi}{2}\,dR\varepsilon' .
\end{equation}
By \eqref{8.5:eq-regular-extension-a} we have $dR\varepsilon'\le(dR)^{-1}$, so
$\bar\beta\le\pi/(2dR)$. This is at most $\tfrac14$, since $d=4$ and $R\ge2\mathsf{m}\ge4$ give
$dR\ge2\pi$.

\emph{The cut-off.} Take $\chi$ from the sites of $\mathbb{T}^{(n)}$ to $[0,1]$ with three
properties: $\chi=1$ on $Q'$; $\chi=0$ at sites within distance $1$ of $\partial Q$ and at sites
outside $Q$; and $|\chi(x)-\chi(x')|\le 2/\mathsf{m}$ for nearest neighbours $x,x'$. For instance,
for $x\in Q$ put
\begin{equation*}
\chi(x):=\min\Bigl\{1,\max\Bigl\{0,\frac{\operatorname{dist}(x,\partial Q)-1}{\mathsf{m}-1}\Bigr\}\Bigr\},
\end{equation*}
and put $\chi(x):=0$ outside $Q$. Since $\mathsf{m}\ge2$, its increments between neighbours are at
most $1/(\mathsf{m}-1)\le 2/\mathsf{m}$.

Define $V^{\natural}(b):=e^{i\chi(b_-)B(b)}$ for bonds $b$ of $Q$, and $V^{\natural}(b):=\mathbf{1}$
for all other bonds. Each $V^{\natural}(b)$ is unitary. It equals
$\mathbf{1}$ off $Q$. For a bond $b$ of $Q'$ we have $\chi(b_-)=1$, hence
$V^{\natural}(b)=V^u(b)$. Thus on $Q'$ the configuration $V^{\natural}$ coincides with the gauge
transform $V^u$ of $V$.

\emph{The plaquette bound.} First let $p\not\subset Q$. Then some site $z$ of $p$ lies outside $Q$,
with a coordinate $z_\mu$ outside the range of $Q$. Every site of $p$ differs from $z$ by at most
one in each coordinate. Hence every site of $p$ that lies in $Q$ has its $\mu$-th coordinate at an
extreme value of that range, and so lies on $\partial Q$. Therefore $\chi$ vanishes at all sites of
$p$, every bond variable of $V^{\natural}$ along $\partial p$ equals $\mathbf{1}$, and
$V^{\natural}(\partial p)=\mathbf{1}$.

Now let $p\subset Q$. Write $p=(b_1,b_2,b_3^{-1},b_4^{-1})$ with base site $x$, and put
$\chi_p:=\chi(x)$ and $B_i:=B(b_i)$. For $s\in[0,1]$ put
\begin{equation*}
W(s):=e^{isB_1}e^{isB_2}e^{-isB_3}e^{-isB_4},\qquad
\Sigma_p:=B_1+B_2-B_3-B_4 .
\end{equation*}
Then $V^{\natural}(\partial p)$ is the same product with $s$ replaced by $\chi(b_{i,-})$ in the
$i$-th factor.

(i) By the first fact and the Lipschitz bound for $e^{itA}$,
$|V^{\natural}(\partial p)-W(\chi_p)|\le\sum_i|\chi(b_{i,-})-\chi_p|\,|B_i|$. Each initial site
$b_{i,-}$ is joined to $x$ by at most two nearest-neighbour steps, so
$|\chi(b_{i,-})-\chi_p|\le 4/\mathsf{m}$. With \eqref{8.5:eq-regular-extension-2} this gives
\begin{equation}\label{8.5:eq-regular-extension-3}
|V^{\natural}(\partial p)-W(\chi_p)|\le 4\cdot\frac{4}{\mathsf{m}}\cdot\bar\beta
=\frac{8\pi dR}{\mathsf{m}}\,\varepsilon' .
\end{equation}

(ii) Let $A_1,\dots,A_4$ be matrices with $|A_i|\le a\le\tfrac14$. Expand
$e^{iA_1}\cdots e^{iA_4}$ and $e^{i(A_1+\dots+A_4)}$ in powers of the $A_i$.

The terms of order zero and one agree. The terms of order two differ by
$-\tfrac12\sum_{i<j}[A_i,A_j]$. This follows from the identity
\begin{equation*}
\tfrac12\Bigl(\sum_iA_i\Bigr)^2=\tfrac12\sum_iA_i^2+\tfrac12\sum_{i<j}(A_iA_j+A_jA_i) .
\end{equation*}
In each of the two expansions, the terms of order $m$ have total norm at most $(4a)^m/m!$. Hence
\begin{equation*}
\bigl|e^{iA_1}\cdots e^{iA_4}-e^{i(A_1+\dots+A_4)}\bigr|
\le\tfrac12\sum_{i<j}|[A_i,A_j]|+2\sum_{m\ge3}\frac{(4a)^m}{m!}
\le 6a^2+15a^2=21a^2 .
\end{equation*}
Here $|[A_i,A_j]|\le 2a^2$ over six pairs, and $(4a)^m\le(4a)^2=16a^2$ for $m\ge3$ since
$4a\le1$. Moreover
$\sum_{m\ge3}1/m!=e-\tfrac52<\tfrac{15}{32}$.

We apply (ii) with $A_i=\pm B_i$ and $a=\bar\beta$. The product $e^{iB_1}e^{iB_2}e^{-iB_3}e^{-iB_4}$
is $V^u(\partial p)$, and $|V^u(\partial p)-\mathbf{1}|=|V(\partial p)-\mathbf{1}|\le\varepsilon'$.
Hence $|e^{i\Sigma_p}-\mathbf{1}|\le\varepsilon'+21\bar\beta^2$.

Since $\Sigma_p$ is Hermitian with $|\Sigma_p|\le4\bar\beta\le1$, the second fact gives
$|\Sigma_p|\le(\pi/2)|e^{i\Sigma_p}-\mathbf{1}|$.

For $s\in[0,1]$, apply (ii) with $A_i=\pm sB_i$ and $a=s\bar\beta\le\bar\beta$, and use
$|e^{is\Sigma_p}-\mathbf{1}|\le s|\Sigma_p|$. Together with $\pi/2\le2$,
$21(1+\pi/2)\le54$, $54\cdot\pi^2/4\le135$ and \eqref{8.5:eq-regular-extension-a}, we obtain
\begin{equation}\label{8.5:eq-regular-extension-4}
\begin{split}
|W(s)-\mathbf{1}|&\le s|\Sigma_p|+21\bar\beta^2\le 2\varepsilon'+54\bar\beta^2\\
&\le(2+135\,d^2R^2\varepsilon')\,\varepsilon'\le 137\,\varepsilon' .
\end{split}
\end{equation}
Taking $s=\chi_p$ in \eqref{8.5:eq-regular-extension-4} and adding
\eqref{8.5:eq-regular-extension-3} gives $|V^{\natural}(\partial p)-\mathbf{1}|\le C_{\mathrm{ext}}\varepsilon'$.
This proves \eqref{8.5:eq-regular-extension-1}.

\emph{Determination by $V$ on $Q$.} The function $\chi$ is fixed, and the comb gauge is fixed up to
the value $u(c)$. Changing $u(c)$ by a constant unitary matrix $h$ changes $u$ by the same constant
$h$ at every site, and hence conjugates every $V^u(b)$ by $h$. The logarithm with spectrum in
$(-\pi,\pi]$ commutes
with unitary conjugation, so $B(b)$ is replaced by $hB(b)h^{-1}$ and $V^{\natural}$ by
$hV^{\natural}h^{-1}$. Thus $V^{\natural}$ is determined by the restriction of $V$ to the bonds of
$Q$ up to one constant rotation.
\end{proof}

Since a constant rotation is a gauge transformation, every gauge-invariant function of
$V^{\natural}$ is a gauge-invariant function of the restriction of $V$ to the bonds of $Q$.

\begin{definition}[The history-free remainder on a box]\label{8.5:def-boxed-remainder}
Work on the unit lattice $\mathbb{T}^{(k_0)}$, $k_0=K-s$, with the test functions $f_i$ and
their supports $S_i=\operatorname{supp} f_i$ as in the two-point witness theorem. For each $i$
let $B(x_i,\rho)\supset S_i$ be the ball of centre $x_i$ and radius $\rho$ fixed for the $i$-th
test function in the two-point witness theorem, with $\rho$ as there.

\emph{The boxes.} Let $Q_i$ and $Q_i'$ be the cubes centred at $x_i$ of half-sides
$\tfrac54\rho$ and $\tfrac98\rho$ respectively. In the notation of the regular extension lemma
(Lemma~\ref{8.5:lem-regular-extension}), applied with $Q:=Q_i$ and $Q':=Q_i'$, the side of $Q$
and the distance from $Q'$ to $\partial Q$ are
\begin{equation}\label{8.5:eq-boxed-remainder-1}
R=\tfrac52\rho,\qquad \mathsf{m}=\tfrac18\rho,\qquad \frac{R}{\mathsf{m}}=20,
\qquad C_{\mathrm{ext}}=137+160\pi d .
\end{equation}

\emph{The window.} A configuration $V$ is \emph{in the window on $Q_i$} if $V$ is real and
\begin{equation}\label{8.5:eq-boxed-remainder-2}
\lvert V(\partial p)-1\rvert\le \varepsilon^{\mathrm{win}}:=C_rL^2\varepsilon_{k_0}
\qquad\text{for every plaquette } p\subset Q_i .
\end{equation}
Here $C_rL^2$ is the constant so denoted in the sourced format of the correlation theorem; the
window has the form $O(L^2)\varepsilon_{k_0}$ of the condition
$\lvert V_j(\partial p)-1\rvert<O(L^2)\varepsilon_j$ stated in
\cite[(2.10) and the lines following it]{B-CMP119}.

\emph{The remainder.} For $V$ in the window on $Q_i$, let $V^{\natural}$ be the real configuration
on all bonds of $\mathbb{T}^{(k_0)}$ given by Lemma~\ref{8.5:lem-regular-extension} for the data
$V$ restricted to the bonds of $Q_i$, with $Q:=Q_i$, $Q':=Q_i'$ and
$\varepsilon':=\varepsilon^{\mathrm{win}}$; and let $U_{k_0}(V^{\natural})$ be the configuration
of minimal action at level $k_0$ with block data $V^{\natural}$, as in
Lemma~\ref{8.5:lem-minimiser-regular}. Put
\begin{equation}\label{8.5:eq-boxed-remainder-a}
\begin{aligned}
\operatorname{Irr}_i(V)&:=\text{the right side of \eqref{8.5:eq-remainder-level-sum-a}
with $f:=f_i$,}\\
&\qquad\text{evaluated at $U_{k_0}:=U_{k_0}(V^{\natural})$,}\\
&\qquad\text{if $V$ is in the window on $Q_i$;}\\
\operatorname{Irr}_i(V)&:=0\qquad\text{otherwise.}
\end{aligned}
\end{equation}
The functional $\operatorname{Irr}_i$ is history-free: it is a function of $V$ alone and involves
no large-field region and no history of the large-field operations. It is gauge invariant, it
depends on $V$ only through the restriction of $V$ to the bonds of $Q_i$, and it vanishes at
$V=\mathbf{1}$. As to the dependence on the restriction: the window condition
\eqref{8.5:eq-boxed-remainder-2} reads only plaquettes in $Q_i$, and, by the last assertion of
Lemma~\ref{8.5:lem-regular-extension}, $V^{\natural}$ is determined by the restriction of $V$ to the
bonds of $Q_i$ up to one constant rotation, so that every gauge-invariant function of
$V^{\natural}$ is a gauge-invariant function of that restriction.

\emph{The regularity size.} Put
\begin{equation}\label{8.5:eq-boxed-remainder-3}
\mathfrak{e}:=C_{\mathfrak{e}}\,\varepsilon_{k_0},\qquad
C_{\mathfrak{e}}:=C_{\mathrm{reg}}\,C_{\mathrm{ext}}\,C_rL^2 .
\end{equation}
This is the value $C_{\mathrm{reg}}\varepsilon'$ of \eqref{8.5:eq-minimiser-regular-a} at
$\varepsilon'=C_{\mathrm{ext}}\varepsilon^{\mathrm{win}}=C_{\mathrm{ext}}C_rL^2\varepsilon_{k_0}$,
the bound that Lemma~\ref{8.5:lem-regular-extension} gives for
$\sup_p\lvert V^{\natural}(\partial p)-1\rvert$.

\emph{The bad event.} In the proof of the two-point witness theorem the bad event attached to
the $i$-th ball is read with its region enlarged by $Q_i$. This does not affect the bound used
there on the number of level-$k$ blocks meeting that region: that bound counts the blocks of a
cube of side $8\rho$ about $x_i$, and $Q_i$, a cube of side $\tfrac52\rho$ about $x_i$, lies
inside it. A box, rather than a neighbourhood of
$S_i$ with holes, is needed in \eqref{8.5:eq-boxed-remainder-a} because a field of small
curvature on a region with holes need not have a small potential.
\end{definition}

\begin{lemma}[The counting bound near a support]\label{8.5:lem-counting-bound}
Let $C_M$ and $D_X$ be as in Lemma~\ref{8.5:lem-single-terms}, and let the units of length be such
that the level-$j$ lattice $\mathbb{T}^{(j)}$ has spacing $t\le1$. Let $\pi_j$ be the partition of
$\mathbb{T}^{(j)}$ into cubes of side $Mt$, each containing $M^4$ points of
$\mathbb{T}^{(j)}$. Let $\mathcal{D}_j$ be the family of domains at level $j$, which are unions of
cubes of $\pi_j$. Let $|X|_{\pi_j}$ be the number of cubes of $\pi_j$ contained in $X$, and let
$d_j(X)$ be as fixed in Notation~\ref{0.2:not-glossary}. Let $\delta$ be a fixed number; it enters
only through the reduced decay rate $(1-\frac14\delta)\kappa$. Put
\begin{equation}\label{8.5:eq-counting-bound-1}
\mathsf{S}_{\kappa}:=\sup_{c}\sum_{X\in\mathcal{D}_j,\;X\ni c}
|X|_{\pi_j}\,(d_j(X)+5)^2\,e^{-(1-\frac14\delta)\kappa d_j(X)}
\end{equation}
(cf.~\cite[(1.26)]{B-CMP116}), the supremum being over the cubes $c$ of $\pi_j$, and
\begin{equation}\label{8.5:eq-counting-bound-2}
C_D:=32\cdot 5^d\,C_M^2M^2\,\mathsf{S}_{\kappa}.
\end{equation}
Both depend on $(d,L,M,\kappa,\delta)$ only. Let $S\subset\mathbb{T}^{(j)}$, let $C_h\ge1$ be a
constant, and let $S^{+}$ be the set of points within sup-norm distance $C_hM$ of $S$, with
$|S^{+}|$ its volume measured
in the units in which a cube of $\pi_j$ has volume $(Mt)^4$. Let $\varrho_j\ge0$ (in the applications,
$\varrho_j$ is the ratio of
Lemma~\ref{8.5:lem-single-terms}) and $\mathfrak{n}_T\ge0$. Let $\mathfrak{b}$ be a real background.
Let $T(X,\cdot,y)$, $X\in\mathcal{D}_j$, $y\in X\cap\mathbb{T}^{(j)}$, be a family (or an unpointed
family $T(X,\cdot)$) such that
\begin{equation}\label{8.5:eq-counting-bound-4}
|T^{\mathrm{vac}}(X,\mathfrak{b},y)|\le 2\mathfrak{n}_T\,e^{-(1-\frac14\delta)\kappa d_j(X)}
\bigl(C_MD_Xt^2\varrho_j\bigr)^2
\end{equation}
and such that $T(X,\cdot,y)$ vanishes unless $X^{\oplus}\cap S\neq\emptyset$; here
$T^{\mathrm{vac}}(X,\mathfrak{b},y):=T(X,\mathfrak{b},y)-T(X,1,y)$ is the vacuum-subtracted family.
Then
\begin{equation}\label{8.5:eq-counting-bound-a}
\sum_{X,y}|T^{\mathrm{vac}}(X,\mathfrak{b},y)|\le C_D\,|S^{+}|\,\mathfrak{n}_T\,\varrho_j^2 .
\end{equation}
Let now $n\ge j$ be a scale, let $\mathsf{N}_n$ be the side of $\mathbb{T}^{(n)}$, and call $X$
wrapping if its own box does not fit in the torus in the sense of
Lemma~\ref{8.5:lem-single-terms}, that is, if $D_X>\frac12\mathsf{N}_nt^{-1}$. Put
\begin{equation*}
d_{\mathrm{wr}}:=\frac{\mathsf{N}_nt^{-1}}{8M}-5,\qquad
\mathsf{S}'_{\kappa}:=\sup_{c}\sum_{X\in\mathcal{D}_j,\;X\ni c}
|X|_{\pi_j}\,e^{-\frac12(1-\frac14\delta)\kappa d_j(X)},
\end{equation*}
and assume that $\mathsf{S}'_{\kappa}<\infty$ and
\begin{equation}\label{8.5:eq-counting-bound-3}
\mathsf{S}'_{\kappa}\,e^{-\frac12(1-\frac14\delta)\kappa d_{\mathrm{wr}}}
\le 16\,C_M^2M^2t^4\varrho_j^2\,\mathsf{S}_{\kappa}.
\end{equation}
Condition \eqref{8.5:eq-counting-bound-3} is a hypothesis of the last two assertions below only; it
is to be checked wherever they are applied, for the scale $n$ and the ratio $\varrho_j$ at hand.
Suppose that \eqref{8.5:eq-counting-bound-4} holds only for the $X$ that are not wrapping, and that
for the wrapping $X$ only the trivial bound
\begin{equation*}
|T^{\mathrm{vac}}(X,\mathfrak{b},y)|\le2\mathfrak{n}_T\,e^{-(1-\frac14\delta)\kappa d_j(X)}
\end{equation*}
holds. Then, under \eqref{8.5:eq-counting-bound-3}, the sum of these trivial bounds over the
wrapping $X$ with $X^{\oplus}\cap S\neq\emptyset$ and over $y$ is at most
$C_D|S^{+}|\mathfrak{n}_T\varrho_j^2$. Moreover, \eqref{8.5:eq-counting-bound-a} holds with $C_D$
read doubled, that is, with $2C_D$ in place of $C_D$.
\end{lemma}

\begin{proof}
Only domains $X$ with $X^{\oplus}\cap S\neq\emptyset$ contribute.

First, every cube of $\pi_j$ has side $Mt\le M\le C_hM$, so each of its points lies within sup-norm
distance $C_hM$ of any other of its points. Hence a cube of $\pi_j$ meeting $S$ lies in $S^{+}$. The
cubes of
$\pi_j$ are disjoint, and each has volume $(Mt)^4$. Hence at most $|S^{+}|(Mt)^{-4}$ of them meet $S$.

Second, $X^{\oplus}$ is $X$ enlarged by two layers of cubes of $\pi_j$. Hence, if
$X^{\oplus}\cap S\neq\emptyset$, there is a
cube $c$ of $\pi_j$ meeting $S$ and a cube $c'\subset X$ of $\pi_j$ lying within two layers of cubes
of $c$. For a given $c$ there are at most $5^d$ such cubes $c'$, since the offset of $c'$ from $c$ is
one of the five integers $-2,\dots,2$ in each of the $d$ coordinate directions. Hence, for every
non-negative function $a$ on domains,
\begin{equation*}
\sum_{X:\,X^{\oplus}\cap S\neq\emptyset}a(X)\le\sum_{c\cap S\neq\emptyset}\;
\sum_{c'\text{ within two layers of }c}\;\sum_{X\ni c'}a(X).
\end{equation*}

Third, the number of points $y$ is $\#(X\cap\mathbb{T}^{(j)})=M^4|X|_{\pi_j}$.

Fourth, $D_X\le 4M(d_j(X)+5)$ by Lemma~\ref{8.5:lem-single-terms}. Therefore
\begin{equation*}
\bigl(C_MD_Xt^2\varrho_j\bigr)^2\le\bigl(4C_MM(d_j(X)+5)\bigr)^2t^4\varrho_j^2 .
\end{equation*}
Together with the third point and the factor $2$ of the hypothesis, this gives
\begin{equation*}
\sum_{y}|T^{\mathrm{vac}}(X,\mathfrak{b},y)|
\le 32\,C_M^2M^6t^4\,\mathfrak{n}_T\varrho_j^2\,|X|_{\pi_j}(d_j(X)+5)^2
e^{-(1-\frac14\delta)\kappa d_j(X)} .
\end{equation*}
In the unpointed case the sum over $y$ is a single term. Since $M^4|X|_{\pi_j}\ge1$, the same
bound holds.

We sum over $X\ni c'$ and use \eqref{8.5:eq-counting-bound-1}. This yields at most
$32\,C_M^2M^6t^4\,\mathfrak{n}_T\varrho_j^2\,\mathsf{S}_{\kappa}$ for each $c'$. We then sum over the
at most $5^d$ cubes $c'$ and the at most $|S^{+}|(Mt)^{-4}$ cubes $c$. This gives
\begin{equation*}
\sum_{X,y}|T^{\mathrm{vac}}(X,\mathfrak{b},y)|
\le |S^{+}|(Mt)^{-4}\cdot5^d\cdot32\,C_M^2M^6t^4\,\mathfrak{n}_T\varrho_j^2\,\mathsf{S}_{\kappa}
= C_D\,|S^{+}|\,\mathfrak{n}_T\varrho_j^2 ,
\end{equation*}
by \eqref{8.5:eq-counting-bound-2}. This is \eqref{8.5:eq-counting-bound-a}.

Consider finally the wrapping domains. Condition \eqref{8.5:eq-counting-bound-3} plays, for the sums
weighted by $|X|_{\pi_j}$ at the reduced rate $(1-\frac14\delta)\kappa$, the role that
\eqref{8.5:eq-renormalised-groups-b} plays for the unweighted sums at rate $\kappa$ in the proof of
Lemma~\ref{8.5:lem-renormalised-groups}. A wrapping $X$ has $D_X>\frac12\mathsf{N}_nt^{-1}$, and since
$D_X\le4M(d_j(X)+5)$ by the fourth point, $4M(d_j(X)+5)>\frac12\mathsf{N}_nt^{-1}$, that is,
$d_j(X)>d_{\mathrm{wr}}$. Write $\kappa_\delta:=(1-\frac14\delta)\kappa$. For a wrapping $X$ we then
have
\begin{equation*}
e^{-\kappa_\delta d_j(X)}
=e^{-\frac12\kappa_\delta d_j(X)}\,e^{-\frac12\kappa_\delta d_j(X)}
\le e^{-\frac12\kappa_\delta d_{\mathrm{wr}}}\,e^{-\frac12\kappa_\delta d_j(X)} .
\end{equation*}
Each summand of a wrapping $X$ is at most $2\mathfrak{n}_T e^{-\kappa_\delta d_j(X)}$, and by the
third point (with a single term in the unpointed case, and $M^4|X|_{\pi_j}\ge1$) the sum over $y$ of
these bounds is at most $2\mathfrak{n}_TM^4|X|_{\pi_j}e^{-\kappa_\delta d_j(X)}$. Hence, for every
cube $c'$ of $\pi_j$, by the definition of $\mathsf{S}'_{\kappa}$ and by
\eqref{8.5:eq-counting-bound-3},
\begin{align*}
\sum_{X\ni c',\;X\text{ wrapping}}2\mathfrak{n}_TM^4|X|_{\pi_j}e^{-\kappa_\delta d_j(X)}
&\le 2\mathfrak{n}_TM^4\,e^{-\frac12\kappa_\delta d_{\mathrm{wr}}}
\sum_{X\ni c'}|X|_{\pi_j}e^{-\frac12\kappa_\delta d_j(X)}\\
&\le 2\mathfrak{n}_TM^4\,\mathsf{S}'_{\kappa}\,e^{-\frac12\kappa_\delta d_{\mathrm{wr}}}
\le 32\,C_M^2M^6t^4\,\mathfrak{n}_T\varrho_j^2\,\mathsf{S}_{\kappa}.
\end{align*}
Arranging the sum over the wrapping $X$ with $X^{\oplus}\cap S\neq\emptyset$ as in the second point
and summing over the at most $5^d$ cubes $c'$ and the at most $|S^{+}|(Mt)^{-4}$ cubes $c$, we obtain
that the sum of the trivial bounds over these $X$ and over $y$ is at most
\begin{equation*}
|S^{+}|(Mt)^{-4}\cdot5^d\cdot32\,C_M^2M^6t^4\,\mathfrak{n}_T\varrho_j^2\,\mathsf{S}_{\kappa}
=C_D\,|S^{+}|\,\mathfrak{n}_T\varrho_j^2 ,
\end{equation*}
by \eqref{8.5:eq-counting-bound-2}, which is the right-hand side of
\eqref{8.5:eq-counting-bound-a}. The $X$ that are not wrapping satisfy
\eqref{8.5:eq-counting-bound-4}, and
the argument above, applied to them alone, bounds their sum by $C_D|S^{+}|\mathfrak{n}_T\varrho_j^2$.
Adding the two sums gives \eqref{8.5:eq-counting-bound-a} with $2C_D$ in place of $C_D$.
\end{proof}

\begin{theorem}[The sup-norm bound on the irrelevant remainder]\label{8.5:thm-remainder-sup}
Assume the following statements, which are the hypotheses of this theorem.
\begin{itemize}
\item[(a)] The correlation theorem, in the setting of the two-point witness theorem: it is read at
$m:=s$, with top level $K'=k_0:=K-s$, the unit lattice being $\mathbb{T}^{(k_0)}$. This includes
its bounds at every level $k\le k_0$, its running-shift statement and its homogeneity lemma.
\item[(b)] The marked-expansion lemma and its corollary.
\item[(c)] The proposition of the two-point witness theorem that defines the irrelevant remainder of
the carried source, and the regularity lemma of that theorem.
\end{itemize}
Let $f$ be Lipschitz with $\operatorname{supp} f=S\subset B(x,\rho)$, where $\rho=L^{s}\mathsf{d}$
and $B(x,\rho)$ is the ball of the setting of the two-point witness theorem. Put
\[
\|f\|_{\sharp}:=\max\{\|f\|_{\infty},\,40M\operatorname{Lip}_{k_0}f\},
\qquad \bar y:=x_{k_0}-4c_2 .
\]
Let $Q$ be the cube of half-side $\tfrac54\rho$ about $x$, and let the window on $Q$ and the size
$\varepsilon^{\mathrm{win}}$ be as in Definition~\ref{8.5:def-boxed-remainder}. Let
$\operatorname{Irr}[f](V)$ be given by \eqref{8.5:eq-boxed-remainder-a}. Let
$\mathfrak{e}:=C_{\mathrm{reg}}C_{\mathrm{ext}}\varepsilon^{\mathrm{win}}$ be the regularity size
\eqref{8.5:eq-minimiser-regular-a} of data of size $C_{\mathrm{ext}}\varepsilon^{\mathrm{win}}$;
thus $\mathfrak{e}=C_{\mathfrak{e}}\varepsilon_{k_0}$, with $\varepsilon_{k_0}$ and the constant
$C_{\mathfrak{e}}$ of Definition~\ref{8.5:def-boxed-remainder}.

Assume \eqref{8.5:eq-local-source-split-a}, \eqref{8.5:eq-local-source-split-b},
\eqref{8.5:eq-block-indicators-a}, \eqref{8.5:eq-block-indicators-b},
\eqref{8.5:eq-renormalised-groups-b} and \eqref{8.5:eq-regular-extension-a}, the last at
$(R,\varepsilon')=(\tfrac52\rho,\varepsilon^{\mathrm{win}})$. Assume moreover
\begin{equation}\label{8.5:eq-remainder-sup-a}
\begin{gathered}
C_{\mathrm{ext}}\varepsilon^{\mathrm{win}}\le a_1;\qquad
\mathfrak{e}\le\mathfrak{a}_{k_0}\le\mathfrak{a}_{\circ};\qquad
x_{k_0}\ge 8c_2+3 .
\end{gathered}
\end{equation}
Here $a_1$ is the threshold of Lemma~\ref{8.5:lem-minimiser-regular}. In the middle condition,
$\mathfrak{e}=C_{\mathfrak{e}}\varepsilon_{k_0}$, $\mathfrak{a}_{k_0}=C_{\mathfrak{a}}\alpha_{*,k_0}$
and $\mathfrak{a}_{\circ}=C_{\mathfrak{a}}\alpha_*$, where $\alpha_{*,k_0}$ is the minimal field
radius at level $k_0$, $\alpha_*$ the minimal field radius at the absolute radii, and $C_*$ the
constant in $\alpha_{*,k_0}=C_*g_{k_0}(\log x_{k_0})^{q}$. After division by $g_{k_0}$ the condition
$\mathfrak{e}\le\mathfrak{a}_{k_0}$ therefore reads
\[
C_{\mathfrak{e}}A_0(\log x_{k_0})^{\mathfrak{p}_0}\le C_{\mathfrak{a}}C_{*}(\log x_{k_0})^{q},
\]
which has degree $\mathfrak{p}_0-q<0$ in $\log x_{k_0}$, and the condition
$\mathfrak{a}_{k_0}\le\mathfrak{a}_{\circ}$ reads $\alpha_{*,k_0}\le\alpha_*$. The last
condition gives $\bar y\ge x_{k_0}/2\ge 1$ and $\log x_{k_0}\ge1$.

Then for every $V$ in the window on $Q$,
\begin{equation}\label{8.5:eq-remainder-sup-1}
|\operatorname{Irr}[f](V)|\le\bigl(\mathsf{I}^P+\mathsf{I}^Q+\mathsf{I}^R\bigr)\,|S^{+}|\,\|f\|_{\sharp},
\end{equation}
where $S^{+}:=S^{+C_hM}$ is the set $S$ enlarged by $C_hM$, as in the setting of the two-point
witness theorem and in the counting bound near a support (Lemma~\ref{8.5:lem-counting-bound}), and
\begin{equation}\label{8.5:eq-remainder-sup-b}
\begin{aligned}
\mathsf{I}^P&:=C_P\,\hat\theta^{\circ}_{k_0-1}\,\mathfrak{e}^2,\qquad
C_P:=\frac{C_{LR}\bigl(C_I'S_I^{\circ}+\tfrac{3}{80M}C_D\bigr)}{\mathfrak{a}_{\circ}^2},\\
\mathsf{I}^Q&:=\frac{1}{\lambda(\kappa_0-6)r}
\Bigl[12C_DC_{\mathfrak{m}}\bar y^{-(\kappa_0-4)/2}
+8C_{\beta}C_tC_{\mathfrak{m}}\mathfrak{e}^2\bar y^{-(\kappa_0-6)/2}\Bigr],\\
\mathsf{I}^R&:=r^{-1}C_D\,\mathfrak{m}_{k_0}.
\end{aligned}
\end{equation}
Here $\kappa_0$ and $C_{\mathfrak{m}}$ are the exponent and the constant that enter
\eqref{8.5:eq-local-source-split-a}.
The letters $\lambda$, $C_{\beta}$, $C_t$ and $K_R$ are constants of the correlation theorem:
$\lambda$ is the lower bound on the increments $x_{j+1}-x_j$ of the inverse running couplings,
$C_{\beta}$ is the flow constant of Definition~\ref{1.2:def-flow-constants}, $C_t$ is the
constant with this name in the correlation theorem, and $K_R$ is the constant in the ratio
$\theta_k=K_R\delta_k/\alpha$ of the fluctuation cut to the field radius, whose running maximum at
the absolute radii is $\hat\theta^{\circ}_k$.
In particular, suppose $\kappa_0\ge7$, $\delta_k=g_kA_1(\log x_k)^{p_0}$ and
$\varepsilon_k=g_kA_0(\log x_k)^{\mathfrak{p}_0}$, where $\delta_k$ is the real fluctuation cut and
$p_0$ denotes here the exponent of the logarithm in it. Then
\begin{equation}\label{8.5:eq-remainder-sup-c}
\begin{split}
|\operatorname{Irr}[f](V)|&\le C_{\star}\,x_{k_0}^{-3/2}(\log x_{k_0})^{p_0+2\mathfrak{p}_0}\,
|S^{+}|\,\|f\|_{\sharp},\\
C_{\star}&:=\frac{2C_PK_RA_1A_0^2C_{\mathfrak{e}}^2}{\alpha_*}
+\frac{34C_DC_{\mathfrak{m}}+12C_{\beta}C_tC_{\mathfrak{m}}A_0^2C_{\mathfrak{e}}^2}{\lambda r}
+\frac{6C_DC_{\mathfrak{m}}}{r}.
\end{split}
\end{equation}
Since $x_{k_0}=g_{k_0}^{-2}$, the bound \eqref{8.5:eq-remainder-sup-c} is one power of $g_{k_0}$,
times a polylogarithm, below the size
$g_{k_0}^2A_0^2(\log x_{k_0})^{2\mathfrak{p}_0}=\varepsilon_{k_0}^2$ of the main insertion.
\end{theorem}

\begin{proof}[Proof of Theorem~\ref{8.5:thm-remainder-sup}: the decomposition]
Fix $V$ in the window on $Q$. Let $V^{\natural}$ be the extension of $V$ entering
\eqref{8.5:eq-boxed-remainder-a}, and put $U:=U_{k_0}(V^{\natural})$.

\emph{Regularity of $U$.} Since $V$ is in the window on $Q$, it is real and satisfies
$|V(\partial p)-1|\le\varepsilon^{\mathrm{win}}$ for $p\subset Q$. The side of $Q$ is
$R=\tfrac52\rho$, and \eqref{8.5:eq-regular-extension-a} holds at
$(R,\varepsilon')=(\tfrac52\rho,\varepsilon^{\mathrm{win}})$. The regular extension lemma
(Lemma~\ref{8.5:lem-regular-extension}) therefore applies. It shows that $V^{\natural}$ is real on all
bonds and satisfies
\[
\sup_p|V^{\natural}(\partial p)-1|\le C_{\mathrm{ext}}\varepsilon^{\mathrm{win}} .
\]
By the first condition of \eqref{8.5:eq-remainder-sup-a}, $C_{\mathrm{ext}}\varepsilon^{\mathrm{win}}\le a_1$.
Lemma~\ref{8.5:lem-minimiser-regular}, on minimisers of small-curvature data, applies with
$\varepsilon'=C_{\mathrm{ext}}\varepsilon^{\mathrm{win}}$. Hence $U$ is real, regular on the whole torus,
and $\mathfrak{e}$-regular at scale $k_0$.

For every $j\le k_0$, $U$ is a real level-$j$ minimal configuration. Indeed, by part (v) of
Proposition~\ref{8.5:prop-local-source-split} (splitting the sourced terms on shrinking radii), the
class of real level-$j$ minimal configurations considered there contains $U_n(W)$ for every
$n\ge j$ and every real $W$ on $\mathbb{T}^{(n)}$ whose minimiser $U_n(W)$ is real and regular. By
the preceding paragraph, $V^{\natural}$ is real with
$\sup_p|V^{\natural}(\partial p)-1|\le C_{\mathrm{ext}}\varepsilon^{\mathrm{win}}\le a_1$, and
$U_{k_0}(V^{\natural})$ is real and regular; we take $n=k_0$ and $W=V^{\natural}$.

Next, $(U,J(U))\restriction X\in\mathcal{U}^{\flat}_j(X)\subset\mathcal{U}_j(X)$ for all $X\in D_j$.
Here $D_j$ is the family of level-$j$ localisation domains, and $J(U)$ is the second member of the
background pair $(\mathbf{U},\mathbf{J})$ that the correlation theorem attaches to $U$.
This rests on three inputs. First, \cite[p.~263]{B-CMP109} states that ``the minimal configurations
$U_j$ satisfying the bound $|\partial U_j-1|<\varepsilon_0\xi^2$ with $\varepsilon_0$ sufficiently
small, satisfy the above conditions'', the conditions there being those that define these background
spaces; the letters $\varepsilon_0$ and $\xi$ are those of \cite{B-CMP109}. Second, the homogeneity
lemma of the correlation theorem applies. Third, $\mathfrak{e}\le\mathfrak{a}_{k_0}$ holds by
\eqref{8.5:eq-remainder-sup-a}.
In particular, every family below that is defined on $\mathcal{U}^{\flat}_j(X)$ or on
$\mathcal{U}_j(X)$ may be evaluated at $(U,J(U))\restriction X$.

\emph{Splitting the $E$-totals.} Let the totals $E^{(j)}(U;w)$ and their derivatives
$\dot{E}^{(j)}[f]$ be those of Lemma~\ref{8.5:lem-remainder-level-sum}. For $X\in D_j$ and
$y\in X$ we put
\[
\dot{Q}^{(j)}(X,\cdot,y)[f]:=\partial_zQ^{(j)}(X,\cdot,y;zf)\big|_{z=0}.
\]
Part (i) of Proposition~\ref{8.5:prop-local-source-split} gives, on
$\mathcal{U}^{\flat}_j(X)\times\mathcal{W}_j(X)$,
\[
E^{(j)}(X,\cdot,y;w)=P^{(j)}(X,\cdot,y;w)+Q^{(j)}(X,\cdot,y;w).
\]
We sum over $X$ and $y$, put
$w=zf$ and differentiate at $z=0$. This gives
\begin{equation}\label{8.5:eq-remainder-sup-2}
\dot{E}^{(j)}[f]=\partial_z\Big|_{z=0}\sum_{X,y}P^{(j)}(X,\cdot,y;zf)
+\sum_{X,y}\dot{Q}^{(j)}(X,\cdot,y)[f].
\end{equation}

\emph{The $P$-total.} By part (v) of Proposition~\ref{8.5:prop-local-source-split}, applied to the real
level-$j$ minimal configuration $U$ and the Lipschitz function $g=f$, the first total on the right of
\eqref{8.5:eq-remainder-sup-2}, evaluated at $U$, equals $\dot{\mathcal{T}}^{(j)}[f](U)$ plus a
constant. By part (d) of the linear response on absolute radii
(Theorem~\ref{8.5:thm-linear-response}), on the same class of configurations,
\[
\sum_{X,y}\dot{P}^{(j)}(X,U,J(U),y)[f]=\dot{\mathcal{T}}^{(j)}[f](U).
\]
Hence the first total equals $\sum_{X,y}\dot{P}^{(j)}(X,U,J(U),y)[f]$ plus a constant independent
of $U$. That constant disappears under the vacuum subtraction
$F^{\mathrm{vac}}=F-F(1)$.

\emph{The flow coefficient.} Put $\mathsf{b}^Q_j:=\partial_{\zeta}\beta[Q^{(j)}_{\zeta}]\big|_{\zeta=0}$.
The running-shift statement of the correlation theorem gives
\[
b_j'(0)=\mathsf{b}^P_j+\mathsf{b}^Q_j .
\]

\emph{The response family.} By part (b) of Theorem~\ref{8.5:thm-linear-response}, the Lipschitz
remainder is defined by
\[
\mathsf{D}^{(j)}(X,\cdot,y)[g]:=\dot{P}^{(j)}(X,\cdot,y)[g]-g(p_y)F^{(j)}(X,\cdot,y).
\]
Therefore
\[
\sum_{X,y}\dot{P}^{(j)}(X,\cdot,y)[f]
=\sum_{y}f(p_y)\sum_{X\ni y}F^{(j)}(X,\cdot,y)+\sum_{X,y}\mathsf{D}^{(j)}(X,\cdot,y)[f].
\]
By part (c) of the same theorem, $F^{(j)}$ is a pointed $E$-family at scale $j$ on the spaces
$\mathcal{U}_j(X)$ and
$\beta[F^{(j)}]=\mathsf{b}^P_j$. Its renormalised pointed group is therefore
\[
\mathfrak{i}_y[F^{(j)}](U)
=\sum_{X\ni y}\bigl(F^{(j)}(X,\cdot,y)\bigr)^{\mathrm{vac}}(U)-\mathsf{b}^P_jA(h^{(j)}_y,U),
\]
the group treated in Lemma~\ref{8.5:lem-renormalised-groups}. Recall that
$I_jf=\sum_{y\in\mathbb{T}^{(j)}}f(p_y)h^{(j)}_y$ and that $A(\cdot,U)$ is linear in its first
argument. Hence
\[
\mathsf{b}^P_jA(I_jf,U)=\sum_yf(p_y)\,\mathsf{b}^P_jA(h^{(j)}_y,U).
\]

\emph{The decomposition.} Insert these identities into \eqref{8.5:eq-remainder-level-sum-a}, taken
at $U_{k_0}=U$ as prescribed by \eqref{8.5:eq-boxed-remainder-a}.

In the $j$-th summand, $(\dot{E}^{(j)}[f])^{\mathrm{vac}}(U)$ splits, by \eqref{8.5:eq-remainder-sup-2}
and the identification of the $P$-total, into three pieces:
\begin{itemize}
\item the vacuum subtraction of $\sum_yf(p_y)\sum_{X\ni y}F^{(j)}$;
\item the vacuum subtraction of $\sum_{X,y}\mathsf{D}^{(j)}[f]$;
\item the vacuum subtraction of $\sum_{X,y}\dot{Q}^{(j)}[f]$.
\end{itemize}
The subtracted term $b_j'(0)A(I_jf,U)$ splits into $\mathsf{b}^P_jA(I_jf,U)$ and
$\mathsf{b}^Q_jA(I_jf,U)$.

The first piece minus $\mathsf{b}^P_jA(I_jf,U)$ equals, by the expression for
$\mathsf{b}^P_jA(I_jf,U)$ above and the formula for $\mathfrak{i}_y[F^{(j)}](U)$,
$\sum_yf(p_y)\mathfrak{i}_y[F^{(j)}](U)$. The second piece is the $j$-th term of the second sum in
$\operatorname{Irr}^P$ below. The third piece minus $\mathsf{b}^Q_jA(I_jf,U)$ is the $j$-th term of
$\operatorname{Irr}^Q$ below.

This proves
\begin{equation}\label{8.5:eq-remainder-sup-3}
\operatorname{Irr}[f](V)=\operatorname{Irr}^P+\operatorname{Irr}^Q+\operatorname{Irr}^R,
\end{equation}
where
\begin{equation}\label{8.5:eq-remainder-sup-4}
\begin{aligned}
\operatorname{Irr}^P&:=\sum_{j\le k_0}\Bigl\{\sum_yf(p_y)\,\mathfrak{i}_y[F^{(j)}](U)
+\sum_{X,y}\bigl(\mathsf{D}^{(j)}(X,\cdot,y)[f]\bigr)^{\mathrm{vac}}(U)\Bigr\},\\
\operatorname{Irr}^Q&:=\sum_{j\le k_0}\Bigl\{\sum_{X,y}\bigl(\dot{Q}^{(j)}(X,\cdot,y)[f]\bigr)^{\mathrm{vac}}(U)
-\mathsf{b}^Q_jA(I_jf,U)\Bigr\},\\
\operatorname{Irr}^R&:=\sum_{j\le k_0}\sum_X\bigl(\dot{R}^{(j)}(X,\cdot)[f]\bigr)^{\mathrm{vac}}(U).
\end{aligned}
\end{equation}
The bounds \eqref{8.5:eq-remainder-sup-1}--\eqref{8.5:eq-remainder-sup-c} are deduced from
\eqref{8.5:eq-remainder-sup-3} by estimating the three parts $\operatorname{Irr}^P$,
$\operatorname{Irr}^Q$ and $\operatorname{Irr}^R$ of \eqref{8.5:eq-remainder-sup-4} in turn; this is
done in the estimates of the three parts of the irrelevant remainder, which complete the proof of
Theorem~\ref{8.5:thm-remainder-sup}.
\end{proof}

\begin{proof}[Proof of Theorem~\ref{8.5:thm-remainder-sup}, continued]\label{8.5:prf-remainder-sup-proof}
It remains to bound the three parts $\operatorname{Irr}^P$, $\operatorname{Irr}^Q$ and $\operatorname{Irr}^R$
of the decomposition of $\operatorname{Irr}[f](V)$ obtained above at the real background
$U=U_{k_0}(V^{\natural})$, and then to derive the final display. Throughout, $j$ runs over the
levels $j\le k_0$. We recall that $U$ is $\mathfrak{e}$-regular at scale $k_0\ge j$, and that
$(\cdot)^{\mathrm{vac}}(U)$ denotes the value at $U$ minus the value at the trivial background.

\smallskip
\noindent\textit{The part $\operatorname{Irr}^P$.} Here the radii are the absolute ones (the system
$\circ$), and we put $\varrho:=\mathfrak{e}/\mathfrak{a}_{\circ}$. The hypothesis
$\mathfrak{e}\le\mathfrak{a}_{k_0}\le\mathfrak{a}_{\circ}$ of
Theorem~\ref{8.5:thm-remainder-sup} gives $\varrho\le1$.

We first treat the frozen part $\sum_{j}\sum_y f(p_y)\,\mathfrak{i}_y[F^{(j)}](U)$. By the linear
response on absolute radii (Theorem~\ref{8.5:thm-linear-response}, parts (a) and (b)), the family
$F^{(j)}$ is a pointed family at level $j$, analytic on the absolute radii, with norm at most
$\mathfrak{n}_j$. Lemma~\ref{8.5:lem-renormalised-groups} applies at $\circ$ with
$\varrho_j=\varrho\le1$ and background $U$. Since $1+\varrho\le2$, it gives
\begin{equation}\label{8.5:eq-remainder-sup-proof-1}
\bigl|\mathfrak{i}_y[F^{(j)}](U)\bigr|\le 2C_I'\,\mathfrak{n}_j\,\nu^3\varrho^2\,t^{4+\hat\alpha}.
\end{equation}
Here $\nu:=k_0-j+1$ and $t:=L^{-(\nu-1)}$ are the level difference and the scale ratio of
Lemma~\ref{8.5:lem-renormalised-groups} with $n=k_0$, so that
$t^{\hat\alpha}=L^{-\hat\alpha(\nu-1)}$.

The value $f(p_y)$ is non-zero for at most $|S^+|\,t^{-4}$ points $y$ of $\mathbb{T}^{(j)}$, and
$|f(p_y)|\le\|f\|_\infty$. Hence the contribution of level $j$ is at most
$2C_I'\|f\|_\infty|S^+|\,\mathfrak{n}_j\nu^3\varrho^2t^{\hat\alpha}$. The totals established in
the proof of Theorem~\ref{8.5:thm-linear-response} give, for $j\le k_0$,
$\mathfrak{n}_j\le C_{LR}\hat\theta^{\circ}_{k_0-1}$. Summing over the levels, that is over
$\nu$, and using $2\sum_{\nu\ge1}\nu^3L^{-\hat\alpha(\nu-1)}=S_I^{\circ}$, we obtain
\begin{equation}\label{8.5:eq-remainder-sup-proof-2}
\Bigl|\sum_{j\le k_0}\sum_y f(p_y)\,\mathfrak{i}_y[F^{(j)}](U)\Bigr|
\le\|f\|_\infty\,|S^+|\cdot C_I'S_I^{\circ}\cdot C_{LR}\,\hat\theta^{\circ}_{k_0-1}\,\varrho^2.
\end{equation}

We next treat the Lipschitz part
$\sum_j\sum_{X,y}\bigl(\mathsf{D}^{(j)}(X,\cdot,y)[f]\bigr)^{\mathrm{vac}}(U)$. By
Theorem~\ref{8.5:thm-linear-response}\,(b), $|\mathsf{D}^{(j)}(X,\cdot,y)[f]|$ is bounded by
$\mathfrak{d}_j\operatorname{Lip}_j(f;X^{\oplus})$ times the exponential factor in $d_j(X)$
displayed in the corresponding display of Theorem~\ref{8.5:thm-linear-response}. Moreover $\mathsf{D}^{(j)}(X,\cdot,y)[f]=0$ if
$f\equiv0$ on $X^{\oplus}$, so the family vanishes unless $X^{\oplus}\cap S\neq\emptyset$.

Lemma~\ref{8.5:lem-single-terms} (the bound \eqref{8.5:eq-single-terms-a}) converts this into the
input of the counting bound, Lemma~\ref{8.5:lem-counting-bound}. For every $X$ whose box fits in
the torus it gives the factor $2\min\{1,(C_MD_Xt^2\varrho)^2\}\le2(C_MD_Xt^2\varrho)^2$, and the
remaining $X$ are covered by the counting bound itself. The counting bound
\eqref{8.5:eq-counting-bound-a}, with
$\mathfrak{n}_T=\mathfrak{d}_j\operatorname{Lip}_j(f)$, then gives
\begin{equation*}
\sum_{X,y}\bigl|\bigl(\mathsf{D}^{(j)}(X,\cdot,y)[f]\bigr)^{\mathrm{vac}}(U)\bigr|
\le C_D\,|S^+|\,\mathfrak{d}_j\operatorname{Lip}_j(f)\,\varrho^2 .
\end{equation*}
Rescaling the Lipschitz modulus from level $k_0$ to level $j$ and using the definition of
$\|f\|_\sharp$, we have
\begin{equation*}
\operatorname{Lip}_jf=L^{-(\nu-1)}\operatorname{Lip}_{k_0}f\le L^{-(\nu-1)}\,\frac{\|f\|_\sharp}{40M}.
\end{equation*}
Since $L$ is an odd integer greater than one, $\sum_{\nu\ge1}L^{-(\nu-1)}=L/(L-1)\le3/2$. With
$\mathfrak{d}_j\le C_{LR}\hat\theta^{\circ}_{k_0-1}$ for $j\le k_0$, summation over the levels gives
\begin{equation}\label{8.5:eq-remainder-sup-proof-3}
\sum_{j\le k_0}\sum_{X,y}\bigl|\bigl(\mathsf{D}^{(j)}(X,\cdot,y)[f]\bigr)^{\mathrm{vac}}(U)\bigr|
\le\|f\|_\sharp\,|S^+|\,\frac{3}{80M}\,C_D\,C_{LR}\,\hat\theta^{\circ}_{k_0-1}\,\varrho^2.
\end{equation}
Adding \eqref{8.5:eq-remainder-sup-proof-2} and \eqref{8.5:eq-remainder-sup-proof-3}, and using
$\|f\|_\infty\le\|f\|_\sharp$ and $\varrho^2=\mathfrak{e}^2/\mathfrak{a}_{\circ}^2$, we obtain
\begin{equation*}
|\operatorname{Irr}^P|\le C_P\,\hat\theta^{\circ}_{k_0-1}\mathfrak{e}^2\,|S^+|\,\|f\|_\sharp
=\mathsf{I}^P|S^+|\,\|f\|_\sharp .
\end{equation*}

\smallskip
\noindent\textit{The part $\operatorname{Irr}^Q$.} If $|z|\,\|f\|_\sharp<r$, then
$zf\in\mathcal{W}_j(X)$ by the definition of the admissible sources at level $j$ (as recalled in
Corollary~\ref{8.5:cor-source-derivative-bis}). Hence
$z\mapsto Q^{(j)}(X,\cdot,y;zf)$ is analytic on the disc $|z|<r/\|f\|_\sharp$. On that disc
it is bounded by $\mathfrak{m}_{j-1}e^{-\kappa d_j(X)}$, the bound of the Q-family in the
correlation theorem taken at $m:=s$. The Cauchy estimate for the derivative at $z=0$ gives
\begin{equation}\label{8.5:eq-remainder-sup-proof-4}
\bigl|\dot{Q}^{(j)}(X,\cdot,y)[f]\bigr|\le\mathfrak{m}_{j-1}\,\|f\|_\sharp\,r^{-1}e^{-\kappa d_j(X)},
\end{equation}
and $\dot{Q}^{(j)}(X,\cdot,y)[f]=0$ unless $X^{+}\cap S\neq\emptyset$.

We apply Lemma~\ref{8.5:lem-single-terms} on the shrinking radii $\flat$, with
$\varrho_j=\mathfrak{e}/\mathfrak{a}_j$. Since $\mathfrak{e}\le\mathfrak{a}_{k_0}$, we have
$\varrho_j\le\mathfrak{a}_{k_0}/\mathfrak{a}_j=\rho_{k_0,j}$. The counting bound
\eqref{8.5:eq-counting-bound-a} then shows that level $j$ contributes at most
$C_D|S^+|\rho_{k_0,j}^2\,\mathfrak{m}_{j-1}\|f\|_\sharp/r$ to the first sum in
$\operatorname{Irr}^Q$.

To sum over $j$ we use the reference-flow comparison of the correlation theorem. It gives
$y_{j-1}\ge\bar y+\lambda(k_0-j+1)$ and $x_j\le x_{k_0}+3\lambda(k_0-j)+2c_2$. The hypothesis
$x_{k_0}\ge8c_2+3$ gives $x_{k_0}\ge7c_2$, hence $14c_2\le2x_{k_0}+3\lambda$. Since
$\bar y=x_{k_0}-4c_2$, the inequality $14c_2\le2x_{k_0}+3\lambda$ is equivalent to the second
step of
\begin{equation*}
x_j\le x_{k_0}+3\lambda(k_0-j)+2c_2\le3\bigl(\bar y+\lambda(k_0-j+1)\bigr).
\end{equation*}
The ratio bound on field radii of the correlation theorem gives
$\rho_{k_0,j}^2\le2x_j/x_{k_0}$. The Q-family bound is
$\mathfrak{m}_{j-1}=C_{\mathfrak{m}}y_{j-1}^{-(\kappa_0-2)/2}$. With $\nu=k_0-j+1\ge1$ as before we
therefore have
\begin{equation*}
\rho_{k_0,j}^2\,\mathfrak{m}_{j-1}\le\frac{6C_{\mathfrak{m}}}{x_{k_0}}\,(\bar y+\lambda\nu)
\,(\bar y+\lambda\nu)^{-(\kappa_0-2)/2}
=\frac{6C_{\mathfrak{m}}}{x_{k_0}}\,(\bar y+\lambda\nu)^{-(\kappa_0-4)/2}.
\end{equation*}

The summation bound of the correlation theorem states that, for $a>1$,
\begin{equation}\label{8.5:eq-remainder-sup-proof-5}
\sum_{\nu\ge1}(\bar y+\lambda\nu)^{-a}\le\frac{\bar y^{\,1-a}}{\lambda(a-1)} .
\end{equation}
We use it at $a=(\kappa_0-4)/2$, which exceeds $1$ exactly when $\kappa_0>6$. Together with
$x_{k_0}\ge\bar y$ it gives
\begin{equation}\label{8.5:eq-remainder-sup-proof-6}
\sum_{j\le k_0}\rho_{k_0,j}^2\,\mathfrak{m}_{j-1}
\le\frac{6C_{\mathfrak{m}}}{x_{k_0}}\sum_{\nu\ge1}(\bar y+\lambda\nu)^{-(\kappa_0-4)/2}
\le\frac{12\,C_{\mathfrak{m}}\,\bar y^{-(\kappa_0-4)/2}}{\lambda(\kappa_0-6)} .
\end{equation}

Next come the counterterms $\mathsf{b}^Q_jA(I_jf,U)$ of $\operatorname{Irr}^Q$, as introduced in
the first part of the proof. The bound on the counterterm coefficients of the
first-order recursion gives $|\mathsf{b}^Q_j|\le v_{j-1}/(2r)$. Part (b) of the lemma on the
level-$j$ interpolants gives
$|A(I_jf,U)|\le4\mathfrak{e}^2\|f\|_\infty|S^+|$; here the interpolant $I_jf$ is supported in
$S^{+2}$. The flow errors of the correlation theorem satisfy
\begin{equation*}
\sum_{j\le k_0}v_{j-1}\le2C_{\beta}C_tC_{\mathfrak{m}}\sum_{\nu\ge1}(\bar y+\lambda\nu)^{-(\kappa_0-4)/2}.
\end{equation*}
By \eqref{8.5:eq-remainder-sup-proof-5} at the same $a$ this gives
\begin{equation*}
\sum_{j\le k_0}v_{j-1}\le\frac{4C_{\beta}C_tC_{\mathfrak{m}}\,\bar y^{-(\kappa_0-6)/2}}{\lambda(\kappa_0-6)} .
\end{equation*}
Consequently,
\begin{equation*}
\sum_{j\le k_0}\bigl|\mathsf{b}^Q_jA(I_jf,U)\bigr|
\le\frac{4\mathfrak{e}^2\|f\|_\infty|S^+|}{2r}\sum_{j\le k_0}v_{j-1}
\le\frac{8C_{\beta}C_tC_{\mathfrak{m}}\,\mathfrak{e}^2\,\bar y^{-(\kappa_0-6)/2}}{\lambda(\kappa_0-6)\,r}
\,|S^+|\,\|f\|_\infty .
\end{equation*}
Together with \eqref{8.5:eq-remainder-sup-proof-6}, multiplied by $C_D|S^+|\|f\|_\sharp/r$, and
with $\|f\|_\infty\le\|f\|_\sharp$, this gives
$|\operatorname{Irr}^Q|\le\mathsf{I}^Q|S^+|\,\|f\|_\sharp$.

\smallskip
\noindent\textit{The part $\operatorname{Irr}^R$.} The derivative $\dot{R}^{(j)}(X,\cdot)[f]$ is
the derivative at $z=0$ of the level-$j$ remainder term at the source $zf$. On
$|z|<r/\|f\|_\sharp$ this term is bounded by $g_j^{\kappa_0}e^{-\kappa d_j(X)}$, by the bound on
the remainder family in the correlation theorem. The Cauchy estimate gives
\begin{equation*}
\bigl|\dot{R}^{(j)}(X,\cdot)[f]\bigr|\le g_j^{\kappa_0}\,\|f\|_\sharp\,r^{-1}e^{-\kappa d_j(X)}.
\end{equation*}
Lemma~\ref{8.5:lem-single-terms} on the shrinking radii and the counting bound
\eqref{8.5:eq-counting-bound-a} then bound the contribution of level $j$ by
$C_D|S^+|\rho_{k_0,j}^2g_j^{\kappa_0}\|f\|_\sharp/r$. The weights $w_{k_0,j}$ of the memory sum
$\mu_{k_0}$ satisfy $w_{k_0,j}\ge\rho_{k_0,j}^2$. With the summation bounds of the correlation
theorem this gives
$\sum_{j}\rho_{k_0,j}^2g_j^{\kappa_0}\le\mu_{k_0}\le\mathfrak{m}_{k_0}$. Hence
\begin{equation*}
|\operatorname{Irr}^R|\le r^{-1}C_D\mathfrak{m}_{k_0}|S^+|\,\|f\|_\sharp
=\mathsf{I}^R|S^+|\,\|f\|_\sharp .
\end{equation*}

The three bounds together are \eqref{8.5:eq-remainder-sup-b}.

\smallskip
\noindent\textit{The final display.} Write $x:=x_{k_0}$ and assume $\kappa_0\ge7$. By
\eqref{8.5:eq-block-indicators-a}, $\hat\theta^{\circ}_{k_0-1}\le2K_R\delta_{k_0}/\alpha_*$, and
\begin{equation*}
\delta_{k_0}=g_{k_0}A_1(\log x)^{p_0}=A_1x^{-1/2}(\log x)^{p_0}.
\end{equation*}
By the choice of $\mathfrak{e}$ in the setting of Theorem~\ref{8.5:thm-remainder-sup} and
$\varepsilon_{k_0}=g_{k_0}A_0(\log x)^{\mathfrak{p}_0}$, we have
$\mathfrak{e}^2\le C_{\mathfrak{e}}^2A_0^2x^{-1}(\log x)^{2\mathfrak{p}_0}$. Hence
\begin{equation*}
\mathsf{I}^P\le\frac{2C_PK_RA_1A_0^2C_{\mathfrak{e}}^2}{\alpha_*}\,
x^{-3/2}(\log x)^{p_0+2\mathfrak{p}_0}.
\end{equation*}

By hypothesis, $\bar y\ge x/2\ge1$. Since $\kappa_0\ge7$, we have $(\kappa_0-4)/2\ge3/2$,
$(\kappa_0-6)/2\ge1/2$ and $\kappa_0-6\ge1$. This gives
$\bar y^{-(\kappa_0-4)/2}\le\bar y^{-3/2}\le2^{3/2}x^{-3/2}$ and
\begin{equation*}
\mathfrak{e}^2\,\bar y^{-(\kappa_0-6)/2}\le\mathfrak{e}^2\,\bar y^{-1/2}
\le\sqrt2\,C_{\mathfrak{e}}^2A_0^2(\log x)^{2\mathfrak{p}_0}x^{-3/2}.
\end{equation*}
With $12\cdot2^{3/2}\le34$ and $8\sqrt2\le12$ these give
\begin{equation*}
\mathsf{I}^Q\le\frac{34C_DC_{\mathfrak{m}}+12C_{\beta}C_tC_{\mathfrak{m}}A_0^2C_{\mathfrak{e}}^2
(\log x)^{2\mathfrak{p}_0}}{\lambda r}\,x^{-3/2}.
\end{equation*}

The reference-flow comparison gives $y_{k_0}\ge\bar y$, and $(\kappa_0-2)/2\ge5/2$. Hence
\begin{equation*}
\mathfrak{m}_{k_0}=C_{\mathfrak{m}}y_{k_0}^{-(\kappa_0-2)/2}\le C_{\mathfrak{m}}\bar y^{-5/2}
\le6C_{\mathfrak{m}}x^{-5/2},
\end{equation*}
and therefore $\mathsf{I}^R\le6C_DC_{\mathfrak{m}}r^{-1}x^{-3/2}$, because $x\ge1$.

Finally, $\log x\ge1$ and the exponents $p_0$ and $\mathfrak{p}_0$ are non-negative. So
$1\le(\log x)^{2\mathfrak{p}_0}\le(\log x)^{p_0+2\mathfrak{p}_0}$. Inserting the three bounds into
\eqref{8.5:eq-remainder-sup-b} gives \eqref{8.5:eq-remainder-sup-c}, with the constant $C_{\star}$
displayed there.
\end{proof}

\begin{theorem}[The sup-norm bound on shrinking radii]\label{8.5:thm-remainder-sup-shrinking}
Work in the setting of Theorem~\ref{8.5:thm-remainder-sup}: the correlation
theorem at $m:=s$ with $K'=k_0$, in the setting of the two-point witness theorem, together with
the two further antecedents named in the statement of Theorem~\ref{8.5:thm-remainder-sup}.
Let $f$ be
Lipschitz with
$\operatorname{supp}f=S\subset B(x,\rho)$, and let $V$
be in the window on $Q$, the cube of half-side $(5/4)\rho$ about $x$, in the sense of
Definition~\ref{8.5:def-boxed-remainder}. Of the
further conditions of Theorem~\ref{8.5:thm-remainder-sup}, assume only
\eqref{8.5:eq-local-source-split-b},
\eqref{8.5:eq-renormalised-groups-b}, \eqref{8.5:eq-regular-extension-a} at
$(R,\varepsilon')=((5/2)\rho,\varepsilon^{\mathrm{win}})$, and
\eqref{8.5:eq-remainder-sup-a}; assume in addition the ceilings
\eqref{8.5:eq-source-quadratic-a-bis} of Lemma~\ref{8.5:lem-source-quadratic-bis}. Then
\begin{equation}\label{8.5:eq-remainder-sup-shrinking-a}
\begin{gathered}
\bigl|\operatorname{Irr}[f](V)\bigr|
\le\bigl[C_{\star}^{\flat}\,\hat\theta_{k_0-1}\,\mathfrak{e}^2+\mathsf{I}^R\bigr]\,
|S^{+}|\,\|f\|_{\sharp},\\
C_{\star}^{\flat}:=C_s\bigl(C_I'S'_I+C_BC_DS'_{\omega}\bigr),
\end{gathered}
\end{equation}
with $\mathsf{I}^R$ as in \eqref{8.5:eq-remainder-sup-b}, $S^{+}$ as in
Theorem~\ref{8.5:thm-remainder-sup}, $C_B$ as in Corollary~\ref{8.5:cor-source-derivative-bis},
$C_s$ as in Lemma~\ref{8.5:lem-source-quadratic-bis}, $C_I'$ as in
Lemma~\ref{8.5:lem-renormalised-groups}, $C_D$ as in Lemma~\ref{8.5:lem-counting-bound}, and
$S'_I$, $S'_{\omega}$ the level sums entering \eqref{8.5:eq-source-quadratic-a-bis}.
\end{theorem}

The first term in the bracket has the natural size $\varepsilon_{k_0}^2$ times
$C_{\theta}(\log x_{k_0})^{a_{\delta}-q}$; here $a_{\delta}$ denotes the exponent of $\log x_k$ in
$\delta_k=g_kA_1(\log x_k)^{a_{\delta}}$.

The proof below bounds the source-dependent terms directly on the shrinking radii through
Corollary~\ref{8.5:cor-source-derivative-bis}. It does not use the splitting into frozen-response
families and the families $Q^{(j)}_{\zeta}$ of the three-family split on which the first two
parts of the proof of
Theorem~\ref{8.5:thm-remainder-sup} rest, nor the absolute-radius response bounds entering
those parts.

\begin{proof}
Put $U:=U_{k_0}(V^{\natural})$. As at the start of the proof of
Theorem~\ref{8.5:thm-remainder-sup}, \eqref{8.5:eq-regular-extension-a} and
\eqref{8.5:eq-remainder-sup-a} make $U$ real and $\mathfrak{e}$-regular at scale $k_0$, with
$\mathfrak{e}\le\mathfrak{a}_{k_0}$.

\emph{Grouping.} By Definition~\ref{8.5:def-boxed-remainder}, $\operatorname{Irr}[f](V)$ is the
right side of \eqref{8.5:eq-remainder-level-sum-a} at $U_{k_0}:=U$. We write that right side in
the grouping of the correlation theorem's format, with the families $\dot S^{(j)}$ and
$\dot D^{(j)}[f]$ of Corollary~\ref{8.5:cor-source-derivative-bis}; by that corollary
$\beta[\dot S^{(j)}]=b_j'(0)$, the coefficient of the counterterm in
\eqref{8.5:eq-remainder-level-sum-a}. This gives
\begin{equation}\label{8.5:eq-remainder-sup-shrinking-1}
\operatorname{Irr}[f](V)=\sum_{j\le k_0}\Bigl\{\sum_y f(p_y)\,
\mathfrak{i}_y[\dot S^{(j)}](U)+\sum_{X,y}\bigl(\dot D^{(j)}(X,\cdot,y)[f]\bigr)^{\mathrm{vac}}(U)
\Bigr\}+\operatorname{Irr}^R,
\end{equation}
where $\operatorname{Irr}^R$ is the second sum of \eqref{8.5:eq-remainder-level-sum-a}:
\begin{equation*}
\operatorname{Irr}^R:=\sum_{j\le k_0}\sum_X\bigl(\dot R^{(j)}(X,\cdot)[f]\bigr)^{\mathrm{vac}}(U).
\end{equation*}

\emph{Source bounds.} By Corollary~\ref{8.5:cor-source-derivative-bis}, which applies under
\eqref{8.5:eq-source-quadratic-a-bis} and \eqref{8.5:eq-local-source-split-b} (with $K'=k_0$),
$\dot S^{(j)}$ is a pointed $E$-family with
$\|\dot S^{(j)}\|^{\flat}\le C_s\hat\theta_{j-1}\mathfrak{a}_j^2$. Moreover
\begin{equation*}
|\dot D^{(j)}(X,\cdot,y)[f]|
\le C_BC_s\hat\theta_{j-1}\mathfrak{a}_j^2\,\omega_{k_0,j}\,\|f\|_{\sharp}\,
e^{-(1-\frac14\delta)\kappa d_j(X)},
\end{equation*}
where $\omega_{k_0,j}$ and the fraction $\delta$ in the exponent are as in
Corollary~\ref{8.5:cor-source-derivative-bis}.

\emph{Lemmas at $\flat$.} We apply the renormalised-group bound
\eqref{8.5:eq-renormalised-groups-a} of Lemma~\ref{8.5:lem-renormalised-groups} and the
single-term bound \eqref{8.5:eq-single-terms-a} of Lemma~\ref{8.5:lem-single-terms}, both at
$\flat$. They are applied with background $U$, which is $\mathfrak{e}$-regular at scale
$n=k_0\ge j$, and with $\varrho_j=\mathfrak{e}/\mathfrak{a}_j$. Here and below $\nu$,
$t=L^{-(\nu-1)}$ (with $\nu-1=k_0-j$) and $\hat\alpha$ are the letters of
Lemma~\ref{8.5:lem-renormalised-groups} at $n=k_0$; this $t$ is not the $t$ of the glossary
(Notation~\ref{0.2:not-glossary}). Since
$\mathfrak{e}\le\mathfrak{a}_{k_0}$, we have $\varrho_j\le\rho_{k_0,j}$, which is the hypothesis
of Lemma~\ref{8.5:lem-renormalised-groups} at $\flat$. In both applications the norm contains
the factor $C_s\hat\theta_{j-1}\mathfrak{a}_j^2$ and the bound the factor $\varrho_j^2$, and
\begin{equation*}
C_s\hat\theta_{j-1}\mathfrak{a}_j^2\cdot\varrho_j^2
=C_s\hat\theta_{j-1}\mathfrak{e}^2,
\end{equation*}
so that the radius $\mathfrak{a}_j$ cancels.

\emph{First sum.} By Lemma~\ref{8.5:lem-renormalised-groups},
\begin{equation*}
|\mathfrak{i}_y[\dot S^{(j)}](U)|
\le C_I'\,C_s\hat\theta_{j-1}\mathfrak{e}^2\,\nu^3(1+\varrho_j)\,t^{4+\hat\alpha}.
\end{equation*}
As in the proof of Theorem~\ref{8.5:thm-remainder-sup}, $f(p_y)\ne0$ for at most $|S^{+}|t^{-4}$
points of $\mathbb{T}^{(j)}$. Since $t^{\hat\alpha}=L^{-\hat\alpha(\nu-1)}$ and
$\varrho_j\le\rho_{k_0,j}$, the level sum satisfies
\begin{equation*}
\sum_{\nu}\nu^3(1+\varrho_j)L^{-\hat\alpha(\nu-1)}\le S'_I.
\end{equation*}
Using $|f(p_y)|\le\|f\|_{\infty}\le\|f\|_{\sharp}$, the first sum of
\eqref{8.5:eq-remainder-sup-shrinking-1} is therefore at most
\begin{equation*}
C_sC_I'S'_I\,\max_{j\le k_0}\hat\theta_{j-1}\,\mathfrak{e}^2\,|S^{+}|\,\|f\|_{\sharp}.
\end{equation*}

\emph{Second sum.} Put
$\mathfrak{n}_T:=C_BC_s\hat\theta_{j-1}\mathfrak{a}_j^2\omega_{k_0,j}\|f\|_{\sharp}$.
Lemma~\ref{8.5:lem-single-terms}, applied to $F:=\dot D^{(j)}(X,\cdot,y)[f]$ with
$\mathfrak{n}=\mathfrak{n}_Te^{-(1-\frac14\delta)\kappa d_j(X)}$, gives for
$(\dot D^{(j)}[f])^{\mathrm{vac}}(U)$ the size factor in the hypothesis of the counting bound
\eqref{8.5:eq-counting-bound-a} of Lemma~\ref{8.5:lem-counting-bound}, with
$\varrho=\varrho_j$; for the wrapping $X$ one uses \eqref{8.5:eq-renormalised-groups-b}. The
counting bound, applied to this family at level $j$, bounds level $j$ by
\begin{equation*}
C_D|S^{+}|\,\mathfrak{n}_T\varrho_j^2
=C_BC_DC_s\,\hat\theta_{j-1}\mathfrak{e}^2\,\omega_{k_0,j}\,|S^{+}|\,\|f\|_{\sharp}.
\end{equation*}
The levels are summed with $\sum_{j}\omega_{k_0,j}=S'_{\omega}$.

\emph{Conclusion.} The quantity $\hat\theta_{j-1}$ is a running maximum, so
$\hat\theta_{j-1}\le\hat\theta_{k_0-1}$ for $j\le k_0$. The two sums of
\eqref{8.5:eq-remainder-sup-shrinking-1} are therefore together at most
$C_{\star}^{\flat}\hat\theta_{k_0-1}\mathfrak{e}^2|S^{+}|\|f\|_{\sharp}$. The part
$\operatorname{Irr}^R$ is bounded by $\mathsf{I}^R|S^{+}|\|f\|_{\sharp}$ exactly as the part of the
same name in the proof of Theorem~\ref{8.5:thm-remainder-sup}. That bound uses, besides the
setting, only Cauchy's estimate for $\dot R^{(j)}$ on the source disc,
Lemma~\ref{8.5:lem-single-terms} at $\flat$, the counting bound \eqref{8.5:eq-counting-bound-a}
with \eqref{8.5:eq-renormalised-groups-b} for the wrapping $X$, and the summation of the level
factors; none of these uses a condition of Theorem~\ref{8.5:thm-remainder-sup} that is not
assumed here. Adding the three contributions gives
\eqref{8.5:eq-remainder-sup-shrinking-a}.
\end{proof}

\begin{corollary}[The natural-size bound on the window]\label{8.5:cor-natural-size}
Let $f_1,f_2$ be test functions of the two-point witness theorem, in its class
$\mathfrak{T}(\mathsf{d};\vartheta)$, with supports at separation $\mathsf{d}$. Let $k_0=K-s$, and let
$\tilde f_i$ be the function built from $f_i$ that enters the main insertion of that theorem. Let
\begin{equation*}
G_i:=N\mathcal{O}^{(k_0)}(\tilde f_i)+\operatorname{Irr}_i,\qquad i=1,2,
\end{equation*}
be the functional of the two-point witness theorem: the sum of the main insertion
$N\mathcal{O}^{(k_0)}(\tilde f_i)$ and the irrelevant remainder $\operatorname{Irr}_i$.
Let $S_i=\operatorname{supp} f_i$ and let $S_i^{+}$ be its enlargement used in
Theorem~\ref{8.5:thm-remainder-sup-shrinking}. Write $\rho:=L^{s}\mathsf{d}$ for the separation
measured in cells of the level-$k_0$ lattice; \emph{the window} is the separation window
$C_W(s)\le\rho<LC_W(s)$ of the two-point witness theorem, a condition on the test functions. Let
$\varepsilon=\varepsilon_{k_0}$ be the
small-field size at level $k_0$, and let $r$ be the parameter of the shrinking radii of the
correlation theorem. Assume:
\begin{itemize}
\item[(a)] the conditions of Theorem~\ref{8.5:thm-remainder-sup-shrinking}, so that
\eqref{8.5:eq-remainder-sup-shrinking-a} holds with $f=f_i$ and $S^{+}=S_i^{+}$;
\item[(a$'$)] at the shrinking radii of Theorem~\ref{8.5:thm-remainder-sup-shrinking} (those of
the correlation theorem) and for fixed $r$, the bracket in
\eqref{8.5:eq-remainder-sup-shrinking-a} satisfies
\begin{equation*}
C_{\star}^{\flat}\,\hat\theta_{k_0-1}\,\mathfrak{e}^2+\mathsf{I}^R\le C_{\mathrm{br}}\,\varepsilon^2
\end{equation*}
with a constant $C_{\mathrm{br}}$ depending on $r$;
\item[(b)] on the window, the bound
\begin{equation*}
|N\mathcal{O}^{(k_0)}(\tilde f_i)|\le C_{\mathrm{ins}}\,\mathfrak{n}(f_i)\,\varepsilon^2
\end{equation*}
with a constant $C_{\mathrm{ins}}$, supplied by the regularity lemma of the two-point witness
theorem;
\item[(c)] $\rho\ge 40M/\vartheta$ on the window.
\end{itemize}
Then at the shrinking radii of Theorem~\ref{8.5:thm-remainder-sup-shrinking} (those of the
correlation theorem), for $r$ fixed, with $C_G:=80\,C_{\mathrm{br}}/\vartheta+C_{\mathrm{ins}}$, on
the window
\begin{equation}\label{8.5:eq-natural-size-1}
|G_i|\le C_G\,\mathfrak{n}(f_i)\,\varepsilon^2 ,\qquad i=1,2.
\end{equation}
\end{corollary}

\begin{proof}
We first treat the remainder $\operatorname{Irr}_i$. Apply hypothesis (a) with $f=f_i$. It bounds
$|\operatorname{Irr}_i|$ by
\begin{equation*}
\bigl[C_{\star}^{\flat}\,\hat\theta_{k_0-1}\,\mathfrak{e}^2+\mathsf{I}^R\bigr]\,
|S_i^{+}|\,\|f_i\|_{\sharp}\le C_{\mathrm{br}}\,\varepsilon^2\,|S_i^{+}|\,\|f_i\|_{\sharp},
\end{equation*}
where the inequality is hypothesis (a$'$).

By hypothesis (c), $\rho\ge 40M/\vartheta$ on the window, and for test functions of the class of
the two-point witness theorem $\|f\|_{\sharp}=\|f\|_{\infty}$ once $\rho\ge 40M/\vartheta$. Hence
$\|f_i\|_{\sharp}=\|f_i\|_{\infty}$. By the definitions of $S_i^{+}$ and of $\mathfrak{n}$ in the
setting of the two-point witness theorem we also
have $|S_i^{+}|\le 80\rho^4$ cells and $\vartheta\rho^4\|f_i\|_{\infty}\le\mathfrak{n}(f_i)$.
Combining these gives
\begin{equation}\label{8.5:eq-natural-size-2}
|S_i^{+}|\,\|f_i\|_{\infty}\le 80\,\rho^4\|f_i\|_{\infty}\le \frac{80\,\mathfrak{n}(f_i)}{\vartheta}.
\end{equation}
By \eqref{8.5:eq-natural-size-2},
$|\operatorname{Irr}_i|\le 80\,C_{\mathrm{br}}\,\vartheta^{-1}\,\mathfrak{n}(f_i)\,\varepsilon^2$
on the window.

By hypothesis (b),
$|N\mathcal{O}^{(k_0)}(\tilde f_i)|\le C_{\mathrm{ins}}\,\mathfrak{n}(f_i)\,\varepsilon^2$ on the
window. Since $G_i$ is the sum of the two pieces, adding the two bounds gives
\eqref{8.5:eq-natural-size-1} with $C_G=80\,C_{\mathrm{br}}/\vartheta+C_{\mathrm{ins}}$.
\end{proof}

The bound \eqref{8.5:eq-natural-size-1} is the bound stated in item~(d) of the source lemma.

\begin{lemma}[Comparison with the branch reading of the remainder]
\label{8.5:lem-branch-comparison}
Let $k_0$ be the level and $i\in\{1,2\}$; let $Q_i$, $Q_i'$, the window on $Q_i$ and
$\operatorname{Irr}_i$ be as in Definition~\ref{8.5:def-boxed-remainder}, with $\operatorname{Irr}_i$
given by \eqref{8.5:eq-boxed-remainder-a}. Let $\mathfrak{e}$ be the regularity size fixed in that
definition. The centre $x_i$, the radius $\rho$ and the support $S_i\subset B(x_i,\rho)$ are those of
the same definition, and $f_i$ is the test function supported in $S_i$; $S_i^{+}$ is read as in
\eqref{8.5:eq-counting-bound-a} with $S=S_i$; and $r$, $E_0$, $C_I$, $C_B$, $\mathfrak{m}$ are the
constants of the correlation theorem's sourced format at level $k_0$.
Assume the following four statements.
\begin{enumerate}
\item[(1)] \emph{Domain-uniform localisation of minimisers.} There are constants $c_{\eta}$,
$B_{\eta}$, $C_{\mathrm{ax}}$, depending on $(d,L)$ only, with the following property, uniformly in
the region $\Omega$ and for every size $\mathfrak{e}'>0$. A real configuration on $\Omega^{+}$ that is
regular at size $\mathfrak{e}'$
and critical for a determining set of level $k_0$ on $\Omega$ is, modulo gauge, within
\begin{equation*}
6B_{\eta}C_{\mathrm{ax}}\,\mathfrak{e}'\,\eta^{2}\,e^{-c_{\eta}R_{\Omega}(y)}
\end{equation*}
in every bare plaquette over the cell of $y$ of the $\Omega$-local minimiser of its own level-$k_0$
averages. Here $R_{\Omega}(y)\ge\operatorname{dist}(y,\Omega^{c})-2M_1$. This is the domain-uniform
form of the localisation of minimisers, which goes back to
\cite[Proposition~9, (190)]{B-CMP102b}.
\item[(2)] The bounds \textup{(c)} and \textup{(d)} of the correlation theorem at level $k_0$.
\item[(3)] The backgrounds considered below satisfy the conditions \textup{(c1)}--\textup{(c3)} of
the correlation theorem on backgrounds, in the form recalled in
Definition~\ref{8.5:def-regular-configurations}, with the size $\mathfrak{a}_n$ replaced
by $\mathfrak{e}$.
\item[(4)] The correlation theorem's bound on single terms and its bound on renormalised groups,
whose proofs read a background only through \textup{(c1)}--\textup{(c3)}; this is the first step of
its lemma on complex backgrounds. Also the Schwarz-lemma estimate used in the proof of part
\textup{(iii)} of its bound on single terms.
\end{enumerate}
Let $(h,V)$ be such that no live label of $h$ meets $Q_i$ and $V$ is in the window on $Q_i$. Let
$\mathfrak{b}^h:=U^{(h)}_{k_0}$ be the background of the branch $h$. It depends on $V$. Through the
determining set $\mathfrak{B}_h$ it also depends on the variables $\mathsf{V}$ that $\mathbf{T}_h$
still integrates in live regions, however far these are from $Q_i$. Let
\begin{equation*}
\operatorname{Irr}^h_i:=\partial_{z_i}\bigl[\mathbf{E}^w+\mathbf{D}^w+\mathbf{R}^w\bigr]_h
\big|_{z=0}\ \text{at}\ \mathfrak{b}^h\ \text{minus its value at}\ \mathbf{1},
\end{equation*}
where only the terms with $X\subset\Lambda_j(h)$ are kept. Then, for
\begin{equation*}
\rho\ \ge\ 64(2M_1+2)\vee 96M
\end{equation*}
and for all values of $\mathsf{V}$,
\begin{equation}\label{8.5:eq-branch-comparison-a}
\begin{split}
\bigl|\operatorname{Irr}^h_i(V)-\operatorname{Irr}_i(V)\bigr|\ \le\ \mathfrak{f}^{\mathrm{irr}}_i
:={}& r^{-1}\|f_i\|_{\sharp}\Bigl[C^{\mathrm{irr}}_{\mathrm{loc}}\,|S_i^{+}|\,e^{-c_{\eta}\rho/64}\\
&\qquad{}+C^{\mathrm{irr}}_{\mathrm{far}}\Bigl(\frac{\rho}{M}\Bigr)^{4}L^{4}
e^{-\kappa''\rho/(128M)}\Bigr].
\end{split}
\end{equation}
The constants in \eqref{8.5:eq-branch-comparison-a} are as follows:
\begin{equation*}
C^{\mathrm{irr}}_{\mathrm{loc}}:=8C_{\mathrm{av}}(12B_{\eta}C_{\mathrm{ax}})^{1/2}
C_{\mathrm{top}}^{(1/2)},\qquad
C_{\mathrm{top}}^{(1/2)}:=C_IE_0S_I^{(1/2)}+C_BC_DE_0S_{\omega}+C_D\mathfrak{m}.
\end{equation*}
Here $C_D$ is the counting constant of \eqref{8.5:eq-counting-bound-a}. The sum $S_I^{(1/2)}$ is the
level sum $S_I$ of the expansion of the sourced format at $n=k_0$, read at the H\"older exponent
$\tfrac12\hat\alpha$; like $S_I$, and for the same reason, it does not depend on $K$. Finally,
$C^{\mathrm{irr}}_{\mathrm{far}}$ is the constant fixed in part (i) of the proof.
\end{lemma}

\begin{proof}
The terms of $\operatorname{Irr}^h_i$ and of $\operatorname{Irr}_i$ that read $f_i$ are labelled by a
level $j$ and a polymer $X$. The box of such a term is $X$ with two layers of cubes of $\pi_j$ of
margin. We put $\nu:=1+k_0-j$ and $t:=L^{j-k_0}$, as in
Definition~\ref{8.5:def-regular-configurations} at $n=k_0$. A term is \emph{far} if its box leaves
$S_i^{+\rho/16}$. It is \emph{near} if its box lies in $S_i^{+\rho/16}$.

\emph{(i) Far terms.} A far term either occurs in one of the two sums and not in the other, or it
occurs in both but is read on different fields. We apply the Cauchy estimate in $z$ to the bounds
\textup{(c)}, \textup{(d)} of hypothesis~(2). This bounds each far term by
\begin{equation*}
(E_0\vee 1)\,\|f_i\|_{\sharp}\,r^{-1}e^{-\kappa''d_j(X)}.
\end{equation*}
The term reads $f_i$ and its box leaves $S_i^{+\rho/16}$. Hence
\begin{equation*}
d_j(X)\ \ge\ \frac{\rho L^{\nu-1}}{16M}-8.
\end{equation*}
At level $j$ at most $(4L\rho L^{\nu-1}/M+3)^{4}$ cubes meet $S_i^{+}$. We anchor each far term at
one of these cubes. The polymers anchored at a given cube are summed by the tree bound
\cite[(1.26)]{B-CMP116} at the rate $\tfrac14\kappa''$. This leads to the sum over the levels
\begin{equation*}
\sum_{\nu}\Bigl(\frac{4L\rho L^{\nu-1}}{M}+3\Bigr)^{4}
e^{-\frac14\kappa''\left[\rho L^{\nu-1}/(16M)-8\right]_{+}}
\ \le\ C\Bigl(\frac{\rho}{M}\Bigr)^{4}L^{4}e^{-\kappa''\rho/(128M)}.
\end{equation*}
The resulting bound on the sum of all far terms has the form
\begin{equation*}
r^{-1}\|f_i\|_{\sharp}\Bigl(\frac{\rho}{M}\Bigr)^{4}L^{4}e^{-\kappa''\rho/(128M)}
\end{equation*}
times a constant. By definition, that
constant is $C^{\mathrm{irr}}_{\mathrm{far}}$. This gives the second term of
\eqref{8.5:eq-branch-comparison-a}.

\emph{(ii) Near terms.} Let $\Omega$ be the union of the cubes of side $M_1$ contained in $Q_i'$. Put
$U:=U_{k_0}(V^{\natural})$, with $V^{\natural}$ the regular extension of
Lemma~\ref{8.5:lem-regular-extension}. We compare $\mathfrak{b}^h$ and $U$.
\begin{itemize}
\item Both are real.
\item Both are critical for determining sets of level $k_0$ on $\Omega$. For $\mathfrak{b}^h$ this
holds because no live label of $h$ meets $Q_i$. For $U$ it holds on the whole lattice.
\item Both are regular at a size at most $\mathfrak{e}$.
\item Their level-$k_0$ averages are gauge equivalent to $V$ on $\Omega^{+}$; for $U$ this is
Lemma~\ref{8.5:lem-regular-extension}.
\end{itemize}
We apply hypothesis~(1) twice, with $\mathfrak{e}'=\mathfrak{e}$, against the same $\Omega$-local
minimiser, namely that of the common level-$k_0$ averages. Let $y$ lie in $S_i^{+\rho/16}$. Its depth
satisfies $R_{\Omega}(y)\ge\rho/16-2M_1-2$. Moreover $\rho/16-2M_1-2\ge\rho/32$, because
$\rho\ge64(2M_1+2)$. Each of $\mathfrak{b}^h$ and $U$ is therefore within
$6B_{\eta}C_{\mathrm{ax}}\mathfrak{e}\eta^{2}e^{-c_{\eta}\rho/32}$ of the same minimiser. By the
triangle inequality, modulo gauge, on $S_i^{+\rho/16}$,
\begin{equation}\label{8.5:eq-branch-comparison-1}
|\mathfrak{b}^h(\partial q)-U(\partial q)|\ \le\ \mathfrak{h}\,\eta^{2},\qquad
\mathfrak{h}:=12B_{\eta}C_{\mathrm{ax}}\,\mathfrak{e}\,e^{-c_{\eta}\rho/32},
\end{equation}
for every bare plaquette $q$.

Fix a near term with box $\square_0$. By \textup{(c1)}, on $\square_0$ the two backgrounds are the
box minimisers $U_j(\square_0,\mathsf{V}^a)$ of their own data, $a\in\{1,2\}$. We take $a=1$ for
$\mathfrak{b}^h$ and $a=2$ for $U$. In the comb gauge centred at $y$ we write
$\mathsf{V}^a=e^{iB^a}$. Here $B^a$ is as in Definition~\ref{8.5:def-regular-configurations} at size
$\mathfrak{e}$, with curvature $F^a$. Put $\mathsf{B}:=B^1-B^2$. Then
\begin{equation*}
|\mathsf{B}(b)|\ \le\ C_{\mathrm{av}}\,d\,(1+|b-y|)\,\mathfrak{h}\,t^{2}.
\end{equation*}
Three facts give this bound. First, a bond variable in the comb gauge is a product of at most
$d(1+|b-y|)$ unitarily conjugated plaquette variables. Second, the differences of the factors add.
Third, the constant $C_{\mathrm{av}}=C_{\mathrm{av}}(d,L)$ carries the bound
\eqref{8.5:eq-branch-comparison-1} on fine plaquettes to level-$j$ plaquettes, by
\cite[p.~263]{B-CMP109}. Put $\mathsf{F}:=F^1-F^2$. For the same reasons
$|\mathsf{F}|\le C_{\mathrm{av}}\mathfrak{h}t^{2}$. Moreover $\mathsf{F}$ is a difference of two
curvatures, each obeying the H\"older bound of Definition~\ref{8.5:def-regular-configurations} with its
constant $C$. Hence
\begin{align*}
|\mathsf{F}(x)-\mathsf{F}(x')|
&\le\min\bigl\{2C_{\mathrm{av}}\mathfrak{h}t^{2},\ 2C\mathfrak{e}t^{2+\hat\alpha}|x-x'|^{\hat\alpha}\bigr\}\\
&\le 2C_{\mathrm{av}}'(\mathfrak{h}\mathfrak{e})^{1/2}t^{2+\frac12\hat\alpha}|x-x'|^{\frac12\hat\alpha}.
\end{align*}
The last line holds for a constant $C_{\mathrm{av}}'$, by bounding the minimum by the geometric mean
of its two members.

Consequently, put $B_w:=B^2+w\mathsf{B}$ and
$\mathsf{w}:=\mathfrak{a}_{k_0}/(2C_{\mathrm{av}}(\mathfrak{h}\mathfrak{e})^{1/2})$. For $|w|<\mathsf{w}$,
the data $B_w$ form a holomorphic family of complex data. They obey \textup{(c3)} at the size
\begin{equation*}
\mathfrak{e}+|w|\,C_{\mathrm{av}}(\mathfrak{h}\mathfrak{e})^{1/2}\ \le\ \mathfrak{a}_{k_0},
\end{equation*}
with the H\"older exponent $\tfrac12\hat\alpha$. This follows from the bounds on $\mathsf{B}$ and
$\mathsf{F}$ above, since $\mathfrak{h}\le(\mathfrak{h}\mathfrak{e})^{1/2}$. Condition \textup{(c2)}
holds along the family.

By hypothesis~(4), the proofs of the bounds on single terms and on renormalised groups read a
background only through \textup{(c1)}--\textup{(c3)}. So every renormalised group in the first row,
and every single term in the second and third rows, of the expansion of the $z_i$-derivative of the
sourced format at level $k_0$ (the expansion whose level sum is $S_I$)
is bounded along the whole disc $|w|<\mathsf{w}$ by its bound $\mathfrak{N}$. For the
groups, $\hat\alpha$ is replaced by $\tfrac12\hat\alpha$. By the Schwarz-lemma estimate of
hypothesis~(4), the value of each such term at $w=1$ differs from its value at $w=0$ by at most
$(2/\mathsf{w})\mathfrak{N}$. No $\mathbf{J}$-variable is interpolated: $J_j(\square_0,\cdot)$ is a
function of the data.

Summed over the near terms, these bounds $\mathfrak{N}$ give at most
$C_{\mathrm{top}}^{(1/2)}|S_i^{+}|\,\|f_i\|_{\sharp}/r$. The near terms therefore contribute at most
\begin{equation*}
\frac{2}{\mathsf{w}}\,C_{\mathrm{top}}^{(1/2)}\,|S_i^{+}|\,\frac{\|f_i\|_{\sharp}}{r}.
\end{equation*}
By the definition of
$\mathfrak{h}$,
\begin{equation*}
\frac{2}{\mathsf{w}}=\frac{4C_{\mathrm{av}}(\mathfrak{h}\mathfrak{e})^{1/2}}{\mathfrak{a}_{k_0}}
=4C_{\mathrm{av}}(12B_{\eta}C_{\mathrm{ax}})^{1/2}\,\frac{\mathfrak{e}}{\mathfrak{a}_{k_0}}\,
e^{-c_{\eta}\rho/64}
\le 4C_{\mathrm{av}}(12B_{\eta}C_{\mathrm{ax}})^{1/2}e^{-c_{\eta}\rho/64}.
\end{equation*}
The last inequality amounts to $\mathfrak{e}\le\mathfrak{a}_{k_0}$. So the near terms contribute at
most
\begin{equation*}
4C_{\mathrm{av}}(12B_{\eta}C_{\mathrm{ax}})^{1/2}C_{\mathrm{top}}^{(1/2)}\,
r^{-1}\|f_i\|_{\sharp}\,|S_i^{+}|\,e^{-c_{\eta}\rho/64},
\end{equation*}
which is at most the first term of \eqref{8.5:eq-branch-comparison-a}, by the definition of
$C^{\mathrm{irr}}_{\mathrm{loc}}$.

Adding the contributions of (i) and (ii) gives \eqref{8.5:eq-branch-comparison-a}, for all values of
$\mathsf{V}$.
\end{proof}

The quantity $\mathfrak{f}^{\mathrm{irr}}_i$ is of the same kind as the error $\mathfrak{f}_i$ of the
source lemma of the two-point witness theorem. It is super-polynomially small in $\bar x_s$ under a
lower bound on $C_0$ that is independent of $g$ and of the same kind as the lower bound on $C_0$ in
the source lemma.
The decomposition $\Phi^h_i=c^0_i+G_i+\mathfrak{L}^h_i$ of the two-point witness theorem holds with
the reading \eqref{8.5:eq-boxed-remainder-a}, provided $\mathfrak{F}^h_i$ is replaced by
$\mathfrak{F}^h_i+\mathbb{E}_h[\operatorname{Irr}^h_i-\operatorname{Irr}_i]$. For every $h$ with no
live label on $Q_i$, the supremum of this replacement is at most
$\mathfrak{f}_i+\mathfrak{f}^{\mathrm{irr}}_i$. For the other $h$ the difference is placed in
$\mathfrak{L}^{\mathrm{near},h}_i$, which is bounded by the suprema of the source lemma.

The conditional covariance inside $\Psi^h$ changes by at most
$4\mathfrak{f}^{\mathrm{irr}}_1\mathfrak{f}^{\mathrm{irr}}_2$ plus cross terms having
$\mathfrak{f}^{\mathrm{irr}}$ as a factor. All of these lie among the terms of that theorem's
remainder built from the errors $\mathfrak{f}_i$, under the same lower bound on $C_0$.
This lemma, and with it hypothesis~(1), is not used in the proofs of the sup bound for
$\operatorname{Irr}_i$ and of the extraction estimate of this section.

\begin{lemma}[Covariances bounded by sup norms]\label{8.5:lem-covariance-sup}
In this lemma $\sup|H|$ denotes the supremum of $|H|$ over the whole domain of a function $H$.
\begin{itemize}
\item[(e1)] Let $(\Omega,\mathcal{F},P)$ be a probability space, with expectation $E$. Let $X,Y$ be
bounded random variables on it, real- or complex-valued, and put
$\operatorname{Cov}(X,Y):=E[XY]-E[X]E[Y]$. Then $\operatorname{Cov}(X,Y)=E[X(Y-EY)]$.
\item[(e2)] Let $k_0<K'$ be integers and $\Omega_{k_0},\dots,\Omega_{K'}$ measurable spaces. Let
$\mathbb{E}_j$, $k_0\le j\le K'-1$, be a tower of positive normalised functionals. By this we mean
that for every $w\in\Omega_{j+1}$ there is a probability measure $\nu_{j,w}$ on $\Omega_j$, such
that $w\mapsto\nu_{j,w}(A)$ is measurable for every measurable $A\subset\Omega_j$, and
$\mathbb{E}_j[G](w)=\int G\,d\nu_{j,w}$ for every bounded measurable $G$ on $\Omega_j$. For a
bounded measurable $G$ on $\Omega_{k_0}$ put
$G^{\langle\rangle}:=\mathbb{E}_{K'-1}\circ\cdots\circ\mathbb{E}_{k_0}[G]$. Then for each
$W\in\Omega_{K'}$ there is a probability measure $\mu_W$ on $\Omega_{k_0}$ with
$G^{\langle\rangle}(W)=\int G\,d\mu_W$ for every such $G$. Consequently, for bounded measurable
$F_1,F_2$ on $\Omega_{k_0}$, the quantity
\begin{equation*}
T^{\langle\rangle}[F_1,F_2](W):=(F_1F_2)^{\langle\rangle}(W)-F_1^{\langle\rangle}(W)\,
F_2^{\langle\rangle}(W)
\end{equation*}
is the covariance of $F_1$ and $F_2$ under $\mu_W$.
\end{itemize}
Put $\operatorname{osc}(H):=\sup_{W,W'}|H(W)-H(W')|$. Then, for every $W\in\Omega_{K'}$ and every
bounded measurable $F$ on $\Omega_{k_0}$,
\begin{equation}\label{8.5:eq-covariance-sup-a}
\begin{gathered}
|\operatorname{Cov}(X,Y)|\le 2\sup|X|\cdot\sup|Y|,\\
|T^{\langle\rangle}[F_1,F_2](W)|\le 2\sup|F_1|\cdot\sup|F_2|,\\
\operatorname{osc}(F^{\langle\rangle})\le 2\sup|F|.
\end{gathered}
\end{equation}
\end{lemma}

\begin{proof}
(e1) We have $E[1]=1$. By linearity of $E$, $E[X(Y-EY)]=E[XY]-E[X]\,E[Y]=\operatorname{Cov}(X,Y)$.
Pointwise on $\Omega$,
\begin{equation*}
|Y-EY|\le|Y|+|EY|\le 2\sup|Y|.
\end{equation*}
Hence
\begin{equation*}
|\operatorname{Cov}(X,Y)|\le E\bigl[|X|\,|Y-EY|\bigr]\le 2\sup|X|\cdot\sup|Y|,
\end{equation*}
which is the first line of \eqref{8.5:eq-covariance-sup-a}.

(e2) For $k_0<j\le K'$ let $G_j:=\mathbb{E}_{j-1}\circ\cdots\circ\mathbb{E}_{k_0}[G]$, a function
on $\Omega_j$. We show by induction on $j$ the following. For each $w\in\Omega_j$ there is a
probability measure $\mu_{j,w}$ on $\Omega_{k_0}$ with $G_j(w)=\int G\,d\mu_{j,w}$ for every bounded
measurable $G$, and $w\mapsto\mu_{j,w}(A)$ is measurable for every measurable $A$.

For $j=k_0+1$ take $\mu_{j,w}=\nu_{k_0,w}$. Suppose the claim holds for some $j<K'$. For
$w\in\Omega_{j+1}$ define
\begin{equation*}
\mu_{j+1,w}(A):=\int\mu_{j,w'}(A)\,d\nu_{j,w}(w').
\end{equation*}
The integrand is bounded and measurable in $w'$ by the inductive hypothesis. We have
$\mu_{j+1,w}(\Omega_{k_0})=1$. Countable additivity follows from that of each $\mu_{j,w'}$ and
monotone convergence. Moreover $w\mapsto\mu_{j+1,w}(A)=\mathbb{E}_j[\mu_{j,\cdot}(A)](w)$ is
measurable, because $\mathbb{E}_j$ maps bounded measurable functions to measurable ones. The
identity
\begin{equation*}
\int G\,d\mu_{j+1,w}=\int\Bigl(\int G\,d\mu_{j,w'}\Bigr)d\nu_{j,w}(w')
\end{equation*}
holds for indicators by definition and for simple $G$ by linearity. It then holds for bounded
measurable $G$ by uniform approximation by simple functions. By the inductive hypothesis its
right-hand side is $\int G_j\,d\nu_{j,w}=G_{j+1}(w)$. This completes the induction. We take
$\mu_W:=\mu_{K',W}$, so that $G^{\langle\rangle}=G_{K'}$ is integration against $\mu_W$.

Applying this with $G=F_1F_2$, $G=F_1$ and $G=F_2$ shows that $T^{\langle\rangle}[F_1,F_2](W)$
equals $\operatorname{Cov}(F_1,F_2)$ under $\mu_W$. The first line of
\eqref{8.5:eq-covariance-sup-a}, applied on $(\Omega_{k_0},\mu_W)$, gives the second line, since
the suprema of $|F_1|$ and $|F_2|$ over $\Omega_{k_0}$ do not depend on $W$.

Finally, for every $W$ we have
$|F^{\langle\rangle}(W)|\le\int|F|\,d\mu_W\le\sup|F|$. Hence, for all $W,W'$,
\begin{equation*}
|F^{\langle\rangle}(W)-F^{\langle\rangle}(W')|\le 2\sup|F|,
\end{equation*}
which is the third line of \eqref{8.5:eq-covariance-sup-a}.
\end{proof}

\begin{definition}[The two sizes on the window]\label{8.5:not-window-sizes}
We work in the setting of the two-point witness theorem at the level $k_0$, with the window on
the box $Q_i$ of Definition~\ref{8.5:def-boxed-remainder}. For each index $i$ the remainder
$\operatorname{Irr}_i$ and the main insertion $N_{k_0}\mathcal{O}^{(k_0)}(\tilde f_i)$ are set equal
to zero off the window, as in the two-point witness theorem. Put $x:=x_{k_0}$ and
\begin{equation}\label{8.5:eq-window-sizes-a}
\begin{gathered}
\mathsf{I}_i:=C_{\star}\,x^{-3/2}(\log x)^{p_0+2\mathfrak{p}_0}\,|S_i^{+}|\,\|f_i\|_{\sharp},\\
\mathsf{O}_i:=C_{\mathcal{O}}\,\varepsilon_{k_0}^{2}\,|S_i^{+}|\,\|f_i\|_{\infty},
\qquad C_{\mathcal{O}}:=\tfrac{8}{3}\,C_{\square}C_{\mathfrak{e}}^{2}.
\end{gathered}
\end{equation}
Here $C_{\star}$, $p_0$ and $\mathfrak{p}_0$ are as in Theorem~\ref{8.5:thm-remainder-sup}: $p_0$
is the exponent of the logarithm in $\delta_k=g_kA_1(\log x_k)^{p_0}$ and $\mathfrak{p}_0$ the
exponent in $\varepsilon_k=g_kA_0(\log x_k)^{\mathfrak{p}_0}$. The two exponents are equal, both
being the exponent $p_0$ of the degree function $p_0(\cdot)$; both letters are retained to match
Theorem~\ref{8.5:thm-remainder-sup}.
$C_{\square}$ is the constant of the regularity lemma of the two-point witness theorem's
construction, and $|S_i^{+}|$ is the number of cells of the unit lattice at level $k_0$ in
$S_i^{+}$. Under the hypotheses of Theorem~\ref{8.5:thm-remainder-sup}, read with $f=f_i$,
$S=S_i$, the ball $B(x_i,\rho)$ and the box $Q=Q_i$, and with $\kappa_0\ge 7$ as in its final
display, one has
\begin{equation}\label{8.5:eq-window-sizes-1}
\sup_V|\operatorname{Irr}_i(V)|\le\mathsf{I}_i,\qquad
\sup_V\bigl|N_{k_0}\mathcal{O}^{(k_0)}(\tilde f_i)(V)\bigr|\le\mathsf{O}_i.
\end{equation}
Moreover
\begin{equation}\label{8.5:eq-window-sizes-2}
|S_i^{+}|\,\|f_i\|_{\infty}\le 80\rho^{4}\|f_i\|_{\infty}\le\frac{80\,\mathfrak{n}(f_i)}{\vartheta},
\end{equation}
and $\|f_i\|_{\sharp}=\|f_i\|_{\infty}$ once $\rho\ge 40M/\vartheta$.
\end{definition}

\begin{proof}
Both functions in \eqref{8.5:eq-window-sizes-1} vanish off the window on $Q_i$, so it suffices
to take $V$ in the window on $Q_i$.

For the first bound, apply Theorem~\ref{8.5:thm-remainder-sup}
with $f=f_i$, $S=S_i$, the ball $B(x_i,\rho)$ and the box
$Q=Q_i$, under the hypotheses listed there. Its final display
\eqref{8.5:eq-remainder-sup-c} bounds $|\operatorname{Irr}_i(V)|$ by the expression defining
$\mathsf{I}_i$ in \eqref{8.5:eq-window-sizes-a}, that is, by $\mathsf{I}_i$.

For the second bound, the two-point witness theorem gives $N_{k_0}\le 4/3$,
$|\tilde f_i|\le 2\|f_i\|_{\infty}$ and $\operatorname{supp}\tilde f_i\subset S_i^{+}$. Write
$\mathcal{O}^{(k_0)}(\tilde f_i)=\sum_p\tilde f_i(x(p))\,a_p$, the sum over the plaquettes $p$ of
level $k_0$, with $a_p:=1-\operatorname{Re}\operatorname{tr}U_{k_0}(\partial p)$ at the background
$U_{k_0}=U_{k_0}(V^{\natural})$ of Definition~\ref{8.5:def-boxed-remainder}; $a_p\ge0$ because
the normalised trace of a unitary matrix has modulus at most one. This background is
$\mathfrak{e}$-regular by Lemma~\ref{8.5:lem-minimiser-regular} and
Definition~\ref{8.5:def-boxed-remainder}. The regularity lemma, applied at regularity size
$\mathfrak{e}$, gives
$\sum_{p\subset c}a_p\le C_{\square}\mathfrak{e}^{2}$ for every cell $c$. Summing over the cells
of $S_i^{+}$,
\begin{equation*}
\bigl|N_{k_0}\mathcal{O}^{(k_0)}(\tilde f_i)(V)\bigr|
\le\frac43\cdot 2\|f_i\|_{\infty}\sum_{c\subset S_i^{+}}\sum_{p\subset c}a_p
\le\frac83\,C_{\square}\,\mathfrak{e}^{2}\,|S_i^{+}|\,\|f_i\|_{\infty}.
\end{equation*}
With $\mathfrak{e}\le C_{\mathfrak{e}}\varepsilon_{k_0}$ (the relation fixed with the constant
$C_{\mathfrak{e}}$ in Definition~\ref{8.5:def-boxed-remainder}), the right side is at most
$\mathsf{O}_i$.

For \eqref{8.5:eq-window-sizes-2}: in the setting of the two-point witness theorem,
$|S_i^{+}|\le 80\rho^{4}$ cells and $\vartheta\rho^{4}\|f\|_{\infty}\le\mathfrak{n}(f)$. Together
these give the chain of inequalities. The same setting gives
$\|f\|_{\sharp}=\max\{\|f\|_{\infty},40M\operatorname{Lip}_{k_0}f\}=\|f\|_{\infty}$ once
$\rho\ge 40M/\vartheta$.
\end{proof}

\begin{theorem}[The second-order extraction estimate]\label{8.5:thm-second-order-extraction}
Assume the following statements:
\begin{itemize}
\item the correlation theorem, applied on the torus $\mathbb{T}_s$, that is with $m:=s$;
\item the sup-norm bound on the irrelevant remainder, Theorem~\ref{8.5:thm-remainder-sup};
\item the bound of covariances by sup norms, Lemma~\ref{8.5:lem-covariance-sup};
\item the two sizes on the window, Notation~\ref{8.5:not-window-sizes}, with the bounds recorded
there.
\end{itemize}
Work in the quantifier prefix of the Gaussian extraction, $k_0$ being the extraction level fixed
in that prefix. Take the measure $\nu^B_{K,s}$, the
set $\mathfrak{P}_s$ of translates, the size functional $\mathfrak{n}$ and the main insertions
$N\mathcal{O}^{(k_0)}(\tilde f_i)$ from the two-point witness theorem. Impose the hypotheses of
Theorem~\ref{8.5:thm-remainder-sup}. Let $\operatorname{Irr}_1,\operatorname{Irr}_2$ be the
irrelevant remainders of Notation~\ref{8.5:not-window-sizes}. For real $x>1$ put
\begin{equation}\label{8.5:eq-second-order-extraction-a}
\mathsf{E}^{\mathrm{irr}}(x):=12800\,\vartheta^{-2}\Bigl[\,2C_{\star}C_{\mathcal{O}}A_0^2\,
x^{-1/2}(\log x)^{p_0+4\mathfrak{p}_0}+C_{\star}^2\,x^{-1}(\log x)^{2p_0+4\mathfrak{p}_0}\Bigr].
\end{equation}
Here $p_0$ is the exponent of the logarithm in the correlation theorem as it enters the size
$\mathsf{I}_i$ of Notation~\ref{8.5:not-window-sizes}, and $\mathfrak{p}_0$ is the exponent of
the logarithm in $\varepsilon_{k_0}=g_{k_0}A_0(\log x_{k_0})^{\mathfrak{p}_0}$.
Let $K$ be arbitrary and $s$ any depth in the range of the Gaussian extraction. Let
$(f_1,f_2)$ be any pair admitted by the two-point witness theorem, with its associated
$\tilde f_1,\tilde f_2$, and let $v\in\mathfrak{P}_s$ be any translate. Take
$G_1=\operatorname{Irr}_1$ and $G_2\in\{N\mathcal{O}^{(k_0)}(\tilde f_2),\operatorname{Irr}_2\}$,
or, symmetrically, $G_2=\operatorname{Irr}_2$ and
$G_1\in\{N\mathcal{O}^{(k_0)}(\tilde f_1),\operatorname{Irr}_1\}$. Then
\begin{equation}\label{8.5:eq-second-order-extraction-1}
\bigl|\operatorname{Cov}_{\nu^B_{K,s}}(G_1^v,G_2^v)\bigr|
\le g_{k_0}^4\,\mathfrak{n}(f_1)\,\mathfrak{n}(f_2)\,\mathsf{E}^{\mathrm{irr}}(x_{k_0}).
\end{equation}
Even the sum of the three covariances is bounded by the same quantity:
\begin{equation}\label{8.5:eq-second-order-extraction-2}
\begin{split}
&\bigl|\operatorname{Cov}(\operatorname{Irr}_1^v,N\mathcal{O}^{(k_0)}(\tilde f_2)^v)\bigr|
+\bigl|\operatorname{Cov}(N\mathcal{O}^{(k_0)}(\tilde f_1)^v,\operatorname{Irr}_2^v)\bigr|
+\bigl|\operatorname{Cov}(\operatorname{Irr}_1^v,\operatorname{Irr}_2^v)\bigr|\\
&\qquad\le g_{k_0}^4\,\mathfrak{n}(f_1)\,\mathfrak{n}(f_2)\,\mathsf{E}^{\mathrm{irr}}(x_{k_0}),
\end{split}
\end{equation}
where every covariance is taken under $\nu^B_{K,s}$. No average over $\mathfrak{P}_s$ is
needed. Since the right-hand sides do not depend on $v$, the same bounds hold for the averages
$\langle\!\langle\cdot\rangle\!\rangle$ over $\mathfrak{P}_s$.

These estimates give the bound of part~(b) of the Gaussian extraction with the constant $C_b$ of
that part, the coefficient of its kernel term, equal to zero. The whole covariance is then an
error term of the type of the first member $g_{k_0}^{\sigma_1}\Pi_r$ of the error bound of
part~(a) of the Gaussian extraction, $\sigma_1$ and $\Pi_r$ being the exponent and the factor of
that member. That is, with $\sigma_1=1$,
\begin{align*}
\mathsf{E}^{\mathrm{irr}}(x_{k_0})
&\le g_{k_0}^{\sigma_1}\cdot12800\,\vartheta^{-2}
\Bigl[2C_{\star}C_{\mathcal{O}}A_0^2(\log x_{k_0})^{p_0+4\mathfrak{p}_0}\\
&\qquad\qquad+C_{\star}^2(\log x_{k_0})^{2p_0+4\mathfrak{p}_0}\Bigr].
\end{align*}
\end{theorem}

\begin{proof}
Fix $K$, $s$, the pair $(f_1,f_2)$ and $v\in\mathfrak{P}_s$, and write $x:=x_{k_0}$. By the
hypotheses of Theorem~\ref{8.5:thm-remainder-sup} we have $\log x\ge1$. Since
$x_{k_0}=g_{k_0}^{-2}$ (Definition~\ref{1.2:def-couplings}), $g_{k_0}=x^{-1/2}$ and
$g_{k_0}^4=x^{-2}$. The field-size parameter is
$\varepsilon_{k_0}=g_{k_0}A_0(\log x)^{\mathfrak{p}_0}$, hence
\begin{equation*}
\varepsilon_{k_0}^2=A_0^2\,x^{-1}(\log x)^{2\mathfrak{p}_0}.
\end{equation*}

\emph{Sup bounds.} By Notation~\ref{8.5:not-window-sizes}, in the form
\eqref{8.5:eq-window-sizes-a}, and Theorem~\ref{8.5:thm-remainder-sup}, for $i=1,2$,
\begin{align*}
\sup|\operatorname{Irr}_i|&\le\mathsf{I}_i
=C_{\star}x^{-3/2}(\log x)^{p_0+2\mathfrak{p}_0}\,|S_i^{+}|\,\|f_i\|_{\sharp},\\
\sup|N\mathcal{O}^{(k_0)}(\tilde f_i)|&\le\mathsf{O}_i
=C_{\mathcal{O}}\,\varepsilon_{k_0}^2\,|S_i^{+}|\,\|f_i\|_{\infty}.
\end{align*}
The constants in these sizes do not depend on the position of the box. The same bounds
therefore hold for the translates $\operatorname{Irr}_i^v$ and
$N\mathcal{O}^{(k_0)}(\tilde f_i)^v$. By the same notation, $\|f_i\|_{\sharp}=\|f_i\|_{\infty}$ and
\begin{equation*}
|S_i^{+}|\,\|f_i\|_{\infty}\le 80\,\mathfrak{n}(f_i)/\vartheta .
\end{equation*}

\emph{Covariances.} The measure $\nu^B_{K,s}$ is a probability measure, and the functions
$G_1^v,G_2^v$ are bounded. Lemma~\ref{8.5:lem-covariance-sup}, in the form
\eqref{8.5:eq-covariance-sup-a}, gives
\begin{equation*}
\bigl|\operatorname{Cov}_{\nu^B_{K,s}}(G_1^v,G_2^v)\bigr|\le 2\sup|G_1^v|\,\sup|G_2^v|.
\end{equation*}
For $G_1=\operatorname{Irr}_1$ and $G_2=N\mathcal{O}^{(k_0)}(\tilde f_2)$, insert the bounds
above. The powers of $x$ combine as $x^{-3/2}x^{-1}=x^{-5/2}$, and those of $\log x$ as
$(\log x)^{p_0+2\mathfrak{p}_0}(\log x)^{2\mathfrak{p}_0}=(\log x)^{p_0+4\mathfrak{p}_0}$. The
factors $|S_1^{+}|\,\|f_1\|_{\sharp}$ and $|S_2^{+}|\,\|f_2\|_{\infty}$ contribute at most
$80^2\vartheta^{-2}\mathfrak{n}(f_1)\mathfrak{n}(f_2)$ by the display above.
Hence
\begin{equation*}
2\mathsf{I}_1\mathsf{O}_2\le 2\cdot6400\,\vartheta^{-2}C_{\star}C_{\mathcal{O}}A_0^2\,
x^{-5/2}(\log x)^{p_0+4\mathfrak{p}_0}\,\mathfrak{n}(f_1)\mathfrak{n}(f_2).
\end{equation*}
The symmetric bound holds for $2\mathsf{O}_1\mathsf{I}_2$. For $G_2=\operatorname{Irr}_2$, in
the same way,
\begin{equation*}
2\mathsf{I}_1\mathsf{I}_2\le 2\cdot6400\,\vartheta^{-2}C_{\star}^2\,
x^{-3}(\log x)^{2p_0+4\mathfrak{p}_0}\,\mathfrak{n}(f_1)\mathfrak{n}(f_2).
\end{equation*}
Write $x^{-5/2}=x^{-2}x^{-1/2}$ and $x^{-3}=x^{-2}x^{-1}$, and use $x^{-2}=g_{k_0}^4$. Adding the
three bounds gives
\begin{align*}
2\mathsf{I}_1\mathsf{O}_2+2\mathsf{O}_1\mathsf{I}_2+2\mathsf{I}_1\mathsf{I}_2
&\le g_{k_0}^4\,\mathfrak{n}(f_1)\mathfrak{n}(f_2)\cdot 12800\,\vartheta^{-2}
\Bigl[2C_{\star}C_{\mathcal{O}}A_0^2x^{-1/2}(\log x)^{p_0+4\mathfrak{p}_0}\\
&\qquad\qquad+C_{\star}^2x^{-1}(\log x)^{2p_0+4\mathfrak{p}_0}\Bigr],
\end{align*}
whose right-hand side is
$g_{k_0}^4\mathfrak{n}(f_1)\mathfrak{n}(f_2)\mathsf{E}^{\mathrm{irr}}(x)$ by
\eqref{8.5:eq-second-order-extraction-a}. This is
\eqref{8.5:eq-second-order-extraction-2}. Each single covariance in
\eqref{8.5:eq-second-order-extraction-1} is one of the three summands, so
\eqref{8.5:eq-second-order-extraction-1} follows.

\emph{Averages.} Since the bound is the same for every $v\in\mathfrak{P}_s$, the triangle
inequality gives
\begin{equation*}
|\langle\!\langle\operatorname{Cov}(G_1^v,G_2^v)\rangle\!\rangle|
\le\langle\!\langle|\operatorname{Cov}(G_1^v,G_2^v)|\rangle\!\rangle
\le g_{k_0}^4\mathfrak{n}(f_1)\mathfrak{n}(f_2)\mathsf{E}^{\mathrm{irr}}(x).
\end{equation*}

\emph{Type of the bound.} In \eqref{8.5:eq-second-order-extraction-a} the first term carries
$x^{-1/2}=g_{k_0}$. The second carries $x^{-1}=g_{k_0}^2\le g_{k_0}$, since $x\ge1$. Both are
multiplied by constants and powers of $\log x$. Replacing $x^{-1}$ by $g_{k_0}$ in the second
term gives the displayed bound on $\mathsf{E}^{\mathrm{irr}}(x_{k_0})$ with $\sigma_1=1$.
\end{proof}

\begin{corollary}[The irrelevant term against the main term]\label{8.5:cor-irrelevant-small}
Assume the setting and the ceilings of the second-order extraction estimate,
Theorem~\ref{8.5:thm-second-order-extraction}, and let $K$, $m$, $s$, the pair $(f_1,f_2)$ and the
level $k_0$ be as there. Assume the following inputs (a)--(d) from the two-point witness theorem,
and the inequalities (e):
\begin{itemize}
\item[(a)] the irrelevant term $R^{\mathrm{irr}}$ of the two-point witness theorem is bounded by the
$\mathfrak{P}_s$-average of three covariances,
\begin{multline*}
R^{\mathrm{irr}}\le\Big\langle\!\Big\langle\,
\big|\operatorname{Cov}_{\nu^B_{K,s}}\big(\operatorname{Irr}_1^v,N\mathcal{O}^{(k_0)}(\tilde f_2)^v\big)\big|
+\big|\operatorname{Cov}_{\nu^B_{K,s}}\big(N\mathcal{O}^{(k_0)}(\tilde f_1)^v,\operatorname{Irr}_2^v\big)\big|\\
+\big|\operatorname{Cov}_{\nu^B_{K,s}}\big(\operatorname{Irr}_1^v,\operatorname{Irr}_2^v\big)\big|
\,\Big\rangle\!\Big\rangle,
\end{multline*}
where $\operatorname{Cov}_{\nu^B_{K,s}}$ is the covariance in $\nu^B_{K,s}$, the superscript $v$
denotes the translate by $v$, and $\langle\!\langle\cdot\rangle\!\rangle$ is the average over the
translates $v\in\mathfrak{P}_s$;
\item[(b)] the main term satisfies $\mathfrak{M}\ge c_M\mathfrak{S}$, where
$\mathfrak{S}=g_{k_0}^4\mathfrak{n}(f_1)\mathfrak{n}(f_2)\rho^{-8}$, $\rho>0$ being the parameter of
the two-point witness theorem that enters $\mathfrak{S}$, with $c_M>0$ and $\mathfrak{n}(f_i)\ge0$;
\item[(c)] $\rho<LC_W(s)$, where $C_W(s)=C_0A_0^2(\log\bar x_s)^{2\mathfrak{p}_0}$;
\item[(d)] $x_{k_0}\ge\underline{x}_s$ and $\bar x_s\le\bar x(\underline{x}_s)$, where $\bar x(\cdot)$
is the envelope of the two-point witness theorem;
\item[(e)] $\underline{x}_s\ge X-c_5$ and $\bar x_s>1$, where $X=g^{-2}$ and $c_5=2c_2$ is the
constant of Definition~\ref{1.2:def-flow-constants}.
\end{itemize}
Put $\mathsf{E}^{\mathrm{irr}}_{\downarrow}(y)=\sup_{x\ge y}\mathsf{E}^{\mathrm{irr}}(x)$, with
$\mathsf{E}^{\mathrm{irr}}$ as in \eqref{8.5:eq-second-order-extraction-a}. If
\begin{equation}\label{8.5:eq-irrelevant-small-a}
\sup_{y\ge X-c_5}\Big(LC_0A_0^2\big(\log\bar x(y)\big)^{2\mathfrak{p}_0}\Big)^{8}
\,\mathsf{E}^{\mathrm{irr}}_{\downarrow}(y)\le\frac{c_M}{32},
\end{equation}
then
\begin{equation*}
R^{\mathrm{irr}}\le g_{k_0}^4\mathfrak{n}(f_1)\mathfrak{n}(f_2)\,\mathsf{E}^{\mathrm{irr}}(x_{k_0})
\le\frac{\mathfrak{M}}{32}.
\end{equation*}
\end{corollary}

\begin{proof}
By Theorem~\ref{8.5:thm-second-order-extraction}, for every translate $v$ the sum of the three
covariances in (a) is at most $g_{k_0}^4\mathfrak{n}(f_1)\mathfrak{n}(f_2)\mathsf{E}^{\mathrm{irr}}(x_{k_0})$.
This bound does not depend on $v$, and an average of numbers each bounded by the same constant is bounded
by that constant; by (a), the first inequality follows.

For the second inequality, write
\begin{equation*}
g_{k_0}^4\mathfrak{n}(f_1)\mathfrak{n}(f_2)\mathsf{E}^{\mathrm{irr}}(x_{k_0})
=\mathfrak{S}\rho^{8}\mathsf{E}^{\mathrm{irr}}(x_{k_0}).
\end{equation*}
If
$\mathsf{E}^{\mathrm{irr}}(x_{k_0})<0$, this is at most $0$, while $\mathfrak{M}\ge c_M\mathfrak{S}\ge0$
by (b); the inequality holds. Let $\mathsf{E}^{\mathrm{irr}}(x_{k_0})\ge0$. By (b),
$\mathfrak{S}\le\mathfrak{M}/c_M$, so
\begin{equation}\label{8.5:eq-irrelevant-small-1}
g_{k_0}^4\mathfrak{n}(f_1)\mathfrak{n}(f_2)\mathsf{E}^{\mathrm{irr}}(x_{k_0})
\le\frac{\mathfrak{M}}{c_M}\,\rho^{8}\mathsf{E}^{\mathrm{irr}}(x_{k_0}),
\end{equation}
and it suffices to show $\rho^8\mathsf{E}^{\mathrm{irr}}(x_{k_0})\le c_M/32$. By (c) and $\rho>0$,
$0<\rho^8<(LC_W(s))^8$. By (e), $\bar x_s>1$, so $\log\bar x_s>0$; since $C_W(s)>\rho/L>0$ by (c),
this gives $C_0A_0^2>0$. By (d) and the
definition of $\mathsf{E}^{\mathrm{irr}}_{\downarrow}$,
$\mathsf{E}^{\mathrm{irr}}(x_{k_0})\le\mathsf{E}^{\mathrm{irr}}_{\downarrow}(\underline{x}_s)$. By (d),
$\bar x_s\le\bar x(\underline{x}_s)$; the logarithm is increasing, and since $\mathfrak{p}_0>0$ in
the setting of Theorem~\ref{8.5:thm-second-order-extraction}, the map $t\mapsto t^{2\mathfrak{p}_0}$
is increasing on $[0,\infty)$. Multiplying by $C_0A_0^2>0$ gives
$C_W(s)\le C_0A_0^2(\log\bar x(\underline{x}_s))^{2\mathfrak{p}_0}$. Hence
\begin{align*}
\rho^8\mathsf{E}^{\mathrm{irr}}(x_{k_0})
&\le\big(LC_W(s)\big)^8\,\mathsf{E}^{\mathrm{irr}}_{\downarrow}(\underline{x}_s)\\
&\le\Big(LC_0A_0^2\big(\log\bar x(\underline{x}_s)\big)^{2\mathfrak{p}_0}\Big)^8
\mathsf{E}^{\mathrm{irr}}_{\downarrow}(\underline{x}_s)
\le\frac{c_M}{32},
\end{align*}
the last step by \eqref{8.5:eq-irrelevant-small-a} at $y=\underline{x}_s$, which is admissible since
$\underline{x}_s\ge X-c_5$ by (e). Inserting this into \eqref{8.5:eq-irrelevant-small-1} gives the
second inequality.
\end{proof}

By \eqref{8.5:eq-second-order-extraction-a}, $\mathsf{E}^{\mathrm{irr}}(x)$ is the sum of a constant
multiple of $x^{-1/2}(\log x)^{p_0+4\mathfrak{p}_0}$ and a constant multiple of
$x^{-1}(\log x)^{2p_0+4\mathfrak{p}_0}$, and the function of $y$ under the supremum in
\eqref{8.5:eq-irrelevant-small-a} is $\mathsf{E}^{\mathrm{irr}}_{\downarrow}(y)$ multiplied by
$(LC_0A_0^2)^8(\log\bar x(y))^{16\mathfrak{p}_0}$; it decreases to zero as $y\to\infty$. Condition
\eqref{8.5:eq-irrelevant-small-a} is therefore a ceiling on the reference coupling $g$, that is, a
lower bound on $X=g^{-2}$, of the same kind as one of the ceilings already collected in $\gamma_W$:
a single ceiling, the same for all $K$, $m$ and all $s$ for which (a)--(e) hold, chosen after $C_0$.
Consequently, in the
bookkeeping of the two-point witness theorem, the condition $C_0\ge(448C_b)^{1/2}$ in the list of
conditions on $C_0$ can be omitted, where $C_b$ is the constant of the irrelevant-term bound of
the Gaussian extraction, of which Theorem~\ref{8.5:thm-second-order-extraction} is the case $C_b=0$;
the bound $|R_K|\le 6/32$ is unchanged.

\begin{corollary}[The estimate on towers of one-step laws]\label{8.5:cor-one-step-towers}
Let the following hold.
\begin{itemize}
\item[(a)] The hypotheses, and hence the conclusion, of the sup-norm bound on the irrelevant
remainder (Theorem~\ref{8.5:thm-remainder-sup}) hold for the test functions $f_1,f_2$ of two
shapes. In
particular \eqref{8.5:eq-regular-extension-a} and \eqref{8.5:eq-remainder-sup-a} hold at the
radius $\rho$ of each shape, where $\rho\le D_{\max}/2$ and $D_{\max}$ is the constant of the
flat bound on the truncated irrelevant remainder, which admits shapes of radii
$\rho\le\rho_\star\le D_{\max}/2$.
\item[(b)] The statement \eqref{8.5:eq-covariance-sup-a} of Lemma~\ref{8.5:lem-covariance-sup}
(covariances bounded by sup norms) holds for towers of positive normalised functionals.
\item[(c)] $T^{\langle h_0\rangle}$ is the truncation along the tower
$\mathbb{E}_{K'-1}\circ\cdots\circ\mathbb{E}_{k_0}$ of the one-step laws $\mathbb{E}_j$,
$k_0\le j\le K'-1$, of the history $h_0$ in which every step is small-field; here $K'-1$ is the
index of the last one-step law of the tower. It is a tower of
positive normalised functionals whose state space is the set of configurations $V$ in the window:
\[
T^{\langle h_0\rangle}[F_1,F_2]:=(F_1F_2)^{\langle h_0\rangle}
-F_1^{\langle h_0\rangle}F_2^{\langle h_0\rangle}.
\]
\end{itemize}
Write $x:=x_{k_0}$ and $\operatorname{Irr}_i:=\operatorname{Irr}[f_i]$ for $i=1,2$, let
$N\mathcal{O}^{(k_0)}(\tilde f_i)$ be the main insertion of $f_i$ as in
Definition~\ref{8.5:not-window-sizes}, and let $\mathsf{E}^{\mathrm{irr}}$ be the function of
\eqref{8.5:eq-second-order-extraction-a}. Then for every pair of shapes, at every distance between
them, and for $G_2\in\{N\mathcal{O}^{(k_0)}(\tilde f_2),\operatorname{Irr}_2\}$,
\begin{equation}\label{8.5:eq-one-step-towers-1}
\bigl|T^{\langle h_0\rangle}[\operatorname{Irr}_1,G_2](\mathbf{1})\bigr|
\le g_{k_0}^4\,\mathfrak{n}(f_1)\mathfrak{n}(f_2)\,\mathsf{E}^{\mathrm{irr}}(x),
\end{equation}
and symmetrically with the roles of the two shapes exchanged. Consequently the flat bound on the
truncated irrelevant remainder along this tower, whose test functions $\varphi_1,\varphi_2$ are
here $f_1,f_2$, holds with its constant $C_b$ equal to $0$ and
with its error function $E^{\flat}(x;\rho_\star)$ taken to be $\mathsf{E}^{\mathrm{irr}}(x)$;
here $\rho_\star$ is the upper bound on the radii of the shapes admitted in the flat bound, and
$\mathsf{E}^{\mathrm{irr}}(x)$ does not depend on $\rho_\star$.
\end{corollary}

\begin{proof}
By hypothesis~(c) and \eqref{8.5:eq-covariance-sup-a} of hypothesis~(b), the quantity
$T^{\langle h_0\rangle}[\operatorname{Irr}_1,G_2](\mathbf{1})$ is a covariance under a probability
measure on the window. Therefore
\[
\bigl|T^{\langle h_0\rangle}[\operatorname{Irr}_1,G_2](\mathbf{1})\bigr|
\le 2\sup|\operatorname{Irr}_1|\cdot\sup|G_2|,
\]
where the suprema are taken over $V$ in the window.

By hypothesis~(a) and Theorem~\ref{8.5:thm-remainder-sup}, $\sup|\operatorname{Irr}_i|\le\mathsf{I}_i$.
The sizes of Definition~\ref{8.5:not-window-sizes} also give
$\sup|N\mathcal{O}^{(k_0)}(\tilde f_i)|\le\mathsf{O}_i$. In the notation of that definition,
\[
\begin{gathered}
\mathsf{I}_i=C_{\star}\,x^{-3/2}(\log x)^{p_0+2\mathfrak{p}_0}\,|S_i^+|\,\|f_i\|_\sharp,\\
\mathsf{O}_i=C_{\mathcal{O}}\,\varepsilon_{k_0}^2\,|S_i^+|\,\|f_i\|_\infty,
\end{gathered}
\]
with $\varepsilon_{k_0}=g_{k_0}A_0(\log x)^{\mathfrak{p}_0}$, and the same definition records
\[
|S_i^+|\,\|f_i\|_\sharp=|S_i^+|\,\|f_i\|_\infty\le 80\,\mathfrak{n}(f_i)/\vartheta .
\]
Inserting this bound and $\varepsilon_{k_0}^2=A_0^2x^{-1}(\log x)^{2\mathfrak{p}_0}$, which holds
since $g_{k_0}^2=x^{-1}$, into the expressions for $\mathsf{I}_i$ and $\mathsf{O}_i$ gives
\begin{equation}\label{8.5:eq-one-step-towers-2}
\begin{aligned}
2\mathsf{I}_1\mathsf{O}_2&\le 2\cdot 6400\,\vartheta^{-2}C_{\star}C_{\mathcal{O}}A_0^2\,
x^{-5/2}(\log x)^{p_0+4\mathfrak{p}_0}\,\mathfrak{n}(f_1)\mathfrak{n}(f_2),\\
2\mathsf{I}_1\mathsf{I}_2&\le 2\cdot 6400\,\vartheta^{-2}C_{\star}^2\,
x^{-3}(\log x)^{2p_0+4\mathfrak{p}_0}\,\mathfrak{n}(f_1)\mathfrak{n}(f_2).
\end{aligned}
\end{equation}
Since $g_{k_0}^4=x^{-2}$, the first right-hand side of \eqref{8.5:eq-one-step-towers-2} is one
half of, and the second is equal to, $g_{k_0}^4\mathfrak{n}(f_1)\mathfrak{n}(f_2)$ times the
corresponding term of $\mathsf{E}^{\mathrm{irr}}(x)$ in \eqref{8.5:eq-second-order-extraction-a}.
Both terms of $\mathsf{E}^{\mathrm{irr}}(x)$ are non-negative, so in either case the right-hand
side is at most $g_{k_0}^4\mathfrak{n}(f_1)\mathfrak{n}(f_2)\mathsf{E}^{\mathrm{irr}}(x)$. This gives
\eqref{8.5:eq-one-step-towers-1}.

The symmetric statement follows from the same argument with the indices $1$ and $2$ exchanged.

The constants $C_{\star}$ and $C_{\mathcal{O}}$ involve no position, so the bound holds for
every pair of shapes and at every distance. The right-hand side of
\eqref{8.5:eq-one-step-towers-1} is the right-hand side of the flat bound with $C_b=0$ and with
$E^{\flat}(x;\rho_\star)$ replaced by $\mathsf{E}^{\mathrm{irr}}(x)$: the term of the flat bound
that carries the factor $C_b$ is not needed. The function $\mathsf{E}^{\mathrm{irr}}(x)$ of
\eqref{8.5:eq-second-order-extraction-a} does not involve $\rho_\star$.

The choice of the admissible source radius $r$ at which the correlation theorem is applied does
not matter either: $\operatorname{Irr}$ is a first derivative in the source variable at $z=0$, so
it does not depend on that radius, and the terms of $C_{\star}$ that depend on $r$ only decrease
as $r$ grows.
\end{proof}

\begin{corollary}[Weak sup, oscillation and truncation bounds]\label{8.5:cor-weak-bounds}
Assume the following two statements as hypotheses:
\begin{itemize}
\item[(1)] the sup-norm bound on the irrelevant remainder, Theorem~\ref{8.5:thm-remainder-sup},
in the form \eqref{8.5:eq-remainder-sup-c}, in its setting and under its hypotheses, for
each of two test functions $f_1,f_2$ satisfying its hypotheses (with supports in balls
$B(x_a,\rho)$ and boxes $Q_a$, $a=1,2$);
\item[(2)] the covariance bounds \eqref{8.5:eq-covariance-sup-a} of
Lemma~\ref{8.5:lem-covariance-sup}.
\end{itemize}
Write $x:=x_{k_0}$; it is not to be confused with the centres $x_a$. Let $\sup$ denote the
supremum over the window on $Q_a$, off
which $\operatorname{Irr}_a$ vanishes, and use the sizes $\mathsf{I}_a$, $\mathsf{O}_a$ of
Notation~\ref{8.5:not-window-sizes}. Let $(\mathbb{E}_j)$ be a tower of positive normalised
functionals, and let $F^{\langle\rangle}$ and $T^{\langle\rangle}$ be as in
Lemma~\ref{8.5:lem-covariance-sup}. Then (a) and (b) hold for $a=1,2$, and (c) holds:
\begin{itemize}
\item[(a)] the weak sup bound holds:
\begin{equation*}
\sup|\operatorname{Irr}_a|
\le 80\,C_{\star}\,\rho^4\,\|f_a\|_{\sharp}\,x^{-3/2}(\log x)^{p_0+2\mathfrak{p}_0};
\end{equation*}
in particular $\sup|\operatorname{Irr}_a|\le\|f_a\|_{\sharp}\,x^{-1}$ whenever the
ceiling
\begin{equation}\label{8.5:eq-weak-bounds-1}
\sup_{y\ge X-c_5}80\,C_{\star}\bigl(LC_0A_0^2(\log\bar x(y))^{2\mathfrak{p}_0}\bigr)^4
\sup_{x'\ge y}x'^{-1/2}(\log x')^{p_0+2\mathfrak{p}_0}\le 1
\end{equation}
holds, with the function $\bar x(\cdot)$ of \eqref{8.5:eq-irrelevant-small-a}; like that
display, it is one condition for every $s$, $K$, $m$;
\item[(b)] the oscillation bound
$\operatorname{osc}(\operatorname{Irr}_a^{\langle\rangle})\le 2\sup|\operatorname{Irr}_a|$
holds;
\item[(c)] for $G_2\in\{N_{k_0}\mathcal{O}^{(k_0)}(\tilde f_2),\operatorname{Irr}_2\}$ and every
$W$, the truncation bound holds:
\begin{equation*}
\bigl|T^{\langle\rangle}[\operatorname{Irr}_1,G_2](W)\bigr|
\le g_{k_0}^4\,\mathfrak{n}(f_1)\,\mathfrak{n}(f_2)\,\mathsf{E}^{\mathrm{irr}}(x),
\end{equation*}
with $\mathsf{E}^{\mathrm{irr}}$ as in \eqref{8.5:eq-second-order-extraction-a}. If moreover
\begin{equation*}
\mathfrak{M}\ge c_M\,g_{k_0}^4\,\mathfrak{n}(f_1)\,\mathfrak{n}(f_2)\,\rho^{-8}
\end{equation*}
and the ceiling \eqref{8.5:eq-irrelevant-small-a} holds with $c_M/64$ in place of $c_M/32$,
then the right-hand side is at most $\mathfrak{M}/64$.
\end{itemize}
The same bounds hold with $f_a$ replaced by $f_a\circ\iota^{-1}$ for every lattice isometry
$\iota$. The constants read $f_a$ only through $\|f_a\|_{\sharp}$ and $|S_a^+|$.
\end{corollary}

\begin{proof}
\emph{(a).} Hypothesis~(1) gives $\sup|\operatorname{Irr}_a|\le\mathsf{I}_a$, where
\begin{equation*}
\mathsf{I}_a=C_{\star}x^{-3/2}(\log x)^{p_0+2\mathfrak{p}_0}\,|S_a^+|\,\|f_a\|_{\sharp}.
\end{equation*}
By Notation~\ref{8.5:not-window-sizes}, $|S_a^+|\le80\rho^4$. Inserting this bound gives the
first inequality. For the second, we use the relations
\begin{equation*}
\rho<LC_W(s)=LC_0A_0^2(\log\bar x_s)^{2\mathfrak{p}_0},\qquad x_{k_0}\ge\underline{x}_s,
\qquad \bar x_s\le\bar x(\underline{x}_s).
\end{equation*}
These relations, together with the fact that $\underline{x}_s$ lies in the range
$y\ge X-c_5$ of the suprema in \eqref{8.5:eq-weak-bounds-1} and
\eqref{8.5:eq-irrelevant-small-a}, hold in the setting of hypothesis~(1); they are the ones
used in the proof of Corollary~\ref{8.5:cor-irrelevant-small}. Take $y:=\underline{x}_s$.
Since the logarithm is increasing, the first and third relations give the first, and the
second relation gives the second, of the inequalities
\begin{gather*}
\rho^4\le\bigl(LC_0A_0^2(\log\bar x(\underline{x}_s))^{2\mathfrak{p}_0}\bigr)^4,\\
x^{-1/2}(\log x)^{p_0+2\mathfrak{p}_0}\le\sup_{x'\ge\underline{x}_s}
x'^{-1/2}(\log x')^{p_0+2\mathfrak{p}_0}.
\end{gather*}
Hence, by \eqref{8.5:eq-weak-bounds-1},
\begin{equation*}
80\,C_{\star}\,\rho^4\,x^{-1/2}(\log x)^{p_0+2\mathfrak{p}_0}\le 1.
\end{equation*}
Multiplying this inequality by $\|f_a\|_{\sharp}x^{-1}$ gives the second inequality.

\emph{(b).} Apply the oscillation bound of Lemma~\ref{8.5:lem-covariance-sup}, contained in
\eqref{8.5:eq-covariance-sup-a}, to $F=\operatorname{Irr}_a$.

\emph{(c).} By Lemma~\ref{8.5:lem-covariance-sup},
$T^{\langle\rangle}[\operatorname{Irr}_1,G_2](W)$ is, for each
$W$, a covariance under a probability measure. Hence the covariance bound of
\eqref{8.5:eq-covariance-sup-a} applies and gives the bound
$2\sup|\operatorname{Irr}_1|\sup|G_2|$. By Notation~\ref{8.5:not-window-sizes}, this is at most
$2\mathsf{I}_1\mathsf{O}_2$ for $G_2=N_{k_0}\mathcal{O}^{(k_0)}(\tilde f_2)$, and at most
$2\mathsf{I}_1\mathsf{I}_2$ for $G_2=\operatorname{Irr}_2$.

Notation~\ref{8.5:not-window-sizes} gives three further facts. First,
$|S_a^+|\|f_a\|_{\infty}\le80\,\mathfrak{n}(f_a)/\vartheta$. Second, since
$\rho\ge 40M/\vartheta$ in the setting of hypothesis~(1), $\|f_a\|_{\sharp}=\|f_a\|_{\infty}$.
Third, $\mathsf{O}_2$ carries the
factor
\begin{equation*}
\varepsilon_{k_0}^2=x^{-1}A_0^2(\log x)^{2\mathfrak{p}_0}.
\end{equation*}
With these three facts,
\begin{align*}
2\mathsf{I}_1\mathsf{O}_2
&\le 12800\,\vartheta^{-2}C_{\star}C_{\mathcal{O}}A_0^2\,x^{-5/2}
(\log x)^{p_0+4\mathfrak{p}_0}\,\mathfrak{n}(f_1)\mathfrak{n}(f_2),\\
2\mathsf{I}_1\mathsf{I}_2
&\le 12800\,\vartheta^{-2}C_{\star}^2\,x^{-3}
(\log x)^{2p_0+4\mathfrak{p}_0}\,\mathfrak{n}(f_1)\mathfrak{n}(f_2).
\end{align*}
Both terms of $\mathsf{E}^{\mathrm{irr}}(x)$ are non-negative; by
\eqref{8.5:eq-second-order-extraction-a}, $\mathsf{E}^{\mathrm{irr}}(x)$ is the sum
\begin{align*}
\mathsf{E}^{\mathrm{irr}}(x)
&=25600\,\vartheta^{-2}C_{\star}C_{\mathcal{O}}A_0^2\,x^{-1/2}(\log x)^{p_0+4\mathfrak{p}_0}\\
&\quad+12800\,\vartheta^{-2}C_{\star}^2\,x^{-1}(\log x)^{2p_0+4\mathfrak{p}_0}.
\end{align*}
Hence
each right-hand side above is at most
$x^{-2}\mathfrak{n}(f_1)\mathfrak{n}(f_2)\mathsf{E}^{\mathrm{irr}}(x)$. Since
$g_{k_0}^4=x^{-2}$, this proves the first bound in~(c).

For the second bound, by the hypothesis
$\mathfrak{M}\ge c_Mg_{k_0}^4\mathfrak{n}(f_1)\mathfrak{n}(f_2)\rho^{-8}$ it suffices to
show that $\rho^8\mathsf{E}^{\mathrm{irr}}(x_{k_0})\le c_M/64$. This follows, as in the proof of
Corollary~\ref{8.5:cor-irrelevant-small}, from the three relations displayed in the proof
of~(a), with the function $\bar x(\cdot)$ of \eqref{8.5:eq-irrelevant-small-a},
together with the monotone majorant $\mathsf{E}^{\mathrm{irr}}_{\downarrow}$, read with
$c_M/64$ in place of $c_M/32$.

Finally, Theorem~\ref{8.5:thm-remainder-sup} holds for every $f$ in its setting. A lattice
isometry $\iota$ carries that setting about $x_a$ to the same setting about $\iota x_a$
(Lemma~\ref{1.1:lem-invariance}) and preserves $\|\cdot\|_{\sharp}$ and $|S_a^+|$. In particular
it holds for $f_a\circ\iota^{-1}$. The constants $C_{\star}$, $C_{\mathcal{O}}$, $A_0$ and
$\vartheta$ do not depend on $f$, so the argument above applies unchanged to
$f_a\circ\iota^{-1}$.
\end{proof}

The bounds (a)--(c) of Corollary~\ref{8.5:cor-weak-bounds} are the necessary core of a
stronger membership statement formulated for the Gaussian extraction. That statement
says that, on the window, $\operatorname{Irr}_a$ lies within $e^{-c_{\eta}\rho}$ of one function
belonging to the Gaussian-extraction class $\mathfrak{F}_{k_0}(\cdot;\cdot,\varrho_{\eta})$. The
membership statement itself is not asserted here, for the following reason.

The part $\operatorname{Irr}^P$, composed with the local background, is analytic at an absolute
radius, by part~(a) of the absolute-radius response theorem. The other two parts behave
differently. The sum $\operatorname{Irr}^Q+\operatorname{Irr}^R$ is of size
$x^{-3/2}/(\lambda r)$, where $\lambda$ and $r$ are the parameters of the level-dependent radius
system. It is known to be analytic only on the $\flat$ tube, that is, on the complex
backgrounds $\mathcal{U}^{\flat}_j$ at the level-dependent (shrinking) radii, and it is not small
of order $e^{-c_{\eta}\rho}$.

Only the bounds (a)--(c), and not the membership statement, are used in the proof of the
Gaussian extraction.

\begin{theorem}[The second-order estimate under an exponent floor]\label{8.5:thm-exponent-floor}
Let $q_0=q_1=q$ be Ba{\l}aban's free integer exponents of \cite{B-CMP119}, which are required there
to be integers greater than $1$. Let $K$, $s$, the level $k_0$, the pair $(f_1,f_2)$ with the
functions $\tilde f_i$ and the support sets $S_i^{+}$ ($i=1,2$), the functionals
$\operatorname{Irr}_i$ and $N\mathcal{O}^{(k_0)}(\tilde f_i)$, the translate $v$ and the measure
$\nu^B_{K,s}$ be as in the second-order extraction estimate
(Theorem~\ref{8.5:thm-second-order-extraction}), and write $x:=x_{k_0}$. Assume:
\begin{itemize}
\item[(a)] the correlation theorem on the torus $\mathbb{T}_s$, that is, with $m:=s$; the
proposition on the main insertion in the construction of the two-point witness theorem; the
regularity lemma at regularity size $\mathfrak{e}$, as it enters
Definition~\ref{8.5:not-window-sizes}; and the hypotheses of the sup-norm bound on shrinking radii
(Theorem~\ref{8.5:thm-remainder-sup-shrinking}) other than the conditions listed in~(b);
\item[(b)] the conditions \eqref{8.5:eq-source-quadratic-a-bis}, \eqref{8.5:eq-local-source-split-b},
\eqref{8.5:eq-renormalised-groups-b}, \eqref{8.5:eq-regular-extension-a} and
\eqref{8.5:eq-remainder-sup-a};
\item[(c)] the floor \textup{(F-$q$)}:
\begin{equation}\label{8.5:eq-exponent-floor-1}
q\ \ge\ p_0+20\mathfrak{p}_0+1 ,
\end{equation}
where $p_0$ is the exponent in $\theta_{k_0}=C_{\theta}(\log x)^{p_0-q}$, the ratio entering the
sup-norm bound on shrinking radii (Theorem~\ref{8.5:thm-remainder-sup-shrinking}), and
$\mathfrak{p}_0$ is the exponent in $\varepsilon_k=g_kA_0(\log x_k)^{\mathfrak{p}_0}$;
\eqref{8.5:eq-exponent-floor-1} is a lower bound on $q$ which reads only these two exponents.
\end{itemize}
Let $\rho$ be the radius parameter of the construction of the two-point witness theorem,
$C_W(s)=C_0A_0^2(\log\bar x_s)^{2\mathfrak{p}_0}$, with its constant $C_0$, the quantity for which
$\rho<LC_W(s)$ there, $\underline{x}_s$ the lower bound of that construction on $x_{k_0}$,
$\bar x_s$ the point at which the logarithm in $C_W(s)$ is evaluated, and $\bar x(\cdot)$ the
function of that construction with
$\bar x_s\le\bar x(\underline{x}_s)$. In the same construction let $\mathfrak{M}$ be the main term,
$c_M$ the constant of its lower bound, and $R^{\mathrm{irr}}$ the irrelevant remainder.
Put
\begin{equation}\label{8.5:eq-exponent-floor-2}
\begin{split}
\mathsf{E}^{\mathrm{irr},\flat}(x) := 12800\,\vartheta^{-2}\Bigl[\,
&2\bigl(2C_{\star}^{\flat}C_{\theta}C_{\mathfrak{e}}^2A_0^2(\log x)^{p_0-q+2\mathfrak{p}_0}
+x\,\mathsf{I}^R\bigr)\,C_{\mathcal{O}}A_0^2(\log x)^{2\mathfrak{p}_0}\\
&+\bigl(2C_{\star}^{\flat}C_{\theta}C_{\mathfrak{e}}^2A_0^2(\log x)^{p_0-q+2\mathfrak{p}_0}
+x\,\mathsf{I}^R\bigr)^2\,\Bigr],
\end{split}
\end{equation}
with $C_{\star}^{\flat}$ as in \eqref{8.5:eq-remainder-sup-shrinking-a} and $\mathsf{I}^R$ as in
\eqref{8.5:eq-remainder-sup-b}, and
$\mathsf{E}^{\mathrm{irr},\flat}_{\downarrow}(y):=\sup_{x\ge y}\mathsf{E}^{\mathrm{irr},\flat}(x)$.
Then each of the three covariances
\begin{gather*}
\operatorname{Cov}_{\nu^B_{K,s}}\bigl(\operatorname{Irr}_1^{v},(N\mathcal{O}^{(k_0)}(\tilde f_2))^{v}\bigr),\\
\operatorname{Cov}_{\nu^B_{K,s}}\bigl((N\mathcal{O}^{(k_0)}(\tilde f_1))^{v},\operatorname{Irr}_2^{v}\bigr),\\
\operatorname{Cov}_{\nu^B_{K,s}}\bigl(\operatorname{Irr}_1^{v},\operatorname{Irr}_2^{v}\bigr)
\end{gather*}
is at most $g_{k_0}^4\,\mathfrak{n}(f_1)\mathfrak{n}(f_2)\,\mathsf{E}^{\mathrm{irr},\flat}(x)$ in
absolute value, and so is the sum of their absolute values. Moreover
$\rho^8\mathsf{E}^{\mathrm{irr},\flat}$ has degree at most
$p_0-q+20\mathfrak{p}_0\le-1$ in $\log x$, and the irrelevant remainder satisfies
$R^{\mathrm{irr}}\le\mathfrak{M}/32$ provided
\begin{equation}\label{8.5:eq-exponent-floor-3}
\sup_{y}\,\bigl(LC_0A_0^2(\log\bar x(y))^{2\mathfrak{p}_0}\bigr)^{8}\,
\mathsf{E}^{\mathrm{irr},\flat}_{\downarrow}(y)\ \le\ \frac{c_M}{32}.
\end{equation}
\end{theorem}

\begin{proof}
Fix $i\in\{1,2\}$. Since $x_k=g_k^{-2}$, we have $g_{k_0}^4=x^{-2}$.

\emph{The irrelevant remainder.} By the sup-norm bound on shrinking radii
(Theorem~\ref{8.5:thm-remainder-sup-shrinking}), whose conditions are those listed in~(b) and whose
remaining hypotheses are assumed in~(a),
applied to $f=f_i$ with support set $S^{+}=S_i^{+}$, we have, uniformly in the background,
\begin{equation*}
|\operatorname{Irr}_i|\le\bigl[C_{\star}^{\flat}\,\hat\theta_{k_0-1}\,\mathfrak{e}^2+\mathsf{I}^R\bigr]
\,|S_i^{+}|\,\|f_i\|_{\sharp}.
\end{equation*}
At level $k_0$ the running maximum satisfies
$\hat\theta_{k_0-1}\le2\theta_{k_0}=2C_{\theta}(\log x)^{p_0-q}$, and the regularity size satisfies
$\mathfrak{e}^2=C_{\mathfrak{e}}^2A_0^2(\log x)^{2\mathfrak{p}_0}/x$. Inserting both,
\begin{equation*}
\sup|\operatorname{Irr}_i|\le x^{-1}\,a\,|S_i^{+}|\,\|f_i\|_{\sharp},\qquad
a:=2C_{\star}^{\flat}C_{\theta}C_{\mathfrak{e}}^2A_0^2(\log x)^{p_0-q+2\mathfrak{p}_0}+x\,\mathsf{I}^R .
\end{equation*}

\emph{The main insertion.} By the two sizes on the window (Definition~\ref{8.5:not-window-sizes},
display \eqref{8.5:eq-window-sizes-a}),
\begin{align*}
\sup|N\mathcal{O}^{(k_0)}(\tilde f_i)|&\le\mathsf{O}_i
=C_{\mathcal{O}}\varepsilon_{k_0}^2\,|S_i^{+}|\,\|f_i\|_\infty,\\
\varepsilon_{k_0}^2&=g_{k_0}^2A_0^2(\log x)^{2\mathfrak{p}_0}=x^{-1}A_0^2(\log x)^{2\mathfrak{p}_0}.
\end{align*}
Hence, with $b:=C_{\mathcal{O}}A_0^2(\log x)^{2\mathfrak{p}_0}$,
\begin{equation*}
\sup|N\mathcal{O}^{(k_0)}(\tilde f_i)|\le x^{-1}\,b\,|S_i^{+}|\,\|f_i\|_\infty .
\end{equation*}

\emph{Sizes of the test functions.} By Definition~\ref{8.5:not-window-sizes},
$\|f_i\|_{\sharp}=\|f_i\|_\infty$ and $|S_i^{+}|\,\|f_i\|_\infty\le80\,\mathfrak{n}(f_i)/\vartheta$. Therefore
\begin{equation*}
\sup|\operatorname{Irr}_i|\le\frac{80\,a}{\vartheta x}\,\mathfrak{n}(f_i),\qquad
\sup|N\mathcal{O}^{(k_0)}(\tilde f_i)|\le\frac{80\,b}{\vartheta x}\,\mathfrak{n}(f_i),
\end{equation*}
and translation by $v$ does not change these suprema.

\emph{The covariances.} By Lemma~\ref{8.5:lem-covariance-sup}, \eqref{8.5:eq-covariance-sup-a}, for
bounded random variables under any probability measure, in particular under $\nu^B_{K,s}$,
$|\operatorname{Cov}(X,Y)|\le2\sup|X|\sup|Y|$. Let $\Sigma$ denote the sum of the absolute values of
the three covariances displayed in the statement. Applying this bound to each of the three
covariances and adding,
\begin{equation*}
\begin{split}
\Sigma
&\le\frac{2\cdot6400}{\vartheta^2x^2}\,(ab+ba+a^2)\,\mathfrak{n}(f_1)\mathfrak{n}(f_2)
=\frac{6400}{\vartheta^2x^2}\,(4ab+2a^2)\,\mathfrak{n}(f_1)\mathfrak{n}(f_2)\\
&=g_{k_0}^4\,\mathfrak{n}(f_1)\mathfrak{n}(f_2)\cdot12800\,\vartheta^{-2}(2ab+a^2)
=g_{k_0}^4\,\mathfrak{n}(f_1)\mathfrak{n}(f_2)\,\mathsf{E}^{\mathrm{irr},\flat}(x),
\end{split}
\end{equation*}
the last equality being \eqref{8.5:eq-exponent-floor-2} written in $a$ and $b$. Each of the three
absolute values is bounded by $\Sigma$.

\emph{Degree.} In the term of $2ab$ that contains $C_{\star}^{\flat}$ the power of $\log x$ has
exponent $p_0-q+4\mathfrak{p}_0$. In the term of $a^2$ that contains $(C_{\star}^{\flat})^2$ it has
exponent $2(p_0-q)+4\mathfrak{p}_0$, which is at most
$p_0-q+4\mathfrak{p}_0$ because $q>p_0$ by \eqref{8.5:eq-exponent-floor-1}. By the choice of $\rho$
in the construction of the two-point witness theorem,
$\rho<LC_W(s)=LC_0A_0^2(\log\bar x_s)^{2\mathfrak{p}_0}$, so $\rho^8$ is bounded by
$(LC_0A_0^2)^8(\log\bar x_s)^{16\mathfrak{p}_0}$. Counting the powers of $\log\bar x_s$ together
with those of $\log x$, the two terms of $\rho^8\mathsf{E}^{\mathrm{irr},\flat}$ just named have
degree at most $p_0-q+20\mathfrak{p}_0$, which is at
most $-1$ by \eqref{8.5:eq-exponent-floor-1}.

\emph{The remainder against the main term.} We run the argument of
Corollary~\ref{8.5:cor-irrelevant-small} with $\mathsf{E}^{\mathrm{irr}}$ replaced by
$\mathsf{E}^{\mathrm{irr},\flat}$. The irrelevant remainder $R^{\mathrm{irr}}$, which in that
argument is bounded by the sum $\Sigma$ of the absolute values of the three covariances, therefore
satisfies
$R^{\mathrm{irr}}\le g_{k_0}^4\mathfrak{n}(f_1)\mathfrak{n}(f_2)\mathsf{E}^{\mathrm{irr},\flat}(x_{k_0})$.
By the lower bound on the main term in the construction of the two-point witness theorem,
$\mathfrak{M}\ge c_Mg_{k_0}^4\mathfrak{n}(f_1)\mathfrak{n}(f_2)\rho^{-8}$, so it suffices that
$\rho^8\mathsf{E}^{\mathrm{irr},\flat}(x_{k_0})\le c_M/32$. The relations $\rho<LC_0A_0^2(\log\bar x_s)^{2\mathfrak{p}_0}$,
$x_{k_0}\ge\underline{x}_s$ and $\bar x_s\le\bar x(\underline{x}_s)$ give
\begin{equation*}
\rho^8\,\mathsf{E}^{\mathrm{irr},\flat}(x_{k_0})
\le\bigl(LC_0A_0^2(\log\bar x(\underline{x}_s))^{2\mathfrak{p}_0}\bigr)^{8}\,
\mathsf{E}^{\mathrm{irr},\flat}_{\downarrow}(\underline{x}_s)\le\frac{c_M}{32}
\end{equation*}
by \eqref{8.5:eq-exponent-floor-3}.

\emph{The floor and the order of choice.}
The floor \textup{(F-$q$)} is a one-sided condition on the exponent $q$, which is chosen at
Stage~0 of the order of choice (Definition~\ref{2.3:not-stages-first}), before
$\kappa,\dots,\gamma$; it reads only the exponents $p_0$ and $\mathfrak{p}_0$, and the correlation
theorem places no upper bound on $q$. The floor creates no cycle in the order of choice, in the
sense of Lemma~\ref{2.4:lem-no-cycle}. The constant
$C_W(s)$, Ba{\l}aban's degree function of the coupling (Definition~\ref{0.2:not-glossary}), $\rho$
and $c_M$ do not depend on $q$. By the
correlation theorem, every quantity depending on $q$ either decreases in $q$ or is a ceiling on
$\gamma$ chosen after $q$. Raising $q$ therefore lowers only $\gamma_J$.
\end{proof}

\begin{corollary}[The array argument from the sup-norm bound]\label{8.5:cor-array-argument}
Let $f_1,f_2$, $\rho$, the boxes $Q_1,Q_2$, the level $k_0$ and the remaining quantities be as
in the array-route theorem, and write $x:=x_{k_0}$. Assume:
\begin{enumerate}
\item[(a)] the array-route theorem, as an implication from its inputs, one of which is its
localisation input for the irrelevant remainder $\operatorname{Irr}_a$, $a=1,2$. This input
places $\operatorname{Irr}_a$, on the level-$k_0$ window, within $C_I\mathfrak{n}(f_a)e^{-c_\eta\rho}$,
with the constant $C_I$ of that input,
of a member of the Gaussian-extraction class $\mathfrak{F}_{k_0}$, and it bounds the truncated
covariance of that member;
\item[(b)] the sup-norm bound on the irrelevant remainder, Theorem~\ref{8.5:thm-remainder-sup},
under its hypotheses;
\item[(c)] the comparison with the branch reading of the remainder,
Lemma~\ref{8.5:lem-branch-comparison};
\item[(d)] the ceiling
\begin{equation}\label{8.5:eq-array-argument-1}
\sup_{y\ge X-c_5}\,80\,C_\star\,y^{-1/2}(\log y)^{p_0}\;\le\;C_{\mathcal{O}}A_0^2 ,
\end{equation}
where $X=g^{-2}$;
\end{enumerate}
together with the ceiling
\eqref{8.5:eq-irrelevant-small-a} of Corollary~\ref{8.5:cor-irrelevant-small}. Then, in the
array-route theorem, the localisation input for the irrelevant remainder may be replaced by (b)
and (c), with the changes (s1)--(s6) made in the proof below. Every other input, condition and
constant of the array-route theorem stands.
\end{corollary}

\begin{proof}
Throughout, $x:=x_{k_0}$. The letters $\varepsilon_A$, $\Gamma$, $R$, $r_A$, $\mathfrak{M}$,
$\sigma_\star$, $e_A$, $\bar e$, $\bar e_2$, $\bar\alpha$, $\hat e$, $\mathfrak{f}$, $C_w$,
$E_1$, $\Pi_1$, $C_{\mathfrak{F}}$, $\eta_A$, $\varepsilon$, $\epsilon_{\mathrm{bg}}$, $C_I$,
$g_c^{(i)}$, $g_c^{\mathrm{loc},(i)}$, $\hat g_c^{(i)}$, $\psi_c$, $a_i$, $\mathsf{f}^\natural$,
the sets $\mathfrak{S}_c$, the variables $V_i$, the clean sets
$\mathrm{Cl}_c$ and the carried quantities $F\mapsto F^{(i)}$, $T^{(i)}$ denote the quantities
so named in the array-route theorem; $X^{(k_0)}(\tilde f)$, $Q_S$, $R_S$, $b_A$ and $\ell$
denote those so named in the Gaussian extraction.

\emph{The places where the localisation input enters.} The array-route theorem uses its
localisation input for the irrelevant remainder at exactly three places:
\begin{itemize}
\item in its positivity step;
\item in the two bounds of its lemma on carried quantities, the first on the carried one-point
functions $\alpha_a^{(i)}$ and the second on the carried truncated covariances $T^{(i)}$;
\item through that lemma, in the quantity $\hat e$ of its comparison proposition, in the quantity
$\bar\alpha$ of its lemma on the clean set, and in its final budget display.
\end{itemize}
We write
\begin{equation}\label{8.5:eq-array-argument-2}
G_a=G_a^X+\operatorname{Irr}_a,\qquad
G_a^X:=N_{k_0}X^{(k_0)}(\tilde f_a)\in\mathfrak{F}_{k_0}\bigl(S_a;C\mathfrak{n}(f_a),\varrho_\eta\bigr),
\end{equation}
with $S_a:=\operatorname{supp}f_a$,
the membership being that of the membership lemma on which the Gaussian extraction rests. If
$\mathcal{O}^{(k_0)}$ is read globally, the defect $\mathfrak{G}_a$ of its localisation joins the
family $\mathfrak{F}'$ of the array-route theorem.

The Gaussian extraction, the construction of the class $\mathfrak{F}_{k_0}$ on which it rests,
the theorem on functions of class members whose part~(A) is used in the comparison proposition
of the array-route theorem,
and the domain-uniform localisation estimate on which Lemma~\ref{8.5:lem-branch-comparison}
rests are all applied to $G_a^X$ only. The remainder $\operatorname{Irr}_a$ is never expanded and
never localised, and it is never required to belong to any class. We now go through the three
places in the order of the array-route theorem.

\emph{(s1) The positivity step.} By the definitions of the two sizes on the window, Notation~\ref{8.5:not-window-sizes} and \eqref{8.5:eq-window-sizes-a},
using $\|f_a\|_\sharp=\|f_a\|_\infty$ (valid since $\rho\ge40M/\vartheta$) and
$\varepsilon_{k_0}=g_{k_0}A_0(\log x)^{\mathfrak{p}_0}$ with $g_{k_0}^2=x^{-1}$, one has
\begin{equation}\label{8.5:eq-array-argument-3}
\frac{\mathsf{I}_a}{\mathsf{O}_a}
=\frac{C_\star\,x^{-1/2}(\log x)^{p_0}}{C_{\mathcal{O}}A_0^2}.
\end{equation}
The ceiling \eqref{8.5:eq-array-argument-1} of hypothesis~(d),
read at $y=x$, gives $\mathsf{I}_a\le\mathsf{O}_a/80$ by \eqref{8.5:eq-array-argument-3}.
Notation~\ref{8.5:not-window-sizes} gives $\sup|N_{k_0}\mathcal{O}^{(k_0)}(\tilde f_a)|\le\mathsf{O}_a$
for the main insertion; when $\mathcal{O}^{(k_0)}$ is read locally,
$X^{(k_0)}(\tilde f_a)=\mathcal{O}^{(k_0)}(\tilde f_a)$, so this is $\sup|G_a^X|\le\mathsf{O}_a$.
Theorem~\ref{8.5:thm-remainder-sup}, in the form \eqref{8.5:eq-remainder-sup-c}, gives
$\sup|\operatorname{Irr}_a|\le\mathsf{I}_a$. Notation~\ref{8.5:not-window-sizes} also gives
$|S_a^+|\,\|f_a\|_\infty\le 80\,\mathfrak{n}(f_a)/\vartheta$. Together these yield
\begin{equation}\label{8.5:eq-array-argument-4}
|G_a|\;\le\;\mathsf{O}_a+\mathsf{I}_a\;\le\;\tfrac{81}{80}\,\mathsf{O}_a
\;\le\;81\,\vartheta^{-1}C_{\mathcal{O}}\,\mathfrak{n}(f_a)\,\varepsilon_{k_0}^2 .
\end{equation}
Hence the sup-norm bound on $G_a$ used in the positivity step of the array-route theorem holds
with the constant
\begin{equation*}
C_G:=81\,C_{\mathcal{O}}/\vartheta .
\end{equation*}
The quantities $\varepsilon_A$, $\Gamma$ and $R$ are then those of the array-route theorem, read
at this value of $C_G$.

In the display of the positivity step, the quantity $\mathfrak{f}$ is replaced by
$\mathfrak{f}+\mathfrak{f}^{\mathrm{irr}}_a$, with $\mathfrak{f}^{\mathrm{irr}}_a$ as in
\eqref{8.5:eq-branch-comparison-a}. Lemma~\ref{8.5:lem-branch-comparison} holds for all values of
the variables $\mathsf{V}$, which is what that display requires. The lemma rests on the
domain-uniform localisation estimate, and this estimate is already among the inputs of the lemma
on carried quantities of the array-route theorem.

Two further readings are made in the array-route theorem.
\begin{itemize}
\item The definition ``$h$ is $k_0$-good if no live label of $h$ is near
$S_1^c\cup S_2^c$'' is read as ``$h$ is $k_0$-good if no live label of $h$ meets
$Q_1^c\cup Q_2^c$''.
\item The box $Q_S\supset S_1^+\cup S_2^+$ of the Gaussian extraction, of size $R_S\le12\rho$,
is taken to contain $Q_1\cup Q_2$. Since $Q_a$ exceeds $B(x_a,\rho)$ by $\rho/4$ on each side,
the bound $R_S\le12\rho$ becomes $R_S\le\tfrac{25}{2}\rho$ at most. The constants of the
Gaussian extraction depend polynomially on $R_S$, so the quantities $b_A$ and $\ell$ change only
by constants.
\end{itemize}
With these readings, $\mathrm{Cl}_c$ is contained in the set of $(h,V)$ for which no live label of
$h$ meets $Q_a^c$ and $V$ is in the window on $Q_a^c$.

\emph{(s2) The first bound of the lemma on carried quantities.} The carried map
$F\mapsto F^{(i)}$ is linear, so by \eqref{8.5:eq-array-argument-2}
\begin{equation*}
\alpha_a^{(i)}=(G_a^X)^{(i)}+(\operatorname{Irr}_a)^{(i)} .
\end{equation*}
The one-step functionals $\mathbb{E}_{i,\omega}$ of the array-route theorem are positive and
normalised, as its lemma on these functionals states. Hence, by fact~(e2) of
Lemma~\ref{8.5:lem-covariance-sup}, $(\operatorname{Irr}_a)^{(i)}$ is, for each argument, the
integral of $\operatorname{Irr}_a$ against a probability measure. This gives
$|(\operatorname{Irr}_a)^{(i)}|\le\sup|\operatorname{Irr}_a|\le\mathsf{I}_a$.

Therefore the first bound holds with $\alpha^{\mathrm{loc},(i)}$ and $\alpha^\flat$ taken to be
the explicit members of $G_a^X$ supplied by the construction of the class $\mathfrak{F}_{k_0}$,
and with the replacements
\begin{equation}\label{8.5:eq-array-argument-5}
\bar e\;\mapsto\;\bar e^X+\mathsf{I}_a,\qquad \bar\alpha\;\mapsto\;\bar\alpha^X+\mathsf{I}_a .
\end{equation}
Here $\bar e^X$ and $\bar\alpha^X$ are the quantities $\bar e$ and $\bar\alpha$ of the
array-route theorem with the member $C_Ie^{-c_\eta\rho}x$ omitted.

We next rewrite the new member. Notation~\ref{8.5:not-window-sizes} gives
\begin{gather*}
\mathsf{I}_a=C_\star x^{-3/2}(\log x)^{p_0+2\mathfrak{p}_0}|S_a^+|\,\|f_a\|_\sharp,\\
\|f_a\|_\sharp=\|f_a\|_\infty,\qquad |S_a^+|\,\|f_a\|_\infty\le80\,\mathfrak{n}(f_a)/\vartheta .
\end{gather*}
Consequently
\begin{equation}\label{8.5:eq-array-argument-6}
\mathsf{I}_a\;\le\;C_w\,\mathfrak{n}(f_a)\,x^{-1}E_1^{\mathrm{irr}}(x),\qquad
E_1^{\mathrm{irr}}(x):=\frac{80\,C_\star}{\vartheta\,C_w}\,x^{-1/2}(\log x)^{p_0+2\mathfrak{p}_0}.
\end{equation}
Since $x^{-1/2}=g_{k_0}$, the quantity $E_1^{\mathrm{irr}}$ is $g_{k_0}$ times a power of
$\log x$. This is the form of the first member $g_{k_0}\Pi_1$ of the quantity $E_1$ of the
array-route theorem.

\emph{(s3) The second bound of the lemma on carried quantities.} The truncated covariance
$T^{(i)}[\cdot,\cdot]$ is bilinear. By \eqref{8.5:eq-array-argument-2},
\begin{equation}\label{8.5:eq-array-argument-7}
\begin{split}
T^{(i)}[G_1,G_2]={}&T^{(i)}[G_1^X,G_2^X]+T^{(i)}[G_1^X,\operatorname{Irr}_2]\\
&+T^{(i)}[\operatorname{Irr}_1,G_2^X]+T^{(i)}[\operatorname{Irr}_1,\operatorname{Irr}_2].
\end{split}
\end{equation}
By fact~(e2) of Lemma~\ref{8.5:lem-covariance-sup}, \eqref{8.5:eq-covariance-sup-a}, each of the
last three terms is a covariance. Its modulus is therefore at most twice the product of the sup
norms of its two entries. The sup norms are at most $\mathsf{O}_a$ for $G_a^X$ and at most
$\mathsf{I}_a$ for $\operatorname{Irr}_a$. Hence the last three terms of
\eqref{8.5:eq-array-argument-7} are together at most
\begin{equation}\label{8.5:eq-array-argument-8}
2\mathsf{O}_1\mathsf{I}_2+2\mathsf{I}_1\mathsf{O}_2+2\mathsf{I}_1\mathsf{I}_2
\;\le\;x^{-2}\,\mathfrak{n}(f_1)\mathfrak{n}(f_2)\,\mathsf{E}^{\mathrm{irr}}(x),
\end{equation}
with $\mathsf{E}^{\mathrm{irr}}$ as in \eqref{8.5:eq-second-order-extraction-a}. The inequality
in \eqref{8.5:eq-array-argument-8} is the estimate of the same three products made in the proof
of Theorem~\ref{8.5:thm-second-order-extraction}, with $g_{k_0}^4=x^{-2}$. It holds pointwise and
on every branch, whether clean or not.

Therefore the second bound holds with $T^{\mathrm{loc}}$ taken to be the local part of
$T^{(i)}[G_1^X,G_2^X]$, and with the replacement
\begin{equation}\label{8.5:eq-array-argument-9}
\bar e_2\;\mapsto\;\bar e_2^X+x^{-2}\,\mathfrak{n}(f_1)\mathfrak{n}(f_2)\,
\mathsf{E}^{\mathrm{irr}}(x),
\end{equation}
where $\bar e_2^X$ is the quantity $\bar e_2$ of the array-route theorem with the member that
came from the localisation input omitted.

The array-route theorem proves $|T^{(r_A)}(W)-\mathfrak{M}|\le\tfrac{3}{32}\mathfrak{M}$ as a sum
of three contributions $\mathfrak{M}/32$. The third of these bounds the terms that contain the
irrelevant remainder, and the array-route theorem obtains it under the floor
$C_0^2\ge448\,C_b$. This third contribution is now supplied by
Corollary~\ref{8.5:cor-irrelevant-small} under the ceiling \eqref{8.5:eq-irrelevant-small-a}.
The floor $C_0^2\ge448\,C_b$ is not needed.

\emph{(s4) The comparison proposition.} By (s1), $g_c^{(i)}\in[\tfrac12,2]$, since $g_c^{(i)}$ is
a positive normalised average of $\mathsf{f}^\natural\in[\tfrac12,2]$. We put
\begin{equation*}
\hat g_c^{(i)}:=g_c^{\mathrm{loc},(i)}+\hat e,\qquad \hat e:=3R\bar e+R^2\bar e_2 .
\end{equation*}
Here $g_c^{\mathrm{loc},(i)}$ is built on $\alpha^{\mathrm{loc}}$ and $T^{\mathrm{loc}}$ of the
$G^X$ alone, and $\bar e$, $\bar e_2$ are the new quantities of \eqref{8.5:eq-array-argument-5}
and \eqref{8.5:eq-array-argument-9}. With these, $g_c^{(i)}\le\hat g_c^{(i)}\le g_c^{(i)}+2\hat e$,
exactly as in the comparison proposition.

The function $\psi_c=\hat g_c^{(i)}+a_i$ contains no term of the irrelevant remainder. It is a
function of the restriction $V_i|_{\mathfrak{S}_c}$, and it depends on it only through members of
the class $\mathfrak{F}_{k_0}$. Consequently the following apply without change, with
$w_a=C\mathfrak{n}(f_a)$ as in the membership lemma:
\begin{itemize}
\item the measurability lemma of the array-route theorem;
\item part~(A) of the theorem on functions of class members named above;
\item the bound
\begin{equation*}
\eta_A\le C_{\mathfrak{F}}\,R\,(w_1+w_2)\,2\varepsilon\,\epsilon_{\mathrm{bg}} .
\end{equation*}
\end{itemize}

\emph{(s5) The lemma on the clean set.} That lemma states that on $\mathrm{Cl}_c$ the carried
integrand is at most $\bar\alpha+\mathfrak{f}$ in modulus, and it takes this from the first bound
of the lemma on carried quantities. It is read at the new $\bar\alpha$ of
\eqref{8.5:eq-array-argument-5} and at the new $\mathfrak{f}$ of (s1), that is, at
$\mathfrak{f}+\mathfrak{f}^{\mathrm{irr}}_a$.

\emph{(s6) The final budget display.} From $R^2x^{-2}=\varepsilon_A^2$ we have $R=\varepsilon_Ax$.
The first budget inequality of the array-route theorem gains the member $56R\mathsf{I}_a$. By
\eqref{8.5:eq-array-argument-6},
\begin{equation}\label{8.5:eq-array-argument-10}
56R\,\mathsf{I}_a\;\le\;56\,\varepsilon_A\cdot80\,\vartheta^{-1}C_\star\,\mathfrak{n}\,
x^{-1/2}(\log x)^{p_0+2\mathfrak{p}_0}.
\end{equation}
The second budget inequality gains
\begin{equation}\label{8.5:eq-array-argument-11}
r_A\bigl(12R\,\mathsf{I}_a+4\varepsilon_A^2\,\mathfrak{n}(f_1)\mathfrak{n}(f_2)\,
\mathsf{E}^{\mathrm{irr}}(x)\bigr).
\end{equation}
Indeed, $2e_A$ contains the member $4r_A\hat e$, with $\hat e=3R\bar e+R^2\bar e_2$. Inserting
the new members $\mathsf{I}_a$ of $\bar e$ and
$x^{-2}\mathfrak{n}(f_1)\mathfrak{n}(f_2)\mathsf{E}^{\mathrm{irr}}(x)$ of $\bar e_2$, and using
$R^2x^{-2}=\varepsilon_A^2$, gives \eqref{8.5:eq-array-argument-11}.

Both budget inequalities are read with these members included, and each is still required to be at
most $\sigma_\star/8$, where $\sigma_\star=\tfrac{3}{64}c_M\vartheta^2\varepsilon_A^2$. We have
$\mathfrak{n}\le10\rho^4$, $\varepsilon_A^{-1}=170\,C_G\rho^4p_0(g_{k_0})^2$, with the degree
function $p_0(\cdot)$, and
$r_A=O(\log\bar x_s)$. With these, each of the two requirements becomes the condition that the
supremum over $y$ of $y^{-1/2}$ times a polylogarithm of $y$, with coefficients fixed by the
constants of the array-route theorem, be at most a constant fixed by the same constants.
This is the final ceiling of the array-route theorem, re-evaluated with
the new members. Its form is unchanged, and it is chosen last. No further input on the
irrelevant remainder enters.

\emph{Conclusion.} (s1)--(s6) exhaust the places at which the array-route theorem uses its
localisation input for the irrelevant remainder. At each of them the input is served by
Theorem~\ref{8.5:thm-remainder-sup} and Lemma~\ref{8.5:lem-branch-comparison}, at the cost of
the ceiling \eqref{8.5:eq-array-argument-1}, the ceiling \eqref{8.5:eq-irrelevant-small-a} and the
re-evaluated final ceiling. All other inputs, conditions and constants are used as they stand.

Lemma~\ref{8.5:lem-branch-comparison} rests on the domain-uniform localisation estimate, which is
already an input of the array-route theorem. Hence the conclusion of the array-route theorem,
the two-sided form of clause~(iv) in every volume by the array route, follows from its inputs
other than the localisation input for the irrelevant remainder, each with the standing it has
there, together with Theorem~\ref{8.5:thm-remainder-sup} and
Lemma~\ref{8.5:lem-branch-comparison}.

The same substitution applies to two by-products derived in the array-route theorem: the bound
$|S^T_2|\le C\rho^8p_0(g_{K-s})^4g^4_{K-s}$ in every volume, with the degree function
$p_0(\cdot)$, and the second proof of clause~(iv) on a fixed torus that it gives.
\end{proof}

\begin{corollary}[Decay of the irrelevant term at polylogarithmic room]
\label{8.5:cor-irrelevant-decay}
Work in the prefix and under the hypotheses of the second-order extraction estimate
(Theorem~\ref{8.5:thm-second-order-extraction}), whose notation $k_0$, $x_{k_0}$,
$\nu^B_{K,s}$, $\mathfrak{n}(f_i)$, $\operatorname{Irr}_i$,
$N\mathcal{O}^{(k_0)}(\tilde f_i)$, $C_{\star}$, $C_{\mathcal{O}}$, $\varepsilon_{k_0}$ and
$\mathsf{E}^{\mathrm{irr}}$ of~\eqref{8.5:eq-second-order-extraction-a} we use. Put
$\mathsf{E}^{\mathrm{irr}}_{\downarrow}(y):=\sup_{x\ge y}\mathsf{E}^{\mathrm{irr}}(x)$.
Let $R^{\mathrm{irr}}$ be the irrelevant remainder of the two-point witness theorem, and let
$\mathfrak{M}$, $c_M$ and $\underline{x}_s\le\bar x_s$ be its main term, the constant of the
lower bound of the main term, and its envelope, so that $x_{k_0}\ge\underline{x}_s$ and, with the
coupling $g_{k_0}$ of~\eqref{8.5:eq-second-order-extraction-a}, in the form used in
Corollary~\ref{8.5:cor-irrelevant-small},
\begin{equation}\label{8.5:eq-irrelevant-decay-1}
\mathfrak{M}\ge c_M\,g_{k_0}^4\,\mathfrak{n}(f_1)\mathfrak{n}(f_2)\,\rho^{-8}.
\end{equation}
We use from the two-point witness theorem that its envelope satisfies
$\bar x_s\le\bar x(\underline{x}_s)$ for a function $\bar x(\cdot)$ such that
$\log\bar x(y)$ is bounded by a power of $\log y$, and that $\underline{x}_s\to\infty$ as
$s\to\infty$ at fixed $g$ and as $g\to0$. Here $\rho$ is the room of the two-point witness
theorem, the parameter whose power $\rho^{-8}$ enters~\eqref{8.5:eq-irrelevant-decay-1}, and
$C_r$ is the constant of the setting of Theorem~\ref{8.5:thm-second-order-extraction} with which
the hypothesis \eqref{8.5:eq-regular-extension-a} of the regular extension lemma is read there.
Let $\rho\ge 40M/\vartheta$ satisfy that hypothesis in the form
\begin{equation}\label{8.5:eq-irrelevant-decay-2}
d^2\,\tfrac{25}{4}\,\rho^2C_rL^2\,\varepsilon_{k_0}\le 1 .
\end{equation}
\begin{itemize}
\item[(a)] The conclusions of Theorem~\ref{8.5:thm-second-order-extraction} hold for this
$\rho$ without change, with constants $C_{\star}$ and $C_{\mathcal{O}}$ that depend neither on
$\rho$ nor on any box; in particular
\begin{equation}\label{8.5:eq-irrelevant-decay-4}
R^{\mathrm{irr}}\le g_{k_0}^4\,\mathfrak{n}(f_1)\mathfrak{n}(f_2)\,\mathsf{E}^{\mathrm{irr}}(x_{k_0}).
\end{equation}
Thus the input of the sharp asymptotic-freedom law holds with vanishing coefficients of the
corrections proportional to the main kernel, and with the right side
of~\eqref{8.5:eq-irrelevant-decay-4} as its single error member, which carries one power of
$g_{k_0}$ beyond $g_{k_0}^4$.
\item[(b)] If moreover $\rho\le(\log\bar x_s)^{c_\rho}$ for some fixed $c_\rho>0$
(arbitrary), then
\begin{equation}\label{8.5:eq-irrelevant-decay-3}
\frac{R^{\mathrm{irr}}}{\mathfrak{M}}
\le\frac{\rho^8\,\mathsf{E}^{\mathrm{irr}}_{\downarrow}(\underline{x}_s)}{c_M}
\le\frac{(\log\bar x_s)^{8c_\rho}\,\mathsf{E}^{\mathrm{irr}}_{\downarrow}(\underline{x}_s)}{c_M},
\end{equation}
and the right side tends to $0$ as $s\to\infty$ at fixed $g$, and as $g\to0$. The same holds
for every $\rho<L\,C_W(s)$, with $(\log\bar x_s)^{8c_\rho}$ replaced by
$(LC_0A_0^2)^8(\log\bar x_s)^{16\mathfrak{p}_0}$.
\end{itemize}
\end{corollary}

\begin{proof}
(a) In Theorem~\ref{8.5:thm-second-order-extraction} the constants $C_{\star}$ and
$C_{\mathcal{O}}$ depend neither on $\rho$ nor on any box, and $\rho$ enters only through the
regular extension lemma (Lemma~\ref{8.5:lem-regular-extension}), whose hypothesis
\eqref{8.5:eq-regular-extension-a} reads \eqref{8.5:eq-irrelevant-decay-2} for the boxes and
plaquette bound used in Theorem~\ref{8.5:thm-second-order-extraction}. Hence the theorem holds
verbatim for every $\rho\ge40M/\vartheta$ satisfying \eqref{8.5:eq-irrelevant-decay-2}, at every
translate $v$, and its summed bound, by which $R^{\mathrm{irr}}$ is bounded as in
Corollary~\ref{8.5:cor-irrelevant-small}, gives \eqref{8.5:eq-irrelevant-decay-4}. This bound
contains no term proportional to the main kernel, so the coefficients of such corrections
vanish. Since $x=g^{-2}$ gives $x^{-1/2}=g$ and $x^{-1}=g^2$,
the two members of $\mathsf{E}^{\mathrm{irr}}(x_{k_0})$ in
\eqref{8.5:eq-second-order-extraction-a} are $g_{k_0}$, respectively $g_{k_0}^2$, times powers
of $\log x_{k_0}$; so the error term carries $g_{k_0}^1$. Since
$\varepsilon_{k_0}=g_{k_0}A_0(\log x_{k_0})^{\mathfrak{p}_0}$ and $g_{k_0}=x_{k_0}^{-1/2}$, and
since $\rho\le(\log\bar x_s)^{c_\rho}$ gives $\rho^2\le(\log\bar x_s)^{2c_\rho}$, on that range
condition \eqref{8.5:eq-irrelevant-decay-2} follows from
\begin{equation*}
\tfrac{25}{4}\,d^2C_rL^2A_0\,(\log\bar x_s)^{2c_\rho}(\log x_{k_0})^{\mathfrak{p}_0}\,
x_{k_0}^{-1/2}\le1 ,
\end{equation*}
which is a ceiling on the coupling in the sense of Definition~\ref{2.3:not-stages-first}: it
sets $x_{k_0}^{-1/2}$ against a product of logarithms of total degree
$2c_\rho+\mathfrak{p}_0$.

(b) Dividing \eqref{8.5:eq-irrelevant-decay-4} by \eqref{8.5:eq-irrelevant-decay-1} gives
$R^{\mathrm{irr}}/\mathfrak{M}\le\rho^8\mathsf{E}^{\mathrm{irr}}(x_{k_0})/c_M$. As
$x_{k_0}\ge\underline{x}_s$, the definition of $\mathsf{E}^{\mathrm{irr}}_{\downarrow}$ gives
$\mathsf{E}^{\mathrm{irr}}(x_{k_0})\le\mathsf{E}^{\mathrm{irr}}_{\downarrow}(\underline{x}_s)$,
which is the first inequality of \eqref{8.5:eq-irrelevant-decay-3}; the second is
$\rho^8\le(\log\bar x_s)^{8c_\rho}$. By \eqref{8.5:eq-second-order-extraction-a},
$\mathsf{E}^{\mathrm{irr}}(x)$ is $x^{-1/2}$ times a power of $\log x$ plus $x^{-1}$ times a
power of $\log x$, so $\mathsf{E}^{\mathrm{irr}}_{\downarrow}(y)$ is bounded by $y^{-1/2}$ times a
power of $\log y$ for large $y$. With $\bar x_s\le\bar x(\underline{x}_s)$ and $\log\bar x(y)$
bounded by a power of $\log y$, the right side of \eqref{8.5:eq-irrelevant-decay-3} is bounded by
$y^{-1/2}$ times a power of $\log y$ at $y=\underline{x}_s$, which tends to $0$ as $y\to\infty$.
Since $\underline{x}_s\to\infty$ as $s\to\infty$ at fixed $g$ and as $g\to0$, the right side
tends to $0$ in both limits. For $\rho<LC_W(s)$, with
$LC_W(s)=LC_0A_0^2(\log\bar x_s)^{2\mathfrak{p}_0}$ as in
Corollary~\ref{8.5:cor-irrelevant-small}, one has
$\rho^8<(LC_0A_0^2)^8(\log\bar x_s)^{16\mathfrak{p}_0}$, and the same argument applies, the
factor $(LC_0A_0^2)^8$ not affecting the limit.
\end{proof}

\begin{remark}[Model: the stated bounds alone do not suffice]\label{8.5:rem-moduli-insufficient}
In the setting of Theorem~\ref{8.5:thm-second-order-extraction}, let $K$, $s$, the test functions $f_1,f_2$ and the functions $\tilde f_1,\tilde f_2$ built from them be as admitted there. Let $\mathsf{r}$ be the source radius of the correlation theorem, whose frozen families are taken at constant sources $\hat w\in D(0,2\mathsf{r})$. Let $q$ be the auxiliary parameter of the shrinking field radii, and put
\begin{equation*}
\Lambda(x):=\frac{x}{\mathsf{r}\,(\log x)^{2q}} .
\end{equation*}
No argument can prove part~(b) of the Gaussian extraction, nor therefore the stronger estimate of Theorem~\ref{8.5:thm-second-order-extraction}, which implies part~(b), if it uses, of the families entering the irrelevant remainder~\eqref{8.5:eq-remainder-level-sum-a}, only the following properties of the correlation theorem:
\begin{itemize}
\item the four clauses of its sourced format (the definition of the running shift $\zeta_k$; the extracted families built on the frozen families; the local families, with their analyticity, locality and freezing at constants; the bounds on the remainder families);
\item the four clauses of its inductive hypothesis at the levels $k\le k_0$ (the comparison of the running inverse couplings; the bound on $b_j(\zeta)-b_j(0)$; the bounds on the pure, $Q$- and frozen families; the bound on the remainder families);
\item the running-shift statement (the recursion for the coefficients $b_{k+1}$);
\item its memory-sum bound;
\item its statement on real values.
\end{itemize}
The reason is that there is a second choice of the level-$k_0$ local families with the following three features.
\begin{itemize}
\item It satisfies all of these properties, with the a~priori bound $E_0$ on the frozen and pure families in the inductive hypothesis of the correlation theorem replaced by $2E_0$.
\item Its irrelevant remainder differs from the true one by a term whose covariance with the main insertion $N\mathcal{O}^{(k_0)}(\tilde f_2)$, computed at the leading order of part~(a) of the Gaussian extraction, has modulus at least a positive constant times $\Lambda(x_{k_0})\rho^{-2}\mathfrak{M}$.
\item The factor $\Lambda(x_{k_0})\rho^{-2}$ is unbounded as $x_{k_0}\to\infty$.
\end{itemize}
Here $\mathfrak{M}:=N^2b_0g^4\mathcal{K}_{K,s}(f_1,f_2)$ is the main term of the two-point witness theorem, and $\rho=L^s\mathsf{d}$, the separation of the supports in cells of $\mathbb{T}^{(k_0)}$.
\end{remark}

\begin{proof}
The model is the family obtained from the level-$k_0$ local families of the correlation theorem by adding, at every point $y$ of the block lattice $\mathbb{T}^{(k_0)}$ of Definition~\ref{1.5:def-block-average}, a multiple, linear in the value of the source at the evaluation plaquette $p_y$, of a discrete second difference of weighted Wilson actions.

\emph{The added functionals.} Let $\alpha_{*,k_0}$ be the minimal field radius at level $k_0$, and let $\mathfrak{a}_{k_0}=C_{\mathfrak{a}}\alpha_{*,k_0}$ be the regularity size at level $k_0$, with $C_{\mathfrak{a}}$ its fixed constant; put
\begin{equation*}
c_{\mathcal{Z}}:=(64d\cdot 2^d)^{-1},\qquad \mathfrak{a}:=\mathfrak{a}_{k_0}=C_{\mathfrak{a}}\alpha_{*,k_0}.
\end{equation*}
For $y\in\mathbb{T}^{(k_0)}$, with $e_\lambda$ the unit vectors of $\mathbb{T}^{(k_0)}$ and $h_y$ the profiles, define
\begin{equation}\label{8.5:eq-moduli-insufficient-1}
\mathcal{Z}_y(U):=c_{\mathcal{Z}}\,\mathfrak{a}^{-2}\sum_{\lambda=1}^{d}
\bigl[2A(h_y,U)-A(h_{y+e_\lambda},U)-A(h_{y-e_\lambda},U)\bigr].
\end{equation}
We write $\mathcal{Z}$ for the pointed family $y\mapsto\mathcal{Z}_y$. It has the following properties.
\begin{itemize}
\item $\mathcal{Z}_y$ is entire and gauge invariant, since each weighted action $A(h,\cdot)$ is.
\item $\mathcal{Z}_y(1)=0$, since $A(h,1)=0$ for every weight $h$.
\item $\mathcal{Z}$ is pointed-covariant: $\mathcal{Z}_{\iota y}(\iota U)=\mathcal{Z}_y(U)$ for every lattice symmetry $\iota$ of $\mathbb{T}^{(k_0)}$.
\item On the complex background space $\mathcal{U}^{\flat}_{k_0}$ at the shrinking radii, the bound $|A(h_{y'},U)|\le 4\cdot2^d\mathfrak{a}^2$ of the correlation theorem on weighted actions holds for each of $y'=y,\,y\pm e_\lambda$. Hence
\begin{equation*}
|\mathcal{Z}_y(U)|\le c_{\mathcal{Z}}\mathfrak{a}^{-2}\cdot d\cdot(2+1+1)\cdot4\cdot2^d\mathfrak{a}^2
=\frac{16\,d\,2^d}{64\,d\,2^d}=\frac14\le1 .
\end{equation*}
\item Let $\beta[\,\cdot\,]$ denote the total coefficient of the weighted actions $A(h_{y'},\cdot)$, $y'\in\mathbb{T}^{(k_0)}$, that the renormalisation prescription \cite[(3.61)]{B-CMP119} extracts from a pointed family (the coefficient entering $b_j$ through the running-shift statement). In each direction $\lambda$ the three weighted actions in~\eqref{8.5:eq-moduli-insufficient-1} carry the coefficients $2,-1,-1$, so
\begin{equation*}
\beta[\mathcal{Z}_y]=c_{\mathcal{Z}}\mathfrak{a}^{-2}\sum_{\lambda=1}^d(2-1-1)=0 .
\end{equation*}
\item On the torus the substitution $y\mapsto y\mp e_\lambda$ gives $\sum_yA(h_{y\pm e_\lambda},U)=\sum_yA(h_y,U)$. Therefore
\begin{equation}\label{8.5:eq-moduli-insufficient-2}
\sum_{y\in\mathbb{T}^{(k_0)}}\mathcal{Z}_y\equiv0 .
\end{equation}
\end{itemize}

\emph{The model family.} Let $X_y$ be the $2^d$-cube hull of $y$, and put
\begin{equation*}
c_1:=\frac{E_0e^{-2^d\kappa}}{2\mathsf{r}} .
\end{equation*}
At level $k_0$ define
\begin{equation*}
\tilde E^{(k_0)}(X_y,\cdot,y;w):=E^{(k_0)}(X_y,\cdot,y;w)+c_1\,w(p_y)\,\mathcal{Z}_y .
\end{equation*}
All other local families are left unchanged. For $w$ with $|w(p_y)|\le2\mathsf{r}$ and $U\in\mathcal{U}^{\flat}_{k_0}$, the added term is bounded by $c_1\cdot2\mathsf{r}\cdot1=E_0e^{-2^d\kappa}$.

The family $\tilde E$ obeys the third clause of the sourced format and the second and third clauses of the inductive hypothesis, with $E_0$ replaced by $2E_0$. We go through the requirements one at a time.
\begin{itemize}
\item It is analytic on $\mathcal{W}_{k_0}(X_y)$, because $\mathcal{Z}_y$ is entire and $w(p_y)$ is linear in $w$.
\item It is local, because the added term reads $w$ only at $p_y$.
\item If $w\equiv\hat w$ on the enlarged domain $X_y^{+}$ of $X_y$ on which the local family reads the source, it freezes at constants to
\begin{equation*}
\tilde E_{\hat w}=E_{\hat w}+c_1\hat w\,\mathcal{Z} ,
\end{equation*}
which is again a pointed $E$-family.
\item The coefficients $b_j$ are unchanged, since $\beta[\mathcal{Z}_y]=0$.
\item By~\eqref{8.5:eq-moduli-insufficient-2} the totals are unchanged. Hence $\tilde P^{\flat(k_0)}_{\zeta}:=P^{\flat(k_0)}_{\zeta}+c_1\zeta\mathcal{Z}$ still corresponds to the pure family $\mathsf{P}^{(k_0)}_{\xi}$.
\item The $Q$-family is unchanged.
\item The difference family is unchanged:
\begin{align*}
\tilde D^{(k_0)}
&=E^{(k_0)}(\cdot;w)+c_1w(p_y)\mathcal{Z}_y-E^{(k_0)}_{w(p_y)}-c_1w(p_y)\mathcal{Z}_y\\
&=D^{(k_0)} .
\end{align*}
\end{itemize}
The remaining listed properties concern only the couplings, the coefficients $b_j$, the $Q$-family, the remainder families and the difference family, none of which has changed, and the frozen families, which are again a pointed $E$-family; hence the model has all of them.

\emph{The two remainders.} In~\eqref{8.5:eq-remainder-level-sum-a} only the level-$k_0$ term $\dot E^{(k_0)}[f]$ changes. Put $w=zf$. Then $\partial_z\bigl(c_1zf(p_y)\mathcal{Z}_y\bigr)=c_1f(p_y)\mathcal{Z}_y$. The vacuum subtraction does not affect this term, since $\mathcal{Z}_y(1)=0$. The terms $b_j'(0)A(I_jf,U_{k_0})$ and $\dot R^{(j)}[f]$ are unchanged.

On $\mathbb{T}^{(k_0)}$ put
\begin{equation*}
(-\Delta f)(p_y):=\sum_{\lambda=1}^d\bigl[2f(p_y)-f(p_{y+e_\lambda})-f(p_{y-e_\lambda})\bigr].
\end{equation*}
Substituting $y\mapsto y\mp e_\lambda$ in the two shifted sums (summation by parts) and using the linearity of $A(\cdot,U)$ in the weight, together with $I_{k_0}g=\sum_yg(p_y)h_y$, we obtain
\begin{align}
\widetilde{\operatorname{Irr}}[f]-\operatorname{Irr}[f]
&=c_1\sum_yf(p_y)\mathcal{Z}_y(U_{k_0})
=c_1c_{\mathcal{Z}}\mathfrak{a}^{-2}\sum_y(-\Delta f)(p_y)A(h_y,U_{k_0})\notag\\
&=c_1c_{\mathcal{Z}}\mathfrak{a}^{-2}\,A\bigl(I_{k_0}(-\Delta f),U_{k_0}\bigr)
=c_1c_{\mathcal{Z}}\mathfrak{a}^{-2}\,\mathcal{O}^{(k_0)}\bigl(I_{k_0}(-\Delta f)\bigr).
\label{8.5:eq-moduli-insufficient-3}
\end{align}
Here $\mathcal{O}^{(k_0)}(h):=A(h,U_{k_0})$ is the curvature field of Definition~\ref{1.3:def-curvature-field} with weight $h$, read on the block field $U_{k_0}$.

\emph{The covariance.} Let $\mathsf{k}(x,y)$ denote the kernel of the leading order of part~(a) of the Gaussian extraction; at that order it is a constant multiple of the power $|x-y|^{-8}$. At that order the covariance of~\eqref{8.5:eq-moduli-insufficient-3}, taken at $f=f_1$, with $N\mathcal{O}^{(k_0)}(\tilde f_2)$ is
\begin{equation*}
c_1c_{\mathcal{Z}}\mathfrak{a}^{-2}\,N b_0 g^4\sum_x\sum_y f_1(x)f_2(y)\,(-\Delta_x\mathsf{k})(x,y),
\end{equation*}
the second difference of~\eqref{8.5:eq-moduli-insufficient-3} having been carried from $f_1$ onto the first argument of the kernel by the same substitution $y\mapsto y\mp e_\lambda$ as above. The single factor $N$ is that of the one insertion $N\mathcal{O}^{(k_0)}(\tilde f_2)$; the difference~\eqref{8.5:eq-moduli-insufficient-3} carries $\mathcal{O}^{(k_0)}$, not $N\mathcal{O}^{(k_0)}$.
In $\mathbb{R}^4$ one has $\Delta|x|^{-a}=a(a-2)|x|^{-a-2}$. Hence
\begin{equation*}
-\Delta|x|^{-8}=-48|x|^{-10}=-48\,|x|^{-2}\,|x|^{-8}\neq0,
\end{equation*}
so the Laplacian does not annihilate the power $|x|^{-8}$. On the supports of $f_1$ and $f_2$ one has therefore $-\Delta_x\mathsf{k}(x,y)=-48\,|x-y|^{-2}\mathsf{k}(x,y)$, and $|x-y|$ is there at most a constant multiple of $\rho$, so that $|x-y|^{-2}$ is at least a positive constant times $\rho^{-2}$. Moreover $c_1c_{\mathcal{Z}}=E_0e^{-2^d\kappa}/(128\,d\,2^d\mathsf{r})$ is a constant multiple of $E_0/\mathsf{r}$. Since $\alpha_{*,k_0}$ is a constant multiple of $g_{k_0}(\log x_{k_0})^{q}$, the factor $\mathfrak{a}^{-2}$ is a constant multiple of $x_{k_0}(\log x_{k_0})^{-2q}$. The modulus of the covariance is therefore at least a positive constant, depending on $N$, times
\begin{equation*}
\frac{E_0}{\mathsf{r}}\,\mathfrak{a}^{-2}\rho^{-2}\,\mathfrak{M}\asymp\Lambda(x_{k_0})\,\rho^{-2}\,\mathfrak{M}.
\end{equation*}
Since $\rho$ is polylogarithmic in $x_{k_0}$, while $\Lambda(x_{k_0})$ grows like $x_{k_0}$ up to powers of $\log x_{k_0}$, the factor $\Lambda(x_{k_0})\rho^{-2}$ tends to infinity as $x_{k_0}\to\infty$.

\emph{Conclusion.} Suppose an argument proved part~(b) of the Gaussian extraction from the listed properties alone. It would then apply to the true family and, with $E_0$ replaced by $2E_0$, equally to the model. The two bounds together would bound the covariance of the difference~\eqref{8.5:eq-moduli-insufficient-3} with $N\mathcal{O}^{(k_0)}(\tilde f_2)$ by the sum of the two right sides of part~(b), each of which under the conditions of that clause is a bounded multiple of $\mathfrak{M}$. This is incompatible with the growth of $\Lambda(x_{k_0})\rho^{-2}$ in the preceding paragraph. The true remainder must therefore be distinguished from the model through its production formula.
\end{proof}

\begin{definition}[The local run, its nested cells and image cells]\label{8.6:not-local-run}
\emph{The run.} In what follows the measure is the measure of Theorem~0, the Wilson measure at
bare coupling $g_0$ on the bare lattice of the torus $\mathbb{T}_m$
(Definition~\ref{1.1:def-wilson-action}). Ba{\l}aban's renormalisation steps $\mathbf{R}T$, in
which $T$ denotes Ba{\l}aban's renormalisation transformation of one step, built on the
block-averaging map of Definition~\ref{1.5:def-block-average}, and $\mathbf{R}$ the large-field
operation of Definition~\ref{1.5:def-densities}, are applied at the levels $k=1,\dots,K'$, where
\begin{equation}\label{8.6:eq-local-run-1}
K' := K-m .
\end{equation}
The sequence of these $K'$ steps is called the \emph{local run}, and $K'$ its \emph{terminal
level}; as in the correlation theorem, $m\ge 0$ is an integer that does not depend on $K$. The
only condition on the torus that enters there is a lower bound on $m$.

\emph{Cells.} For each level $k$, $\mathbb{T}^{(k)}$ denotes the block lattice of level $k$
(Definition~\ref{1.5:def-block-average}), read in the units in which its spacing is $1$. It
carries two level-$k$ partitions: one into the
\emph{$M$-cubes} of level $k$ (closed cubes of side $M$) and one into the \emph{cells} of level
$k$, a cell of level $k$ being a closed cube of side $M\check{R}_k$. Here $\check{R}_k$ is the
cell-size parameter of level $k$ fixed in the correlation theorem; it is nonincreasing in $k$.
By a consequence of the correlation theorem's definition of the local run, every side of a cube
of these partitions is a power of $L$ dividing the side of the torus, and the partitions of
different levels are nested $L$-adically. Each partition is compatible with all the other
partitions, in the sense of \cite{B-CMP119}.

\emph{Image cells.} Let $\mathfrak{c}$ be a cell of level $k'$ and let $k\ge k'$. The image of
$\mathfrak{c}$ at level $k$ is the same subset of the torus, read in the units of
$\mathbb{T}^{(k)}$; in those units it is a closed cube of side $M\check{R}_{k'}L^{-(k-k')}$. We
define
\begin{equation}\label{8.6:eq-local-run-2}
\mathfrak{c}^{(k)} := \text{the cell of level } k \text{ that contains the image of }
\mathfrak{c}.
\end{equation}
This is one cell, for the following reason. The parameters $\check{R}_j$ drop by at most one
factor $L$ per step, that is, $\check{R}_{j+1}\ge L^{-1}\check{R}_j$ for every $j$ with
$k'\le j<k$.
Iterating this over the $k-k'$ steps from level $k'$ to level $k$ gives
$\check{R}_k\ge L^{-(k-k')}\check{R}_{k'}$, and hence
\begin{equation*}
M\check{R}_{k'}L^{-(k-k')}\le M\check{R}_k .
\end{equation*}
The side of the image is therefore at most the side of a cell of level $k$. Since the partitions
are nested, the image lies in a single cell of level $k$, and that cell is
$\mathfrak{c}^{(k)}$. For $k=k'$ the definition gives $\mathfrak{c}^{(k')}=\mathfrak{c}$.
\end{definition}

\begin{definition}[The island representation of the prepared density]
\label{8.6:def-island-representation}
Fix a level $k$. The \emph{prepared density} at level $k$ is the level-$k$ density $\rho_k$ taken
after the preparation of the $k$-th step by Ba{\l}aban's large-field operation $\mathbf{R}$ of
Definition~\ref{1.5:def-densities} and before
the exponentiation \cite[(1.98)]{B-CMP122b}. In the form of \cite[(1.70)--(1.72)]{B-CMP122b} it is a
finite sum over outer data,
\begin{equation}\label{8.6:eq-island-representation-1}
  \rho_k=\sum_{\mathfrak{O}}\rho_{k,\mathfrak{O}},
\end{equation}
whose terms are described in the following items.

\smallskip
\noindent\emph{Outer data.} An \emph{outer datum} $\mathfrak{O}$ consists of two pieces of data:
\begin{itemize}
\item the large-field region $\tilde{Z}_k$ of the $k$-th step, together with its admissible
subsequence and with those of its internal variables that are not integrated at step $k$, as in
\cite[(1.70)--(1.72)]{B-CMP122b};
\item finitely many second-class components $(Y_i,\mathfrak{h}_i)$, $i\in I$, where $I$ is a finite
index set, each with its admissible subsequence and with those of its internal variables that are
not integrated at step $k$.
\end{itemize}
Here $Y_i$ is a region and $\mathfrak{h}_i$ is the history datum that the component carries in
\cite[(1.70)--(1.72)]{B-CMP122b}. The words \emph{first class}, \emph{second class} and
\emph{class component} refer to the
classification of large-field components in \cite[\S1]{B-CMP122a}, conditions (i) and (ii) there.
The \emph{live large-field region} of $\mathfrak{O}$ is the set
\begin{equation}\label{8.6:eq-island-representation-2}
  |\mathfrak{O}|:=\tilde{Z}_k\cup\bigcup_{i\in I}Y_i ,
\end{equation}
as in \cite[\S1]{B-CMP122b}.

\smallskip
\noindent\emph{Admissible island families.} Let an outer datum $\mathfrak{O}$ be given. An
\emph{island} is a first-class component healed at step $k$. Consider a
finite family $\mathfrak{X}=\{(X_a,\mathfrak{h}_a)\}_a$ of islands.
Here each $X_a$ is a region and $\mathfrak{h}_a$ is the history datum that the island carries in
\cite[(1.70)--(1.72)]{B-CMP122b}. We write $a\in\mathfrak{X}$ for the indices of the family. To
$\mathfrak{X}$ are attached the set $\mathcal{Y}(\mathfrak{X})$ of $V$-domains and the boundary
terms $V_{\mathfrak{X}}(Y)$, $Y\in\mathcal{Y}(\mathfrak{X})$, of \cite[(1.72)]{B-CMP122b}
(Ba{\l}aban's letter $V$ for the boundary terms; not to be confused with the block field $V_k$). The
$V$-domains $Y$ are not to be confused with the second-class components $Y_i$ of the outer datum,
which may lie inside them. We
write $Z(\mathfrak{X})^{\sharp}$ for the strip-enlarged region of the family $\mathfrak{X}$, the
strips being those of \cite[(2.1)--(2.3)]{B-CMP119}. The family $\mathfrak{X}$ is
\emph{admissible under $\mathfrak{O}$} if conditions (c1)--(c3) hold; its $V$-domains are then
described by (c4):
\begin{itemize}
\item[(c1)] each $X_a$ is, by itself, a class component satisfying conditions (i) and (ii) of
\cite[\S1]{B-CMP122a}, and is a box of the level-$k$ cells of
Notation~\ref{8.6:not-local-run}, each of whose sides has at most $100$ cells,
as in \cite[(2.3)]{B-CMP119};
\item[(c2)] the regions $X_a$, $a\in\mathfrak{X}$, pairwise share no point;
\item[(c3)] no $X_a$ shares a point with the live large-field region $|\mathfrak{O}|$;
\item[(c4)] every $Y\in\mathcal{Y}(\mathfrak{X})$ has two properties:
the meet $Y\sqcap Z(\mathfrak{X})^{\sharp}$ is a non-empty union of whole strips, and
$Y\setminus Z(\mathfrak{X})^{\sharp}\neq\emptyset$.
\end{itemize}
Conditions (c1)--(c3) hold vacuously for the empty family, so the empty family is admissible under
every outer datum.

\smallskip
\noindent\emph{Island operators.} For a first-class island $X_a$ the operator
\begin{equation}\label{8.6:eq-island-representation-3}
  \mathbf{T}'_k(X_a)=\sum_{\mathfrak{h}}\mathbf{T}'_{k,\mathfrak{h}}(X_a)
\end{equation}
is the normalised integral operation of the island, \cite[(1.71)]{B-CMP122b}. It includes the sum
over the island's internal admissible histories $\mathfrak{h}$, and $\mathbf{T}'_{k,\mathfrak{h}}(X_a)$
denotes the contribution of the history $\mathfrak{h}$. The operator $\mathbf{T}'_k(X_a)$ acts on
functions of the internal variables of $X_a$ and is the identity on all other variables. By (c2),
the operators of distinct islands of $\mathfrak{X}$ act on disjoint sets of variables. They therefore
commute, as in \cite[(2.19)]{B-CMP119}, so the product over $a\in\mathfrak{X}$ below does not depend
on the order of its factors.

\smallskip
\noindent\emph{The $\mathfrak{O}$-part.} For an outer datum $\mathfrak{O}$ the term
$\rho_{k,\mathfrak{O}}$ of \eqref{8.6:eq-island-representation-1} is
\begin{equation}\label{8.6:eq-island-representation-a}
\begin{gathered}
  \rho_{k,\mathfrak{O}}
  =\sum_{\mathfrak{X}\ \text{admissible under}\ \mathfrak{O}}
   \mathbf{O}_{\mathfrak{O},\mathfrak{X}}\bigl[W_{\mathfrak{O}}(\mathfrak{X})\bigr],\\
  W_{\mathfrak{O}}(\mathfrak{X})
  :=\Bigl[\prod_{a\in\mathfrak{X}}\mathbf{T}'_k(X_a)\Bigr]
   \Bigl(\exp\sum_{Y\in\mathcal{Y}(\mathfrak{X})}V_{\mathfrak{X}}(Y)\Bigr).
\end{gathered}
\end{equation}
In \eqref{8.6:eq-island-representation-a}, $\mathbf{O}_{\mathfrak{O},\mathfrak{X}}$ is the factor
that stands outside the bracket of \cite[(1.70)--(1.72)]{B-CMP122b} for the term
$(\mathfrak{O},\mathfrak{X})$. It is the composition of the following operations:
\begin{itemize}
\item multiplication by the characteristic functions of the new small-field domain;
\item multiplication by $\exp A'_k$;
\item the integral operations of the components $\tilde{Z}_k$ and $Y_i$, $i\in I$, of
$\mathfrak{O}$ (the live components);
\item integration over the internal variables of the live components.
\end{itemize}
Here $A'_k$ is the completed small-field action of \cite[\S1]{B-CMP122b}. It is obtained there by
adding to the effective action all the terms missing from it, so that its terms are, by definition,
those of the effective action on the whole lattice. The operator $\mathbf{O}_{\mathfrak{O},\mathfrak{X}}$
is a positive linear operator on non-negative functions of the block field $V_k$, of the internal
variables of $\mathfrak{O}$, and of the outputs of the island operators.

That the prepared density has the form \eqref{8.6:eq-island-representation-1},
\eqref{8.6:eq-island-representation-a} on the family of regions used in this book, with condition
(c3), with the sum over internal histories inside $\mathbf{T}'_k(X_a)$ and with the $V$-domains as
in (c4), is the representation lemma for the bracket of \cite[(1.70)--(1.72)]{B-CMP122b}, one of
the lemmas of the representation and bound of the large-field terms; in the form used in this
book, the identification of the terms of $A'_k$ with those of the effective action on the whole
lattice is part~(i) of the lemma identifying the terms of the completed action.
\end{definition}

\begin{lemma}[Monotonicity of the common prefactor under island removal. Stated; proof incomplete at the step marked $(\ast)$]\label{8.6:lem-prefactor-monotone}
Let $\mathcal{O}$, the families of islands admissible under $\mathcal{O}$, and, for each such family
$\mathfrak{X}$, the operator $\mathbf{O}_{\mathcal{O},\mathfrak{X}}$ be as in the island representation
of the prepared density, Definition~\ref{8.6:def-island-representation} and
\eqref{8.6:eq-island-representation-a}. Let $\mathfrak{X}\supset\mathfrak{X}'$ be two families admissible
under $\mathcal{O}$. Then
\begin{equation}\label{8.6:eq-prefactor-monotone-a}
\mathbf{O}_{\mathcal{O},\mathfrak{X}}\le\mathbf{O}_{\mathcal{O},\mathfrak{X}'}
\end{equation}
as positive operators. Moreover, $\mathbf{O}_{\mathcal{O},\mathfrak{X}}$ and
$\mathbf{O}_{\mathcal{O},\mathfrak{X}'}$ are equal except possibly in the small-field characteristic
functions of the current level $k$ (the top level of the hierarchy) on the regions of the islands of
$\mathfrak{X}\setminus\mathfrak{X}'$. For
$\mathfrak{X}$ these characteristic functions are those of the preparation,
\cite[(1.3)--(1.4)]{B-CMP122a}; for $\mathfrak{X}'$ they are the step's own,
\cite[(2.17), (3.2)]{B-CMP119}; both are taken on the same cubes of side $LM_2\check{R}_k$,
$\check{R}_k$ being the cube-size parameter of level $k$ fixed in \cite[\S1]{B-CMP122a}. The former are
dominated pointwise by the latter: wherever the preparation's characteristic function on such a cube
equals $1$, so does the step's.
\end{lemma}

Stated; the proof below is incomplete at the step marked $(\ast)$.

\begin{proof}
By Definition~\ref{8.6:def-island-representation}, $\mathbf{O}_{\mathcal{O},\mathfrak{X}}$ is the
factor standing outside the bracket in the term indexed by $(\mathcal{O},\mathfrak{X})$ in
\eqref{8.6:eq-island-representation-a}. It is the composition of four operations:
\begin{itemize}
\item multiplication by the characteristic functions of the new small-field domain;
\item multiplication by $\exp A'_k$, where $A'_k$ is the completed small-field action;
\item the integral operations of the live components;
\item integration over the internal variables of the live components.
\end{itemize}
Each of these is a positive linear operator on non-negative functions. We compare
$\mathbf{O}_{\mathcal{O},\mathfrak{X}}$ with $\mathbf{O}_{\mathcal{O},\mathfrak{X}'}$ one factor at a
time.

\emph{The completed action.} In \cite{B-CMP122b} all the missing terms are added to the effective action
in order to complete it, and the completed action is denoted $A'_k$. By that definition, its members are
those of the family taken on the whole lattice. Hence $A'_k$ does not depend on which
islands are healed. The factor $\exp A'_k$ is therefore the same for $\mathfrak{X}$ and for
$\mathfrak{X}'$.

\emph{The coupling profile.} The position-dependent coupling profile, in the sense of
Definition~\ref{1.5:def-domains}, is reset to $1$ on the region of every island healed at this step, in
particular, in the term of $\mathfrak{X}$, on the regions $X_a$ of the islands of
$\mathfrak{X}\setminus\mathfrak{X}'$; this reset is part of the statement fixing the coupling profile
after a region is healed. By that statement the profile coincides for $\mathfrak{X}$ and for
$\mathfrak{X}'$.

\emph{The new domains and the live components.} The new domains coincide term by term as hierarchies of
regions; the characteristic functions they carry coincide off the regions $X_a$ of the islands of
$\mathfrak{X}\setminus\mathfrak{X}'$, which are compared below. So do the
integral operations of the live components and the integrations over their internal variables.

\emph{Members that read an island.} Some members of the term depend on the variables of an island: the
boundary terms; the members that carry the position-dependent coupling profile of
Definition~\ref{1.5:def-domains}, which we write $\varphi$ and call the $\varphi$-members; and the
vacuum constants. Each of these stands inside the bracket of \eqref{8.6:eq-island-representation-a},
not in the common prefactor, by the positivity and bound of the normalised island operation,
\eqref{8.6:eq-island-operation-bound-d}, by the lemma on the $\varphi$-members, and by the
normalisation of the vacuum constants.
These last two statements are part of the representation and bound of the large-field terms. None of
these members, therefore, contributes to $\mathbf{O}_{\mathcal{O},\mathfrak{X}}$ or to
$\mathbf{O}_{\mathcal{O},\mathfrak{X}'}$.

\emph{The remaining factor.} The only factors not yet compared are the top-level small-field
characteristic functions on the regions $X_a$ of the islands of
$\mathfrak{X}\setminus\mathfrak{X}'$. Write $\chi^{\mathrm{prep}}_a$ for the one carried by
$\mathbf{O}_{\mathcal{O},\mathfrak{X}}$ on $X_a$. It is the preparation's characteristic function of
\cite[(1.3)--(1.4)]{B-CMP122a}, on the cubes of side $LM_2\check{R}_k$ of $X_a$. Write
$\chi^{\mathrm{step}}_a$ for the one carried by $\mathbf{O}_{\mathcal{O},\mathfrak{X}'}$. It is the
step's own characteristic function of \cite[(2.17), (3.2)]{B-CMP119}, on the same cubes.

$(\ast)$ \emph{Wherever $\chi^{\mathrm{prep}}_a$ equals $1$, so does $\chi^{\mathrm{step}}_a$.} Two
statements bear on this domination. First, \cite{B-CMP122b} states that the new action satisfies the
induction hypothesis and that the new large-field region equals $\tilde{Z}_k\cup\bigcup_i Y_i$; here, in
the notation of \cite[(1.72)]{B-CMP122b}, $\tilde{Z}_k$ and the regions $Y_i$ are the regions of the
outer datum $(\tilde{Z}_k,\{Y_i\})$, which carry the large-field operations standing in the common
prefactor. The regions $X_a$ of the healed islands are thus not part of the new large-field region, and
the new action, with the characteristic functions it carries on $X_a$, satisfies the induction
hypothesis at level $k$; a characteristic function of the preparation equal to $1$ at a configuration on
$X_a$ at which the step's characteristic function vanishes would be incompatible with that statement.
Second, the coupling profile is reset to $1$ on the healed region, so the thresholds on $X_a$ are the
standard ones of level $k$. It remains to compare the thresholds of
\cite[(1.3)--(1.4)]{B-CMP122a} with those of \cite[(2.17), (3.2)]{B-CMP119} on the cubes of side
$LM_2\check{R}_k$ of $X_a$.

Granting $(\ast)$, we have $\chi^{\mathrm{prep}}_a\le\chi^{\mathrm{step}}_a$ pointwise for every island
$X_a$ of $\mathfrak{X}\setminus\mathfrak{X}'$; both are characteristic functions, so
\begin{equation*}
0\le\chi^{\mathrm{prep}}_a\le\chi^{\mathrm{step}}_a .
\end{equation*}
The products over the islands of $\mathfrak{X}\setminus\mathfrak{X}'$ satisfy the same inequality:
\begin{equation*}
0\le\prod_{a}\chi^{\mathrm{prep}}_a\le\prod_{a}\chi^{\mathrm{step}}_a .
\end{equation*}
Let $F\ge0$. The two compositions differ in exactly one factor: multiplication by
$\prod_a\chi^{\mathrm{prep}}_a$ in $\mathbf{O}_{\mathcal{O},\mathfrak{X}}$ against multiplication by
$\prod_a\chi^{\mathrm{step}}_a$ in $\mathbf{O}_{\mathcal{O},\mathfrak{X}'}$. By the last display the
first multiplication is dominated by the second as positive operators, wherever it stands in the
composition. Every other factor is common to both prefactors, as shown above, and is a positive linear
operator. A positive operator preserves the pointwise order, so composing the same positive operators
before and after two ordered factors in the same position preserves their order. Hence
$\mathbf{O}_{\mathcal{O},\mathfrak{X}}F\le\mathbf{O}_{\mathcal{O},\mathfrak{X}'}F$ for every $F\ge0$.
This is \eqref{8.6:eq-prefactor-monotone-a}. The argument also shows that the two operators can differ
in these characteristic functions and nowhere else.

\emph{On equality.} In \cite[(1.72)]{B-CMP122b} the factor outside the curly bracket is
\begin{equation*}
\chi_k(\cdot)\Bigl(\sum\mathbf{T}''_k(\tilde{Z}_k)\Bigr)\exp A'_k\cdot\prod_i\bigl[\,\cdots\,
\mathbf{T}'_k(Y_i)\bigr],
\end{equation*}
where $\mathbf{T}''_k(\tilde{Z}_k)$ and $\mathbf{T}'_k(Y_i)$ are the large-field operations of
\cite[(1.72)]{B-CMP122b} on $\tilde{Z}_k$ and on the $Y_i$, and $\chi_k(\cdot)$ is the step's
small-field characteristic function.
This factor is indexed by the outer datum $(\tilde{Z}_k,\{Y_i\})$ only. In \cite{B-CMP122b} the new
$k$-th domain, there written $\Omega_k$, is the one defined by the determining set; consequently the
step's small-field characteristic function is placed on the regions of the islands as well. The
preparation's characteristic functions $\chi'(X)$ of \cite[(1.3)--(1.4)]{B-CMP122a} on a region $X$
enter inside $\mathbf{T}'_k(X)$,
by \cite[(1.71)]{B-CMP122b}. Hence equality in \eqref{8.6:eq-prefactor-monotone-a} reduces to the
identification of the argument of $\chi_k(\cdot)$ in that factor with the new domain $\Omega_k$.
\end{proof}

\begin{definition}[Histories of the run, their masses and the annealed measure]
\label{8.6:def-history-measure}
Throughout this definition the complex sources vanish, $z=0$, and all gauge fields are real.
\begin{itemize}
\item[(1)] \emph{Histories.} $\mathfrak{H}$ denotes the finite set of un-resummed histories of the
run.
A history $\omega\in\mathfrak{H}$ records, for every decomposition of unity performed at every step,
the outcome selected in it; $\mathfrak{H}$ is the set of all such records.
In particular, at each step $k$ it records the outer datum of level $k$, the island family
$\mathfrak{X}_k$ of the island representation of the prepared density
(Definition~\ref{8.6:def-island-representation}) together with the internal histories of its
islands, and the outcomes of the preparatory decompositions.
\item[(2)] \emph{Masses.} For $\omega\in\mathfrak{H}$, the \emph{term} of $\omega$ is the summand of
the fully expanded density selected by the outcomes that $\omega$ records, and $\tau_\omega$ is the
total integral of the term of $\omega$ over all its variables. We call $\tau_\omega$ the \emph{mass}
of $\omega$ and also write $\mathfrak{m}(\omega):=\tau_\omega$. Then $\tau_\omega\ge0$ for every
$\omega\in\mathfrak{H}$.
\item[(3)] \emph{Normalisation.} We put
\begin{equation}\label{8.6:eq-history-measure-1}
\mathsf{Z}:=\sum_{\omega\in\mathfrak{H}}\tau_\omega .
\end{equation}
Then $\mathsf{Z}=\int\rho_0$.
\item[(4)] \emph{The annealed measure.} For $E\subset\mathfrak{H}$ we put
\begin{equation}\label{8.6:eq-history-measure-a}
\tilde\nu(E):=\mathsf{Z}^{-1}\sum_{\omega\in E}\tau_\omega .
\end{equation}
\item[(5)] \emph{The measure jointly with the top field.} Fix a level $k$. For
$\omega\in\mathfrak{H}$
write $\tau_{\omega,k}(V)$ for the integral of the level-$k$ term of $\omega$ (the summand of the
expanded density of level $k$ selected by $\omega$) over all variables
other than the top field $V_k$, evaluated at $V_k=V$. For $E\subset\mathfrak{H}$ we put
\begin{equation}\label{8.6:eq-history-measure-2}
\tilde\nu_k(E\times dV):=\mathsf{Z}^{-1}\sum_{\omega\in E}\tau_{\omega,k}(V)\,dV .
\end{equation}
\item[(6)] \emph{Two index spaces.} Two index spaces for the terms of the run are used.
\begin{itemize}
\item[(a)] The index of Ba{\l}aban's expansion, which is also the index of the expansion of the
correlation theorem, keeps only the current large-field data. In it, healed components are summed
out \cite[(0.4)]{B-CMP122a}.
\item[(b)] The un-resummed index $\mathfrak{H}$ of item~(1) keeps every label and never forgets it.
\end{itemize}
The space (b) refines the space (a). An event that names a healed label is an event on (b). The
other events used in what follows, which do not name a healed label, are measurable on (a).
Every $\tilde\nu$-value of an
(a)-measurable event is the same whether it is computed on (a) or on (b).
\end{itemize}
\end{definition}

\begin{proof}
The assertions in items~(2), (3) and~(6) require proof. Items~(1), (4) and~(5) are definitions.

\emph{Item~\textup{(2)}.} Fix $\omega\in\mathfrak{H}$. Every factor of the term of $\omega$ belongs
to one of three kinds:
\begin{itemize}
\item a characteristic function;
\item the exponential of a real action;
\item an integral operation whose density is non-negative at real gauge fields and vanishing
sources.
\end{itemize}
This holds for the factors of the large-field operations and for the densities they produce
\cite{B-CMP122a}, \cite{B-CMP122b}.

A characteristic function takes values in $\{0,1\}$. The exponential of a real number is positive.
Integration against a non-negative density maps non-negative functions to non-negative functions.
Hence the term of $\omega$ is a non-negative function of its variables. Its total integral
$\tau_\omega$ is therefore non-negative.

\emph{Item~\textup{(3)}.} The set $\mathfrak{H}$ indexes the complete assignments of outcomes to the
decompositions of unity performed along the run. Each such decomposition, and each expansion
performed at each step, is an
identity of functions. Summing the terms over all outcomes therefore reproduces the density before
the decomposition.

Each operation of the run preserves total integrals. This holds for the renormalisation
transformation from level $k$ to level $k+1$ \cite[(0.1), (3.1)]{B-CMP119}, and for the
large-field operation for each fixed outer datum \cite[(0.4), (1.102)]{B-CMP122a}.

It follows that $\sum_{\omega\in\mathfrak{H}}\tau_\omega$ equals the total integral of the bare
density $\rho_0$. This proves $\mathsf{Z}=\int\rho_0$. In particular $\mathsf{Z}>0$, since $\rho_0$
is the exponential of a real action.

By item~(2) each summand in \eqref{8.6:eq-history-measure-a} and in
\eqref{8.6:eq-history-measure-1} is non-negative. Hence $\tilde\nu$ is a non-negative set function
on the subsets of $\mathfrak{H}$. It is additive, and by \eqref{8.6:eq-history-measure-1} it
satisfies $\tilde\nu(\mathfrak{H})=1$.

\emph{Item~\textup{(6)}.} Let $\pi$ be the map that assigns to a history $\omega$ in the space (b)
the index in the space (a) obtained by forgetting the labels of healed components. The space (a)
is obtained from the space (b) by summing out the healed components. Hence the mass of an index
$\iota$ in (a) is the sum of the masses $\tau_\omega$ over $\omega\in\pi^{-1}(\iota)$.

An (a)-measurable event is a set $E_a$ of indices of (a). On (b) it is the event $\pi^{-1}(E_a)$.
Grouping the sum in \eqref{8.6:eq-history-measure-a} according to the value of $\pi$ gives
\begin{equation}\label{8.6:eq-history-measure-3}
\tilde\nu\bigl(\pi^{-1}(E_a)\bigr)
=\mathsf{Z}^{-1}\sum_{\iota\in E_a}\ \sum_{\omega\in\pi^{-1}(\iota)}\tau_\omega .
\end{equation}
The inner sum is the mass of $\iota$ in (a). The right-hand side is therefore the value of the
event $E_a$ computed on (a).

An event naming a healed label is defined on the space (b), where that label is retained; it is
not in general of the form $\pi^{-1}(E_a)$.
\end{proof}

\begin{definition}[Healed labels, live labels, seeds and lineages]\label{8.6:def-labels-lineages}
Every large-field label, that is, every component, live or healed, falls under one of the following
three kinds.
\begin{itemize}
\item[(H)] A \emph{healed label} is a pair $(k,X)$ in which $X$ belongs to the family
$\mathfrak{X}_k$ of islands healed at step $k$, taken together with the histories internal to it,
in the sense of Definition~\ref{8.6:def-history-measure}. The \emph{cells} of a healed label
$(k,X)$ are level-$k$ cells.
\item[(L)] A \emph{live label} at level $k'$ is a connected component $C$ of the live large-field
set $|\mathcal{O}_{k'}|$ at level $k'$, the set whose components are live in the sense of this
chapter; that is, $C$ is either a component of the current large-field
region $\widetilde{Z}_{k'}$ at level $k'$ or a component $Y$ of the second class.
\item[(S)] A \emph{seed}, or creation event, at $(j,\square)$ is a component of the cover of the new
large-field region created at step $j$ in Ba{\l}aban's large-field operation \cite{B-CMP122b}, where
$\square$ is a cube of that cover lying in the component. Its
\emph{birth level} is $j$. Seeds of a second, preparatory, type are also admitted: the creation
events of the preparatory type fixed together with the histories of the run in
Definition~\ref{8.6:def-history-measure}.
\end{itemize}
For a label, $\mathcal{Z}(\cdot)$ denotes the set of seeds occurring in its history, and for a seed,
$j(\cdot)$ denotes its birth level; a label formed by merging carries the oldest birth level among
the labels merged, also written $j(\cdot)$ (see below).

\emph{Nesting along a history.} By \cite[(2.1), (2.3)]{B-CMP119}, the sequence of Ba{\l}aban's
domains $\Omega_j,\Lambda_j$ of that paper attached to a history is nested,
\begin{equation*}
\Omega_1\supset\Lambda_1\supset\Omega_2\supset\cdots,
\end{equation*}
and with $Z_j:=\Lambda_j^{c}$ one has $Z_j\subset Z_{j+1}$ along a history, until a step of the
operation $\mathbf{R}$ resets $\Lambda$ on the healed region \cite[\S2]{B-CMP119}. Consequently a
seed created at $(j,\square)$ lies, at every later level up to the step at which its component is
healed, inside a live component that contains the image of $\square$ at that level, and, at the
healing step, inside the healed island.

\emph{Lineage of a live component.} Let $C$ be a live component at level $k'$. Its \emph{lineage} is
the sequence of live components that contain the image of $C$ at the levels $k',k'+1,\dots$, taken
up to the first step $k$ at which that component is healed, that is, absorbed into an island
$X\in\mathfrak{X}_k$ which contains the image of $C$, or up to the final level $K'$ of the
run if it is never healed. The image of $C$ contains, for every cell
$\mathfrak{c}\subset C$, the level-$k$ cell $\mathfrak{c}^{(k)}$ that is the image of $\mathfrak{c}$
at level $k$.

\emph{Lineage of a healed island.} Let $X\in\mathfrak{X}_k$. Its \emph{lineage} consists of its live
ancestor components at the levels $j(X),\dots,k-1$, together with every label merged into them.
Here, as in \cite[(1.84)--(1.86)]{B-CMP122b}, a component formed by merging inherits the oldest
birth level among the components merged, and $j(X)$ is the birth level inherited in this way by
$X$.

\emph{Seeds of many cubes.} A connected seed region consisting of many cubes, all created at the
same step, is one seed, and the quantity $\theta p_0$ enters the weight $\varpi_\theta$ once for it.
The size of such a region enters only through the term $-\kappa_1 d_j(\,\cdot\,)$ of Ba{\l}aban's
large-field bound \cite{B-CMP122b} at the level $j$ considered, $d_j(\cdot)$ being the size
function of the glossary. The
one-point live exponent $\theta$ does not depend on the size of the seed.
\end{definition}

\begin{definition}[Reserve weight of a set of internal histories]\label{8.6:def-reserve-weight}
Let $\beta_0$ be the constant of \cite[(1.87), (1.89)]{B-CMP122b}. It satisfies $\beta_0<1/3$; an example given there is $\beta_0=1/7$. Let $\theta$ be a number in the window
\begin{equation}\label{8.6:eq-reserve-weight-b}
0<\theta<\frac{1-3\beta_0}{2}.
\end{equation}
Since $\beta_0<1/3$, the right-hand side of \eqref{8.6:eq-reserve-weight-b} is positive, so the window
is non-empty. The number $\theta$ is one of the Stage-0 numbers of the order of choice
(cf.\ Definition~\ref{2.3:not-stages-first}); when the reserve weight is used, $\theta$ enters through the
per-seed reserve statement, which supplies a factor $e^{-\theta p_0(g_{j(Z)})}$ at an $\mathbf{R}$-step,
with $Z$ and $j(Z)$ as below. Condition
\eqref{8.6:eq-reserve-weight-b} constrains the auxiliary number $\theta$ only; it is not a hypothesis of
Theorem~\ref{3.1:thm-main}.

Fix an internal history $\mathfrak{h}$, in the sense of the histories of Definition~\ref{8.6:def-labels-lineages}. We write $\mathcal{Z}(\mathfrak{h})$ for the set of large-field regions $Z$ of $\mathfrak{h}$. For $Z\in\mathcal{Z}(\mathfrak{h})$ we write $j(Z)$ for the step at which $Z$ is created. We write $p_0(\cdot)$ for the degree function, so that $p_0(g_{j(Z)})$ is its value at the running coupling $g_{j(Z)}$ of Definition~\ref{1.2:def-couplings} at level $j(Z)$.

For a set $\mathcal{S}$ of internal histories, the \emph{reserve weight} of $\mathcal{S}$ is
\begin{equation}\label{8.6:eq-reserve-weight-a}
\varpi_\theta(\mathcal{S})
:=\sup_{\mathfrak{h}\in\mathcal{S}}
\exp\Bigl(-\theta\sum_{Z\in\mathcal{Z}(\mathfrak{h})}p_0\bigl(g_{j(Z)}\bigr)\Bigr)
\in\mathopen{]}0,1].
\end{equation}
The right-hand side of \eqref{8.6:eq-reserve-weight-a} also makes sense at $\theta=0$, where each
exponential factor equals $1$; we write $\varpi_0\equiv1$ for this case, in which no reserve is read
and the per-seed reserve statement is not used.
\end{definition}

\begin{definition}[Running couplings and the degree function along the run]\label{8.6:not-running-couplings}
Throughout this chapter $x_k=g_k^{-2}$, $0\le k\le K$, are the running inverse couplings of
Definition~\ref{1.2:def-couplings}. The number $g^{-2}$ is the inverse square of the reference
coupling $g$. In this chapter it is always written out as $g^{-2}$, because the letter $X$ is
reserved for large-field regions.

\emph{The degree function.} For a coupling $g'$ the degree function $p_0(\cdot)$ is
\begin{equation*}
  p_0(g')=A_0\bigl(\log g'^{-2}\bigr)^{p_0},
  \qquad\text{so that}\qquad p_0(g_k)=A_0(\log x_k)^{p_0}.
\end{equation*}
Its fixed exponent $p_0$ satisfies the one-sided condition $p_0\ge 5r$, as in
\cite{B-CMP119}. Here $r$ is an exponent of \cite{B-CMP119}, not to be confused with the
extraction depth $r(s)$ or the offset $r_0$ of this chapter. With the same $r$,
\cite{B-CMP119} states the bound
\begin{equation*}
  \check{R}_k\le L(\log x_k)^{r}
\end{equation*}
for its level-$k$ exponent $\check{R}_k$, which enters this chapter through the factors
$L^{4\check{R}_k}$.

\emph{The lower flow bound.} Clause~(0) of Theorem~\ref{3.1:thm-main} gives
\begin{equation}\label{8.6:eq-running-couplings-a}
  x_i-x_j\ \ge\ \lambda^{\sharp}(j-i)-c_5\qquad\text{for all } i\le j\le K .
\end{equation}
Here $\lambda^{\sharp}$ and $c_5=2c_2$ are the flow constants of
Definition~\ref{1.2:def-flow-constants}. By the comparison of the reference and running
couplings,
\begin{equation*}
  g^{-2}\le x_k\qquad\text{for every } k .
\end{equation*}
The bound \eqref{8.6:eq-running-couplings-a} is one-sided.

\emph{Comparison with the correlation theorem.} The following is recorded for comparison and
is not part of the definition. Among the hypotheses of the correlation theorem is the bound
\begin{equation}\label{8.6:eq-running-couplings-1}
  \lambda(j'-j)-2c_2\ \le\ x_j-x_{j'}\ \le\ 3\lambda(j'-j)+2c_2 .
\end{equation}
It is required together with the condition on the last coupling
\begin{equation*}
  x_K\in\mathopen{]}g^{-2},\,g^{-2}+b_+\mathclose{]} .
\end{equation*}
In the two displays above, $\lambda$ and $b_+$ are the window constants fixed in the statement
of the correlation theorem; this $\lambda$ is not the gauge transformation of
Lemma~\ref{1.1:lem-invariance}. The two-sided bound \eqref{8.6:eq-running-couplings-1} is a
hypothesis of the correlation theorem and of the large-field block bound. Its lower half has
the same form as \eqref{8.6:eq-running-couplings-a}, with $\lambda$ in place of
$\lambda^{\sharp}$ and $2c_2=c_5$. From the window \eqref{8.6:eq-running-couplings-1} one has
moreover, for $j\le k'$,
\begin{equation*}
  q_j\le e^{A_0}q_{k'},
\end{equation*}
where $q_j$ is the level-$j$ quantity fixed in this section in the definition of the
large-field regions created at each step; the implication is derived there from
\eqref{8.6:eq-running-couplings-1}.
\end{definition}

\begin{definition}[Positivity and bound of the normalised island operation]
\label{8.6:def-island-operation-bound}
Fix a level $k$, $0\le k\le K$, and the island representation of the level-$k$ density of
Definition~\ref{8.6:def-island-representation}, with its family $\mathfrak{X}$ of islands, its
outer datum $\mathfrak{O}$, and the internal histories of its islands
(Definition~\ref{8.6:def-history-measure}). For an island $X$ and a set $\mathcal{S}$ of
internal histories of $X$, write $\mathbf{T}'_{k,\mathcal{S}}(X)$ for Ba{\l}aban's normalised
operation $\mathbf{T}'_k(X)$ of \cite[\S1, (1.73)]{B-CMP122b} with its sum over internal
histories restricted to $\mathcal{S}$. Thus $\mathbf{T}'_{k,\mathcal{S}}(X)=\mathbf{T}'_k(X)$ when
$\mathcal{S}$ is the set of all internal histories of $X$. We write
$\mathbf{T}'_{k,\mathcal{S}}(X,(U,\sigma))$ for this operation at an argument $(U,\sigma)$, where
$\sigma$ denotes the second component of the argument of $\mathbf{T}'_k(X)$ in
\cite[\S1]{B-CMP122b}, and $\mathbf{T}'_{k,\mathcal{S}}(X,(U,0))$ for its value at
the argument whose second component $\sigma$ vanishes identically.

The localised terms of the representation are denoted $V_{\mathfrak{X}}(Y)$, indexed by the
domains $Y$ of the set $\mathcal{Y}(\mathfrak{X})$. The enlarged large-field region
$Z(\mathfrak{X})^\sharp$, the family $\mathcal{Y}(\mathfrak{X})$ and the distance
$d_\sharp(Y)\in[0,\infty]$ are those of Definition~\ref{8.6:def-island-representation}. For a
family $\mathfrak{X}$ we write $\mathbf{O}_{\mathfrak{O},\mathfrak{X}}$ for
the operation that the representation attaches to the outer datum $\mathfrak{O}$ and the
family $\mathfrak{X}$. For an island $X$ not in $\mathfrak{X}'$, the family obtained by
adjoining $X$ to $\mathfrak{X}'$ is written $\mathfrak{X}'\sqcup X$.

The \emph{island hypotheses at level $k$} are the following statements. Each is referred to by
the name in front of it.
\begin{enumerate}
\item The positivity and comparison bound (comparison with the value at $(U,0)$). For every
first-class island $X$ admissible on its own, every set
$\mathcal{S}$ of internal histories of $X$, and every non-negative bounded function $\Phi$ of the
internal variables of $X$ and of other variables, we have, pointwise in the other variables,
\begin{equation}\label{8.6:eq-island-operation-bound-a}
0\;\le\;\mathbf{T}'_{k,\mathcal{S}}(X)\,\Phi
\;\le\;\bigl(\mathbf{T}'_{k,\mathcal{S}}(X,(U,0))\,1\bigr)\,\sup\nolimits_X\Phi .
\end{equation}
Here $\sup_X$ is the supremum over the internal variables of $X$, restricted to the supports of
the characteristic functions entering $\mathbf{T}'_k(X)$. This is the positivity form of
\cite[(1.73)]{B-CMP122b} at the argument $(U,0)$. Since the densities summed over internal
histories are non-negative, the same inequalities hold for every sub-sum
$\mathbf{T}'_{k,\mathcal{S}}(X)$. At the argument $(U,0)$ the variable $\sigma$ vanishes
identically and no exponential prefactor is present.

\item The island bound at $(U,0)$. For every first-class island $X$ admissible on its
own, pointwise in $V_k$ and in the outer variables,
\begin{equation}\label{8.6:eq-island-operation-bound-b}
\mathbf{T}'_k(X,(U,0))\,1\;\le\;\exp\bigl(-2(1+\beta_0)^{-1}p_0(g_k)\bigr).
\end{equation}
Here $p_0(\cdot)$ is the degree function
(Definition~\ref{8.6:not-running-couplings}), and $\beta_0$ is the exponent parameter of
\cite[(1.89)]{B-CMP122b}. This is \cite[(1.89)]{B-CMP122b}.

\item The decay bound for the localised terms. For every $Y\in\mathcal{Y}(\mathfrak{X})$
with $Y\setminus Z(\mathfrak{X})^\sharp\neq\emptyset$,
\begin{equation}\label{8.6:eq-island-operation-bound-c}
\bigl|V_{\mathfrak{X}}(Y)\bigr|\;\le\;\alpha\,\exp\bigl(-(1+2\beta)\kappa\,d_\sharp(Y)\bigr).
\end{equation}
The bound holds uniformly in the histories, in the first-class internal variables on the
supports, and in $V_k$ on the support of the term. Moreover $V_{\mathfrak{X}}(Y)=0$ if
$d_\sharp(Y)=\infty$. Here $\alpha$ is the constant prefactor and $\beta$ the exponent parameter
of \cite[(1.68)]{B-CMP122b}, and $\kappa$ is the decay parameter of the tuple $\mathfrak{p}$ of
auxiliary parameters (Definition~\ref{2.1:not-prefix}). The bound
has the shape of \cite[(1.68)]{B-CMP122b}, written with $d_\sharp$ and $Z(\mathfrak{X})^\sharp$.

\item The locality statement (for the localised terms). The term $V_{\mathfrak{X}}(Y)$ depends
only on three things: the pair $(\mathbf{U},\mathbf{J})$ of field variables of the
representation restricted to $Y$; the variables of the regions $Y_i\subset Y$ of the
representation; and the internal variables and histories of the islands $X_a\subset Y$.
Furthermore, let $\mathfrak{X}$ and $\mathfrak{X}'$ have the same islands inside $Y$, and let
$Y\in\mathcal{Y}(\mathfrak{X})\cap\mathcal{Y}(\mathfrak{X}')$. Then
\begin{equation}\label{8.6:eq-island-operation-bound-d}
V_{\mathfrak{X}}(Y)\;=\;V_{\mathfrak{X}'}(Y).
\end{equation}
This is the remark headed ``Notice'' in \cite[\S1]{B-CMP122b}, used here as stated there.

\item The bracket on the family of islands. The bracket
\cite[(1.70)--(1.72)]{B-CMP122b} holds on the family $\mathfrak{X}$ of islands of the
representation, with its four side conditions and its constant $K_3(\beta)$.

\item Monotonicity in the family of islands. For every family
$\mathfrak{X}'$ and every island $X$ for which $\mathfrak{X}'\sqcup X$ is a family of the
representation,
\begin{equation*}
\mathbf{O}_{\mathfrak{O},\mathfrak{X}'\sqcup X}\;\le\;\mathbf{O}_{\mathfrak{O},\mathfrak{X}'} .
\end{equation*}
The inequality between operations means that, for every non-negative $\Phi$, pointwise,
\begin{equation*}
\mathbf{O}_{\mathfrak{O},\mathfrak{X}'\sqcup X}\,\Phi\;\le\;\mathbf{O}_{\mathfrak{O},\mathfrak{X}'}\,\Phi .
\end{equation*}

\item The positivity statement (for the terms of the representation). Every term $\tau_h$ of the
representation, $h$ a
history (Definition~\ref{8.6:def-history-measure}), is non-negative at real fields. This
is the positivity of the densities in \cite[\S0]{B-CMP122a}, in \cite[\S1]{B-CMP122b} and in
\cite[(2.17)--(2.22)]{B-CMP119}.

\item Preservation of integrals. Let $T$ be the renormalisation transformation of
\cite[(0.1)]{B-CMP119}, whose kernel
$\delta(\bar V_kV_{k+1}^{-1})$, with $\bar V_k$ the block average of $V_k$ under the map of
Definition~\ref{1.5:def-block-average}, has integral one \cite[(3.1)]{B-CMP119}. Then for every
density $\rho$,
\begin{equation*}
\int dV_{k+1}\,(T\rho)(V_{k+1})=\int dV_k\,\rho(V_k).
\end{equation*}
For each outer datum we also have
\begin{equation*}
\int dV\,(\mathbf{R}\rho)(V)=\int dV\,\rho(V),
\end{equation*}
as in \cite[(0.4), (1.102)]{B-CMP122a}. The Mayer expansion and the identity
\cite[(1.98)]{B-CMP122a} are identities of functions.

\item The geometry conditions and the nesting condition (geometry of islands, strips and
nesting). They state the following.
\begin{itemize}
\item Islands and live components are unions of cells.
\item Every class component, that is, every island and every second-class domain $Y$, is a
rectangular box of cells contained in a cube of side $100M\check R_k$, where $\check R_k$ is the
size parameter of class components at level $k$.
\item Let the strips $S^\sharp$ and $S'^\sharp$ belong to components $S$ and $S'$ that share
no point. Then $S^\sharp$ and $S'^\sharp$ share no point. Moreover $S^\sharp\subset S^{\boxplus}$,
where $S^{\boxplus}$ is the set attached to the component $S$, together with its strip
$S^\sharp$, in the conventions of \cite[(2.1), (2.3)]{B-CMP119}.
\item The number $N(X^\sharp)$ of cubes of side $M$ in the strip $X^\sharp$ of an island $X$
satisfies $N(X^\sharp)\le(101\check R_k)^d$.
\item Along a history, large-field regions are nested until they are healed, with one cell
$\mathfrak{c}^{(j)}$ at each later level $j$.
\end{itemize}
The conventions are those of \cite[(2.1), (2.3)]{B-CMP119} and of
\cite[condition~(i)]{B-CMP122a}, where the regions are rectangular parallelepipeds; the nesting
along a history is stated with the correlation theorem.

\item The tree-graph, touch-connected and anchor counts. These are three counting bounds.
\begin{itemize}
\item The tree-graph count. For every $t\ge0$, the number of cubes of side $M$ met by a tree
graph of length $tM$ is at most $5^d(t+1)$.
\item The touch-connected count. For every integer $n\ge1$, the number of touch-connected
families of $n$ cubes (families connected for the relation of sharing a point) through a given
cube is at most $9^{d(n-1)}$.
\item The anchor count. Let $\square$ be a cube and $D\ge0$ an integer. Put
\begin{equation*}
A_d:=5^d\bigl[(1+3^d)\log 2+2d\log 3\bigr].
\end{equation*}
Then
\begin{equation*}
\#\bigl\{Y\in\mathcal{Y}(\mathfrak{X}):\ \square\text{ is an anchor cube of }Y,\
\lfloor d_\sharp(Y)\rfloor=D\bigr\}\;\le\;2e^{A_d(D+2)} .
\end{equation*}
Here an anchor cube of $Y$ is one of the cubes that the island representation of
Definition~\ref{8.6:def-island-representation} attaches to $Y$ as anchors.
The proof of this bound uses of the anchor cube $\square$ only that the shortest tree graph of
$Y$ meets it.
\end{itemize}

\item The lower flow bound (lower bound on the differences of the running inverse couplings). For
$i\le j\le K$,
\begin{equation*}
x_i-x_j\;\ge\;\lambda^\sharp(j-i)-c_5 .
\end{equation*}
Moreover $g^{-2}\le x_j$ for $0\le j\le K$, where $g$ is the reference coupling
(Definition~\ref{0.2:not-glossary}). The first bound is that of clause~\textup{(0)} of
Theorem~\ref{3.1:thm-main}; the second is stated with the correlation theorem.

\item The reserve per seed. At a step at which the operation $\mathbf{R}$ acts, each
seed $Z$ of an island $X$
(Definition~\ref{8.6:def-labels-lineages}) contributes the reserve factor
$e^{-\theta p_0(g_{j(Z)})}$ to the bound on $\mathbf{T}_k(X)\,1$, where $\mathbf{T}_k(X)$ is the
operation of \cite[\S1]{B-CMP122b} without the normalisation of $\mathbf{T}'_k(X)$; in this
chapter the reserve is read only on the normalised operation $\mathbf{T}'_k(X)$. Here $\theta$ is
an exponent in the range fixed with the reserve weight
(Definition~\ref{8.6:def-reserve-weight}), and $j(Z)$ is the step
at which $Z$ was created. The same holds in the variant of this reserve for linear members. The
inputs are \cite[(1.79)--(1.89)]{B-CMP122b}.

\item The torus-independence statement (independence of the step from the torus exponent). The
data read by the $k$-th step of the
run are $g_k$, the level-$k$ parameter $\varepsilon_k$ of Ba{\l}aban's construction,
$\check R_k$, the partitions of $\mathbb{T}^{(k)}$, and, in the sourced expansion of the
correlation theorem, the frozen running shift $\zeta'_k$. None of these depends on $m$. The
exponent $m$ enters the run only through the stopping level $K'$ of the representation and
through one condition at the lowest levels.
Hence two admissible runs on tori of exponents $\tilde m\le m$ coincide through step $K-m$. For
every event $E$ determined by labels of levels at most $K-m$, the annealed measure
$\tilde\nu(E)$ of Definition~\ref{8.6:def-history-measure} is the same number for both runs.
This statement is stated with the correlation theorem.

\item The arithmetic of the absorptions. We have $p_0\ge 5r$, the one-sided
condition on the exponent of the degree function. We also have
\begin{equation*}
L^{4\check R(x)}\operatorname{polylog}(x)\;\le\;e^{\frac12 c_q p_0(g')} .
\end{equation*}
Here $g'$ is a coupling at which the inequality is used in an absorption of the two-point
witness theorem or in an upper bound on $\gamma$ imposed in this chapter, and $x=g'^{-2}$. The
inequality is required for each of the factors polylogarithmic in $x$ that occur in the bounds
of the two-point witness theorem, each written $\operatorname{polylog}(x)$; $\check R(x)$ is the
size parameter $\check R_k$
of class components written as a function of the inverse coupling; and $c_q$ is the constant of
the construction of the two-point witness theorem with which that construction absorbs these
factors.
These inequalities govern the absorptions in the two-point witness theorem and every upper
bound on $\gamma$ imposed in this chapter. They rest on \cite{B-CMP119}.
\end{enumerate}
\end{definition}

\begin{lemma}[The reserve form of the island operation bound]\label{8.6:lem-reserve-island-bound}
Let $k$ be a level, let $X$ be a large-field region, let $U$ be a configuration for which the
normalised island operation $\mathbf{T}'_k(X,(U,0))1$ of
Definition~\ref{8.6:def-island-operation-bound} is defined, and let $\theta$ lie in the window
\eqref{8.6:eq-reserve-weight-b}. For an internal history $\mathfrak{h}$, write $\mathcal{Z}(\mathfrak{h})$ for
its large-field regions, $j(Z)$ for the step that created $Z\in\mathcal{Z}(\mathfrak{h})$, and
\begin{equation}\label{8.6:eq-reserve-island-bound-1}
  \Theta(\mathfrak{h})=\theta\sum_{Z\in\mathcal{Z}(\mathfrak{h})}p_0(g_{j(Z)}),
\end{equation}
where $p_0(\cdot)$ is the degree function of the coupling (Notation~\ref{0.2:not-glossary}),
so that, by Definition~\ref{8.6:def-reserve-weight}, the reserve weight \eqref{8.6:eq-reserve-weight-a}
of a set $\mathcal{S}$ of internal histories is
$w_\theta(\mathcal{S})=\sup_{\mathfrak{h}\in\mathcal{S}}e^{-\Theta(\mathfrak{h})}$. For such a set
$\mathcal{S}$, let $\mathbf{T}'_{k,\mathcal{S}}(X,(U,0))1$ denote the internal sum defining the
normalised island operation $\mathbf{T}'_k(X,(U,0))1$ of
Definition~\ref{8.6:def-island-operation-bound}, taken over the histories
$\mathfrak{h}\in\mathcal{S}$ only. Assume:
\begin{itemize}
\item[(a)] (per-history factorisation in the reserve statement at an $\mathbf{R}$-step) for every
internal history $\mathfrak{h}$ the contribution of $\mathfrak{h}$ to $\mathbf{T}'_k(X,(U,0))1$ is a
positive number, equal to $e^{-\Theta(\mathfrak{h})}\,e^{-\kappa^{(\theta)}(\mathfrak{h})}$ with
$\kappa^{(\theta)}(\mathfrak{h})$ real, an identity which defines the
$\theta$-stripped exponent $\kappa^{(\theta)}(\mathfrak{h})$ of $\mathfrak{h}$, and
$e^{-\kappa^{(\theta)}(\mathfrak{h})}\le\Phi_\theta(\mathfrak{h})$, where $\Phi_\theta(\mathfrak{h})$ is
the $\theta$-stripped majorant of $\mathfrak{h}$;
\item[(b)] (the $\theta$-stripped form of \cite[(1.79)--(1.89)]{B-CMP122b}) the sum of the
$\theta$-stripped majorants over all internal histories $\mathfrak{h}$, that is, over all internal
admissible sequences of that summation, obeys
\begin{equation}\label{8.6:eq-reserve-island-bound-2}
  \sum_{\mathfrak{h}}\Phi_\theta(\mathfrak{h})
  \le \exp\bigl(-(2-\theta)(1+\beta_0)^{-1}p_0(g_k)\bigr);
\end{equation}
here $\beta_0$ is Ba{\l}aban's exponent of \cite[(1.87), (1.89)]{B-CMP122b}, the letter entering the
window \eqref{8.6:eq-reserve-weight-b}.
\end{itemize}
Then for every set $\mathcal{S}$ of internal histories
\begin{equation}\label{8.6:eq-reserve-island-bound-a}
  \mathbf{T}'_{k,\mathcal{S}}(X,(U,0))1
  \le w_\theta(\mathcal{S})\,\exp\bigl(-(2-\theta)(1+\beta_0)^{-1}p_0(g_k)\bigr).
\end{equation}
\end{lemma}

Hypothesis (a) is the per-history factorisation established with the reserve statement at an
$\mathbf{R}$-step, and hypothesis (b) is the $\theta$-stripped form of the summation cited in it. In that
summation, \cite[\S1]{B-CMP122b} states that ``the summations over the admissible sequences can be
replaced by the factors $\exp O(1)(MR_j)^{-d}|Z_j|$'', the symbols in the quoted sentence being
Ba{\l}aban's. The reserve statement at an $\mathbf{R}$-step is a bound uniform in
the configuration; in \eqref{8.6:eq-reserve-island-bound-a} it is applied only to the internal sum of
the normalised operator $\mathbf{T}'_k$ of Definition~\ref{8.6:def-island-operation-bound}, that is, to a
quantity already divided by its normalising factor.

\begin{proof}
Fix a set $\mathcal{S}$ of internal histories. By definition, $\mathbf{T}'_{k,\mathcal{S}}(X,(U,0))1$
is the sum over $\mathfrak{h}\in\mathcal{S}$ of the contributions of the histories
$\mathfrak{h}$. By hypothesis (a), each contribution is
$e^{-\Theta(\mathfrak{h})}e^{-\kappa^{(\theta)}(\mathfrak{h})}$. Hence
\begin{equation*}
  \mathbf{T}'_{k,\mathcal{S}}(X,(U,0))1
  =\sum_{\mathfrak{h}\in\mathcal{S}}e^{-\Theta(\mathfrak{h})}\,e^{-\kappa^{(\theta)}(\mathfrak{h})}.
\end{equation*}

For $\mathfrak{h}\in\mathcal{S}$ we have
$e^{-\Theta(\mathfrak{h})}\le\sup_{\mathfrak{h}'\in\mathcal{S}}e^{-\Theta(\mathfrak{h}')}=w_\theta(\mathcal{S})$,
by \eqref{8.6:eq-reserve-island-bound-1} and the definition \eqref{8.6:eq-reserve-weight-a}. Each factor
$e^{-\kappa^{(\theta)}(\mathfrak{h})}$ is positive, so bounding $e^{-\Theta(\mathfrak{h})}$ by its supremum over
$\mathcal{S}$ term by term gives
\begin{equation*}
  \mathbf{T}'_{k,\mathcal{S}}(X,(U,0))1
  \le w_\theta(\mathcal{S})\sum_{\mathfrak{h}\in\mathcal{S}}e^{-\kappa^{(\theta)}(\mathfrak{h})}.
\end{equation*}

By hypothesis (a), $e^{-\kappa^{(\theta)}(\mathfrak{h})}\le\Phi_\theta(\mathfrak{h})$ for every
internal history $\mathfrak{h}$; in particular every $\Phi_\theta(\mathfrak{h})$ is positive. Since $\mathcal{S}$
is contained in the set of all internal histories and every $\Phi_\theta(\mathfrak{h})$ is positive, the sum
of the majorants over $\mathcal{S}$ is at most their sum over all internal histories. By hypothesis (b),
that sum is at most the right-hand side of
\eqref{8.6:eq-reserve-island-bound-2}. Therefore
\begin{align*}
  \mathbf{T}'_{k,\mathcal{S}}(X,(U,0))1
  &\le w_\theta(\mathcal{S})\sum_{\mathfrak{h}\in\mathcal{S}}\Phi_\theta(\mathfrak{h})
  \le w_\theta(\mathcal{S})\sum_{\mathfrak{h}}\Phi_\theta(\mathfrak{h})\\
  &\le w_\theta(\mathcal{S})\exp\bigl(-(2-\theta)(1+\beta_0)^{-1}p_0(g_k)\bigr),
\end{align*}
which is \eqref{8.6:eq-reserve-island-bound-a}.
\end{proof}

\begin{lemma}[Later steps do not read earlier healed islands. Stated; proof incomplete at the step marked $(\ast)$]\label{8.6:lem-healed-islands-unread}
Let $0\le k<K$. Write the level-$k$ term in the island representation of the prepared density
(Definition~\ref{8.6:def-island-representation}), with outer datum $\mathcal{O}_k$ and island
datum $\mathfrak{X}_k$.
\begin{enumerate}
\item[(a)] Step $k+1$ maps the level-$k$ terms to level-$(k+1)$ terms by the following operations:
the renormalisation transformation of step $k+1$; the insertion of decompositions of unity in the
new fluctuation variables and the new background variables; integrations; and the preparation for
the operation $\mathbf{R}$ together with healing. The result is a finite sum, over the new labels,
of positive linear operators applied to the level-$k$ term, followed by identities of functions.
\item[(b)] The following holds:
\begin{equation}\label{8.6:eq-healed-islands-unread-a}
\begin{gathered}
\text{the loci and forms of the decompositions of unity of step } k+1,\\
\text{the extension of the operations of the live components, and}\\
\text{the denominators of every later $\mathbf{R}$-quotient}\\
\text{and every later completed action}\\
\text{are determined by } \mathcal{O}_k \text{ and by the universal small-field action}\\
\text{on the current small-field region, and not by the healed islands of } \mathfrak{X}_k .
\end{gathered}
\end{equation}
Here the extension of the operations of the live components is the one of
\cite[(2.20)]{B-CMP119}. Moreover, the denominators of every later $\mathbf{R}$-quotient, and
every later completed action, are functionals of the universal small-field action on the current
small-field region; they do not read the factor
\begin{equation*}
\prod_a \mathbf{T}'(X_a)\, e^{\sum \mathcal{V}}
\end{equation*}
carried by an earlier step, that is, by a step preceding the later step whose $\mathbf{R}$-quotient
or completed action is considered. Here $X_a$ runs over the large-field regions of the term at that
earlier step $j$, $\mathbf{T}'(X_a)$ is the large-field operation $\mathbf{T}'_j(X_a)$ of step $j$
on $X_a$, the level subscript being suppressed in the display,
and $e^{\sum\mathcal{V}}$ is the exponential factor carried
with these regions in the island representation, the sum in the exponent being the one of
Definition~\ref{8.6:def-island-representation}.
\end{enumerate}
\end{lemma}

\begin{proof}
\emph{Part} (a). This is the one-step positivity representation of the renormalisation
transformation, in which one step acting on the level-$k$ terms is written as a finite sum, over
the new labels, of positive linear operators followed by identities of functions; only the
linearity and positivity of the individual operations enter into it.

\emph{Part} (b). In \cite[\S3]{B-CMP119} the level-$k$ term has the large-field region
$Z_k=\Lambda_k^c$, the complement of the term's small-field region $\Lambda_k$. For a term whose
large-field region is $Z_k$, step $k+1$ introduces new restrictions on the new fluctuation and
background fields of the step \cite[\S3]{B-CMP119}. Consequently the loci and the forms of the
decompositions of unity of step $k+1$ are fixed by the large-field data of the term.

By \cite[p.~391]{B-CMP122b}, the healed islands are no longer large-field regions after the step
that heals them. They therefore do not enter the large-field data from which the new restrictions
are placed, and these data constitute the outer datum $\mathcal{O}_k$ of
Definition~\ref{8.6:def-island-representation}. This proves that the
loci and forms of the decompositions of unity of step $k+1$, as asserted
in~\eqref{8.6:eq-healed-islands-unread-a}, are determined by $\mathcal{O}_k$ and not by the
healed islands of $\mathfrak{X}_k$.

For the extension \cite[(2.20)]{B-CMP119} of the operations of the live components, note that a
live component still carries its large-field factor at level $k$, so it is not a healed island of
$\mathfrak{X}_k$; it belongs to the large-field data of the term, which, by the preceding
paragraph, are the outer datum $\mathcal{O}_k$. By \cite[p.~177]{B-CMP122a}, the
operations so extended are determined by the small-field effective actions only, that is, by the
universal small-field action on the current small-field region. Hence the extension named in the
second clause of~\eqref{8.6:eq-healed-islands-unread-a} is determined by $\mathcal{O}_k$ and the
universal small-field action, and not by the healed islands of $\mathfrak{X}_k$.

$(\ast)$ It remains to verify, from \cite[\S3]{B-CMP119}, \cite[p.~177]{B-CMP122a} and
\cite[p.~391]{B-CMP122b}, the third clause of~\eqref{8.6:eq-healed-islands-unread-a} together
with the sentence following that display: the denominators of every later
$\mathbf{R}$-quotient, and every later completed action, are functionals of the universal
small-field action on the current small-field region, and do not read the carried factor
$\prod_a\mathbf{T}'(X_a)\,e^{\sum\mathcal{V}}$ of an earlier step.
\end{proof}

\begin{definition}[One-sided conditions on the auxiliary parameters]
\label{8.6:def-parameter-conditions}
Let $d=4$. Let $\beta$ be the exponent in the factor $e^{-(1+2\beta)\kappa D}$ of the large-field
weight described below, and let $\beta_0$ be Ba{\l}aban's exponent of
\cite[(1.87), (1.89)]{B-CMP122b}, where the example value $\beta_0=1/7$ is given; both are fixed
at the start of the order of choice stated below. Let
$\theta$ be the exponent of the reserve weight of Definition~\ref{8.6:def-reserve-weight}, and let
$g_k$, $x_k=g_k^{-2}$ and the degree function $p_0(\cdot)$ be as in
Notation~\ref{8.6:not-running-couplings}, so that $p_0(g_k)=A_0(\log x_k)^{p_0}$ with the fixed
exponent $p_0\ge 5r$. We put
\begin{equation*}
A_d:=5^d\bigl[(1+3^d)\log 2+2d\log 3\bigr],
\end{equation*}
the constant of the counting bound for families of cubes with a given anchor cube. Let $\alpha$
be the amplitude of the large-field weight, that is, the constant multiplying
$e^{-(1+2\beta)\kappa D}$ in the weight of a family $Y$ of cubes with a given anchor cube for
which $\lfloor d_\sharp(Y)\rfloor=D$; here $d_\sharp(Y)$ is the length, in units of $M$, of a
shortest tree graph meeting the cubes of $Y$. Let
$\check{R}_k$ denote the radius at level $k$ entering the left member of
\eqref{8.6:eq-parameter-conditions-b}, a quantity bounded by a polynomial of degree $r$ in
$\log x_k$, and let
$C_{\mathrm{rm}}$ be the constant of the same name in the statements of this section that bound
the live tail; it is fixed where those statements are stated and is not defined anew here. The
auxiliary parameters are subject to the following three conditions:
\begin{equation}\label{8.6:eq-parameter-conditions-a}
\tfrac12\,\beta\kappa\;\ge\;A_d+5^d+\log 2;
\end{equation}
\begin{equation}\label{8.6:eq-parameter-conditions-b}
\begin{gathered}
C_\delta(d)\,\alpha\,(101\check{R}_k)^d\;\le\;
\Bigl(\frac{2-\theta}{1+\beta_0}-\frac32\Bigr)\,p_0(g_k),
\qquad C_\delta(d):=4(3^d+1)e^{2A_d},\\
\text{for }\theta=0\text{ and for the exponent }\theta\text{ of
Definition~\ref{8.6:def-reserve-weight},}\\
\text{which lies in the window }\eqref{8.6:eq-reserve-weight-b};
\end{gathered}
\end{equation}
\begin{equation}\label{8.6:eq-parameter-conditions-c}
\bigl(\kappa_1 5^{-d}-C_{\mathrm{rm}}3^d\bigr)\,\check{R}_k^{\,d}\;\ge\;d\log 9+\log 2 .
\end{equation}
Condition \eqref{8.6:eq-parameter-conditions-c} enters only as a hypothesis of the statements of
this section that bound the live tail, in which the constant $C_{\mathrm{rm}}$ is fixed. The
window \eqref{8.6:eq-reserve-weight-b}, that is
$0<\theta<(1-3\beta_0)/2$, is a condition on the auxiliary exponent $\theta$ alone; it is non-empty
if and only if $\beta_0<1/3$.

Each of the three conditions is one-sided; none of them bounds $L$; none is a cap in the sense
of Definition~\ref{2.3:def-cap}; none fixes the value of a parameter. They are met in the
following order: $d$, $\beta$, $\beta_0$, $\theta$ are fixed; $\kappa$ is chosen after $M$,
subject to \eqref{8.6:eq-parameter-conditions-a}; $\kappa_1$ is chosen together with $M$,
subject to \eqref{8.6:eq-parameter-conditions-c}; $L$ occurs in none of the three conditions;
$\gamma$ is chosen last, subject
to \eqref{8.6:eq-parameter-conditions-b}, to two further conditions on $\gamma$ and to an
absorption condition, each stated in the lemma of this chapter that uses it. No
condition bounds $K$, $m$, $s$, $k$ or $L$, and no condition bounds a coupling from below.
\end{definition}

\begin{proof}
We verify the assertions made after the three displays: that each condition is one-sided, bounds
none of $L$, $K$, $m$, $s$, $k$, is not a cap and prescribes no value, and that the window
is non-empty exactly when $\beta_0<1/3$.

Condition \eqref{8.6:eq-parameter-conditions-a} is a lower bound on $\kappa$ whose right member
depends on $d$ alone and whose coefficient $\tfrac12\beta$ is fixed before $\kappa$; it is a
lower bound of the kind required in \cite{B-CMP122b}, where $\kappa$ is taken large enough.

In \eqref{8.6:eq-parameter-conditions-b} the coefficient of $p_0(g_k)$ is positive exactly when
$(2-\theta)(1+\beta_0)^{-1}>3/2$. Since $1+\beta_0>0$, this is $2-\theta>\tfrac32(1+\beta_0)$,
and
\begin{equation*}
2-\tfrac32(1+\beta_0)=\frac{1-3\beta_0}{2},
\end{equation*}
so this coefficient is positive if and only if $\theta<(1-3\beta_0)/2$; since $C_\delta(d)>0$ and
the amplitude $\alpha$ and the radius $\check{R}_k$ are positive, the left member of
\eqref{8.6:eq-parameter-conditions-b} is positive, and the condition can be satisfied only if
$\theta<(1-3\beta_0)/2$. For $\theta=0$ this is
$\beta_0<1/3$, which is also the condition for the window \eqref{8.6:eq-reserve-weight-b} to be
non-empty. The condition is of the same kind as the requirement of \cite{B-CMP122b} that
$g_k$ be so small that
\begin{equation*}
-2(1+\beta_0)^{-1}p_0(g_k)+1+2\kappa(100R_k)^d\;\le\;-\tfrac32\,p_0(g_k),
\end{equation*}
where $R_k$ denotes the radius at level $k$ used there; after adding
$2(1+\beta_0)^{-1}p_0(g_k)$ to both members, this reads
\begin{equation*}
1+2\kappa(100R_k)^d\;\le\;\Bigl(\frac{2}{1+\beta_0}-\frac32\Bigr)p_0(g_k),
\end{equation*}
with the coefficient of \eqref{8.6:eq-parameter-conditions-b} at $\theta=0$. The left member of
\eqref{8.6:eq-parameter-conditions-b} is at most $\alpha$ times a polynomial of degree $dr$ in
$\log x_k$, while its right member is a positive multiple of $A_0(\log x_k)^{p_0}$, of degree
$p_0\ge 5r>dr$ because $d=4$. Hence \eqref{8.6:eq-parameter-conditions-b} is a smallness
condition on the coupling, imposed through an upper bound on $\gamma$, the parameter chosen last
in the order stated in the definition above;
it bounds no coupling from below. It is implied by the smallness condition on $g_k$ that
accompanies the counting lemmas for families of cubes with a given anchor cube, a condition of the
same kind as the requirement displayed above, whenever $C_\delta(d)\alpha\le c_1(d)\kappa$, that
is, under an upper bound on $\alpha$; here $c_1(d)$ is the constant of that condition and depends
on $d$ only.

Condition \eqref{8.6:eq-parameter-conditions-c} is a lower bound on Ba{\l}aban's $\kappa_1$ of
\cite{B-CMP122b} relative to the
constant $C_{\mathrm{rm}}$: its right member is fixed by $d$, and its left member is
non-decreasing in $\kappa_1$, since $\check{R}_k\ge0$. It can be met within the
hypotheses of the correlation theorem: a corollary established within the proof of that theorem
gives, for a large-field region $Z$ at step $n$,
\begin{equation*}
\kappa_n(Z)\;\ge\;\sqrt{d}\,c_n|Z|\;\ge\;C^{\flat}_{\mathfrak{d}}M^d\,d_n(Z),
\end{equation*}
where $\kappa_n(Z)$ is the decay exponent attached to $Z$ at step $n$, $|Z|$ is the size of $Z$
as measured in that corollary, $d_n(Z)$ is the quantity $d_j(X)$ of
Notation~\ref{0.2:not-glossary} with $j=n$ and $X=Z$, and $c_n$, $C^{\flat}_{\mathfrak{d}}$
are the constants of that corollary. Thus a lower bound with the constant
$C^{\flat}_{\mathfrak{d}}M^d$, growing with $M$, is
available; this is why $\kappa_1$ is chosen together with $M$.

Finally, \eqref{8.6:eq-parameter-conditions-a} and \eqref{8.6:eq-parameter-conditions-c} are
lower bounds on $\kappa$ and $\kappa_1$ and \eqref{8.6:eq-parameter-conditions-b} is an upper
bound on $\gamma$; none of them contains any of $L$, $K$, $m$, $s$, $k$ as a letter to be
bounded, and
each is met by choosing its letter at the place of the order stated above. None bounds one of
the letters listed in Definition~\ref{2.3:def-cap} from above, and none prescribes a value.
\end{proof}

\begin{lemma}[Coupling comparison and summability of the healing factors]\label{8.6:lem-coupling-comparison}
Let $x_k=g_k^{-2}$ $(0\le k\le K)$ be the running inverse couplings, let
$p_0(g')=A_0(\log g'^{-2})^{p_0}$ be the degree function, and let $\lambda^{\sharp}$ and $c_5$ be
the constants of the flow inequality \eqref{8.6:eq-running-couplings-a}, all as in
Notation~\ref{8.6:not-running-couplings}. Put $X=g^{-2}$.
\begin{itemize}
\item[(a)] \emph{Coupling comparison.} Suppose that every $x\ge X$ satisfies
\begin{equation*}
A_0p_0(\log x)^{p_0-1}\cdot\frac{2c_5}{x}\le 1
\qquad\text{and}\qquad x\ge 2c_5 .
\end{equation*}
Then
\begin{equation}\label{8.6:eq-coupling-comparison-a}
p_0(g_i)\ \ge\ p_0(g_j)-1\qquad\text{for all } 0\le i\le j\le K .
\end{equation}
\item[(b)] \emph{The integral test.} Suppose that
\begin{equation*}
\tfrac32A_0\bigl(\log(X/2)\bigr)^{p_0-1}\ge 2,\qquad \log(X/2)\ge 1,\qquad X\ge 2c_5 .
\end{equation*}
Let $K'$ be an integer with $0\le K'\le K$. Then
\begin{equation}\label{8.6:eq-coupling-comparison-b}
S_*:=\sup_{0\le k''\le K'}\ \sum_{k=k''}^{K'}e^{-\frac32p_0(g_k)}\ \le\ 1+\frac{1}{\lambda^{\sharp}} .
\end{equation}
\end{itemize}
\end{lemma}

Each hypothesis is a lower bound on $X=g^{-2}$ and holds for all sufficiently large $X$. The
bound \eqref{8.6:eq-coupling-comparison-b} involves $\lambda^{\sharp}$ only; it does not depend
on $K$, on $K'$, or on the torus $\mathbb{T}_m$.

\begin{proof}
Both parts use the flow inequality \eqref{8.6:eq-running-couplings-a}, namely
$x_i-x_j\ge\lambda^{\sharp}(j-i)-c_5$ for $0\le i\le j\le K$. They also use the bound
$x_k\ge X$ for every $k$, recorded in Notation~\ref{8.6:not-running-couplings}.

\emph{Part} (a). Let $i\le j\le K$. Since $\lambda^{\sharp}(j-i)\ge0$, the flow inequality gives
\begin{equation}\label{8.6:eq-coupling-comparison-1}
x_i\ \ge\ x_j-c_5 .
\end{equation}
The function $t\mapsto t^{p_0}$ increases, and so does the logarithm. Hence
\eqref{8.6:eq-coupling-comparison-1} yields
\begin{equation*}
p_0(g_i)=A_0(\log x_i)^{p_0}\ge A_0\bigl(\log(x_j-c_5)\bigr)^{p_0} .
\end{equation*}
It remains to compare the right-hand side with $p_0(g_j)=A_0(\log x_j)^{p_0}$.

Apply the mean value theorem to $t\mapsto t^{p_0}$ on the interval
$[\log(x_j-c_5),\log x_j]$. On that interval the derivative $p_0t^{p_0-1}$ is at most
$p_0(\log x_j)^{p_0-1}$. This gives the first inequality in
\begin{equation}\label{8.6:eq-coupling-comparison-2}
\begin{aligned}
(\log x_j)^{p_0}-\bigl(\log(x_j-c_5)\bigr)^{p_0}
&\le p_0(\log x_j)^{p_0-1}\log\frac{x_j}{x_j-c_5}\\
&\le p_0(\log x_j)^{p_0-1}\cdot\frac{2c_5}{x_j} .
\end{aligned}
\end{equation}
For the second inequality, put $s=c_5/x_j$. Since $x_j\ge X$, the hypothesis at $x=x_j$ gives
$x_j\ge2c_5$, that is $0\le s\le\frac12$. Therefore
\begin{equation*}
\log\frac{x_j}{x_j-c_5}=-\log(1-s)\le\frac{s}{1-s}\le 2s=\frac{2c_5}{x_j} .
\end{equation*}

Multiply \eqref{8.6:eq-coupling-comparison-2} by $A_0$ and use the first hypothesis at
$x=x_j\ge X$. This gives
\begin{equation*}
p_0(g_j)-A_0\bigl(\log(x_j-c_5)\bigr)^{p_0}\le A_0p_0(\log x_j)^{p_0-1}\cdot\frac{2c_5}{x_j}\le1 .
\end{equation*}
Combined with the lower bound on $p_0(g_i)$ above, this is
\eqref{8.6:eq-coupling-comparison-a}.

\emph{Part} (b). For $n\ge0$ put $y_n:=X/2+\lambda^{\sharp}n$.

\emph{Lower bound on the couplings.} Let $0\le k\le K'$. The flow inequality with $i=k$,
$j=K'$ gives the first inequality in
\begin{equation}\label{8.6:eq-coupling-comparison-3}
x_k\ \ge\ x_{K'}+\lambda^{\sharp}(K'-k)-c_5\ \ge\ \frac X2+\lambda^{\sharp}(K'-k)=y_{K'-k} .
\end{equation}
For the second, use $x_{K'}\ge X$ together with $c_5\le X/2$, which is the hypothesis
$X\ge2c_5$.

\emph{Monotonicity.} Define
\begin{equation*}
f(u):=e^{-\frac32A_0(\log u)^{p_0}}\qquad (u\ge1).
\end{equation*}
This function decreases, since $u\mapsto(\log u)^{p_0}$ increases on $[1,\infty)$ and $A_0>0$.
By the hypothesis $\log(X/2)\ge1$, every $y_n$ satisfies $y_n\ge y_0=X/2\ge e$. Hence
\eqref{8.6:eq-coupling-comparison-3} gives
$e^{-\frac32p_0(g_k)}=f(x_k)\le f(y_{K'-k})$ for $0\le k\le K'$.

\emph{Reduction to a series.} For every $0\le k''\le K'$, the substitution $n=K'-k$ gives
\begin{equation*}
\sum_{k=k''}^{K'}e^{-\frac32p_0(g_k)}\le\sum_{n=0}^{K'-k''}f(y_n)\le\sum_{n\ge0}f(y_n) .
\end{equation*}
The sum runs over the levels $k\le K'$, that is, toward $K'$ from below. Along this direction,
read from $K'$ downward, the flow inequality makes the inverse couplings grow at least
linearly. This growth is what turns the sum into the series above.

\emph{Comparison with an integral.} Let $n\ge1$. The interval $[y_{n-1},y_n]$ has length
$\lambda^{\sharp}$, and $f$ decreases. Therefore
\begin{equation*}
f(y_n)\le(\lambda^{\sharp})^{-1}\int_{y_{n-1}}^{y_n}f(u)\,du .
\end{equation*}
Summing over $n\ge1$ and using $f(y_0)\le1$, we get
\begin{equation}\label{8.6:eq-coupling-comparison-4}
\sum_{n\ge0}f(y_n)\ \le\ f(y_0)+\frac{1}{\lambda^{\sharp}}\int_{y_0}^{\infty}f(u)\,du
\ \le\ 1+\frac{1}{\lambda^{\sharp}}\int_{y_0}^{\infty}f(u)\,du .
\end{equation}

\emph{Evaluation of the integral.} With the substitution $u=e^v$,
\begin{equation*}
\int_{y_0}^{\infty}f(u)\,du=\int_{\log y_0}^{\infty}e^{-\frac32A_0v^{p_0}+v}\,dv .
\end{equation*}
Put $a:=A_0(\log y_0)^{p_0-1}$. Since $y_0=X/2$, the first hypothesis reads $\frac32a\ge2$,
that is $\frac34a\ge1$.

Let $v\ge\log y_0$. Then $v^{p_0-1}\ge(\log y_0)^{p_0-1}$, and therefore
\begin{equation*}
\tfrac32A_0v^{p_0-1}-1\ \ge\ \tfrac32a-1\ \ge\ \tfrac32a-\tfrac34a=\tfrac34a\ \ge\ 1 .
\end{equation*}
Multiplying by $v\ge0$ gives
\begin{equation*}
-\tfrac32A_0v^{p_0}+v=-v\bigl(\tfrac32A_0v^{p_0-1}-1\bigr)\le-\tfrac34a\,v .
\end{equation*}
Consequently
\begin{equation}\label{8.6:eq-coupling-comparison-5}
\int_{y_0}^{\infty}f(u)\,du\ \le\ \int_{\log y_0}^{\infty}e^{-\frac34a\,v}\,dv
=\frac{e^{-\frac34a\log y_0}}{\frac34a}\ \le\ 1 .
\end{equation}
The last inequality holds because $\frac34a\ge1$ and $\log y_0=\log(X/2)\ge1>0$.

\emph{Conclusion.} Inserting \eqref{8.6:eq-coupling-comparison-5} into
\eqref{8.6:eq-coupling-comparison-4} bounds the sum over $k$ by $1+1/\lambda^{\sharp}$. This
holds for every $0\le k''\le K'$, and taking the supremum over $k''$ gives
\eqref{8.6:eq-coupling-comparison-b}. No quantity in this bound depends on $K$, on $K'$ or
on $m$.
\end{proof}

\begin{lemma}[Removal domination for a healed island]\label{8.6:lem-healed-removal}
Fix a step $k$, an outer datum $\mathfrak{o}$ of the island representation
\eqref{8.6:eq-island-representation-a} of Definition~\ref{8.6:def-island-representation}, all
outer variables, and $V_k$. Every assertion of the lemma is pointwise in these.

For an island $X$ and a set $\mathcal{S}$ of internal histories of $X$ let
$\mathbf{T}'_{k,\mathcal{S}}(X):=\sum_{\mathfrak{h}\in\mathcal{S}}\mathbf{T}'_{k,\mathfrak{h}}(X)$
be the normalised operation of $X$ with its sum over internal histories restricted to
$\mathcal{S}$. Write $\mathbf{T}'_{k,\mathcal{S}}(X,(U,0))$ for the same operation at the
reference configuration $(U,0)$, as in Definition~\ref{8.6:def-island-operation-bound}. Put
$t_X(\mathcal{S}):=\mathbf{T}'_{k,\mathcal{S}}(X,(U,0))1$.

Let $\mathfrak{c}$ be a level-$k$ cell. If some island of an admissible family $\mathfrak{X}$
contains $\mathfrak{c}$, that island is unique, because distinct islands of an admissible family
share no point (condition (c2) of Definition~\ref{8.6:def-island-representation}). Its index is
denoted $a(\mathfrak{c})$. Put
\begin{equation}\label{8.6:eq-healed-removal-a}
\begin{gathered}
E_{\mathfrak{c}}:=\{\mathfrak{X}\ \text{admissible under}\ \mathfrak{o}:\ \text{some island}\
X_a\ \text{of}\ \mathfrak{X}\ \text{contains}\ \mathfrak{c}\},\\
\operatorname{red}_{\mathfrak{c}}(\mathfrak{X}):=\mathfrak{X}\setminus\{X_{a(\mathfrak{c})}\}
\ \ (\mathfrak{X}\in E_{\mathfrak{c}}),\qquad
\operatorname{red}_{\mathfrak{c}}(\mathfrak{X}):=\mathfrak{X}\ \ (\mathfrak{X}\notin E_{\mathfrak{c}}),\\
W^{\mathcal{S},\mathfrak{c}}_{\mathfrak{o}}(\mathfrak{X}):=W_{\mathfrak{o}}(\mathfrak{X})\
\text{with}\ \mathbf{T}'_k(X_{a(\mathfrak{c})})\ \text{replaced by}\
\mathbf{T}'_{k,\mathcal{S}}(X_{a(\mathfrak{c})})\ \ (\mathfrak{X}\in E_{\mathfrak{c}}).
\end{gathered}
\end{equation}

Let $\theta=0$, or let $\theta$ lie in the range of Lemma~\ref{8.6:lem-reserve-island-bound}.
Assume the following.
\begin{itemize}
\item \emph{Positivity.} The positivity form \eqref{8.6:eq-island-operation-bound-a} at $(U,0)$
of \cite[(1.73)]{B-CMP122b} holds: for every island $X$ admissible alone, every set
$\mathcal{S}$ of its internal histories and every non-negative bounded function $\Phi$ of the
internal variables of $X$ and of other variables, $\mathbf{T}'_{k,\mathcal{S}}(X)\Phi$ lies
between $0$ and $t_X(\mathcal{S})\sup_X\Phi$ pointwise in the other variables, the supremum
being over the internal variables of $X$ on the supports of the characteristic functions of
$\mathbf{T}'_{k,\mathcal{S}}(X)$.
\item \emph{Island bound.} If $\theta=0$, the island bound \eqref{8.6:eq-island-operation-bound-b},
which is \cite[(1.89)]{B-CMP122b}, holds for every island admissible alone. If $\theta$ lies in
the range of Lemma~\ref{8.6:lem-reserve-island-bound}, its reserve bound
\eqref{8.6:eq-reserve-island-bound-a} holds for every $\mathcal{S}$.
\item \emph{Boundary terms.} The bound \eqref{8.6:eq-island-operation-bound-c}, the form of
\cite[(1.68)]{B-CMP122b} used here, holds for the terms $V_{\mathfrak{X}}(Y)$ with
$Y\in\mathcal{Y}(\mathfrak{X})$ and $Y\setminus Z(\mathfrak{X})^\sharp\neq\emptyset$, together with
$V_{\mathfrak{X}}(Y)=0$ when $d_\sharp(Y)=\infty$.
\item \emph{Locality.} Assume \eqref{8.6:eq-island-operation-bound-d}. Among internal variables of
islands, $V_{\mathfrak{X}}(Y)$ depends only on those of the islands of $\mathfrak{X}$ contained in
$Y$, together with their histories. Further, $V_{\mathfrak{X}}(Y)=V_{\mathfrak{X}'}(Y)$ whenever
$\mathfrak{X}$ and $\mathfrak{X}'$ have the same islands inside $Y$ and
$Y\in\mathcal{Y}(\mathfrak{X})\cap\mathcal{Y}(\mathfrak{X}')$.
\item \emph{Island representation.} Assume the representation
\eqref{8.6:eq-island-representation-a}, with $W_{\mathfrak{o}}(\mathfrak{X})$ as there, the
admissibility conditions (c1)--(c3) and the domain condition (c4) of
Definition~\ref{8.6:def-island-representation}; each $\mathbf{O}_{\mathfrak{o},\mathfrak{X}}$ is a
positive linear operator. Assume also the monotonicity of the common prefactor,
Lemma~\ref{8.6:lem-prefactor-monotone}, whose proof is incomplete at one step: for admissible
$\mathfrak{X}\supset\mathfrak{X}'$ one has
$\mathbf{O}_{\mathfrak{o},\mathfrak{X}}\le\mathbf{O}_{\mathfrak{o},\mathfrak{X}'}$ as positive
operators.
\item \emph{Box condition.} Every island is a box of level-$k$ cells, each of whose $d$ side
lengths is at most $100$ cells.
\item \emph{Strip condition.} Let $X$ be an island of an admissible family. The strip $X^\sharp$
consists of $N(X^\sharp)\le(101\check{R}_k)^d$ cubes. A cube touching $X^\sharp$ but not contained
in it lies in no other strip. Finally, whenever $\mathfrak{X}'\sqcup\{X\}$ is admissible, so that
$X$ shares no point with the islands of $\mathfrak{X}'$ and with $|\mathfrak{o}|$ by (c2) and (c3),
the strip $X^\sharp$ shares no point with $Z(\mathfrak{X}')^\sharp$.
\item \emph{Domain counting.} The domain-counting lemma holds: for every admissible
$\mathfrak{X}$, every cube $q$ and every integer $D\ge0$, the number of
$Y\in\mathcal{Y}(\mathfrak{X})$ having $q$ as an anchor cube and with
$\lfloor d_\sharp(Y)\rfloor=D$ is at most $2e^{A_d(D+2)}$, where
$A_d=5^d[(1+3^d)\log2+2d\log3]$ is the constant appearing in
\eqref{8.6:eq-parameter-conditions-a}. Moreover, the shortest tree graph $\Gamma_Y$ of a domain
$Y\in\mathcal{Y}(\mathfrak{X})$ meets every cube of $Y\setminus Z(\mathfrak{X})^\sharp$.
\item \emph{Parameters.} The conditions \eqref{8.6:eq-parameter-conditions-a} and
\eqref{8.6:eq-parameter-conditions-b} hold, the latter at the chosen $\theta$. Here
$C_\delta(d)$ denotes the constant of \eqref{8.6:eq-parameter-conditions-b}.
\end{itemize}
Put
\begin{equation}\label{8.6:eq-healed-removal-b}
\epsilon^{\mathrm{rm}}_k(\mathcal{S}):=K_{\mathrm{rm}}\,\varpi_\theta(\mathcal{S})\,
\exp\bigl(-\tfrac32p_0(g_k)\bigr),\qquad K_{\mathrm{rm}}:=10^{4d},
\end{equation}
where $\varpi_\theta$ is the reserve weight \eqref{8.6:eq-reserve-weight-a} and $\varpi_0\equiv1$.
We call $\epsilon^{\mathrm{rm}}_k(\mathcal{S})$ the removal factor and $K_{\mathrm{rm}}$ its entropy
constant.
Then the following two statements hold.

\emph{Single-island domination.} Let $\mathfrak{X}'\notin E_{\mathfrak{c}}$ be admissible, and let
$X\ni\mathfrak{c}$ be an island with $\mathfrak{X}'\sqcup\{X\}$ admissible. Then
\begin{equation}\label{8.6:eq-healed-removal-c}
W^{\mathcal{S},\mathfrak{c}}_{\mathfrak{o}}(\mathfrak{X}'\sqcup\{X\})\le
t_X(\mathcal{S})\,e^{\delta_X}\,W_{\mathfrak{o}}(\mathfrak{X}').
\end{equation}
Here $\delta_X\ge0$ is a number depending only on $X$, $\mathfrak{X}'$ and the boundary terms,
which satisfies
\begin{equation}\label{8.6:eq-healed-removal-e}
\delta_X\le C_\delta(d)\,\alpha\,N(X^\sharp)\le C_\delta(d)\,\alpha\,(101\check{R}_k)^d .
\end{equation}
Both sides of \eqref{8.6:eq-healed-removal-c} are non-negative functions of $V_k$ and of the
outer variables, the internal variables of every island being integrated by its own operation.

\emph{Summed domination.} Call an event $B'$ of island families \emph{$\mathfrak{c}$-blind} if
$1_{B'}(\mathfrak{X})=1_{B'}(\operatorname{red}_{\mathfrak{c}}(\mathfrak{X}))$ for every admissible
$\mathfrak{X}$. For every $\mathfrak{c}$-blind event $B'$,
\begin{equation}\label{8.6:eq-healed-removal-d}
\begin{aligned}
\sum_{\mathfrak{X}\in E_{\mathfrak{c}}\cap B'}\mathbf{O}_{\mathfrak{o},\mathfrak{X}}
\bigl[W^{\mathcal{S},\mathfrak{c}}_{\mathfrak{o}}(\mathfrak{X})\bigr]
&\le\epsilon^{\mathrm{rm}}_k(\mathcal{S})\sum_{\mathfrak{X}'\in B'\setminus E_{\mathfrak{c}}}
\mathbf{O}_{\mathfrak{o},\mathfrak{X}'}\bigl[W_{\mathfrak{o}}(\mathfrak{X}')\bigr]\\
&\le\epsilon^{\mathrm{rm}}_k(\mathcal{S})\sum_{\mathfrak{X}\in B'}
\mathbf{O}_{\mathfrak{o},\mathfrak{X}}\bigl[W_{\mathfrak{o}}(\mathfrak{X})\bigr].
\end{aligned}
\end{equation}
These are inequalities between non-negative functions of $V_k$, all internal variables being
integrated.

Define the $(\mathfrak{o},\mathfrak{c},\mathcal{S})$-pinned part of $\rho_{k,\mathfrak{o}}$ to be
the left member of \eqref{8.6:eq-healed-removal-d} with $B'$ the set of all admissible families.
This part is at most $\epsilon^{\mathrm{rm}}_k(\mathcal{S})\,\rho_{k,\mathfrak{o}}$, pointwise in
$V_k$.

If $\mathfrak{c}$ shares a point with $|\mathfrak{o}|$, then by (c3) no island contains
$\mathfrak{c}$. Hence $E_{\mathfrak{c}}=\emptyset$ and the left members vanish.
\end{lemma}

\begin{proof}
Fix $\mathfrak{X}'$ and $X$ as in the single-island domination and put
$\mathfrak{X}:=\mathfrak{X}'\sqcup\{X\}$. The proof has four steps.

\emph{Step (i): the exponents.} Put
\begin{equation*}
\mathcal{E}_1:=\sum_{Y\in\mathcal{Y}(\mathfrak{X})}V_{\mathfrak{X}}(Y),\qquad
\mathcal{E}_0:=\sum_{Y\in\mathcal{Y}(\mathfrak{X}')}V_{\mathfrak{X}'}(Y).
\end{equation*}
We have $Z(\mathfrak{X})^\sharp=Z(\mathfrak{X}')^\sharp\sqcup X^\sharp$. Moreover $X^\sharp$ shares
no point with $Z(\mathfrak{X}')^\sharp$, by (c2), (c3) and the strip condition.

First consider a domain $Y\in\mathcal{Y}(\mathfrak{X})$ that meets $X^\sharp$. By (c4),
$Y\cap Z(\mathfrak{X})^\sharp$ is a union of whole strips, so $Y$ contains $X^\sharp$.

Next consider a domain $Y$ that does not meet $X^\sharp$. Since
$Z(\mathfrak{X})^\sharp=Z(\mathfrak{X}')^\sharp\sqcup X^\sharp$, one has
$Y\cap Z(\mathfrak{X})^\sharp=Y\cap Z(\mathfrak{X}')^\sharp$ and
$Y\setminus Z(\mathfrak{X})^\sharp=Y\setminus Z(\mathfrak{X}')^\sharp$. Hence $Y$ belongs
to $\mathcal{Y}(\mathfrak{X})$ if and only if it belongs to $\mathcal{Y}(\mathfrak{X}')$. In that
case $X\not\subset Y$, so the islands of $\mathfrak{X}$ and of $\mathfrak{X}'$ inside $Y$ are the
same. The locality hypothesis \eqref{8.6:eq-island-operation-bound-d} then gives
$V_{\mathfrak{X}}(Y)=V_{\mathfrak{X}'}(Y)$.

Put
\begin{align*}
\mathcal{E}^{\mathrm{near}}_1&:=\sum_{Y\in\mathcal{Y}(\mathfrak{X}):\,Y\supseteq X^\sharp}
V_{\mathfrak{X}}(Y),\qquad
\mathcal{E}^{\mathrm{near}}_0:=\sum_{Y\in\mathcal{Y}(\mathfrak{X}'):\,Y\cap X^\sharp\neq\emptyset}
V_{\mathfrak{X}'}(Y),\\
\mathcal{E}^{\mathrm{far}}&:=\sum_{Y\in\mathcal{Y}(\mathfrak{X}'):\,Y\cap X^\sharp=\emptyset}
V_{\mathfrak{X}'}(Y).
\end{align*}
By the two facts above,
\begin{equation}\label{8.6:eq-healed-removal-1}
\mathcal{E}_1=\mathcal{E}^{\mathrm{far}}+\mathcal{E}^{\mathrm{near}}_1,\qquad
\mathcal{E}_0=\mathcal{E}^{\mathrm{far}}+\mathcal{E}^{\mathrm{near}}_0 .
\end{equation}
By \eqref{8.6:eq-island-operation-bound-d}, only the terms $V_{\mathfrak{X}}(Y)$ with
$X\subset Y$ can depend on the internal variables of $X$. The domains in
$\mathcal{E}^{\mathrm{far}}$ and in
$\mathcal{E}^{\mathrm{near}}_0$ are domains of $\mathfrak{X}'$, which has no island $X$. Hence
$\mathcal{E}^{\mathrm{far}}$ and $\mathcal{E}^{\mathrm{near}}_0$ do not depend on the internal
variables of $X$.

Now define
\begin{equation*}
\delta_X:=\sup|\mathcal{E}^{\mathrm{near}}_1|+\sup|\mathcal{E}^{\mathrm{near}}_0|,
\end{equation*}
the suprema being over all variables on the supports. We bound $\delta_X$ by attaching to each
domain a cube of $Y\setminus Z(\cdot)^\sharp$ lying in or touching $X^\sharp$.

Let $Y\in\mathcal{Y}(\mathfrak{X})$ with $Y\supseteq X^\sharp$. By (c4) the set
$Y\setminus Z(\mathfrak{X})^\sharp$ is nonempty. Take a chain of touching cubes of $Y$ from a cube
of $X^\sharp$ to a cube of $Y\setminus Z(\mathfrak{X})^\sharp$. The first cube of the chain outside
$X^\sharp$ touches $X^\sharp$. By the strip condition it lies in no other strip, so it lies off
$Z(\mathfrak{X})^\sharp$. Thus $Y$ contains a cube of $Y\setminus Z(\mathfrak{X})^\sharp$ touching
$X^\sharp$, and there are at most $3^dN(X^\sharp)$ such cubes.

Let $Y\in\mathcal{Y}(\mathfrak{X}')$ meet $X^\sharp$. Then $Y$ contains a cube of $X^\sharp$.
That cube lies off $Z(\mathfrak{X}')^\sharp$, because $X^\sharp$ shares no point with
$Z(\mathfrak{X}')^\sharp$. There are at most $N(X^\sharp)$ such cubes.

In both cases the attached cube lies in $Y\setminus Z(\cdot)^\sharp$. In particular this set is
nonempty, so the bound \eqref{8.6:eq-island-operation-bound-c} applies to $V(Y)$. The
domain-counting lemma is stated at an anchor cube of $Y$, but its proof uses of the anchor cube
only that the shortest tree graph $\Gamma_Y$ of $Y$ meets it; since $\Gamma_Y$ meets every cube of
$Y\setminus Z(\cdot)^\sharp$, the count holds at the attached cube: the domains attached to a given
cube with $\lfloor d_\sharp(Y)\rfloor=D$ number at most $2e^{A_d(D+2)}$.
Domains with $d_\sharp(Y)=\infty$ contribute nothing, since $V(Y)=0$ for them. For the others,
\eqref{8.6:eq-island-operation-bound-c} bounds each term by
$\alpha e^{-(1+2\beta)\kappa D}$, because $d_\sharp(Y)\ge D$.

Summing over the at most $(3^d+1)N(X^\sharp)$ cubes and over $D$, and writing
$2e^{A_d(D+2)}=2e^{2A_d}e^{A_dD}$, we get
\begin{equation}\label{8.6:eq-healed-removal-2}
\begin{aligned}
\delta_X&\le(3^d+1)N(X^\sharp)\sum_{D\ge0}2e^{A_d(D+2)}\,\alpha\,e^{-(1+2\beta)\kappa D}\\
&\le(3^d+1)N(X^\sharp)\,2e^{2A_d}\alpha\sum_{D\ge0}e^{-((1+2\beta)\kappa-A_d)D}\\
&\le4(3^d+1)e^{2A_d}\,\alpha\,N(X^\sharp)=C_\delta(d)\,\alpha\,N(X^\sharp).
\end{aligned}
\end{equation}
The last inequality uses the following estimate. Since $1+2\beta\ge\frac12\beta$, the condition
\eqref{8.6:eq-parameter-conditions-a} (which reads
$\frac12\beta\kappa\ge A_d+5^d+\log2$) gives
\begin{equation*}
(1+2\beta)\kappa-A_d\ge\tfrac12\beta\kappa-A_d\ge5^d+\log2>\log2 .
\end{equation*}
Hence the sum over $D$ is a geometric series with ratio less than $\frac12$, and it is at most
$2$. The constant $4(3^d+1)e^{2A_d}$ is the constant $C_\delta(d)$ of
\eqref{8.6:eq-parameter-conditions-b}. Finally $N(X^\sharp)\le(101\check{R}_k)^d$ by the strip
condition. This proves \eqref{8.6:eq-healed-removal-e}.

\emph{Step (ii): positivity.} Fix the values of all variables other than the internal variables of
$X$. By \eqref{8.6:eq-healed-removal-1},
$e^{\mathcal{E}_1}=e^{\mathcal{E}^{\mathrm{far}}}\,e^{\mathcal{E}^{\mathrm{near}}_1}$. By step (i),
$e^{\mathcal{E}^{\mathrm{far}}}$ is constant in the variables of $X$. The operation
$\mathbf{T}'_{k,\mathcal{S}}(X)$ acts on the internal variables of $X$ and is the identity on the
rest. By linearity it therefore commutes with multiplication by $e^{\mathcal{E}^{\mathrm{far}}}$.

Apply the positivity hypothesis \eqref{8.6:eq-island-operation-bound-a} with
$\Phi:=e^{\mathcal{E}^{\mathrm{near}}_1}$. This $\Phi$ is non-negative, and it is bounded because
$|\mathcal{E}^{\mathrm{near}}_1|\le\delta_X$. We obtain
\begin{equation}\label{8.6:eq-healed-removal-3}
\begin{aligned}
\mathbf{T}'_{k,\mathcal{S}}(X)\bigl(e^{\mathcal{E}_1}\bigr)
&=e^{\mathcal{E}^{\mathrm{far}}}\,\mathbf{T}'_{k,\mathcal{S}}(X)
\bigl(e^{\mathcal{E}^{\mathrm{near}}_1}\bigr)\\
&\le e^{\mathcal{E}^{\mathrm{far}}}\,e^{\sup\mathcal{E}^{\mathrm{near}}_1}\,t_X(\mathcal{S})\\
&\le t_X(\mathcal{S})\,e^{\delta_X}\,e^{\mathcal{E}_0}.
\end{aligned}
\end{equation}
In the last step we used two facts. First, by \eqref{8.6:eq-healed-removal-1},
$\mathcal{E}^{\mathrm{far}}=\mathcal{E}_0-\mathcal{E}^{\mathrm{near}}_0\le
\mathcal{E}_0+\sup|\mathcal{E}^{\mathrm{near}}_0|$. Second,
$\sup\mathcal{E}^{\mathrm{near}}_1\le\sup|\mathcal{E}^{\mathrm{near}}_1|$. The sum of these two
suprema is $\delta_X$.

\emph{Step (iii): the other islands.} The operator $\prod_{b\in\mathfrak{X}'}\mathbf{T}'_k(X_b)$
has three relevant properties. It is positive. It acts on variables disjoint from those of $X$. It
commutes with $\mathbf{T}'_{k,\mathcal{S}}(X)$ \cite[(2.19)]{B-CMP119}.

Apply it to both members of \eqref{8.6:eq-healed-removal-3}. The right member is the non-negative
function $e^{\mathcal{E}_0}$ of the remaining variables times the constant
$t_X(\mathcal{S})e^{\delta_X}$. By the definition of $W_{\mathfrak{o}}$ in
\eqref{8.6:eq-island-representation-a} and of $W^{\mathcal{S},\mathfrak{c}}_{\mathfrak{o}}$ in
\eqref{8.6:eq-healed-removal-a}, we obtain
\begin{align*}
W^{\mathcal{S},\mathfrak{c}}_{\mathfrak{o}}(\mathfrak{X}'\sqcup\{X\})
&=\Bigl[\prod_{b\in\mathfrak{X}'}\mathbf{T}'_k(X_b)\Bigr]\mathbf{T}'_{k,\mathcal{S}}(X)
\bigl(e^{\mathcal{E}_1}\bigr)\\
&\le t_X(\mathcal{S})\,e^{\delta_X}\Bigl[\prod_{b\in\mathfrak{X}'}\mathbf{T}'_k(X_b)\Bigr]
\bigl(e^{\mathcal{E}_0}\bigr)\\
&=t_X(\mathcal{S})\,e^{\delta_X}\,W_{\mathfrak{o}}(\mathfrak{X}').
\end{align*}
This is \eqref{8.6:eq-healed-removal-c}.

\emph{Step (iv): the fibres.} We claim that the map
$\mathfrak{X}\mapsto(\operatorname{red}_{\mathfrak{c}}(\mathfrak{X}),X_{a(\mathfrak{c})})$ is a
bijection from $E_{\mathfrak{c}}\cap B'$ onto the set of pairs $(\mathfrak{X}',X)$ with
$\mathfrak{X}'\in B'\setminus E_{\mathfrak{c}}$, $X\ni\mathfrak{c}$, and
$\mathfrak{X}'\sqcup\{X\}$ admissible.

First, the map lands in this set. Let $\mathfrak{X}\in E_{\mathfrak{c}}\cap B'$. The family
$\operatorname{red}_{\mathfrak{c}}(\mathfrak{X})$ is a sub-family of an admissible family, hence
admissible: (c1) is a condition on each island separately, and (c2), (c3) pass to sub-families.
It does not lie in $E_{\mathfrak{c}}$, since $X_{a(\mathfrak{c})}$ was the only island containing
$\mathfrak{c}$. It lies in $B'$ because $B'$ is $\mathfrak{c}$-blind. Finally,
$\operatorname{red}_{\mathfrak{c}}(\mathfrak{X})\sqcup\{X_{a(\mathfrak{c})}\}=\mathfrak{X}$ is
admissible.

Second, the map is injective, since $\mathfrak{X}$ is recovered as this union.

Third, the map is onto. Given a pair $(\mathfrak{X}',X)$, the family
$\mathfrak{X}:=\mathfrak{X}'\sqcup\{X\}$ is admissible. Its island containing $\mathfrak{c}$ is
$X$, so $\mathfrak{X}\in E_{\mathfrak{c}}$ and $\operatorname{red}_{\mathfrak{c}}(\mathfrak{X})
=\mathfrak{X}'$. By $\mathfrak{c}$-blindness,
$1_{B'}(\mathfrak{X})=1_{B'}(\mathfrak{X}')=1$.

Next we count the islands $X$ that can appear. By the box condition, such an island is a box of
level-$k$ cells with side lengths $\ell_\mu\in\{1,\dots,100\}$. Since it contains $\mathfrak{c}$,
it is determined by these side lengths and by its offsets relative to $\mathfrak{c}$, the offset
in direction $\mu$ being less than $\ell_\mu$ and so taking at most $\ell_\mu$ values. There are
therefore at most
\begin{equation*}
\prod_{\mu=1}^d\Bigl(\sum_{\ell=1}^{100}\ell\Bigr)=5050^d\le10^{4d}=K_{\mathrm{rm}}
\end{equation*}
candidates $X$.

For each pair, Lemma~\ref{8.6:lem-prefactor-monotone} (whose proof is incomplete at one step)
gives $\mathbf{O}_{\mathfrak{o},\mathfrak{X}'\sqcup\{X\}}\le\mathbf{O}_{\mathfrak{o},\mathfrak{X}'}$
as positive operators.

We now reindex the left member of \eqref{8.6:eq-healed-removal-d} by the bijection. We apply the
last operator inequality to the non-negative function
$W^{\mathcal{S},\mathfrak{c}}_{\mathfrak{o}}(\mathfrak{X}'\sqcup\{X\})$. We then apply the positive
linear operator $\mathbf{O}_{\mathfrak{o},\mathfrak{X}'}$ to \eqref{8.6:eq-healed-removal-c} and sum
over $X$. Finally we bound each of the at most $K_{\mathrm{rm}}$ numbers
$t_X(\mathcal{S})e^{\delta_X}$ by their maximum. The result is
\begin{equation}\label{8.6:eq-healed-removal-4}
\begin{aligned}
\sum_{\mathfrak{X}\in E_{\mathfrak{c}}\cap B'}\mathbf{O}_{\mathfrak{o},\mathfrak{X}}
\bigl[W^{\mathcal{S},\mathfrak{c}}_{\mathfrak{o}}(\mathfrak{X})\bigr]
&\le\sum_{\mathfrak{X}'\in B'\setminus E_{\mathfrak{c}}}\mathbf{O}_{\mathfrak{o},\mathfrak{X}'}
\Bigl[\sum_{X\ni\mathfrak{c}}t_X(\mathcal{S})\,e^{\delta_X}\,W_{\mathfrak{o}}(\mathfrak{X}')\Bigr]\\
&\le\Bigl[K_{\mathrm{rm}}\max_X t_X(\mathcal{S})\,e^{\delta_X}\Bigr]
\sum_{\mathfrak{X}'\in B'\setminus E_{\mathfrak{c}}}\mathbf{O}_{\mathfrak{o},\mathfrak{X}'}
\bigl[W_{\mathfrak{o}}(\mathfrak{X}')\bigr].
\end{aligned}
\end{equation}
The inner sum in the middle member runs over the islands $X\ni\mathfrak{c}$ with
$\mathfrak{X}'\sqcup\{X\}$ admissible.

It remains to bound $t_X(\mathcal{S})e^{\delta_X}$ for each candidate $X$.

Suppose $\theta=0$. The positivity hypothesis \eqref{8.6:eq-island-operation-bound-a}, applied
with $\Phi\equiv1$ to a singleton set $\{\mathfrak{h}\}$, gives
$0\le\mathbf{T}'_{k,\{\mathfrak{h}\}}(X)1\le\mathbf{T}'_{k,\{\mathfrak{h}\}}(X,(U,0))1$. So the
summands of $\mathbf{T}'_k(X,(U,0))1$ over internal histories are non-negative. Restricting the
sum to $\mathcal{S}$ therefore gives $t_X(\mathcal{S})\le\mathbf{T}'_k(X,(U,0))1$, and
\eqref{8.6:eq-island-operation-bound-b} bounds the right side by
$\exp(-2(1+\beta_0)^{-1}p_0(g_k))$. Recall that $\varpi_0\equiv1$.

Suppose $\theta$ lies in the range of Lemma~\ref{8.6:lem-reserve-island-bound}. Then
\eqref{8.6:eq-reserve-island-bound-a} applies directly.

In both cases, combining with \eqref{8.6:eq-healed-removal-e} and then with
\eqref{8.6:eq-parameter-conditions-b}, we get
\begin{align*}
t_X(\mathcal{S})\,e^{\delta_X}
&\le\varpi_\theta(\mathcal{S})\exp\bigl(-(2-\theta)(1+\beta_0)^{-1}p_0(g_k)
+C_\delta(d)\,\alpha\,(101\check{R}_k)^d\bigr)\\
&\le\varpi_\theta(\mathcal{S})\,e^{-\frac32p_0(g_k)},
\end{align*}
with $\theta=0$ included. Hence the bracket in the last member of \eqref{8.6:eq-healed-removal-4}
is at most $\epsilon^{\mathrm{rm}}_k(\mathcal{S})$, and the first inequality of
\eqref{8.6:eq-healed-removal-d} follows.

The second inequality of \eqref{8.6:eq-healed-removal-d} follows by adding the terms with
$\mathfrak{X}\in B'\cap E_{\mathfrak{c}}$. Each of them,
$\mathbf{O}_{\mathfrak{o},\mathfrak{X}}[W_{\mathfrak{o}}(\mathfrak{X})]$, is non-negative as the
image of a non-negative function under a positive operator.

The set of all admissible families is trivially $\mathfrak{c}$-blind. For this choice of $B'$ the
right member of \eqref{8.6:eq-healed-removal-d} is
$\epsilon^{\mathrm{rm}}_k(\mathcal{S})\rho_{k,\mathfrak{o}}$ by
\eqref{8.6:eq-island-representation-a}. This is the bound on the pinned part.

If $\mathfrak{c}$ shares a point with $|\mathfrak{o}|$, then by (c3) no island of an admissible
family contains $\mathfrak{c}$. Hence $E_{\mathfrak{c}}=\emptyset$ and the left members are empty
sums.
\end{proof}

\begin{remark}
\leavevmode
\begin{enumerate}
\item[(a)] In the proof of Lemma~\ref{8.6:lem-healed-removal} no lower bound on any activity, no
cluster expansion and no analyticity enter. The
comparison is made term by term between two families that differ in exactly one healed island.
Of Ba{\l}aban's estimates only the bound \eqref{8.6:eq-island-operation-bound-b}, the bound
\eqref{8.6:eq-island-operation-bound-c}, the locality \eqref{8.6:eq-island-operation-bound-d} and
positivity enter; the remaining hypotheses are combinatorial or are conditions on the parameters.
For a healed island, a comparison of Peierls type against a local lower bound
is already contained in Ba{\l}aban's normalised operation $\mathbf{T}'_k$ and in the completed
small-field effective action. This is the purpose of the preparation preceding the large-field
operation $\mathbf{R}$. There the denominators are taken equal to integrals of parts of the
densities, localised in neighbourhoods of the large-field domains and determined by small-field
effective actions only \cite{B-CMP122a}.
\item[(b)] The exponent in $\epsilon^{\mathrm{rm}}_k(\mathcal{S})$ is the full
$-\frac32p_0(g_k)$.
No part of it is spent on a sum over the positions of the other islands; the entropy of the pinned
island is the finite factor $K_{\mathrm{rm}}$.
\item[(c)] The $\mathfrak{c}$-blind event $B'$ may fix the entire configuration of the other
islands of the step. In that case \eqref{8.6:eq-healed-removal-d} is a conditional bound of
Dobrushin--Lanford--Ruelle type.
\end{enumerate}
\end{remark}

\begin{lemma}[One-point, conditional and joint bounds in a positive hard-core gas]
\label{8.6:lem-hard-core-gas}
Let $\mathcal{P}$ be a finite set and let $\iota$ be a reflexive and symmetric relation on
$\mathcal{P}$; we say that $\eta,\eta'\in\mathcal{P}$ are \emph{incompatible} if
$\eta\,\iota\,\eta'$. For $\eta\in\mathcal{P}$ put
$\iota(\eta)=\{\eta'\in\mathcal{P}:\eta'\,\iota\,\eta\}$, and for
$\mathcal{B}\subset\mathcal{P}$ put $\iota(\mathcal{B})=\bigcup_{\eta\in\mathcal{B}}\iota(\eta)$;
by reflexivity $\eta\in\iota(\eta)$. A subset $\mathfrak{F}\subset\mathcal{P}$ is
\emph{compatible} if no two distinct elements of $\mathfrak{F}$ are incompatible. For an
\emph{activity} $\mathsf{z}:\mathcal{P}\to\mathbb{C}$ and $\mathcal{Q}\subset\mathcal{P}$ put
\begin{equation}\label{8.6:eq-hard-core-gas-1}
\Xi(\mathcal{Q})=\sum_{\mathfrak{F}\subset\mathcal{Q}\ \text{compatible}}
\ \prod_{\eta\in\mathfrak{F}}\mathsf{z}(\eta),
\end{equation}
the empty family contributing $1$. For $\eta_1,\dots,\eta_n\in\mathcal{P}$ let
$\phi^{\mathrm{T}}(\eta_1,\dots,\eta_n)=\sum_G(-1)^{|G|}$, the sum over the connected graphs
$G$ on $\{1,\dots,n\}$ each edge $\{i,j\}$ of which satisfies $\eta_i\,\iota\,\eta_j$, with
$|G|$ the number of edges of $G$; the \emph{cluster series} of $\mathsf{z}$ is
\begin{equation}\label{8.6:eq-hard-core-gas-2}
\sum_{n\ge1}\frac1{n!}\sum_{(\eta_1,\dots,\eta_n)\in\mathcal{P}^n}
\phi^{\mathrm{T}}(\eta_1,\dots,\eta_n)\prod_{i=1}^n\mathsf{z}(\eta_i).
\end{equation}
For $\eta_0\in\mathcal{P}$ and $\mathsf{t}\in\mathbb{R}$ let
$\mathsf{z}_{\mathsf{t}}(\eta_0)=(1+\mathsf{t})\,\mathsf{z}(\eta_0)$ and
$\mathsf{z}_{\mathsf{t}}(\eta)=\mathsf{z}(\eta)$ for $\eta\ne\eta_0$, and let
$\Xi_{\mathsf{t}}$ be \eqref{8.6:eq-hard-core-gas-1} formed with $\mathsf{z}_{\mathsf{t}}$.

If $\mathsf{z}\ge0$, then $\Xi(\mathcal{P})\ge1$, and
$\mathbb{P}(\mathfrak{F})=\prod_{\eta\in\mathfrak{F}}\mathsf{z}(\eta)/\Xi(\mathcal{P})$ is a
probability on the compatible subsets of $\mathcal{P}$. For part (d), suppose that a function
$\mathsf{a}:\mathcal{P}\to[0,\infty[$ satisfies
\begin{equation}\label{8.6:eq-hard-core-gas-3}
\sum_{\eta\in\iota(\eta'')}|\mathsf{z}(\eta)|\,e^{\mathsf{a}(\eta)}\le\mathsf{a}(\eta'')
\qquad\text{for every }\eta''\in\mathcal{P},
\end{equation}
and assume the conclusions of the lemma on absolute convergence of the cluster series and the
pinned rooted-tree bound for every activity that satisfies \eqref{8.6:eq-hard-core-gas-3} with
this $\mathsf{a}$; for $\mathsf{z}$ they read: $\Xi(\mathcal{P})\ne0$, $\log\Xi(\mathcal{P})$ is
the absolutely convergent series \eqref{8.6:eq-hard-core-gas-2}, and for every
$\eta_0\in\mathcal{P}$
\begin{equation}\label{8.6:eq-hard-core-gas-4}
\sum_{n\ge1}\frac1{(n-1)!}\sum_{\eta_2,\dots,\eta_n\in\mathcal{P}}
\bigl|\phi^{\mathrm{T}}(\eta_0,\eta_2,\dots,\eta_n)\bigr|\prod_{i=2}^n|\mathsf{z}(\eta_i)|
\ \le\ e^{\mathsf{a}(\eta_0)}.
\end{equation}

Then, for every $\eta_0\in\mathcal{P}$, every compatible $\mathcal{B}\subset\mathcal{P}$ with
$\eta_0\notin\mathcal{B}$ and $\mathbb{P}(\mathfrak{F}\supseteq\mathcal{B})>0$, and all pairwise
distinct $\eta_1,\dots,\eta_n\in\mathcal{P}$: (a)--(c) hold when $\mathsf{z}\ge0$, and (d) holds
when \eqref{8.6:eq-hard-core-gas-3} holds:
\begin{equation}\label{8.6:eq-hard-core-gas-a}
\begin{aligned}
&\text{(a)}\quad \mathbb{P}(\eta_0\in\mathfrak{F})
=\frac{\mathsf{z}(\eta_0)\,\Xi(\mathcal{P}\setminus\iota(\eta_0))}{\Xi(\mathcal{P})}
\le\mathsf{z}(\eta_0),\\
&\text{(b)}\quad \mathbb{P}(\eta_0\in\mathfrak{F}\mid\mathfrak{F}\supseteq\mathcal{B})
\le\mathsf{z}(\eta_0),\\
&\text{(c)}\quad \mathbb{P}(\eta_1,\dots,\eta_n\in\mathfrak{F})\le\prod_{i=1}^n\mathsf{z}(\eta_i),\\
&\text{(d)}\quad \bigl|\partial_{\mathsf{t}}\log\Xi_{\mathsf{t}}(\mathcal{P})
\big|_{\mathsf{t}=0}\bigr|\le|\mathsf{z}(\eta_0)|\,e^{\mathsf{a}(\eta_0)}.
\end{aligned}
\end{equation}
The left member of (b) vanishes if $\eta_0\,\iota\,\eta$ for some $\eta\in\mathcal{B}$; and if
$\mathsf{z}\ge0$, the derivative in (d) equals $\mathbb{P}(\eta_0\in\mathfrak{F})$.
\end{lemma}

\begin{proof}
For a compatible $\mathcal{B}\subset\mathcal{P}$ the map $\mathfrak{F}'\mapsto\mathcal{B}\sqcup\mathfrak{F}'$
is a bijection from the compatible subsets of $\mathcal{P}\setminus\iota(\mathcal{B})$ onto the
compatible subsets of $\mathcal{P}$ containing $\mathcal{B}$. Indeed, if $\mathfrak{F}\supseteq\mathcal{B}$
is compatible, every element of $\mathfrak{F}\setminus\mathcal{B}$ is distinct from, and compatible
with, each element of $\mathcal{B}$, hence lies outside $\iota(\mathcal{B})$; conversely
$\mathcal{B}\cup\mathfrak{F}'$ is compatible for compatible
$\mathfrak{F}'\subset\mathcal{P}\setminus\iota(\mathcal{B})$, and the union is disjoint because
$\mathcal{B}\subset\iota(\mathcal{B})$. Hence
\begin{equation}\label{8.6:eq-hard-core-gas-5}
\sum_{\mathfrak{F}\supseteq\mathcal{B}\ \text{compatible}}\ \prod_{\eta\in\mathfrak{F}}\mathsf{z}(\eta)
=\Bigl(\prod_{\eta\in\mathcal{B}}\mathsf{z}(\eta)\Bigr)\,\Xi(\mathcal{P}\setminus\iota(\mathcal{B})).
\end{equation}

(a) Take $\mathcal{B}=\{\eta_0\}$ in \eqref{8.6:eq-hard-core-gas-5} and divide by
$\Xi(\mathcal{P})$: this is the identity in (a). The compatible subsets of
$\mathcal{P}\setminus\iota(\eta_0)$ are compatible subsets of $\mathcal{P}$, and all terms of
\eqref{8.6:eq-hard-core-gas-1} are non-negative, so
$\Xi(\mathcal{P}\setminus\iota(\eta_0))\le\Xi(\mathcal{P})$, which gives the inequality.

(b) If $\eta_0\,\iota\,\eta$ for some $\eta\in\mathcal{B}$, no compatible family contains
$\mathcal{B}\cup\{\eta_0\}$ and the conditional probability is $0$. Otherwise
$\mathcal{B}\cup\{\eta_0\}$ is compatible, and $\iota(\mathcal{B}\cup\{\eta_0\})
=\iota(\mathcal{B})\cup\iota(\eta_0)$. Apply \eqref{8.6:eq-hard-core-gas-5} to $\mathcal{B}$ and to
$\mathcal{B}\cup\{\eta_0\}$ and take the quotient; the factor
$\prod_{\eta\in\mathcal{B}}\mathsf{z}(\eta)$, which is positive because
$\mathbb{P}(\mathfrak{F}\supseteq\mathcal{B})>0$, cancels, and with
$\mathcal{Q}=\mathcal{P}\setminus\iota(\mathcal{B})$
\begin{equation*}
\mathbb{P}(\eta_0\in\mathfrak{F}\mid\mathfrak{F}\supseteq\mathcal{B})
=\frac{\mathsf{z}(\eta_0)\,\Xi(\mathcal{Q}\setminus\iota(\eta_0))}{\Xi(\mathcal{Q})}.
\end{equation*}
This is the right member of the identity in (a) with $\mathcal{P}$ replaced by $\mathcal{Q}$, and
the argument of (a) bounds it by $\mathsf{z}(\eta_0)$.

(c) Put $\mathcal{B}_0=\emptyset$ and $\mathcal{B}_j=\{\eta_1,\dots,\eta_j\}$. If $\mathcal{B}_n$ is
not compatible, or if $\mathbb{P}(\mathfrak{F}\supseteq\mathcal{B}_j)=0$ for some $j$, the left member
of (c) is $0$, since the event $\mathfrak{F}\supseteq\mathcal{B}_n$ is contained in
$\mathfrak{F}\supseteq\mathcal{B}_j$. Otherwise each $\mathcal{B}_{j-1}$ is compatible, does not
contain $\eta_j$ (the $\eta_i$ are distinct) and has positive probability, and
\begin{equation*}
\mathbb{P}(\mathfrak{F}\supseteq\mathcal{B}_n)
=\prod_{j=1}^n\frac{\mathbb{P}(\mathfrak{F}\supseteq\mathcal{B}_j)}
{\mathbb{P}(\mathfrak{F}\supseteq\mathcal{B}_{j-1})}
=\prod_{j=1}^n\mathbb{P}(\eta_j\in\mathfrak{F}\mid\mathfrak{F}\supseteq\mathcal{B}_{j-1})
\le\prod_{j=1}^n\mathsf{z}(\eta_j),
\end{equation*}
by (b) applied $n$ times.

(d) By the lemma on absolute convergence of the cluster series and the pinned rooted-tree
bound, $\Xi(\mathcal{P})\ne0$ and $\log\Xi(\mathcal{P})$ is the absolutely convergent series
\eqref{8.6:eq-hard-core-gas-2}. In that series for $\mathsf{z}_{\mathsf{t}}$ the product
$\prod_i\mathsf{z}_{\mathsf{t}}(\eta_i)$ equals $(1+\mathsf{t})^{\mathsf{m}}\prod_i\mathsf{z}(\eta_i)$,
where $\mathsf{m}$ is the number of indices $i$ with $\eta_i=\eta_0$; its $\mathsf{t}$-derivative
at $0$ is $\mathsf{m}\prod_i\mathsf{z}(\eta_i)$. For $-1<\mathsf{t}\le0$ we have
$|\mathsf{z}_{\mathsf{t}}|\le|\mathsf{z}|$, so $\mathsf{z}_{\mathsf{t}}$ satisfies
\eqref{8.6:eq-hard-core-gas-3} with the same $\mathsf{a}$, and by the same lemma
$\log\Xi_{\mathsf{t}}(\mathcal{P})$ is the absolutely convergent series
\eqref{8.6:eq-hard-core-gas-2} formed with $\mathsf{z}_{\mathsf{t}}$. Differentiating that series
term by term, for $-1<\mathsf{t}\le0$ the term indexed by $(\eta_1,\dots,\eta_n)$ has absolute value
at most
\begin{equation*}
\frac{\mathsf{m}}{n!}\,\bigl|\phi^{\mathrm{T}}(\eta_1,\dots,\eta_n)\bigr|\prod_{i=1}^n|\mathsf{z}(\eta_i)|,
\end{equation*}
since $0\le\mathsf{m}(1+\mathsf{t})^{\mathsf{m}-1}\le\mathsf{m}$ there. By the symmetrisation carried
out below, the sum of these majorants is $|\mathsf{z}(\eta_0)|$ times the left member of
\eqref{8.6:eq-hard-core-gas-4}, hence finite. The differentiated series therefore converges
uniformly on $]-1,0]$, and the derivative of $\log\Xi_{\mathsf{t}}(\mathcal{P})$ from the left at
$\mathsf{t}=0$ is the termwise derivative. Since $\Xi_{\mathsf{t}}(\mathcal{P})$ is a polynomial in
$\mathsf{t}$ which does not vanish at $\mathsf{t}=0$, the function
$\mathsf{t}\mapsto\log\Xi_{\mathsf{t}}(\mathcal{P})$ is differentiable at $\mathsf{t}=0$, with
derivative $\partial_{\mathsf{t}}\Xi_{\mathsf{t}}(\mathcal{P})/\Xi_{\mathsf{t}}(\mathcal{P})$ at
$\mathsf{t}=0$, so this derivative equals the termwise one. Thus the derivative at $\mathsf{t}=0$
selects the sequences passing through $\eta_0$, each counted with multiplicity $\mathsf{m}$. Since
$\phi^{\mathrm{T}}$ and the product are symmetric in $(\eta_1,\dots,\eta_n)$, the sequences with
$\eta_i=\eta_0$ contribute equally for each $i$, so the $n$-th term becomes
$n/n!=1/(n-1)!$ times the sum over sequences with $\eta_1=\eta_0$:
\begin{equation*}
\partial_{\mathsf{t}}\log\Xi_{\mathsf{t}}(\mathcal{P})\big|_{\mathsf{t}=0}
=\mathsf{z}(\eta_0)\sum_{n\ge1}\frac1{(n-1)!}\sum_{\eta_2,\dots,\eta_n\in\mathcal{P}}
\phi^{\mathrm{T}}(\eta_0,\eta_2,\dots,\eta_n)\prod_{i=2}^n\mathsf{z}(\eta_i).
\end{equation*}
The same symmetrisation applied to the majorants gives the identity used above. Taking absolute
values and applying \eqref{8.6:eq-hard-core-gas-4} gives (d). Finally, splitting
\eqref{8.6:eq-hard-core-gas-1} for $\mathsf{z}_{\mathsf{t}}$ according to whether $\eta_0\in\mathfrak{F}$
and using \eqref{8.6:eq-hard-core-gas-5} with $\mathcal{B}=\{\eta_0\}$,
\begin{equation*}
\Xi_{\mathsf{t}}(\mathcal{P})=\Xi(\mathcal{P}\setminus\{\eta_0\})
+(1+\mathsf{t})\,\mathsf{z}(\eta_0)\,\Xi(\mathcal{P}\setminus\iota(\eta_0)),
\end{equation*}
so the derivative equals $\mathsf{z}(\eta_0)\Xi(\mathcal{P}\setminus\iota(\eta_0))/\Xi(\mathcal{P})$,
which for $\mathsf{z}\ge0$ is $\mathbb{P}(\eta_0\in\mathfrak{F})$ by (a).
\end{proof}

\begin{remark}
For non-negative activities, parts (a)--(c) give upper bounds, joint and conditional, without
any expansion. The removal domination for a healed island, \eqref{8.6:eq-healed-removal-d}, is
parts (a) and (b) carried out on the sum that appears there in Ba{\l}aban's representation of the
effective density $\rho_k$, with three inputs: the terms of that sum are non-negative; the
configuration from which an island has been removed is again a term of the same sum, with the same
factor outside the island, by \eqref{8.6:eq-prefactor-monotone-a}; and the weight of an island
relative to its own removal is at most its activity, which is supplied by the quotient
$\mathbf{T}'$ of Ba{\l}aban's operation $\mathbf{R}$ through the positivity and bound of the
normalised island operation, \eqref{8.6:eq-island-operation-bound-b}.
\end{remark}

\begin{remark}[Bounds on the majorant gas do not bound the history measure]\label{8.6:rem-majorant-gas}
Let $\tilde{\nu}$ be the annealed measure on the histories of the run, with the masses of
Definition~\ref{8.6:def-history-measure}. The bounds of the correlation theorem on the
live-component gas, in which the partition function $\Xi^{\mathrm{maj}}$ of the gas of level $k$,
displayed in~\eqref{8.6:eq-majorant-gas-1} below, satisfies
$\Xi^{\mathrm{maj}}\le\exp\bigl(e^{-p_0(g_k)}|\mathbb{L}_k|\bigr)$ and the sourced partition function
satisfies $|\tilde{Z}|\le e^{E_+^{\sharp}\mathcal{V}}$, together with
\cite[(1.79)]{B-CMP122b} and \cite[(1.89)]{B-CMP122b}, bound the sum of the history terms
$\tau_h$ in the form
\begin{equation}\label{8.6:eq-majorant-gas-1}
\Bigl|\sum_h\tau_h\Bigr|\le\Phi_k\,\Xi^{\mathrm{maj}},\qquad
\Xi^{\mathrm{maj}}=\sum_{\mathfrak{F}}\ \prod_{Z\in\mathfrak{F}}\mathfrak{a}(Z).
\end{equation}
Here $p_0(\cdot)$ is the degree function; $\mathbb{L}_k$ is the block lattice of level $k$ of
Definition~\ref{1.5:def-block-average} and $|\mathbb{L}_k|$ the number of its points;
$\tilde{Z}$ is the sourced partition function
of the correlation theorem and $\tilde{Z}(0)$ its value at zero source; $\mathcal{V}$ is the
volume factor in which the correlation theorem states its bounds, the fourth power of the
linear size of the torus $\mathbb{T}_m$ in the units of those bounds, which grows with $m$;
$E_+^{\sharp}$ and
$C_{\mathrm{low}}$ are the constants of the correlation theorem's upper bound on $|\tilde{Z}|$
stated above and of its lower bound on $\tilde{Z}(0)$ recalled in the proof below;
$\Phi_k$ is a majorant of the small-field factor;
the sum runs over the compatible families $\mathfrak{F}$ of live components; and, for a
component $Z$ created at step $j=j(Z)$, the activity $\mathfrak{a}(Z)\ge0$ is bounded by
$e^{-\kappa_j(Z)-2p_0(g_{j(Z)})}$ times a combinatorial entropy factor, where $\kappa_j(Z)$ is
the exponent attached to the region $Z$, created at step $j$, in the bound of
\cite[\S1]{B-CMP122b} on its large-field factor.
The following hold.
\begin{itemize}
\item[(a)] Part (a) of Lemma~\ref{8.6:lem-hard-core-gas}, applied to $\Xi^{\mathrm{maj}}$,
bounds the occupation weight of each component $Z$ in the majorant gas by $\mathfrak{a}(Z)$.
\item[(b)] This is a statement about the normalisation of the majorant
$\Xi^{\mathrm{maj}}$, not about $\tilde{\nu}$.
\item[(c)] An application of Lemma~\ref{8.6:lem-hard-core-gas} to $\tilde{\nu}$ itself
requires activities of $\tilde{\nu}$ of the form
\begin{equation}\label{8.6:eq-majorant-gas-2}
\frac{\text{mass with $Z$ present, the rest given}}
{\text{mass with the region of $Z$ in small-field form}},
\end{equation}
together with an upper bound on this quotient by one number, the same for every
component $Z$. For healed islands this is
removal domination for a healed island (Lemma~\ref{8.6:lem-healed-removal}). For live
components it is a live-tail statement, asked at each step of the run and for each
large-field region from which the live component descends; the annealed live-tail
statement is not proved and is not used in the proof of
Theorem~\ref{3.1:thm-main}. The proof of clause~(iv) of Theorem~\ref{3.1:thm-main} given in
this book uses no live-tail statement and does not pass through such
activities for live components: there the live part of the large-field bound is controlled
separately for young live regions by the corollary on the rarity of young large-field blocks,
for old healed regions by the healed old-region bound, and, for old regions never healed, by
the un-normalised age pre-payment at the last lattice of the run, prolonged on the same
lattice, under the per-torus depth floor of clause~(iv) of Theorem~\ref{3.1:thm-main}. A bound
on the quotients~\eqref{8.6:eq-majorant-gas-2} for live components would give an alternative
route to the same control; that alternative route is not used in the proof of
Theorem~\ref{3.1:thm-main}.
\end{itemize}
\end{remark}

\begin{proof}
(a) Take for $\mathcal{P}$ the finite set of live components and for the incompatibility
relation the negation of compatibility. Take $\mathfrak{a}$ as the activity; it is
non-negative. Then $\Xi^{\mathrm{maj}}=\Xi(\mathcal{P})$ is the partition function of a
positive hard-core gas in the sense of Lemma~\ref{8.6:lem-hard-core-gas}. Write
$P^{\mathrm{maj}}$ for the occupation measure of the majorant gas, which gives to a compatible
family $\mathfrak{F}$ the weight
\begin{equation*}
P^{\mathrm{maj}}(\mathfrak{F})=\frac{1}{\Xi^{\mathrm{maj}}}\prod_{Z\in\mathfrak{F}}\mathfrak{a}(Z);
\end{equation*}
write $\iota(Z)$ for the set of components incompatible with $Z$; and, for
$\mathcal{Q}\subset\mathcal{P}$, let $\Xi^{\mathrm{maj}}(\mathcal{Q})$
denote the sum~\eqref{8.6:eq-majorant-gas-1} restricted to families contained in
$\mathcal{Q}$, so that $\Xi^{\mathrm{maj}}=\Xi^{\mathrm{maj}}(\mathcal{P})$. Inequality (a)
of~\eqref{8.6:eq-hard-core-gas-a} gives,
for every $Z\in\mathcal{P}$,
\begin{equation*}
P^{\mathrm{maj}}(Z\in\mathfrak{F})
=\frac{\mathfrak{a}(Z)\,\Xi^{\mathrm{maj}}\bigl(\mathcal{P}\setminus\iota(Z)\bigr)}
{\Xi^{\mathrm{maj}}(\mathcal{P})}
\le\mathfrak{a}(Z).
\end{equation*}

(b) Both the numerator and the denominator of $P^{\mathrm{maj}}(Z\in\mathfrak{F})$ are built
from the activities $\mathfrak{a}$, that is, from the majorant in~\eqref{8.6:eq-majorant-gas-1}.
The probability under $\tilde{\nu}$ of the event that $Z$ is present is instead a quotient
whose numerator and denominator are sums of the masses of histories of
Definition~\ref{8.6:def-history-measure}. The bound~\eqref{8.6:eq-majorant-gas-1} controls
such a sum only from above. Transferring the bound of~(a) to $\tilde{\nu}$ would require an
identity expressing the total mass of the histories of
Definition~\ref{8.6:def-history-measure}, which is the denominator of $\tilde{\nu}$, as the
product of a local factor and the partition function of a hard-core gas of components. The
upper bounds quoted above provide no such identity, and the only lower bound available on the
denominator is the global floor
\begin{equation*}
\tilde{Z}(0)\ge e^{-C_{\mathrm{low}}\mathcal{V}}
\end{equation*}
of the correlation theorem, which is not local.
Hence~(a) bounds the majorant's own occupation weights and gives no bound on $\tilde{\nu}$.

(c) Lemma~\ref{8.6:lem-hard-core-gas}, applied with $\tilde{\nu}$ in place of the majorant
gas, uses an
activity for each component. That activity is the mass of a configuration in which the
component is present, relative to the mass of the configuration in which its region is in
small-field form, the rest being held fixed; these are the
quotients~\eqref{8.6:eq-majorant-gas-2}. When in addition the masses are non-negative and
removing the region of the component yields again a term of the same sum with the same
exterior factor, part (a) of the lemma gives the one-point
bound, and part~(b) the conditional bound given the other components, as in the proof of
Lemma~\ref{8.6:lem-healed-removal}; what is needed beyond this is an upper bound
on the quotients~\eqref{8.6:eq-majorant-gas-2} by one number, the same for every component.
The two cases, healed islands and live components, are as stated in~(c).
\end{proof}

\begin{definition}[Pinned island events and events blind to them]\label{8.6:def-blind-events}
Let $\mathcal{H}$ denote the set of histories of the run of
Definition~\ref{8.6:def-history-measure}. Fix a step $k$ and a level-$k$ cell $\mathfrak{c}$. For a
history $h\in\mathcal{H}$ write $\mathfrak{X}_k(h)$ for its family of islands of step $k$, that
is, the islands healed at step $k$, which are the healed labels of step $k$ in the sense of
Definition~\ref{8.6:def-labels-lineages}. The island of $\mathfrak{X}_k(h)$ containing
$\mathfrak{c}$, when there is one, is called the \emph{pinned island}.

\emph{Deletion of the pinned island with its lineage.} The map
$\mathrm{red}^{(k)}_{\mathfrak{c}}\colon\mathcal{H}\to\mathcal{H}$ is defined as follows. Suppose
$\mathfrak{X}_k(h)$ has an island $X$ containing $\mathfrak{c}$, and let $j(X)$ be its birth
level. Then $\mathrm{red}^{(k)}_{\mathfrak{c}}(h)$ is obtained from $h$ by deleting $X$ together
with its lineage in the sense of Definition~\ref{8.6:def-labels-lineages}, namely its internal
history: its live ancestor components at the levels $j(X),\dots,k-1$ and every label merged into
them; that is, at each level of that internal history the large-field structure of the lineage is
replaced by small field. If no island of
$\mathfrak{X}_k(h)$ contains $\mathfrak{c}$, then $\mathrm{red}^{(k)}_{\mathfrak{c}}(h)=h$.

The image $\mathrm{red}^{(k)}_{\mathfrak{c}}(h)$ is again a history. Indeed, the labels
of $h$ at later steps remain admissible after the deletion: in Ba{\l}aban's
construction~\cite{B-CMP122b} a healed island is not a large-field region at any later step, so
no admissibility condition of a later step reads it.

\emph{Pinned island events and blindness.} For a set $\mathcal{S}$ of internal histories, and
for an event $B\subset\mathcal{H}$, we put
\begin{equation}\label{8.6:eq-blind-events-a}
\begin{gathered}
\begin{aligned}
E^{(k)}_{\mathfrak{c},\mathcal{S}}
:=\bigl\{h\in\mathcal{H}:\ &\mathfrak{X}_k(h)\text{ has an island }X\ni\mathfrak{c}\\
&\text{whose internal history lies in }\mathcal{S}\bigr\},
\end{aligned}\\[2pt]
B\ \text{is $(k,\mathfrak{c})$-blind}
\quad:\Longleftrightarrow\quad
\mathbf{1}_B=\mathbf{1}_B\circ\mathrm{red}^{(k)}_{\mathfrak{c}} ,
\end{gathered}
\end{equation}
where $\mathbf{1}_B$ is the indicator function of $B$.
We also put $E^{(k)}_{\mathfrak{c}}:=E^{(k)}_{\mathfrak{c},\mathcal{S}}$, where $\mathcal{S}$ is
the set of all internal histories. Thus an event is $(k,\mathfrak{c})$-blind exactly when it is
invariant under the deletion of the pinned island with its lineage.

\emph{Examples.} Every event determined by the following data is $(k,\mathfrak{c})$-blind:
\begin{itemize}
\item the block field $V_k$, through the joint law $\tilde{\nu}_k$
(Definition~\ref{8.6:def-history-measure});
\item the outer data of the levels $\ge k$;
\item the islands of step $k$ other than the one containing $\mathfrak{c}$, with their
internal histories;
\item the islands of steps other than $k$, and the live labels of any level, provided they are
neither stages of the lineage of the pinned island (Definition~\ref{8.6:def-labels-lineages})
nor merged into that lineage;
\item the preparatory outcomes outside the lineage of the pinned island
(Definition~\ref{8.6:def-labels-lineages}).
\end{itemize}

\emph{Two classes of events that are not blind in general.} The following two classes of events
are not $(k,\mathfrak{c})$-blind in general:
\begin{itemize}
\item events determined by the outer data of all levels;
\item events generated by labels whose regions do not contain $\mathfrak{c}$.
\end{itemize}
For the first class: the earlier live stages and the seeds of the pinned island $X$ are outer
data of the levels $j(X),\dots,k-1$, and $\mathrm{red}^{(k)}_{\mathfrak{c}}$ deletes them. For the
second class: these stages and seeds need not contain the image of $\mathfrak{c}$
(Definition~\ref{8.6:def-labels-lineages}), because $X$ grows by layers before it reaches
$\mathfrak{c}$, and yet they are deleted; likewise, a label whose region is far from
$\mathfrak{c}$ but which later merges into the component that is eventually healed as $X$
belongs to the lineage of $X$, and so it is deleted together with $X$ by
$\mathrm{red}^{(k)}_{\mathfrak{c}}$.
\end{definition}

\begin{lemma}[Conservation of mass along the history tree]\label{8.6:lem-mass-conservation}
Let $\mathfrak{H}$, $\tau_\omega$, $Z$, $\tilde{\nu}$ and $\tilde{\nu}_k$ be as in
Definition~\ref{8.6:def-history-measure}, and let $k''$ denote the last level of the run, so that the
elements of $\mathfrak{H}$ are the level-$k''$ un-resummed histories. For a level $k$ of the run,
write $h_k$ for a level-$k$ un-resummed history, $F^{(h_k)}$ for its level-$k$ term, and
$\mathfrak{m}(h_k):=\int F^{(h_k)}\,dV$ for the total integral of that term.
\begin{itemize}
\item[(a)] For $k<k''$ and every level-$k$ un-resummed history $h_k$,
\begin{equation}\label{8.6:eq-mass-conservation-a}
\sum_{\omega\in\mathfrak{H}:\ \omega\ \text{extends}\ h_k}\tau_\omega=\mathfrak{m}(h_k).
\end{equation}
\item[(b)] Consequently, for an event $E\subset\mathfrak{H}$ determined by labels of levels $\le k$,
$\tilde{\nu}(E)$ is $Z^{-1}$ times the sum of the integrals of the level-$k$ terms over $E$, whatever
steps are performed after level $k$. Moreover $Z=\int\rho_k$ for every $k$.
\item[(c)] In tree form: the histories form a rooted tree on the un-resummed index space of
Definition~\ref{8.6:def-history-measure}, the children of $h_k$ being the histories
$h_{k+1}=(h_k,D_{k+1})$, where $D_{k+1}$ denotes the data produced by step $k+1$. Then
\begin{equation}\label{8.6:eq-mass-conservation-1}
\sum_{h_{k+1}\ \text{child of}\ h_k}\mathfrak{m}(h_{k+1})=\mathfrak{m}(h_k).
\end{equation}
For $j\le k$ and every set $E$ of level-$j$ histories,
\begin{equation}\label{8.6:eq-mass-conservation-2}
\tilde{\nu}_k\bigl(\{h_k:\ h_k\ \text{extends an element of}\ E\}\bigr)=\tilde{\nu}_j(E)=:\tilde{\nu}(E).
\end{equation}
\end{itemize}
\end{lemma}

\begin{proof}
\emph{One step conserves mass.} We prove \eqref{8.6:eq-mass-conservation-1} first. Fix a level-$k$
history $h_k$. Step $k+1$ acts on the term $F^{(h_k)}$ by a sequence of operations. Each operation
either produces new labels or produces none. The children $(h_k,D_{k+1})$ of $h_k$ are indexed by the
joint outcome $D_{k+1}$ of all new labels, and $F^{(h_k,D_{k+1})}$ is the term obtained by keeping
that outcome. It therefore suffices to show that each operation preserves the total integral of the
term it acts on, once the result is summed over the new labels it produces. We take the operations in
the order of the step.

\begin{itemize}
\item[(i)] \emph{The renormalisation transformation.} Let $T$ denote Ba{\l}aban's renormalisation
transformation \cite[(0.1), (3.1)]{B-CMP119}. It preserves total integrals:
\[
\int dV_{k+1}\,(T\rho)(V_{k+1})=\int dV_k\,\rho(V_k).
\]
We apply this to $\rho=F^{(h_k)}$. This is the same integral preservation that is used in the
correlation theorem.

\item[(ii)] \emph{Decompositions of unity.} Each decomposition of unity inserted in the step is a
finite sum of functions adding up to $1$ pointwise, the basic instance being the product, over
cubes $\square$, of the identities $1=\chi^{\kappa}_{\square}+(1-\chi^{\kappa}_{\square})$ of
\cite[(2.17)]{B-CMP119}, \cite[\S2]{B-CMP109} and \cite[(1.22)--(1.23)]{B-CMP122a}. Expanding
produces one term for each outcome. The outcomes are finitely many, and the terms sum to the function
they multiply. Hence the integral of that function equals the sum, over outcomes, of the integrals of
the terms.

\item[(iii)] \emph{Integrations.} Integrating out variables of a term, in whatever order the step
prescribes, leaves its total integral unchanged by Fubini's theorem.

\item[(iv)] \emph{The small-field reorganisation.} This rewrites a function as the same function. It
is an identity, so it changes no integral.

\item[(v)] \emph{The large-field operation.} $\mathbf{R}$ preserves the integral of the partial sum
belonging to each outer datum \cite[(0.3)--(0.5)]{B-CMP122a}, and it does so component by component
\cite[(1.100)--(1.102)]{B-CMP122a}. On the un-resummed index space the pair of large-field regions
$(Z',Z'')$ of \cite[(0.3)--(0.5)]{B-CMP122a} entering $\mathbf{R}$ is kept as a label. There the
same preservation holds for each fixed pair $(Z',Z'')$, by linearity of the operations in the
density they act on.

\item[(vi)] \emph{The Mayer expansion and the exponentiation identity} \cite[(1.98)]{B-CMP122b}.
Both are identities of functions. Hence they change no total integral.
\end{itemize}

Each step of the run maps the level-$k$ term to a finite sum, over the new labels, of positive
linear operators applied to it, followed by identities of functions; this is the structural identity
for one step of the run on the un-resummed index space of Definition~\ref{8.6:def-history-measure}.
Composing (i)--(vi) along this structure and summing over all outcomes $D_{k+1}$ gives
\eqref{8.6:eq-mass-conservation-1}.

\emph{Proof of \textup{(a)}.} For $k\le l\le k''$, let $\mathcal{D}_l(h_k)$ be the set of level-$l$
histories that extend $h_k$. We show
\begin{equation}\label{8.6:eq-mass-conservation-3}
\sum_{h_l\in\mathcal{D}_l(h_k)}\mathfrak{m}(h_l)=\mathfrak{m}(h_k)
\end{equation}
by induction on $l-k$. For $l=k$ the sum has the single term $\mathfrak{m}(h_k)$. Now suppose
\eqref{8.6:eq-mass-conservation-3} holds for some $l<k''$. Every element of $\mathcal{D}_{l+1}(h_k)$
is a child of exactly one element of $\mathcal{D}_l(h_k)$. Grouping the sum over
$\mathcal{D}_{l+1}(h_k)$ by parents and applying \eqref{8.6:eq-mass-conservation-1} to each parent
reduces it to the sum over $\mathcal{D}_l(h_k)$, which equals $\mathfrak{m}(h_k)$ by the induction
hypothesis.

At $l=k''$ the elements of $\mathcal{D}_{k''}(h_k)$ are the histories $\omega\in\mathfrak{H}$ extending
$h_k$. For each such $\omega$ we have $\mathfrak{m}(\omega)=\tau_\omega$, since $\tau_\omega$ is by
definition the total integral of the term $F^{(\omega)}$ of $\omega$. This proves
\eqref{8.6:eq-mass-conservation-a}.

\emph{Proof of \textup{(b)}.} Let $E\subset\mathfrak{H}$ be determined by labels of levels $\le k$.
Let $E_k$ be the set of level-$k$ histories all of whose extensions lie in $E$. Then $E$ is the
disjoint union, over $h_k\in E_k$, of the sets of histories extending $h_k$. By
\eqref{8.6:eq-history-measure-a} and \eqref{8.6:eq-mass-conservation-a},
\[
\tilde{\nu}(E)=Z^{-1}\sum_{h_k\in E_k}\ \sum_{\omega\ \text{extends}\ h_k}\tau_\omega
=Z^{-1}\sum_{h_k\in E_k}\mathfrak{m}(h_k).
\]
The right-hand side involves only level-$k$ terms and the constant $Z$, which equals $\int\rho_0$ by
Definition~\ref{8.6:def-history-measure}; so it does not depend on the steps performed after
level $k$.

Every $\omega\in\mathfrak{H}$ extends exactly one level-$k$ history. Summing
\eqref{8.6:eq-mass-conservation-a} over all $h_k$ therefore gives
\[
Z=\sum_{\omega\in\mathfrak{H}}\tau_\omega=\sum_{h_k}\mathfrak{m}(h_k).
\]
The last sum is the integral of the sum of the level-$k$ terms. These terms are the terms of the
expansion of $\rho_k$, so the sum equals $\int\rho_k$.

\emph{Proof of \textup{(c)}.} The identity \eqref{8.6:eq-mass-conservation-1} was proved above. Let
$j\le k$ and let $E$ be a set of level-$j$ histories. The quantity $\tilde{\nu}_k$ is taken on the set
of level-$k$ histories extending an element of $E$, with the top field integrated out. By
Definition~\ref{8.6:def-history-measure} it is therefore $Z^{-1}$ times the sum of $\mathfrak{m}(h_k)$
over these $h_k$. Each such $h_k$ extends exactly one $h_j\in E$, so
\[
\tilde{\nu}_k\bigl(\{h_k:\ h_k\ \text{extends an element of}\ E\}\bigr)
=Z^{-1}\sum_{h_j\in E}\ \sum_{h_k\in\mathcal{D}_k(h_j)}\mathfrak{m}(h_k).
\]
By \eqref{8.6:eq-mass-conservation-3} with $(j,k)$ in place of $(k,l)$, which was proved by induction
on the difference of the two levels, the inner sum is $\mathfrak{m}(h_j)$. Hence the right-hand side
equals $Z^{-1}\sum_{h_j\in E}\mathfrak{m}(h_j)=\tilde{\nu}_j(E)$. This is
\eqref{8.6:eq-mass-conservation-2}.

The common value does not depend on $k\ge j$. By (b) it agrees with the value of
\eqref{8.6:eq-history-measure-a} on the event of $\mathfrak{H}$ determined by $E$. This justifies the
notation $\tilde{\nu}(E)$. This is the reason why the annealed one-point weight of a label event
decided at level $j$ is a level-$j$ quantity.
\end{proof}

\begin{lemma}[Annealing and conditioning of the removal bound]\label{8.6:lem-annealed-removal}
Assume the following.
\begin{itemize}
\item[(a)] The hypotheses of the removal domination for a healed island,
Lemma~\ref{8.6:lem-healed-removal}, hold. Let $\varepsilon_k(\mathcal{S})$ be the number defined
in \eqref{8.6:eq-healed-removal-b}.
\item[(b)] The step structure of Lemma~\ref{8.6:lem-healed-islands-unread} (stated; proof
incomplete) holds: each step maps the terms of its level to terms of the next
level by a finite sum, over the new labels, of positive linear operators applied to the term,
followed by identities of functions.
\item[(c)] The statement \eqref{8.6:eq-healed-islands-unread-a} of
Lemma~\ref{8.6:lem-healed-islands-unread} (stated; proof incomplete) holds: the later
decompositions of unity, the quotient denominators and the completed actions are determined by
the outer datum and the universal small-field action on the current small-field region, not by
the healed islands.
\end{itemize}
Then the following holds for every $K$, every admissible run, every step $k\le K'$, every
level-$k$ cell $\mathfrak{c}$, every set $\mathcal{S}$ of internal histories and every
$(k,\mathfrak{c})$-blind event $B\subset\Omega$ in the sense of
Definition~\ref{8.6:def-blind-events}:
\begin{equation}\label{8.6:eq-annealed-removal-a}
\tilde\nu\bigl(E^{(k)}_{\mathfrak{c},\mathcal{S}}\cap B\bigr)
\le \varepsilon_k(\mathcal{S})\,\tilde\nu\bigl(B\setminus E^{(k)}_{\mathfrak{c}}\bigr)
\le \varepsilon_k(\mathcal{S})\,\tilde\nu(B).
\end{equation}
Here $\tilde\nu$ is the annealed measure of \eqref{8.6:eq-history-measure-a}, and
$E^{(k)}_{\mathfrak{c},\mathcal{S}}$, $E^{(k)}_{\mathfrak{c}}$ are the events of
\eqref{8.6:eq-blind-events-a}.
\end{lemma}

\begin{proof}
Fix $K$, an admissible run, a step $k\le K'$, a level-$k$ cell $\mathfrak{c}$, a set
$\mathcal{S}$ of internal histories and a $(k,\mathfrak{c})$-blind event $B$.

\emph{Coordinates on the histories.} We write every history $\omega\in\Omega$ as
$\omega=(\pi,\mathcal{O},\mathfrak{X},c)$. The four coordinates are as follows.
\begin{itemize}
\item $\pi$ is the collection of labels of $\omega$ at steps $<k$ that are not stages of the
lineages of the islands of step $k$.
\item $\mathcal{O}=\mathcal{O}_k$ is the outer datum of level $k$.
\item $\mathfrak{X}=\mathfrak{X}_k$ is the step-$k$ island family with its internal histories.
\item $c$ is the collection of labels of $\omega$ at steps $>k$.
\end{itemize}
The healed islands of earlier steps influence the integrand of step $k$ through the terms that
their large-field operations have produced. They do not influence the label structure of
step $k$, and are therefore recorded in $\pi$.

The whole lineage of each island of step $k$ is placed inside the coordinate $\mathfrak{X}$.
This agrees with \cite[(1.1)]{B-CMP122a}, where the whole accumulated operation of a
component is packaged into $\mathbf{T}_k(X)$. The deletion $\mathrm{red}^{(k)}_{\mathfrak{c}}$ of
Definition~\ref{8.6:def-blind-events} removes the pinned island together with its lineage, and
it keeps the later labels of $\omega$. It therefore acts on the coordinate $\mathfrak{X}$ alone,
leaves $\pi$ unchanged, and gives
\begin{equation*}
\mathrm{red}^{(k)}_{\mathfrak{c}}(\pi,\mathcal{O},\mathfrak{X},c)
=(\pi,\mathcal{O},\mathrm{red}_{\mathfrak{c}}(\mathfrak{X}),c),
\end{equation*}
with $\mathrm{red}_{\mathfrak{c}}$ as in \eqref{8.6:eq-healed-removal-a}.

\emph{The masses as positive images of the island representation.} Fix $(\pi,\mathcal{O},c)$.
By hypotheses (b) and (c), there is a positive linear operator
$\mathsf{O}_{c\mid\pi,\mathcal{O}}$. It is the composition, over the steps $k+1,k+2,\dots$ of the
run, of the following operations:
\begin{itemize}
\item the renormalization transformation $T$ of one step;
\item the characteristic functions selected by $c$, placed by the current outer data, which
descend from $\mathcal{O}$ and $c$ and not from $\mathfrak{X}$;
\item the integrations;
\item the quotient denominators and completed actions of the later $\mathbf{R}$-steps, which by
hypothesis (c) are functionals of the universal small-field action on the current small-field
region;
\item the selections of the later island events recorded in $c$, which are selections of
non-negative summands by label.
\end{itemize}
This operator satisfies
\begin{equation}\label{8.6:eq-annealed-removal-1}
\tau_{(\pi,\mathcal{O},\mathfrak{X},c)}
=\int \mathsf{O}_{c\mid\pi,\mathcal{O}}
\bigl(\mathbf{O}_{\mathcal{O},\mathfrak{X}}[W_{\mathcal{O}}(\mathfrak{X})]\bigr),
\end{equation}
where $W_{\mathcal{O}}$ is computed on the integrand that descends from $\pi$, as in the island
representation \eqref{8.6:eq-island-representation-a}.

Three facts enter \eqref{8.6:eq-annealed-removal-1} and its use below.
\begin{itemize}
\item The equality of total integrals is supplied by the conservation of mass along the history
tree, Lemma~\ref{8.6:lem-mass-conservation}, in the form \eqref{8.6:eq-mass-conservation-a}.
\item Lemma~\ref{8.6:lem-healed-removal} holds for every integrand of the inductive form, because
its inputs are uniform in the histories. In particular it holds for the integrand descending
from $\pi$.
\item The terms that later large-field operations produce from an island $X$ are part of the
integrand, and the domination propagates through it.
\end{itemize}

\emph{Blindness in coordinates.} Since $B$ is $(k,\mathfrak{c})$-blind, we have
$1_B=1_B\circ\mathrm{red}^{(k)}_{\mathfrak{c}}$. By the formula for
$\mathrm{red}^{(k)}_{\mathfrak{c}}$ above, there is a function $b$ with
\begin{equation*}
1_B(\pi,\mathcal{O},\mathfrak{X},c)=b\bigl(\pi,\mathcal{O},\mathrm{red}_{\mathfrak{c}}(\mathfrak{X}),c\bigr).
\end{equation*}
For fixed $(\pi,\mathcal{O},c)$ we put
\begin{equation*}
B'=B'_{\pi,\mathcal{O},c}
:=\bigl\{\mathfrak{X}:\ b(\pi,\mathcal{O},\mathrm{red}_{\mathfrak{c}}(\mathfrak{X}),c)=1\bigr\}.
\end{equation*}
The event $B'$ is $\mathfrak{c}$-blind in the sense of \eqref{8.6:eq-healed-removal-d}, that is,
$1_{B'}(\mathfrak{X})=1_{B'}(\mathrm{red}_{\mathfrak{c}}(\mathfrak{X}))$. To see this, take
$\mathfrak{X}\in E_{\mathfrak{c}}$. The island containing $\mathfrak{c}$ is unique, so no island
of $\mathrm{red}_{\mathfrak{c}}(\mathfrak{X})$ contains $\mathfrak{c}$. Hence
$\mathrm{red}_{\mathfrak{c}}(\mathfrak{X})\notin E_{\mathfrak{c}}$, and so
$\mathrm{red}_{\mathfrak{c}}(\mathrm{red}_{\mathfrak{c}}(\mathfrak{X}))=\mathrm{red}_{\mathfrak{c}}(\mathfrak{X})$
by the convention $\mathrm{red}_{\mathfrak{c}}(\mathfrak{X}):=\mathfrak{X}$ off $E_{\mathfrak{c}}$
of \eqref{8.6:eq-healed-removal-a}. For $\mathfrak{X}\notin E_{\mathfrak{c}}$ the same convention
gives the identity directly. Thus $\mathrm{red}_{\mathfrak{c}}\circ\mathrm{red}_{\mathfrak{c}}=\mathrm{red}_{\mathfrak{c}}$.

By the definitions of the events, a history $\omega=(\pi,\mathcal{O},\mathfrak{X},c)$ lies in
$E^{(k)}_{\mathfrak{c},\mathcal{S}}\cap B$ exactly when
$\mathfrak{X}\in E_{\mathfrak{c}}\cap B'_{\pi,\mathcal{O},c}$ and the internal history of the
island containing $\mathfrak{c}$ lies in $\mathcal{S}$. Likewise, $\omega$ lies in
$B\setminus E^{(k)}_{\mathfrak{c}}$ exactly when
$\mathfrak{X}\in B'_{\pi,\mathcal{O},c}\setminus E_{\mathfrak{c}}$.

\emph{The estimate.} We use three inputs.
\begin{itemize}
\item The representation \eqref{8.6:eq-annealed-removal-1}. Its restriction to the internal
histories of the pinned island lying in $\mathcal{S}$ is expressed through
$W^{\mathcal{S},\mathfrak{c}}_{\mathcal{O}}$ of \eqref{8.6:eq-healed-removal-a}.
\item The bound \eqref{8.6:eq-healed-removal-d} of Lemma~\ref{8.6:lem-healed-removal}, applied
to the $\mathfrak{c}$-blind event $B'_{\pi,\mathcal{O},c}$. This bound holds as an inequality
between non-negative functions of $V_k$ and of the outer variables.
\item The monotonicity of the positive operator
$\int\mathsf{O}_{c\mid\pi,\mathcal{O}}(\,\cdot\,)$ on non-negative functions, together with its
linearity. Linearity lets the number $\varepsilon_k(\mathcal{S})$, which does not depend on the
variables on which $\int\mathsf{O}_{c\mid\pi,\mathcal{O}}$ acts, be taken out of the operator.
\end{itemize}
Together these give
\begin{equation}\label{8.6:eq-annealed-removal-2}
\begin{aligned}
\sum_{\omega\in E^{(k)}_{\mathfrak{c},\mathcal{S}}\cap B}\tau_\omega
&=\sum_{\pi,\mathcal{O},c}\int\mathsf{O}_{c\mid\pi,\mathcal{O}}
\Bigl(\sum_{\mathfrak{X}\in E_{\mathfrak{c}}\cap B'}
\mathbf{O}_{\mathcal{O},\mathfrak{X}}\bigl[W^{\mathcal{S},\mathfrak{c}}_{\mathcal{O}}(\mathfrak{X})\bigr]\Bigr)\\
&\le \varepsilon_k(\mathcal{S})\sum_{\pi,\mathcal{O},c}\int\mathsf{O}_{c\mid\pi,\mathcal{O}}
\Bigl(\sum_{\mathfrak{X}'\in B'\setminus E_{\mathfrak{c}}}
\mathbf{O}_{\mathcal{O},\mathfrak{X}'}\bigl[W_{\mathcal{O}}(\mathfrak{X}')\bigr]\Bigr)\\
&=\varepsilon_k(\mathcal{S})\sum_{\omega\in B\setminus E^{(k)}_{\mathfrak{c}}}\tau_\omega .
\end{aligned}
\end{equation}
In \eqref{8.6:eq-annealed-removal-2}, $B'$ stands for $B'_{\pi,\mathcal{O},c}$. The last
equality is \eqref{8.6:eq-annealed-removal-1} read backwards, using the description of
$B\setminus E^{(k)}_{\mathfrak{c}}$ in coordinates given above.

\emph{Conclusion.} By \eqref{8.6:eq-history-measure-a} and
Lemma~\ref{8.6:lem-mass-conservation}, the annealed measure of an event is the sum of the masses
$\tau_\omega$ over that event divided by $Z=\int\rho_k$. Dividing
\eqref{8.6:eq-annealed-removal-2} by $Z$ gives the first inequality of
\eqref{8.6:eq-annealed-removal-a}.

For the second inequality, note that $B\setminus E^{(k)}_{\mathfrak{c}}\subset B$ and that
$\tilde\nu$ is a non-negative measure, so
$\tilde\nu(B\setminus E^{(k)}_{\mathfrak{c}})\le\tilde\nu(B)$. Since
$\varepsilon_k(\mathcal{S})\ge0$, multiplying by $\varepsilon_k(\mathcal{S})$ preserves this
inequality.
\end{proof}

\begin{theorem}[The large-field block bound, healed part]\label{8.6:thm-block-bound-healed}
Let $\tilde{\nu}$ be the annealed measure on the histories of an admissible run,
\eqref{8.6:eq-history-measure-a}, and let $K'$ be the terminal level of the run. For a step $k$
and a level-$k$ cell $\mathfrak{c}$ let $E^{(k)}_{\mathfrak{c}}$ be the event that a first-class
island healed at step $k$ contains $\mathfrak{c}$. For a set $\mathcal{S}$ of internal histories let
$E^{(k)}_{\mathfrak{c},\mathcal{S}}$, and the $(k,\mathfrak{c})$-blind events, be as in
Definition~\ref{8.6:def-blind-events}, \eqref{8.6:eq-blind-events-a}. Let $\theta$ be the reserve
parameter, with either $\theta=0$ or $\theta$ in the window fixed in
Definition~\ref{8.6:def-parameter-conditions}. Let $w_\theta$ be the reserve weight
\eqref{8.6:eq-reserve-weight-a}, with $w_0\equiv 1$. Let $K_H=10^{4d}$ be the constant of
\eqref{8.6:eq-healed-removal-b}. Assume the following.
\begin{itemize}
\item The positivity form of \cite[(1.73)]{B-CMP122b} at the configuration $(U,0)$, for the
normalised island operation $\mathbf{T}'_k$ restricted to sets of internal histories (the operation
of Definition~\ref{8.6:def-island-operation-bound}): for every first-class island $X$ admissible
alone, every set $\mathcal{S}$ of its internal histories and every non-negative bounded function
$\Phi$ of the internal variables of $X$ and of other variables,
\begin{equation*}
0\ \le\ \mathbf{T}'_{k,\mathcal{S}}(X)\Phi
\ \le\ \bigl(\mathbf{T}'_{k,\mathcal{S}}(X,(U,0))1\bigr)\,\sup_X\Phi
\end{equation*}
pointwise in the other variables, where $\sup_X$ is the supremum over the internal variables of $X$
on the supports of the characteristic functions of $\mathbf{T}'_{k,\mathcal{S}}(X)$.
\item The bound \cite[(1.89)]{B-CMP122b} on $\mathbf{T}'_k(X,(U,0))1$ for first-class islands $X$,
in the form supplied by the representation and bound of the large-field terms. For $\theta>0$,
assume also the per-seed reserve at an $\mathbf{R}$-step stated with clause~(i) of
Theorem~\ref{3.1:thm-main} (a reserve factor $e^{-\theta p_0(g_{j(Z)})}$ for each seed $Z$ of birth
level $j(Z)$), from which the reserve form of the island operation bound
(Lemma~\ref{8.6:lem-reserve-island-bound}) follows.
\item The activity bound \cite[(1.68)]{B-CMP122b} in the following sharpened form. For an
admissible family $\mathfrak{X}$ of islands, every polymer $Y$ of the expansion around
$\mathfrak{X}$ that is not contained in the strip $Z(\mathfrak{X})^{\sharp}$ around its large-field
region has activity
\begin{equation*}
|V_{\mathfrak{X}}(Y)|\ \le\ \alpha\,e^{-(1+2\beta)\kappa\,d_{\sharp}(Y)},
\end{equation*}
where $\alpha$ and $\beta$ are the amplitude and the exponent of \cite[(1.68)]{B-CMP122b},
auxiliary parameters of Ba{\l}aban's construction, and $d_{\sharp}(Y)$ is the size of $Y$
relative to the strip entering \cite[(1.68)]{B-CMP122b}; the bound holds uniformly in histories, in
the internal variables of the first-class islands on the supports and in $V_k$ on the support of the
term, and $V_{\mathfrak{X}}(Y)=0$ if $d_{\sharp}(Y)=\infty$.
\item The locality of the activities \cite[\S1]{B-CMP122b}: $V_{\mathfrak{X}}(Y)$ depends only on
the fields restricted to $Y$, on the variables of the polymers contained in $Y$ and on the internal
variables and histories of the islands contained in $Y$; and
$V_{\mathfrak{X}}(Y)=V_{\mathfrak{X}'}(Y)$ whenever $\mathfrak{X}$ and $\mathfrak{X}'$ have the same
islands inside $Y$ and $Y$ is a polymer of both expansions.
\item The representation \cite[(1.70)--(1.72)]{B-CMP122b} of the large-field terms as sums over
admissible families of islands, with the conditions on admissible families; among them, distinct
islands of one step are disjoint and have disjoint lineages, and a live large-field region is
separated from the islands healed before it. Assume,
together with it, the monotonicity of the common prefactor under island removal,
Lemma~\ref{8.6:lem-prefactor-monotone}, \eqref{8.6:eq-prefactor-monotone-a}, whose proof is
incomplete at one step.
\item The geometry of islands \cite[(2.1), (2.3)]{B-CMP119}:
\begin{itemize}
\item islands and live large-field components are unions of cells;
\item every island is a box of cells whose sides are at most $100$ cells long \cite{B-CMP122a};
\item the strips $S^{\sharp}$ of two components that share no point share no point, and the strip
$S^{\sharp}$ of a component lies in its enlarged strip $S^{\boxplus}$;
\item the strip $X^{\sharp}$ of an island $X$ consists of at most $(101\check{R}_k)^d$ cubes of
side $M$;
\item large-field regions nest along a history until they are healed;
\item the cells of successive levels nest, so that a level-$j$ cell $\mathfrak{c}$ determines one
level-$k$ cell $\mathfrak{c}^{(k)}$ for each $k\ge j$.
\end{itemize}
\item Positivity: at real fields every term of the history expansion is non-negative
\cite[(2.17)--(2.22)]{B-CMP119}.
\item Integral preservation \cite[(0.1), (3.1)]{B-CMP119}, \cite[(0.4), (1.102)]{B-CMP122a}:
\begin{itemize}
\item the renormalisation step preserves integrals, and so does the operation $\mathbf{R}$ for each
outer datum;
\item the cluster expansions used are identities of functions.
\end{itemize}
\item The statement that later steps do not read earlier healed islands,
Lemma~\ref{8.6:lem-healed-islands-unread}, \eqref{8.6:eq-healed-islands-unread-a}, whose proof is
incomplete at one step. This includes the statement that each step maps the level-$k$ terms to
level-$(k+1)$ terms by a finite sum of positive linear operators followed by identities of
functions.
\item The lower flow bound of clause~(0) of Theorem~\ref{3.1:thm-main}: for $i\le j\le K$,
\begin{equation*}
x_i-x_j\ \ge\ \lambda^{\sharp}(j-i)-c_5,
\end{equation*}
and $x_k\ge g^{-2}$ for every $k\le K$.
\item Two counting bounds, with $A_d=5^d[(1+3^d)\log 2+2d\log 3]$: for every integer $n\ge0$, the
number of cubes met by a tree graph of length $nM$ is at most $5^d(n+1)$; and, for every integer
$D\ge0$, the number of polymers $Y$ of the expansion around $\mathfrak{X}$ having a given cube among
their anchor cubes and with $\lfloor d_{\sharp}(Y)\rfloor=D$ is at most $2e^{A_d(D+2)}$.
\item The one-sided conditions of Definition~\ref{8.6:def-parameter-conditions} on $\kappa$ and on
$g_k$, for $\theta=0$ and for the reserve parameter $\theta$.
\item The conditions \eqref{8.6:eq-coupling-comparison-a} and \eqref{8.6:eq-coupling-comparison-b}.
\end{itemize}
Then the following hold for every $K$ and every admissible run. They hold uniformly in $K$, in $m$
and in the volume, in position, and in everything the conditioning fixes. Conditional probabilities
$\tilde{\nu}(\,\cdot\mid B)$ are taken for events $B$ with $\tilde{\nu}(B)>0$.
\begin{enumerate}
\item[(a)] For every step $k\le K'$, every level-$k$ cell $\mathfrak{c}$ and every
$(k,\mathfrak{c})$-blind event $B$,
\begin{equation}\label{8.6:eq-block-bound-healed-a}
\tilde{\nu}\bigl(E^{(k)}_{\mathfrak{c}}\bigm| B\bigr)\ \le\ K_H\,e^{-\frac32 p_0(g_k)} .
\end{equation}
The same bound holds conditionally on $V_k$, on the outer datum of step $k$ and on the
configuration of all other islands of step $k$; this quenched form is asserted for healed islands
only, not for the creation of large-field regions. It also holds unconditionally.
\item[(b)] For every set $\mathcal{S}$ of internal histories, with $k$, $\mathfrak{c}$ and $B$ as
in~(a),
\begin{equation}\label{8.6:eq-block-bound-healed-1}
\tilde{\nu}\bigl(E^{(k)}_{\mathfrak{c},\mathcal{S}}\bigm| B\bigr)
\ \le\ K_H\,w_\theta(\mathcal{S})\,e^{-\frac32 p_0(g_k)} .
\end{equation}
For $j\le k$ let $E^{(k)}_{\mathfrak{c},\le j}$ be the event that the island healed at step $k$
containing $\mathfrak{c}$ has a seed of birth level at most $j$. Then
\begin{equation}\label{8.6:eq-block-bound-healed-2}
\tilde{\nu}\bigl(E^{(k)}_{\mathfrak{c},\le j}\bigm| B\bigr)
\ \le\ K_H\,e^{\theta}\,e^{-\theta p_0(g_j)}\,e^{-\frac32 p_0(g_k)} .
\end{equation}
\item[(c)] Let $\mathfrak{c}$ be a level-$j$ cell, and let $S_*$ be the supremum over $k''\le K'$ of
the sums $\sum_{k=k''}^{K'}e^{-\frac32 p_0(g_k)}$, as in \eqref{8.6:eq-coupling-comparison-b}. For
every event $B$ that is $(k,\mathfrak{c}^{(k)})$-blind for every $k\in[j,K']$,
\begin{equation}\label{8.6:eq-block-bound-healed-3}
\tilde{\nu}\Bigl(\bigcup_{k=j}^{K'}E^{(k)}_{\mathfrak{c}^{(k)},\le j}\Bigm| B\Bigr)
\ \le\ K_H\,e^{\theta}S_*\,e^{-\theta p_0(g_j)}
\ \le\ K_H\,e^{\theta}\Bigl(1+\frac{1}{\lambda^{\sharp}}\Bigr)e^{-\theta p_0(g_j)} .
\end{equation}
\item[(d)] A healed label at $(k,\mathfrak{c})$ is an island healed at step $k$ that contains the
level-$k$ cell $\mathfrak{c}$, so that its presence is the event $E^{(k)}_{\mathfrak{c}}$. Let
$(k_1,\mathfrak{c}_1),\dots,(k_n,\mathfrak{c}_n)$ be distinct healed labels, that is, labels of
distinct islands; equal steps are allowed. Let $\mathcal{S}_1,\dots,\mathcal{S}_n$ be sets
of internal histories. For every event $B$ that is $(k_i,\mathfrak{c}_i)$-blind for every $i$,
\begin{equation}\label{8.6:eq-block-bound-healed-4}
\tilde{\nu}\Bigl(\bigcap_{i=1}^{n}E^{(k_i)}_{\mathfrak{c}_i,\mathcal{S}_i}\Bigm| B\Bigr)
\ \le\ \prod_{i=1}^{n}K_H\,w_\theta(\mathcal{S}_i)\,e^{-\frac32 p_0(g_{k_i})} .
\end{equation}
\item[(e)] Let $\mathfrak{c},\mathfrak{c}'$ be level-$k$ cells, and let
$E^{(k)}_{\{\mathfrak{c},\mathfrak{c}'\}}$ be the event that one island healed at step $k$ contains
both. Then
\begin{equation}\label{8.6:eq-block-bound-healed-5}
\begin{aligned}
&\tilde{\nu}\bigl(E^{(k)}_{\{\mathfrak{c},\mathfrak{c}'\}}\bigr)=0
&&\text{if the sup-distance of $\mathfrak{c}$ and $\mathfrak{c}'$ is at least $100$ cells,}\\
&\tilde{\nu}\bigl(E^{(k)}_{\{\mathfrak{c},\mathfrak{c}'\}}\bigr)\le K_H\,e^{-\frac32 p_0(g_k)}
&&\text{at any distance.}
\end{aligned}
\end{equation}
\item[(f)] The bound
\begin{equation}\label{8.6:eq-block-bound-healed-6}
\tilde{\nu}\bigl(E^{(k)}_{\mathfrak{c}}\cap A\bigr)\ \le\ K_H\,e^{-\frac32 p_0(g_k)}\,\tilde{\nu}(A)
\end{equation}
holds for every $(k,\mathfrak{c})$-blind event $A$. In particular it holds for every $A$ in the
$\sigma$-algebra generated by labels (large-field components of any step), of step $k$ or of other
steps, adjacent to $\mathfrak{c}$ or not, that are neither stages of the lineage of the island at
$(k,\mathfrak{c})$ nor merged into it.
This is a conditional sub-multiplicativity for healed labels, with multiplicative constant $1$ and
one-point weight $K_H\,e^{-\frac32 p_0(g_k)}$ per level-$k$ cell.
\end{enumerate}
\end{theorem}

\begin{proof}
\emph{The basic inequality.} Three statements provide the inputs of the annealing and conditioning
of the removal bound (Lemma~\ref{8.6:lem-annealed-removal}).
\begin{itemize}
\item Removal domination for a healed island (Lemma~\ref{8.6:lem-healed-removal}). For $\theta>0$
it is taken with the reserve form of the island operation bound
(Lemma~\ref{8.6:lem-reserve-island-bound}).
\item Conservation of mass along the history tree (Lemma~\ref{8.6:lem-mass-conservation}).
\item Lemma~\ref{8.6:lem-healed-islands-unread}, whose proof is incomplete at one step.
\end{itemize}
Under the hypotheses listed in the theorem these three inputs are available, so
Lemma~\ref{8.6:lem-annealed-removal} applies. Consider any $K$, any admissible run, any step
$k\le K'$, any level-$k$ cell $\mathfrak{c}$, any set $\mathcal{S}$ of internal histories and any
$(k,\mathfrak{c})$-blind event $B$. Then \eqref{8.6:eq-annealed-removal-a} gives
\begin{equation}\label{8.6:eq-block-bound-healed-7}
\tilde{\nu}\bigl(E^{(k)}_{\mathfrak{c},\mathcal{S}}\cap B\bigr)\ \le\ \varepsilon_k(\mathcal{S})\,
\tilde{\nu}(B),
\qquad
\varepsilon_k(\mathcal{S})=K_H\,w_\theta(\mathcal{S})\,e^{-\frac32 p_0(g_k)},
\end{equation}
where $\varepsilon_k(\mathcal{S})$ is as in \eqref{8.6:eq-healed-removal-b}.

\emph{Clauses \textup{(a)} and \textup{(b)}.} Dividing \eqref{8.6:eq-block-bound-healed-7} by
$\tilde{\nu}(B)>0$ gives \eqref{8.6:eq-block-bound-healed-1}.

Lemma~\ref{8.6:lem-annealed-removal} also applies with $\theta=0$, since its hypotheses at
$\theta=0$ are among those assumed. With $\theta=0$, so that $w_0\equiv 1$, and with $\mathcal{S}$
the set of all internal histories, $E^{(k)}_{\mathfrak{c},\mathcal{S}}$ is the event
$E^{(k)}_{\mathfrak{c}}$ of Definition~\ref{8.6:def-blind-events}. Hence
\eqref{8.6:eq-block-bound-healed-1} at $\theta=0$ becomes \eqref{8.6:eq-block-bound-healed-a}.

The quenched form of (a) is Lemma~\ref{8.6:lem-healed-removal} read pointwise. The domination
\eqref{8.6:eq-healed-removal-d} holds at fixed $V_k$, at fixed outer datum of step $k$ and at fixed
configuration $\mathfrak{X}'$ of the other islands of step $k$. At such fixed data it bounds the
total weight of the configurations in which an island containing $\mathfrak{c}$ is added to
$\mathfrak{X}'$ by $\varepsilon_k(\mathcal{S})$ times the weight of $\mathfrak{X}'$ itself. All
weights are non-negative by positivity, and the configuration $\mathfrak{X}'$ is among those
compatible with the conditioning, so the total weight compatible with the conditioning is at least
the weight of $\mathfrak{X}'$. The conditional probability of $E^{(k)}_{\mathfrak{c},\mathcal{S}}$
given these data is therefore at most $\varepsilon_k(\mathcal{S})$; at $\theta=0$ and with
$\mathcal{S}$ the set of all internal histories this is the bound of (a). The unconditional form of
(a) is the case in which $B$ is the sure event.

For the second bound in (b), fix $j\le k$. Let $\mathcal{S}_{\le j}$ be the set of internal
histories having a seed $Z$ of birth level $j(Z)\le j$. Then
$E^{(k)}_{\mathfrak{c},\le j}=E^{(k)}_{\mathfrak{c},\mathcal{S}_{\le j}}$. For such a seed,
Lemma~\ref{8.6:lem-coupling-comparison} under \eqref{8.6:eq-coupling-comparison-a} applies with
$i=j(Z)\le j\le K$ and gives
\begin{equation*}
p_0(g_{j(Z)})\ \ge\ p_0(g_j)-1 .
\end{equation*}
Since $\theta\ge0$, the reserve factor $e^{-\theta p_0(g_{j(Z)})}$ of that seed is therefore at
most $e^{\theta}e^{-\theta p_0(g_j)}$. The reserve factors of the remaining seeds of the history are
at most $1$, because $\theta\ge0$ and $p_0(g_i)\ge0$ for the couplings of the run, as
$\log x_i\ge\log g^{-2}>0$ under the conditions of Lemma~\ref{8.6:lem-coupling-comparison}. By the
definition \eqref{8.6:eq-reserve-weight-a} of the reserve weight this yields
\begin{equation*}
w_\theta(\mathcal{S}_{\le j})\ \le\ e^{\theta}\,e^{-\theta p_0(g_j)} .
\end{equation*}
Inserting this into \eqref{8.6:eq-block-bound-healed-1} with $\mathcal{S}=\mathcal{S}_{\le j}$ gives
\eqref{8.6:eq-block-bound-healed-2}.

\emph{Clause \textup{(c)}.} Let $\mathfrak{c}$ be a level-$j$ cell. If $j>K'$ the union in
\eqref{8.6:eq-block-bound-healed-3} is empty, so let $j\le K'$. For each $k\in[j,K']$ the cell
$\mathfrak{c}^{(k)}$ is a level-$k$ cell and $j\le k$. By hypothesis $B$ is
$(k,\mathfrak{c}^{(k)})$-blind. So \eqref{8.6:eq-block-bound-healed-2} applies at the cell
$\mathfrak{c}^{(k)}$, for every such $k$.

The quantity $S_*$ of \eqref{8.6:eq-coupling-comparison-b} is a supremum over $k''\le K'$. Its
term with $k''=j$ is $\sum_{k=j}^{K'}e^{-\frac32 p_0(g_k)}$, so this sum is at most $S_*$.
Lemma~\ref{8.6:lem-coupling-comparison} under \eqref{8.6:eq-coupling-comparison-b} gives
$S_*\le 1+1/\lambda^{\sharp}$. Subadditivity of $\tilde{\nu}(\,\cdot\mid B)$ and these facts give
\begin{align*}
\tilde{\nu}\Bigl(\bigcup_{k=j}^{K'}E^{(k)}_{\mathfrak{c}^{(k)},\le j}\Bigm| B\Bigr)
&\le \sum_{k=j}^{K'}\tilde{\nu}\bigl(E^{(k)}_{\mathfrak{c}^{(k)},\le j}\bigm| B\bigr)
\le K_H\,e^{\theta}e^{-\theta p_0(g_j)}\sum_{k=j}^{K'}e^{-\frac32 p_0(g_k)}\\
&\le K_H\,e^{\theta}S_*\,e^{-\theta p_0(g_j)}
\le K_H\,e^{\theta}\Bigl(1+\frac{1}{\lambda^{\sharp}}\Bigr)e^{-\theta p_0(g_j)} .
\end{align*}
This is \eqref{8.6:eq-block-bound-healed-3}.

\emph{Clause \textup{(d)}.} Both sides of \eqref{8.6:eq-block-bound-healed-4} are symmetric in the
labels. We may therefore number the labels so that $k_1\le k_2\le\dots\le k_n$. Write
$E_i=E^{(k_i)}_{\mathfrak{c}_i,\mathcal{S}_i}$ and $\varepsilon_i=\varepsilon_{k_i}(\mathcal{S}_i)$.
We prove, by induction on $n$, the multiplied form
\begin{equation}\label{8.6:eq-block-bound-healed-8}
\tilde{\nu}\Bigl(\bigcap_{i=1}^{n}E_i\cap B\Bigr)\ \le\ \prod_{i=1}^{n}\varepsilon_i\;\tilde{\nu}(B).
\end{equation}
It is to hold for every family of $n$ distinct healed labels ordered in this way and every $B$ that
is blind to all of them.

For $n=1$, \eqref{8.6:eq-block-bound-healed-8} is \eqref{8.6:eq-block-bound-healed-7}.

Let $n\ge2$ and put $A=B\cap\bigcap_{i<n}E_i$. We show that $A$ is $(k_n,\mathfrak{c}_n)$-blind.
The event $B$ is $(k_n,\mathfrak{c}_n)$-blind by hypothesis. Each $E_i$ with $i<n$ is
$(k_n,\mathfrak{c}_n)$-blind, by the following case distinction.
\begin{itemize}
\item If $k_i=k_n$, the two islands are distinct islands of one step. By the disjointness condition
on admissible families of islands they are disjoint and have disjoint lineages.
\item If $k_i<k_n$, the $i$-th island is healed at step $k_i$. It is therefore not part of the
lineage of the $n$-th island, which is live until step $k_n$ and is separated from it by the
separation condition on admissible families.
\end{itemize}
Hence $A$, the intersection of these $(k_n,\mathfrak{c}_n)$-blind events, is
$(k_n,\mathfrak{c}_n)$-blind (Definition~\ref{8.6:def-blind-events}).
So \eqref{8.6:eq-block-bound-healed-7}, applied at $(k_n,\mathfrak{c}_n,\mathcal{S}_n)$ with $A$ in
place of $B$, gives
\begin{equation*}
\tilde{\nu}(E_n\cap A)\le\varepsilon_n\,\tilde{\nu}(A).
\end{equation*}
The first $n-1$ labels are distinct healed labels ordered as above, and $B$ is blind to each of them.
The induction hypothesis therefore gives
\begin{equation*}
\tilde{\nu}(A)\le\prod_{i<n}\varepsilon_i\,\tilde{\nu}(B).
\end{equation*}
Since $E_n\cap A=\bigcap_{i\le n}E_i\cap B$, combining the last two inequalities gives
\eqref{8.6:eq-block-bound-healed-8}. Dividing by $\tilde{\nu}(B)>0$ and inserting the value of
$\varepsilon_i$ from \eqref{8.6:eq-block-bound-healed-7} gives \eqref{8.6:eq-block-bound-healed-4}.

\emph{Clause \textup{(e)}.} An island is a box of cells whose sides are at most $100$ cells long.
Two level-$k$ cells at sup-distance at least $100$ cells therefore never lie in one island healed at
step $k$. The event $E^{(k)}_{\{\mathfrak{c},\mathfrak{c}'\}}$ is then empty, so its
$\tilde{\nu}$-measure is $0$.

At any distance, the event $E^{(k)}_{\{\mathfrak{c},\mathfrak{c}'\}}$ is contained in
$E^{(k)}_{\mathfrak{c}}$. Its $\tilde{\nu}$-measure is therefore at
most $K_H\,e^{-\frac32 p_0(g_k)}$, by the unconditional form of (a). This is
\eqref{8.6:eq-block-bound-healed-5}.

\emph{Clause \textup{(f)}.} The event $A$ is $(k,\mathfrak{c})$-blind. We apply
\eqref{8.6:eq-block-bound-healed-7} with $\theta=0$, with $\mathcal{S}$ the set of all internal
histories, and with $A$ in place of $B$; this is \eqref{8.6:eq-block-bound-healed-6}. The events
described in the second sentence of (f) are $(k,\mathfrak{c})$-blind by
Definition~\ref{8.6:def-blind-events}.
\end{proof}

\begin{remark}
Every island healed at step $k$ and containing the level-$k$ cell $\mathfrak{c}$ is a box of cells
whose sides are at most $100$ cells long, and therefore lies in the fixed box formed by the level-$k$
cells at sup-distance at most $99$ cells from $\mathfrak{c}$. Consider a label whose region does not
meet this fixed box. Such a label is neither a stage of the lineage of an island at
$(k,\mathfrak{c})$ nor merged into it. Hence \eqref{8.6:eq-block-bound-healed-6} holds for every $A$
in the $\sigma$-algebra generated by the labels whose regions do not meet this box.

This observation uses that the pinned island is healed. The lineage of a live large-field region is
not confined to a bounded box, and for it later mergers across the complement of such a box are
not excluded.
\end{remark}

\begin{corollary}[Rarity of healed islands near a cell, without sector restriction]
\label{8.6:cor-healed-rarity}
Assume the hypotheses of Lemma~\ref{8.6:lem-healed-removal} in the case $\theta=0$. Let
$k\le K'$ be a step of an admissible run with terminal level $K'$. Fix an outer datum of
step $k$, and write $|\mathcal{O}|$ for its live large-field region. Let $\mathfrak{c}$ be a
level-$k$ cell and let $R_*>0$. Then, for every configuration $V_k$,
\begin{equation}\label{8.6:eq-healed-rarity-1}
\begin{split}
&\tilde\nu\Bigl(\operatorname{dist}\Bigl(\mathfrak{c},\,\textstyle\bigcup_e X_e\Bigr)\le R_*
\Bigm| V_k\Bigr)\\
&\qquad\le K_H\,\#\{\mathfrak{c}' : \operatorname{dist}(\mathfrak{c},\mathfrak{c}')\le R_*\}\,
e^{-\frac32 p_0(g_k)},
\end{split}
\end{equation}
with $K_H=10^{4d}$. Here $\tilde\nu$ is the measure of
Theorem~\ref{8.6:thm-block-bound-healed}. We write $\tilde\nu(\,\cdot\mid V_k)$ for the law
of the step-$k$ islands given $V_k$ and the fixed outer datum. The family $\{X_e\}$ consists
of the first-class islands healed at step $k$, $\mathfrak{c}'$ ranges over level-$k$ cells,
and $p_0(\cdot)$ is the degree function. The bound
\eqref{8.6:eq-healed-rarity-1} holds for every position of $\mathfrak{c}$: it requires no
lower bound on the distance of $\mathfrak{c}$ from $|\mathcal{O}|$. This is the sense in
which no sector restriction is imposed. The bound is uniform in $K$, $m$, the volume and the
position of $\mathfrak{c}$. The number of level-$k$ cells within distance $R_*$ of
$\mathfrak{c}$ does not depend on $\mathfrak{c}$. Consequently, for any fixed level-$k$ cell
$\mathfrak{c}_0$,
\begin{equation*}
\begin{split}
\hat q_k&:=\sup_{\mathfrak{c}}\tilde\nu\Bigl(\operatorname{dist}\Bigl(\mathfrak{c},\,
\textstyle\bigcup_e X_e\Bigr)\le R_*\Bigm| V_k\Bigr)\\
&\le K_H\,\#\{\mathfrak{c}' : \operatorname{dist}(\mathfrak{c}_0,\mathfrak{c}')\le R_*\}\,
e^{-\frac32 p_0(g_k)},
\end{split}
\end{equation*}
pointwise in $V_k$, with a constant independent of the cell.
\end{corollary}

\begin{proof}
Fix $V_k$. Each island is a union of level-$k$ cells. Hence the event
$\operatorname{dist}(\mathfrak{c},\bigcup_e X_e)\le R_*$ occurs only if some first-class island
healed at step $k$ contains a level-$k$ cell $\mathfrak{c}'$ with
$\operatorname{dist}(\mathfrak{c},\mathfrak{c}')\le R_*$. The number of such cells is the
cardinality appearing in \eqref{8.6:eq-healed-rarity-1}. Write $A_{\mathfrak{c}'}$ for the
event that a first-class island healed at step $k$ contains $\mathfrak{c}'$. By
subadditivity of the conditional measure $\tilde\nu(\,\cdot\mid V_k)$, given $V_k$ and the
outer datum,
\begin{equation}\label{8.6:eq-healed-rarity-2}
\tilde\nu\Bigl(\operatorname{dist}\Bigl(\mathfrak{c},\,\textstyle\bigcup_e X_e\Bigr)\le R_*
\Bigm| V_k\Bigr)
\le\sum_{\mathfrak{c}'\,:\,\operatorname{dist}(\mathfrak{c},\mathfrak{c}')\le R_*}
\tilde\nu\bigl(A_{\mathfrak{c}'}\bigm| V_k\bigr).
\end{equation}
We bound each term on the right of \eqref{8.6:eq-healed-rarity-2} separately.

Suppose first that $\mathfrak{c}'$ shares a point with $|\mathcal{O}|$. By the condition of
Lemma~\ref{8.6:lem-healed-removal} that no island contains a cell sharing a point with
$|\mathcal{O}|$, no island contains $\mathfrak{c}'$. The set $E_{\mathfrak{c}'}$ of
Lemma~\ref{8.6:lem-healed-removal} is therefore empty, and the
$(\mathcal{O},\mathfrak{c}',\mathcal{S})$-pinned part of $\rho_{k,\mathcal{O}}$ vanishes.
Hence $A_{\mathfrak{c}'}$ is empty and $\tilde\nu(A_{\mathfrak{c}'}\mid V_k)=0$. This
is why no restriction on the position of $\mathfrak{c}$ relative to $|\mathcal{O}|$ enters the
statement.

Suppose next that $\mathfrak{c}'$ shares no point with $|\mathcal{O}|$. We apply the removal
domination for a healed island, Lemma~\ref{8.6:lem-healed-removal}, in its pointwise form
\eqref{8.6:eq-healed-removal-d}, with the cell $\mathfrak{c}'$ in place of $\mathfrak{c}$.
The datum $\mathcal{S}$ of that lemma enters its bound only through the weight
$w_\theta(\mathcal{S})$ in \eqref{8.6:eq-healed-removal-b}, and $w_0\equiv1$; at $\theta=0$
the bound therefore does not depend on $\mathcal{S}$. The lemma gives, as non-negative
functions of $V_k$ with all internal variables integrated,
that the $(\mathcal{O},\mathfrak{c}',\mathcal{S})$-pinned part of $\rho_{k,\mathcal{O}}$ is at most
$\varepsilon_k(\mathcal{S})\,\rho_{k,\mathcal{O}}$. In the case $\theta=0$,
\eqref{8.6:eq-healed-removal-b} gives $\varepsilon_k(\mathcal{S})=K_H\,e^{-\frac32 p_0(g_k)}$.

At fixed $V_k$ with $\rho_{k,\mathcal{O}}(V_k)>0$, the quotient of the pinned part by
$\rho_{k,\mathcal{O}}$ is $\tilde\nu(A_{\mathfrak{c}'}\mid V_k)$. Hence
\begin{equation}\label{8.6:eq-healed-rarity-3}
\tilde\nu\bigl(A_{\mathfrak{c}'}\bigm| V_k\bigr)\le K_H\,e^{-\frac32 p_0(g_k)}.
\end{equation}
This is the form, conditional on $V_k$ and the outer datum, of the one-point bound
\eqref{8.6:eq-block-bound-healed-a} of the large-field block bound, healed part
(Theorem~\ref{8.6:thm-block-bound-healed}).

In both cases \eqref{8.6:eq-healed-rarity-3} holds. Inserting it into each term of
\eqref{8.6:eq-healed-rarity-2} gives
\eqref{8.6:eq-healed-rarity-1}. Since $V_k$ was arbitrary, the bound holds pointwise in $V_k$.
The translations of the level-$k$ lattice map level-$k$ cells to level-$k$ cells and
preserve distances. Hence the number of level-$k$ cells within distance $R_*$ of
$\mathfrak{c}$ is the same for every $\mathfrak{c}$, and taking the supremum over
$\mathfrak{c}$ in \eqref{8.6:eq-healed-rarity-1} gives the bound for $\hat q_k$.
\end{proof}

\begin{definition}[The live event and its decomposition by healing step]\label{8.6:def-live-event}
Healed labels, live labels, seeds, the birth level $j(Z)$ of a seed $Z$, internal histories and
lineages are those of Definition~\ref{8.6:def-labels-lineages}. The step-$k$ island family
$\mathfrak{X}_k$, the pinned island events $E^{(k)}_{\mathfrak{c},\mathcal{S}}$ and
$(k,\mathfrak{c})$-blindness are those of Definition~\ref{8.6:def-blind-events}, display
\eqref{8.6:eq-blind-events-a}. The reserve weight $w_\theta$ of a set of internal histories is
that of \eqref{8.6:eq-reserve-weight-a}.

Fix a level $k'\le K'$, where $K'$ is the terminal level up to which lineages are followed in
Definition~\ref{8.6:def-labels-lineages}, and a level-$k'$ cell $\mathfrak{c}$. For $k\ge k'$ let
$\mathfrak{c}^{(k)}$ be the level-$k$ cell containing the image of $\mathfrak{c}$. In the
definition of the \emph{live event} $\mathcal{L}^{(k')}_{\mathfrak{c}}$ below, live components of
both kinds of Definition~\ref{8.6:def-labels-lineages}, the second-class components included, and
of any birth level are admitted. For $\omega\in\mathcal{L}^{(k')}_{\mathfrak{c}}$ we write $C$
for the live component of $\omega$ at level $k'$ containing $\mathfrak{c}$; it is unique, distinct
components being disjoint, and we follow its lineage along $\omega$. For integers $k$ with
$k'<k\le K'$ we put
\begin{equation}\label{8.6:eq-live-event-a}
\begin{aligned}
\mathcal{L}^{(k')}_{\mathfrak{c}}&:=\bigl\{\omega:
\text{some live component of $\omega$ at level $k'$ contains $\mathfrak{c}$}\bigr\},\\
\mathcal{L}^{(k')}_{\mathfrak{c},k}&:=\bigl\{\omega\in\mathcal{L}^{(k')}_{\mathfrak{c}}:
\text{the lineage of $C$ is first healed at step $k$}\bigr\},\\
\mathcal{E}^{(k')}_{\mathfrak{c}}(K')&:=\bigl\{\omega\in\mathcal{L}^{(k')}_{\mathfrak{c}}:
\text{the lineage of $C$ is live at level $K'$}\bigr\},\\
\mathcal{S}_{k'}&:=\bigl\{\text{internal histories having a seed of birth level}\le k'\bigr\}.
\end{aligned}
\end{equation}

A set $B$ of histories is \emph{lineage-blind} if two conditions hold. First, $B$ is
$(k,\mathfrak{c}^{(k)})$-blind for every $k$ with $k'<k\le K'$. Second, $B$ is blind to the
terminal live component containing $\mathfrak{c}^{(K')}$: deleting that component together with
its lineage leaves $\mathbf{1}_B$ unchanged.

With these definitions, the following hold.
\begin{enumerate}
\item[(a)] The live event decomposes as a disjoint union:
\begin{equation}\label{8.6:eq-live-event-1}
\mathcal{L}^{(k')}_{\mathfrak{c}}
=\bigsqcup_{k'<k\le K'}\mathcal{L}^{(k')}_{\mathfrak{c},k}\ \sqcup\
\mathcal{E}^{(k')}_{\mathfrak{c}}(K').
\end{equation}
\item[(b)] For every $k$ with $k'<k\le K'$,
\begin{equation}\label{8.6:eq-live-event-2}
\mathcal{L}^{(k')}_{\mathfrak{c},k}\subset E^{(k)}_{\mathfrak{c}^{(k)},\mathcal{S}_{k'}}.
\end{equation}
\item[(c)] Under the hypotheses of Lemma~\ref{8.6:lem-coupling-comparison},
\begin{equation}\label{8.6:eq-live-event-3}
w_\theta(\mathcal{S}_{k'})\le e^{\theta}e^{-\theta p_0(g_{k'})}.
\end{equation}
\end{enumerate}
\end{definition}

\begin{proof}
(a) Let $\omega\in\mathcal{L}^{(k')}_{\mathfrak{c}}$. By Definition~\ref{8.6:def-labels-lineages},
the lineage of $C$ is the sequence of live components containing the image of $C$ at levels
$k',k'+1,\dots$. The sequence runs up to the first step at which that component is healed, or up
to $K'$ if it is never healed. The sequence is well defined for the following reason. Along a
history the large-field regions only grow from one level to the next until the component is
healed: this is the nesting of the regions in \cite[(2.1), (2.3)]{B-CMP119}. Moreover, a step of
Ba{\l}aban's operation $\mathbf{R}$ heals components of each class whole
(Definition~\ref{8.6:def-labels-lineages}; \cite{B-CMP122b}). Hence no part of the
component is healed before the whole of it, and no bound on the lifetime of a region is used.

At level $k'$ the component $C$ is live. So either there is a first step $k$ with
$k'<k\le K'$ at which the lineage is healed, and then $\omega\in\mathcal{L}^{(k')}_{\mathfrak{c},k}$;
or there is none,
and then the lineage is live at $K'$, that is, $\omega\in\mathcal{E}^{(k')}_{\mathfrak{c}}(K')$.
The first healing step is unique, so the sets $\mathcal{L}^{(k')}_{\mathfrak{c},k}$ are pairwise
disjoint. A lineage
healed at some step $k\le K'$ is not live at $K'$, so each $\mathcal{L}^{(k')}_{\mathfrak{c},k}$
is disjoint from
$\mathcal{E}^{(k')}_{\mathfrak{c}}(K')$. This is \eqref{8.6:eq-live-event-1}.

(b) Let $\omega\in\mathcal{L}^{(k')}_{\mathfrak{c},k}$, and let $X\in\mathfrak{X}_k(\omega)$ be the
island as part of
which the lineage of $C$ is healed at step $k$. By Definition~\ref{8.6:def-labels-lineages}, $X$
contains the level-$k$ component of
the lineage of $C$. That component contains the image of $\mathfrak{c}$, by the nesting recalled
in (a). Islands at step $k$ are unions of level-$k$ cells (Definition~\ref{8.6:def-labels-lineages}).
Hence $X$ contains the cell $\mathfrak{c}^{(k)}$.

The internal history of $X$ consists of its live ancestor components at levels $j(X),\dots,k-1$
together with every label merged into them. The component $C$ is one of these ancestors, so the
internal history of $X$ contains $C$. Every seed of $C$ was created at a step not later than the
level $k'$ at which $C$ is live, so every seed of $C$ has birth level at most $k'$. Therefore the
internal history of $X$ lies in $\mathcal{S}_{k'}$. By the definition of the pinned island events
in \eqref{8.6:eq-blind-events-a}, this gives
$\omega\in E^{(k)}_{\mathfrak{c}^{(k)},\mathcal{S}_{k'}}$, which proves
\eqref{8.6:eq-live-event-2}.

(c) Let $Z$ be a seed of birth level $j(Z)\le k'$. Since $k'$ is a level, $k'\le K$. We apply
Lemma~\ref{8.6:lem-coupling-comparison}, display \eqref{8.6:eq-coupling-comparison-a}, with
$i=j(Z)$ and $j=k'$. The comparison it rests on is $x_{j(Z)}\ge x_{k'}-c_5$, which holds in this
direction because $j(Z)\le k'$. The lemma gives
\begin{equation*}
p_0(g_{j(Z)})\ge p_0(g_{k'})-1 .
\end{equation*}
Since $\theta>0$ in the reserve weight \eqref{8.6:eq-reserve-weight-a}, multiplying by $-\theta$
reverses the inequality, and exponentiating gives
\begin{equation*}
e^{-\theta p_0(g_{j(Z)})}\le e^{\theta}e^{-\theta p_0(g_{k'})} .
\end{equation*}
Every internal history in $\mathcal{S}_{k'}$ has such a seed. Inserting this bound for that seed
into the reserve weight \eqref{8.6:eq-reserve-weight-a} gives \eqref{8.6:eq-live-event-3}.
\end{proof}

\begin{proposition}[The healed part of the live event]\label{8.6:prop-live-healed-part}
Assume the statements from which Theorem~\ref{8.6:thm-block-bound-healed} (the large-field block
bound, healed part) and Lemma~\ref{8.6:lem-annealed-removal} are derived, together with the
conditions of those derivations. These include the hypotheses of
Lemma~\ref{8.6:lem-coupling-comparison} and the reserve form \eqref{8.6:eq-reserve-island-bound-a}
of the island operation bound. Assume further the reserve-splitting statement, that is, the
statement invoked for the $\theta$-forms of the island operation bound.

Fix $K$ and an admissible run with terminal level $K'$. Let $k'\le K'$, let $\mathfrak{c}$ be a
level-$k'$ cell, and let $\mathcal{L}^{(k')}_{\mathfrak{c}}$,
$\mathcal{E}^{(k')}_{\mathfrak{c}}(K')$, the cells $\mathfrak{c}^{(k)}$ and the set
$\mathcal{S}_{k'}$ of internal histories be as in Definition~\ref{8.6:def-live-event}. Then for
every set $B$ that is lineage-blind in the sense of Definition~\ref{8.6:def-live-event},
\begin{equation}\label{8.6:eq-live-healed-part-a}
\begin{aligned}
\tilde\nu\Bigl(\bigl(\mathcal{L}^{(k')}_{\mathfrak{c}}\setminus
\mathcal{E}^{(k')}_{\mathfrak{c}}(K')\bigr)\cap B\Bigr)
&\le \sum_{k'<k\le K'}\tilde\nu\bigl(E^{(k)}_{\mathfrak{c}^{(k)},\mathcal{S}_{k'}}\cap B\bigr)\\
&\le K_H\,e^{\theta}\,S_*\,e^{-\theta p_0(g_{k'})}\,\tilde\nu(B).
\end{aligned}
\end{equation}
Here $\tilde\nu$ is the annealed measure \eqref{8.6:eq-history-measure-a}, and
$E^{(k)}_{\mathfrak{c},\mathcal{S}}$ is the pinned island event of \eqref{8.6:eq-blind-events-a}.
The constant $K_H$ is that of \eqref{8.6:eq-healed-removal-b}, and $\theta$ is the reserve
exponent of \eqref{8.6:eq-reserve-island-bound-a}. Finally, $S_*$ is the constant of
\eqref{8.6:eq-coupling-comparison-b}.
\end{proposition}

\begin{proof}
\emph{First inequality.} Definition~\ref{8.6:def-live-event} decomposes
$\mathcal{L}^{(k')}_{\mathfrak{c}}$ by the step at which the lineage of the live component
containing $\mathfrak{c}$ is first healed. For $k'<k\le K'$ we write $\mathcal{H}_k$ for the
subset of $\mathcal{L}^{(k')}_{\mathfrak{c}}$ on which that lineage is first healed at step $k$.
This gives the disjoint union
\begin{equation*}
\mathcal{L}^{(k')}_{\mathfrak{c}}
=\bigsqcup_{k'<k\le K'}\mathcal{H}_k\ \sqcup\ \mathcal{E}^{(k')}_{\mathfrak{c}}(K').
\end{equation*}
Hence $\mathcal{L}^{(k')}_{\mathfrak{c}}\setminus\mathcal{E}^{(k')}_{\mathfrak{c}}(K')$ is the
disjoint union of the sets $\mathcal{H}_k$, $k'<k\le K'$. By additivity of $\tilde\nu$,
\begin{equation*}
\tilde\nu\Bigl(\bigl(\mathcal{L}^{(k')}_{\mathfrak{c}}\setminus
\mathcal{E}^{(k')}_{\mathfrak{c}}(K')\bigr)\cap B\Bigr)
=\sum_{k'<k\le K'}\tilde\nu(\mathcal{H}_k\cap B).
\end{equation*}
By Definition~\ref{8.6:def-live-event}, $\mathcal{H}_k\subset E^{(k)}_{\mathfrak{c}^{(k)},\mathcal{S}_{k'}}$ for
each such $k$, in the notation of \eqref{8.6:eq-blind-events-a}. This inclusion holds because the
island healed at step $k$ contains the cell $\mathfrak{c}^{(k)}$. Its internal history contains
the component, and every seed of that component has birth level at most $k'$. Monotonicity of
$\tilde\nu$ then gives the first inequality in \eqref{8.6:eq-live-healed-part-a}.

\emph{Second inequality.} Fix $k$ with $k'<k\le K'$. Since $B$ is lineage-blind, it is
$(k,\mathfrak{c}^{(k)})$-blind in the sense of \eqref{8.6:eq-blind-events-a}.
Lemma~\ref{8.6:lem-annealed-removal} is applied at step $k$ to the cell $\mathfrak{c}^{(k)}$, the
set $\mathcal{S}_{k'}$ and the set $B$. It gives
\begin{equation*}
\tilde\nu\bigl(E^{(k)}_{\mathfrak{c}^{(k)},\mathcal{S}_{k'}}\cap B\bigr)
\le\varepsilon_k(\mathcal{S}_{k'})\,\tilde\nu(B),
\end{equation*}
where $\varepsilon_k(\mathcal{S})$ is the removal factor of \eqref{8.6:eq-healed-removal-b}.

We next make this bound explicit for the set $\mathcal{S}_{k'}$. The conditional reserve bound
\eqref{8.6:eq-block-bound-healed-a} of Theorem~\ref{8.6:thm-block-bound-healed}, in the reserve
form supplied by \eqref{8.6:eq-reserve-island-bound-a} together with the splitting of the reserve
exponent assumed in the first paragraph of Proposition~\ref{8.6:prop-live-healed-part}, is applied
at step $k$ to the cell $\mathfrak{c}^{(k)}$, the set $\mathcal{S}_{k'}$ and
the $(k,\mathfrak{c}^{(k)})$-blind set $B$. It gives
\begin{equation*}
\tilde\nu\bigl(E^{(k)}_{\mathfrak{c}^{(k)},\mathcal{S}_{k'}}\cap B\bigr)
\le K_H\,\varpi_\theta(\mathcal{S}_{k'})\,e^{-\frac32 p_0(g_k)}\,\tilde\nu(B),
\end{equation*}
where $\varpi_\theta(\mathcal{S})$ denotes the reserve weight of a set $\mathcal{S}$ of internal
histories appearing in \eqref{8.6:eq-block-bound-healed-a}. By Definition~\ref{8.6:def-live-event},
which uses the coupling
comparison of Lemma~\ref{8.6:lem-coupling-comparison}, we have
$\varpi_\theta(\mathcal{S}_{k'})\le e^{\theta}e^{-\theta p_0(g_{k'})}$. Here the comparison is
applied in the direction in which a seed $Z$ of birth level $j(Z)\le k'$ satisfies
$x_{j(Z)}\ge x_{k'}-c_5$, by Lemma~\ref{8.6:lem-coupling-comparison}. Summing over $k$ gives
\begin{equation*}
\sum_{k'<k\le K'}\tilde\nu\bigl(E^{(k)}_{\mathfrak{c}^{(k)},\mathcal{S}_{k'}}\cap B\bigr)
\le K_H\,e^{\theta}e^{-\theta p_0(g_{k'})}
\Bigl(\sum_{k=k'+1}^{K'}e^{-\frac32 p_0(g_k)}\Bigr)\tilde\nu(B).
\end{equation*}
The sum in parentheses is zero if $k'=K'$. Otherwise it is one of the sums over which the
supremum defining $S_*$ in \eqref{8.6:eq-coupling-comparison-b} is taken, namely the one with
$k''=k'+1\le K'$. In either case it is at most $S_*$. This proves the second inequality in
\eqref{8.6:eq-live-healed-part-a}.

\emph{The roles of the two factors.} The exponent $\theta$ in \eqref{8.6:eq-live-healed-part-a}
comes from the reserve of the oldest seed, through the bound on $\varpi_\theta(\mathcal{S}_{k'})$.
The healing step contributes the factor $e^{-\frac32 p_0(g_k)}$. Because $k$ is later than $k'$, this
factor is not itself small in $g_{k'}$. It is, however, summable in $k$ uniformly in $K$: by
Lemma~\ref{8.6:lem-coupling-comparison}, $S_*\le 1+1/\lambda^\sharp$. Thus the reserve supplies
the rate $e^{-\theta p_0(g_{k'})}$, and the summability bound supplies a constant independent of
$K$.

The exponent $\theta$ is split off from the island operation bound once. The remaining exponent
$2-\theta$ stays in the reserve form \eqref{8.6:eq-reserve-island-bound-a} and is not used again
in this proof. The splitting is applied to the normalised island operator, for which an absolute
factor is a relative one.
\end{proof}

\begin{proposition}[Removal domination for a live component. Proof outstanding]
\label{8.6:prop-live-removal}
Let $\mathcal{O}$ be an outer datum at level $k$ and let $C$ be a live component of $\mathcal{O}$.
We use the following notation.
\begin{itemize}
\item $\mathcal{O}\setminus C$ is the outer datum obtained from $\mathcal{O}$ by deleting $C$ and its
lineage, that is, by declaring the field small at every level of the internal history of $C$. It is
an outer datum, since the separation and admissibility conditions hold a fortiori.
\item $\operatorname{hull}(C)$ is the set consisting of $C$ together with the cells touching it.
\item For a set $\mathcal{S}$ of internal histories of $C$, $\rho_k(\mathcal{O};C\restriction\mathcal{S})$
is the $\mathcal{O}$-part of the level-$k$ density in the island representation
\eqref{8.6:eq-island-representation-a}, with the sum over the internal histories of $C$ restricted
to $\mathcal{S}$. For an island family $\mathfrak{X}$ of step $k$,
$\rho_k(\mathcal{O};C\restriction\mathcal{S};\mathfrak{X})$ and $\rho_k(\mathcal{O}\setminus C;\mathfrak{X})$
are the terms indexed by $\mathfrak{X}$.
\item $\int dV_k{\restriction}_{\operatorname{hull}(C)}$ denotes integration over the variables of $V_k$
on $\operatorname{hull}(C)$, all other variables being held fixed.
\end{itemize}
Let $w_\theta(\mathcal{S})$ be the reserve weight \eqref{8.6:eq-reserve-weight-a} of
Definition~\ref{8.6:def-reserve-weight}. Let $C_{\mathrm{rm}}$ be a multiple of $\alpha$ depending on
$d$ only, the constant of the boundary-term correction: $C_{\mathrm{rm}}N(\operatorname{hull}(C))$ is the
bound \eqref{8.6:eq-healed-removal-e} on the boundary-term correction $\delta$ of
Lemma~\ref{8.6:lem-healed-removal}, taken over $\operatorname{hull}(C)$; here
$N(\operatorname{hull}(C))$ is the count entering \eqref{8.6:eq-healed-removal-e}, read on
$\operatorname{hull}(C)$. Let $\kappa_1$ be the coefficient of the additional term
$-\kappa_1 d_k(\cdot)$ in the exponential, in the second sentence of \cite{B-CMP122b} quoted below. Let
$\beta_0$ be Ba{\l}aban's constant entering the term $-2(1+\beta_0)^{-1}p_0(g_k)$ of the exponent of the
bound on $\mathbf{T}'1$ that follows from \cite[(1.79), (1.80)]{B-CMP122b}. Assume the condition
\eqref{8.6:eq-parameter-conditions-c} of Definition~\ref{8.6:def-parameter-conditions}.

Then the following holds for every set $\mathcal{S}$ of internal histories of $C$, every configuration
of $V_k$ off $\operatorname{hull}(C)$ and of the internal variables of the other live components of
$\mathcal{O}$, and every $C$-blind event $B'$ of the step-$k$ island families, that is, every such event
blind to $C$ in the sense in which the event of \eqref{8.6:eq-healed-removal-d} is blind to the pinned
island:
\begin{multline}\label{8.6:eq-live-removal-a}
\int dV_k{\restriction}_{\operatorname{hull}(C)}\sum_{\mathfrak{X}\in B'}
\rho_k(\mathcal{O};C\restriction\mathcal{S};\mathfrak{X})(V_k)\\
\le w_\theta(\mathcal{S})\,
\exp\Bigl(-(2-\theta)(1+\beta_0)^{-1}p_0(g_k)-\kappa_1 d_k(C)
+C_{\mathrm{rm}}N(\operatorname{hull}(C))\Bigr)\\
\times\int dV_k{\restriction}_{\operatorname{hull}(C)}\sum_{\mathfrak{X}\in B'}
\rho_k(\mathcal{O}\setminus C;\mathfrak{X})(V_k).
\end{multline}
\end{proposition}

Proof outstanding.

\emph{Route.} Suppose $C$ is a component that Ba{\l}aban's operation $\mathbf{R}$ heals at step $k$, that
is, a component satisfying conditions (i) and (ii) of \cite{B-CMP122a}. For such a $C$ the inequality
\eqref{8.6:eq-live-removal-a} holds pointwise in $V_k$, without the integral over $\operatorname{hull}(C)$.
This is the content of \eqref{8.6:eq-healed-removal-c} and \eqref{8.6:eq-healed-removal-d} in
Lemma~\ref{8.6:lem-healed-removal}, which is proved as an implication.

A proof has to treat a live component $C$ that violates condition (i) or condition (ii)
of \cite{B-CMP122a}, that is, a component larger than a cube of side $100M\check{R}_k$, with
$\check{R}_k$ the length entering the condition \eqref{8.6:eq-parameter-conditions-c}, or carrying a
seed younger than the prescribed number of levels, or containing an old region that violates
condition (i) at its own scale; components of the second class are included, namely the components for
which the term in which the component is absent can vanish where a top window of the preparation of
route (a) below has failed. Here a seed of $C$ is a
large-field region of the internal history of $C$, read at the step that creates it; it is of either
kind, an ordinary one or a preparatory one in the sense of route (a) below. Fix the exterior. Integrate
out the internal variables of $C$ and the
top field over $\operatorname{hull}(C)$. The task is then to identify the resulting term with the factor
on the right of \eqref{8.6:eq-live-removal-a} times the term in which $C$ is absent.

For a healed component this identification is Ba{\l}aban's $\mathbf{R}$-preparation of
\cite[\S1]{B-CMP122a}. It consists of:
\begin{itemize}
\item pre-integration near $\partial C$, using the minimiser of \cite[Proposition~1]{B-CMP122a};
\item the frame average \cite[(1.100)--(1.102)]{B-CMP122a};
\item the top windows \cite[(1.3)--(1.4)]{B-CMP122a};
\item the normalised operator \cite[(1.71)]{B-CMP122a};
\item the completed action $A'_k$.
\end{itemize}
In \cite[\S1]{B-CMP122a} the preparation is carried out for components satisfying (i) and (ii).

For a component violating them, the top field on part of $C$ may lie outside the window of step $k$.
There the term in which $C$ is absent vanishes. Hence no domination pointwise in $V_k$ can hold, and the
integral over $\operatorname{hull}(C)$ in \eqref{8.6:eq-live-removal-a} is essential.

On the factor attached to the term in which $C$ is present, Ba{\l}aban states in \cite{B-CMP122b}:
\begin{quotation}
Notice that the above inequality [(1.79)] holds for all large field regions, not only for the regions
satisfying the conditions (i), (ii)
\end{quotation}
and
\begin{quotation}
the inequality (1.79) holds for the T-operation connected with an arbitrary large field region. The
inequality (1.80) holds quite generally for such regions, hence also an improved bound (1.89), with the
additional term $-\kappa_1 d_k(X)$ in the exponential. This implies the inequality (2.50) [III]
\end{quotation}
These sentences concern the upper bound on the term in which $C$ is present. The inequality
\eqref{8.6:eq-live-removal-a} requires in addition the comparison of that term with the term in which
$C$ is absent. Either of the following two routes would supply this comparison.

\smallskip
\noindent (a) Run the preparation \cite[\S1, (1.1)--(1.102)]{B-CMP122a} on a component
violating condition (ii). Such a component contains large-field characteristic functions of the
prescribed number of steps before $k$ (the age threshold of condition (ii)), which condition (ii)
excludes \cite{B-CMP122a}, and its regions need not be rectangular
parallelepipeds. The preparation is to be run with the following ingredients:
\begin{itemize}
\item the boundary pre-integration in the axial gauges of \cite[(1.100)]{B-CMP122a};
\item the analytic minimiser of \cite[Proposition~1]{B-CMP122a};
\item the frame average;
\item the top windows \cite[(1.3)--(1.4)]{B-CMP122a}, whose large-field alternatives are entered as new
large-field regions of the preparation, the preparatory seeds, each carrying its own factor
$e^{-p_0(g_k)}$ as in \cite{B-CMP122b};
\item the normalised operator \cite[(1.71)]{B-CMP122a};
\item the completion of the action $A'_k$.
\end{itemize}
For the operation $\mathbf{T}'$ attached to $C$ after this preparation, one has to establish the
inequalities \cite[(1.79), (1.80)]{B-CMP122b}, and from them
\begin{equation*}
\mathbf{T}'1\le\exp\bigl(-2(1+\beta_0)^{-1}p_0(g_k)-\kappa_1 d_k(C)\bigr).
\end{equation*}
This is the bound that the second sentence quoted above states for an arbitrary large-field region. The
three steps of the proof of Lemma~\ref{8.6:lem-healed-removal} then apply without change. They give
\eqref{8.6:eq-live-removal-a} pointwise in $V_k$ on the small-field branch of the preparation. An
induction on the number of preparatory seeds closes the argument.

\smallskip
\noindent (b) Prove, at the single level $k$ and for the hull of one component, the following local
lower bound. For a level $j$, a bond $e$ with hull $W(e)$, and a class of histories fixed outside
$W(e)$, the sum of the $\theta$-weights of the histories $h$ of the class whose seed is at $e$ is at
most $e^{C_{\mathrm{loc}}|W(e)|}$, with a constant $C_{\mathrm{loc}}$, times the sum of
$\int d\tau_{h'}$ over the histories $h'$ no seed of which meets $W(e)$. Here this is required at
level $k$ with $W(e)$ replaced by $\operatorname{hull}(C)$. Combine it with the lemma bounding the
large-field factors cube by cube in the form of \cite[(1.85)]{B-CMP122b}, and with the linear removal
bound, which, where the term $-\kappa_1 d_k(C)$ is not used, pays instead the charge
$C_{\mathrm{rm}}N(\operatorname{hull}(C))$ of \eqref{8.6:eq-live-removal-a}; both are proved as
implications, and through them such a local bound yields the large-field block bound, healed part.

\smallskip
Either way the large-field block bound, healed part, requires no lower bound on the activity of a
healed island, no cluster expansion and no analyticity: only positivity and one normalised operator
bound per live component enter.

\begin{proposition}[The terminal live member]\label{8.6:prop-terminal-live}
Assume the following statements:
\begin{itemize}
\item[(a)] removal domination for a live component, Proposition~\ref{8.6:prop-live-removal}, at the terminal level $K'$; its proof is outstanding;
\item[(b)] conservation of mass along the history tree, Lemma~\ref{8.6:lem-mass-conservation};
\item[(c)] the positivity dichotomy of the large-field operations, in the form that the level-$K'$ integrand remaining after the integration over the hull of a live component, and the masses $\tau_\omega$, are non-negative;
\item[(d)] the geometry of live components: a live component is a touch-connected union of cells; a union $C$ of $n$ cells consists of $N(C)=n\check{R}_{K'}^{\,d}$ $M$-cubes, where $\check{R}_{K'}$ is the quantity of \eqref{8.6:eq-parameter-conditions-c} at level $K'$, so that a level-$K'$ cell consists of $\check{R}_{K'}^{\,d}$ $M$-cubes; and $N(\operatorname{hull}C)\le 3^d n\check{R}_{K'}^{\,d}$;
\item[(e)] two lemmas of the representation and bound of the large-field terms: the lemma by which at most $9^{dn}$ touch-connected unions of $n$ cells contain a given cell, and the lemma bounding $d_{K'}(C)$ below by $5^{-d}N(C)-1$;
\item[(f)] the condition \eqref{8.6:eq-parameter-conditions-c} at level $K'$, and the coupling comparison of Lemma~\ref{8.6:lem-coupling-comparison}.
\end{itemize}
Let $k'\le K'$, let $\mathfrak{c}$ be a level-$k'$ cell, and let $\mathfrak{c}^{(K')}$ be the level-$K'$ cell containing its image. Let $B$ be lineage-blind in the sense of Definition~\ref{8.6:def-live-event}, and determined by the labels of levels $\le K'$ and by $V_{K'}$ off the hull of the terminal live component containing $\mathfrak{c}^{(K')}$. Then
\begin{equation}\label{8.6:eq-terminal-live-a}
\begin{split}
\tilde{\nu}\bigl(\mathcal{E}^{(k')}_{\mathfrak{c}}(K')\cap B\bigr)
&\le 2e^{\kappa_1+\theta}\,e^{-\theta p_0(g_{k'})}\,
e^{-(2-\theta)(1+\beta_0)^{-1}p_0(g_{K'})}\,\tilde{\nu}(B)\\
&\le e^{-\theta p_0(g_{k'})}\,\tilde{\nu}(B).
\end{split}
\end{equation}
The second inequality holds whenever the reference coupling $g$ satisfies the one-sided ceiling
\begin{equation}\label{8.6:eq-terminal-live-1}
2e^{\kappa_1+\theta}\,e^{-(2-\theta)(1+\beta_0)^{-1}(p_0(g)-1)}\le 1,
\end{equation}
a condition requiring $X=g^{-2}$ to be large; it follows from this ceiling together with $p_0(g_{K'})\ge p_0(g)-1$ and $x_{K'}\ge X$.
\end{proposition}

\begin{proof}
\emph{The terminal mass.} The level $K'$ is terminal, so no later step is performed. Conservation of mass along the history tree (Lemma~\ref{8.6:lem-mass-conservation}) applied at $k=K'$ gives
\begin{equation*}
Z=\int dV_{K'}\sum_{\omega}\rho_{K',\omega}(V_{K'}).
\end{equation*}
It also gives that, for an event $E$ determined by labels of levels $\le K'$, the measure $\tilde{\nu}(E)$ of \eqref{8.6:eq-history-measure-a} equals $Z^{-1}$ times the sum of the integrals of the level-$K'$ terms over $E$. By Lemma~\ref{8.6:lem-mass-conservation} the mass $\tau_\omega$ of \eqref{8.6:eq-history-measure-a} is the integral of the level-$K'$ term of $\omega$, so that $Z\,\tilde{\nu}(E)=\sum_{\omega\in E}\tau_\omega$.

\emph{The structure on the terminal event.} Let $\omega\in\mathcal{E}^{(k')}_{\mathfrak{c}}(K')\cap B$. By Definition~\ref{8.6:def-live-event}, the lineage of the live component of $\mathfrak{O}_{k'}(\omega)$ containing $\mathfrak{c}$ is still live at $K'$. Hence the terminal outer datum $\mathfrak{O}$ of $\omega$ contains a live component $C_{\mathrm{fin}}\supseteq\mathfrak{c}^{(K')}$ whose internal history contains the whole lineage, and in particular a seed of birth level $\le k'$. Let $\mathcal{S}_{k'}$ be the set of internal histories with a seed of birth level $\le k'$. The internal history of $C_{\mathrm{fin}}$ lies in $\mathcal{S}_{k'}$. The coupling comparison of Lemma~\ref{8.6:lem-coupling-comparison}, in the form recorded in Definition~\ref{8.6:def-live-event}, bounds the reserve weight \eqref{8.6:eq-reserve-weight-a} of this set:
\begin{equation}\label{8.6:eq-terminal-live-2}
w_\theta(\mathcal{S}_{k'})\le e^{\theta-\theta p_0(g_{k'})}.
\end{equation}

\emph{Decomposition.} We decompose each such $\omega$ by the pair $(\mathfrak{O}',C_{\mathrm{fin}})$ with $\mathfrak{O}':=\mathfrak{O}\setminus C_{\mathrm{fin}}$.

For given labels other than these, the map $\omega\mapsto(\mathfrak{O}',C_{\mathrm{fin}})$ is injective on the histories in $\mathcal{E}^{(k')}_{\mathfrak{c}}(K')\cap B$. The event $B$ is blind to the terminal live component containing $\mathfrak{c}^{(K')}$, that is, to $C_{\mathrm{fin}}$. Deleting $C_{\mathrm{fin}}$ with its lineage therefore leaves $\mathbf{1}_B$ unchanged. The datum $\mathfrak{O}'$ ranges over the outer data none of whose live components contains $\mathfrak{c}^{(K')}$ through the lineage.

Fix $\mathfrak{O}'$, the labels other than those of $C_{\mathrm{fin}}$, the configuration of $V_{K'}$ off $\operatorname{hull}(C_{\mathrm{fin}})$, and the internal variables of the other live components. Since $B$ is determined by $V_{K'}$ off the hull, the factor $\mathbf{1}_B$ is constant in the variables integrated next.

We first integrate over $V_{K'}\restriction_{\operatorname{hull}(C_{\mathrm{fin}})}$ and over the internal variables of $C_{\mathrm{fin}}$. We then apply Proposition~\ref{8.6:prop-live-removal}, whose proof is outstanding, at $k=K'$. It is applied with the live component $C=C_{\mathrm{fin}}$, with the set of internal histories $\mathcal{S}_{k'}$, and with the section of $B$ as the $C$-blind event of the step-$K'$ island families. By \eqref{8.6:eq-live-removal-a}, the integral is bounded by
\begin{equation*}
w_\theta(\mathcal{S}_{k'})\,
e^{-(2-\theta)(1+\beta_0)^{-1}p_0(g_{K'})-\kappa_1 d_{K'}(C_{\mathrm{fin}})
+C_{\mathrm{rm}}N(\operatorname{hull}C_{\mathrm{fin}})}
\end{equation*}
times the corresponding integral of the term with outer datum $\mathfrak{O}'$.

We then integrate over the remaining variables. By the positivity dichotomy of the large-field operations, the remaining integrand is non-negative, and it is common to both sides. Fubini's theorem therefore carries the inequality through this integration. Summing over $C_{\mathrm{fin}}\ni\mathfrak{c}^{(K')}$ and inserting \eqref{8.6:eq-terminal-live-2}, we obtain
\begin{equation}\label{8.6:eq-terminal-live-3}
\sum_{\omega\in\mathcal{E}^{(k')}_{\mathfrak{c}}(K')\cap B}\tau_\omega
\le\sum_{(\mathfrak{O}',\dots)\in B}
\Bigl[\sum_{C\ni\mathfrak{c}^{(K')}}\varepsilon(C)\Bigr]\tau_{(\mathfrak{O}',\dots)},
\end{equation}
where
\begin{equation*}
\varepsilon(C):=e^{\theta-\theta p_0(g_{k'})}\,
\exp\bigl(-(2-\theta)(1+\beta_0)^{-1}p_0(g_{K'})-\kappa_1d_{K'}(C)
+C_{\mathrm{rm}}N(\operatorname{hull}C)\bigr).
\end{equation*}

\emph{The count of shapes.} By hypothesis~(d), a live component $C\ni\mathfrak{c}^{(K')}$ is a touch-connected union of cells. Let $n\ge1$ be its number of cells. By the first lemma in~(e), applied at cell scale, there are at most $9^{dn}$ such shapes through the given cell $\mathfrak{c}^{(K')}$. By~(d), $N(C)=n\check{R}_{K'}^{\,d}$. The second lemma in~(e) then gives $d_{K'}(C)\ge5^{-d}n\check{R}_{K'}^{\,d}-1$. By~(d) again, $N(\operatorname{hull}C)\le3^dn\check{R}_{K'}^{\,d}$. Hence
\begin{equation*}
-\kappa_1d_{K'}(C)+C_{\mathrm{rm}}N(\operatorname{hull}C)
\le\kappa_1-\bigl(\kappa_15^{-d}-C_{\mathrm{rm}}3^d\bigr)n\check{R}_{K'}^{\,d}.
\end{equation*}
The condition \eqref{8.6:eq-parameter-conditions-c} at level $K'$ states
\begin{equation*}
\bigl(\kappa_15^{-d}-C_{\mathrm{rm}}3^d\bigr)\check{R}_{K'}^{\,d}\ge d\log9+\log2 .
\end{equation*}
It therefore gives
\begin{equation*}
9^{dn}e^{-(\kappa_15^{-d}-C_{\mathrm{rm}}3^d)n\check{R}_{K'}^{\,d}}\le 2^{-n},
\end{equation*}
and so
\begin{equation*}
\begin{split}
\sum_{C\ni\mathfrak{c}^{(K')}}e^{-\kappa_1d_{K'}(C)+C_{\mathrm{rm}}N(\operatorname{hull}C)}
&\le e^{\kappa_1}\sum_{n\ge1}9^{dn}
e^{-(\kappa_15^{-d}-C_{\mathrm{rm}}3^d)n\check{R}_{K'}^{\,d}}\\
&\le e^{\kappa_1}\sum_{n\ge1}2^{-n}\le 2e^{\kappa_1}.
\end{split}
\end{equation*}
Consequently the bracket in \eqref{8.6:eq-terminal-live-3} is at most
\begin{equation*}
2e^{\kappa_1+\theta}\,e^{-\theta p_0(g_{k'})}\,
e^{-(2-\theta)(1+\beta_0)^{-1}p_0(g_{K'})}.
\end{equation*}

\emph{Division by $Z$.} The bound on the bracket is uniform in $(\mathfrak{O}',\dots)$. The masses are non-negative by the positivity dichotomy of the large-field operations, and the histories $(\mathfrak{O}',\dots)$ on the right of \eqref{8.6:eq-terminal-live-3} form a sub-family of those contributing to $Z\,\tilde{\nu}(B)$ by Lemma~\ref{8.6:lem-mass-conservation}. Hence
\begin{equation*}
\sum_{(\mathfrak{O}',\dots)\in B}\tau_{(\mathfrak{O}',\dots)}
\le\sum_{\omega\in B}\tau_\omega
=Z\,\tilde{\nu}(B).
\end{equation*}
The left side of \eqref{8.6:eq-terminal-live-3} equals $Z\,\tilde{\nu}(\mathcal{E}^{(k')}_{\mathfrak{c}}(K')\cap B)$ by the first paragraph. Dividing by $Z$ gives the first inequality of \eqref{8.6:eq-terminal-live-a}.

\emph{The ceiling.} For the second inequality we use $x_{K'}\ge X$, the lower bound on the running inverse couplings along the run in clause~(0) of Theorem~\ref{3.1:thm-main}, and, in the form supplied by the coupling comparison of Lemma~\ref{8.6:lem-coupling-comparison}, $p_0(g_{K'})\ge p_0(g)-1$. The coefficient $(2-\theta)(1+\beta_0)^{-1}$ is positive, so
\begin{equation*}
e^{-(2-\theta)(1+\beta_0)^{-1}p_0(g_{K'})}\le e^{-(2-\theta)(1+\beta_0)^{-1}(p_0(g)-1)}.
\end{equation*}
By \eqref{8.6:eq-terminal-live-1}, the factor $2e^{\kappa_1+\theta}$ is then absorbed:
\begin{equation*}
2e^{\kappa_1+\theta}\,e^{-(2-\theta)(1+\beta_0)^{-1}p_0(g_{K'})}\le1.
\end{equation*}
This gives the second inequality of \eqref{8.6:eq-terminal-live-a}.
\end{proof}

\begin{theorem}[The large-field block bound, live part]\label{8.6:thm-block-bound-live}
Let $\tilde{\nu}$ be the annealed measure on histories $h$ of
\eqref{8.6:eq-history-measure-a}. For events $A$, $B$ and a number $c\ge0$, the inequality
$\tilde{\nu}(A\mid B)\le c$ means $\tilde{\nu}(A\cap B)\le c\,\tilde{\nu}(B)$. Cells $\mathfrak{c}$,
labels, live components $C$, their seeds and birth levels, and lineage-blind and $C$-blind
events are as in Definition~\ref{8.6:def-live-event}. The event that a live component at level
$k'$ contains $\mathfrak{c}$ is $\mathcal{L}^{(k')}_{\mathfrak{c}}$ of \eqref{8.6:eq-live-event-a}.
The constant $K_H=10^{4d}$ is that of Theorem~\ref{8.6:thm-block-bound-healed}, and $\theta$,
$\beta_0$, $\kappa_1$, $\check{R}_k$ and the degree function $p_0(\cdot)$ are as there. Let $K'$
denote the terminal level of the longest admissible run, the one with
$m=m_{\mathbb{T}}(g)$. Assume:
\begin{itemize}
\item[($\alpha$)] the statements and conditions from which
Theorem~\ref{8.6:thm-block-bound-healed} is proved hold. In particular its conclusions are
available, and so are the conditions \eqref{8.6:eq-coupling-comparison-a} and
\eqref{8.6:eq-coupling-comparison-b} of Lemma~\ref{8.6:lem-coupling-comparison};
\item[($\beta$)] the $\theta$-reserve statement named among the antecedents of
Theorem~\ref{8.6:thm-block-bound-healed} for its $\theta$-forms holds;
\item[($\gamma$)] the torus-independence statement of the correlation theorem holds: the $k$-th
step of an admissible run reads only data independent of $m$, namely the
coupling $g_k$, the small-field parameter $\varepsilon_k$, the radius $\check{R}_k$, the
partitions of the level-$k$ lattice and the frozen running shift. The parameter $m$ enters only
through the terminal level and through the lower threshold condition of the run. Consequently
two admissible runs with parameters $m'\le m$ coincide
through step $K-m$, and for every event $F$ determined by labels of levels $\le K-m$ the number
$\tilde{\nu}(F)$ is the same in both;
\item[($\delta$)] the removal domination for a live component,
Proposition~\ref{8.6:prop-live-removal}, whose proof is outstanding, holds at the single level
$K'$;
\item[($\varepsilon$)] the conditions \eqref{8.6:eq-parameter-conditions-c} hold.
\end{itemize}
Then statements (a), (c) and (d) below hold; statement (b), which needs different hypotheses, is
stated after them.
\begin{enumerate}
\item[(a)] For every $K$, every level $k'\le K'$, every level-$k'$ cell $\mathfrak{c}$ and every
lineage-blind event $B$ determined by labels of levels $\le K'$,
\begin{equation}\label{8.6:eq-block-bound-live-a}
\tilde{\nu}\bigl(\mathcal{L}^{(k')}_{\mathfrak{c}}\bigm| B\bigr)\le C_L\,e^{-\theta p_0(g_{k'})},
\qquad C_L:=K_He^{\theta}\Bigl(1+\frac{1}{\lambda^\sharp}\Bigr)+1,
\end{equation}
uniformly in $K$, in the volume and in the position of $\mathfrak{c}$.
\item[(c)] For every level $j$, every cube $\square$ of the level-$j$ partition, consisting of
$n_\square$ cells, and every kind of large-field creation,
\begin{equation}\label{8.6:eq-block-bound-live-1}
\tilde{\nu}\bigl(h\ \text{creates a seed at}\ (j,\square)\bigr)\le n_\square\,C_L\,
e^{-\theta p_0(g_j)}.
\end{equation}
This is the creation-event bound in the form $C_qe^{-c_qp_0(g_j)}$ used by the Gaussian
extraction, with $C_q=n_\square C_L$ and $c_q=\theta$, at every level $j$, coarser or finer than
the level of the observable.
\item[(d)] Let $k_0\le K'$ be a level and $\mathcal{N}$ a set of points of the torus. Let
$n_k(\mathcal{N})$ be the number of level-$k$ cells meeting $\mathcal{N}$. Let $\Xi_{\mathcal{N}}$
be the union of two events: that some label of $h$ live at level $k_0$ meets $\mathcal{N}$, and
that $V_{k_0}$ leaves the level-$k_0$ small-field window somewhere on $\mathcal{N}$. Then
\begin{equation}\label{8.6:eq-block-bound-live-2}
\tilde{\nu}(\Xi_{\mathcal{N}})\le\sum_{\substack{\mathfrak{c}\ \text{level-}k_0\ \text{cell}\\
\mathfrak{c}\cap\mathcal{N}\neq\emptyset}}\tilde{\nu}\bigl(\mathcal{L}^{(k_0)}_{\mathfrak{c}}\bigr)
\le n_{k_0}(\mathcal{N})\,C_L\,e^{-\theta p_0(g_{k_0})}.
\end{equation}
Let $\mathfrak{l}\ge0$ be an integer. Let
$\Xi_{\mathcal{N}}^{\mathfrak{l}}$ be the union of $\Xi_{\mathcal{N}}$ with the event that some
label healed at a step $k\in[k_0-\mathfrak{l},k_0]$ meets $\mathcal{N}$. Then
\begin{equation}\label{8.6:eq-block-bound-live-3}
\tilde{\nu}\bigl(\Xi_{\mathcal{N}}^{\mathfrak{l}}\bigr)\le n_{k_0}(\mathcal{N})\,C_L\,
e^{-\theta p_0(g_{k_0})}
+e^{3/2}K_H\sum_{k=\max(0,k_0-\mathfrak{l})}^{k_0}n_k(\mathcal{N})\,e^{-\frac32p_0(g_{k_0})}.
\end{equation}
The second summand uses hypothesis~($\alpha$) only.
When $\mathcal{N}$ is one of the neighbourhoods $N_i$ of the two-point witness theorem and $k_0$
the level of its observable, $\Xi_{\mathcal{N}}$ is the contact event $\mathcal{B}_i$ of that
theorem, and in the form of $\mathcal{B}_i$ that carries in addition the clause on healed labels
it is $\Xi_{\mathcal{N}}^{\mathfrak{l}}$ with $\mathfrak{l}=L\check{R}_{k_0}+c_L$, where $c_L$ is
the constant of that theorem; neither carries a restriction on birth levels, the lifetime entering
the estimate and not the event.
\end{enumerate}
Now fix a level $k'$ and assume, in place of ($\alpha$), ($\gamma$), ($\delta$) and
($\varepsilon$):
\begin{itemize}
\item[($\zeta$)] Proposition~\ref{8.6:prop-live-removal}, whose proof is outstanding, holds at
the level $k'$ itself;
\item[($\eta$)] the conservation of mass, Lemma~\ref{8.6:lem-mass-conservation}, holds, and so do
the positivity statement and the two estimates on the large-field terms named among the
antecedents of Proposition~\ref{8.6:prop-terminal-live}, the geometry of the large-field regions
used there, and the condition \eqref{8.6:eq-coupling-comparison-a} of
Lemma~\ref{8.6:lem-coupling-comparison};
\item[($\varepsilon'$)] the conditions \eqref{8.6:eq-parameter-conditions-c} hold with
$\frac12\kappa_1$ in place of $\kappa_1$;
\end{itemize}
together with ($\beta$). Then statement (b) holds at that level.
\begin{enumerate}
\item[(b)] Fix a level-$k'$ cell $\mathfrak{c}$, a level $j$ and a real number $D$. Let $B$ be
any event that is $C$-blind and determined by labels of levels $\le k'$ and by $V_{k'}$ off the
hull of $C$. Then
\begin{equation}\label{8.6:eq-block-bound-live-4}
\begin{split}
&\tilde{\nu}\bigl(\text{a live component $C$ at level $k'$ contains $\mathfrak{c}$, has a seed of
birth level $\le j$,}\\
&\qquad\text{and has } d_{k'}(C)\ge D\bigm| B\bigr)\\
&\quad\le2e^{\kappa_1+\theta}\,e^{-\theta p_0(g_j)}\,
e^{-(2-\theta)(1+\beta_0)^{-1}p_0(g_{k'})}\,e^{-\frac12\kappa_1D}.
\end{split}
\end{equation}
In particular, for two level-$k'$ cells $\mathfrak{c}$, $\mathfrak{c}'$ at distance $D$ cells,
\begin{equation}\label{8.6:eq-block-bound-live-5}
\begin{split}
&\tilde{\nu}\bigl(\text{one live component at level $k'$ contains $\mathfrak{c}$ and
$\mathfrak{c}'$}\bigm| B\bigr)\\
&\quad\le2e^{2\kappa_1+\theta}\,e^{-\theta p_0(g_{k'})}\,
e^{-(2-\theta)(1+\beta_0)^{-1}p_0(g_{k'})}\,
e^{-\frac12\kappa_1\,5^{-d}\check{R}_{k'}^{\,d}D}.
\end{split}
\end{equation}
\end{enumerate}
\end{theorem}

\begin{proof}
\emph{Proof of (a).} Fix $K$, a level $k'\le K'$, a level-$k'$ cell $\mathfrak{c}$ and a
lineage-blind event $B$ determined by labels of levels $\le K'$.

We first fix the run. We state and prove \eqref{8.6:eq-block-bound-live-a} in the longest
admissible run, with $m=m_{\mathbb{T}}(g)$ and terminal level $K'$, and nothing is lost thereby.
By Lemma~\ref{8.6:lem-mass-conservation}, $\tilde{\nu}$
of an event determined by labels of levels $\le k$ equals $Z^{-1}$ times the sum of the integrals
of the level-$k$ terms over that event, whatever later steps are performed. By
hypothesis~($\gamma$), two admissible runs with parameters $m'\le m$ coincide through step $K-m$.
Hence, for an event determined by labels of levels $\le K-m$, its $\tilde{\nu}$-value in any
admissible run with parameter $m$ is the same number as in the longest admissible run, and the
left member of \eqref{8.6:eq-block-bound-live-a} does not depend on where the run stops.

This choice is a device of the proof. The measure $\tilde{\nu}$ is fixed; only the level at which
its representation stops moves. In particular, when the bound is applied to events of the
two-point witness theorem, which are determined by labels of levels not exceeding the witness
depth, it is read in the longest admissible run down to that run's own terminal level. The choice
of $m$ made in the two-point witness theorem, its measure $\tilde{\nu}$, its lower threshold
condition and its depth floor are left unchanged.

We next split the live event. Definition~\ref{8.6:def-live-event} decomposes the live event by
healing step into its healed part $\mathcal{L}^{(k')}_{\mathfrak{c}}\setminus
\mathcal{E}^{(k')}_{\mathfrak{c}}(K')$ and its terminal member
$\mathcal{E}^{(k')}_{\mathfrak{c}}(K')$. Since
$\mathcal{L}^{(k')}_{\mathfrak{c}}\subset(\mathcal{L}^{(k')}_{\mathfrak{c}}\setminus
\mathcal{E}^{(k')}_{\mathfrak{c}}(K'))\cup\mathcal{E}^{(k')}_{\mathfrak{c}}(K')$,
\begin{equation}\label{8.6:eq-block-bound-live-6}
\tilde{\nu}\bigl(\mathcal{L}^{(k')}_{\mathfrak{c}}\cap B\bigr)\le
\tilde{\nu}\bigl((\mathcal{L}^{(k')}_{\mathfrak{c}}\setminus
\mathcal{E}^{(k')}_{\mathfrak{c}}(K'))\cap B\bigr)
+\tilde{\nu}\bigl(\mathcal{E}^{(k')}_{\mathfrak{c}}(K')\cap B\bigr).
\end{equation}

For the healed part we use Proposition~\ref{8.6:prop-live-healed-part}. Its inputs are the list
of Theorem~\ref{8.6:thm-block-bound-healed} and the $\theta$-reserve statement, which are
hypotheses ($\alpha$) and ($\beta$), and it applies to every lineage-blind $B$. It bounds the
healed part by
the sum, over the healing steps $k$ with $k'<k\le K'$, of the measures of the corresponding
healed-island events intersected with $B$. It then bounds that sum by
\begin{equation*}
K_He^{\theta}S_*\,e^{-\theta p_0(g_{k'})}\,\tilde{\nu}(B).
\end{equation*}
By Lemma~\ref{8.6:lem-coupling-comparison}, whose conditions are part of ($\alpha$), the quantity
$S_*$ of \eqref{8.6:eq-coupling-comparison-b} satisfies $S_*\le1+1/\lambda^\sharp$. Therefore
\begin{equation}\label{8.6:eq-block-bound-live-7}
\tilde{\nu}\bigl((\mathcal{L}^{(k')}_{\mathfrak{c}}\setminus
\mathcal{E}^{(k')}_{\mathfrak{c}}(K'))\cap B\bigr)\le K_He^{\theta}
\Bigl(1+\frac{1}{\lambda^\sharp}\Bigr)e^{-\theta p_0(g_{k'})}\,\tilde{\nu}(B).
\end{equation}

For the terminal member we use Proposition~\ref{8.6:prop-terminal-live}. Its inputs are:
\begin{itemize}
\item the removal domination of Proposition~\ref{8.6:prop-live-removal} at level $K'$, which is
hypothesis~($\delta$) (the proof of Proposition~\ref{8.6:prop-live-removal} is outstanding);
\item Lemma~\ref{8.6:lem-mass-conservation};
\item the positivity statement, the geometry of the large-field regions, and the two estimates on
the large-field terms, all of which are among the inputs of
Theorem~\ref{8.6:thm-block-bound-healed} and hence part of ($\alpha$);
\item conditions \eqref{8.6:eq-parameter-conditions-c} and
\eqref{8.6:eq-coupling-comparison-a}, which are hypotheses ($\varepsilon$) and ($\alpha$).
\end{itemize}
Our $B$ is lineage-blind and determined by labels of levels $\le K'$, so it is of the kind
admitted there. The second inequality of Proposition~\ref{8.6:prop-terminal-live} absorbs the
constant $2e^{\kappa_1+\theta}$ into $e^{-(2-\theta)(1+\beta_0)^{-1}p_0(g_{K'})}$ by a ceiling
condition on $\gamma$ read at the terminal coupling $g_{K'}$, as in that proposition. It gives
\begin{equation}\label{8.6:eq-block-bound-live-8}
\tilde{\nu}\bigl(\mathcal{E}^{(k')}_{\mathfrak{c}}(K')\cap B\bigr)\le
e^{-\theta p_0(g_{k'})}\,\tilde{\nu}(B).
\end{equation}
Inserting \eqref{8.6:eq-block-bound-live-7} and \eqref{8.6:eq-block-bound-live-8} into
\eqref{8.6:eq-block-bound-live-6} gives
\begin{equation*}
\tilde{\nu}\bigl(\mathcal{L}^{(k')}_{\mathfrak{c}}\cap B\bigr)\le
\Bigl(K_He^{\theta}\Bigl(1+\frac{1}{\lambda^\sharp}\Bigr)+1\Bigr)
e^{-\theta p_0(g_{k'})}\,\tilde{\nu}(B),
\end{equation*}
which is \eqref{8.6:eq-block-bound-live-a}. The summand $1$ in $C_L$ is exactly the contribution
\eqref{8.6:eq-block-bound-live-8} of the terminal member. All constants are independent of $K$,
of the volume and of the position of $\mathfrak{c}$, and so is the bound.

\emph{Proof of (b).} Fix $k'$, $\mathfrak{c}$, $j$, $D$ and $B$ as in (b). The bound
\eqref{8.6:eq-block-bound-live-4} results from the computation that proves the first inequality of
Proposition~\ref{8.6:prop-terminal-live}, carried out at the level $k'$ instead of $K'$. The
removal domination is taken from hypothesis~($\zeta$), at level $k'$, for the live component $C$
at level $k'$ that contains $\mathfrak{c}$ and has a seed of birth level $\le j$. The remaining
inputs of that computation are those listed in ($\eta$), among them the geometry of the
large-field regions and \eqref{8.6:eq-coupling-comparison-a}, together with ($\beta$). The reserve
of the seed of birth level $\le j$ takes the place of that of the
oldest seed of the terminal member, so that the factor $e^{-\theta p_0(g_j)}$ replaces
$e^{-\theta p_0(g_{k'})}$. The factor supplied by hypothesis~($\zeta$) carries
$e^{-(2-\theta)(1+\beta_0)^{-1}p_0(g_{k'})}$ at the level $k'$, together with
$e^{-\kappa_1d_{k'}(C)}$. We write
\begin{equation*}
e^{-\kappa_1d_{k'}(C)}=e^{-\frac12\kappa_1d_{k'}(C)}\,e^{-\frac12\kappa_1d_{k'}(C)}.
\end{equation*}
On the event $d_{k'}(C)\ge D$ the first factor is at most $e^{-\frac12\kappa_1D}$, and it is kept.
The second factor plays, in the computation of Proposition~\ref{8.6:prop-terminal-live}, the part
played there by $e^{-\kappa_1d(C)}$; accordingly the conditions
\eqref{8.6:eq-parameter-conditions-c} are read with $\frac12\kappa_1$ in place of $\kappa_1$,
which is hypothesis~($\varepsilon'$). The
computation then yields the constant $2e^{\kappa_1+\theta}$ of that proposition. This gives
\eqref{8.6:eq-block-bound-live-4}.

For the two-cell form we use the geometry of the large-field regions assumed in ($\eta$): if two
level-$k'$ cells $\mathfrak{c}$,
$\mathfrak{c}'$ at distance $D$ cells lie in one live component $C$ at level $k'$, then
\begin{equation*}
d_{k'}(C)\ge D':=5^{-d}(D+1)\check{R}_{k'}^{\,d}-1.
\end{equation*}
Every seed of a component live at level $k'$ has birth level $\le k'$. The event in
\eqref{8.6:eq-block-bound-live-5} is therefore contained in the event of
\eqref{8.6:eq-block-bound-live-4} with $j:=k'$ and with $D'$ in place of $D$. Moreover,
\begin{equation*}
e^{-\frac12\kappa_1D'}=e^{\frac12\kappa_1}\,e^{-\frac12\kappa_1\,5^{-d}\check{R}_{k'}^{\,d}}\,
e^{-\frac12\kappa_1\,5^{-d}\check{R}_{k'}^{\,d}D}\le e^{\frac12\kappa_1}\,
e^{-\frac12\kappa_1\,5^{-d}\check{R}_{k'}^{\,d}D},
\end{equation*}
and $2e^{\kappa_1+\theta}e^{\frac12\kappa_1}\le2e^{2\kappa_1+\theta}$ because $\kappa_1>0$.
Substituting these into \eqref{8.6:eq-block-bound-live-4} gives
\eqref{8.6:eq-block-bound-live-5}. For healed labels the corresponding two-cell statement is the
clause of Theorem~\ref{8.6:thm-block-bound-healed} according to which two cells $100$ apart never
lie in one island.

\emph{Proof of (c).} Let a seed be created at $(j,\square)$. By
\cite[condition~(ii)]{B-CMP122a}, a component is never a class component at the birth step
of its seed. The seed lies in $\square$; it is therefore part of a live component at level $j$
that contains at least one cell $\mathfrak{c}\subset\square$. If the seed is merged at birth with
an existing component, that component inherits
an older oldest seed, whose degree $p_0$ is larger. This does not affect the argument, because the
event $\mathcal{L}^{(j)}_{\mathfrak{c}}$ places no condition on birth levels. Hence
\begin{equation*}
\{h\ \text{creates a seed at}\ (j,\square)\}\subset\bigcup_{\mathfrak{c}\subset\square}
\mathcal{L}^{(j)}_{\mathfrak{c}}.
\end{equation*}
The sure event is lineage-blind and determined by labels of levels $\le K'$, and $\tilde{\nu}$
has total mass one by Lemma~\ref{8.6:lem-mass-conservation}. Statement (a) with $k'=j$, which is
at most $K'$ because a creation at level $j$ occurs at a step of the run and the run ends at level
$K'$, and with $B$ the sure event therefore gives
$\tilde{\nu}(\mathcal{L}^{(j)}_{\mathfrak{c}})\le C_Le^{-\theta p_0(g_j)}$ for each of the
$n_\square$ cells $\mathfrak{c}$ of $\square$. Summing over these cells gives
\eqref{8.6:eq-block-bound-live-1}. Only the hypotheses of (a) were used. The bound holds at every
level $j$, both coarser and finer than any fixed level of an observable.

\emph{Proof of (d).} We first treat the field clause. The small-field characteristic function
$\chi_{k_0}$ is a factor of every term of the level-$k_0$ density \cite[(2.18)]{B-CMP119}.
Consider the event that $V_{k_0}$ leaves the level-$k_0$ window somewhere on $\mathcal{N}$ while
no label live at level $k_0$ meets $\mathcal{N}$. On this event every term vanishes, so the event
is $\tilde{\nu}$-null. Hence $\tilde{\nu}(\Xi_{\mathcal{N}})$ equals $\tilde{\nu}$ of the label
clause.

Next the label clause. A label live at level $k_0$ that meets $\mathcal{N}$ is part of a live
component at level $k_0$ containing a level-$k_0$ cell that meets $\mathcal{N}$. The label clause
is therefore contained in the union of $\mathcal{L}^{(k_0)}_{\mathfrak{c}}$ over the level-$k_0$
cells $\mathfrak{c}$ meeting $\mathcal{N}$. This gives the first inequality of
\eqref{8.6:eq-block-bound-live-2}. The second
follows from (a) with $k'=k_0\le K'$ and $B$ the sure event, applied to each of the
$n_{k_0}(\mathcal{N})$ cells. No restriction on the birth level enters: the object bounded is the
live component at level
$k_0$ itself, whatever the birth levels of its seeds. This includes small old components kept alive
by later creations.

Finally the healed addendum. A label healed at a step $k\in[k_0-\mathfrak{l},k_0]$ that meets
$\mathcal{N}$ contains a level-$k$ cell $\mathfrak{c}$ meeting $\mathcal{N}$. It therefore lies
in the event $E^{(k)}_{\mathfrak{c}}$ that a first-class island healed at step $k$ contains
$\mathfrak{c}$, the pinned island event of \eqref{8.6:eq-blind-events-a} with no restriction on
the reserve datum. The one-point bound of Theorem~\ref{8.6:thm-block-bound-healed},
with $B$ the sure event, gives
$\tilde{\nu}(E^{(k)}_{\mathfrak{c}})\le K_He^{-\frac32p_0(g_k)}$. Since $k\le k_0\le K$,
\eqref{8.6:eq-coupling-comparison-a} gives $p_0(g_k)\ge p_0(g_{k_0})-1$, hence
\begin{equation*}
e^{-\frac32p_0(g_k)}\le e^{3/2}\,e^{-\frac32p_0(g_{k_0})}.
\end{equation*}
Summing over the at most $n_k(\mathcal{N})$ cells at each step $k$ and over the steps $k$ with
$\max(0,k_0-\mathfrak{l})\le k\le k_0$, and adding \eqref{8.6:eq-block-bound-live-2}, gives
\eqref{8.6:eq-block-bound-live-3}. This summand used Theorem~\ref{8.6:thm-block-bound-healed} and
Lemma~\ref{8.6:lem-coupling-comparison} only, that is hypothesis~($\alpha$), and not
Proposition~\ref{8.6:prop-live-removal}.
\end{proof}

\begin{remark}[Use of the bounds in the two-point witness theorem]
Let $\mathcal{N}$, $k_0$ and $\mathfrak{l}=L\check{R}_{k_0}+c_L$ be as at the end of statement (d)
of Theorem~\ref{8.6:thm-block-bound-live}. In the arithmetic of the two-point witness theorem the
cell counts satisfy
\begin{equation*}
n_k(\mathcal{N})\le\Bigl(\frac{8\varrho L^{k_0-k}}{\check{R}_k}+2\Bigr)^4,
\end{equation*}
with $\varrho$ the radius of the neighbourhoods in that theorem, and $n_{k_0}(\mathcal{N})$ is at
most $(8\varrho+2)^4$ times a factor polylogarithmic in $x$, where $x$ is the inverse squared
coupling entering that theorem's absorption. The position counts are absorbed there as follows.
For the healed member that theorem uses
\begin{equation*}
L^{4\mathfrak{l}}\operatorname{polylog}(x)\le e^{\frac12c_qp_0(g_{k_0})},\qquad c_q=\tfrac32;
\end{equation*}
this is a ceiling on $\gamma_W$, imposed after $L$ (so that the thresholds move with $L$), and it
can be met because the exponent $p_0$ of the degree function exceeds $r_{\mathrm{pr}}$ and
\begin{align*}
L^{4\mathfrak{l}}&=\exp\bigl(4(L\check{R}_{k_0}+c_L)\log L\bigr)\\
&\le\exp\bigl(4L\log L\,(\log x)^{r_{\mathrm{pr}}}+c\bigr)
\end{align*}
for a constant $c$, where $r_{\mathrm{pr}}$ is the power of $\log x$ bounding $\check{R}_{k_0}$ in
that estimate. With these absorptions the two-point witness theorem's bound on its contact event
$\mathcal{B}_i$ applies as stated there, with $c_q=\theta$ for the live member and
$c_q=\tfrac32$ for the healed member.

The two-point witness theorem uses bounds of the form $\tilde{\nu}(A\mid B)$, uniformly over
events $B$ generated by labels whose collars are disjoint from that of the pinned label. For
healed labels the conditional clause of Theorem~\ref{8.6:thm-block-bound-healed} gives them with
constant one. It imposes no restriction on $B$ beyond blindness to the pinned label and its
lineage. For live labels they hold in the same form by statement (b) of
Theorem~\ref{8.6:thm-block-bound-live}, which rests on
Proposition~\ref{8.6:prop-live-removal} at the level of the label, whose proof is outstanding. This
holds for $B$ determined by labels of levels not exceeding the level of the pinned label; it is not
asserted for $B$ involving far labels of later levels.
\end{remark}

\begin{remark}[Why global bounds cannot localise the terminal member]\label{8.6:rem-no-global-route}
We use the notation fixed with the live event \eqref{8.6:eq-live-event-a}: the normalised history measure
$\tilde\nu$, the level $k'$ and the cell $\mathfrak{c}$ indexing that event, the outer data, their live components
and the hulls $\operatorname{hull}(C)$ of these. Let $K'=K-m_{\mathbb{T}}(g)$ be the terminal level, and let
$\mathcal{E}^{(k')}_{\mathfrak{c}}(K')$ be the live event of \eqref{8.6:eq-live-event-a} at that level.
\begin{enumerate}
\item[(a)] Suppose the numerator of $\tilde\nu(\mathcal{E}^{(k')}_{\mathfrak{c}}(K'))$ is bounded from above, and
its denominator from below, by the global bounds on the history sums. The resulting bound on the quotient
carries a factor that grows with the volume. No choice of the depth floor removes this factor.
\item[(b)] Each of the absolute bounds on the history sums listed in the proof below compares with the true total
mass only globally. Turning such a bound into a bound on $\tilde\nu$ with a volume-free constant requires a
normalised local ratio of the kind stated in Proposition~\ref{8.6:prop-live-removal},
\eqref{8.6:eq-live-removal-a}, whose proof is outstanding.
\item[(c)] Bounds by the supremum of the action over the support of $\chi_k$ lose against the small
factor $e^{-\theta c''p_0(g_k)^2}$, with $c''$ the constant of the required gain, for every $\theta$.
\item[(d)] At the level $K'$ and for a live component $C$ of the terminal outer datum, removal domination
\eqref{8.6:eq-live-removal-a} (Proposition~\ref{8.6:prop-live-removal}, whose proof is outstanding) is a
deterministic inequality between two local integrals over $V_{K'}\restriction\operatorname{hull}(C)$. The two
integrals have common
exterior data, and the inequality is required uniformly in the exterior. The volume enters none of its
constants. What has to be paid is uniformity over all exterior configurations admissible off
$\operatorname{hull}(C)$.
\end{enumerate}
\end{remark}

\begin{proof}
We write $\mathcal{L}$ for the linear size, in sites, of the lattice on $\mathbb{T}_m$ on which the history sums
are taken, $\mathcal{L}=n_TL^m$ with $n_T$ free, so that the number of sites is $\mathcal{L}^4=n_T^4L^{4m}$.
The letters $M^{(\theta)}_h$, $E$, $B$ and $t$ below are local to this proof; they are unrelated to the
glossary's $M$, $E_\pm$ and $B^\sharp$.

(a) Let $M^{(\theta)}_h$ be the majorant weights of the histories $h$. Let $\Theta$ be the penalty exponent
carried on the event $\mathcal{E}^{(k')}_{\mathfrak{c}}(K')$ by the seeds of the lineage, in the sense of the
decomposition \eqref{8.6:eq-live-event-a}; this $\Theta$ is unrelated to the stress-tensor density
$\Theta_{\nu\sigma}$. By the absolute upper
bounds on the history sums, the numerator of the terminal member is at most
\begin{equation*}
e^{-\Theta}\sum_h M^{(\theta)}_h\le e^{-\Theta}\,e^{E_+^{(\theta)}\mathcal{L}^4}.
\end{equation*}
The correlation theorem uses the bound of the same type $|\tilde Z|\le e^{E_+^{\sharp}\mathcal{L}^4}$, where
$\tilde Z$ is the sourced partition function of the correlation theorem and $\tilde Z(0)$ its value at zero source.

The denominator is bounded from below in one of two ways. The first is the floor
$\tilde Z(0)\ge e^{-C_{\mathrm{low}}\mathcal{L}^4}$ used in the correlation theorem. The second is the lower bound
$e^{-E_-|T|}$ of \cite[Corollary~3]{B-CMP119}, where $T$ is the lattice of that corollary and the constant
$E_-$, like the companion constant of the upper bound there, is independent of the parameter $\eta$ of that
corollary and of $T$, but depends on $g_k$.

Writing $E_+$ for $E_+^{(\theta)}$ or $E_+^{\sharp}$, whichever of the upper bounds above is used, the quotient is
therefore bounded only by
\begin{equation*}
e^{(E_++C_{\mathrm{low}})\mathcal{L}^4}\,e^{-\Theta}.
\end{equation*}
This depends on the volume through $\mathcal{L}^4$, and $n_T$ is free. On the other side, a lineage surviving
any number of levels, in particular to the terminal level, has
\begin{equation*}
\Theta\ge\theta\,p_0(g)\cdot(\text{number of its seeds}),
\end{equation*}
and this lower bound is unrelated to $\mathcal{L}^4$. Hence no choice of the depth floor cures the volume
dependence.

(b) Consider the absolute bounds on the history sums: the bounds of the family from which the two upper bounds
used in (a) are taken, these two included; the representation and bound of the large-field terms;
the weighted bound described in the next paragraph, together with the estimate that accompanies it; and the
torus bound of the correlation theorem.
Each of them compares with the true total only globally.

The weighted bound shows this most clearly. It carries the factor $(1-\theta_W)$, with $\theta_W$ a parameter of
that bound, on the whole action, and the constant $\mathsf{K}_{\mathrm{cell}}$ once for each cell of the volume.

Now let $E$ be an event determined by a creation at step $j$. Mark the terms of this creation by a factor
$e^t$, with
\begin{equation*}
t\le c'p_0(g_j)
\end{equation*}
for a constant $c'$, and let $Z_t$ be the total mass so modified. A bound $Z_t\le e^{B}Z$, with $B$ a constant
(not the event of Proposition~\ref{8.6:prop-terminal-live}), gives
\begin{equation*}
\tilde\nu(E)\le 2e^{B-t}.
\end{equation*}
The only constant available from the global bounds, however, is $B=(E_+^{\sharp}+C_{\mathrm{low}})\mathcal{L}^4$. To
localise $B$, the dependence of the history sum on $t$ must sit in finitely many local terms. That is, one needs
a normalised local ratio: removal domination for a live component
(Proposition~\ref{8.6:prop-live-removal}, whose proof is outstanding), or its local-normalisation form per step
and per parent. Bounds on the majorant gas supply no such
ratio (compare Remark~\ref{8.6:rem-majorant-gas}).

(c) Let $\tilde\square$ be a cube, the restriction to which is written $\big|_{\tilde\square}$. Let $C_B$ and
$n_\square$ be the constants of the small-field bound on the action over $\tilde\square$ displayed next, let $c''$
be the exponent of the required small factor $e^{-\theta c''p_0(g_k)^2}$, and let $c_A$ be the constant entering
the comparison of exponents below. On $\operatorname{supp}\chi_k$
one has
\begin{equation*}
x_k\,A(U_k(V))\big|_{\tilde\square}\le C_B^2\,n_\square\,p_0(g_k)^2 .
\end{equation*}
Consequently the average $\langle\cdot\rangle$ of the local weight satisfies
\begin{equation*}
\big\langle e^{\theta x_kA|_{\tilde\square}}\big\rangle\le e^{\theta C_B^2n_\square p_0(g_k)^2}.
\end{equation*}
Set against the required factor $e^{-\theta c''p_0(g_k)^2}$, this bound loses by the constant
$C_B^4n_\square/c_A>1$, for every $\theta$. The reason is as follows. The small-field window is $p_0(g_k)$
standard deviations wide, and a bound by the supremum cannot see that typical fields are of size $g_k$. The
constant $\mathsf{K}_{\mathrm{cell}}$ of the weighted bound, and the termwise majorants of the second of the two
upper bounds used in (a), are bounds of this kind.

The two-point witness theorem avoids bounds of this kind through the law of total covariance. The price of
doing so is the local rarity of live labels at the top level, that is, an inequality of the form
\eqref{8.6:eq-live-removal-a} at a single level; in the generality of
Proposition~\ref{8.6:prop-live-removal}, the proof of that inequality is outstanding.

(d) By construction, the event $\mathcal{E}^{(k')}_{\mathfrak{c}}(K')$ sits at $K'=K-m_{\mathbb{T}}(g)$. No
further step is taken there, and the total mass is directly $Z=\int dV_{K'}\sum\rho_{K'}$.

At this level, removal domination \eqref{8.6:eq-live-removal-a} for a live component $C$ is a deterministic
inequality between two local integrals over $V_{K'}\restriction\operatorname{hull}(C)$ with common exterior
data, required uniformly in the exterior. The infrared content at the top scale lies entirely in the exterior
configuration: the torons, that is, the zero-action (flat) configurations of the torus at the top scale, and the
volume $n_T^4L^{4m}$. That configuration is integrated afterwards against a
common non-negative integrand.

The volume therefore enters no constant:
\begin{itemize}
\item the constants $K_H$, $C_\delta$ and $A_d$ entering these local integrals depend on the dimension $d$ only;
\item the constant $S_*$ entering these local integrals is bounded through the lower flow bound, independently of
the volume;
\item the constant $C_L$ entering these local integrals likewise does not depend on the volume;
\item the factor $2e^{\kappa_1}$ in Proposition~\ref{8.6:prop-terminal-live} arises from counting the shapes of
the live component. That proposition is proved as an implication from removal domination, whose proof is
outstanding.
\end{itemize}
What has to be paid is that removal domination is required uniformly over all exterior configurations admissible
off $\operatorname{hull}(C)$.
\end{proof}

\begin{definition}[Creations, blobs, quiet lifetime and the age threshold]\label{8.6:def-creations-blobs}
The following organises the run step by step. It is an equivalent alternative to the organisation by
histories in Definition~\ref{8.6:def-history-measure}. The running couplings $g_j$, $x_j=g_j^{-2}$ and
the degree function $p_0(\cdot)$ along the run are those of Definition~\ref{8.6:not-running-couplings}.

\emph{The step.} For each level $k$ we write
\begin{equation}\label{8.6:eq-creations-blobs-1}
\rho_{k+1}:=\mathbf{R}_{k+1}\,\mathcal{T}_{k+1}\,\rho_k .
\end{equation}
Here $\mathcal{T}_{k+1}$ is Ba{\l}aban's one-step renormalization transformation \cite{B-CMP119}
followed by the step-$(k+1)$ decompositions of unity of every \emph{kind} $\iota$. The kinds form
Ba{\l}aban's finite list: they are the kinds of the decompositions of \cite[(2.17)]{B-CMP119}, of the
restrictions on the fluctuation field of \cite[\S2]{B-CMP109} and \cite[(2.21)]{B-CMP119}, of the
characteristic functions $\chi^{(0)}$, $\chi^{(1)}$ internal to $\mathbf{R}$ of
\cite[(1.22)--(1.23)]{B-CMP122a}, and of the preparatory list of
\cite[\S1]{B-CMP122b}. At a step $j$ each such decomposition is of the form
$1=\chi^{\iota}_{j,\square}+(1-\chi^{\iota}_{j,\square})$, $\square$ running over the cubes of the
step-$j$ partition of kind $\iota$.

\emph{Creations.} Let $j$ be a level, $\iota$ a kind and $h_{j-1}$ a level-$(j-1)$ history. Let
$\square$ be a cube of the step-$j$ partition of kind $\iota$ that lies in the \emph{creation zone} of
$h_{j-1}$ for the kind $\iota$, that is, in the part of the lattice on which, given $h_{j-1}$, the
step-$j$ decomposition of kind $\iota$ is performed.
A \emph{creation of kind $\iota$ at
$(j,\square)$} is the choice of the term $1-\chi^{\iota}_{j,\square}$ in the step-$j$ data
$D_j$. We write $E^{\iota}_{j,\square}$ for the set of level-$j$ histories that contain this choice.

\emph{Conditional weights.} Let $E$ be a set of level-$j$ histories and let $h_{j-1}$ be a level-$(j-1)$
history. The masses $m(h)$ are those of the history tree of Lemma~\ref{8.6:lem-mass-conservation}, and
$h_j$ runs over the children of $h_{j-1}$. We put
\begin{equation}\label{8.6:eq-creations-blobs-2}
P(E\mid h_{j-1}):=\sum_{\substack{h_j\in E\\ h_j\ \text{a child of}\ h_{j-1}}}
\frac{m(h_j)}{m(h_{j-1})} .
\end{equation}

\emph{Sites.} A \emph{site} is a pair $(j,\square)$ taken together with a kind. Creations of different
kinds at the same $(j,\square)$ are different sites.

\emph{The region rules.} The large-field regions of Ba{\l}aban's procedure obey the following region
rules (a)--(d), with the numbers \eqref{8.6:eq-creations-blobs-a}. They
are used here as stated in \cite[\S1]{B-CMP122a} and \cite[\S1]{B-CMP122b}.
\begin{itemize}
\item[(a)] (The blob bound.) Let $n$ be a later level at which a creation at $(j,\square)$ is present (is
still part of the large-field region at level $n$). The \emph{blob} of the creation at level $n$ is the
cube $S^{n-j}(\square)$ that \cite[\S1]{B-CMP122b} attaches to $\square$. As stated there,
$S^{n-j}(Z)=\bigcup_{\square\subset Z}S^{n-j}(\square)$ for a large-field region $Z$; the cube of
side $MR_n$ in the centre of $S^{n-j}(\square)$, where $R_n$ is the radius that
\cite[\S1]{B-CMP122b} fixes at level $n$, contains $\square$; and $S^{n-j}(Z)$ is contained in the
union of enlargements, fixed there, of the cubes $\square\subset Z$. The blob is a cube containing the
image of $\square$ on the level-$n$ lattice, and its side obeys the bound in the first line of
\eqref{8.6:eq-creations-blobs-a}, with $\check R_n$ as described after that display. A connected
component of the large-field region present at level $n$ is the union of the
blobs of its creations, together with the cubes of the layers of its creation regions, these layers
being those formed in \cite[\S1]{B-CMP122a}.
\item[(b)] (The lifetime bound.) Let $Z$ be a component that is quiet from level $n$ on. Then, by
\cite[(1.80)--(1.81)]{B-CMP122b}, $Z$ is healed within $\Delta(Z)$ further steps, where
$\Delta(Z)$ obeys the first bound in the second line of
\eqref{8.6:eq-creations-blobs-a}.
Here $d'_n(Z)$ is the quantity attached to $Z$ at level $n$ in \cite[(1.80)--(1.81)]{B-CMP122b}
through which that bound depends on $Z$; it satisfies the second
bound in the second line of \eqref{8.6:eq-creations-blobs-a}.
\item[(c)] (Merging; cancelling.) Components merge and creations cancel according to the rules of
\cite[(1.84)]{B-CMP122b} and \cite[\S1]{B-CMP122a}.
\item[(d)] (Monotonicity.) Being live is an up-set in the configuration of creations: if a
component is live under a set of creations, it is live under every larger set of creations.
\end{itemize}

\emph{Numbers.} We put
\begin{equation}\label{8.6:eq-creations-blobs-a}
\begin{gathered}
\operatorname{side}\bigl(\text{blob at level } n\bigr)\le C_G M\check R_n,\\
\Delta(Z)\le 2\log_L d'_n(Z)+\check R_n+2,\qquad
d'_n(Z)\le C_G\cdot\#\{\text{blobs of }Z\},\\
T_0:=2+2\log_L(2C_G),\qquad T_1(x):=\check R(x)+T_0,\\
q_j:=C_q\,e^{-c_q\,p_0(g_j)} .
\end{gathered}
\end{equation}
Here $\check R(\cdot)$ is a radius depending on the running inverse coupling, and $\check R_n=\check R(x_n)$
is its value at level
$n$. The constant $C_G$ is a constant of the cube construction of \cite[\S1]{B-CMP122b}; $C_q$ and
$c_q$ are constants that do not depend on the level, so that $q_j$ depends on $j$ only through $g_j$.

\emph{Age.} The \emph{age} of a creation at $(j,\square)$, read at an observation level $k'$, is $k'-j$.
The creation is \emph{young} at level $k'$ if $k'-j\le T_1(x_{k'})$, and \emph{old} otherwise.
\end{definition}

\begin{proposition}[Conditional rarity of simultaneous creations. Proof outstanding]
\label{8.6:prop-creation-rarity}
There are constants $C_{\mathrm{cr}}\ge 1$ and $C_q,c_q>0$, depending only on $L$, $M$, $\mathfrak{p}$
and $N$, with the following property. Let $q_j$ be as in \eqref{8.6:eq-creations-blobs-a}, formed
with these constants $C_q$, $c_q$. For every $K$, every $m$, every level $j$ with $1\le j\le K$,
every parent history $h_{j-1}$ in the larger of the two spaces of histories used in this chapter
(and hence, a fortiori, in the smaller one), every integer $n\ge1$, and every family of $n$ pairwise distinct
cubes $\square_1,\dots,\square_n$ of the partitions used at step $j$, carrying kinds
$\mathfrak{k}_1,\dots,\mathfrak{k}_n$, such that each $\square_i$ lies in the creation zone of
kind $\mathfrak{k}_i$ determined by $h_{j-1}$, we have
\begin{equation}\label{8.6:eq-creation-rarity-a}
\sum_{\substack{h_j \text{ child of } h_{j-1}\\
h_j\in\bigcap_{i\le n}E^{\mathfrak{k}_i}_{j,\square_i}}} m(h_j)
\;\le\;\bigl(C_{\mathrm{cr}}\,q_j\bigr)^{n}\, m(h_{j-1}).
\end{equation}
Here $E^{\mathfrak{k}}_{j,\square}$ is the event that step $j$ creates a large-field region of
kind $\mathfrak{k}$ at the cube $\square$, in the sense of
Definition~\ref{8.6:def-creations-blobs}, and $m(h)$ is the mass of the history $h$, that is,
the total integral of its term, as in Lemma~\ref{8.6:lem-mass-conservation}. The case $n=1$ of
\eqref{8.6:eq-creation-rarity-a} is called the one-cube case.
\end{proposition}

Proof outstanding.

\medskip
\noindent\textit{Route.} (R1) \emph{The quotient of the large-field operation as an estimate.}
Fix the step $j$, the parent $h_{j-1}$ and the sites $(\mathfrak{k}_i,\square_i)$, $i\le n$.
Let $\chi^{\mathfrak{k}}_{\square}$ denote the small-field characteristic function of kind
$\mathfrak{k}$ on $\square$, so that $1-\chi^{\mathfrak{k}}_{\square}$ is the indicator of the
corresponding large-field condition. Form the quotient that defines Ba{\l}aban's operation
$\mathbf{R}$ in \cite[(0.3)]{B-CMP122a}. Its numerator is the integrand of the children of
$h_{j-1}$ multiplied by
\[
\prod_{i\le n}\bigl(1-\chi^{\mathfrak{k}_i}_{\square_i}\bigr),
\]
with the large-field operations acting there. Its denominator is, as described in
\cite{B-CMP122a}, the integral of parts of the densities localised in neighbourhoods of the cubes
$\square_i$ and determined by the small-field effective actions alone. This quotient is to be
used as an estimate at the current step only, and not as a change of representation: nothing is
carried to later steps. The reasons for which \cite{B-CMP122a} restricts $\mathbf{R}$ to
regions satisfying its conditions (i) and (ii) concern the representation carried forward (the
small factor does not control the further steps) and the rectangular geometry of the regions
involved; neither arises for an estimate confined to one step. For the numerator the bounds of
\cite{B-CMP122b} hold for general large-field regions, per cube in the form
\cite[(1.85)]{B-CMP122b}. For the denominator a proof must establish positivity, non-emptiness of
the domain of integration, and a lower bound for the local small-field integral next to an
arbitrary large-field region: a young region, a large one, an irregular one, and one adjacent to
a live region. In \cite[\S\S1--3, (1.100)--(1.102)]{B-CMP122a} this lower bound is established
for regions satisfying conditions (i) and (ii). The small-field window is $p_0(g_j)$ standard
deviations wide. The mismatch factor of \cite{B-CMP122b} is $\exp\bigl(O(1)\bigr)$ raised to the
number of cells of the large-field region of step $j$ that lie in the small-field region of
step $j$; it costs $\exp(O(1))$ per such cell, and it is affordable per cube against the gain of
order $p_0(g_j)^2$ because the exponent of the degree function $p_0(\cdot)$ satisfies a one-sided
lower bound. The
logarithm of the quotient is then to be expanded in the polymer format of the representation and
bound of the large-field terms, under the convergence condition of part~(d) of
\eqref{8.6:eq-hard-core-gas-a}. Each cube $\square_i$ is given an auxiliary complex parameter, its
mark, entering the activities of this expansion; there is one mark for each of the $n$ cubes. The
conditional probability of $\bigcap_{i\le n}E^{\mathfrak{k}_i}_{j,\square_i}$
given $h_{j-1}$ is then read off from the derivatives with respect to the marks at $0$, by
part~(d) of Lemma~\ref{8.6:lem-hard-core-gas}.

(R2) \emph{The local exponential moment.} Equivalently, for creations of the basic kind
(Definition~\ref{8.6:def-creations-blobs}),
\eqref{8.6:eq-creation-rarity-a} follows from the local exponential moment
\eqref{8.6:eq-exponential-moment-a}. That moment bound is to be proved for arbitrary pasts by the
fluctuation expansion of \cite[\S\S3--4]{B-CMP109}. In that expansion the local tilt is absorbed
into the covariance and the coupling on the neighbourhoods $\tilde\square_i$ of the cubes
$\square_i$. This is a bounded relative change on $n\cdot n_{\square}$ cells, where $n_{\square}$
is the number of cells in one neighbourhood $\tilde\square_i$.

(R3) \emph{Localisation in the volume.} On either route, condition on the block field $V_{j'}$,
where $j'=j+\lceil\frac14\log_L x_j\rceil+C$ with $C$ a constant. This
localises the estimate in the volume.

\medskip
The inequality \eqref{8.6:eq-creation-rarity-a} concerns one step. It is annealed in the fields:
the field variables of the children are integrated. It is conditional only on the large-field
labels of the past, that is, on $h_{j-1}$, and on no field configuration. The cubes
$\square_1,\dots,\square_n$ may be adjacent. The small factor per cube that adjacency requires is
of the kind supplied by the extraction \cite[(1.77)--(1.78)]{B-CMP122b}, which is performed
before the regions are coarsened. It is not of the kind supplied by the per-region bound
\cite[(1.85)]{B-CMP122b}, which yields a gain of order $p_0(g_j)^2$ per unit of linear extent.
Summed over all parents $h_{j-1}$ whose creation zone of kind $\mathfrak{k}_1$ contains
$\square_1$, the one-cube case bounds the total mass of the creation event
at $(\mathfrak{k}_1,\square_1)$ by $C_{\mathrm{cr}}q_j$ times the total mass; by
Lemma~\ref{8.6:lem-mass-conservation} the latter is the same at every level.

The variant that conditions also on the field cannot be expected to hold uniformly; the
obstruction is the following, which is an indication and is not proved here. In that variant the
probability of a creation at $\square$ at step $k+1$ is conditioned on $h_k$ and on the field
off $\tilde\square$, and is required to be at most $q$ uniformly. Suppose the
exterior field lies in the small-field window but near its edge. The lattice harmonic
(minimiser) extension into $\tilde\square$ then satisfies the maximum principle, and this forces
an interior curvature of order $\varepsilon_k$. Here $\varepsilon_k$ is the small-field threshold
at level $k$. This curvature is much larger than the creation threshold of the next step,
$\varepsilon_{k+1}/(C_B C_{\mathrm{av}} L^2)$, where $C_B$ and $C_{\mathrm{av}}$ are the
constants of the curvature bound for the block configuration and of the averaging estimate. The
creation is therefore likely. For this reason \eqref{8.6:eq-creation-rarity-a} integrates the
field and conditions on labels only.

This mechanism does not conflict with Proposition~\ref{8.6:prop-live-removal}, whose proof is
outstanding. That proposition conditions on the field off the hull of a live component at the
same level, where no threshold of a next step enters. The mechanism above, by contrast, concerns
the field on the boundary of $\tilde\square$ at plaquette range, compared with a threshold of the
next step.

Removal domination for a live component at level $j$ (Proposition~\ref{8.6:prop-live-removal},
whose proof is outstanding) can be applied to a newborn component. Its lineage consists of the
component itself. The outer datum with the component removed is the datum of the parent without
that creation, and the other creations are events blind to the component. Integrated over the
exterior field, this application yields the one-cube case at level $j$. Applied $n$ times, it
yields \eqref{8.6:eq-creation-rarity-a}. As families of statements, however, neither implies the
other. The one-point bounds of this chapter use removal domination only at the one terminal level,
and at the level itself for the conditional two-cell decay. The route through the present
proposition uses \eqref{8.6:eq-creation-rarity-a} at every level up to the level of observation.
Both rest on the same local ratio estimate. This is a lower bound for the local small-field
integral next to the region, compared with the integral carrying the large-field condition, as
described in (R1).

\begin{lemma}[Peeling creations level by level]\label{8.6:lem-peeling}
In the notation of Definition~\ref{8.6:def-creations-blobs} and
Lemma~\ref{8.6:lem-mass-conservation}, write $\sigma(h_{j-1})$ for the $\sigma$-algebra of events
determined by the history up to level $j-1$, $E_s$ for the creation event of a site $s$ (a site in the
sense of Definition~\ref{8.6:def-creations-blobs}, with its level), and $\tilde\nu$ for the annealed
measure on histories. Assume that the conditional rarity of simultaneous creations,
Proposition~\ref{8.6:prop-creation-rarity} (whose proof is outstanding), holds for every $n\ge1$ in
the form \eqref{8.6:eq-creation-rarity-a}, with the constant $C_{\mathrm{cr}}\ge1$ and the numbers
$q_j$ given there. Let $S$ be a finite family of pairwise distinct sites, let $j_1<\dots<j_r$ be the
levels that occur among them, and let $n_t$ be the number of sites of $S$ at level $j_t$. Let $A$ be
an event in $\sigma(h_{j_1-1})$. Then
\begin{equation}\label{8.6:eq-peeling-1}
\tilde\nu\Bigl(A\cap\bigcap_{s\in S}E_s\Bigr)
\le\tilde\nu(A)\prod_{t\le r}\bigl(C_{\mathrm{cr}}\,q_{j_t}\bigr)^{n_t}.
\end{equation}
\end{lemma}

\begin{proof}
Let $m(h)=\int\tau_h\,dV$ and $Z=\int\rho_k$, the same for every $k$, be as in the conservation of
mass along the history tree (Lemma~\ref{8.6:lem-mass-conservation}). For $1\le t\le r$ let $S_t$ be
the set of sites of $S$ at level $j_t$. Put $A_0:=A$ and, for $1\le t\le r$,
\begin{equation*}
A_t:=A\cap\bigcap_{u\le t}\ \bigcap_{s\in S_u}E_s .
\end{equation*}
We prove by induction on $t$ that
\begin{equation*}
\tilde\nu(A_t)\le\tilde\nu(A)\prod_{u\le t}\bigl(C_{\mathrm{cr}}q_{j_u}\bigr)^{n_u}.
\end{equation*}
For $t=0$ this is an equality. For $t=r$ it is \eqref{8.6:eq-peeling-1}.

Let $1\le t\le r$ and assume the bound for $t-1$. The event $A_{t-1}$ is
$\sigma(h_{j_t-1})$-measurable: $A\in\sigma(h_{j_1-1})\subset\sigma(h_{j_t-1})$ because
$j_1\le j_t$, and for $u<t$ each $E_s$ with $s\in S_u$ is determined by the labels of level
$j_u\le j_t-1$. The event $A_t=A_{t-1}\cap\bigcap_{s\in S_t}E_s$ is determined by the labels of
levels $\le j_t$. Lemma~\ref{8.6:lem-mass-conservation}, in the form
\eqref{8.6:eq-mass-conservation-a}, gives $\tilde\nu(A_t)$ as $Z^{-1}$ times the sum of $m(h_{j_t})$
over the level-$j_t$ histories in $A_t$. We group these histories according to their parents
$h_{j_t-1}$, which lie in $A_{t-1}$, and obtain
\begin{equation}\label{8.6:eq-peeling-2}
\tilde\nu(A_t)=\sum_{h_{j_t-1}\in A_{t-1}}Z^{-1}
\sum_{\substack{h_{j_t}\ \text{child of}\ h_{j_t-1}\\ h_{j_t}\in\bigcap_{s\in S_t}E_s}}m(h_{j_t}).
\end{equation}
Whenever $m(h_{j_t-1})>0$, the inner sum is
$m(h_{j_t-1})\,P\bigl(\bigcap_{s\in S_t}E_s\bigm| h_{j_t-1}\bigr)$. Here $P(B\mid h_{j_t-1})$
denotes the total mass of the children of $h_{j_t-1}$ lying in $B$, divided by $m(h_{j_t-1})$, and
$Z^{-1}m(h_{j_t-1})=\tilde\nu(h_{j_t-1})$ by the same lemma.

For each parent $h_{j_t-1}$ we apply the hypothesis \eqref{8.6:eq-creation-rarity-a} with $n=n_t$
to the $n_t$ distinct sites of $S_t$. It bounds the inner sum in \eqref{8.6:eq-peeling-2} by
$(C_{\mathrm{cr}}q_{j_t})^{n_t}m(h_{j_t-1})$. Summing over $h_{j_t-1}\in A_{t-1}$ and applying
Lemma~\ref{8.6:lem-mass-conservation} once more at level $j_t-1$, we get
\begin{equation*}
\tilde\nu(A_t)\le(C_{\mathrm{cr}}q_{j_t})^{n_t}\,\tilde\nu(A_{t-1})
\le\tilde\nu(A)\prod_{u\le t}(C_{\mathrm{cr}}q_{j_u})^{n_u},
\end{equation*}
where the last step is the induction hypothesis.

Thus the product over levels comes from conditioning on the past, and the product within a level
comes from the hypothesis \eqref{8.6:eq-creation-rarity-a} applied to the $n_t$ distinct sites of one
level; nothing is conditioned on the future.
\end{proof}

\begin{lemma}[Reaching cells through chains of creation blobs]\label{8.6:lem-blob-chains}
Let $Z$ be a component present at level $n$. Components present at a level, creations, blobs,
$n$-cells, $n$-images, the $n$-distance $\operatorname{dist}_n$ (the distance in the maximum norm
over the four coordinates, in units of the spacing of the level-$n$ lattice), the step-$j$
partitions and the constants $C_G$, $\check{R}_n$, $\check{R}$ are those of
Definition~\ref{8.6:def-creations-blobs}. Cubes are closed, and two blobs touch when they have a
common point.
\begin{itemize}
\item[(a)] If $Z$ meets the $n$-cell $e$, then some creation $(j'',\square'')$ of $Z$ has $e$ within
$n$-distance $C_G M\check{R}_n$ of the $n$-image of $\square''$. Given a point $x$ of $e$ that lies in
the blob of $(j'',\square'')$ at level $n$, the cube $\square''$ is one of at most
\begin{equation*}
\bigl(2C_G M\check{R}_n L^{n-j''}+2\bigr)^4\le (4C_G M\check{R})^4 L^{4(n-j'')}
\end{equation*}
cubes of the step-$j''$ partition; the count is made for each such point $x$ of $e$.
\item[(b)] If $Z$ meets the $n$-cells $e$ and $e'$, and $\operatorname{dist}_n(e,e')=D$, then $Z$
contains $N\ge\lceil D/(2C_G M\check{R}_n)\rceil$ distinct creations $(j_i,\square_i)$,
$1\le i\le N$, whose blobs at level $n$ form a chain from $e$ to $e'$: the first blob meets $e$, the
last meets $e'$, and consecutive blobs touch. Consequently, given $\square_i$, the cube
$\square_{i+1}$ is one of at most $(4C_G M\check{R})^4L^{4(n-j_{i+1})}$ cubes of the step-$j_{i+1}$
partition.
\end{itemize}
\end{lemma}

\begin{proof}
Everything follows from the blob rule of Definition~\ref{8.6:def-creations-blobs}
(see \eqref{8.6:eq-creations-blobs-a}), taken from \cite[\S1]{B-CMP122b}. We use it in the following
form. For a creation at $(j,\square)$ and a level $n$ at which it is present, the blob is a cube of
side at most $C_G M\check{R}_n$ that contains the $n$-image of $\square$. It is the cube that
\cite[\S1]{B-CMP122b} attaches to $\square$ and to the number $n-j$ of steps, so it is determined by
$(j,\square)$ and $n$. A component present at level $n$ is the union of the blobs of its creations;
the cubes of the layers of its creation regions, which the blob rule also lists, lie inside these
blobs (Definition~\ref{8.6:def-creations-blobs}). Besides the blob rule we use two facts about
$n$-images. By Definition~\ref{1.5:def-block-average}, one spacing of the level-$j$ lattice is
$L^{-(n-j)}$ spacings of the level-$n$ lattice; so, by Definition~\ref{8.6:def-creations-blobs}, the
$n$-images of the cubes of the step-$j$ partition are cubes of side at least $L^{-(n-j)}$, hence of
volume at least $L^{-4(n-j)}$, with pairwise disjoint interiors. In the maximum norm, two points of
a cube of side $b$ are at distance at most $b$, and a cube of side $b$ containing a point $y$ lies
in the cube of side $2b$ centred at $y$.

(a) Since $Z$ is the union of the blobs of its creations and meets $e$, some creation
$(j'',\square'')$ of $Z$ has a blob $B''$ that meets $e$. Pick $x\in e\cap B''$. The $n$-image of
$\square''$ also lies in $B''$, and the side of $B''$ is at most $C_G M\check{R}_n$. Hence $x$, and with
it $e$, lies within $n$-distance $C_G M\check{R}_n$ of that image.

For the count, let $x$ be a point of $e$ lying in $B''$. Since $B''$ contains $x$ and has side at
most $C_G M\check{R}_n$, it lies in the cube $Q(x)$ of side $2C_G M\check{R}_n$ centred at $x$. So does
the $n$-image of $\square''$. The $n$-images of the cubes of the step-$j''$ partition that lie in
$Q(x)$ have disjoint interiors and volume at least $L^{-4(n-j'')}$ each, while $Q(x)$ has volume
$(2C_G M\check{R}_n)^4$. Hence $\square''$ is one of at most
$(2C_G M\check{R}_nL^{n-j''})^4\le(2C_G M\check{R}_nL^{n-j''}+2)^4$ cubes of that partition.

For the second inequality put $a=2C_G M\check{R}_nL^{n-j''}$. The blob $B''$ contains the $n$-image of
$\square''$, of side at least $L^{-(n-j'')}$, and has side at most $C_G M\check{R}_n$. Hence
$C_G M\check{R}_nL^{n-j''}\ge1$, that is $a\ge2$, and so $a+2\le 2a$. Together with the bound
$\check{R}_n\le\check{R}$ of Definition~\ref{8.6:def-creations-blobs}, this gives
$(a+2)^4\le(2a)^4\le(4C_G M\check{R})^4L^{4(n-j'')}$.

(b) Form the graph whose vertices are the creations of $Z$ present at level $n$, with an edge between
two creations whenever their blobs touch. Suppose this graph had two nonempty vertex classes with no
edge between them. Then $Z$ would be the union of two subsets, each a union of blobs, with no blob of
one touching a blob of the other, that is, the union of two disjoint nonempty closed sets. This
contradicts the fact that $Z$ is a component, so the graph is connected.

By part (a), some creation has a blob meeting $e$, and some creation has a blob meeting $e'$.
Connectedness gives a path in the graph from the first to the second. Choose such a path
$(j_1,\square_1),\dots,(j_N,\square_N)$ with $N$ as small as possible. Its creations are then
distinct: a repetition could be cut out, giving a shorter path. Write $B_i$ for the blob of
$(j_i,\square_i)$ at level $n$.

Pick points $x_0\in e\cap B_1$ and $x_N\in e'\cap B_N$, and for $1\le i<N$ a common point
$y_i\in B_i\cap B_{i+1}$. Put $y_0=x_0$ and $y_N=x_N$. Consecutive points $y_{i-1},y_i$ both lie in
$B_i$, a cube of side at most $C_G M\check{R}_n$. The triangle inequality then gives
\begin{equation*}
D\le \operatorname{dist}_n(x_0,x_N)
\le\sum_{i=1}^{N}\operatorname{dist}_n(y_{i-1},y_i)
\le N\,C_G M\check{R}_n .
\end{equation*}
Hence $N\ge D/(C_G M\check{R}_n)\ge D/(2C_G M\check{R}_n)$. Since $N$ is an integer,
$N\ge\lceil D/(2C_G M\check{R}_n)\rceil$.

Finally, fix $i<N$ and let $c_i$ be the centre of $B_i$. By the blob rule, $B_i$ and hence $c_i$ are
determined by $(j_i,\square_i)$ and $n$. The common point $y_i$ of $B_i$ and $B_{i+1}$ lies within
$n$-distance $\tfrac12 C_G M\check{R}_n$ of $c_i$. Every point of $B_{i+1}$ lies within $n$-distance
$C_G M\check{R}_n$ of $y_i$. So $B_{i+1}$, and with it the $n$-image of $\square_{i+1}$, lies in the
cube of side $3C_G M\check{R}_n$ centred at $c_i$. As in part (a), the $n$-images of the cubes of the
step-$j_{i+1}$ partition lying in this cube number at most
$(3C_G M\check{R}_nL^{n-j_{i+1}})^4$. By $\check{R}_n\le\check{R}$ this is at most
$(4C_G M\check{R})^4L^{4(n-j_{i+1})}$, which is the count in (b).
\end{proof}

\begin{theorem}[Young labels and two-cell decay from the creation bound]
\label{8.6:thm-young-two-cell}
Let $K$, $m$ and an observation level $k'\le K$ be given. Let $\tilde{\nu}$ be the normalised measure
on histories of Lemma~\ref{8.6:lem-mass-conservation}.
Creations, sites, the events $E^{\iota}_{j,\square}$ (a creation of kind $\iota$ at the cube
$\square$ of the step-$j$ partition), blobs, the blob constant $C_G$, the factors
$\check{R}_k=\check{R}(x_k)$, the age threshold $T_1(\cdot)$, the numbers
$q_j=C_q e^{-c_q p_0(g_j)}$ and ages are those of Definition~\ref{8.6:def-creations-blobs}. The constants
$C_{\mathrm{cr}}\ge1$, $C_q$, $c_q>0$ are those of Proposition~\ref{8.6:prop-creation-rarity}
(whose proof is outstanding). We put
$E_{j,\square}:=\bigcup_{\iota}E^{\iota}_{j,\square}$, the union over the finitely many kinds. We write
$C_\kappa$ for a number not smaller than $C_{\mathrm{cr}}$ times the number of kinds; it is the
constant of the creation bound summed over the kinds in (a) below. We put
$\kappa_c(x):=\tfrac12 c_q p_0(x^{-1/2})$. Assume:
\begin{enumerate}
\item[(1)] the conditional rarity of simultaneous creations, Proposition~\ref{8.6:prop-creation-rarity}
(whose proof is outstanding), that is, the bound \eqref{8.6:eq-creation-rarity-a} for every number of
pairwise distinct cubes;
\item[(2)] the blob rule of Definition~\ref{8.6:def-creations-blobs}, as displayed in
\eqref{8.6:eq-creations-blobs-a};
\item[(3)] the positivity dichotomy of the large-field operations;
\item[(4)] conservation of mass along the history tree, Lemma~\ref{8.6:lem-mass-conservation};
\item[(5)] the reference bounds \eqref{8.6:eq-running-couplings-a} of
Notation~\ref{8.6:not-running-couplings};
\item[(6)] $g\le\gamma_{\mathrm{LF}}$, where $\gamma_{\mathrm{LF}}>0$ is a coupling ceiling such that
for every $y\ge\gamma_{\mathrm{LF}}^{-2}-4c_2$ the following hold:
\begin{equation}\label{8.6:eq-young-two-cell-1}
C_\kappa\bigl(4C_G M\check{R}(y)\bigr)^{8}L^{4T_1(y)}\bigl(T_1(y)+1\bigr)e^{A_0}C_{\mathrm{cr}}
\le e^{\frac14 c_q p_0(y^{-1/2})},
\end{equation}
and $y$ exceeds the thresholds beyond which the functions of $y$ entering
\eqref{8.6:eq-young-two-cell-1} and the next display are monotone, and
\begin{equation}\label{8.6:eq-young-two-cell-2}
\sum_{\nu\ge0}e^{-\frac12 c_q p_0((y+\lambda\nu)^{-1/2})}\le 2 ,
\end{equation}
with $\lambda$ the flow constant of \eqref{8.6:eq-running-couplings-a}, and
\begin{equation}\label{8.6:eq-young-two-cell-6}
C_q^{1/4}\le\bigl(4C_GM\check{R}(y)\bigr)^4C_{\mathrm{cr}},\qquad
C_q^{3/4}e^{-\frac14 c_qp_0(y^{-1/2})}\le\tfrac12 ;
\end{equation}
these two are one-sided conditions on $y$, like \eqref{8.6:eq-young-two-cell-1}, and are used in
the proofs of (c) and (d).
\end{enumerate}
Under these assumptions the following hold; the bounds of (c) and (d) are collected in
\eqref{8.6:eq-young-two-cell-a} at the end of the statement.
\begin{enumerate}
\item[(a)] (Creation form.) For every site $(j,\square,\iota)$,
$\tilde{\nu}(E^{\iota}_{j,\square})\le C_{\mathrm{cr}}q_j$; summed over the kinds,
$\tilde{\nu}(E_{j,\square})\le C_\kappa q_j$.
\item[(b)] (Joint form.) Let $S$ be a finite family of pairwise distinct sites
$s=(j(s),\square_s,\iota_s)$, and write $E_s:=E^{\iota_s}_{j(s),\square_s}$. Then
\begin{equation*}
\tilde{\nu}\Bigl(\bigcap_{s\in S}E_s\Bigr)\le\prod_{s\in S}C_{\mathrm{cr}}q_{j(s)} .
\end{equation*}
\item[(c)] (Young label form.) Let $1\le j\le k''\le k'$, let $e$ be a cell of level $k''$, and put
$T_1:=T_1(x_{k''})$. Let $Y_{k''}(e;j)$ be the event that $e$ meets a component present at level
$k''$ all of whose creations have levels $j''$ with $\max(j,k''-T_1)\le j''\le k''$. Then the first
bound in \eqref{8.6:eq-young-two-cell-a} holds.
\item[(d)] (Two-cell form for all-young components.) Let $e,e'$ be cells of level $k'$ with
$\operatorname{dist}_{k'}(e,e')=D$. Put $T_1:=T_1(x_{k'})$ and
$J_0:=\max\{1,\lceil D/(2C_G M\check{R}_{k'})\rceil\}$. Let $Y_{k'}(e,e')$ be the event that some
component
present at level $k'$, all of whose creations are young at level $k'$, meets both $e$ and $e'$. No
assumption is made on how the component arose from its creations; in particular it may have arisen
by merging. Put
\begin{equation}\label{8.6:eq-young-two-cell-3}
\eta:=\sum_{\max\{1,k'-T_1\}\le j\le k'}\bigl(4C_G M\check{R}_{k'}\bigr)^4L^{4(k'-j)}C_\kappa q_j ,
\qquad
C':=2\cdot 4^4C_G^4\,(T_1+1)L^{4T_1}C_\kappa e^{A_0};
\end{equation}
the number $C'$ depends on $k'$ only through $T_1=T_1(x_{k'})$.
Then $\eta\le q_{k'}^{3/4}$, and the second bound in \eqref{8.6:eq-young-two-cell-a} holds.
\end{enumerate}
The bounds of (c) and (d) are
\begin{equation}\label{8.6:eq-young-two-cell-a}
\begin{gathered}
\tilde{\nu}\bigl(Y_{k''}(e;j)\bigr)\le q_{k''}^{3/4}\le q_{k'}^{3/4}e^{A_0},\\
\tilde{\nu}\bigl(Y_{k'}(e,e')\bigr)\le\sum_{J\ge J_0}\eta^J\le 2\eta^{J_0}
\le C'\bigl(M\check{R}_{k'}\bigr)^4q_{k'}\,e^{-\kappa_c(x_{k'})(J_0-1)} .
\end{gathered}
\end{equation}
\end{theorem}

Concerning the existence of a ceiling $\gamma_{\mathrm{LF}}$ as in hypothesis (6): the conditions
\eqref{8.6:eq-young-two-cell-1} and \eqref{8.6:eq-young-two-cell-2} have the form
$\Phi(y)\le c$ for a constant $c$, with $\Phi$ decreasing to $0$ as $y\to\infty$, since the exponent
$p_0$ of the degree function $p_0(\cdot)$ exceeds the exponent $r$ in the bound
$\check{R}(y)\le L(\log y)^{r}$ of Notation~\ref{8.6:not-running-couplings}.
The second condition of \eqref{8.6:eq-young-two-cell-6} is of the same form, since
$p_0(y^{-1/2})\to\infty$ as $y\to\infty$. The first condition of \eqref{8.6:eq-young-two-cell-6} has
the form $c\le\Psi(y)$ with $\Psi(y)=(4C_GM\check{R}(y))^4C_{\mathrm{cr}}$. Since
\eqref{8.6:eq-creation-rarity-a} (Proposition~\ref{8.6:prop-creation-rarity}, whose proof is
outstanding) remains true when its constant is replaced by a larger number, hypothesis (1) holds with
$C_{\mathrm{cr}}$ replaced by $\max\{C_{\mathrm{cr}},C_q^{1/4}\}$, and with $C_\kappa$ enlarged
accordingly; with that choice the first condition of \eqref{8.6:eq-young-two-cell-6} holds for
every $y$ with $4C_GM\check{R}(y)\ge1$.

\begin{proof}
We first record four facts used throughout.
\begin{itemize}
\item By Lemma~\ref{8.6:lem-mass-conservation}, $\tilde{\nu}$ is normalised by $\int\rho_k$, which is
the same at every level $k$. Hence the set of all histories has $\tilde{\nu}$-measure $1$.
\item The weights defining $\tilde{\nu}$ are nonnegative by assumption~(3). So $\tilde{\nu}$ is monotone
and subadditive, and the union bounds below are legitimate.
\item The reference bounds \eqref{8.6:eq-running-couplings-a} give $q_j\le e^{A_0}q_k$ whenever
$j\le k$.
\item By Notation~\ref{8.6:not-running-couplings}, $x_k\ge g^{-2}$ for every $k$, and
$g^{-2}\ge\gamma_{\mathrm{LF}}^{-2}\ge\gamma_{\mathrm{LF}}^{-2}-4c_2$ because $g\le\gamma_{\mathrm{LF}}$.
Hence \eqref{8.6:eq-young-two-cell-1} and \eqref{8.6:eq-young-two-cell-6} may be read at $y=x_k$ for
every level $k$, where $\check{R}(x_k)=\check{R}_k$ and $p_0(x_k^{-1/2})=p_0(g_k)$.
\end{itemize}
Condition \eqref{8.6:eq-young-two-cell-2} is part of the ceiling $\gamma_{\mathrm{LF}}$; the proofs
of (a)--(d) below use only \eqref{8.6:eq-young-two-cell-1} and \eqref{8.6:eq-young-two-cell-6}.

\emph{Proof of (a).} The set of all histories is determined by the history up to level $j-1$, that
is, it belongs to the class $\sigma(h_{j-1})$ of Lemma~\ref{8.6:lem-peeling}. Apply
Lemma~\ref{8.6:lem-peeling} to $A$ equal to this set and to the family consisting of the single site
$(j,\square,\iota)$. That lemma is proved
from Proposition~\ref{8.6:prop-creation-rarity}, whose proof is outstanding, and from
Lemma~\ref{8.6:lem-mass-conservation}. Since $\tilde{\nu}(A)=1$, it gives
$\tilde{\nu}(E^{\iota}_{j,\square})\le C_{\mathrm{cr}}q_j$. Summing over the kinds
gives $\tilde{\nu}(E_{j,\square})\le C_\kappa q_j$.

\emph{Proof of (b).} Group $S$ by levels $j_1<\dots<j_r$, with $n_t$ sites at level $j_t$. Apply
Lemma~\ref{8.6:lem-peeling} with $A$ the set of all histories, which has $\tilde{\nu}$-measure $1$. It
gives
\begin{equation*}
\tilde{\nu}\Bigl(\bigcap_{s\in S}E_s\Bigr)\le\prod_{t\le r}\bigl(C_{\mathrm{cr}}q_{j_t}\bigr)^{n_t},
\end{equation*}
which is the product over $s\in S$ of $C_{\mathrm{cr}}q_{j(s)}$.

\emph{Proof of (c).} Suppose $Y_{k''}(e;j)$ occurs, and let $Z$ be such a component present at level
$k''$ meeting $e$. By the first assertion of Lemma~\ref{8.6:lem-blob-chains}, some creation
$(j'',\square'')$ of $Z$ has $e$ within $k''$-distance $C_G M\check{R}_{k''}$ of the $k''$-image of
$\square''$. Moreover $\square''$ is one of at most $(4C_G M\check{R}_{k''})^4L^{4(k''-j'')}$ cubes of
the step-$j''$ partition, and $\max(j,k''-T_1)\le j''\le k''$ by the definition of the event. Hence
\begin{equation}\label{8.6:eq-young-two-cell-4}
Y_{k''}(e;j)\subset\bigcup_{j''=\max(j,k''-T_1)}^{k''}\ \bigcup_{\square''}E_{j'',\square''},
\end{equation}
where, for each $j''$, the cube $\square''$ ranges over those cubes. By subadditivity and (a), and then
using $k''-j''\le T_1$, the fact that there are at most $T_1+1$ values of $j''$, and
$q_{j''}\le e^{A_0}q_{k''}$, we get
\begin{equation}\label{8.6:eq-young-two-cell-5}
\begin{aligned}
\tilde{\nu}\bigl(Y_{k''}(e;j)\bigr)
&\le\sum_{j''}\bigl(4C_G M\check{R}_{k''}\bigr)^4L^{4(k''-j'')}C_\kappa q_{j''}\\
&\le(T_1+1)\bigl(4C_G M\check{R}_{k''}\bigr)^4L^{4T_1}C_\kappa e^{A_0}q_{k''} .
\end{aligned}
\end{equation}
Write $q_{k''}=C_q^{1/4}e^{-\frac14c_qp_0(g_{k''})}\,q_{k''}^{3/4}$. By the first inequality in
\eqref{8.6:eq-young-two-cell-6} and then by \eqref{8.6:eq-young-two-cell-1}, both read at $y=x_{k''}$,
\begin{align*}
&(T_1+1)\bigl(4C_GM\check{R}_{k''}\bigr)^4L^{4T_1}C_\kappa e^{A_0}C_q^{1/4}\\
&\qquad\le C_\kappa\bigl(4C_GM\check{R}_{k''}\bigr)^8L^{4T_1}(T_1+1)e^{A_0}C_{\mathrm{cr}}
\le e^{\frac14c_qp_0(g_{k''})} .
\end{align*}
Multiplying by $e^{-\frac14c_qp_0(g_{k''})}q_{k''}^{3/4}$ shows that the last expression in
\eqref{8.6:eq-young-two-cell-5} is at most
$q_{k''}^{3/4}$. Finally $k''\le k'$, so $q_{k''}\le e^{A_0}q_{k'}$. Hence
$q_{k''}^{3/4}\le e^{3A_0/4}q_{k'}^{3/4}\le q_{k'}^{3/4}e^{A_0}$. This is the first line of
\eqref{8.6:eq-young-two-cell-a}.

\emph{Proof of (d).} Suppose $Y_{k'}(e,e')$ occurs, and let $Z$ be such an all-young component meeting
$e$ and $e'$. By the second assertion of Lemma~\ref{8.6:lem-blob-chains}, $Z$ contains
$J\ge\lceil D/(2C_GM\check{R}_{k'})\rceil$ distinct creations $(j_1,\square_1),\dots,(j_J,\square_J)$
whose $k'$-blobs form a chain from $e$ to $e'$, consecutive blobs touching. If $D>0$, then $J\ge1$, so
$J\ge J_0$. If $D=0$, we use instead the single creation of $Z$ near $e$ given by the first assertion
of that lemma, a chain of length $J=1=J_0$. In either case $J\ge J_0$.

The first blob meets $e$, and blobs have side at most $C_G M\check{R}_{k'}$ and contain the image of
their cube. So, by the count in the first assertion of that lemma, $\square_1$ is one of at most
$(4C_G M\check{R}_{k'})^4L^{4(k'-j_1)}$ cubes. Given $\square_i$, the cube $\square_{i+1}$ is one of at
most $(4C_G M\check{R}_{k'})^4L^{4(k'-j_{i+1})}$ cubes. Since $Z$ is all-young, every level satisfies
$\max\{1,k'-T_1\}\le j_i\le k'$.

Hence $Y_{k'}(e,e')$ is contained in the union, over $J\ge J_0$ and over $J$-tuples of pairwise distinct
sites $(j_i,\square_i,\iota_i)$ obeying these constraints, of the events
$\bigcap_{i\le J}E^{\iota_i}_{j_i,\square_i}$. For each tuple, (b) gives
\begin{equation*}
\tilde{\nu}\Bigl(\bigcap_{i\le J}E^{\iota_i}_{j_i,\square_i}\Bigr)
\le\prod_{i\le J}C_{\mathrm{cr}}q_{j_i}.
\end{equation*}

Now sum over the tuples, dropping the distinctness constraint, which only enlarges the sum. Sum first
over the last site given the previous one, then successively over the earlier ones. Each summation over
one site, its kind included, is bounded by the number $\eta$ of \eqref{8.6:eq-young-two-cell-3}. Hence
$\tilde{\nu}(Y_{k'}(e,e'))\le\sum_{J\ge J_0}\eta^J$.

The argument of \eqref{8.6:eq-young-two-cell-5} and of the lines following it, with $k''=k'$ and
$j=1$, applies
verbatim to $\eta$, whose range of levels is the same; with \eqref{8.6:eq-young-two-cell-1} and the
first inequality in \eqref{8.6:eq-young-two-cell-6} read at $y=x_{k'}$ it gives
\begin{equation*}
\eta\le(T_1+1)\bigl(4C_G M\check{R}_{k'}\bigr)^4L^{4T_1}C_\kappa e^{A_0}q_{k'}\le q_{k'}^{3/4}.
\end{equation*}
By the second inequality in \eqref{8.6:eq-young-two-cell-6} read at $y=x_{k'}$, and since
$\kappa_c(x_{k'})=\tfrac12c_qp_0(g_{k'})$,
\begin{equation*}
q_{k'}^{3/4}=C_q^{3/4}e^{-\frac14c_qp_0(g_{k'})}\,e^{-\frac12c_qp_0(g_{k'})}
\le\tfrac12e^{-\kappa_c(x_{k'})} .
\end{equation*}
Hence $\eta\le\frac12$ and $\eta\le e^{-\kappa_c(x_{k'})}$. The first gives
$\sum_{J\ge J_0}\eta^J=\eta^{J_0}/(1-\eta)\le2\eta^{J_0}$.

Finally write $2\eta^{J_0}=2\eta\cdot\eta^{J_0-1}$. By the definition of $C'$ in
\eqref{8.6:eq-young-two-cell-3}, the first bound on $\eta$ is
$2\eta\le C'(M\check{R}_{k'})^4q_{k'}$. Since $J_0\ge1$ and $\eta\le e^{-\kappa_c(x_{k'})}$, one has
$\eta^{J_0-1}\le e^{-\kappa_c(x_{k'})(J_0-1)}$.
This is the second line of \eqref{8.6:eq-young-two-cell-a}.
\end{proof}

\begin{corollary}[Young-member bounds for the witness events]\label{8.6:cor-young-witness-bounds}
Assume the following.
\begin{itemize}
\item[(i)] The conditional rarity of simultaneous creations holds in the form
\eqref{8.6:eq-creation-rarity-a} of Proposition~\ref{8.6:prop-creation-rarity}, whose proof is
outstanding.
\item[(ii)] The hypotheses of Theorem~\ref{8.6:thm-young-two-cell} other than
Proposition~\ref{8.6:prop-creation-rarity} hold. Under \textup{(i)} and \textup{(ii)} that theorem
supplies its young label form and its two-cell bound \eqref{8.6:eq-young-two-cell-a} for all-young
components.
\end{itemize}

\emph{Notation.} The measure on histories $h$ with respect to which
Theorem~\ref{8.6:thm-young-two-cell} is stated is denoted $\tilde\nu$. The level-$j$ creation rate
of Proposition~\ref{8.6:prop-creation-rarity} is $q_j$; it satisfies $q_j\le C_qe^{-c_qp_0(g_j)}$,
where $p_0(\cdot)$ is the degree function, as in \eqref{8.6:eq-creation-rarity-a}, with the
constants $C_q,c_q$ of that proposition; the
level-$j$ bound supplied by the young label form is $\tilde q_j$. The constants $C_\kappa$, $C_G$,
$C'$ and the lifetime parameters $\check R_j$ are those of Theorem~\ref{8.6:thm-young-two-cell};
$C_\kappa$ is the constant with which that theorem, using hypothesis \textup{(i)} with $n=1$, bounds
the $\tilde\nu$-measure of a creation in a single cube of level $j$ by $C_\kappa q_j$. We write
$\check R$ for a number $\check R\ge1$ with $\check R_j\le\check R$ at all levels $j$ occurring
below. The rate of Theorem~\ref{8.6:thm-young-two-cell} is the function
\begin{equation*}
\kappa_c(x)=\tfrac12c_qp_0(x^{-1/2}).
\end{equation*}
The level $k_0$, the regions $N_i$ attached to the supports, the scale $\varrho$ and the constants
$c_L$, $C''$, $c'$ are those of the two-point witness theorem; $\kappa_W$ denotes the decay constant
with which that theorem bounds the crossing event of its large-field part and the live spanning
piece of its localised remainder. Ages and youth are understood as in
Theorem~\ref{8.6:thm-young-two-cell}.

Then the following hold.
\begin{itemize}
\item[(a)] Let $\mathcal{E}_i$ be the event that some large-field region of $h$ created at a level at
least $k_0-\check R-c_L$ meets $N_i$. Then
\begin{equation}\label{8.6:eq-young-witness-bounds-1}
\tilde\nu(\mathcal{E}_i)\le 2\sum_{j=k_0-\check R-c_L}^{k_0}\bigl(8\varrho L^{k_0-j}\bigr)^4\tilde q_j
=:\mathfrak{q}_W\le C''\varrho^4e^{-c'p_0(g_{k_0})};
\end{equation}
this defines $\mathfrak{q}_W$, the quantity so denoted in the two-point witness theorem. The event
$\mathcal{E}_i$ is the young member, restricted by creation level, of the label clause of the
exceptional event of the two-point witness theorem, and \eqref{8.6:eq-young-witness-bounds-1} is
the bound $\mathfrak{q}_W$ that the two-point witness theorem, and the first item of the exceptional
event of its large-field part, use.
\item[(b)] Let $r$, $D\ge1$, $\mathfrak{l}$ and $C$ be the depth, the box side (in lattice units of
the level concerned), the window length and the window constant of the clean event of the Gaussian
extraction. Let $\mathcal{C}$ denote that clean event and
$\mathcal{C}^c$ its complement, that is, the union over $j\in[k_0,k_0+r]$, over levels
$j'\in[j-\mathfrak{l}-C,j]$, and over cubes of level $j'$ lying under the level-$j$ box of side $D$
enlarged by $100M\check R_j$, of the event that a large-field region of any kind is created at
level $j'$ in that cube. We write $\mathcal{F}$ for the family of these creation events, so that
$\mathcal{C}^c$ is the union of the members of $\mathcal{F}$. The weight that the Gaussian
extraction assigns to $\mathcal{C}^c$ is at most $\sum_{E\in\mathcal{F}}\tilde\nu(E)$, and so is
$\tilde\nu(\mathcal{C}^c)$. Assume the absorption inequality
\begin{equation}\label{8.6:eq-young-witness-bounds-2}
L^{4(\mathfrak{l}+C)}(M\check R)^4(\mathfrak{l}+C+1)\le e^{\frac14c_qp_0(g_{k_0})}.
\end{equation}
This is the absorption condition of the Gaussian extraction, of the same form as the threshold
conditions of Theorem~\ref{8.6:thm-young-two-cell}: a one-sided condition on the coupling
$g_{k_0}$. Assume also the range condition of the Gaussian extraction together with its depth
restriction $r\le C_r\,p_0(g_{k_0})/\log L$, where $C_r$ is
the constant of that restriction, under which, by the arithmetic of the Gaussian extraction,
\begin{equation}\label{8.6:eq-young-witness-bounds-3}
p_0(g_{j'})\ge\tfrac12p_0(g_{k_0})
\end{equation}
for all levels $j'$ occurring in $\mathcal{C}^c$. Then
\begin{equation}\label{8.6:eq-young-witness-bounds-4}
\begin{split}
\sum_{E\in\mathcal{F}}\tilde\nu(E)
&\le C_\kappa(1+r)(\mathfrak{l}+C+1)(D+200M\check R)^4L^{4(\mathfrak{l}+C)}\max_{j'}q_{j'}\\
&\le 201^4C_\kappa C_q\,D^4(1+r)\,e^{-\frac14c_qp_0(g_{k_0})}.
\end{split}
\end{equation}
Hence both $\tilde\nu(\mathcal{C}^c)$ and the weight that the Gaussian extraction assigns to
$\mathcal{C}^c$ are at most the right-hand side of \eqref{8.6:eq-young-witness-bounds-4}.
In particular the weight bound of the Gaussian extraction for $\mathcal{C}^c$, a constant times
$D^4(1+r)$ times an exponential in $-p_0(g_{k_0})$, holds with the constant $201^4C_\kappa C_q$ and
the rate $\tfrac14c_q$. This bound uses neither the young label form of
Theorem~\ref{8.6:thm-young-two-cell} nor any survival statement for old regions; it uses only
hypothesis \textup{(i)} with $n=1$.
\item[(c)] Fix a family of at most $\varrho^8$ pairs of level-$k_0$ cells, each pair at mutual
distance at least $\varrho/4$; for a pair we write $D'$ for this distance, the distance denoted $D$
in the two-cell bound. Let $\mathcal{S}$ be the event that for some pair of the family one connected
large-field component, all of whose creations are young, contains both cells of the pair. Then
\begin{equation}\label{8.6:eq-young-witness-bounds-5}
\tilde\nu(\mathcal{S})\le C'\varrho^8(M\check R)^4q_{k_0}\,e^{\kappa_c(x_{k_0})}
\exp\Bigl(-\frac{\kappa_c(x_{k_0})\,\varrho}{8C_GM\check R}\Bigr),
\end{equation}
and, whenever $\tfrac12c_qp_0(g_{k_0})\ge C_G\kappa_W/50$, a one-sided condition on $g_{k_0}$, the
rate satisfies
\begin{equation*}
\frac{\kappa_c(x_{k_0})\varrho}{8C_GM\check R}\ge\frac{\kappa_W\varrho}{400M\check R}.
\end{equation*}
The event $\mathcal{S}$ covers two young pieces: the crossing event of the large-field part of the
two-point witness theorem, and the live spanning piece of its localised remainder.
\end{itemize}
\end{corollary}

\begin{proof}
\emph{Part (a).} Suppose a large-field region created at a level at least $k_0-\check R-c_L$ meets
$N_i$ at a level $j\le k_0$. The creation level of a component is the level of its oldest creation
\cite[(1.85)--(1.86)]{B-CMP122b}. Hence every creation entering the region took place at a level in
$[k_0-\check R-c_L,j]$, and $j$ ranges over $[k_0-\check R-c_L,k_0]$. Every creation of the region
is young at the level at which the region meets $N_i$. Consequently $\mathcal{E}_i$ is contained in
the union, over these meeting levels $j$ and over the level-$j$ cells lying under a level-$k_0$
cell of $N_i$, of the events to which the young label form of
Theorem~\ref{8.6:thm-young-two-cell} applies. Summing the young label form over these cells and
over the meeting levels gives the first inequality in \eqref{8.6:eq-young-witness-bounds-1}. The
sum is $\mathfrak{q}_W$ by definition. At meeting level $j$ the number of level-$j$ cells lying
under the level-$k_0$ cells of $N_i$ is what the factor $2(8\varrho L^{k_0-j})^4$ counts; this
factor and the last inequality of \eqref{8.6:eq-young-witness-bounds-1} are the display and the
arithmetic by which the two-point witness theorem bounds $\mathfrak{q}_W$.

\emph{Part (b).} By definition, $\mathcal{C}^c$ is the union of the members of $\mathcal{F}$,
creation events in single cubes, so its measure is at most $\sum_{E\in\mathcal{F}}\tilde\nu(E)$;
it remains to bound this sum. For a cube of level $j'$, hypothesis \textup{(i)}
with $n=1$, in the form in which Theorem~\ref{8.6:thm-young-two-cell} uses it, bounds the measure
of a creation in that cube by $C_\kappa q_{j'}$.

We count the terms. There are $1+r$ values of $j$, and for each of them $\mathfrak{l}+C+1$ values
of $j'$. The enlarged box has side $D+200M\check R_j\le D+200M\check R$ in level-$j$ units. Since
$j-j'\le\mathfrak{l}+C$,
it contains at most $(D+200M\check R)^4L^{4(\mathfrak{l}+C)}$ cubes of level $j'$. Bounding each
$q_{j'}$ by $\max_{j'}q_{j'}$ gives the first line of \eqref{8.6:eq-young-witness-bounds-4}.

For the second line, use $D\ge1$ and $M\check R\ge1$ (since $M=L^{m'}\ge1$ and $\check R\ge1$),
which give $D+200M\check R\le201\,D\,M\check R$. Then \eqref{8.6:eq-young-witness-bounds-2} bounds
\begin{equation*}
(\mathfrak{l}+C+1)(D+200M\check R)^4L^{4(\mathfrak{l}+C)}
\le201^4D^4e^{\frac14c_qp_0(g_{k_0})}.
\end{equation*}
The creation rate satisfies $q_{j'}\le C_qe^{-c_qp_0(g_{j'})}$. Together with
\eqref{8.6:eq-young-witness-bounds-3} this gives
\begin{equation*}
\max_{j'}q_{j'}\le C_qe^{-\frac12c_qp_0(g_{k_0})}.
\end{equation*}
Multiplying these two bounds by $C_\kappa(1+r)$ gives the second line of
\eqref{8.6:eq-young-witness-bounds-4}.

\emph{Part (c).} Apply the two-cell bound \eqref{8.6:eq-young-two-cell-a} of
Theorem~\ref{8.6:thm-young-two-cell} at level $k'=k_0$. For two level-$k_0$ cells at distance $D'$,
it bounds the measure of the event that one all-young connected component contains both by
\begin{equation*}
C'(M\check R)^4q_{k_0}\,e^{-\kappa_c(x_{k_0})(n_0-1)},\qquad
n_0=\Bigl\lceil\frac{D'}{2C_GM\check R}\Bigr\rceil,
\end{equation*}
where every chain of distinct young creations joining the two cells has at least $n_0$ members.
Since $D'\ge\varrho/4$, we have $n_0\ge\varrho/(8C_GM\check R)$, and therefore
\begin{equation*}
e^{-\kappa_c(x_{k_0})(n_0-1)}
\le e^{\kappa_c(x_{k_0})}\exp\Bigl(-\frac{\kappa_c(x_{k_0})\varrho}{8C_GM\check R}\Bigr).
\end{equation*}
Summing over the at most $\varrho^8$ pairs of the family gives
\eqref{8.6:eq-young-witness-bounds-5}. Since $x_{k_0}=g_{k_0}^{-2}$, we have
$\kappa_c(x_{k_0})=\tfrac12c_qp_0(g_{k_0})$; under the condition stated in (c) this is at least
$C_G\kappa_W/50$, and multiplying by $\varrho/(8C_GM\check R)$ shows that the rate
$\kappa_c(x_{k_0})\varrho/(8C_GM\check R)$ so obtained is at least $\kappa_W\varrho/(400M\check R)$,
the rate with which the two-point witness theorem bounds the crossing event and the live spanning
piece of its localised remainder. Both of these, in their young part, are events of the form
$\mathcal{S}$.
\end{proof}

\begin{lemma}[Survival bound for long-lived creations. Proof outstanding]\label{8.6:lem-survival-bound}
We use the sites, creations, blobs, live components and the numbers $q_j$ of
Definition~\ref{8.6:def-creations-blobs} and~\eqref{8.6:eq-creations-blobs-a}, and the measure
$\tilde\nu$ on histories with respect to which Lemma~\ref{8.6:lem-peeling} is stated: the
measure on histories in which the conditional probability of a child given its parent is the
ratio of the weight of the child to the weight of the parent.
There are a constant $C_{\mathrm{S}}$ and a function $\eta_{\mathrm{S}}(\cdot)$ of the inverse
squared coupling, neither depending on the site $(j',\square')$ or on $a$, such that the
following holds. For every site $(j',\square')$ and every $a\ge 0$ with $k':=j'+a\le K$, let
$E_{j',\square'}$ be the event that the creation at $(j',\square')$ occurs. Put
\begin{equation*}
y_a:=x_{j'+\lceil a/2\rceil},
\qquad
N_\dagger(y):=\bigl\lfloor e^{\,c_q\,p_0(y^{-1/2})/128}\bigr\rfloor .
\end{equation*}
Then
\begin{equation}\label{8.6:eq-survival-bound-1}
\begin{split}
&\tilde\nu\bigl\{E_{j',\square'}\ \text{and the material of $\square'$ belongs to a component
live at level $k'$}\bigr\}\\
&\qquad\le C_{\mathrm{S}}\, q_{j'}\cdot
\max\bigl\{e^{-(4\log L+1)a},\ e^{-\eta_{\mathrm{S}}(y_a)\,N_\dagger(y_a)}\bigr\}.
\end{split}
\end{equation}
Here $p_0(\cdot)$ is the degree function of the coupling and $c_q$ is the constant in the
definition of $q_j$.
\end{lemma}

Proof outstanding.

\emph{Route.} A proof would have to establish~\eqref{8.6:eq-survival-bound-1} by
combinatorics on the region rules~\eqref{8.6:eq-creations-blobs-a} (among them the quiet
lifetime of \cite[(1.81)]{B-CMP122b}), together with
Lemma~\ref{8.6:lem-peeling}. Lemma~\ref{8.6:lem-peeling} is proved as an implication from
Proposition~\ref{8.6:prop-creation-rarity}, whose proof is outstanding, and the tree structure
of the measure on histories. The steps are the
following.
\begin{enumerate}
\item[(s1)] By the monotonicity rule in~\eqref{8.6:eq-creations-blobs-a}, being live at level
$k'$ is an up-set in the creation configuration. The event on the left
of~\eqref{8.6:eq-survival-bound-1} is therefore a union over the minimal families of creations
that sustain the component. The $\tilde\nu$-measure of the joint occurrence of each such family
is bounded by Lemma~\ref{8.6:lem-peeling}.
\item[(s2)] Each minimal sustaining family is organised by a spanning tree of Kruskal type,
whose vertices are the creations of the family, each standing for its blob, and whose edges are
ordered by the level at which the blobs they join first merge.
\item[(s3)] The pairs formed by a creation of the family and a level at which its blob is
present are charged injectively to creation events, so that no creation event is charged by
two such pairs.
\item[(s4)] The small factors gained from the creations in the first half of the interval of
levels from $j'$ to $k'$ are evaluated at the mid-life inverse coupling $y_a$, under an
auxiliary condition on which this step depends.
\item[(s5)] A truncation at $N_\dagger(y_a)$, monotone in the level, is applied.
\end{enumerate}

The lemma enters only the per-step organisation of the large-field block bound. In the
organisation in which that bound is read at the terminal level, the contribution of old live
components is absorbed by a spare small factor available at the step at which the component
heals.

\begin{proposition}[Reduction of creation rarity to a local exponential moment]
\label{8.6:prop-exponential-moment}
Let $0\le k<K$, let $h_k$ be a level-$k$ history with non-negative term $\tau_{h_k}$ and mass
$m(h_k)=\int\tau_{h_k}\,dV_k>0$, the total integral of $\tau_{h_k}$ as in
Lemma~\ref{8.6:lem-mass-conservation}, and
let $\langle\cdot\rangle_{h_k}$ be the average with respect to $\tau_{h_k}\,dV_k/m(h_k)$. Let
$\square$ be a cube of the level-$(k+1)$ partition of \cite[(2.17)]{B-CMP119} lying above the
interior of the small-field region $\Omega_k(h_k)$ of $h_k$, and let $\tilde\square$ be its
enlargement in that partition. Put
\begin{equation*}
  E=E^{\mathrm{basic}}_{k+1,\square}
  =\Bigl\{V_{k+1}:\ \sup_{p\in\square}\bigl|U_{k+1,\square}(V_{k+1},\partial p)-\mathbf{1}\bigr|
  \ge\varepsilon_{k+1}\eta_{k+1}^2\Bigr\},
\end{equation*}
where $\varepsilon_{k+1},\eta_{k+1}>0$ are the numbers of Ba{\l}aban's basic small-field
condition at level $k+1$ in \cite[(2.17)]{B-CMP119}, $\varepsilon_{k+1}$ the threshold and
$\eta_{k+1}^2$ the factor multiplying it, and $U_{k+1,\square}(V_{k+1},\cdot)$ is
Ba{\l}aban's level-$(k+1)$ field on
$\square$ built from the block variables $V_{k+1}$. Let $f_{h_k}$ be the image of $\tau_{h_k}$
under Ba{\l}aban's renormalisation transformation, that is, the push-forward of
$\tau_{h_k}\,dV_k$ under the block-averaging map of Definition~\ref{1.5:def-block-average},
and put
\begin{equation*}
  P(E\mid h_k)=\frac{\int_E f_{h_k}}{\int f_{h_k}} .
\end{equation*}
By Lemma~\ref{8.6:lem-mass-conservation}, $P(E\mid h_k)$ is the ratio
$\sum_{h_{k+1}}m(h_{k+1})/m(h_k)$ over the children $h_{k+1}$ of $h_k$ that contain the
creation of basic kind at $(k+1,\square)$ in the sense of
Definition~\ref{8.6:def-creations-blobs}, the conditional probability of
Proposition~\ref{8.6:prop-creation-rarity}, whose proof is outstanding.
For a configuration $V_k$, with $U_k(V_k)$ the level-$k$ configuration determined by $V_k$ in
Ba{\l}aban's construction, whose plaquette variables enter the effective Wilson action, and for
a plaquette $p\subset\tilde\square$ of its lattice, put
$F_p:=U_k(V_k)(\partial p)-\mathbf{1}$ and $|F_p|^2:=\operatorname{tr}F_p^*F_p$, and let
$A(\cdot)|_{\tilde\square}$ be the Wilson action restricted to the plaquettes in
$\tilde\square$. Since $U_k(V_k)(\partial p)$ is unitary, $\operatorname{tr}\mathbf{1}=1$ and
$\operatorname{tr}U^*=\overline{\operatorname{tr}U}$, one has
$|F_p|^2=2-2\operatorname{Re}\operatorname{tr}U_k(V_k)(\partial p)$, hence
\begin{equation}\label{8.6:eq-exponential-moment-5}
  A(U_k(V_k))|_{\tilde\square}
  =\sum_{p\subset\tilde\square}\bigl[1-\operatorname{Re}\operatorname{tr}U_k(V_k)(\partial p)\bigr]
  =c_A\sum_{p\subset\tilde\square}|F_p|^2,
  \qquad c_A=\tfrac12 .
\end{equation}
Assume the following two hypotheses.
\begin{itemize}
\item[(i)] The bound of \cite[Theorem~1]{B-CMP102b} (see also \cite{B-CMP109}) on the block
field in terms of the block variables: there are $\alpha_B,C_B>0$ such
that for every $V$ with
$\mathfrak{m}:=\max_{p'\subset\tilde\square}|V(\partial p')-\mathbf{1}|\le\alpha_B$, $p'$
running over the plaquettes of the level-$(k+1)$ lattice in $\tilde\square$,
\begin{equation*}
  \sup_{p\in\square}\bigl|U_{k+1,\square}(V,\partial p)-\mathbf{1}\bigr|
  \le C_B\,\eta_{k+1}^2\,\mathfrak{m}.
\end{equation*}
\item[(ii)] The averaging energy inequality of \cite{B-CMP95}, used here as a hypothesis in
the following linearised form; its form for the non-abelian block-averaging map is not proved
here. There is a number $\delta\ge0$ such that for every
$V_{k+1}$ and every $V_k$ in the fibre of the block-averaging map over $V_{k+1}$, with
$\mathfrak{m}$ computed from $V=V_{k+1}$,
\begin{equation*}
  \mathfrak{m}^2\le(1+\delta)\sum_{p\subset\tilde\square}|F_p|^2 .
\end{equation*}
\end{itemize}
Assume moreover that
\begin{equation*}
  \varepsilon_{k+1}\le C_B\,\alpha_B,
  \qquad
  x_k\,\varepsilon_{k+1}^2\ge(1+\delta)\,p_0(g_{k+1})^2 .
\end{equation*}
Put $t_1=\varepsilon_{k+1}/C_B$ and $c''=c_A/C_B^2$. Then:
\begin{itemize}
\item[(a)] for $0<\theta\le\tfrac12$,
\begin{equation}\label{8.6:eq-exponential-moment-1}
  P\bigl(E^{\mathrm{basic}}_{k+1,\square}\bigm| h_k\bigr)
  \le e^{-\theta c''p_0(g_{k+1})^2}\,
  \Bigl\langle\exp\bigl(\theta x_k\,A(U_k(V_k))|_{\tilde\square}\bigr)\Bigr\rangle_{h_k};
\end{equation}
\item[(b)] likewise, for an integer $n\ge1$ and pairwise distinct cubes
$\square_1,\dots,\square_n$ of the level-$(k+1)$ partition, each as above, and
$0<\theta\le\tfrac12$,
\begin{equation}\label{8.6:eq-exponential-moment-2}
  P\Bigl(\bigcap_{i\le n}E^{\mathrm{basic}}_{k+1,\square_i}\Bigm| h_k\Bigr)
  \le e^{-n\theta c''p_0(g_{k+1})^2}\,
  \Bigl\langle\exp\Bigl(\theta x_k\sum_{i\le n}A(U_k(V_k))|_{\tilde\square_i}\Bigr)
  \Bigr\rangle_{h_k};
\end{equation}
\item[(c)] if, with $n_\square$ the number of cells of $\tilde\square$, there are
$0<\theta_0\le\tfrac12$
and $C_M$ independent of $k$, $K$, $m$, $h_k$ and the position of the cubes such that
\begin{equation}\label{8.6:eq-exponential-moment-a}
  \Bigl\langle\exp\Bigl(\theta_0x_k\sum_{i\le n}A(U_k(V_k))|_{\tilde\square_i}\Bigr)
  \Bigr\rangle_{h_k}\le e^{C_Mn\cdot n_\square},
\end{equation}
then
\begin{equation}\label{8.6:eq-exponential-moment-3}
  P\Bigl(\bigcap_{i\le n}E^{\mathrm{basic}}_{k+1,\square_i}\Bigm| h_k\Bigr)
  \le\bigl(q'_{k+1}\bigr)^n,
  \qquad q'_{k+1}=e^{C_Mn_\square-\theta_0c''p_0(g_{k+1})^2},
\end{equation}
which, when $\theta_0$ and $C_M$ can be taken independent of $n$ and
$q'_{k+1}\le e^{-p_0(g_{k+1})}$, gives the bound \eqref{8.6:eq-creation-rarity-a} of
Proposition~\ref{8.6:prop-creation-rarity}, whose proof is outstanding, for the basic kind at
level $k+1$ with each of its three constants equal to $1$.
\end{itemize}
\end{proposition}

Parts (a)--(c) are proved below from hypotheses (i) and (ii) and the two conditions on
$\varepsilon_{k+1}$; the bound \eqref{8.6:eq-exponential-moment-a} itself, the local
exponential moment of the effective Wilson action density, of order one per cell, is not
proved. The proposition is not used in the proof of
Theorem~\ref{3.1:thm-main}; it reduces Proposition~\ref{8.6:prop-creation-rarity}, whose proof
is outstanding, to \eqref{8.6:eq-exponential-moment-a} for the basic kind.

\begin{proof}
\emph{The block deviation on $E$.} Let $V_{k+1}\in E$ and let $\mathfrak{m}$ be computed from
$V=V_{k+1}$. If $\mathfrak{m}>\alpha_B$, then $\mathfrak{m}\ge t_1$, since
$\alpha_B\ge\varepsilon_{k+1}/C_B$ by the first condition on $\varepsilon_{k+1}$ and
$\varepsilon_{k+1}/C_B=t_1$. If $\mathfrak{m}\le\alpha_B$, then
hypothesis (i) and the definition of $E$ give
\begin{equation*}
  \varepsilon_{k+1}\eta_{k+1}^2
  \le\sup_{p\in\square}\bigl|U_{k+1,\square}(V_{k+1},\partial p)-\mathbf{1}\bigr|
  \le C_B\,\eta_{k+1}^2\,\mathfrak{m},
\end{equation*}
hence $\mathfrak{m}\ge\varepsilon_{k+1}/C_B=t_1$. Thus $\mathfrak{m}\ge t_1$ on $E$.

\emph{The action on $E$.} Let $V_{k+1}\in E$ and let $V_k$ be any configuration in the fibre
over $V_{k+1}$. The identity \eqref{8.6:eq-exponential-moment-5}, the inequality of
hypothesis (ii), the bound $\mathfrak{m}\ge t_1\ge0$ just
proved, $x_k>0$ and the second condition on $\varepsilon_{k+1}$ give
\begin{equation}\label{8.6:eq-exponential-moment-4}
\begin{aligned}
  x_k\,A(U_k(V_k))|_{\tilde\square}
  &=c_A\,x_k\sum_{p\subset\tilde\square}|F_p|^2
  \ge\frac{c_A\,x_k\,\mathfrak{m}^2}{1+\delta}
  \ge\frac{c_A}{C_B^2}\,\frac{x_k\,\varepsilon_{k+1}^2}{1+\delta}\\
  &\ge\frac{c_A}{C_B^2}\,p_0(g_{k+1})^2=c''p_0(g_{k+1})^2 .
\end{aligned}
\end{equation}

\emph{Proof of (a).} Write $V_{k+1}=V_{k+1}(V_k)$ for the block average of $V_k$. Since
$f_{h_k}$ is the push-forward of $\tau_{h_k}\,dV_k$ under the block-averaging
map, writing $\mathbf{1}_E$ for the indicator function of $E$,
\begin{equation*}
  \int_E f_{h_k}=\int\mathbf{1}_E\bigl(V_{k+1}(V_k)\bigr)\,\tau_{h_k}(V_k)\,dV_k,
  \qquad\int f_{h_k}=m(h_k),
\end{equation*}
the second identity by Lemma~\ref{8.6:lem-mass-conservation}. By
\eqref{8.6:eq-exponential-moment-4}, for $\theta>0$ and every $V_k$ with $V_{k+1}(V_k)\in E$,
\begin{equation*}
  \exp\Bigl(\theta\bigl(x_kA(U_k(V_k))|_{\tilde\square}-c''p_0(g_{k+1})^2\bigr)\Bigr)\ge1,
\end{equation*}
and the left side is positive for every $V_k$. Hence
\begin{equation*}
  \mathbf{1}_E\bigl(V_{k+1}(V_k)\bigr)
  \le e^{-\theta c''p_0(g_{k+1})^2}\exp\bigl(\theta x_kA(U_k(V_k))|_{\tilde\square}\bigr)
\end{equation*}
for every $V_k$. Integrating against the non-negative density
$\tau_{h_k}$ and dividing by $m(h_k)>0$ gives
\eqref{8.6:eq-exponential-moment-1}.

\emph{Proof of (b).} For each $i$ the two preceding steps, applied to $\square_i$, give
\begin{equation*}
  \mathbf{1}_{E_i}\bigl(V_{k+1}(V_k)\bigr)
  \le\exp\Bigl(\theta\bigl(x_kA(U_k(V_k))|_{\tilde\square_i}-c''p_0(g_{k+1})^2\bigr)\Bigr)
\end{equation*}
with $E_i=E^{\mathrm{basic}}_{k+1,\square_i}$. The indicator of the
intersection is the product of the $n$ indicators, so
\begin{equation*}
  \prod_{i\le n}\mathbf{1}_{E_i}\bigl(V_{k+1}(V_k)\bigr)
  \le e^{-n\theta c''p_0(g_{k+1})^2}
  \exp\Bigl(\theta x_k\sum_{i\le n}A(U_k(V_k))|_{\tilde\square_i}\Bigr);
\end{equation*}
no disjointness of the $\tilde\square_i$ is used. Integrating against $\tau_{h_k}\,dV_k/m(h_k)$
as in (a) gives \eqref{8.6:eq-exponential-moment-2}.

\emph{Proof of (c).} Since $0<\theta_0\le\tfrac12$, (b) applies with
$\theta=\theta_0$. Inserting \eqref{8.6:eq-exponential-moment-a} into the right side of
\eqref{8.6:eq-exponential-moment-2} at $\theta=\theta_0$ gives
\begin{equation*}
  P\Bigl(\bigcap_{i\le n}E_i\Bigm|h_k\Bigr)
  \le e^{-n\theta_0c''p_0(g_{k+1})^2}\,e^{C_Mn\cdot n_\square}
  =\bigl(e^{C_Mn_\square-\theta_0c''p_0(g_{k+1})^2}\bigr)^n=\bigl(q'_{k+1}\bigr)^n,
\end{equation*}
which is \eqref{8.6:eq-exponential-moment-3}.
\end{proof}

Since $q'_{k+1}=e^{C_Mn_\square-\theta_0c''p_0(g_{k+1})^2}$, one has
$q'_{k+1}\le e^{-p_0(g_{k+1})}$ whenever
$\theta_0c''p_0(g_{k+1})^2-C_Mn_\square\ge p_0(g_{k+1})$, a condition met under a ceiling on
$\gamma$ in the sense of Definition~\ref{2.3:not-stages-first}. For the decompositions of unity
of fluctuation kind (Definition~\ref{8.6:def-creations-blobs}), the analogue of
\eqref{8.6:eq-exponential-moment-a} is a bound on the
average $\langle e^{\theta Y^2/(2\sigma^2)}\rangle$ in the fibre of the step, where $Y$ is the
fluctuation variable restricted by that decomposition and $\sigma^2$ its variance under the
step's Gaussian measure, the change of
covariance being of rank one; when the history $h_k$ carries no live large-field region the
analogue is a statement of the type of the localisation of the sourced step, and for a general
history it is of the type of \eqref{8.6:eq-exponential-moment-a}.
In \cite[(1.77)--(1.78)]{B-CMP122b} the absolute part of the per-cube extraction is taken and
the normalising integral that remains is estimated by $1$; the bound
\eqref{8.6:eq-exponential-moment-a} is a bound on that normalising integral with the factor
$\exp\bigl(\theta_0x_k\sum_{i\le n}A(U_k(V_k))|_{\tilde\square_i}\bigr)$ retained.

\begin{remark}[The infrared depth of the per-step route and the cluster-expansion comparison]
\label{8.6:rem-infrared-depth}
Throughout, $\tilde{\nu}$ denotes the probability measure on histories and fields defined by the
effective densities of Definition~\ref{1.5:def-densities}, normalised by the total integral.
$E_{k+1,\square}$ denotes the event that a large-field region is created at step $k+1$ in the
cube $\square$, and $V_K$ denotes the field at the top level $K$, which lives on the unit lattice.
\begin{enumerate}
\item[(a)] \emph{Infrared depth.} The following is a heuristic orientation, not an assertion of
this remark. The flux of a unit-lattice field through a plaquette is at most $\pi$, however large
the field, and block averaging dilutes flux. On this heuristic, such a field forces on the
level-$k$ field constrained by it an additional curvature of order $L^{-2(K-k)}$ per plaquette,
and for a constant $C_{\mathrm{IR}}(L)$ depending on $L$ this additional curvature is below one
half of the large-field threshold at level $k$ as soon as
\begin{equation}\label{8.6:eq-infrared-depth-1}
K-k\;\ge\;\tfrac14\log_L x_k+C_{\mathrm{IR}}(L).
\end{equation}
On the same heuristic, at levels satisfying \eqref{8.6:eq-infrared-depth-1} a proof of
conditional rarity of simultaneous creations (Proposition~\ref{8.6:prop-creation-rarity}, whose
proof is outstanding) may condition on $V_K$ and is thereby localised in the volume, while within
the last $\frac14\log_L x_k$ levels below the top the large deviations of the top-level theory
enter such a proof. If a proof of Proposition~\ref{8.6:prop-creation-rarity} covered only the
steps $j$ with
$K-j\ge s_{\mathrm{IR}}(g)$, where
\[
s_{\mathrm{IR}}(g):=\bigl\lceil\tfrac14\log_L g^{-2}\bigr\rceil+C_{\mathrm{IR}}(L),
\]
then every statement with depth floor $s$ that reads large-field labels up to level
$K-s+r$, where $r$ is its extraction depth (Notation~\ref{0.2:not-glossary}), would need
its depth floor $s$ replaced by
\[
\max\{s,\ s_{\mathrm{IR}}(g)+r\}.
\]
This is a one-sided condition depending on $g$; no hypothesis of Theorem~\ref{3.1:thm-main} and
no statement of this chapter carries such a floor. Such a
floor has the form $\frac14\log_L g^{-2}+C$ of the depth in clause (ii-d) of
Theorem~\ref{3.1:thm-main}, and it is coarser than the depth
$O(\log_L\log g^{-2})\vee m_{\mathbb{T}}(g)$ at which the two-point witness theorem reads
large-field labels. Nothing is
asserted here about Proposition~\ref{8.6:prop-creation-rarity} within the last
$s_{\mathrm{IR}}(g)$ levels; item~(a) only identifies where a proof of it must use the top-level
theory.

Removal domination for a live component (Proposition~\ref{8.6:prop-live-removal}, whose proof is
outstanding), placed at the terminal level where no further step is taken, is a comparison at one
level, uniform in the exterior, with the infrared content held fixed and common to both sides; it
carries no floor of this kind, and it is the formulation adopted in this chapter.
\item[(b)] \emph{Comparison with the cluster-expansion route.} Let $\Xi$ be the partition
function of the polymer representation of the large-field terms at level $k$, after its Mayer
step. For polymers $X_0'$ whose core $\operatorname{core}X_0'$ (a set of cubes fixed by the
representation and bound of the large-field terms) contains the cube $\mathfrak{c}$,
let $F_{\mathfrak{c}}(X_0')$ be the activity of $X_0'$ in that representation, pinned at
$\mathfrak{c}$ (the pinned activity), and let
$\Xi_{\operatorname{comp}(X_0')}$ be the partition function of the polymers compatible with
$X_0'$. Let $\mathcal{Y}_{\mathfrak{c}}/\mathcal{Y}$ be the ratio of the large-field sum restricted
to families in which the cube $\mathfrak{c}$ lies in an island to the full large-field sum. For a
set of cubes, $N(\cdot)$ denotes the number of its cubes, as in \eqref{8.6:eq-live-removal-a}, and
$K_{\mathrm{KP}}(d)$ is the constant of the cluster-expansion lemma that accompanies the
representation and bound of the large-field terms in its polymer form.
Under the parameter conditions \eqref{8.6:eq-parameter-conditions-a}, the cluster-expansion route
gives
\begin{equation}\label{8.6:eq-infrared-depth-2}
\begin{aligned}
\frac{\mathcal{Y}_{\mathfrak{c}}}{\mathcal{Y}}
&=\sum_{X_0':\ \operatorname{core}X_0'\ni\mathfrak{c}}
F_{\mathfrak{c}}(X_0')\,\frac{\Xi_{\operatorname{comp}(X_0')}}{\Xi}
\le\sum_{X_0':\ \operatorname{core}X_0'\ni\mathfrak{c}}
|F_{\mathfrak{c}}(X_0')|\,e^{N(\operatorname{core}X_0')}\\
&\le K_{\mathrm{KP}}(d)\cdot10^{4d}\cdot e^{-\frac32 p_0(g_k)}.
\end{aligned}
\end{equation}
The right-hand side is $K_{\mathrm{KP}}(d)$ times the quantity $\varepsilon_k(\mathcal{S})$ of
Lemma~\ref{8.6:lem-healed-removal} at $\theta=0$ (where $w_0\equiv1$). This route gives the
one-point bound and the bound conditional on the presence of given polymers. It does not give the
bound conditional on the presence of given islands unless a lower bound on the pinned activities
is available. The positivity route of Lemma~\ref{8.6:lem-healed-removal} has no such limitation.
\item[(c)] \emph{Second-class components.} Components prepared at step $k$ and returned to the
expansion together with a large-field restriction inherited from the preparatory step of
\cite[\S1]{B-CMP122b} (a preparatory seed) are live, they fall under part $(\beta)$ of the
majorant proposition accompanying the representation and bound of the large-field terms, and
they require the hull integral of
\eqref{8.6:eq-live-removal-a} (Proposition~\ref{8.6:prop-live-removal}, whose proof is
outstanding) in the same way as the other live components. No separate statement is made for
them.
\end{enumerate}
\end{remark}

\begin{proof}
\emph{Item~(a).} By the positivity statement, in the form of the tree representation of the
effective densities by non-negative terms, $\tilde{\nu}$ is a positive measure that disintegrates
along the levels of the effective densities, from level $k+1$ up to the top level $K$. The law of
total probability then gives
\[
\tilde{\nu}(E_{k+1,\square})
=\mathbb{E}_{\tilde{\nu}}\bigl[\tilde{\nu}\bigl(E_{k+1,\square}\bigm| V_K,\ \text{top-level data}\bigr)\bigr].
\]
The discussion of flux dilution in item~(a), and of the levels at which conditioning on $V_K$
localises a proof, is heuristic and is not used below. The display shows only how a proof
localised in the volume would be organised: the conditional probability inside the expectation
would be estimated for each fixed $V_K$, and the estimate integrated afterwards against
$\tilde{\nu}$.

The remaining assertion of item~(a) follows from the definition of $s_{\mathrm{IR}}(g)$. A proof
valid only
for $K-j\ge s_{\mathrm{IR}}(g)$ gives nothing at the levels above $K-s_{\mathrm{IR}}(g)$, so a
statement with depth floor $s$ that reads large-field labels up to level $K-s+r$ needs
$K-s+r\le K-s_{\mathrm{IR}}(g)$ as well; that is, its depth floor $s$ must be replaced by
$\max\{s,\ s_{\mathrm{IR}}(g)+r\}$.

Proposition~\ref{8.6:prop-live-removal}, whose proof is outstanding, compares two integrals over
the field on the hull of a live component $C$ at one and the same level, with common exterior
data, uniformly in the
exterior. The infrared content of the top scale is contained in that exterior configuration, and
it is held fixed and common to both sides. No condition on the depth therefore enters that
comparison.

\emph{Item~(b).} The first line of \eqref{8.6:eq-infrared-depth-2} is obtained as follows. Apply
the polymer representation of the large-field terms at level $k$, after its Mayer step, and
collect the terms according to the polymer $X_0'$ whose core contains $\mathfrak{c}$. This is the
pinned-polymer form of the one-point quotient.

The first inequality in \eqref{8.6:eq-infrared-depth-2} comes from the bound on ratios of
partition functions in the cluster-expansion lemma that accompanies the representation and bound
of the large-field terms, used in the following form. The Koteck\'y--Preiss
condition~\cite{KP86} of Lemma~\ref{8.6:lem-hard-core-gas}(d) holds for the majorant
weights of these activities with $a(\cdot)=N(\operatorname{core}\,\cdot\,)$, by part $(\beta)$ of
the majorant proposition accompanying the representation and bound of the large-field terms, in
which that condition is established for these weights. Under that condition the
cluster-expansion lemma bounds $|\Xi_{\operatorname{comp}(X_0')}/\Xi|$ by
$e^{N(\operatorname{core}X_0')}$.

The second inequality comes from the proof of the activity estimate accompanying the
representation and bound of the large-field terms, with one change: the factor
$e^{-\frac12p_0(g_k)}$ extracted there is replaced by
the entropy factor $10^{4d}$ of the pinned island. The estimate is combined with the
corollary accompanying the representation and bound of the large-field terms in its polymer
form, and with the parameter conditions
\eqref{8.6:eq-parameter-conditions-a}.

Since the constant $K_H$ of Lemma~\ref{8.6:lem-healed-removal} equals $10^{4d}$ and $w_0\equiv1$
there, the right-hand side equals
$K_{\mathrm{KP}}(d)\,\varepsilon_k(\mathcal{S})$ at $\theta=0$. Thus the two routes give bounds of
the same form. This route reads the probabilities as derivatives of the logarithm of a marked
partition function with respect to the marking parameter, evaluated at zero. Marking given
polymers in this way yields the one-point bound and the bound conditional on the presence of those
polymers.

The form conditional on the presence of given islands would require a lower bound on the pinned
activities, which after the Mayer step are signed. Lemma~\ref{8.6:lem-healed-removal} works instead
with the non-negative island weights of a positive hard-core gas. There the conditional bound of
Lemma~\ref{8.6:lem-hard-core-gas}(b) holds given any compatible family, and no lower bound is
needed.

\emph{Item~(c).} Consider the second-class components of \cite[\S1]{B-CMP122b}, prepared at step
$k$ and returned with such a preparatory seed. They are live components in the sense of
Proposition~\ref{8.6:prop-live-removal}, whose proof is outstanding. For such
a component $Y$, the operator $\mathbf{T}'_k(Y)$ exists in the normalised form
\cite[(1.71)]{B-CMP122b}, and \cite[(1.79)]{B-CMP122b} holds for it with what is there called the
new factor $\exp(-p_0(g_j))$, $j$ being the step index in the notation of that display.

The preparatory large-field function that makes $Y$ second-class may, however, be the one
attached to the top window of \cite[(1.3)--(1.4)]{B-CMP122b}. In that case the term in which $Y$
is absent can vanish
at configurations where the term in which $Y$ is present does not. A pointwise ratio of the two
terms is then unavailable, and the comparison must be made after integrating the field over the
hull of $Y$, as in \eqref{8.6:eq-live-removal-a}. Since that display already contains the integral
over the hull of the component, no separate statement is needed for second-class
components.
\end{proof}

\begin{definition}[Creation events, live-cell events and the one-point bounds]
\label{8.7:def-creation-live-events}
Fix $K$ and an admissible run of Ba{\l}aban's procedure with terminal level $K'$. Write $\Omega$
for the set of its label histories $\omega$ (the histories $h$ of the glossary, written $\omega$
in this chapter), $\tau_{\omega}>0$ for the mass of the history
$\omega$, $Z=\sum_{\omega\in\Omega}\tau_{\omega}$ for the total mass, and $\tilde{\nu}$ for the
normalised measure on label events,
\begin{equation*}
\tilde{\nu}(E)=Z^{-1}\sum_{\omega\in E}\tau_{\omega},\qquad E\subset\Omega .
\end{equation*}
The bounds below are unconditioned: they concern $\tilde{\nu}$ itself, not $\tilde{\nu}$
conditioned on an event.

\emph{Creation events.} Let $j$ be a step, $\kappa$ a kind (here $\kappa$ indexes the kinds of
characteristic functions of the step; it is distinct from the auxiliary parameter $\kappa$ of
$\mathfrak{p}$), and $\square$ a cube of the step-$j$
decomposition of unity of kind $\kappa$, with characteristic function
$\chi^{\kappa}_{j,\square}$. The \emph{creation event} at $(j,\square)$ of kind $\kappa$ is
\begin{equation*}
E^{\kappa}_{j,\square}:=\bigl\{\omega\in\Omega:\ \text{the step-$j$ decomposition of $\omega$
chose the term } 1-\chi^{\kappa}_{j,\square}\bigr\},
\end{equation*}
that is, the event that a large-field region is created at step $j$ at the cube $\square$.
We write $n_{\square}$ for the number of level-$j$ cells of $\square$.

\emph{Live-cell events.} For a level $k$ and a level-$k$ cell $\square$, let $|\mathcal{O}_k(\omega)|$
be the union of the large-field components of $\omega$ that are live at level $k$. The
\emph{live-cell event} is
\begin{equation*}
\mathfrak{lv}_k(\square):=\bigl\{\omega\in\Omega:\ \square\subset|\mathcal{O}_k(\omega)|\bigr\},
\end{equation*}
also written $L^{(k)}_{\square}$ (an event, not a power of the blocking factor): the event that
$\square$ lies in a component live at level $k$, whatever the step at which that
component was created and whatever its kind, components of the second class included.

\emph{The one-point bounds.} They are the two inequalities
\begin{equation}\label{8.7:eq-creation-live-events-a}
\begin{gathered}
\sum_{\omega\in E^{\kappa}_{j,\square}}\tau_{\omega}\le q_j\, Z,\qquad
q_j:=C_q\, e^{-c_q\, p_0(g_j)},\\
\tilde{\nu}\bigl(\mathfrak{lv}_k(\square)\bigr)\le q^{lv}_k
:=C_{lv}\exp\bigl(-c_{lv}\, p_0(g_k)\bigr),
\end{gathered}
\end{equation}
where $p_0(\cdot)$ is the degree function and $g_j$, $g_k$ are the running couplings at levels $j$
and $k$. By the definition of $\tilde{\nu}$, the first line of
\eqref{8.7:eq-creation-live-events-a} is the inequality
$\tilde{\nu}(E^{\kappa}_{j,\square})\le q_j$.
\begin{itemize}
\item[(a)] We say that \emph{rarity of young large-field blocks} holds if there are constants
$C_q,c_q>0$, depending only on $L$, $M$, $\mathfrak{p}$ and the rank $N$ of $\mathrm{SU}(N)$,
such that the first line of
\eqref{8.7:eq-creation-live-events-a} holds for every $K$, every admissible run, every step $j$,
every kind $\kappa$ and every cube $\square$ as above.
\item[(b)] We say that the \emph{one-cell live bound at level $k$} holds if there are constants
$C_{lv}\ge1$ and $c_{lv}>0$, which read only $d$, $L$, $M$, $\mathfrak{p}$, $\kappa_1$, $\theta$
and further fixed constants of the construction, and are uniform in $K$, in $m$, in $N_T$ and
in the position of $\square$,
such that the second line of \eqref{8.7:eq-creation-live-events-a} holds for every $K$, every
admissible run with terminal level $K'\ge k$ and every level-$k$ cell $\square$.
\end{itemize}
\end{definition}

\begin{lemma}[A creation event lies in the live-cell events it meets]\label{8.7:lem-creation-in-live}
Let $E^{\kappa}_{j,\square}$ be the creation event and $L^{(j)}_{\mathfrak{c}}$ the live-cell event of
Definition~\ref{8.7:def-creation-live-events}, and let $n_{\square}$ be the number of level-$j$ cells of $\square$.
For every $K$, every torus size $N_T$, every level $j$, every admissible run with terminal level $K'\ge j$,
every kind $\kappa$ and every $\kappa$-cube $\square$,
\begin{equation}\label{8.7:eq-creation-in-live-a}
E^{\kappa}_{j,\square}\;\subset\;\bigcup_{\substack{\mathfrak{c}\ \text{a level-}j\ \text{cell}\\ \mathfrak{c}\cap\square\neq\emptyset}}
L^{(j)}_{\mathfrak{c}} ,
\end{equation}
a union over $n_{\square}$ cells; there is one such cell if the $\kappa$-cube is smaller than a cell, which is the case
when $LM_2\le M$. Consequently, if at level $j$ the live-cell one-point bound of
Definition~\ref{8.7:def-creation-live-events} holds with constants $C_{lv},c_{lv}$, that is,
\begin{equation}\label{8.7:eq-creation-in-live-1}
\sum_{\omega\in L^{(j)}_{\mathfrak{c}}}\tau_{\omega}\;\le\;q^{lv}_j\,Z,
\qquad q^{lv}_j:=C_{lv}\,e^{-c_{lv}\,p_0(g_j)},
\end{equation}
for every level-$j$ cell $\mathfrak{c}$, then the per-block weight bound of
Definition~\ref{8.7:def-creation-live-events} holds at level $j$, with
\begin{equation}\label{8.7:eq-creation-in-live-2}
C_q:=n_{\square}\,C_{lv},\qquad c_q:=c_{lv}.
\end{equation}
\end{lemma}

\begin{proof}
Let $\omega\in E^{\kappa}_{j,\square}$, so that the step-$j$ decomposition of $\omega$ chose the factor
$1-\chi^{\kappa}_{j,\square}$; a large-field region is thereby created at step $j$ at the cube $\square$. By
clause~(ii) of the classification of large-field regions stated in \cite[p.~175]{B-CMP122a}, a region is never, at
the step at which it is created, in the class of regions that no longer carry their large-field factor; a region
created at step $j$ is therefore live at level $j$. Thus the region created at $(j,\square)$ is part of a
component that is live at level $j$ and contains the level-$j$ cells of $\square$. This holds for every kind
$\kappa$. Therefore there is a level-$j$ cell $\mathfrak{c}$ meeting $\square$ that lies in a live component at
level $j$, that is, $\omega\in L^{(j)}_{\mathfrak{c}}$. This proves \eqref{8.7:eq-creation-in-live-a}. The cells
meeting $\square$ are the $n_{\square}$ level-$j$ cells of $\square$, and a single cell if the $\kappa$-cube is
smaller than a cell, as it is when $LM_2\le M$.

Now assume \eqref{8.7:eq-creation-in-live-1}. Every mass $\tau_{\omega}$ is positive. Summing $\tau_{\omega}$ over
the left-hand side of \eqref{8.7:eq-creation-in-live-a} is therefore bounded by summing over the union on the
right, and this in turn by the sum over the $n_{\square}$ cells of the sums over the single events:
\begin{align*}
\sum_{\omega\in E^{\kappa}_{j,\square}}\tau_{\omega}
&\le\sum_{\substack{\mathfrak{c}\ \text{a level-}j\ \text{cell}\\ \mathfrak{c}\cap\square\neq\emptyset}}
\ \sum_{\omega\in L^{(j)}_{\mathfrak{c}}}\tau_{\omega}
\;\le\; n_{\square}\,q^{lv}_j\,Z\\
&= n_{\square}\,C_{lv}\,e^{-c_{lv}\,p_0(g_j)}\,Z .
\end{align*}
With the constants \eqref{8.7:eq-creation-in-live-2} the right-hand side is
$C_q\,e^{-c_q\,p_0(g_j)}\,Z=q_j\,Z$, which is the per-block weight bound at level $j$.
\end{proof}

The inclusion \eqref{8.7:eq-creation-in-live-a} and the implication hold level by level. The live-cell one-point
bound is established below, in the theorem on the live-cell one-point bound, for every $K$ and every
$j\le K-u_{\gamma}(g)$, with $C_{lv}=C_L$ and $c_{lv}=\theta$; by the lemma the per-block weight bound then holds at
the same levels, with $C_q=n_{\square}C_L$ and $c_q=\theta$. The quantifiers over $K$, over admissible runs with
$K'\ge j$, over $N_T$, over positions and over kinds are unchanged; the only change is the restriction of the level
to $j\le K-u_{\gamma}(g)$.

\begin{lemma}[First-moment bound for an extensive mark]\label{8.7:lem-first-moment}
Let $\Omega$ be a finite set carrying a positive mass distribution $\tau$: masses $\tau_{\omega}\ge 0$ for
$\omega\in\Omega$, and $\tau(E):=\sum_{\omega\in E}\tau_{\omega}$ for $E\subset\Omega$. Put
$Z:=\tau(\Omega)$, assume $Z>0$, let $\tilde\nu:=\tau/Z$ be the normalised measure, and write
$\mathbb{E}_{\tilde\nu}$ for the expectation under $\tilde\nu$. Let $\Theta\colon\Omega\to[0,\infty[$,
and let $\mathfrak{M}$, $\mathfrak{m}$ be real numbers. Subsets of $\Omega$ are called events; the
infimum over the empty set is $+\infty$ and $e^{-\infty}=0$. Assume
\begin{equation}\label{8.7:eq-first-moment-a}
\begin{aligned}
&(\mathrm{G}\!\uparrow)\qquad
&&\tau(E)\le \exp\Bigl(-\inf_{E}\Theta\Bigr)\,\mathfrak{M}\qquad\text{for every } E\subset\Omega,\\
&(\mathrm{G}\!\downarrow)\qquad
&&Z\ge\mathfrak{m}>0.
\end{aligned}
\end{equation}
Then the following hold.
\begin{enumerate}
\item[(i)] With $c_J:=e^{1/2}/(1-e^{-1/2})$,
\begin{equation}\label{8.7:eq-first-moment-1}
\mathbb{E}_{\tilde\nu}\bigl[e^{\Theta/2}\bigr]\le c_J\,\frac{\mathfrak{M}}{\mathfrak{m}},
\qquad
\mathbb{E}_{\tilde\nu}[\Theta]\le 2\log\Bigl(c_J\,\frac{\mathfrak{M}}{\mathfrak{m}}\Bigr).
\end{equation}
If $(\mathrm{G}\!\uparrow)$ is available termwise, in the sense that there are numbers
$\mathfrak{M}_{\omega}$, $\omega\in\Omega$, with $\tau_{\omega}\le e^{-\Theta(\omega)}\mathfrak{M}_{\omega}$
for every $\omega$ and $\sum_{\omega}\mathfrak{M}_{\omega}\le\mathfrak{M}$, then
\begin{equation}\label{8.7:eq-first-moment-2}
\mathbb{E}_{\tilde\nu}\bigl[e^{\Theta}\bigr]\le\frac{\mathfrak{M}}{\mathfrak{m}},
\qquad
\mathbb{E}_{\tilde\nu}[\Theta]\le\log\frac{\mathfrak{M}}{\mathfrak{m}}.
\end{equation}
\item[(ii)] Let $(A_c)_{c\in\mathcal{C}}$ be a family of events indexed by a nonempty finite set
$\mathcal{C}$, and let $t>0$ be such that $\sum_{c\in\mathcal{C}}\mathbf{1}_{A_c}(\omega)\le\Theta(\omega)/t$
for every $\omega\in\Omega$. Then
\begin{equation}\label{8.7:eq-first-moment-3}
\frac{1}{|\mathcal{C}|}\sum_{c\in\mathcal{C}}\tilde\nu(A_c)
\le\frac{2\log(c_J\mathfrak{M}/\mathfrak{m})}{t\,|\mathcal{C}|}.
\end{equation}
If moreover $B_c\supset A_c$, $c\in\mathcal{C}$, are events such that $\tilde\nu(B_c)$ does not depend on
$c$, then for every $c\in\mathcal{C}$
\begin{equation}\label{8.7:eq-first-moment-4}
\tilde\nu(B_c)\le\max_{c'\in\mathcal{C}}\tilde\nu\bigl(B_{c'}\setminus A_{c'}\bigr)
+\frac{2\log(c_J\mathfrak{M}/\mathfrak{m})}{t\,|\mathcal{C}|}.
\end{equation}
\item[(iii)] (Sharpness.) For all real $V,\mathsf{B},t>0$ and every integer $n\ge\mathsf{B}V/t$ there are
a finite set $\mathcal{C}$ with $|\mathcal{C}|=n$, data $(\Omega,\tau,\Theta)$ and events
$(A_c)_{c\in\mathcal{C}}$, invariant under a group of permutations of $\mathcal{C}$ acting transitively on
$\mathcal{C}$ (the group acts on $\Omega$, preserves $\tau$ and $\Theta$, and every $\sigma$ in it carries
$A_c$ onto $A_{\sigma(c)}$), such that $(\mathrm{G}\!\uparrow)$ and $(\mathrm{G}\!\downarrow)$ hold with
$\mathfrak{M}/\mathfrak{m}=e^{\mathsf{B}V}$, $\sum_{c}\mathbf{1}_{A_c}=\Theta/t$, and
$\tilde\nu(A_c)=\lfloor\mathsf{B}V/t\rfloor/n$ for every $c\in\mathcal{C}$. Thus from positivity of
$\tau$, $(\mathrm{G}\!\uparrow)$, $(\mathrm{G}\!\downarrow)$, the reserve $t$ per marked index (an index $c$
with $\omega\in A_c$) and symmetry, nothing better than \textup{(ii)} follows, up to the factor $2$ and
the term $\log c_J$.
\end{enumerate}
\end{lemma}

\begin{proof}
\emph{Preliminary.} Since $Z>0$, the set $\Omega$ is nonempty, and $(\mathrm{G}\!\uparrow)$ applied to
$E=\Omega$ gives $0<Z\le e^{-\inf_{\Omega}\Theta}\mathfrak{M}$; hence $\mathfrak{M}>0$. For $E=\emptyset$
both sides of $(\mathrm{G}\!\uparrow)$ vanish by the convention on the empty infimum.

\emph{(i).} Slice $\Omega$ according to the values of $\Theta$: for $l\in\{0,1,2,\dots\}$ put
\begin{equation*}
E_l:=\{\omega\in\Omega:\ l\le\Theta(\omega)<l+1\}.
\end{equation*}
Since $\Theta\ge0$, the sets $E_l$ are pairwise disjoint and their union is $\Omega$; since $\Omega$ is
finite, only finitely many of them are nonempty. On $E_l$ we have $\inf_{E_l}\Theta\ge l$, so
$(\mathrm{G}\!\uparrow)$ gives $\tau(E_l)\le e^{-l}\mathfrak{M}$; for empty $E_l$ this holds because
$\tau(E_l)=0$ and $\mathfrak{M}>0$. On $E_l$ also $e^{\Theta/2}\le e^{(l+1)/2}$. Therefore
\begin{equation*}
\begin{aligned}
\sum_{\omega\in\Omega}\tau_{\omega}e^{\Theta(\omega)/2}
&=\sum_{l\ge0}\ \sum_{\omega\in E_l}\tau_{\omega}e^{\Theta(\omega)/2}
\le\sum_{l\ge0}e^{(l+1)/2}\,\tau(E_l)
\le\sum_{l\ge0}e^{(l+1)/2}e^{-l}\,\mathfrak{M}\\
&=e^{1/2}\sum_{l\ge0}e^{-l/2}\,\mathfrak{M}
=\frac{e^{1/2}}{1-e^{-1/2}}\,\mathfrak{M}=c_J\,\mathfrak{M}.
\end{aligned}
\end{equation*}
Dividing by $Z$ and using $Z\ge\mathfrak{m}$ from $(\mathrm{G}\!\downarrow)$ gives the first inequality
of \eqref{8.7:eq-first-moment-1}. Since the exponential is convex, Jensen's inequality under the
probability measure $\tilde\nu$ gives
\begin{equation*}
e^{\mathbb{E}_{\tilde\nu}[\Theta]/2}\le\mathbb{E}_{\tilde\nu}\bigl[e^{\Theta/2}\bigr]
\le c_J\,\frac{\mathfrak{M}}{\mathfrak{m}},
\end{equation*}
and taking logarithms and multiplying by $2$ gives the second inequality of
\eqref{8.7:eq-first-moment-1}.

For the termwise variant, the hypothesis gives $\tau_{\omega}e^{\Theta(\omega)}\le\mathfrak{M}_{\omega}$ for
every $\omega$, so
\begin{equation*}
\sum_{\omega\in\Omega}\tau_{\omega}e^{\Theta(\omega)}\le\sum_{\omega\in\Omega}\mathfrak{M}_{\omega}
\le\mathfrak{M};
\end{equation*}
dividing by $Z\ge\mathfrak{m}$ gives $\mathbb{E}_{\tilde\nu}[e^{\Theta}]\le\mathfrak{M}/\mathfrak{m}$, and
Jensen's inequality, $e^{\mathbb{E}_{\tilde\nu}[\Theta]}\le\mathbb{E}_{\tilde\nu}[e^{\Theta}]$, gives the
second inequality of \eqref{8.7:eq-first-moment-2} after taking logarithms.

\emph{(ii).} By linearity of the expectation, the pointwise hypothesis
$\sum_{c}\mathbf{1}_{A_c}\le\Theta/t$ and its monotonicity, and then \eqref{8.7:eq-first-moment-1},
\begin{equation*}
\sum_{c\in\mathcal{C}}\tilde\nu(A_c)
=\mathbb{E}_{\tilde\nu}\Bigl[\sum_{c\in\mathcal{C}}\mathbf{1}_{A_c}\Bigr]
\le\frac{\mathbb{E}_{\tilde\nu}[\Theta]}{t}
\le\frac{2\log(c_J\mathfrak{M}/\mathfrak{m})}{t};
\end{equation*}
dividing by $|\mathcal{C}|$ gives \eqref{8.7:eq-first-moment-3}. For the second claim fix
$c\in\mathcal{C}$. Since $\tilde\nu(B_{c'})$ does not depend on $c'$, it equals its average over
$c'\in\mathcal{C}$; since $A_{c'}\subset B_{c'}$, the set $B_{c'}$ is the disjoint union of
$B_{c'}\setminus A_{c'}$ and $A_{c'}$. Hence
\begin{equation*}
\begin{aligned}
\tilde\nu(B_c)
&=\frac{1}{|\mathcal{C}|}\sum_{c'\in\mathcal{C}}\tilde\nu(B_{c'})
=\frac{1}{|\mathcal{C}|}\sum_{c'\in\mathcal{C}}
\Bigl[\tilde\nu\bigl(B_{c'}\setminus A_{c'}\bigr)+\tilde\nu(A_{c'})\Bigr]\\
&\le\max_{c'\in\mathcal{C}}\tilde\nu\bigl(B_{c'}\setminus A_{c'}\bigr)
+\frac{1}{|\mathcal{C}|}\sum_{c'\in\mathcal{C}}\tilde\nu(A_{c'}),
\end{aligned}
\end{equation*}
and \eqref{8.7:eq-first-moment-3} bounds the last term, which gives \eqref{8.7:eq-first-moment-4}.

\emph{(iii).} Let $V,\mathsf{B},t>0$ and let $n\ge\mathsf{B}V/t$ be an integer. Put
$\nu:=\lfloor\mathsf{B}V/t\rfloor$, so that $0\le\nu\le\mathsf{B}V/t\le n$ and $t\nu\le\mathsf{B}V$. Let
$\mathcal{C}:=\mathbb{Z}/n$, let $\Omega$ be the set of subsets of $\mathcal{C}$ of cardinality $\nu$
(nonempty, because $\nu\le n$), and put
\begin{equation*}
\tau_{\omega}:=1,\qquad \Theta(\omega):=t\nu,\qquad A_c:=\{\omega\in\Omega:\ c\in\omega\},
\qquad \mathfrak{m}:=Z,\qquad \mathfrak{M}:=e^{\mathsf{B}V}Z .
\end{equation*}
The group of translations of $\mathbb{Z}/n$ acts transitively on $\mathcal{C}$ and acts on $\Omega$ by
translating subsets; it preserves cardinality, hence maps $\Omega$ to itself, and it preserves $\tau\equiv1$
and the constant $\Theta$. Translation by $a$ carries $A_c$ onto $A_{c+a}$, since
$c\in\omega$ if and only if $c+a\in\omega+a$.

$(\mathrm{G}\!\downarrow)$ holds because $\mathfrak{m}=Z=|\Omega|>0$, and
$\mathfrak{M}/\mathfrak{m}=e^{\mathsf{B}V}$. For $(\mathrm{G}\!\uparrow)$, let $E\subset\Omega$ be nonempty
(the empty set was treated above); then $\inf_E\Theta=t\nu$, and since $t\nu\le\mathsf{B}V$ we have
$\mathfrak{M}=e^{\mathsf{B}V}Z\ge e^{t\nu}Z$, so
\begin{equation*}
\tau(E)\le Z=e^{-t\nu}\cdot e^{t\nu}Z\le e^{-\inf_E\Theta}\,\mathfrak{M}.
\end{equation*}
For every $\omega\in\Omega$, $\sum_{c}\mathbf{1}_{A_c}(\omega)=|\omega|=\nu=\Theta(\omega)/t$. Finally,
translation by $a$ is a bijection of $A_c$ onto $A_{c+a}$ and $\tau\equiv1$, so $\tilde\nu(A_c)$ does not
depend on $c$; and by linearity of the expectation
\begin{equation*}
\sum_{c\in\mathcal{C}}\tilde\nu(A_c)
=\mathbb{E}_{\tilde\nu}\Bigl[\sum_{c\in\mathcal{C}}\mathbf{1}_{A_c}\Bigr]=\nu,
\end{equation*}
so $\tilde\nu(A_c)=\nu/n=\lfloor\mathsf{B}V/t\rfloor/n$ for every $c$. For these data the right-hand side
of \eqref{8.7:eq-first-moment-3} equals
\begin{equation*}
\frac{2\log\bigl(c_J e^{\mathsf{B}V}\bigr)}{t\,n}=\frac{2\mathsf{B}V+2\log c_J}{t\,n},
\end{equation*}
while the left-hand side equals $\lfloor\mathsf{B}V/t\rfloor/n$; this is the comparison stated at the end
of \textup{(iii)}.
\end{proof}

\begin{remark}[Density bounds against the comparison of one event]\label{8.7:rem-density-versus-event}
The lemma on the conversion of absolute per-object large-field factors into probabilities of label
events lists, in its clause~(iii), three ways to turn an absolute per-object factor
into a probability. Here an absolute factor is a bound of the type $(\mathrm{G}\!\uparrow)$ of
\eqref{8.7:eq-first-moment-a} on the unnormalised mass $\tau(E)$ of an event, not yet divided by
$Z$. The three ways listed there are:
\begin{itemize}
\item[(R)] removal of the object;
\item[(G)] global comparison of the mass of the event with the total mass, through constants whose
logarithm is a volume, for instance $(E_+^{\sharp}+C_{\mathrm{low}})N^d$, with $N$ the side of the
torus in terminal-lattice units;
\item[(S)] symmetrisation.
\end{itemize}
The first-moment bound for an extensive mark, Lemma~\ref{8.7:lem-first-moment}, is a fourth device:
\begin{itemize}
\item[(M)] global comparison of the extensive mark $\Theta$, then Jensen's inequality, then division by
the number of objects.
\end{itemize}

To compare (G) and (M), let $\mathsf{B}$ be a free-energy constant per unit of a volume $V$, and
suppose that the constants in $(\mathrm{G}\!\uparrow)$ and $(\mathrm{G}\!\downarrow)$ of
\eqref{8.7:eq-first-moment-a} satisfy $\mathfrak{M}/\mathfrak{m}=e^{\mathsf{B}V}$. Let $t>0$.

Let $E\subset\Omega$ be an event on which $\Theta\ge t$. Then $\inf_E\Theta\ge t$, and
$(\mathrm{G}\!\uparrow)$ together with $(\mathrm{G}\!\downarrow)$ gives
\begin{equation}\label{8.7:eq-density-versus-event-1}
\tilde\nu(E)=\frac{\tau(E)}{Z}\le e^{-t}\,\frac{\mathfrak{M}}{\mathfrak{m}}=e^{\mathsf{B}V-t}.
\end{equation}
This bound is useless when no bound on $V$ is available, since its right side exceeds $1$ as soon as
$\mathsf{B}V\ge t$.

Now let $(A_c)_{c\in\mathcal{C}}$ satisfy the hypothesis of clause~(ii) of
Lemma~\ref{8.7:lem-first-moment}, that is,
$\sum_{c\in\mathcal{C}}\mathbf{1}_{A_c}\le\Theta/t$ pointwise. Taking the expectation of this
inequality under $\tilde\nu$ and applying clause~(i) of that lemma, we obtain
\begin{equation}\label{8.7:eq-density-versus-event-2}
\frac{1}{|\mathcal{C}|}\sum_{c\in\mathcal{C}}\tilde\nu(A_c)
\le\frac{\mathbb{E}_{\tilde\nu}[\Theta]}{t\,|\mathcal{C}|}
\le\frac{2\log\bigl(c_J e^{\mathsf{B}V}\bigr)}{t\,|\mathcal{C}|}
=\frac{2(\mathsf{B}V+\log c_J)}{t\,|\mathcal{C}|}.
\end{equation}
Thus (M) bounds the density of the marked objects. The bound \eqref{8.7:eq-density-versus-event-2}
is useful exactly when the marked objects
are many per unit of the volume in which $\mathsf{B}$ is counted.

In Ba{\l}aban's scheme $\mathsf{B}$ is counted per site of the terminal lattice, at level $K'$. All
free energy of the finer levels is matched exactly, numerator against denominator, by the
renormalised vacuum energies; this is ultraviolet stability, in the form stated in
\cite[Corollary~3]{B-CMP119}. Accordingly $V=N^d$, the number of terminal sites of the torus.

The level-$j$ cells are the cubes of side $M\check{R}_j$ of the level-$j$ lattice. They number
$L^{d(K'-j)}/(M\check{R}_j)^d$ per terminal site. Hence, when $\mathcal{C}$ is the family of all
level-$j$ cells,
\begin{equation*}
|\mathcal{C}|=\frac{N^d L^{d(K'-j)}}{(M\check{R}_j)^d}.
\end{equation*}
Substituting $V=N^d$ and this value of $|\mathcal{C}|$ into the right side of
\eqref{8.7:eq-density-versus-event-2} gives
\begin{equation}\label{8.7:eq-density-versus-event-3}
\frac{2}{t}\,\bigl(\mathsf{B}+N^{-d}\log c_J\bigr)\,(M\check{R}_j)^d\,L^{-d(K'-j)} .
\end{equation}
\end{remark}

\begin{definition}[The terminal-live reserve and its global bound]\label{8.7:def-terminal-reserve}
Fix $K$ and a run whose terminal level is $K'$, and put $N:=N_T L^{K-K'}$, the side of the
torus measured in units of the terminal level. For a label history $\omega\in\Omega$ let
$|\mathcal{O}_{K'}(\omega)|$ be the union of the large-field components of the history
$\omega$ that are live at the terminal level $K'$, and let $\mathcal{Z}^{\mathrm{live}}(\omega)$
be the set of seeds of the first type in the internal histories of these components. For
$Z\in\mathcal{Z}^{\mathrm{live}}(\omega)$ let $j(Z)$ be the step at which $Z$ was created
and $d'_{j(Z)}(Z)$ the linear size of $Z$ measured in cubes of that step. The
\emph{terminal-live reserve} of $\omega$ is
\begin{equation}\label{8.7:eq-terminal-reserve-a}
  \Theta^{\mathrm{live}}(\omega)
  :=\sum_{Z\in\mathcal{Z}^{\mathrm{live}}(\omega)}
  \theta'\gamma_0A_1\,p_0(g_{j(Z)})^{2}\,\bigl(d'_{j(Z)}(Z)+1\bigr),
\end{equation}
where $p_0(\cdot)$ is the degree function.

Let $E_+^{\mathrm{lin}}$ be the upper free-energy constant $E_+^{\sharp}$ of the
correlation theorem, evaluated at the coefficients from which the per-seed reserve, the
summand of \eqref{8.7:eq-terminal-reserve-a}, has been removed. It does not depend on $K$,
on $N_T$ or on position, and it depends on $(g,m)$ only through $x_{K'}$. The
\emph{global bound} of the terminal-live reserve is the statement that, for every
$E\subset\Omega$,
\begin{equation}\label{8.7:eq-terminal-reserve-b}
  \sum_{\omega\in E}\tau_{\omega}
  \le\exp\Bigl(-\inf_{\omega\in E}\Theta^{\mathrm{live}}(\omega)\Bigr)\,
  e^{E_+^{\mathrm{lin}}N^4}.
\end{equation}
\end{definition}

Each summand of \eqref{8.7:eq-terminal-reserve-a} is linear in the size of the seed and
quadratic in the value $p_0(g_{j(Z)})$ of the degree function; it has the form of the member
which, in the notation of
\cite[(1.85)]{B-CMP122b}, reads $\gamma_0A_1p_0^2(g_{j+1})(d'_{j+1}(Z_i^{j+1})+1)$. In the
construction of \cite{B-CMP122b}, the large-field factors of each component combine
into one exponential factor, and in that factor the volumes of the large-field sets are
divided by the volume of a cube, so that they enter cube by cube.

\begin{proof}[Proof of \eqref{8.7:eq-terminal-reserve-b}]
For a history $\omega\in\Omega$ write $M_{\omega}$ for its majorant at the coefficients
from which the per-seed reserve has been removed. The reserve bound in its global form,
taken with the linear reserve \eqref{8.7:eq-terminal-reserve-a} per seed, states that for
every $E\subset\Omega$
\begin{equation*}
  \sum_{\omega\in E}\tau_{\omega}
  \le\exp\Bigl(-\inf_{\omega\in E}\Theta^{\mathrm{live}}(\omega)\Bigr)
  \sum_{\omega\in\Omega}M_{\omega}.
\end{equation*}
The upper free-energy bound of the correlation theorem, applied at the same reserve-stripped
coefficients, at which its constant $E_+^{\sharp}$ is $E_+^{\mathrm{lin}}$ by the definition
of $E_+^{\mathrm{lin}}$, states that
\begin{equation*}
  \sum_{\omega\in\Omega}M_{\omega}\le e^{E_+^{\mathrm{lin}}N^4}.
\end{equation*}
Inserting the second inequality into the right-hand side of the first gives
\eqref{8.7:eq-terminal-reserve-b} for every $E\subset\Omega$.
\end{proof}

In
\eqref{8.7:eq-terminal-reserve-b} the reserve enters only through
$\Theta^{\mathrm{live}}$, that is, summed over the seeds of the components that are live
at the terminal level $K'$, and the infimum is taken over the histories of $E$.

With the flat reserve (one reserve per component) in place of
\eqref{8.7:eq-terminal-reserve-a}, the same two bounds give the upper bound on the
numerator of the terminal member, the member of the expansion attached to the terminal
level $K'$; \eqref{8.7:eq-terminal-reserve-b} has the form of that upper bound, with the
flat reserve replaced by the linear
reserve \eqref{8.7:eq-terminal-reserve-a} (one reserve per seed, linear in its size).
Where it is used, it enters only as an absolute factor on majorants, compared with the
global lower bound $\tilde{Z}(0)\ge e^{-C_{\mathrm{low}}N^4}$ of the correlation theorem
for the normalising sum $\tilde{Z}(0)$ at zero source. It is not used relative to any
local quantity.

\begin{remark}[Healed seeds admit no cutoff-free reserve]\label{8.7:rem-healed-seeds}
The reserve functional $\Theta^{\mathrm{live}}$ of Definition~\ref{8.7:def-terminal-reserve} sums only over
the set $\mathcal{Z}^{\mathrm{live}}(\omega)$ of seeds of Definition~\ref{8.7:def-terminal-reserve}, attached to the components of
$|\mathcal{O}_{K'}(\omega)|$, that is, to the components live at the terminal level $K'$. No seed of a
healed large-field region enters it. For the bound by moduli that yields
$(\mathrm{G}\!\uparrow)$ this restriction cannot be lifted, as follows. Let $\Theta$ be a functional that
also counts the seeds of healed regions, each weighted by a mark $e^{t}$ with $t>0$, and call the
marked sum the sum over histories $\omega$ of the masses $\tau_{\omega}$ with the factor $e^{t}$
inserted for every seed counted by $\Theta$. Then bounding every
term of the marked sum by its modulus does not give the global upper bound
$(\mathrm{G}\!\uparrow)$ of Lemma~\ref{8.7:lem-first-moment}, displayed in
\eqref{8.7:eq-first-moment-a}, with
\begin{equation*}
\mathfrak{M}=\exp\bigl(E_+\cdot(\text{number of terminal sites})\bigr)
\end{equation*}
and an upper free-energy constant $E_+$ independent of $K$.
The reserve of a healed seed is carried instead on the kernel of each $\mathbf{R}$-term.
\end{remark}

\begin{proof}
For $X$ a large-field region created at step $j$, write $R^{(j)}(X,U)$ for the exponentiated
$\mathbf{R}$-term localised on $X$, and $R^{(j)}(X,1)$ for its value at the trivial configuration
(its vacuum part). Write $e^{(j)}(X)$ for the vacuum-energy term attached to $X$ in the torus bound of
the correlation theorem. The torus bound exhibits $R^{(j)}(X,1)$ as one of the contributions to
$e^{(j)}(X)$.

A mark $e^{t}$ placed on the seeds of healed regions multiplies the operators $\mathbf{T}'$ of these
regions. Hence it multiplies the exponentiated $\mathbf{R}$-terms built from them, and hence their vacuum
parts $R^{(j)}(X,1)$. Each such vacuum part is therefore changed by $(e^{t}-1)$ times itself.

On the other hand, the counterterm $E$ in the bare density $\rho_0=\exp[-A-E]$ of
Definition~\ref{1.5:def-densities} is that of the unmarked
theory, and the mark leaves it unchanged. The change of the vacuum parts is therefore not compensated.
Let $n_j$ be the number of level-$j$ cells of the torus; these are cubes of side $M\check{R}_j$ in
level-$j$ units, with $\check{R}_j=R(\min_{i\le j}x_i)$ the running radius. Let $\beta_0$ be the
parameter of Ba{\l}aban's large-field bounds that enters the exponent of the large-field factor. Then
the majorant of the marked sum acquires, for each level $j$, the factor
\begin{equation}\label{8.7:eq-healed-seeds-1}
\exp\Bigl[(e^{t}-1)\cdot O\bigl(e^{-2(1+\beta_0)^{-1}p_0(g_j)}\bigr)\cdot n_j\Bigr].
\end{equation}

The torus contains $L^{d(K'-j)}(M\check{R}_j)^{-d}$ level-$j$ cells per site of the terminal lattice.
Per terminal site, the exponent in \eqref{8.7:eq-healed-seeds-1} is therefore of the order of
\begin{equation}\label{8.7:eq-healed-seeds-2}
L^{d(K'-j)}\,(M\check{R}_j)^{-d}\,e^{\,t-2(1+\beta_0)^{-1}p_0(g_j)}.
\end{equation}

This quantity is unbounded in $K$ for every fixed $t>0$. To see this, let $K'-j$ grow. The factor
$L^{d(K'-j)}$ grows geometrically in $K'-j$, whereas the number $p_0(g_j)$, the degree function
$p_0(g')=A_0(\log g'^{-2})^{p_0}$ evaluated at the running coupling $g_j$, is polylogarithmic in
$K'-j$ by the reference window of Definition~\ref{1.2:def-flow-constants}. The difference $K'-j$
between the level at which a healed region is created and the terminal
level is not bounded independently of $K$. Compare \cite[p.~175]{B-CMP122a}, which states:
``it controls a large number of steps, but this number is a small fraction of the total number of
steps''.

The conclusion holds a fortiori for the mark $t=\theta p_0(g_j)$, with $\theta$ the reserve exponent.
For this mark the exponent $t$ in \eqref{8.7:eq-healed-seeds-2} itself grows with $p_0(g_j)$.

Consequently the majorant of the moduli of the marked sum exceeds every bound of the form
$\exp(E_+\cdot\text{number of terminal sites})$ with $E_+$ independent of $K$. The argument by moduli
therefore cannot give $(\mathrm{G}\!\uparrow)$ with a $K$-free $E_+$ for a functional $\Theta$ that
counts healed seeds. This is why \eqref{8.7:eq-terminal-reserve-a} and the global bound
\eqref{8.7:eq-terminal-reserve-b} are stated for $\Theta^{\mathrm{live}}$, read on terminal-live seeds
only.

Nothing is lost by this restriction. A healed seed carries its reserve on the kernel of each
$\mathbf{R}$-term. This is the quantity used by the summation of boundary terms at finer zones; it is
read on the normalised operator $\mathbf{T}'$.
\end{proof}

\begin{lemma}[Translation covariance of the label histories. Stated; proof incomplete at the step marked $(\ast)$]\label{8.7:lem-translation-covariance}
Let $k$ be a level, and let $a$ be a translation of the torus $\mathbb{T}_m$ by a multiple of the side of
the level-$k$ cells, the level-$k$ cells being the $M\check{R}_k$-cubes of the block lattice
$\mathbb{T}^{(k)}$ of Definition~\ref{1.5:def-block-average}. Consider the label histories of
Definition~\ref{8.7:def-creation-live-events} at zero complex source. For a label history $\omega$,
write $\omega_{\le k}$ for its part consisting of the steps up to level $k$, taken before the labels
of these steps are resummed; these are the level-$k$ un-resummed histories. For a level-$k$
un-resummed history $\omega_{\le k}$, let $a\cdot\omega_{\le k}$ be the history obtained by
translating by $a$ every region, every cube and every label of $\omega_{\le k}$. Write
$\tau_{\le k}(\omega_{\le k})$ for the mass that the steps up to level $k$ assign to
$\omega_{\le k}$. Then
$\omega_{\le k}\mapsto a\cdot\omega_{\le k}$ is a bijection of the set of level-$k$ un-resummed
histories onto itself, and for every such history
\begin{equation}\label{8.7:eq-translation-covariance-a}
\tau_{\le k}(a\cdot\omega_{\le k})=\tau_{\le k}(\omega_{\le k}).
\end{equation}
Consequently, with the live-cell events $\mathfrak{lv}_k(\mathfrak{c})$ and the normalised measure
$\tilde{\nu}$ of Definition~\ref{8.7:def-creation-live-events}, for every level-$k$ cell
$\mathfrak{c}$ and every such $a$,
\begin{equation}\label{8.7:eq-translation-covariance-1}
\tilde{\nu}\bigl(\mathfrak{lv}_k(\mathfrak{c}+a)\bigr)=\tilde{\nu}\bigl(\mathfrak{lv}_k(\mathfrak{c})\bigr),
\end{equation}
and the translations of $\mathbb{T}_m$ by multiples of the level-$k$ cell side act transitively on
the level-$k$ cells.

This lemma asserts nothing about events of levels larger than $k$ under translations by multiples
of the level-$k$ cell side. The event $\mathcal{E}_{\mathfrak{c}}$ that the cell $\mathfrak{c}$ is
still live at the terminal level (its never-healed fate) is such an event. The second clause of
part (ii) of the lemma on the global constants $\mathfrak{M}$, $\mathfrak{m}$ (the global upper
bound with reserve $(\mathrm{G}\!\uparrow)$ and the global floor $(\mathrm{G}\!\downarrow)$) is
designed for exactly this situation. There, the event $B_c=\mathfrak{lv}_k(\mathfrak{c})$ is
covariant by \eqref{8.7:eq-translation-covariance-1}, whereas the event $A_c=\mathcal{E}_{\mathfrak{c}}$
need not be.
\end{lemma}

\begin{proof}
\emph{The weight.} The translation $a$ moves the level-$k$ lattice by an integer number of its
spacings. One spacing of $\mathbb{T}^{(k)}$ is $L^{k}$ spacings of the bare lattice, so $a$ is a
translation of the bare lattice by an integer number of bare spacings, hence a lattice isometry.
Since the complex source is zero here, $z=0$, the weight is the Wilson weight, and by
Lemma~\ref{1.1:lem-invariance} the Wilson weight and the product Haar measure are invariant under it.

\emph{The partitions.} Consider first the cube partitions used by step $k$. In
$\mathbb{T}^{(k-1)}$ they are the partitions into $LM\check{R}_k$-cubes and into
$L^2M_2\check{R}_k$-cubes of \cite[\S3]{B-CMP119}. One spacing of $\mathbb{T}^{(k)}$ is $L$ spacings of
$\mathbb{T}^{(k-1)}$, so these are the partitions of $\mathbb{T}^{(k)}$ into $M\check{R}_k$-cubes
and into $LM_2\check{R}_k$-cubes, where the auxiliary parameters satisfy $LM_2\le M$.
For the steps before $k$, the cube partitions are nested $L$-adically in those of step $k$: this
is the nesting condition, through its consequence that the cube partitions of Ba{\l}aban's
successive steps are nested $L$-adically. Hence every cube
partition used by the steps up to $k$ has, in units of $\mathbb{T}^{(k)}$, a side that is a power
of $L$ and divides $M\check{R}_k$. A translation by a
multiple of $M\check{R}_k$ level-$k$ spacings is therefore a translation by a multiple of the side
of each of these partitions. Consequently $a$ maps each of them onto itself.

$(\ast)$ \emph{The rules of a step.} Every rule of a step is stated relative to these cubes or to
the component. This covers:
\begin{itemize}
\item the block average and its contours;
\item the axial trees;
\item the thresholds;
\item the layers;
\item the class clause with its two alternatives;
\item the frame average in Ba{\l}aban's operation $\mathbf{R}$;
\item the modifications of the step made by the correlation theorem.
\end{itemize}

\emph{Conclusion.} The weight and the Haar measure are invariant under $a$, and $a$ maps every
partition used by the steps up to $k$ onto itself. Every rule of these steps refers only to those
partitions or to the component, so translating by $a$ all regions, cubes and labels of a level-$k$
un-resummed history produces a level-$k$ un-resummed history. For the same reason, the translated
history carries the same mass, which is \eqref{8.7:eq-translation-covariance-a}. The map by $-a$ is
also a translation by a multiple of the level-$k$ cell side, and it inverts
$\omega_{\le k}\mapsto a\cdot\omega_{\le k}$, which is therefore a bijection.

\emph{The consequences.} The event $\mathfrak{lv}_k(\mathfrak{c})$ says that $\mathfrak{c}$ lies in a
component live at level $k$; it depends on a history only through its level-$k$ un-resummed part.
By the definition of $\tilde{\nu}$, its $\tilde{\nu}$-probability is $Z^{-1}$ times the sum of the
masses $\tau_{\omega}$ of the label histories $\omega$ in it. Group these histories according to
their level-$k$ un-resummed part; the steps after $k$ being decompositions of unity, the sum of
$\tau_{\omega}$ over the histories $\omega$ with a given part $\omega_{\le k}$ is
$\tau_{\le k}(\omega_{\le k})$. Hence the $\tilde{\nu}$-probability of
$\mathfrak{lv}_k(\mathfrak{c})$ is $Z^{-1}$ times the sum of the masses $\tau_{\le k}$ of the
level-$k$ un-resummed histories in it. The bijection $\omega_{\le k}\mapsto a\cdot\omega_{\le k}$
translates components together with cells. It therefore carries the level-$k$ un-resummed histories
in $\mathfrak{lv}_k(\mathfrak{c})$ onto those in $\mathfrak{lv}_k(\mathfrak{c}+a)$ and preserves their
masses by \eqref{8.7:eq-translation-covariance-a}. This gives \eqref{8.7:eq-translation-covariance-1}.

For transitivity, recall that the level-$k$ cells are the cubes of the partition of
$\mathbb{T}^{(k)}$ into $M\check{R}_k$-cubes. Any two of them differ by a translation whose
components are multiples of $M\check{R}_k$ level-$k$ spacings, that is, by a translation of the
kind considered.
\end{proof}

\begin{definition}[Constants of the density bound and the depth function]\label{8.7:not-depth-function}
Fix $g$. We write
\[
K':=K-m_{\mathbb{T}}(g)
\]
for the terminal level of the longest admissible run, and $N:=N_T L^{m_{\mathbb{T}}(g)}$. At that run we put
\[
\mathsf{B}(g):=E_+^{\mathrm{lin}}+C_{\mathrm{low}},
\]
where $E_+^{\mathrm{lin}}$ is the linear upper free-energy constant of the correlation theorem and
$C_{\mathrm{low}}$ is the constant of its floor, the lower bound
\[
\int\rho^{0}_{K'}\ge e^{-C_{\mathrm{low}}N^{d}}
\]
on the integral of the level-$K'$ effective density at zero source.

The number $\mathsf{B}(g)$ is finite. It depends on $g$ only through the bound
\[
x_{K'}\le X+b_++3\lambda m_{\mathbb{T}}(g)+2c_2,
\]
and it does not depend on $K$, on $N_T$ or on the position. No further property of $\mathsf{B}(g)$ is
used, and in particular its form as a function of $g$ is not used.

We further put
\[
C_{\mathrm{cnt}}:=(C_G+2)^d5^d,\qquad t_j:=\theta'\gamma_0A_1\bigl(p_0(g_j)-1\bigr)^2 .
\]
Here $p_0(\cdot)$ is the degree function. We also put
\[
n_j:=\Bigl(\frac{N L^{K'-j}}{M\check{R}_j}\Bigr)^{d},
\]
with $\check{R}_j=R(\min_{i\le j}x_i)$; this quotient is an integer, and $n_j$ is the number of
level-$j$ cells of the torus. Finally
\[
\bar{x}(u):=X+b_++3\lambda u+2c_2 ,
\]
where $b_+$ and $\lambda$ are the constants of the reference bound for the running couplings, by which
$x_{K-u}\le\bar{x}(u)$ for every depth $u$.

With $c_J$ the constant of Lemma~\ref{8.7:lem-first-moment}, the first-moment error term is
\begin{equation}\label{8.7:eq-depth-function-a}
\mathfrak{e}_j:=\frac{2C_{\mathrm{cnt}}\bigl(\mathsf{B}(g)+\log c_J\bigr)\,(M\check{R}_j)^d\,
L^{-d(K'-j)}}{t_j}.
\end{equation}

The depth function and the least admissible depth are
\begin{equation}\label{8.7:eq-depth-function-b}
\begin{split}
\mathfrak{F}_{\gamma}(u)&:=d\bigl(u-m_{\mathbb{T}}(g)\bigr)\log L-\theta A_0\bigl(\log\bar{x}(u)\bigr)^{p_0}
-d\log\Bigl(ML\bigl(\log\bar{x}(u)\bigr)^{r_{\mathrm{pr}}}\Bigr)\\
&\qquad-\log\bigl(2C_{\mathrm{cnt}}(\mathsf{B}(g)+\log c_J)\bigr)
+\log\bigl(\theta'\gamma_0A_1(p_0(g)-1)^2\bigr),\\
u_{\gamma}(g)&:=\min\bigl\{u_0\in\mathbb{Z},\ u_0\ge m_{\mathbb{T}}(g):\ \mathfrak{F}_{\gamma}(u)\ge0
\text{ for all integers }u\ge u_0\bigr\}.
\end{split}
\end{equation}
Here $u$ and $u_0$ range over integers, depths below the cutoff; $p_0$ in the exponent is the fixed
exponent of the degree function, and $r_{\mathrm{pr}}$ is the degree of $R(\cdot)$.
\end{definition}

\begin{proof}[Existence of the minimum in \eqref{8.7:eq-depth-function-b}]
In \eqref{8.7:eq-depth-function-b} the first term of $\mathfrak{F}_{\gamma}(u)$ is linear in $u$, with slope $d\log L>0$.

The second and third terms depend on $u$ only through $\log\bar{x}(u)$. Expanding the third term,
\[
d\log\Bigl(ML\bigl(\log\bar{x}(u)\bigr)^{r_{\mathrm{pr}}}\Bigr)
=d\log(ML)+d\,r_{\mathrm{pr}}\log\log\bar{x}(u).
\]
The quantity $\bar{x}(u)$ is affine in $u$ with slope $3\lambda>0$. Hence both $(\log\bar{x}(u))^{p_0}$ and $\log\log\bar{x}(u)$ are polylogarithmic in $u$, and each divided by $u$ tends to $0$ as $u\to\infty$. The last two terms do not depend on $u$.

Therefore $\mathfrak{F}_{\gamma}(u)/u\to d\log L>0$, so $\mathfrak{F}_{\gamma}(u)\to+\infty$. Consequently there is $u_0\ge m_{\mathbb{T}}(g)$ with $\mathfrak{F}_{\gamma}(u)\ge0$ for all $u\ge u_0$. The set in \eqref{8.7:eq-depth-function-b} is thus non-empty, and it is bounded below by $m_{\mathbb{T}}(g)$.

Since $u$ and $u_0$ range over integers, this set is a non-empty set of integers bounded below, and
therefore it has a least element.
\end{proof}

\begin{theorem}[Rarity of young large-field blocks]\label{8.7:thm-young-block-rarity}
Let $\Omega$ be the set of label histories of Ba{\l}aban's construction on the torus, with masses
$\tau_{\omega}$, total mass $Z=\sum_{\omega\in\Omega}\tau_{\omega}$ and normalised measure
$\tilde{\nu}=\tau/Z$, and let $L^{(j)}_{\mathfrak{c}}$ and $E^{\kappa}_{j,\square}$ be the live-cell
events and the creation events of Definition~\ref{8.7:def-creation-live-events}. Let $g>0$ be at
most the one-sided ceilings on the coupling under which hypotheses \textup{(IV)}, \textup{(V)} and
\textup{(X)} below hold,
let $K$ be a number of renormalisation steps, let $N_T\ge1$, put $K'=K-m_{\mathbb{T}}(g)$ and
$N=N_TL^{m_{\mathbb{T}}(g)}$ (here $N$ denotes the torus side in level-$K'$ units, not the rank of
the gauge group), and let $j\le K'$ be a level. Let $(\Omega,\tau)$ be those of any admissible run
with terminal level $\ge j$; by \textup{(I)} and \textup{(II)} below the masses considered here do
not depend on this choice. The level-$j$ cells $\mathfrak{c}$ are the
cubes of side $M\check{R}_j$; their number $n_j$, the constants $C_{\mathrm{cnt}}$, $t_j$,
$\mathsf{B}(g)$ and the function $\bar{x}(u)$ are those of Notation~\ref{8.7:not-depth-function}.
Assume the following.
\begin{enumerate}
\item[(I)] Positivity and the tree lemma: $\tau_{\omega}\ge0$ for every $\omega\in\Omega$, and for
every level $k$ the $\tilde{\nu}$-mass of an event determined by the steps of levels $\le k$ is a
level-$k$ quantity, that is, it does not change when the run is continued beyond level $k$.
\item[(II)] Coincidence of admissible runs: an admissible run stops the construction at its
terminal level $K-m$, with $m_{\mathbb{T}}(g)\le m\le K$; two admissible runs coincide through the
last step of the shorter one; the longest admissible run has $m=m_{\mathbb{T}}(g)$ and terminal
level $K'$.
\item[(III)] The fate decomposition: for every level-$j$ cell $\mathfrak{c}$ the event
$L^{(j)}_{\mathfrak{c}}$ is the disjoint union of the events $H_k$, $j<k\le K'$, that the component
containing $\mathfrak{c}$ is healed at step $k$, and of the event
$\mathcal{E}_{\mathfrak{c}}:=\mathcal{E}^{(j)}_{\mathfrak{c}}(K')$ that its lineage, the sequence
of the components containing it at the levels after $j$, is still live at
level $K'$; and if a component present at level $j$ has a lineage ending in a component $C'$ of
$|\mathcal{O}_{K'}(\omega)|$, its internal history, in particular its creation regions, is part of
the internal history of $C'$.
\item[(IV)] The healed-later bound, which bounds the part of the weight of a live cell whose
component is healed after step $j$, under the one-sided conditions of that statement, among them
\textup{(X)(c)} and \textup{(X)(d)} below: in its form without conditioning, for every level-$j$
cell $\mathfrak{c}$,
\[
\tilde{\nu}\bigl(L^{(j)}_{\mathfrak{c}}\setminus\mathcal{E}_{\mathfrak{c}}\bigr)
\le K_He^{\theta}S_*\,e^{-\theta p_0(g_j)} .
\]
\item[(V)] The terminal-live reserve bound~\eqref{8.7:eq-terminal-reserve-b} of
Definition~\ref{8.7:def-terminal-reserve}, for the reserve functional
$\Theta^{\mathrm{live}}$ of~\eqref{8.7:eq-terminal-reserve-a}, in the form of hypothesis
$(\mathrm{G}\!\uparrow)$ of~\eqref{8.7:eq-first-moment-a} with $\Theta=\Theta^{\mathrm{live}}$ and
$\mathfrak{M}=e^{E_+^{\mathrm{lin}}N^4}$.
\item[(VI)] The floor of the correlation theorem: $Z\ge e^{-C_{\mathrm{low}}N^4}$, with
$C_{\mathrm{low}}$ depending on $x_{K'}$ only.
\item[(VII)] Blob covering: for every level $n$, a component present at level $n$ is covered by the
$n$-blobs of the cells of its creation regions, the $n$-blob of a cell being a cube of side at most
$C_GM\check{R}_n$ containing the image of that cell at level $n$; more precisely, it is covered by
the $n$-blobs of the cells of the creation regions of its seeds of the first kind (the seeds
entering $\mathcal{Z}^{\mathrm{live}}$ in Definition~\ref{8.7:def-terminal-reserve}), because a seed
of the second, preparatory kind lies inside an already existing component and adds layers to it
only through the layer-adding operation of the region rule, and these layers are covered by the
blobs of the cells of the component's creation regions of the first kind. This rests on the region
rule of Ba{\l}aban's large-field operation~\cite{B-CMP122b}.
\item[(VIII)] Cell count along trees and nesting: a tree of length $\ell$ cell-sides meets at most
$5^d(\ell+1)$ cells; the image of a level-$j$ cell at a level $n\ge j$ lies in one level-$n$ cell.
\item[(IX)] The inclusion~\eqref{8.7:eq-creation-in-live-a} of
Lemma~\ref{8.7:lem-creation-in-live}: $E^{\kappa}_{j,\square}$ is contained in the union of the
events $L^{(j)}_{\mathfrak{c}}$ over the $n_{\square}$ level-$j$ cells $\mathfrak{c}$ meeting
$\square$.
\item[(X)] Inequalities of the coupling flow: (a) $x_{K-u}\le\bar{x}(u)$ for every depth $u=K-j$
with $j$ as in part~(d) below; (b) $X=g^{-2}\le x_k$ for every $k$; (c) $p_0(g_i)\ge p_0(g_j)-1$
for every level $i\le j$; (d) $S_*\le 1+1/\lambda^{\sharp}$.
\item[(XI)] Translation covariance of the label histories at level $j$, as in
Lemma~\ref{8.7:lem-translation-covariance} (stated; proof incomplete).
\end{enumerate}
Then, under \textup{(I)}--\textup{(X)}:
\begin{enumerate}
\item[(a)] (without \textup{(XI)}) with the sum over all level-$j$ cells of the torus and
$\mathfrak{e}_j$ as in~\eqref{8.7:eq-depth-function-a},
\begin{equation}\label{8.7:eq-young-block-rarity-a}
\bar{q}_j:=n_j^{-1}\sum_{\mathfrak{c}}\tilde{\nu}\bigl(L^{(j)}_{\mathfrak{c}}\bigr)
\le K_He^{\theta}S_*\,e^{-\theta p_0(g_j)}+\mathfrak{e}_j ;
\end{equation}
\item[(b)] if \textup{(XI)} holds, then for every level-$j$ cell $\mathfrak{c}$
\begin{equation}\label{8.7:eq-young-block-rarity-1}
\tilde{\nu}\bigl(L^{(j)}_{\mathfrak{c}}\bigr)\le K_He^{\theta}S_*\,e^{-\theta p_0(g_j)}
+\mathfrak{e}_j ;
\end{equation}
\item[(c)] for every kind $\kappa$ and every $\kappa$-cube $\square$ of step $j$, the average of
$\tilde{\nu}(E^{\kappa}_{j,\square+a})$ over the $n_j$ translations $a$ of the torus by multiples
of the level-$j$ cell side is at most $n_{\square}$ times the right member
of~\eqref{8.7:eq-young-block-rarity-a}; if \textup{(XI)} holds, then
$\tilde{\nu}(E^{\kappa}_{j,\square})$ is at most $n_{\square}$ times the right member
of~\eqref{8.7:eq-young-block-rarity-1};
\item[(d)] if
\begin{equation}\label{8.7:eq-young-block-rarity-b}
K-j\ge u_{\gamma}(g),
\end{equation}
with $u_{\gamma}(g)$ as in~\eqref{8.7:eq-depth-function-b}, then
$\mathfrak{e}_j\le e^{-\theta p_0(g_j)}$; hence, if moreover \textup{(XI)} holds, for every
admissible run with terminal level
$\ge j$, the live-cell one-point bound of Definition~\ref{8.7:def-creation-live-events} holds at
level $j$ with $C_{lv}=C_L:=K_He^{\theta}(1+1/\lambda^{\sharp})+1$ and $c_{lv}=\theta$, and the
per-block creation bound of Definition~\ref{8.7:def-creation-live-events} holds at level $j$ with
$C_q=n_{\square}C_L$ and $c_q=\theta$ (both in the form~\eqref{8.7:eq-creation-live-events-a}),
uniformly in $K$, $m$, $N_T$, the position of the cell and of the cube, and the kind $\kappa$.
\end{enumerate}
\end{theorem}

\begin{proof}
\emph{The run.} The events $L^{(j)}_{\mathfrak{c}}$ and $E^{\kappa}_{j,\square}$ are determined by
the steps of levels $\le j$. By hypothesis~(I) their $\tilde{\nu}$-masses are level-$j$
quantities, and by hypothesis~(II) any two admissible runs with terminal level $\ge j$ coincide
through level $j$; hence these masses do not depend on where the run stops. The measure is fixed
and only the level at which its representation stops varies, so we evaluate all masses on the
longest admissible run, whose terminal level is $K'=K-m_{\mathbb{T}}(g)\ge j$. On that run the torus
has side $N=N_TL^{m_{\mathbb{T}}(g)}$ in level-$K'$ units.

\emph{Fates.} Fix a level-$j$ cell $\mathfrak{c}$. By hypothesis~(III), grouping the healed fates
$H_k$, $j<k\le K'$, into one event,
\begin{equation}\label{8.7:eq-young-block-rarity-2}
L^{(j)}_{\mathfrak{c}}=\bigl(L^{(j)}_{\mathfrak{c}}\setminus\mathcal{E}_{\mathfrak{c}}\bigr)
\sqcup\mathcal{E}_{\mathfrak{c}},
\qquad \mathcal{E}_{\mathfrak{c}}=\mathcal{E}^{(j)}_{\mathfrak{c}}(K').
\end{equation}

\emph{Healed later.} By hypothesis~(IV), for every level-$j$ cell $\mathfrak{c}$,
\begin{equation}\label{8.7:eq-young-block-rarity-3}
\tilde{\nu}\bigl(L^{(j)}_{\mathfrak{c}}\setminus\mathcal{E}_{\mathfrak{c}}\bigr)
\le K_He^{\theta}S_*\,e^{-\theta p_0(g_j)} .
\end{equation}

\emph{The count.} For $\omega\in\Omega$ let $N^{\mathrm{live}}_j(\omega)$ be the number of level-$j$
cells $\mathfrak{c}$ with $\omega\in\mathcal{E}_{\mathfrak{c}}$, and for a region $Z$ let
$|Z|_{\mathrm{cells}}$ denote the number of its cells. Let $\omega\in\mathcal{E}_{\mathfrak{c}}$.
Since $\mathcal{E}_{\mathfrak{c}}\subset L^{(j)}_{\mathfrak{c}}$, the cell $\mathfrak{c}$ lies in a
component $C$ present at level $j$, and by the definition of $\mathcal{E}_{\mathfrak{c}}$ the
lineage of $C$ ends in a component $C'$ of $|\mathcal{O}_{K'}(\omega)|$. By the last assertion of
hypothesis~(III), the internal history of $C$, in particular its creation regions, all of birth
level $\le j$, is part of the internal history of $C'$. By hypothesis~(VII) at $n=j$, in the form
that covers $C$ by the blobs of the cells of the creation regions of its seeds of the first kind,
the cell $\mathfrak{c}$ lies in the $j$-blob of a cell $\square''$ of such a creation region; since
this region belongs to the internal history of $C'$, it is a region
$Z\in\mathcal{Z}^{\mathrm{live}}(\omega)$ with $j(Z)\le j$. The layers added to $C$ by seeds of the
preparatory kind need no separate treatment, since by the same hypothesis they are covered by these
blobs. The $j$-blob of $\square''$ is a cube of side at most
$C_GM\check{R}_j$, and a level-$j$ cell has side $M\check{R}_j$; hence the blob meets at most
$(C_G+2)^d$ level-$j$ cells. Every cell counted in $N^{\mathrm{live}}_j(\omega)$ is among the cells
met by one of these blobs; summing over the cells $\square''$ of the regions $Z$ we obtain
\begin{equation}\label{8.7:eq-young-block-rarity-4}
N^{\mathrm{live}}_j(\omega)
\le (C_G+2)^d\sum_{\substack{Z\in\mathcal{Z}^{\mathrm{live}}(\omega)\\ j(Z)\le j}}
|Z|_{\mathrm{cells}}
\le (C_G+2)^d5^d\sum_{\substack{Z\in\mathcal{Z}^{\mathrm{live}}(\omega)\\ j(Z)\le j}}
\bigl(d'_{j(Z)}(Z)+1\bigr).
\end{equation}
The second inequality is hypothesis~(VIII) at cell scale: the shortest tree through the cells of
$Z$ has length $d'_{j(Z)}(Z)$ cell-sides, so it meets, and $Z$ has, at most
$5^d(d'_{j(Z)}(Z)+1)$ cells. For $j(Z)\le j$, hypothesis~(X)(c) gives
$p_0(g_{j(Z)})\ge p_0(g_j)-1$, hence
$\theta'\gamma_0A_1p_0(g_{j(Z)})^2\ge\theta'\gamma_0A_1(p_0(g_j)-1)^2=t_j$, and each summand of
$\Theta^{\mathrm{live}}(\omega)$ in~\eqref{8.7:eq-terminal-reserve-a} belonging to such a $Z$ is at
least $t_j(d'_{j(Z)}(Z)+1)$. The summands belonging to the regions with $j(Z)>j$ are non-negative.
With $C_{\mathrm{cnt}}=(C_G+2)^d5^d$, \eqref{8.7:eq-young-block-rarity-4} therefore gives
\begin{equation}\label{8.7:eq-young-block-rarity-5}
N^{\mathrm{live}}_j(\omega)\le C_{\mathrm{cnt}}\,\Theta^{\mathrm{live}}(\omega)/t_j
\qquad(\omega\in\Omega).
\end{equation}

\emph{The first moment.} We apply Lemma~\ref{8.7:lem-first-moment} to $(\Omega,\tau)$ with
$\Theta:=\Theta^{\mathrm{live}}$. Hypothesis $(\mathrm{G}\!\uparrow)$
of~\eqref{8.7:eq-first-moment-a} is hypothesis~(V), with $\mathfrak{M}:=e^{E_+^{\mathrm{lin}}N^4}$;
hypothesis $(\mathrm{G}\!\downarrow)$ is hypothesis~(VI), with
$\mathfrak{m}:=e^{-C_{\mathrm{low}}N^4}$. By part~(i) of that lemma, with
$\mathsf{B}=E_+^{\mathrm{lin}}+C_{\mathrm{low}}$ at the longest run,
\begin{equation}\label{8.7:eq-young-block-rarity-6}
\mathbb{E}_{\tilde\nu}\bigl[\Theta^{\mathrm{live}}\bigr]
\le 2\log\bigl(c_J\mathfrak{M}/\mathfrak{m}\bigr)
=2\bigl(\mathsf{B}N^4+\log c_J\bigr).
\end{equation}
We use part~(ii) with $\mathcal{C}$ the set of the $n_j$ level-$j$ cells,
$A_{\mathfrak{c}}:=\mathcal{E}_{\mathfrak{c}}$ and $t:=t_j/C_{\mathrm{cnt}}$; its hypothesis is
\eqref{8.7:eq-young-block-rarity-5}, since
$\sum_{\mathfrak{c}}\mathbf{1}_{\mathcal{E}_{\mathfrak{c}}}(\omega)=N^{\mathrm{live}}_j(\omega)$.
By linearity of $\mathbb{E}_{\tilde\nu}$,
$\sum_{\mathfrak{c}}\tilde{\nu}(\mathcal{E}_{\mathfrak{c}})=\mathbb{E}_{\tilde\nu}[N^{\mathrm{live}}_j]
\le\mathbb{E}_{\tilde\nu}[\Theta^{\mathrm{live}}]/t$, and with~\eqref{8.7:eq-young-block-rarity-6}
\begin{equation}\label{8.7:eq-young-block-rarity-7}
n_j^{-1}\sum_{\mathfrak{c}}\tilde{\nu}(\mathcal{E}_{\mathfrak{c}})
\le\frac{2C_{\mathrm{cnt}}(\mathsf{B}N^4+\log c_J)}{t_jn_j}
\le\frac{2C_{\mathrm{cnt}}(\mathsf{B}+\log c_J)N^4}{t_jn_j}=\mathfrak{e}_j .
\end{equation}
The middle inequality holds because $N\ge1$ and $\log c_J>0$. For the last equality, $d=4$ gives
\[
n_j=\Bigl(\frac{NL^{K'-j}}{M\check{R}_j}\Bigr)^{d}=\frac{N^4L^{d(K'-j)}}{(M\check{R}_j)^d},
\]
so the factor $N^4$ cancels, and with $K'=K-m_{\mathbb{T}}(g)$ the last member
of~\eqref{8.7:eq-young-block-rarity-7} is
\[
2C_{\mathrm{cnt}}\bigl(\mathsf{B}(g)+\log c_J\bigr)\,(M\check{R}_j)^d\,
L^{-d(K-m_{\mathbb{T}}(g)-j)}/t_j=\mathfrak{e}_j ,
\]
the quantity $\mathfrak{e}_j$ of~\eqref{8.7:eq-depth-function-a}; in particular it does not depend
on $N_T$.
By~\eqref{8.7:eq-young-block-rarity-2}, $\tilde{\nu}(L^{(j)}_{\mathfrak{c}})$ is the sum of the two
masses on its right; averaging over $\mathfrak{c}$ and using~\eqref{8.7:eq-young-block-rarity-3}
and~\eqref{8.7:eq-young-block-rarity-7} gives~\eqref{8.7:eq-young-block-rarity-a}. This proves~(a).

\emph{Position.} Assume hypothesis~(XI). The translation by a multiple $a$ of the level-$j$ cell
side is a mass-preserving bijection of the histories which carries $L^{(j)}_{\mathfrak{c}}$ onto
$L^{(j)}_{\mathfrak{c}+a}$, so $\tilde{\nu}(L^{(j)}_{\mathfrak{c}})$ does not depend on the
level-$j$ cell $\mathfrak{c}$. The second clause of part~(ii) of
Lemma~\ref{8.7:lem-first-moment}, applied with $B_{\mathfrak{c}}:=L^{(j)}_{\mathfrak{c}}$, then
applies: each $\tilde{\nu}(L^{(j)}_{\mathfrak{c}})$ equals the average $\bar{q}_j$, and
\eqref{8.7:eq-young-block-rarity-1} follows from~\eqref{8.7:eq-young-block-rarity-a}. This
proves~(b).

\emph{Creation.} Let $\kappa$ be a kind and $\square$ a $\kappa$-cube of step $j$. By
hypothesis~(IX), $\tilde{\nu}(E^{\kappa}_{j,\square})$ is at most the sum of
$\tilde{\nu}(L^{(j)}_{\mathfrak{c}})$ over the $n_{\square}$ level-$j$ cells $\mathfrak{c}$ meeting
$\square$. Under~(XI) each term is bounded by~\eqref{8.7:eq-young-block-rarity-1}, which gives the
second assertion of~(c). For the first, the cells meeting $\square+a$ are the translates
$\mathfrak{c}'+a$ of the $n_{\square}$ cells $\mathfrak{c}'$ meeting $\square$, so
\[
n_j^{-1}\sum_a\tilde{\nu}\bigl(E^{\kappa}_{j,\square+a}\bigr)
\le\sum_{\mathfrak{c}'\cap\square\neq\emptyset}\,
n_j^{-1}\sum_a\tilde{\nu}\bigl(L^{(j)}_{\mathfrak{c}'+a}\bigr)
=n_{\square}\,\bar{q}_j ,
\]
because for fixed $\mathfrak{c}'$ the cell $\mathfrak{c}'+a$ runs once through every level-$j$ cell
as $a$ runs through the $n_j$ translations. By~\eqref{8.7:eq-young-block-rarity-a} this proves~(c).

\emph{The floor.} Assume~\eqref{8.7:eq-young-block-rarity-b} and put $u:=K-j\ge u_{\gamma}(g)$.
Hypothesis~(X)(a) gives $x_j=x_{K-u}\le\bar{x}(u)$. With $\check{R}_j=R(\min_{i\le j}x_i)$ this
yields $\check{R}_j\le R(x_j)\le L(\log\bar{x}(u))^{r_{\mathrm{pr}}}$, and with the degree function
$p_0(g_j)=A_0(\log x_j)^{p_0}\le A_0(\log\bar{x}(u))^{p_0}$. By hypothesis~(X)(b),
$p_0(g)=A_0(\log X)^{p_0}\le p_0(g_j)$, whence $t_j\ge\theta'\gamma_0A_1(p_0(g)-1)^2$. Finally
$K'-j=u-m_{\mathbb{T}}(g)$. Inserting these bounds into~\eqref{8.7:eq-depth-function-a} gives
\[
\mathfrak{e}_je^{\theta p_0(g_j)}
\le\frac{2C_{\mathrm{cnt}}(\mathsf{B}(g)+\log c_J)\,
\bigl(ML(\log\bar{x}(u))^{r_{\mathrm{pr}}}\bigr)^d}
{\theta'\gamma_0A_1(p_0(g)-1)^2}\,
L^{-d(u-m_{\mathbb{T}}(g))}\,e^{\theta A_0(\log\bar{x}(u))^{p_0}},
\]
and the logarithm of the right member is at most $-\mathfrak{F}_{\gamma}(u)$, with
$\mathfrak{F}_{\gamma}$ the depth function of~\eqref{8.7:eq-depth-function-b}; since
$u\ge u_{\gamma}(g)$, the definition of $u_{\gamma}(g)$ there gives $-\mathfrak{F}_{\gamma}(u)\le0$.
Hence $\mathfrak{e}_j\le e^{-\theta p_0(g_j)}$. Assume now moreover~(XI).
With~\eqref{8.7:eq-young-block-rarity-1} and hypothesis~(X)(d),
\[
\tilde{\nu}\bigl(L^{(j)}_{\mathfrak{c}}\bigr)
\le\bigl(K_He^{\theta}(1+1/\lambda^{\sharp})+1\bigr)e^{-\theta p_0(g_j)}
=C_L\,e^{-\theta p_0(g_j)}
\]
for every level-$j$ cell $\mathfrak{c}$, and by~(c)
$\tilde{\nu}(E^{\kappa}_{j,\square})\le n_{\square}C_L\,e^{-\theta p_0(g_j)}$ for every kind and
every $\kappa$-cube of step $j$. Since $\tilde{\nu}=\tau/Z$, these are the two one-point bounds
of~\eqref{8.7:eq-creation-live-events-a} at level $j$ with the constants stated in~(d). By the first
part of the proof they hold on every admissible run with terminal level $\ge j$; the constants
$C_L$ and $\theta$ do not involve $K$, $m$, $N_T$, the position or the kind, and $C_q=n_{\square}C_L$
involves the kind only through the number $n_{\square}$ of level-$j$ cells of the $\kappa$-cube.
This proves~(d).
\end{proof}

\begin{remark}[Scope and constants of the rarity bound]\label{8.7:rem-rarity-scope}
We comment on the scope of the rarity of young large-field blocks
(Theorem~\ref{8.7:thm-young-block-rarity}) and on its constants. Here $B\subset\Omega$ denotes a
conditioning event on the set $\Omega$ of label histories; $B=\Omega$ means that no conditioning is
imposed.
\begin{enumerate}
\item[(r1)] The constants of part~(d) of Theorem~\ref{8.7:thm-young-block-rarity} are those of the
large-field block bound, healed part, in its one-point form. In particular, consider the additive $1$
in the constant of that bound,
\begin{equation*}
C_L=K_He^{\theta}\bigl(1+1/\lambda_{\sharp}\bigr)+1 .
\end{equation*}
In that bound this summand absorbed a contribution that is controlled there under a live-tail
hypothesis at the terminal level $K'$. Here the same summand absorbs
$\mathfrak{e}_j\,e^{\theta p_0(g_j)}$, where $\mathfrak{e}_j$ is the first-moment error term of part~(a).

Theorem~\ref{8.7:thm-young-block-rarity} is that one-point bound specialised to $B=\Omega$, with two
hypotheses replaced. The hypotheses removed are the live-tail hypothesis at level $K'$ and
the condition on $\kappa_1$ used there. In their place stand:
\begin{itemize}
\item the terminal-live reserve bound \eqref{8.7:eq-terminal-reserve-b};
\item the global floor on the total mass;
\item the side bound carrying the constant $C_G$ (the constant that enters the counting constant
$C_{\mathrm{cnt}}=(C_G+2)^d5^d$);
\item the hypothesis labelled (C1) among those of Theorem~\ref{8.7:thm-young-block-rarity};
\item for part~(b) only, translation covariance of the label histories
\eqref{8.7:eq-translation-covariance-a}.
\end{itemize}
These are taken together with the additional depth condition \eqref{8.7:eq-young-block-rarity-b}.
\item[(r2)] Every kind $\kappa$ of characteristic function $\chi^{\kappa}_{j,\square}$ is covered,
local or not. This includes the kind lettered $\mathrm{R}$ in this chapter. The reason is
that the proof of Theorem~\ref{8.7:thm-young-block-rarity} reads only the labels, that is, the
creation events $E^{\kappa}_{j,\square}$ and the fates of cells along a label history.
\item[(r3)] The proof uses no bound on lifetimes and no decay in the size of components. It uses no
entropy sum over birth levels or over shapes. Every other component is accounted for either in the
number of level-$j$ cells that are, under the history $\omega$, in the never-healed fate, that is,
still live at the terminal level $K'$, or in the proposition of this section on the healed fates
$H_k$, whose conditions are among the hypotheses of Theorem~\ref{8.7:thm-young-block-rarity}.
This covers large, old, refreshed, merged and adjacent components, and
components of the second class.
\item[(r4)] The proposition on the healed fates $H_k$ admits conditioning events $B\neq\Omega$ that
do not depend on the creation history of a cell.
The first-moment bound for an extensive mark (Lemma~\ref{8.7:lem-first-moment}) does not admit them.
Its lower-bound hypothesis, the global floor $Z\ge\mathfrak{m}>0$ on the total mass
$Z=\tau(\Omega)$, provides no lower
bound on $\tilde{\nu}(B)$. Theorem~\ref{8.7:thm-young-block-rarity} is therefore stated for
$B=\Omega$ only. This is the case in which it is applied elsewhere in this book.
\item[(r5)] The rate $c_q$ of the one-point bound may be taken to be any fixed number in
$(0,\theta]$. The least admissible depth $u_{\gamma}(g)$ scales with $c_q$. In the applications only
the $\gamma$-ceilings on the coupling, of the kind appearing among the hypotheses of
Theorem~\ref{8.7:thm-young-block-rarity}, move; this is possible because $p_0>r_{\mathrm{pr}}$, where
$p_0$ is the exponent of the degree function $p_0(\cdot)$ and $r_{\mathrm{pr}}$ the degree of
$R(\cdot)$. Nothing else changes.
\end{enumerate}
\end{remark}

\begin{corollary}[The young-member bound at the witness depth]\label{8.7:cor-young-member}
Assume part~(d) of Theorem~\ref{8.7:thm-young-block-rarity} (rarity of young large-field blocks),
in the range of $g$, $K$, $N_T$ and of levels $j\le K'=K-m_{\mathbb{T}}(g)$ of that theorem. Let
$s$ be the witness depth and $k_0=K-s$ the witness level. Let $\mathfrak{l}$ be the lifetime, the
number of levels below the witness level over which the young member of the additive assembly of the
two-point witness theorem collects creation events, and $\mathfrak{B}_1(j)$, $\mathfrak{B}_2(j)$ the
block families near the two balls that are read by that young member. Let $C_{\mathrm{ext}}$ be the
constant in the range of levels read by the Gaussian extraction, and $r(s)$ the extraction depth. Put
$C_L=K_He^{\theta}(1+1/\lambda^{\sharp})+1$. Then the following hold.
\begin{itemize}
\item[(a)] If $s\ge u_{\gamma}(g)$, then for every $j\in[k_0-\mathfrak{l},k_0]$ and every
$(\kappa,\square)\in\mathfrak{B}_1(j)\cup\mathfrak{B}_2(j)$ the per-block weight bound
of~\eqref{8.7:eq-creation-live-events-a} holds in the form
\begin{equation}\label{8.7:eq-young-member-1}
\sum_{\omega\in E^{\kappa}_{j,\square}}\tau_{\omega}\le q_j\,Z,
\qquad q_j=n_{\square}\,C_L\,e^{-\theta p_0(g_j)},
\end{equation}
that is, with $c_q=\theta$. These are the inputs read by the young member of the additive assembly in
its creation-event form. They enter it through the arithmetic of part~(b) of the proposition on the
additive assembly, which, together with the lower flow bound and the $\gamma$-ceiling used there,
applies verbatim with $c_q=\theta$.
\item[(b)] If $s\ge u_{\gamma}(g)$, then for every level-$k_0$ cell $\square$
\begin{equation}\label{8.7:eq-young-member-2}
\tilde{\nu}\big(\mathfrak{lv}_{k_0}(\square)\big)\le C_L\,e^{-\theta p_0(g_{k_0})},
\end{equation}
that is, the live-cell bound of~\eqref{8.7:eq-creation-live-events-a} at level $k_0$ with
$(C_{lv},c_{lv})=(C_L,\theta)$. This is the input read by the additive assembly in its live-cell form.
\item[(c)] If $s-r(s)\ge u_{\gamma}(g)$, then for every level $j\le k_0+r(s)$ and every level-$j$ cell
$\square$
\begin{equation}\label{8.7:eq-young-member-3}
\tilde{\nu}\big(\mathfrak{lv}_{j}(\square)\big)\le C_L\,e^{-\theta p_0(g_j)}.
\end{equation}
In particular this holds at every level $j\in[k_0-\mathfrak{l}-C_{\mathrm{ext}},\,k_0+r(s)]$ read by
the Gaussian extraction.
\end{itemize}
\end{corollary}

\begin{proof}
By the definition of $u_{\gamma}(g)$ in~\eqref{8.7:eq-depth-function-b}, the inequality
$\mathfrak{F}_{\gamma}(u)\ge0$ holds for every $u\ge u_{\gamma}(g)$. Hence the depth condition
\eqref{8.7:eq-young-block-rarity-b} of part~(d) of Theorem~\ref{8.7:thm-young-block-rarity} is met at
a level $j$ as soon as $K-j\ge u_{\gamma}(g)$. We use this in all three parts; the levels at which
part~(d) is applied below are taken in its range $j\le K'$.

(a) Let $j\in[k_0-\mathfrak{l},k_0]$. Then $j\le k_0$, so
\begin{equation*}
K-j\ge K-k_0=s\ge u_{\gamma}(g).
\end{equation*}
Thus the single condition $s\ge u_{\gamma}(g)$ gives $K-j\ge u_{\gamma}(g)$ at every level of the
interval. Part~(d) of
Theorem~\ref{8.7:thm-young-block-rarity}, applied at each such $j$, gives the per-block weight bound
for every kind $\kappa$ and every cube $\square$, with $q_j=n_{\square}C_Le^{-\theta p_0(g_j)}$. In
particular it gives the bound for the pairs $(\kappa,\square)\in\mathfrak{B}_1(j)\cup\mathfrak{B}_2(j)$,
which is~\eqref{8.7:eq-young-member-1}. Since $\tilde{\nu}=\tau/Z$, the bound can equivalently be read
as $\tilde{\nu}(E^{\kappa}_{j,\square})\le q_j$.

(b) Take $j=k_0$. Then $K-k_0=s\ge u_{\gamma}(g)$, so part~(d) of
Theorem~\ref{8.7:thm-young-block-rarity} applies at level $k_0$. It gives the live-cell bound at that
level, for every level-$k_0$ cell, with the constants $(C_{lv},c_{lv})=(C_L,\theta)$. These are the
values that the live-cell bound of~\eqref{8.7:eq-creation-live-events-a} itself names. This
proves~\eqref{8.7:eq-young-member-2}.

(c) The highest level read by the Gaussian extraction is $k_0+r(s)$, and at that level
\begin{equation*}
K-\big(k_0+r(s)\big)=s-r(s)\ge u_{\gamma}(g).
\end{equation*}
Now let $j\le k_0+r(s)$. Then
\begin{equation*}
K-j\ge s-r(s)\ge u_{\gamma}(g),
\end{equation*}
so every such level has depth at least that of the highest level read, and the condition is met there
a fortiori. Part~(d) of Theorem~\ref{8.7:thm-young-block-rarity}, applied at each such $j$, gives the
live-cell bound at level $j$ with $(C_{lv},c_{lv})=(C_L,\theta)$, which
is~\eqref{8.7:eq-young-member-3}. The interval $[k_0-\mathfrak{l}-C_{\mathrm{ext}},\,k_0+r(s)]$
consists of levels $j\le k_0+r(s)$, so the bound holds throughout it.
\end{proof}

\begin{lemma}[Translation averages suffice for the witness]\label{8.7:lem-translation-average}
Let $s$ be the witness depth of the two-point witness theorem and $k_0=K-s$ the witness level. Let
$\mathcal{V}_{k_0}$ be the group of translations of $\mathbb{T}_m$ by vectors of the level-$k_0$ lattice
(these are the depth-$s$ translations), and for a
function $F$ on $\mathcal{V}_{k_0}$ write
$\langle\!\langle F\rangle\!\rangle=|\mathcal{V}_{k_0}|^{-1}\sum_{v\in\mathcal{V}_{k_0}}F(v)$ and
$\mathbb{P}_v(\cdot)$ for the corresponding uniform probability. Assume:
\begin{itemize}
\item[(A)] the translation-average identity used in the Gaussian extraction of the two-point witness
theorem: with
$f^v(x)=f(x-v)$,
\begin{equation*}
S^{T}_{2;K,m}(f_1,f_2)=\langle\!\langle S^{T}_{2;K,m}(f_1^{v},f_2^{v})\rangle\!\rangle ,
\end{equation*}
so that every $\tilde\nu$-quantity entering the bounds of that extraction may be replaced by its
$\langle\!\langle\cdot\rangle\!\rangle$-average;
\item[(B)] clauses (a) and (c) of the rarity of young large-field blocks,
Theorem~\ref{8.7:thm-young-block-rarity}.
\end{itemize}
For $Y\subset\mathbb{T}_m$ let $n_{k_0}(Y)=\max_{v\in\mathcal{V}_{k_0}}\#\{\mathfrak{c}:\mathfrak{c}\text{ meets }Y+v\}$,
the count running over level-$k_0$ cells. Then
\begin{equation}\label{8.7:eq-translation-average-1}
\begin{aligned}
\Big\langle\!\!\Big\langle\sum_{\mathfrak{c}\text{ meets }Y+v}\tilde\nu\big(L^{(k_0)}_{\mathfrak{c}}\big)\Big\rangle\!\!\Big\rangle
&=\sum_{\mathfrak{c}}\tilde\nu\big(L^{(k_0)}_{\mathfrak{c}}\big)\,
\mathbb{P}_v(\mathfrak{c}\text{ meets }Y+v)\\
&=\big\langle\!\big\langle\#\{\mathfrak{c}\text{ meeting }Y+v\}\big\rangle\!\big\rangle\,\bar q_{k_0}
\;\le\; n_{k_0}(Y)\,\bar q_{k_0},
\end{aligned}
\end{equation}
with $\bar q_{k_0}$ as in \eqref{8.7:eq-young-block-rarity-a}. Consequently, clauses (a) and (c) of
Theorem~\ref{8.7:thm-young-block-rarity} serve both additive large-field bounds of the witness
construction with no input beyond (A) and the two clauses named in (B); in particular the
translation covariance of the
label histories \eqref{8.7:eq-translation-covariance-a} is not used.
\end{lemma}

\begin{proof}
\emph{The identity \eqref{8.7:eq-translation-average-1}.} Write the sum over cells meeting $Y+v$ as
$\sum_{\mathfrak{c}}\tilde\nu(L^{(k_0)}_{\mathfrak{c}})\,\mathbf{1}[\mathfrak{c}\text{ meets }Y+v]$. The weights
$\tilde\nu(L^{(k_0)}_{\mathfrak{c}})$ do not depend on $v$, so averaging over $v$ replaces the indicator by
$\mathbb{P}_v(\mathfrak{c}\text{ meets }Y+v)$; this is the first equality.

This probability is the same for every level-$k_0$ cell. Let $\mathfrak{c}$ and $\mathfrak{c}'$ be two cells, and let
$a$ be the cell translation with $\mathfrak{c}'=\mathfrak{c}+a$. The cells are cubes whose vertices lie in the
level-$k_0$ lattice and which tile the torus by translation,
so $a$ translates by an integer multiple of the level-$k_0$ spacing in each coordinate direction, and
hence $a\in\mathcal{V}_{k_0}$. Now $\mathfrak{c}'$ meets $Y+v$ if and only if $\mathfrak{c}$ meets $Y+(v-a)$.
The map $v\mapsto v-a$ is a bijection of $\mathcal{V}_{k_0}$ and preserves the uniform measure. Hence
$\mathbb{P}_v(\mathfrak{c}'\text{ meets }Y+v)=\mathbb{P}_v(\mathfrak{c}\text{ meets }Y+v)$. Call this common
value $\varpi$.

By the definition of $\bar q_{k_0}$ in \eqref{8.7:eq-young-block-rarity-a}, the second member of
\eqref{8.7:eq-translation-average-1},
$\sum_{\mathfrak{c}}\tilde\nu(L^{(k_0)}_{\mathfrak{c}})\,\mathbb{P}_v(\mathfrak{c}\text{ meets }Y+v)$, equals
$\varpi\,n_{k_0}\bar q_{k_0}$, where $n_{k_0}$, without argument, is the number of level-$k_0$ cells of the
torus, as in \eqref{8.7:eq-young-block-rarity-a}. By linearity of the average,
\begin{equation*}
\langle\!\langle\#\{\mathfrak{c}\text{ meeting }Y+v\}\rangle\!\rangle
=\sum_{\mathfrak{c}}\mathbb{P}_v(\mathfrak{c}\text{ meets }Y+v)=n_{k_0}\,\varpi .
\end{equation*}
This proves the second equality. The last inequality holds because
$\#\{\mathfrak{c}\text{ meeting }Y+v\}\le n_{k_0}(Y)$ for every $v$, by the definition of $n_{k_0}(Y)$.

\emph{The opening step of the additive large-field bound at the witness level.} By (A) every
$\tilde\nu$-quantity in that step may be replaced by its $\langle\!\langle\cdot\rangle\!\rangle$-average.
In each item of its opening step the one-point masses have the form $n_i\,q^{lv}_{k_0}$, where $n_i$ is
the position-free cell count that this step attaches to ball $i$, that is, a bound, uniform in the
position, for the number of level-$k_0$ cells meeting the relevant set $Y_i$; thus
$n_{k_0}(Y_i)\le n_i$ in the notation above. The quantities bounded there by $n_i\,q^{lv}_{k_0}$, through the
one-point bounds \eqref{8.7:eq-creation-live-events-a}, are sums of $\tilde\nu(L^{(k_0)}_{\mathfrak{c}})$ over the
level-$k_0$ cells meeting a translate of $Y_i$. By \eqref{8.7:eq-translation-average-1} their
$\langle\!\langle\cdot\rangle\!\rangle$-averages are bounded by $n_i\,\bar q_{k_0}$, so after this replacement
these bounds hold with
\begin{equation*}
q^{lv}_{k_0}:=\bar q_{k_0},
\end{equation*}
which clause (a) of Theorem~\ref{8.7:thm-young-block-rarity} controls.

\emph{The later steps.} Every later step of the additive large-field bound at the witness level does one
of two things.
\begin{itemize}
\item[(1)] It is linear in these masses. Then the average passes through it term by term.
\item[(2)] It bounds a second factor by its supremum over positions. For functions $F\ge0$ and $G$ on
$\mathcal{V}_{k_0}$ one has $FG\le F\,\sup_{v}G$ pointwise in $v$, and therefore
\begin{equation*}
\langle\!\langle FG\rangle\!\rangle\le\Big(\sup_v G\Big)\langle\!\langle F\rangle\!\rangle .
\end{equation*}
\end{itemize}
The main term of the two-point witness theorem, the one isolated by the Gaussian extraction, is
invariant under depth-$s$ translations, which is the form in which hypothesis (A) is
used in the Gaussian extraction.

\emph{The blocks at levels $j<k_0$.} In the additive bound for blocks of levels $j<k_0$, enlarge each block
family
$\mathfrak{B}_i(j)$ to the family of all step-$j$ cubes contained in those level-$k_0$ cells that meet a
cube of $\mathfrak{B}_i(j)$, and denote the enlarged family
again by $\mathfrak{B}_i(j)$. Enlarging the family only increases the sum, its terms being non-negative;
and since a union of whole level-$k_0$ cells contains every class of step-$j$ cubes modulo
$\mathcal{V}_{k_0}$ equally often, the $v$-average of the sum over this family is $\#\mathfrak{B}_i(j)$ times the
average over the torus of the quantity bounded in clause (c) of Theorem~\ref{8.7:thm-young-block-rarity}.

\emph{Conclusion.} Both additive bounds therefore use only the averaged quantities controlled by
clauses (a) and (c) of Theorem~\ref{8.7:thm-young-block-rarity}, after the replacement by translation
averages that (A) permits. Neither the
position-uniform clause of that theorem nor the covariance \eqref{8.7:eq-translation-covariance-a} enters.
\end{proof}

\begin{remark}[The two depth conditions as the coupling tends to zero]\label{8.7:rem-depth-conditions}
Let $\mathfrak{F}_{\gamma}$ and $u_{\gamma}(g)$ be as in \eqref{8.7:eq-depth-function-b} and
$\mathfrak{e}_j$ as in \eqref{8.7:eq-depth-function-a}
(Notation~\ref{8.7:not-depth-function}), and write $k_0=K-s$ for the level of the witness depth $s$,
with $s\ge m_{\mathbb{T}}(g)$, so that $k_0\le K'$. Both conditions below are lower bounds on the
depth $s$ depending on $g$; neither bounds any quantity from above, and neither is a hypothesis of
Theorem~0.
\begin{itemize}
\item[(a)] The depth condition \eqref{8.7:eq-young-block-rarity-b}, $s-r(s)\ge u_{\gamma}(g)$, is for
each $g$ a lower bound on $s$: the number $u_{\gamma}(g)$ is finite. As $g\to0$, provided
$\mathsf{B}(g)>0$ and $\log\mathsf{B}(g)=o\bigl((\log g^{-2})^{p_0}\bigr)$,
\begin{equation}\label{8.7:eq-depth-conditions-1}
u_{\gamma}(g)=m_{\mathbb{T}}(g)+\frac{\theta A_0}{d\log L}\,(\log g^{-2})^{p_0}\,\bigl(1+o(1)\bigr).
\end{equation}
\item[(b)] The budget inequality of the additive assembly at the witness level,
\begin{equation}\label{8.7:eq-depth-conditions-2}
L^{8(\mathfrak{l}(\bar{x}_s)+1)}\,\mathfrak{p}^{\mathrm{add}}_{\mathrm{tot}}(\bar{x}_s)\,\bar{x}_s^{2}\,
q^{lv}_{k_0}\le\frac{\vartheta^{4}c_M}{64\,C^{\mathrm{add}}_{\mathrm{tot}}},
\end{equation}
in which $q^{lv}_{k_0}$ is the bound of Theorem~\ref{8.7:thm-young-block-rarity} at level $k_0$,
the sum of a healed member and the never-healed member $\mathfrak{e}_{k_0}$, holds as soon as each
member satisfies it with $128$ in place of $64$. For the healed member this is exactly the condition
that the additive assembly imposes on that member. For the never-healed member it is equivalent to
\begin{equation}\label{8.7:eq-depth-conditions-a}
\begin{split}
d\bigl(s-m_{\mathbb{T}}(g)\bigr)\ge{}&8\bigl(\mathfrak{l}(\bar{x}_s)+1\bigr)+2\log_L\bar{x}_s\\
&+\log_L\frac{128\,C^{\mathrm{add}}_{\mathrm{tot}}\,\mathfrak{p}^{\mathrm{add}}_{\mathrm{tot}}(\bar{x}_s)\cdot
2C_{\mathrm{cnt}}\bigl(\mathsf{B}(g)+\log c_J\bigr)\bigl(M\check{R}_{k_0}\bigr)^{d}}
{\vartheta^{4}c_M\,t_{k_0}},
\end{split}
\end{equation}
where the lifetime $\mathfrak{l}(\bar{x})$ of the additive assembly is $L\,R(\bar{x})$ plus a
constant of that assembly, a polynomial of degree $r_{\mathrm{pr}}<p_0$ in $\log\bar{x}$.
In the organisation of Corollary~\ref{8.7:cor-young-member}, where the number of blocks read carries
the factor $L^{4\mathfrak{l}}$, the right-hand side
of \eqref{8.7:eq-depth-conditions-a} acquires the further term $4\mathfrak{l}$. Assume the proviso
of (a); assume $p_0>1$, which is needed for the member $2\log_L\bar{x}_s$, linear in
$\log\bar{x}_s$; assume
$\log\mathfrak{p}^{\mathrm{add}}_{\mathrm{tot}}(\bar{x})=o\bigl((\log\bar{x})^{p_0}\bigr)$ as
$\bar{x}\to\infty$; and assume that at the level $k_0$ the two bounds
\begin{equation*}
M\check{R}_{k_0}\le ML(\log\bar{x}_s)^{r_{\mathrm{pr}}},\qquad
t_{k_0}\ge\theta'\gamma_0A_1\bigl(p_0(g)-1\bigr)^{2}
\end{equation*}
hold. It is through these two bounds that $\check{R}_{k_0}$ and $t_{k_0}$ enter $\mathfrak{F}_{\gamma}$.
Then for every $\varepsilon>0$ and all $g$ small enough the right-hand side of
\eqref{8.7:eq-depth-conditions-a}, with or without the further term $4\mathfrak{l}$, is at most
$\varepsilon(\log\bar{x}_s)^{p_0}$ for every $s$ with $m_{\mathbb{T}}(g)\le s\le K$. Moreover, for
$g$ small enough, every $s$ with $s-r(s)\ge u_{\gamma}(g)$ satisfies
\eqref{8.7:eq-depth-conditions-a}, with or without that term; in this sense
\eqref{8.7:eq-depth-conditions-a} is weaker than \eqref{8.7:eq-young-block-rarity-b} as $g\to0$.
\item[(c)] If the depth of clause (iii)(c) of Theorem~0 is required to satisfy
\eqref{8.7:eq-young-block-rarity-b} in addition to $s\ge s_W(g)$, the admissible depths are those
with
\begin{equation*}
s\ge s_W(g)\quad\text{and}\quad s-r(s)\ge u_{\gamma}(g),
\end{equation*}
that is, if the extraction depth $r(s)$ takes one value $r$ for all $s\ge s_W(g)$, the floor
$s_W(g)$ becomes $\max\{s_W(g),u_{\gamma}(g)+r\}$, and
the admissible separations shrink to $\mathsf{d}\lesssim L^{-u_{\gamma}(g)}$. The
conclusion of that clause then remains a statement at all sufficiently short distances, for every
cutoff and every volume; the floor becomes part of the conclusion of clause (iii)(c).
\end{itemize}
\end{remark}

\begin{proof}
(a) For $u\ge m_{\mathbb{T}}(g)$, dividing the definition of $\mathfrak{F}_{\gamma}(u)$ by
$d\log L>0$ shows that $\mathfrak{F}_{\gamma}(u)\ge0$ if and only if
\begin{equation}\label{8.7:eq-depth-conditions-3}
\begin{split}
u-m_{\mathbb{T}}(g)\ge\frac{1}{d\log L}\Bigl[&\theta A_0(\log\bar{x}(u))^{p_0}
+d\log\bigl(ML(\log\bar{x}(u))^{r_{\mathrm{pr}}}\bigr)\\
&+\log\bigl(2C_{\mathrm{cnt}}(\mathsf{B}(g)+\log c_J)\bigr)
-\log\bigl(\theta'\gamma_0A_1(p_0(g)-1)^{2}\bigr)\Bigr].
\end{split}
\end{equation}
For fixed $g$ the first member of $\mathfrak{F}_{\gamma}(u)$ is linear in $u$ with slope
$d\log L>0$, and the other members are polylogarithmic in $\bar{x}(u)$, which is affine in $u$.
Hence $\mathfrak{F}_{\gamma}(u)\to+\infty$ as $u\to+\infty$, the set of $u_0\ge m_{\mathbb{T}}(g)$ with
$\mathfrak{F}_{\gamma}(u)\ge0$ for all $u\ge u_0$ is non-empty, and $u_{\gamma}(g)$ is finite.
Consequently \eqref{8.7:eq-young-block-rarity-b} is a lower bound on $s$ for each $g$.

For the asymptotics recall $X=g^{-2}$, so $\bar{x}(u)=g^{-2}+b_++3\lambda u+2c_2$. Write
$c=\theta A_0/(d\log L)$ and $\Phi_g(u)$ for the right-hand side of
\eqref{8.7:eq-depth-conditions-3}. The number $m_{\mathbb{T}}(g)=11+m'+\log_L R(g)$ is of order
$\log\log g^{-2}$, the function $R$ being a polynomial of degree $r_{\mathrm{pr}}$ in the logarithm
of the inverse square coupling (we write $R(g)$ for its value at $g^{-2}$, as in the glossary entry
for $m_{\mathbb{T}}$, and $R(\bar{x})$ for its value at $\bar{x}$); in particular
$m_{\mathbb{T}}(g)\le g^{-2}/2$ for $g$ small. Since
$p_0(g)=A_0(\log g^{-2})^{p_0}\to\infty$, the member $-\log(\theta'\gamma_0A_1(p_0(g)-1)^2)$ does not
depend on $u$, is of order $\log\log g^{-2}$, and is $\le0$ for $g$ small. The member
$\log(2C_{\mathrm{cnt}}(\mathsf{B}(g)+\log c_J))$ does not depend on $u$. Since $\mathsf{B}(g)>0$ and
$\log c_J>0$ (as $c_J>e^{1/2}>1$), one has
$\log c_J\le\mathsf{B}(g)+\log c_J\le2\max\{\mathsf{B}(g),\log c_J\}$, so by the proviso
$\log\mathsf{B}(g)=o\bigl((\log g^{-2})^{p_0}\bigr)$ this member is
$o\bigl((\log g^{-2})^{p_0}\bigr)$.

First let $m_{\mathbb{T}}(g)\le u\le g^{-2}$. Since $\lambda u\ge0$,
\begin{equation*}
g^{-2}+b_++2c_2\le\bar{x}(u)\le(1+3\lambda)g^{-2}+b_++2c_2,
\end{equation*}
so $\log\bar{x}(u)=\log g^{-2}+O(1)$ uniformly in such $u$, and
$(\log\bar{x}(u))^{p_0}=(\log g^{-2})^{p_0}(1+o(1))$ uniformly; the member
$d\log(ML(\log\bar{x}(u))^{r_{\mathrm{pr}}})$ is $O(\log\log g^{-2})$ uniformly. With the two
$u$-independent members above, $\Phi_g(u)=c(\log g^{-2})^{p_0}(1+o(1))$ uniformly for
$m_{\mathbb{T}}(g)\le u\le g^{-2}$. Fix $\varepsilon\in(0,1)$ and put
\begin{equation*}
u_{\pm}=m_{\mathbb{T}}(g)+(1\pm\varepsilon)\,c\,(\log g^{-2})^{p_0};
\end{equation*}
for $g$ small, $u_{\pm}\le g^{-2}$. At $u=u_-$ one has, for $g$ small,
\begin{equation*}
u_--m_{\mathbb{T}}(g)=(1-\varepsilon)c(\log g^{-2})^{p_0}<\Phi_g(u_-),
\end{equation*}
so $\mathfrak{F}_{\gamma}(u_-)<0$; as $u_-\ge m_{\mathbb{T}}(g)$, the definition of $u_{\gamma}(g)$ gives
$u_{\gamma}(g)>u_-$. For $u_+\le u\le g^{-2}$ one has, for $g$ small and uniformly in $u$,
\begin{equation*}
u-m_{\mathbb{T}}(g)\ge(1+\varepsilon)c(\log g^{-2})^{p_0}>\Phi_g(u),
\end{equation*}
so $\mathfrak{F}_{\gamma}(u)\ge0$.

Next let $u\ge g^{-2}$. Then $\bar{x}(u)\le(1+3\lambda)u+b_++2c_2\le(2+3\lambda)u$ for $g$ small,
hence $\log\bar{x}(u)\le2\log u$ for $g$ small. Dropping the non-positive member
$-\log(\theta'\gamma_0A_1(p_0(g)-1)^2)$ and using $\log g^{-2}\le\log u$ in the bound for the
$\mathsf{B}$-member, one finds a constant $C'$ with $\Phi_g(u)\le C'(\log u)^{p_0}$ for all
$u\ge g^{-2}$ and $g$ small. On the other hand $u-m_{\mathbb{T}}(g)\ge u/2$, since
$m_{\mathbb{T}}(g)\le g^{-2}/2\le u/2$. As $(\log u)^{p_0}/u\to0$, one has
$u/2>C'(\log u)^{p_0}$ for all $u\ge g^{-2}$ once $g$ is small, so $\mathfrak{F}_{\gamma}(u)\ge0$ for
all $u\ge g^{-2}$.

Together, $\mathfrak{F}_{\gamma}\ge0$ on $[u_+,\infty[$, so $u_{\gamma}(g)\le u_+$, while
$u_{\gamma}(g)>u_-$. Since $\varepsilon\in(0,1)$ is arbitrary, \eqref{8.7:eq-depth-conditions-1}
follows.

(b) Both members of $q^{lv}_{k_0}$ are non-negative, and the factor
\begin{equation*}
L^{8(\mathfrak{l}(\bar{x}_s)+1)}\,\mathfrak{p}^{\mathrm{add}}_{\mathrm{tot}}(\bar{x}_s)\,\bar{x}_s^{2}
\end{equation*}
is positive. If this factor times each member is at most
$\vartheta^4c_M/(128\,C^{\mathrm{add}}_{\mathrm{tot}})$, then the factor times their sum
$q^{lv}_{k_0}$ is at most twice this bound, that is, at most
$\vartheta^4c_M/(64\,C^{\mathrm{add}}_{\mathrm{tot}})$; this is the first assertion.
Since $K-k_0=s$, \eqref{8.7:eq-depth-function-a} reads
\begin{equation*}
\mathfrak{e}_{k_0}=\frac{2C_{\mathrm{cnt}}(\mathsf{B}(g)+\log c_J)\,(M\check{R}_{k_0})^{d}}{t_{k_0}}\,
L^{-d(s-m_{\mathbb{T}}(g))}.
\end{equation*}
Insert this into the inequality for the never-healed member,
\begin{equation*}
L^{8(\mathfrak{l}(\bar{x}_s)+1)}\,\mathfrak{p}^{\mathrm{add}}_{\mathrm{tot}}(\bar{x}_s)\,\bar{x}_s^{2}\,
\mathfrak{e}_{k_0}\le\frac{\vartheta^4c_M}{128\,C^{\mathrm{add}}_{\mathrm{tot}}}.
\end{equation*}
Multiplying by $L^{d(s-m_{\mathbb{T}}(g))}$ and by
$128C^{\mathrm{add}}_{\mathrm{tot}}/(\vartheta^4c_M)$, and taking $\log_L$ of both sides, gives
\eqref{8.7:eq-depth-conditions-a}; every step is reversible.

For the comparison, write $y_s=\log\bar{x}_s$ and $D(s)$ for the right-hand side of
\eqref{8.7:eq-depth-conditions-a}, with or without the further term $4\mathfrak{l}(\bar{x}_s)$.
Since $\bar{x}_s\ge g^{-2}+b_++2c_2\ge g^{-2}/2$ for $g$ small and all $s\ge0$, $y_s\to\infty$ as
$g\to0$ uniformly in $s$, and $\log g^{-2}\le y_s+\log 2$. We bound $D(s)$ member by member. The
members $8(\mathfrak{l}(\bar{x}_s)+1)$ and $4\mathfrak{l}(\bar{x}_s)$ are polynomials of degree
$r_{\mathrm{pr}}$ in $y_s$, and $2\log_L\bar{x}_s=2y_s/\log L$. In the last member,
$\log_L(256\,C^{\mathrm{add}}_{\mathrm{tot}}C_{\mathrm{cnt}}/(\vartheta^4c_M))$ is a constant; by the
first of the two displayed bounds, $d\log_L(M\check{R}_{k_0})\le d\log_L(MLy_s^{r_{\mathrm{pr}}})$,
which is $O(\log y_s)$; by the second,
$-\log_L t_{k_0}\le-\log_L(\theta'\gamma_0A_1(p_0(g)-1)^2)\le0$ for $g$ small;
$\log_L\mathfrak{p}^{\mathrm{add}}_{\mathrm{tot}}(\bar{x}_s)=o(y_s^{p_0})$ by the proviso on
$\mathfrak{p}^{\mathrm{add}}_{\mathrm{tot}}$; and
\begin{equation*}
\log_L\bigl(\mathsf{B}(g)+\log c_J\bigr)\le\log_L2+\max\{\log_L\mathsf{B}(g),\log_L\log c_J\},
\end{equation*}
which by the proviso of (a) is $o\bigl((\log g^{-2})^{p_0}\bigr)$, hence $o(y_s^{p_0})$ uniformly in
$s$. Since $\max\{r_{\mathrm{pr}},1\}<p_0$, for every $\varepsilon>0$ there is $g_{\varepsilon}>0$
such that $D(s)\le\varepsilon y_s^{p_0}$ for all $g<g_{\varepsilon}$ and all $s$ with
$m_{\mathbb{T}}(g)\le s\le K$; this is the penultimate assertion of (b).

Now let $s-r(s)\ge u_{\gamma}(g)$. As $r(s)\ge0$, $s\ge u_{\gamma}(g)$, and by
\eqref{8.7:eq-depth-conditions-1}
\begin{equation*}
d\bigl(s-m_{\mathbb{T}}(g)\bigr)\ge d\bigl(u_{\gamma}(g)-m_{\mathbb{T}}(g)\bigr)
\ge\frac{\theta A_0}{2\log L}\,(\log g^{-2})^{p_0}
\end{equation*}
for $g$ small. If $s\le g^{-2}$, then $\bar{x}_s\le(2+3\lambda)g^{-2}$ for $g$ small, so
$y_s\le\log g^{-2}+\log(2+3\lambda)$ and $y_s^{p_0}\le2(\log g^{-2})^{p_0}$ for $g$ small; with
$\varepsilon=\theta A_0/(4\log L)$ this gives
\begin{equation*}
D(s)\le\frac{\theta A_0}{2\log L}\,(\log g^{-2})^{p_0}\le d\bigl(s-m_{\mathbb{T}}(g)\bigr).
\end{equation*}
If $s\ge g^{-2}$, then
$\bar{x}_s\le(2+3\lambda)s$ and $y_s\le2\log s$ for $g$ small, so with $\varepsilon\le1$ one has
$D(s)\le2^{p_0}(\log s)^{p_0}$, while $d(s-m_{\mathbb{T}}(g))\ge s/2$ because
$m_{\mathbb{T}}(g)\le g^{-2}/2\le s/2$; as $(\log s)^{p_0}/s\to0$, $D(s)\le s/2$ for all
$s\ge g^{-2}$ once $g$ is small. In both cases \eqref{8.7:eq-depth-conditions-a} holds.

(c) The condition $s-r(s)\ge u_{\gamma}(g)$ reads $s\ge u_{\gamma}(g)+r(s)$; together with
$s\ge s_W(g)$ it is the displayed condition. If $r(s)=r$ for all $s\ge s_W(g)$, the displayed
condition is $s\ge\max\{s_W(g),u_{\gamma}(g)+r\}$. By (a), $u_{\gamma}(g)$ is finite for every $g$ and
independent of $K$ and of $m$; the condition is then a lower bound on $s$ only, so the admissible
depths still comprise all sufficiently large $s$, that is, all sufficiently short separations, for
every cutoff and every volume.
\end{proof}

\begin{remark}[What the first moment does not reach at shallow levels]\label{8.7:rem-shallow-levels}
Let $K'=K-m_{\mathbb{T}}(g)$ be the terminal level of Theorem~\ref{8.7:thm-young-block-rarity}
(Notation~\ref{8.7:not-depth-function}). Let $u_{\gamma}(g)$ be the least admissible depth
and $\mathfrak{e}_j$ the first-moment error term, both fixed in
Notation~\ref{8.7:not-depth-function} (see \eqref{8.7:eq-depth-function-b} and
\eqref{8.7:eq-depth-function-a}). Consider the levels
\begin{equation*}
K-u_{\gamma}(g)<j\le K-m_{\mathbb{T}}(g).
\end{equation*}
At these levels parts (a)--(c) of Theorem~\ref{8.7:thm-young-block-rarity} hold, but
$\mathfrak{e}_j$ is not bounded by $e^{-\theta p_0(g_j)}$. At $j=K'$ the error term is of the
order of
\begin{equation*}
\frac{\mathsf{B}(g)\,(M\check{R}_{K'})^{d}}{t_{K'}} .
\end{equation*}
This quantity is small only as a negative power of $\log g^{-2}$. It is so if twice the exponent
$p_0$ of the degree function $p_0(\cdot)$ exceeds $d\,r_{\mathrm{pr}}$, with $r_{\mathrm{pr}}$ the
degree of $R(\cdot)$ in $\log g^{-2}$, plus the degree of $\mathsf{B}(g)$ as a power of
$\log g^{-2}$.

The inputs used cannot do better; this is part~(iii) of the first-moment bound for an extensive
mark, Lemma~\ref{8.7:lem-first-moment}. What these levels require is the rarity of the
never-healed fate of a level-$j$ cell $\mathfrak{c}$, that is of the live event
$\mathcal{E}^{(j)}_{\mathfrak{c}}(K')$ that $\mathfrak{c}$ is still live, not yet healed, at the
terminal level, at bounded depth below the terminal level.
That rarity amounts to one of two things: a moderate-deviation bound for the terminal
field itself at level $K'$, or one further block constraint.

The same inputs do not give the age-weighted one-point bound at the witness level $k_0=K-s$.
That bound controls
\begin{equation*}
\sum_{a>\mathfrak{l}} L^{8(a+1)}\,q^{=a},
\end{equation*}
where $\mathfrak{l}$ is the lifetime threshold of that bound, the age above which the sum runs,
and $q^{=a}$ is the one-point quantity of age exactly $a$ entering that bound.

Any statement that imposes the depth condition \eqref{8.7:eq-young-block-rarity-b}, in
particular in the form $s-r(s)\ge u_{\gamma}(g)$, uses no level of the displayed range.
\end{remark}

\begin{proof}
Let $j$ lie in the displayed range. Then $j\le K'$. Theorem~\ref{8.7:thm-young-block-rarity}
is stated for every $g$ below its ceilings, every $K$, every $N_T\ge 1$ and every level
$j\le K'$, so its parts (a)--(c) hold at $j$. On the other hand $K-j<u_{\gamma}(g)$, so the
depth condition $K-j\ge u_{\gamma}(g)$ of part~(d), displayed in
\eqref{8.7:eq-young-block-rarity-b}, fails at $j$. Thus $j$ lies at a depth $K-j$ below the
least admissible depth $u_{\gamma}(g)$ of \eqref{8.7:eq-depth-function-b}; we make the size of
$\mathfrak{e}_j$ explicit at $j=K'$.

At $j=K'$, the definition of $\mathfrak{e}_j$ in \eqref{8.7:eq-depth-function-a} makes
$\mathfrak{e}_{K'}$ of the order of $\mathsf{B}(g)(M\check{R}_{K'})^d/t_{K'}$. We compare the
three factors with powers of $\log g^{-2}$.
In this comparison $\log g_{K'}^{-2}$ and $\log\min_{i\le K'}x_i$ are counted on the scale of
$\log g^{-2}$, the scale on which the order of magnitude of $\mathfrak{e}_{K'}$ is stated.
\begin{itemize}
\item The reserve per marked cell is $t_{K'}=\theta'\gamma_0A_1\bigl(p_0(g_{K'})-1\bigr)^2$.
Here the degree function is $p_0(g')=A_0(\log g'^{-2})^{p_0}$. So $t_{K'}$ grows as the
$2p_0$-th power of the logarithm.
\item $M$ is fixed and $R(\cdot)$ has degree $r_{\mathrm{pr}}$. So $(M\check{R}_{K'})^d$ grows
as the $d\,r_{\mathrm{pr}}$-th power.
\item $\mathsf{B}(g)$ contributes its own degree.
\end{itemize}
The quotient is therefore a power of $\log g^{-2}$ with exponent
$d\,r_{\mathrm{pr}}+\text{(degree of }\mathsf{B})-2p_0$. This exponent is negative exactly under
the stated condition on $2p_0$. By contrast,
\begin{equation*}
e^{-\theta p_0(g_{K'})}=\exp\bigl(-\theta A_0(\log g_{K'}^{-2})^{p_0}\bigr)
\end{equation*}
decays faster than every power of the logarithm as the logarithm grows. So a quantity that is
small only as a power of the logarithm is not bounded by it.

The bound on $\mathfrak{e}_j$ comes from a global upper bound with reserve and a global floor,
that is, from hypotheses of the form $(\mathrm{G}\!\uparrow)$ and $(\mathrm{G}\!\downarrow)$ of
Lemma~\ref{8.7:lem-first-moment}. Part~(iii) of that lemma applies to these inputs. From them
one obtains a first moment of the extensive mark and no better bound.

For the age-weighted bound, consider first the error term. It is blind to age: the error term
$\mathfrak{e}_j$ of Notation~\ref{8.7:not-depth-function} carries no dependence on the age $a$,
so it does not compensate the factor $L^{8a}$. Consider next the one-point bound at level
$k_0-a$, whose rate is $e^{-\theta p_0(g_{k_0-a})}$. This rate does not compensate $L^{8a}$
either. Finally, suppose the ages are marked in the extensive mark of
Lemma~\ref{8.7:lem-first-moment}. Then the $\tilde{\nu}$-expectation of the total age over the
marked cells is bounded by $\mathsf{B}(g)N^{d}$ divided by the reserve carried per unit of age,
where $N=N_TL^{K-K'}$ is the side of the torus in units of the terminal level
(Notation~\ref{8.7:not-depth-function}). This
is a first moment. By part~(iii) of Lemma~\ref{8.7:lem-first-moment}, these inputs give no
exponential tail in the age, and such a tail is what the factor $L^{8(a+1)}$ would require.
\end{proof}

\begin{lemma}[Reflection positivity and the chessboard estimate on any torus]\label{8.7:lem-chessboard-torus}
Let
\begin{equation*}
\mathbb{T}=\prod_{\mu=1}^{d}\mathbb{Z}/n_{\mu}\mathbb{Z}
\end{equation*}
be a discrete torus with arbitrary periods $n_{\mu}\ge 3$; odd periods are allowed. A bond is an ordered
pair $(x,y)$ of nearest-neighbour sites. Let $G$ be a compact matrix group. We take its elements to be
unitary matrices. This loses no generality: averaging a Hermitian inner product over the Haar measure of
$G$ gives an inner product for which $G$ is unitary, so $G$ is conjugate to a group of unitary matrices,
and conjugating every bond variable by one fixed matrix leaves the action $S$ defined below unchanged.
A configuration $U$ assigns to every bond an element $U(x,y)$ of $G$, with $U(y,x)=U(x,y)^{-1}$, and we
write $U_b=U(x,y)$ for $b=(x,y)$. For a plaquette $p$, $U(\partial p)$ is the ordered product of the bond
variables along the boundary of $p$. Let $\beta\ge 0$, let $\operatorname{tr}$ be the trace in any
normalisation, and put
\begin{equation*}
S(U):=\beta\sum_{p}\bigl(1-\operatorname{Re}\operatorname{tr}U(\partial p)\bigr),\qquad
d\mu_{\mathrm{un}}:=e^{-S}\prod_{b}dU_b ,
\end{equation*}
where $dU_b$ is the normalised Haar measure of $G$ and the product runs over the bonds, each taken with
one orientation. Let $\mu$ be the probability measure $\mu_{\mathrm{un}}/\int d\mu_{\mathrm{un}}$ and
$\langle\,\cdot\,\rangle$ its expectation.

For $c\in\frac12\mathbb{Z}$ let $\vartheta$ be the reflection $x_1\mapsto 2c-x_1$ of $\mathbb{T}$, the
other coordinates being fixed. Its fixed hyperplanes are $x_1=c$ and $x_1=c+n_1/2$; each of them either
passes through sites (when its coordinate is an integer) or bisects bonds (when its coordinate is a
half-integer), and for $n_1$ odd there is one of each type. Let $\Lambda_+$ be the set of bonds both of
whose ends lie in the closed arc $x_1\in[c,c+n_1/2]$, let $\mathfrak{A}_+$ be the set of bounded
measurable complex functions of $U\restriction\Lambda_+$, and put
\begin{equation*}
(\vartheta^{*}U)(x,y):=U(\vartheta x,\vartheta y),\qquad (\theta F)(U):=\overline{F(\vartheta^{*}U)}.
\end{equation*}
\begin{itemize}
\item[(a)] For every $F\in\mathfrak{A}_+$,
\begin{equation}\label{8.7:eq-chessboard-torus-1}
\int F\,\theta F\,d\mu_{\mathrm{un}}\ \ge\ 0 ;
\end{equation}
no gauge invariance of $F$ is required.
\item[(b)] Let $s$ be a positive integer with $s\mid n_{\mu}$ for all $\mu$, put $r_{\mu}:=n_{\mu}/s$, let
$c^{\mu}\in\frac12\mathbb{Z}$ and $c^{\mu}_i:=c^{\mu}+i\,s/2$ for $i\in\mathbb{Z}/2r_{\mu}\mathbb{Z}$. The
hyperplanes $x_{\mu}=c^{\mu}_i$ cut $\mathbb{T}$ into the $n_c:=\prod_{\mu}2r_{\mu}$ closed tiles
\begin{equation*}
\Delta_m:=\bigl\{x:\ x_{\mu}\in[c^{\mu}_{m_{\mu}},c^{\mu}_{m_{\mu}+1}]\ \text{for all }\mu\bigr\},\qquad
m\in\prod_{\mu}\mathbb{Z}/2r_{\mu}\mathbb{Z},
\end{equation*}
an even number $2r_{\mu}$ of them in each direction. Let $\vartheta^{(\mu)}_i$ be the reflection
$x_{\mu}\mapsto 2c^{\mu}_i-x_{\mu}$, let
$\Phi^{(\mu)}_l:=\vartheta^{(\mu)}_l\circ\dots\circ\vartheta^{(\mu)}_1$ for $0\le l\le 2r_{\mu}-1$
($\Phi^{(\mu)}_0$ the identity), and
$\Phi_m:=\Phi^{(1)}_{m_1}\circ\dots\circ\Phi^{(d)}_{m_d}$. Let $F\ge 0$ be a bounded measurable function
of the bonds with both ends in one closed tile, which we index as $\Delta_0$ (the family of hyperplanes is
unchanged by $i\mapsto i+j$), and let $F^{[m]}(U):=F(\Phi_m^{*}U)$, where
$(\Phi^{*}U)(x,y):=U(\Phi x,\Phi y)$; thus $F^{[m]}$ is the copy of $F$ transported by successive tile
reflections. Then
\begin{equation}\label{8.7:eq-chessboard-torus-2}
\langle F\rangle\ \le\ \Bigl\langle\prod_{\text{tiles }m}F^{[m]}\Bigr\rangle^{1/n_c}.
\end{equation}
\item[(c)] Let $s$ be as in \textup{(b)}. For
$a\in(s\mathbb{Z})^d$ and an event $A$, let $A+a:=\{U:\ U(\cdot+a,\cdot+a)\in A\}$. For every event
$A$ of the bonds of a box with at most $\lfloor s/2\rfloor+1$ sites per side,
\begin{equation}\label{8.7:eq-chessboard-torus-3}
\mu(A)\ \le\ \mu\Bigl(\bigcap_{a\in(s\mathbb{Z})^d/\sim}(A+a)\Bigr)^{1/n_c},\qquad
n_c=\prod_{\nu=1}^{d}\frac{2n_{\nu}}{s},
\end{equation}
where $a$ runs over $(s\mathbb{Z})^d$ modulo the periods of $\mathbb{T}$.
\end{itemize}
\end{lemma}

\begin{proof}
\emph{Proof of \textup{(a)}.} The elements of $G$ are unitary, so $U^{-1}=U^{\dagger}$. The bonds split
into four disjoint classes: $\Lambda_+\setminus\Lambda_0$; its mirror image
$\vartheta(\Lambda_+\setminus\Lambda_0)$; the set $\Lambda_0$ of bonds lying inside a fixed hyperplane
through sites, each of which is fixed by $\vartheta$ together with its ends; and the set $C$ of bonds
$(x,\vartheta x)$ bisected by a fixed hyperplane of bond type, each of which is reversed by $\vartheta$.
Indeed, a bond not bisected by a fixed hyperplane has both ends in one of the two closed arcs
$[c,c+n_1/2]$, $[c+n_1/2,c+n_1]$, and if its ends lie in both arcs they lie in the intersection of the
arcs, which consists of the fixed hyperplanes through sites.

The plaquettes split accordingly into four classes. First, those all of whose bonds lie in $\Lambda_+$
but not all in $\Lambda_0$; we write $S_+$ for the sum of the terms
$\beta(1-\operatorname{Re}\operatorname{tr}U(\partial p))$ over them, a real function of
$U\restriction\Lambda_+$. Second, their mirror images; the sum of their terms is
$\theta S_+=S_+\circ\vartheta^{*}$. Third, those all of whose bonds lie in $\Lambda_0$; the sum of their
terms is a real function $S_0(U_0)$ of $U_0:=U\restriction\Lambda_0$. Fourth, the crossing plaquettes,
those containing bonds of $C$: a plaquette containing no bond of $C$ has its corners in one closed arc
and so belongs to one of the first three classes. A crossing plaquette has boundary
\begin{equation*}
(x,y),\ (y,\vartheta y),\ (\vartheta y,\vartheta x),\ (\vartheta x,x),\qquad b=(x,y)\in\Lambda_+ .
\end{equation*}
Let $\Lambda_{\times}$ be the set of the bonds $b$ so arising and $n_{\times}$ its cardinality, which is
the number of crossing plaquettes. For $b=(x,y)\in\Lambda_{\times}$ put
\begin{equation*}
A_b:=U_b,\qquad W_b:=(\vartheta^{*}U)_b=U(\vartheta x,\vartheta y),\qquad h_x:=U(x,\vartheta x),
\end{equation*}
the last for every site $x$ of the arc $[c,c+n_1/2]$ that is an end of a bond of $C$. Since
$U(\vartheta y,\vartheta x)=W_b^{-1}=W_b^{\dagger}$ and $U(\vartheta x,x)=h_x^{-1}$, cyclicity of the
trace gives
\begin{equation*}
\operatorname{tr}U(\partial p)=\operatorname{tr}\bigl(A_b h_y W_b^{\dagger}h_x^{-1}\bigr)
=\operatorname{tr}\bigl(h_x^{-1}A_b h_y W_b^{\dagger}\bigr).
\end{equation*}
Put $A^{h}_b:=h_x^{-1}A_bh_y$ and
\begin{equation*}
k(A,W):=\prod_{b\in\Lambda_{\times}}\exp\bigl(\beta\operatorname{Re}\operatorname{tr}(A_bW_b^{\dagger})\bigr),
\end{equation*}
so that the crossing plaquettes contribute the factor $e^{-\beta n_{\times}}k(A^{h},W)$ to $e^{-S}$.
For any matrices $M$, $W$,
\begin{equation*}
\operatorname{Re}\operatorname{tr}(MW^{\dagger})
=\frac12\sum_{i,l}\bigl(M_{il}\overline{W_{il}}+\overline{M_{il}}\,W_{il}\bigr),
\end{equation*}
a sum of terms $P(M)\overline{P(W)}$ with $P$ a matrix entry or its conjugate, each with the coefficient
$\frac12\ge 0$ (multiplied by the positive normalisation of the trace). Expanding the exponential series
of $\beta$ times such a sum and multiplying over $b$ gives
\begin{equation}\label{8.7:eq-chessboard-torus-4}
k(A,W)=\sum_{\alpha}c_{\alpha}P_{\alpha}(A)\overline{P_{\alpha}(W)},\qquad c_{\alpha}\ge 0,
\end{equation}
with $P_{\alpha}$ monomials in the entries of the $A_b$ and their conjugates. The series converges
absolutely and uniformly, since the entries of elements of the compact group $G$ are bounded.

For $g=(g_x)_x$ indexed like $h$, write $(hg)_x:=h_xg_x$ and $W^{g}_b:=g_x^{-1}W_bg_y$ for
$b=(x,y)\in\Lambda_{\times}$. Then $A^{hg}_b=g_x^{-1}A^{h}_bg_y$, and by unitarity and cyclicity of the
trace
\begin{equation*}
\operatorname{tr}\bigl(A^{hg}_b(W^{g}_b)^{\dagger}\bigr)
=\operatorname{tr}\bigl(g_x^{-1}A^{h}_bg_yg_y^{-1}W_b^{\dagger}g_x\bigr)
=\operatorname{tr}\bigl(A^{h}_bW_b^{\dagger}\bigr),
\end{equation*}
so $k(A^{hg},W^{g})=k(A^{h},W)$ for every $g$. Let $dh$ and $dg$ be products of normalised Haar measures
and $K(A,W):=\int dh\,k(A^{h},W)$. Integrating the last identity over $h$ and using the invariance of
$dh$ under $h\mapsto hg$ gives $K(A,W)=K(A,W^{g})$ for every $g$, hence
\begin{equation*}
K(A,W)=\int dg\,K(A,W^{g})=\int dh\int dg\,k(A^{h},W^{g})
=\sum_{\alpha}c_{\alpha}\varphi_{\alpha}(A)\overline{\varphi_{\alpha}(W)},
\end{equation*}
where $\varphi_{\alpha}(A):=\int dh\,P_{\alpha}(A^{h})$; the last step inserts
\eqref{8.7:eq-chessboard-torus-4} with $A^{h}$, $W^{g}$ in place of $A$, $W$, interchanges sum and
integrals by uniform convergence, and uses that $W^{g}$ is formed from $W$ as $A^{h}$ is from $A$.

Now write $V:=U\restriction(\Lambda_+\setminus\Lambda_0)$, so that $F(U)=F(V,U_0)$ and $A$ is a part of
$V$ (the bonds of $\Lambda_{\times}$ do not lie in $\Lambda_0$). Change the variables on
$\vartheta(\Lambda_+\setminus\Lambda_0)$ to $W:=(\vartheta^{*}U)\restriction(\Lambda_+\setminus\Lambda_0)$;
this is a relabelling of bonds composed with inversion on the bonds whose orientation $\vartheta$
reverses, and the Haar measure is invariant under both, so the product of the $dU_b$ over the mirror
bonds becomes $dW$. Since $\vartheta$ fixes the ends of the bonds of $\Lambda_0$,
$(\vartheta^{*}U)\restriction\Lambda_+=(W,U_0)$, hence $\theta F=\overline{F(W,U_0)}$ and
$S_+\circ\vartheta^{*}=S_+(W,U_0)$. The variables $h$ on $C$ enter only through $k$. Therefore
\begin{align*}
\int F\,\theta F\,d\mu_{\mathrm{un}}
&=e^{-\beta n_{\times}}\int dU_0\,e^{-S_0(U_0)}\int dV\int dW\\
&\qquad\times F(V,U_0)e^{-S_+(V,U_0)}\,\overline{F(W,U_0)e^{-S_+(W,U_0)}}\,K(A,W)\\
&=e^{-\beta n_{\times}}\int dU_0\,e^{-S_0(U_0)}\sum_{\alpha}c_{\alpha}
\Bigl|\int dV\,F\,e^{-S_+}\varphi_{\alpha}\Bigr|^{2}\ \ge\ 0 ,
\end{align*}
using $c_{\alpha}\ge0$ and that $S_0$ and $S_+$ are real. All four combinations of types of the two fixed
hyperplanes are covered by the same computation: $\Lambda_0$ is empty (and the $U_0$-integral absent)
when both fixed hyperplanes bisect bonds, $C$ and $\Lambda_{\times}$ are empty (and $K\equiv 1$) when
both pass through sites, and for one hyperplane of each type both occur.

\emph{Proof of \textup{(b)}: one direction.} Fix a direction, say $\mu=1$, and write $c_i=c^1_i$,
$r=r_1$, $\vartheta_i=\vartheta^{(1)}_i$, $\Phi_l=\Phi^{(1)}_l$, and $\Delta_l$ for the closed slab
$x_1\in[c_l,c_{l+1}]$. The reflection $\vartheta_i$ has the fixed hyperplanes $x_1=c_i$ and
$x_1=c_i+n_1/2=c_{i+r}$, and it permutes the family of hyperplanes, $\vartheta_i(c_l)=c_{2i-l}$; in
particular it maps $\Delta_l$ onto $\Delta_{2i-l-1}$. Induction on $l$ from $\Phi_1=\vartheta_1$ gives, in the
first coordinate,
\begin{equation*}
\Phi_l(x)_1=x_1+\tfrac{ls}{2}\quad(l\text{ even}),\qquad
\Phi_l(x)_1=2c+\tfrac{(l+1)s}{2}-x_1\quad(l\text{ odd}),
\end{equation*}
since $\vartheta_l$ applied to $2c+ls/2-x_1$ gives $x_1+ls/2$, and applied to $x_1+(l-1)s/2$ gives
$2c+(l+1)s/2-x_1$. We use these formulas to define $\Phi_l$ for all $l\in\mathbb{Z}$; they are unchanged
under $l\mapsto l+2r$ because $rs=n_1$, and $\Phi_l$ maps $\Delta_0$ onto $\Delta_l$. Moreover
$2c_i-\Phi_l(x)_1$ equals $2c-x_1+(2i-l)s/2$ for $l$ even and $x_1+(2i-l-1)s/2$ for $l$ odd, that is,
\begin{equation}\label{8.7:eq-chessboard-torus-5}
\vartheta_i\circ\Phi_l=\Phi_{2i-l-1}.
\end{equation}

Let $a_0,\dots,a_{2r-1}$ be bounded nonnegative measurable functions of the bonds with both ends in the
slab $\Delta_0$ (indices in $\mathbb{Z}/2r\mathbb{Z}$), let $a^{[l]}(U):=a(\Phi_l^{*}U)$, and put
\begin{equation*}
\mathcal{G}(a_0,\dots,a_{2r-1}):=\Bigl\langle\prod_{l}a_l^{[l]}\Bigr\rangle .
\end{equation*}
By \eqref{8.7:eq-chessboard-torus-5}, for real $a$ we have $\theta_i(a^{[l]})=a^{[2i-l-1]}$, where
$\theta_i$ is the operation $\theta$ for the reflection $\vartheta_i$. Fix a cut $i$ and let $\Pi_+$ be
the product of the factors $a_l^{[l]}$ with $l=i,\dots,i+r-1$, the tiles in the arc $[c_i,c_{i+r}]$, and
$\Pi_-$ the product of the others; $\Pi_+$ and $\theta_i\Pi_-$ depend only on bonds with both ends in
$[c_i,c_{i+r}]$. Since $\vartheta_i$ is an involution, $\Pi_-=\theta_i(\theta_i\Pi_-)$, so
$\mathcal{G}(a)=\langle\Pi_+\,\theta_i(\theta_i\Pi_-)\rangle$. By \textup{(a)} (with $c$ replaced by
$c_i$, and after division by $\int d\mu_{\mathrm{un}}>0$) the sesquilinear form
$(F,F')\mapsto\langle F\,\theta_iF'\rangle$ is positive semidefinite on the functions of the bonds with
both ends in $[c_i,c_{i+r}]$, and the Cauchy--Schwarz inequality for it gives, for every cut $i$,
\begin{equation}\label{8.7:eq-chessboard-torus-6}
\mathcal{G}(a)^2\le\langle\Pi_+\theta_i\Pi_+\rangle\,\langle(\theta_i\Pi_-)\,\Pi_-\rangle
=\mathcal{G}(a')\,\mathcal{G}(a''),
\end{equation}
where $a'$ agrees with $a$ on the tiles $i,\dots,i+r-1$ and is continued by its mirror image,
$a'_{2i-1-l}:=a_l$ for those $l$, and $a''$ is formed in the same way from the half of $a$ on the other
side of the cut. By \textup{(a)}, $\mathcal{G}(a')\ge0$ and $\mathcal{G}(a'')\ge0$.

Let $\|a\|:=\mathcal{G}(a,\dots,a)^{1/2r}$. We show
$|\mathcal{G}(a_0,\dots,a_{2r-1})|\le\prod_l\|a_l\|$. Both sides are positively homogeneous of degree one
in each $a_l$. If some $\|a_l\|=0$, replace every $a_l$ by $a_l+\varepsilon$, for which
$\|a_l+\varepsilon\|\ge\varepsilon>0$, and let $\varepsilon\downarrow0$ at the end, both sides being
continuous in $\varepsilon$. So we may assume $\|a_l\|=1$ for all $l$. Among the finitely many sequences
with entries in $\{a_0,\dots,a_{2r-1}\}$ let $\mathsf{m}$ be the maximum of $|\mathcal{G}|$; then
$\mathsf{m}\ge\mathcal{G}(a_0,\dots,a_0)=1$. Choose a maximiser
$\tilde a=(\tilde a_0,\dots,\tilde a_{2r-1})$ whose longest cyclic run of equal entries has the largest
length $\ell$ among all maximisers, and suppose $\ell<2r$; let the run occupy
the places $j,\dots,j+\ell-1$. If $\ell\le r$, cut at $c_j$, an end of the run: the half
$j,\dots,j+r-1$ contains the run, and $\tilde a'$ has the run on the places $j-\ell,\dots,j+\ell-1$, of
length $2\ell$. If $\ell>r$, cut at $c_j$ again: the half $j,\dots,j+r-1$ lies inside the run and
$\tilde a'$ is constant. In both cases $\tilde a'$ has a run of length at least $\min(2\ell,2r)>\ell$,
and by \eqref{8.7:eq-chessboard-torus-6}
\begin{equation*}
\mathsf{m}^2=|\mathcal{G}(\tilde a)|^2\le\mathcal{G}(\tilde a')\,\mathcal{G}(\tilde a'')
\le|\mathcal{G}(\tilde a')|\,\mathsf{m},
\end{equation*}
so $|\mathcal{G}(\tilde a')|\ge\mathsf{m}$ and $\tilde a'$, whose entries lie in the same set, is again a
maximiser, contradicting the choice of $\tilde a$. Hence $\ell=2r$, a maximiser is constant, equal to
$(a_k,\dots,a_k)$ for some $k$, and $\mathsf{m}=\mathcal{G}(a_k,\dots,a_k)=\|a_k\|^{2r}=1$. This proves
$|\mathcal{G}(a_0,\dots,a_{2r-1})|\le\prod_l\|a_l\|$. With
$a_0=F$ and $a_l=1$ for $l\ne0$, and $\|1\|=1$,
\begin{equation*}
\langle F\rangle=\mathcal{G}(F,1,\dots,1)\le\|F\|=\Bigl\langle\prod_{l}F^{[l]}\Bigr\rangle^{1/2r_1}.
\end{equation*}

\emph{Proof of \textup{(b)}: all directions.} The argument above uses only the first coordinate, and
\textup{(a)} holds with any coordinate in the role of $x_1$; so the one-direction inequality holds in
every direction $\mu$ for bounded nonnegative functions of the bonds with both ends in the slab
$x_{\mu}\in[c^{\mu}_0,c^{\mu}_1]$. The product $\Pi^{(1)}:=\prod_{l}F^{[(l,0,\dots,0)]}$ is such a
function for $\mu=2$: it depends on the bonds with both ends in the closed slab
$x_{\nu}\in[c^{\nu}_0,c^{\nu}_1]$, $\nu\ge2$. The maps $\Phi^{(\mu)}_l$ for different $\mu$ act on
different coordinates and commute, so the direction-two copies of $\Pi^{(1)}$ are products of the
$F^{[m]}$ with $m=(m_1,m_2,0,\dots,0)$. Iterating over $\mu=2,\dots,d$ and using the monotonicity of
$t\mapsto t^{1/2r_{\mu}}$,
\begin{equation*}
\langle F\rangle\le\langle\Pi^{(1)}\rangle^{1/2r_1}\le\dots\le
\Bigl\langle\prod_{m}F^{[m]}\Bigr\rangle^{1/\prod_{\mu}2r_{\mu}},
\end{equation*}
which is \eqref{8.7:eq-chessboard-torus-2}.

\emph{Proof of \textup{(c)}.} A closed tile $[c,c+s/2]$ with $c\in\mathbb{Z}$ contains the sites
$c,\dots,c+\lfloor s/2\rfloor$; so, choosing each $c^{\mu}$ as the least coordinate of the box in
direction $\mu$, the box lies in the closed tile $\Delta_0$, and $F:=\mathbf{1}_A$ satisfies the
hypotheses of \textup{(b)}. In each direction $\vartheta_{i+1}\vartheta_i$ is the translation by $s$,
since $2c_{i+1}-(2c_i-x_{\mu})=x_{\mu}+s$; hence for a tile $m$ all of whose indices are even,
$\Phi_m$ is the translation $x\mapsto x+a(m)$ with $a(m):=s\,m/2$, and $F^{[m]}=\mathbf{1}_{A+a(m)}$. As
$m$ runs over the tiles with all indices even, $a(m)$ runs over $(s\mathbb{Z})^d$ modulo the periods,
each class once. The factors $F^{[m]}$ of the other tiles take values in $[0,1]$ and are dropped, so by
\eqref{8.7:eq-chessboard-torus-2}
\begin{equation*}
\mu(A)=\langle\mathbf{1}_A\rangle\le\Bigl\langle\prod_{m}F^{[m]}\Bigr\rangle^{1/n_c}
\le\mu\Bigl(\bigcap_{a\in(s\mathbb{Z})^d/\sim}(A+a)\Bigr)^{1/n_c},
\end{equation*}
with $n_c=\prod_{\nu}2r_{\nu}=\prod_{\nu}2n_{\nu}/s$.

For one type of fixed hyperplane, part~(a) is the reflection positivity of the Wilson measure in
\cite{OSe78} and \cite{Sei82}, and part~(b) is the chessboard estimate of \cite{FILS78}; the proof
above uses in addition only that the two fixed hyperplanes of a reflection need not be of the
same type. Part~(c) keeps only the translates by $s$, twice the side of a tile. When $s$ is odd, for
instance a power of the odd blocking factor $L$, the side $s/2$ of a tile is a half-integer and the
cutting hyperplanes $x_{\mu}=c^{\mu}+i\,s/2$, $c^{\mu}\in\mathbb{Z}$, alternate between hyperplanes
through sites ($i$ even) and hyperplanes bisecting bonds ($i$ odd).
\end{proof}

\begin{lemma}[The renormalisation operation as a resampling kernel. Stated; proof incomplete at the step marked $(\ast)$]
\label{8.7:lem-resampling-kernel}
Let $\Omega$ be the set of label histories of the run, let $\tau_{\omega}>0$ be the mass of the history
$\omega\in\Omega$ and $Z=\sum_{\omega}\tau_{\omega}$, and let $\tilde{\nu}=\tau/Z$, as in
Definition~\ref{8.7:def-creation-live-events}. Let $\mu$ be the normalised Wilson measure at the bare
coupling $g_0$ (Definition~\ref{1.1:def-wilson-action}).
\begin{itemize}
\item[(a)] Each operation that makes up the composite $\mathbf{R}T$ of Ba{\l}aban's step at level $k$
is a Markov kernel on pairs (labels, fields) whose total mass is one at every value of its argument.
These operations are: the averaging $T$, the sharp decompositions of unity, the
Faddeev--Popov insertions, the operation $\mathbf{R}'$ of \cite{B-CMP122a}, and the closing operation
of \cite{B-CMP122b}. Consider a term of the large-field expansion of $\rho_k$ on which $\mathbf{R}'$
acts, with class components $X_1,\dots,X_n$ and, for each $i$, the part $\Lambda_i$ of the
integration domain attached to $X_i$, as in \cite[(1.100)]{B-CMP122a}; let $V_{\Lambda_i}$ be the
minimiser of \cite[Proposition~1]{B-CMP122a}. On this term the operation $\mathbf{R}'$
discards $V_k\restriction\Lambda_i$ and draws it afresh as $V'V_{\Lambda_i}$. Here $V'$ is
distributed by a normalised density that is a function of $V_k\restriction(X_i\setminus\Lambda_i)$
only.
\item[(b)] The same holds for the modified large-field operation used in the run of the correlation
theorem.
\item[(c)] Consequently $\tilde{\nu}$ is the law of the label history of a Markov chain in the level
index, started from $V_0=U$ distributed by $\mu$. Write $P$ for the law of this chain. Then for every
$E\subset\Omega$ the function $G_E(U):=P(E\mid V_0=U)$ takes values in $[0,1]$, and
\begin{equation}\label{8.7:eq-resampling-kernel-1}
\tilde{\nu}(E)=\int d\mu(U)\,G_E(U).
\end{equation}
\end{itemize}
\end{lemma}

Stated; the proof below is incomplete at the step marked $(\ast)$.

\begin{proof}
The proof is given for the operations as Ba{\l}aban defines them, and then for the modified operation.

\emph{The transformations preceding the integration.} \cite{B-CMP122a} states:
``All the above transformations preserve the $k$th density $\rho_k$, they change only the representation
of this density''. These transformations are the sharp decompositions of unity and the Faddeev--Popov
insertions; their Markov-kernel property is established in the paragraph on the remaining operations
below.

\emph{The integration over $\Lambda$.} \cite{B-CMP122a} continues:
``Now we define the next operation, which changes the density \dots\ The next operation is an
integration of this term with respect to the field variables $V_k$ over the domain $\Lambda$. We obtain
a new expression, which we consider as a function of all field variables $V_k$, but independent of
$V_k\restriction\Lambda$''. The explicit form is \cite[(1.100), p.~201]{B-CMP122a}. Let $X_i$ be
a class component of the term and $\Lambda_i$ the part of the integration domain attached to it. Then
$X_i$ carries the factor
\begin{equation}\label{8.7:eq-resampling-kernel-2}
\frac{1}{N_i}\sum_{\dots,\,G_i,\,T_i}\chi_{k,\Lambda_i}
\cdot\frac{\delta_{G_i}(V'_k)\,\chi(\Lambda_i)\,
\exp\bigl[-g_k^{-2}A\bigl(\zeta_i,U_{k,X_i}(V'_kV_{\Lambda_i})\bigr)\bigr]}
{\int d(V'\restriction\Lambda_i)\,\delta_{G_i}(V')\,\chi(\Lambda_i)\,
\exp\bigl[-g_k^{-2}A\bigl(\zeta_i,U_{k,X_i}(V'V_{\Lambda_i})\bigr)\bigr]}.
\end{equation}
This factor multiplies the expression integrated over $\Lambda_i$. In this display:
\begin{itemize}
\item the sum runs over the label data of the component $X_i$ in \cite[(1.100)]{B-CMP122a}, among
them $G_i$ and $T_i$; $N_i^{-1}$ is the numerical prefactor of \cite[(1.100)]{B-CMP122a},
$\delta_{G_i}$ is the
gauge-fixing $\delta$-function attached to $G_i$, and $\chi_{k,\Lambda_i}$ and $\chi(\Lambda_i)$ are
the characteristic functions of that display;
\item $\zeta_i$ is the weight in the action of that display, not a complex source;
\item $V_{\Lambda_i}$ is the minimiser of \cite[Proposition~1]{B-CMP122a} determined by the data on
the part of the large-field region outside $\Lambda$;
\item $U_{k,X_i}$ is the background configuration local to $X_i$.
\end{itemize}
Moreover \cite[(1.102)]{B-CMP122a} states
\begin{equation}\label{8.7:eq-resampling-kernel-3}
\int dV_k\,(\mathbf{R}'\rho_k)=\int dV_k\,\rho_k .
\end{equation}
We read \eqref{8.7:eq-resampling-kernel-2} as follows. The integrand of the denominator is
non-negative, being a product of a $\delta$-function, a characteristic function and an exponential,
and its integral over $V'\restriction\Lambda_i$ is strictly positive, the domain cut out by
$\delta_{G_i}\chi(\Lambda_i)$ being non-empty \cite{B-CMP122a}. The quotient has denominator equal to
the integral of its numerator over $V'\restriction\Lambda_i$. It is therefore a probability density in
$V'$. It is a function of $V_k\restriction(X_i\setminus\Lambda_i)$ only. Thus, on the term with class
components $X_1,\dots,X_n$, the operation $\mathbf{R}'$ does two things. It forgets
$V_k\restriction\Lambda_i$,
as the quoted sentence says. It then puts back $V'V_{\Lambda_i}$, with $V'$ drawn from this normalised
density. For each fixed value of the remaining variables this is a Markov kernel of total mass one.
Integrating over those variables gives \eqref{8.7:eq-resampling-kernel-3}.

\emph{The closing operation.} \cite{B-CMP122b} removes the factor $\delta_{G\cap X}(V'_k)$ ``by a
reversed Faddeev--Popov procedure \dots\ equivalent, but not equal''. In the reading just given, this
operation applies inside $X_i$ a gauge transformation distributed by Haar measure. That is a Markov
kernel of total mass one.

\emph{The remaining operations.}
\begin{itemize}
\item The averaging $T$ is a $\delta$-kernel: the block variables are a deterministic function of the
finer ones.
\item In the sharp decompositions of unity, each label is a deterministic function of the current
variables. Each such decomposition is therefore again a $\delta$-kernel, now in the label variable.
\item The Faddeev--Popov insertions of \cite[(1.6)]{B-CMP122a} are normalised one-site densities, that
is, gauge transformations drawn at random.
\end{itemize}
Hence every operation of $\mathbf{R}T$ is a Markov kernel on (labels, fields) that conserves mass
pointwise in its argument. This is the content of the lemma on conservation of total mass under
$\mathbf{R}T$, read fibre by fibre; \eqref{8.7:eq-resampling-kernel-3}, stated in
\cite[(1.102)]{B-CMP122a}, is its instance for $\mathbf{R}'$.
This proves (a).

$(\ast)$ The run of the correlation theorem uses a modified large-field operation. That its factors
have, component by component, the form \eqref{8.7:eq-resampling-kernel-2} is not verified here;
granted this, the preceding paragraphs apply without change and give (b).

\emph{The chain.} Compose the kernels of (a) and (b) level by level, starting from $V_0=U$ distributed
by $\mu$. The result is a Markov chain in the level index on (labels, fields). The bare density is the
unnormalised Wilson weight, which is positive, and all kernels of (a) and (b) have non-negative
weights. The mass $\tau_{\omega}$ of a history is the integral of the part of the
density carrying the labels of $\omega$. Since every kernel has mass one at every argument, $Z$ is the
normalising constant of the unnormalised Wilson weight, that is, of $\mu$, and $\tau_{\omega}/Z$ is the
probability that the chain runs through the
history $\omega$. Summing over $\omega\in E$ gives \eqref{8.7:eq-resampling-kernel-1}, with
$G_E(U)=P(E\mid V_0=U)\in[0,1]$. This proves (c).
\end{proof}

We note what the lemma does and does not give for label events.

A label event $E$ that lies past an $\mathbf{R}$-step is not an event of the bare field. An inequality
$\tilde{\nu}(E)\le\mu(F)$ with a bare-field event $F$ therefore does not follow from the lemma. What
does follow is that the chessboard estimate, Lemma~\ref{8.7:lem-chessboard-torus}(b), can be applied
to the function $G_E$ of \eqref{8.7:eq-resampling-kernel-1}. Let $E$ be attached to a cube $\square$
(for instance $E=E^{\kappa}_{j,\square}$ of Definition~\ref{8.7:def-creation-live-events}), let
$\square_a$, $1\le a\le n_c$, be its translates over the $n_c$ tiles of
Lemma~\ref{8.7:lem-chessboard-torus}, and let $E_a$ be the corresponding translates of $E$. Suppose,
in addition, that:
\begin{itemize}
\item $G_E$ is measurable with respect to the variables of a single tile, and
\item the translates $E_a$ use disjoint resampling variables.
\end{itemize}
Then $\prod_a G_{E_a}=P(\bigcap_a E_a\mid U)$, and
\begin{equation}\label{8.7:eq-resampling-kernel-4}
\tilde{\nu}(E)\le\tilde{\nu}\Bigl(\bigcap_a E_a\Bigr)^{1/n_c}.
\end{equation}
Both members of \eqref{8.7:eq-resampling-kernel-4} are probabilities, under the same measure
$\tilde{\nu}$, of events of the same run. On $\bigcap_a E_a$ the event occurs at every translate
$\square_a$.

Take now $E=E^{\kappa}_{j,\square}$. For the right-hand side, stop the run at level $j$. Take a seed in
$\mathcal{Z}^{\mathrm{live}}$ of Definition~\ref{8.7:def-terminal-reserve} containing $n'$ of the cubes
$\square_a$, pairwise at distance at least the tile period $s$ of
Lemma~\ref{8.7:lem-chessboard-torus} (not the depth) less a fixed length. Its linear size in
cubes, $d'_j$ of Definition~\ref{8.7:def-terminal-reserve}, satisfies $d'_j+1\ge n'$. On
$\bigcap_aE_a$, with the run stopped at level $j$, each $\square_a$ has been created at the terminal
level $j$, is therefore live there, and lies in one such seed, created at step $j$. Summing the
summands of the reserve \eqref{8.7:eq-terminal-reserve-a} of Definition~\ref{8.7:def-terminal-reserve}
over these seeds, whose counts $n'$ add up to $n_c$, gives
\begin{equation}\label{8.7:eq-resampling-kernel-5}
\inf_{\omega\in\bigcap_a E_a}\Theta^{\mathrm{live}}(\omega)\ge t'_j\,n_c,\qquad
t'_j:=\theta'\gamma_0A_1\,p_0(g_j)^2,
\end{equation}
where $t'_j$ is the per-seed coefficient of the summands of $\Theta^{\mathrm{live}}$. This covers the
kinds whose seeds are counted in $\Theta^{\mathrm{live}}$. For the remaining kinds $\kappa$ of
Definition~\ref{8.7:def-creation-live-events}, those created by the preparatory operations, the
translates lie in distinct class components, and each contributes a reserve equal to $\theta$ times
the degree function $p_0(\cdot)$, with no factor that grows with the size.

If the two hypotheses above held, \eqref{8.7:eq-resampling-kernel-4},
\eqref{8.7:eq-resampling-kernel-5} and the global bound \eqref{8.7:eq-terminal-reserve-b} of
Definition~\ref{8.7:def-terminal-reserve} would give, at depth zero (the letter $s$ denoting here, as
in Lemma~\ref{8.7:lem-chessboard-torus}, the tile period of that lemma and not the depth below the
cutoff),
\begin{equation}\label{8.7:eq-resampling-kernel-6}
\tilde{\nu}(E^{\kappa}_{j,\square})
\le\exp\bigl(-2^{-d}\bigl[t'_j-\mathsf{B}(x_j)\,s^{d}\bigr]\bigr).
\end{equation}
Requiring the bracket to be large would impose an upper bound on $\gamma$. That bound comes from
comparing the degree $2p_0$ of $t'_j$ in $\log g_j^{-2}$ with $d\,r_{\mathrm{pr}}$ plus the degree of
$\mathsf{B}$. The measurability hypothesis above is not established here, and
\eqref{8.7:eq-resampling-kernel-6} is not used below.

\begin{remark}[The fluctuation test reads the global background field]\label{8.7:rem-nonlocal-test}
We work in Ba{\l}aban's renormalisation step as it is set up in \cite[\S3]{B-CMP119}, and use its
letters: $\Gamma_{k+1}$ and $\Lambda_{k+1}$ are domains of that step, $\delta_k$ is the threshold
entering its characteristic functions, and $R_k$ is the number in the weight $e^{-R_k}$ of
\cite[(3.18)--(3.19)]{B-CMP119}. The characteristic functions of the second kind described below
determine at level $k$ a region, which \cite[(3.20)]{B-CMP119} denotes by the same letter with the
index $k+1$. We call it the \emph{region of the second kind at level $k$} and reserve the letter
$R_k$ for the number. By \cite[(3.20)]{B-CMP119} this region enters $\Lambda_{k+1}$.
Further, $A_k(b)$ is the
variable of the fluctuation field on the bond $b$. Two kinds of characteristic functions are tested
at a cube $\square'$.
\begin{itemize}
\item Functions of the \emph{first kind}, which are local in the pair $(V_k,V_{k+1})$ through the
configuration $U_{k+1,\square'}$ that \cite[\S3]{B-CMP119} attaches to the cube $\square'$.
\item Functions of the \emph{second kind}, namely the characteristic functions
$\chi(\{|A_k(b)|<\delta_k\})$ of \cite[(3.16)]{B-CMP119}, which are local in $A_k$.
\end{itemize}
A \emph{label history} $\omega\in\Omega$ records which term of each decomposition of unity was chosen
at each cube and each step, and a \emph{label event} is a set of label histories. For a label event
$E$, $G_E(U)$ is the conditional probability of $E$ given the bare field $U$. Here $\mu$ is the Wilson
measure, and $\tilde{\nu}$ is the normalised measure on label histories. A function of $U$ is
\emph{tile-measurable} if it depends only on the bare bond variables in one tile, that is, in one
cell of the periodic decomposition of the torus used in the chessboard estimate of
Lemma~\ref{8.7:lem-chessboard-torus}.

The following statements hold.
\begin{enumerate}
\item[(a)] Regarded as a function of the sequence of fields
$V_{k+1}\restriction\Gamma_{k+1},V_k\restriction\Gamma_k,\dots$ of the step, the characteristic
function of the second kind at $\square'$ depends on $V_{k+1}$ on all of $\Gamma_{k+1}$. The
non-local part of this dependence enters the tested variable $A_k(b)$ with weight at most
$O(1)e^{-R_k}\delta_k$ \cite[(3.18)--(3.19)]{B-CMP119}. Hence at most one
of the two kinds is a strictly local function in any one set of variables.
\item[(b)] A characteristic function of the second kind is not a local function of $V_{k+1}$; in
particular, the locality hypothesis on the characteristic functions $\chi^{\kappa}_{j,\square}$
fails, as stated, for $\kappa$ the second kind.
\item[(c)] Every label event of a level later than $k$ is defined through tests that inherit this
dependence on $V_{k+1}$. Hence the construction does not present $G_E$, for such an event $E$, as a
tile-measurable function, and for no such event other than $\emptyset$ and $\Omega$ does
tile-measurability of $G_E$ follow from the construction.
\item[(d)] The thresholds are sharp, so the smallness of the non-local weight in (a) does not by
itself give an approximation of $G_E$ in the supremum norm by a tile-measurable function. The test
of the second kind can be bracketed from both sides by local
tests with the thresholds $\delta_k(1\mp e^{-R_k})$, but this leaves a window of undetermined outcome
at every bond of every finer level under the tile. Such an approximation requires a bound on the
$\mu$-mass of that window, that is, an
anti-concentration bound, which is a lower-regularity statement about $\tilde{\nu}$ that is not
proved; otherwise a renormalisation step whose test of the second kind is local is needed.
\end{enumerate}
In particular, an argument for rarity of young large-field blocks at levels of small depth $s$ below
the cutoff cannot rest on tile-measurability of the conditional probabilities $G_E$, by (c) and (d);
such an argument is needed only at those levels.
\end{remark}

\begin{proof}
\emph{The background field is global.} In \cite[(3.11)]{B-CMP119} the fluctuation field is defined on
the set $B(\Gamma_{k+1})$ of \cite[(3.11)]{B-CMP119} by
\begin{equation*}
V'_k:=V_k\bigl(V^{(k)}\bigr)^{-1}.
\end{equation*}
The background field $V^{(k)}$ entering this definition is described at
\cite[(3.10)--(3.11)]{B-CMP119} as follows (the sets $\Omega_j$ in this sentence are Ba{\l}aban's
domains, not the set of label histories, and $M^k$ is the map of \cite[\S3]{B-CMP119} that takes the
minimal configuration $U_{k+1}$ to the background field $V^{(k)}$, not a power of the constant $M$):
\begin{quote}
$V^{(k)}=M^k(U_{k+1})$, where the minimal configuration $U_{k+1}$ is determined by the sequence of
domains $\{\Omega_{k+1},\Omega_k,\dots\}$, and the corresponding sequence of fields
$\{V_{k+1}\restriction\Gamma_{k+1},V_k\restriction\Gamma_k,\dots\}$
\end{quote}
Thus $V^{(k)}$ is determined by the whole sequence of domains and fields. It is a global background,
not a function of the fields near one cube.

\emph{The second decomposition of unity.} After \cite[(3.15)]{B-CMP119}, where $\chi'_k$ denotes the
characteristic functions of the first decomposition of unity, Ba{\l}aban writes:
\begin{quote}
the characteristic functions $\chi'_k$ depend now on the field $V^{(k)}$, hence on the new field
variables $V_{k+1}$ on the whole set $\Gamma_{k+1}$. We want to have characteristic functions
depending on the fluctuation field only, hence we introduce another decomposition of unity
\end{quote}
This further decomposition is \cite[(3.16)]{B-CMP119}. Its characteristic functions are
$\chi(\{|A_k(b)|<\delta_k\})$, which are the functions of the second kind. Later in the same section,
of the characteristic function written in the fluctuation variables, \cite[\S3]{B-CMP119} states:
\begin{quote}
The last function above depends nonlocally on $V_{k+1}$, through the function $V^{(k)}$
\end{quote}

\emph{Proof of (a).} A function of the first kind is local in $(V_k,V_{k+1})$ through the
configuration $U_{k+1,\square'}$ that \cite[\S3]{B-CMP119} attaches to $\square'$. A function of the
second kind is local in $A_k$, that is, in the fluctuation field
$V'_k$. By the displayed definition, $V'_k$ involves $V^{(k)}$, and by the first quoted sentence
$V^{(k)}$ involves $V_{k+1}$ on all of $\Gamma_{k+1}$. Hence, as a function of the sequence of fields
$V_{k+1}\restriction\Gamma_{k+1},V_k\restriction\Gamma_k,\dots$, the
characteristic function of the second kind at $\square'$ reads $V_{k+1}$ on all of $\Gamma_{k+1}$;
this is the content of the third quoted sentence. By \cite[(3.18)--(3.19)]{B-CMP119}, the non-local
part of this dependence enters the tested variable $A_k(b)$ with weight at most
$O(1)e^{-R_k}\delta_k$. Consequently no single set of variables makes both kinds strictly local: in
the pair $(V_k,V_{k+1})$ the second kind is not, as just shown; and in the fluctuation variables the
first kind is not, since by the second quoted sentence the characteristic functions $\chi'_k$ depend
on $V^{(k)}$, hence on $V_{k+1}$ on all of $\Gamma_{k+1}$, and not on the fluctuation field alone.

\emph{Proof of (b).} By (a), a characteristic function of the second kind at $\square'$ depends on
$V_{k+1}$ on all of $\Gamma_{k+1}$, so it is not a local function of $V_{k+1}$. This is the failure of
the locality hypothesis on the characteristic functions $\chi^{\kappa}_{j,\square}$ for $\kappa$ the
second kind.

\emph{Proof of (c).} The region of the second kind at level $k$ is determined by the characteristic
functions of the second kind at level $k$, so by (a) it depends on $V_{k+1}$ on all of
$\Gamma_{k+1}$; by \cite[(3.20)]{B-CMP119} it enters the domain
$\Lambda_{k+1}$. Through $\Lambda_{k+1}$ it enters every domain tested at a later level, every
condition sorting the later domains into the classes used by the later label events, and every later
resampling step, the renormalisation
operation being read as a resampling kernel in Lemma~\ref{8.7:lem-resampling-kernel}, whose proof is
incomplete at one step. Each of these therefore depends on the choices of the second kind made at
level $k$, and by (a) these choices depend on $V_{k+1}$ on all of $\Gamma_{k+1}$. Every label event of
a later level is defined through these domains, classes and resampling steps, so its defining tests
inherit the dependence, and the expression for $G_E(U)$ given by the construction involves the bare
variables outside the tile. Whether $G_E$ is nevertheless constant on the set of bare fields that
agree in the tile depends on whether the non-local part of a tested variable moves it across its
threshold with positive conditional probability, that is, on the conditional law of the tested
variables near their thresholds, the information discussed in (d). For $E=\emptyset$ and
$E=\Omega$ the function $G_E$ is constant; for no other later label event does the construction give
tile-measurability.

\emph{Proof of (d).} The characteristic functions $\chi(\{|A_k(b)|<\delta_k\})$ have sharp thresholds:
a shift of the tested variable, however small, can change the value of the test. Hence the smallness
of the non-local part of $A_k(b)$ does not by itself give a tile-measurable function
$G_E^{\mathrm{loc}}$ close to $G_E$ in the supremum norm. By (a), the non-local part of the tested
variable $A_k(b)$ has size at most
$O(1)e^{-R_k}\delta_k$. Replace the tested variable by its part that depends only on the fields under
the tile, and the threshold $\delta_k$ by $\delta_k(1-e^{-R_k})$ and by $\delta_k(1+e^{-R_k})$. The
first of these two local tests implies the test of the second kind, which in turn implies the second
local test, so the test of the second kind is bracketed from both sides. The two brackets
disagree on a window, which occurs
at every bond of every finer level under the tile. The event $E$ and the events obtained from it by
replacing each test of the second kind by either bracket differ only on histories in which some tested
variable lies in its window; so an approximation of $G_E$ in the supremum norm by the conditional
probabilities with local tests amounts to a bound on the conditional probability of the window.
Controlling the $\mu$-mass of this window needs an
anti-concentration bound for the tested variable near its threshold, that is, a lower-regularity
statement about $\tilde{\nu}$. The alternative is a renormalisation step whose test of the second kind
is local.
The second quoted sentence shows why \cite[(3.16)]{B-CMP119} is introduced: the fluctuation integral
requires a restriction that, as a function of the fluctuation field, does not involve $V_{k+1}$. The
test written in $A_k$ provides exactly that restriction, and (a) shows that the same test, read as a
function of the sequence of fields $V_{k+1}\restriction\Gamma_{k+1},V_k\restriction\Gamma_k,\dots$,
is not local in $V_{k+1}$.

The final assertion of the remark follows from (c) and (d). By (c) the construction does not give
tile-measurability of $G_E$; by (d) an approximation of $G_E$ by a tile-measurable function requires
the anti-concentration bound, which is not proved, and the renormalisation step used has a test of the
second kind that is not local.
\end{proof}

\begin{definition}[The annealed history measure and the large-field events]\label{8.8:not-history-measure}
\emph{The measure.} Ba{\l}aban's procedure $(\mathbf{R}T)^k$, with $T$ his renormalisation
transformation, the large-field operation
$\mathbf{R}$ and the effective densities $\rho_k$ of Definition~\ref{1.5:def-densities}, is run as
a Markov chain started from a bare configuration $U$ distributed according to
$\rho_0\,dU/\int\rho_0$, and every outcome of every decomposition of unity introduced along the way
is kept. A \emph{history} $\omega$ is the resulting sequence of outcomes, with no resummation;
$\Omega$ denotes the finite set of histories, and $\tau_\omega$ the total integral of the term of
$\omega$. By the positivity of the history terms, $\tau_\omega\ge 0$; by the conservation of the
total integral along the chain,
\begin{equation}\label{8.8:eq-history-measure-1}
  Z:=\sum_{\omega\in\Omega}\tau_\omega=\int\rho_0\,dU=\int\rho_k\,dV_k
  \qquad\text{for every }k .
\end{equation}
The \emph{annealed history measure} is defined, for $E\subset\Omega$, by
\begin{equation}\label{8.8:eq-history-measure-2}
  \tilde\nu(E):=Z^{-1}\sum_{\omega\in E}\tau_\omega .
\end{equation}
Jointly with the top field, we write $\tilde\nu_k(E\times dV)$ for the joint law of the history
$\omega\in E$ and the field $V_k$ of the chain after $k$ steps, whose marginal on $\Omega$ is
$\tilde\nu$.

The operation $\mathbf{R}$ re-draws the restriction $V_k\restriction\Lambda$ of the field to the
region $\Lambda$ on which it acts from a normalised reference density
\cite[(1.100)--(1.102)]{B-CMP122a}. Accordingly, under $\tilde\nu$ the field $V_k$ at level $k$ is
the chain's field, and not the $k$-fold block average of $U$ in the sense of
Definition~\ref{1.5:def-block-average}. Nothing in what follows requires the two to coincide.

\emph{Three functionals.} Fix a level $k$. Let $\square$ be a cube of side $LM_2R_k$ of the
partition of the level-$k$ lattice $\mathbb{T}^{(k)}$ of Definition~\ref{1.5:def-block-average},
and $\square'\supset\square$ the cube of side $LM_2R_{k+1}$ of the
next partition; $\square^{\sim}$, $\square^{\sim4}$, $\square'^{\sim2}$ are the layered enlargements
of \cite{B-CMP119}.
Put $\varepsilon_k=g_k\,p_0(g_k)$, with the degree function $p_0(g)=A_0(\log g^{-2})^{p_0}$, and
$\delta_k=(A_1/A_0)\varepsilon_k$; let $\eta$ be the fine spacing in level-$k$ units. Following
\cite[(2.16)--(2.17)]{B-CMP119}, with $U_{k,\square}$ the local minimiser,
\begin{equation}\label{8.8:eq-history-measure-3}
  a_k(\square):=(\varepsilon_k\eta^2)^{-1}\sup_{p\subset\square^{\sim}}
  \bigl|U_{k,\square}(V_k,\partial p)-1\bigr| ,
\end{equation}
and $a_k(\square):=+\infty$ where $V_k\restriction\square^{\sim4}$ is not regular enough for
$U_{k,\square}$ to exist. Following \cite[(3.3)--(3.4)]{B-CMP119}, in the axial gauge, with the
background $V^{(k)}_{\square'}:=M^k\bigl(U_{k+1,\square'}(V_{k+1})\bigr)$, in which $M^k$ is the
map of the cited displays and not a power of the constant $M$,
\begin{equation}\label{8.8:eq-history-measure-4}
  \varphi_k(\square'):=(2\delta_k)^{-1}\sup_{b\in(\square'^{\sim2})^{(k)*}}
  \bigl|V_k(b)\bigl(V^{(k)}_{\square'}(b)\bigr)^{-1}-1\bigr| ,
\end{equation}
where $(\square'^{\sim2})^{(k)*}$ denotes the set of bonds $b$ of the level-$k$ lattice contained in
$\square'^{\sim2}$; and, following \cite[(3.10)--(3.11), (3.15)--(3.16)]{B-CMP119}, with $A_k$ the
scaled fluctuation variable about the global background,
\begin{equation}\label{8.8:eq-history-measure-5}
  \psi_k(\square'):=(\delta_k/g_k)^{-1}\sup_{b\in(\square'^{\sim2})^{(k)*}}|A_k(b)| .
\end{equation}

\emph{Kinds of large-field creation.} The kind $\mathrm{S}1$ at step $k$ is the condition
$a_k(\square)\ge1$; the kind $\mathrm{S}2$ at step $k+1$ is $\varphi_k(\square')\ge1$; the kind
$\mathrm{S}3$ at step $k+1$ is $\psi_k(\square')\ge1$; each is imposed only on the cubes on which
Ba{\l}aban's construction tests it. The kinds $\mathrm{S}4,\dots,\mathrm{S}10$ are those defined by
the characteristic functions of
\cite[(1.22)--(1.24), (1.27), (1.28), (1.55), (1.75), (1.82), (1.88)]{B-CMP122a}.

\emph{Events.} For a kind $\kappa$ of large-field creation among those listed above (as a
superscript on $E$ and $\chi$ the letter $\kappa$ indexes kinds and is not the auxiliary parameter
$\kappa$), let
$\chi^{\kappa}_{k,\square}$ be its characteristic function in the step-$k$ decomposition of unity,
and put
\begin{equation}\label{8.8:eq-history-measure-6}
  E^{\kappa}_{k,\square}:=\bigl\{\omega\in\Omega:\ \text{the step-}k\text{ decomposition of }
  \omega\text{ chose }1-\chi^{\kappa}_{k,\square}\bigr\} .
\end{equation}
Graded at threshold $\lambda$, $E^{\kappa}_{k,\square}(\lambda)$ is the intersection of
$E^{\kappa}_{k,\square}$ with the event that the functional of the kind $\kappa$ is at least
$\lambda$. The field event $F_{k,\square}(\lambda):=\{a_k(\square)\ge\lambda\}$ is an event of
the chain's field $V_k$, whether or not $\square$ is tested. Finally, for a cell $\mathfrak{c}$
of level $k$, $L^{(k)}_{\mathfrak{c}}$ is the event that a connected component of the large-field
structure $\mathcal{O}_k(\omega)$ of the history at level $k$ (the letter $\mathcal{O}$ with
subscript and argument is reserved for this object) which is live at level $k$ (still carrying its
large-field factor there) contains $\mathfrak{c}$.
\end{definition}

\begin{lemma}[Graded rarity in creation and field form, reduced to live cells]
\label{8.8:lem-graded-rarity}
Let $\tilde{\nu}$ be the normalised annealed measure on histories of
Definition~\ref{8.8:not-history-measure}. For an event that involves the scale-$k$ field,
$\tilde{\nu}$ is read jointly with that field, as in that definition. For a level $k$, let
$\square$ range over the cubes of the level-$k$ partition and $\mathfrak{c}$ over the
level-$k$ cells. Let $\kappa$ range over the kinds $\mathrm{S}0,\dots,\mathrm{S}10$ of
large-field creation. We use the following events of
Definition~\ref{8.8:not-history-measure}, which we recall.
\begin{itemize}
\item[] $E^{\kappa}_{k,\square}(\lambda)$, for $\lambda\ge 1$, is the creation event of kind
$\kappa$ at $(k,\square)$, graded at the multiple $\lambda$ of its threshold;
$E^{\kappa}_{k,\square}=E^{\kappa}_{k,\square}(1)$ is the creation event itself.
\item[] $F_{k,\square}(\lambda)$ is the field event $\{a_k(\square)\ge\lambda\}$, in which
$a_k(\square)$ is Ba{\l}aban's small-field functional, normalised as in that definition.
Thus $F_{k,\square}(\lambda)$ is the event that the scale-$k$ field reaches or exceeds
$\lambda$ times its threshold somewhere in $\square$.
\item[] $L^{(k)}_{\mathfrak{c}}$ is the event that some live component at level $k$
contains $\mathfrak{c}$.
\end{itemize}
The \emph{creation form} of graded rarity is the bound
\begin{equation}\label{8.8:eq-graded-rarity-b}
\sup_{\kappa,\square}\tilde{\nu}\bigl(E^{\kappa}_{k,\square}(\lambda)\bigr)
\le C e^{-c\lambda^{2}p_0(g_k)}\qquad\text{for all $k$ and all $\lambda\ge 1$,}
\end{equation}
and the \emph{field form} is the bound
\begin{equation}\label{8.8:eq-graded-rarity-c}
\sup_{\square}\tilde{\nu}\bigl(F_{k,\square}(\lambda)\bigr)\le C e^{-c\lambda^{2}p_0(g_k)}
\qquad\text{for all $k$ and all $\lambda\ge 1$.}
\end{equation}
Fix a level $k$, and let $n_{\square}$ be an upper bound, valid for every cube $\square$ of
level $k$, for the number of level-$k$ cells that meet $\square$. Then the following hold.
\begin{itemize}
\item[(i)] Up to a $\tilde{\nu}$-null set,
\begin{equation*}
F_{k,\square}(1)\subset\bigcup_{\mathfrak{c}\cap\square\neq\emptyset}L^{(k)}_{\mathfrak{c}} .
\end{equation*}
\item[(ii)] For every kind $\kappa$,
\begin{equation*}
E^{\kappa}_{k,\square}\subset\bigcup_{\mathfrak{c}\cap\square\neq\emptyset}L^{(k)}_{\mathfrak{c}} .
\end{equation*}
\item[(iii)] Suppose that every level-$k$ cell $\mathfrak{c}$ satisfies
\begin{equation}\label{8.8:eq-graded-rarity-a}
\tilde{\nu}\bigl(L^{(k)}_{\mathfrak{c}}\bigr)\le C_L e^{-\theta p_0(g_k)} .
\end{equation}
Then \eqref{8.8:eq-graded-rarity-b} and \eqref{8.8:eq-graded-rarity-c} hold at level $k$
for $\lambda=1$, with $C=n_{\square}C_L$ and $c=\theta$.
\end{itemize}
\end{lemma}

The field form \eqref{8.8:eq-graded-rarity-c} bounds the probability that the scale-$k$
field reaches or exceeds $\lambda$ times its threshold somewhere in $\square$. The creation
form \eqref{8.8:eq-graded-rarity-b} at $\lambda=1$ is the input of the bound on rarity of
young large-field blocks.

\begin{proof}
\emph{(i).} Let a history $\omega$, together with a scale-$k$ field, lie in
$F_{k,\square}(1)$, and suppose that no live component at level $k$ contains a cell that
meets $\square$. Let $x$ be a point of $\square$ at which the field reaches or exceeds its
threshold.
Every live component is a union of level-$k$ cells. If $x$ lay in a live component, the
cell of that component containing $x$ would meet $\square$. Hence $x$ lies in the
complement of the live region.

Suppose first that $x$ lies in territory healed at step $k$. On that territory the
small-field conditions are tested again \cite[p.~378]{B-CMP122b}, and a connected component
on which they fail is placed in the second class of that test, the class of components that
keep their large-field factor; components of this class are live at level $k$. The
component containing $x$ fails the test, so
it is live, and its cell containing $x$ meets $\square$. This contradicts the assumption,
so $x$ does not lie in such territory.

Otherwise the characteristic function $\chi_k$ is a factor of the term of $\omega$ there,
because by \cite[(2.18)]{B-CMP119} $\chi_k$ is a factor of every term on the complement of
the live region. Since $\chi_k$ vanishes where $a_k(\square)\ge 1$, the joint
weight of $\omega$ and the field vanishes on the set of such pairs. That set is therefore
$\tilde{\nu}$-null, which proves (i).

Inclusion (i) is also one of the clauses of the large-field block bound, the clause stating
that the field event is contained, up to a $\tilde{\nu}$-null set, in the union of the
live-cell events of the cells meeting the cube.

\emph{(ii).} Inclusion (ii) is the creation clause of the large-field block bound, which
states, for every kind $\kappa$ of large-field creation, that the creation event
$E^{\kappa}_{k,\square}$ is contained in the union of the events $L^{(k)}_{\mathfrak{c}}$
over the level-$k$ cells $\mathfrak{c}$ meeting $\square$.

\emph{(iii).} By (ii), subadditivity of
$\tilde{\nu}$ and \eqref{8.8:eq-graded-rarity-a}, for every $\kappa$ and every $\square$,
\begin{equation*}
\tilde{\nu}\bigl(E^{\kappa}_{k,\square}\bigr)
\le\sum_{\mathfrak{c}\cap\square\neq\emptyset}\tilde{\nu}\bigl(L^{(k)}_{\mathfrak{c}}\bigr)
\le n_{\square}C_L e^{-\theta p_0(g_k)} .
\end{equation*}
By (i), the null set being invisible to $\tilde{\nu}$, the same chain of inequalities holds
with $F_{k,\square}(1)$ in place of $E^{\kappa}_{k,\square}$. With $C=n_{\square}C_L$ and
$c=\theta$, and since $\lambda^2=1$ when $\lambda=1$, these two bounds are
\eqref{8.8:eq-graded-rarity-b} and \eqref{8.8:eq-graded-rarity-c} at level $k$ for
$\lambda=1$.
\end{proof}

\begin{lemma}[Boundedness of the grading on tested cubes]\label{8.8:lem-grading-bounded}
Let $k\ge 1$ and let $\omega\in\Omega$ be any history. Put
\begin{equation*}
\bar\lambda_1:=2B_3L^2(1+\beta_0)^{3},\qquad
\bar\lambda_2:=O(L^2)\,\frac{A_0}{A_1},\qquad
\bar\lambda_3:=2+o(1),
\end{equation*}
where $O(L^2)$ is the constant of that order in \cite[p.~247]{B-CMP119}, and the term $o(1)$ is the size,
relative to $\delta_k$, of the deviation of the background from the identity bounded in
\cite[(3.19)]{B-CMP119}. Then:
\begin{enumerate}
\item[(a)] on $E^{\mathrm{S}1}_{k,\square}$ one has $a_k(\square)<\bar\lambda_1$;
\item[(b)] on $E^{\mathrm{S}2}_{k+1,\square'}$ one has $\varphi_k(\square')<\bar\lambda_2$;
\item[(c)] on $E^{\mathrm{S}3}_{k+1,\square'}$ one has $\psi_k(\square')<\bar\lambda_3$.
\end{enumerate}
Hence, for $\kappa\in\{\mathrm{S}1,\mathrm{S}2,\mathrm{S}3\}$, the graded creation event
$E^{\kappa}_{k,\square}(\lambda)$ of Definition~\ref{8.8:not-history-measure} is empty for
$\lambda\ge\bar\lambda_{\kappa}$ (for $\kappa=\mathrm{S}2,\mathrm{S}3$ with $(k,\square)$ read as
$(k+1,\square')$). Consequently, for these kinds, the creation form \eqref{8.8:eq-graded-rarity-b} of
graded rarity for all $\lambda\ge 1$ is equivalent to its case $\lambda=1$: if
\begin{equation*}
\tilde{\nu}\bigl(E^{\kappa}_{k,\square}\bigr)\le C e^{-\theta p_0(g_k)},
\end{equation*}
then
\begin{equation}\label{8.8:eq-grading-bounded-1}
\tilde{\nu}\bigl(E^{\kappa}_{k,\square}(\lambda)\bigr)
\le C\exp\Bigl(-\frac{\theta}{\bar\lambda_{\kappa}^{2}}\,\lambda^{2}p_0(g_k)\Bigr)
\qquad\text{for every }\lambda\ge 1 .
\end{equation}
\end{lemma}

\begin{proof}
\emph{The kind} $\mathrm{S}1$. In \cite[\S3]{B-CMP119} the condition of kind $\mathrm{S}1$ is tested only on
the complement $(\bar Z_{k-1}^{\sim 4})^{c}$ of the enlargement of $\bar Z_{k-1}$ by four layers of cubes, and
there the characteristic function $\chi_{k-1}$ is a factor of the term. Ba{\l}aban states in
\cite[\S3]{B-CMP119} that ``the restrictions introduced by the characteristic function $\chi_k$ imply that
the new fields $V_{k+1}$ satisfy the bounds $|V_{k+1}(\partial p')-1|<2L^2\varepsilon_k$''. Applied one
level lower, the factor $\chi_{k-1}$ gives
\begin{equation*}
|V_k(\partial p')-1|<2L^2\varepsilon_{k-1}
\end{equation*}
on the cubes where $\mathrm{S}1$ is tested. By the theorem on the minimiser invoked at
\cite[(3.12)]{B-CMP119}, in the form in which it is used there, the local minimiser $U_{k,\square}$ determined by these
block fields then satisfies, on the plaquettes $p$ of $\square$,
\begin{equation*}
|U_{k,\square}(\partial p)-1|<2B_3L^2(1+\beta_0)\,\varepsilon_{k-1}\eta^{2}.
\end{equation*}
By \cite[(2.8)]{B-CMP119}, $\varepsilon_{k-1}\le(1+\beta_0)^2\varepsilon_k$. Inserting this into the
previous bound gives
\begin{equation*}
|U_{k,\square}(\partial p)-1|<2B_3L^2(1+\beta_0)^{3}\,\varepsilon_k\eta^{2}
=\bar\lambda_1\,\varepsilon_k\eta^{2},
\end{equation*}
which is the bound $a_k(\square)<\bar\lambda_1$ of~(a).

\emph{The kind} $\mathrm{S}2$. The condition of kind $\mathrm{S}2$ is tested where the characteristic
functions $\chi_k$, $\chi_{k+1}$ and $\chi_{\mathrm{Ax}}$ all hold, $\chi_{\mathrm{Ax}}$ denoting the
characteristic function of the axial-gauge condition of \cite[\S3]{B-CMP119}. These three restrictions give
\begin{equation*}
\bigl|V_k\,(V^{(k)}_{\square'})^{-1}-1\bigr|<O(L^2)\,\varepsilon_k
\qquad\text{on }\square'^{\sim 2},
\end{equation*}
with the constant $O(L^2)$ of \cite[p.~247]{B-CMP119}. Since $\delta_k=(A_1/A_0)\varepsilon_k$, the
right-hand side equals $O(L^2)(A_0/A_1)\,\delta_k=\bar\lambda_2\,\delta_k$, which is the bound
$\varphi_k(\square')<\bar\lambda_2$ of~(b).

\emph{The kind} $\mathrm{S}3$. The condition of kind $\mathrm{S}3$ is tested under the characteristic
function $\chi'_k$ of the small-field region of step $k+1$ in \cite[\S3]{B-CMP119}, that is, under
\begin{equation*}
\bigl|e^{ig_kA_k}V^{(k)}(V^{(k)}_{\square'})^{-1}-1\bigr|<2\delta_k,
\end{equation*}
while by \cite[(3.19)]{B-CMP119} the background alone satisfies
$|V^{(k)}(V^{(k)}_{\square'})^{-1}-1|\ll\delta_k$. Together these give $\psi_k(\square')<2+o(1)$, the term
$o(1)$ being the size of this background deviation relative to $\delta_k$, so that~(c) holds.

\emph{Emptiness.} In each of the three cases the functional that defines the grade is bounded above by
$\bar\lambda_{\kappa}$ on the ungraded event; hence no history in $E^{\kappa}_{k,\square}$ reaches the grade
$\lambda\ge\bar\lambda_{\kappa}$, and $E^{\kappa}_{k,\square}(\lambda)=\emptyset$ for those $\lambda$.

\emph{The last sentence.} Let $\lambda\ge 1$. If $\lambda\ge\bar\lambda_{\kappa}$, then
$\tilde{\nu}(E^{\kappa}_{k,\square}(\lambda))=\tilde{\nu}(\emptyset)=0$ and \eqref{8.8:eq-grading-bounded-1}
holds. If $1\le\lambda<\bar\lambda_{\kappa}$, then $\lambda^2/\bar\lambda_{\kappa}^2\le 1$, so
\begin{equation*}
\frac{\theta}{\bar\lambda_{\kappa}^{2}}\,\lambda^{2}\le\theta .
\end{equation*}
Multiplying by $p_0(g_k)\ge 0$ and exponentiating, and using
$E^{\kappa}_{k,\square}(\lambda)\subset E^{\kappa}_{k,\square}(1)=E^{\kappa}_{k,\square}$, we get
\begin{equation*}
\tilde{\nu}\bigl(E^{\kappa}_{k,\square}(\lambda)\bigr)\le\tilde{\nu}\bigl(E^{\kappa}_{k,\square}\bigr)
\le Ce^{-\theta p_0(g_k)}
\le C\exp\Bigl(-\frac{\theta}{\bar\lambda_{\kappa}^{2}}\,\lambda^{2}p_0(g_k)\Bigr),
\end{equation*}
which is \eqref{8.8:eq-grading-bounded-1}. The converse implication is the case $\lambda=1$ of
\eqref{8.8:eq-graded-rarity-b}.
\end{proof}

\begin{remark}[Scope of the grading]
(i) Lemma~\ref{8.8:lem-grading-bounded} shows that the grading by $\lambda$ is compatible with Ba{\l}aban's
hierarchy of thresholds: the window of the previous level caps the exceedance at the next level at a
constant multiple of the threshold. Thus the rate constant $c$ of \eqref{8.8:eq-graded-rarity-b} may be
taken to be $\theta/\bar\lambda_{\kappa}^{2}$; for $\kappa=\mathrm{S}1$ it depends on
$(L,B_3,\beta_0,\theta)$ only, and it is independent of the volume, of the position of the cube and of the
cutoff.

(ii) For $k=0$ (the kind $\mathrm{S}0$, bare plaquettes, where the grade ranges up to $2/\varepsilon_0$) and
for the preparatory kinds whose a priori range is polylogarithmic (for the kind $\mathrm{S}8$ the a priori
bound is the one of \cite{B-CMP122a}), the grade is not bounded a priori, and a graded bound is obtained
only through the statement that the exponential reserve of a large-field region created at grade $\lambda$
is $\lambda^2$ times the ungraded reserve. This comes from the same pointwise extraction as in
\cite[p.~380]{B-CMP122b}: a cube of a region created at level $j$ failing the test at $\lambda\varepsilon_j$
forces a block plaquette of size at least $B_3^{-1}\lambda\varepsilon_j$, and the Wilson action pays
$\gamma_0(B_3^{-1}\lambda\varepsilon_j)^2/(2g_j^2)$. The surplus in the exponent is read by none of the
later inductive bounds of the large-field expansion; in particular it is not an input of the statement on
the rarity of young large-field blocks.

(iii) The grading has content only in the field form \eqref{8.8:eq-graded-rarity-c} inside live regions,
where $V_k$ is an average of older unrestricted fields and $a_k$ is not bounded a priori. This is the
regime that an estimate of the live cells a fixed number $\iota$ of levels further, at level $k+\iota$, by
the graded field event at grade $\tfrac12L^{2\iota}\lambda$ would require; the argument of this chapter does
not use the grading there.
\end{remark}

\begin{definition}[Bare data, members and attainable levels]\label{8.8:not-bare-data-members}
The following objects are used throughout the rest of this chapter.
\begin{itemize}
\item[(a)] A \emph{bare datum} is a pair $(x_0,\mathsf{N})$ consisting of a bare inverse coupling $x_0$
and a number $\mathsf{N}$ of bare sites per side. A bare datum fixes the Wilson measure, which
we denote by $\mu$ (Definition~\ref{1.1:def-wilson-action}).
\item[(b)] The \emph{reference flow} starts from $x_0$ and is defined by
\begin{equation*}
x_{k+1}:=x_k-b_{k+1}(0),
\end{equation*}
where $b_{k+1}(0)$ is the value at zero source of the sourced coupling increment $b_{k+1}$ of the
correlation theorem. The reference flow does not depend on the torus.
\item[(c)] A \emph{member} over the bare datum $(x_0,\mathsf{N})$ is a pair $(g,N_T)$ such that
$g\le\gamma_J$ and $X:=g^{-2}<x_0$, and such that, with $K(x_0,g)$ the largest $k$ for which
$x_j>X$ holds for every $j\le k$, we have
\begin{equation*}
\mathsf{N}=N_T\,L^{K(x_0,g)} .
\end{equation*}
The renormalisation steps applied for a member (its admissible runs) are carried from the bare
level $0$ to level $K-m$, where $K=K(x_0,g)$. The integer $m$ obeys
$m\ge m_{\mathbb{T}}(g)$, the lower bound given by the torus exponent floor $m_{\mathbb{T}}$
of the glossary, and the torus read at level $K-m$ has $N_{\mathrm{site}}:=N_T L^m$ sites per side.
\item[(d)] With $K=K(x_0,g)$, the reference flow satisfies
\begin{equation}\label{8.8:eq-bare-data-members-1}
\lambda(j-i)-2c_2\;\le\;x_i-x_j\;\le\;3\lambda(j-i)+2c_2\qquad(i<j\le K),
\end{equation}
together with
\begin{equation*}
x_K\in\,]X,\,X+b_+] .
\end{equation*}
Here $\lambda$ (the flow slope), $c_2$ and $b_+$ are the constants with which the correlation
theorem states these bounds; the flow constants are fixed in
Definition~\ref{1.2:def-flow-constants}. From here on the letter $\lambda$ denotes the flow slope
of \eqref{8.8:eq-bare-data-members-1}; it is not the grading parameter of the creation events
used earlier in the chapter.
\item[(e)] The coupling $\gamma_*$ is the least of the following ceilings on the coupling: those
under which the theorem on rarity of young large-field blocks holds, and the ceiling
\eqref{8.8:eq-coupling-ceilings-a} below. We put
\begin{equation*}
X_*:=\gamma_*^{-2},\qquad K_*:=K(x_0,\gamma_*),\qquad
c_{\mathrm{gap}}:=\lceil 2c_2/\lambda\rceil+1 .
\end{equation*}
\item[(f)] A level $k^{\diamond}<K_*$ is \emph{attainable} if
\begin{equation}\label{8.8:eq-bare-data-members-2}
x_{k^{\diamond}+1}<\min_{j\le k^{\diamond}}x_j .
\end{equation}
\item[(g)] Let $u_{\gamma}$ be the depth function of the theorem on rarity of young large-field
blocks, and $\mathfrak{F}_{\gamma}$ the family of functions of the coupling from which it is
built. We put
\begin{equation}\label{8.8:eq-bare-data-members-3}
\bar u(g):=\sup\{\,u_{\gamma}(g') : g\le g'\le\gamma_*\,\}.
\end{equation}
The quantity $\bar u(g)$ is finite and polylogarithmic of degree $\mathfrak{p}_0$ in $\log X$.
\end{itemize}
\end{definition}

\begin{proof}
Items (b), (d) and (g) contain assertions; the remaining items are definitions.

That the reference flow of (b) does not depend on the torus is the content of the lemma on
torus-independence of the reference flow, which is used here as stated there.

The bounds \eqref{8.8:eq-bare-data-members-1} in (d), and the inclusion of $x_K$ in
$]X,X+b_+]$, are the flow bounds of the correlation theorem. They apply here with
$K=K(x_0,g)$, the first-passage level of (c).

For (g), set $X'=g'^{-2}$. The condition $g\le g'\le\gamma_*$ then reads $X_*\le X'\le X$, so the
supremum in \eqref{8.8:eq-bare-data-members-3} is taken over $X'$ in this interval. Each member
of the family $\mathfrak{F}_{\gamma}$ is monotone or logarithmic in $X'$. This is the ground on
which $\bar u(g)$ is finite and polylogarithmic of degree $\mathfrak{p}_0$ in $\log X$.
\end{proof}

\begin{lemma}[Independence of the steps from the target coupling. Stated; proof incomplete at the step marked $(\ast)$]\label{8.8:lem-target-independence}
Let $(x_0,\mathsf{N})$ be a bare datum in the sense of Notation~\ref{8.8:not-bare-data-members}, and let $k\ge 0$ be a step of a run over $(x_0,\mathsf{N})$ whose block lattices $\mathbb{T}^{(j)}$ carry the $L$-adically nested partitions anchored at the origin of the torus. The $k$-th step $\mathbf{R}T$ of that run reads the following objects:
\begin{itemize}
\item $g_k$, the parameters $\varepsilon_k$, $\delta_k$ of Ba{\l}aban's $k$-th step, and the parameter $\delta'_k$ of that step, defined together with $\varepsilon_k$ and $\delta_k$ in \cite[\S2--\S3]{B-CMP119} and \cite[\S1]{B-CMP122a};
\item the running radius $\check{R}_k=R(\min_{i\le k}x_i)$;
\item the bare count $\mathsf{N}$ and the per-scale integers $\mathsf{N}_0$ of Ba{\l}aban's step, the latter functions of $(x_i)_{i\le k}$;
\item the $L$-adically nested partitions of $\mathbb{T}^{(k)}$ anchored at the origin of the torus.
\end{itemize}
All of these are determined by $(x_0,\mathsf{N})$ and by $k$. The target coupling $g$ of a member $(g,N_T)$ over $(x_0,\mathsf{N})$ enters a run only through two things. The first is the run's stopping level. The second is the additive constant $E$ in
\begin{equation*}
\rho_0=\exp\bigl[-g_0^{-2}A-E\bigr],
\end{equation*}
and this constant cancels in $\tilde{\nu}$. Hence two runs over the same bare datum, for any two members over it, coincide through the last step of the shorter run. Moreover, for every event $\mathcal{E}$ of levels at most that step, $\tilde{\nu}(\mathcal{E})$ is the same number in both runs.
\end{lemma}

Stated; the proof below is incomplete at the step marked $(\ast)$.

\begin{proof}
Fix a member $(g,N_T)$ over $(x_0,\mathsf{N})$, and a run of that member stopping at level $K(x_0,g)-m$ for some $m\ge m_{\mathbb{T}}(g)$.

$(\ast)$ Each object read by the $k$-th step $\mathbf{R}T$ is defined in \cite[\S2--\S3]{B-CMP119} and \cite[\S1]{B-CMP122a} from two data: the couplings $(g_j)_{j\le k}$, and the partitions of the block lattices. The partition of $\mathbb{T}^{(k)}$ used at step $k$ is stated in \cite[\S2--\S3]{B-CMP119} to be compatible with all other partitions. By hypothesis, the partitions used at step $k$ are the $L$-adically nested partitions of $\mathbb{T}^{(k)}$ anchored at the origin of the torus. They are therefore determined by the bare datum $(x_0,\mathsf{N})$ and by $k$.

Next we show that the couplings $(g_j)_{j\le k}$ depend on $x_0$ and $k$ only. By Notation~\ref{8.8:not-bare-data-members}, the reference flow is defined recursively, starting from $x_0$, by
\begin{equation*}
x_{j+1}=x_j-b_{j+1}(0).
\end{equation*}
This flow is independent of the torus. The recursion does not involve $X=g^{-2}$: in Notation~\ref{8.8:not-bare-data-members} the target coupling enters only after the flow is defined, through the first passage $K(x_0,g)$, the largest $k'$ with $x_j>X$ for all $j\le k'$, and through the conditions on members and on their admissible runs. Hence $(x_j)_{j\le k}$, and with it $g_k=x_k^{-1/2}$, is a function of $x_0$ and $k$.

The remaining objects of the list are functions of these. First, $\check{R}_k=R(\min_{i\le k}x_i)$ is a function of $(x_i)_{i\le k}$. Second, $\varepsilon_k=g_k\,p_0(g_k)$ is a function of $g_k$ through the degree function $p_0(\cdot)$. Third, $\delta_k=(A_1/A_0)\,\varepsilon_k$. Fourth, by the step marked $(\ast)$, $\delta'_k$ and the per-scale integers $\mathsf{N}_0$ are defined from $(g_j)_{j\le k}$ and the partitions, and so are functions of $(x_0,\mathsf{N})$ and $k$. Finally, the bare count $\mathsf{N}$ is part of the bare datum. The whole list read by the $k$-th step is therefore determined by $(x_0,\mathsf{N})$ and $k$.

Two quantities of the run remain in which $g$ could enter. The first is the stopping level $K(x_0,g)-m$. The second is the constant $E$ in $\rho_0=\exp[-g_0^{-2}A-E]$. This constant enters $\rho_0$ only as the factor $e^{-E}$, and it cancels in the ratio
\begin{equation*}
\tilde{\nu}(\mathcal{E})=Z^{-1}\sum_{\omega\in\mathcal{E}}\tau_{\omega}
\end{equation*}
of Notation~\ref{8.8:not-history-measure}.

Now let $(g,N_T)$ and $(g',N'_T)$ be two members over $(x_0,\mathsf{N})$. Fix an admissible run of each, and let $k_1$ be the last step of the shorter of the two runs. For each $k\le k_1$, the $k$-th step of either run reads the same objects, by the preceding paragraphs. The two runs therefore coincide through step $k_1$.

By Notation~\ref{8.8:not-history-measure}, $\tilde{\nu}$ is the normalised measure on the histories of the procedure $\mathbf{R}T$ iterated from $U\sim\rho_0\,dU/\int\rho_0$. The initial law $\rho_0\,dU/\int\rho_0$ does not depend on $E$, and the first $k_1$ steps of the two runs coincide. The measure $\tilde{\nu}$ is therefore fixed by the initial law and by the steps, and the stopping level of a run moves only the level at which the representation of $\tilde{\nu}$ ends; an event of levels at most $k_1$ is determined by the first $k_1$ steps. Hence, for every event $\mathcal{E}$ of levels at most $k_1$, $\tilde{\nu}(\mathcal{E})$ is the same number in both runs.
\end{proof}

\begin{lemma}[Re-targeting a member to an attainable level]\label{8.8:lem-retargeting}
Let $(x_0,\mathsf{N})$ be a bare datum and $(g,N_T)$ a member over it, in the sense of
Notation~\ref{8.8:not-bare-data-members}. Put $X:=g^{-2}$ and $K:=K(x_0,g)$. Let $\gamma_*$,
$X_*=\gamma_*^{-2}$ and $K_*=K(x_0,\gamma_*)$ be as in
Notation~\ref{8.8:not-bare-data-members}, and assume $K<K_*$.
\begin{itemize}
\item[(a)] The level $K$ is attainable. For every level $k_1$ with
$K\le k_1\le K_*-c_{\mathrm{gap}}-1$ there is an attainable level in
$[k_1,\,k_1+c_{\mathrm{gap}}]$ (indeed in $[k_1,\,k_1+c_{\mathrm{gap}}-1]$).
\item[(b)] Let $k^{\diamond}$ be an attainable level with $K\le k^{\diamond}<K_*$, and assume that
$L^{k^{\diamond}-K}$ divides $N_T$. Put $X^{\diamond}:=x_{k^{\diamond}+1}$, so that
$X_*<X^{\diamond}\le X$, and put $g^{\diamond}:=(X^{\diamond})^{-1/2}$ and
$N_T^{\diamond}:=N_T L^{-(k^{\diamond}-K)}$. Then $(g^{\diamond},N_T^{\diamond})$ is a member
over the same bare datum $(x_0,\mathsf{N})$, with $K(x_0,g^{\diamond})=k^{\diamond}$ and
$g\le g^{\diamond}<\gamma_*$. Moreover $m_{\mathbb{T}}(g^{\diamond})\le m_{\mathbb{T}}(g)$, and
the longest run of $(g^{\diamond},N_T^{\diamond})$ stops at the level
$k^{\diamond}-m_{\mathbb{T}}(g^{\diamond})$, which satisfies
\begin{equation*}
k^{\diamond}-m_{\mathbb{T}}(g^{\diamond})\ \ge\ K-m_{\mathbb{T}}(g)+(k^{\diamond}-K),
\end{equation*}
with $N^{\diamond}=N_T^{\diamond}L^{m_{\mathbb{T}}(g^{\diamond})}$ sites per side.
\item[(c)] The number of levels between $K$ and $K_*$ satisfies
\begin{equation}\label{8.8:eq-retargeting-a}
K_*-K\ \ge\ \frac{X-X_*-b_+-2c_2}{3\lambda}.
\end{equation}
\end{itemize}
\end{lemma}

\begin{proof}
Throughout, $(x_k)$ is the reference flow of Notation~\ref{8.8:not-bare-data-members}, and a
level $k$ is attainable when $x_{k+1}<\min_{j\le k}x_j$; this is the form in which
attainability is used below.

\emph{(a).} By the definition of $K=K(x_0,g)$ as the largest $k$ with $x_j>X$ for all $j\le k$,
we have $x_j>X$ for every $j\le K$, so $X<\min_{j\le K}x_j$; and, $K$ being the largest such
index, the condition fails at $K+1$, which can only happen through $x_{K+1}\le X$. Hence
\begin{equation*}
x_{K+1}\ \le\ X\ <\ \min_{j\le K}x_j ,
\end{equation*}
and $K$ is attainable.

Now let $K\le k_1\le K_*-c_{\mathrm{gap}}-1$. Put $\mu_1:=\min_{j\le k_1}x_j$ and choose
$j_1\le k_1$ with $\mu_1=x_{j_1}$. We apply the lower bound on the decrements of the reference
flow recorded in Notation~\ref{8.8:not-bare-data-members}, $x_i-x_j\ge\lambda(j-i)-2c_2$, to the
member $\gamma_*$, whose first passage is $K_*$; with $i=j_1$ and $j=j_1+n$ it gives, for
$n\ge 0$ with $j_1+n\le K_*$,
\begin{equation*}
x_{j_1}-x_{j_1+n}\ \ge\ \lambda n-2c_2 ,
\end{equation*}
and by the choice of $c_{\mathrm{gap}}$ the right-hand side is strictly positive for
$n\ge c_{\mathrm{gap}}$. Take $n:=k_1+c_{\mathrm{gap}}-j_1$. Since $j_1\le k_1$ we have
$n\ge c_{\mathrm{gap}}$, and $j_1+n=k_1+c_{\mathrm{gap}}\le K_*-1$, so the bound applies and
yields $x_{k_1+c_{\mathrm{gap}}}<x_{j_1}=\mu_1$. Consequently the set of indices $i>k_1$ with
$x_i<\mu_1$ is non-empty, since it contains $k_1+c_{\mathrm{gap}}$; let $k_2$ be its least
element. Then
\begin{equation*}
k_1\ <\ k_2\ \le\ k_1+c_{\mathrm{gap}}\ \le\ K_*-1 .
\end{equation*}
For $k_1<i<k_2$ the minimality of $k_2$ gives $x_i\ge\mu_1$, and $\mu_1>x_{k_2}$ by the choice
of $k_2$. The minimum of $x_j$ over $j\le k_1$ is $\mu_1$ and every $x_i$ with $k_1<i\le k_2-1$
is at least $\mu_1$; therefore $\min_{j\le k_2-1}x_j=\mu_1>x_{k_2}$. With
$k^{\diamond}:=k_2-1$ this reads $x_{k^{\diamond}+1}<\min_{j\le k^{\diamond}}x_j$, so
$k^{\diamond}$ is attainable, and $k^{\diamond}\in[k_1,\,k_1+c_{\mathrm{gap}}-1]$ by the display
above.

\emph{(b).} Since $k^{\diamond}$ is attainable, we have
\begin{equation*}
x_j\ \ge\ \min_{i\le k^{\diamond}}x_i\ >\ x_{k^{\diamond}+1}\ =\ X^{\diamond}
\qquad(j\le k^{\diamond}),
\end{equation*}
whereas
$x_{k^{\diamond}+1}=X^{\diamond}$ is not strictly greater than $X^{\diamond}$. Hence the largest
$k$ with $x_j>X^{\diamond}$ for all $j\le k$ is $k^{\diamond}$, that is,
$K(x_0,g^{\diamond})=k^{\diamond}$.

We verify $X_*<X^{\diamond}\le X$. Since $k^{\diamond}+1\le K_*=K(x_0,\gamma_*)$, the definition
of $K_*$ gives $x_{k^{\diamond}+1}>X_*$, that is, $X^{\diamond}>X_*$. If $k^{\diamond}=K$, then
$X^{\diamond}=x_{K+1}\le X$ by part (a). If $k^{\diamond}>K$, then $K+1\le k^{\diamond}$ and
attainability of $k^{\diamond}$ gives
\begin{equation*}
X^{\diamond}\ <\ \min_{j\le k^{\diamond}}x_j\ \le\ x_{K+1}\ \le\ X .
\end{equation*}
In both cases $X^{\diamond}\le X$.
Taking inverse square roots, $g\le g^{\diamond}<\gamma_*$.

Membership of $(g^{\diamond},N_T^{\diamond})$ over $(x_0,\mathsf{N})$: we have
$X^{\diamond}\le X<x_0$; since $\gamma_*$ is the coupling of a member (it was used as such in
part (a)), $\gamma_*\le\gamma_J$, and $g^{\diamond}<\gamma_*$ gives $g^{\diamond}\le\gamma_J$;
$N_T^{\diamond}$ is an integer because $L^{k^{\diamond}-K}$ divides $N_T$; and
\begin{equation*}
\mathsf{N}\ =\ N_T L^{K}\ =\ N_T L^{-(k^{\diamond}-K)}L^{k^{\diamond}}
\ =\ N_T^{\diamond}L^{k^{\diamond}}\ =\ N_T^{\diamond}L^{K(x_0,g^{\diamond})} .
\end{equation*}

For the torus exponent we use the torus size floor of
Notation~\ref{8.8:not-bare-data-members},
\begin{equation*}
N_{\mathbb{T}}(g,m)\ =\ C_{\mathbb{T}}\sup_{u}L^{-u}\,
R\bigl(X+b_++4c_2+3\lambda(m+u)\bigr),\qquad X=g^{-2}.
\end{equation*}
This quantity is nondecreasing in $X$; since $X=g^{-2}$ decreases as $g$ increases, the floor
$m_{\mathbb{T}}$ determined by $N_{\mathbb{T}}$ is nonincreasing in $g$. As $g\le g^{\diamond}$,
we obtain $m_{\mathbb{T}}(g^{\diamond})\le m_{\mathbb{T}}(g)$.

Finally, by Notation~\ref{8.8:not-bare-data-members} the admissible runs of the member
$(g^{\diamond},N_T^{\diamond})$ stop at the levels $K(x_0,g^{\diamond})-m$ with
$m\ge m_{\mathbb{T}}(g^{\diamond})$, with $N_T^{\diamond}L^m$ sites per side; the longest is the
one with $m=m_{\mathbb{T}}(g^{\diamond})$. It stops at
$k^{\diamond}-m_{\mathbb{T}}(g^{\diamond})$ with $N^{\diamond}=N_T^{\diamond}
L^{m_{\mathbb{T}}(g^{\diamond})}$ sites per side, and, by
$m_{\mathbb{T}}(g^{\diamond})\le m_{\mathbb{T}}(g)$,
\begin{equation*}
k^{\diamond}-m_{\mathbb{T}}(g^{\diamond})\ \ge\ k^{\diamond}-m_{\mathbb{T}}(g)
\ =\ K-m_{\mathbb{T}}(g)+(k^{\diamond}-K).
\end{equation*}

\emph{(c).} By the definition of $K$ we have $x_K>X$. At the first passage $K_*$ of $\gamma_*$
the flow satisfies $x_{K_*}\le X_*+b_+$, with $b_+$ the flow constant of
Notation~\ref{8.8:not-bare-data-members}. Hence
\begin{equation*}
x_K-x_{K_*}\ >\ X-X_*-b_+ .
\end{equation*}
The upper bound $x_i-x_j\le 3\lambda(j-i)+2c_2$ on the decrements of the reference flow recorded
in Notation~\ref{8.8:not-bare-data-members}, applied with $i=K$ and $j=K_*$, gives
$x_K-x_{K_*}\le 3\lambda(K_*-K)+2c_2$. Combining the two bounds,
$3\lambda(K_*-K)>X-X_*-b_+-2c_2$, which implies \eqref{8.8:eq-retargeting-a}.
\end{proof}

\begin{remark}[Re-targeting in the parametrisation of Theorem~\ref{3.1:thm-main}]
In the parametrisation of Theorem~\ref{3.1:thm-main}, a run is given by the triple
$(g_0,g,m)$ of bare coupling, reference coupling and torus exponent, with $K$ steps and a torus
of $2L^m$ sites per side at the end. Here $K$ is the number of renormalisation steps of
Theorem~\ref{3.1:thm-main}, as in the glossary, and the first passage of
Lemma~\ref{8.8:lem-retargeting} is written $K(x_0,g)$ throughout this remark. Re-targeting to an
attainable level $k^{\diamond}$ as in
Lemma~\ref{8.8:lem-retargeting}(b) acts as
\begin{equation*}
(g_0,g,m)\ \longmapsto\ \bigl(g_0,\ g^{\diamond},\ m-(k^{\diamond}-K(x_0,g))\bigr),
\end{equation*}
and the hypotheses of Notation~\ref{2.2:not-hypotheses} read: (H4) $g^{\diamond}\le\gamma_H$;
(H6) $g_0\le g_-(g)\le g_-(g^{\diamond})$; (H5)
$m-(k^{\diamond}-K(x_0,g))\ge m_{\mathbb{T}}(g_-(g^{\diamond}))$. The bare lattice has
$2L^{m+K}$ sites per side in both descriptions. The first-passage convention of
Definition~\ref{1.2:def-flow-constants}, $\min\{k:\ x_k\le g^{-2}+B^{\sharp}\}$, differs from
$K(x_0,g)$ by one index; the argument of the proof above applies to it in the same way.

The two parametrisations correspond as follows. The run of Theorem~\ref{3.1:thm-main} with $K$
steps, $2L^m$ sites per side at the end and terminal coupling approximately $g$ is the longest
run of the member $(\tilde g,N_T)$ whose coupling $\tilde g$ and terminal count $N_T$ satisfy
\begin{equation*}
K(x_0,\tilde g)-m_{\mathbb{T}}(\tilde g)=K,\qquad N_T=2L^{\,m-m_{\mathbb{T}}(\tilde g)} .
\end{equation*}
Here $\tilde g$ is the coupling $m_{\mathbb{T}}(\tilde g)$ levels beyond $K$, a number of
levels of doubly logarithmic order, which lies inside the margin of the ceiling $(\gamma_G)$ on
$g$; in the notation $N_T=n_0L^a$ one has $n_0=2$ and $a=m-m_{\mathbb{T}}(\tilde g)$.
\end{remark}

\begin{lemma}[Logarithmic growth of the global stability exponent]\label{8.8:lem-stability-exponent}
Let $(g,N_T)$ be a member over a bare datum $(x_0,\mathsf{N})$ in the sense of
Notation~\ref{8.8:not-bare-data-members}, with $X=g^{-2}$. Let $K'$ be the terminal level of an
admissible run of the member, put $k=K'$, and let $N_k$ be the number of level-$k$ sites. Let
$\nu=\dim\mathfrak{g}$ be the dimension of the Lie algebra of $\mathrm{SU}(N)$ (throughout this
lemma $N$ is the group parameter of $\mathrm{SU}(N)$), write $1$ for the identity matrix, let
$|\cdot|$ be the
operator norm of $N\times N$ matrices, and let $c_G>0$, depending only on $N$, be such that for every
$r\in(0,1]$ the Haar measure of $\{W\in\mathrm{SU}(N):|W-1|<r\}$ is at least $(c_Gr)^{\nu}$. Assume:
\begin{itemize}
\item[\textup{(a)}] clause \textup{(i-b)} of Theorem~\ref{3.1:thm-main}. We use its bound
$E_+\le C(L,\mathfrak{p},\gamma,N)(1+\log g_k^{-2})$, and its lower bound in the integrated form
\begin{equation}\label{8.8:eq-stability-exponent-1}
\int\rho_k\;\ge\;e^{-E_-N_k}\int\chi_k\,e^{-g_k^{-2}A(U_k(V))}\,dV ,
\end{equation}
where $E_-$ is the constant of that lower bound; here $V$ runs over the level-$k$ bond
configurations, $dV$ is the product of Haar measures over the bonds, and $U_k(V)$ is
Ba{\l}aban's minimising configuration on the fine lattice determined by $V$;
\item[\textup{(a$'$)}] the constant $E_-$ of \eqref{8.8:eq-stability-exponent-1} satisfies
$E_-\le C_-(1+\log g_k^{-2})$ with $C_-=C_-(L,\mathfrak{p},\gamma,N)\ge0$;
\item[\textup{(b)}] the majorant bound for the flat reserve weight. For the majorants
$\mathfrak{M}^{(\theta)}_k(h)$ of the terms $h$ of the level-$k$ expansion that remain after the flat
reserve weight is split off, $\sum_h\mathfrak{M}^{(\theta)}_k(h)\le e^{E_+^{(\theta)}N_k}$ holds with
$E_+^{(\theta)}\le C_+(1+\log g_k^{-2})$, where $C_+=C_+(L,\mathfrak{p},\gamma,N)\ge0$;
\item[\textup{(c)}] the linear-reserve statement. The linear reserve weight is unspent, on the same
terms as
the flat reserve weight in~(b): the bound of~(b) holds for the majorants of the terms $h$ that remain
after the linear reserve weight is split off, with an exponent $E_+^{\mathrm{lin}}$ in place of
$E_+^{(\theta)}$ that obeys the same bound $E_+^{\mathrm{lin}}\le C_+(1+\log g_k^{-2})$;
\item[\textup{(d)}] the terminal inverse coupling of the run satisfies
\begin{equation}\label{8.8:eq-stability-exponent-2}
x_{K'}\;\le\;X+b_++3\lambda\,m_{\mathbb{T}}(g)+2c_2\;\le\;2X ,
\end{equation}
where $b_+$, $c_2$, $\lambda$ are the flow constants of the reference inequality of
Notation~\ref{8.8:not-bare-data-members} and $m_{\mathbb{T}}(g)$ is the torus exponent floor;
\item[\textup{(e)}] the degree function $p_0(\cdot)$ satisfies $p_0(g_k)\ge4B_3$, with $B_3$ as below.
\end{itemize}
Let $Z=\int\rho_k$ be the normaliser. Define $C_{\mathrm{low}}$ by $Z=e^{-C_{\mathrm{low}}N_k}$ and put
$\mathsf{B}(g):=E_+^{\mathrm{lin}}+C_{\mathrm{low}}$. Then
\begin{equation}\label{8.8:eq-stability-exponent-3}
C_{\mathrm{low}}\;\le\;E_-+C_dB_3^2+\nu d\,\log\frac{1}{c_Gg_k},
\end{equation}
with $C_d$ depending on the dimension only and $B_3$ the constant of Ba{\l}aban's bound on the
minimiser. Moreover there is a constant $C_{\mathsf{B}}$, depending only on $(L,\mathfrak{p},\gamma,N)$,
such that
\begin{equation}\label{8.8:eq-stability-exponent-4}
\mathsf{B}(g)\;\le\;C_{\mathsf{B}}\,(1+\log g^{-2}).
\end{equation}
\end{lemma}

\begin{proof}
\emph{Upper half.} By hypothesis~(b), the majorants that remain after the flat reserve weight is split
off obey the bound of~(b) with an exponent $E_+^{(\theta)}\le C_+(1+\log g_k^{-2})$. By
hypothesis~(c), the same bound holds with the linear reserve weight split off in place of the flat
one, with the exponent $E_+^{\mathrm{lin}}$. This has the shape of the bound on $E_+$ in clause (i-b)
(hypothesis~(a)). The logarithmic form is that of the constant $O(\log g_j^{-2})$, in the notation of
\cite[p.~383]{B-CMP122b}, stated there for the last of the constants considered, in place of $O(1)$
for the others. Hence $E_+^{\mathrm{lin}}\le C_+(1+\log g_k^{-2})$, which is the upper half.

\emph{Lower half.} By \eqref{8.8:eq-stability-exponent-1},
\begin{equation*}
Z=\int\rho_k\;\ge\;e^{-E_-N_k}\int\chi_k\,e^{-g_k^{-2}A(U_k(V))}\,dV .
\end{equation*}
The integrand on the right is nonnegative, so the integral is bounded below by its restriction to the set
\begin{equation*}
\mathcal{V}:=\{V:\ |V(b)-1|<g_k\ \text{for every bond } b\}.
\end{equation*}
The measure $dV$ is the product of Haar measures over the bonds, and each site carries $d$ bonds.
An admissible run stops at $K'=K(x_0,g)-m$ for some $m\ge m_{\mathbb{T}}(g)$
(Notation~\ref{8.8:not-bare-data-members}), so $K'\le K(x_0,g)$, and the definition of $K(x_0,g)$
there gives
\begin{equation*}
g_k^{-2}=x_{K'}>X=g^{-2}\ge1,
\end{equation*}
since $g\le\gamma_J\le1$. Hence $g_k<1$, and the choice of
$c_G$ with $r=g_k$ shows that the mass of $\mathcal{V}$ is at least $(c_Gg_k)^{\nu d}$ per site. In total
it is at least $(c_Gg_k)^{\nu dN_k}$.

Let $V\in\mathcal{V}$ and let $p'$ be a plaquette of the level-$k$ lattice, with bonds
$b_1,\dots,b_4$. Then $V(\partial p')=v_1v_2v_3v_4$, where $v_i:=V(b_i)^{\sigma_i}$ with
$\sigma_i\in\{\pm1\}$. For unitary $W$ one has $|W^{-1}-1|=|W^{-1}(1-W)|=|W-1|$, so each
$|v_i-1|<g_k$. The telescoping identity
\begin{equation*}
v_1v_2v_3v_4-1=(v_1-1)v_2v_3v_4+(v_2-1)v_3v_4+(v_3-1)v_4+(v_4-1),
\end{equation*}
together with the fact that unitary factors have norm one, gives $|V(\partial p')-1|<4g_k$.

We compare this with the small-field threshold. Since $\varepsilon_k=g_k\,p_0(g_k)$, with $p_0(\cdot)$
the degree function, we have $\varepsilon_k/B_3=g_k\,p_0(g_k)/B_3$, so the inequality
$4g_k\le\varepsilon_k/B_3$ is the inequality $4B_3\le p_0(g_k)$. By hypothesis~(e) this inequality
holds, so that $4g_k\le\varepsilon_k/B_3$. Hence every
configuration of $\mathcal{V}$ lies in the small-field region, so that $\chi_k=1$ on $\mathcal{V}$.
Moreover, for every plaquette $p$ of the fine lattice, whose spacing in level-$k$ units is $\eta$,
\begin{equation*}
|U_k(\partial p)-1|<4B_3\,g_k\,\eta^2 .
\end{equation*}
This is the bound on the minimiser quoted in \cite{B-CMP119} as Theorem~1 of the reference numbered~16
there, applied to the configuration $V$.

For $W\in\mathrm{SU}(N)$ and $\operatorname{tr}1=1$ one has
$\operatorname{tr}\big((W-1)^*(W-1)\big)=2\big(1-\operatorname{Re}\operatorname{tr}W\big)$. Hence
$1-\operatorname{Re}\operatorname{tr}W\le\tfrac12|W-1|^2$. Inserting the bound on
$|U_k(\partial p)-1|$, each plaquette of the fine lattice contributes at most
$\tfrac12(4B_3g_k\eta^2)^2=8B_3^2g_k^2\eta^4$. The fine lattice has $\eta^{-d}N_k$ sites and hence
$\tfrac{d(d-1)}{2}\eta^{-d}N_k$ plaquettes; since $\eta^{4-d}=1$ for $d=4$, summing over the
plaquettes gives
\begin{equation*}
A(U_k(V))\le C_dB_3^2g_k^2N_k ,\qquad\text{that is,}\qquad g_k^{-2}A(U_k(V))\le C_dB_3^2N_k
\quad\text{on }\mathcal{V},
\end{equation*}
with $C_d:=4d(d-1)$, which depends on the dimension only. Therefore
\begin{align*}
Z&\ge e^{-E_-N_k}\,e^{-C_dB_3^2N_k}\,(c_Gg_k)^{\nu dN_k}\\
&=\exp\Big(-N_k\Big[E_-+C_dB_3^2+\nu d\log\frac{1}{c_Gg_k}\Big]\Big).
\end{align*}
Comparing with $Z=e^{-C_{\mathrm{low}}N_k}$ gives \eqref{8.8:eq-stability-exponent-3}.

\emph{Dependence on $g$.} By hypothesis~(d), that is by \eqref{8.8:eq-stability-exponent-2}, the
terminal inverse coupling satisfies $x_{K'}\le 2X$.
As shown in the lower half, $K'\le K(x_0,g)$ and hence $x_{K'}>X$. Write
$\ell_k:=\log g_k^{-2}=\log x_{K'}$.
Then
\begin{equation*}
\log X<\ell_k\le\log2+\log X .
\end{equation*}
We also have
\begin{equation*}
\log\frac{1}{c_Gg_k}=\log\frac{1}{c_G}+\frac12\,\ell_k .
\end{equation*}
Put
\begin{equation*}
A_{\mathsf{B}}:=C_++C_-+C_dB_3^2+\nu d\max\{\log c_G^{-1},0\}.
\end{equation*}
The upper half,
\eqref{8.8:eq-stability-exponent-3} and the bound on $E_-$ in hypothesis~(a$'$) give
\begin{align*}
\mathsf{B}(g)&\le C_+(1+\ell_k)+C_-(1+\ell_k)+C_dB_3^2+\nu d\log\frac1{c_G}+\frac{\nu d}2\,\ell_k\\
&\le\Big(A_{\mathsf{B}}+\frac{\nu d}{2}\Big)(1+\ell_k),
\end{align*}
where the last step uses $\ell_k\ge0$, which holds since $\ell_k>\log X$ and
$\log X=\log g^{-2}\ge0$ because $g\le\gamma_J\le1$.

Finally, because $\log X\ge0$,
\begin{equation*}
1+\ell_k\le1+\log2+\log X\le(1+\log2)(1+\log X).
\end{equation*}
Since $C_+\ge0$ and $C_-\ge0$, the number $A_{\mathsf{B}}+\nu d/2$ is nonnegative, and multiplying the
last inequality by it and combining with the previous bound shows that
\eqref{8.8:eq-stability-exponent-4} holds with
\begin{equation*}
C_{\mathsf{B}}:=(1+\log2)\Big(A_{\mathsf{B}}+\frac{\nu d}{2}\Big).
\end{equation*}
The factor coming from the terminal level is thus part of $C_{\mathsf{B}}$. The constants $C_+$, $C_-$,
$C_d$, $B_3$, $\nu$, $d$ and $c_G$ depend only on $(L,\mathfrak{p},\gamma,N)$, so $C_{\mathsf{B}}$ depends
only on these.
\end{proof}

\begin{definition}[The two coupling ceilings]\label{8.8:def-coupling-ceilings}
Let $d=4$, let $C_{\mathsf{B}}$ be the constant of Lemma~\ref{8.8:lem-stability-exponent}, and let
$X$, $X_*$, $b_+$, $c_2$, $\lambda$, $c_{\mathrm{gap}}$, $\bar n$ and the torus exponent floor
$m_{\mathbb{T}}(\cdot)$ be as in Notation~\ref{8.8:not-bare-data-members}, with
$X_*=\gamma_*^{-2}$. Let $p_0(\cdot)$ be the degree function,
$p_0(g)=A_0(\log g^{-2})^{\mathfrak{p}_0}$, and let
$\bar u(g):=\sup_{g\le g''\le\gamma_*}u_{\gamma}(g'')$ be the supremum of the depth function
$u_{\gamma}$ of the glossary over $[g,\gamma_*]$.

The first ceiling is the following condition on $\gamma_*$: for every $X'\ge X_*$,
with $g':=X'^{-1/2}$, and for every $p\ge p_0(\gamma_*)-1$,
\begin{equation}\label{8.8:eq-coupling-ceilings-a}
\begin{gathered}
C_{\mathsf{B}}(1+\log X')\bigl(\bar n\,L^{c_{\mathrm{gap}}}\,L^{m_{\mathbb{T}}(g')}\bigr)^{d}
\le \tfrac14\,\theta'\gamma_0A_1\bigl(p_0(g')-2\bigr)^2,\\
\tfrac14\,\theta'\gamma_0A_1(p-1)^2\ge\theta p .
\end{gathered}
\end{equation}

The second ceiling is the following condition on a number $\gamma_G$ with $0<\gamma_G\le\gamma_*$:
\begin{equation}\label{8.8:eq-coupling-ceilings-b}
X-X_*-b_+-2c_2\ \ge\ 3\lambda\bigl(\bar u(g)+2c_{\mathrm{gap}}+1\bigr)
\qquad\text{for every } X\ge\gamma_G^{-2},\ g=X^{-1/2}.
\end{equation}

Both conditions are one-sided. Condition \eqref{8.8:eq-coupling-ceilings-a} holds for every
$\gamma_*$ sufficiently small, once $L$, $M$, $\bar n$, $\theta'$ and $A_1$ have been fixed; it
is therefore a ceiling on $\gamma_*$, chosen after these. Condition
\eqref{8.8:eq-coupling-ceilings-b} holds for every $\gamma_G$ sufficiently small, once $\gamma_*$
has been fixed; it is therefore a ceiling on $g$, chosen after $\gamma_*$.
\end{definition}

\begin{proof}[Proof that the two ceilings can be met]
Consider the first inequality in \eqref{8.8:eq-coupling-ceilings-a}. By
Notation~\ref{8.8:not-bare-data-members}, the torus exponent floor satisfies
$L^{m_{\mathbb{T}}(g')}\le L\,C_{\mathbb{T}}\,R(g')$, where $C_{\mathbb{T}}$ is the constant
fixed there and the radius $R(g')$ grows
like the power $r_{\mathrm{pr}}$ of $\log X'$. Hence, apart from the factors
$C_{\mathsf{B}}$, $\bar n^{d}$, $L^{d c_{\mathrm{gap}}}$ and $(L\,C_{\mathbb{T}})^{d}$, which do
not depend on $X'$, the left member grows like the power $d\,r_{\mathrm{pr}}+1$ of $\log X'$: one
power from the factor $1+\log X'$, and $d\,r_{\mathrm{pr}}$ powers from the $d$-th power of
$R(g')$.

The right member is $\tfrac14\theta'\gamma_0A_1\bigl(A_0(\log X')^{\mathfrak{p}_0}-2\bigr)^2$; its
coefficient $\tfrac14\theta'\gamma_0A_1A_0^2$ is positive, and it grows like the power
$2\mathfrak{p}_0$ of $\log X'$. The conditions $\mathfrak{p}_0\ge 5r_{\mathrm{pr}}$ and
$r_{\mathrm{pr}}\ge 2$ of the order of choice (Notation~\ref{2.3:not-stages-first},
Notation~\ref{2.3:not-stages-middle} and Notation~\ref{2.3:not-stages-last}), together with
$d=4$, give
\begin{equation*}
2\mathfrak{p}_0\ \ge\ 10\,r_{\mathrm{pr}}\ >\ 4\,r_{\mathrm{pr}}+1\ =\ d\,r_{\mathrm{pr}}+1,
\end{equation*}
the strict inequality because $6\,r_{\mathrm{pr}}\ge 12>1$. The ratio of the left member to the
right member therefore tends to $0$ as $X'\to\infty$. So there is a number $\bar X$, depending on
$L$, $M$, $\bar n$, $\theta'$ and $A_1$ (and on the constants fixed before them), such that the
first inequality holds for every $X'\ge\bar X$. It then holds for every $X'\ge X_*$ whenever
$X_*\ge\bar X$, that is, whenever $\gamma_*\le\bar X^{-1/2}$.

Consider the second inequality in \eqref{8.8:eq-coupling-ceilings-a}. The difference
$\tfrac14\theta'\gamma_0A_1(p-1)^2-\theta p$ is a quadratic polynomial in $p$ with positive
leading coefficient $\tfrac14\theta'\gamma_0A_1$. It is therefore nonnegative for every $p$
exceeding its larger real root, and for every $p$ if it has no real root. Moreover
$p_0(\gamma_*)-1=A_0(\log X_*)^{\mathfrak{p}_0}-1$ tends to $+\infty$ as $X_*\to\infty$. Hence
the second inequality holds for every $p\ge p_0(\gamma_*)-1$ once $X_*$ is large enough. Taking
$X_*$ large enough for both inequalities, condition \eqref{8.8:eq-coupling-ceilings-a} holds for
every $\gamma_*$ below a threshold fixed after $(L,M,\bar n,\theta',A_1)$.

Consider \eqref{8.8:eq-coupling-ceilings-b}, with $\gamma_*$, and hence $X_*$, fixed. The left
member $X-X_*-b_+-2c_2$ is $X$ minus a quantity that does not depend on $X$. The right member is
$3\lambda$ times $\bar u(g)+2c_{\mathrm{gap}}+1$. Here $\bar u(g)$ is polylogarithmic in $X$, and
$c_{\mathrm{gap}}$ does not depend on $X$. The left member therefore exceeds the right member for
every $X$ larger than some number $\hat X$, which depends on $\gamma_*$. Any $\gamma_G$ with
$\gamma_G^{-2}\ge\hat X$ satisfies \eqref{8.8:eq-coupling-ceilings-b}. Since $\hat X$ depends on
$\gamma_*$, this is a ceiling on $g$ fixed after $\gamma_*$.
\end{proof}

\emph{The fate decomposition.} Let $(g,N_T)$ be a member over a bare datum $(x_0,\mathsf{N})$, in the sense of Notation~\ref{8.8:not-bare-data-members}. Let $j$ be a level and $\mathfrak{c}$ a level-$j$ cell. Let $L^{(j)}_{\mathfrak{c}}$ be the event that a live component at level $j$ contains $\mathfrak{c}$.

Fix an admissible run of $(g,N_T)$ containing level $j$, and let $K^{\diamond}>j$ be a level of this run or of any run, of a member over the same bare datum, that prolongs it. The \emph{lineage} of this component consists of the live components at levels $k>j$ that contain it. For $j<k\le K^{\diamond}$ let $H_k$ be the event that the lineage of this component is first healed at step $k$. Let $\mathcal{E}^{(j)}_{\mathfrak{c}}(K^{\diamond})$, the \emph{live event}, be the event that the lineage is healed at no step $k$ with $j<k\le K^{\diamond}$, that is, that it is still live at level $K^{\diamond}$. On $\mathcal{E}^{(j)}_{\mathfrak{c}}(K^{\diamond})$, the internal history of the component that is live at level $K^{\diamond}$ contains a seed $Z$ of the first kind whose birth level $j(Z)$ satisfies $j(Z)\le j$.

The large-field regions of a history only grow until the component is healed as a whole. Hence the first healing step of the lineage, when there is one, is unique, and
\begin{equation}\label{8.8:eq-induction-base-a}
L^{(j)}_{\mathfrak{c}}
=\bigsqcup_{j<k\le K^{\diamond}}H_k\ \sqcup\ \mathcal{E}^{(j)}_{\mathfrak{c}}(K^{\diamond}).
\end{equation}

\begin{lemma}[Base case: volume scale or coupling room]\label{8.8:lem-induction-base}
Let $g\le\gamma_G$, where $\gamma_G$ is fixed by the ceiling \eqref{8.8:eq-coupling-ceilings-b} of Definition~\ref{8.8:def-coupling-ceilings}. Let $(g,N_T)$ be a member over $(x_0,\mathsf{N})$ with $N_T=n_0L^a$ and $n_0\le\bar n$, and put $K:=K(x_0,g)$ and $X:=g^{-2}$. Put
\begin{equation*}
\hat\kappa:=\min\{K+a,\ K_*-1\},
\end{equation*}
and let $k^{\dagger}$ be the largest attainable level $\le\hat\kappa$. Let $(g^{\dagger},N_T^{\dagger})$ be the member given by part~(b) of Lemma~\ref{8.8:lem-retargeting} at $k^{\diamond}=k^{\dagger}$, and put $X^{\dagger}:=(g^{\dagger})^{-2}$. Put
\begin{equation*}
K^{\dagger\prime}:=k^{\dagger}-m_{\mathbb{T}}(g^{\dagger}),\qquad
N^{\dagger}:=N_T^{\dagger}L^{m_{\mathbb{T}}(g^{\dagger})}.
\end{equation*}

For a level $j$ put
\begin{equation*}
t_j:=\theta'\gamma_0A_1\bigl(p_0(g_j)-1\bigr)^2 .
\end{equation*}

Assume the two ceilings \eqref{8.8:eq-coupling-ceilings-a} and \eqref{8.8:eq-coupling-ceilings-b}. Assume also the following statements.
\begin{enumerate}
\item[(h1)] Lemma~\ref{8.8:lem-target-independence}, independence of the steps from the target coupling, whose proof is incomplete at one step.
\item[(h2)] Lemma~\ref{8.8:lem-retargeting}, re-targeting a member to an attainable level.
\item[(h3)] Independence of the stopping level: the $\tilde\nu$-mass of $L^{(j)}_{\mathfrak{c}}$ does not depend on the level at which the run stops.
\item[(h4)] The terminal reserve bound, in the form used below. For a member whose longest run has $N=N_TL^{m}$ sites per side, in the sense of Notation~\ref{8.8:not-bare-data-members}, and an event $E$ of histories of that run on which the terminal reserve $\Theta^{\mathrm{term}}$ is at least $t$, it gives
\begin{equation*}
\sum_{\omega\in E}\tau_{\omega}\le e^{-t}e^{E_+^{\mathrm{lin}}N^d}.
\end{equation*}
\item[(h5)] The lower half of clause (i-b) of Theorem~\ref{3.1:thm-main}, the floor on the normaliser: for the same member, with $\Omega$ the set of histories of its longest run,
\begin{equation*}
\sum_{\omega\in\Omega}\tau_{\omega}\ge e^{-C_{\mathrm{low}}N^d}.
\end{equation*}
\item[(h6)] Lemma~\ref{8.8:lem-stability-exponent}, logarithmic growth of the global stability exponent.
\item[(h7)] The comparison of reserves at the birth coupling, together with the continuity of the flow, in the two forms used here. First, if the internal history of a component contains a seed $Z$ of the first kind with birth level $j(Z)\le j$, then
\begin{equation*}
\Theta^{\mathrm{term}}\ge\theta'\gamma_0A_1\,p_0(g_{j(Z)})^2(d'+1)\ge t_j .
\end{equation*}
Here $d'+1$ is the factor with which the comparison of reserves bounds the terminal reserve of a first-kind seed; both inequalities of the display are part of the assumed statement, and below only the resulting bound $\Theta^{\mathrm{term}}\ge t_j$ is used.
Second,
\begin{equation*}
p_0(g_j)\ge p_0(g_{k^{\dagger}})-1\ge p_0(g^{\dagger})-1 .
\end{equation*}
\item[(h8)] In case \textup{(C)} only: parts (b) and (d) of the live-event average theorem; only part (d) is invoked below. Part (d), in the form used here, gives the following for a member $(g',N_T')$ over $(x_0,\mathsf{N})$ and a level $j$ with $K(x_0,g')-j\ge u_{\gamma}(g')$. Let $n_j$ be the number of level-$j$ cells and $K'$ the terminal level of the longest run of $(g',N_T')$. Then the average over the level-$j$ cells satisfies
\begin{equation*}
n_j^{-1}\sum_{\mathfrak{c}}\tilde\nu\bigl(\mathcal{E}^{(j)}_{\mathfrak{c}}(K')\bigr)
\le\mathfrak{e}_j\le e^{-\theta p_0(g_j)} .
\end{equation*}
\item[(h9)] For the last assertion of \textup{(C)}, two further statements. The first is the healed-part bound
\begin{equation*}
\tilde\nu\bigl(L^{(j)}_{\mathfrak{c}}\setminus\mathcal{E}^{(j)}_{\mathfrak{c}}(K^{\diamond})\bigr)
\le K_He^{\theta}S_*e^{-\theta p_0(g_j)} .
\end{equation*}
The second is the second clause of part (ii) of the first-moment lemma.
\item[(h10)] For the fixed-cell form of the last assertion of \textup{(C)}: the covariance statement at level $j$, for the family of events $L^{(j)}_{\mathfrak{c}}$ over the level-$j$ cells.
\item[(h11)] For the cell-averaged form of the last assertion of \textup{(C)}: the cell-average convention, under which bounds on $\tilde\nu(L^{(j)}_{\mathfrak{c}})$ are asserted only for the average $\langle\!\langle\cdot\rangle\!\rangle$ over the level-$j$ cells.
\end{enumerate}

Then $k^{\dagger}$ exists, $K^{\dagger\prime}\ge K-m_{\mathbb{T}}(g)$, $K^{\dagger\prime}$ is the terminal level of the longest run of $(g^{\dagger},N_T^{\dagger})$, and for every level $j\le K-m_{\mathbb{T}}(g)$ the following hold.
\begin{enumerate}
\item[(V)] Suppose $K+a\le K_*-1$. The prolonged run then reaches the volume scale, the level $K+a$ at which the factor $L^{a}$ of $N_T$ is used up by re-targeting, before the coupling reaches $\gamma_*$. This is the case of every torus with $a<K_*-K$; recall that always, by part (c) of Lemma~\ref{8.8:lem-retargeting},
\begin{equation*}
K_*-K\ge\frac{X-X_*-b_+-2c_2}{3\lambda}.
\end{equation*}
In this case $N_T^{\dagger}\le\bar nL^{c_{\mathrm{gap}}}$, and for every level-$j$ cell $\mathfrak{c}$
\begin{equation}\label{8.8:eq-induction-base-1}
\tilde\nu\bigl(\mathcal{E}^{(j)}_{\mathfrak{c}}(K^{\dagger\prime})\bigr)
\le\exp\bigl(\mathsf{B}(g^{\dagger})(N^{\dagger})^d-t_j\bigr)
\le e^{-\theta p_0(g_j)} .
\end{equation}
\item[(C)] Suppose $K+a\ge K_*$. The torus is then larger than in (V), and coupling room is what is spent. In this case $k^{\dagger}-K\ge\bar u(g)$. The average over the level-$j$ cells satisfies
\begin{equation}\label{8.8:eq-induction-base-2}
n_j^{-1}\sum_{\mathfrak{c}}\tilde\nu\bigl(\mathcal{E}^{(j)}_{\mathfrak{c}}(K^{\dagger\prime})\bigr)
\le\mathfrak{e}_j\le e^{-\theta p_0(g_j)} .
\end{equation}
Hence, with the two statements of (h9),
\begin{equation*}
\tilde\nu\bigl(L^{(j)}_{\mathfrak{c}}\bigr)
\le K_He^{\theta}S_*e^{-\theta p_0(g_j)}+e^{-\theta p_0(g_j)} .
\end{equation*}
This holds for every $\mathfrak{c}$ under (h10). Under (h11) it holds in the cell-averaged form $\langle\!\langle\cdot\rangle\!\rangle$.
\end{enumerate}
\end{lemma}

\begin{proof}
The proof has two preliminary steps, common to both cases, followed by one step for each case.

\emph{The level $k^{\dagger}$.} Part (c) of Lemma~\ref{8.8:lem-retargeting} gives $K_*-K\ge(X-X_*-b_+-2c_2)/(3\lambda)$. Since $X=g^{-2}\ge\gamma_G^{-2}$, the ceiling \eqref{8.8:eq-coupling-ceilings-b} applies. Dividing it by $3\lambda$ shows that this quotient is at least $\bar u(g)+2c_{\mathrm{gap}}+1$. Hence
\begin{equation}\label{8.8:eq-induction-base-3}
K_*-K\ge\bar u(g)+2c_{\mathrm{gap}}+1 .
\end{equation}

By part (a) of Lemma~\ref{8.8:lem-retargeting}, $K$ is attainable. We have $K\le K+a$, and $K\le K_*-1$ by \eqref{8.8:eq-induction-base-3}, so $K\le\hat\kappa$. The attainable levels $\le\hat\kappa$ therefore form a non-empty finite set, and its largest element $k^{\dagger}$ exists, with $k^{\dagger}\ge K$.

Put $k_1:=\max\{K,\hat\kappa-c_{\mathrm{gap}}\}$. Then $\hat\kappa-c_{\mathrm{gap}}\le K_*-1-c_{\mathrm{gap}}$. Moreover $K\le K_*-c_{\mathrm{gap}}-1$ by \eqref{8.8:eq-induction-base-3}. Hence $k_1\in[K,K_*-c_{\mathrm{gap}}-1]$. Two cases arise.
\begin{itemize}
\item If $k_1=K$, then $K\ge\hat\kappa-c_{\mathrm{gap}}$, and $K$ itself is an attainable level in $[\hat\kappa-c_{\mathrm{gap}},\hat\kappa]$.
\item If $k_1=\hat\kappa-c_{\mathrm{gap}}$, part (a) of Lemma~\ref{8.8:lem-retargeting} gives an attainable level between $k_1$ and $k_1+c_{\mathrm{gap}}=\hat\kappa$.
\end{itemize}
In either case, by the maximality of $k^{\dagger}$,
\begin{equation}\label{8.8:eq-induction-base-4}
\hat\kappa-k^{\dagger}\le c_{\mathrm{gap}} .
\end{equation}

\emph{The prolonged member.} Since $k^{\dagger}\le\hat\kappa\le K+a$, we have $k^{\dagger}-K\le a$. Hence $L^{k^{\dagger}-K}$ divides $N_T=n_0L^a$. Also $K\le k^{\dagger}\le K_*-1$. So part (b) of Lemma~\ref{8.8:lem-retargeting} applies at $k^{\diamond}=k^{\dagger}$, and gives the following:
\begin{itemize}
\item $(g^{\dagger},N_T^{\dagger})$ is a member over the same bare datum $(x_0,\mathsf{N})$;
\item $K(x_0,g^{\dagger})=k^{\dagger}$, and $g\le g^{\dagger}<\gamma_*$;
\item $N_T^{\dagger}=N_TL^{-(k^{\dagger}-K)}$, and $m_{\mathbb{T}}(g^{\dagger})\le m_{\mathbb{T}}(g)$.
\end{itemize}
Consequently $K^{\dagger\prime}=k^{\dagger}-m_{\mathbb{T}}(g^{\dagger})\ge K-m_{\mathbb{T}}(g)$, since $k^{\dagger}\ge K$. By Notation~\ref{8.8:not-bare-data-members}, the admissible runs of $(g^{\dagger},N_T^{\dagger})$ stop at $k^{\dagger}-m$ with $m\ge m_{\mathbb{T}}(g^{\dagger})$, on a lattice of $N_T^{\dagger}L^{m}$ sites per side; so its longest run ends at $K^{\dagger\prime}$, on a lattice of $N^{\dagger}$ sites per side.

Now consider levels $j\le K-m_{\mathbb{T}}(g)\le K^{\dagger\prime}$. We use Lemma~\ref{8.8:lem-target-independence}, whose proof is incomplete at one step, together with the independence of the stopping level (h3). They show that the events $L^{(j)}_{\mathfrak{c}}$ at these levels are those of the original member, and that so are their $\tilde\nu$-masses. The fate decomposition \eqref{8.8:eq-induction-base-a} may therefore be taken at $K^{\diamond}=K^{\dagger\prime}$.

\emph{Case (V): $\hat\kappa=K+a$.} Then $N_T^{\dagger}=n_0L^{a-(k^{\dagger}-K)}$. Here $a-(k^{\dagger}-K)=\hat\kappa-k^{\dagger}\le c_{\mathrm{gap}}$ by \eqref{8.8:eq-induction-base-4}, and $n_0\le\bar n$, so $N_T^{\dagger}\le\bar nL^{c_{\mathrm{gap}}}$. The longest run of $(g^{\dagger},N_T^{\dagger})$ ends on a lattice of $N^{\dagger}$ sites per side, and $L^{m_{\mathbb{T}}(g^{\dagger})}$ is polylogarithmic in $X^{\dagger}$.

On $\mathcal{E}^{(j)}_{\mathfrak{c}}(K^{\dagger\prime})$, consider the component that is live at the terminal level. Its internal history contains a seed $Z$ of the first kind with birth level $j(Z)\le j$ (see the remarks before \eqref{8.8:eq-induction-base-a}). By the first form of (h7), $\Theta^{\mathrm{term}}\ge\theta'\gamma_0A_1p_0(g_{j(Z)})^2(d'+1)\ge t_j$ on this event.

We apply the terminal reserve bound (h4) with $E:=\mathcal{E}^{(j)}_{\mathfrak{c}}(K^{\dagger\prime})$ and $t:=t_j$, and the floor (h5), both for the member $(g^{\dagger},N_T^{\dagger})$. Since $\tilde\nu(E)$ is the total integral of the histories in $E$ divided by the normaliser, we obtain
\begin{equation*}
\tilde\nu\bigl(\mathcal{E}^{(j)}_{\mathfrak{c}}(K^{\dagger\prime})\bigr)
\le\frac{e^{-t_j}\,e^{E_+^{\mathrm{lin}}(N^{\dagger})^d}}{e^{-C_{\mathrm{low}}(N^{\dagger})^d}}
=\exp\bigl(\mathsf{B}(g^{\dagger})(N^{\dagger})^d-t_j\bigr).
\end{equation*}
The last equality is the definition $\mathsf{B}(g^{\dagger})=E_+^{\mathrm{lin}}+C_{\mathrm{low}}$ of Lemma~\ref{8.8:lem-stability-exponent}. This is the first inequality of \eqref{8.8:eq-induction-base-1}.

Since $g^{\dagger}<\gamma_*$, we have $X^{\dagger}\ge X_*$. Moreover $N^{\dagger}\le\bar nL^{c_{\mathrm{gap}}}L^{m_{\mathbb{T}}(g^{\dagger})}$. Lemma~\ref{8.8:lem-stability-exponent} and the ceiling \eqref{8.8:eq-coupling-ceilings-a} at $X'=X^{\dagger}$ therefore give
\begin{equation*}
\mathsf{B}(g^{\dagger})(N^{\dagger})^d
\le\tfrac14\theta'\gamma_0A_1\bigl(p_0(g^{\dagger})-2\bigr)^2
\le\tfrac14t_j .
\end{equation*}
The last inequality uses $p_0(g_j)\ge p_0(g_{k^{\dagger}})-1\ge p_0(g^{\dagger})-1$, the second form of (h7). Hence
\begin{equation*}
\tilde\nu\bigl(\mathcal{E}^{(j)}_{\mathfrak{c}}(K^{\dagger\prime})\bigr)\le e^{-\frac34t_j}.
\end{equation*}
By the second clause of \eqref{8.8:eq-coupling-ceilings-a}, applied with $p=p_0(g_j)$, we have $\frac14t_j\ge\theta p_0(g_j)$. That clause applies: since $j\le K$, we have $x_j>X$, so that
\begin{equation*}
g_j<g\le g^{\dagger}<\gamma_* ;
\end{equation*}
as the degree function $p_0(\cdot)$ increases when its argument decreases, $p_0(g_j)\ge p_0(\gamma_*)\ge p_0(\gamma_*)-1$. Hence $\frac34t_j\ge\theta p_0(g_j)$, which proves the second inequality of \eqref{8.8:eq-induction-base-1}.

No covariance, no first moment and no averaging over cells has entered, so the bound holds for every cell as it stands. The coupling $g^{\dagger}$ may be far above $g$, up to $\gamma_*$. This does no harm: both $\mathsf{B}(g^{\dagger})$ and $N^{\dagger}$ decrease as the coupling increases, while $t_j$ is read at the birth coupling.

\emph{Case (C): $\hat\kappa=K_*-1$.} By \eqref{8.8:eq-induction-base-4} and \eqref{8.8:eq-induction-base-3},
\begin{equation*}
k^{\dagger}-K\ge K_*-1-c_{\mathrm{gap}}-K\ge\bar u(g)+c_{\mathrm{gap}}\ge\bar u(g)\ge u_{\gamma}(g^{\dagger}).
\end{equation*}
The last inequality holds because $\bar u(g)$ is the supremum of $u_{\gamma}$ over $[g,\gamma_*]$, and $g\le g^{\dagger}<\gamma_*$. This proves the first assertion of (C).

For $j\le K$ it follows that
\begin{equation*}
K(x_0,g^{\dagger})-j=k^{\dagger}-j\ge u_{\gamma}(g^{\dagger}).
\end{equation*}
Part (d) of the live-event average theorem, hypothesis (h8), applies to the member $(g^{\dagger},N_T^{\dagger})$, whose longest run ends at $K^{\dagger\prime}$. It bounds the cell average of the masses $\tilde\nu(\mathcal{E}^{(j)}_{\mathfrak{c}}(K^{\dagger\prime}))$ of the live events as in \eqref{8.8:eq-induction-base-2}. The bound \eqref{8.8:eq-induction-base-2} is a bound in the prolonged run, in which, by the second preliminary step, the fate decomposition \eqref{8.8:eq-induction-base-a} of $L^{(j)}_{\mathfrak{c}}$ is taken.

The passage from this cell average to a fixed cell is made under the covariance statement at level $j$ (h10). It applies to the family of events $L^{(j)}_{\mathfrak{c}}$, which is covariant, and not to the family $\mathcal{E}^{(j)}_{\mathfrak{c}}(K^{\dagger\prime})$, which is not. This passage, combined with the fate decomposition \eqref{8.8:eq-induction-base-a} at $K^{\diamond}=K^{\dagger\prime}$, the healed-part bound of (h9) and the second statement of (h9), gives the last assertion of (C). Under the cell-average convention (h11), the same combination gives it in the cell-averaged form.
\end{proof}

\begin{corollary}[The flat reserve at the volume scale]\label{8.8:cor-flat-reserve}
Assume the following statements:
\begin{enumerate}
\item[(a)] clause (i-d) of Theorem~\ref{3.1:thm-main}, with its reserve weight $w_{\theta}$ and
the two reserve statements of that clause;
\item[(b)] the lower stability bound of clause (i-b) of Theorem~\ref{3.1:thm-main}, with its
constant $C_{\mathrm{low}}$;
\item[(c)] the proposition reducing the live-cell bound \eqref{8.8:eq-graded-rarity-a} to the fate
decomposition \eqref{8.8:eq-induction-base-a};
\item[(d)] the torus lemma of the induction;
\item[(e)] rarity of young large-field blocks.
\end{enumerate}
Assume moreover the following flat ceiling: for every $X'\ge X_*$, with $g':=X'^{-1/2}$,
\begin{equation}\label{8.8:eq-flat-reserve-1}
C_{\mathsf{B}}\,(1+\log X')\,\bigl(\bar n\,L^{c_{\mathrm{gap}}}\,L^{m_{\mathbb{T}}(g')}\bigr)^{d}
\le \tfrac12\,\theta\,\bigl(p_0(g')-2\bigr).
\end{equation}
Consider the volume-scale case (Case~V) of Lemma~\ref{8.8:lem-induction-base}. Let $j$,
$\mathfrak{c}$, $K^{\dagger\prime}$ and $N^{\dagger}$ be as in that lemma. Then the live event
$\mathcal{E}^{(j)}_{\mathfrak{c}}(K^{\dagger\prime})$ of the fate decomposition
\eqref{8.8:eq-induction-base-a} satisfies
\begin{equation}\label{8.8:eq-flat-reserve-2}
\tilde{\nu}\bigl(\mathcal{E}^{(j)}_{\mathfrak{c}}(K^{\dagger\prime})\bigr)
\le e^{\theta/2}\,e^{-\frac12\theta p_0(g_j)} ,
\end{equation}
and the live-cell bound \eqref{8.8:eq-graded-rarity-a} holds at level $j$ with the constants
\begin{equation}\label{8.8:eq-flat-reserve-3}
c_{\mathrm{lv}}=\tfrac12\theta,\qquad C_{\mathrm{lv}}=K_H\,e^{\theta}\,S_*+e^{\theta/2}.
\end{equation}
\end{corollary}

\begin{proof}
We use the reserve weight $w_{\theta}$ of clause (i-d) of Theorem~\ref{3.1:thm-main}. Its two
reserve statements are applied termwise, in the same way as in the proof of
Lemma~\ref{8.8:lem-stability-exponent}.

For a history in the live event $\mathcal{E}^{(j)}_{\mathfrak{c}}(K^{\dagger\prime})$, let $Z$ be the
seed region of the live component containing $\mathfrak{c}$, and let $j(Z)$ be the level at which
$Z$ was created. The reserve weight of such a history satisfies
\begin{equation}\label{8.8:eq-flat-reserve-4}
w_{\theta}\le e^{-\theta p_0(g_{j(Z)})}\le e^{-\theta(p_0(g_j)-1)} .
\end{equation}

We now bound the measure of the event. The reserve statements of clause (i-d) are applied to each
term on the torus of side $N^{\dagger}$. The reserve weight is split off, and the sum of the
remaining majorants is bounded by $e^{E_+^{(\theta)}(N^{\dagger})^{d}}$, by the second reserve
statement of clause (i-d), cited as in Lemma~\ref{8.8:lem-stability-exponent}. The normaliser
of $\tilde{\nu}$ is bounded
below by the lower stability bound of clause (i-b), which contributes the factor
$e^{C_{\mathrm{low}}(N^{\dagger})^{d}}$. Together with \eqref{8.8:eq-flat-reserve-4} this gives
\begin{equation}\label{8.8:eq-flat-reserve-5}
\tilde{\nu}\bigl(\mathcal{E}^{(j)}_{\mathfrak{c}}(K^{\dagger\prime})\bigr)
\le \exp\Bigl(\bigl(E_+^{(\theta)}+C_{\mathrm{low}}\bigr)(N^{\dagger})^{d}
-\theta\bigl(p_0(g_j)-1\bigr)\Bigr).
\end{equation}

The right member of \eqref{8.8:eq-flat-reserve-5} is at most
$e^{\theta/2}e^{-\frac12\theta p_0(g_j)}$ exactly when
\begin{equation}\label{8.8:eq-flat-reserve-6}
\bigl(E_+^{(\theta)}+C_{\mathrm{low}}\bigr)(N^{\dagger})^{d}\le \tfrac12\theta\bigl(p_0(g_j)-1\bigr).
\end{equation}
Indeed, subtracting $\theta(p_0(g_j)-1)$ from both sides of \eqref{8.8:eq-flat-reserve-6} gives the
exponent bound $\tfrac12\theta-\tfrac12\theta p_0(g_j)$.

It remains to derive \eqref{8.8:eq-flat-reserve-6} from the flat ceiling
\eqref{8.8:eq-flat-reserve-1}. Put $g':=g^{\dagger}$ and $X':=(g^{\dagger})^{-2}$, where
$g^{\dagger}$ is the coupling of the member fixed in Lemma~\ref{8.8:lem-induction-base}. As in
Case~V of the proof of that lemma, $X'\ge X_*$, so that \eqref{8.8:eq-flat-reserve-1} applies at
this $X'$. By Lemma~\ref{8.8:lem-stability-exponent}, together with the second reserve statement
of clause (i-d), whose exponent $E_+^{(\theta)}$ has the shape of clause (i-b),
$E_+^{(\theta)}+C_{\mathrm{low}}\le C_{\mathsf{B}}(1+\log X')$; as in Case~V of the proof of
Lemma~\ref{8.8:lem-induction-base},
\begin{equation*}
N^{\dagger}\le \bar n\,L^{c_{\mathrm{gap}}}\,L^{m_{\mathbb{T}}(g')},\qquad
p_0(g')-2\le p_0(g_j)-1 .
\end{equation*}
Multiplying the bound on $E_+^{(\theta)}+C_{\mathrm{low}}$ by the $d$-th power of the bound on
$N^{\dagger}$ and applying \eqref{8.8:eq-flat-reserve-1}, the left member of
\eqref{8.8:eq-flat-reserve-6} is at most $\tfrac12\theta(p_0(g')-2)$; by the inequality
$p_0(g')-2\le p_0(g_j)-1$ this is at most $\tfrac12\theta(p_0(g_j)-1)$.
Hence \eqref{8.8:eq-flat-reserve-6} holds, and \eqref{8.8:eq-flat-reserve-2} follows.

The flat ceiling \eqref{8.8:eq-flat-reserve-1} is one-sided, and it is met as soon as $X_*$ is large
enough. To see this, compare the growth of the two members in $\log X'$.
\begin{itemize}
\item The left member is of degree $d\,r_{\mathrm{pr}}+1$ in $\log X'$: the factor
$(L^{m_{\mathbb{T}}(g')})^{d}=(L^{11+m'}R(g'))^{d}$ contributes degree $d\,r_{\mathrm{pr}}$, and
the factor $1+\log X'$ one more.
\item The right member is of degree $\mathfrak{p}_0$, since $p_0(g')=A_0(\log X')^{\mathfrak{p}_0}$.
\end{itemize}
With $d=4$ the left degree is $4r_{\mathrm{pr}}+1$. Because $r_{\mathrm{pr}}\ge 2$ and
$\mathfrak{p}_0\ge 5r_{\mathrm{pr}}$, we have
\begin{equation*}
4r_{\mathrm{pr}}+1<5r_{\mathrm{pr}}\le\mathfrak{p}_0 .
\end{equation*}
Hence the right member of \eqref{8.8:eq-flat-reserve-1} exceeds the left member for all
$X'\ge X_*$ once $X_*$ is large enough; the condition bounds $X_*$ from below only.

The remaining steps are those of the proof of Lemma~\ref{8.8:lem-induction-base} in Case~V. Here
\eqref{8.8:eq-flat-reserve-2} serves as the bound on the event
$\mathcal{E}^{(j)}_{\mathfrak{c}}(K^{\dagger\prime})$ of the fate decomposition
\eqref{8.8:eq-induction-base-a}. These steps use hypotheses (c), (d) and (e). The healed summands
$H_k$ of \eqref{8.8:eq-induction-base-a}, bounded as in Lemma~\ref{8.8:lem-induction-base},
contribute $K_H e^{\theta}S_*$ to the constant, and \eqref{8.8:eq-flat-reserve-2} contributes
$e^{\theta/2}$. This gives the live-cell bound \eqref{8.8:eq-graded-rarity-a} at level $j$ with the
constants \eqref{8.8:eq-flat-reserve-3}. Apart from hypotheses (a)--(e), the flat ceiling
\eqref{8.8:eq-flat-reserve-1}, Lemma~\ref{8.8:lem-stability-exponent}, through which the reserve
statements are cited termwise, and the comparisons of Case~V of the proof of
Lemma~\ref{8.8:lem-induction-base} used above, the argument uses none of the further inputs of
Lemma~\ref{8.8:lem-induction-base}; in particular the reserve statement in linear form, with
exponent $E_+^{\mathrm{lin}}$, is not used.

Finally, the rate $c_{\mathrm{lv}}=\tfrac12\theta$ is adequate for the later uses of the live-cell
bound: the statements that consume it admit any fixed rate $c_q\in(0,\theta]$.
\end{proof}

\begin{remark}[What re-targeting changes; a divisibility restriction]\label{8.8:rem-divisibility}
We set out what re-targeting a member to an attainable level (Lemma~\ref{8.8:lem-retargeting}) and the base of the induction (Lemma~\ref{8.8:lem-induction-base}, which is proved as an implication from the statements it names) change, and what they leave unchanged. Case~V (the volume scale) and Case~C (coupling room) are the two cases of Lemma~\ref{8.8:lem-induction-base}.
\begin{itemize}
\item[(i)] Let $K:=K(x_0,g)$, and let $N^{\dagger}$ be the terminal number of sites per side of the prolonged run, as in Lemma~\ref{8.8:lem-induction-base}. In Case~V the steps between the level $K$ and the volume scale are Ba{\l}aban steps: they are the renormalisation steps of another member of the same family over the bare datum $(x_0,\mathsf{N})$, and no new estimate enters. The final system is finite-dimensional in the sense that, with the degree function $p_0(\cdot)$,
\begin{equation*}
N^{\dagger d}\,\log g^{-2}\ll p_0(g)^2 .
\end{equation*}
\item[(ii)] Two statements made at a fixed terminal volume, that is, for a fixed number of sites per side at the terminal level, are not affected. The first is the assertion, made with the large-field block bound, healed part, that no choice of the depth floor cures the obstruction treated there. The second is the sharpness assertion of the fixed-volume sharpness lemma, which is likewise stated at a fixed terminal volume.
\item[(iii)] The ceiling \eqref{8.8:eq-coupling-ceilings-b} is used in Case~C only; Case~V needs $g<\gamma_*$ and nothing else. The following statements are affected: the two-point witness theorem, the Gaussian extraction, and the additive assembly theorem. In each of them the $g$-dependent condition of the depth floor is replaced by the ceiling $\gamma_W\le\gamma_G$. That condition is $s-r(s)\ge u_{\gamma}(g)$, together with the corresponding depth condition of the two-point witness theorem. On the tori of side $2L^m$ of \textup{(H1)} the ceiling $\gamma_W\le\gamma_G$ suffices as it stands. The other depth floors, those not named in this item, are not affected.
\item[(iv)] \emph{A divisibility restriction.} Write the terminal torus count as $N_T=n_0L^{a}$ with $a\ge0$ and $L\nmid n_0$. Lemma~\ref{8.8:lem-induction-base} is stated under the restriction
\begin{equation}\label{8.8:eq-divisibility-a}
N_T=n_0L^{a},\qquad L\nmid n_0,\qquad n_0\le\bar n ;
\end{equation}
If $n_0>\bar n$, the prolongation stops on a lattice of $n_0L^{m_{\mathbb{T}}(g^{\dagger})}$ sites per side, where $g^{\dagger}$ is the coupling of the last re-targeted member; Case~V is then not available, and Case~C needs $a\ge K_*-K$, or at least $a\ge\bar u(g)+c_{\mathrm{gap}}$. For smaller $a$ the depth floor survives, for a large-field region created at step $j$, in the form
\begin{equation*}
a+(K-j)\ge u_{\gamma}(g).
\end{equation*}
This is a restriction of the freedom to take any terminal count $N_T\ge1$. It is not a restriction of \textup{(H1)}, and it is not a cap on the volume.
\end{itemize}
\end{remark}

\begin{proof}
\emph{(i).} We first show that the intermediate steps are Ba{\l}aban steps. By Lemma~\ref{8.8:lem-retargeting}\,(b), the re-targeted member $(g^{\dagger},N_T^{\dagger})$ of Lemma~\ref{8.8:lem-induction-base} is a member over the same bare datum $(x_0,\mathsf{N})$ as $(g,N_T)$. Its steps are therefore steps of Ba{\l}aban's construction for another member of the same family, so no new estimate is used.

We next show finite-dimensionality. The torus exponent is $m_{\mathbb{T}}(g)=11+m'+\log_L R(g)$ and $M=L^{m'}$, so that
\begin{equation*}
L^{m_{\mathbb{T}}(g)}\asymp L^{11}MR(g).
\end{equation*}
Hence the minimal admissible torus has a number of sites per side of the order $L^{11}MR(g)$. We have $N^{\dagger}=N_T^{\dagger}L^{m_{\mathbb{T}}(g^{\dagger})}$, where $N_T^{\dagger}\le\bar nL^{c_{\mathrm{gap}}}$ by the choice of $k^{\dagger}$ in Lemma~\ref{8.8:lem-induction-base}, and $m_{\mathbb{T}}(g^{\dagger})\le m_{\mathbb{T}}(g)$ by Lemma~\ref{8.8:lem-retargeting}\,(b); therefore
\begin{equation*}
N^{\dagger}\le\bar nL^{c_{\mathrm{gap}}}L^{m_{\mathbb{T}}(g)} .
\end{equation*}
The left-hand side $N^{\dagger d}\log g^{-2}$ is then polylogarithmic in $g^{-2}$. The ceiling \eqref{8.8:eq-coupling-ceilings-a}, read at $X'=g^{-2}$ (so that $g'=g$), which is admissible since $g<\gamma_*$, bounds a multiple of
\begin{equation*}
(1+\log g^{-2})\,\bigl(\bar n L^{c_{\mathrm{gap}}}L^{m_{\mathbb{T}}(g)}\bigr)^d
\end{equation*}
by a multiple of $(p_0(g)-2)^2$; together with the bound on $N^{\dagger}$ this gives an inequality of the displayed form.

\emph{(ii).} Both statements concern a fixed terminal volume. Lemma~\ref{8.8:lem-retargeting} changes the terminal volume, since $N_T^{\diamond}=N_TL^{-(k^{\diamond}-K)}$. In those statements the longest admissible run was taken inside one member, with $m=m_{\mathbb{T}}(g)$ at fixed $g$. The bare datum admits longer runs, namely those of the re-targeted members of Lemma~\ref{8.8:lem-retargeting}. Neither statement speaks about these longer runs.

\emph{(iii).} The condition \eqref{8.8:eq-coupling-ceilings-b} is read only in Case~C of Lemma~\ref{8.8:lem-induction-base}. In Case~V that lemma uses only $g<\gamma_*$.

Lemma~\ref{8.8:lem-induction-base} holds for every $g\le\gamma_G$, with no depth condition on the coupling; this is why, in the statements listed in item~(iii), the $g$-dependent condition of the depth floor can be replaced by the ceiling $\gamma_W\le\gamma_G$. On the tori of side $2L^m$ of \textup{(H1)} the terminal count satisfies \eqref{8.8:eq-divisibility-a} with $n_0=2$, as shown in~(iv) below, so no further condition enters there.

\emph{(iv).} Let $N_T=n_0L^{a}$ with $L\nmid n_0$ and $n_0>\bar n$. The re-targeting of Lemma~\ref{8.8:lem-retargeting}\,(b) requires the hypothesis $L^{k^{\diamond}-K}\mid N_T$. Only powers of $L$ can therefore be divided out of $N_T$, and the prolongation stops on a lattice of $n_0L^{m_{\mathbb{T}}(g^{\dagger})}$ sites per side.

Case~V fails. The count $N^{\dagger}$ contains the factor $n_0$, so the volume is free through $n_0$ and the inequality of~(i) is not available.

Case~C needs $a\ge K_*-K$, or at least $a\ge\bar u(g)+c_{\mathrm{gap}}$. The second condition already gives clause~(d) of the theorem with the two-rate bound on which Case~C of Lemma~\ref{8.8:lem-induction-base} rests.

For smaller $a$, that theorem gives its two-rate bound with the depth counted from $K+a$; for a large-field region created at step $j$ this is the surviving form of the floor displayed in item~(iv).

This is why \eqref{8.8:eq-divisibility-a} imposes $n_0\le\bar n$. \textup{(H1)} prescribes tori of side $2L^m$, so there $n_0=2$; the restriction \eqref{8.8:eq-divisibility-a} is therefore not a restriction of \textup{(H1)}. The exponent $a$ is free, so \eqref{8.8:eq-divisibility-a} is not a cap on the volume.

Treating an arbitrary $L$-free part $n_0>\bar n$ is the live-tail problem at the terminal level; it is an open problem off the path of Theorem~\ref{3.1:thm-main}, and nothing in this remark contributes to it.
\end{proof}

\begin{theorem}[Rarity of live cells at every admissible depth]\label{8.8:thm-live-cell-rarity}
We use the normalised annealed history measure $\tilde{\nu}$, the live-cell events
$L^{(j)}_{\mathfrak{c}}$, the creation events $E^{\kappa}_{j,\square}$ and $E^{\kappa}_{j,\square}(\lambda)$
and the field events $F_{j,\square}(\lambda)$ of Definition~\ref{8.8:not-history-measure}, and the bare
data, members, attainable levels and admissible runs of Definition~\ref{8.8:not-bare-data-members}.
Let $\theta$ be the rate number, fixed at Stage~0 of the order of choice
(Definition~\ref{2.3:not-stages-first}). For a kind $\kappa$ and a $\kappa$-cube $\square$ of step $j$, let
$n_{\square}$ be the number of level-$j$ cells meeting $\square$. Let $K_H=K_H(d)$, depending on $d$ only,
be a bound on the number of boxes that contain a given cell, the boxes being those that enter the first
of the two conditions of \cite[\S1]{B-CMP122a} defining
the components healed at a step; let $\lambda^{\sharp}$ be the flow slope of
Definition~\ref{0.2:not-glossary}, and put
\begin{equation*}
  C_L:=K_H\,e^{\theta}\Bigl(1+\frac{1}{\lambda^{\sharp}}\Bigr)+1 .
\end{equation*}
Assume the conditions on $\theta$ and on the auxiliary parameters $\mathfrak{p}$ under which the
large-field block bound, healed part, of (a) is stated, and assume moreover that, for every bare
datum and every level $k$ of every run over it,
\begin{equation}\label{8.8:eq-live-cell-rarity-9}
  (2-\theta)(1+\beta_0)^{-1}p_0(g_k)-C_{\delta}(d)\,\alpha\,(101\,\check{R}_k)^{d}
  \ge\tfrac32\,p_0(g_k),
\end{equation}
where $C_{\delta}(d)$ is the strip constant, $\alpha$ the radius parameter of $\mathfrak{p}$ and
$\check{R}_k=R(\min_{i\le k}x_i)$ the running radius; this is the form in which the condition on
$\theta$ of the healed-part proposition of (a1) below is used in the proof; assume also the ceilings
on $\gamma$ under which
Lemma~\ref{8.8:lem-stability-exponent} holds, and the two coupling ceilings
\eqref{8.8:eq-coupling-ceilings-a} and \eqref{8.8:eq-coupling-ceilings-b} of
Definition~\ref{8.8:def-coupling-ceilings}. Assume finally the following statements:
\begin{enumerate}
\item[(a)] the large-field block bound, healed part, and clause (i-d) of
  Theorem~\ref{3.1:thm-main}; the flow window, that is the two-sided bound on $x_i-x_j$ stated in
  Definition~\ref{8.8:not-bare-data-members}; and the bound on the terminal reserve
  and the lower bound of clause (i-b) of Theorem~\ref{3.1:thm-main}, both also
  among the hypotheses of Lemma~\ref{8.8:lem-induction-base};
\item[(a1)] the statement on the reference flow $(x_k)$ that stands beside the flow window among the
  premises of the \emph{healed-part proposition}, the proposition of this chapter bounding the
  $\tilde{\nu}$-mass of the part of $L^{(j)}_{\mathfrak{c}}$ that is healed before the terminal level
  of the run;
\item[(a2)] the two lemmas of this chapter that are premises of the healed-part proposition
  besides the inputs listed in (a), (a1) and (b)--(e);
\item[(b)] the non-negativity of the total integrals $\tau_{\omega}$ and the identity
  $Z=\int\rho_k\,dV_k$ for every $k$, both stated in Definition~\ref{8.8:not-history-measure};
\item[(c)] for every kind $\kappa$, every step $j$ and every $\kappa$-cube $\square$ of step $j$: on
  $E^{\kappa}_{j,\square}$ every level-$j$ cell meeting $\square$ lies in a component that is live at
  level $j$; and the large-field regions of a history only grow until the component containing them
  is healed as a whole, so that, if the live component containing a level-$j$ cell, followed through
  the later steps, is healed at a step $k>j$, the component healed at step $k$ (an \emph{island} of
  step $k$) contains the level-$k$ cell containing that cell;
\item[(d)] every operation of the chain after a step $k$ is a positive linear operator whose placement
  is determined by the outer data and not by the islands healed at step $k$;
\item[(e)] there is a constant $S_*\le 1+1/\lambda^{\sharp}$, depending on $\lambda^{\sharp}$ only,
  such that for every level $j$: the set $\mathcal{S}_j$ of internal histories (of a component: the
  sequences of its large-field data at the earlier steps) each of which contains a seed of birth level
  at most $j$, that is, a large-field region created at a step at most $j$, has reserve weight
  \begin{equation*}
    w_{\theta}(\mathcal{S}_j)\le e^{\theta}e^{-\theta p_0(g_j)},
  \end{equation*}
  and
  \begin{equation*}
    \sum_{k>j}e^{-\frac32 p_0(g_k)}\le S_* ;
  \end{equation*}
\item[(f)] the hypotheses from which Lemma~\ref{8.8:lem-induction-base} is derived, and, in the second
  of the two cases of that lemma only, the covariance statement at level $j$ or the cell-averaged
  statement of this chapter (either suffices), which transfer a bound on the average over the
  level-$j$ cells to the individual events $L^{(j)}_{\mathfrak{c}}$;
\item[(g)] Lemma~\ref{8.8:lem-target-independence} (stated; proof incomplete at one step), and the
  torus-independence statement (the step at scale $k$ does not depend on the torus exponent $m$,
  that is on the stopping level $K-m$); and
  Lemma~\ref{8.8:lem-stability-exponent}.
\end{enumerate}
Fix $\bar n\ge 1$ before $\gamma$ in the order of choice. Then for every $g\le\gamma_G$, every
$x_0>g^{-2}$, with $K=K(x_0,g)$, every $N_T=n_0L^{a}$ with $1\le n_0\le\bar n$ and $a\ge 0$, every
admissible run of the member $(g,N_T)$, every level $j$ of it ($j\le K-m_{\mathbb{T}}(g)$), every level-$j$
cell $\mathfrak{c}$, every kind $\kappa$ and every $\kappa$-cube $\square$ of step $j$,
\begin{align}
  \tilde{\nu}\bigl(L^{(j)}_{\mathfrak{c}}\bigr)
    &\le C_L\,e^{-\theta p_0(g_j)}, \label{8.8:eq-live-cell-rarity-1}\\
  \tilde{\nu}\bigl(E^{\kappa}_{j,\square}\bigr)
    &\le n_{\square}C_L\,e^{-\theta p_0(g_j)}, \label{8.8:eq-live-cell-rarity-a}\\
  \tilde{\nu}\bigl(F_{j,\square}(1)\bigr)
    &\le n_{\square}C_L\,e^{-\theta p_0(g_j)}, \label{8.8:eq-live-cell-rarity-2}
\end{align}
and, for $\kappa\in\{\mathrm{S}1,\mathrm{S}2,\mathrm{S}3\}$, $j\ge1$ and all $\lambda\ge1$,
\begin{equation}\label{8.8:eq-live-cell-rarity-3}
  \tilde{\nu}\bigl(E^{\kappa}_{j,\square}(\lambda)\bigr)
    \le n_{\square}C_L\,e^{-(\theta/\bar\lambda_{\kappa}^{2})\lambda^{2}p_0(g_j)} ,
\end{equation}
with $\bar\lambda_{\kappa}$ the bounds of Lemma~\ref{8.8:lem-grading-bounded}. The constants are the
displayed ones, uniform in $K$, in $a$ (the volume), in the position and in the kind; no condition on
$K-j$ is imposed. On any bounded range $1\le\lambda\le\bar\lambda$ the field form follows as well:
\begin{equation}\label{8.8:eq-live-cell-rarity-4}
  \tilde{\nu}\bigl(F_{j,\square}(\lambda)\bigr)
    \le n_{\square}C_L\,e^{-(\theta/\bar\lambda^{2})\lambda^{2}p_0(g_j)} .
\end{equation}
In the notation of Theorem~\ref{3.1:thm-main}: under hypotheses (H1)--(H6) of
Definition~\ref{2.2:not-hypotheses}, with $g\le\gamma_G$, where $\gamma_G\le\gamma_J$ is a threshold on
the reference coupling, fixed after $\gamma_J$ in the order of choice
(Definition~\ref{8.8:def-coupling-ceilings}), with $K$ the number of renormalisation steps of
Theorem~\ref{3.1:thm-main}, the same bounds hold at every level $j\le K$ of the run
on $\mathbb{T}_m$, for every $m\ge m_{\mathbb{T}}(g_-(g))$; the tori of (H1) are those with $n_0=2$.
\end{theorem}

\begin{proof}
Throughout, for a level $k>j$ we write $\mathfrak{c}^{(k)}$ for the level-$k$ cell containing the
level-$j$ cell $\mathfrak{c}$.

\emph{Step 1: a creation leaves its cells live; decomposition by fate.} Hypothesis (c) gives that a
creation of any kind at $(j,\square)$ leaves the cells of $\square$ in a live component at level $j$.
(In Ba{\l}aban's construction this holds for the following reasons. For the kinds tested by the
three fundamental small-field characteristic functions of step $j$,
$\square\subset P_j\cup Q_j\cup R_j\subset\Lambda_j^{c}$, where $P_j$, $Q_j$, $R_j$ are the regions on
which these three tests fail and $\Lambda_j$ is the small-field region of step $j$ in the notation of
\cite[\S1]{B-CMP122a}; and the component so created does not belong, at its birth step, to the class of
components healed at that step, by the second of the two conditions defining that class in
\cite[\S1]{B-CMP122a}, which requires that in the prescribed number of preceding renormalisation steps no
new large-field regions were created inside the component. For the preparatory kinds, the component in
which the large branch was chosen is excluded from the region $Z_j$ of \cite[\S1]{B-CMP122a}, that is, it
is not healed at step $j$.) Consequently
\begin{equation}\label{8.8:eq-live-cell-rarity-5}
  E^{\kappa}_{j,\square}\subset\bigcup_{\mathfrak{c}\cap\square\neq\emptyset}L^{(j)}_{\mathfrak{c}},
\end{equation}
a union of $n_{\square}$ events, and by Lemma~\ref{8.8:lem-graded-rarity} the same inclusion holds for
$F_{j,\square}(1)$ up to a $\tilde{\nu}$-null set.

We next decompose $L^{(j)}_{\mathfrak{c}}$ according to the fate of the lineage of $\mathfrak{c}$, that
is, of the live component containing $\mathfrak{c}$ followed through the later steps together with its
internal history. By hypothesis (c), the large-field regions of a history only grow until the component
is healed as a whole. Hence, for any later level $K^{\diamond}$ of any run prolonging the given one, the
event $L^{(j)}_{\mathfrak{c}}$ is the disjoint union \eqref{8.8:eq-induction-base-a} of the events
$\mathcal{H}_k$, $k\in\,]j,K^{\diamond}]$, where $\mathcal{H}_k$ is the event that the lineage is first
healed at step $k$,
and of the event $\mathcal{E}^{(j)}_{\mathfrak{c}}(K^{\diamond})$ that it is still live at level
$K^{\diamond}$. By the identity $Z=\int\rho_k\,dV_k$, valid for every $k$ (hypothesis (b)), the
$\tilde{\nu}$-mass of $L^{(j)}_{\mathfrak{c}}$ does not depend on the level at which the run stops. By
Lemma~\ref{8.8:lem-target-independence} (stated; proof incomplete at one step), together with the
torus-independence statement of hypothesis (g) (the step at scale $k$ does not depend on the torus
exponent $m$, that is on the stopping level $K-m$), the level $K^{\diamond}$ may be
taken to be the terminal level $K^{\dagger\prime}$ of the prolonged run of
Lemma~\ref{8.8:lem-induction-base}.

\emph{Step 2: healing attempts in which all tests pass; the conditional bound.} Fix the step $k$, the
outer datum $\mathcal{O}$, the field $V_k$ and all outer variables. By
\cite[(1.70)--(1.72)]{B-CMP122b}, the $\mathcal{O}$-part of the density before exponentiation is
\begin{equation*}
  \sum_{\mathfrak{X}}\mathbf{O}_{\mathcal{O},\mathfrak{X}}\bigl[W_{\mathcal{O}}(\mathfrak{X})\bigr],
  \qquad
  W_{\mathcal{O}}(\mathfrak{X})
  =\Bigl[\prod_{X'\in\mathfrak{X}}\mathbf{T}'_k(X')\Bigr]
   \exp\sum_{Y\in\mathcal{Y}(\mathfrak{X})}V_{\mathfrak{X}}(Y),
\end{equation*}
where $\mathfrak{X}$ is the family of islands healed at this step. Let $X$ be an island
containing $\mathfrak{c}^{(k)}$ and $\mathcal{S}$ a set of its internal histories; write
$\mathbf{T}'_{k,\mathcal{S}}(X)$ for the part of $\mathbf{T}'_k(X)$ coming from the histories in
$\mathcal{S}$, write $\mathcal{V}_1:=\sum_{Y\in\mathcal{Y}(\mathfrak{X})}V_{\mathfrak{X}}(Y)$ for the
exponent, and $\mathcal{V}_0$ for the same sum with $X$ removed from
the family. The inputs of hypothesis (a) give the following.
\begin{itemize}
\item[(P1)] By the locality of the terms $V_{\mathfrak{X}}(Y)$ in \cite[\S1]{B-CMP122b},
  $\mathcal{V}_1$ and
  $\mathcal{V}_0$ differ only in the $Y$ that meet the strip $X^{\sharp}$ of $X$ (the layer around $X$
  so denoted in \cite[\S1]{B-CMP122b}); by
  \cite[(1.68)]{B-CMP122b}
  and the bound on the number of such $Y$, the sum of these differences is at most $\delta_X$, with
  \begin{equation*}
    \delta_X\le C_{\delta}(d)\,\alpha\,(101\,\check{R}_k)^{d},
  \end{equation*}
  $\alpha$ being the radius parameter of $\mathfrak{p}$.
\item[(P2)] By \cite[(1.73)]{B-CMP122b}, $\mathbf{T}'_{k,\mathcal{S}}(X)$ has a non-negative
  density at the configuration $(U,0)$ (background $U$, zero fluctuation
  field); hence, with $t_X(\mathcal{S}):=\mathbf{T}'_{k,\mathcal{S}}(X,(U,0))1$,
  \begin{equation*}
    \mathbf{T}'_{k,\mathcal{S}}(X)\bigl(e^{\mathcal{V}_1}\bigr)
    \le t_X(\mathcal{S})\,e^{\delta_X}e^{\mathcal{V}_0}.
  \end{equation*}
\item[(P3)] The operators of the other islands are positive and commute with
  $\mathbf{T}'_{k,\mathcal{S}}(X)$.
\item[(P4)] At most $K_H$ boxes contain a given cell, by the first of the two conditions of
  \cite[\S1]{B-CMP122a} defining the components healed at a step.
\item[(P5)] By \cite[(1.89)]{B-CMP122b}, with the reserve of clause (i-d) of
  Theorem~\ref{3.1:thm-main} split off once,
  \begin{equation*}
    t_X(\mathcal{S})\le w_{\theta}(\mathcal{S})\exp\bigl(-(2-\theta)(1+\beta_0)^{-1}p_0(g_k)\bigr).
  \end{equation*}
\end{itemize}
Write $\rho_{k,\mathcal{O}}$ for the $\mathcal{O}$-part of $\rho_k$, and
$\rho^{(\mathfrak{c}^{(k)},\mathcal{S})}_{k,\mathcal{O}}$ for its part pinned at
$(\mathcal{O},\mathfrak{c}^{(k)},\mathcal{S})$, that is, the sum of the terms in which an island healed
at step $k$ contains $\mathfrak{c}^{(k)}$ and has its internal history in $\mathcal{S}$. In each such
term we apply (P2) to the factor of $X$, which is legitimate inside the product by (P3); this
replaces the factor of $X$ and the exponent $\mathcal{V}_1$ by $t_X(\mathcal{S})e^{\delta_X}$ and the
exponent
$\mathcal{V}_0$ of the term in which $X$ is absent; $t_X(\mathcal{S})$ is bounded by (P5), $\delta_X$ by
(P1), and the number of islands through which $\mathfrak{c}^{(k)}$ can be pinned by (P4).
By \eqref{8.8:eq-live-cell-rarity-9}, the resulting exponent
$-(2-\theta)(1+\beta_0)^{-1}p_0(g_k)+\delta_X$ is at most $-\tfrac32p_0(g_k)$, and this gives,
pointwise in $V_k$,
\begin{equation}\label{8.8:eq-live-cell-rarity-6}
  \rho^{(\mathfrak{c}^{(k)},\mathcal{S})}_{k,\mathcal{O}}\le\varepsilon_k(\mathcal{S})\,\rho_{k,\mathcal{O}},
  \qquad
  \varepsilon_k(\mathcal{S}):=K_H\,w_{\theta}(\mathcal{S})\,e^{-\frac32 p_0(g_k)} .
\end{equation}
In this comparison the numerator is controlled by the positivity of the Wilson action and of the
quadratic forms \cite[(1.79)]{B-CMP122b}; the denominator by the Gaussian ball at the conditional
minimiser, \cite[(1.2)--(1.10)]{B-CMP122b}, which is already contained in $\mathbf{T}'$; the comparison is
made term by term between two histories that differ in exactly one island, and no lower bound on any
activity is used.

Let $E^{(k)}_{\mathfrak{c}^{(k)},\mathcal{S}}$ be the event that an island healed at step $k$ contains
$\mathfrak{c}^{(k)}$ and has its internal history in $\mathcal{S}$. By hypothesis (d), every later
operation of the chain is a positive linear operator placed by the outer data and not by the healed
islands, so it acts on both sides of \eqref{8.8:eq-live-cell-rarity-6} in the same way and preserves the
inequality. Integrating over $V_k$ and the outer variables, summing over $\mathcal{O}$ and dividing by
$Z$ (hypothesis (b)), we obtain
\begin{equation}\label{8.8:eq-live-cell-rarity-7}
  \tilde{\nu}\bigl(E^{(k)}_{\mathfrak{c}^{(k)},\mathcal{S}}\bigr)\le\varepsilon_k(\mathcal{S}).
\end{equation}
The bounds \eqref{8.8:eq-live-cell-rarity-6} and \eqref{8.8:eq-live-cell-rarity-7} are, respectively,
the two intermediate statements
through which the healed-part proposition named in hypothesis (a1) is proved; they are distinct from
the two lemmas of hypothesis (a2).

\emph{Step 3: healing attempts in which a preparatory test fails; charge at the same step.} If at a
healing attempt one of the seven preparatory decompositions takes its large branch, no comparison is
made: the component stays live, its internal history acquires a preparatory seed, and the factor
$e^{-p_0(g_j)}$ enters the last product of \cite[(1.79)]{B-CMP122b}. By
\cite[(1.83)ff.]{B-CMP122b} the bookkeeping restarts with
\begin{equation*}
  \mathfrak{a}_{j+1}(D)=p_0(g_j)-O(1)\,M^{d}R_{j+1}^{d+1}d'_{j+1}(D),
\end{equation*}
where $D$ is the region created by the large branch, a small domain in the sense of
\cite[\S1]{B-CMP122b}. Here $\mathfrak{a}_{j+1}(D)$ is the amortisation exponent written
$\kappa_{j+1}(Z)$ in \cite[(1.83)ff.]{B-CMP122b}; we write $\mathfrak{a}$ because $\kappa$ denotes a kind
here, and $D$ to keep the unsubscripted $Z$ for the normaliser, the subscripted $Z_j$, $Z_i$ of
Step~1 and below being the region sequences so written in \cite{B-CMP122a,B-CMP122b}. Further,
$R_{j+1}$ is the radius so denoted in that display at step $j+1$ (a radius, not the region $R_j$ of
Step~1), and $d'_{j+1}(D)$ is the quantity of \cite[(1.80)--(1.88)]{B-CMP122b} attached to the domain
$D$ at step $j+1$, through which the amortisation exponent depends on the size of $D$. As stated in
\cite[(1.83)ff.]{B-CMP122b}, $\mathfrak{a}_{j+1}(D)$ satisfies \cite[(1.80)]{B-CMP122b} because $D$ is a
small domain.
All such histories are summed inside $\mathbf{T}'_k(X)$ at the step $k$ at which the lineage is finally
healed: by \cite[(1.71)]{B-CMP122b} the sum runs over the sequences
$\{\Omega^{c}_i\cap X,\ Z_i\cap X\}$, where $\Omega_i$ denotes the small-field regions so written in
\cite[(1.71)]{B-CMP122b}, unrelated to the history space $\Omega$, and $Z_i$ the regions of the same
display. The bound
(P5) is uniform in these histories. Thus the bound is obtained at the healing step itself, without
passing to a coarser level or enlarging the event.

\emph{Step 4: induction in time.} On $\mathcal{H}_k$ the island healed at step $k$ contains
$\mathfrak{c}^{(k)}$, by the nesting of regions in hypothesis (c), and its internal history
contains a seed of birth level at most $j$; that is,
$\mathcal{H}_k\subset E^{(k)}_{\mathfrak{c}^{(k)},\mathcal{S}_j}$
for the set $\mathcal{S}_j$ of hypothesis (e), for which
$w_{\theta}(\mathcal{S}_j)\le e^{\theta}e^{-\theta p_0(g_j)}$ by the first condition of hypothesis (e).
Applying \eqref{8.8:eq-live-cell-rarity-7} at each $k\in\,]j,K^{\dagger\prime}]$ and summing over $k$,
with the second condition of hypothesis (e),
\begin{equation}\label{8.8:eq-live-cell-rarity-8}
\begin{aligned}
  \tilde{\nu}\bigl(L^{(j)}_{\mathfrak{c}}\setminus\mathcal{E}^{(j)}_{\mathfrak{c}}(K^{\dagger\prime})\bigr)
  &\le K_H\,e^{\theta}e^{-\theta p_0(g_j)}\sum_{k>j}e^{-\frac32 p_0(g_k)}\\
  &\le K_H\,e^{\theta}S_*\,e^{-\theta p_0(g_j)},
  \qquad S_*\le 1+\frac{1}{\lambda^{\sharp}} .
\end{aligned}
\end{equation}
The bound \eqref{8.8:eq-live-cell-rarity-8} is the statement of the healed-part proposition named in
hypothesis (a1), whose premises are the statements of hypotheses (a), (a1), (a2) and (b)--(e).
The rate is the reserve of the oldest seed at its birth coupling, and the healing step supplies a
summability that does not depend on $K$; the constants are the same at every level because nothing is
iterated.

It remains to bound the last event of the decomposition \eqref{8.8:eq-induction-base-a}. In the first
case of Lemma~\ref{8.8:lem-induction-base}, that lemma gives
\begin{equation*}
  \tilde{\nu}\bigl(\mathcal{E}^{(j)}_{\mathfrak{c}}(K^{\dagger\prime})\bigr)\le e^{-\theta p_0(g_j)}
\end{equation*}
for every cell. In the second case it gives this bound for the average over the level-$j$ cells, and
the covariance statement at level $j$, or the cell-averaged statement, of hypothesis (f) carries it
over to
$L^{(j)}_{\mathfrak{c}}$: the contribution of the last event to $\tilde{\nu}(L^{(j)}_{\mathfrak{c}})$ is
again at most $e^{-\theta p_0(g_j)}$.

\emph{Step 5: constants and the three forms.} The numbers $K_H$ and $C_{\delta}$ depend on $d$
only; $S_*$ depends on $\lambda^{\sharp}$ only; $\theta$ is fixed at Stage~0; and the bound of
Lemma~\ref{8.8:lem-induction-base} does not depend on the volume. Adding the bound of Step~4 on the
last event of \eqref{8.8:eq-induction-base-a} to
\eqref{8.8:eq-live-cell-rarity-8} and using $S_*\le1+1/\lambda^{\sharp}$ gives
\begin{equation*}
  \tilde{\nu}\bigl(L^{(j)}_{\mathfrak{c}}\bigr)
  \le\Bigl[K_H\,e^{\theta}\bigl(1+1/\lambda^{\sharp}\bigr)+1\Bigr]e^{-\theta p_0(g_j)}
  =C_L\,e^{-\theta p_0(g_j)},
\end{equation*}
which is \eqref{8.8:eq-live-cell-rarity-1}, that is, the live-cell bound
\eqref{8.8:eq-graded-rarity-a} at level $j$, with constant $C_L$ and rate $\theta$. The creation
form \eqref{8.8:eq-live-cell-rarity-a} and the
field form \eqref{8.8:eq-live-cell-rarity-2} follow from \eqref{8.8:eq-live-cell-rarity-1} by Step~1:
each of the two events is contained, up to a $\tilde{\nu}$-null set, in a union of $n_{\square}$ events
$L^{(j)}_{\mathfrak{c}}$. The graded form \eqref{8.8:eq-live-cell-rarity-3} follows from
\eqref{8.8:eq-live-cell-rarity-a} by Lemma~\ref{8.8:lem-grading-bounded}, applied with constant
$n_{\square}C_L$. Finally, for $1\le\lambda\le\bar\lambda$ the event $F_{j,\square}(\lambda)$ is contained
in $F_{j,\square}(1)$, and, since $p_0(g_j)\ge0$ and $\lambda^{2}/\bar\lambda^{2}\le1$,
\begin{equation*}
  e^{-\theta p_0(g_j)}\le e^{-(\theta/\bar\lambda^{2})\lambda^{2}p_0(g_j)} ;
\end{equation*}
so \eqref{8.8:eq-live-cell-rarity-4} follows from \eqref{8.8:eq-live-cell-rarity-2}.
\end{proof}

\begin{remark}[Large-field kinds enter only through live cells]\label{8.8:rem-creation-kinds}
Ba{\l}aban's large-field creations come in the kinds $\mathrm{S}0,\dots,\mathrm{S}10$, where
$\mathrm{S}0$ and $\mathrm{S}1$ have one and the same indicator. In the proof of
Theorem~\ref{8.8:thm-live-cell-rarity} nothing specific to a kind is estimated. A kind
enters in only two ways:
\begin{itemize}
\item through the cube that becomes live when its indicator is switched on;
\item through Ba{\l}aban's own factor for that indicator inside \cite[(1.79)]{B-CMP122a}.
\end{itemize}
Every one of these factors is graded in the multiple $\lambda$ of the threshold by the
extraction used in the proof of Theorem~\ref{8.8:thm-live-cell-rarity}.
\end{remark}

\begin{proof}
For each kind we list four items, in this order:
\begin{itemize}
\item Ba{\l}aban's indicator, with its reference;
\item the cubes or fields on which the indicator is tested;
\item whether the creation is seeded at the step itself or in the preparation of the
large-field data in \cite{B-CMP122a};
\item Ba{\l}aban's factor, graded in $\lambda$.
\end{itemize}
In each case the indicator decides only which cube becomes live. The factor is the one
stated at the cited equation.

\emph{Kinds $\mathrm{S}0$ and $\mathrm{S}1$.} The indicator is $1-\chi_k(\square)$, that is,
the event $a_k(\square)\ge 1$; it is the indicator of \cite[(1.1), (2.17), (3.2)]{B-CMP119},
and the test is repeated in \cite[p.~378]{B-CMP122b}. At the step it is tested on the new
field at level $k$, held fixed, off $\bar Z_{k-1}^{\sim4}$; at the repeated test of
\cite[p.~378]{B-CMP122b} it is tested on the healed region $X^{\sim6}$.
These kinds are seeded at the step itself. Their factor comes from the Wilson action and is,
per cube,
\begin{equation*}
\exp\bigl(-\gamma_0\,(2g_k^2)^{-1}\,(B_3^{-1}\lambda\varepsilon_k)^2\bigr).
\end{equation*}
In the proof of Theorem~\ref{8.8:thm-live-cell-rarity} such a creation either opens a live
component or joins an existing one, and its factor sits in the first product of
\cite[(1.79)]{B-CMP122a}. If it arises at the repeated test of \cite[p.~378]{B-CMP122b}, it is
treated, in the proof of Theorem~\ref{8.8:thm-live-cell-rarity}, with the creations that arise
at a repeated test rather than at the step.

\emph{Kind $\mathrm{S}2$.} The indicator is $\varphi_{k-1}(\square')\ge1$, from
\cite[(3.3)]{B-CMP119}. It is tested on the region $(B^k(P^1_k))^{\sim-1}$ on which
$\varphi_{k-1}$ is evaluated, in the notation of \cite[(3.3)]{B-CMP119}, and it is seeded at
the step itself. Its factor comes from the quadratic form and is
\begin{equation*}
\exp\bigl(-\tfrac12\gamma_0\,g^{-2}(2\lambda\delta)^2\bigr),
\end{equation*}
with $g$ and $\delta$ the coupling and the threshold of the step as in \cite[(3.3)]{B-CMP119}.
It enters the proof in the same way as $\mathrm{S}0$ and $\mathrm{S}1$.

\emph{Kind $\mathrm{S}3$.} The indicator is $\psi_{k-1}(\square')\ge1$, from
\cite[(3.16)]{B-CMP119}. It is tested on $\Omega_k^{\sim-1}$, where $\Omega_k$ is Ba{\l}aban's
region of \cite[(3.16)]{B-CMP119} on which $\psi_{k-1}$ is tested (not the set of histories of
Notation~\ref{8.8:not-history-measure}), and it is seeded at the step itself.
Its factor is $\exp(-\gamma_0 g^{-2}(\lambda\delta)^2)$, with $g$ and $\delta$ the coupling and
the threshold of the step as in \cite[(3.16)]{B-CMP119}. It enters the proof in the same
way as $\mathrm{S}0$ and $\mathrm{S}1$.

The remaining kinds all arise in the preparation of the large-field data. In the proof of
Theorem~\ref{8.8:thm-live-cell-rarity} they enter at the step that treats the prepared data.

\emph{Kind $\mathrm{S}4$.} The indicator is $1-\chi_k^{(n)}$, from
\cite[(1.22)--(1.24)]{B-CMP122a}. It is tested on the backgrounds on the old layers of one
component of the class of regions considered in \cite[(1.22)--(1.24)]{B-CMP122a}. Its factor
is at most $\exp(-A_1^2(\log g_j^{-2})^{2p_0})$.

\emph{Kinds $\mathrm{S}5$, $\mathrm{S}6$ and $\mathrm{S}7$.} Their indicators and tested
objects are as follows.
\begin{itemize}
\item $\mathrm{S}5$: the indicator is $1-\chi'_j$, from \cite[(1.27)]{B-CMP122a}. It is tested
on the boundary-layer field $V_j$ compared with the background.
\item $\mathrm{S}6$: the indicator is the one introduced in \cite[(1.28)]{B-CMP122a}. It is
tested on the deferred variable $A_j$ on $Z_{j+1}\cap\Omega_{j+1}$, where $Z_{j+1}$ and
$\Omega_{j+1}$ are Ba{\l}aban's regions in the notation of \cite[(1.28)]{B-CMP122a}.
\item $\mathrm{S}7$: the indicator is $1-\chi'^{(j)}$, from \cite[(1.55)]{B-CMP122a}. It is
tested on $B_j$, the argument of $\chi'^{(j)}$ in \cite[(1.55)]{B-CMP122a}.
\end{itemize}
All three kinds have the factor $\exp(-\gamma_0A_1^2p_1^2(g_j))$ as an upper bound, where
$p_1(\cdot)$ denotes the function of the coupling so lettered in \cite[(1.27)]{B-CMP122a}, the
letter also standing for its exponent, as with $p_0$.

\emph{Kind $\mathrm{S}8$.} The indicator is $1-\chi_{k,\Lambda}$, from
\cite[(1.75)]{B-CMP122a}. It is tested on $V_k$ restricted to $Z\cap\Lambda^c$, where $Z$ is
the large-field region and $\Lambda$ the region of \cite[(1.75)]{B-CMP122a} (this $Z$ is not
the normaliser of Notation~\ref{8.8:not-history-measure}). Its factor is
\begin{equation*}
\exp\bigl(-\tfrac14 g_k^{-2}\,(O(1)B_3M^2)^{-1}\varepsilon_k^2\bigr).
\end{equation*}

\emph{Kind $\mathrm{S}9$.} The indicator is $1-\chi'$, from \cite[(1.82)]{B-CMP122a}. It comes
in two cases and is tested on the variable $B'$ on the set $\mathbb{B}_0$, in the notation of
\cite[(1.82)]{B-CMP122a}. The larger of the factors of the two cases is
\begin{equation}\label{8.8:eq-creation-kinds-1}
\exp\Bigl(-\tfrac12\gamma_0\,
\bigl[6(d+3)\bigl(100M(L+1)N^{\beta_0}R_j\bigr)^{d+2}\bigr]^{-1}A_1^2p_1^2\Bigr),
\end{equation}
with the constants as in \cite[(1.82)]{B-CMP122a}, under the condition
$2p_1-(d+5)r_0>p_0$ on the exponents of \cite[(1.82)]{B-CMP122a}; here $p_1$ is the exponent
of that function, and $r_0$ denotes the exponent so lettered in
\cite[(1.82)]{B-CMP122a}, not the depth offset $r_0$ of the witness construction fixed in
Notation~\ref{0.2:not-glossary}.

\emph{Kind $\mathrm{S}10$.} The indicator is the one introduced in \cite[(1.88)]{B-CMP122a},
with threshold $\tfrac12\varepsilon_h$. It is tested on the oldest layer $V_h$, $h$ denoting
the index of that layer in \cite[(1.88)]{B-CMP122a} (not a history). Its factor is the factor
of $\mathrm{S}4$.

\emph{The gauge and auxiliary variables.} These are not a kind. The Faddeev--Popov identity
\cite[(1.5)]{B-CMP119} is exact, so no large-field indicator is attached to them.

In every case listed, the kind contributes to the proof of
Theorem~\ref{8.8:thm-live-cell-rarity} only two things: which cube becomes live when its
indicator is switched on, and the factor displayed for it. Each factor is graded in $\lambda$
by the extraction used in that proof. No estimate depends on the kind beyond these two
items, which is the assertion.
\end{proof}

\begin{corollary}[Rarity of young large-field blocks]\label{8.8:cor-young-block-rarity}
Assume the antecedents of Theorem~\ref{8.8:thm-live-cell-rarity} (rarity of live cells at every
admissible depth), namely the statements listed as hypotheses of that theorem, under the
conditions listed there. Among these conditions is the coupling ceiling
\eqref{8.8:eq-coupling-ceilings-b} of Definition~\ref{8.8:def-coupling-ceilings}.

Let $\Omega$ be the finite set of histories, and let $\tau_{\omega}$ be the total integral of the
term of the history $\omega\in\Omega$. Let $Z=\sum_{\omega\in\Omega}\tau_{\omega}$, and let
$E^{\kappa}_{j,\square}\subset\Omega$ be the event that a large-field region of kind $\kappa$ is
created at level $j$ in the cube $\square$. These objects are as in
Lemma~\ref{8.8:lem-graded-rarity}, and the normalised annealed measure is
$\tilde{\nu}(E)=Z^{-1}\sum_{\omega\in E}\tau_{\omega}$. Let $C_L$ and $\theta$ be the constants of
Theorem~\ref{8.8:thm-live-cell-rarity}, let $K_H$ be the number of boxes per cell, and let
$n_{\square}$ be the number of cells per cube. Put
\begin{equation}\label{8.8:eq-young-block-rarity-1}
\begin{gathered}
C_q:=n_{\square}C_L=n_{\square}\bigl(K_He^{\theta}(1+1/\lambda^{\sharp})+1\bigr),\qquad
c_q:=\theta,\\
q_j:=C_q\,e^{-c_q\,p_0(g_j)} .
\end{gathered}
\end{equation}
Then $C_q,c_q>0$ depend on $(L,M,\mathfrak{p},N)$ only. Moreover, the bound
\begin{equation}\label{8.8:eq-young-block-rarity-2}
\sum_{\omega\in E^{\kappa}_{j,\square}}\tau_{\omega}\;\le\;q_j\,Z
\end{equation}
holds for all of the following data:
\begin{itemize}
\item every $g\le\gamma_G$ and every $K$;
\item every terminal torus count $N_T=n_0L^a$ with $n_0\le\bar n$;
\item every admissible run and every level $j$ of it;
\item every kind $\kappa$ and every cube $\square$.
\end{itemize}
Equivalently, the live-cell bound \eqref{8.8:eq-graded-rarity-a} holds at level $j$ with the
constants $C_{\mathrm{lv}}=C_L$ and $c_{\mathrm{lv}}=\theta$.
\end{corollary}

\begin{proof}
Fix $g\le\gamma_G$, $K$, a terminal torus count $N_T=n_0L^a$ with $n_0\le\bar n$, an admissible
run, a level $j$ of it, a kind $\kappa$ and a cube $\square$. All of these are data to which
Theorem~\ref{8.8:thm-live-cell-rarity} applies: its antecedents are our hypotheses, and
$g\le\gamma_G$ places $g$ below the ceiling \eqref{8.8:eq-coupling-ceilings-b} under which that
theorem is stated.

The bound \eqref{8.8:eq-live-cell-rarity-a} of that theorem bounds
$\tilde{\nu}(E^{\kappa}_{j,\square})$ by $q_j$, with the constants
\eqref{8.8:eq-young-block-rarity-1}. The constants $n_{\square}$, $C_L$, $K_H$ and $\theta$
entering \eqref{8.8:eq-young-block-rarity-1} are those of
Theorem~\ref{8.8:thm-live-cell-rarity}, and $\lambda^{\sharp}$ is the flow slope of
Definition~\ref{1.2:def-flow-constants};
each of these reads $(L,M,\mathfrak{p},N)$ only, so $C_q$ and $c_q$ do, and both are positive.

By the definition of $\tilde{\nu}$,
\[
Z\,\tilde{\nu}(E^{\kappa}_{j,\square})=\sum_{\omega\in E^{\kappa}_{j,\square}}\tau_{\omega} .
\]
Since $Z=\int\rho_0>0$, where $\rho_0$ is the effective density at level $0$ of
Definition~\ref{1.5:def-densities}, we may multiply the bound
$\tilde{\nu}(E^{\kappa}_{j,\square})\le q_j$ by $Z$ without changing the direction of the
inequality. This gives \eqref{8.8:eq-young-block-rarity-2}.

The equivalent form stated is the form \eqref{8.8:eq-graded-rarity-a} at level $j$ with the
constants $C_{\mathrm{lv}}=C_L$ and $c_{\mathrm{lv}}=\theta$.
\end{proof}

The bound \eqref{8.8:eq-young-block-rarity-2} is used in two places. In the proof of the two-point
witness theorem it is used at the levels $j$ with $k_0-\mathfrak{l}\le j\le k_0$, where $k_0$ is
the level fixed in that proof and $\mathfrak{l}$ is its depth range, on the block families
$\mathfrak{B}_1(j)\cup\mathfrak{B}_2(j)$ introduced there. In the Gaussian extraction it is used
at the levels $j\le k_0+r(s)$, where $r(s)$ is the extraction depth. All of these are levels of an
admissible run, so the corollary applies at each of them. Consequently the young large-field
event $\mathrm{LF}^{\mathrm{young}}_i$ of the $i$-th block family of those arguments satisfies
\[
\tilde{\nu}(\mathrm{LF}^{\mathrm{young}}_i)
\;\le\; n^{\mathrm{blk}}_i\,\max_j q_j
\;\le\; \mathfrak{q}^{\mathrm{add}}_W .
\]
Here $n^{\mathrm{blk}}_i$ is the number of blocks of the $i$-th family and
$\mathfrak{q}^{\mathrm{add}}_W$ is the rarity budget of those arguments. The second inequality
holds under the smallness condition on the coupling imposed there, and the bound holds at every
admissible depth $s$.

The bound \eqref{8.8:eq-young-block-rarity-2} carries no factor decaying in the age. For a region
created at level $j\le k_0$, hence of age $k_0-j$ at level $k_0$, the fate bound used in the proof
of Theorem~\ref{8.8:thm-live-cell-rarity} yields the factor $e^{-\theta p_0(g_j)}$ and no factor
$L^{-4(k_0-j)}$; the only decay in the age is the decay of $e^{-\theta p_0(g_j)}$ in $j$.

\begin{definition}[The rim copy subtracted at a healing and its apparent excess]
\label{8.8:def-rim-copy}
Let $Z$ be a large-field region which is healed by Ba{\l}aban's large-field operation
$\mathbf{R}_k$ at step $k$, and let $j$ be a level whose terms are restored at this healing.
The terms of level $j$ are restored on the patch $P_j$. The values that were previously
missing are added at the block field $U_k$ and subtracted, as in \cite{B-CMP122b}. Among
the subtracted values are the constituents of the rim lying along the healed segment of the
region's own old rim $\partial\Lambda_j^{\circledR}$ at level $j$, a segment which lies in
the interior after the healing; here $\Lambda_j^{\circledR}$ is the region $\Lambda_j$ of
\cite{B-CMP119} with one layer of width $MR_j$ removed, and $\partial\Lambda_j^{\circledR}$
is its rim. These constituents enter with undifferentiated weight. The
\emph{rim copy subtracted at the healing} is the family of these subtracted rim constituents.

Let $n>j$ be a level and put $t=L^{j-n}$. Suppose that the rim copy is valued in a zone of
level $n$, cell by cell. Each cell of level $n$ that meets the healed segment then carries
$t^{-3}$ boundary bonds of the healed rim, and each of these bonds has size of order
$\mathsf{u}_j\rho_{n,j}^2t^4$. The resulting value of the rim copy per cell of level $n$ is
\begin{equation}\label{8.8:eq-rim-copy-1}
C\,\mathsf{u}_j\,\sup|\mathsf{h}|\,\rho_{n,j}^2L^{-(n-j)} .
\end{equation}
Here $C$ is a constant, $\mathsf{h}$ is the test function of the vector package,
$\mathsf{u}_j$ is its unit size at level $j$, and $\rho_{n,j}$ is the scale ratio between
levels $j$ and $n$.

Let $\hat{\alpha}$ be the load exponent of the vector package, for which the load entry
$\varsigma^V_{1+n-j}$ is of order $L^{-2\hat{\alpha}(n-j)}$. The ratio of
\eqref{8.8:eq-rim-copy-1} to the load entry $\varsigma^V_{1+n-j}\mathsf{u}_j$ of the vector
package is of order $C\rho_{n,j}^2\sup|\mathsf{h}|\,L^{(2\hat{\alpha}-1)(n-j)}$. The
\emph{apparent excess} of the rim copy at level $n$ is the factor
\begin{equation}\label{8.8:eq-rim-copy-2}
L^{(2\hat{\alpha}-1)(n-j)}
\end{equation}
in this ratio, and the vector package shows that this factor is not summable over $n>j$.

In what follows, three ways of removing this excess are considered:
\begin{itemize}
\item[(a)] $n-j$ is bounded by a constant;
\item[(b)] the age $k-j$ is confined to a window;
\item[(c)] the vector package is given an additional load entry indexed by the age $k-j$.
\end{itemize}
\end{definition}

\begin{proof}[Derivation of the sizes \eqref{8.8:eq-rim-copy-1} and \eqref{8.8:eq-rim-copy-2}]
With the count of bonds per cell of level $n$ and the size per bond recorded above, the
product of the count and the size per bond is
\begin{equation*}
t^{-3}\cdot\mathsf{u}_j\rho_{n,j}^2t^4=\mathsf{u}_j\rho_{n,j}^2t .
\end{equation*}
Since $t=L^{j-n}$, this equals $\mathsf{u}_j\rho_{n,j}^2L^{-(n-j)}$. Multiplying by
$\sup|\mathsf{h}|$, from the test function, and by the constant $C$ gives
\eqref{8.8:eq-rim-copy-1}.

In the vector package the load entry $\varsigma^V_{1+n-j}$ is of order
$L^{-2\hat{\alpha}(n-j)}$. The factor $\mathsf{u}_j$ occurs both in
\eqref{8.8:eq-rim-copy-1} and in $\varsigma^V_{1+n-j}\mathsf{u}_j$, so it cancels in the
ratio. The ratio of \eqref{8.8:eq-rim-copy-1} to $\varsigma^V_{1+n-j}\mathsf{u}_j$ is
therefore of order
\begin{equation*}
C\sup|\mathsf{h}|\,\rho_{n,j}^2\,L^{-(n-j)}L^{2\hat{\alpha}(n-j)}
=C\sup|\mathsf{h}|\,\rho_{n,j}^2\,L^{(2\hat{\alpha}-1)(n-j)} ,
\end{equation*}
whose factor depending on $n-j$ through the power of $L$ is \eqref{8.8:eq-rim-copy-2}.
\end{proof}

\begin{lemma}[Location of the healed rim segment]\label{8.8:lem-healed-rim-location}
Let $j$ be a level. Let $Z_j:=\Lambda_j^c$, $\Omega_j$ and $\Omega_{j+1}$ denote the unprimed sets
of the current label, in the convention fixed in the statement of the correlation theorem. Let
$\Lambda_j^{\circledR}$ be the set $\Lambda_j$ with one layer of width $M\check{R}_j$ (one
$M\check{R}_j$-layer) removed, as in
\cite{B-CMP119}, and let $\partial\Lambda_j^{\circledR}$ be its own rim. Assume the radius condition
$\check{R}_j\le L\check{R}_{j+1}$ of the correlation theorem. Then, with distances measured in
level-$j$ units,
\begin{equation}\label{8.8:eq-healed-rim-location-1}
\operatorname{dist}\bigl(\Lambda_j^c,\Omega_{j+1}\bigr)\ \ge\ 8LM\check{R}_{j+1}\ \ge\ 8M\check{R}_j
\ >\ M\check{R}_j+1 ,
\end{equation}
and consequently
\begin{equation}\label{8.8:eq-healed-rim-location-2}
\partial\Lambda_j^{\circledR}\subset\Omega_j\setminus\Omega_{j+1},
\end{equation}
where $\Omega_j$ and $\Omega_{j+1}$ are the unprimed sets of the current label. In particular, at a
healing in the sense of Definition~\ref{8.8:def-rim-copy}, the healed segment of the old own rim
$\partial\Lambda_j^{\circledR}$ lies in $\Omega_j\setminus\Omega_{j+1}$, inside the component $C$
treated at that healing. No cap in the sense of Definition~\ref{2.3:def-cap} is
placed on $M$, and no lower bound is placed on $L$.
\end{lemma}

\begin{proof}
The sets $Z_j=\Lambda_j^c$, $\Omega_j$ and $\Omega_{j+1}$ are the unprimed sets of the current
label. This is the convention fixed in the statement of the correlation theorem, and it is used
here as stated there. As in the correlation theorem, the radius $\check{R}_j$ is the radius written
$R_j$ in \cite{B-CMP119}. Accordingly, the layer widths $MR_j$ and $LMR_{k+1}$ of
\cite{B-CMP119} are written $M\check{R}_j$ and $LM\check{R}_{k+1}$ below.

\emph{First inequality.} In the construction of \cite{B-CMP119}, the set $Z_k$ is surrounded first
by four layers of $LM\check{R}_{k+1}$-cubes, then by one more such layer, then by one more, and then
by two more; only after these layers have been added is $\Omega_{k+1}$ defined. Take $k=j$. Then
$Z_j=\Lambda_j^c$ is separated from $\Omega_{j+1}$ by $4+1+1+2=8$ layers of
$LM\check{R}_{j+1}$-cubes, each of width $LM\check{R}_{j+1}$ in level-$j$ units; hence
$\operatorname{dist}(\Lambda_j^c,\Omega_{j+1})\ge 8LM\check{R}_{j+1}$.
This is the first inequality of \eqref{8.8:eq-healed-rim-location-1}.

\emph{Second inequality.} Multiply the radius condition $\check{R}_j\le L\check{R}_{j+1}$ by $8M>0$.
This gives $8M\check{R}_j\le 8LM\check{R}_{j+1}$, which is the second inequality.

The statement of the correlation theorem contains the same comparison in the form
$8L\check{R}_{j+1}\ge 8\check{R}_j$, that is, with the common factor $M$ divided out. Neither
inequality uses an upper bound on $M$ or a lower bound on $L$.

\emph{Third inequality.} The strict inequality $8M\check{R}_j>M\check{R}_j+1$ is equivalent to
$7M\check{R}_j>1$. The width $M\check{R}_j$ of an $M\check{R}_j$-layer of level-$j$ cells is a
positive integer number of level-$j$ units, so $M\check{R}_j\ge 1$ and $7M\check{R}_j\ge 7>1$.
This completes the chain \eqref{8.8:eq-healed-rim-location-1}.

\emph{The inclusion.} By definition, $\Lambda_j^{\circledR}$ is $\Lambda_j$ with one layer of width
$M\check{R}_j$ (one $M\check{R}_j$-layer) removed. Its rim $\partial\Lambda_j^{\circledR}$ therefore
lies in $\Lambda_j$, at distance at most $M\check{R}_j+1$ from $\Lambda_j^c$ in level-$j$ units.

By \eqref{8.8:eq-healed-rim-location-1}, every point of $\Omega_{j+1}$ is at distance greater than
$M\check{R}_j+1$ from $\Lambda_j^c$. Hence $\partial\Lambda_j^{\circledR}\cap\Omega_{j+1}=\emptyset$.

In Ba{\l}aban's construction \cite{B-CMP119}, $\Lambda_j\subset\Omega_j$. Hence
$\partial\Lambda_j^{\circledR}\subset\Omega_j$, and \eqref{8.8:eq-healed-rim-location-2} follows,
with $\Omega_j$, $\Omega_{j+1}$ the unprimed sets of the current label.

The argument applies to every point of $\partial\Lambda_j^{\circledR}$. It therefore applies in
particular to the points in the treated component $C$, among them the healed segment of the old
own rim of Definition~\ref{8.8:def-rim-copy}.
\end{proof}

\begin{definition}[Counting, indexing and reading level of a subtracted copy]
\label{8.8:def-copy-readings}
Let a large-field region $Z$ be healed by Ba{\l}aban's operation $\mathbf{R}_k$, let $j$ be a
level at which terms are restored on the patch $P_j$, and consider one of the values
subtracted at this healing along the healed segment of the old own rim
$\partial\Lambda_j^{\circledR}$, that is, a rim copy in the sense of
Definition~\ref{8.8:def-rim-copy}; in what follows we call it a \emph{subtracted copy}.
The following conventions fix how a subtracted copy is counted, indexed and read.

\smallskip
\emph{(a) Zones of the determining set.} \cite[(2.2)]{B-CMP119} sets
\begin{equation}\label{8.8:eq-copy-readings-1}
  \Gamma_0=\Omega_1^{c},\qquad
  \Gamma_j=\Omega_j^{(j)}\setminus\Omega_{j+1}^{(j)},\quad\dots,\qquad
  \mathbf{B}=\bigcup_j\Gamma_j ,
\end{equation}
and calls $\mathbf{B}$ the determining set; the $\Gamma_j$ are its zones.
In this book the superscript $(j)$ in \eqref{8.8:eq-copy-readings-1} is read as follows:
$\Omega_j^{(j)}$ is the set $\Omega_j$ regarded as a union of level-$j$ cells, so that
$\Gamma_j$ is the set of level-$j$ cells of the $j$-th zone. We assume that this is the
meaning of the superscript fixed in \cite[Sect.~0]{B-CMP109}, to which
\cite[(2.2)]{B-CMP119} refers; every convention below that refers to cells of the
determining set is conditional on this reading.

\smallskip
\emph{(b) The cells in which a subtracted copy is counted.} \cite[(1.69)]{B-CMP122b} states
\begin{equation}\label{8.8:eq-copy-readings-2}
  |F(Y,(U,J))|\le O(1)\sum_j|\Gamma^0_j\cap Y|
  \qquad\text{if $Y$ is a component of $Z^{\sim}$,}
\end{equation}
where $\{\Gamma^0_j\}$ denotes the old determining set, obtained after the last
renormalisation transformation, and $Z^{\sim}$ is the enlarged large-field region of
\cite[(1.69)]{B-CMP122b}; here $F(Y,(U,J))$ denotes the expression of
\cite[(1.69)]{B-CMP122b} attached to the domain $Y$, a function of the bond variables $U$
and of the further variable $J$ of that paper. The text following
\cite[(1.69)]{B-CMP122b} states further, of the values restored at the healing
(cf.\ Definition~\ref{8.8:def-rim-copy}), that they may be
resummed in domains $Y$, that the resummed expressions satisfy the
bounds \cite[(1.68), (1.69)]{B-CMP122b}, and that they are subtracted from the
corresponding terms $V(Y)$ of \cite{B-CMP122b}, which are not the block fields $V_k$. In
the correlation theorem each term, with localisation domain $X$ and localisation level $m$,
is charged once, at an anchor
cube $\square(X)\in\pi_m$ of $X\cap\Omega_m\cap\Omega^{c}_{m+1}$, which is a cell of
zone $m$ of the old determining set $\{\Gamma^0_j\}$ counted in
\eqref{8.8:eq-copy-readings-2}; the constant charged per cell, at a cell of zone $n\le k$
counted in level-$n$ units as the sum in \eqref{8.8:eq-copy-readings-2} counts, is the
constant $\mathsf{K}_{\mathrm{cell}}(k,n)$ of the charging rule, formed from the last column
of the table of per-cell sizes of the correlation theorem summed over its entries.

We adopt the following convention: a subtracted copy enters the count on the right-hand
side of
\eqref{8.8:eq-copy-readings-2}. The cells of that count are the cells of the old
determining set $\{\Gamma^0_j\}$, that is, the zone cells of the pre-attempt (unprimed)
large-field data, and zone $j$ is counted in level-$j$ units. This choice of cells rests
on \cite[(1.69)]{B-CMP122b} and the sentence following it, together with
\cite[(2.2)]{B-CMP119} in the reading (a), and it is a convention of this book. The charging
rule of the correlation theorem,
the weighted count of old cells used with it, and the corollary recalled in (e) below use
the same cells.

\smallskip
\emph{(c) The index of a subtracted copy.} A \emph{use} is a triple $u=(k',n,v)$: a
unit $v$ read at level $k'$ in zone $n\le k'$, where the zone is determined by the
point or support of $v$ lying in $\Omega_n\setminus\Omega_{n+1}$; this is the definition of
a use in terms of which the theorem on the occurrences of uses, recalled at the end of this
paragraph, is stated. The definition of a use takes its sets $\Omega_n$ from the
correlation theorem; we assume that they are the unprimed (pre-attempt) sets $\Omega_n$
of (b) and not the current ones, and convention (c) is conditional on this identification.
A subtracted copy lies in the unprimed
set $\Omega_j\setminus\Omega_{j+1}$. It is therefore the use $(k,j,\cdot)$, and its
\emph{dial} is $\lambda_{j,k}$, the parameter attached to a use read at level $k$ in
zone $j$. Clause (d) of that theorem, in its statement on the occurrences of a use,
then supplies the factor
\begin{equation}\label{8.8:eq-copy-readings-3}
  q^{(k-j-1)_+},
\end{equation}
with $q$ the geometric forgetting ratio. This factor multiplies the term with index $n=j$
that is already present in the rim estimate of the vector package.

\smallskip
\emph{(d) The level at which the size of a subtracted copy is read.} The table of
per-cell sizes of the correlation theorem is indexed, at a point $y$, by the cube
$\tilde{\square}\in\pi_n$ of the current zone containing $y$. Its indices are
\begin{equation}\label{8.8:eq-copy-readings-4}
  \nu=1+n-j,\qquad t=L^{j-n}.
\end{equation}
The charging rule recalled in (b) assigns to a charged value the last column of this table
per unit $n$-cell. The index $n$ of \eqref{8.8:eq-copy-readings-4} is the level of the
current-zone cube at which a size is read; it is not the zone index $n$ of a use in (c),
nor the zone index $n$ of a cell of the count in (b).
For a subtracted copy each of the latter two equals $j$, and the former is fixed below to
be $k$.

At the site of $\mathbf{R}_k$ at which the subtracted values are placed, the placement
clause of the correlation theorem fixes the following vocabulary:
$\mathfrak{B}:=\mathbf{B}''_k$, the current cover-and-window structure at step $k$; levels
of points are read from $\mathfrak{B}$; cells,
levels and zones are those of $\mathfrak{B}$; and, inside this clause and only there,
$Z_j$, $\Omega_j$ denote the sets of the current large-field data, not the pre-attempt sets
of (b) and (c). The weight clause attached to this vocabulary names the
configuration of a subtracted copy: the weight
\begin{equation}\label{8.8:eq-copy-readings-5}
  \mathsf{L}_0^{2(k_a-m)}
\end{equation}
is carried on the old ring of the top zone of pre-attempt level
$m$ with $k_{0a}<m\le k_a$, where $k_{0a}$ and $k_a$ are the base and top levels of the
running attempt $\mathsf{a}$, and $\mathsf{L}_0$ is the window base of the correlation
theorem ($3\le\mathsf{L}_0<\tfrac12 L$), distinct from the constant $L_0$ of (H2).
For a subtracted copy the pre-attempt level is $m=j$ and the top
level is $k_a=k$, so that the weight carried is $\mathsf{L}_0^{2(k-j)}$. The far-core
clause of the correlation theorem reads the smooth far
cores of the subtracted copies at
\begin{equation*}
  (n,\mathsf{W})=(k,3).
\end{equation*}

We adopt the following convention: the size of a subtracted copy is read at $n=k$. In
\eqref{8.8:eq-copy-readings-4} this means
\begin{equation}\label{8.8:eq-copy-readings-6}
  \nu=1+k-j,\qquad t=L^{j-k}.
\end{equation}
Thus the level of $\mathbf{B}''_k$ is the reading scale of a subtracted copy.

\smallskip
\emph{(e) Scope of the absorption of window weights.} The corollary of the correlation
theorem on running windows has a last clause; here the running window is the window of the
running attempt $\mathsf{a}$, whose cell weights $w(\Delta)$ are not frozen, as opposed to
the frozen windows
of live arrests. That clause absorbs a window weight into the
per-cell count of old cells. The weight is the window weight $\bar{\mathsf{w}}^3$ of the
pointed units on window cells, which satisfies
\begin{equation}\label{8.8:eq-copy-readings-7}
  \bar{\mathsf{w}}^3\le 27\,\mathsf{L}_0^{6}\,L^{de},
\end{equation}
where $d=4$ and $e$ is the weight exponent of the cell, the number of levels by which the
level of the cell exceeds $j$ (a cell of level $k$ has $e=k-j$).
The clause states that a current cell $\Delta$ of the running window carrying the weight
$w(\Delta)=\mathsf{L}_0^{2e}$ contains at least $L^{de}$ cells of the old determining set
$\{\Gamma^0_j\}$ that \eqref{8.8:eq-copy-readings-2} and the charging rule of (b) count.
That clause is stated for the pieces of pointed units that are plain interior units,
under the hypothesis on plain interior units that accompanies it in the correlation
theorem. It contains no statement about the cell at which a value subtracted at a healing
is charged. Accordingly, convention (b) is not derived from it in this definition.

\smallskip
The indices in (c) and (d) are consistent.
The subtracted copy lies along the healed segment of $\partial\Lambda_j^{\circledR}$. By the
location of the healed rim segment (Lemma~\ref{8.8:lem-healed-rim-location}), this segment
lies in $\Omega_j\setminus\Omega_{j+1}$ of the unprimed large-field data, inside the
treated component, the component of the large-field data treated at this healing. In the
definition of a use recalled in (c), the zone of a unit
is the index $n$ with its point or support in $\Omega_n\setminus\Omega_{n+1}$. Hence the
zone of the copy is $n=j$, its level of reading is $k$, and it is the use $(k,j,\cdot)$
with dial $\lambda_{j,k}$, as stated in (c). Finally, substituting the reading index $n=k$
into \eqref{8.8:eq-copy-readings-4} gives $\nu=1+k-j$ and $t=L^{j-k}$, which is
\eqref{8.8:eq-copy-readings-6}.
\end{definition}

\begin{proposition}[The healed-rim bound without excess]\label{8.8:prop-healed-rim-bound}
Let $0\le j\le k$, let $Z$ be a large-field region healed by $\mathbf{R}_k$, and let the rim copy
along the healed segment of the old own rim $\partial\Lambda_j^{\circledR}$ be as in
Definition~\ref{8.8:def-rim-copy}. Let $\Delta$ be an old zone-$j$ cell of $\Gamma^0_j$, that is,
a unit cell of the level-$j$ lattice in $\Omega_j\setminus\Omega_{j+1}$ of the label before the
healing, met by the healed segment. Suppose that $\Delta$ lies on the top zone's old ring of
pre-attempt level $j$ of the running attempt $\mathsf{a}$ (in the sense of the correlation
theorem), whose top level is $k_a=k$ and whose base level satisfies $k_{0a}<j$. Put
\begin{equation*}
e:=k-j\ge 0,\qquad t:=L^{j-k}=L^{-e}\le 1 .
\end{equation*}
Assume the following.
\begin{itemize}
\item[(a)] The copy is counted, indexed and read as in Definition~\ref{8.8:def-copy-readings}:
its modulus enters the count per old zone-$j$ cell of $\Gamma^0_j$, counted in level-$j$
units; its dial index is $j$; and its reading level is the level $k$ of $\mathbf{B}''_k$, so that
it is read at scale $n=k$ with $\nu=1+e$ and $t=L^{j-k}$.
\item[(b)] The per-old-cell charging of window weights, which the correlation theorem provides
for the weighted mass of pieces of pointed units (in the sense of that theorem), applies to a
value subtracted at a healing of $Z$ by $\mathbf{R}_k$: if such a value is read in a cell of the
running window of $\mathbf{B}''_k$, its modulus is charged per old zone-$j$ cell of $\Gamma^0_j$
with the window weight absorbed into that count; if it is read in a cell of the frozen window of
a live arrest $\mathfrak{a}$ (in the sense of the correlation theorem), its weighted modulus is
charged extensively, per pre-arrest zone-$j$ cell, inside $\mathsf{P}(\mathfrak{a})$. (Only the
first alternative is used in the proof below, $\Delta$ being placed as above; the second concerns
healed rims lying inside a live arrest.)
\item[(c)] The coupling $g$ appearing in the scale-ratio bound of the correlation theorem
satisfies $g\le 1$.
\end{itemize}
Put
\begin{equation}\label{8.8:eq-healed-rim-bound-b}
C_w:=27\mathsf{L}_0^6\,(1+o(1))\,(1+2c_2)(1+3\lambda),
\end{equation}
where $1+o(1)$ is the factor in the same scale-ratio bound of the correlation theorem.
Then:
\begin{itemize}
\item[(1)] the number of rim constituents of level $j$ pinned in $\Delta$, those of the own rim
together with those of the copy, is at most $2C_{\mathfrak{J}}'dM^3$;
\item[(2)] for a rim constituent supported on a set $X\subset\tilde{\square}^2$, where
$\tilde{\square}$ is the scale-$k$ cube of its point $y$, the size of its copy at level $k$ and
window weight $\bar{\mathsf{w}}$, relative to the own-rim entry at level $j$ and weight $1$, is at
most $C_w$;
\item[(3)] for a rim constituent supported on a set $X\not\subset\tilde{\square}^2$ and $e\ge1$,
its decay supplies the extra factor $\exp(-\frac{15}{16}\kappa e)$, and
\begin{equation}\label{8.8:eq-healed-rim-bound-c}
\begin{aligned}
&\exp(-\tfrac{15}{16}\kappa e)<t^5=L^{-5e},\\
&t^5\cdot 27\mathsf{L}_0^6(L/3)^{3e}\bigl(1+(3\lambda e+2c_2)g^2\bigr)\\
&\qquad=27\mathsf{L}_0^6\,(27L^2)^{-e}\bigl(1+(3\lambda e+2c_2)g^2\bigr)
\le 27\mathsf{L}_0^6(1+2c_2)(1+3\lambda);
\end{aligned}
\end{equation}
\item[(4)] the constant $C_w$ depends only on $\mathsf{L}_0$, $\lambda$ and $c_2$, and is uniform in
$e=k-j$, in $k$ and in $K$. In particular, charged per old zone-$j$ cell, neither the count (1)
nor the bounds (2), (3) carry a factor $t^{-3}$, and no factor $L^{(2\hat{\alpha}-1)(k-j)}$
arises, $\hat{\alpha}$ being the load exponent of the vector package.
\end{itemize}
\end{proposition}

\begin{proof}
The proof has three steps: the count per old cell, the size per constituent at reading level
$k$, and an elementary inequality that bounds the resulting factor uniformly in $e$.

\emph{Step 1: the count per old zone-$j$ cell.} By Lemma~\ref{8.8:lem-healed-rim-location} the
healed segment of $\partial\Lambda_j^{\circledR}$ lies in $\Omega_j\setminus\Omega_{j+1}$ of the
label before the healing, inside the treated component. Hence the cells $\Delta$ of the
statement are old zone-$j$ cells of $\Gamma^0_j$, and by hypothesis~(a) these are the cells in
which the count is taken. The rim constituents of level $j$ pinned in $\Delta$ are those of the
own-rim entry of the table of load entries of the vector package: at most
$C_{\mathfrak{J}}'dM^3$ bonds of $\partial\Lambda_j^{\circledR}$ meet $\Delta$, the own-rim entry
being at most $C_{\mathfrak{J}}'M^3\mathsf{u}_j$ in zone $n=j$. By
Definition~\ref{8.8:def-rim-copy}, the copy subtracted at the healing consists of the rim
constituents along the healed segment, so it duplicates these very constituents. Therefore
the count per old zone-$j$ cell at most doubles, which is the factor $2$ in~(1).

No factor $t^{-3}$ appears, because the count is per old zone-$j$ cell and not per cell of level
$k$ of $\mathbf{B}''_k$. A cell of level $k$ would contain $L^{3(k-j)}=t^{-3}$ of the old
zone-$j$ cells met by the rim, and counting per cell of level $k$ is what produces the apparent
excess described in Definition~\ref{8.8:def-rim-copy}.

\emph{Step 2: the size per constituent at reading level $k$.}

By hypothesis~(a), a rim constituent of the copy is read at scale $n=k$. The correlation
theorem's lemma on admissible fields is stated at a general scale $n\ge j$, with
$t=L^{j-n}$ and $\tilde{\square}$ the scale-$n$ cube of $y$. It reads as follows. Let $F(X,\cdot)$
have, on the domain of fields specified there, the two properties required of functionals in
that lemma, which we denote by \textup{(f1)} and \textup{(f2)}, and let
$|F|\le\mathfrak{n}$. Let $\mathfrak{b}$ be an admissible field at scale $n\ge j$. Put
\begin{equation*}
D_X:=\operatorname{diam}_jX+5M\le 4M\,(d_j(X)+5),
\end{equation*}
where $\operatorname{diam}_j$ is the diameter in level-$j$ units. Then part~(i) of that lemma
states: if $X\subset\tilde{\square}^2$, then
\begin{equation}\label{8.8:eq-healed-rim-bound-1}
|F(X,\mathfrak{b})-F(X,\mathbb{1})|\le 2\mathfrak{n}\,\bigl(C_MD_X\,t^2\rho_{n,j}\bigr)^2 .
\end{equation}

The first part of the vector package's lemma on the operations $\mathfrak{J}[F]$, $F^{\sharp}$,
$\mathsf{S}[F]$ states the following about an E-family $F$, that is, a family of functionals of
the kind on which these operations are defined in the vector package, with $|F|\le\mathfrak{n}$:
\begin{itemize}
\item[(i)] these operations are linear in $F$;
\item[(ii)] they are analytic in every parameter in which $F$ is analytic;
\item[(iii)] they have \textup{(f1)}, \textup{(f2)} on the radii of $F$ and vanish to second
order at $\mathbb{1}$;
\item[(iv)] $\mathfrak{J}[F]$ satisfies
\begin{equation*}
\sup\exp\bigl(\tfrac{15}{8}\kappa d_j(X)\bigr)\bigl|\mathfrak{J}[F](X,\cdot;b;\epsilon)\bigr|
\le C_{\mathfrak{J}}\mathfrak{n}\,|\epsilon| ,
\end{equation*}
where $b$ and $\epsilon$ are the further arguments of $\mathfrak{J}[F]$ in the vector package and
the supremum is taken over the field argument;
\item[(v)] $F^{\sharp}$ is a pointed E-family typed by $V$, and $\mathsf{S}[F]$ a pointed E-family
typed by $V\otimes V$, in the sense of the vector package; they are bounded by
$(1+8C_{\mathfrak{J}})\mathfrak{n}$ and $4C_{\mathfrak{J}}\mathfrak{n}$ respectively in the flat
norm $\|\cdot\|^{\flat}$, at the rate $\kappa_{\sharp}:=\frac{15}{8}\kappa$.
\end{itemize}
This supplies the rim constituent's properties \textup{(f1)}, \textup{(f2)}, and a bound
$\mathfrak{n}$ on its modulus of the type
\begin{equation*}
C_{\mathfrak{J}}\mathsf{u}_j\sup|\mathsf{h}|\,e^{-\frac{15}{8}\kappa d_j(X)} .
\end{equation*}
The own-rim entry of the vector package applies this pair of statements at $t=1$. It does so
only because the own rim lies in zone $j$; the statements themselves carry a general $t$.

\emph{The window weight.} In both cases of the dichotomy below the copy is read at $n=k$, at the
window weight of its cell. At the localisation site of $\mathbf{R}_k$ in the correlation theorem,
where the window structure is $\mathbf{B}''_k$, the window-weight rule of that theorem assigns
the weight $\mathsf{L}_0^{2(k_a-m)}$ on the top zone's old ring of pre-attempt level $m$ with
$k_{0a}<m\le k_a$. Here $m=j$ and $k_a=k$, so the running-window cell of $\mathbf{B}''_k$ that
contains $\Delta$ carries a weight $w(\Delta)\le\mathsf{L}_0^{2e}$.

Ba{\l}aban's auxiliary window base $\mathsf{L}_0$ satisfies $3\le\mathsf{L}_0<\frac12L$ and
$\mathsf{L}_0^2\le L/3$; these are among the parameter conditions of the correlation theorem. In
particular $\mathsf{L}_0^{2e}\le(L/3)^e$.

The entries of pointed units carry the weight
$\bar{\mathsf{w}}_v:=3w^{\vee}(\Delta_y)\le 3\mathsf{L}_0^{2(e+1)}$. Here $w^{\vee}(\Delta_y)$ is
the window weight attached to the cell $\Delta_y$ of the point $y$, and passing from $w$ to
$w^{\vee}$ raises the exponent $e$ by at most one. This weight enters the per-cell charge as
$\bar{\mathsf{w}}^3$. Using $\mathsf{L}_0^{6e}=(\mathsf{L}_0^2)^{3e}\le(L/3)^{3e}$, we obtain
\begin{equation}\label{8.8:eq-healed-rim-bound-2}
\bar{\mathsf{w}}^3\le(3w^{\vee})^3\le 27\mathsf{L}_0^{6(e+1)}\le 27\mathsf{L}_0^6\,(L/3)^{3e}.
\end{equation}

\emph{The scale ratio.} The correlation theorem bounds the scale ratio by
\begin{equation}\label{8.8:eq-healed-rim-bound-3}
\rho_{n,j}^2\le(1+o(1))\bigl(1+(3\lambda(n-j)+2c_2)g^2\bigr).
\end{equation}

We split the constituents according to the dichotomy of the lemma on admissible fields at scale
$n=k$, where $\tilde{\square}$ is the scale-$k$ cube of $y$.

\emph{Case (ii-a): $X\subset\tilde{\square}^2$.} Part~(i) of the lemma on admissible fields,
\eqref{8.8:eq-healed-rim-bound-1} with $n=k$, applies with this $\mathfrak{n}$.
It bounds the constituent's marginal modulus by
\begin{equation*}
2C_{\mathfrak{J}}\mathsf{u}_j\sup|\mathsf{h}|\,\bigl(C_MD_X\,t^2\rho_{k,j}\bigr)^2
e^{-\frac{15}{8}\kappa d_j(X)} .
\end{equation*}

\emph{Case (ii-b): $X\not\subset\tilde{\square}^2$.} This is the tail class of the correlation
theorem. It is handled by the device used there for a set $X\not\subset\square^{\sim\nu}$. For
such a set $d_j(X)\ge\nu-1$, and the device bounds $2\mathfrak{n}\exp(-\kappa d_j(X))$ by
$2\mathfrak{n}\exp(-\frac12\kappa d_j(X))\,t^5$. This is the inequality
$\exp(-\frac12\kappa(n-j))<(L^jL^{-n})^5$ used in \cite{B-CMP119}. In the correlation theorem it
appears as the floor $\kappa\ge 10\log L$, one of the lower bounds on $\kappa$ among that
theorem's parameter conditions.

At scale $n=k$ we have $\nu=1+e$. Part~(ii) of the lemma on admissible fields gives, for such
$X$, $d_j(X)\ge t^{-1}$. Since $t^{-1}=L^e\ge e$, this yields
\begin{equation*}
d_j(X)\ge t^{-1}\ge \nu-1=e .
\end{equation*}
Write the decay factor supplied above as
\begin{equation*}
\exp\bigl(-\tfrac{15}{8}\kappa d_j(X)\bigr)
=\exp\bigl(-\tfrac{15}{16}\kappa d_j(X)\bigr)\,\exp\bigl(-\tfrac{15}{16}\kappa d_j(X)\bigr) .
\end{equation*}
The first factor is kept for the sum over $X$, as the entries of the table of load entries of the
vector package do. The second factor is at most $\exp(-\frac{15}{16}\kappa e)$, by $d_j(X)\ge e$.

For $e\ge1$ we have $\exp(-\frac{15}{16}\kappa e)<\exp(-\frac12\kappa e)$, since $\kappa>0$ and
$\frac{15}{16}>\frac12$. Further, $\exp(-\frac12\kappa e)\le L^{-5e}$, by $\kappa\ge10\log L$.
This gives the first line of \eqref{8.8:eq-healed-rim-bound-c}.

The factor $t^5$ does the work of the factor $t^4$ of case~(ii-a), with one power to spare.
Multiplying it by the weight bound \eqref{8.8:eq-healed-rim-bound-2} and by the second factor of
the scale-ratio bound \eqref{8.8:eq-healed-rim-bound-3}, one has
$t^5(L/3)^{3e}=L^{-5e}L^{3e}27^{-e}=(27L^2)^{-e}$; this is the equality in the second line of
\eqref{8.8:eq-healed-rim-bound-c}. The inequality in that line follows from the elementary
inequality of Step~3, applied with $a=27L^2$, and is proved there.
No new floor is used: $\kappa\ge 10\log L$ is one of the correlation theorem's own parameter
conditions. For $e=0$ nothing is needed: then $t=1$, and the constituent is the own-rim entry
itself.

\emph{The ratio in case (ii-a).} Compare the bound of case~(ii-a) at level $k$ and weight
$\bar{\mathsf{w}}$ with the own-rim entry, which is the same bound at level $j$, $t=1$ and
weight $1$. The factors $2C_{\mathfrak{J}}\mathsf{u}_j\sup|\mathsf{h}|$, $(C_MD_X)^2$ and
$e^{-\frac{15}{8}\kappa d_j(X)}$ are common to both. The own-rim entry of the vector package
carries the scale-ratio factor $1$; accordingly, in the denominator below, $\rho_{j,j}^2$ is
replaced by $1$. Combining \eqref{8.8:eq-healed-rim-bound-2},
$t^4=L^{-4e}$ and \eqref{8.8:eq-healed-rim-bound-3} at $n=k$, and using
$(L/3)^{3e}L^{-4e}=(27L)^{-e}$, the copy's size relative to the own-rim entry is at most
\begin{equation}\label{8.8:eq-healed-rim-bound-a}
\begin{split}
\bar{\mathsf{w}}^3\,t^4\,\frac{\rho_{k,j}^2}{\rho_{j,j}^2}
&\le 27\mathsf{L}_0^6\,(L/3)^{3e}L^{-4e}\,
(1+o(1))\bigl(1+(3\lambda e+2c_2)g^2\bigr)\\
&=27\mathsf{L}_0^6\,(1+o(1))\,(27L)^{-e}\bigl(1+(3\lambda e+2c_2)g^2\bigr).
\end{split}
\end{equation}

\emph{Step 3: the factor in \eqref{8.8:eq-healed-rim-bound-a} and
\eqref{8.8:eq-healed-rim-bound-c} is bounded uniformly in $e$.} We use the following elementary
inequality. For real $x\ge0$, $b\ge0$ and $a>1$ with $\ln a\ge1$,
\begin{equation}\label{8.8:eq-healed-rim-bound-4}
a^{-x}(1+bx)\le 1+b .
\end{equation}
If $x\le1$, then $bx\le b$, and since $a^{-x}\le1$ we get $a^{-x}(1+bx)\le1+bx\le1+b$. If $x\ge1$,
then $1+bx\le x+bx=(1+b)x$, so
\begin{equation*}
a^{-x}(1+bx)\le(1+b)\,x\,a^{-x}\le(1+b)\sup_{x\ge1}x\,a^{-x}\le\frac{1+b}{\exp(1)\,\ln a}
\le 1+b .
\end{equation*}
Here we used that $y\mapsto y\,a^{-y}$ on $y>0$ attains its maximum $1/(\exp(1)\ln a)$ at
$y=1/\ln a$, and that $\exp(1)\ln a\ge1$ once $\ln a\ge1$. The condition $\ln a\ge1$ holds
whenever $a\ge3$.

We apply \eqref{8.8:eq-healed-rim-bound-4} with $x=e$, with $a=27L$ in
\eqref{8.8:eq-healed-rim-bound-a} and $a=27L^2$ in \eqref{8.8:eq-healed-rim-bound-c}. Both values
are at least $3$. We take
\begin{equation*}
b=\frac{3\lambda g^2}{1+2c_2g^2}\le 3\lambda ,
\end{equation*}
where the bound uses $1+2c_2g^2\ge1$ and hypothesis~(c). With this choice,
$1+(3\lambda e+2c_2)g^2=(1+2c_2g^2)(1+be)$. Hence, for both values of $a$,
\begin{equation*}
\begin{split}
a^{-e}\bigl(1+(3\lambda e+2c_2)g^2\bigr)
&=(1+2c_2g^2)\,a^{-e}(1+be)\\
&\le(1+2c_2g^2)(1+3\lambda)\le(1+2c_2)(1+3\lambda),
\end{split}
\end{equation*}
the last step by hypothesis~(c). The linear growth in $e$ is thus absorbed by $a^{-e}$.

Inserted in \eqref{8.8:eq-healed-rim-bound-a}, this shows that the ratio of case~(ii-a) is at most
$C_w$ as defined in \eqref{8.8:eq-healed-rim-bound-b}, which is~(2). Inserted in the second line
of \eqref{8.8:eq-healed-rim-bound-c}, it gives the bound
$27\mathsf{L}_0^6(1+2c_2)(1+3\lambda)$ stated there; together with the first line, established in
Step~2, this is~(3).

\emph{Conclusion.} The constant $C_w$ depends only on $\mathsf{L}_0$, $\lambda$ and $c_2$.
$\mathsf{L}_0$ is the base of the window weights, a number fixed before $\gamma$ in the order of
choice of the constants; $\lambda$ and $c_2$ are the constants of the scale-ratio bound
\eqref{8.8:eq-healed-rim-bound-3}. None of these depends on $e=k-j$, on $k$ or on $K$, which gives
the uniformity in~(4).

By Step~1 the count per old zone-$j$ cell is at most $2C_{\mathfrak{J}}'dM^3$, with no factor
$t^{-3}$. By hypothesis~(b), in its running-window alternative ($\Delta$ lying inside a cell of
the running window of the running attempt), the modulus read at the window weight of that cell is
charged per old zone-$j$ cell of $\Gamma^0_j$ with that weight absorbed. By Steps~2 and~3 each
constituent's weighted size at reading level $k$ exceeds the own-rim entry by at most $C_w$,
respectively by at most $27\mathsf{L}_0^6(1+2c_2)(1+3\lambda)$ in the tail class.

Neither bound involves $L^{(2\hat{\alpha}-1)(k-j)}$. Such a factor would require the count per
cell of level $k$, with its factor $t^{-3}$, which hypothesis~(a) replaces by the count per old
zone-$j$ cell. This proves~(4).
\end{proof}

\begin{proof}[Proof of Proposition~\ref{8.8:prop-healed-rim-bound}, continued]
\label{8.8:prf-rim-per-cell}
We keep the notation of the preceding part of the proof: $\Delta$ is an old zone-$j$ cell of
$\Gamma^0_j$, that is, a unit cell of the level-$j$ lattice in $\Omega_j\setminus\Omega_{j+1}$, met
by the healed segment of $\partial\Lambda_j^{\circledR}$. The scale ratio is
$t=L^{j-n}\le 1$ at reading level $n$, and $e$ denotes the exponent of the window weight.

\emph{The contribution per old zone-$j$ cell.} We first take the reading level $n=k$ of the
reading-level clause of Definition~\ref{8.8:def-copy-readings}.
\begin{itemize}
\item Count. The rim constituents pinned in $\Delta$ are those of the rim entry of $T^V_j$, that
is, of the term of the vector package that bounds the constituents of the own rim
$\partial\Lambda_j^{\circledR}$ placed in zone $j$. The healed copy duplicates exactly these
constituents, so their number per old zone-$j$ cell is at most doubled.
\item Size. Each constituent of the copy is read at level $k$ with the window weight of its
cell. Its size is at most $C_w$ times its size in the rim entry, which is read at level $j$ with
weight $1$. For constituents with $X\subset\tilde{\square}^2$ this is the ratio bound
\eqref{8.8:eq-healed-rim-bound-a} together with the uniform bound
\eqref{8.8:eq-healed-rim-bound-b}. For constituents in the tail class
$X\not\subset\tilde{\square}^2$ it is the tail bound \eqref{8.8:eq-healed-rim-bound-c}, in which
the factor $t^5$ takes the place of $t^4$.
\end{itemize}
The rim entry bounds the constituents pinned in $\Delta$ by $C_{\mathfrak{J}}'M^3\mathsf{u}_j
\sup|\mathsf{h}|$. Multiplying by the factor $2$ of the count and the factor $C_w$ of the size,
the healed copy contributes, per old zone-$j$ cell, at most
\begin{equation}\label{8.8:eq-rim-per-cell-a}
2\,C_w\,C_{\mathfrak{J}}'M^3\,\mathsf{u}_j\sup|\mathsf{h}|
\;\le\; C_{\mathrm{rim}}\,\varsigma^V_1\,\mathsf{u}_j\sup|\mathsf{h}| ,
\end{equation}
for a constant $C_{\mathrm{rim}}$ depending only on
$(d,M,C_{\mathfrak{J}},\kappa,\mathsf{L}_0,\lambda,c_2)$; in particular $C_{\mathrm{rim}}$ does not
depend on $j$, $k$ or $K$.

The left member of \eqref{8.8:eq-rim-per-cell-a} is the bound of the rim entry of $T^V_j$
multiplied by the pure number $2C_w$. The contribution therefore lies inside the rim entry of
$T^V_j$, provided the rim constant in $C_L^V$ is multiplied by $2C_w$.

By \eqref{8.8:eq-healed-rim-bound-b}, $C_w$ depends only on $\mathsf{L}_0$, $\lambda$, $c_2$, and
it is uniform in $e=k-j$, in $k$ and in $K$. Hence the enlarged rim constant $C_{\mathrm{rim}}$ is
a number depending on $(d,M,C_{\mathfrak{J}},\kappa,\mathsf{L}_0,\lambda,c_2)$. It is fixed before
$\gamma$, and it enters only the existing one-sided ceilings on $\gamma$, namely those in which the
rim constant of $C_L^V$ already appears.

The age $k-j$ is not paid by \eqref{8.8:eq-rim-per-cell-a}. It is paid, at dial index $j$ as
the indexing clause of Definition~\ref{8.8:def-copy-readings} prescribes, by the factor
$q^{(k-j-1)_+}$ already present in the age bound of the vector package; this is carried out
below under the heading on the age.

No factor $L^{(2\hat{\alpha}-1)(k-j)}$ arises. Such a factor would be formed as follows:
\begin{itemize}
\item the number, of order $t^{-3}=L^{3(k-j)}$, of level-$j$ rim bonds of
$\partial\Lambda_j^{\circledR}$ meeting one level-$k$ cell,
\item multiplied by the size $t^4$ per bond, giving $t$ per level-$k$ cell,
\item compared with the load entry
$\varsigma^V_{1+k-j}\asymp L^{-2\hat{\alpha}(k-j)}$.
\end{itemize}
The quotient $t\,L^{2\hat{\alpha}(k-j)}$ equals $L^{(2\hat{\alpha}-1)(k-j)}$. This product does
not form here, for two reasons:
\begin{itemize}
\item its first factor is a count per level-$k$ cell, whereas the counting clause of
Definition~\ref{8.8:def-copy-readings} counts per old zone-$j$ cell;
\item its comparison entry $\varsigma^V_{1+k-j}$ belongs to the dial index $n=k$, whereas the
indexing clause of Definition~\ref{8.8:def-copy-readings} fixes the dial index $j$.
\end{itemize}

\emph{The placements.} The healed levels are all $j\ge j(Z)-1$; this is the statement of the
vector package on the levels restored at a healing. The window-weight clause of the correlation
theorem, read at the site of $\mathbf{R}_k$ on the cover $\mathbf{B}''_k$, distinguishes
three placements of a point of $(\Omega_j\setminus\Omega_{j+1})\cap C$ under
$\mathbf{B}''_k$. For each we list the reading level $n$ and the weight exponent $e$:
\begin{equation}\label{8.8:eq-rim-per-cell-b}
\begin{gathered}
\text{(P1)}\quad n=k,\qquad e=k-j;\\
\text{(P2)}\quad n\in\,]k_{0a},k],\quad e=n-k_{0a}\le n-j,\quad\text{or}\quad
\text{in }\mathsf{C}:\ \text{weight }1,\ e=0;\\
\text{(P3)}\quad n=\text{level of the frozen cell},\qquad e\le n-j .
\end{gathered}
\end{equation}

(P1) The point lies in the running attempt $\mathsf{a}$, on the old ring of its top zone. The
weight there is $\mathsf{L}_0^{2(k_a-j)}$ on the top zone's old ring of a pre-attempt level in
$]k_{0a},k_a]$; here that level is $j$, $k_a=k$ and $k_{0a}=k_0(g_k)$, so $j\in\,]k_{0a},k]$,
$n=k$ and $e=k-j$. This is the case treated by \eqref{8.8:eq-healed-rim-bound-a} and
\eqref{8.8:eq-healed-rim-bound-b}.

(P2) The point lies in the running attempt $\mathsf{a}$, in its strata. The weight is
$\mathsf{L}_0^{2e}$, with $e$ equal to the level of the window cell minus $k_{0a}$, on a window
cell $\Delta'$ of $\mathsf{a}$ under which the pre-attempt label has levels $\le k_{0a}$; here the
label means the sets $Z_j$, $\Omega_j$, together with the cells and their levels, that the
correlation theorem attaches to the cover. These cells are the strata of levels
$]k_{0a},k_a[$ and the top shell off $\Omega_{k_{0a}+1}$. A healed level $j\le k_{0a}$ therefore
lies in one of two places:
\begin{itemize}
\item under a stratum cell of some level $n\in\,]k_{0a},k]$, with $e=n-k_{0a}$, and $e\le n-j$
because $j\le k_{0a}$;
\item in the core $\mathsf{C}$, of weight $w(\Delta')=1$, with $n$ the level of its cell and
$e=0$.
\end{itemize}

(P3) The point lies inside a live arrest $\mathfrak{a}\subset C$; we write $k_{\mathfrak{a}}$ for
its top level, which does not exceed the level bound for live arrests in the correlation theorem,
and $X_{\mathfrak{a}}\subset Z_{k_{\mathfrak{a}}}\cap C$. For $j<k_{\mathfrak{a}}$ the correlation
theorem's inclusion of rings gives
\[
\partial\Lambda_j^{\circledR}\subset Z_{j+1}\subset X_{\mathfrak{a}} .
\]
On a live piece the correlation theorem assigns the frozen weight $w_{\mathfrak{a}}$, so the ring
lies under the arrest's frozen window with weight $w_{\mathfrak{a}}$. The reading level $n$ is the
level of the frozen cell, and the weight exponent is likewise $\le n-j$.

\emph{The ratio bound in each placement.} In each placement of
\eqref{8.8:eq-rim-per-cell-b} the weight exponent is at most $n-j$ for the reading level $n$. Since
the bounds $\bar{\mathsf{w}}^3\le 27\mathsf{L}_0^{6(e+1)}$ and $\mathsf{L}_0^{2e}\le (L/3)^e$ are
increasing in $e$, the ratio bound applies verbatim with $e:=n-j$ and $t=L^{j-n}$. Using
$(L/3)^{3(n-j)}L^{-4(n-j)}=[27L]^{-(n-j)}$ and the elementary bound behind
\eqref{8.8:eq-healed-rim-bound-b}, we obtain
\begin{equation}\label{8.8:eq-rim-per-cell-1}
\begin{split}
\bar{\mathsf{w}}^3\,t^4\rho^2_{n,j}
&\le 27\mathsf{L}_0^6\,(1+o(1))\,[27L]^{-(n-j)}
\bigl(1+(3\lambda(n-j)+2c_2)g^2\bigr)\\
&\le C_w .
\end{split}
\end{equation}
In the tail class, $t^4$ is replaced by $t^5$ as in \eqref{8.8:eq-healed-rim-bound-c}. So no
factor $L^{(2\hat{\alpha}-1)(\cdot)}$ forms in any placement: the arithmetic of
\eqref{8.8:eq-rim-per-cell-1} does not depend on the placement.

\emph{The charge.} What does depend on the placement is the clause that charges the weighted
modulus to old zone-$j$ cells.

In (P1) and (P2) the count per old zone-$j$ cell is the running-window branch of the charge
hypothesis of Proposition~\ref{8.8:prop-healed-rim-bound}. That branch is the running-window
clause of the corollary of the correlation theorem on weighted window cells, applied here to the
subtracted value of the healed copy: a current cell
$\Delta'$ of the running window carrying the weight $w(\Delta')=\mathsf{L}_0^{2e}$ holds at least
$L^{de}$ cells of the old determining set $\{\Gamma^0_j\}$.

In (P3) the current $\{\Gamma^0_j\}$ is the frozen, coarser label. There the frozen-window branch
of the same corollary places the weighted mass extensively in the domain
$\mathsf{P}(\mathfrak{a})$ of the arrest, per pre-arrest cell. This goes through the weighted
per-old-cell clause of the correlation theorem for frozen windows of live arrests; the remark
accompanying that per-old-cell clause notes that a count per current cell would instead
carry a factor $\check{R}^{d}$, with $\check{R}$ the radius letter of the correlation theorem.
Accordingly, in (P3) the charge is not per cell of $\Gamma^0_j$ but extensive, per pre-arrest
zone-$j$ cell inside $\mathsf{P}(\mathfrak{a})$. This is the frozen-window branch of the charge
hypothesis of Proposition~\ref{8.8:prop-healed-rim-bound}, which is used here for the subtracted
value of the healed copy. With this charge, \eqref{8.8:eq-rim-per-cell-a} holds per pre-arrest
zone-$j$ cell in (P3), and per old zone-$j$ cell of $\Gamma^0_j$ in (P1) and (P2).

\emph{The age.} Under the count at zone $j$, the indexing clause of
Definition~\ref{8.8:def-copy-readings} fixes the dial $\lambda_{j,k}$: the copy, pinned in
$\Omega_j\setminus\Omega_{j+1}$, is the use $(k,j,\cdot)$. The lateness $k-j$ is then paid as
follows:
\begin{itemize}
\item clause (d) of the large-field theorem on the occurrence of old regions, combined with the
geometric forgetting of old regions, gives the factor $q^{(k-j-1)_+}$, $q$ being the forgetting
ratio;
\item this factor enters through the age bound of the vector package,
\[
\epsilon^V_k\sum_{n\le k}(k-n+1)\,q^{(k-n-1)_+}\,T^V_n ,
\]
in its term with $n=j$, namely $\epsilon^V_k\,(k-j+1)\,q^{(k-j-1)_+}\,T^V_j$;
\item by \eqref{8.8:eq-rim-per-cell-a}, the healed copy lies inside the rim entry of $T^V_j$
in that term.
\end{itemize}
This uses the age bound of the vector package in the form in which a use at level $k$ of a rim
unit of zone $j$ is admitted among its terms. That form is a hypothesis of
Proposition~\ref{8.8:prop-healed-rim-bound}. Without it, placing the own rim in zone $j$ admits
that rim unit to its own level and to no later one, which excludes exactly the use $(k,j,\cdot)$.
The age step depends on this hypothesis.

\emph{Constants.} Every condition used is one-sided. No constant of the vector package changes
except two:
\begin{itemize}
\item the rim constant in $C_L^V$, by the pure factor $2C_w$ of \eqref{8.8:eq-rim-per-cell-a};
\item the decay rate in the tail class $X\not\subset\tilde{\square}^2$, halved to
$\tfrac{15}{16}\kappa$ inside the rate accounting of the terms of the vector package, as in
\eqref{8.8:eq-healed-rim-bound-c}.
\end{itemize}
The only floor used is the correlation theorem's own parameter floor $\kappa\ge 10\log L$; no
new floor is introduced. No cap in the sense of Definition~\ref{2.3:def-cap} is used.
\end{proof}

\begin{remark}[Charging subtracted values to old zone cells]\label{8.8:rem-old-cell-charging}
Let $Z$ be a large-field region healed by the large-field operation $\mathbf{R}_k$, let $C$ be the
treated component, and let $j\ge j(Z)-1$ be a healed level. By
Definition~\ref{8.8:def-copy-readings}, the copy of a level-$j$ term subtracted at this healing is
\emph{read} in a cell of level $k$ of the running window of $\mathbf{B}''_k$, with weight
$\mathsf{L}_0^{2(k-j)}$. Its modulus is \emph{charged} per old zone-$j$ cell of $\Gamma^0_j$, with
the window weight absorbed. The healed-rim bound without excess
(Proposition~\ref{8.8:prop-healed-rim-bound}) and the healed-rim bound per old zone cell,
\eqref{8.8:eq-rim-per-cell-a}, pass from the first of these descriptions to the second. The
remainder of this remark states exactly what that passage uses.

\smallskip
(a) \emph{The per-cell clause and its two branches.} In the correlation theorem, the passage from
a reading in a running-window cell to a charge per old cell is the running-window per-old-cell
clause of its corollary on pointed units. That clause states the following. A current cell
$\Delta$ of the running window carrying the weight $w(\Delta)=\mathsf{L}_0^{2e}$, where $e\ge0$ is
its weight exponent, contains at least $L^{de}$ cells of the old determining set $\{\Gamma^0_j\}$,
and these are the cells counted in \cite[(1.69)]{B-CMP122b} and in the per-cell clause of the
correlation theorem. The clause is stated for the weighted mass of the pieces of pointed units
that are plain interior units, under the condition of the correlation theorem on such units. It is
not stated for a value subtracted at a healing.

For copies lying inside a live arrest $\mathfrak{a}$, the correlation theorem proves the analogous
per-old-cell bound only for the frozen windows of live arrests, again for the weighted mass of
pointed pieces and not for a subtracted value. In that case the weighted mass is bounded
extensively over $\mathsf{P}(\mathfrak{a})$, per pre-arrest cell.

Accordingly, Proposition~\ref{8.8:prop-healed-rim-bound} and the per-cell bound
\eqref{8.8:eq-rim-per-cell-a} are proved as implications. At this step the hypothesis they assume
consists of the following clause, which has two branches, together with the form of the
vector package specified in item~(c) below.
\begin{itemize}
\item[(i)] Consider a copy in the first or the second placement listed with
\eqref{8.8:eq-rim-per-cell-b}. Let $k_{0\mathsf{a}}$ denote the base level of the running attempt
$\mathsf{a}$, whose top level is $k$. In the first placement the copy lies on the old ring of the
top zone of the running attempt, and $j\in\,]k_{0\mathsf{a}},k]$. In the second placement it lies
under a stratum cell or in the core $\mathsf{C}$ of the running attempt, and
$j\le k_{0\mathsf{a}}$. For such a copy, the
running-window per-old-cell clause holds with the weighted mass of a pointed piece replaced by the
modulus of the subtracted copy.
\item[(ii)] Consider a copy in the third placement, that is, inside a live arrest
$\mathfrak{a}\subset C$ of top level $k_{\mathfrak{a}}\le h$, where $h$ is the level bound of the
live arrest, with $X_{\mathfrak{a}}\subset Z_{k_{\mathfrak{a}}}\cap C$, where
$Z_{k_{\mathfrak{a}}}$ is the large-field region at level $k_{\mathfrak{a}}$, under its frozen
window with the frozen weight
$w_{\mathfrak{a}}$. For such a copy, the extensive bound over $\mathsf{P}(\mathfrak{a})$ per
pre-arrest zone-$j$ cell holds with the same replacement.
\end{itemize}

\smallskip
(b) \emph{The third placement.} In the third placement the current $\{\Gamma^0_j\}$ is the frozen,
coarser label. Branch~(ii) therefore contains two further statements.
\begin{itemize}
\item[(1)] The first is that such a copy is counted at a definite zone index $i$ in the quantity
$\lambda_{i,k}$ that the one-lattice Lipschitz step attaches to a use at level $k$ of a term
counted at zone $i$, where $i=j$ or $i=n'$, the level of the frozen cell. Item~(c) treats $i=j$;
for $i=n'$ the quantity entering item~(c) is $\lambda_{n',k}$ in place of $\lambda_{j,k}$.
\item[(2)] The second is an extensive entry of the type $\mathsf{P}(\mathfrak{a})$ among the
per-level entries of the vector package, whose entries are otherwise per cell.
\end{itemize}
The size estimate is the same in all three placements: in each of them the weight exponent is at
most $n-j$ for the reading level $n$. Since the window weight $\mathsf{L}_0^{2e}$ increases with
$e$, the ratio bound \eqref{8.8:eq-healed-rim-bound-b} therefore applies with $e$ replaced by
$n-j$, and no factor
$L^{(2\hat{\alpha}-1)(\cdot)}$ arises in any placement.

\smallskip
(c) \emph{The age step.} For the count at zone $j$, the quantity of item~(b)(1) is
$\lambda_{j,k}$. The lateness $k-j$ is paid by the age factor $q^{(k-j-1)_+}$ of the one-lattice
Lipschitz step, $q$ denoting here the ratio of geometric forgetting of old regions, through the
term $i=j$ of the remainder sum of the vector package
\begin{equation*}
\epsilon^V_k\sum_{i\le k}(k-i+1)\,q^{(k-i-1)_+}\,T^V_i .
\end{equation*}
This step uses the per-level entries $T^V_i$ of the vector package in the following form,
which is part of the hypothesis stated in item~(a).
The table of entries contains a rim entry for the healed patch, counted per old zone-$j$ cell and
read at the level of its window. The rim clause of the package admits a use at level $k$ of a
zone-$j$ rim unit of a healed patch. The entry of the package for the operation
$\mathbf{R}$ names the clause of item~(a).

\smallskip
(d) \emph{Conventions of citation.} In \cite[(2.2)]{B-CMP119}, as recalled in
Definition~\ref{8.8:def-copy-readings}, the zones of the determining set are given by
\begin{equation*}
\Gamma_0=\Omega_1^c,\qquad \Gamma_j=\Omega_j^{(j)}\setminus\Omega_{j+1}^{(j)},\qquad \dots,\qquad
\mathbf{B}=\bigcup_j\Gamma_j ,
\end{equation*}
the set $\mathbf{B}$ being called the determining set, with notation as in
\cite[Sect.~0]{B-CMP109}. In this notation the superscript $(j)$ of
$\Omega_j^{(j)}$ is read here as indicating that $\Omega_j$ is taken as a union of cells of
level $j$; the argument uses this reading. The regions $\Omega_1,\Omega_2,\dots$ through which the
one-lattice Lipschitz step defines the age factor used in item~(c) are identified with the regions
of the same names in the correlation theorem. The reading level $n$ of a copy is the level of the
corresponding window cell in the placement clause of the correlation theorem.

\smallskip
(e) \emph{The two levels of a copy.} In the charge, the zone in which a copy is counted is its old
zone $j$. The difference between the counting zone and $j$ is therefore zero, so the load entry
compared with, $\varsigma^V_{1+(\text{counting zone}-j)}$, is $\varsigma^V_1$, as in
\eqref{8.8:eq-rim-per-cell-a}. The level $n=k$ is the reading level of the same copy. It enters
the size of the copy through the scale ratio $t=L^{j-k}$, which satisfies $t\le1$ because $j\le k$,
and, together with the window weight, through the ratio bound
\eqref{8.8:eq-healed-rim-bound-b}; it can only make the size of the copy smaller. The two indices
are different and do not conflict. Finally, with $N_0(g_k)$ the age threshold at level $k$, the
window $k-j<N_0(g_k)$ is where the first placement occurs; it contributes no factor to the
bound.
\end{remark}

\begin{definition}[Territory age and seed age of a live region]\label{8.9:def-seed-age}
Let $K'=K-m_{\mathbb{T}}(g)$ be the terminal level, $k_0$ the witness level, and $\tilde{\nu}$ the
annealed normalised large-field measure on histories $h$. For a level $k\le K'$, let
$\mathcal{O}_k(h)$ be the live large-field region of $h$ at level $k$ and $|\mathcal{O}_k(h)|$
the set of points it covers. Here a large-field region is live at level $k$ when it still carries
its large-field factor at that level.

\emph{Territory age.} For a live component $Z$ of $\mathcal{O}_k(h)$, $j(Z)$ is the level of
the oldest layer of its territory; $\operatorname{age}_k(Z):=k-j(Z)$ is its
\emph{territory age} at level $k$.

\emph{Live-cell events.} For a level-$k$ cell $\square$ and $a\ge 0$ we put
\begin{equation}\label{8.9:eq-seed-age-1}
\begin{gathered}
\mathfrak{lv}_k(\square):=\{h:\ \square\subset|\mathcal{O}_k(h)|\},\\
\mathfrak{lv}^{=a}_k(\square):=\{h:\ \square\subset Z\text{ for a live component }Z\\
\qquad\text{of }\mathcal{O}_k(h)\text{ with }\operatorname{age}_k(Z)=a\},
\end{gathered}
\end{equation}
so that $\mathfrak{lv}_k(\square)$ is the disjoint union of the events
$\mathfrak{lv}^{=a}_k(\square)$ over $a\ge0$.

\emph{Live-cell probabilities.} We put
$q^{=a}_k:=\max_{\square}\tilde{\nu}(\mathfrak{lv}^{=a}_k(\square))$. The
\emph{age-weighted live-cell inequality at level $k$} is the assertion
\begin{equation}\label{8.9:eq-seed-age-a}
\begin{gathered}
\sum_{a>\mathfrak{l}(x_k)}L^{8(a+1)}\,q^{=a}_k\ \le\ A_{\mathrm{lv}}\,
L^{8(\mathfrak{l}(x_k)+1)}\,q^{\mathrm{lv}}_k,\qquad\text{where}\\
q^{\mathrm{lv}}_k:=C_{\mathrm{lv}}\,e^{-c_{\mathrm{lv}}\,p_0(g_k)},\qquad
(C_{\mathrm{lv}},c_{\mathrm{lv}})=(C_L,\theta),\\
C_L:=K_He^{\theta}\bigl(1+1/\lambda^{\sharp}\bigr)+1,\qquad
\mathfrak{l}:=L\check{R}_{k_0}+c_L,
\end{gathered}
\end{equation}
in which $A_{\mathrm{lv}}\ge0$ is either a constant depending on $d$, $L$, $M$ and further fixed
constants, or a polynomial in $\log x_k$. Here $q^{\mathrm{lv}}_k$ is the live rate,
$x_k=g_k^{-2}$, $p_0(\cdot)$ is
the degree function, $\theta$ is the reserve parameter, $\check{R}_{k_0}$ is
Ba{\l}aban's number $R$ at level $k_0$, $K_H$ is a constant of the live rate, and $c_L$ is a
fixed additive constant of the lifetime count $\mathfrak{l}$, distinct from $C_L$.

\emph{Seed age.} Let $h$ be a history and $C$ a live component of $\mathcal{O}_k(h)$. Its
\emph{internal history}
$\mathfrak{h}$ consists of the live ancestor components of $C$ together with every label merged
into them, and its \emph{seed set} $\mathcal{Z}(\mathfrak{h})$
consists of the seeds of the first kind,
which are the components of the cover of a new large-field region \cite{B-CMP122b}, and
of the preparatory seeds. Each seed $Y$ has a birth level, and on seeds we
write $j(Y)$ for it; on live components $j$ keeps the meaning fixed above. We put
\begin{equation}\label{8.9:eq-seed-age-2}
j_{\mathrm{seed}}(C):=\min_{Y\in\mathcal{Z}(\mathfrak{h})}j(Y),\qquad
\operatorname{age}^{\mathrm{seed}}_k(C):=k-j_{\mathrm{seed}}(C).
\end{equation}
Thus $j_{\mathrm{seed}}(C)$ is the index that \cite{B-CMP122b} attaches to a large-field
region, namely the index of a first large-field region contained in it, read on the live
component $C$. Merged regions inherit the oldest
birth index \cite[(1.85)--(1.86)]{B-CMP122b}. For $a\ge0$ and a level-$k$ cell $\square$ we put
\begin{equation*}
\begin{gathered}
\mathfrak{lv}^{\ge a,\mathrm{seed}}_k(\square):=\{h:\ \square\subset C\text{ for a live component }
C\\
\qquad\text{of }\mathcal{O}_k(h)\text{ with }\operatorname{age}^{\mathrm{seed}}_k(C)\ge a\}.
\end{gathered}
\end{equation*}

With these definitions, for every history $h$ every live component $C$ of $\mathcal{O}_k(h)$
satisfies
\begin{equation}\label{8.9:eq-seed-age-3}
\operatorname{age}_k(C)\le\operatorname{age}^{\mathrm{seed}}_k(C),
\end{equation}
and for every $a\ge0$ and every level-$k$ cell $\square$,
\begin{equation}\label{8.9:eq-seed-age-4}
\mathfrak{lv}^{\ge a,\mathrm{seed}}_k(\square)\supset\mathfrak{lv}^{=a}_k(\square).
\end{equation}
The bounds on live-cell probabilities of this chapter are stated for the events
$\mathfrak{lv}^{\ge a,\mathrm{seed}}_k(\square)$; by \eqref{8.9:eq-seed-age-4} a tail bound in the
seed age dominates the slice in the territory age, and so controls the
quantities $q^{=a}_k$ entering \eqref{8.9:eq-seed-age-a}.
\end{definition}

\begin{proof}
By the definitions of the two ages, inequality \eqref{8.9:eq-seed-age-3} is equivalent to
$j(C)\ge j_{\mathrm{seed}}(C)$. In words, the oldest layer of the territory of $C$ was created no
earlier than the birth of the oldest seed in $\mathcal{Z}(\mathfrak{h})$.

A point of the territory of $C$ leaves the small-field
region in one of two ways only. In the first way the point
lies inside a new large-field region; in the second it lies in a layer added to a live region. In
the first case the components of the cover of that new region are seeds of the first kind in
$\mathcal{Z}(\mathfrak{h})$. Their birth level is the level of that creation, and it is at least
$j_{\mathrm{seed}}(C)$ by \eqref{8.9:eq-seed-age-2}. In the second case the live region to which
the layer is added is an ancestor component of $C$. The layer is added at or after the birth of
that region, hence at or after the birth of the oldest seed. So every layer of the territory of
$C$ is created at or after the birth of the oldest seed.

If a region is healed and later created again, then by the corresponding bound of the scaffold
construction for old regions the re-creation can only make the oldest layer of the territory
younger; in particular it does not move $j(C)$ below $j_{\mathrm{seed}}(C)$. This proves
$j(C)\ge j_{\mathrm{seed}}(C)$, that is, \eqref{8.9:eq-seed-age-3}.

For \eqref{8.9:eq-seed-age-4}, let $h\in\mathfrak{lv}^{=a}_k(\square)$. Then
$\square\subset C$ for a live component $C$ of $\mathcal{O}_k(h)$ with
$\operatorname{age}_k(C)=a$. By \eqref{8.9:eq-seed-age-3},
$\operatorname{age}^{\mathrm{seed}}_k(C)\ge a$. Hence $h\in\mathfrak{lv}^{\ge
a,\mathrm{seed}}_k(\square)$.

Finally, by \eqref{8.9:eq-seed-age-4} and monotonicity of $\tilde{\nu}$,
$\tilde{\nu}(\mathfrak{lv}^{=a}_k(\square))\le\tilde{\nu}(\mathfrak{lv}^{\ge
a,\mathrm{seed}}_k(\square))$ for every $\square$. Taking the maximum over $\square$, every bound
on $\max_\square\tilde{\nu}(\mathfrak{lv}^{\ge a,\mathrm{seed}}_k(\square))$ is also a bound on
$q^{=a}_k$.
\end{proof}

We use the notions of Definition~\ref{8.9:def-seed-age}:
\begin{itemize}
\item a history $h$ with its live large-field region $\mathcal{O}_k(h)$ at level $k$;
\item a live component $C$ of $\mathcal{O}_k(h)$ with its internal history $\mathfrak{h}$ and its seed set
$\mathcal{Z}(\mathfrak{h})$, whose seeds are of the first kind (components of the cover of a new
large-field region) or preparatory;
\item the birth level $j(Z)$ of a seed $Z$;
\item the seed age $\operatorname{age}^{\mathrm{seed}}_k(C)=k-j_{\mathrm{seed}}(C)$, where
$j_{\mathrm{seed}}(C)=\min_{Z\in\mathcal{Z}(\mathfrak{h})}j(Z)$;
\item the territory age $\operatorname{age}_k(C)$, which satisfies
$\operatorname{age}_k(C)\le\operatorname{age}^{\mathrm{seed}}_k(C)$.
\end{itemize}

For a live component $W$ of the history at level $n$ we write $\mathcal{Z}(W)$ for its seed set and $Z_{\mathrm{old}}(W)$ for its oldest seed, that is, a seed of minimal birth level.

Ba{\l}aban's number $R$ at level $n$ is the number written $\check{R}_n$ in the description of the run in the correlation theorem; it is non-increasing in $n$, as stated there, and \cite[\S1]{B-CMP122a} states that ``in some steps the number $R_k$ decreases by the factor $L^{-1}$''. We write $d'_n(Z)$ for the linear size of a domain $Z$ measured in cubes of side $M\check{R}_n$. We write $S(Z)$ for Ba{\l}aban's enlargement of $Z$ in \cite[\S1]{B-CMP122b}: the cover of $Z$ by cubes of the next level, taken with ten layers.

\emph{The survival clock.} Let $W$ be a live component of the large-field region at level $n$. If $W$ does not belong to the class of components described in the healing rule below, its \emph{survival clock} $\Phi_n(W)$ is the integer that Ba{\l}aban writes $K$ in the sentences quoted below; outside quotations we write $\Phi_n(W)$ for it. In the words of \cite[\S1]{B-CMP122b}, for a domain $Z$ created at level $j$,
\begin{quote}
``the number $K$ is the smallest positive integer having the property that the domain $S^K(Z)$, considered as a domain in the lattice of the scale $L^{-(j+K)}$, satisfies the conditions (i), (ii), with $N=R_j$''.
\end{quote}
Here (i), (ii) are Ba{\l}aban's conditions of \cite[\S1]{B-CMP122a}. If $W$ satisfies the conditions (i), (ii) at level $n$, we put $\Phi_n(W):=0$, in accordance with the healing rule below. The lemma below uses four rules for this clock, stated next. The birth rule is among the region rules stated with the large-field block bound in this chapter; the quiet rule follows from the definition of $\Phi$; the two cases of the event rule and the healing rule are taken from the sentences of \cite[\S1]{B-CMP122b} and \cite[\S1]{B-CMP122a} quoted with them and are used as stated there.

\emph{Birth rule.} For a new region $Z$ born at level $j$ put $K_{\mathrm{init}}(Z):=\Phi_j(Z)$. Then
\begin{equation}\label{8.9:eq-age-lemma-1}
K_{\mathrm{init}}(Z)\le 2\log_L\bigl(d'_j(Z)+1\bigr)+\check{R}_j+C_S ,
\end{equation}
with $C_S$ an absolute constant. This is the bound of \cite[(1.81)]{B-CMP122b} and the discussion following it, which runs as follows.
\begin{itemize}
\item The domain $Z^{(n-j)}$ reached from $Z$ after $n-j$ steps satisfies
$d'_n(Z^{(n-j)})\le L^{-(n-j)/2}d'_j(Z)$.
\item Let $n_0$ be the last $n$ with $d'_n(Z^{(n-j)})>0$. Then, in the words of
\cite[\S1]{B-CMP122b}, ``doing at most $R_j$ further steps we obtain a domain satisfying both
conditions (i), (ii). Thus $K\le n_0-j+R_j$''.
\end{itemize}

The constant $C_S$ absorbs the additive $2$ in the cruder form $K_{\mathrm{init}}(Z)\le 2\log_L d'_j(Z)+\check{R}_j+2$ of this bound, and the constant of the ten layers in $d'$.

\emph{Quiet rule.} Suppose that at step $n+1$ no large field is created in $S(W)$, neither in $W$ nor in the added layers $S(W)\setminus W$, and that $W$ joins no other piece. In the words of \cite[\S1]{B-CMP122b}, ``In the first case no large fields were introduced in the last step, hence $Z=S(Z_0)$''. Then
\begin{equation}\label{8.9:eq-age-lemma-2}
\Phi_{n+1}\bigl(S(W)\bigr)=\Phi_n(W)-1 .
\end{equation}
Indeed $S^{\Phi_n(W)-1}(S(W))=S^{\Phi_n(W)}(W)$, and $\Phi$ is defined as a smallest integer.

\emph{Event rule, preparatory case.} This case is described in \cite[\S1]{B-CMP122b} as follows:
\begin{quote}
``$Z=S(Z_0)$, $Z_0$ satisfies the conditions (i), (ii), and a new large field was introduced in the preparatory operations \ldots\ we have the new factor $\exp(-p_0(g_j))$ \ldots\ it is contained in a cube of the size $100MR_{j+1}$, hence $K=R_{j+1}$ for $Z$''.
\end{quote}
In this case a new preparatory seed $P$ is born at step $n+1$, and
\begin{equation}\label{8.9:eq-age-lemma-3}
\Phi_{n+1}\bigl(S(W)\bigr)=\check{R}_{n+1}\le K_{\mathrm{init}}(P):=\check{R}_{n+1}+1 .
\end{equation}

\emph{Event rule, join case.} This case is described in \cite[\S1]{B-CMP122b} as follows:
\begin{quote}
``In the second case $Z$ is obtained from some number of components of $Z_j$ and some number of new large field regions, joined together into the one component of $Z_{j+1}$ by the operations of the last step, in particular by adding layers''.
\end{quote}
\cite[\S1]{B-CMP122b} bounds the clock of a join by an induction ``with respect to the number of domains'', whose step is the following rule for two domains:
\begin{quote}
``The domains $X$, $Y$ determine the corresponding indices $K_1$, $K_2$ \ldots\ we may assume \ldots\ $K_1\le K_2$. The domain $S^{K_1}(X)$ satisfies the conditions (i), (ii) \ldots\ applying $n_1$ times the operation $S$ to the last domain, where $n_1$ is a rather small number, e.g., $n_1<10$, we obtain the domain $S^{K_2+n_1}(Y)$ containing $S^{K_2+n_1}(X)$ \ldots\ and that $K\le K_2+n_1+R_{j+1}$''.
\end{quote}

Let the level-$(n+1)$ component $W'$ be the join of the following pieces:
\begin{itemize}
\item $p\ge0$ old pieces $S(W_1),\dots,S(W_p)$, the piece $S(W_i)$ entering the two-domain rule with
the index $\Phi_n(W_i)-1$ (quiet rule) or with its own event value;
\item $q\ge0$ new regions $X_1,\dots,X_q$.
\end{itemize}

Here $p+q\ge2$, or $p=0$. Put
\begin{equation*}
E_{n+1}:=n_1+\check{R}_{n+1}+2 .
\end{equation*}
The additive $2$ allows for rounding in Ba{\l}aban's bound, and $n_1$ is his absolute ``rather small number''. Iterating this two-domain rule over a spanning tree of the $p+q$ pieces, as in the induction of \cite[\S1]{B-CMP122b}, one is led to the bound
\begin{equation}\label{8.9:eq-age-lemma-4}
\Phi_{n+1}(W')\le\max\Bigl\{\max_{1\le i\le p}\bigl(\Phi_n(W_i)-1\bigr),\;
\max_{1\le l\le q}K_{\mathrm{init}}(X_l)\Bigr\}+(p+q-1)\,E_{n+1},
\end{equation}
each maximum being taken over the pieces present. The bound \eqref{8.9:eq-age-lemma-4} is used below as stated, and its use is the step marked $(\ast)$ in the proof of Lemma~\ref{8.9:lem-age-lemma}.

\emph{Healing rule.} The following are equivalent:
\begin{itemize}
\item $\Phi_n(W)=0$;
\item $W$ satisfies the conditions (i), (ii) at level $n$;
\item $W$ belongs to the class of components on whose union the operation $\mathbf{R}$ acts at that step.
\end{itemize}

The class is the one of \cite[\S1]{B-CMP122a}: ``we consider the class of components such that each satisfies \ldots\ (i) \ldots\ (ii) \ldots\ Let us denote the union of the above class of components by $Z$''. The operation $\mathbf{R}$ is defined on this $Z$. In \cite[\S1]{B-CMP122b}, ``the domain $X$ in the definition (1.71) satisfies the assumption of the statement with $K=0$'' (see \cite[(1.71)]{B-CMP122b}).

Consequently a component with $\Phi=0$ is healed at that step: it is absorbed by $\mathbf{R}$ and leaves the live region \cite[\S1]{B-CMP122b}, by the description of the live large-field region given with the large-field block bound. The exception is when the preparatory case of the event rule occurs. The component then stays live, with a new preparatory seed and with $\Phi$ equal to the value in \eqref{8.9:eq-age-lemma-3}.

Hence, along a live lineage, $\Phi\ge1$ at every level at which the component is live and not in the class. Moreover, every level is a step of the quiet rule or a step of the event rule.

In \cite[\S1]{B-CMP122b} the parameter of condition (ii), written $N$ there, is fixed at the value $\check{R}_j$, the value at the start of the clock. Suppose the healing step were instead to use the value $\check{R}_n$ at the current level $n$. Since $\check{R}$ is non-increasing, condition (ii) becomes weaker and healing occurs no later. Lemma~\ref{8.9:lem-age-lemma} then holds a fortiori. The lemma does not use this alternative.

\begin{lemma}[The age lemma. Stated; proof incomplete at the step marked $(\ast)$]
\label{8.9:lem-age-lemma}
Let the number of renormalisation steps $K$ be given, consider an admissible run in the sense of the description of the run in the correlation theorem, and let $k$ be a level. Let $C$ be a component that is live at level $k$, with internal history $\mathfrak{h}$ and seed set $\mathcal{Z}(\mathfrak{h})$ (seeds of both kinds), and let $Z_{\mathrm{old}}$ be its oldest seed. Then
\begin{equation}\label{8.9:eq-age-lemma-5}
\operatorname{age}^{\mathrm{seed}}_k(C)
\le\operatorname{age}^{\mathrm{seed}}_k(C)+\Phi_k(C)+E(Z_{\mathrm{old}})
\le\sum_{Z\in\mathcal{Z}(\mathfrak{h})}\mathfrak{a}(Z),
\end{equation}
where the charge of a seed $Z$ is
\begin{equation}\label{8.9:eq-age-lemma-a}
\mathfrak{a}(Z):=K_{\mathrm{init}}(Z)+E(Z),\qquad E(Z):=n_1+\check{R}_{j(Z)}+2 .
\end{equation}
With the birth rule \eqref{8.9:eq-age-lemma-1}, for seeds of the first kind
\begin{equation}\label{8.9:eq-age-lemma-6}
\mathfrak{a}(Z)\le 2\log_L\bigl(d'_{j(Z)}(Z)+1\bigr)+2\check{R}_{j(Z)}+n_1+C_S+4 .
\end{equation}
For preparatory seeds, using $\check{R}_{j+1}\le\check{R}_j$,
\begin{equation}\label{8.9:eq-age-lemma-7}
\mathfrak{a}(Z)\le 2\check{R}_{j(Z)}+n_1+4 .
\end{equation}
In particular, $\operatorname{age}_k(C)\le\sum_{Z\in\mathcal{Z}(\mathfrak{h})}\mathfrak{a}(Z)$ for the territory age of Definition~\ref{8.9:def-seed-age}.
\end{lemma}

Stated; the proof below is incomplete at the step marked $(\ast)$. The step marked $(\ast)$ is the appeal to \eqref{8.9:eq-age-lemma-4}, taken by statement from \cite[\S1]{B-CMP122b}.

\begin{proof}
For a live component $W$ of the history at its level $n$, put
\begin{equation*}
\Psi(W):=\sum_{Z\in\mathcal{Z}(W)}\mathfrak{a}(Z).
\end{equation*}
We argue by induction along the levels of the history of $C$. The invariant for $W$ is the inequality
\begin{equation}\label{8.9:eq-age-lemma-8}
\operatorname{age}^{\mathrm{seed}}_n(W)+\Phi_n(W)+E\bigl(Z_{\mathrm{old}}(W)\bigr)\le\Psi(W),
\end{equation}
to hold for every live component $W$ of the history at its level $n$. In the steps of the induction below, for the components $W$, $W_i$, $W'$ of the history we abbreviate $\operatorname{age}^{\mathrm{seed}}_n(W)$ to $\operatorname{age}_n(W)$; the territory age does not occur in these steps.

\emph{Birth.} Let $W$ be a newborn component at level $j$, made of $q\ge1$ new regions $X_1,\dots,X_q$ and no old piece. This is the join case of the event rule with $p=0$, or $q=1$ with no join.
\begin{itemize}
\item Its seed age is $0$, since every seed is born at level $j$.
\item By \eqref{8.9:eq-age-lemma-4} with $p=0$ (the input of the step marked $(\ast)$ below), and
for $q=1$ by the definition $K_{\mathrm{init}}(X_1)=\Phi_j(W)$,
\begin{equation*}
\Phi_j(W)\le\max_lK_{\mathrm{init}}(X_l)+(q-1)E_j\le\sum_lK_{\mathrm{init}}(X_l)+(q-1)E_j .
\end{equation*}
\item Since $j(X_l)=j$ for every $l$, we have $E_j=E(X_l)$ for every $l$. The oldest seed is one of
the $X_l$, so
\begin{equation*}
E(Z_{\mathrm{old}})+(q-1)E_j\le\sum_lE(X_l).
\end{equation*}
\end{itemize}

Adding these,
\begin{equation*}
\operatorname{age}_j(W)+\Phi_j(W)+E(Z_{\mathrm{old}})
\le\sum_l\bigl(K_{\mathrm{init}}(X_l)+E(X_l)\bigr)=\Psi(W),
\end{equation*}
which is the invariant \eqref{8.9:eq-age-lemma-8} for $W$.

\emph{Quiet step.} Under the quiet rule \eqref{8.9:eq-age-lemma-2}, passing from $W$ at level $n$ to $S(W)$ at level $n+1$:
\begin{itemize}
\item the seed age increases by $1$;
\item the clock decreases by $1$;
\item the seed set and the oldest seed are unchanged.
\end{itemize}

Both members of \eqref{8.9:eq-age-lemma-8} are therefore unchanged, and the invariant is preserved.

\emph{Preparatory refresh.} Consider the preparatory case of the event rule at step $n+1$. A new seed $P$ is born, the seed set becomes $\mathcal{Z}(W)\sqcup\{P\}$, and the oldest seed is unchanged because $P$ is born at the latest level. The seed age increases by $1$. Let $W'=S(W)$ be the component at level $n+1$. By \eqref{8.9:eq-age-lemma-3}, $\Phi_{n+1}(W')=\check{R}_{n+1}=K_{\mathrm{init}}(P)-1$. Since $\Phi_n(W)\ge0$, we obtain
\begin{align*}
\operatorname{age}_{n+1}(W')+\Phi_{n+1}(W')+E(Z_{\mathrm{old}})
&=(\operatorname{age}_n(W)+1)+\bigl(K_{\mathrm{init}}(P)-1\bigr)+E(Z_{\mathrm{old}})\\
&=\operatorname{age}_n(W)+K_{\mathrm{init}}(P)+E(Z_{\mathrm{old}})\\
&\le\bigl(\operatorname{age}_n(W)+\Phi_n(W)+E(Z_{\mathrm{old}})\bigr)+K_{\mathrm{init}}(P)\\
&\le\Psi(W)+\mathfrak{a}(P)=\Psi(W').
\end{align*}
In the last line we used the invariant for $W$ and $K_{\mathrm{init}}(P)\le\mathfrak{a}(P)$.

\emph{Join.} Consider the join case of the event rule at step $n+1$, with $p\ge1$ old pieces and $q\ge0$ new regions $X_1,\dots,X_q$, where $p+q\ge2$.
\begin{itemize}
\item Each old piece $W_i$ satisfies the invariant \eqref{8.9:eq-age-lemma-8} for $W_i$ at level $n$
and has oldest seed $Z_i$.
\item Let $i_0$ be an index with $j(Z_{i_0})$ minimal. The new regions are born at level $n+1$ and
every $Z_i$ at a level at most $n$, so the oldest seed of $W'$ is $Z_{i_0}$.
\end{itemize}

\emph{The age.} Since $\operatorname{age}_n(W_i)=n-j(Z_i)$, we get
\begin{equation*}
\operatorname{age}_{n+1}(W')=\max_i\operatorname{age}_n(W_i)+1 .
\end{equation*}

\emph{The clock.} $(\ast)$ By inequality \eqref{8.9:eq-age-lemma-4} of the join case of the event rule,
\begin{equation*}
\Phi_{n+1}(W')\le\max\Bigl\{\max_i\bigl(\Phi_n(W_i)-1\bigr),\;\max_lK_{\mathrm{init}}(X_l)\Bigr\}
+(p+q-1)E_{n+1}.
\end{equation*}
Both inner maxima are non-negative: $\Phi_n(W_i)\ge1$ by the healing rule, and, if $q\ge1$, $K_{\mathrm{init}}(X_l)\ge1$; if $q=0$ the second maximum is absent. Hence the outer maximum is at most their sum. We bound the resulting terms in three groups.
\begin{enumerate}
\item Since $\Phi_n(W_i)$ and $\operatorname{age}_n(W_i)$ are all non-negative,
\begin{align*}
\max_i\bigl(\Phi_n(W_i)-1\bigr)+\max_i\operatorname{age}_n(W_i)+1
&\le\max_i\Phi_n(W_i)+\max_i\operatorname{age}_n(W_i)\\
&\le\sum_i\bigl(\Phi_n(W_i)+\operatorname{age}_n(W_i)\bigr).
\end{align*}
If both maxima are attained at one index the right side contains that summand; otherwise it contains
the two summands at which they are attained.
\item Trivially,
\begin{equation*}
\max_lK_{\mathrm{init}}(X_l)\le\sum_lK_{\mathrm{init}}(X_l).
\end{equation*}
\item Every seed $Z$ present has $j(Z)\le n+1$, and $\check{R}$ is non-increasing. Hence
\begin{equation*}
E_{n+1}=n_1+\check{R}_{n+1}+2\le n_1+\check{R}_{j(Z)}+2=E(Z)
\end{equation*}
for every such $Z$. The seeds $Z_i$ with $i\ne i_0$ and the $X_l$ give exactly $(p-1)+q$ such terms, so
\begin{equation*}
(p+q-1)E_{n+1}\le\sum_{i\ne i_0}E(Z_i)+\sum_lE(X_l).
\end{equation*}
\end{enumerate}

\emph{Combining.} Adding these bounds and using the invariant for each $W_i$,
\begin{align*}
\operatorname{age}_{n+1}(W')+\Phi_{n+1}(W')+E(Z_{i_0})
&\le\sum_i\bigl[\operatorname{age}_n(W_i)+\Phi_n(W_i)+E(Z_i)\bigr]
+\sum_l\bigl[K_{\mathrm{init}}(X_l)+E(X_l)\bigr]\\
&\le\sum_i\Psi(W_i)+\sum_l\mathfrak{a}(X_l)=\Psi(W').
\end{align*}
The last equality holds because the seed sets of the pieces are disjoint, and their union with $\{X_1,\dots,X_q\}$ is $\mathcal{Z}(W')$.

\emph{A piece staying live by a refresh.} Suppose an old piece $W_i$ is in the class at level $n$, and stays live only by a preparatory refresh at the same step at which it joins. It then enters the rule with the value $\check{R}_{n+1}\le K_{\mathrm{init}}(P_i)$ in place of $\Phi_n(W_i)-1$, and with the new seed $P_i$. The same lines give the bound with $\mathfrak{a}(P_i)$ added on the right, and $\Psi(W')$ contains this term.

\emph{Conclusion.} By the healing rule, the steps of the quiet rule, the preparatory case and the join case of the event rule exhaust the levels of a live lineage. The invariant therefore holds for $C$ at level $k$, which gives
\begin{equation*}
\operatorname{age}^{\mathrm{seed}}_k(C)\le\Psi(C)-\Phi_k(C)-E(Z_{\mathrm{old}})\le\Psi(C).
\end{equation*}
The first inequality of \eqref{8.9:eq-age-lemma-5} holds because $\Phi_k(C)\ge0$ and $E(Z_{\mathrm{old}})\ge0$.

\emph{The bounds on the charges.} For a seed of the first kind, insert \eqref{8.9:eq-age-lemma-1} into \eqref{8.9:eq-age-lemma-a}:
\begin{equation*}
\mathfrak{a}(Z)\le 2\log_L\bigl(d'_{j(Z)}(Z)+1\bigr)+2\check{R}_{j(Z)}+n_1+C_S+2 ,
\end{equation*}
and a fortiori \eqref{8.9:eq-age-lemma-6}. For a preparatory seed $P$ born at step $n+1$ we have $j(P)\le n+1$ (whichever endpoint of the step is taken as the birth level). By \eqref{8.9:eq-age-lemma-3} and since $\check{R}$ is non-increasing,
\begin{equation*}
K_{\mathrm{init}}(P)=\check{R}_{n+1}+1\le\check{R}_{j(P)}+1 ,
\end{equation*}
so $\mathfrak{a}(P)\le2\check{R}_{j(P)}+n_1+3$, and a fortiori \eqref{8.9:eq-age-lemma-7}.

\emph{The territory age.} Finally, the territory age $\operatorname{age}_k(C)$ of Definition~\ref{8.9:def-seed-age} satisfies $\operatorname{age}_k(C)\le\operatorname{age}^{\mathrm{seed}}_k(C)$, which together with \eqref{8.9:eq-age-lemma-5} gives the last assertion.
\end{proof}

We note three features of Lemma~\ref{8.9:lem-age-lemma}.

First, consider how each term of the right side of \eqref{8.9:eq-age-lemma-5} is used.
\begin{itemize}
\item Each seed pays one initial clock. Its birth size enters once, logarithmically, through
\eqref{8.9:eq-age-lemma-1}.
\item Each seed also pays one join extension $E$. The oldest seed of a piece that loses its oldest
status is charged exactly once; the final oldest seed is never charged.
\item Ba{\l}aban's argument with $n_1$ is what makes the join extension independent of the sizes of the
pieces: the piece with the smaller index already satisfies the conditions (i), (ii) when it is
absorbed, so only $n_1$ further enlargements and $\check{R}_{n+1}$ further steps are added,
independently of its size.
\end{itemize}

Second, no entropy, no probability and no field enters. The lemma is a statement about the label structure of one history.

Third, let $\mathfrak{l}=L\check{R}_{k_0}+c_L$ be the lifetime threshold of the annealed live-tail statement, a statement that is not proved and is not used in the proof of any clause of Theorem~0, with $k_0$ the witness level and $c_L$ the constant of that statement. For one seed of bounded birth size, $\mathfrak{a}(Z)$ is at most $2\check{R}_{j(Z)}$ plus a constant independent of $j(Z)$, a bounded multiple of one $\check{R}$, whereas $\mathfrak{l}$ carries the factor $L$. Hence an age exceeding $\mathfrak{l}$ forces $\sum_Z\mathfrak{a}(Z)>L\check{R}_{k_0}+c_L$. This happens in one of three ways:
\begin{itemize}
\item there are many seeds;
\item seeds are born large, with logarithm of the birth size of order $L\check{R}_{k_0}$;
\item seeds are born so deep that $\check{R}_{j(Z)}$ is much larger than $\check{R}_{k_0}$.
\end{itemize}

Each of these alternatives is charged by the reserve weights of the annealed live-tail statement.

\begin{lemma}[The age lemma on the correlation theorem's clock]\label{8.9:lem-age-lemma-correlation}
Consider the large-field labels of a run of Ba{\l}aban's operation $\mathbf{R}$ read on the clock
of the correlation theorem (its count of the operations that remain for a live region), and assume
the following statements about that clock, established in the large-field part of the proof of the
correlation theorem.
\begin{itemize}
\item[(a)] \emph{Events.} An \emph{event} of the lineage of a component (the component itself
together with its live ancestor components at all earlier levels, with the labels merged into
them, as in Definition~\ref{8.9:def-seed-age}) is a pair consisting of a level and a component of
the lineage at that level at which at least one of the following occurs: the component is born, is
joined, is hit by an $\mathrm{I}$-creation (the creation of a region of the domain class that the
correlation theorem's clock denotes by $\mathrm{I}$), or is re-created; several such occurrences on
one component at one level count as one event. The level of an event is the level in this pair.
\item[(b)] \emph{Event-free epochs.} If a component $C$ was born, joined, hit by an
$\mathrm{I}$-creation or re-created at level $e$ and is event-free during $(e,e+t]$, then
\begin{equation*}
t\le 5+\frac{2\operatorname{diam}_e(C)}{L}+\check{R}_e ,
\end{equation*}
and, writing $\operatorname{diam}_{e+u}$ for the diameter at level $e+u$ of the successor of $C$
in its lineage at that level, the diameters along such an epoch obey
\begin{equation*}
\operatorname{diam}_{e+u}\le L^{-\lfloor u/2\rfloor}\operatorname{diam}_e+45 ,\qquad 0\le u\le t .
\end{equation*}
\item[(c)] \emph{Re-creation.} A component living after $\mathbf{R}$ at level $n$ with
$\mathsf{K}_n(Z)=0$, where $\mathsf{K}_n(Z)$ is the remaining count of the correlation theorem's
clock for its region $Z$ at level $n$, was re-created at $n$; a re-creation is an event at which a
new seed of the lineage is born.
\item[(d)] \emph{The join rule.} At an event the diameter of the component grows by at most the
diameters of the constituents it acquires at that event plus $C_J$ for each of them, where
$C_J\ge0$ is the join-layer constant of the clock; this applies to a join of touching or nearby
domains, to new regions and to joined old components alike, the component formed at a join
growing from one of the joined old components and acquiring the others.
\end{itemize}
Let $C$ be a component live at level $k$, with internal history $\mathfrak{h}$ and seed set
$\mathcal{Z}=\mathcal{Z}(\mathfrak{h})$ as in Definition~\ref{8.9:def-seed-age}, and let
$e_1<\dots<e_r\le k$ be the event levels of its lineage, $e_1$ being the birth level of the oldest
constituent, that is, the birth level $j_{\mathrm{seed}}(C)$ of the oldest seed of $C$. Then
\begin{equation}\label{8.9:eq-age-lemma-correlation-a}
\begin{split}
&\operatorname{age}^{\mathrm{seed}}_k(C)\le\sum_{Z\in\mathcal{Z}}\mathfrak{a}_J(Z),\\
&\mathfrak{a}_J(Z):=2\Bigl(5+\check{R}_{j(Z)}+\frac{360}{L}\Bigr)
+\frac{8}{L}\bigl(\operatorname{diam}_{j(Z)}(Z)+C_J\bigr).
\end{split}
\end{equation}
\end{lemma}

\begin{proof}
Let $\mathcal{E}$ be the set of events of the lineage of $C$ at levels at most $k$. For
$\xi\in\mathcal{E}$ let $e(\xi)$ be its level, so that $e_1<\dots<e_r$ are the
values of $e$ on $\mathcal{E}$; let $C_\xi$ be the component of the lineage that carries
$\xi$, taken after the event, and put
$D_\xi:=\operatorname{diam}_{e(\xi)}(C_\xi)$. The successors of
$C_\xi$ at the levels above $e(\xi)$ are ancestors of $C$, hence live up to level
$k$. If one of them takes part in an event of $\mathcal{E}$ at a level in $(e(\xi),k]$,
let $\xi^+$ be the first such event and put $t_\xi:=e(\xi^+)-e(\xi)$;
otherwise put $t_\xi:=k-e(\xi)$. There is exactly one event without a successor
event, which we denote $\bar\xi$: an event of highest level has none, and if two events
had none, the successors of their components would both be $C$ at level $k$, whereas two distinct
components of the lineage become one component only at a join, which is an event. For
$\xi\in\mathcal{E}$ let $\xi^-$ be the set of events $\psi$ with
$\psi^+=\xi$. Every event other than $\bar\xi$ lies in exactly one of the sets
$\xi^-$, and $e(\psi)<e(\xi)$ for $\psi\in\xi^-$.

\emph{Epoch lengths.} The epoch of $\xi$ is the interval
$(e(\xi),e(\xi)+t_\xi]$. Each epoch is event-free for the successor of
$C_\xi$, and this successor lives throughout it, since it is an ancestor of $C$, which is
live at level $k$. At the left end $e(\xi)$ the component $C_\xi$ was born, joined,
hit by an $\mathrm{I}$-creation or re-created, by hypothesis (a). We call the epoch of
$\xi$ \emph{long} if $t_\xi\ge2$ and \emph{short} if $t_\xi\le1$. For a
long epoch, hypothesis (b), applied with $e=e(\xi)$ and $t=t_\xi$, gives
\begin{equation*}
t_\xi\le 5+\frac{2D_\xi}{L}+\check{R}_{e(\xi)} ;
\end{equation*}
for a short epoch $t_\xi\le1\le5+\check{R}_{e(\xi)}$, since
$\check{R}_{e(\xi)}\ge0$; this covers in particular $\bar\xi$ when
$e(\bar\xi)=k$, where $t_{\bar\xi}=0$. By Definition~\ref{8.9:def-seed-age} the
seed age is measured from the birth level of the oldest seed, which is $e_1$; hence
$\operatorname{age}^{\mathrm{seed}}_k(C)=k-e_1$. Let $\xi_1$ be an event of level $e_1$.
The chain $\xi_1,\xi_1^+,(\xi_1^+)^+,\dots$ has strictly increasing levels
and ends at $\bar\xi$; the levels of two consecutive members $\psi,\psi^+$ of the chain differ
by $t_\psi$, and $t_{\bar\xi}=k-e(\bar\xi)$. So $k-e_1$ is the sum of $t_\psi$ over the members
$\psi$ of the chain, which is at most $\sum_{\xi\in\mathcal{E}}t_\xi$ since every
$t_\xi\ge0$. Therefore
\begin{equation}\label{8.9:eq-age-lemma-correlation-1}
\operatorname{age}^{\mathrm{seed}}_k(C)\le\sum_{\xi\in\mathcal{E}}
\bigl(5+\check{R}_{e(\xi)}\bigr)
+\frac{2}{L}\sum_{\xi\in\mathcal{E},\ t_\xi\ge2} D_\xi .
\end{equation}

\emph{Diameters.} Fix $\xi\in\mathcal{E}$. The old pieces entering the event
$\xi$ are the successors, just before the event, of the components $C_\psi$ with
$\psi\in\xi^-$, one for each such $\psi$. Between two events the diameter shrinks as in
the second display of hypothesis (b): applied to the epoch of $\psi$, it bounds the diameter of
the piece coming from $C_\psi$ just before $\xi$ by
$L^{-\lfloor t_\psi/2\rfloor}D_\psi+45$. At the event the diameter grows, by hypothesis (d), by at
most the diameters of the constituents acquired there plus $C_J$ for each, starting from the piece
from which $C_\xi$ grows (no such piece when $\xi^-$ is empty). The acquired
constituents are the seeds born at $\xi$, with their birth diameters, and, when
$|\xi^-|\ge2$, all old pieces but one, that is $|\xi^-|-1$ of them, each carrying
its own diameter, already shrunk in the same way. Writing $(x)_+:=\max\{x,0\}$, this gives
\begin{equation*}
D_\xi\le\sum_{\psi\in\xi^-}L^{-\lfloor t_\psi/2\rfloor}D_\psi+b_\xi,
\end{equation*}
where
\begin{equation*}
b_\xi:=45\,|\xi^-|+C_J\bigl(|\xi^-|-1\bigr)_+
+\sum_{Z\in\mathcal{Z}\ \text{born at}\ \xi}\bigl(\operatorname{diam}_{j(Z)}(Z)+C_J\bigr).
\end{equation*}
For $\xi\in\mathcal{E}$ let $\mathcal{E}(\xi)$ be the set of events $\psi'$ whose
chain $\psi',\psi'^+,\dots$ passes through $\xi$ (so $\xi\in\mathcal{E}(\xi)$),
and for $\psi'\in\mathcal{E}(\xi)$ let $\pi(\psi',\xi)$ be the product of the
factors $L^{-\lfloor t_\psi/2\rfloor}$ over the members $\psi\neq\xi$ of that chain up to
$\xi$, the empty product being $1$. Since each event has at most one successor event,
$\mathcal{E}(\xi)$ is the disjoint union of $\{\xi\}$ and the sets
$\mathcal{E}(\psi)$, $\psi\in\xi^-$, and
$\pi(\psi',\xi)=\pi(\psi',\psi)L^{-\lfloor t_\psi/2\rfloor}$ for
$\psi'\in\mathcal{E}(\psi)$. By induction on $e(\xi)$ the last display therefore gives
\begin{equation}\label{8.9:eq-age-lemma-correlation-3}
D_\xi\le\sum_{\psi'\in\mathcal{E}(\xi)}b_{\psi'}\,\pi(\psi',\xi),
\qquad \xi\in\mathcal{E}.
\end{equation}
The contraction is thus composed epoch by epoch along the chain, and an epoch with $t_\psi\le1$
contributes the factor $1$. Summing \eqref{8.9:eq-age-lemma-correlation-3} over the events with
long epochs and exchanging the order of summation,
\begin{equation*}
\sum_{\xi\in\mathcal{E},\ t_\xi\ge2}D_\xi\le\sum_{\psi'\in\mathcal{E}}
b_{\psi'}\sum_{\xi:\ \psi'\in\mathcal{E}(\xi),\ t_\xi\ge2}
\pi(\psi',\xi).
\end{equation*}
Fix $\psi'$; the events $\xi$ of the inner sum lie on the chain of $\psi'$. For each of
them the product $\pi(\psi',\xi)$ contains, for every earlier member of the chain with a
long epoch, one factor at most $L^{-1}$, because $\lfloor t_\psi/2\rfloor\ge1$ when $t_\psi\ge2$,
and all its other factors are at most $1$. Hence the inner sum is at most
\begin{equation*}
1+L^{-1}+L^{-2}+\cdots=\frac{L}{L-1}\le 2 ,
\end{equation*}
the last inequality because $L\ge2$. It remains to bound $\sum_{\xi}b_\xi$. First,
$\sum_\xi|\xi^-|=|\mathcal{E}|-1$, since every event other than $\bar\xi$
lies in exactly one set $\xi^-$. Second, each seed of $\mathcal{Z}$ is born at most once,
so it enters the last sum in $b_\xi$ for at most one $\xi$. Third, fix once and for all a linear
order of the seeds of $\mathcal{Z}$ that refines the order of their birth levels; call a seed older
than another if it precedes it in this order, and call the first seed of a component in this order
its oldest seed. At an event with
$|\xi^-|\ge2$ the old pieces have disjoint seed sets, exactly one of them contains the
oldest seed of $C_\xi$, and so the oldest seeds of the other $|\xi^-|-1$ pieces
lose the status of oldest seed at $\xi$; a seed that has lost that status lies afterwards
in a component together with an older seed, so it cannot lose it a second time. Hence, since
$C_J\ge0$ and diameters are non-negative,
\begin{equation*}
\sum_\xi\bigl(|\xi^-|-1\bigr)_+\le|\mathcal{Z}|,\qquad
C_J|\mathcal{Z}|\le\sum_{Z\in\mathcal{Z}}\bigl(\operatorname{diam}_{j(Z)}(Z)+C_J\bigr).
\end{equation*}
Altogether we obtain
\begin{equation}\label{8.9:eq-age-lemma-correlation-2}
\sum_{\xi\in\mathcal{E},\ t_\xi\ge2}D_\xi\le 90\,|\mathcal{E}|
+4\sum_{Z\in\mathcal{Z}}\bigl(\operatorname{diam}_{j(Z)}(Z)+C_J\bigr).
\end{equation}

\emph{Events versus seeds.} Every event of the lineage is either the birth of at least one seed of
the lineage or a join of at least two old pieces. Indeed, a birth or an $\mathrm{I}$-creation is
the birth of a seed, and a re-creation is an event with a new seed by hypothesis (c); the remaining
events are joins. We charge each event to one seed. An event at which a seed is born is charged to
one seed born there. A join of at least two old pieces is charged to the oldest seed of a piece
that loses its status of oldest seed at that join; such a piece exists, and a seed loses that
status at most once, as shown in the diameter step above. Thus each seed is charged with at most
one join, and in addition with at most the event of its own birth. This is the same charging as in
the age lemma (Lemma~\ref{8.9:lem-age-lemma}, whose proof is incomplete at one step). Consequently
$|\mathcal{E}|\le2|\mathcal{Z}|$, each seed $Z$ is charged with at most two events $\xi$,
and each of them satisfies $e(\xi)\ge j(Z)$, hence
$\check{R}_{e(\xi)}\le\check{R}_{j(Z)}$ because $\check{R}_n$ is non-increasing in $n$.

\emph{Collecting.} Inserting \eqref{8.9:eq-age-lemma-correlation-2} into
\eqref{8.9:eq-age-lemma-correlation-1} and using
$\tfrac{2}{L}\cdot90\,|\mathcal{E}|=\tfrac{180}{L}\,|\mathcal{E}|$,
\begin{equation*}
\operatorname{age}^{\mathrm{seed}}_k(C)\le\sum_{\xi\in\mathcal{E}}
\Bigl(5+\check{R}_{e(\xi)}+\frac{180}{L}\Bigr)
+\frac{8}{L}\sum_{Z\in\mathcal{Z}}\bigl(\operatorname{diam}_{j(Z)}(Z)+C_J\bigr).
\end{equation*}
Group the first sum according to the seed to which each event is charged. Each summand
$5+\check{R}_{e(\xi)}+180/L$ is positive, each seed receives at most two summands, and each
summand charged to $Z$ is at most $5+\check{R}_{j(Z)}+180/L$. Hence, using $180/L\le360/L$,
\begin{align*}
\operatorname{age}^{\mathrm{seed}}_k(C)
&\le\sum_{Z\in\mathcal{Z}}\Bigl[2\Bigl(5+\check{R}_{j(Z)}+\frac{180}{L}\Bigr)
+\frac{8}{L}\bigl(\operatorname{diam}_{j(Z)}(Z)+C_J\bigr)\Bigr]\\
&\le\sum_{Z\in\mathcal{Z}}\Bigl[2\Bigl(5+\check{R}_{j(Z)}+\frac{360}{L}\Bigr)
+\frac{8}{L}\bigl(\operatorname{diam}_{j(Z)}(Z)+C_J\bigr)\Bigr]\\
&=\sum_{Z\in\mathcal{Z}}\mathfrak{a}_J(Z),
\end{align*}
which is \eqref{8.9:eq-age-lemma-correlation-a}.
\end{proof}

\begin{remark}
The charge $\mathfrak{a}_J(Z)$ of \eqref{8.9:eq-age-lemma-correlation-a} has the same form as the
charge \eqref{8.9:eq-age-lemma-a} of the age lemma (Lemma~\ref{8.9:lem-age-lemma}, whose proof is
incomplete at one step). Each seed pays the number $\check{R}$ at its birth level plus a constant,
and a term measuring its size at birth. Here that term is linear in the birth diameter, with the
small coefficient $8/L$, whereas there it is logarithmic; either form of the size term is paid by
the reserve inequality for the large-field weights of this chapter. Under the induction invariant
of the correlation theorem's clock, which bounds the large-field exponent of a live region $Z$ at
level $n$ from below by a term proportional to its remaining count $\mathsf{K}_n(Z)$, the
remaining waiting time of a live region is already prepaid in that exponent. Only the past of the
lineage, up to level $k$, enters the argument, and the seeds' charges $\mathfrak{a}_J(Z)$ carry it.
\end{remark}

\begin{definition}[Flat and linear reserve weights per seed]\label{8.9:def-reserve-weights}
Let $C$ be a live component of the live large-field region $\mathcal{O}_k(h)$ of a history $h$ at
a level $k$, let $\mathcal{S}$ be a non-empty set of internal histories
$\mathfrak{h}$ of $C$, and for $\mathfrak{h}\in\mathcal{S}$ let $\mathcal{Z}(\mathfrak{h})$ be
its seed set and $j(Z)$ the birth level of a seed $Z\in\mathcal{Z}(\mathfrak{h})$, as in
Definition~\ref{8.9:def-seed-age}. Let $p_0(\cdot)$ be the degree function, and let $\gamma_0$
and $A_1$ be Ba{\l}aban's positive constants. For a seed $Z$ of the first kind in the sense of
Definition~\ref{8.9:def-seed-age}, born at level $j(Z)$, let $d'_{j(Z)}(Z)\ge0$ be its linear
size, counted in cubes of side $M\check{R}_{j(Z)}$; for preparatory seeds see convention~(a)
below.

We recall the \emph{flat reserve weight} of $\mathcal{S}$ at the parameter $\theta$:
\begin{equation}\label{8.9:eq-reserve-weights-1}
w_\theta(\mathcal{S}):=\sup_{\mathfrak{h}\in\mathcal{S}}
\exp\Bigl(-\theta\sum_{Z\in\mathcal{Z}(\mathfrak{h})}p_0(g_{j(Z)})\Bigr).
\end{equation}
The \emph{linear reserve} of a seed $Z$ at the parameter $\theta'$ is the quantity
\begin{equation*}
\theta'\gamma_0A_1\,p_0(g_{j(Z)})^2\bigl(d'_{j(Z)}(Z)+1\bigr),
\end{equation*}
linear in the size of the seed and quadratic in $p_0$. Its form is that of the member
$\gamma_0A_1\,p_0(g_{j+1})^2\bigl(d'_{j+1}(Z^{j+1}_i)+1\bigr)$ which
\cite[(1.85)]{B-CMP122b} attaches to each component $Z^{j+1}_i$ of a large-field region
created at step $j+1$; in \cite[\S1]{B-CMP122b}, for each such component the large-field
factors yield an exponential factor. The \emph{linear reserve weight} of $\mathcal{S}$, the
splitting of the linear-reserve parameter $\theta'$ into two parts $\theta'_A\ge0$ and
$\theta'_{\mathrm{age}}\ge0$, and the \emph{combined reserve weight} of $\mathcal{S}$ are
\begin{equation}\label{8.9:eq-reserve-weights-a}
\begin{aligned}
w^{\mathrm{lin}}_{\theta'}(\mathcal{S})&:=\sup_{\mathfrak{h}\in\mathcal{S}}
\exp\Bigl(-\theta'\gamma_0A_1\sum_{Z\in\mathcal{Z}(\mathfrak{h})}
p_0(g_{j(Z)})^2\bigl(d'_{j(Z)}(Z)+1\bigr)\Bigr),\\
\theta'&=\theta'_A+\theta'_{\mathrm{age}},\\
w^{\natural}(\mathcal{S})&:=w_\theta(\mathcal{S})\,
w^{\mathrm{lin}}_{\theta'_{\mathrm{age}}}(\mathcal{S})\le w_\theta(\mathcal{S}).
\end{aligned}
\end{equation}
The following conventions accompany these definitions.
\begin{itemize}
\item[(a)] A preparatory seed carries the flat reserve $\theta p_0(g_{j(Z)})$ of
\eqref{8.9:eq-reserve-weights-1}. It has no size member, and in
$w^{\mathrm{lin}}_{\theta'}(\mathcal{S})$ it is given $d'_{j(Z)}(Z):=0$; its per-seed scaffold
charge $\mathfrak{a}(Z)$ has no size term either.
\item[(b)] The splitting of $\theta'$ in \eqref{8.9:eq-reserve-weights-a} is a matter of
bookkeeping only. The part $\theta'_A$ serves the two uses made of the linear reserve by the
absorption proposition for block and boundary terms: the absorption of the
product of its absorption constant with the number of blocks, and the absorption of the
boundary terms. The part $\theta'_{\mathrm{age}}$ serves the age lemma.
\item[(c)] Accordingly, the absorption condition of the absorption proposition for block and
boundary terms is read with $\theta'_A$ in place of $\theta'$:
\begin{equation}\label{8.9:eq-reserve-weights-2}
2C_{\mathrm{abs}}\,d(\mathfrak{g})\,d\cdot O(1)\,M^d\,\check{R}(y)^{d+1}
\le\theta'_A\gamma_0A_1^2\,p_0(y)^2,
\end{equation}
where $C_{\mathrm{abs}}$ is the absorption constant of the absorption proposition for block and
boundary terms and $d(\mathfrak{g})$ is the constant so denoted in that statement; here $y$
stands for an inverse-square coupling $g'^{-2}$, and $p_0(y)$ and $\check{R}(y)$ denote the
degree function and Ba{\l}aban's number $R$ evaluated at the coupling $g'$ with $g'^{-2}=y$.
The condition \eqref{8.9:eq-reserve-weights-2} is of the same kind as with $\theta'$: in
$\log y$ its left side has degree $(d+1)r$, where $r$ is the exponent of $\check{R}(y)$, against
degree $2p_0$ on the right, and the ratio of the left side to the right side decreases to zero
as $y\to\infty$.
\end{itemize}
\end{definition}

\begin{proof}
We prove the inequality $w^{\natural}(\mathcal{S})\le w_\theta(\mathcal{S})$ in
\eqref{8.9:eq-reserve-weights-a}. Fix $\mathfrak{h}\in\mathcal{S}$ and
$Z\in\mathcal{Z}(\mathfrak{h})$. We have $\theta'_{\mathrm{age}}\ge0$, the constants
$\gamma_0$ and $A_1$ are positive, $p_0(g_{j(Z)})^2\ge0$, and
$d'_{j(Z)}(Z)+1\ge1$, since $d'_{j(Z)}(Z)$ is a number of cubes for a seed of the first kind
and equals $0$ for a preparatory seed by convention~(a). Hence
\begin{equation*}
\theta'_{\mathrm{age}}\gamma_0A_1\,p_0(g_{j(Z)})^2\bigl(d'_{j(Z)}(Z)+1\bigr)\ge0 .
\end{equation*}
Summing over $Z\in\mathcal{Z}(\mathfrak{h})$, the exponent in the definition of
$w^{\mathrm{lin}}_{\theta'_{\mathrm{age}}}(\mathcal{S})$ is at most $0$ for every
$\mathfrak{h}\in\mathcal{S}$, so each exponential is at most $1$, and since $\mathcal{S}$ is
non-empty the supremum satisfies
$0<w^{\mathrm{lin}}_{\theta'_{\mathrm{age}}}(\mathcal{S})\le1$. The weight
$w_\theta(\mathcal{S})$ is a supremum of positive numbers, hence
$w_\theta(\mathcal{S})>0$. Multiplying the inequality
$w^{\mathrm{lin}}_{\theta'_{\mathrm{age}}}(\mathcal{S})\le1$ by $w_\theta(\mathcal{S})$ gives
$w^{\natural}(\mathcal{S})\le w_\theta(\mathcal{S})$.
\end{proof}

\begin{lemma}[The age-payment inequality under a coupling ceiling]\label{8.9:lem-age-payment}
Let $b\ge 1$ be fixed and kept symbolic. Let $\theta$ and $\theta'_{\mathrm{age}}$ be the positive
reserve fractions of the weights $w_\theta$ and $w^{\mathrm{lin}}_{\theta'_{\mathrm{age}}}$ of
Definition~\ref{8.9:def-reserve-weights}. Let $A_0$, $A_1$, $\gamma_0$ be Ba{\l}aban's constants,
and let $p_0$ be the exponent of the degree function $p_0(\cdot)$, so that
$p_0(g')=A_0(\log g'^{-2})^{p_0}$. Let $n_1$ and $C_S$ be the constants of the age lemma
(Lemma~\ref{8.9:lem-age-lemma}). Let $R(y)$ be Ba{\l}aban's number $R$ written as a function of
$y=g'^{-2}$, so that $\check{R}_j=R(x_j)$ at every level $j$, and let $r$ be the exponent in the
bound
\begin{equation}\label{8.9:eq-age-payment-1}
R(y)\le L(\log y)^{r}
\end{equation}
of \cite{B-CMP119}. Put $X=g^{-2}$.
\begin{itemize}
\item[(a)] Consider the condition that, for every $y\ge X$,
\begin{equation}\label{8.9:eq-age-payment-a}
\begin{gathered}
(b+1)\ln L\,\bigl(2R(y)+n_1+C_S+4\bigr)\le \theta A_0(\log y)^{p_0},\\
2(b+1)\le \theta'_{\mathrm{age}}\gamma_0A_1A_0^2(\log y)^{2p_0}.
\end{gathered}
\end{equation}
There is a number
\begin{equation*}
\gamma_{\mathrm{age}}=\gamma_{\mathrm{age}}(b,L,\theta,\theta'_{\mathrm{age}},A_0,A_1,\gamma_0,
p_0,r,n_1,C_S)>0
\end{equation*}
such that \eqref{8.9:eq-age-payment-a} holds whenever $g\le\gamma_{\mathrm{age}}$. This
condition is one ceiling on $g$. It is the same for every $s$, $K$ and $m$, it is uniform in the
level, and it is not a cap in the sense of Definition~\ref{2.3:def-cap}. In the order of choice of
Notations~\ref{2.3:not-stages-first}--\ref{2.3:not-stages-last} it is a ceiling on $g$ read after
$L$ and after the constants on which $\gamma_{\mathrm{age}}$ depends; its threshold depends on $L$,
which that order allows.
\item[(b)] Assume \eqref{8.9:eq-age-payment-a}. Let $Z$ be a seed at any level, and write
$d'(Z)=d'_{j(Z)}(Z)$, with $d'(Z)=0$ for the seeds that carry no size, as in
Definition~\ref{8.9:def-reserve-weights}. Then
\begin{equation}\label{8.9:eq-age-payment-b}
\theta p_0(g_{j(Z)})+\theta'_{\mathrm{age}}\gamma_0A_1\,p_0(g_{j(Z)})^2\bigl(d'(Z)+1\bigr)
\ge (b+1)\ln L\cdot\mathfrak{a}(Z),
\end{equation}
where $\mathfrak{a}(Z)$ is the charge \eqref{8.9:eq-age-lemma-a}.
\end{itemize}
\end{lemma}

\begin{proof}
\emph{Part (a).} By \eqref{8.9:eq-age-payment-1}, the first inequality of
\eqref{8.9:eq-age-payment-a} follows from
\begin{equation*}
(b+1)\ln L\,\bigl(2L(\log y)^{r}+n_1+C_S+4\bigr)\le\theta A_0(\log y)^{p_0}.
\end{equation*}
As a function of $\log y$, its left side has degree $r$ and its right side has degree $p_0$. The
exponents $r$ and $p_0$ are positive, and the exponent of the degree function satisfies the
one-sided condition $p_0\ge 5r$ of the order of choice
(Notations~\ref{2.3:not-stages-first}--\ref{2.3:not-stages-last}); hence $p_0>r$. Therefore, as
$y\to\infty$, the quotient of the left side by $\theta A_0(\log y)^{p_0}$ tends to $0$.

In the second inequality of \eqref{8.9:eq-age-payment-a}, the left side has degree $0$ in
$\log y$ and the right side has degree $2p_0>0$. Hence, as $y\to\infty$, the quotient of $2(b+1)$
by $\theta'_{\mathrm{age}}\gamma_0A_1A_0^2(\log y)^{2p_0}$ tends to $0$.

So there is $y_1>1$ such that both quotients are at most $1$ for every $y\ge y_1$. This $y_1$
depends only on $b$, $L$, $\theta$, $\theta'_{\mathrm{age}}$, $A_0$, $A_1$, $\gamma_0$, $p_0$,
$r$, $n_1$ and $C_S$. Put $\gamma_{\mathrm{age}}=y_1^{-1/2}$. If $g\le\gamma_{\mathrm{age}}$, then
$X=g^{-2}\ge y_1$, so every $y\ge X$ satisfies $y\ge y_1$, and both inequalities of
\eqref{8.9:eq-age-payment-a} hold.

The condition reads only $g$, through the range $y\ge X$, together with the constants listed. It
therefore involves neither $s$, $K$, $m$ nor a level, and it bounds no constant other than $g$;
in particular it is not a cap in the sense of Definition~\ref{2.3:def-cap}. Since
$\gamma_{\mathrm{age}}$ depends only on $L$ and the other constants listed, the ceiling is read
after them in the order of choice.

\emph{Part (b).} Let $Z$ be a seed and put $y=x_{j(Z)}$. By the lower bound for the running
inverse couplings of Definition~\ref{1.2:def-couplings}, $x_j\ge X$ at every level $j$; hence
$y\ge X$ and \eqref{8.9:eq-age-payment-a} applies at this $y$.
By the definition of the degree function,
\begin{equation*}
p_0(g_{j(Z)})=A_0(\log x_{j(Z)})^{p_0}=A_0(\log y)^{p_0},\qquad
p_0(g_{j(Z)})^2=A_0^2(\log y)^{2p_0}.
\end{equation*}

For the first member of the left side of \eqref{8.9:eq-age-payment-b}, the first inequality of
\eqref{8.9:eq-age-payment-a} gives
\begin{equation}\label{8.9:eq-age-payment-2}
\theta p_0(g_{j(Z)})\ge(b+1)\ln L\,\bigl(2R(y)+n_1+C_S+4\bigr).
\end{equation}

For the second member, we use the second inequality of \eqref{8.9:eq-age-payment-a}, then
$\ln u\le u$ for $u\ge 1$, then $\ln(d'(Z)+1)=\ln L\cdot\log_L(d'(Z)+1)$:
\begin{equation}\label{8.9:eq-age-payment-3}
\begin{aligned}
\theta'_{\mathrm{age}}\gamma_0A_1\,p_0(g_{j(Z)})^2\bigl(d'(Z)+1\bigr)
&=\theta'_{\mathrm{age}}\gamma_0A_1A_0^2(\log y)^{2p_0}\bigl(d'(Z)+1\bigr)\\
&\ge 2(b+1)\bigl(d'(Z)+1\bigr)\\
&\ge 2(b+1)\ln\bigl(d'(Z)+1\bigr)\\
&=(b+1)\ln L\cdot 2\log_L\bigl(d'(Z)+1\bigr).
\end{aligned}
\end{equation}

Adding \eqref{8.9:eq-age-payment-2} and \eqref{8.9:eq-age-payment-3} gives
\begin{equation*}
\begin{split}
\theta p_0(g_{j(Z)})&+\theta'_{\mathrm{age}}\gamma_0A_1\,p_0(g_{j(Z)})^2\bigl(d'(Z)+1\bigr)\\
&\ge(b+1)\ln L\,\Bigl(2\log_L\bigl(d'(Z)+1\bigr)+2R(y)+n_1+C_S+4\Bigr).
\end{split}
\end{equation*}

It remains to bound the charge. By the age lemma (Lemma~\ref{8.9:lem-age-lemma}), whose proof is
incomplete at one step, a seed of the first kind satisfies
\begin{equation*}
\mathfrak{a}(Z)\le 2\log_L\bigl(d'_{j(Z)}(Z)+1\bigr)+2\check{R}_{j(Z)}+n_1+C_S+4.
\end{equation*}
For the remaining seeds, the bound of that lemma has no size term, and they are given
$d'(Z)=0$, as in the statement. Since $\check{R}_{j(Z)}=R(x_{j(Z)})=R(y)$, the parenthesis on the
right of the last display is at least $\mathfrak{a}(Z)$. This proves \eqref{8.9:eq-age-payment-b}.
\end{proof}

\begin{remark}
The two members of the left side of \eqref{8.9:eq-age-payment-b} are the per-seed reserves of
$w_\theta$ and of $w^{\mathrm{lin}}_{\theta'_{\mathrm{age}}}$ in
Definition~\ref{8.9:def-reserve-weights}.

For the charge $\mathfrak{a}_J(Z)$ of \eqref{8.9:eq-age-lemma-correlation-a}, used by the age lemma
on the correlation theorem's clock (Lemma~\ref{8.9:lem-age-lemma-correlation}), the ceiling is of
the same kind, with two replacements in \eqref{8.9:eq-age-payment-a}. The member
$2R(y)+n_1+C_S+4$ is replaced by the member of $\mathfrak{a}_J(Z)$ that carries $\check{R}$. The
number $2(b+1)$ is replaced by $(b+1)\cdot 8\ln L/L$ times the factor that multiplies $8/L$ in the
size member of \eqref{8.9:eq-age-lemma-correlation-a}.
\end{remark}

\begin{theorem}[Decay in age of live large-field regions]\label{8.9:thm-live-age-decay}
Fix a number $b\ge 1$, a reserve parameter $\beta_0\ge0$, a reserve parameter $\theta$ in the window
$\mathopen{]}0,(1-3\beta_0)/2\mathclose{[}$, and a level $k_0$. We use the following objects.
\begin{itemize}
\item The internal history $\mathfrak{h}$ of a live component, its seed set
$\mathcal{Z}(\mathfrak{h})$, the birth level $j(Z)$ of a seed $Z$, the birth level
$j_{\mathrm{seed}}(C)$ of the oldest seed, the seed age
$\operatorname{age}^{\mathrm{seed}}_k(C)$ and the territory age, all as in
Definition~\ref{8.9:def-seed-age}.
\item The reserve weights $w_\theta(\mathcal{S})$, $w^{\mathrm{lin}}_{\theta'}(\mathcal{S})$,
$w^{\natural}(\mathcal{S})$ and the split $\theta'=\theta'_A+\theta'_{\mathrm{age}}$ of
Definition~\ref{8.9:def-reserve-weights}, display~\eqref{8.9:eq-reserve-weights-a}.
\item The per-seed charge $\mathfrak{a}(Z)$ of display~\eqref{8.9:eq-age-lemma-a}.
\item The events $\mathfrak{lv}^{=a}_k(\square)$, the probabilities $q^{=a}_k$, the live rate
$q^{\mathrm{lv}}_k$ with its constant $C_L$, and the lifetime
$\mathfrak{l}=L\check{R}_{k_0}+c_L$, all of display~\eqref{8.9:eq-seed-age-a}.
\item The event $\mathfrak{lv}^{\ge a,\mathrm{seed}}_k(\square)$: the level-$k$ cell $\square$
lies in a live component of the live region $\mathcal{O}_k(h)$ whose seed age is at least $a$.
\end{itemize}
A \emph{touch-connected} union of level-$k$ cells is a union of level-$k$ cells that is connected
when two cells are counted as adjacent whenever their closures meet.
For a touch-connected union $C$ of level-$k$ cells and a set $\mathcal{S}$ of internal histories
of $C$, $L_{C,\mathcal{S}}$ denotes the lineage event that
$C$ is a live component of $\mathcal{O}_k(h)$ whose internal history lies in $\mathcal{S}$, and
$L_C$ denotes the lineage event of $C$ in which every internal history of $C$ is admitted. The
annealed normalised large-field measure on histories is $\tilde{\nu}$; it has total mass one.
Assume the following.
\begin{enumerate}
\item[(a)] For every $K$, every admissible run with terminal level at least $k_0$, every
touch-connected union $C$ of level-$k_0$ cells and every set $\mathcal{S}$ of internal histories
of $C$,
\begin{equation}\label{8.9:eq-live-age-decay-a}
\begin{gathered}
\tilde{\nu}(L_{C,\mathcal{S}})\le \varepsilon^{\natural}_{k_0}(C,\mathcal{S})\,
\bigl(1-\tilde{\nu}(L_C)\bigr),\\
\varepsilon^{\natural}_k(C,\mathcal{S}):=w^{\natural}(\mathcal{S})\,
\exp\Bigl(-\frac{2-\theta}{1+\beta_0}\,p_0(g_k)-\kappa_1 d_k(C)
+C_{\mathrm{rm}}N(\operatorname{hull}C)\Bigr).
\end{gathered}
\end{equation}
Here $d_k(C)$ and $N(\operatorname{hull}C)$ are the size and hull functionals of $C$ at level $k$,
$\kappa_1>0$ is the auxiliary parameter fixed in the order of choice
(Definition~\ref{2.3:not-stages-middle}), proportional to $M^d$, and $C_{\mathrm{rm}}\ge0$ is a
constant. Hypothesis~(a) is the case in which $B$ is the sure event and $k=k_0$ of the following
statement at level $k$, which carries the reserve weight $w^{\natural}$ in place of the flat
reserve weight $w_\theta$: for every $K$, every admissible run with terminal level at least $k$,
every touch-connected union $C$ of level-$k$ cells, every set $\mathcal{S}$ of internal histories
of $C$ and every event $B$ of levels at most $k$ that is blind to the lineage of $C$,
\begin{equation*}
\tilde{\nu}(L_{C,\mathcal{S}}\cap B)\le
\varepsilon^{\natural}_k(C,\mathcal{S})\,\tilde{\nu}(B\setminus L_C).
\end{equation*}
The sure event is trivially such an event for $k=k_0$, and for it
$\tilde{\nu}(B\setminus L_C)=1-\tilde{\nu}(L_C)$.
\item[(b)] The age lemma, Lemma~\ref{8.9:lem-age-lemma} (stated; proof incomplete at the step
marked $(\ast)$).
In particular, for every $K$, every admissible run, every level $k$ and every component $C$ that
is live at level $k$ with internal history $\mathfrak{h}$,
\begin{equation*}
\operatorname{age}^{\mathrm{seed}}_k(C)\le \sum_{Z\in\mathcal{Z}(\mathfrak{h})}\mathfrak{a}(Z).
\end{equation*}
\item[(c)] The coupling ceiling~\eqref{8.9:eq-age-payment-a} for this $b$, and the lower bound
$x_j\ge g^{-2}$ at every level $j$ on the running inverse couplings $x_j=g_j^{-2}$ of
Definition~\ref{1.2:def-couplings}.
\item[(d)] The largeness condition on $\kappa_1$ of the shape sum, in the form
\begin{equation}\label{8.9:eq-live-age-decay-1}
\sum_{n\ge1}9^{dn}\,e^{-(\kappa_1 5^{-d}-C_{\mathrm{rm}}3^d)\,n\,\check{R}_{k_0}^d}\le 2 .
\end{equation}
\item[(e)] The first of the two conditions on the flow of the running couplings $g_k$ of
Definition~\ref{1.2:def-couplings}.
\end{enumerate}
Let $K$ be arbitrary. Take any admissible run with terminal level at least $k_0$, any
$N_T\ge1$, any level-$k_0$ cell $\square$ and any $a\ge1$. Then the following hold.
\begin{enumerate}
\item[(i)] With $C_{\mathrm{old}}:=2e^{\kappa_1}$,
\begin{equation}\label{8.9:eq-live-age-decay-b}
q^{\ge a}_{k_0}(\square):=\tilde{\nu}\bigl(\mathfrak{lv}^{\ge a,\mathrm{seed}}_{k_0}(\square)\bigr)
\le C_{\mathrm{old}}\,L^{-(b+1)a}\,
e^{-(2-\theta)(1+\beta_0)^{-1}p_0(g_{k_0})}.
\end{equation}
This holds uniformly in $K$, $m$, $N_T$ and the position of $\square$. A fortiori the same bound
holds for $q^{=a}_{k_0}=\max_{\square}\tilde{\nu}(\mathfrak{lv}^{=a}_{k_0}(\square))$.
\item[(ii)] Suppose (a)--(e) hold with $b=8$. Then the age-weighted live-cell bound of
display~\eqref{8.9:eq-seed-age-a} holds at level $k_0$ in the following form:
\begin{equation}\label{8.9:eq-live-age-decay-2}
\begin{aligned}
\sum_{a>\mathfrak{l}}L^{8(a+1)}q^{=a}_{k_0}
&\le C_{\mathrm{old}}(L-1)^{-1}L^{8-\mathfrak{l}}\,e^{-(2-\theta)(1+\beta_0)^{-1}p_0(g_{k_0})}\\
&\le A_{\mathrm{lv}}\,L^{8(\mathfrak{l}+1)}\,q^{\mathrm{lv}}_{k_0},
\qquad A_{\mathrm{lv}}:=\frac{C_{\mathrm{old}}}{(L-1)C_L}.
\end{aligned}
\end{equation}
The constant $A_{\mathrm{lv}}$ depends on $L$ through the factor $(L-1)^{-1}$, on $(d,M)$
through $\kappa_1$, which is proportional to $M^d$, and on the constants entering $C_L$.
Suppose instead (a)--(e) hold with $b=4$. Then for every index $i$ and every number $n_i\ge0$
attached to it,
\begin{equation}\label{8.9:eq-live-age-decay-3}
\sum_{a>\mathfrak{l}}L^{4(a+1)}n_i\,q^{=a}_{k_0}
\le n_i\,C_{\mathrm{old}}(L-1)^{-1}L^{4-\mathfrak{l}}\,
e^{-(2-\theta)(1+\beta_0)^{-1}p_0(g_{k_0})},
\end{equation}
a bound that carries the factor $L^{-\mathfrak{l}}$.
\end{enumerate}
Now consider the following second form of hypothesis~(a): for every $K$, every admissible run
with terminal level at least $k_0+c_0$, every touch-connected union $C$ of level-$k_0$ cells and
every set $\mathcal{S}$ of internal histories of $C$, the bound of~(a) holds with the following
replacements made in the definition of $\varepsilon^{\natural}_k$
in~\eqref{8.9:eq-live-age-decay-a}, read at $k=k_0$:
\begin{equation*}
\begin{aligned}
N(\operatorname{hull}C)&\mapsto C_{\mathrm{zone}}3^d n(C)\check{R}_k^d,\\
\kappa_1 d_k(C)&\mapsto \kappa_1 d_k(C)-a_{\mathrm{cl}}n(C)-\log 2,\\
w^{\natural}&\mapsto 2e^{\kappa_1}w^{\natural}.
\end{aligned}
\end{equation*}
Here $n(C)$ is the number of level-$k$ cells of $C$, and $C_{\mathrm{zone}}$, $a_{\mathrm{cl}}$
are fixed constants, the same here and in the condition~\eqref{8.9:eq-live-age-decay-10} below.
Assume this second form in place of (a), together with the coupling ceiling
\begin{equation}\label{8.9:eq-live-age-decay-11}
g\le\gamma_{\mathrm{cl}}
\end{equation}
on the reference coupling $g$, where $\gamma_{\mathrm{cl}}$ is a fixed constant; the proof below
uses~\eqref{8.9:eq-live-age-decay-11} only through the second form of hypothesis~(a). Assume
moreover the depth condition that the cutoff $K$ admits an admissible
run with terminal level at least $k_0+c_0$; the conclusion below is asserted for every such $K$.
Assume in place of (d) the strengthened largeness condition on $\kappa_1$, and that, for every
level-$k_0$ cell $\square$,
\begin{equation}\label{8.9:eq-live-age-decay-10}
\sum_{C\ni\square}2e^{\kappa_1}\cdot2\,
e^{-\kappa_1 d_{k_0}(C)+a_{\mathrm{cl}}n(C)
+C_{\mathrm{rm}}C_{\mathrm{zone}}3^dn(C)\check{R}_{k_0}^d}\le 4e^{2\kappa_1},
\end{equation}
the sum running over touch-connected unions $C$ of level-$k_0$ cells containing $\square$. Then,
for every such $K$, (i) and (ii) hold, in every admissible run with terminal level at least $k_0$,
with $C_{\mathrm{old}}:=4e^{2\kappa_1}$.
\end{theorem}

\begin{proof}
\emph{Proof of (i).} Fix $\square$ and $a\ge1$. For a touch-connected union $C$ of level-$k_0$
cells put
\begin{equation*}
\mathcal{S}^{(\ge a)}_C:=\Bigl\{\mathfrak{h}\ \text{an internal history of}\ C:\ 
\min_{Z\in\mathcal{Z}(\mathfrak{h})}j(Z)\le k_0-a\Bigr\}.
\end{equation*}

\emph{The cover.} Let $h\in\mathfrak{lv}^{\ge a,\mathrm{seed}}_{k_0}(\square)$. The support
$|\mathcal{O}_{k_0}(h)|$ of the live region $\mathcal{O}_{k_0}(h)$ is a union of level-$k_0$
cells. Hence the live component of
$\mathcal{O}_{k_0}(h)$ containing $\square$ is a touch-connected union $C$ of level-$k_0$ cells
containing $\square$. Its seed age is at least $a$, that is,
$j_{\mathrm{seed}}(C)=\min_{Z\in\mathcal{Z}(\mathfrak{h})}j(Z)\le k_0-a$ for its internal
history $\mathfrak{h}$ (Definition~\ref{8.9:def-seed-age}). Thus
$h\in L_{C,\mathcal{S}^{(\ge a)}_C}$, and
\begin{equation}\label{8.9:eq-live-age-decay-4}
\mathfrak{lv}^{\ge a,\mathrm{seed}}_{k_0}(\square)\subset
\bigcup_{C\ni\square}L_{C,\mathcal{S}^{(\ge a)}_C},
\end{equation}
the union running over touch-connected unions $C$ of level-$k_0$ cells containing $\square$.

\emph{Application of hypothesis~(a).} Apply~\eqref{8.9:eq-live-age-decay-a} with
$\mathcal{S}=\mathcal{S}^{(\ge a)}_C$. Since $\tilde{\nu}$ is normalised,
$1-\tilde{\nu}(L_C)\le1$. Hence
\begin{equation}\label{8.9:eq-live-age-decay-5}
\tilde{\nu}\bigl(L_{C,\mathcal{S}^{(\ge a)}_C}\bigr)\le
\varepsilon^{\natural}_{k_0}\bigl(C,\mathcal{S}^{(\ge a)}_C\bigr).
\end{equation}

\emph{The age factor.} Let $\mathfrak{h}\in\mathcal{S}^{(\ge a)}_C$. Then the component $C$ is
live at level $k_0$ with $\operatorname{age}^{\mathrm{seed}}_{k_0}(C)\ge a$. By the age lemma
(hypothesis~(b); Lemma~\ref{8.9:lem-age-lemma}, stated; proof incomplete at the step marked
$(\ast)$),
\begin{equation*}
\sum_{Z\in\mathcal{Z}(\mathfrak{h})}\mathfrak{a}(Z)\ge
\operatorname{age}^{\mathrm{seed}}_{k_0}(C)\ge a .
\end{equation*}
By hypothesis~(c), $x_{j(Z)}\ge g^{-2}$ for every seed $Z$, so that
Lemma~\ref{8.9:lem-age-payment} applies at $y=x_{j(Z)}$; its conclusion is the age-payment
inequality~\eqref{8.9:eq-age-payment-b} with $d'(Z)=d'_{j(Z)}(Z)$. Summing it over
$Z\in\mathcal{Z}(\mathfrak{h})$ gives
\begin{equation}\label{8.9:eq-live-age-decay-6}
\begin{aligned}
&\theta\sum_{Z}p_0(g_{j(Z)})
+\theta'_{\mathrm{age}}\gamma_0A_1\sum_{Z}p_0(g_{j(Z)})^2\bigl(d'(Z)+1\bigr)\\
&\qquad\ge(b+1)\ln L\sum_{Z}\mathfrak{a}(Z)\ge(b+1)\,a\ln L .
\end{aligned}
\end{equation}
For each $\mathfrak{h}\in\mathcal{S}^{(\ge a)}_C$, consider the product of the exponential whose
supremum defines $w_\theta$ and the exponential whose supremum defines
$w^{\mathrm{lin}}_{\theta'_{\mathrm{age}}}$. By~\eqref{8.9:eq-live-age-decay-6} this product is at
most $L^{-(b+1)a}$ for every $\mathfrak{h}\in\mathcal{S}^{(\ge a)}_C$. Taking the supremum over
$\mathfrak{h}\in\mathcal{S}^{(\ge a)}_C$, with $w^{\natural}$ as in
Definition~\ref{8.9:def-reserve-weights}, gives
\begin{equation}\label{8.9:eq-live-age-decay-7}
w^{\natural}\bigl(\mathcal{S}^{(\ge a)}_C\bigr)\le L^{-(b+1)a}.
\end{equation}

\emph{The shape sum.} Group the unions $C\ni\square$ by their number $n\ge1$ of cells. We use
the counting bounds of the shape sum of the large-field block bound: for every $n\ge1$ the number
of touch-connected unions of $n$ level-$k_0$ cells containing a given cell is at most $9^{dn}$,
and every such $C$ satisfies
\begin{equation*}
d_{k_0}(C)\ge 5^{-d}n\check{R}_{k_0}^d-1,\qquad
N(\operatorname{hull}C)\le 3^d n\check{R}_{k_0}^d .
\end{equation*}
Hence for each of them
\begin{equation*}
-\kappa_1 d_{k_0}(C)\le -\kappa_1\bigl(5^{-d}n\check{R}_{k_0}^d-1\bigr),\qquad
C_{\mathrm{rm}}N(\operatorname{hull}C)\le C_{\mathrm{rm}}3^dn\check{R}_{k_0}^d .
\end{equation*}
The first inequality uses $\kappa_1>0$ and the second $C_{\mathrm{rm}}\ge0$.
Hence, by~\eqref{8.9:eq-live-age-decay-1},
\begin{equation}\label{8.9:eq-live-age-decay-8}
\begin{aligned}
\sum_{C\ni\square}e^{-\kappa_1 d_{k_0}(C)+C_{\mathrm{rm}}N(\operatorname{hull}C)}
&\le\sum_{n\ge1}9^{dn}\,
e^{-\kappa_1(5^{-d}n\check{R}_{k_0}^d-1)+C_{\mathrm{rm}}3^dn\check{R}_{k_0}^d}\\
&=e^{\kappa_1}\sum_{n\ge1}9^{dn}\,
e^{-(\kappa_1 5^{-d}-C_{\mathrm{rm}}3^d)\,n\,\check{R}_{k_0}^d}
\le 2e^{\kappa_1}.
\end{aligned}
\end{equation}

\emph{Collection.} Combine the cover~\eqref{8.9:eq-live-age-decay-4}, subadditivity of
$\tilde{\nu}$, the bound~\eqref{8.9:eq-live-age-decay-5}, the definition of
$\varepsilon^{\natural}_{k_0}$ in~\eqref{8.9:eq-live-age-decay-a},
the weight bound~\eqref{8.9:eq-live-age-decay-7} and the shape sum~\eqref{8.9:eq-live-age-decay-8}.
This gives
\begin{align*}
q^{\ge a}_{k_0}(\square)
&\le\sum_{C\ni\square}\varepsilon^{\natural}_{k_0}\bigl(C,\mathcal{S}^{(\ge a)}_C\bigr)\\
&\le L^{-(b+1)a}\,e^{-(2-\theta)(1+\beta_0)^{-1}p_0(g_{k_0})}
\sum_{C\ni\square}e^{-\kappa_1 d_{k_0}(C)+C_{\mathrm{rm}}N(\operatorname{hull}C)}\\
&\le L^{-(b+1)a}\,e^{-(2-\theta)(1+\beta_0)^{-1}p_0(g_{k_0})}\cdot 2e^{\kappa_1},
\end{align*}
which is~\eqref{8.9:eq-live-age-decay-b} with $C_{\mathrm{old}}=2e^{\kappa_1}$.

\emph{Uniformity.} Every factor on the right of the last display is free of $K$, $m$, $N_T$ and
the position of $\square$: hypothesis~(a) holds with the same rate for every $K$ and every
admissible run, $\kappa_1$ and $C_{\mathrm{rm}}$ are constants, and the remaining factors are
$L^{-(b+1)a}$ and the bound~\eqref{8.9:eq-live-age-decay-1}.

\emph{Territory age.} By Definition~\ref{8.9:def-seed-age} the territory age of a live component
does not exceed its seed age. Hence
$\mathfrak{lv}^{=a}_{k_0}(\square)\subset\mathfrak{lv}^{\ge a,\mathrm{seed}}_{k_0}(\square)$ for
every $\square$. Taking the maximum over $\square$, the bound~\eqref{8.9:eq-live-age-decay-b}
holds for $q^{=a}_{k_0}$.

\emph{Proof of (ii).} First compare the exponential rates. Since $\beta_0\ge0$ and $\theta$ lies
in the window $\mathopen{]}0,(1-3\beta_0)/2\mathclose{[}$, we have
$(2-\theta)/(1+\beta_0)>3/2$: as $1+\beta_0>0$, this is equivalent to $4-2\theta>3+3\beta_0$, that
is, to $2\theta<1-3\beta_0$. Also
\begin{equation*}
\theta<\frac{1-3\beta_0}{2}\le\frac12<\frac32 .
\end{equation*}
Because $p_0(g_{k_0})\ge0$ and, by display~\eqref{8.9:eq-seed-age-a},
$q^{\mathrm{lv}}_{k_0}=C_Le^{-\theta p_0(g_{k_0})}$, it follows that
\begin{equation}\label{8.9:eq-live-age-decay-9}
e^{-(2-\theta)(1+\beta_0)^{-1}p_0(g_{k_0})}\le e^{-\theta p_0(g_{k_0})}
=q^{\mathrm{lv}}_{k_0}/C_L .
\end{equation}
Next, with $b=8$, part~(i) gives
\begin{equation*}
q^{=a}_{k_0}\le C_{\mathrm{old}}L^{-9a}e^{-(2-\theta)(1+\beta_0)^{-1}p_0(g_{k_0})}.
\end{equation*}
Summing over $a\ge\mathfrak{l}+1$ requires the geometric series
\begin{equation*}
\sum_{a\ge\mathfrak{l}+1}L^{8(a+1)}L^{-9a}
=L^8\sum_{a\ge\mathfrak{l}+1}L^{-a}
=L^{8-\mathfrak{l}}(L-1)^{-1}.
\end{equation*}
This yields the first inequality of~\eqref{8.9:eq-live-age-decay-2}. For the second, write
$L^{8-\mathfrak{l}}=L^{-9\mathfrak{l}}\,L^{8(\mathfrak{l}+1)}$ and apply~\eqref{8.9:eq-live-age-decay-9}:
\begin{equation*}
\begin{aligned}
C_{\mathrm{old}}(L-1)^{-1}L^{8-\mathfrak{l}}\,e^{-(2-\theta)(1+\beta_0)^{-1}p_0(g_{k_0})}
&\le C_{\mathrm{old}}(L-1)^{-1}L^{-9\mathfrak{l}}\,L^{8(\mathfrak{l}+1)}\,
q^{\mathrm{lv}}_{k_0}/C_L .
\end{aligned}
\end{equation*}
Since $L^{-9\mathfrak{l}}\le1$, the right-hand side is at most
$A_{\mathrm{lv}}L^{8(\mathfrak{l}+1)}q^{\mathrm{lv}}_{k_0}$, with
$A_{\mathrm{lv}}=C_{\mathrm{old}}/((L-1)C_L)$. This proves~\eqref{8.9:eq-live-age-decay-2}.

With $b=4$ the argument is the same. Part~(i) gives
\begin{equation*}
q^{=a}_{k_0}\le C_{\mathrm{old}}L^{-5a}e^{-(2-\theta)(1+\beta_0)^{-1}p_0(g_{k_0})},
\end{equation*}
and the geometric series is now
\begin{equation*}
\sum_{a>\mathfrak{l}}L^{4(a+1)}L^{-5a}=L^{4-\mathfrak{l}}(L-1)^{-1}.
\end{equation*}
Multiplying by $n_i\ge0$ gives~\eqref{8.9:eq-live-age-decay-3}.

\emph{The second form of hypothesis~(a).} By the depth condition there is an admissible run with
terminal level at least $k_0+c_0$; fix one. The cover~\eqref{8.9:eq-live-age-decay-4}, the bound
$1-\tilde{\nu}(L_C)\le1$ and the weight bound~\eqref{8.9:eq-live-age-decay-7} are
unchanged, and the second form of hypothesis~(a) bounds
$\tilde{\nu}(L_{C,\mathcal{S}^{(\ge a)}_C})$ by its rate in place
of~\eqref{8.9:eq-live-age-decay-5}. In that rate
\begin{equation*}
e^{-(\kappa_1 d_{k_0}(C)-a_{\mathrm{cl}}n(C)-\log2)}
=2\,e^{-\kappa_1 d_{k_0}(C)+a_{\mathrm{cl}}n(C)},
\end{equation*}
and the factors $e^{a_{\mathrm{cl}}n(C)}$ so produced are absorbed in the
condition~\eqref{8.9:eq-live-age-decay-10}. Hence, by~\eqref{8.9:eq-live-age-decay-7}
and~\eqref{8.9:eq-live-age-decay-10},
\begin{align*}
q^{\ge a}_{k_0}(\square)
&\le L^{-(b+1)a}\,e^{-(2-\theta)(1+\beta_0)^{-1}p_0(g_{k_0})}\\
&\qquad\times\sum_{C\ni\square}2e^{\kappa_1}\cdot2\,
e^{-\kappa_1 d_{k_0}(C)+a_{\mathrm{cl}}n(C)
+C_{\mathrm{rm}}C_{\mathrm{zone}}3^dn(C)\check{R}_{k_0}^d}\\
&\le L^{-(b+1)a}\,e^{-(2-\theta)(1+\beta_0)^{-1}p_0(g_{k_0})}\cdot4e^{2\kappa_1}.
\end{align*}
This is (i) with $C_{\mathrm{old}}=4e^{2\kappa_1}$ in the fixed run, whose terminal level is at
least $k_0+c_0$. The event $\mathfrak{lv}^{\ge a,\mathrm{seed}}_{k_0}(\square)$ is of levels
at most $k_0$, and so is $\mathfrak{lv}^{=a}_{k_0}(\square)$. By the run-independence of
low-level events (the $\tilde{\nu}$-mass of every event of levels at most $k_0$ is the same in
every admissible run with terminal level at least $k_0$), the bound obtained in the fixed run
holds in every admissible run with terminal level at least $k_0$; the bound for
$q^{=a}_{k_0}$ then follows as before. The proof of (ii) uses only (i), the
comparison~\eqref{8.9:eq-live-age-decay-9} and the geometric series. Hence (ii) holds with
$C_{\mathrm{old}}=4e^{2\kappa_1}$ and $A_{\mathrm{lv}}=C_{\mathrm{old}}/((L-1)C_L)$ for this
$C_{\mathrm{old}}$.
\end{proof}

\begin{remark}[Remarks on the allocation of the reserve]\label{8.9:rem-reserve-allocation}
We use the notation of Definition~\ref{8.9:def-reserve-weights}, Lemma~\ref{8.9:lem-age-payment}
and Theorem~\ref{8.9:thm-live-age-decay}. Throughout, $b\ge 1$ is the fixed, otherwise arbitrary,
exponent of the age-payment inequality \eqref{8.9:eq-age-payment-a}, and $\mathfrak{l}$ is the lifetime
entering part~(ii) of Theorem~\ref{8.9:thm-live-age-decay}. For a touch-connected union $C$ of
level-$k_0$ cells and $a\ge1$, let $\mathcal{S}^{(\ge a)}_C$ be the set of internal histories
$\mathfrak{h}$ of $C$ with $\min_{Z\in\mathcal{Z}(\mathfrak{h})} j(Z)\le k_0-a$.
\begin{enumerate}
\item[(a)] With the weight $L^{ba}$ in place of $L^{8(a+1)}$, the series in part~(ii) of
Theorem~\ref{8.9:thm-live-age-decay} converges for every $b\ge1$, the exponent of part~(i) being
otherwise arbitrary. The factor $L^{8(\mathfrak{l}+1)}$ allowed on the right-hand side
of part~(ii) is not drawn upon: the bound carries the factor $L^{-\mathfrak{l}}$ instead.
\item[(b)] The rate $e^{-\theta p_0(g_{k_0})}$ in part~(ii) comes from the factor of
$\varepsilon^{\natural}_{k_0}$ in \eqref{8.9:eq-live-age-decay-a} that depends on the level $k_0$
only, not from the flat
reserve of the oldest seed. The whole reserve $w^{\natural}$ is therefore available for the age.
If one preferred, the rate could also be taken from the oldest seed. Its birth level $j$ satisfies
$\theta p_0(g_j)\ge\theta\,(p_0(g_{k_0})-1)$ by the comparison of the degree function
$p_0(\cdot)$ between the birth level and level $k_0$, which is one of the hypotheses of
Theorem~\ref{8.9:thm-live-age-decay}.
\item[(c)] The flat weight $w_\theta$ alone does not suffice, and the linear reserve
$w^{\mathrm{lin}}_{\theta'}$ of \eqref{8.9:eq-reserve-weights-a} is needed, for the following reason.
A single seed $Z$ created at its birth level $j=j(Z)$ with linear size $d'=d'_j(Z)$ remains live,
by the lifetime bound for a single seed in terms of its linear size, for approximately
$2\log_L d'+\check{R}_j$ further levels, as long as no newly created region joins it; here
$\check{R}_j$ is Ba{\l}aban's number $\check{R}_n$ at the birth level $n=j$. At
level $k_0$ it is small, since linear sizes shrink by a factor $L^{-1/2}$ per level
\cite[(1.81)]{B-CMP122b}. Hence neither the flat exponent $\theta p_0$ nor the
term $-\kappa_1 d_{k_0}(C)$ at level $k_0$ in the exponent of $\varepsilon^{\natural}_{k_0}$
depends on $d'$. The linear reserve does depend on $d'$, and it overpays: it is of order
$p_0(g_j)^2\,(d'+1)$, against a requirement of order $\ln(d'+1)$.

In a per-cube bookkeeping, in which each cube of a large-field region counts as one creation and
carries its own factor, the flat reserve per cube would suffice. Write $c_q\,p_0(g_n)$ for the flat
reserve of one cube created at level $n$, and $C'$ for the summand of the lifetime count of such a
region that is a constant, independent of $s$. The requirement for a region of $s$ cubes created
at level $n$ then reads
\begin{equation}\label{8.9:eq-reserve-allocation-1}
s\,c_q\,p_0(g_n)\;\ge\;(b+1)\ln L\,\bigl(2\check{R}_n+C'+2\log_L(5^d s+1)\bigr).
\end{equation}
This alternative bookkeeping is not used. The large-field quantities entering the proof of
Theorem~\ref{8.9:thm-live-age-decay} are indexed by seed regions, not by cubes, and this is why
the linear reserve of Definition~\ref{8.9:def-reserve-weights} is used.
\item[(d)] No part of the reserve is spent twice. The flat fraction $\theta$ is used by no other
factor of $\varepsilon^{\natural}_{k_0}$ when $B=\Omega$. The fraction $\theta'_{\mathrm{age}}$ is
disjoint from the fraction $\theta'_A$ by the split in \eqref{8.9:eq-reserve-weights-a}. The margin
in $\kappa_1$ and the constant $C_{\mathrm{rm}}$ are exactly those consumed by the shape sum in the
proof of Theorem~\ref{8.9:thm-live-age-decay}, under the same lower bound on $\kappa_1$.
\item[(e)] Conditional forms. Assume the annealed live-tail statement with weight
$w^{\natural}$, display \eqref{8.9:eq-live-age-decay-a}, for every event $B$ allowed in that
statement and not only for $B=\Omega$ (a statement that is not proved in this book and is not used
in the proof of Theorem~\ref{3.1:thm-main}). Then the proof of part~(i) gives, for every level-$k_0$
cell $\square$, every $a\ge1$, and every event $B$ of levels $\le k_0$ that is $C$-lineage-blind
for every touch-connected union $C$ of level-$k_0$ cells containing $\square$,
\begin{equation}\label{8.9:eq-reserve-allocation-2}
\tilde{\nu}\bigl(\mathfrak{lv}^{\ge a,\mathrm{seed}}_{k_0}(\square)\cap B\bigr)
\;\le\;C_{\mathrm{old}}\,L^{-(b+1)a}\,e^{-(2-\theta)p_0(g_{k_0})/(1+\beta_0)}\,\tilde{\nu}(B).
\end{equation}
This form is not needed where Theorem~\ref{8.9:thm-live-age-decay} is applied, since that
application uses only $B=\Omega$.
\end{enumerate}
\end{remark}

\begin{proof}[Proof of items \textup{(a)}, \textup{(b)} and \textup{(e)}]
(a) Write $\mathcal{E}$ for the factor $e^{-(2-\theta)(1+\beta_0)^{-1}p_0(g_{k_0})}$.
Part~(i) of Theorem~\ref{8.9:thm-live-age-decay}, display~\eqref{8.9:eq-live-age-decay-b}, gives
$q^{\ge a}_{k_0}(\square)\le C_{\mathrm{old}}L^{-(b+1)a}\mathcal{E}$, and the same bound holds
a fortiori for $q^{=a}_{k_0}$. Hence, for the weight $L^{ba}$, the series of part~(ii) satisfies
\begin{equation*}
\sum_{a>\mathfrak{l}}L^{ba}\,q^{=a}_{k_0}
\;\le\;C_{\mathrm{old}}\,\mathcal{E}\sum_{a\ge\mathfrak{l}+1}L^{-a}
\;=\;C_{\mathrm{old}}\,\mathcal{E}\,\frac{L^{-\mathfrak{l}}}{L-1}.
\end{equation*}
This is finite for every $b$. With $b=8$ and the weight $L^{8(a+1)}$ of part~(ii), part~(i) gives
$L^{8(a+1)}L^{-9a}=L^{8}L^{-a}$, and the same summation yields the factor
$L^{8-\mathfrak{l}}(L-1)^{-1}$. So the factor $L^{8(\mathfrak{l}+1)}$ appearing in
part~(ii) is an allowance that the bound does not draw upon; the bound itself decays like
$L^{-\mathfrak{l}}$.

(b) In the proof of part~(i) the bound $w^{\natural}(\mathcal{S}^{(\ge a)}_C)\le L^{-(b+1)a}$ uses
only the age lemma and the age-payment inequality of Lemma~\ref{8.9:lem-age-payment}, applied to
the flat and the linear exponents of $w^{\natural}$ together; so the whole of $w^{\natural}$ is
spent on the age.

The factor $\mathcal{E}$ in \eqref{8.9:eq-live-age-decay-b} is the factor
$\exp(-(2-\theta)(1+\beta_0)^{-1}p_0(g_k))$ of $\varepsilon^{\natural}_k$ at $k=k_0$. In part~(ii)
it is bounded by $e^{-\theta p_0(g_{k_0})}$.

For the alternative, let $Z_{\mathrm{old}}$ be a seed of $\mathfrak{h}$ whose birth level
$j=j(Z_{\mathrm{old}})$ is minimal. The comparison $\theta p_0(g_j)\ge\theta(p_0(g_{k_0})-1)$ gives
\begin{equation*}
e^{-\theta p_0(g_j)}\;\le\;e^{\theta}\,e^{-\theta p_0(g_{k_0})}.
\end{equation*}
The flat factor of the oldest seed therefore yields the rate $e^{-\theta p_0(g_{k_0})}$ as well,
up to the constant $e^{\theta}$.

(e) Fix $\square$ and $a\ge1$, and repeat the proof of part~(i) with $B$ in place of $\Omega$.
Intersecting with $B$ the inclusion of $\mathfrak{lv}^{\ge a,\mathrm{seed}}_{k_0}(\square)$ in the
union of the events $L_{C,\mathcal{S}^{(\ge a)}_C}$ used there reduces the left-hand side of
\eqref{8.9:eq-reserve-allocation-2} to a sum over the touch-connected unions $C$ of level-$k_0$
cells containing $\square$. For each such $C$ the event $B$ is of levels $\le k_0$ and
$C$-lineage-blind by assumption, so the annealed live-tail statement with weight $w^{\natural}$
applies to $C$, $\mathcal{S}^{(\ge a)}_C$ and $B$. The only change is that
$\tilde{\nu}(B\setminus L_C)$ replaces $\tilde{\nu}(\Omega\setminus L_C)\le1$, and
\begin{equation*}
\tilde{\nu}\bigl(L_{C,\mathcal{S}^{(\ge a)}_C}\cap B\bigr)
\;\le\;\varepsilon^{\natural}_{k_0}\bigl(C,\mathcal{S}^{(\ge a)}_C\bigr)\,\tilde{\nu}(B\setminus L_C)
\;\le\;\varepsilon^{\natural}_{k_0}\bigl(C,\mathcal{S}^{(\ge a)}_C\bigr)\,\tilde{\nu}(B),
\end{equation*}
where the last step uses $B\setminus L_C\subset B$. All other factors do not involve $B$ and are
as in the proof of part~(i): the age factor $w^{\natural}(\mathcal{S}^{(\ge a)}_C)\le L^{-(b+1)a}$,
and the bound $2e^{\kappa_1}=C_{\mathrm{old}}$ given by the shape sum. Collecting them gives
\eqref{8.9:eq-reserve-allocation-2}.
\end{proof}

\begin{proposition}[The healed old-region bound]\label{8.9:prop-healed-old-region}
Fix an integer $b\ge1$ and the reserve parameters $\beta_0$ and
$\theta\in\left]0,(1-3\beta_0)/2\right[$. Split the linear-reserve fraction as
$\theta'=\theta'_A+\theta'_{\mathrm{age}}$, where $\theta'_{\mathrm{age}}$ lies in the range in which
the linear-in-size reserve per seed of Definition~\ref{8.9:def-reserve-weights} is available. Let
$w_\theta$, $w^{\mathrm{lin}}_{\theta'}$ and
$w^{\natural}=w_\theta\, w^{\mathrm{lin}}_{\theta'_{\mathrm{age}}}$ be the reserve weights of
\eqref{8.9:eq-reserve-weights-a}, and let $\mathfrak{a}(Z)$ be the per-seed charge of
\eqref{8.9:eq-age-lemma-a}. Let $\tilde{\nu}$ be the annealed normalised large-field measure on
histories, $k_0$ the witness level and $K'=K-m_{\mathbb{T}}(g)$ the terminal level. For a level-$k_0$
cell $\square$ let $\mathcal{E}^{(k_0)}_{\square}(K')$, the \emph{live event}, be the event that the
lineage of the level-$k_0$
component containing $\square$ is healed at no step $k$ with $k_0<k\le K'$, and let $K_H$ be the
constant of the annealed bound on healed islands, which is proved from the inputs of
hypothesis~(i) below and recalled in the proof. Assume the
following.
\begin{itemize}
\item[(i)] The inputs of the large-field block bound, healed part, where its flat-reserve strip
bound, the bound
\begin{equation}\label{8.9:eq-healed-old-region-11}
\mathbf{T}'_{k,\mathcal{S}}(X,(U,0))1
\le w_\theta(\mathcal{S})\,
\exp\bigl(-(2-\theta)(1+\beta_0)^{-1}p_0(g_k)\bigr)
\end{equation}
for every set $\mathcal{S}$ of internal histories, is replaced by~(ii).
\item[(ii)] For every set $\mathcal{S}$ of internal histories,
\begin{equation}\label{8.9:eq-healed-old-region-1}
\mathbf{T}'_{k,\mathcal{S}}(X,(U,0))1
\le w^{\natural}(\mathcal{S})\,
\exp\bigl(-(2-\theta)(1+\beta_0)^{-1}p_0(g_k)\bigr).
\end{equation}
\item[(iii)] The age lemma, Lemma~\ref{8.9:lem-age-lemma}, whose proof is incomplete at one step.
\item[(iv)] The coupling ceiling of Lemma~\ref{8.9:lem-age-payment} for the given $b$, in the
following form: for every $y\ge g^{-2}$,
\begin{equation}\label{8.9:eq-healed-old-region-2}
\begin{gathered}
(b+1)\ln L\,\bigl(2R(y)+n_1+C_S+4\bigr)\le \theta A_0(\log y)^{p_0},\\
(b+1)\ln L\,\bigl(2R(y)+n_1+C_S+4\bigr)
\le \tfrac12\,\theta'_{\mathrm{age}}\gamma_0A_1A_0^2(\log y)^{2p_0},\\
2(b+1)\le \tfrac12\,\theta'_{\mathrm{age}}\gamma_0A_1A_0^2(\log y)^{2p_0}.
\end{gathered}
\end{equation}
\item[(v)] (First condition on the coupling flow $k\mapsto g_k$.) For every component live at
level $k_0$ and every internal history of it, the oldest seed
$Z_{\mathrm{old}}$ of that history satisfies
$p_0(g_{j(Z_{\mathrm{old}})})\ge p_0(g_{k_0})-1$.
\item[(vi)] (Second condition on the coupling flow $k\mapsto g_k$.) There is a number $S_*$,
independent of $K$, such that
\begin{equation}\label{8.9:eq-healed-old-region-3}
\sum_{k_0<k\le K'}e^{-\frac32 p_0(g_k)}\le S_*\le 1+\frac{1}{\lambda^\sharp},
\end{equation}
and the constant $C_L$ of \eqref{8.9:eq-seed-age-a} satisfies $K_He^{\theta}S_*\le C_L$.
\item[(vii)] The condition on Ba{\l}aban's parameter $\kappa$ assumed together with the inputs of
hypothesis~(i).
\item[(viii)] The standing hypothesis at the parameter $\theta$ used in the fourth step of the lemma
on island classes recalled in the proof; it is non-vacuous for every
$\theta\in\left]0,(1-3\beta_0)/2\right[$.
\end{itemize}
Then for every $K$, every admissible run, every $N_T$, every level-$k_0$ cell $\square$ and every
$a\ge1$,
\begin{equation}\label{8.9:eq-healed-old-region-4}
\begin{aligned}
\tilde{\nu}\bigl(\mathfrak{lv}^{\ge a,\mathrm{seed}}_{k_0}(\square)
\setminus\mathcal{E}^{(k_0)}_{\square}(K')\bigr)
&\le K_He^{\theta}S_*\,L^{-(b+1)a}\,e^{-\theta p_0(g_{k_0})}\\
&\le K_He^{\theta}\Bigl(1+\frac{1}{\lambda^\sharp}\Bigr)L^{-(b+1)a}\,e^{-\theta p_0(g_{k_0})}.
\end{aligned}
\end{equation}
Hence, for $b=8$, with the lifetime $\mathfrak{l}$ and the live rate
$q^{\mathrm{lv}}_{k_0}=C_Le^{-\theta p_0(g_{k_0})}$ of \eqref{8.9:eq-seed-age-a}, and using
$K_He^{\theta}S_*\le C_L$ from hypothesis~(vi),
\begin{equation}\label{8.9:eq-healed-old-region-5}
\begin{aligned}
\sum_{a>\mathfrak{l}}L^{8(a+1)}\,
\tilde{\nu}\bigl(\mathfrak{lv}^{=a}_{k_0}(\square)\setminus\mathcal{E}^{(k_0)}_{\square}(K')\bigr)
&\le (L-1)^{-1}L^{8-\mathfrak{l}}\,q^{\mathrm{lv}}_{k_0}\\
&\le (L-1)^{-1}C_L^{-1}L^{8(\mathfrak{l}+1)}\,q^{\mathrm{lv}}_{k_0}.
\end{aligned}
\end{equation}
Moreover, for every internal history $\mathfrak{h}$ whose seed set $\mathcal{Z}=\mathcal{Z}(\mathfrak{h})$,
with oldest seed $Z_{\mathrm{old}}$, satisfies $\sum_{Z\in\mathcal{Z}}\mathfrak{a}(Z)\ge a$,
\begin{equation}\label{8.9:eq-healed-old-region-6}
\begin{aligned}
&\theta\sum_{Z\in\mathcal{Z},\,Z\ne Z_{\mathrm{old}}}p_0(g_{j(Z)})
+\theta'_{\mathrm{age}}\gamma_0A_1\sum_{Z\in\mathcal{Z}}p_0(g_{j(Z)})^2\bigl(d'_{j(Z)}(Z)+1\bigr)\\
&\qquad\ge (b+1)\ln L\sum_{Z\in\mathcal{Z}}\mathfrak{a}(Z)\ \ge\ (b+1)\,a\ln L.
\end{aligned}
\end{equation}
\end{proposition}

The left-hand
sides of \eqref{8.9:eq-healed-old-region-2} have degrees $r$, $r$ and $0$ in $\log y$, the
right-hand sides degrees $p_0$, $2p_0$ and $2p_0$, as in Lemma~\ref{8.9:lem-age-payment}.

Hypothesis~(ii) is the strip bound, per internal history, of the normalised operator $\mathbf{T}'$,
in which the reserve stripped along an internal history $\mathfrak{h}$ is the flat reserve plus the
$\theta'_{\mathrm{age}}$-linear reserve,
\begin{equation*}
\Theta(\mathfrak{h}):=\theta\sum_{Z\in\mathcal{Z}(\mathfrak{h})}p_0(g_{j(Z)})
+\theta'_{\mathrm{age}}\gamma_0A_1\sum_{Z\in\mathcal{Z}(\mathfrak{h})}
p_0(g_{j(Z)})^2\bigl(d'_{j(Z)}(Z)+1\bigr);
\end{equation*}
it asks that the majorants of the terms of $\mathbf{T}'$, after the factor $e^{-\Theta(\mathfrak{h})}$
has been taken out along each internal history, still sum to at most
$e^{-(2-\theta)(1+\beta_0)^{-1}p_0(g_k)}$, the largeness of $p_0$ in Ba{\l}aban's construction being
used together with the remaining part $(1-\theta')\gamma_0A_1p_0(g_{j(Z)})^2(d'_{j(Z)}(Z)+1)$ of the
linear term.

The reserve is allocated as follows. The rate $e^{-\theta p_0(g_{k_0})}$ in
\eqref{8.9:eq-healed-old-region-4} comes from the flat member of the oldest seed alone, through
hypothesis~(v); the factor $e^{-\frac32 p_0(g_k)}$ of the healing step $k>k_0$ gives only the
summability \eqref{8.9:eq-healed-old-region-3}, independent of $K$, and never the rate. The age factor
$L^{-(b+1)a}$ comes from the flat members of all the other seeds and from the
$\theta'_{\mathrm{age}}$-linear members of all seeds; the part of $\mathfrak{a}(Z_{\mathrm{old}})$
carrying the number $R$ is paid from the linear member of the oldest seed. This is the content of
\eqref{8.9:eq-healed-old-region-6}.

\begin{proof}
\emph{The allocation inequality \eqref{8.9:eq-healed-old-region-6}.} Let $\mathfrak{h}$ have seed
set $\mathcal{Z}$ with $\sum_{Z}\mathfrak{a}(Z)\ge a$. Fix $Z\in\mathcal{Z}$ and put
$y:=x_{j(Z)}$; the inverse couplings satisfy $x_j\ge g^{-2}$ at every level, as in the statement of
Lemma~\ref{8.9:lem-age-payment}, so $y\ge g^{-2}$ and
\eqref{8.9:eq-healed-old-region-2} applies at $y$, the number $\check{R}_{j(Z)}$ being $R(y)$ as in
the derivation of \eqref{8.9:eq-age-payment-b}. By the bounds on the per-seed charge in the age
lemma (Lemma~\ref{8.9:lem-age-lemma}, whose proof is incomplete at one step),
\begin{equation*}
\mathfrak{a}(Z)\le 2\log_L\bigl(d'_{j(Z)}(Z)+1\bigr)+2R(y)+n_1+C_S+4,
\end{equation*}
where for a preparatory seed $d'_{j(Z)}(Z):=0$, as in Definition~\ref{8.9:def-reserve-weights}, so
that the logarithmic term vanishes. Since
\begin{equation*}
(b+1)\ln L\cdot 2\log_L(d'+1)=2(b+1)\ln(d'+1)\le 2(b+1)(d'+1),
\end{equation*}
we obtain
\begin{equation}\label{8.9:eq-healed-old-region-7}
(b+1)\ln L\,\mathfrak{a}(Z)\le (b+1)\ln L\,\bigl(2R(y)+n_1+C_S+4\bigr)
+2(b+1)\bigl(d'_{j(Z)}(Z)+1\bigr).
\end{equation}
By the definition of the degree function, $p_0(g_{j(Z)})=A_0(\log y)^{p_0}$ and
$p_0(g_{j(Z)})^2=A_0^2(\log y)^{2p_0}$. The third inequality of
\eqref{8.9:eq-healed-old-region-2}, multiplied by $d'_{j(Z)}(Z)+1$, bounds the last term of
\eqref{8.9:eq-healed-old-region-7} by
$\tfrac12\theta'_{\mathrm{age}}\gamma_0A_1p_0(g_{j(Z)})^2(d'_{j(Z)}(Z)+1)$. If $Z\ne Z_{\mathrm{old}}$,
the first inequality of \eqref{8.9:eq-healed-old-region-2} bounds the first term on the right of
\eqref{8.9:eq-healed-old-region-7} by $\theta p_0(g_{j(Z)})$. If $Z=Z_{\mathrm{old}}$, the second
inequality of \eqref{8.9:eq-healed-old-region-2} bounds it by
$\tfrac12\theta'_{\mathrm{age}}\gamma_0A_1p_0(g_{j(Z)})^2$, which is at most the same quantity
multiplied by $d'_{j(Z)}(Z)+1\ge1$. Hence
\begin{equation*}
\begin{aligned}
(b+1)\ln L\,\mathfrak{a}(Z)&\le \theta p_0(g_{j(Z)})
+\tfrac12\theta'_{\mathrm{age}}\gamma_0A_1p_0(g_{j(Z)})^2\bigl(d'_{j(Z)}(Z)+1\bigr)
&&(Z\ne Z_{\mathrm{old}}),\\
(b+1)\ln L\,\mathfrak{a}(Z)&\le \theta'_{\mathrm{age}}\gamma_0A_1p_0(g_{j(Z)})^2
\bigl(d'_{j(Z)}(Z)+1\bigr)
&&(Z=Z_{\mathrm{old}}).
\end{aligned}
\end{equation*}
Summing over $Z\in\mathcal{Z}$, and dropping the factor $\frac12$ in the linear members of the seeds
$Z\ne Z_{\mathrm{old}}$ (they are nonnegative), gives the first inequality of
\eqref{8.9:eq-healed-old-region-6}; the second follows from $\sum_Z\mathfrak{a}(Z)\ge a$ and
$\ln L>0$.

\emph{Decomposition by the healing step.} On
$\mathfrak{lv}^{\ge a,\mathrm{seed}}_{k_0}(\square)$ the cell $\square$ lies in a component $C$ live
at level $k_0$. Off $\mathcal{E}^{(k_0)}_{\square}(K')$ the lineage of $C$ is healed at some step $k$
with $k_0<k\le K'$; let $H_k$ be the event that it is first healed at step $k$. Then
\begin{equation}\label{8.9:eq-healed-old-region-8}
\mathfrak{lv}^{\ge a,\mathrm{seed}}_{k_0}(\square)\setminus\mathcal{E}^{(k_0)}_{\square}(K')
=\bigcup_{k_0<k\le K'}\Bigl(H_k\cap\mathfrak{lv}^{\ge a,\mathrm{seed}}_{k_0}(\square)\Bigr).
\end{equation}
We follow the argument by which the corresponding estimate with the flat weight $w_\theta$ in place
of $w^{\natural}$ is derived from the inputs of hypothesis~(i). On $H_k$ the island $X$ healed at
step $k$ (the component, in its class, of the large-field region healed at step $k$; such components
have at most $100$ cells to a side) contains the
level-$k$ cell $\mathfrak{c}:=\square^{(k)}$ containing $\square$, and its internal history
$\mathfrak{h}_X$ contains the internal history $\mathfrak{h}_C$ of $C$. Let
$\mathcal{S}^{(k,a)}$ be the set of internal histories of islands that contain a level-$k_0$ live
ancestor containing $\square$ whose seed age at level $k_0$ is at least $a$. For $k_0<k\le K'$ let
$E^{(k)}_{\mathfrak{c},\mathcal{S}}$ be the event that the island healed at step $k$ and containing
$\mathfrak{c}$ has its internal history in $\mathcal{S}$. By what was just said,
\begin{equation*}
H_k\cap\mathfrak{lv}^{\ge a,\mathrm{seed}}_{k_0}(\square)
\subset E^{(k)}_{\mathfrak{c},\mathcal{S}^{(k,a)}}.
\end{equation*}

\emph{The annealed bound on a healed island.} The annealed bound on
$\tilde{\nu}(E^{(k)}_{\mathfrak{c},\mathcal{S}})$ is proved from the inputs of hypothesis~(i), the
flat-reserve strip bound \eqref{8.9:eq-healed-old-region-11} among them, through the lemma on island
classes in four steps. We rerun
that proof with \eqref{8.9:eq-healed-old-region-1} in place of the flat-reserve strip
bound and obtain
\begin{equation}\label{8.9:eq-healed-old-region-9}
\tilde{\nu}\bigl(E^{(k)}_{\mathfrak{c},\mathcal{S}^{(k,a)}}\bigr)
\le K_H\,w^{\natural}\bigl(\mathcal{S}^{(k,a)}\bigr)\,e^{-\frac32 p_0(g_k)}.
\end{equation}
Indeed, the first three steps of the lemma on island classes do not involve the strip bound and are
unchanged. Its fourth step bounds the class-restricted term $t_X(\mathcal{S})$ of that lemma, multiplied
by the correction factor $e^{\delta_X}$ of that step, through
\begin{equation*}
t_X(\mathcal{S})e^{\delta_X}
\le w_\theta(\mathcal{S})\,
e^{-(2-\theta)(1+\beta_0)^{-1}p_0(g_k)+C_\delta\alpha(101\check{R}_k)^d}
\le w_\theta(\mathcal{S})\,e^{-\frac32 p_0(g_k)},
\end{equation*}
the first inequality being the flat-reserve strip bound \eqref{8.9:eq-healed-old-region-11}. Reading
\eqref{8.9:eq-healed-old-region-1}
there instead gives the same chain with $w^{\natural}(\mathcal{S})$ in place of
$w_\theta(\mathcal{S})$; the second inequality is hypothesis~(viii), which is unchanged. Finally,
the annealing in the bound on
$\tilde{\nu}(E^{(k)}_{\mathfrak{c},\mathcal{S}})$ is linear in its rate $\varepsilon_k(\mathcal{S})$,
in which the weight enters as a factor, so $w^{\natural}(\mathcal{S}^{(k,a)})$ passes through to
\eqref{8.9:eq-healed-old-region-9}.

\emph{The age factor.} Let $\mathfrak{h}_X\in\mathcal{S}^{(k,a)}$, and let $C$ be a level-$k_0$ live
ancestor containing $\square$ with seed age at least $a$, whose internal history $\mathfrak{h}_C$ is
contained in $\mathfrak{h}_X$. Then $\mathcal{Z}(\mathfrak{h}_X)\supseteq\mathcal{Z}(\mathfrak{h}_C)$.
Every seed contributes a nonnegative amount to the reserve exponent $\Theta$, so extra seeds only
increase it and
\begin{equation*}
\Theta(\mathfrak{h}_X)\ge\Theta(\mathfrak{h}_C).
\end{equation*}
Since $C$ is live at level $k_0$ with
$\operatorname{age}^{\mathrm{seed}}_{k_0}(C)\ge a$, the age lemma (Lemma~\ref{8.9:lem-age-lemma},
whose proof is incomplete at one step) at level $k_0$ gives
$\sum_{Z\in\mathcal{Z}(\mathfrak{h}_C)}\mathfrak{a}(Z)\ge a$. Let $Z_{\mathrm{old}}$ be the oldest
seed of $\mathfrak{h}_C$. Separating its flat member and applying
\eqref{8.9:eq-healed-old-region-6} to $\mathfrak{h}_C$,
\begin{equation*}
\Theta(\mathfrak{h}_C)\ge \theta p_0(g_{j(Z_{\mathrm{old}})})+(b+1)\,a\ln L .
\end{equation*}
By hypothesis~(v), multiplied by $\theta>0$,
$\theta p_0(g_{j(Z_{\mathrm{old}})})\ge\theta(p_0(g_{k_0})-1)$. Therefore the
contribution $\exp(-\Theta(\mathfrak{h}_X))$ of $\mathfrak{h}_X$ to $w^{\natural}$ satisfies
\begin{equation*}
\exp\bigl(-\Theta(\mathfrak{h}_X)\bigr)
\le e^{-\theta p_0(g_{j(Z_{\mathrm{old}})})}L^{-(b+1)a}
\le e^{\theta}e^{-\theta p_0(g_{k_0})}L^{-(b+1)a}.
\end{equation*}
Taking the supremum over $\mathfrak{h}_X\in\mathcal{S}^{(k,a)}$,
\begin{equation}\label{8.9:eq-healed-old-region-10}
w^{\natural}\bigl(\mathcal{S}^{(k,a)}\bigr)\le e^{\theta}e^{-\theta p_0(g_{k_0})}L^{-(b+1)a}.
\end{equation}

\emph{Summation over the healing step.} By \eqref{8.9:eq-healed-old-region-8}, the inclusion into
$E^{(k)}_{\mathfrak{c},\mathcal{S}^{(k,a)}}$, \eqref{8.9:eq-healed-old-region-9} and
\eqref{8.9:eq-healed-old-region-10},
\begin{equation*}
\tilde{\nu}\bigl(\mathfrak{lv}^{\ge a,\mathrm{seed}}_{k_0}(\square)
\setminus\mathcal{E}^{(k_0)}_{\square}(K')\bigr)
\le K_He^{\theta}e^{-\theta p_0(g_{k_0})}L^{-(b+1)a}\sum_{k_0<k\le K'}e^{-\frac32 p_0(g_k)}.
\end{equation*}
By hypothesis~(vi), \eqref{8.9:eq-healed-old-region-3}, the sum is at most $S_*\le 1+1/\lambda^\sharp$.
This proves both inequalities of \eqref{8.9:eq-healed-old-region-4}; every factor is independent of
$K$, of the run, of $N_T$ and of the position of $\square$.

\emph{The series.} A component of territory age $a$ at level $k_0$ has seed age at least $a$
(Definition~\ref{8.9:def-seed-age}), so
$\mathfrak{lv}^{=a}_{k_0}(\square)\subset\mathfrak{lv}^{\ge a,\mathrm{seed}}_{k_0}(\square)$. Take
$b=8$ in \eqref{8.9:eq-healed-old-region-4}. Since
$\sum_{a\ge\mathfrak{l}+1}L^{-a}=L^{-\mathfrak{l}}(L-1)^{-1}$,
\begin{equation*}
\sum_{a>\mathfrak{l}}L^{8(a+1)}L^{-9a}=L^{8}\sum_{a\ge\mathfrak{l}+1}L^{-a}
=\frac{L^{8-\mathfrak{l}}}{L-1},
\end{equation*}
and since $K_He^{\theta}S_*\le C_L$ by hypothesis~(vi), so that
\begin{equation*}
K_He^{\theta}S_*e^{-\theta p_0(g_{k_0})}\le C_Le^{-\theta p_0(g_{k_0})}=q^{\mathrm{lv}}_{k_0},
\end{equation*}
we have
\begin{equation*}
\sum_{a>\mathfrak{l}}L^{8(a+1)}\,
\tilde{\nu}\bigl(\mathfrak{lv}^{=a}_{k_0}(\square)\setminus\mathcal{E}^{(k_0)}_{\square}(K')\bigr)
\le K_He^{\theta}S_*e^{-\theta p_0(g_{k_0})}\,\frac{L^{8-\mathfrak{l}}}{L-1}
\le (L-1)^{-1}L^{8-\mathfrak{l}}\,q^{\mathrm{lv}}_{k_0}.
\end{equation*}
This is the first inequality of \eqref{8.9:eq-healed-old-region-5}. Writing
$L^{8-\mathfrak{l}}=L^{-9\mathfrak{l}}\,L^{8(\mathfrak{l}+1)}$, the right-hand member of
\eqref{8.9:eq-healed-old-region-5} is the middle one with the factor $L^{-9\mathfrak{l}}$ replaced by
$C_L^{-1}$; the second inequality of \eqref{8.9:eq-healed-old-region-5} thus amounts to the
comparison $C_L\le L^{9\mathfrak{l}}$ of the constants of \eqref{8.9:eq-seed-age-a}.
\end{proof}

Proposition~\ref{8.9:prop-healed-old-region} bounds the part of the old-age series carried by live
regions that are healed at a later step. Its proof uses only the inputs of the large-field block
bound, healed part, with \eqref{8.9:eq-healed-old-region-1} in place of the flat-reserve strip bound,
and the age lemma (Lemma~\ref{8.9:lem-age-lemma}, whose proof is incomplete at one step), under the
coupling ceiling~(iv) and the conditions~(v) and~(vi) on the coupling flow. It involves
no union over regions of unbounded extent, since every island is a component of its class and has at
most $100$ cells to a side, and it imposes no lower bound on the depth beyond those already fixed.

\begin{remark}[Old regions never healed by the terminal level]\label{8.9:rem-never-healed}
Fix $K$ and consider the longest admissible run, whose terminal level is
$K'=K-m_{\mathbb{T}}(g)$; fix a level-$k_0$ cell $\square$ and an integer $a\ge1$. We write
$\square^{(K')}$ for the level-$K'$ cell
containing $\square$. Proposition~\ref{8.9:prop-healed-old-region} (the healed old-region bound)
estimates the part of the event $\mathfrak{lv}^{\ge a,\mathrm{seed}}_{k_0}(\square)$ on which the
lineage of the level-$k_0$ live component containing $\square$ is healed at some step
$k\in\,]k_0,K']$. We write $\mathcal{E}^{(k_0)}_{\square}(K')$, the live event, for the set of
histories in which that lineage is not healed at any such step. The part that remains is the event
\begin{equation}\label{8.9:eq-never-healed-1}
\begin{split}
\mathfrak{lv}^{\ge a,\mathrm{seed}}_{k_0}(\square)\cap\mathcal{E}^{(k_0)}_{\square}(K')
=\{&\text{the level-}k_0\text{ component of }\square\text{ has seed age}\ge a\text{ at }k_0,\\
&\text{and its lineage is live at the terminal level }K'\}.
\end{split}
\end{equation}
The contribution of this event to the series over the age $a$ whose healed part is bounded in
Proposition~\ref{8.9:prop-healed-old-region} is the never-healed old-age tail of
clause~(iv) of Theorem~\ref{3.1:thm-main}.

Let $\mathcal{S}^{(\ge a)}_{k_0}$ be the class of internal histories $\mathfrak{h}$ of
touch-connected unions of level-$K'$ cells such that $\mathfrak{h}$ contains a level-$k_0$ live
ancestor $C\ni\square$ with $\operatorname{age}^{\mathrm{seed}}_{k_0}(C)\ge a$. Let $\theta$ and
$\theta'_{\mathrm{age}}$ lie in the windows of Theorem~\ref{8.9:thm-live-age-decay}, let $b$ be
the exponent of the coupling ceiling of Proposition~\ref{8.9:prop-healed-old-region}, and assume
the hypotheses of Theorem~\ref{8.9:thm-live-age-decay} other than the annealed live-tail
statement at level $k_0$ (among them the age lemma and the coupling ceiling). Suppose in addition
that the annealed live-tail statement, in the notation
of Theorem~\ref{8.9:thm-live-age-decay} with the rate $\varepsilon^{\natural}_k$ of
\eqref{8.9:eq-live-age-decay-a}, holds at level $k=K'$ for $B$ equal to the sure event. Then
\begin{equation}\label{8.9:eq-never-healed-2}
\begin{split}
\tilde{\nu}\bigl(\mathfrak{lv}^{\ge a,\mathrm{seed}}_{k_0}(\square)\cap\mathcal{E}^{(k_0)}_{\square}(K')\bigr)
&\le\sum_{C_{\mathrm{fin}}\ni\square^{(K')}}
\varepsilon^{\natural}_{K'}\bigl(C_{\mathrm{fin}},\mathcal{S}^{(\ge a)}_{k_0}\bigr)\\
&\le 2e^{\kappa_1+\theta}\,L^{-(b+1)a}\,e^{-\theta p_0(g_{k_0})}\,
e^{-(2-\theta)(1+\beta_0)^{-1}p_0(g_{K'})},
\end{split}
\end{equation}
where the sum runs over the touch-connected unions $C_{\mathrm{fin}}$ of level-$K'$ cells
containing $\square^{(K')}$. The annealed live-tail statement at level $K'$ is not proved; it is
an open statement on which no clause of Theorem~\ref{3.1:thm-main} depends; accordingly
\eqref{8.9:eq-never-healed-2} holds only under that statement. On the reference torus
the event \eqref{8.9:eq-never-healed-1} is not estimated through any live-tail statement: under
the depth floor displayed in clause~(iv) of Theorem~\ref{3.1:thm-main}, its contribution is
bounded directly, before normalisation, by the age factor accumulated up to the last lattice of
the run and prolonged on that lattice; this is part of the proof of the two-point witness on every torus,
Theorem~\ref{8.18:thm-torus-witness}, in this chapter.
\end{remark}

\begin{proof}
The identity \eqref{8.9:eq-never-healed-1} unwinds the definitions. The event
$\mathfrak{lv}^{\ge a,\mathrm{seed}}_{k_0}(\square)$ requires that $\square$ lie in a level-$k_0$
live component of seed age at least $a$. The event $\mathcal{E}^{(k_0)}_{\square}(K')$ requires that the
lineage of that component not be healed at any step up to $K'$, which means that it is still
live at level $K'$.

We now prove the first inequality of \eqref{8.9:eq-never-healed-2}. On the event
\eqref{8.9:eq-never-healed-1} the lineage of the level-$k_0$ component of $\square$ is live at
level $K'$. Hence there is a live component $C_{\mathrm{fin}}$ at level $K'$, a touch-connected
union of level-$K'$ cells, which contains $\square^{(K')}$ and whose internal history contains a
level-$k_0$ live ancestor containing $\square$ of seed age at least $a$. That internal history
therefore belongs to $\mathcal{S}^{(\ge a)}_{k_0}$. The event is thus contained in the union,
over such $C_{\mathrm{fin}}$, of the lineage events
$L_{C_{\mathrm{fin}},\mathcal{S}^{(\ge a)}_{k_0}}$.

Let $B$ be the sure event, the set of all histories; it is lineage-blind for every
$C_{\mathrm{fin}}$. The annealed live-tail statement at level $K'$ with this $B$ therefore gives
\begin{equation*}
\tilde{\nu}\bigl(L_{C_{\mathrm{fin}},\mathcal{S}^{(\ge a)}_{k_0}}\bigr)
\le\varepsilon^{\natural}_{K'}\bigl(C_{\mathrm{fin}},\mathcal{S}^{(\ge a)}_{k_0}\bigr)\,
\tilde{\nu}\bigl(B\setminus L_{C_{\mathrm{fin}}}\bigr)
\le\varepsilon^{\natural}_{K'}\bigl(C_{\mathrm{fin}},\mathcal{S}^{(\ge a)}_{k_0}\bigr).
\end{equation*}
The last step holds because $\tilde{\nu}$ is normalised. Summing over $C_{\mathrm{fin}}$ gives
the first inequality.

We now prove the second inequality. We first bound the reserve weight of the class
$\mathcal{S}^{(\ge a)}_{k_0}$. In the proof of Proposition~\ref{8.9:prop-healed-old-region} the
reserve weight $w^{\natural}$ of the class of island internal histories that contain a
level-$k_0$ live ancestor containing $\square$ of seed age at least $a$ at level $k_0$ is bounded
using only this defining property, the age lemma at level $k_0$, the allocation of the reserve
that follows from the coupling ceiling, and the flow condition on the degree function
$p_0(\cdot)$ which gives $p_0(g_{j(Z_{\mathrm{old}})})\ge p_0(g_{k_0})-1$ for the oldest seed
$Z_{\mathrm{old}}$. Every $\mathfrak{h}\in\mathcal{S}^{(\ge a)}_{k_0}$ has the same defining
property, so the same argument, applied to $\mathcal{S}^{(\ge a)}_{k_0}$, gives
\begin{equation}\label{8.9:eq-never-healed-3}
w^{\natural}\bigl(\mathcal{S}^{(\ge a)}_{k_0}\bigr)
\le e^{\theta}\,e^{-\theta p_0(g_{k_0})}\,L^{-(b+1)a}.
\end{equation}

The rate $\varepsilon^{\natural}_{K'}(C_{\mathrm{fin}},\mathcal{S})$ of
\eqref{8.9:eq-live-age-decay-a} carries, beside the reserve weight $w^{\natural}(\mathcal{S})$,
the factor $e^{-(2-\theta)(1+\beta_0)^{-1}p_0(g_{K'})}$ that the bound on the normalised
large-field operator at level $K'$ contributes. By \eqref{8.9:eq-live-age-decay-b}, read at
level $K'$ for the class $\mathcal{S}=\mathcal{S}^{(\ge a)}_{k_0}$, the sum over the
touch-connected unions $C_{\mathrm{fin}}$ of level-$K'$ cells containing $\square^{(K')}$
satisfies
\begin{equation*}
\sum_{C_{\mathrm{fin}}\ni\square^{(K')}}
\varepsilon^{\natural}_{K'}\bigl(C_{\mathrm{fin}},\mathcal{S}^{(\ge a)}_{k_0}\bigr)
\le 2e^{\kappa_1}\,w^{\natural}\bigl(\mathcal{S}^{(\ge a)}_{k_0}\bigr)\,
e^{-(2-\theta)(1+\beta_0)^{-1}p_0(g_{K'})}.
\end{equation*}
Inserting \eqref{8.9:eq-never-healed-3} gives the product
\begin{equation*}
2e^{\kappa_1}\cdot e^{\theta}e^{-\theta p_0(g_{k_0})}L^{-(b+1)a}
\cdot e^{-(2-\theta)(1+\beta_0)^{-1}p_0(g_{K'})},
\end{equation*}
which is the right-hand side of \eqref{8.9:eq-never-healed-2}.
\end{proof}

\begin{remark}[First-moment bounds are blind to age]\label{8.9:rem-first-moment-blind}
Let $k_0$ be the witness level, let $K'=K-m_{\mathbb{T}}(g)$ be the terminal level of the longest run,
let $\tilde{\nu}$ be the annealed normalised large-field measure on histories, and, for a level-$k_0$
cell $\square$, let $\mathcal{E}^{(k_0)}_{\square}(K')$ be the live event of
Remark~\ref{8.9:rem-never-healed}, namely that the
lineage of the level-$k_0$ component of $\square$ is live at the level $K'$. The first-moment lemma
bounds $\tilde{\nu}$-expectations of functionals that are additive over the components live at the
terminal level and linear in the reserve of each such component. Its inputs are positivity, the two
hypotheses of the first-moment lemma, a reserve attached to each seed, and symmetry. For $a\ge 0$, the
\emph{density of age $a$} at $\square$ is the $\tilde{\nu}$-probability of the event that the
level-$k_0$ component of $\square$ has seed age $a$ at level $k_0$ and
$\mathcal{E}^{(k_0)}_{\square}(K')$ holds;
a bound is \emph{age-blind} if it uses no decay of these densities in $a$. From these inputs alone no
bound on
\begin{equation}\label{8.9:eq-first-moment-blind-1}
  \sum_{a} L^{8a}\cdot\bigl(\text{density of age }a\bigr)
\end{equation}
better than an age-blind one follows. The sharpness clause of the first-moment lemma states that a
first moment gives no exponential age tail; the proof below exhibits the mechanism. The obstruction is
a property of the weight $L^{8\cdot\operatorname{age}}$, not of the first-moment lemma.
\end{remark}

\begin{proof}
\emph{What the first-moment lemma controls.} In the letters of the first-moment lemma, $\Theta(X)$ is
the reserve of a component $X$, $\mathsf{B}$ is the volume constant, $N^4$ is the volume, and $t$,
$c_J$, $\eta$ are its remaining parameters; expectations $\mathbb{E}$ without subscript are also taken
under $\tilde{\nu}$. The first-moment lemma bounds $\mathbb{E}_{\tilde{\nu}}[F]$ for functionals that
assign to a history $h$ the number
\begin{equation*}
  F(h)=\sum_{X\ \text{live at the terminal level}} f(X),
  \qquad f(X)\le \frac{\Theta(X)}{t},
\end{equation*}
where the reserve of a component is the sum of the reserves of its seeds,
\begin{equation*}
  \Theta(X)=\sum_{Z\in X}\operatorname{reserve}(Z).
\end{equation*}
Such an $F$ is additive in the components live at the terminal level and linear in the reserve
$\Theta(X)$ of each of them. In the proof of the first-moment lemma the volume factor $N^4$ cancels only
because
\begin{equation*}
  \mathbb{E}[\Theta]\le 2\bigl(\mathsf{B}N^4+\log c_J\bigr),
\end{equation*}
a bound linear in the volume.

\emph{What the age series requires.} The age series for the never-healed old regions requires a bound on
\begin{equation}\label{8.9:eq-first-moment-blind-2}
  \frac{1}{n_{k_0}}\,\mathbb{E}_{\tilde{\nu}}\Bigl[\,\sum_{\mathfrak{c}}
  \mathbf{1}_{\mathcal{E}^{(k_0)}_{\mathfrak{c}}(K')}\,
  L^{8\cdot\operatorname{age}_{k_0}(\mathfrak{c})}\Bigr],
\end{equation}
where $n_{k_0}$ is the normalising factor of the age series, $\mathbf{1}$ denotes an indicator, and the
sum runs over the level-$k_0$ cells $\mathfrak{c}$.
By the age lemma (Lemma~\ref{8.9:lem-age-lemma}), whose proof is incomplete at one step, the seed age of
a live component is bounded by the sum over its seeds of the per-seed charges $\mathfrak{a}(Z)$ of
\eqref{8.9:eq-age-lemma-a}. This bound is sharp for every kind of seed that the age lemma treats: a chain
of successive re-creations realises an age of the order of the number of re-creations times
Ba{\l}aban's number $\check{R}_j$ at the levels $j$ of the chain. Hence
$L^{8\cdot\operatorname{age}(X)}$ can be as large as
\begin{equation*}
  L^{8\sum_{Z\in X}\mathfrak{a}(Z)}=\prod_{Z\in X} e^{8\ln L\cdot\mathfrak{a}(Z)}.
\end{equation*}
This quantity is exponential in the seed structure of $X$, and it is not bounded by $\Theta(X)/t$. The
summand of \eqref{8.9:eq-first-moment-blind-2} is therefore not of the form that the first-moment lemma
controls.

\emph{Taming the weight by the reserve returns the volume to the exponent.} Suppose the weight is
dominated by the reserve, $L^{8\cdot\operatorname{age}}\le e^{\eta\Theta(X)}$, under a ceiling condition.
Then \eqref{8.9:eq-first-moment-blind-2} leads to
\begin{equation*}
  \mathbb{E}_{\tilde{\nu}}\Bigl[\,\sum_{X}\Theta(X)\,e^{\eta\Theta(X)}\Bigr]
  \quad\text{or}\quad
  \mathbb{E}\bigl[e^{\eta\Theta}\bigr].
\end{equation*}
By the exponential-moment clause of the first-moment lemma, the two hypotheses of the first-moment lemma
together bound these expectations only by quantities of the type $e^{\mathsf{B}N^4}$. In such a bound
the volume appears in the exponent.

Higher moments, of order $p>1$, do not help either. They satisfy
\begin{equation*}
  \mathbb{E}\bigl[\Theta^p\bigr]\lesssim \bigl(2\mathsf{B}N^4\bigr)^p ,
\end{equation*}
and after division by $n_{k_0}$ they leave a factor $N^{4(p-1)}$.

\emph{Counting old seeds instead of old cells.} Count old seeds at their birth level $j_0=k_0-a$ instead
of old cells. This does give a density that is free of the age. The position entropy $L^{4a}$ of a
seed of level $j_0$ under a level-$k_0$ cell cancels against the depth gain $L^{-4a}$ carried by the
seed factor $\mathfrak{e}_{j_0}$. The resulting density, however, is age-free, not age-decaying. Moreover,
tying a counted seed to the component of $\square$ is a joint event.

\emph{Conclusion.} The sharpness clause of the first-moment lemma is thus met directly. From positivity,
the two hypotheses of the first-moment lemma, a reserve per seed and symmetry, no bound on
\eqref{8.9:eq-first-moment-blind-1} better than age-blind follows. The obstruction lies in the weight
$L^{8\cdot\operatorname{age}}$ and not in the first-moment lemma.
\end{proof}

\begin{definition}[The terminal lattice, transports and plaquette defects]\label{8.10:not-terminal-lattice}
Throughout this section $d\ge 3$ denotes the dimension. In the application $d=4$, and we keep the
letter $d$ in the counts. We put $c_d:=d(d-1)/2$.

\emph{Lattice and cubes.} Let $\mathfrak{N}\ge 3$ be an integer. The \emph{terminal unit torus of
side $\mathfrak{N}$} is $\mathsf{T}:=(\mathbb{Z}/\mathfrak{N}\mathbb{Z})^d$, and its elements are
\emph{sites}. The vectors $e_\mu$, $1\le\mu\le d$, are the unit vectors. A \emph{$k$-cube} is a set
\begin{equation*}
\Delta=y+\{0,1\}^J:=\Bigl\{\,y+\sum_{\mu\in J}t_\mu e_\mu : t\in\{0,1\}^J\Bigr\},
\qquad y\in\mathsf{T},\ J\subset\{1,\dots,d\},\ |J|=k .
\end{equation*}
Its \emph{faces} are the cubes contained in it. The $0$-cubes are the sites, the $1$-cubes are the
\emph{bonds} and the $2$-cubes are the \emph{plaquettes}. A \emph{cubical complex} is a finite set
$\Xi$ of cubes that is closed under taking faces. We write $\Xi_k$ for its set of $k$-cubes, and
$(\Xi_0,\Xi_1)$ is its graph. For a set $Y$ of sites, $\mathcal{F}(Y)$ is the \emph{full complex}
on $Y$, that is, the set of all cubes $\Delta\subset Y$.

\emph{Distances.} For sites $z,z'$ and finite sets of sites $Y_1,Y_2$ put
\begin{align*}
|z-z'|_\infty&:=\max_{\mu}\,\min\{|a| : a\in\mathbb{Z},\ a\equiv z_\mu-z'_\mu \ \mathrm{mod}\ \mathfrak{N}\},\\
\operatorname{dist}_\infty(Y_1,Y_2)&:=\min\{|y_1-y_2|_\infty : y_1\in Y_1,\ y_2\in Y_2\}.
\end{align*}
When distances are taken, a cube is identified with its vertex set. For a plaquette $p$ and an integer
$r\ge 0$ let $\mathcal{U}_r(p)$ be the set of sites $z$ with $|z-y|_\infty\le r$ for some vertex $y$
of $p$. We say that a cube $\Delta$ \emph{lies in} $\mathcal{U}_r(p)$, and write
$\Delta\subset\mathcal{U}_r(p)$, if every vertex of $\Delta$ does.

\emph{Group and norms.} Let $N\ge 2$. The group is $G:=\mathrm{SU}(N)$, with Lie algebra
$\mathfrak{g}$, $\dim\mathfrak{g}=N^2-1$, and $\|\cdot\|$ is the operator norm on
$\mathbb{C}^{N\times N}$. For unitary $\Phi,\Psi,\Gamma$ and unitaries $\Phi_1,\dots,\Phi_n$,
$\Psi_1,\dots,\Psi_n$ one has
\begin{align}
\|\Phi\Psi-1\|&\le\|\Phi-1\|+\|\Psi-1\|, \label{8.10:eq-terminal-lattice-1}\\
\|\Gamma\Phi\Gamma^{-1}-1\|&=\|\Phi-1\|,\qquad \|\Phi^{-1}-1\|=\|\Phi-1\|,
\label{8.10:eq-terminal-lattice-2}\\
\|\Phi_1\cdots\Phi_n-\Psi_1\cdots\Psi_n\|&\le\sum_{i=1}^n\|\Phi_i-\Psi_i\| .
\label{8.10:eq-terminal-lattice-3}
\end{align}

\emph{Configurations and transports.} Let $\mathcal{V}$ be the product of copies of $\mathrm{SU}(N)$
indexed by the bonds of $\mathsf{T}$, and let $dV$ be the product of the Haar
probability measures. For a bond $b=\{y,y+e_\mu\}$ the variable $V_b$ is attached to the
orientation $(y,y+e_\mu)$, and we put $V_{(y+e_\mu,y)}:=V_b^{-1}$. An \emph{edge path} is a sequence
$\sigma=(y_0,\dots,y_n)$ of sites in which consecutive sites are joined by bonds. Its
\emph{transport} is
\begin{equation*}
T_\sigma(V):=V_{(y_0,y_1)}\cdots V_{(y_{n-1},y_n)},
\end{equation*}
and the empty product is $1$. The path is a \emph{loop at} $y_0$ if $y_n=y_0$. Transport is
multiplicative under concatenation of edge paths. For a plaquette $p=y+\{0,1\}^{\{\mu,\nu\}}$ with
$\mu<\nu$, $\partial p$ is the loop $(y,\,y+e_\mu,\,y+e_\mu+e_\nu,\,y+e_\nu,\,y)$. We write
$V(\partial p):=T_{\partial p}(V)$, and the \emph{plaquette defect} is
\begin{equation}\label{8.10:eq-terminal-lattice-4}
F_p(V):=\|V(\partial p)-1\|\in[0,2].
\end{equation}
This number does not depend on the corner at which the loop is started or on its orientation.

\emph{Gauge transformations.} A gauge transformation $\lambda:\mathsf{T}\to\mathrm{SU}(N)$
acts by $(V^\lambda)_{(y,z)}:=\lambda(y)V_{(y,z)}\lambda(z)^{-1}$. For every edge path $\sigma$ and
every plaquette $p$,
\begin{equation}\label{8.10:eq-terminal-lattice-5}
T_\sigma(V^\lambda)=\lambda(y_0)\,T_\sigma(V)\,\lambda(y_n)^{-1},\qquad F_p(V^\lambda)=F_p(V),
\end{equation}
and the map $V\mapsto V^\lambda$ preserves $dV$.

\emph{Cells.} Let $\mathfrak{s}\ge 2$ be an integer dividing $\mathfrak{N}$. Put
$\mathfrak{N}':=\mathfrak{N}/\mathfrak{s}$. The
cells are indexed by $c\in\mathsf{C}:=(\mathbb{Z}/\mathfrak{N}'\mathbb{Z})^d$, and we put
\begin{equation*}
S(c):=\mathfrak{s}c+\{0,\dots,\mathfrak{s}-1\}^d,\qquad
\bar S(c):=\mathfrak{s}c+\{0,\dots,\mathfrak{s}\}^d,\qquad
\overline{\mathcal{Q}}_c:=\mathcal{F}(\bar S(c)).
\end{equation*}
The sets $S(c)$, $c\in\mathsf{C}$, partition $\mathsf{T}$. For $X\subset\mathsf{C}$ put
\begin{equation*}
S(X):=\bigcup_{c\in X}S(c),\quad \bar S(X):=\bigcup_{c\in X}\bar S(c),\quad
\overline{\mathcal{F}}(X):=\bigcup_{c\in X}\overline{\mathcal{Q}}_c,\quad n(X):=|X|.
\end{equation*}
Two cells \emph{touch} if their indices differ by a non-zero element of $\{-1,0,1\}^d$ modulo
$\mathfrak{N}'$. The distance $\operatorname{dist}_{\blacksquare}$ is the sup-distance of indices on
the index torus:
\begin{equation*}
\operatorname{dist}_{\blacksquare}(c,c'):=\max_\mu\,\min\{|a| : a\in\mathbb{Z},\
a\equiv c_\mu-c'_\mu\ \mathrm{mod}\ \mathfrak{N}'\},
\end{equation*}
and $\operatorname{dist}_{\blacksquare}(c,X):=\min_{c'\in X}\operatorname{dist}_{\blacksquare}(c,c')$.
A set $X$ is \emph{touch-connected} if the touch-graph on $X$ is connected. For an integer $q\ge 0$
put $\operatorname{hull}_q(X):=\{c : \operatorname{dist}_{\blacksquare}(c,X)\le q\}$. Then
\begin{equation}\label{8.10:eq-terminal-lattice-6}
\operatorname{hull}_q(\operatorname{hull}_{q'}(X))=\operatorname{hull}_{q+q'}(X),
\end{equation}
and we have the counts
\begin{equation}\label{8.10:eq-terminal-lattice-7}
n(\operatorname{hull}_q(X))\le(2q+1)^d\,n(X),\qquad |\bar S(X)|\le(\mathfrak{s}+1)^d\,n(X).
\end{equation}
The complex $\overline{\mathcal{F}}(X)$ has at most $d(\mathfrak{s}+1)^dn(X)$ bonds and at most
$c_d(\mathfrak{s}+1)^dn(X)$ plaquettes. Finally,
$\bar S(c)\subset S(\operatorname{hull}_1(\{c\}))$.
\end{definition}

\begin{proof}
We verify the assertions made above.

\emph{Well-definedness of the orientation.} Because $\mathfrak{N}\ge 3$, a bond $b$ determines the
site $y$ and the direction $\mu$ with $b=\{y,y+e_\mu\}$. Suppose that also $b=\{y',y'+e_\nu\}$. If
$y'=y$, then $e_\mu\equiv e_\nu$ modulo $\mathfrak{N}$, and hence $\mu=\nu$. Otherwise $y'=y+e_\mu$
and $y'+e_\nu=y$. Then $e_\mu+e_\nu\equiv 0$ modulo $\mathfrak{N}$, which says $2\equiv 0$ if
$\mu=\nu$ and $1\equiv0$ if $\mu\neq\nu$; both are excluded for $\mathfrak{N}\ge3$. So the
orientation $(y,y+e_\mu)$ attached to $V_b$ is well defined.

\emph{The norm relations.} Since $\|\Phi\|=1$,
\begin{equation*}
\|\Phi\Psi-1\|=\|\Phi(\Psi-1)+(\Phi-1)\|\le\|\Psi-1\|+\|\Phi-1\|,
\end{equation*}
which is \eqref{8.10:eq-terminal-lattice-1}. Next, $\Gamma\Phi\Gamma^{-1}-1=\Gamma(\Phi-1)\Gamma^{-1}$,
and the operator norm is invariant under multiplication by unitaries on either side. This gives the
first identity of \eqref{8.10:eq-terminal-lattice-2}. Likewise $\Phi^{-1}-1=\Phi^{-1}(1-\Phi)$ gives
the second. For \eqref{8.10:eq-terminal-lattice-3} we use the telescoping identity
\begin{equation*}
\Phi_1\cdots\Phi_n-\Psi_1\cdots\Psi_n=\sum_{i=1}^n\Phi_1\cdots\Phi_{i-1}(\Phi_i-\Psi_i)\Psi_{i+1}\cdots\Psi_n .
\end{equation*}
In this identity each product of unitaries has norm one.

\emph{Range of $F_p$.} The matrix $V(\partial p)$ is unitary, so
$\|V(\partial p)-1\|\le\|V(\partial p)\|+1=2$. This gives
\eqref{8.10:eq-terminal-lattice-4}.

\emph{Multiplicativity.} Let $\sigma=(y_0,\dots,y_n)$ and $\sigma'=(y_n,\dots,y_{n+n'})$, and let
$\sigma\sigma'$ be their concatenation. Then $T_{\sigma\sigma'}(V)=T_\sigma(V)T_{\sigma'}(V)$ by the
definition of $T$ as an ordered product.

\emph{Independence of corner and orientation.} Write
$V(\partial p)=\Phi_1\Phi_2\Phi_3\Phi_4$, where the $\Phi_i$ are the four oriented bond variables
of $\partial p$. Starting the loop at another corner replaces this product by a cyclic permutation
$\Phi_{j+1}\cdots\Phi_4\Phi_1\cdots\Phi_j$, and
\begin{equation*}
\Phi_{j+1}\cdots\Phi_4\Phi_1\cdots\Phi_j=\Gamma^{-1}\,V(\partial p)\,\Gamma,
\qquad \Gamma:=\Phi_1\cdots\Phi_j .
\end{equation*}
Reversing an edge path $\sigma=(y_0,\dots,y_n)$ to $\bar\sigma=(y_n,\dots,y_0)$ gives, because
$V_{(z,y)}=V_{(y,z)}^{-1}$,
\begin{equation*}
T_{\bar\sigma}(V)=V_{(y_{n-1},y_n)}^{-1}\cdots V_{(y_0,y_1)}^{-1}=T_\sigma(V)^{-1}.
\end{equation*}
So a change of corner or of orientation, or both, replaces $V(\partial p)$ by a conjugate of itself
or of its inverse. By \eqref{8.10:eq-terminal-lattice-2} the norm $\|V(\partial p)-1\|$ is
unchanged.

\emph{Gauge transformations.} The rule $(V^\lambda)_{(y,z)}=\lambda(y)V_{(y,z)}\lambda(z)^{-1}$ is
consistent with $V_{(z,y)}=V_{(y,z)}^{-1}$, since
\begin{equation*}
\lambda(z)V_{(y,z)}^{-1}\lambda(y)^{-1}=\bigl(\lambda(y)V_{(y,z)}\lambda(z)^{-1}\bigr)^{-1}.
\end{equation*}
In the product $T_\sigma(V^\lambda)$ the interior factors $\lambda(y_i)^{-1}\lambda(y_i)$,
$0<i<n$, cancel. This leaves $\lambda(y_0)T_\sigma(V)\lambda(y_n)^{-1}$, the first identity of
\eqref{8.10:eq-terminal-lattice-5}. For a loop $\partial p$ at $y$ it gives
$V^\lambda(\partial p)=\lambda(y)V(\partial p)\lambda(y)^{-1}$, and the second identity follows from
\eqref{8.10:eq-terminal-lattice-2}. The map $V\mapsto V^\lambda$ acts bond by bond: on the factor of
the bond $\{y,y+e_\mu\}$ it acts by $V_b\mapsto\lambda(y)V_b\lambda(y+e_\mu)^{-1}$, a left and a right
multiplication by fixed elements of $\mathrm{SU}(N)$. Haar probability measure on the compact group
$\mathrm{SU}(N)$ is invariant under both. Hence each factor of the product measure is preserved,
and so is $dV$.

\emph{The cells partition the torus.} Since $\mathfrak{s}\mathfrak{N}'=\mathfrak{N}$, the site
$\mathfrak{s}c$ is well defined for $c\in\mathsf{C}$. Every residue $z_\mu$ modulo $\mathfrak{N}$ can
be written in exactly one way as $\mathfrak{s}c_\mu+\xi_\mu$ with $c_\mu$ a residue modulo
$\mathfrak{N}'$ and $\xi_\mu\in\{0,\dots,\mathfrak{s}-1\}$. Hence each site lies in exactly one
$S(c)$.

\emph{Composition of hulls.} The inclusion of the left side of \eqref{8.10:eq-terminal-lattice-6}
in the right side follows from the triangle inequality for $\operatorname{dist}_{\blacksquare}$.
That inequality holds coordinatewise, because representatives of residues add. For the converse, let
$\operatorname{dist}_{\blacksquare}(c,c')\le q+q'$ with $c'\in X$. Choose for each $\mu$ a
representative $a_\mu\equiv c_\mu-c'_\mu$ with $|a_\mu|\le q+q'$, and put
$\eta_\mu:=\operatorname{sign}(a_\mu)\min\{|a_\mu|,q'\}$. Then $|\eta_\mu|\le q'$ and
$|a_\mu-\eta_\mu|\le q$. So the cell $c'+\eta$ satisfies
$\operatorname{dist}_{\blacksquare}(c'+\eta,c')\le q'$ and
$\operatorname{dist}_{\blacksquare}(c,c'+\eta)\le q$. Hence
$c\in\operatorname{hull}_q(\operatorname{hull}_{q'}(X))$.

\emph{Counts.} We have
\begin{equation*}
\operatorname{hull}_q(X)=\bigcup_{c'\in X}\{c:\operatorname{dist}_{\blacksquare}(c,c')\le q\}.
\end{equation*}
Every $c$ in the set indexed by $c'$ is of the form $c'+a$ with $a\in\{-q,\dots,q\}^d$, so that set
has at most $(2q+1)^d$ elements. This gives the first bound of \eqref{8.10:eq-terminal-lattice-7}.
The set $\bar S(c)$ is the image of $\{0,\dots,\mathfrak{s}\}^d$, so
$|\bar S(c)|\le(\mathfrak{s}+1)^d$. Summing over $c\in X$ gives the second bound.

Every bond contained in $\bar S(c)$ has the form $y+\{0,1\}^{\{\mu\}}$ with $y\in\bar S(c)$ and
$1\le\mu\le d$. So $\overline{\mathcal{Q}}_c$ has at most $d\,|\bar S(c)|\le d(\mathfrak{s}+1)^d$
bonds. Similarly, every plaquette contained in $\bar S(c)$ has the form
$y+\{0,1\}^{\{\mu,\nu\}}$ with $y\in\bar S(c)$ and $\mu<\nu$, and there are $c_d$ such pairs. So
$\overline{\mathcal{Q}}_c$ has at most $c_d(\mathfrak{s}+1)^d$ plaquettes. Since
$\overline{\mathcal{F}}(X)$ is the union of the $\overline{\mathcal{Q}}_c$, $c\in X$, summing over
$c\in X$ gives the bounds on its bonds and plaquettes.

\emph{The closed cell lies in the hull.} Let $z=\mathfrak{s}c+t$ with $t\in\{0,\dots,\mathfrak{s}\}^d$.
For each $\mu$ put $\varepsilon_\mu:=1$ if $t_\mu=\mathfrak{s}$ and $\varepsilon_\mu:=0$ otherwise, and
put $\xi_\mu:=t_\mu-\mathfrak{s}\varepsilon_\mu\in\{0,\dots,\mathfrak{s}-1\}$. Then
$z=\mathfrak{s}(c+\varepsilon)+\xi\in S(c+\varepsilon)$. Since $\varepsilon\in\{0,1\}^d$, we have
$\operatorname{dist}_{\blacksquare}(c+\varepsilon,c)\le1$, so
$z\in S(\operatorname{hull}_1(\{c\}))$.
\end{proof}

\begin{lemma}[Fullness of closed-cube complexes]\label{8.10:lem-cube-fullness}
Let $\mathbb{T}$ be the terminal lattice of Notation~\ref{8.10:not-terminal-lattice}, periodic of
side $\mathfrak{N}$, and let $\ell\ge 2$ be an integer dividing $\mathfrak{N}$;
put $\mathfrak{n}:=\mathfrak{N}/\ell$, and index cells by
$c\in\mathsf{C}:=(\mathbb{Z}/\mathfrak{n}\mathbb{Z})^d$. Put
$\overline{S}(c):=\ell c+\{0,\dots,\ell\}^d$ and, for $X\subset\mathsf{C}$,
$\overline{S}(X):=\bigcup_{c\in X}\overline{S}(c)$. A \emph{cube} is a set
$y+\{0,1\}^J$ with $y\in\mathbb{T}$ and $J\subset\{1,\dots,d\}$, namely the set of sites
$y+\sum_{\mu\in J}\epsilon_\mu e_\mu$ with $\epsilon_\mu\in\{0,1\}$; these sites are its
vertices, and $e_\mu$ is the unit vector in direction $\mu$. For a set $Y$ of sites,
$\mathsf{K}(Y)$ is the set of cubes all of whose vertices lie in $Y$. Put
$\overline{Q}_c:=\mathsf{K}(\overline{S}(c))$ and
$\overline{\mathsf{K}}(X):=\bigcup_{c\in X}\overline{Q}_c$. Then for every
$X\subset\mathsf{C}$
\begin{equation}\label{8.10:eq-cube-fullness-1}
\overline{\mathsf{K}}(X)=\mathsf{K}\bigl(\overline{S}(X)\bigr);
\end{equation}
that is, a cube all of whose vertices lie in $\overline{S}(X)$ lies in $\overline{Q}_c$ for
some $c\in X$.
\end{lemma}

\begin{proof}
Coordinates $z_\mu$ of sites are represented in $\{0,\dots,\mathfrak{N}-1\}$, and cell
indices are read modulo $\mathfrak{n}$. For each $c\in X$ we have
$\overline{S}(c)\subset\overline{S}(X)$, so
$\overline{Q}_c\subset\mathsf{K}(\overline{S}(X))$. This is the inclusion
$\overline{\mathsf{K}}(X)\subset\mathsf{K}(\overline{S}(X))$, and it remains to prove the
reverse one.

Let $\Gamma=y+\{0,1\}^J$ be a cube all of whose vertices lie in $\overline{S}(X)$.

\emph{Sites.} For a site $z$ and a direction $\mu$ put
\begin{equation*}
\mathcal{I}_\mu(z):=\{\lfloor z_\mu/\ell\rfloor\}\cup\{z_\mu/\ell-1:\ \ell\mid z_\mu\}.
\end{equation*}
The coordinate $z_\mu$ lies in $[\ell c_\mu,\ell c_\mu+\ell]$ exactly when either
$z_\mu\in[\ell c_\mu,\ell c_\mu+\ell-1]$, that is $c_\mu=\lfloor z_\mu/\ell\rfloor$, or
$z_\mu=\ell c_\mu+\ell$, that is $\ell\mid z_\mu$ and $c_\mu=z_\mu/\ell-1$. Hence
$z\in\overline{S}(c)$ if and only if $c_\mu\in\mathcal{I}_\mu(z)$ for every $\mu$.

\emph{Cubes.} We have $\Gamma\subset\overline{S}(c)$ if and only if every vertex of $\Gamma$
lies in $\overline{S}(c)$. For $\mu\in J$ the vertices of $\Gamma$ take both values $y_\mu$
and $y_\mu+1$ in coordinate $\mu$, and both must lie in $[\ell c_\mu,\ell c_\mu+\ell]$
(recall $\ell\ge2$). This excludes $y_\mu=\ell c_\mu+\ell$, so
$y_\mu\in[\ell c_\mu,\ell c_\mu+\ell-1]$, i.e.\ $c_\mu=\lfloor y_\mu/\ell\rfloor$.
Conversely, this value of $c_\mu$ gives $y_\mu$ and $y_\mu+1$ in the interval. For
$\mu\notin J$ all vertices have coordinate $y_\mu$, and the condition is
$c_\mu\in\mathcal{I}_\mu(y)$. Thus $\Gamma\subset\overline{S}(c)$ if and only if
$c_\mu=\lfloor y_\mu/\ell\rfloor$ for $\mu\in J$ and $c_\mu\in\mathcal{I}_\mu(y)$ for
$\mu\notin J$.

\emph{A distinguished vertex.} For $\mu\in J$ put $\delta_\mu:=1$ if $\ell\mid y_\mu$ and
$\delta_\mu:=0$ otherwise. Let $v:=y+\sum_{\mu\in J}\delta_\mu e_\mu$, which is a vertex of
$\Gamma$.

Take $\mu\in J$. If $\ell\mid y_\mu$, then $v_\mu=y_\mu+1$ is not divisible by $\ell$,
since $\ell\ge2$. Hence
\begin{equation*}
\mathcal{I}_\mu(v)=\{\lfloor (y_\mu+1)/\ell\rfloor\}=\{\lfloor y_\mu/\ell\rfloor\}.
\end{equation*}
If $\ell\nmid y_\mu$, then $v_\mu=y_\mu$ and again
$\mathcal{I}_\mu(v)=\{\lfloor y_\mu/\ell\rfloor\}$.

For $\mu\notin J$ we have $v_\mu=y_\mu$, and so $\mathcal{I}_\mu(v)=\mathcal{I}_\mu(y)$.

Comparing with the two characterisations above, $v\in\overline{S}(c)$ if and only if
$\Gamma\subset\overline{S}(c)$.

\emph{Conclusion.} Since $v$ is a vertex of $\Gamma$, it lies in
$\overline{S}(X)=\bigcup_{c\in X}\overline{S}(c)$. Hence $v\in\overline{S}(c)$ for some
$c\in X$, and for this $c$ we get $\Gamma\subset\overline{S}(c)$. This means
\begin{equation*}
\Gamma\in\mathsf{K}(\overline{S}(c))=\overline{Q}_c\subset\overline{\mathsf{K}}(X),
\end{equation*}
which proves \eqref{8.10:eq-cube-fullness-1}.
\end{proof}

\begin{lemma}[Lattice steps and cell distance]\label{8.10:lem-steps-cells}
We use the notation of Notation~\ref{8.10:not-terminal-lattice}: the cells $S(c)$ of side $\ell$
of the terminal lattice, their indices $c$ read modulo $\mathfrak{n}$, the cell distance
$\operatorname{dist}_{\blacksquare}$, and the norm $|\cdot|_\infty$. Let $k\ge 0$ be
an integer, and let $z$ and $w$ be points with $|z-w|_\infty\le k\ell$ and $w\in S(c)$. Then
$z\in S(c'')$ for some cell index $c''$ with $\operatorname{dist}_{\blacksquare}(c'',c)\le k$.
\end{lemma}

\begin{proof}
The argument is carried out one coordinate at a time. Fix a direction $\mu$.

Since $|z-w|_\infty\le k\ell$, we may choose a representative $a$ of $z_\mu-w_\mu$ with $|a|\le
k\ell$. Since $w\in S(c)$, the coordinate $w_\mu$ has a representative in the interval
$[\ell c_\mu,\ \ell c_\mu+\ell-1]$; we write $w_\mu$ also for this representative. Adding $a$ to
it gives
\begin{equation*}
\ell(c_\mu-k)\;\le\; w_\mu+a\;\le\; \ell(c_\mu+k)+\ell-1 .
\end{equation*}
Put $c''_\mu=\lfloor (w_\mu+a)/\ell\rfloor$. Dividing the last display by $\ell$ and taking
integer parts yields the integer inequalities
\begin{equation*}
c_\mu-k\;\le\; c''_\mu\;\le\; c_\mu+k .
\end{equation*}
By the definition of the integer part,
\begin{equation*}
\ell c''_\mu\;\le\; w_\mu+a\;\le\; \ell c''_\mu+\ell-1 .
\end{equation*}
The number $w_\mu+a$ is a representative of $z_\mu$. Hence the $\mu$-th coordinate of $z$
satisfies the defining condition of the cell of index $c''$.

We now let $\mu$ range over all directions and read each $c''_\mu$ modulo $\mathfrak{n}$. This
defines a cell index $c''$ with $z\in S(c'')$. For every $\mu$, the coordinate $c''_\mu$ lies in
$[c_\mu-k,\ c_\mu+k]$, read modulo $\mathfrak{n}$. This is the bound
$\operatorname{dist}_{\blacksquare}(c'',c)\le k$.
\end{proof}

We use the sets $\bar S(c)$ and $\bar Q_c$ attached to a cell index $c$, the set $\bar K(X)$
attached to a region $X$, and the notion of touching cells. If the region $X$ is
touch-connected, then the graph of $\bar K(X)$ is connected. The reason
is twofold: each $\bar Q_c$ has a connected graph, and $\bar S(c)\cap\bar S(c')\neq\emptyset$
whenever the cells $c$ and $c'$ touch.

\begin{definition}[Lifting shapes from the torus to the lattice]\label{8.10:def-torus-lifting}
We use the set $\mathsf{C}$ of cells, the number $\mathfrak{n}$ of residue classes of cell indices
in each direction, the cell side $\ell$, the side $\mathfrak{N}$ of the terminal torus, and, for a
set $Y\subset\mathsf{C}$ of cells, the vertex set $\bar{S}(Y)$, the complex $\bar{K}(Y)$ of closed
cubes and the closed cubes $\bar{Q}_c$, $c\in Y$, in the notation of
Notation~\ref{8.10:not-terminal-lattice} and Lemma~\ref{8.10:lem-cube-fullness}. The $\mu$-th
index of a cell $c$ is a residue class $c_\mu$ mod $\mathfrak{n}$. A set $X\subset\mathsf{C}$ is
called a shape.

A shape $X$ satisfies the \emph{lifting condition} if in every direction $\mu$ some residue class
mod $\mathfrak{n}$ contains no $\mu$-th index of a cell of $X$.

Let $X$ satisfy the lifting condition. For each $\mu$ fix an integer $o_\mu$ whose class mod
$\mathfrak{n}$ contains no $\mu$-th index of a cell of $X$; the $\mu$-th indices of $X$ then lie in
the $\mathfrak{n}-1$ consecutive classes of $o_\mu+1,\dots,o_\mu+\mathfrak{n}-1$. The
\emph{lifted} $\mu$-th index of $c\in X$ is the unique integer of
$\{o_\mu+1,\dots,o_\mu+\mathfrak{n}-1\}$ in the class $c_\mu$.
Let $a_\mu$ and $b_\mu$ be the least and the largest lifted $\mu$-th index of a cell of $X$, so
that $b_\mu-a_\mu\le\mathfrak{n}-2$, and let
\begin{equation*}
B(X):=\prod_\mu\{a_\mu,\dots,b_\mu\}
\end{equation*}
be the box of cells spanned coordinatewise by $X$ in lifted indices, regarded as a set of cells by
reduction of the indices mod $\mathfrak{n}$, which is injective on $B(X)$ since
$b_\mu-a_\mu\le\mathfrak{n}-2$. The box $B(X)$ depends on the choice of the integers $o_\mu$; one
such choice is fixed for each shape. Then $\bar{S}(B(X))$ is the reduction mod $\mathfrak{N}$ of
the set of integer points
\begin{equation}\label{8.10:eq-torus-lifting-1}
\prod_\mu\{\ell a_\mu,\dots,\ell b_\mu+\ell\}\subset\mathbb{Z}^d,
\end{equation}
reduction mod $\mathfrak{N}$ is injective on the set \eqref{8.10:eq-torus-lifting-1},
$\bar{S}(B(X))\supset\bar{S}(X)$, and the complexes $\bar{K}(B(X))\supset\bar{K}(X)$ lift
injectively to $\mathbb{Z}^d$, each cube to a closed unit cube of $\mathbb{Z}^d$. For a complex $K$
lifted in this way, $|K|\subset\mathbb{R}^d$ denotes the union of the closed unit cubes of $K$.

The lifting condition restricts the position of a shape relative to the torus; it bounds none of
the parameters.
\end{definition}

\begin{proof}[Proof of the assertions in Definition~\ref{8.10:def-torus-lifting}]
Fix $\mu$. The set \eqref{8.10:eq-torus-lifting-1} consists in direction $\mu$ of
$\ell(b_\mu-a_\mu)+\ell+1$ consecutive integers, and since $b_\mu-a_\mu\le\mathfrak{n}-2$,
\begin{equation*}
\ell(b_\mu-a_\mu)+\ell+1\le\ell(\mathfrak{n}-1)+1=\mathfrak{N}-\ell+1\le\mathfrak{N}-1,
\end{equation*}
where the equality is $\mathfrak{N}=\ell\mathfrak{n}$ and the last inequality holds because the cell
side $\ell$ of Notation~\ref{8.10:not-terminal-lattice} satisfies $\ell\ge2$. Distinct integers in
a set of at most $\mathfrak{N}-1$ consecutive integers are distinct
mod $\mathfrak{N}$, so reduction mod $\mathfrak{N}$ maps the product
\eqref{8.10:eq-torus-lifting-1} injectively onto $\bar{S}(B(X))$ in the torus.

Every cell of $X$ has its lifted indices in $\{a_\mu,\dots,b_\mu\}$, so $X\subset B(X)$; hence
$\bar{S}(X)\subset\bar{S}(B(X))$ and $\bar{K}(X)\subset\bar{K}(B(X))$.

By Lemma~\ref{8.10:lem-cube-fullness}, $\bar{K}(B(X))=K(\bar{S}(B(X)))$, so every cube of
$\bar{K}(B(X))$ has all its vertices in $\bar{S}(B(X))$. Lift these vertices into
\eqref{8.10:eq-torus-lifting-1} by the inverse of the reduction map. Two vertices of the cube whose
$\mu$-th coordinates differ by $1$ mod $\mathfrak{N}$ have lifted $\mu$-th coordinates whose
difference is congruent to $\pm1$ mod $\mathfrak{N}$. Both lie in an interval of at most
$\mathfrak{N}-1$ consecutive integers, so the difference has absolute value at most
$\mathfrak{N}-2$. Of the integers congruent to $\pm1$ mod $\mathfrak{N}$, only $\pm1$ have
absolute value at most $\mathfrak{N}-2$, and therefore the lifted coordinates differ by exactly
$1$. Coordinates
that are equal in the torus have equal lifts, by injectivity. The lifted vertices are therefore
the vertices of a unit cube of $\mathbb{Z}^d$, and distinct vertices have distinct lifts; that is,
the cube lifts injectively to $\mathbb{Z}^d$. Since all vertices are lifted by the inverse of one
injective map, the lifts of the cubes of $\bar{K}(B(X))$, and of its subcomplex $\bar{K}(X)$,
together form an injective lift of these complexes to $\mathbb{Z}^d$.
\end{proof}

\begin{definition}[Edge-loop homotopy and the fundamental group]\label{8.10:def-edge-loop-homotopy}
Let $X$ be a cubical complex, whose cells are cubes of dimensions $0,\dots,d$. We write $X_j$ for
the set of its $j$-dimensional cubes, so that $X_0$, $X_1$, $X_2$ are its vertices, edges and
plaquettes. The \emph{graph of $X$} is the graph with vertex set $X_0$ and edge set $X_1$, and
$|X|$ denotes the realisation of $X$, the union of its cubes.

Let $v\in X_0$. A \emph{loop at $v$} in the graph of $X$ is a sequence
$\gamma=(y_0,y_1,\dots,y_n)$ of vertices with $y_0=y_n=v$ in which consecutive vertices are the
endpoints of an edge of $X_1$. The constant loop is $(v)$, with $n=0$. More generally, an
\emph{edge path} from $y$ to $y'$ is a sequence of vertices from $y$ to $y'$ in which consecutive
vertices are the endpoints of an edge of $X_1$; $\overline{\alpha}$ denotes the reverse of an edge
path $\alpha$, and juxtaposition denotes concatenation. The loop $\gamma$
determines the path in $|X|$ that traverses its edges in order at constant speed. For a
plaquette $P\in X_2$ with corner $y$, a \emph{boundary circuit of $P$ at $y$} is the loop
$(y,a_1,a_2,a_3,y)$ that runs once around the four edges of $P$ in either orientation.

Two loops at $v$ are \emph{combinatorially homotopic in $X$} if one arises from the other by
finitely many of the following moves:
\begin{itemize}
\item[(E1)] insertion or deletion of a backtrack, $(\dots,y,z,y,\dots)\leftrightarrow(\dots,y,\dots)$;
\item[(E2)] insertion or deletion, at a vertex $y$ of the loop, of a boundary circuit at $y$ of a
plaquette $P\in X_2$ having $y$ as a corner (either orientation).
\end{itemize}
Concatenation of loops at $v$ is compatible with these moves. The classes therefore form a group
$\pi_1^{\square}(X,v)$ under concatenation. Its identity is the class of $(v)$. The inverse of the
class of a loop is the class of the reversed loop: the concatenation of a loop of length $n$ with
its reverse is reduced to $(v)$ by $n$ successive deletions of backtracks.

A loop is \emph{null-homotopic in $X$} if it is combinatorially homotopic in $X$ to $(v)$.
Suppose $\gamma$ is reduced to the constant loop by moves (E1) together with (E2)-moves across
plaquettes $q_1,\dots,q_A$, listed with repetition. Such a reduction is called a
\emph{null-homotopy of $\gamma$ across $(q_1,\dots,q_A)$}, and $A$ is its \emph{area}.

Let $V$ be a field given at least on $X_1$, with plaquette transports $V(\partial P)$ and
plaquette defects $F_P(V)$ as in Notation~\ref{8.10:not-terminal-lattice}. Then $V$ is
\emph{flat on $X$} if $V(\partial P)=1$ for every $P\in X_2$. For a number $\delta$, $V$ is
\emph{$\delta$-flat on $X$} if $F_P(V)\le\delta$ for every $P\in X_2$.

The following holds and is proved below. For a finite cubical complex $X\subset\mathbb{Z}^d$ and
$v\in X_0$, the map that sends the class of a loop to the homotopy class of its path is an
isomorphism
\begin{equation}\label{8.10:eq-edge-loop-homotopy-1}
\pi_1^{\square}(X,v)\;\cong\;\pi_1(|X|,v).
\end{equation}
\end{definition}

\begin{proof}
\emph{Reduction to the component of $v$.} Every loop at $v$ lies in the path component of $v$ in
$|X|$. So does every plaquette used by a move (E2), since such a plaquette has a vertex of the loop
as a corner. The cubes of $X$ contained in this component form a cubical complex. Both sides of
\eqref{8.10:eq-edge-loop-homotopy-1} are unchanged if $X$ is replaced by it. We may therefore
assume that $|X|$ is path connected.

\emph{Cell structure.} $|X|$ is a finite CW complex whose open $k$-cells are the relative
interiors of the $k$-cubes of $X$. Its $1$-skeleton $|X|^{(1)}$ is the realisation of the graph
of $X$. Its $2$-skeleton $|X|^{(2)}$ arises from $|X|^{(1)}$ by attaching one $2$-cell for each
$P\in X_2$, along the path of a boundary circuit of $P$.

\emph{The graph.} Let $G$ be the set of loops at $v$ modulo moves (E1) alone. Under concatenation
it is a group, by the argument given for $\pi_1^{\square}$. We write $[\eta]$ for the class of a
loop $\eta$ in $G$. By the edge-path description of the fundamental group of a graph, every loop
in $|X|^{(1)}$ at $v$ is homotopic relative to its endpoints to the path of a loop at $v$.
Moreover, by the same description, the paths of two loops are homotopic relative to the endpoints
in $|X|^{(1)}$ if and only if the loops are related by finitely many moves (E1). Hence
$G\cong\pi_1(|X|^{(1)},v)$.

\emph{Attaching the $2$-cells.} Since $|X|$ is path connected, so is its $1$-skeleton
$|X|^{(1)}$. For each $P\in X_2$ fix a boundary circuit $c_P$ at a corner $y_P$, and an edge path
$\alpha_P$ from $v$ to $y_P$; such a path exists because $|X|^{(1)}$ is path connected. We use
the theorem on the effect of attaching $2$-cells on the fundamental group: if a space $Y$ is
obtained from a path-connected space $Z$ by attaching $2$-cells along loops $\phi_\beta$, the
inclusion induces a surjection $\pi_1(Z,v)\to\pi_1(Y,v)$ whose kernel is the normal subgroup
generated by the classes of the loops $\gamma_\beta\phi_\beta\overline{\gamma_\beta}$, where
$\gamma_\beta$ is a path from $v$ to the base point of $\phi_\beta$. Applied with
$Z=|X|^{(1)}$ and $Y=|X|^{(2)}$, it shows that the inclusion induces a surjection
$\pi_1(|X|^{(1)},v)\to\pi_1(|X|^{(2)},v)$. Its kernel corresponds to the normal subgroup
$\mathcal{N}\subset G$ generated by the classes $[\alpha_P c_P\overline{\alpha_P}]$.

We claim that $\mathcal{N}$ is the set of finite products of classes $[\alpha c\overline{\alpha}]$,
where $c$ is any boundary circuit of any plaquette $P$ at any corner $y$, and $\alpha$ is any
edge path from $v$ to $y$. Three facts give this.
\begin{itemize}
\item[(a)] Let $\beta$ run along $c_P$ from $y$ to $y_P$. The circuit of the same orientation at
$y$ is obtained from $\beta c_P\overline{\beta}$ by deleting backtracks.
\item[(b)] Reversing $c$ inverts its class.
\item[(c)] The identity
\[
[\eta][\alpha c\overline{\alpha}][\overline{\eta}]=[(\eta\alpha)c\,\overline{\eta\alpha}]
\]
holds for every loop $\eta$ at $v$.
\end{itemize}
By (a), (b) and (c) (the last with $\eta=\alpha\beta\overline{\alpha_P}$), each such class is a
conjugate in $G$ of a generator or of its inverse, so it lies in
$\mathcal{N}$. By (c), the set of such classes is closed under conjugation, and by (b) it is
closed under inversion. This proves the claim.

\emph{Higher cells.} $|X|$ arises from $|X|^{(2)}$ by attaching cells of dimension at least three.
By the cellular approximation theorem, the inclusion of the $2$-skeleton of a CW complex induces
an isomorphism of fundamental groups; hence $\pi_1(|X|^{(2)},v)\cong\pi_1(|X|,v)$.

\emph{Identification with the moves.} It remains to show that two loops $\eta,\eta'$ at $v$ are
combinatorially homotopic in $X$ if and only if $[\eta']\in\mathcal{N}[\eta]$.

A move (E1) does not change the class in $G$. Suppose $\eta'$ arises from $\eta$ by inserting a
circuit $c$ at a vertex $y$. Write $\eta=\eta_1\eta_2$, where $\eta_1$ runs from $v$ to $y$. Then
$\eta'=\eta_1c\eta_2$. The loop $\eta_1c\overline{\eta_1}\eta_1\eta_2$ reduces to $\eta'$ by
deletion of backtracks, so
\[
[\eta']=[\eta_1c\overline{\eta_1}]\,[\eta]\in\mathcal{N}[\eta].
\]
A deletion is the reverse of an insertion.

Conversely, let $[\eta']=n[\eta]$ with $n\in\mathcal{N}$. By the claim and (b), $n$ is a product
of $r$ classes $[\alpha c\overline{\alpha}]$. We argue by induction on $r$. It suffices to show
that $\alpha c\overline{\alpha}\eta$ is combinatorially homotopic to $\eta$. This loop arises by
one move (E2) from $\alpha\overline{\alpha}\eta$, inserting $c$ at the endpoint $y$ of $\alpha$.
The loop $\alpha\overline{\alpha}\eta$ reduces to $\eta$ by deletions of backtracks.

Therefore $\pi_1^{\square}(X,v)=G/\mathcal{N}$, and
\[
G/\mathcal{N}\cong\pi_1(|X|^{(2)},v)\cong\pi_1(|X|,v).
\]
The composite isomorphism sends the class of a loop to the homotopy class of its path. This is
\eqref{8.10:eq-edge-loop-homotopy-1}.
\end{proof}

\begin{lemma}[Transports under homotopy moves]\label{8.10:lem-homotopy-moves}
Transports are unchanged by moves of type \textup{(E1)} of edge-loop homotopy
(Definition~\ref{8.10:def-edge-loop-homotopy}). Here the transport $T_\gamma(V)$ of a
configuration $V$ on $X$ along an edge path $\gamma$, the plaquette variable $V(\partial q)$ of a
plaquette $q$ of $X$, flatness of $V$ on $X$, and the plaquette functional $F_q(V)$ are those of
the definition of transports in this chapter. Suppose a loop $\gamma$ based at the vertex $v$
admits a null-homotopy in $X$ across the plaquettes $(q_1,\dots,q_A)$. Then for every
configuration $V$
\begin{equation}\label{8.10:eq-homotopy-moves-1}
\lVert T_\gamma(V)-1\rVert\le\sum_{i\le A}F_{q_i}(V).
\end{equation}
In particular $T_\gamma(V)=1$ if $V$ is flat on $X$.
\end{lemma}

\begin{proof}
We write $T_\sigma$ for $T_\sigma(V)$ when the configuration is fixed, and we use that the
transport along a concatenation of paths is the product of the transports along the pieces,
in the order of traversal. From the definition of transports we use two properties of the
plaquette functional: $\lVert V(\partial q)-1\rVert\le F_q(V)$ for every plaquette $q$ of $X$,
and $F_q(V)=0$ for every such $q$ when $V$ is flat on $X$. The norm $\lVert\cdot\rVert$ is the
matrix norm of that definition; we use that it is invariant under left and right
multiplication by unitary matrices.

\emph{Moves of type \textup{(E1)}.} Such a move inserts or deletes a backtrack $y\to z\to y$
along a bond $(y,z)$. Write the loop as $\sigma\cdot\tau$, with $\sigma$ running from $v$ to
$y$ and the backtrack inserted at $y$. Since $V_{(y,z)}V_{(z,y)}=1$, the transport of the
modified loop is
\begin{equation*}
T_\sigma V_{(y,z)}V_{(z,y)}T_\tau=T_\sigma T_\tau,
\end{equation*}
the transport of the original loop. Hence moves of type \textup{(E1)} leave the transport
unchanged, which is the first assertion.

\emph{Moves of type \textup{(E2)}.} Such a move inserts or deletes the boundary $\partial q$ of a
plaquette $q$ of $X$, traversed in one of its two orientations. Write the loop as
$\sigma\cdot\tau$ with $\sigma$ running from $v$ to $y$, and insert $\partial q$ at $y$. The
transport of the new loop is
\begin{equation*}
T_\sigma V(\partial q)^{\pm1}T_\tau
=\bigl[T_\sigma V(\partial q)^{\pm1}T_\sigma^{-1}\bigr]\cdot T_\sigma T_\tau .
\end{equation*}
Thus the new transport is the old transport $T_\sigma T_\tau$ multiplied on the left by the
unitary $W=T_\sigma V(\partial q)^{\pm1}T_\sigma^{-1}$, which is a conjugate of
$V(\partial q)^{\pm1}$ by the unitary $T_\sigma$. Since
$V(\partial q)^{-1}-1=-V(\partial q)^{-1}\bigl(V(\partial q)-1\bigr)$, the unitary invariance of
the norm and the first property of $F_q(V)$ recalled above give
\begin{equation*}
\lVert W-1\rVert=\lVert V(\partial q)^{\pm1}-1\rVert=\lVert V(\partial q)-1\rVert\le F_q(V),
\end{equation*}
so $W$ lies within $F_q(V)$ of $1$. Consequently
\begin{equation*}
\lVert (W-1)\,T_\sigma T_\tau\rVert=\lVert W-1\rVert\le F_q(V),
\end{equation*}
the equality because $T_\sigma T_\tau$ is unitary and the norm is unitarily invariant. So
under a move of type \textup{(E2)} across $q$, in either direction, the transport changes in
norm by at most $F_q(V)$.

\emph{Telescoping.} Let $\gamma=\gamma_0,\gamma_1,\dots,\gamma_n=(v)$ be the sequence of loops
of the null-homotopy, each obtained from the previous one by a single move, the moves of type
\textup{(E2)} being across $q_1,\dots,q_A$. The transport of the constant loop $(v)$ is $1$. By
the triangle inequality,
\begin{equation*}
\lVert T_\gamma-1\rVert\le\sum_{j=1}^{n}\lVert T_{\gamma_{j-1}}-T_{\gamma_j}\rVert .
\end{equation*}
The terms coming from moves of type \textup{(E1)} vanish by the first step, and the term coming
from the move across $q_i$ is at most $F_{q_i}(V)$ by the second step. This
gives~\eqref{8.10:eq-homotopy-moves-1}. If $V$ is flat on $X$, each $F_{q_i}(V)$ vanishes by the
second property of the plaquette functional recalled above, and
\eqref{8.10:eq-homotopy-moves-1} gives $T_\gamma(V)=1$.
\end{proof}

\begin{lemma}[Flat fields on simply connected complexes are pure gauge]\label{8.10:lem-pure-gauge}
Let $X$ be a complex whose graph, with vertex set $X_0$, is connected, let $v\in X_0$, and let
$V$ be a configuration on the edges of $X$ that is flat on $X$.
Assume that every loop at $v$ is null-homotopic in $X$.
For $y\in X_0$ put
\begin{equation}\label{8.10:eq-pure-gauge-1}
u(y):=T_\sigma(V),\qquad \sigma \text{ any edge path in } X \text{ from } v \text{ to } y .
\end{equation}
Then $u$ is well defined on $X_0$, and
\begin{equation}\label{8.10:eq-pure-gauge-2}
V_{(y,z)}=u(y)^{-1}u(z)\qquad\text{for every edge } (y,z) \text{ of } X .
\end{equation}
\end{lemma}

\begin{proof}
Here $T_\sigma(V)$ is the transport of $V$ along the edge path $\sigma$. Recall that transports
are multiplicative under concatenation of composable edge paths,
$T_{\sigma\cdot\tau}(V)=T_\sigma(V)\,T_\tau(V)$, and that reversing a path replaces its
transport by the inverse.

Let $y\in X_0$. Since the graph of $X$ is connected, there is at least one edge path in $X$
from $v$ to $y$. Let $\sigma$ and $\sigma'$ be two such paths. Then $\sigma$ followed by
$\sigma'$ reversed is a loop at $v$ in $X$, and by hypothesis it is null-homotopic in $X$.
Since $V$ is flat on $X$, Lemma~\ref{8.10:lem-homotopy-moves} gives that the transport of $V$
along this loop equals $1$. By the multiplicativity and the behaviour under reversal recalled
above, that transport is $T_\sigma(V)\,T_{\sigma'}(V)^{-1}$. Hence
\begin{equation*}
T_\sigma(V)\,T_{\sigma'}(V)^{-1}=1,
\end{equation*}
that is, $T_\sigma(V)=T_{\sigma'}(V)$. The right side of \eqref{8.10:eq-pure-gauge-1} therefore
does not depend on the choice of $\sigma$, and $u$ is well defined on $X_0$.

Now let $(y,z)$ be an edge of $X$, and let $\sigma$ be an edge path in $X$ from $v$ to $y$.
Then $\sigma\cdot(y,z)$ is an edge path in $X$ from $v$ to $z$. By the well-definedness just
proved and by multiplicativity,
\begin{equation*}
u(z)=T_{\sigma\cdot(y,z)}(V)=T_\sigma(V)\,V_{(y,z)}=u(y)\,V_{(y,z)} .
\end{equation*}
Multiplying on the left by $u(y)^{-1}$ gives \eqref{8.10:eq-pure-gauge-2}.
\end{proof}

\begin{lemma}[Flat field with prescribed holonomy]\label{8.10:lem-prescribed-holonomy}
Let $X$ be a cubical complex whose graph is connected, let $T$ be a spanning tree of that graph,
let $v\in X_0$, and let $\Phi\colon \pi_1^{\square}(X,v)\to G$ be a homomorphism, with
$\pi_1^{\square}(X,v)$ the edge-loop group of Definition~\ref{8.10:def-edge-loop-homotopy}.
For an edge $e=(y,z)$ of $X$ let $c_e$ be the loop at $v$ that follows the path in $T$ from $v$
to $y$, then the edge $e$, then the path in $T$ from $z$ to $v$. Put $W_e=1$ for every edge $e$
of $T$, and $W_{(y,z)}:=\Phi([c_{(y,z)}])$ for every other edge $(y,z)$ of $X$. Then
\begin{equation}\label{8.10:eq-prescribed-holonomy-1}
  T_\gamma(W)=\Phi([\gamma])\qquad\text{for every loop $\gamma$ at $v$ in $X$;}
\end{equation}
in particular $W$ is flat on $X$.
\end{lemma}

\begin{proof}
We first show that $W_e=\Phi([c_e])$ for every edge $e$ of $X$. For an edge not in $T$ this is
the definition of $W_e$. Let $e$ be an edge of $T$. Then the loop $c_e$ runs entirely in the tree
$T$, and it reduces to the constant loop $(v)$ by moves (E1) of
Definition~\ref{8.10:def-edge-loop-homotopy}. Hence $[c_e]$ is the identity of
$\pi_1^{\square}(X,v)$, and since $\Phi$ is a homomorphism, $\Phi([c_e])=1=W_e$.

Now let $\gamma=(v=y_0,y_1,\dots,y_n=v)$ be a loop at $v$. Consider the concatenation
\begin{equation*}
  c_{(y_0,y_1)}\,c_{(y_1,y_2)}\cdots c_{(y_{n-1},y_n)} .
\end{equation*}
For $1\le i\le n-1$, the loop $c_{(y_{i-1},y_i)}$ ends with the path in $T$ from $y_i$ to $v$,
and the next loop $c_{(y_i,y_{i+1})}$ begins with the path in $T$ from $v$ to $y_i$; moreover
the first loop begins, and the last ends, with the path in $T$ from $v$ to $v=y_0=y_n$. These tree
segments are removed by moves (E1), and what remains is $\gamma$. Thus the concatenation reduces
to $\gamma$ by moves (E1), and therefore
\begin{equation*}
  [\gamma]=\prod_{i=1}^{n}\,[c_{(y_{i-1},y_i)}]
\end{equation*}
in $\pi_1^{\square}(X,v)$. Applying the homomorphism $\Phi$ and the identity
$W_e=\Phi([c_e])$ of the first paragraph, we obtain
\begin{equation*}
  T_\gamma(W)=\prod_{i=1}^{n}W_{(y_{i-1},y_i)}
  =\prod_{i=1}^{n}\Phi\bigl([c_{(y_{i-1},y_i)}]\bigr)=\Phi([\gamma]),
\end{equation*}
which is \eqref{8.10:eq-prescribed-holonomy-1}.

Finally let $P\in X_2$, and let $\gamma_P$ be the loop at $v$ that follows the path in $T$ from
$v$ to a corner of $P$, then the boundary $\partial P$, then the same path in $T$ back to $v$.
This loop is null-homotopic in $X$, so $[\gamma_P]$ is the identity and, by
\eqref{8.10:eq-prescribed-holonomy-1}, $T_{\gamma_P}(W)=\Phi([\gamma_P])=1$. Since $W=1$ on
every edge of $T$, the transport $T_{\gamma_P}(W)$ equals $W(\partial P)$. Hence
$W(\partial P)=1$ for every $P\in X_2$, that is, $W$ is flat on $X$.
\end{proof}

\begin{lemma}[Lattice Stokes bound for a rectangle]\label{8.10:lem-stokes-rectangle}
Let $\mu\neq\nu$ be two coordinate directions, let $A,B\ge 1$ be integers, let $y$ be a lattice point, and put
\begin{equation*}
  \Pi:=\{\,y+s e_\mu+t e_\nu \;:\; 0\le s\le A,\ 0\le t\le B\,\}.
\end{equation*}
Let $\gamma$ be the boundary loop of $\Pi$ based at $y$: $A$ steps $+e_\mu$, then $B$ steps $+e_\nu$, then $A$ steps $-e_\mu$, then $B$ steps $-e_\nu$. Let $P(\Pi)$ be the set of the $AB$ plaquettes
\begin{equation*}
  (y+s e_\mu+t e_\nu)+\{0,1\}^{\{\mu,\nu\}},\qquad 0\le s<A,\quad 0\le t<B .
\end{equation*}
Then, for every configuration $V$, the transport $T_\gamma(V)$ is an ordered product of $AB$ factors $g_i V(\partial p_i)^{\pm1} g_i^{-1}$, with $g_i\in G$, one factor for each $p_i\in P(\Pi)$. Consequently
\begin{equation*}
  \|T_\gamma(V)-1\|\le \sum_{p\in P(\Pi)} F_p(V),
\end{equation*}
and $\gamma$ admits a null-homotopy in $K(\Pi)$ across the $AB$ plaquettes of $P(\Pi)$.
\end{lemma}

Here $T_\gamma(V)$, $F_p(V)$, the moves and the notion of a null-homotopy across a sequence of plaquettes are those of Definition~\ref{8.10:def-edge-loop-homotopy} and Lemma~\ref{8.10:lem-homotopy-moves}. The group $G$ is the gauge group $\mathrm{SU}(N)$, in which the bond variables of $V$ take values, and $\|\cdot\|$ is the operator norm. Finally, $K(\Pi)$ is the cubical complex formed by the lattice points, bonds and plaquettes contained in $\Pi$. The transport along a concatenation of paths is the product of the transports in the order of traversal, and the transport along a reversed path is the inverse.

\begin{proof}
For a lattice point $z$ and an integer $B'\ge1$, write $\eta_{B'}(z)$ for the boundary loop, based at $z$, of the $1\times B'$ rectangle
\begin{equation*}
  \{z+s e_\mu+t e_\nu : 0\le s\le 1,\ 0\le t\le B'\}.
\end{equation*}
This loop takes one step $+e_\mu$, then $B'$ steps $+e_\nu$, then one step $-e_\mu$, then $B'$ steps $-e_\nu$. Put $p_z:=z+\{0,1\}^{\{\mu,\nu\}}$, so that $\eta_1(z)=\partial p_z$.

\emph{(a) Rectangles of width one.} We show by induction on $B'$ that $T_{\eta_{B'}(z)}(V)$ is an ordered product of $B'$ factors $g V(\partial p)^{\pm1}g^{-1}$, $g\in G$, one for each plaquette $p_{z+te_\nu}$, $0\le t<B'$. For $B'=1$ the single factor is $V(\partial p_z)$. Let $B'\ge2$ and consider the concatenated loop
\begin{equation*}
  \partial p_z\cdot\bigl[(z,z+e_\nu)\cdot\eta_{B'-1}(z+e_\nu)\cdot(z+e_\nu,z)\bigr].
\end{equation*}
It passes through
\begin{equation*}
  z,\ z+e_\mu,\ z+e_\mu+e_\nu,\ z+e_\nu,\ z,\ z+e_\nu,\ z+e_\mu+e_\nu,\ \dots,
\end{equation*}
and then follows $\eta_{B'-1}(z+e_\nu)$ back to $z+e_\nu$ before stepping to $z$. The loop contains the backtrack $z+e_\nu\to z\to z+e_\nu$; we delete it. The result then contains the backtrack $z+e_\mu+e_\nu\to z+e_\nu\to z+e_\mu+e_\nu$; we delete it as well. What remains is exactly $\eta_{B'}(z)$.

By Lemma~\ref{8.10:lem-homotopy-moves}, transports are unchanged by these moves. Since $V_{(z+e_\nu,z)}=V_{(z,z+e_\nu)}^{-1}$, we obtain
\begin{equation*}
  T_{\eta_{B'}(z)}(V)
  =V(\partial p_z)\,V_{(z,z+e_\nu)}\,T_{\eta_{B'-1}(z+e_\nu)}(V)\,V_{(z,z+e_\nu)}^{-1}.
\end{equation*}
By the induction hypothesis, $T_{\eta_{B'-1}(z+e_\nu)}(V)$ is an ordered product of $B'-1$ factors $gV(\partial p)^{\pm1}g^{-1}$, one for each $p_{z+te_\nu}$ with $1\le t<B'$. Conjugation by $h:=V_{(z,z+e_\nu)}\in G$ conjugates each factor, replacing $g$ by $hg$. Together with the factor $V(\partial p_z)$ this gives the assertion for $B'$.

\emph{(b) Induction on the width.} Keep $y$ and $B$ fixed. For $1\le A'\le A$, write $\Pi_{A'}$ for the rectangle of width $A'$ and height $B$ at $y$, and $\partial\Pi_{A'}$ for its boundary loop, formed as $\gamma$ is with $A'$ in place of $A$. Thus $\partial\Pi_A=\gamma$ and $\partial\Pi_1=\eta_B(y)$, so that (a) settles $A'=1$.

Let $A'\ge2$. Put $y':=y+(A'-1)e_\mu$, let $\xi$ be the path of $A'-1$ steps $+e_\mu$ from $y$ to $y'$, and let $\bar\xi$ be its reverse. Consider
\begin{equation}\label{8.10:eq-stokes-rectangle-1}
  \xi\cdot\eta_B(y')\cdot\bar\xi\cdot\partial\Pi_{A'-1}.
\end{equation}
The loop $\partial\Pi_{A'-1}$ begins with $\xi$. Deleting the $A'-1$ successive backtracks of $\bar\xi\cdot\xi$ leaves three pieces in order: $\xi$; then $\eta_B(y')$; then the rest of $\partial\Pi_{A'-1}$. That rest consists of $B$ steps $+e_\nu$ from $y'$, then $A'-1$ steps $-e_\mu$, then $B$ steps $-e_\nu$ down to $y$.

The last side of $\eta_B(y')$ descends $B$ steps from $y'+Be_\nu$ to $y'$. It is immediately followed by the ascending side from $y'$ to $y'+Be_\nu$, and deleting these $B$ successive backtracks leaves the following loop. It takes the $A'-1$ steps of $\xi$, one step $+e_\mu$ to $y+A'e_\mu$, and $B$ steps $+e_\nu$. It then takes one step $-e_\mu$ to $y'+Be_\nu$, followed by $A'-1$ steps $-e_\mu$ to $y+Be_\nu$. It ends with $B$ steps $-e_\nu$. This is $\partial\Pi_{A'}$.

By Lemma~\ref{8.10:lem-homotopy-moves} and the rules for concatenation and reversal,
\begin{equation}\label{8.10:eq-stokes-rectangle-2}
  T_{\partial\Pi_{A'}}(V)=T_\xi(V)\,T_{\eta_B(y')}(V)\,T_\xi(V)^{-1}\,T_{\partial\Pi_{A'-1}}(V).
\end{equation}
By (a), $T_{\eta_B(y')}(V)$ is an ordered product of $B$ factors of the required form, one for each plaquette $(y+(A'-1)e_\mu+te_\nu)+\{0,1\}^{\{\mu,\nu\}}$, $0\le t<B$. Conjugation by $T_\xi(V)\in G$ preserves this form. By the induction hypothesis, $T_{\partial\Pi_{A'-1}}(V)$ is an ordered product of $(A'-1)B$ such factors, one for each plaquette with $0\le s<A'-1$, $0\le t<B$. Hence $T_{\partial\Pi_{A'}}(V)$ is an ordered product of $A'B$ factors, one for each plaquette of $P(\Pi_{A'})$. For $A'=A$ this is the product representation of $T_\gamma(V)$.

\emph{The norm bound.} Let $h_1,\dots,h_n\in G$. Each $h_j$ is unitary, so left multiplication by a product of the $h_j$ preserves the operator norm. Hence
\begin{equation*}
  \|h_1\cdots h_n-1\|
  \le\|h_1\cdots h_{n-1}(h_n-1)\|+\|h_1\cdots h_{n-1}-1\|
  =\|h_n-1\|+\|h_1\cdots h_{n-1}-1\|,
\end{equation*}
and by induction on $n$, $\|h_1\cdots h_n-1\|\le\sum_j\|h_j-1\|$. We apply this to the product representation of $T_\gamma(V)$. For $g\in G$ and unitary $W$ we have
\begin{equation*}
  \|gW^{\pm1}g^{-1}-1\|=\|W^{\pm1}-1\|=\|W-1\|.
\end{equation*}
Therefore
\begin{equation*}
  \|T_\gamma(V)-1\|\le\sum_{p\in P(\Pi)}\|V(\partial p)-1\|.
\end{equation*}
For each $p\in P(\Pi)$, the loop $\partial p$ is null-homotopic across the single plaquette $p$. Lemma~\ref{8.10:lem-homotopy-moves} then gives $\|V(\partial p)-1\|\le F_p(V)$, and the asserted bound follows.

\emph{The null-homotopy.} The identities above read as sequences of moves, and every path occurring in them lies in $\Pi$, hence in $K(\Pi)$.

In (a), the loop $\eta_{B'}(z)$ is obtained from $\partial p_z\cdot[(z,z+e_\nu)\cdot\eta_{B'-1}(z+e_\nu)\cdot(z+e_\nu,z)]$ by two backtrack deletions, so the reverse moves turn $\eta_{B'}(z)$ into that loop. Removing the plaquette loop $\partial p_z$ across $p_z$ leaves $(z,z+e_\nu)\cdot\eta_{B'-1}(z+e_\nu)\cdot(z+e_\nu,z)$. By induction, $\eta_{B'-1}(z+e_\nu)$ is null-homotopic across $p_{z+te_\nu}$, $1\le t<B'$. The backtrack $(z,z+e_\nu)\cdot(z+e_\nu,z)$ that remains is then deleted. This gives a null-homotopy of $\eta_{B'}(z)$ across the $B'$ plaquettes $p_{z+te_\nu}$, $0\le t<B'$.

In (b), the reverse of the backtrack deletions turns $\partial\Pi_{A'}$ into \eqref{8.10:eq-stokes-rectangle-1}. The null-homotopy of $\eta_B(y')$ across its $B$ plaquettes, followed by the deletion of the backtracks $\xi\cdot\bar\xi$, leaves $\partial\Pi_{A'-1}$. By induction, $\partial\Pi_{A'-1}$ is null-homotopic across the $(A'-1)B$ plaquettes of $P(\Pi_{A'-1})$. For $A'=A$, this is a null-homotopy of $\gamma$ in $K(\Pi)$ across the $AB$ plaquettes of $P(\Pi)$.
\end{proof}

\begin{lemma}[The ladder identity]\label{8.10:lem-ladder-identity}
Let $V$ be a configuration on the terminal lattice, with the transports $T(\cdot)$, the plaquette
variables $V(\partial p)$ and the plaquette defects $F_p(V)=\|V(\partial p)-1\|$ of
Definition~\ref{8.10:not-terminal-lattice}. Let $P=f_1\cdots f_n$ be an edge path with vertices
$v_0,\dots,v_n$, the bond $f_k$ running from $v_{k-1}$ to $v_k$. Let $e$ be a lattice unit vector, let
$P'=f_1'\cdots f_n'$ be the translated path, with vertices $v_k+e$, let $c_k:=(v_k,v_k+e)$,
$0\le k\le n$, be the rungs, and let $\sigma_k$, $1\le k\le n$, be the plaquette with boundary
$f_k\,c_k\,f_k'^{-1}\,c_{k-1}^{-1}$. Then
\begin{equation}\label{8.10:eq-ladder-identity-1}
T(P)\,V_{c_n}\,T(P')^{-1}\,V_{c_0}^{-1}
\end{equation}
is a product of conjugates of the elements $V(\partial\sigma_k)^{\pm1}$, $1\le k\le n$, and therefore
\begin{equation}\label{8.10:eq-ladder-identity-2}
\bigl\|T(P)\,V_{c_n}\,T(P')^{-1}\,V_{c_0}^{-1}-1\bigr\|\le\sum_{k=1}^{n}F_{\sigma_k}(V).
\end{equation}
If $P$ is a loop, then
\begin{equation}\label{8.10:eq-ladder-identity-3}
\|T(P)-1\|\le\sum_{k=1}^{n}F_{\sigma_k}(V)+\|T(P')-1\|.
\end{equation}
\end{lemma}

\begin{proof}
We prove the first assertion by induction on $n$. For $n=1$ the element
\eqref{8.10:eq-ladder-identity-1} equals $V_{f_1}V_{c_1}V_{f_1'}^{-1}V_{c_0}^{-1}$, which is the
holonomy of $\partial\sigma_1$ read from the vertex $v_0$; it is therefore a conjugate of
$V(\partial\sigma_1)^{\pm1}$.

Let $n>1$, put $P_{n-1}:=f_1\cdots f_{n-1}$, and let $P_{n-1}'$ be its translate. Since
$T(P)=T(P_{n-1})V_{f_n}$ and $T(P')=T(P_{n-1}')V_{f_n'}$, inserting
$V_{c_{n-1}}^{-1}\,T(P_{n-1})^{-1}\,T(P_{n-1})\,V_{c_{n-1}}=1$ gives
\begin{align*}
T(P)\,V_{c_n}\,T(P')^{-1}\,V_{c_0}^{-1}
&=T(P_{n-1})\,V_{f_n}V_{c_n}V_{f_n'}^{-1}\,T(P_{n-1}')^{-1}\,V_{c_0}^{-1}\\
&=T(P_{n-1})\bigl[V_{f_n}V_{c_n}V_{f_n'}^{-1}V_{c_{n-1}}^{-1}\bigr]T(P_{n-1})^{-1}\\
&\qquad\cdot\bigl[T(P_{n-1})\,V_{c_{n-1}}\,T(P_{n-1}')^{-1}\,V_{c_0}^{-1}\bigr].
\end{align*}
The first bracket is the holonomy of $\partial\sigma_n$ read from $v_{n-1}$, hence a conjugate of
$V(\partial\sigma_n)^{\pm1}$, and so is its conjugate by $T(P_{n-1})$. The second bracket is the
element \eqref{8.10:eq-ladder-identity-1} for the path $P_{n-1}$, whose rungs are $c_0,\dots,c_{n-1}$
and whose plaquettes are $\sigma_1,\dots,\sigma_{n-1}$; by the induction hypothesis it is a product of
conjugates of the $V(\partial\sigma_k)^{\pm1}$, $k\le n-1$. This proves the first assertion.

For \eqref{8.10:eq-ladder-identity-2}: by the induction just given, the element
\eqref{8.10:eq-ladder-identity-1} equals $\Xi_n\Xi_{n-1}\cdots\Xi_1$, one factor for each $k$, where
$\Xi_k=Y_k\,V(\partial\sigma_k)^{\varepsilon_k}\,Y_k^{-1}$ with $Y_k\in\mathrm{SU}(N)$ and
$\varepsilon_k\in\{1,-1\}$. For unitary $\Phi,\Psi$ one has $\Phi\Psi-1=\Phi(\Psi-1)+(\Phi-1)$, so,
the norm being invariant under multiplication by unitary matrices,
$\|\Phi\Psi-1\|\le\|\Phi-1\|+\|\Psi-1\|$; by induction on the number of factors,
$\|\Xi_n\cdots\Xi_1-1\|\le\sum_k\|\Xi_k-1\|$. Moreover, for unitary $Y$ and $\Xi$,
\begin{gather*}
\|Y\Xi Y^{-1}-1\|=\|Y(\Xi-1)Y^{-1}\|=\|\Xi-1\|,\\
\|\Xi^{-1}-1\|=\|\Xi^{-1}(1-\Xi)\|=\|\Xi-1\|.
\end{gather*}
Hence $\|\Xi_k-1\|=\|V(\partial\sigma_k)-1\|=F_{\sigma_k}(V)$, and \eqref{8.10:eq-ladder-identity-2} follows.

If $P$ is a loop, then $v_n=v_0$, so $c_n=c_0$. Put $\Pi:=T(P)\,V_{c_0}\,T(P')^{-1}\,V_{c_0}^{-1}$;
then
\begin{equation*}
T(P)=\Pi\cdot V_{c_0}\,T(P')\,V_{c_0}^{-1},\qquad \|\Pi-1\|\le\sum_{k=1}^{n}F_{\sigma_k}(V),
\end{equation*}
the bound being \eqref{8.10:eq-ladder-identity-2}. By the product inequality and the conjugation
invariance just used,
\begin{equation*}
\|T(P)-1\|\le\|\Pi-1\|+\bigl\|V_{c_0}T(P')V_{c_0}^{-1}-1\bigr\|
=\|\Pi-1\|+\|T(P')-1\|,
\end{equation*}
which is \eqref{8.10:eq-ladder-identity-3}.
\end{proof}

\begin{definition}[Couplings and thresholds at the terminal level]\label{8.10:not-terminal-couplings}
In this section $g>0$ is the reference coupling of Theorem~\ref{3.1:thm-main}, and
\begin{equation*}
X:=g^{-2}.
\end{equation*}

\emph{Terminal level.} The bare spacing is $L^{-K}$ with $K\ge 1$. Of the $K$ renormalisation
steps, $m\ge 0$ are not run; the number $m$ does not depend on $K$, and the run ends at the
\emph{terminal level}
\begin{equation*}
K':=K-m .
\end{equation*}
The volume is $\mathfrak{N}=N_{\mathrm{T}}L^{m}$, with $N_{\mathrm{T}}$ an integer on which no
condition is placed.
We write $g_{K'}$ for the running coupling at level $K'$ (Definition~\ref{1.2:def-couplings}), and
\begin{equation}\label{8.10:eq-terminal-couplings-1}
x:=x_{K'}:=g_{K'}^{-2}.
\end{equation}

\emph{Degree function.} For $x'\ge e$ we put
\begin{equation}\label{8.10:eq-terminal-couplings-2}
p_0(x'):=A_0\,(\log x')^{p_0},
\end{equation}
where $A_0>0$ is Ba{\l}aban's constant and the exponent, denoted by the same letter, satisfies
$p_0\ge 5r$, where $r>0$ is Ba{\l}aban's exponent. For a coupling $g'$ we write
$p_0(g'):=p_0(g'^{-2})$. In particular
\begin{equation*}
p_0(g_{K'})=p_0(x).
\end{equation*}

\emph{Small-field threshold.} The small-field threshold at the terminal level $K'$ is
\begin{equation}\label{8.10:eq-terminal-couplings-3}
\varepsilon:=\varepsilon_{K'}:=g_{K'}\,p_0(g_{K'})=x^{-1/2}\,p_0(x).
\end{equation}

\emph{The integer $\check{R}_{K'}$ and the length $\ell$.} We fix a positive integer
$\check{R}_{K'}$ with
\begin{equation*}
\check{R}_{K'}\le L(\log x)^{r},
\end{equation*}
and we put $\ell:=M\check{R}_{K'}$. Then
\begin{equation}\label{8.10:eq-terminal-couplings-4}
\ell=M\check{R}_{K'}\le \ell_*(x):=ML(\log x)^{r}.
\end{equation}

\emph{Inputs from clause~\textup{(0)}.} Clause~(0) of Theorem~\ref{3.1:thm-main} supplies two
inequalities, which this chapter uses as inputs. The first is
\begin{equation}\label{8.10:eq-terminal-couplings-5}
X\le x_k\qquad\text{for every level }k .
\end{equation}
The second is
\begin{equation}\label{8.10:eq-terminal-couplings-6}
p_0(g_j)\ge p_0(g_k)-1\qquad\text{for all levels } j\le k\le K' .
\end{equation}
Clause~(0) provides the inequality \eqref{8.10:eq-terminal-couplings-6} under a condition on the
couplings, the \emph{coupling floor}. Wherever \eqref{8.10:eq-terminal-couplings-6} is used
as an input, that floor is part of the input.

\emph{Ba{\l}aban's constants.} The constants
\begin{equation*}
\beta_0\in\bigl(0,\tfrac13\bigr),\qquad \kappa_1\ge 0,\qquad \gamma_0>0,\qquad A_1>0,
\end{equation*}
and the function $d_{K'}(\cdot)\ge 0$, are those of \cite[(1.79)--(1.89)]{B-CMP122b}.

\emph{Reserve fractions.} The flat reserve fraction is a number
\begin{equation*}
\theta\in\Bigl(0,\tfrac{1-3\beta_0}{2}\Bigr) ,
\end{equation*}
and the linear reserve fraction is a number $\theta'\in(0,1)$. The first is used in the flat
reserve estimate of this chapter, the second in the linear reserve estimate.
\end{definition}

\begin{definition}[Histories, live components, fibres and the deletion map]\label{8.10:def-histories-fibres}
We use the terminal level $K'$, the terminal lattice with its cells and its configurations $V$,
and the running couplings $g_j$, as fixed in Definitions~\ref{8.10:not-terminal-lattice}
and~\ref{8.10:not-terminal-couplings}. We write $\mathcal{V}$ for the space of configurations $V$
and $dV$ for the measure on it with which the representation recalled in item~(a) is written.

Items (a)--(f) recall, by statement, the objects of the representation of Ba{\l}aban's run down
to the terminal level by un-resummed histories; item (g) defines freeness of a component from
the kept structure.

\begin{itemize}
\item[(a)] \emph{Histories and their densities.} The set $\mathcal{W}$ of \emph{un-resummed
histories} of the run is a finite set. To each $\omega\in\mathcal{W}$ are attached two objects.
The first is a $\sigma$-finite measure space $(\mathcal{U}_\omega,du)$ of internal variables that
are not integrated; it is a finite product of compact groups, each with its Haar measure, and
Euclidean spaces, each with Lebesgue measure. The second is a measurable function
\begin{equation*}
\rho_\omega:\mathcal{U}_\omega\times\mathcal{V}\to[0,\infty)
\end{equation*}
whose total integral is finite:
\begin{equation*}
\tau_\omega:=\int\rho_\omega\,du\,dV<\infty .
\end{equation*}
That $\rho_\omega\ge0$ is the non-negativity of the history densities; it
rests on \cite{B-CMP122a} and \cite{B-CMP122b}. The representation also gives
$\mathsf{Z}>0$, where $\mathsf{Z}$ and the probability $\tilde\nu$ on subsets
$E\subset\mathcal{W}$ are defined by
\begin{equation*}
\mathsf{Z}:=\sum_{\omega\in\mathcal{W}}\tau_\omega,
\qquad
\tilde\nu(E):=\mathsf{Z}^{-1}\sum_{\omega\in E}\tau_\omega .
\end{equation*}

\item[(b)] \emph{Terminal live components.} A union of cells is \emph{touch-connected} if it is
connected for the relation of touching between cells; its \emph{touch-components} are its
maximal touch-connected sub-unions. Each $\omega\in\mathcal{W}$ has a \emph{terminal live
region}, which is a union of cells. The touch-components of the terminal live region of
$\omega$ are the \emph{terminal live components} of $\omega$.

\item[(c)] \emph{Internal histories and seeds.} Let $C$ be a touch-connected union of cells.
We put
\begin{equation*}
\mathcal{L}_C:=\{\omega\in\mathcal{W}:\ C\text{ is a terminal live component of }\omega\}.
\end{equation*}
\begin{itemize}
\item[$\cdot$] For $\omega\in\mathcal{L}_C$, the \emph{internal history}
$\operatorname{hist}_C(\omega)$ of $C$ in $\omega$ lies in a finite set $\mathfrak{H}(C)$.
\item[$\cdot$] Each $\mathfrak{h}\in\mathfrak{H}(C)$ carries a finite non-empty set
$\mathcal{Z}(\mathfrak{h})$ of \emph{seeds} $\eta$.
\item[$\cdot$] Each seed $\eta$ has a \emph{birth level} $j(\eta)\in\{1,\dots,K'\}$.
\item[$\cdot$] Each seed is of one of two types. A seed $\eta$ of the \emph{first type}, with
$j=j(\eta)$, is a creation component in the sense of \cite[p.~380]{B-CMP122b}: a component of
the cover, by the $MR_j$-cubes of that paper (cubes of the side fixed there for step $j$), of
the large-field region created at step $j$. Otherwise the seed is \emph{preparatory}.
\item[$\cdot$] We write $\mathcal{Z}^1(\mathfrak{h})\subset\mathcal{Z}(\mathfrak{h})$ for the
set of seeds of the first type.
\item[$\cdot$] For $\eta\in\mathcal{Z}^1(\mathfrak{h})$, the number
$\mathfrak{d}(\eta)\in\{0,1,2,\dots\}$ is the linear size $d'_{j(\eta)}(\eta)$ of the creation
component $\eta$ at its birth level. This is the quantity $d'_j(Z^{(i)}_j)$ of
\cite[(1.79)]{B-CMP122b}.
\end{itemize}

\item[(d)] \emph{Sets of internal histories and their weights.} For
$\mathcal{S}\subset\mathfrak{H}(C)$ we put
\begin{equation*}
\mathcal{L}_{C,\mathcal{S}}:=\{\omega\in\mathcal{L}_C:\ \operatorname{hist}_C(\omega)\in
\mathcal{S}\}.
\end{equation*}
For a real number $\theta\ge0$, and with the degree function $p_0(\cdot)$, we also put
\begin{equation*}
\varpi_\theta(\mathcal{S}):=\sup_{\mathfrak{h}\in\mathcal{S}}
\exp\Bigl(-\theta\sum_{\eta\in\mathcal{Z}(\mathfrak{h})}p_0\bigl(g_{j(\eta)}\bigr)\Bigr)
\in(0,1].
\end{equation*}

\item[(e)] \emph{The deletion map and fibres.} The \emph{deletion map}
\begin{equation*}
\operatorname{red}_C:\mathcal{L}_C\to\mathcal{W}\setminus\mathcal{L}_C
\end{equation*}
is defined on $\omega\in\mathcal{L}_C$ as follows: it replaces $C$, together with its lineage
in $\omega$, by small field, and it leaves every other label of $\omega$ unchanged. Since $C$ is
a component at the terminal level $K'$, there is no later label that the replacement could
affect. We put
$\Lambda_C:=\operatorname{red}_C(\mathcal{L}_C)$. The \emph{fibre} of $\lambda\in\Lambda_C$ is
\begin{equation*}
\mathcal{F}(\lambda,C):=\operatorname{red}_C^{-1}(\lambda).
\end{equation*}
The representation of item (a) provides the following two properties, for every
$\lambda\in\Lambda_C$.
\begin{itemize}
\item[(i)] The map $\operatorname{hist}_C$ is injective on $\mathcal{F}(\lambda,C)$.
\item[(ii)] For every $\omega\in\mathcal{F}(\lambda,C)$ the space of internal variables
factorises as
\begin{equation*}
\mathcal{U}_\omega=\mathcal{U}^C_\omega\times\mathcal{U}_\lambda ,
\end{equation*}
where $\mathcal{U}^C_\omega$ carries the internal variables of $C$.
\end{itemize}

\item[(f)] \emph{Kept structure.} For $\lambda\in\Lambda_C$, the \emph{kept structure}
$\mathfrak{K}(\lambda)$ is the set of cells of the terminal live region of $\lambda$.

\item[(g)] \emph{Freeness of a component from the kept structure.} Let $C$ be a touch-connected
union of cells, let $\lambda\in\Lambda_C$, and let $q\ge0$ be an integer. With
$\operatorname{hull}_q(C)$ the $q$-hull of $C$, we say that $C$ is
\emph{free at range $q$ from the kept structure of} $\lambda$, and we write
$\mathsf{K}_{\mathrm{free}}(C;\lambda;q)$, when
$\operatorname{hull}_q(C)\cap\mathfrak{K}(\lambda)=\emptyset$; that is,
\begin{equation}\label{8.10:eq-histories-fibres-1}
\mathsf{K}_{\mathrm{free}}(C;\lambda;q)\ :\Longleftrightarrow\
\operatorname{hull}_q(C)\cap\mathfrak{K}(\lambda)=\emptyset .
\end{equation}
\end{itemize}
\end{definition}

\begin{definition}[Domination of a fibre by the hole majorant. Proof outstanding]
\label{8.10:def-fibre-domination}
We use the touch-connected unions $C$ of cells, the set $\Lambda_C$ of fibre labels $\lambda$
attached to $C$, the fibre $\mathcal{F}(\lambda,C)$, the set $\mathfrak{H}(C)$ of histories
internal to $C$ and the map $\omega\mapsto\operatorname{hist}_C(\omega)$, which assigns to a
member $\omega$ of the fibre its history on $C$, as in
Definition~\ref{8.10:def-histories-fibres}. For $\lambda\in\Lambda_C$ we write
$\mathcal{U}_\lambda$ for the space of the variables $\xi$ of the fibre outside $C$, and
$\mathcal{V}$ for the space of the block fields $V$. For $\omega\in\mathcal{F}(\lambda,C)$ we
write $\mathcal{U}^C_\omega$ for the space of the variables $u$ internal to $C$, and
$\rho_\omega(u,\xi,V)$ for the density of $\omega$. Let $\theta'>0$ be the fraction, and
$\gamma_0$ the constant, appearing in the exponent of the linear-decay statement; let $\theta$
be the parameter of the weight factor $\varpi_\theta$ and of the factor
$e^{-(2-\theta)\underline{p}_0}$ below; and let $K'$ be a level. For a set $\mathcal{S}$ of
histories, $\varpi_\theta(\mathcal{S})$ is the weight factor attached to $\mathcal{S}$. For
$\mathfrak{h}\in\mathfrak{H}(C)$, $\mathcal{Z}^1(\mathfrak{h})$ is the set of creation
components of $\mathfrak{h}$. For $\mathfrak{z}\in\mathcal{Z}^1(\mathfrak{h})$, $j(\mathfrak{z})$
is the level at which $\mathfrak{z}$ is created and $\mathfrak{d}(\mathfrak{z})$ is the quantity
that the linear-decay statement attaches to $\mathfrak{z}$ in its exponential factor.
The number $d_{K'}(C)$ is the function $d_j(\cdot)$ of the glossary,
Notation~\ref{0.2:not-glossary}, at $j=K'$, evaluated on $C$. Let $0<\theta'_\tau\le\theta'$ be
the part of the fraction $\theta'$ that is allotted to the present estimate; the index $\tau$
marks this allotment and is carried by the penalty $\operatorname{pen}^\tau_\theta$ below. Put
\[
  \underline{p}_0:=\min_{1\le j\le K'}p_0(g_j),
\]
where $p_0(\cdot)$ is the degree function.

\emph{Fibre domination at $(\theta,\theta')$ and level $K'$} is the following statement. For
every touch-connected union $C$ of cells and every $\lambda\in\Lambda_C$ there exist a
measurable function
\[
  M_{\lambda,C}:\mathcal{U}_\lambda\times\mathcal{V}\to[0,\infty[
  \quad\text{with}\quad
  \int M_{\lambda,C}\,d\xi\,dV<\infty,
\]
and a number $\delta_C\ge0$ depending on $C$ only, with the following property. For every
non-empty $\mathcal{S}\subset\mathfrak{H}(C)$ and every $(\xi,V)$,
\begin{equation}\label{8.10:eq-fibre-domination-1}
  \sum_{\substack{\omega\in\mathcal{F}(\lambda,C)\\ \operatorname{hist}_C(\omega)\in\mathcal{S}}}
  \int_{\mathcal{U}^C_\omega}\rho_\omega(u,\xi,V)\,du
  \le \operatorname{pen}^\tau_\theta(C,\mathcal{S})\,e^{\delta_C}\,M_{\lambda,C}(\xi,V).
\end{equation}
Here
\begin{equation}\label{8.10:eq-fibre-domination-2}
\begin{aligned}
  \operatorname{pen}^\tau_\theta(C,\mathcal{S})
  &:=\varpi_\theta(\mathcal{S})\,r_{\theta'_\tau}(\mathcal{S})\,
     \exp\bigl(-(2-\theta)\underline{p}_0-\kappa_1 d_{K'}(C)\bigr),\\
  r_{\theta'_\tau}(\mathcal{S})
  &:=\sup_{\mathfrak{h}\in\mathcal{S}}\ \prod_{\mathfrak{z}\in\mathcal{Z}^1(\mathfrak{h})}
     \exp\Bigl(-\theta'_\tau\gamma_0A_1^2\,p_0\bigl(g_{j(\mathfrak{z})}\bigr)^2
     \bigl(\mathfrak{d}(\mathfrak{z})+1\bigr)\Bigr).
\end{aligned}
\end{equation}
Since $\theta'_\tau\le\theta'$, writing $\theta'_\tau$ in place of $\theta'$ in the factor
$r_{\theta'_\tau}$ only weakens the linear-decay factor.

The function $M_{\lambda,C}$ is the \emph{hole majorant}. It is obtained in two steps from the
term of the history in which $C$ is present: the operator of $C$ is replaced by $1$, and the
factors of the remaining integrand that depend on $C$ are replaced by their supremum over the
variables internal to $C$. Thus, on the level-$k$ region of $C$, the restriction $V_k|_C$ is no
longer weighted. Off $C$ the integrand is that of the history in which $C$ is present, whose
background near $C$ treats $C$ as a hole. There is one majorant for the whole
fibre, and it does not depend on the internal history of $C$.

\emph{The hole-correction bound} is the following statement. For every touch-connected union
$C$ of cells, the number $\delta_C$ in \eqref{8.10:eq-fibre-domination-1} satisfies
\begin{equation}\label{8.10:eq-fibre-domination-3}
  \delta_C\le C_{\delta'}\,\alpha\,\mathcal{N}_M\bigl(\operatorname{hull}_1(C)\bigr)
  \le C_{\delta'}\,\alpha\,3^d\check{R}_{K'}^{\,d}\,\#C .
\end{equation}
In this inequality:
\begin{itemize}
\item $\mathcal{N}_M(\cdot)$ is the number of cubes of side $M$, and each cell contains
  $\check{R}_{K'}^{\,d}$ of them;
\item $\#C$ is the number of cells of $C$;
\item $C_{\delta'}$ and $\alpha$ are constants of \cite{B-CMP122b}, $\alpha$ being
  Ba{\l}aban's constant that carries the logarithms;
\item $\alpha\le\alpha_1(1+\log x)^{a_\alpha}$, with $a_\alpha\in\{0,1\}$, with $\alpha_1$
  depending on $(d,N,L,M)$ only, and with $x$ the inverse square coupling, in the notation
  $x_k=g_k^{-2}$ of the glossary, at the level at which $\alpha$ is read.
\end{itemize}
\end{definition}

In what follows both statements are used as hypotheses, under the names fibre domination and
the hole-correction bound.

\emph{Route.} A proof of fibre domination would have to bound the sum over the fibre in
\eqref{8.10:eq-fibre-domination-1} by the penalty \eqref{8.10:eq-fibre-domination-2} times the
single majorant $M_{\lambda,C}$. It would derive this from the following inputs:
\begin{itemize}
\item the positivity statement for the densities $\rho_\omega$;
\item the zeroth-order bound on the densities;
\item the linear-decay statement, which provides for each creation component the factor of
  $r_{\theta'_\tau}$ with $\theta'$ in place of $\theta'_\tau$. These two bounds are derived from
  Ba{\l}aban's estimates under one-sided upper bounds on the couplings;
\item a comparison of the format of the large-field region $C$ with the regular format;
\item Ba{\l}aban's bound \cite[(1.79)]{B-CMP122b}, applied to $C$ as an arbitrary live region.
  This bound is stated in \cite[(1.79)]{B-CMP122b} and \cite[p.~383]{B-CMP122b} for all
  large-field regions.
\end{itemize}
The form with one majorant per fibre, independent of the internal history of $C$, requires in
addition that the small-field conditions (windows) carried by the histories in which $C$ is
present be dropped in a collar around $C$, consistently with the format of the majorant; this
requirement is read from \cite[(1.70)--(1.72)]{B-CMP122b}. A proof of the hole-correction
bound would have to show that $\delta_C$ is at most $C_{\delta'}\alpha$ times the number of
$M$-cubes of $\operatorname{hull}_1(C)$, with the constants $C_{\delta'}$ and $\alpha$ of
\cite{B-CMP122b}, and then to count these cubes as in the second inequality of
\eqref{8.10:eq-fibre-domination-3}.

\begin{definition}[The small-field windows and their localisation. Proof outstanding]
\label{8.10:def-window-structure}
The terms of Ba{\l}aban's expansion at the terminal level carry small-field characteristic
functions. This definition fixes the seven statements about these characteristic functions, and
about the region rules that place them, that are used in this chapter. The statements refer to the
following objects:
\begin{itemize}
\item the histories, their kept structures $\mathfrak{K}(\cdot)$, the classes $\Lambda_C$ and the
fibres $F(\lambda,C)$ of Definition~\ref{8.10:def-histories-fibres};
\item the terminal-level threshold $\varepsilon$ and length $\ell$ of
Notation~\ref{8.10:not-terminal-couplings}.
\end{itemize}
Throughout, $\mathcal{V}$ is the space of terminal-level configurations $V=(V_b)_b$, indexed by
bonds $b$. For a plaquette $p$, $F_p(V)$ is the plaquette functional of $V$ at $p$, which measures
the deviation of the plaquette variable $V(\partial p)$ from the identity. For an integer $r\ge1$,
$\mathcal{N}_r(p)$ is the neighbourhood of radius $r$ of $p$ in the terminal-level lattice.
Distances between cells are cell-distances. Statements (a)--(g) are used in this chapter as
hypotheses; their proof is outstanding, and each use names this definition and the letter of the
statement.
\begin{enumerate}
\item[(a)] \emph{Localisation of the windows.} Each history $\eta$ has a finite set
$\mathcal{P}_\eta$ of plaquettes, its \emph{window index set}. For each $p\in\mathcal{P}_\eta$ there
is a measurable, gauge-invariant function $\chi^\eta_p\colon\mathcal{V}\to[0,1]$ that depends on $V$
only through $(V_b)_{b\subset\mathcal{N}_{r_W+1}(p)}$. Here $r_W$ is a positive integer, depending
on the coupling only through $\ell$, with
\begin{equation}\label{8.10:eq-window-structure-1}
r_W\le c_r\,\ell ,
\end{equation}
and $c_r\ge1$ depends only on $d$, $L$, $M$ and the constants of Ba{\l}aban's block construction.
Of $r_W$ only the bound \eqref{8.10:eq-window-structure-1} is used.
\item[(b)] \emph{The windows lie in the small-field region.} There is $c_W\ge1$, depending only on
$d$, $N$, $L$ and $M$, such that for every history $\eta$, every $p\in\mathcal{P}_\eta$ and every
$V\in\mathcal{V}$,
\begin{equation}\label{8.10:eq-window-structure-2}
\chi^\eta_p(V)>0\quad\Longrightarrow\quad F_p(V)\le c_W\,\varepsilon .
\end{equation}
\item[(c)] \emph{Regular fields lie in the windows.} There is $B_W\ge1$, depending only on $d$, $N$,
$L$, $M$ and the constants of Ba{\l}aban's construction, with the following property. Let $\eta$ be
a history, $p\in\mathcal{P}_\eta$, $V\in\mathcal{V}$ and $\delta\ge0$. Suppose that
$F_{p'}(V)\le\delta$ for every plaquette $p'\subset\mathcal{N}_{r_W+1}(p)$, and that
$B_W\delta\le\varepsilon$. Then
\begin{equation}\label{8.10:eq-window-structure-3}
\chi^\eta_p(V)=1 .
\end{equation}
A number $\mu\in\,]0,1[$ is fixed once and for all, and we put
$\varepsilon_\mu:=(1-\mu)\varepsilon/B_W$. We say that \emph{$p$ is in window with margin $\mu$
at $V$} if $F_{p'}(V)\le\varepsilon_\mu$ for all plaquettes $p'\subset\mathcal{N}_{r_W+1}(p)$.
\item[(d)] \emph{Placement of the windows and factorisation of the densities.} There is an integer
$p_W\ge1$, depending only on $d$ and $L$, with the following two properties for every history
$\lambda$. First, every plaquette $p$ all of whose vertices lie in cells at cell-distance $>p_W$ from
$\mathfrak{K}(\lambda)$ belongs to $\mathcal{P}_\lambda$. Second, the terminal density $\rho_\lambda$
of $\lambda$ factorises, for every complex source $w$ and every $V\in\mathcal{V}$, as
\begin{equation}\label{8.10:eq-window-structure-4}
\rho_\lambda(w,V)
=\Bigl[\prod_{p\in\mathcal{P}_\lambda}\chi^\lambda_p(V)\Bigr]\,\widehat{\rho}_\lambda(w,V),
\end{equation}
with $\widehat{\rho}_\lambda\ge0$ measurable.
\item[(e)] \emph{Agreement of the windows on fibres.} Let $C$ be a set of cells, $\lambda\in\Lambda_C$
and $\omega\in F(\lambda,C)$, with $\Lambda_C$ and $F(\lambda,C)$ as in
Definition~\ref{8.10:def-histories-fibres}. Let $p$ be a plaquette all of whose vertices lie in
cells at cell-distance $>p_W$ from $C$. Then $p\in\mathcal{P}_\omega$ if and only if
$p\in\mathcal{P}_\lambda$, and in that case
\begin{equation}\label{8.10:eq-window-structure-5}
\chi^\omega_p=\chi^\lambda_p .
\end{equation}
\item[(f)] \emph{Factorisation of the fibre functional.} Let $\mathsf{C}$ denote the set of all
cells, let $C$ be a set of cells and let $\lambda\in\Lambda_C$; we write $\mathsf{M}_{\lambda,C}$ for
the fibre functional of $\lambda$ over $C$. A \emph{core cell} of $C$ is a cell of $C$ at
cell-distance $\ge b_W$ from $\mathsf{C}\setminus C$, where $b_W\ge1$ is an integer depending only
on $d$ and $L$. For every $\lambda\in\Lambda_C$ there is a set
$\mathcal{P}_{\lambda,C}\subset\mathcal{P}_\lambda$ such that
\begin{equation}\label{8.10:eq-window-structure-6}
\mathsf{M}_{\lambda,C}(w,V)
=\Bigl[\prod_{p\in\mathcal{P}_{\lambda,C}}\chi^\lambda_p(V)\Bigr]\,
\widehat{\mathsf{M}}_{\lambda,C}(w,V),
\end{equation}
with $\widehat{\mathsf{M}}_{\lambda,C}\ge0$ measurable. Moreover, $\mathcal{P}_{\lambda,C}$ contains
no plaquette with a vertex in $\mathsf{S}(c)$ for any core cell $c$ of $C$, where $\mathsf{S}(c)$
is the set of vertices attached to the cell $c$.
\item[(g)] \emph{Positivity in the absence of kept structure.} Let $C$ be a set of cells and
$\lambda\in\Lambda_C$ with $\mathfrak{K}(\lambda)=\emptyset$, and let $w$ be a complex source in the
set $\mathcal{U}_\lambda$ of complex sources admitted for $\lambda$. Then
\begin{equation}\label{8.10:eq-window-structure-7}
\widehat{\mathsf{M}}_{\lambda,C}(w,V)>0\qquad\text{for $dV$-almost every }V .
\end{equation}
\end{enumerate}
\end{definition}

Proof outstanding.

\smallskip
\noindent\emph{Route.} A proof has to establish statements (a)--(g) for the small-field
characteristic functions and region rules of Ba{\l}aban's construction at the terminal level. These
objects are defined in \cite[(2.17)--(2.18), \S3]{B-CMP119}, \cite[(1.3)--(1.4)]{B-CMP122a} and
\cite[(1.70)--(1.72)]{B-CMP122b}. The proof has the following parts.
\begin{itemize}
\item For (a): by \cite[(2.17)]{B-CMP119}, the window at a test cube is a function of the local
minimiser. The local minimiser is a function of the field on a neighbourhood of side comparable to
the side of the cube. It remains to bound the radius of that neighbourhood by $c_r\ell$.
\item For (b): the starting point is the small-field condition
\begin{equation*}
|V_{k+1}(\partial p')-1|<2L^2\varepsilon_k
\end{equation*}
of \cite{B-CMP119}, in which $p'$ is a plaquette of the lattice at level $k+1$ and
$\varepsilon_k$ is Ba{\l}aban's small-field threshold at step $k$. It has to be converted into the
bound \eqref{8.10:eq-window-structure-2} on $F_p$, with $c_W$ depending only on $d$, $N$, $L$ and
$M$.
\item For (c): the input is the regularity of the local minimiser in terms of the regularity of
the data, \cite[Theorem~1]{B-CMP102b}. From it the constant $B_W$ has to be extracted.
\item For (d): the region rules place the windows relative to the kept structure
$\mathfrak{K}(\lambda)$ at a distance of four layers of the cubes used by the region rules. This
gives $p_W$ of the size $4L+O(1)$. The factorisation \eqref{8.10:eq-window-structure-4} of
$\rho_\lambda$ also has to be established.
\item For (e): it has to be shown that, at cell-distance $>p_W$ from $C$, the windows of a history
$\omega\in F(\lambda,C)$ are placed by the same rules, and with the same functions, as those of
$\lambda$.
\item For (f): the input is the statement that, on the level-$k$ region of $C$, the field $V_k$
restricted to $C$ is no longer weighted. Consequently no window factor remains on a plaquette
meeting $\mathsf{S}(c)$ for a core cell $c$.
\item For (g): it has to be shown that, when $\mathfrak{K}(\lambda)=\emptyset$, the factors remaining
in $\mathsf{M}_{\lambda,C}$ after the windows are removed are positive wherever those windows hold.
The restriction to $\mathfrak{K}(\lambda)=\emptyset$ in (g) cannot be dropped. Suppose the kept
structure of $\lambda$ contains a large-field region created at step $K'$. Then
$\widehat{\mathsf{M}}_{\lambda,C}$ carries a large-field characteristic function of the top field,
and a positivity statement of the form \eqref{8.10:eq-window-structure-7} cannot hold for all such
$\lambda$.
\end{itemize}

\begin{definition}[The interior--exterior split of the thickened hull]\label{8.10:def-hull-split}
We use the cells, the closed cube $\bar{Q}_c$ of a cell $c$ regarded as a complex, the vertex sets
$\bar{S}(X)$, the closed-cube complex $\bar{K}(X)$ of a set $X$ of cells and the full complexes $K(S)$
of Lemma~\ref{8.10:lem-cube-fullness}. We also use
the configurations $V=(V_b)_b$ of the terminal lattice of Notation~\ref{8.10:not-terminal-lattice}. For a
complex $\mathsf{X}$ we write $\mathsf{X}_1$ for its bonds and $\mathsf{X}_2$ for its plaquettes.

Fix integers $p_1\ge1$ and $p_2\ge1$: $p_1$ is the thickness of the hull and $p_2$ the width of the
collar, both counted in cells. Let $C$ be a touch-connected union of cells, and let
$\operatorname{hull}_{p_1}(C)$ and $\operatorname{hull}_{p_1+p_2}(C)$ be its hulls of thickness $p_1$
and $p_1+p_2$. Put
\begin{equation*}
  R_{\mathrm{int}}:=\operatorname{hull}_{p_1}(C),\qquad
  R_{\mathrm{col}}:=\operatorname{hull}_{p_1+p_2}(C)\setminus R_{\mathrm{int}},
\end{equation*}
and
\begin{equation}\label{8.10:eq-hull-split-1}
  \mathsf{R}:=\bar{K}(R_{\mathrm{int}}),\quad
  \mathsf{N}:=\bar{K}\bigl(\operatorname{hull}_{p_1+p_2}(C)\bigr),\quad
  \mathsf{A}:=\bar{K}(R_{\mathrm{col}}),\quad
  \mathsf{I}:=\mathsf{R}\cap\mathsf{A}.
\end{equation}
Put $\mathsf{N}_0:=\bar{S}(\operatorname{hull}_{p_1+p_2}(C))$. Then $\mathsf{R}$, $\mathsf{N}$, $\mathsf{A}$
and $\mathsf{I}$ are the full complexes
\begin{equation}\label{8.10:eq-hull-split-2}
  K\bigl(\bar{S}(R_{\mathrm{int}})\bigr),\quad K(\mathsf{N}_0),\quad
  K\bigl(\bar{S}(R_{\mathrm{col}})\bigr),\quad
  K\bigl(\bar{S}(R_{\mathrm{int}})\cap\bar{S}(R_{\mathrm{col}})\bigr),
\end{equation}
respectively, and $\mathsf{N}=\mathsf{R}\cup\mathsf{A}$.

\emph{Interior and exterior bonds.} The \emph{interior bonds} are the bonds of $\mathsf{R}_1$. This is
the closed-cube convention: by \eqref{8.10:eq-hull-split-2}, $\mathsf{R}_1$ consists of every bond both
of whose endpoints lie in $\bar{S}(R_{\mathrm{int}})$. All other bonds of the terminal lattice are
\emph{exterior}. We put
\begin{equation*}
  V^{\mathrm{in}}:=(V_b)_{b\in\mathsf{R}_1}\in\mathcal{V}^{\mathrm{in}}:=\mathrm{SU}(N)^{\mathsf{R}_1},\qquad
  V^{\mathrm{ext}}:=(V_b)_{b\notin\mathsf{R}_1}\in\mathcal{V}^{\mathrm{ext}}.
\end{equation*}
Here $\mathcal{V}^{\mathrm{ext}}$ is the product of copies of $\mathrm{SU}(N)$ over the exterior
bonds. The product Haar measure then factorises as $dV=dV^{\mathrm{in}}\,dV^{\mathrm{ext}}$. The
domination bound of this chapter, being pointwise in every variable not attached to $C$, is
indifferent to this split.

\emph{Straddling plaquettes, the exterior collar complex, the filling plaquettes.} Put
\begin{align*}
  \mathrm{Str}&:=\{p\in\mathsf{A}_2\setminus\mathsf{R}_2:\ p\text{ has at least one edge in }\mathsf{R}_1\},\\
  \mathsf{Ex}&:=\{\eta\in\mathsf{A}:\ \text{no edge of }\eta\text{ lies in }\mathsf{R}_1\},\\
  \mathsf{F}&:=\mathsf{R}_2\sqcup\mathrm{Str}.
\end{align*}
We call $\mathrm{Str}$ the \emph{straddling plaquettes}, and $\mathsf{Ex}$ the \emph{exterior collar
complex}. The complex $\mathsf{Ex}$ contains every vertex of $\mathsf{A}$, since a vertex has no edge.
The set $\mathsf{F}$ consists of the plaquettes on which the filling problem of this chapter is posed;
the union defining it is disjoint because $\mathrm{Str}\subset\mathsf{A}_2\setminus\mathsf{R}_2$.
\end{definition}

\begin{proof}
Apply Lemma~\ref{8.10:lem-cube-fullness} with $X=R_{\mathrm{int}}$, with
$X=\operatorname{hull}_{p_1+p_2}(C)$ and with $X=R_{\mathrm{col}}$. Each is a set of cells. This gives
$\mathsf{R}=K(\bar{S}(R_{\mathrm{int}}))$, $\mathsf{N}=K(\mathsf{N}_0)$ and
$\mathsf{A}=K(\bar{S}(R_{\mathrm{col}}))$.

For $\mathsf{I}$ we unwind the definition of the full complex. A cube lies in $K(S)$ exactly when all
its vertices lie in $S$. Hence a cube lies in both $K(\bar{S}(R_{\mathrm{int}}))$ and
$K(\bar{S}(R_{\mathrm{col}}))$ exactly when all its vertices lie in
$\bar{S}(R_{\mathrm{int}})\cap\bar{S}(R_{\mathrm{col}})$. Therefore
\begin{equation*}
  \mathsf{I}=\mathsf{R}\cap\mathsf{A}
  =K\bigl(\bar{S}(R_{\mathrm{int}})\cap\bar{S}(R_{\mathrm{col}})\bigr).
\end{equation*}
This proves \eqref{8.10:eq-hull-split-2}.

For $\mathsf{N}=\mathsf{R}\cup\mathsf{A}$, recall from Lemma~\ref{8.10:lem-cube-fullness} that
$\bar{K}(X)$ is made of the closed cubes $\bar{Q}_c$ with $c\in X$. Consequently
$\bar{K}(X\cup Y)=\bar{K}(X)\cup\bar{K}(Y)$ for any two sets of cells $X$ and $Y$. By the definition
of the hull of a given thickness, the hull of thickness $p_1$ is contained in the hull of thickness
$p_1+p_2$. Hence
$\operatorname{hull}_{p_1+p_2}(C)=R_{\mathrm{int}}\cup R_{\mathrm{col}}$, and so
$\mathsf{N}=\bar{K}(R_{\mathrm{int}})\cup\bar{K}(R_{\mathrm{col}})=\mathsf{R}\cup\mathsf{A}$.
\end{proof}

\begin{lemma}[Decomposition of the collar plaquettes]\label{8.10:lem-collar-plaquettes}
We use the complexes $\mathsf{R}$, $\mathsf{N}$, $\mathsf{A}$ and $\mathsf{I}$ of
Definition~\ref{8.10:def-hull-split}, with the block side $\ell$, the block distance
$\operatorname{dist}_{\blacksquare}$, the hulls $\operatorname{hull}_r$ and the lattice
$\mathbb{T}$ of this chapter, and the closed-cube complexes $\bar{K}(\cdot)$, vertex sets
$\bar{S}(\cdot)$ and closed cubes $\bar{Q}_c$ appearing in
Lemma~\ref{8.10:lem-cube-fullness}. For a complex $\mathsf{X}$ we write $\mathsf{X}_0$,
$\mathsf{X}_1$, $\mathsf{X}_2$ for its sets of vertices, bonds and plaquettes. A bond is
\emph{interior} if it lies in $\mathsf{R}_1$ and \emph{exterior} otherwise. Let
$\mathrm{Str}$ be the set of plaquettes of $\mathsf{A}$ not in $\mathsf{R}$ that have an
interior edge, and $\mathsf{Ex}_2$ the set of plaquettes of $\mathsf{A}$ not in
$\mathsf{R}$ without interior edge. Let $\mathsf{Ex}$ be the set of cells of $\mathsf{A}$ none
of whose edges lies in $\mathsf{R}_1$; since the edges of a plaquette of $\mathsf{R}$ lie in
$\mathsf{R}_1$, the plaquettes of $\mathsf{Ex}$ form the set $\mathsf{Ex}_2$ just defined,
$\mathsf{Ex}_1=\mathsf{A}_1\setminus\mathsf{R}_1$, and $\mathsf{Ex}_0$ is its set of
vertices. Let $\mathsf{F}=\mathsf{R}_2\cup\mathrm{Str}$. Finally $V$ denotes the collar
configuration, $F_P$ the plaquette factor of a plaquette $P$, and $V^{\mathrm{ext}}$ the
restriction of $V$ to exterior bonds.
\begin{itemize}
\item[(a)] $\mathsf{N}_2=\mathsf{R}_2\sqcup\mathrm{Str}\sqcup\mathsf{Ex}_2$. Every cube of
$\mathbb{T}$ containing a bond of $\mathsf{R}_1$ lies in
$\bar{K}(\operatorname{hull}_1(R_{\mathrm{int}}))\subset\mathsf{N}$. Every plaquette of
$\mathbb{T}$ with an edge in $\mathsf{R}_1$ lies in $\mathsf{F}$.
\item[(b)] Every exterior bond of $\mathsf{N}$ lies in $\mathsf{Ex}_1$, and
$\mathsf{Ex}_0=\mathsf{A}_0$.
\item[(c)] Each $P\in\mathrm{Str}$ has exactly one maximal run of exterior edges in its
boundary circuit. This run has two or three edges, and its two endpoints lie in
$\mathsf{I}_0$. For every plaquette $P$ without interior edge, in particular for every
$P\in\mathsf{Ex}_2$, $F_P$ is a function of $V^{\mathrm{ext}}$. For $P\in\mathsf{F}$,
$F_P$ is a function of $V$.
\end{itemize}
\end{lemma}

\begin{proof}
(a) We have $\mathsf{N}=\mathsf{R}\cup\mathsf{A}$, so every plaquette of $\mathsf{N}$ lies
in $\mathsf{R}_2$ or is a plaquette of $\mathsf{A}$ not in $\mathsf{R}$. A plaquette of the
second kind either has an edge in $\mathsf{R}_1$, and then it lies in $\mathrm{Str}$, or it
has no such edge, and then it lies in $\mathsf{Ex}_2$. The three sets are disjoint by their
definitions. This proves the first assertion.

For the second, let $\{y,z\}\in\mathsf{R}_1$. Then $\{y,z\}\subset\bar{Q}_c$ for some
$c\in R_{\mathrm{int}}$. Let $y'$ be a vertex of a cube of $\mathbb{T}$ containing
$\{y,z\}$. Then $|y'-y|_\infty\le1$, and for each coordinate $\nu$,
\begin{equation*}
y'_\nu\in[\ell c_\nu-1,\ \ell c_\nu+\ell+1].
\end{equation*}
Hence $y'\in\bar{S}(c'')$ for some block $c''$ with
$\operatorname{dist}_{\blacksquare}(c'',c)\le1$. Such a $c''$ satisfies
\begin{equation*}
c''\in\operatorname{hull}_1(R_{\mathrm{int}})\subset\operatorname{hull}_{p_1+p_2}(C).
\end{equation*}
So every vertex of the cube lies in $\bar{S}(\operatorname{hull}_1(R_{\mathrm{int}}))$. By
fullness of closed-cube complexes (Lemma~\ref{8.10:lem-cube-fullness}), applied with
$X=\operatorname{hull}_1(R_{\mathrm{int}})$, the cube lies in
$\bar{K}(\operatorname{hull}_1(R_{\mathrm{int}}))$. By the inclusion of hulls just displayed,
this complex is contained in $\mathsf{N}$.

For the third assertion, let $P$ be a plaquette of $\mathbb{T}$ with an edge in
$\mathsf{R}_1$. Then $P$ is a face of a cube of $\mathbb{T}$ containing that edge. By the
previous paragraph $P$ lies in $\mathsf{N}_2$. Since $P$ has an interior edge, the first
assertion places it in $\mathsf{R}_2$ or in $\mathrm{Str}$, that is, in $\mathsf{F}$.

(b) Let $b$ be an exterior bond of $\mathsf{N}$. Then $b\notin\mathsf{R}$, and since
$\mathsf{N}=\mathsf{R}\cup\mathsf{A}$ we get $b\in\mathsf{A}_1$. The only edge of $b$ is $b$
itself, which is not in $\mathsf{R}_1$. So $b$ is a bond of $\mathsf{A}$ with no interior
edge, that is, $b\in\mathsf{Ex}_1$. A vertex has no edge, so every vertex of $\mathsf{A}$
belongs to $\mathsf{Ex}$, and $\mathsf{Ex}_0=\mathsf{A}_0$.

(c) Let $P\in\mathrm{Str}$. An interior edge of $P$ lies in $\mathsf{R}_1$, so both its
endpoints lie in $\mathsf{R}_0$. Since $P\in\mathsf{A}$, all vertices of $P$ lie in
$\mathsf{A}_0$. Because $\mathsf{I}=\mathsf{R}\cap\mathsf{A}$, the endpoints of interior
edges of $P$ therefore lie in $\mathsf{I}_0$.

We now count the vertices of $P$ in $\mathsf{R}_0$. By Lemma~\ref{8.10:lem-cube-fullness},
$\mathsf{R}=\bar{K}(R_{\mathrm{int}})=K(\bar{S}(R_{\mathrm{int}}))$. So a cell of
$\mathbb{T}$ all of whose vertices lie in $\mathsf{R}_0$ belongs to $\mathsf{R}$.
Since $P$ has an interior edge, at least two of its four vertices lie in $\mathsf{R}_0$.

\emph{Exactly two vertices in $\mathsf{R}_0$.} These two vertices are adjacent, because $P$
has an interior edge. Each of the other three edges of the boundary circuit has an endpoint
outside $\mathsf{R}_0$, so it is exterior. These three edges are consecutive and form the
unique maximal run of exterior edges. Its endpoints are the two endpoints of the interior
edge, which lie in $\mathsf{I}_0$.

\emph{Exactly three vertices in $\mathsf{R}_0$.} The two edges of $P$ joining these three
vertices have both endpoints in $\mathsf{R}_0$, so they are interior by the fullness just
recalled. The complementary two edges contain the fourth vertex, which is not in
$\mathsf{R}_0$, so they are exterior. They form the unique maximal run of exterior edges,
with two edges. Its endpoints are endpoints of interior edges and hence lie in
$\mathsf{I}_0$.

\emph{Four vertices in $\mathsf{R}_0$.} This case cannot occur: fullness would put $P$ in
$\mathsf{R}_2$, contrary to $P\notin\mathsf{R}$.

For the last sentence of (c), note that $F_P$ depends on $V$ through $V(\partial P)$, which
involves only the bonds of $P$. If $P$ has no interior edge, all its bonds are exterior,
so $F_P$ is a function of $V^{\mathrm{ext}}$. This applies in particular to
$P\in\mathsf{Ex}_2$, which by definition has no interior edge. For $P\in\mathsf{F}$, the
bonds of $P$ are bonds on which $V$ is given, so $F_P$, being a function of
$V(\partial P)$, is a function of $V$.
\end{proof}

\begin{definition}[The two local integrals and the admissible exteriors]
\label{8.10:def-admissible-exteriors}
Let $C$ be the set whose thickened hull is split in Definition~\ref{8.10:def-hull-split}, and let
$\Lambda_C$ be the index set attached to $C$. For $\lambda\in\Lambda_C$ let $\mathcal{U}_\lambda$ be
the parameter domain and $d\xi$ the measure it carries, $\rho_\lambda$ the density, $\tau_\lambda$
the total weight of $\lambda$, that is, the integral of $\rho_\lambda$ over all its variables, and
$\mathcal{P}_\lambda$ the set of windows attached to $\lambda$. Let $M_{\lambda,C}$ be the hole
majorant of Definition~\ref{8.10:def-fibre-domination}, whose proof is outstanding. Let
$\chi^\lambda_p$, $p\in\mathcal{P}_\lambda$, be the window functions, $N_r(p)$ the neighbourhoods
of $p$ and $r_W$ the localisation radius, as in Definition~\ref{8.10:def-window-structure}, whose
proof is outstanding. Call a bond \emph{interior} if it belongs to the interior part of the
interior--exterior split of Definition~\ref{8.10:def-hull-split}, and \emph{exterior} otherwise;
write $V=(V^{\mathrm{in}},V^{\mathrm{ext}})$ for the corresponding splitting of the bond variables
into the variables of interior and of exterior bonds. Fix $\lambda\in\Lambda_C$ and
$\xi\in\mathcal{U}_\lambda$.
\begin{itemize}
\item[(a)] The \emph{two local integrals} are
\begin{equation}\label{8.10:eq-admissible-exteriors-1}
\begin{aligned}
a_\lambda(V^{\mathrm{ext}},\xi)
&:=\int\rho_\lambda(\xi,V^{\mathrm{in}},V^{\mathrm{ext}})\,dV^{\mathrm{in}},\\
m_\lambda(V^{\mathrm{ext}},\xi)
&:=\int M_{\lambda,C}(\xi,V^{\mathrm{in}},V^{\mathrm{ext}})\,dV^{\mathrm{in}}.
\end{aligned}
\end{equation}
They are non-negative measurable functions of $(V^{\mathrm{ext}},\xi)$, and
\begin{equation}\label{8.10:eq-admissible-exteriors-2}
\int a_\lambda\,dV^{\mathrm{ext}}\,d\xi=\tau_\lambda,
\qquad
\int m_\lambda\,dV^{\mathrm{ext}}\,d\xi<\infty .
\end{equation}
\item[(b)] The \emph{held window set} is
\begin{equation*}
\mathcal{P}^{\mathrm{held}}_\lambda
:=\bigl\{p\in\mathcal{P}_\lambda:\ \text{every bond } b\subset N_{r_W+1}(p)
\text{ is exterior}\bigr\}.
\end{equation*}
For $p\in\mathcal{P}^{\mathrm{held}}_\lambda$ the window function $\chi^\lambda_p$ is a function of
$V^{\mathrm{ext}}$ alone. The windows $p\in\mathcal{P}^{\mathrm{held}}_\lambda$ are called held
because their positivity condition involves exterior variables only and is kept, in (c), as a
condition on $V^{\mathrm{ext}}$.
\item[(c)] The set of \emph{admissible exteriors} is
\begin{equation}\label{8.10:eq-admissible-exteriors-3}
\operatorname{Adm}(\lambda;C):=\bigl\{V^{\mathrm{ext}}:\ \chi^\lambda_p(V^{\mathrm{ext}})>0
\text{ for every } p\in\mathcal{P}^{\mathrm{held}}_\lambda\bigr\}.
\end{equation}
\end{itemize}
\end{definition}

\begin{proof}
For (a): the integrands $\rho_\lambda$ and $M_{\lambda,C}$ are non-negative measurable functions of
$(\xi,V^{\mathrm{in}},V^{\mathrm{ext}})$. Tonelli's theorem, applied to each of them on the product of
the space of $(V^{\mathrm{ext}},\xi)$ with the space of $V^{\mathrm{in}}$, gives that the partial
integrals $a_\lambda$ and $m_\lambda$ in \eqref{8.10:eq-admissible-exteriors-1} are non-negative
measurable functions of $(V^{\mathrm{ext}},\xi)$, and that their integrals over
$(V^{\mathrm{ext}},\xi)$ equal the integrals of $\rho_\lambda$ and of $M_{\lambda,C}$ over all the
variables $(\xi,V^{\mathrm{in}},V^{\mathrm{ext}})$. The first of these full integrals equals
$\tau_\lambda$ by the definition of the total weight $\tau_\lambda$, and the second is finite by the
integrability of the hole majorant $M_{\lambda,C}$ of Definition~\ref{8.10:def-fibre-domination},
whose proof is outstanding; this is \eqref{8.10:eq-admissible-exteriors-2}.

For (b): let $p\in\mathcal{P}^{\mathrm{held}}_\lambda$. By the localisation property of the windows in
Definition~\ref{8.10:def-window-structure}, whose proof is outstanding, the window function
$\chi^\lambda_p$ depends only on the variables of the bonds contained in $N_{r_W+1}(p)$. Each such bond
is exterior by the definition of $\mathcal{P}^{\mathrm{held}}_\lambda$, so $\chi^\lambda_p$ is a
function of $V^{\mathrm{ext}}$ alone. Consequently every condition
$\chi^\lambda_p(V^{\mathrm{ext}})>0$ in \eqref{8.10:eq-admissible-exteriors-3} is a condition on
$V^{\mathrm{ext}}$ only, and $\operatorname{Adm}(\lambda;C)$ is a well-defined set of exterior
configurations.
\end{proof}

\begin{definition}[Collar functionals, margined fillings and the good set]\label{8.10:def-good-set}
Let $\mathsf{R}$, $\mathsf{N}$, $\mathsf{A}$, $\mathsf{I}$ be the sets of Definition~\ref{8.10:def-hull-split}.
Their sets of vertices, bonds and plaquettes carry the subscripts $0$, $1$, $2$.
Let $\mathsf{N}_2=\mathsf{R}_2\sqcup\mathrm{Str}\sqcup\mathsf{Ex}_2$, the set $\mathsf{F}$ and the plaquette
functions $F_p$ be as in Lemma~\ref{8.10:lem-collar-plaquettes}.
Configurations are written $V=(V^{\mathrm{in}},V^{\mathrm{ext}})$, with $V^{\mathrm{in}}$ the interior and
$V^{\mathrm{ext}}$ the exterior bond variables. We write $\mathcal{V}^{\mathrm{in}}$ for the space of interior
configurations and $\mathcal{V}^{\mathrm{ext}}$ for the space of exterior configurations.
Let $\operatorname{Adm}(\lambda;C)$ be the set of admissible exteriors of
Definition~\ref{8.10:def-admissible-exteriors}.
Fix a margin index $\mu$ and its margin $\varepsilon_\mu$.

\begin{enumerate}
\item[(a)] \emph{Near collar plaquettes.} A plaquette $p'\in\mathsf{Ex}_2$ belongs to
$\mathsf{Ex}_2^{\mathrm{near}}$ if every vertex of $p'$ lies within sup-distance $2r_W+3$ of $\mathsf{R}_0$.
The \emph{near-collar margin set} is
\begin{equation}\label{8.10:eq-good-set-1}
G^{\mathrm{mar}}_{\mu}:=\bigl\{V^{\mathrm{ext}}:\ F_{p'}(V^{\mathrm{ext}})\le\varepsilon_\mu
\text{ for every }p'\in\mathsf{Ex}_2^{\mathrm{near}}\bigr\}.
\end{equation}

\item[(b)] \emph{Cell-smoothed curvature.} Let $\mathfrak{a}\in R_{\mathrm{col}}$ be a collar cell, of side
$\ell$ and with set of sites $S(\mathfrak{a})$, and let
$\mu'<\nu'$ be two directions.
Let $P_{\mathfrak{a},\mu'\nu'}$ be the set of plaquettes $p=y+\{0,1\}^{\{\mu',\nu'\}}$ with
$y\in S(\mathfrak{a})$ and $p\in\mathsf{Ex}_2$; it has at most $\ell^d$ elements.
Fix unitaries $\mathfrak{t}_p(V^{\mathrm{ext}})$, $p\in P_{\mathfrak{a},\mu'\nu'}$, depending measurably on
$V^{\mathrm{ext}}$.
For instance, one may take transports along fixed exterior edge paths from the corner $\ell\mathfrak{a}$, or
$\mathfrak{t}_p\equiv 1$. We put
\begin{equation}\label{8.10:eq-good-set-2}
f_{\mathfrak{a},\mu'\nu'}(V^{\mathrm{ext}}):=\ell^{-d}\Bigl\|\sum_{p\in P_{\mathfrak{a},\mu'\nu'}}
\mathfrak{t}_p\bigl(V(\partial p)-1\bigr)\mathfrak{t}_p^{-1}\Bigr\|.
\end{equation}
For a number $f$ we put
\begin{equation}\label{8.10:eq-good-set-3}
G^{\mathrm{av}}_{f}:=\bigl\{V^{\mathrm{ext}}:\ f_{\mathfrak{a},\mu'\nu'}(V^{\mathrm{ext}})\le f
\text{ for all }\mathfrak{a}\in R_{\mathrm{col}}\text{ and all }\mu'<\nu'\bigr\}.
\end{equation}
No gauge invariance of $f_{\mathfrak{a},\mu'\nu'}$ is claimed or used.

\item[(c)] \emph{The threshold.} The number $f_*>0$ is the threshold in the collar condition on exteriors,
which requires the cell-smoothed curvature of $V^{\mathrm{ext}}$ on the collar to be at most $f_*$; we write
$G_{f_*}(C)$ for the set of exteriors that satisfy the collar condition.
Only the positivity of $f_*$ is used below.

\item[(d)] \emph{Collar energy.} We put
\begin{equation}\label{8.10:eq-good-set-5}
\mathsf{Q}_{\mathrm{col}}(V^{\mathrm{ext}}):=\sum_{p\in\mathsf{Ex}_2}F_p(V^{\mathrm{ext}})^2.
\end{equation}
For $E>0$ we put $G^{\mathrm{en}}_E:=\{\mathsf{Q}_{\mathrm{col}}\le E\}$. This set is not one of the
conditions defining the good set in (f).

\item[(e)] \emph{Fillings.} For $V=(V^{\mathrm{in}},V^{\mathrm{ext}})$ the \emph{lattice filling energy} is
\begin{equation*}
\mathsf{E}(V):=\sum_{p\in\mathsf{F}}F_p(V)^2 .
\end{equation*}
We call $V^{\mathrm{in}}$ a \emph{margin-$\mu$ filling} of $V^{\mathrm{ext}}$ if $F_p(V)\le\varepsilon_\mu$ for
every $p\in\mathsf{F}$. For $\mathfrak{E}\ge0$ we put
\begin{equation}\label{8.10:eq-good-set-6}
\operatorname{Fill}_{\mu,\mathfrak{E}}(C):=\bigl\{V^{\mathrm{ext}}\in G^{\mathrm{mar}}_{\mu}:\
\mathsf{E}(V^{\mathrm{in}},V^{\mathrm{ext}})\le\mathfrak{E}\text{ for some margin-}\mu\text{ filling }
V^{\mathrm{in}}\text{ of }V^{\mathrm{ext}}\bigr\}
\end{equation}
and
\begin{equation}\label{8.10:eq-good-set-7}
E^{\mu}_{\mathrm{fill}}(V^{\mathrm{ext}}):=\inf\bigl\{\mathsf{E}(V^{\mathrm{in}},V^{\mathrm{ext}}):\
V^{\mathrm{in}}\text{ a margin-}\mu\text{ filling of }V^{\mathrm{ext}}\bigr\}\in[0,\infty],
\end{equation}
with the convention that the infimum of the empty set is $+\infty$.

\item[(f)] \emph{The good set.} The good set is used for those $\lambda$ for which the freeness condition
$K_{\mathrm{free}}(C;\lambda;p_1+p_2+p_W)$ of the region $C$ for $\lambda$ at distance $p_1+p_2+p_W$
holds. It is
\begin{equation}\label{8.10:eq-good-set-8}
G_{\lambda,C}(\mathfrak{E}):=\operatorname{Adm}(\lambda;C)\cap G^{\mathrm{av}}_{f_*}\cap
\operatorname{Fill}_{\mu,\mathfrak{E}}(C).
\end{equation}
\end{enumerate}

These objects have the following properties.
\begin{itemize}
\item[(i)] Each $f_{\mathfrak{a},\mu'\nu'}$ is a measurable function of $V^{\mathrm{ext}}$, and
\begin{equation}\label{8.10:eq-good-set-9}
f_{\mathfrak{a},\mu'\nu'}\le\ell^{-d}\sum_{p\in P_{\mathfrak{a},\mu'\nu'}}F_p
\le\max_{p\in P_{\mathfrak{a},\mu'\nu'}}F_p .
\end{equation}
\item[(ii)] The sets $G^{\mathrm{mar}}_\mu$ and $\operatorname{Fill}_{\mu,\mathfrak{E}}(C)$ are closed.
\item[(iii)] The function $E^{\mu}_{\mathrm{fill}}$ is lower semicontinuous.
\item[(iv)] The good set satisfies
\begin{equation*}
G_{\lambda,C}(\mathfrak{E})\subset G^{\mathrm{mar}}_\mu\cap G^{\mathrm{av}}_{f_*}.
\end{equation*}
In particular, $G_{\lambda,C}(\mathfrak{E})$ lies inside the set $G_{f_*}(C)$ of (c).
\end{itemize}
\end{definition}

\begin{proof}
\emph{(i).} By Lemma~\ref{8.10:lem-collar-plaquettes}(c), a plaquette $p\in\mathsf{Ex}_2$ has no interior edge.
Hence $V(\partial p)$ is a product of exterior bond variables and depends continuously on $V^{\mathrm{ext}}$.
The unitaries $\mathfrak{t}_p$ depend measurably on $V^{\mathrm{ext}}$.
The sum in \eqref{8.10:eq-good-set-2} is therefore a measurable function of $V^{\mathrm{ext}}$, and so is its norm.
Conjugation by a unitary preserves the norm, so each summand
$\mathfrak{t}_p(V(\partial p)-1)\mathfrak{t}_p^{-1}$ has norm $F_p$.
The triangle inequality gives the first inequality in \eqref{8.10:eq-good-set-9}.
The second holds because $P_{\mathfrak{a},\mu'\nu'}$ has at most $\ell^d$ elements: a sum of at most $\ell^d$
terms, multiplied by $\ell^{-d}$, is at most the largest term.

\emph{(ii).} Each $F_p$ is a continuous function of $V$, and hence so is $\mathsf{E}$.
$G^{\mathrm{mar}}_\mu$ is the intersection of the finitely many sets $\{F_{p'}\le\varepsilon_\mu\}$,
$p'\in\mathsf{Ex}_2^{\mathrm{near}}$. Each of these is the preimage of a closed interval under a continuous
function, so $G^{\mathrm{mar}}_\mu$ is closed.
For $\mathfrak{E}\ge0$ let
\begin{equation*}
\Phi_{\mu,\mathfrak{E}}:=\bigl\{(V^{\mathrm{in}},V^{\mathrm{ext}}):\ F_p(V)\le\varepsilon_\mu
\text{ for all }p\in\mathsf{F},\ \mathsf{E}(V)\le\mathfrak{E}\bigr\}.
\end{equation*}
This set is closed for the same reason.
The space $\mathcal{V}^{\mathrm{in}}$ is a finite product of copies of $\mathrm{SU}(N)$, hence compact.
Let $\pi:\mathcal{V}^{\mathrm{in}}\times\mathcal{V}^{\mathrm{ext}}\to\mathcal{V}^{\mathrm{ext}}$ be the
projection; since $\mathcal{V}^{\mathrm{in}}$ is compact, $\pi$ maps closed sets to closed sets.
To see this, let $V^{\mathrm{ext}}$ lie outside the projection $\pi(\Phi_{\mu,\mathfrak{E}})$.
Then every point of the compact fibre $\mathcal{V}^{\mathrm{in}}\times\{V^{\mathrm{ext}}\}$ has a product
neighbourhood disjoint from $\Phi_{\mu,\mathfrak{E}}$.
Finitely many of these cover the fibre.
The intersection of their second factors is a neighbourhood of $V^{\mathrm{ext}}$ disjoint from
$\pi(\Phi_{\mu,\mathfrak{E}})$.
By \eqref{8.10:eq-good-set-6},
\begin{equation*}
\operatorname{Fill}_{\mu,\mathfrak{E}}(C)=G^{\mathrm{mar}}_\mu\cap\pi(\Phi_{\mu,\mathfrak{E}}),
\end{equation*}
which is closed.

\emph{(iii).} We claim that $\{E^{\mu}_{\mathrm{fill}}\le\mathfrak{E}\}=\pi(\Phi_{\mu,\mathfrak{E}})$ for every
$\mathfrak{E}\ge0$.

The inclusion $\supset$ follows directly from \eqref{8.10:eq-good-set-7}.

For $\subset$, suppose $E^{\mu}_{\mathrm{fill}}(V^{\mathrm{ext}})\le\mathfrak{E}<\infty$.
Then the set of margin-$\mu$ fillings of $V^{\mathrm{ext}}$ is non-empty.
It is closed in the compact space $\mathcal{V}^{\mathrm{in}}$, being the intersection of the sets
$\{V^{\mathrm{in}}: F_p(V^{\mathrm{in}},V^{\mathrm{ext}})\le\varepsilon_\mu\}$, $p\in\mathsf{F}$.
The continuous function $\mathsf{E}(\cdot,V^{\mathrm{ext}})$ therefore attains its infimum on this set.
The minimiser $V^{\mathrm{in}}$ satisfies $(V^{\mathrm{in}},V^{\mathrm{ext}})\in\Phi_{\mu,\mathfrak{E}}$, which
proves the claim.

The sublevel sets of $E^{\mu}_{\mathrm{fill}}$ at levels $\mathfrak{E}\ge0$ are thus closed by (ii).
At negative levels they are empty.
A function with values in $[0,\infty]$ all of whose sublevel sets are closed is lower semicontinuous.

\emph{(iv).} By \eqref{8.10:eq-good-set-6}, $\operatorname{Fill}_{\mu,\mathfrak{E}}(C)\subset
G^{\mathrm{mar}}_\mu$. Intersecting \eqref{8.10:eq-good-set-8} with this inclusion gives
\begin{equation*}
G_{\lambda,C}(\mathfrak{E})\subset G^{\mathrm{mar}}_\mu\cap G^{\mathrm{av}}_{f_*}.
\end{equation*}
The condition defining $G^{\mathrm{av}}_{f_*}$ is the transcription
\eqref{8.10:eq-good-set-2}--\eqref{8.10:eq-good-set-3} of the collar condition of (c),
so $G^{\mathrm{mar}}_\mu\cap G^{\mathrm{av}}_{f_*}\subset G_{f_*}(C)$, and hence
$G_{\lambda,C}(\mathfrak{E})\subset G_{f_*}(C)$.
\end{proof}

\begin{lemma}[The good set places a margined competitor in every window]\label{8.10:lem-margined-competitor}
Assume the inner-window clause and the localisation clause of
Definition~\ref{8.10:def-window-structure}, which fixes the small-field windows $\chi^\lambda_p$;
both clauses are stated there without proof, and their proof is outstanding.
Assume moreover the lower bound on the collar width $p_2$,
\begin{equation}\label{8.10:eq-margined-competitor-1}
(p_2-1)\ell\ \ge\ 2r_W+3,\qquad p_2\ge 3 .
\end{equation}
Let $V^{\mathrm{ext}}\in G_{\lambda,C}(\mathfrak{E})$, the good set of
Definition~\ref{8.10:def-good-set}, let $V^{\mathrm{in}}$ be a margin-$\mu$ filling with energy
$\mathsf{E}\le\mathfrak{E}$, in the sense of Definition~\ref{8.10:def-good-set}, and put
$V:=(V^{\mathrm{in}},V^{\mathrm{ext}})$. Then
\begin{itemize}
\item[(i)] every $p\in\mathcal{P}_\lambda\setminus\mathcal{P}^{\mathrm{held}}_\lambda$ has
$\chi^\lambda_p(V)=1$;
\item[(ii)] every $p\in\mathcal{P}^{\mathrm{held}}_\lambda$ has
$\chi^\lambda_p(V)=\chi^\lambda_p(V^{\mathrm{ext}})>0$;
\item[(iii)] $E^\mu_{\mathrm{fill}}(V^{\mathrm{ext}})\le\mathfrak{E}$.
\end{itemize}
In words: the set of interior configurations in $\mathcal{V}^{\mathrm{in}}$ satisfying the window
constraints of the absent term contains the competitor $V^{\mathrm{in}}$, with margin $\mu$ on every
window it influences, all held windows positive, and filling energy at most $\mathfrak{E}$.
Nothing is asserted about exterior configurations outside $G_{\lambda,C}(\mathfrak{E})$.
\end{lemma}

\begin{proof}
\emph{Clause \textup{(i)}.} Let $p\in\mathcal{P}_\lambda\setminus\mathcal{P}^{\mathrm{held}}_\lambda$.
Since $p$ is not held, some bond $b\subset N_{r_W+1}(p)$ is interior. Hence some vertex $y_0$ of $p$
satisfies $|y_0-w_0|_\infty\le r_W+1$ for an endpoint $w_0$ of $b$ lying in
$\mathsf{R}_0=\bar S(R_{\mathrm{int}})$.

Let $p'\subset N_{r_W+1}(p)$ be a plaquette and $z$ a vertex of $p'$. Then
$|z-y_1|_\infty\le r_W+1$ for some vertex $y_1$ of $p$, and $|y_1-y_0|_\infty\le 1$ because $y_0$ and
$y_1$ are vertices of the same plaquette. By the triangle inequality,
\begin{equation*}
|z-w_0|_\infty\ \le\ |z-y_1|_\infty+|y_1-y_0|_\infty+|y_0-w_0|_\infty
\ \le\ (r_W+1)+1+(r_W+1)\ =\ 2r_W+3 .
\end{equation*}
Now $w_0\in\bar S(c)\subset S(\operatorname{hull}_1\{c\})$ for some $c\in R_{\mathrm{int}}$; say
$w_0\in S(c')$ with $\operatorname{dist}_{\blacksquare}(c',c)\le 1$. By
\eqref{8.10:eq-margined-competitor-1}, $2r_W+3\le(p_2-1)\ell$, so $|z-w_0|_\infty\le(p_2-1)\ell$ with
$p_2-1$ a non-negative integer and $w_0\in S(c')$. Lemma~\ref{8.10:lem-steps-cells}, applied with
$k=p_2-1$, gives $z\in S(c'')$ for some $c''$ with $\operatorname{dist}_{\blacksquare}(c'',c')\le p_2-1$.
Since $\operatorname{dist}_{\blacksquare}(c',c)\le 1$ and $c\in R_{\mathrm{int}}$, the cell $c''$ lies
within cell distance $p_2$ of $R_{\mathrm{int}}$, that is,
\begin{equation*}
c''\in\operatorname{hull}_{p_2}(R_{\mathrm{int}})=\operatorname{hull}_{p_1+p_2}(C),
\end{equation*}
and therefore $z\in\mathsf{N}_0$.

As $z$ was an arbitrary vertex of $p'$, all vertices of $p'$ lie in $\mathsf{N}_0$. Hence
$p'\in\mathsf{N}_2$, by the lemma on the layers $\mathsf{N}_0$ and $\mathsf{N}_2$ of the collar,
which places in $\mathsf{N}_2$ every plaquette all of whose vertices lie in $\mathsf{N}_0$. By
Lemma~\ref{8.10:lem-collar-plaquettes}(a),
$\mathsf{N}_2=\mathsf{R}_2\sqcup\mathrm{Str}\sqcup\mathsf{Ex}_2$.
\begin{itemize}
\item If $p'\in\mathsf{F}$, then $F_{p'}(V)\le\varepsilon_\mu$, because $V^{\mathrm{in}}$ is a
margin-$\mu$ filling.
\item If $p'\in\mathsf{Ex}_2$, then all vertices of $p'$ lie within distance $2r_W+3$ of
$\mathsf{R}_0$, by the bound on $|z-w_0|_\infty$ above. Hence $p'\in\mathsf{Ex}_2^{\mathrm{near}}$.
By Lemma~\ref{8.10:lem-collar-plaquettes}(c), $F_{p'}$ is a function of $V^{\mathrm{ext}}$ alone,
since $p'$ has no interior edge. Definition~\ref{8.10:def-good-set} gives
$G_{\lambda,C}(\mathfrak{E})\subset\operatorname{Fill}_{\mu,\mathfrak{E}}(C)\subset G^{\mathrm{mar}}_\mu$,
and therefore $F_{p'}(V)=F_{p'}(V^{\mathrm{ext}})\le\varepsilon_\mu$.
\end{itemize}
So every plaquette $p'\subset N_{r_W+1}(p)$ has $F_{p'}(V)\le\varepsilon_\mu$, and consequently
\begin{equation*}
B_W\,F_{p'}(V)\ \le\ B_W\varepsilon_\mu\ =\ (1-\mu)\varepsilon\ \le\ \varepsilon .
\end{equation*}
The inner-window clause of Definition~\ref{8.10:def-window-structure}, whose proof is outstanding,
now gives $\chi^\lambda_p(V)=1$.

\emph{Clause \textup{(ii)}.} Let $p\in\mathcal{P}^{\mathrm{held}}_\lambda$. By the localisation
clause of Definition~\ref{8.10:def-window-structure}, whose proof is outstanding,
$\chi^\lambda_p$ depends only on exterior bonds. Hence
$\chi^\lambda_p(V)=\chi^\lambda_p(V^{\mathrm{ext}})$. By Definition~\ref{8.10:def-good-set},
$V^{\mathrm{ext}}\in G_{\lambda,C}(\mathfrak{E})\subset\operatorname{Adm}(\lambda;C)$, and
admissibility gives $\chi^\lambda_p(V^{\mathrm{ext}})>0$.

\emph{Clause \textup{(iii)}.} The inequality
$E^\mu_{\mathrm{fill}}(V^{\mathrm{ext}})\le\mathfrak{E}$ is the definition of
$E^\mu_{\mathrm{fill}}$ in Definition~\ref{8.10:def-good-set}.
\end{proof}

In this lemma and its proof the set $R_{\mathrm{int}}$ of cells is the one in the interior--exterior split of the thickened hull (Definition~\ref{8.10:def-hull-split}), where the complexes
\[
\mathsf{R}=\bar{K}(R_{\mathrm{int}}),\qquad \mathsf{A}=\bar{K}(R_{\mathrm{col}}),\qquad \mathsf{I}=\mathsf{R}\cap\mathsf{A}
\]
are fixed. The notations $K(\cdot)$, $\bar{K}(\cdot)$, $\bar{S}(\cdot)$ and $\bar{Q}_c$ are those of the fullness lemma for closed-cube complexes (Lemma~\ref{8.10:lem-cube-fullness}). The cell side $\ell$, the collar $R_{\mathrm{col}}$, the hulls $\operatorname{hull}_p(\cdot)$, the depth $p_2$ and the relation ``touches'' between cells are likewise those of Definition~\ref{8.10:def-hull-split}.

Unit cubes (closed, of any dimension $0,\dots,d$) are those of Lemma~\ref{8.10:lem-cube-fullness}. For a complex $\mathsf{K}$ of closed cubes, $\mathsf{K}_0$ denotes the set of its vertices.

We regard
\[
R:=R_{\mathrm{int}}
\]
as the closed subset of the torus, or of $\mathbb{R}^d$ when the cells are placed in $\mathbb{R}^d$, that is covered by its closed cells; this is the set $\bar{S}(R_{\mathrm{int}})$ of Lemma~\ref{8.10:lem-cube-fullness}. We write $R^{\circ}$ for its interior and
\[
\partial R:=R\setminus R^{\circ}
\]
for its boundary. We write $W$ for the complement of $R^{\circ}$, which is a closed set.

\begin{lemma}[Structure of the hull's boundary]\label{8.10:lem-hull-boundary}
With these notations the following hold.
\begin{itemize}
\item[(a)] Every unit cube $D$ satisfies $D\subset R$ or $D\subset W$. Moreover, $D\cap R^{\circ}\neq\emptyset$ if and only if every unit $d$-cube containing $D$ is contained in $R$.
\item[(b)] The set $\partial R$ is a union of unit cubes. Put
\begin{equation*}
\mathsf{R}^{\mathrm{deep}}:=\{D\in\mathsf{R}: D\cap R^{\circ}\neq\emptyset\},
\end{equation*}
the \emph{deep} cubes, and
\begin{equation*}
\mathsf{R}^{\partial}:=\mathsf{R}\setminus\mathsf{R}^{\mathrm{deep}}
=\{D\in\mathsf{R}: D\subset\partial R\},
\end{equation*}
the \emph{boundary} cubes. Then $\mathsf{R}^{\mathrm{deep}}$ is the complement in $\mathsf{R}$ of a subcomplex. That subcomplex is
\[
\mathsf{R}^{\partial}=K(\partial R)=\mathsf{R}\cap K(W).
\]
\item[(c)] Statements (a) and (b) hold also for the cell structure in place of the unit structure. The cells are the closed cubes of side $\ell$ with corners in $\ell\mathbb{Z}^d$.
\item[(d)] A vertex $y\in\mathsf{R}_0$ is deep, that is $y\in R^{\circ}$, if and only if
\begin{equation}\label{8.10:eq-hull-boundary-1}
y+\{-1,0,1\}^d\subset\bar{S}(R_{\mathrm{int}}).
\end{equation}
The vertices of $\mathsf{R}_0$ that are not deep are exactly those of $\mathsf{I}_0=\mathsf{R}_0\cap\mathsf{A}_0$. This holds for every $p_2\ge1$.
\end{itemize}
\end{lemma}

\begin{proof}
We call the relative interior $\operatorname{relint}D$ of a unit cube $D$ an \emph{open unit face}.

The open unit faces partition space. A closed cube is the union of the relative interiors of its faces. Hence $R$, being a union of closed unit $d$-cubes, is a union of open unit faces.

\emph{(a)} Suppose first that $D\not\subset R$. Then $\operatorname{relint}D$ is an open unit face that is not among those making up $R$. Indeed, if it were one of them, then $D$, the closure of $\operatorname{relint}D$, would lie in the closed set $R$. Since the open unit faces partition space, $\operatorname{relint}D\cap R=\emptyset$.

No point of $D$ then lies in $R^{\circ}$. A point of $D$ in $R^{\circ}$ would have a neighbourhood contained in $R$. That neighbourhood meets $\operatorname{relint}D$, because $D$ is the closure of $\operatorname{relint}D$. This is impossible, so $D\subset W$.

Suppose next that $D\subset R$, and let $q\in\operatorname{relint}D$. The closed unit $d$-cubes through $q$ are exactly the unit $d$-cubes containing $D$, and near $q$ they cover space. Hence $q\in R^{\circ}$ if and only if every closed unit $d$-cube through $q$ lies in $R$. This condition does not depend on the choice of $q$ in $\operatorname{relint}D$. Therefore either $\operatorname{relint}D\subset R^{\circ}$ or $\operatorname{relint}D\cap R^{\circ}=\emptyset$.

In the latter case no boundary point of $D$ lies in $R^{\circ}$ either. Suppose such a point had a neighbourhood contained in $R$. Each unit $d$-cube containing $D$ meets that neighbourhood in a set with non-empty interior. So the interior of that $d$-cube, an open unit face, meets $R$ and is therefore one of the open unit faces of $R$. As $R$ is closed, the $d$-cube lies in $R$. Then every unit $d$-cube containing $D$ lies in $R$, and the criterion just proved gives $\operatorname{relint}D\subset R^{\circ}$, a contradiction. Hence in this case $D\cap R^{\circ}=\emptyset$, that is $D\subset W$.

This proves the dichotomy. For the equivalence, suppose first that $D\cap R^{\circ}\neq\emptyset$. Then $D\not\subset W$, so $D\subset R$ by what was shown. By the last paragraph, $\operatorname{relint}D\subset R^{\circ}$. By the criterion at points of $\operatorname{relint}D$, every unit $d$-cube containing $D$ lies in $R$.

Conversely, suppose every unit $d$-cube containing $D$ lies in $R$. There is at least one such cube, so $D\subset R$. The same criterion then gives $\operatorname{relint}D\subset R^{\circ}$, and in particular $D\cap R^{\circ}\neq\emptyset$.

\emph{(b)} By (a), every open unit face contained in $R$ is either contained in $R^{\circ}$ or disjoint from it. Hence $R^{\circ}$ and $\partial R$ are unions of open unit faces.

Let $\operatorname{relint}D\subset\partial R$. Then $D\subset R$, because $R$ is closed. Also $\operatorname{relint}D\cap R^{\circ}=\emptyset$, so, by the second case in the proof of (a), $D\cap R^{\circ}=\emptyset$. Thus $D\subset\partial R$. Hence $\partial R=R\cap W$ is a union of closed unit cubes.

By Lemma~\ref{8.10:lem-cube-fullness} and the identification $R=\bar{S}(R_{\mathrm{int}})$ made above,
\[
\mathsf{R}=\bar{K}(R_{\mathrm{int}})=K(\bar{S}(R_{\mathrm{int}}))=K(R).
\]
Since $\partial R=R\cap W$, a unit cube lies in $\partial R$ exactly when it lies in $R$ and in $W$. Hence $K(\partial R)=\mathsf{R}\cap K(W)$.

By definition, a cube of $\mathsf{R}$ is deep if and only if it meets $R^{\circ}$. A cube $D\in\mathsf{R}$ that is not deep satisfies $D\subset R$ and $D\cap R^{\circ}=\emptyset$, that is $D\subset\partial R$. Conversely, a cube contained in $\partial R$ does not meet $R^{\circ}$. This proves
\[
\mathsf{R}^{\partial}=\{D\in\mathsf{R}:D\subset\partial R\}=K(\partial R).
\]
Faces of non-deep cubes of $\mathsf{R}$ lie in $\partial R$ and are therefore non-deep. So $\mathsf{R}^{\partial}$ is a subcomplex, and $\mathsf{R}^{\mathrm{deep}}$ is its complement in $\mathsf{R}$.

\emph{(c)} The proofs of (a) and (b) go through word for word with the closed cells of side $\ell$ (of any dimension) in place of the closed unit cubes, and with the complex of closed cells of $R_{\mathrm{int}}$ in place of $\mathsf{R}$. The facts about the unit structure on which those proofs rest hold for the cell structure, since $R_{\mathrm{int}}$ is a set of cells:
\begin{itemize}
\item the relative interiors of the closed cells partition space;
\item $R$ is a union of closed top-dimensional cells;
\item near a point of $\operatorname{relint}D$, the top-dimensional cells containing $D$ are exactly those through that point, and they cover space.
\end{itemize}

\emph{(d)} Apply (a) with $D=\{y\}$. It gives $y\in R^{\circ}$ if and only if all $2^d$ unit $d$-cubes having $y$ as a vertex lie in $R$, that is, lie in $\mathsf{R}=K(R)$.

By Lemma~\ref{8.10:lem-cube-fullness}, a cube lies in $\mathsf{R}=\bar{K}(R_{\mathrm{int}})$ if and only if all its vertices lie in $\bar{S}(R_{\mathrm{int}})$. The union of the vertex sets of the $2^d$ unit $d$-cubes at $y$ is $y+\{-1,0,1\}^d$. This proves the criterion~\eqref{8.10:eq-hull-boundary-1}.

Now let $y\in\mathsf{R}_0$ be not deep. By the criterion, some unit $d$-cube $D\ni y$ is not in $\mathsf{R}$. The cube $D$ lies in a unique cell $c'$, and $c'\notin R_{\mathrm{int}}$; moreover
\[
y\in D\subset\bar{S}(c').
\]
Since $y\in\mathsf{R}_0$, there is a cell $c\in R_{\mathrm{int}}$ with $y\in\bar{S}(c)$. So $c'$ touches $c$, and
\[
c'\in\operatorname{hull}_1(R_{\mathrm{int}})\setminus R_{\mathrm{int}}\subset R_{\mathrm{col}};
\]
here the inclusion is that of Definition~\ref{8.10:def-hull-split}, valid for every $p_2\ge1$.
Hence $y\in\mathsf{A}_0$.

Conversely, let $y\in\mathsf{R}_0\cap\mathsf{A}_0$, so that $y\in\bar{S}(c')$ for some $c'\in R_{\mathrm{col}}$. A unit $d$-cube lies in exactly one cell. Hence the unit $d$-cubes of $\bar{Q}_{c'}$ having $y$ as a vertex are not in $\mathsf{R}$. By the criterion, $y$ is not deep.

Both inclusions together give that the non-deep vertices of $\mathsf{R}_0$ are exactly those of $\mathsf{I}_0=\mathsf{R}_0\cap\mathsf{A}_0$, for every $p_2\ge1$.
\end{proof}

\begin{definition}[Linked, unlinked and box-like components]\label{8.10:def-linked-components}
Let $C$ be as in Definition~\ref{8.10:def-hull-split}, with widths $p_1,p_2$. Assume the
lifting condition of Definition~\ref{8.10:def-torus-lifting} for
$\operatorname{hull}_{p_1+p_2+1}(C)$. Then the complexes $\mathsf{N}$, $\mathsf{R}$,
$\mathsf{A}$ of Definition~\ref{8.10:def-hull-split} and the exterior complex $\mathsf{Ex}$
associated with the interior--exterior split of that definition are complexes in
$\mathbb{Z}^d$; an edge is called \emph{interior} or \emph{exterior} in the sense of the same
split. We write $\mathsf{A}_0$ for the set of vertices of
$\mathsf{A}$. The \emph{graph} of $\mathsf{Ex}$ is the graph formed by its vertices and
edges. Loops, combinatorial homotopy and the edge-loop fundamental group $\pi_1^{\square}$
are those of Definition~\ref{8.10:def-edge-loop-homotopy}.

\begin{enumerate}
\item[(a)] \emph{Unlinked.} $C$ is \emph{unlinked} (at widths $p_1,p_2$) if at least one of
the following two conditions holds.
\begin{itemize}
\item[$(\mathrm{U}_1)$] Every loop in the graph of $\mathsf{Ex}$ is null-homotopic in
$\mathsf{Ex}$. Equivalently, each connected component of $\mathsf{Ex}$, isolated vertices
included, has trivial $\pi_1^{\square}$.
\item[$(\mathrm{U}_2)$] The graph of $\mathsf{Ex}$ is connected and, for a vertex
$v\in\mathsf{A}_0$, both of the following hold:
\begin{itemize}
\item every loop at $v$ in $\mathsf{N}$ is combinatorially homotopic in $\mathsf{N}$ to a
loop in $\mathsf{Ex}$;
\item every loop in $\mathsf{Ex}$ that is null-homotopic in $\mathsf{N}$ is null-homotopic
in $\mathsf{Ex}$.
\end{itemize}
Equivalently, the homomorphism sending the class of a loop at $v$ in $\mathsf{Ex}$ to its
class in $\mathsf{N}$,
\begin{equation*}
\pi_1^{\square}(\mathsf{Ex},v)\longrightarrow\pi_1^{\square}(\mathsf{N},v),
\end{equation*}
is an isomorphism.
\end{itemize}

\item[(b)] \emph{Linked.} $C$ is \emph{linked} if there exist the following three objects:
\begin{itemize}
\item an edge loop $\gamma$ all of whose edges are exterior;
\item a null-homotopy of $\gamma$ in the cubical complex $K(\mathbb{T})$ of the torus
$\mathbb{T}$ of Definition~\ref{8.10:def-torus-lifting};
\item an exterior configuration $V^{\mathrm{ext}}\in\mathcal{V}^{\mathrm{ext}}$ that is flat
on every plaquette of $\mathbb{T}$ without interior edge;
\end{itemize}
such that the parallel transport $T_\gamma(V^{\mathrm{ext}})$ along $\gamma$ satisfies
$T_\gamma(V^{\mathrm{ext}})\neq 1$, that is, it differs from the identity element of
$\mathrm{SU}(N)$. Under the lifting condition, $\gamma$ is null-homotopic in $K(\mathbb{T})$
if and only if it is null-homotopic in $\bar K(B)$ for some box of cells $B$ containing the
vertices of $\gamma$.

\item[(c)] \emph{Box-like.} $C$ is \emph{box-like} if
\begin{equation*}
R_{\mathrm{int}}=\operatorname{hull}_{p_1}(C)
\end{equation*}
is a box of cells $\prod_\mu\{a_\mu,\dots,b_\mu\}$ in lifted indices, in the sense of
Definition~\ref{8.10:def-torus-lifting}. This holds, for example, when $C$ is a single cell
or when $C$ is itself a box.
\end{enumerate}
\end{definition}

An unlinked $C$ is not linked; this is item~(i) of the lemma on unlinked components proved
in this section. The converse is not
asserted in this book. The exact criterion for an exterior configuration flat in the sense
of~(b) to extend to a configuration flat on every plaquette of $\mathbb{T}$ is an extension
condition on the holonomy representations of the components of $\mathsf{Ex}$; it is not
stated precisely here, and neither it nor the converse is used in what follows.

\begin{definition}[The obstruction group of the filling problem]\label{8.10:def-obstruction-group}
We use the torus $\mathbb{T}$ of this section, its cubical complex $K(\mathbb{T})$, and the sets
$\mathsf{R}_0$, $\mathsf{R}_1$, $\mathsf{F}$ and the closed complex
$\overline{K}(\operatorname{hull}_1 R_{\mathrm{int}})$ attached to the region $C$, as in
Lemma~\ref{8.10:lem-collar-plaquettes}.

\emph{The cochain groups.} For $q\ge 1$ let $C^q_\sharp$ be the group of $\mathbb{Z}$-valued functions
on the $q$-cubes of $\mathbb{T}$ that have at least one edge in $\mathsf{R}_1$, and let $C^0_\sharp$
be the group of $\mathbb{Z}$-valued functions on $\mathsf{R}_0$. By
Lemma~\ref{8.10:lem-collar-plaquettes}(a) these cubes lie in
$\overline{K}(\operatorname{hull}_1 R_{\mathrm{int}})$; $C^1_\sharp$ is the group of functions on
$\mathsf{R}_1$, and $C^2_\sharp$ is the group of functions on $\mathsf{F}$.

\emph{The coboundary.} Let $e$ be extension by zero to all cubes of $\mathbb{T}$, $r$ restriction to
the cubes indexing $C^{q+1}_\sharp$ when $d$ acts on $C^q_\sharp$, and $d$ the cubical coboundary
of $K(\mathbb{T})$: for a
function $\chi$ on vertices, a function $a$ on bonds and a plaquette $p$ at $x$ with directions
$i<j$,
\begin{align*}
(d\chi)(x,x+e_i) &= \chi(x+e_i)-\chi(x),\\
(da)(p) &= a(x,x+e_i)+a(x+e_i,x+e_i+e_j)\\
&\qquad -a(x+e_j,x+e_i+e_j)-a(x,x+e_j);
\end{align*}
for a function $F$ on plaquettes, $(dF)(c)$ is the alternating sum of the values of $F$ on the six
$2$-faces of the $3$-cube $c$; and in every higher degree $d$ is the cubical coboundary of
$K(\mathbb{T})$ in the same way. We put $d_\sharp := r\circ d\circ e$, and we write $C_\sharp$ for
the sequence $(C^q_\sharp,d_\sharp)_{q\ge0}$; in degrees $q\ge1$ it is a cochain complex,
and only its segment $C^1_\sharp\to C^2_\sharp\to C^3_\sharp$ enters the
definitions that follow.

\emph{The obstruction group and its rank.} We define
\begin{equation}\label{8.10:eq-obstruction-group-1}
\mathfrak{O}^\sharp(C) := H^2(C_\sharp;\mathbb{Z})
=\ker\bigl(d_\sharp : C^2_\sharp\to C^3_\sharp\bigr)\big/ d_\sharp\bigl(C^1_\sharp\bigr),
\end{equation}
\begin{equation}\label{8.10:eq-obstruction-group-2}
\operatorname{lk}^\sharp(C) := \operatorname{rank}_{\mathbb{Z}}\mathfrak{O}^\sharp(C)
=\dim_{\mathbb{R}} H^2(C_\sharp\otimes\mathbb{R})\le|\mathsf{F}|.
\end{equation}
For a real vector space $\mathfrak{t}$, the complex $C_\sharp\otimes\mathfrak{t}$ is the linear
caricature of the filling problem: its unknowns live on $\mathsf{R}_1$ and its curvature on
$\mathsf{F}$.

\emph{The comparison group.} The following group is determined by the closed unit cubes of the
closed region $R$ of Lemma~\ref{8.10:lem-hull-boundary} alone. With $\mathsf{R}$ its complex of
closed unit cubes and $\mathsf{R}^\partial\subset\mathsf{R}$ its subcomplex of boundary
cubes, as in Lemma~\ref{8.10:lem-hull-boundary}(b), we put
\begin{equation*}
\mathfrak{O}(C) := H^2(\mathsf{R},\mathsf{R}^\partial;\mathbb{Z}),\qquad
\operatorname{lk}(C) := \operatorname{rank}_{\mathbb{Z}}\mathfrak{O}(C),
\end{equation*}
the relative cubical cochains being the functions on the deep cubes $\mathsf{R}^{\mathrm{deep}}$
(Lemma~\ref{8.10:lem-hull-boundary}(b)). This group belongs to a coarser version of the filling
problem, in which the bonds of $\partial R$ are prescribed. The word \emph{ball-like} is used only
informally, for $\operatorname{lk}(C)=\operatorname{lk}^\sharp(C)=0$; the sufficient condition
operative in the lemma on the filling problem is that $C$ be unlinked, in the sense of that lemma.
\end{definition}

\begin{proof}[Proof of the assertions in Definition~\ref{8.10:def-obstruction-group}]
\emph{The complex in degrees $q\ge1$.} Let $q\ge1$ and $f\in C^q_\sharp$. The function $ef$ is
supported on $q$-cubes having at least one edge in $\mathsf{R}_1$. By the formulas for $d$, the
value $(d(ef))(c)$ at a $(q+1)$-cube $c$ is a signed sum of the values of $ef$ on the $q$-faces
of $c$; it vanishes unless $c$ has a face in the support of $ef$. Such a face has an edge in
$\mathsf{R}_1$, and that edge is an edge of $c$. Hence $d(ef)$ is supported on $(q+1)$-cubes with
at least one edge in $\mathsf{R}_1$, which are the cubes indexing $C^{q+1}_\sharp$. Restricting to
these cubes and extending by zero therefore changes nothing:
\begin{equation}\label{8.10:eq-obstruction-group-3}
e\,r\,d\,e = d\,e \qquad\text{on } C^q_\sharp,\ q\ge1 .
\end{equation}
Consequently, for $q\ge1$ and $f\in C^q_\sharp$, using \eqref{8.10:eq-obstruction-group-3} in the
middle step and $d\circ d=0$ for the cubical coboundary of $K(\mathbb{T})$ in the last,
\begin{equation*}
d_\sharp d_\sharp f = r\,d\,(e\,r\,d\,e)f = r\,d\,d\,e f = 0 .
\end{equation*}
So $d_\sharp\circ d_\sharp=0$ in degrees $q\ge1$. The group $H^2(C_\sharp;\mathbb{Z})$
in \eqref{8.10:eq-obstruction-group-1} involves only the maps
$C^1_\sharp\to C^2_\sharp\to C^3_\sharp$, for which this has been shown; hence
$d_\sharp(C^1_\sharp)\subset\ker(d_\sharp: C^2_\sharp\to C^3_\sharp)$ and the quotient is
defined.

\emph{The rank.} The torus $\mathbb{T}$ is finite, so each $C^q_\sharp$ is a free abelian group of
finite rank, with basis the indicator functions of the cubes indexing it. For the cochain complex
$(C^q_\sharp,d_\sharp)_{q\ge1}$ of finitely generated free abelian groups, whose second cohomology
group is $H^2(C_\sharp;\mathbb{Z})$, the universal coefficient theorem gives
$H^2(C_\sharp\otimes\mathbb{R})\cong H^2(C_\sharp;\mathbb{Z})\otimes\mathbb{R}$, the torsion term
vanishing because $\mathbb{R}$ is torsion-free. The rank of a finitely generated abelian group
equals the real dimension of its tensor product with $\mathbb{R}$; this is the equality in
\eqref{8.10:eq-obstruction-group-2}.

\emph{The bound.} $H^2(C_\sharp\otimes\mathbb{R})$ is the quotient of the subspace
$\ker d_\sharp$ of $C^2_\sharp\otimes\mathbb{R}$ by a subspace. Its dimension is therefore at most
$\dim_{\mathbb{R}}(C^2_\sharp\otimes\mathbb{R})$, the number of plaquettes indexing $C^2_\sharp$.
By Lemma~\ref{8.10:lem-collar-plaquettes}(a) every plaquette of $\mathbb{T}$ with an edge in
$\mathsf{R}_1$ lies in $\mathsf{F}$, so this number is at most $|\mathsf{F}|$. This proves the
inequality in \eqref{8.10:eq-obstruction-group-2}.
\end{proof}

\begin{definition}[The spanned box and its exterior subcomplexes]\label{8.10:def-spanned-box}
Let $C$ be the set of cells of the obstruction group of the filling problem
(Definition~\ref{8.10:def-obstruction-group}) and $p_1,p_2$ the two integers fixed with it, and write
$\operatorname{hull}_r(\cdot)$ for
the hull of width $r$ of a set of cells. Let $R$, $W$, $R^{\circ}$, the complex $\mathsf{R}$ and its deep cubes
$\mathsf{R}^{\mathrm{deep}}$ be as in the structure of the hull's boundary
(Lemma~\ref{8.10:lem-hull-boundary}). Let $R_{\mathrm{int}}$ be the set of interior cells attached to the
region $R$, with set of
sites $\bar{S}(R_{\mathrm{int}})$. We write $\mathsf{R}_0$ for the set of vertices and $\mathsf{R}_1$ for the
set of interior edges attached to the region $R$ together with its hull; an edge in $\mathsf{R}_1$ is called an
\emph{interior edge}. A \emph{deep vertex} is a
vertex belonging to $\mathsf{R}^{\mathrm{deep}}$, that is, a site lying in $R^{\circ}$. Finally,
$\mathsf{Ex}$ denotes the exterior complex attached to the region $R$.

Assume that $\operatorname{hull}_{p_1+p_2+1}(C)$ satisfies the patch condition of lifting shapes from the torus
to the lattice (Definition~\ref{8.10:def-torus-lifting}), and let
\begin{equation*}
B:=B\bigl(\operatorname{hull}_{p_1+p_2+1}(C)\bigr)=\prod_{\mu}\{a_\mu,\dots,b_\mu\}
\end{equation*}
be the box of cells spanned by $\operatorname{hull}_{p_1+p_2+1}(C)$ there. Put
$\mathbb{B}:=\bar{K}(B)=K(\bar{S}(B))$. This is a full complex on a lattice box, and it lifts injectively to
$\mathbb{Z}^d$ by Definition~\ref{8.10:def-torus-lifting}. We set
\begin{align}
E_B&:=\{\kappa\in\mathbb{B}:\ \kappa\cap R^{\circ}=\emptyset\},\notag\\
\mathbb{E}_B&:=\{\kappa\in\mathbb{B}:\ \text{no edge of }\kappa\text{ lies in }\mathsf{R}_1,
\text{ and }\kappa\text{ is not a deep vertex}\},\label{8.10:eq-spanned-box-1}
\end{align}
and we call the \emph{boundary layer} of $\bar{S}(B)$ the set
\begin{equation}\label{8.10:eq-spanned-box-2}
\partial\bar{S}(B):=\bigl\{z\in\bar{S}(B):\ z_\mu\in\{\ell a_\mu,\ \ell b_\mu+\ell\}\ \text{for some }\mu\bigr\}.
\end{equation}
The exterior complex $\mathsf{Ex}$ satisfies $\mathsf{Ex}\subset\mathbb{E}_B$ by its construction.
Moreover the following hold.
\begin{enumerate}
\item[(i)] $E_B$ and $\mathbb{E}_B$ are subcomplexes of $\mathbb{B}$.
\item[(ii)] $\mathbb{E}_B\subset E_B$.
\item[(iii)] $K(\partial\bar{S}(B))\subset\mathbb{E}_B$.
\end{enumerate}
\end{definition}

\begin{proof}
(i) First consider $E_B$, and let $\kappa\in E_B$. By parts (a) and (b) of
Lemma~\ref{8.10:lem-hull-boundary}, $\kappa$ lies in $W$, and so do its faces. A face of $\kappa$ is a subset
of $\kappa$, so it does not meet $R^{\circ}$. It belongs to $\mathbb{B}$ because $\mathbb{B}$ is a complex.
Hence $E_B$ is closed under taking faces.

Now consider $\mathbb{E}_B$. Let $\kappa\in\mathbb{E}_B$ and let $\tau$ be a face of $\kappa$. Every edge of
$\tau$ is an edge of $\kappa$, so no edge of $\tau$ lies in $\mathsf{R}_1$. It remains to show that $\tau$ is
not a deep vertex. If $\tau=\kappa$, this holds by \eqref{8.10:eq-spanned-box-1}. Otherwise $\kappa$ has
dimension at least one. Suppose $\tau$ were a deep vertex $y$. By part (d) of
Lemma~\ref{8.10:lem-hull-boundary}, a cube of dimension at least one with a deep vertex $y$ has an interior
edge at $y$. Then $\kappa$ would have an edge in $\mathsf{R}_1$, contrary to $\kappa\in\mathbb{E}_B$. Hence
$\tau\in\mathbb{E}_B$, and $\mathbb{E}_B$ is a subcomplex.

(ii) Let $\kappa\in\mathbb{E}_B$. Suppose first that $\kappa$ is a vertex. It is not a deep vertex, so it is
not a site of $R^{\circ}$, and therefore $\kappa\cap R^{\circ}=\emptyset$.

Suppose next that $\kappa$ has dimension at least one, and assume that $\kappa$ meets $R^{\circ}$. By part (a)
of Lemma~\ref{8.10:lem-hull-boundary}, $\kappa$ lies in $R$. Then, by the definitions of $\mathsf{R}_0$ and
$\mathsf{R}_1$, all its vertices lie in
$\bar{S}(R_{\mathrm{int}})$ and it has interior edges. This contradicts $\kappa\in\mathbb{E}_B$. Hence in both
cases $\kappa\cap R^{\circ}=\emptyset$, that is, $\kappa\in E_B$.

(iii) We have $\operatorname{hull}_1(R_{\mathrm{int}})\subset\operatorname{hull}_{p_1+p_2+1}(C)$.
The box $B$ is spanned by $\operatorname{hull}_{p_1+p_2+1}(C)$, so
\begin{equation*}
\operatorname{hull}_1(R_{\mathrm{int}})\subset\operatorname{hull}_{p_1+p_2+1}(C)\subset B .
\end{equation*}
Consequently every site of $\bar{S}(R_{\mathrm{int}})$ is at sup-distance at least $\ell\ge 2$ from the
boundary layer \eqref{8.10:eq-spanned-box-2}. Therefore the cubes of $K(\partial\bar{S}(B))$ have no vertex
in $\mathsf{R}_0$, hence no edge in $\mathsf{R}_1$ and no deep vertex.

Now let $\kappa\in K(\partial\bar{S}(B))$. Since $\partial\bar{S}(B)\subset\bar{S}(B)$, we have
$\kappa\in\mathbb{B}$. By the preceding paragraph, no edge of $\kappa$ lies in $\mathsf{R}_1$ and $\kappa$ is
not a deep vertex. Thus $\kappa\in\mathbb{E}_B$.
\end{proof}

\begin{definition}[The core bound and the allotment of the linear reserve; the core bound: proof
outstanding]\label{8.10:def-core-bound}
We use the terminal level $K'$ and the couplings of Notation~\ref{8.10:not-terminal-couplings}, and
the live components and their histories of Definition~\ref{8.10:def-histories-fibres}. For a live
component $C$ at level $K'$:
\begin{itemize}
\item $n(C)$ is the size of $C$;
\item $\mathfrak{H}(C)$ is the family of histories of $C$;
\item for $\mathfrak{h}\in\mathfrak{H}(C)$, $\mathcal{Z}^{1}(\mathfrak{h})$ is the set of seeds of the
first kind of $\mathfrak{h}$.
\end{itemize}
A seed $\zeta$ of the first kind is one connected creation region, of arbitrary size. It is created
at step $j(\zeta)$, and $\mathfrak{d}(\zeta)$ denotes its extent.

\emph{The core bound.} For every live component $C$ at level $K'$ and every
$\mathfrak{h}\in\mathfrak{H}(C)$,
\begin{equation}\label{8.10:eq-core-bound-1}
n(C)\le C_G'\sum_{\zeta\in\mathcal{Z}^{1}(\mathfrak{h})}\bigl(\mathfrak{d}(\zeta)+1\bigr),
\qquad C_G':=2^{d+1}\cdot 219^{d};
\end{equation}
in particular $\mathcal{Z}^{1}(\mathfrak{h})\neq\emptyset$. The bound \eqref{8.10:eq-core-bound-1} is
taken as a hypothesis in what follows.

A bound of $n(C)$ by a fixed constant times the number of all seeds of $\mathfrak{h}$ is not used,
since a seed of the first kind may be arbitrarily large, its size being carried by
$\mathfrak{d}(\zeta)+1$ and by the linear-reserve bound rather than by a count of seeds.

\emph{Budgets.} For $\mathfrak{h}\in\mathfrak{H}(C)$ we put
\begin{equation}\label{8.10:eq-core-bound-2}
B_\tau(C,\mathfrak{h}):=\theta'_\tau\gamma_0A_1^{2}
\sum_{\zeta\in\mathcal{Z}^{1}(\mathfrak{h})}p_0\bigl(g_{j(\zeta)}\bigr)^{2}\bigl(\mathfrak{d}(\zeta)+1\bigr).
\end{equation}
Here $\theta'_\tau$ is the part of $\theta'$ fixed below, $\gamma_0$ is the constant denoted by this
letter in the linear-reserve bound, $g_{j(\zeta)}$ is the running coupling at step $j(\zeta)$, and
$p_0(\cdot)$ is the degree function. With this choice the reserve factor of a set
$\mathcal{S}\subset\mathfrak{H}(C)$ of histories is
\begin{equation*}
r_{\theta'_\tau}(\mathcal{S})=\sup_{\mathfrak{h}\in\mathcal{S}}e^{-B_\tau(C,\mathfrak{h})}.
\end{equation*}
Put $p_0':=p_0(x)-1$, with $x$ as in Notation~\ref{8.10:not-terminal-couplings}. Let $c_A$ be the
constant introduced under this letter in the lower bound stated later in this chapter, when that
constant is positive, and $c_A:=1$ otherwise; in either case $c_A>0$. We define
\begin{equation}\label{8.10:eq-core-bound-3}
B^\flat_\tau(C):=\frac{\theta'_\tau\gamma_0A_1^{2}\,p_0'^{\,2}\,n(C)}{C_G'},
\qquad
\mathfrak{E}_*(C):=\frac{B^\flat_\tau(C)}{4c_Ax}.
\end{equation}

\emph{The allotment of the linear reserve.} The linear-reserve bound holds for every fraction
$\theta'\in\,]0,1[$ of the linear reserve under two one-sided ceilings on the coupling, namely the
conditions \cite[(1.82), (1.88)]{B-CMP122b}, written in the letters of this book and read at
$\theta'$. In the
order of choice of Notation~\ref{2.3:not-stages-first}, $\theta'$ is one number fixed at Stage~0;
each statement that uses $\theta'$ uses only its own part, and these parts are disjoint. Two
disjoint parts $\theta'_A$ and $\theta'_{\mathrm{age}}$ of $\theta'$ are used elsewhere in this
chapter. The budgets \eqref{8.10:eq-core-bound-2} and \eqref{8.10:eq-core-bound-3} require a part
$\theta'_\tau>0$ disjoint from these two. Accordingly $\theta'\in\,]0,1[$ is fixed as
\begin{equation}\label{8.10:eq-core-bound-4}
\theta'=\theta'_A+\theta'_{\mathrm{age}}+\theta'_\tau ,
\end{equation}
with these $\theta'_A$, $\theta'_{\mathrm{age}}$ and $\theta'_\tau>0$. The two ceilings of the
linear-reserve bound are read at this
$\theta'$. They are conditions of the kind ``$p_0$ large, $\gamma$ small'' of Ba{\l}aban's
construction, and they are one-sided in the coupling.

No other quantity depends on $\theta'_\tau$. The following do not enter \eqref{8.10:eq-core-bound-2}
or \eqref{8.10:eq-core-bound-3}: the flat reserve, the factor $e^{-(2-\theta)\underline{p}_0}$, and
the factor $e^{-\kappa_1d_{K'}(C)}$ of the penalty $\mathrm{pen}^\tau_\theta$. That penalty is the one
appearing in the domination bound of Definition~\ref{8.10:def-fibre-domination}, whose proof is
outstanding, and $\theta$ and $\underline{p}_0$ are as there.
\end{definition}

\emph{Route.} A proof would have to establish \eqref{8.10:eq-core-bound-1} for every live component
$C$ at level $K'$ and every history $\mathfrak{h}$ of the family $\mathfrak{H}(C)$. This family
contains the histories of Ba{\l}aban's construction \cite{B-CMP122b} as a special case. The
proof would work from Ba{\l}aban's rules for the creation of large-field regions, under a one-sided
ceiling on the coupling and a fixed numerical one-sided lower bound on the blocking factor. The
bound would be obtained by comparing the size $n(C)$ with the extents
$\mathfrak{d}(\zeta)+1$ of the seeds of the first kind of $\mathfrak{h}$, with the constant $C_G'$ of
\eqref{8.10:eq-core-bound-1}, together with one further estimate on the histories of that family
used jointly with this comparison.

\begin{definition}[The one-sided coupling and hull-width conditions]\label{8.10:def-one-sided-conditions}
Let $x$ and $X$ be the inverse squared couplings of Notation~\ref{8.10:not-terminal-couplings}. Let
$C$ be a core,
and $\theta'_\tau$ the number entering the core bound, both as in Definition~\ref{8.10:def-core-bound};
the core bound stated there is a bound whose proof is outstanding. Let $\mu\in(0,1)$ be an auxiliary
pure number of this chapter, on which
the constant $c_L'$ and the threshold $x_L$ of item~(c) below depend. We
write $p_0(x)$ for the value of the degree function $p_0(\cdot)$ at the coupling $g'=x^{-1/2}$,
so that $p_0(x)=A_0(\log x)^{p_0}$. In what follows, $p_1$ denotes the hull thickness and $p_2$
the collar width. The following conditions are imposed, and each of them is one-sided.
\begin{enumerate}
\item[(a)] \emph{The hull condition.} The hull thickness satisfies the lower bound
\begin{equation*}
p_1\ \ge\ p_W(d,L)+1,
\end{equation*}
where $p_W(d,L)$ is a number depending on $(d,L)$ only, attached to the small-field window
structures of Definition~\ref{8.10:def-window-structure}: by the
localisation of these structures stated there, a statement whose proof is outstanding, a plaquette
on which they may differ has a vertex in a cell at
cell-distance at most $p_W$ from the core $C$. Write $\mathsf{R}_2$ for the region about the core $C$
into which this condition places these plaquettes, $\operatorname{hull}_{p_W+1}(C)$ for the
cell hull of radius $p_W+1$ about $C$, $\mathsf{I}$ for the interior region of the hull
construction about $C$, and $\bar{\mathsf{P}}(\mathsf{I})$ for its closed plaquette set. This
condition places inside $\mathsf{R}_2$ every plaquette on which the window structures may differ,
those of the datum $\Lambda$ of a fibre on the one hand and those of the histories present in
that fibre on the other, the fibre and its datum $\Lambda$ being as fixed earlier in this chapter.
Indeed, let $p$ be a plaquette with a vertex in a cell at
cell-distance at most $p_W$ from $C$. Then all vertices of $p$ lie in cells of
$\operatorname{hull}_{p_W+1}(C)\subset \mathsf{I}$, and hence $p$ lies in
$\bar{\mathsf{P}}(\mathsf{I})$. The condition enters the hypotheses of the statement that supplies
the lower bound on the Laplace kernel; it is not
invoked in any proof given in this section.
\item[(b)] \emph{The collar condition.} The collar width satisfies the lower bounds
\begin{equation*}
(p_2-1)\,\ell\ \ge\ 2r_W+3\qquad\text{and}\qquad p_2\ \ge\ 3,
\end{equation*}
where $\ell$ is the side of a cell and $r_W$ a length attached to the window structures of
Definition~\ref{8.10:def-window-structure} (their localisation, stated there, is a statement whose
proof is outstanding); the letter
$\ell$ here is not the length $\ell(n,R,g)$ of the glossary, Notation~\ref{0.2:not-glossary}.
Both bounds hold whenever $p_2\ge 2c_r+3$, where $c_r$ is a constant of this chapter.
\item[(c)] \emph{The Laplace-floor condition.} This condition reads
\begin{equation}\label{8.10:eq-one-sided-conditions-1}
\mu\,p_0(x)\ \ge\ c_L',
\end{equation}
that is, $x\ge x_L$, where
\begin{equation}\label{8.10:eq-one-sided-conditions-2}
x_L:=\inf\bigl\{x_0\ge e:\ \mu A_0(\log x')^{p_0}\ge c_L'\ \text{for all } x'\ge x_0\bigr\}.
\end{equation}
Here $c_L'\ge 1$ is the constant of the lower bound on the Laplace kernel. It depends on
$(d,N,L,M,p_1,p_2,\mu)$ and on the constants of Ba{\l}aban's construction. Condition
\eqref{8.10:eq-one-sided-conditions-1} is the floor on the window margin of the Laplace kernel.
It is implied by $X\ge x_L$ through clause~(0) of Theorem~\ref{3.1:thm-main}.
\item[(d)] \emph{The degree condition.} This condition bears on the exponents of Ba{\l}aban's
construction and reads
\begin{equation*}
d\,r+1\ <\ 2p_0 ,
\end{equation*}
where $r$ is Ba{\l}aban's exponent, the one in the condition $p_0\ge 5r$ (not the extraction depth
$r(s)$ of the two-point witness theorem). The inequality is
strict, and no equality case occurs. For $d=4$ the condition is implied by $p_0\ge 5r$ whenever
$r>1/6$, since then
\begin{equation*}
2p_0\ \ge\ 10r\ =\ 4r+6r\ >\ 4r+1 .
\end{equation*}
In particular, the condition holds in every case of Ba{\l}aban's exponents used in this book.
\item[(e)] \emph{The coupling-smallness condition.} This condition reads $x\ge x_\tau$, where
$x_\tau$ is the threshold fixed later in this chapter. It is a smallness condition on the coupling
given $(d,N,L,M,p_1,p_2,\mu,\theta'_\tau)$. It is implied by $X\ge x_\tau$ through clause~(0) of
Theorem~\ref{3.1:thm-main}.
\item[(f)] \emph{The second threshold condition.} This condition reads $x\ge x_{\mathrm{hol}}$, where
$x_{\mathrm{hol}}$ is the threshold fixed later in this chapter. It is implied by
$X\ge x_{\mathrm{hol}}$.
\item[(g)] \emph{The shape condition.} This is a condition on the shape of the regions
considered. It is a hypothesis on a shape and not a condition on the parameters.
\item[(h)] \emph{Inherited conditions.} Two conditions are inherited from the statement that
supplies the core bound stated in Definition~\ref{8.10:def-core-bound}, a statement whose proof is
outstanding: the condition $L\ge 5$ and a further hypothesis of that statement. In addition, a
statement of this book that this chapter uses by statement, as it uses clause~(0) of
Theorem~\ref{3.1:thm-main}, brings with it the hypotheses under which it is proved.
\end{enumerate}
No upper bound on $L$, $M$, $K$, $m$, the torus-size integer $N_T$ of this chapter, the number
$n(C)$ entering the core
bound stated in Definition~\ref{8.10:def-core-bound} (a bound whose proof is outstanding), the number
of seeds in
the constructions of this chapter, the ages or the depths
occurs anywhere in the conditions imposed in this chapter. No lower bound on the coupling occurs
there either.
\end{definition}

\begin{definition}[Standard facts of algebraic topology]\label{8.10:not-topology-facts}
The following seven facts are theorems of algebraic topology. They are cited by statement, and
their proofs are not reproduced.
\begin{itemize}
\item[(1)] For a finite CW pair, the cellular cochain complex computes singular cohomology.
Two finite CW structures on the same pair give isomorphic cohomology.
\item[(2)] For a pair of spaces, singular cohomology has the long exact cohomology sequence.
\item[(3)] A contractible space $S$ has $H^q(S)=0$ for every $q\ge 1$.
\item[(4)] Let $S$ be a finite CW complex. There is an isomorphism
\begin{equation*}
H^1(S;\mathbb{Z})\cong\operatorname{Hom}\bigl(H_1(S;\mathbb{Z}),\mathbb{Z}\bigr),
\end{equation*}
and this group is free abelian. For cochain complexes of finitely generated free abelian groups,
the ranks of the cohomology groups over $\mathbb{Z}$ equal the dimensions over $\mathbb{R}$ of
the corresponding real cohomology groups. A non-zero class of $H^1(S;\mathbb{R})$ pairs
non-trivially with some integral $1$-cycle. Hence, when the $1$-skeleton of $S$ is a graph, a
non-zero class of $H^1(S;\mathbb{R})$ pairs non-trivially with some edge loop.
\item[(5)] Let $X\subset\mathbb{Z}^d$ be a finite cubical complex whose geometric realisation is
simply connected. Then every edge loop in $X$ admits a null-homotopy by the elementary moves of
Definition~\ref{8.10:def-edge-loop-homotopy}; this follows from the isomorphism of the
edge-loop group with the fundamental group stated there.
\item[(6)] For $d\ge 3$ the sphere $S^{d-1}$ satisfies $H^1(S^{d-1};\mathbb{R})=0$. Moreover, for
$d\ge 3$ the geometric realisation of the boundary of a lattice box in $\mathbb{Z}^d$ is
homeomorphic to the sphere $S^{d-1}$.
\item[(7)] Two further facts are used.
\begin{itemize}
\item[(a)] The nerve theorem holds for finite covers by subcomplexes all of whose non-empty
intersections are contractible.
\item[(b)] Alexander duality holds.
\end{itemize}
Both facts in \textup{(7)} are used only in a remark of this chapter.
\end{itemize}
Facts \textup{(1)}--\textup{(7)} are cited by statement from two sources:
A.~Hatcher, \emph{Algebraic Topology} (Propositions~1.14 and~1.26 with Corollary~4.12,
Corollary~0.20, Theorem~3.44, the last being Alexander duality), and A.~Bj\"orner,
\emph{Topological methods}, Handbook of Combinatorics (Theorem~10.6, the nerve theorem).
\end{definition}

\begin{lemma}[The linear filling problem]\label{8.10:lem-linear-filling}
Let $\mathfrak{t}$ be a finite-dimensional real vector space, and let $\mathsf{R}_1$, $\mathsf{F}$, the
complex $C_\sharp$, its differential $d_\sharp$ and the number $\mathrm{lk}^\sharp(C)$ be as in
Definition~\ref{8.10:def-obstruction-group}. A bond of $\mathsf{R}_1$ is called \emph{interior}, and the
other bonds of $\mathbb{T}$ are called \emph{exterior}. Consider $\mathfrak{t}$-valued $1$-cochains
\begin{equation*}
A=A^{\mathrm{ext}}\oplus A^{\mathrm{in}}
\end{equation*}
on the bonds of $\mathbb{T}$, with $A^{\mathrm{ext}}$ on the exterior bonds and $A^{\mathrm{in}}$ on
$\mathsf{R}_1$. For a map $F\colon\mathsf{F}\to\mathfrak{t}$ let $F\cup dA^{\mathrm{ext}}$ denote the
$2$-cochain that equals $F$ on $\mathsf{F}$ and equals $d(A^{\mathrm{ext}}\oplus 0)$ on the plaquettes
without interior edge.
\begin{enumerate}
\item[(i)] For fixed $A^{\mathrm{ext}}$, the set
\begin{equation}\label{8.10:eq-linear-filling-1}
\bigl\{(dA)|_{\mathsf{F}} : A^{\mathrm{in}}\ \text{arbitrary}\bigr\}
\end{equation}
equals the set of all $F\colon\mathsf{F}\to\mathfrak{t}$ such that $F\cup dA^{\mathrm{ext}}$ is closed
on every $3$-cube having an interior edge, if and only if $\mathrm{lk}^\sharp(C)=0$.
\item[(ii)] Consequently, if $\mathrm{lk}^\sharp(C)=0$, every exterior datum $A^{\mathrm{ext}}$ together
with any target curvature $F$ on $\mathsf{F}$ compatible in the sense of \textup{(i)} is realised by
some $A^{\mathrm{in}}$: in this linear setting, fillability with prescribed curvature is equivalent to
the cocycle condition.
\end{enumerate}
\end{lemma}

\begin{proof}
Fix $A^{\mathrm{ext}}$ and put $\Phi_0:=d(A^{\mathrm{ext}}\oplus 0)$. For a $2$-cochain $\Psi$ on
$\mathsf{F}$ let $e\Psi$ be its extension by zero to all plaquettes of $\mathbb{T}$. The value of $dA$ on
a plaquette depends only on the values of $A$ on the bonds of its boundary; hence on a plaquette without
interior edge we have $dA=d(A^{\mathrm{ext}}\oplus 0)=\Phi_0$ for every $A^{\mathrm{in}}$.

Let $F\colon\mathsf{F}\to\mathfrak{t}$ satisfy the cocycle condition of (i), and put
\begin{equation}\label{8.10:eq-linear-filling-2}
\Phi:=F-\Phi_0|_{\mathsf{F}}\quad\text{on }\mathsf{F}.
\end{equation}
On every plaquette, $e\Phi$ and $F\cup dA^{\mathrm{ext}}-\Phi_0$ agree: on $\mathsf{F}$ both equal
$F-\Phi_0$ by \eqref{8.10:eq-linear-filling-2}, and on a plaquette without interior edge both vanish,
because there $F\cup dA^{\mathrm{ext}}$ and $\Phi_0$ both equal $d(A^{\mathrm{ext}}\oplus 0)$. Thus
$F\cup dA^{\mathrm{ext}}$ differs from $e\Phi$ by the coboundary $\Phi_0$, and since $d\Phi_0=dd(A^{\mathrm{ext}}\oplus0)=0$,
\begin{equation}\label{8.10:eq-linear-filling-3}
d(e\Phi)=d\bigl(F\cup dA^{\mathrm{ext}}\bigr).
\end{equation}
By hypothesis the right-hand side vanishes on every $3$-cube with an interior edge; so $d(e\Phi)$
vanishes there, that is, $\Phi\in\ker d_\sharp\subset C^2_\sharp\otimes\mathfrak{t}$. Read backwards,
the same computation shows that for every $\Phi\in\ker d_\sharp$ the map
$F:=\Phi+\Phi_0|_{\mathsf{F}}$ satisfies the cocycle condition, since then
$F\cup dA^{\mathrm{ext}}=e\Phi+\Phi_0$ and \eqref{8.10:eq-linear-filling-3} holds. Hence the set of
$F$ satisfying the cocycle condition is $\Phi_0|_{\mathsf{F}}+\ker d_\sharp$.

Next, by linearity of $d$,
\begin{equation*}
(dA)|_{\mathsf{F}}=\Phi_0|_{\mathsf{F}}+\bigl(d(0\oplus A^{\mathrm{in}})\bigr)|_{\mathsf{F}}
=\Phi_0|_{\mathsf{F}}+d_\sharp A^{\mathrm{in}},
\qquad A^{\mathrm{in}}\in C^1_\sharp\otimes\mathfrak{t}.
\end{equation*}
Therefore $F=(dA)|_{\mathsf{F}}$ for $A=A^{\mathrm{ext}}\oplus A^{\mathrm{in}}$ if and only if
$\Phi=d_\sharp(A^{\mathrm{in}})$ with $A^{\mathrm{in}}\in C^1_\sharp\otimes\mathfrak{t}$, and the set
\eqref{8.10:eq-linear-filling-1} is $\Phi_0|_{\mathsf{F}}+d_\sharp(C^1_\sharp\otimes\mathfrak{t})$.

Conversely, every $(dA)|_{\mathsf{F}}$ satisfies the cocycle condition: for $F=(dA)|_{\mathsf{F}}$ the
cochain $F\cup dA^{\mathrm{ext}}$ equals $dA$ on every plaquette (on the plaquettes without interior
edge because $dA=\Phi_0$ there), and $d(dA)=0$. This is the inclusion
$d_\sharp(C^1_\sharp\otimes\mathfrak{t})\subset\ker d_\sharp$.

Both sets in (i) are thus translates of subspaces by the same vector $\Phi_0|_{\mathsf{F}}$, so they
coincide if and only if
\begin{equation*}
\ker d_\sharp=d_\sharp\bigl(C^1_\sharp\otimes\mathfrak{t}\bigr),
\end{equation*}
that is, if and only if $H^2(C_\sharp\otimes\mathfrak{t})=0$. By the standard facts of algebraic
topology collected in Definition~\ref{8.10:not-topology-facts}, this holds if and only if
$\mathrm{lk}^\sharp(C)=0$. This proves (i).

For (ii), let $\mathrm{lk}^\sharp(C)=0$. By (i) the two sets coincide for the given
$A^{\mathrm{ext}}$, so every $F$ compatible with $A^{\mathrm{ext}}$ in the sense of (i) lies in the set
\eqref{8.10:eq-linear-filling-1}, i.e.\ $F=(dA)|_{\mathsf{F}}$ for some $A^{\mathrm{in}}$.
\end{proof}

\begin{lemma}[Ranks and cohomological descriptions of the obstruction groups]
\label{8.10:lem-obstruction-ranks}
Let $C$ be as in Definition~\ref{8.10:def-obstruction-group}. The following objects are those fixed in that
definition and in Lemma~\ref{8.10:lem-hull-boundary}:
\begin{itemize}
\item the regions $R_{\mathrm{int}}$ and $R$, the interior $R^{\circ}$ and the boundary $\partial R$ of $R$;
\item the set $\mathsf{R}$ of closed unit cubes of $R$, its boundary cubes $\mathsf{R}^{\partial}$, and its
deep cubes, namely the cubes of $\mathsf{R}$ meeting $R^{\circ}$;
\item the cells of the cell structure referred to in part~(c) of Lemma~\ref{8.10:lem-hull-boundary};
\item the set $\mathsf{R}_0$ of vertices on which the degree-$0$ cochains of $C_\sharp$ live, and the set
$\mathsf{R}_1$ of interior edges;
\item the cochain complex $C_\sharp$ with coboundary $d_\sharp$, and the plaquette set $\mathsf{F}$;
\item the groups $\mathfrak{O}(C)$ and $\mathfrak{O}^\sharp(C)$, with their ranks
$\operatorname{lk}(C)$ and $\operatorname{lk}^\sharp(C)$.
\end{itemize}
We also use the hulls $\operatorname{hull}_r(\cdot)$, $r\ge1$, the parameters $p_1$ and $p_2$, the number
$\ell$ entering the plaquette count in~(a), the cubical complex $\bar K(X)$ of a region $X$, and the
patching condition on a region, with the meanings fixed earlier in this chapter.
The box $\mathbb{B}$ and its subcomplexes $\mathbb{E}_B$ and $E_B$ are those of
Definition~\ref{8.10:def-spanned-box}. We write $n(\cdot)$ for the number of cells, and
$c_d=\binom{d}{2}=d(d-1)/2$, so that $c_4=6$; a $d$-dimensional cell has $c_d\cdot2^{d-2}$
two-dimensional faces.
\begin{itemize}
\item[(a)] One has
\begin{align*}
\operatorname{lk}(C)&\le c_d\,n(R_{\mathrm{int}})\le c_d(2p_1+1)^d\,n(C),\\
\operatorname{lk}^\sharp(C)&\le|\mathsf{F}|\le c_d(\ell+1)^d(2p_1+3)^d\,n(C).
\end{align*}
For $d=4$ the first line gives $\operatorname{lk}(C)\le 6(2p_1+1)^4\,n(C)$.
\item[(b)] Assume that the patching condition holds on $\operatorname{hull}_{p_1+p_2+1}(C)$. Then
\begin{align*}
\mathfrak{O}^\sharp(C)&\cong H^1(\mathbb{E}_B;\mathbb{Z})
\cong\operatorname{Hom}\bigl(H_1(|\mathbb{E}_B|;\mathbb{Z}),\mathbb{Z}\bigr),\\
\mathfrak{O}(C)&\cong H^1(E_B;\mathbb{Z})
\cong\operatorname{Hom}\bigl(H_1(|E_B|;\mathbb{Z}),\mathbb{Z}\bigr).
\end{align*}
Both groups are free abelian. Moreover
\begin{equation*}
\operatorname{lk}^\sharp(C)=\dim H^1(\mathbb{E}_B;\mathbb{R}),\qquad
\operatorname{lk}(C)=\dim H^1(E_B;\mathbb{R}).
\end{equation*}
\item[(c)] For $q\ge1$ let $C^q_{\mathrm{deep}}$ be the space of functions on the deep $q$-cubes. In
degrees $q\ge1$, $C_{\mathrm{deep}}$ is a subcomplex of $C_\sharp$. The quotient in degree $q\ge1$ is
the space $\mathrm{Sh}^q$ of functions on those $q$-cubes that have an edge in $\mathsf{R}_1$ and do
not meet $R^{\circ}$. All the complexes in this item are taken in degrees $q\ge1$ and set equal to zero
in degree $0$; this does not change $H^2$, which reads only the degrees $1$, $2$ and $3$. Consequently
\begin{equation*}
\operatorname{lk}^\sharp(C)\le\operatorname{lk}(C)+\dim H^2(\mathrm{Sh}\otimes\mathbb{R}).
\end{equation*}
Equality $\operatorname{lk}^\sharp(C)=\operatorname{lk}(C)$ is neither proved nor used here.
\end{itemize}
\end{lemma}

\begin{proof}
We use the following standard facts of algebraic topology, which are among those of
Notation~\ref{8.10:not-topology-facts}; the numbering (1)--(4) below is used in this proof only:
\begin{itemize}
\item[(1)] the cohomology of a finite cubical or CW pair is computed by its relative cellular cochains,
so that its real dimension in degree $q$ is at most the number of relative $q$-cells;
\item[(2)] a contractible space has vanishing cohomology in positive degrees;
\item[(3)] every pair has a long exact cohomology sequence;
\item[(4)] for a finite complex $X$, $H^1(X;\mathbb{Z})\cong\operatorname{Hom}(H_1(X;\mathbb{Z}),\mathbb{Z})$,
and this group is free abelian.
\end{itemize}

\emph{Part (a).} By part~(c) of Lemma~\ref{8.10:lem-hull-boundary}, the pair $(R,\partial R)$ is a CW
pair also for the cell structure. Its relative $2$-cochains are therefore the functions on the deep
$2$-faces of cells.

We count these faces. A cell has $c_d\cdot2^{d-2}$ two-faces. A deep $2$-face is counted once by each of
the $2^{d-2}$ cells containing it. Summing over the $n(R_{\mathrm{int}})$ cells, we find at most
$c_d\,n(R_{\mathrm{int}})$ deep $2$-faces.

The real cohomology of the pair is computed by these relative cochains. Hence
\begin{equation*}
\dim H^2(R,\partial R;\mathbb{R})\le c_d\,n(R_{\mathrm{int}}).
\end{equation*}
The left side is $\operatorname{lk}(C)$. The cell count of the hulls gives
$n(R_{\mathrm{int}})\le(2p_1+1)^d\,n(C)$, and this proves the first line. For $d=4$ one has $c_4=6$,
which gives the stated special case.

The first inequality of the second line, $\operatorname{lk}^\sharp(C)\le|\mathsf{F}|$, is part of
Definition~\ref{8.10:def-obstruction-group}. For the second inequality we use part~(a) of the lemma of
this chapter that places every cube with an interior edge inside
$\bar K(\operatorname{hull}_1R_{\mathrm{int}})$, together
with the cell count of the hulls. They give
\begin{align*}
|\mathsf{F}|
&\le\#\{\text{plaquettes of }\bar K(\operatorname{hull}_1R_{\mathrm{int}})\}
\le c_d(\ell+1)^d\,n(\operatorname{hull}_1R_{\mathrm{int}})\\
&\le c_d(\ell+1)^d(2p_1+3)^d\,n(C).
\end{align*}

\emph{Part (b).} The relative cochain complex of the CW pair $(\mathbb{B},\mathbb{E}_B)$ consists of
the functions on the cubes of $\mathbb{B}$ that do not belong to $\mathbb{E}_B$. By
Definition~\ref{8.10:def-spanned-box}, these are the cubes of $\mathbb{B}$ having an edge in
$\mathsf{R}_1$, together with the deep vertices.

By the same part~(a) of the lemma used in the proof of~(a) above, every cube with an interior edge lies in
$\bar K(\operatorname{hull}_1R_{\mathrm{int}})\subset\mathbb{B}$. Hence, in degrees $\ge1$, the relative
cochain complex of $(\mathbb{B},\mathbb{E}_B)$ is $C_\sharp$. The coboundaries also agree, since both are
computed in the ambient cubical complex.

The two complexes differ only in degree $0$: there the relative complex has the deep vertices and
$C_\sharp$ has $\mathsf{R}_0$. A second cohomology group involves only the degrees $1$, $2$ and~$3$, so
this difference does not affect it. Therefore
\begin{equation*}
\mathfrak{O}^\sharp(C)=H^2(C_\sharp;\mathbb{Z})\cong H^2(\mathbb{B},\mathbb{E}_B;\mathbb{Z}).
\end{equation*}

The realisation $|\mathbb{B}|$ is a box, hence contractible. By facts~(1) and~(2),
$H^1(\mathbb{B})=H^2(\mathbb{B})=0$. The long exact sequence of the pair
$(\mathbb{B},\mathbb{E}_B)$ contains the exact segment
\begin{equation*}
H^1(\mathbb{B})\longrightarrow H^1(\mathbb{E}_B)\longrightarrow H^2(\mathbb{B},\mathbb{E}_B)
\longrightarrow H^2(\mathbb{B}).
\end{equation*}
Both ends vanish, so the middle map is an isomorphism. This gives
$\mathfrak{O}^\sharp(C)\cong H^1(\mathbb{E}_B;\mathbb{Z})$.

By fact~(4), $H^1(\mathbb{E}_B;\mathbb{Z})$ is isomorphic to
$\operatorname{Hom}(H_1(|\mathbb{E}_B|;\mathbb{Z}),\mathbb{Z})$ and is free abelian.

The identification of the relative cochain complex with $C_\sharp$ in degrees $\ge1$ is a statement about
cochains, so it persists after tensoring with $\mathbb{R}$. The same exact segment with real
coefficients then gives $H^2(C_\sharp\otimes\mathbb{R})\cong H^1(\mathbb{E}_B;\mathbb{R})$. By
Definition~\ref{8.10:def-obstruction-group}, $\operatorname{lk}^\sharp(C)=\dim H^2(C_\sharp\otimes\mathbb{R})$,
and so $\operatorname{lk}^\sharp(C)=\dim H^1(\mathbb{E}_B;\mathbb{R})$.

The same argument applies to $E_B$. The relative cochains of $(\mathbb{B},E_B)$ are the functions on the
cubes of $\mathbb{B}$ meeting $R^{\circ}$. By part~(a) of Lemma~\ref{8.10:lem-hull-boundary}, such cubes
lie in $\mathsf{R}$, so they are exactly the deep cubes. The functions on the deep cubes are the
relative cochains of $(\mathsf{R},\mathsf{R}^{\partial})$, with the same coboundary computed in the
ambient complex. Hence
\begin{equation*}
H^2(\mathbb{B},E_B;\mathbb{Z})\cong H^2(\mathsf{R},\mathsf{R}^{\partial};\mathbb{Z})=\mathfrak{O}(C).
\end{equation*}

The exact segment above, with $E_B$ in place of $\mathbb{E}_B$, then gives
$\mathfrak{O}(C)\cong H^1(E_B;\mathbb{Z})$. By fact~(4) this group is isomorphic to
$\operatorname{Hom}(H_1(|E_B|;\mathbb{Z}),\mathbb{Z})$ and is free abelian. Repeating the argument with
real coefficients gives
\begin{equation*}
\operatorname{lk}(C)=\dim H^2(\mathsf{R},\mathsf{R}^{\partial};\mathbb{R})=\dim H^1(E_B;\mathbb{R}).
\end{equation*}

\emph{Part (c).} Let $D$ be a closed unit cube meeting $R^{\circ}$. By part~(a) of
Lemma~\ref{8.10:lem-hull-boundary}, $D\subset R$, so $D$ lies in $\mathsf{R}$. If its dimension is at
least one, it has interior edges. Thus every deep $q$-cube with $q\ge1$ has an edge in $\mathsf{R}_1$, and
$C^q_{\mathrm{deep}}\subset C^q_\sharp$ for $q\ge1$.

Let $\phi$ be a function supported on deep cubes. Its coboundary $d_\sharp\phi$ is supported on cubes
having a deep face. Such a cube meets $R^{\circ}$, because its deep face does, so it is deep. Hence
$d_\sharp$ maps $C_{\mathrm{deep}}$ into itself in degrees $\ge1$, and $C_{\mathrm{deep}}$ is a
subcomplex of $C_\sharp$ there.

For $q\ge1$, $C^q_\sharp$ is the space of functions on the $q$-cubes having an edge in $\mathsf{R}_1$
(Definition~\ref{8.10:def-obstruction-group}; see also the proof of~(b)). $C^q_{\mathrm{deep}}$ is the
space of functions on those of them that meet $R^{\circ}$. The
quotient is therefore the space of functions on the $q$-cubes having an edge in $\mathsf{R}_1$ and not
meeting $R^{\circ}$, which is $\mathrm{Sh}^q$. Its coboundary is the one induced by $d_\sharp$.

Take all three complexes in degrees $\ge1$, setting them equal to zero in degree $0$; this does not
change their second cohomology, which reads only the degrees $1$, $2$ and $3$. Tensor the short exact
sequence of complexes
\begin{equation*}
0\longrightarrow C_{\mathrm{deep}}\longrightarrow C_\sharp\longrightarrow\mathrm{Sh}\longrightarrow0
\end{equation*}
with $\mathbb{R}$. Its long exact cohomology sequence contains the exact segment
\begin{equation*}
H^2(C_{\mathrm{deep}}\otimes\mathbb{R})\longrightarrow H^2(C_\sharp\otimes\mathbb{R})
\longrightarrow H^2(\mathrm{Sh}\otimes\mathbb{R}),
\end{equation*}
in which only the degrees $1$, $2$ and $3$ enter.

The complex $C_{\mathrm{deep}}$ is the relative cochain complex of $(\mathsf{R},\mathsf{R}^{\partial})$ in
degrees $\ge1$, so $H^2(C_{\mathrm{deep}})=\mathfrak{O}(C)$. Exactness at the middle term bounds the
dimension of $H^2(C_\sharp\otimes\mathbb{R})$ by the dimension of the image of the first map plus that of
the third space. This gives
\begin{equation*}
\operatorname{lk}^\sharp(C)\le\dim H^2(C_{\mathrm{deep}}\otimes\mathbb{R})
+\dim H^2(\mathrm{Sh}\otimes\mathbb{R})=\operatorname{lk}(C)+\dim H^2(\mathrm{Sh}\otimes\mathbb{R}).
\end{equation*}
\end{proof}

\begin{remark}[A count of exterior classes by Alexander duality]\label{8.10:rem-duality-count}
Let $C$, $p_1$ and $n(C)$ be as in Lemma~\ref{8.10:lem-obstruction-ranks}, let $\mathsf{R}$ be
the hull of $C$ whose cells are counted in Lemma~\ref{8.10:lem-obstruction-ranks}(a), and write
$|\mathsf{R}|\subset\mathbb{R}^d$ for the union of
its closed cells. The \emph{touch-graph} of $\mathsf{R}$ has the closed cells of $\mathsf{R}$ as
vertices, and an edge between two distinct cells whenever they intersect.

The homology classes of loops in the exterior $\mathbb{R}^d\setminus|\mathsf{R}|$, that is, the
elements of $\tilde H_1(\mathbb{R}^d\setminus|\mathsf{R}|;\mathbb{Z})$, may be counted by the
minimal number of generators of
\begin{equation}\label{8.10:eq-duality-count-1}
H^{d-2}(|\mathsf{R}|;\mathbb{Z})\cong\tilde H_1(\mathbb{R}^d\setminus|\mathsf{R}|;\mathbb{Z}),
\end{equation}
where the isomorphism is Alexander duality (Notation~\ref{8.10:not-topology-facts}).
This number is bounded by
\begin{equation*}
3^{d(d-2)}(2p_1+1)^d\,n(C),
\end{equation*}
a bound obtained from the nerve theorem (Notation~\ref{8.10:not-topology-facts}) applied to the
cover of $|\mathsf{R}|$ by its closed cells; by Helly's property of boxes (a finite family of
pairwise intersecting boxes has a common point), the nerve of this cover
is the flag complex of the touch-graph.

Suppose $|\mathsf{R}|$ is a manifold with boundary. Then, by Lefschetz duality
(Notation~\ref{8.10:not-topology-facts}), the rank of the group
\eqref{8.10:eq-duality-count-1} agrees with the rank $\operatorname{lk}(C)$ of
Lemma~\ref{8.10:lem-obstruction-ranks}. When $|\mathsf{R}|$ is not a manifold with boundary
the two counts may differ.

The two dualities are the facts of Notation~\ref{8.10:not-topology-facts}, and none of the
statements of this remark is used elsewhere in this book.
\end{remark}

\begin{proof}[Proof of the bound]
The closed cells of $\mathsf{R}$ are closed lattice cubes of one side length, and by
Lemma~\ref{8.10:lem-obstruction-ranks}(a) their number is at most $(2p_1+1)^d\,n(C)$.

Every non-empty intersection of closed cells is a box, hence convex and contractible. The nerve
theorem (Notation~\ref{8.10:not-topology-facts}) therefore shows that $|\mathsf{R}|$ is homotopy
equivalent to the nerve $\mathcal{N}$ of the cover of $|\mathsf{R}|$ by its closed cells, so that
$H^{d-2}(|\mathsf{R}|;\mathbb{Z})\cong H^{d-2}(\mathcal{N};\mathbb{Z})$. By Helly's property of
boxes, a set of cells spans a simplex of $\mathcal{N}$ exactly when its members touch pairwise,
that is, $\mathcal{N}$ is the flag complex of the touch-graph.

Let $\#\mathcal{N}_{d-2}$ be the number of $(d-2)$-simplices of $\mathcal{N}$. The group
$H^{d-2}(\mathcal{N};\mathbb{Z})$ is a quotient of a subgroup of the cochain group
$C^{d-2}(\mathcal{N};\mathbb{Z})$, which is free abelian of rank $\#\mathcal{N}_{d-2}$; a
subgroup of a free abelian group of finite rank is free abelian of at most that rank, so its
quotient is generated by at most $\#\mathcal{N}_{d-2}$ elements.

A $(d-2)$-simplex is a set of $d-1$ pairwise touching cells. Fix one of its cells $\square$; each
of the other $d-2$ cells touches $\square$. A closed cube of the lattice meets exactly $3^d-1$
other cubes of the same lattice, namely those whose integer coordinates differ from those of
$\square$ by at most one in every direction. Hence at most $(3^d-1)^{d-2}\le 3^{d(d-2)}$
simplices of dimension $d-2$ contain $\square$, and
\begin{equation*}
\#\mathcal{N}_{d-2}\le 3^{d(d-2)}\cdot\#\{\text{cells of }\mathsf{R}\}
\le 3^{d(d-2)}(2p_1+1)^d\,n(C).
\end{equation*}
Together with the isomorphism \eqref{8.10:eq-duality-count-1} this gives the bound.
\end{proof}

\begin{lemma}[Filling from a small relative gauge]\label{8.10:lem-gauge-filling}
We use the objects $\mathsf{N}$, $\mathsf{A}_0$, $\mathsf{R}_0$, $\mathsf{R}_1$, $\mathsf{R}_2$,
$\mathsf{I}_0$, $\mathrm{Str}$, $\mathcal{V}^{\mathrm{ext}}$ and $F_P$ as in
Lemma~\ref{8.10:lem-collar-plaquettes}, and we write $G$ for the gauge group. Let
$V^{\mathrm{ext}}\in\mathcal{V}^{\mathrm{ext}}$ and $\eta\ge0$, and suppose that $u\colon\mathsf{A}_0\to G$
satisfies
\begin{equation}\label{8.10:eq-gauge-filling-1}
\bigl\|u(v)\,V^{\mathrm{ext}}_{(v,v')}\,u(v')^{-1}-1\bigr\|\le\eta
\end{equation}
for every exterior edge $(v,v')$ of $\mathsf{N}$ that is an edge of some plaquette of $\mathrm{Str}$.
Define $\tilde u\colon\mathsf{R}_0\to G$ by $\tilde u:=u$ on
$\mathsf{I}_0=\mathsf{R}_0\cap\mathsf{A}_0$ and $\tilde u:=1$ on
$\mathsf{R}_0\setminus\mathsf{A}_0$, and put $V^{\mathrm{in}}_{(v,v')}:=\tilde u(v)^{-1}\tilde u(v')$ for
$(v,v')\in\mathsf{R}_1$. Then
\begin{equation*}
F_P(V^{\mathrm{in}},V^{\mathrm{ext}})=0\quad(P\in\mathsf{R}_2),\qquad
F_P(V^{\mathrm{in}},V^{\mathrm{ext}})\le3\eta\quad(P\in\mathrm{Str}).
\end{equation*}
\end{lemma}

\begin{proof}
Write $V$ for the configuration equal to $V^{\mathrm{in}}$ on $\mathsf{R}_1$ and to $V^{\mathrm{ext}}$
on the exterior edges. All group elements involved are unitary, and the norm is invariant under
multiplication by unitaries on either side. In particular, reversing an edge inverts its variable, and
$\|\Gamma^{-1}-1\|=\|\Gamma(\Gamma^{-1}-1)\|=\|1-\Gamma\|$ for unitary $\Gamma$. Hence
\eqref{8.10:eq-gauge-filling-1} holds for both orientations of each edge concerned, and the formula
defining $V^{\mathrm{in}}$ is consistent under reversal.

\emph{Plaquettes of $\mathsf{R}_2$.} Let $P\in\mathsf{R}_2$, and let $v_0,v_1,\dots,v_4=v_0$ be the
vertices of its boundary circuit. The transport around $\partial P$ is
\begin{equation*}
V(\partial P)=\prod_{i=1}^{4}\tilde u(v_{i-1})^{-1}\tilde u(v_i)=\tilde u(v_0)^{-1}\tilde u(v_4)=1,
\end{equation*}
since the product telescopes. Hence $F_P(V^{\mathrm{in}},V^{\mathrm{ext}})=0$.

\emph{Plaquettes of $\mathrm{Str}$.} Let $P\in\mathrm{Str}$. By
Lemma~\ref{8.10:lem-collar-plaquettes}(c), the boundary circuit $\partial P$ has exactly one maximal
run $\sigma$ of exterior edges. The run has $|\sigma|\in\{2,3\}$ edges, and its two endpoints lie in
$\mathsf{I}_0$. Write
\begin{equation*}
\sigma=(v_0,v_1,\dots,v_{|\sigma|}),\qquad v_0,\;v_{|\sigma|}\in\mathsf{I}_0 .
\end{equation*}
The remaining part of $\partial P$ is a path of edges of $\mathsf{R}_1$ from $v_{|\sigma|}$ to $v_0$.
By the telescoping of the previous case, its transport is $\tilde u(v_{|\sigma|})^{-1}\tilde u(v_0)$.
Let $T_\sigma$ be the transport of $V^{\mathrm{ext}}$ along $\sigma$. Reading $\partial P$ from
$v_{|\sigma|}$, we obtain
\begin{equation*}
V(\partial P)=\tilde u(v_{|\sigma|})^{-1}\tilde u(v_0)\,T_\sigma .
\end{equation*}
Since $\tilde u=u$ on $\mathsf{I}_0$, conjugating by the unitary $\tilde u(v_{|\sigma|})$ gives
\begin{equation*}
F_P(V^{\mathrm{in}},V^{\mathrm{ext}})=\bigl\|u(v_0)\,T_\sigma\,u(v_{|\sigma|})^{-1}-1\bigr\|.
\end{equation*}

By Lemma~\ref{8.10:lem-collar-plaquettes}(b), every exterior edge of $\mathsf{N}$ lies in
$\mathsf{Ex}_1$, and $\mathsf{Ex}_0=\mathsf{A}_0$. Hence every vertex $v_i$ of $\sigma$ lies in
$\mathsf{A}_0$, and $u(v_i)$ is defined. Inserting $u(v_i)^{-1}u(v_i)$ between consecutive factors of
$T_\sigma$ gives
\begin{equation*}
u(v_0)\,T_\sigma\,u(v_{|\sigma|})^{-1}=\prod_{i=1}^{|\sigma|}\Gamma_i,\qquad
\Gamma_i:=u(v_{i-1})\,V^{\mathrm{ext}}_{(v_{i-1},v_i)}\,u(v_i)^{-1}.
\end{equation*}
Each edge $(v_{i-1},v_i)$ is an exterior edge of $\mathsf{N}$ and an edge of the plaquette
$P\in\mathrm{Str}$. By \eqref{8.10:eq-gauge-filling-1}, each $\Gamma_i$ is therefore a unitary with
$\|\Gamma_i-1\|\le\eta$. From the identity
\begin{equation*}
\Gamma_1\cdots\Gamma_{|\sigma|}-1=\sum_{i=1}^{|\sigma|}\Gamma_1\cdots\Gamma_{i-1}(\Gamma_i-1)
\end{equation*}
and the unitary invariance of the norm, we get
\begin{equation*}
F_P(V^{\mathrm{in}},V^{\mathrm{ext}})\le\sum_{i=1}^{|\sigma|}\|\Gamma_i-1\|\le|\sigma|\,\eta\le3\eta. \qedhere
\end{equation*}
\end{proof}

\begin{lemma}[Small relative gauge from a flat collar]\label{8.10:lem-collar-gauge}
Assume the patching hypothesis on the collar $\mathsf{Ex}$ and condition $(\mathrm{U}_1)$ of
Definition~\ref{8.10:def-linked-components}. Let $\mathfrak{F}$ be a spanning forest of the graph of
$\mathsf{Ex}$, with a root chosen in each connected component. For an edge
$e=(y,z)\in\mathsf{Ex}_1\setminus\mathfrak{F}$ let $c_e$ be the fundamental cycle of $e$: the path in
$\mathfrak{F}$ from the root to $y$, followed by $e$, followed by the path in $\mathfrak{F}$ from $z$
to the root. Let $k_e$ be the least area of a null-homotopy of $c_e$ in $\mathsf{Ex}$, that is, the
least integer $A$ such that $c_e$ admits a null-homotopy across plaquettes $(q_1,\dots,q_A)$ of
$\mathsf{Ex}_2$ in the sense of Lemma~\ref{8.10:lem-homotopy-moves}. Put
\begin{equation}\label{8.10:eq-collar-gauge-1}
\mathfrak{a}(\mathsf{Ex}):=\min_{\mathfrak{F}}\ \max_{e\in\mathsf{Ex}_1\setminus\mathfrak{F}} k_e ,
\end{equation}
which is finite under $(\mathrm{U}_1)$, with the convention $\mathfrak{a}(\mathsf{Ex}):=1$ when the
graph of $\mathsf{Ex}$ is a forest, so that $\mathsf{Ex}_1\subset\mathfrak{F}$ and the maximum in
\eqref{8.10:eq-collar-gauge-1} is over the empty set. Suppose that the configuration
$V^{\mathrm{ext}}$ of the patching hypothesis is $\delta$-flat on $\mathsf{Ex}$, so that
$F_q(V^{\mathrm{ext}})\le\delta$ for
every plaquette $q\in\mathsf{Ex}_2$, with $F_q$ and the transports $T_\gamma$ as in
Lemma~\ref{8.10:lem-homotopy-moves}. Let
$\mathfrak{F}$ realise the minimum in \eqref{8.10:eq-collar-gauge-1} and, for every $y$ in the vertex
set $\mathsf{A}_0$ fixed in the patching hypothesis, to which the vertices of the graph of
$\mathsf{Ex}$ belong, let
$u(y):=T_{\mathrm{root}\rightsquigarrow y}(V^{\mathrm{ext}})$ be the transport of $V^{\mathrm{ext}}$
along the path in $\mathfrak{F}$ from the root of the component of $y$ to $y$. Then for every edge
$e=(y,z)$ of $\mathsf{Ex}$
\begin{equation}\label{8.10:eq-collar-gauge-2}
\bigl\|\,u(y)\,V^{\mathrm{ext}}_e\,u(z)^{-1}-1\,\bigr\|\le\mathfrak{a}(\mathsf{Ex})\,\delta .
\end{equation}
\end{lemma}

\begin{proof}
We treat forest edges and the remaining edges separately.

Let $e=(y,z)$ be an edge of $\mathfrak{F}$, with $z$ the child of $y$. The path in $\mathfrak{F}$
from the root to $z$ is the path from the root to $y$ followed by $e$, so that
\begin{equation*}
T_{\mathrm{root}\rightsquigarrow z}(V^{\mathrm{ext}})
=T_{\mathrm{root}\rightsquigarrow y}(V^{\mathrm{ext}})\,V^{\mathrm{ext}}_e ,
\end{equation*}
that is, $u(z)=u(y)V^{\mathrm{ext}}_e$. Hence $u(y)V^{\mathrm{ext}}_e u(z)^{-1}=1$, the left side of
\eqref{8.10:eq-collar-gauge-2} vanishes, and the bound holds. In particular, if
$\mathsf{Ex}_1\subset\mathfrak{F}$, every edge is of this kind.

Let now $e=(y,z)\in\mathsf{Ex}_1\setminus\mathfrak{F}$. The inverse of the transport along the path in
$\mathfrak{F}$ from the root to $z$ is the transport along the same path traversed from $z$ to the
root. Therefore
\begin{equation*}
u(y)\,V^{\mathrm{ext}}_e\,u(z)^{-1}
=T_{\mathrm{root}\rightsquigarrow y}(V^{\mathrm{ext}})\,V^{\mathrm{ext}}_e\,
T_{z\rightsquigarrow\mathrm{root}}(V^{\mathrm{ext}})
=T_{c_e}(V^{\mathrm{ext}}),
\end{equation*}
the transport of $V^{\mathrm{ext}}$ around the fundamental cycle $c_e$, which is a loop at the root.
By the definition of $k_e$, the loop $c_e$ admits a null-homotopy in $\mathsf{Ex}$ across $k_e$
plaquettes $q_1,\dots,q_{k_e}$ of $\mathsf{Ex}_2$. Since $\mathfrak{F}$ realises the minimum in
\eqref{8.10:eq-collar-gauge-1}, we have $k_e\le\mathfrak{a}(\mathsf{Ex})$. Since $V^{\mathrm{ext}}$ is
$\delta$-flat on $\mathsf{Ex}$, each of these plaquettes satisfies $F_{q_i}(V^{\mathrm{ext}})\le\delta$.
Lemma~\ref{8.10:lem-homotopy-moves}, applied to the loop $c_e$, the null-homotopy across
$(q_1,\dots,q_{k_e})$ and the configuration $V^{\mathrm{ext}}$, gives
\begin{equation*}
\bigl\|T_{c_e}(V^{\mathrm{ext}})-1\bigr\|
\le\sum_{i=1}^{k_e}F_{q_i}(V^{\mathrm{ext}})
\le k_e\,\delta
\le\mathfrak{a}(\mathsf{Ex})\,\delta .
\end{equation*}
This is \eqref{8.10:eq-collar-gauge-2} for $e$, and the proof is complete.
\end{proof}

\begin{lemma}[Comb gauge on a box]\label{8.10:lem-comb-gauge}
Let
\begin{equation*}
\mathcal{Q}=\prod_{i=1}^{d}[\alpha_i,\beta_i]\subset\mathbb{Z}^d
\end{equation*}
be a lattice box with sides $s_i:=\beta_i-\alpha_i\ge 1$, and put $\bar s:=\max_i s_i$. Let a configuration
$V$ be given on all edges of $K(\mathcal{Q})$, the lattice cell complex of $\mathcal{Q}$, and suppose that
$F_\sigma\le\eta$ for every plaquette $\sigma\subset\mathcal{Q}$; here $V_e$, the transport $T(P)$ along an
edge path $P$, and $F_\sigma$ are as in the ladder identity (Lemma~\ref{8.10:lem-ladder-identity}). Then
there is a function $u$ on the vertices of $\mathcal{Q}$ such that
\begin{equation}\label{8.10:eq-comb-gauge-1}
\bigl\|u(x)\,V_{(x,y)}\,u(y)^{-1}-1\bigr\|\le (d-1)\,\bar s\,\eta
\qquad\text{for every edge }(x,y)\subset\mathcal{Q}.
\end{equation}
\end{lemma}

\begin{proof}
Write $o:=(\alpha_1,\ldots,\alpha_d)$ for the lower corner of $\mathcal{Q}$. For a vertex $x$ of
$\mathcal{Q}$ let $P(x)$ be the comb path from $o$ to $x$: it runs first in direction $1$ from $o$ to
$(x_1,\alpha_2,\ldots,\alpha_d)$, then in direction $2$ to $(x_1,x_2,\alpha_3,\ldots,\alpha_d)$, and so
on, the last leg running in direction $d$ to $x$. Every vertex of $P(x)$ has its $\nu$-th coordinate
between $\alpha_\nu$ and $x_\nu$, so $P(x)$ lies in $\mathcal{Q}$. Put $u(x):=T(P(x))$.

Let $e=(x,x+e_\mu)$ be an edge of $\mathcal{Q}$. Since the transport of a concatenation is the product of
the transports and the transport of a reversed path is the inverse,
\begin{equation*}
u(x)\,V_e\,u(x+e_\mu)^{-1}=T\bigl(P(x)\cdot e\cdot P(x+e_\mu)^{-1}\bigr).
\end{equation*}
The vertices $x$ and $x+e_\mu$ have the same coordinates of index $\ne\mu$, and the $\mu$-th coordinate of
$x+e_\mu$ exceeds that of $x$ by one. Hence the two combs share the initial segment $P_0$ from $o$ to
\begin{equation*}
a:=(x_1,\ldots,x_\mu,\alpha_{\mu+1},\ldots,\alpha_d).
\end{equation*}
After $P_0$, the comb $P(x)$ continues by the staircase $S_x$ from $a$ to $x$ in the directions
$\mu+1,\ldots,d$, while $P(x+e_\mu)$ continues by the edge $\bar e:=(a,a+e_\mu)$ followed by the
translated staircase $S_x+e_\mu$ from $a+e_\mu$ to $x+e_\mu$. Therefore the loop above is
\begin{equation*}
P_0\cdot\bigl[S_x\cdot e\cdot(S_x+e_\mu)^{-1}\cdot\bar e^{-1}\bigr]\cdot P_0^{-1},
\end{equation*}
and its transport is
\begin{equation*}
T(P_0)\,\bigl[T(S_x)\,V_e\,T(S_x+e_\mu)^{-1}\,V_{\bar e}^{-1}\bigr]\,T(P_0)^{-1}.
\end{equation*}

We apply Lemma~\ref{8.10:lem-ladder-identity} with $P=S_x$, with vertices $v_0=a,\ldots,v_n=x$, with the
unit vector $e_\mu$, and so with $P'=S_x+e_\mu$. The rungs are $c_k=(v_k,v_k+e_\mu)$; the first rung is
$c_0=\bar e$ and the last rung is $c_n=e$. The plaquette $\sigma_k$ has the vertices $v_{k-1}$, $v_k$,
$v_k+e_\mu$, $v_{k-1}+e_\mu$. The vertices $v_k$ lie in $\mathcal{Q}$, since their $\nu$-th coordinates
are between $\alpha_\nu$ and $x_\nu$. The vertices $v_k+e_\mu$ lie in $\mathcal{Q}$ as well, since their
$\mu$-th coordinate is $x_\mu+1\le\beta_\mu$, the edge $e$ being in $\mathcal{Q}$. Hence every
$\sigma_k\subset\mathcal{Q}$ and $F_{\sigma_k}\le\eta$. The ladder identity then gives
\begin{equation*}
\bigl\|T(S_x)\,V_e\,T(S_x+e_\mu)^{-1}\,V_{\bar e}^{-1}-1\bigr\|\le\sum_{k=1}^{n}F_{\sigma_k}\le n\,\eta .
\end{equation*}
When $\mu=d$ the staircase is empty, $n=0$, $a=x$, $\bar e=e$, and the bracket equals $1$.

The length of $S_x$ is
\begin{equation*}
n=\sum_{\nu>\mu}(x_\nu-\alpha_\nu)\le\sum_{\nu>\mu}s_\nu\le(d-1)\,\bar s .
\end{equation*}
Finally, $T(P_0)$ is unitary, and for unitary $W$ the identity $WXW^{-1}-1=W(X-1)W^{-1}$ shows that
$\|WXW^{-1}-1\|=\|X-1\|$. So the transport of the loop, being the conjugate by $T(P_0)$ of the unitary
in brackets, is within $(d-1)\bar s\eta$ of $1$. This is \eqref{8.10:eq-comb-gauge-1} for the edge
$(x,x+e_\mu)$.

For the same edge traversed in the opposite direction, with $V_{(x+e_\mu,x)}=V_{(x,x+e_\mu)}^{-1}$, the
gauge-transformed variable is $W^{-1}$, where $W:=u(x)V_eu(x+e_\mu)^{-1}$ is unitary. The identity
$W^{-1}-1=-W^{-1}(W-1)$ gives $\|W^{-1}-1\|=\|W-1\|$, and the same bound follows.
\end{proof}

Throughout Lemma~\ref{8.10:lem-box-filling} and its proof the following conventions are in force. A configuration on a set of lattice edges assigns to each edge $b$ of the set, oriented in the positive coordinate direction, an element $V_b\in\mathrm{SU}(N)$; for the reversed edge $b^{-1}$ we put $V_{b^{-1}}:=V_b^{-1}$. For an edge path $\gamma=b_1b_2\cdots b_n$ we write $T(\gamma):=V_{b_1}V_{b_2}\cdots V_{b_n}$, the empty product being $1$. For a plaquette $\sigma$ we write $V(\partial\sigma)$ for $T$ of its boundary loop, and we put $F_\sigma:=\|V(\partial\sigma)-1\|$, as in Lemma~\ref{8.10:lem-stokes-rectangle} and Lemma~\ref{8.10:lem-ladder-identity}. Here $\|\cdot\|$ is the matrix norm of those two lemmas. The only property of it used below is that $\|gAh\|=\|A\|$ for all matrices $A$ and all $g,h\in\mathrm{SU}(N)$.

For $X,Y,g,X_j,Y_j\in\mathrm{SU}(N)$ this property gives
\begin{equation}\label{8.10:eq-box-filling-1}
\|X^{-1}-Y^{-1}\|=\|X-Y\|,\qquad \|XY^{-1}-1\|=\|X-Y\|,\qquad \|gXg^{-1}-1\|=\|X-1\|.
\end{equation}
These follow from the identities $X^{-1}-Y^{-1}=X^{-1}(Y-X)Y^{-1}$, $XY^{-1}-1=(X-Y)Y^{-1}$ and $gXg^{-1}-1=g(X-1)g^{-1}$. The same property gives the telescoping bound
\begin{equation}\label{8.10:eq-box-filling-2}
\|X_1X_2\cdots X_n-Y_1Y_2\cdots Y_n\|\le\sum_{j=1}^{n}\|X_j-Y_j\|.
\end{equation}
This follows from the identity
\begin{equation*}
X_1\cdots X_n-Y_1\cdots Y_n=\sum_{j=1}^{n}X_1\cdots X_{j-1}\,(X_j-Y_j)\,Y_{j+1}\cdots Y_n .
\end{equation*}
Changing the base point of the boundary loop of $\sigma$ conjugates $V(\partial\sigma)$ by an element of $\mathrm{SU}(N)$, and reversing its orientation inverts $V(\partial\sigma)$. Hence, by \eqref{8.10:eq-box-filling-1}, $F_\sigma$ depends on neither choice.

\begin{lemma}[Rough filling of a box from face data]\label{8.10:lem-box-filling}
Let $d\ge 3$, and let $\delta\ge 0$. Let $\mathcal{Q}=\prod_{i=1}^{d}[\alpha_i,\beta_i]$ be a lattice box, with integers $\alpha_i<\beta_i$ whose sides satisfy $s_i:=\beta_i-\alpha_i\ge 2$ for all $i$. Put $\bar s:=\max_i s_i$. Let $\partial\mathcal{Q}$ denote the union of the faces of $\mathcal{Q}$, that is, the points of $\mathcal{Q}$ with some coordinate $x_i\in\{\alpha_i,\beta_i\}$. A lattice edge or plaquette contained in $\mathcal{Q}$ is said to lie in $\partial\mathcal{Q}$ if it is contained in a face of $\mathcal{Q}$.

Let a configuration be given on the edges lying in $\partial\mathcal{Q}$ such that every plaquette $\sigma$ lying in $\partial\mathcal{Q}$ has $F_\sigma\le\delta$. Then it extends to a configuration on all lattice edges contained in $\mathcal{Q}$ such that every plaquette $\sigma\subset\mathcal{Q}$ has
\begin{equation*}
F_\sigma\le 13\,\bar s^{2}\,\delta .
\end{equation*}
\end{lemma}

\begin{proof}
\noindent\textit{Decomposition of the box.}
Write $\mathcal{Q}=\mathcal{Q}'\times[\alpha_d,\beta_d]$ with $\mathcal{Q}':=\prod_{i<d}[\alpha_i,\beta_i]$, and write points as $(z',t)$ with $z'\in\mathcal{Q}'$ and $t\in[\alpha_d,\beta_d]$. Let $\partial\mathcal{Q}'$ be the union of the faces of $\mathcal{Q}'$. An edge or plaquette of $\mathcal{Q}'$ lies in $\partial\mathcal{Q}'$, written $e'\subset\partial\mathcal{Q}'$, if it is contained in a face of $\mathcal{Q}'$. Edges of the form $e'\times\{t\}$ are called horizontal, and edges $\{z'\}\times[t,t+1]$ vertical; vertical edges are oriented upwards.

The faces $\{x_d=\alpha_d\}$ and $\{x_d=\beta_d\}$ of $\mathcal{Q}$ contain no vertical edge. A face $\{x_i=c\}$ with $i<d$ contains $e'\times\{t\}$ if and only if $e'\subset\{x_i=c\}$, and it contains $\{z'\}\times[t,t+1]$ if and only if $z'_i=c$. Hence the edges lying in $\partial\mathcal{Q}$ are:
\begin{itemize}
\item the horizontal edges $e'\times\{t\}$ with $t\in\{\alpha_d,\beta_d\}$ (bottom and lid);
\item the horizontal edges $e'\times\{t\}$ with $e'\subset\partial\mathcal{Q}'$ (cylinder);
\item the vertical edges $\{z'\}\times[t,t+1]$ with $z'\in\partial\mathcal{Q}'$ (cylinder).
\end{itemize}
The remaining edges are the horizontal edges $e'\times\{t\}$ with $e'\not\subset\partial\mathcal{Q}'$ and $\alpha_d<t<\beta_d$, and the vertical edges over points $z'\in\mathcal{Q}'\setminus\partial\mathcal{Q}'$.

In the same way, the plaquettes lying in $\partial\mathcal{Q}$ are those of the bottom $\mathcal{Q}'\times\{\alpha_d\}$, those of the lid $\mathcal{Q}'\times\{\beta_d\}$, and the cylinder plaquettes. The cylinder plaquettes are the $p'\times\{t\}$ with $p'\subset\partial\mathcal{Q}'$ and the $e'\times[t,t+1]$ with $e'\subset\partial\mathcal{Q}'$. The other plaquettes of $\mathcal{Q}$ are of two types:
\begin{itemize}
\item[(a)] $p'\times\{t\}$ with $p'\not\subset\partial\mathcal{Q}'$ and $\alpha_d<t<\beta_d$;
\item[(b)] $e'\times[t,t+1]$ with $e'\not\subset\partial\mathcal{Q}'$.
\end{itemize}

\smallskip
\noindent\textit{Step 1 (column gauge).}
For a map $\lambda$ from the vertices of $\mathcal{Q}$ to $\mathrm{SU}(N)$ and a configuration $V$, put $(\lambda\cdot V)_{[x,y]}:=\lambda(x)V_{[x,y]}\lambda(y)^{-1}$. For a plaquette $\sigma$ whose boundary loop is based at $x$ we then have $(\lambda\cdot V)(\partial\sigma)=\lambda(x)V(\partial\sigma)\lambda(x)^{-1}$. By \eqref{8.10:eq-box-filling-1}, $F_\sigma$ is therefore the same for $V$ and for $\lambda\cdot V$.

Consequently, if $V^{\circ}$ denotes the given data, the transformed data $\lambda\cdot V^{\circ}$ on the edges lying in $\partial\mathcal{Q}$ satisfy the hypothesis of the lemma. Suppose $W$ extends $\lambda\cdot V^{\circ}$ to all edges of $\mathcal{Q}$ with the asserted bounds. Then $\lambda^{-1}\cdot W$, where $\lambda^{-1}$ is the pointwise inverse, extends $V^{\circ}$ and has the same values of $F_\sigma$. It therefore suffices to prove the lemma for the transformed data, which we choose as follows.

Let $\lambda\equiv 1$, except at the cylinder vertices $(z',t)$ with $z'\in\partial\mathcal{Q}'$ and $t>\alpha_d$. There we set $\lambda(z',t):=T(\text{column})$, the transport of $V^{\circ}$ along the vertical column from $(z',\alpha_d)$ up to $(z',t)$. With the empty product at $t=\alpha_d$, this formula also covers $t=\alpha_d$, where $\lambda=1$.

For a vertical cylinder edge $b=\{z'\}\times[t,t+1]$ we have $\lambda(z',t+1)=\lambda(z',t)V^{\circ}_b$. Hence
\begin{equation*}
(\lambda\cdot V^{\circ})_b=\lambda(z',t)\,V^{\circ}_b\,\bigl(\lambda(z',t)V^{\circ}_b\bigr)^{-1}=1 .
\end{equation*}
So every transformed vertical cylinder edge equals $1$. From now on $V$ denotes the transformed data, and the configuration to be constructed extends it.

For $e'=[x',y']\subset\partial\mathcal{Q}'$ and $\alpha_d\le t\le\beta_d$, let $\Gamma_t(e'):=V_{e'\times\{t\}}$. The plaquette $e'\times[t,t+1]$ lies in $\partial\mathcal{Q}$. Its boundary loop based at $(x',t)$ runs along $e'\times\{t\}$, up $\{y'\}\times[t,t+1]$, back along $e'\times\{t+1\}$ and down $\{x'\}\times[t,t+1]$. Its holonomy is therefore
\begin{equation*}
\Gamma_t(e')\cdot 1\cdot\Gamma_{t+1}(e')^{-1}\cdot 1 .
\end{equation*}
By the hypothesis and the second identity of \eqref{8.10:eq-box-filling-1}, $\|\Gamma_t(e')-\Gamma_{t+1}(e')\|\le\delta$. Summing these bounds over the $t-\alpha_d\le s_d$ layers between $\alpha_d$ and $t$, and over the $\beta_d-t\le s_d$ layers between $t$ and $\beta_d$, gives, for $\alpha_d\le t\le\beta_d$,
\begin{equation}\label{8.10:eq-box-filling-3}
\|\Gamma_t(e')-\Gamma_{\alpha_d}(e')\|\le s_d\,\delta,\qquad
\|\Gamma_t(e')-\Gamma_{\beta_d}(e')\|\le s_d\,\delta .
\end{equation}

For every edge $e'$ of $\mathcal{Q}'$ put
\begin{equation*}
V^{\downarrow}_{e'}:=V_{e'\times\{\alpha_d\}},\qquad V^{\uparrow}_{e'}:=V_{e'\times\{\beta_d\}} .
\end{equation*}
Both are given data, since bottom and lid edges lie in $\partial\mathcal{Q}$. When $e'\subset\partial\mathcal{Q}'$ we have $V^{\downarrow}_{e'}=\Gamma_{\alpha_d}(e')$ and $V^{\uparrow}_{e'}=\Gamma_{\beta_d}(e')$.

\smallskip
\noindent\textit{Step 2 (the extension).}
We assign values to the remaining edges as follows.
\begin{itemize}
\item For $e'\not\subset\partial\mathcal{Q}'$ and $\alpha_d<t<\beta_d$, set $V_{e'\times\{t\}}:=V^{\downarrow}_{e'}$.
\item For $z'\in\mathcal{Q}'\setminus\partial\mathcal{Q}'$ and $\alpha_d\le t\le\beta_d-2$, set $V_{\{z'\}\times[t,t+1]}:=1$.
\item The topmost vertical edges $\{z'\}\times[\beta_d-1,\beta_d]$ with $z'\in\mathcal{Q}'\setminus\partial\mathcal{Q}'$ receive values $\eta(z')\in\mathrm{SU}(N)$, defined in Step~4.
\end{itemize}
For $z'\in\partial\mathcal{Q}'$ we put $\eta(z'):=1$; this is the value of the transformed cylinder edge $\{z'\}\times[\beta_d-1,\beta_d]$ found in Step~1. Since $s_d\ge 2$, we have $\beta_d-1>\alpha_d$. By the description of the remaining edges above, every edge of $\mathcal{Q}$ not lying in $\partial\mathcal{Q}$ now receives exactly one value, and the data on $\partial\mathcal{Q}$ are unchanged.

A plaquette lying in $\partial\mathcal{Q}$ is contained in a face, so all its edges lie in that face and carry given data. Such plaquettes therefore keep $F_\sigma\le\delta$. It remains to bound the plaquettes of types (a) and (b).

\smallskip
\noindent\textit{Step 3 (types (a) and (b) below the top layer).}
(a) Let $p'\not\subset\partial\mathcal{Q}'$ and $\alpha_d<t<\beta_d$. The holonomy of $p'\times\{t\}$ is the ordered product, along the boundary loop of $p'$ and with inverses for reversed edges, of the four edge values at height $t$. Each of these is $V^{\downarrow}_{e'}$ if $e'\not\subset\partial\mathcal{Q}'$ (Step~2), or $\Gamma_t(e')$ if $e'\subset\partial\mathcal{Q}'$ (cylinder data).

The holonomy of the bottom plaquette $p'\times\{\alpha_d\}$, taken along the same loop, is the same ordered product with every factor replaced by its bottom value $V^{\downarrow}_{e'}$. For the edges $e'\subset\partial\mathcal{Q}'$ this bottom value is $\Gamma_{\alpha_d}(e')$. So the two products differ at most in the factors belonging to edges $e'\subset\partial\mathcal{Q}'$, of which there are at most four. For each such factor, \eqref{8.10:eq-box-filling-3} and the first identity of \eqref{8.10:eq-box-filling-1} (for reversed edges) give $\|\Gamma_t(e')^{\pm1}-\Gamma_{\alpha_d}(e')^{\pm1}\|\le s_d\delta$. By \eqref{8.10:eq-box-filling-2} the two holonomies therefore differ in norm by at most $4s_d\delta$.

The bottom plaquette lies in $\partial\mathcal{Q}$, so its holonomy is within $\delta$ of $1$. By the triangle inequality,
\begin{equation*}
\|V(\partial(p'\times\{t\}))-1\|\le(4s_d+1)\,\delta\le 5\bar s\,\delta .
\end{equation*}

(b) Let $e'=[x',y']\not\subset\partial\mathcal{Q}'$ and $\alpha_d\le t\le\beta_d-2$. The lower horizontal edge of $e'\times[t,t+1]$ carries $V^{\downarrow}_{e'}$: for $t=\alpha_d$ this is the definition of $V^{\downarrow}_{e'}$, and for $t>\alpha_d$ it is Step~2. The upper horizontal edge, at height $t+1$ with $\alpha_d<t+1\le\beta_d-1$, also carries $V^{\downarrow}_{e'}$ by Step~2.

The two vertical edges $\{x'\}\times[t,t+1]$ and $\{y'\}\times[t,t+1]$ carry $1$. Over points of $\mathcal{Q}'\setminus\partial\mathcal{Q}'$ this holds by Step~2, because $t\le\beta_d-2$; over points of $\partial\mathcal{Q}'$ it holds by Step~1. The holonomy based at $(x',t)$ is therefore
\begin{equation*}
V^{\downarrow}_{e'}\cdot 1\cdot(V^{\downarrow}_{e'})^{-1}\cdot 1=1 .
\end{equation*}

\smallskip
\noindent\textit{Step 4 (the top layer).}
For $z'\in\mathcal{Q}'$ put $o(z'):=(\alpha_1,z'_2,\dots,z'_{d-1})\in\partial\mathcal{Q}'$, and let $\gamma_{z'}$ be the straight path from $o(z')$ to $z'$ in direction $e_1$; it consists of $z'_1-\alpha_1\le s_1$ edges. Let $T^{\downarrow}(z')$ be the ordered product of the bottom values $V^{\downarrow}$ along $\gamma_{z'}$, and $T^{\uparrow}(z')$ that of the lid values $V^{\uparrow}$. Put
\begin{equation*}
\hat\eta(z'):=T^{\downarrow}(z')^{-1}\,T^{\uparrow}(z'),\qquad
\eta(z'):=\hat\eta(z')\quad\text{for } z'\in\mathcal{Q}'\setminus\partial\mathcal{Q}' .
\end{equation*}

Let $e'=[x',y']\not\subset\partial\mathcal{Q}'$ with $y'=x'+e_i$, $1\le i\le d-1$. Consider the top plaquette $e'\times[\beta_d-1,\beta_d]$. Its lower edge lies at height $\beta_d-1>\alpha_d$, because $s_d\ge 2$, and so carries $V^{\downarrow}_{e'}$ by Step~2. Its upper edge is the lid edge, carrying $V^{\uparrow}_{e'}$. Its vertical edges carry $\eta(x')$ and $\eta(y')$. The holonomy based at $(x',\beta_d-1)$ and running along $e'$ first is
\begin{equation*}
\operatorname{hol}=V^{\downarrow}_{e'}\,\eta(y')\,(V^{\uparrow}_{e'})^{-1}\,\eta(x')^{-1} .
\end{equation*}

Now $\eta$ and $\hat\eta$ differ only at points of $\partial\mathcal{Q}'$, where $\eta=1$. By \eqref{8.10:eq-box-filling-2} and the first identity of \eqref{8.10:eq-box-filling-1}, replacing $\eta$ by $\hat\eta$ changes $\operatorname{hol}$ in norm by at most
\begin{equation}\label{8.10:eq-box-filling-4}
\mathbf{1}_{x'\in\partial\mathcal{Q}'}\,\|\hat\eta(x')-1\|+\mathbf{1}_{y'\in\partial\mathcal{Q}'}\,\|\hat\eta(y')-1\| .
\end{equation}

Consider the main term
\begin{equation*}
\Xi:=V^{\downarrow}_{e'}\,\hat\eta(y')\,(V^{\uparrow}_{e'})^{-1}\,\hat\eta(x')^{-1} .
\end{equation*}
Insert $\hat\eta(y')=T^{\downarrow}(y')^{-1}T^{\uparrow}(y')$ and $\hat\eta(x')^{-1}=T^{\uparrow}(x')^{-1}T^{\downarrow}(x')$, and put
\begin{equation*}
\Phi^{\downarrow}:=T^{\downarrow}(x')\,V^{\downarrow}_{e'}\,T^{\downarrow}(y')^{-1},\qquad
\Phi^{\uparrow}:=T^{\uparrow}(x')\,V^{\uparrow}_{e'}\,T^{\uparrow}(y')^{-1} .
\end{equation*}
Then
\begin{equation*}
\begin{aligned}
T^{\downarrow}(x')\,\Xi\,T^{\downarrow}(x')^{-1}
&=T^{\downarrow}(x')V^{\downarrow}_{e'}T^{\downarrow}(y')^{-1}\,
T^{\uparrow}(y')(V^{\uparrow}_{e'})^{-1}T^{\uparrow}(x')^{-1}\\
&=\Phi^{\downarrow}\,(\Phi^{\uparrow})^{-1} .
\end{aligned}
\end{equation*}
By the third and second identities of \eqref{8.10:eq-box-filling-1},
\begin{equation*}
\|\Xi-1\|=\|\Phi^{\downarrow}(\Phi^{\uparrow})^{-1}-1\|=\|\Phi^{\downarrow}-\Phi^{\uparrow}\| .
\end{equation*}

If $i=1$, then $o(y')=o(x')$ and $\gamma_{y'}$ is $\gamma_{x'}$ followed by $e'$. Hence $T^{\downarrow}(y')=T^{\downarrow}(x')V^{\downarrow}_{e'}$ and $T^{\uparrow}(y')=T^{\uparrow}(x')V^{\uparrow}_{e'}$, so $\Phi^{\downarrow}=\Phi^{\uparrow}=1$ and $\Xi=1$.

Let $2\le i\le d-1$. Let $\bar e$ be the edge from $o(x')$ to $o(y')=o(x')+e_i$; it lies in the face $\{x_1=\alpha_1\}$ of $\mathcal{Q}'$, so $\bar e\subset\partial\mathcal{Q}'$. We apply Lemma~\ref{8.10:lem-ladder-identity} in the bottom face $\mathcal{Q}'\times\{\alpha_d\}$, with the following data:
\begin{itemize}
\item the path there called $P$ is the bottom copy of $\gamma_{x'}$;
\item the translation vector is $e_i$, so the translated path is the bottom copy of $\gamma_{y'}$ (both are straight in direction $e_1$, with $n:=x'_1-\alpha_1=y'_1-\alpha_1$ edges);
\item the first rung is the bottom copy of $\bar e$ and the last rung is the bottom copy of $e'$.
\end{itemize}
All edges involved are bottom edges carrying the values $V^{\downarrow}$. By that lemma, $\Phi^{\downarrow}(V^{\downarrow}_{\bar e})^{-1}$ is a product of conjugates of the holonomies, or their inverses, of the $n\le s_1$ bottom plaquettes between the bottom copies of $\gamma_{x'}$ and $\gamma_{y'}$. These plaquettes are contained in $\mathcal{Q}'\times\{\alpha_d\}$, because $y'\in\mathcal{Q}'$, and so they lie in $\partial\mathcal{Q}$. Hence
\begin{equation*}
\|\Phi^{\downarrow}(V^{\downarrow}_{\bar e})^{-1}-1\|\le s_1\delta,
\end{equation*}
and by the second identity of \eqref{8.10:eq-box-filling-1}, $\|\Phi^{\downarrow}-V^{\downarrow}_{\bar e}\|\le s_1\delta$.

The same argument in the lid $\mathcal{Q}'\times\{\beta_d\}$ gives $\|\Phi^{\uparrow}-V^{\uparrow}_{\bar e}\|\le s_1\delta$. Since $\bar e\subset\partial\mathcal{Q}'$, \eqref{8.10:eq-box-filling-3} at $t=\beta_d$ gives
\begin{equation*}
\|V^{\downarrow}_{\bar e}-V^{\uparrow}_{\bar e}\|=\|\Gamma_{\alpha_d}(\bar e)-\Gamma_{\beta_d}(\bar e)\|\le s_d\delta .
\end{equation*}
By the triangle inequality, in both cases
\begin{equation}\label{8.10:eq-box-filling-5}
\|\Xi-1\|\le(2s_1+s_d)\,\delta .
\end{equation}

\smallskip
\noindent\textit{Rim corrections.}
It remains to bound the terms \eqref{8.10:eq-box-filling-4}, that is, $\|\hat\eta(z')-1\|$ for $z'\in\partial\mathcal{Q}'$. By the second identity of \eqref{8.10:eq-box-filling-1},
\begin{equation*}
\|\hat\eta(z')-1\|=\|T^{\downarrow}(z')^{-1}T^{\uparrow}(z')-1\|=\|T^{\uparrow}(z')-T^{\downarrow}(z')\| .
\end{equation*}
Write $z'=(z'_1,z'')$ with $z''=(z_2,\dots,z_{d-1})\in\mathcal{Q}'':=\prod_{2\le i\le d-1}[\alpha_i,\beta_i]$; since $d\ge 3$, $\mathcal{Q}''$ is a box of dimension $d-2\ge 1$. The edges of $\gamma_{z'}$ point in direction $e_1$ and have last coordinates $z''$. No face $\{x_1=c\}$ contains such an edge, and a face $\{x_i=c\}$ with $2\le i\le d-1$ contains it if and only if $z_i=c$. Hence two situations can occur.

If every edge of $\gamma_{z'}$ lies in $\partial\mathcal{Q}'$, then edgewise, by \eqref{8.10:eq-box-filling-3},
\begin{equation*}
\|V^{\downarrow}_e-V^{\uparrow}_e\|=\|\Gamma_{\alpha_d}(e)-\Gamma_{\beta_d}(e)\|\le s_d\delta .
\end{equation*}
By \eqref{8.10:eq-box-filling-2}, applied to at most $s_1$ edges, $\|\hat\eta(z')-1\|\le s_1s_d\delta$.

Otherwise $\gamma_{z'}$ is non-empty and $z''$ lies in the interior of $\mathcal{Q}''$. Since $z'\in\partial\mathcal{Q}'$, this forces $z'_1\in\{\alpha_1,\beta_1\}$, and $z'_1=\alpha_1$ would make $\gamma_{z'}$ empty. Conversely, if $z'_1=\beta_1$ and $z''$ is interior to $\mathcal{Q}''$, then no edge of $\gamma_{z'}$ lies in $\partial\mathcal{Q}'$. So this situation occurs exactly when $z'$ lies on the face $\{x_1=\beta_1\}$ with $z''$ interior to $\mathcal{Q}''$; then $\gamma_{z'}$ crosses $\mathcal{Q}'$.

For every $\zeta''\in\mathcal{Q}''$ define the belt loop $\mathfrak{l}(\zeta'')$ as the concatenation of four paths:
\begin{itemize}
\item the bottom copy of the straight path from $(\alpha_1,\zeta'')$ to $(\beta_1,\zeta'')$;
\item the vertical column at $(\beta_1,\zeta'')$, from height $\alpha_d$ to height $\beta_d$;
\item the lid copy of that straight path, traversed backwards;
\item the vertical column at $(\alpha_1,\zeta'')$, traversed downwards.
\end{itemize}
This loop lies in $\partial\mathcal{Q}$: its four pieces lie in the bottom, in the face $\{x_1=\beta_1\}$, in the lid and in the face $\{x_1=\alpha_1\}$. Both columns lie over points of $\partial\mathcal{Q}'$, so their edges equal $1$ by Step~1. For $z'=(\beta_1,z'')$ the straight path is $\gamma_{z'}$, so
\begin{equation*}
T(\mathfrak{l}(z''))=T^{\downarrow}(z')\,T^{\uparrow}(z')^{-1} .
\end{equation*}
By the second identity of \eqref{8.10:eq-box-filling-1},
\begin{equation*}
\|\hat\eta(z')-1\|=\|T^{\downarrow}(z')-T^{\uparrow}(z')\|=\|T(\mathfrak{l}(z''))-1\| .
\end{equation*}

Let $\zeta''=(\zeta_2,\zeta_3,\dots,\zeta_{d-1})\in\mathcal{Q}''$ with $\zeta_2>\alpha_2$. The translate of $\mathfrak{l}(\zeta'')$ by $-e_2$ is $\mathfrak{l}(\zeta''-e_2)$. We apply the loop case of Lemma~\ref{8.10:lem-ladder-identity} with the loop $\mathfrak{l}(\zeta'')$ and translation vector $-e_2$. The rung plaquettes are the plaquettes between each of the $2s_1+2s_d$ edges of $\mathfrak{l}(\zeta'')$ and its translate:
\begin{itemize}
\item for the bottom edges, bottom plaquettes;
\item for the column at $x_1=\beta_1$, plaquettes of the face $\{x_1=\beta_1\}$;
\item for the lid edges, lid plaquettes;
\item for the column at $x_1=\alpha_1$, plaquettes of the face $\{x_1=\alpha_1\}$.
\end{itemize}
Since $\zeta_2-1\ge\alpha_2$, all of them are contained in $\mathcal{Q}$ and lie in $\partial\mathcal{Q}$, so each has $F_\sigma\le\delta$. The lemma gives
\begin{equation*}
\|T(\mathfrak{l}(\zeta''))-1\|\le(2s_1+2s_d)\,\delta+\|T(\mathfrak{l}(\zeta''-e_2))-1\| .
\end{equation*}
Starting from $\zeta''=z''$ and applying this bound $z_2-\alpha_2\le s_2$ times, we get
\begin{equation*}
\|T(\mathfrak{l}(z''))-1\|\le s_2(2s_1+2s_d)\,\delta+\|T(\mathfrak{l}(\alpha_2,z_3,\dots,z_{d-1}))-1\| .
\end{equation*}

The loop $\mathfrak{l}(\alpha_2,z_3,\dots,z_{d-1})$ moves only in the coordinates $1$ and $d$. Since $d\ge 3$, the coordinate $2$ is distinct from both, so every vertex of this loop has $x_2=\alpha_2$; that is, the loop lies in the face $\{x_2=\alpha_2\}$ of $\mathcal{Q}$. Put $y:=(\alpha_1,\alpha_2,z_3,\dots,z_{d-1},\alpha_d)$. This loop is the boundary loop from $y$ of the rectangle of Lemma~\ref{8.10:lem-stokes-rectangle} with $\mu=1$, $\nu=d$, $A=s_1$ and $B=s_d$. Indeed, it consists of $s_1$ steps $+e_1$ in the bottom, $s_d$ steps $+e_d$ up the column at $x_1=\beta_1$, $s_1$ steps $-e_1$ in the lid, and $s_d$ steps $-e_d$ down the column at $x_1=\alpha_1$. The $s_1s_d$ plaquettes of that rectangle are contained in the face $\{x_2=\alpha_2\}$, so each has $F_p\le\delta$. By Lemma~\ref{8.10:lem-stokes-rectangle},
\begin{equation*}
\|T(\mathfrak{l}(\alpha_2,z_3,\dots,z_{d-1}))-1\|\le s_1s_d\,\delta .
\end{equation*}
Altogether, in this situation,
\begin{equation*}
\|\hat\eta(z')-1\|\le\bigl(2s_2(s_1+s_d)+s_1s_d\bigr)\delta\le 5\bar s^{2}\delta .
\end{equation*}
In the first situation we found $\|\hat\eta(z')-1\|\le s_1s_d\delta\le 5\bar s^2\delta$ as well, so
\begin{equation}\label{8.10:eq-box-filling-6}
\|\hat\eta(z')-1\|\le 5\bar s^{2}\,\delta\qquad(z'\in\partial\mathcal{Q}') .
\end{equation}

\smallskip
\noindent\textit{Conclusion.}
Let $e'\not\subset\partial\mathcal{Q}'$, and consider the top plaquette $e'\times[\beta_d-1,\beta_d]$ of type (b). Its holonomy differs from $\Xi$ by at most \eqref{8.10:eq-box-filling-4}. Combining this with \eqref{8.10:eq-box-filling-5} and \eqref{8.10:eq-box-filling-6}, the holonomy is within
\begin{equation*}
(2s_1+s_d)\,\delta+2\cdot 5\bar s^{2}\delta\le(3\bar s+10\bar s^{2})\,\delta\le 13\bar s^{2}\delta
\end{equation*}
of $1$; the last inequality holds because $\bar s\ge 2$ gives $3\bar s\le 3\bar s^{2}$.

Together with Step~3, where plaquettes of type (a) are within $5\bar s\delta\le 13\bar s^2\delta$ of $1$ and plaquettes of type (b) below the top layer have holonomy $1$, every plaquette of $\mathcal{Q}$ not lying in $\partial\mathcal{Q}$ has $F_\sigma\le 13\bar s^{2}\delta$. Those lying in $\partial\mathcal{Q}$ have $F_\sigma\le\delta\le 13\bar s^{2}\delta$.

Finally we undo the column gauge. As shown in Step~1, applying $\lambda^{-1}$ to the configuration just constructed gives an extension of the original data with the same values of $F_\sigma$.
\end{proof}

We use the objects attached to a touch-connected union of cells $C$ in the definitions and lemmas of this section:
\begin{itemize}
\item $\mathsf{Ex}$ is the collar, with plaquette set $\mathsf{Ex}_2$.
\item $\mathsf{A}$ is the cell complex of the collar region. Its plaquettes belong to $\mathsf{Ex}_2$, and its
vertex set $\mathsf{A}_0$ carries the gauge maps of Lemmas~\ref{8.10:lem-gauge-filling}
and~\ref{8.10:lem-collar-gauge}. $\mathsf{N}$ is the neighbourhood complex of $C$ on whose exterior edges
$V^{\mathrm{ext}}$ is given, with vertex set $\mathsf{N}_0$.
\item $\mathsf{R}_0,\mathsf{R}_1,\mathsf{R}_2$ are the interior vertices, edges and plaquettes.
\item $\mathrm{Str}$ is the set of straddling plaquettes.
\item $\mathcal{V}^{\mathrm{ext}}$ and $\mathcal{V}^{\mathrm{in}}$ are the exterior and interior configurations.
\item $F_P(V)$ is the curvature functional of the plaquette $P$ in the configuration $V$. It is the norm of
the difference between the holonomy of $P$ and $1$. $F_P(V^{\mathrm{in}},V^{\mathrm{ext}})$ is the same
quantity for the configuration equal to $V^{\mathrm{in}}$ on interior edges and to $V^{\mathrm{ext}}$ on
exterior edges.
\item $\varepsilon_\mu$, $G^{\mathrm{mar}}_\mu$, $\mathsf{E}$, $E^{\mu}_{\mathrm{fill}}$ and
$\mathrm{Fill}_{\mu,\mathfrak{E}}(C)$ are as in Definition~\ref{8.10:def-good-set}. So are
$\mathsf{F}=\mathsf{R}_2\sqcup\mathrm{Str}$, the plaquette set on which a margin-$\mu$ filling is
required to lie within the margin, and $\mathsf{Ex}_2^{\mathrm{near}}\subset\mathsf{Ex}_2$, the set of
near collar plaquettes entering the definition of $G^{\mathrm{mar}}_\mu$.
\end{itemize}

The index $\mu$ always denotes the margin. Coordinate directions are indexed by $i$, $j$, $k$, and $e_i$
denotes the unit lattice vector in direction $i$.

\begin{theorem}[Quantitative filling of unlinked and box hulls from a nearly flat collar]
\label{8.10:thm-quantitative-filling}
Let $C$ be a touch-connected union of cells. Let $p_1\ge 1$ and $p_2\ge 1$ be integers such that the
patch condition (the hypothesis of that name in Lemma~\ref{8.10:lem-collar-gauge}) holds for the hull
$\mathrm{hull}_{p_1+p_2+1}(C)$ of $C$ of width $p_1+p_2+1$. Let $\hat{\mathfrak{a}}\ge 1$ be such
that one of the following holds.
\begin{itemize}
\item[($\alpha$)] Condition (U$_1$) of Definition~\ref{8.10:def-linked-components} holds, and
$\hat{\mathfrak{a}}\ge 3\,\mathfrak{a}(\mathsf{Ex})$, where $\mathfrak{a}(\mathsf{Ex})$ is the number
defined in Lemma~\ref{8.10:lem-collar-gauge}.
\item[($\beta$)] $C$ is box-like in the sense of Definition~\ref{8.10:def-linked-components}, and
$\hat{\mathfrak{a}}$ satisfies the bound displayed below.
\end{itemize}
In case ($\beta$) the bound on $\hat{\mathfrak{a}}$ and the meaning of $\bar s$ are
\begin{equation*}
\hat{\mathfrak{a}}\ \ge\ 2+65(d-1)(\bar s+2)^3,\qquad
\bar s:=\ell\cdot\max_{i}\,(b_i-a_i+1).
\end{equation*}
Here $\bar s$ is the largest lattice side of the site box
\begin{equation*}
\mathsf{Q}:=\bar S(R_{\mathrm{int}})=\prod_{i=1}^{d}\{\ell a_i,\dots,\ell b_i+\ell\}.
\end{equation*}
In this display $\ell$ is the lattice side of a cell, and $R_{\mathrm{int}}$ is the box of interior cells of
the box-like union $C$, whose cell indices range over $\prod_{i=1}^{d}\{a_i,\dots,b_i\}$. The map
$\bar S(\cdot)$ assigns to a union of cells its closed site box, as in
Definition~\ref{8.10:def-linked-components}.
Let $\delta\ge 0$ and $V^{\mathrm{ext}}\in\mathcal{V}^{\mathrm{ext}}$ satisfy
\begin{equation}\label{8.10:eq-quantitative-filling-1}
F_{p'}(V^{\mathrm{ext}})\le\delta\qquad\text{for every }p'\in\mathsf{Ex}_2 .
\end{equation}
Then there is $V^{\mathrm{in}}\in\mathcal{V}^{\mathrm{in}}$ with
\begin{equation}\label{8.10:eq-quantitative-filling-2}
F_P(V^{\mathrm{in}},V^{\mathrm{ext}})=0\quad(P\in\mathsf{R}_2),\qquad
F_P(V^{\mathrm{in}},V^{\mathrm{ext}})\le\hat{\mathfrak{a}}\,\delta\quad(P\in\mathrm{Str}).
\end{equation}
Suppose moreover that
\begin{equation}\label{8.10:eq-quantitative-filling-3}
\hat{\mathfrak{a}}\,\delta\le\varepsilon_\mu
\end{equation}
and
\begin{equation}\label{8.10:eq-quantitative-filling-4}
\hat{\mathfrak{a}}^2\delta^2\cdot|\mathrm{Str}|\le\mathfrak{E}.
\end{equation}
Then the following hold:
\begin{itemize}
\item $V^{\mathrm{in}}$ is a margin-$\mu$ filling;
\item $V^{\mathrm{ext}}\in G^{\mathrm{mar}}_\mu$;
\item $\mathsf{E}(V^{\mathrm{in}},V^{\mathrm{ext}})\le\mathfrak{E}$;
\item $V^{\mathrm{ext}}\in\mathrm{Fill}_{\mu,\mathfrak{E}}(C)$, with
\begin{equation*}
E^{\mu}_{\mathrm{fill}}(V^{\mathrm{ext}})\le\hat{\mathfrak{a}}^2\delta^2|\mathrm{Str}|.
\end{equation*}
\end{itemize}
\end{theorem}

\begin{proof}
For a map $u$ on a set of vertices into the gauge group and an edge $(y,z)$ with both endpoints in that
set, we call
\begin{equation*}
\|u(y)V_{(y,z)}u(z)^{-1}-1\|
\end{equation*}
the \emph{defect} of $(y,z)$ under $u$.

We use two elementary facts about unitaries $X,Y,Z$.

First, $\|X^{-1}-1\|=\|X^{-1}(1-X)\|=\|1-X\|$. Since $V_{(z,y)}=V_{(y,z)}^{-1}$, it follows that an edge
and its reverse have the same defect.

Second, from the identity
\begin{equation*}
XYZ-1=(X-1)YZ+(Y-1)Z+(Z-1)
\end{equation*}
we get
\begin{equation}\label{8.10:eq-quantitative-filling-5}
\|XYZ-1\|\le\|X-1\|+\|Y-1\|+\|Z-1\|.
\end{equation}

\emph{Case ($\alpha$).} Condition \eqref{8.10:eq-quantitative-filling-1} states that $V^{\mathrm{ext}}$
is $\delta$-flat on $\mathsf{Ex}$. The patch condition of Lemma~\ref{8.10:lem-collar-gauge} and
condition (U$_1$) hold by assumption. Hence
Lemma~\ref{8.10:lem-collar-gauge} gives a map $u$ on $\mathsf{A}_0$ under which every edge of
$\mathsf{Ex}$ has defect at most $\mathfrak{a}(\mathsf{Ex})\,\delta$.

By part~(b) of the lemma of this chapter describing the exterior edges of $\mathsf{N}$, this applies in
particular to every exterior edge of $\mathsf{N}$. Lemma~\ref{8.10:lem-gauge-filling}, applied with
$\eta:=\mathfrak{a}(\mathsf{Ex})\,\delta$, gives $V^{\mathrm{in}}\in\mathcal{V}^{\mathrm{in}}$ with
$F_P(V^{\mathrm{in}},V^{\mathrm{ext}})=0$ for $P\in\mathsf{R}_2$ and
\begin{equation*}
F_P(V^{\mathrm{in}},V^{\mathrm{ext}})\le 3\,\mathfrak{a}(\mathsf{Ex})\,\delta\le\hat{\mathfrak{a}}\,\delta
\qquad(P\in\mathrm{Str}),
\end{equation*}
the last inequality because $\hat{\mathfrak{a}}\ge 3\,\mathfrak{a}(\mathsf{Ex})$. This is
\eqref{8.10:eq-quantitative-filling-2}.

\emph{Case ($\beta$): preliminaries.} Let $\mathsf{Q}'$ be the site box $\mathsf{Q}$ enlarged by one
lattice layer:
\begin{equation*}
\mathsf{Q}':=\prod_{i=1}^{d}\{\ell a_i-1,\dots,\ell b_i+\ell+1\}.
\end{equation*}
We have $\mathsf{Q}'\subset\mathsf{N}_0$, since $p_2\ell\ge 1$. The lattice sides of $\mathsf{Q}'$ are
\begin{equation*}
s'_i=\ell(b_i-a_i+1)+2\in[4,\bar s+2].
\end{equation*}
We also write
\begin{equation*}
\mathsf{Q}^\circ:=\prod_{i=1}^{d}\{\ell a_i+1,\dots,\ell b_i+\ell-1\}.
\end{equation*}

Membership in a box is decided coordinatewise. Consider a plaquette with vertices $x$, $x+e_j$, $x+e_k$,
$x+e_j+e_k$.
\begin{itemize}
\item If two diagonal vertices $x$ and $x+e_j+e_k$ lie in $\mathsf{Q}$, then the $j$-th and $k$-th
coordinates of both lie in the ranges of $\mathsf{Q}$. The other coordinates of all four vertices
coincide with those of $x$. Hence all four vertices lie in $\mathsf{Q}$.
\item So a plaquette with two vertices in $\mathsf{Q}$ has two adjacent vertices in $\mathsf{Q}$, that
is, an edge in $K(\mathsf{Q})_1=\mathsf{R}_1$. Here $K(\mathsf{Q})$ denotes the complex of vertices,
edges and plaquettes of the box $\mathsf{Q}$, with skeleta $K(\mathsf{Q})_0$ and $K(\mathsf{Q})_1$, as in
Lemma~\ref{8.10:lem-comb-gauge}.
\end{itemize}
Consequently a plaquette with no edge in $\mathsf{R}_1$ has at most one vertex in $\mathsf{Q}$. That
vertex lies in $\mathsf{Q}\setminus\mathsf{Q}^\circ$: a vertex $t\in\mathsf{Q}^\circ$ has all its
lattice neighbours $t\pm e_j$ in $\mathsf{Q}$, so a plaquette containing $t$ would have an edge in
$\mathsf{R}_1$.

The vertices of $\mathsf{Q}'\setminus\mathsf{Q}^\circ$ lie in $\bar S(R_{\mathrm{col}})$, the closed site
box of the collar cells $R_{\mathrm{col}}$ of the box-like $C$
(Definition~\ref{8.10:def-linked-components}): the closed box of the cell of $R_{\mathrm{col}}$ adjacent
across the relevant face contains them. Let $q$ be a plaquette
of $K(\mathsf{Q}')$ without interior edge. Its vertices lie in $\mathsf{Q}'\setminus\mathsf{Q}^\circ$.
Hence $q$ lies in $\mathsf{A}$, hence in $\mathsf{Ex}_2$, and by \eqref{8.10:eq-quantitative-filling-1}
\begin{equation}\label{8.10:eq-quantitative-filling-6}
F_q(V^{\mathrm{ext}})\le\delta\qquad\text{for every plaquette $q$ of $K(\mathsf{Q}')$ with no edge in
}\mathsf{R}_1 .
\end{equation}

\emph{Step (1): a rough filling of $\mathsf{Q}'$.} An edge lying in $\partial\mathsf{Q}'$ has both
endpoints with a common coordinate extreme for $\mathsf{Q}'$, that is, equal to $\ell a_i-1$ or to
$\ell b_i+\ell+1$. Such a coordinate is outside the range of $\mathsf{Q}$. So both endpoints lie outside
$\mathsf{Q}$ and the edge is exterior. In particular $V^{\mathrm{ext}}$ is defined on every edge lying
in $\partial\mathsf{Q}'$.

The plaquettes lying in $\partial\mathsf{Q}'$ therefore have no interior edge. By
\eqref{8.10:eq-quantitative-filling-6} they satisfy $F_\sigma(V^{\mathrm{ext}})\le\delta$.

All sides of $\mathsf{Q}'$ are at least $4\ge 2$, and its largest side is at most $\bar s+2$. We apply
Lemma~\ref{8.10:lem-box-filling} to the box $\mathsf{Q}'$, with the data $V^{\mathrm{ext}}$ on the edges
lying in $\partial\mathsf{Q}'$. It gives a configuration $W$ on all edges of $K(\mathsf{Q}')$ with two
properties: $W$ equals $V^{\mathrm{ext}}$ on the edges lying in $\partial\mathsf{Q}'$, and every
plaquette $\sigma$ of $K(\mathsf{Q}')$ satisfies $F_\sigma(W)\le 13(\bar s+2)^2\delta$.

\emph{Step (2): a gauge on $\mathsf{Q}'$.} Lemma~\ref{8.10:lem-comb-gauge}, applied to $W$ on
$\mathsf{Q}'$ with $\eta:=13(\bar s+2)^2\delta$ and with largest side at most $\bar s+2$, gives a map $u$
on the vertices of $\mathsf{Q}'$ such that
\begin{equation*}
\|u(x)W_{(x,y)}u(y)^{-1}-1\|\le\eta':=13(d-1)(\bar s+2)^3\delta
\end{equation*}
for every edge $(x,y)$ of $K(\mathsf{Q}')$. On the edges lying in $\partial\mathsf{Q}'$ we have
$W=V^{\mathrm{ext}}$. Hence every such edge has defect at most $\eta'$ under $u$, computed with
$V^{\mathrm{ext}}$.

\emph{Step (3): transfer to the faces of $\mathsf{Q}$.} The vertices of $\mathsf{Q}'$ not in $\mathsf{Q}$
are exactly those with some coordinate extreme for $\mathsf{Q}'$, that is, the vertices of
$\partial\mathsf{Q}'$.

For $t\in\mathsf{Q}\setminus\mathsf{Q}^\circ$ let $n_1(t)$ be the outward normal $\pm e_i$ for the least
$i$ with $t_i\in\{\ell a_i,\ \ell b_i+\ell\}$. The sign is $-$ if $t_i=\ell a_i$ and $+$ if
$t_i=\ell b_i+\ell$. Then $t+n_1(t)$ has a coordinate extreme for $\mathsf{Q}'$, so it is a vertex of
$\partial\mathsf{Q}'$. The edge $(t+n_1(t),t)$ has an endpoint outside $\mathsf{Q}$, so it is exterior.

We put
\begin{equation*}
u'(t):=u(t+n_1(t))\,V^{\mathrm{ext}}_{(t+n_1(t),t)}\qquad(t\in\mathsf{Q}\setminus\mathsf{Q}^\circ),
\end{equation*}
and $u':=u$ on the vertices of $\partial\mathsf{Q}'$. At the remaining vertices of $\mathsf{A}_0$, if
any, $u'$ is given the value $1$. These values do not enter any of the defects estimated below. Only the
values of $u'$ at vertices of straddling plaquettes enter the bounds on $\mathrm{Str}$ obtained below from
Lemma~\ref{8.10:lem-gauge-filling}, and these vertices lie in
$(\mathsf{Q}\setminus\mathsf{Q}^\circ)\cup\partial\mathsf{Q}'$ (Step~(4)). We write
$V$ for $V^{\mathrm{ext}}$ in this step.

Let $t\in\mathsf{Q}\setminus\mathsf{Q}^\circ$ and let $n$ be an outward normal at $t$, that is, the
coordinate of $t$ in the direction of $n$ is extreme for $\mathsf{Q}$ on the side of $n$. Write
$n_1:=n_1(t)$.

If $n=n_1$, then, since $u'(t+n_1)=u(t+n_1)$,
\begin{equation*}
u'(t)V_{(t,t+n_1)}u'(t+n_1)^{-1}
=u(t+n_1)V_{(t+n_1,t)}V_{(t,t+n_1)}u(t+n_1)^{-1}=1 .
\end{equation*}

Otherwise $n$ and $n_1$ are in different coordinate directions. A coordinate of $t$ cannot be extreme on
both sides, because the sides of $\mathsf{Q}$ are positive. Let $q$ be the plaquette with vertices $t$,
$t+n_1$, $t+n$, $t+n_1+n$. Each of $t+n_1$, $t+n$, $t+n_1+n$ has a coordinate outside the range of
$\mathsf{Q}$. So $q$ has only the vertex $t$ in $\mathsf{Q}$, and in particular no edge in
$\mathsf{R}_1$. All four vertices lie in $\mathsf{Q}'$. By the paragraph preceding
\eqref{8.10:eq-quantitative-filling-6}, $q\in\mathsf{Ex}_2$, and by
\eqref{8.10:eq-quantitative-filling-6}, $F_q(V)\le\delta$.

Write the holonomy of $q$ from $t+n_1$ as
\begin{equation*}
H:=V_{(t+n_1,t)}V_{(t,t+n)}V_{(t+n,t+n+n_1)}V_{(t+n+n_1,t+n_1)} .
\end{equation*}
Using $V_{(z,y)}=V_{(y,z)}^{-1}$, we have
\begin{equation*}
H\,V_{(t+n_1,t+n_1+n)}\,V_{(t+n_1+n,t+n)}=V_{(t+n_1,t)}V_{(t,t+n)} .
\end{equation*}
Since $u'(t+n)=u(t+n)$, this gives
\begin{align*}
u'(t)V_{(t,t+n)}u'(t+n)^{-1}
&=\bigl[u(t+n_1)Hu(t+n_1)^{-1}\bigr]\cdot
\bigl[u(t+n_1)V_{(t+n_1,t+n_1+n)}u(t+n_1+n)^{-1}\bigr]\\
&\qquad\cdot\bigl[u(t+n_1+n)V_{(t+n_1+n,t+n)}u(t+n)^{-1}\bigr].
\end{align*}
This is a product of three unitaries.
\begin{itemize}
\item The first lies within $\delta$ of $1$, because $F_q(V)\le\delta$ and conjugation by a unitary
preserves the norm.
\item The edge $(t+n_1,t+n_1+n)$ lies in $\partial\mathsf{Q}'$: both endpoints share the coordinate of
$t+n_1$ in the direction of $n_1$, which is extreme for $\mathsf{Q}'$. So the second factor lies within
$\eta'$ of $1$ by Step~(2).
\item The edge $(t+n_1+n,t+n)$ lies in $\partial\mathsf{Q}'$: both endpoints share the coordinate of
$t+n$ in the direction of $n$, which is extreme for $\mathsf{Q}'$. So the third factor also lies within
$\eta'$ of $1$ by Step~(2).
\end{itemize}
By \eqref{8.10:eq-quantitative-filling-5} the defect of $(t,t+n)$ under $u'$ is at most $\delta+2\eta'$.

Together with the case $n=n_1$ we obtain
\begin{equation}\label{8.10:eq-quantitative-filling-7}
\|u'(t)V^{\mathrm{ext}}_{(t,t+n)}u'(t+n)^{-1}-1\|\le\delta+2\eta'
\end{equation}
for every $t\in\mathsf{Q}\setminus\mathsf{Q}^\circ$ and every outward normal $n$ at $t$.

\emph{Step (4): the filling.} Let $P$ be a straddling plaquette. It has exactly one edge $(s,t)$ in
$\mathsf{Q}$. This edge lies in a face of $\mathsf{Q}$ whose outward normal $n$ is the direction in which
$P$ leaves $\mathsf{Q}$. The exterior run of $P$ is
\begin{equation*}
t\to t+n\to s+n\to s,
\end{equation*}
and $s,t\in\mathsf{Q}\setminus\mathsf{Q}^\circ$ with $n$ an outward normal at both. Under $u'$ the three
exterior edges have the following defects.
\begin{itemize}
\item The edge $(t,t+n)$ has defect at most $\delta+2\eta'$, by \eqref{8.10:eq-quantitative-filling-7}.
\item The edge $(t+n,s+n)$ lies in $\partial\mathsf{Q}'$, and $u'=u$ at its endpoints. So its defect is
at most $\eta'$ by Step~(2).
\item The edge $(s+n,s)$ is the reverse of $(s,s+n)$, so it has the same defect. That defect is at most
$\delta+2\eta'$, by \eqref{8.10:eq-quantitative-filling-7}.
\end{itemize}

We apply Lemma~\ref{8.10:lem-gauge-filling} with the map $u'$ on $\mathsf{A}_0$. Its hypothesis is needed
only on the exterior edges of straddling plaquettes, with the respective bounds just listed, and its
proof gives $F_P$ at most the sum of the three defects of the exterior run of $P$. It yields
$V^{\mathrm{in}}\in\mathcal{V}^{\mathrm{in}}$ with $F_P(V^{\mathrm{in}},V^{\mathrm{ext}})=0$ for
$P\in\mathsf{R}_2$ and, for $P\in\mathrm{Str}$,
\begin{equation*}
F_P(V^{\mathrm{in}},V^{\mathrm{ext}})
\le 2\delta+5\eta'
=\bigl(2+65(d-1)(\bar s+2)^3\bigr)\delta
\le\hat{\mathfrak{a}}\,\delta .
\end{equation*}
This is \eqref{8.10:eq-quantitative-filling-2} in case ($\beta$).

\emph{The remaining assertions.} Let $V^{\mathrm{in}}$ be as in \eqref{8.10:eq-quantitative-filling-2},
in either case, and assume \eqref{8.10:eq-quantitative-filling-3} and
\eqref{8.10:eq-quantitative-filling-4}.

Then $F_P(V^{\mathrm{in}},V^{\mathrm{ext}})\le\varepsilon_\mu$ on $\mathsf{F}=\mathsf{R}_2\sqcup\mathrm{Str}$.
So $V^{\mathrm{in}}$ is a margin-$\mu$ filling.

Next, $\mathsf{Ex}_2^{\mathrm{near}}\subset\mathsf{Ex}_2$, and $\hat{\mathfrak{a}}\ge1$ together with
\eqref{8.10:eq-quantitative-filling-3} gives $\delta\le\hat{\mathfrak{a}}\,\delta\le\varepsilon_\mu$.
With \eqref{8.10:eq-quantitative-filling-1} this gives $V^{\mathrm{ext}}\in G^{\mathrm{mar}}_\mu$.

Finally, by \eqref{8.10:eq-quantitative-filling-2} and \eqref{8.10:eq-quantitative-filling-4},
\begin{equation*}
\mathsf{E}(V^{\mathrm{in}},V^{\mathrm{ext}})
=\sum_{P\in\mathrm{Str}}F_P(V^{\mathrm{in}},V^{\mathrm{ext}})^2
\le\hat{\mathfrak{a}}^2\delta^2|\mathrm{Str}|\le\mathfrak{E}.
\end{equation*}
Thus $V^{\mathrm{in}}$ is a margin-$\mu$ filling of $V^{\mathrm{ext}}$ with energy at most
$\hat{\mathfrak{a}}^2\delta^2|\mathrm{Str}|\le\mathfrak{E}$. By the definitions of $E^{\mu}_{\mathrm{fill}}$
and $\mathrm{Fill}_{\mu,\mathfrak{E}}(C)$ in Definition~\ref{8.10:def-good-set}, this gives
$E^{\mu}_{\mathrm{fill}}(V^{\mathrm{ext}})\le\hat{\mathfrak{a}}^2\delta^2|\mathrm{Str}|$ and
$V^{\mathrm{ext}}\in\mathrm{Fill}_{\mu,\mathfrak{E}}(C)$.
\end{proof}

\begin{remark}
The conditions \eqref{8.10:eq-quantitative-filling-1}, \eqref{8.10:eq-quantitative-filling-3} and
\eqref{8.10:eq-quantitative-filling-4} select exterior configurations. They bound none of the parameters
of the construction.

The number $\hat{\mathfrak{a}}$ grows with the extent and the aspect of the hull. In case ($\beta$) this
is explicit. In case ($\alpha$) it happens through $\mathfrak{a}(\mathsf{Ex})$, whose growth is discussed
in a remark below.

The number $\hat{\mathfrak{a}}$ is not used as a threshold of the good set of
Definition~\ref{8.10:def-good-set}.
\end{remark}

\begin{lemma}[Exact flat filling of unlinked hulls, openness and counts]\label{8.10:lem-flat-filling}
Let $C$ be a touch-connected union of cells, let $p_1\ge 1$ and $p_2\ge 1$, and assume that the
patching condition holds for $\operatorname{hull}_{p_1+p_2+1}(C)$. The words \emph{linked},
\emph{unlinked} and \emph{box-like} and condition $(\mathrm{U}_1)$ are those of
Definition~\ref{8.10:def-linked-components}. Margin-$\mu$ fillings, the collar energy
$\mathsf{E}$ and the sets $\mathrm{Fill}_{\mu,\mathfrak{E}}(C)$, $G^{\mathrm{mar}}_\mu$,
$G^{\mathrm{av}}_f$, $G^{\mathrm{en}}_E$ are those of Definition~\ref{8.10:def-good-set}.
The neighbourhood complex $\mathsf{N}$ of $C$, which contains the set $\mathsf{F}$ of plaquettes
to be filled; its interior part $\mathsf{R}$, with vertex, edge and plaquette sets
$\mathsf{R}_0,\mathsf{R}_1,\mathsf{R}_2$; the exterior collar $\mathsf{Ex}$, with plaquette set
$\mathsf{Ex}_2$; the vertex set $\mathsf{A}_0$; the set $\mathrm{Str}$ of straddling plaquettes;
the space $\mathcal{V}^{\mathrm{ext}}$ of exterior fields; the defects $F_P$, the transports
$T_\gamma$ and the groups $\pi_1^\square$ are those fixed for $C$ earlier in this section.
\begin{enumerate}
\item[(i)] \textup{(Exact flat filling.)} If $C$ is unlinked and $V^{\mathrm{ext}}$ is flat on
$\mathsf{Ex}$, there is $V^{\mathrm{in}}$ with $F_P(V^{\mathrm{in}},V^{\mathrm{ext}})=0$ for
every $P\in\mathsf{F}$. So $\mathsf{E}(V^{\mathrm{in}},V^{\mathrm{ext}})=0$, $V^{\mathrm{in}}$ is a
margin-$\mu$ filling for every $\mu$, and $V^{\mathrm{ext}}\in\mathrm{Fill}_{\mu,0}(C)$ whenever
$V^{\mathrm{ext}}\in G^{\mathrm{mar}}_\mu$. This last condition is automatic if $V^{\mathrm{ext}}$
is flat on $\mathsf{Ex}_2^{\mathrm{near}}$.

Consequently an unlinked $C$ is not linked. Indeed, if $V^{\mathrm{ext}}$ is flat on every
plaquette of $\mathbb{T}$ without interior edge, then the configuration
$(V^{\mathrm{in}},V^{\mathrm{ext}})$ is flat on every plaquette of $\mathbb{T}$, and $T_\gamma=1$
for every loop $\gamma$ null-homotopic in $K(\mathbb{T})$.
\item[(ii)] \textup{(Openness.)} For $\eta>0$ the set
\begin{equation}\label{8.10:eq-flat-filling-1}
\mathcal{O}_\eta:=\bigl\{V^{\mathrm{ext}}\in\mathcal{V}^{\mathrm{ext}}:\ \text{some }
V^{\mathrm{in}}\text{ has }F_P(V^{\mathrm{in}},V^{\mathrm{ext}})<\eta\ \text{for all }
P\in\mathsf{F}\bigr\}
\end{equation}
is open in $\mathcal{V}^{\mathrm{ext}}$ and contains $V^{\mathrm{ext}}\equiv 1$. Let
$\mathfrak{E}>0$, $f>0$ and $E>0$, and put
\begin{equation}\label{8.10:eq-flat-filling-2}
\eta:=\min\Bigl(\varepsilon_\mu,\bigl(\tfrac{\mathfrak{E}}{|\mathsf{F}|+1}\bigr)^{1/2}\Bigr),
\qquad
\eta':=\min\Bigl(\varepsilon_\mu,\,f,\,\bigl(\tfrac{E}{|\mathsf{Ex}_2|+1}\bigr)^{1/2}\Bigr).
\end{equation}
Then the set
$\mathrm{Fill}_{\mu,\mathfrak{E}}(C)\cap G^{\mathrm{av}}_f\cap G^{\mathrm{en}}_E$ contains the
non-empty set
\begin{equation}\label{8.10:eq-flat-filling-3}
\mathcal{O}_\eta\cap\bigl\{V^{\mathrm{ext}}:\ F_{p'}(V^{\mathrm{ext}})<\eta'\ \text{for all }
p'\in\mathsf{Ex}_2\bigr\}.
\end{equation}
This set is open in $\mathcal{V}^{\mathrm{ext}}$ and has positive Haar measure, and if $C$ is
unlinked it contains every $V^{\mathrm{ext}}$ flat on $\mathsf{Ex}$.
\item[(iii)] \textup{(Box-like components.)} If $C$ is box-like, then condition
$(\mathrm{U}_1)$ holds.
\item[(iv)] \textup{(Counts.)}
\begin{equation}\label{8.10:eq-flat-filling-4}
\operatorname{lk}(C)\le c_d(2p_1+1)^d\,n(C),\qquad
\operatorname{lk}^\sharp(C)\le|\mathsf{F}|\le c_d(\ell+1)^d(2p_1+3)^d\,n(C).
\end{equation}
Here the numbers $c_d$ and $\ell$ are those of Lemma~\ref{8.10:lem-obstruction-ranks}.
\end{enumerate}
\end{lemma}

\begin{proof}[Proof of \textup{(i)}--\textup{(iv)}]
\emph{Proof of \textup{(i)}.} We say that $V^{\mathrm{in}}$ \emph{fills} $V^{\mathrm{ext}}$ if
$F_P(V^{\mathrm{in}},V^{\mathrm{ext}})=0$ for every $P\in\mathsf{F}$. For a gauge transformation
$u$, write $(\cdot)^u$ for the transformed field.

The filling problem is gauge covariant. If $V^{\mathrm{in}}$ fills $(V^{\mathrm{ext}})^u$, then
$(V^{\mathrm{in}})^{u^{-1}}$ fills $V^{\mathrm{ext}}$. Hence, to fill $V^{\mathrm{ext}}$, it
suffices to fill a gauge transform of it.

Since $C$ is unlinked, condition $(\mathrm{U}_1)$ or condition $(\mathrm{U}_2)$ of
Definition~\ref{8.10:def-linked-components} holds. We treat the two cases separately.

\emph{Case $(\mathrm{U}_1)$.} The hypotheses of Lemma~\ref{8.10:lem-collar-gauge} are the
patching condition, assumed in the statement, and $(\mathrm{U}_1)$. Flatness of
$V^{\mathrm{ext}}$ on $\mathsf{Ex}$ is $\delta$-flatness with $\delta=0$.
Lemma~\ref{8.10:lem-collar-gauge}, with the constant $\mathfrak{a}(\mathsf{Ex})$ defined there,
therefore gives $u$ on $\mathsf{A}_0$ such that
\begin{equation*}
\|u(y)V^{\mathrm{ext}}_e u(z)^{-1}-1\|\le\mathfrak{a}(\mathsf{Ex})\cdot 0=0
\end{equation*}
for every edge $e=(y,z)$ of $\mathsf{Ex}$. So $u$ has zero defect on every edge of
$\mathsf{Ex}$, and in particular on the exterior edges of $\mathsf{N}$.

This is the hypothesis of Lemma~\ref{8.10:lem-gauge-filling} with $\eta=0$. That lemma yields a
field $V^{\mathrm{in}}$ on $\mathsf{R}_1$ with $F_P(V^{\mathrm{in}},V^{\mathrm{ext}})=0$ for every
$P\in\mathsf{R}_2$, and with
$F_P(V^{\mathrm{in}},V^{\mathrm{ext}})\le 3\cdot 0=0$ for every $P\in\mathrm{Str}$. Hence
$F\equiv 0$ on $\mathsf{F}$.

\emph{Case $(\mathrm{U}_2)$.} Let $T_{\mathrm{ex}}$ be a spanning tree of the connected graph of
$\mathsf{Ex}$. Its vertex set is $\mathsf{A}_0$, by part~(b) of the structure lemma for the
neighbourhood complex $\mathsf{N}$. Fix $v\in\mathsf{A}_0$.

We first build a spanning tree of the graph of $\mathsf{N}$. The graph of $\mathsf{R}$ is
connected, because $R_{\mathrm{int}}$ is touch-connected. It meets $\mathsf{A}_0$ in
$\mathsf{I}_0\neq\emptyset$. Hence $T_{\mathrm{ex}}\cup\mathsf{R}_1$ is a connected spanning
subgraph of the graph of $\mathsf{N}$. A tree contained in a connected graph can be enlarged,
using edges of that graph, to a spanning tree of it. We therefore extend $T_{\mathrm{ex}}$ by
interior edges to a spanning tree $T$ of the graph of $\mathsf{N}$. Since only interior edges
are added, an exterior edge of $\mathsf{N}$ that lies in $T$ lies in $T_{\mathrm{ex}}$.

Next we normalise $V^{\mathrm{ext}}$ on $T_{\mathrm{ex}}$. Let $u(y)$, for $y\in\mathsf{A}_0$, be
the transport of $V^{\mathrm{ext}}$ along the path in $T_{\mathrm{ex}}$ from $v$ to $y$, and put
$u:=1$ at the other vertices. After this gauge transformation $V^{\mathrm{ext}}\equiv 1$ on the
edges of $T_{\mathrm{ex}}$. The transformed field is still flat on $\mathsf{Ex}$. By gauge
covariance it suffices to fill it, so from now on we assume $V^{\mathrm{ext}}\equiv 1$ on
$T_{\mathrm{ex}}$.

Define $\Phi:\pi_1^\square(\mathsf{N},v)\to G$ by
\begin{equation*}
\Phi([\gamma]):=T_{\gamma'}(V^{\mathrm{ext}}),
\end{equation*}
where $\gamma'$ is any loop at $v$ in $\mathsf{Ex}$ homotopic in $\mathsf{N}$ to $\gamma$. We
check that this is well defined and a homomorphism.
\begin{itemize}
\item Such a loop $\gamma'$ exists by $(\mathrm{U}_2)$.
\item Let $\gamma_1,\gamma_2$ be loops at $v$ representing the same class, and let $\gamma'_i$ be
a choice for $\gamma_i$. Then $\gamma'_1$, $\gamma_1$, $\gamma_2$ and $\gamma'_2$ are all
homotopic in $\mathsf{N}$. By $(\mathrm{U}_2)$, $\gamma'_1$ and $\gamma'_2$ are then homotopic in
$\mathsf{Ex}$, so $\gamma'_2$ is
obtained from $\gamma'_1$ by moves across plaquettes of $\mathsf{Ex}$. Since
$V^{\mathrm{ext}}$ is flat on $\mathsf{Ex}$, Lemma~\ref{8.10:lem-homotopy-moves} gives
$T_{\gamma'_1}(V^{\mathrm{ext}})=T_{\gamma'_2}(V^{\mathrm{ext}})$. So $\Phi$ is well defined.
\item If $\gamma'_1,\gamma'_2$ represent $[\gamma_1],[\gamma_2]$, then the concatenation
$\gamma'_1\gamma'_2$ is a loop at $v$ in $\mathsf{Ex}$ homotopic in $\mathsf{N}$ to
$\gamma_1\gamma_2$. The transport along a concatenation of loops at $v$ is the ordered product
of the transports. Hence $\Phi$ is a homomorphism.
\end{itemize}

Let $W$ be the field on $\mathsf{N}$ given by Lemma~\ref{8.10:lem-prescribed-holonomy}, applied
with $X=\mathsf{N}$, whose graph is connected, with the spanning tree $T$ and with this $\Phi$.
Thus $W\equiv 1$ on the edges of $T$, and $W_{(y,z)}=\Phi([c_{(y,z)}])$ on the other edges.
Moreover $T_\gamma(W)=\Phi([\gamma])$ for every loop $\gamma$ at $v$, and $W$ is flat on
$\mathsf{N}$.

We show that $W$ agrees with $V^{\mathrm{ext}}$ on every exterior edge $e$ of $\mathsf{N}$. By
the remark after the construction of $T$, either $e\in T_{\mathrm{ex}}$ or $e\notin T$.
\begin{itemize}
\item If $e\in T_{\mathrm{ex}}$, then $W_e=1=V^{\mathrm{ext}}_e$.
\item Let $e=(y,z)\notin T$. The endpoints $y,z$ lie in $\mathsf{A}_0$, the vertex set of
$T_{\mathrm{ex}}$. Since $T\supset T_{\mathrm{ex}}$ and both are trees, the unique path in $T$
from $v$ to $y$ is the path in $T_{\mathrm{ex}}$, and likewise for $z$. Hence $c_e$ is a loop in
$\mathsf{Ex}$. It is therefore an admissible choice of $\gamma'$ for its own class, and
\begin{equation*}
W_e=\Phi([c_e])=T_{c_e}(V^{\mathrm{ext}})=V^{\mathrm{ext}}_e .
\end{equation*}
The last equality holds because every edge of $c_e$ other than $e$ is an edge of
$T_{\mathrm{ex}}$, where $V^{\mathrm{ext}}$ equals $1$.
\end{itemize}
Thus $W$ extends $V^{\mathrm{ext}}$. Put $V^{\mathrm{in}}:=W$ restricted to $\mathsf{R}_1$. Then
$(V^{\mathrm{in}},V^{\mathrm{ext}})$ coincides with $W$ on $\mathsf{N}$ and is flat there. Since
$\mathsf{N}\supset\mathsf{F}$, we get $F_P(V^{\mathrm{in}},V^{\mathrm{ext}})=0$ for every
$P\in\mathsf{F}$.

\emph{Energy, margin and membership.} In both cases we have a $V^{\mathrm{in}}$ with
$F_P(V^{\mathrm{in}},V^{\mathrm{ext}})=0$ for every $P\in\mathsf{F}$. By
Definition~\ref{8.10:def-good-set}, the collar energy $\mathsf{E}(V^{\mathrm{in}},V^{\mathrm{ext}})$
therefore vanishes, and $V^{\mathrm{in}}$ is a margin-$\mu$ filling for every $\mu$. Hence
$V^{\mathrm{ext}}\in\mathrm{Fill}_{\mu,0}(C)$ as soon as $V^{\mathrm{ext}}\in G^{\mathrm{mar}}_\mu$,
and by the definition of $G^{\mathrm{mar}}_\mu$ in Definition~\ref{8.10:def-good-set} this holds in
particular when $V^{\mathrm{ext}}$ is flat on $\mathsf{Ex}_2^{\mathrm{near}}$.

\emph{The consequence.} Suppose $V^{\mathrm{ext}}$ is flat on every plaquette of $\mathbb{T}$
without interior edge. The plaquettes of $\mathsf{Ex}$ have no interior edge, so
$V^{\mathrm{ext}}$ is flat on $\mathsf{Ex}$ and the construction above applies; let
$V^{\mathrm{in}}$ be the resulting field. Every
plaquette of $\mathbb{T}$ either has an interior edge or has none.
\begin{itemize}
\item If it has an interior edge, it lies in $\mathsf{F}$, by part~(a) of the structure lemma
for the neighbourhood complex $\mathsf{N}$. There $F_P(V^{\mathrm{in}},V^{\mathrm{ext}})=0$.
\item If it has none, the configuration on it is $V^{\mathrm{ext}}$, which is flat there by
hypothesis.
\end{itemize}
Hence $(V^{\mathrm{in}},V^{\mathrm{ext}})$ is flat on every plaquette of $\mathbb{T}$.
Lemma~\ref{8.10:lem-homotopy-moves}, applied with $X=K(\mathbb{T})$, then gives $T_\gamma=1$ for
every loop $\gamma$ null-homotopic in $K(\mathbb{T})$.

\emph{Proof of \textup{(ii)}.} Write $G^{\mathsf{R}_1}$ for the space of interior fields
$V^{\mathrm{in}}$, one element of $G$ for each edge of $\mathsf{R}_1$. Each
$F_P(V^{\mathrm{in}},V^{\mathrm{ext}})$ is a continuous function of the finitely many bond
variables on the boundary of $P$, and $\mathsf{F}$ is finite. Hence the set
\begin{equation*}
\mathcal{U}_\eta:=\bigl\{(V^{\mathrm{in}},V^{\mathrm{ext}})\in G^{\mathsf{R}_1}\times
\mathcal{V}^{\mathrm{ext}}:\ F_P(V^{\mathrm{in}},V^{\mathrm{ext}})<\eta\ \text{for all }
P\in\mathsf{F}\bigr\}
\end{equation*}
is a finite intersection of open sets, hence open. The set $\mathcal{O}_\eta$ of
\eqref{8.10:eq-flat-filling-1} is the image of $\mathcal{U}_\eta$ under the projection
$(V^{\mathrm{in}},V^{\mathrm{ext}})\mapsto V^{\mathrm{ext}}$, and a projection of a product onto a
factor is an open map; so $\mathcal{O}_\eta$ is open in $\mathcal{V}^{\mathrm{ext}}$. The
configuration equal to $1$ on every edge is flat, so $F_P(1,1)=0<\eta$ for every $P\in\mathsf{F}$;
taking $V^{\mathrm{in}}\equiv 1$ shows that $V^{\mathrm{ext}}\equiv 1$ lies in $\mathcal{O}_\eta$.

Now let $\eta,\eta'$ be as in \eqref{8.10:eq-flat-filling-2}, and let $V^{\mathrm{ext}}$ lie in
the set \eqref{8.10:eq-flat-filling-3}. Choose $V^{\mathrm{in}}$ with
$F_P(V^{\mathrm{in}},V^{\mathrm{ext}})<\eta$ for all $P\in\mathsf{F}$. Then
$F_P(V^{\mathrm{in}},V^{\mathrm{ext}})<\varepsilon_\mu$ for every $P\in\mathsf{F}$, and
\begin{equation*}
\sum_{P\in\mathsf{F}}F_P(V^{\mathrm{in}},V^{\mathrm{ext}})^2\le|\mathsf{F}|\,\eta^2
\le\frac{|\mathsf{F}|\,\mathfrak{E}}{|\mathsf{F}|+1}<\mathfrak{E}.
\end{equation*}
Likewise $F_{p'}(V^{\mathrm{ext}})<\varepsilon_\mu$ and $F_{p'}(V^{\mathrm{ext}})<f$ for every
$p'\in\mathsf{Ex}_2$, and
\begin{equation*}
\sum_{p'\in\mathsf{Ex}_2}F_{p'}(V^{\mathrm{ext}})^2\le|\mathsf{Ex}_2|\,\eta'^2
\le\frac{|\mathsf{Ex}_2|\,E}{|\mathsf{Ex}_2|+1}<E.
\end{equation*}
By Definition~\ref{8.10:def-good-set}, these inequalities give that $V^{\mathrm{in}}$ is a
margin-$\mu$ filling with collar energy below $\mathfrak{E}$, that
$V^{\mathrm{ext}}\in G^{\mathrm{mar}}_\mu$, and that
$V^{\mathrm{ext}}\in G^{\mathrm{av}}_f\cap G^{\mathrm{en}}_E$; hence
$V^{\mathrm{ext}}\in\mathrm{Fill}_{\mu,\mathfrak{E}}(C)\cap G^{\mathrm{av}}_f\cap G^{\mathrm{en}}_E$.

The second set in the intersection \eqref{8.10:eq-flat-filling-3} is
open, since each $F_{p'}$ is a continuous function of $V^{\mathrm{ext}}$ and $\mathsf{Ex}_2$ is
finite; so \eqref{8.10:eq-flat-filling-3} is open in $\mathcal{V}^{\mathrm{ext}}$. It contains
$V^{\mathrm{ext}}\equiv 1$, which lies in $\mathcal{O}_\eta$ and has $F_{p'}=0<\eta'$; so it is
non-empty. The space $\mathcal{V}^{\mathrm{ext}}$ is a finite product of copies of the compact
group $G$, and its Haar measure is the product of normalised Haar measures. A non-empty open
subset of a compact group has positive Haar measure: finitely many of its translates cover the
group, Haar measure is translation invariant and the group has measure one. Hence
\eqref{8.10:eq-flat-filling-3} has positive Haar measure.

Finally let $C$ be unlinked and $V^{\mathrm{ext}}$ flat on $\mathsf{Ex}$. By part~(i) there is
$V^{\mathrm{in}}$ with $F_P(V^{\mathrm{in}},V^{\mathrm{ext}})=0<\eta$ for every $P\in\mathsf{F}$,
so $V^{\mathrm{ext}}\in\mathcal{O}_\eta$; and $F_{p'}(V^{\mathrm{ext}})=0<\eta'$ for every
plaquette $p'\in\mathsf{Ex}_2$ of $\mathsf{Ex}$. So $V^{\mathrm{ext}}$ lies in
\eqref{8.10:eq-flat-filling-3}.

\emph{Proof of \textup{(iii)}.} Let $C$ be box-like. Condition $(\mathrm{U}_1)$ asks that every
connected component of $\mathsf{Ex}$, isolated vertices included, have trivial $\pi_1^\square$
(Definition~\ref{8.10:def-linked-components}). It follows from three facts: the topological
statement on box-like unions of cells; the fact of homotopy theory
that if a pair $(X,A)$ has the homotopy extension property and the inclusion
$A\hookrightarrow X$ is a homotopy equivalence, then $A$ is a deformation retract of $X$; and
the simple connectivity $\pi_1(S^{d-1})=1$ of the sphere $S^{d-1}$, which holds since $d=4$. The
three facts are combined in the proof of part~\textup{(iii)} given after the present proof.

\emph{Proof of \textup{(iv)}.} The two inequalities \eqref{8.10:eq-flat-filling-4} are part~(a)
of Lemma~\ref{8.10:lem-obstruction-ranks}, whose hypotheses are those of the present lemma.
\end{proof}

\begin{proof}[Proof of part \textup{(iii)} of Lemma~\ref{8.10:lem-flat-filling}]
We use the objects of Lemma~\ref{8.10:lem-flat-filling} and the notation fixed with it. Cells
have side $\ell$ in lattice units, the closed box of the cell $c$ being the site set
$\prod_\nu\{\ell c_\nu,\dots,\ell c_\nu+\ell\}$, and $\bar S(\cdot)$ is the union of the closed
boxes of a set of cells; $\operatorname{hull}_{p}(\cdot)$ is the enlargement of a set of cells
that appears in the hypothesis of that lemma. Further, $\mathsf{R}$ is the complex spanned by the
closed boxes of the cells of the box $R_{\mathrm{int}}$ of part~(iii), with edge set
$\mathsf{R}_1$; $\mathsf{N}$ is the complex, with site set $\mathsf{N}_0$, spanned by the closed
boxes of the cells of $\operatorname{hull}_{p_2}R_{\mathrm{int}}$; the annulus complex
$\mathsf{A}$ is spanned by the set $\mathsf{A}_0$ of sites of the closed boxes of the cells of
$\operatorname{hull}_{p_2}R_{\mathrm{int}}$ outside $R_{\mathrm{int}}$; and the exterior complex
$\mathsf{Ex}$ consists of the cubes of $\mathsf{A}$ having no edge in $\mathsf{R}_1$.
Throughout, coordinate directions are indexed by $\nu\in\{1,\dots,d\}$ and $e_\nu$ is the unit
vector in direction $\nu$. For a set $S$ of sites we write $\operatorname{cx}(S)$ for the cubical
complex spanned by $S$, the complex used in part~(i) of Lemma~\ref{8.10:lem-flat-filling}: all
lattice cubes (of every dimension, side $1$) whose vertices lie in $S$; and we write
$|\cdot|$ for geometric realisation in $\mathbb{R}^d$.

\emph{The boxes.} Write the box of part~(iii) as
$R_{\mathrm{int}}=\prod_{\nu}\{a_\nu,\dots,b_\nu\}$, and put
\begin{equation*}
\begin{gathered}
Q:=\bar S(R_{\mathrm{int}})=\prod_\nu\{\ell a_\nu,\dots,\ell b_\nu+\ell\},\qquad
Q^\circ:=\prod_\nu\{\ell a_\nu+1,\dots,\ell b_\nu+\ell-1\},\\
Q^+:=\prod_\nu\{\ell(a_\nu-p_2),\dots,\ell(b_\nu+p_2)+\ell\}
=\bar S(\operatorname{hull}_{p_2}R_{\mathrm{int}})=\mathsf{N}_0 .
\end{gathered}
\end{equation*}
Then $\mathsf{R}=\operatorname{cx}(Q)$ and $\mathsf{N}=\operatorname{cx}(Q^+)$.

\emph{The annulus.} We have $\mathsf{A}_0=Q^+\setminus Q^\circ$. Indeed, let $x\in Q^\circ$; if
$x$ lies in the closed box $\prod_\nu\{\ell c_\nu,\dots,\ell c_\nu+\ell\}$ of a cell $c$, then
$\ell c_\nu\le x_\nu\le\ell b_\nu+\ell-1$ forces $c_\nu\le b_\nu$ and
$\ell c_\nu+\ell\ge x_\nu\ge\ell a_\nu+1$ forces $c_\nu\ge a_\nu$, so $c\in R_{\mathrm{int}}$: a site
of $Q^\circ$ lies in closed boxes of cells of $R_{\mathrm{int}}$ only. Conversely let
$x\in Q^+\setminus Q^\circ$; then some coordinate has $x_\nu\le\ell a_\nu$ or
$x_\nu\ge\ell b_\nu+\ell$. In the first case
$\ell(a_\nu-p_2)\le x_\nu\le\ell a_\nu$, and the intervals $[\ell c,\ell c+\ell]$ with
$a_\nu-p_2\le c\le a_\nu-1$ cover $[\ell(a_\nu-p_2),\ell a_\nu]$, so one may choose such a
$c_\nu$ with $x_\nu$ in its interval; the second case is symmetric, with
$b_\nu+1\le c_\nu\le b_\nu+p_2$. Choosing the other coordinates of $c$ in the same way within
$\{a_{\nu'}-p_2,\dots,b_{\nu'}+p_2\}$, we obtain a cell $c$ of
$\operatorname{hull}_{p_2}R_{\mathrm{int}}$ outside $R_{\mathrm{int}}$ (a collar cell) whose closed
box contains $x$. Consequently $\mathsf{A}=\operatorname{cx}(Q^+\setminus Q^\circ)$.

\emph{The exterior complex.} Let $\Delta$ be a cube with vertices in $Q^+\setminus Q^\circ$. Since
$Q$ and $\Delta$ are products over the coordinates, $\Delta\cap Q$ is the product of the
intersections of the coordinate factors, each of which is the whole factor, one point, or empty;
so $\Delta\cap Q$ is a face of $\Delta$ or empty. An edge of $\Delta$ lies in $\mathsf{R}_1$, the
edges of $\mathsf{R}=\operatorname{cx}(Q)$, exactly when both its endpoints lie in $Q$; hence
$\Delta$ has an edge in $\mathsf{R}_1$ exactly when $\Delta\cap Q$ has dimension $\ge1$. The cubes
of $\mathsf{Ex}$ are therefore those with $\Delta\cap Q$ empty or a single vertex $v$, and in the
latter case $v\in Q\setminus Q^\circ$. For $v\in Q\setminus Q^\circ$ let $\mathrm{Out}(v)$ be the set
of pairs $(\nu,s_\nu)$ with $s_\nu=-1$ if $v_\nu=\ell a_\nu$ and $s_\nu=+1$ if
$v_\nu=\ell b_\nu+\ell$ (the outward directions at $v$; non-empty since $v\notin Q^\circ$), and
let $\mathrm{Co}(v)$ be the cube
\begin{equation*}
v+\sum_{(\nu,s_\nu)\in\mathrm{Out}(v)}\{0,s_\nu\}e_\nu
\end{equation*}
with all its faces. Every vertex of $\mathrm{Co}(v)$ other than $v$ has a coordinate moved outward
from $Q$, hence lies outside $Q$, and (as $p_2\ell\ge1$) inside $Q^+$; so every face of
$\mathrm{Co}(v)$ meets $Q$ in $\{v\}$ at most. A cube meeting $Q$ in $\{v\}$ alone can extend from
$v$ only in outward directions, so it is a face of $\mathrm{Co}(v)$. Hence
\begin{equation*}
\mathsf{Ex}=\mathrm{Sh}\cup\bigcup_{v\in Q\setminus Q^\circ}\mathrm{Co}(v),
\qquad \mathrm{Sh}:=\operatorname{cx}(Q^+\setminus Q).
\end{equation*}

\emph{Realisation.} Let $|Q^+|_{\mathbb{R}}:=\prod_\nu[\ell(a_\nu-p_2),\ell(b_\nu+p_2)+\ell]$ and
$\mathcal{I}:=\prod_\nu\mathopen{]}\ell a_\nu-1,\ell b_\nu+\ell+1\mathclose{[}$. A closed unit cube
avoids $\mathcal{I}$ if and only if its vertices avoid $Q$ (the cube meets $\mathcal{I}$ iff in every
coordinate one of its two values lies in $\{\ell a_\nu,\dots,\ell b_\nu+\ell\}$); and a point of
$|Q^+|_{\mathbb{R}}$ outside $\mathcal{I}$ lies in the smallest lattice cube containing it, which
avoids $\mathcal{I}$ because the faces of $\mathcal{I}$ lie in integer hyperplanes. Hence
$|\mathrm{Sh}|=|Q^+|_{\mathbb{R}}\setminus\mathcal{I}$. Moreover
$\bar{\mathcal{I}}\subset\operatorname{int}|Q^+|_{\mathbb{R}}$, since $p_2\ell\ge2$.

\emph{Retraction onto $|\mathrm{Sh}|$.} For each $v$, $|\mathrm{Co}(v)|$ meets
$|\mathrm{Sh}|\cup\bigcup_{v'\ne v}|\mathrm{Co}(v')|$ exactly in the union of the faces of the top
cube of $\mathrm{Co}(v)$ not containing $v$. These faces have all vertices outside $Q$, so they lie
in $|\mathrm{Sh}|$, and their union is star-shaped about the far corner
$v+\sum_{(\nu,s_\nu)\in\mathrm{Out}(v)}s_\nu e_\nu$, hence contractible. The pair formed by the
cube and this subcomplex has the homotopy extension property and, both being contractible, its
inclusion is a homotopy equivalence; by the deformation-retraction criterion of
Notation~\ref{8.10:not-topology-facts} the cube deformation retracts onto the union of these
faces. These deformations
are stationary on the
overlaps, which lie in the far faces; extending them by the identity on $|\mathrm{Sh}|$, they
paste over the finite closed cover to a deformation retraction of $|\mathsf{Ex}|$ onto
$|\mathrm{Sh}|$. Finally $\mathcal{I}$ and $|Q^+|_{\mathbb{R}}$ are concentric boxes with
$\bar{\mathcal{I}}$ in the interior, so radial projection from their common centre deformation
retracts $|\mathrm{Sh}|=|Q^+|_{\mathbb{R}}\setminus\mathcal{I}$ onto
$\partial|Q^+|_{\mathbb{R}}\cong S^{d-1}$.

\emph{Conclusion.} Thus $|\mathsf{Ex}|$ is connected and
$\pi_1(|\mathsf{Ex}|)\cong\pi_1(S^{d-1})=1$, as $d\ge3$. By the identification of the edge-loop
group with the fundamental group (Definition~\ref{8.10:def-edge-loop-homotopy}), every loop in
the graph of $\mathsf{Ex}$ is null-homotopic in $\mathsf{Ex}$, which is the conclusion of
part~(iii).
\end{proof}

\begin{lemma}[Flux bound for margined fillings]\label{8.10:lem-flux-bound}
Let $\gamma$ be an edge loop all of whose edges are exterior, and suppose that $\gamma$ admits
a null-homotopy in $K(\mathbb{T})$ across the plaquettes $(q_1,\dots,q_A)$. If
$V^{\mathrm{in}}$ is a margin-$\mu$ filling of $V^{\mathrm{ext}}$ in the sense of
Definition~\ref{8.10:def-good-set}, then
\begin{equation}\label{8.10:eq-flux-bound-1}
\bigl\|T_\gamma(V^{\mathrm{ext}})-1\bigr\|
\;\le\;
\varepsilon_\mu\cdot\#\{\,i : q_i\in\mathsf{F}\,\}
\;+\;\sum_{i\,:\,q_i\notin\mathsf{F}} F_{q_i}(V^{\mathrm{ext}}).
\end{equation}
\end{lemma}

\begin{proof}
Let $V:=(V^{\mathrm{in}},V^{\mathrm{ext}})$ be the configuration that equals
$V^{\mathrm{in}}$ on the interior edges and $V^{\mathrm{ext}}$ on the exterior edges. Every
edge of $\gamma$ is exterior, so the transport $T_\gamma$ reads only the values of the
configuration on exterior edges. Hence
\begin{equation*}
T_\gamma(V)=T_\gamma(V^{\mathrm{ext}}).
\end{equation*}

The loop $\gamma$ admits a null-homotopy in $K(\mathbb{T})$ across $(q_1,\dots,q_A)$. We
apply Lemma~\ref{8.10:lem-homotopy-moves} with the complex $K(\mathbb{T})$ and the
configuration $V$, and obtain
\begin{equation*}
\bigl\|T_\gamma(V^{\mathrm{ext}})-1\bigr\|
=\bigl\|T_\gamma(V)-1\bigr\|
\le\sum_{i\le A}F_{q_i}(V).
\end{equation*}
We split this sum according to whether $q_i\in\mathsf{F}$ or not.

Suppose first that $q_i\in\mathsf{F}$. Since $V^{\mathrm{in}}$ is a margin-$\mu$ filling of
$V^{\mathrm{ext}}$, every plaquette in $\mathsf{F}$ carries $F_{q_i}(V)\le\varepsilon_\mu$.
These indices therefore contribute at most
$\varepsilon_\mu\cdot\#\{\,i : q_i\in\mathsf{F}\,\}$.

Suppose next that $q_i\notin\mathsf{F}$. By
Lemma~\ref{8.10:lem-collar-plaquettes}\,(a), $q_i$ has no interior edge. By
Lemma~\ref{8.10:lem-collar-plaquettes}\,(c), $F_P$ is a function of $V^{\mathrm{ext}}$ for
every plaquette $P$ without interior edge. Hence $F_{q_i}(V)=F_{q_i}(V^{\mathrm{ext}})$.

Adding the two contributions gives~\eqref{8.10:eq-flux-bound-1}.
\end{proof}

\begin{remark}[Coherent flux on long box hulls. Stated; proof incomplete at the step marked $(\ast)$]
\label{8.10:rem-coherent-flux}
Lemma~\ref{8.10:lem-flat-filling} leaves open a quantitative question. Let $C$ be unlinked, and
suppose that its collar is flat in the following sense:
\begin{itemize}
\item $F_{p'}(V^{\mathrm{ext}})\le\delta$ for every $p'\in\mathsf{Ex}_2$, with $\delta$ a fixed
multiple of $\varepsilon_\mu$;
\item every cell average is at most $f$;
\item no bound is imposed on the collar energy.
\end{itemize}
Does it follow that $V^{\mathrm{ext}}\in\mathrm{Fill}_{\mu,\mathfrak{E}}(C)$, with constants
independent of the extent of $C$? Granted part~(a), whose proof is incomplete at the step marked
$(\ast)$, part~(c) below shows that it does not, through a coherent flux carried by a long box.

Let $d=4$, and let $C$ be box-like (as in case~$(\beta)$ of
Theorem~\ref{8.10:thm-quantitative-filling}), with site box $Q$ of lattice side $T$ in the
directions $1,2,3$ and lattice side $h$ in direction $4$, with
\begin{equation*}
T\ge 2h\ge 4 .
\end{equation*}
Choose lattice coordinates in which $Q=\{0,\dots,T\}^3\times\{0,\dots,h\}$.
Let $\mathcal{L}$ be the boundary loop of the lattice rectangle $[-2,T+2]\times[-2,h+2]$ in the
$(x_2,x_4)$-plane at a central position $(x_1^0,x_3^0)$. All vertices of $\mathcal{L}$ lie outside
$Q$, so all its edges are exterior. By Lemma~\ref{8.10:lem-stokes-rectangle}, $\mathcal{L}$ has a
null-homotopy across the $(T+4)(h+4)$ plaquettes of that rectangle. Of these plaquettes, at most
\begin{equation*}
\mathfrak{n}_{\mathrm{in}}:=(T+2)(h+2)
\end{equation*}
have an interior edge, and at most
\begin{equation*}
\mathfrak{n}_{\mathrm{out}}:=4(T+h)+16
\end{equation*}
have none.
\begin{enumerate}
\item[(a)] There are an absolute constant $c_{\mathrm{coh}}>0$ and, for every $\delta\in\,]0,1]$,
an exterior configuration $V^{\mathrm{ext}}_\delta$ with values in a fixed maximal torus of
$\mathrm{SU}(N)$, constructed below, with the following two properties. First,
$F_p(V^{\mathrm{ext}}_\delta)\le\delta$ for every plaquette $p$ without interior edge, hence every
cell average is at most $\delta$. Second,
\begin{equation}\label{8.10:eq-coherent-flux-1}
\|T_{\mathcal{L}}(V^{\mathrm{ext}}_\delta)-1\|
\ge\min\Bigl(1,\;\frac{c_{\mathrm{coh}}\,\delta\,T^2}{\log(T/h)}\Bigr).
\end{equation}
\item[(b)] Let $V^{\mathrm{ext}}$ be an exterior configuration with
$F_p(V^{\mathrm{ext}})\le\delta$ for every plaquette $p$ without interior edge, and which
satisfies \eqref{8.10:eq-coherent-flux-1} in place of $V^{\mathrm{ext}}_\delta$. If
\begin{equation}\label{8.10:eq-coherent-flux-2}
\min\Bigl(1,\;\frac{c_{\mathrm{coh}}\,\delta\,T^2}{\log(T/h)}\Bigr)
>\mathfrak{n}_{\mathrm{in}}\,\varepsilon_\mu+\mathfrak{n}_{\mathrm{out}}\,\delta ,
\end{equation}
then $V^{\mathrm{ext}}\notin\mathrm{Fill}_{\mu,\mathfrak{E}}(C)$ for every $\mathfrak{E}\ge0$. This
holds although $V^{\mathrm{ext}}$ lies in the set $G^{\mathrm{av}}_\delta$ of exterior
configurations all of whose cell averages are at most $\delta$, and is pointwise $\delta$-flat on
the collar.
\item[(c)] Granted~(a), fix the ratio $\varepsilon_\mu/\delta$; an example is $\delta$ equal to
$\varepsilon_\mu$ divided by a constant depending on $d,N,L,M$ only. Then
\eqref{8.10:eq-coherent-flux-2} holds whenever
\begin{equation}\label{8.10:eq-coherent-flux-3}
\frac{T}{h\log(T/h)}>\frac{3\varepsilon_\mu/\delta+5}{c_{\mathrm{coh}}}
\qquad\text{and}\qquad
\Bigl(3+\frac{5\delta}{\varepsilon_\mu}\Bigr)\,T\,h\,\varepsilon_\mu<1 .
\end{equation}
For fixed $h$, the first condition holds for all large $T$, and the second for all $T$ below a
bound that grows without limit as $\varepsilon_\mu\to0$ at fixed ratio $\varepsilon_\mu/\delta$;
hence values of $T$ satisfying both conditions exist once $\varepsilon_\mu$ is small enough in
terms of that ratio, $h$ and $c_{\mathrm{coh}}$. No hypothesis on $C$ bounds the extent $T$ of a
component. Consequently no inclusion
\begin{equation*}
\{\text{pointwise $\delta$-flat collar}\}\cap G^{\mathrm{av}}_f\cap G^{\mathrm{en}}_E
\subset\mathrm{Fill}_{\mu,\mathfrak{E}},
\end{equation*}
with $\delta$ a fixed fraction of $\varepsilon_\mu$, where $G^{\mathrm{av}}_f$ is the set of
exterior configurations all of whose cell averages are at most $f$ and $G^{\mathrm{en}}_E$ the set
of exterior configurations with collar energy at most $E$, holds uniformly in the extent of
unlinked components. In particular, no size-uniform inclusion of this form can serve as a
hypothesis.
\end{enumerate}
Accordingly, membership of the exterior configuration in the fillable set of $C$ is imposed as an
explicit condition in the definition of the good configurations for every $C$. Configurations of
the coherent-flux type of part~(a) are among those whose weight the typicality estimate must
control. The topological statements are unaffected by the above:
\begin{itemize}
\item the exact flat filling and the openness assertions of Lemma~\ref{8.10:lem-flat-filling};
\item Theorem~\ref{8.10:thm-quantitative-filling}, in which membership in the fillable set is
automatic with a constant that depends on the extent of $C$. No size-uniform version of that
constant is used.
\end{itemize}
Part~(a) is stated; its proof below is incomplete at the step marked $(\ast)$. Parts~(b) and~(c)
are proved, part~(c) as an implication from part~(a).
\end{remark}

\begin{proof}[Proof of part~(a), incomplete at the step marked $(\ast)$]
Fix a maximal torus of $\mathrm{SU}(N)$ and a nonzero element $H_0$ of its Lie algebra. Let $\tau$
be a tent function of width $T$. Let $\alpha$ be the winding form of the cross-section rectangle of
$Q$ in the $(x_2,x_4)$-plane, averaged over the points of that rectangle. For a constant
$c\asymp\delta T^2/\log(T/h)$ put
\begin{equation*}
\eta:=c\,\tau(x_1)\,\tau(x_3)\,\alpha(x_2,x_4),
\end{equation*}
a $1$-form. For each exterior edge $e$ put
\begin{equation*}
A_e:=\int_e\eta,\qquad V_e:=\exp(A_eH_0),
\end{equation*}
and let $V^{\mathrm{ext}}_\delta$ be the configuration $(V_e)$. Reversing an edge changes the sign
of $A_e$ and inverts $V_e$, so the definition is consistent with orientation.

All $V_e$ are exponentials of real multiples of the single element $H_0$, so they commute. A
product of them is therefore the exponential of the sum of the exponents. Hence, for every edge
loop $\gamma$ made of exterior edges,
\begin{equation*}
T_\gamma(V^{\mathrm{ext}}_\delta)=\exp\Bigl(\Bigl(\oint_\gamma\eta\Bigr)H_0\Bigr).
\end{equation*}
This applies in particular to $\gamma=\mathcal{L}$ and to $\gamma=\partial p$ for every plaquette
$p$ without interior edge. Both assertions of~(a) are thus statements about the integrals of $\eta$
over $\mathcal{L}$ and over the boundaries of such plaquettes. The construction is the lattice form
of one in the abelian continuum setting with the averaged winding form~$\alpha$.

$(\ast)$ It remains to transcribe to the lattice two estimates of planar potential theory on the
averaged winding form: a bound of $|\alpha|$ by a constant multiple of $\log(T/h)/T$, and the
closedness $d\alpha=0$ outside the cross-section rectangle. From these,
$F_p(V^{\mathrm{ext}}_\delta)\le\delta$ on plaquettes without interior edge and
\eqref{8.10:eq-coherent-flux-1} are to follow, the constant $c_{\mathrm{coh}}$ being absolute. The
proof stops at this step.
\end{proof}

\begin{proof}[Proof of part~(b)]
Suppose, for contradiction, that $V^{\mathrm{ext}}\in\mathrm{Fill}_{\mu,\mathfrak{E}}(C)$ for some
$\mathfrak{E}\ge0$. Then $V^{\mathrm{ext}}$ has a margin-$\mu$ filling $V^{\mathrm{in}}$.

Apply Lemma~\ref{8.10:lem-stokes-rectangle} to the directions $2$ and $4$, with side lengths
$A=T+4$, $B=h+4$ and corner $y=(x_1^0,-2,x_3^0,-2)$. Its boundary loop is $\mathcal{L}$, and it
supplies a null-homotopy of
$\mathcal{L}$ across the $(T+4)(h+4)$ plaquettes $q_1,\dots,q_{(T+4)(h+4)}$ of the rectangle. All
edges of $\mathcal{L}$ are exterior, so Lemma~\ref{8.10:lem-flux-bound} applies to
$\gamma=\mathcal{L}$ with this null-homotopy, and gives
\begin{equation*}
\|T_{\mathcal{L}}(V^{\mathrm{ext}})-1\|
\le\varepsilon_\mu\,\#\{i:q_i\in\mathsf{F}\}+\sum_{i:\,q_i\notin\mathsf{F}}F_{q_i}(V^{\mathrm{ext}}).
\end{equation*}
The number of $q_i$ in $\mathsf{F}$ is at most the number of plaquettes of the rectangle with an
interior edge, hence at most $\mathfrak{n}_{\mathrm{in}}$. The remaining $q_i$ have no interior
edge; there are at most
$\mathfrak{n}_{\mathrm{out}}$ of them, and each satisfies $F_{q_i}(V^{\mathrm{ext}})\le\delta$ by
hypothesis. Therefore
\begin{equation*}
\|T_{\mathcal{L}}(V^{\mathrm{ext}})-1\|
\le\mathfrak{n}_{\mathrm{in}}\,\varepsilon_\mu+\mathfrak{n}_{\mathrm{out}}\,\delta .
\end{equation*}
Together with \eqref{8.10:eq-coherent-flux-1}, this contradicts
\eqref{8.10:eq-coherent-flux-2}. Hence $V^{\mathrm{ext}}\notin\mathrm{Fill}_{\mu,\mathfrak{E}}(C)$,
and since $\mathfrak{E}\ge0$ was arbitrary, this holds for every $\mathfrak{E}$.

The two remaining properties follow from the hypothesis. Since $F_p(V^{\mathrm{ext}})\le\delta$ on
every plaquette without interior edge, every cell average is at most $\delta$, so
$V^{\mathrm{ext}}\in G^{\mathrm{av}}_\delta$, and $V^{\mathrm{ext}}$ is pointwise $\delta$-flat on
the collar.
\end{proof}

\begin{proof}[Proof of part~(c)]
Put $\Phi:=c_{\mathrm{coh}}\delta T^2/\log(T/h)$; note that
$\log(T/h)\ge\log2>0$. From $T\ge4$ and $h\ge2$ we get $T+2\le\tfrac32T$ and $h+2\le2h$, hence
\begin{equation*}
\mathfrak{n}_{\mathrm{in}}\le 3Th .
\end{equation*}
From $h\le T/2$ and $16\le4T$ we get
\begin{equation*}
\mathfrak{n}_{\mathrm{out}}\le 4T+2T+4T=10T .
\end{equation*}
The right side of \eqref{8.10:eq-coherent-flux-2} is therefore at most
$3Th\,\varepsilon_\mu+10T\delta$.

We first show $\Phi>3Th\,\varepsilon_\mu+10T\delta$. This inequality is equivalent to
\begin{equation*}
c_{\mathrm{coh}}\,T/\log(T/h)>3(\varepsilon_\mu/\delta)h+10 .
\end{equation*}
Since $h\ge2$ gives $10\le5h$, we have
$3(\varepsilon_\mu/\delta)h+10\le(3\varepsilon_\mu/\delta+5)h$. So the inequality follows from the
first condition in \eqref{8.10:eq-coherent-flux-3}.

Next we show $1>3Th\,\varepsilon_\mu+10T\delta$. Using $10T\le5Th$,
\begin{equation*}
3Th\,\varepsilon_\mu+10T\delta\le\Bigl(3+\frac{5\delta}{\varepsilon_\mu}\Bigr)Th\,\varepsilon_\mu ,
\end{equation*}
which is less than $1$ by the second condition in \eqref{8.10:eq-coherent-flux-3}.

Both members of the minimum thus exceed the right side of \eqref{8.10:eq-coherent-flux-2}, so
\eqref{8.10:eq-coherent-flux-2} holds.

For fixed $h$, the first condition of \eqref{8.10:eq-coherent-flux-3} holds for all large $T$,
because $T/(h\log(T/h))\to\infty$ as $T\to\infty$. The second condition reads
\begin{equation*}
T<\Bigl(\Bigl(3+\frac{5\delta}{\varepsilon_\mu}\Bigr)h\,\varepsilon_\mu\Bigr)^{-1},
\end{equation*}
and its right side tends to infinity as $\varepsilon_\mu\to0$ with $\varepsilon_\mu/\delta$ fixed.
Hence values of $T$ satisfying both conditions exist once $\varepsilon_\mu$ is small enough in terms
of $\varepsilon_\mu/\delta$, $h$ and $c_{\mathrm{coh}}$.

For such $T$, let $V^{\mathrm{ext}}_\delta$ be as in~(a). It has a pointwise $\delta$-flat collar
and lies in $G^{\mathrm{av}}_\delta\subset G^{\mathrm{av}}_f$ for every $f\ge\delta$. Its collar
energy is some finite number; since $E$ is a free parameter of the inclusion, we may take $E$ at
least that number, and then $V^{\mathrm{ext}}_\delta$ lies in $G^{\mathrm{en}}_E$. By~(b), it lies
in no $\mathrm{Fill}_{\mu,\mathfrak{E}}(C)$. This is the assertion of~(c).
\end{proof}

\begin{remark}[Scope of the filling results]\label{8.10:rem-filling-scope}
The following points delimit the scope of Theorem~\ref{8.10:thm-quantitative-filling} and
Lemma~\ref{8.10:lem-flat-filling}.

\begin{enumerate}
\item With $\delta=0$ in Theorem~\ref{8.10:thm-quantitative-filling}, that theorem together
with part~(i) of Lemma~\ref{8.10:lem-flat-filling} states the following. For every unlinked
component $C$, and in particular for every box-like one, every exterior configuration that
is flat on $\mathsf{Ex}$ is fillable within the window, with zero filling energy. By part~(ii)
of Lemma~\ref{8.10:lem-flat-filling}, the intersection of $\mathrm{Fill}_{\mu,\mathfrak{E}}(C)$
with the averaging good set $G^{\mathrm{av}}_{f}$ and the energy good set $G^{\mathrm{en}}_{E}$,
\begin{equation*}
\mathrm{Fill}_{\mu,\mathfrak{E}}(C)\cap G^{\mathrm{av}}_{f}\cap G^{\mathrm{en}}_{E},
\end{equation*}
contains an open neighbourhood of the exteriors that are flat on $\mathsf{Ex}$. The statement
is about this neighbourhood. It does not say that the set of flat exteriors itself has
interior. For components $C$ with linking count $\operatorname{lk}^{\sharp}(C)\ge 1$, the
proposition of this chapter on linked components shows that the opposite holds.

\item Theorem~\ref{8.10:thm-quantitative-filling} works in the regime
$\delta\lesssim\varepsilon_\mu/\hat{\mathfrak{a}}$, with the margin scale $\varepsilon_\mu$, and
this regime is atypical. Under the
law defined by $\mathcal{M}$, a typical collar plaquette has size of order $g$, and the
collar plaquettes are incoherent. That such exteriors are fillable with margin and with small
filling energy, with high probability, is the typicality statement for fillability. That
statement belongs to the typicality input for the family $G\supset\mathrm{Fill}$. An
interior regularity estimate is an input of that typicality statement. It is used on the
typicality side, and it is not an input of the old-member statement, that is, of the lemma
of this chapter that is formulated on the domain $G_{\lambda,C}(\mathfrak{E})$ for every $C$.

\item No result of this section is a hypothesis of the old-member statement. That statement
is formulated on the domain $G_{\lambda,C}(\mathfrak{E})$ for every $C$.
\end{enumerate}
\end{remark}

\begin{remark}[Rough collars and the energy clause of the good set]\label{8.10:rem-rough-collars}
In the abelian caricature, a Bianchi-type identity on the unit $3$-cubes that meet both
$\mathsf{F}$ and $\mathsf{Ex}_2$ bounds $\mathsf{E}(V)$ from below by a positive multiple of
the squared exterior curvature flux entering those $3$-cubes. The fluxes through several
exterior faces of one $3$-cube may cancel. The bound is therefore on the net flux, and not
on the sum $\sum' F^2$ of the squared curvatures over the exterior faces of those $3$-cubes
taken separately.

Now consider exteriors admissible in the weaker sense given by the pointwise bound
$F\le c_W\varepsilon$ and the bound $f_*$ on cell averages alone, without the energy clause.
Such exteriors may force fillings of lattice energy of order
$\varepsilon^2$ per interface site. The $x$-multiple of this energy is of order $p_0(x)^2$
per site, with $p_0(\cdot)$ the degree function, and no linear reserve pays for it.

These exteriors are excluded from $\mathrm{Fill}_{\mu,\mathfrak{E}_*}$ by the energy clause.
They carry their own action under the law defined by $M$, and so they belong to
the self-penalised part of the complement. This is why $\mathrm{Fill}$, and not
$G^{\mathrm{av}}_{f_*}$ alone, enters the good set $G_{\lambda,C}(\mathfrak{E})$ of
Definition~\ref{8.10:def-good-set}, even for unlinked components.
\end{remark}

The following Laplace-type lower bound at the terminal level compares the window-constrained local small-field integral $a_\lambda$, taken around an arbitrary margined filling of the hole, with the local integral $m_\lambda$ of the majorant, at the cost of the absolute filling action.

\begin{proposition}[The window-constrained Laplace lower bound; proof outstanding]
\label{8.10:prop-window-lower-bound}
There are constants $C_L\ge 0$, $c_H\ge 0$, $c_A>0$ and $c_L'\ge 1$ with the following property. They depend only on $(d,N,L,M,\mu,p_1,p_2)$ and on the constants of Ba{\l}aban's construction fixed at Stage~0 of the order of choice
(Definition~\ref{2.3:not-stages-first}).

Assume the hull-width conditions (C-$p_1$) and (C-$p_2$) of Definition~\ref{8.10:def-one-sided-conditions}. Assume also the one-sided coupling condition of the same definition, which reads
\begin{equation*}
\mu\, p_0(x)\ \ge\ c_L' ,
\qquad\text{equivalently}\qquad x\ \ge\ x_L ,
\end{equation*}
with $p_0(\cdot)$ the degree function and $x$ the inverse coupling of Definition~\ref{8.10:def-one-sided-conditions}. Then the following holds for every $K$, $m$ and every value of the terminal parameter $N_T$, and for every touch-connected union $C$ of terminal cells. Let $\lambda\in\Lambda_C$ be a history in the sense of Definition~\ref{8.10:def-histories-fibres} satisfying $K_{\mathrm{free}}(C;\lambda;p_1+p_2+p_W)$ in the sense of the same definition, where $p_W$ is the localisation radius of the windows of Definition~\ref{8.10:def-window-structure} (the window bounds stated there have proofs outstanding). Let $\xi\in\mathcal{U}_\lambda$, let $\mathfrak{E}\ge 0$, and let $V^{\mathrm{ext}}\in G_{\lambda,C}(\mathfrak{E})$, the good set of Definition~\ref{8.10:def-good-set}. Then
\begin{equation}\label{8.10:eq-window-lower-bound-1}
\begin{split}
m_\lambda(V^{\mathrm{ext}},\xi)\ \le\ 
&\exp\Bigl(C_L\,|\mathsf{N}_0| + c_H\,|\mathsf{R}_1|\,\log x\\
&\qquad\quad + c_A\, x\, E^{\mu}_{\mathrm{fill}}(V^{\mathrm{ext}})\Bigr)\,
a_\lambda(V^{\mathrm{ext}},\xi).
\end{split}
\end{equation}
The notation is as follows.
\begin{itemize}
\item $m_\lambda$ and $a_\lambda$ are the two local integrals of Definition~\ref{8.10:def-admissible-exteriors}.
\item $\mathcal{U}_\lambda$ is the space of internal variables of $\lambda$.
\item $E^{\mu}_{\mathrm{fill}}$ is the absolute filling action with margin $\mu$ of Definition~\ref{8.10:def-good-set}.
\item $\mathsf{N}$ is the neighbourhood of $C$ outside which the factors of the two local integrands agree; $\mathsf{N}_0$ is the set of sites near $C$, within $\mathsf{N}$, on which the boundary terms and the differences of normalisation are localised, and $|\mathsf{N}_0|$ is its size.
\item $\mathsf{R}_1$ is the region of bonds around $C$, containing the bonds of the cells of $C$, on which the comparison with the majorant's integral $m_\lambda$ is made, and $|\mathsf{R}_1|$ is its size.
\end{itemize}
\end{proposition}

Proof outstanding.

\medskip\noindent\emph{Route.}
A proof would have to be built from Ba{\l}aban's one-level small-field analysis at the terminal level. It would use the following inputs:
\begin{itemize}
\item the regularity theory for the variational problem \cite[Theorem~1]{B-CMP102b};
\item the form of the effective action \cite[\S1]{B-CMP109};
\item Ba{\l}aban's propagator bounds;
\item the boundary terms of \cite{B-CMP119}.
\end{itemize}
It would have to establish three things, uniformly in the exterior configuration on $G_{\lambda,C}(\mathfrak{E})$ and in the internal variables $\xi$ under $K_{\mathrm{free}}$.

\begin{enumerate}
\item[(i)] \emph{Laplace form.} Write the integrand of $a_\lambda$ in exponential coordinates on the space $\mathcal{V}^{\mathrm{in}}$ of interior bond variables. In these coordinates it must be $e^{-x\mathcal{A}_{\mathrm{eff}}}$ times the window factors of Definition~\ref{8.10:def-window-structure}. Here $\mathcal{A}_{\mathrm{eff}}$ is an effective action of the form described in \cite[\S1]{B-CMP109}, with Hessian of order one per coordinate and a controlled remainder.

\item[(ii)] \emph{Reference volumes.} The kernel of this step is the lemma on the Gaussian volume around a margined filling, stated below in this section. The base point is the interior configuration $V^{\mathrm{in}}$ that completes $V^{\mathrm{ext}}$ by a margined filling. It need not be a critical point of $\mathcal{A}_{\mathrm{eff}}$; this costs nothing, since the linear term of the action averages to zero over a cube symmetric about $V^{\mathrm{in}}$. Take the integral over the cube of side $x^{-1/2}$ in each coordinate around $V^{\mathrm{in}}$. This cube lies inside every window, by the margin $\mu$ and the one-sided coupling condition. Indeed, a bond fluctuation of the size of the cube side, $x^{-1/2}$, moves each plaquette variable by $O(x^{-1/2})$. This is at most $\mu\varepsilon/B_W$ when $\mu\,p_0(x)$ is large, with $\varepsilon$ the small-field parameter of the windows of Definition~\ref{8.10:def-window-structure} and $B_W$ a constant attached to those windows. This step is exactly the content of the one-sided coupling condition, and it removes the need to estimate the fraction of volume occupied by a window.

The cube integral gives the Gaussian volume
\begin{equation*}
x^{-\dim\mathcal{V}^{\mathrm{in}}/2}\,e^{-O(1)\dim\mathcal{V}^{\mathrm{in}}}
\end{equation*}
times $e^{-x\mathcal{A}_{\mathrm{eff}}(V^{\mathrm{in}})}$, where $\dim\mathcal{V}^{\mathrm{in}}$ is the number of coordinates. This must be compared with the reference integral of the majorant. That integral carries Haar measure of total mass one on the bonds of the cells of $C$, whose variables are unweighted there. Off $C$ it carries the integrand of the hole background. The mismatch is
\begin{equation*}
\tfrac12 (N^2-1)\log x + O(1)
\end{equation*}
per unweighted bond. Inequality \eqref{8.10:eq-window-lower-bound-1} charges this mismatch on all of $\mathsf{R}_1$ through the term $c_H|\mathsf{R}_1|\log x$. Charging it there, rather than on the bonds of $C$ alone, removes the need for a Laplace upper bound on the majorant's own integral over the bonds of $\mathsf{R}_1$ outside $C$. This over-charge is a power of $x$ whose degree is bounded per cell of $\mathsf{R}_1$, and is affordable where \eqref{8.10:eq-window-lower-bound-1} is used later in this section; the mismatch computed above arises per unweighted bond, and the unweighted bonds are those of the cells of $C$.

The action excess charged is the absolute filling action of $V^{\mathrm{in}}$. It is comparable to $c_A\,x\,E^{\mu}_{\mathrm{fill}}(V^{\mathrm{ext}})$, since the minimum on the hole side is nonnegative. Because the charge is absolute rather than relative to the hole side's minimum, no lower bound on the difference between the filling action and that minimum is needed near the boundary of the interior region, where this difference may vanish.

\item[(iii)] \emph{Localisation.} Consider the factors attached to variables outside the neighbourhood $\mathsf{N}$ of $C$. The factors of the majorant $M_{\lambda,C}$ in the domination inequality of Definition~\ref{8.10:def-fibre-domination}, whose proof is outstanding, must agree there with those of the terminal small-field density $\rho_\lambda$ attached to $\lambda$; this agreement rests on the window format of Definition~\ref{8.10:def-window-structure}, and on the hull-width condition at width $p_1$ of Definition~\ref{8.10:def-one-sided-conditions}. The boundary terms of \cite{B-CMP119}, together with Ba{\l}aban's localisation bounds near $C$, and the differences of normalisation near $C$ are of order one per site of $\mathsf{N}_0$, which gives the term $C_L|\mathsf{N}_0|$; whether some of these terms are instead of order $\log x$ per site, and then belong with the logarithmic charge, is a point a proof would have to settle. A proof would also have to establish uniformity in $\xi$ beyond the radius in $K_{\mathrm{free}}$: the internal variables of $\lambda$ located beyond that radius enter only through screened, far contributions, and it is the bound on these far contributions, uniform in the internal variables of a neighbouring component, that is needed here. Finally, a single majorant is used for each fibre in the domination inequality of Definition~\ref{8.10:def-fibre-domination}, whose proof is outstanding.
\end{enumerate}

\begin{lemma}[The Laplace kernel around a margined competitor]\label{8.10:lem-competitor-kernel}
Let $D\ge 1$ be an integer, and let $x>0$, $\varepsilon'>0$ and $\mu'\in\,]0,1[$. Let
$\mathfrak{m}\ge1$ be an integer and let
$\eta_1,\dots,\eta_{\mathfrak{m}}\in\mathbb{R}^D$ have $\ell^1$-norms bounded by $c_\eta$:
\begin{equation*}
\sum_{j=1}^{D}|\eta_{i,j}|\le c_\eta\qquad(i=1,\dots,\mathfrak{m}).
\end{equation*}
Put
\begin{equation*}
\mathcal{W}:=\{z\in\mathbb{R}^D:\ |\eta_i\cdot z|\le\varepsilon'\ \text{for all } i\},
\end{equation*}
and let $z_0\in\mathbb{R}^D$ satisfy $|\eta_i\cdot z_0|\le(1-\mu')\varepsilon'$ for all $i$. Put
$r:=x^{-1/2}$ (in this lemma and its proof only) and assume the one-sided condition
$c_\eta\, r\le\mu'\varepsilon'$. Let
\begin{equation*}
\mathcal{Q}:=\{z\in\mathbb{R}^D:\ \max_{j}|z_j-z_{0,j}|\le r\},
\end{equation*}
and let $\Phi:\mathcal{W}\to\mathbb{R}$ be measurable and such that, for some $v\in\mathbb{R}^D$ and
some $\Lambda\ge0$,
\begin{equation*}
\Phi(z)\le\Phi(z_0)+v\cdot(z-z_0)+\tfrac12\Lambda|z-z_0|^2\qquad\text{for } z\in\mathcal{Q}.
\end{equation*}
Then $\mathcal{Q}\subset\mathcal{W}$, and
\begin{equation}\label{8.10:eq-competitor-kernel-1}
\int_{\mathcal{W}}e^{-x\Phi(z)}\,dz\;\ge\;\int_{\mathcal{Q}}e^{-x\Phi(z)}\,dz
\;\ge\;e^{-x\Phi(z_0)}\,e^{-\Lambda D/2}\,2^D\,x^{-D/2}.
\end{equation}
\end{lemma}

\begin{proof}
Let $z\in\mathcal{Q}$. For each $i$,
\begin{equation*}
|\eta_i\cdot(z-z_0)|\le\sum_{j=1}^D|\eta_{i,j}|\,|z_j-z_{0,j}|\le c_\eta\, r\le\mu'\varepsilon',
\end{equation*}
by the bound on the $\ell^1$-norm of $\eta_i$, the definition of $\mathcal{Q}$, and the one-sided
condition $c_\eta r\le\mu'\varepsilon'$. Together with $|\eta_i\cdot z_0|\le(1-\mu')\varepsilon'$
this gives
\begin{equation*}
|\eta_i\cdot z|\le|\eta_i\cdot z_0|+|\eta_i\cdot(z-z_0)|\le(1-\mu')\varepsilon'+\mu'\varepsilon'
=\varepsilon'
\end{equation*}
for every $i$, so $z\in\mathcal{W}$. Hence $\mathcal{Q}\subset\mathcal{W}$; in particular $\Phi$ is
defined on $\mathcal{Q}$. The integrand $e^{-x\Phi}$ is measurable and non-negative on
$\mathcal{W}$, so the integral over $\mathcal{W}$ is at least the integral over the subset
$\mathcal{Q}$; this is the first inequality in \eqref{8.10:eq-competitor-kernel-1}.

For the second inequality, note that on $\mathcal{Q}$ each coordinate satisfies
$|z_j-z_{0,j}|\le r$, so
\begin{equation*}
|z-z_0|^2=\sum_{j=1}^D|z_j-z_{0,j}|^2\le Dr^2=\frac{D}{x}.
\end{equation*}
Multiplying the hypothesis on $\Phi$ by $x>0$ and using $\Lambda\ge0$, we obtain for $z\in\mathcal{Q}$
\begin{equation*}
x\Phi(z)\le x\Phi(z_0)+x\,v\cdot(z-z_0)+\tfrac12\Lambda x\cdot\frac{D}{x}
=x\Phi(z_0)+x\,v\cdot(z-z_0)+\frac{\Lambda D}{2},
\end{equation*}
and therefore
\begin{equation}\label{8.10:eq-competitor-kernel-2}
\int_{\mathcal{Q}}e^{-x\Phi(z)}\,dz\ge e^{-x\Phi(z_0)}\,e^{-\Lambda D/2}
\int_{\mathcal{Q}}e^{-x\,v\cdot(z-z_0)}\,dz .
\end{equation}
It remains to bound the last integral from below. The cube $\mathcal{Q}$ has volume
$\operatorname{vol}(\mathcal{Q})=(2r)^D$, and $dz/\operatorname{vol}(\mathcal{Q})$ is a probability
measure on $\mathcal{Q}$. The function $z\mapsto -x\,v\cdot(z-z_0)$ is bounded on $\mathcal{Q}$,
hence integrable, and the exponential function is convex; Jensen's inequality gives
\begin{equation*}
\int_{\mathcal{Q}}e^{-x\,v\cdot(z-z_0)}\,dz\ge\operatorname{vol}(\mathcal{Q})\,
\exp\Bigl(-x\,\operatorname{vol}(\mathcal{Q})^{-1}\int_{\mathcal{Q}}v\cdot(z-z_0)\,dz\Bigr).
\end{equation*}
The cube $\mathcal{Q}$ is symmetric about $z_0$: the reflection $z\mapsto 2z_0-z$ maps $\mathcal{Q}$
onto itself, preserves Lebesgue measure and changes the sign of $v\cdot(z-z_0)$. Hence
$\int_{\mathcal{Q}}v\cdot(z-z_0)\,dz=0$, the exponential on the right equals one, and
\begin{equation*}
\int_{\mathcal{Q}}e^{-x\,v\cdot(z-z_0)}\,dz\ge\operatorname{vol}(\mathcal{Q})=(2r)^D=2^D x^{-D/2}.
\end{equation*}
Multiplying this by $e^{-x\Phi(z_0)}e^{-\Lambda D/2}>0$ and inserting it into
\eqref{8.10:eq-competitor-kernel-2} gives the second inequality in
\eqref{8.10:eq-competitor-kernel-1}.
\end{proof}

\begin{remark}[The parameters of Lemma~\ref{8.10:lem-competitor-kernel} in its application]
In the application of Lemma~\ref{8.10:lem-competitor-kernel} in this chapter, the parameters are
specialised as follows. The dimension is $D=\dim\mathfrak{g}\cdot|\mathsf{R}_1|$, where
$\mathfrak{g}$ is the Lie algebra of $\mathrm{SU}(N)$ and $\mathsf{R}_1$ is the bond set of the
competitor construction, with one block of $\dim\mathfrak{g}$ coordinates per bond. The vectors
$\eta_i$ are the linearised window functionals of the window condition, written
in exponential coordinates around the competitor; their $\ell^1$-norms are bounded by a constant
depending on $(d,N,L,M)$ only. The number $\varepsilon'$ corresponds to $\varepsilon/B_W$, where
$\varepsilon$ is the threshold and $B_W$ the constant of the window condition, and $\mu'$
corresponds to the margin parameter $\mu$ of that condition. The condition
$c_\eta x^{-1/2}\le\mu'\varepsilon'$ then reads $c_\eta\le\mu\,p_0(g)/B_W$, where $g$ is the
coupling with $g^{-2}=x$ and $p_0(\cdot)$ is the degree function of
Notation~\ref{0.2:not-glossary}; this is the condition imposed on the constant $c'_L$ of the window,
with $c'_L:=B_W c_\eta$. The factor $x^{-D/2}$, compared with Haar volume one on the unweighted
bonds, is the term $c_H\log x$ per bond of the bound for the competitor, $c_H$ denoting the
constant of that logarithmic term; the factor $e^{-\Lambda D/2}$ and the Haar density on
$\mathcal{Q}$ are bounded by a constant per coordinate and are contained in the constant $C_L$ of
that bound.
\end{remark}

\begin{proposition}[The old-member statement. Stated; proof incomplete at the step marked
$(\ast)$]\label{8.10:prop-old-member}
Assume the following hypotheses:
\begin{itemize}
\item the representation hypothesis on the terminal densities;
\item the single-majorant bound at the terminal level $K'$ with fractions $\theta,\theta'$;
\item the boundary-factor bound, which gives, for every $C$ as below, with $n(C)$ the number
of terminal cells of $C$, and for the factor
$e^{\delta_C}$ in the domination inequality of Definition~\ref{8.10:def-fibre-domination}
\textup{(}an inequality whose proof is outstanding\textup{)},
\begin{equation}\label{8.10:eq-old-member-1}
e^{\delta_C}\le \exp\bigl(C_\delta'\,\alpha\,3^d\,\check{R}_{K'}^{\,d}\,n(C)\bigr);
\end{equation}
\item the properties of the small-field windows and of the region formats stated in
Definition~\ref{8.10:def-window-structure}; the proof of these properties is outstanding;
\item the two one-sided hull-width conditions and the one-sided coupling floor of
Definition~\ref{8.10:def-one-sided-conditions}.
\end{itemize}
Let $K$, $m$ and the terminal-level parameter $N_T$ be given. Let $C$ be a touch-connected union
of terminal cells. Let
$\lambda\in\Lambda_C$ be a fibre satisfying the freeness condition
$\mathrm{K}_{\mathrm{free}}(C;\lambda;p_1+p_2+p_W)$ of
Definition~\ref{8.10:def-histories-fibres}, at hull width $p_1+p_2+p_W$, where $p_1$, $p_2$ are
the widths of the hull-width conditions of Definition~\ref{8.10:def-one-sided-conditions} and
$p_W$ is a further fixed width; let $\xi\in\mathcal{U}_\lambda$ be a value
of the internal variables of $\lambda$, let $\mathfrak{E}\ge 0$, and let
$V^{\mathrm{ext}}\in G_{\lambda,C}(\mathfrak{E})$, the good set of
Definition~\ref{8.10:def-good-set}. Then
\begin{equation}\label{8.10:eq-old-member-2}
\int e^{\delta_C}\,M_{\lambda,C}(\xi,V^{\mathrm{in}},V^{\mathrm{ext}})\,dV^{\mathrm{in}}
\le e^{\Gamma_C(\mathfrak{E})}
\int \rho_\lambda(\xi,V^{\mathrm{in}},V^{\mathrm{ext}})\,dV^{\mathrm{in}},
\end{equation}
where
\begin{equation}\label{8.10:eq-old-member-3}
\Gamma_C(\mathfrak{E}):=C_\delta'\,\alpha\,3^d\,\check{R}_{K'}^{\,d}\,n(C)
+C_L\,|\mathsf{N}_0|+c_H\,|\mathsf{R}_1|\,\log x+c_A\,x\,\mathfrak{E},
\end{equation}
with $C_L$, $c_H$, $c_A$, $\mathsf{N}_0$, $\mathsf{R}_1$ as in the window-constrained Laplace
lower bound, Proposition~\ref{8.10:prop-window-lower-bound}, whose proof is outstanding, and $x$
the coupling variable of the
coupling floor of Definition~\ref{8.10:def-one-sided-conditions}. The bound
\eqref{8.10:eq-old-member-2} holds uniformly in $K$, $m$, $N_T$ and the position of $C$:
$\Gamma_C(\mathfrak{E})$ depends only on $n(C)$, $x$, $p_1$, $p_2$, $\mathfrak{E}$, the displayed
constants and the parameter $\ell$ of Proposition~\ref{8.10:prop-window-lower-bound}. Nothing is
asserted for
$V^{\mathrm{ext}}\notin G_{\lambda,C}(\mathfrak{E})$,
nor for fibres $\lambda$ violating $\mathrm{K}_{\mathrm{free}}(C;\lambda;p_1+p_2+p_W)$.
If Proposition~\ref{8.10:prop-window-lower-bound} and \eqref{8.10:eq-old-member-2} hold only for
$dV^{\mathrm{ext}}$-almost every $V^{\mathrm{ext}}\in G_{\lambda,C}(\mathfrak{E})$, this suffices
wherever the proposition is used, since the good set enters only as a domain of integration in the
exterior field.
\end{proposition}

\begin{proof}
Fix $K$, $m$, $N_T$, $C$, $\lambda$, $\xi$, $\mathfrak{E}$ and $V^{\mathrm{ext}}$ as in the
statement. Let $a_\lambda$ and $m_\lambda$ be the two local integrals of
Definition~\ref{8.10:def-admissible-exteriors} at $(\xi,V^{\mathrm{ext}})$. Thus $a_\lambda$ is
the integral over $V^{\mathrm{in}}$ on the right of \eqref{8.10:eq-old-member-2}, that is, the
local integral of the density $\rho_\lambda$ of the term in which $C$ is absent. Likewise
$m_\lambda$ is the integral over $V^{\mathrm{in}}$ of the hole majorant $M_{\lambda,C}$ (the net
majorant of the term in which $C$ is present) on the left
of \eqref{8.10:eq-old-member-2}, taken without the factor $e^{\delta_C}$.

First, the filling energy of the exterior is bounded by $\mathfrak{E}$. By
Definition~\ref{8.10:def-good-set} the good set is
\begin{equation*}
G_{\lambda,C}(\mathfrak{E})=\operatorname{Adm}(\lambda;C)\cap G^{\mathrm{av}}_{f_*}\cap
\mathrm{Fill}_{\mu,\mathfrak{E}}(C).
\end{equation*}
Hence $V^{\mathrm{ext}}\in
\mathrm{Fill}_{\mu,\mathfrak{E}}(C)$, and the definition of the margined fillings in the same
definition gives
\begin{equation}\label{8.10:eq-old-member-4}
E^{\mu}_{\mathrm{fill}}(V^{\mathrm{ext}})\le\mathfrak{E}.
\end{equation}

$(\ast)$ Next, we invoke the window-constrained Laplace lower bound,
Proposition~\ref{8.10:prop-window-lower-bound}, whose proof is outstanding. We apply it at the
fibre $\lambda$, the internal variables $\xi$ and the exterior $V^{\mathrm{ext}}$, whose filling
energy obeys \eqref{8.10:eq-old-member-4}, under the conditions assumed above. It gives
\begin{equation}\label{8.10:eq-old-member-5}
m_\lambda\le\exp\bigl(C_L\,|\mathsf{N}_0|+c_H\,|\mathsf{R}_1|\,\log x
+c_A\,x\,\mathfrak{E}\bigr)\,a_\lambda .
\end{equation}

Finally, we multiply \eqref{8.10:eq-old-member-5} by $e^{\delta_C}$. The left-hand side becomes
$e^{\delta_C}m_\lambda$, which is the left-hand side of \eqref{8.10:eq-old-member-2}. On the
right-hand side we bound $e^{\delta_C}$ by the boundary-factor bound
\eqref{8.10:eq-old-member-1}, which is one of the hypotheses. The two exponentials multiply,
and we obtain
\begin{equation*}
e^{\delta_C}m_\lambda\le
\exp\bigl(C_\delta'\,\alpha\,3^d\,\check{R}_{K'}^{\,d}\,n(C)\bigr)
\exp\bigl(C_L\,|\mathsf{N}_0|+c_H\,|\mathsf{R}_1|\,\log x+c_A\,x\,\mathfrak{E}\bigr)\,a_\lambda .
\end{equation*}
By \eqref{8.10:eq-old-member-3} the exponent on the right is $\Gamma_C(\mathfrak{E})$. This is
\eqref{8.10:eq-old-member-2}.

No constant entering \eqref{8.10:eq-old-member-1}, \eqref{8.10:eq-old-member-4} or
\eqref{8.10:eq-old-member-5} depends on $K$, $m$, $N_T$ or the position of $C$. The bound is
therefore uniform in these quantities, as stated.
\end{proof}

Four points are established above, or in the statements cited, without appeal to
Proposition~\ref{8.10:prop-window-lower-bound}.
\begin{enumerate}
\item[\textup{(a)}] Consider $V^{\mathrm{ext}}\in G_{\lambda,C}(\mathfrak{E})$. By
Lemma~\ref{8.10:lem-margined-competitor}, the window-constraint set, over the interior fields
$V^{\mathrm{in}}$, of the term in which $C$ is absent is non-empty. That lemma is proved as an
implication from the
window properties stated in Definition~\ref{8.10:def-window-structure}; the proof of those
properties is outstanding. The
constraint set contains a competitor that has margin $\mu$ on every window it influences. All held
windows are positive, and the filling energy of the competitor is at most $\mathfrak{E}$. Now take
$\mathfrak{E}=\mathfrak{E}_*(C)$. By the definition of the energy allowance $\mathfrak{E}_*(C)$,
the filling action charged on $G_{\lambda,C}(\mathfrak{E}_*(C))$ is
\begin{equation*}
c_A\,x\,\mathfrak{E}_*(C)=\tfrac14\,B^{\flat}_{\tau}(C),
\end{equation*}
that is, one quarter of the budget $B^{\flat}_{\tau}(C)$ that the allotted part
$\theta'_\tau$ of the linear creation reserve assigns to $C$; this lies inside the allowance
drawn from $\theta'_\tau$ against which the costs collected in \eqref{8.10:eq-old-member-3} are
paid.
\item[\textup{(b)}] Proposition~\ref{8.10:prop-window-lower-bound} and
\eqref{8.10:eq-old-member-2} compare the true local integral of the absent term with the integral
of the hole majorant $M_{\lambda,C}$. That majorant carries its own exterior factor, and the
comparison is not made against the number one. At the top level there is a mismatch between the
Haar measure of total mass one on the integrated bonds of $C$ and a Gaussian whose width is of
the order of the coupling at the terminal level.
This mismatch is displayed as the cost $c_H\,|\mathsf{R}_1|\,\log x$ in
\eqref{8.10:eq-old-member-3}. It is paid, with the other costs collected there, from the allotment
$\theta'_\tau$, and it is not absorbed into a constant depending on $d$ alone.
\item[\textup{(c)}] Nothing is used pointwise off the good set. The bound
\eqref{8.10:eq-old-member-2} is asserted only for $V^{\mathrm{ext}}\in G_{\lambda,C}(\mathfrak{E})$.
No quantitative use is made of the clause $G^{\mathrm{av}}_{f_*}$ of the good set, and no property
of the complement of the good set enters.
\item[\textup{(d)}] The bound is uniform: no constant depends on $K$, $m$, $N_T$ or the position
of $C$. The exterior field and the internal variables $\xi$ of $\lambda$ are held fixed, and the
infrared content of the top scale stays in the held exterior.
\end{enumerate}

\begin{remark}[Fibres with kept live structure near the component. Stated; proof incomplete at the step marked $(\ast)$]\label{8.10:rem-kept-structure}
The content of this remark is not used in what follows.

Let $C$ and $\lambda$ be as in the old-member statement, Proposition~\ref{8.10:prop-old-member}
(stated; proof incomplete), except that $\lambda$ violates the condition
$K_{\mathrm{free}}(C;\lambda;p_1+p_2+p_W)$ of that proposition, so that the fibre of $\lambda$ may
carry kept live structure within $p_1+p_2+p_W$ cells of $C$. The symbols $\mathfrak{E}$,
$\Gamma_C(\mathfrak{E})$, $G^{\mathrm{av}}$, $G^{\mathrm{mar}}$, $R_{\mathrm{col}}$, $\mathrm{Fill}$,
$\check{R}_{K'}$ and the creation step $K'$ are those of Proposition~\ref{8.10:prop-old-member};
the kept structure $\mathfrak{K}(\lambda)$ of $\lambda$, its hull and the argument $x$ of the
degree function $p_0(\cdot)$ are as fixed earlier in this chapter. For such $\lambda$
one expects the following (this expectation is the step marked $(\ast)$): the domination bound of
Proposition~\ref{8.10:prop-old-member} (stated; proof incomplete) holds with the following three
changes.
\begin{itemize}
\item[(a)] The collar clauses $G^{\mathrm{av}}$ and $G^{\mathrm{mar}}$ are taken on the collar
region $R_{\mathrm{col}}$ with the cells of $\operatorname{hull}_{p_W}(\mathfrak{K}(\lambda))$,
the hull of width $p_W$ cells about the kept structure $\mathfrak{K}(\lambda)$ of $\lambda$,
removed. The plaquettes of the removed cells
are unwindowed by $\lambda$ itself, by the format convention for $\lambda$.
\item[(b)] The filling clause $\mathrm{Fill}$ is unchanged.
\item[(c)] The exponent $\Gamma_C(\mathfrak{E})$ is augmented by a kept-structure term. This
term is of the order of
\begin{equation*}
c\,p_0(x)^2\,\bigl(M\check{R}_{K'}\bigr)^{-4}
\end{equation*}
for each kept component created at step $K'$ within the collar; here $c$ is a constant and
$p_0(x)$ is the degree function $p_0(\cdot)$ read at the argument $x$.
\end{itemize}
The term in (c) is the collar curvature that such a kept component forces. It can be paid only
from a reserve term in the exponent that is quadratic in the degree function and available once
for each cell at which a kept component is created, and only under a lower bound on $M$.

The step $(\ast)$ requires the last of the hypotheses listed in
Proposition~\ref{8.10:prop-old-member} (stated; proof incomplete) to hold uniformly in the
unwindowed internal variables of a kept component of the kind described in (c). This uniformity,
together with the forced collar curvature of (c), is not proved in this book.
\end{remark}

\begin{lemma}[Sizes of the hull complexes]\label{8.10:lem-hull-sizes}
Let the cube set $C$ and its cardinality $n(C)$, the hull widths $p_1$ and $p_2$, the length $\ell$,
and the sets $\mathsf{N}_0$, $\mathsf{N}_2$, $\mathsf{R}_1$, $\mathrm{Str}$, $\mathsf{F}$,
$\mathsf{Ex}_2$ of the hull complexes, together with the argument $x$ of the ceiling $\ell_*(x)$
and its exponent $r_{\mathrm{pr}}$, be as fixed in Notation~\ref{8.10:not-terminal-lattice}. Let
$c_d$ be the constant depending on $d$ only that appears in
Lemma~\ref{8.10:lem-obstruction-ranks}\textup{(a)}. Then
\begin{align*}
|\mathsf{N}_0| &\le (2p_1+2p_2+1)^d(\ell+1)^d\,n(C),\\
|\mathsf{R}_1| &\le d\,(2p_1+1)^d(\ell+1)^d\,n(C),\\
|\mathrm{Str}| \le |\mathsf{F}| &\le c_d\,(2p_1+3)^d(\ell+1)^d\,n(C),\\
|\mathsf{Ex}_2|,\ |\mathsf{N}_2| &\le c_d\,(2p_1+2p_2+1)^d(\ell+1)^d\,n(C),
\end{align*}
and
\begin{equation*}
\ell\le \ell_*(x)=ML(\log x)^{r_{\mathrm{pr}}}.
\end{equation*}
\end{lemma}

\begin{proof}
The bounds on $|\mathsf{N}_0|$, $|\mathsf{R}_1|$, $|\mathsf{Ex}_2|$ and $|\mathsf{N}_2|$, and the
ceiling $\ell\le\ell_*(x)$ with its exponent $r_{\mathrm{pr}}$, are the counts recorded in
Notation~\ref{8.10:not-terminal-lattice}, where these sets and the length $\ell$ are fixed.

For the plaquette sets we use Lemma~\ref{8.10:lem-collar-plaquettes}\textup{(a)}, the
decomposition of the collar plaquettes; the inequality $|\mathrm{Str}|\le|\mathsf{F}|$ is
obtained from that statement.

Finally, the bound
\begin{equation*}
|\mathsf{F}|\le c_d\,(\ell+1)^d(2p_1+3)^d\,n(C)
\end{equation*}
is the second inequality of Lemma~\ref{8.10:lem-obstruction-ranks}\textup{(a)}, the lemma on the
ranks and cohomological descriptions of the obstruction groups. Combined with the preceding
paragraph, it gives the third line of the display in the statement.
\end{proof}

\begin{definition}[The cost polynomial and the coupling-smallness condition]
\label{8.10:def-cost-polynomial}
Let $\ell_*(x)=ML(\log x)^{r_{\mathrm{pr}}}$ be the bound on $\ell$ of
Lemma~\ref{8.10:lem-hull-sizes}, in which $p_1,p_2$ are the widths entering the size
bounds of that lemma. Here $C_G',C_\delta',\alpha_1,C_L,c_H$ are positive constants
and $a_\alpha$ is an exponent; all of them are fixed in this chapter and do not
depend on $x$. For $x\ge e$ put
\begin{equation}\label{8.10:eq-cost-polynomial-1}
\begin{split}
\mathfrak{c}(x):=C_G'\cdot\Bigl[\,
&C_\delta'\,\alpha_1(1+\log x)^{a_\alpha}\,3^dL^d(\log x)^{d\,r_{\mathrm{pr}}}
+C_L(2p_1+2p_2+1)^d\bigl(\ell_*(x)+1\bigr)^d\\
&+c_H\,d\,(2p_1+1)^d\bigl(\ell_*(x)+1\bigr)^d\log x
+2d\log 3+\log 2+2\,\Bigr].
\end{split}
\end{equation}
This is a positive combination of powers of $\log x$ with top exponent at most
$d\,r_{\mathrm{pr}}+1$. We call $\mathfrak{c}$ the \emph{cost polynomial}.

Write $p_0(\cdot)$ for the degree function, read as a function of the inverse coupling,
$p_0(x)=A_0(\log x)^{\mathfrak{p}_0}$, where $\mathfrak{p}_0$ is the fixed exponent of the
degree function (written $p_0$ in Notation~\ref{0.2:not-glossary}). The exponents are
assumed to satisfy the one-sided degree condition $d\,r_{\mathrm{pr}}+1<2\mathfrak{p}_0$.
The constants $\theta_\tau'$ and $\gamma_0$ are the positive constants that enter,
together with the auxiliary parameter $A_1$, the core bound $B^\flat_\tau$ of
Definition~\ref{8.10:def-core-bound}, whose proof is outstanding. Put
\begin{equation}\label{8.10:eq-cost-polynomial-2}
\begin{split}
x_\tau:=\inf\Bigl\{x_0\ge e:\ \text{for all }x'\ge x_0,\ 
&\mathfrak{c}(x')\le\tfrac12\,\theta_\tau'\gamma_0A_1^2\bigl(p_0(x')-1\bigr)^2\\
&\text{and }p_0(x')\ge 2\Bigr\}.
\end{split}
\end{equation}
The \emph{coupling-smallness condition} is
\begin{equation}\label{8.10:eq-cost-polynomial-3}
x\ge x_\tau ,
\end{equation}
for the inverse coupling $x$ of Notation~\ref{8.10:not-terminal-couplings}. The
threshold $x_\tau$ is defined by the infimum \eqref{8.10:eq-cost-polynomial-2}
alone; no value is assigned to it. We claim that, under the degree condition,
$x_\tau<\infty$, and that \eqref{8.10:eq-cost-polynomial-3} holds whenever the
reference inverse coupling $X=g^{-2}$ satisfies $X\ge x_\tau$.
\end{definition}

\begin{proof}
The term $C_G'c_H\,d\,(2p_1+1)^d\log x$ occurs with positive coefficient in the
expansion of \eqref{8.10:eq-cost-polynomial-1}, so the top exponent of
$\mathfrak{c}$, which is at most $d\,r_{\mathrm{pr}}+1$, is at least $1$. By the degree
condition $d\,r_{\mathrm{pr}}+1<2\mathfrak{p}_0$ assumed above, $2\mathfrak{p}_0>1$, so the
exponent $\mathfrak{p}_0$ is positive. Expanding the square,
\begin{equation*}
\tfrac12\,\theta_\tau'\gamma_0A_1^2\bigl(p_0(x')-1\bigr)^2
=\tfrac12\,\theta_\tau'\gamma_0A_1^2\bigl(A_0^2(\log x')^{2\mathfrak{p}_0}
-2A_0(\log x')^{\mathfrak{p}_0}+1\bigr),
\end{equation*}
a combination of powers of $\log x'$ with top exponent $2\mathfrak{p}_0$ and
positive leading coefficient $\tfrac12\theta_\tau'\gamma_0A_1^2A_0^2$. By the
degree condition, this top exponent exceeds $d\,r_{\mathrm{pr}}+1$, which bounds the
top exponent of $\mathfrak{c}$ in \eqref{8.10:eq-cost-polynomial-1}. For real
exponents $a<b$ one has $(\log x')^a/(\log x')^b\to0$ as $x'\to\infty$. Dividing both
sides of the first inequality in \eqref{8.10:eq-cost-polynomial-2} by
$(\log x')^{2\mathfrak{p}_0}$, the right-hand side tends to
$\tfrac12\theta_\tau'\gamma_0A_1^2A_0^2>0$, because $\mathfrak{p}_0>0$, while
$\mathfrak{c}(x')(\log x')^{-2\mathfrak{p}_0}\to0$, because each term of
\eqref{8.10:eq-cost-polynomial-1} is bounded, for $x'\ge e$, by a constant times
$(\log x')^{d\,r_{\mathrm{pr}}+1}$, the top exponent of $\mathfrak{c}$ being at most
$d\,r_{\mathrm{pr}}+1$, and $d\,r_{\mathrm{pr}}+1<2\mathfrak{p}_0$. There is therefore an
$x_1\ge e$ such that, for all $x'>x_1$,
\begin{equation*}
\mathfrak{c}(x')(\log x')^{-2\mathfrak{p}_0}
<\tfrac14\,\theta_\tau'\gamma_0A_1^2A_0^2
<\tfrac12\,\theta_\tau'\gamma_0A_1^2\bigl(p_0(x')-1\bigr)^2(\log x')^{-2\mathfrak{p}_0};
\end{equation*}
hence the first inequality in
\eqref{8.10:eq-cost-polynomial-2} holds for all $x'>x_1$.

Since $A_0>0$ and $\mathfrak{p}_0>0$, the function $p_0(x')=A_0(\log x')^{\mathfrak{p}_0}$
tends to $\infty$ as $x'\to\infty$. Enlarging $x_1$ if necessary, we may also
assume $p_0(x')\ge2$ for all $x'>x_1$. Thus the set of $x'\ge e$ at which either
condition in \eqref{8.10:eq-cost-polynomial-2} fails is contained in $[e,x_1]$, so
it is bounded.

Every $x_0>x_1$ therefore belongs to the set whose infimum defines $x_\tau$. That
set is non-empty, and $x_\tau\le x_1<\infty$.

Finally, suppose $X\ge x_\tau$. Clause~(0) of Theorem~\ref{3.1:thm-main} gives $x\ge X$,
hence $x\ge X\ge x_\tau$. This is \eqref{8.10:eq-cost-polynomial-3}.
\end{proof}

\begin{lemma}[The budget inequality]\label{8.10:lem-budget-inequality}
Let $C$ be as in the old-member statement, Proposition~\ref{8.10:prop-old-member}, whose proof
is incomplete at one step, with $L_C\neq\emptyset$. Put
$p_0'=p_0(x)-1$, where $p_0(\cdot)$ is the degree function. Assume the following.
\begin{itemize}
\item[(a)] For every level $j\le K'$,
$p_0(g_j)\ge p_0(x)-1$.
\item[(b)] For every history $\mathfrak{h}\in\mathfrak{H}(C)$ of $C$, its seeds
$\zeta\in\mathcal{Z}^1(\mathfrak{h})$ satisfy
\begin{equation*}
\sum_{\zeta\in\mathcal{Z}^1(\mathfrak{h})}\bigl(\mathfrak{d}(\zeta)+1\bigr)\ \ge\ \frac{n(C)}{C_G'}.
\end{equation*}
\item[(c)] The exponent $\delta_C$ has the form assumed for it among the hypotheses of
Proposition~\ref{8.10:prop-old-member}, so that $\Gamma_C(\cdot)$ is the function defined
there.
\item[(d)] The coupling-smallness condition $x\ge x_\tau$ of
Definition~\ref{8.10:def-cost-polynomial} holds.
\end{itemize}
Then for every $\mathfrak{h}\in\mathfrak{H}(C)$
\begin{equation}\label{8.10:eq-budget-inequality-1}
\Gamma_C\bigl(\mathfrak{E}_*(C)\bigr)+(2d\log3+\log2)\,n(C)+(2-\theta)
\ \le\ \tfrac34 B^\flat_\tau(C)\ \le\ \tfrac34 B_\tau(C,\mathfrak{h}).
\end{equation}
Consequently, for every non-empty $\mathcal{S}\subset\mathfrak{H}(C)$,
\begin{multline}\label{8.10:eq-budget-inequality-2}
\operatorname{pen}^\tau_\theta(C,\mathcal{S})\,e^{\Gamma_C(\mathfrak{E}_*(C))}
\le w_\theta(\mathcal{S})\,\exp\Bigl(-(2-\theta)(\underline{p}_0+1)-\kappa_1 d_{K'}(C)\\
-(2d\log3+\log2)\,n(C)-\tfrac14 B^\flat_\tau(C)\Bigr).
\end{multline}
\end{lemma}

\begin{proof}
Fix $C$ as in the statement and $\mathfrak{h}\in\mathfrak{H}(C)$. The quantities
\begin{equation*}
B^\flat_\tau(C)=\frac{\theta'_\tau\gamma_0A_1^2p_0'^2\,n(C)}{C_G'},
\qquad
\mathfrak{E}_*(C)=\frac{B^\flat_\tau(C)}{4c_Ax}
\end{equation*}
are those of Definition~\ref{8.10:def-core-bound}.
The function $\Gamma_C(\cdot)$ is that of the old-member statement,
Proposition~\ref{8.10:prop-old-member}, whose proof is incomplete at one step.

\emph{The lower bound on $B_\tau(C,\mathfrak{h})$.}
By hypothesis~(d) the two inequalities of Definition~\ref{8.10:def-cost-polynomial}
hold at $x$. In particular $p_0(x)\ge2$, so $p_0'=p_0(x)-1\ge1>0$.

Every seed $\zeta\in\mathcal{Z}^1(\mathfrak{h})$ has creation level $j(\zeta)\le K'$.
Hypothesis~(a) therefore applies with $j=j(\zeta)$ and gives
\begin{equation*}
p_0\bigl(g_{j(\zeta)}\bigr)\ \ge\ p_0'\ >\ 0 .
\end{equation*}
Both sides are positive, so $p_0(g_{j(\zeta)})^2\ge p_0'^2$ for every such seed.
We insert this lower bound, seed by seed, into the definition of $B_\tau(C,\mathfrak{h})$.
Then we apply hypothesis~(b). This gives
\begin{equation}\label{8.10:eq-budget-inequality-3}
B_\tau(C,\mathfrak{h})
\ \ge\ \theta'_\tau\gamma_0A_1^2p_0'^2\sum_{\zeta\in\mathcal{Z}^1(\mathfrak{h})}
\bigl(\mathfrak{d}(\zeta)+1\bigr)
\ \ge\ \frac{\theta'_\tau\gamma_0A_1^2p_0'^2\,n(C)}{C_G'}
\ =\ B^\flat_\tau(C).
\end{equation}
Multiplying \eqref{8.10:eq-budget-inequality-3} by $\tfrac34$ gives the right-hand
inequality of \eqref{8.10:eq-budget-inequality-1}.

Now let $\mathcal{S}\subset\mathfrak{H}(C)$ be non-empty. The bound
\eqref{8.10:eq-budget-inequality-3} holds for every $\mathfrak{h}\in\mathcal{S}$, and
$B^\flat_\tau(C)$ does not depend on $\mathfrak{h}$. Hence the factor
$r_{\theta'_\tau}(\mathcal{S})$ entering $\operatorname{pen}^\tau_\theta(C,\mathcal{S})$ in
Definition~\ref{8.10:def-fibre-domination} satisfies
\begin{equation}\label{8.10:eq-budget-inequality-4}
r_{\theta'_\tau}(\mathcal{S})\ \le\ e^{-B^\flat_\tau(C)}.
\end{equation}

\emph{The upper bound on the cost.}
We use the following four facts:
\begin{itemize}
\item the sizes of the hull complexes given by Lemma~\ref{8.10:lem-hull-sizes};
\item the bounds $\check{R}_{K'}\le L(\log x)^{r_{\mathrm{pr}}}$ and
$\alpha\le\alpha_1(1+\log x)^{a_\alpha}$ on the two parameters $\check{R}_{K'}$ and
$\alpha$ that enter $\Gamma_C$;
\item the inequality $n(C)\ge1$;
\item the identity $c_Ax\,\mathfrak{E}_*(C)=\tfrac14 B^\flat_\tau(C)$, which is the
definition of $\mathfrak{E}_*(C)$ multiplied by $c_Ax$.
\end{itemize}
Insert these into the expression for $\Gamma_C(\mathfrak{E})$ that hypothesis~(c) provides, at
$\mathfrak{E}=\mathfrak{E}_*(C)$. Every term other than the term
$c_Ax\,\mathfrak{E}$, together with $(2d\log3+\log2)\,n(C)+(2-\theta)$, is then bounded
by $n(C)\,\mathfrak{c}(x)/C_G'$. Here $\mathfrak{c}$ is the cost polynomial of
Definition~\ref{8.10:def-cost-polynomial}, and $n(C)\ge1$ is used to bound the constant
$2-\theta$. The term $c_Ax\,\mathfrak{E}$ of $\Gamma_C(\mathfrak{E})$ equals
$\tfrac14 B^\flat_\tau(C)$ at $\mathfrak{E}=\mathfrak{E}_*(C)$. Hence
\begin{equation*}
\Gamma_C\bigl(\mathfrak{E}_*(C)\bigr)+(2d\log3+\log2)\,n(C)+(2-\theta)
\ \le\ \frac{n(C)\,\mathfrak{c}(x)}{C_G'}+\tfrac14 B^\flat_\tau(C).
\end{equation*}
By hypothesis~(d), $\mathfrak{c}(x)\le\tfrac12\theta'_\tau\gamma_0A_1^2(p_0(x)-1)^2$,
that is, $\mathfrak{c}(x)\le\tfrac12\theta'_\tau\gamma_0A_1^2p_0'^2$. Therefore
\begin{equation*}
\frac{n(C)\,\mathfrak{c}(x)}{C_G'}
\ \le\ \frac12\,\frac{\theta'_\tau\gamma_0A_1^2p_0'^2\,n(C)}{C_G'}
\ =\ \tfrac12 B^\flat_\tau(C).
\end{equation*}
Combining the last two displays gives
\begin{equation}\label{8.10:eq-budget-inequality-5}
\Gamma_C\bigl(\mathfrak{E}_*(C)\bigr)+(2d\log3+\log2)\,n(C)+(2-\theta)
\ \le\ \tfrac12 B^\flat_\tau(C)+\tfrac14 B^\flat_\tau(C)
\ =\ \tfrac34 B^\flat_\tau(C).
\end{equation}
This is the left-hand inequality of \eqref{8.10:eq-budget-inequality-1}.

\emph{The consequence.}
Rearranging \eqref{8.10:eq-budget-inequality-5} gives
\begin{equation*}
e^{\Gamma_C(\mathfrak{E}_*(C))}
\ \le\ \exp\Bigl(\tfrac34 B^\flat_\tau(C)-(2d\log3+\log2)\,n(C)-(2-\theta)\Bigr).
\end{equation*}
By the definition of $\operatorname{pen}^\tau_\theta(C,\mathcal{S})$ in
Definition~\ref{8.10:def-fibre-domination}, the factor
$r_{\theta'_\tau}(\mathcal{S})$ multiplies $w_\theta(\mathcal{S})$ and the factor
$\exp\bigl(-(2-\theta)\underline{p}_0-\kappa_1d_{K'}(C)\bigr)$. We bound
$r_{\theta'_\tau}(\mathcal{S})$ by \eqref{8.10:eq-budget-inequality-4} and multiply by
the last display. The factor $w_\theta(\mathcal{S})$, the term
$-(2-\theta)\underline{p}_0$ and the term $-\kappa_1d_{K'}(C)$ are left unchanged.
The exponents of $B^\flat_\tau(C)$ combine as
$-B^\flat_\tau(C)+\tfrac34B^\flat_\tau(C)=-\tfrac14B^\flat_\tau(C)$. The term
$-(2-\theta)$ coming from \eqref{8.10:eq-budget-inequality-5} joins
$-(2-\theta)\underline{p}_0$ to give $-(2-\theta)(\underline{p}_0+1)$. The result is
\eqref{8.10:eq-budget-inequality-2}.
\end{proof}

\begin{remark}[Charging the costs to the linear creation reserve]\label{8.10:rem-reserve-allocation}
Let $C$ be as in the budget inequality (Lemma~\ref{8.10:lem-budget-inequality}), and let
$\mathfrak{h}$ be a history of $C$. To each seed $\eta$ of $\mathfrak{h}$ are attached its
creation step $j(\eta)$, the coupling $g_{j(\eta)}$ at that step, and a quantity
$\mathfrak{d}(\eta)$ entering its linear factor below; to $C$ is attached a step $j(C)$. The
seeds are either of the first kind or preparatory. The proof of
clause~(i) of Theorem~\ref{3.1:thm-main} leaves the following four factors
unspent, per history $\mathfrak{h}$ of $C$.
\begin{itemize}
\item[(R1)] The \emph{seed factor}, supplied by the seed-reserve statement of the proof of
clause~(i): for every seed $\eta$, of either kind, the factor
\begin{equation*}
e^{-\theta p_0(g_{j(\eta)})},
\end{equation*}
where $p_0(\cdot)$ is the degree function.
\item[(R2)] The \emph{linear factor}, supplied by the linear creation reserve statement of the
proof of clause~(i): for every seed $\eta$ of the first kind only, the factor
\begin{equation*}
e^{-\theta'\gamma_0A_1^2\,p_0(g_{j(\eta)})^2\,(\mathfrak{d}(\eta)+1)};
\end{equation*}
here $\gamma_0$ and $\theta'$ are the constants of the linear creation reserve and $A_1$ is the
auxiliary parameter of that name. A preparatory seed carries no linear factor, only its unit
of flat reserve. The coefficient $\theta'$ has the parts $\theta'_\tau$, $\theta'_A$ and
$\theta'_{\mathrm{age}}$. Of the linear factor, only the part $\theta'_\tau$ is used in this
remark.
\item[(R3)] The \emph{flat factor}:
\begin{equation*}
e^{-(2-\theta)p_0(g_{j(C)})}\le e^{-(2-\theta)\underline{p}_0},
\end{equation*}
where $\underline{p}_0$ is the lower bound for the degree function that appears in
Lemma~\ref{8.10:lem-budget-inequality}.
\item[(R4)] The \emph{distance factor}: $e^{-\kappa_1 d_{K'}(C)}$, where $d_{K'}(C)$ is the
quantity of that name in the exponent of the consequence stated in
Lemma~\ref{8.10:lem-budget-inequality}.
\end{itemize}
The costs of the old-member statement (Proposition~\ref{8.10:prop-old-member}, stated; proof
incomplete at the step marked $(\ast)$) are as follows; the constants $c_H$, $C_L$, the sets
$\mathsf{R}_1$, $\mathsf{N}_0$ and the variable $x$ are those of the exponent $\Gamma_C$ of
that proposition.
\begin{itemize}
\item[(c1)] The top-level mismatch between Haar and Gaussian measures, $c_H|\mathsf{R}_1|\log x$,
and the losses of order one per site, $C_L|\mathsf{N}_0|$, coming from window margins,
Jacobians, the Hessian and the matching of the exterior factor.
\item[(c2)] The filling action $c_{Ax}\mathfrak{E}_*(C)=\tfrac14 B^\flat_\tau(C)$, where
$\mathfrak{E}_*(C)$ is the allotment of Definition~\ref{8.10:def-core-bound}.
\item[(c3)] The exponent $\delta_C$ of the factor $e^{\delta_C}$ in
Proposition~\ref{8.10:prop-old-member}, in the form fixed by the hypothesis of that
proposition that specifies $\delta_C$.
\end{itemize}
These costs, together with the shape entropy $(2d\log3+\log2)\,n(C)$ and the additive constant
$(2-\theta)$ that enter the consequence stated in Lemma~\ref{8.10:lem-budget-inequality}, are
all charged to the part $\theta'_\tau$ of the linear factor (R2), through the core bound stated
in Definition~\ref{8.10:def-core-bound}; the proof of that bound is outstanding.
This leaves $\tfrac14B^\flat_\tau(C)$ of that part unspent, and that quarter remains available
to the typicality estimate stated later in this chapter.

The factors (R1), (R3), (R4) and the parts $\theta'_A$, $\theta'_{\mathrm{age}}$ of (R2) pay
for none of these costs; (R3) and (R4) pass unchanged into the consequence of
Lemma~\ref{8.10:lem-budget-inequality}. Consequently $\theta$ is unchanged, and the weight
$w_\theta$ of Lemma~\ref{8.10:lem-budget-inequality} appears in that lemma's consequence exactly,
with exponent $\theta$. Nothing is added to the removal constant $C_{\mathrm{rm}}$, and no
factor is used twice.

Under the funding by (R2), the degrees to be compared, in the logarithm of the inverse square
coupling, are $d\,r_{\mathrm{pr}}+1$ on the side of the removal costs, $r_{\mathrm{pr}}$ being
the exponent of the chapter's degree comparison, and $2p_0$ on the side of the linear factor,
where $p_0$ is the fixed exponent of the degree function. The comparison is strict in every
regime of exponents considered:
\begin{equation*}
d\,r_{\mathrm{pr}}+1<2p_0 .
\end{equation*}
\end{remark}

\begin{proof}
We first establish the identity in (c2). By Definition~\ref{8.10:def-core-bound},
$\mathfrak{E}_*(C)=B^\flat_\tau(C)/(4c_{Ax})$. Multiplying by $c_{Ax}$ gives
\begin{equation}\label{8.10:eq-reserve-allocation-1}
c_{Ax}\,\mathfrak{E}_*(C)=\tfrac14\,B^\flat_\tau(C).
\end{equation}

Lemma~\ref{8.10:lem-budget-inequality} gives
\begin{equation}\label{8.10:eq-reserve-allocation-2}
\begin{aligned}
\Gamma_C(\mathfrak{E}_*(C))+(2d\log3+\log2)\,n(C)+(2-\theta)
&\le \tfrac34\,B^\flat_\tau(C)\\
&\le \tfrac34\,B_\tau(C,\mathfrak{h}).
\end{aligned}
\end{equation}
Here $B_\tau(C,\mathfrak{h})$ is the exponent furnished by the part $\theta'_\tau$ of the linear
factor (R2) along $\mathfrak{h}$, and $B^\flat_\tau(C)$ is its lower bound independent of
$\mathfrak{h}$, given by the core bound stated in Definition~\ref{8.10:def-core-bound} (an input
of that lemma; the proof of the bound is outstanding). In $\Gamma_C(\mathfrak{E}_*(C))$ the
old-member statement collects the costs (c1), (c2) with the value
\eqref{8.10:eq-reserve-allocation-1}, and (c3). By \eqref{8.10:eq-reserve-allocation-2}, these
costs, the shape entropy and the additive constant together use at most three quarters of
$B^\flat_\tau(C)$. Hence $\tfrac14B^\flat_\tau(C)$ is left unspent, which is the allocation
asserted.

Finally, since $p_0(g')=A_0(\log g'^{-2})^{p_0}$, the exponent of (R2), which is quadratic in
$p_0(g_{j(\eta)})$, has degree $2p_0$ in $\log g_{j(\eta)}^{-2}$; against it the removal costs
enter with degree $d\,r_{\mathrm{pr}}+1$, and the comparison $d\,r_{\mathrm{pr}}+1<2p_0$ is
strict in every regime of exponents considered.
\end{proof}

\begin{remark}[The distance factor together with the linear reserve. Stated; proof incomplete at the step marked $(\ast)$]\label{8.10:rem-distance-factor}
Let $\theta'$ be the fraction of the linear reserve that is spent; it is fixed in the allotment of the linear reserve in Definition~\ref{8.10:def-core-bound} (printed there with proof outstanding), where it carries an index $\tau$ that we suppress. Using the linear reserve with fraction $\theta'$ means the following: for the coefficient $\gamma_0$ of the recursion, the $\kappa$-recursion of \cite[(1.80)--(1.88)]{B-CMP122b} is run on the coefficient
\begin{equation*}
(1-\theta')\gamma_0 .
\end{equation*}

For a live region $X$ at level $k$, write $d_k(X)$ for Ba{\l}aban's distance function of $X$ at that level. The assertion is that this use of the reserve does not remove the clause of the basic bound on a live large-field region that carries the additional factor
\begin{equation*}
e^{-\kappa_1 d_k(X)}
\end{equation*}
for a general live region $X$. That is, the bound with coefficient $(1-\theta')\gamma_0$ retains the factor $e^{-\kappa_1 d_k(X)}$ for a general live region.

Stated; the proof below is incomplete at the step marked $(\ast)$.

The old-member statement does not use, for its region $C$, the distance factor $e^{-\kappa_1 d_{K'}(C)}$, where $d_{K'}(C)$ is the distance function of $C$ at the level $K'$ at which that statement considers $C$. The present remark is therefore not used in its proof. It concerns only those uses of the bounds in which both factors are kept. The clause carrying $e^{-\kappa_1 d_k(X)}$ is itself used here in the form stated in \cite{B-CMP122b} for arbitrary regions.
\end{remark}

\begin{proof}
Run the $\kappa$-recursion of \cite[(1.80)--(1.88)]{B-CMP122b} on the coefficient $(1-\theta')\gamma_0$. The clause with the additional factor $e^{-\kappa_1 d_k(X)}$ is to be derived from two corollaries of this book; the first is used as a whole and the second through its item~(c).

$(\ast)$ One of these two corollaries is stated with a single scaling factor. The derivation of the clause in the present proof applies it with two unequal scaling factors; that a corollary stated with one scaling factor may be applied with two unequal ones is not established here.

Granting this application, the two corollaries give that the bound obtained from the recursion with coefficient $(1-\theta')\gamma_0$ still carries the additional factor $e^{-\kappa_1 d_k(X)}$ for a general live region $X$.
\end{proof}

\begin{proposition}[Non-fillable exteriors with flat collar. Stated; proof incomplete at the step marked $(\ast)$]\label{8.10:prop-non-fillable}
Let $p_1\ge 1$ and $p_2\ge 3$. Let $C$ be a touch-connected union of cells, and let $\lambda\in\Lambda_C$ satisfy $K_{\mathrm{free}}(C;\lambda;p_1+p_2+p_W)$. Suppose we are given the following data.
\begin{itemize}
\item An edge loop $\gamma$ all of whose edges are exterior.
\item A null-homotopy of $\gamma$ in $K(\mathbb{T})$ across plaquettes $(q_1,\dots,q_A)$, where $A\ge 1$.
\item An exterior configuration $V^{*}_{\mathrm{ext}}\in\mathcal{V}^{\mathrm{ext}}$ such that $F_p(V^{*}_{\mathrm{ext}})=0$ for every plaquette $p$ of $\mathbb{T}$ having no interior edge, and such that
\begin{equation*}
h_0:=\bigl\|T_\gamma(V^{*}_{\mathrm{ext}})-1\bigr\|>0 .
\end{equation*}
\end{itemize}
For any $E>0$ put
\begin{equation}\label{8.10:eq-non-fillable-1}
\eta_*:=\min\biggl(\frac{h_0}{8A},\ \frac{\varepsilon_\mu}{2},\ f_*,\
\Bigl(\frac{E}{|\mathsf{Ex}_2|+1}\Bigr)^{1/2}\biggr)>0,
\end{equation}
and let $\mathcal{Y}$ be the set of those $V^{\mathrm{ext}}\in\mathcal{V}^{\mathrm{ext}}$ that satisfy both
\begin{equation}\label{8.10:eq-non-fillable-2}
\bigl\|T_\gamma(V^{\mathrm{ext}})-1\bigr\|>\frac{h_0}{2}
\end{equation}
and $F_p(V^{\mathrm{ext}})<\eta_*$ for every plaquette $p$ of $\mathbb{T}$ with no interior edge. We use the two local integrals $a_\lambda(V^{\mathrm{ext}},w)$ and $m_\lambda(V^{\mathrm{ext}},w)$, $w\in\mathcal{U}_\lambda$, of Definition~\ref{8.10:def-admissible-exteriors}. Then the following hold.
\begin{enumerate}
\item[(a)] The set $\mathcal{Y}$ is open in $\mathcal{V}^{\mathrm{ext}}$, contains $V^{*}_{\mathrm{ext}}$, and has positive Haar measure. Moreover, for every $V^{\mathrm{ext}}\in\mathcal{Y}$ and every $V^{\mathrm{in}}$, some $q_i$ lies in $\mathsf{F}$, and
\begin{equation}\label{8.10:eq-non-fillable-3}
\max_{i:\,q_i\in\mathsf{F}}F_{q_i}(V^{\mathrm{in}},V^{\mathrm{ext}})>\frac{3h_0}{8A}.
\end{equation}
\item[(b)] Assume the outer window bound of Definition~\ref{8.10:def-window-structure}, whose proof is outstanding, and the format statement for the density $\rho_\lambda$. Assume also the one-sided floor
\begin{equation}\label{8.10:eq-non-fillable-4}
c_W\,\varepsilon\le\frac{3h_0}{8A}.
\end{equation}
Then $a_\lambda(V^{\mathrm{ext}},w)=0$ for every $V^{\mathrm{ext}}\in\mathcal{Y}$ and every $w\in\mathcal{U}_\lambda$. Also $\mathcal{Y}\cap\mathrm{Fill}_{\mu,\mathfrak{E}}(C)=\emptyset$ for every $\mathfrak{E}\ge 0$.
\item[(c)] The set $\mathcal{Y}$ satisfies
\begin{equation*}
\mathcal{Y}\subset G^{\mathrm{mar}}_{\mu}\cap G^{\mathrm{av}}_{f_*}\cap G^{\mathrm{en}}_{E}.
\end{equation*}
If, in addition, the inner window bound and the localisation of the windows of Definition~\ref{8.10:def-window-structure} hold, then $\mathcal{Y}\subset\mathrm{Adm}(\lambda;C)$ as well (Definition~\ref{8.10:def-admissible-exteriors}). Under these two hypotheses $\mathcal{Y}$ therefore lies inside every clause of the good set of Definition~\ref{8.10:def-good-set} except the filling clause. Unconditionally, it lies inside $G_{f_*}(C)$ of Definition~\ref{8.10:def-good-set}, and inside the collar-energy clause $G^{\mathrm{en}}_E$ for the chosen $E$.
\item[(d)] Assume the inner window bound of Definition~\ref{8.10:def-window-structure}, and the format statement for $M_{\lambda,C}$, which gives $M_{\lambda,C}=\bigl[\prod_{p\in\mathcal{P}_M}\chi_p\bigr]\cdot\widehat{M}$ with $\widehat{M}\ge0$. Suppose that $\lambda$ satisfies the positivity statement for $\widehat{M}$ (for instance, $\mathfrak{K}(\lambda)=\emptyset$), which yields, for every $w\in\mathcal{U}_\lambda$, $\widehat{M}(w,\cdot)>0$ for $dV$-almost every argument. Suppose also that $V^{*}_{\mathrm{ext}}$ admits an extension $V^{*}\in\mathcal{V}$ with $F_{p'}(V^{*})=0$ for every plaquette $p'$ contained in $N_{r_W+1}(p)$, the neighbourhood of radius $r_W+1$ of $p$, for every $p\in\mathcal{P}_M$. Then there is an open neighbourhood $\mathcal{N}_{\mathrm{ext}}$ of $V^{*}_{\mathrm{ext}}$ in $\mathcal{V}^{\mathrm{ext}}$ with the following property: for every $w\in\mathcal{U}_\lambda$, one has $m_\lambda(V^{\mathrm{ext}},w)>0$ for $dV^{\mathrm{ext}}$-almost every $V^{\mathrm{ext}}\in\mathcal{Y}\cap\mathcal{N}_{\mathrm{ext}}$. The set $\mathcal{Y}\cap\mathcal{N}_{\mathrm{ext}}$ has positive Haar measure. Positivity of $m_\lambda(V^{\mathrm{ext}},w)$ at every $V^{\mathrm{ext}}\in\mathcal{Y}$ is not proved here. The reason is that, for $V^{\mathrm{ext}}\in\mathcal{Y}$ far from $V^{*}_{\mathrm{ext}}$, the windows $\chi_p$, $p\in\mathcal{P}_M$, that straddle the interior and the exterior bonds depend on $V^{\mathrm{ext}}$ directly.
\item[(e)] The union $C$ is linked in the sense of Definition~\ref{8.10:def-linked-components}. Suppose moreover that $\gamma$ lies in $\mathsf{N}$ (all its vertices lie in $\mathsf{N}_0$) and that all $q_i$ lie in $\mathsf{N}_2$. Then, under the patching hypothesis, neither condition $(\mathrm{U}_1)$ nor condition $(\mathrm{U}_2)$ of Definition~\ref{8.10:def-linked-components} holds.
\end{enumerate}
\end{proposition}

\begin{proof}
Throughout, a configuration $V\in\mathcal{V}$ is written $V=(V^{\mathrm{in}},V^{\mathrm{ext}})$, where $\mathcal{V}=\mathcal{V}^{\mathrm{in}}\times\mathcal{V}^{\mathrm{ext}}$.

\emph{(a).} By Lemma~\ref{8.10:lem-collar-plaquettes}(c), $F_p$ is a function of $V^{\mathrm{ext}}$ for every plaquette $p$ without interior edge. Since all edges of $\gamma$ are exterior, $T_\gamma$ involves exterior bonds only, so it too is a function of $V^{\mathrm{ext}}$. Both are continuous functions of $V^{\mathrm{ext}}$. Hence $\mathcal{Y}$ is defined by finitely many strict inequalities between continuous functions of $V^{\mathrm{ext}}$, and is therefore open.

The point $V^{*}_{\mathrm{ext}}$ lies in $\mathcal{Y}$ for two reasons. First, $\|T_\gamma(V^{*}_{\mathrm{ext}})-1\|=h_0>h_0/2$ because $h_0>0$. Second, $F_p(V^{*}_{\mathrm{ext}})=0<\eta_*$ for every plaquette $p$ without interior edge. Here $\eta_*>0$ because each entry of the minimum in \eqref{8.10:eq-non-fillable-1} is positive. Since $\mathcal{Y}$ is a non-empty open subset of the compact group $\mathcal{V}^{\mathrm{ext}}$, it has positive Haar measure.

Now let $V^{\mathrm{ext}}\in\mathcal{Y}$, let $V^{\mathrm{in}}$ be arbitrary, and put $V=(V^{\mathrm{in}},V^{\mathrm{ext}})$. Lemma~\ref{8.10:lem-homotopy-moves} applies to the given null-homotopy of $\gamma$ across $(q_1,\dots,q_A)$. Together with \eqref{8.10:eq-non-fillable-2}, it gives
\begin{equation*}
\frac{h_0}{2}<\bigl\|T_\gamma(V)-1\bigr\|
\le\sum_{i:\,q_i\in\mathsf{F}}F_{q_i}(V)+\sum_{i:\,q_i\notin\mathsf{F}}F_{q_i}(V).
\end{equation*}
By Lemma~\ref{8.10:lem-collar-plaquettes}(a), every plaquette of $\mathbb{T}$ with an edge in $\mathsf{R}_1$ lies in $\mathsf{F}$. So a plaquette not in $\mathsf{F}$ has no interior edge.

Consider first an index $i$ with $q_i\notin\mathsf{F}$. Then $F_{q_i}(V)=F_{q_i}(V^{\mathrm{ext}})$, and this is $<\eta_*$ because $V^{\mathrm{ext}}\in\mathcal{Y}$; moreover $\eta_*\le h_0/(8A)$ by \eqref{8.10:eq-non-fillable-1}. There are at most $A$ such indices. If there is at least one, the second sum is $<A\cdot h_0/(8A)=h_0/8$. If there are none, the second sum is $0<h_0/8$. In either case,
\begin{equation*}
\sum_{i:\,q_i\in\mathsf{F}}F_{q_i}(V)>\frac{h_0}{2}-\frac{h_0}{8}=\frac{3h_0}{8}.
\end{equation*}
In particular the index set $\{i:q_i\in\mathsf{F}\}$ is non-empty, so some $q_i$ lies in $\mathsf{F}$. The sum has at most $A$ terms, so the largest of them exceeds $3h_0/(8A)$. This is \eqref{8.10:eq-non-fillable-3}.

\emph{(b).} Let $V^{\mathrm{ext}}\in\mathcal{Y}$, let $V^{\mathrm{in}}$ be arbitrary, and put $V=(V^{\mathrm{in}},V^{\mathrm{ext}})$. By~(a) there is a plaquette $q\in\mathsf{F}$ among the $q_i$ with $F_q(V)>3h_0/(8A)$. By the floor \eqref{8.10:eq-non-fillable-4}, $3h_0/(8A)\ge c_W\varepsilon$, so $F_q(V)>c_W\varepsilon$. The outer window bound of Definition~\ref{8.10:def-window-structure}, whose proof is outstanding, states that positivity of the window at a plaquette forces $F$ at that plaquette to be at most $c_W\varepsilon$. Read in contrapositive form, it therefore gives $\chi^\lambda_q(V)=0$.

Next we show that $q\in\mathcal{P}_\lambda$. The plaquette $q$ has an interior edge. The endpoints of that edge lie in
\begin{equation*}
\bar S(R_{\mathrm{int}})\subset S(\mathrm{hull}_1R_{\mathrm{int}})=S(\mathrm{hull}_{p_1+1}C).
\end{equation*}
The other vertices of $q$ are within sup-distance $1\le\ell$ of these endpoints. Lemma~\ref{8.10:lem-steps-cells} with $k=1$ therefore places them in $S(c'')$ for cells $c''$ at cell-distance at most $1$ from cells of $\mathrm{hull}_{p_1+1}C$. Hence all vertices of $q$ lie in $S(\mathrm{hull}_{p_1+2}C)$.

By $K_{\mathrm{free}}(C;\lambda;p_1+p_2+p_W)$, a cell of $\mathrm{hull}_{p_1+2}(C)$ is at cell-distance at least
\begin{equation*}
(p_1+p_2+p_W+1)-(p_1+2)=p_2+p_W-1>p_W
\end{equation*}
from $\mathfrak{K}(\lambda)$. The last inequality holds because $p_2\ge3$. The format statement for $\rho_\lambda$, which enters here as a hypothesis of~(b), then gives $q\in\mathcal{P}_\lambda$, and $\chi^\lambda_q$ is a factor of $\rho_\lambda$.

Since $\chi^\lambda_q(V)=0$, we get $\rho_\lambda(w,V^{\mathrm{in}},V^{\mathrm{ext}})=0$. The plaquette $q$ depends on $V^{\mathrm{in}}$, but one such plaquette exists for every $V^{\mathrm{in}}$. So this holds for all $V^{\mathrm{in}}$ and all $w\in\mathcal{U}_\lambda$, and consequently $a_\lambda(V^{\mathrm{ext}},w)=0$.

For the second claim, suppose $V^{\mathrm{ext}}\in\mathcal{Y}$ admitted a filling of margin $\mu$. On $\mathsf{F}$ it would satisfy
\begin{equation*}
F\le\varepsilon_\mu=\frac{(1-\mu)\varepsilon}{B_W}\le\varepsilon\le c_W\varepsilon\le\frac{3h_0}{8A},
\end{equation*}
where the two middle inequalities are $(1-\mu)/B_W\le1$ and $c_W\ge1$, the latter being part of the outer window bound of Definition~\ref{8.10:def-window-structure}, whose proof is outstanding, and the last step is \eqref{8.10:eq-non-fillable-4}. This contradicts \eqref{8.10:eq-non-fillable-3}. Hence $\mathcal{Y}\cap\mathrm{Fill}_{\mu,\mathfrak{E}}(C)=\emptyset$ for every $\mathfrak{E}\ge0$.

\emph{(c).} Let $V^{\mathrm{ext}}\in\mathcal{Y}$. Every $p'\in\mathsf{Ex}_2$ has no interior edge, so $F_{p'}(V^{\mathrm{ext}})<\eta_*$. Using the entries of the minimum \eqref{8.10:eq-non-fillable-1} one at a time, we obtain three consequences.
\begin{itemize}
\item Since $\eta_*\le\varepsilon_\mu/2$, we have $F<\varepsilon_\mu/2$ on $\mathsf{Ex}_2^{\mathrm{near}}$. This is the clause $G^{\mathrm{mar}}_\mu$.
\item Since $\eta_*^2\le E/(|\mathsf{Ex}_2|+1)$, we have
\begin{equation*}
\mathsf{Q}_{\mathrm{col}}<|\mathsf{Ex}_2|\,\eta_*^2\le E .
\end{equation*}
This is the clause $G^{\mathrm{en}}_E$.
\item Since $P_{\mathfrak{a},\mu'\nu'}\subset\mathsf{Ex}_2$ and $\eta_*\le f_*$, we have
\begin{equation*}
f_{\mathfrak{a},\mu'\nu'}\le\max_{P_{\mathfrak{a},\mu'\nu'}}F<\eta_*\le f_* .
\end{equation*}
This is the clause $G^{\mathrm{av}}_{f_*}$.
\end{itemize}

For the admissibility clause, let $p\in\mathcal{P}^{\mathrm{held}}_\lambda$. Every bond $b\subset N_{r_W+1}(p)$ is exterior, so every plaquette $p'\subset N_{r_W+1}(p)$ has no interior edge, and
\begin{equation*}
F_{p'}(V^{\mathrm{ext}})<\eta_*\le\frac{\varepsilon_\mu}{2}\le\frac{\varepsilon}{B_W}.
\end{equation*}
By the localisation of the windows of Definition~\ref{8.10:def-window-structure}, whose proof is outstanding and which enters here as a hypothesis of~(c), $\chi^\lambda_p$ is a function of the bonds of $N_{r_W+1}(p)$, all of which are exterior, hence a function of the exterior configuration $V^{\mathrm{ext}}$ alone. The inner window bound of Definition~\ref{8.10:def-window-structure}, whose proof is outstanding, applied with $\delta:=\varepsilon/B_W$, gives $\chi^\lambda_p(V^{\mathrm{ext}})=1>0$. As this holds for every $p\in\mathcal{P}^{\mathrm{held}}_\lambda$, we conclude $V^{\mathrm{ext}}\in\mathrm{Adm}(\lambda;C)$.

\emph{(d).} By the format statement for $M_{\lambda,C}$, which is a hypothesis of~(d),
\begin{equation*}
M_{\lambda,C}=\Bigl[\prod_{p\in\mathcal{P}_M}\chi_p\Bigr]\cdot\widehat{M},\qquad \widehat{M}\ge0 .
\end{equation*}
Let $\mathcal{N}$ be the set of $V\in\mathcal{V}$ with
\begin{equation*}
F_{p'}(V)<\frac{\varepsilon}{B_W}\quad\text{for every plaquette } p'\subset N_{r_W+1}(p),\ p\in\mathcal{P}_M .
\end{equation*}
This set is open, being defined by finitely many strict inequalities between continuous functions. It contains $V^{*}$, where these $F$ vanish. On $\mathcal{N}$, every $\chi_p$ with $p\in\mathcal{P}_M$ equals $1$, by the inner window bound of Definition~\ref{8.10:def-window-structure}, whose proof is outstanding, applied with $\delta:=\varepsilon/B_W$.

In the product topology of $\mathcal{V}=\mathcal{V}^{\mathrm{in}}\times\mathcal{V}^{\mathrm{ext}}$, the set $\mathcal{N}$ contains a product $\mathcal{N}_{\mathrm{in}}\times\mathcal{N}_{\mathrm{ext}}$ of open neighbourhoods of $V^{*}|_{\mathsf{R}_1}$ and of $V^{*}_{\mathrm{ext}}$. For $V^{\mathrm{ext}}\in\mathcal{N}_{\mathrm{ext}}$ and every $w\in\mathcal{U}_\lambda$,
\begin{equation*}
m_\lambda(V^{\mathrm{ext}},w)\ge\int_{\mathcal{N}_{\mathrm{in}}}\widehat{M}(w,V^{\mathrm{in}},V^{\mathrm{ext}})\,dV^{\mathrm{in}} .
\end{equation*}

By the positivity statement for $\widehat{M}$, which is a hypothesis of~(d), $\widehat{M}(w,\cdot)>0$ for $dV$-almost every argument. By Fubini's theorem, for $dV^{\mathrm{ext}}$-almost every $V^{\mathrm{ext}}$ the function $\widehat{M}(w,\cdot,V^{\mathrm{ext}})$ is positive $dV^{\mathrm{in}}$-almost everywhere. For such $V^{\mathrm{ext}}\in\mathcal{N}_{\mathrm{ext}}$, the integral over the non-empty open set $\mathcal{N}_{\mathrm{in}}$ is positive, because $\mathcal{N}_{\mathrm{in}}$ has positive $dV^{\mathrm{in}}$-measure. Hence $m_\lambda(V^{\mathrm{ext}},w)>0$ for $dV^{\mathrm{ext}}$-almost every $V^{\mathrm{ext}}\in\mathcal{Y}\cap\mathcal{N}_{\mathrm{ext}}$.

Finally, $\mathcal{Y}\cap\mathcal{N}_{\mathrm{ext}}$ is open, as the intersection of $\mathcal{Y}$, open by~(a), with the open set $\mathcal{N}_{\mathrm{ext}}$, and it contains $V^{*}_{\mathrm{ext}}$. Being a non-empty open subset of the compact group $\mathcal{V}^{\mathrm{ext}}$, it has positive Haar measure.

\emph{(e).} Four facts hold together:
\begin{itemize}
\item $V^{*}_{\mathrm{ext}}$ is flat on every plaquette without interior edge;
\item $\gamma$ is exterior;
\item $\gamma$ is null-homotopic in $K(\mathbb{T})$;
\item $T_\gamma(V^{*}_{\mathrm{ext}})\ne1$, since $h_0>0$.
\end{itemize}
By Definition~\ref{8.10:def-linked-components}, $C$ is linked.

Now suppose that $\gamma$ lies in $\mathsf{N}$. Its edges are exterior bonds of $\mathsf{N}$. By Lemma~\ref{8.10:lem-collar-plaquettes}(b), they are therefore edges of $\mathsf{Ex}$. The configuration $V^{*}_{\mathrm{ext}}$ is flat on $\mathsf{Ex}_2$, because the plaquettes of $\mathsf{Ex}_2$ have no interior edge. If $\gamma$ were null-homotopic in $\mathsf{Ex}$, Lemma~\ref{8.10:lem-homotopy-moves} would give $T_\gamma(V^{*}_{\mathrm{ext}})=1$, contradicting $h_0>0$. So $\gamma$ is not null-homotopic in $\mathsf{Ex}$, and $(\mathrm{U}_1)$ fails.

$(\ast)$ If all $q_i$ lie in $\mathsf{N}_2$, then $\gamma$ is null-homotopic in $\mathsf{N}$; the passage from the given null-homotopy of $\gamma$ across $(q_1,\dots,q_A)$, all of which lie in $\mathsf{N}_2$, to a null-homotopy of $\gamma$ in $\mathsf{N}$ is the step not carried out here. Since $\gamma$ is not null-homotopic in $\mathsf{Ex}$, the injectivity in $(\mathrm{U}_2)$ fails.
\end{proof}

\begin{proposition}[Instances from a non-zero obstruction group]\label{8.10:prop-obstruction-instances}
Let $p_1\ge 1$ and $p_2\ge 3$. Let $C$ be a touch-connected union of cells such that the patch
condition holds for $\operatorname{hull}_{p_1+p_2+1}(C)$, and suppose that
$\operatorname{lk}^{\sharp}(C)\ge 1$. Let $\mathbb{B}$ be the spanned box and $\mathbb{E}_B$ the
exterior subcomplex of the closed-cube split, as in Definition~\ref{8.10:def-spanned-box}. Then
the following hold.
\begin{itemize}
\item[(i)] There are a closed real $1$-cochain $\xi$ on $\mathbb{E}_B$ that is not a coboundary,
an edge loop $\gamma$ in the graph of $\mathbb{E}_B$ (so that every edge of $\gamma$ is
exterior) with $\langle \xi,\gamma\rangle=1$, and a null-homotopy of $\gamma$ in $\mathbb{B}$
across plaquettes $q_1,\dots,q_A$ of $\mathbb{B}$, for some $A\ge 1$.
\item[(ii)] There is $V^{\flat}\in\mathcal{V}^{\mathrm{ext}}$ (here and below in the notation of
Proposition~\ref{8.10:prop-non-fillable}) with $F_p(V^{\flat})=0$ for every
plaquette $p$ of $\mathbb{T}$ without interior edge and
$T_\gamma(V^{\flat})=\operatorname{diag}(-1,-1,1,\dots,1)$; in particular
$h_0:=\|T_\gamma(V^{\flat})-1\|=2$.
\item[(iii)] The data $\gamma$, $(q_1,\dots,q_A)$ and
$V^{*}_{\mathrm{ext}}:=V^{\flat}$ satisfy the hypotheses of
Proposition~\ref{8.10:prop-non-fillable} with $h_0=2$. Assume in addition the outer window bound,
the inner window bound, the localisation of the windows and the format condition in $\lambda$, in
the form in which these enter Proposition~\ref{8.10:prop-non-fillable}, and assume
Proposition~\ref{8.10:prop-non-fillable} (stated; proof incomplete at the step marked $(\ast)$).
Then, applied to the data of (i) and (ii), it gives the following:
for every $E>0$ and every
$\lambda\in\Lambda_C$ satisfying $K_{\mathrm{free}}(C;\lambda;p_1+p_2+p_W)$, the open set $U$ of
Proposition~\ref{8.10:prop-non-fillable} formed with these data has positive measure and
satisfies
\begin{equation*}
U\subset \mathrm{Adm}\cap G^{\mathrm{mar}}_{\mu}\cap G^{\mathrm{av}}_{f_*}\cap
G^{\mathrm{en}}_{E},
\qquad
U\cap \mathrm{Fill}_{\mu,\mathfrak{E}}(C)=\emptyset\ \text{ for every }\mathfrak{E};
\end{equation*}
and if moreover $c_W\varepsilon\le 3/(4A)$, that is $x\ge x_{\mathrm{hol}}(C)$ with
\begin{equation}\label{8.10:eq-obstruction-instances-1}
x_{\mathrm{hol}}(C):=\inf\Bigl\{x_0\ge e:\ \tfrac43\,c_W A\,A_0(\log x')^{p_0}\le x'^{1/2}
\ \text{ for all } x'\ge x_0\Bigr\}<\infty ,
\end{equation}
then $a_\lambda\equiv 0$ on $U$. Here $c_W$ is the constant of the outer and inner window
bounds, and $x$ and $\varepsilon$ are the coupling variable and the parameter so denoted in
Proposition~\ref{8.10:prop-non-fillable}.
\end{itemize}
The floor $x_{\mathrm{hol}}(C)$ depends on $C$; such a $C$-dependent coupling floor is legitimate
when the proposition is used to exhibit a counterexample.
\end{proposition}

\begin{proof}
(i) By Lemma~\ref{8.10:lem-obstruction-ranks}(b), which applies because the patch condition
holds for $\operatorname{hull}_{p_1+p_2+1}(C)$, the real cohomology group
$H^1(\mathbb{E}_B;\mathbb{R})$ has dimension $\operatorname{lk}^{\sharp}(C)\ge 1$. A non-zero
class in it is represented by a closed real $1$-cochain $\xi$ on $\mathbb{E}_B$ that is not a
coboundary. By the facts collected in Notation~\ref{8.10:not-topology-facts}, a closed real
$1$-cochain on a cubical complex that is not a coboundary has non-zero pairing with some edge loop
of that complex; let $\gamma$ be an edge loop of $\mathbb{E}_B$ with
$\langle\xi,\gamma\rangle\neq 0$. Replacing $\xi$ by $\xi/\langle\xi,\gamma\rangle$, which is
again closed and again not a coboundary, we obtain $\langle\xi,\gamma\rangle=1$. The edges of
$\mathbb{E}_B$ are exterior by Definition~\ref{8.10:def-spanned-box}, so every edge of $\gamma$
is exterior.

The realisation $|\mathbb{B}|$ is a box, hence contractible. By the facts collected in
Notation~\ref{8.10:not-topology-facts}, an edge loop in a cubical complex with contractible
realisation admits a null-homotopy in that complex across finitely many of its plaquettes; so
$\gamma$ admits a null-homotopy in $\mathbb{B}$ across finitely many plaquettes
$q_1,\dots,q_A$ of $\mathbb{B}$. The number $A$ is at least $1$: a null-homotopy across no
plaquette would give $T_\gamma(V)=1$ for every field $V$, whereas the field $V^{\flat}$ of (ii),
whose construction does not use the null-homotopy, has
$T_\gamma(V^{\flat})\neq 1$.

(ii) We write $\bar S(B)$ for the lattice box whose cubical complex is
$K(\bar S(B))=\mathbb{B}$, as in Definition~\ref{8.10:def-spanned-box}, and
\begin{equation*}
\bar S(B)^{\circ}:=\prod_{\nu}\{\ell a_\nu+1,\dots,\ell b_\nu+\ell-1\}
\end{equation*}
for its set of strictly interior sites, with $\ell$, $a_\nu$, $b_\nu$ the integers that fix
$\bar S(B)$ in Definition~\ref{8.10:def-spanned-box}. The boundary complex
$K(\partial\bar S(B))$ is contained
in $\mathbb{E}_B$ (Definition~\ref{8.10:def-spanned-box}), and
$H^1(K(\partial\bar S(B));\mathbb{R})=0$ by the facts on the cohomology of the boundary complex
of a box collected in Notation~\ref{8.10:not-topology-facts}. Hence the restriction of $\xi$ to
$K(\partial\bar S(B))$, which is closed, is a coboundary: there is a real $0$-cochain $\psi$ on
the vertices of $\partial\bar S(B)$ with
\begin{equation*}
\xi\big|_{K(\partial\bar S(B))}=d\psi .
\end{equation*}
Put $\tilde\psi:=\psi$ on $\partial\bar S(B)$ and $\tilde\psi:=0$ on the sites of $\mathbb{T}$
outside $\bar S(B)$. Define $\tilde\xi$ on the exterior bonds of $\mathbb{T}$ by
\begin{itemize}
\item[] $\tilde\xi:=\xi$ on the edges of $\mathbb{E}_B$ (bonds of the first kind), and
\item[] $\tilde\xi:=d\tilde\psi$ on the exterior edges both of whose endpoints lie outside
$\bar S(B)^{\circ}$ (bonds of the second kind).
\end{itemize}
A bond of both kinds is an edge of $K(\partial\bar S(B))$, where $\xi=d\psi=d\tilde\psi$; so the
two prescriptions agree.

Every exterior bond of $\mathbb{T}$ receives a value. Indeed, let $e$ be an exterior bond. If $e$
has an endpoint in $\bar S(B)^{\circ}$, then $e$ lies in $K(\bar S(B))=\mathbb{B}$, because one
step from a site strictly inside a lattice box stays inside the closed box; being exterior, $e$
then lies in $\mathbb{E}_B$, and is of the first kind. If both endpoints of $e$ lie outside
$\bar S(B)^{\circ}$, then $e$ is of the second kind.

Next, $d\tilde\xi(p)=0$ for every plaquette $p$ of $\mathbb{T}$ without interior edge; all
edges of such a $p$ are exterior, so $d\tilde\xi(p)$ is defined. Suppose first that $p$ has a
vertex in $\bar S(B)^{\circ}$. Two steps from a strictly interior site reach at most the
boundary layer, so all vertices of $p$ lie in $\bar S(B)$ and $p\in\mathbb{B}$. The plaquette
$p$ has no interior edge and is not a vertex, in particular not a deep vertex; hence
$p\in\mathbb{E}_B$ by Definition~\ref{8.10:def-spanned-box}, and
$d\tilde\xi(p)=d\xi(p)=0$ because $\xi$ is closed. Suppose now that no vertex of $p$ lies in
$\bar S(B)^{\circ}$. Then every edge of $p$ is of the second kind, and
$d\tilde\xi(p)=dd\tilde\psi(p)=0$. In all of this $\bar S(B)$ is regarded as a subset of
$\mathbb{T}$; that it is embedded injectively in the torus uses $\ell\ge 2$, through the patch
condition.

Now let $\hat H:=i\operatorname{diag}(1,-1,0,\dots,0)$, an element of the Lie algebra
$\mathfrak{su}(N)$ of $\mathrm{SU}(N)$, and put
\begin{equation*}
V^{\flat}_e:=\exp\bigl(\pi\tilde\xi_e\hat H\bigr)
\end{equation*}
on the exterior bonds $e$ of $\mathbb{T}$. These matrices commute with one another, so for every
plaquette $p$ without interior edge the ordered product around $\partial p$ collapses to the
exponential of the oriented sum:
\begin{equation*}
V^{\flat}(\partial p)=\exp\bigl(\pi\,(d\tilde\xi)(p)\,\hat H\bigr)=\exp(0)=1 ,
\end{equation*}
and hence $F_p(V^{\flat})=0$. In the same way, since $\gamma$ runs in $\mathbb{E}_B$, where
$\tilde\xi=\xi$, and $\langle\xi,\gamma\rangle=1$,
\begin{align*}
T_\gamma(V^{\flat})
&=\exp\bigl(\pi\langle\xi,\gamma\rangle\hat H\bigr)
=\exp(\pi\hat H)
=\operatorname{diag}\bigl(e^{i\pi},e^{-i\pi},1,\dots,1\bigr)\\
&=\operatorname{diag}(-1,-1,1,\dots,1)\in\mathrm{SU}(N),
\end{align*}
so that $T_\gamma(V^{\flat})-1=\operatorname{diag}(-2,-2,0,\dots,0)$ and
$\|T_\gamma(V^{\flat})-1\|=2$.

(iii) The hypotheses of Proposition~\ref{8.10:prop-non-fillable} (stated; proof incomplete at
the step marked $(\ast)$) are: $p_1\ge1$, $p_2\ge3$; $C$ a touch-connected union of cells and
$\lambda\in\Lambda_C$ with $K_{\mathrm{free}}(C;\lambda;p_1+p_2+p_W)$; an edge loop all of whose
edges are exterior; a null-homotopy of it in $K(\mathbb{T})$ across plaquettes
$(q_1,\dots,q_A)$ with $A\ge1$; and a field in $\mathcal{V}^{\mathrm{ext}}$ with $F_p=0$ on every
plaquette of $\mathbb{T}$ having no interior edge and with $h_0>0$. The first two are assumed;
the loop $\gamma$ and the plaquettes $q_1,\dots,q_A$ with $A\ge1$ are provided by (i); the field
$V^{*}_{\mathrm{ext}}:=V^{\flat}$ is provided by (ii), with $h_0=2>0$. The null-homotopy of (i)
lies in $\mathbb{B}=K(\bar S(B))$, which is contained in $K(\mathbb{T})$ since $\bar S(B)$ is
embedded in $\mathbb{T}$; and the proof of Proposition~\ref{8.10:prop-non-fillable} uses of the
null-homotopy only that it lies in $K(\mathbb{T})$. Since
Proposition~\ref{8.10:prop-non-fillable} is a hypothesis of (iii), it applies verbatim to these
data with $h_0=2$: under the outer window
bound, the inner window bound, the localisation of the windows and the format condition in
$\lambda$, it gives, for every $E>0$ and every $\lambda\in\Lambda_C$ with
$K_{\mathrm{free}}(C;\lambda;p_1+p_2+p_W)$, that $U$ is open, of positive measure, contained in
$\mathrm{Adm}\cap G^{\mathrm{mar}}_{\mu}\cap G^{\mathrm{av}}_{f_*}\cap G^{\mathrm{en}}_{E}$ and
disjoint from $\mathrm{Fill}_{\mu,\mathfrak{E}}(C)$ for every $\mathfrak{E}$, and that
$a_\lambda\equiv0$ on $U$ when $c_W\varepsilon\le 3/(4A)$, that is when
$x\ge x_{\mathrm{hol}}(C)$.

Finally, $x_{\mathrm{hol}}(C)<\infty$: since $(\log x')^{p_0}x'^{-1/2}\to0$ as $x'\to\infty$,
there is $x_0\ge e$ with
$\tfrac43\,c_WA\,A_0(\log x')^{p_0}x'^{-1/2}\le 1$ for all $x'\ge x_0$, so the set in
\eqref{8.10:eq-obstruction-instances-1} is non-empty.
\end{proof}

\begin{remark}
Proposition~\ref{8.10:prop-obstruction-instances} is stated with $\operatorname{lk}^{\sharp}(C)$,
the rank of the obstruction group of the split used in this chapter, and shows that every such
$C$ is linked, in the sense of the definition of linked unions of cells in this chapter, when the
patch condition holds. For the counting argument of this chapter that uses these ranks,
Lemma~\ref{8.10:lem-obstruction-ranks}(a) bounds $\operatorname{lk}(C)$ by a multiple of $n(C)$
with a constant independent of $\ell$, whereas for $\operatorname{lk}^{\sharp}(C)$ it gives only
$\operatorname{lk}^{\sharp}(C)\le|\mathsf{F}|$, whose bound in terms of $n(C)$ grows with $\ell$.
Equality of the two ranks is not proved; it is expected for hulls whose boundary has a product
collar. A thickened
$2$-torus carries, besides abelian classes, obstructions through commutators of exterior
holonomies; such a union of cells is linked, although these obstructions are counted by neither
rank.
\end{remark}

\begin{proposition}[The hollow-shell instance]\label{8.10:prop-hollow-shell}
We use the cells of side $\ell$, the site set $S(c)$ of a cell $c$ and its union $S(X)$ over a union
$X$ of cells, the set $\bar S(X)$, the cell distance
$\operatorname{dist}_{\blacksquare}$, the $k$-hull $\operatorname{hull}_k(X)$ (the union of the cells
at cell distance at most $k$ from $X$), and the set $\mathsf{C}$ of all cells, as in
Lemma~\ref{8.10:lem-steps-cells}. We use the regions $\mathsf{R}_0$, $\mathsf{R}_2$, $\mathsf{N}_0$,
$\mathsf{N}_2$, $\mathsf{N}$ and the interior and exterior edges as fixed in this section, together
with the lemma of this section which places in $\mathsf{R}_2$ (resp.\ $\mathsf{N}_2$) every plaquette
whose four vertices lie in $\mathsf{R}_0$ (resp.\ $\mathsf{N}_0$). Let $p_1\ge 1$ and $p_2\ge 3$ be
integers as in Proposition~\ref{8.10:prop-non-fillable}, and let $b_W$, $r_W$, $p_W$, $c_W$ be the
constants of the window structure of Definition~\ref{8.10:def-window-structure}. The sets
$\Lambda_C$, $\mathcal{G}_{f_*}(C)$, $\mathcal{N}_{\mathrm{ext}}$, the set $\mathcal{G}$ with its
clauses (among them $\mathrm{Fill}$), the condition $K_{\mathrm{free}}$, the local integrals and the
holonomy smallness condition are those of Proposition~\ref{8.10:prop-non-fillable}.

Let $\mathbb{T}$ be the terminal torus, of side $\mathfrak{N}$ in lattice units, on which
Proposition~\ref{8.10:prop-non-fillable} is stated. Write its sites as $(y,t)$,
with $y\in(\mathbb{Z}/\mathfrak{N}\mathbb{Z})^{d-1}$ (directions $1,\dots,d-1$) and
$t\in\mathbb{Z}/\mathfrak{N}\mathbb{Z}$ (direction $d$). Lift coordinates to
$]-\mathfrak{N}/2,\mathfrak{N}/2]$, and let $|\cdot|_\infty$ and the sup-distance refer to the lifted
coordinates. Put $k_0:=\lceil (r_W+2)/\ell\rceil$. Let $\mathfrak{b}$ and $\mathfrak{R}$ be integers,
and let $\mathfrak{N}$ satisfy
\begin{equation}\label{8.10:eq-hollow-shell-1}
\mathfrak{b}\ge b_W+k_0,\qquad \mathfrak{R}\ge\mathfrak{a},\qquad
\mathfrak{N}\ge 2\mathfrak{R}+2(p_1+\mathfrak{b}+p_2+3)\ell+3,
\end{equation}
where $\delta:=(p_1+\mathfrak{b}+1)\ell$ and $\mathfrak{a}:=\delta+2$.
Put $\mathfrak{D}:=\{y:|y|_\infty\le\mathfrak{R}\}$. A \emph{wall plaquette} is a plaquette
$(z,0)+\{0,1\}^{\{i,d\}}$, with $1\le i\le d-1$, such that exactly one of $z$, $z+e_i$ lies in
$\mathfrak{D}$. Let $W_0$ be the set of vertices of wall plaquettes, and put
\begin{equation*}
C_0:=\{c\in\mathsf{C}: S(c)\cap W_0\neq\emptyset\},\qquad
C_{\mathrm{shell}}:=\operatorname{hull}_{\mathfrak{b}}(C_0).
\end{equation*}
Let
\begin{equation*}
\mathsf{h}:=\operatorname{diag}\bigl(e^{i\pi/3},e^{-i\pi/3},1,\dots,1\bigr)\in\mathrm{SU}(N),
\end{equation*}
so that $\|\mathsf{h}-1\|=2\sin(\pi/6)=1$. Define the configuration $V^*$ by
$V^*_{b'}:=\mathsf{h}$ on the bonds $b'$ joining $(y,0)$ to $(y,1)$ with $y\in\mathfrak{D}$, and
$V^*_{b'}:=1$ on every other bond. Let $V^*_{\mathrm{ext}}$ be the restriction of $V^*$ to the
exterior bonds. Put $C:=C_{\mathrm{shell}}$ and
\begin{equation*}
R_{\mathrm{int}}=\operatorname{hull}_{p_1}(C_{\mathrm{shell}})
=\operatorname{hull}_{p_1+\mathfrak{b}}(C_0).
\end{equation*}
Recall that a core cell of $C$ is a cell of $C$ at cell distance at least $b_W$ from
$\mathsf{C}\setminus C$. Assume the following.
\begin{itemize}
\item[(a)] The small-field window structure of Definition~\ref{8.10:def-window-structure} holds,
including its inner window bound. This structure is assumed here; its proof is outstanding.
\item[(b)] The format condition on the set of plaquettes $\mathcal{P}_M$ holds in the following
form: if each vertex $v$ of a plaquette $p$ lies in $S(c_v)$ for some core cell $c_v$ of $C$,
depending on $v$, then $p\notin\mathcal{P}_M$.
\item[(c)] For the positivity assertion in item~(7) only: $\lambda$ satisfies the strengthened format
condition. This holds, for instance, when the set $\mathfrak{K}(\lambda)$ attached to $\lambda$ is
empty.
\item[(d)] (Realisability, taken from the cover rule.) $C_{\mathrm{shell}}$ is a live component of
some history, and $\Lambda_{C_{\mathrm{shell}}}$ contains an element $\lambda$ with
$\mathfrak{K}(\lambda)=\emptyset$.
\end{itemize}
Then the following hold.
\begin{enumerate}
\item[(1)] The set $W_0$ is connected under lattice adjacency, so $C_0$ and $C_{\mathrm{shell}}$ are
touch-connected. Every site $(y,t)$ of
$\bar S(\operatorname{hull}_{p_1+\mathfrak{b}+p_2+1}(C_0))$ satisfies
\begin{equation*}
|y|_\infty\le\mathfrak{R}+1+(p_1+\mathfrak{b}+p_2+2)\ell,\qquad
|t|\le 1+(p_1+\mathfrak{b}+p_2+2)\ell .
\end{equation*}
The coordinate-patch condition holds for $\operatorname{hull}_{p_1+p_2+1}(C_{\mathrm{shell}})$.
\item[(2)] $F_p(V^*)=0$ unless $p$ is a wall plaquette, and
$V^*(\partial p)\in\{\mathsf{h},\mathsf{h}^{-1}\}$ on wall plaquettes. Further:
\begin{itemize}
\item $F_p(V^*_{\mathrm{ext}})=0$ for every plaquette $p$ without interior edge;
\item every cell at cell distance at most $k_0$ from $C_0$ is a core cell of $C_{\mathrm{shell}}$;
\item under \textup{(b)}, $F_{p'}(V^*)=0$ for every $p\in\mathcal{P}_M$ and every plaquette
$p'$ contained in the $(r_W+1)$-neighbourhood $N_{r_W+1}(p)$ of $p$.
\end{itemize}
\item[(3)] Every site of $\bar S(R_{\mathrm{int}})$ lies within sup-distance $\delta$ of $W_0$. Every
site within sup-distance $(p_1+\mathfrak{b}+p_2)\ell$ of $W_0$ lies in
$S(\operatorname{hull}_{p_1+\mathfrak{b}+p_2}(C_0))\subset\mathsf{N}_0$.
\item[(4)] Let $\Pi$ be the meridian square in the $(1,d)$-plane,
\begin{equation*}
\Pi:=\bigl\{((\mathfrak{R}+s)e_1,t):-\mathfrak{a}\le s\le\mathfrak{a},\
-\mathfrak{a}\le t\le\mathfrak{a}\bigr\},
\end{equation*}
and let $\gamma$ be its boundary loop. Then:
\begin{itemize}
\item $|P(\Pi)|=4\mathfrak{a}^2$;
\item every vertex of $\gamma$ is at sup-distance at least $\delta+1$ from $W_0$, lies outside
$\mathsf{R}_0$, and every edge of $\gamma$ is exterior;
\item all vertices of $\Pi$ lie in $\mathsf{N}_0$, and $P(\Pi)\subset\mathsf{N}_2$;
\item $\gamma$ admits a null-homotopy in $K(\Pi)\subset\mathsf{N}$ across the $4\mathfrak{a}^2$
plaquettes of $P(\Pi)$.
\end{itemize}
\item[(5)] Exactly one edge of $\gamma$ carries $\mathsf{h}$ under $V^*$. Hence
$T_\gamma(V^*_{\mathrm{ext}})$ is a conjugate of $\mathsf{h}^{\pm1}$, and
$h_0:=\|T_\gamma(V^*_{\mathrm{ext}})-1\|=1$.
\item[(6)] Let $x=g'^{-2}$ be the inverse square of the coupling $g'$ at which the degree function
$p_0(\cdot)$ is evaluated in the holonomy smallness condition of
Proposition~\ref{8.10:prop-non-fillable}, so that $A_0(\log x)^{p_0}=p_0(g')$, and let
$r_{\mathrm{pr}}$ be the fixed exponent of the bound $\ell\le ML(\log x)^{r_{\mathrm{pr}}}$ on the
cell side. For $C_{\mathrm{shell}}$ the holonomy smallness condition reads
\begin{equation*}
c_W\,x^{-1/2}A_0(\log x)^{p_0}\le\frac{3}{32\,\mathfrak{a}^2},
\end{equation*}
and it holds whenever $x\ge x_{\mathrm{hol}}$, where
\begin{equation}\label{8.10:eq-hollow-shell-2}
\begin{split}
x_{\mathrm{hol}}:=\inf\Bigl\{x_0\ge e:\ &\tfrac{32}{3}\,c_WA_0(\log x')^{p_0}
\bigl((p_1+\mathfrak{b}+1)ML(\log x')^{r_{\mathrm{pr}}}+2\bigr)^2\le x'^{1/2}\\
&\text{for all }x'\ge x_0\Bigr\},
\end{split}
\end{equation}
and $x_{\mathrm{hol}}<\infty$.
\item[(7)] All hypotheses of Proposition~\ref{8.10:prop-non-fillable}, whose proof is incomplete at
one step, hold for $C_{\mathrm{shell}}$ with $A=4\mathfrak{a}^2$ and $h_0=1$; moreover
$\gamma\subset\mathsf{N}$, and every plaquette $q_i$ of the null-homotopy lies in $\mathsf{N}_2$.
Consequently, by that proposition, $C_{\mathrm{shell}}$ is linked, and neither of the two unlinking
conditions of that proposition holds. Let
$\lambda\in\Lambda_{C_{\mathrm{shell}}}$ satisfy
$K_{\mathrm{free}}(C_{\mathrm{shell}};\lambda;p_1+p_2+p_W)$. Then:
\begin{itemize}
\item there is an open set $U$ of exterior configurations, of positive Haar measure, contained in
$\mathcal{G}_{f_*}(C_{\mathrm{shell}})$ and in the set defined by every clause of $\mathcal{G}$
other than $\mathrm{Fill}$;
\item under the holonomy smallness condition, the local integral $a_\lambda$ with the component
$C_{\mathrm{shell}}$ absent vanishes identically on $U$;
\item under \textup{(a)} and \textup{(b)}, if $\lambda$ satisfies \textup{(c)}, the local integral of
the majorant is positive at almost every point of $U\cap\mathcal{N}_{\mathrm{ext}}$, an open set of
positive measure containing $V^*_{\mathrm{ext}}$.
\end{itemize}
\end{enumerate}
\end{proposition}

The conditions \eqref{8.10:eq-hollow-shell-1} are one-sided conditions on the auxiliary integers
$\mathfrak{b}$ and $\mathfrak{R}$, and on the side $\mathfrak{N}$ of a torus large enough to contain
the shape. The proposition exhibits one shape, and an open set of exterior configurations, on one such
torus. It bounds no parameter of the order of choice.

\begin{proof}
\emph{Item \textup{(1)}.} Let $(z,0)+\{0,1\}^{\{i,d\}}$ be a wall plaquette with, say,
$z\in\mathfrak{D}$ and $z+e_i\notin\mathfrak{D}$. Then $|z+e_i|_\infty\ge\mathfrak{R}+1$, and $z$,
$z+e_i$ differ by one in one coordinate, so $|z|_\infty=\mathfrak{R}$. The same holds with the roles of
$z$ and $z+e_i$ exchanged. Conversely, if $|y|_\infty=\mathfrak{R}$, some coordinate satisfies
$y_i=\pm\mathfrak{R}$, and $y\pm e_i\notin\mathfrak{D}$.

Hence the $y$-parts of the sites of $W_0$ consist of two kinds of points:
\begin{itemize}
\item the discrete sphere $\{|y|_\infty=\mathfrak{R}\}$;
\item sites of sup norm $\mathfrak{R}+1$ adjacent to that sphere.
\end{itemize}
Their $t$-coordinates are $0$ and $1$. The discrete sphere is connected, since $d-1\ge 2$. The sites
of sup norm $\mathfrak{R}+1$ in $W_0$ are attached to it by lattice edges, and so are the two layers
$t=0$ and $t=1$. So $W_0$ is connected under lattice adjacency, and therefore $C_0$ and
$C_{\mathrm{shell}}=\operatorname{hull}_{\mathfrak{b}}(C_0)$ are touch-connected.

A site of $\bar S(c'')$ with $\operatorname{dist}_{\blacksquare}(c'',c)\le k$, where $S(c)$ contains
some $w\in W_0$, is within sup-distance $k\ell+\ell$ of $w$. Take
$k=p_1+\mathfrak{b}+p_2+1$, and recall that the $y$-parts of $W_0$ have sup norm at most
$\mathfrak{R}+1$ and the $t$-coordinates of $W_0$ lie in $\{0,1\}$. This gives the two bounds of
item~(1) for every site of $\bar S(\operatorname{hull}_{p_1+\mathfrak{b}+p_2+1}(C_0))$. Since
\begin{equation*}
\operatorname{hull}_{p_1+\mathfrak{b}+p_2+1}(C_0)=\operatorname{hull}_{p_1+p_2+1}(C_{\mathrm{shell}}),
\end{equation*}
the coordinate-patch condition for $\operatorname{hull}_{p_1+p_2+1}(C_{\mathrm{shell}})$ holds by the
choice of $\mathfrak{N}$ in \eqref{8.10:eq-hollow-shell-1}.

\emph{Item \textup{(2)}.} A plaquette contains a bond joining $(y,0)$ to $(y,1)$ only if it is of the
form $(z,0)+\{0,1\}^{\{i,d\}}$ with $1\le i\le d-1$. Write $[P]$ for $1$ if $P$ holds and $0$
otherwise. The holonomy of such a plaquette from $(z,0)$, taken along $e_i$ first, is
\begin{equation*}
1\cdot\mathsf{h}^{[z+e_i\in\mathfrak{D}]}\cdot 1\cdot\mathsf{h}^{-[z\in\mathfrak{D}]}
=\mathsf{h}^{[z+e_i\in\mathfrak{D}]-[z\in\mathfrak{D}]}.
\end{equation*}
This equals $1$ unless exactly one of $z$, $z+e_i$ lies in $\mathfrak{D}$. In that case the plaquette
is a wall plaquette and the holonomy is $\mathsf{h}$ or $\mathsf{h}^{-1}$. Every other plaquette meets
only bonds carrying $1$. Hence $F_p(V^*)=0$ unless $p$ is a wall plaquette, and
$V^*(\partial p)\in\{\mathsf{h},\mathsf{h}^{-1}\}$ on wall plaquettes.

Every wall plaquette has its four vertices in $W_0$. Moreover
\begin{equation*}
W_0\subset S(C_0)\subset\bar S(R_{\mathrm{int}})=\mathsf{R}_0 .
\end{equation*}
By the lemma of this section which places in $\mathsf{R}_2$ every plaquette whose four vertices lie in
$\mathsf{R}_0$, every wall plaquette lies in $\mathsf{R}_2$. Hence $F_p(V^*)=0$ for every plaquette
$p$ without interior edge. That is, $V^*_{\mathrm{ext}}$ is flat on all such plaquettes.

Next, the cells of $C_0$ are core cells of $C_{\mathrm{shell}}$. Indeed, a cell at cell distance at
most $b_W-1\le\mathfrak{b}$ from $C_0$ lies in
$\operatorname{hull}_{\mathfrak{b}}(C_0)=C_{\mathrm{shell}}$. So cells outside $C_{\mathrm{shell}}$
are at cell distance at least $\mathfrak{b}+1>b_W$ from $C_0$.

More is true. Let a cell be at cell distance at most $k_0$ from $C_0$. Since $k_0\le\mathfrak{b}$, it
lies in $C_{\mathrm{shell}}$. Its cell distance from $\mathsf{C}\setminus C_{\mathrm{shell}}$ is at
least $\mathfrak{b}+1-k_0$, which is at least $b_W$ by \eqref{8.10:eq-hollow-shell-1}. So it is a core
cell.

Now let $p'$ be a plaquette with $F_{p'}(V^*)\neq0$; by the above, $p'$ is a wall plaquette with
vertices in $W_0$. Suppose $p'\subset N_{r_W+1}(p)$ for some plaquette $p$. Then every vertex of $p$
is within sup-distance $r_W+2$ of $W_0\subset S(C_0)$, and $r_W+2\le k_0\ell$ by the definition of
$k_0$. Apply Lemma~\ref{8.10:lem-steps-cells} with $k=k_0$, $z$ a vertex of $p$, and $w\in W_0$ in
$S(c)$ for some $c\in C_0$. It places each vertex of $p$ in $S(c'')$ for a cell $c''$ at cell distance
at most $k_0$ from $C_0$, which is a core cell. By hypothesis~(b), $p\notin\mathcal{P}_M$.

Hence $F_{p'}(V^*)=0$ for every $p\in\mathcal{P}_M$ and every $p'\subset N_{r_W+1}(p)$. This is the
extension hypothesis of the positivity part of Proposition~\ref{8.10:prop-non-fillable}, whose proof
is incomplete at one step, and it holds with $V^*$ itself.

\emph{Item \textup{(3)}.} Let $c''$ be a cell with $\operatorname{dist}_{\blacksquare}(c'',c)\le
p_1+\mathfrak{b}$, where $S(c)$ contains some $w\in W_0$. A site of $\bar S(c'')$ is within
sup-distance $(p_1+\mathfrak{b})\ell+\ell=\delta$ of $w$. Since
$R_{\mathrm{int}}=\operatorname{hull}_{p_1+\mathfrak{b}}(C_0)$, every site of $\bar S(R_{\mathrm{int}})$
is within sup-distance $\delta$ of $W_0$.

Now let $z$ be within sup-distance $(p_1+\mathfrak{b}+p_2)\ell$ of $W_0$. Apply
Lemma~\ref{8.10:lem-steps-cells} with $k=p_1+\mathfrak{b}+p_2$. It gives
$z\in S(\operatorname{hull}_{p_1+\mathfrak{b}+p_2}(C_0))$, and this set is contained in
$\mathsf{N}_0$.

\emph{Item \textup{(4)}.} The set $\Pi$ is the rectangle of Lemma~\ref{8.10:lem-stokes-rectangle} with
the following data: base point $((\mathfrak{R}-\mathfrak{a})e_1,-\mathfrak{a})$, $\mu=1$, $\nu=d$, and
side lengths $2\mathfrak{a}$ and $2\mathfrak{a}$. So $|P(\Pi)|=(2\mathfrak{a})^2=4\mathfrak{a}^2$.

Every vertex of $\gamma$ is at sup-distance at least $\mathfrak{a}-1=\delta+1>\delta$ from $W_0$.
There are three cases.
\begin{itemize}
\item For $|t|=\mathfrak{a}$: the $t$-coordinates of $W_0$ are $0$ and $1$.
\item For $s=+\mathfrak{a}$: $|y|_\infty=\mathfrak{R}+\mathfrak{a}$, while the $y$-parts of $W_0$
have sup norm at most $\mathfrak{R}+1$.
\item For $s=-\mathfrak{a}$: $|y|_\infty=\mathfrak{R}-\mathfrak{a}\ge0$, while the $y$-parts of $W_0$
have sup norm at least $\mathfrak{R}$. This is the cavity of the shell, and uses
$\mathfrak{R}\ge\mathfrak{a}$.
\end{itemize}
By item~(3), the vertices of $\gamma$ therefore lie outside $\bar S(R_{\mathrm{int}})=\mathsf{R}_0$,
and the edges of $\gamma$ are exterior.

Every vertex $((\mathfrak{R}+s)e_1,t)$ of $\Pi$, with $|s|,|t|\le\mathfrak{a}$, is within sup-distance
$\mathfrak{a}$ of $(\mathfrak{R}e_1,0)$. This point is a wall vertex: $\mathfrak{R}e_1\in\mathfrak{D}$
and $\mathfrak{R}e_1+e_1\notin\mathfrak{D}$, so it is a vertex of the wall plaquette
$(\mathfrak{R}e_1,0)+\{0,1\}^{\{1,d\}}$. Since $p_2\ge3$ and $\ell\ge1$, we have $(p_2-1)\ell\ge2$,
and therefore
\begin{equation*}
\mathfrak{a}=(p_1+\mathfrak{b}+1)\ell+2\le(p_1+\mathfrak{b}+p_2)\ell .
\end{equation*}
By item~(3), all vertices of $\Pi$ lie in $\mathsf{N}_0$. By the lemma of this section which places in
$\mathsf{N}_2$ every plaquette whose four vertices lie in $\mathsf{N}_0$, $P(\Pi)\subset\mathsf{N}_2$.
Lemma~\ref{8.10:lem-stokes-rectangle} then gives a null-homotopy of $\gamma$ in
$K(\Pi)\subset\mathsf{N}$ across the $4\mathfrak{a}^2$ plaquettes of $P(\Pi)$.

\emph{Item \textup{(5)}.} The edges of $\gamma$ in direction $d$ are the bonds joining
$((\mathfrak{R}\pm\mathfrak{a})e_1,t)$ to $((\mathfrak{R}\pm\mathfrak{a})e_1,t+1)$, for
$-\mathfrak{a}\le t<\mathfrak{a}$. Such an edge carries $\mathsf{h}$ under $V^*$ if and only if $t=0$
and $(\mathfrak{R}\pm\mathfrak{a})e_1\in\mathfrak{D}$. Now
$|(\mathfrak{R}-\mathfrak{a})e_1|_\infty=\mathfrak{R}-\mathfrak{a}\le\mathfrak{R}$, whereas
$|(\mathfrak{R}+\mathfrak{a})e_1|_\infty=\mathfrak{R}+\mathfrak{a}>\mathfrak{R}$. So the condition
holds for $(\mathfrak{R}-\mathfrak{a})e_1$ only.

The edges of $\gamma$ in direction $1$ lie at $|t|=\mathfrak{a}\ge2$ and carry $1$. So exactly one edge
of $\gamma$ carries $\mathsf{h}$. All edges of $\gamma$ are exterior by item~(4), so
$T_\gamma(V^*_{\mathrm{ext}})$ is defined. It is a conjugate of $\mathsf{h}^{\pm1}$, and
$h_0=\|T_\gamma(V^*_{\mathrm{ext}})-1\|=1$.

\emph{Item \textup{(6)}.} For the shell, the holonomy smallness condition reads as displayed in
item~(6). Here $c_W\ge1$ is the constant of the window structure of
Definition~\ref{8.10:def-window-structure}, a structure whose proof is outstanding, and, by the
bound $\ell\le ML(\log x)^{r_{\mathrm{pr}}}$ on the cell side,
\begin{equation*}
\mathfrak{a}=(p_1+\mathfrak{b}+1)\ell+2\le(p_1+\mathfrak{b}+1)ML(\log x)^{r_{\mathrm{pr}}}+2 .
\end{equation*}
Let $x_{\mathrm{hol}}$ be defined by \eqref{8.10:eq-hollow-shell-2}.
The left-hand side of the inequality in \eqref{8.10:eq-hollow-shell-2} grows like a power of
$\log x'$, so the set is non-empty and $x_{\mathrm{hol}}<\infty$. The holonomy smallness condition is
implied by $x\ge x_{\mathrm{hol}}$. This is a lower bound on $x$. Through the coupling statement of
clause~(0) of Theorem~\ref{3.1:thm-main} it becomes a lower bound on $X$. No value is assigned to
$x_{\mathrm{hol}}$.

\emph{Item \textup{(7)}.} We check the hypotheses of Proposition~\ref{8.10:prop-non-fillable}, whose
proof is incomplete at one step, for $C=C_{\mathrm{shell}}$ and the integers $p_1\ge1$, $p_2\ge3$.
\begin{itemize}
\item $C_{\mathrm{shell}}$ is a touch-connected union of cells, by item~(1). The coordinate-patch
condition for $\operatorname{hull}_{p_1+p_2+1}(C_{\mathrm{shell}})$ also holds, by item~(1).
\item $\gamma$ is an edge loop all of whose edges are exterior, by item~(4).
\item $\gamma$ has a null-homotopy in $K(\mathbb{T})$ across the $A=4\mathfrak{a}^2$ plaquettes
$q_1,\dots,q_A$ of $P(\Pi)$, with all $q_i\in\mathsf{N}_2$ and $\gamma\subset K(\Pi)\subset\mathsf{N}$,
by item~(4).
\item $V^*_{\mathrm{ext}}$ is flat on every plaquette of $\mathbb{T}$ with no interior edge, by
item~(2).
\item $h_0=1>0$, by item~(5).
\end{itemize}

Proposition~\ref{8.10:prop-non-fillable} therefore gives the following. $C_{\mathrm{shell}}$ is linked,
and neither of the two unlinking conditions holds. For every $\lambda\in\Lambda_{C_{\mathrm{shell}}}$
with $K_{\mathrm{free}}(C_{\mathrm{shell}};\lambda;p_1+p_2+p_W)$ there is an open set $U$ of exterior
configurations, of positive Haar measure, with the following properties:
\begin{itemize}
\item $U$ is contained in $\mathcal{G}_{f_*}(C_{\mathrm{shell}})$ and in the set defined by every
clause of $\mathcal{G}$ other than $\mathrm{Fill}$;
\item under the holonomy smallness condition, which item~(6) reduces to $x\ge x_{\mathrm{hol}}$, the
local integral $a_\lambda$ vanishes identically on $U$.
\end{itemize}

Finally, the extension hypothesis of the positivity part of Proposition~\ref{8.10:prop-non-fillable}
holds with $V^*$ itself, by item~(2), whose last assertion uses hypothesis~(b). Under
hypotheses~(a) and~(b), the positivity part of that proposition therefore applies to every such
$\lambda$ satisfying~(c). It gives that the local integral of the majorant is positive at almost every
point of the open set $U\cap\mathcal{N}_{\mathrm{ext}}$, which contains $V^*_{\mathrm{ext}}$ and has
positive measure.
\end{proof}

Hypothesis~(d) is taken from the cover rule and is not proved here; it is the only point at which
realisability of $C_{\mathrm{shell}}$ as a live component of a history enters. The sums over shapes in
which the old-member statement is applied run over all touch-connected unions of cells, whether or not
they are realisable as live components. So $C_{\mathrm{shell}}$ enters those sums in either case.

\begin{corollary}[Failure of the domination on the curvature-defined good set]
\label{8.10:cor-domination-failure}
We work in the setting of the old-member statement, Proposition~\ref{8.10:prop-old-member}, at the
terminal level $K'$. There the top gauge field is integrated over the interior bond variables
$V^{\mathrm{in}}$, with the exterior configuration $V^{\mathrm{ext}}$ and all labels held fixed. For a label
$\lambda\in\Lambda_C$ with $K_{\mathrm{free}}(C;\lambda;p_1+p_2+p_W)$ and a fixed
$w\in\mathcal{U}_\lambda$ we put
\begin{align*}
m_\lambda(V^{\mathrm{ext}})&:=\int e^{\delta_C}M_{\lambda,C}(w,V^{\mathrm{in}},V^{\mathrm{ext}})
\,dV^{\mathrm{in}},\\
a_\lambda(V^{\mathrm{ext}})&:=\int \rho_\lambda(w,V^{\mathrm{in}},V^{\mathrm{ext}})\,dV^{\mathrm{in}},
\end{align*}
the left and the right integral in the inequality of Proposition~\ref{8.10:prop-old-member}:
$m_\lambda$ integrates the net majorant of the term in which $C$ is present, and $a_\lambda$ the
integrand of the term in which $C$ and its lineage are absent.
For a set $G$ of exterior configurations and a number $\Gamma<\infty$, we say that the
\emph{good-exterior domination holds on $G$ with constant $\Gamma$} if
\begin{equation}\label{8.10:eq-domination-failure-1}
m_\lambda(V^{\mathrm{ext}})\le e^{\Gamma}\,a_\lambda(V^{\mathrm{ext}})
\qquad\text{for every } V^{\mathrm{ext}}\in G .
\end{equation}
We further write:
\begin{itemize}
\item $\mathrm{Adm}$ for the set of admissible exterior configurations;
\item $G^{\mathrm{av}}_{f_*}(C)$ for the set of exterior configurations whose collar curvature, averaged
over each cell of the collar of $C$, is at most $f_*$;
\item $G^{\mathrm{mar}}_\mu$ and $G^{\mathrm{en}}_E$ for the margin-good set at margin $\mu$ and the
energy-good set at allowance $E$;
\item $\mathrm{Fill}(C)$ for the set of exterior configurations admitting an in-window filling of the
hull of $C$ with the fixed margin and with lattice filling energy inside the allowance displayed in
the setting of Proposition~\ref{8.10:prop-old-member}.
\end{itemize}

Assume the following statements:
\begin{itemize}
\item the three statements of this section on the exterior configuration, on the interior
configuration and on localisation;
\item the three statements of this section on the form of $\rho_\lambda$, of the majorant
$M_{\lambda,C}$ and of its enlarged version, each under one common hypothesis of this section.
\end{itemize}
Then the following hold.
\begin{enumerate}
\item[(i)] Assume moreover the reading of the cover rule under which the hollow shell
$C_{\mathrm{shell}}$ of Proposition~\ref{8.10:prop-hollow-shell} is realisable as a terminal
component.
For $C=C_{\mathrm{shell}}$, assumed to satisfy also the hypotheses of the domain class of
this section that fixes the exterior set $\mathcal{N}_{\mathrm{ext}}$, with the label $\lambda$ and the
parameters $p_1\ge1$, $p_2\ge3$
of Proposition~\ref{8.10:prop-non-fillable}, there is no finite $\Gamma$ for which the good-exterior
domination holds, pointwise in $V^{\mathrm{ext}}$, on
$G^{\mathrm{av}}_{f_*}(C_{\mathrm{shell}})\cap\mathrm{Adm}$.
More precisely, there is a subset of $G^{\mathrm{av}}_{f_*}(C_{\mathrm{shell}})\cap\mathrm{Adm}$ of
positive measure on which the right member of \eqref{8.10:eq-domination-failure-1} is $0$, that is
$a_\lambda=0$, while the left member $m_\lambda$ is positive. This holds for every $p_1$ and every
$p_2$. In particular, the failure is not removed by enlarging $p_1$.
\item[(ii)] Let $C$ satisfy the hypotheses of Proposition~\ref{8.10:prop-non-fillable} together with
the hypotheses of the domain class of this section that fixes the exterior set
$\mathcal{N}_{\mathrm{ext}}$. Then Proposition~\ref{8.10:prop-old-member} (stated; proof incomplete
at the step marked $(\ast)$) cannot be extended from its good set
$G_{\lambda,C}(\mathfrak{E})$ to $\mathrm{Adm}\cap G^{\mathrm{av}}_{f_*}(C)$. It also cannot be extended
to
\begin{equation*}
\mathrm{Adm}\cap G^{\mathrm{av}}_{f_*}(C)\cap G^{\mathrm{mar}}_\mu\cap G^{\mathrm{en}}_E .
\end{equation*}
On neither set does the good-exterior domination hold with any finite $\Gamma$.
More precisely, the set $U\cap\mathcal{N}_{\mathrm{ext}}$, where $U$ is the set of exterior
configurations provided by Proposition~\ref{8.10:prop-non-fillable} (stated; proof incomplete at the
step marked $(\ast)$), lies in the second of these sets
and outside $\mathrm{Fill}(C)$, and on it, a set of positive measure,
\begin{equation*}
0=a_\lambda<e^{-\Gamma}\,m_\lambda\qquad\text{for every finite }\Gamma .
\end{equation*}
\end{enumerate}
\end{corollary}

\begin{proof}
\emph{Item} (i). Fix $p_1$ and $p_2$. By the assumed reading of the cover rule, $C_{\mathrm{shell}}$ is
a terminal component. By Proposition~\ref{8.10:prop-hollow-shell}, it satisfies the hypotheses of
Proposition~\ref{8.10:prop-non-fillable} (stated; proof incomplete at the step marked $(\ast)$). Let
$U$ be the set of exterior configurations provided there, and let $\mathcal{N}_{\mathrm{ext}}$ be the
set of exterior configurations fixed by the hypotheses of the domain class of this section, which
$C_{\mathrm{shell}}$ satisfies by assumption.

By Proposition~\ref{8.10:prop-non-fillable} (stated; proof incomplete at the step marked $(\ast)$),
the intersection
$U\cap\mathcal{N}_{\mathrm{ext}}$ has two properties. First, it lies in
\begin{equation*}
\mathrm{Adm}\cap G^{\mathrm{av}}_{f_*}(C_{\mathrm{shell}})\cap G^{\mathrm{mar}}_\mu\cap G^{\mathrm{en}}_E
\;\subset\;\mathrm{Adm}\cap G^{\mathrm{av}}_{f_*}(C_{\mathrm{shell}}) .
\end{equation*}
Second, it lies outside $\mathrm{Fill}(C_{\mathrm{shell}})$.

By Proposition~\ref{8.10:prop-non-fillable} (stated; proof incomplete at the step marked $(\ast)$),
no interior configuration $V^{\mathrm{in}}$ completes a member $V^{\mathrm{ext}}$ of
$U\cap\mathcal{N}_{\mathrm{ext}}$ to an in-window filling of the hull: such a filling across
$q_1,\dots,q_A$ would force the holonomy $T_\gamma$ around the loop $\gamma$ of that proposition to be
close to the identity, whereas $\|T_\gamma(V^{\mathrm{ext}})-1\|>h_0/2$ on $U$. Under the six
statements assumed above, the integrand $\rho_\lambda$ then vanishes identically in $V^{\mathrm{in}}$,
so that $a_\lambda(V^{\mathrm{ext}})=0$. The wall of $C_{\mathrm{shell}}$ (the shell of
Proposition~\ref{8.10:prop-hollow-shell}) sits where $M_{\lambda,C}$ carries neither window nor action
weight; under the same statements the integrand of $m_\lambda$ is positive on a set of interior
configurations of positive Haar measure, so that $m_\lambda(V^{\mathrm{ext}})>0$. By
Proposition~\ref{8.10:prop-non-fillable} (stated; proof incomplete at the step marked $(\ast)$), the
set $U\cap\mathcal{N}_{\mathrm{ext}}$ has positive measure; at every point of it
\begin{equation}\label{8.10:eq-domination-failure-2}
0=a_\lambda(V^{\mathrm{ext}})<e^{-\Gamma}\,m_\lambda(V^{\mathrm{ext}})
\qquad\text{for every finite }\Gamma .
\end{equation}
At every such point, \eqref{8.10:eq-domination-failure-2} contradicts
\eqref{8.10:eq-domination-failure-1}, whatever the finite constant $\Gamma$. Hence the domination
fails, pointwise in $V^{\mathrm{ext}}$, on a subset of
$G^{\mathrm{av}}_{f_*}(C_{\mathrm{shell}})\cap\mathrm{Adm}$ of positive measure. On that subset its
right member vanishes and its left member is positive.

Since $p_1$ and $p_2$ were arbitrary, the construction exists for every $p_1$ and $p_2$. In particular
no enlargement of $p_1$ restores the domination.

\emph{Item} (ii). The argument is exactly that of item (i), with $C_{\mathrm{shell}}$ replaced by an
arbitrary $C$ satisfying the hypotheses of Proposition~\ref{8.10:prop-non-fillable} and those of the
domain class of this section that fixes $\mathcal{N}_{\mathrm{ext}}$. In place of the shell of
Proposition~\ref{8.10:prop-hollow-shell}, the argument uses the wall of the hull of $C$, which, as the
shell does for $C=C_{\mathrm{shell}}$, sits where $M_{\lambda,C}$ carries neither window nor action
weight.

By that proposition (stated; proof incomplete at the step marked $(\ast)$), the set
$U\cap\mathcal{N}_{\mathrm{ext}}$
lies in the set
\begin{equation*}
\mathrm{Adm}\cap G^{\mathrm{av}}_{f_*}(C)\cap G^{\mathrm{mar}}_\mu\cap G^{\mathrm{en}}_E ,
\end{equation*}
hence also in $\mathrm{Adm}\cap G^{\mathrm{av}}_{f_*}(C)$, and it lies outside $\mathrm{Fill}(C)$. As in
item (i), under the six statements assumed above, \eqref{8.10:eq-domination-failure-2} holds there, for
every finite $\Gamma$, on a set of positive measure.

An extension of Proposition~\ref{8.10:prop-old-member} to either of the two sets would assert
\eqref{8.10:eq-domination-failure-1} at every point of that set with a finite constant. That is
excluded on this set of positive measure.
\end{proof}

\begin{remark}
The curvature-defined form of the domination comes with an assertion about the exterior configurations
outside its good set, namely those whose collar curvature is at least $f_*$ on some cell. The assertion
is that these configurations are penalised by their own action, of size at least
$x_{K'}\,f_*^{2}\,(M\check{R}_{K'})^{d}$, where $x_{K'}$ is the inverse coupling at the terminal level
$K'$ and $\check{R}_{K'}$ is the cell parameter of the setting of
Proposition~\ref{8.10:prop-old-member}.

This assertion describes the complement of $G^{\mathrm{av}}_{f_*}(C)$, and also the energy-bad and
margin-bad classes, the complements of $G^{\mathrm{en}}_E$ and $G^{\mathrm{mar}}_\mu$: these classes
carry their own action under the law defined by the majorant $M_{\lambda,C}$, which is weighted by the
action off the core of $C$.

It does not, however, describe the complement of any set on which a good-exterior domination
\eqref{8.10:eq-domination-failure-1} can hold. Indeed the complement of $\mathrm{Fill}(C)$ contains the
set $U$ of Proposition~\ref{8.10:prop-non-fillable} (stated; proof incomplete at the step marked
$(\ast)$). The members of $U$ are $\eta_*$-flat on the
collar, and the configuration $V^{*}_{\mathrm{ext}}$ there is exactly flat. They carry collar action
smaller than any prescribed $E$. Meanwhile the wall of the hull of $C$ (for $C=C_{\mathrm{shell}}$, the
shell of Proposition~\ref{8.10:prop-hollow-shell}) sits where $M_{\lambda,C}$
carries neither window nor action weight.

For this reason the fillability condition, and not curvature alone, enters the definition of the good
set. A bound on the weight of $U$ under that law is among the requirements on the exterior
configurations outside the good set; these requirements are displayed in this chapter with reasons and
are not proved.
\end{remark}

\begin{remark}[The complement of the good set and its pricing. Proof outstanding]
\label{8.10:rem-good-set-complement}
We state a requirement on the typicality input of the first route of the old-member statement,
that is, on the hypothesis of that route which supplies lower bounds for probabilities under the
hole law defined below. The
requirement is not proved, and it is not part of the old-member statement.

Let $C$ be the region of the old-member statement, with its core and its collar. Let
$\mathcal{V}^{\mathrm{ext}}$ be the space of exterior configurations $V^{\mathrm{ext}}$. Let
$P_{\mathcal{M}}$ be the probability on $\mathcal{V}^{\mathrm{ext}}$ with density proportional to
the weight $m_\lambda(\cdot,w)$ at the parameter $\lambda$ of the old-member statement (not the
gauge transformation of Lemma~\ref{1.1:lem-invariance}) and the source $w$, taken either for
fixed $\lambda$ and $w$ or mixed over $\lambda$. This is the law defined by the data
$\mathcal{M}$ of the old-member statement (not to be confused with $M=L^{m'}$), and we call it
the \emph{hole law}.

Write $G$
for the good set $G_{\lambda,C}(\mathfrak{E}_*(C))$ at the energy threshold $\mathfrak{E}_*(C)$.
The typicality input must supply the lower bound
\begin{equation}\label{8.10:eq-good-set-complement-1}
P_{\mathcal{M}}\bigl(G_{\lambda,C}(\mathfrak{E}_*(C))\bigr)\ \ge\ e^{-\Lambda_C},
\qquad \Lambda_C\ \le\ \tfrac14 B^{\flat}_\tau(C).
\end{equation}
Here $\Lambda_C$ is the exponent conceded on the good set; it is to be kept inside the unspent
quarter $\tfrac14 B^{\flat}_\tau(C)$ of the reserve $B^{\flat}_\tau(C)$ of the old-member statement.
The bound is required uniformly in $K$, $m$, the parameter $N_T$ of the setting of the old-member
statement, the position of $C$ and $\lambda$, over those fibres of the mixture (the sets of fixed
$\lambda$) on which the freeness condition $K_{\mathrm{free}}$ holds. The typicality input must in
addition cover the fibres violating $K_{\mathrm{free}}$; these concern the cells whose distance
from kept structure, that is, from the structure kept near $C$ in the setting of the old-member
statement, is at most the sum $p_1+p_2+p_W$ of the lengths $p_1$, $p_2$ and $p_W$ (this is the
rule on kept structure near $C$ used in the route below).

The complement to be priced is
\begin{equation}\label{8.10:eq-good-set-complement-2}
\begin{gathered}
G^{c}=\mathrm{Adm}^{c}\cup\bigl(G^{\mathrm{av}}_{f_*}\bigr)^{c}
\cup\mathrm{Fill}_{\mu,\mathfrak{E}_*}^{c},\\
\mathrm{Fill}^{c}\ \supseteq\ \bigl(G^{\mathrm{mar}}_{\mu}\bigr)^{c}
\cup\bigl\{E^{\mu}_{\mathrm{fill}}>\mathfrak{E}_*\bigr\}\cup\mathfrak{NF}.
\end{gathered}
\end{equation}
In this display $G$ is the intersection of the admissibility event $\mathrm{Adm}$, the curvature
event $G^{\mathrm{av}}_{f_*}$ at the threshold $f_*$, and the fillability event
$\mathrm{Fill}=\mathrm{Fill}_{\mu,\mathfrak{E}_*}$ at the margin $\mu$ and the energy threshold
$\mathfrak{E}_*$; $G^{\mathrm{mar}}_{\mu}$ is the margin event and $E^{\mu}_{\mathrm{fill}}$ the
filling energy at margin $\mu$. The set $\mathfrak{NF}$ is
\begin{equation*}
\mathfrak{NF}:=\bigl\{V^{\mathrm{ext}}:\ \text{no margin-}\mu\text{ filling exists}\bigr\}.
\end{equation*}
Under the respective hypotheses, $\mathfrak{NF}$ contains the linked class of the propositions on
linking and null-homotopy. It also contains the coherent-flux class of the remark that introduces
the coherent-flux exteriors.
\end{remark}

\medskip\noindent\emph{Route.}
A proof of \eqref{8.10:eq-good-set-complement-1} would have to price under $P_{\mathcal{M}}$ each
part of \eqref{8.10:eq-good-set-complement-2}: the curvature-bad, margin-bad and energy-bad
collars through their own collar action; the event $\mathrm{Adm}^{c}$, which by the format
statement for $\mathcal{M}$ has $P_{\mathcal{M}}$-measure zero when every window held by
$\lambda$ off the core of $C$ (a window in the sense of the old-member statement, not the
reference window of Definition~\ref{1.2:def-flow-constants}) is a window of $\mathcal{M}$; the
non-fillability of each
independent abelian linking class of the collar; the coherent-flux class; the remaining
obstructions to a margin-$\mu$ filling, among them torsion classes, which also obstruct for
$\mathrm{SU}(N)\supset\mathrm{U}(1)$, and obstructions of commutator type; and the collar
curvature forced by kept structure near $C$, which is to be excluded from the curvature test and
from the filling and priced from a quadratic reserve. It would then have to show that the total
exponent stays within $\tfrac14 B^{\flat}_\tau(C)$ under a one-sided smallness condition on the
coupling.

\begin{proposition}[Typicality of fillable exteriors under the hole law. Proof outstanding]
\label{8.10:prop-hole-law-typicality}
Let $C$ be a touch-connected union of terminal cells, let $\lambda$ be a fibre variable of $C$, and
let $w\in\mathcal{U}_\lambda$ be the held variables of the fibre $\lambda$. Let
$\mathcal{V}^{\mathrm{ext}}$ be the space of exterior
configurations $V^{\mathrm{ext}}$, with its Haar measure $dV^{\mathrm{ext}}$. Write
$m_\lambda(V^{\mathrm{ext}},w)$ for the integral of the hole integrand
$M_{\lambda,C}(w,V^{\mathrm{in}},V^{\mathrm{ext}})$ over the interior configurations
$V^{\mathrm{in}}$.

Let $P_M$ be the probability on $\mathcal{V}^{\mathrm{ext}}$ whose density is proportional to
$m_\lambda(\cdot,w)$. This is meant either for fixed $\lambda$ and $w$, or, in the mixture
reading, for the law mixed over $\lambda$ with the mixture's weights. The measure $P_M$ is the law
defined by $M_{\lambda,C}$; we call it the \emph{hole law}.

Recall the good set of exteriors of $C$ on the fibre $\lambda$ at filling energy
$\mathfrak{E}$:
\begin{equation*}
G_{\lambda,C}(\mathfrak{E})
=\mathrm{Adm}\cap G^{\mathrm{av}}_{f_*}\cap\mathrm{Fill}_{\mu,\mathfrak{E}},
\end{equation*}
where the three sets on the right are the following:
\begin{itemize}
\item[(a)] $\mathrm{Adm}$ is the admissible set;
\item[(b)] $G^{\mathrm{av}}_{f_*}$ is the set of exteriors whose cell averages of collar curvature
are below $f_*$;
\item[(c)] $\mathrm{Fill}_{\mu,\mathfrak{E}}$ is the set of exteriors admitting a margin-$\mu$
filling whose filling energy $E^\mu_{\mathrm{fill}}$ is at most $\mathfrak{E}$.
\end{itemize}
Let $\mathfrak{E}_*(C)$ be the energy threshold of fillings for $C$. Let $B^\flat_\tau(C)$ be the
flat reserve of $C$; the quarter $\tfrac14 B^\flat_\tau(C)$ of it is left unspent by the reserve
allocation of the old-member statement.

Then there is $\Lambda_C$ with
\begin{equation}\label{8.10:eq-hole-law-typicality-1}
P_M\bigl(G_{\lambda,C}(\mathfrak{E}_*(C))\bigr)\ \ge\ e^{-\Lambda_C},
\qquad \Lambda_C\ \le\ \tfrac14 B^\flat_\tau(C).
\end{equation}
This holds uniformly in $K$, $m$, $N_T$, the position of $C$, and $\lambda$, on the fibres $\lambda$
satisfying $K_{\mathrm{free}}(C;\lambda;p_1+p_2+p_W)$; that condition says that no kept structure of
$\lambda$ lies within cell-distance $p_1+p_2+p_W$ of $C$. In addition, a lower bound of the form
$e^{-\Lambda_C}$ on $P_M\bigl(G_{\lambda,C}(\mathfrak{E}_*(C))\bigr)$ holds on the fibres $\lambda$
violating $K_{\mathrm{free}}(C;\lambda;p_1+p_2+p_W)$, namely those having kept structure in cells
within cell-distance $p_1+p_2+p_W$ of $C$, with the collar clauses of the good set modified as in
the rider to the old-member statement, which is stated there with its proof incomplete, and with
the cost of the kept structure near $C$ priced from a quadratic reserve.

The complement of $G:=G_{\lambda,C}(\mathfrak{E}_*(C))$ whose probability is to be bounded is
described by
\begin{equation}\label{8.10:eq-hole-law-typicality-a}
\begin{gathered}
G^c=\mathrm{Adm}^c\cup\bigl(G^{\mathrm{av}}_{f_*}\bigr)^c
\cup\mathrm{Fill}_{\mu,\mathfrak{E}_*}^c,\\
\mathrm{Fill}^c\supseteq\bigl(G^{\mathrm{mar}}_\mu\bigr)^c
\cup\bigl\{E^\mu_{\mathrm{fill}}>\mathfrak{E}_*\bigr\}\cup\mathfrak{N}\mathfrak{F},
\end{gathered}
\end{equation}
where $\mathfrak{E}_*=\mathfrak{E}_*(C)$ and $G^{\mathrm{mar}}_\mu$ is the margin good set.
The last set is
\begin{equation*}
\mathfrak{N}\mathfrak{F}:=\bigl\{V^{\mathrm{ext}}\in\mathcal{V}^{\mathrm{ext}}:
\text{no margin-}\mu\text{ filling of }V^{\mathrm{ext}}\text{ exists}\bigr\}.
\end{equation*}
Under the hypotheses of the statements named in (i) and (ii) below, the construction behind (ii)
being carried out in the abelian continuum caricature only, $\mathfrak{N}\mathfrak{F}$ contains
the following classes:
\begin{itemize}
\item[(i)] the linked class of the proposition on non-fillable collar-flat exteriors and of
its instances from the obstruction group;
\item[(ii)] the coherent-flux class of the remark on coherent flux.
\end{itemize}
\end{proposition}

Proof outstanding.

\medskip
\noindent\textit{Route.}
A proof has to price each part of the complement \eqref{8.10:eq-hole-law-typicality-a} under
$P_M$, within the allowance $\tfrac14 B^\flat_\tau(C)$ of \eqref{8.10:eq-hole-law-typicality-1}.
The following reasons for the price of each part are available.

\begin{enumerate}
\item[(1)] \emph{Curvature-, margin- and energy-bad collars.} Here and below $x$ is the inverse
square coupling at which the degree function $p_0(\cdot)$ is read, with $x\ge X$; here
$p_0(x)=A_0(\log x)^{\mathfrak{p}_0}$, $\mathfrak{p}_0$ denoting its exponent; and
$\ell\le\ell_*(x)$ is the cell side in lattice units. Under $P_M$ curvature-, margin- and
energy-bad collars are penalised by their own collar action: at the rate $x f_*^2\ell^d$ per bad
cell, a positive multiple of $p_0(x)$, and by action of type $x\varepsilon_\mu^2$ per violated
plaquette.

The set $\mathrm{Adm}^c$ is $P_M$-null whenever, off the core of $C$, every plaquette whose
window $\lambda$ holds is windowed by $M_{\lambda,C}$; this follows from the product format of
$M_{\lambda,C}$.

\item[(2)] \emph{The linked class.} Each independent abelian linking class $\gamma$ of the
collar carries an exterior transport $T_\gamma$. Apply the flux bound for margined fillings to a
margin-$\mu$ filling, along the null-homotopy furnished by the instances from the obstruction
group; let $A_\gamma$ be the area of that null-homotopy. On $\mathrm{Fill}$ this gives
\begin{equation*}
\|T_\gamma-1\|\le A_\gamma\varepsilon_\mu .
\end{equation*}

Under $P_M$ the bonds of the core of $C$ are unweighted. The exterior action terms that could
correlate $T_\gamma$ with the rest are supported on loops encircling the hole $\mathsf{R}$. Hence
$T_\gamma$ is, to leading order, Haar-distributed over $\mathrm{SU}(N)$, or over a coset of a
torus in $\mathrm{SU}(N)$.

For small $\eta>0$ the set $\{T:\|T-1\|\le\eta\}$ has Haar measure comparable to
$\eta^{\dim\mathfrak{g}}$.
Taking $\eta=A_\gamma\varepsilon_\mu$ for each independent class, the expected cost is
\begin{equation}\label{8.10:eq-hole-law-typicality-b}
\begin{split}
-\log P_M(\text{all classes fillable})
\ \lesssim\ &(\text{number of independent classes})\cdot\dim\mathfrak{g}\\
&\cdot\Bigl(\tfrac12\log x-\mathfrak{p}_0\log\log x+O\bigl(\log(A_\gamma B_W)\bigr)\Bigr).
\end{split}
\end{equation}

The number of classes is at most the lattice-scale count $\operatorname{lk}^\sharp(C)$. Part (a)
of the lemma on the ranks of the obstruction groups proves
\begin{equation*}
\operatorname{lk}^\sharp(C)\le c_d(2p_1+3)^d\bigl(\ell_*(x)+1\bigr)^d n(C).
\end{equation*}
The number of classes is expected to be the cell-scale count $\operatorname{lk}(C)$, with
\begin{equation*}
\operatorname{lk}(C)\le c_d(2p_1+1)^d n(C);
\end{equation*}
the equality of the two ranks is not proved.

Each class contributes degree one in $\log x$, and $\ell_*(x)$ has degree $r_{\mathrm{pr}}$ in
$\log x$. So even with the lattice-scale count the cost has degree at most
$d\,r_{\mathrm{pr}}+1$ in $\log x$ per cell of $C$. The unspent $\tfrac14 B^\flat_\tau(C)$ has
degree $2\mathfrak{p}_0$ per cell, and the degree hypothesis $2\mathfrak{p}_0>d\,r_{\mathrm{pr}}+1$
is strict. The cost is therefore absorbable under a further coupling floor, of the same
one-sided kind as the floor under which the budget of the old-member statement holds. With the
cell-scale count the cost has degree one in $\log x$ per cell.

Torsion classes, which also obstruct for $\mathrm{SU}(N)\supset\mathrm{U}(1)$, and
commutator-type obstructions are not counted here.

\item[(3)] \emph{The coherent-flux class.} Let the hull of $C$ be a long box of length $T$ and
width $h$. The exteriors of the remark on coherent flux are then collar-flat and cheap in collar
action per cell, so they are not covered by item~(1).

Consider the abelian Gaussian caricature of $P_M$: a massless Gaussian field of width
$x^{-1/2}$ off the hull.
Let $\mathcal{L}$ be a central cross-section loop, of perimeter comparable to $T+h$. The flux
$\oint_{\mathcal{L}}A$ has variance of order $x^{-1}(T+h)$, by the perimeter law of the free
abelian theory. Non-fillability at margin $\mu$ requires
\begin{equation*}
\Bigl|\oint_{\mathcal{L}}A\Bigr|\ \gtrsim\ Th\,\varepsilon_\mu-O\bigl((T+h)x^{-1/2}\bigr)
\ \asymp\ Th\,x^{-1/2}p_0(x)(1-\mu)/B_W .
\end{equation*}
The Gaussian tail of the caricature then prices the class, per cross-section and for some
constant $c>0$, at
\begin{equation*}
-\log P\ \gtrsim\ c\,Th^2p_0(x)^2(1-\mu)^2/B_W^2 ,
\end{equation*}
$P$ denoting the law of the caricature. This price is far inside the allowance. This is a
caricature, not a theorem about the law defined by $M_{\lambda,C}$.

\item[(4)] \emph{Kept structure near $C$.} Consider a kept component created at level $K'$
within the collar, on a fibre violating $K_{\mathrm{free}}(C;\lambda;p_1+p_2+p_W)$. It forces a
collar curvature of order $p_0(x)^2(M\check{R})^{-4}x^{-1}$, up to a constant, where $M=L^{m'}$.
This curvature must be excluded from the curvature test and from the filling. It must be priced
from a quadratic reserve.
\end{enumerate}

In words, a proof has to show the following. The cost $\Lambda_C$ in
\eqref{8.10:eq-hole-law-typicality-1} must account for, besides the typicality of the collar
curvature, the Haar cost of fillability, of degree one in $\log x$ per independent linking class,
together with the coherent-flux and kept-structure clauses. It may be established either in the
mixture reading, or fibre by fibre, the per-fibre typicality statement being restated with the
collar clauses taken off kept structure. The arithmetic above shows that these costs are of lower
degree in $\log x$ than the reserves, under a one-sided smallness condition on the coupling.
Items (1)--(4) are reasons, not proofs.

\begin{remark}[Splitting the exterior integral along the fillable good set]
\label{8.10:rem-good-set-split}
Fix the level $K'$, a touch-connected union $C$ of terminal cells, and a non-empty set
$\mathcal{S}$ of histories of $C$. For a fibre $\lambda$ of $C$ put
\begin{equation*}
G_{\lambda,C}(\mathfrak{E}_*(C)) := \mathrm{Adm}\cap G^{\mathrm{av}}_{f_*}
\cap\mathrm{Fill}_{\mu,\mathfrak{E}_*(C)},
\qquad \Gamma_C:=\Gamma_C(\mathfrak{E}_*(C)),
\end{equation*}
and write
\begin{equation*}
m_\lambda(V^{\mathrm{ext}},w)=\int M_{\lambda,C}(w,V^{\mathrm{in}},V^{\mathrm{ext}})\,dV^{\mathrm{in}},
\qquad
a_\lambda(V^{\mathrm{ext}},w)=\int \rho_\lambda(w,V^{\mathrm{in}},V^{\mathrm{ext}})\,dV^{\mathrm{in}}
\end{equation*}
for the two local integrals: $m_\lambda$, the hole-law density (the local integral of the
majorant's hole integrand), and $a_\lambda$, the local integral of the effective density with $C$
absent. Here $V^{\mathrm{in}}$ and $V^{\mathrm{ext}}$ are the interior and exterior
configurations, and $w$ denotes the further variables of the fibre, ranging over a set
$\mathcal{U}_\lambda$ and integrated against $dw$. Let $P_M$ be the probability on
$\mathcal{V}^{\mathrm{ext}}$ with density proportional to
$m_\lambda(\cdot,w)$ (the hole law), for fixed $\lambda$ and $w$ or mixed over $\lambda$ (the
mixture reading). Assume the following.
\begin{itemize}
\item[(a)] The old-member statement (good-exterior removal domination), taken at
$\mathfrak{E}=\mathfrak{E}_*(C)$,
which is proved as an implication from the bound of $m_\lambda$ by a multiple of $a_\lambda$
on the good set and from the bound on the increment $\delta_C$: for every $K$, $m$, $N_T$ and
every position of $C$, every fibre $\lambda$ with
$K_{\mathrm{free}}(C;\lambda;p_1+p_2+p_W)$, every $w\in\mathcal{U}_\lambda$ and every
$V^{\mathrm{ext}}\in G_{\lambda,C}(\mathfrak{E}_*(C))$, uniformly in $K$, $m$, $N_T$ and the
position of $C$, with $\Gamma_C$ reading only $n(C)$, $\ell$, $x$, $p_1$, $p_2$,
$\mathfrak{E}_*(C)$ and fixed constants,
\begin{equation}\label{8.10:eq-good-set-split-1}
\int e^{\delta_C}M_{\lambda,C}(w,V^{\mathrm{in}},V^{\mathrm{ext}})\,dV^{\mathrm{in}}
\le e^{\Gamma_C}\int\rho_\lambda(w,V^{\mathrm{in}},V^{\mathrm{ext}})\,dV^{\mathrm{in}} .
\end{equation}
\item[(b)] The budget lemma, in its consequence for the penalty factor
$\operatorname{pen}^\tau_\theta(C,\mathcal{S})$:
\begin{equation}\label{8.10:eq-good-set-split-2}
\operatorname{pen}^\tau_\theta(C,\mathcal{S})\,e^{\Gamma_C}
\le w_\theta(\mathcal{S})\exp\Bigl(-(2-\theta)(\underline{p}_0+1)-\kappa_1 d_{K'}(C)
-(2d\log 3+\log 2)\,n(C)-\tfrac14 B^\flat_\tau(C)\Bigr).
\end{equation}
\item[(c)] The comparison of the running couplings at the seeds with the coupling at level $K'$,
in the form $\underline{p}_0+1\ge p_0(g_{K'})$.
\item[(d)] The numerator form, read with the seeds' linear reserve factors as well as their
weight factors: the exterior integrals of both local integrals are integrals
against the common non-negative measure $dV^{\mathrm{ext}}\,dw$, with non-negative integrands.
\item[(e)] The typicality input of Proposition~\ref{8.10:prop-hole-law-typicality}, whose
proof is outstanding, delivered with
\begin{equation}\label{8.10:eq-good-set-split-3}
\Lambda_C\le\tfrac14 B^\flat_\tau(C):
\end{equation}
for each fibre $\lambda$ with $K_{\mathrm{free}}(C;\lambda;p_1+p_2+p_W)$, uniformly in $K$,
$m$, $N_T$, the position of $C$ and such $\lambda$,
$P_M\bigl(G_{\lambda,C}(\mathfrak{E}_*(C))\bigr)\ge e^{-\Lambda_C}$; in the mixture reading,
the same lower bound for the event $\widehat{G}$ defined below under the hole mixture.
\end{itemize}
Then the following hold.
\begin{itemize}
\item[(i)] For every fibre $\lambda$ with $K_{\mathrm{free}}(C;\lambda;p_1+p_2+p_W)$, with
$G:=G_{\lambda,C}(\mathfrak{E}_*(C))$,
\begin{equation}\label{8.10:eq-good-set-split-4}
\iint_{V^{\mathrm{ext}}\in G} e^{\delta_C}\,m_\lambda\,dV^{\mathrm{ext}}\,dw
\le e^{\Gamma_C}\iint_{V^{\mathrm{ext}}\in G} a_\lambda\,dV^{\mathrm{ext}}\,dw ,
\end{equation}
and the factor $\operatorname{pen}^\tau_\theta(C,\mathcal{S})\,e^{\Gamma_C}$ is bounded
by~\eqref{8.10:eq-good-set-split-2}.
\item[(ii)] With the cost $\Lambda_C$ of~(e),
\begin{equation}\label{8.10:eq-good-set-split-5}
\begin{aligned}
&\operatorname{pen}^\tau_\theta(C,\mathcal{S})\,e^{\Gamma_C+\Lambda_C}\\
&\quad\le w_\theta(\mathcal{S})\exp\Bigl(-(2-\theta)(\underline{p}_0+1)-\kappa_1 d_{K'}(C)
-(2d\log3+\log2)\,n(C)\Bigr)\,e^{\Lambda_C-\frac14 B^\flat_\tau(C)}\\
&\quad\le w_\theta(\mathcal{S})\exp\Bigl(-(2-\theta)(1+\beta_0)^{-1}p_0(g_{K'})
-\kappa_1 d_{K'}(C)-(2d\log3+\log2)\,n(C)\Bigr).
\end{aligned}
\end{equation}
\item[(iii)] In the mixture reading, put
\begin{equation*}
\widehat{G}:=\bigl\{(\lambda,V^{\mathrm{ext}},\cdot):K_{\mathrm{free}}(C;\lambda;p_1+p_2+p_W)
\text{ and }V^{\mathrm{ext}}\in G_{\lambda,C}(\mathfrak{E}_*(C))\bigr\}.
\end{equation*}
Then, the essential infimum being taken with respect to the hole mixture,
\begin{equation}\label{8.10:eq-good-set-split-6}
\operatorname*{ess\,inf}_{\widehat{G}}\frac{a_\lambda}{m_\lambda}\ge e^{\delta_C-\Gamma_C}.
\end{equation}
\end{itemize}
On the complement of $G_{\lambda,C}(\mathfrak{E}_*(C))$, decomposed in
\eqref{8.10:eq-hole-law-typicality-a} as the union of $\mathrm{Adm}^c$,
$(G^{\mathrm{av}}_{f_*})^c$ and $\mathrm{Fill}^c$, the last of which contains the linked class
and the coherent-flux class of exteriors admitting no margin-$\mu$ filling, and on the fibres
violating
$K_{\mathrm{free}}(C;\lambda;p_1+p_2+p_W)$, nothing follows from~(a)--(d): there the bound must
be delivered by the typicality input at the costs priced in
Proposition~\ref{8.10:prop-hole-law-typicality}, whose proof is outstanding (for the Haar cost of
fillability see~\eqref{8.10:eq-hole-law-typicality-b}). The exterior bound at level $K'$
required of $C$ then follows from~(a), that input and~(d) by the argument of the first part of
the rarity proposition for terminal components; off the good set and on the fibres violating
$K_{\mathrm{free}}$ no pointwise domination is available, so the probabilistic ingredient remains
genuinely needed there. In the mixture reading, the bound at level $K'$ required of the
ratio $a_\lambda/m_\lambda$ follows from~(iii) and the typicality of $\widehat{G}$ under the hole
mixture at the same costs. The Jensen and exponential-moment forms of that bound require
$\mathbb{E}[-\log(a_\lambda/m_\lambda)]<\infty$ under the hole mixture; this fails where
$a_\lambda=0<m_\lambda$ on a set of positive measure, as in the hollow-shell instance, unless the
expectation is restricted likewise to $\widehat{G}$.

The coupling floor entering~(b) is the one of the budget lemma; it rests on the comparison
$d\,r_{\mathrm{pr}}+1<2\mathfrak{p}_0$ between degrees in $\log x$. This comparison is strict.
The smallness for $C$ obtained from~\eqref{8.10:eq-good-set-split-5} together with the
typicality input, by the remaining steps of the rarity theorem for terminal components and of
its corollary, is at most the smallness
$\varepsilon_{K'}(C,\mathcal{S})$ required of $C$ at level $K'$, for every value
$C_{\mathrm{rm}}\ge 0$ of the removal constant.

For fibres violating $K_{\mathrm{free}}(C;\lambda;p_1+p_2+p_W)$ the old-member statement
supplies no pair $(G,\Gamma)$. One possible reorganisation is to remove the maximal
cluster of live components chained to $C$ at touch-distance at most $c_K$, a fixed chaining
constant, so that
$K_{\mathrm{free}}$ holds by maximality. It is not carried out here; it would have to be
compatible with the injectivity and the lineage-independence of the summation over histories.

The self-penalisation of bad exteriors by their own collar action is used only for
$(G^{\mathrm{av}}_{f_*})^c$ and for the energy- and margin-bad classes, as in part~(iii) of the
corollary to the hollow-shell instance. The pointwise domination on the good set is used only in
the form~(a). This remark does not provide the bound of $m_\lambda$ by a multiple of $a_\lambda$
on the good set, on
which~(a) depends. It provides no statement at a level $k<K'$, and neither the
label-conditional nor the joint form of the estimate.
\end{remark}

\begin{proof}
(i) Fix a fibre $\lambda$ satisfying $K_{\mathrm{free}}(C;\lambda;p_1+p_2+p_W)$ and put
$G:=G_{\lambda,C}(\mathfrak{E}_*(C))$. For every $w\in\mathcal{U}_\lambda$ and every
$V^{\mathrm{ext}}\in G$, hypothesis~(a) gives~\eqref{8.10:eq-good-set-split-1}. The factor
$e^{\delta_C}$ does not depend on $V^{\mathrm{in}}$, so the left member
of~\eqref{8.10:eq-good-set-split-1} is
$e^{\delta_C}m_\lambda(V^{\mathrm{ext}},w)$. Its right member is
$e^{\Gamma_C}a_\lambda(V^{\mathrm{ext}},w)$. By~(d), both members are integrated against the same
non-negative measure $dV^{\mathrm{ext}}\,dw$, and the integrands are non-negative.

By Tonelli's theorem the iterated integrals over $V^{\mathrm{in}}$ and over
$(V^{\mathrm{ext}},w)$ therefore exist in $[0,\infty]$ and do not depend on the order of
integration. Integrating the pointwise inequality over $\{V^{\mathrm{ext}}\in G\}$ preserves it,
and the constant $e^{\Gamma_C}$ comes out of the integral. This
is~\eqref{8.10:eq-good-set-split-4}. In the numerator form the term indexed by
$(C,\mathcal{S})$ carries the factor $\operatorname{pen}^\tau_\theta(C,\mathcal{S})$; after
\eqref{8.10:eq-good-set-split-4} the bound carries
$\operatorname{pen}^\tau_\theta(C,\mathcal{S})\,e^{\Gamma_C}$, which is bounded by~(b), that is,
by~\eqref{8.10:eq-good-set-split-2}.

(ii) Multiplying~\eqref{8.10:eq-good-set-split-2} by $e^{\Lambda_C}$ and writing the factor
$e^{-\frac14B^\flat_\tau(C)}$ of its right member together with $e^{\Lambda_C}$ gives the first
inequality of~\eqref{8.10:eq-good-set-split-5}. For the second, two facts are used.

First, \eqref{8.10:eq-good-set-split-3} gives $\Lambda_C-\frac14B^\flat_\tau(C)\le0$, so
$e^{\Lambda_C-\frac14B^\flat_\tau(C)}\le1$. Second, hypothesis~(c) gives
$\underline{p}_0+1\ge p_0(g_{K'})$. Together with $(1+\beta_0)^{-1}\le1$, this replaces the term
$-(2-\theta)(\underline{p}_0+1)$ in the exponent by $-(2-\theta)(1+\beta_0)^{-1}p_0(g_{K'})$.
The terms $-\kappa_1d_{K'}(C)$ and $-(2d\log3+\log2)\,n(C)$ are unchanged.

(iii) Let $(\lambda,V^{\mathrm{ext}},\cdot)\in\widehat{G}$. Then $\lambda$ satisfies
$K_{\mathrm{free}}(C;\lambda;p_1+p_2+p_W)$ and $V^{\mathrm{ext}}\in G_{\lambda,C}(\mathfrak{E}_*(C))$,
so hypothesis~(a) applies. As in~(i), it reads
$e^{\delta_C}m_\lambda\le e^{\Gamma_C}a_\lambda$, whence
$a_\lambda/m_\lambda\ge e^{\delta_C-\Gamma_C}$ wherever $m_\lambda>0$.

The hole mixture has density proportional to $m_\lambda$. Hence the part of $\widehat{G}$ on which
$m_\lambda=0$ is a null set for it, and the lower bound holds almost everywhere on $\widehat{G}$.
This is~\eqref{8.10:eq-good-set-split-6}.
\end{proof}

\begin{definition}[Hole, fixed exterior and window constraints]\label{8.11:not-hole-exterior-windows}
Throughout the remainder of this chapter the following objects are fixed.

\emph{The fine lattice and the zone.} The dimension is $d=4$. We write $k_o$ for the observation
level and $T_\eta$ for the fine lattice of spacing $\eta=L^{-k_o}$, and we put $n:=\eta^{-1}$, the
number of fine spacings per top unit. The zone is $H=R^{+}_{k_o}$, the zone at level $k_o$; it is a
box of top blocks whose sides $D_1,\dots,D_4$ satisfy
\begin{equation*}
  D_W/\varrho_H\;\le\; D_i\;\le\; D_W,\qquad i=1,\dots,4,
\end{equation*}
with $\varrho_H\le 2$. The bound $\varrho_H\le2$ holds because the sides of the zone differ by at
most the size $2M\check R\hat r$ of its core, while
\begin{equation*}
  D_W\;\ge\;8LM\check R\check r
\end{equation*}
by the size bounds of the zone.

\emph{The exterior.} The configuration $e$ is the top exterior held fixed; it is admissible and
isolated, in the sense of the standing scope inherited throughout this chapter.

\emph{The support hypothesis on the exterior.} We assume that $e=V^{\sharp}{\restriction}\mathfrak{E}$
is the restriction of a support
configuration $V^{\sharp}$ to the exterior's data set $\mathfrak{E}$, and that
$w^{\sharp}:=V^{\sharp}{\restriction}\mathfrak{h}$, the restriction of $V^{\sharp}$ to the set
$\mathfrak{h}$ of hole data, is a support point: every window of the weight passes strictly at
$w^{\sharp}$. In particular the set $S(e)$ of support points of $e$ is non-empty. We put
$\varepsilon_1:=\varepsilon_1(V^{\sharp})$, the regularity of $V^{\sharp}$. The domain
$\mathsf{K}$ of hole data and the constants $\varepsilon_W$, $\hat N_W$, $\Theta_W$, $C_E$ (the
block count and the energy constants of the zone) are those fixed with the zone. For hole data
$x\in\mathsf{K}$, $e\oplus x$ denotes the configuration on $\mathfrak{E}\cup\mathfrak{h}$ assembled
from the exterior $e$ and the hole data $x$, as fixed with the zone, and
\begin{equation*}
  \mathbf{A}(x):=A\bigl(U(e\oplus x)\bigr).
\end{equation*}

\emph{The clamp width.} We put $w:=M_2R_{k_o}$, where $R_{k_o}$ is the size parameter of level
$k_o$; then $w\ge2$. In this chapter $w$ denotes the clamp width, not the complex source
$w=\sum_i z_if_i$.

\emph{Straddling sites.} $\mathfrak{S}$ is the set of straddling sites. For $a\in\mathfrak{S}$ we
put $g_a(x):=f_a(e\oplus x)$, where $f_a$ is the window functional of the site $a$ (fixed with the
zone), measured in units of $\varepsilon_{k_o}\eta^2$. The test region
$\tilde a$ of $a$ meets $H$ only at depth at most $w$, that is $\tilde a\cap H\subset\bar
S_1=\{\text{depth}\le w\}$, and $a^{\sim4}$ is the dependence region of $a$.

\emph{The young family.} $\mathfrak{Y}$ is a young family (in the sense of the age vocabulary of
this chapter) consisting of sites $\square'$ of level
$k_o-1$ with $\square'^{\sim4}\subset H$. For $\square'\in\mathfrak{Y}$ we put
$g_{\square'}(x):=f_{k_o-1,\square'}(e\oplus x)$, where $f_{k_o-1,\square'}$ is the
level-$(k_o-1)$ window functional of $\square'$, measured in units of
\begin{equation*}
  \varepsilon_{k_o-1}\,\eta_{(k_o-1)}^2\;=\;\varepsilon_{k_o-1}L^2\eta^2 .
\end{equation*}

\emph{The constrained class and its minimiser.} For a threshold $\theta$ we put
\begin{equation}\label{8.11:eq-hole-exterior-windows-1}
  \widehat{\mathcal{W}}^{\theta}
  :=\bigl\{x\in\mathsf{K}\;:\;g_a(x)\le\theta\ \text{for all } a\in\mathfrak{S}\cup\mathfrak{Y}\bigr\},
  \qquad
  \hat m(\theta):=\min_{\widehat{\mathcal{W}}^{\theta}}\mathbf{A},
\end{equation}
and we let $\hat x=\hat x^{\theta}$ denote any minimiser of $\mathbf{A}$ on
$\widehat{\mathcal{W}}^{\theta}$. For $\theta\ge1$ the set $\widehat{\mathcal{W}}^{\theta}$ is
compact and contains $w^{\sharp}$ (every window passes strictly at the support point
$w^{\sharp}$), and $\mathbf{A}$ is continuous; hence the minimum in
\eqref{8.11:eq-hole-exterior-windows-1} is attained and $\hat x$ exists. If $\mathfrak{Y}=\emptyset$,
a minimiser $\hat x$ is a straddling-only minimiser and is written $\tilde x$. If $\mathfrak{Y}$ is
the first young row, then $\hat x$ is a canonical comparison point.

\emph{Coordinates and signed depth.} We unwrap a neighbourhood of $H$ and use fine sites
$x\in\mathbb{Z}^4$, with the sites of $H$ forming the box
\begin{equation*}
  \prod_{i=1}^{4}\,[0,\mathsf{D}_i],\qquad \mathsf{D}_i:=D_i\,n.
\end{equation*}
The signed depth of a site is
\begin{equation}\label{8.11:eq-hole-exterior-windows-2}
  \delta(x):=\eta\,\min_{i}\,\min\bigl(x_i,\ \mathsf{D}_i-x_i\bigr),
\end{equation}
measured in top units; it is negative outside $H$. For $c\in\eta\mathbb{Z}$ we put
$\Omega_c:=\{x:\delta(x)\ge c\}$, which is a box, and $\Sigma_c:=\partial\Omega_c$. A cell lies in
a region if all its corners do, and for a plaquette $p$ we put
$\delta(p):=\min\{\delta(x):x\ \text{a corner of}\ p\}$. The block depth used elsewhere in this
chapter differs from $\delta$ by less
than one block; this difference plays no role in what follows.

\emph{Standing scope.} Since the zone is isolated and the cube carrying the constrained class is
strictly smaller than the torus at every level, the box $\Omega_{-(3w+2)}$ embeds in $T_\eta$
without wrapping and meets no other structure retained by the expansion; this is assumed
throughout.
\end{definition}

\begin{definition}[The first young row covers a depth band]\label{8.11:def-young-row-cover}
Let the zone $H=R^{+}_{k_o}$, the clamp width $w=M_2R_{k_o}$, the signed depth $\delta$, the young
family $\mathfrak{Y}$ and, for each young cube $\square'\in\mathfrak{Y}$, its test region
$\tilde{\square}'$ be as in Definition~\ref{8.11:not-hole-exterior-windows}. For real numbers
$c\le c'$ we write $\{\delta\le c'\}$ for the set of points of signed depth at most $c'$, and
$\{c\le\delta\le c'\}$ for the set of points whose signed depth lies in $[c,c']$.

We say that the young family $\mathfrak{Y}$ \emph{covers a depth band} if there are real numbers
$y_-$ and $y_+$ with
\begin{equation*}
y_-\le w-\frac{1}{L},\qquad y_+\ge w,
\end{equation*}
such that
\begin{equation}\label{8.11:eq-young-row-cover-a}
\{y_-\le\delta\le y_+\}\cap H\;\subset\;\bigcup_{\square'\in\mathfrak{Y}}\tilde{\square}'
\;\subset\;\{\delta\le y_+\}.
\end{equation}
The inclusion \eqref{8.11:eq-young-row-cover-a} is a hypothesis on the configuration; it is
assumed, and cited by this label, wherever the first young row is used below. The first young
row is the young family of Definition~\ref{8.11:not-hole-exterior-windows} at which the canonical
comparison point $\hat x$ is taken.

When $\mathfrak{Y}$ is the first young row and $r_R=1$, the case in which the radii of
Definition~\ref{8.11:not-hole-exterior-windows} do not jump between the levels $k_o-1$ and $k_o$,
one puts $\lambda_y=w/L$ and takes the depth range $(y_-,y_+)$ as follows:
\begin{equation}\label{8.11:eq-young-row-cover-1}
\lambda_y=\frac{w}{L},\qquad y_-=3\lambda_y,\qquad y_+=w+\frac{5w}{L}.
\end{equation}
For $\mathfrak{Y}=\emptyset$ we put $y_+:=w$. In every case we put
\begin{equation*}
\varpi_y:=y_+-w .
\end{equation*}

Two alternative choices of the depth range are treated later in this chapter: in the first the
radii of Definition~\ref{8.11:not-hole-exterior-windows} jump between the levels $k_o-1$ and $k_o$,
$R_{k_o-1}=L\,R_{k_o}$, and one takes
$y_-=3w$; in the second the young rows are aligned with cubes, and one takes
\begin{equation*}
y_+=2w+\frac{w}{L}.
\end{equation*}
\end{definition}

\begin{definition}[Units, young cap and exterior amplitude]\label{8.11:not-units-cap-amplitude}
We keep the objects of Definition~\ref{8.11:not-hole-exterior-windows}: the top level $k_o$, the fine
spacing $\eta=L^{-k_o}$, the zone $H=R^{+}_{k_o}$, the admissible and isolated top exterior $e$,
the clamp width $w=M_2R_{k_o}$, the threshold $\theta$, and, for each level-$k_o$ site $a$, its
window functional $f_a$, its test region $\tilde a$ and its dependence region $a^{\sim 4}$. We also
use the signed depth $\delta$ of the coordinates unwrapped around $H$, the boxes
$\Omega_c=\{\delta\ge c\}$ and their boundaries $\Sigma_c=\partial\Omega_c$. For each level $k$,
$\varepsilon_k$ is the constant for which the level-$k$ windows are measured in units of
$\varepsilon_k$ times the square of the level-$k$ spacing, and $\vartheta_1$ is a fixed small
relative error.

\emph{Units.} We put $\varepsilon:=\varepsilon_{k_o}$. For a real number $X$, the expression
$X\,[\mathrm{u}]$ stands for $X\varepsilon\eta^2$; this is the exterior unit, the unit of the
level-$k_o$ windows. We put $\theta':=\theta(1+\vartheta_1)$.

\emph{Young cap and exterior amplitude.} We put
\begin{equation}\label{8.11:eq-units-cap-amplitude-a}
\begin{gathered}
\mathsf{T}:=\frac{\varepsilon_{k_o-1}L^2}{\varepsilon_{k_o}}\in\Bigl[\frac{L^2}{1+\beta_0},\,L^2\Bigr],
\\
s_e:=\theta^{-1}\max\bigl\{f_a(e)\;:\;a\text{ is a level-}k_o\text{ site},
\\
a^{\sim4}\cap H=\emptyset,\ \tilde a\cap\Sigma_{-(3w+1)}\neq\emptyset\bigr\}
\\
\in\bigl[0,\theta^{-1}\bigr].
\end{gathered}
\end{equation}
Since $\varepsilon_{k_o}>0$, the inclusion of $\mathsf{T}$ in the displayed interval is the same as
the bound $\varepsilon_{k_o}/(1+\beta_0)\le\varepsilon_{k_o-1}\le\varepsilon_{k_o}$ on consecutive
units. A young window at
level $k_o-1$ is measured in the unit $\varepsilon_{k_o-1}L^2\eta^2$, which equals
$\mathsf{T}\varepsilon\eta^2$; so $\mathsf{T}$ is the young cap in exterior units. A bound $X$ in the
young unit is the bound $X\mathsf{T}\,[\mathrm{u}]$ in the exterior unit. Accordingly, we put
$\theta_y:=\theta(1+\vartheta_1)\mathsf{T}=\theta'\mathsf{T}$; this is the threshold $\theta'$ for the
young windows, expressed in $[\mathrm{u}]$.

Every site $a$ entering the maximum in \eqref{8.11:eq-units-cap-amplitude-a} has dependence region
$a^{\sim4}$ disjoint from $H$. Its window functional reads only the configuration on $a^{\sim4}$
\cite[(1.2)]{B-CMP119}, hence only the exterior $e$. Therefore $s_e$ is a function of
$e$ alone; it is the amplitude of the exterior as seen by these windows.

\emph{Norms and plaquette logarithms.} The symbol $\|\cdot\|$ denotes the operator norm. For a
plaquette $p$ we put $\Phi(p):=\frac1i\log U(\partial p)$. By the plaquette-logarithm comparison,
\begin{equation}\label{8.11:eq-units-cap-amplitude-1}
\|U(\partial p)-1\|\le\|\Phi(p)\|\le2\,\|U(\partial p)-1\|.
\end{equation}

\emph{Further small letters.} We put
\begin{equation*}
\varepsilon_{\min}:=\min(\varepsilon_{k_o-1},\varepsilon_{k_o}),\qquad
\varepsilon_a:=\tfrac12\,\varepsilon_W\Theta_W.
\end{equation*}
Here $\varepsilon_a$ is the class bound in top units, formed from the energy constants
$\varepsilon_W$, $\Theta_W$ of the zone.

\emph{Comparison with the amplitude of a kinked exterior.} The following comparison is stated for
orientation; nothing in the definitions above depends on it. Suppose that the exterior $e$ is the
kinked exterior of amplitude $\sigma$, and let $o_{\gamma}(1)$ denote a quantity that is small
under the one-sided upper bounds on the coupling. By the exterior comparison statement read in the
reverse direction, and by the admissibility lemma, with $B_3$ the constant of the admissibility
lemma, $s_e$ is comparable with $\sigma/\theta$ in the form
\begin{equation}\label{8.11:eq-units-cap-amplitude-2}
\sigma\le\theta s_e+\vartheta_1+o_{\gamma}(1),\qquad \theta s_e\le\sigma B_3\bigl(1+o(1)\bigr).
\end{equation}
The first inequality comes from the exterior comparison statement: three thin
layers away from $H$ the field is the exterior's own, of curvature $\sigma\varepsilon$, up to the
tail of \cite[(190)]{B-CMP102b}. The second inequality is the admissibility lemma. Under
the amplitude form of that lemma one has instead
\begin{equation*}
\theta s_e\le\sigma(1+\kappa_{\mathrm{sp}})(1+\vartheta_1),
\end{equation*}
where $\kappa_{\mathrm{sp}}$ is the relative constant of the amplitude form of the admissibility
lemma.
\end{definition}

\begin{lemma}[Energy bounds below the support value]\label{8.11:lem-energy-below-support}
Let the zone $H=R^{+}_{k_o}$, the exterior $e$, the support point $w^{\sharp}$, the domain $K$ of hole
data, the regularity $\varepsilon_1$, the block count $\hat N_W$ and the energy constant
$C_E=3B_3^2+2$, where $B_3$ is the regularity constant of the support configuration (see
\textup{(E3)} below), be as in Definition~\ref{8.11:not-hole-exterior-windows}, and let $\varepsilon_a$
be the class bound of Definition~\ref{8.11:not-units-cap-amplitude}. Write
$\mathbf{A}(x)=A(U(e\oplus x))$ and $U^{\sharp}=U(e\oplus w^{\sharp})$, so that
$\mathbf{A}(w^{\sharp})=A(U^{\sharp})$. Let $A=A_{\mathrm{in}}+A_{\mathrm{out}}$ be the splitting of
the Wilson action into its inner part $A_{\mathrm{in}}$ and its outer part $A_{\mathrm{out}}$, the
latter carried by the plaquettes lying beyond the comparison length $\ell_W$ of \textup{(E3)} below.
Let $\Omega_{-(3w+2)}$ be the box of
fine sites of signed depth at least $-(3w+2)$, and let
$\mathfrak{e}(U(\partial q))=1-\operatorname{Re}\operatorname{tr}U(\partial q)$ be the energy
of the plaquette $q$, that is, its term in the Wilson action of
Definition~\ref{1.1:def-wilson-action}. Throughout, the one-sided smallness condition involving
$\gamma_W$ under which the
zone and its windows are constructed is in force. Assume:
\begin{enumerate}
\item[(E1)] (energy, class and interiority) for every $x\in K$ with
$\mathbf{A}(x)\le\mathbf{A}(w^{\sharp})$ one has (a)
$A_{\mathrm{in}}(U(e\oplus x))\le C_E\varepsilon_1^2\hat N_W$, (b)
$\|U(e\oplus x)(\partial p)-1\|\le\varepsilon_a\eta^2$ for every fine plaquette $p$, and (c)
$x\in\operatorname{int}K$;
\item[(E2)] (the energy bound at the support point and the far expansion of the outer action,
together with the regularity lemma of the zone) one has
$A_{\mathrm{in}}(U^{\sharp})\le 3B_3^2\varepsilon_1^2\hat N_W$, and for every
$x\in K$,
\begin{equation*}
A_{\mathrm{out}}(U^{\sharp})-A_{\mathrm{out}}(U(e\oplus x))\le 2\varepsilon_1^2\hat N_W ;
\end{equation*}
\item[(E3)] (comparison away from the hole) $U^{\sharp}$ is $B_3\varepsilon_1$-regular, and for $x$
as in \textup{(E1)}, beyond the comparison length $\ell_W$, the configuration $U(e\oplus x)$
differs from $U^{\sharp}$ by the correction field $\mathcal{H}'$ of the comparison, which is
exponentially small, so small that $U(e\oplus x)$ is $2B_3\varepsilon_1$-regular there;
\item[(E4)] (energy per block) if a configuration $U$ is $b$-regular on a block, then
$\sum_{q}\mathfrak{e}(U(\partial q))\le 3b^2$, the sum running over the plaquettes $q$ of that
block;
\item[(E5)] $3w+2\le D_W$.
\end{enumerate}
Then
\begin{equation}\label{8.11:eq-energy-below-support-a}
\mathbf{A}(w^{\sharp})-\mathbf{A}(x)\le C_E\varepsilon_1^2\hat N_W
\qquad\text{for every } x\in K ,
\end{equation}
and, for every $x\in K$ with $\mathbf{A}(x)\le\mathbf{A}(w^{\sharp})$,
\begin{equation}\label{8.11:eq-energy-below-support-b}
\sum_{q\subset\Omega_{-(3w+2)}}\mathfrak{e}\bigl(U(e\oplus x)(\partial q)\bigr)
\le 5C_E\varepsilon_1^2\hat N_W .
\end{equation}
\end{lemma}

\begin{proof}
\emph{The bound \eqref{8.11:eq-energy-below-support-a}.} Fix $x\in K$ and write $U=U(e\oplus x)$.
By the additivity $A=A_{\mathrm{in}}+A_{\mathrm{out}}$,
\begin{equation*}
\mathbf{A}(w^{\sharp})-\mathbf{A}(x)=A(U^{\sharp})-A(U)
=\bigl[A_{\mathrm{in}}(U^{\sharp})-A_{\mathrm{in}}(U)\bigr]
+\bigl[A_{\mathrm{out}}(U^{\sharp})-A_{\mathrm{out}}(U)\bigr].
\end{equation*}
We bound $A_{\mathrm{in}}(U^{\sharp})$ from above, $A_{\mathrm{in}}(U)$ from below, and the outer
difference from above. First, $A_{\mathrm{in}}(U^{\sharp})\le 3B_3^2\varepsilon_1^2\hat N_W$ by
\textup{(E2)}. Second,
$A_{\mathrm{in}}(U)\ge0$. Indeed $A_{\mathrm{in}}(U)$ is a sum of terms
$1-\operatorname{Re}\operatorname{tr}U(\partial p)$ of the Wilson action of
Definition~\ref{1.1:def-wilson-action}. Each of these is non-negative, because
$|\operatorname{tr}V|\le1$ for $V\in\mathrm{SU}(N)$ in the normalisation $\operatorname{tr}1=1$.
Third,
\textup{(E2)} applies because $x\in K$; it gives
$A_{\mathrm{out}}(U^{\sharp})-A_{\mathrm{out}}(U)\le2\varepsilon_1^2\hat N_W$. Adding the three
bounds,
\begin{equation*}
\mathbf{A}(w^{\sharp})-\mathbf{A}(x)
\le 3B_3^2\varepsilon_1^2\hat N_W-0+2\varepsilon_1^2\hat N_W
=(3B_3^2+2)\varepsilon_1^2\hat N_W
=C_E\varepsilon_1^2\hat N_W .
\end{equation*}
No hypothesis on $x$ other than $x\in K$ has been used, so this holds for every $x\in K$.

\emph{The bound \eqref{8.11:eq-energy-below-support-b}.} Let $x\in K$ satisfy
$\mathbf{A}(x)\le\mathbf{A}(w^{\sharp})$, and split the sum over plaquettes
$q\subset\Omega_{-(3w+2)}$ into an inner part and an outer part, according to the splitting of the
action.

For the inner part, each energy $\mathfrak{e}(U(e\oplus x)(\partial q))$ is a term of the Wilson
action and hence non-negative, as shown above for the terms of $A_{\mathrm{in}}$, and the inner
plaquettes $q\subset\Omega_{-(3w+2)}$ are
among the plaquettes whose terms make up $A_{\mathrm{in}}$. Hence the inner part is at most
$A_{\mathrm{in}}(U(e\oplus x))$, and \textup{(E1)}(a) applies to $x$ and gives
$A_{\mathrm{in}}(U(e\oplus x))\le C_E\varepsilon_1^2\hat N_W$.

For the outer part, the plaquettes of the outer part lie beyond $\ell_W$, and \textup{(E3)} says
that there the configuration $U(e\oplus x)$
differs from the $B_3\varepsilon_1$-regular configuration $U^{\sharp}$ by the exponentially small
field $\mathcal{H}'$, so that $U(e\oplus x)$ is $2B_3\varepsilon_1$-regular on the blocks containing
them. Grouping the outer plaquettes by the blocks containing them and applying
\textup{(E4)} with $b=2B_3\varepsilon_1$, each such block carries energy at most
\begin{equation*}
3(2B_3\varepsilon_1)^2=12B_3^2\varepsilon_1^2 .
\end{equation*}
These blocks lie in $\Omega_{-(3w+2)}$. Since $3w+2\le D_W$ by \textup{(E5)}, they are among the
blocks counted by $\hat N_W$. The outer part is therefore at most
$12B_3^2\varepsilon_1^2\hat N_W$. Moreover
\begin{equation*}
12B_3^2\le 12B_3^2+8=4C_E ,
\end{equation*}
so the outer part is at most $4C_E\varepsilon_1^2\hat N_W$.

Adding the inner and the outer part gives
\begin{equation*}
C_E\varepsilon_1^2\hat N_W+4C_E\varepsilon_1^2\hat N_W=5C_E\varepsilon_1^2\hat N_W ,
\end{equation*}
which is \eqref{8.11:eq-energy-below-support-b}.
\end{proof}

\begin{definition}[Lipschitz drop of the constrained minimum value]\label{8.11:def-value-lipschitz}
We use the hole, the fixed exterior and the window constraints of
Notation~\ref{8.11:not-hole-exterior-windows}: the action $\mathbf{A}(x)$ as a function of the hole
data $x$, the straddling sites $\mathfrak{S}$, the young family $\mathfrak{Y}$, and, for
$a\in\mathfrak{S}\cup\mathfrak{Y}$, the window functional $g_a$ of the site $a$. Recall from
Notation~\ref{8.11:not-hole-exterior-windows} that, for a threshold $\theta$,
$\widehat{\mathcal{W}}^{\theta}$ denotes the constrained class of hole data $x$ with
$g_a(x)\le\theta$ for every $a\in\mathfrak{S}\cup\mathfrak{Y}$, and that
$\hat m(\theta)=\min_{\widehat{\mathcal{W}}^{\theta}}\mathbf{A}$.

Let $\vartheta\in(0,1]$ and $\Lambda>0$; here $\vartheta$ denotes an increment of the threshold,
not the separation-window parameter. We say that the constrained minimum value has
\emph{Lipschitz drop} $\Lambda$ at $\theta$ over the increment $\vartheta$, and write that
$(\mathrm{VL})_{\theta}[\vartheta,\Lambda]$ holds, if
\begin{equation}\label{8.11:eq-value-lipschitz-a}
\hat m(\theta)-\hat m(\theta+\vartheta)\le\vartheta\,\Lambda .
\end{equation}
The function $\theta\mapsto\hat m(\theta)$ is non-increasing, since
$\widehat{\mathcal{W}}^{\theta}\subset\widehat{\mathcal{W}}^{\theta'}$ for $\theta\le\theta'$; in
particular the left-hand side of \eqref{8.11:eq-value-lipschitz-a} is non-negative, and
\eqref{8.11:eq-value-lipschitz-a} is a one-sided bound on it.
\end{definition}

\begin{proof}
We verify the monotonicity of $\hat m$ asserted in Definition~\ref{8.11:def-value-lipschitz}.
Let $\theta\le\theta'$ be thresholds at which $\hat m$ is defined in
Notation~\ref{8.11:not-hole-exterior-windows}. If $x\in\widehat{\mathcal{W}}^{\theta}$, then for every
$a\in\mathfrak{S}\cup\mathfrak{Y}$ we have $g_a(x)\le\theta\le\theta'$, so $x$ satisfies every
constraint defining $\widehat{\mathcal{W}}^{\theta'}$, and $x\in\widehat{\mathcal{W}}^{\theta'}$. Hence
$\widehat{\mathcal{W}}^{\theta}\subset\widehat{\mathcal{W}}^{\theta'}$.

Let $\hat x^{\theta}$ be a minimiser of $\mathbf{A}$ on $\widehat{\mathcal{W}}^{\theta}$, so that
$\mathbf{A}(\hat x^{\theta})=\hat m(\theta)$. By the inclusion just proved,
$\hat x^{\theta}\in\widehat{\mathcal{W}}^{\theta'}$, and therefore
\begin{equation*}
\hat m(\theta')=\min_{\widehat{\mathcal{W}}^{\theta'}}\mathbf{A}\le\mathbf{A}(\hat x^{\theta})=\hat m(\theta).
\end{equation*}
Thus $\hat m$ is non-increasing. Taking $\theta'=\theta+\vartheta$ with $\vartheta>0$ gives
$\hat m(\theta)-\hat m(\theta+\vartheta)\ge0$, which is the non-negativity of the left-hand side
of \eqref{8.11:eq-value-lipschitz-a}.
\end{proof}

\begin{remark}[Lipschitz drop from the second-derivative room lemma]
Let $\tau$ be the room parameter of the second-derivative room lemma (not the layer thickness
used later in this chapter), let $q(x',x_1)$ be the quantity of that lemma depending on its points
$x'$ and $x_1$ (not the auxiliary parameter $q$ of $\mathfrak{p}$), and let $\Lambda_{\mathbf{A}}$
be the constant of that lemma. If the second-derivative room condition at the threshold $\theta$
with the parameter $\tau$, in the sense of the second-derivative room lemma, holds and
$\tau>2q(x',x_1)$, then
\begin{equation*}
(\mathrm{VL})_{\theta}\bigl[\vartheta,\ \Lambda_{\mathbf{A}}/(\tau-2q(x',x_1))\bigr]
\end{equation*}
holds for every $\vartheta\in(0,1]$. This implication is taken here by statement from the lemma of
this chapter that derives the Lipschitz drop from the second-derivative room lemma. It applies
without change to the family of constraints indexed by $\mathfrak{S}\cup\mathfrak{Y}$, because its
proof uses the window functionals $g_a$, $a\in\mathfrak{S}\cup\mathfrak{Y}$, only through
property~(iii) of the window functionals in the window-functional lemma and through the
second-derivative room condition for that family.
\end{remark}

\begin{lemma}[Abundance of good thresholds]\label{8.11:lem-good-thresholds}
Let $J=[\theta_a,\theta_b]\subset[1,2]$, and write $|J|=\theta_b-\theta_a$ for its length. Let
$0<\vartheta\le |J|/2$, and put
\begin{equation}\label{8.11:eq-good-thresholds-1}
\Lambda_J:=\frac{4\,C_E\,\varepsilon_1^2\,\hat N_W}{|J|}.
\end{equation}
Then the set
\begin{equation}\label{8.11:eq-good-thresholds-2}
G_J(e,\vartheta):=\bigl\{\theta\in[\theta_a,\theta_b-\vartheta]\;:\;
(\mathrm{VL})_{\theta}[\vartheta,\Lambda_J]\ \text{holds}\bigr\}
\end{equation}
has Lebesgue measure at least $|J|/4$. Here $(\mathrm{VL})_{\theta}[\vartheta,\Lambda]$ is the property of
Definition~\ref{8.11:def-value-lipschitz}, and the argument $e$ is the exterior datum entering the
configurations $U(e\oplus x)$ of Lemma~\ref{8.11:lem-energy-below-support}, on which the
classes $\widehat{\mathcal{W}}^{\theta}$ and their minimum values $\hat m(\theta)$ depend.
The proof below uses only the bound \eqref{8.11:eq-energy-below-support-a} and the monotonicity of
$\theta\mapsto\widehat{\mathcal{W}}^{\theta}$; it uses neither the second-derivative room lemma, nor
any near-affinity property, nor any multiplier, nor any other property of the functionals that define
the classes.
\end{lemma}

\begin{proof}
Since $0<\vartheta\le|J|/2$ and $J\subset[1,2]$, we have $0<\vartheta\le 1/2$, so
$\vartheta\in(0,1]$ and Definition~\ref{8.11:def-value-lipschitz} applies to $\vartheta$; moreover
$|J|\ge 2\vartheta>0$.

The function $\theta\mapsto\hat m(\theta)$ is non-increasing, because
$\widehat{\mathcal{W}}^{\theta}\subset\widehat{\mathcal{W}}^{\theta'}$ for $\theta\le\theta'$ and
$\hat m(\theta)$ is the minimum of $\mathbf{A}$ over $\widehat{\mathcal{W}}^{\theta}$. We bound it at
the two ends of $J$. Since $\theta_a\ge 1$, the support point satisfies
$w^{\sharp}\in\widehat{\mathcal{W}}^{1}\subset\widehat{\mathcal{W}}^{\theta_a}$, whence
$\hat m(\theta_a)\le\mathbf{A}(w^{\sharp})$. At the other end, $\hat m(\theta_b)$ is a minimum of
$\mathbf{A}$ over a subset of the domain of hole data $x$ of
Lemma~\ref{8.11:lem-energy-below-support}, so it is at least the minimum of $\mathbf{A}$ over that
domain. The bound
\eqref{8.11:eq-energy-below-support-a} of Lemma~\ref{8.11:lem-energy-below-support} gives
$\mathbf{A}(w^{\sharp})-\mathbf{A}(x)\le C_E\varepsilon_1^2\hat N_W$ for every $x$ in that domain.
Together,
\begin{equation}\label{8.11:eq-good-thresholds-3}
\mathbf{A}(w^{\sharp})-C_E\varepsilon_1^2\hat N_W\;\le\;\hat m(\theta_b)\;\le\;\hat m(\theta)
\;\le\;\hat m(\theta_a)\;\le\;\mathbf{A}(w^{\sharp}),
\qquad \theta\in J.
\end{equation}
In particular $\hat m$ is a bounded monotone function on $J$. It is therefore Lebesgue measurable and
integrable on $J$.

Put $f(\theta):=\hat m(\theta)-\hat m(\theta+\vartheta)$ for $\theta\in[\theta_a,\theta_b-\vartheta]$.
The substitution $\theta\mapsto\theta+\vartheta$ gives
\begin{equation*}
\int_{\theta_a}^{\theta_b-\vartheta}\hat m(\theta+\vartheta)\,d\theta
=\int_{\theta_a+\vartheta}^{\theta_b}\hat m .
\end{equation*}
Because $\theta_a+\vartheta\le\theta_b-\vartheta$, the integral of $\hat m$ over
$[\theta_a+\vartheta,\theta_b-\vartheta]$ occurs in both terms and cancels. Hence
\begin{equation}\label{8.11:eq-good-thresholds-4}
\begin{split}
\int_{\theta_a}^{\theta_b-\vartheta}f(\theta)\,d\theta
&=\int_{\theta_a}^{\theta_a+\vartheta}\hat m-\int_{\theta_b-\vartheta}^{\theta_b}\hat m
\le\vartheta\,\bigl[\hat m(\theta_a)-\hat m(\theta_b)\bigr]
\le\vartheta\,C_E\varepsilon_1^2\hat N_W .
\end{split}
\end{equation}
The first inequality holds because $\hat m\le\hat m(\theta_a)$ on $[\theta_a,\theta_a+\vartheta]$ and
$\hat m\ge\hat m(\theta_b)$ on $[\theta_b-\vartheta,\theta_b]$, by monotonicity. The second inequality
is \eqref{8.11:eq-good-thresholds-3}.

By monotonicity, $f\ge 0$. By definition of $(\mathrm{VL})_{\theta}[\vartheta,\Lambda_J]$ in
\eqref{8.11:eq-value-lipschitz-a}, the set $G_J(e,\vartheta)$ is
$\{\theta\in[\theta_a,\theta_b-\vartheta]:f(\theta)\le\vartheta\Lambda_J\}$, which is measurable.
Its complement in $[\theta_a,\theta_b-\vartheta]$ is the set where $f>\vartheta\Lambda_J$. Since
$f\ge0$, Markov's inequality applies to $f$, and together with \eqref{8.11:eq-good-thresholds-4} and
\eqref{8.11:eq-good-thresholds-1} it gives
\begin{equation*}
\bigl|\{f>\vartheta\Lambda_J\}\bigr|
\le\frac{1}{\vartheta\Lambda_J}\int_{\theta_a}^{\theta_b-\vartheta}f
\le\frac{C_E\varepsilon_1^2\hat N_W}{\Lambda_J}=\frac{|J|}{4}.
\end{equation*}
The interval $[\theta_a,\theta_b-\vartheta]$ has length $|J|-\vartheta\ge|J|/2$. Therefore
\begin{equation*}
\bigl|G_J(e,\vartheta)\bigr|\;\ge\;\frac{|J|}{2}-\frac{|J|}{4}\;=\;\frac{|J|}{4}. \qedhere
\end{equation*}
\end{proof}

\begin{definition}[The released problem and its free depth]\label{8.11:def-released-problem}
Let $y_+$ be the upper end of the depth band covered by the first young row, as fixed in
Definition~\ref{8.11:def-young-row-cover}. Let $\delta(\cdot)$ be the signed depth, and let
$\hat x$, the canonical comparison point, be a minimiser of $\mathbf{A}$ over the constrained
class $\widehat{\mathcal{W}}^{\theta}$ at threshold $\theta$. Here $\mathbf{A}(x)$ is the action
regarded as a function of the hole data $x$, which range over the domain $\mathfrak{K}$ of hole data
over which $\mathbf{A}$ is minimised.

\emph{The released class.} For $r\in L^{-1}\mathbb{N}$ put $d_0:=y_+ + r$. Let
$\mathfrak{D}_{\mathrm{near}}$ be the set of hole data bonds, of any level, that have an endpoint
block at depth less than $d_0$, and let $\mathfrak{D}_{\mathrm{far}}$ be the set of all other hole
data bonds. The \emph{released class} is
\begin{equation*}
  \mathfrak{L}(\hat x):=\{\,x\in\mathfrak{K} : x_c=\hat x_c \text{ for every }
  c\in\mathfrak{D}_{\mathrm{near}}\,\}.
\end{equation*}
Thus a point of $\mathfrak{L}(\hat x)$ agrees with $\hat x$ on the near bonds and is
unconstrained on the far bonds. The \emph{released problem} is the minimisation of $\mathbf{A}$
over $\mathfrak{L}(\hat x)$. A minimiser $\check x$ of this problem is fixed in
\eqref{8.11:eq-released-problem-a} below.

\emph{The tail error.} Let $K_w$, $C_{\delta_0}$, $\delta_0>0$, $\sigma_K$, $C_{tg}$ and $D_z$ be
the constants of the window-functional and gauge lemmas, and let
$\varepsilon_{\min}=\min(\varepsilon_{k_o-1},\varepsilon_{k_o})$. For
$r\in L^{-1}\mathbb{N}$ put
\begin{equation}\label{8.11:eq-released-problem-1}
  \vartheta_r:=K_w\,C_{\delta_0}\,\sigma_K\,\varepsilon_{\min}^{-1}\,e^{-\delta_0 r}.
\end{equation}
The relative error $\vartheta_\Delta$, the parameter $\varepsilon>0$ that enters the condition
below together with $\vartheta_\Delta$, and the two one-sided ceilings on the coupling that accompany
$\vartheta_\Delta$ and $\vartheta_t$ respectively, are taken by statement from the
constrained-minimisation estimates of this chapter. We put
\begin{equation*}
  \vartheta_t:=4K_wC_{\delta_0}C_{tg}\,D_z^2\,\vartheta_\Delta ,
  \qquad\text{and require}\qquad \vartheta_t\le\vartheta_1 .
\end{equation*}

\emph{Tail length and free depth.} Let $\check x$ be a minimiser of $\mathbf{A}$ over
$\mathfrak{L}(\hat x)$ and, for $\Lambda>0$, define $r_0(\Lambda)$, $r_1$ and $d_1$ by
\begin{equation}\label{8.11:eq-released-problem-a}
  \begin{gathered}
    \check x\in\operatorname{arg\,min}_{\mathfrak{L}(\hat x)}\mathbf{A},\qquad
    r_0(\Lambda):=\min\bigl\{\,r\in L^{-1}\mathbb{N} :
      \vartheta_r\,\Lambda\le\tfrac12\,\vartheta_\Delta^2\,\varepsilon^2\,\bigr\},\\
    r_1:=r_0(\Lambda)+\tfrac{2}{L},\qquad d_1:=y_+ + r_1 .
  \end{gathered}
\end{equation}
We call $r_0(\Lambda)$ the \emph{tail length} and $d_1$ the \emph{free depth}.

By \eqref{8.11:eq-released-problem-1}, the condition
$\vartheta_r\Lambda\le\frac12\vartheta_\Delta^2\varepsilon^2$ reads
\begin{equation*}
  e^{-\delta_0 r}\ \le\ \frac{\vartheta_\Delta^2\,\varepsilon_{\min}\,\varepsilon^2}
  {2K_wC_{\delta_0}\sigma_K\Lambda}.
\end{equation*}
Since $\delta_0>0$, this is equivalent to
\begin{equation*}
  r\ \ge\ \delta_0^{-1}\log\frac{2K_wC_{\delta_0}\sigma_K\Lambda}
  {\vartheta_\Delta^2\,\varepsilon_{\min}\,\varepsilon^2}.
\end{equation*}
Let $\lceil u\rceil_L$ denote the least element of $L^{-1}\mathbb{N}$ not smaller than $u$. Then
\begin{equation}\label{8.11:eq-released-problem-2}
  r_0(\Lambda)=\Bigl\lceil\,\delta_0^{-1}\log\frac{2K_wC_{\delta_0}\sigma_K\Lambda}
  {\vartheta_\Delta^2\,\varepsilon_{\min}\,\varepsilon^2}\Bigr\rceil_L .
\end{equation}

Now take $\Lambda=\Lambda_J$, the constant of Lemma~\ref{8.11:lem-good-thresholds}, for an
interval $J$ of length $|J|=4\bar q$, where $\bar q$ is a threshold parameter fixed outside this
definition. Then
$\Lambda_J=C_E\varepsilon_1^2\hat N_W/\bar q$, and the argument of the logarithm in
\eqref{8.11:eq-released-problem-2} equals
\begin{equation*}
  \frac{2K_wC_{\delta_0}\sigma_K\,C_E\,\varepsilon_1^2\,\hat N_W}
  {\bar q\,\vartheta_\Delta^2\,\varepsilon_{\min}\,\varepsilon^2}.
\end{equation*}
Hence $r_0(\Lambda_J)$ is the rounding up into $L^{-1}\mathbb{N}$ of $\delta_0^{-1}$ times the
logarithm of this quantity, which is linear in the block count $\hat N_W$ of the zone and carries
the factor $\varepsilon_1^2\varepsilon_{\min}^{-1}\varepsilon^{-2}$, the remaining factors being the
constants displayed above.
\end{definition}

\begin{lemma}[Release and squeeze at the constrained minimiser]\label{8.11:lem-release-squeeze}
We work in the setting of Notation~\ref{8.11:not-hole-exterior-windows}: the domain $\mathsf{K}$ of the
hole data $x$, the held exterior $e$ and its data set $\mathfrak{E}$, the support point $w^{\sharp}$, the
action $\mathbf{A}(x)=A(U(e\oplus x))$ of the hole data, the window functionals $g_a$ of the straddling
sites $a\in\mathfrak{S}$ and of the young cubes $\square'\in\mathfrak{Y}$, the classes
$\widehat{\mathcal{W}}^{\theta}$ with minimum values $\hat m(\theta)$, the stencil $D_c$ of a data bond
$c$, and the set $\mathfrak{S}_{\mathfrak{E}}$ on which relative gauges are normalised. We also use the
released problem of Definition~\ref{8.11:def-released-problem}: the depth $d_0=y_++r$, the near and far
hole data bonds $\mathfrak{D}_{\mathrm{near}}$, $\mathfrak{D}_{\mathrm{far}}$, the released class
$\mathfrak{L}(\hat x)$, its minimiser $\check x$, the numbers $\vartheta_r$, $\vartheta_\Delta$,
$\vartheta_t$, $r_0(\Lambda)$, $r_1$, $d_1$ of \eqref{8.11:eq-released-problem-a}, and the scale
$\varepsilon$ entering the definitions of $\vartheta_\Delta$ and $r_0(\Lambda)$ there. Assume the
following hypotheses.
\begin{enumerate}
\item[(a)] The energy bounds below the support value, Lemma~\ref{8.11:lem-energy-below-support}: for
every $x\in\mathsf{K}$ with $\mathbf{A}(x)\le\mathbf{A}(w^{\sharp})$ one has $x\in\operatorname{int}\mathsf{K}$,
together with the energy and class bounds stated there.
\item[(b)] The free-directions statement: if $x\in\operatorname{int}\mathsf{K}$ and
$\partial\mathbf{A}/\partial x_c(x)=0$ for every data bond $c$ in a set $\mathfrak{F}$, then
$U(e\oplus x)$ satisfies the lattice Yang--Mills equation exactly at every fine bond lying in no
stencil $D_c$ with $c\notin\mathfrak{F}$.
\item[(c)] The window-functional comparison: for a site $a$ of level $i$ and data
$x,x'\in\mathsf{K}$ in a common block gauge with bond discrepancies $\theta(b')\le\alpha_P$ (not to be
confused with the threshold $\theta$ below),
\begin{equation}\label{8.11:eq-release-squeeze-1}
|g_a(x)-g_a(x')|\le\frac{K_w}{\varepsilon_i}\sum_{b'\subset a^{\sim 4}}
e^{-\delta_0\operatorname{dist}(\tilde a,b')}\,\theta(b'),
\end{equation}
the distance measured in blocks of level $i$, hence at least the distance in top blocks; here
$C_{\delta_0}$ denotes the supremum over sites $a$ of
$\sum_{b'\subset a^{\sim 4}}e^{-\delta_0\operatorname{dist}(\tilde a,b')}$.
\item[(d)] The path-curvature estimate: if $x_0,x\in\mathsf{K}$ are joined inside $\mathsf{K}$ by the
geodesic $x_0\exp(t\xi)$ of a common gauge with $|\xi_b|\le\alpha_P$, and
$\varphi(t):=\mathbf{A}(x_0\exp(t\xi))$, then for every data block plaquette $p'$, with
$\Delta b(p')$ the discrepancy between the block data of $x_0$ and of $x$ at $p'$, with
$\vartheta_{CS}$ the small relative error, and with $C_{PE}$ and $\hat\varepsilon_K$ the constant
and the small parameter of this estimate,
\begin{equation}\label{8.11:eq-release-squeeze-2}
\int_0^1(1-t)\varphi''(t)\,dt\ \ge\ \tfrac12(1-\vartheta_{CS})|\Delta b(p')|^2
-C_{PE}\,\hat\varepsilon_K\sum_b|\xi_b|^2 .
\end{equation}
\item[(e)] The relative-gauge statement: for data $x,x'$ there is a relative gauge transformation
$\lambda$, equal to $1$ on $\mathfrak{S}_{\mathfrak{E}}$, with
$|\log x_b(\lambda\cdot x')_b^{-1}|\le C_{tg}D_z^2\,\pi(x,x')$ for every bond $b$, where
$\pi(x,x')$ is the discrepancy parameter of the pair built from the block-plaquette discrepancies
$\Delta b(p')$ of (d), so that a bound on $|\Delta b(p')|$ uniform over the data block plaquettes
$p'$ is a bound on $\pi(x,x')$; one puts
$\sigma_K:=2C_{tg}D_z^2C_{\mathrm{avg}}\varepsilon_W$, where $C_{\mathrm{avg}}\varepsilon_W$ is the bound
on $\|\partial x(P)-1\|$ over the block plaquettes $P$ of data $x\in\mathsf{K}$; and for every fine
plaquette $p$
\begin{equation}\label{8.11:eq-release-squeeze-3}
\|U(e\oplus x)(\partial p)\,U(e\oplus x')(\partial p)^{-1}-1\|
\le K_wC_{\delta_0}\sup_b|\log x_bx_b'^{-1}|\,\eta^2 .
\end{equation}
\end{enumerate}
Assume moreover the one-sided smallness condition on $\varepsilon_W$ under which
Lemma~\ref{8.11:lem-energy-below-support} holds, and the two one-sided smallness conditions under
which $\vartheta_\Delta$ and $\vartheta_t$ are fixed in Definition~\ref{8.11:def-released-problem}.

Let $\theta\in[1,2]$ and $\Lambda>0$, let $\hat x$ be a minimiser of $\mathbf{A}$ over
$\widehat{\mathcal{W}}^{\theta}$, put $r:=r_0(\Lambda)$, and assume $(\mathrm{VL})_{\theta}[\vartheta_r,\Lambda]$
in the sense of Definition~\ref{8.11:def-value-lipschitz}. Then:
\begin{enumerate}
\item[(i)] $\check x$ exists, $\check x\in\operatorname{int}\mathsf{K}$,
\begin{equation*}
\mathbf{A}(\check x)\le\mathbf{A}(\hat x)\le\mathbf{A}(w^{\sharp}),
\end{equation*}
and $U(e\oplus\check x)$ satisfies the lattice Yang--Mills equation exactly at every fine bond with an
endpoint in a block of depth $\ge d_0+1/L$; in particular at every bond of $\Omega_{d_1}$;
\item[(ii)] for every fine plaquette $p$ of $T_\eta$,
\begin{equation}\label{8.11:eq-release-squeeze-4}
\|U(e\oplus\hat x)(\partial p)\,U(e\oplus\check x)(\partial p)^{-1}-1\|\le\vartheta_t\,[\mathrm{u}].
\end{equation}
\end{enumerate}
\end{lemma}

\begin{proof}
The proof has five steps: existence of $\check x$ and the Yang--Mills equation; membership of
$\check x$ in $\widehat{\mathcal{W}}^{\theta+\vartheta_r}$; the value gap from
$(\mathrm{VL})_{\theta}[\vartheta_r,\Lambda]$; the squeeze of the block data; the transfer to the fine
plaquettes.

\emph{Step 1: existence, interiority and the Yang--Mills equation.} The released class
$\mathfrak{L}(\hat x)$ is compact, so a minimiser $\check x$ of $\mathbf{A}$ over it exists. Since
$\hat x\in\mathfrak{L}(\hat x)$, and since $\hat m(\theta)\le\mathbf{A}(w^{\sharp})$ because the support
point $w^{\sharp}$ lies in $\widehat{\mathcal{W}}^{1}\subset\widehat{\mathcal{W}}^{\theta}$ (the classes
increase with the threshold, as noted in Definition~\ref{8.11:def-value-lipschitz}), we have
\begin{equation}\label{8.11:eq-release-squeeze-5}
\mathbf{A}(\check x)\le\mathbf{A}(\hat x)=\hat m(\theta)\le\mathbf{A}(w^{\sharp}).
\end{equation}
By \eqref{8.11:eq-release-squeeze-5} both $\check x$ and $\hat x$ satisfy the hypothesis of
Lemma~\ref{8.11:lem-energy-below-support}, hypothesis (a), which places them in
$\operatorname{int}\mathsf{K}$. The class $\mathfrak{L}(\hat x)$ fixes the coordinates $x_c$ with
$c\in\mathfrak{D}_{\mathrm{near}}$ and leaves the coordinates with $c\in\mathfrak{D}_{\mathrm{far}}$ free;
since $\check x$ minimises $\mathbf{A}$ over this class and is an interior point of $\mathsf{K}$, we get
$\partial\mathbf{A}/\partial x_c(\check x)=0$ for every $c\in\mathfrak{D}_{\mathrm{far}}$. Hypothesis (b)
with $\mathfrak{F}=\mathfrak{D}_{\mathrm{far}}$ then shows that $U(e\oplus\check x)$ satisfies the
lattice Yang--Mills equation exactly at every fine bond lying in no stencil of a data bond of
$\mathfrak{E}\cup\mathfrak{D}_{\mathrm{near}}$. A stencil of a bond of $\mathfrak{D}_{\mathrm{near}}$ is
a block pair one of whose blocks has depth $<d_0$, so both blocks of the pair have depth $\le d_0$;
hence a fine bond with an endpoint in a block of depth $\ge d_0+1/L$ lies in no such stencil, nor in a
stencil of a bond of $\mathfrak{E}$, whose blocks lie in the held exterior, at depth below $d_0$.
Finally, by \eqref{8.11:eq-released-problem-a},
\begin{equation*}
d_1=y_++r_1=y_++r_0(\Lambda)+\tfrac{2}{L}=d_0+\tfrac{2}{L},
\end{equation*}
so every site of $\Omega_{d_1}$ lies in a block of depth $\ge d_0+1/L$, and every bond of
$\Omega_{d_1}$ is among the bonds just described. This proves (i).

\emph{Step 2: the released minimiser lies in a slightly larger class.} Present the pair
$(\hat x,\check x)$ in the relative gauge of hypothesis (e), built on a breadth-first forest rooted on
$\mathfrak{S}_{\mathfrak{E}}$. The two data agree on $\mathfrak{D}_{\mathrm{near}}$, so there
$\lambda=1$ and the bond discrepancy $\theta_\lambda$ vanishes; on $\mathfrak{D}_{\mathrm{far}}$ one has
$\theta_\lambda\le\sigma_K$. Now consider the two kinds of sites of $\mathfrak{S}\cup\mathfrak{Y}$.
\begin{itemize}
\item For $a\in\mathfrak{S}$, a far bond $b'\subset a^{\sim 4}$ satisfies
$\operatorname{dist}(\tilde a,b')\ge d_0-w\ge r$, because $d_0=y_++r$ and $y_+\ge w$.
\item For $\square'\in\mathfrak{Y}$, the test region satisfies $\tilde\square'\subset\{\delta\le y_+\}$
by \eqref{8.11:eq-young-row-cover-a} of Definition~\ref{8.11:def-young-row-cover}, so a far bond $b'$
satisfies $\operatorname{dist}(\tilde\square',b')\ge(d_0-y_+)L\ge r$, the distance counted in the young
cube's own blocks.
\end{itemize}
In the sum \eqref{8.11:eq-release-squeeze-1} of hypothesis (c) only far bonds contribute, each with
discrepancy at most $\sigma_K$ and at distance at least $r$; with $\varepsilon_i\ge\varepsilon_{\min}$ and
$g_a(\hat x)\le\theta$ (as $\hat x\in\widehat{\mathcal{W}}^{\theta}$) this gives, for all
$a\in\mathfrak{S}\cup\mathfrak{Y}$,
\begin{equation}\label{8.11:eq-release-squeeze-6}
g_a(\check x)\le g_a(\hat x)+\frac{K_w}{\varepsilon_{\min}}\,\sigma_KC_{\delta_0}e^{-\delta_0 r}
\le\theta+\vartheta_r ,
\end{equation}
the last inequality being the definition of $\vartheta_r$ in \eqref{8.11:eq-released-problem-a}.
Hence $\check x\in\widehat{\mathcal{W}}^{\theta+\vartheta_r}$, and by \eqref{8.11:eq-release-squeeze-5}
\begin{equation}\label{8.11:eq-release-squeeze-7}
\hat m(\theta+\vartheta_r)\le\mathbf{A}(\check x)\le\mathbf{A}(\hat x)=\hat m(\theta).
\end{equation}

\emph{Step 3: the value gap.} By $(\mathrm{VL})_{\theta}[\vartheta_r,\Lambda]$,
\eqref{8.11:eq-value-lipschitz-a}, we have $\hat m(\theta)-\hat m(\theta+\vartheta_r)\le\vartheta_r\Lambda$.
Together with \eqref{8.11:eq-release-squeeze-7}, and with the definition of $r_0(\Lambda)$ as the least
$r\in L^{-1}\mathbb{N}$ with $\vartheta_r\Lambda\le\tfrac12\vartheta_\Delta^2\varepsilon^2$ (recall
$r=r_0(\Lambda)$), this gives
\begin{equation}\label{8.11:eq-release-squeeze-8}
\Delta_{\mathrm{val}}:=\mathbf{A}(\hat x)-\mathbf{A}(\check x)
\le\hat m(\theta)-\hat m(\theta+\vartheta_r)\le\vartheta_r\Lambda
\le\tfrac12\vartheta_\Delta^2\varepsilon^2 .
\end{equation}

\emph{Step 4: the squeeze of the block data.} Put $\xi:=\log(\check x^{-1}(\lambda\cdot\hat x))$. By
Step~2, $\xi$ is supported on $\mathfrak{D}_{\mathrm{far}}$ and $|\xi_b|\le\sigma_K$. Let
$x_t:=\check x\exp(t\xi)$ and $\varphi(t):=\mathbf{A}(x_t)$. Since $\xi$ vanishes off
$\mathfrak{D}_{\mathrm{far}}$, Step~1 gives
\begin{equation*}
\varphi'(0)=\sum_{c\in\mathfrak{D}_{\mathrm{far}}}\langle\partial_c\mathbf{A}(\check x),\xi_c\rangle=0 .
\end{equation*}
The geodesic $x_t$ stays in $\mathsf{K}$: by Lemma~\ref{8.11:lem-energy-below-support} both endpoints
lie in the set of data whose block configuration $\tilde V$ has regularity
$\varepsilon_1(\tilde V)\le\tfrac12C_{\mathrm{avg}}\varepsilon_W$, so for every block plaquette $P$ and
$t\in[0,1]$
\begin{align*}
\|\partial x_t(P)-1\|
&\le(1-t)\|\partial\check x(P)-1\|+t\|\partial\hat x(P)-1\|+O(\sigma_K^2)\\
&\le\tfrac12C_{\mathrm{avg}}\varepsilon_W+O(\sigma_K^2)\le C_{\mathrm{avg}}\varepsilon_W ,
\end{align*}
the last inequality by the smallness condition on $\varepsilon_W$. By Taylor's formula with integral
remainder and $\varphi'(0)=0$, $\Delta_{\mathrm{val}}=\int_0^1(1-t)\varphi''(t)\,dt$. Applying
hypothesis (d), \eqref{8.11:eq-release-squeeze-2}, to this geodesic, with
$\sum_b|\xi_b|^2\le N_{\mathfrak{h}}\sigma_K^2$, where $N_{\mathfrak{h}}$ is the number of hole data
bonds, we obtain for every data block plaquette $p'$
\begin{equation*}
\Delta_{\mathrm{val}}\ge\tfrac12(1-\vartheta_{CS})|\Delta b(p')|^2
-C_{PE}\,\hat\varepsilon_K N_{\mathfrak{h}}\sigma_K^2 .
\end{equation*}
Combined with \eqref{8.11:eq-release-squeeze-8}, the smallness condition fixing $\vartheta_\Delta$ yields
\begin{equation}\label{8.11:eq-release-squeeze-9}
|\Delta b(p')|\le2\vartheta_\Delta\varepsilon\qquad\text{for every data block plaquette }p'.
\end{equation}

\emph{Step 5: transfer to the fine plaquettes.} By \eqref{8.11:eq-release-squeeze-9} the bound
$2\vartheta_\Delta\varepsilon$ on $|\Delta b(p')|$ holds uniformly over the data block plaquettes $p'$,
so the pair $(\check x,\hat x)$ satisfies $\pi(\check x,\hat x)\le2\vartheta_\Delta\varepsilon$ in
hypothesis (e). The relative-gauge bound of (e) then gives, for every bond $b$,
\begin{equation*}
|\xi_b|\le2C_{tg}D_z^2\vartheta_\Delta\varepsilon\,(1+O(\sigma_K)),
\end{equation*}
and the plaquette comparison \eqref{8.11:eq-release-squeeze-3}, applied with this bound on the bond
logarithms, gives, for every fine plaquette $p$ of $T_\eta$,
\begin{equation*}
\|U(e\oplus\hat x)(\partial p)\,U(e\oplus\check x)(\partial p)^{-1}-1\|
\le2K_wC_{\delta_0}C_{tg}D_z^2\vartheta_\Delta\varepsilon\,(1+O(\sigma_K))\,\eta^2 ,
\end{equation*}
which is the bound \eqref{8.11:eq-release-squeeze-4} of (ii), with
$\vartheta_t=4K_wC_{\delta_0}C_{tg}D_z^2\vartheta_\Delta$ as fixed in
Definition~\ref{8.11:def-released-problem}.
\end{proof}

\begin{remark}[Bounds inherited by the released minimiser]\label{8.11:rem-inherited-bounds}
Let $\theta\in[1,2]$ and $\hat x$ be as in
Lemma~\ref{8.11:lem-release-squeeze}, assume the hypotheses of that lemma, and let $\check x$ be
the released minimiser it provides. Write $\check U:=U(e\oplus\check x)$, $\hat U:=U(e\oplus\hat x)$
and $\check\Phi(p):=\frac1i\log\check U(\partial p)$ for a fine plaquette $p$. Assume further the
following plaquette bounds in the zone: for every $x\in K$, with $U:=U(e\oplus x)$,
\begin{enumerate}
\item[(1)] every fine plaquette $p\subset\tilde a$ with $a$ pure, that is, a level-$k_o$ site with
$a^{\sim4}\cap H=\emptyset$, has $\|U(\partial p)-1\|\le(f_a(e)+\vartheta_1)\,[\mathrm{u}]$, and
the test regions $\tilde a$ of the pure sites cover $\{\delta\le-3w\}$;
\item[(2)] if $x$ satisfies the windows of all straddling sites $a\in\mathfrak{S}$ at threshold
$\theta$, then every fine plaquette $p\subset\{\delta\le w\}$ has
$\|U(\partial p)-1\|\le\theta'\,[\mathrm{u}]$;
\item[(3)] if $x$ satisfies the window of a young cube $\square'\in\mathfrak{Y}$ at threshold
$\theta$, then every fine plaquette $p\subset\tilde\square'$ has
$\|U(\partial p)-1\|\le\theta_y\,[\mathrm{u}]$.
\end{enumerate}
\begin{enumerate}
\item[(a)] The value-Lipschitz property $(\mathrm{VL})_{\theta}[\vartheta_r,\Lambda]$ enters
Lemma~\ref{8.11:lem-release-squeeze} as a hypothesis. All calibrations that serve only to
establish that property enter neither that lemma nor the bounds below. These are the
threshold parameters used in the proof of that property (among them $\bar q$) and the auxiliary
estimates of that proof: a second-derivative estimate, a linear accumulation bound, a geometric
decay condition, a trace condition and an a priori bound.
\item[(b)] Releasing the far data violates no young window: the only windows on young cubes
imposed in the class $\widehat{\mathcal{W}}^{\theta}$ are those of the young family
$\mathfrak{Y}$, and they are treated exactly like the straddling windows.
\item[(c)] The plaquette bounds that hold at $\hat x$ pass to the released minimiser $\check x$ up to
$\vartheta_t$. Put
\begin{equation}\label{8.11:eq-inherited-bounds-a}
\theta'':=\theta'+2\vartheta_t,\qquad \theta_y'':=\theta_y+2\vartheta_t,\qquad
\sigma_e:=s_e\theta+\vartheta_1 .
\end{equation}
Then, writing $1+o$ for the fixed number $1+\varepsilon_a\eta^2$, the following bounds hold:
\begin{equation}\label{8.11:eq-inherited-bounds-1}
\begin{aligned}
\|\check\Phi(p)\|&\le\theta''(1+o)\,[\mathrm{u}] &&\text{for } p\subset\{\delta\le w\},\\
\|\check\Phi(p)\|&\le\theta_y''(1+o)\,[\mathrm{u}] &&\text{for } p\subset\{\delta\le y_+\},\\
\|\check\Phi(p)\|&\le\sigma_e(1+o)\,[\mathrm{u}] &&\text{for } p \text{ a plaquette of }
\Sigma_{-(3w+1)} .
\end{aligned}
\end{equation}
The factor $1+o$ is the comparison constant between $\|\check U(\partial p)-1\|$ and
$\|\check\Phi(p)\|$ on the plaquettes considered: we assume, as a further hypothesis of this
remark, that $\|\check\Phi(p)\|\le(1+o)\,\|\check U(\partial p)-1\|$ there. It is absorbed into
the small relative
errors $\vartheta$ of the estimates that use \eqref{8.11:eq-inherited-bounds-1}.
\item[(d)] We impose the one-sided ceiling on $g$
\begin{equation}\label{8.11:eq-inherited-bounds-b}
r_1\le \frac{w}{12},
\end{equation}
with $r_1$ as in \eqref{8.11:eq-released-problem-a}.
\end{enumerate}
\end{remark}

\begin{proof}
(a) Lemma~\ref{8.11:lem-release-squeeze} assumes $(\mathrm{VL})_{\theta}[\vartheta_r,\Lambda]$
and uses it only to bound the value difference $\mathbf{A}(\hat x)-\mathbf{A}(\check x)$ by
$\vartheta_r\Lambda$. Nothing used to establish that property is invoked in the proof of the lemma
or in the arguments for (b)--(d) below.

(b) In the proof of Lemma~\ref{8.11:lem-release-squeeze} the window bound after release is
obtained for every $a\in\mathfrak{S}\cup\mathfrak{Y}$ by one argument. For a young cube
$\square'\in\mathfrak{Y}$, the young-row cover of Definition~\ref{8.11:def-young-row-cover} gives
$\tilde\square'\subset\{\delta\le y_+\}$, so a far bond $b'$ lies at least
$(d_0-y_+)L\ge r_0(\Lambda)$ blocks of the cube's own scale from $\tilde\square'$. This is the
same lower bound $r_0(\Lambda)$ that separates a far bond from the test region $\tilde a$ of a
straddling site $a\in\mathfrak{S}$, where $\operatorname{dist}(\tilde a,b')\ge d_0-w\ge r_0(\Lambda)$.
The same estimate
therefore gives the young windows of $\check x$ at threshold $\theta+\vartheta_r$, exactly as for
the straddling windows. No other window on a young cube is imposed, so none can be violated.

(c) Since $\hat x\in\operatorname{int}K$ by the interiority statement for the constrained and
released minimisers, and $\hat x\in\widehat{\mathcal{W}}^{\theta}$, the point
$\hat x$ satisfies the straddling
and the young windows at threshold $\theta$. These are the premises of hypotheses (2) and (3)
above, and by them:
\begin{itemize}
\item every fine plaquette $p\subset\{\delta\le w\}$ has
$\|\hat U(\partial p)-1\|\le\theta'\,[\mathrm{u}]$;
\item every fine plaquette $p\subset\tilde\square'$, $\square'\in\mathfrak{Y}$, has
$\|\hat U(\partial p)-1\|\le\theta_y\,[\mathrm{u}]$.
\end{itemize}
By the young-row cover of Definition~\ref{8.11:def-young-row-cover}, every fine plaquette $p$ in
$\{\delta\le y_+\}$ lies in $\{\delta\le w\}$ or in $\tilde\square'$ for some
$\square'\in\mathfrak{Y}$. Since $\theta_y=\theta'\mathsf{T}$ and $\mathsf{T}\ge1$
(Notation~\ref{8.11:not-units-cap-amplitude}), we have $\theta'\le\theta_y$, so every such
plaquette satisfies $\|\hat U(\partial p)-1\|\le\theta_y\,[\mathrm{u}]$.

Let $p$ be such a plaquette. We have
\begin{equation*}
\check U(\partial p)-\hat U(\partial p)=\bigl(1-\hat U(\partial p)\check U(\partial p)^{-1}\bigr)
\check U(\partial p).
\end{equation*}
Because $\check U(\partial p)$ is unitary, the left-hand side has the same operator norm as
$1-\hat U(\partial p)\check U(\partial p)^{-1}$.
Part (ii) of Lemma~\ref{8.11:lem-release-squeeze} bounds the latter by $\vartheta_t\,[\mathrm{u}]$.
Hence, in $\{\delta\le w\}$,
\begin{equation*}
\|\check U(\partial p)-1\|\le \|\check U(\partial p)-\hat U(\partial p)\|+\|\hat U(\partial p)-1\|
\le(\theta'+\vartheta_t)\,[\mathrm{u}]\le\theta''\,[\mathrm{u}],
\end{equation*}
and in the same way $\|\check U(\partial p)-1\|\le(\theta_y+\vartheta_t)\,[\mathrm{u}]\le
\theta_y''\,[\mathrm{u}]$ in $\{\delta\le y_+\}$.

For the exterior bound we use part (i) of Lemma~\ref{8.11:lem-release-squeeze}, which gives
$\check x\in\operatorname{int}K\subset K$. Hypothesis (1) above then
applies to $\check x$ directly, with no use of part (ii).

Let $p$ be a plaquette of $\Sigma_{-(3w+1)}$. Then $p\subset\{\delta\le-3w\}$, so $p\subset\tilde a$
for some pure $a$, and $\tilde a\cap\Sigma_{-(3w+1)}$ contains $p$, hence is non-empty. The
definition of $s_e$
in \eqref{8.11:eq-units-cap-amplitude-a} therefore gives $f_a(e)\le s_e\theta$. Consequently
\begin{equation*}
\|\check U(\partial p)-1\|\le(s_e\theta+\vartheta_1)\,[\mathrm{u}]=\sigma_e\,[\mathrm{u}] .
\end{equation*}
Applying the comparison $\|\check\Phi(p)\|\le(1+o)\,\|\check U(\partial p)-1\|$ assumed in
item (c) to these three bounds gives
\eqref{8.11:eq-inherited-bounds-1}.

(d) The depth $r_1$ of \eqref{8.11:eq-released-problem-a} is of order $\delta_0^{-1}$ times the sum
of $\log\varepsilon^{-1}$ and the logarithm of a polynomial in the size of the hole. The clamp width
is $w=M_2R_{k_o}$ with $R_{k_o}\ge(\log g^{-2})^{r}$ for the exponent $r\ge2$ of \cite{B-CMP119}
(this $r$ is not the tail length $r_0(\Lambda)$ of Lemma~\ref{8.11:lem-release-squeeze}). Thus
\eqref{8.11:eq-inherited-bounds-b} sets a quantity of that order against
\begin{equation*}
\frac{w}{12}=\frac{M_2R_{k_o}}{12}\ge\frac{M_2(\log g^{-2})^{r}}{12},
\end{equation*}
and it is imposed as a one-sided ceiling on $g$. The fraction $1/12$ is
numerical and enters only through the inequality \eqref{8.11:eq-inherited-bounds-b} itself.
\end{proof}

\begin{definition}[Cochains on a box and boundary cells]\label{8.11:not-box-cochains}
Let $d\ge 3$, let $\alpha_i<\beta_i$ be integers for $i=1,\dots,d$, and put
\begin{equation*}
\Omega=\prod_{i=1}^{d}[\alpha_i,\beta_i]\cap\mathbb{Z}^d,\qquad
\ell_i:=\beta_i-\alpha_i\ge 4,\qquad \ell:=\min_i\ell_i .
\end{equation*}
Fix a real number $\varrho$, the aspect bound of the box, and assume $\max_i\ell_i\le\varrho\,\ell$.

\emph{Cochains.} Let $\mathsf{V}$ be a finite-dimensional real normed space (in the applications
$\mathbb{R}$, or the Lie algebra $\mathfrak{g}$ with the operator norm). We write $\mathsf{C}^k(\Omega)$
for the space of
functions on the oriented unit $k$-cells of $\Omega$ with values in $\mathsf{V}$. Every operator on
these spaces used below has a real scalar kernel, acts componentwise on $\mathsf{V}$, and obeys the
triangle inequality in $\|\cdot\|_{\mathsf{V}}$. We use index notation $f_{\mu_1\dots\mu_k}(x)$ for
the value of $f$ on the cell with lowest corner $x$ spanned by $e_{\mu_1},\dots,e_{\mu_k}$, with
$f_{\mu_1\dots\mu_k}$ antisymmetric in its indices, so that reversing the orientation of a cell
changes the sign of the value. The exterior derivative $d$ (no confusion with the dimension $d$ will
arise) and the forward and backward differences
$\nabla_\mu$, $\nabla^-_\mu$ are those of the Glossary (Notation~\ref{0.2:not-glossary}):
\begin{equation*}
(\nabla_\mu f)(x)=f(x+e_\mu)-f(x),\qquad (\nabla^-_\mu f)(x)=f(x)-f(x-e_\mu),
\end{equation*}
and on $1$-cochains $(d\xi)_{\mu\nu}=\nabla_\mu\xi_\nu-\nabla_\nu\xi_\mu$. We put
$\langle f,g\rangle:=\sum_{c}\langle f(c),g(c)\rangle$, the sum running over the unit $k$-cells
$c\subset\Omega$, each counted once with one orientation. We write $d^*$ for the adjoint of $d$ on
$\Omega$ with respect to this pairing. Since $d$ has a real scalar kernel, $d^*$ is the operator
with the transposed kernel, acting componentwise.

\emph{Boundary cells.} A unit cell of $\Omega$ is a \emph{$\partial$-cell} if it lies in one of the
faces $\{x_i=\alpha_i\}$ or $\{x_i=\beta_i\}$ of $\Omega$. We write $\partial\Omega$ for the set of
sites of $\Omega$ lying in some face. We put
\begin{equation*}
\mathsf{C}^k_0(\Omega):=\{f\in\mathsf{C}^k(\Omega):\ f(c)=0\ \text{for every $\partial$-cell } c\}.
\end{equation*}
We write $\mathfrak{b}(\Omega)$ for the set of $\partial$-bonds that lie in exactly one face of
$\Omega$. For $b'\in\mathfrak{b}(\Omega)$ we write $p(b')$ for the plaquette of the second assertion
of \eqref{8.11:eq-box-cochains-a}. We write $E_0$ for the extension by zero of a function on the
$\partial$-bonds to an element of $\mathsf{C}^1(\Omega)$. For a function $\lambda$ on the sites of
$\partial\Omega$
and a $\partial$-bond $b'$ with initial point $b'_-$ and final point $b'_+$, we put
\begin{equation*}
(d_\partial\lambda)(b'):=\lambda(b'_+)-\lambda(b'_-).
\end{equation*}
This is defined because both endpoints of $b'$ lie in a face, and hence in $\partial\Omega$.

\emph{Facts.} First, $d\,\mathsf{C}^k_0(\Omega)\subset\mathsf{C}^{k+1}_0(\Omega)$. Second, every
plaquette of $\mathbb{Z}^d$ that contains a bond of $\Omega$ which is not a $\partial$-bond is a
plaquette of $\Omega$ and is not a $\partial$-plaquette. Third, the following two assertions hold:
\begin{equation}\label{8.11:eq-box-cochains-a}
\begin{gathered}
\langle\Psi,d\xi\rangle=\sum_{b\ \text{not a $\partial$-bond}}\langle(d^*\Psi)(b),\xi(b)\rangle
\qquad\text{for } \Psi\in\mathsf{C}^2(\Omega),\ \xi\in\mathsf{C}^1_0(\Omega),\\
\text{where}\quad (d^*\Psi)(b)=(d^*\Psi)_\nu(x)=-\sum_{\mu}\nabla^-_\mu\Psi_{\mu\nu}(x)\\
\text{for every bond $b=\langle x,x+e_\nu\rangle$ of $\Omega$ that is not a $\partial$-bond};\\
b'\in\mathfrak{b}(\Omega)\ \Longrightarrow\ \text{$b'$ lies in exactly one plaquette $p(b')$ of
$\Omega$}\\
\text{that is not a $\partial$-plaquette.}
\end{gathered}
\end{equation}
In the first assertion the formula for $(d^*\Psi)(b)$ is the full lattice codifferential of
$\mathbb{Z}^d$, read at bonds that are not $\partial$-bonds. The plaquette $p(b')$ of the second
assertion is spanned by $e_k$ and
the inward $\pm e_i$, where $b'=\langle x,x+e_k\rangle$ and the face containing $b'$ is
perpendicular to $e_i$. Finally, a $\partial$-bond lying in two faces lies in no plaquette of
$\Omega$ that is not a $\partial$-plaquette.
\end{definition}

\begin{proof}
\emph{The inclusion.} For $f\in\mathsf{C}^k_0(\Omega)$ and a unit $(k+1)$-cell $c$ of $\Omega$, the
value $(df)(c)$ is a signed sum of the values of $f$ on the $k$-faces of $c$. Let $c$ be a
$\partial$-cell, lying in a face $F$ of $\Omega$. Then every $k$-face of $c$ lies in $F$ and is a
$\partial$-cell. Hence $f$ vanishes on all of them, and $(df)(c)=0$.

\emph{Plaquettes at a bond that is not a $\partial$-bond.} Let $b=\langle x,x+e_\nu\rangle$ be a bond
of $\Omega$ that is not a $\partial$-bond. For $i\neq\nu$ the $i$-th coordinate is constant on $b$.
So $b$ lies in $\{x_i=\alpha_i\}$ or $\{x_i=\beta_i\}$ exactly when $x_i\in\{\alpha_i,\beta_i\}$.
Hence $\alpha_i<x_i<\beta_i$ for every $i\neq\nu$.

The plaquettes of $\mathbb{Z}^d$ containing $b$ are those spanned by $e_\nu$ and $\pm e_\mu$ with
$\mu\neq\nu$. Their vertices are $x$, $x+e_\nu$, $x\pm e_\mu$ and $x+e_\nu\pm e_\mu$. Since
$\alpha_\mu<x_\mu<\beta_\mu$, we have $x_\mu\pm1\in[\alpha_\mu,\beta_\mu]$, so all these vertices lie
in $\Omega$. Such a plaquette lies in no face:
\begin{itemize}
\item for $i\notin\{\mu,\nu\}$ its $i$-coordinate is $x_i$, which is strictly between $\alpha_i$
and $\beta_i$;
\item its $\mu$- and $\nu$-coordinates each take two different values.
\end{itemize}

\emph{The first assertion of \eqref{8.11:eq-box-cochains-a}.} We have
\begin{equation*}
\langle\Psi,d\xi\rangle=\sum_{p\subset\Omega}\langle\Psi(p),(d\xi)(p)\rangle .
\end{equation*}
Here $(d\xi)(p)$ is a
signed sum of the values of $\xi$ on the four bonds of $p$. Since $\xi\in\mathsf{C}^1_0(\Omega)$,
only bonds that are not $\partial$-bonds contribute. Exchanging the two sums gives
\begin{equation*}
\langle\Psi,d\xi\rangle=\sum_{b\ \text{not a $\partial$-bond}}\Bigl\langle\sum_{p\subset\Omega,\ p\ni b}
\pm\Psi(p),\ \xi(b)\Bigr\rangle .
\end{equation*}
The bracket is $(d^*\Psi)(b)$. By the preceding paragraph, the inner sum runs over all $2(d-1)$
plaquettes of $\mathbb{Z}^d$ containing $b$.

We now identify its value. Take $b=\langle x,x+e_\nu\rangle$. For $\mu\neq\nu$ and a site $y$,
\begin{equation*}
(d\xi)_{\mu\nu}(y)=\xi_\nu(y+e_\mu)-\xi_\nu(y)-\xi_\mu(y+e_\nu)+\xi_\mu(y).
\end{equation*}
The term $\xi_\nu(x)$ therefore occurs in two places:
\begin{itemize}
\item with coefficient $+1$ at $y=x-e_\mu$;
\item with coefficient $-1$ at $y=x$.
\end{itemize}
These are the two $(\mu,\nu)$-plaquettes containing $b$. Each unordered plaquette is counted once
in the pairing, and by antisymmetry $\Psi_{\mu\nu}(d\xi)_{\mu\nu}=\Psi_{\nu\mu}(d\xi)_{\nu\mu}$.
Hence the coefficient of $\xi_\nu(x)$ is
\begin{equation*}
\sum_{\mu}\bigl(\Psi_{\mu\nu}(x-e_\mu)-\Psi_{\mu\nu}(x)\bigr)=-\sum_\mu\nabla^-_\mu\Psi_{\mu\nu}(x),
\end{equation*}
using $\Psi_{\nu\nu}=0$.

\emph{The second assertion of \eqref{8.11:eq-box-cochains-a}.} Let $b'=\langle x,x+e_k\rangle\in
\mathfrak{b}(\Omega)$, lying in its one face. This face is perpendicular to some $e_i$ with $i\neq k$,
because the $k$-coordinate takes two values on $b'$. Say $x_i=\alpha_i$; the case $x_i=\beta_i$ is
symmetric, with $-e_i$ inward. Since $b'$ lies in no other face, $\alpha_j<x_j<\beta_j$ for
$j\notin\{i,k\}$. The plaquettes of $\mathbb{Z}^d$ containing $b'$ are spanned by $e_k$ and
$\pm e_j$ with $j\neq k$. We treat them in three cases.
\begin{itemize}
\item For $j\notin\{i,k\}$: both plaquettes are plaquettes of $\Omega$, since $x_j\pm1$ lies in
$[\alpha_j,\beta_j]$. Their $i$-coordinate is identically $\alpha_i$, so they lie in the face
$\{x_i=\alpha_i\}$ and are $\partial$-plaquettes.
\item For $j=i$ with $-e_i$: this plaquette has a vertex with $i$-coordinate $\alpha_i-1$, so it is
not a plaquette of $\Omega$.
\item For $j=i$ with $+e_i$: its vertices have $i$-coordinates $\alpha_i$ and $\alpha_i+1\le\beta_i$
(because $\ell_i\ge4$), so it is a plaquette of $\Omega$. It lies in no face: its $i$- and
$k$-coordinates take two values each, and its $j$-coordinates with $j\notin\{i,k\}$ are strictly
between $\alpha_j$ and $\beta_j$.
\end{itemize}
So $p(b')$ is this last plaquette, and it is the only one.

\emph{A $\partial$-bond in two faces.} Let $b'=\langle x,x+e_k\rangle$ lie in two faces. They are
perpendicular to $e_i$ and $e_j$ with $i\neq j$, because $\alpha_i\neq\beta_i$. Both $i$ and $j$
differ from $k$. A plaquette of $\Omega$ containing $b'$ is spanned by $e_k$ and $\pm e_m$ with
$m\neq k$.
\begin{itemize}
\item If $m\neq i$, its $i$-coordinate is constant and equal to $x_i\in\{\alpha_i,\beta_i\}$, so
it lies in a face perpendicular to $e_i$.
\item If $m=i$, then $m\neq j$, and by the same argument it lies in a face perpendicular to $e_j$.
\end{itemize}
In either case it is a $\partial$-plaquette.
\end{proof}

\begin{lemma}[Relative Hodge projection on a box]\label{8.11:lem-relative-projection}
Let $\Omega$, the cochain spaces $\mathsf{C}^k(\Omega)$, $\mathsf{C}^k_0(\Omega)$, the pairing
$\langle\cdot,\cdot\rangle$, the adjoint $d^*$, the set $\mathfrak{b}(\Omega)$ of boundary bonds
lying in exactly one face, the plaquettes $p(b')$, the extension by zero $E_0$ and the boundary
difference $d_\partial$ be as in Definition~\ref{8.11:not-box-cochains}. Let $\Pi^0$ be the
orthogonal projection of $\mathsf{C}^2(\Omega)$ onto $d\mathsf{C}^1_0(\Omega)$.
\begin{itemize}
\item[(a)] For $\Psi\in\mathsf{C}^2(\Omega)$, $\Pi^0\Psi$ is the unique element $d\zeta\in
d\mathsf{C}^1_0(\Omega)$ with
\begin{equation*}
\langle d\zeta,d\xi\rangle=\langle\Psi,d\xi\rangle\qquad\text{for all }\xi\in\mathsf{C}^1_0(\Omega).
\end{equation*}
It depends only on the restriction of $d^*\Psi$ to the bonds that are not boundary cells, and
only on the restriction of $\Psi$ to the plaquettes that are not boundary cells.
\item[(b)] If $A\in\mathsf{C}^1(\Omega)$ and $a$ is the restriction of $A$ to the boundary bonds,
then $(1-\Pi^0)dA=(1-\Pi^0)dE_0a$.
\item[(c)] For every function $\lambda$ on the boundary sites of $\Omega$,
$(1-\Pi^0)dE_0(d_\partial\lambda)=0$.
\item[(d)] For every function $a$ on the boundary bonds, on the plaquettes that are not boundary
cells
\begin{equation*}
dE_0a=\sum_{b'\in\mathfrak{b}(\Omega)}\pm\,a(b')\,\mathbf{1}_{p(b')},
\end{equation*}
where $a(b')\mathbf{1}_{p(b')}$ denotes the $2$-cochain with value $a(b')$ on $p(b')$ and $0$ on
every other plaquette, and the sign is the incidence sign of $b'$ in the boundary of $p(b')$.
\end{itemize}
\end{lemma}

\begin{proof}
(a) The subspace $d\mathsf{C}^1_0(\Omega)$ is finite-dimensional. By definition of the
orthogonal projection, $\Pi^0\Psi$ is the element $d\zeta$ of this subspace for which
$\Psi-d\zeta$ is orthogonal to every $d\xi$, $\xi\in\mathsf{C}^1_0(\Omega)$. This is the displayed
condition. If two elements of the subspace satisfy it, their difference lies in the subspace and
is orthogonal to it, hence vanishes; so the element is unique.

By Definition~\ref{8.11:not-box-cochains}, every plaquette containing a bond that is not a boundary
cell is a plaquette of $\Omega$ and is not a boundary cell. Consequently, by
\eqref{8.11:eq-box-cochains-a}, for $\xi\in\mathsf{C}^1_0(\Omega)$,
\begin{equation*}
\langle\Psi,d\xi\rangle=\sum_{b\ \text{not a boundary cell}}\langle(d^*\Psi)(b),\xi(b)\rangle,
\end{equation*}
with $(d^*\Psi)(b)$ the full lattice codifferential. Hence if $\Psi,\Psi'$ have the same $d^*$ on
the bonds that are not boundary cells, the right-hand sides of the displayed condition agree for
every $\xi$. The same $d\zeta$ then satisfies both conditions, and uniqueness gives
$\Pi^0\Psi=\Pi^0\Psi'$.

The faces of a boundary cell are boundary cells, so $d\mathsf{C}^1_0(\Omega)\subset
\mathsf{C}^2_0(\Omega)$. Thus every $d\xi$ with $\xi\in\mathsf{C}^1_0(\Omega)$ vanishes on the
boundary plaquettes. Therefore $\langle\Psi,d\xi\rangle$ is a sum over the plaquettes that are
not boundary cells only. If $\Psi$ and $\Psi'$ agree there, the same uniqueness argument gives
$\Pi^0\Psi=\Pi^0\Psi'$.

(b) The cochain $E_0a$ equals $a$ on the boundary bonds, so $A-E_0a$ vanishes there. Thus
$A-E_0a\in\mathsf{C}^1_0(\Omega)$, and so $dA-dE_0a\in d\mathsf{C}^1_0(\Omega)$. The projection
$\Pi^0$ is the identity on its range, so $(1-\Pi^0)(dA-dE_0a)=0$.

(c) Let $E_0\lambda$ be $\lambda$ extended by $0$ to all sites of $\Omega$. For a boundary bond
$b'$, both endpoints $b'_-$ and $b'_+$ are boundary sites, since the faces of a boundary cell are
boundary cells. Hence
\begin{equation*}
(dE_0\lambda)(b')=\lambda(b'_+)-\lambda(b'_-)=(d_\partial\lambda)(b')=(E_0d_\partial\lambda)(b').
\end{equation*}
So $E_0d_\partial\lambda-dE_0\lambda\in\mathsf{C}^1_0(\Omega)$. Applying $d$ and using $dd=0$,
\begin{equation*}
dE_0d_\partial\lambda=d\bigl(E_0d_\partial\lambda-dE_0\lambda\bigr)\in d\mathsf{C}^1_0(\Omega),
\end{equation*}
and $1-\Pi^0$ annihilates it.

(d) Let $p$ be a plaquette that is not a boundary cell. Then $(dE_0a)(p)$ is the sum, over the
bonds $b$ of the boundary of $p$, of $\pm(E_0a)(b)$, each sign being the incidence sign of $b$ in
that boundary. Only boundary bonds contribute, because $E_0a$ vanishes elsewhere.

By \eqref{8.11:eq-box-cochains-a}, a boundary bond lying in two faces lies in no plaquette that
is not a boundary cell. So every contributing bond $b'$ belongs to $\mathfrak{b}(\Omega)$. By the
same display, each $b'\in\mathfrak{b}(\Omega)$ lies in exactly one such plaquette, namely $p(b')$,
and therefore $p(b')=p$. Hence
\begin{equation*}
(dE_0a)(p)=\sum_{b'\in\mathfrak{b}(\Omega):\,p(b')=p}\pm\,a(b'),
\end{equation*}
which is the asserted identity evaluated at $p$.
\end{proof}

\begin{lemma}[Reflection onto the doubled torus]\label{8.11:lem-torus-reflection}
Let $\Omega$, the cochain spaces $\mathsf{C}^k(\Omega)$, $\mathsf{C}^k_0(\Omega)$ and the $\partial$-cells be as in
Definition~\ref{8.11:not-box-cochains}, and let $\Pi^0$ be the orthogonal projection of $\mathsf{C}^2(\Omega)$
onto $d\mathsf{C}^1_0(\Omega)$, as in Lemma~\ref{8.11:lem-relative-projection}. Translate so that
$\alpha_i=0$ for $i=1,\dots,d$; thus $\Omega=\prod_{i=1}^d[0,\ell_i]\cap\mathbb{Z}^d$. Let
$\mathbb{T}:=\prod_{i=1}^d\mathbb{Z}/(2\ell_i)$ be the \emph{doubled torus} of $\Omega$ (not to be confused with
the torus $\mathbb{T}_m$ of Definition~\ref{1.1:def-lattice}), whose cells contain those of $\Omega$. For each $i$
let $R_i$ be the reflection $x_i\mapsto -x_i$ of $\mathbb{T}$, the other coordinates being fixed, and let
$T_i:=-R_i^{*}$ on cochains of $\mathbb{T}$, where $R_i^{*}$ is the cellular pull-back by $R_i$, with the
orientation signs. For $\Psi\in\mathsf{C}^2(\Omega)$ let $\tilde\Psi$ be the $2$-cochain on $\mathbb{T}$ such that
\begin{itemize}
\item[(1)] $\tilde\Psi=\Psi$ on the plaquettes of $\Omega$ that are not $\partial$-cells;
\item[(2)] $T_i\tilde\Psi=\tilde\Psi$ for $i=1,\dots,d$;
\item[(3)] $\tilde\Psi=0$ on every plaquette lying in a hyperplane $\{x_i=0\}$ or $\{x_i=\ell_i\}$.
\end{itemize}
Then $\Pi^0\Psi=(\Pi_{\mathbb{T}}\tilde\Psi)|_{\Omega}$, where $\Pi_{\mathbb{T}}$ is the orthogonal projection of
$\mathsf{C}^2(\mathbb{T})$ onto $d\mathsf{C}^1(\mathbb{T})$ and $|_{\Omega}$ denotes restriction to the cells of
$\Omega$.
\end{lemma}

\begin{proof}
Below, \emph{the hyperplanes} are the $2d$ sets $\{x_i=0\}$, $\{x_i=\ell_i\}$ of $\mathbb{T}$. Since the faces of
$\Omega$ are $\{x_i=0\}$ and $\{x_i=\ell_i\}$, a cell of $\Omega$ is a $\partial$-cell exactly when it lies in a
hyperplane. The inner product $\langle\cdot,\cdot\rangle_{\mathbb{T}}$ on cochains of $\mathbb{T}$ is the sum over
the cells of $\mathbb{T}$, as for $\Omega$; for a set $P$ of plaquettes, $\|\chi\|_P^2$ denotes the sum over
$p\in P$ of the summands of $\langle\chi,\chi\rangle_{\mathbb{T}}$, and $\|\chi\|_{\mathbb{T}}^2:=
\langle\chi,\chi\rangle_{\mathbb{T}}$. Let $P_0$ be the set of plaquettes of $\Omega$ that are not $\partial$-cells.
For $I\subset\{1,\dots,d\}$ put $R_I:=\prod_{i\in I}R_i$ and $T_I:=\prod_{i\in I}T_i$. The $R_i$ commute and
$R_i^2=1$, so the $T_i$ commute, $T_i^2=1$, $T_iT_I=T_{I\triangle\{i\}}$, and the $T_I$ form the group generated
by $T_1,\dots,T_d$.

\emph{$\tilde\Psi$ is well defined.} In each coordinate, $\mathbb{Z}/(2\ell_i)$ is the union of $[0,\ell_i]$ and its
image $[\ell_i,2\ell_i]$ under $x_i\mapsto-x_i$, and every unit interval or point lies in one of the two, which meet
only at $0$ and $\ell_i$. Hence the $2^d$ reflected copies $R_I(\Omega)$ tile $\mathbb{T}$ cell by cell and two
distinct copies overlap only in the hyperplanes. Every plaquette of $\mathbb{T}$ not lying in a hyperplane is
therefore $R_I(q)$ for exactly one $I$ and one $q\in P_0$, and (2) prescribes $\tilde\Psi(R_I(q))$ as
$(-1)^{|I|}\Psi(q)$ times the orientation sign; (3) prescribes $0$ on the remaining plaquettes, which $R_i$
permutes among themselves. No plaquette receives two prescriptions, so $\tilde\Psi$ is well defined.

\emph{An invariant primitive.} Each $R_i$ is an automorphism of the cell complex of $\mathbb{T}$, so $R_i^{*}$, and
hence $T_i$, is an isometry of $\langle\cdot,\cdot\rangle_{\mathbb{T}}$ commuting with $d$; and $T_i$ fixes
$\tilde\Psi$ by (2). Being invertible and mapping $d\mathsf{C}^1(\mathbb{T})$ into itself, $T_i$ maps this subspace
and its orthogonal complement onto themselves, so $T_i\Pi_{\mathbb{T}}=\Pi_{\mathbb{T}}T_i$. Choose
$\tilde\xi_0\in\mathsf{C}^1(\mathbb{T})$ with $d\tilde\xi_0=\Pi_{\mathbb{T}}\tilde\Psi$. Then
\begin{equation*}
d(T_i\tilde\xi_0)=T_i\Pi_{\mathbb{T}}\tilde\Psi=\Pi_{\mathbb{T}}T_i\tilde\Psi=\Pi_{\mathbb{T}}\tilde\Psi ,
\end{equation*}
and by iteration $d(T_I\tilde\xi_0)=\Pi_{\mathbb{T}}\tilde\Psi$ for every $I$. Averaging over the group,
\begin{equation}\label{8.11:eq-torus-reflection-1}
\tilde\xi_1:=2^{-d}\sum_{I\subset\{1,\dots,d\}}T_I\tilde\xi_0
\end{equation}
satisfies $d\tilde\xi_1=\Pi_{\mathbb{T}}\tilde\Psi$, and $T_i\tilde\xi_1=\tilde\xi_1$ because
$I\mapsto I\triangle\{i\}$ permutes the subsets.

\emph{Invariant $1$-cochains vanish on the hyperplanes.} Let $\tilde\eta\in\mathsf{C}^1(\mathbb{T})$ satisfy
$T_i\tilde\eta=\tilde\eta$ for all $i$. A bond $b$ in $\{x_i=0\}$ or $\{x_i=\ell_i\}$ has direction different from
$i$ and endpoints fixed by $R_i$ (as $-\ell_i\equiv\ell_i$ modulo $2\ell_i$), so $R_i$ fixes $b$ with its
orientation and $\tilde\eta(b)=(T_i\tilde\eta)(b)=-\tilde\eta(b)$, i.e.\ $\tilde\eta(b)=0$. All bonds of a plaquette
lying in a hyperplane lie in it, so $d\tilde\eta=0$ on such plaquettes. For $\tilde\eta=\tilde\xi_1$: the
$\partial$-bonds of $\Omega$ lie in hyperplanes, so $\xi_1:=\tilde\xi_1|_{\Omega}\in\mathsf{C}^1_0(\Omega)$. The
coboundary at a plaquette of $\Omega$ reads only bonds of $\Omega$, hence
\begin{equation}\label{8.11:eq-torus-reflection-2}
d\xi_1=(d\tilde\xi_1)|_{\Omega}=(\Pi_{\mathbb{T}}\tilde\Psi)|_{\Omega}.
\end{equation}

\emph{Minimality.} Let $\tilde\xi\in\mathsf{C}^1(\mathbb{T})$ be invariant under all $T_i$, and
$\xi:=\tilde\xi|_{\Omega}$. Then $\chi:=\tilde\Psi-d\tilde\xi$ is invariant, since $T_i$ commutes with $d$, and
$\chi=0$ on plaquettes lying in a hyperplane by (3) and the preceding paragraph. For $q\in P_0$, invariance under
$T_I$ gives $\chi(R_I(q))=\pm\chi(q)$, so the two summands coincide, and
$\chi(q)=\Psi(q)-(d\xi)(q)$ by (1). Summing over the tiling,
\begin{equation}\label{8.11:eq-torus-reflection-3}
\|\tilde\Psi-d\tilde\xi\|_{\mathbb{T}}^2=2^d\,\|\Psi-d\xi\|_{P_0}^2 .
\end{equation}
Every $\xi\in\mathsf{C}^1_0(\Omega)$ has such an invariant extension: it is built as $\tilde\Psi$ was, from the bonds
of $\Omega$ not lying in a hyperplane, with the value $0$ on bonds in the hyperplanes, and it agrees with $\xi$ on
the $\partial$-bonds, where $\xi$ vanishes. As $\Pi_{\mathbb{T}}\tilde\Psi=d\tilde\xi_1$ is the element of
$d\mathsf{C}^1(\mathbb{T})$ closest to $\tilde\Psi$, \eqref{8.11:eq-torus-reflection-3} applied to $\tilde\xi_1$ and
to the invariant extension $\tilde\xi$ of an arbitrary $\xi\in\mathsf{C}^1_0(\Omega)$ yields
\begin{equation*}
2^d\|\Psi-d\xi_1\|_{P_0}^2=\|\tilde\Psi-d\tilde\xi_1\|_{\mathbb{T}}^2
\le\|\tilde\Psi-d\tilde\xi\|_{\mathbb{T}}^2=2^d\|\Psi-d\xi\|_{P_0}^2 .
\end{equation*}
Since $d\mathsf{C}^1_0(\Omega)\subset\mathsf{C}^2_0(\Omega)$ (Definition~\ref{8.11:not-box-cochains}), $d\xi$ vanishes
on the $\partial$-plaquettes, and $\|\Psi-d\xi\|^2$ over all plaquettes of $\Omega$ equals $\|\Psi-d\xi\|_{P_0}^2$
plus the sum of the summands of $\Psi$ over the $\partial$-plaquettes, which does not depend on $\xi$. So $\xi_1$
minimises $\|\Psi-d\xi\|$ over $\mathsf{C}^1_0(\Omega)$, and $d\xi_1$ is the element of $d\mathsf{C}^1_0(\Omega)$
closest to $\Psi$, namely $\Pi^0\Psi$. With \eqref{8.11:eq-torus-reflection-2},
$\Pi^0\Psi=(\Pi_{\mathbb{T}}\tilde\Psi)|_{\Omega}$.
\end{proof}

\begin{lemma}[Derivative bounds for the torus Green's function]\label{8.11:lem-torus-kernel}
Let $d\ge3$ (in this book $d=4$; the argument uses only $d\ge3$). Let
$\mathbb{T}=\prod_{i=1}^{d}\mathbb{Z}/(2\ell_i)$ be the doubled torus of
Lemma~\ref{8.11:lem-torus-reflection}, and let $\ell$ and $\varrho_{\mathbb{T}}$ be the length
scale and the aspect bound of the underlying box, so that
$\ell\le\ell_i\le\varrho_{\mathbb{T}}\ell$ for
$i=1,\dots,d$. Write $(\nabla_\mu\phi)(x)=\phi(x+e_\mu)-\phi(x)$,
$(\nabla^-_\nu\phi)(x)=\phi(x)-\phi(x-e_\nu)$ and
$\Delta=\sum_{\mu}\nabla^-_\mu\nabla_\mu$, so that
$(\Delta\phi)(x)=\sum_{e=\pm e_\mu}(\phi(x+e)-\phi(x))$. Let $G=G_{\mathbb{Z}^d}$ be the
Green's function of $-\Delta$ on $\mathbb{Z}^d$, let $Q_s(x)$ be the $\ell^\infty$-cube of
radius $s$ about $x$, and let $|x|$ be the $\ell^\infty$-norm, on $\mathbb{T}$ the
$\ell^\infty$-distance to $0$. Assume the following three statements.
\begin{itemize}
\item[(a)] The continuum kernel: there is a constant $a_d$, depending only on $d$, such that
$G(x)=a_d|x|_2^{2-d}+O(|x|^{-d})$.
\item[(b)] The gradient estimate for discrete harmonic functions: there is a constant
$C_\nabla$, depending only on $d$, such that if $u$ is discrete-harmonic on the lattice ball
of radius $r\ge2$ about $x$, then $|u(x+e)-u(x)|\le C_\nabla r^{-1}\max|u|$, the maximum
taken over the ball, for every unit vector $e$.
\item[(c)] The mean-value inequality for discrete subharmonic functions: there is a constant
$C_{MV}$, depending only on $d$, such that if $\phi\ge0$ and $\Delta\phi\ge0$ on
$Q_{2s}(x_0)$, then $\phi(x_0)\le C_{MV}s^{-d}\sum_{Q_{2s+1}(x_0)}\phi$.
\end{itemize}
Let $G_{\mathbb{T}}$ be the mean-zero solution of
$-\Delta G_{\mathbb{T}}=\mathbf{1}_0-|\mathbb{T}|^{-1}$ on $\mathbb{T}$. Then for all
$\mu,\nu$ and all $x\in\mathbb{T}$,
\begin{equation}\label{8.11:eq-torus-kernel-1}
|\nabla_\mu G_{\mathbb{T}}(x)|\le C_{G1}(1+|x|)^{1-d},\qquad
|\nabla_\mu\nabla^-_\nu G_{\mathbb{T}}(x)|\le C_{G2}(1+|x|)^{-d},
\end{equation}
with constants $C_{G1},C_{G2}$ depending only on $d$ and $\varrho_{\mathbb{T}}$.
\end{lemma}

\begin{proof}
Throughout, $C$ and $c$ denote positive constants depending only on $d$,
$\varrho_{\mathbb{T}}$ and the
constants $a_d$, $C_\nabla$, $C_{MV}$ of (a)--(c); their values change from line to line.
We write $\bar\phi=|\mathbb{T}|^{-1}\sum_{\mathbb{T}}\phi$ for the mean of a function on
$\mathbb{T}$. The argument has five steps.

If $\ell$ is smaller than a suitable constant depending only on $d$, then every
$\ell_i\le\varrho_{\mathbb{T}}\ell$ is bounded by a constant depending only on $d$ and
$\varrho_{\mathbb{T}}$, so only finitely many tori $\mathbb{T}$ occur; for each of them
$G_{\mathbb{T}}$ is a fixed function on a finite set, and \eqref{8.11:eq-torus-kernel-1} holds
with $C_{G1},C_{G2}$ the largest of the finitely many ratios
$|\nabla_\mu G_{\mathbb{T}}(x)|(1+|x|)^{d-1}$ and
$|\nabla_\mu\nabla^-_\nu G_{\mathbb{T}}(x)|(1+|x|)^{d}$ over these tori. We assume henceforth
that $\ell$ is at least such a constant, large enough that all cubes and balls used below have
radius at least $2$ and at least a fixed multiple of $\ell$.

\emph{Step (1): the bounds on $\mathbb{Z}^d$.} Put $\varphi(x)=a_d|x|_2^{2-d}$. Since
$|x|\le|x|_2\le\sqrt d\,|x|$, the partial derivatives of $\varphi$ of first and second order
satisfy $|\partial\varphi(y)|\le C|y|^{1-d}$ and $|\partial^2\varphi(y)|\le C|y|^{-d}$ for
$y\ne0$. For $|x|\ge3$ the points $x$, $x+e_\mu$, $x-e_\nu$, $x+e_\mu-e_\nu$ and the segments
joining them stay at $\ell^\infty$-distance at least $|x|/3$ from $0$; by Taylor's formula,
$\nabla_\mu\varphi(x)$ is a first partial derivative of $\varphi$ at a point of the segment
$[x,x+e_\mu]$, and $\nabla_\mu\nabla^-_\nu\varphi(x)$ is bounded by the supremum of the second
partial derivatives over the unit square spanned at $x-e_\nu$. Hence
$|\nabla_\mu\varphi(x)|\le C|x|^{1-d}$ and $|\nabla_\mu\nabla^-_\nu\varphi(x)|\le C|x|^{-d}$.
By (a), $G-\varphi=O(|x|^{-d})$, and a first or second difference of this error is bounded by
at most four of its values at points within distance $2$ of $x$, hence is $O(|x|^{-d})$; this
error is already of the size of the second-order bound, and a fortiori of the first. The
finitely many points with $|x|\le2$ contribute finitely many finite values. Therefore
\begin{equation}\label{8.11:eq-torus-kernel-2}
|\nabla_\mu G(x)|\le C(1+|x|)^{1-d},\qquad |\nabla_\mu\nabla^-_\nu G(x)|\le C(1+|x|)^{-d},
\end{equation}
and in the same way $|G(x)|\le C(1+|x|)^{2-d}$ by (a).

\emph{Step (2): the parametrix.} Let $\chi_1$ be a smooth function on $\mathbb{R}$ equal to
$1$ on $[-\tfrac14,\tfrac14]$ and supported in $[-\tfrac12,\tfrac12]$, and put
$\chi(x)=\prod_i\chi_1(x_i/\ell)$. Since $\ell\le\ell_i$, the support of $\chi$ lies in a
fundamental domain of $\mathbb{T}$, and $P:=\chi G$ is a function on $\mathbb{T}$. By the
discrete Leibniz rule
\begin{equation*}
\nabla_\mu(\chi G)(x)=\chi(x+e_\mu)\nabla_\mu G(x)+G(x)\nabla_\mu\chi(x),
\end{equation*}
applied twice,
\begin{equation*}
-\Delta P=\chi\,(-\Delta G)+f=\mathbf{1}_0+f,
\end{equation*}
where we used $\chi(0)=1$, and $f$ is a sum of products $(\nabla^j\chi)(\nabla^{2-j}G)$,
$j=1,2$, of differences of $\chi$ and of $G$ evaluated at points at distance at most $1$ from
$x$. The differences of $\chi$ vanish unless $\ell/4-1\le|x|\le\ell/2+1$, so $f$ is supported
there. On that set $|\nabla^j\chi|\le C\ell^{-j}$ by the mean value theorem for $\chi_1$, and
by \eqref{8.11:eq-torus-kernel-2} and $|G|\le C(1+|x|)^{2-d}$ we have
$|\nabla^{2-j}G|\le C\ell^{2-j-d}$; hence $|f|\le C\ell^{-d}$. A difference of $f$ is a sum of
products in which either $\chi$ receives one more difference ($|\nabla^{j+1}\chi|\le
C\ell^{-j-1}$) or $G$ does; in the latter case, for $j=1$, the factor is $\nabla\nabla^-G$,
bounded by \eqref{8.11:eq-torus-kernel-2}. Hence $|\nabla f|\le C\ell^{-d-1}$. Put
$H:=G_{\mathbb{T}}-P$. Then
\begin{equation*}
-\Delta H=g:=-|\mathbb{T}|^{-1}-f,\qquad \sum_{\mathbb{T}}g=0,
\end{equation*}
the second identity because the sum over $\mathbb{T}$ of the Laplacian of any function
vanishes. Note $|g|\le|\mathbb{T}|^{-1}+C\ell^{-d}\le C\ell^{-d}$ and $|\nabla g|=|\nabla f|\le
C\ell^{-d-1}$, since $|\mathbb{T}|=\prod_i2\ell_i\ge(2\ell)^d$.

\emph{Step (3): the $L^2$ bound.} The eigenvalues of $-\Delta$ on $\mathbb{T}$ are
$\sum_i2(1-\cos(\pi n_i/\ell_i))$, and the smallest non-zero one is
$2(1-\cos(\pi/\max_i\ell_i))\ge c(\varrho_{\mathbb{T}}\ell)^{-2}$, because
$\max_i\ell_i\le\varrho_{\mathbb{T}}\ell$.
Since $H-\bar H$ and $g$ are orthogonal to the constants and $-\Delta(H-\bar H)=g$,
\begin{equation*}
\sum_{\mathbb{T}}|H-\bar H|^2\le C(\varrho_{\mathbb{T}}\ell)^4\sum_{\mathbb{T}}g^2
\le C(\varrho_{\mathbb{T}}\ell)^4\bigl(2|\mathbb{T}|^{-1}+2\ell^d\cdot C\ell^{-2d}\bigr),
\end{equation*}
using that $f$ is supported in a set of at most $C\ell^d$ points. With
$|\mathbb{T}|\ge(2\ell)^d$ this gives
\begin{equation}\label{8.11:eq-torus-kernel-3}
|\mathbb{T}|^{-1}\sum_{\mathbb{T}}|H-\bar H|^2\le C\ell^{2(2-d)}.
\end{equation}

\emph{Step (4): the local correction.} Fix $x_0\in\mathbb{T}$ and put
$Q':=Q_{\lfloor\ell/4\rfloor}(x_0)$, embedded in $\mathbb{Z}^d$ (its side is at most
$2\ell_i$); through this embedding $g$ and $H$ are read on $Q'$, and $H$ on all of
$\mathbb{Z}^d$ periodically. Define on $\mathbb{Z}^d$
\begin{equation*}
v(x):=\bigl(G*(g\mathbf{1}_{Q'})\bigr)(x)=\sum_{y\in Q'}G(x-y)g(y).
\end{equation*}
For $x\in Q'$, since $|g|\le C\ell^{-d}$ and
\begin{equation*}
\sum_{|z|\le\ell/2}(1+|z|)^{2-d}\le C\ell^2,\qquad\sum_{|z|\le\ell/2}(1+|z|)^{1-d}\le C\ell,
\end{equation*}
the bounds of Step (1) give $|v(x)|\le C\ell^{2-d}$ and $|\nabla v(x)|\le C\ell^{1-d}$. For
the second differences,
\begin{equation*}
\nabla_\mu\nabla^-_\nu v(x)=\sum_{y\in Q'}\nabla_\mu\nabla^-_\nu G(x-y)\bigl(g(y)-g(x)\bigr)
+g(x)\sum_{y\in Q'}\nabla_\mu\nabla^-_\nu G(x-y).
\end{equation*}
Let $x\in Q_{\lfloor\ell/8\rfloor}(x_0)$. In the first sum, $|g(y)-g(x)|\le
C|x-y|\,\ell^{-d-1}$, by summing $|\nabla g|\le C\ell^{-d-1}$ along a lattice path from $x$ to
$y$ inside the cube $Q'$; with \eqref{8.11:eq-torus-kernel-2} the first sum is at most
\begin{equation*}
\sum_{y\in Q',\,y\ne x}C|x-y|^{-d}\,|x-y|\,C\ell^{-d-1}
\le C\ell^{-d-1}\sum_{0<|z|\le\ell/2}|z|^{1-d}\le C\ell^{-d}.
\end{equation*}
In the
second sum, $\nabla_\mu$ acting on $x$ in $G(x-y)$ is $h(y-e_\mu)-h(y)$ with
$h(y)=\nabla^-_\nu G(x-y)$; summing over each line of $Q'$ in the direction $e_\mu$ the sum
telescopes to two sums of $\pm\nabla^-_\nu G(x-y)$ over the two faces of $Q'$ orthogonal to
$e_\mu$. These faces have at most $C\ell^{d-1}$ points, at $\ell^\infty$-distance at least
$\lfloor\ell/4\rfloor-\lfloor\ell/8\rfloor\ge c\ell$ from $x$, where
$|\nabla^-_\nu G|\le C\ell^{1-d}$; the face sums are therefore at most $C$, and the second
term is at most $|g(x)|\,C\le C\ell^{-d}$. Hence $|\nabla_\mu\nabla^-_\nu v(x)|\le C\ell^{-d}$
for $x\in Q_{\lfloor\ell/8\rfloor}(x_0)$.

Put $u:=H-\bar H-v$. At every $x\in Q'$ we have $-\Delta H(x)=g(x)$ and $-\Delta
v(x)=g(x)\mathbf{1}_{Q'}(x)=g(x)$, so $u$ is harmonic in $Q'$. For any function $u$,
\begin{equation*}
\Delta(u^2)(x)=2u(x)\Delta u(x)+\sum_{e=\pm e_\mu}\bigl(u(x+e)-u(x)\bigr)^2,
\end{equation*}
so $u^2\ge0$ and $\Delta(u^2)\ge0$ wherever $u$ is harmonic. Apply (c) to $\phi=u^2$ at each
point $z\in Q_{\ell/16}(x_0)$, with a radius $s\ge c\ell$ such that
$Q_{2s+1}(z)\subset Q'$:
\begin{equation*}
\sup_{Q_{\ell/16}(x_0)}|u|\le C\Bigl(\ell^{-d}\sum_{Q'}u^2\Bigr)^{1/2}.
\end{equation*}
Now $\sum_{Q'}u^2\le2\sum_{\mathbb{T}}|H-\bar H|^2+2\sum_{Q'}v^2$; by
\eqref{8.11:eq-torus-kernel-3} and $|\mathbb{T}|\le(2\varrho_{\mathbb{T}}\ell)^d$ the first
term is at most
$C\ell^d\ell^{2(2-d)}$, and by $|v|\le C\ell^{2-d}$ on $Q'$ so is the second. Hence
$\sup_{Q_{\ell/16}(x_0)}|u|\le C\ell^{2-d}$. Applying (b) to $u$ on the ball of radius
comparable to $\ell$ about each point of $Q_{\ell/32}(x_0)$, contained in
$Q_{\ell/16}(x_0)$, gives $|\nabla u|\le C\ell^{1-d}$ on $Q_{\ell/32}(x_0)$, in particular
$|\nabla u(x_0)|\le C\ell^{1-d}$. The differences $\nabla^-_\nu u$ are harmonic where $u$ is
harmonic at both points involved; applying (b) to $\nabla^-_\nu u$ on a ball about $x_0$ of
radius comparable to $\ell$ inside $Q_{\ell/32}(x_0)$ gives $|\nabla_\mu\nabla^-_\nu
u(x_0)|\le C\ell^{-1}\ell^{1-d}=C\ell^{-d}$.

\emph{Step (5): conclusion.} Near $x_0$ we have $H-\bar H=u+v$, and differences of $H$ and of
$H-\bar H$ coincide. By Step (4),
\begin{equation*}
|\nabla_\mu H(x_0)|\le C\ell^{1-d},\qquad|\nabla_\mu\nabla^-_\nu H(x_0)|\le C\ell^{-d},
\end{equation*}
and since $x_0$ was arbitrary these hold on all of $\mathbb{T}$. Recall
$G_{\mathbb{T}}=P+H$. If $|x|\le\ell/4-2$, then $\chi=1$ at all points entering the first and
second differences at $x$, so there $P=G$ and the differences of $G_{\mathbb{T}}$ are those of
$G$ plus those of $H$; by \eqref{8.11:eq-torus-kernel-2} and $\ell\ge c(1+|x|)$ we obtain
\eqref{8.11:eq-torus-kernel-1}. If $|x|>\ell/4-2$ (and always
$|x|\le\max_i\ell_i\le\varrho_{\mathbb{T}}\ell$
on $\mathbb{T}$), the Leibniz rule and the bounds of Steps (1) and (2) give
$|\nabla_\mu P(x)|\le C\ell^{1-d}$ and $|\nabla_\mu\nabla^-_\nu P(x)|\le C\ell^{-d}$, so all of
$\nabla P$, $\nabla H$ are at most $C\ell^{1-d}$ and all of $\nabla\nabla^-P$,
$\nabla\nabla^-H$ are at most $C\ell^{-d}$. Since $1+|x|\le2\varrho_{\mathbb{T}}\ell$,
\begin{equation*}
C\ell^{1-d}\le C\varrho_{\mathbb{T}}^{d-1}(1+|x|)^{1-d},\qquad
C\ell^{-d}\le C\varrho_{\mathbb{T}}^{d}(1+|x|)^{-d},
\end{equation*}
which is \eqref{8.11:eq-torus-kernel-1} in this range as well, with constants depending only
on $d$ and $\varrho_{\mathbb{T}}$.
\end{proof}

\begin{corollary}[Decay of the relative Hodge projection kernels]\label{8.11:cor-projection-kernels}
Let the box $\Omega$ with side lengths $\ell_i$ and aspect bound $\varrho$, the cochain spaces
$\mathsf{C}^k(\Omega)$, the codifferential $d^*$ and the
$\partial$-cells be as in Notation~\ref{8.11:not-box-cochains}. Let $\Pi^0$ be the relative Hodge
projection of Lemma~\ref{8.11:lem-relative-projection}. For a cell $c$ write $x(c)$ for its
barycentre. There are constants $C_\Pi$ and $C_\Pi'$, depending only on $d$ and $\varrho$, with the
following property. For every non-$\partial$ plaquette $p_0$ of $\Omega$ there are real kernels
$\mathcal{K}_\Omega(p_0,\cdot)$ on the non-$\partial$ plaquettes and $\mathcal{J}_\Omega(p_0,\cdot)$ on
the non-$\partial$ bonds of $\Omega$ such that, for every $\Psi\in\mathsf{C}^2(\Omega)$,
\begin{equation}\label{8.11:eq-projection-kernels-a}
\begin{aligned}
&\text{(a)}\quad (\Pi^0\Psi)(p_0)=\sum_{p\ \text{non-}\partial}\mathcal{K}_\Omega(p_0,p)\,\Psi(p),\\
&\phantom{\text{(a)}\quad}|\mathcal{K}_\Omega(p_0,p)|\le C_\Pi\,\bigl(1+|x(p_0)-x(p)|_\infty\bigr)^{-d},\\
&\text{(b)}\quad (\Pi^0\Psi)(p_0)=\sum_{b\ \text{non-}\partial}\mathcal{J}_\Omega(p_0,b)\,(d^*\Psi)(b),\\
&\phantom{\text{(b)}\quad}|\mathcal{J}_\Omega(p_0,b)|\le C_\Pi'\,\bigl(1+|x(p_0)-x(b)|_\infty\bigr)^{1-d}.
\end{aligned}
\end{equation}
\end{corollary}

\begin{proof}
The assertions are invariant under integer translations of $\Omega$, so we may assume $\alpha_i=0$ for
all $i$, as in Lemma~\ref{8.11:lem-torus-reflection}. Let $\mathbb{T}=\prod_i\mathbb{Z}/(2\ell_i)$,
the reflections $R_i$, the operators $T_i$ and the invariant extension $\tilde\Psi$ be as in
Lemma~\ref{8.11:lem-torus-reflection}. That lemma gives
\begin{equation}\label{8.11:eq-projection-kernels-1}
\Pi^0\Psi=(\Pi_{\mathbb{T}}\tilde\Psi)\big|_{\Omega}.
\end{equation}

\textit{A formula for $\Pi_{\mathbb{T}}$.}
For functions $F,f$ on $\mathbb{T}$ put $(F\ast f)(x):=\sum_{y\in\mathbb{T}}F(x-y)f(y)$. We apply this
to cochains componentwise in index notation. Difference operators commute with convolution:
$\nabla_\mu(F\ast f)=(\nabla_\mu F)\ast f=F\ast\nabla_\mu f$, and likewise for $\nabla^-_\mu$.

On $\mathbb{T}$ we have $(d^*\tilde\Psi)_\nu=-\sum_\lambda\nabla^-_\lambda\tilde\Psi_{\lambda\nu}$. For
every $h$ the sum $\sum_{y\in\mathbb{T}}\nabla^-_\lambda h(y)$ vanishes by telescoping. Hence each
component of $d^*\tilde\Psi$ has mean zero.

Let $G_{\mathbb{T}}$ be the function of Lemma~\ref{8.11:lem-torus-kernel}, and put
$\zeta_\nu:=G_{\mathbb{T}}\ast(d^*\tilde\Psi)_\nu$. Since
$-\Delta G_{\mathbb{T}}=\mathbf{1}_0-|\mathbb{T}|^{-1}$ and $d^*\tilde\Psi$ has mean zero,
\begin{equation*}
-\Delta\zeta=d^*\tilde\Psi.
\end{equation*}
Moreover $d^*\zeta=G_{\mathbb{T}}\ast d^*d^*\tilde\Psi=0$. By the lattice Weitzenb\"ock identity, which
states that $d^*d+dd^*=-\Delta$ on $1$-cochains of $\mathbb{T}$ with $\Delta$ acting componentwise,
and since $d^*\zeta=0$, we get $d^*d\zeta=-\Delta\zeta=d^*\tilde\Psi$.

Consequently, for every $\xi\in\mathsf{C}^1(\mathbb{T})$,
\begin{equation*}
\langle d\zeta,d\xi\rangle=\langle d^*d\zeta,\xi\rangle=\langle d^*\tilde\Psi,\xi\rangle
=\langle\tilde\Psi,d\xi\rangle .
\end{equation*}
So $\tilde\Psi-d\zeta$ is orthogonal to $d\mathsf{C}^1(\mathbb{T})$. Since $d\zeta$ lies in
$d\mathsf{C}^1(\mathbb{T})$, we conclude $d\zeta=\Pi_{\mathbb{T}}\tilde\Psi$.

Now write $(d\zeta)_{\mu\nu}=\nabla_\mu\zeta_\nu-\nabla_\nu\zeta_\mu$ and move the differences onto
$G_{\mathbb{T}}$. This gives
\begin{equation}\label{8.11:eq-projection-kernels-2}
\begin{aligned}
(\Pi_{\mathbb{T}}\tilde\Psi)_{\mu\nu}
&=(\nabla_\mu G_{\mathbb{T}})\ast(d^*\tilde\Psi)_\nu-(\nabla_\nu G_{\mathbb{T}})\ast(d^*\tilde\Psi)_\mu\\
&=-\sum_\lambda\Bigl[(\nabla_\mu\nabla^-_\lambda G_{\mathbb{T}})\ast\tilde\Psi_{\lambda\nu}
-(\nabla_\nu\nabla^-_\lambda G_{\mathbb{T}})\ast\tilde\Psi_{\lambda\mu}\Bigr].
\end{aligned}
\end{equation}

\textit{Collecting the images.}
For $S\subset\{1,\dots,d\}$ put $\sigma_S:=\prod_{i\in S}R_i$; these are $2^d$ maps. As in the proof of
Lemma~\ref{8.11:lem-torus-reflection}, the copies $\sigma_S\Omega$ tile $\mathbb{T}$ and overlap only in
the hyperplanes $x_i\in\{0,\ell_i\}$. These hyperplanes are mapped to themselves by every $R_j$.

Let $q$ be a plaquette of $\mathbb{T}$ lying in no such hyperplane. Then $q=\sigma_S p$ for exactly one
$S$ and exactly one plaquette $p$ of $\Omega$ lying in no face of $\Omega$, that is, a non-$\partial$
plaquette. The relations $T_i\tilde\Psi=\tilde\Psi$, with $T_i=-R_i^*$ carrying orientation signs,
give $\tilde\Psi(\sigma_S p)=\epsilon_S(p)\Psi(p)$ with $\epsilon_S(p)\in\{\pm1\}$. On the hyperplane
plaquettes $\tilde\Psi=0$.

Let $z$ be the lowest corner of $p_0$ and $\mu,\nu$ its directions. For a plaquette $q$ of
$\mathbb{T}$, let $\mathcal{K}_{\mathbb{T}}(p_0,q)$ be the coefficient of $\tilde\Psi(q)$ in the second
line of \eqref{8.11:eq-projection-kernels-2} at $z$. By \eqref{8.11:eq-projection-kernels-1},
assertion (a) then holds with
\begin{equation*}
\mathcal{K}_\Omega(p_0,p):=\sum_{S}\epsilon_S(p)\,\mathcal{K}_{\mathbb{T}}(p_0,\sigma_S p).
\end{equation*}

Suppose $q$ has lowest corner $y'$ and directions $\alpha,\beta$. Then $\tilde\Psi(q)$ enters the first
sum in \eqref{8.11:eq-projection-kernels-2} only if $\nu\in\{\alpha,\beta\}$, and the second only if
$\mu\in\{\alpha,\beta\}$. Hence $\mathcal{K}_{\mathbb{T}}(p_0,q)$ is a sum of at most two terms of the
form $\pm(\nabla_\mu\nabla^-_\lambda G_{\mathbb{T}})(z-y')$ or
$\pm(\nabla_\nu\nabla^-_\lambda G_{\mathbb{T}})(z-y')$.
By Lemma~\ref{8.11:lem-torus-kernel} we get
\begin{equation*}
|\mathcal{K}_{\mathbb{T}}(p_0,q)|\le 2C_{G2}\,(1+|z-y'|)^{-d},
\end{equation*}
where $|\cdot|$ is the $\ell^\infty$-distance to $0$ on $\mathbb{T}$.

\textit{Distances to images.}
Let $z,y\in\Omega$ and $S$ be given. We claim that in each coordinate $i$ the torus distance between
$z_i$ and $(\sigma_S y)_i$ is at least $|z_i-y_i|$.

First let $i\notin S$. Since $|z_i-y_i|\le\ell_i$, this distance is
$\min(|z_i-y_i|,\,2\ell_i-|z_i-y_i|)$, and $2\ell_i-|z_i-y_i|\ge|z_i-y_i|$.

Next let $i\in S$. Then $0\le z_i+y_i\le 2\ell_i$, and the distance is
$\min(z_i+y_i,\,2\ell_i-(z_i+y_i))$. We have $|z_i+y_i|\ge|z_i-y_i|$ because $z_i,y_i\ge0$. Also
\begin{equation*}
2\ell_i-(z_i+y_i)=(\ell_i-z_i)+(\ell_i-y_i)\ge|z_i-y_i|,
\end{equation*}
because both summands are non-negative and their difference is $y_i-z_i$. This proves the claim, and
taking the maximum over $i$ gives $|z-\sigma_S y|\ge|z-y|_\infty$.

If $p$ has lowest corner $y$, then $\sigma_S p$ has lowest corner $y'=\sigma_S y-\eta$ with
$\eta\in\{0,1\}^d$, where $\eta_i=1$ exactly when $i\in S$ and $e_i$ is a direction of the cell. The
reason is that $R_i$ then reverses the edge of the cell in direction $e_i$.

If $|v|\le1$, then $1+|u+v|\le 2(1+|u|)$. So a unit shift of the argument costs a factor at most $2^d$
in $(1+|\cdot|)^{-d}$. Applying this to the shift $\eta$, and using the claim, we get
\begin{equation*}
(1+|z-y'|)^{-d}\le 2^d(1+|z-\sigma_S y|)^{-d}\le 2^d(1+|z-y|_\infty)^{-d}.
\end{equation*}

Finally, $x(p_0)-z$ and $x(p)-y$ have coordinates in $\{0,\tfrac12\}$. Hence
$1+|x(p_0)-x(p)|_\infty\le 2(1+|z-y|_\infty)$, which costs a further factor $2^d$. Summing over the
$2^d$ sets $S$ yields (a) with $C_\Pi:=2^{3d+1}C_{G2}$. This constant depends only on $(d,\varrho)$ by
Lemma~\ref{8.11:lem-torus-kernel}. Before the passage to barycentres the bound holds with the lowest
corners $z,y$ in place of $x(p_0),x(p)$; the two forms of the bound differ by a factor at most $2^d$.

\textit{Part (b).}
We use the first line of \eqref{8.11:eq-projection-kernels-2}. As noted in the proof of
Lemma~\ref{8.11:lem-torus-reflection}, each $T_i$ is an isometry commuting with $d$. It therefore
commutes with $d^*$, and $T_id^*\tilde\Psi=d^*\tilde\Psi$.

A bond in a hyperplane $x_i\in\{0,\ell_i\}$ is fixed by $R_i$ together with its orientation. At such a
bond, $d^*\tilde\Psi$ therefore equals its own negative, and so it vanishes.

Now let $b$ be a non-$\partial$ bond of $\Omega$. The value $(d^*\tilde\Psi)(b)$ reads $\tilde\Psi$
only on the plaquettes containing $b$. By the first of the facts \eqref{8.11:eq-box-cochains-a}, these
are non-$\partial$ plaquettes of $\Omega$. None of them lies in a hyperplane $x_i\in\{0,\ell_i\}$, so on
them $\tilde\Psi=\Psi$. Again by the first of the facts \eqref{8.11:eq-box-cochains-a}, the same full
lattice codifferential is $(d^*\Psi)(b)$. Hence
$(d^*\tilde\Psi)(b)=(d^*\Psi)(b)$, and by invariance
$(d^*\tilde\Psi)(\sigma_S b)=\pm(d^*\Psi)(b)$.

As for plaquettes, every bond of $\mathbb{T}$ outside the hyperplanes is $\sigma_S b$ for exactly one
$S$ and one non-$\partial$ bond $b$ of $\Omega$. Thus (b) holds with $\mathcal{J}_\Omega(p_0,b)$ equal to
the signed sum over $S$ of the coefficients of $(d^*\tilde\Psi)(\sigma_S b)$ in
\eqref{8.11:eq-projection-kernels-2} at $z$.

A bond of direction $\alpha$ enters only if $\alpha\in\{\mu,\nu\}$, and then, since $\mu\ne\nu$, in
exactly one term. That term has the form $\pm(\nabla_\mu G_{\mathbb{T}})(z-y')$ or
$\pm(\nabla_\nu G_{\mathbb{T}})(z-y')$, where $y'=\sigma_S y-\eta$ and $\eta\in\{0,e_\alpha\}$. By
Lemma~\ref{8.11:lem-torus-kernel} it is bounded by $C_{G1}(1+|z-y'|)^{1-d}$.

The distance argument above, now with exponent $d-1$, applies unchanged. The unit shift and the passage
to barycentres each cost a factor at most $2^{d-1}\le2^d$, and there are $2^d$ sets $S$. This gives (b)
with $C_\Pi':=2^{3d}C_{G1}$, a constant depending only on $(d,\varrho)$.
\end{proof}

\begin{lemma}[The axial homotopy operator]\label{8.11:lem-axial-homotopy}
Let $\Omega_*\subset\mathbb{Z}^d$ be a box, that is, a product of $d$ integer intervals, let
$z_0\in\Omega_*$ have coordinates $z_{0,1},\dots,z_{0,d}$, and let $\mathsf{C}^k(\Omega_*)$ be the
$\mathsf{V}$-valued $k$-cochains on $\Omega_*$, defined as in
Notation~\ref{8.11:not-box-cochains} with $\Omega_*$ in place of $\Omega$. Cochains are evaluated
on positively oriented cells. For a cell $\sigma$ we write $x_\sigma$ for its lowest corner and
$D(\sigma)$ for its set of directions, so that in the index notation
$\omega(\sigma)=\omega_{\nu_1\dots\nu_k}(x_\sigma)$ with $D(\sigma)=\{\nu_1<\dots<\nu_k\}$.

For integers $a,b$ and a function $f$ on $\mathbb{Z}$ put
\begin{equation*}
\mathbf{s}_a^b[f]:=\sum_{t=a}^{b-1}f(t)\quad\text{if }b\ge a,\qquad
\mathbf{s}_a^b[f]:=-\sum_{t=b}^{a-1}f(t)\quad\text{if }b<a.
\end{equation*}
Fix $j\in\{1,\dots,d\}$, and let
$\sigma$ be a cell of $\Omega_*$ with $j\notin D(\sigma)$. Let $\pi_j\sigma$ be the cell with
directions $D(\sigma)$ and lowest corner $x_\sigma+(z_{0,j}-x_{\sigma,j})e_j$. For an integer $t$, let
$\sigma_j[t]$ be the cell with directions $D(\sigma)\cup\{j\}$ and lowest corner
$x_\sigma+(t-x_{\sigma,j})e_j$. Let $n_j(\sigma)$ be the number of $\mu\in D(\sigma)$ with $\mu<j$.
Because $\sigma$ and $z_0$ lie in $\Omega_*$, so do $\pi_j\sigma$ and every $\sigma_j[t]$ with $t$
between $z_{0,j}$ and $x_{\sigma,j}-1$, or between $x_{\sigma,j}$ and $z_{0,j}-1$.

Define $P_j\colon\mathsf{C}^k(\Omega_*)\to\mathsf{C}^k(\Omega_*)$ for $k\ge0$, and
$K_j\colon\mathsf{C}^k(\Omega_*)\to\mathsf{C}^{k-1}(\Omega_*)$ for $k\ge1$, by
\begin{align*}
(P_j\omega)(\sigma)&:=\omega(\pi_j\sigma),\\
(K_j\omega)(\sigma)&:=(-1)^{n_j(\sigma)}\,
\mathbf{s}_{z_{0,j}}^{x_{\sigma,j}}\bigl[t\mapsto\omega(\sigma_j[t])\bigr]
\end{align*}
if $j\notin D(\sigma)$. Both vanish on cells $\sigma$ with $j\in D(\sigma)$. We put $K_j:=0$ on
$\mathsf{C}^0(\Omega_*)$. Then:
\begin{enumerate}
\item[(a)] on $\mathsf{C}^k(\Omega_*)$, $k\ge0$, one has $dK_j+K_jd=1-P_j$ and $P_jd=dP_j$;
\item[(b)] the operator
\begin{equation*}
K:=K_d+K_{d-1}P_d+\dots+K_1P_2\cdots P_d
\end{equation*}
satisfies
\begin{equation}\label{8.11:eq-axial-homotopy-a}
dK+Kd=1\qquad\text{on }\mathsf{C}^k(\Omega_*),\ k\ge1;
\end{equation}
\item[(c)] let $k\ge1$. For every $(k-1)$-cell $\sigma$ of $\Omega_*$ there are at most
$|x_\sigma-z_0|_1$ $k$-cells, each contained in the smallest box containing $\sigma$ and $z_0$,
and coefficients in $\{0,\pm1\}$ depending only on $\sigma$ and $z_0$, such that for every
$\omega\in\mathsf{C}^k(\Omega_*)$ the value $(K\omega)(\sigma)$ is the combination of the values of
$\omega$ on these cells with these coefficients; in particular
\begin{equation*}
\|(K\omega)(\sigma)\|_{\mathsf{V}}\le\sum_{c}\|\omega(c)\|_{\mathsf{V}},
\end{equation*}
the sum running over these cells $c$.
\end{enumerate}
\end{lemma}

\begin{proof}
We use the coboundary on positively oriented cells in the index notation:
\begin{equation*}
(d\omega)_{\nu_1\dots\nu_{k+1}}(y)
=\sum_{i=1}^{k+1}(-1)^{i-1}\nabla_{\nu_i}\omega_{\nu_1\dots\widehat{\nu_i}\dots\nu_{k+1}}(y).
\end{equation*}
Here $\nabla_\mu f(y)=f(y+e_\mu)-f(y)$.

\emph{One-dimensional telescoping.} Splitting into the cases $b\ge a$ and $b<a$ gives two
identities:
\begin{equation}\label{8.11:eq-axial-homotopy-1}
\mathbf{s}_a^{b+1}[f]-\mathbf{s}_a^{b}[f]=f(b),\qquad
\mathbf{s}_a^{b}\bigl[t\mapsto g(t+1)-g(t)\bigr]=g(b)-g(a).
\end{equation}
If $b\ge a$, the first identity adds the term $t=b$. If $b+1\le a$, it reads
$-\sum_{b+1}^{a-1}f+\sum_{b}^{a-1}f=f(b)$. The second identity is a telescoping sum in both cases.

\emph{Proof of (a).} Let $\sigma$ be a $k$-cell with $D(\sigma)=\{\nu_1<\dots<\nu_k\}$, lowest corner
$y$, and $n:=\#\{i:\nu_i<j\}$. Write $y^t:=y+(t-y_j)e_j$. Write $\nu$ for the string
$\nu_1\dots\nu_k$, $\nu\setminus\nu_i$ for the string with $\nu_i$ deleted, and $\nu\cup j$ for
the string with $j$ inserted in increasing position.

First suppose $j\notin D(\sigma)$. In $D(\sigma)\cup\{j\}$ the index $j$ stands in position $n+1$.
An index $\nu_i$ stands in position $i$ if $\nu_i<j$, and in position $i+1$ if $\nu_i>j$. Hence
$(K_jd\omega)(\sigma)$ equals
\begin{equation*}
(-1)^{n}\mathbf{s}_{z_{0,j}}^{y_j}\Bigl[t\mapsto(-1)^{n}\nabla_j\omega_{\nu}(y^t)
+\sum_{i}\epsilon_i\nabla_{\nu_i}\omega_{(\nu\cup j)\setminus\nu_i}(y^t)\Bigr].
\end{equation*}
Here $\epsilon_i=(-1)^{i-1}$ if $\nu_i<j$ and $\epsilon_i=(-1)^{i}$ if $\nu_i>j$. By the second
identity of \eqref{8.11:eq-axial-homotopy-1}, the first term contributes
$\omega_\nu(y)-\omega_\nu(y^{z_{0,j}})=\omega(\sigma)-(P_j\omega)(\sigma)$.

Next consider
\begin{equation*}
(dK_j\omega)(\sigma)=\sum_i(-1)^{i-1}\nabla_{\nu_i}(K_j\omega)_{\nu\setminus\nu_i}(y).
\end{equation*}
The cell with directions $\nu\setminus\nu_i$ has $n_j$ equal to $n-1$ if $\nu_i<j$, and to $n$ if
$\nu_i>j$. Since $\nu_i\ne j$, the shift by $e_{\nu_i}$ leaves the $j$-coordinate unchanged, so
$\nabla_{\nu_i}$ passes inside $\mathbf{s}$. The term of index $i$ therefore carries the sign
$(-1)^{i-1+n-1}$ or $(-1)^{i-1+n}$. In $K_jd\omega$ the same term carries $(-1)^{n+i-1}$ or
$(-1)^{n+i}$ respectively. The two signs are opposite, and these terms cancel. For $k=0$ there are
no such terms and $dK_j=0$. In both cases $(dK_j+K_jd)\omega(\sigma)=\omega(\sigma)-(P_j\omega)(\sigma)$.

Now suppose $j=\nu_{n+1}$. Then $(P_j\omega)(\sigma)=0$, and $(K_jd\omega)(\sigma)=0$ by definition.
In $(dK_j\omega)(\sigma)$ the terms with $\nu_i\ne j$ vanish, because $K_j\omega$ vanishes on cells
extending in direction $j$. The remaining term is
\begin{equation*}
(-1)^{n}\bigl[(K_j\omega)_{\nu\setminus j}(y+e_j)-(K_j\omega)_{\nu\setminus j}(y)\bigr]
=\bigl(\mathbf{s}_{z_{0,j}}^{y_j+1}-\mathbf{s}_{z_{0,j}}^{y_j}\bigr)
\bigl[t\mapsto\omega_\nu(y^t)\bigr].
\end{equation*}
Here the cell with directions $\nu\setminus j$ has $n_j=n$, so the factor $(-1)^n$ cancels the
sign $(-1)^{n_j}$ in the definition of $K_j$, and the cells summed over have lowest corners $y^t$
in both terms, because $(y+e_j)^t=y^t$.
By the first identity of \eqref{8.11:eq-axial-homotopy-1} this equals $\omega_\nu(y)=\omega(\sigma)$.
This proves $dK_j+K_jd=1-P_j$.

We turn to $P_jd=dP_j$. If $j\in D(\sigma)$, then $(P_jd\omega)(\sigma)=0$. In $(dP_j\omega)(\sigma)$ only
the term $\nabla_j(P_j\omega)_{\nu\setminus j}(y)$ survives. It vanishes, since
$(P_j\omega)_{\nu\setminus j}(y')=\omega_{\nu\setminus j}(y'+(z_{0,j}-y'_j)e_j)$ does not depend on
$y'_j$. If $j\notin D(\sigma)$, then $(P_jd\omega)(\sigma)$ and $(dP_j\omega)(\sigma)$ are both equal to
$\sum_i(-1)^{i-1}\nabla_{\nu_i}\omega_{\nu\setminus\nu_i}(y^{z_{0,j}})$. The reason is that a shift
by $e_{\nu_i}$, $\nu_i\ne j$, commutes with replacing the $j$-coordinate by $z_{0,j}$.

\emph{Proof of (b).} By (a), $P_l$ commutes with $d$ for every $l$. Applying (a) on
$\mathsf{C}^k(\Omega_*)$, $k\ge1$ (the operators $P_l$ preserve the degree), we get
\begin{equation*}
dK+Kd=\sum_{j=1}^d(dK_j+K_jd)P_{j+1}\cdots P_d
=\sum_{j=1}^d\bigl(P_{j+1}\cdots P_d-P_j\cdots P_d\bigr)=1-P_1\cdots P_d.
\end{equation*}
Here the empty product for $j=d$ is $1$. Let $\sigma$ have positive dimension, and let $l$ be the
least element of $D(\sigma)$. The projections $\pi_1,\dots,\pi_{l-1}$ preserve the set of
directions. Hence $(P_1\cdots P_d\psi)(\sigma)$ is either $0$ or equal to
$(P_l\cdots P_d\psi)(\pi_{l-1}\cdots\pi_1\sigma)$, and the latter vanishes because the cell extends
in direction $l$. Thus $P_1\cdots P_d=0$ on $\mathsf{C}^k(\Omega_*)$, $k\ge1$, and
\eqref{8.11:eq-axial-homotopy-a} follows.

\emph{Proof of (c).} By the same unwinding, $(P_{j+1}\cdots P_d\omega)(c)$ is either $0$ or equal
to $\omega(c')$. Here $c'$ has the directions of $c$, and its lowest corner is that of $c$ with
coordinates $j+1,\dots,d$ replaced by $z_{0,j+1},\dots,z_{0,d}$. Hence
$(K_jP_{j+1}\cdots P_d\omega)(\sigma)$ is a sum, with coefficients $\pm1$ or $0$ determined by
$\sigma$ and $z_0$, of values of
$\omega$ at no more than $|x_{\sigma,j}-z_{0,j}|$ cells. These cells agree with $\sigma$ in the
coordinates $i<j$. In coordinate $j$ they occupy $[t,t+1]$ with $t$ between $z_{0,j}$ and
$x_{\sigma,j}$. In the coordinates $i>j$, where $\sigma$ has no extent unless the term vanishes,
they sit at $z_{0,i}$.

Each such cell lies in the smallest box containing $\sigma$ and $z_0$. Summing over $j$ gives at
most $|x_\sigma-z_0|_1$ cells. The triangle inequality in $\|\cdot\|_{\mathsf{V}}$ then gives the
bound.
\end{proof}

\begin{definition}[Path--surface systems]\label{8.11:def-path-surface}
Let $\Omega_1\subset\Omega_*$ be two boxes in the sense of Notation~\ref{8.11:not-box-cochains}.
As there, $\partial\Omega_1$ denotes the set of sites of $\Omega_1$ that lie in a face of
$\Omega_1$, that is, the sites of $\Omega_1$ at which some coordinate takes its least or its
greatest value on $\Omega_1$; and $\mathfrak{b}(\Omega_1)$ denotes the set
of bonds of $\Omega_1$ that lie in exactly one face of $\Omega_1$. A \emph{lattice path} from
a site $y$ to a site $z$ is a finite sequence of sites in which consecutive sites are nearest
neighbours, beginning at $y$ and ending at $z$; it is identified with the integral $1$-chain
formed by the sum of its bonds, each oriented in the direction of traversal. An
\emph{integral cellular $2$-chain} in $\Omega_*$ is a finite formal integer combination of
oriented plaquettes of $\Omega_*$, a plaquette with the opposite orientation of $p$ being
identified with $-p$. Fix once and for all one orientation of each plaquette of $\Omega_*$;
an integral cellular $2$-chain is then uniquely a finite sum $S=\sum_p n_p\,p$ with
$n_p\in\mathbb{Z}$, over the plaquettes $p$ of $\Omega_*$ so oriented, and the modulus $|n_p|$
of the multiplicity of $p$ in $S$ does not depend on the orientations chosen.

A \emph{path--surface system} for the pair $\Omega_1\subset\Omega_*$ consists of the data
\textup{(a)}--\textup{(c)} below; we write it
\begin{equation*}
\mathfrak{Sy}=\bigl(\{\gamma_x\},\{S(b')\}\bigr),
\end{equation*}
the base site $o$ being understood:
\begin{itemize}
\item[(a)] a base site $o\in\Omega_*$;
\item[(b)] for every site $x\in\partial\Omega_1$, a lattice path $\gamma_x$ in $\Omega_*$ from
$o$ to $x$;
\item[(c)] for every bond $b'=\langle x,x'\rangle\in\mathfrak{b}(\Omega_1)$, oriented from $x$
to $x'$, an integral cellular $2$-chain $S(b')$ in $\Omega_*$ with
\begin{equation}\label{8.11:eq-path-surface-1}
\partial S(b')=\gamma_x+b'-\gamma_{x'} .
\end{equation}
\end{itemize}
The right-hand side of \eqref{8.11:eq-path-surface-1} is a $1$-cycle, since its boundary is
$(x-o)+(x'-x)-(x'-o)=0$.

In sums over the plaquettes of $S(b')$, and in its area, each plaquette is counted with the
modulus of its multiplicity: if $S(b')=\sum_p n_p(b')\,p$ in the form above and $F$ is a
function on the (unoriented) plaquettes of $\Omega_*$, then
\begin{equation}\label{8.11:eq-path-surface-2}
\sum_{p\in S(b')}F(p):=\sum_p |n_p(b')|\,F(p),
\qquad
|S(b')|:=\sum_p |n_p(b')| .
\end{equation}
\end{definition}

\begin{proposition}[Trace representation with nonabelian defects]\label{8.11:prop-trace-representation}
Let $\Omega_1\subset\Omega_*$ be boxes with the cochain conventions of~\eqref{8.11:eq-box-cochains-a}.
Let $\Phi\in\mathsf{C}^2(\Omega_*;\mathsf{V})$ and $z_0\in\Omega_1$. Let $K$ be the axial homotopy
operator of Lemma~\ref{8.11:lem-axial-homotopy} centred at $z_0$, and put
\begin{equation*}
\rho_2:=K(d\Phi),\qquad N_1:=(d^*\Phi)\big|_{\text{non-}\partial\text{ bonds of }\Omega_1}.
\end{equation*}
Let $\mathfrak{Sy}=(\{\gamma_x\},\{S(b')\})$ be a path--surface system for
$\Omega_1\subset\Omega_*$ in the sense of Definition~\ref{8.11:def-path-surface}. Let
$p_0=p_{\mu\nu}(z_0)$ be a plaquette of $\Omega_1$ none of whose bonds is a $\partial$-bond
of $\Omega_1$. For a cell $c$ with lowest corner $x_c$ put $\rho(c):=1+|z_0-x_c|_\infty$. Then
\begin{equation}\label{8.11:eq-trace-representation-a}
\begin{split}
\|\Phi(p_0)\|\le{}& C_\Pi\sum_{b'\in\mathfrak{b}(\Omega_1)}\rho(p(b'))^{-d}
\sum_{p\in S(b')}\bigl[\|\Phi(p)\|+\|\rho_2(p)\|\bigr]
+C_\Pi\sum_{p\subset\Omega_1}\rho(p)^{-d}\|\rho_2(p)\|\\
&+\|\rho_2(p_0)\|+C_\Pi'\sum_{b\subset\Omega_1}\rho(b)^{1-d}\|N_1(b)\|,
\end{split}
\end{equation}
where $C_\Pi$, $C_\Pi'$ are the constants of Corollary~\ref{8.11:cor-projection-kernels}
for the box $\Omega_1$.
\end{proposition}

\begin{proof}
Since $\Phi$ is a $2$-cochain, the identity $dK+Kd=1$ of Lemma~\ref{8.11:lem-axial-homotopy},
recorded in~\eqref{8.11:eq-axial-homotopy-a}, applies on $\mathsf{C}^2(\Omega_*)$. With
$\tilde A:=K\Phi\in\mathsf{C}^1(\Omega_*;\mathsf{V})$ it gives, on $\Omega_*$,
\begin{equation}\label{8.11:eq-trace-representation-1}
\Phi=dK\Phi+Kd\Phi=d\tilde A+\rho_2 .
\end{equation}
Every bond of a plaquette of $\Omega_1$ is a bond of $\Omega_1$, so the restriction of
$d\tilde A$ to $\Omega_1$ is $d$ applied to the restriction of $\tilde A$. Let $\Pi^0$ be the
relative Hodge projection of $\mathsf{C}^2(\Omega_1)$ of Lemma~\ref{8.11:lem-relative-projection}.
Writing $\Phi=\Pi^0\Phi+(1-\Pi^0)\Phi$ and inserting~\eqref{8.11:eq-trace-representation-1}
into the second term, we obtain on $\Omega_1$
\begin{equation}\label{8.11:eq-trace-representation-2}
\Phi=\Pi^0\Phi+(1-\Pi^0)d\tilde A+(1-\Pi^0)\rho_2 .
\end{equation}

For a $1$-cochain $B$ and a lattice path $\gamma$ write $B(\gamma):=\sum_{b\in\gamma}\pm B(b)$,
the sign being $+$ when $\gamma$ traverses $b$ along its orientation; $B$ is evaluated on
integral $1$-chains by linearity. Let $\tilde a$ be the restriction of $\tilde A$ to the
$\partial$-bonds of $\Omega_1$, let $\lambda(x):=\tilde A(\gamma_x)$ for every site $x$ of
$\partial\Omega_1$, and put $a':=\tilde a-d_\partial\lambda$. Part~(b) of
Lemma~\ref{8.11:lem-relative-projection}, applied to the restriction of $\tilde A$ to
$\Omega_1$, whose $\partial$-values are $\tilde a$, gives
$(1-\Pi^0)d\tilde A=(1-\Pi^0)dE_0\tilde a$. Part~(c) gives
$(1-\Pi^0)dE_0(d_\partial\lambda)=0$. Since $E_0$ and $d$ are linear,
\begin{equation}\label{8.11:eq-trace-representation-3}
(1-\Pi^0)d\tilde A=(1-\Pi^0)dE_0a' .
\end{equation}

Let $b'=\langle x,x'\rangle\in\mathfrak{b}(\Omega_1)$, oriented from $x$ to $x'$. Then
$(d_\partial\lambda)(b')=\lambda(x')-\lambda(x)$, and hence
\begin{equation*}
a'(b')=\tilde A(\gamma_x)+\tilde A(b')-\tilde A(\gamma_{x'})=\tilde A(\gamma_x+b'-\gamma_{x'})
=\tilde A(\partial S(b')) .
\end{equation*}
Write $S(b')=\sum_p n_p\,p$ with integer multiplicities $n_p$. By the definition of the
coboundary, $\tilde A(\partial p)=(d\tilde A)(p)$ for every plaquette $p$; this is Stokes'
formula for cochains. Since $S(b')$ lies in $\Omega_*$, where~\eqref{8.11:eq-trace-representation-1}
holds, we get
\begin{equation}\label{8.11:eq-trace-representation-4}
a'(b')=\sum_p n_p\,(d\tilde A)(p)=\sum_p n_p\,(\Phi-\rho_2)(p).
\end{equation}
Counting plaquettes with $|n_p|$, as in Definition~\ref{8.11:def-path-surface}, this yields
$\|a'(b')\|\le\sum_{p\in S(b')}[\|\Phi(p)\|+\|\rho_2(p)\|]$.

We now evaluate~\eqref{8.11:eq-trace-representation-2} at $p_0$. A $\partial$-plaquette lies in
a face of $\Omega_1$, and so do its bonds; since no bond of $p_0$ is a $\partial$-bond, $p_0$ is
a non-$\partial$ plaquette. Its base point is $z_0$, so $1+|x_{p_0}-x_c|_\infty=\rho(c)$ for
every cell $c$. Four evaluations are needed.

First, by part~(d) of Lemma~\ref{8.11:lem-relative-projection},
$(dE_0a')(p_0)=\sum_{b'\in\mathfrak{b}}\pm a'(b')\mathbf{1}_{p(b')}(p_0)$. Each $b'$ is a bond
of $p(b')$ and a $\partial$-bond, whereas $p_0$ has no $\partial$-bond. Hence
$p_0\neq p(b')$ for all $b'$, and $(dE_0a')(p_0)=0$.

Second, Corollary~\ref{8.11:cor-projection-kernels}(a), recorded
in~\eqref{8.11:eq-projection-kernels-a}, expresses $(\Pi^0dE_0a')(p_0)$ as a sum over
non-$\partial$ plaquettes $p$ of $\mathcal{K}_{\Omega_1}(p_0,p)(dE_0a')(p)$. Inserting
part~(d) of Lemma~\ref{8.11:lem-relative-projection} on these plaquettes gives
\begin{equation*}
(\Pi^0dE_0a')(p_0)=\sum_{b'\in\mathfrak{b}(\Omega_1)}\pm\mathcal{K}_{\Omega_1}(p_0,p(b'))\,a'(b').
\end{equation*}

Third, Corollary~\ref{8.11:cor-projection-kernels}(b) gives
$(\Pi^0\Phi)(p_0)=\sum_{b}\mathcal{J}_{\Omega_1}(p_0,b)N_1(b)$, the sum running over
non-$\partial$ bonds $b$.

Fourth, Corollary~\ref{8.11:cor-projection-kernels}(a) gives
$(\Pi^0\rho_2)(p_0)=\sum_{p}\mathcal{K}_{\Omega_1}(p_0,p)\rho_2(p)$, the sum running over
non-$\partial$ plaquettes $p$.

Collecting these four evaluations, \eqref{8.11:eq-trace-representation-2}
and~\eqref{8.11:eq-trace-representation-3} give
\begin{equation}\label{8.11:eq-trace-representation-5}
\begin{split}
\Phi(p_0)={}&\sum_b\mathcal{J}_{\Omega_1}(p_0,b)N_1(b)
-\sum_{b'\in\mathfrak{b}(\Omega_1)}\pm\mathcal{K}_{\Omega_1}(p_0,p(b'))\,a'(b')\\
&+\rho_2(p_0)-\sum_p\mathcal{K}_{\Omega_1}(p_0,p)\rho_2(p).
\end{split}
\end{equation}
The kernels are real scalars acting componentwise, so the triangle inequality in
$\|\cdot\|_{\mathsf{V}}$ applies term by term. We use
$|\mathcal{K}_{\Omega_1}(p_0,p)|\le C_\Pi\rho(p)^{-d}$ and
$|\mathcal{J}_{\Omega_1}(p_0,b)|\le C_\Pi'\rho(b)^{1-d}$, together with the bound on
$\|a'(b')\|$ following~\eqref{8.11:eq-trace-representation-4}. Finally we enlarge the sums over
non-$\partial$ cells to all cells of $\Omega_1$. This gives~\eqref{8.11:eq-trace-representation-a}.
\end{proof}

\begin{remark}
Suppose $d\Phi=0$ on $\Omega_*$ and $N_1=0$; this is the lattice Maxwell case in $\Omega_1$.
Then $\rho_2=0$, and \eqref{8.11:eq-trace-representation-4}
with~\eqref{8.11:eq-trace-representation-5} give the exact identity
\begin{equation*}
\Phi(p_0)=-\sum_{b'\in\mathfrak{b}(\Omega_1)}\pm\mathcal{K}_{\Omega_1}(p_0,p(b'))\,\Phi(S(b')),
\end{equation*}
where $\Phi(S(b'))=\sum_p n_p\Phi(p)$ is the flux of $\Phi$ through $S(b')$. Thus the field at
$p_0$ is the kernel $\mathcal{K}_{\Omega_1}$ applied to the trace potential $a'$, in the gauge
defined by the paths $\gamma_x$.
\end{remark}

\begin{definition}[An explicit path--surface system between boxes]\label{8.11:def-explicit-system}
Boxes, cells and $\partial$-cells are as in Notation~\ref{8.11:not-box-cochains}, and so is, for a box
$\Omega$, the set $\mathfrak{b}(\Omega)$ of $\partial$-bonds of $\Omega$ lying in exactly one face of
$\Omega$; in particular $d\ge 3$.

\emph{Conventions.} If $\Omega=\prod_{i=1}^d[\alpha_i',\beta_i']\cap\mathbb{Z}^d$ is a box, a
coordinate $x_i$ of a site $x$ is called \emph{$\Omega$-extreme} if $x_i\in\{\alpha_i',\beta_i'\}$.
The word always refers to the box named.

A lattice path is identified with the integral $1$-chain formed by its bonds, each oriented in the
direction of traversal. A sum of paths, each starting where the preceding one ends, is their
concatenation. For sites $u,v$ that differ in at most one coordinate, $[u\to v]$ is the straight path
from $u$ to $v$ (the empty path if $u=v$), and $-\beta$ is the bond $\beta$ with reversed orientation.

A plaquette spanned by $e_j$ and $e_k$ is called a plaquette $(j,k)$.

\emph{The setting.} Let $\Omega_1=\prod_{i=1}^d[\alpha_i,\beta_i]\cap\mathbb{Z}^d$ be a box in the
sense of Notation~\ref{8.11:not-box-cochains}. For
$\tau\in\mathbb{N}$, the thickness of the layer in units of the fine lattice, let
$\Omega_*=\prod_{i=1}^d[\alpha_i-\tau,\beta_i+\tau]\cap\mathbb{Z}^d$. Put
\begin{equation*}
\ell_i:=\beta_i-\alpha_i,\qquad \ell:=\min_i\ell_i,\qquad
\ell_{\max}:=\max_i\ell_i .
\end{equation*}
Since $\Omega_1$ is such a box, $\ell_i\ge4$ for every $i$ and $\ell_{\max}\le\varrho\,\ell$. Assume
further $\tau\le\ell$. The \emph{layer} $\mathcal{L}$ is the set of cells of $\Omega_*$ that are
not interior to $\Omega_1$; the cells of $\partial\Omega_1$ belong to $\mathcal{L}$.

Let $h_z\ge 2$, and let $p_{\mu\nu}(z_0)\subset\Omega_1$, the \emph{centre plaquette}, be a plaquette
all of whose corners are at $\ell^\infty$-distance at least $h_z$ from $\partial\Omega_1$. We write
$(z_0)_i$ for the $i$-th coordinate of $z_0$. After a permutation of the
coordinate axes and reflections of coordinates, the face of $\Omega_1$ nearest to $z_0$ is
$\{x_1=\alpha_1\}$, so that
\begin{equation*}
h_z\le (z_0)_1-\alpha_1=\operatorname{dist}_\infty(z_0,\partial\Omega_1)\le \ell_1/2 .
\end{equation*}
Put $s_z:=(\alpha_1,(z_0)_2,\dots,(z_0)_d)$ and
\begin{equation*}
o:=s_z-\tau e_1 .
\end{equation*}
The site $o$ lies in the \emph{primary face} $\{x_1=\alpha_1-\tau\}$ of $\Omega_*$. The
\emph{opposite faces} are $\{x_1=\beta_1\}$ of $\Omega_1$ and $\{x_1=\beta_1+\tau\}$ of $\Omega_*$.
For a cell $c$ put $\rho(c):=1+|z_0-x_c|_\infty$, where $x_c$ is the lowest corner of $c$, as in
Notation~\ref{8.11:not-box-cochains}.

\emph{Paths on $\partial\Omega_*$.} A site $y\in\partial\Omega_*$ is \emph{of type A} if
$y_1=\alpha_1-\tau$ or if $y_i$ is $\Omega_*$-extreme for some $i\ge 2$. For such $y$ and
$1\le j\le d$ put
\begin{equation*}
y^{[j]}:=(o_1,y_2,\dots,y_j,o_{j+1},\dots,o_d),
\end{equation*}
so that $y^{[1]}=o$. Define the legs
\begin{equation*}
P_y^{(j)}:=[y^{[j-1]}\to y^{[j]}]\ \ (2\le j\le d),\qquad P_y^{(1)}:=[y^{[d]}\to y],
\end{equation*}
and the path
\begin{equation}\label{8.11:eq-explicit-system-1}
P_y:=P_y^{(2)}+P_y^{(3)}+\dots+P_y^{(d)}+P_y^{(1)}.
\end{equation}
Thus $P_y$ runs from $o$ through $(o_1,y_2,o_3,\dots,o_d)$, $(o_1,y_2,y_3,o_4,\dots,o_d)$, \dots,
$(o_1,y_2,\dots,y_d)$, and then along $e_1$ to $y$. The first $d-1$ legs lie in the primary face. The
last leg lies in a face orthogonal to $e_i$ for some $i\ge 2$, or is empty.

A site $y\in\partial\Omega_*$ not of type A is \emph{of type B}; it lies in the open opposite face.
For every site $u$ of the opposite face of $\Omega_*$ put
$u^{\downarrow}:=(u_1,\dots,u_{d-1},\alpha_d-\tau)$. If $y$ is of type B, $y^{\downarrow}$ is of type A
and we set
\begin{equation*}
P_y:=P_{y^{\downarrow}}+[y^{\downarrow}\to y],
\end{equation*}
the last segment running along $e_d$.

The \emph{rim} of the opposite face of $\Omega_*$ consists of its sites with some $\Omega_*$-extreme
coordinate $y_i$, $i\ge 2$. It is divided into the bottom rim $\{x_d=\alpha_d-\tau\}$, the top rim
$\{x_d=\beta_d+\tau\}$, and the side rims $\{x_i\ \text{$\Omega_*$-extreme}\}$, $2\le i\le d-1$.

\emph{Chains on $\partial\Omega_*$.} To every $\partial$-bond $\beta=\langle y,y+e_k\rangle$ of
$\Omega_*$ we assign an integral $2$-chain $S_*(\beta)\subset\partial\Omega_*$ with
\begin{equation}\label{8.11:eq-explicit-system-2}
\partial S_*(\beta)=P_y+\beta-P_{y+e_k},
\end{equation}
and we put $S_*(-\beta):=-S_*(\beta)$.

Let $\sigma$ be a straight path in $\partial\Omega_*$ from $u$ to $v$, all of whose sites are of
type A and all of whose bonds are $\partial$-bonds. We extend $S_*$ additively to such paths:
$S_*(\sigma):=\sum_{\beta''\in\sigma}S_*(\beta'')$, where each $\beta''$ is oriented along $\sigma$.

The chains $S_*(\beta)$ are as follows.
\begin{itemize}
\item[(A/A, $k=1$)] Both endpoints are of type A and $k=1$. The loop
$P_y+\beta-P_{y+e_1}$ is degenerate, and $S_*(\beta):=0$.
\item[(A/A, $k\ge2$)] Both endpoints are of type A and $k\ge 2$. Then
$S_*(\beta):=\sum_{j=k+1}^d\mathrm{St}_j(\beta)+\mathrm{St}_1(\beta)$. Here $\mathrm{St}_j(\beta)$
($j>k$) is the strip of the $|y_j-o_j|$ plaquettes $(k,j)$ in the primary face between the
corresponding legs $P^{(j)}_y$ and $P^{(j)}_{y+e_k}$ of the two paths. $\mathrm{St}_1(\beta)$ is the
strip of plaquettes $(1,k)$ at the sites $(t,y_2,\dots,y_d)$, $t$
between $o_1$ and $y_1$. The chain satisfies
\begin{equation}\label{8.11:eq-explicit-system-3}
|S_*(\beta)|\le|y-o|_1 .
\end{equation}
\item[(B/B)] Both endpoints are of type B. If $k=d$, put $S_*(\beta):=0$. If $k<d$, put
$\beta^{\downarrow}:=\langle y^{\downarrow},y^{\downarrow}+e_k\rangle$ and let $D(\beta)$ be the strip
of plaquettes $(k,d)$ in the opposite face above $\beta^{\downarrow}$, up to $\beta$. Then
$S_*(\beta):=D(\beta)+S_*(\beta^{\downarrow})$, and
$|S_*(\beta)|\le(d+1)(\ell_{\max}+2\tau)$.
\item[(mixed)] The bond $\beta$ lies in the opposite face, and exactly one endpoint lies on its rim.
\begin{itemize}
\item If $k=d$ and $y$ lies on the bottom rim, so that $\beta=\langle y,y+e_d\rangle$ with
$y=(y+e_d)^{\downarrow}$, put $S_*(\beta):=0$.
\item If $k<d$, let $y'\in\{y,y+e_k\}$ be the endpoint on a side rim and put
$\beta^{\downarrow}:=\langle y^{\downarrow},y^{\downarrow}+e_k\rangle$. Then
\begin{equation}\label{8.11:eq-explicit-system-4}
P_y+\beta-P_{y+e_k}=\partial S_*(\beta^{\downarrow})+\partial D(\beta)
+\sum_{\beta''\in[y'^{\downarrow}\to y']}\partial S_*(\beta'')
\end{equation}
when $y'=y+e_k$, and the same identity holds with the sign of the last sum reversed when $y'=y$.
The $\beta''$ are A/A bonds of the rim column, and $D(\beta)$ is as in the case (B/B). We put
\begin{equation*}
S_*(\beta):=S_*(\beta^{\downarrow})+D(\beta)\pm S_*([y'^{\downarrow}\to y']),
\end{equation*}
with the sign of \eqref{8.11:eq-explicit-system-4}.
\item If $k=d$ and $y+e_d$ lies on the top rim, the column $[y^{\downarrow}\to y+e_d]$ is slid
inside the opposite face to the side rim $\{x_2=\alpha_2-\tau\}$, and A/A loops are telescoped
along the bottom rim, the rim column and the top rim. Precisely, put
\begin{gather*}
u_0:=y^{\downarrow},\qquad u_1:=y+e_d,\\
u_0':=u_0-(y_2-\alpha_2+\tau)e_2,\qquad u_1':=u_1-(y_2-\alpha_2+\tau)e_2 .
\end{gather*}
Let $\mathrm{Sh}(\beta)$ be the sheet of plaquettes $(2,d)$ in the opposite face swept by the
column when it is slid from $[u_0'\to u_1']$ to $[u_0\to u_1]$. Then
\begin{align*}
S_*(\beta):=\mathrm{Sh}(\beta)&+S_*([u_0'\to u_1'])\\
&+S_*([u_1'\to u_1])-S_*([u_0'\to u_0]).
\end{align*}
\end{itemize}
In the last two cases $|S_*(\beta)|\le c_d(\ell_{\max}+2\tau)^2$, and the number of bonds $\beta$ of
the mixed kind is at most $c_d(\ell_{\max}+2\tau)^{d-2}$; here $c_d$ is a constant depending only on
$d$, the same in the two bounds.
\end{itemize}

\emph{The system.} For $x\in\partial\Omega_1$ let $\iota(x)$ be the least $i$ such that $x_i$ is
$\Omega_1$-extreme. Let $\nu(x):=\pm e_{\iota(x)}$ be the outward unit vector, namely $-e_{\iota(x)}$
if $x_{\iota(x)}=\alpha_{\iota(x)}$ and $+e_{\iota(x)}$ if $x_{\iota(x)}=\beta_{\iota(x)}$. Put
$\pi(x):=x+\tau\nu(x)\in\partial\Omega_*$, let $n_x:=[\pi(x)\to x]$ be the straight path from
$\pi(x)$ to $x$, and set
\begin{equation*}
\gamma_x:=P_{\pi(x)}+n_x .
\end{equation*}

Let $b'=\langle x,x'\rangle\in\mathfrak{b}(\Omega_1)$ with $x'=x+e_k$.
\begin{itemize}
\item[Case I] $\nu(x)=\nu(x')=\pm e_\iota$. Put $\pi b':=\langle\pi(x),\pi(x')\rangle$ and let
$N(b')$, the \emph{normal strip}, be the strip of the $\tau$ plaquettes $(\iota,k)$ between
$\pi b'$ and $b'$. Its boundary is
\begin{equation*}
\partial N(b')=n_x+b'-n_{x'}-\pi b'.
\end{equation*}
Set
\begin{equation*}
S(b'):=N(b')+S_*(\pi b').
\end{equation*}
\item[Case II] $\nu(x)\ne\nu(x')$. Then $b'$ lies in exactly one face, orthogonal to $e_m$ say. Its
endpoints differ only in the $k$-th coordinate, so exactly one of them has its $k$-th coordinate
$\Omega_1$-extreme. We call that endpoint $x$ and the other $x'$, read $b'$ from $x$ to $x'$, and
replace $S(b')$ by its negative for the opposite orientation. Then $k<m$, $\iota(x)=k$ and
$\iota(x')=m$.

The \emph{corner path} $c(b')$ on $\partial\Omega_*$ runs from $\pi(x)$ outward along $e_m$, in
$\tau$ steps within the face orthogonal to $e_k$, to the edge of $\Omega_*$. It then runs in the
direction $-\nu(x)$, parallel to $e_k$, in $\tau+1$ steps within the face orthogonal to $e_m$, to
$\pi(x')$. The loop
$n_x+b'-n_{x'}-c(b')$ bounds the planar rectangle $\mathrm{Sq}(b')$ of $(\tau+1)\tau$ plaquettes
$(k,m)$, all in $\mathcal{L}$. Set
\begin{equation*}
S(b'):=\mathrm{Sq}(b')+\sum_{\beta\in c(b')}S_*(\beta),
\end{equation*}
each $\beta$ oriented along $c(b')$. The boundary of the last sum telescopes to
$P_{\pi(x)}+c(b')-P_{\pi(x')}$.
\end{itemize}

\emph{Classes of bonds.} $\mathfrak{b}_{II}$ is the set of Case II bonds. $\mathfrak{b}_{\mathrm{far}}$
is the set of bonds of $\mathfrak{b}(\Omega_1)$ with an endpoint in the opposite face of $\Omega_1$.
$\mathfrak{b}_{\mathrm{near}}$ is the set of the remaining bonds of $\mathfrak{b}(\Omega_1)$.
\end{definition}

\begin{proof}
We show that the objects above are well defined and that the stated containments, boundary identities
and bounds hold. It follows that $\mathfrak{Sy}=(\{\gamma_x\},\{S(b')\})$, with base site $o$, is a
path--surface system in the sense of Definition~\ref{8.11:def-path-surface}.

\emph{The paths $P_y$.} Let $y$ be of type A. The points $y^{[j]}$ have first coordinate
$o_1=\alpha_1-\tau$ and all coordinates within the ranges of $\Omega_*$. Hence the legs
$P^{(2)}_y,\dots,P^{(d)}_y$ lie in the primary face. If $y_1=o_1$, the leg $P^{(1)}_y$ is empty.
Otherwise $y_i$ is $\Omega_*$-extreme for some $i\ge2$. Since $P^{(1)}_y$ keeps $y_2,\dots,y_d$ fixed,
it lies in the face $\{x_i=y_i\}$. Thus $P_y\subset\partial\Omega_*$ runs from $o$ to $y$.

Let $y$ be of type B. Then $y_1\ne\alpha_1-\tau$ and no $y_i$, $i\ge2$, is extreme, so
$y_1=\beta_1+\tau$. The site $y^{\downarrow}$ has the extreme coordinate $\alpha_d-\tau$ with $d\ge2$,
so it is of type A. The segment $[y^{\downarrow}\to y]$ lies in the opposite face.

\emph{The sums $S_*(\sigma)$.} Let $\sigma=[u\to v]$ be as in the definition. Each bond
$\beta''=\langle a,a'\rangle$ of $\sigma$, oriented along $\sigma$, satisfies
\begin{equation*}
\partial S_*(\beta'')=P_a+\beta''-P_{a'} .
\end{equation*}
For a negatively oriented bond this follows from \eqref{8.11:eq-explicit-system-2} and
$S_*(-\beta)=-S_*(\beta)$. Summing along $\sigma$, the intermediate paths cancel, and
$\partial S_*(\sigma)=P_u+\sigma-P_v$.

\emph{Case (A/A, $k=1$).} We have $y^{[d]}=(y+e_1)^{[d]}$, and $y_1\ge o_1$. Hence
$P^{(1)}_{y+e_1}=P^{(1)}_y+\beta$ as chains, and all other legs agree. Therefore
$P_{y+e_1}=P_y+\beta$, the loop $P_y+\beta-P_{y+e_1}$ is the zero chain, and $S_*=0$ satisfies
\eqref{8.11:eq-explicit-system-2}.

\emph{Case (A/A, $k\ge2$).} For $j<k$ the legs $P^{(j)}_y$ and $P^{(j)}_{y+e_k}$ coincide. The $k$-th
legs start at the same point and end at $y^{[k]}$ and $y^{[k]}+e_k$; as chains,
$P^{(k)}_{y+e_k}-P^{(k)}_y$ is the bond $\langle y^{[k]},y^{[k]}+e_k\rangle$ in direction $e_k$ at
$y^{[k]}$. For $j>k$ and for $j=1$ the leg $P^{(j)}_{y+e_k}$ is the translate of $P^{(j)}_y$ by $e_k$.

Let $r_j$ and $r_j'$ be the bonds in direction $e_k$ at the initial and at the final point of
$P^{(j)}_y$. Orient $\mathrm{St}_j(\beta)$ so that
\begin{equation*}
\partial\,\mathrm{St}_j(\beta)=P^{(j)}_y+r_j'-P^{(j)}_{y+e_k}-r_j .
\end{equation*}
If $P^{(j)}_y$ is empty, then $\mathrm{St}_j(\beta)=0$ and $r_j=r_j'$. The final point of each leg is
the initial point of the next, so $r_{k+1}=\langle y^{[k]},y^{[k]}+e_k\rangle$, $r_j'=r_{j+1}$ for
$k<j<d$, $r_d'=r_1$, and $r_1'=\beta$; when $k=d$, $r_{k+1}$ is to be read as $r_1$.
Summing,
\begin{align*}
\partial\Bigl(\sum_{j=k+1}^d\mathrm{St}_j(\beta)+\mathrm{St}_1(\beta)\Bigr)
&=\sum_{j\in\{k+1,\dots,d,1\}}\bigl(P^{(j)}_y-P^{(j)}_{y+e_k}\bigr)\\
&\quad-r_{k+1}+\beta .
\end{align*}
The right-hand side equals $P_y+\beta-P_{y+e_k}$.

The strips $\mathrm{St}_j(\beta)$, $j>k$, lie in the primary face. We show that $\mathrm{St}_1(\beta)$
lies in $\partial\Omega_*$. If
$y_1=o_1$ it is empty. If $y_1\ne o_1$, then $y$ and $y+e_k$ have a common $\Omega_*$-extreme
coordinate $i\notin\{1,k\}$, and $\mathrm{St}_1(\beta)$, which keeps every coordinate other than $1$
and $k$ fixed,
lies in the face $\{x_i=y_i\}$. To see that such an $i$ exists, note first that $\beta$ is not in the
primary face, because $y_1\ne o_1$.
\begin{itemize}
\item If $\beta$ lies in a side face orthogonal to $e_i$, take that $i$; then $i\ne1$, and $i\ne k$
because $\beta$ has direction $e_k$.
\item If $\beta$ lies in the opposite face, then $y+e_k$, being of type A with
$(y+e_k)_1=\beta_1+\tau$, has an extreme coordinate $i\ge2$. If $i\ne k$, $y$ shares it. If $k$ is the
only such coordinate, then $y$'s own extreme coordinate $i'\ge2$ is $\ne k$, because the side
$\ell_k+2\tau$ exceeds $1$ and so $y_k$ and $y_k+1$ cannot both be extreme; and $y+e_k$ shares
$i'$.
\end{itemize}
The strips $\mathrm{St}_j(\beta)$, $j>k$, contain $|y_j-o_j|$ plaquettes each, and
$\mathrm{St}_1(\beta)$ contains $|y_1-o_1|$.
Hence $|S_*(\beta)|\le\sum_{j\ne k}|y_j-o_j|\le|y-o|_1$, which is
\eqref{8.11:eq-explicit-system-3}. Since every coordinate difference inside $\Omega_*$ is at most
$\ell_{\max}+2\tau$, this also gives $|S_*(\beta)|\le d(\ell_{\max}+2\tau)$.

\emph{Case (B/B).} Both endpoints have first coordinate $\beta_1+\tau$, so $k\ge2$.
\begin{itemize}
\item If $k=d$, then $(y+e_d)^{\downarrow}=y^{\downarrow}$ and $y_d\ge\alpha_d-\tau$. Hence
$P_{y+e_d}=P_y+\beta$, and $S_*=0$ satisfies \eqref{8.11:eq-explicit-system-2}.
\item If $k<d$, both endpoints of $\beta^{\downarrow}$ have the extreme coordinate $\alpha_d-\tau$. So
$\beta^{\downarrow}$ is an A/A bond with $k\ge2$, and $S_*(\beta^{\downarrow})$ is defined. The strip
$D(\beta)$ lies in the opposite face and, suitably oriented, satisfies
\begin{equation*}
\partial D(\beta)=[y^{\downarrow}\to y]+\beta-[y^{\downarrow}+e_k\to y+e_k]-\beta^{\downarrow}.
\end{equation*}
Adding $\partial S_*(\beta^{\downarrow})=P_{y^{\downarrow}}+\beta^{\downarrow}-P_{y^{\downarrow}+e_k}$
and using the definition of $P$ at type-B sites gives \eqref{8.11:eq-explicit-system-2}.

$D(\beta)$ has $y_d-\alpha_d+\tau\le\ell_{\max}+2\tau$ plaquettes. By the A/A bound,
$|S_*(\beta^{\downarrow})|\le d(\ell_{\max}+2\tau)$. Hence
$|S_*(\beta)|\le(d+1)(\ell_{\max}+2\tau)$.
\end{itemize}

\emph{Mixed bonds: the cases are exhaustive.} Let one endpoint be of type A and the other of type B.
A type-B site lies in the opposite face and in no other face of $\Omega_*$, and a $\partial$-bond lies
in a face containing both its endpoints. Hence $\beta$ and both its endpoints lie in the opposite face,
so $k\ge2$. Since the endpoints differ only in the $k$-th
coordinate, and the type-B endpoint has no extreme coordinate $i\ge2$, the $k$-th coordinate is the
only extreme coordinate $i\ge 2$ of the type-A endpoint.
\begin{itemize}
\item If $k<d$, the type-A endpoint $y'$ lies on a side rim.
\item If $k=d$, either $y_d=\alpha_d-\tau$, in which case $y$ is on the bottom rim and
$y=(y+e_d)^{\downarrow}$, or $(y+e_d)_d=\beta_d+\tau$, in which case $y+e_d$ is on the top rim. The
other two possibilities would put $y$ or $y+e_d$ outside $\Omega_*$.
\end{itemize}
These are the three mixed cases of the definition.

Together with the cases above this covers every $\partial$-bond. A bond with a type-B endpoint has
that endpoint in the open opposite face. If its direction is $e_1$, its other endpoint has first
coordinate $\beta_1+\tau-1$ and no extreme coordinate, so the bond is not a $\partial$-bond.

\emph{Mixed bonds: the identities and bounds.}
\begin{itemize}
\item $k=d$, $y$ on the bottom rim. Here $P_{y+e_d}=P_y+[y\to y+e_d]=P_y+\beta$, and $S_*=0$.
\item $k<d$. All sites of the column $[y'^{\downarrow}\to y']$ have the extreme $k$-th coordinate of
$y'$. So they are of type A, and its bonds are A/A bonds with direction $e_d$, $d\ge2$. By the
telescoping above,
\begin{equation*}
\partial S_*([y'^{\downarrow}\to y'])=P_{y'^{\downarrow}}+[y'^{\downarrow}\to y']-P_{y'} .
\end{equation*}
Let $y'=y+e_k$. Then
\begin{align*}
P_y&=P_{y^{\downarrow}}+[y^{\downarrow}\to y],\\
P_{y+e_k}&=P_{y^{\downarrow}+e_k}+[y^{\downarrow}+e_k\to y+e_k]
-\partial S_*([y'^{\downarrow}\to y']).
\end{align*}
Inserting these and the boundaries of $S_*(\beta^{\downarrow})$ and $D(\beta)$ from the (B/B)
computation gives \eqref{8.11:eq-explicit-system-4}. If $y'=y$, the same substitution is made in
$P_y$, and the sign of the last sum is reversed.

For the size: $|S_*(\beta^{\downarrow})|\le d(\ell_{\max}+2\tau)$ and
$|D(\beta)|\le\ell_{\max}+2\tau$. The column has at most $\ell_{\max}+2\tau$ bonds, and for each of
them $|S_*(\beta'')|\le d(\ell_{\max}+2\tau)$. Since $\ell_{\max}+2\tau\ge1$, the total satisfies
$|S_*(\beta)|\le(2d+1)(\ell_{\max}+2\tau)^2$.
\item $k=d$, $y+e_d$ on the top rim. Since $y$ is of type B, $y_2\ne\alpha_2-\tau$. All sites of
$[u_0'\to u_0]$ lie on the bottom rim, all sites of $[u_1'\to u_1]$ on the top rim, and all sites of
$[u_0'\to u_1']$ on the side rim $\{x_2=\alpha_2-\tau\}$. So all are of type A, and all bonds are
$\partial$-bonds with directions $e_2$ or $e_d$, where $2\ne d$. The sheet $\mathrm{Sh}(\beta)$ lies
in the opposite face and, oriented suitably, satisfies
\begin{equation*}
\partial\,\mathrm{Sh}(\beta)=[u_0'\to u_0]+[u_0\to u_1]-[u_1'\to u_1]-[u_0'\to u_1'] .
\end{equation*}
Adding
\begin{equation*}
\partial S_*([u_0'\to u_1'])+\partial S_*([u_1'\to u_1])-\partial S_*([u_0'\to u_0]),
\end{equation*}
computed by
the telescoping above, leaves $P_{u_0}+[u_0\to u_1]-P_{u_1}$. This equals $P_y+\beta-P_{y+e_d}$,
because $P_y=P_{u_0}+[u_0\to y]$.

For the size: $\mathrm{Sh}(\beta)$ has $(y_2-\alpha_2+\tau)(\ell_d+2\tau)\le(\ell_{\max}+2\tau)^2$
plaquettes. The three segments contain at most $3(\ell_{\max}+2\tau)$ bonds, each contributing at most
$d(\ell_{\max}+2\tau)$. So $|S_*(\beta)|\le(3d+1)(\ell_{\max}+2\tau)^2$.
\end{itemize}
Every mixed bond has an endpoint on the rim of the opposite face. The rim consists of at most
$2(d-1)(\ell_{\max}+2\tau+1)^{d-2}$ sites, and each site carries at most $2d$ bonds. Since
$\ell_{\max}+2\tau+1\le2(\ell_{\max}+2\tau)$, the number of mixed bonds is at most
$d(d-1)2^{d}(\ell_{\max}+2\tau)^{d-2}$. Both bounds of the definition therefore hold with
$c_d:=\max\{3d+1,\,d(d-1)2^d\}$, which depends only on $d$.

\emph{The paths $\gamma_x$.} For $x\in\partial\Omega_1$, the coordinate $x_{\iota(x)}$ is
$\Omega_1$-extreme. So $\pi(x)$ has the $\Omega_*$-extreme coordinate
$\alpha_{\iota(x)}-\tau$ or $\beta_{\iota(x)}+\tau$, and lies in $\partial\Omega_*$. The path
$n_x$ lies in $\Omega_*$. Thus $\gamma_x=P_{\pi(x)}+n_x$ is a lattice path in $\Omega_*$ from $o$ to
$x$.

\emph{Case I.} Since $\ell_k\ge 4$, $x_k$ and $x'_k=x_k+1$ are not both $\Omega_1$-extreme; as
$x_\iota$ and $x'_\iota$ are both extreme, $k\ne\iota$. As $\nu(x)=\nu(x')$, we
have $\pi(x')=\pi(x)+e_k$, so $\pi b'$ is a bond in the face of $\Omega_*$ orthogonal to $e_\iota$
that contains $\pi(x)$, a
$\partial$-bond of $\Omega_*$. The strip $N(b')$ is swept by $n_x$ under translation by $e_k$; it
consists of $\tau$ plaquettes $(\iota,k)$, and its boundary is as stated. By
\eqref{8.11:eq-explicit-system-2} for $\pi b'$,
\begin{equation*}
\partial S(b')=n_x+b'-n_{x'}-\pi b'+P_{\pi(x)}+\pi b'-P_{\pi(x')}=\gamma_x+b'-\gamma_{x'} .
\end{equation*}

\emph{Case II.} The endpoints agree in every coordinate except the $k$-th. Let $\mathcal{I}$ be the
set of indices $i\ne k$ such that the $i$-th coordinate of the endpoints, which is the same for both,
is $\Omega_1$-extreme. The set of $\Omega_1$-extreme coordinates of each endpoint is $\mathcal{I}$ or
$\mathcal{I}\cup\{k\}$. The two sets cannot both contain $k$, since $\ell_k>1$. If neither does, both
sets equal $\mathcal{I}$; then $\iota$ takes the same value at the two endpoints, the endpoints agree
in that coordinate, so $\nu(x)=\nu(x')$, and we are in Case I. Hence exactly one endpoint, $x$, has
$x_k$ extreme.

The bond $b'$ lies in the face orthogonal to $e_i$ exactly for $i\in\mathcal{I}$. Since
$b'\in\mathfrak{b}(\Omega_1)$, we get $\mathcal{I}=\{m\}$. Then $\iota(x')=m$ and
$\iota(x)=\min\{k,m\}$. The condition $\nu(x)\ne\nu(x')$ forces $\iota(x)\ne m$, so $k<m$ and
$\iota(x)=k$.

Now $\pi(x)=x+\tau\nu(x)$ has $\Omega_*$-extreme $k$-th coordinate and $m$-th coordinate $x_m$. The
$\tau$ steps along the outward direction $\nu(x')$ keep the $k$-th coordinate extreme, so they lie in
the face orthogonal to $e_k$. They make the $m$-th coordinate $\Omega_*$-extreme, and so end on the
edge of $\Omega_*$. The following steps in the direction $-\nu(x)$ keep the $m$-th coordinate extreme,
so they lie in the face orthogonal to $e_m$. Since $x'_k$ is one unit inside $x_k$, they reach
$x'+\tau\nu(x')=\pi(x')$ after $\tau+1$ steps.

The loop $n_x+b'-n_{x'}-c(b')$ lies in the plane through $x$ spanned by $e_k$ and $e_m$. It bounds the
rectangle of extent $\tau+1$ in the $k$-direction and $\tau$ in the $m$-direction, which is
$\mathrm{Sq}(b')$ with $(\tau+1)\tau$ plaquettes $(k,m)$. Each of them has its $m$-th coordinates in
the closed interval with endpoints $x_m\in\{\alpha_m,\beta_m\}$ and $x_m+\tau\nu(x')_m$, so
it is not interior to $\Omega_1$ and belongs to $\mathcal{L}$.

Every bond $\beta$ of $c(b')$ is a $\partial$-bond of $\Omega_*$. By
\eqref{8.11:eq-explicit-system-2} and $S_*(-\beta)=-S_*(\beta)$, the sum of
$\partial S_*(\beta)$ along $c(b')$ telescopes to $P_{\pi(x)}+c(b')-P_{\pi(x')}$. Hence
\begin{equation*}
\partial S(b')=n_x+b'-n_{x'}-c(b')+P_{\pi(x)}+c(b')-P_{\pi(x')}=\gamma_x+b'-\gamma_{x'} .
\end{equation*}
If $b'$ is oriented from $x'$ to $x$, both sides change sign.

All chains $S(b')$ are integral cellular $2$-chains in $\Omega_*$. Together with the paths
$\gamma_x$ and the base site $o\in\Omega_*$, they satisfy the requirements of
Definition~\ref{8.11:def-path-surface}.
\end{proof}

\begin{lemma}[Geometry of the explicit path--surface system]\label{8.11:lem-system-geometry}
Let the boxes $\Omega_1=\prod_i[\alpha_i,\beta_i]\subset\Omega_*$, the layer thickness $\tau$, the
layer $\mathcal{L}$, the lengths $\ell$, $\ell_{\max}$, the point $z_0$, the height $h_z$, the
function $\rho$ on cells and the pair $(\{\gamma_x\},\{S(b')\})$ be those of
Definition~\ref{8.11:def-explicit-system}, and write $\mathfrak{b}:=\mathfrak{b}(\Omega_1)$, with
its classes $\mathfrak{b}_{II}$, $\mathfrak{b}_{\mathrm{near}}$, $\mathfrak{b}_{\mathrm{far}}$ as
fixed there. For $b'=\langle x,x'\rangle\in\mathfrak{b}$ with $x'=x+e_k$, and, in Case~II, with
$e_m$ the normal of the face of $\Omega_1$ containing $b'$, so that $k<m$, as in
Definition~\ref{8.11:def-explicit-system}, put
$S^{\mathrm{lay}}(b'):=N(b')$ and $S^{*}(b'):=S_*(\pi b')$ if $b'$ is in Case~I, and
$S^{\mathrm{lay}}(b'):=\mathrm{Sq}(b')$ and $S^{*}(b'):=\sum_{\beta\in c(b')}S_*(\beta)$ if $b'$
is in Case~II; write $\rho:=\rho(p(b'))$, where $p(b')$ is the plaquette of $\Omega_1$ containing
$b'$ and not contained in $\partial\Omega_1$.
\begin{itemize}
\item[(a)] $(\{\gamma_x\},\{S(b')\})$ is a path--surface system in $\Omega_*$ in the sense of
Definition~\ref{8.11:def-path-surface}, with $S(b')=S^{\mathrm{lay}}(b')+S^{*}(b')$,
$S^{\mathrm{lay}}(b')\in\{N(b'),\mathrm{Sq}(b')\}$ contained in $\mathcal{L}$, and
$S^{*}(b')\subset\partial\Omega_*$.
\end{itemize}
Index the plaquettes of $N(b')$ by $u\in\{0,\dots,\tau-1\}$, the plaquette with index $u$
spanning the normal distances $[u,u+1]$ from $\Omega_1$, and those of $\mathrm{Sq}(b')$ by
$(u_k,u_m)\in\{-1,\dots,\tau-1\}\times\{0,\dots,\tau-1\}$, the plaquette with index
$(u_k,u_m)$ spanning the signed outward distances $[u_k,u_k+1]$ in the direction $e_k$ and
$[u_m,u_m+1]$ in the direction $e_m$; put $u:=\max(u_k,u_m)$. Then the corner of such a
plaquette nearest to $\Omega_1$ lies at $\ell^\infty$-distance $u$ from $\Omega_1$. For
$b'\in\mathfrak{b}_{\mathrm{near}}$, every bond $\beta$ of $\partial\Omega_*$ whose surface
$S_*(\beta)$ enters $S^{*}(b')$, that is $\beta=\pi b'$ in Case~I and $\beta\in c(b')$ in
Case~II, is of type A/A: both its endpoints are of type~A for the paths $P_y$ on
$\partial\Omega_*$ used in Definition~\ref{8.11:def-explicit-system}. Finally there are constants
$c_1$, $c_2$ and, in (d), $c_d$, all depending only on $d$, such that, for all integers
$0\le u_1\le u_2\le\tau$,
\begin{equation}\label{8.11:eq-system-geometry-a}
\begin{aligned}
&\text{(b)}\ \ \#\{b'\in\mathfrak{b}:\ p\in S^{\mathrm{lay}}(b')\}\le 3
   \quad\text{for every plaquette } p \text{ of } \mathcal{L};\\
&\text{(c)}\ \ \#\{p\in N(b'):\ u_1\le u<u_2\}=u_2-u_1,\qquad
   \#\{p\in\mathrm{Sq}(b'):\ u<u_2\}\le(u_2+1)u_2,\\
&\phantom{\text{(c)}\ \ }\#\{p\in\mathrm{Sq}(b'):\ u_1\le u<u_2\}\le(u_2-u_1)(2u_2+1);\\
&\text{(d)}\ \ |S^{*}(b')|\le 2d\rho+2\tau\ \ \text{(Case I, } b'\in\mathfrak{b}_{\mathrm{near}}),\\
&\phantom{\text{(d)}\ \ }|S^{*}(b')|\le(2\tau+1)(2d\rho+4\tau)\ \
   \text{(Case II, } b'\in\mathfrak{b}_{\mathrm{near}});\\
&\phantom{\text{(d)}\ \ }\rho\ge\ell/2\ \text{ for } b'\in\mathfrak{b}_{\mathrm{far}},\qquad
   \sum_{b'\in\mathfrak{b}_{\mathrm{far}}}|S^{*}(b')|\le c_d(\ell_{\max}+2\tau)^d;\\
&\text{(e)}\ \ \sum_{b'\in\mathfrak{b}}\rho^{-d}\le\frac{c_1}{h_z},\qquad
   \sum_{b'\in\mathfrak{b}}\rho^{1-d}\le c_1\bigl(1+\log^{+}(\ell_{\max}/h_z)\bigr),\\
&\phantom{\text{(e)}\ \ }\sum_{b'\in\mathfrak{b}}\rho^{-2d}\le c_1h_z^{-(d+1)};\\
&\phantom{\text{(e)}\ \ }\sum_{b'\in\mathfrak{b}_{II}}\rho^{-d}\le\frac{c_2}{h_z^{2}},\qquad
   \sum_{b'\in\mathfrak{b}_{II}}\rho^{1-d}\le\frac{c_2}{h_z},\\
&\phantom{\text{(e)}\ \ }\sum_{b'\in\mathfrak{b}_{II}}\rho^{-2d}\le c_2h_z^{-(d+2)}.
\end{aligned}
\end{equation}
\end{lemma}

\begin{proof}
Throughout, $c_d$ denotes a constant depending only on $d$, not necessarily the same at each
occurrence; the constant $c_d$ in (d) of the statement is the value it takes at the end of the
estimate of the far bonds below. We recall from Definition~\ref{8.11:def-explicit-system} that
every corner of the base plaquette at $z_0$, in particular $z_0$, lies at $\ell^\infty$-distance
at least $h_z\ge2$ from $\partial\Omega_1$, that the face of $\Omega_1$ nearest to $z_0$ is
$\{x_1=\alpha_1\}$, that
\begin{equation*}
h_z\le z_{0,1}-\alpha_1=\operatorname{dist}_\infty(z_0,\partial\Omega_1)\le\ell_1/2,
\end{equation*}
that $s_z=(\alpha_1,z_{0,2},\dots,z_{0,d})$ and $o=s_z-\tau e_1$, and that the
opposite faces are $\{x_1=\beta_1\}$ of $\Omega_1$ and $\{x_1=\beta_1+\tau\}$ of $\Omega_*$.

(a) This holds by construction. The path $\gamma_x=P_{\pi(x)}+n_x$ runs in $\Omega_*$ from $o$
to $x$. In Case~I, $\partial N(b')=n_x+b'-n_{x'}-\pi b'$ and
$\partial S_*(\pi b')=P_{\pi(x)}+\pi b'-P_{\pi(x')}$, so that
\begin{equation*}
\partial S(b')=P_{\pi(x)}+n_x+b'-n_{x'}-P_{\pi(x')}=\gamma_x+b'-\gamma_{x'}.
\end{equation*}
In Case~II, $\partial\,\mathrm{Sq}(b')=n_x+b'-n_{x'}-c(b')$, and the boundary of
$\sum_{\beta\in c(b')}S_*(\beta)$ telescopes to $P_{\pi(x)}+c(b')-P_{\pi(x')}$; the sum is again
$\gamma_x+b'-\gamma_{x'}$. The normal strip $N(b')$ lies between $\pi b'$ and $b'$, hence in
$\mathcal{L}$, the rectangle $\mathrm{Sq}(b')$ lies in $\mathcal{L}$, and every $S_*(\beta)$ lies in
$\partial\Omega_*$; hence so does $S^{*}(b')$.

(b) Let $p$ be a plaquette of $\mathcal{L}$ spanned by $e_i$ and $e_{i'}$. If $p$ lies in
$N(b')$, then one of $e_i$, $e_{i'}$ is the normal direction of the strip and the other is the
direction of $b'$, and $b'$ is the bond obtained by sliding $p$ along the normal direction onto
$\partial\Omega_1$. Since the normal direction is $e_i$ or $e_{i'}$, there are at most two such
bonds. If $p$ lies in $\mathrm{Sq}(b')$, then $b'$ is determined by the plane of $p$, by the
codimension-two edge of $\Omega_1$ beside which $p$ lies, namely the one on which the coordinates
in the two directions of the plane of $p$ are extreme (the rectangles $\mathrm{Sq}(b')$ of the
bonds along one such edge being disjoint translates of one another along the edge), and by the
rule $k<m$, which decides which of the two directions of the plane is the direction of $b'$;
there is at most one such bond. In total $p$ lies in $S^{\mathrm{lay}}(b')$ for at most three
bonds $b'$.

(c) In $N(b')$ the plaquette with index $u$ spans the normal distances $[u,u+1]$ from
$\Omega_1$, and all its coordinates other than the normal one lie in the corresponding intervals
$[\alpha_j,\beta_j]$. For a point all of whose coordinates but the normal one lie in these
intervals, the $\ell^\infty$-distance from $\Omega_1$ equals its normal distance; hence the
nearest corner lies at distance $u$. In $\mathrm{Sq}(b')$ the outward distances of the nearest
corner of the plaquette with index $(u_k,u_m)$ in the directions $e_k$ and $e_m$ are
$u_k^{+}:=\max(u_k,0)$ and $u_m$, and all other coordinates lie inside the corresponding
intervals; its $\ell^\infty$-distance from $\Omega_1$ is therefore
$\max(u_k^{+},u_m)=\max(u_k,u_m)=u$, because $u_m\ge0$. It remains to count lattice points.
The strip contains exactly one plaquette for each $u\in\{0,\dots,\tau-1\}$, which gives
$\#\{p\in N(b'):u_1\le u<u_2\}=u_2-u_1$ for $0\le u_1\le u_2\le\tau$. For an integer $v$ with
$0\le v\le\tau$, a plaquette of $\mathrm{Sq}(b')$ has $u<v$ exactly when
$u_k\in\{-1,\dots,v-1\}$ and $u_m\in\{0,\dots,v-1\}$, and these ranges lie within the index
ranges because $v\le\tau$; so the set of such plaquettes has exactly $(v+1)v$ elements. With
$v=u_2$ this gives the second count, and with $v=u_1$ and $v=u_2$
\begin{equation*}
\#\{p\in\mathrm{Sq}(b'):u_1\le u<u_2\}=(u_2+1)u_2-(u_1+1)u_1=(u_2-u_1)(u_2+u_1+1),
\end{equation*}
which is at most $(u_2-u_1)(2u_2+1)$.

(d) Let $b'=\langle x,x'\rangle\in\mathfrak{b}_{\mathrm{near}}$. No endpoint of $b'$ lies in
the opposite face of $\Omega_1$. Since $\pi(x)=x+\tau\nu(x)$, the point $\pi(x)$ lies in the
opposite face of $\Omega_*$ only if $x_1=\beta_1$; so $\pi(x)$, $\pi(x')$ and the corner path
$c(b')$ avoid the open opposite face of $\Omega_*$, and every bond $\beta$ whose surface
$S_*(\beta)$ is a summand of $S^{*}(b')$ is of type A/A. For such a bond
$\beta=\langle y,y+e_k\rangle$ the construction of
$S_*(\beta)$ gives $|S_*(\beta)|\le|y-o|_1$. In Case~I, $S^{*}(b')=S_*(\pi b')$ with $y=\pi(x)$,
and
\begin{align*}
|\pi(x)-o|_1&\le|x-s_z|_1+2\tau
\le d\bigl(|x-z_0|_\infty+|z_0-s_z|_\infty\bigr)+2\tau\\
&\le 2d|x-z_0|_\infty+2\tau\le 2d\rho+2\tau.
\end{align*}
The first inequality holds because $\pi(x)-o=(x-s_z)+\tau\nu(x)+\tau e_1$; the second because
$|v|_1\le d|v|_\infty$ and by the triangle inequality; the third because
$|z_0-s_z|_\infty=z_{0,1}-\alpha_1=\operatorname{dist}_\infty(z_0,\partial\Omega_1)$, which is at
most $|x-z_0|_\infty$ since $x\in\partial\Omega_1$; the fourth because $x$ is a corner of $p(b')$,
so that $|x-z_0|_\infty\le\rho$. In Case~II, $S^{*}(b')$ is the sum of the surfaces $S_*(\beta)$
over the $2\tau+1$ bonds $\beta$ of $c(b')$; passing along $c(b')$ from $\pi(x)$ adds at most
$2\tau$ to the above bound for $|y-o|_1$, so each of these surfaces has at most $2d\rho+4\tau$
plaquettes, and $|S^{*}(b')|\le(2\tau+1)(2d\rho+4\tau)$.

Now let $b'\in\mathfrak{b}_{\mathrm{far}}$. Then
\begin{equation*}
\operatorname{dist}_\infty(z_0,\{x_1=\beta_1\})=\beta_1-z_{0,1}=\ell_1-(z_{0,1}-\alpha_1),
\end{equation*}
which is at least $\ell_1/2\ge\ell/2$ because $z_{0,1}-\alpha_1\le\ell_1/2$; since $b'$ has an
endpoint in this face, $\rho\ge\ell/2$. There are at most $c_d\ell_{\max}^{d-1}$ far bonds. The
bonds $\beta$ whose surfaces $S_*(\beta)$ are summands of the surfaces $S^{*}(b')$ are one for
each far bond in Case~I and $2\tau+1$
for each far bond in Case~II, and there are at most $c_d\ell_{\max}^{d-2}$ far bonds in Case~II.
By the construction of the surfaces $S_*(\beta)$, these bonds have
$|S_*(\beta)|\le(d+1)(\ell_{\max}+2\tau)$, except the at most $c_d(\ell_{\max}+2\tau)^{d-2}$
bonds of mixed type, which lie in the opposite face of $\Omega_*$ with exactly one endpoint on its
rim, and for which $|S_*(\beta)|\le c_d(\ell_{\max}+2\tau)^2$. Summing, and using
$2\tau+1\le\ell_{\max}+2\tau$,
\begin{align*}
\sum_{b'\in\mathfrak{b}_{\mathrm{far}}}|S^{*}(b')|
&\le c_d\bigl(\ell_{\max}^{d-1}+(2\tau+1)\ell_{\max}^{d-2}\bigr)(d+1)(\ell_{\max}+2\tau)\\
&\quad+c_d(\ell_{\max}+2\tau)^{d-2}(\ell_{\max}+2\tau)^{2}\\
&\le c_d(\ell_{\max}+2\tau)^d.
\end{align*}

(e) For an integer $j\ge1$ let $\mathfrak{b}(j)$ be the set of $b'\in\mathfrak{b}$
with $\rho(p(b'))\in[j,j+1)$; these sets exhaust $\mathfrak{b}$, since $\rho\ge1$.
We count them. Every $b'\in\mathfrak{b}$ lies in one of the $2d$ faces of $\Omega_1$ and has one
of the $d-1$ directions parallel to that face. Fix a face and a direction. The plaquettes $p(b')$
of the bonds $b'$ with this face and this direction are translates of one another along the face;
hence the coordinate normal to the face of the point of $p(b')$ at which $\rho$ is evaluated takes
one value for all of them, and its distance from the corresponding coordinate of $z_0$ is a number
$D^{\perp}\ge0$ depending only on the face and the direction. Writing $D^{\parallel}(b')$ for the
largest of the distances between the remaining $d-1$ coordinates of that point and those of $z_0$,
we have $\rho(p(b'))=1+\max(D^{\perp},D^{\parallel}(b'))$, and the map sending $b'$ to this point
is injective into a translate of the $(d-1)$-dimensional integer lattice of the hyperplane of the
face. Consequently: if $j\le D^{\perp}$, no such bond lies in $\mathfrak{b}(j)$; if
$D^{\perp}<j\le D^{\perp}+1$, which happens for exactly one integer $j$, the bonds of
$\mathfrak{b}(j)$ have $D^{\parallel}(b')<j$, and there are at most $c_dj^{\,d-1}$ of them; if
$j>D^{\perp}+1$, they have $D^{\parallel}(b')\in[j-1,j)$, so their points lie in the trace, on the
hyperplane of the face, of an $\ell^\infty$-shell of thickness one around $z_0$ at radius $j$, and
there are at most $c_dj^{\,d-2}$ of them. We call the integer $j$ with
$D^{\perp}<j\le D^{\perp}+1$ the tangent shell of the face and the direction. Summing over the
$2d(d-1)$ faces and directions, we obtain
$\#\mathfrak{b}(j)\le c_dj^{\,d-2}$ for every $j$ outside a set $E$ of at most $2d(d-1)$
tangent shells, and $\#\mathfrak{b}(j)\le c_dj^{\,d-1}$ for $j\in E$.

Moreover $\rho\ge h_z$ on $\mathfrak{b}$,
because $z_0$ lies at $\ell^\infty$-distance at least $h_z$ from $\partial\Omega_1\supset b'$;
so $\mathfrak{b}(j)$ is empty unless $j+1>h_z$, that is, unless
$j\ge j_0$, where $j_0$ is the least integer exceeding $h_z-1$; as $h_z\ge2$, we have
$j_0\ge1$ and $j_0>h_z-1\ge h_z/2$. In particular every $j\in E$ with $\mathfrak{b}(j)$ non-empty
satisfies $j\ge j_0$.

A Case~II bond $b'=\langle x,x+e_k\rangle$ lies in the face perpendicular to $e_m$ and has $x_k$
extreme, so it touches the codimension-two edge of $\Omega_1$ on which $x_k$ and $x_m$ are extreme.
For each of the at most $c_d$ such edges, the coordinates in the directions $e_k$ and $e_m$ of the
point of $p(b')$ at which $\rho$ is evaluated take one value for all these bonds, and the same count,
now in the remaining $d-2$ coordinates, gives
$\#(\mathfrak{b}(j)\cap\mathfrak{b}_{II})\le c_dj^{\,d-3}$ for every $j$ outside a set $E_{II}$ of
at most $c_d$ tangent shells, and $\#(\mathfrak{b}(j)\cap\mathfrak{b}_{II})\le c_dj^{\,d-2}$ for
$j\in E_{II}$; again every $j\in E_{II}$ with $\mathfrak{b}(j)\cap\mathfrak{b}_{II}$ non-empty
satisfies $j\ge j_0$.

Summing the series uses, for $a\ge2$ and integers $j_0\ge1$,
\begin{equation*}
\sum_{j\ge j_0}j^{-a}\le j_0^{-a}+\int_{j_0}^{\infty}r^{-a}\,dr\le 2j_0^{1-a}.
\end{equation*}
On $\mathfrak{b}(j)$ we have $\rho\ge j$. Therefore the shells $j\notin E$ contribute to the first
and third sums at most
\begin{align*}
\sum_{j\notin E}\sum_{b'\in\mathfrak{b}(j)}\rho^{-d}&\le c_d\sum_{j\ge j_0}j^{\,d-2}j^{\,-d}
 \le 2c_d\,j_0^{-1}\le 4c_d\,h_z^{-1},\\
\sum_{j\notin E}\sum_{b'\in\mathfrak{b}(j)}\rho^{-2d}&\le c_d\sum_{j\ge j_0}j^{\,-(d+2)}
 \le 2c_d\,j_0^{-(d+1)}\le 2^{d+2}c_d\,h_z^{-(d+1)}.
\end{align*}
For the middle sum we also bound $\rho$ from above: an endpoint of $b'$ and the point $z_0$ both
lie in $\Omega_1$, and $p(b')$ has side one, so $\rho\le\ell_{\max}+2$ and
$\mathfrak{b}(j)$ is empty for $j>\ell_{\max}+2$. Since
$\ell_{\max}\ge\ell_1\ge 2h_z\ge4$, we have $\ell_{\max}+2\le2\ell_{\max}$ and
$\log(\ell_{\max}/h_z)=\log^{+}(\ell_{\max}/h_z)$; hence the shells $j\notin E$ contribute to the
middle sum at most
\begin{align*}
\sum_{j\notin E}\sum_{b'\in\mathfrak{b}(j)}\rho^{1-d}&\le c_d\sum_{j_0\le j\le\ell_{\max}+2}j^{\,-1}
 \le c_d\Bigl(\frac1{j_0}+\log\frac{\ell_{\max}+2}{j_0}\Bigr)\\
&\le c_d\Bigl(1+\log\frac{4\ell_{\max}}{h_z}\Bigr)
 \le c_d(1+\log4)\bigl(1+\log^{+}(\ell_{\max}/h_z)\bigr).
\end{align*}
The tangent shells $j\in E$ are at most $2d(d-1)$ in number, carry at most $c_dj^{\,d-1}$ bonds
each, and, when non-empty, have $j\ge j_0\ge h_z/2$; they contribute at most
\begin{equation*}
c_dj_0^{-1}\le 2c_d\,h_z^{-1},\qquad c_d,\qquad c_dj_0^{-(d+1)}\le 2^{d+1}c_d\,h_z^{-(d+1)}
\end{equation*}
to the sums of $\rho^{-d}$, $\rho^{1-d}$ and $\rho^{-2d}$ respectively, since
$j^{\,d-1}j^{\,-d}=j^{-1}$, $j^{\,d-1}j^{\,1-d}=1$ and $j^{\,d-1}j^{\,-2d}=j^{-(d+1)}$.
For the Case~II bonds the count $c_dj^{\,d-3}$ gives in the same way, for the shells
$j\notin E_{II}$,
\begin{align*}
\sum_{j\notin E_{II}}\sum_{b'\in\mathfrak{b}(j)\cap\mathfrak{b}_{II}}\rho^{-d}
&\le c_d\sum_{j\ge j_0}j^{\,-3}\le 2c_d\,j_0^{-2}\le 8c_d\,h_z^{-2},\\
\sum_{j\notin E_{II}}\sum_{b'\in\mathfrak{b}(j)\cap\mathfrak{b}_{II}}\rho^{1-d}
&\le c_d\sum_{j\ge j_0}j^{\,-2}\le 4c_d\,h_z^{-1},\\
\sum_{j\notin E_{II}}\sum_{b'\in\mathfrak{b}(j)\cap\mathfrak{b}_{II}}\rho^{-2d}
&\le c_d\sum_{j\ge j_0}j^{\,-(d+3)}\le 2c_d\,j_0^{-(d+2)}\le 2^{d+3}c_d\,h_z^{-(d+2)},
\end{align*}
while the at most $c_d$ tangent shells $j\in E_{II}$, with at most $c_dj^{\,d-2}$ bonds each and
$j\ge j_0\ge h_z/2$ when non-empty, contribute at most
\begin{equation*}
c_dj_0^{-2}\le 4c_d\,h_z^{-2},\qquad c_dj_0^{-1}\le 2c_d\,h_z^{-1},\qquad
c_dj_0^{-(d+2)}\le 2^{d+2}c_d\,h_z^{-(d+2)}
\end{equation*}
to the sums of $\rho^{-d}$, $\rho^{1-d}$ and $\rho^{-2d}$ over $\mathfrak{b}_{II}$ respectively.
Adding the two contributions to each sum, and taking for $c_1$ the largest of the resulting
constants in the first three bounds and for $c_2$ the largest of those in the last three, all
depending only on $d$, gives (e).

The logarithm in the second bound of (e) has the following origin. By (d) and the shell count,
the sum $\sum_{b'\in\mathfrak{b}}\rho^{-d}|S^{*}(b')|$ is, shell by shell, of the order of
\begin{equation*}
\sum_{j}j^{\,d-2}\cdot j^{\,-d}\cdot j=\sum_j j^{-1},
\end{equation*}
so that every dyadic range of $j$ between $h_z$ and $\ell$ carries comparable weight.
\end{proof}

\begin{lemma}[Nonabelian defects in the axial gauge]\label{8.11:lem-nonabelian-defects}
Let $\Omega_*$ be a box of the fine lattice and let $\Omega_1\subset\Omega_*$. Write $\|\cdot\|$ for
the operator norm and $\|\cdot\|_{HS}$ for the unnormalised Hilbert--Schmidt norm,
$\|A\|_{HS}=(\sum_{i,j}|A_{ij}|^2)^{1/2}$, of $N\times N$ matrices, so that
$\|A\|_{HS}\le N^{1/2}\|A\|$. Let $U$ be a configuration such that
\begin{enumerate}
\item[(1)] $\|U(\partial p)-1\|\le\varepsilon_a\eta^2$ for every plaquette $p$ within fine distance
$8$ of $\Omega_*$;
\item[(2)] $\operatorname{diam}_1(\Omega_*)\,\varepsilon_a\eta^2\le 1/(64C_{ax})$;
\item[(3)] $U$ satisfies the lattice Yang--Mills equation exactly at every bond of $\Omega_1$.
\end{enumerate}
Fix $z_0\in\Omega_1$ and put $U$ in the complete axial gauge $\mathfrak{G}$ of $\Omega_*$ centred at
$z_0$: $U_b=1$ on the tree formed by all bonds in direction $d$, the bonds in direction $d-1$ in the
hyperplane $x_d=z_{0,d}$, and so on, down to the bonds in direction $1$ in the line
$x_2=z_{0,2},\dots,x_d=z_{0,d}$. Let $\Phi(p):=\frac1i\log U(\partial p)$, with $\partial p$ based at the
lowest corner $x_p$ of $p$; in this gauge $\Phi$ is a $\mathfrak{g}$-valued $2$-cochain. Then for every
$3$-cell $c$ of $\Omega_*$ and every bond $b$ of $\Omega_1$, with $x$ the lowest corner of $c$,
respectively of $b$,
\begin{equation}\label{8.11:eq-nonabelian-defects-1}
\|(d\Phi)(c)\|_{HS},\ \|(d^*\Phi)(b)\|_{HS}\le C_{ND}\,N\,(|x-z_0|_1+1)\,\varepsilon_a^2\eta^4,
\qquad C_{ND}=C(d).
\end{equation}
Consequently, with $K$ the axial homotopy operator of Lemma~\ref{8.11:lem-axial-homotopy} centred at
$z_0$ and $\rho_2:=K(d\Phi)$, for every plaquette $p$ of $\Omega_*$
\begin{equation}\label{8.11:eq-nonabelian-defects-2}
\|\rho_2(p)\|_{HS}\le C_{ND}\,N\,(|x_p-z_0|_1+1)^2\,\varepsilon_a^2\eta^4 .
\end{equation}
\end{lemma}

\begin{proof}
For a point $x$ of $\Omega_*$ let $Q_8(x)$ be the neighbourhood of $x$ on which the lattice calculus of
this chapter is formulated, and let $a$ and $m(x)$ be its two size parameters: every bond
$b\subset Q_8(x)$ has $|b_--x|_1\le 8d$, the plaquettes of $Q_8(x)$ lie within fine distance $8$ of
$\Omega_*$, and so $\|\Phi(p)\|_{HS}\le m(x)\le a$ for them, where $a$
is the supremum of $\|\Phi(p)\|_{HS}$ over the plaquettes $p$ within fine distance $8$ of $\Omega_*$.
We use three statements of that lattice calculus. The axial-gauge bound: in
the complete axial gauge of a box centred at $y$ the tree bonds equal $1$ and, by the nonabelian Stokes
formula, $\|U_b-1\|\le|b_--y|_1\sup_p\|U(\partial p)-1\|$; in the local axial gauge on $Q_8(x)$ one has
$\|U_b-1\|\le C_{ax}\,m(x)$ for the bonds of $Q_8(x)$. The Yang--Mills bound: in the local axial gauge on
$Q_8(x)$, $\|(d^*\Phi)(b)\|_{HS}\le 2(d-1)C_{YM}\,a\,m(x)$ at a bond $b$ with lowest corner $x$ at which
the Yang--Mills equation holds. The Bianchi bound: in the same gauge,
$\|(d\Phi)(c)\|_{HS}\le C_{BI}\,a\,m(x)$ at every $3$-cell $c$ with lowest corner $x$.

\emph{The gauge $\mathfrak{G}$ near $x$.} Put $\xi_{\mathfrak{G}}(x):=\max_{b\subset Q_8(x)}\|U_b-1\|$,
computed in $\mathfrak{G}$. The axial-gauge bound with $y=z_0$, hypothesis~(1) and
$|b_--z_0|_1\le|x-z_0|_1+8d$ for $b\subset Q_8(x)$ give
\begin{equation}\label{8.11:eq-nonabelian-defects-3}
\xi_{\mathfrak{G}}(x)\le(|x-z_0|_1+8d)\,\varepsilon_a\eta^2 .
\end{equation}

\emph{Comparison with the local gauge.} Fix $x$ and let $U^{\mathrm{loc}}=\lambda\cdot U$ be the local
axial gauge on $Q_8(x)$, where $\lambda$ is a gauge transformation normalised by $\lambda(x)=1$. This
normalisation is admissible: a constant gauge transformation conjugates every $U_b$ and every
$\Phi(p)$ by one fixed unitary matrix, which preserves $\|U_b-1\|$, $\|\Phi(p)\|_{HS}$ and the
condition $U_b=1$ on the tree, so it changes none of the three bounds above. By definition
$U^{\mathrm{loc}}_b=\lambda(b_-)U_b\lambda(b_+)^{-1}$, that is,
$U^{\mathrm{loc}}_b\lambda(b_+)=\lambda(b_-)U_b$, whence
\begin{equation*}
\lambda(b_+)-\lambda(b_-)=(1-U^{\mathrm{loc}}_b)\,\lambda(b_+)+\lambda(b_-)\,(U_b-1).
\end{equation*}
Since $\lambda(b_\pm)$ are unitary, for $b\subset Q_8(x)$ the local axial-gauge bound and the definition
of $\xi_{\mathfrak{G}}$ give
\begin{equation*}
\|\lambda(b_+)-\lambda(b_-)\|\le\|U^{\mathrm{loc}}_b-1\|+\|U_b-1\|
\le C_{ax}m(x)+\xi_{\mathfrak{G}}(x).
\end{equation*}
Every $y$ within two steps of $x$ is joined to $x$ by a path of at most two bonds of $Q_8(x)$;
telescoping along it from $\lambda(x)=1$ gives
\begin{equation}\label{8.11:eq-nonabelian-defects-4}
\|\lambda(y)-1\|\le 2\bigl(C_{ax}m(x)+\xi_{\mathfrak{G}}(x)\bigr).
\end{equation}
Let $\Phi^{\mathrm{loc}}(p):=\frac1i\log U^{\mathrm{loc}}(\partial p)$. As $\partial p$ is based at $x_p$,
$U^{\mathrm{loc}}(\partial p)=\lambda(x_p)U(\partial p)\lambda(x_p)^{-1}$, and the logarithm commutes
with conjugation, so $\Phi^{\mathrm{loc}}(p)=\operatorname{Ad}(\lambda(x_p))\Phi(p)$. In particular
$\|\Phi^{\mathrm{loc}}(p)\|_{HS}=\|\Phi(p)\|_{HS}$, so $a$ and $m(x)$ are the same for $\Phi$ and for
$\Phi^{\mathrm{loc}}$. Writing $\mu:=\lambda(x_p)$,
\begin{equation*}
\mu\Phi(p)\mu^{-1}-\Phi(p)=(\mu-1)\Phi(p)\mu^{-1}+\Phi(p)(\mu^{-1}-1),
\end{equation*}
and $\|\mu^{-1}-1\|=\|\mu^{-1}(1-\mu)\|=\|\mu-1\|$; with
$\|AB\|_{HS}\le\|A\|\,\|B\|_{HS}$ this gives
\begin{equation}\label{8.11:eq-nonabelian-defects-5}
\|\Phi(p)-\Phi^{\mathrm{loc}}(p)\|_{HS}\le 2\,\|\lambda(x_p)-1\|\,\|\Phi(p)\|_{HS}.
\end{equation}

\emph{The defects.} The coboundary $(d\Phi)(c)$ is a signed sum of the $6$ values of $\Phi$ at the
faces of $c$, and
$(d^*\Phi)(b)$ a signed sum of the $2(d-1)$ values at the plaquettes containing $b$; all these
plaquettes lie in $Q_8(x)$ and are based within two steps of $x$. For each of them
\eqref{8.11:eq-nonabelian-defects-5}, \eqref{8.11:eq-nonabelian-defects-4} and
$\|\Phi(p)\|_{HS}\le m(x)$ bound the difference by $4(C_{ax}m(x)+\xi_{\mathfrak{G}}(x))\,m(x)$. With
the Bianchi bound for $\Phi^{\mathrm{loc}}$ and $m(x)\le a$,
\begin{align*}
\|(d\Phi)(c)\|_{HS}
&\le\|(d\Phi^{\mathrm{loc}})(c)\|_{HS}+6\cdot4\bigl(C_{ax}m(x)+\xi_{\mathfrak{G}}(x)\bigr)m(x)\\
&\le\bigl[C_{BI}a+24C_{ax}a+24\xi_{\mathfrak{G}}(x)\bigr]\,m(x).
\end{align*}
Likewise, the Yang--Mills bound for $\Phi^{\mathrm{loc}}$, applicable at $b$ by hypothesis~(3), gives
\begin{equation*}
\|(d^*\Phi)(b)\|_{HS}\le 2(d-1)\bigl[C_{YM}a+4C_{ax}a+4\xi_{\mathfrak{G}}(x)\bigr]\,m(x).
\end{equation*}
Now insert $a,\,m(x)\le 2N^{1/2}\varepsilon_a\eta^2$. This follows from hypothesis~(1), which controls
every plaquette entering the supremum $a$, together with $\|A\|_{HS}\le N^{1/2}\|A\|$ and the bound
$\|\log Y\|\le 2\|Y-1\|$ for a unitary matrix $Y$ whose principal logarithm is defined, applied to
$Y=U(\partial p)$. For the latter bound: $Y$ and $\log Y$ are diagonal in one orthonormal basis, an
eigenvalue $e^{i\theta}$ of $Y$ with $|\theta|<\pi$ corresponding to the eigenvalue $i\theta$ of
$\log Y$; since $\sin t\ge 2t/\pi$ for $0\le t\le\pi/2$,
\begin{equation*}
|e^{i\theta}-1|=2|\sin(\theta/2)|\ge 2|\theta|/\pi,
\qquad\text{so}\qquad |\theta|\le\tfrac{\pi}{2}\,|e^{i\theta}-1|\le 2\,|e^{i\theta}-1|,
\end{equation*}
and the operator norm of a normal matrix is the largest modulus of its eigenvalues. Insert also
\eqref{8.11:eq-nonabelian-defects-3}. Using $|x-z_0|_1+8d\le 8d\,(|x-z_0|_1+1)$ and $N^{1/2}\le N$,
\begin{align*}
\|(d\Phi)(c)\|_{HS}&\le\bigl[4C_{BI}N+96C_{ax}N+48N^{1/2}(|x-z_0|_1+8d)\bigr]
\varepsilon_a^2\eta^4\\
&\le(4C_{BI}+96C_{ax}+384d)\,N\,(|x-z_0|_1+1)\,\varepsilon_a^2\eta^4,\\
\|(d^*\Phi)(b)\|_{HS}&\le 2(d-1)\bigl[4C_{YM}N+16C_{ax}N+8N^{1/2}(|x-z_0|_1+8d)\bigr]
\varepsilon_a^2\eta^4\\
&\le 2(d-1)(4C_{YM}+16C_{ax}+64d)\,N\,(|x-z_0|_1+1)\,\varepsilon_a^2\eta^4 .
\end{align*}
This is \eqref{8.11:eq-nonabelian-defects-1} with $C_{ND}$ the larger of the two constants, which are
built from $d$ and the lattice-calculus constants $C_{BI}$, $C_{YM}$, $C_{ax}$; this is the
$C_{ND}=C(d)$ of the statement.

\emph{The homotopy defect.} By Lemma~\ref{8.11:lem-axial-homotopy}, $(K\omega)(p)$ is a signed sum of
values of $\omega$ over at most $|x_p-z_0|_1$ cells, all contained in the box spanned by $p$ and $z_0$,
hence in $\Omega_*$, each with lowest corner coordinatewise nearer to $z_0$ than $x_p$. Applied to
$\omega=d\Phi$, each of these at most $|x_p-z_0|_1$ three-cells $c$ has $|x_c-z_0|_1\le|x_p-z_0|_1$, so
by \eqref{8.11:eq-nonabelian-defects-1} and the triangle inequality
\begin{equation*}
\|\rho_2(p)\|_{HS}\le|x_p-z_0|_1\,C_{ND}N(|x_p-z_0|_1+1)\varepsilon_a^2\eta^4
\le C_{ND}N(|x_p-z_0|_1+1)^2\varepsilon_a^2\eta^4,
\end{equation*}
which is \eqref{8.11:eq-nonabelian-defects-2}.
\end{proof}

\begin{remark}[Why the gauge is centred at $z_0$]
Let $\rho(p)$ denote $1$ plus the $\ell^\infty$-distance of $p$ to the centre. In a complete axial gauge
centred at a point far from a given plaquette, the comparison above gives defects near that plaquette
of order $\operatorname{diam}_1(\Omega_*)\,\varepsilon_a^2\eta^4$, and the Calder\'on--Zygmund sum
$\sum_p\rho(p)^{-d}\|\rho_2(p)\|_{HS}$ then carries a logarithm of the size of $\Omega_*$ in fine units,
which contains $\log\eta^{-1}$: a factor depending on the cutoff, which the smallness of
$\varepsilon_a$ cannot absorb, since $k_o\log L$ can be of order $g^{-2}$. With the gauge centred at
$z_0$, \eqref{8.11:eq-nonabelian-defects-2} makes the shell at distance $\rho$ contribute of order
$\rho^{d-1}\rho^{-d}\rho^2$, and the sum is of order $\operatorname{diam}_1(\Omega_*)^2$: no logarithm of
the mesh occurs. The expansion of $\Phi$ to first order in the bond potentials (the logarithms of the
bond variables $U_b$), read in one global gauge, would give defects proportional to the square of the
distance to the centre, losing a factor of the size of $\Omega_*$ in fine units against the covariant
comparison of the proof, which gives defects proportional to the distance.
\end{remark}

\begin{definition}[The three members of the interior bound]\label{8.11:not-bound-members}
Fix $\theta\in[1,2]$ and $\Lambda>0$. Let $r_1=r_0(\Lambda)+2/L$ and $d_1=y_++r_1$ be the tail length
and the free depth of Definition~\ref{8.11:def-released-problem}, see
\eqref{8.11:eq-released-problem-a}. Here $y_+$ and $\varpi_y=y_+-w$ are the upper end of the depth
band and the excess of Definition~\ref{8.11:def-young-row-cover}, see
\eqref{8.11:eq-young-row-cover-a}, and $w$ is the clamp width.

A \emph{reference depth} is a number
\begin{equation*}
c_*\in\{-(3w+1),\ w-1\}.
\end{equation*}
For a reference depth $c_*$ we put $t:=d_1-c_*$ and $t_c:=w-c_*$. Since
$w-(-(3w+1))=4w+1$ and $w-(w-1)=1$, we have $t_c\in\{4w+1,\,1\}$.

Let $\theta''$ and $\theta_y''$ be the plaquette bounds inherited by the released minimiser
$\check x$, and
$\sigma_e=s_e\theta+\vartheta_1$ the exterior amplitude bound, fixed in
Remark~\ref{8.11:rem-inherited-bounds}, see
\eqref{8.11:eq-inherited-bounds-a}. The exterior amplitude $s_e$ is the one of
Notation~\ref{8.11:not-units-cap-amplitude}, see \eqref{8.11:eq-units-cap-amplitude-a}. We put
\begin{equation*}
\sigma_*:=
\begin{cases}
\sigma_e=s_e\theta+\vartheta_1 & \text{if } c_*=-(3w+1),\\
\theta'' & \text{if } c_*=w-1.
\end{cases}
\end{equation*}

Recall that $H=R^{+}_{k_o}$ is the zone, $D_W$ its width and $\varrho_H$ its aspect bound, and that
$\Omega_c$ is the box $\{\delta\ge c\}$. In the present section $\varrho$ denotes the number
$2\varrho_H$ (and not the parameter $\varrho$ of the Glossary); it bounds the aspect ratio of
$\Omega_{d_1}$ whenever $4\varrho_H d_1\le D_W$, see \eqref{8.11:eq-bound-members-2} below. That
inequality is derived below under the comparison
\begin{equation}\label{8.11:eq-bound-members-3}
\varrho_H\le L\,\frac{M}{M_2}
\end{equation}
between the aspect bound of the zone and the constants $L$, $M$, $M_2$; \eqref{8.11:eq-bound-members-3}
is a condition of the aspect-ratio bound, which is asserted only where it holds.

The following constants are used.
\begin{itemize}
\item $C_\Pi$ and $C_\Pi'$ are the constants of Corollary~\ref{8.11:cor-projection-kernels}, see
\eqref{8.11:eq-projection-kernels-a}, taken at the dimension $d$ and at the aspect bound $\varrho$
just fixed.
\item $c_1$ and $c_2$ are the shell-sum constants of clause~(e) of
Lemma~\ref{8.11:lem-system-geometry}, see \eqref{8.11:eq-system-geometry-a}, and $c_d$ is the
area constant of clause~(d) of the same lemma.
\item The constants $c_3$, $c_4$, $c_5$ are defined by
\begin{equation}\label{8.11:eq-bound-members-1}
c_3:=2dc_1(1+\log(4\varrho_H))+2^{d}c_d(\varrho+2)^{d},\qquad c_4:=4c_1+12dc_2,\qquad
c_5:=48c_2 .
\end{equation}
\end{itemize}
Throughout the present section $c_1,\dots,c_5$ denote the constants just fixed in this list and in
\eqref{8.11:eq-bound-members-1}; none of them is a flow constant of the Glossary
(Notation~\ref{0.2:not-glossary}).

We write $\log^+x:=\max\{\log x,0\}$ for $x>0$. The following quantities also enter:
\begin{itemize}
\item $\varepsilon:=\varepsilon_{k_o}$ and $\varepsilon_a$ are the exterior unit parameter and the
class bound in top units of Notation~\ref{8.11:not-units-cap-amplitude};
\item $\varepsilon_1$ is the regularity of the support configuration $V^{\sharp}$ of the zone;
\item $\hat N_W$ is the block count of the zone, and $C_E$ its energy constant;
\item $N\ge 2$ is the integer of the gauge group $\mathrm{SU}(N)$ of the Glossary
(Notation~\ref{0.2:not-glossary});
\item $C_{\mathrm{nl}}=C_{\mathrm{nl}}(d,\varrho_H)$ is a constant depending only on $d$ and
$\varrho_H$.
\end{itemize}

For $h>0$ the three members of the interior bound at reading height $h$ are
\begin{equation}\label{8.11:eq-bound-members-a}
\begin{aligned}
\mathfrak{B}(h)&:=C_\Pi\Bigl\{\frac{2c_1\bigl[t_c\theta''+\varpi_y\theta_y''\bigr]}{h}\\
&\qquad+\frac{4c_2\bigl[(t+1)^2\theta''+\varpi_y(2\varpi_y+2r_1+1)\theta_y''\bigr]}{h^{2}}\\
&\qquad+\sigma_*\Bigl[c_3+2dc_1\log^+\frac{D_W}{h}+c_4\,\frac{t}{h}
+c_5\Bigl(\frac{t}{h}\Bigr)^{2}\Bigr]\Bigr\},\\
\kappa_{\mathrm{col}}(h)&:=C_\Pi\Bigl[\frac{3072\,(c_1+c_2)\,N\,C_E\,\hat N_W\,(r_1+1)}{h^{5}}
\Bigr]^{1/2}\,\frac{\varepsilon_1}{\varepsilon},\\
\vartheta_{\mathrm{nl}}(h)&:=C_{\mathrm{nl}}\,N\,(D_W+8w)^{2}\,\frac{\varepsilon_a^{2}}{\varepsilon}
\Bigl[1+\log^+\frac{D_W}{h}+\Bigl(\frac{t}{h}\Bigr)^{2}\Bigr].
\end{aligned}
\end{equation}
The member $\mathfrak{B}(h)$ depends on the reference depth $c_*$ through $t$, $t_c$ and
$\sigma_*$, and $\vartheta_{\mathrm{nl}}(h)$ through $t$; $\kappa_{\mathrm{col}}(h)$ does not
depend on $c_*$.
\end{definition}

\begin{proof}[On the aspect ratio of $\Omega_{d_1}$]
By the definition of the zone, the sides of the box $H$ have lengths at most $D_W$ and at least
$D_W/\varrho_H$; and $\Omega_{d_1}=\{\delta\ge d_1\}$ arises from $H$ by moving each face of $H$
inward by $d_1$.
Since $d_1=y_++r_1\ge w>0$, every side of $\Omega_{d_1}$ has length at most $D_W$ and at least
$D_W/\varrho_H-2d_1$. The bound on the aspect ratio therefore follows from the inequality
\begin{equation}\label{8.11:eq-bound-members-2}
4\varrho_H d_1\le D_W ,
\end{equation}
for \eqref{8.11:eq-bound-members-2} gives $2d_1\le D_W/(2\varrho_H)$, hence every side of
$\Omega_{d_1}$ has length at least $D_W/(2\varrho_H)$, and the ratio of the longest to the
shortest side is at most $D_W\big/\bigl(D_W/(2\varrho_H)\bigr)=2\varrho_H=\varrho$.

We now derive \eqref{8.11:eq-bound-members-2} from the two bounds
\begin{equation*}
D_W\ge 8L\,\frac{M}{M_2}\,w,\qquad d_1\le 2w ,
\end{equation*}
together with the comparison \eqref{8.11:eq-bound-members-3}. The first is the lower bound on the
width $D_W$ of the zone fixed with the zone itself. For the
second, recall that
$d_1=y_++r_1$ is the free depth of Definition~\ref{8.11:def-released-problem}. By
Definition~\ref{8.11:def-young-row-cover}, see \eqref{8.11:eq-young-row-cover-a}, $y_+=w+5w/L$,
or $y_+=w$ when the young family $\mathfrak{Y}$ is empty; and by the condition $r_1\le w/12$
stated in Remark~\ref{8.11:rem-inherited-bounds},
\begin{equation*}
d_1=y_++r_1\le w\Bigl(1+\frac{5}{L}+\frac{1}{12}\Bigr)\le 2w,
\end{equation*}
where the last inequality is a one-sided lower bound on the blocking factor $L$ alone, absorbed in
the standing hypothesis $L\ge L_0$ (Lemma~\ref{2.4:lem-stage-zero-data}). The second bound gives
$4\varrho_H d_1\le 8\varrho_H w$, and the first bounds $D_W$ from below by $8L(M/M_2)w$. Where the
comparison \eqref{8.11:eq-bound-members-3} holds, multiplying it by $8w>0$ gives
$8\varrho_H w\le 8L(M/M_2)w$, and therefore
\begin{equation*}
4\varrho_H d_1\le 8\varrho_H w\le 8L\,\frac{M}{M_2}\,w\le D_W ,
\end{equation*}
which is \eqref{8.11:eq-bound-members-2}.
\end{proof}

\begin{theorem}[Curvature bound deep inside the hole]\label{8.11:thm-deep-curvature}
Assume the following statements, each under the one-sided smallness conditions with which it
is stated:
\begin{itemize}
\item the plaquette bounds supplied by the window functionals: the exterior, straddling and
young-cube bounds on $\|U(e\oplus x)(\partial p)-1\|$ for data $x\in K$;
\item the free-direction statement: if $x\in\operatorname{int}K$ and
$\partial\mathbf{A}/\partial x_c(x)=0$ for every data bond $c$ of a set $\mathfrak{F}$, then
$U(e\oplus x)$ is exact lattice Yang--Mills at every fine bond lying in no stencil of a data
bond $c\notin\mathfrak{F}$;
\item Lemma~\ref{8.11:lem-energy-below-support} (energy bounds below the support value);
\item the Lipschitz bound for the window functionals as functions of the data;
\item the lower bound for the second variation of $\mathbf{A}$ along geodesics of a common
gauge;
\item the relative gauge bound, and the comparison of plaquette variables of two data sets
presented in a common gauge;
\item the lattice calculus of the axial gauge: the bound on the non-tree bond variables by the
nonabelian Stokes formula, the bounds on the Yang--Mills and Bianchi defects, the identity
$dd^*+d^*d=-\Delta_{\mathbb{Z}^d}$ componentwise, and the mean-value inequality for
nonnegative subharmonic functions;
\item the size bounds on the zone width $D_W$ and on the block count $\hat N_W$;
\item the asymptotics of the Green's function of $\mathbb{Z}^d$, $d\ge3$, and the gradient
bound for discrete harmonic functions on a lattice ball;
\item the covering display \eqref{8.11:eq-young-row-cover-a} of
Definition~\ref{8.11:def-young-row-cover}.
\end{itemize}
The smallness conditions in question are those of the statement on the plaquette bounds
supplied by the window functionals and those under
which Lemma~\ref{8.11:lem-energy-below-support} and Lemma~\ref{8.11:lem-release-squeeze} are
stated; the smallness condition of Lemma~\ref{8.11:lem-energy-below-support} implies the
inequality
\begin{equation*}
4(D_W+8w)\,\varepsilon_a\eta\le \frac{1}{64C_{ax}} .
\end{equation*}
Let $e$ be an admissible isolated exterior with $S(e)\neq\emptyset$, let $\Lambda>0$, let
$\theta\in[1,2]$ be
such that $(\mathrm{VL})_{\theta}[\vartheta_{r_0(\Lambda)},\Lambda]$ holds
(Definition~\ref{8.11:def-value-lipschitz}, \eqref{8.11:eq-value-lipschitz-a}), with
$r_0(\Lambda)$ as in \eqref{8.11:eq-released-problem-a}, and let $\hat x$ be any minimiser of
$\mathbf{A}$ over $\widehat{\mathcal{W}}^{\theta}(e)$. Then for either choice of the reference
depth $c_*$ of Definition~\ref{8.11:not-bound-members}, every $h$ with $h\ge r_1$ and $hn\ge4$,
where $n=\eta^{-1}$ is the number of fine lattice spacings per unit of depth,
and every fine plaquette $p$ with $\delta(p)\ge d_1+h$,
\begin{equation}\label{8.11:eq-deep-curvature-a}
\|U(e\oplus\hat x)(\partial p)-1\|
\le\bigl[\mathfrak{B}(h)+\kappa_{\mathrm{col}}(h)+\vartheta_{\mathrm{nl}}(h)+2\vartheta_t\bigr]
\,\varepsilon_{k_o}\eta^2 ,
\end{equation}
where $r_1$, $d_1$ and $\vartheta_t$ are as in \eqref{8.11:eq-released-problem-a} and
$\mathfrak{B}(h)$, $\kappa_{\mathrm{col}}(h)$, $\vartheta_{\mathrm{nl}}(h)$ as in
\eqref{8.11:eq-bound-members-a}. The constants $C_\Pi,c_1,\dots,c_5$ depend on $d$ and
$\varrho_H$ only.
\end{theorem}

\begin{proof}
Throughout, $\varepsilon:=\varepsilon_{k_o}$ and $X\,[\mathrm{u}]:=X\,\varepsilon\eta^2$ as in
Definition~\ref{8.11:not-units-cap-amplitude}, and $\eta=n^{-1}$ is the fine spacing in the units
in which depths are measured, so that a displacement by $u$ fine lattice units changes the
depth by $u\eta$.

\smallskip\noindent
Step~0 (reduction to the released minimiser). Let $\check x$ be a minimiser of $\mathbf{A}$
over the released class $\mathfrak{L}(\hat x)$ of \eqref{8.11:eq-released-problem-a} (it exists
by part~(i) of Lemma~\ref{8.11:lem-release-squeeze}), and put
$\hat U:=U(e\oplus\hat x)$, $\check U:=U(e\oplus\check x)$. The hypotheses of
Lemma~\ref{8.11:lem-release-squeeze} hold: $\theta\in[1,2]$, $\hat x$ minimises $\mathbf{A}$
over $\widehat{\mathcal{W}}^{\theta}(e)$, $(\mathrm{VL})_{\theta}[\vartheta_r,\Lambda]$ holds
with $r=r_0(\Lambda)$, the antecedents of that lemma (the free-direction statement,
Lemma~\ref{8.11:lem-energy-below-support}, the Lipschitz bound for the window functionals, the
second-variation bound and the relative gauge bound) are among our hypotheses, and so are its
smallness conditions. By part~(ii) of that lemma,
$\|\hat U(\partial p)\check U(\partial p)^{-1}-1\|\le\vartheta_t\,[\mathrm{u}]$ for every fine
plaquette $p$. We write
\begin{equation*}
\hat U(\partial p)-1=\bigl(\hat U(\partial p)\check U(\partial p)^{-1}-1\bigr)\check U(\partial p)
+\bigl(\check U(\partial p)-1\bigr);
\end{equation*}
since $\check U(\partial p)$ is unitary, the triangle inequality gives
\begin{equation*}
\|\hat U(\partial p)-1\|\le\vartheta_t\,[\mathrm{u}]+\|\check U(\partial p)-1\|,
\end{equation*}
and in particular the same bound with $2\vartheta_t$ in place of $\vartheta_t$, which is the
form in which this term enters \eqref{8.11:eq-deep-curvature-a}.
It therefore suffices to prove, for $p$ as in the statement,
\begin{equation}\label{8.11:eq-deep-curvature-1}
\|\check U(\partial p)-1\|\le\bigl[\mathfrak{B}(h)+\kappa_{\mathrm{col}}(h)
+\vartheta_{\mathrm{nl}}(h)\bigr]\,[\mathrm{u}] .
\end{equation}
We collect the properties of $\check U$ that are used. By part~(i) of
Lemma~\ref{8.11:lem-release-squeeze}, $\check U$ is exact lattice Yang--Mills at every bond of
$\Omega_1:=\Omega_{d_1}$, $\check x\in\operatorname{int}K$ and
$\mathbf{A}(\check x)\le\mathbf{A}(w^{\sharp})$. Since $\check x\in K$ and
$\mathbf{A}(\check x)\le\mathbf{A}(w^{\sharp})$,
Lemma~\ref{8.11:lem-energy-below-support} applies with $x=\check x$: $\|\check U(\partial q)-1\|\le
\varepsilon_a\eta^2$ for every fine plaquette $q$, and $A_{\mathrm{in}}(\check U)\le
C_E\varepsilon_1^2\hat N_W$. Finally, with $\check\Phi(q):=(1/i)\log\check U(\partial q)$ and
$\theta''$, $\theta_y''$, $\sigma_e$ as in \eqref{8.11:eq-inherited-bounds-a},
Remark~\ref{8.11:rem-inherited-bounds} gives $\|\check\Phi(q)\|\le\theta''\,[\mathrm{u}]$ for
$q$ in $\{\delta\le w\}$, $\|\check\Phi(q)\|\le\theta_y''\,[\mathrm{u}]$ for $q$ in
$\{\delta\le y_+\}$, and $\|\check\Phi(q)\|\le\sigma_e\,[\mathrm{u}]$ for the plaquettes $q$ of
$\Sigma_{-(3w+1)}$. That remark states these bounds with the factor $1+\varepsilon_a\eta^2$,
which converts $\|\check U(\partial q)-1\|$ into $\|\check\Phi(q)\|$; this factor is absorbed
in the small relative errors $\vartheta$ of the bracket in \eqref{8.11:eq-deep-curvature-a},
and we write the bounds without it.

\smallskip\noindent
Step~1 (set-up). Put $\Omega_*:=\Omega_{c_*}$ and $\tau:=tn$, the thickness of the layer
between $\Omega_1$ and $\Omega_*$ in fine units. Here
\begin{equation*}
t=d_1-c_*=t_c+\varpi_y+r_1\ge t_c\ge1,
\end{equation*}
because $\varpi_y=y_+-w\ge0$ by \eqref{8.11:eq-young-row-cover-a} and $t_c\in\{4w+1,1\}$; in
particular $\tau=tn>0$, and since $\tau$ is a number of fine lattice units, $\tau\ge1$.
Let $p_0:=p$ and let $z_0$ be its lowest
corner. Since $\delta(p)\ge d_1+h$, all corners of $p_0$ are at $\ell^\infty$-distance at least
$h_z:=hn$ from $\partial\Omega_1$; in particular $h_z\ge4$, so that $p_0$ is a plaquette of
$\Omega_1$ none of whose bonds is a $\partial$-bond of $\Omega_1$. Below only the weaker
inequality $h_z\ge hn/2$ is used, which leaves room for shifts by one lattice unit.

We put $\check U$ in the complete axial gauge $\mathfrak{G}$ of $\Omega_*$ centred at $z_0$, as
in Lemma~\ref{8.11:lem-nonabelian-defects}. The hypotheses of that lemma hold: the bound
$\|\check U(\partial q)-1\|\le\varepsilon_a\eta^2$ holds for every fine plaquette (Step~0),
hence within fine distance $8$ of $\Omega_*$; $\check U$ is exact Yang--Mills at every bond of
$\Omega_1\subset\Omega_*$ (Step~0); and $\mathsf{L}:=\operatorname{diam}_1(\Omega_*)$ satisfies
$\mathsf{L}\le4(D_W+8w)n$, so that
\begin{equation*}
\mathsf{L}\,\varepsilon_a\eta^2\le4(D_W+8w)\,n\,\varepsilon_a\eta^2
=4(D_W+8w)\,\varepsilon_a\eta\le\frac{1}{64C_{ax}}
\end{equation*}
by the smallness condition displayed in the statement. Let $\Phi:=\check\Phi$ in the gauge
$\mathfrak{G}$, a $\mathfrak{g}$-valued $2$-cochain on $\Omega_*$; let $\rho_2$ and $N_1$ be the
homotopy defect and the Yang--Mills defect of $\Phi$ as in
Proposition~\ref{8.11:prop-trace-representation}; and let
$(\{\gamma_x\},\{S(b')\})$ be the path--surface system of
Lemma~\ref{8.11:lem-system-geometry}, with $S(b')=S^{\mathrm{lay}}(b')+S^{*}(b')$, where
$S^{\mathrm{lay}}(b')\subset\mathcal{L}$ and $S^{*}(b')\subset\partial\Omega_*$ (part~(a) of
that lemma). Write $\sum_{\mathfrak{b}}$ for the sum over $b'\in\mathfrak{b}(\Omega_1)$ and
$\rho:=\rho(p(b'))$ in such sums; the bonds of Case~I of
Lemma~\ref{8.11:lem-system-geometry} have $S^{\mathrm{lay}}(b')=N(b')$, those of Case~II
form $\mathfrak{b}_{II}$ and have $S^{\mathrm{lay}}(b')=\mathrm{Sq}(b')$.
Proposition~\ref{8.11:prop-trace-representation}, in the form
\eqref{8.11:eq-trace-representation-a}, gives
\begin{equation}\label{8.11:eq-deep-curvature-2}
\|\Phi(p_0)\|\le T_{\mathrm{lay}}+T_*+T_{\mathrm{nl}},
\end{equation}
where
\begin{equation*}
T_{\mathrm{lay}}:=C_\Pi\sum_{\mathfrak{b}}\rho^{-d}\sum_{q\in S^{\mathrm{lay}}(b')}\|\Phi(q)\|,
\qquad
T_*:=C_\Pi\sum_{\mathfrak{b}}\rho^{-d}\sum_{q\in S^{*}(b')}\|\Phi(q)\|,
\end{equation*}
and $T_{\mathrm{nl}}$ is the sum of the terms of \eqref{8.11:eq-trace-representation-a} that
carry $\rho_2$ or $N_1$. The norms $\|\Phi(q)\|$ are gauge invariant, since a gauge
transformation conjugates $\check U(\partial q)$, hence $\Phi(q)$, by a unitary; so the bounds
of Step~0 apply to them.

\smallskip\noindent
Step~2 (the layer outside the collar). The boxes $\Omega_c$ are concentric with equal offsets,
so a site at $\ell^\infty$-distance $u$ (in fine units) from $\Omega_1$ has
$\delta=d_1-u\eta$. By part~(c) of Lemma~\ref{8.11:lem-system-geometry}, the corner of a layer
plaquette of index $u$ nearest to $\Omega_1$ is at $\ell^\infty$-distance $u$ from $\Omega_1$;
hence such a plaquette lies in $\{\delta\le c\}$ if and only if $u\ge(d_1-c)n$. We split
$\mathcal{L}$ into three parts. The part $\mathcal{L}_c$ consists of the plaquettes with
$u\ge(r_1+\varpi_y)n$; since $d_1-r_1-\varpi_y=y_+-\varpi_y=w$, they lie in $\{\delta\le w\}$,
where $\|\Phi\|\le\theta''\,[\mathrm{u}]$. The part $\mathcal{L}_y$ consists of those with
$r_1n\le u<(r_1+\varpi_y)n$; since $d_1-r_1=y_+$, they lie in $\{\delta\le y_+\}$, where
$\|\Phi\|\le\theta_y''\,[\mathrm{u}]$. The collar $\mathcal{L}_{\mathrm{coll}}$ consists of
those with $u<r_1n$; it is treated in Step~3.

Bonds of Case~I. The indices of $N(b')$ run over $\{0,\dots,\tau-1\}$, and by part~(c)
of Lemma~\ref{8.11:lem-system-geometry} the number with $u_1\le u<u_2$ is $u_2-u_1$. Since
\begin{equation*}
t-r_1-\varpi_y=d_1-c_*-r_1-y_++w=w-c_*=t_c,
\end{equation*}
the surface $N(b')$ has $t_cn$ plaquettes in
$\mathcal{L}_c$ and $\varpi_yn$ in $\mathcal{L}_y$. By part~(e) of
Lemma~\ref{8.11:lem-system-geometry} and $h_z\ge hn/2$,
$\sum_{\mathfrak{b}}\rho^{-d}\le c_1/h_z\le2c_1/(hn)$; hence the contribution of
$\mathcal{L}_c\cup\mathcal{L}_y$ to the inner sums over $N(b')$, multiplied by $\rho^{-d}$ and
summed, is at most
\begin{equation*}
\frac{2c_1}{hn}\,\bigl[t_cn\,\theta''+\varpi_yn\,\theta_y''\bigr]\,[\mathrm{u}]
=\frac{2c_1\,[t_c\theta''+\varpi_y\theta_y'']}{h}\,[\mathrm{u}] .
\end{equation*}
Bonds of Case~II. By part~(c) of Lemma~\ref{8.11:lem-system-geometry}, $\mathrm{Sq}(b')$
has at most $(\tau+1)\tau\le(t+1)^2n^2$ plaquettes and, taking
$u_1=r_1n$, $u_2=(r_1+\varpi_y)n$, at most
\begin{equation*}
\varpi_yn\,\bigl(2(r_1+\varpi_y)n+1\bigr)\le\varpi_y(2\varpi_y+2r_1+1)\,n^2
\end{equation*}
of them lie in $\mathcal{L}_y$. We bound the plaquettes of $\mathcal{L}_c$ by $\theta''$ and
those of $\mathcal{L}_y$ by $\theta_y''$. By part~(e),
$\sum_{\mathfrak{b}_{II}}\rho^{-d}\le c_2/h_z^2\le4c_2/(hn)^2$, so this contribution is at
most
\begin{equation*}
\frac{4c_2\,[(t+1)^2\theta''+\varpi_y(2\varpi_y+2r_1+1)\theta_y'']}{h^2}\,[\mathrm{u}] .
\end{equation*}
Multiplying
by $C_\Pi$, the part of $T_{\mathrm{lay}}$ coming from $\mathcal{L}_c\cup\mathcal{L}_y$ is at
most
\begin{equation}\label{8.11:eq-deep-curvature-3}
C_\Pi\Bigl\{\frac{2c_1[t_c\theta''+\varpi_y\theta_y'']}{h}
+\frac{4c_2[(t+1)^2\theta''+\varpi_y(2\varpi_y+2r_1+1)\theta_y'']}{h^2}\Bigr\}\,[\mathrm{u}] .
\end{equation}

\smallskip\noindent
Step~3 (the collar, by energy). Consider the pairs $(b',q)$ with
$q\in S^{\mathrm{lay}}(b')\cap\mathcal{L}_{\mathrm{coll}}$. By part~(c) of
Lemma~\ref{8.11:lem-system-geometry}, for a bond of Case~I there are $r_1n$ such $q$
(indices $0\le u<r_1n$), and for a bond of $\mathfrak{b}_{II}$ at most $(r_1n+1)r_1n$. By part~(b)
of the same lemma, each plaquette of $\mathcal{L}$ lies in $S^{\mathrm{lay}}(b')$ for at most
three bonds $b'$. The Cauchy--Schwarz inequality over the pairs therefore gives
\begin{equation*}
\sum_{(b',q)}\rho^{-d}\|\Phi(q)\|
\le\Bigl[\sum_{\mathfrak{b}}\rho^{-2d}\,r_1n
+\sum_{\mathfrak{b}_{II}}\rho^{-2d}(r_1n+1)\,r_1n\Bigr]^{1/2}
\Bigl[3\sum_{q\in\mathcal{L}_{\mathrm{coll}}}\|\Phi(q)\|^2\Bigr]^{1/2}.
\end{equation*}
First bracket. Part~(e) of Lemma~\ref{8.11:lem-system-geometry} with $d=4$ gives
$\sum_{\mathfrak{b}}\rho^{-8}\le c_1h_z^{-5}$ and $\sum_{\mathfrak{b}_{II}}\rho^{-8}\le
c_2h_z^{-6}$. With $h_z\ge hn/2$ and $n^{-1}=\eta$,
\begin{equation*}
\sum_{\mathfrak{b}}\rho^{-8}r_1n\le32c_1r_1h^{-5}n^{-4},\qquad
\sum_{\mathfrak{b}_{II}}\rho^{-8}(r_1n+1)r_1n\le64c_2(r_1+\eta)\,r_1h^{-6}n^{-4}.
\end{equation*}
Since $r_1\le h$ and $\eta\le h/4$ (the latter is $hn\ge4$), we have
$(r_1+\eta)r_1h^{-6}\le(1+\tfrac14)r_1h^{-5}$, so the first bracket is at most
\begin{equation*}
(32c_1+80c_2)\,r_1h^{-5}\eta^4\le128(c_1+c_2)(r_1+1)\,h^{-5}\eta^4 .
\end{equation*}
Second bracket. For every plaquette $q$, $\|\Phi(q)\|\le2\|\check U(\partial q)-1\|$
(Definition~\ref{8.11:not-units-cap-amplitude}), and by the comparison of the plaquette variable
with the plaquette energy, $\|\check U(\partial q)-1\|^2\le2N\,\mathfrak{e}(\check U(\partial q))$;
thus
\begin{equation*}
\|\Phi(q)\|^2\le4\|1-\check U(\partial q)\|^2\le8N\,\mathfrak{e}(\check U(\partial q)).
\end{equation*}
The collar lies in the zone $H$, and $H$ lies in the region whose plaquette energies make up
$A_{\mathrm{in}}$. Hence, by Step~0,
\begin{equation*}
3\sum_{q\in\mathcal{L}_{\mathrm{coll}}}\|\Phi(q)\|^2\le24N\,A_{\mathrm{in}}(\check U)
\le24N\,C_E\varepsilon_1^2\hat N_W .
\end{equation*}
Taking the square roots of the two bounds, multiplying them, multiplying by $C_\Pi$, and using
$128\cdot24=3072$, we find that the collar part of $T_{\mathrm{lay}}$ is at most
\begin{equation*}
C_\Pi\bigl[3072(c_1+c_2)\,N C_E\hat N_W(r_1+1)\,h^{-5}\bigr]^{1/2}\varepsilon_1\eta^2
=\kappa_{\mathrm{col}}(h)\,[\mathrm{u}],
\end{equation*}
by the definition of $\kappa_{\mathrm{col}}(h)$ in \eqref{8.11:eq-bound-members-a}, which
carries the factor $\varepsilon_1/\varepsilon$.

\smallskip\noindent
Step~4 (the reference surface: the logarithm). Every plaquette of $S^{*}(b')$ lies in
$\partial\Omega_*=\Sigma_{c_*}$, where $\|\Phi\|\le\sigma_*\,[\mathrm{u}]$: for
$c_*=-(3w+1)$ this is the bound by $\sigma_e=\sigma_*$ of Step~0; for $c_*=w-1$ it is the bound
by $\theta''=\sigma_*$, since $\Sigma_{w-1}\subset\{\delta\le w\}$. Hence
$T_*\le C_\Pi\sigma_*\sum_{\mathfrak{b}}\rho^{-d}|S^{*}(b')|\,[\mathrm{u}]$, and we bound the
last sum using parts~(d) and~(e) of Lemma~\ref{8.11:lem-system-geometry}, with $\ell$ and
$\ell_{\max}$ as there. We use $\ell_{\max}\le D_Wn$, $h_z\ge hn/2$, $\tau=tn$ and $\tau\ge1$
(Step~1), which give
\begin{equation*}
\frac{\tau}{h_z}\le\frac{2t}{h},\qquad\frac{2\tau+1}{h_z}\le\frac{6t}{h},\qquad
\frac{(2\tau+1)\tau}{h_z^2}\le12\Bigl(\frac{t}{h}\Bigr)^2,\qquad
\frac{\ell_{\max}}{h_z}\le\frac{2D_W}{h}.
\end{equation*}
For $b'\in\mathfrak{b}_{\mathrm{near}}$, part~(d) gives $|S^{*}(b')|\le2d\rho+2\tau$ for bonds
of Case~I and $|S^{*}(b')|\le(2\tau+1)(2d\rho+4\tau)$ for bonds of Case~II.
With the sums of part~(e),
\begin{align*}
\sum_{\mathfrak{b}_{\mathrm{near}}}\rho^{-d}|S^{*}(b')|
&\le\sum_{\mathfrak{b}}\rho^{-d}(2d\rho+2\tau)
+\sum_{\mathfrak{b}_{II}}\rho^{-d}(2\tau+1)(2d\rho+4\tau)\\
&\le2dc_1\Bigl(1+\log^+\frac{\ell_{\max}}{h_z}\Bigr)+\frac{2c_1\tau}{h_z}
+\frac{2dc_2(2\tau+1)}{h_z}+\frac{4c_2(2\tau+1)\tau}{h_z^2}\\
&\le2dc_1\Bigl(1+\log^+\frac{2D_W}{h}\Bigr)+\frac{4c_1t}{h}+\frac{12dc_2t}{h}
+48c_2\Bigl(\frac{t}{h}\Bigr)^2 .
\end{align*}
For $b'\in\mathfrak{b}_{\mathrm{far}}$, part~(d) gives $\rho\ge\ell/2$ and
$\sum_{\mathfrak{b}_{\mathrm{far}}}|S^{*}(b')|\le c_d(\ell_{\max}+2\tau)^d$, whence
\begin{equation*}
\sum_{\mathfrak{b}_{\mathrm{far}}}\rho^{-d}|S^{*}(b')|
\le(\ell/2)^{-d}c_d(\ell_{\max}+2\tau)^d\le2^dc_d(2\varrho_H+2)^d .
\end{equation*}
The last inequality amounts to $\ell_{\max}+2\tau\le(2\varrho_H+2)\,\ell$; here
$\ell_{\max}\le2\varrho_H\,\ell$, since $2\varrho_H$ bounds the aspect ratio of $\Omega_{d_1}$
(Definition~\ref{8.11:not-bound-members}), and the remaining part of the inequality is
$\tau\le\ell$.
Since $\varrho_H\ge1$, being an upper bound for an aspect ratio, we have
$\log2\le\log4\varrho_H$; with $\log^+(2D_W/h)\le\log2+\log^+(D_W/h)$, the definitions
$c_3=2dc_1(1+\log4\varrho_H)+2^dc_d(2\varrho_H+2)^d$, $c_4=4c_1+12dc_2$, $c_5=48c_2$ of
Definition~\ref{8.11:not-bound-members} give
\begin{equation}\label{8.11:eq-deep-curvature-4}
\sum_{\mathfrak{b}}\rho^{-d}|S^{*}(b')|
\le c_3+2dc_1\log^+\frac{D_W}{h}+c_4\frac{t}{h}+c_5\Bigl(\frac{t}{h}\Bigr)^2,
\end{equation}
and hence
\begin{equation*}
T_*\le C_\Pi\sigma_*\Bigl[c_3+2dc_1\log^+\frac{D_W}{h}+c_4\frac{t}{h}
+c_5\Bigl(\frac{t}{h}\Bigr)^2\Bigr]\,[\mathrm{u}] .
\end{equation*}
Together with \eqref{8.11:eq-deep-curvature-3}, this is $\mathfrak{B}(h)\,[\mathrm{u}]$ by the
definition of $\mathfrak{B}(h)$ in \eqref{8.11:eq-bound-members-a}.

\smallskip\noindent
Step~5 (defects). By Lemma~\ref{8.11:lem-nonabelian-defects}, whose hypotheses were checked
in Step~1, and since the operator norm is at most the Hilbert--Schmidt norm, we have
$\|\rho_2(q)\|\le C_{ND}N(|x_q-z_0|_1+1)^2\varepsilon_a^2\eta^4$, where $x_q$ is the lowest
corner of $q$. For $q\subset\Omega_*$,
$|x_q-z_0|_1\le\mathsf{L}$, so $\|\rho_2\|\le C_{ND}N(\mathsf{L}+1)^2\varepsilon_a^2\eta^4$ on
$\Omega_*$; for a cell $c$ with lowest corner $x_c$ and $\rho(c)=1+|z_0-x_c|_\infty$ we have
$|x_c-z_0|_1+1\le d\,\rho(c)$, so $\|\rho_2\|\le C_{ND}N(d\rho(c))^2\varepsilon_a^2\eta^4$
there; and $\|N_1(b)\|\le C_{ND}N\,d\rho(b)\,\varepsilon_a^2\eta^4$ at every bond $b$ of
$\Omega_1$. The four kinds of terms
of $T_{\mathrm{nl}}$ are bounded as follows.

(i) The $\rho_2$-part of the first sum of \eqref{8.11:eq-trace-representation-a} is at most
$C_\Pi\sum_{\mathfrak{b}}\rho^{-d}|S(b')|\cdot\sup_{\Omega_*}\|\rho_2\|$. Here
$|S(b')|=|S^{\mathrm{lay}}(b')|+|S^{*}(b')|$, with $|N(b')|=\tau$ and
$|\mathrm{Sq}(b')|\le(\tau+1)\tau$ by part~(c) of Lemma~\ref{8.11:lem-system-geometry}; as in
Step~2, $\sum_{\mathfrak{b}}\rho^{-d}\tau\le2c_1t/h$ and
$\sum_{\mathfrak{b}_{II}}\rho^{-d}(\tau+1)\tau\le4c_2(t+1)^2/h^2$. With
\eqref{8.11:eq-deep-curvature-4}, this term is at most
\begin{align*}
&C_\Pi\Bigl[\frac{2c_1t}{h}+\frac{4c_2(t+1)^2}{h^2}+c_3+2dc_1\log^+\frac{D_W}{h}
+c_4\frac{t}{h}+c_5\Bigl(\frac{t}{h}\Bigr)^2\Bigr]\\
&\qquad{}\times C_{ND}N(\mathsf{L}+1)^2\varepsilon_a^2\eta^4 .
\end{align*}
(ii) There are at most $c_d\rho'^{\,d-1}$ plaquettes of $\Omega_1$ at parameter $\rho'$,
and $\rho'\le\mathsf{L}$; since $d=4$,
\begin{equation*}
C_\Pi\sum_{q\subset\Omega_1}\rho^{-d}C_{ND}N d^2\rho^2\varepsilon_a^2\eta^4
\le C_\Pi C_{ND}c_dd^2N\sum_{\rho'\le\mathsf{L}}\rho'\,\varepsilon_a^2\eta^4
\le C_\Pi C_{ND}c_dd^2N\mathsf{L}^2\varepsilon_a^2\eta^4 .
\end{equation*}
(iii) At $p_0$, whose lowest corner is $z_0$,
$\|\rho_2(p_0)\|\le C_{ND}N\varepsilon_a^2\eta^4$.

(iv) In the same way as in~(ii),
\begin{equation*}
C_\Pi'\sum_{b\subset\Omega_1}\rho^{1-d}C_{ND}N d\rho\,\varepsilon_a^2\eta^4
\le C_\Pi'C_{ND}c_ddN\mathsf{L}^2\varepsilon_a^2\eta^4 .
\end{equation*}
For every $t\ge0$ we have $t/h\le1+(t/h)^2$; and since $t\ge1$ (Step~1),
$(t+1)^2/h^2\le4(t/h)^2$. Also $1\le\mathsf{L}$, so $(\mathsf{L}+1)^2\le4\mathsf{L}^2$ and
$1\le\mathsf{L}^2$. Adding (i)--(iv), there is a constant $C'$ depending on $d$ and
$\varrho_H$ only (through $C_\Pi$, $C_\Pi'$, $C_{ND}$, $c_1,\dots,c_5$, $c_d$) with
\begin{equation*}
T_{\mathrm{nl}}\le C'N\mathsf{L}^2\varepsilon_a^2\eta^4
\Bigl[1+\log^+\frac{D_W}{h}+\Bigl(\frac{t}{h}\Bigr)^2\Bigr].
\end{equation*}
Since $\mathsf{L}\le4(D_W+8w)n$ and $n\eta=1$, we have $\mathsf{L}^2\eta^4\le16(D_W+8w)^2\eta^2$,
and therefore, taking the constant $C_{\mathrm{nl}}=C(d,\varrho_H)$ of
\eqref{8.11:eq-bound-members-a} to be $16C'$,
\begin{equation*}
T_{\mathrm{nl}}\le C_{\mathrm{nl}}N(D_W+8w)^2\frac{\varepsilon_a^2}{\varepsilon}
\Bigl[1+\log^+\frac{D_W}{h}+\Bigl(\frac{t}{h}\Bigr)^2\Bigr]\varepsilon\eta^2
=\vartheta_{\mathrm{nl}}(h)\,[\mathrm{u}] .
\end{equation*}

\smallskip\noindent
Step~6 (conclusion). By Definition~\ref{8.11:not-units-cap-amplitude},
$\|\check U(\partial p_0)-1\|\le\|\Phi(p_0)\|$. By \eqref{8.11:eq-deep-curvature-2}, Steps~2
and~3 (which together bound $T_{\mathrm{lay}}$), Step~4 and Step~5,
\begin{equation*}
\|\Phi(p_0)\|\le\bigl[\mathfrak{B}(h)+\kappa_{\mathrm{col}}(h)
+\vartheta_{\mathrm{nl}}(h)\bigr]\,[\mathrm{u}],
\end{equation*}
which is \eqref{8.11:eq-deep-curvature-1}. By Step~0 this gives
\eqref{8.11:eq-deep-curvature-a}. The constants $C_\Pi$ and $C_\Pi'$ are those of
Proposition~\ref{8.11:prop-trace-representation} at $(d,2\varrho_H)$, where $2\varrho_H$
bounds the aspect ratio of $\Omega_{d_1}$,
$c_1,c_2$ those of Lemma~\ref{8.11:lem-system-geometry}, and $c_3,c_4,c_5$ are built from them
and from $c_d$ and $\varrho_H$; so they depend on $d$ and $\varrho_H$ only.
\end{proof}

\begin{remark}[Scope of the deep curvature bound]\label{8.11:rem-deep-bound-scope}
We use the notation of Theorem~\ref{8.11:thm-deep-curvature} and of Notation~\ref{8.11:not-bound-members}.
The right-hand side of \eqref{8.11:eq-deep-curvature-a} is the sum of the three members
\eqref{8.11:eq-bound-members-a} and $2\vartheta_t$, times the exterior unit $\varepsilon_{k_o}\eta^2$.
\begin{enumerate}
\item[(i)] Two things are used about the minimiser $\hat x$. The first is its minimality, in the form
$\hat m(\theta)=\mathbf{A}(\hat x)$ and in the form $\mathbf{A}(\hat x)\le\mathbf{A}(w^{\sharp})$.
The second is its feasibility, the feasibility hypothesis among the antecedents of
Theorem~\ref{8.11:thm-deep-curvature}.
Nothing else is used. Consequently the validity of
\eqref{8.11:eq-deep-curvature-a} for every minimiser of $\mathbf{A}$ over
$\widehat{\mathcal{W}}^{\theta}(e)$ is automatic.
\item[(ii)] No property of the forces on the layer $\mathcal{L}$ enters the bound. This covers
their size, their sign, their support inside the layer, their divergence and any further structure.
The layer enters only through the norms $\|\Phi(q)\|$ of the logarithms of the plaquette variables
at the layer plaquettes $q$.
Where the bounds $\theta''$, $\theta_y''$ satisfied by $\check x$ apply, these enter through supremum
bounds. On the remaining part of the layer they enter through an $\ell^2$ bound.
\item[(iii)] The bound \eqref{8.11:eq-deep-curvature-a} holds at every fine plaquette $p$ of
$H=R^{+}_{k_o}$ with $\delta(p)\ge d_1+h$, including those in the inner rings and in the core of the
zone, the innermost parts of $H$, where $\check x$ is free. No separate analysis of the minimiser
in the inner rings or in the core is needed.
The bound improves with the depth $\delta(p)$.
\item[(iv)] The constants $C_\Pi,c_1,\dots,c_5$ of $\mathfrak{B}(h)$ do not depend on
$N$. An explicit factor $N$ enters only through $\kappa_{\mathrm{col}}(h)$ and
$\vartheta_{\mathrm{nl}}(h)$.
\item[(v)] The two reference depths behave differently.
  \begin{itemize}
  \item For $c_*=-(3w+1)$ the amplitude $\sigma_*$ multiplying the logarithm in $\mathfrak{B}(h)$ is
  $s_e$, and there
  is an additive member $(4w/h)\theta''$. This is the form for heights $h$ of order at least $w$.
  \item For $c_*=w-1$ the amplitude $\sigma_*$ multiplying $\log(D_W/h)$ is $\theta''$, and there is
  no member of
  order $w/h$. This is the form for small heights $r_1\le h\ll w$, used in the small-height
  corollary.
  \end{itemize}
\item[(vi)] The two members $\kappa_{\mathrm{col}}(h)$ and $\vartheta_{\mathrm{nl}}(h)$ are kept
small by the ceilings
\begin{equation}\label{8.11:eq-deep-bound-scope-a}
\begin{gathered}
\kappa_{\mathrm{col}}(h)\le 1,\\
\vartheta_{\mathrm{nl}}(h)\le\vartheta_1 ,
\end{gathered}
\end{equation}
the first called the collar ceiling and the second the non-linear ceiling. Fix a number $\xi>0$
and put $h=\xi w$. Then the collar ceiling is implied by
\begin{equation}\label{8.11:eq-deep-bound-scope-1}
C(d,N,L,M/M_2,c_0,\beta_0)\,\xi^{-5}\,\frac{r_1+1}{w}\le 1 ,
\end{equation}
for a constant depending only on $d,\varrho_H,N,L,M/M_2,c_0,\beta_0,\hat r$ and the constants
$C_B,C_E,C_{LF}$, written $C(d,N,L,M/M_2,c_0,\beta_0)$ and made explicit in the proof.
This condition is one-sided and is a condition on $g$, chosen after $L$ and $M$. It is uniform in
the level, and it holds for $g$ small. The set of $g^{-1}$ for which it holds contains a
half-line, so the condition is non-empty. It bounds no size, level, age, nor $M$, $M/M_2$ or $L$
from above. The non-linear ceiling bounds a polynomial in the size of the hole times the small
factor $\varepsilon_a^2/\varepsilon$ by a constant. It is of the same kind as the smallness
condition among the hypotheses of Theorem~\ref{8.11:thm-deep-curvature} that contains the gauge
condition $4(D_W+8w)\varepsilon_a\eta\le1/(64C_{ax})$, and as its analogue for the block-mean
model.
\end{enumerate}
\end{remark}

\begin{proof}
We take the items in order. Each is a reading of the statement of
Theorem~\ref{8.11:thm-deep-curvature}, of its proof, or of the definitions
\eqref{8.11:eq-bound-members-a}.

\emph{Item} (i). In the proof of Theorem~\ref{8.11:thm-deep-curvature} the minimiser $\hat x$
appears in two places. It appears in the reduction of \eqref{8.11:eq-deep-curvature-a} to a bound
on $\check U=U(e\oplus\check x)$ for the released minimiser $\check x$, which accounts for the
member $2\vartheta_t$. It also appears in the comparison of actions with the support point
$w^{\sharp}$. Of $\hat x$, these use only three facts:
\begin{itemize}
\item the identity $\hat m(\theta)=\mathbf{A}(\hat x)$, used in the reduction;
\item the inequality $\mathbf{A}(\hat x)\le\mathbf{A}(w^{\sharp})$;
\item its feasibility, the feasibility hypothesis of Theorem~\ref{8.11:thm-deep-curvature}.
\end{itemize}
The first two are minimality of $\hat x$ over $\widehat{\mathcal{W}}^{\theta}(e)$, and the third is
a hypothesis of the theorem. None of them singles out a particular minimiser. Hence the conclusion
holds for every minimiser.

\emph{Item} (ii). In the proof, the layer contributes to $\|\Phi(p)\|$ only through the sums
$\sum\rho^{-d}\|\Phi(q)\|$ over layer plaquettes $q$. The layer is split into three parts.
\begin{itemize}
\item On $\mathcal{L}_c$ the bound $\|\Phi(q)\|\le\theta''[\mathrm{u}]$ is used.
\item On $\mathcal{L}_y$ the bound $\|\Phi(q)\|\le\theta_y''[\mathrm{u}]$ is used.
\item On the collar $\mathcal{L}_{\mathrm{coll}}$, Cauchy--Schwarz reduces the contribution to
$\sum_{q\in\mathcal{L}_{\mathrm{coll}}}\|\Phi(q)\|^2$. This sum is bounded by $8N$ times the interior
action $A_{\mathrm{in}}(\check U)$ of the zone, so that three times it, as it enters the
Cauchy--Schwarz bound, is at most $24N\,A_{\mathrm{in}}(\check U)$.
\end{itemize}
The first two are supremum bounds and the third is an $\ell^2$ bound. The forces on the layer are
used at no point.

\emph{Item} (iii). The hypothesis of Theorem~\ref{8.11:thm-deep-curvature} on $p$ is only that $p$
is a fine plaquette with $\delta(p)\ge d_1+h$. Of $\check U$ the proof uses only the following
inputs: that $\check U$ is exact Yang--Mills at every bond of $\Omega_{d_1}$, which contains the
inner rings and the core; the bound $\|\check U(\partial q)-1\|\le\varepsilon_a\eta^2$ for every
$q$; the interior energy bound $A_{\mathrm{in}}(\check U)\le C_E\varepsilon_1^2\hat N_W$; and the
bounds on $\|\Phi\|$ that $\check x$ satisfies on the layer and on the reference surface. None of
these inputs distinguishes a plaquette of the inner rings or of the core from any other fine plaquette $p$ with
$\delta(p)\ge d_1+h$. A fine plaquette there is treated exactly as any
other.

For the improvement with depth, let $p$ be given with $\delta(p)-d_1\ge r_1$ and
$(\delta(p)-d_1)n\ge 4$. Then every $h$ with $r_1\le h\le\delta(p)-d_1$ and $hn\ge4$ is admissible.
Each member of \eqref{8.11:eq-bound-members-a} is non-increasing in $h$, as follows.
\begin{itemize}
\item In $\mathfrak{B}(h)$, $h$ enters through $1/h$, $1/h^2$, $\log^{+}(D_W/h)$, $t/h$ and
$(t/h)^2$.
\item $\kappa_{\mathrm{col}}(h)$ is proportional to $h^{-5/2}$.
\item In $\vartheta_{\mathrm{nl}}(h)$, $h$ enters through $\log^{+}(D_W/h)$ and $(t/h)^2$.
\end{itemize}
So the largest admissible $h$, which grows with $\delta(p)$, gives the smallest bound.

\emph{Item} (iv). By the statement of Theorem~\ref{8.11:thm-deep-curvature}, $C_\Pi$ and
$c_1,\dots,c_5$ depend on $d$ and $\varrho_H$ only. In the proof the norms are operator norms and
the kernels of $\Pi^0$ are scalar, so no explicit factor $N$ arises in the constants of
$\mathfrak{B}(h)$. An explicit factor $N$ enters in two places:
\begin{itemize}
\item in the collar estimate, through
$\|\Phi\|^2\le4\|1-U\|^2\le 8N\,\mathfrak{e}(U(\partial q))$, which produces the factor $N$ in
$\kappa_{\mathrm{col}}(h)$;
\item in the defect bounds on $\rho_2$ and $N_1$, which carry $C_{ND}N$ and produce the factor $N$
in $\vartheta_{\mathrm{nl}}(h)$.
\end{itemize}

\emph{Item} (v). By definition $t_c=w-c_*$ and $t=d_1-c_*$. For $c_*=w-1$ we have $t_c=1$.
Then the member $2c_1t_c\theta''/h$ of $\mathfrak{B}(h)$ is $2c_1\theta''/h$, which contains no
factor $w$. The plaquettes of the reference surface $\Sigma_{w-1}$ lie in
$\{\delta\le w\}$. There $\theta''$ bounds $\|\Phi\|$, so $\sigma_*=\theta''$ may be taken, and
$\theta''$ is the coefficient of the logarithm. For $c_*=-(3w+1)$ we have $t_c=4w+1$, so the
member $2c_1t_c\theta''/h$ of $\mathfrak{B}(h)$ is the additive member of order $(4w/h)\theta''$
stated in~(v). The coefficient $\sigma_*$ of the logarithm is the exterior amplitude $s_e$. Only
for $h$ of order at least $w$ is that additive member bounded by a multiple of $\theta''$.

\emph{Item} (vi). By the size bounds of the zone, with $d=4$,
\begin{equation*}
\hat N_W\le C_B\bigl[(c_0+1)L^4+3^4\bigr]D_W^4 ,
\qquad D_W\le 17LM\check R\hat r .
\end{equation*}
Since $w=M_2R_{k_o}$, dividing by $w^4$ and by $w$ respectively gives
\begin{equation*}
\frac{\hat N_W}{w^4}\le C_B\bigl[(c_0+1)L^4+3^4\bigr]\Bigl(\frac{D_W}{w}\Bigr)^4,
\qquad \frac{D_W}{w}\le 17L\,\frac{M}{M_2}\,\frac{\check R}{R_{k_o}}\,\hat r .
\end{equation*}
Here $\check R/R_{k_o}\in\{1,L\}$, so $D_W/w$ does not depend on the cutoff.

Two further inputs are needed. The flow bounds give
$\varepsilon_1/\varepsilon\le C_{LF}(1+\beta_0)^{c_0+1}$, and $C_E$ is the energy constant of the
zone, the constant in the bound $A_{\mathrm{in}}(\check U)\le C_E\varepsilon_1^2\hat N_W$ used in
the collar estimate.

Now put $h=\xi w$. Then $\hat N_W(r_1+1)/h^5=\xi^{-5}(\hat N_W/w^4)(r_1+1)/w$. Squaring
$\kappa_{\mathrm{col}}(h)$ therefore gives
\begin{equation}\label{8.11:eq-deep-bound-scope-2}
\begin{split}
\kappa_{\mathrm{col}}(\xi w)^2
&\le C_\Pi^2\,3072(c_1+c_2)\,N\,C_E\,C_B\bigl[(c_0+1)L^4+3^4\bigr]\\
&\quad\times\Bigl(17L\,\frac{M}{M_2}\,\frac{\check R}{R_{k_o}}\,\hat r\Bigr)^4
C_{LF}^2(1+\beta_0)^{2(c_0+1)}\;\xi^{-5}\,\frac{r_1+1}{w}.
\end{split}
\end{equation}
Denote by $C(d,N,L,M/M_2,c_0,\beta_0)$ the product of the factors preceding $\xi^{-5}$ in
\eqref{8.11:eq-deep-bound-scope-2}, with $\check R/R_{k_o}$ replaced by its larger possible value
$L$; besides the displayed arguments it contains $\varrho_H$ (through $C_\Pi$, $c_1$, $c_2$),
$\hat r$ and the constants $C_B$, $C_E$, $C_{LF}$. Then \eqref{8.11:eq-deep-bound-scope-1} gives
$\kappa_{\mathrm{col}}(\xi w)^2\le1$, which is the collar ceiling.

The condition \eqref{8.11:eq-deep-bound-scope-1} is an upper bound on $(r_1+1)/w$ only. Hence it is
one-sided and involves $g$ after $L$ and $M$ have been fixed. By the lower flow bound,
\begin{equation*}
\frac{r_1+1}{w}\le C(L,M)\,\frac{\log g^{-1}}{(\log g^{-2})^{r}}
\end{equation*}
uniformly in the level, where $r\ge2$ is the exponent in Ba{\l}aban's condition $p_0\ge5r$ on the
exponent $p_0$ of the degree function $p_0(\cdot)$ \cite{B-CMP119}. The right-hand side
tends to $0$ as $g\to0$. So \eqref{8.11:eq-deep-bound-scope-1} holds for all $g$ small enough, and
the set of admissible $g^{-1}$ contains a half-line. Since the condition is an upper bound on a
quantity that tends to zero with $g$, it bounds no size, level or age, and none of $M$, $M/M_2$
or $L$, from above.

Finally, by \eqref{8.11:eq-bound-members-a},
\begin{equation*}
\vartheta_{\mathrm{nl}}(h)=C_{\mathrm{nl}}N(D_W+8w)^2\,\frac{\varepsilon_a^2}{\varepsilon}
\bigl[1+\log^{+}(D_W/h)+(t/h)^2\bigr].
\end{equation*}
This is a polynomial in the size of the hole times $\varepsilon_a^2/\varepsilon$. The non-linear
ceiling bounds it by the constant $\vartheta_1$, which is the form stated in~(vi).
\end{proof}

Throughout, the notation is that of Notation~\ref{8.11:not-hole-exterior-windows}. There, $e$ is the exterior and $x$ a hole datum. The map $x\mapsto\sigma(x)$ assigns to $x$ its top averages. For a straddling site $a\in\mathfrak{S}$ the quantity $g_a(x)=f_a(e\oplus x)$ is the value of its window, measured in units $\varepsilon_{k_o}\eta^2$. The sets $\tilde a$ and $a^{\sim 4}$ are the test region and the dependence region of the site $a$. Finally, $\delta$ is the signed depth and $w=M_2R_{k_o}\ge 2$ is the clamp width.

A plaquette-by-plaquette bound in the clamped layer can be obtained by comparing the minimiser $U(e\oplus x)$ of the full problem with the straddling local minimisers. The full problem is constrained by the level-$(k_o-1)$ data of the outer ring, whereas a local minimiser reads only the top averages $\sigma(x)$. That comparison is justified at points $x$ that are critical along the fibres of $\sigma$. At the constrained minimiser $\hat x$, however, the fibres through $\{\delta\le w\}$ are not free near a tight young cube.

Where the clamped layer enters an estimate only through fluxes against a kernel that is flat at block scale, a weaker bound suffices. Block-averaged fluxes are the plaquette variables of the top data, and the straddling windows bound these directly, with no question about the fibres. The following lemma records this bound.

\begin{lemma}[Block curvature bound in the clamped layer]\label{8.11:lem-clamped-block}
Assume the following.
\begin{itemize}
\item The exterior $e$ is admissible.
\item $\theta$ is the threshold of the constrained class of
Notation~\ref{8.11:not-hole-exterior-windows}.
\item $x$ lies in the set of hole data of Notation~\ref{8.11:not-hole-exterior-windows}, and
$g_a(x)\le\theta$ for every $a\in\mathfrak{S}$.
\end{itemize}
Let $P$ be a plaquette of the block lattice of level $k_o$ (a top block plaquette) whose four blocks lie in $\{\delta\le w\}$. Then the top configuration $e\oplus\sigma(x)$ satisfies
\begin{equation}\label{8.11:eq-clamped-block-a}
\bigl\|(e\oplus\sigma(x))(\partial P)-1\bigr\|<2\theta\varepsilon_{k_o}.
\end{equation}
\end{lemma}

\begin{proof}
Write $V:=e\oplus\sigma(x)$ for the top configuration.

\emph{Step 1: the four blocks of $P$ lie in one test region.} By Notation~\ref{8.11:not-hole-exterior-windows}, the test regions $\tilde a$, with $a$ pure or straddling, cover $\{\delta\le w\}$ with overlaps of width at least $w\ge 2$ blocks. Consequently every two-by-two array of blocks contained in $\{\delta\le w\}$ lies in a single test region $\tilde a$. The four blocks of $P$ lie in $\{\delta\le w\}$ and form such an array. Hence they lie in a single test region $\tilde a$, where $a$ is either pure or straddling. Fix such an $a$.

\emph{Step 2: the local minimiser at $a$.} Let
\begin{equation*}
U_a:=U_{k_o,a}\bigl(V|_{a^{\sim 4}}\bigr)
\end{equation*}
be the local minimiser of level $k_o$ at the site $a$ with data $V$ on $a^{\sim 4}$, in the sense of \cite[(1.2), (2.17)]{B-CMP119}.

First, $U_a$ is $\theta\varepsilon_{k_o}$-regular on $\tilde a$. This is what the bound $f_a(e\oplus x)\le\theta$ on the window of $a$ expresses. For $a\in\mathfrak{S}$ that bound is the hypothesis
\begin{equation*}
f_a(e\oplus x)=g_a(x)\le\theta .
\end{equation*}
For $a$ pure, the bound is the one supplied by the admissibility of $e$, which is a hypothesis of the lemma.

Second, $U_a$ has top averages equal to $V=e\oplus\sigma(x)$ on $a^{\sim 4}$. Indeed, $U_a$ belongs to its own constrained class: it minimises among configurations whose top averages on $a^{\sim 4}$ are the prescribed data $V|_{a^{\sim 4}}$.

\emph{Step 3: the block plaquette.} The configuration $U_a$ is $\theta\varepsilon_{k_o}$-regular on $\tilde a$ (Step~2). The region $\tilde a$ contains the four blocks of $P$ (Step~1), and the block averages of $U_a$ on these blocks are given by $V$ (second part of Step~2). Applied as stated to $U_a$ on $\tilde a$, \cite[Proposition~2]{B-CMP98} gives
\begin{equation*}
\|V(\partial P)-1\|<2\theta\varepsilon_{k_o},
\end{equation*}
which is \eqref{8.11:eq-clamped-block-a}.

The proof uses neither the estimate \cite[(1.19)]{B-CMP119}, nor the minimiser $U(e\oplus x)$ of the full problem, nor any criticality of $x$ along the fibres of $\sigma$.
\end{proof}

\begin{lemma}[{The signed block-averaging identity; \cite[(45)--(49)]{B-CMP98}}]
\label{8.11:lem-signed-averaging}
Let $d=4$ and let $n$ be the block side.

\emph{Blocks and the averaged flux.} For a block site $y$ let $B(y)$ be the block of $y$, the set of
the $n^d$ lattice sites attached to $y$. Fix two distinct directions $\mu\neq\nu$ and let
$P=P_{\mu\nu}(y)$ be the block plaquette at $y$ in the plane of $e_\mu,e_\nu$. For a site $x$ let
$(P)_x$ be the translate of $P$ with lower corner $x$, and let $\Delta(P)$ be the neighbourhood of
$P$. For a plaquette $p$ let $\Phi(p)=\frac1i\log U(\partial p)$, as in
Notation~\ref{8.11:not-units-cap-amplitude}. Write $\Phi_{\mu\nu}(z)$ for $\Phi$ of the plaquette
with lower corner $z$ in the plane of $e_\mu,e_\nu$. Put
\begin{equation}\label{8.11:eq-signed-averaging-1}
\mathsf{A}_P:=n^{-d}\sum_{x\in B(y)}\ \sum_{p\subset(P)_x}\Phi(p)
=n^{-d}\sum_{z}\omega_y(z)\,\Phi_{\mu\nu}(z),
\end{equation}
where
\begin{equation}\label{8.11:eq-signed-averaging-2}
\omega_y(z):=\#\bigl\{(x,s):\ x\in B(y),\ s=s_\mu e_\mu+s_\nu e_\nu,\
0\le s_\mu,s_\nu<n,\ x+s=z\bigr\}.
\end{equation}
The weight $\omega_y$ is the product of two tents, in the directions $\mu$ and $\nu$, and of the
indicator of the range of the block $B(y)$ in the other directions. Moreover, with the sum taken
over all block sites $y$,
\begin{equation}\label{8.11:eq-signed-averaging-3}
\sum_{y}\omega_y(z)=n^2\qquad\text{for every site } z.
\end{equation}

\emph{The signed bound.} Let $U$ satisfy $\|U(\partial p)-1\|\le\alpha\eta^2$ for every plaquette
$p$ in $\Delta(P)$, with $\eta$ as in Notation~\ref{8.11:not-units-cap-amplitude}, where
$\alpha\le c_2$, $c_2$ being the smallness threshold of the unsigned block-averaging estimate. Let
$\mathsf{A}_P^{\mathrm{loc}}$ denote $\mathsf{A}_P$ computed in the local axial gauge at $P$, and let
$\bar U(\partial P)$ be the holonomy around $\partial P$ of the block-averaged configuration $\bar U$
(Definition~\ref{1.5:def-block-average}). Then
\begin{equation}\label{8.11:eq-signed-averaging-a}
\bigl\|\mathsf{A}_P^{\mathrm{loc}}\bigr\|\le\bigl\|\bar U(\partial P)-1\bigr\|+C_0''\alpha^2 ,
\end{equation}
where $C_0''$ is the constant of the unsigned block-averaging estimate and $\|\cdot\|$ is the
operator norm.
\end{lemma}

In \cite[(45)--(49)]{B-CMP98} the following three relations are stated; we write them with the
block side $n$ in place of $L$. There $p$ is a plaquette, $p'$ a block plaquette with block site
$y_0$, $(p')_x$ the translate of $p'$ with lower corner $x$, $c$ runs over the bonds of
$\partial p'$ and $c_-$ is the initial site of the bond $c$; $V_0$ is the configuration and
$\bar V_0$ its block average, $A(\partial p)$ is the flux through $p$ and $A(\Gamma_{c,x})$ the
flux along Ba{\l}aban's path $\Gamma_{c,x}$, the letter $A$ denoting here the field of
\cite{B-CMP98} and not the Wilson action; $\alpha_0$ is the smallness parameter of the
configurations there. The relations are
\begin{align*}
&\Bigl|\bar V_0(\partial p')-1-i\sum_{c\subset\partial p'}\ \sum_{x\in B(c_-)}
n^{-d}A(\Gamma_{c,x})\Bigr|<O(1)\,(n^2\alpha_0)^2,\\
&\sum_{c\subset\partial p'}\ \sum_{x\in B(c_-)}n^{-d}A(\Gamma_{c,x})
=\sum_{x\in B(y_0)}n^{-d}\sum_{p\subset(p')_x}A(\partial p),\\
&\bigl|V_0(\partial p)-1-iA(\partial p)\bigr|\le O(1)\,n^2\alpha_0^2,
\end{align*}
which are \cite[(47)]{B-CMP98}, \cite[(48)]{B-CMP98} and \cite[(49)]{B-CMP98}, in this order.

The identity \cite[(48)]{B-CMP98} thus rewrites the sum of the averaged fluxes along the paths
$\Gamma_{c,x}$ as the block average over $x\in B(y_0)$ of the sums
$\sum_{p\subset(p')_x}A(\partial p)$, a sum of the same shape as the middle member of
\eqref{8.11:eq-signed-averaging-1} with $A(\partial p)$ in place of $\Phi(p)$.

The display \eqref{8.11:eq-signed-averaging-a} is the reading of these three statements for the
signed sum $\mathsf{A}_P$ itself, in the local axial gauge at $P$.

The second equality in \eqref{8.11:eq-signed-averaging-1} regroups the double sum according to the site $z=x+s$: the plaquettes $p\subset(P)_x$ are those with lower corner $x+s$, $0\le s_\mu,s_\nu<n$, in the plane of $e_\mu,e_\nu$, so $\Phi_{\mu\nu}(z)$ occurs exactly $\omega_y(z)$ times.

For \eqref{8.11:eq-signed-averaging-3}, fix $z$. Each of the $n^2$ admissible values of $s$ determines $x=z-s$. That site lies in exactly one block $B(y)$, so each value of $s$ contributes exactly one triple $(y,x,s)$.

In the unsigned form of the block-averaging estimate it is the modulus of exactly the sum $\mathsf{A}_P$ that is bounded; \eqref{8.11:eq-signed-averaging-a} is the signed form of that estimate.

\begin{lemma}[Variation of the projection kernel]\label{8.11:lem-kernel-variation}
Let $p_0$ be a non-$\partial$ plaquette of $\Omega$, and let $\mathcal{K}_\Omega(p_0,\cdot)$ be the kernel of
$\Pi^0$ of Corollary~\ref{8.11:cor-projection-kernels}, display~\eqref{8.11:eq-projection-kernels-a}.
There is a constant $C_\Pi''=C(d,\varrho)$, depending only on $d$ and $\varrho$ and not on $\Omega$ or
$p_0$, such that for every coordinate direction $j$, $1\le j\le d$, with $e_j$ the unit lattice vector in
direction $j$, and all non-$\partial$ plaquettes $p$, $p+e_j$ of $\Omega$
\begin{equation}\label{8.11:eq-kernel-variation-a}
\bigl|\mathcal{K}_\Omega(p_0,p)-\mathcal{K}_\Omega(p_0,p+e_j)\bigr|
\le C_\Pi''\,\bigl(1+|x(p_0)-x(p)|_\infty\bigr)^{-(d+1)} .
\end{equation}
In particular, if the barycentre of $p_0$ is the centre of $\Omega$, then for a plaquette $p$ of $\Omega$
one has $1+|x(p_0)-x(p)|_\infty=\rho(p)$, with $\rho$ one plus the $\ell^\infty$-distance to the centre,
and the right side of \eqref{8.11:eq-kernel-variation-a} is $C_\Pi''\rho(p)^{-(d+1)}$.
\end{lemma}

\begin{proof}
\emph{Harmonicity.} Let $G_{\mathbb{T}}$ be the mean-zero solution on the doubled torus $\mathbb{T}$ of
$-\Delta G_{\mathbb{T}}=\mathbf{1}_{\{0\}}-|\mathbb{T}|^{-1}$, as in Lemma~\ref{8.11:lem-torus-kernel},
where $\mathbf{1}_{\{0\}}$ is the indicator function of the origin.
Differences commute with $\Delta$, so
\begin{equation*}
-\Delta\,\nabla_\mu\nabla^-_\nu G_{\mathbb{T}}
=\nabla_\mu\nabla^-_\nu\bigl(\mathbf{1}_{\{0\}}-|\mathbb{T}|^{-1}\bigr)
=\nabla_\mu\nabla^-_\nu\mathbf{1}_{\{0\}} ,
\end{equation*}
since differences of the constant $|\mathbb{T}|^{-1}$ vanish. The right-hand side is supported in the
unit neighbourhood of $0$; hence $\nabla_\mu\nabla^-_\nu G_{\mathbb{T}}$ is harmonic on $\mathbb{T}$ off
that neighbourhood.

\emph{Third differences.} Let $|x|$ be the $\ell^\infty$-distance of $x$ to $0$ on $\mathbb{T}$. If
$|x|\ge4$, every $y$ in the ball of radius $|x|/2$ about $x$ satisfies $|y|\ge|x|/2\ge2$, so the ball
lies in the harmonicity region, and Lemma~\ref{8.11:lem-torus-kernel} gives there
\begin{equation*}
|\nabla_\mu\nabla^-_\nu G_{\mathbb{T}}(y)|\le C_{G2}(1+|x|/2)^{-d}\le 2^dC_{G2}(1+|x|)^{-d}.
\end{equation*}
The interior gradient estimate for discrete harmonic functions (a difference at the centre of a ball of
radius $r$ is bounded by $Cr^{-1}$ times the supremum on the ball), applied on this ball with
$r=|x|/2$, and $2/|x|\le\tfrac52(1+|x|)^{-1}$ for $|x|\ge4$, give
\begin{equation}\label{8.11:eq-kernel-variation-1}
|\nabla_j\nabla_\mu\nabla^-_\nu G_{\mathbb{T}}(x)|\le C(1+|x|)^{-(d+1)} ,
\end{equation}
with $C=C(d,\varrho)$.
If $|x|<4$, the triangle inequality and Lemma~\ref{8.11:lem-torus-kernel} bound the left side by
$2C_{G2}\le 2\cdot5^{d+1}C_{G2}(1+|x|)^{-(d+1)}$, so, after enlarging $C$,
\eqref{8.11:eq-kernel-variation-1} holds on all of $\mathbb{T}$.

\emph{Images.} By the proof of Corollary~\ref{8.11:cor-projection-kernels}, $\mathcal{K}_\Omega(p_0,p)$ is
obtained from the kernel of $\Pi_{\mathbb{T}}$, whose entries are $\pm\nabla_\mu\nabla^-_\lambda
G_{\mathbb{T}}$ by the formula for $\Pi_{\mathbb{T}}\tilde\Psi$ there, by restricting to $\Omega$ and
collecting the $2^d$ images of $p$. We collect the images of $p$ and of $p+e_j$ in the same way; under
each reflection the translate by $e_j$ goes to a translate by $\pm e_j$, so the difference in
\eqref{8.11:eq-kernel-variation-a} is a sum of at most $2^{d+1}d$ terms (for each of the $2^d$ images,
at most the $2d$ summands of the formula for $(\Pi_{\mathbb{T}}\tilde\Psi)_{\mu\nu}$ there), each a value of
$\pm\nabla_j\nabla_\mu\nabla^-_\lambda G_{\mathbb{T}}$ at a unit shift of the torus displacement from
$x(p_0)$ to an image of $x(p)$. By that proof, this displacement has torus distance at least
$|x(p_0)-x(p)|_\infty$, and unit shifts cost a factor at most $2^d$ in $(1+|x|)^{-d}$, hence at most
$2^{d+1}$ in $(1+|x|)^{-(d+1)}$. With \eqref{8.11:eq-kernel-variation-1}, each term is at most
$2^{d+1}C(1+|x(p_0)-x(p)|_\infty)^{-(d+1)}$, and summing gives \eqref{8.11:eq-kernel-variation-a} with
$C_\Pi''=2^{2d+2}d\,C$, which depends only on $d$ and $\varrho$.
\end{proof}

\begin{lemma}[Weighted flux through a block region]\label{8.11:lem-weighted-flux}
Let $U$, the complete axial gauge $\mathfrak{G}$ of $\Omega_*$ centred at $z_0$ and the
$\mathfrak{g}$-valued $2$-cochain $\Phi$ be as in Lemma~\ref{8.11:lem-nonabelian-defects}.
Let $\varepsilon_a$ and $\eta$ also be as there, and put $\mathsf{L}:=\operatorname{diam}_1(\Omega_*)$.
Put $V:=\bar U^{k_o}(U)$, the top block average of $U$. Let the block side $n$, the blocks $B(y)$ of
fine sites, the block plaquettes $P_{\alpha\beta}(y)$ and the tent weights $\omega_y$ be as in the
signed block-averaging identity (Lemma~\ref{8.11:lem-signed-averaging}).

Fix directions $\alpha\neq\beta$, a real number $\mathsf{d}$, a finite set $\mathcal{R}$ of fine
sites, and a function $W\colon\mathcal{R}\to\mathbb{R}$, extended by $0$ outside $\mathcal{R}$.
In this lemma $\mathsf{d}$ denotes this bound on block data, not the separation of supports.
Write $\Phi(z):=\Phi_{\alpha\beta}(z)$ for the value of $\Phi$ on the plaquette in the
$(\alpha,\beta)$-plane with lowest corner $z$. Let $\mathcal{Y}$ be the set of top block sites $y$
such that
\begin{equation*}
\operatorname{supp}\omega_y\subset\mathcal{R}
\qquad\text{and}\qquad
\|V(\partial P_{\alpha\beta}(y))-1\|\le\mathsf{d}.
\end{equation*}
For $z\in\mathcal{R}$ put
\begin{equation*}
\operatorname{osc}W(z):=\max\{|W(z)-W(z')| : z'\in\mathcal{R},\ |z-z'|_\infty<2n\},
\end{equation*}
and let
\begin{equation*}
\partial\mathcal{R}:=\{z\in\mathcal{R} : \omega_y(z)>0 \text{ for some } y\notin\mathcal{Y}\}.
\end{equation*}
Then
\begin{equation}\label{8.11:eq-weighted-flux-1}
\begin{split}
\Bigl\|\sum_{z\in\mathcal{R}}W(z)\Phi_{\alpha\beta}(z)\Bigr\|
&\le n^{-2}\mathsf{d}'\sum_{z\in\mathcal{R}}\bigl[|W(z)|+\operatorname{osc}W(z)\bigr]
+\sum_{z\in\mathcal{R}}\operatorname{osc}W(z)\,\|\Phi(z)\|\\
&\quad+2\sum_{z\in\partial\mathcal{R}}\bigl[|W(z)|+\operatorname{osc}W(z)\bigr]\|\Phi(z)\|,
\end{split}
\end{equation}
where
\begin{equation*}
\mathsf{d}':=\mathsf{d}+C_{\mathrm{av}}(1+\mathsf{L}\eta)\varepsilon_a^2
\end{equation*}
and $C_{\mathrm{av}}$ is the constant fixed in the proof.
\end{lemma}

\begin{proof}
\emph{Preliminaries.} For each $y\in\mathcal{Y}$ fix a site $y'\in B(y)$ and put
$\overline{W}_y:=W(y')$.

For $x\in B(y)$ we have $\omega_y(x)\ge1$; this is seen by taking $s=0$ in the definition of
$\omega_y$. Hence $B(y)\subset\operatorname{supp}\omega_y\subset\mathcal{R}$ for $y\in\mathcal{Y}$,
and in particular $y'\in\mathcal{R}$.

If $y\in\mathcal{Y}$ and $\omega_y(z)>0$, then $z=x+s$ with $x\in B(y)$ and
$0\le s_\alpha,s_\beta\le n-1$. Since $x$ and $y'$ both lie in the block $B(y)$ of side $n$, we get
$|z-y'|_\infty\le|s|_\infty+|x-y'|_\infty\le 2n-2<2n$. Consequently, for every $z\in\mathcal{R}$ with
$\omega_y(z)>0$,
\begin{equation}\label{8.11:eq-weighted-flux-2}
|\overline{W}_y|\le|W(z)|+|W(z)-W(y')|\le|W(z)|+\operatorname{osc}W(z).
\end{equation}

Finally, Lemma~\ref{8.11:lem-signed-averaging} states that $\sum_y\omega_y(z)=n^2$ for every $z$,
the sum running over all top block sites.

\emph{Decomposition.} The function $z\mapsto n^{-2}\sum_{y\in\mathcal{Y}}\overline{W}_y\omega_y(z)$
vanishes outside $\mathcal{R}$, because $\operatorname{supp}\omega_y\subset\mathcal{R}$ for
$y\in\mathcal{Y}$. Adding and subtracting it therefore gives
\begin{equation}\label{8.11:eq-weighted-flux-3}
\begin{split}
\sum_{z\in\mathcal{R}}W(z)\Phi(z)
&=n^{-2}\sum_{y\in\mathcal{Y}}\overline{W}_y\sum_{z}\omega_y(z)\Phi(z)\\
&\quad+\sum_{z\in\mathcal{R}}\Bigl[W(z)-n^{-2}\sum_{y\in\mathcal{Y}}\overline{W}_y\,\omega_y(z)\Bigr]
\Phi(z).
\end{split}
\end{equation}

\emph{The bracket.} Let $z\in\mathcal{R}\setminus\partial\mathcal{R}$. Every $y$ with
$\omega_y(z)>0$ lies in $\mathcal{Y}$, so $\sum_{y\in\mathcal{Y}}\omega_y(z)=n^2$, and the bracket
equals
\begin{equation*}
n^{-2}\sum_{y}\omega_y(z)\bigl(W(z)-\overline{W}_y\bigr).
\end{equation*}
For each $y$ contributing here, $y'\in\mathcal{R}$ and $|z-y'|_\infty<2n$, so
$|W(z)-\overline{W}_y|\le\operatorname{osc}W(z)$. The weights $n^{-2}\omega_y(z)$ are non-negative
with sum $1$, hence the bracket has modulus at most $\operatorname{osc}W(z)$.

Now let $z\in\partial\mathcal{R}$. Since $\sum_{y\in\mathcal{Y}}\omega_y(z)\le n^2$, the bracket has
modulus at most
\begin{equation*}
|W(z)|+\max\{|\overline{W}_y| : y\in\mathcal{Y},\ \omega_y(z)>0\}.
\end{equation*}
By \eqref{8.11:eq-weighted-flux-2} this is at most
$2|W(z)|+\operatorname{osc}W(z)\le2[|W(z)|+\operatorname{osc}W(z)]$; if no such $y$ exists, the
bound is $|W(z)|$. Since $\partial\mathcal{R}\subset\mathcal{R}$, the second sum in
\eqref{8.11:eq-weighted-flux-3} has norm at most
\begin{equation*}
\sum_{z\in\mathcal{R}}\operatorname{osc}W(z)\,\|\Phi(z)\|
+2\sum_{z\in\partial\mathcal{R}}\bigl[|W(z)|+\operatorname{osc}W(z)\bigr]\|\Phi(z)\|.
\end{equation*}

\emph{The first sum: change of gauge.} Fix $y\in\mathcal{Y}$ and let $P:=P_{\alpha\beta}(y)$. By the
definition of the block-averaged flux in Lemma~\ref{8.11:lem-signed-averaging},
\begin{equation*}
\sum_z\omega_y(z)\Phi(z)=n^d\mathsf{A}^{\mathfrak{G}}_P,
\end{equation*}
where the superscript records that $\Phi$ is computed in the gauge $\mathfrak{G}$. The signed
block-averaging identity is stated in the local axial gauge at $P$ on $\Delta(P)$. This gauge is
$U^{\mathrm{loc}}=g\cdot U^{\mathfrak{G}}$ for a gauge transformation $g$ with $g(y)=1$, so that
\begin{equation*}
U^{\mathrm{loc}}_b=g(b_-)\,U^{\mathfrak{G}}_b\,g(b_+)^{-1}.
\end{equation*}
Rearranging $U^{\mathrm{loc}}_b g(b_+)=g(b_-)U^{\mathfrak{G}}_b$ gives
\begin{equation*}
g(b_-)-g(b_+)=(U^{\mathrm{loc}}_b-1)\,g(b_+)-g(b_-)\,(U^{\mathfrak{G}}_b-1).
\end{equation*}
The values of $g$ are unitary, so, as in the proof of Lemma~\ref{8.11:lem-nonabelian-defects},
\begin{equation*}
\|g(b_+)-g(b_-)\|\le\|U^{\mathrm{loc}}_b-1\|+\|U^{\mathfrak{G}}_b-1\|
\le(2dn+\mathsf{L})\,\varepsilon_a\eta^2 .
\end{equation*}
Here the first term is bounded by means of \cite[(46)]{B-CMP98}, and the second by the bond bound
of the complete axial gauge $\mathfrak{G}$ used in Lemma~\ref{8.11:lem-nonabelian-defects}.
Summing these differences along a path of at most $2dn$ bonds from $y$, and using $g(y)=1$, we get
at the sites in question
\begin{equation}\label{8.11:eq-weighted-flux-4}
\|g-1\|\le2dn(2dn+\mathsf{L})\,\varepsilon_a\eta^2\le C(1+\mathsf{L}\eta)\,\varepsilon_a .
\end{equation}

\emph{The first sum: the two block fluxes.} For every plaquette $p$ under consideration,
$\|U(\partial p)-1\|\le\varepsilon_a\eta^2$. This bound is unchanged by a gauge transformation,
which conjugates $U(\partial p)$. Since $\Phi(p)=(1/i)\log U(\partial p)$, we have
$\|\Phi(p)\|\le2\varepsilon_a\eta^2$. There are $n^{d}\cdot n^2$ pairs $(x,s)$ in the definition of
$\omega_y$, so $\sum_z\omega_y(z)=n^{d+2}$. Hence, in either gauge,
\begin{equation}\label{8.11:eq-weighted-flux-5}
n^{-d}\sum_z\omega_y(z)\|\Phi(z)\|\le n^2\cdot2\varepsilon_a\eta^2=2\varepsilon_a .
\end{equation}
As in the proof of Lemma~\ref{8.11:lem-nonabelian-defects},
$\Phi^{\mathrm{loc}}(p)=\operatorname{Ad}(g(x_p))\Phi(p)$. Writing
$gXg^{-1}-X=(g-1)Xg^{-1}+X(g^{-1}-1)$ and using $\|g^{-1}-1\|=\|g-1\|$ for unitary $g$, we obtain
\begin{equation*}
\|\Phi(p)-\Phi^{\mathrm{loc}}(p)\|\le2\|g(x_p)-1\|\,\|\Phi(p)\|.
\end{equation*}
Inserting this into the definition of $\mathsf{A}_P$ and using \eqref{8.11:eq-weighted-flux-5} gives
\begin{equation*}
\|\mathsf{A}^{\mathfrak{G}}_P-\mathsf{A}^{\mathrm{loc}}_P\|\le2\|g-1\|\cdot2\varepsilon_a .
\end{equation*}

We now apply the bound \eqref{8.11:eq-signed-averaging-a} of
Lemma~\ref{8.11:lem-signed-averaging}, in the local axial gauge at $P$ with $\alpha=\varepsilon_a$.
Since $\bar U(\partial P)=V(\partial P)$ and $y\in\mathcal{Y}$, it gives
\begin{equation*}
\|\mathsf{A}^{\mathrm{loc}}_P\|\le\|V(\partial P)-1\|+C_0''\varepsilon_a^2\le\mathsf{d}+C_0''\varepsilon_a^2 .
\end{equation*}
Together with \eqref{8.11:eq-weighted-flux-4} we conclude
\begin{equation*}
\|\mathsf{A}^{\mathfrak{G}}_P\|\le\mathsf{d}+C_0''\varepsilon_a^2+4C(1+\mathsf{L}\eta)\varepsilon_a^2
\le\mathsf{d}',
\end{equation*}
with $C_{\mathrm{av}}:=C_0''+4C$.

\emph{Conclusion.} By \eqref{8.11:eq-weighted-flux-2} applied to each $z\in B(y)$, we have
$|\overline{W}_y|\le|W(z)|+\operatorname{osc}W(z)$. Averaging over the $n^d$ sites of $B(y)$ gives
\begin{equation*}
|\overline{W}_y|\le n^{-d}\sum_{z\in B(y)}\bigl[|W(z)|+\operatorname{osc}W(z)\bigr].
\end{equation*}
The first sum in \eqref{8.11:eq-weighted-flux-3} therefore has norm at most
\begin{equation*}
n^{-2}\sum_{y\in\mathcal{Y}}|\overline{W}_y|\,n^d\,\|\mathsf{A}^{\mathfrak{G}}_{P_{\alpha\beta}(y)}\|
\le n^{-2}\mathsf{d}'\sum_{y\in\mathcal{Y}}\sum_{z\in B(y)}\bigl[|W(z)|+\operatorname{osc}W(z)\bigr].
\end{equation*}
The blocks $B(y)$ are disjoint and, for $y\in\mathcal{Y}$, contained in $\mathcal{R}$, so the right
side is at most $n^{-2}\mathsf{d}'\sum_{z\in\mathcal{R}}[|W(z)|+\operatorname{osc}W(z)]$. Adding the
bound for the second sum proves \eqref{8.11:eq-weighted-flux-1}.
\end{proof}

\begin{theorem}[The deep curvature bound through block data]\label{8.11:thm-deep-curvature-blocks}
Assume the following:
\begin{itemize}
\item[(a)] the extension part and the regularity part of the feasibility input of
Theorem~\ref{8.11:thm-deep-curvature}, and no other part of that input;
\item[(b)] the block curvature bound in the clamped layer, Lemma~\ref{8.11:lem-clamped-block};
\item[(c)] the signed block-averaging identity, Lemma~\ref{8.11:lem-signed-averaging};
\item[(d)] the inputs of Theorem~\ref{8.11:thm-deep-curvature} other than its feasibility input
and other than \eqref{8.11:eq-young-row-cover-a}, among them the energy bounds below the support
value, Lemma~\ref{8.11:lem-energy-below-support};
\item[(e)] the covering of a depth band by the first young row, \eqref{8.11:eq-young-row-cover-a}.
\end{itemize}
Let $e$, $\theta$ and $\hat x$ be as in Theorem~\ref{8.11:thm-deep-curvature}, under the ceilings
of that theorem, and take the reference depth $c_*=-(3w+1)$. Assume in addition the one-sided
condition
\begin{equation}\label{8.11:eq-deep-curvature-blocks-7}
C_{av}(1+\mathsf{L}\eta)\varepsilon_a^2\le\vartheta_1\varepsilon_{k_o},
\end{equation}
a ceiling of the same kind as $\vartheta_{\mathrm{nl}}(h)\le\vartheta_1$ in
\eqref{8.11:eq-deep-bound-scope-a}, where $C_{av}$ is the constant of
Lemma~\ref{8.11:lem-weighted-flux}, $\mathsf{L}$ is the $\ell^1$-diameter of the box
$\Omega_{c_*}$ and $\eta$ is the fine spacing in top units. Let $h$ satisfy $h\ge r_1$ and
$hn\ge4$, as there, and in addition $h\ge4$, with $h$ counted in top blocks; the last condition
holds at $h=w/2$, since $w=M_2R_{k_o}=L^2M_1R_{k_o}\ge8$. Put
\begin{equation}\label{8.11:eq-deep-curvature-blocks-1}
C_o:=\frac{2^{d+3}d\,C_\Pi''}{C_\Pi},\qquad
\kappa_{\mathrm{blk}}(h):=C_{\mathrm{blk}}
\Bigl[N\,C_E\,\hat N_W\bigl(h^{-5}+t\,h^{-6}+t^2h^{-8}\bigr)\Bigr]^{1/2}
\frac{\varepsilon_1}{\varepsilon_{k_o}},
\end{equation}
where $C_{\mathrm{blk}}$ depends on $d$ and $\varrho_H$ only. Then for every fine plaquette $p$
with $\delta(p)\ge d_1+h$ the bound \eqref{8.11:eq-deep-curvature-a} holds with the following two
changes. First, in the two members $2c_1t_c\theta''/h$ and $4c_2(t+1)^2\theta''/h^2$ of
$\mathfrak{B}(h)$ in \eqref{8.11:eq-bound-members-a}, the quantity $\theta''$ of
\eqref{8.11:eq-inherited-bounds-a} is replaced by $2\theta(1+\vartheta_1)(1+C_o/h)$, so that
these two members become
\begin{equation}\label{8.11:eq-deep-curvature-blocks-2}
\frac{2c_1t_c}{h}\cdot2\theta(1+\vartheta_1)\Bigl(1+\frac{C_o}{h}\Bigr)
\qquad\text{and}\qquad
\frac{4c_2(t+1)^2}{h^2}\cdot2\theta(1+\vartheta_1)\Bigl(1+\frac{C_o}{h}\Bigr);
\end{equation}
all other members of $\mathfrak{B}(h)$ are unchanged. Second, the additional error
$\kappa_{\mathrm{blk}}(h)$ is added inside the bracket of \eqref{8.11:eq-deep-curvature-a}.
\end{theorem}

The proof below uses neither the comparison, in the clamped layer, of the global minimiser with the
straddling local minimisers, nor the regularity statement for the straddling sites, which is
distinct from the regularity part of the feasibility input, nor the justification of
that comparison at points critical along the fibres of the top averages.

\begin{proof}
We use the notation of the proof of Theorem~\ref{8.11:thm-deep-curvature}: $\check x$ is the
released minimiser of \eqref{8.11:eq-released-problem-a} and $\check U:=U(e\oplus\check x)$;
$\Omega_1:=\Omega_{d_1}$ and $\Omega_*:=\Omega_{c_*}$; $\tau:=tn$; $p_0:=p$, with lowest corner
$z_0$; $h_z:=hn$; below we use $h_z\ge hn/2$ in the shell sums and, from $h\ge4$, $h_z\ge4n$ in
the oscillation bound; $\mathfrak{G}$ is the centred axial gauge
and $\Phi:=\check\Phi$ the logarithm of the plaquette variables of $\check U$ in $\mathfrak{G}$;
$\rho_2$ and $N_1$ are as in Proposition~\ref{8.11:prop-trace-representation}; the path--surface
system is that of Lemma~\ref{8.11:lem-system-geometry}. The fine spacing in top units is
$\eta=n^{-1}$. The layer $\mathcal{L}$ between $\Omega_1$ and $\Omega_*$ is split, according to the
index $u$ of part (c) of Lemma~\ref{8.11:lem-system-geometry}, into $\mathcal{L}_c$
($u\ge(r_1+\varpi_y)n$), $\mathcal{L}_y$ ($r_1n\le u<(r_1+\varpi_y)n$) and the collar
$\mathcal{L}_{\mathrm{coll}}$ ($u<r_1n$); as shown there, $\mathcal{L}_c$ lies in
$\{\delta\le w\}$. Bonds $b'$ with $S^{\mathrm{lay}}(b')=N(b')$ are called bonds of Case~I,
those of $\mathfrak{b}_{II}$, with $S^{\mathrm{lay}}(b')=\mathrm{Sq}(b')$, bonds of Case~II.
Throughout, $\rho$ attached to a bond $b'$ means $\rho(p(b'))$, and $X\lesssim Y$ means
$X\le CY$ with $C$ depending on $d$ and $\varrho_H$ only; $C_\Pi$, $C_\Pi''$, $c_1$, $c_2$ are
such constants.

\emph{The identity.} The proof of Proposition~\ref{8.11:prop-trace-representation} gives, with
signs $\varsigma_{b'},\varsigma_q\in\{-1,+1\}$ and $\mathcal{K}$ the kernel of $\Pi^0$ occurring
there, the identity
\begin{equation}\label{8.11:eq-deep-curvature-blocks-3}
\Phi(p_0)=-\sum_{b'}\varsigma_{b'}\,\mathcal{K}(p_0,p(b'))\sum_{q\in S(b')}
\varsigma_q\,(\Phi-\rho_2)(q)
+\bigl(\Pi^0\Phi-\Pi^0\rho_2+\rho_2\bigr)(p_0).
\end{equation}
The proof of Theorem~\ref{8.11:thm-deep-curvature} bounds $\|\Phi(p_0)\|$ from this identity by
taking norms term by term; the result is the sum of the layer term, the reference-surface term and
the defect term. We keep all of that proof except the bound of the part of the layer term in which
$q\in S^{\mathrm{lay}}(b')\cap\mathcal{L}_c$ and the summand is $\Phi(q)$; there the proof of
Theorem~\ref{8.11:thm-deep-curvature} puts $\|\Phi(q)\|\le\theta''\,\varepsilon_{k_o}\eta^2$ and
obtains the members $2c_1t_c\theta''/h$ and $4c_2(t+1)^2\theta''/h^2$ of $\mathfrak{B}(h)$. The
terms with $\rho_2(q)$, those with $q\in S^*(b')$, those with $q$ in $\mathcal{L}_y$ or
$\mathcal{L}_{\mathrm{coll}}$, and the last term of \eqref{8.11:eq-deep-curvature-blocks-3} are
bounded exactly as there. In place of the removed bound we estimate the signed sum
\begin{equation}\label{8.11:eq-deep-curvature-blocks-4}
\mathsf{F}_c:=\sum_{b'}\varsigma_{b'}\,\mathcal{K}(p_0,p(b'))
\sum_{q\in S^{\mathrm{lay}}(b')\cap\mathcal{L}_c}\varsigma_q\,\Phi(q)
\end{equation}
without taking norms inside it.

\emph{Grouping of the bonds.} (I) The bonds of Case~I are grouped by the face $(i,\pm)$ of
$\Omega_1$ on which they lie and by their direction $k$. On a group the signs $\varsigma_{b'}$,
$\varsigma_q$ are constant, with product $\varsigma$, and $q$ runs over the plaquettes of
orientation $(i,k)$ based at the sites $z$ of a $d$-dimensional slab $\mathcal{R}$, the product
of the set of positions of Case~I bonds in the face with the normal range of length $t_cn$
belonging to $\mathcal{L}_c$; we write $b'(z)$ for the Case~I bond with the plaquette at $z$ in
$N(b'(z))$. Put $W(z):=\varsigma\,\mathcal{K}(p_0,p(b'(z)))$. Then $W$ is constant along the normal
direction. If $z,z'\in\mathcal{R}$ and $|z-z'|_\infty<2n$, the plaquettes $p(b'(z))$ and
$p(b'(z'))$ are joined by at most $2dn$ unit translations, along which $\rho$ stays at least
$\rho-2n\ge\rho/2$, because $\rho\ge h_z\ge4n$, which follows from $h\ge4$. Applying the
variation of the projection kernel, Lemma~\ref{8.11:lem-kernel-variation}, at each translation
gives
\begin{equation}\label{8.11:eq-deep-curvature-blocks-5}
\operatorname{osc}W(z)\le2dn\cdot C_\Pi''\,\rho^{-(d+1)}\cdot2^{d+1}
=2^{d+2}d\,C_\Pi''\,n\,\rho^{-(d+1)}.
\end{equation}
(II) The bonds of Case~II are grouped by the edge $(k,m,\pm,\pm)$ on which they lie. On a group
$q$ runs over the plaquettes of orientation $(k,m)$ based at the sites of the set $\mathcal{R}$,
the product of the positions $x''$ along the edge with the part of the rectangle where
$\max(u_k,u_m)\ge(r_1+\varpi_y)n$; put $W(z):=\varsigma\,\mathcal{K}(p_0,p(b'(x'')))$ for $z$
above $x''$. Here $W$ is constant in the two directions of the rectangle and varies only along the
edge, so the argument just given yields \eqref{8.11:eq-deep-curvature-blocks-5} also in this case.
With this notation $\mathsf{F}_c$ is the sum over all groups of
$\sum_{z\in\mathcal{R}}W(z)\Phi_{\alpha\beta}(z)$, with $(\alpha,\beta)=(i,k)$ in (I) and
$(\alpha,\beta)=(k,m)$ in (II), and $\|\mathsf{F}_c\|$ is at most the sum of the norms of these
group sums.

\emph{The block data.} The configuration is $\check U$. Its top averages on $\{\delta\le w\}$ are
$e\oplus\sigma(\check x)=e\oplus\sigma(\hat x)$: indeed $\check x=\hat x$ on the near hole data
bonds $\mathfrak{D}_{\mathrm{near}}$, and the top averages on $\{\delta\le w\}$ depend only on
these by the iterate structure of the block-averaging map \cite[(69)]{B-CMP98}. The block curvature
bound in the clamped layer, Lemma~\ref{8.11:lem-clamped-block}, applied at $\hat x$ (which lies in
$\widehat{\mathcal{W}}^{\theta}(e)$, the class constrained by the straddling windows at threshold
$\theta$), therefore gives $\|\partial V(P_{\alpha\beta}(y))-1\|<2\theta\varepsilon_{k_o}$ for the
top average $V$ of $\check U$ at every top block site $y$ with
$\operatorname{supp}\omega_y\subset\mathcal{R}\subset\{\delta\le w\}$. We apply the weighted flux
through a block region, Lemma~\ref{8.11:lem-weighted-flux}, to each group, with its directions
$\alpha,\beta$, its set $\mathcal{R}$, its weight $W$, and $\mathsf{d}:=2\theta\varepsilon_{k_o}$;
by what was just said, every block site $y$ with $\operatorname{supp}\omega_y\subset\mathcal{R}$
belongs to $\mathcal{Y}$. The lemma bounds each group sum by a main term and two error terms, which
we treat in turn.

\emph{The main term.} By the decay of the projection kernels \eqref{8.11:eq-projection-kernels-a},
which is the bound giving the factor $C_\Pi\rho^{-d}$ in
Proposition~\ref{8.11:prop-trace-representation}, we have $|W|\le C_\Pi\rho^{-d}$. Together with
\eqref{8.11:eq-deep-curvature-blocks-5} and $n/\rho\le n/h_z\le2/h$,
\begin{align*}
|W|+\operatorname{osc}W
&\le C_\Pi\rho^{-d}+2^{d+2}d\,C_\Pi''\,n\,\rho^{-(d+1)}\\
&\le C_\Pi\Bigl(1+\frac{2^{d+3}d\,C_\Pi''}{C_\Pi\,h}\Bigr)\rho^{-d}
=C_\Pi\Bigl(1+\frac{C_o}{h}\Bigr)\rho^{-d}.
\end{align*}
In (I) each bond contributes $t_cn$ sites of $\mathcal{R}$, and
$\sum_{b'\in\mathfrak{b}(\Omega_1)}\rho^{-d}\le c_1/h_z$ as in the bound, in the proof of
Theorem~\ref{8.11:thm-deep-curvature}, of the part of the layer term bounded pointwise by
$\theta''$ and $\theta_y''$; summing the main term of
Lemma~\ref{8.11:lem-weighted-flux} over the groups of Case~I and using $h_z\ge hn/2$ and
$n^{-2}=\eta^2$,
\begin{equation*}
n^{-2}\mathsf{d}'\sum_{\mathrm{(I)}}\sum_{z\in\mathcal{R}}\bigl(|W|+\operatorname{osc}W\bigr)
\le n^{-2}\mathsf{d}'\Bigl(1+\frac{C_o}{h}\Bigr)\frac{t_cn\,C_\Pi c_1}{h_z}
\le\Bigl(1+\frac{C_o}{h}\Bigr)2C_\Pi c_1\frac{t_c}{h}\,\mathsf{d}'\eta^2.
\end{equation*}
In (II) each edge position carries at most $(\tau+1)\tau\le(t+1)^2n^2$ sites of $\mathcal{R}$, and
$\sum_{\mathfrak{b}_{II}}\rho^{-d}\le c_2/h_z^2$ as in the same bound; hence
\begin{equation*}
n^{-2}\mathsf{d}'\sum_{\mathrm{(II)}}\sum_{z\in\mathcal{R}}\bigl(|W|+\operatorname{osc}W\bigr)
\le n^{-2}\mathsf{d}'\Bigl(1+\frac{C_o}{h}\Bigr)\frac{(\tau+1)\tau\,C_\Pi c_2}{h_z^2}
\le\Bigl(1+\frac{C_o}{h}\Bigr)4C_\Pi c_2\Bigl(\frac{t+1}{h}\Bigr)^2\mathsf{d}'\eta^2.
\end{equation*}
These are the two members that the proof of Theorem~\ref{8.11:thm-deep-curvature} obtains from the
part of the layer term bounded pointwise by $\theta''$, with $\theta''\varepsilon_{k_o}$ replaced by
$(1+C_o/h)\,\mathsf{d}'$. Here
$\mathsf{d}'=\mathsf{d}+C_{av}(1+\mathsf{L}\eta)\varepsilon_a^2$ is the quantity of
Lemma~\ref{8.11:lem-weighted-flux}, with $C_{av}$ its constant and $\mathsf{L}$ the
$\ell^1$-diameter of $\Omega_*$. By the condition \eqref{8.11:eq-deep-curvature-blocks-7}
assumed in the theorem, $C_{av}(1+\mathsf{L}\eta)\varepsilon_a^2\le\vartheta_1\varepsilon_{k_o}$,
so that $\mathsf{d}'=2\theta\varepsilon_{k_o}(1+o_{\gamma}(1))$,
and, since $\theta\ge1$,
\begin{equation*}
\mathsf{d}'\le2\theta\varepsilon_{k_o}+\vartheta_1\varepsilon_{k_o}
\le2\theta(1+\vartheta_1)\varepsilon_{k_o}.
\end{equation*}
So the main terms are bounded by the two members \eqref{8.11:eq-deep-curvature-blocks-2}, times
$C_\Pi\,\varepsilon_{k_o}\eta^2$, which is the first change asserted.

\emph{The energy bound.} Each error term is estimated by the Cauchy--Schwarz inequality against
\begin{equation}\label{8.11:eq-deep-curvature-blocks-6}
\sum_{q\subset\Omega_*}\|\Phi(q)\|^2\le8N\cdot5C_E\,\varepsilon_1^2\hat N_W.
\end{equation}
This follows from $\|\Phi(q)\|^2\le8N\,\mathfrak{e}(\check U(\partial q))$, used in the collar
bound of the proof of Theorem~\ref{8.11:thm-deep-curvature}, and from
\eqref{8.11:eq-energy-below-support-b}, which applies because $\mathcal{L}_c$ reaches $3w+1$ top
blocks outside $H$. As in the collar bound, each pair (group, site) corresponds to a pair $(b',q)$
with $q\in S^{\mathrm{lay}}(b')$, and by part (b) of Lemma~\ref{8.11:lem-system-geometry} each
plaquette $q$ occurs in at most three such pairs, so that the sum of $\|\Phi\|^2$ over all groups
and sites is at most three times the left side of \eqref{8.11:eq-deep-curvature-blocks-6}. We
use $d=4$, $h_z\ge hn/2$, and the series bounds of part (e) of
Lemma~\ref{8.11:lem-system-geometry} in the form used in the collar bound:
\begin{equation*}
\sum_{\mathfrak{b}}\rho^{-2d-2}\le c_1h_z^{-(d+3)},\quad
\sum_{\mathfrak{b}_{II}}\rho^{-2d-2}\le c_2h_z^{-(d+4)},\quad
\sum_{\mathfrak{b}}\rho^{-2d}\le c_1h_z^{-(d+1)},\quad
\sum_{\mathfrak{b}_{II}}\rho^{-2d}\le c_2h_z^{-(d+2)},
\end{equation*}
with $\mathfrak{b}=\mathfrak{b}(\Omega_1)$.

\emph{The oscillation error.} By \eqref{8.11:eq-deep-curvature-blocks-5},
$(\operatorname{osc}W)^2\lesssim n^2\rho^{-2(d+1)}$. Hence the square of the coefficient of the
Cauchy--Schwarz inequality for $\sum\operatorname{osc}W\,\|\Phi\|$ is, in (I) and (II)
respectively,
\begin{align*}
\sum_{\mathfrak{b}}n^2\rho^{-2(d+1)}\,t_cn&\lesssim t_cn^3\,(hn)^{-7}=t_ch^{-7}\eta^4,\\
\sum_{\mathfrak{b}_{II}}n^2\rho^{-2(d+1)}(\tau+1)^2
&\lesssim(t+1)^2n^4\,(hn)^{-8}\lesssim t^2h^{-8}\eta^4,
\end{align*}
where $t+1\le2t$ because $t\ge t_c=4w+1\ge1$. The two coefficients are thus
$\lesssim(t_ch^{-7})^{1/2}\eta^2$ and $\lesssim t\,h^{-4}\eta^2$.

\emph{The rim error.} By the main-term bound and $h\ge4$, $|W|+\operatorname{osc}W\le
C_\Pi(1+C_o/4)\rho^{-d}$, so the rim term of Lemma~\ref{8.11:lem-weighted-flux} is
$\lesssim\sum_{\partial\mathcal{R}}\rho^{-d}\|\Phi\|$. Since every block site $y$ with
$\operatorname{supp}\omega_y\subset\mathcal{R}$ lies in $\mathcal{Y}$, the rim
$\partial\mathcal{R}$ lies within $2n$ of the boundary of $\mathcal{R}$: the weights $\omega_y$
are tents of width less than $2n$ in the directions $\alpha,\beta$, and in the other directions
$\mathcal{R}$ may be misaligned with the blocks. In (I) the rim consists of the two normal
end-slabs, at most $4n$ normal positions per bond, and of the lateral rims, the face positions
within $2n$ of the boundary of the set of positions, over the full normal range $t_cn$; in (II)
it consists of the perimeter strips of the rectangle, at most $8(\tau+1)n$ sites per edge
position. The squares of the Cauchy--Schwarz coefficients are
\begin{align*}
\sum_{\mathfrak{b}}\rho^{-2d}\cdot4n&\lesssim n\,(hn)^{-5}=h^{-5}\eta^4,\\
2n\,h_z^{-(d+2)}\cdot t_cn&\lesssim t_cn^2(hn)^{-6}=t_ch^{-6}\eta^4,\\
\sum_{\mathfrak{b}_{II}}\rho^{-2d}\cdot8(\tau+1)n&\lesssim(t+1)n^2(hn)^{-6}\lesssim t\,h^{-6}\eta^4,
\end{align*}
so the coefficients are $\lesssim h^{-5/2}\eta^2$, $\lesssim(t_ch^{-6})^{1/2}\eta^2$ and
$\lesssim(t\,h^{-6})^{1/2}\eta^2$.

\emph{Summation of the errors.} Each of the five coefficients is multiplied by the square root of
three times the right side of \eqref{8.11:eq-deep-curvature-blocks-6}, that is by
$(120\,N\,C_E\,\hat N_W)^{1/2}\varepsilon_1$. A sum of five square roots is at most
$\sqrt5$ times the square root of the sum, and since $t_c\le t$ and, for $h\ge1$,
$t_ch^{-7}\le t\,h^{-6}$,
\begin{equation*}
t_ch^{-7}+t^2h^{-8}+h^{-5}+t_ch^{-6}+t\,h^{-6}\le h^{-5}+3\,t\,h^{-6}+t^2h^{-8}.
\end{equation*}
Therefore the sum of all error terms over all groups is at most
\begin{equation*}
C_{\mathrm{blk}}\bigl[NC_E\hat N_W(h^{-5}+t\,h^{-6}+t^2h^{-8})\bigr]^{1/2}\varepsilon_1\eta^2
=\kappa_{\mathrm{blk}}(h)\,\varepsilon_{k_o}\eta^2
\end{equation*}
for a suitable $C_{\mathrm{blk}}$ depending
on $d$ and $\varrho_H$ only, by \eqref{8.11:eq-deep-curvature-blocks-1}.

\emph{Conclusion.} Inserting the bounds for $\mathsf{F}_c$ into the proof of
Theorem~\ref{8.11:thm-deep-curvature} in place of its pointwise bound by $\theta''$ on
$\mathcal{L}_c$, and keeping all other bounds of that proof, including the reduction from
$\hat x$ to $\check x$ with its error $2\vartheta_t$ and the final comparison
$\|\check U(\partial p_0)-1\|\le\|\Phi(p_0)\|$, gives \eqref{8.11:eq-deep-curvature-a} with the
two members \eqref{8.11:eq-deep-curvature-blocks-2} and with the additional error
$\kappa_{\mathrm{blk}}(h)$.
\end{proof}

\begin{remark}
Theorem~\ref{8.11:thm-deep-curvature-blocks} is used under the ceiling
\begin{equation}\label{8.11:eq-deep-curvature-blocks-a}
\kappa_{\mathrm{blk}}(h)\le1.
\end{equation}
For $h$ and $t$ comparable to $w$ it reads $C\,w^{-1}\le1$, with $C$ depending on
$d,N,L,M/M_2$ and on the further constants on which the corresponding form of the ceiling
$\kappa_{\mathrm{col}}(h)\le1$ of \eqref{8.11:eq-deep-bound-scope-a} depends. Comparing the
definitions \eqref{8.11:eq-deep-curvature-blocks-1} and \eqref{8.11:eq-bound-members-a}, for $h$
and $t$ comparable to $w$ the quantity $\kappa_{\mathrm{blk}}(h)$ is at most a constant multiple of
$\kappa_{\mathrm{col}}(h)$, since $r_1+1\ge1$; so \eqref{8.11:eq-deep-curvature-blocks-a} is
implied by the ceiling $\kappa_{\mathrm{col}}(h)\le1$ up to a constant.

As a consequence, the first and the third of the corollaries of
Theorem~\ref{8.11:thm-deep-curvature} that follow hold with the constant $C_1^{\flat}$ doubled in
its first two members, and without the comparison of the global minimiser with the straddling local
minimisers and without its justification at points critical along the fibres of the top averages;
the second corollary likewise, in whose setting that comparison is justified in any case. The fourth corollary,
with reference depth $c_*=w-1$, still reads the bound pointwise on $\Sigma_{w-1}$; at $\tilde x$ it
holds as stated, while at $\hat x$ it rests on the justification of the comparison at points
critical along the fibres. It is applied at $\hat x$ nowhere in this book: the gap analysis in
which it enters can be run on Theorem~\ref{8.11:thm-deep-curvature-blocks}, where the members
$w/h$ and $(w/h)^2$ must be at most a fixed fraction of the young cap $\mathsf{T}$, the fraction
being determined by the contraction constant and the constant of that analysis, and this holds as
soon as $h$ is at least a constant multiple of $w/L$, the constant being that of the gap analysis.
\end{remark}

In what follows, the height $h$ enters the definitions \eqref{8.11:eq-bound-members-a} of
$\mathfrak{B}(h)$, $\kappa_{\mathrm{col}}(h)$ and $\vartheta_{\mathrm{nl}}(h)$ only through the factors
$h^{-1}$, $h^{-2}$, $h^{-5/2}$ and $\log^{+}(D_W/h)$. Each of these factors is non-increasing in $h$,
and their coefficients are non-negative. Hence $\mathfrak{B}$, $\kappa_{\mathrm{col}}$ and
$\vartheta_{\mathrm{nl}}$ are non-increasing in $h$, and a bound at height $h$ holds at every greater
depth.

\begin{corollary}[Logarithmic bound beyond twice the clamp width]\label{8.11:cor-logarithmic-bound}
Assume the hypotheses of Theorem~\ref{8.11:thm-deep-curvature} and the condition
\eqref{8.11:eq-inherited-bounds-b}. With $\sigma_e$ as in Remark~\ref{8.11:rem-inherited-bounds} and
$\theta'=\theta(1+\vartheta_1)$ as in Notation~\ref{8.11:not-units-cap-amplitude}, put
\begin{equation*}
s_e^{+}:=\frac{\sigma_e}{\theta'}=\frac{s_e\theta+\vartheta_1}{\theta(1+\vartheta_1)}\le s_e+\vartheta_1 ,
\end{equation*}
where the inequality holds because $\sigma_e=s_e\theta+\vartheta_1$, $\theta\ge1$ and $\vartheta_1\ge0$.
Let $\mathfrak{Y}$ be the first young row, with the depth band of
Definition~\ref{8.11:def-young-row-cover}, display \eqref{8.11:eq-young-row-cover-a}, taken with
$y_+=w+5w/L$. Assume
\begin{equation}\label{8.11:eq-logarithmic-bound-a}
\frac{5}{L}+\frac{r_1}{w}\le\frac12 ,
\end{equation}
together with the two ceilings of \eqref{8.11:eq-deep-bound-scope-a} at height $w/2$, namely
$\kappa_{\mathrm{col}}(w/2)\le1$ and $\vartheta_{\mathrm{nl}}(w/2)\le\vartheta_1$. Then every fine
plaquette $p$ with $\delta(p)\ge 2w$ obeys the bound below. This covers every fine plaquette of $H$
at depth at least $2w$, the inner rings and the core included:
\begin{equation}\label{8.11:eq-logarithmic-bound-1}
\|U(e\oplus\hat x)(\partial p)-1\|
\le\Bigl[C_1^{\flat}+C_2^{\flat}\,\frac{\mathsf{T}}{L}
+C_3^{\flat}\,s_e^{+}\log\frac{D_W}{w}\Bigr]\,\theta(1+\vartheta_1)\,\varepsilon_{k_o}\eta^2,
\qquad \frac{\mathsf{T}}{L}\le L,
\end{equation}
where
\begin{equation}\label{8.11:eq-logarithmic-bound-b}
\begin{aligned}
C_1^{\flat}&:=C_\Pi\bigl(18c_1+484c_2+c_3+10c_4+100c_5+2dc_1\log 2\bigr)(1+2\vartheta_t)+2,\\
C_2^{\flat}&:=C_\Pi(20c_1+120c_2)(1+2\vartheta_t),\qquad
C_3^{\flat}:=2dC_\Pi c_1 .
\end{aligned}
\end{equation}
These constants are functions of $(d,\varrho_H)$ only.
\end{corollary}

\begin{proof}
Fix a fine plaquette $p$ with $\delta(p)\ge 2w$ and put $h:=\delta(p)-d_1$. By
\eqref{8.11:eq-released-problem-a}, $d_1=y_++r_1=w+5w/L+r_1$. Hence, by
\eqref{8.11:eq-logarithmic-bound-a},
\begin{equation*}
h\ge 2w-w-\frac{5w}{L}-r_1=w\Bigl(1-\frac5L-\frac{r_1}{w}\Bigr)\ge\frac w2\ge r_1 .
\end{equation*}
The last inequality holds because \eqref{8.11:eq-logarithmic-bound-a} gives $r_1/w\le 1/2$.
Moreover $hn\ge4$, with $n$ as in Theorem~\ref{8.11:thm-deep-curvature}. Indeed, the clamp width
satisfies $w\ge8$, as stated in Theorem~\ref{8.11:thm-deep-curvature-blocks}, so $h\ge w/2\ge4$;
and $n\ge1$.

We apply Theorem~\ref{8.11:thm-deep-curvature} at this $p$ and this $h$ with the reference depth
$c_*:=-(3w+1)$. For this choice $\sigma_*=\sigma_e$, $t_c=w-c_*$ and $t=d_1-c_*$, and
\eqref{8.11:eq-deep-curvature-a} gives
\begin{equation*}
\|U(e\oplus\hat x)(\partial p)-1\|
\le\bigl[\mathfrak{B}(h)+\kappa_{\mathrm{col}}(h)+\vartheta_{\mathrm{nl}}(h)+2\vartheta_t\bigr]
\varepsilon_{k_o}\eta^2 .
\end{equation*}

\emph{The thicknesses.} Since $w\ge8$, in particular $w\ge2$, we have $t_c=4w+1\le 9w/2$. By
\eqref{8.11:eq-logarithmic-bound-a}, $5w/L+r_1\le w/2$, so
\begin{equation*}
t=4w+1+\frac{5w}{L}+r_1\le 4w+1+\frac w2\le 5w,\qquad t+1\le\frac{11w}{2},
\end{equation*}
again using $w\ge2$. Further, $\varpi_y=y_+-w=5w/L$, and $2\varpi_y+2r_1+1\le w+1\le w+w/2$. For every
height $h'\ge w/2$, in particular for $h'=w/2$, these give
\begin{gather*}
\frac{t_c}{h'}\le 9,\qquad \frac{(t+1)^2}{h'^2}\le 121,\qquad \frac{t}{h'}\le 10,\qquad
\frac{\varpi_y}{h'}\le\frac{10}{L},\\
\frac{\varpi_y(2\varpi_y+2r_1+1)}{h'^2}\le\frac{10}{L}\cdot\frac{w+w/2}{w/2}=\frac{30}{L}.
\end{gather*}
The zone width satisfies $D_W\ge 8L(M/M_2)w\ge w$, as noted in
Notation~\ref{8.11:not-bound-members}, where the constants of \eqref{8.11:eq-bound-members-a} are
fixed. Hence
\begin{equation*}
\log^{+}\frac{D_W}{h'}\le\log^{+}\frac{2D_W}{w}=\log2+\log\frac{D_W}{w}.
\end{equation*}

\emph{The member $\mathfrak{B}$.} By monotonicity, $\mathfrak{B}(h)\le\mathfrak{B}(w/2)$. Inserting the
bounds above into the definition \eqref{8.11:eq-bound-members-a} at height $w/2$ gives
\begin{align*}
\mathfrak{B}(w/2)\le C_\Pi\Bigl\{&(18c_1+484c_2)\theta''+(20c_1+120c_2)\frac{\theta_y''}{L}\\
&+\sigma_e\Bigl[c_3+2dc_1\log2+10c_4+100c_5+2dc_1\log\frac{D_W}{w}\Bigr]\Bigr\}.
\end{align*}
By \eqref{8.11:eq-inherited-bounds-a}, $\theta''=\theta'+2\vartheta_t\le\theta'(1+2\vartheta_t)$,
because $\theta'\ge1$; likewise $\theta_y''=\theta'\mathsf{T}+2\vartheta_t\le\theta'\mathsf{T}(1+2\vartheta_t)$,
because $\theta'\mathsf{T}\ge\mathsf{T}\ge1$ by \eqref{8.11:eq-units-cap-amplitude-a}. In the
$\sigma_e$-member we use two bounds. For the constant part, $\theta s_e\le 1$ by
Notation~\ref{8.11:not-units-cap-amplitude} and $\theta\ge1$, so
\begin{equation*}
\sigma_e=\theta s_e+\vartheta_1\le 1+\vartheta_1\le\theta'\le\theta'(1+2\vartheta_t).
\end{equation*}
For the logarithmic part we write $\sigma_e=\theta' s_e^{+}$ and note that $\log(D_W/w)\ge0$, again
because $D_W\ge w$. Together,
\begin{equation*}
\mathfrak{B}(h)\le\theta'\Bigl[C_1^{\flat}-2+C_2^{\flat}\,\frac{\mathsf{T}}{L}
+C_3^{\flat}\,s_e^{+}\log\frac{D_W}{w}\Bigr],
\end{equation*}
with the constants \eqref{8.11:eq-logarithmic-bound-b}.

\emph{The remaining members.} By monotonicity and the two ceilings assumed at height $w/2$,
$\kappa_{\mathrm{col}}(h)\le\kappa_{\mathrm{col}}(w/2)\le1$ and
$\vartheta_{\mathrm{nl}}(h)\le\vartheta_{\mathrm{nl}}(w/2)\le\vartheta_1$. The relative errors
$\vartheta_1$ and $\vartheta_t$ are small under the hypotheses of
Theorem~\ref{8.11:thm-deep-curvature}; in particular $\vartheta_1+2\vartheta_t\le1$. Hence
\begin{equation*}
\kappa_{\mathrm{col}}(h)+\vartheta_{\mathrm{nl}}(h)+2\vartheta_t\le 2\le 2\theta',
\end{equation*}
the last inequality since $\theta'\ge1$. This supplies the summand $2$ in $C_1^{\flat}$. Adding the
two estimates and recalling $\theta'=\theta(1+\vartheta_1)$ gives the bound in
\eqref{8.11:eq-logarithmic-bound-1}.

Finally, $\mathsf{T}\le L^2$ by \eqref{8.11:eq-units-cap-amplitude-a}, so $\mathsf{T}/L\le L$.
\end{proof}

\begin{remark}[Form of the bound beyond twice the clamp width]
The bracket in \eqref{8.11:eq-logarithmic-bound-1} contains no cutoff quantity. By
Notation~\ref{8.11:not-bound-members}, $D_W/w\ge 8L(M/M_2)$; by
Remark~\ref{8.11:rem-deep-bound-scope}, which takes it from one of the inputs of
Theorem~\ref{8.11:thm-deep-curvature}, $D_W/w$ is bounded above by $17L(M/M_2)$ times a factor free
of the cutoff. Consequently neither
$\check R$ as a function of $g$, nor $\eta$, nor the number $K$ of renormalisation steps enters the
bracket, and the bracket is $o(L^2)\,(1+\log(M/M_2))$ as $L\to\infty$.

Compare \eqref{8.11:eq-logarithmic-bound-1} with a bound at depth $2w$ of the form
\begin{equation*}
C(d,N)\Bigl[1+s_e\log\frac{D_W}{w}\Bigr]\theta(1+\vartheta_1)\,\varepsilon_{k_o}\eta^2 .
\end{equation*}
The logarithmic member of \eqref{8.11:eq-logarithmic-bound-1} has this form: its coefficient is
$C_3^{\flat}s_e^{+}\le C_3^{\flat}(s_e+\vartheta_1)$, linear in $s_e$ up to the relative error
$\vartheta_1$, with $C_3^{\flat}$ depending on $(d,\varrho_H)$ only.
The additive member is $C_1^{\flat}+C_2^{\flat}\mathsf{T}/L\le C_1^{\flat}+C_2^{\flat}L$, in place of a
constant $C(d,N)$.

Theorem~\ref{8.11:thm-deep-curvature-blocks} applies at heights $h\ge w/2\ge4$. Under its ceiling
\eqref{8.11:eq-deep-curvature-blocks-a} at height $w/2$, the argument above, run with that theorem,
gives a bound of the form \eqref{8.11:eq-logarithmic-bound-1} in which the members $18c_1$ and
$484c_2$ of $C_1^{\flat}$ carry the factor $2(1+2C_o/w)$, with $C_o$ as in
Theorem~\ref{8.11:thm-deep-curvature-blocks}, in place of $(1+2\vartheta_t)$, and in which the
summand $2$ of $C_1^{\flat}$ is replaced by $3$.
\end{remark}

\begin{corollary}[The bound without young windows]\label{8.11:cor-no-young-windows}
Assume the hypotheses under which the curvature bound \eqref{8.11:eq-deep-curvature-a} deep
inside the hole is proved, and the bounds
\eqref{8.11:eq-inherited-bounds-b} inherited by the released minimiser. Suppose further
that:
\begin{itemize}
\item the young family is empty, $\mathfrak{Y}=\emptyset$;
\item the value-Lipschitz property $(\mathrm{VL})_{\theta}[\vartheta,\Lambda]$ of
\eqref{8.11:eq-value-lipschitz-a} holds for the family of straddling sites
$\mathfrak{S}$;
\item $r_1\le w/2$;
\item the two conditions of \eqref{8.11:eq-deep-bound-scope-a} hold at reading height
$w/2$.
\end{itemize}
Let $\tilde x$ be the straddling-only minimiser. Then every fine plaquette $p$ with
$\delta(p)\ge 2w$ obeys
\begin{equation}\label{8.11:eq-no-young-windows-1}
\bigl\|U(e\oplus\tilde x)(\partial p)-1\bigr\|
\le\Bigl[C_1^{\flat}+C_3^{\flat}\,s_e^{+}\log\frac{D_W}{w}\Bigr]\,
\theta(1+\vartheta_1)\,\varepsilon_{k_o}\eta^2 .
\end{equation}
Here $C_1^{\flat}$ and $C_3^{\flat}$ are the constants of
\eqref{8.11:eq-logarithmic-bound-b}.
\end{corollary}

\begin{proof}
We run the proof of Corollary~\ref{8.11:cor-logarithmic-bound} with the young family
empty. The point is then the straddling-only minimiser $\tilde x$, and the value-Lipschitz
property is used for the straddling family. Since $\mathfrak{Y}=\emptyset$ there is no
young row, and the proof of Corollary~\ref{8.11:cor-logarithmic-bound} is run with
$\varpi_y=0$ in place of $\varpi_y=5w/L$.

Put $h:=\delta(p)-d_1$. In the proof of
Corollary~\ref{8.11:cor-logarithmic-bound} the lower bound for $h$ reads
$h\ge 2w-w-\varpi_y-r_1$. With $\varpi_y=0$ and the hypothesis $r_1\le w/2$ this gives
\begin{equation*}
h\ \ge\ 2w-w-r_1\ \ge\ \frac w2\ \ge\ r_1 .
\end{equation*}
Here $r_1\le w/2$ does the work that the geometric condition of
Corollary~\ref{8.11:cor-logarithmic-bound} did there.

The quantities $\mathfrak{B}(h)$, $\kappa_{\mathrm{col}}(h)$ and
$\vartheta_{\mathrm{nl}}(h)$ are non-increasing in $h$. A bound at height $w/2$ therefore
holds at height $h\ge w/2$. By hypothesis, the conditions of
\eqref{8.11:eq-deep-bound-scope-a} hold at height $w/2$.

With $c_*=-(3w+1)$ the thicknesses are $t_c=w-c_*$ and $t=d_1-c_*$. Since $w\ge 2$,
\begin{equation*}
t_c\ =\ 4w+1\ \le\ \frac{9w}{2},
\end{equation*}
and since $d_1\le w+r_1\le 3w/2$ (here $\varpi_y=0$ and $r_1\le w/2$),
\begin{equation*}
t\ \le\ 4w+1+\frac w2\ \le\ 5w,\qquad t+1\ \le\ \frac{11w}{2}.
\end{equation*}
Together with $h\ge w/2$ these give
\begin{align*}
&\frac{t_c}{h}\le 9, \qquad \frac{(t+1)^2}{h^2}\le 121, \qquad \frac{t}{h}\le 10,\\
&\log^{+}\frac{D_W}{h}\le \log 2+\log\frac{D_W}{w}.
\end{align*}
The ratios $\varpi_y/h$ and $\varpi_y(2\varpi_y+2r_1+1)/h^2$ are now zero. The two members
proportional to $\varpi_y$, which in Corollary~\ref{8.11:cor-logarithmic-bound} produced
the term $C_2^{\flat}\,\mathsf{T}/L$, therefore vanish.

We insert these estimates in $\mathfrak{B}(w/2)$ and use $\theta''\le\theta'(1+2\vartheta_t)$.
We also use
\begin{equation*}
\kappa_{\mathrm{col}}+\vartheta_{\mathrm{nl}}+2\vartheta_t\ \le\ 2\ \le\ 2\theta' .
\end{equation*}
This gives \eqref{8.11:eq-no-young-windows-1}, with the constants $C_1^{\flat}$ and
$C_3^{\flat}$ of \eqref{8.11:eq-logarithmic-bound-b}.
\end{proof}

In the following corollary, $g_{\square'}$ is the window functional of the young cube $\square'$, and
$\widetilde{\mathcal{W}}^{\theta}$ is the class constrained by the straddling windows alone at threshold
$\theta$. We write $\tilde m(\theta)$ for the minimum of $\mathbf{A}$ over $\widetilde{\mathcal{W}}^{\theta}$,
and a straddling-only minimiser $\tilde x$ is a point of $\widetilde{\mathcal{W}}^{\theta}$ at which
this minimum is attained. The objects $\widehat{\mathcal{W}}^{\theta}$, $\hat m(\theta)$ and $\hat x$
are the corresponding objects when the young windows are imposed as well.

\begin{corollary}[Young windows inactive under a scale jump]\label{8.11:cor-scale-jump}
Assume the hypotheses of Corollary~\ref{8.11:cor-no-young-windows}. Assume moreover that the radii
at the two levels under consideration satisfy $R_{k_o-1}=L\,R_{k_o}$. Assume also that the first young
row (Definition~\ref{8.11:def-young-row-cover}) tests the depth band $\{3w\le\delta\le(L+5)w\}$, that
is, the test region $\tilde{\square}'$ of every cube $\square'$ of that row lies in this band. Let
$\rho$ be a number, used in this corollary only and unrelated to the contraction letter of
clause~(v-d), such that
\begin{equation}\label{8.11:eq-scale-jump-a}
  \bigl[C_1^{\flat}+C_3^{\flat}\,s_e^{+}\log(D_W/w)\bigr](1+\vartheta_1)+\vartheta_1\mathsf{T}
  \le(1-\rho)\,\mathsf{T}.
\end{equation}
Then every straddling-only minimiser $\tilde x$ passes the first young row with room $\rho\theta$,
that is, $g_{\square'}(\tilde x)\le\theta-\rho\theta$ for every cube $\square'$ of that row.
Moreover $\hat m(\theta)=\tilde m(\theta)$, and every constrained minimiser $\hat x$ is a
straddling-only minimiser. Finally, the bound of Corollary~\ref{8.11:cor-no-young-windows} holds at
$\hat x$: for every fine plaquette $p$ with $\delta(p)\ge 2w$,
\begin{equation}\label{8.11:eq-scale-jump-1}
  \|U(e\oplus\hat x)(\partial p)-1\|
  \le\bigl[C_1^{\flat}+C_3^{\flat}\,s_e^{+}\log(D_W/w)\bigr]\,
  \theta(1+\vartheta_1)\,\varepsilon_{k_o}\eta^2 .
\end{equation}
\end{corollary}

\begin{proof}
Fix a straddling-only minimiser $\tilde x$ and a cube $\square'$ of the first young row. The bound
\cite[(1.19)]{B-CMP119}, in the form used in this chapter for the young-window functional, applied at
$\tilde x$, gives
\begin{equation}\label{8.11:eq-scale-jump-2}
  g_{\square'}(\tilde x)
  \le\frac{\sup_{p\subset\tilde{\square}'}\|U(e\oplus\tilde x)(\partial p)-1\|}
          {\varepsilon_{k_o-1}L^2\eta^2}+\vartheta_1 .
\end{equation}
By hypothesis $\tilde{\square}'\subset\{\delta\ge 3w\}$, and $\{\delta\ge 3w\}\subset\{\delta\ge 2w\}$.
Hence every fine plaquette $p\subset\tilde{\square}'$ satisfies $\delta(p)\ge 2w$.
Corollary~\ref{8.11:cor-no-young-windows}, whose hypotheses are assumed, therefore bounds the numerator
in \eqref{8.11:eq-scale-jump-2} by
\begin{equation*}
  \bigl[C_1^{\flat}+C_3^{\flat}\,s_e^{+}\log(D_W/w)\bigr]\,
  \theta(1+\vartheta_1)\,\varepsilon_{k_o}\eta^2 .
\end{equation*}
By the definition of the young cap $\mathsf{T}$, the denominator is
$\varepsilon_{k_o-1}L^2\eta^2=\mathsf{T}\,\varepsilon_{k_o}\eta^2$. This is the cap at the scale
$R_{k_o-1}=L\,R_{k_o}$, expressed in the exterior unit $\varepsilon_{k_o}\eta^2$.
By \eqref{8.11:eq-scale-jump-a}, the right-hand side of \eqref{8.11:eq-scale-jump-2} is therefore
at most $\theta(1-\rho)$. Since $\square'$ was an arbitrary cube of the first young row,
$\tilde x$ passes that row with room $\rho\theta$.

Consequently $\tilde x$ satisfies the young windows in addition to the straddling windows, so
$\tilde x\in\widehat{\mathcal{W}}^{\theta}$. The class $\widehat{\mathcal{W}}^{\theta}$ is obtained
from $\widetilde{\mathcal{W}}^{\theta}$ by imposing further constraints, so
$\widehat{\mathcal{W}}^{\theta}\subset\widetilde{\mathcal{W}}^{\theta}$. A minimum over a smaller set
is not below the minimum over the larger one, so $\hat m(\theta)\ge\tilde m(\theta)$. On the other hand,
$\tilde x\in\widehat{\mathcal{W}}^{\theta}$ gives $\hat m(\theta)\le\mathbf{A}(\tilde x)=\tilde m(\theta)$.
Thus the minima coincide.

Let $\hat x$ be a constrained minimiser. Then $\hat x\in\widehat{\mathcal{W}}^{\theta}\subset
\widetilde{\mathcal{W}}^{\theta}$ and $\mathbf{A}(\hat x)=\hat m(\theta)=\tilde m(\theta)$. A minimiser
over the smaller set that attains the larger set's minimum minimises over the larger set, so
$\hat x$ is a straddling-only minimiser. Corollary~\ref{8.11:cor-no-young-windows} applies to every
straddling-only minimiser, hence to $\hat x$, and this gives \eqref{8.11:eq-scale-jump-1}.
\end{proof}

\begin{remark}[The cube-aligned young row]\label{8.11:rem-cube-aligned}
The covering hypothesis \eqref{8.11:eq-young-row-cover-a} on the first young row is a displayed
hypothesis. Corollary~\ref{8.11:cor-logarithmic-bound} uses it with $y_+=w+5w/L$. Here we display a
second arrangement, without deciding which arrangement the young family $\mathfrak{Y}$ realises. In the
\emph{cube-aligned} arrangement, the cubes of level $k_o-1$ tile the zone $H$ flush with its boundary
$\partial H$. In this arrangement the following hold.
\begin{itemize}
\item[(a)] The first row of young cubes $\square'$ whose dependence regions $\square'^{\sim 4}$ lie in
$H$ is the band $\{w\le\delta\le 2w\}$. Its test band is $\{w-w/L\le\delta\le 2w+w/L\}$.
Consequently $y_+=2w+w/L$ and $\varpi_y=w(1+1/L)$.
\item[(b)] The free depth is $d_1=2w+w/L+r_1$. For every $h$ admitted by
Theorem~\ref{8.11:thm-deep-curvature} and every fine plaquette $p$ with $\delta(p)\ge d_1+h$, the
right-hand side of \eqref{8.11:eq-deep-curvature-a} is, in the exterior unit $\varepsilon_{k_o}\eta^2$,
of the form
\begin{equation}\label{8.11:eq-cube-aligned-1}
C\Bigl[1+\frac{w}{h}\Bigl(1+\frac{w}{h}\Bigr)\mathsf{T}+s_e^{+}\log\frac{D_W}{h}\Bigr]\theta' ,
\end{equation}
where $\theta'=\theta(1+\vartheta_1)$ and $C$ is a constant.
\item[(c)] At $h\asymp w$ the young member $(w/h)(1+w/h)\mathsf{T}$ of the bracket in
\eqref{8.11:eq-cube-aligned-1} is of order $L^2$; it is not $o(L^2)$.
\end{itemize}
In the cube-aligned arrangement the arguments of this section do not establish the bound of
Corollary~\ref{8.11:cor-logarithmic-bound} at $\hat x$: they yield \eqref{8.11:eq-cube-aligned-1}, by
(b), and the corollary of Theorem~\ref{8.11:thm-deep-curvature} for the straddling-only minimiser
$\tilde x$, and neither of these is that bound.
\end{remark}

\begin{proof}
The cubes of level $k_o-1$ tile $H$ flush with $\partial H$, so their rows are bands between level sets
of the signed depth $\delta$, and the first row and its test band are the bands named in (a). Hence the
outer inclusion of \eqref{8.11:eq-young-row-cover-a} holds with $y_+=2w+w/L$, and the definition
$\varpi_y=y_+-w$ gives
\begin{equation*}
\varpi_y = w+\frac{w}{L} = w\Bigl(1+\frac1L\Bigr).
\end{equation*}
This proves (a).

For (b), we apply Theorem~\ref{8.11:thm-deep-curvature} with this value of $y_+$. The free depth
becomes $d_1=2w+w/L+r_1$. For every fine plaquette
$p$ with $\delta(p)\ge d_1+h$, the right-hand side of \eqref{8.11:eq-deep-curvature-a} then takes the
form \eqref{8.11:eq-cube-aligned-1}.

For (c), let $h\asymp w$. Then the factor $(w/h)(1+w/h)$ is bounded above and below by constants, so the
young member $(w/h)(1+w/h)\mathsf{T}$ is of the order of $\mathsf{T}$, that is, of order $L^2$; in
particular it is not $o(L^2)$.
\end{proof}

\begin{corollary}[The bound at small reading heights]\label{8.11:cor-small-heights}
Let $e$, $\theta$ and $\hat x$ be as in Theorem~\ref{8.11:thm-deep-curvature}, under the
hypotheses of that theorem. Let $h$ satisfy
$r_1+2\le h\le w$, and assume the ceilings $(\gamma\text{-col})_h$ and
$(\gamma\text{-nl})_h$ of \eqref{8.11:eq-deep-bound-scope-a}. Then every fine plaquette $p$
with $\delta(p)\ge d_1+h$ satisfies
\begin{equation}\label{8.11:eq-small-heights-a}
\begin{split}
\|U(e\oplus\hat x)(\partial p)-1\|
&\le C_4^{\flat}\Bigl\{1+\log\frac{D_W}{h}
+\Bigl[\frac{\varpi_y}{h}+\Bigl(\frac{\varpi_y}{h}\Bigr)^{2}\Bigr]\mathsf{T}\Bigr\}\\
&\qquad\times\theta(1+\vartheta_1)\,\varepsilon_{k_o}\eta^{2},
\end{split}
\end{equation}
where $C_4^{\flat}=C(d,\varrho_H)$ depends on $d$ and $\varrho_H$ only.
\end{corollary}

\begin{proof}
We apply Theorem~\ref{8.11:thm-deep-curvature} with the reference depth $c_*=w-1$.
We write $\theta'=\theta(1+\vartheta_1)$ and $\theta_y=\theta'\mathsf{T}$.

\emph{The parameters at $c_*=w-1$.} The thicknesses of Theorem~\ref{8.11:thm-deep-curvature}
are $t_c=w-c_*$ and $t=d_1-c_*$. At $c_*=w-1$ we get $t_c=1$, and, since $d_1=y_++r_1$ by
\eqref{8.11:eq-released-problem-a} and $\varpi_y=y_+-w$,
\begin{equation*}
t=y_++r_1-w+1=\varpi_y+r_1+1 .
\end{equation*}
As $h\ge r_1+2$, this gives $t+1=\varpi_y+r_1+2\le\varpi_y+h$. The amplitude attached to this
reference depth is $\sigma_*=\theta''$.

\emph{Application of the deep curvature bound.} Since $h\ge r_1+2>r_1$, the bound
\eqref{8.11:eq-deep-curvature-a} of Theorem~\ref{8.11:thm-deep-curvature} applies at height $h$
with $c_*=w-1$. It gives
\begin{equation*}
\|U(e\oplus\hat x)(\partial p)-1\|
\le\bigl[\mathfrak{B}(h)+\kappa_{\mathrm{col}}(h)+\vartheta_{\mathrm{nl}}(h)+2\vartheta_t\bigr]
\varepsilon_{k_o}\eta^{2}.
\end{equation*}
By \eqref{8.11:eq-bound-members-a}, with $t_c=1$ and $\sigma_*=\theta''$, the first member is
\begin{equation*}
\begin{split}
\mathfrak{B}(h)=C_\Pi\Bigl\{&\frac{2c_1[\theta''+\varpi_y\theta_y'']}{h}
+\frac{4c_2[(t+1)^2\theta''+\varpi_y(2\varpi_y+2r_1+1)\theta_y'']}{h^2}\\
&+\theta''\Bigl[c_3+2dc_1\log^+\frac{D_W}{h}+c_4\frac{t}{h}
+c_5\Bigl(\frac{t}{h}\Bigr)^{2}\Bigr]\Bigr\}.
\end{split}
\end{equation*}
The constants $\theta''$ and $\theta_y''$ attached to the released minimiser $\check x$ of
\eqref{8.11:eq-released-problem-a} satisfy $\theta''\le\theta'$ and $\theta_y''\le\theta_y$, by
part~(iii) of the statement that introduces the released problem together with that display.

\emph{Elementary inequalities.} We use four facts.
\begin{enumerate}
\item From $t+1\le\varpi_y+h$ we get $t/h\le(t+1)/h\le1+\varpi_y/h$. Since
$(a+b)^2\le2a^2+2b^2$, this gives
\begin{equation*}
\Bigl(\frac{t}{h}\Bigr)^{2}\le\frac{(t+1)^2}{h^2}\le2+2\Bigl(\frac{\varpi_y}{h}\Bigr)^{2}.
\end{equation*}
\item Since $2r_1+1\le2h$,
\begin{equation*}
\frac{\varpi_y(2\varpi_y+2r_1+1)}{h^2}\le2\Bigl(\frac{\varpi_y}{h}\Bigr)^{2}
+2\,\frac{\varpi_y}{h}.
\end{equation*}
\item $h\ge r_1+2\ge2$, so $2/h\le1$.
\item In the set-up of Theorem~\ref{8.11:thm-deep-curvature} the aspect ratio of the box
$\Omega_{d_1}$ is bounded by $2\varrho_H$ through the inequality $4\varrho_Hd_1\le D_W$.
Since $d_1=w+\varpi_y+r_1$ with $\varpi_y\ge0$ and $r_1\ge0$, we have $d_1\ge w$; as
$\varrho_H\ge1$ and $h\le w$,
\begin{equation*}
D_W\ge4\varrho_Hd_1\ge4w\ge h,
\end{equation*}
so that $\log^+(D_W/h)=\log(D_W/h)\ge0$.
\end{enumerate}

\emph{The first member.} We insert these facts and $\theta''\le\theta'$,
$\theta_y''\le\theta_y$ into $\mathfrak{B}(h)$ and sort the terms into three groups.

The terms carrying $\theta''$ and no power of $\varpi_y$ are bounded by
\begin{equation*}
\theta'\bigl[c_1+8c_2+c_3+c_4+2c_5+2dc_1\log(D_W/h)\bigr].
\end{equation*}

The terms carrying $\theta''$ and a power of $\varpi_y/h$ come from the bounds on
$(t+1)^2/h^2$, $t/h$ and $(t/h)^2$. They sum to
\begin{equation*}
\theta'\bigl[c_4\varpi_y/h+(8c_2+2c_5)(\varpi_y/h)^2\bigr].
\end{equation*}
Since $\mathsf{T}\ge1$ we have $\theta'\le\theta_y$, so this $\varpi_y$-part is dominated by
the young member:
\begin{equation*}
\theta'\bigl[c_4\tfrac{\varpi_y}{h}+(8c_2+2c_5)\bigl(\tfrac{\varpi_y}{h}\bigr)^2\bigr]
\le\theta'\mathsf{T}\bigl[c_4\tfrac{\varpi_y}{h}
+(8c_2+2c_5)\bigl(\tfrac{\varpi_y}{h}\bigr)^2\bigr].
\end{equation*}

The terms carrying $\theta_y''$ are bounded, by the second fact, by
\begin{equation*}
\theta'\mathsf{T}\bigl[(2c_1+8c_2)\varpi_y/h+8c_2(\varpi_y/h)^2\bigr].
\end{equation*}

Adding the three groups,
\begin{equation*}
\begin{split}
\mathfrak{B}(h)\le C_\Pi\,\theta'\Bigl\{&c_1+8c_2+c_3+c_4+2c_5+2dc_1\log\frac{D_W}{h}\\
&+\Bigl[(2c_1+8c_2+c_4)\frac{\varpi_y}{h}
+(16c_2+2c_5)\Bigl(\frac{\varpi_y}{h}\Bigr)^{2}\Bigr]\mathsf{T}\Bigr\}.
\end{split}
\end{equation*}
Put
\begin{equation*}
C'=C_\Pi\max\{c_1+8c_2+c_3+c_4+2c_5,\ 2dc_1,\ 2c_1+8c_2+c_4,\ 16c_2+2c_5\}.
\end{equation*}
Then
\begin{equation*}
\begin{split}
\mathfrak{B}(h)
&\le C'\Bigl\{1+\log\frac{D_W}{h}
+\Bigl[\frac{\varpi_y}{h}+\Bigl(\frac{\varpi_y}{h}\Bigr)^{2}\Bigr]\mathsf{T}\Bigr\}\,\theta'.
\end{split}
\end{equation*}

\emph{The remaining members.} By $(\gamma\text{-col})_h$ we have $\kappa_{\mathrm{col}}(h)\le1$,
and by $(\gamma\text{-nl})_h$ we have $\vartheta_{\mathrm{nl}}(h)\le\vartheta_1$. Since
$\theta\ge1$ and $\vartheta_1\ge0$, both are at most $\theta'$. The relative error
$\vartheta_t$ of Theorem~\ref{8.11:thm-deep-curvature} is kept small by the ceiling
$(\gamma\text{-t})$ under which that theorem is stated; in particular $\vartheta_t\le1$, so
$2\vartheta_t\le2\theta'$. The brace in \eqref{8.11:eq-small-heights-a} is at least $1$,
because $\log(D_W/h)\ge0$ and the bracket is non-negative. Hence these three members add at
most $4\theta'$ times the brace.

\emph{Conclusion.} This proves \eqref{8.11:eq-small-heights-a} with $C_4^{\flat}=C'+4$. The
constants $C_\Pi,c_1,\dots,c_5$ depend on $d$ and $\varrho_H$ only, by
Theorem~\ref{8.11:thm-deep-curvature}, and hence so does $C_4^{\flat}$.

\emph{The case $\tilde x$.} At the straddling-only minimiser $\tilde x$ we have $\varpi_y=0$.
The argument above, run with $\varpi_y=0$, has $t=r_1+1$, and the terms carrying $\theta_y''$
or a power of $\varpi_y/h$ are absent; the bound reads
\begin{equation*}
\|U(e\oplus\tilde x)(\partial p)-1\|
\le C_4^{\flat}\bigl[1+\log(D_W/h)\bigr]\theta'\varepsilon_{k_o}\eta^{2}
\end{equation*}
at every height $h\ge r_1+2$ for which $(\gamma\text{-col})_h$ holds. An example is
$h=w/L^{2}$, admissible as soon as $w/L^2\ge r_1+2$; by
Remark~\ref{8.11:rem-deep-bound-scope} the ratio $(r_1+1)/w$ is small for $g$ small after $L$
and $M$, so this holds for such $g$.

By \eqref{8.11:eq-bound-members-a}, $\kappa_{\mathrm{col}}(h)$ is proportional to $h^{-5/2}$.
Hence $\kappa_{\mathrm{col}}(w/L^2)=L^{5}\kappa_{\mathrm{col}}(w)$, and
$(\gamma\text{-col})_{w/L^2}$ is the condition $L^{10}\kappa_{\mathrm{col}}(w)^2\le1$. By
Remark~\ref{8.11:rem-deep-bound-scope}, this is again a one-sided smallness condition on $g$,
imposed after $L$ and $M$ are fixed.

Consequently, at $\tilde x$ the young cap $\theta\mathsf{T}$ is respected at all heights
$h\ge h_\flat:=w/L^2$ under the floor condition
\begin{equation*}
C_4^{\flat}\bigl(1+\log(L^2D_W/w)\bigr)(1+\vartheta_1)^2\le(1-\rho)\mathsf{T}
\end{equation*}
on $\mathsf{T}$, in which $\rho\in[0,1)$ is the margin that this floor condition leaves below
$\mathsf{T}$. Indeed, for $h\ge h_\flat$ we have $\log(D_W/h)\le\log(L^2D_W/w)$, so the
right-hand side of the bound at $\tilde x$ satisfies
\begin{equation*}
\begin{split}
C_4^{\flat}\bigl[1+\log(D_W/h)\bigr]\theta'\varepsilon_{k_o}\eta^{2}
&\le(1-\rho)\,\theta\mathsf{T}(1+\vartheta_1)^{-1}\varepsilon_{k_o}\eta^{2}\\
&\le\theta\mathsf{T}\,\varepsilon_{k_o}\eta^{2}.
\end{split}
\end{equation*}

Since $d_1=w+r_1$ when $\varpi_y=0$, the bound covers every plaquette with
$\delta\ge w+r_1+h_\flat$. The only depths at which $\tilde x$ can fail the constraint of the
first young row therefore form the sliver $w<\delta<w+r_1+h_\flat$, of width $w/L^2+r_1$.
The present bound confines the question to this sliver and does not decide whether $\tilde x$
fails there.
\end{proof}

\begin{corollary}[Box holes of arbitrary extent]\label{8.11:cor-arbitrary-extent}
Let the hole $H$ be a box of aspect at most $\varrho_H$ and of width $D$, for which the
inputs $(\mathrm{I}1)$--$(\mathrm{I}8)$ of Theorem~\ref{8.11:thm-deep-curvature} hold. Let $e$,
$\theta$, $\hat x$, $c_*$, $h$ and the fine plaquette $p$ be as in
Theorem~\ref{8.11:thm-deep-curvature}. Assume the remaining hypotheses of that theorem,
namely $(\mathrm{I}9)$, $(\mathrm{CEN})$, $(\gamma\text{-}\mathrm{loc})$ and
$(\gamma\text{-}t)$. Suppose that the ceilings
\eqref{8.11:eq-deep-bound-scope-a} on $\kappa_{\mathrm{col}}(h)$ and
$\vartheta_{\mathrm{nl}}(h)$, and the smallness conditions $(\gamma\text{-}\Delta)$ and
$(\gamma\text{-}\mathrm{W})$ of Theorem~\ref{8.11:thm-deep-curvature}, hold with the width
$D_W$ replaced by $D$. Then:
\begin{itemize}
\item[(a)] The bound \eqref{8.11:eq-deep-curvature-a} holds with the width $D_W$ and the block
count $\hat N_W$ of the zone replaced by the width $D$ and the block count of $H$. The
ceilings \eqref{8.11:eq-deep-bound-scope-a} read at $D$ bound $\kappa_{\mathrm{col}}(h)$
by $1$ and $\vartheta_{\mathrm{nl}}(h)$ by $\vartheta_1$. With these bounds, the bound is
logarithmic in the extent: the extent enters $\mathfrak{B}(h)$ only through
$\log^{+}(D/h)$, with coefficient $2dc_1C_\Pi\sigma_*$. For the reference depth
$c_*=-(3w+1)$ the amplitude $\sigma_*$ is $s_e$
(Remark~\ref{8.11:rem-deep-bound-scope}). Thus the bound grows at most logarithmically in the
smaller of the lateral extent and the depth of the hole, which for a box of aspect at most
$\varrho_H$ is comparable to $D/\varrho_H$.
\item[(b)] The four conditions above are ceilings on $g$ that depend on $D$. Consider an
unbounded family of extents $D$ at fixed $g$. Along such a family the conditions are caps
in the sense of Definition~\ref{2.3:def-cap}, now on $D$: they bound $D$ from above, and
nothing is asserted for that family.
\item[(c)] The constants $C_\Pi$ and $c_1,\dots,c_5$ depend on $d$ and $\varrho_H$ only, and
they deteriorate as $\varrho_H$ grows. Some positive power of $\varrho_H$ cannot be removed
from this dependence.
\end{itemize}
\end{corollary}

\begin{proof}
The proof of Theorem~\ref{8.11:thm-deep-curvature} uses nothing about the hole $H$ beyond two
facts: $H$ is a box of aspect at most $\varrho_H$, and the inputs
$(\mathrm{I}1)$--$(\mathrm{I}8)$ hold for $H$.

The size of $H$ enters that proof in four places only:
\begin{itemize}
\item[(1)] the energy bound $C_E\varepsilon_1^2\hat N_W$, taken from $(\mathrm{I}3)$ in the
reduction to the released minimiser $\check x$;
\item[(2)] the smallness of the centred axial gauge $\mathfrak{G}$, which is the inequality
\begin{equation*}
4(D_W+8w)\,\varepsilon_a\eta\le\frac{1}{64\,C_{ax}}
\end{equation*}
and is contained in $(\gamma\text{-}\mathrm{W})$;
\item[(3)] the length bound $\ell_{\max}\le D_W n$ of the path--surface system, which produces
the member $2dc_1C_\Pi\sigma_*\log^{+}(D_W/h)$ of the bound on the contribution of the
reference surface;
\item[(4)] the bound $\mathsf{L}\le4(D_W+8w)n$ on the diameter used for the defect terms.
\end{itemize}
Each of these is a statement about a box of the given width and aspect.

All of them hold verbatim for a box of width $D$ satisfying $(\mathrm{I}1)$--$(\mathrm{I}8)$.
In them $D_W$ is replaced by $D$ throughout, and $\hat N_W$ by the block count of that box.
This gives the first assertion of (a), under $(\gamma\text{-}\Delta)$,
$(\gamma\text{-}\mathrm{W})$ and the ceilings \eqref{8.11:eq-deep-bound-scope-a}, all read at
$D$. The ceilings give $\kappa_{\mathrm{col}}(h)\le1$ and
$\vartheta_{\mathrm{nl}}(h)\le\vartheta_1$. Hence the members carrying $\hat N_W$ and
$\mathsf{L}$ are bounded independently of $D$. Of the quantities read at $D$ there remains
only the member $2dc_1C_\Pi\sigma_*\log^{+}(D/h)$ of $\mathfrak{B}(h)$. This is the
logarithmic dependence asserted in (a). For the reference depth $c_*=-(3w+1)$ the amplitude
$\sigma_*$ is $s_e$, by Remark~\ref{8.11:rem-deep-bound-scope}.

Each of these conditions bounds $g$ from above by a quantity that depends on $D$. If $D$
ranges over an unbounded set while $g$ is held fixed, the conditions bound $D$ from above.
They are then caps on $D$, in the sense of Definition~\ref{2.3:def-cap}, and no conclusion is
drawn. This gives (b).

The constants $C_\Pi$, $c_1,c_2$ and $c_3,c_4,c_5$ depend on $d$ and $\varrho_H$ only
(Theorem~\ref{8.11:thm-deep-curvature}). Their growth with $\varrho_H$ is given by the lemma
of this chapter that estimates $C_\Pi$, the shell-sum constants $c_1,c_2$ and the derived
constants $c_3,c_4,c_5$ in terms of the aspect. That some positive power of $\varrho_H$ cannot be
removed from the dependence of the bound on the aspect is shown by the proposition on the
dependence of the bound on the aspect, which we take by statement. This gives (c).
\end{proof}

\begin{proposition}[Model: two-sided logarithmic profile behind the clamp]
\label{8.11:prop-two-sided-profile}
The model of this proposition is the block-mean model, whose data we recall. All its objects
live on the two-dimensional lattice $\Lambda_2$, with sites $z$, coordinates $(x_1,x_4)$ and unit
vector $e_4$ in the $x_4$-direction. In this model:
\begin{itemize}
\item $\Lambda_2$ is partitioned into square blocks of side $n$ in fine lattice units;
\item there is a hole made of $P_{bl}\times(P_{bl}-1)$ blocks, whose block boundary is denoted
$\partial^{bl}$;
\item the clamp $\mathsf{S}$ is the union of $\partial^{bl}$ and the set of sites of the hole of
block depth at most $w$, where $w$ is the clamp width;
\item the target is $\mathcal{T}=\eta^2T_P$, where $T_P$ is the exterior profile of the model and
$\eta>0$ is its scale parameter;
\item $M_{B_2}$ denotes the mean over a block $B_2$ of $\Lambda_2$, and $\Delta_2$ denotes the
lattice Laplacian of $\Lambda_2$.
\end{itemize}
Put $a:=w+r_c$, where $r_c\ge0$ is an integer and $a\le 13w/12+1$. Let $\mathfrak{R}_a^{\circ}$,
$\mathfrak{R}_a$ and $\partial\mathfrak{R}_a$ be the interior region, the closed region and its
boundary of the block-mean model, formed with the depth parameter of that model replaced by $a$.
Assume
$P_{bl}\ge 56w+4$.

Let $\psi\colon\Lambda_2\to\mathbb{R}$ be a function, and let $s,\tau,\kappa_j\ge0$ be
numbers; here $s$, $\tau$ and $\kappa_j$ are amplitudes of the model, fixed in this proposition
only. Assume:
\begin{itemize}
\item[(i)] $M_{B_2}\psi=s\,M_{B_2}\mathcal{T}$ on every block $B_2$ exterior to the hole;
\item[(ii)] $|\nabla\psi|\le\tau\eta^2$ on every nearest-neighbour pair of sites in $\mathsf{S}$;
\item[(iii)] $\Delta_2\psi=0$ on $\mathfrak{R}_a^{\circ}$;
\item[(iv)] (small trace jump across the collar) $|\psi(z)-\psi(z')|\le\kappa_j w\eta$ whenever
$z\in\partial\mathfrak{R}_a$ and $z'$ is the normal projection of $z$ on the last site row (or
column) of $\mathsf{S}$.
\end{itemize}
Let $y$ be a block depth with $2w\le y\le P_{bl}/8$. Put $m_y:=(y-a)n$, and let $z_y$ be the
site of the junction column $x_1=0$ at fine height $m_y$ above the bottom row of
$\partial\mathfrak{R}_a$. Then
\begin{equation}\label{8.11:eq-two-sided-profile-1}
\begin{aligned}
c_K'\,s\Bigl[\log\frac{P_*}{m_y+1}-C_K'\Bigr]-C_e(\tau+\kappa_j+s)
&\le \eta^{-2}\bigl[\psi(z_y+e_4)-\psi(z_y)\bigr]\\
&\le C_K''\,s\Bigl[1+\log^{+}\frac{P}{m_y+1}\Bigr]+C_e(\tau+\kappa_j+s).
\end{aligned}
\end{equation}
Here:
\begin{itemize}
\item $C_e:=40\,C_\nabla$ is absolute, where $C_\nabla$ is the absolute constant of the lattice
gradient estimate for harmonic functions on $\mathbb{Z}^2$;
\item $c_K'$, $C_K'$, $C_K''$ are the constants of part~(iii) of the strip-solution theorem for
the block-mean model, and $P$, $P_*$ are the lengths entering that theorem;
\item the ratio $P_*/(m_y+1)$ is bounded above and below by constant multiples of
$P_{bl}/(y-a)$, so that the block side $n$ cancels.
\end{itemize}
\end{proposition}

\begin{proof}
Throughout, write
\begin{equation}\label{8.11:eq-two-sided-profile-2}
\Gamma_a:=(w+3)\tau+\kappa_j w+(a+2)s .
\end{equation}

\emph{Step 1: the block boundary of the hole.} We apply the boundary comparison of the
block-mean model on the block boundary of the hole without change, with the amplitude $s$ in the
place of the amplitude there. It
gives
\begin{equation}\label{8.11:eq-two-sided-profile-3}
|\psi-s\mathcal{T}|\le(2\tau+s)\eta\qquad\text{on }\partial^{bl}.
\end{equation}

\emph{Step 2: transfer to $\partial\mathfrak{R}_a$.} Let $z\in\partial\mathfrak{R}_a$, and let $z'$
be its normal projection on the last site row (or column) of $\mathsf{S}$. We combine four
estimates.
\begin{itemize}
\item The bound \eqref{8.11:eq-two-sided-profile-3} holds on $\partial^{bl}$.
\item Along a chain through the strip, made of $wn+1$ nearest-neighbour pairs of sites in
$\mathsf{S}$, each pair contributes a variation of $\psi$ of at most $\tau\eta^2$, by
hypothesis~(ii).
\item The jump from $z'$ to $z$ is at most $\kappa_j w\eta$, by hypothesis~(iv).
\item Along a wall chain of $an+n$ steps, each step changes $s\mathcal{T}$ by at most $s\eta^2$,
because $|\nabla_1\mathcal{T}|\le\eta^2$.
\end{itemize}
Summing these four contributions gives
\begin{equation}\label{8.11:eq-two-sided-profile-4}
|\psi-s\mathcal{T}|\le\Gamma_a\,\eta\qquad\text{on }\partial\mathfrak{R}_a ,
\end{equation}
with $\Gamma_a$ as in \eqref{8.11:eq-two-sided-profile-2}.

\emph{Step 3: the maximum principle on $\mathfrak{R}_a$.} Let $u_a$ be the function that is
harmonic on $\mathfrak{R}_a^{\circ}$ and equal to $T_P$ on $\partial\mathfrak{R}_a$. The function
$\psi-s\eta^2u_a$ is harmonic on $\mathfrak{R}_a^{\circ}$, because $\psi$ is harmonic there by
hypothesis~(iii) and $u_a$ is harmonic there by construction. On $\partial\mathfrak{R}_a$ it
equals $\psi-s\mathcal{T}$, because $\mathcal{T}=\eta^2T_P$. This is the maximum-principle
comparison of the block-mean model: by the maximum principle and
\eqref{8.11:eq-two-sided-profile-4},
\begin{equation}\label{8.11:eq-two-sided-profile-5}
|\psi-s\eta^2u_a|\le\Gamma_a\,\eta\qquad\text{on }\mathfrak{R}_a .
\end{equation}

\emph{Step 4: comparison with the strip solution.} Put
\begin{equation*}
H_a:=(P_{bl}-2a-1)n+1 ,
\end{equation*}
and let $u$ be the strip solution of height $H_a$ of the strip-solution theorem. The comparison
of the block-mean model with the strip solution, which uses part~(ii) of that theorem together
with the symmetry columns, gives
\begin{equation}\label{8.11:eq-two-sided-profile-6}
|u_a-u|\le 2an\qquad\text{on }\mathfrak{R}_a .
\end{equation}
Put $\phi:=\psi-s\eta^2u$. Then $\phi$ is harmonic on $\mathfrak{R}_a^{\circ}$. Moreover, by
\eqref{8.11:eq-two-sided-profile-5} and \eqref{8.11:eq-two-sided-profile-6},
\begin{equation}\label{8.11:eq-two-sided-profile-7}
|\phi|\le(\Gamma_a+2as)\,\eta\qquad\text{on }\mathfrak{R}_a .
\end{equation}

\emph{Step 5: the gradient at $z_y$.} We have $m_y\ge2$. Also $m_y\le P/8$, and this is less
than the distance from $z_y$ to the walls and to the top row. Hence the $\mathbb{Z}^2$-ball of
radius $m_y-1$ about $z_y$ lies in $\mathfrak{R}_a^{\circ}$, and $\phi$ is harmonic on it. The
lattice gradient estimate for harmonic functions on $\mathbb{Z}^2$, applied to $\phi$ on this
ball together with \eqref{8.11:eq-two-sided-profile-7}, gives
\begin{equation}\label{8.11:eq-two-sided-profile-8}
\bigl|\phi(z_y+e_4)-\phi(z_y)\bigr|
\le\frac{C_\nabla(\Gamma_a+2as)\,\eta}{m_y-1}
\le\frac{2C_\nabla(\Gamma_a+2as)\,\eta^2}{y-a}.
\end{equation}

\emph{Step 6: the error in terms of the amplitudes.} From $y\ge2w$, $a\le 13w/12+1$ and $w\ge2$
we get
\begin{equation*}
y-a\ge 2w-\frac{13w}{12}-1\ge\frac{5w}{12}.
\end{equation*}
Now $\Gamma_a+2as=(w+3)\tau+\kappa_jw+(3a+2)s$. Dividing by $y-a\ge 5w/12$ and using $w\ge2$
gives the first inequality below. The second uses $a\le 13w/12+1$ and $w\ge2$, which yield
$(3a+2)/w\le 13/4+5/w\le 8$:
\begin{equation*}
\frac{\Gamma_a+2as}{y-a}
\le\frac{12}{5}\Bigl[\frac52\tau+\kappa_j+\frac{(3a+2)s}{w}\Bigr]
\le\frac{12}{5}\Bigl[\frac52\tau+\kappa_j+8s\Bigr].
\end{equation*}
Inserting this into \eqref{8.11:eq-two-sided-profile-8} gives
\begin{equation}\label{8.11:eq-two-sided-profile-9}
\begin{aligned}
\eta^{-2}\bigl|\phi(z_y+e_4)-\phi(z_y)\bigr|
&\le C_\nabla\Bigl[12\tau+\frac{24}{5}\kappa_j+\frac{192}{5}s\Bigr]\\
&\le 40\,C_\nabla(\tau+\kappa_j+s)=C_e(\tau+\kappa_j+s).
\end{aligned}
\end{equation}

\emph{Step 7: the strip profile.} By the definition of $\phi$,
\begin{equation*}
\eta^{-2}\bigl[\psi(z_y+e_4)-\psi(z_y)\bigr]
=s\bigl[u(z_y+e_4)-u(z_y)\bigr]+\eta^{-2}\bigl[\phi(z_y+e_4)-\phi(z_y)\bigr].
\end{equation*}
Part~(iii) of the strip-solution theorem applies at fine height $m=m_y$, since $m_y\le H_a/4-1$;
this holds because $y\le P_{bl}/8$ and $P_{bl}\ge 56w+4$. It bounds the increment
$u(z_y+e_4)-u(z_y)$ of the strip solution below by $c_K'[\log(P_*/(m_y+1))-C_K']$ and above by
$C_K''[1+\log^{+}(P/(m_y+1))]$. We multiply these bounds by $s\ge0$ and add the bound
\eqref{8.11:eq-two-sided-profile-9} on the $\phi$-term. This gives
\eqref{8.11:eq-two-sided-profile-1}.
\end{proof}

\begin{remark}[Consistency with the logarithmic lower bounds]\label{8.11:rem-lower-bound-consistency}
In this remark the \emph{five lower bounds} are the following statements: the lower bound in the
block-mean model, with constant $c^{\natural}$; the lower bound
read at an arbitrary depth $t_+$, with constant $2c_K/3$; and three further lower bounds of the
same shape, with constants $c_K/8$, $c_K/9$ and $c_{\mathrm{b}}$. Here $c^{\natural}$, $c_K$ and
$c_{\mathrm{b}}$ are the constants of these statements. Each of the five lower bounds concerns a
configuration that is pinned to the kinked exterior, clamped on the collar and on the inner clamp
set of the respective statement, and free beyond the clamp. It asserts that every number $B$
bounding the curvature of such a configuration in the region lying between the clamped set and
the free part satisfies
\begin{equation*}
B\ \ge\ c\,s\log(D/w)-C(\tau) ,
\end{equation*}
where $s$ is the amplitude of the exterior, $D$ is the outer cutoff of the logarithm in the
respective bound (in the block-mean model, $D=P_{bl}$), and $c$ is one of the five constants
\begin{equation}\label{8.11:eq-lower-bound-consistency-1}
c\in\Bigl\{\,c^{\natural},\ \tfrac{2c_K}{3},\ \tfrac{c_K}{8},\ \tfrac{c_K}{9},\ c_{\mathrm{b}}\Bigr\}.
\end{equation}
The constant $C(\tau)$ may depend on the clamp bound $\tau$.
\begin{enumerate}
\item[(a)] Consider the block-mean model of Proposition~\ref{8.11:prop-two-sided-profile}. With
exactly the information that the proof of Theorem~\ref{8.11:thm-deep-curvature} uses, the
normal--tangential curvature on the junction column at block depth $2w$ is of size
$s\log(P_{bl}/w)$, from below and from above, up to an additive error $C(\tau+\kappa_j+s)$,
with $C$ depending only on the constants $c_K'$, $C_K'$, $C_K''$ and $C_e$ of the model. The
lower bound $c\,s\log(D/w)-C(\tau)$ is therefore attained from above, with the same amplitude $s$,
the same two cutoffs, and an additive error of the size of the clamp.
\item[(b)] The bound \eqref{8.11:eq-deep-curvature-a} and the five lower bounds are consistent if
and only if
\begin{equation*}
C_3^{\flat}=2d\,C_\Pi c_1\ \ge\ c
\end{equation*}
for each of the five constants $c$ in \eqref{8.11:eq-lower-bound-consistency-1}.
\item[(c)] On the right side of \eqref{8.11:eq-deep-curvature-a}, the only logarithm is
$\log(D_W/h)$. Its inner cutoff is the reading height $h$.
\item[(d)] As $h\downarrow r_1$, the bound \eqref{8.11:eq-deep-curvature-a} makes no claim that the
curvature behind the clamp is bounded by the constant $1$ in exterior units, that is, by
$\varepsilon_{k_o}\eta^2$.
\item[(e)] The outer cutoff of the logarithm is the width $D_W$ of the box $H$.
\item[(f)] The coefficient of the logarithm grows with the amplitude of the exterior, in
proportion to $s_e$.
\end{enumerate}
\end{remark}

\begin{proof}
\emph{(a).} The model proposition, Proposition~\ref{8.11:prop-two-sided-profile}, is a statement
in the block-mean model. Its hypotheses correspond to the information used in the proof of
Theorem~\ref{8.11:thm-deep-curvature} as follows.
\begin{itemize}
\item Hypothesis (i), the pinning of the block means to $s\,\mathcal{T}$ on exterior blocks,
corresponds to the pure windows on the reference surface $\Sigma_{c_*}$.
\item Hypothesis (ii), the clamp $|\nabla\psi|\le\tau\eta^2$ on pairs in $\mathsf{S}$, corresponds
to the clamp condition among the inputs of Theorem~\ref{8.11:thm-deep-curvature}.
\item Hypothesis (iv), the small trace jump across the collar, corresponds to the energy bound
on the collar in the proof of Theorem~\ref{8.11:thm-deep-curvature}. Alternatively, it corresponds
to the young cap $\mathsf{T}$ times the thickness $\varpi_y=5w/L$ of the first young row, where
$y_+=w+5w/L$ as in Corollary~\ref{8.11:cor-logarithmic-bound}; this gives $\kappa_j=5\mathsf{T}/L$
at $\hat x$.
\item Hypothesis (iii), harmonicity on $\mathfrak{R}_a^{\circ}$, corresponds to the freedom of the
released minimiser $\check x$, which is exactly Yang--Mills on $\Omega_{d_1}$.
\end{itemize}

Take $y=2w$ in Proposition~\ref{8.11:prop-two-sided-profile}. Its condition
$2w\le y\le P_{bl}/8$ holds because $P_{bl}\ge 56w+4$. With $m_y=(2w-a)n$, the model proposition
gives
\begin{equation}\label{8.11:eq-lower-bound-consistency-2}
\begin{split}
c_K'\,s\bigl[\log\bigl(P_*/(m_y+1)\bigr)-C_K'\bigr]-C_e(\tau+\kappa_j+s)
&\le \eta^{-2}\bigl[\psi(z_y+e_4)-\psi(z_y)\bigr]\\
&\le C_K''\,s\bigl[1+\log^{+}\bigl(P/(m_y+1)\bigr)\bigr]+C_e(\tau+\kappa_j+s).
\end{split}
\end{equation}
Here $P$ is the fine-lattice length of the model that enters the upper member of the display in
Proposition~\ref{8.11:prop-two-sided-profile}; in the proof of that proposition the inequality
$m_y\le P/8$ places $m_y$ below the distance from $z_y$ to the walls and to the top row of the
hole.
The middle term is the normal--tangential curvature at the site $z_y$ of the junction column.

We now compare the cutoffs with $P_{bl}/w$. Since $a\ge w$ we have $y-a\le w$. Since
$a\le 13w/12+1$ and $w\ge2$ we have $y-a\ge 5w/12$. Hence $P_{bl}/(y-a)$ lies between $P_{bl}/w$
and $12P_{bl}/(5w)$. The model proposition states that $P_*/(m_y+1)$ and $P_{bl}/(y-a)$ are
comparable up to absolute factors, the block side $n$ cancelling. Applying this comparability in
the lower member gives $\log(P_{bl}/w)$ there up to an additive constant; in the upper member,
$P/(m_y+1)$ is compared with $P_{bl}/(y-a)$ in the same way, and $\log^{+}(P/(m_y+1))$ equals
$\log(P_{bl}/w)$ up to an additive constant. After multiplication by $s$ these additive constants
are bounded by a constant multiple of $s$. The
two-sided bound \eqref{8.11:eq-lower-bound-consistency-2} therefore has the same amplitude $s$ on
both sides, the same outer cutoff $P_{bl}$ and inner cutoff $w$, and an additive error bounded by
$C(\tau+\kappa_j+s)$, with $C$ depending only on $c_K'$, $C_K'$, $C_K''$ and $C_e$; this error is
of the size of the clamp.

\emph{(b).} Let the exterior be kinked at amplitude $\sigma$. By
\eqref{8.11:eq-units-cap-amplitude-a},
\begin{equation*}
\sigma_e=\theta s_e+\vartheta_1\ \ge\ \sigma-o_{\gamma}(1).
\end{equation*}
Hence the logarithmic member of the bound at depth $2w$ satisfies
\begin{equation*}
C_3^{\flat}\sigma_e\log(D_W/w)\ \ge\ C_3^{\flat}\bigl(\sigma-o_{\gamma}(1)\bigr)\log(D_W/w).
\end{equation*}
This member is the one in \eqref{8.11:eq-logarithmic-bound-b}, which is deduced from
\eqref{8.11:eq-deep-curvature-a}. Against it stand the five lower bounds, each of the form
$c\,\sigma\log(D_W/w)-C$ with $c$ one of the constants \eqref{8.11:eq-lower-bound-consistency-1}.
There is no contradiction between the two if and only if
\begin{equation*}
C_3^{\flat}=2d\,C_\Pi c_1\ \ge\ c
\end{equation*}
for each of these five values; the first equality is the definition of $C_3^{\flat}$ in
\eqref{8.11:eq-logarithmic-bound-b}. This is statement (b).

Two facts bear on the constants of this comparison. In the block-mean model, the corresponding
comparison is between the constants of the two members of
\eqref{8.11:eq-lower-bound-consistency-2}: there the inequality $C_K''\ge c_K'$ is
the two-sidedness of the bound on the strip solution of the block-mean model (part (iii) of the
strip-solution theorem used in the proof of the model proposition,
Proposition~\ref{8.11:prop-two-sided-profile}), which the model proposition uses in its lower
and upper members. For the bound \eqref{8.11:eq-deep-curvature-a}, $C_\Pi c_1$ majorises the
$\ell^1$-mass, per dyadic shell, of the
kernel whose exact two-dimensional section is the reflected two-dimensional kernel $K^R_y$ of the
kernel-bound lemma of the localisation analysis; in that lemma the kernel bound of the first kind
majorises the kernel bound of the second kind. Neither fact is a comparison of
$2d\,C_\Pi c_1$ with the five constants \eqref{8.11:eq-lower-bound-consistency-1}; statement (b)
is the equivalence proved above.

\emph{(c).} The right side of \eqref{8.11:eq-deep-curvature-a} contains no logarithm other than
$\log(D_W/h)$. In particular it contains neither of the following:
\begin{itemize}
\item $\log D_W$ alone, hence no $\check R(g)$;
\item $\log n=k_o\log L$, which the centring of the axial gauge $\mathfrak{G}$ removes.
\end{itemize}
Let $\varepsilon$ and $\tau_{\mathrm{d}}$ denote the two small parameters that enter the choice of
the free depth $d_1$ and the ceilings, the second distinct from the clamp bound $\tau$.
The quantities $r_1$, $\delta_0$, $\log\varepsilon^{-1}$, $\log\tau_{\mathrm{d}}^{-1}$ and
$\Lambda$ enter only
through the free depth $d_1$ and through the ceilings. The constants $B_3$ (which enters through
$C_E=3B_3^2+2$), $K_w$ and $C_E$ enter
only through the ceilings. The integer $N$ of $\mathrm{SU}(N)$ does not enter the main member
$\mathfrak{B}(h)$: it occurs
only in $\kappa_{\mathrm{col}}(h)$ and $\vartheta_{\mathrm{nl}}(h)$, which the ceilings
$\kappa_{\mathrm{col}}(h)\le1$ and $\vartheta_{\mathrm{nl}}(h)\le\vartheta_1$ bound.

The inner cutoff of the logarithm is the reading height $h$. This is the same convention as in the
lower bound read at an arbitrary depth $t_+$ (the one with constant $2c_K/3$), where the
logarithm's inner cutoff is the depth at which the bound is read.

\emph{(d).} Let $h\downarrow r_1$. There are two cases, according to the choice of the reference
depth.
\begin{itemize}
\item For $c_*=-(3w+1)$, the bracket in \eqref{8.11:eq-deep-curvature-a} carries the additive
member $(4w/h)\theta''$, and so grows like $w/h$.
\item For $c_*=w-1$, it carries $\theta''\log(D_W/h)$ and $\kappa_{\mathrm{col}}(h)$, and both
increase as $h$ decreases; compare Corollary~\ref{8.11:cor-small-heights}.
\end{itemize}
Hence \eqref{8.11:eq-deep-curvature-a} never asserts that the curvature behind the clamp is
bounded by the constant $1$ in exterior units (compare the corollary of the block-mean treatment
of the lower bounds that concerns such a profile behind the clamp, and the remark accompanying
it). It is compatible with a flux jump of size
$s\log(D_W/w)$ across the inner face of the clamp. It is also compatible with the saturated profile
$(2s/\pi)\log(D_W/\ell_0)$, where $\ell_0$ is the inner cutoff of that profile, one block behind
that face. There $h<r_1$, and
Theorem~\ref{8.11:thm-deep-curvature} asserts nothing.

\emph{(e).} The logarithm in \eqref{8.11:eq-deep-curvature-a} is $\log^{+}(D_W/h)$, with $D_W$ the
width of the box $H$. In the lower bound with constant $c_{\mathrm{b}}$ the outer cutoff is
\begin{equation*}
\min\bigl(\text{lateral width of the region},\ \text{its depth}/(2\sqrt3)\bigr).
\end{equation*}
This cutoff agrees with $D_W$ up to the aspect bound $\varrho_H$.

\emph{(f).} In \eqref{8.11:eq-logarithmic-bound-b}, the logarithmic member is
$C_3^{\flat}s_e^{+}\log(D_W/w)$ times $\theta(1+\vartheta_1)\varepsilon_{k_o}\eta^2$, where
$C_3^{\flat}$ depends on $(d,\varrho_H)$ only. Its coefficient is therefore proportional to $s_e$,
and so it grows with the amplitude of the exterior.
\end{proof}

\begin{remark}[Interior windows and the blocking-factor floor]\label{8.11:rem-blocking-floor}
Assume the hypotheses of Corollary~\ref{8.11:cor-logarithmic-bound}, among them the geometric
condition \eqref{8.11:eq-logarithmic-bound-a}. Let $\rho$, with $0<\rho<15/16$, be the room
parameter, as in
\eqref{8.11:eq-scale-jump-a}, so that a
window of threshold $\theta$ has \emph{room} $\rho\theta$ when its value is at most
$(1-\rho)\theta$. Write $a$ for the exponent in the zone-width relation of this chapter
\begin{equation*}
8L^{a}(M/M_2)\le D_W/w\le 17L^{a}(M/M_2)
\end{equation*}
between the zone width $D_W$ and the clamp width $w$. Consider a hole window of level
$i\le k_o-1$ whose test region lies in the box
$\Omega_{2w}=\{\delta\ge 2w\}$, and write $f$ for its value at the canonical comparison point
$\hat x$.
\begin{enumerate}
\item[(a)] The value $f$ obeys
\begin{equation}\label{8.11:eq-blocking-floor-1}
f\le\bigl[C_1^{\flat}+C_2^{\flat}L+C_3^{\flat}s_e^{+}\log(D_W/w)\bigr]
\,\theta(1+\vartheta_1)^2(1+\beta_0)L^{-2}+\vartheta_1 .
\end{equation}
\item[(b)] Suppose that $\vartheta_1\le\theta/16$ and
\begin{equation*}
s_e^{+}\log(D_W/w)\le\log\bigl(17\,\hat r\,L^{a}(M/M_2)\bigr);
\end{equation*}
the latter holds, by the zone-width relation, whenever $0\le s_e^{+}\le 1$, $\hat r\ge 1$ and
$D_W/w\ge 1$.
Then the window has room $\rho\theta$ whenever
\begin{equation}\label{8.11:eq-blocking-floor-a}
\bigl[C_1^{\flat}+C_2^{\flat}L+C_3^{\flat}\log\bigl(17\,\hat r\,L^{a}(M/M_2)\bigr)\bigr]
(1+\beta_0)(1+\vartheta_1)^3\le\Bigl(\frac{15}{16}-\rho\Bigr)L^2 .
\end{equation}
\item[(c)] If $M/M_2$ is a power of $L$ whose exponent is fixed before $L$ is chosen, then
\eqref{8.11:eq-blocking-floor-a} is a one-sided lower bound on $L$. If instead $M/M_2$ is
chosen after $L$, with no exponent fixed in advance, then
\eqref{8.11:eq-blocking-floor-a} is a one-sided upper bound on $M/M_2$ at fixed $L$,
\begin{equation*}
M/M_2\le L^{-a}\exp[\Theta(L^2)],
\end{equation*}
where $\Theta(L^2)$ denotes a quantity bounded above and below by positive multiples of $L^2$ for
all sufficiently large $L$; every fixed power of $L$, taken as the value of $M/M_2$, satisfies
this bound for all $L$ beyond a threshold depending on that power.
\end{enumerate}
Condition \eqref{8.11:eq-blocking-floor-a} is not a hypothesis of any statement of this book.
\end{remark}

\begin{proof}
(a) Every plaquette $p$ of the test region satisfies $\delta(p)\ge 2w$, so
Corollary~\ref{8.11:cor-logarithmic-bound} applies to $p$. It gives
\begin{equation*}
\|U(e\oplus\hat x)(\partial p)-1\|
\le\bigl[C_1^{\flat}+C_2^{\flat}\,\mathsf{T}/L+C_3^{\flat}s_e^{+}\log(D_W/w)\bigr]
\,\theta(1+\vartheta_1)\,\varepsilon_{k_o}\eta^2 .
\end{equation*}
Since $\mathsf{T}/L\le L$, the member $C_2^{\flat}\mathsf{T}/L$ is at most $C_2^{\flat}L$.
By the window bound of this chapter, which reads \cite[(1.19)]{B-CMP119} with relative error
$\vartheta_1$ and which we apply to hole windows of level $i\le k_o-1$ at the
configuration $U(e\oplus\hat x)$ of the canonical comparison point $\hat x$,
the value of the window satisfies
\begin{equation}\label{8.11:eq-blocking-floor-2}
f\le\frac{\sup_{p}\|U(e\oplus\hat x)(\partial p)-1\|}{\varepsilon_{k_o-1}L^2\eta^2}+\vartheta_1 ,
\end{equation}
the supremum being taken over the plaquettes $p$ of the test region. Inserting the plaquette bound
just obtained into \eqref{8.11:eq-blocking-floor-2} gives
\begin{equation*}
f\le\bigl[C_1^{\flat}+C_2^{\flat}L+C_3^{\flat}s_e^{+}\log(D_W/w)\bigr]
\,\theta(1+\vartheta_1)\,\frac{\varepsilon_{k_o}}{\varepsilon_{k_o-1}}\,L^{-2}+\vartheta_1 .
\end{equation*}
The bound \eqref{8.11:eq-blocking-floor-1} follows from this inequality and the comparison of
the exterior unit with the unit of level $k_o-1$,
\begin{equation*}
\varepsilon_{k_o}\le(1+\beta_0)(1+\vartheta_1)\,\varepsilon_{k_o-1},
\end{equation*}
which is the comparison of consecutive small-field thresholds along the flow of the running
couplings, and which supplies the factor $(1+\beta_0)$ and the second factor $(1+\vartheta_1)$.

(b) It suffices to show $f\le(1-\rho)\theta$. Multiplying
\eqref{8.11:eq-blocking-floor-a} by $\theta L^{-2}(1+\vartheta_1)^{-1}$ and using the
inequality for $s_e^{+}\log(D_W/w)$ assumed in (b) in \eqref{8.11:eq-blocking-floor-1}, we obtain
\begin{equation*}
f\le\Bigl(\frac{15}{16}-\rho\Bigr)\frac{\theta}{1+\vartheta_1}+\vartheta_1
\le\Bigl(\frac{15}{16}-\rho\Bigr)\theta+\frac{\theta}{16}=(1-\rho)\theta ,
\end{equation*}
where the middle step uses $\vartheta_1\ge0$ and the assumption $\vartheta_1\le\theta/16$; this
is the only use of the margin $1/16$ in \eqref{8.11:eq-blocking-floor-a}.

(c) By Corollary~\ref{8.11:cor-logarithmic-bound}, the constants
$C_1^{\flat},C_2^{\flat},C_3^{\flat}$ of \eqref{8.11:eq-logarithmic-bound-b} depend on $d$ and
$\varrho_H$ only. The left-hand side of \eqref{8.11:eq-blocking-floor-a} involves only $L$, these
constants, $\hat r$, $a$, $M/M_2$, $\beta_0$ and $\vartheta_1$; the multiplicative constant of
the general window theorem of this chapter does not enter it.

First let $M/M_2$ be a power of $L$ whose exponent is fixed before $L$, as in the order of choice
of Definitions~\ref{2.3:not-stages-first}, \ref{2.3:not-stages-middle}
and~\ref{2.3:not-stages-last}. The numbers $\hat r$, $a$, $\beta_0$,
$\vartheta_1$ and $\rho$ are independent of $L$. Then
$\log(17\hat rL^{a}(M/M_2))$ is $O(\log L)$, the bracket is $O(L+\log L)$, and since the factor
$(1+\beta_0)(1+\vartheta_1)^3$ does not grow with $L$, the left-hand side of
\eqref{8.11:eq-blocking-floor-a} is $O(L+\log L)$. Since $\rho<15/16$, the right-hand side is a
fixed positive multiple of $L^2$. Hence \eqref{8.11:eq-blocking-floor-a} holds for all $L$
beyond a threshold, and it imposes no upper bound on $L$.

Now let $M/M_2$ be chosen after $L$, with no exponent fixed in advance, as in \cite{B-CMP119}.
Solving \eqref{8.11:eq-blocking-floor-a} for the
logarithm gives an
upper bound $\log(M/M_2)\le -a\log L+\Theta(L^2)$, that is, $M/M_2\le L^{-a}\exp[\Theta(L^2)]$.
By the first case, every fixed power of $L$, taken as the value of $M/M_2$, satisfies
\eqref{8.11:eq-blocking-floor-a} for all $L$ beyond a threshold depending on that power; hence
beyond that threshold the range of admissible values of $M/M_2$ is non-empty.
\end{proof}

\begin{corollary}[A good threshold for every exterior]\label{8.11:cor-threshold-existence}
Let $\bar q$ and $\bar e$ be the threshold parameters of the skin comparison theorem, with
$0<\bar q\le 1/8$ (the bound under which the interval $J$ below lies in $[1,2]$) and
$\bar e\le 1/8$, and put
\begin{equation}\label{8.11:eq-threshold-existence-1}
J:=[1+4\bar q,\ 1+8\bar q],\qquad
\Lambda_J:=\frac{C_E\,\varepsilon_1^2\,\hat N_W}{\bar q}.
\end{equation}
Let $r\in L^{-1}\mathbb{N}$ be such that $\vartheta_r\le 2\bar q$, where $\vartheta_r$ is the relative error
of the released problem, \eqref{8.11:eq-released-problem-a}. Then the following hold.
\begin{itemize}
\item[(a)] For every exterior $e$ whose set $S(e)$, attached to $e$ in the skin comparison theorem,
is non-empty, there is a threshold $\theta(e)\in J$ such
that $(\mathrm{VL})_{\theta(e)}[\vartheta_r,\Lambda_J]$ holds, in the sense of
Definition~\ref{8.11:def-value-lipschitz}. Every such $\theta(e)$ satisfies
$\theta(e)-1\ge 2\bar q$, which is the only condition on the threshold in the segment inequality
lemma, and $\theta(e)\le 2$, which is within the range of
thresholds allowed by the threshold-range lemma.
\item[(b)] The shift $t_1:=16\bar q$, which does not depend on $e$, satisfies
$t_1(1-\bar e)>\theta-1$ for every $\theta\le 1+8\bar q$, in particular for $\theta=\theta(e)$ for every
exterior $e$ as in \textup{(a)}. With this shift, the one-sided smallness conditions on $g$ of the
skin comparison theorem in which the shift enters as $8\bar q$ are to be read with $16\bar q$ in
place of $8\bar q$.
\end{itemize}
The tail length belonging to $\Lambda_J$ is $r_0(\Lambda_J)$ of \eqref{8.11:eq-released-problem-a},
with $\varepsilon$, $\vartheta_\Delta$ and the constants as fixed there:
\begin{equation}\label{8.11:eq-threshold-existence-2}
r_0(\Lambda_J)=\Bigl\lceil \delta_0^{-1}\log\frac{2K_wC_{\delta_0}\sigma_K\,C_E\,\varepsilon_1^2\,\hat N_W}
{\bar q\,\vartheta_\Delta^2\,\varepsilon_{\min}\,\varepsilon^2}\Bigr\rceil ;
\end{equation}
in particular $r_0(\Lambda_J)$ is of order $\delta_0^{-1}$ times the logarithm of a quantity bounded
polynomially in the number of hole data bonds and in $\varepsilon^{-1}$.
\end{corollary}

\begin{proof}
(a) The interval $J$ has length $|J|=4\bar q>0$, and $J\subset[1,2]$ because $8\bar q\le 1$. The
constant of Lemma~\ref{8.11:lem-good-thresholds} for this interval is
$4C_E\varepsilon_1^2\hat N_W/|J|=C_E\varepsilon_1^2\hat N_W/\bar q$, which is $\Lambda_J$ of
\eqref{8.11:eq-threshold-existence-1}. The number $\vartheta_r$ is positive, being a product of positive
constants and of $e^{-\delta_0 r}$, and $\vartheta_r\le 2\bar q=|J|/2$ by hypothesis. Thus the hypotheses
of Lemma~\ref{8.11:lem-good-thresholds} hold with $\vartheta=\vartheta_r$, for every exterior $e$ with
$S(e)\neq\emptyset$. That lemma gives that the set
\begin{equation*}
G_J(e,\vartheta_r)=\bigl\{\theta\in[1+4\bar q,\,1+8\bar q-\vartheta_r]:
(\mathrm{VL})_{\theta}[\vartheta_r,\Lambda_J]\text{ holds}\bigr\}
\end{equation*}
has Lebesgue measure at least $|J|/4=\bar q>0$. In particular it is not empty, and we choose
$\theta(e)\in G_J(e,\vartheta_r)\subset J$. Any $\theta(e)\in J$ satisfies
$\theta(e)-1\ge 4\bar q\ge 2\bar q$, the condition on the threshold in the segment inequality lemma,
and $\theta(e)\le 1+8\bar q\le 2$, within the range allowed by the threshold-range lemma.

(b) The transport step of the skin comparison theorem imposes exactly one condition linking the
shift $t_1$ to the threshold, namely
$t_1(1-\bar e)>\theta-1$. Take $\theta=1+8\bar q$ and $t_1=16\bar q$. Expanding,
\begin{equation*}
(1-16\bar q)(1+8\bar q)+16\bar q\,\bar e
=1-8\bar q-128\bar q^2+16\bar q\,\bar e
\le 1-6\bar q-128\bar q^2\le 1-6\bar q ,
\end{equation*}
where $16\bar q\,\bar e\le 2\bar q$ was used, by $\bar e\le 1/8$. Writing the left side as
$\theta-t_1\theta+t_1\bar e$ and using $t_1\theta=t_1+t_1(\theta-1)=16\bar q+128\bar q^2$, the first
inequality of the last display is the inequality
\begin{equation*}
\theta-1\le t_1(1-\bar e)-6\bar q,
\end{equation*}
so $t_1(1-\bar e)>\theta-1$ at $\theta=1+8\bar q$, since $\bar q>0$. For $\theta\le 1+8\bar q$ the right
side $\theta-1$ is not larger, while $t_1(1-\bar e)$ does not depend on $\theta$; hence the condition
holds for every $\theta\le 1+8\bar q$, in particular for every $\theta(e)\in J$. The shift $t_1$ is the
same for all exteriors. The remaining sentence of \textup{(b)} is a convention on the smallness
conditions on $g$ of the skin comparison theorem, which this corollary uses as stated there with
$16\bar q$ in place of $8\bar q$; the argument above concerns only the transport condition.

Finally, \eqref{8.11:eq-threshold-existence-2} is the formula for $r_0(\Lambda)$ in
\eqref{8.11:eq-released-problem-a}, the least $r\in L^{-1}\mathbb{N}$ with
$\vartheta_r\Lambda\le\tfrac12\vartheta_\Delta^2\varepsilon^2$, written out at $\Lambda=\Lambda_J$ from
\eqref{8.11:eq-threshold-existence-1}; its size is the one stated with
\eqref{8.11:eq-released-problem-a} for $\Lambda=\Lambda_J$ and $|J|=4\bar q$.
\end{proof}

We add two comments on the use of Corollary~\ref{8.11:cor-threshold-existence}.
The skin comparison theorem is applied at the threshold $\theta=1+4\bar q$, the left endpoint of $J$;
its comparison of
numerator and denominator is made for each exterior $e$ separately, so a threshold depending on $e$ may
be used in it, one exterior at a time, while the shift $t_1$ stays the same for all exteriors.

Under an alternative hypothesis no selection of the threshold is needed: if the room condition
stated with Definition~\ref{8.11:def-value-lipschitz} holds at
a threshold $\theta$ with a margin exceeding $2q$, where $q=q(x',x_1)$ is the quantity attached there
to the comparison pair $(x',x_1)$, then the implication stated with
Definition~\ref{8.11:def-value-lipschitz} gives
$(\mathrm{VL})_\theta[\vartheta,\Lambda]$ for every $\vartheta\le1$, at that
$\theta$, with $\Lambda$ equal to the constant $\Lambda_{\mathbf{A}}$ of that implication divided by the
excess of the margin over $2q$; every threshold at which this room condition holds is then good.
The value-Lipschitz property is
not obtained from a positive room margin alone; the analysis of the thresholds at which the room
vanishes shows where such a derivation fails. Lemma~\ref{8.11:lem-good-thresholds} makes no use of
the room: it produces thresholds at which the value-Lipschitz property holds outside a set of measure
at most $|J|/4$, whatever the room there.

\begin{remark}[The sliver behind the first young row]\label{8.11:rem-young-sliver}
We use the first young row as in Definition~\ref{8.11:def-young-row-cover}, with the covering
property \eqref{8.11:eq-young-row-cover-a}, in the case $\lambda_y=w/L$, $y_-=3\lambda_y$ and
$y_+=w+5w/L$. Then
\begin{equation}\label{8.11:eq-young-sliver-1}
  \varpi_y=y_+-w=\frac{5w}{L}.
\end{equation}
The base point of the half-skin bound at the canonical comparison point is identified as a
constrained minimiser, that is, a point of the form $\hat x$, by a corollary whose proof concludes
this from the fact that $\hat x$ lies
in the closure of a set of configurations defined by the windows at threshold $\theta$. For that
conclusion every window of the outer ring of the zone, the windows of the second young row among
them, has to pass at $\hat x$. The second young row tests the rows
\begin{equation*}
  w+\frac{3w}{L}\le\delta\le 2w+\frac{5w}{L}.
\end{equation*}
Two ranges of depth inside this set need separate attention.
\begin{itemize}
\item[(a)] The first young row makes the rows $w+3w/L<\delta\le y_+$ pass only at its own cap
$\theta_y=\theta(1+\vartheta_1)\mathsf{T}$. On these rows the second row therefore reads at most
$\theta(1+\vartheta_1)^2+\vartheta_1$, not at most $\theta$. This is a sliver of relative size
$\vartheta_1$. It can be removed by imposing the first young row at the threshold
$\theta(1-3\vartheta_1)$.
\item[(b)] The rows $y_+<\delta<2w$ lie outside two regions. They are not in the region of the
regularity condition among the hypotheses of Theorem~\ref{8.11:thm-deep-curvature-blocks}; the
deep-curvature problem of that theorem begins at $\delta=y_+$. They are not in the region of the
window condition in the proof of that corollary, which begins at $\delta=2w$.
\end{itemize}
The curvature bound \eqref{8.11:eq-deep-curvature-a} gives the following on the rows in~(b). Let
$\rho$, with $0<\rho<1$, be the room fraction of the floor condition on the young cap displayed in
\eqref{8.11:eq-young-sliver-4} below. Enlarging $C_4^{\flat}$ if necessary, which only weakens the
bound \eqref{8.11:eq-small-heights-a}, we assume $C_4^{\flat}\ge 1/4$. Put
\begin{equation}\label{8.11:eq-young-sliver-2}
  C_c:=\frac{4C_4^{\flat}}{1-\rho}.
\end{equation}
Take $h=C_c\varpi_y$, and assume the hypotheses of Corollary~\ref{8.11:cor-small-heights} at this
height, among them $r_1+2\le h\le w$. Assume also the floor condition on the young cap at this
height,
\begin{equation}\label{8.11:eq-young-sliver-4}
  C_4^{\flat}\bigl[1+\log(D_W/h)\bigr](1+\vartheta_1)
  \le\tfrac12(1-\rho)(1-\vartheta_1)\,\mathsf{T}.
\end{equation}
This is a one-sided lower bound on $\mathsf{T}$.
Then every plaquette $p$ with
\begin{equation}\label{8.11:eq-young-sliver-3}
  \delta(p)\ \ge\ y_++r_1+C_c\varpi_y\ =\ y_++r_1+\frac{5C_cw}{L}
\end{equation}
satisfies
\begin{equation*}
  \|U(e\oplus\hat x)(\partial p)-1\|\le(1-\rho)\,\theta\,\mathsf{T}\,\varepsilon_{k_o}\eta^2 .
\end{equation*}
This is the young cap with room to spare. If $L\ge 10(1+C_c)$ and $r_1<w/2$, the lowest row
$\delta=y_++r_1+C_c\varpi_y$ of \eqref{8.11:eq-young-sliver-3} lies inside the range $(y_+,2w)$.

The rows of the sliver
\begin{equation*}
  y_+<\delta<y_++r_1+C_c\varpi_y
\end{equation*}
are not covered by this bound.

Theorem~\ref{8.11:thm-deep-curvature-blocks}, whose reference depth $c_*$ lies outside the box and
which does not depend on one of the inputs of Corollary~\ref{8.11:cor-small-heights}, gives the
same sliver up to the constant: its two young members, of sizes $w/h$ and $(w/h)^2$, are at most a
fixed fraction of $(1-\rho)\mathsf{T}$ once $h$ is at least a fixed multiple of $w/L$.
\end{remark}

\begin{proof}
\emph{The depth at which the bound applies.} Corollary~\ref{8.11:cor-small-heights} is taken with
the reference depth $c_*=w-1$. Its thickness is $t=d_1-c_*$, where $d_1$ is the free depth of
\eqref{8.11:eq-released-problem-a}. The corollary records $t=\varpi_y+r_1+1$. Hence
\begin{equation*}
  d_1=c_*+t=(w-1)+(\varpi_y+r_1+1)=y_++r_1 .
\end{equation*}
The corollary applies to plaquettes with $\delta(p)\ge d_1+h$. For $h=C_c\varpi_y$ this is
condition \eqref{8.11:eq-young-sliver-3}. The second form of that condition follows from
\eqref{8.11:eq-young-sliver-1}.

\emph{The young member.} For $h=C_c\varpi_y$, definition \eqref{8.11:eq-young-sliver-2} gives
\begin{equation*}
  \frac{\varpi_y}{h}=\frac1{C_c}=\frac{1-\rho}{4C_4^{\flat}} .
\end{equation*}
Since $C_4^{\flat}\ge1/4\ge(1-\rho)/4$, we have $(1-\rho)^2/(16C_4^{\flat})\le(1-\rho)/4$, and
therefore
\begin{align*}
  C_4^{\flat}\Bigl[\frac{\varpi_y}{h}+\Bigl(\frac{\varpi_y}{h}\Bigr)^2\Bigr]
  &=\frac{1-\rho}{4}+\frac{(1-\rho)^2}{16\,C_4^{\flat}}\\
  &\le\frac{1-\rho}{4}+\frac{1-\rho}{4}
   =\frac{1-\rho}{2} .
\end{align*}
This is the condition under which the young member is small.

\emph{The total bound.} Write $[\mathrm{u}]=\varepsilon_{k_o}\eta^2$. The bound
\eqref{8.11:eq-small-heights-a} splits into two parts. The first is
$C_4^{\flat}[1+\log(D_W/h)]\,\theta(1+\vartheta_1)[\mathrm{u}]$. The second is the young member,
\begin{equation*}
  C_4^{\flat}\Bigl[\frac{\varpi_y}{h}+\Bigl(\frac{\varpi_y}{h}\Bigr)^2\Bigr]
  \mathsf{T}\,\theta(1+\vartheta_1)[\mathrm{u}]
  \le\frac{1-\rho}{2}\,\mathsf{T}\,\theta(1+\vartheta_1)[\mathrm{u}] .
\end{equation*}
By the floor condition \eqref{8.11:eq-young-sliver-4}, multiplied by $\theta[\mathrm{u}]$, the first
part is at most
\begin{equation*}
  \tfrac12(1-\rho)(1-\vartheta_1)\,\theta\,\mathsf{T}[\mathrm{u}] .
\end{equation*}
Adding the two parts,
\begin{align*}
  \|U(e\oplus\hat x)(\partial p)-1\|
  &\le\tfrac12(1-\rho)(1-\vartheta_1)\,\theta\,\mathsf{T}[\mathrm{u}]
   +\tfrac12(1-\rho)(1+\vartheta_1)\,\theta\,\mathsf{T}[\mathrm{u}]\\
  &=(1-\rho)\,\theta\,\mathsf{T}[\mathrm{u}]
\end{align*}
on the rows \eqref{8.11:eq-young-sliver-3}. The same
choice $h=C_c\varpi_y$ works at every deeper row, since the corollary asks only
$\delta(p)\ge d_1+h$.

\emph{The position of the rows.} Suppose $L\ge 10(1+C_c)$ and $r_1<w/2$. Then
$5(1+C_c)w/L\le w/2$, so by \eqref{8.11:eq-young-sliver-1}
\begin{equation*}
  y_++r_1+C_c\varpi_y=w+\frac{5(1+C_c)w}{L}+r_1\le\frac{3w}{2}+r_1<2w .
\end{equation*}
Since $r_1\ge0$ and $C_c\varpi_y>0$, the same row is larger than $y_+$. Hence the lowest row of
\eqref{8.11:eq-young-sliver-3} lies inside $(y_+,2w)$.

\emph{The sliver.} At depths $y_+<\delta<y_++r_1+C_c\varpi_y$ no height $h\ge C_c\varpi_y$ with
$\delta\ge d_1+h$ is available. The computation above therefore does not apply there.
\end{proof}

\begin{remark}[Model: caps on a band do not propagate]\label{8.11:rem-band-caps}
The model is the harmonic half-plane. Consider the closed upper half-plane
$\{(s,h):\ s\in\mathbb{R},\ h\ge 0\}$ (here $s$ is a real coordinate of the model, not the depth
below the cutoff) and the bounded harmonic function $u$ on its interior whose
boundary trace at $h=0$ is
\begin{equation}\label{8.11:eq-band-caps-1}
  \varphi(s)=\mathsf{T}\cdot\max\bigl(-b,\min(s,b)\bigr),\qquad s\in\mathbb{R},
\end{equation}
with $b>0$ and $\mathsf{T}$ the young cap in exterior units. Write $F$ for the field whose
components the young caps bound. In the model the role of $F$ is played by the gradient
$(\partial_s u,\partial_h u)$ of $u$. Its tangential component at $h=0$ is $\varphi'$, which
equals $\mathsf{T}$ on $(-b,b)$ and $0$ outside $[-b,b]$. In the model the capped band is the
layer $\{(s,h):\ s\in\mathbb{R},\ -b\le h\le 0\}$ of thickness $b$ below the interface $h=0$,
and the region $h>0$ beyond it is free. The trace \eqref{8.11:eq-band-caps-1} is reachable
across the band with both components at most $\mathsf{T}$: the function
$v(s,h)=\varphi(s)\,(h+b)/b$ on the layer vanishes at $h=-b$ and equals $\varphi$ at $h=0$, and
at every $s\ne\pm b$ the field $F=(\partial_s v,\partial_h v)$ it carries there has the
components $\varphi'(s)(h+b)/b$ and $\varphi(s)/b$, both of absolute value at most $\mathsf{T}$
because $|\varphi'|\le\mathsf{T}$ and $|\varphi|\le\mathsf{T}b$. In this model the information
that $|F|\le\mathsf{T}$ holds componentwise on a band of thickness $b$, with $F$ free beyond the
band, does not imply that $|F|\le\mathsf{T}$ holds just beyond the band: at height $h>0$ the
normal derivative of $u$, taken towards the boundary, is
\begin{equation}\label{8.11:eq-band-caps-2}
  -\partial_h u(s,h)=\frac{\mathsf{T}}{2\pi}\,
  \log\frac{(s+b)^2+h^2}{(s-b)^2+h^2},
\end{equation}
and at $s=b$ this exceeds $\mathsf{T}$ as soon as
\begin{equation}\label{8.11:eq-band-caps-3}
  h<2b\,\bigl(e^{2\pi}-1\bigr)^{-1/2}=2e^{-\pi}\bigl(1-e^{-2\pi}\bigr)^{-1/2}\,b .
\end{equation}
The constant $2e^{-\pi}(1-e^{-2\pi})^{-1/2}$ exceeds $2e^{-\pi}$ by less than one part in a
thousand.

The model is a caricature. Its purpose is to locate a mechanism, and nothing below rests on it
as an estimate. It shows the following. Behind any capped band there is a sliver, of thickness
proportional to $b$, on which the bound $|F|\le\mathsf{T}$ does not follow from the caps alone;
in the model it is the strip
\begin{equation*}
  0<h<2e^{-\pi}\bigl(1-e^{-2\pi}\bigr)^{-1/2}\,b
\end{equation*}
over the point $s=b$ of the interface. Adding the next
young row of cubes to the young family $\mathfrak{Y}$ enlarges $b$ to $b+w$, where $w$ is
the clamp width, and the sliver grows with it. Write $c\,b$, with a constant $c>0$, for the
thickness of the sliver behind a band of thickness $b$; in the model
$c=2e^{-\pi}(1-e^{-2\pi})^{-1/2}$. Call a round of the scheme described below the imposition of
further young rows until the sliver behind the current band is covered, and write $b_n$ for the
band thickness after $n$ rounds. The band after a round contains the band before it together
with the sliver behind it, so
\begin{equation}\label{8.11:eq-band-caps-4}
  b_{n+1}\ge(1+c)\,b_n ,
\end{equation}
and $b_n\to\infty$.
Consequently the scheme ``impose further young rows as constraints until the remaining
windows pass'' (a window passes when the cap inequality holds on it at the point read)
diverges, and it diverges at every ring of outer windows.

What closes the scheme is a statement about the minimiser itself. Let $\hat x$ be the canonical
comparison point. Let $\varpi_y=y_+-w$, with $y_+$ the upper end of
the depth range of the first young row. Let $r_1$ and $C_c$ be the depth offset and the
constant that fix the width $r_1+C_c\varpi_y$ of the sliver behind the first young row in
Remark~\ref{8.11:rem-young-sliver}. The statement needed is:
\begin{equation}\label{8.11:eq-band-caps-5}
  \begin{gathered}
  \text{at } \hat x, \text{ no plaquette of the first young row that lies within depth}\\
  r_1+C_c\varpi_y \text{ of the row's inner face is binding.}
  \end{gathered}
\end{equation}
This is the inactivity statement for the first young row. It is not proved in this remark.
Suppose \eqref{8.11:eq-band-caps-5} holds. Then three consequences follow.
\begin{itemize}
  \item[(a)] $\hat x$ is free there, up to the squeeze term.
  \item[(b)] The estimate used in Remark~\ref{8.11:rem-young-sliver} to bound $F$ on the rows
  beyond the sliver applies with $y_+$ lowered.
  \item[(c)] The rows up to depth $2w$ are read at the logarithmic scale of the corollary of
  that estimate which Remark~\ref{8.11:rem-young-sliver} uses to bound $F$ on the rows beyond
  the sliver.
\end{itemize}
Statement \eqref{8.11:eq-band-caps-5} is equivalent to the assertion that the young core (the
central part of the young region, whose radius in the kink caricature is given below) rounds
and stays at the clamp's face. In the kink caricature, the
caricature in which the exterior field is a kink of amplitude $\sigma$ (the exterior kink
amplitude of that caricature), the core has radius
\begin{equation}\label{8.11:eq-band-caps-6}
  h_1=D_W\,e^{-\pi\theta\mathsf{T}/(2\sigma)}\ll w/L^2 ,
\end{equation}
with $D_W$ the width of the zone and $\theta$ the threshold.

The same corollary, applied at the straddling-only minimiser $\tilde x$, supports
\eqref{8.11:eq-band-caps-5}. Indeed, $\tilde x$ exceeds the cap at most in the sliver
\begin{equation*}
  w<\delta<w+r_1+w/L^2 ,
\end{equation*}
where $\delta$ is the signed depth. It does not prove \eqref{8.11:eq-band-caps-5} at $\hat x$,
because no comparison principle between $\tilde x$ and $\hat x$ is available in
Ba{\l}aban's setting.
\end{remark}

\begin{proof}[Verification of \eqref{8.11:eq-band-caps-2} and \eqref{8.11:eq-band-caps-3} in the
model]
Write $z=s+ih$ with $h>0$. Put
\begin{equation*}
  G(z)=\frac{i\mathsf{T}}{\pi}\,\log\frac{z+b}{z-b},
\end{equation*}
with the principal branch of the logarithm.

\emph{$G$ is holomorphic.} Since
\begin{equation*}
  \operatorname{Im}\frac{z+b}{z-b}=\frac{-2bh}{|z-b|^2}<0,
\end{equation*}
the quotient stays in the open lower half-plane. Hence $G$ is holomorphic for $h>0$.

\emph{The real part of $G$.} Let $\arg(z\pm b)\in(0,\pi)$. Then
\begin{equation*}
  \operatorname{Re}G=\frac{\mathsf{T}}{\pi}\bigl(\arg(z-b)-\arg(z+b)\bigr)\in(0,\mathsf{T}).
\end{equation*}
As $h\downarrow0$ this tends to $\mathsf{T}$ for $|s|<b$ and to $0$ for $|s|>b$. So
$\operatorname{Re}G$ is the bounded harmonic extension of $\varphi'=\mathsf{T}\,1_{(-b,b)}$.

\emph{Identification with $\partial_s u$.} The function $u$ is the Poisson integral of the
bounded Lipschitz function $\varphi$. Differentiating under the integral shows that
$\partial_s u$ is the Poisson integral of $\varphi'$. Hence $\partial_s u=\operatorname{Re}G$.

\emph{Identification of the imaginary part.} By the Cauchy--Riemann equations,
$\partial_s u-i\partial_h u$ is holomorphic. It has the same real part as $G$, so the two
functions differ by an imaginary constant. To fix the constant, note that $\varphi$ is odd and
the Poisson kernel is even in $s$. Hence $u(0,h)=0$ for every $h$, and so $\partial_h u(0,h)=0$.
Also $\operatorname{Im}G(ih)=0$, because $|ih+b|=|ih-b|$. The constant is therefore zero. It
follows that
\begin{equation*}
  -\partial_h u=\operatorname{Im}G
  =\frac{\mathsf{T}}{\pi}\log\frac{|z+b|}{|z-b|},
\end{equation*}
which is \eqref{8.11:eq-band-caps-2}.

\emph{The threshold.} At $s=b$, \eqref{8.11:eq-band-caps-2} reads
\begin{equation*}
  -\partial_h u(b,h)=\frac{\mathsf{T}}{2\pi}\log\Bigl(1+\frac{4b^2}{h^2}\Bigr).
\end{equation*}
This exceeds $\mathsf{T}$ if and only if $4b^2/h^2>e^{2\pi}-1$, that is, if and only if
$h<2b\,(e^{2\pi}-1)^{-1/2}$. Finally,
\begin{equation*}
  2b\,(e^{2\pi}-1)^{-1/2}=2e^{-\pi}b\,(1-e^{-2\pi})^{-1/2}.
\end{equation*}
This is \eqref{8.11:eq-band-caps-3}.
\end{proof}

\begin{remark}[Box holes and reentrant edges]\label{8.11:rem-reentrant-edges}
Lemma~\ref{8.11:lem-torus-reflection} (reflection onto the doubled torus) applies only when the hole is a box.
The logarithmic dependence on the extent in Corollary~\ref{8.11:cor-arbitrary-extent} (box holes of
arbitrary extent) is likewise stated only for a box hole of bounded aspect.
For unions of boxes having reentrant edges, the statement on power laws at reentrant edges gives a power
law in the extent in place of the logarithm. The lower bound at reentrant edges says the following for an
L-shaped hole of the class treated there: every constant in a bound that includes the edges is at least
$c\,\mathsf{p}^{a'}$, with an exponent $a'<1/3$, where $c$ is a constant of that statement and
$\mathsf{p}$ is the parameter of that statement.
Hence the logarithmic form of Corollary~\ref{8.11:cor-arbitrary-extent} is not to be expected at a reentrant
edge, even at depth $2w$, where $w$ is the clamp width. Any extension of that corollary to unions of boxes
therefore has to carry either a condition on the shape of the union or a power of the extent in place of
the logarithm.
\end{remark}

\begin{proof}
Lemma~\ref{8.11:lem-torus-reflection} is stated and proved for a box hole only.

We now turn to reentrant edges. Let $D_u$ be the extent of the union, and let $\varphi>\pi$ be the interior
opening angle at a reentrant edge. Let $r_e$ be the distance to that edge. The leading corner mode at that
edge has radial profile $r_e^{\pi/\varphi-1}$. Normalising the mode at the scale $D_u$ of the union and
reading it at $r_e\asymp w$ gives the factor
\begin{equation}\label{8.11:eq-reentrant-edges-1}
\Bigl(\frac{D_u}{w}\Bigr)^{1-\pi/\varphi}.
\end{equation}
This is a power of $D_u/w$ and does not involve the lattice cutoff. Hence a dependence on the extent that
is only logarithmic cannot persist at a reentrant edge; the power-law bound itself is the content of the
statement on power laws at reentrant edges recalled above.
\end{proof}

\begin{remark}[Notions of non-triviality for the curvature two-point function]
\label{8.12:rem-nontriviality-notions}
The word \emph{non-trivial} is used in this book only in the sense of
Definition~\ref{3.1:def-non-trivial}. This remark lists the notions of non-triviality that can be
formulated for the curvature observables of Theorem~\ref{3.1:thm-main}, namely smeared action
densities with pairwise separated supports. Where a notion excludes a class of candidate limits,
it says which.

Fix a point $x_0$ of physical space (of $\mathbb{T}_m$, or of $\mathbb{R}^4$ in the
infinite-volume limit) and a \emph{reference pair} $(\varphi_1,\varphi_2)$ of test functions,
chosen so that the dilates below belong to the class $\mathfrak{T}(\mathsf{d};\vartheta)$ of the
two-point witness theorem. Their \emph{sup-normalised dilates} at scale $\mathsf{d}>0$ are
\begin{equation}\label{8.12:eq-nontriviality-notions-1}
  f_i^{(\mathsf{d})}(x):=\varphi_i\bigl((x-x_0)/\mathsf{d}\bigr),\qquad i=1,2,
\end{equation}
so that $\|f_i^{(\mathsf{d})}\|_\infty=\|\varphi_i\|_\infty$ for every $\mathsf{d}$.

We write $S^T_{2;K}$ for the connected two-point function at $K$ renormalisation steps, with the
torus exponent $m$ suppressed from the notation. We write $S^T_2$ for the connected two-point
function in the limit considered ($S^T_{2;m}$ or $S^T_{2;\infty}$).

With this normalisation the kernel $\mathcal{K}_s(f_1^{(\mathsf{d})},f_2^{(\mathsf{d})})$ is
bounded above and below by positive constants independent of $\mathsf{d}$; we say that it is
\emph{scale-free}. The statement that the two-point function is ``of order $g^4$'' means that
$S^T_2(f_1^{(\mathsf{d})},f_2^{(\mathsf{d})})$ is bounded above and below by constant multiples
of $g^4$, with no power of $\mathsf{d}$.

By a \emph{scale-covariant assignment} of dimension $\Delta$ we mean a translation-invariant
two-point function $G$ of the action density with
$G(\mathsf{d}u,\mathsf{d}v)=\mathsf{d}^{-2\Delta}G(u,v)$.

The notions are the following.
\begin{itemize}
\item[(n0)] \emph{Flow.} The running coupling is controlled from both sides, uniformly in the
  cutoff, at every scale, and the effective action at level $k$ is $-g_k^{-2}A$ plus irrelevant
  terms. This is the content of clauses (0) and (i) of Theorem~\ref{3.1:thm-main}. It involves
  neither the rarity of young large-field blocks nor any bound uniform in the volume.
\item[(n1)] \emph{Neither free nor degenerate.} As $\mathsf{d}\downarrow0$, with no rate
  specified,
  \begin{equation}\label{8.12:eq-nontriviality-notions-2}
    \limsup_{K\to\infty}S^T_{2;K}\bigl(f_1^{(\mathsf{d})},f_2^{(\mathsf{d})}\bigr)\to0 ,
  \end{equation}
  and there is a pair $(f_1,f_2)$ with $S^T_2(f_1,f_2)\neq0$ in the limit considered.
  This excludes a free Maxwell field at any charge and a degenerate limit. It does not exclude a
  scale-covariant assignment of dimension $\Delta<4$, and it fixes no
  rate, so it gives no logarithm.
\item[(n2$\flat$)] \emph{Two-sided logarithmic law.} There are positive constants $c,C,B,\lambda$
  and a scale $\mathsf{d}_W>0$ such that for
  $\mathsf{d}\le\mathsf{d}_W$
  \begin{equation}\label{8.12:eq-nontriviality-notions-3}
  \begin{split}
    c\bigl(g^{-2}+B\log_L(1/\mathsf{d})\bigr)^{-2}
    &\le\frac{S^T_2(f_1^{(\mathsf{d})},f_2^{(\mathsf{d})})}
            {\|\varphi_1\|_\infty\|\varphi_2\|_\infty}\\
    &\le C\bigl(g^{-2}+\lambda\log_L(1/\mathsf{d})\bigr)^{-2}.
  \end{split}
  \end{equation}
  This excludes every scale-covariant assignment and every free Maxwell field, and it gives
  decay of the effective coupling at scale $\mathsf{d}$ to zero, logarithmically, with bounds on
  both sides.
\item[(n2)] \emph{The form of clause \textup{(iv)} of Theorem~\ref{3.1:thm-main}.} This is
  \textup{(n2$\flat$)} together with the structure
  \begin{equation}\label{8.12:eq-nontriviality-notions-4}
  \begin{split}
    S^T_{2;K}\bigl(f_1^{(\mathsf{d})},f_2^{(\mathsf{d})}\bigr)
    &=b_0\,g^4_{K-s(\mathsf{d})}\,
      \mathcal{K}_s\bigl(f_1^{(\mathsf{d})},f_2^{(\mathsf{d})}\bigr)\,(1+R),\\
    |R|&\le\tfrac12 .
  \end{split}
  \end{equation}
  Here $b_0=(N^2-1)/2$ is exact up to the factor $1+R$.
\item[(n2$\sharp$)] \emph{Sharp asymptotic freedom.} As $\mathsf{d}\downarrow0$ one has $R\to0$,
  and
  \begin{equation}\label{8.12:eq-nontriviality-notions-5}
    S^T_2\bigl(f_1^{(\mathsf{d})},f_2^{(\mathsf{d})}\bigr)\,(\log 1/\mathsf{d})^2
    \to b_0\hat{\mathcal{K}}(\varphi_1,\varphi_2)/(2\beta_0)^2 ,
  \end{equation}
  where $\hat{\mathcal{K}}=\hat{\mathcal{K}}(\varphi_1,\varphi_2)$ denotes the continuum kernel of
  the two-point witness theorem and $\beta_0$ denotes the one-loop coefficient of perturbation
  theory, in the convention in which $g^{-2}$ at scale $\mathsf{d}$ grows like
  $2\beta_0\log(1/\mathsf{d})$.
  For the dilates $\mathsf{d}^{-4}f_i^{(\mathsf{d})}$, normalised in $L^1$ instead of in sup
  norm, the two-point function is $\mathsf{d}^{-8}S^T_2(f_1^{(\mathsf{d})},f_2^{(\mathsf{d})})$
  by bilinearity, so, written for these $L^1$-normalised dilates, the left-hand side of
  \eqref{8.12:eq-nontriviality-notions-5} is
  \begin{equation*}
    \mathsf{d}^{8}\,S^T_2\bigl(\mathsf{d}^{-4}f_1^{(\mathsf{d})},
      \mathsf{d}^{-4}f_2^{(\mathsf{d})}\bigr)\,(\log 1/\mathsf{d})^2 .
  \end{equation*}
  The notion \textup{(n2$\sharp$)} is agreement with the leading prediction of perturbative
  renormalisation theory; Theorem~\ref{3.1:thm-main} does not assert it. Among the statements that
  clause (iv) does not assert is the sharp form $b_0g^4(1+O(g^\sigma))$ with some
  $\sigma>0$. Clause (0) asserts no one-loop slope.
\item[(n3)] \emph{Non-degeneracy.} There is $f$ with $S^T_2(\mathfrak{r}f,f)>0$. Here
  $\mathfrak{r}$ is the reflection in the hyperplane of reflection positivity
  (Lemma~\ref{1.1:lem-invariance}), acting on test
  functions by composition, and $f$ is supported on one side of that hyperplane, at positive
  distance from it. By reflection positivity this is equivalent to
  $\dim\mathcal{H}^{\mathrm{curv}}\ge2$ (the corollary stated beside the two-point witness
  theorem).
\item[(n4)] \emph{Other notions used for gauge theories.} These are: an area law, or positive
  string tension; a mass gap; a scattering matrix different from the identity; Borel
  summability; a non-vanishing connected four-point function of the action density; and the
  property of not being a Gaussian connection with isotropic covariance. The first four are not
  asserted in this book. The fifth does not serve as a criterion here, since free Maxwell fields
  of any charge already have non-vanishing connected correlations of observables quadratic in the
  field (see the proof). The last follows from \textup{(n2)} by the statement that no
  reflection-positive Gaussian connection with isotropic covariance satisfies \textup{(n2)}; it
  refines the exclusion of free Maxwell fields and is not stated separately. The bare property of
  not being Gaussian is among the statements clause (iv) does not assert.
\end{itemize}
\end{remark}

\begin{proof}
\emph{Scale-freeness.} The action density has dimension four, and the free kernel is
$|x-y|^{-8}$. The change of variables $x=x_0+\mathsf{d}u$, $y=x_0+\mathsf{d}v$ gives
\begin{equation*}
  \iint\varphi_1\Bigl(\frac{x-x_0}{\mathsf{d}}\Bigr)\varphi_2\Bigl(\frac{y-x_0}{\mathsf{d}}\Bigr)
  \frac{dx\,dy}{|x-y|^{8}}
  =\iint\varphi_1(u)\varphi_2(v)\frac{du\,dv}{|u-v|^{8}} .
\end{equation*}
The factor $\mathsf{d}^8$ of the Jacobian cancels the factor $\mathsf{d}^{-8}$ of the kernel, so
the right-hand side is free of $\mathsf{d}$. For the kernel $\mathcal{K}_s$ of the two-point
witness theorem, the kernel bounds accompanying that theorem give, for the reference pairs
considered there,
\begin{equation*}
  \mathcal{K}_s\bigl(f_1^{(\mathsf{d})},f_2^{(\mathsf{d})}\bigr)
  \in\bigl[c_{\mathcal{K}}\vartheta^2,\,100\,C_{\mathcal{K}}\bigr]
  \cdot\|\varphi_1\|_\infty\|\varphi_2\|_\infty .
\end{equation*}
This interval does not depend on $\mathsf{d}$.

\emph{Scale-covariant assignments.} Let $G$ be a scale-covariant assignment of dimension
$\Delta$. The smeared two-point function is
$S^T_2(f_1,f_2)=\iint f_1(x)f_2(y)\,G(x,y)\,dx\,dy$. The change of variables above then gives
\begin{equation}\label{8.12:eq-nontriviality-notions-6}
  S^T_2\bigl(f_1^{(\mathsf{d})},f_2^{(\mathsf{d})}\bigr)=c_\Delta\,\mathsf{d}^{8-2\Delta},
  \qquad c_\Delta:=S^T_2\bigl(\varphi_1(\cdot-x_0),\varphi_2(\cdot-x_0)\bigr).
\end{equation}
The case $\Delta=4$ is the scale-free case just computed.

\emph{What \textup{(n1)} excludes.} A degenerate limit has $S^T_2\equiv0$ and violates the second
condition of \textup{(n1)}. A free Maxwell field at charge $e$ gives the value
$e^4b_0\hat{\mathcal{K}}(\varphi_1,\varphi_2)(1+o(1))$, constant in $\mathsf{d}$. This is the free
Maxwell computation accompanying the two-point witness theorem, with $\hat{\mathcal{K}}$ the
continuum kernel. The constant $\hat{\mathcal{K}}(\varphi_1,\varphi_2)$ is positive for reference
pairs of the class fixed above, by the lower bound of the continuum kernel stated with that
theorem; so the value does not tend to zero, and it violates
\eqref{8.12:eq-nontriviality-notions-2}.

\textup{(n1)} does not exclude an assignment \eqref{8.12:eq-nontriviality-notions-6} with
$\Delta<4$. There $c_\Delta\mathsf{d}^{8-2\Delta}\to0$, which is compatible with an upper bound
and with positivity at one scale. Nor does \textup{(n1)} give a logarithm, since it fixes no rate.

\emph{What \textup{(n2$\flat$)} excludes.} Let \eqref{8.12:eq-nontriviality-notions-6} hold.
\begin{itemize}
\item If $\Delta<4$, the right-hand side tends to zero (like a positive power of $\mathsf{d}$ if
  $c_\Delta\neq0$).
  This is faster than the left-hand member of \eqref{8.12:eq-nontriviality-notions-3}, which
  decays only like an inverse square logarithm. So the lower bound fails for small
  $\mathsf{d}$: the rate is wrong.
\item If $\Delta=4$, the value is the constant $c_4$. For $c_4>0$ the upper bound, which
  tends to zero, fails for small $\mathsf{d}$; for $c_4\le0$ the positive lower bound fails.
\item If $\Delta>4$: for $c_\Delta>0$ the value is unbounded and the upper bound fails; for
  $c_\Delta\le0$ the positive lower bound fails.
\end{itemize}
The free Maxwell value is constant in $\mathsf{d}$, and is excluded in the same way as the case
$\Delta=4$.

By Definition~\ref{3.1:def-non-trivial}, the normalised two-point function is the fourth power
of an effective coupling $g_{\mathrm{eff}}(\mathsf{d})$ at scale $\mathsf{d}$ times a
dilation-invariant kernel factor, and for the dilates fixed above this factor is scale-free in
the sense fixed before the list of notions (bounded above and below by positive constants
independent of $\mathsf{d}$). Dividing
\eqref{8.12:eq-nontriviality-notions-3} by this factor and taking fourth roots gives positive
constants $c',C'$, independent of $\mathsf{d}$, with
\begin{equation*}
  c'\bigl(g^{-2}+B\log_L(1/\mathsf{d})\bigr)^{-1/2}
  \le g_{\mathrm{eff}}(\mathsf{d})
  \le C'\bigl(g^{-2}+\lambda\log_L(1/\mathsf{d})\bigr)^{-1/2}.
\end{equation*}
Hence the effective coupling tends to zero logarithmically, with bounds on both sides.

\emph{\textup{(n3)}.} By reflection positivity, \textup{(n3)} is equivalent to
$\dim\mathcal{H}^{\mathrm{curv}}\ge2$. This equivalence is a corollary stated beside the two-point
witness theorem.

\emph{\textup{(n4)}.} A non-vanishing connected four-point function of the action density does
not distinguish the limit from a free Maxwell field. The reason is that free Maxwell fields of
any charge already have non-vanishing
connected correlations of observables quadratic in the field; this is the computation of the
cumulants of quadratic observables of independent free Maxwell fields.

By the statement that no reflection-positive Gaussian connection with isotropic covariance
satisfies \textup{(n2)}, the property of not being such a connection is implied by
\textup{(n2)}; it is a refinement of the exclusion of free Maxwell fields.
\end{proof}

\begin{definition}[Decomposition of the two-point function into main term and remainders]
\label{8.12:not-two-point-decomposition}
Let $s$ be the depth below the cutoff and put $k_0:=K-s$, the observation level. Let $f_1,f_2$ be
test functions with supports $S_1,S_2$, and let $S^T_{2;K}=S^T_{2;K,m}$ (we suppress $m$) be the
connected two-point function of the observables $\mathcal{O}(f_1)$ and $\mathcal{O}(f_2)$ of
Definition~\ref{1.3:def-curvature-field}. This is the covariance of $\mathcal{O}(f_1)$ and
$\mathcal{O}(f_2)$ under the lattice measure of Definition~\ref{1.1:def-wilson-action}. We use the
following notation.
\begin{itemize}
\item[(a)] $h$ is a history at level $k_0$, and $V=V_{k_0}$ is the block field at level $k_0$
(Definition~\ref{1.5:def-block-average}). We write $\tilde{\nu}$ for the \emph{history law}: the law
of $(h,V)$ on the level-$k_0$ history space in the representation of the lattice expectation over
that space used in the proof of the two-point witness theorem. $\mathbb{E}_h$ is the positive
normalised plain integral over the live components of $h$.
\item[(b)] $\Phi_i^h:=\mathrm{E}[\mathcal{O}(f_i)\mid h,V]$ is the conditional mean. $\Psi^h$ is the
conditional covariance of $\mathcal{O}(f_1)$ and $\mathcal{O}(f_2)$ given $(h,V)$.
\item[(c)] By parts (a)--(d) of the source lemma, in the form in which it enters the proof of the
two-point witness theorem, the conditional mean decomposes as
\begin{equation*}
\Phi_i^h=c^0_i+G_i(V)+\mathfrak{L}^h_i(V),
\qquad
G_i:=N_{k_0}\,\mathcal{O}^{(k_0)}(\tilde f_i)+\mathrm{Irr}_i .
\end{equation*}
Here $c^0_i$ is a constant, and $\mathcal{O}^{(k_0)}$, $\tilde f_i$ and $\mathrm{Irr}_i$ are the
objects so denoted in the source lemma. The function $G_i$ does not depend on the
history. On the window of the two-point witness theorem it satisfies
$|G_i|\le C_G\,\mathfrak{n}(f_i)\,\varepsilon^2$, where $\varepsilon$ is the small parameter of that
theorem and $C_G$ is a constant of that theorem. A corollary in the proof of the two-point witness
theorem bounds the function $\Xi_i$ attached to $f_i$ that is so denoted there by
$|\Xi_i|\le C_X\,\mathfrak{n}(f_i)/\vartheta$ everywhere, not only on the window, with $C_X$ a
constant of that theorem.
\item[(d)] By the same parts of the source lemma, the live part is
$\mathfrak{L}^h_i=\mathfrak{L}^{\mathrm{near},h}_i+\mathfrak{F}^h_i$. It vanishes unless some live
component of $h$, created at any step, lies near $S_i$. Its far part always satisfies
$\sup|\mathfrak{F}^h_i|\le\mathfrak{f}_i$. We write $\mathfrak{L}_i$ for $\mathfrak{L}^h_i$ regarded
as a function of $(h,V)$.
\item[(e)] $\Psi^{\mathrm{span},h}$ is the term of $\Psi^h$ so denoted in the source lemma, and
$\mathfrak{A}^w_h$ is the sourced effective action on the history $h$, with complex source
$w=z_1f_1+z_2f_2$; $\partial_{z_j}\mathfrak{A}^w_h|_0$ denotes its derivative in $z_j$ evaluated at
$z_1=z_2=0$.
\item[(f)] The main term is $\mathfrak{M}:=N^2b_0\,g^4\,\mathcal{K}_s(f_1,f_2)$. The remainders are
\begin{align*}
R^{\mathfrak{L}}&:=\operatorname{Cov}_{\tilde{\nu}}(\mathfrak{L}_1,G_2)
+\operatorname{Cov}_{\tilde{\nu}}(G_1,\mathfrak{L}_2)
+\operatorname{Cov}_{\tilde{\nu}}(\mathfrak{L}_1,\mathfrak{L}_2),\\
R^{\Psi,\mathrm{live}}&:=\mathrm{E}_{\tilde{\nu}}\operatorname{Cov}_{\mathbb{E}_h}
\bigl(\partial_{z_1}\mathfrak{A}^w_h|_0,\partial_{z_2}\mathfrak{A}^w_h|_0\bigr),
\qquad
R^{\mathrm{loc}}:=\mathrm{E}_{\tilde{\nu}}\Psi^{\mathrm{span},h}.
\end{align*}
The remainders $R^{\mathrm{prof}},R^{GX},R^{\mathrm{irr}}$ are those of the Gaussian extraction, for
which
\begin{equation}\label{8.12:eq-two-point-decomposition-1}
\operatorname{Cov}_{\tilde{\nu}}(G_1,G_2)=\mathfrak{M}+R^{\mathrm{prof}}+R^{GX}+R^{\mathrm{irr}} .
\end{equation}
\end{itemize}
With this notation,
\begin{equation}\label{8.12:eq-two-point-decomposition-a}
\begin{gathered}
S^T_{2;K}=\mathrm{E}_{\tilde{\nu}}[\Psi^h]+\operatorname{Cov}_{\tilde{\nu}}(\Phi_1^h,\Phi_2^h),\\
S^T_{2;K}-\mathfrak{M}=R^{\mathrm{prof}}+R^{GX}+R^{\mathrm{irr}}+R^{\mathrm{loc}}
+R^{\mathfrak{L}}+R^{\Psi,\mathrm{live}} .
\end{gathered}
\end{equation}
In the proof of the two-point witness theorem the six remainders are bounded, each under the
hypotheses of that theorem; assembled there, these bounds give
$|S^T_{2;K}-\mathfrak{M}|\le 6\mathfrak{M}/32$.
\end{definition}

\begin{proof}
\emph{The first identity.} Write $F_i:=\mathcal{O}(f_i)$. Each $F_i$ is bounded, since the torus is
finite and the plaquette terms are bounded. The first identity is total covariance on the
level-$k_0$ history space. It rests on the representation of the lattice expectation over that
space used in the proof of the two-point witness theorem: for every bounded function $F$ of the
configuration,
\begin{equation*}
\mathrm{E}[F]=\mathrm{E}_{\tilde{\nu}}\bigl[\mathrm{E}[F\mid h,V]\bigr],
\end{equation*}
where $\mathrm{E}[F\mid h,V]$ is the normalised conditional integral at $(h,V)$ and the outer
expectation is taken under the history law $\tilde{\nu}$. Applying this to $F_1F_2$, $F_1$ and $F_2$
gives
\begin{equation*}
\operatorname{Cov}(F_1,F_2)
=\mathrm{E}_{\tilde{\nu}}\bigl[\mathrm{E}[F_1F_2\mid h,V]\bigr]
-\mathrm{E}_{\tilde{\nu}}[\Phi_1^h]\,\mathrm{E}_{\tilde{\nu}}[\Phi_2^h].
\end{equation*}
Since $\operatorname{Cov}(F_1,F_2\mid h,V)=\mathrm{E}[F_1F_2\mid h,V]-\Phi_1^h\Phi_2^h$, adding and
subtracting $\mathrm{E}_{\tilde{\nu}}[\Phi_1^h\Phi_2^h]$ yields
\begin{equation*}
\operatorname{Cov}(F_1,F_2)
=\mathrm{E}_{\tilde{\nu}}\bigl[\operatorname{Cov}(F_1,F_2\mid h,V)\bigr]
+\operatorname{Cov}_{\tilde{\nu}}(\Phi_1^h,\Phi_2^h).
\end{equation*}
The left side is $S^T_{2;K}$ and the conditional covariance is $\Psi^h$. This gives the first line
of \eqref{8.12:eq-two-point-decomposition-a}.

\emph{The conditional-mean term.} Covariance is bilinear and vanishes whenever one argument is
constant. Substituting $\Phi_i^h=c^0_i+G_i+\mathfrak{L}^h_i$ therefore gives
\begin{align*}
\operatorname{Cov}_{\tilde{\nu}}(\Phi_1^h,\Phi_2^h)
&=\operatorname{Cov}_{\tilde{\nu}}(G_1,G_2)+\operatorname{Cov}_{\tilde{\nu}}(\mathfrak{L}_1,G_2)
+\operatorname{Cov}_{\tilde{\nu}}(G_1,\mathfrak{L}_2)
+\operatorname{Cov}_{\tilde{\nu}}(\mathfrak{L}_1,\mathfrak{L}_2)\\
&=\operatorname{Cov}_{\tilde{\nu}}(G_1,G_2)+R^{\mathfrak{L}} .
\end{align*}
By the Gaussian extraction, \eqref{8.12:eq-two-point-decomposition-1} holds. Hence
\begin{equation*}
\operatorname{Cov}_{\tilde{\nu}}(\Phi_1^h,\Phi_2^h)
=\mathfrak{M}+R^{\mathrm{prof}}+R^{GX}+R^{\mathrm{irr}}+R^{\mathfrak{L}} .
\end{equation*}

\emph{The conditional-covariance term.} By parts (a)--(d) of the source lemma, in the form in which
it enters the proof of the two-point witness theorem, for every history $h$,
\begin{equation*}
\Psi^h=\Psi^{\mathrm{span},h}
+\operatorname{Cov}_{\mathbb{E}_h}\bigl(\partial_{z_1}\mathfrak{A}^w_h|_0,
\partial_{z_2}\mathfrak{A}^w_h|_0\bigr).
\end{equation*}
Taking $\mathrm{E}_{\tilde{\nu}}$ of both sides and using the definitions in (f), we obtain
$\mathrm{E}_{\tilde{\nu}}\Psi^h=R^{\mathrm{loc}}+R^{\Psi,\mathrm{live}}$.

\emph{Conclusion.} Insert the two preceding results into the first line of
\eqref{8.12:eq-two-point-decomposition-a} and subtract $\mathfrak{M}$. This gives the second line.
\end{proof}

\begin{remark}[Remainders sensitive to rarity and their weight on old live regions]
\label{8.12:rem-rarity-weights}
We work at the observation level $k_0=K-s$. We use the decomposition of
Notation~\ref{8.12:not-two-point-decomposition}, in which:
\begin{itemize}
\item $\tilde{\nu}$ is the history law at level $k_0$;
\item $\mathbb{E}_h$ is the positive normalised plain integral of the live components of a history $h$;
\item the connected two-point function is written as in
\eqref{8.12:eq-two-point-decomposition-a}, as the main term $\mathfrak{M}$ plus the remainders
$R^{\mathrm{prof}}$, $R^{GX}$, $R^{\mathrm{irr}}$, $R^{\mathrm{loc}}$, $R^{\mathfrak{L}}$ and
$R^{\Psi,\mathrm{live}}$.
\end{itemize}
A live component $Z$ created at step $j(Z)$ has age $a=k_0-j(Z)$ at level $k_0$. For $i=1,2$ let
$A_i(h)$ be the age of the oldest live component of $h$ near the support of $f_i$. Let $A(h)$ denote
the age parameter of the history $h$ in terms of which the two-slot bound
\eqref{8.12:eq-rarity-weights-3} below is stated; it satisfies $A_1(h)+A_2(h)\le 2A(h)$.
The members of the decomposition, together with the two bad-event probabilities
$\mathfrak{q}_W=\tilde{\nu}(\text{bad event})$ and $\mathfrak{q}_{\times}$, use a normalised rarity
bound for large-field regions as follows.
\begin{enumerate}
\item[(1)] $\mathfrak{M}$ and $R^{\mathrm{prof}}$ use no rarity bound. Their bounds rest on the
kernel statement for $\mathcal{K}_s$ and on the lemma of this chapter which, together with that
kernel statement, supplies the estimates of $\mathfrak{M}$ and $R^{\mathrm{prof}}$.
\item[(2)] $R^{GX}$, the first remainder of the Gaussian extraction, uses rarity through the weight
of the unclean event $\mathcal{C}^c$. This weight is a one-point quantity, at every level
$j\le k_0+r(s)$, where $r(s)$ is the extraction depth of clause~(iv) of
Theorem~\ref{3.1:thm-main}, near the extraction box of
the Gaussian extraction. It enters through the small-piece statement that accompanies the
localisation of the sourced step. The integrand on the unclean event is bounded, using only the
$\eta^2$-law part of the bound on the quantities $X_i$ of
Notation~\ref{8.12:not-two-point-decomposition}:
\begin{equation}\label{8.12:eq-rarity-weights-1}
|X_1X_2|\le C_X^2\,\mathfrak{n}(f_1)\,\mathfrak{n}(f_2)/\vartheta^2 .
\end{equation}
An old live component receives no extra weight, so the bound is blind to age. The one-point
rarity bound used at each such level is obtained as an implication from the corollary on the
rarity of young large-field blocks.
\item[(3)] $R^{\mathrm{irr}}$, the second remainder of the Gaussian extraction, is treated exactly
as $R^{GX}$. It has a bounded integrand and no weight on age.
\item[(4)] $R^{\mathrm{loc}}$ uses no rarity bound. It consists of deterministic spanning terms,
present for every history $h$, and is bounded by the correlation theorem.
\item[(5)] $\mathfrak{q}_W$ and $\mathfrak{q}_{\times}$ use rarity as a one-point live weight,
for live components of any birth. They multiply factors that are bounded in sup norm, and they
carry no weight on age. The bound used is the one-point rarity bound at level $k_0$, obtained from
the corollary on the rarity of young large-field blocks; it is an input of the additive theorem.
\item[(6)] The parts $\operatorname{Cov}(\mathfrak{L}_1,G_2)$ and
$\operatorname{Cov}(G_1,\mathfrak{L}_2)$ of $R^{\mathfrak{L}}$ use rarity. On the event that the
oldest near live age is $A_i=A_i(h)$, the near live part satisfies
\begin{equation}\label{8.12:eq-rarity-weights-2}
|\mathfrak{L}^{\mathrm{near}}_i|\le\mathfrak{s}_i^{(A_i)},
\qquad
\mathfrak{s}_i^{(A_i)}:=\mathfrak{s}_i^{\sharp}\,L^{4(A_i-\mathfrak{l})}.
\end{equation}
An old live component of age $a$ therefore enters with weight $L^{4a}$, from one source slot.
These parts rest on the age-weighted one-point rarity bound at level $k_0$ supplied by the first
step of the additive theorem.
\item[(7)] The part $\operatorname{Cov}(\mathfrak{L}_1,\mathfrak{L}_2)$ of $R^{\mathfrak{L}}$, and
$R^{\Psi,\mathrm{live}}$, use rarity. On one history $h$ their integrand on the rare event is
bounded by the product of the two slot sizes, and
\begin{equation}\label{8.12:eq-rarity-weights-3}
\mathfrak{s}_1^{(A_1)}\mathfrak{s}_2^{(A_2)}
\le\mathfrak{s}_1^{\sharp}\mathfrak{s}_2^{\sharp}\,L^{8(A(h)-\mathfrak{l})}.
\end{equation}
An old live component of age $a$ therefore enters with weight $L^{8a}$, from two slots. These
parts rest on the same age-weighted bound, supplied by the second step of the additive theorem.
\end{enumerate}
The slot size $\mathfrak{s}^{(a)}$ is of order $L^{4a}$ on the event that the age is $a$, and the
slot is unsigned.
\end{remark}

\begin{proof}
Inequality \eqref{8.12:eq-rarity-weights-1} is the product of the bounds
$|X_i|\le C_X\mathfrak{n}(f_i)/\vartheta$, $i=1,2$, which hold everywhere and are stated in
Notation~\ref{8.12:not-two-point-decomposition}. Items (2)--(5) list, for each member of
\eqref{8.12:eq-two-point-decomposition-a} and for the two bad-event probabilities, which rarity
bound its estimate uses; no inequality is asserted there beyond \eqref{8.12:eq-rarity-weights-1}.
Inequality \eqref{8.12:eq-rarity-weights-2} is the rider to the source lemma, in the form in which
it enters the additive theorem.

We derive \eqref{8.12:eq-rarity-weights-3} from \eqref{8.12:eq-rarity-weights-2}. Write
$A_i=A_i(h)$. We have $A_1+A_2\le 2A(h)$ and $L>1$. Hence
\[
L^{4(A_1-\mathfrak{l})}\,L^{4(A_2-\mathfrak{l})}
=L^{4(A_1+A_2-2\mathfrak{l})}
\le L^{8(A(h)-\mathfrak{l})} .
\]
Multiplying by $\mathfrak{s}_1^{\sharp}\mathfrak{s}_2^{\sharp}$ gives
\eqref{8.12:eq-rarity-weights-3}.

We now explain the slot size. Let $\mathsf{U}$ be the multi-scale minimiser configuration, and
write $a_p(\mathsf{U})$ for its plaquette term at $p$ in the energy part of the slot. Inside a live
component $Z$, the minimiser is regular only at the birth scale $j(Z)-1$, by part~(iii) of the
structure statement for the multi-scale minimiser. The regularity lemma then gives, for every
witness cell $\Delta$ inside $Z$,
\begin{equation}\label{8.12:eq-rarity-weights-4}
\sum_{p\subset\Delta}a_p(\mathsf{U})\le C_{\blacksquare}\,\Theta_Z,
\qquad
\Theta_Z=L^{4(k_0-j(Z)+1)}.
\end{equation}
Outside live regions the same sum is bounded by $C_{\square}\varepsilon^2$, with $\varepsilon$ the
small-field parameter.

Consider next the collar classes and the classes attached to live $\mathbf{B}$-operations of the
source lemma, part~(b). Their sum is of the same form, and we substitute $n=k_0-j'$:
\[
\sum_{j'=j(Z)-1}^{k_0}L^{4(k_0-j')}
=\sum_{n=0}^{k_0-j(Z)+1}L^{4n}
\le\frac{\Theta_Z}{1-L^{-4}}
\le 2\,\Theta_Z .
\]
The last inequality holds because $L^{-4}\le\tfrac12$.

With $a=k_0-j(Z)$ we have $\Theta_Z=L^{4(a+1)}$. Hence $\mathfrak{s}^{(a)}$ is comparable to
$L^{4a}$, and this is the size of one slot on the event that the age is $a$.

The slot is unsigned. The energy part, the sum of $f\,a_p(\mathsf{U})$ over plaquettes, is
nonnegative. The terms of the collar and live-$\mathbf{B}$ classes are not. No lower bound on
$\mathbb{E}_h\bigl[\sum a_p\bigr]$ is proved.
\end{proof}

\begin{proposition}[Old live regions enter only through first and second slot moments]
\label{8.12:prop-old-region-moments}
Let $k_0=K-s$, the history law $\tilde{\nu}$ at level $k_0$, the positive normalised plain integrals
$\mathbb{E}_h$ of the live components of a history $h$, the parts $\Phi_i^h$, $G_i$,
$\mathfrak{L}^h_i=\mathfrak{L}^{\mathrm{near},h}_i+\mathfrak{F}^h_i$, $\Psi^h$, the sourced
effective action $\mathfrak{A}^w_h$, the main term $\mathfrak{M}$ and the remainders of the
decomposition \eqref{8.12:eq-two-point-decomposition-a} be as in
Notation~\ref{8.12:not-two-point-decomposition}. We write $\operatorname{Cov}$ and $E$ for
covariance and expectation under $\tilde{\nu}$, and
$\operatorname{Cov}_{\mathbb{E}_h}$ for
\begin{equation*}
\operatorname{Cov}_{\mathbb{E}_h}\bigl(\partial_{z_1}\mathfrak{A}^w_h|_0,
\partial_{z_2}\mathfrak{A}^w_h|_0\bigr).
\end{equation*}
For an integer $a\ge0$ put
\begin{equation*}
\begin{split}
q^{=a}_{k_0}:=\max_{\square}\tilde{\nu}\bigl(&\square \text{ lies in a live component at level }
k_0\\
&\text{whose oldest layer is }k_0-a\bigr),
\end{split}
\end{equation*}
the maximal probability that a cell lies in a live component whose oldest large-field region
was created at step $k_0-a$, that is, has age $a$ at level $k_0$. Let $\mathfrak{l}$ be the age
threshold separating young from old, $q^{\mathrm{lv}}_{k_0}$ the young live weight and
$A_{\mathrm{lv}}$ a constant, and consider the \emph{age series}
\begin{equation}\label{8.12:eq-old-region-moments-a}
\sum_{a>\mathfrak{l}}L^{8(a+1)}q^{=a}_{k_0}\le A_{\mathrm{lv}}L^{8(\mathfrak{l}+1)}q^{\mathrm{lv}}_{k_0}.
\end{equation}
For $i=1,2$ let $B_i$ be the event that some live region of $h$, of any age, lies near the
support of $f_i$ (``near'' in the sense of the source lemma); on $B_i$ let
$\mathsf{A}_i=\mathsf{A}_i(h)$ be the age at level $k_0$ of the oldest such region, and put
$\mathfrak{s}_i^{(a)}:=\mathfrak{s}_i^{\sharp}L^{4(a-\mathfrak{l})}$, where
$\mathfrak{s}_i^{\sharp}$ is the young value of the slot size. Assume:
\begin{itemize}
\item[(i)] the source lemma, the Gaussian extraction, the corollary bounding uniformly the
history-free quantities of the Gaussian extraction, and the corollary on rarity of young
large-field blocks, each in the form in which it enters
Notation~\ref{8.12:not-two-point-decomposition} and
Remark~\ref{8.12:rem-rarity-weights}; in particular $\mathfrak{L}^{\mathrm{near}}_i$ vanishes off
$B_i$, $|\mathfrak{L}^{\mathrm{near}}_i|\le\mathfrak{s}_i^{(\mathsf{A}_i)}$ on $B_i$, and
$\sup|\mathfrak{F}_i|\le\mathfrak{f}_i$;
\item[(ii)] the first- and second-moment bounds and the three covariance inequalities of the
additive theorem for the live slots, where $\mathfrak{q}^{\mathrm{add}}$ is the small constant
of that theorem: whenever \eqref{8.12:eq-old-region-moments-a} holds,
\begin{equation}\label{8.12:eq-old-region-moments-1}
\begin{aligned}
&E|\mathfrak{L}^{\mathrm{near}}_i|\le\mathfrak{s}_i^{\sharp}\mathfrak{q}^{\mathrm{add}}
\quad(i=1,2),\\
&E\bigl[\mathfrak{s}_1^{(\mathsf{A}_1)}\mathfrak{s}_2^{(\mathsf{A}_2)}\mathbf{1}_{B_1\cap B_2}\bigr]
\le\mathfrak{s}_1^{\sharp}\mathfrak{s}_2^{\sharp}\mathfrak{q}^{\mathrm{add}},
\end{aligned}
\end{equation}
and, in all cases,
\begin{equation}\label{8.12:eq-old-region-moments-b}
\begin{aligned}
|\operatorname{Cov}(\mathfrak{L}_1,G_2)|&\le 2\sup|G_2|\cdot E|\mathfrak{L}_1|,\\
|\operatorname{Cov}(\mathfrak{L}_1,\mathfrak{L}_2)|&\le E|\mathfrak{L}_1\mathfrak{L}_2|
+E|\mathfrak{L}_1|\,E|\mathfrak{L}_2|,\\
E\bigl|\operatorname{Cov}_{\mathbb{E}_h}\bigr|&\le 2E\bigl[(\mathfrak{s}_1^{(\mathsf{A}_1)}
\mathbf{1}_{B_1}+\mathfrak{f}_1)(\mathfrak{s}_2^{(\mathsf{A}_2)}\mathbf{1}_{B_2}
+\mathfrak{f}_2)\bigr],
\end{aligned}
\end{equation}
and the first line with the indices $1,2$ exchanged;
\item[(iii)] the split $q^{=a}_{k_0}=q^{\mathrm{heal},=a}+q^{\mathrm{end},=a}$ according to
whether the live component is healed at some later level or remains live to the terminal level
(the nonnegative \emph{healed} and \emph{end} slices), and the healed old-region bound, valid at
every level $k_0$
with no depth floor and with $A^{\mathrm{heal}}_{\mathrm{lv}}$ independent of $K$, $m$ and $N_T$,
\begin{equation}\label{8.12:eq-old-region-moments-2}
\sum_{a>\mathfrak{l}}L^{8(a+1)}q^{\mathrm{heal},=a}\le
A^{\mathrm{heal}}_{\mathrm{lv}}L^{8(\mathfrak{l}+1)}q^{\mathrm{lv}}_{k_0}.
\end{equation}
\end{itemize}
Then:
\begin{itemize}
\item[(a)] \eqref{8.12:eq-old-region-moments-a} implies
\begin{equation}\label{8.12:eq-old-region-moments-3}
\sum_{a>\mathfrak{l}}L^{4(a-\mathfrak{l})}q^{=a}_{k_0}\le A_{\mathrm{lv}}q^{\mathrm{lv}}_{k_0},
\qquad
\sum_{a>\mathfrak{l}}L^{8(a-\mathfrak{l})}q^{=a}_{k_0}\le A_{\mathrm{lv}}q^{\mathrm{lv}}_{k_0};
\end{equation}
\item[(b)] the age series \eqref{8.12:eq-old-region-moments-a}, that is, the numbers
$q^{=a}_{k_0}$ with $a>\mathfrak{l}$ carrying an age-dependent weight, enters the remainder
$S^T_{2;K}-\mathfrak{M}$ only through $R^{\mathfrak{L}}$ (all three of its covariances) and
$R^{\Psi,\mathrm{live}}$, and nowhere else; the other members depend on live regions at most
through age-blind one-cell probabilities; in $R^{\mathfrak{L}}$ and $R^{\Psi,\mathrm{live}}$ the
first-moment members
$\operatorname{Cov}(\mathfrak{L}_1,G_2)$, $\operatorname{Cov}(G_1,\mathfrak{L}_2)$ carry the weight
$L^{4a}$ on a live component of age $a$, and the second-moment members
$\operatorname{Cov}(\mathfrak{L}_1,\mathfrak{L}_2)$, $E\operatorname{Cov}_{\mathbb{E}_h}$ carry the
weight $L^{8a}$;
\item[(c)] \eqref{8.12:eq-old-region-moments-a} holds, with
$A_{\mathrm{lv}}=A^{\mathrm{heal}}_{\mathrm{lv}}+A^{\mathrm{end}}_{\mathrm{lv}}$, as soon as the end
slices satisfy the bound \eqref{8.12:eq-old-region-moments-2} with $q^{\mathrm{end},=a}$ and a
constant $A^{\mathrm{end}}_{\mathrm{lv}}$ in place of $q^{\mathrm{heal},=a}$ and
$A^{\mathrm{heal}}_{\mathrm{lv}}$ (the old-age tail statement); that
statement is not proved in this book, and clause~(iv) of Theorem~\ref{3.1:thm-main} does not rest
on it.
\end{itemize}
\end{proposition}

\begin{proof}
(a) Both inequalities follow from \eqref{8.12:eq-old-region-moments-a} by elementary arithmetic.
Dividing
\eqref{8.12:eq-old-region-moments-a} by $L^{8(\mathfrak{l}+1)}>0$ gives the second inequality of
\eqref{8.12:eq-old-region-moments-3}, since $L^{8(a+1)}/L^{8(\mathfrak{l}+1)}=L^{8(a-\mathfrak{l})}$.
For $a>\mathfrak{l}$ and $L\ge1$ we have $L^{4(a-\mathfrak{l})}\le L^{8(a-\mathfrak{l})}$, and
$q^{=a}_{k_0}\ge0$, being a maximum of probabilities; comparing term by term gives the first.

(b) By \eqref{8.12:eq-two-point-decomposition-a}, $S^T_{2;K}-\mathfrak{M}$ is the sum of the six
remainders of that decomposition. By Remark~\ref{8.12:rem-rarity-weights}, under hypothesis (i):
$\mathfrak{M}$ and $R^{\mathrm{prof}}$ do not depend on the probabilities of live regions;
$R^{\mathrm{loc}}$
consists of deterministic spanning terms, the same for every history; $R^{GX}$ and
$R^{\mathrm{irr}}$ depend on them only through the one-cell probability of the unclean event
$\mathcal{C}^c$, with an integrand bounded on that event independently of any age, by the
uniform bound on the history-free quantities; and the bad-event probabilities
$\mathfrak{q}_W$, $\mathfrak{q}_{\times}$ are one-cell probabilities of lying in a live region
created at any step, which multiply sup-bounded factors. None of these carries a weight
depending on the age of a live component,
so an age-dependent weight on $q^{=a}_{k_0}$, $a>\mathfrak{l}$, can arise only in
$R^{\mathfrak{L}}$ and $R^{\Psi,\mathrm{live}}$.

By hypothesis (i),
\begin{equation*}
|\mathfrak{L}_i|\le|\mathfrak{L}^{\mathrm{near}}_i|+|\mathfrak{F}_i|
\le\mathfrak{s}_i^{(\mathsf{A}_i)}\mathbf{1}_{B_i}+\mathfrak{f}_i.
\end{equation*}
The first line of
\eqref{8.12:eq-old-region-moments-b} bounds $\operatorname{Cov}(\mathfrak{L}_1,G_2)$, and with
indices exchanged $\operatorname{Cov}(G_1,\mathfrak{L}_2)$, by the first moment
$E|\mathfrak{L}_i|\le E[\mathfrak{s}_i^{(\mathsf{A}_i)}\mathbf{1}_{B_i}]+\mathfrak{f}_i$. Here
\begin{equation*}
E\bigl[\mathfrak{s}_i^{(\mathsf{A}_i)}\mathbf{1}_{B_i}\bigr]
=\sum_{a}\mathfrak{s}_i^{\sharp}L^{4(a-\mathfrak{l})}\,\tilde{\nu}(B_i,\ \mathsf{A}_i=a),
\end{equation*}
so the event that the oldest near live region has age $a$ enters with the weight
$L^{4(a-\mathfrak{l})}=L^{-4\mathfrak{l}}L^{4a}$; under \eqref{8.12:eq-old-region-moments-a} the
near part is bounded by the first inequality of \eqref{8.12:eq-old-region-moments-1}.

The second line of \eqref{8.12:eq-old-region-moments-b} bounds
$\operatorname{Cov}(\mathfrak{L}_1,\mathfrak{L}_2)$ by
$E|\mathfrak{L}_1\mathfrak{L}_2|+E|\mathfrak{L}_1|E|\mathfrak{L}_2|$, and the third line bounds
$R^{\Psi,\mathrm{live}}=E\operatorname{Cov}_{\mathbb{E}_h}$; with the bound on $|\mathfrak{L}_i|$
above, both are controlled by
\begin{equation*}
E\bigl[(\mathfrak{s}_1^{(\mathsf{A}_1)}\mathbf{1}_{B_1}+\mathfrak{f}_1)
(\mathfrak{s}_2^{(\mathsf{A}_2)}\mathbf{1}_{B_2}+\mathfrak{f}_2)\bigr]
\end{equation*}
and by products of first moments. Expanding the product, the
terms containing
$\mathfrak{f}_i$ are first moments or age-free, while the term
$E[\mathfrak{s}_1^{(\mathsf{A}_1)}\mathfrak{s}_2^{(\mathsf{A}_2)}\mathbf{1}_{B_1\cap B_2}]$ is the
product of the
two slots on one history. On $B_1\cap B_2$ put $\mathsf{A}(h):=\max\{\mathsf{A}_1,\mathsf{A}_2\}$;
since
$4(\mathsf{A}_1-\mathfrak{l})+4(\mathsf{A}_2-\mathfrak{l})\le8(\mathsf{A}(h)-\mathfrak{l})$ and
$L\ge1$,
\begin{equation*}
\mathfrak{s}_1^{(\mathsf{A}_1)}\mathfrak{s}_2^{(\mathsf{A}_2)}\le\mathfrak{s}_1^{\sharp}
\mathfrak{s}_2^{\sharp}L^{8(\mathsf{A}(h)-\mathfrak{l})},
\end{equation*}
so age $a$ enters with the weight $L^{8(a-\mathfrak{l})}=L^{-8\mathfrak{l}}L^{8a}$; under
\eqref{8.12:eq-old-region-moments-a} this second moment is bounded by the second inequality of
\eqref{8.12:eq-old-region-moments-1}. The two series of \eqref{8.12:eq-old-region-moments-3}
carry exactly these two weights.

(c) All terms are nonnegative by hypothesis (iii), so the series of
\eqref{8.12:eq-old-region-moments-a} is the sum of the healed and end series. The healed series
is bounded by \eqref{8.12:eq-old-region-moments-2}; adding the bound of the same shape on the end
series, with constant $A^{\mathrm{end}}_{\mathrm{lv}}$, assumed in (c)
gives \eqref{8.12:eq-old-region-moments-a} with
$A_{\mathrm{lv}}=A^{\mathrm{heal}}_{\mathrm{lv}}+A^{\mathrm{end}}_{\mathrm{lv}}$. The bound on the
end slices enters only as the hypothesis of (c); it is not proved here.
\end{proof}

\begin{remark}[Old live regions affect the upper and lower bounds alike]
\label{8.12:rem-two-sided-dependence}
We use the decomposition \eqref{8.12:eq-two-point-decomposition-a} of the two-point function, in
the notation of Notation~\ref{8.12:not-two-point-decomposition}: $\tilde{\nu}$ is the history law at the
observation level $k_0=K-s$, $\mathbb{E}_{\tilde{\nu}}$ the expectation under $\tilde{\nu}$,
$\mathbb{E}_h$ the positive normalised plain integral of the live
components of the history $h$, $\mathfrak{L}_i=\mathfrak{L}^h_i$ the live part of $\Phi_i^h$,
$G_i$ its history-free part, and $\mathfrak{A}=\mathfrak{A}^w_h$ the sourced effective action,
with $\partial_i\mathfrak{A}$ its derivative with respect to the coefficient $z_i$ of the
complex source $w=z_1f_1+z_2f_2$. By Proposition~\ref{8.12:prop-old-region-moments}, old live
regions enter the remainders of \eqref{8.12:eq-two-point-decomposition-a} only through the four
quantities
\begin{equation}\label{8.12:eq-two-sided-dependence-1}
\begin{gathered}
  \operatorname{Cov}_{\tilde{\nu}}(\mathfrak{L}_1,G_2),\qquad
  \operatorname{Cov}_{\tilde{\nu}}(G_1,\mathfrak{L}_2),\qquad
  \operatorname{Cov}_{\tilde{\nu}}(\mathfrak{L}_1,\mathfrak{L}_2),\\
  \mathbb{E}_{\tilde{\nu}}\bigl[\operatorname{Cov}_{\mathbb{E}_h}(\partial_1\mathfrak{A},
  \partial_2\mathfrak{A})\bigr],
\end{gathered}
\end{equation}
the first three making up $R^{\mathfrak{L}}$ and the fourth entering $R^{\Psi,\mathrm{live}}$. The
second arises from the first by exchanging the indices $1$ and $2$, so it suffices to discuss the
first. We say that the live slot near the support of $f_i$ is \emph{old} if the live component
containing it was created more than $\mathfrak{l}$ steps before the level $k_0$.
\begin{itemize}
\item[(a)] Each quantity in \eqref{8.12:eq-two-sided-dependence-1} is controlled only in
absolute value, and enters
\begin{equation*}
  S^T_{2;K}-\mathfrak{M}
  =R^{\mathrm{prof}}+R^{GX}+R^{\mathrm{irr}}+R^{\mathrm{loc}}+R^{\mathfrak{L}}
  +R^{\Psi,\mathrm{live}}
\end{equation*}
with no sign. Hence the upper bound
\begin{equation*}
  S^T_{2;K}\le \tfrac{8}{3}\,b_0g^4\mathcal{K}_s(f_1,f_2)
\end{equation*}
requires control of $R^{\mathfrak{L}}+R^{\Psi,\mathrm{live}}$ from above exactly as the lower
bound
\begin{equation*}
  \tfrac{2}{9}\,b_0g^4\mathcal{K}_s(f_1,f_2)\le S^T_{2;K}
\end{equation*}
requires control of it from below, and neither
control is obtained without $\mathbb{E}_{\tilde{\nu}}$-moments of the slots on the event that
they are old.
\item[(b)] No one-sided argument replaces these moments: no sign argument bounds
$\operatorname{Cov}_{\tilde{\nu}}(\mathfrak{L}_1,G_2)$ from above without the first absolute
moment $\mathbb{E}_{\tilde{\nu}}|\mathfrak{L}_1|$, and none bounds
$\operatorname{Cov}_{\tilde{\nu}}(\mathfrak{L}_1,\mathfrak{L}_2)$ from either side with first
moments alone.
\item[(c)] The members of the remainder that depend on rarity only through the one-point rarity
of large-field blocks, level by level (Remark~\ref{8.12:rem-rarity-weights}), do not depend on
the age of the live components. These
are the weight of the unclean event $\mathcal{C}^c$ of the Gaussian extraction, the
probabilities $\mathfrak{q}_W$ and $\mathfrak{q}_{\times}$, and the first and second moments
of the young live member, that is, of the slots of age at most $\mathfrak{l}$, on which
$\mathfrak{s}_i^{(a)}=\mathfrak{s}_i^{\sharp}$. They are proved as an implication from rarity
of young large-field blocks, through its corollary, at every admissible depth $s$ on the tori of
hypothesis \textup{(H1)} of Notation~\ref{2.2:not-hypotheses}, with no depth floor depending
on $g$. The contribution of old live components that are healed at a later level is controlled
by the healed old-region bound. The contribution of the
never-healed old-age tail is the subject of the old-age tail statement, which is not proved;
clause~\textup{(iv)} of Theorem~\ref{3.1:thm-main} does not rest on it.
\end{itemize}
\end{remark}

\begin{proof}
(a) The four quantities in \eqref{8.12:eq-two-sided-dependence-1} are estimated by the
inequalities of \eqref{8.12:eq-old-region-moments-b}, and each of those inequalities bounds
the absolute value of the quantity; none of them asserts a sign. The quantities themselves
carry no sign either. The first is the covariance of the unsigned functional $\mathfrak{L}_1$
with the bounded functional $G_2$, and the second is the same with the indices $1$ and $2$
exchanged. The third is the covariance of the two unsigned slot
functionals $\mathfrak{L}_1$ and $\mathfrak{L}_2$. The fourth is the $\tilde{\nu}$-mean of the
covariance under $\mathbb{E}_h$ of the two unsigned slot functionals $\partial_1\mathfrak{A}$
and $\partial_2\mathfrak{A}$. Since $R^{\mathfrak{L}}$ and $R^{\Psi,\mathrm{live}}$ are sums of
such quantities, they enter the identity for $S^T_{2;K}-\mathfrak{M}$ with no sign.
Consequently an upper bound on $S^T_{2;K}$ requires an upper bound on
$R^{\mathfrak{L}}+R^{\Psi,\mathrm{live}}$, and a lower bound on $S^T_{2;K}$ requires a lower
bound on it. By Proposition~\ref{8.12:prop-old-region-moments}, the only available control of
these quantities on the old event passes through the $\mathbb{E}_{\tilde{\nu}}$-moments of the
slots there. The first-moment members carry weight $L^{4a}$ on a slot of age $a$, and the
second-moment members carry weight $L^{8a}$.

(b) We take the three cases in turn. First, write
\begin{equation*}
  \operatorname{Cov}_{\tilde{\nu}}(\mathfrak{L}_1,G_2)
  =\mathbb{E}_{\tilde{\nu}}\bigl[\mathfrak{L}_1(G_2-\mathbb{E}_{\tilde{\nu}}G_2)\bigr].
\end{equation*}
To bound this from above without
the first absolute moment $\mathbb{E}_{\tilde{\nu}}|\mathfrak{L}_1|$, one would need both
$\mathfrak{L}_1$ and $G_2-\mathbb{E}_{\tilde{\nu}}G_2$ to carry a definite sign. Neither does:
$\mathfrak{L}_1$ is unsigned, and $G_2-\mathbb{E}_{\tilde{\nu}}G_2$ is a centred functional.

Second, consider the upper bound for the second covariance, which reads
\begin{equation*}
  \operatorname{Cov}_{\tilde{\nu}}(\mathfrak{L}_1,\mathfrak{L}_2)
  =\mathbb{E}_{\tilde{\nu}}[\mathfrak{L}_1\mathfrak{L}_2]
  -\mathbb{E}_{\tilde{\nu}}\mathfrak{L}_1\,\mathbb{E}_{\tilde{\nu}}\mathfrak{L}_2 .
\end{equation*}
The inequality
\begin{equation*}
  \operatorname{Cov}_{\tilde{\nu}}(\mathfrak{L}_1,\mathfrak{L}_2)
  \le \mathbb{E}_{\tilde{\nu}}[\mathfrak{L}_1\mathfrak{L}_2]
\end{equation*}
drops the term
$-\mathbb{E}_{\tilde{\nu}}\mathfrak{L}_1\,\mathbb{E}_{\tilde{\nu}}\mathfrak{L}_2$ only if
$\mathfrak{L}_i\ge0$. In any case $\mathbb{E}_{\tilde{\nu}}[\mathfrak{L}_1\mathfrak{L}_2]$ is
itself a second moment. An upper bound with first moments only would need the sign condition
$\mathfrak{L}_i\le0$, and this condition fails, because the energy excess contained in
$\mathfrak{L}_i$ is non-negative.

Third, consider the lower bound for the same covariance. The inequality
\begin{equation*}
  \operatorname{Cov}_{\tilde{\nu}}(\mathfrak{L}_1,\mathfrak{L}_2)
  \ge -\mathbb{E}_{\tilde{\nu}}\mathfrak{L}_1\,\mathbb{E}_{\tilde{\nu}}\mathfrak{L}_2,
\end{equation*}
which uses first moments only, would need $\mathfrak{L}_i\ge0$. This positivity is spoiled by
the collar terms of the live slot, which enter at the same order $L^{4a}$.

The fourth quantity in \eqref{8.12:eq-two-sided-dependence-1} is again a covariance of two
unsigned slot functionals. It is controlled by the third inequality of
\eqref{8.12:eq-old-region-moments-b}, whose right-hand side is a second moment of the near slot
sizes and the far live bounds $\mathfrak{f}_i$.

(c) The members in question depend on rarity only through the one-point rarity of large-field
blocks, level by level, as stated in Remark~\ref{8.12:rem-rarity-weights}. On the
event that the near live slot has age at most $\mathfrak{l}$, the slot size satisfies
$\mathfrak{s}_i^{(a)}=\mathfrak{s}_i^{\sharp}$, so no age weight appears. Their bounds are
therefore those of the corollary on rarity of young large-field blocks, which holds at every
admissible depth $s$ on the tori of hypothesis \textup{(H1)}, with no depth floor depending on
$g$; this proves them as an implication from rarity of young large-field blocks, through that
corollary. By
Proposition~\ref{8.12:prop-old-region-moments}, the age tail enters the remainder only at
$R^{\mathfrak{L}}$ and $R^{\Psi,\mathrm{live}}$. There the old live components that are healed
at a later level are controlled by the healed old-region bound, as stated in
Proposition~\ref{8.12:prop-old-region-moments}; what remains is the never-healed old-age tail.
Its contribution is the subject of the old-age tail statement, which is not proved;
clause~\textup{(iv)} of Theorem~\ref{3.1:thm-main} does not rest on it.
\end{proof}

\begin{remark}[A crude bound on old live regions is the rarity statement itself]
\label{8.12:rem-crude-bound}
Let $k_0=K-s$ be the observation level, $\tilde{\nu}$ the history law at level $k_0$ and
$E_{\tilde{\nu}}$ the expectation with respect to it. Let $f_1,f_2$ be the two test functions of
the two-point witness theorem, with supports $S_1,S_2$, and let $\mathfrak{L}_1,\mathfrak{L}_2$
be the live parts of the corresponding conditional means. By
Proposition~\ref{8.12:prop-old-region-moments}, old live regions enter only through the first
and second slot moments $E_{\tilde{\nu}}|\mathfrak{L}_1|$ and
$E_{\tilde{\nu}}|\mathfrak{L}_1\mathfrak{L}_2|$. A \emph{crude bound} on the share of these moments
carried by old live regions bounds the contribution of the event that a live component of age
$a>\mathfrak{l}$ lies near $S_i$ (for the second moment: near both $S_1$ and $S_2$) by the
supremum of the integrand on that event times the $\tilde{\nu}$-probability of that event.

For a crude bound to be of use for a continuum statement, the $\tilde{\nu}$-probability that an old
live component of age $a$ lies near $S_i$ (near both $S_1$ and $S_2$) must decay like $L^{-4a}$
(like $L^{-8a}$) times $e^{-c\,p_0(g_{k_0})}$ for some constant $c>0$, where $p_0(\cdot)$ is
Ba{\l}aban's degree function of the coupling. This requirement is the
inequality at level $k_0$ on the age series of live components through which the age tail enters
Proposition~\ref{8.12:prop-old-region-moments}, that is, the rarity of old live regions in the
normalised measure $\tilde{\nu}$, whose part concerning the regions whose lineage is still live at
the terminal level of the prolonged run (the never-healed regions) is the old-age tail statement,
which is not proved in
this book. A crude bound therefore does not avoid the never-healed old-age tail; it is that tail.
\end{remark}

\begin{proof}
We write $a_p(U)=1-\operatorname{Re}\operatorname{tr}U(\partial p)$ for the contribution of the
plaquette $p$ to the action. Since the trace is normalised so that the identity matrix has trace
$1$, the normalised trace of an
element of $\mathrm{SU}(N)$ has modulus at most $1$, hence $0\le a_p(U)\le 2$ for every
plaquette.

\emph{The supremum without interior regularity.} If no regularity of the configuration inside the
live component is used, the only available bound on the energy in a witness cell $\Delta$ is the
last clause of the per-cell energy bound for the multi-scale minimiser, which rests on
$a_p\le 2$:
\begin{equation}\label{8.12:eq-crude-bound-1}
\sum_{p\subset\Delta}a_p\le 12\,L^{4k_0}\qquad\text{for every witness cell }\Delta .
\end{equation}
The right-hand side of \eqref{8.12:eq-crude-bound-1} depends on the cutoff: $k_0=K-s$ tends to
infinity as $K\to\infty$ at fixed depth $s$. A supremum of this size gives no bound that survives
the continuum limit.

\emph{The supremum with birth-scale regularity.} Inside a live component $Z$ the multi-scale
minimiser is regular at the birth scale of $Z$; this is the regularity that the construction of
the witness does provide (clause (iii) of the statement on the form of the minimiser inside live
components). With it, the supremum of one live source slot on the event that the oldest near live
component has age $a$ is, by the bound for live source slots that accompanies the last clause of
the per-cell energy bound,
\begin{equation}\label{8.12:eq-crude-bound-2}
C_{\mathrm{live}}\,L^{4(a+1)}\,P_{\mathrm{src}}\,\varsigma^4\,
\frac{\mathfrak{n}(f_i)}{\vartheta}\qquad(i=1,2),
\end{equation}
where $C_{\mathrm{live}}$ is the constant of that bound for one live source slot,
$P_{\mathrm{src}}$ its source prefactor and $\varsigma$ a parameter of that bound, entering
to the fourth power; for the product of two slots the supremum is the product of the two bounds
\eqref{8.12:eq-crude-bound-2}, namely
\begin{equation*}
C_{\mathrm{live}}^2\,L^{8(a+1)}\,P_{\mathrm{src}}^2\,\varsigma^8\,
\frac{\mathfrak{n}(f_1)\,\mathfrak{n}(f_2)}{\vartheta^2}.
\end{equation*}

\emph{The mass.} Multiplying the supremum by the $\tilde{\nu}$-probability of the event, the
old-event share of
$E_{\tilde{\nu}}|\mathfrak{L}_i|$ is controlled, summed over ages $a>\mathfrak{l}$, only if
the $\tilde{\nu}$-probability of the event that an old live component of age $a$ lies near $S_i$
decays like $L^{-4a}$ times
$e^{-c\,p_0(g_{k_0})}$, compensating the factor $L^{4(a+1)}$ of \eqref{8.12:eq-crude-bound-2};
the old-event share of $E_{\tilde{\nu}}|\mathfrak{L}_1\mathfrak{L}_2|$ is controlled only if
the $\tilde{\nu}$-probability of the event that an old live component of age $a$ lies near both
$S_1$ and $S_2$ decays like
$L^{-8a}$ times $e^{-c\,p_0(g_{k_0})}$, compensating $L^{8(a+1)}$. With
$q^{=a}_{k_0}$ the largest $\tilde{\nu}$-probability, over cells $\square$, that $\square$ lies in a
live component at level $k_0$ whose oldest layer is $k_0-a$, and $q^{\mathrm{lv}}_{k_0}$ the
young-live weight, this requirement is the age-series inequality
\begin{equation*}
\sum_{a>\mathfrak{l}}L^{8(a+1)}\,q^{=a}_{k_0}
\le A_{\mathrm{lv}}\,L^{8(\mathfrak{l}+1)}\,q^{\mathrm{lv}}_{k_0},
\end{equation*}
in which $A_{\mathrm{lv}}$ is the constant of the inequality; it is the inequality through which
the age tail enters Proposition~\ref{8.12:prop-old-region-moments}. It is a
statement of rarity in the normalised measure $\tilde{\nu}$, and for the regions whose lineage is
still live at the terminal level of the prolonged run (the never-healed regions) it is the
old-age tail statement,
which is not proved in this book and on which clause (iv) of
Theorem~\ref{3.1:thm-main} does not rest. Thus the crude bound requires exactly the never-healed old-age tail.

\emph{Per-cube small factors do not replace it.} The small factors carried per created cube in
Ba{\l}aban's large-field construction, \cite[(1.79),(1.85)]{B-CMP122b}, namely
$e^{-\kappa_j(Z)-2p_0(g_{j(Z)})}$ and $\gamma_0A_1\,p_0(g_{j(Z)})^2(d'+1)$ (here $j(Z)$ is the
step at which the live component $Z$ is created, $\kappa_j(Z)$ is the exponent of
\cite[(1.79)]{B-CMP122b} for $Z$, and $\gamma_0$ is the constant and $d'$ the
integer of \cite[(1.85)]{B-CMP122b}), together with the weights
$w_{\theta}$, $w^{\mathrm{lin}}_{\theta'}$ and $\varepsilon^{+}(D)$ of the history expansion
($\varepsilon^{+}(D)$ being the small factor attached to a created cube $D$ of the history
expansion), are bounds on un-normalised branch integrals, that is, on the Haar integral of one
term (a branch) of the history expansion, not divided by the partition function:
$\mathbf{T}(X)1$ is at most the supremum of the integrand times
the Haar mass, which is $1$. They become weights of the normalised history law $\tilde{\nu}$ only
through a lower bound on the mass of the alternative in which the same region is absent, taken at
the true scale of the region and with the action matched to the infimum on the side where the
region is present, up to the gain those factors provide. Such a lower bound is, by definition, the
removed-side bound of the annealed live-tail statement, equivalently the half-skin bound at the
canonical comparison point; the annealed live-tail statement is not proved in this book. The
matching-free lower bound, namely the Haar measure of the unrestricted set of configurations of
the region times $e^{-x\sup\mathbf{A}}$, with $x$ the inverse coupling at the scale of the region,
loses at least $\tfrac12 c_*\,p_0(g_j)^2$ per seed created at step $j$, $c_*>0$ being the loss
constant of the proposition on the matching-free lower bound. The loss exceeds the whole gain:
the total gain in the exponent supplied by the small factors listed above is at most
$I\,p_0(g_{j'})^2$ per cube created at step $j'$, for a constant $I$ independent of the cube.
Hence the product of the matching-free lower bound with the
weight has the wrong sign in the exponent, by a margin proportional to the volume, uniformly in the
depth and in the cutoff. By the corollary of that proposition, with the minimal-exceedance lower
bound the required inequality is false with any weight at most $1$; and by the proposition on
keeping the region, if the region is not removed the cell series controlled by the summability
lemma for the history expansion diverges. Consequently no crude bound per created cube converts
into a bound for $\tilde{\nu}$.

\emph{Other upper-bound devices.} None of the following gives the two-point statement of clause
(iv) of Theorem~\ref{3.1:thm-main}. For test functions $f_1,f_2$ write $S^T_{2;K,m}(f_1,f_2)$ for
the connected two-point function of $\mathcal{O}(f_1)$ and $\mathcal{O}(f_2)$ (the case $n=2$ of
$S^T_{n;K,m}$).
\begin{itemize}
\item The Cauchy--Schwarz inequality
\begin{equation*}
|\operatorname{Cov}(\mathcal{O}(f_1),\mathcal{O}(f_2))|
\le(\operatorname{Var}\mathcal{O}(f_1)\,\operatorname{Var}\mathcal{O}(f_2))^{1/2}
\end{equation*}
involves the variances at coincident points, which diverge with the cutoff by the proposition on
coincident-point variances.
\item The Cauchy--Schwarz inequality furnished by reflection positivity of the lattice measure,
\begin{equation*}
|S^T_{2;K,m}(f_1,\iota f_2)|^2
\le S^T_{2;K,m}(f_1,\iota f_1)\,S^T_{2;K,m}(f_2,\iota f_2),
\end{equation*}
where $\iota$ denotes a reflection of Lemma~\ref{1.1:lem-invariance} under which the lattice
measure is reflection positive, only moves the problem to reflected pairs at the same scale.
\item The torus bound of the correlation theorem, taken at $m:=s$, and the per-ball volume-free
bound are uniform in the cutoff but do not see the factor $g^4$: in the torus bound the whole
large-field and history content enters through the global two-sided stability quotient, which is
uniform in $K$, depends on the volume and does not see $g^4$; in the per-ball volume-free bound the
constant $C^\infty_n(\mathsf{d};R)$ of clause (iii-b) of Theorem~\ref{3.1:thm-main} is only
non-increasing in $\mathsf{d}$, no decay in $\mathsf{d}$ being asserted. At best these statements
bound $S^T_{2;K,m}(f_1^{(\mathsf{d})},f_2^{(\mathsf{d})})$, for test
functions $f_1^{(\mathsf{d})},f_2^{(\mathsf{d})}$ supported in balls of radius $\mathsf{d}$ at
mutual distance of order $\mathsf{d}$, as in Definition~\ref{3.1:def-non-trivial}, by a constant
independent
of $\mathsf{d}$; that is the behaviour of the free field, not a decay to zero.
\item Jensen's inequality, or the convexity of the logarithm of the partition function (the
convexity proposition for local free energies), would have to give two-sided bounds on local free
energies to a precision finer than their leading term, which is the two-point witness theorem
itself.
\end{itemize}
\end{proof}

\begin{remark}[The global stability quotient and its volume dependence]\label{8.12:rem-global-quotient}
In this remark $N_T$ denotes the number of sites per side of the terminal lattice, the lattice on which
the run stops, so that $N_T^4$ is the number of its cells. The letter $N_T$ is a volume parameter and
not the rank of the gauge group. Here $\tilde{\nu}$ is the history law, and $\Theta$ is the total
reserve of a terminal internal history. We write $\rho_k$ for the effective density of
Definition~\ref{1.5:def-densities} at the terminal level $k$, $Z_{\mathrm{norm}}$ for the normalising
integral of $\tilde{\nu}$, and $C_{\mathrm{low}}$ for the constant of the lower bound on
$Z_{\mathrm{norm}}$ supplied with the correlation theorem. We call \emph{the global device} the
estimation of a probability under $\tilde{\nu}$ by the global quotient of (a) times the supremum of
$e^{-\Theta}$ on the event, as in (b). The following hold.
\begin{itemize}
\item[(a)] The one normalisation of $\tilde{\nu}$ available without further input is global. The
unnormalised total mass and the normalising integral $Z_{\mathrm{norm}}$ satisfy
\begin{equation}\label{8.12:eq-global-quotient-1}
\int\rho_k\le e^{E_+N_T^4},\qquad Z_{\mathrm{norm}}\ge e^{-C_{\mathrm{low}}N_T^4},
\end{equation}
and their quotient is bounded by $e^{\mathsf{B}N_T^4}$, where $\mathsf{B}:=E_++C_{\mathrm{low}}$.
\item[(b)] For every event $E$ of terminal-live structure,
\begin{equation}\label{8.12:eq-global-quotient-2}
\tilde{\nu}(E)\le e^{\mathsf{B}N_T^4}\,\sup_E e^{-\Theta}.
\end{equation}
\item[(c)] For each fixed terminal volume there is an age threshold $a_{N_T}$, proportional to
$\mathsf{B}N_T^4/\varrho$, beyond which terminal ages of live components are exponentially rare in
$\tilde{\nu}$, globally on the terminal lattice.
\item[(d)] No bound uniform in $N_T$ on the series $\sum_{a>\mathfrak{l}}L^{8(a-\mathfrak{l})}q^{=a}$
of Proposition~\ref{8.12:prop-old-region-moments} follows from the inputs of the global device.
\end{itemize}
Where the terminal cell count is bounded by a function independent of $g$, this normalisation
suffices.
\end{remark}

\begin{proof}
\emph{The quotient (a).} By the stability bound after integration, clause (i-b) of
Theorem~\ref{3.1:thm-main}, the effective density $\rho_k$ has total mass at most $e^{E_+N_T^4}$.
This is the first inequality in
\eqref{8.12:eq-global-quotient-1}. The normalising integral obeys the lower bound
$Z_{\mathrm{norm}}\ge e^{-C_{\mathrm{low}}N_T^4}$ supplied with the correlation theorem. This is the
second inequality there.

Dividing the first inequality by the second gives
\begin{equation*}
\frac{\int\rho_k}{Z_{\mathrm{norm}}}\le e^{(E_++C_{\mathrm{low}})N_T^4}.
\end{equation*}
Write $\mathsf{B}=E_++C_{\mathrm{low}}$; then the right side is the global stability quotient
$e^{\mathsf{B}N_T^4}$. The dependence of $\mathsf{B}$ on the coupling is that of the degree statement
for $\mathsf{B}$, with log-degree constant $C_{\mathsf{B}}$, taken here as stated together with the
lower bound on $Z_{\mathrm{norm}}$. The exponent is proportional to $N_T^4$, the
number of cells of the lattice on which the run stops.

\emph{The event bound (b).} Let $E$ be an event of terminal-live structure. The reserve-stripped
termwise upper bound bounds the unnormalised mass of the terms contributing to $E$ by
$e^{E_+N_T^4}\sup_E e^{-\Theta}$. Dividing by
$Z_{\mathrm{norm}}$ and using the lower bound in \eqref{8.12:eq-global-quotient-1} gives
\eqref{8.12:eq-global-quotient-2}. The two inequalities of \eqref{8.12:eq-global-quotient-1} and the
bound \eqref{8.12:eq-global-quotient-2} are the global inequalities that enter the lemma of (d) below
as inputs.

\emph{The age threshold (c).} Fix the terminal volume. Let $E_a$ be the event that some live component,
live at the terminal level, has terminal age at least $a$; this is an event of terminal-live
structure. By the linear age--reserve relation of the age lemma, $\Theta\ge\varrho a$
on $E_a$. Hence $\sup_{E_a}e^{-\Theta}\le e^{-\varrho a}$. Inserting this in
\eqref{8.12:eq-global-quotient-2} gives the Chebyshev bound
\begin{equation}\label{8.12:eq-global-quotient-3}
\tilde{\nu}(E_a)\le e^{\mathsf{B}N_T^4-\varrho a}=e^{-\varrho\,(a-\mathsf{B}N_T^4/\varrho)}.
\end{equation}
The right side is less than one once $a$ exceeds $\mathsf{B}N_T^4/\varrho$, a threshold $a_{N_T}$
proportional to $\mathsf{B}N_T^4/\varrho$, and decays exponentially in $a-a_{N_T}$ beyond it. The
threshold depends on the volume and grows like $N_T^4$.

In the form of the corollary bounding the probability of a live component of terminal age at least
$A$, the same mechanism reads
\begin{equation*}
\tilde{\nu}\bigl(\text{some live component of terminal age}\ \ge A\bigr)\le 2q^A e^{\mathsf{B}N_T^4},
\end{equation*}
where, in this display only, $q^{A}$ is the $A$-th power of a number $q$, $0<q<1$, fixed in that
corollary and unrelated to the probabilities $q^{=a}$ below.
The factor $e^{\mathsf{B}N_T^4}$ is again the global quotient. The bound is informative only for $A$
beyond a threshold growing linearly in $N_T^4$. This proves (c) for the old slices of the live event
$\mathcal{E}$, that is, for live components whose lineage is still live at the terminal level.

\emph{No uniformity in the volume (d).} The lemma exhibiting a model for the global device settles this
point. Consider a translation-invariant model with exactly one old cell of the live event, of age
$a_0=\lfloor\mathsf{B}N_T^4/\varrho\rfloor$, which exceeds $\mathfrak{l}$ for $N_T$ large. That lemma
shows that the model satisfies every input of the global
device:
\begin{itemize}
\item positivity;
\item the global inequalities \eqref{8.12:eq-global-quotient-1} and \eqref{8.12:eq-global-quotient-2};
\item linear cell counts;
\item the linear age--reserve relation of the age lemma;
\item the live-probability bound in the form used by the global device;
\item the further proposition of this chapter taken as an input of the global device;
\item every pair or second-moment identity derivable from these.
\end{itemize}
In the model,
\begin{equation}\label{8.12:eq-global-quotient-4}
L^{8(a_0-\mathfrak{l})}\,q^{\mathcal{E},=a_0}
=\frac{L^{8(a_0-\mathfrak{l})}}{n}\longrightarrow\infty
\qquad (N_T\to\infty),
\end{equation}
where $q^{\mathcal{E},=a}$ denotes the maximal probability that a cell lies in a live component of
exact age $a$ whose lineage is still live at the terminal level (the superscript $\mathcal{E}$ marks
this live event), and, in this proof
only, $n$ is the parameter of the model fixed in that lemma. The left side of
\eqref{8.12:eq-global-quotient-4} is the term of index $a_0$ of the series
\begin{equation*}
\sum_{a>\mathfrak{l}}L^{8(a-\mathfrak{l})}q^{\mathcal{E},=a},
\end{equation*}
the part of the age series for components whose lineage is live at the terminal level. The event
defining $q^{\mathcal{E},=a}$ is one of the two parts into which the event defining $q^{=a}$ splits,
so $q^{\mathcal{E},=a}\le q^{=a}$. Hence the series
$\sum_{a>\mathfrak{l}}L^{8(a-\mathfrak{l})}q^{=a}$ of
Proposition~\ref{8.12:prop-old-region-moments} dominates this term and is therefore unbounded in the
model as $N_T\to\infty$. So no bound on that
series uniform in $N_T$ can be derived from the inputs listed above. This proves (d).

\emph{Bounded terminal cell count.} If the terminal cell count $N_T^4$ is bounded by a
function independent of $g$, then the exponent $\mathsf{B}N_T^4$ is at most that function times
$\mathsf{B}$, and its dependence on the coupling is that of $\mathsf{B}$, governed by the degree
statement for $\mathsf{B}$.
\end{proof}

\begin{remark}[Non-negativity from reflection positivity and strict positivity at one scale]
\label{8.12:rem-positivity-one-scale}
Let a cutoff $K$, a torus $\mathbb{T}_m$ and a scale $\mathsf{d}_0$ be given, with depth
$s_0=s(\mathsf{d}_0)$ and observation level $k_0=K-s_0$. We write $\tilde{\nu}$ for the history
law at level $k_0$, and $\mathfrak{l}=\mathfrak{l}(x_{k_0})$ for the age threshold separating young
from old live components. We write $K'$ for the terminal level of the longest run, and say that a
live component is \emph{live at the end} if its lineage is still live at level $K'$. We write $N_T$
for the number of sites per side of the terminal lattice of this run, the lattice on which it stops;
this lattice has $N_T^4$ cells, and the global stability quotient of
Remark~\ref{8.12:rem-global-quotient} is $e^{\mathsf{B}N_T^4}$. We write $\mathsf{d}_W(g)$ for the
witness scale at coupling $g$ of the re-targeting corollary of this chapter, the scales
$\mathsf{d}\le\mathsf{d}_W(g)$ being those treated there, and $\mathsf{d}_W(g,m)$ for the
corresponding volume-dependent scale on the fixed torus $\mathbb{T}_m$. For a link-plane reflection
$\theta$ of $\mathbb{T}_m$ (one of the reflections of Lemma~\ref{1.1:lem-invariance}; in this
remark $\theta$ denotes the reflection and not a reserve fraction) and a test function $f$,
$\theta f$ is the reflected test function. The pair $(\theta f,f)$ is a \emph{reflected separated
pair} if $f$ is supported in one of the two open half-tori into which the reflection planes of
$\theta$ divide $\mathbb{T}_m$; in particular the supports of $\theta f$ and $f$ are then separated.
We write $\mathfrak{M}$ for the main term and $\mathfrak{R}$ for the remainder of the two-point
witness theorem at depth $s_0$. The \emph{end member} of $\mathfrak{R}$ is the member of
$\mathfrak{R}$ carried by old live components that are live at the end.
\begin{itemize}
\item[(a)] For every reflected separated pair, at every cutoff, volume, scale and position,
\begin{equation}\label{8.12:eq-positivity-one-scale-1}
S^T_{2;K,m}(\theta f,f)\ge 0 ,
\end{equation}
and the same inequality holds in every limit. It involves no large-field tail and no
volume-dependent constant, and it is the whole content of the non-negativity of the connected
two-point function.
\item[(b)] Strict positivity at the scale $\mathsf{d}_0$ is the inequality
\begin{equation}\label{8.12:eq-positivity-one-scale-2}
\mathfrak{M}-|\mathfrak{R}|>0\qquad\text{at depth } s_0=s(\mathsf{d}_0).
\end{equation}
The end member of $\mathfrak{R}$ at depth $s_0$ is a sum over the ages
\begin{equation}\label{8.12:eq-positivity-one-scale-3}
\mathfrak{l}<a\le k_0-1=K-s_0-1 .
\end{equation}
This range grows with the cutoff. The contribution of old components live at the end is therefore
present at every scale held fixed as $K\to\infty$.
\item[(c)] The age-blind end density bounds the end member by a volume-free quantity, but this
quantity diverges with $K$ at fixed depth. The global exponential moment caps the ages at a
threshold $a_{N_T}$ proportional to $\mathsf{B}N_T^4/\varrho$; this cap does not depend on the
cutoff, but it does depend on the volume.
\item[(d)] Three conditions are equivalent:
\begin{itemize}
\item[--] the end member is dominated at one scale uniformly in $K$;
\item[--] the terminal cell count $N_T^4$ is controlled;
\item[--] either the tori are those of the case of the re-targeting corollary of this chapter in
which $N_T^4$ admits a bound free of $g$, or the torus is fixed and its volume enters the constant.
\end{itemize}
Accordingly the end member is dominated on the tori of that case, indeed at every scale
$\mathsf{d}\le\mathsf{d}_W(g)$. It is also dominated on each fixed torus below a
volume-dependent scale $\mathsf{d}_W(g,m)$. It is not dominated uniformly in the volume.
\item[(e)] No depth window and no choice of test functions keeps old never-healed live components
away from the balls carrying the supports of the test functions by combinatorics, except the
window $k_0\le\mathfrak{l}(x_{k_0})+1$. That window has no continuum content.
\item[(f)] On the terminal side nothing combinatorial keeps the level-$k_0$ ancestor of a
component live at the end away from these balls.
\end{itemize}
\end{remark}

\begin{proof}
\emph{Item} (a). By reflection positivity of the lattice measure, which is the link-plane
reflection positivity of the Wilson action of Osterwalder and Seiler \cite{OSe78},
\eqref{8.12:eq-positivity-one-scale-1} holds for every reflected separated pair.
This property of the lattice measure holds at each fixed $K$ and $m$, and it is obtained
without the renormalisation-group expansion. Consequently \eqref{8.12:eq-positivity-one-scale-1}
holds at every cutoff, volume, scale and position, carries no large-field tail, and has no constant
depending on the volume. A limit of non-negative numbers is non-negative, so the inequality holds
in every limit of the connected two-point functions. Nothing beyond
\eqref{8.12:eq-positivity-one-scale-1} is contained in the statement that the connected two-point
function is non-negative.

\emph{Item} (b). By the two-point witness theorem, $S^T_{2;K,m}$ at depth $s_0$ is the main term
$\mathfrak{M}$ plus the remainder $\mathfrak{R}$. Strict positivity at $\mathsf{d}_0$ is therefore
the inequality \eqref{8.12:eq-positivity-one-scale-2}.

The end member of $\mathfrak{R}$ at depth $s_0$ sums over the old ages
\eqref{8.12:eq-positivity-one-scale-3}. At fixed $s_0$ the upper end $K-s_0-1$ tends to infinity
with $K$. The range is not truncated by any of the mechanisms that control live components, for the
following four reasons.
\begin{itemize}
\item[--] The estimates on a live component's waiting time charge for the wait and bound the size
of the component during it; these are the wait bound entering the correlation theorem, a second
wait bound in flat form, and the arrest estimates. None of them caps the component's age.
\item[--] Quiet live components that were born large, or that wrap around the torus, fail the
condition of the first domain class at every level of the run.
\item[--] Chains of refreshments keep the index of the oldest creation step, so a refreshment does
not lower the age.
\item[--] By the positivity of the history law, every age carries positive $\tilde{\nu}$-weight.
\end{itemize}
Hence the contribution of old components live at the end is present at every scale held fixed as
$K\to\infty$.

\emph{Item} (c). Part (d) of the end-density theorem gives the age-blind end density
\begin{equation*}
\mathfrak{e}_{k_0}=\mathfrak{E}L^{-4(K'-k_0)} ,
\end{equation*}
which contains no volume. Summing it against the weight $L^{8a}$ of age $a$ over the ages
$a\le k_0$ gives a geometric sum, which lies between $L^{8k_0}$ and $L^{8k_0}L^8/(L^8-1)$. With
$k_0=K-s$ this yields
\begin{align*}
\sum_{a\le k_0}L^{8a}\,\mathfrak{e}_{k_0}
&\asymp \mathfrak{E}\,L^{8(K-s)}L^{-4(K'-K+s)}\\
&=\mathfrak{E}\,L^{8K-12s}L^{-4(K'-K)} ,
\end{align*}
and this tends to infinity with $K$ at fixed $s$. The bound is thus volume-free but diverges with
the cutoff.

The one normalisation available without an age-dependent input is the global stability quotient
$e^{\mathsf{B}N_T^4}$ of Remark~\ref{8.12:rem-global-quotient}. The exponential moment it yields
caps the ages at $a_{N_T}$, proportional to $\mathsf{B}N_T^4/\varrho$. This cap is free of the
cutoff and depends on the volume through $N_T^4$.

\emph{Item} (d). By item (b), domination of the end member at one scale uniformly in $K$ requires
control of the full range \eqref{8.12:eq-positivity-one-scale-3} as $K\to\infty$. By item (c),
the volume-free route through $\mathfrak{e}_{k_0}$ diverges with $K$. The remaining cap,
$a_{N_T}$, is uniform in $K$ exactly when $N_T^4$ is controlled.

This control holds in two situations. The first is the case of the re-targeting corollary of this
chapter in which $N_T^4$ admits a bound free of $g$. The second is a fixed torus, where the volume
enters the constant. This gives the equivalences of item (d).

On the tori of the first case the end member is dominated, and this holds at every scale
$\mathsf{d}\le\mathsf{d}_W(g)$. On each fixed torus it is dominated below the volume-dependent
scale $\mathsf{d}_W(g,m)$. Uniformly in the volume it is not dominated.

One cannot trade one uniformity for the other. A bound uniform in $K$ but not in $m$ is the finite-volume
witness of Corollary~\ref{8.12:cor-finite-volume-witness} on a fixed torus, uniform in $K$ with the
volume in the depth floor. A bound uniform in $m$ but not in $K$ holds at each fixed cutoff,
because the bound through $\mathfrak{e}_{k_0}$ in item (c) is volume-free and finite for fixed
$K$. Such a bound carries no continuum content.

\emph{Item} (e). We treat in turn the only combinatorial exclusion, the size of old components at a
cutoff-uniform depth, and the separation geometry.

\emph{The only combinatorial exclusion.} It is $k_0\le\mathfrak{l}(x_{k_0})+1$: the observation
level lies within $\mathfrak{l}+1$ levels of the bare lattice, where no component can be older than
$k_0-1\le\mathfrak{l}$. There the separations satisfy
$\mathsf{d}\le C_W(s)L^{\mathfrak{l}+1}\cdot L^{-K}$, with $s=K-k_0$, where $L^{-K}$ is the bare
spacing. These separations shrink to zero physical size as $K\to\infty$. They also violate the
ultraviolet depth condition $K-s\ge k_W$ as soon as $k_W>\mathfrak{l}+1$. No continuum statement
lives in this window.

\emph{At any cutoff-uniform depth.} The age lemma, whose age bound is
\eqref{8.12:eq-volume-ranges-d}, says that an old component with few seeds was born large: its
size $d'$ at birth, in the sense of the age lemma, and its age $a$ at level $k_0$ satisfy
\begin{equation*}
2\log_L(d'+1)\gtrsim a.
\end{equation*}
A quiet region shrinks by a factor $\sqrt{L}$ per level. Hence at the witness scale an old quiet
component is a speck of linear size at most of order $L^{-a/2}M\check{R}$ cells, and it
fits inside either ball. Chains of refreshments behave in the same way.

The test functions are deterministic and the live regions are random, so no choice of supports
avoids them. Shallow depths, with $\mathsf{d}$ near $\mathsf{d}_W(g)$, do not help, since the ages
still range up to $K-s-1$. Deep depths do not help either, for the same reason.

\emph{Separation geometry.} A single component meeting both balls has diameter at least the
separation of the balls in level-$k_0$ cells. That separation is at least $C_W(s)\gg 100M\check{R}$
cells. Components that were born large, or that arose by merging, are admissible at that size
(they fail the condition of the first domain class). Their un-normalised weight is
$e^{-\kappa d_{k_0}(Z)}$, where $d_{k_0}(Z)$ is at least of order the separation divided by
$M\check{R}$.

Converting even this locality factor to a bound in $\tilde{\nu}$ fails, by the proposition that
gives the matching-free bound: that bound forfeits at least of order $p_0(g)^2$ per seed and gains
$\kappa$ per unit length, while $\kappa<Cp_0(g)^2$ for small $g$, where $p_0(\cdot)$ is the degree
function. So not even the sub-event in which one component spans the gap is excluded, either
combinatorially or trivially.

Moreover, the member of $\mathfrak{R}$ with a live source slot near each of the two balls, the
two-slot member, is carried also by two components, one near each ball. In the proof of
Proposition~\ref{8.12:prop-old-region-moments} the joint bad event of the two slots, $B_1\cap B_2$
in the notation of that proof, is bounded by the probability of one of $B_1$, $B_2$. No joint
rarity of the two events is used there, and none is available.

\emph{Item} (f). A component is live at the end if it is live at the terminal level $K'$ of the
longest run. Nothing combinatorial keeps the level-$k_0$ ancestor of such a component away from the
balls. The terminal torus has $N_T^4$ cells, and the balls are not special positions: the
covariance and expectation bounds used for the terminal side average over positions and do not
exclude any.
\end{proof}

\begin{definition}[Prolonged runs, the two volume ranges and the age bound]\label{8.12:not-volume-ranges}
In this section the letter $N$, with or without decorations ($N$, $N_T$, $N^{\dagger}$,
$N_T^{\dagger}$, $\mathsf{N}$), denotes a number of lattice sites per side of a torus; the integer
$N$ of $\mathrm{SU}(N)$ does not occur in this section.

\emph{Bare data and members.} A \emph{bare datum} is a pair $(x_0,\mathsf{N})$ consisting of a bare
inverse coupling $x_0$ and a number $\mathsf{N}$ of bare sites per side; it fixes the Wilson measure
of Definition~\ref{1.1:def-wilson-action}. A \emph{member} over the bare datum $(x_0,\mathsf{N})$
is a pair $(g,N_T)$ such that, with $K(x_0,g)$ the first-passage index $K(g_0,g)$ of
Definition~\ref{1.2:def-flow-constants} for the bare coupling $g_0$ with $g_0^{-2}=x_0$,
\begin{equation}\label{8.12:eq-volume-ranges-1}
X:=g^{-2}<x_0,\qquad K:=K(x_0,g),\qquad \mathsf{N}=N_TL^{K}.
\end{equation}
The admissible runs of the member stop at a level $K-m_{\mathrm{st}}$ with
$m_{\mathrm{st}}\ge m_{\mathbb{T}}(g)$, on a lattice with $N=N_TL^{m_{\mathrm{st}}}$ sites per side;
we call $m_{\mathrm{st}}$ the \emph{stopping depth} of the run. The data $(g_0,g,m)$ of
Theorem~\ref{3.1:thm-main} are the member
$(g_V,N_T)$ over a bare datum with $x_0=g_0^{-2}$, with
\begin{equation*}
N_T=2L^{\,m-m_{\mathbb{T}}(g_V)},
\end{equation*}
where $g_V$ lies within a fixed margin of $g$, the margin of the comparison between the data of
Theorem~\ref{3.1:thm-main} and members. We put $n_0:=2$ and call
\begin{equation}\label{8.12:eq-volume-ranges-2}
a:=m-m_{\mathbb{T}}(g_V)
\end{equation}
the \emph{torus excess}; thus $N_T=n_0L^{a}$.

\emph{The coupling threshold.} Let $\gamma_*$ be the minimum of $\gamma_J$, of the coupling ceilings
of the theorem that uses the global stability quotient of
Remark~\ref{8.12:rem-global-quotient}, and of a further coupling ceiling, written $\gamma_V$. We put
\begin{equation*}
X_*:=\gamma_*^{-2},\qquad K_*:=K(x_0,\gamma_*),
\end{equation*}
and we have $K_*\ge K$. The quantities $\bar{u}(g)$, $c_{\mathrm{gap}}$ and $\bar{n}$ are those of
the lemma on prolonged runs recalled below.

\emph{Consistency of runs over one bare datum.} We use the statement on the consistency of runs
over one bare datum: no object of the
$k$-th step depends on the target coupling, on $K$ or on the stopping depth $m_{\mathrm{st}}$; two
runs over the same bare datum
coincide through the last step of the shorter one; and for an event $E$ of lower levels the
number $\tilde{\nu}(E)$ is the same for both runs.

\emph{Prolonged runs.} By the lemma on prolonged runs, prolongation on the same lattice to a level
attainable in the sense of that lemma,
\begin{equation*}
k^{\dagger}\le\hat{\kappa}:=\min\{K+a,\;K_*-1\}
\end{equation*}
gives a member $(g^{\dagger},N_T^{\dagger})$ over the same bare datum $(x_0,\mathsf{N})$, with
\begin{equation*}
g\le g^{\dagger}<\gamma_*,\qquad m_{\mathbb{T}}(g^{\dagger})\le m_{\mathbb{T}}(g),
\end{equation*}
with terminal level
\begin{equation*}
K^{\dagger\prime}:=k^{\dagger}-m_{\mathbb{T}}(g^{\dagger})\ge K-m_{\mathbb{T}}(g),
\end{equation*}
and with $N^{\dagger}:=N_T^{\dagger}L^{m_{\mathbb{T}}(g^{\dagger})}$ sites per side at level
$K^{\dagger\prime}$. We write $X^{\dagger}:=(g^{\dagger})^{-2}$. The same lemma gives, with
$\lambda$ and $b_+$ the constants of the flow of the inverse coupling that enter its statement (the
letter $\lambda$ here does not denote a gauge transformation) and $c_2$ the
flow constant of Definition~\ref{1.2:def-flow-constants},
\begin{equation}\label{8.12:eq-volume-ranges-c}
K_*-K\ \ge\ \frac{X-X_*-b_+-2c_2}{3\lambda}.
\end{equation}

\emph{The two volume ranges.} We distinguish
\begin{equation}\label{8.12:eq-volume-ranges-a}
\text{Case V:}\quad K+a\le K_*-1,\qquad\qquad \text{Case C:}\quad K+a\ge K_* .
\end{equation}
Case V is the range in which the prolonged run reaches the volume scale, the level $K+a$, before
the coupling reaches $\gamma_*$; it consists of every torus with $a<K_*-K$. In Case V the level
$k^{\dagger}$ of the lemma on prolonged runs satisfies $\hat{\kappa}-k^{\dagger}\le c_{\mathrm{gap}}$, and
$N_T^{\dagger}\le\bar{n}L^{c_{\mathrm{gap}}}$, whence
\begin{equation}\label{8.12:eq-volume-ranges-3}
(N^{\dagger})^4\le\bigl(\bar{n}L^{c_{\mathrm{gap}}}L^{m_{\mathbb{T}}(g^{\dagger})}\bigr)^4 .
\end{equation}
The right side of \eqref{8.12:eq-volume-ranges-3} is bounded by a function of
$(X^{\dagger},L,M,\bar{n})$ which decreases with the coupling and is polylogarithmic of degree
$4r_{\mathrm{pr}}$ in $\log X^{\dagger}$.

\emph{The exponent of the global quotient.} With a constant $E_+^{\mathrm{lin}}$, independent of
$K$, which plays the role of the stability constant $E_+$ in the termwise bound
\eqref{8.12:eq-volume-ranges-7} below, and $C_{\mathrm{low}}$ the constant of the lower bound on the
partition function, we put
\begin{equation}\label{8.12:eq-volume-ranges-4}
\mathsf{B}(g):=E_+^{\mathrm{lin}}+C_{\mathrm{low}}\ \le\ C_{\mathsf{B}}\bigl(1+\log g^{-2}\bigr).
\end{equation}
The inequality in \eqref{8.12:eq-volume-ranges-4} is the degree bound for $\mathsf{B}$, the content
of the lemma bounding $E_+^{\mathrm{lin}}+C_{\mathrm{low}}$ linearly in $\log g^{-2}$.

\emph{The age bound.} For a large-field region $Z$ we write $j(Z)$ for the step at which it was
created. By the age lemma, for every run, every level $k$, and every live component $C$ at level
$k$ with internal history $\mathfrak{h}$ and seed set $\mathcal{Z}(\mathfrak{h})$ (the set of seeds
of $\mathfrak{h}$, each seed being a large-field region $Z$ created at the step $j(Z)$),
\begin{equation}\label{8.12:eq-volume-ranges-d}
\operatorname{age}_k(C)\ \le\ \sum_{Z\in\mathcal{Z}(\mathfrak{h})}\mathfrak{a}(Z),
\end{equation}
where the per-seed age allowance satisfies
\begin{equation*}
\mathfrak{a}(Z)\le 2\log_L\bigl(d'_{j(Z)}(Z)+1\bigr)+2R_{j(Z)}+n_1+C_S+4
\end{equation*}
for a seed of the first kind, and $\mathfrak{a}(Z)\le 2R_{j(Z)}+n_1+4$ for a preparatory seed;
these are the two kinds of seeds of Ba{\l}aban's rules for large-field regions
\cite{B-CMP122a,B-CMP122b}. Here
$d'_{j}(Z)$ denotes the size of the region $Z$ measured at step $j$, $R_j$ the radius attached
to step $j$ by those rules (not the function $R(\cdot)$ entering $m_{\mathbb{T}}$), and $n_1$,
$C_S$ are constants of the rules by which large-field regions are formed.

\emph{The reserves and their split.} By the lemma on the reserves of histories, for each history
$\mathfrak{h}$ the flat reserve weight
$w_{\theta}$ carries $\theta\,p_0(g_{j(Z)})$ per seed $Z\in\mathcal{Z}(\mathfrak{h})$, and the
linear reserve weight $w^{\mathrm{lin}}_{\theta'}$ carries
$\theta'\gamma_0A_1\,p_0(g_{j(Z)})^2\bigl(d'_{j(Z)}(Z)+1\bigr)$ per seed, $p_0(\cdot)$ being the
degree function and $\gamma_0$ the constant of the linear reserve; the two reserves are available
simultaneously for every history. The total reserve of a history is
\begin{equation}\label{8.12:eq-volume-ranges-5}
\Theta(\mathfrak{h}):=\sum_{Z\in\mathcal{Z}(\mathfrak{h})}
\Bigl[\theta\,p_0(g_{j(Z)})+\theta'\gamma_0A_1\,p_0(g_{j(Z)})^2\bigl(d'(Z)+1\bigr)\Bigr],
\end{equation}
with $d'(Z):=d'_{j(Z)}(Z)$. The linear fraction is split as
\begin{equation}\label{8.12:eq-volume-ranges-6}
\theta'=\theta'_V+\theta'_{\mathrm{age}},
\end{equation}
$\theta'_V$ being the share assigned to the absorption of the volume and
$\theta'_{\mathrm{age}}$ the share assigned to the age. The two shares are fixed positive
fractions of $\theta'$; every comparison made with them in this chapter is a comparison of
polylogarithmic degrees and holds for any fixed positive fractions.

\emph{The termwise reserve-stripped bound.} Consider a member whose longest run ends at level
$K'$ on $N$ sites per side, and let $\mathcal{S}$ be any set of terminal internal histories. For
a history $h$ of the run write $\mathfrak{h}_h$ for its terminal internal history. We use the
termwise reserve-stripped bound contained in the correlation theorem, together with the two
statements by which the reserves are stripped from the termwise bound, valid with the constant
$E_+^{\mathrm{lin}}$ independent of $K$ and
with the Haar measures normalised to total mass one (Definition~\ref{1.1:def-wilson-action}), so
that the Haar factor of the bound equals one,
\begin{equation}\label{8.12:eq-volume-ranges-7}
\sum_{h:\,\mathfrak{h}_h\in\mathcal{S}}\int d\tau_h\ \le\
\sup_{\mathfrak{h}\in\mathcal{S}}e^{-\Theta(\mathfrak{h})}\cdot e^{E_+^{\mathrm{lin}}N^4},
\end{equation}
together with the lower bound $\mathcal{N}\ge e^{-C_{\mathrm{low}}N^4}$ for the partition function
$\mathcal{N}$ of the member's measure, also part of the correlation theorem. Then
\begin{equation}\label{8.12:eq-volume-ranges-b}
\tilde{\nu}\bigl(\{\mathfrak{h}\in\mathcal{S}\}\bigr)\ \le\
\exp\Bigl(\mathsf{B}N^4-\inf_{\mathfrak{h}\in\mathcal{S}}\Theta(\mathfrak{h})\Bigr).
\end{equation}
\end{definition}

\begin{proof}[Proof of \eqref{8.12:eq-volume-ranges-b}; reduction of the Case V bound]
By the definition of the history law $\tilde{\nu}$, the number
$\tilde{\nu}(\{\mathfrak{h}\in\mathcal{S}\})$ is the sum of the integrals
$\int d\tau_h$ over the histories $h$ with $\mathfrak{h}_h\in\mathcal{S}$, divided by the
partition function $\mathcal{N}$. By \eqref{8.12:eq-volume-ranges-7} the numerator is at most
$\sup_{\mathcal{S}}e^{-\Theta}\,e^{E_+^{\mathrm{lin}}N^4}$, and by the lower bound on
$\mathcal{N}$ we have $\mathcal{N}^{-1}\le e^{C_{\mathrm{low}}N^4}$. Since
$\sup_{\mathcal{S}}e^{-\Theta}=e^{-\inf_{\mathcal{S}}\Theta}$, we obtain
\begin{equation*}
\tilde{\nu}\bigl(\{\mathfrak{h}\in\mathcal{S}\}\bigr)
\le e^{C_{\mathrm{low}}N^4}\,e^{E_+^{\mathrm{lin}}N^4}\,e^{-\inf_{\mathcal{S}}\Theta}
=\exp\Bigl(\bigl(E_+^{\mathrm{lin}}+C_{\mathrm{low}}\bigr)N^4-\inf_{\mathcal{S}}\Theta\Bigr),
\end{equation*}
and $E_+^{\mathrm{lin}}+C_{\mathrm{low}}=\mathsf{B}$ by \eqref{8.12:eq-volume-ranges-4}. This is
\eqref{8.12:eq-volume-ranges-b}.

For the bound in Case V, recall $m_{\mathbb{T}}(g')=11+m'+\log_LR(g')$ and $M=L^{m'}$, so that
$L^{m_{\mathbb{T}}(g^{\dagger})}=L^{11}M\,R(g^{\dagger})$. Inserting this into
\eqref{8.12:eq-volume-ranges-3} gives
\begin{equation*}
(N^{\dagger})^4\le\bigl(\bar{n}\,L^{c_{\mathrm{gap}}+11}M\,R(g^{\dagger})\bigr)^4 .
\end{equation*}
The right side depends on $g^{\dagger}$ only through $R(g^{\dagger})$, that is, through
$X^{\dagger}=(g^{\dagger})^{-2}$, and otherwise only on $L$, $M$, $\bar{n}$ and the fixed
constant $c_{\mathrm{gap}}$. That this function of $(X^{\dagger},L,M,\bar{n})$ decreases with the
coupling and, $R$ being of logarithmic degree $r_{\mathrm{pr}}$, is polylogarithmic of degree
$4r_{\mathrm{pr}}$ in $\log X^{\dagger}$, is the content of the lemma bounding the terminal cell
count in Case V by a decreasing polylogarithmic function of the coupling, used here as stated.
\end{proof}

\begin{proposition}[Rarity of never-healed old regions on weakly coupled tori]
\label{8.12:prop-weak-torus-rarity}
We use the bare datum $(x_0,\mathsf{N})$, its members $(g,N_T)$ with $X:=g^{-2}$ and first passage
$K=K(x_0,g)$, their runs and prolonged runs, the letters $\gamma_*$, $X_*=\gamma_*^{-2}$, $K_*$,
$\bar{n}$, $c_{\mathrm{gap}}$ and the two volume ranges, the ceiling $\gamma_G$ on the member
coupling, the torus excess $a$ and the multiplicity $n_0$ of a member, the constant $\gamma_0$ of the
linear reserve, and the constants $n_1$, $C_S$ of the region rules that enter the age allowance, all
as fixed in Notation~\ref{8.12:not-volume-ranges}; $\tilde{\nu}$ is the history law at the observation
level $k_0$. For an internal history $\mathfrak{h}$ of a live component, $\mathcal{Z}(\mathfrak{h})$ is
its seed set. A seed $Z$ is born at level $j(Z)$, is either of the first kind or preparatory, has size
parameter $d'(Z):=d'_{j(Z)}(Z)$ at its birth level ($d'(Z):=0$ for a preparatory seed), and has age
allowance $\mathfrak{a}(Z)$. Let $\theta>0$ and $\theta'>0$ be the flat and the linear reserve
fractions, and let $\theta'_V,\theta'_{\mathrm{age}}>0$ be fixed with
$\theta'_V+\theta'_{\mathrm{age}}\le\theta'$. The total reserve of an internal history is
\begin{equation}\label{8.12:eq-weak-torus-rarity-1}
\Theta(\mathfrak{h}):=\sum_{Z\in\mathcal{Z}(\mathfrak{h})}
\Bigl[\theta\,p_0(g_{j(Z)})+\theta'\gamma_0A_1\,p_0(g_{j(Z)})^2\bigl(d'(Z)+1\bigr)\Bigr],
\end{equation}
with the degree function $p_0(g')=A_0(\log g'^{-2})^{p_0}$.

Let $(g,N_T=n_0L^a)$, with $g\le\gamma_G$ and $n_0\le\bar{n}$, be a member in the first volume range
$K+a\le K_*-1$ of \eqref{8.12:eq-volume-ranges-a}; let $k_0\le K-m_{\mathbb{T}}(g)$ be a level and
$\square$ a level-$k_0$ cell; and let $k^{\dagger}$, $g^{\dagger}$, $K^{\dagger\prime}$, $N^{\dagger}$
be the level, coupling, terminal level and number of sites per side of the prolonged run of this
member in the first volume range, as fixed in Notation~\ref{8.12:not-volume-ranges}, with
$X^{\dagger}:=(g^{\dagger})^{-2}$. Assume the following.
\begin{enumerate}
\item[(a)] (Independence of the run.) No object of step $k$ of Ba{\l}aban's construction depends on
the target coupling, the cutoff index $K$ or the torus exponent $m$; two runs over the same bare datum
coincide through the last step of the shorter one, and for an event of a level not beyond that step
$\tilde{\nu}$ is the same number in both runs.
\item[(b)] (Independence from the torus exponent.) The history law $\tilde{\nu}$ of events of level
$k$ does not depend on the torus exponent $m$ at which an admissible run passing level $k$ stops.
\item[(c)] (Prolongation.) Prolongation on the same lattice to an attainable level
$k^{\dagger}\le\hat{\kappa}:=\min\{K+a,K_*-1\}$ gives a member $(g^{\dagger},N_T^{\dagger})$ over the
same bare datum with $g\le g^{\dagger}<\gamma_*$, $m_{\mathbb{T}}(g^{\dagger})\le m_{\mathbb{T}}(g)$,
terminal level $K^{\dagger\prime}:=k^{\dagger}-m_{\mathbb{T}}(g^{\dagger})\ge K-m_{\mathbb{T}}(g)$, and
$N^{\dagger}:=N_T^{\dagger}L^{m_{\mathbb{T}}(g^{\dagger})}$ sites per side at $K^{\dagger\prime}$.
\item[(d)] (Reserve-stripped bound at the terminal level, termwise.) For every member $(g',N'_T)$
whose longest run ends at level $K'$ on $N'$ sites per side, and every set $\mathcal{S}$ of terminal
internal histories of it, the bound \eqref{8.12:eq-volume-ranges-b} holds in the form
\begin{equation*}
\tilde{\nu}\bigl(\{\mathfrak{h}\in\mathcal{S}\}\bigr)\le
\exp\Bigl(\mathsf{B}(g')N'^{\,4}-\inf_{\mathcal{S}}\Theta\Bigr),
\end{equation*}
where $\mathsf{B}(g')$ is the exponent of the global stability quotient of the member and, for each
history, the flat and the linear reserve of every seed enter $\Theta$ in
\eqref{8.12:eq-weak-torus-rarity-1} simultaneously.
\item[(e)] (The age lemma, \eqref{8.12:eq-volume-ranges-d}.) For every run, every level $k$ and every
live component $C$ at level $k$ with internal history $\mathfrak{h}$,
$\operatorname{age}_k(C)\le\sum_{Z\in\mathcal{Z}(\mathfrak{h})}\mathfrak{a}(Z)$, where
\begin{equation*}
\mathfrak{a}(Z)\le 2\log_L\bigl(d'(Z)+1\bigr)+2R_{j(Z)}+n_1+C_S+4
\end{equation*}
for a seed of the first kind and $\mathfrak{a}(Z)\le 2R_{j(Z)}+n_1+4$ for a preparatory seed;
here $R_j=R(x_j)$, where $R(\cdot)$ is the enlargement radius of Ba{\l}aban's region rules written
as a function of the inverse square coupling (the function entering $m_{\mathbb{T}}$; it is distinct
from the ball radius $R$ of the observables).
\item[(f)] (Degree bound on the stability exponent.) $\mathsf{B}(g')\le C_{\mathsf{B}}(1+\log g'^{-2})$
for every member coupling $g'$.
\item[(g)] (Comparison of degrees along the flow.) Every seed $Z$ of a live component at level $k_0$
satisfies $p_0(g_{j(Z)})\ge p_0(g_{k_0})-1$, and $p_0(g_{j(Z)})\ge p_0(g^{\dagger})-1$ whenever
$x_{j(Z)}>X^{\dagger}$.
\item[(h)] (Ordering of the flow.) Every level $j\le k_0$ has $x_j>X^{\dagger}$.
\end{enumerate}
Assume moreover the following one-sided conditions. For every $X'\ge X_*$, with $g':=X'^{-1/2}$,
\begin{equation}\label{8.12:eq-weak-torus-rarity-a}
C_{\mathsf{B}}(1+\log X')\bigl(\bar{n}L^{c_{\mathrm{gap}}}L^{m_{\mathbb{T}}(g')}\bigr)^{d}
\le\tfrac14\,\theta'_V\gamma_0A_1\bigl(p_0(g')-2\bigr)^2 ;
\end{equation}
and for every $y\ge X$,
\begin{equation}\label{8.12:eq-weak-torus-rarity-b}
\begin{gathered}
9\ln L\,\bigl(2R(y)+n_1+C_S+4\bigr)\le\theta A_0(\log y)^{p_0},\\
9\ln L\,\bigl(2R(y)+n_1+C_S+4\bigr)\le\tfrac12\,\theta'_{\mathrm{age}}\gamma_0A_1A_0^2(\log y)^{2p_0},\\
18\le\tfrac12\,\theta'_{\mathrm{age}}\gamma_0A_1A_0^2(\log y)^{2p_0}.
\end{gathered}
\end{equation}
As functions of $\log X'$ the two members of \eqref{8.12:eq-weak-torus-rarity-a} have degrees
$d\,r_{\mathrm{pr}}+1=4r_{\mathrm{pr}}+1$ and $2p_0\ge10r_{\mathrm{pr}}$, where $r_{\mathrm{pr}}$ is the
logarithmic degree of $R(\cdot)$; thus \eqref{8.12:eq-weak-torus-rarity-a} is an upper bound on
$\gamma_*$, imposed after $L$, $M$, $\bar{n}$, $\theta'_V$, $A_1$. In
\eqref{8.12:eq-weak-torus-rarity-b} the left members have degrees $r_{\mathrm{pr}}$, $r_{\mathrm{pr}}$,
$0$ in $\log y$ against $p_0$, $2p_0$, $2p_0$ on the right; it is an upper bound on $g$. These
conditions are imposed together with the remaining one-sided conditions on which hypotheses (d) and
(g) rest: a condition involving the exponent $\kappa$, the window for the flat reserve fraction
$\theta$, the window of the linear reserve, and the comparison of degrees along the flow stated in
hypothesis (g); none of them is among the hypotheses
\textup{(H1)}--\textup{(H7)} of Notation~\ref{2.2:not-hypotheses}.

Put
\begin{equation*}
\mathcal{E}^{\dagger}_{\square}:=\bigl\{h:\ \square\text{ lies in a component }C\text{ of }h
\text{ live at }k_0\text{ whose lineage is live at }K^{\dagger\prime}\bigr\},
\end{equation*}
the live event of $\square$, and $\operatorname{age}(C):=k_0-(\text{the oldest layer index of }C)$.
Then for every integer $\mathsf{a}\ge1$
\begin{equation}\label{8.12:eq-weak-torus-rarity-c}
\tilde{\nu}\bigl(\mathcal{E}^{\dagger}_{\square}\cap\{\operatorname{age}(C)=\mathsf{a}\}\bigr)
\le e^{\theta}e^{-\theta p_0(g_{k_0})}L^{-9\mathsf{a}} .
\end{equation}
Hence, with $q^{\mathcal{E}^{\dagger},=\mathsf{a}}_{k_0}$ the maximum over level-$k_0$ cells
$\square$ of the left member of \eqref{8.12:eq-weak-torus-rarity-c},
$q^{\mathcal{E}^{\dagger},=\mathsf{a}}_{k_0}\le e^{\theta}e^{-\theta p_0(g_{k_0})}L^{-9\mathsf{a}}$ and
\begin{equation*}
\begin{split}
\sum_{\mathsf{a}>\mathfrak{l}}L^{8(\mathsf{a}+1)}q^{\mathcal{E}^{\dagger},=\mathsf{a}}_{k_0}
&\le e^{\theta}(L-1)^{-1}L^{8-\mathfrak{l}}e^{-\theta p_0(g_{k_0})}\\
&\le A^{\mathcal{E}}_{\mathrm{lv}}\,L^{8(\mathfrak{l}+1)}\,q^{\mathrm{lv}}_{k_0},
\qquad A^{\mathcal{E}}_{\mathrm{lv}}:=\frac{e^{\theta}}{(L-1)C_{\mathrm{lv}}},
\end{split}
\end{equation*}
where $q^{\mathrm{lv}}_{k_0}=C_{\mathrm{lv}}e^{-c_{\mathrm{lv}}p_0(g_{k_0})}$, with
$c_{\mathrm{lv}}\le\theta$, is the young live weight of \eqref{8.12:eq-old-region-moments-a}. These
bounds hold uniformly in $K$, $g_0$, the position of $\square$, and the member throughout the first
volume range, without any depth floor.
\end{proposition}

\begin{proof}
\emph{The same measure.} By hypothesis (c) the prolonged run of the member $(g^{\dagger},N_T^{\dagger})$
runs over the same bare datum as the member $(g,N_T)$ and has terminal level
\begin{equation*}
K^{\dagger\prime}\ge K-m_{\mathbb{T}}(g)\ge k_0 ,
\end{equation*}
the second inequality being the choice of $k_0$. By hypotheses (a) and (b), a run of $(g,N_T)$ passing
level $k_0$ and the prolonged run coincide through level $k_0$, and $\tilde{\nu}$ gives every event of
level $k_0$ the same mass in both. Hence the level-$k_0$ events and their $\tilde{\nu}$-masses are
those of the original member, and the fate of the live lineage of $\square$ may be decided at the
prolonged terminal level $K^{\dagger\prime}$ fixed in Notation~\ref{8.12:not-volume-ranges}.
Let $\mathrm{LV}^{=\mathsf{a}}_{k_0}(\square)$
be the event that $\square$ lies in a component
live at $k_0$ of age $\mathsf{a}$, and $\mathrm{HL}^{=\mathsf{a}}_{k_0}(\square)$ its part on which the
lineage of that component is healed for the first time at some level $k$ with
$k_0<k\le K^{\dagger\prime}$. Either the lineage is so healed or it is live at $K^{\dagger\prime}$, and
not both; thus
\begin{equation*}
\mathrm{LV}^{=\mathsf{a}}_{k_0}(\square)=\mathrm{HL}^{=\mathsf{a}}_{k_0}(\square)\ \sqcup\
\bigl(\mathcal{E}^{\dagger}_{\square}\cap\{\operatorname{age}(C)=\mathsf{a}\}\bigr).
\end{equation*}
The first part is the object of the healed old-region bound in the prolonged run; the estimate below
concerns the second part only. The constants $K_H$, $e^{\theta}$ and
$S_*=\sum_{k>k_0}e^{-\frac32p_0(g_k)}$ of the healed old-region bound are bounded uniformly in
$g^{\dagger}$: indeed
\begin{equation*}
S_*\le1+(\lambda^{\sharp})^{-1}\sum_{n\ge0}
e^{-\frac32A_0\left(\log(X_*+\lambda^{\sharp}n-c_5)\right)^{p_0}},
\end{equation*}
and the right member is an absolute number once $X_*$ exceeds a threshold; this is the summability of
the degree function along the flow, read on the prolonged flow, all of whose couplings lie below
$\gamma_*$. So the healed old-region bound holds in the prolonged run as stated, with constants
depending on $\gamma_*$.

\emph{The event as a set of terminal histories.} On
$\mathcal{E}^{\dagger}_{\square}\cap\{\operatorname{age}(C)=\mathsf{a}\}$ let $C'$ be the component
live at $K^{\dagger\prime}$ that contains the descendant of $C$, and let $\mathfrak{h}_C$ be the
internal history of $C$ up to level $k_0$. The internal history $\mathfrak{h}_{C'}$ of $C'$ satisfies
$\mathcal{Z}(\mathfrak{h}_{C'})\supseteq\mathcal{Z}(\mathfrak{h}_C)$, since the descendant of a live
component inherits its seeds; extra seeds only add nonnegative reserves to
\eqref{8.12:eq-weak-torus-rarity-1}. Let $\mathcal{S}_{\square,\mathsf{a}}$ be the set of terminal
internal histories so arising. It is a set of terminal histories of the member
$(g^{\dagger},N_T^{\dagger})$: whether $h$ contributes is a property of $h$ alone, namely which
component holds $\square$ at $k_0$, its oldest layer, and whether its lineage reaches
$K^{\dagger\prime}$. It is the event that the lineage of the component holding $\square$ is live at
$K^{\dagger\prime}$, refined by the age,
and
\begin{equation*}
\mathcal{E}^{\dagger}_{\square}\cap\{\operatorname{age}(C)=\mathsf{a}\}\subseteq
\{\mathfrak{h}\in\mathcal{S}_{\square,\mathsf{a}}\}.
\end{equation*}

\emph{The reserve on $\mathcal{S}_{\square,\mathsf{a}}$.} Let
$\mathfrak{h}\in\mathcal{S}_{\square,\mathsf{a}}$, with seed set
$\mathcal{Z}\supseteq\mathcal{Z}_C:=\mathcal{Z}(\mathfrak{h}_C)$; let $Z_{\mathrm{old}}\in\mathcal{Z}_C$
be the oldest seed, and pick any seed $Z_1\in\mathcal{Z}_C$ of the first kind (a live component has
one). All terms of \eqref{8.12:eq-weak-torus-rarity-1} are nonnegative. We bound $\Theta(\mathfrak{h})$
from below by three disjoint shares.

The first share is the fraction $\theta'_V$ of the linear reserve of $Z_1$. Since
$j(Z_1)\le k_0$, hypothesis (h) gives $x_{j(Z_1)}>X^{\dagger}$.
Then hypothesis (g) gives
\begin{equation*}
p_0(g_{j(Z_1)})\ge p_0(g^{\dagger})-1\ge p_0(g^{\dagger})-2 ,
\end{equation*}
and with $d'(Z_1)+1\ge1$,
\begin{equation*}
\theta'_V\gamma_0A_1\,p_0(g_{j(Z_1)})^2\bigl(d'(Z_1)+1\bigr)\ge
\theta'_V\gamma_0A_1\bigl(p_0(g^{\dagger})-2\bigr)^2=:t_V .
\end{equation*}

The second share is the flat reserve of $Z_{\mathrm{old}}$. The oldest seed is born at or below level
$k_0$, so hypothesis (g) gives
$\theta p_0(g_{j(Z_{\mathrm{old}})})\ge\theta\bigl(p_0(g_{k_0})-1\bigr)$.

The third share consists of the flat reserves of all $Z\in\mathcal{Z}_C\setminus\{Z_{\mathrm{old}}\}$
and the fraction $\theta'_{\mathrm{age}}$ of the linear reserves of all $Z\in\mathcal{Z}_C$. By
hypothesis (e) at level $k_0$, applied to $C$ with history $\mathfrak{h}_C$,
$\mathsf{a}=\operatorname{age}(C)\le\sum_{Z\in\mathcal{Z}_C}\mathfrak{a}(Z)$. Fix
$Z\in\mathcal{Z}_C$ and put $y:=x_{j(Z)}$, so that $p_0(g_{j(Z)})=A_0(\log y)^{p_0}$ and
$R_{j(Z)}=R(y)$; the conditions \eqref{8.12:eq-weak-torus-rarity-b} apply at this $y$, since
$j(Z)\le k_0<K$ and $K=K(x_0,g)$ is the first passage of the flow to $X$. For a seed of the first kind,
$9\ln L\cdot2\log_L(d'(Z)+1)=18\ln(d'(Z)+1)$, so
\begin{equation*}
9\ln L\,\mathfrak{a}(Z)\le18\ln\bigl(d'(Z)+1\bigr)+9\ln L\,\bigl(2R(y)+n_1+C_S+4\bigr).
\end{equation*}
The size term is paid from the linear reserve: by $\ln u\le u$ and the third line of
\eqref{8.12:eq-weak-torus-rarity-b},
\begin{equation*}
18\ln\bigl(d'(Z)+1\bigr)\le18\bigl(d'(Z)+1\bigr)
\le\tfrac12\,\theta'_{\mathrm{age}}\gamma_0A_1\,p_0(g_{j(Z)})^2\bigl(d'(Z)+1\bigr).
\end{equation*}
The radius term $9\ln L(2R(y)+n_1+C_S+4)$ is paid, for $Z\ne Z_{\mathrm{old}}$, from the flat reserve
by the first line of \eqref{8.12:eq-weak-torus-rarity-b}, being at most $\theta p_0(g_{j(Z)})$; and
for $Z=Z_{\mathrm{old}}$ from its linear reserve by the second line, being at most
\begin{equation*}
\tfrac12\,\theta'_{\mathrm{age}}\gamma_0A_1\,p_0(g_{j(Z)})^2\le
\tfrac12\,\theta'_{\mathrm{age}}\gamma_0A_1\,p_0(g_{j(Z)})^2\bigl(d'(Z)+1\bigr).
\end{equation*}
For a preparatory seed
$d'(Z)=0$, its allowance bound $2R(y)+n_1+4$ has neither the size term nor $C_S$, and it is paid by
the same members as displayed. Each linear reserve is used at most with total fraction
$\frac12\theta'_{\mathrm{age}}+\frac12\theta'_{\mathrm{age}}$; summing over $Z\in\mathcal{Z}_C$,
\begin{equation*}
\theta\!\!\sum_{Z\in\mathcal{Z}_C\setminus\{Z_{\mathrm{old}}\}}\!\!p_0(g_{j(Z)})
+\theta'_{\mathrm{age}}\gamma_0A_1\sum_{Z\in\mathcal{Z}_C}p_0(g_{j(Z)})^2\bigl(d'(Z)+1\bigr)
\ge9\ln L\sum_{Z\in\mathcal{Z}_C}\mathfrak{a}(Z)\ge9\mathsf{a}\ln L .
\end{equation*}

The three shares are disjoint: they use different seeds or different fixed fractions of one seed's
linear reserve, and $\theta'_V+\theta'_{\mathrm{age}}\le\theta'$; this holds also when
$Z_1=Z_{\mathrm{old}}$, since the first share uses the linear and the second the flat reserve of that
seed. Hence
\begin{equation}\label{8.12:eq-weak-torus-rarity-2}
\inf_{\mathcal{S}_{\square,\mathsf{a}}}\Theta\ge t_V+\theta\bigl(p_0(g_{k_0})-1\bigr)+9\mathsf{a}\ln L .
\end{equation}

\emph{The global quotient at the prolonged terminal level.} The longest run of
$(g^{\dagger},N_T^{\dagger})$ ends at $K^{\dagger\prime}$ on $N^{\dagger}$ sites per side. Hypothesis
(d) for this member and $\mathcal{S}=\mathcal{S}_{\square,\mathsf{a}}$, together with
\eqref{8.12:eq-weak-torus-rarity-2}, gives
\begin{equation}\label{8.12:eq-weak-torus-rarity-3}
\tilde{\nu}\bigl(\{\mathfrak{h}\in\mathcal{S}_{\square,\mathsf{a}}\}\bigr)\le
\exp\Bigl(\mathsf{B}(g^{\dagger})N^{\dagger\,4}-t_V-\theta\bigl(p_0(g_{k_0})-1\bigr)
-9\mathsf{a}\ln L\Bigr).
\end{equation}
In the first volume range $N^{\dagger}\le\bar{n}L^{c_{\mathrm{gap}}}L^{m_{\mathbb{T}}(g^{\dagger})}$ by
\eqref{8.12:eq-volume-ranges-a}. Since $g^{\dagger}<\gamma_*$ by hypothesis (c), $X^{\dagger}>X_*$, and
\eqref{8.12:eq-weak-torus-rarity-a} applies at $X'=X^{\dagger}$. With hypothesis (f) and $d=4$,
\begin{equation*}
\begin{split}
\mathsf{B}(g^{\dagger})N^{\dagger\,4}
&\le C_{\mathsf{B}}(1+\log X^{\dagger})
\bigl(\bar{n}L^{c_{\mathrm{gap}}}L^{m_{\mathbb{T}}(g^{\dagger})}\bigr)^{d}\\
&\le\tfrac14\,\theta'_V\gamma_0A_1\bigl(p_0(g^{\dagger})-2\bigr)^2=\tfrac14\,t_V\le t_V .
\end{split}
\end{equation*}
Hence the exponent in \eqref{8.12:eq-weak-torus-rarity-3} is at most
$\theta-\theta p_0(g_{k_0})-9\mathsf{a}\ln L$, and
\begin{equation*}
\tilde{\nu}\bigl(\{\mathfrak{h}\in\mathcal{S}_{\square,\mathsf{a}}\}\bigr)\le
e^{\theta}e^{-\theta p_0(g_{k_0})}L^{-9\mathsf{a}} .
\end{equation*}
By the inclusion of the second paragraph and the
identification of level-$k_0$ masses in the first, this is \eqref{8.12:eq-weak-torus-rarity-c}.

\emph{The series.} The right member of \eqref{8.12:eq-weak-torus-rarity-c} does not depend on
$\square$, so it bounds $q^{\mathcal{E}^{\dagger},=\mathsf{a}}_{k_0}$. Further,
\begin{equation*}
\sum_{\mathsf{a}>\mathfrak{l}}L^{8(\mathsf{a}+1)}L^{-9\mathsf{a}}
=L^{8}\sum_{\mathsf{a}>\mathfrak{l}}L^{-\mathsf{a}}=\frac{L^{8-\mathfrak{l}}}{L-1},
\end{equation*}
which gives the first inequality of the series. For the second, $c_{\mathrm{lv}}\le\theta$ gives
\begin{equation*}
e^{-\theta p_0(g_{k_0})}\le e^{-c_{\mathrm{lv}}p_0(g_{k_0})}=q^{\mathrm{lv}}_{k_0}/C_{\mathrm{lv}},
\end{equation*}
and $L^{8-\mathfrak{l}}\le L^{8(\mathfrak{l}+1)}$ for the threshold $\mathfrak{l}\ge0$; together these
give the bound by $A^{\mathcal{E}}_{\mathrm{lv}}L^{8(\mathfrak{l}+1)}q^{\mathrm{lv}}_{k_0}$.

No average over translations, no first moment and no covariance bound enters; the bound is uniform in
the position of $\square$ as it stands. The coupling $g^{\dagger}$ may lie far above $g$, up to
$\gamma_*$. This does not affect the bound: $\mathsf{B}(g^{\dagger})$ and $N^{\dagger}$ decrease with
the coupling, the rate $p_0(g_{k_0})$ is read at $g_{k_0}\le g$, and the age reserve is read at the
birth couplings of the seeds.
\end{proof}

\begin{corollary}[The age series and the two-point witness on weakly coupled tori]
\label{8.12:cor-weak-torus-witness}
Let $g$ be a reference coupling and $N_T$ a number of sites per side of the terminal lattice, with
$(g,N_T)$ in the weakly coupled range, Case~V of the volume ranges
\eqref{8.12:eq-volume-ranges-a}. Assume, for every level $k_0\le K-m_{\mathbb{T}}(g)$:
\begin{itemize}
\item[(a)] the hypotheses of Proposition~\ref{8.12:prop-weak-torus-rarity}, so that
\eqref{8.12:eq-weak-torus-rarity-c} holds for every integer $\mathsf{a}\ge1$; in particular we
use its prolonged run, with prolonged terminal level $K^{\dagger\prime}$, and the young-live weight
$q^{\mathrm{lv}}_{k_0}=C_{\mathrm{lv}}e^{-c_{\mathrm{lv}}p_0(g_{k_0})}$, where $p_0(\cdot)$ is the
degree function and $c_{\mathrm{lv}}\le\theta$;
\item[(b)] the healed old-region bound holds at level $k_0$ for that prolonged run, with its
constant $A^{\mathrm{heal}}_{\mathrm{lv}}$ read at the threshold coupling $\gamma_*$, that is,
\begin{equation}\label{8.12:eq-weak-torus-witness-1}
\sum_{\mathsf{a}>\mathfrak{l}}L^{8(\mathsf{a}+1)}\,q^{\mathrm{heal},=\mathsf{a}}_{k_0}
\le A^{\mathrm{heal}}_{\mathrm{lv}}\,L^{8(\mathfrak{l}+1)}\,q^{\mathrm{lv}}_{k_0},
\end{equation}
where $q^{\mathrm{heal},=\mathsf{a}}_{k_0}$ is the maximum over level-$k_0$ cells $\square$ of the
$\tilde{\nu}$-mass of the event that $\square$ lies in a component live at $k_0$ of age
$\mathsf{a}$ whose lineage is first healed at some level $k$ with $k_0<k\le K^{\dagger\prime}$;
\item[(c)] the first input at level $k_0$ of the additive theorem, supplied by the corollary on
rarity of young large-field blocks; here the additive theorem is the statement whose moment and
covariance bounds on the near live source slots control the remainders $R^{\mathfrak{L}}$ and
$R^{\Psi,\mathrm{live}}$ of the two-point witness theorem.
\end{itemize}
Assume moreover (d): the moment and covariance bounds of the additive theorem hold, from its first
input and, as second input, from \eqref{8.12:eq-old-region-moments-a}, at every level
$k_0\le K-m_{\mathbb{T}}(g)$, under its one-sided ceiling on the reference coupling.
Then the old-region moment bound \eqref{8.12:eq-old-region-moments-a} holds at every level
$k_0\le K-m_{\mathbb{T}}(g)$ with
\begin{equation}\label{8.12:eq-weak-torus-witness-2}
A_{\mathrm{lv}}:=A^{\mathrm{heal}}_{\mathrm{lv}}+\frac{e^{\theta}}{(L-1)\,C_{\mathrm{lv}}}.
\end{equation}
Consequently, under the additive theorem's one-sided ceiling on the reference coupling, its moment
and covariance bounds, the assembly of the remainder bounds in the proof of the two-point witness
theorem, and the bound $|R_K|\le\tfrac14$ all hold on the tori of Case~V, with no term added to the
witness depth $s_W(g)$.
\end{corollary}

\begin{proof}
Fix a level $k_0\le K-m_{\mathbb{T}}(g)$, a level-$k_0$ cell $\square$ and an age
$\mathsf{a}>\mathfrak{l}$. The fate decomposition of the live lineage of $\square$, taken at the
prolonged terminal level $K^{\dagger\prime}\ge k_0$ as in the proof of
Proposition~\ref{8.12:prop-weak-torus-rarity}, writes the event that $\square$ lies in a component
live at $k_0$ of age $\mathsf{a}$ as the disjoint union of two events. In the first, the lineage
of that component is first healed at some level $k$ with $k_0<k\le K^{\dagger\prime}$. The second
is $\mathcal{E}^{\dagger}_{\square}\cap\{\operatorname{age}(C)=\mathsf{a}\}$, where $C$ is the
component live at $k_0$ containing $\square$ and we write
$\mathcal{E}^{\dagger}_{\square}$ for the live event that $\square$ lies in a component live at
$k_0$ whose lineage is still live at $K^{\dagger\prime}$. Hence the
$\tilde{\nu}$-mass of the event is the sum of the masses of the two parts. Each part's mass is
bounded by its maximum over cells. Taking the maximum over $\square$ therefore gives
\begin{equation*}
q^{=\mathsf{a}}_{k_0}\le q^{\mathrm{heal},=\mathsf{a}}_{k_0}
+q^{\mathcal{E}^{\dagger},=\mathsf{a}}_{k_0},
\end{equation*}
where $q^{=\mathsf{a}}_{k_0}$ is the quantity of \eqref{8.12:eq-old-region-moments-a} and
$q^{\mathcal{E}^{\dagger},=\mathsf{a}}_{k_0}$ denotes the maximum over level-$k_0$ cells $\square$
of $\tilde{\nu}(\mathcal{E}^{\dagger}_{\square}\cap\{\operatorname{age}(C)=\mathsf{a}\})$.
The first event is the object of the healed old-region bound in the prolonged run, whose constants
are uniform in the prolonged run's coupling $g^{\dagger}\le\gamma_*$ by the argument in the proof
of Proposition~\ref{8.12:prop-weak-torus-rarity}; this is why (b) is stated with
$A^{\mathrm{heal}}_{\mathrm{lv}}$ read at $\gamma_*$.

From \eqref{8.12:eq-weak-torus-rarity-c}, which is uniform in the position of $\square$,
Proposition~\ref{8.12:prop-weak-torus-rarity} derives the summed bound, written here with the
constant denoted $A^{\mathcal{E}}_{\mathrm{lv}}$:
\begin{equation}\label{8.12:eq-weak-torus-witness-3}
\sum_{\mathsf{a}>\mathfrak{l}}L^{8(\mathsf{a}+1)}q^{\mathcal{E}^{\dagger},=\mathsf{a}}_{k_0}
\le A^{\mathcal{E}}_{\mathrm{lv}}\,L^{8(\mathfrak{l}+1)}\,q^{\mathrm{lv}}_{k_0},
\qquad A^{\mathcal{E}}_{\mathrm{lv}}:=\frac{e^{\theta}}{(L-1)C_{\mathrm{lv}}}.
\end{equation}
Multiply the inequality for $q^{=\mathsf{a}}_{k_0}$ by $L^{8(\mathsf{a}+1)}$ and sum over
$\mathsf{a}>\mathfrak{l}$. Adding \eqref{8.12:eq-weak-torus-witness-1} and
\eqref{8.12:eq-weak-torus-witness-3} then yields \eqref{8.12:eq-old-region-moments-a} with
$A_{\mathrm{lv}}=A^{\mathrm{heal}}_{\mathrm{lv}}+A^{\mathcal{E}}_{\mathrm{lv}}$, which is the
constant \eqref{8.12:eq-weak-torus-witness-2}. Since $k_0$ was arbitrary, this holds at every
level $k_0\le K-m_{\mathbb{T}}(g)$.

By Proposition~\ref{8.12:prop-old-region-moments}, the age series enters the remainder of the
two-point witness theorem only at $R^{\mathfrak{L}}$ and $R^{\Psi,\mathrm{live}}$, and there only
through the bound \eqref{8.12:eq-old-region-moments-a}, which is the second input of the additive
theorem; it has just been derived at every $k_0\le K-m_{\mathbb{T}}(g)$. The first input is (c).
So, by (d), the moment and covariance bounds of the additive theorem hold on the tori of Case~V
under its one-sided ceiling on the reference coupling, and they control $R^{\mathfrak{L}}$ and
$R^{\Psi,\mathrm{live}}$. The bounds on the other remainders in the proof of the two-point witness
theorem do not involve the age series and are unchanged on these tori; the assembly of the
remainder bounds in that proof therefore gives $|R_K|\le\tfrac14$ under the same ceiling. Since
(a)--(c) and \eqref{8.12:eq-weak-torus-rarity-c} hold at every level $k_0\le K-m_{\mathbb{T}}(g)$
with no depth floor, no term is added to $s_W(g)$.
\end{proof}

\begin{remark}[Extent and physical meaning of the weakly coupled tori]\label{8.12:rem-weak-torus-extent}
We use the letters of Notation~\ref{8.12:not-volume-ranges}: the bare datum $(x_0,\mathsf{N})$, the
member coupling $g_V$ attached there to the reference coupling $g$, the first-passage level
$K=K(x_0,g_V)$, the torus excess $a=m-m_{\mathbb{T}}(g_V)$, the fixed threshold $\gamma_*$ with
$X_*=\gamma_*^{-2}$, $K_*=K(x_0,\gamma_*)$, and the constants $b_+$, $c_2$, $\lambda$ entering
\eqref{8.12:eq-volume-ranges-c}. Case~V of \eqref{8.12:eq-volume-ranges-a} (the weakly coupled
tori) is not a cap on any constant. It is a range
of members, described by a one-sided condition, and it is the range on which
Proposition~\ref{8.12:prop-weak-torus-rarity}, proved there as an implication from the statements
it names, holds.
\begin{itemize}
\item[(a)] Case~V, that is \eqref{8.12:eq-volume-ranges-a}, holds if and only if $a<K_*-K$.
It contains every member whose torus exponent satisfies
\begin{equation}\label{8.12:eq-weak-torus-extent-1}
m\;\le\; m_{\mathbb{T}}(g_V)
+\Big\lfloor \frac{g_V^{-2}-\gamma_*^{-2}-b_+-2c_2}{3\lambda}\Big\rfloor-1 .
\end{equation}
This is a one-sided condition on $m$ in which the cutoff $K$ does not appear. Its right
member grows like $g^{-2}/(3\lambda)$ as $g\downarrow0$.
\item[(b)] Case~V contains the minimal admissible torus, $a=0$, of every member with
$g_V^{-2}>X_*+b_++2c_2$.
\item[(c)] For every fixed $m_0$ there is a threshold $\eta(m_0)>0$ such that Case~V contains
$(g,m_0)$ for every $g$ below $\eta(m_0)$. Equivalently, Case~V may be read as a ceiling on
$g$ that depends on the volume, in place of a range of $m$.
\item[(d)] Case~V is a property of the bare lattice $(x_0,\mathsf{N})$ alone: the complete
renormalisation-group run, from the cutoff to the scale of the torus itself, stays below
$\gamma_*$. In particular it does not depend on the reference labelling. In physical terms, read
through the change-of-units identity, clause \textup{(v-d)(iii)} of Theorem~\ref{3.1:thm-main}
and the two-sided bounds on the running inverse couplings $x_k$ given by the flow envelopes
$\underline{x}_s$ and $\bar{x}_s$, this reads as follows. Let $\ell_*$ be the length at which the
running coupling first reaches $\gamma_*$. Since $\gamma_*$ is fixed, $\ell_*$ is a fixed
fraction, in the logarithmic sense, of the dynamical scale (the length at which the running
coupling becomes of order one). Case~V then says that the torus is smaller than
$\ell_*$; such tori are sometimes called sub-strong-coupling, or femto-universe, tori.
\item[(e)] In Case~C (the complement of Case~V) the torus is larger than $\ell_*$, and the physics
at the scale of the torus is strongly coupled. Case~C is exactly the infinite-volume direction of
clause \textup{(ii-$\mathbb{R}^4$)} of Theorem~\ref{3.1:thm-main}: a sequence $m_j\to\infty$ at
fixed bare coupling $g_0$ and fixed member coupling $g_V$ leaves Case~V after finitely many $j$.
\end{itemize}
\end{remark}

\begin{proof}
(a) The integers $K$, $a$ and $K_*$ satisfy $K+a\le K_*-1$ if and only if $a<K_*-K$. This is the
first assertion.

For the second, put
\begin{equation*}
Q:=\frac{g_V^{-2}-X_*-b_+-2c_2}{3\lambda}.
\end{equation*}
With $X=g_V^{-2}$, the member's inverse square coupling in the letters of
\eqref{8.12:eq-volume-ranges-c}, the inequality \eqref{8.12:eq-volume-ranges-c} reads
$K_*-K\ge Q$. The left
side is an integer and $\lfloor Q\rfloor\le Q$, so $K_*-K\ge\lfloor Q\rfloor$. Now suppose that $m$
satisfies \eqref{8.12:eq-weak-torus-extent-1}. Then
\begin{equation*}
a=m-m_{\mathbb{T}}(g_V)\le\lfloor Q\rfloor-1<\lfloor Q\rfloor\le K_*-K,
\end{equation*}
so the member lies in Case~V.

The right member of \eqref{8.12:eq-weak-torus-extent-1} involves only $g_V$ and the fixed
constants $\gamma_*$, $b_+$, $c_2$, $\lambda$, $m'$, $L$; in particular the cutoff $K$ does not
enter. The floor term differs from $g_V^{-2}/(3\lambda)$ by a quantity bounded independently of
$g_V$. The remaining term is $m_{\mathbb{T}}(g_V)=11+m'+\log_L R(g_V)$. Since $R(\cdot)$ has
log-degree $r_{\mathrm{pr}}$, the quantity $R(g_V)$ grows at most like the power $r_{\mathrm{pr}}$
of $\log g_V^{-2}$, so $\log_L R(g_V)=O(\log\log g_V^{-2})$, which is of lower order in
$g_V^{-2}$. Finally, $g_V\downarrow0$ as $g\downarrow0$. This gives the stated growth.

(b) For $a=0$ the inequality \eqref{8.12:eq-volume-ranges-a} reads $K\le K_*-1$, that is,
$K_*-K\ge1$. If $g_V^{-2}>X_*+b_++2c_2$, then $Q>0$. By the proof of (a), $K_*-K$ is an integer
with $K_*-K\ge Q>0$, hence $K_*-K\ge1$, and the member with $a=0$ lies in Case~V.

(c) Fix $m_0$. By (a), the right member of \eqref{8.12:eq-weak-torus-extent-1} tends to $+\infty$
as $g\downarrow0$. Hence there is $\eta(m_0)>0$ such that $m_0$ does not exceed the right
member for every $g<\eta(m_0)$. By (a), every such $(g,m_0)$ lies in Case~V. Read with $m_0$
fixed and $g$ varying, this is a ceiling on $g$ depending on $m_0$.

(d) By Notation~\ref{8.12:not-volume-ranges} we have $\mathsf{N}=N_TL^K$ and
$N_T=2L^{m-m_{\mathbb{T}}(g_V)}=2L^a$, so
\begin{equation}\label{8.12:eq-weak-torus-extent-2}
\mathsf{N}=2L^{K+a}.
\end{equation}
Thus $K+a=\log_L(\mathsf{N}/2)$ is the level at which the block lattice of
Definition~\ref{1.5:def-block-average} has two sites per side, that is, the scale of the torus
itself. The index $K_*=K(x_0,\gamma_*)$ depends only on $x_0$ and on the fixed $\gamma_*$. By
\eqref{8.12:eq-weak-torus-extent-2}, the inequality \eqref{8.12:eq-volume-ranges-a} becomes
\begin{equation*}
\log_L(\mathsf{N}/2)\le K_*-1 .
\end{equation*}
This is a condition on $(x_0,\mathsf{N})$ alone. It says that the run reaches the scale of the
torus at a level below the first passage of the running coupling to $\gamma_*$. Neither $g$ nor
$m$ enters it, so it is independent of the reference labelling.

Measured in bare spacings, a block at level $k$ has side $L^k$. The inequality above therefore
reads $\mathsf{N}<2L^{K_*}$: the torus side is smaller than twice the side of a block at the
level at which the running coupling first reaches $\gamma_*$. Read through the change-of-units
identity, clause \textup{(v-d)(iii)} of Theorem~\ref{3.1:thm-main} and the two-sided bounds on
the running inverse couplings given by the flow envelopes $\underline{x}_s$ and $\bar{x}_s$, this
is the comparison of the torus with $\ell_*$ stated in (d).

(e) Take a sequence with $g_0$ and $g_V$ held fixed and $m_j\to\infty$. Then $K$ and $K_*$
do not depend on $j$, while $a_j=m_j-m_{\mathbb{T}}(g_V)\to\infty$. Hence $a_j\ge K_*-K$ for all
but finitely many $j$. By (a), these members are in Case~C.
\end{proof}

\begin{proposition}[Rarity of never-healed old regions above a volume-dependent depth floor. Stated; proof incomplete at the step marked $(\ast)$]\label{8.12:prop-depth-floor-rarity}
We keep the notation of Notation~\ref{8.12:not-volume-ranges}. Let $(g,N_T)$, $N_T=n_0L^{a}$, be a member in
Case~C of \eqref{8.12:eq-volume-ranges-a}, that is, $K+a\ge K_*$. Assume:
\begin{itemize}
\item the statements from which Proposition~\ref{8.12:prop-weak-torus-rarity} is derived;
\item the prolongation lemma in the coupling range $K+a\ge K_*$, which provides an attainable level
$k^{\dagger}$ with $k^{\dagger}-K\ge\bar{u}(g)+c_{\mathrm{gap}}$ and the prolonged member
$(g^{\dagger},N_T^{\dagger})$ over the same bare datum;
\item the end-density bound in Case~C, which supplies the age-blind end density $\mathfrak{e}_{k_0}$ (used
only in the variant following the proof);
\item the ceiling \eqref{8.12:eq-weak-torus-rarity-b}.
\end{itemize}
Prolong the run on the same lattice by $k^{\dagger}-K\ge\bar{u}(g)+c_{\mathrm{gap}}$ levels, by the
prolongation lemma in the coupling range. Let $g^{\dagger}$ and $K^{\dagger\prime}=k^{\dagger}-m_{\mathbb{T}}(g^{\dagger})$ be the coupling and the
terminal level of the prolonged member $(g^{\dagger},N_T^{\dagger})$. Its terminal lattice has
\begin{equation}\label{8.12:eq-depth-floor-rarity-1}
N_C:=N_T^{\dagger}L^{m_{\mathbb{T}}(g^{\dagger})}=n_0L^{a-(k^{\dagger}-K)}L^{m_{\mathbb{T}}(g^{\dagger})}
\end{equation}
sites per side, a number that grows with $a$. We write $N_C=N_C(m)$ when its dependence on the torus
exponent $m$, through $a$ and $k^{\dagger}$, is to be stressed.

Let $s\ge m_{\mathbb{T}}(g)$ be an integer, the witness depth. Put $k_0:=K-s$ and let $\square$ be a level-$k_0$
cell. Let the live event $\mathcal{E}^{\dagger}_{\square}$ and $\operatorname{age}(C)$ be as in
Proposition~\ref{8.12:prop-weak-torus-rarity}. For an integer $\mathsf{a}\ge1$, let
$\mathcal{S}_{\square,\mathsf{a}}$ be the set of terminal internal histories of the prolonged member that
arise on $\mathcal{E}^{\dagger}_{\square}\cap\{\operatorname{age}(C)=\mathsf{a}\}$. Then, for every such
$\mathsf{a}$,
\begin{equation}\label{8.12:eq-depth-floor-rarity-2}
\tilde{\nu}(\mathcal{S}_{\square,\mathsf{a}})\le\exp\Bigl(\mathsf{B}(\gamma_*)N_C^4
-\theta\bigl(p_0(g_{k_0})-1\bigr)-9\ln L\,\bigl(\mathsf{a}+s-m_{\mathbb{T}}(g)+k^{\dagger}-K\bigr)\Bigr).
\end{equation}
Suppose moreover that $s$ satisfies the volume-dependent depth floor displayed below. The second line of the
display defines $s_W(g,m)$; there $s_W(g)$ is the depth threshold of the two-point witness theorem.
\begin{equation}\label{8.12:eq-depth-floor-rarity-a}
\begin{gathered}
9\ln L\,\bigl(s-m_{\mathbb{T}}(g)+k^{\dagger}-K\bigr)\ge C_{\mathsf{B}}(1+\log X_*)\,N_C(m)^4,\\
s_W(g,m):=\max\bigl\{s_W(g),\ \text{the least } s \text{ satisfying the first line}\bigr\}.
\end{gathered}
\end{equation}
Then, with $A^{\mathcal{E}}_{\mathrm{lv}}=e^{\theta}/((L-1)C_{\mathrm{lv}})$ as in
Proposition~\ref{8.12:prop-weak-torus-rarity}, and with $q^{\mathcal{E},=\mathsf{a}}_{k_0}$ the quantity
$q^{\mathcal{E}\dagger,=\mathsf{a}}_{k_0}$ of Proposition~\ref{8.12:prop-weak-torus-rarity} for the
prolonged run,
\begin{equation}\label{8.12:eq-depth-floor-rarity-3}
\sum_{\mathsf{a}>\mathfrak{l}}L^{8(\mathsf{a}+1)}q^{\mathcal{E},=\mathsf{a}}_{k_0}
\le A^{\mathcal{E}}_{\mathrm{lv}}L^{8(\mathfrak{l}+1)}q^{\mathrm{lv}}_{k_0}.
\end{equation}
Consequently, on every torus of the family of members described in Notation~\ref{8.12:not-volume-ranges},
the bound \eqref{8.12:eq-old-region-moments-a} holds for $s\ge s_W(g,m)$, and hence so does the two-point
witness theorem. These conclusions hold uniformly in $K$, in $g_0$ and in the position of $\square$. They are
not uniform in $m$.
\end{proposition}

Stated; proof incomplete at the step marked $(\ast)$.

The first line of \eqref{8.12:eq-depth-floor-rarity-a} is a one-sided lower bound on the witness depth $s$.
For every $(g,m)$ its solutions form a half-line in $s$, and it bounds none of the constants. It depends on the
volume only through $N_C(m)^4$, and $N_C(m)^4$ is comparable to
\begin{equation*}
L^{4(a-(K_*-K))}\,L^{4m_{\mathbb{T}}(\gamma_*)},
\end{equation*}
where $a$ is the torus excess of Notation~\ref{8.12:not-volume-ranges}.
This is exponential in the Case-C excess $a-(K_*-K)$. As $m$ decreases into Case~V, that is, into the range
$K+a\le K_*-1$ of \eqref{8.12:eq-volume-ranges-a}, the floor passes to
$s_W(g)$, and Proposition~\ref{8.12:prop-weak-torus-rarity} needs no depth floor.

\begin{proof}
Steps (1)--(3) below are those of the proof of Proposition~\ref{8.12:prop-weak-torus-rarity}, which is proved
there as an implication from the statements listed with it. We recall them in the form needed here, for the
member $(g,N_T)$ in Case~C and its prolonged member $(g^{\dagger},N_T^{\dagger})$.

(1) \emph{Same measure.} Two runs over the same bare datum coincide through the last step of the shorter
one, and the terminal-level comparison of Notation~\ref{8.12:not-volume-ranges} gives
$K^{\dagger\prime}\ge K-m_{\mathbb{T}}(g)\ge k_0$. Hence the level-$k_0$ events and their
$\tilde{\nu}$-masses are those of the original member. The fate decomposition of the live lineage of
$\square$ may be taken at $K^{\dagger\prime}$: the event that $\square$ lies in a live component of age
$\mathsf{a}$ at $k_0$ is the disjoint union of two events.
\begin{itemize}
\item The lineage is first healed at some $k\in\,]k_0,K^{\dagger\prime}]$. This part is the object of the
healed old-region bound in the prolonged run, and its constants depend on the couplings only through the
ceiling $\gamma_*$.
\item The event $\mathcal{E}^{\dagger}_{\square}\cap\{\operatorname{age}=\mathsf{a}\}$.
\end{itemize}

(2) \emph{The event as a set of terminal histories.} On
$\mathcal{E}^{\dagger}_{\square}\cap\{\operatorname{age}(C)=\mathsf{a}\}$ the terminal-live component that
contains the descendant of $C$ has an internal history whose seed set contains $\mathcal{Z}(\mathfrak{h}_C)$.
Here $\mathfrak{h}_C$ is the history of $C$ up to $k_0$. Membership in $\mathcal{S}_{\square,\mathsf{a}}$ is a
property of the history $h$, so $\mathcal{S}_{\square,\mathsf{a}}$ is a set of terminal histories of the member
$(g^{\dagger},N_T^{\dagger})$.

(3) \emph{The reserve.} For $\mathfrak{h}\in\mathcal{S}_{\square,\mathsf{a}}$ the total reserve
$\Theta(\mathfrak{h})$ contains three disjoint shares:
\begin{itemize}
\item[(i)] the fraction $\theta'_V$ of the linear member of a seed of $C$ of the first kind, that is, a seed
that is not preparatory and whose allowance $\mathfrak{a}$ carries a size term; every live component has
such a seed;
\item[(ii)] the flat member of the oldest seed $Z_{\mathrm{old}}$ of $C$, which is at least
$\theta(p_0(g_{k_0})-1)$, because that seed is born at or below $k_0$;
\item[(iii)] the flat members of the seeds in $\mathcal{Z}(\mathfrak{h}_C)$ other than $Z_{\mathrm{old}}$,
together with the fraction $\theta'_{\mathrm{age}}$ of the linear members of all seeds in
$\mathcal{Z}(\mathfrak{h}_C)$. By \eqref{8.12:eq-weak-torus-rarity-b} this share is at least
$9\ln L\sum_{Z\in\mathcal{Z}(\mathfrak{h}_C)}\mathfrak{a}(Z)$.
\end{itemize}

(4) $(\ast)$ \emph{The terminal age and the global quotient.} In this step two inputs are used beyond the
range in which they are displayed: the ceiling \eqref{8.12:eq-weak-torus-rarity-b} is applied to seeds
created after level $K$, whose birth couplings lie between $g$ and $\gamma_*$, and the stability exponent of
the prolonged member is read at $\gamma_*$ in place of $g^{\dagger}$. Share (i) is nonnegative and is not used
here.
Consider a component of age $\mathsf{a}$ at $k_0$ that is live at $K^{\dagger\prime}$. The seed set of the
terminal-live component containing its descendant contains $\mathcal{Z}(\mathfrak{h}_C)$ (step (2)), so its
oldest layer index is at most $k_0-\mathsf{a}$ and its age at level $K^{\dagger\prime}$ is at least
$\mathsf{a}+(K^{\dagger\prime}-k_0)$. Apply the age lemma, in the form \eqref{8.12:eq-volume-ranges-d}, at
level $K^{\dagger\prime}$ in place of $k_0$:
\begin{equation*}
\sum_{Z\in\mathcal{Z}(\mathfrak{h})}\mathfrak{a}(Z)\ge\mathsf{a}+K^{\dagger\prime}-k_0.
\end{equation*}
The allocation of share (iii), applied to every seed of $\mathcal{Z}(\mathfrak{h})$ with
\eqref{8.12:eq-weak-torus-rarity-b} used at birth couplings up to $\gamma_*$, enlarges share (iii) to the flat
members of the seeds in $\mathcal{Z}(\mathfrak{h})$ other than $Z_{\mathrm{old}}$ together with the fraction
$\theta'_{\mathrm{age}}$ of the linear members of all seeds in $\mathcal{Z}(\mathfrak{h})$, and gives that
this share is at least $9\ln L\sum_{Z\in\mathcal{Z}(\mathfrak{h})}\mathfrak{a}(Z)$.
Hence share (iii) is at least
$9\ln L\,(\mathsf{a}+K^{\dagger\prime}-k_0)$. Since $K^{\dagger\prime}=k^{\dagger}-m_{\mathbb{T}}(g^{\dagger})$
and $k_0=K-s$, and since the terminal-level comparison of Notation~\ref{8.12:not-volume-ranges} gives
$m_{\mathbb{T}}(g^{\dagger})\le m_{\mathbb{T}}(g)$,
\begin{equation*}
K^{\dagger\prime}-k_0=\bigl(k^{\dagger}-m_{\mathbb{T}}(g^{\dagger})\bigr)-(K-s)
\ge s-m_{\mathbb{T}}(g)+(k^{\dagger}-K).
\end{equation*}
Adding share (ii), we obtain
\begin{equation*}
\inf_{\mathcal{S}_{\square,\mathsf{a}}}\Theta
\ge\theta\bigl(p_0(g_{k_0})-1\bigr)+9\ln L\,\bigl(\mathsf{a}+s-m_{\mathbb{T}}(g)+k^{\dagger}-K\bigr).
\end{equation*}
We apply \eqref{8.12:eq-volume-ranges-b} termwise to the prolonged member, whose longest run ends at
$K^{\dagger\prime}$ on $N_C$ sites per side; it gives the exponent $\mathsf{B}(g^{\dagger})N_C^4$ of the
prolonged member. In \eqref{8.12:eq-depth-floor-rarity-2} and in the floor
\eqref{8.12:eq-depth-floor-rarity-a} the exponent is read at the threshold coupling, as
$\mathsf{B}(\gamma_*)N_C^4$, and this reading is part of the step marked $(\ast)$. With it, the two
preceding bounds give \eqref{8.12:eq-depth-floor-rarity-2}.

By the bound of Notation~\ref{8.12:not-volume-ranges} on the degree of $\mathsf{B}$, we have
$\mathsf{B}(\gamma_*)\le C_{\mathsf{B}}(1+\log X_*)$. Under the first line of
\eqref{8.12:eq-depth-floor-rarity-a} it follows that
\begin{equation*}
\exp\Bigl(\mathsf{B}(\gamma_*)N_C^4-9\ln L\,\bigl(s-m_{\mathbb{T}}(g)+k^{\dagger}-K\bigr)\Bigr)\le1.
\end{equation*}
Write $e^{-\theta(p_0(g_{k_0})-1)}=e^{\theta}e^{-\theta p_0(g_{k_0})}$ and
$e^{-9\mathsf{a}\ln L}=L^{-9\mathsf{a}}$. By step (2),
\begin{equation*}
\tilde{\nu}\bigl(\mathcal{E}^{\dagger}_{\square}\cap\{\operatorname{age}(C)=\mathsf{a}\}\bigr)
\le\tilde{\nu}(\mathcal{S}_{\square,\mathsf{a}}),
\end{equation*}
and the right-hand side of \eqref{8.12:eq-depth-floor-rarity-2} does not depend on the cell $\square$; hence
it bounds $q^{\mathcal{E},=\mathsf{a}}_{k_0}$. Using the geometric sum
\begin{equation*}
\sum_{\mathsf{a}>\mathfrak{l}}L^{8(\mathsf{a}+1)}L^{-9\mathsf{a}}
=L^{8}\sum_{\mathsf{a}\ge\mathfrak{l}+1}L^{-\mathsf{a}}=\frac{L^{8-\mathfrak{l}}}{L-1},
\end{equation*}
we obtain
\begin{equation*}
\begin{split}
\sum_{\mathsf{a}>\mathfrak{l}}L^{8(\mathsf{a}+1)}q^{\mathcal{E},=\mathsf{a}}_{k_0}
&\le e^{\theta}(L-1)^{-1}L^{8-\mathfrak{l}}e^{-\theta p_0(g_{k_0})}\\
&\quad\times\exp\Bigl(\mathsf{B}(\gamma_*)N_C^4-9\ln L\,\bigl(s-m_{\mathbb{T}}(g)+k^{\dagger}-K\bigr)\Bigr)\\
&\le e^{\theta}(L-1)^{-1}L^{8-\mathfrak{l}}e^{-\theta p_0(g_{k_0})}.
\end{split}
\end{equation*}

(5) \emph{The series.} Since $c_{\mathrm{lv}}\le\theta$,
\begin{equation*}
e^{-\theta p_0(g_{k_0})}\le e^{-c_{\mathrm{lv}}p_0(g_{k_0})}=q^{\mathrm{lv}}_{k_0}/C_{\mathrm{lv}}.
\end{equation*}
Moreover $L^{8-\mathfrak{l}}\le L^{8(\mathfrak{l}+1)}$. Hence the right-hand side of the last display of
step (4) is at most
$A^{\mathcal{E}}_{\mathrm{lv}}L^{8(\mathfrak{l}+1)}q^{\mathrm{lv}}_{k_0}$, which is
\eqref{8.12:eq-depth-floor-rarity-3}. No translation average and no moment is used, so the bound is uniform in
the position of $\square$. It is also uniform in $K$ and $g_0$.

\emph{The consequence.} The healed part of step (1) is bounded by the healed old-region bound. The end part is
bounded by \eqref{8.12:eq-depth-floor-rarity-3} for every $s$ satisfying the first line of
\eqref{8.12:eq-depth-floor-rarity-a}. Together these give \eqref{8.12:eq-old-region-moments-a} for
$s\ge s_W(g,m)$, and with it the two-point witness theorem. The floor depends on $m$ through $N_C(m)$, so the
uniformity does not extend to $m$.
\end{proof}

A variant splits the age series at a threshold age and bounds the terms below it by the age-blind end density
$\mathfrak{e}_{k_0}$. It gives a floor of the same form, slightly smaller; either floor may be used.

\begin{corollary}[The two-point witness in finite volume]\label{8.12:cor-finite-volume-witness}
Let $N=2$ and work within the regime of clause~(iii) of Theorem~\ref{3.1:thm-main}. The test class
$\mathfrak{T}(\mathsf{d};\vartheta)$ and the letters $b_0$, $\mathcal{K}_s$, $\zeta'$, $C_W(s)$,
$s(\mathsf{d})$, $k_W$, $c_{\mathcal{K}}$, $C_{\mathcal{K}}$ are exactly as in clause~(iv) of
Theorem~\ref{3.1:thm-main}.

Let $\gamma_*$ be the minimum of $\gamma_J$, of the ceilings on the coupling carried by the statement
that supplies the age-blind end density $\mathfrak{e}_{k_0}$ entering
Proposition~\ref{8.12:prop-depth-floor-rarity}, and of the ceiling
\eqref{8.12:eq-weak-torus-rarity-a}. The constant $\gamma_W$ is fixed last, subject to
\begin{equation*}
\gamma_W\le\min\bigl\{\gamma_G,\ \gamma^{\mathrm{add}}_{\mathrm{src}},\
\text{the threshold in \eqref{8.12:eq-weak-torus-rarity-b}}\bigr\},
\end{equation*}
where $\gamma^{\mathrm{add}}_{\mathrm{src}}$ is the ceiling on $g$ carried by the additive theorem in
(5) below, and $\gamma_G$ is a further ceiling on $g$ carried by the inputs listed below.

A triple $(g_0,g,m)$ of bare coupling, reference coupling and torus exponent, with $K=K(g_0,g)$, is
called \emph{weakly coupled} if the following holds: its first-passage
run, prolonged on the same lattice, reaches the minimal admissible torus before its running
coupling reaches $\gamma_*$. This is the first of the two cases of
\eqref{8.12:eq-volume-ranges-a}. A sufficient condition
(Remark~\ref{8.12:rem-weak-torus-extent}), with $g_V$, $b_+$ and $\lambda$ as in the statement
containing \eqref{8.12:eq-volume-ranges-a}--\eqref{8.12:eq-volume-ranges-c}, is
\begin{equation}\label{8.12:eq-finite-volume-witness-1}
m\le m_{\mathbb{T}}(g_V)
+\Bigl\lfloor\frac{g_V^{-2}-\gamma_*^{-2}-b_+-2c_2}{3\lambda}\Bigr\rfloor-1 .
\end{equation}

Let $s_W(g)$ be the depth floor of the two-point witness theorem. It is fixed by exactly the three
conditions
\begin{equation*}
s-r(s)\ge s_*,\qquad s\ge m_{\mathbb{T}}(g),\qquad 2B^{\sharp}\bigl(r(s)+1\bigr)\le\underline{x}_s .
\end{equation*}
Put
\begin{equation}\label{8.12:eq-finite-volume-witness-2}
s_W(g,m):=
\begin{cases}
s_W(g) & \text{if $(g_0,g,m)$ is weakly coupled},\\
\max\bigl\{s_W(g),\ \text{the least $s$ with \eqref{8.12:eq-depth-floor-rarity-a}}\bigr\}
& \text{otherwise}.
\end{cases}
\end{equation}
Consider the displayed one-sided conditions
\begin{align}
&g\le\gamma_W,\label{8.12:eq-finite-volume-witness-3}\\
&K-s(\mathsf{d})\ge k_W,\label{8.12:eq-finite-volume-witness-4}\\
&s_W(g,m)\le s(\mathsf{d}).\label{8.12:eq-finite-volume-witness-5}
\end{align}
Condition \eqref{8.12:eq-finite-volume-witness-5} has the following properties:
\begin{itemize}
\item for each $(g,m)$ it defines a half-line in $s$;
\item $s_W(g,m)$ is finite for every $(g,m)$ and non-decreasing in $m$;
\item it reads $g$ through $\log g^{-2}$ and $m$ through $N_C(m)^4$, where $N_C(m)$ is the number of
sites per side of the terminal lattice of the prolonged run;
\item it is not a cap in the sense of Definition~\ref{2.3:def-cap};
\item it imposes on $m$ nothing beyond the torus clause of (H5).
\end{itemize}
Let $\mathsf{d}_W(g,m)$ be defined from $s_W(g,m)$ in the same way as clause~(iv) of
Theorem~\ref{3.1:thm-main} defines $\mathsf{d}_W(g)$ from $s_W(g)$.

Assume the following inputs:
\begin{itemize}
\item[(1)] the correlation theorem, applied with $m:=s$;
\item[(2)] the Gaussian extraction in all three parts, together with localisation of the sourced
step, from which its first and third parts follow with the one-point weights of~(6);
\item[(3)] the source lemma and the continuum kernel;
\item[(4)] clause~(0) and clause~(ii-a) of Theorem~\ref{3.1:thm-main}, the operator theorem at free
current, and the three further inputs listed with the two-point witness theorem beside these;
\item[(5)] the additive theorem of the witness construction. At the observation level
$k_0:=K-s(\mathsf{d})$ its two inputs are the bound on young live regions and the age-series bound
\eqref{8.12:eq-old-region-moments-a};
\item[(6)] the corollary on rarity of young large-field blocks, and the healed old-region bound;
\item[(7)] Proposition~\ref{8.12:prop-weak-torus-rarity}, and
Proposition~\ref{8.12:prop-depth-floor-rarity}, whose proof is incomplete at one step;
\item[(8)] clauses~(ii-$\mathbb{T}$-a) and (ii-$\mathbb{T}$-b) of Theorem~\ref{3.1:thm-main}, and
reflection positivity of the lattice measure.
\end{itemize}

Then the following holds for every $g$ with \eqref{8.12:eq-finite-volume-witness-3}, every
$g_0\in(0,g_-(g)]$ with $K=K(g_0,g)$, every torus exponent $m$ of the family fixed by the torus
clause of (H5), every $\mathsf{d}$ with
\eqref{8.12:eq-finite-volume-witness-4} and \eqref{8.12:eq-finite-volume-witness-5}, with
$s:=s(\mathsf{d})$, and every $(f_1,f_2)\in\mathfrak{T}(\mathsf{d};\vartheta)$:
\begin{equation}\label{8.12:eq-finite-volume-witness-6}
S^T_{2;K,m}(f_1,f_2)=\zeta'^{\,2}_{K-s}\,b_0\,g_{K-s}^4\,\mathcal{K}_s(f_1,f_2)
\bigl(1+R_K(f_1,f_2)\bigr).
\end{equation}
Here $|R_K|\le\frac12$ (indeed $|R_K|\le\frac14$) and $\zeta'_{K-s}\in[\frac23,\frac43]$. Moreover
\begin{equation*}
c_{\mathcal{K}}\|f_1\|_\infty\|f_2\|_\infty\le\mathcal{K}_s(f_1,f_2)
\le C_{\mathcal{K}}\|f_1\|_\infty\|f_2\|_\infty ,
\end{equation*}
and $|\zeta'_{K-s}-1|$ is at most a constant times $g_{K-s}$ times a fixed power of
$\log g_{K-s}^{-2}$. Hence
\begin{equation}\label{8.12:eq-finite-volume-witness-7}
\tfrac29\,b_0\,g_{K-s}^4\,\mathcal{K}_s(f_1,f_2)\le S^T_{2;K,m}(f_1,f_2)
\le\tfrac83\,b_0\,g_{K-s}^4\,\mathcal{K}_s(f_1,f_2).
\end{equation}
This holds uniformly in $K$, $g_0$ and the position of the supports, and uniformly in $m$ over the
weakly coupled members.

Consequences. Let $\ell_g$, $g_{\infty,s}$, $\mathsf{d}_W(g)$, the pairs
$(f_1^{(\mathsf{d})},f_2^{(\mathsf{d})})$ and $g_{\mathrm{phys}}(\mathsf{d})$ be as in clause~(iv)
of Theorem~\ref{3.1:thm-main}. Consider the subsequences of clause~(ii-$\mathbb{T}$-b) and the tuned
family of clause~(ii-$\mathbb{T}$-a) on every torus of side $2L^{m_0}\ell_g$. Along them every limit
point $S^T_{2;m_0}$ obeys
\begin{equation}\label{8.12:eq-finite-volume-witness-8}
\tfrac29\,b_0\,g_{\infty,s}^4\,\mathcal{K}_s(f_1,f_2)\le S^T_{2;m_0}(f_1,f_2)
\le\tfrac83\,b_0\,g_{\infty,s}^4\,\mathcal{K}_s(f_1,f_2).
\end{equation}
Restrict to separations $\mathsf{d}\le\mathsf{d}_W(g)$ on weakly coupled tori, and to
$\mathsf{d}\le\mathsf{d}_W(g,m_0)$ on the others. Let $f^{\mathrm{refl}}$ denote the reflection of a
test function $f$ in the sense of reflection positivity of the lattice measure. Then:
\begin{itemize}
\item[(a)] $S^T_{2;m_0}(f^{\mathrm{refl}},f)>0$ for every $f\neq0$ with
$(f^{\mathrm{refl}},f)\in\mathfrak{T}(\mathsf{d};\vartheta)$. The smeared curvature is therefore not
almost surely constant in any limit state on the torus.
\item[(b)] The function $\mathsf{d}\mapsto S^T_{2;m_0}(f_1^{(\mathsf{d})},f_2^{(\mathsf{d})})$
tends to $0$ and dominates every power of $\mathsf{d}$. No scale-covariant law, of any dimension,
reproduces it.
\item[(c)] No limit state on the torus is the law of a free Maxwell field at any charge.
\item[(d)] $g_{\mathrm{phys}}(\mathsf{d})>0$, and $g_{\mathrm{phys}}(\mathsf{d})$ tends to
$0$ logarithmically, at the rate of the two-sided bounds on the running coupling.
\end{itemize}

This corollary asserts nothing about the limits $S^T_{n;\infty}$ on $\mathbb{R}^4$ of clause~(ii)
of Theorem~\ref{3.1:thm-main}. In particular it does not assert
$\dim\mathcal{H}^{\mathrm{curv}}\ge2$, nor
(b), (c) or (d) on $\mathbb{R}^4$. The reason is that along $m_j\to\infty$ the admissible
separations $\mathsf{d}_W(g,m_j)$ decrease to $0$, so no fixed test pair remains admissible. The
corollary also asserts nothing that clause~(iv) of Theorem~\ref{3.1:thm-main} lists as not asserted.
\end{corollary}

\begin{proof}
\emph{The depth floor.} Let $(g_0,g,m)$ be weakly coupled. Then
$s_W(g,m)=s_W(g)$ is fixed by three conditions, none of which reads $m$.

Let $(g_0,g,m)$ not be weakly coupled. In the condition \eqref{8.12:eq-depth-floor-rarity-a}, the
left member $9\ln L\,(s-m_{\mathbb{T}}(g)+k^{\dagger}-K)$ grows by $9\ln L>0$ with each unit of $s$.
The right member $C_{\mathsf{B}}(1+\log X_*)N_C(m)^4$ does not contain $s$. So the set of $s$
satisfying \eqref{8.12:eq-depth-floor-rarity-a} is a half-line, and its least element is finite.
The maximum in \eqref{8.12:eq-finite-volume-witness-2} is then finite as well, and
\eqref{8.12:eq-finite-volume-witness-5} is a half-line in $s$.

Fix $(g_0,g)$. By \eqref{8.12:eq-volume-ranges-a} and
Remark~\ref{8.12:rem-weak-torus-extent}, the weakly coupled members are those whose torus exponent
lies below a threshold. On them the floor equals $s_W(g)$. Above the threshold the floor is a
maximum with $s_W(g)$, and its second entry reads $m$ only through $N_C(m)^4$. By the form of
$N_C(m)$ fixed with \eqref{8.12:eq-depth-floor-rarity-a}, $N_C(m)$ grows with the torus excess,
while the left member of \eqref{8.12:eq-depth-floor-rarity-a} does not read $m$; so the least $s$
satisfying \eqref{8.12:eq-depth-floor-rarity-a} is non-decreasing in $m$. Since
$s_W(g,m)\ge s_W(g)$ in both cases and the weakly coupled members are those below the threshold in
$m$, $s_W(g,m)$ is non-decreasing in $m$.

\emph{Weakly coupled members.} Fix $g$ with \eqref{8.12:eq-finite-volume-witness-3}, $g_0$, $K$, a
weakly coupled $m$, and $\mathsf{d}$ with \eqref{8.12:eq-finite-volume-witness-4} and
\eqref{8.12:eq-finite-volume-witness-5}. Put $s=s(\mathsf{d})$ and $k_0=K-s$. By
\eqref{8.12:eq-finite-volume-witness-5}, $s\ge s_W(g)$; the second condition defining $s_W(g)$ then
gives $s\ge m_{\mathbb{T}}(g)$, that is, $k_0\le K-m_{\mathbb{T}}(g)$.

The ceilings \eqref{8.12:eq-weak-torus-rarity-a} and \eqref{8.12:eq-weak-torus-rarity-b} hold. The
first holds because $\gamma_*$ does not exceed it; the second holds because
$g\le\gamma_W$ and $\gamma_W$ does not exceed it. Corollary~\ref{8.12:cor-weak-torus-witness}
derives, from the healed old-region bound and Proposition~\ref{8.12:prop-weak-torus-rarity} (both
in (6), (7)), the age-series bound \eqref{8.12:eq-old-region-moments-a}. It holds on every weakly
coupled member at every $k_0\le K-m_{\mathbb{T}}(g)$, in particular at $k_0=K-s$.

By the same corollary, this bound, together with the bound on young live regions from the corollary
on rarity of young large-field blocks, gives the two inputs of the additive theorem in (5). The
additive theorem then holds at $k_0$, since $g\le\gamma^{\mathrm{add}}_{\mathrm{src}}$. Since
\eqref{8.12:eq-finite-volume-witness-3}, \eqref{8.12:eq-finite-volume-witness-4} and
$s\ge s_W(g)$ are the three floors of the two-point witness theorem, that theorem, with inputs
(1)--(4), yields \eqref{8.12:eq-finite-volume-witness-6}
with $|R_K|\le\frac14$.

\emph{Members that are not weakly coupled.} Now \eqref{8.12:eq-finite-volume-witness-5} gives both
$s\ge s_W(g)$ and \eqref{8.12:eq-depth-floor-rarity-a}. Hence again $k_0=K-s\le K-m_{\mathbb{T}}(g)$.

Under \eqref{8.12:eq-depth-floor-rarity-a}, Proposition~\ref{8.12:prop-depth-floor-rarity}, whose
proof is incomplete at one step, supplies the end part of the age series at $k_0$. This is the part
contributed by old live regions that are never healed. The healed old-region bound supplies the
healed part. The age-series bound \eqref{8.12:eq-old-region-moments-a} is obtained by adding the two
parts; its constant is the sum of the two constants.

The corollary on rarity of young large-field blocks supplies the bound on young live regions. The
additive theorem therefore holds at $k_0$. As in the weakly coupled case, the three floors of the
two-point witness theorem are met, and that theorem gives \eqref{8.12:eq-finite-volume-witness-6}
with $|R_K|\le\frac14$.

\emph{The display and the two-sided law.} In both cases the two-point witness theorem also gives
the following:
\begin{itemize}
\item $\zeta'_{K-s}\in[\frac23,\frac43]$;
\item the two bounds on $\mathcal{K}_s$;
\item the bound on $|\zeta'_{K-s}-1|$ by a constant times $g_{K-s}$ times a fixed power of
$\log g_{K-s}^{-2}$.
\end{itemize}
In particular $|R_K|\le\frac14\le\frac12$, and so $1+R_K\in[\frac12,\frac32]$. From
$\zeta'_{K-s}\in[\frac23,\frac43]$ we get $\zeta'^{\,2}_{K-s}\in[\frac49,\frac{16}9]$. Also
$b_0=(N^2-1)/2=\frac32>0$, and
$\mathcal{K}_s\ge c_{\mathcal{K}}\|f_1\|_\infty\|f_2\|_\infty\ge0$. Multiplying these ranges in
\eqref{8.12:eq-finite-volume-witness-6} gives the factors
\begin{equation*}
\tfrac49\cdot\tfrac12=\tfrac29,\qquad\tfrac{16}9\cdot\tfrac32=\tfrac83 ,
\end{equation*}
which is \eqref{8.12:eq-finite-volume-witness-7}.

None of the constants entering \eqref{8.12:eq-finite-volume-witness-7} depends on $K$, on $g_0$ or
on the position of the supports. On weakly coupled members the hypotheses
\eqref{8.12:eq-finite-volume-witness-3}--\eqref{8.12:eq-finite-volume-witness-5} do not read $m$,
because the floor there is $s_W(g)$. Moreover, by Corollary~\ref{8.12:cor-weak-torus-witness} the
age-series bound holds on weakly coupled members with the constant
\begin{equation*}
A_{\mathrm{lv}}=A^{\mathrm{heal}}_{\mathrm{lv}}+\frac{e^{\theta}}{(L-1)C_{\mathrm{lv}}},
\end{equation*}
which does not depend on $m$; hence the inputs of the additive theorem, and so $|R_K|\le\frac14$,
hold there with constants free of $m$. The letters are those of clause~(iv) of
Theorem~\ref{3.1:thm-main}. This gives
the uniformity in $m$ over the weakly coupled members. Beyond them the floor reads $m$.

\emph{Consequences.} These are drawn from \eqref{8.12:eq-finite-volume-witness-7} together with
clause~(0) of Theorem~\ref{3.1:thm-main}, clauses~(ii-$\mathbb{T}$-a) and (ii-$\mathbb{T}$-b),
reflection positivity of the lattice measure, and the continuum kernel.

The bounds \eqref{8.12:eq-finite-volume-witness-7} are uniform in $K$, $g_0$ and the position of the
supports. They therefore pass to every limit point along the subsequences of
clause~(ii-$\mathbb{T}$-b) and along the tuned family of clause~(ii-$\mathbb{T}$-a), with
$g_{\infty,s}$ in place of $g_{K-s}$. This is \eqref{8.12:eq-finite-volume-witness-8}.

For (a), take $f\neq0$ with $(f^{\mathrm{refl}},f)\in\mathfrak{T}(\mathsf{d};\vartheta)$. Since
$\|f^{\mathrm{refl}}\|_\infty=\|f\|_\infty$, we have
$\mathcal{K}_s(f^{\mathrm{refl}},f)\ge c_{\mathcal{K}}\|f\|_\infty^2>0$. Also $b_0>0$, and
$g_{\infty,s}>0$ as in (d). So the left member of \eqref{8.12:eq-finite-volume-witness-8} is
positive, and $S^T_{2;m_0}(f^{\mathrm{refl}},f)>0$. A connected two-point function vanishes if one
of its two observables is almost surely constant. Hence the smeared curvature is not almost surely
constant.

For (b), insert the pairs $(f_1^{(\mathsf{d})},f_2^{(\mathsf{d})})$ in
\eqref{8.12:eq-finite-volume-witness-8}, and control $\mathcal{K}_s$ along them by the continuum
kernel. Then $S^T_{2;m_0}$ is bounded above and below by multiples of $g^4_{\infty,s(\mathsf{d})}$,
and it tends to $0$ with the logarithmic rate of (d). It therefore dominates every power of
$\mathsf{d}$. A scale-covariant law is a constant times a power $\mathsf{d}^{\delta}$:
\begin{itemize}
\item for $\delta\le0$ it does not tend to $0$;
\item for $\delta>0$ it is eventually dominated by our function.
\end{itemize}
So no such law reproduces it.

For (c), as $K\to\infty$ the depth $s$ of admissible separations is unbounded above. Hence on every
torus there are two admissible separations at which the values of $g^4_{\infty,s}$, say
$g_1^4>12\,g_2^4$, have ratio larger than $12=\frac83\big/\frac29$. Then
\eqref{8.12:eq-finite-volume-witness-8} gives, for the quotient $S^T_{2;m_0}/(b_0\mathcal{K}_s)$ at
the two separations,
\begin{equation*}
\text{quotient at the first}\ \ge\tfrac29\,g_1^4>\tfrac83\,g_2^4\ \ge\ \text{quotient at the second}.
\end{equation*}
So the quotient depends on the separation. This is excluded for a free Maxwell field at any charge
in the sense of Definition~\ref{3.1:def-non-trivial}.

For (d), the positivity of $g_{\mathrm{phys}}(\mathsf{d})$ and its logarithmic decay at the two-sided
rate are the two-sided bounds on the running coupling, read at separation $\mathsf{d}$.
\end{proof}

\begin{proposition}[Volume-uniform rarity of never-healed old regions. Proof outstanding]
\label{8.12:prop-uniform-old-rarity}
Let $s$ be a witness depth, $k_0=K-s$ the observation level, $\tilde{\nu}$ the history law at
$k_0$, $\mathfrak{l}$ the age threshold and $q^{\mathrm{lv}}_{k_0}$ the young-live weight, all as in
Proposition~\ref{8.12:prop-old-region-moments}. For an integer $\mathsf{a}>\mathfrak{l}$, let
$q^{=\mathsf{a}}_{k_0}$ be the maximum over cells $\square$ of the $\tilde{\nu}$-probability that
$\square$ lies in a live component at level $k_0$ whose oldest layer was created at level
$k_0-\mathsf{a}$. Split it by the fate of that component,
\begin{equation*}
q^{=\mathsf{a}}_{k_0}=q^{\mathrm{heal},=\mathsf{a}}_{k_0}+q^{\mathcal{E},=\mathsf{a}}_{k_0},
\end{equation*}
where $q^{\mathcal{E},=\mathsf{a}}_{k_0}$ is the part on the live event, on which the component is
never healed,
that is, its lineage is still live at the terminal level $K^{\dagger\prime}$ of the prolonged run
of \eqref{8.12:eq-volume-ranges-a}, in either of the two volume ranges there, Case~V and Case~C.
Then there is a constant $A^{\mathcal{E}}_{\mathrm{lv}}$, independent of $m$ (that is, of the
volume), of $K$, of $g_0$, of the position and of the level $k_0$, such that on every torus
$\mathbb{T}_m$ of the family of tori of Theorem~\ref{3.1:thm-main}, for $g\le\gamma_W$ and for
every witness depth $s\ge s_W(g)$, with no volume-dependent depth floor imposed on the witness
depth,
\begin{equation}\label{8.12:eq-uniform-old-rarity-1}
\sum_{\mathsf{a}>\mathfrak{l}} L^{8(\mathsf{a}+1)}\, q^{\mathcal{E},=\mathsf{a}}_{k_0}
\le A^{\mathcal{E}}_{\mathrm{lv}}\, L^{8(\mathfrak{l}+1)}\, q^{\mathrm{lv}}_{k_0}.
\end{equation}
In particular, in Case~C the constant $A^{\mathcal{E}}_{\mathrm{lv}}$ is uniform in the number
$N_C$ of sites per side of the terminal lattice of the prolonged run.
\end{proposition}

Proof outstanding.

\noindent\emph{Route.}
By the split above, the age-series bound \eqref{8.12:eq-old-region-moments-a} is obtained by
adding \eqref{8.12:eq-uniform-old-rarity-1} to the healed old-region bound, which is the same
inequality for the healed parts $q^{\mathrm{heal},=\mathsf{a}}_{k_0}$ with a constant
$A^{\mathrm{heal}}_{\mathrm{lv}}$ in place of $A^{\mathcal{E}}_{\mathrm{lv}}$ and which holds at
every level with no depth floor, uniformly in $K$ and $m$; only
\eqref{8.12:eq-uniform-old-rarity-1} is the content of the present proposition. By
Proposition~\ref{8.12:prop-old-region-moments}, old live regions enter the two-point witness
only through first and second slot moments, with weights $L^{4\mathsf{a}}$ and
$L^{8\mathsf{a}}$. In Case~V the never-healed parts $q^{\mathcal{E},=\mathsf{a}}_{k_0}$ are
bounded without a volume-dependent depth floor
by the rarity bound \eqref{8.12:eq-weak-torus-rarity-c} on weakly coupled tori; what remains is
the range Case~C, uniformly in $N_C$. For each fixed torus, Proposition~\ref{8.12:prop-depth-floor-rarity}, whose
proof is incomplete at one step, gives \eqref{8.12:eq-uniform-old-rarity-1} only under the
volume-dependent depth floor \eqref{8.12:eq-depth-floor-rarity-a}. That floor reads the volume
through $N_C^4$. A proof of the present proposition must therefore reach
\eqref{8.12:eq-uniform-old-rarity-1} without that floor. Two reductions are available.

The first reduction: the bound follows from the annealed live-tail statement, in a modified form,
in its one-cell instance, the instance in which the region on
which the statement is taken is the whole domain. The annealed live-tail statement is not proved
and is not an input of Theorem~\ref{3.1:thm-main}. That
instance is reduced to two inputs: the input treating unbounded unions (the
unbounded-union depth statement), and a further bound,
itself reduced to one estimate whose leading input is the half-skin bound at the canonical comparison point.
What must be established on this route is therefore the half-skin bound at
the canonical comparison point, together with the case of unbounded unions.

The second reduction: the bound follows from a statement on the total never-healed part at one
fixed level, which is in turn reduced to a single estimate at that same level.

\begin{remark}[Relation to clause~\textup{(iv)}]
The volume-uniform form of the finite-volume witness of
Corollary~\ref{8.12:cor-finite-volume-witness} is that corollary with the depth floor $s_W(g,m)$
equal to $s_W(g)$ for every $m$. Along the route of
Proposition~\ref{8.12:prop-old-region-moments}, on which old live regions enter the witness only
through the age series, this form is obtained exactly when \eqref{8.12:eq-old-region-moments-a}
holds uniformly in the volume, that is, exactly when \eqref{8.12:eq-uniform-old-rarity-1} holds
in Case~C uniformly in $N_C$. On the two reductions above this comes down to a single one-cell
statement on the whole domain: on the first reduction, the half-skin bound at the canonical
comparison point together with the case of unbounded unions; on the second, the single estimate at
one fixed level.
Clause~(iv) of Theorem~\ref{3.1:thm-main} does not rest on the present proposition. In finite
volume it is Corollary~\ref{8.12:cor-finite-volume-witness}, in which the never-healed old parts
$q^{\mathcal{E},=\mathsf{a}}_{k_0}$ are controlled, under the per-torus depth floor
\eqref{8.12:eq-depth-floor-rarity-a}, by the rarity bound \eqref{8.12:eq-weak-torus-rarity-c} in
Case~V and by Proposition~\ref{8.12:prop-depth-floor-rarity}, whose proof is incomplete at one
step, in Case~C; the volume dependence of
that floor is removed by the volume-transfer theorem of this chapter, and not by
\eqref{8.12:eq-uniform-old-rarity-1}. The present proposition, like the annealed live-tail
statement, is not an input of Theorem~\ref{3.1:thm-main}.
\end{remark}

\begin{remark}[Infinite-volume bounds along the age series need uniform old-region rarity]
\label{8.12:rem-infinite-volume-dependence}
We use the notions (n0)--(n3) of non-triviality for the curvature two-point function fixed in
Remark~\ref{8.12:rem-nontriviality-notions}, the decomposition
\eqref{8.12:eq-two-point-decomposition-a} of the connected two-point function into the main term
$\mathfrak{M}$ and six remainders, the old-region moments \eqref{8.12:eq-old-region-moments-a} and
\eqref{8.12:eq-old-region-moments-b}, and the two ranges of volumes, Case~V and Case~C, of
\eqref{8.12:eq-volume-ranges-a}. Write $\mathsf{a}$ for the age of a live large-field region at the
observation level $k_0=K-s$. Along the route of Proposition~\ref{8.12:prop-old-region-moments}, on
which old live regions enter only through the age series, the following hold on $\mathbb{R}^4$.
Clause~(iv) of Theorem~\ref{3.1:thm-main} does not take this route: it is obtained by the volume
transfer of Corollary~\ref{8.22:cor-transfer-depth-floor}, and its consequences for the limits on
$\mathbb{R}^4$, $\dim\mathcal{H}^{\mathrm{curv}}\ge2$ among them, are drawn after
Corollary~\ref{8.22:cor-reference-torus-kernel}.
\begin{itemize}
\item[(a)] The upper bound (n1),
\begin{equation*}
S^T_{2;m}(f_1,f_2)\le C\,g(\mathsf{d})^4\,\mathcal{K}_{s(\mathsf{d})}(f_1,f_2)
\quad\text{for all small }\mathsf{d},\ \text{with }C\text{ independent of }m,
\end{equation*}
where $g(\mathsf{d})$ is the running coupling at the physical scale $\mathsf{d}$, depends, through
the second-moment members of \eqref{8.12:eq-old-region-moments-b} (the covariance of the two live
source slots and the conditional-covariance term), on the old-age tail statement in Case~C in its
form uniform in the volume (Proposition~\ref{8.12:prop-uniform-old-rarity}, whose proof is
outstanding).
\item[(b)] The strict positivity (n3), namely that $S^T_{2;\infty}$, evaluated at one test
function $f$ and its reflection in a link plane, is strictly positive, depends on a lower bound,
uniform in the volume, at one scale, on these members; reflection positivity gives only $\ge 0$.
\item[(c)] The two-sided notions (n2$\flat$) and (n2) on $\mathbb{R}^4$ depend on both sides of
these members. The notion (n2$\sharp$) is not asserted in this book.
\item[(d)] Independent of the never-healed old-age tail, the theory on $\mathbb{R}^4$ has (n0),
the non-negativity given by reflection positivity, and the existence and Osterwalder--Schrader
content of clause~(ii) of Theorem~\ref{3.1:thm-main} with the space $\mathcal{H}^{\mathrm{curv}}$.
\item[(e)] No exchange or iteration of the limits $K\to\infty$, $m\to\infty$ removes the
dependence in (a) and (b).
\end{itemize}
\end{remark}

\begin{proof}
\emph{(a).} By Remark~\ref{8.12:rem-two-sided-dependence}, the upper side of
$|S^T_{2;K,m}-\mathfrak{M}|$ contains the covariance
$\operatorname{Cov}_{\tilde{\nu}}(\mathfrak{L}^h_1,\mathfrak{L}^h_2)$ of the live parts of the
two conditional means, and the term
$\mathbb{E}_{\tilde{\nu}}\operatorname{Cov}_{\mathbb{E}_h}(\partial_1\mathfrak{A}^w_h,
\partial_2\mathfrak{A}^w_h)$ coming from the conditional covariance part $\Psi^h$, where
$\partial_i$ denotes the derivative in the source coefficient $z_i$ of $w=z_1f_1+z_2f_2$, as in
\eqref{8.12:eq-old-region-moments-b}. The
second-moment member of these terms in the additive theorem carries the weight
$L^{8\mathsf{a}}$ and has
no sign. A bound on it uniform in the volume is the moment bound
\eqref{8.12:eq-old-region-moments-a} uniform in the volume, and this contains the old-age tail
statement in Case~C; this is the content of
Proposition~\ref{8.12:prop-uniform-old-rarity}, whose proof is outstanding.

No global device supplies this bound. The global stability quotient
$e^{\mathsf{B}N_T^4}$ of Remark~\ref{8.12:rem-global-quotient}, $N_T$ the number of sites per side
of the terminal lattice of the run, gives no bound uniform in $N_T$; and the translation-invariant
model of that remark, which satisfies every input of the global device and contains one old end
cell whose weighted end probability is unbounded as $N_T\to\infty$, shows the same for the
second-moment member when it is taken with a set of two cells in place of one. Moreover, by the
scale-freeness of the end density per unit physical volume, at the creation level the depth gain
$L^{-4\mathsf{a}}$ of the end density is cancelled exactly by the placement entropy
$L^{4\mathsf{a}}$. A crude bound by the
supremum of the integrand times the $\tilde{\nu}$-mass of the old event is the same rarity
statement, by Remark~\ref{8.12:rem-crude-bound}. Finally, neither the correlation theorem nor the
per-ball volume-free bound carries the factor $g^4$, so they cannot give (n1).
The first-moment members of the additive theorem, of weight $L^{4\mathsf{a}}$, are not bounded
uniformly in the volume here either; and the one available route for them, re-measuring the
birth-scale size $\Theta_Z$ of a live component $Z$ by its own-scale cell count, would shrink but
not remove the end member, in which the product of the two slots remains.

\emph{(b).} Strict positivity at one scale requires a lower bound, uniform in the volume, at that
scale on the lower side of the same members, by Remarks~\ref{8.12:rem-two-sided-dependence}
and~\ref{8.12:rem-positivity-one-scale}. Reflection positivity of the lattice measure gives only
the non-negativity of $S^T_{2;K,m}$ at every test function and its reflection in a link plane.
No volume monotonicity and no correlation inequality is known
for $\mathrm{SU}(N)$ beyond reflection positivity and the chessboard estimate. Therefore strict
positivity on a torus does not pass to the limit $m\to\infty$; nor can semicontinuity help, since
a limit of positive numbers may vanish. Hence, along this route, even the non-degeneracy (n3) of
the theory on $\mathbb{R}^4$ depends on the old-age tail statement restricted to Case~C.

\emph{(c).} The notions (n2$\flat$) and (n2) contain both an upper and a lower bound of the kind
treated in (a) and (b), so by the two preceding paragraphs they depend on both sides of the
members listed in (a).

\emph{(d).} None of the items listed in (d) passes through the members of (a). What they provide
on $\mathbb{R}^4$ is a reconstructed Hilbert space with vacuum $\Omega$ and Hamiltonian $H\ge0$,
whose dimension these items alone do not show to exceed one, together with a
renormalisation-group trajectory
along which the coupling is non-zero and controlled from both sides at every scale.

\emph{(e).} The change-of-units identity, clause~(v-d)(iii) of Theorem~\ref{3.1:thm-main},
identifies the datum
$(m_0,\hat{x},\mathsf{d})$ with $(m_0+s,\hat{x}_s,L^{-s}\mathsf{d})$ on the same lattice, with
$\hat{x}_s$ the canonical comparison point at depth $s$ of that clause. Case~V is
a property of the bare datum and is invariant under this identification, so the identity cannot
move a statement from Case~C into Case~V. If instead one lets $m_j\to\infty$ with $g_j\downarrow0$
so as to remain in Case~V, the physical side length of the torus stays below $\ell_*$, the
physical length at which the running coupling reaches $\gamma_*$; this is not an infinite-volume
limit. Iterated limits, $K\to\infty$ on each torus followed by $m\to\infty$ by the compactness that
the per-ball volume-free bound provides, meet the same obstruction: the window
$\mathsf{d}_W(g,m)$ below which the bounds of the two-point witness theorem hold on
$\mathbb{T}_m$ at coupling $g$ (the scale of depth $s_W(g,m)$ in the notation of
Proposition~\ref{8.12:prop-uniform-old-rarity}, whose proof is outstanding) satisfies
$\mathsf{d}_W(g,m)\downarrow0$ as $m\to\infty$.
\end{proof}

\begin{definition}[Model: the two-tier box-clamp lattice model]\label{8.13:def-two-tier-model}
The model is a real scalar field on a periodic two-dimensional strip. The values on the bottom and top rows are prescribed. The field minimises the nearest-neighbour Dirichlet energy subject to bounds, imposed bond by bond, on the differences of the field across the bonds lying in two bands of rows at the bottom and, by reflection, two bands at the top: a lower tier with bound $\tau$ and a young tier above it with bound $c$. Equivalently, it is the model with the lower tier alone (the straddling clamp, defined in (a) below) to which one further band of bond-by-bond bounds, invariant under lateral translations, has been added. The number bounding the modulus of the difference across a bond is called the \emph{cap} of that bond. This use of the word is local to the model and is unrelated to Definition~\ref{2.3:def-cap}.

\emph{The strip.} Let $P\ge 8$ be an even integer, $N:=2P$, and let $H$ be a positive integer. Put
\begin{equation*}
\Lambda:=\mathbb{Z}_N\times\{0,1,\dots,H\}.
\end{equation*}
For $x\in\mathbb{Z}_N$ let $\|x\|:=\operatorname{dist}(x,N\mathbb{Z})$. The bonds of $\Lambda$ are of two kinds:
\begin{itemize}
\item the horizontal bonds $\{(x,y),(x+1,y)\}$, for $x\in\mathbb{Z}_N$ and $0\le y\le H$;
\item the vertical bonds $\{(x,y),(x,y+1)\}$, for $0\le y\le H-1$.
\end{itemize}
The vertical bond $\{(x,y),(x,y+1)\}$ is called the vertical bond $y$ of column $x$, or the vertical bond $(x,y)$. For $\psi:\Lambda\to\mathbb{R}$ we write
\begin{equation*}
F_y(x):=\psi(x,y+1)-\psi(x,y),
\end{equation*}
and
\begin{equation*}
\Delta_x\psi(x,y):=\psi(x+1,y)+\psi(x-1,y)-2\psi(x,y).
\end{equation*}
The energy of $\psi$ is
\begin{equation*}
E(\psi):=\tfrac12\sum_{\text{nearest-neighbour bonds}}(\nabla\psi)^2,
\end{equation*}
where $\nabla\psi$ on a bond is the difference of $\psi$ across it.

\emph{The datum.} Fix $\sigma>0$ and put $\varphi_0(x):=\sigma\|x\|$. The graph of $\varphi_0$ is a lattice V, with bottom tip at $x\equiv0$ and top tip at $x\equiv P$. Since $\|x\|=x$ and $\|x+P\|=P-x$ for $0\le x\le P$, one has $|\nabla\varphi_0|=\sigma$ on every horizontal bond, and
\begin{equation}\label{8.13:eq-two-tier-model-1}
\begin{gathered}
\Delta_x\varphi_0=2\sigma\,\mathbf{1}_{x\equiv0}-2\sigma\,\mathbf{1}_{x\equiv P},\\
\varphi_0(-x)=\varphi_0(x),\qquad\varphi_0(x+N)=\varphi_0(x),\qquad
\varphi_0(x+P)=\sigma P-\varphi_0(x).
\end{gathered}
\end{equation}
The datum may be replaced by a \emph{mollified datum}. In this mollified variant $\varphi_0$ denotes the mollified datum, and every occurrence of $\varphi_0$ in what follows refers to it; the role of the tip column $\{0\}$ is played by the tip set
\begin{equation*}
\mathcal{T}:=\{x:\ \|x\|\le t_0,\ \Delta_x\varphi_0(x)>0\},
\end{equation*}
where $t_0$ is the half-width of the mollification, a fixed parameter; $\mathcal{T}$ is the set of columns within distance $t_0$ of the bottom tip at which the lateral second difference of the mollified datum is positive, and it consists of a fixed number $\ell_{\mathrm{mn}}$ of columns. The role of the top tip column $\{P\}$ is played by the image of $\mathcal{T}$ at $P$. The statements made about the model in what follows are worded so that, for the mollified datum, the only change is this substitution.

\emph{Dirichlet rows.} We require
\begin{equation*}
\psi(\cdot,0)=\psi(\cdot,H)=\varphi_0.
\end{equation*}

\emph{Constraints.} Fix $\tau$ and $c$ with $0<\sigma<\tau<c$, a positive integer $W$, and an integer $\bar W\ge4$. Put $W':=W+\bar W$.
\begin{itemize}
\item[(a)] \emph{Lower tier} (the straddling clamp). We require
\begin{equation*}
|F_y(x)|\le\tau\quad\text{for }0\le y\le W-1,
\end{equation*}
and
\begin{equation*}
|\psi(x+1,y)-\psi(x,y)|\le\tau\quad\text{on the rows }0\le y\le W.
\end{equation*}
\item[(b)] \emph{Young tier} (the clamp of the first held young row, with cap $c$). We require
\begin{equation*}
|F_y(x)|\le c\quad\text{for }W\le y\le W'-1,
\end{equation*}
and
\begin{equation*}
|\psi(x+1,y)-\psi(x,y)|\le c\quad\text{on the rows }W+1\le y\le W'.
\end{equation*}
\item[(c)] \emph{Mirror constraints.} On the rows $H-W'\le y\le H$ we impose the images of the constraints (a) and (b) under the reflection $y\mapsto H-y$ of $\Lambda$.
\item[(d)] All other bonds are \emph{free}, that is, they carry no constraint.
\end{itemize}
The caps so defined are invariant under translations in $x$ and under $y\mapsto H-y$.

\emph{Minimisers.} Let $\mathcal{K}_c$ be the feasible set, that is, the set of $\psi:\Lambda\to\mathbb{R}$ that satisfy the Dirichlet rows and the constraints (a)--(c). The energy $E$ has a unique minimiser over $\mathcal{K}_c$, by the existence and uniqueness lemma for minimisers of $E$ under caps that act bond by bond and are invariant under translations in $x$, applied to this family of caps; we denote it $\psi^{(c)}$. We put
\begin{equation*}
\tilde\psi:=\psi^{(\infty)},\qquad \hat\psi:=\psi^{(\tau')},
\end{equation*}
where $\tau'$ is a fixed value of the young cap $c$ (so that $\tau'>\tau$), namely the value attached to the comparison point $\hat x$. Thus $\tilde\psi$ is the minimiser of the model with no young tier, that is, with the constraints (b) and their mirror images removed, the model attached to the comparison point $\tilde x$; and $\hat\psi$ is the minimiser of the model attached to the comparison point $\hat x$.

\emph{Notation at a minimiser.} Let $\psi=\psi^{(c)}$.
\begin{itemize}
\item The \emph{exit trace} of the lower tier is $g(x):=\psi(x,W)$.
\item The \emph{lateral mean} is $A_y(x):=\tfrac12[\psi(x-1,y)+\psi(x+1,y)]$, so that
\begin{equation*}
\Delta_x\psi(x,y)=2[A_y(x)-\psi(x,y)].
\end{equation*}
\item The \emph{lower lift deficit} is defined for $0\le y\le W$ by
\begin{equation}\label{8.13:eq-two-tier-model-2}
f(x,y):=\psi(x,y)-\varphi_0(x)-\tau y\le0.
\end{equation}
The inequality follows from feasibility and the Dirichlet row. Indeed, $\psi(x,y)=\varphi_0(x)+\sum_{0\le y'<y}F_{y'}(x)$, and each summand is at most $\tau$ by (a).
\item The \emph{young lift deficit} is defined for $W\le y\le W'$ by
\begin{equation}\label{8.13:eq-two-tier-model-3}
f^{y}(x,y):=\psi(x,y)-g(x)-c(y-W)\le0.
\end{equation}
Here the superscript is a name and not the row variable. The inequality follows from feasibility and row $W$. Indeed, $\psi(x,y)=g(x)+\sum_{W\le y'<y}F_{y'}(x)$, and each summand is at most $c$ by (b).
\item For a vertical bond $(x,y)$ in either tier, let $\lambda^{+}(x,y),\lambda^{-}(x,y)\ge0$ be the Karush--Kuhn--Tucker multipliers of the constraints $F_y(x)\le\text{cap}$ and $-F_y(x)\le\text{cap}$. The multiplier of the bond is
\begin{equation*}
\mu(x,y):=\lambda^{+}(x,y)-\lambda^{-}(x,y),
\end{equation*}
and $\mu:=0$ on free bonds.
\item A vertical bond of a tier is \emph{up-tight} if $F=+\text{cap}$ and \emph{down-tight} if $F=-\text{cap}$, where the cap is $\tau$ in the lower tier and $c$ in the young tier.
\item An \emph{up-run} on column $x$ in a tier is a maximal set of consecutive up-tight bonds of that tier in column $x$.
\item A column $x$ is \emph{lower-full} if $F_y(x)=\tau$ for all $0\le y\le W-1$. Then $g(x)=\varphi_0(x)+\tau W$.
\item A column $x$ is \emph{young-full} if $F_y(x)=c$ for all $W\le y\le W'-1$.
\end{itemize}
\end{definition}

\begin{remark}[Model: dictionary between the model and the block variables]\label{8.13:rem-model-dictionary}
The model is the two-tier box-clamp lattice model of Definition~\ref{8.13:def-two-tier-model}.
It is a model and not Ba{\l}aban's step. Its field lives on the periodic strip
$\Lambda=\mathbb{Z}_N\times\{0,1,\dots,H\}$, $N=2P$, with the kinked datum
$\varphi_0(x)=\sigma\|x\|$, where $\|x\|=\operatorname{dist}(x,N\mathbb{Z})$. Its constraints are
the caps of Definition~\ref{8.13:def-two-tier-model}: $\tau$ on the lower tier and the young cap $c$
on the young band.

We set out how the letters of the model are read in the block variables at a level $k_o$.
Put $n=L^{k_o}$ and $\eta=1/n$, and measure fields of the model in the unit $\varepsilon_{k_o}\eta^2$.
Lengths on the block side are measured in top blocks.
Put $\mathsf{w}:=M_2R_{k_o}$, the height of the clamp in top blocks; here $R_{k_o}$ denotes the
level-$k_o$ length entering the zone of the model.

\begin{itemize}
\item[(a)] \emph{Period.} We have $N=2P_{\mathrm{bl}}\,n$, that is, $P=P_{\mathrm{bl}}\,n$, so that one
top block spans $n$ lattice steps of the strip.
Here $P_{\mathrm{bl}}$ is the lateral half-width, in top blocks, of the face of the zone in which
the model is read. In the zone of the proposition on the cutoff-free lateral width, whose
half-width is $D_W$, one has $P_{\mathrm{bl}}=D_W$, and that proposition gives, with its
exponent $a$ of the zone-width bound,
\begin{equation*}
\frac{P_{\mathrm{bl}}}{\mathsf{w}}=\frac{D_W}{\mathsf{w}}
\le 17\,L^{a}\,\frac{M}{M_2}\,\frac{\check R}{R_{k_o}} ,
\end{equation*}
a bound that does not involve the cutoff; here $\check R$ is the parameter of that proposition's
zone which the bound compares with $R_{k_o}$.

\item[(b)] \emph{Heights.} We have $W=\mathsf{w}n$, $\bar W=\varpi_y n$ and $W'=y_+n$.
Here $y_+=\mathsf{w}+\varpi_y$, so that $W'=W+\bar W$. In the geometry of the first young
band, $\varpi_y=\mathsf{w}+\mathsf{w}/L$; under the alternative convention for the width of the
young band, $\varpi_y=\mathsf{w}/L$.

\item[(c)] \emph{Caps and slope.} We have $\tau=\theta'$, where $\theta'=\theta(1+\vartheta_1)$
with $\theta\in[1,2]$.
The datum slope is $\sigma=s_e\theta'$, where $s_e<1$ is the relative kink slope of the exterior.
The young cap is
\begin{equation*}
c=\tau':=\theta_y=\theta'\mathsf{T},\qquad \mathsf{T}=\varepsilon_{k_o-1}L^2/\varepsilon_{k_o}.
\end{equation*}
Since $\varepsilon_{k_o-1}\in[\varepsilon_{k_o}/(1+\beta_0),\varepsilon_{k_o}]$, with $\beta_0$
the constant of this inclusion (distinct from the coupling increments $\beta_{k+1}$), this gives
\begin{equation*}
\frac{\tau'}{\tau}=\mathsf{T}\in\Bigl[\frac{L^2}{1+\beta_0},\,L^2\Bigr],
\qquad
\frac{\tau'}{\sigma}=\frac{\mathsf{T}}{s_e}.
\end{equation*}

\item[(d)] \emph{Depths.} A lattice height $\mathsf{m}$ above the inner face of the clamp corresponds
to the top depth $\delta=\mathsf{w}+\mathsf{m}/n$.

\item[(e)] \emph{Bonds and plaquettes.} By the lemma identifying the model with the
$x_2x_3$-invariant abelian four-dimensional box-clamp problem, the model is read as that
problem. In that reading the vertical bond corresponds to the
normal--tangential plaquette. The horizontal bond corresponds to the tangential--tangential
plaquette.

\item[(f)] \emph{Young windows.} Each young window of the first held young row $\mathfrak{Y}_1$ has
a test region. This region has lateral extent $(1+2/L)\mathsf{w}$ and consists of the rows in
$[W-\mathsf{w}n/L,\,W']$. Call a bond \emph{young-tight} if the young-cap constraint of
Definition~\ref{8.13:def-two-tier-model} is active on it, that is, if $c$ is attained on it.
A window is tight if and only if some bond of its test region is young-tight. The rows in
$[W-\mathsf{w}n/L,\,W-1]$ consist of lower-tier bonds, whose cap is $\tau<c$. These bonds are
therefore never young-tight.
\end{itemize}
\end{remark}

\begin{definition}[Model: the mild-fringe hypothesis]\label{8.13:def-mild-fringe}
The model is the two-tier box-clamp lattice model of
Definition~\ref{8.13:def-two-tier-model}, which is a model statement. It consists of fields on the
periodic strip $\Lambda=\mathbb{Z}_N\times\{0,1,\dots,H\}$, $N=2P$, with the kinked datum
$\varphi_0(x)=\sigma\|x\|$, where $\|x\|=\operatorname{dist}(x,N\mathbb{Z})$, with clamp height $W$
and with young-cap parameter $c$; for each $c$, $\psi^{(c)}$ denotes the minimiser of the Dirichlet
energy $E(\psi)$ over the feasible set $\mathcal{K}_c$ with young cap $c$, and $\tau'$ denotes the
young-cap value at which the doubly-clamped minimiser $\psi^{(\tau')}$ is taken. Every
letter below has the meaning fixed in that model or in this definition, not the meaning fixed in the
glossary; this concerns in particular $g$, $E$, $F$, $S$, $N$, $P$, $H$, $W$, $\sigma$, $c$,
$\tau'$ and $\kappa$.

Let $\tilde{\psi}=\psi^{(\infty)}$ be the straddling-only minimiser, that is, the minimiser of
$E(\psi)$ over the feasible set $\mathcal{K}_\infty$, in which the young-cap constraint is absent. Put
\begin{equation}\label{8.13:eq-mild-fringe-1}
\tilde g:=\tilde{\psi}(\cdot,W),\qquad
\tilde F_y(x):=\tilde{\psi}(x,y+1)-\tilde{\psi}(x,y).
\end{equation}
Let $\tilde S$ be the half-width of the lower-full block of $\tilde{\psi}$ near the bottom tip of the
kink. By the block lemma of the two-tier model, which is proved in that model, the lower-full
columns of $\tilde{\psi}$ that lie in $[-P/2,P/2]$ form an interval:
\begin{equation}\label{8.13:eq-mild-fringe-2}
\{\text{lower-full columns of }\tilde{\psi}\}\cap[-P/2,P/2]=[-\tilde S,\tilde S],
\qquad \tilde S\in\{-1,0,1,\dots\}.
\end{equation}
The value $\tilde S=-1$ corresponds to the empty interval. By the half-period symmetry of the
model, the top tip carries the mirror block.

Fix $\kappa\in(0,\tfrac12]$ with $\kappa W\in\mathbb{N}$ and $\kappa W\ge1$. The
\emph{mild-fringe hypothesis with lateral fraction $\kappa$} is the conjunction of the following two
conditions.
\begin{enumerate}
\item[(i)] (Block.) $\tilde S\ge 2\kappa W$.
\item[(ii)] (No strong column off the kink.) For every $x\in\mathbb{Z}_N$ with
\begin{equation*}
\kappa W\le\|x\|\le P-\kappa W
\end{equation*}
one has
\begin{equation}\label{8.13:eq-mild-fringe-3}
|\tilde F_W(x)|\le\tau'/2 .
\end{equation}
\end{enumerate}
\end{definition}

\begin{remark}[Model: equivalent readings of the mild-fringe hypothesis]\label{8.13:rem-fringe-readings}
The model is the constrained minimisation of the Dirichlet energy $E$ on the periodic strip
$\Lambda$, with kinked datum of slope $\sigma$, straddling cap $\tau$, young cap $\tau'=\mathsf{T}\tau$,
clamp height $W$ and half-period $P$; $\tilde{\psi}$ denotes its straddling-only minimiser. The
mild-fringe hypothesis of Definition~\ref{8.13:def-mild-fringe} (a statement about this model)
concerns $\tilde{\psi}$.
We write $\tilde g=\tilde{\psi}(\cdot,W)$ for its exit trace and
$\tilde F_y(x)=\tilde{\psi}(x,y+1)-\tilde{\psi}(x,y)$ for its vertical bond differences. We write
$\tilde{\mu}(x,y)$ for the multiplier of the vertical bond $(x,y)$ at $\tilde{\psi}$, and $\Delta_x$
for the lateral second difference.

Clause (i) of Definition~\ref{8.13:def-mild-fringe} says that $\tilde g$ has slope exactly
$\pm\sigma$ on $\|z\|<2\kappa W$ and on $\|z-P\|<2\kappa W$. This is the only consequence of (i)
used in the proof of the model theorem on inactivity at full band width.

Clause (ii) says that for $\kappa W\le\|x\|\le P-\kappa W$, that is, beyond lateral distance
$\kappa W$ from either tip, the first free bond above the straddling clamp is pulled by at most
half the young cap.

By the classification of lower columns of the straddling-only minimiser, a lower column is of one
of four kinds:
\begin{itemize}
\item full (a block column), with $\tilde F_W(x)=\tau+\mu_x$, where $\mu_x$ denotes the multiplier
attached to the full column at $x$ in that classification;
\item a tent pole, with
\begin{equation*}
\tilde F_W(x)=\tau+\tilde{\mu}(x,W-1)-\Delta_x\tilde g(x)>\tau ;
\end{equation*}
\item harmonic, with
\begin{equation*}
\tilde F_W(x)=\tilde F_{W-1}(x)-\Delta_x\tilde g(x)\le\tau+2\sigma ;
\end{equation*}
\item an inverted pole.
\end{itemize}
In the harmonic case $\tilde F_W(x)\le\tau'/2$ holds automatically once $\mathsf{T}\ge 2+4s_e$. Indeed,
$\sigma=s_e\tau$ and $\tau'=\mathsf{T}\tau$, so that $\tau+2\sigma=(1+2s_e)\tau\le\mathsf{T}\tau/2$.
Hence clause (ii) is equivalent to the following statement:
\begin{equation}\label{8.13:eq-fringe-readings-1}
\begin{gathered}
\text{no full (block) column, no tent pole and no inverted pole}\\
\text{at lateral distance at least }\kappa W\text{ from the tips}\\
\text{carries a straddling multiplier with }
|\tilde{\mu}(x,W-1)|\ge \tfrac{1}{2}\tau'-\tau-2\sigma .
\end{gathered}
\end{equation}

Apply at height $0$ the kernel representation of the vertical differences of the strip harmonic
extension with symmetric data. By it, \eqref{8.13:eq-fringe-readings-1} is in turn equivalent to the
following statement: $\tilde g$ carries no near-full-strength kink coherent over the lateral scales
\begin{equation*}
\bigl[\,1,\ N e^{-\tau'/(12\sigma)+O(1)}\,\bigr]
\end{equation*}
around any point beyond lateral distance $\kappa W$ from the tips. Since $\tau'/\sigma=\mathsf{T}/s_e$,
this is a regularity statement for the fan described by $\tilde g$ at a lateral scale exponentially
small in $\mathsf{T}/s_e$.

In the language of Ba{\l}aban's large-field analysis, the hypothesis has two parts. First, at the
straddling-only minimiser the excess pull of the straddling windows over their cap is below half
the young cap everywhere beyond lateral distance $\kappa$ times the clamp height from the foot of the
exterior kink. Second, the rigidly transported strip is at least $2\kappa$ clamp heights wide.
\end{remark}

\begin{definition}[Model: the one-sided conditions of the model instance]\label{8.13:def-model-conditions}
The model is the two-tier box-clamp lattice model of
Definition~\ref{8.13:def-two-tier-model}, a model setting and not Ba{\l}aban's step. It is posed on
the periodic strip $\Lambda=\mathbb{Z}_N\times\{0,1,\dots,H\}$, with $N=2P$, and carries the kinked
Dirichlet datum $\varphi_0(x)=\sigma\|x\|$. The heights $W$, $\bar W$, $W'$, $H$, the slope
$\sigma$, the straddling cap $\tau$ and the young cap $\tau'$ are those of that definition. Let
$\kappa$ be the lateral fraction of the mild-fringe hypothesis of
Definition~\ref{8.13:def-mild-fringe}, which is a hypothesis of the model. We put
\begin{equation*}
  \beta_\kappa=\frac{24\tau}{\kappa},\qquad
  B_\beta=\sigma\Bigl[C_K^{0}+6\log^{+}\frac{N}{\bar W}\Bigr],\qquad
  h_{\sharp}=h_{\sharp}(\tau'),\qquad \mathfrak{h}=\mathfrak{h}(\tau'),\qquad C_d=144 .
\end{equation*}
Here $B_\beta$ coincides with the column bound of \eqref{8.13:eq-no-full-column-a}, and
$h_{\sharp}(c)$, $\mathfrak{h}(c)$ are the hot-box scale and the cooling height of
\eqref{8.13:eq-hot-box-a}, as functions of the young cap $c$, here specialised to $c=\tau'$.
The constant $C_T$ is that of \eqref{8.13:eq-kink-pull-decay-a}, the needle bound
$m_{\mathrm{ndl}}(c)$ is that of \eqref{8.13:eq-short-needle-a}, and $h_\rho$ is the margin height of
\eqref{8.13:eq-inactivity-a}; all of these are statements of the model. The
\emph{one-sided conditions of the model instance} are the following four conditions.

The lateral floor:
\begin{equation}\label{8.13:eq-model-conditions-a}
  \sigma\Bigl[C_T+12\log^{+}\frac{N}{2\kappa W-1}\Bigr]+\beta_\kappa\le\frac{\tau'}{2}.
\end{equation}
The top-of-band floor:
\begin{equation}\label{8.13:eq-model-conditions-b}
  B_\beta\le\frac{\tau'}{2}-3\sigma .
\end{equation}
The fit of the hot box and of the needle:
\begin{equation}\label{8.13:eq-model-conditions-c}
\begin{aligned}
  &\text{(a)}\quad \Bigl(\frac{\tau'}{\sigma}\Bigr)^{2}\frac{\bar W^{3}}{W'h_{\sharp}^{2}}
     >110592\,C_d\Bigl(1+\frac{\beta_\kappa+19\sigma}{\tau'}\Bigr),\\
  &\phantom{\text{(a)}}\quad N\,\frac{\tau'}{\sigma}\,\frac{\bar W^{2}}{W'h_{\sharp}^{2}}
     >2304\,C_d\Bigl(1+\frac{\beta_\kappa+19\sigma}{\tau'}\Bigr),\\
  &\text{(b)}\quad h_{\sharp}\Bigl[\frac{128\,C_d(\tau'+\beta_\kappa+19\sigma)}{\sigma}\Bigr]^{1/2}
     \le W,\\
  &\text{(c)}\quad m_{\mathrm{ndl}}(\tau')\le\tfrac12\min(W,\bar W)-2,\qquad
     9h_{\sharp}<\bar W,\qquad h_{\sharp}\le\kappa W,\\
  &\text{(d)}\quad h_\rho\le\bar W-1 .
\end{aligned}
\end{equation}
The ranges of the model instance:
\begin{equation}\label{8.13:eq-model-conditions-d}
\begin{aligned}
  &0<\sigma<\tau,\qquad \tau'>\tau+19\sigma,\qquad \tau'\bar W>\tau W,\qquad \bar W\ge4,
   \qquad \kappa W\in\mathbb{N},\\
  &P\ge 8\ \text{even},\qquad N=2P\ge4(W+\bar W),\\
  &H\ge 2W'+\bar W+2,\qquad H\ge4W+1,\qquad W+9h_{\sharp}<\tfrac12 H .
\end{aligned}
\end{equation}

The following readings use the dictionary between the model and the block variables
(Remark~\ref{8.13:rem-model-dictionary}, a statement about the model). Write
$\mathsf{w}$ for the clamp height in top blocks, so that $W=\mathsf{w}n$. In the block variables
\eqref{8.13:eq-model-conditions-a} reads
\begin{equation*}
  s_e\Bigl[C_T+12\log^{+}\frac{P_{\mathrm{bl}}}{\kappa\mathsf{w}-1/(2n)}\Bigr]
  +\frac{24}{\kappa}\le\frac{\mathsf{T}}{2}.
\end{equation*}
A typical instance of this reading is the pair of conditions $\kappa\mathsf{T}\ge96$ and
\begin{equation*}
  s_e\Bigl[C_T+12\log\frac{P_{\mathrm{bl}}\mathsf{T}}{96\,\mathsf{w}}\Bigr]\le\frac{\mathsf{T}}{4}.
\end{equation*}
Condition \eqref{8.13:eq-model-conditions-b} reads
\begin{equation*}
  s_e\Bigl[C_K^{0}+3+6\log^{+}\frac{2P_{\mathrm{bl}}}{\varpi_y}\Bigr]\le\frac{\mathsf{T}}{2}.
\end{equation*}
All parts of \eqref{8.13:eq-model-conditions-c} have one shape. For instance, the first
inequality of part (c) takes, in the block variables and with the cube-root expression of
\eqref{8.13:eq-short-needle-a} in the place of $m_{\mathrm{ndl}}(\tau')$, the form
\begin{equation*}
  \biggl[128\,C_d\Bigl(\frac{\mathsf{T}}{s_e}+\frac{24}{\kappa s_e}+19\Bigr)\mathsf{w}\,
  P_{\mathrm{bl}}^{2}\,e^{C_T/6}\,e^{-(\mathsf{T}-24/\kappa)/(6s_e)}\biggr]^{1/3}
  \le\tfrac12\min(\mathsf{w},\varpi_y)-\frac{2}{n}.
\end{equation*}
For a constant $C$, this is the condition
\begin{equation*}
  \frac{\mathsf{T}}{6s_e}\ge 2\log\frac{P_{\mathrm{bl}}}{\mathsf{w}}
  +3\log\frac{2\mathsf{w}}{\min(\mathsf{w},\varpi_y)}
  +\log\Bigl(\frac{\mathsf{T}}{s_e}+\frac{24}{\kappa s_e}+19\Bigr)+C .
\end{equation*}
Every part is therefore of the form
\begin{equation*}
  s_e\biggl[C'+C''\log\frac{P_{\mathrm{bl}}}{\mathsf{w}}+\log\frac{\mathsf{T}}{s_e}
  +\Bigl(\log\frac{\mathsf{w}}{\varpi_y}\Bigr)^{+}\biggr]\le c'\,\mathsf{T},
\end{equation*}
with absolute constants $C'$, $C''$, $c'$. By Remark~\ref{8.13:rem-model-dictionary} (a statement
about the model), in the
zone on which the model is read the ratio $P_{\mathrm{bl}}/\mathsf{w}=D_W/\mathsf{w}$ is at most
$17L^{a}(M/M_2)(\check R/R)$, with $a$ the exponent of the width bound of that zone, and is free
of the cutoff. Moreover $\mathsf{T}\ge L^{2}/(1+\beta_0)$ by the same remark about the model. Under
the order of
choice of the constants (Definitions~\ref{2.3:not-stages-first}--\ref{2.3:not-stages-last},
one-sided by Lemma~\ref{2.4:lem-no-cycle}), these conditions are therefore one-sided floors on $L$
for given $M/M_2$. They are of exactly the same kind as the floor
\begin{equation*}
  C_4^{\flat}\bigl(1+\log(L^{2}D_W/\mathsf{w})\bigr)(1+\vartheta_1)^{2}\le(1-\rho)\mathsf{T}
\end{equation*}
of the corollary whose condition it is, with $C_4^{\flat}$ the constant and $\rho$ the parameter
of that corollary. They
can be met: their left-hand sides are
$O(\log L+\log(M/M_2))$, while $\mathsf{T}$ is of order $L^{2}$. No numerical value of $L$ is
proposed.

In the block variables the first three conditions of \eqref{8.13:eq-model-conditions-d} read
\begin{equation*}
  s_e<1,\qquad \mathsf{T}>1+19s_e,\qquad \mathsf{T}\varpi_y>\mathsf{w},
\end{equation*}
and these hold automatically. The condition $N\ge4(W+\bar W)$ reads
$P_{\mathrm{bl}}\ge2(\mathsf{w}+\varpi_y)$: the zone is far wider than the band. The condition
$H\ge4W+1$ places the bond at depth $2W$ inside the lower half of the strip, as used in the
corollary on the bond at depth $2W$, a statement of the model. The condition $W+9h_{\sharp}<H/2$
says that the extent of the box normal
to the face exceeds that of the band by the band's width. In the dictionary the zone on which
the model is read is a box of bounded aspect ratio, with sides $D_i$ between $D_W/\varrho_H$ and
$D_W$, $\varrho_H$ the aspect bound of the zone; the dictionary gives
$D_*=\min(P,(H-2W)/\sqrt3)$.

None of these conditions bounds any of $L$, $K$, $k_o$, $M$, $M/M_2$, $R$, $\eta$, $\theta$,
$s_e$, $g$, or the parameter $N_T$ of the block description, from the side excluded by the order
of choice; the lateral fraction
$\kappa$ is a free parameter
of the mild-fringe hypothesis, and \eqref{8.13:eq-model-conditions-a} fixes it at
order $\mathsf{T}^{-1}$.
\end{definition}

We work in the model of Definition~\ref{8.13:def-two-tier-model}, in the form displayed there; this is a model statement. Throughout this lemma and the statements that follow it we write $\psi:=\psi^{(c)}$.

All statements are made near the bottom tip, that is, on the columns with $\|x\|\le P/2$, and near the bottom band, that is, on the rows $W,\dots,W'$. The corresponding statements at the top tip and in the top band follow from the symmetries of the model: the map $(x,\psi)\mapsto(x+P,\sigma P-\psi)$, which exchanges up and down, and the map $y\mapsto H-y$.

The following facts about the model's minimiser are used throughout:
\begin{itemize}
\item the comparison principle, $|\psi(x+1,y)-\psi(x,y)|\le\sigma$ for all $(x,y)$, and hence $|\Delta_x\psi|\le 2\sigma$;
\item the bond equation \eqref{8.13:eq-opposite-bonds-2}, recalled below;
\item the lower and young lift deficits satisfy $f\le 0$ and $f^{y}\le0$;
\item $\psi\ge\varphi_0$ on the columns with $\|x\|\le P/2$.
\end{itemize}
The comparison principle, the bond equation with the sign rules of its multipliers, and the vanishing of the horizontal multipliers are items of the minimiser theorem for lattice models with $x$-invariant componentwise caps, of which the model of Definition~\ref{8.13:def-two-tier-model} is an instance for every $c$. The two symmetries above and the lower bound $\psi\ge\varphi_0$ are statements about the model of Definition~\ref{8.13:def-two-tier-model}, and the signs of $f$ and $f^{y}$ are the sign bounds for the lower and young lift deficits established for that model. We use all of them as stated there; each of them is a model statement.

For each constrained vertical bond $y$ let $\mathrm{cap}_y\in\{\tau,c\}$ be the bound on $|F_y(x)|$ carried by the row $y$; it depends on $y$ only. Let $\lambda^{\pm}(x,y)\ge0$ be the multipliers of the constraints $\pm F_y(x)\le \mathrm{cap}_y$; they satisfy complementary slackness. Put $\mu(x,y)=\lambda^{+}(x,y)-\lambda^{-}(x,y)$, and set $\mu(x,y)=0$ on free bonds $y$. Thus $\mu>0$ forces $F=+\mathrm{cap}$, $\mu<0$ forces $F=-\mathrm{cap}$, and $|F|<\mathrm{cap}$ forces $\mu=0$. The horizontal multipliers vanish, because $|\psi(x+1,y)-\psi(x,y)|\le\sigma<\tau$ and $\tau$ is at most every horizontal cap. The bond equation then reads, at every non-Dirichlet site,
\begin{equation}\label{8.13:eq-opposite-bonds-2}
\Delta\psi(x,y)=\mu(x,y-1)-\mu(x,y).
\end{equation}

\begin{lemma}[Model: no adjacent opposite tight bonds]\label{8.13:lem-opposite-bonds}
This is a statement about the two-tier box-clamp lattice model of
Definition~\ref{8.13:def-two-tier-model}. In it, $\psi=\psi^{(c)}$ is the minimiser of the
Dirichlet energy $E$ over the feasible set $\mathcal{K}_c$ on the strip
$\Lambda=\mathbb{Z}_N\times\{0,1,\dots,H\}$, with the datum built from $\varphi_0$ as fixed
there. In $\mathcal{K}_c$ the vertical bond difference $F_y(x)=\psi(x,y+1)-\psi(x,y)$ of every
constrained bond $y$ satisfies $|F_y(x)|\le\mathrm{cap}_y$, where the cap $\mathrm{cap}_y\in\{\tau,c\}$
depends only on the row $y$, the young cap $c$ being carried by the young tier of bonds
$W,\dots,W'-1$. The slope and the two caps satisfy $\sigma<\tau<c$ and $c\bar W>\tau W$.

Let $x$ be a column and let the vertical bonds $y$ and $y+1$ of that column both be
constrained. The two bonds may belong to either tier; in particular the interface pair $W-1$, $W$ is included. Then
\begin{equation*}
F_y(x)=+\mathrm{cap}_y,\quad F_{y+1}(x)=-\mathrm{cap}_{y+1}
\end{equation*}
is impossible, and so is
\begin{equation*}
F_y(x)=-\mathrm{cap}_y,\quad F_{y+1}(x)=+\mathrm{cap}_{y+1}.
\end{equation*}
\end{lemma}

\begin{proof}
The argument uses only the comparison principle, the bond equation with the sign rules of its multipliers, and $\mathrm{cap}_y,\mathrm{cap}_{y+1}\ge\tau>\sigma$.

Write the four-neighbour Laplacian at $(x,y+1)$ as its vertical part plus its lateral part:
\begin{equation}\label{8.13:eq-opposite-bonds-1}
\Delta\psi(x,y+1)=\bigl(F_{y+1}(x)-F_y(x)\bigr)+\Delta_x\psi(x,y+1).
\end{equation}

\emph{First case.} Suppose $F_y(x)=+\mathrm{cap}_y$ and $F_{y+1}(x)=-\mathrm{cap}_{y+1}$. We have $|\Delta_x\psi|\le2\sigma$. Both caps are at least $\tau$, and $\sigma<\tau$. Hence \eqref{8.13:eq-opposite-bonds-1} gives
\begin{equation*}
\Delta\psi(x,y+1)=-\mathrm{cap}_{y+1}-\mathrm{cap}_y+\Delta_x\psi(x,y+1)\le-2\tau+2\sigma<0.
\end{equation*}

We now obtain the opposite sign from the multipliers.
\begin{itemize}
\item Since $F_y(x)=+\mathrm{cap}_y\ne-\mathrm{cap}_y$, complementary slackness gives $\lambda^-(x,y)=0$. Hence $\mu(x,y)=\lambda^+(x,y)\ge0$.
\item Since $F_{y+1}(x)=-\mathrm{cap}_{y+1}\ne+\mathrm{cap}_{y+1}$, we get $\lambda^+(x,y+1)=0$. Hence $\mu(x,y+1)=-\lambda^-(x,y+1)\le0$.
\end{itemize}

The bond equation \eqref{8.13:eq-opposite-bonds-2} at $(x,y+1)$ now gives
\begin{equation*}
\Delta\psi(x,y+1)=\mu(x,y)-\mu(x,y+1)\ge0,
\end{equation*}
a contradiction.

\emph{Second case.} This is the sign reverse of the first. Suppose $F_y(x)=-\mathrm{cap}_y$ and $F_{y+1}(x)=+\mathrm{cap}_{y+1}$. Then \eqref{8.13:eq-opposite-bonds-1} gives
\begin{equation*}
\Delta\psi(x,y+1)=\mathrm{cap}_{y+1}+\mathrm{cap}_y+\Delta_x\psi(x,y+1)\ge2\tau-2\sigma>0.
\end{equation*}

On the other hand:
\begin{itemize}
\item $F_y(x)\ne+\mathrm{cap}_y$ gives $\lambda^+(x,y)=0$, hence $\mu(x,y)\le0$;
\item $F_{y+1}(x)\ne-\mathrm{cap}_{y+1}$ gives $\lambda^-(x,y+1)=0$, hence $\mu(x,y+1)\ge0$.
\end{itemize}

So the bond equation \eqref{8.13:eq-opposite-bonds-2} at $(x,y+1)$ gives $\Delta\psi(x,y+1)=\mu(x,y)-\mu(x,y+1)\le0$, again a contradiction.
\end{proof}

\emph{Ends of young runs.} Fix a column $x$. A \emph{young up-run} is a maximal set of consecutive bonds $p,\dots,q-1$, contained in $\{W,\dots,W'-1\}$, on which $F_y(x)=c$ for $y=p,\dots,q-1$. At each end that lies inside the young tier, a young up-run has a non-tight neighbouring bond.
\begin{itemize}
\item \emph{Lower end.} Let $p\ge W+1$. Maximality gives $F_{p-1}(x)\ne c$. Lemma~\ref{8.13:lem-opposite-bonds} (second case, with $y=p-1$) gives $F_{p-1}(x)\ne-c$. Hence $|F_{p-1}(x)|<c$.
\item \emph{Upper end.} Let $q\le W'-1$. Maximality gives $F_q(x)\ne c$, and the first case of the lemma (with $y=q-1$) gives $F_q(x)\ne -c$. Hence $|F_q(x)|<c$.
\end{itemize}

Consequently $\mu(x,p-1)=0$, respectively $\mu(x,q)=0$, in these cases.

Moreover $\mu(x,W')=0$ always. Indeed, bond $W'$ is free, since $W'\le H-W'-1$ by the one-sided condition~\eqref{8.13:eq-model-conditions-d} of the model instance, Definition~\ref{8.13:def-model-conditions}, itself a model statement.

On the run itself, $F=c\ne-c$ gives $\lambda^-=0$, so $\mu=\lambda^+\ge0$.

In summary, for a young up-run $p,\dots,q-1$ on the column $x$,
\begin{equation}\label{8.13:eq-opposite-bonds-3}
\begin{aligned}
&|F_{p-1}(x)|<c \text{ and } \mu(x,p-1)=0 &&\text{if } p\ge W+1,\\
&|F_{q}(x)|<c \text{ and } \mu(x,q)=0 &&\text{if } q\le W'-1,\\
&\mu(x,W')=0,\\
&\mu(x,y)=\lambda^{+}(x,y)\ge0 &&\text{for } p\le y\le q-1.
\end{aligned}
\end{equation}

The mirror statements hold for \emph{young down-runs}, the maximal sets of consecutive young bonds $p,\dots,q-1$ on which $F_y(x)=-c$ for $y=p,\dots,q-1$. For such a run:
\begin{itemize}
\item $|F_{p-1}(x)|<c$ and $\mu(x,p-1)=0$ if $p\ge W+1$;
\item $|F_q(x)|<c$ and $\mu(x,q)=0$ if $q\le W'-1$;
\item $\mu(x,W')=0$;
\item on the run, $\mu=-\lambda^-\le0$.
\end{itemize}

The proofs are the same, with the two cases of the lemma exchanged.

\begin{lemma}[Model: no floating tight runs]\label{8.13:lem-floating-runs}
The model is the constrained Dirichlet problem of this chapter on the periodic strip
$\Lambda=\mathbb{Z}_N\times\{0,\dots,H\}$: $\psi=\psi^{(c)}$ is the minimiser of the Dirichlet energy
$E(\psi)$ over the feasible set $\mathcal{K}_c$. Recall that in $\mathcal{K}_c$ the vertical bond
difference $F_y(x)=\psi(x,y+1)-\psi(x,y)$ satisfies $|F_y(x)|\le \operatorname{cap}_y$. Here
$\operatorname{cap}_y=\tau$ on the lower tier of bonds $0\le y\le W-1$, and
$\operatorname{cap}_y=c$, the young cap, on the young tier $W\le y\le W'-1$. On a column $x$ an
\emph{up-run} (resp.\ \emph{down-run}) of a tier is a maximal interval $[p,q-1]$ of bonds of that
tier on which $F_y(x)=+\operatorname{cap}_y$ (resp.\ $F_y(x)=-\operatorname{cap}_y$). A bond $y$ on
the column $x$ is \emph{tight} if $|F_y(x)|=\operatorname{cap}_y$, an \emph{up-bond} if
$F_y(x)=+\operatorname{cap}_y$ and a \emph{down-bond} if $F_y(x)=-\operatorname{cap}_y$. The
statement concerns the bottom band; the top band follows from it by the symmetries
$(x,\psi)\mapsto(x+P,\sigma P-\psi)$, which exchanges up and down, and $y\mapsto H-y$. Then:
\begin{itemize}
\item[(a)] every young up-run $[p,q-1]$ has $p=W$ or $q=W'$, and so does every young down-run;
\item[(b)] every lower up-run or down-run $[p,q-1]\subset[0,W-1]$ has $p=0$ or $q=W$.
\end{itemize}
\end{lemma}

\begin{proof}
We write $A_y(x)=\tfrac12[\psi(x-1,y)+\psi(x+1,y)]$ for the lateral neighbour mean and
$\Delta_x\psi(x,y)=\psi(x+1,y)+\psi(x-1,y)-2\psi(x,y)$ for the lateral second difference. Thus
\begin{equation}\label{8.13:eq-floating-runs-1}
\Delta_x\psi(x,y)=2\bigl(A_y(x)-\psi(x,y)\bigr),\qquad
\Delta\psi(x,y)=F_y(x)-F_{y-1}(x)+\Delta_x\psi(x,y),
\end{equation}
where $\Delta$ is the four-neighbour Laplacian. Recall that the minimality condition of the model
reads
\begin{equation}\label{8.13:eq-floating-runs-2}
\Delta\psi(x,y)=\mu(x,y-1)-\mu(x,y).
\end{equation}
In it $\mu=\lambda^+-\lambda^-$ with multipliers
$\lambda^\pm\ge0$, and $\lambda^+(x,y)=0$ unless $F_y(x)=\operatorname{cap}_y$, while
$\lambda^-(x,y)=0$ unless $F_y(x)=-\operatorname{cap}_y$. Hence $\mu(x,y)=0$ on a bond that is not
tight, $\mu=\lambda^+\ge0$ on an up-run, and $\mu=-\lambda^-\le0$ on a down-run.

\emph{Young up-runs.} Suppose $[p,q-1]$ is a young up-run on the column $x$ with
$W+1\le p<q\le W'-1$; we derive a contradiction.

The bond $p-1$ is young. It is not an up-bond, by maximality. It is not a down-bond either, by the
model lemma on opposite tight bonds, Lemma~\ref{8.13:lem-opposite-bonds}, which is a model statement
and excludes $F_{p-1}=-c$, $F_p=+c$. Hence $|F_{p-1}(x)|<c$ and $\mu(x,p-1)=0$. The same argument
gives $|F_q(x)|<c$ and $\mu(x,q)=0$.

At $(x,p)$, \eqref{8.13:eq-floating-runs-1} gives
$\Delta\psi(x,p)=(c-F_{p-1}(x))+\Delta_x\psi(x,p)$ with $c-F_{p-1}(x)>0$. By
\eqref{8.13:eq-floating-runs-2},
$\Delta\psi(x,p)=-\mu(x,p)\le0$. So $\Delta_x\psi(x,p)<0$, that is, $\psi(x,p)>A_p(x)$.

At $(x,q)$, $\Delta\psi(x,q)=(F_q(x)-c)+\Delta_x\psi(x,q)$ with $F_q(x)-c<0$. By
\eqref{8.13:eq-floating-runs-2},
$\Delta\psi(x,q)=\mu(x,q-1)\ge0$. So $\Delta_x\psi(x,q)>0$, that is, $\psi(x,q)<A_q(x)$.

On the run, $\psi(x,q)-\psi(x,p)=\sum_{y=p}^{q-1}F_y(x)=c(q-p)$. The bonds $p,\dots,q-1$ of the
neighbouring columns $x\pm1$ are young, so feasibility gives $F_y(x\pm1)\le c$ for each of them, and
\begin{equation*}
A_q(x)-A_p(x)=\tfrac12\sum_{\pm}\bigl[\psi(x\pm1,q)-\psi(x\pm1,p)\bigr]\le c(q-p).
\end{equation*}
Therefore
\begin{equation*}
\psi(x,q)=\psi(x,p)+c(q-p)>A_p(x)+c(q-p)\ge A_p(x)+\bigl(A_q(x)-A_p(x)\bigr)=A_q(x),
\end{equation*}
which contradicts $\psi(x,q)<A_q(x)$.

\emph{Young down-runs.} Let $F_y(x)=-c$ for $p\le y\le q-1$, with $W+1\le p<q\le W'-1$. As before,
$|F_{p-1}(x)|<c$ and $|F_q(x)|<c$, so $\mu(x,p-1)=\mu(x,q)=0$.

At $(x,p)$, $\Delta\psi(x,p)=(-c-F_{p-1}(x))+\Delta_x\psi(x,p)$ with $-c-F_{p-1}(x)<0$. By
\eqref{8.13:eq-floating-runs-2},
$\Delta\psi(x,p)=-\mu(x,p)\ge0$. So $\Delta_x\psi(x,p)>0$, that is, $\psi(x,p)<A_p(x)$.

At $(x,q)$, $\Delta\psi(x,q)=(F_q(x)+c)+\Delta_x\psi(x,q)$ with $F_q(x)+c>0$. By
\eqref{8.13:eq-floating-runs-2},
$\Delta\psi(x,q)=\mu(x,q-1)\le0$. So $\Delta_x\psi(x,q)<0$, that is, $\psi(x,q)>A_q(x)$.

Now $\psi$ drops by exactly $c$ per step on the run, while $F_y(x\pm1)\ge-c$ gives
$A_q(x)-A_p(x)\ge-c(q-p)$. Hence
\begin{equation*}
\psi(x,q)=\psi(x,p)-c(q-p)<A_p(x)-c(q-p)\le A_q(x),
\end{equation*}
a contradiction.

\emph{Lower runs.} Let $[p,q-1]\subset[0,W-1]$ be a lower up-run or down-run with
$1\le p<q\le W-1$. The bonds $p-1,\dots,q$ at $x$ and the bonds $p,\dots,q-1$ at $x\pm1$ are all
lower bonds, with cap $\tau$. Lemma~\ref{8.13:lem-opposite-bonds}, which is a model statement and
applies to constrained bonds of either tier, and maximality give $|F_{p-1}(x)|<\tau$ and
$|F_q(x)|<\tau$, so $\mu(x,p-1)=\mu(x,q)=0$.

For a lower up-run, \eqref{8.13:eq-floating-runs-1} gives
$\Delta\psi(x,p)=(\tau-F_{p-1}(x))+\Delta_x\psi(x,p)$ with $\tau-F_{p-1}(x)>0$, and
\eqref{8.13:eq-floating-runs-2} gives $\Delta\psi(x,p)=-\mu(x,p)\le0$; so $\Delta_x\psi(x,p)<0$,
that is, $\psi(x,p)>A_p(x)$. Likewise
$\Delta\psi(x,q)=(F_q(x)-\tau)+\Delta_x\psi(x,q)$ with $F_q(x)-\tau<0$, and
$\Delta\psi(x,q)=\mu(x,q-1)\ge0$; so $\Delta_x\psi(x,q)>0$, that is, $\psi(x,q)<A_q(x)$. On the run
$\psi(x,q)-\psi(x,p)=\tau(q-p)$, while the lower bonds of the neighbouring columns satisfy
$F_y(x\pm1)\le\tau$, so $A_q(x)-A_p(x)\le\tau(q-p)$. Hence
\begin{equation*}
\psi(x,q)=\psi(x,p)+\tau(q-p)>A_p(x)+\tau(q-p)\ge A_q(x),
\end{equation*}
which contradicts $\psi(x,q)<A_q(x)$. For a lower down-run the signs reverse: the same two
computations give $\psi(x,p)<A_p(x)$ and $\psi(x,q)>A_q(x)$, while $\psi$ drops by exactly $\tau$ per
step on the run and $F_y(x\pm1)\ge-\tau$ gives $A_q(x)-A_p(x)\ge-\tau(q-p)$, so
\begin{equation*}
\psi(x,q)=\psi(x,p)-\tau(q-p)<A_p(x)-\tau(q-p)\le A_q(x),
\end{equation*}
a contradiction.
\end{proof}

\begin{lemma}[Model: hanging young runs and full columns]\label{8.13:lem-hanging-runs}
We work in the model of this section. It is a minimisation problem on the strip
$\Lambda=\mathbb{Z}_N\times\{0,\dots,H\}$. The field $\psi=\psi^{(c)}$ minimises the Dirichlet energy $E(\psi)$
over the feasible set $\mathcal{K}_c$; among the constraints defining $\mathcal{K}_c$, the vertical bond
differences $F_y(x)=\psi(x,y+1)-\psi(x,y)$ of the young bonds $W\le y\le W'-1$ satisfy
$|F_y(x)|\le c$, and the lower bonds $0\le y\le W-1$ carry the straddling cap $\tau$.
We write $A_y(x)=\tfrac12[\psi(x-1,y)+\psi(x+1,y)]$ and
$\Delta_x\psi(x,y)=2[A_y(x)-\psi(x,y)]$. The multipliers of the vertical bonds are
$\mu=\lambda^{+}-\lambda^{-}$, with $\lambda^{\pm}\ge0$ vanishing unless the bond is tight at $\pm$ its
cap, and the stationarity condition of the constrained minimiser reads
\begin{equation}\label{8.13:eq-hanging-runs-1}
\Delta\psi(x,y)=\mu(x,y-1)-\mu(x,y).
\end{equation}
A young up-run (down-run) of column $x$ is a maximal interval $[p,q-1]$ of young bonds on which
$F_y(x)=c$ (respectively $F_y(x)=-c$). Column $x$ is young-full if it is of the full-column type
among the run types of this section; below we use only that its young bonds $W,\dots,W'-1$ form
an up-run $[W,W'-1]$, the case $p=W$ complementary to (a). Bond $W'$ carries no active
constraint, so $\mu(x,W')=0$. Then:
\begin{itemize}
\item[(a)] If a young up-run is $[p,W'-1]$ with $p\ge W+1$, then $\psi(x,y)>A_y(x)$ for every
$p\le y\le W'$, and
\begin{equation}\label{8.13:eq-hanging-runs-2}
F_{W'}(x)=c+\mu(x,W'-1)-\Delta_x\psi(x,W')>c.
\end{equation}
\item[(b)] If a young down-run is $[p,W'-1]$ with $p\ge W+1$, then $F_{W'}(x)<-c$.
\item[(c)] If column $x$ is young-full, then
\begin{equation}\label{8.13:eq-hanging-runs-3}
F_{W'}(x)\ge c-\Delta_x\psi(x,W')\ge c-2\sigma.
\end{equation}
\end{itemize}
\end{lemma}

\begin{proof}
(a) We first treat the site $(x,p)$ by the argument used at the lower end of a run in the proof of
the model lemma, Lemma~\ref{8.13:lem-floating-runs}. That argument uses only the site $(x,p)$ and
the hypothesis $p\ge W+1$, so it applies here. At this site
$\Delta\psi(x,p)=(c-F_{p-1})+\Delta_x\psi(x,p)$, and $c-F_{p-1}>0$ because $p-1\ge W$ is a young
bond outside the maximal run. On the other hand \eqref{8.13:eq-hanging-runs-1} gives
\begin{equation*}
\Delta\psi(x,p)=\mu(x,p-1)-\mu(x,p)=-\mu(x,p)\le0 .
\end{equation*}
Hence $\Delta_x\psi(x,p)<0$, that is, $\psi(x,p)>A_p(x)$.

Now let $p\le y\le W'$. The bonds $p,\dots,y-1$ of column $x$ lie in the run, so
$\psi(x,y)=\psi(x,p)+c(y-p)$. The bonds $p,\dots,y-1$ of the neighbouring columns $x\pm1$ are young
bonds, each with $F\le c$, so $\psi(x\pm1,y)-\psi(x\pm1,p)\le c(y-p)$. Averaging over the two
neighbours gives $A_y(x)\le A_p(x)+c(y-p)$. Therefore
\begin{equation*}
\psi(x,y)-A_y(x)\ge\psi(x,p)-A_p(x)>0 .
\end{equation*}

At the site $(x,W')$ we have $F_{W'-1}=c$, so
$\Delta\psi(x,W')=(F_{W'}-c)+\Delta_x\psi(x,W')$. Since bond $W'$ is free and bond $W'-1$ is tight
at $+c$, \eqref{8.13:eq-hanging-runs-1} gives
\begin{equation*}
\Delta\psi(x,W')=\mu(x,W'-1)-\mu(x,W')=\mu(x,W'-1)\ge0 .
\end{equation*}
Equating the two expressions yields the identity in \eqref{8.13:eq-hanging-runs-2}. The case
$y=W'$ of the bound just proved gives $\Delta_x\psi(x,W')=2[A_{W'}(x)-\psi(x,W')]<0$. Hence
$F_{W'}(x)>c$.

(b) This is the mirror image of (a). At $(x,p)$ we have
$\Delta\psi(x,p)=(-c-F_{p-1})+\Delta_x\psi(x,p)$ with $-c-F_{p-1}<0$, while
$\Delta\psi(x,p)=-\mu(x,p)\ge0$. Hence $\psi(x,p)<A_p(x)$.

For $p\le y\le W'$ we have $\psi(x,y)=\psi(x,p)-c(y-p)$. The neighbours' young bonds satisfy
$F\ge-c$, so $A_y(x)\ge A_p(x)-c(y-p)$, and therefore $\psi(x,y)<A_y(x)$. In particular
$\Delta_x\psi(x,W')>0$.

At $(x,W')$ we have
\begin{equation*}
\Delta\psi(x,W')=(F_{W'}+c)+\Delta_x\psi(x,W')=\mu(x,W'-1)\le0 .
\end{equation*}
Hence
\begin{equation*}
F_{W'}(x)=-c+\mu(x,W'-1)-\Delta_x\psi(x,W')<-c .
\end{equation*}

(c) For a young-full column, bond $W'-1$ is tight at $+c$ and bond $W'$ is free. Exactly as in (a),
\begin{equation*}
F_{W'}(x)-c+\Delta_x\psi(x,W')=\mu(x,W'-1)\ge0 .
\end{equation*}
Together with the model's bound $\Delta_x\psi(x,W')\le2\sigma$ on the lateral second difference
this gives \eqref{8.13:eq-hanging-runs-3}.

In words: a young constraint active at the top of the band occurs only together with a free bond
immediately above the band, and that bond is pulled beyond the young cap.
\end{proof}

\begin{lemma}[Model: lower tent poles are laterally concave]\label{8.13:lem-tent-poles-concave}
The model is the strip model of this section. The field $\psi=\psi^{(c)}$ is the minimiser of the
Dirichlet energy $E$ over the feasible set $\mathcal{K}_c$ with young cap $c$ on the strip $\Lambda$.
Lower bonds, their cap $\tau$, and the lower up-runs and down-runs of a column are as fixed in the
set-up of the strip model. We write $g=\psi(\cdot,W)$ for the exit trace.

Let $x$ be a column.
\begin{itemize}
\item[(a)] If a lower up-run on $x$ is $[p,W-1]$ with $p\ge1$, then $\psi(x,y)>A_y(x)$ for
$p\le y\le W$; in particular $\Delta_x g(x)<0$.
\item[(b)] If a lower down-run on $x$ is $[p,W-1]$ with $p\ge1$, then $\Delta_x g(x)>0$.
\end{itemize}
\end{lemma}

\begin{proof}
We use the multipliers $\mu=\lambda^{+}-\lambda^{-}$ of the vertical bonds and the stationarity
condition of the model, both fixed in its set-up. Here $\lambda^{\pm}\ge0$, and for a lower bond $y$
the multiplier $\lambda^{\pm}(x,y)$ vanishes unless $F_y(x)=\pm\tau$ (for a young bond, unless
$F_y(x)=\pm c$). The stationarity condition reads
\[
\Delta\psi(x,y)=\mu(x,y-1)-\mu(x,y),
\qquad
\Delta\psi(x,y)=\bigl(F_y(x)-F_{y-1}(x)\bigr)+\Delta_x\psi(x,y).
\]
Finally, since $\Delta_x\psi(x,y)=\psi(x-1,y)+\psi(x+1,y)-2\psi(x,y)$, we have
$\Delta_x\psi(x,y)=2[A_y(x)-\psi(x,y)]$.

(a) Since $1\le p\le W-1$, the bond $p-1$ is a lower bond, so $|F_{p-1}(x)|\le\tau$. By maximality
of the run, $F_{p-1}(x)\ne+\tau$. Since bond $p$ has $F_p(x)=+\tau$,
Lemma~\ref{8.13:lem-opposite-bonds}, a statement proved in the model, excludes
$F_{p-1}(x)=-\tau$. Hence
$|F_{p-1}(x)|<\tau$, and therefore $\mu(x,p-1)=0$.

Bond $p$ sits at its upper cap, so $\lambda^{-}(x,p)=0$ and $\mu(x,p)\ge0$. Stationarity at
$(x,p)$ gives
\[
(\tau-F_{p-1}(x))+\Delta_x\psi(x,p)=-\mu(x,p)\le0 .
\]
As $\tau-F_{p-1}(x)>0$, this forces $\Delta_x\psi(x,p)<0$, that is, $\psi(x,p)>A_p(x)$.

Now let $p\le y\le W$. The bonds $p,\dots,y-1$ of column $x$ lie in the run, so
$\psi(x,y)=\psi(x,p)+\tau(y-p)$. On the neighbouring columns $x\pm1$, the bonds $p,\dots,y-1$ are
lower bonds, because $y-1\le W-1$. Each therefore has $F_j(x\pm1)\le\tau$, which gives
\[
A_y(x)=A_p(x)+\sum_{j=p}^{y-1}\tfrac12\bigl[F_j(x-1)+F_j(x+1)\bigr]\le A_p(x)+\tau(y-p).
\]
Hence
\[
\psi(x,y)-A_y(x)\ge\psi(x,p)-A_p(x)>0 .
\]
Taking $y=W$ gives
\[
\Delta_xg(x)=2[A_W(x)-\psi(x,W)]<0 .
\]
The young bond $W$ above the run enters none of these steps.

(b) This is the mirror image of (a). As before, bond $p-1$ is a lower bond. By maximality,
$F_{p-1}(x)\ne-\tau$. Since $F_p(x)=-\tau$, Lemma~\ref{8.13:lem-opposite-bonds}, a statement
proved in the model, excludes $F_{p-1}(x)=+\tau$. Hence $|F_{p-1}(x)|<\tau$ and
$\mu(x,p-1)=0$.

Bond $p$ sits at its lower cap, so $\lambda^{+}(x,p)=0$ and $\mu(x,p)\le0$. Stationarity at
$(x,p)$ gives
\[
(-\tau-F_{p-1}(x))+\Delta_x\psi(x,p)=-\mu(x,p)\ge0 .
\]
As $-\tau-F_{p-1}(x)<0$, this forces $\Delta_x\psi(x,p)>0$, that is, $\psi(x,p)<A_p(x)$.

For $p\le y\le W$ we have $\psi(x,y)=\psi(x,p)-\tau(y-p)$. The neighbouring columns' lower bonds
satisfy $F_j(x\pm1)\ge-\tau$, so $A_y(x)\ge A_p(x)-\tau(y-p)$. Hence
\[
\psi(x,y)-A_y(x)\le\psi(x,p)-A_p(x)<0 .
\]
At $y=W$ this gives $\Delta_xg(x)=2[A_W(x)-\psi(x,W)]>0$.
\end{proof}

\begin{lemma}[Model: anchored young runs sit on tip columns]\label{8.13:lem-anchored-runs}
The model is the lattice obstacle problem on the periodic strip
$\Lambda=\mathbb{Z}_N\times\{0,\dots,H\}$ in which $\psi=\psi^{(c)}$ is the minimiser of the
Dirichlet energy $E(\psi)$ over the feasible set $\mathcal{K}_c$. Among the constraints defining
$\mathcal{K}_c$ are the following:
\begin{itemize}
\item the bottom row carries the kinked datum, $\psi(\cdot,0)=\varphi_0$;
\item the vertical bond differences $F_y(x)=\psi(x,y+1)-\psi(x,y)$ of the \emph{lower} bonds
$0\le y\le W-1$ satisfy $|F_y(x)|\le\tau$;
\item those of the \emph{young} bonds $W\le y\le W'-1$ satisfy $|F_y(x)|\le c$.
\end{itemize}
The young cap exceeds the lower cap, $c>\tau$.
We write $g(x)=\psi(x,W)$ and
\[
\Delta_x g(x)=g(x+1)+g(x-1)-2g(x),
\]
and likewise $\Delta_x\psi(x,y)$ and $\Delta_x\varphi_0(x)$ for the second differences in the
column index. A \emph{young up-run} of column $x$ is a maximal
interval $[p,q-1]$ of young bond indices on which $F_y(x)=c$. A \emph{lower up-run} is a maximal
interval of lower bond indices on which $F_y(x)=\tau$. A bond $y$ of column $x$ is
\emph{up-tight} if $F_y(x)$ equals the cap of that bond, that is, $\tau$ for a lower bond and $c$
for a young bond.

Let $x$ be a column with $\|x\|\le P/2$, and suppose that $x$ carries a young up-run $[W,q-1]$
with $W+1\le q\le W'-1$. The run is anchored at the inner face of the clamp and does not reach
the top of the young band. Then:
\begin{itemize}
\item[(a)] $\Delta_x g(x)>0$;
\item[(b)] the lower bond $W-1$ of column $x$ is up-tight, with $\mu(x,W-1)>c-\tau>0$;
\item[(c)] column $x$ is lower-full: $g(x)=\varphi_0(x)+\tau W$;
\item[(d)] $\Delta_x\varphi_0(x)>0$.
\end{itemize}
In other words, $x$ lies in the tip set $\mathcal{T}$ of the kink. For the datum
$\varphi_0(x)=\sigma\|x\|$ this means $x\equiv0\pmod N$, the kink's tip column.
\end{lemma}

We use the Karush--Kuhn--Tucker conditions of the model in the sign convention fixed in the
set-up of the model in this section. There is a multiplier $\mu(x,y)$ for each vertical
bond such that, at every site $(x,y)$ of the strip at which the stationarity condition of the
minimisation is imposed, and in particular at the sites $(x,q)$ and $(x,W)$ used below,
\begin{equation}\label{8.13:eq-anchored-runs-1}
\Delta\psi(x,y)
=\bigl(F_y(x)-F_{y-1}(x)\bigr)+\Delta_x\psi(x,y)
=\mu(x,y-1)-\mu(x,y).
\end{equation}
The multipliers obey three rules:
\begin{itemize}
\item $\mu(x,y)=0$ if $|F_y(x)|$ is strictly below the cap of the bond;
\item $\mu(x,y)\ge0$ if the bond is up-tight;
\item $\mu(x,y)>0$ forces the upper multiplier $\lambda^{+}(x,y)$ to be positive, and hence forces
$F_y(x)$ to equal the upper cap of the bond.
\end{itemize}
For a column $x'$ we also use the young and lower lift deficits
\[
f^{y}(x',y)=\psi(x',y)-g(x')-c(y-W),\qquad
f(x',y)=\psi(x',y)-\varphi_0(x')-\tau y .
\]

\begin{proof}
(a) Consider the site $(x,q)$. Since $q\le W'-1$, bond $q$ is young. It lies outside the run,
and it satisfies $|F_q(x)|<c$, so $\mu(x,q)=0$. Bond $q-1$ belongs to the run, so
$F_{q-1}(x)=c$ and $\mu(x,q-1)\ge0$. Hence \eqref{8.13:eq-anchored-runs-1} at $(x,q)$ reads
\[
(F_q(x)-c)+\Delta_x\psi(x,q)=\mu(x,q-1)\ge0,
\]
and therefore $\Delta_x\psi(x,q)\ge c-F_q(x)>0$.

Along the run every bond $W,\dots,q-1$ of column $x$ has difference $c$, so
$\psi(x,q)=g(x)+c(q-W)$. In the columns $x\pm1$ the bonds $W,\dots,q-1$ are young bonds, each
with difference at most $c$. Summing these differences gives
$\psi(x\pm1,q)\le g(x\pm1)+c(q-W)$, that is, $f^{y}(x\pm1,q)\le0$. Inserting the three
expressions into the definition of $\Delta_x$, the terms $c(q-W)$ cancel, and
\[
\Delta_x\psi(x,q)=\Delta_x g(x)+f^{y}(x+1,q)+f^{y}(x-1,q)\le\Delta_x g(x).
\]
Combined with $\Delta_x\psi(x,q)>0$ this gives $\Delta_x g(x)>0$.

(b) Consider the site $(x,W)$. Since $q\ge W+1$, bond $W$ belongs to the run, so $F_W(x)=c$.
Bond $W$ is up-tight, so $\mu(x,W)\ge0$. Moreover $\Delta_x\psi(x,W)=\Delta_x g(x)$ by
definition of $g$. Thus \eqref{8.13:eq-anchored-runs-1} at $(x,W)$ reads
\[
\bigl(c-F_{W-1}(x)\bigr)+\Delta_x g(x)=\mu(x,W-1)-\mu(x,W).
\]
Since $\mu(x,W)\ge0$, we get $\mu(x,W-1)\ge c-F_{W-1}(x)+\Delta_x g(x)$. The lower constraint
gives $F_{W-1}(x)\le\tau$, and (a) gives $\Delta_x g(x)>0$. Hence
\[
\mu(x,W-1)>c-\tau>0 .
\]
By the third rule for the multipliers, $\lambda^{+}(x,W-1)>0$, and so $F_{W-1}(x)=\tau$. Bond
$W-1$ is therefore up-tight.

(c) By (b), bond $W-1$ lies in a lower up-run of column $x$. Since $W-1$ is the top lower bond,
this run has the form $[p,W-1]$. By the model lemma, Lemma~\ref{8.13:lem-floating-runs}, applied
in the lower tier, there are two cases: either $p=0$, or the run hangs from the top of the lower
tier with $p\ge1$. Suppose $p\ge1$. Then the model lemma,
Lemma~\ref{8.13:lem-tent-poles-concave}, applied to the lower up-run $[p,W-1]$, gives
$\Delta_x g(x)<0$. This contradicts (a). Hence $p=0$.

All lower bonds of column $x$ then have difference $\tau$. Since $\psi(x,0)=\varphi_0(x)$, we get
$\psi(x,y)=\varphi_0(x)+\tau y$ for $0\le y\le W$. In particular $g(x)=\varphi_0(x)+\tau W$.

(d) In the columns $x\pm1$ every lower bond has difference at most $\tau$, and the bottom value
is $\varphi_0(x\pm1)$. Hence $f(x\pm1,W)\le0$. Write $g(x\pm1)=\varphi_0(x\pm1)+\tau W+f(x\pm1,W)$
and use (c) for $g(x)$; the terms $\tau W$ cancel in $\Delta_x$, and
\[
\Delta_x g(x)=\Delta_x\varphi_0(x)+f(x+1,W)+f(x-1,W)\le\Delta_x\varphi_0(x).
\]
By (a), $\Delta_x\varphi_0(x)>0$.

It remains to identify $x$. Let $\varphi_0(x)=\sigma\|x\|$ and $\|x\|\le P/2$, where $N=2P$. For
$\|x\|\le P/2$ the function $\|\cdot\|$ is affine on $\{x-1,x,x+1\}$ unless $x\equiv0\pmod N$,
where $\Delta_x\|x\|=2$. So $\Delta_x\|x\|=0$ off the tip column, and $\Delta_x\varphi_0(x)>0$
forces $x\equiv0\pmod N$.
\end{proof}

\begin{remark}
Lemma~\ref{8.13:lem-anchored-runs} is the analogue, one tier up, of the corresponding statement
for the lower band alone. In the model with the lower band alone the sign of $\Delta_x\varphi_0$
is read off the datum directly. Here the structure of the lower tier supplies it, through the two
model lemmas, Lemmas~\ref{8.13:lem-floating-runs} and~\ref{8.13:lem-tent-poles-concave}.

The lemma says that the lattice admits no young constraint anchored at the face of the clamp
anywhere off the tip of the kink, whatever the values of the caps. In particular:
\begin{itemize}
\item there is none on the columns of the block of the model other than the tip column, since on
those columns $\Delta_x g=\Delta_x\varphi_0=0$;
\item there is none above the fringe of the model, on whose columns the lower runs reaching the
top lower bond hang from it as tent poles; the tops of these poles are laterally concave by the
model lemma, Lemma~\ref{8.13:lem-tent-poles-concave};
\item there is none above the fan of the model, on whose columns the multiplier of the top lower
bond vanishes, $\mu(x,W-1)=0<c-\tau$.
\end{itemize}

In the continuum the corresponding statement is false: active regions anchored at the base,
with multiplier vanishing at their free boundary, are allowed there. What forbids them on the
lattice is the strict inequality $|F_q(x)|<c$ at the end of the run.

Suppose moreover that the tangential slope is at most $\sigma$ everywhere, with $\sigma$ much
smaller than $c$. Then a ramp, that is, a trace whose tangential slope attains the young cap
over a lateral stretch, is infeasible in the model with the young band. Of that mechanism only
the kink's own logarithmic singularity survives, and it can act only through items (b)--(d) of
the model lemma on that singularity stated later in this section.
\end{remark}

In the following lemma and its proof, $\psi=\psi^{(c)}$ is the minimiser of the Dirichlet energy
$E(\psi)$ over the feasible set $\mathcal{K}_c$ of the model instance, on the periodic strip
$\Lambda=\mathbb{Z}_N\times\{0,\dots,H\}$, and $F_y(x)=\psi(x,y+1)-\psi(x,y)$. The proof uses
the following constraints of the feasible set:
the lower bonds $0\le y\le W-1$ of every column satisfy $|F_y|\le\tau$;
the young bonds $W\le y\le W'-1$, where $W'=W+\bar W$, satisfy $|F_y|\le c$;
and $\psi(\cdot,0)=\varphi_0$. Beyond them it uses two further properties of the model instance:
the lateral slope bound $|\psi(x\pm1,W)-\psi(x,W)|\le\sigma$, and the lower bound
$\psi\ge\varphi_0$ on the columns with $\|x\|\le P/2$. It also uses the optimality conditions
and two model lemmas.

We write $g(x)=\psi(x,W)$ for the exit trace. The following objects are those used in the model
lemmas, Lemma~\ref{8.13:lem-opposite-bonds} and Lemma~\ref{8.13:lem-tent-poles-concave}:
\begin{itemize}
\item the lateral second difference $\Delta_x$;
\item the multipliers $\mu(x,y)=\lambda^{+}(x,y)-\lambda^{-}(x,y)$ of the vertical bonds.
Here $\lambda^{+}(x,y)>0$ only if bond $y$ of column $x$ is at its upper cap, and
$\lambda^{-}(x,y)>0$ only if it is at its lower cap;
\item the optimality condition at a site;
\item the up- and down-runs of a column. These are maximal strings of consecutive bonds of one
tier at the upper, respectively lower, cap of that tier.
\end{itemize}
Thus a young down-run $[p,q-1]$ of column $x$ has $F_y(x)=-c$ for $p\le y\le q-1$.

\begin{lemma}[Model: no anchored young down-run]\label{8.13:lem-anchored-down-run}
The model is the lattice variational problem of the model instance whose one-sided conditions
are fixed in Definition~\ref{8.13:def-model-conditions}, with minimiser $\psi=\psi^{(c)}$ and
the conventions above. Then, for $\|x\|\le P/2$, no young down-run $[p,q-1]$ of column $x$
has $p=W$.
\end{lemma}

\begin{proof}
Suppose that $[W,q-1]$ is a young down-run of column $x$.

\emph{The bond below the run.} At $(x,W)$ we have $F_W=-c$ and $\psi(x,W)=g(x)$, so the
optimality condition reads
\begin{equation*}
(-c-F_{W-1})+\Delta_x g(x)=\mu(x,W-1)-\mu(x,W).
\end{equation*}
Since $c>0$, bond $W$ is not at its upper cap. Hence $\lambda^{+}(x,W)=0$ and
$\mu(x,W)\le0$. We use $-F_{W-1}\le\tau$ and the bound $\Delta_x g(x)\le2\sigma$. The latter
holds because, by the lateral slope bound of the model instance, the lateral differences
$|\psi(x\pm1,W)-\psi(x,W)|$ are at most $\sigma$. With these,
\begin{equation*}
\mu(x,W-1)\le -c-F_{W-1}+\Delta_x g(x)\le -c+\tau+2\sigma<0.
\end{equation*}
The last inequality is the condition $c>\tau+2\sigma$ of the model instance, displayed in
\eqref{8.13:eq-model-conditions-d} of Definition~\ref{8.13:def-model-conditions}. Thus
$\lambda^{-}(x,W-1)>0$, and bond $W-1$ is at its lower cap: $F_{W-1}(x)=-\tau$.

Let $[p,W-1]$ be the lower down-run containing bond $W-1$.

Suppose first that $p=0$. Then the column falls by $\tau$ at each step from the bottom row, so
\begin{equation*}
\psi(x,W)=\varphi_0(x)-\tau W<\varphi_0(x).
\end{equation*}
This contradicts the lower bound $\psi\ge\varphi_0$ of the model instance on the columns with
$\|x\|\le P/2$.

Hence $p\ge1$. We apply the model lemma, Lemma~\ref{8.13:lem-tent-poles-concave}, to the lower
down-run $[p,W-1]$ with $p\ge1$. It gives
\begin{equation}\label{8.13:eq-anchored-down-run-1}
\Delta_x g(x)>0 .
\end{equation}

\emph{Case $q\le W'-1$.} Here bond $q$ is a young bond. By maximality of the run,
$F_q\ne-c$. Bonds $q-1$ and $q$ are adjacent constrained bonds, and $F_{q-1}=-c$. So the model
lemma, Lemma~\ref{8.13:lem-opposite-bonds}, excludes $F_q=+c$. Hence $|F_q|<c$, so
$\mu(x,q)=0$ and $F_q+c>0$. Bond $q-1$ is at its lower cap, so $\lambda^{+}(x,q-1)=0$ and
$\mu(x,q-1)\le0$.

The optimality condition at $(x,q)$ is therefore
\begin{equation*}
(F_q+c)+\Delta_x\psi(x,q)=\mu(x,q-1)\le 0 .
\end{equation*}
Since $F_q+c>0$, this gives $\Delta_x\psi(x,q)<0$.

On the other hand, along the run $\psi(x,q)=g(x)-c(q-W)$. The young bonds $W,\dots,q-1$ of
the columns $x\pm1$ are each at least $-c$, so $\psi(x\pm1,q)\ge g(x\pm1)-c(q-W)$. Hence
\begin{equation*}
\Delta_x\psi(x,q)
=\psi(x+1,q)+\psi(x-1,q)-2\psi(x,q)
\ge\Delta_x g(x)>0,
\end{equation*}
by \eqref{8.13:eq-anchored-down-run-1}. This is a contradiction.

\emph{Case $q=W'$.} Here the run fills the young band. The lower bonds of column $x$ are each
at most $\tau$, so $g(x)\le\varphi_0(x)+\tau W$. Therefore
\begin{equation*}
\psi(x,W')=g(x)-c\bar W\le\varphi_0(x)+\tau W-c\bar W<\varphi_0(x).
\end{equation*}
The last inequality is the condition $c\bar W>\tau W$ of the model instance in
\eqref{8.13:eq-model-conditions-d}. This contradicts the lower bound $\psi\ge\varphi_0$ of the
model instance on $\|x\|\le P/2$.

Both cases are impossible, which proves the lemma.
\end{proof}

\begin{lemma}[Model: census of young-tight bonds]\label{8.13:lem-tight-census}
This lemma concerns the model of this section: the minimiser $\psi=\psi^{(c)}$ of the Dirichlet
energy $E(\psi)$ over the feasible set $\mathcal{K}_c$ with young cap $c$, on the periodic strip
$\Lambda=\mathbb{Z}_N\times\{0,\dots,H\}$, $N=2P$. Its bottom row carries the Dirichlet datum
$\varphi_0$. The vertical bond differences $F_y(x)=\psi(x,y+1)-\psi(x,y)$ of the lower tier
$0\le y\le W-1$ are bounded by $\tau$ in absolute value. The bonds of the bottom young band
$W\le y\le W'-1$, $W'=W+\bar W$ (and of its mirror image at the top), are bounded by the young
cap $c$.

A young bond is \emph{tight} if its constraint is saturated. A young vertical bond is up-tight
if $F_y(x)=c$ and down-tight if $F_y(x)=-c$. A young \emph{up-run} (\emph{down-run}) of
column $x$ is a maximal interval $[p,q-1]$ of rows, $W\le p<q\le W'$, on which every bond is
up-tight (down-tight). Column $x$ is \emph{young-full} (up-full, down-full) if it carries the
young up-run (down-run) $[W,W'-1]$. It is \emph{lower-full} if every bond of its lower tier is
at the cap, $F_y(x)=\tau$ for $0\le y\le W-1$, that is, if
$\psi(x,y)=\varphi_0(x)+\tau y$ for $0\le y\le W$.

Then no young horizontal bond is tight (the lateral differences satisfy
$|\nabla_x\psi|\le\sigma<c$). Near the bottom
tip, that is for $\|x\|\le P/2$, every young-tight vertical bond of the bottom band belongs to
exactly one of the following:
\begin{itemize}
\item[$(\mathrm{N})$] a tip needle: an up-run $[W,W+m^{*})$ with $m^{*}\le\bar W-1$, on a
lower-full column $x\in\mathcal{T}$ (in the lattice variant of the model whose tip set is the
tip column of the kink, $x=0$);
\item[$(\mathrm{F})$] a young-full up-column; then $F_{W'}(x)\ge c-2\sigma$;
\item[$(\mathrm{P}^{+})$] a hanging up-pole $[p,W'-1]$ with $p\ge W+1$; then $F_{W'}(x)>c$;
\item[$(\mathrm{P}^{-})$] a hanging down-pole $[p,W'-1]$ with $p\ge W+1$; then $F_{W'}(x)<-c$.
\end{itemize}
Near the top tip the same holds with up and down exchanged: the needles are down-runs on
$P+\mathcal{T}$, and the full columns are down-full. The statement for the top band is the
mirror image of the above in the vertical direction.
\end{lemma}

\begin{proof}
A young horizontal bond is not tight: by the lateral slope bound of the model the lateral
differences satisfy $|\nabla_x\psi|\le\sigma$, and $\sigma<c$ because the caps of the model
satisfy $c>\tau+2\sigma$.

Let $\|x\|\le P/2$. A tight vertical bond of the bottom band in column $x$ lies in exactly one
maximal run $[p,q-1]$, which is either an up-run or a down-run.

Consider first up-runs. By Lemma~\ref{8.13:lem-floating-runs}, a lemma of the same model, every
young up-run has $p=W$ or $q=W'$. Three cases remain.
\begin{itemize}
\item $p=W$ and $q<W'$. The column carries an up-run $[W,q-1]$ with $q\le W'-1$ and
$\|x\|\le P/2$. These are exactly the hypotheses of Lemma~\ref{8.13:lem-anchored-runs}, a lemma
of the same model. By its parts (c) and (d), column $x$ is lower-full and
$x\in\mathcal{T}$ (in the lattice variant of the model named in the statement, $x=0$). Writing
$q=W+m^{*}$, the condition $q\le W'-1$ reads $m^{*}\le\bar W-1$. This is case $(\mathrm{N})$.
\item $p=W$ and $q=W'$. Column $x$ is young-full. Part (c) of
Lemma~\ref{8.13:lem-hanging-runs}, a lemma of the same model, gives
\begin{equation*}
F_{W'}(x)\ge c-\Delta_x\psi(x,W')\ge c-2\sigma .
\end{equation*}
This is case $(\mathrm{F})$.
\item $p>W$ and $q=W'$. The run is $[p,W'-1]$ with $p\ge W+1$. Part (a) of
Lemma~\ref{8.13:lem-hanging-runs}, a lemma of the same model, gives $F_{W'}(x)>c$. This is case
$(\mathrm{P}^{+})$.
\end{itemize}

Consider next down-runs. The case $p=W$ is excluded by
Lemma~\ref{8.13:lem-anchored-down-run}, a lemma of the same model, since $\|x\|\le P/2$.

If $p>W$ and $q=W'$, the run is $[p,W'-1]$ with $p\ge W+1$. Part (b) of
Lemma~\ref{8.13:lem-hanging-runs}, a lemma of the same model, gives $F_{W'}(x)<-c$. This is
case $(\mathrm{P}^{-})$.

The remaining case, a floating down-run with $p>W$ and $q<W'$, is excluded by
Lemma~\ref{8.13:lem-floating-runs}, a lemma of the same model.

The four cases are told apart by the sign of the run, by whether $p=W$, and by whether $q=W'$.
Since the maximal run containing a given bond is unique, the bond belongs to exactly one case.

The statements near the top tip and for the top band follow from the bottom-tip statement by
the symmetries of the model: its invariance exchanging the two tips (columns near $0$ with
columns near $P$) together with up- and down-runs, and its invariance under the reflection of
the strip in the vertical direction, which exchanges the two bands.
\end{proof}

Consequently, imposing the young cap $c$ on the band $W\le y\le W'-1$ can make a young
constraint active beyond the tip needle
only through $(\mathrm{F})$ or $(\mathrm{P}^{\pm})$. Case $(\mathrm{F})$ is a young-full
column, along which the kink is transported rigidly through the whole band. Cases
$(\mathrm{P}^{\pm})$ are poles hanging from the top of the band. Each of these cases requires
the first free row above the band to be pulled by at least $c-2\sigma$, that is,
$|F_{W'}(x)|\ge c-2\sigma$.

\begin{definition}[Model: the straddling-only minimiser above the clamp]\label{8.13:not-straddling-minimiser}
Everything in this definition takes place in a model, not in Ba{\l}aban's step. The model is the
two-tier box-clamp lattice model of Definition~\ref{8.13:def-two-tier-model}, which is a model
statement, on the periodic strip $\Lambda=\mathbb{Z}_N\times\{0,1,\dots,H\}$ with $N=2P$. We take
the instance of that model that satisfies the one-sided conditions of
Definition~\ref{8.13:def-model-conditions}, also a model statement. In it, $\tilde{\psi}$ is the
straddling-only minimiser of the model. We write $\tilde g:=\tilde{\psi}(\cdot,W)$ for its trace on
row $W$, and for every row $y$ with $0\le y\le H-1$ we put
$\tilde F_y(x):=\tilde{\psi}(x,y+1)-\tilde{\psi}(x,y)$.

The \emph{free strip} of $\tilde{\psi}$ consists of the rows $W,\dots,H-W$. Its height is
$H':=H-2W$. By the condition \eqref{8.13:eq-model-conditions-d} of the model instance,
\begin{equation*}
H'\ge 2\bar W+4 .
\end{equation*}
By the harmonicity of the straddling-only minimiser in its free strip, a model statement of this
chapter, the field $\tilde{\psi}$ is harmonic for the four-neighbour Laplacian $\Delta$ at every
site of the rows $W+1,\dots,H-W-1$, and it is symmetric under the reflection $y\mapsto H-y$ of the
strip, so that $\tilde{\psi}(\cdot,H-W)=\tilde g$ as well.

We apply the strip bound above a Lipschitz trace with $Y=W$. It identifies the function
$(x,m)\mapsto\tilde{\psi}(x,W+m)$, $0\le m\le H'$, with the strip harmonic extension
$\mathcal{U}_{N,H'}[\tilde g,\tilde g]$, whose data are symmetric. The kernel representation of
vertical differences, applied to this extension, then gives the exact identity
\begin{equation}\label{8.13:eq-straddling-minimiser-1}
\tilde F_{W+m}(x)=\bigl((\nabla\tilde g)\ast K^{R}_m\bigr)(x)
:=\sum_{z}\nabla\tilde g(z)\,K^{R}_m(x-z),\qquad 0\le m\le H'-1 .
\end{equation}
Here $\nabla\tilde g(z):=\tilde g(z+1)-\tilde g(z)$. The functions $K^{R}_m$ are the real kernels of
the strip $\mathfrak{S}_{N,H'}$ given by that representation, and they satisfy
$K^{R}_m(1-z)=-K^{R}_m(z)$.

The logarithmic bound in the strip bound above a Lipschitz trace gives, for all $x$ and
$0\le m\le H'-1$,
\begin{equation}\label{8.13:eq-straddling-minimiser-a}
\bigl|\tilde F_{W+m}(x)\bigr|\le b(m)
:=\sigma\Bigl[C_K^{0}+6\log^{+}\Bigl(\frac{N}{2m_*+1}\Bigr)\Bigr].
\end{equation}
In this bound $m_*:=\min(m,\,H'-1-m)$ and $C_K^{0}=6+2\sqrt{3}$.

Finally we put $T(x):=\|x\|=\operatorname{dist}(x,N\mathbb{Z})$, so that $\varphi_0=\sigma T$. The
difference $\nabla T(z):=T(z+1)-T(z)$ then takes the values
\begin{equation*}
\nabla T(z)=
\begin{cases}
+1, & 0\le z\le P-1,\\
-1, & -P\le z\le -1,
\end{cases}
\end{equation*}
with $z$ read modulo $N$. For $0\le m\le H'-1$ we put
\begin{equation}\label{8.13:eq-straddling-minimiser-2}
D_T(x,m):=\bigl((\nabla T)\ast K^{R}_m\bigr)(x)=\sum_{z}\nabla T(z)\,K^{R}_m(x-z).
\end{equation}
\end{definition}

\begin{lemma}[Model: the kink pulls upward near the tip]\label{8.13:lem-kink-pull-sign}
The model is the one fixed in Definition~\ref{8.13:not-straddling-minimiser}, itself a statement of
the model. It consists of the periodic strip $\mathfrak{S}_{N,H'}$ of period $N=2P$ and height
$H'=H-2W$, with rows numbered $0,\dots,H'$; the datum $T(x)=\|x\|$, the distance from $x$ to
$N\mathbb{Z}$; and the pull $D_T(x,m)=((\nabla T)\ast K^{R}_m)(x)$, $0\le m\le H'-1$, where
$K^{R}_m$ are the kernels of the strip representation. In this model, for all $x$ with
$\|x\|\le P/2$ and all $m$ with $0\le m\le \lfloor H'/2\rfloor-1$,
\[
D_T(x,m)\ge 0 .
\]
\end{lemma}

\begin{proof}
We write $\Delta$ for the four-neighbour Laplacian on the strip,
\[
\Delta f(x,y)=f(x+1,y)+f(x-1,y)+f(x,y+1)+f(x,y-1)-4f(x,y),
\]
and $\Delta_x$ for the lateral second difference, $\Delta_x f(x)=f(x+1)+f(x-1)-2f(x)$.

Let $\Psi:=\mathcal{U}_{N,H'}[T,T]$ be the strip harmonic extension of the symmetric datum.
It is the unique function on the strip with $\Delta\Psi=0$ at every site of rows $1,\dots,H'-1$
and with $\Psi(x,0)=\Psi(x,H')=T(x)$ for all $x$. For $0\le m\le H'-1$ put
\[
F_T(x,m):=\Psi(x,m+1)-\Psi(x,m).
\]
By the kernel representation of vertical differences in the strip, applied to the datum $T$
on both boundary rows, $D_T(x,m)=F_T(x,m)$. It therefore suffices to prove $F_T\ge0$ on the
stated set.

\emph{The extension lies above the datum.}
The function $T-P/2$ is odd about $x=P/2$, that is,
\[
T(P/2+u)+T(P/2-u)=P\qquad\text{for every integer }u .
\]
To see this, put $s=P/2+u$. Then $T(P-s)$ is the distance from $s$ to $P+2P\mathbb{Z}$, while
$T(s)$ is the distance from $s$ to $2P\mathbb{Z}$. The point $s$ lies between two consecutive
multiples of $P$, one in each of these sets, so the two distances add up to $P$.

Consequently the function $(x,y)\mapsto P-\Psi(P-x,y)$ is harmonic on rows $1,\dots,H'-1$ and
equals $T(x)$ on rows $0$ and $H'$. By uniqueness it coincides with $\Psi$, so
$\Psi(P/2+u,y)+\Psi(P/2-u,y)=P$; at $u=0$ this gives $\Psi(P/2,y)=P/2$.

Since $T$ is even, $(x,y)\mapsto\Psi(-x,y)$ has the same boundary data as $\Psi$. Hence it
equals $\Psi$, and $\Psi(-P/2,y)=P/2$ as well. Thus
\[
\Psi(\pm P/2,y)=P/2=T(\pm P/2)\qquad\text{for all }y .
\]

Let $\Lambda_1:=\{(x,y):\|x\|\le P/2,\ 0\le y\le H'\}$. Its boundary sites are those of rows $0$
and $H'$ and of the columns $x\equiv\pm P/2$. Its interior sites are those with
$\|x\|\le P/2-1$ and $1\le y\le H'-1$, and all four neighbours of an interior site lie in
$\Lambda_1$.

Put $\Psi^{\circ}:=\Psi-T$, with $T$ regarded as constant in $y$. By the boundary data and the
identity $\Psi(\pm P/2,y)=T(\pm P/2)$, $\Psi^{\circ}$ vanishes at every boundary site of
$\Lambda_1$.

At an interior site, $\Delta\Psi^{\circ}=\Delta\Psi-\Delta T=-\Delta_xT(x)$, since $T$ does not
depend on $y$. On the representative interval $[-P/2,P/2]$ we have $T(x)=|x|$, so
$\Delta_xT(x)=2\cdot\mathbf{1}_{x\equiv0}$ for $\|x\|\le P/2-1$. Hence
$\Delta\Psi^{\circ}=-2\cdot\mathbf{1}_{x\equiv0}\le0$ at the interior sites: $\Psi^{\circ}$ is
superharmonic there.

Now let the minimum of $\Psi^{\circ}$ over $\Lambda_1$ be attained at an interior site.
Superharmonicity gives $4\Psi^{\circ}$ at that site $\ge$ the sum of the four neighbouring
values, and each of these is at least the minimum. So all four equal the minimum, in particular
the right-hand neighbour. Repeating with increasing $x$ reaches the column $x\equiv P/2$ in
finitely many steps. So the minimum is attained at a boundary site, where $\Psi^{\circ}=0$.
Therefore $\Psi^{\circ}\ge0$ on $\Lambda_1$.

\emph{Boundary values of the vertical differences.}
For $\|x\|\le P/2$ we have
\[
F_T(x,0)=\Psi(x,1)-T(x)=\Psi^{\circ}(x,1)\ge0 .
\]
Moreover, since the columns $x\equiv\pm P/2$ are flat, $F_T(\pm P/2,m)=P/2-P/2=0$ for
$0\le m\le H'-1$.

\emph{The discrete minimum argument.}
For $1\le m\le H'-2$ both rows $m$ and $m+1$ lie among the rows $1,\dots,H'-1$, where $\Psi$ is
harmonic. Hence $\Delta F_T(x,m)=\Delta\Psi(x,m+1)-\Delta\Psi(x,m)=0$, that is,
\begin{equation}\label{8.13:eq-kink-pull-sign-1}
4F_T(x,m)=F_T(x+1,m)+F_T(x-1,m)+F_T(x,m+1)+F_T(x,m-1),\qquad 1\le m\le H'-2 .
\end{equation}

The function $(x,y)\mapsto\Psi(x,H'-y)$ has the same boundary data as $\Psi$, so it equals
$\Psi$ by uniqueness. Therefore, for $0\le m\le H'-1$,
\begin{equation}\label{8.13:eq-kink-pull-sign-2}
F_T(x,H'-1-m)=\Psi(x,m)-\Psi(x,m+1)=-F_T(x,m).
\end{equation}

Let $\bar m:=\lfloor H'/2\rfloor$. Definition~\ref{8.13:not-straddling-minimiser}, a statement of
the model, gives $H'\ge2\bar W+4\ge4$, so $\bar m\ge2$. Put
\[
\Lambda_2:=\{(x,m):\|x\|\le P/2,\ 0\le m\le\bar m-1\},\qquad
\Lambda_2^{\circ}:=\{(x,m):\|x\|\le P/2-1,\ 1\le m\le\bar m-1\}.
\]
The sites of $\Lambda_2\setminus\Lambda_2^{\circ}$ lie in row $0$ or in the columns
$x\equiv\pm P/2$, and there $F_T\ge0$ by the boundary values just established.

Suppose that $\min\{F_T(x,m):(x,m)\in\Lambda_2^{\circ}\}=:-a<0$, attained at $(x_0,m_0)$. Then
$F_T\ge-a$ on all of $\Lambda_2$.

Consider first the case $m_0\le\bar m-2$. The four neighbours of $(x_0,m_0)$ lie in
$\Lambda_2$, and \eqref{8.13:eq-kink-pull-sign-1} applies because
$1\le m_0\le\bar m-2\le H'-2$. Since each neighbour is $\ge-a$ and their sum is $-4a$, all four
equal $-a$. The values on $\Lambda_2\setminus\Lambda_2^{\circ}$ are $\ge0>-a$, so none of the
neighbours lies there. In particular $(x_0,m_0+1)$ is again a minimiser in $\Lambda_2^{\circ}$.
Repeating this step, one reaches row $\bar m-1$. In every case there is therefore a point
$(x_1,\bar m-1)\in\Lambda_2^{\circ}$, namely $x_1=x_0$, with $F_T(x_1,\bar m-1)=-a$.

At this point \eqref{8.13:eq-kink-pull-sign-1} applies, since $1\le\bar m-1\le H'-2$. Three of
its neighbours, $(x_1\pm1,\bar m-1)$ and $(x_1,\bar m-2)$, lie in $\Lambda_2$ and carry values
$\ge-a$. The fourth neighbour, $(x_1,\bar m)$, is handled by \eqref{8.13:eq-kink-pull-sign-2}:
\begin{itemize}
\item If $H'$ is odd, then $H'=2\bar m+1$. Taking $m=\bar m$ in
\eqref{8.13:eq-kink-pull-sign-2} gives $F_T(x_1,\bar m)=-F_T(x_1,\bar m)$, so
$F_T(x_1,\bar m)=0$.
\item If $H'$ is even, then $H'=2\bar m$. Taking $m=\bar m-1$ in
\eqref{8.13:eq-kink-pull-sign-2} gives $F_T(x_1,\bar m)=-F_T(x_1,\bar m-1)=a$.
\end{itemize}
In both cases $F_T(x_1,\bar m)\ge0$, and so
\begin{equation}\label{8.13:eq-kink-pull-sign-3}
-4a=4F_T(x_1,\bar m-1)=\sum_{\text{four neighbours}}F_T\ge3(-a)+0 ,
\end{equation}
that is, $a\le0$. This is a contradiction.

Hence $F_T\ge0$ on $\Lambda_2^{\circ}$, and $F_T\ge0$ on $\Lambda_2\setminus\Lambda_2^{\circ}$ by
the boundary values. Thus $F_T(x,m)\ge0$ for $\|x\|\le P/2$ and $0\le m\le\bar m-1$. Since
$D_T=F_T$, the lemma follows.
\end{proof}

\begin{lemma}[Model: decay of the kink's pull]\label{8.13:lem-kink-pull-decay}
The model is the strip model fixed in Notation~\ref{8.13:not-straddling-minimiser}, itself a
statement made within that model and not about the lattice theory: the periodic strip
$\mathbb{Z}_N\times\{0,\dots,H\}$, of period $N=2P$, together
with the free strip $\mathfrak{S}_{N,H'}$ of the straddling-only minimiser $\tilde{\psi}$, of height
$H'=H-2W$. Write $K^{R}_m$, $0\le m\le H'-1$, for the real kernels of $\mathfrak{S}_{N,H'}$ through
which vertical differences of strip harmonic extensions with symmetric data are represented. Put
$T(x)=\|x\|$, the distance from $x$ to $N\mathbb{Z}$, and $\nabla T(z)=T(z+1)-T(z)$. Define
$D_T(x,m)=((\nabla T)\ast K^{R}_m)(x)$, the convolution being taken over one period. In this lemma
$N$ denotes the period and $m$ a row index.

Then for every $x\in\mathbb{Z}_N$ with $\|x\|\le P/2$ and every integer $m$ with
$0\le m\le (H'-1)/2$,
\begin{equation*}
|D_T(x,m)|\le C_T+12\,\log^{+}\!\Big(\frac{N}{2\max(\|x\|,m)+1}\Big),
\end{equation*}
where
\begin{equation}\label{8.13:eq-kink-pull-decay-a}
C_T:=2C_K^{0}+18+12\log 3 ,
\end{equation}
and $C_K^{0}$ is the constant in the $\ell^1$ bound for the kernels recalled in the proof.
\end{lemma}

\begin{proof}
Fix $m$ as in the statement and write $K:=K^{R}_m$, regarded as an $N$-periodic function on
$\mathbb{Z}$. We use two facts about these kernels. The first comes from the kernel representation
of vertical differences in the strip $\mathfrak{S}_{N,H'}$: $K$ is real and odd about $\tfrac12$,
\begin{equation*}
K(1-z)=-K(z).
\end{equation*}
The second consists of three of the kernel bounds, which hold for $0\le y\le (H'-1)/2$ and hence for
$y=m$:
\begin{itemize}
\item[(b)] $|K^{R}_y(z)|\le 3/\|z-\tfrac12\|$;
\item[(c)] $\|K^{R}_y\|_{\ell^1}\le C_K^{0}+6\log^{+}(N/(2y+1))$;
\item[(f)] $\sum_{\|z-\frac12\|>\rho}|K^{R}_y(z)|\le 6[1+\log^{+}(N/(2\rho))]$ for every
half-integer $\rho\ge\tfrac12$.
\end{itemize}

\emph{Reduction to $0\le x\le P/2$.} Since $T$ is even, $\nabla T(-z-1)=-\nabla T(z)$. Substituting
$z=-z'-1$ and then using the oddness of $K$ gives
\begin{equation*}
D_T(-x,m)=\sum_{z'}\big(-\nabla T(z')\big)K\big(1-(x-z')\big)=\sum_{z'}\nabla T(z')K(x-z')=D_T(x,m).
\end{equation*}
So $D_T(\cdot,m)$ is even, and we may take $0\le x\le P/2$.

\emph{Rewriting.} On the period $-P\le z\le P-1$ we have $\nabla T=+1$ for $0\le z\le P-1$ and
$\nabla T=-1$ for $-P\le z\le -1$. With $\zeta=x-z$ this gives
\begin{equation*}
D_T(x,m)=\sum_{z=0}^{P-1}K(x-z)-\sum_{z=-P}^{-1}K(x-z)
=\sum_{\zeta=x-P+1}^{x}K(\zeta)-\sum_{\zeta=x+1}^{x+P}K(\zeta).
\end{equation*}
Because $2x\le P$, the first sum splits over $\zeta\in[x-P+1,-x]$ and $\zeta\in[1-x,x]$. The map
$\zeta\mapsto 1-\zeta$ sends $[1-x,x]$ onto itself and changes the sign of $K$. Hence the sum over
$[1-x,x]$ equals its own negative and vanishes. On $[x-P+1,-x]$ the substitution $\zeta'=1-\zeta$
gives the range $[x+1,P-x]$, so this part equals $-\sum_{\zeta'=x+1}^{P-x}K(\zeta')$. We split the
second sum at $P-x$ and obtain
\begin{equation}\label{8.13:eq-kink-pull-decay-1}
D_T(x,m)=-2\sum_{\zeta=x+1}^{P-x}K(\zeta)-\sum_{\zeta=P-x+1}^{P+x}K(\zeta).
\end{equation}

\emph{The last sum.} It has $2x$ terms. For $\zeta\in[P-x+1,P+x]$ the point $\zeta-\tfrac12$ lies at
distance at least $P-x+\tfrac12$ from both $0$ and $2P=N$. Therefore
$\|\zeta-\tfrac12\|\ge P-x+\tfrac12\ge P/2$. By (b), each term has modulus at most $6/P$, and the
sum has modulus at most $12x/P\le 6$.

\emph{The main sum.} For $\zeta\in[x+1,P-x]$ one has $\|\zeta-\tfrac12\|\ge x+\tfrac12$. These
$\zeta$ are distinct modulo $N$. Since $N$ is even, $\|\zeta-\tfrac12\|$ is a half-integer, so the
condition $\|\zeta-\tfrac12\|\ge x+\tfrac12$ is the same as $\|\zeta-\tfrac12\|>x-\tfrac12$. Hence
\begin{equation*}
\Big|\sum_{\zeta=x+1}^{P-x}K(\zeta)\Big|\le\sum_{\|\zeta-\frac12\|\ge x+\frac12}|K(\zeta)|
\le\min\Big\{\|K\|_{\ell^1},\ \sum_{\|\zeta-\frac12\|>x-\frac12}|K(\zeta)|\Big\}.
\end{equation*}
By (c), and by (f) with $\rho=x-\tfrac12$ (admissible when $x\ge1$), this minimum is at most
\begin{equation*}
\min\Big\{C_K^{0}+6\log^{+}\tfrac{N}{2m+1},\ 6+6\log^{+}\tfrac{N}{2x-1}\Big\}.
\end{equation*}
For $x=0$ only the first entry is used. For $x\ge1$ we have $(2x+1)/(2x-1)\le3$ and
$\log^{+}(ab)\le\log^{+}a+\log^{+}b$. Hence
\begin{equation*}
\log^{+}\tfrac{N}{2x-1}\le\log 3+\log^{+}\tfrac{N}{2x+1}.
\end{equation*}

\emph{Conclusion.} Put $\ell_m=\log^{+}(N/(2m+1))$ and $\ell_x=\log^{+}(N/(2x+1))$. For $a,b\ge0$,
$\min\{a+\alpha,b+\beta\}\le\max\{a,b\}+\min\{\alpha,\beta\}$, and $\max\{a,b\}\le a+b$. We apply
this with $a=C_K^{0}$ and $b=6+6\log3$ to the main sum. Inserting the result and the bound on the
last sum into \eqref{8.13:eq-kink-pull-decay-1} gives
\begin{equation*}
|D_T(x,m)|\le 6+2\big(C_K^{0}+6+6\log 3\big)+12\min\{\ell_m,\ell_x\}.
\end{equation*}
The constant part equals $C_T$ of \eqref{8.13:eq-kink-pull-decay-a}. The function
$t\mapsto\log^{+}(N/(2t+1))$ is non-increasing, so
\begin{equation*}
\min\{\ell_m,\ell_x\}=\log^{+}\big(N/(2\max(x,m)+1)\big).
\end{equation*}
For $x=0$ the same bound follows from $2(C_K^{0}+6\ell_m)\le C_T+12\ell_m$. This proves the lemma.
Item (a) of the kernel bounds, $|K^{R}_y(z)|\le\min(1,\sqrt3/(2y+1))$, would sharpen
$\max(\|x\|,m)$ to a joint norm of $(x,m)$; this is not needed.
\end{proof}

\begin{lemma}[Model: two-sided pull bound near the tip]\label{8.13:lem-tip-pull-bound}
This lemma concerns the strip model of Notation~\ref{8.13:not-straddling-minimiser}. The model is
the Dirichlet energy on the periodic strip $\mathbb{Z}_N\times\{0,\dots,H\}$, $N=2P$, with kinked
datum $\varphi_0(x)=\sigma\|x\|$; in it $\tilde{\psi}$ is the straddling-only minimiser,
$\tilde g=\tilde{\psi}(\cdot,W)$ is its exit trace, $\tilde F_y(x)=\tilde{\psi}(x,y+1)-\tilde{\psi}(x,y)$,
$H'=H-2W$, and $K^{R}_m$ are the kernels of the free strip $\mathfrak{S}_{N,H'}$.
Let $\kappa\in(0,\tfrac12]$ with $\kappa W$ an integer and $\kappa W\ge 1$. Assume part~(i) of the
mild-fringe hypothesis, Definition~\ref{8.13:def-mild-fringe}, which is a model hypothesis:
$\tilde S\ge 2\kappa W$. Put $\beta_\kappa:=24\tau/\kappa$. Then for $\|x\|\le\kappa W$ and
$0\le m\le\lfloor H'/2\rfloor-1$,
\begin{equation}\label{8.13:eq-tip-pull-bound-a}
-\beta_\kappa\;\le\;\tilde F_{W+m}(x)\;\le\;
\sigma\Bigl[C_T+12\log^{+}\frac{N}{2\max(\|x\|,m)+1}\Bigr]+\beta_\kappa ,
\end{equation}
where $C_T$ is the constant of \eqref{8.13:eq-kink-pull-decay-a}.
\end{lemma}

\begin{proof}
Write $\nabla h(z)=h(z+1)-h(z)$ and $T(x)=\|x\|$, so that $\varphi_0=\sigma T$, $\nabla T=+1$ on
$0\le z\le P-1$ and $\nabla T=-1$ on $-P\le z\le -1$. Define the slope deficit $\tilde{\delta}$ on
$\mathbb{Z}_N$ by $\nabla\tilde g=\sigma\nabla T-\tilde{\delta}$.

\emph{Signs and symmetry of $\tilde{\delta}$.} The comparison property of the model minimisers,
stated with their symmetries (the property from which their lower bound $\psi\ge\varphi_0$ on
$\|x\|\le P/2$ is derived), gives $\tilde{\delta}\ge0$ on $0\le z<P/2$.
The minimiser $\tilde{\psi}$ is even in $x$. Hence
\[
\nabla\tilde g(-1-z)=\tilde g(z)-\tilde g(z+1)=-\nabla\tilde g(z).
\]
The same identity holds for $T$, and therefore $\tilde{\delta}(-1-z)=-\tilde{\delta}(z)$.
The map $z\mapsto-1-z$ sends $[0,P/2)$ onto $[-P/2,0)$, so $\tilde{\delta}\le0$ on
$-P/2\le z<0$. From $\tilde{\psi}(x+P,y)=\sigma P-\tilde{\psi}(x,y)$ we get
$\nabla\tilde g(z+P)=-\nabla\tilde g(z)$. Moreover $T(z+P)=P-T(z)$ gives
$\nabla T(z+P)=-\nabla T(z)$. Hence $\tilde{\delta}(z+P)=-\tilde{\delta}(z)$.

\emph{Total mass.} By the same symmetries the quarter columns are flat, so
$\tilde g(P/2)=\sigma P/2$. Telescoping over $0\le z<P/2$ gives
\[
\sum_{0\le z<P/2}\tilde{\delta}(z)=\frac{\sigma P}{2}-\bigl[\tilde g(P/2)-\tilde g(0)\bigr]
=\tilde g(0)-\varphi_0(0).
\]
This lies in $[0,\tau W]$. It is nonnegative by the lower bound $\tilde{\psi}\ge\varphi_0$ on
$\|x\|\le P/2$. It is at most $\tau W$ by the nonpositivity $f\le0$ of the lower lift deficit.
Since $\tilde{\delta}\ge0$ there, $\sum_{0\le z<P/2}|\tilde{\delta}(z)|\le\tau W$. The bijections
$z\mapsto-1-z$ from $[0,P/2)$ onto $[-P/2,0)$, and $z\mapsto z+P$ from $[-P/2,0)$ onto $[P/2,P)$
and from $[0,P/2)$ onto $[-P,-P/2)$, change only the sign of $\tilde{\delta}$. Thus $|\tilde{\delta}|$
has the same sum on each of the four quarters, and
\begin{equation}\label{8.13:eq-tip-pull-bound-1}
\sum_{z\in\mathbb{Z}_N}|\tilde{\delta}(z)|\le 4\tau W .
\end{equation}

\emph{Vanishing near the tips.} The lower-full columns in $[-P/2,P/2]$ are those of
$[-\tilde S,\tilde S]$, and at a lower-full column $\tilde g=\varphi_0+\tau W$. Let $\|z\|<2\kappa W$,
with representative $|z|\le2\kappa W-1$. Then $z$ and $z+1$ lie in $[-2\kappa W,2\kappa W]$, which
is contained in $[-\tilde S,\tilde S]$ by hypothesis~(i). Hence $\nabla\tilde g(z)=\sigma\nabla T(z)$,
that is, $\tilde{\delta}(z)=0$. If $\|z-P\|<2\kappa W$, then $\tilde{\delta}(z)=-\tilde{\delta}(z-P)=0$.
Note also that $2\kappa W\le\tilde S\le P/2$, so $\kappa W\le P/2$.

\emph{The deficit's pull.} Let $\|x\|\le\kappa W$ and $z\in\operatorname{supp}\tilde{\delta}$.
Then $\|z\|\ge2\kappa W$, so $\|x-z\|\ge\kappa W$, and consequently
$\|x-z-\tfrac12\|\ge\kappa W-\tfrac12>0$. For every such $z$ the lateral part of the kernel bound
for $K^{R}_m$ bounds $|K^{R}_m(x-z)|$ by $3/(\kappa W-\tfrac12)$. Combining this with
\eqref{8.13:eq-tip-pull-bound-1} and with $\kappa W-\tfrac12\ge\kappa W/2$ (as $\kappa W\ge1$):
\begin{equation}\label{8.13:eq-tip-pull-bound-2}
|(\tilde{\delta}*K^{R}_m)(x)|\le\frac{4\tau W\cdot 3}{\kappa W-\tfrac12}\le\frac{24\tau}{\kappa}
=\beta_\kappa .
\end{equation}

\emph{Conclusion.} By Notation~\ref{8.13:not-straddling-minimiser}, a model statement,
$\tilde F_{W+m}=(\nabla\tilde g)*K^{R}_m$ for $0\le m\le H'-1$. Inserting
$\nabla\tilde g=\sigma\nabla T-\tilde{\delta}$ and $D_T(x,m)=((\nabla T)*K^{R}_m)(x)$ gives
\[
\tilde F_{W+m}(x)=\sigma D_T(x,m)-(\tilde{\delta}*K^{R}_m)(x).
\]
Let $\|x\|\le\kappa W\le P/2$ and $0\le m\le\lfloor H'/2\rfloor-1$. This range of $m$ lies in
$[0,H'-1]$ and satisfies $m\le (H'-1)/2$.
By the model lemma Lemma~\ref{8.13:lem-kink-pull-sign}, $D_T(x,m)\ge0$. Since $\sigma>0$,
\eqref{8.13:eq-tip-pull-bound-2} gives the lower bound in \eqref{8.13:eq-tip-pull-bound-a}.
By the model lemma Lemma~\ref{8.13:lem-kink-pull-decay},
\[
\sigma D_T(x,m)\le\sigma\Bigl[C_T+12\log^{+}\frac{N}{2\max(\|x\|,m)+1}\Bigr],
\]
and together with \eqref{8.13:eq-tip-pull-bound-2} this gives the upper bound.
\end{proof}

Without hypothesis~(i), the vertical part of the kernel bound gives only
$|(\tilde{\delta}*K^{R}_m)(x)|\le 4\sqrt3\,\tau W/(2m+1)$ for every $x$. This confines the bonds whose
difference exceeds the young cap to heights $m$ of order $W/\mathsf{T}$, and it carries no lateral
information.

\begin{remark}[Model: lateral confinement and the mild-fringe hypothesis]
\label{8.13:rem-lateral-confinement}
In the model of the present section, namely the Dirichlet energy $E$ on the periodic strip
$\Lambda=\mathbb{Z}_N\times\{0,\dots,H\}$, with $N=2P$ and $\|x\|$ the distance from $x$ to
$N\mathbb{Z}$, and with the kinked Dirichlet datum $\varphi_0(x)=\sigma\|x\|$, the object of this
remark is the straddling-only minimiser $\tilde{\psi}$ of $E$, with straddling clamp
height $W$, block half-width $\tilde S$, exit trace $\tilde g=\tilde{\psi}(\cdot,W)$ and vertical
bond differences $\tilde F_y(x)=\tilde{\psi}(x,y+1)-\tilde{\psi}(x,y)$. The straddling cap is
$\tau$ and the young cap is $\tau'$.

The mild-fringe hypothesis with parameter $\kappa$ (Definition~\ref{8.13:def-mild-fringe},
a model statement) has two clauses. Clause (i) asks that $\tilde S\ge 2\kappa W$. Clause (ii)
asks that $|\tilde F_W(x)|\le\tau'/2$ whenever $\kappa W\le\|x\|\le P-\kappa W$.

\emph{Clause \textup{(i)}.} Fix $\kappa\le\tfrac12$. Clause (i) follows from the rigid-transport
statement: in the regime of that statement the block of $\tilde{\psi}$ has half-width at least
$W$, whence
\begin{equation*}
\tilde S\;\ge\;W\;\ge\;2\kappa W .
\end{equation*}
The rigid-transport statement is not proved in this book, and clause (i) is used here under that
statement.

The model theorem on inactivity at full band width uses clause (i) only with $\kappa$ a constant
multiple of $1/\mathsf{T}$. There $\kappa$ is subject to the one-sided condition
$\kappa\mathsf{T}\ge 96$ of \eqref{8.13:eq-model-conditions-a}. In that range clause (i) asks only
for a block whose half-width is of order $W/\mathsf{T}$, which is far weaker than
$\tilde S\ge W$.

\emph{Clause \textup{(ii)}.} Clause (ii) is not proved. It is the half of a free-boundary problem
that concerns the fan. The rigid-transport
statement is the block half of the same problem.

One part of clause (ii) follows from clause (i) and the conditions
\eqref{8.13:eq-model-conditions-a}. On the range $\kappa W\le\|x\|\le\tilde S/2$ the argument of
the model Lemma~\ref{8.13:lem-kink-pull-decay}, together with clause (i), gives
\begin{equation}\label{8.13:eq-lateral-confinement-1}
|\tilde F_W(x)|\;\le\;\sigma\Bigl[C_T+12\log^{+}\frac{N}{2\kappa W+1}\Bigr]
+\frac{24\,\tau W}{\tilde S},
\end{equation}
and under \eqref{8.13:eq-model-conditions-a} the right-hand side of
\eqref{8.13:eq-lateral-confinement-1} is at most $\tau'/2$.
What remains of clause (ii) is the edge zone $\tilde S/2<\|x\|$. This zone consists of the outer
half of the block, the fringe and the fan.

\emph{Why a hypothesis of this kind is needed in the present argument.} The height-confinement
corollary confines the hot young bonds at the comparison point $\tilde{x}$ in height; in the
model the confining height is of order $W/\mathsf{T}$ above the straddling clamp. The argument of
the model Lemma~\ref{8.13:lem-kink-pull-decay} re-derives this confinement as follows.

The slope deficit $\tilde\delta$, defined by $\nabla\tilde g=\sigma\nabla T-\tilde\delta$ with
$T(x)=\|x\|$, has total mass $4\tau W$ on the fan. By the kernel bound for $K^{R}_m$, its
contribution at height $m$ satisfies, for every $x$,
\begin{equation}\label{8.13:eq-lateral-confinement-2}
|(\tilde\delta\ast K^{R}_m)(x)|\;\le\;\frac{4\sqrt3\,\tau W}{2m+1}.
\end{equation}
The right-hand side is at most $\tau'/4$ as soon as $2m+1\ge 16\sqrt3\,\tau W/\tau'$, and
$\tau W/\tau'=W/\mathsf{T}$.

Below that height, however, the confinement does not extend laterally: the bound
\eqref{8.13:eq-lateral-confinement-2} holds at every $x$ and says nothing about where the hot bonds
lie. Suppose one knows only that $\tilde g$ is Lipschitz with constant $\sigma$. The classification
of the columns of the straddling band then admits tent poles, which are columns that are not full
and on which
\begin{equation*}
\tilde F_W(x)=\tau+\mu(x,W-1)-\Delta_x\tilde g(x)>\tau ,
\end{equation*}
with $\mu(x,W-1)$ the multiplier of the vertical bond at $(x,W-1)$ and $\Delta_x$ the lateral
second difference. The classification bounds the multiplier at the top of a tent pole only by
$2\sigma$ times the length of the pole. Consequently the lattice lemmas of the model do not exclude
a tent pole of length at least of order $\tau'/(2\sigma)$ lattice rows, placed anywhere in the fan,
that carries $\tilde F_W>\tau'$.

The energy comparison in the proof of the model theorem on inactivity at full band width needs the
hot set of $\tilde{\psi}$ to be confined laterally. Its competitor clips $\tilde{\psi}$ where
$\tilde{\psi}$ is hot, and a laterally spread hot set carries a clipping cost that is not
controlled. The mild-fringe hypothesis is exactly this lateral confinement.

It is not among the hypotheses of Theorem~\ref{3.1:thm-main}
(Notation~\ref{2.2:not-hypotheses}). It is not a cap in the sense of
Definition~\ref{2.3:def-cap}, and it is not a floor of the order of choice
(Notation~\ref{2.3:not-stages-first}). It is a property of the model's straddling-only minimiser,
and no statement other than the model theorem on inactivity at full band width uses it.
\end{remark}

\begin{lemma}[Model: the band is cool away from the tips]\label{8.13:lem-cool-band}
The model is the one of this section: $\tilde{\psi}$ is the straddling-only minimiser
$\psi^{(\infty)}$, that is, the minimiser of the Dirichlet energy
$E(\psi)=\frac12\sum(\nabla\psi)^2$ over the feasible set $\mathcal{K}_c$ at young cap
$c=\infty$, on the periodic strip $\Lambda=\mathbb{Z}_N\times\{0,\dots,H\}$. Its vertical bond
differences are $\tilde F_y(x)=\tilde{\psi}(x,y+1)-\tilde{\psi}(x,y)$, and $\|x\|$ is the
distance from $x$ to $N\mathbb{Z}$, where $N=2P$. Assume both clauses of the mild-fringe
hypothesis of the model, Definition~\ref{8.13:def-mild-fringe}, and the lateral floor
condition \eqref{8.13:eq-model-conditions-a} of Definition~\ref{8.13:def-model-conditions}.
Then for all $x$ with $\kappa W\le\|x\|\le P-\kappa W$ and all $m$ with
$0\le m\le\lfloor H'/2\rfloor-1$,
\begin{equation}\label{8.13:eq-cool-band-1}
  \bigl|\tilde F_{W+m}(x)\bigr|\le \frac{\tau'}{2}.
\end{equation}
\end{lemma}

\begin{proof}
Put $k:=\lfloor H'/2\rfloor$ and
\begin{equation}\label{8.13:eq-cool-band-2}
  B:=\max\Bigl\{\frac{\tau'}{2},\;
  \sigma\Bigl[C_T+12\log^{+}\Bigl(\frac{N}{2\kappa W-1}\Bigr)\Bigr]+\beta_\kappa\Bigr\}.
\end{equation}
For $m=0$ the bound \eqref{8.13:eq-cool-band-1} is clause (ii) of the mild-fringe hypothesis
(Definition~\ref{8.13:def-mild-fringe}). If $k=1$ there is nothing further to prove, so let
$k\ge2$. Write $G(x,m):=\tilde F_{W+m}(x)$.

The proof of the model lemma, Lemma~\ref{8.13:lem-kink-pull-sign}, uses two properties of the
pull of the kink. In the same way, $G$ satisfies the four-neighbour mean-value equation at
every $(x,m)$ with $1\le m\le H'-2$, and $G(x,H'-1-m)=-G(x,m)$.

The two arcs are
\begin{equation*}
  \Omega:=\{\kappa W\le\|x\|\le P-\kappa W\}\times\{1\le m\le k-1\}.
\end{equation*}
The sites outside $\Omega$ adjacent to it lie in three sets.
\begin{itemize}
\item[(a)] The lateral columns $\|x\|=\kappa W-1$ and $\|x-P\|=\kappa W-1$, with
$0\le m\le k-1$.
\item[(b)] The bottom row $m=0$.
\item[(c)] The row $m=k$.
\end{itemize}

\emph{Lateral columns.} Let $\|x\|=\kappa W-1$ and $0\le m\le k-1$. The model lemma,
Lemma~\ref{8.13:lem-tip-pull-bound}, holds under clause (i) of
Definition~\ref{8.13:def-mild-fringe}, and its bound \eqref{8.13:eq-tip-pull-bound-a} gives
\begin{equation*}
  -\beta_\kappa\le G(x,m)\le
  \sigma\Bigl[C_T+12\log^{+}\Bigl(\frac{N}{2\max(\|x\|,m)+1}\Bigr)\Bigr]+\beta_\kappa.
\end{equation*}
Since $2\max(\|x\|,m)+1\ge2\kappa W-1$ and $\log^{+}$ is non-decreasing, while $\sigma\ge0$
and $C_T\ge0$, this yields
\begin{equation}\label{8.13:eq-cool-band-3}
  |G(x,m)|\le\sigma\Bigl[C_T+12\log^{+}\Bigl(\frac{N}{2\kappa W-1}\Bigr)\Bigr]
  +\beta_\kappa\le B.
\end{equation}
Now let $\|x-P\|=\kappa W-1$. Then $x=x'+P$ with $\|x'\|=\kappa W-1$. The symmetry of the
model under translation by the half-period gives $\tilde F_y(x'+P)=-\tilde F_y(x')$, so
\eqref{8.13:eq-cool-band-3} holds at $x$ as well.

\emph{Bottom row.} For $\kappa W\le\|x\|\le P-\kappa W$, clause (ii) of
Definition~\ref{8.13:def-mild-fringe} gives $|G(x,0)|\le\tau'/2\le B$.

\emph{Row $k$.} If $H'=2k+1$, then $H'-1-k=k$, so antisymmetry gives $G(\cdot,k)=-G(\cdot,k)$,
that is $G(\cdot,k)=0$. If $H'=2k$, then $H'-1-k=k-1$, so $G(\cdot,k)=-G(\cdot,k-1)$.

\emph{Maximum argument for $G-B$.} Suppose $a:=\max_\Omega G>B$, and let it be attained at
$(x_0,m_0)$. Note that $B\ge\tau'/2>0$, since the young cap $\tau'$ is positive.

If $m_0\le k-2$, all four neighbours of $(x_0,m_0)$ lie in $\Omega$ or in the sets (a) and (b).
Their values are therefore at most $a$, and at most $B<a$ when they lie in (a) or (b). The
mean-value equation at $(x_0,m_0)$ then forces all four neighbours to equal $a$. In
particular, $(x_0,m_0+1)$ is a point of $\Omega$ at which $G=a$. Repeating this step, we find
that $G(x_0,k-1)=a$.

At $(x_0,k-1)$ the neighbours $(x_0\pm1,k-1)$ and $(x_0,k-2)$ have values at most $a$, by the
bounds just listed. The neighbour $(x_0,k)$ has value $0$ or $-a$. The mean-value equation
therefore gives
\begin{equation*}
  4a\le3a+0\qquad\text{or}\qquad4a\le3a-a,
\end{equation*}
so $a\le0<B$, a contradiction. Hence $G\le B$ on $\Omega$.

\emph{Minimum argument for $B-G$.} The same argument applied to $-G$ shows $G\ge-B$ on $\Omega$.
This is legitimate because $-G$ satisfies the same mean-value equation and antisymmetry, and
the bounds on the sets (a) and (b) are bounds on $|G|$. At row $k-1$ one obtains
$-4a\ge-3a+0$ or $-4a\ge-3a+a$.

Thus $|G|\le B$ on $\Omega$. By \eqref{8.13:eq-model-conditions-a} the second entry of the
maximum in \eqref{8.13:eq-cool-band-2} is at most $\tau'/2$, so $B=\tau'/2$. Together with the
case $m=0$, this proves \eqref{8.13:eq-cool-band-1}.
\end{proof}

\begin{corollary}[Model: the hot box]\label{8.13:cor-hot-box}
The setting is that of the model fixed in Notation~\ref{8.13:not-straddling-minimiser}:
$\tilde{\psi}$ is the straddling-only minimiser of the Dirichlet energy
$E(\psi)=\tfrac12\sum(\nabla\psi)^2$ on the periodic strip
$\Lambda=\mathbb{Z}_N\times\{0,\dots,H\}$, $N=2P$, with the kinked Dirichlet datum
$\varphi_0(x)=\sigma\|x\|$ of slope $\sigma>0$, where $\|x\|$ is the distance from $x$ to
$N\mathbb{Z}$. Here $\tilde F_y(x)=\tilde{\psi}(x,y+1)-\tilde{\psi}(x,y)$ are its vertical bond
differences, $W$ is the clamp height and $\bar W$ the height of the young band. The
\emph{bottom band} is the set of vertical bonds $(x,W+m)$ with $0\le m\le \bar W-1$. The
constants $C_K^{0}$, $\beta_\kappa$, $C_T$ are those of the model displays
\eqref{8.13:eq-straddling-minimiser-a}, \eqref{8.13:eq-tip-pull-bound-a} and
\eqref{8.13:eq-kink-pull-decay-a}.

For $c>0$ let
\[
A(c):=\bigl\{(x,W+m):\ 0\le m\le \bar W-1,\ \tilde F_{W+m}(x)>c\bigr\}
\]
be the set of up-hot bonds of the bottom band, and put
\begin{equation}\label{8.13:eq-hot-box-a}
\begin{aligned}
h_{\sharp}(c)&:=\max\Bigl\{1,\ \tfrac12\Bigl[N\exp\Bigl(-\tfrac{1}{12}
\Bigl(\frac{c-\beta_\kappa}{\sigma}-C_T\Bigr)\Bigr)-1\Bigr]\Bigr\},\\
\mathfrak{h}(c)&:=\max\Bigl\{0,\ \tfrac12\Bigl[N\exp\Bigl(-\tfrac{1}{6}
\Bigl(\frac{c}{\sigma}-C_K^{0}-6\Bigr)\Bigr)-1\Bigr]\Bigr\}.
\end{aligned}
\end{equation}
Assume the mild-fringe hypothesis of the model with lateral fraction $\kappa$ and the
lateral floor condition \eqref{8.13:eq-model-conditions-a} of the model instance
(Definition~\ref{8.13:def-model-conditions}). Then for every $c\ge\tau'$ the following hold.
\begin{enumerate}
\item[(a)] $A(c)\subset Q(c)$, where
\[
Q(c):=\bigl\{(x,W+m):\ \|x\|<\kappa W,\ \max(\|x\|,m)<h_{\sharp}(c)\bigr\},
\]
and on $A(c)$
\[
0<\tilde F_{W+m}(x)-c\le 12\sigma
\log\frac{2h_{\sharp}(c)+1}{2\max(\|x\|,m)+1}.
\]
\item[(b)] $\tilde F_{W+m}(x)\ge -c$ for all $x$ with $\|x\|\le P/2$ and all
$0\le m\le\bar W-1$.
\item[(c)] $\tilde F_{W+m}(x)\le c-6\sigma$ for all $x$ and all $m$ with
$\mathfrak{h}(c)\le m\le (H'-1)/2$, in particular for all bonds of the bottom band with
$m\ge\mathfrak{h}(c)$; and $\mathfrak{h}(c)\le h_{\sharp}(c)$.
\item[(d)] The down-hot set $\{\tilde F<-c\}$ of the bottom band is the image of $A(c)$
under $x\mapsto x+P$. The hot sets of the top band are the mirror images of those of the
bottom band under the reflection in the vertical coordinate $y$.
\end{enumerate}
\end{corollary}

\begin{proof}
\emph{Preliminaries.} By condition \eqref{8.13:eq-model-conditions-d} of the model
instance, $\bar W\le (H'-1)/2$. Hence every bond of the bottom band satisfies
$2m\le 2\bar W-2\le H'-3$. Since $\lfloor H'/2\rfloor\ge (H'-1)/2$, this gives
$m\le\lfloor H'/2\rfloor-1$. Therefore the bound \eqref{8.13:eq-tip-pull-bound-a} of the
model lemma, Lemma~\ref{8.13:lem-tip-pull-bound}, and the bound of the model lemma,
Lemma~\ref{8.13:lem-cool-band}, both apply on the whole bottom band.

The constant $\beta_\kappa=24\tau/\kappa$ is non-negative, and $C_T>0$. Since
$\sigma>0$ and $\log^{+}\ge0$, condition \eqref{8.13:eq-model-conditions-a} gives
$\sigma C_T+\beta_\kappa\le\tau'/2$. In particular $\beta_\kappa\le\tau'/2<\tau'$, and for
$c\ge\tau'$
\begin{equation}\label{8.13:eq-hot-box-1}
c-\beta_\kappa-\sigma C_T\ \ge\ c-\tau'/2\ \ge\ \tau'/2\ >\ 0 .
\end{equation}
We also use the symmetries of the model, as stated in the symmetry statement of the
model: antisymmetry under lateral translation by $P$, namely
$\tilde F_y(x+P)=-\tilde F_y(x)$, and reflection in the vertical coordinate $y$.
We write $\nu:=\max(\|x\|,m)$.

(a) Let $(x,W+m)$ lie in the bottom band. There are three cases.

\emph{Away from the tips}, $\kappa W\le\|x\|\le P-\kappa W$. Here the model lemma,
Lemma~\ref{8.13:lem-cool-band}, gives $|\tilde F_{W+m}(x)|\le\tau'/2$. Since
$\tau'/2<\tau'\le c$, the bond is not hot.

\emph{Near the top tip}, $\|x-P\|<\kappa W$. Put $x'=x-P$, so that $\|x'\|<\kappa W$. The
lower side of the model bound \eqref{8.13:eq-tip-pull-bound-a} gives
$\tilde F_{W+m}(x')\ge-\beta_\kappa$.
By the antisymmetry, $\tilde F_{W+m}(x)=-\tilde F_{W+m}(x')\le\beta_\kappa$. Since
$\beta_\kappa<\tau'\le c$, the bond is not hot.

\emph{Near the bottom tip}, $\|x\|<\kappa W$. The upper side of the model bound
\eqref{8.13:eq-tip-pull-bound-a} gives
\begin{equation}\label{8.13:eq-hot-box-2}
\tilde F_{W+m}(x)-c\ \le\ 12\sigma\log^{+}\frac{N}{2\nu+1}
-\bigl(c-\beta_\kappa-\sigma C_T\bigr).
\end{equation}
Suppose first that $2\nu+1\ge N$. Then the logarithm vanishes, and by
\eqref{8.13:eq-hot-box-1} the right side is at most $-\tau'/2<0$, so the bond is not hot.
On $A(c)$ we therefore have $2\nu+1<N$, and $\log^{+}$ may be replaced by $\log$ in
\eqref{8.13:eq-hot-box-2}.

If the maximum defining $h_{\sharp}(c)$ in \eqref{8.13:eq-hot-box-a} is attained by its
second entry, then
\[
2h_{\sharp}(c)+1=N\exp\Bigl(-\frac{c-\beta_\kappa-\sigma C_T}{12\sigma}\Bigr),
\qquad\text{that is}\qquad
c-\beta_\kappa-\sigma C_T=12\sigma\log\frac{N}{2h_{\sharp}(c)+1}.
\]
Substituting this into \eqref{8.13:eq-hot-box-2} yields
\[
\tilde F_{W+m}(x)-c\le 12\sigma\log\frac{2h_{\sharp}(c)+1}{2\nu+1}.
\]
The right side is positive only if $\nu<h_{\sharp}(c)$.

If $h_{\sharp}(c)=1$ is the floor value, then the second entry is at most $1$. This means
$N\exp\bigl(-(c-\beta_\kappa-\sigma C_T)/(12\sigma)\bigr)\le 3$, that is,
$c-\beta_\kappa-\sigma C_T\ge 12\sigma\log(N/3)$. Then \eqref{8.13:eq-hot-box-2} gives
\[
\tilde F_{W+m}(x)-c\le 12\sigma\log\frac{3}{2\nu+1}
=12\sigma\log\frac{2h_{\sharp}(c)+1}{2\nu+1}.
\]
This is positive only for $\nu=0$, that is, only for $\nu<1=h_{\sharp}(c)$.

In each case a hot bond has $\|x\|<\kappa W$ and $\nu<h_{\sharp}(c)$, and it satisfies the
displayed bound. This proves $A(c)\subset Q(c)$ together with the bound on $A(c)$.

(b) Let $\|x\|\le P/2$ and $0\le m\le\bar W-1$. If $\|x\|\le\kappa W$, the lower side of
the model bound \eqref{8.13:eq-tip-pull-bound-a} gives
\[
\tilde F_{W+m}(x)\ \ge\ -\beta_\kappa\ \ge\ -\tau'\ \ge\ -c,
\]
using $\beta_\kappa\le\tau'$ from
\eqref{8.13:eq-model-conditions-a} as above. If $\kappa W\le\|x\|\le P/2$, then
$\kappa W\le P/2$, so $P/2\le P-\kappa W$. Hence
the model lemma, Lemma~\ref{8.13:lem-cool-band}, applies and gives
$\tilde F_{W+m}(x)\ge-\tau'/2\ge-c$.

(c) For $0\le m\le(H'-1)/2$, the bound \eqref{8.13:eq-straddling-minimiser-a} of the
model notation gives
\[
|\tilde F_{W+m}(x)|\le b(m)\le\sigma\Bigl[C_K^{0}+6\log^{+}\frac{N}{2m+1}\Bigr]
\qquad\text{for all }x .
\]
By the preliminaries, this range contains the bottom band. Next,
\eqref{8.13:eq-hot-box-1} together with $\beta_\kappa\ge0$ gives $c\ge\tau'\ge2\sigma C_T$.
Since $C_T\ge 2C_K^{0}+18$ by the model bound \eqref{8.13:eq-kink-pull-decay-a}, it
follows that
$c/\sigma-C_K^{0}-6\ge0$. Consequently the following conditions are equivalent:
\begin{align*}
\sigma\Bigl[C_K^{0}+6\log^{+}\frac{N}{2m+1}\Bigr]\le c-6\sigma
&\iff \log\frac{N}{2m+1}\le\frac16\Bigl(\frac{c}{\sigma}-C_K^{0}-6\Bigr)\\
&\iff 2m+1\ge N\exp\Bigl(-\frac16\Bigl(\frac{c}{\sigma}-C_K^{0}-6\Bigr)\Bigr).
\end{align*}
For integers $m\ge0$, the last inequality is equivalent to $m\ge\mathfrak{h}(c)$ by
\eqref{8.13:eq-hot-box-a}. Hence $\tilde F_{W+m}(x)\le c-6\sigma$ for all $x$ and all
$m\ge\mathfrak{h}(c)$ with $m\le(H'-1)/2$, in particular in the bottom band.

For the comparison of the two scales: the first entries of the two maxima in
\eqref{8.13:eq-hot-box-a} satisfy $0\le1$, and, since the exponential is increasing, the
second entry of $\mathfrak{h}(c)$ is at most that of $h_{\sharp}(c)$ as soon as
\[
\frac16\Bigl(\frac{c}{\sigma}-C_K^{0}-6\Bigr)\ge
\frac1{12}\Bigl(\frac{c-\beta_\kappa}{\sigma}-C_T\Bigr).
\]
This inequality is equivalent to
\[
\frac{c}{\sigma}\ \ge\ 2C_K^{0}+12-C_T-\frac{\beta_\kappa}{\sigma}.
\]
Since $C_T\ge 2C_K^{0}+18$, the right side is at most $-6-\beta_\kappa/\sigma<0$, whereas
$c/\sigma>0$. So $\mathfrak{h}(c)\le h_{\sharp}(c)$.

(d) By the antisymmetry $\tilde F_y(x+P)=-\tilde F_y(x)$ of the symmetry statement of the
model, we have
$\tilde F_{W+m}(x+P)<-c$ if and only if $\tilde F_{W+m}(x)>c$. Hence the down-hot set of
the bottom band is $A(c)+P$. The assertion on the top band is the reflection symmetry of
the model in the vertical coordinate $y$, from the same symmetry statement.
\end{proof}

For orientation only: in the block units of the model instance, with $n=L^{k_o}$,
$\theta_y=\theta'\mathsf{T}$ and $P_{\mathrm{bl}}$ the zone half-width in blocks, the two
scales at the young cap are
\[
\frac{h_{\sharp}(\theta_y)}{n}=P_{\mathrm{bl}}\,e^{C_T/12}\,
e^{-(\mathsf{T}-24/\kappa)/(12s_e)}\,(1+o(1)),
\qquad
\frac{\mathfrak{h}(\theta_y)}{n}=P_{\mathrm{bl}}\,e^{1+C_K^{0}/6}\,e^{-\mathsf{T}/(6s_e)},
\]
where $o(1)$ tends to zero as $n\to\infty$.
These are the lattice counterparts of the continuum scale
$D_W\,e^{-\pi\theta\mathsf{T}/(2s_e)}$, where $D_W$ is the width parameter of the zone and
$\theta$ is the straddling cap before the factor $1+\vartheta_1$, so that
$\theta'=\theta(1+\vartheta_1)$; the pair of lattice constants $12$, $6$ stands in
place of $\pi/2$. The right-hand sides do not involve the cutoff. Since $\mathsf{T}/s_e$
is of order $L^2$ and $P_{\mathrm{bl}}$ is comparable to $D_W$, both scales are
exponentially small in $\mathsf{T}/s_e$ relative to $P_{\mathrm{bl}}$. Hence they are much
smaller than $W/(nL^2)$ under condition \eqref{8.13:eq-model-conditions-c}.

\begin{lemma}[Model: the variational inequality]\label{8.13:lem-variational-inequality}
This lemma is a statement about the two-tier box-clamp lattice model of
Definition~\ref{8.13:def-two-tier-model}, which is a model and not Ba{\l}aban's step.

In that model the fields are real functions $\psi$ on the periodic strip
$\Lambda=\mathbb{Z}_N\times\{0,1,\dots,H\}$. They satisfy the Dirichlet conditions
$\psi(\cdot,0)=\psi(\cdot,H)=\varphi_0$, and their energy is
$E(\psi)=\tfrac12\sum(\nabla\psi)^2$, summed over the nearest-neighbour bonds of $\Lambda$.

The feasible set $\mathcal{K}_c$ is defined by bounds on bond differences:
\begin{itemize}
\item on the bonds of the lower tier, the bound is $\tau$;
\item on the bonds of the young tier, the bound is the young cap $c$;
\item the mirror constraints near the top row are imposed as well.
\end{itemize}
The set $\mathcal{K}_\infty$ is obtained by dropping the young tier. Further:
\begin{itemize}
\item $\psi^{(c)}$ is the minimiser of $E$ over $\mathcal{K}_c$;
\item $\tilde{\psi}=\psi^{(\infty)}$ is the minimiser of $E$ over $\mathcal{K}_\infty$.
\end{itemize}

For two fields $\psi,v$ on $\Lambda$ write
\begin{equation*}
\langle\nabla\psi,\nabla v\rangle=\sum\nabla\psi\,\nabla v,
\qquad
\|\nabla v\|^2=\langle\nabla v,\nabla v\rangle,
\end{equation*}
where both sums run over the nearest-neighbour bonds.

Let $c\ge\tau'$, and let $\psi^{c}\in\mathcal{K}_c$ be an arbitrary field. It is written without
parentheses, to distinguish an arbitrary competitor from the minimiser $\psi^{(c)}$. Then
\begin{equation}\label{8.13:eq-variational-inequality-a}
\tfrac12\bigl\|\nabla\bigl(\psi^{(c)}-\tilde{\psi}\bigr)\bigr\|^2
\le E\bigl(\psi^{(c)}\bigr)-E(\tilde{\psi})
\le E\bigl(\psi^{c}\bigr)-E(\tilde{\psi}).
\end{equation}
Moreover, the field $v:=\psi^{(c)}-\tilde{\psi}$ vanishes on the Dirichlet rows $y=0$ and $y=H$.
\end{lemma}

\begin{proof}
Let $c\ge\tau'$. The set $\mathcal{K}_c$ is obtained from $\mathcal{K}_\infty$ by imposing the
young-tier constraints in addition. Hence $\mathcal{K}_c\subset\mathcal{K}_\infty$, and in
particular $\psi^{(c)}\in\mathcal{K}_\infty$.

Put $v:=\psi^{(c)}-\tilde{\psi}$. Both fields equal $\varphi_0$ on the rows $y=0$ and $y=H$, so
$v$ vanishes there. This proves the last assertion.

Since $E$ is quadratic, expanding $(\nabla\tilde{\psi}+t\nabla v)^2$ bond by bond gives, for every
real $t$,
\begin{equation}\label{8.13:eq-variational-inequality-1}
E(\tilde{\psi}+tv)=E(\tilde{\psi})+t\,\langle\nabla\tilde{\psi},\nabla v\rangle
+\tfrac{t^2}{2}\,\|\nabla v\|^2 .
\end{equation}
At $t=1$ this reads
\begin{equation*}
E\bigl(\psi^{(c)}\bigr)-E(\tilde{\psi})
=\langle\nabla\tilde{\psi},\nabla v\rangle+\tfrac12\|\nabla v\|^2 .
\end{equation*}
By \eqref{8.13:eq-variational-inequality-1}, the function $t\mapsto E(\tilde{\psi}+tv)$ has
derivative $\langle\nabla\tilde{\psi},\nabla v\rangle$ at $t=0$.

We show that this derivative is non-negative. The set $\mathcal{K}_\infty$ is defined by
affine equalities (the Dirichlet rows) and by bounds on absolute values of linear functionals
of $\psi$ (the bond differences). It is therefore convex. Since $\tilde{\psi}$ and $\psi^{(c)}$
both lie in $\mathcal{K}_\infty$, so does
\begin{equation*}
\tilde{\psi}+tv=(1-t)\tilde{\psi}+t\psi^{(c)}
\qquad\text{for every } t\in[0,1].
\end{equation*}
By minimality of $\tilde{\psi}$ on $\mathcal{K}_\infty$, we have
$E(\tilde{\psi}+tv)\ge E(\tilde{\psi})$ for $t\in[0,1]$. With
\eqref{8.13:eq-variational-inequality-1}, for $t\in(0,1]$,
\begin{equation*}
\langle\nabla\tilde{\psi},\nabla v\rangle+\tfrac{t}{2}\|\nabla v\|^2
=\frac{E(\tilde{\psi}+tv)-E(\tilde{\psi})}{t}\ge0 .
\end{equation*}
Letting $t\downarrow0$ gives $\langle\nabla\tilde{\psi},\nabla v\rangle\ge0$.

Inserting this into the identity at $t=1$ yields
$E(\psi^{(c)})-E(\tilde{\psi})\ge\tfrac12\|\nabla v\|^2$, which is the first inequality in
\eqref{8.13:eq-variational-inequality-a}.

For the second inequality: $\psi^{(c)}$ minimises $E$ over $\mathcal{K}_c$ and
$\psi^{c}\in\mathcal{K}_c$, so $E(\psi^{(c)})\le E(\psi^{c})$. Subtracting $E(\tilde{\psi})$
from both sides gives the second inequality in \eqref{8.13:eq-variational-inequality-a}.
\end{proof}

\begin{definition}[Model: the clipped competitor]\label{8.13:def-clipped-competitor}
The model of this definition is the minimisation of the discrete Dirichlet energy, defined below,
over feasible fields on the periodic strip $\Lambda=\mathbb{Z}_N\times\{0,\dots,H\}$, with $N=2P$;
here $\|x\|$ is the distance from $x$ to $N\mathbb{Z}$.
For a field $u$ on $\Lambda$, $\|\nabla u\|^2$ denotes the sum over the bonds of
$\Lambda$ of the squared differences of $u$, and the Dirichlet energy is
$E(u)=\tfrac12\|\nabla u\|^2$. The field $\tilde\psi$ is the minimiser of $E$ over
$\mathcal{K}_\infty$ as in the model statement Lemma~\ref{8.13:lem-variational-inequality}, with
vertical bond differences $\tilde F_y(x)=\tilde\psi(x,y+1)-\tilde\psi(x,y)$. The bottom band is the
set of rows $W+m$ with $0\le m\le\bar W-1$.

Let $c>0$ be a young cap. Define $\mathsf{d}=\mathsf{d}_c$ on $\Lambda$ by the column recursion
\begin{align*}
\mathsf{d}(x,y)&:=0 &&\text{for } y\le W,\\
\mathsf{d}(x,y+1)&:=\bigl(\mathsf{d}(x,y)+\tilde F_y(x)-c\bigr)^{+}
&&\text{for } W\le y\le\lfloor H/2\rfloor-1,\\
\mathsf{d}(x,y)&:=0 &&\text{for } y>\lfloor H/2\rfloor .
\end{align*}
Write $\nabla_y\mathsf{d}(x,y):=\mathsf{d}(x,y+1)-\mathsf{d}(x,y)$, and put
\begin{align*}
\mathsf{D}(x,y)&:=\mathsf{d}(x,y)-\mathsf{d}(x-P,y) &&\text{for } y\le\lfloor H/2\rfloor,\\
\mathsf{D}(x,y)&:=\mathsf{D}(x,H-y) &&\text{for } y>\lfloor H/2\rfloor .
\end{align*}
The \emph{clipped competitor} is $\psi^{c}:=\tilde\psi-\mathsf{D}$. Its vertical bond differences
are written $F^{c}_y(x)=\psi^{c}(x,y+1)-\psi^{c}(x,y)$.

The clipped competitor has the following properties. Property (e1) holds for every $c>0$.
Property (e2) holds for $c\ge\tau'$ under the hypotheses of the model statement
Corollary~\ref{8.13:cor-hot-box} and part (c) of the model condition
\eqref{8.13:eq-model-conditions-c} of Definition~\ref{8.13:def-model-conditions}; the latter
enters (e2) only to exclude hot bonds above the bottom band. Property (e3) holds for $c\ge\tau'$
under these hypotheses and the model condition \eqref{8.13:eq-model-conditions-d}.
\begin{itemize}
\item[(e1)] $\mathsf{d}\ge0$. For $W\le y\le\lfloor H/2\rfloor-1$ the increment
$\nabla_y\mathsf{d}(x,y)$ is given by the second relation in the first line of
\eqref{8.13:eq-clipped-competitor-a}. At every such vertical bond at whose endpoints
$\mathsf{D}=\mathsf{d}$, the bond difference $F^{c}_y$ is given by the third relation in that line.
\item[(e2)] The column $\mathsf{d}(x,\cdot)$ vanishes identically unless the column $x$ meets
$A(c)$. Consequently $\mathsf{d}$ is supported in the set on the second line of
\eqref{8.13:eq-clipped-competitor-a}. Moreover
$\mathsf{d}\equiv0$ on the columns near the top tip, $\|x-P\|<\kappa W$. Hence $\mathsf{d}$ and
$\mathsf{d}(\cdot-P,\cdot)$ have disjoint supports, $\mathsf{D}=\mathsf{d}$ near the bottom tip,
$\|x\|<\kappa W$, and $\mathsf{D}=-\mathsf{d}(\cdot-P,\cdot)$ near the top tip.
\item[(e3)] On a column $x$ meeting $A(c)$, the function $\mathsf{d}$ increases only at bonds of
$A(c)$ (the hot bonds), at heights $m<h_\sharp(c)$. It never exceeds
\[
\mathsf{d}_{\max}(x):=\sum_{0\le m<h_\sharp(c)}\bigl(\tilde F_{W+m}(x)-c\bigr)^{+},
\]
and $\mathsf{d}_{\max}(x)\le48\sigma h_\sharp(c)$. At the rows $W+m$ with $m\ge\mathfrak{h}(c)$
each step lowers $\mathsf{d}$ by at least $\min(\mathsf{d},6\sigma)$, so $\mathsf{d}=0$ from height
$W+\mathfrak{h}(c)+8h_\sharp(c)$ on, and this height is at most $W+9h_\sharp(c)$.
The support of $\mathsf{d}$ lies in the box
\[
\mathsf{R}(c):=\bigl\{(x,y):\ \|x\|<h_\sharp(c),\ W+1\le y\le W+9h_\sharp(c)\bigr\}.
\]
This box lies inside the bottom band and below height $H/2$. Consequently $\mathsf{D}$ is well
defined and consists of four congruent pieces with pairwise disjoint supports. These bounds are
collected in the third line of \eqref{8.13:eq-clipped-competitor-a}.
\end{itemize}
In formulas:
\begin{equation}\label{8.13:eq-clipped-competitor-a}
\begin{gathered}
\mathsf{d}\ge0,\qquad
\nabla_y\mathsf{d}(x,y)=\max\bigl(\tilde F_y(x)-c,\,-\mathsf{d}(x,y)\bigr),\qquad
F^{c}_y=\min\bigl(c,\tilde F_y+\mathsf{d}\bigr)\ \text{where }\mathsf{D}=\mathsf{d};\\
\operatorname{supp}\mathsf{d}\subset\bigl\{\|x\|<\min\bigl(h_\sharp(c),\kappa W\bigr)\bigr\}
\times\{y\ge W+1\};\\
\mathsf{d}(x,\cdot)\le\mathsf{d}_{\max}(x)\le48\sigma h_\sharp(c),\qquad
\operatorname{supp}\mathsf{d}\subset\mathsf{R}(c),\qquad
\|\nabla\mathsf{D}\|^2=4\|\nabla\mathsf{d}\|^2 .
\end{gathered}
\end{equation}
\end{definition}

\begin{proof}[Proof of properties \textup{(e1)}--\textup{(e3)}]
\emph{(e1).} We have $\mathsf{d}(x,y)=0$ for $y\le W$ and for $y>\lfloor H/2\rfloor$, and every
other value is a positive part.
Hence $\mathsf{d}\ge0$ by induction on $y$. For $W\le y\le\lfloor H/2\rfloor-1$,
\[
\nabla_y\mathsf{d}(x,y)=\max\bigl(\mathsf{d}+\tilde F_y-c,\,0\bigr)-\mathsf{d}
=\max\bigl(\tilde F_y(x)-c,\,-\mathsf{d}(x,y)\bigr).
\]
At a bond with $\mathsf{D}=\mathsf{d}$ at both endpoints, $F^{c}_y=\tilde F_y-\nabla_y\mathsf{d}$.
By the identity just proved this equals $\tilde F_y-\max(\tilde F_y-c,-\mathsf{d})$, which is
$\min(c,\tilde F_y+\mathsf{d})$.

\emph{(e2).} Since $-\mathsf{d}\le0$, property (e1) shows that $\mathsf{d}(x,y+1)>\mathsf{d}(x,y)$
only if $\tilde F_y(x)>c$. Starting from $\mathsf{d}(x,W)=0$, the column $\mathsf{d}(x,\cdot)$
therefore stays zero unless $\tilde F_y(x)>c$ at some $y$ with $W\le y\le\lfloor H/2\rfloor-1$.

Within the bottom band, the bonds with $\tilde F_{W+m}(x)>c$ are exactly the points of $A(c)$.
Above the bottom band there are none, for the following reason. The function $h_\sharp$ of
\eqref{8.13:eq-hot-box-a}, in the model Corollary~\ref{8.13:cor-hot-box}, is non-increasing in its
argument, so $h_\sharp(c)\le h_\sharp(\tau')$.
Together with part (c) of the model condition \eqref{8.13:eq-model-conditions-c}, this gives
\begin{equation}\label{8.13:eq-clipped-competitor-1}
9h_\sharp(c)\le9h_\sharp(\tau')<\bar W .
\end{equation}
Part (c) of the model Corollary~\ref{8.13:cor-hot-box} gives $\mathfrak{h}(c)\le h_\sharp(c)$.
Hence every row $W+m$ with $m\ge\bar W$ has $m\ge\mathfrak{h}(c)$. The same part of that
corollary then gives $\tilde F_{W+m}(x)\le c-6\sigma<c$ on such rows.

By part (a) of the model Corollary~\ref{8.13:cor-hot-box}, $A(c)\subset Q(c)$. So a column
meeting $A(c)$ has $\|x\|<\kappa W$ and $\|x\|\le\max(\|x\|,m)<h_\sharp(c)$. Since
$\mathsf{d}=0$ for $y\le W$, the second line of \eqref{8.13:eq-clipped-competitor-a} follows.

Within the bottom band, on the columns near the top tip, $\|x-P\|<\kappa W$, the lateral bound
near the tips gives $\tilde F\le\beta_\kappa$, as in the proof of part (a) of the model
Corollary~\ref{8.13:cor-hot-box}. The lateral floor condition of the model
Definition~\ref{8.13:def-model-conditions} gives $\beta_\kappa\le\tau'/2<\tau'$, and $\tau'\le c$
by hypothesis, so $\tilde F<c$ there. Above the bottom band $\tilde F<c$ for all $x$ by the
preceding paragraphs. Hence no bond on these columns raises $\mathsf{d}$, and $\mathsf{d}\equiv0$
on these columns.

If $\|x\|<\kappa W$, then $\|(x-P)-P\|=\|x-N\|$, which equals $\|x\|$ and so is less than
$\kappa W$. So $x-P$ lies near the top tip,
$\mathsf{d}(x-P,\cdot)\equiv0$, and $\mathsf{D}=\mathsf{d}$ there. If $\|x-P\|<\kappa W$, then
$\mathsf{d}(x,\cdot)\equiv0$ and $\mathsf{D}=-\mathsf{d}(x-P,\cdot)$. Thus $\mathsf{d}$ is
supported where $\|x\|<\kappa W$ and $\|x-P\|\ge\kappa W$. The function $\mathsf{d}(\cdot-P,\cdot)$
is supported where $\|x-P\|<\kappa W$ and $\|x\|\ge\kappa W$. These two supports are disjoint.

\emph{(e3).} By (e1), $\nabla_y\mathsf{d}>0$ only where $\tilde F_y(x)>c$, and there
$\nabla_y\mathsf{d}=\tilde F_y(x)-c$. By (e2) such bonds are points $(x,W+m)$ of
$A(c)\subset Q(c)$, so $m<h_\sharp(c)$. The steps with $\nabla_y\mathsf{d}\le0$ do not increase
$\mathsf{d}$. Therefore $\mathsf{d}(x,y)\le\mathsf{d}_{\max}(x)$.

Consider a term of $\mathsf{d}_{\max}(x)$ that is not zero. Then $(x,W+m)\in A(c)$, and part (a)
of the model Corollary~\ref{8.13:cor-hot-box} bounds it by
\[
12\sigma\log\frac{2h_\sharp+1}{2\max(\|x\|,m)+1}\le12\sigma\log\frac{2h_\sharp+1}{2m+1}.
\]
The right side is non-negative for $m<h_\sharp$, so the bound also holds for the vanishing terms.
The function $t\mapsto\log\bigl((2h_\sharp+1)/(2t+1)\bigr)$ is decreasing and non-negative on
$[0,h_\sharp]$. Hence each logarithm $\log\bigl((2h_\sharp+1)/(2m+1)\bigr)$ with
$1\le m<h_\sharp$ is at most the integral of this function
over $[m-1,m]$, and these intervals lie disjointly in $[0,h_\sharp]$. The substitution
$u=(2t+1)/(2h_\sharp+1)$ and the identity $\int_0^1\log(1/u)\,du=1$ give
\[
\int_0^{h_\sharp}\log\frac{2h_\sharp+1}{2t+1}\,dt
=\frac{2h_\sharp+1}{2}\int_{1/(2h_\sharp+1)}^{1}\log\frac1u\,du\le h_\sharp+\tfrac12 .
\]
Keeping the term $m=0$ separately, and using $h_\sharp\ge1$ in the last step, we obtain
\begin{align*}
\mathsf{d}_{\max}(x)
&\le\sum_{0\le m<h_\sharp}12\sigma\log\frac{2h_\sharp+1}{2m+1}\\
&\le12\sigma\Bigl[\log(2h_\sharp+1)+\int_0^{h_\sharp}\log\frac{2h_\sharp+1}{2t+1}\,dt\Bigr]\\
&\le12\sigma\bigl[\log(2h_\sharp+1)+h_\sharp+\tfrac12\bigr]\le48\sigma h_\sharp .
\end{align*}

For $m\ge\mathfrak{h}(c)$, part (c) of the model Corollary~\ref{8.13:cor-hot-box} gives
$\tilde F_{W+m}(x)\le c-6\sigma$. By (e1),
\[
\nabla_y\mathsf{d}=-\min\bigl(c-\tilde F,\,\mathsf{d}\bigr)\le-\min(6\sigma,\mathsf{d}).
\]
So above height $\mathfrak{h}(c)$ every step decreases $\mathsf{d}$ by at least
$\min(\mathsf{d},6\sigma)$. Once $\mathsf{d}$ vanishes it stays zero, because then
$\nabla_y\mathsf{d}=\max(\tilde F-c,0)=0$. Since $\mathsf{d}\le48\sigma h_\sharp=6\sigma\cdot
8h_\sharp$, the function $\mathsf{d}$ vanishes from height $W+\mathfrak{h}(c)+8h_\sharp(c)$ on.
By $\mathfrak{h}(c)\le h_\sharp(c)$, this height is at most $W+9h_\sharp(c)$. Together with (e2),
this gives $\operatorname{supp}\mathsf{d}\subset\mathsf{R}(c)$.

By \eqref{8.13:eq-clipped-competitor-1}, $\mathsf{R}(c)$ lies inside the bottom band. By the model
condition \eqref{8.13:eq-model-conditions-d} it lies below height $H/2$. Hence, for
$y\le\lfloor H/2\rfloor$, $\mathsf{D}$ consists of the piece $\mathsf{d}$ on $\mathsf{R}(c)$ and
the piece $-\mathsf{d}(\cdot-P,\cdot)$ on the translate of $\mathsf{R}(c)$ by $P$ in $x$. For
$y>\lfloor H/2\rfloor$, $\mathsf{D}$ consists of the mirror images of these two pieces under
$y\mapsto H-y$, which lie above height $H/2$. The four pieces have pairwise disjoint supports.
Each is congruent to $\mathsf{d}$ up to sign, by a translation by $P$, the reflection
$y\mapsto H-y$, or both. Therefore $\|\nabla\mathsf{D}\|^2=4\|\nabla\mathsf{d}\|^2$.
\end{proof}

\begin{lemma}[Model: clipping preserves the tangential Lipschitz bound]\label{8.13:lem-clipping-lipschitz}
The model is that of the clipped competitor, Definition~\ref{8.13:def-clipped-competitor}, which
is a model statement. It works on the periodic strip $\Lambda$ of that definition, with the
field $\tilde{\psi}$, its vertical bond differences $\tilde F_y$, the young cap $c$, the heights
$W$ and $H$, the half-period $P$, the clipping function $\mathsf{d}=\mathsf{d}_c$, its extension
$\mathsf{D}$, and the clipped competitor $\psi^{c}=\tilde{\psi}-\mathsf{D}$. Let $\sigma$ be the
slope constant of the model, for which the comparison principle for $\tilde{\psi}$ gives
$|\tilde{\psi}(x+1,y)-\tilde{\psi}(x,y)|\le\sigma$. Then
\begin{equation}\label{8.13:eq-clipping-lipschitz-1}
|\psi^{c}(x+1,y)-\psi^{c}(x,y)|\le\sigma\qquad\text{for all }(x,y)\in\Lambda .
\end{equation}
Hence $|\nabla_x\mathsf{d}|\le2\sigma$ wherever $\mathsf{d}$ is defined, where
$\nabla_x\mathsf{d}(x,y)=\mathsf{d}(x+1,y)-\mathsf{d}(x,y)$.
\end{lemma}

\begin{proof}
We treat the rows $y\ge W$ near the bottom tip. There, by the support statement in
\eqref{8.13:eq-clipped-competitor-a} for the clipped competitor (a model statement),
$\mathsf{D}=\mathsf{d}$. Elsewhere $\psi^{c}$ equals $\tilde{\psi}$, and the bound is the
comparison principle, or $\psi^{c}$ is a symmetric image of the field treated below.

Fix $x$. For the rows $y$ at which $\mathsf{d}$ is defined we put
\begin{equation*}
e(y):=\tilde{\psi}(x+1,y)-\tilde{\psi}(x,y),\qquad
\delta(y):=\mathsf{d}(x+1,y)-\mathsf{d}(x,y).
\end{equation*}
The letter $\delta$ is local to this proof. By the comparison principle,
$e(y)\in[-\sigma,\sigma]$. Where $\mathsf{D}=\mathsf{d}$ at both columns $x$ and $x+1$, the
definition $\psi^{c}=\tilde{\psi}-\mathsf{D}$ gives
\begin{equation}\label{8.13:eq-clipping-lipschitz-2}
\psi^{c}(x+1,y)-\psi^{c}(x,y)=e(y)-\delta(y).
\end{equation}
Since $\tilde F_y(x)=\tilde{\psi}(x,y+1)-\tilde{\psi}(x,y)$, we have
\begin{equation}\label{8.13:eq-clipping-lipschitz-3}
\tilde F_y(x+1)-\tilde F_y(x)=e(y+1)-e(y).
\end{equation}

We prove $e(y)-\delta(y)\in[-\sigma,\sigma]$ by induction on $y\ge W$. The column recursion
of the clipped competitor sets $\mathsf{d}(\cdot,W)=0$, so $\delta(W)=0$. Hence
$e(W)-\delta(W)=e(W)\in[-\sigma,\sigma]$.

Suppose the claim holds at $y$, and let $y+1\le\lfloor H/2\rfloor$. Put
\begin{equation*}
a_0:=\mathsf{d}(x,y)+\tilde F_y(x)-c,\qquad a_1:=\mathsf{d}(x+1,y)+\tilde F_y(x+1)-c.
\end{equation*}
The recursion $\mathsf{d}(x,y+1)=(\mathsf{d}(x,y)+\tilde F_y(x)-c)^{+}$, applied at the columns
$x$ and $x+1$, gives $\delta(y+1)=a_1^{+}-a_0^{+}$. Subtracting the two definitions and using
\eqref{8.13:eq-clipping-lipschitz-3} gives
\begin{equation}\label{8.13:eq-clipping-lipschitz-4}
a_1-a_0=\delta(y)+e(y+1)-e(y).
\end{equation}
There are four cases.

If $a_0\le0$ and $a_1\le0$, then $\delta(y+1)=0$. Hence $e(y+1)-\delta(y+1)=e(y+1)$, and
$|e(y+1)|\le\sigma$.

If $a_0>0$ and $a_1>0$, then $\delta(y+1)=a_1-a_0$. By \eqref{8.13:eq-clipping-lipschitz-4},
\begin{equation*}
e(y+1)-\delta(y+1)=e(y)-\delta(y)\in[-\sigma,\sigma],
\end{equation*}
by the induction hypothesis.

If $a_1>0\ge a_0$, then $\delta(y+1)=a_1$, and $0<\delta(y+1)\le a_1-a_0$. By
\eqref{8.13:eq-clipping-lipschitz-4},
\begin{equation*}
e(y+1)-\delta(y+1)\in\bigl[e(y)-\delta(y),\,e(y+1)\bigr)\subset[-\sigma,\sigma],
\end{equation*}
because both endpoints lie in $[-\sigma,\sigma]$: the left one by the induction hypothesis and
the right one by the comparison principle.

If $a_0>0\ge a_1$, then $\delta(y+1)=-a_0$, and $a_1-a_0\le\delta(y+1)<0$. By
\eqref{8.13:eq-clipping-lipschitz-4},
\begin{equation*}
e(y+1)-\delta(y+1)\in\bigl(e(y+1),\,e(y)-\delta(y)\bigr]\subset[-\sigma,\sigma],
\end{equation*}
again because both endpoints lie in $[-\sigma,\sigma]$. This completes the induction.

By \eqref{8.13:eq-clipping-lipschitz-2}, the claim just proved is
\eqref{8.13:eq-clipping-lipschitz-1} near the bottom tip. The induction used only the column
recursion, so the claim holds at every $x$. The triangle inequality then gives
\begin{equation*}
|\nabla_x\mathsf{d}(x,y)|=|\delta(y)|\le|e(y)|+|e(y)-\delta(y)|\le2\sigma
\end{equation*}
for $y\ge W$. For $y\le W$ the function $\mathsf{d}$ vanishes.
\end{proof}

\begin{lemma}[Model: feasibility and energy of the competitor]\label{8.13:lem-competitor-energy}
We work in the strip model of this chapter. Fields $\psi$ live on the periodic strip
$\Lambda=\mathbb{Z}_N\times\{0,\dots,H\}$, $N=2P$, and $\|x\|$ denotes the distance from $x$ to
$N\mathbb{Z}$. For two fields we write $\langle\nabla\psi,\nabla\chi\rangle$ for the sum over the
nearest-neighbour bonds of $\Lambda$ of the products of the differences of $\psi$ and $\chi$ across the
bond, and we put $\|\nabla\psi\|^2=\langle\nabla\psi,\nabla\psi\rangle$. The Dirichlet energy is
$E(\psi)=\tfrac12\|\nabla\psi\|^2$, and $\mathcal{K}_c$ is the feasible set at young cap $c$. The field
$\tilde{\psi}$ is the minimiser of $E$ over $\mathcal{K}_\infty$, and
$\tilde F_y(x)=\tilde{\psi}(x,y+1)-\tilde{\psi}(x,y)$.

Let $c\ge\tau'$. Assume the standing hypotheses of the model under which
Lemma~\ref{8.13:lem-cool-band}, Corollary~\ref{8.13:cor-hot-box} and the facts
\eqref{8.13:eq-clipped-competitor-a} hold. Let $\mathsf{d}=\mathsf{d}_c$, $\mathsf{D}$ and
$\psi^{c}=\tilde{\psi}-\mathsf{D}$ be as in Definition~\ref{8.13:def-clipped-competitor}. Then
\begin{equation}\label{8.13:eq-competitor-energy-a}
\psi^{c}\in\mathcal{K}_c,\qquad
E(\psi^{c})-E(\tilde{\psi})=\tfrac12\|\nabla\mathsf{D}\|^2=2\|\nabla\mathsf{d}\|^2 .
\end{equation}
\end{lemma}

\begin{proof}
Write $F^{c}_y(x)=\psi^{c}(x,y+1)-\psi^{c}(x,y)$. Then
$F^{c}_y=\tilde F_y-\nabla_y\mathsf{D}$, where
$\nabla_y\mathsf{D}(x,y)=\mathsf{D}(x,y+1)-\mathsf{D}(x,y)$.

\emph{Feasibility.} We check the constraints of $\mathcal{K}_c$ region by region.

First, the rows $y\le W$ and $y\ge H-W$. The Dirichlet rows and the lower tiers lie in these rows, and
$\mathsf{D}=0$ there. Indeed, $\mathsf{d}(x,y)=0$ for $y\le W$, so $\mathsf{D}(x,y)=0$ for $y\le W$.
For $y\ge H-W$ we have $\mathsf{D}(x,y)=\mathsf{D}(x,H-y)$ with $H-y\le W$, so again
$\mathsf{D}(x,y)=0$. Hence these rows are untouched.

Next, the young vertical bonds. By the model definition and its facts
\eqref{8.13:eq-clipped-competitor-a}:
\begin{itemize}
\item $\mathsf{d}$ is supported in the columns with $\|x\|<\kappa W$;
\item $\mathsf{d}(\cdot-P,\cdot)$ is supported in the columns with $\|x-P\|<\kappa W$;
\item $\mathsf{D}=\mathsf{d}$ near the bottom tip ($\|x\|<\kappa W$) and
$\mathsf{D}=-\mathsf{d}(\cdot-P,\cdot)$ near the top tip ($\|x-P\|<\kappa W$).
\end{itemize}

Near the bottom tip, the same facts give $F^{c}_y=\min(c,\tilde F_y+\mathsf{d})\le c$. Since
$\mathsf{d}\ge0$, also $F^{c}_y\ge\min(c,\tilde F_y)$. By part~(b) of the model corollary,
Corollary~\ref{8.13:cor-hot-box}, $\tilde F_y\ge-c$ there, and therefore $F^{c}_y\ge-c$.

On the columns with $\kappa W\le\|x\|\le P-\kappa W$ we have $\mathsf{D}=0$, so $F^{c}=\tilde F$. By the
model lemma, Lemma~\ref{8.13:lem-cool-band}, $\tilde F\in[-\tau'/2,\tau'/2]$ there, and this interval
lies in $[-c,c]$ because $c\ge\tau'$.

Near the top tip, the symmetry $\tilde F_y(x)=-\tilde F_y(x-P)$ of $\tilde{\psi}$ under the half-period
shift gives
\begin{align*}
F^{c}_y(x)&=\tilde F_y(x)+\nabla_y\mathsf{d}(x-P,y)\\
&=-\bigl[\tilde F_y(x-P)-\nabla_y\mathsf{d}(x-P,y)\bigr]\\
&=-F^{c}_y(x-P)\in[-c,c].
\end{align*}
The last step uses that $x-P$ lies near the bottom tip, where $\mathsf{D}=\mathsf{d}$ and the previous
case applies.

The young vertical bonds of the top band follow from those of the bottom band by the reflection symmetry
in $y$, $F^{c}_{H-1-y}=-F^{c}_y$. For the $\mathsf{D}$-part this is the identity
$\mathsf{D}(x,H-y)=\mathsf{D}(x,y)$.

Finally, the young horizontal bonds. By the model lemma, Lemma~\ref{8.13:lem-clipping-lipschitz},
\[
|\psi^{c}(x+1,y)-\psi^{c}(x,y)|\le\sigma<c .
\]
Hence $\psi^{c}\in\mathcal{K}_c$.

\emph{Energy.} Since $E$ is quadratic,
\[
E(\tilde{\psi}-\mathsf{D})-E(\tilde{\psi})
=-\langle\nabla\tilde{\psi},\nabla\mathsf{D}\rangle+\tfrac12\|\nabla\mathsf{D}\|^2 .
\]
Let $\Delta$ be the graph Laplacian of the strip, $(\Delta\psi)(u)=\sum_{v\sim u}(\psi(v)-\psi(u))$.
Summation by parts on the finite graph $\Lambda$, which produces no boundary terms, gives
\[
\langle\nabla\tilde{\psi},\nabla\mathsf{D}\rangle
=\sum_{u\in\Lambda}\mathsf{D}(u)\sum_{v\sim u}\bigl(\tilde{\psi}(u)-\tilde{\psi}(v)\bigr)
=-\sum_{u\in\Lambda}\mathsf{D}(u)\,(\Delta\tilde{\psi})(u).
\]
By \eqref{8.13:eq-clipped-competitor-a}, $\mathsf{D}$ is supported in the rows
$W+1,\dots,H-W-1$. Every site of these rows has four neighbours, so $\Delta$ is the $4$-neighbour
Laplacian there, and $\tilde{\psi}$ is harmonic there. Hence
$\langle\nabla\tilde{\psi},\nabla\mathsf{D}\rangle=0$ and
$E(\psi^{c})-E(\tilde{\psi})=\tfrac12\|\nabla\mathsf{D}\|^2$.

By \eqref{8.13:eq-clipped-competitor-a}, $\mathsf{D}$ consists of four disjoint congruent pieces, so
$\|\nabla\mathsf{D}\|^2=4\|\nabla\mathsf{d}\|^2$. This gives \eqref{8.13:eq-competitor-energy-a}.
\end{proof}

\begin{lemma}[Model: the local energy budget]\label{8.13:lem-energy-budget}
This lemma is a statement about the strip model of this section: on the periodic strip $\Lambda$, the
field $\psi^{(c)}$ minimises the quadratic energy $E$ over the feasible set $\mathcal{K}_c$ at young
cap $c$, and $\tilde\psi$ minimises $E$ over $\mathcal{K}_\infty$, as in the model lemma,
Lemma~\ref{8.13:lem-variational-inequality}. For a function $u$ on $\Lambda$ we write
$\|\nabla u\|^2$ for the sum, over all bonds of $\Lambda$, of the square of the difference of $u$
along the bond.

Assume the mild-fringe hypothesis of the model, Definition~\ref{8.13:def-mild-fringe}, and the
conditions \eqref{8.13:eq-model-conditions-a}, \eqref{8.13:eq-model-conditions-c} and
\eqref{8.13:eq-model-conditions-d} of the model, Definition~\ref{8.13:def-model-conditions}. Let
$c\ge\tau'$. Put $v:=\psi^{(c)}-\tilde\psi$, and let $\mathsf{d}_c$ be the clipping function of the
model, Definition~\ref{8.13:def-clipped-competitor}. Then
\begin{equation}\label{8.13:eq-energy-budget-1}
\|\nabla v\|^2\le 4\|\nabla\mathsf{d}_c\|^2\le 4\Delta(c),
\end{equation}
where, with $h_{\sharp}(c)$ the hot-box scale of \eqref{8.13:eq-hot-box-a},
\begin{equation}\label{8.13:eq-energy-budget-a}
\Delta(c):=C_d\,\sigma\,(c+\beta_\kappa+19\sigma)\,h_{\sharp}(c)^2,\qquad C_d:=144.
\end{equation}
\end{lemma}

\begin{proof}
Every statement and every display cited in this proof belongs to the strip model.
Write $\mathsf{d}=\mathsf{d}_c$ and $h_{\sharp}=h_{\sharp}(c)$. Put
$\nabla_x\mathsf{d}(x,y):=\mathsf{d}(x+1,y)-\mathsf{d}(x,y)$ and
$\nabla_y\mathsf{d}(x,y):=\mathsf{d}(x,y+1)-\mathsf{d}(x,y)$. These are the differences along the
horizontal and the vertical bonds. By the definition in \eqref{8.13:eq-hot-box-a}, $h_{\sharp}\ge1$.

\emph{The first inequality.} Let $\psi^{c}=\tilde\psi-\mathsf{D}$ be the clipped competitor of
Definition~\ref{8.13:def-clipped-competitor}. Under the hypotheses of the present lemma, the model
lemma, Lemma~\ref{8.13:lem-competitor-energy}, gives $\psi^{c}\in\mathcal{K}_c$ and
\eqref{8.13:eq-competitor-energy-a}, namely
\[
E(\psi^{c})-E(\tilde\psi)=\tfrac12\|\nabla\mathsf{D}\|^2=2\|\nabla\mathsf{d}\|^2 .
\]
Since $c\ge\tau'$ and $\psi^{c}\in\mathcal{K}_c$, we may insert $\psi^{c}$ as the competitor in
\eqref{8.13:eq-variational-inequality-a}. This gives
\[
\tfrac12\|\nabla v\|^2\le E(\psi^{c})-E(\tilde\psi)=2\|\nabla\mathsf{d}\|^2 .
\]
Multiplying by $2$ gives the first inequality in \eqref{8.13:eq-energy-budget-1}.

\emph{Localisation.} By the support statement in \eqref{8.13:eq-clipped-competitor-a},
\[
\operatorname{supp}\mathsf{d}\subset\mathsf{R}(c)
:=\{\|x\|<h_{\sharp}\}\times\{W+1\le y\le W+9h_{\sharp}\}.
\]
So a bond contributes to $\|\nabla\mathsf{d}\|^2$ only if one of its endpoints lies in
$\mathsf{R}(c)$. The columns $x$ with $\|x\|<h_{\sharp}$ form a set of consecutive columns. There are
at most $2h_{\sharp}+1$ of them, and $2h_{\sharp}+1\le3h_{\sharp}$.

\emph{Horizontal bonds.} The rows meeting $\mathsf{R}(c)$ number at most $9h_{\sharp}$. In each such
row, the horizontal bonds with an endpoint in one of these at most $2h_{\sharp}+1$ consecutive columns
number at most $2h_{\sharp}+2$. So at most $(2h_{\sharp}+2)\cdot9h_{\sharp}$ horizontal bonds
meet $\mathsf{R}(c)$. On each of them $(\nabla_x\mathsf{d})^2\le4\sigma^2$, by the model lemma,
Lemma~\ref{8.13:lem-clipping-lipschitz}. Since $2h_{\sharp}+2\le4h_{\sharp}$, the horizontal bonds
contribute at most
\[
36\sigma^2h_{\sharp}(2h_{\sharp}+2)\le144\sigma^2h_{\sharp}^2 .
\]

\emph{Vertical bonds: increasing steps.} Fix a column $x$ with $\|x\|<h_{\sharp}$. By
\eqref{8.13:eq-clipped-competitor-a},
\[
\nabla_y\mathsf{d}(x,y)=\max\bigl(\tilde F_y(x)-c,\,-\mathsf{d}(x,y)\bigr)
\quad\text{and}\quad \mathsf{d}\ge0 .
\]
Hence a step increases $\mathsf{d}$ exactly when $\tilde F_y(x)>c$, and then
$\nabla_y\mathsf{d}=(\tilde F_y(x)-c)^{+}$. By the height statement in
\eqref{8.13:eq-clipped-competitor-a}, $\mathsf{d}$ increases only at bonds $(x,W+m)$ of the hot set
$A(c)$, with $m<h_{\sharp}$. Part~(a) of the model corollary, Corollary~\ref{8.13:cor-hot-box}, gives
on these bonds
\[
0<\tilde F-c\le12\sigma\log\frac{2h_{\sharp}+1}{2m+1}.
\]
Each $m$ occurs at most once in the column. The function
$t\mapsto\log^2\bigl((2h_{\sharp}+1)/(2t+1)\bigr)$ is non-negative and non-increasing on
$[0,h_{\sharp}]$. Therefore the term $m=0$ equals $\log^2(2h_{\sharp}+1)$. For $1\le m<h_{\sharp}$,
the term at $m$ is at most the integral of this function over $[m-1,m]\subset[0,h_{\sharp}]$. Hence
\begin{align*}
\sum\bigl((\tilde F-c)^{+}\bigr)^2
&\le\sum_{0\le m<h_{\sharp}}\Bigl[12\sigma\log\frac{2h_{\sharp}+1}{2m+1}\Bigr]^2\\
&\le144\sigma^2\Bigl[\log^2(2h_{\sharp}+1)
+\int_0^{h_{\sharp}}\log^2\frac{2h_{\sharp}+1}{2t+1}\,dt\Bigr].
\end{align*}
Substitute $u=(2t+1)/(2h_{\sharp}+1)$ and use $\int_0^1\log^2(1/u)\,du=2$:
\[
\int_0^{h_{\sharp}}\log^2\frac{2h_{\sharp}+1}{2t+1}\,dt
=\frac{2h_{\sharp}+1}{2}\int_{1/(2h_{\sharp}+1)}^{1}\log^2\frac1u\,du\le2h_{\sharp}+1 .
\]
Since $\log^2t\le t$ for $t\ge1$, we also have $\log^2(2h_{\sharp}+1)\le2h_{\sharp}+1$. With
$2\le2h_{\sharp}$, the increasing steps of the column therefore contribute at most
\[
144\sigma^2\bigl[\log^2(2h_{\sharp}+1)+2h_{\sharp}+1\bigr]
\le144\sigma^2(4h_{\sharp}+2)\le864\sigma^2h_{\sharp}.
\]

\emph{Vertical bonds: decreasing steps.} On a step that decreases $\mathsf{d}$, the formula above gives
\[
|\nabla_y\mathsf{d}|=\min(\mathsf{d},\,c-\tilde F)\le c+\beta_\kappa .
\]
The last inequality holds because $\tilde F\ge-\beta_\kappa$ on these columns, by part~(b) of the model
corollary, Corollary~\ref{8.13:cor-hot-box}.

Next we bound the total decrease. The function $\mathsf{d}$ vanishes at the bottom of the column and
stays non-negative. So at every height, the total decrease so far is at most the total increase so
far. The total increase of the column satisfies
\[
\sum(\tilde F-c)^{+}\le48\sigma h_{\sharp},
\]
by the height statement in \eqref{8.13:eq-clipped-competitor-a}. The sum of the squares of the
decreasing steps is at most their maximum times their sum, so it is at most
$48\sigma h_{\sharp}(c+\beta_\kappa)$. Steps on which $\mathsf{d}$ is constant contribute nothing.

\emph{Summation.} Adding the increasing and decreasing steps, the vertical bonds of one column
contribute at most
\[
\sigma h_{\sharp}\bigl[864\sigma+48(c+\beta_\kappa)\bigr].
\]
There are at most $3h_{\sharp}$ columns, so all vertical bonds together contribute at most
\[
3\sigma h_{\sharp}^2\bigl[864\sigma+48(c+\beta_\kappa)\bigr]
=144\sigma h_{\sharp}^2\,[18\sigma+c+\beta_\kappa].
\]
Adding the horizontal contribution $144\sigma^2h_{\sharp}^2$ gives
\[
\|\nabla\mathsf{d}\|^2\le144\sigma h_{\sharp}^2\,[19\sigma+c+\beta_\kappa]=\Delta(c).
\]
This is the second inequality in \eqref{8.13:eq-energy-budget-1}.
\end{proof}

\begin{remark}[Model: monotonicity and locality of the energy budget]\label{8.13:rem-budget-monotone}
This remark concerns the model of Lemma~\ref{8.13:lem-energy-budget}, whose notation it uses: the
Dirichlet energy $E$, the feasible set $\mathcal{K}_c$ at young cap $c$, the minimiser $\psi^{(c)}$
of $E$ over $\mathcal{K}_c$, and the minimiser $\tilde{\psi}$ under the straddling cap alone.
We write $\tilde F$ for the vertical bond differences of
$\tilde{\psi}$. The energy budget $\Delta(c)$ is the quantity defined in
\eqref{8.13:eq-energy-budget-a} of the model lemma, Lemma~\ref{8.13:lem-energy-budget}.
\begin{enumerate}
\item[(i)] The budget $\Delta(c)$ is decreasing in $c$ on $c\ge\tau'$. Indeed, by the definition
of the hot-box scale $h_{\sharp}(c)$, the budget $\Delta(c)$ is a $c$-independent multiple of
\begin{equation*}
(c+\beta_\kappa+19\sigma)\,e^{-(c-\beta_\kappa)/(6\sigma)},
\end{equation*}
whose derivative in $c$ is
\begin{equation*}
e^{-(c-\beta_\kappa)/(6\sigma)}\Bigl(1-\frac{c+\beta_\kappa+19\sigma}{6\sigma}\Bigr)
=-\,e^{-(c-\beta_\kappa)/(6\sigma)}\,\frac{c+\beta_\kappa+13\sigma}{6\sigma}.
\end{equation*}
This is negative for every $c>-\beta_\kappa-13\sigma$, hence on the whole range $c\ge\tau'$ of
positive caps, since $\beta_\kappa=24\tau/\kappa$ and $\sigma$ are positive; in particular the
expression decreases once $c\ge 6\sigma$. Hence the budget at every cap $c\ge\tau'$
is at most the budget at $\tau'$.
\item[(ii)] Under the dictionary of this chapter, with $\theta_y=\tau'$ the young cap, we have
$\Delta(\theta_y)\asymp\sigma\,\theta_y\,h_{\sharp}(\theta_y)^2$, measured in fine lattice
units. Here $\asymp$ means that the two sides are bounded by constant multiples of each
other. Every use of this budget in this chapter compares it with a quantity of the same
scaling, and the conclusions drawn are therefore free of the cutoff.
\item[(iii)] The competitor is local. The clipping function $\mathsf{d}_c$ from which it is
built is carried by its support box $\mathsf{R}(c)$, which has of the order of
$h_{\sharp}(c)\times 9h_{\sharp}(c)$ sites and lies at the tip. The energy argument is
therefore local. Its lateral locality is exactly the point at which the mild-fringe
hypothesis of the model lemma, Lemma~\ref{8.13:lem-energy-budget}, with lateral fraction
$\kappa$, enters.
\item[(iv)] Let $c_{\max}$ be the maximum of $\tilde F$ over the bottom band. For every
$c\ge c_{\max}$ we have $\tilde{\psi}\in\mathcal{K}_c$, $\psi^{(c)}=\tilde{\psi}$ and
$\mathsf{d}_c\equiv0$. For such caps the model lemma has nothing to prove.
\end{enumerate}
Items (i) and (iv) together read
\begin{equation}\label{8.13:eq-budget-monotone-a}
\begin{gathered}
\Delta(c)\le\Delta(\tau')\quad\text{for all } c\ge\tau',\\
\tilde{\psi}\in\mathcal{K}_c,\quad \psi^{(c)}=\tilde{\psi},\quad \mathsf{d}_c\equiv0
\quad\text{for all } c\ge c_{\max}.
\end{gathered}
\end{equation}
\end{remark}

\begin{lemma}[Model: a transported segment costs cubic energy]\label{8.13:lem-segment-cost}
This lemma concerns the two-tier box-clamp lattice model of
Definition~\ref{8.13:def-two-tier-model}, which is a model and not Ba{\l}aban's step. Its fields
are real functions on the periodic strip $\Lambda=\mathbb{Z}_N\times\{0,1,\dots,H\}$, $N=2P$,
with the kinked Dirichlet datum $\varphi_0(x)=\sigma\|x\|$ and the Dirichlet energy
$E(\psi)=\tfrac12\|\nabla\psi\|^2$. Here $\|\nabla\psi\|^2$ is the sum, over all bonds of
$\Lambda$, of the squared differences of $\psi$ along them. The vertical bond differences are
$F_y(x)=\psi(x,y+1)-\psi(x,y)$.

Let $c\ge\tau'$. Let $\psi^{(c)}$ be the minimiser of $E$ over the feasible set $\mathcal{K}_c$
with young cap $c$, and let $\tilde\psi$ be the minimiser of $E$ over $\mathcal{K}_\infty$.
Write $F^{(c)}_y$ and $\tilde F_y$ for their vertical bond differences, and put
$v:=\psi^{(c)}-\tilde\psi$. Let $x_0\in\mathbb{Z}_N$, let $Y_1,Y_2$ be integers with
$W\le Y_1<Y_2\le W'$, and let $\gamma>0$. Suppose that one of the following holds:
\begin{itemize}
\item[(a)] for every $y$ with $Y_1\le y<Y_2$, the young bond $(x_0,y)$ is up-tight at
$\psi^{(c)}$ and $\tilde F_y(x_0)\le c-\gamma$;
\item[(b)] for every $y$ with $Y_1\le y<Y_2$, the young bond $(x_0,y)$ is down-tight at
$\psi^{(c)}$ and $\tilde F_y(x_0)\ge -c+\gamma$.
\end{itemize}
Put $\Gamma:=\gamma(Y_2-Y_1)$. Then
\begin{equation}\label{8.13:eq-segment-cost-1}
\|\nabla v\|^2\ \ge\ \min\Bigl\{\frac{\Gamma}{8\sigma},\,N\Bigr\}\cdot\frac{\Gamma^2}{16\,Y_2}.
\end{equation}
\end{lemma}

\begin{proof}
Consider case (a) first. For $Y_1\le y<Y_2$, up-tightness of the bond $(x_0,y)$ at
$\psi^{(c)}$ means that its difference sits at the young cap, so $F^{(c)}_y(x_0)=c$. The
hypothesis then gives $F^{(c)}_y(x_0)-\tilde F_y(x_0)\ge c-(c-\gamma)=\gamma$. Since
$v=\psi^{(c)}-\tilde\psi$, the differences telescope along the column $x_0$:
\begin{equation*}
v(x_0,Y_2)-v(x_0,Y_1)=\sum_{Y_1\le y<Y_2}\bigl[F^{(c)}_y-\tilde F_y\bigr](x_0)
\ \ge\ \gamma\,(Y_2-Y_1)=\Gamma .
\end{equation*}
Note that $\Gamma>0$. If both $|v(x_0,Y_1)|$ and $|v(x_0,Y_2)|$ were smaller than $\Gamma/2$,
their difference would be smaller than $\Gamma$. Hence there is $Y_0\in\{Y_1,Y_2\}$ with
$|v(x_0,Y_0)|\ge\Gamma/2$.

By the comparison principle of the model, applied to both fields $\psi^{(c)}$ and $\tilde\psi$,
the row $v(\cdot,Y_0)$ is $2\sigma$-Lipschitz with respect to the periodic distance $\|\cdot\|$.
Put
\begin{equation*}
I:=\Bigl\{x\in\mathbb{Z}_N:\ \|x-x_0\|\le\frac{\Gamma}{8\sigma}\Bigr\}.
\end{equation*}
For $x\in I$ we have
\begin{equation*}
|v(x,Y_0)|\ \ge\ |v(x_0,Y_0)|-2\sigma\|x-x_0\|\ \ge\ \frac{\Gamma}{2}-\frac{\Gamma}{4}
=\frac{\Gamma}{4}.
\end{equation*}
Write $r:=\Gamma/(8\sigma)$. The set $I$ has $\min\{2\lfloor r\rfloor+1,\,N\}$ elements. Since
$2\lfloor r\rfloor+1\ge r$ for every $r\ge0$, it follows that $|I|\ge\min\{\Gamma/(8\sigma),N\}$.

Now fix $x\in I$. The function $v$ vanishes on the Dirichlet rows of the model, so $v(x,0)=0$.
In particular $Y_0\ge1$, because $|v(x_0,Y_0)|\ge\Gamma/2>0$. Write
$\nabla_yv(x,y)=v(x,y+1)-v(x,y)$. Then $v(x,Y_0)=\sum_{0\le y<Y_0}\nabla_yv(x,y)$, and the
Cauchy--Schwarz inequality along the column gives $v(x,Y_0)^2\le Y_0\sum_{y<Y_0}(\nabla_yv(x,y))^2$.
Using $Y_0\le Y_2$ and $|v(x,Y_0)|\ge\Gamma/4$, we obtain
\begin{equation}\label{8.13:eq-segment-cost-2}
\sum_{0\le y<Y_0}\bigl(\nabla_yv(x,y)\bigr)^2\ \ge\ \frac{v(x,Y_0)^2}{Y_0}
\ \ge\ \frac{\Gamma^2}{16\,Y_2}.
\end{equation}
For distinct $x\in I$ the vertical bonds occurring in \eqref{8.13:eq-segment-cost-2} are
distinct bonds of $\Lambda$, and $\|\nabla v\|^2$ contains all of their squared differences
(together with further non-negative terms). Summing \eqref{8.13:eq-segment-cost-2} over $x\in I$
therefore gives
\begin{equation*}
\|\nabla v\|^2\ \ge\ |I|\,\frac{\Gamma^2}{16\,Y_2}
\ \ge\ \min\Bigl\{\frac{\Gamma}{8\sigma},\,N\Bigr\}\cdot\frac{\Gamma^2}{16\,Y_2},
\end{equation*}
which is \eqref{8.13:eq-segment-cost-1}.

In case (b) the signs are reversed. Down-tightness gives $F^{(c)}_y(x_0)=-c$, and the hypothesis
gives $F^{(c)}_y(x_0)-\tilde F_y(x_0)\le -c-(-c+\gamma)=-\gamma$ for $Y_1\le y<Y_2$. Hence
$v(x_0,Y_2)-v(x_0,Y_1)\le-\Gamma$. From this point the argument above uses only
$|v(x_0,Y_2)-v(x_0,Y_1)|\ge\Gamma$ and absolute values, so it applies verbatim and yields
\eqref{8.13:eq-segment-cost-1}.
\end{proof}

\begin{corollary}[Model: no young-full column]\label{8.13:cor-no-full-column}
This is a statement about the strip model whose one-sided conditions are those of
Definition~\ref{8.13:def-model-conditions}: the minimiser
$\psi^{(c)}$ of the Dirichlet energy $E$ over the feasible set $\mathcal{K}_c$ with young cap $c$,
compared with the straddling-only minimiser $\tilde{\psi}$ of
Notation~\ref{8.13:not-straddling-minimiser}. A \emph{young-full} column is a column of the
full-column run type in the model's classification of runs into needles, full columns and
hanging poles; such a column is either up-full or down-full. Put
\begin{equation}\label{8.13:eq-no-full-column-a}
B_\beta:=\sigma\bigl[C_K^{0}+6\log^{+}(N/\bar W)\bigr].
\end{equation}
Assume the mild-fringe condition of the model at lateral fraction $\kappa$ and the one-sided
conditions of Definition~\ref{8.13:def-model-conditions}, in particular the top-of-band
floor~\eqref{8.13:eq-model-conditions-b} and the hot-box and needle-fit
conditions~\eqref{8.13:eq-model-conditions-c}. Then for every cap $c\ge\tau'$ the minimiser
$\psi^{(c)}$ has no young-full column.
\end{corollary}

\begin{proof}
Fix $c\ge\tau'$ and put $v:=\psi^{(c)}-\tilde{\psi}$. By the model statement,
Notation~\ref{8.13:not-straddling-minimiser}, $H'\ge2\bar W+4$, so that $\bar W\le(H'-1)/2$;
moreover $\bar W\ge4$ in the model instance.

\emph{Up-full columns.} Suppose that $x_0$, anywhere on the strip, is an up-full column. Take
$Y_1:=W+\lceil\bar W/2\rceil$ and $Y_2:=W'$. Since $\bar W\ge4$ we have $W\le Y_1<Y_2\le W'$,
and the bonds $(x_0,y)$, $Y_1\le y<Y_2$, are up-tight at $\psi^{(c)}$ because the column is full
(young bonds and up- and down-tightness in the sense of the model's set-up).
For such $y$ put $m:=y-W$, so that $m\ge\bar W/2$ and $m\le\bar W-1<H'-1$. The logarithmic strip
bound of the model statement, Notation~\ref{8.13:not-straddling-minimiser}, gives
$\tilde F_y(x_0)\le b(m)$, with $b(m)$ and the index $m_*$ as
in~\eqref{8.13:eq-straddling-minimiser-a}. Since $m\ge\bar W/2$ and $\bar W\le(H'-1)/2$, we
have $2m_*+1\ge\bar W$; since $\log^{+}$ is non-decreasing and $N/(2m_*+1)\le N/\bar W$, it
follows that $b(m)\le B_\beta$. By the top-of-band floor~\eqref{8.13:eq-model-conditions-b}
and $c\ge\tau'$,
\begin{equation*}
\tilde F_y(x_0)\le b(m)\le B_\beta\le \tfrac{\tau'}{2}-3\sigma\le\tfrac{\tau'}{2}\le\tfrac{c}{2},
\qquad Y_1\le y<Y_2 .
\end{equation*}
Thus $\tilde F_y(x_0)\le c-\gamma$ with $\gamma:=c/2>0$. Since $\bar W\ge4$,
\begin{equation*}
Y_2-Y_1=\bar W-\lceil\bar W/2\rceil=\lfloor\bar W/2\rfloor\ge\tfrac{\bar W-1}{2}\ge\tfrac{\bar W}{3},
\qquad
\Gamma:=\gamma(Y_2-Y_1)=\tfrac{c}{2}\lfloor\bar W/2\rfloor\ge\tfrac{c\bar W}{6}.
\end{equation*}
The hypotheses of the model lemma, Lemma~\ref{8.13:lem-segment-cost}, hold with these
$Y_1,Y_2,\gamma$. Its lower bound increases with $\Gamma$; with $\Gamma\ge c\bar W/6$,
$16\cdot36=576$ and $Y_2=W'$ it gives
\begin{equation}\label{8.13:eq-no-full-column-1}
\|\nabla v\|^2\ge\min\Bigl\{\frac{c\bar W}{48\sigma},\,N\Bigr\}\cdot
\frac{c^2\bar W^2}{576\,W'} .
\end{equation}
On the other hand the model lemma, Lemma~\ref{8.13:lem-energy-budget}, with $C_d=144$ gives
\begin{equation}\label{8.13:eq-no-full-column-2}
\|\nabla v\|^2\le4\Delta(c)=576\,\sigma\,(c+\beta_\kappa+19\sigma)\,h_{\sharp}(c)^2 .
\end{equation}
We show that the right-hand side of \eqref{8.13:eq-no-full-column-1} exceeds that of
\eqref{8.13:eq-no-full-column-2}, a contradiction. Note that
$c+\beta_\kappa+19\sigma=c\,\bigl(1+(\beta_\kappa+19\sigma)/c\bigr)$. If the
minimum in \eqref{8.13:eq-no-full-column-1} equals $c\bar W/(48\sigma)$, then, since
$48\cdot576\cdot576=110592\cdot144$, division by $\sigma c\,h_{\sharp}(c)^2$ shows that the
required strict inequality is equivalent to
\begin{equation}\label{8.13:eq-no-full-column-3}
\Bigl(\frac{c}{\sigma}\Bigr)^2\frac{\bar W^3}{W'h_{\sharp}(c)^2}
>110592\,C_d\Bigl(1+\frac{\beta_\kappa+19\sigma}{c}\Bigr).
\end{equation}
If the minimum equals $N$, then, since $576\cdot576=2304\cdot144$, the same division shows that
it is equivalent to
\begin{equation}\label{8.13:eq-no-full-column-4}
N\,\frac{c}{\sigma}\,\frac{\bar W^2}{W'h_{\sharp}(c)^2}
>2304\,C_d\Bigl(1+\frac{\beta_\kappa+19\sigma}{c}\Bigr).
\end{equation}
For $c=\tau'$, \eqref{8.13:eq-no-full-column-3} and \eqref{8.13:eq-no-full-column-4} are the two
inequalities displayed in~\eqref{8.13:eq-model-conditions-c}. They persist for every
$c\ge\tau'$. The right-hand sides decrease in $c$, because $1+(\beta_\kappa+19\sigma)/c$ does.
For the left-hand sides, $h_{\sharp}(c)^2=\Delta(c)/\bigl(C_d\sigma(c+\beta_\kappa+19\sigma)\bigr)$,
so
\begin{align*}
\Bigl(\frac{c}{\sigma}\Bigr)^2\frac{\bar W^3}{W'h_{\sharp}(c)^2}
&=\frac{C_d\,c^2(c+\beta_\kappa+19\sigma)\bar W^3}{\sigma\,W'\,\Delta(c)},\\
N\frac{c}{\sigma}\frac{\bar W^2}{W'h_{\sharp}(c)^2}
&=\frac{C_d\,N c\,(c+\beta_\kappa+19\sigma)\bar W^2}{W'\,\Delta(c)},
\end{align*}
and these increase in $c$. Indeed their numerators increase, and $\Delta(c)$ decreases on
$c\ge\tau'$ by item~(i) of the model remark, Remark~\ref{8.13:rem-budget-monotone}. Hence both
\eqref{8.13:eq-no-full-column-3} and \eqref{8.13:eq-no-full-column-4} hold at $c$, whichever
value the minimum takes. This contradicts \eqref{8.13:eq-no-full-column-1} together with
\eqref{8.13:eq-no-full-column-2}. The two conditions of~\eqref{8.13:eq-model-conditions-c} are
one-sided inequalities of the model instance and are not evaluated here. Hence no column is
up-full.

\emph{Down-full columns.} A down-full column $x_0$ with $\|x_0\|\le P/2$, near the bottom tip,
is excluded by the model lemma, Lemma~\ref{8.13:lem-anchored-down-run}: its young down-tight
bonds start at the clamp's face $y=W$, and, by that lemma, no young down-run anchored at the
clamp's face exists on a column $x_0$ with $\|x_0\|\le P/2$.

For a down-full column anywhere we use the mirror argument. The bound of the model statement,
Notation~\ref{8.13:not-straddling-minimiser}, is a bound on $|\tilde F_{W+m}|$. The estimates
above therefore give
\begin{equation*}
\tilde F_y(x_0)\ge -b(m)\ge -B_\beta\ge -\tfrac{c}{2}=-c+\gamma,\qquad Y_1\le y<Y_2,
\end{equation*}
with the same $Y_1$, $Y_2$, $\gamma$ and $\Gamma$. The down case of the model lemma,
Lemma~\ref{8.13:lem-segment-cost}, then
yields \eqref{8.13:eq-no-full-column-1} again. The comparison with
\eqref{8.13:eq-no-full-column-2} is unchanged and gives the same contradiction. Hence
$\psi^{(c)}$ has no young-full column.
\end{proof}

\begin{corollary}[Model: the tip needle is short]\label{8.13:cor-short-needle}
We work in the strip model studied in the present section. In it, $\psi^{(c)}$ is the minimiser of
the Dirichlet energy $E$ on the periodic strip $\Lambda$ over the feasible set $\mathcal{K}_c$ with
young cap $c$, and $\tilde{\psi}$ is the minimiser with the straddling cap only. Assume the conditions
\eqref{8.13:eq-model-conditions-c} and \eqref{8.13:eq-model-conditions-d} of the model instance,
Definition~\ref{8.13:def-model-conditions}, and let $c\ge\tau'$. Let $[W,W+m^{*})$ be a tip needle at
$\psi^{(c)}$, in the sense of type $(\mathrm{N})$ of the model census,
Lemma~\ref{8.13:lem-tight-census}. Then $m^{*}\le m_{\mathrm{ndl}}(c)$, where
\begin{equation}\label{8.13:eq-short-needle-a}
m_{\mathrm{ndl}}(c):=\max\bigl\{2\mathfrak{h}(c)+2,\ m_e(c)\bigr\},\qquad
m_e(c):=\Bigl[128\,C_d\,\frac{c+\beta_\kappa+19\sigma}{\sigma}\,W\,h_{\sharp}(c)^2\Bigr]^{1/3}.
\end{equation}
Moreover, for all $c\ge\tau'$,
\begin{equation*}
m_{\mathrm{ndl}}(c)\le m_{\mathrm{ndl}}(\tau')\le \tfrac12\min(W,\bar W)-2 .
\end{equation*}
\end{corollary}

\begin{proof}
Fix $c\ge\tau'$ and a tip needle $[W,W+m^{*})$ on a column $x_0$. Put $v:=\psi^{(c)}-\tilde{\psi}$.
By type $(\mathrm{N})$ of the model census (Lemma~\ref{8.13:lem-tight-census}), the bonds $(x_0,y)$
with $W\le y<W+m^{*}$ are young and up-tight at $\psi^{(c)}$, and $m^{*}\le\bar W-1$.

If $m^{*}\le 2\mathfrak{h}(c)+1$, then $m^{*}\le m_{\mathrm{ndl}}(c)$ and there is nothing to prove.

Assume from now on that $m^{*}>2\mathfrak{h}(c)+1$. Since $\mathfrak{h}(c)\ge0$ and $m^{*}$ is an
integer, we have $m^{*}\ge2$. Put
\begin{equation*}
Y_1:=W+\lceil m^{*}/2\rceil,\qquad Y_2:=W+m^{*}.
\end{equation*}
Then $Y_1-W\ge m^{*}/2>\mathfrak{h}(c)+\tfrac12$. Next, $Y_1<Y_2$ because
$\lceil m^{*}/2\rceil<m^{*}$ for $m^{*}\ge2$. Finally, $Y_2\le W+\bar W-1<W'$. Hence
$W\le Y_1<Y_2\le W'$.

Every bond $(x_0,y)$ with $Y_1\le y<Y_2$ has the form $y=W+m$ with $m\ge\mathfrak{h}(c)$. By item (c)
of the model hot-box corollary, Corollary~\ref{8.13:cor-hot-box}, with $\mathfrak{h}(c)$ as in
\eqref{8.13:eq-hot-box-a}, we get $\tilde F_y(x_0)\le c-6\sigma$ for these bonds. That bound holds
uniformly in $x$, so the needle's column requires no lateral information.

The hypotheses of the up case of the model segment-cost lemma, Lemma~\ref{8.13:lem-segment-cost},
therefore hold
with $\gamma=6\sigma$. Its quantity is $\Gamma=6\sigma(Y_2-Y_1)=6\sigma\lfloor m^{*}/2\rfloor$. We
claim $\Gamma\ge2\sigma m^{*}$. For even $m^{*}$ we have $\lfloor m^{*}/2\rfloor=m^{*}/2\ge m^{*}/3$.
For odd $m^{*}\ge3$ we have $\lfloor m^{*}/2\rfloor=(m^{*}-1)/2\ge m^{*}/3$.

The right side $\min\{\Gamma/(8\sigma),N\}\,\Gamma^2/(16Y_2)$ of that lemma is increasing in $\Gamma$.
We also have $m^{*}\le\bar W-1<N$ by \eqref{8.13:eq-model-conditions-d}, so
$\min\{m^{*}/4,N\}=m^{*}/4$. Hence
\begin{equation*}
\|\nabla v\|^2\ge\min\Bigl\{\frac{m^{*}}{4},N\Bigr\}\frac{4\sigma^2 (m^{*})^{2}}{16(W+m^{*})}
=\frac{m^{*}}{4}\cdot\frac{\sigma^2(m^{*})^{2}}{4(W+m^{*})}.
\end{equation*}

The model budget lemma, Lemma~\ref{8.13:lem-energy-budget}, bounds $\|\nabla v\|^2$ from above by
$4\Delta(c)$, with $\Delta(c)$ as in \eqref{8.13:eq-energy-budget-a}. Combining the two bounds gives
\begin{equation}\label{8.13:eq-short-needle-1}
\frac{\sigma^2(m^{*})^{3}}{16(W+m^{*})}\le\|\nabla v\|^2\le4\Delta(c)
=4C_d\,\sigma\,(c+\beta_\kappa+19\sigma)\,h_{\sharp}(c)^2 .
\end{equation}

There are two cases.

(1) Suppose $m^{*}\le W$. Then $W+m^{*}\le2W$, and \eqref{8.13:eq-short-needle-1} gives
\begin{equation*}
\frac{\sigma^2(m^{*})^{3}}{32W}\le4C_d\,\sigma\,(c+\beta_\kappa+19\sigma)\,h_{\sharp}(c)^2 .
\end{equation*}
Multiplying by $32W/\sigma^2$, this is
\begin{equation*}
(m^{*})^{3}\le128\,C_d\,\frac{c+\beta_\kappa+19\sigma}{\sigma}\,W\,h_{\sharp}(c)^2=m_e(c)^3,
\end{equation*}
that is, $m^{*}\le m_e(c)\le m_{\mathrm{ndl}}(c)$.

(2) Suppose $m^{*}>W$. Then $W+m^{*}<2m^{*}$, and \eqref{8.13:eq-short-needle-1} gives
$\sigma^2(m^{*})^{2}/32\le4\Delta(c)$, that is,
\begin{equation*}
m^{*}\le h_{\sharp}(c)\Bigl[128\,C_d\,\frac{c+\beta_\kappa+19\sigma}{\sigma}\Bigr]^{1/2}.
\end{equation*}
Item (i) of the model remark, Remark~\ref{8.13:rem-budget-monotone}, states that $\Delta$ is
decreasing on $c\ge\tau'$. Since $\Delta(c)=C_d\sigma(c+\beta_\kappa+19\sigma)h_{\sharp}(c)^2$, the
product $(c+\beta_\kappa+19\sigma)h_{\sharp}(c)^2$ is at most its value at $\tau'$. Hence the right
side above is at most $h_{\sharp}(\tau')[128C_d(\tau'+\beta_\kappa+19\sigma)/\sigma]^{1/2}$. By
condition (b) of \eqref{8.13:eq-model-conditions-c}, this is at most $W$, which contradicts
$m^{*}>W$. So this case is empty.

Monotonicity in $c$. Since $\sigma>0$, the exponential in the definition of $\mathfrak{h}(c)$ in
\eqref{8.13:eq-hot-box-a} decreases in $c$, so $\mathfrak{h}$ is decreasing. Moreover
$m_e(c)^3=128\,W\Delta(c)/\sigma^2$, which is proportional to
$(c+\beta_\kappa+19\sigma)h_{\sharp}(c)^2$. By item (i) of the model remark,
Remark~\ref{8.13:rem-budget-monotone},
$m_e(c)^3$ decreases on $c\ge\tau'$. Both entries of the maximum in \eqref{8.13:eq-short-needle-a}
therefore decrease, and $m_{\mathrm{ndl}}(c)\le m_{\mathrm{ndl}}(\tau')$ for $c\ge\tau'$. Condition (c)
of \eqref{8.13:eq-model-conditions-c} is the inequality
$m_{\mathrm{ndl}}(\tau')\le\frac12\min(W,\bar W)-2$.

The down-needles at the top tip, and the mirror images in the top band, are treated in the same way.
One uses the down case of the model segment-cost lemma, Lemma~\ref{8.13:lem-segment-cost}. In place
of item (c) of the model hot-box corollary, Corollary~\ref{8.13:cor-hot-box}, one uses the lower-side
bound $\tilde F_{W+m}(x)\ge-c+6\sigma$ for $m\ge\mathfrak{h}(c)$, which is the lower side of the
model's bound on the vertical differences $\tilde F$ of $\tilde{\psi}$ at the top tip.
\end{proof}

In the dictionary of the model instance (Definition~\ref{8.13:def-model-conditions}), at the cap
$c=\tau'=\theta_y$ and with $n=L^{k_o}$, the two entries of \eqref{8.13:eq-short-needle-a} read
\begin{equation*}
\frac{m_e(\theta_y)}{n}=\Bigl[128\,C_d\Bigl(\frac{\mathsf{T}}{s_e}+\frac{24}{\kappa s_e}+19\Bigr)
\frac{W}{n}\Bigl(\frac{h_{\sharp}(\theta_y)}{n}\Bigr)^2\Bigr]^{1/3},
\qquad \frac{2\mathfrak{h}(\theta_y)+2}{n}.
\end{equation*}
Both are independent of the cutoff: $W/n$ does not depend on $n$, and $h_{\sharp}$ is proportional to
$n$, so the product $W\,h_{\sharp}(\theta_y)^2$ under the cube root is proportional to $n^3$. Under
\eqref{8.13:eq-model-conditions-c} both are much smaller than $W/(nL^2)$. Thus the needle's height is
at most of the order of
\begin{equation*}
\max\Bigl\{\bigl[L^2\,(W/n)\,P_{\mathrm{bl}}^2\,e^{-\mathsf{T}/(6s_e)}\bigr]^{1/3},\
P_{\mathrm{bl}}\,e^{-\mathsf{T}/(6s_e)}\Bigr\}\ \text{blocks}.
\end{equation*}

\begin{lemma}[Model: continuity of the minimiser in the cap]\label{8.13:lem-cap-continuity}
This lemma is a statement about the two-tier box-clamp lattice model of
Definition~\ref{8.13:def-two-tier-model}, and about nothing else. In that model, fields are
real-valued functions $\psi\in\mathbb{R}^{\Lambda}$ on the periodic strip
$\Lambda=\mathbb{Z}_N\times\{0,1,\dots,H\}$
and have Dirichlet energy $E(\psi)=\frac12\sum(\nabla\psi)^2$, the sum running over
nearest-neighbour bonds. For $0<\sigma<\tau<c$, the set $\mathcal{K}_c$ consists of the fields
that equal the datum $\varphi_0(x)=\sigma\|x\|$ on the rows $y=0$ and $y=H$, obey the
lower-tier cap $\tau$ and the young-tier cap $c$ on the vertical and horizontal bond
differences, and obey the mirror constraints. The field $\psi^{(c)}$ is the minimiser of $E$
over $\mathcal{K}_c$, unique by the comparison principle for the two-tier model, and
$\tilde{\psi}=\psi^{(\infty)}$ is the minimiser
without the young tier.

Then the map $c\mapsto\psi^{(c)}$ is continuous on $(\tau+2\sigma,\infty)$, with
$\mathbb{R}^{\Lambda}$ carrying any norm. Moreover, $\psi^{(c)}=\tilde{\psi}$ for
$c\ge c_{\max}$, where $c_{\max}$ is the maximum over the bottom band of the vertical bond
differences $\tilde F$ of $\tilde{\psi}$, as in part~(iv) of
Remark~\ref{8.13:rem-budget-monotone}, a remark about the same model.
\end{lemma}

\begin{proof}
Since $\mathbb{R}^{\Lambda}$ is finite-dimensional, all norms on it are equivalent, so it
suffices to prove continuity for one of them.

\emph{Properties of the feasible sets.} Each $\mathcal{K}_c$ is cut out by finitely many
linear equalities (the two Dirichlet rows) and finitely many non-strict inequalities
$|\ell(\psi)|\le b$, where $\ell$ is a linear functional and $b\in\{\tau,c\}$. Hence
$\mathcal{K}_c$ is closed and convex. The only constraints that depend on $c$ are those of
the young tier and their mirror images, all of the form $|\ell(\psi)|\le c$. Therefore
$\mathcal{K}_{c'}\subset\mathcal{K}_c$ whenever $c'<c$.

Put $\Phi_0(x,y):=\varphi_0(x)$. This field lies in every $\mathcal{K}_c$:
\begin{itemize}
\item its vertical differences vanish;
\item its horizontal differences are $\pm\sigma$, because $|\nabla\varphi_0|=\sigma$, and
$\sigma<\tau<c$ places them within every cap;
\item it equals $\varphi_0$ on the Dirichlet rows.
\end{itemize}
Hence $E(\psi^{(c)})\le E(\Phi_0)$ for every $c$. Each bond satisfies
$(\nabla\psi)^2\le 2E(\psi)$, and every site $(x,y)$ is joined to $(x,0)$ by at most $H$
vertical bonds. Therefore
\[
|\psi^{(c)}(x,y)|\le|\varphi_0(x)|+H\sqrt{2E(\Phi_0)} ,
\]
so the fields $\psi^{(c)}$ lie in a bounded set. Consequently every sequence
$\psi^{(c_k)}$ has limit points, and it converges to $\psi^{(c)}$ as soon as every one of
its limit points equals $\psi^{(c)}$.

\emph{Right-continuity.} Let $c\in(\tau+2\sigma,\infty)$ and let $c_k\downarrow c$. Let
$\psi^{*}$ be a limit point of $\psi^{(c_k)}$. For each fixed $j$ and all $k\ge j$, nesting
gives $\psi^{(c_k)}\in\mathcal{K}_{c_k}\subset\mathcal{K}_{c_j}$. Since $\mathcal{K}_{c_j}$ is
closed, $\psi^{*}\in\mathcal{K}_{c_j}$ for every $j$. Thus
$\psi^{*}\in\bigcap_j\mathcal{K}_{c_j}=\mathcal{K}_c$, because a young-tier difference
bounded by every $c_j$ is bounded by $c$. Since $\psi^{(c)}\in\mathcal{K}_c\subset
\mathcal{K}_{c_k}$, we have $E(\psi^{(c_k)})\le E(\psi^{(c)})$ for every $k$. By continuity
of $E$, taken along the subsequence converging to $\psi^{*}$,
\[
E(\psi^{*})=\lim E(\psi^{(c_k)})\le E(\psi^{(c)}).
\]
So $\psi^{*}$ minimises $E$ over $\mathcal{K}_c$. By the comparison principle for the
two-tier model the minimiser of $E$ over $\mathcal{K}_c$ is unique, so $\psi^{*}=\psi^{(c)}$.

\emph{Left-continuity.} Let $c_k\uparrow c$ with $c_k\in(\tau+2\sigma,c)$. Nesting gives
$\psi^{(c_k)}\in\mathcal{K}_c$, and since $\mathcal{K}_c$ is closed, every limit point
$\psi^{*}$ of $\psi^{(c_k)}$ lies in $\mathcal{K}_c$. Put
\begin{equation}\label{8.13:eq-cap-continuity-1}
\psi_k:=\Phi_0+\frac{c_k}{c}\bigl(\psi^{(c)}-\Phi_0\bigr).
\end{equation}
We check that $\psi_k\in\mathcal{K}_{c_k}$.
\begin{itemize}
\item Since $\Phi_0$ has vanishing vertical differences, the vertical differences of
$\psi_k$ are those of $\psi^{(c)}$ multiplied by $c_k/c\le1$. In the young tier they are
therefore at most $(c_k/c)\,c=c_k$ in absolute value, and in the lower tier at most
$(c_k/c)\,\tau\le\tau$.
\item The horizontal differences of $\psi_k$ are convex combinations, with weights
$1-c_k/c$ and $c_k/c$, of the horizontal differences $\pm\sigma$ of $\Phi_0$ and those of
$\psi^{(c)}$. Every cap is a symmetric interval, so these differences lie inside every cap.
\item On the Dirichlet rows both $\Phi_0$ and $\psi^{(c)}$ equal $\varphi_0$, hence so
does $\psi_k$.
\end{itemize}
Hence $E(\psi^{(c_k)})\le E(\psi_k)$. As $c_k/c\to1$ we have $\psi_k\to\psi^{(c)}$, so
$E(\psi_k)\to E(\psi^{(c)})$ by continuity of $E$. Along the subsequence converging to
$\psi^{*}$ this gives $E(\psi^{*})\le E(\psi^{(c)})$. The uniqueness of the minimiser, given
by the comparison principle for the two-tier model, then yields $\psi^{*}=\psi^{(c)}$.

Together the two parts show that $c\mapsto\psi^{(c)}$ is continuous on
$(\tau+2\sigma,\infty)$.

\emph{The cap beyond $c_{\max}$.} Part~(iv) of Remark~\ref{8.13:rem-budget-monotone}, a
statement about the same model, gives $\tilde{\psi}\in\mathcal{K}_c$ for $c\ge c_{\max}$.
The set $\mathcal{K}_c$ is contained in the feasible set without the young tier, and
$\tilde{\psi}$ minimises $E$ over that larger set. Hence $\tilde{\psi}$ minimises $E$ over
$\mathcal{K}_c$, and the uniqueness of the minimiser, given by the comparison principle for
the two-tier model, yields $\psi^{(c)}=\tilde{\psi}$.
\end{proof}

\begin{proposition}[Model: no hanging young pole]\label{8.13:prop-no-hanging-pole}
We work in the strip model: on the periodic strip
$\Lambda=\mathbb{Z}_N\times\{0,\dots,H\}$ the field $\psi^{(c)}$ is the unique minimiser of
the Dirichlet energy $E(\psi)=\frac12\sum(\nabla\psi)^2$ over the closed convex feasible set
$\mathcal{K}_c$ at young cap $c$. Its elements carry the kinked Dirichlet datum
$\varphi_0(x)=\sigma\|x\|$, and every bond difference, horizontal and vertical, lies within
its cap: the straddling cap $\tau$ on the lower bonds, those with endpoints in rows at most $W$
or at least $H-W$, and the young cap $c$ on the bonds of the young band of height $\bar W$
above row $W$ and of its mirror band at the top of the strip. The field $\tilde{\psi}$ is the
straddling-only
minimiser, and $W'=W+\bar W$.

Assume the one-sided conditions
\eqref{8.13:eq-model-conditions-a}--\eqref{8.13:eq-model-conditions-d} of the model instance
(Definition~\ref{8.13:def-model-conditions}, which belongs to the model). Then for every
$c\ge\tau'$ the minimiser $\psi^{(c)}$ has no run of type $(\mathrm{P}^{+})$ or
$(\mathrm{P}^{-})$, in either band and at either tip. Here the run types are the hanging poles
of the model lemma, Lemma~\ref{8.13:lem-tight-census}.
\end{proposition}

\begin{proof}
We write $F^{(c)}_y(x):=\psi^{(c)}(x,y+1)-\psi^{(c)}(x,y)$ and
$\tilde F_y(x):=\tilde{\psi}(x,y+1)-\tilde{\psi}(x,y)$. The argument is written for the bottom
band near the bottom tip. The top tip is treated in the same way with up and down exchanged,
and the top band is the mirror image under $y\mapsto H-y$, as in the model lemma,
Lemma~\ref{8.13:lem-tight-census}. Put
\begin{equation*}
G:=\{c\ge\tau' : \psi^{(c)}\ \text{has no run of type}\ (\mathrm{P}^{+})\ \text{or}\ (\mathrm{P}^{-})\}.
\end{equation*}

\emph{Large caps.} By the model lemma, Lemma~\ref{8.13:lem-cap-continuity}, we have
$\psi^{(c)}=\tilde{\psi}$ for $c\ge c_{\max}$, where $c_{\max}$ is the constant of
\eqref{8.13:eq-budget-monotone-a}.

Let $c\ge\tau'$ with $c\ge c_{\max}$, and suppose that $\tilde{\psi}$ had a run of type
$(\mathrm{P}^{+})$ or $(\mathrm{P}^{-})$ at a column $x$. Parts (a) and (b) of the model lemma,
Lemma~\ref{8.13:lem-hanging-runs}, give $|\tilde F_{W'}(x)|>c\ge\tau'$.

On the other hand $W'=W+\bar W$, so the bound \eqref{8.13:eq-straddling-minimiser-a} of the
model Notation~\ref{8.13:not-straddling-minimiser}, taken with $m=\bar W$, gives
$|\tilde F_{W'}(x)|\le b(\bar W)$. The free strip of $\tilde{\psi}$ has height
$H'\ge 2\bar W+4$, so $H'-1-\bar W\ge\bar W$. The smaller of $\bar W$ and $H'-1-\bar W$ is
therefore $\bar W$. Since $\log^{+}$ is nondecreasing,
\begin{equation*}
b(\bar W)=\sigma\bigl[C_K^{0}+6\log^{+}\bigl(N/(2\bar W+1)\bigr)\bigr]
\le\sigma\bigl[C_K^{0}+6\log^{+}(N/\bar W)\bigr]=B_\beta ,
\end{equation*}
where $B_\beta$ is the quantity \eqref{8.13:eq-no-full-column-a}. By the model condition
\eqref{8.13:eq-model-conditions-b},
\begin{equation*}
B_\beta\le\tau'/2-3\sigma\le\tau'/2 .
\end{equation*}
Hence $|\tilde F_{W'}(x)|\le\tau'/2<\tau'$, which contradicts $|\tilde F_{W'}(x)|>\tau'$.
Therefore every $c\ge\tau'$ with $c\ge c_{\max}$ lies in $G$.

\emph{The bound on $G$.} Let $c\in G$. By the model lemma, Lemma~\ref{8.13:lem-tight-census},
no young horizontal bond is tight. Every young-tight vertical bond is of one of the types
$(\mathrm{N})$, $(\mathrm{F})$, $(\mathrm{P}^{+})$, $(\mathrm{P}^{-})$. The types
$(\mathrm{P}^{\pm})$ are absent because $c\in G$. Type $(\mathrm{F})$ is absent by the model
corollary, Corollary~\ref{8.13:cor-no-full-column}.

By the model corollary, Corollary~\ref{8.13:cor-short-needle}, every tip needle has height
$m^{*}\le m_{\mathrm{ndl}}(c)\le m_{\mathrm{ndl}}(\tau')$, with $m_{\mathrm{ndl}}$ as in
\eqref{8.13:eq-short-needle-a}. Hence the young-tight set of $\psi^{(c)}$ consists of tip
needles only:
\begin{itemize}
\item at the bottom tip, up-needles on $\mathcal{T}$ of height at most
$m_{\mathrm{ndl}}(\tau')$;
\item at the top tip, down-needles;
\item in the top band, the mirror images of these.
\end{itemize}

A site therefore touches an active constraint only if it is an endpoint of one of the
following:
\begin{itemize}
\item a tight lower bond, with endpoints in rows at most $W$ or at least $H-W$;
\item a needle bond, with endpoints in rows at most $W+m_{\mathrm{ndl}}(\tau')$ or at least
$H-W-m_{\mathrm{ndl}}(\tau')$.
\end{itemize}
Set $Y:=W+m_{\mathrm{ndl}}(\tau')+1$. No site of the rows $Y,\dots,H-Y$ touches an active
constraint. By the bond equation for the minimiser, $\psi^{(c)}$ is therefore harmonic for
the four-neighbour Laplacian $\Delta$ at every site of these rows.

We now apply the strip bound above a Lipschitz trace. Its data are the traces
$\psi^{(c)}(\cdot,Y)=\psi^{(c)}(\cdot,H-Y)$; these coincide by the mirror symmetry of the model
under $y\mapsto H-y$. By the comparison principle, their Lipschitz constant is at most
$\sigma$. By uniqueness of the Dirichlet problem on the strip, the restriction of
$\psi^{(c)}$ to the rows $Y,\dots,H-Y$ is the strip harmonic extension
$\mathcal{U}_{N,H-2Y}[g_Y,g_Y]$ of $g_Y=\psi^{(c)}(\cdot,Y)$. The strip bound then gives, with
$m_*:=\min(m,H-2Y-1-m)$,
\begin{equation}\label{8.13:eq-no-hanging-pole-a}
\begin{gathered}
Y:=W+m_{\mathrm{ndl}}(\tau')+1,\\
|F^{(c)}_{Y+m}(x)|\le\sigma\bigl[C_K^{0}+6\log^{+}\bigl(N/(2m_*+1)\bigr)\bigr]
\quad\text{for all } x \text{ and } 0\le m\le H-2Y-1 .
\end{gathered}
\end{equation}

We evaluate \eqref{8.13:eq-no-hanging-pole-a} at the top bond of the band. Put
$m:=\bar W-m_{\mathrm{ndl}}(\tau')-1$, so that $Y+m=W+\bar W=W'$. By the model corollary,
Corollary~\ref{8.13:cor-short-needle}, $m_{\mathrm{ndl}}(\tau')\le\bar W/2-2$, hence
$m\ge\bar W/2+1$.

The model condition \eqref{8.13:eq-model-conditions-d} gives $H\ge 2W'+\bar W+2$. Using this
and then $m_{\mathrm{ndl}}(\tau')\le\bar W/2-2$,
\begin{align*}
H-2Y-1-m&=H-Y-W'-1=H-W'-W-m_{\mathrm{ndl}}(\tau')-2\\
&\ge W'+\bar W-W-m_{\mathrm{ndl}}(\tau')=2\bar W-m_{\mathrm{ndl}}(\tau')\\
&\ge\tfrac32\bar W+2 .
\end{align*}
Thus $0\le m\le H-2Y-1$ and $m_*\ge\bar W/2$, so $2m_*+1\ge\bar W+1$. From
\eqref{8.13:eq-no-hanging-pole-a}, the monotonicity of $\log^{+}$, the definition
\eqref{8.13:eq-no-full-column-a} of $B_\beta$ and \eqref{8.13:eq-model-conditions-b}, for every
$x$ we get
\begin{equation}\label{8.13:eq-no-hanging-pole-1}
|F^{(c)}_{W'}(x)|\le\sigma\bigl[C_K^{0}+6\log^{+}\bigl(N/(\bar W+1)\bigr)\bigr]\le B_\beta
\le\tau'/2<c .
\end{equation}

The mirror bond of the top band is $F^{(c)}_{H-W'-1}$. We have $H-W'-1=Y+m'$ with
$m'=H-2Y-1-m$, and $\min(m',H-2Y-1-m')=m_*$. Hence \eqref{8.13:eq-no-hanging-pole-1} holds for
it as well.

\emph{The supremum argument.} Suppose that $B:=[\tau',\infty)\setminus G$ is non-empty, and
let $c^{*}:=\sup B$. By the large-cap step every element of $B$ is smaller than $c_{\max}$, so
$c^{*}\le c_{\max}<\infty$. Since $B\subset[\tau',\infty)$, also $c^{*}\ge\tau'$.

For fixed $x_0$ the map $c\mapsto F^{(c)}_{W'}(x_0)$ is a difference of two coordinates of
$\psi^{(c)}$. By the model lemma, Lemma~\ref{8.13:lem-cap-continuity}, it is continuous in
$c$.

Suppose first that $c^{*}\in B$. Then $\psi^{(c^{*})}$ has a pole. By the model lemma,
Lemma~\ref{8.13:lem-hanging-runs}, this gives a column $x_0$ with
$|F^{(c^{*})}_{W'}(x_0)|>c^{*}$. Every $c>c^{*}$ lies in $G$, so
\eqref{8.13:eq-no-hanging-pole-1} gives $|F^{(c)}_{W'}(x_0)|\le B_\beta\le\tau'/2$. Letting
$c\downarrow c^{*}$ yields $|F^{(c^{*})}_{W'}(x_0)|\le\tau'/2$, which is smaller than
$\tau'\le c^{*}$, a
contradiction.

Suppose next that $c^{*}\notin B$. Then there are $c_k\in B$ with $c_k\uparrow c^{*}$, and
$\psi^{(c_k)}$ has a pole at a column $x_k$. There are finitely many columns, and the two
bands and tips are handled alike, so we may choose a subsequence along which $x_k=x_0$ and
the band is fixed. By the model lemma, Lemma~\ref{8.13:lem-hanging-runs},
$|F^{(c_k)}_{W'}(x_0)|>c_k$. Letting
$k\to\infty$ yields $|F^{(c^{*})}_{W'}(x_0)|\ge c^{*}\ge\tau'$. But $c^{*}\in[\tau',\infty)$
and $c^{*}\notin B$, so $c^{*}\in G$, and \eqref{8.13:eq-no-hanging-pole-1} gives
$|F^{(c^{*})}_{W'}(x_0)|\le\tau'/2<\tau'$. This is a contradiction.

Hence $B$ is empty and $G=[\tau',\infty)$.
\end{proof}

\begin{remark}[Model: the continuity argument in words]\label{8.13:rem-continuity-argument}
The model is the one of Proposition~\ref{8.13:prop-no-hanging-pole}. For each young cap
$c\ge\tau'$, $\psi^{(c)}$ is the minimiser of the Dirichlet energy $E(\psi)$ over the feasible
set with young cap $c$. The model theorem, Proposition~\ref{8.13:prop-no-hanging-pole}, states
that no $\psi^{(c)}$ has a run of type $(\mathrm{P}^{+})$ or $(\mathrm{P}^{-})$, a hanging pole.

Its proof is a continuity argument in the cap. One lets $c$ decrease from $+\infty$ down to
$\tau'$. For large caps there is no pole, and the first cap at which the young constraint binds
at the top of the band must be one at which the bound on $F_{W'}$ is violated. Three
ingredients exclude this.

First, the lattice classification of young-tight runs says what could bind at the top of the
band. Only poles can bind there, and a pole needs $|F_{W'}|>c$.

Second, the energy budget for needles and full columns keeps the rest of the active set inside
the first $m_{\mathrm{ndl}}(\tau')$ rows of each band, uniformly in $c$.

Third, above an exit trace with Lipschitz constant at most $\sigma$, the vertical bond
difference at height at least $\bar W/2$ obeys the logarithmic strip bound, whose value is at
most $B_\beta$.

Hence along the whole family of caps the top of the band is never pulled hard enough for a
pole to appear. The underlying principle is that the young constraint cannot act at depth
$2W$ except through a trace whose Lipschitz constant it cannot raise. For depths inside the
young band this principle is supplemented by the needle bound $m_{\mathrm{ndl}}(\tau')$.
\end{remark}

\begin{theorem}[Model: inactivity at full band width]\label{8.13:thm-inactivity}
The model is the young-row strip model of this section: on the periodic strip
$\Lambda=\mathbb{Z}_N\times\{0,\dots,H\}$ with the Dirichlet energy $E(\psi)$, and for a young cap $c$,
$\psi^{(c)}$ is the minimiser of $E$ over the feasible set $\mathcal{K}_c$, in which every young bond is
capped by $c$. We work at the young cap $c=\tau'$. Assume the one-sided conditions
\eqref{8.13:eq-model-conditions-a}--\eqref{8.13:eq-model-conditions-d} of
Definition~\ref{8.13:def-model-conditions}, and assume the mild-fringe hypothesis of
Definition~\ref{8.13:def-mild-fringe}. Let $\hat{\psi}:=\psi^{(\tau')}$, let
$m_{\mathrm{ndl}}:=m_{\mathrm{ndl}}(\tau')$ with $m_{\mathrm{ndl}}(\cdot)$ as in
\eqref{8.13:eq-short-needle-a}, and put $Y:=W+m_{\mathrm{ndl}}+1$. Then the following hold.
\begin{enumerate}
\item[(i)] (The active set.) The young constraints active at $\hat{\psi}$ are exactly the following: no
horizontal young bond; and vertical young bonds only on the tip columns $\mathcal{T}$ (in the lattice
variant of the model whose tip set is the single column $x\equiv0$, only on that column) within the
first $m_{\mathrm{ndl}}$ rows above the
clamp's inner face, forming one up-run $[W,W+m^{*})$ per tip column with $m^{*}\le m_{\mathrm{ndl}}$;
together with their symmetric images (the down-needles at the top tip, and the mirror images in the
top band). In particular every young bond at height at least $m_{\mathrm{ndl}}$ above the inner face,
and every young bond off the tip columns, carries an inactive young constraint.
\item[(ii)] (Harmonicity.) $\hat{\psi}$ is $\Delta$-harmonic at every site of the rows $Y,\dots,H-Y$.
With $u(x,m):=\hat{\psi}(x,Y+m)$ and $g_Y:=\hat{\psi}(\cdot,Y)$ one has
$u=\mathcal{U}_{N,H-2Y}[g_Y,g_Y]$ and $\operatorname{Lip} g_Y\le\sigma$.
\item[(iii)] (Room.) For all $x$ and all $m$ with $0\le m\le H-2Y-1$, writing
$m_*:=\min\{m,\,H-2Y-1-m\}$,
\begin{equation}\label{8.13:eq-inactivity-1}
\bigl|\hat{\psi}(x,Y+m+1)-\hat{\psi}(x,Y+m)\bigr|
\le\sigma\Bigl[C_K^{0}+6\log^{+}\Bigl(\frac{N}{2m_*+1}\Bigr)\Bigr],
\end{equation}
and $|\hat{\psi}(x+1,y)-\hat{\psi}(x,y)|\le\sigma$. Hence, for every $\rho\in(0,\tfrac12)$, every young
bond at height $h\ge h_\rho$ above the inner face obeys $|\hat F|\le(1-\rho)\tau'$, where
\begin{equation}\label{8.13:eq-inactivity-a}
h_\rho:=m_{\mathrm{ndl}}+1+\mathfrak{h}_\rho,\qquad
\mathfrak{h}_\rho:=\max\Bigl\{0,\ \tfrac12\Bigl[N\,
e^{-((1-\rho)\tau'/\sigma-C_K^{0})/6}-1\Bigr]\Bigr\}.
\end{equation}
Part~(d) of \eqref{8.13:eq-model-conditions-c} states $h_\rho\le\bar W-1$.
\end{enumerate}
\end{theorem}

\begin{proof}
Here a young constraint is called active at $\hat{\psi}$ when its bond is tight, that is, when the
difference of $\hat{\psi}$ across it equals $\tau'$ in absolute value.

\emph{Clause (i).} By the model proposition, Proposition~\ref{8.13:prop-no-hanging-pole}, applied at
the cap $c=\tau'$, the field $\hat{\psi}$ has no run of type $(\mathrm{P}^{+})$ or $(\mathrm{P}^{-})$,
in either band and at either tip. By the model corollary, Corollary~\ref{8.13:cor-no-full-column}, at
the cap $c=\tau'$, the field $\hat{\psi}$ has no young-full column, that is, no run of type
$(\mathrm{F})$. By the model lemma, Lemma~\ref{8.13:lem-tight-census}, every young-tight vertical bond
of the bottom band near the bottom tip belongs to exactly one of the types $(\mathrm{N})$,
$(\mathrm{F})$, $(\mathrm{P}^{+})$, $(\mathrm{P}^{-})$. With $(\mathrm{F})$, $(\mathrm{P}^{+})$ and
$(\mathrm{P}^{-})$ excluded, the remaining young-tight vertical bonds are tip needles, up-runs
$[W,W+m^{*})$ on columns of $\mathcal{T}$; the same holds for their symmetric images, the down-needles
at the top tip and the mirror images in the top band. By the model corollary,
Corollary~\ref{8.13:cor-short-needle}, at the cap $c=\tau'$ every such needle has
$m^{*}\le m_{\mathrm{ndl}}(\tau')=m_{\mathrm{ndl}}$. For the horizontal young bonds, the comparison
principle of the model gives $|\hat{\psi}(x+1,y)-\hat{\psi}(x,y)|\le\sigma<\tau'$, so no horizontal
young bond is tight. This proves the description of the active set. Since a needle $[W,W+m^{*})$
occupies only the heights $0,\dots,m^{*}-1\le m_{\mathrm{ndl}}-1$ above the inner face, every young bond
at height at least $m_{\mathrm{ndl}}$, and every young bond off the tip columns, carries an
inactive young constraint.

\emph{Clauses (ii) and (iii).} These follow as in the derivation of the strip bound
\eqref{8.13:eq-no-hanging-pole-a} of the model proposition, Proposition~\ref{8.13:prop-no-hanging-pole},
carried out at $c=\tau'$. By clause~(i) the young-tight
set of $\hat{\psi}$ consists of tip needles only. A site is an endpoint of a bond with active
constraint only if it is an endpoint of a tight lower bond, which lies in the rows $\le W$ or
$\ge H-W$, or of a needle bond, which lies in the rows $\le W+m_{\mathrm{ndl}}$ or
$\ge H-W-m_{\mathrm{ndl}}$. Hence no site of the rows $Y,\dots,H-Y$ is such an endpoint, and the bond
equation of the model makes $\hat{\psi}$ $\Delta$-harmonic at each of them. The boundary data
$\hat{\psi}(\cdot,Y)=\hat{\psi}(\cdot,H-Y)=g_Y$ satisfy $\operatorname{Lip}g_Y\le\sigma$ by the
comparison principle; by uniqueness of the Dirichlet problem on the strip,
$u=\mathcal{U}_{N,H-2Y}[g_Y,g_Y]$. This is clause~(ii). The strip bound above a Lipschitz trace,
applied to $u$, gives \eqref{8.13:eq-inactivity-1} for all $x$ and $0\le m\le H-2Y-1$, and the
horizontal bound is the comparison principle once more.

For the last assertion of clause~(iii), let a young bond lie at height $h\ge h_\rho$ above the inner
face, so that its lower row is $W+h=Y+m$ with $m:=h-m_{\mathrm{ndl}}-1\ge\mathfrak{h}_\rho$. Then
\begin{equation}\label{8.13:eq-inactivity-2}
2m+1\ \ge\ 2\mathfrak{h}_\rho+1\ \ge\ N\,e^{-((1-\rho)\tau'/\sigma-C_K^{0})/6}.
\end{equation}
For $h\le\bar W$, the condition \eqref{8.13:eq-model-conditions-d} gives $m_*=m$. Put
$a:=(1-\rho)\tau'/\sigma-C_K^{0}$. By \eqref{8.13:eq-model-conditions-b} and $\log^{+}\ge0$ we have
\begin{equation*}
\sigma C_K^{0}\le B_\beta\le\tau'/2-3\sigma<(1-\rho)\tau',
\end{equation*}
since $\rho<\tfrac12$; hence $a>0$.
By \eqref{8.13:eq-inactivity-2}, $N/(2m_*+1)\le e^{a/6}$, so $\log^{+}(N/(2m_*+1))\le a/6$, and
\eqref{8.13:eq-inactivity-1} gives
\begin{equation*}
|\hat F|\ \le\ \sigma\bigl[C_K^{0}+a\bigr]\ =\ (1-\rho)\tau'.
\end{equation*}
\end{proof}

\begin{corollary}[Model: logarithmic bound at depth two clamp-widths]\label{8.13:cor-depth-log-bound}
The model is the one of Theorem~\ref{8.13:thm-inactivity}, the model theorem on inactivity at full
band width: the field $\hat{\psi}=\psi^{(\tau')}$ at young cap $\tau'$ on the periodic strip
$\mathbb{Z}_N\times\{0,\dots,H\}$, under the hypotheses of that theorem and the one-sided
conditions of the model instance of Definition~\ref{8.13:def-model-conditions}. Put
$m_{\mathrm{ndl}}:=m_{\mathrm{ndl}}(\tau')$ and $Y:=W+m_{\mathrm{ndl}}+1$. For every row $y$ with
$Y\le y\le H-Y-1$ and every $x$,
\begin{equation}\label{8.13:eq-depth-log-bound-1}
  \bigl|\hat{\psi}(x,y+1)-\hat{\psi}(x,y)\bigr|
  \le \sigma\Bigl[C_K^{0}+6\log^{+}\Bigl(\frac{N}{2(y-Y)_*+1}\Bigr)\Bigr],
\end{equation}
where $(y-Y)_*:=\min(y-Y,\,H-Y-1-y)$. Read through the dictionary between the model and the
block variables (Remark~\ref{8.13:rem-model-dictionary}), the vertical bond difference at row $y$
is the normal--tangential plaquette component $\nabla_4\hat{\psi}$ at depth $w+(y-W)/n$, and the
horizontal bond difference is the tangential--tangential component $\nabla_1\hat{\psi}$, both in
the field unit $\varepsilon_{k_o}\eta^2$. In particular, at depth $2w$, that is, at the row
$y=2W$,
\begin{equation}\label{8.13:eq-depth-log-bound-2}
\begin{aligned}
  |\nabla_4\hat{\psi}|
  &\le \sigma\Bigl[C_K^{0}+6\log^{+}\Bigl(\frac{N}{W+1}\Bigr)\Bigr]
   \le s_e\Bigl[C_K^{0}+6\log\frac{2P_{\mathrm{bl}}}{w}\Bigr]\,\theta'\varepsilon_{k_o}\eta^2,\\
  |\nabla_1\hat{\psi}| &\le s_e\,\theta'\varepsilon_{k_o}\eta^2 .
\end{aligned}
\end{equation}
At every depth $w+(y-W)/n\ge w+(m_{\mathrm{ndl}}+1)/n$ with $y\le H-Y-1$ the same two bounds
hold with $\log(2P_{\mathrm{bl}}/w)$ replaced by $\log^{+}\bigl(N/(2(y-Y)_*+1)\bigr)$.
\end{corollary}

The constants $C_K^{0}$ and $6$ in \eqref{8.13:eq-depth-log-bound-2} are absolute, the coefficient
of the logarithm does not involve $\theta'$, and the dictionary of
Remark~\ref{8.13:rem-model-dictionary}, which belongs to the model, bounds $P_{\mathrm{bl}}/w$
independently of the cutoff. The logarithmic bound at the canonical comparison point $\hat{x}$
obtained beyond the young band requires the young band to have height at most $W/2$, so that
depth $2w$ lies beyond the band; \eqref{8.13:eq-depth-log-bound-2} holds at depth $2w$ also when
this depth lies inside the young band, as it does when $\varpi_y=w+w/L$.

\begin{proof}
Let $Y\le y\le H-Y-1$ and put $m:=y-Y$, so that $0\le m\le H-2Y-1$. By part~(ii) of
Theorem~\ref{8.13:thm-inactivity}, the model theorem on inactivity at full band width,
$\hat{\psi}$ is harmonic at every site of the rows $Y,\dots,H-Y$. By part~(iii) of the same
model theorem, applied with this $m$, for all $x$
\begin{equation*}
  \bigl|\hat{\psi}(x,Y+m+1)-\hat{\psi}(x,Y+m)\bigr|
  \le \sigma\Bigl[C_K^{0}+6\log^{+}\Bigl(\frac{N}{2m_*+1}\Bigr)\Bigr],
\end{equation*}
and $|\hat{\psi}(x+1,y)-\hat{\psi}(x,y)|\le\sigma$. Here $m_*$ is the quantity of part~(iii),
which at $m=y-Y$ is $(y-Y)_*=\min(m,\,H-2Y-1-m)$. Since $Y+m=y$, the first inequality is
\eqref{8.13:eq-depth-log-bound-1}.

By the dictionary of Remark~\ref{8.13:rem-model-dictionary}, which belongs to the model,
$\sigma=s_e\theta'$, field values are measured in the unit $\varepsilon_{k_o}\eta^2$, and height
$y-W$ above the clamp's inner face corresponds to depth $w+(y-W)/n$. Hence
\eqref{8.13:eq-depth-log-bound-1} gives, at every row $Y\le y\le H-Y-1$, that is, at every depth
at least $w+(m_{\mathrm{ndl}}+1)/n$ in this range of rows,
\begin{equation*}
  |\nabla_4\hat{\psi}|
  \le s_e\Bigl[C_K^{0}+6\log^{+}\Bigl(\frac{N}{2(y-Y)_*+1}\Bigr)\Bigr]\theta'\varepsilon_{k_o}\eta^2,
\end{equation*}
and the horizontal bound of part~(iii) gives $|\nabla_1\hat{\psi}|\le\sigma
=s_e\theta'\varepsilon_{k_o}\eta^2$. This is the last assertion of the corollary.

Now take $y=2W$. Then $y-Y=2W-(W+m_{\mathrm{ndl}}+1)=W-m_{\mathrm{ndl}}-1$, and by condition~(c)
of \eqref{8.13:eq-model-conditions-c}, one of the conditions of the model instance,
\begin{equation*}
  y-Y=W-m_{\mathrm{ndl}}-1\ge \tfrac{W}{2}+1 .
\end{equation*}
In particular $y\ge Y$. By \eqref{8.13:eq-model-conditions-d}, another of the conditions of the
model instance, $y-Y\le H-Y-1-y$, so $(y-Y)_*=y-Y$ and $y\le H-Y-1$. Thus $y=2W$ lies in the
range of \eqref{8.13:eq-depth-log-bound-1}, and
\begin{equation*}
  2(y-Y)_*+1=2(y-Y)+1\ge W+3\ge W+1 .
\end{equation*}
Since $\log^{+}$ is non-decreasing, $\log^{+}\bigl(N/(2(y-Y)_*+1)\bigr)\le\log^{+}\bigl(N/(W+1)\bigr)$,
and \eqref{8.13:eq-depth-log-bound-1} yields the first inequality of
\eqref{8.13:eq-depth-log-bound-2}. For the second, the dictionary gives $N=2P_{\mathrm{bl}}n$ and
$W=wn$, hence
\begin{equation*}
  \frac{N}{W+1}<\frac{N}{W}=\frac{2P_{\mathrm{bl}}}{w},
\end{equation*}
so that $\log^{+}\bigl(N/(W+1)\bigr)\le\log(2P_{\mathrm{bl}}/w)$; together with
$\sigma=s_e\theta'$ and the field unit $\varepsilon_{k_o}\eta^2$ this is the second inequality.
The bound on $|\nabla_1\hat{\psi}|$ at $y=2W$ is the horizontal bound of part~(iii) read through
the dictionary, as above.
\end{proof}

\begin{corollary}[Model: inactivity read in the block variables]\label{8.13:cor-inactivity-blocks}
The model is that of the model theorem, Theorem~\ref{8.13:thm-inactivity}: the minimisation of
the Dirichlet energy $E(\psi)$ on the periodic strip
$\Lambda=\mathbb{Z}_N\times\{0,\dots,H\}$, with the kinked Dirichlet datum $\varphi_0(x)=\sigma\|x\|$,
over the feasible set $\mathcal{K}_{\tau'}$ at young cap $\tau'$, with minimiser
$\hat{\psi}=\psi^{(\tau')}$. It is read, through the model
dictionary of Remark~\ref{8.13:rem-model-dictionary}, as the $x_2x_3$-invariant abelian
four-dimensional box-clamp problem at the canonical comparison point $\hat{x}$, for the exterior
of relative kink slope $s_e<1$. Here $w$ is the clamp width in top blocks ($W=wn$,
$n=L^{k_o}$, $\eta=1/n$), $\delta$ denotes the depth below the top face, in top blocks, so that
height $m$ above the clamp's inner face $\delta=w$ is depth $\delta=w+m/n$ (by the model
dictionary, Remark~\ref{8.13:rem-model-dictionary}), and the young band above the inner face has
width $\varpi_y=w+w/L$.

Assume the mild-fringe hypothesis, with lateral fraction $\kappa$, of the model theorem,
Theorem~\ref{8.13:thm-inactivity}. Assume also the one-sided conditions of the model
instance, Definition~\ref{8.13:def-model-conditions}. Then the following hold.
\begin{itemize}
\item[(a)] At $\hat{x}$ the young windows of the first held young row $\mathfrak{Y}_1$ bind only
on plaquettes of the tip columns of the kink lying at depths $w\le\delta<w+m_{\mathrm{ndl}}/n$,
where $m_{\mathrm{ndl}}=m_{\mathrm{ndl}}(\tau')$ and
\begin{equation}\label{8.13:eq-inactivity-blocks-1}
\begin{split}
\frac{m_{\mathrm{ndl}}}{n}\le\max\Bigl\{&2P_{\mathrm{bl}}\,\exp\bigl(1+C_K^{0}/6\bigr)\,
\exp\bigl(-\mathsf{T}/(6s_e)\bigr)+\frac{2}{n},\\
&\Bigl[128\cdot144\cdot\Bigl(\frac{\mathsf{T}}{s_e}+\frac{24}{\kappa s_e}+19\Bigr)
\,w\,P_{\mathrm{bl}}^2\\
&\qquad\cdot \exp\bigl(C_T/6\bigr)\,\exp\bigl(-(\mathsf{T}-24/\kappa)/(6s_e)\bigr)\Bigr]^{1/3}\Bigr\}.
\end{split}
\end{equation}
Moreover, for every $\rho\in(0,\tfrac12)$, on the depths
$\{w+h_\rho/n\le\delta\le 2w+w/L\}$ the normal--tangential components of the configuration
$U(e\oplus\hat{x})$ formed from the kinked exterior $e$ and the comparison point $\hat{x}$
are bounded in absolute value by $(1-\rho)\theta_y\varepsilon_{k_o}\eta^2$. On the same depths the
tangential--tangential components are bounded in absolute value by
$s_e\theta'\varepsilon_{k_o}\eta^2$. Here $h_\rho$ is the height of
\eqref{8.13:eq-inactivity-a}, and
\begin{equation}\label{8.13:eq-inactivity-blocks-2}
\frac{h_\rho}{n}\le\frac{m_{\mathrm{ndl}}}{n}+\frac1n
+P_{\mathrm{bl}}\,\exp\bigl(C_K^{0}/6\bigr)\,\exp\bigl(-(1-\rho)\mathsf{T}/(6s_e)\bigr).
\end{equation}
\item[(b)] Suppose instead that the young band is the thin band $(w,w+w/L]$, so that
$\varpi_y=w/L$. Then the model Theorem~\ref{8.13:thm-inactivity} holds verbatim with
$\bar W=wn/L$, the conditions \eqref{8.13:eq-model-conditions-b} and
\eqref{8.13:eq-model-conditions-c} being read with $\varpi_y=w/L$. In the variables of the
dictionary, the condition \eqref{8.13:eq-model-conditions-b} then reads
\begin{equation}\label{8.13:eq-inactivity-blocks-3}
s_e\Bigl[C_K^{0}+3+6\log^{+}\Bigl(\frac{2P_{\mathrm{bl}}L}{w}\Bigr)\Bigr]\le\frac{\mathsf{T}}{2}.
\end{equation}
The conditions \eqref{8.13:eq-model-conditions-b} and \eqref{8.13:eq-model-conditions-c}, read
with $\varpi_y=w/L$, remain one-sided lower bounds on $L$ with coefficients free of the cutoff.
\end{itemize}
\end{corollary}

\begin{proof}
\emph{The dictionary.} We use from the model dictionary, Remark~\ref{8.13:rem-model-dictionary},
which is a model statement, the following identifications: $N=2P_{\mathrm{bl}}n$, $W=wn$ and
$\bar W=\varpi_y n$, and $W'=y_+n$ with $y_+=w+\varpi_y=2w+w/L$; $\tau=\theta'$,
$\sigma=s_e\theta'$ and $\tau'=\theta_y=\theta'\mathsf{T}$, so that $\tau/\sigma=1/s_e$ and
$\tau'/\sigma=\mathsf{T}/s_e$; height $m$ above the clamp's inner face corresponds to depth
$\delta=w+m/n$; a vertical bond corresponds to a normal--tangential plaquette, and a horizontal
bond to a tangential--tangential plaquette; differences of $\psi$ are measured in units of
$\varepsilon_{k_o}\eta^2$. Moreover, a window of $\mathfrak{Y}_1$ binds if and only if some bond
of its test region, which consists of the rows $W-wn/L\le y\le W'$, is young-tight; the rows
$W-wn/L\le y\le W-1$ carry lower-tier bonds, with cap $\tau<\tau'$, and these bonds are never
young-tight.

\emph{Where the windows bind.} By the model theorem, Theorem~\ref{8.13:thm-inactivity}(i), no
horizontal young bond is active at $\hat{\psi}$. Every young bond at height at least
$m_{\mathrm{ndl}}$ above the inner face is inactive at $\hat{\psi}$, that is, its young constraint
holds there with strict inequality; the same holds for every young bond off the tip
columns $\mathcal{T}$. Hence the young-tight bonds of a test region are vertical bonds on the tip
columns in the rows $W\le y<W+m_{\mathrm{ndl}}$. By the dictionary these are normal--tangential
plaquettes of the tip columns at depths $w\le\delta<w+m_{\mathrm{ndl}}/n$.

\emph{The bound \eqref{8.13:eq-inactivity-blocks-1}.} By the model corollary,
Corollary~\ref{8.13:cor-short-needle}, and \eqref{8.13:eq-short-needle-a},
\[
m_{\mathrm{ndl}}=\max\{2\mathfrak{h}(\tau')+2,\ m_e(\tau')\},
\]
where
\[
m_e(\tau')=\Bigl[128\,C_d\,\frac{\tau'+\beta_\kappa+19\sigma}{\sigma}\,W\,h_{\sharp}(\tau')^2\Bigr]^{1/3}.
\]
Since $C_d=144$, $\beta_\kappa=24\tau/\kappa$ and $\tau/\sigma=1/s_e$, we have
\[
\frac{\tau'+\beta_\kappa+19\sigma}{\sigma}=\frac{\mathsf{T}}{s_e}+\frac{24}{\kappa s_e}+19 .
\]
Using $W=wn$, this gives
\[
\frac{m_e(\tau')}{n}=\Bigl[128\cdot144\cdot\Bigl(\frac{\mathsf{T}}{s_e}+\frac{24}{\kappa s_e}+19\Bigr)
\,w\,\Bigl(\frac{h_{\sharp}(\tau')}{n}\Bigr)^2\Bigr]^{1/3}.
\]
Next we insert $N=2P_{\mathrm{bl}}n$, $\tau'/\sigma=\mathsf{T}/s_e$ and
$\beta_\kappa/\sigma=24/(\kappa s_e)$ into the definitions of $h_{\sharp}(c)$ and
$\mathfrak{h}(c)$ in the model display \eqref{8.13:eq-hot-box-a}, at $c=\tau'$. This gives
\[
\Bigl(\frac{h_{\sharp}(\tau')}{n}\Bigr)^2\le P_{\mathrm{bl}}^2\,\exp\bigl(C_T/6\bigr)\,
\exp\bigl(-(\mathsf{T}-24/\kappa)/(6s_e)\bigr)
\]
and
\[
\frac{2\mathfrak{h}(\tau')}{n}\le 2P_{\mathrm{bl}}\,\exp\bigl(1+C_K^{0}/6\bigr)\,
\exp\bigl(-\mathsf{T}/(6s_e)\bigr).
\]
Dividing $m_{\mathrm{ndl}}$ by $n$ and combining the last three displays gives
\eqref{8.13:eq-inactivity-blocks-1}.

\emph{The bounds on the depths above $h_\rho$.} Fix $\rho\in(0,\tfrac12)$. By the model theorem,
Theorem~\ref{8.13:thm-inactivity}(iii), every young bond at height $h\ge h_\rho$ above the inner
face obeys $|\hat F|\le(1-\rho)\tau'=(1-\rho)\theta_y$. The horizontal differences of
$\hat{\psi}$ are bounded by $\sigma=s_e\theta'$. Here
\[
h_\rho=m_{\mathrm{ndl}}+1+\mathfrak{h}_\rho .
\]
The young bonds of these heights fill the depths
\[
w+h_\rho/n\le\delta\le y_+=2w+w/L .
\]
This range
is non-empty, because \eqref{8.13:eq-model-conditions-c} displays $h_\rho\le\bar W-1$. By the
dictionary, the vertical bonds there give the normal--tangential bound
$(1-\rho)\theta_y\varepsilon_{k_o}\eta^2$. The horizontal bonds give the tangential--tangential
bound $s_e\theta'\varepsilon_{k_o}\eta^2$. Next,
\[
\mathfrak{h}_\rho=\max\Bigl\{0,\tfrac12\bigl[N\exp\bigl(-((1-\rho)\tau'/\sigma-C_K^{0})/6\bigr)-1\bigr]\Bigr\}
\le\tfrac12\,N\exp\bigl(-((1-\rho)\tau'/\sigma-C_K^{0})/6\bigr).
\]
Inserting $N=2P_{\mathrm{bl}}n$ and $\tau'/\sigma=\mathsf{T}/s_e$ gives
\[
\mathfrak{h}_\rho\le P_{\mathrm{bl}}\,n\,\exp\bigl(C_K^{0}/6\bigr)\,
\exp\bigl(-(1-\rho)\mathsf{T}/(6s_e)\bigr).
\]
Dividing $h_\rho=m_{\mathrm{ndl}}+1+\mathfrak{h}_\rho$ by $n$ gives
\eqref{8.13:eq-inactivity-blocks-2}. This proves (a).

\emph{The thin band.} With $\varpi_y=w/L$ the dictionary gives $\bar W=\varpi_y n=wn/L$. The
model theorem is then applied verbatim with this $\bar W$, the conditions
\eqref{8.13:eq-model-conditions-b} and \eqref{8.13:eq-model-conditions-c} being read with
$\varpi_y=w/L$. In the variables of the dictionary, \eqref{8.13:eq-model-conditions-b} reads
\[
s_e\bigl[C_K^{0}+3+6\log^{+}(2P_{\mathrm{bl}}/\varpi_y)\bigr]\le\mathsf{T}/2 .
\]
Since $2P_{\mathrm{bl}}/\varpi_y=2P_{\mathrm{bl}}L/w$, this is
\eqref{8.13:eq-inactivity-blocks-3}. Since $P_{\mathrm{bl}}/w$ is free of the cutoff by the
dictionary, so is $\log^{+}(2P_{\mathrm{bl}}L/w)$; since moreover $\mathsf{T}$ is comparable to
$L^2$ by the dictionary, the conditions
\eqref{8.13:eq-model-conditions-b} and \eqref{8.13:eq-model-conditions-c} read with
$\varpi_y=w/L$ remain one-sided lower bounds on $L$ with coefficients free of the cutoff. This
proves (b).
\end{proof}

We compare clause (a) of Corollary~\ref{8.13:cor-inactivity-blocks} with the inactivity statement
for the first held young row of full width at $\hat{x}$ in Ba{\l}aban's setting, whose content we
state in full. Let $r_1$ be the width of a collar of Ba{\l}aban's construction, and let $C$ be a
constant. The statement is that at $\hat{x}$ the young windows of $\mathfrak{Y}_1$ bind only on
plaquettes within $r_1+Cw/L^2$ behind the inner face $\delta=w$; equivalently,
that $U(e\oplus\hat{x})$ is at most $(1-\rho)\theta_y\varepsilon_{k_o}\eta^2$ on the depths
\[
\{w+r_1+Cw/L^2\le\delta\le 2w+w/L\}.
\]
In the model this statement holds, with a sliver narrower than the one it allows. Apart from the
lattice terms $2/n$ and $1/n$, the bounds \eqref{8.13:eq-inactivity-blocks-1} and
\eqref{8.13:eq-inactivity-blocks-2} are exponentially small in $\mathsf{T}/s_e$; in the second
term of \eqref{8.13:eq-inactivity-blocks-1} this comes with a cube root of $w$. Here
$\mathsf{T}/s_e\ge\mathsf{T}$, and $\mathsf{T}$ is comparable to $L^2$ by the dictionary. In
particular both right-hand sides are small compared with $w/L^2$, the width of the sliver allowed
there. The sliver in Ba{\l}aban's setting has width $r_1+Cw/L^2$, and the collar width $r_1$ does
not occur in the model.

The test regions of the next row of cubes, on which one tests that the young binding does not
pass to that row, cover the depth band $[2w-w/L,\,3w+w/L]$. In the model this band lies inside
the range of rows of the model theorem, Theorem~\ref{8.13:thm-inactivity}(iii).

\begin{remark}[Model: the mollified kink]\label{8.13:rem-mollified-kink}
The model is the strip model of this section: on
$\Lambda=\mathbb{Z}_N\times\{0,1,\dots,H\}$, the minimiser $\psi^{(c)}$ of the Dirichlet
energy $E$ over the feasible set $\mathcal{K}_c$ with its lower and young clamped tiers. In
this model the kinked datum $\varphi_0=\sigma T$, $T(x)=\|x\|$, is replaced by the mollified
datum $\varphi_0=\sigma T^{\mathrm{mol}}$, where $T^{\mathrm{mol}}$ is the mollified profile;
in the statements of this section $\varphi_0$ then denotes the mollified datum. Its kink is
spread over the tip set
\begin{equation*}
\mathcal{T}=\{x:\ \|x\|\le t_0,\ \Delta_x\varphi_0(x)>0\},
\end{equation*}
which consists of $\ell_m n$ columns, where $\ell_m$ is the width of the mollification
measured in blocks of $n$ columns; here $t_0$ is the half-width of the mollification.
In Ba{\l}aban's construction the kink that this datum represents is resolved over
$\ell_0=O(1+\theta'/s_e)$ blocks; the tip set $\mathcal{T}$ reproduces this resolution. With
the mollified datum the following substitutions are made.
\begin{itemize}
\item In the model lemma, Lemma~\ref{8.13:lem-anchored-runs}, conclusion (d) reads
$x\in\mathcal{T}$.
\item Needles may stand on each tip column. The model corollary,
Corollary~\ref{8.13:cor-short-needle}, applies to each tip column separately, with the same
bound $m_{\mathrm{ndl}}$ as in \eqref{8.13:eq-short-needle-a}.
\item In step (b) of the proof of the model proposition,
Proposition~\ref{8.13:prop-no-hanging-pole}, and in the model theorem,
Theorem~\ref{8.13:thm-inactivity}, the set on which the field is not harmonic is
$\mathcal{T}\times[W,W+m_{\mathrm{ndl}}]$. The row $Y$ of
\eqref{8.13:eq-no-hanging-pole-a} is unchanged.
\item The bounds of the model lemmas, Lemma~\ref{8.13:lem-kink-pull-sign} and
Lemma~\ref{8.13:lem-kink-pull-decay}, acquire the additive correction
\begin{equation}\label{8.13:eq-mollified-kink-1}
2\sigma|\mathcal{T}|\cdot\min\Bigl\{\frac{\sqrt{3}}{2m+1},\
\frac{3}{\|x\|-t_0}\Bigr\}.
\end{equation}
This correction comes from the convolution
$(\nabla T^{\mathrm{mol}}-\nabla T)\ast K^{R}_m$ and is estimated by items (a) and (b) of the
lemma on kernel bounds for the real odd strip kernel $K^{R}_m$. For
$\max(\|x\|,m)\ge|\mathcal{T}|$ the
correction \eqref{8.13:eq-mollified-kink-1} is absorbed into the constant $C_T$ of
\eqref{8.13:eq-kink-pull-decay-a}. Otherwise it is absorbed into the hot box of the model
corollary, Corollary~\ref{8.13:cor-hot-box}, by replacing $h_{\sharp}$ of
\eqref{8.13:eq-hot-box-a} with $\max(h_{\sharp},2t_0)$.
\end{itemize}
With these substitutions every model statement of this section holds for the mollified
datum in the same form as for the kinked datum.
\end{remark}

\begin{remark}[From the model to the canonical comparison point]\label{8.13:rem-model-transfer}
The model of Theorem~\ref{8.13:thm-inactivity} is the $x_2x_3$-invariant abelian four-dimensional
box-clamp problem on the strip $\Lambda$ of Theorem~\ref{8.13:thm-inactivity}. In it a vertical bond
stands for a normal--tangential plaquette and a horizontal bond for a tangential--tangential
plaquette. Corollary~\ref{8.13:cor-inactivity-blocks}, also a model statement, reads the model
theorem in the block variables.

Ba{\l}aban's canonical comparison point $\hat{x}$, a configuration of block variables of the
four-dimensional non-abelian theory, is formed by the straddling family $\mathfrak{S}$ together
with the first young row $\mathfrak{Y}_1$, whose windows are held at the cap $\theta$ (the
straddling cap being $\theta'=\theta(1+\vartheta_1)$).
At $\hat{x}$ the conclusion of Corollary~\ref{8.13:cor-inactivity-blocks} is not proved in this
book; it is proved as an implication from the following two statements.
\begin{itemize}
\item[(a)] The mild-fringe hypothesis of the model theorem, Theorem~\ref{8.13:thm-inactivity}, read
for Ba{\l}aban's $\tilde{x}$, the comparison point without the young row.
\item[(b)] The statements that carry a conclusion about the model's minimisers over to
Ba{\l}aban's block variables at $\tilde{x}$ and $\hat{x}$. These are the following.
\begin{itemize}
\item[(1)] The passage from the linearised to the non-linear problem at the cap $\theta$, with the
contraction of \cite[Proposition~1]{B-CMP122a} as template.
\item[(2)] The localisation of the problem to the box of the zone, in which the remaining faces
enter only as spectators.
\item[(3)] The block-handle structure. The windows constrain local minimisers of the block data.
The fine field is determined by the block variables. The clamp is the straddling cap
$\theta'$, up to a relative correction of order $\vartheta_1$.
\item[(4)] For the tangential--tangential components of the conclusion, the corresponding
tangential--tangential statement for the block means at $\hat{x}$.
\end{itemize}
\end{itemize}
Items (1)--(4) are used here as hypotheses of the implication; they are not proved.

No comparison principle between $\tilde{x}$ and $\hat{x}$ is available in Ba{\l}aban's setting:
the order argument that compares the minimisers of the model fails once the field is replaced by
its block means, and no such principle is supplied here. Instead, the model theorem,
Theorem~\ref{8.13:thm-inactivity}, establishes in the model, by energy and continuity rather than
by order, the statement that a comparison principle would deliver.
\end{remark}

\begin{remark}[Model: a laterally delocalised hot set]\label{8.13:rem-delocalised-hot-set}
This remark concerns the model of Theorem~\ref{8.13:thm-inactivity}, with the minimisers
$\psi^{(c)}$, $\tilde{\psi}$ (for $\tilde{x}$) and $\hat{\psi}=\psi^{(\tau')}$, $\tau'=\theta_y$
(for $\hat{x}$), fixed there.

By the model lemma, Lemma~\ref{8.13:lem-tight-census}, a young constraint can be active beyond the
tip needle $(\mathrm{N})$ in only two ways:
\begin{itemize}
\item as $(\mathrm{F})$, a young-full column;
\item as $(\mathrm{P}^{\pm})$, a pole hanging from the top of the band.
\end{itemize}
That lemma does not use the mild-fringe hypothesis of Definition~\ref{8.13:def-mild-fringe}, a
model statement. Both alternatives require the first free row
above the band to be pulled by at least $c-2\sigma$, where $\sigma=s_e\theta'$. At $\hat{\psi}$,
where $c=\tau'=\theta_y$, this is
\begin{equation*}
|\hat{F}_{W'}(x)|\ge\theta_y-2s_e\theta'=\theta_y-2\sigma.
\end{equation*}

The model theorem, Theorem~\ref{8.13:thm-inactivity}, excludes both alternatives given the budget
\begin{equation*}
\|\nabla(\hat{\psi}-\tilde{\psi})\|^2\le 4\Delta(\tau'),
\end{equation*}
with $\Delta(c)$ the budget defined in \eqref{8.13:eq-energy-budget-a}. The model lemma,
Lemma~\ref{8.13:lem-energy-budget}, supplies this budget at $c=\tau'$ under the conditions of the
model of Theorem~\ref{8.13:thm-inactivity} and the mild-fringe hypothesis. That hypothesis
asks that the hot set of $\tilde{x}$ sit above the foot of the kink.

When the mild-fringe hypothesis fails, the hot set of $\tilde{x}$ is shallow and laterally
delocalised. This happens in either of two ways.
\begin{itemize}
\item The block of $\tilde{\psi}$ has half-width $\tilde{S}<2\kappa W$, contrary to clause~(i) of
the model hypothesis of Definition~\ref{8.13:def-mild-fringe}; that is, the block is narrower
than of the order of $C\,W/(nL^{2})$ blocks, with $C$ a constant.
\item There is a strong fringe: a column $x$ with $\kappa W\le\|x\|\le P-\kappa W$ whose first free
bond above the clamp is pulled by more than $\theta_y/2$, that is,
$|\tilde{F}_W(x)|>\tau'/2=\theta_y/2$, contrary to clause~(ii) of the model hypothesis of
Definition~\ref{8.13:def-mild-fringe}. Equivalently, there is a straddling multiplier of order at
least $\theta_y/2$ away from the kink.
\end{itemize}
For the straddling-only minimiser this is the one configuration not covered by the model theorem,
Theorem~\ref{8.13:thm-inactivity}.

Even in this configuration, the model lemma, Lemma~\ref{8.13:lem-anchored-runs}, forbids anchored
young activity away from the tip columns. Propagation would therefore still have to take the form
$(\mathrm{F})$ or $(\mathrm{P}^{\pm})$. What fails is only the energy accounting.

Whether the straddling-only minimiser $\tilde{\psi}$ can violate the mild-fringe hypothesis is one
half of the free-boundary problem of the saturation statement: the half set in the region of the
straddling band beside the kink and at its lateral edge, where the lateral confinement of the hot
set of $\tilde{x}$ is decided. The model constrains such a violation. Write
$\tilde{g}=\tilde{\psi}(\cdot,W)$ for the exit trace, $T(x)=\|x\|$, and $D_T$ for the pull of $T$
through the strip representation, so that $\sigma T$ is the exit trace of the kink transported
rigidly to the row $W$ and $\sigma D_T$ its pull. In the model, lowering the rigid exit
trace at columns where the rigid pull $\sigma D_T$ exceeds $\tau$
raises the energy $E$ to first order.

Consider finally a young active set of depth at least of the order of $W/L$ rows into the band,
that is, of relative depth at least of the order of $W/(nL)$.
It would require $(\mathrm{F})$ or $(\mathrm{P}^{\pm})$, and hence
$|\hat{F}_{W'}|\ge\theta_y-2\sigma$ at the top of the band. In the model this is impossible
whenever the budget holds.
\end{remark}

\begin{remark}[Model: scope of the model theorem]\label{8.13:rem-model-scope}
The model is the one of the model theorem, Theorem~\ref{8.13:thm-inactivity}. On the periodic strip
$\Lambda=\mathbb{Z}_N\times\{0,\dots,H\}$, the Dirichlet energy $E(\psi)$ is minimised with the kinked
datum $\varphi_0(x)=\sigma\|x\|$, where $\sigma=s_e\theta'$, over the feasible set $\mathcal{K}_c$ cut
out by clamps at the straddling cap $\tau$ and at the young cap $c$, with minimiser $\psi^{(c)}$.
We describe the scope of that theorem and of the model lemmas on which it rests.

\emph{Geometry of the clamps.} The clamps are componentwise throughout. Each vertical bond difference
$F_y(x)=\psi(x,y+1)-\psi(x,y)$ below the young band is constrained by the straddling cap $\tau$, and
each in the young band by the young cap $c$, bond by bond and with no coupling between bonds. This
box geometry is the one for which $\mathcal{K}_c$ is a lattice, that is, closed under pointwise maximum
and minimum. It is also the one for which the lateral slope bound
\begin{equation*}
|\psi^{(c)}(x+1,y)-\psi^{(c)}(x,y)|\le\sigma
\end{equation*}
holds at every cap $c$.

\emph{Class of data.} The field outside the model strip (the exterior, whose relative kink slope is
$s_e$) enters the model only through the slope $\sigma=s_e\theta'$ and through
the class of the kinked datum: the sharp V-shaped datum $\varphi_0$, or its mollified form treated in the
model remark on the mollified kink, Remark~\ref{8.13:rem-mollified-kink}. General $\sigma$-Lipschitz data
with many points of convexity are not covered by the model lemmas, Lemma~\ref{8.13:lem-anchored-runs}
and Lemma~\ref{8.13:lem-kink-pull-decay}, for two reasons.
\begin{itemize}
\item By part~(d) of the model lemma, Lemma~\ref{8.13:lem-anchored-runs}, every column $x$ with
$\Delta_x\varphi_0(x)>0$ may carry a needle.
\item The model lemma on the decay of the kink's pull, Lemma~\ref{8.13:lem-kink-pull-decay}, uses the
V-shaped datum near the hot box $Q(c)$.
\end{itemize}
The model theorem is a statement for the kinked class.

\emph{Profile of the vertical bound.} Below height $m_{\mathrm{ndl}}$ the needle sits at the cap. At heights
$h\in(2m_{\mathrm{ndl}},\bar W)$ the bound on vertical bond differences is
\begin{equation*}
\sigma\bigl[C_K^{0}+6\log^{+}\bigl(N/(h+1)\bigr)\bigr].
\end{equation*}
For small $h$ this is a logarithm of the width $N$ measured in fine lattice units; the bound decreases
to the order of $\sigma[C_K^{0}+6\log^{+}(N/\bar W)]$ only at heights $h$ comparable with $W$.

\emph{The half-disc picture.} Consider the heuristic picture in which the young active set is a half-disc
of some radius $h_1$. It is centred where the kink's tip meets the inner face of the clamp, and beyond it
the field is harmonic with the kink rounded at scale $h_1$. The model theorem, Theorem~\ref{8.13:thm-inactivity},
refines this picture. On the lattice the active set is not a half-disc but a needle on the tip column,
whose height is at most of the order of the scale $h_1$ of the half-disc picture. The neighbouring columns
follow the tip column to within $\sigma$, by the lateral slope bound. By the model lemma,
Lemma~\ref{8.13:lem-anchored-runs}, these
columns are close to their caps but do not attain them, and the rounding of the kink takes place in this
neighbourhood where no constraint is active. Above the needle the field is exactly the harmonic extension
of a $\sigma$-Lipschitz trace. Thus, under the hypothesis of the model theorem,
Theorem~\ref{8.13:thm-inactivity}, no further active young constraints occur above the needle.
\end{remark}

\begin{definition}[Model: truncated half-space problems with young caps]
\label{8.14:def-half-space-problems}
The model of this definition is a convex minimisation problem for real vector fields on the
half-lattice $\mathbb{T}\times\{0,1,2,\dots\}$: a lateral field on each row $y\ge0$, whose
transverse part is called the trace of the row, and a level field between consecutive rows.
The problem has a quadratic energy, the trace of row $0$ is held fixed, an arbitrary closed
convex constraint acts near the bottom, and caps act on a finite band of levels above it.

\emph{The lateral torus and its operators.} Let $N$ be a positive integer; in this
definition $N$ is the side of the lateral torus, not the rank of the gauge group. Put
\begin{equation*}
  \mathbb{T}:=(\mathbb{Z}/N)^3,\qquad \mathfrak{H}:=\ell^2(\mathbb{T};\mathbb{R}^3).
\end{equation*}
Let $\nabla$ be the lattice gradient on scalars $\ell^2(\mathbb{T};\mathbb{R})$, and let
$\mathfrak{H}_T\subset\mathfrak{H}$ be the space of transverse fields, so that
\begin{equation*}
  \mathfrak{H}=\operatorname{ran}\nabla\oplus\mathfrak{H}_T\oplus\{\text{constants}\}.
\end{equation*}
Let $P$ be the projection of $\mathfrak{H}$ onto $\mathfrak{H}_T$ along this decomposition,
and let $P^{\perp}:=1-P$. Let $\Delta_{\parallel}$ be the lattice Laplacian of $\mathbb{T}$,
and put $\lambda:=-\Delta_{\parallel}$. The operator $\lambda$ acts componentwise on all of
$\mathfrak{H}$ and on the scalars $\ell^2(\mathbb{T};\mathbb{R})$. On $\mathfrak{H}_T$ it
coincides with $\operatorname{curl}^*\operatorname{curl}$, where $\operatorname{curl}$ is the
lateral curl and $\operatorname{curl}^*$ its adjoint, and
$\|\operatorname{curl}\alpha\|^2=\langle P\alpha,\lambda P\alpha\rangle$ for
$\alpha\in\mathfrak{H}$. By functional calculus in $\lambda\ge0$ we put
\begin{equation}\label{8.14:eq-half-space-problems-1}
  \mathsf{v}:=\tfrac12\bigl[-\lambda+(\lambda^2+4\lambda)^{1/2}\bigr],\qquad
  \mathsf{C}:=1-\mathsf{v},\qquad
  \mathsf{U}:=\lambda+\mathsf{v}.
\end{equation}
Squaring the definition of $\mathsf{v}$ gives $\mathsf{v}^2+\lambda\mathsf{v}=\lambda$.
Hence
\begin{equation*}
  (1-\mathsf{v})(1+\lambda+\mathsf{v})=1+\lambda-\lambda\mathsf{v}-\mathsf{v}^2=1,
\end{equation*}
that is, $\mathsf{C}=(1+\lambda+\mathsf{v})^{-1}=(1+\mathsf{U})^{-1}$.

The operators $\mathsf{v}$ and $\mathsf{C}$ are convolutions by scalar kernels acting on
each component, and they commute with $P$ and with the translations of $\mathbb{T}$. For
$j\ge1$ let $K_j$ be the kernel of $\mathsf{C}^j$, so that
$(\mathsf{C}^jf)(x)=\sum_{z\in\mathbb{T}}K_j(z)f(x-z)$ componentwise. Then $K_j=K_1^{*j}$,
and
\begin{equation*}
  K_j\ge0,\qquad \sum_{z\in\mathbb{T}}K_j(z)=1,\qquad K_j(-z)=K_j(z),\qquad
  K_i*K_j=K_{i+j}.
\end{equation*}
We further put
\begin{equation}\label{8.14:eq-half-space-problems-2}
  C_N:=\max\operatorname{spec}\bigl(\mathsf{C}^{-1}\bigr),\qquad C':=\tfrac12(1+C_N).
\end{equation}
For $\varphi\in\mathfrak{H}$ we put
\begin{equation}\label{8.14:eq-half-space-problems-3}
  N(\varphi):=\tfrac12\langle P\varphi,\mathsf{C}^{-1}P\varphi\rangle
  +\tfrac12\|P^{\perp}\varphi\|^2 .
\end{equation}
For $\beta\in\mathfrak{H}_T$ we put $d_0(\beta):=-\mathsf{v}\beta$.
We write $V_+$ for the free value and $V_1$ for the value with one capped level, both as
fixed in the exit-value model theorem, and we put $W:=V_1-V_+$.

\emph{Configurations and energy.} A configuration consists of the traces
$\beta_y\in\mathfrak{H}_T$, $y\ge0$, which are the transverse parts $\beta_y=P\alpha_y$ of
row fields $\alpha_y\in\mathfrak{H}$, and the level fields $E_y\in\mathfrak{H}$, $y\ge0$.
The level field $E_y$ sits between the rows $y$ and $y+1$, and
\begin{equation*}
  \beta_{y+1}=\beta_y+PE_y\qquad(y\ge0).
\end{equation*}
A field $A^0\in\mathfrak{H}$ is given, and the bottom trace is held at
$\beta_0=PA^0$. The energy of a configuration is
\begin{equation}\label{8.14:eq-half-space-problems-4}
  \Phi:=\sum_{y\ge0}\tfrac12\|E_y\|^2+\sum_{y\ge1}\tfrac12\langle\beta_y,\lambda\beta_y\rangle .
\end{equation}
By the identity $\|\operatorname{curl}\alpha\|^2=\langle P\alpha,\lambda P\alpha\rangle$, the
second sum equals $\sum_{y\ge1}\frac12\|\operatorname{curl}\alpha_y\|^2$.

\emph{Lower constraints.} Fix an integer $a_0\ge0$. Let $\mathcal{L}_{\mathrm{low}}$ be an
arbitrary closed convex set of the variables carried by the plaquettes within the rows
$0\le y\le a_0$, that is, of the levels $(E_y)_{y<a_0}$ and the traces
$(\beta_y)_{y\le a_0}$. The lower constraint is that these variables lie in
$\mathcal{L}_{\mathrm{low}}$.

\emph{Young caps.} Let $\mathfrak{B}\subset\mathfrak{H}$ be closed and convex, contain $0$,
and be invariant under the translations of $\mathbb{T}$. In the items of the exit-value model
theorem that use the box, the cap on the level $y$ is the box
\begin{equation}\label{8.14:eq-half-space-problems-5}
  \mathfrak{B}_y=\bigl\{E\in\mathfrak{H}:\ |E(x)|\le\tau_y\ \text{for all } x\in\mathbb{T}\bigr\}
\end{equation}
with cap height $\tau_y\ge0$, where $|E(x)|$ is a fixed norm of the vector
$E(x)\in\mathbb{R}^3$, the one used in the exit-value model theorem. Each $\mathfrak{B}_y$ is
closed, convex, contains $0$ and is invariant under the translations of $\mathbb{T}$; in those
items the set $\mathfrak{B}$ in the cap on the level $y$ is $\mathfrak{B}_y$.

\emph{The truncated problems.} For an integer $a\ge a_0$, let $P_a$ be the problem of
minimising $\Phi$ of \eqref{8.14:eq-half-space-problems-4} over configurations subject to
three conditions: $\beta_0=PA^0$, the lower constraint, and the young caps
\begin{equation*}
  E_y\in\mathfrak{B}\qquad(a_0\le y<a),
\end{equation*}
with $\mathfrak{B}_y$ in place of $\mathfrak{B}$ on the level $y$ in the items that use the box.
Let $m(a)$ be the value of $P_a$ and $X_a$ its minimiser, and put
\begin{equation}\label{8.14:eq-half-space-problems-6}
  \Delta m(a):=m(a+1)-m(a)\ge0 .
\end{equation}
The inequality in \eqref{8.14:eq-half-space-problems-6} holds because the constraints of
$P_{a+1}$ are those of $P_a$ together with the cap $E_a\in\mathfrak{B}$ on the level $a$
($E_a\in\mathfrak{B}_a$ in the items that use the box). The admissible set of
$P_{a+1}$ is therefore contained in that of $P_a$.
\end{definition}

\begin{lemma}[Model: Poisson operator identities, uniform in lateral size]
\label{8.14:lem-poisson-identities}
The model is that of Definition~\ref{8.14:def-half-space-problems}, which is a model statement and
not Ba{\l}aban's renormalisation step. It consists of the truncated half-space problems with young
caps over the lateral torus $\mathbb{T}=(\mathbb{Z}/N)^3$, where $N$ is the lateral size of
the model (not the rank of the gauge group). Its level fields lie in
$\mathfrak{H}=\ell^2(\mathbb{T};\mathbb{R}^3)$, and $\mathfrak{H}_T$ is the transverse part of
$\mathfrak{H}$.

Let $e_1,e_2,e_3$ be the unit vectors of the three lateral directions, and let
$\lambda=-\Delta_{\parallel}$ be the lateral lattice Laplacian,
$\Delta_{\parallel}u(x)=\sum_{k=1}^{3}[u(x+e_k)+u(x-e_k)-2u(x)]$, acting componentwise on
$\mathfrak{H}$ and on $\ell^2(\mathbb{T};\mathbb{R})$. It is self-adjoint and non-negative. By the
functional calculus of $\lambda$ put
\begin{equation*}
\mathsf{v}:=\tfrac12\bigl[-\lambda+(\lambda^2+4\lambda)^{1/2}\bigr],\qquad
\mathsf{C}:=1-\mathsf{v},\qquad \mathsf{U}:=\lambda+\mathsf{v}.
\end{equation*}
Let $K_1$ be the kernel of $\mathsf{C}$, so that $\mathsf{C}u(x)=\sum_{z}K_1(x-z)u(z)$. Put
$d_0(\beta):=-\mathsf{v}\beta$ for $\beta\in\mathfrak{H}$. Then the following hold.
\begin{enumerate}
\item[(a)] The operator $\mathsf{v}$ is the non-negative root of
\begin{equation}\label{8.14:eq-poisson-identities-a}
\mathsf{v}^2+\lambda\mathsf{v}=\lambda .
\end{equation}
\item[(b)] The operator $\mathsf{C}$ is invertible, and
\begin{equation}\label{8.14:eq-poisson-identities-b}
(1-\mathsf{v})(1+\lambda+\mathsf{v})=1,\qquad \mathsf{C}=(1+\mathsf{U})^{-1},\qquad
\mathsf{C}^{-1}=2+\lambda-\mathsf{C}.
\end{equation}
\item[(c)] The following identity holds:
\begin{equation}\label{8.14:eq-poisson-identities-c}
\mathsf{U}\mathsf{C}=\mathsf{v}.
\end{equation}
Hence $\mathsf{v}\lambda^{-1}=(\lambda+\mathsf{v})^{-1}$ on mean-zero data and on transverse data.
\item[(d)] The following identity holds:
\begin{equation}\label{8.14:eq-poisson-identities-d}
\mathsf{v}+\mathsf{C}=1.
\end{equation}
\item[(e)] For every $\beta\in\mathfrak{H}$,
\begin{equation}\label{8.14:eq-poisson-identities-e}
(1+\mathsf{U})d_0(\beta)=-\mathsf{U}\beta,\qquad
\tfrac12\langle d_0(\beta),(1+\mathsf{U})d_0(\beta)\rangle=\tfrac12\langle\beta,\lambda\beta\rangle .
\end{equation}
\item[(f)] In the sense of quadratic forms,
\begin{equation}\label{8.14:eq-poisson-identities-f}
1-\mathsf{C}^2=\mathsf{v}(1+\mathsf{C}),\qquad \mathsf{v}\le 1-\mathsf{C}^2\le 2\mathsf{v},\qquad
1-\mathsf{C}\le 1-\mathsf{C}^2 .
\end{equation}
\item[(g)] (Poisson.) Mode by mode, for given $u_0$ the recurrence
$u_{y+1}+u_{y-1}-(2+\lambda)u_y=0$ for $y\ge1$ has exactly one solution bounded in $y$, namely
\begin{equation}\label{8.14:eq-poisson-identities-g}
u_y=\mathsf{C}^{y}u_0\qquad(y\ge0).
\end{equation}
Moreover $\mathsf{C}=e^{-\eta}$ with $\eta\ge0$ and $2\cosh\eta-2=\lambda$. The kernel $K_1(x-z)$
is the harmonic measure of $(z,0)$ on row $0$ for the random walk on $\mathbb{T}\times\mathbb{Z}$
that at each step makes one of the eight moves $(\pm e_k,0)$ $(k=1,2,3)$, $(0,\pm1)$, each with
probability $1/8$, started at $(x,1)$. In particular $K_1\ge0$,
$\sum_z K_1(z)=1$ and $\mathsf{C}1=1$.
\item[(h)] Write $\mathsf{v}(s):=\tfrac12[-s+(s^2+4s)^{1/2}]$ for a number $s\ge0$.
The number $C_N$, the largest eigenvalue of $\mathsf{C}^{-1}$ on
$\ell^2(\mathbb{T};\mathbb{R})$, is bounded independently of $N$. Indeed the spectrum of
$\lambda$ on $\ell^2(\mathbb{T};\mathbb{R})$ is the set of the numbers
$\sum_{k=1}^{3}4\sin^2(\pi\xi_k/N)$, $\xi\in(\mathbb{Z}/N)^3$, which lies in $[0,12]$ for every
$N$, and $1+s+\mathsf{v}(s)$, the value of $\mathsf{C}^{-1}$ at the spectral point $s$, is
increasing in $s$; hence
\begin{equation}\label{8.14:eq-poisson-identities-h}
C_N\le 1+12+\mathsf{v}(12)=7+4\sqrt3\qquad\text{for every }N .
\end{equation}
\end{enumerate}
\end{lemma}

\begin{proof}
All operators involved are real functions of the single self-adjoint operator $\lambda$. Hence
they commute with one another and with lattice translations. We use the same letter $\lambda$ for
the operator and for a point of its spectrum: each identity is an identity between functions of
$\lambda$ on the spectrum, and is checked at one point $\lambda\ge0$ of the spectrum at a time.

Eigenvalues of $\lambda$ are computed from the functions
$x\mapsto N^{-3/2}\exp(2\pi i\,\xi\cdot x/N)$, $\xi\in(\mathbb{Z}/N)^3$, which form an
orthonormal eigenbasis of $\ell^2(\mathbb{T};\mathbb{C})$. Indeed, applied to such a function,
$-\Delta_{\parallel}$ multiplies it by
\begin{equation*}
\sum_{k=1}^{3}\bigl(2-2\cos(2\pi\xi_k/N)\bigr)
=\sum_{k=1}^{3}4\sin^2(\pi\xi_k/N) .
\end{equation*}
The spectrum is therefore contained in $[0,12]$. The eigenvalue is $0$ only for $\xi=0$, that is,
on the constants.

Preliminary: $0\le\mathsf{v}<1$. Read as a quadratic equation for the unknown $\mathsf{v}$,
equation \eqref{8.14:eq-poisson-identities-a} has the two roots
$\tfrac12[-\lambda\pm(\lambda^2+4\lambda)^{1/2}]$, and $\mathsf{v}$ is the one with the plus sign.
It is non-negative because $(\lambda^2+4\lambda)^{1/2}\ge\lambda$, and it is less than $1$ because
$\lambda^2+4\lambda<(\lambda+2)^2$. Consequently $0<\mathsf{C}\le1$.

(a) By the preliminary, $\mathsf{v}$ is a root of \eqref{8.14:eq-poisson-identities-a}, read as a
quadratic equation, and it is the non-negative root.

(b) Expanding and using \eqref{8.14:eq-poisson-identities-a},
\begin{equation*}
(1-\mathsf{v})(1+\lambda+\mathsf{v})=1+\lambda-\lambda\mathsf{v}-\mathsf{v}^2=1 .
\end{equation*}
Thus $\mathsf{C}=1-\mathsf{v}$ is invertible with $\mathsf{C}^{-1}=1+\lambda+\mathsf{v}$. Since
$\mathsf{U}=\lambda+\mathsf{v}$, this inverse equals $1+\mathsf{U}$. Since $\mathsf{v}=1-\mathsf{C}$,
it also equals $2+\lambda-\mathsf{C}$.

(c) Expanding and using \eqref{8.14:eq-poisson-identities-a},
\begin{equation*}
\mathsf{U}\mathsf{C}=(\lambda+\mathsf{v})(1-\mathsf{v})
=\lambda+\mathsf{v}-(\lambda\mathsf{v}+\mathsf{v}^2)=\lambda+\mathsf{v}-\lambda=\mathsf{v}.
\end{equation*}
The same computation shows $\mathsf{v}(\lambda+\mathsf{v})=\lambda$. The constants form the kernel
of $\lambda$, so on their orthogonal complement $\lambda$ is invertible, and so is
$\lambda+\mathsf{v}\ge\lambda$. Dividing $\mathsf{v}(\lambda+\mathsf{v})=\lambda$ by $\lambda$ and by
$\lambda+\mathsf{v}$, which commute, gives $\mathsf{v}\lambda^{-1}=(\lambda+\mathsf{v})^{-1}$ on
mean-zero data. Transverse data are orthogonal to the constant fields in the orthogonal
decomposition of $\mathfrak{H}$ into gradients, transverse fields and constants of the model of
Definition~\ref{8.14:def-half-space-problems}. The identity therefore holds on
transverse data as well.

(d) This is the definition $\mathsf{C}=1-\mathsf{v}$.

(e) By (b),
\begin{equation*}
(1+\mathsf{U})d_0(\beta)=-(1+\lambda+\mathsf{v})\mathsf{v}\beta .
\end{equation*}
Moreover $(1+\lambda+\mathsf{v})\mathsf{v}=\mathsf{v}+\lambda\mathsf{v}+\mathsf{v}^2=\mathsf{v}+\lambda$
by \eqref{8.14:eq-poisson-identities-a}. Hence
$(1+\mathsf{U})d_0(\beta)=-(\mathsf{v}+\lambda)\beta=-\mathsf{U}\beta$. Since $\mathsf{v}$ is
self-adjoint and commutes with $\mathsf{U}$,
\begin{equation*}
\tfrac12\langle d_0(\beta),(1+\mathsf{U})d_0(\beta)\rangle
=\tfrac12\langle-\mathsf{v}\beta,-\mathsf{U}\beta\rangle
=\tfrac12\langle\beta,\mathsf{v}(\lambda+\mathsf{v})\beta\rangle
=\tfrac12\langle\beta,\lambda\beta\rangle .
\end{equation*}

(f) First, $1-\mathsf{C}^2=(1-\mathsf{C})(1+\mathsf{C})=\mathsf{v}(1+\mathsf{C})$ by (d). The
operators $\mathsf{v}$ and $\mathsf{C}$ commute, $\mathsf{v}\ge0$, and $0\le\mathsf{C}\le1$. Hence
$0\le\mathsf{v}\mathsf{C}\le\mathsf{v}$, which gives
$\mathsf{v}\le\mathsf{v}(1+\mathsf{C})\le2\mathsf{v}$. Finally
\begin{equation*}
1-\mathsf{C}=\mathsf{v}\le\mathsf{v}(1+\mathsf{C})=1-\mathsf{C}^2 .
\end{equation*}

(g) Fix a point $\lambda$ of the spectrum. The characteristic equation of the scalar recurrence is
$r^2-(2+\lambda)r+1=0$. The numbers $1+\lambda+\mathsf{v}$ and $1-\mathsf{v}=\mathsf{C}$ satisfy
\begin{equation*}
(1+\lambda+\mathsf{v})+(1-\mathsf{v})=2+\lambda,\qquad (1+\lambda+\mathsf{v})(1-\mathsf{v})=1,
\end{equation*}
the second by (b). They are therefore its two roots, $\mathsf{C}$ and $\mathsf{C}^{-1}$.

If $\lambda>0$, then $\mathsf{v}>0$ and $\mathsf{C}<1$. The general solution is
$c_+\mathsf{C}^y+c_-\mathsf{C}^{-y}$ with constants $c_\pm$, and it is bounded only if $c_-=0$, so
$u_y=\mathsf{C}^yu_0$. If $\lambda=0$, then $\mathsf{C}=1$ is a double root. The general solution
is $c_++c_-y$, and it is bounded only if $c_-=0$, so again $u_y=u_0=\mathsf{C}^yu_0$.
Conversely, in both cases $\mathsf{C}^yu_0$ is of the general form with $c_-=0$ and $c_+=u_0$, so it
solves the recurrence, and it is bounded because $0<\mathsf{C}\le1$.

Since $0<\mathsf{C}\le1$ we may write $\mathsf{C}=e^{-\eta}$ with $\eta\ge0$. Then
$2\cosh\eta=\mathsf{C}+\mathsf{C}^{-1}=2+\lambda$.

For the harmonic measure, write the recurrence of item (g) as an
equation for $u(x,y)=u_y(x)$. Using $\lambda u_y=6u_y-\sum_k[u_y(\cdot+e_k)+u_y(\cdot-e_k)]$, it
reads
\begin{equation*}
8u(x,y)=\sum_{k=1}^3\bigl[u(x+e_k,y)+u(x-e_k,y)\bigr]+u(x,y+1)+u(x,y-1),\qquad y\ge1 .
\end{equation*}
So $u$ is harmonic at every point of the rows $y\ge1$ for the random walk
$(\mathsf{Y}_n)_{n\ge0}$ on $\mathbb{T}\times\mathbb{Z}$ which at each step makes one of the eight
moves $(\pm e_k,0)$, $k=1,2,3$, and $(0,\pm1)$, each with probability $1/8$.

Let $\mathsf{n}_0$ be the first $n$ with $\mathsf{Y}_n$ in row $0$. The vertical coordinate of
$\mathsf{Y}$ changes at each step with probability $1/4$, by $\pm1$ with equal probabilities.
Observed at its times of change it is a simple random walk on $\mathbb{Z}$, which is recurrent.
Hence $\mathsf{n}_0<\infty$ almost surely.

Now take $u_y=\mathsf{C}^yu_0$ for a given $u_0$. This field is bounded, because
$\|\mathsf{C}^yu_0\|_\infty\le\|\mathsf{C}^yu_0\|_2\le\|u_0\|_2$. Therefore
$u(\mathsf{Y}_{n\wedge\mathsf{n}_0})$ is a bounded martingale for the walk started at $(x,1)$.
Optional stopping, followed by bounded convergence as $n\to\infty$, gives
\begin{equation*}
\mathsf{C}u_0(x)=u(x,1)
=\sum_{z\in\mathbb{T}}\Pr\nolimits_{(x,1)}\bigl(\mathsf{Y}_{\mathsf{n}_0}=(z,0)\bigr)\,u_0(z).
\end{equation*}
By translation invariance the probability on the right depends only on $x-z$. It is $K_1(x-z)$,
the harmonic measure of $(z,0)$ on row $0$ for the walk started one row up. Hence $K_1\ge0$ and
$\sum_zK_1(z)=1$. In agreement with this, $\lambda1=0$ gives $\mathsf{v}1=0$, that is,
$\mathsf{C}1=1$.

(h) By (b), $\mathsf{C}^{-1}=1+\lambda+\mathsf{v}(\lambda)$. For $\lambda>0$,
\begin{equation*}
\mathsf{v}(\lambda)=\frac{2\lambda}{\lambda+(\lambda^2+4\lambda)^{1/2}}
=\frac{2}{1+(1+4/\lambda)^{1/2}},
\end{equation*}
and $\mathsf{v}(0)=0$. So $\mathsf{v}$, and with it $1+\lambda+\mathsf{v}(\lambda)$, is
increasing in $\lambda\ge0$. The spectrum of $-\Delta_{\parallel}$ on $(\mathbb{Z}/N)^3$ is
contained in $[0,12]$ for every $N$, as computed at the start. Moreover
$\mathsf{v}(12)=\tfrac12(-12+\sqrt{192})=-6+4\sqrt3$. Therefore
\begin{equation*}
C_N\le1+12+\mathsf{v}(12)=13+(-6+4\sqrt3)=7+4\sqrt3
\end{equation*}
for every $N$.
\end{proof}

\begin{lemma}[Model: the innovations identity]\label{8.14:lem-innovations-identity}
This lemma is a statement in the model set-up of Definition~\ref{8.14:def-half-space-problems}, which is
itself a model and not Ba{\l}aban's step. That set-up consists of a half-lattice with lateral torus
$\mathbb{T}$ and rows $y\ge 0$, and level fields $E_y\in\mathfrak{H}=\ell^2(\mathbb{T};\mathbb{R}^3)$
between the rows $y$ and $y+1$. The transverse traces $\beta_y\in\mathfrak{H}_T$ satisfy
$\beta_{y+1}=\beta_y+PE_y$. Here $\lambda=-\Delta_{\parallel}$ acts componentwise on $\mathfrak{H}$.

For a row $r$ and levels $(E_y)_{y\ge r}$ above it, with traces $\beta_{y+1}=\beta_y+PE_y$ starting
from $\beta_r$, the \emph{energy above row $r$} is
\[
\sum_{y\ge r}\tfrac12\|E_y\|^2+\sum_{y\ge r+1}\tfrac12\langle\beta_y,\lambda\beta_y\rangle .
\]
The \emph{free tail} from a row is the configuration above it of least energy, given the trace on that row.

For $\gamma\in\mathfrak{H}_T$ put $V_+(\gamma):=\tfrac12\langle\gamma,\mathsf{v}\gamma\rangle$. We use the
free pull $d_0(\beta)=-\mathsf{v}\beta$ of $\beta\in\mathfrak{H}_T$, which enters
\eqref{8.14:eq-poisson-identities-e}, and the defect form, given for $E\in\mathfrak{H}$ by
\[
N(E)=\tfrac12\langle PE,\mathsf{C}^{-1}PE\rangle+\tfrac12\|P^{\perp}E\|^2 .
\]
Here $\mathsf{v}$ and $\mathsf{C}=1-\mathsf{v}$ are the operators of the model lemma,
Lemma~\ref{8.14:lem-poisson-identities}.

Then for every $\beta\in\mathfrak{H}_T$ and every $E\in\mathfrak{H}$, with $\beta':=\beta+PE$,
\begin{equation}\label{8.14:eq-innovations-identity-a}
\tfrac12\|E\|^2+\tfrac12\langle\beta',\lambda\beta'\rangle+V_+(\beta')
=V_+(\beta)+N\bigl(d_0(\beta)-E\bigr).
\end{equation}

Consequently, if a configuration consists of levels $E_r,\dots,E_{s-1}$ above a row $r$ followed by
the free tail from row $s$, its energy above row $r$ is
\begin{equation}\label{8.14:eq-innovations-identity-1}
V_+(\beta_r)+\sum_{r\le y<s}N\bigl(d_0(\beta_y)-E_y\bigr).
\end{equation}
For an arbitrary configuration of finite energy, the partial sums
\[
V_+(\beta_r)+\sum_{r\le y<s}N\bigl(d_0(\beta_y)-E_y\bigr)
\]
increase to its energy above row $r$ as $s\to\infty$.
\end{lemma}

\begin{proof}
The operators $\mathsf{v}$, $\mathsf{C}$ and $\mathsf{U}=\lambda+\mathsf{v}$ are functions of $\lambda$.
They are therefore self-adjoint, commute with one another and commute with $P$, as fixed in the model
set-up of Definition~\ref{8.14:def-half-space-problems}. Since
$\lambda\ge 0$, we have $\mathsf{C}^{-1}=1+\lambda+\mathsf{v}\ge 1$ by \eqref{8.14:eq-poisson-identities-b}
of the model lemma, Lemma~\ref{8.14:lem-poisson-identities}. In particular $N\ge 0$.

\emph{The free tail and its value.}
Fix a row $r$ and a trace $\gamma\in\mathfrak{H}_T$, and minimise the energy above row $r$ among
configurations with $\beta_r=\gamma$. The component $P^{\perp}E_y$ of a level enters the energy only through
$\tfrac12\|P^{\perp}E_y\|^2$ and does not change any trace. It therefore vanishes at the minimiser, whose
levels are $E_y=\beta_{y+1}-\beta_y$. For $i\ge 1$ the terms of the energy containing $\beta_{r+i}$ are
\[
\tfrac12\|\beta_{r+i}-\beta_{r+i-1}\|^2+\tfrac12\|\beta_{r+i+1}-\beta_{r+i}\|^2
+\tfrac12\langle\beta_{r+i},\lambda\beta_{r+i}\rangle .
\]
Their gradient in $\beta_{r+i}$ is $(2+\lambda)\beta_{r+i}-\beta_{r+i-1}-\beta_{r+i+1}$. The traces of the
minimiser therefore solve the recursion of \eqref{8.14:eq-poisson-identities-g} with initial value $\gamma$.
Finite energy selects the decaying mode, which by \eqref{8.14:eq-poisson-identities-g} is
$\beta_{r+i}=\mathsf{C}^i\gamma$. This is the decaying lattice-harmonic extension of $\gamma$, and its levels
are
\[
E_{r+i}=(\mathsf{C}-1)\mathsf{C}^i\gamma=-\mathsf{v}\mathsf{C}^i\gamma .
\]
Its energy is thus a quadratic form in $\gamma$ given by a function of $\lambda$. We write this least
energy as $\tfrac12\langle\gamma,\mathsf{v}_+\gamma\rangle$, where $\mathsf{v}_+\ge 0$ is a function of
$\lambda$; it is non-negative because the energy is. We show that $\mathsf{v}_+=\mathsf{v}$.

Split off the first level. The energy above row $r$ equals
\[
\tfrac12\|E_r\|^2+\tfrac12\langle\beta_{r+1},\lambda\beta_{r+1}\rangle+\text{(energy above row $r+1$)}.
\]
Minimising first over the levels above row $r+1$ replaces the last term by
$\tfrac12\langle\beta_{r+1},\mathsf{v}_+\beta_{r+1}\rangle$. The component $P^{\perp}E_r$ costs
$\tfrac12\|P^{\perp}E_r\|^2$ and changes nothing above, so it vanishes at the minimum. Writing
$e=PE_r\in\mathfrak{H}_T$ and $\mathsf{U}_+:=\lambda+\mathsf{v}_+$, we obtain the Bellman equation
\[
\tfrac12\langle\gamma,\mathsf{v}_+\gamma\rangle
=\min_{e\in\mathfrak{H}_T}\Bigl[\tfrac12\|e\|^2+\tfrac12\langle\gamma+e,\mathsf{U}_+(\gamma+e)\rangle\Bigr].
\]
The bracket is a strictly convex quadratic function of $e$. Its gradient $e+\mathsf{U}_+(\gamma+e)$ vanishes
exactly at
\[
e=-\mathsf{U}_+(1+\mathsf{U}_+)^{-1}\gamma .
\]
There $\gamma+e=(1+\mathsf{U}_+)^{-1}\gamma$, and the value is
\[
\tfrac12\bigl\langle\gamma,(\mathsf{U}_+^2+\mathsf{U}_+)(1+\mathsf{U}_+)^{-2}\gamma\bigr\rangle
=\tfrac12\bigl\langle\gamma,\mathsf{U}_+(1+\mathsf{U}_+)^{-1}\gamma\bigr\rangle .
\]
Both sides of the Bellman equation are quadratic forms $\tfrac12\langle\gamma,A\gamma\rangle$ on
$\mathfrak{H}_T$ whose operators $A$ are self-adjoint functions of $\lambda$; equality for all $\gamma$
therefore gives
$\mathsf{v}_+=\mathsf{U}_+(1+\mathsf{U}_+)^{-1}$, that is,
$\mathsf{v}_+(1+\lambda+\mathsf{v}_+)=\lambda+\mathsf{v}_+$, that is,
\[
\mathsf{v}_+^2+\lambda\mathsf{v}_+=\lambda .
\]
On each eigenspace of $\lambda$ the scalar quadratic equation obtained from the last display has exactly
one non-negative root, since the product of its roots is minus the eigenvalue, which is $\le 0$. As
$\mathsf{v}_+\ge 0$, it follows that
$\mathsf{v}_+=\mathsf{v}$, the non-negative root of \eqref{8.14:eq-poisson-identities-a}. Consequently
$\mathsf{U}_+=\mathsf{U}$, and $(1+\mathsf{U})^{-1}=\mathsf{C}$ by \eqref{8.14:eq-poisson-identities-b}.
The minimising first level is
\[
e=-\mathsf{U}\mathsf{C}\gamma=-\mathsf{v}\gamma=d_0(\gamma)
\]
by \eqref{8.14:eq-poisson-identities-c}. Thus the free tail from a row of trace $\gamma$ has energy
$V_+(\gamma)$.

\emph{Direct evaluation of the same value.}
Summing the energy of the explicit tail $\beta_{r+i}=\mathsf{C}^i\gamma$, $E_{r+i}=-\mathsf{v}\mathsf{C}^i\gamma$
as two geometric series gives
\[
\sum_{i\ge 0}\tfrac12\|\mathsf{v}\mathsf{C}^i\gamma\|^2
+\sum_{i\ge 1}\tfrac12\langle\mathsf{C}^i\gamma,\lambda\mathsf{C}^i\gamma\rangle
=\tfrac12\bigl\langle\gamma,(\mathsf{v}^2+\lambda\mathsf{C}^2)(1-\mathsf{C}^2)^{-1}\gamma\bigr\rangle .
\]
This equals $\tfrac12\langle\gamma,\mathsf{v}\gamma\rangle$, because
$\mathsf{v}(1-\mathsf{C}^2)=\mathsf{v}^2+\lambda\mathsf{C}^2$. To see this, note that $\mathsf{C}=1-\mathsf{v}$
gives $1-\mathsf{C}^2=2\mathsf{v}-\mathsf{v}^2$. Using \eqref{8.14:eq-poisson-identities-a} in the form
$\mathsf{v}^2=\lambda-\lambda\mathsf{v}$ at each occurrence of $\mathsf{v}^2$, we get
\begin{align*}
\mathsf{v}(2\mathsf{v}-\mathsf{v}^2)
&=2(\lambda-\lambda\mathsf{v})-\mathsf{v}(\lambda-\lambda\mathsf{v})
=2\lambda-3\lambda\mathsf{v}+\lambda\mathsf{v}^2
=2\lambda-3\lambda\mathsf{v}+\lambda^2-\lambda^2\mathsf{v},\\
\mathsf{v}^2+\lambda(1-\mathsf{v})^2
&=(\lambda-\lambda\mathsf{v})+\lambda-2\lambda\mathsf{v}+\lambda\mathsf{v}^2
=2\lambda-3\lambda\mathsf{v}+\lambda^2-\lambda^2\mathsf{v}.
\end{align*}

\emph{The identity \eqref{8.14:eq-innovations-identity-a}.}
Put $p:=PE$, so that $\beta'=\beta+p$ and $\|E\|^2=\|p\|^2+\|P^{\perp}E\|^2$. Then
$\tfrac12\langle\beta',\lambda\beta'\rangle+V_+(\beta')=\tfrac12\langle\beta+p,\mathsf{U}(\beta+p)\rangle$.
Expanding with $\mathsf{U}$ self-adjoint and using $\mathsf{U}-\mathsf{v}=\lambda$, the left side minus
$V_+(\beta)$ equals
\begin{align*}
&\tfrac12\|P^{\perp}E\|^2+\tfrac12\|p\|^2+\tfrac12\langle\beta+p,\mathsf{U}(\beta+p)\rangle
-\tfrac12\langle\beta,\mathsf{v}\beta\rangle\\
&\qquad=\tfrac12\|P^{\perp}E\|^2+\tfrac12\langle p,(1+\mathsf{U})p\rangle+\langle p,\mathsf{U}\beta\rangle
+\tfrac12\langle\beta,\lambda\beta\rangle .
\end{align*}
On the right, $d_0=d_0(\beta)=-\mathsf{v}\beta$ is transverse, because $\mathsf{v}$ commutes with $P$ and
$\beta\in\mathfrak{H}_T$. Hence $P(d_0-E)=d_0-p$ and $P^{\perp}(d_0-E)=-P^{\perp}E$. With
$\mathsf{C}^{-1}=1+\mathsf{U}$ from \eqref{8.14:eq-poisson-identities-b}, the right side minus $V_+(\beta)$
equals
\begin{align*}
&\tfrac12\langle d_0-p,(1+\mathsf{U})(d_0-p)\rangle+\tfrac12\|P^{\perp}E\|^2\\
&\qquad=\tfrac12\langle p,(1+\mathsf{U})p\rangle-\langle p,(1+\mathsf{U})d_0\rangle
+\tfrac12\langle d_0,(1+\mathsf{U})d_0\rangle+\tfrac12\|P^{\perp}E\|^2\\
&\qquad=\tfrac12\langle p,(1+\mathsf{U})p\rangle+\langle p,\mathsf{U}\beta\rangle
+\tfrac12\langle\beta,\lambda\beta\rangle+\tfrac12\|P^{\perp}E\|^2 .
\end{align*}
In the last step we used $(1+\mathsf{U})d_0=-\mathsf{U}\beta$ and
$\tfrac12\langle d_0,(1+\mathsf{U})d_0\rangle=\tfrac12\langle\beta,\lambda\beta\rangle$, both from
\eqref{8.14:eq-poisson-identities-e}. The two sides agree, which proves
\eqref{8.14:eq-innovations-identity-a}.

\emph{Summation.}
Apply \eqref{8.14:eq-innovations-identity-a} with $\beta=\beta_y$ and $E=E_y$, so that $\beta'=\beta_{y+1}$.
This gives
\[
\tfrac12\|E_y\|^2+\tfrac12\langle\beta_{y+1},\lambda\beta_{y+1}\rangle
=V_+(\beta_y)-V_+(\beta_{y+1})+N\bigl(d_0(\beta_y)-E_y\bigr).
\]
Summing over $r\le y<s$, the differences telescope. We obtain
\[
\sum_{r\le y<s}\Bigl[\tfrac12\|E_y\|^2+\tfrac12\langle\beta_{y+1},\lambda\beta_{y+1}\rangle\Bigr]+V_+(\beta_s)
=V_+(\beta_r)+\sum_{r\le y<s}N\bigl(d_0(\beta_y)-E_y\bigr).
\]
If the configuration continues by the free tail from row $s$, then $V_+(\beta_s)$ is the energy above row
$s$, by the first part of the proof. The left side is then the energy above row $r$, which proves
\eqref{8.14:eq-innovations-identity-1}.

For an arbitrary configuration of finite energy, the right side of the last display is the partial sum of
\eqref{8.14:eq-innovations-identity-1} up to $s$. It is non-decreasing in $s$ because $N\ge 0$. Since
$V_+(\beta_s)$ is the least energy above row $s$, the energy above row $r$ minus this partial sum equals
the energy above row $s$ minus $V_+(\beta_s)$. This difference therefore lies between $0$ and the energy
above row $s$. The latter is the tail of a convergent series and tends to $0$. Hence the partial sums
increase to the energy above row $r$.
\end{proof}

\begin{remark}[Model: the truncated problem as a finite convex control problem]
We apply the lemma to the truncated problems $P_a$ of Definition~\ref{8.14:def-half-space-problems}, a model
statement, with lower constraints on the rows $\le a_0$ and capped levels $E_y\in\mathfrak{B}$ for
$a_0\le y<a$. The levels $E_y$ with $y\ge a$ are unconstrained there, so minimising over them replaces the
energy above row $a$ by $V_+(\beta_a)$. By \eqref{8.14:eq-innovations-identity-1} applied with $r=a_0$ and
$s=a$, the problem $P_a$ is then the following control problem: minimise the energy of the levels below
$a_0$ and of the rows $1,\dots,a_0$, together with the remaining terms, that is,
\begin{multline*}
\sum_{y<a_0}\tfrac12\|E_y\|^2+\sum_{1\le y\le a_0}\tfrac12\langle\beta_y,\lambda\beta_y\rangle\\
+V_+(\beta_{a_0})+\sum_{a_0\le y<a}N\bigl(d_0(\beta_y)-E_y\bigr),
\end{multline*}
over the variables of the lower constraints and over $E_y\in\mathfrak{B}$ for $a_0\le y<a$. The free tail
above row $a$ contributes $0$ to the sum of the terms $N(\cdot)$.

Feasibility of the lower constraints alone makes $P_a$ feasible: take $E_y:=0$ on the levels
$a_0\le y<a$, which is allowed since $0\in\mathfrak{B}$. By the same identity, the objective equals
\[
\sum_{y<a}\tfrac12\|E_y\|^2+\sum_{1\le y\le a}\tfrac12\langle\beta_y,\lambda\beta_y\rangle+V_+(\beta_a),
\]
in finitely many real variables, since $\mathbb{T}$ is finite. The traces are affine in the levels, and
$\lambda$ and $\mathsf{v}$ are non-negative. The objective is therefore strictly convex through the terms
$\tfrac12\|E_y\|^2$, and it is coercive, being at least $\sum_{y<a}\tfrac12\|E_y\|^2$. It is minimised over
a non-empty closed convex set. Hence the minimiser $X_a$ exists and is unique in the level variables
$(E_y)$.
\end{remark}

\begin{theorem}[Model: the exit-value theorem]\label{8.14:thm-exit-value}
The model is the family of truncated half-space problems with young caps of
Definition~\ref{8.14:def-half-space-problems}; it is a convex minimisation problem on the rows
$y\ge0$ over the lateral torus $\mathbb{T}$, and not Ba{\l}aban's renormalisation step. We use
the setting and notation of Definition~\ref{8.14:def-half-space-problems} and of
Lemmas~\ref{8.14:lem-poisson-identities} and~\ref{8.14:lem-innovations-identity} (model
statements). These include the traces $\beta_y\in\mathfrak{H}_T$ and the levels
$E_y\in\mathfrak{H}$, the projection $P$ and $P^{\perp}=1-P$, and the problems $P_a$
($a\ge a_0$) with their values $m(a)$, minimisers $X_a$ and increments
$\Delta m(a)=m(a+1)-m(a)$. They also include the young caps $\mathfrak{B}\subset\mathfrak{H}$,
which are closed, convex, contain $0$ and are invariant under lateral translations. Where the
box is named, $\mathfrak{B}$ is the set of $E\in\mathfrak{H}$ with $|E_k(x)|\le\tau_k$ for all
$x\in\mathbb{T}$ and $k=1,2,3$, with the cap heights $\tau_1,\tau_2,\tau_3$ of
Definition~\ref{8.14:def-half-space-problems}. We write $\beta^{(a)}_y$ and $E^{(a)}_y$ for the
traces and levels of $X_a$. The operators $\mathsf{v}=\tfrac12[-\lambda+(\lambda^2+4\lambda)^{1/2}]$
and $\mathsf{C}=1-\mathsf{v}$ are the functions of the componentwise operator
$\lambda=-\Delta_{\parallel}$ fixed there. They act on each component as convolutions by
symmetric scalar kernels and commute with each other, with $P$ and with lateral translations.
$K_1$ is the kernel of $\mathsf{C}$, $C_N$ is the maximum of $\mathsf{C}^{-1}$ over the spectrum
of $\lambda$, $C'=\tfrac12(1+C_N)$, and
\begin{equation*}
N(w)=\tfrac12\langle Pw,\mathsf{C}^{-1}Pw\rangle+\tfrac12\|P^{\perp}w\|^2,
\qquad d_0(\beta)=-\mathsf{v}\beta .
\end{equation*}
$V_+(\gamma)$ is the least energy above a row of trace $\gamma$, attained by the free tail,
and $V_+(\gamma)=\tfrac12\langle\gamma,\mathsf{v}\gamma\rangle$
(Lemma~\ref{8.14:lem-innovations-identity}, a model statement). $V_1(\gamma)$ is the least
energy above a row of trace $\gamma$ among configurations whose first level above that row
lies in $\mathfrak{B}$, and $W=V_1-V_+$. Then the following hold.
\begin{enumerate}
\item[(a)] For every $\beta\in\mathfrak{H}_T$,
\begin{equation*}
V_1(\beta)=V_+(\beta)+\min_{E\in\mathfrak{B}}N(d_0(\beta)-E),
\quad\text{so}\quad W(\beta)=\min_{E\in\mathfrak{B}}N(d_0(\beta)-E),
\end{equation*}
and the minimum is attained at a unique $E_0\in\mathfrak{B}$. The \emph{multiplier}
\begin{equation*}
\mu:=\nabla N(d_0(\beta)-E_0)=\mathsf{C}^{-1}P(d_0(\beta)-E_0)-P^{\perp}E_0
\end{equation*}
satisfies $\langle\mu,E-E_0\rangle\le 0$ for all $E\in\mathfrak{B}$, that is, $\mu$ lies in
the normal cone of $\mathfrak{B}$ at $E_0$. For the box this gives complementarity:
$\mu_k(x)>0$ implies $E_{0,k}(x)=\tau_k$, and $\mu_k(x)<0$ implies $E_{0,k}(x)=-\tau_k$.
Moreover $P^{\perp}(E_0+\mu)=0$, so $\mathcal{D}:=E_0+\mu\in\mathfrak{H}_T$, and
\begin{equation}\label{8.14:eq-exit-value-2}
W(\beta)=\tfrac12\langle\mu,(\mathsf{C}P+P^{\perp})\mu\rangle;
\end{equation}
in particular
\begin{equation}\label{8.14:eq-exit-value-a}
\frac{\|\mu\|^2}{2C_N}\le W(\beta)\le\frac{\|\mu\|^2}{2}.
\end{equation}
\item[(b)] For $\beta\in\mathfrak{H}_T$, with $E_0,\mu,\mathcal{D}$ as in \textup{(a)} and
$\beta'=\beta+PE_0$, one has $d_0(\beta')=\mathsf{C}\mathcal{D}$. Let $a>a_0$ and
$E=E^{(a)}_{a-1}$. Then $E$ is the minimiser $E_0$ of \textup{(a)} for
$\beta=\beta^{(a)}_{a-1}$; with $\mu$ its multiplier, the exit pull $e_a:=E^{(a)}_a$ satisfies
$e_a=d_0(\beta^{(a)}_a)=\mathsf{C}(E+\mu)$, and the free tail of $X_a$ is
\begin{equation}\label{8.14:eq-exit-value-b}
E^{(a)}_{a+j}=\mathsf{C}^j e_a\qquad(j\ge 0).
\end{equation}
For every $f\in\mathfrak{H}$ and every convex $\varphi:\mathbb{R}\to\mathbb{R}$,
\begin{equation*}
\sum_{x,k}\varphi\big((\mathsf{C}f)_k(x)\big)\le\sum_{x,k}\varphi\big(f_k(x)\big).
\end{equation*}
\item[(c)] For every $a\ge a_0$,
\begin{equation}\label{8.14:eq-exit-value-c}
W\big(\beta^{(a+1)}_a\big)\le\Delta m(a)\le W\big(\beta^{(a)}_a\big).
\end{equation}
\item[(d)] Let $a\ge a_0$, $\beta=\beta^{(a+1)}_a$, $E_0=E^{(a+1)}_a$ (the minimiser of
\textup{(a)} for $\beta$, by \textup{(b)}), $\mu$ its multiplier, and
$\beta'=\beta+PE_0=\beta^{(a+1)}_{a+1}$. Then
\begin{equation}\label{8.14:eq-exit-value-3}
W(\beta')\le W(\beta)-\tfrac12\langle P^{\perp}\mu,(1-\mathsf{C}^2)P^{\perp}\mu\rangle,
\end{equation}
and, since $\Delta m(a+1)\le W(\beta')$ and $W(\beta)\le\Delta m(a)$ by \textup{(c)},
\begin{equation*}
\Delta m(a+1)\le\Delta m(a)-\tfrac12\langle P^{\perp}\mu,(1-\mathsf{C}^2)P^{\perp}\mu\rangle;
\end{equation*}
in particular $\Delta m$ is non-increasing on $a\ge a_0$.
\item[(e)] Let $\mathfrak{B}$ be the box, and let $a$, $\beta$, $E_0$, $\mu$, $\beta'$ be as
in \textup{(d)}. For $(x,k)$ with $\mu_k(x)\neq0$ put
\begin{equation*}
\kappa_1(x,k)=\sum_{x'}K_1(x-x')\big[\tau_k-\operatorname{sgn}(\mu_k(x))\,E_{0,k}(x')\big]
\in[0,2\tau_k].
\end{equation*}
Then, with the sums below taken over the pairs $(x,k)$ with $\mu_k(x)\ne0$,
\begin{align*}
W(\beta')&\le W(\beta)-\tfrac12\langle P^{\perp}\mu,\mathsf{C}\mathsf{v}P^{\perp}\mu\rangle
-\tfrac12\sum_{x,k}|\mu_k(x)|\min\Big(\kappa_1(x,k),\frac{|\mu_k(x)|}{2C'}\Big)\\
&\le W(\beta)-\tfrac12\sum_{x,k}|\mu_k(x)|\min\Big(\kappa_1(x,k),\frac{|\mu_k(x)|}{2C'}\Big),
\end{align*}
and by \textup{(c)} the same bound holds with $\Delta m(a+1)$ and $\Delta m(a)$ in place of
$W(\beta')$ and $W(\beta)$.
\item[(f)] Let $a>a_0$. Then $e_a\in\mathfrak{B}$ if and only if $\Delta m(a)=0$. In that case
$\Delta m(a')=0$ and $X_{a'}=X_a$ for every $a'\ge a$, and $E^{(a)}_y\in\mathfrak{B}$ for every
$y\ge a_0$. For the box, with $E=E^{(a)}_{a-1}$ as in \textup{(b)}, put
\begin{equation*}
\kappa_d=\min_{x,k,\pm}\sum_{x'}K_1(x-x')\big[\tau_k\mp E_k(x')\big];
\end{equation*}
then $e_a\in\mathfrak{B}$ as soon as $(2C_N\Delta m(a-1))^{1/2}\le\kappa_d$.
\end{enumerate}
\end{theorem}

\begin{proof}[Proof of Theorem~\ref{8.14:thm-exit-value}, clauses (a)--(c)]
We first collect facts of the model. By the model Poisson identities,
Lemma~\ref{8.14:lem-poisson-identities}, $\mathsf{C}=(1+\lambda+\mathsf{v})^{-1}$, so
$\mathsf{C}^{-1}=1+\lambda+\mathsf{v}$. Since $\lambda\ge0$ and
$(\lambda^2+4\lambda)^{1/2}\ge\lambda$, we have $\mathsf{v}\ge0$; hence $\mathsf{C}\le1$,
$\mathsf{C}^{-1}\ge1$, $C_N\ge1$, and
\begin{equation}\label{8.14:eq-exit-value-4}
C_N^{-1}\le\mathsf{C}\le 1 .
\end{equation}
As $P$ is an orthogonal projection commuting with the self-adjoint operator
$\mathsf{C}^{-1}$, we have $\langle Pw,\mathsf{C}^{-1}Pw\rangle=\langle w,\mathsf{C}^{-1}Pw\rangle$,
so that
\begin{equation}\label{8.14:eq-exit-value-5}
N(w)=\tfrac12\langle w,(\mathsf{C}^{-1}P+P^{\perp})w\rangle,
\qquad N(w)\ge\tfrac12\|Pw\|^2+\tfrac12\|P^{\perp}w\|^2=\tfrac12\|w\|^2 .
\end{equation}
Finally, for $\beta\in\mathfrak{H}_T=\operatorname{ran}P$, since $\mathsf{v}$ commutes with
$P$, we have $Pd_0(\beta)=d_0(\beta)$ and $P^{\perp}d_0(\beta)=0$.

In this proof, for a configuration $X$ (its traces and levels) and $r\ge 0$, we write
$\mathrm{en}_r(X)$ for the \emph{energy above row $r$} and $\mathrm{en}(X)$ for the
\emph{energy} of $X$:
\begin{equation}\label{8.14:eq-exit-value-1}
\mathrm{en}_r(X)=\sum_{y\ge r}\tfrac12\|E_y\|^2
+\sum_{y\ge r+1}\tfrac12\langle\beta_y,\lambda\beta_y\rangle,
\qquad \mathrm{en}(X)=\mathrm{en}_0(X).
\end{equation}

\emph{Clause (a).} Let $\beta\in\mathfrak{H}_T$ and consider a configuration above a row $r$
with trace $\beta$ at row $r$ and level $E_r=E\in\mathfrak{B}$; then $\beta_{r+1}=\beta'$ with
$\beta'=\beta+PE$, and by \eqref{8.14:eq-exit-value-1}
\begin{equation*}
\mathrm{en}_r=\tfrac12\|E\|^2+\tfrac12\langle\beta',\lambda\beta'\rangle+\mathrm{en}_{r+1}.
\end{equation*}
Since $\mathrm{en}_{r+1}\ge V_+(\beta')$, with equality when the configuration above row $r+1$
is the free tail from $\beta'$, the identity \eqref{8.14:eq-innovations-identity-a} of the
model innovations identity, Lemma~\ref{8.14:lem-innovations-identity}, gives
$\mathrm{en}_r\ge V_+(\beta)+N(d_0(\beta)-E)$, with equality in the same case. Minimising
first over the configuration above row $r+1$ and then over $E\in\mathfrak{B}$ gives
$V_1(\beta)=V_+(\beta)+\inf_{E\in\mathfrak{B}}N(d_0(\beta)-E)$. Since $\mathbb{T}$ is finite,
$\mathfrak{H}$ is finite-dimensional; by \eqref{8.14:eq-exit-value-5} the map
$E\mapsto N(d_0(\beta)-E)$ is a quadratic function whose quadratic part is positive definite,
hence strictly convex and tending to infinity at infinity; $\mathfrak{B}$ is non-empty, closed
and convex. So the infimum is a minimum, attained at a unique $E_0\in\mathfrak{B}$, and
$W(\beta)=V_1(\beta)-V_+(\beta)=N(d_0(\beta)-E_0)$.

By \eqref{8.14:eq-exit-value-5}, $\nabla N(w)=\mathsf{C}^{-1}Pw+P^{\perp}w$. At
$w=d_0(\beta)-E_0$, using $P^{\perp}d_0(\beta)=0$, this is
$\mu=\mathsf{C}^{-1}P(d_0(\beta)-E_0)-P^{\perp}E_0$. For $E\in\mathfrak{B}$ and
$t\in[0,1]$, the field $E_t=E_0+t(E-E_0)$ lies in $\mathfrak{B}$ by convexity, so
$g(t)=N(d_0(\beta)-E_t)$ satisfies $g(t)\ge g(0)$; $g$ is differentiable and
$g'(0)=-\langle\mu,E-E_0\rangle$, whence $\langle\mu,E-E_0\rangle\le0$.

For the box, suppose $\mu_k(x)>0$ and $E_{0,k}(x)<\tau_k$. The field $E$ equal to $E_0$
except that $E_k(x)=\tau_k$ lies in $\mathfrak{B}$, and
$\langle\mu,E-E_0\rangle=\mu_k(x)(\tau_k-E_{0,k}(x))>0$, a contradiction; so
$E_{0,k}(x)=\tau_k$. Likewise, if $\mu_k(x)<0$ and $E_{0,k}(x)>-\tau_k$, setting
$E_k(x)=-\tau_k$ gives $\langle\mu,E-E_0\rangle=\mu_k(x)(-\tau_k-E_{0,k}(x))>0$; so
$E_{0,k}(x)=-\tau_k$.

Since $P^{\perp}\mathsf{C}^{-1}P=\mathsf{C}^{-1}P^{\perp}P=0$, we get
$P^{\perp}\mu=-P^{\perp}E_0$, that is $P^{\perp}(E_0+\mu)=0$, and $\mathcal{D}=E_0+\mu$ lies
in $\operatorname{ran}P=\mathfrak{H}_T$. Further $P\mu=\mathsf{C}^{-1}P(d_0(\beta)-E_0)$, so
$P(d_0(\beta)-E_0)=\mathsf{C}P\mu$, while $P^{\perp}(d_0(\beta)-E_0)=-P^{\perp}E_0=P^{\perp}\mu$.
Therefore, $\mathsf{C}$ being self-adjoint and commuting with $P$,
\begin{align*}
W(\beta)=N(d_0(\beta)-E_0)
&=\tfrac12\langle\mathsf{C}P\mu,\mathsf{C}^{-1}\mathsf{C}P\mu\rangle+\tfrac12\|P^{\perp}\mu\|^2\\
&=\tfrac12\langle P\mu,\mathsf{C}P\mu\rangle+\tfrac12\|P^{\perp}\mu\|^2
=\tfrac12\langle\mu,(\mathsf{C}P+P^{\perp})\mu\rangle,
\end{align*}
which is \eqref{8.14:eq-exit-value-2}. By \eqref{8.14:eq-exit-value-4} and $C_N\ge1$, the
last expression lies between $(2C_N)^{-1}(\|P\mu\|^2+\|P^{\perp}\mu\|^2)$ and
$\tfrac12(\|P\mu\|^2+\|P^{\perp}\mu\|^2)$, and $\|P\mu\|^2+\|P^{\perp}\mu\|^2=\|\mu\|^2$; this
is \eqref{8.14:eq-exit-value-a}.

\emph{Clause (b).} Let $\beta'=\beta+PE_0$. By \eqref{8.14:eq-poisson-identities-d},
$\mathsf{v}=1-\mathsf{C}$, so $-\mathsf{v}PE_0=-PE_0+\mathsf{C}PE_0$ and
\begin{equation*}
d_0(\beta')=-\mathsf{v}\beta-\mathsf{v}PE_0=(d_0(\beta)-PE_0)+\mathsf{C}PE_0
=\mathsf{C}P\mu+\mathsf{C}PE_0=\mathsf{C}P\mathcal{D}=\mathsf{C}\mathcal{D},
\end{equation*}
where we used $d_0(\beta)-PE_0=P(d_0(\beta)-E_0)=\mathsf{C}P\mu$ (from clause (a) and
$Pd_0(\beta)=d_0(\beta)$) and $P\mathcal{D}=\mathcal{D}$.

Next we represent $P_a$ through the model innovations identity. For a configuration $X$ of
finite energy put
\begin{align*}
\Psi(X)&=\sum_{y<a_0}\tfrac12\|E_y\|^2+\sum_{1\le y\le a_0}\tfrac12\langle\beta_y,\lambda\beta_y\rangle,\\
\Phi_a(X)&=\Psi(X)+V_+(\beta_{a_0})+\sum_{a_0\le y<a}N(d_0(\beta_y)-E_y).
\end{align*}
Then $\mathrm{en}(X)=\Psi(X)+\mathrm{en}_{a_0}(X)$, and by
Lemma~\ref{8.14:lem-innovations-identity} (a model statement) the partial sums of
$V_+(\beta_{a_0})+\sum_{y\ge a_0}N(d_0(\beta_y)-E_y)$ increase to $\mathrm{en}_{a_0}(X)$; so
\begin{equation}\label{8.14:eq-exit-value-6}
\mathrm{en}(X)=\Phi_a(X)+\sum_{y\ge a}N(d_0(\beta_y)-E_y)\ge\Phi_a(X),
\end{equation}
and, by the same lemma applied with $r=a_0$ and $s=a$, equality holds when the levels of $X$
above row $a$ form the free tail from $\beta_a$. If $X$ is feasible for $P_a$, let
$\widetilde X$ agree with $X$ in the lower variables and in the levels $E_y$, $y<a$, and be
the free tail above row $a$. Since $a\ge a_0$, the lower constraints and the young caps of
$P_a$ read only these data, so $\widetilde X$ is feasible for $P_a$, and
$\mathrm{en}(\widetilde X)=\Phi_a(\widetilde X)=\Phi_a(X)$. Hence $\Phi_a(X)\ge m(a)$ for every
finite-energy configuration feasible for $P_a$. Applied to $X=X_a$ this gives
\begin{equation*}
m(a)=\mathrm{en}(X_a)\le\Phi_a(X_a)\le\mathrm{en}(X_a),
\end{equation*}
so that
\begin{equation}\label{8.14:eq-exit-value-7}
m(a)=\Phi_a(X_a)=\min\{\Phi_a(X):\ X\ \text{feasible for}\ P_a,\ \mathrm{en}(X)<\infty\},
\end{equation}
and by \eqref{8.14:eq-exit-value-6} the sum $\sum_{y\ge a}N(d_0(\beta^{(a)}_y)-E^{(a)}_y)$
vanishes. By \eqref{8.14:eq-exit-value-5} each term vanishes only at zero argument, so
$E^{(a)}_y=d_0(\beta^{(a)}_y)$ for all $y\ge a$. Then
\begin{equation*}
\beta^{(a)}_{y+1}=\beta^{(a)}_y+Pd_0(\beta^{(a)}_y)=(1-\mathsf{v})\beta^{(a)}_y
=\mathsf{C}\beta^{(a)}_y\qquad(y\ge a),
\end{equation*}
so $\beta^{(a)}_{a+j}=\mathsf{C}^j\beta^{(a)}_a$ and,
$\mathsf{v}$ and $\mathsf{C}$ commuting,
$E^{(a)}_{a+j}=-\mathsf{v}\mathsf{C}^j\beta^{(a)}_a=\mathsf{C}^jd_0(\beta^{(a)}_a)$. With
$e_a=E^{(a)}_a=d_0(\beta^{(a)}_a)$ this is \eqref{8.14:eq-exit-value-b}.

Now let $a>a_0$, and for $E'\in\mathfrak{B}$ let $X^{E'}$ agree with $X_a$ in the lower
variables and in the levels $E_y$, $y<a-1$, have level $a-1$ equal to $E'$, and be the free
tail above row $a$. The lower constraints read rows $y\le a_0\le a-1$ and levels below
$a_0$, which are unchanged, and the young caps of $P_a$ hold at levels $a_0,\dots,a-2$ as for
$X_a$ and at level $a-1$ because $E'\in\mathfrak{B}$; so $X^{E'}$ is feasible for $P_a$. The
rows up to $a-1$ of $X^{E'}$ are those of $X_a$, so, by the equality case
of \eqref{8.14:eq-exit-value-6},
\begin{equation*}
\mathrm{en}(X^{E'})=\Phi_a(X^{E'})=\Phi_{a-1}(X_a)+N\big(d_0(\beta^{(a)}_{a-1})-E'\big),
\end{equation*}
and by \eqref{8.14:eq-exit-value-7} the same formula with $E'=E^{(a)}_{a-1}$ gives $m(a)$.
Minimality of $m(a)$ therefore yields
\begin{equation*}
N\big(d_0(\beta^{(a)}_{a-1})-E^{(a)}_{a-1}\big)\le N\big(d_0(\beta^{(a)}_{a-1})-E'\big)
\qquad\text{for all } E'\in\mathfrak{B}.
\end{equation*}
By the uniqueness in clause (a), $E=E^{(a)}_{a-1}$ is the minimiser $E_0$
for $\beta=\beta^{(a)}_{a-1}$, and
\begin{equation}\label{8.14:eq-exit-value-8}
m(a)=\Phi_{a-1}(X_a)+W\big(\beta^{(a)}_{a-1}\big)\qquad(a>a_0).
\end{equation}
Since $\beta^{(a)}_a=\beta^{(a)}_{a-1}+PE$, the first identity of clause (b), applied with
$\beta=\beta^{(a)}_{a-1}$, gives $e_a=d_0(\beta^{(a)}_a)=\mathsf{C}(E+\mu)$.

For the convexity inequality: $\mathsf{C}$ acts on each component as convolution by $K_1$,
so $(\mathsf{C}f)_k(x)=\sum_{x'}K_1(x-x')f_k(x')$, and by
Lemma~\ref{8.14:lem-poisson-identities} (a model statement) $K_1\ge0$ and $\sum_zK_1(z)=1$;
hence $\sum_xK_1(x-x')=\sum_{x'}K_1(x-x')=1$. For convex $\varphi$, Jensen's inequality for the
probability weights $K_1(x-\cdot)$ gives
$\varphi((\mathsf{C}f)_k(x))\le\sum_{x'}K_1(x-x')\varphi(f_k(x'))$; summing over $x$ and using
$\sum_xK_1(x-x')=1$,
\begin{equation*}
\sum_{x,k}\varphi\big((\mathsf{C}f)_k(x)\big)
\le\sum_{x',k}\varphi\big(f_k(x')\big)\sum_xK_1(x-x')=\sum_{x',k}\varphi\big(f_k(x')\big).
\end{equation*}

\emph{Clause (c).} Let $a\ge a_0$. For the upper bound let $E^{*}\in\mathfrak{B}$ be the
minimiser of clause (a) for $\beta=\beta^{(a)}_a$, and let $X'$ agree with $X_a$ in the lower
variables and in the levels $E_y$, $y<a$, have level $a$ equal to $E^{*}$, and be the free tail
above row $a+1$. It satisfies the lower constraints and the young caps at levels
$a_0,\dots,a-1$ as $X_a$ does, and the cap at level $a$ since $E^{*}\in\mathfrak{B}$; so it is
feasible for $P_{a+1}$. Its rows up to $a$ are those of $X_a$, so by the equality case
of \eqref{8.14:eq-exit-value-6}, by \eqref{8.14:eq-exit-value-7} and by clause (a),
\begin{equation*}
m(a+1)\le\mathrm{en}(X')=\Phi_{a+1}(X')=\Phi_a(X_a)+N\big(d_0(\beta^{(a)}_a)-E^{*}\big)
=m(a)+W\big(\beta^{(a)}_a\big).
\end{equation*}
For the lower bound, $a+1>a_0$, so \eqref{8.14:eq-exit-value-8} with $a+1$ in place of $a$
gives $m(a+1)=\Phi_a(X_{a+1})+W(\beta^{(a+1)}_a)$: the top capped level of $P_{a+1}$ is chosen
to minimise $N(d_0(\beta^{(a+1)}_a)-\cdot\,)$ over $\mathfrak{B}$, and nothing after it costs
anything. The configuration $X_{a+1}$ has finite energy and satisfies the lower constraints and
the young caps at levels $a_0,\dots,a-1$, so it is feasible for $P_a$, and by
\eqref{8.14:eq-exit-value-7} $\Phi_a(X_{a+1})\ge m(a)$. Hence
$m(a+1)\ge m(a)+W(\beta^{(a+1)}_a)$. The two bounds together are
\eqref{8.14:eq-exit-value-c}.

This proves clauses (a)--(c); clauses (d)--(f) are derived from them in a separate proof.
\end{proof}

\begin{theorem}[Model: concavity, absorption and exactness]\label{8.14:prf-concavity-absorption}
The model is the capped-level variational problem of the model statements
Lemma~\ref{8.14:lem-poisson-identities}, Lemma~\ref{8.14:lem-innovations-identity} and
Theorem~\ref{8.14:thm-exit-value}. Its configurations are stacks of level fields $E_y$ in
$\mathfrak{H}=\ell^2(\mathbb{T};\mathbb{R}^3)$ on the lateral torus $\mathbb{T}$, with row fields in
$\mathfrak{H}_T$ related by $\beta_{y+1}=\beta_y+PE_y$. Here $a_0$ denotes the lowest young level,
so that $P_a$ is considered for $a\ge a_0$, as in the model theorem,
Theorem~\ref{8.14:thm-exit-value}. In the truncated problem $P_a$ the levels
$a_0\le y<a$ are constrained to the convex, closed, translation-invariant set $\mathfrak{B}$ of
young caps, and the levels from $a$ on are free. We write $m(a)$ for the value of $P_a$, $X_a$ for
its minimiser, $\beta^{(a)}_y$ and $E^{(a)}_y$ for the row and level fields of $X_a$,
$\Delta m(a)=m(a+1)-m(a)$, and $e_a=E^{(a)}_a=d_0(\beta^{(a)}_a)$ for the exit pull of $P_a$. The
operators $\mathsf{v}$, $\mathsf{C}=1-\mathsf{v}$, the kernel $K_1$ of $\mathsf{C}$, the form $N$,
the functional $W=V_1-V_+$, the constants $C_N$ and $C'=\tfrac12(1+C_N)$ and the multiplier of a
capped level are those of these three statements. By the \emph{box} we mean
\[
  \mathfrak{B}=\{E\in\mathfrak{H}:\ |E_k(x)|\le\tau_k\ \text{for all } x\in\mathbb{T},\ k=1,2,3\},
\]
where $\tau_1,\tau_2,\tau_3\ge0$ are the half-widths of the box. The lettering (d)--(f) continues
that of parts (a)--(c) of the model theorem, Theorem~\ref{8.14:thm-exit-value}.
\begin{itemize}
\item[(d)] \emph{Concavity.} Let $a\ge a_0$. Put $\beta:=\beta^{(a+1)}_a$ and $E_0:=E^{(a+1)}_a$,
and let $\mu$ be the multiplier of $E_0$. Then $\beta^{(a+1)}_{a+1}=\beta+PE_0$, and
\begin{equation}\label{8.14:eq-concavity-absorption-a}
\begin{split}
  \Delta m(a+1)\le W(\beta^{(a+1)}_{a+1})
  &\le W(\beta)-\tfrac12\langle P^{\perp}\mu,(1-\mathsf{C}^2)P^{\perp}\mu\rangle\\
  &\le \Delta m(a)-\tfrac12\langle P^{\perp}\mu,(1-\mathsf{C}^2)P^{\perp}\mu\rangle,
\end{split}
\end{equation}
where $1-\mathsf{C}^2=\mathsf{v}(1+\mathsf{C})\ge0$. In particular
$0\le\Delta m(a+1)\le\Delta m(a)$ for every $a\ge a_0$.
\item[(e)] \emph{Absorption.} Let $\mathfrak{B}$ be the box, and let $a$, $\beta$, $E_0$, $\mu$
be as in (d). At every $(x,k)$ with $\mu_k(x)\neq0$, put $s_k(x):=\operatorname{sign}\mu_k(x)$
and
\begin{equation}\label{8.14:eq-concavity-absorption-1}
  \kappa_1(x,k):=\sum_{x'\in\mathbb{T}}K_1(x-x')\bigl[\tau_k-s_k(x)E_{0,k}(x')\bigr].
\end{equation}
Then $\kappa_1(x,k)\in[0,2\tau_k]$, and
\begin{equation}\label{8.14:eq-concavity-absorption-b}
\begin{split}
  \Delta m(a+1)\le W(\beta^{(a+1)}_{a+1})
  &\le W(\beta)-\frac12\sum_{(x,k):\,\mu_k(x)\neq0}|\mu_k(x)|
    \min\Bigl\{\kappa_1(x,k),\frac{|\mu_k(x)|}{2C'}\Bigr\}\\
  &\le \Delta m(a)-\frac12\sum_{(x,k):\,\mu_k(x)\neq0}|\mu_k(x)|
    \min\Bigl\{\kappa_1(x,k),\frac{|\mu_k(x)|}{2C'}\Bigr\}.
\end{split}
\end{equation}
Both gains also hold at once, with $\mathsf{C}\mathsf{v}$ in place of the operator
$\mathsf{v}(1+\mathsf{C})$ of \eqref{8.14:eq-concavity-absorption-a}:
\begin{equation}\label{8.14:eq-concavity-absorption-2}
\begin{split}
  W(\beta^{(a+1)}_{a+1})\le W(\beta)
  &-\tfrac12\langle P^{\perp}\mu,\mathsf{C}\mathsf{v}P^{\perp}\mu\rangle\\
  &-\frac12\sum_{(x,k):\,\mu_k(x)\neq0}|\mu_k(x)|
    \min\Bigl\{\kappa_1(x,k),\frac{|\mu_k(x)|}{2C'}\Bigr\}.
\end{split}
\end{equation}
\item[(f)] \emph{Exactness.} Let $a\ge a_0$. Then $e_a\in\mathfrak{B}$ if and only if
$\Delta m(a)=0$. In that case $\Delta m(a')=0$, $m(a')=m(a)$ and $X_{a'}=X_a$ for every
$a'\ge a$; thus $X_a$ meets every young cap, and we say that the \emph{front has stopped} at $a$.
If $\mathfrak{B}$ is the box and $a>a_0$, let $E$ and $\mu$ now denote the field and the multiplier
of the top capped level $a-1$ of $P_a$. Then
\begin{equation}\label{8.14:eq-concavity-absorption-c}
  e_a\in\mathfrak{B}\quad\text{as soon as}\quad \bigl(2C_N\,\Delta m(a-1)\bigr)^{1/2}\le\kappa_d,
\end{equation}
where
\begin{equation}\label{8.14:eq-concavity-absorption-7}
  \kappa_d:=\min_{x,k,\pm}\ \sum_{x'\in\mathbb{T}}K_1(x-x')\bigl[\tau_k\mp E_k(x')\bigr].
\end{equation}
\end{itemize}
\end{theorem}

\begin{proof}
Throughout we use the following facts.
\begin{itemize}
\item The operators $\mathsf{v}$ and $\mathsf{C}$ are real functions of the self-adjoint operator
$\lambda$. Hence they are self-adjoint, and they commute with $P$ and $P^{\perp}$.
\item $\mathsf{v}\ge0$, because $\mathsf{v}$ is the non-negative root of
$\mathsf{v}^2+\lambda\mathsf{v}=\lambda$ in the model lemma,
Lemma~\ref{8.14:lem-poisson-identities}.
\item By the same model lemma, $\mathsf{C}^{-1}=1+\lambda+\mathsf{v}\ge1$. Consequently
$0<\mathsf{C}\le1$, $C_N\ge1$, and $N(w)\ge0$ for every $w$.
\item By the kernel representation of $\mathsf{C}$ in the same model lemma,
$(\mathsf{C}f)_k(x)=\sum_{x'}K_1(x-x')f_k(x')$ with
$K_1\ge0$ and $\sum_{x'}K_1(x-x')=1$.
\end{itemize}
The following consequences of the model theorem, Theorem~\ref{8.14:thm-exit-value}, are used
repeatedly.
\begin{itemize}
\item By its part (a) (display \eqref{8.14:eq-exit-value-a}), for every row field
$\beta\in\mathfrak{H}_T$, $W(\beta)=\min_{E\in\mathfrak{B}}N(d_0(\beta)-E)\ge0$.
\item By its part (c) (display \eqref{8.14:eq-exit-value-c}), for every $a\ge a_0$,
\begin{equation}\label{8.14:eq-concavity-absorption-3}
  W(\beta^{(a+1)}_a)\le\Delta m(a)\le W(\beta^{(a)}_a).
\end{equation}
\end{itemize}

\emph{Proof of (d).} Fix $a\ge a_0$. The top capped level of $P_{a+1}$ is $a$, and $a+1>a_0$. By
part (b) of the model theorem, Theorem~\ref{8.14:thm-exit-value}, the tail of $X_{a+1}$ above
row $a$ is $V_1$-optimal given $\beta$. Hence $E_0$ is the unique minimiser of
$N(d_0(\beta)-\cdot\,)$ over $\mathfrak{B}$, and $W(\beta)=N(d_0(\beta)-E_0)$. By part (a) of the
same model theorem, the multiplier is
\[
  \mu=\mathsf{C}^{-1}P(d_0(\beta)-E_0)-P^{\perp}E_0,\qquad P^{\perp}(E_0+\mu)=0.
\]
The row relation gives $\beta':=\beta+PE_0=\beta^{(a+1)}_{a+1}$.

First, $\mathsf{C}E_0\in\mathfrak{B}$. For $z\in\mathbb{T}$ write $E_0(\cdot-z)$ for the lateral
translate $x\mapsto E_0(x-z)$ of $E_0$. Then
\[
  \mathsf{C}E_0=\sum_{z\in\mathbb{T}}K_1(z)\,E_0(\cdot-z) .
\]
This is a convex combination of lateral translates of $E_0$, since $K_1\ge0$ and $\sum_zK_1(z)=1$.
Each translate $E_0(\cdot-z)$ lies in $\mathfrak{B}$, because $\mathfrak{B}$ is translation
invariant, and $\mathfrak{B}$ is convex and closed. Hence $\mathsf{C}E_0\in\mathfrak{B}$. Since
$W(\beta')=\min_{E\in\mathfrak{B}}N(d_0(\beta')-E)$, this gives
\[
  W(\beta')\le N\bigl(d_0(\beta')-\mathsf{C}E_0\bigr).
\]

Next we compute the two components of $d_0(\beta')-\mathsf{C}E_0$. Since $\beta\in\mathfrak{H}_T$
and $\mathsf{v}$ commutes with $P$, we have $d_0(\beta)=-\mathsf{v}\beta\in\mathfrak{H}_T$ and
\[
  d_0(\beta')=-\mathsf{v}\beta-\mathsf{v}PE_0=d_0(\beta)-\mathsf{v}PE_0 .
\]
Both terms are transverse. Using the identity $\mathsf{v}+\mathsf{C}=1$ of the model lemma,
Lemma~\ref{8.14:lem-poisson-identities} (display \eqref{8.14:eq-poisson-identities-d}), and then
$P^{\perp}E_0=-P^{\perp}\mu$, we obtain
\[
  P\bigl(d_0(\beta')-\mathsf{C}E_0\bigr)=d_0(\beta)-\mathsf{v}PE_0-\mathsf{C}PE_0=d_0(\beta)-PE_0,
\]
\[
  P^{\perp}\bigl(d_0(\beta')-\mathsf{C}E_0\bigr)=-\mathsf{C}P^{\perp}E_0=\mathsf{C}P^{\perp}\mu .
\]
By the definition of $N$, therefore,
\begin{equation}\label{8.14:eq-concavity-absorption-4}
  N\bigl(d_0(\beta')-\mathsf{C}E_0\bigr)
  =\tfrac12\bigl\langle d_0(\beta)-PE_0,\mathsf{C}^{-1}(d_0(\beta)-PE_0)\bigr\rangle
  +\tfrac12\|\mathsf{C}P^{\perp}\mu\|^2 .
\end{equation}

On the other hand, $P(d_0(\beta)-E_0)=d_0(\beta)-PE_0$ and
$P^{\perp}(d_0(\beta)-E_0)=-P^{\perp}E_0=P^{\perp}\mu$. Hence
\[
  W(\beta)=N(d_0(\beta)-E_0)
  =\tfrac12\bigl\langle d_0(\beta)-PE_0,\mathsf{C}^{-1}(d_0(\beta)-PE_0)\bigr\rangle
  +\tfrac12\|P^{\perp}\mu\|^2 .
\]
The first term on the right of \eqref{8.14:eq-concavity-absorption-4} thus equals
$W(\beta)-\tfrac12\|P^{\perp}\mu\|^2$. Since $\mathsf{C}$ is self-adjoint and commutes with
$P^{\perp}$,
\[
  \tfrac12\|\mathsf{C}P^{\perp}\mu\|^2=\tfrac12\langle P^{\perp}\mu,\mathsf{C}^2P^{\perp}\mu\rangle .
\]
Consequently
\begin{equation}\label{8.14:eq-concavity-absorption-5}
\begin{split}
  W(\beta')\le N\bigl(d_0(\beta')-\mathsf{C}E_0\bigr)
  &=W(\beta)-\tfrac12\|P^{\perp}\mu\|^2+\tfrac12\|\mathsf{C}P^{\perp}\mu\|^2\\
  &=W(\beta)-\tfrac12\langle P^{\perp}\mu,(1-\mathsf{C}^2)P^{\perp}\mu\rangle .
\end{split}
\end{equation}

Now apply \eqref{8.14:eq-concavity-absorption-3}. Its upper member at $a+1$ reads
$\Delta m(a+1)\le W(\beta^{(a+1)}_{a+1})=W(\beta')$. Its lower member at $a$ reads
$W(\beta)=W(\beta^{(a+1)}_a)\le\Delta m(a)$. Together with \eqref{8.14:eq-concavity-absorption-5}
these give \eqref{8.14:eq-concavity-absorption-a}. By the identity
\eqref{8.14:eq-poisson-identities-f} of the model lemma, Lemma~\ref{8.14:lem-poisson-identities},
$1-\mathsf{C}^2=\mathsf{v}(1+\mathsf{C})\in[\mathsf{v},2\mathsf{v}]$, which is non-negative because
$\mathsf{v}\ge0$. Finally, the lower member of \eqref{8.14:eq-concavity-absorption-3} at $a+1$
gives $\Delta m(a+1)\ge W(\beta^{(a+2)}_{a+1})\ge0$. Hence $0\le\Delta m(a+1)\le\Delta m(a)$.

\emph{Proof of (e).} Let $\mathfrak{B}$ be the box, and fix $(x,k)$ with $\mu_k(x)\neq0$; write
$s=s_k(x)$. By the complementarity in part (a) of the model theorem,
Theorem~\ref{8.14:thm-exit-value} ($\mu_k(x)>0$ implies $E_{0,k}(x)=\tau_k$, and $\mu_k(x)<0$
implies $E_{0,k}(x)=-\tau_k$), we have
$E_{0,k}(x)=s\tau_k$. Since $s^2=1$ and $\sum_{x'}K_1(x-x')=1$,
\[
  (\mathsf{C}E_0)_k(x)=s\sum_{x'}K_1(x-x')\,sE_{0,k}(x')=s\bigl[\tau_k-\kappa_1(x,k)\bigr],
\]
with $\kappa_1(x,k)$ as in \eqref{8.14:eq-concavity-absorption-1}. Since $|E_{0,k}(x')|\le\tau_k$,
each bracket in \eqref{8.14:eq-concavity-absorption-1} lies in $[0,2\tau_k]$. As $K_1\ge0$ with
total mass $1$, it follows that $\kappa_1(x,k)\in[0,2\tau_k]$. Define $\eta\in\mathfrak{H}$ by
\[
  \eta_k(x):=s_k(x)\min\Bigl\{\kappa_1(x,k),\frac{|\mu_k(x)|}{2C'}\Bigr\}
  \ \text{ if }\mu_k(x)\neq0,\qquad \eta_k(x):=0\ \text{ otherwise},
\]
and put $E':=\mathsf{C}E_0+\eta$.

Then $E'\in\mathfrak{B}$. Where $\mu_k(x)\neq0$,
\[
  s_k(x)E'_k(x)=\tau_k-\kappa_1(x,k)+\min\Bigl\{\kappa_1(x,k),\frac{|\mu_k(x)|}{2C'}\Bigr\}
  \in\bigl[\tau_k-\kappa_1(x,k),\tau_k\bigr]\subset[-\tau_k,\tau_k],
\]
because $\kappa_1(x,k)\le2\tau_k$. Elsewhere $E'_k(x)=(\mathsf{C}E_0)_k(x)$, and
$\mathsf{C}E_0\in\mathfrak{B}$ by the proof of (d).

Put $\bar w:=d_0(\beta')-\mathsf{C}E_0$, so that $d_0(\beta')-E'=\bar w-\eta$. The form $N$ is
quadratic, $N(w)=\tfrac12\langle w,(\mathsf{C}^{-1}P+P^{\perp})w\rangle$, with gradient
$\nabla N(w)=\mathsf{C}^{-1}Pw+P^{\perp}w$. By the componentwise computation in the proof of (d),
and since $P\mu=\mathsf{C}^{-1}(d_0(\beta)-PE_0)$ by the formula for $\mu$,
\[
  \nabla N(\bar w)=\mathsf{C}^{-1}\bigl(d_0(\beta)-PE_0\bigr)+\mathsf{C}P^{\perp}\mu
  =P\mu+\mathsf{C}P^{\perp}\mu=\mu-(1-\mathsf{C})P^{\perp}\mu .
\]
Exactly, because $N$ is quadratic,
\[
  N(\bar w-\eta)=N(\bar w)-\langle\nabla N(\bar w),\eta\rangle+N(\eta)
  =N(\bar w)-\langle\mu,\eta\rangle+\langle(1-\mathsf{C})P^{\perp}\mu,\eta\rangle+N(\eta).
\]
We bound the last three terms in turn.
\begin{itemize}
\item Since $\eta_k(x)$ has the sign of $\mu_k(x)$ where $\mu_k(x)\neq0$ and vanishes elsewhere,
$\langle\mu,\eta\rangle=\sum_{x,k}|\mu_k(x)||\eta_k(x)|$.
\item The operator $1-\mathsf{C}=\mathsf{v}$ is positive semidefinite and satisfies
$1-\mathsf{C}\le1$. The Cauchy--Schwarz inequality for the form it defines therefore gives
\[
  \langle(1-\mathsf{C})P^{\perp}\mu,\eta\rangle
  \le\tfrac12\langle P^{\perp}\mu,(1-\mathsf{C})P^{\perp}\mu\rangle
  +\tfrac12\langle\eta,(1-\mathsf{C})\eta\rangle
  \le\tfrac12\langle P^{\perp}\mu,(1-\mathsf{C})P^{\perp}\mu\rangle+\tfrac12\|\eta\|^2 .
\]
\item Since $\mathsf{C}^{-1}\le C_N$ and $C_N\ge1$,
\[
  N(\eta)\le\tfrac12C_N\|P\eta\|^2+\tfrac12\|P^{\perp}\eta\|^2\le\tfrac12C_N\|\eta\|^2 .
\]
\end{itemize}
Since $\tfrac12+\tfrac12C_N=C'$, the second and third bounds together contribute
$C'\|\eta\|^2=\sum_{x,k}C'\eta_k(x)^2$. Inserting the value of $N(\bar w)$ from
\eqref{8.14:eq-concavity-absorption-5} and using $E'\in\mathfrak{B}$, we obtain
\begin{equation}\label{8.14:eq-concavity-absorption-6}
\begin{split}
  W(\beta')\le N(\bar w-\eta)\le W(\beta)
  &-\tfrac12\bigl\langle P^{\perp}\mu,\bigl[(1-\mathsf{C}^2)-(1-\mathsf{C})\bigr]P^{\perp}\mu\bigr\rangle\\
  &-\sum_{x,k}\bigl[|\mu_k(x)||\eta_k(x)|-C'\eta_k(x)^2\bigr].
\end{split}
\end{equation}

The bracket $(1-\mathsf{C}^2)-(1-\mathsf{C})=\mathsf{C}(1-\mathsf{C})=\mathsf{C}\mathsf{v}$ is
non-negative: $\mathsf{C}$ and $\mathsf{v}$ are commuting non-negative operators, and
$1-\mathsf{C}\le1-\mathsf{C}^2$ by \eqref{8.14:eq-poisson-identities-f} of the model lemma,
Lemma~\ref{8.14:lem-poisson-identities}. As for the
sum, $|\eta_k(x)|\le|\mu_k(x)|/(2C')$ gives $C'\eta_k(x)^2\le\tfrac12|\mu_k(x)||\eta_k(x)|$, and
hence
\[
  |\mu_k(x)||\eta_k(x)|-C'\eta_k(x)^2\ge\tfrac12|\mu_k(x)|
  \min\Bigl\{\kappa_1(x,k),\frac{|\mu_k(x)|}{2C'}\Bigr\}
\]
where $\mu_k(x)\neq0$, while both sides vanish elsewhere. Keeping the bracket in
\eqref{8.14:eq-concavity-absorption-6} gives \eqref{8.14:eq-concavity-absorption-2}. Dropping the
non-negative bracket gives the second member of \eqref{8.14:eq-concavity-absorption-b}. The first
and last members of \eqref{8.14:eq-concavity-absorption-b} follow from
\eqref{8.14:eq-concavity-absorption-3}, exactly as in the proof of (d): the upper member at $a+1$
bounds $\Delta m(a+1)$ by $W(\beta')$, and the lower member at $a$ bounds $W(\beta)$ by
$\Delta m(a)$.

\emph{Proof of (f).} Fix $a\ge a_0$. Suppose first that $e_a\in\mathfrak{B}$. The choice
$E:=e_a=d_0(\beta^{(a)}_a)$ in $W(\beta^{(a)}_a)=\min_{E\in\mathfrak{B}}N(d_0(\beta^{(a)}_a)-E)$
gives $N(0)=0$. Since $W\ge0$, it follows that $W(\beta^{(a)}_a)=0$. By
\eqref{8.14:eq-concavity-absorption-3},
\[
  0\le W(\beta^{(a+1)}_a)\le\Delta m(a)\le W(\beta^{(a)}_a)=0,
\]
so $\Delta m(a)=0$.

Conversely, suppose $\Delta m(a)=0$. Every $P_{a'}$, $a'\ge a_0$, minimises the same objective,
the energy of the configuration, which by the model lemma,
Lemma~\ref{8.14:lem-innovations-identity}, equals the lower part of the energy plus
$V_+(\beta_{a_0})+\sum_{y\ge a_0}N(d_0(\beta_y)-E_y)$; the problems differ only in the set
$a_0\le y<a'$ of levels constrained to $\mathfrak{B}$. This objective is strictly convex and
coercive in the level variables, so the minimiser of each $P_{a'}$ is unique in them. In
particular $P_{a+1}$ carries, in addition to the caps of $P_a$, the young cap at level $a$. Hence
$X_{a+1}$ is admissible for $P_a$, and its objective value is $m(a+1)=m(a)$, so $X_{a+1}$ is
optimal for $P_a$. By the uniqueness just stated, $X_{a+1}=X_a$ in the level variables. In
particular $e_a=E^{(a)}_a=E^{(a+1)}_a$. This field lies in $\mathfrak{B}$, since level $a$ is capped
in $P_{a+1}$. So $e_a\in\mathfrak{B}$.

Now let $\Delta m(a)=0$. By (d), $0\le\Delta m(a')\le\Delta m(a)$ for all $a'\ge a$. Hence
$\Delta m(a')=0$ for all $a'\ge a$, and $m(a')=m(a)$. For $a'\ge a$, $P_{a'}$ has the same
objective as $P_a$ and more caps. So $X_{a'}$ is admissible for $P_a$ with value $m(a')=m(a)$,
hence optimal for $P_a$, and therefore $X_{a'}=X_a$ in the level variables by uniqueness. Each
$X_{a'}$ meets the caps at levels $a_0\le y<a'$. Consequently $X_a$ meets every young cap: the
front has stopped.

For the box criterion let $a>a_0$, and let $E$, $\mu$ be the field and the multiplier of the top
capped level $a-1$ of $P_a$. By part (b) of the model theorem, Theorem~\ref{8.14:thm-exit-value}
(display \eqref{8.14:eq-exit-value-b} and the line before it), $e_a=\mathsf{C}(E+\mu)$. For each
sign $\pm$, each $x$ and each $k$,
\[
  \pm(\mathsf{C}E)_k(x)=\sum_{x'}K_1(x-x')\bigl(\pm E_k(x')\bigr)
  =\tau_k-\sum_{x'}K_1(x-x')\bigl[\tau_k\mp E_k(x')\bigr],
\]
and $\pm(\mathsf{C}\mu)_k(x)\le(\mathsf{C}|\mu_k|)(x)$ because $K_1\ge0$. Hence
\[
  \pm(e_a)_k(x)=\pm(\mathsf{C}E)_k(x)\pm(\mathsf{C}\mu)_k(x)
  \le\tau_k-\sum_{x'}K_1(x-x')\bigl[\tau_k\mp E_k(x')\bigr]+(\mathsf{C}|\mu_k|)(x).
\]

It remains to bound the last term. Since $(\mathsf{C}|\mu_k|)(x)$ is a convex combination of
values of $|\mu_k|$, we have $(\mathsf{C}|\mu_k|)(x)\le\|\mu\|_\infty\le\|\mu\|_2$. The multiplier
$\mu$ belongs to the $V_1$-optimal level given $\beta^{(a)}_{a-1}$, by part (b) of
Theorem~\ref{8.14:thm-exit-value}. The lower bound $\|\mu\|^2/(2C_N)\le W$ of part (a) of that
model theorem therefore gives $\|\mu\|_2\le(2C_NW(\beta^{(a)}_{a-1}))^{1/2}$. The lower member of
\eqref{8.14:eq-concavity-absorption-3} at $a-1\ge a_0$ gives
$W(\beta^{(a)}_{a-1})\le\Delta m(a-1)$. Altogether
\[
  (\mathsf{C}|\mu_k|)(x)\le\|\mu\|_\infty\le\|\mu\|_2
  \le\bigl(2C_N\,W(\beta^{(a)}_{a-1})\bigr)^{1/2}\le\bigl(2C_N\,\Delta m(a-1)\bigr)^{1/2}.
\]
If $(2C_N\Delta m(a-1))^{1/2}\le\kappa_d$, then for every $x$, $k$ and sign $\pm$,
\[
  \pm(e_a)_k(x)\le\tau_k-\kappa_d+\bigl(2C_N\,\Delta m(a-1)\bigr)^{1/2}\le\tau_k .
\]
Hence $|(e_a)_k(x)|\le\tau_k$ for all $x$ and $k$, that is, $e_a\in\mathfrak{B}$. This is
\eqref{8.14:eq-concavity-absorption-c}.
\end{proof}

\begin{corollary}[Model: monotone increments and a stopping height]\label{8.14:cor-stopping-height}
The model is the one of the model exit-value theorem, Theorem~\ref{8.14:thm-exit-value}.
Its ingredients are the following.
\begin{itemize}
\item The level fields form $\mathfrak{H}=\ell^2(\mathbb{T};\mathbb{R}^3)$ on the lateral torus
$\mathbb{T}$. On $\mathfrak{H}$ act the transverse projection $P$ and its complement $P^{\perp}$, and
the one-step Poisson operator $\mathsf{C}=1-\mathsf{v}$, whose kernel $K_1$ is non-negative with unit
sums. We put $C_N=\max\mathsf{C}^{-1}$ and $C'=\tfrac12(1+C_N)$.
\item The quadratic defect form is $N$, the value functions are $V_+$, $V_1$ and $W=V_1-V_+$, and the
free pull is $d_0(\beta)=-\mathsf{v}\beta$.
\item The set of young caps $\mathfrak{B}$ is convex, closed and translation-invariant.
\item For $a\ge a_0$, with $a_0$ the base level fixed in Theorem~\ref{8.14:thm-exit-value}, there are
truncated problems $P_a$, with value $m(a)$, minimiser $X_a$ and
increment $\Delta m(a)=m(a+1)-m(a)$.
\end{itemize}
As there, $\beta^{(a)}_{a'}$ is the row-$a'$ argument of the value functions along $X_a$, and
$E^{(a)}_{a'}$ is the field of $X_a$ at level $a'$. For $a>a_0$, $\mu^{(a)}$ is the multiplier of the
top capped level $a-1$ of $P_a$. We put
\begin{equation*}
u_a:=W\bigl(\beta^{(a)}_{a-1}\bigr)\qquad(a>a_0).
\end{equation*}
\begin{itemize}
\item[(E1)] For every $a>a_0$,
\begin{equation}\label{8.14:eq-stopping-height-a}
\begin{gathered}
m(a)-m(a_0)\ \ge\ (a-a_0)\,\Delta m(a-1)\ \ge\ (a-a_0)\,\Delta m(a),\\
u_a\ \le\ \Delta m(a-1)\ \le\ \frac{m(a)-m(a_0)}{a-a_0}.
\end{gathered}
\end{equation}
\item[(E2)] Let $a>a_0$ and $\mu:=\mu^{(a)}$. Then
\begin{equation}\label{8.14:eq-stopping-height-1}
\begin{gathered}
u_a=\tfrac12\bigl\langle\mu,(\mathsf{C}P+P^{\perp})\mu\bigr\rangle,\qquad
\frac{\|\mu\|^2}{2C_N}\le u_a\le\frac{\|\mu\|^2}{2},\\
u_{a+1}\ \le\ \Delta m(a)\ \le\ W\bigl(\beta^{(a)}_a\bigr).
\end{gathered}
\end{equation}
If moreover $\mathfrak{B}$ is the box $\{E:\ |E_k(x)|\le\tau_k\ \text{for all }x,k\}$, with the cap
heights $\tau_k>0$ of Theorem~\ref{8.14:thm-exit-value}, and, for $\mu_k(x)\neq0$, $\kappa_1(x,k)$ is
the one-step Poisson deficit of $E^{(a)}_{a-1}$, defined as in the model statement on concavity,
absorption and exactness, then
\begin{equation}\label{8.14:eq-stopping-height-b}
\begin{split}
u_{a+1}\ &\le\ \Delta m(a)\ \le\ W\bigl(\beta^{(a)}_a\bigr)\\
&\le\ u_a-\frac12\sum_{x,k}|\mu_k(x)|\,\min\Bigl(\kappa_1(x,k),\frac{|\mu_k(x)|}{2C'}\Bigr).
\end{split}
\end{equation}
Keep $\mathfrak{B}$ the box. Let $\kappa_d>0$ and $s_\infty>0$; the constant $\kappa_d$ plays the role
of a uniform lower bound, on the range of levels considered, for the minimal one-step deficit of the
top capped field that enters the box criterion of display \eqref{8.14:eq-concavity-absorption-c}. Put
$\lambda_*:=\min(\kappa_d/s_\infty,\,1/(2C'))$.
\begin{itemize}
\item[(a)] Suppose that, at a level $a>a_0$, the minimisers satisfy the one-sided conditions
$\kappa_1(x,k)\ge\kappa_d$ whenever $\mu^{(a)}_k(x)\neq0$, and $\|\mu^{(a)}\|_\infty\le s_\infty$.
Then
\begin{equation}\label{8.14:eq-stopping-height-2}
u_{a+1}\le(1-\lambda_*)\,u_a .
\end{equation}
\item[(b)] Let $a_1>a_0$, put
\begin{equation*}
\mathfrak{n}:=\log^+\bigl(2C_Nu_{a_1}/\kappa_d^2\bigr),
\end{equation*}
and let $j$ be the least integer with $j\ge\lambda_*^{-1}\mathfrak{n}$, so that
$j\le\lambda_*^{-1}\mathfrak{n}+1$. Suppose the conditions of (a) hold at the levels
$a_1,\dots,a_1+j-1$. Suppose also that
\begin{equation*}
\kappa_d\le\min_{x,k,\pm}\sum_{x'}K_1(x-x')\bigl[\tau_k\mp E_k(x')\bigr],
\end{equation*}
where $E$ is the top capped field of $P_{a_1+j+1}$. Then the front has stopped at level $a_1+j+1$,
in the sense of the model statement
on concavity, absorption and exactness: $\Delta m$ vanishes from that level on, and
$X_{a'}=X_{a_1+j+1}$ for all $a'\ge a_1+j+1$.
\end{itemize}
\end{itemize}
The conditions in the second part of (E2) are hypotheses on the minimisers, which are not known
explicitly. They delimit
the range of levels on which \eqref{8.14:eq-stopping-height-2} and the stopping height are asserted.
\end{corollary}

\begin{proof}
Every appeal below is to the model exit-value theorem, Theorem~\ref{8.14:thm-exit-value}, or to the
model statement on concavity, absorption and exactness. Both are statements about the model.

\emph{(E1).} Let $a'\ge a_0$, and let $\beta:=\beta^{(a'+1)}_{a'}$. By display
\eqref{8.14:eq-exit-value-b}, the $V_1$-optimal step from $\beta$ is $E_0:=E^{(a'+1)}_{a'}$; let $\mu$
denote its multiplier, and put $\beta':=\beta+PE_0=\beta^{(a'+1)}_{a'+1}$. Apply the concavity
inequality \eqref{8.14:eq-concavity-absorption-a} to this $\beta$. By display
\eqref{8.14:eq-exit-value-c} at $a'+1$ (upper bound)
and at $a'$ (lower bound), it gives
\begin{equation*}
\Delta m(a'+1)\le W(\beta')\le W(\beta)\le\Delta m(a').
\end{equation*}
The middle inequality holds because $\langle P^{\perp}\mu,(1-\mathsf{C}^2)P^{\perp}\mu\rangle\ge0$:
the operators $\lambda=-\Delta_{\parallel}\ge0$ and
$\mathsf{v}=\tfrac12[-\lambda+(\lambda^2+4\lambda)^{1/2}]\ge0$ commute, so
$\mathsf{C}=(1+\lambda+\mathsf{v})^{-1}$ satisfies $0\le\mathsf{C}\le1$ and $1-\mathsf{C}^2$ is
positive semi-definite.
Thus $\Delta m$ is non-increasing on $a'\ge a_0$. Telescoping the definition of $\Delta m$ gives
\begin{equation*}
m(a)-m(a_0)=\sum_{a_0\le a'<a}\Delta m(a').
\end{equation*}
Each summand has $a'\le a-1$, so it is at least $\Delta m(a-1)$, which in turn is at least
$\Delta m(a)$. This proves the first line of \eqref{8.14:eq-stopping-height-a}; dividing by
$a-a_0>0$ gives the last inequality of the second line. Finally, the lower member of
\eqref{8.14:eq-exit-value-c} at $a-1\ge a_0$ reads $W(\beta^{(a)}_{a-1})\le\Delta m(a-1)$, that is,
$u_a\le\Delta m(a-1)$.

\emph{(E2).} Let $a>a_0$ and $\beta:=\beta^{(a)}_{a-1}$. By display \eqref{8.14:eq-exit-value-b}, the
tail of $P_a$ above row $a-1$ is $V_1$-optimal given $\beta$. Hence the minimiser of
$N(d_0(\beta)-\cdot)$ over $\mathfrak{B}$ is the top capped field $E^{(a)}_{a-1}$ of $P_a$, and its
multiplier is $\mu=\mu^{(a)}$. Display \eqref{8.14:eq-exit-value-a} then gives
$W(\beta)=\frac12\langle\mu,(\mathsf{C}P+P^{\perp})\mu\rangle$ and
$\|\mu\|^2/(2C_N)\le W(\beta)\le\|\mu\|^2/2$. This is the first line of
\eqref{8.14:eq-stopping-height-1}.

We now prove the three inequalities of \eqref{8.14:eq-stopping-height-b}; the first two are also the
second line of \eqref{8.14:eq-stopping-height-1} and do not use the box.
\begin{itemize}
\item The first is the lower member of \eqref{8.14:eq-exit-value-c} at $a$, namely
$\Delta m(a)\ge W(\beta^{(a+1)}_a)=u_{a+1}$.
\item The second is the upper member of \eqref{8.14:eq-exit-value-c} at $a$.
\item For the third, which uses the box, apply the absorption inequality
\eqref{8.14:eq-concavity-absorption-b} for the
box to $\beta=\beta^{(a)}_{a-1}$. Its $V_1$-optimal step is $E^{(a)}_{a-1}$, with multiplier $\mu$, as
just shown, and the next row is $\beta'=\beta^{(a)}_a$. The inequality gives
\begin{equation*}
W\bigl(\beta^{(a)}_a\bigr)\le
W(\beta)-\frac12\sum_{x,k}|\mu_k(x)|\min\Bigl(\kappa_1(x,k),\frac{|\mu_k(x)|}{2C'}\Bigr),
\end{equation*}
with $\kappa_1$ the deficit of $E^{(a)}_{a-1}$, and $W(\beta)=u_a$.
\end{itemize}

\emph{(E2)(a).} Let $\mu_k(x)\neq0$. Then $|\mu_k(x)|\le s_\infty$ and $\kappa_1(x,k)\ge\kappa_d$. Hence
$\kappa_1(x,k)\ge|\mu_k(x)|\,\kappa_d/s_\infty$, and therefore
\begin{equation*}
\min\Bigl(\kappa_1(x,k),\frac{|\mu_k(x)|}{2C'}\Bigr)\ \ge\
|\mu_k(x)|\min\Bigl(\frac{\kappa_d}{s_\infty},\frac1{2C'}\Bigr)=\lambda_*|\mu_k(x)|.
\end{equation*}
Where $\mu_k(x)=0$, the corresponding summand of \eqref{8.14:eq-stopping-height-b} vanishes. So the
subtracted sum in \eqref{8.14:eq-stopping-height-b} is at least $\frac12\lambda_*\|\mu\|^2$. By the
upper bound in \eqref{8.14:eq-stopping-height-1}, this is at least $\lambda_*u_a$. Therefore
\eqref{8.14:eq-stopping-height-b} gives $u_{a+1}\le u_a-\lambda_*u_a$.

\emph{(E2)(b).} Since $\lambda_*\le1/(2C')=1/(1+C_N)<1$, the factor $1-\lambda_*$ lies in $[0,1)$.
Moreover $u_a\ge0$, because $N\ge0$. Iterating \eqref{8.14:eq-stopping-height-2} over the levels
$a_1,\dots,a_1+j-1$, and using $1-\lambda_*\le e^{-\lambda_*}$, we obtain
\begin{equation*}
u_{a_1+j}\le(1-\lambda_*)^j\,u_{a_1}\le e^{-\lambda_*j}\,u_{a_1}.
\end{equation*}
Two cases arise.
\begin{itemize}
\item If $\mathfrak{n}=0$, then $j=0$ and $2C_Nu_{a_1}\le\kappa_d^2$.
\item If $\mathfrak{n}>0$, then $e^{-\lambda_*j}\le e^{-\mathfrak{n}}=\kappa_d^2/(2C_Nu_{a_1})$.
\end{itemize}
In both cases $2C_Nu_{a-1}\le\kappa_d^2$ with $a:=a_1+j+1$.

Since $a-1>a_0$, \eqref{8.14:eq-stopping-height-b} at level $a-1$ gives
$\Delta m(a-1)\le W(\beta^{(a-1)}_{a-1})\le u_{a-1}$, the subtracted sum being non-negative.
Hence
\begin{equation*}
\bigl(2C_N\Delta m(a-1)\bigr)^{1/2}\le\kappa_d.
\end{equation*}
By the displayed hypothesis of (b), the minimal one-step deficit of the top capped field of $P_a$ is
at least $\kappa_d$. The
box criterion of the exactness display \eqref{8.14:eq-concavity-absorption-c} therefore gives that the
exit pull $e_a$ lies in $\mathfrak{B}$. By the same display, $\Delta m$ then vanishes from $a$ on, and
$X_{a'}=X_a$ for every $a'\ge a$. That is, the front has stopped at level $a=a_1+j+1$.
\end{proof}

\begin{lemma}[Model: exit pull from exit-row curvature]\label{8.14:lem-exit-pull}
The model is that of Lemma~\ref{8.14:lem-poisson-identities} and Theorem~\ref{8.14:thm-exit-value}.
Fields are arranged in rows $y$ over the lateral torus $\mathbb{T}$. The level $E_y\in\mathfrak{H}$
between row $y$ and row $y+1$ carries the trace $\beta_y\in\mathfrak{H}_T$ of row $y$ to the trace
$\beta_{y+1}=\beta_y+PE_y$ of row $y+1$. The truncated problem $P_a$ restricts the levels
$a_0\le y<a$ (the \emph{capped} levels) to the convex set $\mathfrak{B}$ and leaves the levels above
row $a$ free; here $a_0$ is the row at which the capped levels begin, and the constraints on the
lower variables involve rows $\le a_0$ only. Its minimiser is $X_a$, and
$\beta^{(a)}_y$ denotes the trace of $X_a$ on row $y$. The operators $\lambda=-\Delta_{\parallel}$,
$\mathsf{v}$, $\mathsf{C}=1-\mathsf{v}$ and $\mathsf{U}=\lambda+\mathsf{v}$ are those of
Lemma~\ref{8.14:lem-poisson-identities}.

Let $a>a_0$ and put $\beta:=\beta^{(a)}_a$, the trace of $X_a$ on row $a$, its exit row. Let
$B_a:=\operatorname{curl}\beta$ be the tangential
curvature of $X_a$ on its exit row, which is gauge-invariant. Let $(E,\mu)$ be the field and the
multiplier of the level $a-1$, and let $e_a$ be the exit pull of $P_a$. Then the following hold.
\begin{enumerate}
\item[(i)] The exit pull satisfies
\begin{equation}\label{8.14:eq-exit-pull-1}
e_a=-\mathsf{v}\beta=-(\lambda+\mathsf{v})^{-1}\operatorname{curl}^*B_a .
\end{equation}
\item[(ii)] The multiplier satisfies
\begin{equation}\label{8.14:eq-exit-pull-2}
\mu=e_a-E-\operatorname{curl}^*B_a .
\end{equation}
This is the stationarity of the energy in the tangential bonds of row $a$.
\item[(iii)] The sum of field and multiplier satisfies
\begin{equation}\label{8.14:eq-exit-pull-3}
E+\mu=e_a-\operatorname{curl}^*B_a\in\mathfrak{H}_T .
\end{equation}
The relation \eqref{8.14:eq-exit-value-b} of the model theorem, Theorem~\ref{8.14:thm-exit-value},
then reads $\mathsf{v}e_a=-\mathsf{C}\operatorname{curl}^*B_a$. This is consistent with (i) by
$\mathsf{C}(\lambda+\mathsf{v})=\mathsf{v}$, which is \eqref{8.14:eq-poisson-identities-c} of the
model lemma, Lemma~\ref{8.14:lem-poisson-identities}.
\end{enumerate}
By (i), the exit pull is the order-zero Riesz image of the exit row's tangential curvature, namely
its image under the operator $-(\lambda+\mathsf{v})^{-1}\operatorname{curl}^*$. By (ii), the
multiplier is local in $(e_a,E,B_a)$.
\end{lemma}

\begin{proof}
\emph{(i).} By the model theorem, Theorem~\ref{8.14:thm-exit-value}, in its part
\eqref{8.14:eq-exit-value-b}, the exit pull is the free pull of the exit-row trace. Hence
\[
e_a=d_0(\beta)=-\mathsf{v}\beta .
\]
The trace $\beta$ lies in $\mathfrak{H}_T$. On $\mathfrak{H}_T$ the operator $\lambda$ acts as
$\operatorname{curl}^*\operatorname{curl}$, so
\[
\lambda\beta=\operatorname{curl}^*\operatorname{curl}\beta=\operatorname{curl}^*B_a .
\]
The operator $\lambda$ is invertible on $\mathfrak{H}_T$. On such data the model lemma,
Lemma~\ref{8.14:lem-poisson-identities}, gives
$\mathsf{v}\lambda^{-1}=(\lambda+\mathsf{v})^{-1}$ through \eqref{8.14:eq-poisson-identities-c}.
Therefore
\[
e_a=-\mathsf{v}\lambda^{-1}\lambda\beta=-\mathsf{v}\lambda^{-1}\operatorname{curl}^*B_a
=-(\lambda+\mathsf{v})^{-1}\operatorname{curl}^*B_a ,
\]
which is \eqref{8.14:eq-exit-pull-1}.

\emph{(ii).} Put $\beta_-:=\beta^{(a)}_{a-1}$, the trace of $X_a$ on row $a-1$. Since $E$ is the
level $a-1$, we have $\beta=\beta_-+PE$. The level $a-1$ is the top capped level of $P_a$. By
\eqref{8.14:eq-exit-value-a} of the model theorem, Theorem~\ref{8.14:thm-exit-value}, applied to this
level with trace $\beta_-$ and field $E$, the multiplier is
\[
\mu=\mathsf{C}^{-1}P\bigl(d_0(\beta_-)-E\bigr)-P^{\perp}E .
\]
We have $d_0(\beta_-)=-\mathsf{v}\beta_-$, and by the model lemma,
Lemma~\ref{8.14:lem-poisson-identities}, $\mathsf{C}^{-1}=1+\lambda+\mathsf{v}=1+\mathsf{U}$.
This gives the first equality of
\begin{align*}
\mathsf{C}^{-1}(-\mathsf{v}\beta_--PE)
&=-(1+\lambda+\mathsf{v})\mathsf{v}\beta_--(1+\mathsf{U})PE\\
&=-\mathsf{U}\beta_--PE-\mathsf{U}PE\\
&=-\mathsf{U}\beta-PE .
\end{align*}
The second equality holds because
$(1+\lambda+\mathsf{v})\mathsf{v}=\mathsf{v}+\lambda\mathsf{v}+\mathsf{v}^2$. This equals
$\mathsf{v}+\lambda=\mathsf{U}$ by the identity $\mathsf{v}^2+\lambda\mathsf{v}=\lambda$, which is
\eqref{8.14:eq-poisson-identities-a}. The third equality uses $\beta_-+PE=\beta$.

Insert this into the expression for $\mu$ and use $P+P^{\perp}=1$. Then
\[
\mu=-\lambda\beta-\mathsf{v}\beta-PE-P^{\perp}E=-\lambda\beta-\mathsf{v}\beta-E .
\]
By the proof of (i), $\lambda\beta=\operatorname{curl}^*B_a$ and $-\mathsf{v}\beta=e_a$. Hence
$\mu=-\operatorname{curl}^*B_a+e_a-E$, which is \eqref{8.14:eq-exit-pull-2}.

\emph{(iii).} Adding $E$ to both sides of \eqref{8.14:eq-exit-pull-2} gives the identity in
\eqref{8.14:eq-exit-pull-3}. The right-hand side lies in $\mathfrak{H}_T$, since
$e_a=-\mathsf{v}\beta$ and $\operatorname{curl}^*B_a=\lambda\beta$ with $\beta\in\mathfrak{H}_T$,
and $\mathsf{v}$ and $\lambda$ preserve $\mathfrak{H}_T$.

The relation \eqref{8.14:eq-exit-value-b} of the model theorem, Theorem~\ref{8.14:thm-exit-value},
is $e_a=\mathsf{C}(E+\mu)$. By
\eqref{8.14:eq-exit-pull-3} it becomes
\[
e_a=\mathsf{C}e_a-\mathsf{C}\operatorname{curl}^*B_a .
\]
Since $1-\mathsf{C}=\mathsf{v}$, this is $\mathsf{v}e_a=-\mathsf{C}\operatorname{curl}^*B_a$.

It remains to compare with (i). Apply $\mathsf{v}$ to \eqref{8.14:eq-exit-pull-1}. This gives
\[
\mathsf{v}e_a=-\mathsf{v}(\lambda+\mathsf{v})^{-1}\operatorname{curl}^*B_a .
\]
The operators $\mathsf{C}$, $\mathsf{v}$ and $\lambda+\mathsf{v}$ are functions of $\lambda$ and
commute. The identity $\mathsf{C}(\lambda+\mathsf{v})=\mathsf{v}$ of the model lemma,
Lemma~\ref{8.14:lem-poisson-identities}, in the form \eqref{8.14:eq-poisson-identities-c}, gives
$\mathsf{v}(\lambda+\mathsf{v})^{-1}=\mathsf{C}$. So the right-hand side equals
$-\mathsf{C}\operatorname{curl}^*B_a$, in agreement with the form of \eqref{8.14:eq-exit-value-b}
just obtained.
\end{proof}

\begin{remark}[Model: the lateral cutoff in the stopping height]\label{8.14:rem-cutoff-in-height}
The model is the exit-value model of Lemma~\ref{8.14:lem-poisson-identities} and
Corollary~\ref{8.14:cor-stopping-height}. Its level fields lie in
$\mathfrak{H}=\ell^2(\mathbb{T};\mathbb{R}^3)$ on the lateral torus $\mathbb{T}=(\mathbb{Z}/N)^3$.
Here $N$ is the side of the lateral torus of the model, not the rank of the gauge group.
$P$ and $P^{\perp}$ are the projections onto transverse fields and onto their complement, and
$\lambda=-\Delta_{\parallel}$. The free-value operator is
$\mathsf{v}=\tfrac12[-\lambda+(\lambda^2+4\lambda)^{1/2}]$ and the one-step Poisson operator is
$\mathsf{C}=1-\mathsf{v}$. We write $C_N=\max\mathsf{C}^{-1}$, $C'=\tfrac12(1+C_N)$, and
$\mathfrak{B}$ for the set of young caps. The truncated problem is $P_a$, with value $m(a)$,
minimiser $X_a$, increment $\Delta m(a)$, top multiplier $\mu$ and exit pull $e_a$. The
one-step Poisson deficit is $\kappa_1$, with minimum $\kappa_d$ on the range considered, and
the top-level defect is $u_a$.

Every statement of the model invoked below is a model statement. In this remark $N$ plays the
role of the lateral cutoff. We ask how the stopping height of
Corollary~\ref{8.14:cor-stopping-height} depends on it.

The model is read in the lattice problem through the following dictionary:
\begin{equation*}
N=2P_{\mathrm{bl}}\,n,\qquad W_y=\varpi_y\,n,\qquad
\tau_y=\theta'\mathsf{T},\qquad \ell=s_e\theta'.
\end{equation*}
Here $n$ is the cutoff variable (in this remark $n$ is not the order of a Schwinger function);
$2P_{\mathrm{bl}}$ is the ratio of the lateral side $N$ to $n$;
$W_y$ is the band height, with $\varpi_y$ free of $n$; $\tau_y$ is the cap and $\ell$ the unit
of the lift bound of the model theorem on inactivity at full band width (below, the model
inactivity theorem); and $\theta'$ and
$s_e$ are the two constants of the dictionary
through which $\tau_y$ and $\ell$ are expressed. The last two relations are taken in the unit in
which the dictionary measures field values, so that $\tau_y=\ell\mathsf{T}/s_e$.
The level of the lattice problem at which the model is read is $k_o$, and $n=L^{k_o}$, so that
at that level
\begin{equation*}
\log N=\log(2P_{\mathrm{bl}})+k_o\log L .
\end{equation*}
A quantity is called \emph{cutoff-free} if it is free of $n$ when read in this dictionary.
Two kinds of quantity are allowed in a cutoff-free bound:
\begin{itemize}
\item ratios such as $N/W_y=2P_{\mathrm{bl}}/\varpi_y$;
\item floors of the kind recorded in the corollary of the model theorem on inactivity at full
band width, in which $n$ cancels against $N$.
\end{itemize}
A bare $\log N=\log(2P_{\mathrm{bl}})+k_o\log L$ is not cutoff-free.

The quantity at issue is the following. We consider a range of levels on which
$\kappa_1\ge\kappa_d$ on the support of $\mu$ and $\|\mu\|_\infty\le s_\infty$. On such a range
the geometric decay \eqref{8.14:eq-stopping-height-b} of the model corollary,
Corollary~\ref{8.14:cor-stopping-height}, combined with the stopping criterion
\eqref{8.14:eq-concavity-absorption-c} of the model statement on concavity, absorption and
exactness, stops the front within $H_{\mathrm{stop}}$ levels of the first level $a_1$ of the
range, where the exit-value stopping height and its factors are
\begin{equation}\label{8.14:eq-cutoff-in-height-a}
\begin{gathered}
H_{\mathrm{stop}}:=\mathfrak{h}\,\mathfrak{n}+1,\qquad
\mathfrak{h}:=\lambda_*^{-1}=\max\Bigl(\frac{s_\infty}{\kappa_d},\,2C'\Bigr),\\
\mathfrak{n}:=\log^{+}\Bigl(\frac{2C_Nu_{a_1}}{\kappa_d^2}\Bigr).
\end{gathered}
\end{equation}
The equality in the definition of $\mathfrak{h}$ is the definition
$\lambda_*=\min(\kappa_d/s_\infty,1/(2C'))$ of the decay rate, inverted.

Read in the dictionary, ``cutoff-uniform'' means that $H_{\mathrm{stop}}/W_y$ is bounded by a
cutoff-free quantity. The standard is set by the model theorem on inactivity at full band
width. Its bound on short segments, which we write $\bar a$ (this letter keeps it apart from
the base level $a_0$ of the truncated problems used below), with $b_y$ the floor of that bound,
gives
\begin{equation*}
\frac{\bar a}{W_y}\le
\frac{\max\bigl(b_y,\;N e^{-(\tau_y/(3\ell)-C_K^0-6)/6}\bigr)}{\varpi_y\,n},
\end{equation*}
and $n$ cancels against $N=2P_{\mathrm{bl}}n$, as the corollary of that theorem records.

The deficit $\kappa_d$ does not carry the cutoff. What follows in this paragraph is a
description of sizes, not a proved statement. Suppose the field $E_k$ of a top level $k$ is
saturated on a lateral ball of radius $R$, that is $E_k=\varsigma\tau_k$ there, where $\tau_k$ is
the cap of level $k$ and $\varsigma$ is the factor of the model by which a saturated field is a
multiple of its cap, and that it is unsaturated by a fixed fraction
outside it. Then $\kappa_1$ at the centre of the ball is $\tau_k$ times the $K_1$-mass outside
the ball. The lattice harmonic measure from height one has tails $K_1(z)\asymp|z|^{-4}$ in three
lateral dimensions. The $K_1$-mass outside the ball is therefore $\asymp 1/R$, and
$\kappa_1\asymp\tau/R$, with $\tau$ the size of the caps.

A heuristic reading of the model, itself a description, proceeds in four steps:
\begin{itemize}
\item it takes $\kappa_d\asymp\tau/R_*$, for a radius $R_*\asymp W_ye^{-c/c_B}$;
\item it takes $s_\infty\lesssim c_B\tau\log N$, where $c_B$ is the constant of a weak
tangential bound $\|B_a\|_\infty\le c_B\tau$ and $c$ is a constant;
\item it takes $u_{a_1}\lesssim n_{\mathrm{hot}}s_\infty^2$, where $n_{\mathrm{hot}}$ is the
number of hot fine sites of a level, that is of site--component pairs in the support of the
top multiplier;
\item it infers a stopping height $\asymp R_*(s_\infty/\tau)[\log n_{\mathrm{hot}}
+\log(R_*s_\infty/\tau)]$, which against the band height reads
$e^{c/c_B}\gtrsim(\log N)^2$.
\end{itemize}
From $\kappa_d\asymp\tau/R$ one gets $\mathfrak{h}\asymp R\,(s_\infty/\tau)$. The $e$-folding
height is thus proportional to the lateral radius of the saturated patch. This is the
scale-covariant dependence: the free field above a patch of radius $R$ decays over a height
$\asymp R$. In the setting of the model inactivity theorem, the saturated set under a kink is a
half-disc, and $\bar a$ is comparable to its radius.

We measure this against $W_y$ with $R=\varrho_*n$, where $\varrho_*$ is cutoff-free, for
instance $R_*\asymp W_ye^{-c/c_B}$ above. Then
\begin{equation*}
\frac{\mathfrak{h}}{W_y}\asymp\frac{\varrho_*}{\varpi_y}\cdot\frac{s_\infty}{\tau},
\end{equation*}
which is cutoff-free exactly when $s_\infty/\tau$ is. So the cutoff does not sit in the dyadic
structure of the deficit.

The supremum of the top multiplier is the next point. Let $a>a_0$, where $a_0$ is the base
level of the truncated problems of Corollary~\ref{8.14:cor-stopping-height}, let $B_a$ be the
curl of the exit-row trace, and let $E$ be the field of level $a-1$. The field $E$ lies in the
set $\mathfrak{B}$ of young caps; writing $\tau_{\max}$ for the largest cap, we have
$\|E\|_\infty\le\tau_{\max}$.

By the model lemma, Lemma~\ref{8.14:lem-exit-pull}, part (ii), we have
$\mu=e_a-E-\operatorname{curl}^*B_a$. Hence, pointwise on $\mathbb{T}$,
$|\mu|\le\tau_{\max}+|\operatorname{curl}^*B_a|+|e_a|$.

By part (i) of the same lemma, $e_a=-(\lambda+\mathsf{v})^{-1}\operatorname{curl}^*B_a$. The
operator $(\lambda+\mathsf{v})^{-1}\operatorname{curl}^*$ commutes with lateral translations, so
it is convolution on $\mathbb{T}$ with a kernel. Bounding a convolution in the supremum norm by
the $\ell^1(\mathbb{T})$-norm of the kernel times the supremum of the datum, and taking for
$s_\infty$ the supremum of $|\mu|$, we obtain
\begin{equation}\label{8.14:eq-cutoff-in-height-b}
s_\infty\le\tau_{\max}+\|\operatorname{curl}^*B_a\|_\infty+\|e_a\|_\infty,\qquad
\|e_a\|_\infty\le\mathfrak{k}_N\,\|B_a\|_\infty .
\end{equation}
Here $\mathfrak{k}_N$ is the $\ell^1(\mathbb{T})$-norm of the convolution kernel of
$(\lambda+\mathsf{v})^{-1}\operatorname{curl}^*$.

The symbol of this operator is homogeneous of degree zero at small lateral momenta, since
$\mathsf{v}\sim\lambda^{1/2}$ there. It is therefore a lattice operator of Riesz type. Such
kernels have $\ell^1$-norm $\asymp 1+\log N$ on $(\mathbb{Z}/N)^3$. The status of the two sides
differs:
\begin{itemize}
\item The upper bound $\mathfrak{k}_N\le C(1+\log N)$, with $C$ a constant (in this remark $C$
denotes such a constant, not the same at each occurrence), is of the kind produced by the
lattice Fourier analysis on which the theorem giving the strip profiles rests. It
is not derived here, and no proved statement of this book depends on it.
\item The lower bound, $\gtrsim\log N$ at data of lattice scale, is not proved here for
$(\mathbb{Z}/N)^3$. For the two-dimensional strip, the bound on the curvature of the free
continuation into a hole has a kernel whose $\ell^1$-norm is of the order of the logarithm of
the strip's lateral size, from above and from below.
\end{itemize}

Against data that are smooth at a lateral scale $\sigma$, the oddness of the kernel leaves only
$1+\log(N/\sigma)$; this is the mechanism of the lemma, in the setting of the model inactivity
theorem, on data smooth at a given scale. But
$B_a$ is the exit row
of a truncated problem. It lies on the free boundary, where no smoothing scale is available, and
the multipliers just below it are objects of lattice scale. A lattice kink is an admissible
datum: the cutoff-free corollary of the model inactivity theorem allows the sharp lattice kink.

So, with a weak tangential bound $\|B_a\|_\infty\le c_B\tau$,
\eqref{8.14:eq-cutoff-in-height-b} gives
\begin{equation*}
\frac{s_\infty}{\tau}\lesssim 1+c_B\,(C+\mathfrak{k}_N)\asymp 1+c_B\log N .
\end{equation*}
This is the first logarithm of the heuristic reading above, and its exact source is the Riesz
image in Lemma~\ref{8.14:lem-exit-pull}\,(i).

This logarithm is present at a front even in the scalar model $\widehat{M}$. There, part (iii)
of the theorem on $\widehat{M}$ bounds the normal field at the levels $m=0,1,\dots$ steps above
the last clamped level; at $m=0$, one level above the last clamped level, the bound is
$[C_K^0+6\log N]\ell$, a bare $\log N$, which is not cutoff-free. The theorem on
$\widehat{M}$ notes that its bound has this logarithmic shape at $m=0,1,\dots$, as it must.

Part (ii) of the model lemma, Lemma~\ref{8.14:lem-exit-pull}, reads, in the letters of
$\widehat{M}$ and at what that model calls an up-active top pair,
\begin{equation*}
\mu(Y-1)=F_Y-F_{Y-1}+\Delta_x\hat\psi=F_Y-\tau_y+\Delta_x\hat\psi .
\end{equation*}
Here $Y$ is the top row of $\widehat{M}$, $F_Y$ its exit field, $F_{Y-1}$ the field one row
below, equal there to the cap $\tau_y$, $\Delta_x$ the lateral Laplacian and $\hat\psi$ the lift
function of $\widehat{M}$. The identity shows that the top multiplier of the problem truncated
there is as large as the exit field
$F_Y$ allows. At a kink of the datum, $F_Y$ is of the size of the kernel's $\ell^1$-norm times
$\ell$, that is $\asymp\ell\log N$, by the two-sided bound on the curvature of the free
continuation into a hole of the two-dimensional strip. We
use that two-sided size only as an indication that the logarithm at distance zero is attained,
not as a proved lower bound in the setting of the model inactivity theorem.

Granting that size, $F_Y$ exceeds $\tau_y=\ell\mathsf{T}/s_e$ at every level $k_o$ with
$\log N\gtrsim\mathsf{T}/s_e$, so that the top multiplier is unbounded in $k_o$, while the
model inactivity theorem stops the front with a cutoff-free bound. The only a priori bounds on
$s_\infty$ available to
\eqref{8.14:eq-stopping-height-b} are thus of two kinds:
\begin{itemize}
\item bounds that carry $\log N$, obtained through the model lemma,
Lemma~\ref{8.14:lem-exit-pull}\,(i);
\item extensive bounds, namely $\|\mu\|_\infty\le(2C_Nu_a)^{1/2}$. These come from
\eqref{8.14:eq-exit-value-a}, through the lower bound $u_a\ge\|\mu\|_2^2/(2C_N)$ of
\eqref{8.14:eq-stopping-height-b} and $\|\mu\|_\infty\le\|\mu\|_2$.
\end{itemize}

The count of $e$-folds is the last point. The number
$\mathfrak{n}=\log^{+}(2C_Nu_{a_1}/\kappa_d^2)$
compares two quantities of different nature. On one side stands $u_{a_1}$, which can be bounded
in two ways:
\begin{itemize}
\item By \eqref{8.14:eq-exit-value-a} and \eqref{8.14:eq-stopping-height-b}, $u_{a_1}$ is
comparable to $\|\mu^{(a_1)}\|_2^2$. This is a sum over the support of $\mu^{(a_1)}$, the hot
fine sites of a whole level. Hence $u_{a_1}\le\tfrac12n_{\mathrm{hot}}s_\infty^2$, where
$n_{\mathrm{hot}}\le 3N^3$ is the number of site--component pairs in that support.
\item By \eqref{8.14:eq-stopping-height-a}, $u_{a_1}\le\mathcal{V}/(a_1-a_0)$, where, for any
integer $h_0\ge a_1-a_0$, since $m$ is non-decreasing ($\Delta m(a)\ge u_{a+1}\ge0$ by
\eqref{8.14:eq-stopping-height-b}),
\begin{equation*}
\mathcal{V}:=m(a_1)-m(a_0)\le\mathcal{V}_{h_0}:=m(a_0+h_0)-m(a_0).
\end{equation*}
The latter is an extensive value budget.
\end{itemize}
On the other side stands $\kappa_d^2$, the threshold of one site.

In the dictionary we use the following sizes, all of which are descriptions:
\begin{itemize}
\item $\mathcal{V}_{h_0}\le\mathfrak{e}_{h_0}$, with $\mathfrak{e}_{h_0}$ of the kind given by
the energy-budget proposition of the auxiliary setting of the model inactivity theorem, namely
$\mathfrak{e}_{h_0}\asymp N^3\,b\,\tau_y^2(s_e\bar\lambda/\mathsf{T})^2$, where $b$ is the number
of levels entering that proposition and $\bar\lambda$ is of the order of the quadratic mean,
over these $b$ levels, of the strip profile entering it;
\item $b=\beta_{\mathrm{hot}}n$, where
$\beta_{\mathrm{hot}}\asymp(H_{\mathrm{b}}/n)e^{-c_{\mathrm{hot}}\mathsf{T}}$ is cutoff-free,
$H_{\mathrm{b}}$ being the height entering the energy-budget proposition and $c_{\mathrm{hot}}$
a constant;
\item $a_1-a_0=\varpi''n$, with $\varpi''$ cutoff-free;
\item $\kappa_d=c_\kappa\tau_y/R$, with $R=\varrho_*n$ and $c_\kappa$ a constant.
\end{itemize}
Inserting them gives
\begin{equation}\label{8.14:eq-cutoff-in-height-c}
\begin{gathered}
\frac{2C_Nu_{a_1}}{\kappa_d^2}\lesssim
\Bigl[\frac{96\,C_NP_{\mathrm{bl}}^3\beta_{\mathrm{hot}}\varrho_*^2
(s_e\bar\lambda/\mathsf{T})^2}{c_\kappa^2\,\varpi''}\Bigr]\,n^5,\\
\text{so that}\qquad
\mathfrak{n}=5k_o\log L+\log(\text{a cutoff-free quantity}).
\end{gathered}
\end{equation}

This is the limitation of every value or energy device that is expressed by the lemma of the
exit analysis according to which a value budget resolves a pointwise event only up to the number
of fine sites. There, a single fine plaquette carrying the value $2\theta_y$, twice the
threshold $\theta_y$ of the exit analysis, costs an amount smaller than the value budget by a
factor at least the number of fine sites. What the exit-value model adds
to the other devices of the exit analysis is the geometric decay
\eqref{8.14:eq-stopping-height-b}. That
decay converts the power loss of that lemma into a logarithm, and the logarithm is a
logarithm of the cutoff.

A bookkeeping that does not use $s_\infty$ confirms this, and we prove it. Assume
$\kappa_1\ge\kappa_d$ on the support of $\mu=\mu^{(a)}$. The model corollary,
Corollary~\ref{8.14:cor-stopping-height},
applies the absorption inequality \eqref{8.14:eq-concavity-absorption-b} to
$\beta=\beta^{(a)}_{a-1}$, with $\beta'=\beta^{(a)}_a$, and gives
\begin{equation*}
u_{a+1}\le\Delta m(a)\le W(\beta^{(a)}_a)
\le u_a-\tfrac12\sum|\mu|\min\Bigl(\kappa_1,\frac{|\mu|}{2C'}\Bigr).
\end{equation*}
We split the sum according to whether $|\mu|\ge 2C'\kappa_d$:
\begin{itemize}
\item where $|\mu|\ge2C'\kappa_d$, the minimum is at least $\kappa_d$;
\item where $|\mu|<2C'\kappa_d$, the minimum is at least $\min(\kappa_d,|\mu|/(2C'))=|\mu|/(2C')$.
\end{itemize}
Let $A$ and $B$ be the sums of $\mu^2$ over the first and the second set. On the first set we
use $\sum_{|\mu|\ge2C'\kappa_d}|\mu|\ge A^{1/2}$. The absorbed amount is therefore at least
$\tfrac12\kappa_dA^{1/2}+B/(4C')$.

By \eqref{8.14:eq-stopping-height-b}, $A+B=\|\mu\|_2^2\ge 2u_a$. Two cases are possible:
\begin{itemize}
\item if $A\ge u_a$, the first term is at least $\tfrac12\kappa_du_a^{1/2}$;
\item if $A<u_a$, then $B>u_a$, and the second term is at least $u_a/(4C')$.
\end{itemize}
Hence
\begin{equation}\label{8.14:eq-cutoff-in-height-1}
u_{a+1}\le u_a-\min\Bigl(\tfrac12\kappa_du_a^{1/2},\;\frac{u_a}{4C'}\Bigr).
\end{equation}

While $u_a^{1/2}\ge2C'\kappa_d$, the minimum in \eqref{8.14:eq-cutoff-in-height-1} is
$\tfrac12\kappa_du_a^{1/2}$. Since $2C'\ge1$, we have $u_a^{1/2}\ge\kappa_d/4$, and
\begin{equation*}
u_{a+1}\le u_a-\tfrac12\kappa_du_a^{1/2}\le\bigl(u_a^{1/2}-\kappa_d/4\bigr)^2 .
\end{equation*}
So $u^{1/2}$ decreases by $\kappa_d/4$ per level, and this phase lasts at most
$4u_{a_1}^{1/2}/\kappa_d$ levels.

Once $u_a^{1/2}\le2C'\kappa_d$, the minimum is $u_a/(4C')$, and
$u_{a+1}\le(1-1/(4C'))u_a$. The decay is now geometric, with a rate that depends only on $C'$.
It brings $u$ from $(2C'\kappa_d)^2$ below the threshold $\kappa_d^2/(2C_N)$ of
\eqref{8.14:eq-concavity-absorption-c} in at most $j_0$ levels. Here $j_0$ is the least
integer with
\begin{equation*}
j_0\ge\frac{\log(8C'^2C_N)}{-\log(1-1/(4C'))},
\end{equation*}
a number that depends on $C'$ and $C_N$ only, and not on $\kappa_d$ or on the data. The stopping
height is therefore at most
$4u_{a_1}^{1/2}/\kappa_d+j_0$.

This bound contains no logarithm and no $s_\infty$. But
$4u_{a_1}^{1/2}/\kappa_d=4(u_{a_1}/\kappa_d^2)^{1/2}$ is proportional to $n^{5/2}$ by
\eqref{8.14:eq-cutoff-in-height-c}: the power loss described above appears in it openly.

The square of $\log N$ in the heuristic reading is therefore the product of the two factors
\eqref{8.14:eq-cutoff-in-height-b} and \eqref{8.14:eq-cutoff-in-height-c}:
\begin{equation*}
\mathfrak{h}\,\mathfrak{n}\asymp R\,(1+c_B\log N)\,(5k_o\log L+C)
\qquad\text{against}\qquad W_y\propto R ,
\end{equation*}
and both factors contain $k_o\log L$. The factor
$\log n_{\mathrm{hot}}+\log(R_*s_\infty/\tau)$ of that reading is $\mathfrak{n}$ written with
the first bound $u_{a_1}\le\tfrac12n_{\mathrm{hot}}s_\infty^2$ and with $\kappa_d\asymp\tau/R_*$.
\end{remark}

\begin{lemma}[Model: the one-row ceiling]\label{8.14:lem-one-row-ceiling}
The model is that of the model theorem, Theorem~\ref{8.14:thm-exit-value} (one capped level
followed by free levels, the caps forming a box). Its level
fields are $\mathfrak{H}=\ell^2(\mathbb{T};\mathbb{R}^3)$ on the lateral torus $\mathbb{T}$, with
transverse part $\mathfrak{H}_T$, orthogonal projection $P$ onto transverse fields and
$P^{\perp}=1-P$. Its operators $\lambda=-\Delta_{\parallel}$, $\mathsf{v}$ and
$\mathsf{C}=1-\mathsf{v}$, with the kernel $K_1$ of $\mathsf{C}$, are those of
Lemma~\ref{8.14:lem-poisson-identities}. Its quadratic defect form is
\[
N(w)=\tfrac12\langle Pw,\mathsf{C}^{-1}Pw\rangle+\tfrac12\|P^{\perp}w\|^2 .
\]
Its caps form the box
$\mathfrak{B}=\{E\in\mathfrak{H}:|E_k(x)|\le\tau_k\ \text{for all }x,k\}$ with half-widths
$\tau_k$ as in Theorem~\ref{8.14:thm-exit-value}. Its free pull is $d_0(\beta)=-\mathsf{v}\beta$,
and $W=V_1-V_+$ is the excess of the one-capped-level value over the free value.
Let $C_N=\max\mathsf{C}^{-1}$ and $C'=\tfrac12(1+C_N)$.

Let $\beta\in\mathfrak{H}_T$, let $E_0\in\mathfrak{B}$ be its $V_1$-optimal step and $\mu$ the
multiplier of the capped level, and put $\beta'=\beta+PE_0$. Let $\kappa_1$ be the one-step
Poisson deficit of the absorption bound~\eqref{8.14:eq-concavity-absorption-b}. That is, for
$(x,k)$ with $\mu_k(x)\ne0$ and $s=\operatorname{sign}\mu_k(x)$,
\[
\kappa_1(x,k)=\sum_{x'}K_1(x-x')\bigl[\tau_k-sE_{0,k}(x')\bigr],
\]
and $\kappa_1(x,k)=0$ where $\mu_k(x)=0$; such pairs contribute nothing to the sums below.
Then
\begin{equation}\label{8.14:eq-one-row-ceiling-1}
\begin{split}
W(\beta)-\langle P^{\perp}\mu,\mathsf{v}P^{\perp}\mu\rangle
-\sum_{x,k}|\mu_k(x)|\,\kappa_1(x,k)
&\le W(\beta')\\
&\le W(\beta)-\max\Bigl\{\tfrac12\langle P^{\perp}\mu,\mathsf{v}(1+\mathsf{C})P^{\perp}\mu\rangle,\\
&\qquad\tfrac12\sum_{x,k}|\mu_k(x)|\,
\min\Bigl(\kappa_1(x,k),\frac{|\mu_k(x)|}{2C'}\Bigr)\Bigr\}.
\end{split}
\end{equation}
\end{lemma}

\begin{proof}
\emph{The upper member.} The upper member of~\eqref{8.14:eq-one-row-ceiling-1} states that
$W(\beta')$ is at most each of the two quantities obtained by subtracting from $W(\beta)$ one of
the two terms inside the maximum. These two bounds are the concavity
bound~\eqref{8.14:eq-concavity-absorption-a} and the absorption
bound~\eqref{8.14:eq-concavity-absorption-b} of the model proposition,
Proposition~\ref{8.14:prf-concavity-absorption} (concavity, absorption and exactness). The concavity
bound~\eqref{8.14:eq-concavity-absorption-a} reads
\[
W(\beta')\le W(\beta)-\tfrac12\langle P^{\perp}\mu,(1-\mathsf{C}^2)P^{\perp}\mu\rangle .
\]
By the identity $1-\mathsf{C}^2=\mathsf{v}(1+\mathsf{C})$ of~\eqref{8.14:eq-poisson-identities-f}
in Lemma~\ref{8.14:lem-poisson-identities}, a model lemma, its right-hand side is $W(\beta)$ minus
the first term inside the maximum. The absorption
bound~\eqref{8.14:eq-concavity-absorption-b} reads
\[
W(\beta')\le W(\beta)-\tfrac12\sum_{x,k}|\mu_k(x)|\,
\min\Bigl(\kappa_1(x,k),\frac{|\mu_k(x)|}{2C'}\Bigr),
\]
whose right-hand side is $W(\beta)$ minus the second term inside the maximum. Since $W(\beta')$
is at most both right-hand sides, it is at most $W(\beta)$ minus the larger of the two
subtracted terms, which is the upper member.

\emph{The lower member: the objects.} By the model theorem,
Theorem~\ref{8.14:thm-exit-value}, $W(\beta)=\min_{E\in\mathfrak{B}}N(d_0(\beta)-E)$, and the
minimum is attained at the unique point $E_0$. The multiplier is
\[
\mu=\nabla N(d_0(\beta)-E_0).
\]
Complementarity holds: $\mu_k(x)>0$ implies $E_{0,k}(x)=\tau_k$, and $\mu_k(x)<0$ implies
$E_{0,k}(x)=-\tau_k$. Since $P$ projects onto transverse fields, $\beta'=\beta+PE_0$ lies in
$\mathfrak{H}_T$. The same theorem at $\beta'$ therefore gives
$W(\beta')=\min_{E'\in\mathfrak{B}}N(d_0(\beta')-E')$.

Fix any $E'\in\mathfrak{B}$ and put $\delta=E'-\mathsf{C}E_0$ and
$\bar w=d_0(\beta')-\mathsf{C}E_0$, so that $d_0(\beta')-E'=\bar w-\delta$.

\emph{The lower member: the expansion.} The form $N$ is the quadratic form of the operator
$P\mathsf{C}^{-1}P+P^{\perp}$:
\[
N(w)=\tfrac12\bigl\langle w,(P\mathsf{C}^{-1}P+P^{\perp})w\bigr\rangle .
\]
This operator is
self-adjoint because $P$ is an orthogonal projection and $\mathsf{C}^{-1}$, a function of
$\lambda$, is self-adjoint. Hence $\nabla N(w)=(P\mathsf{C}^{-1}P+P^{\perp})w$, and expanding the
quadratic form gives
\begin{equation}\label{8.14:eq-one-row-ceiling-2}
N(\bar w-\delta)=N(\bar w)-\langle\nabla N(\bar w),\delta\rangle+N(\delta).
\end{equation}
The proof of~\eqref{8.14:eq-concavity-absorption-a} and~\eqref{8.14:eq-concavity-absorption-b}
in the model proposition, Proposition~\ref{8.14:prf-concavity-absorption}, computes, for
this $\bar w$,
\[
N(\bar w)=W(\beta)-\tfrac12\langle P^{\perp}\mu,(1-\mathsf{C}^2)P^{\perp}\mu\rangle,
\qquad
\nabla N(\bar w)=\mu-(1-\mathsf{C})P^{\perp}\mu .
\]
Inserting these into~\eqref{8.14:eq-one-row-ceiling-2} gives
\begin{equation}\label{8.14:eq-one-row-ceiling-3}
\begin{split}
N(d_0(\beta')-E')={}&W(\beta)-\tfrac12\langle P^{\perp}\mu,(1-\mathsf{C}^2)P^{\perp}\mu\rangle\\
&-\langle\mu,\delta\rangle+\langle(1-\mathsf{C})P^{\perp}\mu,\delta\rangle+N(\delta).
\end{split}
\end{equation}

\emph{First estimate: $\langle\mu,\delta\rangle\le\sum_{x,k}|\mu_k(x)|\kappa_1(x,k)$.}
Summands with $\mu_k(x)=0$ vanish, so only $(x,k)$ with $\mu_k(x)\ne0$ contribute. Fix such a
pair and let $s=\operatorname{sign}\mu_k(x)$. By complementarity, $E_{0,k}(x)=s\tau_k$. Since
$(\mathsf{C}E_0)_k(x)=\sum_{x'}K_1(x-x')E_{0,k}(x')$, $\sum_{x'}K_1(x-x')=1$ and $s^2=1$, we have
\[
(\mathsf{C}E_0)_k(x)
=s\Bigl[\tau_k-\sum_{x'}K_1(x-x')\bigl(\tau_k-sE_{0,k}(x')\bigr)\Bigr]
=s\bigl(\tau_k-\kappa_1(x,k)\bigr).
\]
Writing $\mu_k(x)=s|\mu_k(x)|$, we obtain
\begin{align*}
\langle\mu,\delta\rangle
&=\sum_{\mu_k(x)\ne0}\mu_k(x)\bigl[E'_k(x)-s(\tau_k-\kappa_1(x,k))\bigr]\\
&=\sum_{\mu_k(x)\ne0}|\mu_k(x)|\bigl[sE'_k(x)-\tau_k+\kappa_1(x,k)\bigr].
\end{align*}
Since $E'\in\mathfrak{B}$, we have $sE'_k(x)\le|E'_k(x)|\le\tau_k$. Each bracket is therefore at
most $\kappa_1(x,k)$, which is the first estimate.

\emph{Second estimate.} By the Cauchy--Schwarz inequality and $ab\le\tfrac12a^2+\tfrac12b^2$,
\[
\langle(1-\mathsf{C})P^{\perp}\mu,\delta\rangle
\ge-\tfrac12\|(1-\mathsf{C})P^{\perp}\mu\|^2-\tfrac12\|\delta\|^2 .
\]
By Lemma~\ref{8.14:lem-poisson-identities}, a model lemma, $\mathsf{C}^{-1}=1+\lambda+\mathsf{v}$,
where $\lambda\ge0$ and $\mathsf{v}\ge0$ is the non-negative root; hence $\mathsf{C}^{-1}\ge1$.
It follows that
\[
N(\delta)\ge\tfrac12\|P\delta\|^2+\tfrac12\|P^{\perp}\delta\|^2=\tfrac12\|\delta\|^2 .
\]
Since $1-\mathsf{C}$ is self-adjoint,
$\|(1-\mathsf{C})P^{\perp}\mu\|^2=\langle P^{\perp}\mu,(1-\mathsf{C})^2P^{\perp}\mu\rangle$.

\emph{Conclusion.} Insert both estimates into~\eqref{8.14:eq-one-row-ceiling-3}. The terms
$\pm\tfrac12\|\delta\|^2$ cancel, leaving
\begin{align*}
N(d_0(\beta')-E')\ge{}&W(\beta)
-\tfrac12\bigl\langle P^{\perp}\mu,\bigl[(1-\mathsf{C}^2)+(1-\mathsf{C})^2\bigr]P^{\perp}\mu\bigr\rangle\\
&-\sum_{x,k}|\mu_k(x)|\,\kappa_1(x,k).
\end{align*}
The operator in the middle term simplifies:
\[
\tfrac12\bigl[(1-\mathsf{C}^2)+(1-\mathsf{C})^2\bigr]
=\tfrac12\bigl(1-\mathsf{C}^2+1-2\mathsf{C}+\mathsf{C}^2\bigr)=1-\mathsf{C}=\mathsf{v},
\]
where the last equality is the identity $\mathsf{v}+\mathsf{C}=1$ of
Lemma~\ref{8.14:lem-poisson-identities}, a model lemma. Hence, for every $E'\in\mathfrak{B}$,
\[
N(d_0(\beta')-E')\ge W(\beta)-\langle P^{\perp}\mu,\mathsf{v}P^{\perp}\mu\rangle
-\sum_{x,k}|\mu_k(x)|\,\kappa_1(x,k).
\]
Minimising over $E'\in\mathfrak{B}$ and using
$W(\beta')=\min_{E'\in\mathfrak{B}}N(d_0(\beta')-E')$ gives the lower member
of~\eqref{8.14:eq-one-row-ceiling-1}.
\end{proof}

\begin{remark}[Model: invisibility of deficits beyond one row]\label{8.14:rem-one-row-deficit}
This remark concerns the layered quadratic model of this section and only that model. Each row carries a level in
$\mathfrak{H}=\ell^2(\mathbb{T};\mathbb{R}^3)$ over the lateral torus $\mathbb{T}=(\mathbb{Z}/N)^3$. The free value is
$V_+(\gamma)=\tfrac12\langle\gamma,\mathsf{v}\gamma\rangle$. The value $V_1$ is the value with one level capped in the box
$\mathfrak{B}$ of young caps, and $W=V_1-V_+$, as in the model lemma, Lemma~\ref{8.14:lem-innovations-identity}. We write $\tau_k$
for the cap of the $k$-th component, so that $|E'_k(x)|\le\tau_k$ for every $E'\in\mathfrak{B}$. A \emph{hot site} is a pair
$(x,k)$ in the support of the multiplier $\mu$. At a hot site we write $\sigma_k(x)$ for the sign of $\mu_k(x)$.

The thesis of this remark is that the absorption cannot be sharpened to a scale-invariant deficit per dyadic band within the
one-row method. We argue this in three steps.

\emph{First step: the drop of $W$ along one row is bounded above and below by one-row quantities.} Let
$\beta\in\mathfrak{H}_T$, let $(E_0,\mu)$ be its
$V_1$-optimal step and multiplier, and put $\beta'=\beta+PE_0$. By the model lemma on the one-row ceiling,
Lemma~\ref{8.14:lem-one-row-ceiling}, rearranged,
\begin{equation}\label{8.14:eq-one-row-deficit-1}
\begin{split}
&\tfrac12\langle P^{\perp}\mu,\mathsf{v}(1+\mathsf{C})P^{\perp}\mu\rangle
+\tfrac12\sum_{x,k}|\mu_k(x)|\min\bigl(\kappa_1(x,k),|\mu_k(x)|/(2C')\bigr)\\
&\qquad\le W(\beta)-W(\beta')
\le\langle P^{\perp}\mu,\mathsf{v}P^{\perp}\mu\rangle+\sum_{x,k}|\mu_k(x)|\kappa_1(x,k).
\end{split}
\end{equation}
Thus the vectorial part of the drop lies between $\tfrac12\langle P^{\perp}\mu,\mathsf{v}(1+\mathsf{C})P^{\perp}\mu\rangle$
and $\langle P^{\perp}\mu,\mathsf{v}P^{\perp}\mu\rangle$, a factor of at most $2$. The box part is at most
$\sum_{x,k}|\mu_k(x)|\kappa_1(x,k)$, where $\kappa_1$ is the one-step Poisson deficit.

Whatever base point one expands the defect form $N$ around, the linear gain at a hot site, measured against $\mathsf{C}E_0$,
satisfies
\begin{equation}\label{8.14:eq-one-row-deficit-2}
\mu_k(x)\bigl(E'_k(x)-(\mathsf{C}E_0)_k(x)\bigr)\le|\mu_k(x)|\,\kappa_1(x,k)
\qquad\text{for every } E'\in\mathfrak{B}.
\end{equation}
This is the first step of the proof of the lower member in the model lemma, Lemma~\ref{8.14:lem-one-row-ceiling}: by
complementarity $(E_0)_k(x)=\sigma_k(x)\tau_k$ at every hot site, hence
$(\mathsf{C}E_0)_k(x)=\sigma_k(x)\bigl(\tau_k-\kappa_1(x,k)\bigr)$ there, so the left side of
\eqref{8.14:eq-one-row-deficit-2} equals
$|\mu_k(x)|\bigl(\sigma_k(x)E'_k(x)-\tau_k+\kappa_1(x,k)\bigr)$, and $\sigma_k(x)E'_k(x)\le\tau_k$. The box forbids pushing
past the cap. A deficit $\kappa_j$ at a scale $j>1$ is
therefore invisible to the one-row sandwich \eqref{8.14:eq-one-row-deficit-1}.

To see this directly, let $\kappa_j(x,k)$ denote the scale-$j$ Poisson deficit, so that on the support of $\mu$
\[
(\mathsf{C}^jE_0)_k(x)=\sigma_k(x)\bigl(\tau_k-\kappa_j(x,k)\bigr).
\]
Try $E'=\mathsf{C}^jE_0+\eta'$ with $|\eta'_k(x)|\le\kappa_j(x,k)$. At a hot site the gain is largest for
$\eta'_k(x)=\sigma_k(x)\kappa_j(x,k)$. Since
$(\mathsf{C}^jE_0-\mathsf{C}E_0)_k(x)=-\sigma_k(x)\bigl(\kappa_j(x,k)-\kappa_1(x,k)\bigr)$ there, the gain is
\begin{align*}
\mu_k(x)\bigl(E'_k(x)-(\mathsf{C}E_0)_k(x)\bigr)
&=|\mu_k(x)|\bigl[\kappa_j(x,k)-\bigl(\kappa_j(x,k)-\kappa_1(x,k)\bigr)\bigr]\\
&=|\mu_k(x)|\kappa_1(x,k),
\end{align*}
which is the same as before. Off the support of $\mu$ the extra displacement $(\mathsf{C}-\mathsf{C}^j)E_0$ does not enter
the linear gain $\langle\mu,E'-\mathsf{C}E_0\rangle$, since $\mu$ vanishes there; it only costs.

\emph{Second step: comparing rows $j$ apart.} To see $\kappa_j$ one must compare rows $j$ apart. Let $m(a)$ be the value of
the truncated problem $P_a$ and $X_a$ its minimiser. Write $\beta^{(b)}_a$ for the row-$a$ trace of $X_b$. Write $W_j=V_j-V_+$,
where $V_j$ is the value with $j$ capped levels. Put
\[
M_j(a)=m(a+j)-m(a).
\]
The model lemma, Lemma~\ref{8.14:lem-innovations-identity}, applied with $j$ capped levels, by the argument that gives
\eqref{8.14:eq-exit-value-c} in the model exit-value theorem, yields
\begin{equation}\label{8.14:eq-one-row-deficit-3}
W_j\bigl(\beta^{(a+j)}_a\bigr)\le M_j(a)\le W_j\bigl(\beta^{(a)}_a\bigr).
\end{equation}
With $a+j$ in place of $a$, the upper bound for $M_j(a+j)$ is $W_j(\beta^{(a+j)}_{a+j})$. The bounds for $M_j(a+j)$ and for
$M_j(a)$ therefore both involve the minimiser $X_{a+j}$.

Consider $j$ rows $\beta_0=\beta,\beta_1,\dots,\beta_j$ with levels $E^{(0)},\dots,E^{(j-1)}$ and
$\beta_{l+1}=\beta_l+PE^{(l)}$. The natural competitor is the $j$-row shift
\[
(\beta_i,E^{(i)})_{i<j}\longmapsto\bigl(\beta_j+\mathsf{C}^j(\beta_i-\beta),\ \mathsf{C}^jE^{(i)}\bigr).
\]
It does not reproduce the transverse defects. Let $\xi_l=d_0(\beta_l)-PE^{(l)}$. Since
$\beta_{l+1}=\beta_l+PE^{(l)}=\mathsf{C}\beta_l-\xi_l$, induction gives
\[
\beta_l=\mathsf{C}^l\beta-\sum_{l'<l}\mathsf{C}^{l-1-l'}\xi_{l'}.
\]
Inserting this, the defects of the shifted configuration are
\begin{equation}\label{8.14:eq-one-row-deficit-4}
\tilde \xi_i=\mathsf{C}^j\xi_i+\mathsf{v}\sum_{l<j}\mathsf{C}^{j-1-l}\xi_l .
\end{equation}
The second member of \eqref{8.14:eq-one-row-deficit-4} is the propagated deviation of the trace from free evolution. It
produces unsigned cross terms of the size of the signal.

For $j=1$ it collapses to $\tilde \xi_0=(\mathsf{C}+\mathsf{v})\xi_0=\xi_0$. This is why the concavity step of the model
statement on concavity, absorption and exactness is exact. For $j>1$ this
competitor
gives no comparison of $M_j(a+j)$ with $M_j(a)$. A concavity statement over several rows would need a different competitor,
and none is given in this section.

\emph{Third step: a finer deficit would not remove the cutoff.} Such a statement would not help in any case. By the model
remark, Remark~\ref{8.14:rem-cutoff-in-height}, the rate of \eqref{8.14:eq-stopping-height-b} is already the scale-covariant
one. A per-dyadic-band deficit could at best remove the factor $s_\infty/\tau$ of \eqref{8.14:eq-cutoff-in-height-b} from
$\mathfrak{h}$, with $\tau$ the cap height entering there. A corollary later in this section does exactly that on the
charged class, by other means. It could not remove the count $\mathfrak{n}$ of
\eqref{8.14:eq-cutoff-in-height-c}, and that count is where the term $5k_o\log L$ sits, $k_o$ being the level entering
that count.
\end{remark}

\begin{remark}[Model: ring freedom absorbs only cutoff-free constants]\label{8.14:rem-ring-freedom}
The model here is the one of this section. It consists of the truncated variational problems
$P_a$ for level fields in $\mathfrak{H}=\ell^2(\mathbb{T};\mathbb{R}^3)$ on the lateral torus
$\mathbb{T}=(\mathbb{Z}/N)^3$, with the young caps $\mathfrak{B}$ imposed on the capped levels. The
front of capped levels of this model stops within the height $H_{\mathrm{EV}}$ of
\eqref{8.14:eq-cutoff-in-height-a}, by the model statements Corollary~\ref{8.14:cor-stopping-height}
and Remark~\ref{8.14:rem-cutoff-in-height}, on the range fixed in the latter.

Let $j$ be the seed level, let $c_0$ be the number of levels for which the observation level is
$k_o=j+c_0+1$, and let $c_\sharp$ be the offset for which the rings of the ladder are the levels $i$
with $j+c_\sharp+1\le i\le k_o-1$. One might hope that it suffices for the front to stop within
these $c_0-c_\sharp$ rings. The arithmetic below shows that this hope absorbs only constants free
of $g$.

In this model the freedom to let the front cross the $c_0-c_\sharp-1$ rings between the seed level
$j$ and the observation level $k_o=j+c_0+1$ absorbs constants independent of $g$. It absorbs
nothing that grows with $k_o$ or with $g$. Requiring the stopping height not to exceed the depth
available across these rings would impose on $c_0$ the floor
\begin{equation}\label{8.14:eq-ring-freedom-1}
c_0\ \ge\ c_\sharp+\log_L(j+c_0+1)+O(1),
\end{equation}
which depends on $j$. Hence no constant $c_0$ satisfies this condition for every seed level $j$.
Under it the observation level recedes logarithmically from the seed, and the form of the
fixed-depth condition of this chapter, in which $k_o-j$ is a fixed number of levels, is changed.

\emph{The available depth.} The ladder of rings between the seed level and the observation level
lets the front cross $c_0-c_\sharp-1$ rings. For $j+c_\sharp+1\le i\le k_o-1$, ring $i$ consists
of level-$i$ cubes and is $L^{k_o-1-i}$ times deeper than the outer ring. We use this geometry as
it is stated in part~(c) of the existence proposition for the ladder in this chapter.
Summing the geometric series over the rings, the depth available to the
front is
\begin{equation}\label{8.14:eq-ring-freedom-2}
D_{\mathrm{avail}}\ \le\ C_r\,\varpi_y\,n\,L^{c_0-c_\sharp}
\end{equation}
fine levels. Here $C_r$ is free of $n$: it counts the rows per ring and absorbs the geometric sum.

Even if the young caps grow inward by factors $\mathsf{T}_i\le L^2$ per ring (the factor
$L^{-2(k_o-1-i)}$ in part~(ii) of the proposition on the ladder), this growth can only shorten the
front. We nevertheless take the same
cap on every ring, which is the case most favourable to the stopping height.

\emph{The required depth.} By \eqref{8.14:eq-stopping-height-b} (Corollary~\ref{8.14:cor-stopping-height},
a model statement), the front stops within the height $H_{\mathrm{EV}}$ of
\eqref{8.14:eq-cutoff-in-height-a}, on the range of Remark~\ref{8.14:rem-cutoff-in-height} on which
$\kappa_1\ge\kappa_d$ and $\|\mu\|_\infty\le s_\infty$. By the two factors
\eqref{8.14:eq-cutoff-in-height-b} and
\eqref{8.14:eq-cutoff-in-height-c} of the model statement Remark~\ref{8.14:rem-cutoff-in-height},
up to factors free of $n$,
\begin{equation}\label{8.14:eq-ring-freedom-3}
H_{\mathrm{EV}}=\mathfrak{h}\,\mathfrak{n}+1\ \asymp\
\varrho_*\,n\,\bigl(1+c_B\log N\bigr)\bigl(5k_o\log L+C_{\mathfrak{n}}\bigr),
\qquad N=2P_{\mathrm{bl}}\,n.
\end{equation}
Here $c_B$ is the $n$-free constant of the factor \eqref{8.14:eq-cutoff-in-height-b}, and
$C_{\mathfrak{n}}$ is the $n$-free additive constant of the factor
\eqref{8.14:eq-cutoff-in-height-c}. Insert \eqref{8.14:eq-ring-freedom-2} and
\eqref{8.14:eq-ring-freedom-3} into the condition $H_{\mathrm{EV}}\le D_{\mathrm{avail}}$ and divide
by $C_r\varpi_y n$. Since $n=L^{k_o}$, we have $\log N=\log(2P_{\mathrm{bl}})+k_o\log L$, and the
condition then reads
\begin{equation}\label{8.14:eq-ring-freedom-4}
L^{c_0-c_\sharp}\ \ge\ \frac{\varrho_*}{C_r\varpi_y}\,
\bigl(1+c_B(\log 2P_{\mathrm{bl}}+k_o\log L)\bigr)\bigl(5k_o\log L+C_{\mathfrak{n}}\bigr).
\end{equation}

\emph{The floor.} In \eqref{8.14:eq-ring-freedom-4} we drop the logarithm in the first factor,
keeping only its term $1$. We bound the second factor below by $5k_o\log L$. The condition then
implies
\begin{equation*}
L^{c_0-c_\sharp}\ \ge\ \frac{5\varrho_*}{C_r\varpi_y}\,k_o\log L .
\end{equation*}
Taking $\log_L$ of both sides gives
\begin{equation}\label{8.14:eq-ring-freedom-5}
c_0-c_\sharp\ \ge\ \log_L k_o+\log_L\log L+\log_L\frac{5\varrho_*}{C_r\varpi_y}.
\end{equation}
Both factors on the right of \eqref{8.14:eq-ring-freedom-4} grow linearly in $k_o$. Either of
them alone forces a term $\log_L k_o$ in the floor; in \eqref{8.14:eq-ring-freedom-5} this term
comes from the second factor. Keeping both factors instead of dropping the logarithm in the first
gives $2\log_L k_o$ in place of $\log_L k_o$ in \eqref{8.14:eq-ring-freedom-5}.

Now $k_o=j+c_0+1$, so \eqref{8.14:eq-ring-freedom-5} gives \eqref{8.14:eq-ring-freedom-1}, with the
$O(1)$ term independent of $j$. The seed level $j$ ranges over the whole hierarchy. For any fixed
$c_0$, the term $\log_L(j+c_0+1)$ is unbounded in $j$, while the left side of
\eqref{8.14:eq-ring-freedom-1} is fixed. Hence no constant $c_0$ satisfies
\eqref{8.14:eq-ring-freedom-1} for all $j$.

\emph{Comparison with the stopping condition.} The analysis of the stopping condition at fixed
$c_0$ in this chapter leads to a floor of the same shape,
$c_0\ge c_\sharp+C\log_L\log P_{\mathrm{bl}}(g)$ with a constant $C$. In the present model
the term $\log_L k_o$ takes the place of $\log_L\log P_{\mathrm{bl}}(g)$. Neither floor is imposed
as a condition on $c_0$. The inequality \eqref{8.14:eq-ring-freedom-1} is the reason that the
stopping height of the model is not asserted to fit within a fixed number of rings.

\emph{Conclusion.} The ring freedom therefore absorbs constants independent of $g$: by part~(b) of
the statement of this chapter on the absorption of constants by additional levels, such a constant
$B^\natural$ is absorbed by
\begin{equation*}
\Bigl\lceil \log\bigl(4B^\natural\cdots\bigr)\big/\log\bigl(L^2/(1+\beta_0)\bigr)\Bigr\rceil
\end{equation*}
levels, the dots standing for the further factors displayed in that statement, and $\beta_0$ being
the parameter of that statement. It absorbs nothing that grows with $k_o$ or with $g$. By part~(c)
of the same statement, a constant that grows with $\log P_{\mathrm{bl}}$ is absorbed by no fixed
number of levels, and the arithmetic above shows that the same holds here for $\log n$.
\end{remark}

\begin{remark}[Model: the double logarithm belongs to the method]\label{8.14:rem-double-logarithm}
This remark compares the one-row method of the model theorem on the exit value, specialised to a
scalar field with one lateral direction of size $N$, with the scalar model in which the model
theorem on inactivity at full band width, called the inactivity theorem below, is proved: a
scalar field with one lateral direction of size $N$, the front of active levels advancing from a
clamped level. The model theorem on the exit value holds verbatim for scalar fields in any lateral
dimension, as observed in this section for its setting. Its member for a scalar field with one
lateral direction differs from the scalar model of the inactivity theorem only in the mirror band,
as also observed there: a half-space against a symmetric strip.

We use the letters $\tau_y$, $\ell$, $b_y$, $a_0$, $W_y$, $C_K^0$, $\Lambda_N$, $\lambda_N$,
$\hat\psi$ of the inactivity theorem. We also use the letters $n$,
$k_o$, $s_\infty$, $\kappa_d$, $\mathfrak{h}$, $\mathfrak{n}$ of the model remark,
Remark~\ref{8.14:rem-cutoff-in-height}. We say that a front quantity is \emph{cutoff-free} when
its bound does not depend on $n$, and \emph{cutoff-dependent} when it grows like the logarithm of
the lateral size $N$, hence without bound as $k_o$ grows.

\emph{The front stops cutoff-free in the model.} By parts (b) and (c) of the inactivity theorem,
the tight young levels lie within distance $a_0$ of the clamped
level, where
\begin{equation*}
a_0\le\max\Bigl(b_y,\;N\,e^{-(\tau_y/(3\ell)-C_K^0-6)/6}\Bigr).
\end{equation*}
Under the floor condition of the corollary to that theorem, $a_0$ is a fraction of $W_y$. This
bound is cutoff-free, since $n$ cancels in it.

\emph{The ingredients of the stopping height are cutoff-dependent at that front.} The stopping
height of the model display \eqref{8.14:eq-stopping-height-b} uses two ingredients at the front.
Both are cutoff-dependent there.

The exit field one level above a clamped level is bounded by $[C_K^0+6\log N]\ell$. This is part
(iii) of the theorem on the scalar model $\widehat{M}$, read at the first level above the clamped
level. The remark following the inactivity theorem states that this
quantity is cutoff-dependent at that level and the next ones, as it must be. At the level of the
kernel the logarithm is two-sided by the free-continuation bound: on the two-dimensional strip the
kernel of the free continuation has $\ell^1$-norm comparable, above and below, to the logarithm
of the lateral size.

Hence the ratio $s_\infty/\tau$ of the supremum of the top multiplier to the threshold $\tau$,
entering the model display \eqref{8.14:eq-cutoff-in-height-b}, is
cutoff-dependent. This ratio enters through part (ii) of the model lemma,
Lemma~\ref{8.14:lem-exit-pull}, which expresses the top multiplier through the exit pull, the
field of the level below and the curvature of the exit row.

The count in the model display \eqref{8.14:eq-cutoff-in-height-c} compares an energy with the
threshold of one pair. It does so in any dimension.

Consequently no sharpening of the ingredients of \eqref{8.14:eq-stopping-height-b} themselves can
make that stopping height cutoff-free. They are cutoff-dependent front quantities in a model in
which cutoff-free stopping is a theorem, namely the inactivity theorem.

\emph{What the inactivity theorem uses instead is integrated over the height.} Part (a) of that
model theorem, the lift, bounds the potential climbed over $k$ levels: for every lateral site $x$,
every level $p$ and every number $k$ of levels in the range of that theorem,
\begin{equation}\label{8.14:eq-double-logarithm-1}
|\hat\psi(x,p+k)-\hat\psi(x,p)|\le\ell\bigl[\Lambda_N(k+b_y)+\Lambda_N(b_y)\bigr].
\end{equation}
Here the lift functions satisfy
\begin{equation}\label{8.14:eq-double-logarithm-2}
\Lambda_N(t)\le t\,\lambda_N(t),\qquad
\lambda_N(t)=C_K^0+6+6\log^+(N/t).
\end{equation}
Inserting $t=k+b_y$ into \eqref{8.14:eq-double-logarithm-2}, the logarithm attached to a climb of
length $k$ in \eqref{8.14:eq-double-logarithm-1} is $\log^+\bigl(N/(k+b_y)\bigr)$. This integrates
the cutoff-dependent front values into a cutoff-free length.

Moreover the lift is a comparison: it rests on the lattice minimum and maximum $\psi\wedge\Psi$,
$\psi\vee\Psi$ of two configurations $\psi$, $\Psi$, as in the comparison lemma of the scalar
model. That comparison in turn rests on order, namely on the column-wise sign structure of the
Karush--Kuhn--Tucker conditions in the comparison-principle lemmas for the exit problem. For
$1$-forms this sign structure meets the constraint on the multiplier charge $q=\nabla^*\mu$.

So the mechanism by which the model stops its front cutoff-free is exactly the one that this
charge constraint excludes for vector fields. The front-by-front value functional of the model
theorem on the exit value is exactly the kind of quantity that is cutoff-dependent even where the
true stopping is cutoff-free.

\emph{Where the cutoff enters.} The cutoff enters the stopping height
\eqref{8.14:eq-stopping-height-b} in two places.
\begin{itemize}
\item[(1)] Through \eqref{8.14:eq-cutoff-in-height-b}, the supremum of the top multiplier. By
parts (i) and (ii) of the model lemma, Lemma~\ref{8.14:lem-exit-pull}, the cutoff-dependent part
of the top multiplier is the exit pull, the Riesz image of the tangential curvature of the exit
row. It is of size $\log N$ at a free boundary, a logarithm attained at the level of the kernel in
the sense described above.
\item[(2)] Through \eqref{8.14:eq-cutoff-in-height-c}, the count $\mathfrak{n}$ of e-folds. This
count compares an extensive budget with the threshold of one site, and equals $5k_o\log L$ up to
an additive term independent of $n$.
\end{itemize}
The cutoff does not enter through the scale $\kappa_d$ of the one-step Poisson deficit. Nor does
it enter through the constants of the model theorem on the exit value in
\eqref{8.14:eq-poisson-identities-h}.

The first of these enters the stopping height as a factor, through $\mathfrak{h}$. The second
enters it as the number $\mathfrak{n}$ of e-folds. Either one alone already exceeds what the ring
freedom can absorb; see the model remark, Remark~\ref{8.14:rem-ring-freedom}.

Finally, consider a sharpening of the absorption to a scale-invariant deficit per dyadic band.
By the model lemma, Lemma~\ref{8.14:lem-one-row-ceiling}, it cannot enter the one-row method. At
best it addresses the scale $\kappa_d$ of the deficit or the supremum in (1), and never the count
in (2).
\end{remark}

\begin{lemma}[Model: concavity defect as multiplier-charge energy]\label{8.14:lem-charge-energy}
The model is that of Lemma~\ref{8.14:lem-poisson-identities}, the model statement on Poisson
operator identities: real level fields on a three-dimensional lateral torus, acted on componentwise
by minus the lateral Laplacian and by its one-step Poisson operator. Its ingredients are the
following.
\begin{itemize}
\item The lateral torus $\mathbb{T}$ is a three-dimensional discrete torus, and a level field is an
element of $\mathfrak{H}=\ell^2(\mathbb{T};\mathbb{R}^3)$.
\item $P$ is the orthogonal projection of $\mathfrak{H}$ onto its transverse part
$\mathfrak{H}_T$, and $P^{\perp}=1-P$.
\item $\lambda=-\Delta_{\parallel}$, where $\Delta_{\parallel}$ is the lateral lattice Laplacian,
acting componentwise on level fields.
\item $\mathsf{v}=\tfrac12\bigl[-\lambda+(\lambda^2+4\lambda)^{1/2}\bigr]$ is the free-value
operator, and $\mathsf{C}=1-\mathsf{v}=(1+\lambda+\mathsf{v})^{-1}$ is the one-step Poisson
operator.
\item $P_a$ are the truncated capped problems, with values $m(a)$ and increments $\Delta m(a)$.
\end{itemize}
Let $e_j$ denote the unit vector of the $j$-th lateral direction, $j=1,2,3$. Write
$(\nabla\varphi)_j(x)=\varphi(x+e_j)-\varphi(x)$ for the forward lateral gradient of a
scalar $\varphi$, and write $\partial_j^-f(x)=f(x)-f(x-e_j)$. Put
$\operatorname{div}^-\nu=\sum_j\partial_j^-\nu_j$ and $\nabla^*=-\operatorname{div}^-$.
For $\mu\in\mathfrak{H}$ put $q:=\nabla^*\mu\in\ell^2(\mathbb{T};\mathbb{R})$, the lateral
co-differential of $\mu$; then $\sum_xq(x)=0$. Let $\lambda:=-\Delta_{\parallel}$ act on mean-zero
scalars, and let $\mathsf{v},\mathsf{C}$ denote the same functions of this operator. Then
\begin{equation}\label{8.14:eq-charge-energy-1}
\tfrac12\bigl\langle P^{\perp}\mu,(1-\mathsf{C}^2)P^{\perp}\mu\bigr\rangle
=\tfrac12\bigl\langle q,(1-\mathsf{C}^2)\lambda^{-1}q\bigr\rangle
=\tfrac12\bigl\langle q,(1+\mathsf{C})(\lambda+\mathsf{v})^{-1}q\bigr\rangle .
\end{equation}
Moreover, $(1+\mathsf{C})(\lambda+\mathsf{v})^{-1}$ is a positive multiplier. It is bounded above
and below by absolute constant multiples of $\min(\lambda^{-1},\lambda^{-1/2})$. Thus the right side
of \eqref{8.14:eq-charge-energy-1} is comparable to the square of the discrete Sobolev norm of $q$
of order $-1$ on the modes with $\lambda\ge1$ (lattice scales) and of order $-\tfrac12$ on the
modes with $\lambda\le1$ (large scales).

Let $\mu$ be the multiplier field of the top capped level of one of the truncated problems of the
upward variational problem. Then $q$ coincides with the multiplier charge $(\delta\mu)_4$ of that
level, as defined in the Gauss-law lemma of the exit construction, up to a sign convention that is
immaterial in the quadratic form. Consequently, writing
\begin{equation*}
\|q\|^2_{(1+\mathsf{C})(\lambda+\mathsf{v})^{-1}}
:=\langle q,(1+\mathsf{C})(\lambda+\mathsf{v})^{-1}q\rangle ,
\end{equation*}
the concavity inequality \eqref{8.14:eq-concavity-absorption-a} of the model statement on
concavity, absorption and exactness reads
\begin{equation}\label{8.14:eq-charge-energy-2}
\Delta m(a+1)\le\Delta m(a)
-\tfrac12\,\bigl\|q^{(a+1)}\bigr\|^2_{(1+\mathsf{C})(\lambda+\mathsf{v})^{-1}} .
\end{equation}
Here $q^{(a+1)}$ is the charge of the top level of $P_{a+1}$.
\end{lemma}

\begin{proof}
\emph{Adjointness and the mean of $q$.} For a scalar $\varphi$ and a level field $\nu$, summation by
parts on $\mathbb{T}$ gives
\begin{equation*}
\langle\nabla\varphi,\nu\rangle
=\sum_x\sum_j\bigl(\varphi(x+e_j)-\varphi(x)\bigr)\nu_j(x)
=\sum_x\varphi(x)\sum_j\bigl(\nu_j(x-e_j)-\nu_j(x)\bigr)
=\langle\varphi,-\operatorname{div}^-\nu\rangle .
\end{equation*}
So $\nabla^*=-\operatorname{div}^-$ is the adjoint of the forward gradient, and
$\nabla^*\nabla=-\Delta_{\parallel}$ on scalars. Since the gradient of a constant vanishes, we get
$\sum_xq(x)=\langle1,\nabla^*\mu\rangle$, which equals $\langle\nabla1,\mu\rangle=0$. Hence
$\lambda^{-1}q$ is defined, $\lambda$ acting on mean-zero scalars.

\emph{Hodge decomposition on $\mathbb{T}$.} Let $\bar\mu$ be the mean of $\mu$, regarded as a constant
field, and put $\varphi:=\lambda^{-1}\nabla^*\mu=\lambda^{-1}q$. Then
$\nabla^*(\mu-\nabla\varphi)=q-\lambda\varphi=0$, and $\nabla^*\bar\mu=0$. The field
$\mu-\nabla\varphi-\bar\mu$ therefore has mean zero and is annihilated by $\nabla^*$, so it lies in
$\mathfrak{H}_T$. This gives
\begin{equation*}
\mu=P\mu+\nabla\varphi+\bar\mu,\qquad P^{\perp}\mu=\nabla\varphi+\bar\mu .
\end{equation*}
The field $\nabla\varphi+\bar\mu$ is orthogonal to $\mathfrak{H}_T$. Indeed, for $t\in\mathfrak{H}_T$
we have $\langle\nabla\varphi,t\rangle=\langle\varphi,\nabla^*t\rangle=0$, and
$\langle\bar\mu,t\rangle=0$ because $t$ has mean zero. Also, a constant is orthogonal to every
gradient: $\langle\bar\mu,\nabla\psi\rangle=\langle\nabla^*\bar\mu,\psi\rangle=0$.

\emph{A general identity.} Let $f$ be a real function of the Laplacian, acting on level fields
componentwise. Since the lateral torus is flat, $f(\lambda)$ commutes with lateral translations.
Hence $f(\lambda)\nabla=\nabla f(\lambda)$. Since $\lambda$ annihilates constants,
$f(\lambda)\bar\mu=f(0)\bar\mu$. Expanding $\langle P^{\perp}\mu,f(\lambda)P^{\perp}\mu\rangle$ with
$P^{\perp}\mu=\nabla\varphi+\bar\mu$, the two cross terms are
$\langle\bar\mu,\nabla f(\lambda)\varphi\rangle$ and $f(0)\langle\nabla\varphi,\bar\mu\rangle$. Both
are pairings of a constant with a gradient, so they vanish. Using $\nabla^*\nabla=\lambda$,
$\varphi=\lambda^{-1}q$, and the fact that functions of $\lambda$ are self-adjoint and commute, we get
\begin{equation}\label{8.14:eq-charge-energy-3}
\begin{aligned}
\bigl\langle P^{\perp}\mu,f(\lambda)P^{\perp}\mu\bigr\rangle
&=f(0)\|\bar\mu\|^2+\langle\nabla\varphi,\nabla f(\lambda)\varphi\rangle
=f(0)\|\bar\mu\|^2+\langle\varphi,\lambda f(\lambda)\varphi\rangle\\
&=f(0)\|\bar\mu\|^2+\langle q,f(\lambda)\lambda^{-1}q\rangle .
\end{aligned}
\end{equation}

\emph{The choice $f=1-\mathsf{C}^2$.} Since $\mathsf{v}(0)=0$, we have $\mathsf{C}=1$ on constants,
so $f(0)=0$. Identity \eqref{8.14:eq-poisson-identities-f} and identity
\eqref{8.14:eq-poisson-identities-c} of the model lemma, Lemma~\ref{8.14:lem-poisson-identities},
give respectively $1-\mathsf{C}^2=\mathsf{v}(1+\mathsf{C})$ and
$\mathsf{v}\lambda^{-1}=(\lambda+\mathsf{v})^{-1}$ on mean-zero data. Hence
\begin{equation*}
f(\lambda)\lambda^{-1}=\mathsf{v}(1+\mathsf{C})\lambda^{-1}=(1+\mathsf{C})(\lambda+\mathsf{v})^{-1}.
\end{equation*}
Inserting this into \eqref{8.14:eq-charge-energy-3} and multiplying by $\tfrac12$ gives both
equalities in \eqref{8.14:eq-charge-energy-1}.

\emph{Size of the multiplier.} Since $\mathsf{C}=(1+\lambda+\mathsf{v})^{-1}\in(0,1]$, we have
$1+\mathsf{C}\in[1,2]$. Since $\lambda,\mathsf{v}\ge0$, we have
$\lambda+\mathsf{v}\in[\max(\lambda,\mathsf{v}),2\max(\lambda,\mathsf{v})]$.

Next we bound $\mathsf{v}$. By identity \eqref{8.14:eq-poisson-identities-a} of the same model lemma,
$\mathsf{v}^2\le\mathsf{v}^2+\lambda\mathsf{v}=\lambda$, so $\mathsf{v}\le\sqrt\lambda$. The same
identity gives $\mathsf{v}=\lambda/(\lambda+\mathsf{v})$, whence
\begin{equation*}
\mathsf{v}\ge\frac{\lambda}{\lambda+\sqrt\lambda}=\frac{\sqrt\lambda}{1+\sqrt\lambda}
\ge\tfrac12\min(\sqrt\lambda,1).
\end{equation*}
The last step uses $1+\sqrt\lambda\le2$ when $\sqrt\lambda\le1$, and
$\sqrt\lambda/(1+\sqrt\lambda)\ge\tfrac12$ when $\sqrt\lambda\ge1$. Thus
$\mathsf{v}\in[\tfrac12\min(\sqrt\lambda,1),\sqrt\lambda]$.

These bounds give $\max(\lambda,\mathsf{v})\le\max(\lambda,\sqrt\lambda)$. In the other direction,
\begin{equation*}
\max(\lambda,\mathsf{v})\ge\max\bigl(\lambda,\tfrac12\min(\sqrt\lambda,1)\bigr)
\ge\tfrac12\max(\lambda,\sqrt\lambda).
\end{equation*}
For $\lambda\ge1$ this holds because the left side is $\lambda$. For $\lambda<1$ it holds because the
left side is at least $\tfrac12\sqrt\lambda$. Since
$1/\max(\lambda,\sqrt\lambda)=\min(\lambda^{-1},\lambda^{-1/2})$, we conclude that on mean-zero
scalars
\begin{equation*}
\tfrac12\min(\lambda^{-1},\lambda^{-1/2})
\le(1+\mathsf{C})(\lambda+\mathsf{v})^{-1}
\le4\min(\lambda^{-1},\lambda^{-1/2}) .
\end{equation*}
The minimum is $\lambda^{-1}$ for $\lambda\ge1$ (lattice scales) and $\lambda^{-1/2}$ for
$\lambda\le1$ (large scales).

\emph{Identification with the multiplier charge.} In the upward variational problem only the normal
plaquettes carry multipliers, namely those spanned by the normal direction $4$ and a lateral
direction $k$. The
multiplier field of the level at normal coordinate $y$ is therefore the level field
$\mu=(\mu_{4k}(\cdot,y))_{k=1,2,3}$. Reversing the orientation of a plaquette reverses the sign of its
multiplier, $\mu_{j4}=-\mu_{4j}$. Hence the multiplier charge of the level, as defined in the
Gauss-law lemma of the exit construction, is
\begin{equation*}
(\delta\mu)_4=\sum_j\partial_j^-\mu_{j4}=-\operatorname{div}^-(\mu_{4\cdot})=\nabla^*\mu=q .
\end{equation*}
A different sign convention for the charge replaces $q$ by $-q$. This leaves the quadratic form in
\eqref{8.14:eq-charge-energy-1} unchanged.

\emph{Consequence.} In \eqref{8.14:eq-concavity-absorption-a}, of the model statement on
concavity, absorption and exactness, the increment $\Delta m(a+1)$ is bounded by $\Delta m(a)$
minus the concavity defect $\tfrac12\langle P^{\perp}\mu,(1-\mathsf{C}^2)P^{\perp}\mu\rangle$. Here
$\mu$ is the multiplier of the top capped level of $P_{a+1}$. By \eqref{8.14:eq-charge-energy-1} and
the identification just made, this defect equals
$\tfrac12\|q^{(a+1)}\|^2_{(1+\mathsf{C})(\lambda+\mathsf{v})^{-1}}$. This is
\eqref{8.14:eq-charge-energy-2}.
\end{proof}

\begin{corollary}[Model: geometric decay on the charged class]\label{8.14:cor-charged-decay}
The model is that of the model theorem, Theorem~\ref{8.14:thm-exit-value}. It consists of the
truncated problems $P_a$ on level fields $\mathfrak{H}=\ell^2(\mathbb{T};\mathbb{R}^3)$ over the
lateral torus $\mathbb{T}$. The young caps form a set $\mathfrak{B}$, and $P_a$ has value $m(a)$,
increment $\Delta m(a)=m(a+1)-m(a)$ and one-capped-level defect $W=V_1-V_+$. The operators
$\lambda=-\Delta_{\parallel}$, $\mathsf{v}$ and $\mathsf{C}=1-\mathsf{v}$ are those fixed there.

Let $a_1<a_2$. Write $u_a=W(\beta^{(a)}_{a-1})$ for the top-level defect of $P_a$, as in the
model corollary, Corollary~\ref{8.14:cor-stopping-height}. Write $\mu^{(a)}$ for the multiplier of
the top capped level of $P_a$, and $q^{(a)}=\nabla^*\mu^{(a)}$ for its charge, as in the model
lemma, Lemma~\ref{8.14:lem-charge-energy}. Suppose that for $a_1\le a<a_2$ the one-sided
hypothesis
\begin{equation}\label{8.14:eq-charged-decay-a}
(\mathrm{CH})_{\theta_q,R_q}:\qquad
\tfrac12\bigl\langle q^{(a)},(1+\mathsf{C})(\lambda+\mathsf{v})^{-1}q^{(a)}\bigr\rangle
\;\ge\;\frac{\theta_q}{R_q}\,u_a
\end{equation}
holds for some $\theta_q\in(0,1]$ and some $R_q\ge1$. Then
\begin{equation}\label{8.14:eq-charged-decay-1}
u_{a+1}\le\Bigl(1-\frac{\theta_q}{R_q}\Bigr)u_a\qquad(a_1\le a<a_2).
\end{equation}
No sup bound on the multiplier is assumed, and so no $s_\infty$ enters. No Poisson deficit is
assumed either, so neither $\kappa_d$ nor $\kappa_1$ enters. No box structure is used:
$\mathfrak{B}$ may be any convex, translation-invariant set of young caps.
\end{corollary}

\begin{proof}
Fix $a$ with $a_1\le a<a_2$. The argument has four steps.

\emph{First step.} We use the lower member of the increment bounds
\eqref{8.14:eq-exit-value-c} of the model theorem, Theorem~\ref{8.14:thm-exit-value}, at the index
$a$. It gives $\Delta m(a)\ge W(\beta^{(a+1)}_a)$. The right-hand side is $u_{a+1}$ by definition,
so $u_{a+1}\le\Delta m(a)$.

\emph{Second step.} The upper member of the same bounds \eqref{8.14:eq-exit-value-c}, again at $a$,
gives $\Delta m(a)\le W(\beta^{(a)}_a)$.

\emph{Third step.} Put $\beta:=\beta^{(a)}_{a-1}$ and $\beta':=\beta^{(a)}_a$. By the statement
\eqref{8.14:eq-exit-value-b} of the model theorem, Theorem~\ref{8.14:thm-exit-value}, the step from
$\beta$ that is optimal for $V_1$ is the top capped level of $P_a$. Its multiplier is therefore
$\mu^{(a)}$, and the field it adds to $\beta$ carries $\beta$ to $\beta'$. Thus $\beta$ and $\beta'$
stand in the relation assumed in the concavity inequality \eqref{8.14:eq-concavity-absorption-a}
of the model statement on concavity, absorption and exactness. That inequality gives
\begin{equation*}
W(\beta^{(a)}_a)\le W(\beta^{(a)}_{a-1})
-\tfrac12\bigl\langle P^{\perp}\mu^{(a)},(1-\mathsf{C}^2)P^{\perp}\mu^{(a)}\bigr\rangle .
\end{equation*}
Here $W(\beta^{(a)}_{a-1})=u_a$ by definition.

\emph{Fourth step.} By the model lemma, Lemma~\ref{8.14:lem-charge-energy}, applied to
$\mu=\mu^{(a)}$ with $q=q^{(a)}$,
\begin{equation*}
\tfrac12\bigl\langle P^{\perp}\mu^{(a)},(1-\mathsf{C}^2)P^{\perp}\mu^{(a)}\bigr\rangle
=\tfrac12\bigl\langle q^{(a)},(1+\mathsf{C})(\lambda+\mathsf{v})^{-1}q^{(a)}\bigr\rangle .
\end{equation*}
By \eqref{8.14:eq-charged-decay-a}, the right-hand side is at least $(\theta_q/R_q)\,u_a$.

\emph{Conclusion.} Chaining the four steps gives
\begin{align*}
u_{a+1}\le\Delta m(a)\le W(\beta^{(a)}_a)
&\le u_a-\tfrac12\bigl\langle q^{(a)},(1+\mathsf{C})(\lambda+\mathsf{v})^{-1}q^{(a)}\bigr\rangle\\
&\le\Bigl(1-\frac{\theta_q}{R_q}\Bigr)u_a ,
\end{align*}
which is \eqref{8.14:eq-charged-decay-1}.
\end{proof}

\begin{remark}[Model: order without charge, absorption with charge]\label{8.14:rem-order-absorption}
The model is that of the model remark, Remark~\ref{8.14:rem-cutoff-in-height}, and of the model
lemma, Lemma~\ref{8.14:lem-charge-energy}. On the lateral torus
$\mathbb{T}=(\mathbb{Z}/N)^3$ we have the level fields $\mathfrak{H}=\ell^2(\mathbb{T};\mathbb{R}^3)$,
their transverse part $\mathfrak{H}_T$, the projection $P$ onto transverse fields and its
complement $P^{\perp}$, the lateral Laplacian $\lambda=-\Delta_{\parallel}$, and the operators
$\mathsf{v}$ and $\mathsf{C}=1-\mathsf{v}$ built from $\lambda$. A front climbs level by level
through the truncated problems $P_a$, whose caps $\mathfrak{B}$ are convex and
translation-invariant. For each $a$:
\begin{itemize}
\item $m(a)$ is the value of $P_a$ and $\Delta m(a)$ its increment;
\item $u_a$ is the top-level defect of $P_a$;
\item $\mu^{(a)}$ is the multiplier of the top capped level of $P_a$;
\item $q^{(a)}=\nabla^*\mu^{(a)}$ is its charge, the lateral co-differential.
\end{itemize}
All of this is as in the model lemma, Lemma~\ref{8.14:lem-charge-energy}.

For $\mu\in\mathfrak{H}$ we write $\bar\mu$ for its mean, which is a constant field. For
$q\in\ell^2(\mathbb{T};\mathbb{R})$ of mean zero we put
\begin{equation}\label{8.14:eq-order-absorption-1}
\|q\|^2_{(1+\mathsf{C})(\lambda+\mathsf{v})^{-1}}
:=\langle q,(1+\mathsf{C})(\lambda+\mathsf{v})^{-1}q\rangle .
\end{equation}
By Lemma~\ref{8.14:lem-charge-energy}, a model statement, the operator
$(1+\mathsf{C})(\lambda+\mathsf{v})^{-1}$ is a positive multiplier. Within absolute factors it is
comparable to $\min(\lambda^{-1},\lambda^{-1/2})$. The quantity
\eqref{8.14:eq-order-absorption-1} is therefore the $H^{-1}$-norm of $q$ at lattice scales and
its $H^{-1/2}$-norm at large scales. We call it the charge energy.

\emph{(a) Order without charge, absorption with charge.}
The dichotomy proposition for the axial potential shows that the scalar order structure exists
exactly on the charge-free class $q\equiv0$. By Lemma~\ref{8.14:lem-charge-energy}, a model
statement, the concavity defect that Jensen's inequality absorbs for free at each level is
exactly one half of the charge energy of that level. The two mechanisms available in the model
therefore divide the problem between them:
\begin{itemize}
\item order on $\ker\nabla^*$, where the charge-free transfer proposition carries the lift of
the model theorem on inactivity at full band width, with constants independent of the cutoff
index $n$ of the model remark, Remark~\ref{8.14:rem-cutoff-in-height};
\item absorption on $(\ker\nabla^*)^{\perp}$, at a rate set by the $H^{-1/2}$-energy of $q$.
\end{itemize}

The vanishing of the concavity defect when $P^{\perp}\mu$ is a constant field is the end
$q\equiv0$ of this division. Indeed, by the proof of the model lemma,
Lemma~\ref{8.14:lem-charge-energy}, $P^{\perp}\mu=\nabla\lambda^{-1}q+\bar\mu$, so $P^{\perp}\mu$
is constant exactly when $q=0$; and then the defect
$\tfrac12\langle q,(1+\mathsf{C})(\lambda+\mathsf{v})^{-1}q\rangle$ vanishes.
The other end is the statement of the charged-needle lemma: a laterally bounded active set of a
single orientation, a needle, is charged. The hypothesis \eqref{8.14:eq-charged-decay-a} of the
model corollary, Corollary~\ref{8.14:cor-charged-decay}, is the quantitative form of that
statement.

The following is a heuristic description and is used in no proof. Consider a multiplier field
whose charge lives at a lateral scale, written $R$ in this paragraph only. That is, the lateral
Fourier transform $\hat q(\xi)$ of $q$, with $\xi$ the lateral momentum in the dual torus of
$\mathbb{T}$, is concentrated where $|\xi|\asymp 1/R$, as for a single-orientation needle of
radius $R$ varying on that scale.
There $\lambda\asymp R^{-2}$, so $(1+\mathsf{C})(\lambda+\mathsf{v})^{-1}\asymp R$ and
$\lambda^{-1}\asymp R^2$. By the computation in the proof of the model lemma,
Lemma~\ref{8.14:lem-charge-energy}, $P^{\perp}\mu-\bar\mu=\nabla\lambda^{-1}q$, so
\begin{equation*}
\|P^{\perp}\mu-\bar\mu\|^2=\langle q,\lambda^{-1}q\rangle\asymp R^2\|q\|^2 .
\end{equation*}
Hence
\begin{equation*}
\langle q,(1+\mathsf{C})(\lambda+\mathsf{v})^{-1}q\rangle\asymp R\,\|q\|^2
\asymp \frac{\|P^{\perp}\mu-\bar\mu\|^2}{R}.
\end{equation*}
In this description $\theta_q\asymp\|P^{\perp}\mu-\bar\mu\|^2/\|\mu\|^2$, the charged fraction
of the mass of the multiplier, and $R_q\asymp R$.

The hypothesis \eqref{8.14:eq-charged-decay-a} is a hypothesis on the minimisers of the
truncated problems, which are not known explicitly. It is of the same kind as
\eqref{8.14:eq-stopping-height-b}. It is not a hypothesis of any clause of
Theorem~\ref{3.1:thm-main}.

\emph{(b) The cost of charge per level.}
The heuristic that charged needles cost value has an exact per-level form in the model. By
Lemma~\ref{8.14:lem-charge-energy}, a model statement,
\begin{equation}\label{8.14:eq-order-absorption-2}
\Delta m(a+1)\le \Delta m(a)-\tfrac12\,\|q^{(a+1)}\|^2_{(1+\mathsf{C})(\lambda+\mathsf{v})^{-1}} .
\end{equation}
Thus at every level that the front climbs with a charged top multiplier, the value increment
drops by at least one half of the charge energy of the new top level.

\emph{(c) What the decay on the charged class removes, and what it does not.}
Suppose \eqref{8.14:eq-charged-decay-a} holds for $a_1\le a<a_2$. The model corollary,
Corollary~\ref{8.14:cor-charged-decay}, then gives $u_{a+1}\le(1-\theta_q/R_q)\,u_a$ on that
range, with no sup bound on $\mu$. The decay rate is therefore $\theta_q/R_q$, and the e-folding
height is $\mathfrak{h}=R_q/\theta_q$. The bound \eqref{8.14:eq-cutoff-in-height-b} of the model
remark, Remark~\ref{8.14:rem-cutoff-in-height}, no longer enters. The e-folding height is now
the radius $R_q$ divided by the dimensionless charged fraction $\theta_q$. As a fraction of the
band height $W_y$ of the model remark, Remark~\ref{8.14:rem-cutoff-in-height}, it is independent
of $n$ whenever $R_q/n$ is. The count $\mathfrak{n}$ of
e-folds in \eqref{8.14:eq-cutoff-in-height-c} is unaffected.

\emph{(d) The missing bound.}
A proof free of the cutoff must bound the collar in the sup norm. The scalar descent does this on
the charge-free class. For $q\ne0$ it meets the term
\begin{equation*}
\mathcal{J}_k=\partial_k D-T_k,\qquad D(\cdot,y)=\sum_{y'<y}q(\cdot,y') .
\end{equation*}
Here $q(\cdot,y')$ is the charge of the level at height $y'$, so that $D(\cdot,y)$ is the charge
integrated up the column below height $y$; $\partial_k$ is the lateral difference in a lateral
direction $k$; and $T_k$ is the term of the scalar descent with which $\partial_k D$ is compared,
$\mathcal{J}_k$ being their difference. Thus $\partial_k D$ is the lateral
$k$-gradient of the column-integrated charge. By the remark on the size of the charge term, it is
pointwise as large as the signal of the descent.

The only mechanism in the model that pays for charge is the increment inequality
\eqref{8.14:eq-order-absorption-2}. It pays
\begin{equation*}
\sum_a\tfrac12\,\|q^{(a)}\|^2_{(1+\mathsf{C})(\lambda+\mathsf{v})^{-1}}
\le\sum_a\Delta m(a)=M_h ,
\end{equation*}
where the sums run over $h$ consecutive levels above a base level $a_0$ (the letters $h$ and
$a_0$ are used in this paragraph only), and $M_h=m(a_0+h)-m(a_0)$ is the value budget of these
levels, the sum of their increments.
This payment is in the $H^{-1/2}$-norm, and it is extensive. A lift with a charge error, that is,
a form of the lift referred to in (a) that admits a nonzero charge, would
require the sup norm of the Riesz potential $\partial_k\Delta_{\parallel}^{-1}D$ to be controlled
by this budget. That is again a comparison of energy with sup norm, now for the charge. This
bound is not proved.
\end{remark}

\begin{lemma}[Square-summability of the iterated Poisson kernel]\label{8.14:lem-kernel-square-sum}
Let $N\ge 1$ be an integer and let $\mathbb{T}=(\mathbb{Z}/N)^3$ be the lateral torus of side $N$ (as
in the model of Lemma~\ref{8.14:lem-poisson-identities}). Let $\lambda=-\Delta_{\parallel}$, where
$\Delta_{\parallel}u(x)=\sum_{k=1}^{3}[u(x+e_k)+u(x-e_k)-2u(x)]$ is the lateral lattice Laplacian on
$\ell^2(\mathbb{T};\mathbb{R})$, and $x\pm e_k$ denote the two neighbours of $x$ in the $k$-th
lateral coordinate direction. Put
\begin{equation*}
\mathsf{v}=\tfrac12\bigl[-\lambda+(\lambda^2+4\lambda)^{1/2}\bigr],\qquad \mathsf{C}=1-\mathsf{v},
\end{equation*}
and for $j\ge 1$ let $K_j$ be the convolution kernel of $\mathsf{C}^j$ on $\mathbb{T}$, so that
$\mathsf{C}^ju=K_j*u$. Then for $1\le J\le N$
\begin{equation}\label{8.14:eq-kernel-square-sum-1}
\|K_J\|^2_{\ell^2(\mathbb{T})}=K_{2J}(0)\le C_{\mathrm{hm}}\,J^{-3},
\qquad C_{\mathrm{hm}}:=\frac{185}{16}+27e^{-3}.
\end{equation}
\end{lemma}

\begin{proof}
\emph{The identity.} The operator $\mathsf{C}^J$ is a real function of the lateral Laplacian, which
commutes with the reflection $z\mapsto -z$ of $\mathbb{T}$; hence $K_J$ is real and even. Since
$\mathsf{C}^J\mathsf{C}^J=\mathsf{C}^{2J}$, we have $K_J*K_J=K_{2J}$. Therefore
\begin{equation*}
\sum_{z\in\mathbb{T}}K_J(z)^2=\sum_{z\in\mathbb{T}}K_J(z)K_J(-z)=(K_J*K_J)(0)=K_{2J}(0).
\end{equation*}
The characters $x\mapsto e^{2\pi i\xi\cdot x/N}$, $\xi=(\xi_1,\xi_2,\xi_3)$ with
$\xi_k\in(-N/2,N/2]\cap\mathbb{Z}$, diagonalise $-\Delta_{\parallel}$ with eigenvalues
\begin{equation*}
\lambda(\xi)=\sum_{k=1}^{3}4\sin^2(\pi\xi_k/N),
\end{equation*}
so the kernel of the multiplier $\mathsf{C}^{2J}$ evaluated at the origin is
\begin{equation}\label{8.14:eq-kernel-square-sum-2}
K_{2J}(0)=N^{-3}\sum_{\xi}\mathsf{C}(\lambda(\xi))^{2J}.
\end{equation}

\emph{Bound on each mode.} By the model lemma, Lemma~\ref{8.14:lem-poisson-identities}, one has
$\mathsf{C}=(1+\lambda+\mathsf{v})^{-1}$ (indeed
$(1-\mathsf{v})(1+\lambda+\mathsf{v})=1+\lambda-\lambda\mathsf{v}-\mathsf{v}^2=1$, because
$\mathsf{v}^2+\lambda\mathsf{v}=\lambda$ by the definition of $\mathsf{v}$ as the non-negative root), so
$\mathsf{C}\ge 0$. Since $1-t\le e^{-t}$ for every real $t$, also $\mathsf{C}=1-\mathsf{v}\le e^{-\mathsf{v}}$.
By the size estimate in the proof of the model lemma, Lemma~\ref{8.14:lem-charge-energy},
$\mathsf{v}\ge\frac12\min(\sqrt{\lambda},1)$; there it follows from $\mathsf{v}\le\sqrt{\lambda}$ (a
consequence of $\mathsf{v}^2+\lambda\mathsf{v}=\lambda$) and, for $\lambda>0$,
\begin{equation*}
\mathsf{v}=\frac{\lambda}{\lambda+\mathsf{v}}\ge\frac{\lambda}{\lambda+\sqrt{\lambda}}
=\frac{\sqrt{\lambda}}{1+\sqrt{\lambda}}\ge\tfrac12\min(\sqrt{\lambda},1),
\end{equation*}
while at $\lambda=0$ both $\mathsf{v}$ and $\frac12\min(\sqrt{\lambda},1)$ vanish.
Hence, mode by mode,
\begin{equation}\label{8.14:eq-kernel-square-sum-3}
\mathsf{C}^{2J}\le e^{-2J\mathsf{v}}\le e^{-J\min(\sqrt{\lambda},1)}\le e^{-J\sqrt{\lambda}}+e^{-J},
\end{equation}
the last step because $e^{-J\min(\sqrt{\lambda},1)}$ equals one of the two non-negative terms on the
right.

\emph{Lower bound on $\sqrt{\lambda}$.} Since $\sin\pi t\ge 2t$ for $t\in[0,\frac12]$ (concavity of
$\sin\pi t$ on this interval, with equality at both ends) and $|\xi_k|/N\le\frac12$, for each
$k\in\{1,2,3\}$ we have
\begin{equation*}
\sqrt{\lambda(\xi)}\ge 2\,|\sin(\pi\xi_k/N)|=2\sin(\pi|\xi_k|/N)\ge 4|\xi_k|/N,
\end{equation*}
and therefore $\sqrt{\lambda(\xi)}\ge 4|\xi|_\infty/N$, where $|\xi|_\infty=\max_k|\xi_k|$.

\emph{The sum over modes.} Put $\eta:=4J/N$. Regarding the set of modes $\xi$ as a subset of
$\mathbb{Z}^3$ and grouping by $m=|\xi|_\infty$, we obtain
\begin{align*}
\sum_{\xi}e^{-J\sqrt{\lambda(\xi)}}
&\le\sum_{m\ge 0}\#\{\xi\in\mathbb{Z}^3:|\xi|_\infty=m\}\,e^{-\eta m}
\le 1+\sum_{m\ge 1}(24m^2+2)e^{-\eta m}\\
&\le 1+26\sum_{m\ge 1}m^2e^{-\eta m}
\le 1+26\int_0^\infty(t+1)^2e^{-\eta t}\,dt\\
&=1+26\bigl(2\eta^{-3}+2\eta^{-2}+\eta^{-1}\bigr).
\end{align*}
Here the count uses $\#\{\xi\in\mathbb{Z}^3:|\xi|_\infty=m\}=(2m+1)^3-(2m-1)^3=24m^2+2$ for $m\ge 1$
and the value $1$ at $m=0$, and the next step uses $24m^2+2\le 26m^2$ for $m\ge 1$. The integral
comparison rests on the bound, for each $m\ge 1$,
\begin{equation*}
m^2e^{-\eta m}\le\int_{m-1}^{m}(t+1)^2e^{-\eta t}\,dt,
\end{equation*}
valid because $(t+1)^2\ge m^2$ and $e^{-\eta t}\ge e^{-\eta m}$ for $t\in[m-1,m]$. Summing this over
$m\ge 1$ gives the integral bound, and the last equality is the evaluation
$\int_0^\infty(t^2+2t+1)e^{-\eta t}\,dt=2\eta^{-3}+2\eta^{-2}+\eta^{-1}$.

\emph{Conclusion.} The term $e^{-J}$ in \eqref{8.14:eq-kernel-square-sum-3} contributes
$N^{-3}\cdot N^3e^{-J}=e^{-J}$ to \eqref{8.14:eq-kernel-square-sum-2}, since there are $N^3$ modes.
Inserting $\eta^{-3}=N^3/(64J^3)$, $\eta^{-2}=N^2/(16J^2)$, $\eta^{-1}=N/(4J)$ gives
\begin{align*}
K_{2J}(0)&\le N^{-3}+26\Bigl[\frac{2N^3}{64J^3}+\frac{2N^2}{16J^2}+\frac{N}{4J}\Bigr]N^{-3}+e^{-J}\\
&=N^{-3}+\frac{13}{16}J^{-3}+\frac{13}{4}N^{-1}J^{-2}+\frac{13}{2}N^{-2}J^{-1}+e^{-J}\\
&\le\Bigl[1+\frac{13}{16}+\frac{13}{4}+\frac{13}{2}\Bigr]J^{-3}+e^{-J}
=\frac{185}{16}J^{-3}+e^{-J},
\end{align*}
where the inequality uses $J\le N$, so that $N^{-3}$, $N^{-1}J^{-2}$ and $N^{-2}J^{-1}$ are each at most
$J^{-3}$. Finally, the function $t\mapsto t^3e^{-t}$ on $[0,\infty)$ has derivative $(3t^2-t^3)e^{-t}$
and attains its maximum at $t=3$, so $J^3e^{-J}\le 27e^{-3}$, that is, $e^{-J}\le 27e^{-3}J^{-3}$.
Together with the identity proved first, this gives \eqref{8.14:eq-kernel-square-sum-1} with
$C_{\mathrm{hm}}=\frac{185}{16}+27e^{-3}$.
\end{proof}

\begin{theorem}[Model: exactness beyond a collar]\label{8.14:thm-collar-exactness}
The model is the variational model of the model exit-value theorem,
Theorem~\ref{8.14:thm-exit-value}. Its level fields are the elements $E=(E_k(x))$, with
$x\in\mathbb{T}$ and $k\in\{1,2,3\}$, of $\mathfrak{H}=\ell^2(\mathbb{T};\mathbb{R}^3)$
on the lateral torus $\mathbb{T}=(\mathbb{Z}/N)^3$. The one-step Poisson operator
$\mathsf{C}$ acts componentwise as convolution with its kernel $K_1$, and $K_J$ is the
kernel of $\mathsf{C}^J$. We put $C_N=\max\mathsf{C}^{-1}$. The young cap is the box
$\mathfrak{B}=\{E:\ |E_k(x)|\le\tau_k\ \text{for all }x,k\}$. For each $a$, $P_a$ is the
truncated problem, $m(a)$ is its value and $X_a$ its minimiser, and $W=V_1-V_+$. The
constant $C_{\mathrm{hm}}=185/16+27e^{-3}$ is that of Lemma~\ref{8.14:lem-kernel-square-sum}.

Let $a>a_0$ and $1\le J\le N$. Let $(E,\mu)$ be the field and the multiplier of the top
capped level of $P_a$. Let $u_a=W(\beta^{(a)}_{a-1})$, and let the scale-$J$ Poisson
deficit of the top capped level be
\begin{equation}\label{8.14:eq-collar-exactness-a}
\kappa_J(a):=\min_{x,k,\pm}\ \sum_{x'\in\mathbb{T}}K_J(x-x')\bigl[\tau_k\mp E_k(x')\bigr].
\end{equation}
Suppose that
\begin{equation}\label{8.14:eq-collar-exactness-1}
C_{\mathrm{hm}}^{1/2}J^{-3/2}\,(2C_Nu_a)^{1/2}\le\kappa_J(a).
\end{equation}
This holds in particular if
\begin{equation}\label{8.14:eq-collar-exactness-2}
C_{\mathrm{hm}}^{1/2}J^{-3/2}
\Bigl[\frac{2C_N\,(m(a)-m(a_0))}{a-a_0}\Bigr]^{1/2}\le\kappa_J(a).
\end{equation}
Then $X_a$ satisfies the young cap at every level $\ge a+J-1$. Consequently $X_a$ is the
minimiser of every problem obtained from $P_a$ by capping any set of levels $\ge a+J-1$,
and those caps carry no multiplier there.
\end{theorem}

\begin{proof}
\emph{The free tail.} By part (b) of the model exit-value theorem,
Theorem~\ref{8.14:thm-exit-value}, the tail of $P_a$ above row $a-1$ consists of the top
capped level and free levels. The exit pull is $e_a=\mathsf{C}(E+\mu)$, and the free tail
of $X_a$ is given by \eqref{8.14:eq-exit-value-b}:
\begin{equation}\label{8.14:eq-collar-exactness-3}
e_{a+j}=\mathsf{C}^j e_a=\mathsf{C}^{j+1}(E+\mu),\qquad j\ge0.
\end{equation}
Satisfying the young cap at level $a+j$ means $e_{a+j}\in\mathfrak{B}$.

\emph{Propagation upward.} The same part of
Theorem~\ref{8.14:thm-exit-value} records that $K_1\ge0$ and
$\sum_{x'}K_1(x-x')=1$. Hence for $F\in\mathfrak{B}$ we have
\[
|(\mathsf{C}F)_k(x)|\le\sum_{x'}K_1(x-x')\,|F_k(x')|\le\tau_k ,
\]
so $\mathsf{C}$ maps $\mathfrak{B}$ into itself. By \eqref{8.14:eq-collar-exactness-3},
$e_{a+j}=\mathsf{C}^{\,j-J+1}e_{a+J-1}$ for $j\ge J-1$. Therefore, if
$e_{a+J-1}\in\mathfrak{B}$, then $e_{a+j}\in\mathfrak{B}$ for all $j\ge J-1$. It remains
to show that $e_{a+J-1}\in\mathfrak{B}$.

\emph{The field part.} The kernel $K_J$ is the $J$-fold convolution of $K_1$. It is
therefore non-negative, and $\sum_{x'}K_J(x-x')=1$. We argue as in the box criterion of the
model theorem on concavity, absorption and exactness (see
\eqref{8.14:eq-concavity-absorption-c}), now at scale $J$. By
\eqref{8.14:eq-collar-exactness-3} with $j=J-1$,
\[
\pm(e_{a+J-1})_k(x)=\pm(\mathsf{C}^JE)_k(x)\pm(\mathsf{C}^J\mu)_k(x).
\]
Since $K_J$ has total mass one,
\[
\sum_{x'}K_J(x-x')\bigl[\tau_k\mp E_k(x')\bigr]=\tau_k\mp(\mathsf{C}^JE)_k(x).
\]
Since $K_J\ge0$, we also have $\pm(\mathsf{C}^J\mu)_k(x)\le(\mathsf{C}^J|\mu_k|)(x)$. Hence
\begin{equation}\label{8.14:eq-collar-exactness-4}
\pm(e_{a+J-1})_k(x)\le\tau_k-\sum_{x'}K_J(x-x')\bigl[\tau_k\mp E_k(x')\bigr]
+(\mathsf{C}^J|\mu_k|)(x).
\end{equation}

\emph{The multiplier part.} By the Cauchy--Schwarz inequality, and since the
$\ell^2(\mathbb{T})$-norm of $x'\mapsto K_J(x-x')$ equals $\|K_J\|_2$,
\[
(\mathsf{C}^J|\mu_k|)(x)=\sum_{x'}K_J(x-x')\,|\mu_k(x')|
\le\|K_J\|_2\,\|\mu_k\|_2 .
\]
Lemma~\ref{8.14:lem-kernel-square-sum} gives $\|K_J\|_2^2\le C_{\mathrm{hm}}J^{-3}$ for
$1\le J\le N$, and $\|\mu_k\|_2\le\|\mu\|_2$. By part (b) of
Theorem~\ref{8.14:thm-exit-value}, the top capped level of $P_a$ is the $V_1$-optimal step
from $\beta^{(a)}_{a-1}$, and $\mu$ is its multiplier. The lower bound in
\eqref{8.14:eq-exit-value-a} then gives $\|\mu\|_2^2\le 2C_N\,W(\beta^{(a)}_{a-1})=2C_Nu_a$.
Altogether,
\begin{equation}\label{8.14:eq-collar-exactness-5}
(\mathsf{C}^J|\mu_k|)(x)\le C_{\mathrm{hm}}^{1/2}J^{-3/2}\|\mu\|_2
\le C_{\mathrm{hm}}^{1/2}J^{-3/2}(2C_Nu_a)^{1/2}.
\end{equation}

\emph{Conclusion of the cap.} By \eqref{8.14:eq-collar-exactness-a}, the sum in
\eqref{8.14:eq-collar-exactness-4} is at least $\kappa_J(a)$ for every $x$, $k$ and choice
of sign. Inserting \eqref{8.14:eq-collar-exactness-5} gives
\[
\pm(e_{a+J-1})_k(x)\le\tau_k-\kappa_J(a)+C_{\mathrm{hm}}^{1/2}J^{-3/2}(2C_Nu_a)^{1/2}
\le\tau_k ,
\]
where the last step is hypothesis \eqref{8.14:eq-collar-exactness-1}. Thus
$e_{a+J-1}\in\mathfrak{B}$, and by the propagation step $X_a$ satisfies the young cap at
every level $\ge a+J-1$.

\emph{The second hypothesis.} Since $a>a_0$, the model corollary on monotone increments and
a stopping height, Corollary~\ref{8.14:cor-stopping-height}, gives in
\eqref{8.14:eq-stopping-height-a} the bound
\[
u_a\le\frac{m(a)-m(a_0)}{a-a_0}.
\]
Since $t\mapsto t^{1/2}$ is increasing, \eqref{8.14:eq-collar-exactness-2} implies
\eqref{8.14:eq-collar-exactness-1}.

\emph{The consequence.} Consider a problem obtained from $P_a$ by capping a set of levels
$\ge a+J-1$. By what has been shown, $X_a$ is feasible for this enlarged problem. It is
optimal for the smaller problem $P_a$, whose feasible set contains that of the enlarged
problem. Hence $X_a$ is the minimiser of the enlarged problem, and the added caps carry no
multiplier there.
\end{proof}

\begin{corollary}[Model: smallness of the exit over-pull]\label{8.14:cor-over-pull-small}
The setting is the model of Theorem~\ref{8.14:thm-exit-value} (the model theorem) and its notation.
Its level fields form $\mathfrak{H}=\ell^2(\mathbb{T};\mathbb{R}^3)$ over the lateral torus
$\mathbb{T}=(\mathbb{Z}/N)^3$, and $N(\cdot)$ is the quadratic defect form
\begin{equation*}
N(w)=\tfrac12\langle Pw,\mathsf{C}^{-1}Pw\rangle+\tfrac12\|P^{\perp}w\|^2 .
\end{equation*}
The young caps form the set $\mathfrak{B}$, which is closed, convex and invariant under lateral
translations; when $\mathfrak{B}$ is the box, further statements are made below. The truncated
problems $P_a$ have values
$m(a)$, exit pulls $e_a$ and free tails $e_{a+j}$. We write $\|\cdot\|_2$ and $\|\cdot\|_\infty$ for
the $\ell^2$ and supremum norms over $(x,k)$, and $\operatorname{dist}_2$,
$\operatorname{dist}_{\ell^\infty}$ for the corresponding distances to $\mathfrak{B}$.

Let $a>a_0$. Let $\pi_{\mathfrak{B}}$ denote the $\ell^2$-nearest-point map onto $\mathfrak{B}$; for
the box it is componentwise clipping. Put $S_a:=e_a-\pi_{\mathfrak{B}}(e_a)$, the exit over-pull. Let
$u_a=W(\beta^{(a)}_{a-1})$ be the top-level defect of the model corollary,
Corollary~\ref{8.14:cor-stopping-height}, and let
$C_{\mathrm{hm}}=185/16+27e^{-3}$.
\begin{enumerate}
\item[(i)] We have
\begin{equation}\label{8.14:eq-over-pull-small-1}
\|S_a\|_2^2\le 2W(\beta^{(a)}_a)\le 2u_a\le \frac{2\,(m(a)-m(a_0))}{a-a_0}.
\end{equation}
\item[(ii)] For every integer $j\ge1$,
$\operatorname{dist}_{\ell^\infty}(e_{a+j},\mathfrak{B})\le\|\mathsf{C}^jS_a\|_\infty$, and for
$1\le j\le N$ moreover
\begin{equation}\label{8.14:eq-over-pull-small-2}
\|\mathsf{C}^jS_a\|_\infty
\le C_{\mathrm{hm}}^{1/2}\,j^{-3/2}\Bigl[\frac{2\,(m(a)-m(a_0))}{a-a_0}\Bigr]^{1/2}.
\end{equation}
For the box, $\operatorname{dist}_{\ell^\infty}(e_{a+j},\mathfrak{B})$ is the largest amount by which
a component of $e_{a+j}$ exceeds its cap in absolute value.
\end{enumerate}
\end{corollary}

\begin{proof}
(i) Recall that $\mathsf{C}^{-1}=1+\lambda+\mathsf{v}$, where $\lambda=-\Delta_{\parallel}\ge0$ and
$\mathsf{v}\ge0$. Hence $\mathsf{C}^{-1}\ge1$ as an operator, and
$\langle Pw,\mathsf{C}^{-1}Pw\rangle\ge\|Pw\|^2$. Since $\|Pw\|^2+\|P^{\perp}w\|^2=\|w\|^2$, we
obtain $N(w)\ge\frac12\|w\|^2$ for every $w\in\mathfrak{H}$.

We next identify $W(\beta^{(a)}_a)$. By the model theorem, Theorem~\ref{8.14:thm-exit-value},
display~\eqref{8.14:eq-exit-value-a} gives
$W(\beta)=\min_{E\in\mathfrak{B}}N(d_0(\beta)-E)$. By \eqref{8.14:eq-exit-value-b} of the same model
theorem, the exit pull is $e_a=d_0(\beta^{(a)}_a)$. Therefore
\begin{equation*}
W(\beta^{(a)}_a)=\min_{E\in\mathfrak{B}}N(e_a-E)\ge\tfrac12\min_{E\in\mathfrak{B}}\|e_a-E\|_2^2
=\tfrac12\operatorname{dist}_2(e_a,\mathfrak{B})^2=\tfrac12\|S_a\|_2^2 .
\end{equation*}
The last equality holds because $\pi_{\mathfrak{B}}(e_a)$ is the $\ell^2$-nearest point of the closed
convex set $\mathfrak{B}$ to $e_a$. This is the first inequality of
\eqref{8.14:eq-over-pull-small-1}.

For the second inequality, we apply the concavity inequality \eqref{8.14:eq-concavity-absorption-a}
of the model statement~\ref{8.14:prf-concavity-absorption} (concavity, absorption and exactness)
with $\beta:=\beta^{(a)}_{a-1}$. It
applies here because, by \eqref{8.14:eq-exit-value-b} of the model theorem,
Theorem~\ref{8.14:thm-exit-value}, the $V_1$-optimal step from $\beta^{(a)}_{a-1}$ is the top capped
level of $P_a$. The configuration reached after that step is therefore $\beta'=\beta^{(a)}_a$. The
inequality gives
\begin{equation*}
W(\beta^{(a)}_a)\le W(\beta^{(a)}_{a-1})
-\tfrac12\langle P^{\perp}\mu,(1-\mathsf{C}^2)P^{\perp}\mu\rangle ,
\end{equation*}
with $\mu$ the multiplier of that level. Since $\mathsf{C}=(1+\lambda+\mathsf{v})^{-1}$ satisfies
$0<\mathsf{C}\le1$, the operator $1-\mathsf{C}^2$ is non-negative. Hence
$W(\beta^{(a)}_a)\le W(\beta^{(a)}_{a-1})=u_a$.

The third inequality is $u_a\le(m(a)-m(a_0))/(a-a_0)$ for $a>a_0$. This is
\eqref{8.14:eq-stopping-height-a} of the model corollary,
Corollary~\ref{8.14:cor-stopping-height}.

(ii) By \eqref{8.14:eq-exit-value-b} of the model theorem, Theorem~\ref{8.14:thm-exit-value}, the
free tail of $P_a$ is $e_{a+j}=\mathsf{C}^je_a$ for $j\ge0$. Splitting $e_a$ gives
\begin{equation*}
e_{a+j}=\mathsf{C}^j\pi_{\mathfrak{B}}(e_a)+\mathsf{C}^jS_a .
\end{equation*}
We check that $\mathsf{C}^j\pi_{\mathfrak{B}}(e_a)\in\mathfrak{B}$. The operator $\mathsf{C}$ acts
by convolution in the lateral variable with the kernel $K_1$. This kernel is non-negative and has
total mass one (the facts behind the Jensen inequality stated in the same model theorem). Hence
$\mathsf{C}f$
is a convex combination of lateral translates of $f$. Since $\mathfrak{B}$ is convex, closed and
translation-invariant, $\mathsf{C}$ maps $\mathfrak{B}$ into itself, and so does $\mathsf{C}^j$.
Therefore
\begin{equation*}
\operatorname{dist}_{\ell^\infty}(e_{a+j},\mathfrak{B})
\le\|e_{a+j}-\mathsf{C}^j\pi_{\mathfrak{B}}(e_a)\|_\infty=\|\mathsf{C}^jS_a\|_\infty .
\end{equation*}

Next we bound the right-hand side. For each $(x,k)$ we have
$(\mathsf{C}^jS_a)_k(x)=\sum_{x'}K_j(x-x')(S_a)_k(x')$. The Cauchy--Schwarz inequality gives
\begin{equation*}
|(\mathsf{C}^jS_a)_k(x)|\le\|K_j\|_{\ell^2(\mathbb{T})}\,\|(S_a)_k\|_2
\le\|K_j\|_{\ell^2(\mathbb{T})}\,\|S_a\|_2 .
\end{equation*}
For $1\le j\le N$, the square-summability of the iterated Poisson kernel
(Lemma~\ref{8.14:lem-kernel-square-sum}) gives
$\|K_j\|_{\ell^2(\mathbb{T})}\le C_{\mathrm{hm}}^{1/2}j^{-3/2}$. Inserting the bound on $\|S_a\|_2$
from \eqref{8.14:eq-over-pull-small-1} yields \eqref{8.14:eq-over-pull-small-2}.

Finally, for the box, componentwise clipping is a nearest point in $\ell^\infty$ as well. Hence
$\operatorname{dist}_{\ell^\infty}(e_{a+j},\mathfrak{B})$ is the largest overshoot of a component
over its cap in absolute value.
\end{proof}

\begin{remark}[Model: cutoff-free bounds outside the collar]\label{8.14:rem-outside-collar}
The model is the one of the model Theorem~\ref{8.14:thm-collar-exactness} and the model
Corollary~\ref{8.14:cor-over-pull-small}. Its ingredients are:
\begin{itemize}
\item level fields in $\mathfrak{H}=\ell^2(\mathbb{T};\mathbb{R}^3)$ on the lateral torus
$\mathbb{T}=(\mathbb{Z}/N)^3$;
\item the young caps $\mathfrak{B}$, the box with caps $\tau_k$;
\item the one-step Poisson operator $\mathsf{C}=(1+\lambda+\mathsf{v})^{-1}$ with kernel $K_1$,
and the kernels $K_j$ of $\mathsf{C}^j$;
\item for $a>a_0$, the truncated problem $P_a$, whose levels from $a$ on are free, with value
$m(a)$, minimiser $X_a$, exit pull $e_a$ and free tail $e_{a+j}=\mathsf{C}^je_a$;
\item the exit over-pull $S_a=e_a-\pi_{\mathfrak{B}}(e_a)$.
\end{itemize}
All statements below concern this model.

We read its quantities in the dictionary that expresses them through one parameter $n$, which
carries the whole dependence on the cutoff.
A quantity is called \emph{free of $n$}, or \emph{cutoff-free}, if, read in this dictionary, it
is bounded independently of $n$. The dictionary takes
\begin{equation*}
N=2P_{\mathrm{bl}}n,\qquad J=\varpi' n,\qquad a-a_0=\varpi'' n ,
\end{equation*}
and, with $M_h=m(a_0+h)-m(a_0)$ the value budget of $h$ levels above $a_0$ and
$\mathfrak{e}_h$ an energy budget of the kind given by the energy proposition accompanying the
model theorem on inactivity at full band width,
\begin{equation*}
m(a)-m(a_0)\le M_h\le\mathfrak{e}_h
\asymp 48P_{\mathrm{bl}}^3n^3\cdot\beta_{\mathrm{hot}}n\cdot
\tau_y^2\bigl(s_e\bar\lambda/\mathsf{T}\bigr)^2 .
\end{equation*}
Here $P_{\mathrm{bl}}$ is the number, free of $n$, for which $N=2P_{\mathrm{bl}}n$;
$\varpi',\varpi''$ are free of $n$; $\tau_y$ is the young cap; $s_e$ and $\bar\lambda$ are
constants free of $n$ entering the energy budget; and $\beta_{\mathrm{hot}}$ is the fraction of
levels, free of $n$, for which $\beta_{\mathrm{hot}}n$ is the number of levels counted by the
budget, with $\beta_{\mathrm{hot}}^{1/2}$ proportional to
$e^{-c\mathsf{T}/2}$ for a constant $c>0$ free of $n$. These relations are descriptions. No
estimate of the model depends on
them; only the numerical readings below use them, with $\mathfrak{e}_h$ replaced by the
displayed expression.

First, freedom from $n$ beyond the collar. By the model Corollary~\ref{8.14:cor-over-pull-small},
part (ii), for $J\le j\le N$,
\begin{equation*}
\operatorname{dist}_{\ell^\infty}(e_{a+j},\mathfrak{B})
\le C_{\mathrm{hm}}^{1/2}j^{-3/2}\Bigl[\frac{2(m(a)-m(a_0))}{a-a_0}\Bigr]^{1/2}
\le C_{\mathrm{hm}}^{1/2}J^{-3/2}\Bigl[\frac{2(m(a)-m(a_0))}{a-a_0}\Bigr]^{1/2}.
\end{equation*}
In the dictionary the square of the right-hand side is at most
\begin{equation*}
C_{\mathrm{hm}}(\varpi'n)^{-3}\cdot
\frac{96P_{\mathrm{bl}}^3n^4\beta_{\mathrm{hot}}\tau_y^2(s_e\bar\lambda/\mathsf{T})^2}{\varpi''n}
=\tau_y^2\,\frac{96C_{\mathrm{hm}}P_{\mathrm{bl}}^3\beta_{\mathrm{hot}}
(s_e\bar\lambda/\mathsf{T})^2}{\varpi'^3\varpi''} .
\end{equation*}
Thus, beyond a collar of $\varpi'n$ fine levels above any truncation at height at least
$\varpi''n$ into the band, the truncated minimiser overshoots the young caps in the supremum
norm by at most
\begin{equation}\label{8.14:eq-outside-collar-1}
\tau_y\Bigl[\frac{96C_{\mathrm{hm}}P_{\mathrm{bl}}^3\beta_{\mathrm{hot}}
(s_e\bar\lambda/\mathsf{T})^2}{\varpi'^3\varpi''}\Bigr]^{1/2}.
\end{equation}
This bound is free of $n$, because the powers $n^3\cdot n\cdot n^{-3}\cdot n^{-1}$ cancel. It
is small through the factor $\beta_{\mathrm{hot}}^{1/2}\propto e^{-c\mathsf{T}/2}$.

The root-mean-square overshoot per site at the exit level, averaged over the $3N^3$ components
of a level field, is also free of $n$. Part (i) of the
same corollary gives
\begin{equation*}
\|S_a\|_2^2\le\frac{2(m(a)-m(a_0))}{a-a_0}
\le\frac{96P_{\mathrm{bl}}^3n^3\beta_{\mathrm{hot}}\tau_y^2(s_e\bar\lambda/\mathsf{T})^2}{\varpi''} .
\end{equation*}
Since $3N^3=24P_{\mathrm{bl}}^3n^3$, it follows that
\begin{equation*}
\Bigl(\frac{\|S_a\|_2^2}{3N^3}\Bigr)^{1/2}
\le\tau_y\Bigl[\frac{4\beta_{\mathrm{hot}}(s_e\bar\lambda/\mathsf{T})^2}{\varpi''}\Bigr]^{1/2}.
\end{equation*}

Second, exactness beyond the collar. Let $\tau_{\min}$ be the smallest of the caps $\tau_k$,
let $0<\theta_J\le1$, and let $\kappa_J(a)$ be the scale-$J$ Poisson deficit
\eqref{8.14:eq-collar-exactness-a} of the top capped level of $P_a$. Consider the hypothesis
\begin{equation}\label{8.14:eq-outside-collar-a}
(\mathrm{NS})_J:\qquad \kappa_J(a)\ge\theta_J\tau_{\min}.
\end{equation}
It says that the top capped level of $P_a$ is not saturated, in Poisson mean at scale $J$,
beyond the fraction $1-\theta_J$. It is a one-sided hypothesis on the minimiser, which is not
known, of a problem of the kind treated in part (b) of the model theorem on inactivity at full
band width. It excludes a saturated patch of lateral radius of order $J$ or larger; in the
description by a saturated patch of lateral radius $R_*$, it corresponds to $R_*<J$.

Under \eqref{8.14:eq-outside-collar-a} the condition of the model
Theorem~\ref{8.14:thm-collar-exactness} reduces to a floor on $a-a_0$; when that floor holds,
the overshoot beyond the collar is exactly zero. The second sufficient condition of that
theorem is
\begin{equation*}
C_{\mathrm{hm}}^{1/2}J^{-3/2}
\Bigl[\frac{2C_N(m(a)-m(a_0))}{a-a_0}\Bigr]^{1/2}\le\kappa_J(a).
\end{equation*}
Since $m(a)-m(a_0)\le M_h$, this condition holds when
\begin{equation*}
C_{\mathrm{hm}}^{1/2}J^{-3/2}\Bigl[\frac{2C_NM_h}{a-a_0}\Bigr]^{1/2}\le\theta_J\tau_{\min},
\end{equation*}
that is, when
\begin{equation}\label{8.14:eq-outside-collar-2}
a-a_0\ge\frac{2C_NC_{\mathrm{hm}}M_h}{\theta_J^2\tau_{\min}^2J^3}.
\end{equation}
In the dictionary all caps equal the young cap $\tau_y$, so that $\tau_{\min}=\tau_y$, and
\eqref{8.14:eq-outside-collar-2} reads
\begin{equation*}
a-a_0\ge\Bigl[\frac{96C_NC_{\mathrm{hm}}P_{\mathrm{bl}}^3\beta_{\mathrm{hot}}
(s_e\bar\lambda/\mathsf{T})^2}{\theta_J^2\varpi'^3}\Bigr]\cdot n .
\end{equation*}
Once the truncation is $\varpi''n$ levels into the band, this becomes the floor
\begin{equation}\label{8.14:eq-outside-collar-3}
\frac{96C_NC_{\mathrm{hm}}P_{\mathrm{bl}}^3\beta_{\mathrm{hot}}
(s_e\bar\lambda/\mathsf{T})^2}{\theta_J^2\varpi'^3}\le\varpi'' ,
\end{equation}
which is free of $n$ and small with $\beta_{\mathrm{hot}}$. The index $n$ cancels exactly as in
the cutoff-free corollary to the model theorem on inactivity at full band width: the $n^4$ from
$M_h$ cancels against the $n\cdot n^3$ from $(a-a_0)J^3$.

The readings \eqref{8.14:eq-outside-collar-1}--\eqref{8.14:eq-outside-collar-3} use only the
items (a)--(d) of the model statement on concavity, absorption and exactness, the statement
containing \eqref{8.14:eq-concavity-absorption-c}, and the increment
inequality
\eqref{8.14:eq-stopping-height-a} of the model Corollary~\ref{8.14:cor-stopping-height}. They
use neither the absorption statement nor the deficit $\kappa_d$, the supremum bound
$s_\infty$ or the decay rate $\lambda_*$.

Third, where the cutoff sits. The cutoff therefore sits in the collar $[a,a+J-1)$, that is, in
the first $J-1$ free levels above the truncation, a number proportional to $n$. There
the only controls are the bound of a supremum by an $\ell^2$ norm,
\begin{equation*}
\operatorname{dist}_{\ell^\infty}(e_{a+j},\mathfrak{B})\le\|S_a\|_2 ,
\end{equation*}
where, by the bound on $\|S_a\|_2^2$ above, $\|S_a\|_2$ is at most $n^{3/2}$ times a factor free
of $n$, and the pointwise criterion \eqref{8.14:eq-concavity-absorption-c}
of the model statement on concavity, absorption and exactness. That criterion carries the two
factors entering the stopping height $H_{\mathrm{EV}}=\mathfrak{h}\mathfrak{n}+1$ of the model
Remark~\ref{8.14:rem-cutoff-in-height}: the supremum bound $s_\infty$ on the top multiplier,
through $\mathfrak{h}$, and the count $\mathfrak{n}$ of e-folds.

This is the obstruction met in three further places:
\begin{itemize}
\item the collar transfer, which needs a bound in the supremum norm on the $r_1$-collar of a
release surface lying inside a tested band;
\item the requirement that one closed shell of windows be exactly inactive at the canonical
comparison point $\hat{x}$;
\item the loaded form of the collar transfer.
\end{itemize}
The collar of the truncated problem is that shell. The model bounds of this section make every
statement beyond the collar free of $n$, and no statement inside it.

The collar statement does not by itself provide the input of the ladder argument over window
families. A truncated minimiser that satisfies every window beyond a gap, but not the window of
the gap's own row, does not lie in the segment set $\mathcal{S}$, for two reasons: the
membership lemma requires exact membership, and, by part (ii) of the exit proposition, a
centre that violates a window by at least $\nu_\flat$ leaves the lower part of the skin empty.
The exception is the case in which the violation at the gap is small enough for the tolerance
inequality of the ladder argument.
\end{remark}

\begin{proposition}[Truncated minimisers as admissible sandwich centres]
\label{8.14:prop-admissible-centres}
We use the objects of the sandwich theorem of this chapter: the chart at $e$ with its action
$\mathbf{A}(e,\cdot)$ and its regular domain; the family $W(e,\cdot)$ of windows, with the seed
window $f_{\square}$, the other windows being called non-seed windows; the set $S(e)$ of points
at which every window of $W(e,\cdot)$ reads less than $1$, and its closure
$\operatorname{cl}S(e)$; the segment set $\mathcal{S}$; the event set $\mathcal{E}$; the shift
map $\Xi$ and the map $\mathcal{M}$ carrying points of the chart at $e$ to points of the chart
at $\Xi e$; the numerator $\operatorname{Num}(e)$ and the denominator
$\operatorname{Den}(\Xi e)$; the function $\Psi$ whose supremum enters the bound on the
numerator; the non-seed threshold $\theta$ and the constants $\bar q$, $t_1$, $\bar e$,
$\Gamma_{\square}$, $\Gamma^{\kappa}$, $C'$, $\varepsilon$, $\mathfrak{d}$; and the sandwich
constant $\epsilon_{\kappa}$. Assume the following.
\begin{itemize}
\item The sandwich theorem, one of whose two inequalities, called below the upper sandwich
inequality, asserts $\operatorname{Num}(e)\le\epsilon_{\kappa}\operatorname{Den}(\Xi e)$ at its
minimising centre $x_0$.
\item The scaffold comparison in all zones.
\item The near-convexity lemma for windows, with its part~(i), which, along the path
$s\mapsto y\exp(s\log(y^{-1}x))$ from a point $y$ at which a window reads at most $b$ to a point
$x$ at which it reads at most $1$, bounds its reading by $(1-s)b+s+2s(1-s)\bar q$, and its
part~(ii), which bounds the Haar measure of the contraction of a set about a centre.
\item The second-order lemma, the path lemma, and the two further lemmas of the exit-value
construction used with them in the second-order step of the sandwich theorem.
\item The nine bounds for the shifted chart, among them the window bound used in the first
claim of the localised volume proposition, which carries a bound on the reading of a non-seed
window at $(e,y)$ to a bound on its reading at $(\Xi e,\mathcal{M}y)$, and the bound on the Haar
measure of the shifted contraction.
\item The basic identity of the exit-value construction and its fourth input bound, in the forms
used in the proof of the sandwich theorem.
\item The three claims of the localised volume proposition.
\item The segment lemma, in the form in which its assertion on the segment issuing from a
minimising centre is read below for a family of windows.
\end{itemize}
All of these are taken with the seed-window margin $\kappa$ at least the floor
$\kappa_{\flat}$ of the sandwich theorem, and under the following ceilings on the parameters, in
the sense of Notation~\ref{2.3:not-stages-first}: the two ceilings on $\gamma$ at the margin
$\kappa$ under which the sandwich theorem is stated, the second of which is the smallness
condition on the action bound used in the third claim of the localised volume proposition; the
ceiling on $\gamma$ of the scaffold comparison; the $F$-, $J$- and $\Lambda$-ceilings on
$\gamma$; the near-zone condition; and the growth ceiling on $\gamma$.

Let $\mathcal{W}'$ be any family of non-seed windows of $W(e,\cdot)$, each window
$f\in\mathcal{W}'$ being imposed at a threshold $\theta_f$ with
$1+2\bar q\le\theta_f\le\theta$ (for instance $\widehat{\mathcal{W}}_n$, the base window family
$\mathfrak{S}$ together with the windows of index at most $n$ at the graded thresholds, or
$\mathcal{W}^{\ge i}$), so that the set
$\{\mathcal{W}'\text{ holds}\}$ of points at which every window of $\mathcal{W}'$ reads at most
its threshold satisfies
$\{\mathcal{W}'\text{ holds}\}\supseteq\mathcal{S}\supseteq\operatorname{cl}S(e)$.
Let $X\in\mathsf{M}(\mathcal{W}')$, a minimiser of $\mathbf{A}(e,\cdot)$ over
$\{\mathcal{W}'\text{ holds}\}$, satisfy
\begin{itemize}
\item[(i)] $f_{\square}(e,X)\le1-\kappa$, and
\item[(ii)] $f(e,X)\le\theta+\nu_*$ for every non-seed window $f\notin\mathcal{W}'$, where
$\nu_*\ge0$ obeys
\begin{equation}\label{8.14:eq-admissible-centres-a}
(1-t_1)(\theta+\nu_*)+t_1\bar e<1 .
\end{equation}
\end{itemize}
Then the upper sandwich inequality holds with $X$ in place of $x_0$:
\begin{equation}\label{8.14:eq-admissible-centres-1}
\operatorname{Num}(e)\le\epsilon_{\kappa}\operatorname{Den}(\Xi e).
\end{equation}
\end{proposition}

For $\theta=1+4\bar q$, $t_1=16\bar q$ and $\bar e\le\tfrac18$ one has $t_1\bar e\le2\bar q$, hence
\begin{align*}
(1-t_1)(\theta+\nu_*)+t_1\bar e
&\le(1-16\bar q)(1+4\bar q+\nu_*)+2\bar q\\
&=1-10\bar q-64\bar q^{2}+(1-16\bar q)\nu_* .
\end{align*}
The choice $\nu_*:=10\bar q$ therefore satisfies \eqref{8.14:eq-admissible-centres-a}, the
right-hand side becoming $1-224\bar q^{2}<1$; so does every $\nu_*\ge0$ with
$\nu_*<(10\bar q+64\bar q^{2})/(1-16\bar q)$.

\begin{proof}
The proof of the upper sandwich inequality in the sandwich theorem consists of a numerator step,
bounding $\operatorname{Num}(e)$ from above, a denominator step, bounding
$\operatorname{Den}(\Xi e)$ from below, and a ratio step; we follow it and at each point say
where $X$ enters.

\emph{Numerator step.} Let $x\in\mathcal{E}\cap\operatorname{cl}S(e)$, put
$\xi:=\log(X^{-1}x)$ and $x_s:=X\exp(s\xi)$ for $0\le s\le1$. Let $f\in\mathcal{W}'$. Since $X$
lies in $\{\mathcal{W}'\text{ holds}\}$, $f(e,X)\le\theta_f$; since $x\in\operatorname{cl}S(e)$,
the definition of $S(e)$ gives $f(e,x)\le1$. Part~(i) of the near-convexity lemma therefore gives
\begin{equation}\label{8.14:eq-admissible-centres-2}
f(e,x_s)\le(1-s)\theta_f+s\cdot1+2s(1-s)\bar q\le\theta_f ,
\end{equation}
where the last inequality holds because the difference of its two sides is
\begin{equation*}
\theta_f-\bigl[(1-s)\theta_f+s+2s(1-s)\bar q\bigr]
=s\bigl(\theta_f-1-2(1-s)\bar q\bigr)\ge s\bigl(\theta_f-1-2\bar q\bigr)\ge0,
\end{equation*}
by $\theta_f\ge1+2\bar q$. Thus the whole segment $\{x_s:0\le s\le1\}$ lies in
$\{\mathcal{W}'\text{ holds}\}$. This is the assertion of the segment lemma with $\mathcal{S}$
replaced by $\{\mathcal{W}'\text{ holds}\}$: that assertion uses only the bound at the centre
for the windows imposed, the bound $f(e,x)\le1$ at the far end, part~(i) of the near-convexity
lemma, and the minimality of the centre over a set containing the segment, and each of these is
available here, the first because $X\in\{\mathcal{W}'\text{ holds}\}$, the last because $X$
minimises over that set.

Windows not in $\mathcal{W}'$ are no constraint of the family $\mathcal{W}'$. Along the segment
they read at most $\theta+\nu_*<2$: for a non-seed window $f\notin\mathcal{W}'$, part~(i) of the
near-convexity lemma with hypothesis~(ii) and $f(e,x)\le1$ gives
$f(e,x_s)\le(1-s)(\theta+\nu_*)+s+2s(1-s)\bar q$, and this is at most $\theta+\nu_*$ because
$\theta+\nu_*\ge\theta\ge1+2\bar q$, by the computation above with $\theta_f$ replaced by
$\theta+\nu_*$. Hence $x_s$ stays in the regular domain of the chart exactly as the segment lemma
asserts for its segment; we use here the regularity assertion of the segment lemma read for the
present segment, and nothing beyond it.

Put $\varphi(s):=\mathbf{A}(e,x_s)$. Since $X$ minimises $\mathbf{A}(e,\cdot)$ over
$\{\mathcal{W}'\text{ holds}\}$ and the segment lies in that set,
$\varphi(s)\ge\mathbf{A}(e,X)=\varphi(0)$ for $0\le s\le1$, and therefore $\varphi'(0)\ge0$.
These are the two properties of $\varphi$ on which the second-order step of the sandwich theorem
rests; that step, which applies the second-order lemma symmetrically with its parameter equal to
$\tfrac43\kappa$, together with the further second-order bound among the inputs of
the sandwich theorem, gives
\begin{equation*}
\mathbf{A}(e,x)\ge\mathbf{A}(e,X)+\Gamma^{\kappa}
\qquad\text{whenever}\qquad f_{\square}(e,x)\ge1\ge f_{\square}(e,X)+\kappa ,
\end{equation*}
the second of these two inequalities being hypothesis~(i). With this lower bound on the action in
place of the one at $x_0$, the numerator step of the sandwich theorem gives
\begin{equation}\label{8.14:eq-admissible-centres-3}
\operatorname{Num}(e)\le\operatorname{Haar}\bigl(S(e)\bigr)\,
\exp\Bigl(-\frac{\mathbf{A}(e,X)+\Gamma^{\kappa}}{g^{2}}+\sup\Psi\Bigr),
\end{equation}
with $\Psi$ as in the sandwich theorem.

\emph{Denominator step.} For $x\in S(e)$ put $\xi_x:=\log(X^{-1}x)$ and
$x_\sigma:=X\exp(\sigma\xi_x)$, where $\sigma$ is the contraction ratio of the denominator step
of the sandwich theorem, and let
$\mathcal{C}^{X}_{\sigma}(S(e))$ be the contraction of $S(e)$ about the centre $X$, whose points
are the $x_\sigma$. Put
\begin{equation*}
\mathfrak{J}:=\mathcal{M}\bigl(\mathcal{C}^{X}_{\sigma}(S(e))\bigr).
\end{equation*}
We verify the first claim of the localised volume proposition for the centre $X$, with the
centre's reading $\theta+\nu_*$ in place of $\theta$ for the windows outside $\mathcal{W}'$.

For a non-seed window $f\notin\mathcal{W}'$, part~(i) of the near-convexity lemma with
hypothesis~(ii) and $f(e,x)<1$ (definition of $S(e)$) gives
\begin{equation}\label{8.14:eq-admissible-centres-4}
f(e,x_\sigma)\le(1-\sigma)(\theta+\nu_*)+\sigma\cdot1+2\sigma\bar q\le\theta+\nu_* ,
\end{equation}
the last inequality being equivalent to $\sigma(\theta+\nu_*-1-2\bar q)\ge0$, which holds since
$\sigma\ge0$, $\nu_*\ge0$ and $\theta\ge1+2\bar q$. For $f\in\mathcal{W}'$ the window part of
the first claim applies verbatim and gives $f(e,x_\sigma)\le\theta_f\le\theta$. Thus every
non-seed window reads at most $\theta+\nu_*$ at $(e,x_\sigma)$, and the window bound for the
shifted chart gives, for every non-seed window $f$,
\begin{equation}\label{8.14:eq-admissible-centres-5}
f(\Xi e,\mathcal{M}x_\sigma)\le(1-t_1)(\theta+\nu_*)+t_1\bar e<1,
\end{equation}
the strict inequality being the choice \eqref{8.14:eq-admissible-centres-a} of $\nu_*$. Hence
every non-seed window holds at $(\Xi e,\mathcal{M}x_\sigma)$. The part of the first claim
concerning the seed window follows verbatim from hypothesis~(i) and $\kappa\ge\kappa_{\flat}$,
as in the sandwich theorem, and the remaining part of the first claim holds verbatim.

The second claim (the bound on the Haar measure of the contraction about an arbitrary centre,
from part~(ii) of the near-convexity lemma and the Haar bound for the shifted chart) and the
third claim,
\begin{equation}\label{8.14:eq-admissible-centres-6}
\mathbf{A}(\Xi e,x')\le\mathbf{A}(e,X)+\frac{\kappa^{2}\Gamma_{\square}}{4}
+2C'\varepsilon\mathfrak{d}g\qquad\text{for }x'\in\mathfrak{J},
\end{equation}
whose proof consists of derivative bounds along the $\sigma$-path issuing from the centre, the
path staying regular by the window part of the first claim, together with the smallness
condition on the action bound along that path, read the centre only through its window readings.
They use neither the minimality of the centre nor its membership in $\mathcal{S}$: this is the
observation, made with the proposition giving the denominator step of the sandwich theorem, that
this step never uses the minimality of the centre. The window readings of $X$ are those
established in \eqref{8.14:eq-admissible-centres-4}, \eqref{8.14:eq-admissible-centres-5} and
in the seed part above, so the second and third claims hold with $X$ as centre, and the
denominator step of the sandwich theorem bounds $\operatorname{Den}(\Xi e)$ from below through
$\mathfrak{J}$ exactly as for $x_0$.

\emph{Ratio.} The ratio step of the sandwich theorem now applies verbatim to
\eqref{8.14:eq-admissible-centres-3} and to the lower bound for $\operatorname{Den}(\Xi e)$
obtained from the second claim and \eqref{8.14:eq-admissible-centres-6}; the terms
$\mathbf{A}(e,X)$ cancel, and, using $\varepsilon=g\,p_0(g)$ with $p_0(\cdot)$ the degree
function, the exponent of the ratio is
\begin{equation*}
-\frac{\Gamma^{\kappa}-\kappa^{2}\Gamma_{\square}/4-2C'\varepsilon\mathfrak{d}g}{g^{2}}
=-\frac{15}{4}\,\frac{\kappa^{2}\Gamma_{\square}}{g^{2}}+2C'\mathfrak{d}\,p_0(g) ,
\end{equation*}
which is the exponent of the sandwich theorem. This gives \eqref{8.14:eq-admissible-centres-1}.
\end{proof}

\begin{corollary}[The fibrewise sandwich under negative room]\label{8.14:cor-negative-room}
Let $\hat{x}$ be the canonical comparison point and assume
$\hat{x}\in\mathsf{M}(\widehat{\mathcal{W}})$, that is, $\hat{x}$ is a minimiser under the window
family $\widehat{\mathcal{W}}$. Assume that the upper window threshold $\theta$
satisfies $\theta\le 2$. Here $e$ is the base datum of the fibre, and $e\oplus x$ is the
configuration of the chart at the point $x$ over $e$. The value
$f(e,\hat{x})$ of a window $f$ at $\hat{x}$ is called its \emph{reading} there. Further, $y_+$ is the
exit row, $\delta(p)$ is the row coordinate of a fine plaquette $p$, and $r_1$ is the collar
width. Finally, $\theta_y$ is the exit-row threshold and $\beta_0$ is the per-level growth
constant. Let $\nu_*\ge 0$ obey the tolerance inequality \eqref{8.14:eq-admissible-centres-a}.
Let $\nu'_*$ be the largest $\nu'$ with
\begin{equation}\label{8.14:eq-negative-room-1}
2\nu'(1+\vartheta_1)+3\vartheta_1\le\nu_*,
\end{equation}
where $\vartheta_1\ge 0$ is the exponentially small tail constant,
and let $0\le\nu'\le\nu'_*$. Let $j$ be the level fixed in the setting of part~(iv) of the
proposition in hypothesis~(b1) below. Put $k_o=j+c_{\mathrm{obs}}+1$, where the observation depth
$c_{\mathrm{obs}}$ is at least the threshold depth $c_\sharp$. The seed window is $f_{\square}$;
every other window is called \emph{non-seed}. The sites at which a window is read are of three
kinds: the outer-ring sites beyond the exit row $y_+$, the second-row sites, and the sites at levels
$i\le k_o-2$. Assume the following.
\begin{itemize}
\item[(a)] Proposition~\ref{8.14:prop-admissible-centres} (truncated minimisers as admissible
sandwich centres).
\item[(b)] The following statements, on which hypothesis~(c) rests:
\begin{itemize}
\item[(b1)] the proposition that bounds the window readings at $\hat{x}$ through the collar
bound, in its parts (i), (ii) and (iv);
\item[(b2)] the statements from which parts (i) and (ii) of the proposition in (b1) are proved:
part~(ii) of a lemma that this proof applies twice, a maximum principle and an estimate that
this proof applies together with it, and the two further statements of the list of inputs of
that proposition.
\end{itemize}
\item[(c)] The proof of parts (i) and (ii) of the proposition in (b1), re-run with the
collar-room parameter set to $-2\nu'\le 0$ in place of a non-negative value (each step of that
proof enters the collar bound affinely, so the proof keeps its form when the room is negative),
yields the following
from the exit bound \eqref{8.14:eq-negative-room-a} below. At $\hat{x}$, every non-seed window
outside $\widehat{\mathcal{W}}$ reads at most $(1+\nu')\theta$ at outer-ring sites beyond $y_+$, and
at most $(1+\nu')\theta(1+\vartheta_1)+\vartheta_1$ at second-row sites. At levels $i\le k_o-2$ it
reads at most
\begin{equation}\label{8.14:eq-negative-room-2}
(1+\nu')\theta(1+\vartheta_1)^2(1+\beta_0)^{k_o-1-i}L^{-2(k_o-1-i)}+\vartheta_1\le\theta .
\end{equation}
The last inequality is part of what the re-run yields. Part~(iv) of the same proposition, applied
with the exit bound
\eqref{8.14:eq-negative-room-a} in place of the exit bound of that proposition, yields
$f_{\square}(\hat{x})\le\tfrac12$ (with $k_o$ and $c_{\mathrm{obs}}\ge c_\sharp$ as above);
in this application the factor $1+\nu'$ of \eqref{8.14:eq-negative-room-a} is at most $2$ and lies
inside the margin of the skin bound used in the proof of part~(iv).
\item[(d)] The remaining hypotheses of Proposition~\ref{8.14:prop-admissible-centres} hold for the
window family $\widehat{\mathcal{W}}$ and the sandwich parameter $\tfrac12$. That is,
$\widehat{\mathcal{W}}$ is a family of non-seed windows imposed at thresholds in
$[1+2\bar q,\theta]$, and the set on which it holds contains $\mathcal{S}$. Moreover $\tfrac12$ is
at least the floor $\kappa_\flat$ of that proposition, and the ceilings under which that
proposition is stated hold at this parameter.
\end{itemize}
Suppose that the exit bound holds with negative room $-\nu'$. That is, every fine plaquette $p$
with $y_+<\delta(p)\le y_++r_1+1$ obeys
\begin{equation}\label{8.14:eq-negative-room-a}
\bigl\|U(e\oplus\hat{x})(\partial p)-1\bigr\|\le(1+\nu')\,\theta_y .
\end{equation}
Then the following hold.
\begin{itemize}
\item[(1)] Every non-seed window outside $\widehat{\mathcal{W}}$ reads at most $\theta+\nu_*$
at $\hat{x}$.
\item[(2)] The seed window satisfies $f_{\square}(\hat{x})\le\tfrac12$.
\item[(3)] The fibrewise sandwich of Proposition~\ref{8.14:prop-admissible-centres} holds with
centre $\hat{x}$ and sandwich parameter $\tfrac12$. In that proposition's notation
$\operatorname{Num}$, $\operatorname{Den}$, $\Xi$, $\epsilon_{1/2}$ it reads
\begin{equation}\label{8.14:eq-negative-room-3}
\operatorname{Num}(e)\le\epsilon_{1/2}\,\operatorname{Den}(\Xi e).
\end{equation}
This holds although $\hat{x}$ need not lie in $\mathcal{S}$.
\end{itemize}
\end{corollary}

\begin{proof}
\emph{Clause (1).} Let $f$ be a non-seed window of the fibre over $e$ that is not in
$\widehat{\mathcal{W}}$. Hypothesis~(c) is applicable, since its input is the exit bound
\eqref{8.14:eq-negative-room-a}, which we have assumed. It bounds the reading of $f$ at $\hat{x}$
by one of three numbers, according to the kind of site. We compare each with $\theta+\nu_*$.
Two facts are used throughout. Since $\theta\le 2$, we have $\nu'\theta\le 2\nu'$. Since
$\nu'\le\nu'_*$ and $\nu'_*$ obeys \eqref{8.14:eq-negative-room-1}, so does $\nu'$.

At outer-ring sites beyond $y_+$ the bound is
\begin{equation*}
(1+\nu')\theta=\theta+\nu'\theta\le\theta+2\nu'\le\theta+\nu_* .
\end{equation*}
Here the last step uses $2\nu'\le 2\nu'(1+\vartheta_1)+3\vartheta_1\le\nu_*$. This holds because
$\vartheta_1\ge 0$ and $\nu'$ obeys \eqref{8.14:eq-negative-room-1}.

At second-row sites we expand the product and use $\theta\le 2$ twice:
\begin{align*}
(1+\nu')\theta(1+\vartheta_1)+\vartheta_1
&=\theta+\nu'\theta(1+\vartheta_1)+\theta\vartheta_1+\vartheta_1\\
&\le\theta+2\nu'(1+\vartheta_1)+3\vartheta_1\le\theta+\nu_* .
\end{align*}
The last inequality is \eqref{8.14:eq-negative-room-1} for $\nu'$.

At levels $i\le k_o-2$, hypothesis~(c) gives the bound \eqref{8.14:eq-negative-room-2}. This is
at most $\theta$, and hence at most $\theta+\nu_*$, because $\nu_*\ge 0$. This proves clause~(1).

\emph{Clause (2).} By the last statement of hypothesis~(c), whose only input is the exit bound
\eqref{8.14:eq-negative-room-a}, which we have assumed, the seed window satisfies
$f_{\square}(\hat{x})\le\tfrac12$.

\emph{Clause (3).} We apply Proposition~\ref{8.14:prop-admissible-centres} with
$\mathcal{W}':=\widehat{\mathcal{W}}$, $X:=\hat{x}$ and $\kappa:=\tfrac12$. Here $X$ and $\kappa$
are that proposition's centre variable and sandwich parameter, in its notation. Its requirements
on the centre are three; its remaining hypotheses are hypothesis~(d).
\begin{itemize}
\item[--] $X\in\mathsf{M}(\mathcal{W}')$. This is our assumption
$\hat{x}\in\mathsf{M}(\widehat{\mathcal{W}})$.
\item[--] Its condition~(i), $f_{\square}(X)\le 1-\kappa$. For $\kappa=\tfrac12$ this reads
$f_{\square}(\hat{x})\le\tfrac12$, which is clause~(2).
\item[--] Its condition~(ii). Every non-seed window $f\notin\widehat{\mathcal{W}}$ must satisfy
$f(e,\hat{x})\le\theta+\nu_*$, with $\nu_*\ge 0$ obeying \eqref{8.14:eq-admissible-centres-a}.
The bound is clause~(1). The condition on $\nu_*$ is one of our assumptions.
\end{itemize}
Membership of the centre in $\mathcal{S}$ is not among the hypotheses of
Proposition~\ref{8.14:prop-admissible-centres}. That proposition therefore applies to $\hat{x}$,
and its conclusion gives \eqref{8.14:eq-negative-room-3}, the fibrewise sandwich. This holds
whether or not $\hat{x}\in\mathcal{S}$.
\end{proof}

\begin{remark}[Model: exponentially small defects within the tolerance]\label{8.14:rem-defect-tolerance}
The model of this remark is that of the model statement Remark~\ref{8.14:rem-outside-collar}. It
consists of the truncated problem $P_a$, truncated at height $a$ into the band, its minimiser $X_a$,
and the comparison of the exit level of $X_a$ with the young caps $\mathfrak{B}$. The bounds of that
model are compared here with the tolerance of Proposition~\ref{8.14:prop-admissible-centres} and
Corollary~\ref{8.14:cor-negative-room}, both of which are proved as implications from the statements
listed there.

\emph{(a) What the tolerance changes.} At the canonical comparison point $\hat{x}$, the exit bound
with positive room $\rho$ requires $\rho\ge 8\vartheta_1+8\vartheta_t$. The exit bound at the canonical
comparison point with room $\rho=0$ rests on exact membership of its comparison point in the segment
set $\mathcal{S}$, which the exact-membership lemma supplies.
Proposition~\ref{8.14:prop-admissible-centres} and Corollary~\ref{8.14:cor-negative-room} show that
the fibrewise sandwich is obtained already under room $-\nu'_*$ at $\hat{x}$ itself, which need not
lie in $\mathcal{S}$, where $\nu'_*\asymp\bar q$. This is the hypothesis
\eqref{8.14:eq-negative-room-a}: the exit collar may exceed the young cap by the relative amount
$\nu'_*$.

More generally, it suffices to have any statement of the following form: some truncated minimiser
violates the remaining non-seed windows by at most $\nu_*$, in the sense of hypothesis (ii) of
Proposition~\ref{8.14:prop-admissible-centres}, with $\nu_*$ subject to
\eqref{8.14:eq-admissible-centres-a}. The exactness of the exact-membership lemma is therefore
stricter than the half-skin bound at the canonical comparison point needs.

This is consistent with the emptiness statement for violating centres. Let $K_w$, $C_{\delta_0}$ and
$\varepsilon_{\min}$ be the constants of that statement. A centre that violates a window by
$\nu\ge\nu_\flat$, where
\begin{equation*}
  \nu_\flat=48\bar q+\frac{2K_wC_{\delta_0}\,g}{\varepsilon_{\min}},
\end{equation*}
leaves empty the set over which the lower bound of the half-skin bound is taken, so that this lower
bound is void. For violations between $\nu_*$ and $\nu_\flat$, the one-sided bounds of clause (e) of
the comparison statement listed among the antecedents of
Proposition~\ref{8.14:prop-admissible-centres} decide neither way.

\emph{(b) Where it helps.} Some bounds on the collar have natural defects of a size $\bar\vartheta$
exponentially small in $R$ or in $r_1$; here $R$ is the radius that enters
\eqref{8.14:eq-cutoff-in-height-c}, and $r_1$ is the width of the collar in
\eqref{8.14:eq-negative-room-a}. Three sources of such defects occur:
\begin{itemize}
\item the exponential tail terms in the list of antecedents of
Proposition~\ref{8.14:prop-admissible-centres};
\item the tail constant $\vartheta_A\le\vartheta_1$ of the maximum principle used in the proof of
Corollary~\ref{8.14:cor-negative-room};
\item the constant $\vartheta_t$ of the lemma applied twice in that proof.
\end{itemize}
Since $\bar\vartheta\ll\bar q$, a bound of the collar by $(1+C\bar\vartheta)\theta_y$, for a constant
$C$, suffices where a bound by $\theta_y$ exactly was required.

\emph{(c) Where it does not help.} It does not help bounds obtained from the energy estimate.

The smallness free of the cutoff $n$ of the model statement Remark~\ref{8.14:rem-outside-collar}, in
its cutoff-free overshoot bound beyond the collar, is free of $g$. It is $e^{-c\mathsf{T}/2}$ times a
polynomial in $P_{\mathrm{bl}}$, $s_e/\mathsf{T}$ and the further constants entering it.

The tolerance $\nu_*\tau_y$, by contrast, is small with $g$, because $\bar q\to0$ as $g\to0$. Two
facts show this. First, the floor $\kappa_\flat'$ of the sandwich theorem for minimisers in the
segment set is $O(\bar q+1/p_0)$, and the condition on it closes as $g\to0$. Second, the bracket in one
of the ceilings under which Proposition~\ref{8.14:prop-admissible-centres} holds reads, with
$\sigma_P$ and $\mathfrak{d}$ the constants of that ceiling,
\begin{equation*}
  \frac{16\,\bar q\,\sigma_P\,p_0\,\mathfrak{d}}{g}
  =O\bigl(g\,p_0^3\cdot\mathrm{polylog}(g^{-2})\cdot\mathfrak{d}\bigr).
\end{equation*}

By the model statement Remark~\ref{8.14:rem-outside-collar}, in whose notation the following display
is written, the overshoot beyond the collar is at most $\tau_y$ times the left side of
\begin{equation}\label{8.14:eq-defect-tolerance-1}
  \Bigl[\frac{96\,C_{\mathrm{hm}}P_{\mathrm{bl}}^3\,\beta_{\mathrm{hot}}
  \,(s_e\bar\lambda/\mathsf{T})^2}{\varpi'^3\varpi''}\Bigr]^{1/2}\le\nu_* .
\end{equation}
The energy bound would fall within the tolerance under the floor \eqref{8.14:eq-defect-tolerance-1}.
That floor fails as $g\to0$ at fixed $L$. Adopted as a condition, it would bound $g$ from below, so
it is not adopted. Inequality \eqref{8.14:eq-defect-tolerance-1} is the precise point at which the
energy bound fails to meet the tolerance.

Moreover, inside the collar the tolerance $\nu_*\tau_y$ is still a threshold at a single site. The
pointwise criterion for stopping the front is used together with \eqref{8.14:eq-stopping-height-b}
in the model statement Remark~\ref{8.14:rem-cutoff-in-height}. Under the tolerance, the count of
levels in that criterion becomes
\begin{equation*}
  \log\bigl(2u_{a_1}/(\nu_*\tau)^2\bigr),
\end{equation*}
and this count contains the terms
\begin{equation*}
  3k_o\log L+2\log(1/\bar q).
\end{equation*}
This is the form of \eqref{8.14:eq-cutoff-in-height-c} with $\bar q$ in place of $1/R$.

Proposition~\ref{8.14:prop-admissible-centres} therefore lowers what has to be proved. An exact
free-boundary identity is replaced by a quantitative bound in supremum: on one shell, the overshoot
over the young cap is at most $\nu_*$ times the cap. This bound is of the same kind as the exit bound,
with negative room. The proposition does not change the fact that the bound is in supremum, and it is
this supremum nature that the cutoff-free overshoot bound beyond the collar of the model statement
Remark~\ref{8.14:rem-outside-collar} isolates.
\end{remark}

\begin{remark}[Model: mean square against supremum on the collar]\label{8.14:rem-collar-supremum}
This remark concerns the model of the present section. Its ingredients are the truncated
problems $P_a$ for fields on the levels above the lateral torus $\mathbb{T}=(\mathbb{Z}/N)^3$, with
values $m(a)$, minimisers $X_a$, increments $\Delta m(a)$, top-level defects $u_a$, exit pulls
$e_a$ and exit over-pulls $S_a=e_a-\pi_{\mathfrak{B}}(e_a)$ against the young caps $\mathfrak{B}$,
as in Corollary~\ref{8.14:cor-stopping-height} and Corollary~\ref{8.14:cor-over-pull-small}
(both model). A level carries $3N^3$ fine sites. We write $n$ for the cutoff
parameter of the model; $N$ is proportional to $n$. A quantity is called \emph{cutoff-free}
if it is bounded independently of $n$.

\emph{The statement not proved.} Consider the stopping property, namely that one closed
shell of windows is exactly inactive at the canonical comparison point $\hat{x}$, in the
following cutoff-uniform form: the front of Remark~\ref{8.14:rem-cutoff-in-height} (model) stops
within a number of levels whose ratio to the width $W_y$, in levels, of the band of young
levels is bounded independently of $n$.
This form is not proved by the route of this section.

\emph{Where the cutoff enters.} The route controls the exit over-pull of a truncated
minimiser independently of $n$ in two senses:
\begin{itemize}
\item in root mean square, by Corollary~\ref{8.14:cor-over-pull-small}\,(i) (model);
\item in supremum beyond a collar of a number of levels proportional to $n$, by
Corollary~\ref{8.14:cor-over-pull-small}\,(ii) and by Theorem~\ref{8.14:thm-collar-exactness}
(both model).
\end{itemize}
Proposition~\ref{8.14:prop-admissible-centres} admits truncated minimisers as sandwich
centres. Even at its tolerance \eqref{8.14:eq-admissible-centres-a} (overshoot at most
$\nu_*$ times the cap, with $\nu_*\asymp 10\bar q$), it requires a supremum bound on the
over-pull on the collar itself. The passage from root mean square to supremum on a level of
$3N^3\propto n^3$ fine sites is where the cutoff enters.

The only supremum device of the route is the pointwise stopping criterion driven by
\eqref{8.14:eq-stopping-height-b}. It pays the cutoff as the product
$\mathfrak{h}\,\mathfrak{n}$ in the stopping height \eqref{8.14:eq-cutoff-in-height-a},
through the two factors analysed in Remark~\ref{8.14:rem-cutoff-in-height} (model):
\begin{itemize}
\item the e-folding height $\mathfrak{h}$, bounded through \eqref{8.14:eq-cutoff-in-height-b};
\item the e-fold count $\mathfrak{n}$, estimated in \eqref{8.14:eq-cutoff-in-height-c}.
\end{itemize}
For the first factor, the e-folding height $\mathfrak{h}$ is proportional to $s_\infty/\tau$. Here
$\tau$ denotes the common order of magnitude of the cap levels $\tau_k$ of the box
$\mathfrak{B}$, and $s_\infty$ is
the supremum bound on the top multiplier. Under a
tangential bound $\|B_a\|_\infty\le c_B\tau$ on the exit row one has
\begin{equation*}
\frac{s_\infty}{\tau}\lesssim 1+c_B\bigl(C+\mathfrak{k}_N\bigr),
\end{equation*}
where $C$ is a constant independent of $N$ and $\mathfrak{k}_N$ is the $\ell^1(\mathbb{T})$-norm
of the convolution kernel of $(\lambda+\mathsf{v})^{-1}\operatorname{curl}^*$; by a lattice Fourier
estimate for convolution kernels whose symbol is homogeneous of degree zero, not reproduced
here, $\mathfrak{k}_N$ is of order at most $1+\log N$. Lemma~\ref{8.14:lem-exit-pull} (model)
identifies $e_a$ as the
image $-(\lambda+\mathsf{v})^{-1}\operatorname{curl}^*B_a$ of the tangential curvature $B_a$ of
the exit row under this operator, whose symbol is homogeneous of degree zero at small lateral
momenta; the exit row lies on a free boundary, where no smoothing scale is available. The only
a priori bounds on $s_\infty$ available to the route are this one, of logarithmic size in units
of $\tau$, and the extensive bound $\|\mu\|_\infty\le(2C_Nu_a)^{1/2}$. The available bound has
this logarithmic form at the front even in the scalar model of the model theorem on inactivity
at full band width.

For the second factor, the e-fold count is
\begin{equation*}
\mathfrak{n}=5k_o\log L+O(1),
\end{equation*}
where $k_o$ is the observation level and the $O(1)$ term is independent of $n$; here one uses
that, under the identification of the model's parameters with those of the lattice problem used
in Remark~\ref{8.14:rem-outside-collar} (model), $\log n$ equals $k_o\log L$ up to a constant
independent of $n$. This count compares an extensive quantity with a one-site threshold.

\emph{Why the route does not remove these factors.} Consider the class of arguments in which
the monotone quantity is a value increment or an energy ($u_a$, $\Delta m(a)$, or
$m(a+j)-m(a)$), and the terminal criterion is pointwise feasibility of one shell. Within
this class the two factors are not removed, for three reasons.
\begin{itemize}
\item[(1)] The one-row ceiling, Lemma~\ref{8.14:lem-one-row-ceiling} (model), pins the gain
per row to the one-step Poisson deficit $\kappa_1$.
\item[(2)] The second factor is intrinsic to the comparison of an extensive quantity with a
pointwise one, whatever the rate. The bookkeeping of that comparison without $s_\infty$
trades the logarithm for a factor $n^{5/2}$.
\item[(3)] By Remark~\ref{8.14:rem-ring-freedom} (model), no fixed number of rings absorbs
either logarithm.
\end{itemize}
The obstruction is not a property of the model: see Remark~\ref{8.14:rem-double-logarithm}
(model).

Five natural questions about the route are answered as follows.
\begin{itemize}
\item[(a)] A Poisson-weighted deficit summed over heights does not remove the cutoff.
Remark~\ref{8.14:rem-cutoff-in-height} (model) shows that the dyadic scale of the deficit is
not where the cutoff enters. As a consequence of Lemma~\ref{8.14:lem-one-row-ceiling} (model),
the linear gain at a site of the support of $\mu$ is at most $|\mu_k(x)|\kappa_1(x,k)$ for
every admissible
competitor, so a deficit at scale $j>1$ is invisible to the one-row comparison. Even if it were
visible, it could remove at best the factor $s_\infty/\tau$ from $\mathfrak{h}$, never the
count $\mathfrak{n}$.
\item[(b)] A better lower bound for the absorption term enters only through $s_\infty$, that
is, through the first factor. That factor is removable on the charged class (multiplier charge
$q=\nabla^*\mu\neq0$) by
Corollary~\ref{8.14:cor-charged-decay} (model) under hypothesis
\eqref{8.14:eq-charged-decay-a}. It is genuine on the charge-free class, where a comparison
(order) structure is available instead, by the transfer of the lift of the model theorem on
inactivity at full band width to the charge-free class.
\item[(c)] Confining the front to a fixed number of rings does not help. By
Remark~\ref{8.14:rem-ring-freedom} (model), the required number of rings would grow like
$\log_L k_o$.
\item[(d)] The cutoff-uniform form is proved for everything outside the collar, independently
of $n$, by Theorem~\ref{8.14:thm-collar-exactness} and
Corollary~\ref{8.14:cor-over-pull-small} (both model). The size of the resulting bounds is
given in Remark~\ref{8.14:rem-outside-collar} (model).
\item[(e)] The product $\mathfrak{h}\,\mathfrak{n}$ of the two logarithmic factors (the double
logarithm of Remark~\ref{8.14:rem-double-logarithm}, model) belongs to the method: a value
functional read front by
front, together with a pointwise exit criterion. It sits at the passage from mean square to
supremum on the collar. It is not a property of the scalar member of the model, in which the
front stops within a fraction of the band width $W_y$ bounded independently of $n$, by the
model theorem on inactivity at full band width.
\end{itemize}

\emph{What a cutoff-free proof of the stopping property that uses no comparison (order)
argument must contain.} It must contain a
supremum statement on one shell, obtained from a quantity integrated over heights. The
prototype is part (a) of the model theorem on inactivity at full band width. In the form it
takes on the charge-free class, that statement bounds the climb of the axial potential over $j$
consecutive levels by $\ell_V[\Lambda_N(j+b_y)+\Lambda_N(b_y)]$, where $\Lambda_N$ is the lift
function of the scalar model, $b_y$ is the level offset and $\ell_V$ the field scale of that
statement. Since $\Lambda_N(t)\le t\lambda_N(t)$ with $\lambda_N(t)=C_K^0+6+6\log^{+}(N/t)$,
the logarithm attached to a climb of $j$ levels is $\log(N/j)$. The required
statement must be of this kind but valid for charged configurations.

By the partition of the problem into the charge-free class and its complement, described after
Corollary~\ref{8.14:cor-charged-decay} (model), such a statement is the column-wise sign
(order) descent for the axial potential,
with the error term $\partial_k\sum_{y'<y}q(\cdot,y')$ paid, where $y,y'$ run over heights,
$\partial_k$ is the lateral difference in the direction $k$, and $q=\nabla^*\mu$ is the
multiplier charge.
The only available statement that pays for charge is
Lemma~\ref{8.14:lem-charge-energy} (model). It pays
$\tfrac12\langle q,(1+\mathsf{C})(\lambda+\mathsf{v})^{-1}q\rangle$ per level out of the
value budget $M_h:=m(a_0+h)-m(a_0)$ of $h$ levels above the base level $a_0$ of the
truncations of Corollary~\ref{8.14:cor-stopping-height} (model), that is, in energy and not in
supremum. The missing link
is the statement that, for the top levels of the truncated problems, the supremum
\begin{equation*}
\sup_x\Bigl|\partial_k\Delta_{\parallel}^{-1}\sum_{y'<y}q(\cdot,y')\Bigr|(x)
\end{equation*}
is controlled by the energy
\begin{equation*}
\sum_a\tfrac12\bigl\langle q^{(a)},(1+\mathsf{C})(\lambda+\mathsf{v})^{-1}q^{(a)}\bigr\rangle
\le M_h .
\end{equation*}
This is an energy-to-supremum statement for the Riesz potential of the charge integrated over
heights. By the mechanism behind the second factor, it cannot hold without further structure:
a single charged fine needle (a laterally bounded active set of one orientation, at fine scale)
is invisible in $M_h$.

Two further structures are available and not used on this route.
\begin{itemize}
\item[($\alpha$)] The charge of a minimiser is not arbitrary: the proposition that activity is
never self-sustained, together with its companion lemma for the vector problem, ties the sign
pattern of $q$ to the field.
\item[($\beta$)] Under block caps the windows read block averages of the field rather than its
values at fine sites. A fine needle of amplitude
below a quantity of order $L^2$ times the cap is then invisible to the windows as well, so
the needles that matter are coherent on blocks, of lateral size at least $L$ fine units.
This does not make the count independent of $n$, since the block volume $L^4$ is much
smaller than $n^3$. It does align what the energy sees with what the windows see up to a
factor $(n/L)^3$ instead of $n^3$, and it leaves the nature of the estimate unchanged.
\end{itemize}
No device is supplied here. The cutoff-free target is a supremum bound on one shell, and along
the route of this section it is exactly the passage from mean square to supremum there.
\end{remark}

\begin{lemma}[Model: the optimality system inside the capped band]\label{8.15:lem-band-optimality}
We work in the model of the model theorem on the exit value: its unknowns are
square-summable sequences $(E_y)_{y\ge0}$ of lateral $1$-forms $E_y\in\mathfrak{H}$ on a finite
torus. The transversal potentials $\beta_y\in\mathfrak{H}_T$ obey $\beta_{y+1}=\beta_y+PE_y$, where
$P$ and $P^{\perp}$ are the transversal and longitudinal projections. We write $\nabla^{*}$ for the
lattice divergence of the model; then $\mathfrak{H}_T=\operatorname{ran}P\subset\ker\nabla^{*}$.
The energy is
\[
  J(E)=\sum_{y\ge0}\tfrac12\|E_y\|^2+\sum_{y\ge1}\tfrac12\langle\beta_y,\lambda\beta_y\rangle ,
\]
with $\lambda\ge0$ the magnetic operator of the model acting on $\mathfrak{H}_T$, so that
$\|\operatorname{curl}\alpha\|^2=\langle P\alpha,\lambda P\alpha\rangle$ for lateral $1$-forms
$\alpha$, where $\operatorname{curl}$ is the lattice curl of the model. The model fixes a level
$a_0$, the bottom of the capped band; the lower constraints of the model read only the rows
$\le a_0$, that is, the fields $E_y$ with $y<a_0$. For $a_0\le a\le\infty$ the
problem $P_a$ requires in addition $E_y\in\mathfrak{B}$, the feasible box of a capped level, for
$a_0\le y<a$; thus $P_\infty$ caps all levels $y\ge a_0$. In the model $\mathfrak{B}$ is the box
of cap $\tau>0$; it is closed and convex and contains a neighbourhood of $0$ in the norm of
$\mathfrak{H}$. For $E\in\mathfrak{B}$ let
$\mathcal{N}_{\mathfrak{B}}(E)$ be the set of all $\xi'\in\mathfrak{H}$ such that
$\langle\xi',E'-E\rangle\le0$ for every $E'\in\mathfrak{B}$; this is the normal cone of
$\mathfrak{B}$ at $E$.

Fix $a$ with $a_0\le a\le\infty$ and let $(E_y,\beta_y)_{y\ge0}$ be the minimiser of $P_a$, which
is unique by the well-posedness of the model. Put
\[
  \mathcal{D}_y:=-\sum_{y'>y}\lambda\beta_{y'}\quad(y\ge a_0),\qquad \mu_y:=\mathcal{D}_y-E_y .
\]
Then the following hold.
\begin{enumerate}
\item[(i)] $\mathcal{D}_y\in\mathfrak{H}_T$.
\item[(ii)] $\mu_y\in\mathcal{N}_{\mathfrak{B}}(E_y)$ for $a_0\le y<a$, and $\mu_y=0$ for $y\ge a$.
\item[(iii)] The fields satisfy
\begin{equation}\label{8.15:eq-band-optimality-a}
  \beta_{y+1}-\beta_y=PE_y,\qquad
  \mathcal{D}_y-\mathcal{D}_{y-1}=\lambda\beta_y\quad(y>a_0),
\end{equation}
and hence, for $y>a_0$,
\begin{equation}\label{8.15:eq-band-optimality-1}
  \mathcal{D}_{y+1}-2\mathcal{D}_y+\mathcal{D}_{y-1}
  =\lambda PE_y=\operatorname{curl}^{*}\operatorname{curl}E_y .
\end{equation}
\item[(iv)] Let $\operatorname{clip}_\tau$ be the nearest-point projection onto the box
$\mathfrak{B}$ of cap $\tau$, and put
$\operatorname{shrink}_\tau:=\mathrm{id}-\operatorname{clip}_\tau$. Then on capped levels
$E_y=\operatorname{clip}_\tau(\mathcal{D}_y)$ and $\mu_y=\operatorname{shrink}_\tau(\mathcal{D}_y)$.
\item[(v)] $P^{\perp}E_y=-P^{\perp}\mu_y$; thus the longitudinal field exists only through the
multipliers. The charge at level $y$ is
\[
  q_y:=\nabla^{*}\mu_y=-\nabla^{*}E_y ,
\]
and this $q_y$ is the charge that enters the exit-charge statement of this chapter.
\end{enumerate}
\end{lemma}

\begin{proof}
The minimiser has finite energy, so $\sum_y\|E_y\|^2\le2J(E)<\infty$ and $\|E_y\|\to0$ as
$y\to\infty$. We use repeatedly that $\lambda\beta_{y'}\in\mathfrak{H}_T$, because
$\beta_{y'}\in\mathfrak{H}_T$ and $\lambda$ acts on $\mathfrak{H}_T$; hence
$P\lambda\beta_{y'}=\lambda\beta_{y'}$, and since $P$ is an orthogonal projection,
$\langle Ph,\lambda\beta_{y'}\rangle=\langle h,\lambda\beta_{y'}\rangle$ for $h\in\mathfrak{H}$.
The recursion of the transversal potentials gives
\[
  \beta_{y'}=\beta_0+\sum_{y''<y'}PE_{y''} ,
\]
with $\beta_0$ fixed by the data of the model.

We call a level $y'\ge a_0$ \emph{interior} if $E_{y'}$ is an interior point of its constraint
set, that is, if there is $\delta>0$ such that $E_{y'}+\eta$ satisfies the constraint of level
$y'$ for every $\eta\in\mathfrak{H}$ with $\|\eta\|<\delta$. Every free level $y'\ge a$ is
interior. All sufficiently large levels are interior. If $a<\infty$, this holds because
they are free. If $a=\infty$, we use that $\mathfrak{B}$ contains a ball
$\{\eta:\|\eta\|<2\delta\}$ about $0$ for some $\delta>0$, as recalled above. Choose $\bar y$ with
$\|E_{y'}\|<\delta$ for $y'\ge\bar y$; for such $y'$ and $\|\eta\|<\delta$ we have
$\|E_{y'}+\eta\|<2\delta$, so $E_{y'}+\eta\in\mathfrak{B}$.

We now vary the minimiser along directions of finite energy. Let $y\ge a_0$, let $a'>y$ be an
interior level, with its $\delta$, and let $h\in\mathfrak{H}$ be such that $E_y+h$ satisfies the
constraint of level $y$; for $y\ge a$ every $h$ is allowed. For $\epsilon\in\mathbb{R}$ let
$E^\epsilon$ coincide with $E$ except at two levels, where
\[
  E^\epsilon_y:=E_y+\epsilon h,\qquad E^\epsilon_{a'}:=E_{a'}-\epsilon Ph .
\]
Since $P^2=P$, the formula for $\beta_{y'}$ gives $\beta^\epsilon_{y'}=\beta_{y'}+\epsilon Ph$ for
$y<y'\le a'$ and $\beta^\epsilon_{y'}=\beta_{y'}$ for all other $y'$. Let $0<\epsilon<1$ and
$\epsilon\|h\|<\delta$.
Then $E^\epsilon_y=(1-\epsilon)E_y+\epsilon(E_y+h)$ satisfies the constraint of level $y$, by
convexity of $\mathfrak{B}$ when $y<a$. Further $E^\epsilon_{a'}$ satisfies the constraint of
level $a'$, because
$\|\epsilon Ph\|<\delta$. The lower constraints read only $E_{y''}$ with $y''<a_0\le y$ and are
unchanged, and no other level's cap is touched. So $E^\epsilon$ is feasible. Only the terms of $J$
at the levels $y$, $a'$ and $y<y'\le a'$ change; since $\lambda$ is self-adjoint,
\begin{align*}
  J(E^\epsilon)-J(E)
  &=\epsilon\Bigl[\langle h,E_y\rangle-\langle Ph,E_{a'}\rangle
    +\sum_{y<y'\le a'}\langle Ph,\lambda\beta_{y'}\rangle\Bigr]\\
  &\quad+\tfrac{\epsilon^2}{2}\bigl[\|h\|^2+\|Ph\|^2+(a'-y)\langle Ph,\lambda Ph\rangle\bigr] .
\end{align*}
Minimality gives $J(E^\epsilon)\ge J(E)$. We divide by $\epsilon>0$, let $\epsilon\to0$, and use
$\langle Ph,\lambda\beta_{y'}\rangle=\langle h,\lambda\beta_{y'}\rangle$ and
$\langle Ph,E_{a'}\rangle=\langle h,PE_{a'}\rangle$. This yields
\begin{equation}\label{8.15:eq-band-optimality-2}
  \Bigl\langle h,\,E_y+\sum_{y<y'\le a'}\lambda\beta_{y'}-PE_{a'}\Bigr\rangle\ge0 .
\end{equation}
If $-h$ is allowed as well, equality holds in \eqref{8.15:eq-band-optimality-2}.

The series defining $\mathcal{D}_y$ converges. The torus is finite, so $\lambda$ is a bounded
operator, and the minimiser has finite energy, so $\sum_y\langle\beta_y,\lambda\beta_y\rangle<\infty$.
If $a<\infty$, then by parts (a) and (b) of the model theorem on the exit value the potentials on
the free levels $y'\ge a$ form the tail $\beta_{a+j}=\mathsf{C}^{j}\beta_a$, where
\[
  \mathsf{C}=1-\mathsf{v}=(1+\lambda+\mathsf{v})^{-1}\in(0,1]
\]
and $\mathsf{v}\in[0,1)$ is the
non-negative root of $\mathsf{v}^2+\lambda\mathsf{v}-\lambda=0$; the operator $\mathsf{C}$ is
strictly less than one on every mode on which $\lambda>0$, so the terms $\lambda\beta_{y'}$
decay geometrically. If $a=\infty$, let $a'<a''$ be interior levels, the $\delta$ of $a'$ being
$\delta'$. Apply \eqref{8.15:eq-band-optimality-2} with $y=a'$, with $a''$ in place of $a'$, and
with $\pm h$, $\|h\|<\delta'$. Equality results, and by linearity in $h$
\[
  \sum_{a'<y'\le a''}\lambda\beta_{y'}=PE_{a''}-E_{a'} .
\]
The right side tends to $0$ as $a',a''\to\infty$, since $\|E_{y'}\|\to0$. So the partial sums
are Cauchy in the finite-dimensional space $\mathfrak{H}_T$, and the series converges.

We now derive the optimality conditions. In \eqref{8.15:eq-band-optimality-2} let
$a'\to\infty$ through interior levels. Then $PE_{a'}\to0$, and the partial sums tend to
$-\mathcal{D}_y$. Hence $\langle h,E_y-\mathcal{D}_y\rangle\ge0$ for every allowed $h$. In this
sense the partial gradient of $J$ in $E_y$ along variations of finite energy is
\[
  E_y+P\sum_{y'>y}\lambda\beta_{y'}=E_y-\mathcal{D}_y .
\]
For $a_0\le y<a$ and $E'\in\mathfrak{B}$ take $h=E'-E_y$. This gives
$\langle\mu_y,E'-E_y\rangle\le0$, that is, $\mu_y=-(E_y-\mathcal{D}_y)\in
\mathcal{N}_{\mathfrak{B}}(E_y)$. For $y\ge a$ every $h$ and $-h$ are allowed, so
$\langle h,E_y-\mathcal{D}_y\rangle=0$ for all $h$. Thus $E_y=\mathcal{D}_y$ and $\mu_y=0$. This
proves (ii).

For (i), we use again that $\lambda\beta\in\mathfrak{H}_T$. The space $\mathfrak{H}_T$ is a closed
subspace, so the convergent sum $\mathcal{D}_y$ lies in $\mathfrak{H}_T$.

For (iii), the second identity in \eqref{8.15:eq-band-optimality-a} is the definition of
$\mathcal{D}$. For $y>a_0$ both $\mathcal{D}_y$ and $\mathcal{D}_{y-1}$ are defined, and
\[
  \mathcal{D}_y-\mathcal{D}_{y-1}
  =-\sum_{y'>y}\lambda\beta_{y'}+\sum_{y'>y-1}\lambda\beta_{y'}=\lambda\beta_y .
\]
The first identity in \eqref{8.15:eq-band-optimality-a} is the kinematic recursion of the model.
For $y>a_0$, apply the second identity at $y+1$ and at $y$ and subtract. This gives
$\lambda\beta_{y+1}-\lambda\beta_y$, which equals $\lambda PE_y$ by the first identity. This is
the first equality in \eqref{8.15:eq-band-optimality-1}. The second equality is the
identification, in the model, of $\lambda P$ with $\operatorname{curl}^{*}\operatorname{curl}$ on
lateral $1$-forms; it is the operator form of the relation
$\|\operatorname{curl}\alpha\|^2=\langle P\alpha,\lambda P\alpha\rangle$ defining $\lambda$.

For (iv), let $a_0\le y<a$. By (ii),
\[
  \langle\mathcal{D}_y-E_y,E'-E_y\rangle\le0\quad\text{for all }E'\in\mathfrak{B}.
\]
This is the variational characterisation of the nearest point of the closed convex set
$\mathfrak{B}$ to $\mathcal{D}_y$. Hence $E_y=\operatorname{clip}_\tau(\mathcal{D}_y)$, and
$\mu_y=\mathcal{D}_y-E_y=\operatorname{shrink}_\tau(\mathcal{D}_y)$.

For (v), we have $P^{\perp}\mathcal{D}_y=0$ by (i). Applying $P^{\perp}$ to
$\mu_y=\mathcal{D}_y-E_y$ gives $P^{\perp}\mu_y=-P^{\perp}E_y$. In the same way, since
$\mathcal{D}_y\in\mathfrak{H}_T\subset\ker\nabla^{*}$, we obtain $\nabla^{*}\mu_y=-\nabla^{*}E_y$.
\end{proof}

\begin{proposition}[Model: convexity of the multiplier norm in the height]
\label{8.15:prop-multiplier-convexity}
The model is the exit-value model of the exit-value model theorem; the notation is that of the
model lemma, Lemma~\ref{8.15:lem-band-optimality}. For $a_0\le a\le\infty$,
$(E_y,\beta_y)_{y\ge a_0}$ is the minimiser of the problem $P_a$. In $P_a$ the levels
$a_0\le y<a$ are capped. We work in the box case of the model: on a capped level the lateral
$1$-form $E_y$ on the finite torus is constrained to the feasible set $\mathfrak{B}$, which is
the componentwise box $|E_{y,k}(x)|\le\tau_k$ for every lateral site $x$ and direction $k$.
In $P_\infty$ every level $y\ge a_0$ is capped. The remaining objects are those of
Lemma~\ref{8.15:lem-band-optimality}: $\lambda=\operatorname{curl}^*\operatorname{curl}+\nabla\nabla^*$,
the transversal and longitudinal projections $P$ and $P^{\perp}$,
$\mathcal{D}_y:=-\sum_{y'>y}\lambda\beta_{y'}$, the multiplier $\mu_y:=\mathcal{D}_y-E_y$ and the
charge $q_y:=\nabla^*\mu_y$. Norms and pairings are $\ell^2$ over sites and directions, and
$x'\sim x$ means that $x'$ is a nearest neighbour of $x$. We put
$f(y):=\tfrac12\|\mu_y\|^2$, and we let $s_{y,k}(x)\in\{-1,0,1\}$ be the sign of
$\mu_{y,k}(x)$.

\emph{Second differences.} Let $y-1,y,y+1$ be levels that are capped in $P_a$, with
$y-1\ge a_0$. Then
\begin{equation}\label{8.15:eq-multiplier-convexity-1}
f(y+1)-2f(y)+f(y-1)\ \ge\ \mathsf{A}_y+\|q_y\|^2,
\end{equation}
where
\begin{equation}\label{8.15:eq-multiplier-convexity-2}
\mathsf{A}_y:=\sum_{x,k}|\mu_{y,k}(x)|\sum_{x'\sim x}\bigl[\tau_k-s_{y,k}(x)\,E_{y,k}(x')\bigr]
\ \ge\ 0 .
\end{equation}

\emph{The problem $P_\infty$.} For $P_\infty$ the function $f$ is convex and non-increasing on
the levels $y\ge a_0$. Moreover, for every $y_1\ge a_0$,
\begin{equation}\label{8.15:eq-multiplier-convexity-3}
\sum_{y>y_1}(y-y_1)\bigl(\mathsf{A}_y+\|q_y\|^2\bigr)\ \le\ \tfrac12\|\mu_{y_1}\|^2 .
\end{equation}

\emph{Localised form.} Let $\chi\ge0$ be a lateral weight, and put
$f_\chi(y):=\tfrac12\sum_x\chi(x)|\mu_y(x)|^2$. Then, for $y$ as in the first part,
\begin{equation}\label{8.15:eq-multiplier-convexity-4}
\begin{split}
f_\chi(y+1)-2f_\chi(y)+f_\chi(y-1)\ \ge\ &\sum_{x,k}\chi(x)|\mu_{y,k}(x)|
\sum_{x'\sim x}\bigl[\tau_k-s_{y,k}(x)E_{y,k}(x')\bigr]\\
&+\sum_x\chi(x)\,q_y(x)^2-\sum|\nabla\chi|\,|\mu_y|\,|q_y| .
\end{split}
\end{equation}
\end{proposition}

\begin{proof}
\emph{Convexity of the level functional.} For $\tau>0$ put
$\varphi_\tau(t):=\tfrac12\bigl((|t|-\tau)_+\bigr)^2$ for $t\in\mathbb{R}$. Here
$\operatorname{clip}_\tau(t):=\max(-\tau,\min(\tau,t))$ and
$\operatorname{shrink}_\tau(t):=t-\operatorname{clip}_\tau(t)$.

The function $\varphi_\tau$ is $C^1$ with $\varphi_\tau'=\operatorname{shrink}_\tau$. Indeed,
$\varphi_\tau'(t)=0$ for $|t|\le\tau$, and $\varphi_\tau'(t)=\operatorname{sign}(t)(|t|-\tau)$ for
$|t|>\tau$. This derivative is continuous and non-decreasing, so $\varphi_\tau$ is convex.

By Lemma~\ref{8.15:lem-band-optimality}, a model lemma, on a capped level we have
$E_y=\operatorname{clip}_\tau(\mathcal{D}_y)$ and $\mu_y=\operatorname{shrink}_\tau(\mathcal{D}_y)$,
componentwise with cap $\tau_k$. Hence $|\mu_{y,k}(x)|=(|\mathcal{D}_{y,k}(x)|-\tau_k)_+$, and
\begin{equation*}
f(y)=\sum_{x,k}\varphi_{\tau_k}\bigl(\mathcal{D}_{y,k}(x)\bigr).
\end{equation*}
The levels $y\pm1$ are capped as well, so $f(y\pm1)$ has the same representation. The tangent
inequality $\varphi_{\tau_k}(b)\ge\varphi_{\tau_k}(a)+\varphi_{\tau_k}'(a)(b-a)$ is applied at
$a=\mathcal{D}_{y,k}(x)$ and $b=\mathcal{D}_{y\pm1,k}(x)$, and summed over $x,k$. This gives
\begin{equation*}
f(y\pm1)-f(y)\ \ge\ \langle\mu_y,\mathcal{D}_{y\pm1}-\mathcal{D}_y\rangle .
\end{equation*}
We add the two inequalities. Since $y>a_0$, the second-difference identity that the model
lemma, Lemma~\ref{8.15:lem-band-optimality}, derives from \eqref{8.15:eq-band-optimality-a} gives
$\mathcal{D}_{y+1}-2\mathcal{D}_y+\mathcal{D}_{y-1}=\lambda PE_y$. Hence
\begin{equation*}
f(y+1)-2f(y)+f(y-1)\ \ge\ \langle\mu_y,\lambda PE_y\rangle .
\end{equation*}

\emph{The charge.} We write $\lambda PE_y=\lambda E_y-\lambda P^{\perp}E_y$. By the model lemma,
Lemma~\ref{8.15:lem-band-optimality}, $\mathcal{D}_y$ lies in the transversal space
$\mathfrak{H}_T$. Hence $P^{\perp}\mathcal{D}_y=0$, and $\mu_y=\mathcal{D}_y-E_y$ gives
$P^{\perp}E_y=-P^{\perp}\mu_y$. Therefore
\begin{equation*}
\langle\mu_y,\lambda PE_y\rangle
=\langle\mu_y,\lambda E_y\rangle-\langle\mu_y,\lambda P^{\perp}E_y\rangle
=\langle\mu_y,\lambda E_y\rangle+\langle\mu_y,\lambda P^{\perp}\mu_y\rangle .
\end{equation*}

Next we show that $\lambda P^{\perp}=\nabla\nabla^*$. The inputs are four facts:
$\lambda=\operatorname{curl}^*\operatorname{curl}+\nabla\nabla^*$, $\operatorname{curl}P^{\perp}=0$,
$\nabla^*P=0$, and $\lambda$ annihilates constant $1$-forms.
\begin{itemize}
\item The first two facts give $\lambda P^{\perp}=\nabla\nabla^*P^{\perp}$.
\item The third fact says that $\nabla\nabla^*$ vanishes on the range of $P$.
\item For a constant $1$-form $c$ we have
\begin{equation*}
0=\langle c,\lambda c\rangle=\|\operatorname{curl}c\|^2+\|\nabla^*c\|^2,
\end{equation*}
so $\nabla\nabla^*c=0$.
\end{itemize}
In the model every lateral $1$-form is the sum of an element of the range of $P$, an element of
the range of $P^{\perp}$ and a constant $1$-form, so the operator $1-P^{\perp}$ maps into the
span of the range of $P$ and the constants.
Hence $\nabla\nabla^*P^{\perp}=\nabla\nabla^*$. It follows that the last term above is
\begin{equation*}
\langle\mu_y,\nabla\nabla^*\mu_y\rangle=\|\nabla^*\mu_y\|^2=\|q_y\|^2 .
\end{equation*}

\emph{Absorption.} In components,
\begin{equation*}
\langle\mu,\lambda E\rangle=\sum_{x,k}\mu_k(x)\sum_{x'\sim x}\bigl[E_k(x)-E_k(x')\bigr].
\end{equation*}
Let $x$ lie in the support of $\mu_{y,k}$. There $E_{y,k}(x)=s_{y,k}(x)\tau_k$ and
$\mu_{y,k}(x)=s_{y,k}(x)|\mu_{y,k}(x)|$, with $s_{y,k}(x)^2=1$. Hence
\begin{equation*}
\mu_{y,k}(x)\bigl[E_{y,k}(x)-E_{y,k}(x')\bigr]
=|\mu_{y,k}(x)|\bigl[\tau_k-s_{y,k}(x)E_{y,k}(x')\bigr].
\end{equation*}
Sites outside the support contribute nothing, so $\langle\mu_y,\lambda E_y\rangle=\mathsf{A}_y$.
Since the level $y$ is capped, $|E_{y,k}(x')|\le\tau_k$ for every site $x'$. Each bracket in
\eqref{8.15:eq-multiplier-convexity-2} is therefore non-negative, and so $\mathsf{A}_y\ge0$.
Collecting the three steps proves \eqref{8.15:eq-multiplier-convexity-1}.

\emph{The problem $P_\infty$.} Every level $y\ge a_0$ is capped, so
\eqref{8.15:eq-multiplier-convexity-1} holds at every $y>a_0$. Its right-hand side is
non-negative, so $f$ is convex on the levels $y\ge a_0$. Also $f\ge0$.

In the exit-value model the front is finite: there is a finite level, the front, from which on
the multipliers $\mu_y$ vanish; this is the remark on the finiteness of the front in the section
of the exit-value model theorem. Hence $f$ vanishes from the front on.

Suppose $f(y_0+1)>f(y_0)$ for some $y_0\ge a_0$. The increments $f(y+1)-f(y)$ are non-decreasing
in $y$, so for every $y>y_0$
\begin{equation*}
f(y)\ge f(y_0+1)+(y-y_0-1)\bigl(f(y_0+1)-f(y_0)\bigr)>0 .
\end{equation*}
This contradicts the vanishing of $f$ beyond the front, so $f$ is non-increasing.

For the summation by parts, put $\delta_y:=f(y+1)-f(y)$, so that
$f(y+1)-2f(y)+f(y-1)=\delta_y-\delta_{y-1}$. For $Y>y_1\ge a_0$,
\begin{equation}\label{8.15:eq-multiplier-convexity-5}
\begin{split}
\sum_{y_1<y\le Y}(y-y_1)(\delta_y-\delta_{y-1})
&=(Y-y_1)\delta_Y-\sum_{y=y_1}^{Y-1}\delta_y\\
&=(Y-y_1)\bigl[f(Y+1)-f(Y)\bigr]-f(Y)+f(y_1).
\end{split}
\end{equation}
Choose $Y$ beyond both the front and $y_1$. Then $f(Y)=f(Y+1)=0$. For $y>Y$ we have $\mu_y=0$,
so $\mathsf{A}_y=0$ and $q_y=0$. Bounding each second difference from below by
\eqref{8.15:eq-multiplier-convexity-1} in \eqref{8.15:eq-multiplier-convexity-5} yields
\eqref{8.15:eq-multiplier-convexity-3}.

\emph{Localised form.} Write $f_\chi(y)=\sum_{x,k}\chi(x)\varphi_{\tau_k}(\mathcal{D}_{y,k}(x))$.
Multiply the tangent inequality at $(x,k)$ by $\chi(x)\ge0$ and add. The second difference of
$f_\chi$ is then at least
\begin{equation*}
\sum_x\chi(x)\bigl\langle\mu_y(x),(\mathcal{D}_{y+1}-2\mathcal{D}_y+\mathcal{D}_{y-1})(x)\bigr\rangle .
\end{equation*}

By the model lemma, Lemma~\ref{8.15:lem-band-optimality}, this second difference of
$\mathcal{D}$ equals
$\operatorname{curl}^*\operatorname{curl}E_y$. The operator
$\operatorname{curl}^*\operatorname{curl}=\lambda-\nabla\nabla^*$ is local. By the charge step,
\begin{equation*}
\operatorname{curl}^*\operatorname{curl}E_y=\lambda PE_y=\lambda E_y+\nabla\nabla^*\mu_y .
\end{equation*}

The term with $\lambda E_y$ is treated as in the absorption step with the weight $\chi(x)$. It
gives the first sum in \eqref{8.15:eq-multiplier-convexity-4}.

For the other term, $\nabla^*$ is the adjoint of $\nabla$, so
\begin{equation*}
\sum_x\chi\langle\mu_y,\nabla(\nabla^*\mu_y)\rangle
=\sum_x\bigl(\nabla^*(\chi\mu_y)\bigr)(\nabla^*\mu_y)
=\sum_x\chi\,q_y^2+\sum_x\bigl([\nabla^*,\chi]\mu_y\bigr)q_y .
\end{equation*}
Here $[\nabla^*,\chi]\mu:=\nabla^*(\chi\mu)-\chi\nabla^*\mu$. The commutator term is bounded in
absolute value by $\sum|\nabla\chi||\mu_y||q_y|$. This gives
\eqref{8.15:eq-multiplier-convexity-4}.
\end{proof}

\begin{remark}
The inequality admits three readings.

First, in the scalar treatment of the exit problem the charge appears as an unsigned
derivative term in the component equations for $E_k$.
Paired with $\mu$ itself, it is a square; this is the only pairing that keeps the Gauss
constraint. Inequality \eqref{8.15:eq-multiplier-convexity-3} is an $\ell^2$ form of the
statement that charged activity costs: the height-weighted total charge above a level $y_1$ is
bounded by the multiplier norm on that level.

Second, for weights other than $\mu$ the cross term $\langle\nabla^*\nu,\nabla^*\mu\rangle$ is
unsigned. This holds for any other element $\nu$ of the normal cone, and for $\ell^p$ weights
with $p\ne2$. That only $p=2$ pairs with the energy is thus the Gauss law.

Third, Proposition~\ref{8.15:prop-multiplier-convexity} is the second-order, single-minimiser
counterpart of clauses (d)--(e) of the exit-value model theorem. Like them it is global in $x$,
or local with a commutator term, and it gives no bound in the supremum norm. If a test field is
paired with the whole co-closed $2$-form $\hat{F}+\Lambda\hat{F}/\tau$,
the minimising curvature $\hat{F}$ plus its multiplier $2$-form built from the multiplier
magnitude $\Lambda\ge0$, normal bonds included,
no sign of the charge is needed: the Gauss law is then used as one of the equations rather than
treated as a source.
\end{remark}

\begin{definition}[Capped lattice data on a cube and their defects]\label{8.15:def-capped-data}
\emph{Lattice and forms.} We work on the lattice $\mathbb{Z}^4$ with unit vectors $e_1,\dots,e_4$ and
forward differences $\partial_i f(z):=f(z+e_i)-f(z)$. The arguments below do not depend on the
dimension; the number $4$ enters only as the exponent of the volume, through the power $r^4$ in
\eqref{8.15:eq-capped-data-c}. Let $\mathsf{V}$, the colour space, be a
finite-dimensional real vector space with a Euclidean inner product $\langle\cdot,\cdot\rangle$ and
norm $|\cdot|$. All forms are $\mathsf{V}$-valued.
\begin{itemize}
\item A $1$-form $\xi$ assigns a vector $\xi_i(z)$ to each bond $(z,i)$.
\item A $2$-form $X$ assigns a vector $X_{ij}(z)=-X_{ji}(z)$ to each plaquette $p=(z,ij)$ with $i<j$.
\end{itemize}
In both cases $z$ is the \emph{base site} of the cell. The exterior difference is
$(d\xi)_{ij}:=\partial_i\xi_j-\partial_j\xi_i$. The boundary of $p=(z,ij)$ is the set of four bonds
\[
\partial p:=\{(z,i),\,(z+e_i,j),\,(z+e_j,i),\,(z,j)\},
\]
and we write $|\xi|_{\partial p}$ for the largest of the numbers $|\xi(b)|$, $b\in\partial p$.

\emph{Cubes and the spaces $\mathfrak{X}_r$.} Fix $c\in\mathbb{Z}^4$ and put
$Q_r=Q_r(c):=\{z\in\mathbb{Z}^4:|z-c|_\infty\le r\}$. A cell (bond or plaquette) is \emph{in} $Q_r$ if
its base site lies in $Q_r$. For $2$-forms we put
\[
\|X\|_r^2:=\sum_{p\text{ in }Q_r}|X(p)|^2,\qquad
\langle X,Y\rangle:=\sum_p\langle X(p),Y(p)\rangle,
\]
the second sum being taken whenever one of the two forms is finitely supported. We let
$\mathfrak{X}_r$ be the space of $1$-forms supported on the bonds having at least one endpoint in
$Q_{r-1}$. The two endpoints of a bond differ by a unit vector, so both endpoints of such a bond lie
in $Q_r$; the bonds lying in the boundary layer $Q_r\setminus Q_{r-1}$ are excluded. For
$\xi\in\mathfrak{X}_r$ the form $d\xi$ is supported on plaquettes in $Q_r$. Indeed, if
$(d\xi)(p)\neq0$, then $\xi$ is non-zero on some bond of $\partial p$, and that bond has an endpoint,
a vertex of $p$, in $Q_{r-1}$. Every vertex of $p$ differs from
its base site by a vector with entries in $\{0,1\}$, and therefore the base site lies in $Q_r$.

\emph{Constant $2$-forms.} A \emph{constant} $2$-form $F_0$ is given by one vector
$F_{0,ij}\in\mathsf{V}$ for each orientation $i<j$. We write $F_0(p):=F_{0,ij}$ for plaquettes $p$
of orientation $ij$, and $|F_0|:=\max_{i<j}|F_{0,ij}|$. For every finitely supported $1$-form $\xi$,
\[
\langle F_0,d\xi\rangle=\sum_{i<j}\Bigl\langle F_{0,ij},
\sum_{z}\bigl(\partial_i\xi_j(z)-\partial_j\xi_i(z)\bigr)\Bigr\rangle=0,
\]
because the sum over $z\in\mathbb{Z}^4$ of a forward difference of a finitely supported function
telescopes to zero.

\emph{Capped data.} Let $R\ge1$ be an integer and let $\tau>0$. \emph{Capped lattice data on $Q_{3R}$
with cap $\tau$} consist of the following objects:
\begin{itemize}
\item a $1$-form $A$;
\item a $2$-form $\hat{F}$;
\item a function $\Lambda\ge0$ on plaquettes;
\item a non-decreasing function $\varpi:[1,R]\to[0,1]$, the \emph{defect function}.
\end{itemize}
These objects are given on $Q_{3R}$, and for every plaquette $p$ in $Q_{3R}$ they satisfy
\begin{equation}\label{8.15:eq-capped-data-a}
|\hat{F}(p)|\le\tau,\qquad
\Lambda(p)=0\ \text{ unless }\ |\hat{F}(p)|=\tau.
\end{equation}
The cap in the first condition is imposed on plaquettes of every orientation. The second condition
is the complementarity of the multiplier magnitude $\Lambda$ with the cap. We set
\[
\mu(p):=\frac{\Lambda(p)\hat{F}(p)}{\tau}
\]
and call the $2$-form $\mu$ the \emph{multiplier $2$-form}.

The data are further required to satisfy, for every $1\le r\le R$, the two \emph{defect bounds at
scale $r$}:
\begin{equation}\label{8.15:eq-capped-data-b}
\begin{gathered}
|\hat{F}(p)-(dA)(p)|\le\varpi(r)\,\tau\qquad\text{for every plaquette }p\text{ in }Q_{3r},\\
|\langle\hat{F}+\mu,d\xi\rangle|\le\frac{\varpi(r)}{r}\sum_{p\text{ in }Q_{3r}}
\bigl(|\hat{F}(p)|+\Lambda(p)\bigr)\,|\xi|_{\partial p}
\qquad\text{for every }\xi\in\mathfrak{X}_{3r}.
\end{gathered}
\end{equation}
The first bound measures how far $\hat{F}$ is from being the curvature $dA$. The second measures how
far $\hat{F}+\mu$ is from being orthogonal to the exact forms $d\xi$, $\xi\in\mathfrak{X}_{3r}$. Since
$d\xi$ is supported on plaquettes in $Q_{3r}\subset Q_{3R}$, the pairing on the left is a finite sum
that involves only the values of the data on $Q_{3R}$.

\emph{The abelian case.} This is the case in which $\mathsf{V}=\mathbb{R}$, $\hat{F}=d\hat{A}$ for a
$1$-form $\hat{A}$, and $\varpi\equiv0$. With $A=\hat{A}$, the first bound in
\eqref{8.15:eq-capped-data-b} then holds with $\varpi\equiv0$. The second bound, imposed for all
$1\le r\le R$, states that $\langle\hat{F}+\mu,d\xi\rangle=0$ for every $\xi\in\mathfrak{X}_{3R}$.
If $|\hat{F}(p)|=\tau$, then $\hat{F}(p)/\tau$ is the sign of $\hat{F}(p)$. The second bound is
therefore the multiplier identity of the capped abelian minimisation problem in the following
configuration:
\begin{itemize}
\item every bond with an endpoint in $Q_{3R-1}$ is free, that is, not among the bonds whose values
are held fixed in the minimisation;
\item the cap function of the problem, which bounds $|\hat{F}(p)|$ plaquette by plaquette, equals
$\tau$ on every plaquette in $Q_{3R}$.
\end{itemize}

\emph{Hybrid excess.} For a constant $2$-form $F_0$ we define the \emph{hybrid excess} and the
\emph{normalised excess}
\begin{equation}\label{8.15:eq-capped-data-c}
\mathcal{Y}(r;F_0):=\sum_{p\text{ in }Q_r}\Bigl[\tfrac12|\hat{F}(p)-F_0(p)|^2+\tau\Lambda(p)\Bigr],
\qquad
y(r;F_0):=\frac{\mathcal{Y}(3r;F_0)}{\tau^2r^4}.
\end{equation}
\end{definition}

\begin{lemma}[Gauge-invariant translation potentials]\label{8.15:lem-translation-potentials}
We work on $\mathbb{Z}^4$ with the conventions of Definition~\ref{8.15:def-capped-data}, of which we
recall what is used here.
\begin{itemize}
\item $\partial_\mu f(z):=f(z+e_\mu)-f(z)$.
\item A $1$-form $\xi$ takes values $\xi_\mu(z)$ on bonds $(z,\mu)$.
\item A $2$-form $X$ takes values $X_{\mu\nu}(z)=-X_{\nu\mu}(z)$ on plaquettes with base site $z$; in
particular $X_{\nu\nu}=0$.
\item $(d\xi)_{\mu\nu}:=\partial_\mu\xi_\nu-\partial_\nu\xi_\mu$, and for a function $\lambda$ on sites
we write $(d\lambda)_\nu:=\partial_\nu\lambda$.
\item A cell is in $Q_s$ if its base site is, and $\|X\|_s^2$ is the sum of $|X(p)|^2$ over the
plaquettes $p$ in $Q_s$.
\end{itemize}
Let $r\ge1$ be an integer, let $A$ be a $1$-form on $Q_{3r}$ (in the application, $1\le r\le R$ and
$A$ is the $1$-form of the data of Definition~\ref{8.15:def-capped-data}), and
put $F':=dA$ on $Q_{3r}$. Let $\kappa\ge1$ be an integer, let $K$ be the uniform probability on
$\mathsf{K}:=[-\kappa,\kappa]^4\cap\mathbb{Z}^4$, and for a $2$-form $X$ put
$X_K(p):=\sum_{h}K(h)X(p+h)$, where $p+h$ is the translate of $p$ by $h$. For $h\in\mathsf{K}$ let
$h^{(i)}:=(h_1,\dots,h_i,0,\dots,0)$, so that $h^{(0)}=0$. For a $2$-form $X$ define
$\beta_{h,\nu}[X](z):=\sum_{i=1}^4\sigma_i$, where
\begin{equation*}
\sigma_i:=
\begin{cases}
\sum_{t=0}^{h_i-1}X_{i\nu}\bigl(z+h^{(i-1)}+te_i\bigr), & h_i>0,\\[2pt]
-\sum_{t=1}^{|h_i|}X_{i\nu}\bigl(z+h^{(i-1)}-te_i\bigr), & h_i<0,\\[2pt]
0, & h_i=0.
\end{cases}
\end{equation*}
Put $\beta_K[X]:=\sum_hK(h)\beta_h[X]$, $\beta_h:=\beta_h[F']$ and $\beta_K:=\beta_K[F']$.

Then, wherever the cells involved lie in $Q_{3r}$, the following hold. The $1$-form $\beta_K[X]$ is
linear in $X$, and $\beta_K[F_0]=0$ for every constant $2$-form $F_0$. Moreover, for every constant
$2$-form $F_0$ and every $s$,
\begin{equation}\label{8.15:eq-translation-potentials-a}
\begin{aligned}
&\text{(t1)}\quad d\beta_h=F'(\cdot+h)-F',\qquad d\beta_K=F'_K-F';\\
&\text{(t2)}\quad \beta_K[F']=\beta_K[F'-F_0];\\
&\text{(t3)}\quad |\beta_{K,\nu}(z)|\le4\kappa\,\sup|F'-F_0|;\\
&\text{(t4)}\quad \Bigl(\sum_{b\text{ in }Q_s}|\beta_K(b)|^2\Bigr)^{1/2}\le8\kappa\,\|F'-F_0\|_{s+\kappa}.
\end{aligned}
\end{equation}
In \textup{(t3)} the supremum is over the plaquettes whose base site lies within
$|\cdot|_\infty$-distance $\kappa$ of $z$.
\end{lemma}

\begin{proof}
\emph{The one-step identity.} By definition of $d$,
\begin{equation*}
F'_{i\nu}(z)=\partial_iA_\nu(z)-\partial_\nu A_i(z)=A_\nu(z+e_i)-A_\nu(z)-\partial_\nu A_i(z).
\end{equation*}
Hence
\begin{equation}\label{8.15:eq-translation-potentials-1}
A_\nu(z+e_i)-A_\nu(z)=F'_{i\nu}(z)+\partial_\nu A_i(z).
\end{equation}
For $i=\nu$ this reads $A_\nu(z+e_\nu)-A_\nu(z)=\partial_\nu A_\nu(z)$, since $F'_{\nu\nu}=0$.
Replacing $z$ by $z-e_i$ in \eqref{8.15:eq-translation-potentials-1} gives
\begin{equation}\label{8.15:eq-translation-potentials-2}
A_\nu(z-e_i)-A_\nu(z)=-F'_{i\nu}(z-e_i)-\partial_\nu A_i(z-e_i).
\end{equation}

\emph{Telescoping along the staircase.} Go from $z$ to $z+h$ through the points
$z+h^{(1)},\dots,z+h^{(4)}=z+h$. Put $z^{(i-1)}:=z+h^{(i-1)}$. The $i$-th leg runs from $z^{(i-1)}$
to $z+h^{(i)}=z^{(i-1)}+h_ie_i$.

If $h_i>0$, write
\begin{equation*}
A_\nu(z^{(i-1)}+h_ie_i)-A_\nu(z^{(i-1)})
=\sum_{t=0}^{h_i-1}\bigl[A_\nu(z^{(i-1)}+(t+1)e_i)-A_\nu(z^{(i-1)}+te_i)\bigr].
\end{equation*}
Applying \eqref{8.15:eq-translation-potentials-1} at each $z^{(i-1)}+te_i$ turns this into
\begin{equation*}
\sigma_i+\sum_{t=0}^{h_i-1}\partial_\nu A_i(z^{(i-1)}+te_i).
\end{equation*}

If $h_i<0$, write
\begin{equation*}
A_\nu(z^{(i-1)}+h_ie_i)-A_\nu(z^{(i-1)})
=\sum_{t=1}^{|h_i|}\bigl[A_\nu(z^{(i-1)}-te_i)-A_\nu(z^{(i-1)}-(t-1)e_i)\bigr].
\end{equation*}
Applying \eqref{8.15:eq-translation-potentials-2} at the points $z^{(i-1)}-(t-1)e_i$ turns this into
\begin{equation*}
\sigma_i-\sum_{t=1}^{|h_i|}\partial_\nu A_i(z^{(i-1)}-te_i).
\end{equation*}

If $h_i=0$, the leg contributes $0=\sigma_i$.

Let $\lambda_h(z)$ be the corresponding signed sum of the $A_i$: it is the sum over $i$ of
$\sum_{t=0}^{h_i-1}A_i(z+h^{(i-1)}+te_i)$ if $h_i>0$, of
$-\sum_{t=1}^{|h_i|}A_i(z+h^{(i-1)}-te_i)$ if $h_i<0$, and of $0$ if $h_i=0$. The difference
$\partial_\nu$ commutes with translations. Summing the four legs therefore gives
\begin{equation}\label{8.15:eq-translation-potentials-3}
A_\nu(z+h)-A_\nu(z)=\beta_{h,\nu}(z)+\partial_\nu\lambda_h(z),
\qquad\text{that is,}\qquad A(\cdot+h)-A=\beta_h+d\lambda_h.
\end{equation}

\emph{Proof of \textup{(t1)}.} Apply $d$ to \eqref{8.15:eq-translation-potentials-3}. First,
$d(A(\cdot+h))=(dA)(\cdot+h)=F'(\cdot+h)$, because $d$ commutes with translations. Second,
$(dd\lambda_h)_{\mu\nu}=\partial_\mu\partial_\nu\lambda_h-\partial_\nu\partial_\mu\lambda_h=0$, because
forward differences commute. Hence $d\beta_h=F'(\cdot+h)-F'$. Multiplying by $K(h)$ and summing over
$h$, with $\sum_hK(h)=1$, gives $d\beta_K=F'_K-F'$.

\emph{Proof of \textup{(t2)}.} Each $\sigma_i$ is linear in $X$, so $\beta_h[X]$ and $\beta_K[X]$
are linear in $X$.

Let $X\equiv F_0$ be constant. If $h_i>0$, then $\sigma_i$ has $h_i$ terms equal to $F_{0,i\nu}$. If
$h_i<0$, then $\sigma_i=-|h_i|F_{0,i\nu}=h_iF_{0,i\nu}$. Hence
$\beta_{h,\nu}[F_0]=\sum_ih_iF_{0,i\nu}$, which is odd in $h$. The cube $\mathsf{K}$ is symmetric
under $h\mapsto-h$ and $K$ is uniform, so $K$ is even. Therefore $\beta_K[F_0]=0$.

By linearity,
\begin{equation*}
\beta_K[F']=\beta_K[F'-F_0]+\beta_K[F_0]=\beta_K[F'-F_0].
\end{equation*}

\emph{Proof of \textup{(t3)}.} Put $X:=F'-F_0$. By \textup{(t2)}, $\beta_K=\sum_hK(h)\beta_h[X]$.

The sum $\sigma_i$ has $|h_i|$ terms. Each is $\pm X_{i\nu}(y)$ with $y=z+h^{(i-1)}\pm te_i$ and
$0\le t\le|h_i|$. The coordinates of $y-z$ are $h_1,\dots,h_{i-1},\pm t,0,\dots,0$, so
$|y-z|_\infty\le\kappa$. Thus $\beta_h[X]$ has $|h|_1:=\sum_{i=1}^4|h_i|$ terms, each bounded by the
supremum in \textup{(t3)}. Since $|h|_1\le4\kappa$ for $h\in\mathsf{K}$, we get
\begin{equation*}
|\beta_{K,\nu}(z)|\le\sum_hK(h)\,|h|_1\sup|X|\le4\kappa\sup|F'-F_0|.
\end{equation*}

\emph{Proof of \textup{(t4)}.} Again put $X:=F'-F_0$ and use \textup{(t2)}. Fix $h$ and $\nu$. As
just shown, the function $z\mapsto\beta_{h,\nu}[X](z)$ on $Q_s$ is a sum of $|h|_1$ functions of
the form $z\mapsto\pm X_{i\nu}(z+a)$ with $|a|_\infty\le\kappa$.

For a function $f$ on the sites of $Q_s$ put
$\|f\|_{\ell^2(Q_s)}:=\bigl(\sum_{z\in Q_s}|f(z)|^2\bigr)^{1/2}$.
For $z\in Q_s$ the point $z+a$ lies in $Q_{s+\kappa}$, and $z\mapsto z+a$ is injective. The
$\ell^2(Q_s)$-norm of such a function is therefore at most the $\ell^2$-norm of $X$ over the
plaquettes of orientation $\{i,\nu\}$ in $Q_{s+\kappa}$. This is at most $\|X\|_{s+\kappa}$, and the
function vanishes when $i=\nu$.

By Minkowski's inequality,
\begin{equation*}
\|\beta_{K,\nu}\|_{\ell^2(Q_s)}\le\sum_hK(h)\,|h|_1\,\|X\|_{s+\kappa}\le4\kappa\,\|X\|_{s+\kappa}.
\end{equation*}
Summing the squares over the four components $\nu$ gives
\begin{equation*}
\sum_{b\text{ in }Q_s}|\beta_K(b)|^2=\sum_{\nu=1}^4\|\beta_{K,\nu}\|_{\ell^2(Q_s)}^2
\le4\,(4\kappa)^2\,\|X\|_{s+\kappa}^2.
\end{equation*}
Taking square roots yields \textup{(t4)}.
\end{proof}

\begin{remark}
By construction, $\beta_h$ and $\beta_K$ depend on $A$ only through $F'=dA$. They are therefore
unchanged when $A$ is replaced by $A+d\lambda$ for a function $\lambda$ on sites, because
$dd\lambda=0$. In this sense $\beta_K$ is gauge invariant. By \textup{(t1)}, $\beta_K$ is a potential
for the difference between $F'_K$, an average of translates of $F'$, and $F'$ itself. In the model
theorem on the exit value this difference is the averaging operator $\mathsf{C}$ applied to the datum
of that model, minus the datum; \eqref{8.15:eq-translation-potentials-a}\,(t1) thus supplies that
difference with a gauge-invariant potential.
\end{remark}

\begin{lemma}[Multiplier mass and oscillation bounded by the excess]\label{8.15:lem-multiplier-mass}
Let $R\ge 1$ be an integer and $\tau>0$. Let $A$, $\hat{F}$, $\Lambda$, $\varpi$ be data on
$Q_{3R}$ as in Definition~\ref{8.15:def-capped-data}, satisfying \eqref{8.15:eq-capped-data-a}
and \eqref{8.15:eq-capped-data-b}. Let $\mu$ be the multiplier $2$-form and $\mathcal{Y}$ the
hybrid excess of that definition, see \eqref{8.15:eq-capped-data-c}. Let $4\le r\le R$, let
$\kappa$ be an integer with $1\le\kappa\le r-2$, and put $\varepsilon:=\kappa/r$ and
$\varpi:=\varpi(r)$. Assume $\varpi\le 1/32$. Let $F_0$ be a constant $2$-form with
$|F_0|\le\tfrac34\tau$, and put $D':=dA-F_0$. Let $K$ be the uniform probability on
$\mathsf{K}=[-\kappa,\kappa]^4\cap\mathbb{Z}^4$ and $X_K(p):=\sum_h K(h)X(p+h)$, as in
Lemma~\ref{8.15:lem-translation-potentials}. Assume
\begin{equation}\label{8.15:eq-multiplier-mass-a}
(2\kappa+1)^{-4}\,\|D'\|_{3r}^2\le(\tau/16)^2 .
\end{equation}
Then, with an absolute constant $C_J$,
\begin{equation}\label{8.15:eq-multiplier-mass-b}
\sum_{p\text{ in }Q_r}\Big[\frac{\tau}{16}\,\Lambda(p)+\frac14\,|D'(p)-D'_K(p)|^2\Big]
\le C_J\big[\varepsilon\,\mathcal{Y}(3r;F_0)+\varpi^2\tau^2r^4\big].
\end{equation}
\end{lemma}

\begin{proof}
All cubes $Q_s$ are taken about the centre $c$ of Definition~\ref{8.15:def-capped-data}. A
plaquette or bond is in $Q_s$ when its base site is. We use $\varepsilon\le(r-2)/r<1$ and
$\varpi\le 1$ throughout. Since $\kappa\le r-2$, we have $Q_{2r+1+\kappa}\subset Q_{3r}$. The
letter $C$ denotes an absolute constant whose value may change from one occurrence to the
next.

\emph{The cutoff and the test field.} Put
\begin{equation*}
\zeta(z):=\min\big(1,(2r-|z-c|_\infty)_+/r\big).
\end{equation*}
Then $0\le\zeta\le 1$, $\zeta=1$ on $Q_r$, and $\zeta=0$ off $Q_{2r-1}$. Since $|\cdot|_\infty$
is $1$-Lipschitz for $|\cdot|_1$ and $t\mapsto\min(1,t_+/r)$ is $1/r$-Lipschitz, we have
$|\zeta(z)-\zeta(z')|\le|z-z'|_1/r$. In particular $|\partial_\mu\zeta|\le 1/r$. For a plaquette
$p$ we write $\zeta(p)$ for the value of $\zeta$ at its base site.

Let $\beta_K$ be the averaged translation potential of
Lemma~\ref{8.15:lem-translation-potentials}, built from $F':=dA$ on $Q_{3r}$ and the integer
$\kappa$. Since $\beta_K$ is linear in $F'$ and vanishes for constant $F'$
(see \eqref{8.15:eq-translation-potentials-a}), it is unchanged when $F'$ is replaced by
$F'-F_0=D'$.

Define $\xi_\mu(z):=\zeta(z)\,\beta_{K,\mu}(z)$. Since $\zeta$ vanishes off $Q_{2r-1}$, $\xi$ is
supported on bonds whose base site lies in $Q_{2r-1}$. Hence
$\xi\in\mathfrak{X}_{2r}\subset\mathfrak{X}_{3r}$, and $d\xi$ is supported on plaquettes in
$Q_{2r}$.

By its definition as an average of staircase sums, $\beta_{K,\nu}(z)$ depends only on the
values of $F'$ on plaquettes within $|\cdot|_\infty$-distance $\kappa$ of $z$. Below, $\beta_K$
is used only at base sites in $Q_{2r+1}$. It therefore depends only on the values of $D'$ on
plaquettes in
$Q_{2r+1+\kappa}\subset Q_{3r}$, where all the statements of
Lemma~\ref{8.15:lem-translation-potentials} apply.

\emph{Product rule.} We have
\begin{equation*}
\partial_\mu(\zeta\beta_{K,\nu})(z)
=\zeta(z)\,\partial_\mu\beta_{K,\nu}(z)+\partial_\mu\zeta(z)\,\beta_{K,\nu}(z+e_\mu).
\end{equation*}
Hence $(d\xi)_{\mu\nu}(z)=\zeta(z)(d\beta_K)_{\mu\nu}(z)+\gamma_{\mu\nu}(z)$, where
\begin{equation*}
\gamma_{\mu\nu}(z):=\partial_\mu\zeta(z)\,\beta_{K,\nu}(z+e_\mu)
-\partial_\nu\zeta(z)\,\beta_{K,\mu}(z+e_\nu).
\end{equation*}
If $\partial_\mu\zeta(z)\ne 0$, then $z$ or $z+e_\mu$ lies in $Q_{2r-1}$, so both lie in
$Q_{2r}$. Thus $\gamma$ is supported on plaquettes in $Q_{2r}$, and only values of $\beta_K$ on
bonds with base site in $Q_{2r}$ enter it.

By the relation $d\beta_K=F'_K-F'$ of \eqref{8.15:eq-translation-potentials-a}, valid at
plaquettes in $Q_{2r}$ since their translates by $h\in\mathsf{K}$ lie in
$Q_{2r+\kappa}\subset Q_{3r}$, we get $d\beta_K=D'_K-D'$ there. Indeed $F'=D'+F_0$, and
$(F_0)_K=F_0$ because $F_0$ is constant and $K$ is a probability. Therefore
\begin{equation}\label{8.15:eq-multiplier-mass-1}
d\xi=\zeta\,(D'_K-D')+\gamma .
\end{equation}

\emph{The identity.} Put $\mathfrak{n}:=\hat{F}-dA$ and
$\mathfrak{d}:=\langle\hat{F}+\mu,d\xi\rangle$. Then $D'=\hat{F}-F_0-\mathfrak{n}$. Since $\xi$
is finitely supported, $\langle F_0,d\xi\rangle=0$ (Definition~\ref{8.15:def-capped-data}). Hence
\begin{equation*}
\langle D'+\mu,d\xi\rangle=\langle\hat{F}+\mu,d\xi\rangle-\langle F_0,d\xi\rangle
-\langle\mathfrak{n},d\xi\rangle=\mathfrak{d}-\langle\mathfrak{n},d\xi\rangle .
\end{equation*}
Inserting \eqref{8.15:eq-multiplier-mass-1} gives
$\langle D'+\mu,\zeta(D'_K-D')+\gamma\rangle=\mathfrak{d}-\langle\mathfrak{n},d\xi\rangle$, that is,
\begin{equation}\label{8.15:eq-multiplier-mass-c}
\sum_p\zeta\langle\mu,D'-D'_K\rangle+\sum_p\zeta\langle D',D'-D'_K\rangle
=\langle D'+\mu,\gamma\rangle-\mathfrak{d}+\langle\mathfrak{n},d\xi\rangle .
\end{equation}
All terms are supported on plaquettes in $Q_{2r}$.

\emph{Step (i): the multipliers, with a good sign and a rate.} Let $p$ be in $Q_{2r}$. By
Cauchy--Schwarz, with $K$ a probability, and since the plaquettes $p+h$, $h\in\mathsf{K}$, are
distinct and lie in $Q_{3r}$, we have by \eqref{8.15:eq-multiplier-mass-a}
\begin{equation*}
|D'_K(p)|^2\le\sum_hK(h)|D'(p+h)|^2\le(2\kappa+1)^{-4}\|D'\|_{3r}^2\le(\tau/16)^2 .
\end{equation*}
Averaging $\hat{F}=F_0+D'+\mathfrak{n}$ gives $\hat{F}_K=F_0+D'_K+\mathfrak{n}_K$. By
\eqref{8.15:eq-capped-data-b}, $|\mathfrak{n}|\le\varpi\tau$ on $Q_{3r}$, so
$|\mathfrak{n}_K(p)|\le\varpi\tau\le\tau/32$. Hence
\begin{equation*}
|\hat{F}_K(p)|\le\big(\tfrac34+\tfrac1{16}+\tfrac1{32}\big)\tau<\tfrac78\tau .
\end{equation*}
By the complementary slackness in \eqref{8.15:eq-capped-data-a}, $|\hat{F}(p)|=\tau$ wherever
$\Lambda(p)>0$. Hence $\langle\mu(p),\hat{F}(p)\rangle=\Lambda(p)\tau$ and
$|\mu(p)|=\Lambda(p)$ at every $p$. Since $D'-D'_K=(\hat{F}-\hat{F}_K)-(\mathfrak{n}-\mathfrak{n}_K)$,
and $2\varpi\le 1/16$,
\begin{align*}
\langle\mu,D'-D'_K\rangle
&=\langle\mu,\hat{F}-\hat{F}_K\rangle-\langle\mu,\mathfrak{n}-\mathfrak{n}_K\rangle
\ge\Lambda(\tau-|\hat{F}_K|)-2\varpi\tau\Lambda\\
&\ge\big(\tfrac18-\tfrac1{16}\big)\tau\Lambda=\tfrac1{16}\tau\Lambda .
\end{align*}

\emph{Step (ii): the field, a variance.} For vectors $a,b$ we have
\begin{equation*}
\langle a,a-b\rangle=\tfrac12|a-b|^2+\tfrac12|a|^2-\tfrac12|b|^2 .
\end{equation*}
Apply this with $a=D'(p)$ and
$b=D'(p+h)$. Sum against $K(h)\zeta(p)$. In the last sum substitute $q=p+h$ and use
$K(-h)=K(h)$; with $\zeta_K(q):=\sum_hK(h)\zeta(q+h)$ this gives
\begin{equation*}
\sum_p\zeta\langle D',D'-D'_K\rangle
=\frac12\sum_{p,h}K(h)\zeta(p)|D'(p)-D'(p+h)|^2+\frac12\sum_q(\zeta-\zeta_K)(q)|D'(q)|^2 .
\end{equation*}
The function $\zeta-\zeta_K$ vanishes off $Q_{2r-1+\kappa}\subset Q_{3r}$. By the Lipschitz bound
for $\zeta$,
\begin{equation*}
|\zeta-\zeta_K|\le\sum_hK(h)|h|_1/r\le 4\kappa/r=4\varepsilon .
\end{equation*}
Jensen's inequality gives $|D'(p)-D'_K(p)|^2\le\sum_hK(h)|D'(p)-D'(p+h)|^2$. Combining this with
step (i), the left side of \eqref{8.15:eq-multiplier-mass-c} is at least
\begin{equation}\label{8.15:eq-multiplier-mass-2}
\sum_p\zeta\Big[\frac{\tau\Lambda}{16}+\frac12|D'-D'_K|^2\Big]-2\varepsilon\|D'\|_{3r}^2 .
\end{equation}

\emph{Step (iii): the commutators.} From the definition of $\gamma$ and $|\partial\zeta|\le 1/r$,
\begin{equation*}
|\gamma_{\mu\nu}(z)|^2\le\frac2{r^2}\big(|\beta_{K,\nu}(z+e_\mu)|^2
+|\beta_{K,\mu}(z+e_\nu)|^2\big).
\end{equation*}
Each bond value of $\beta_K$ occurs in at most six of the $\gamma$'s, and only bonds with base
site in $Q_{2r}$ occur. Apply Cauchy--Schwarz and the $\ell^2$ bound on $\beta_K$ in
\eqref{8.15:eq-translation-potentials-a}, with $s=2r$ and $\|D'\|_{2r+\kappa}\le\|D'\|_{3r}$.
Since $12^{1/2}\cdot 8\le 28$, this gives, with $\|\gamma\|:=\|\gamma\|_{2r}$ (recall that
$\gamma$ is supported on plaquettes in $Q_{2r}$),
\begin{equation*}
\|\gamma\|\le\frac1r\,12^{1/2}\cdot 8\kappa\,\|D'\|_{3r}\le 28\varepsilon\|D'\|_{3r},
\qquad
|\langle D',\gamma\rangle|\le 28\varepsilon\|D'\|_{3r}^2 .
\end{equation*}

Next we bound $\gamma$ pointwise. For $p$ in $Q_{3r}$, by \eqref{8.15:eq-capped-data-b} and the
cap condition in \eqref{8.15:eq-capped-data-a},
\begin{equation*}
|D'(p)|\le|dA(p)-\hat{F}(p)|+|\hat{F}(p)|+|F_0|\le\varpi\tau+\tau+\tfrac34\tau\le 2\tau .
\end{equation*}
Here the cap is used at plaquettes of every orientation. The sup-norm bound on $\beta_K$ in
\eqref{8.15:eq-translation-potentials-a}, at bonds with base site in $Q_{2r+1}$, involves only
plaquettes within $|\cdot|_\infty$-distance $\kappa$ of $Q_{2r+1}$, hence plaquettes in
$Q_{2r+1+\kappa}\subset Q_{3r}$. It gives $|\beta_K|\le 4\kappa\cdot 2\tau$ at these bonds, and
so $|\gamma|\le(2/r)\cdot 4\kappa\cdot 2\tau=16\varepsilon\tau$. With $|\mu|=\Lambda$,
\begin{equation*}
|\langle\mu,\gamma\rangle|\le 16\varepsilon\,\tau\sum_{p\text{ in }Q_{2r}}\Lambda(p).
\end{equation*}
This multiplier commutator is not absorbed into the left side. It is bounded by the multiplier
mass on $Q_{2r}$, which is paid for by the excess on the larger cube $Q_{3r}$.

\emph{Step (iv): the defects.} Let $n_r$ be the number of plaquettes in $Q_{2r}$. Then
$n_r\le 6(4r+1)^4\le 6\cdot 5^4r^4$.

By \eqref{8.15:eq-multiplier-mass-1}, $d\xi$ is supported in $Q_{2r}$ with
$|d\xi|\le\zeta|D'-D'_K|+|\gamma|$, and $|\mathfrak{n}|\le\varpi\tau$ on $Q_{3r}$. Apply
Cauchy--Schwarz, using $\zeta^2\le\zeta$. Then use
\begin{equation*}
ab\le\tfrac14b^2+a^2,\qquad
28\varepsilon\varpi\tau n_r^{1/2}\|D'\|_{3r}\le\varepsilon\|D'\|_{3r}^2+196\,\varepsilon\varpi^2\tau^2n_r,
\end{equation*}
together with $\varepsilon\le 1$. This gives
\begin{align*}
|\langle\mathfrak{n},d\xi\rangle|
&\le\varpi\tau\sum_{p\text{ in }Q_{2r}}\big(\zeta|D'-D'_K|+|\gamma|\big)
\le\varpi\tau\,n_r^{1/2}\Big[\Big(\sum_p\zeta|D'-D'_K|^2\Big)^{1/2}+\|\gamma\|\Big]\\
&\le\frac14\sum_p\zeta|D'-D'_K|^2+\varepsilon\|D'\|_{3r}^2+C\varpi^2\tau^2r^4 .
\end{align*}

Now apply \eqref{8.15:eq-capped-data-b} to $\xi\in\mathfrak{X}_{3r}$, and use
$|\hat{F}|\le\tau$ and $|\xi|_{\partial p}\le|\beta_K|_{\partial p}$. The quantity
$|\xi|_{\partial p}$ vanishes unless a bond of $\partial p$ has base site in $Q_{2r-1}$, which
forces $p$ to be in $Q_{2r}$. The bonds of $\partial p$ for such $p$ have base sites in
$Q_{2r+1}$, and each bond lies in at most six plaquettes. Therefore, by Cauchy--Schwarz and the
$\ell^2$ bound in \eqref{8.15:eq-translation-potentials-a} with $s=2r+1$
($2r+1+\kappa\le 3r$, $6^{1/2}\le 4$),
\begin{equation*}
\sum_{p\text{ in }Q_{2r}}|\beta_K|_{\partial p}
\le n_r^{1/2}\,6^{1/2}\Big(\sum_{b\text{ in }Q_{2r+1}}|\beta_K(b)|^2\Big)^{1/2}
\le n_r^{1/2}\cdot 4\cdot 8\kappa\|D'\|_{3r}.
\end{equation*}
The sup-norm bound gives $|\beta_K|_{\partial p}\le 8\kappa\tau$ for $p$ in $Q_{2r}$. Hence,
using $\varpi\le 1$ in the last term,
\begin{align*}
|\mathfrak{d}|
&\le\frac{\varpi}{r}\sum_{p\text{ in }Q_{2r}}(\tau+\Lambda)|\beta_K|_{\partial p}
\le\frac{\varpi}{r}\Big[\tau n_r^{1/2}\cdot 4\cdot 8\kappa\|D'\|_{3r}
+8\kappa\tau\sum_{p\text{ in }Q_{2r}}\Lambda\Big]\\
&\le\varepsilon\|D'\|_{3r}^2+C\varepsilon\varpi^2\tau^2r^4
+8\varepsilon\,\tau\sum_{p\text{ in }Q_{2r}}\Lambda .
\end{align*}
Here we used
\begin{equation*}
32\varepsilon\varpi\tau n_r^{1/2}\|D'\|_{3r}
\le\varepsilon\|D'\|_{3r}^2+256\,\varepsilon\varpi^2\tau^2n_r .
\end{equation*}

\emph{Step (v): the field by the excess.} Since $D'=(\hat{F}-F_0)-\mathfrak{n}$, and
$\frac12\|\hat{F}-F_0\|_{3r}^2\le\mathcal{Y}(3r;F_0)$ by \eqref{8.15:eq-capped-data-c}, we have
\begin{equation}\label{8.15:eq-multiplier-mass-3}
\|D'\|_{3r}^2\le 2\|\hat{F}-F_0\|_{3r}^2+2\|\mathfrak{n}\|_{3r}^2
\le 4\mathcal{Y}(3r;F_0)+C\varpi^2\tau^2r^4 .
\end{equation}
In the last step we used $|\mathfrak{n}|\le\varpi\tau$ on the at most $6(6r+1)^4\le 6\cdot 7^4r^4$
plaquettes of $Q_{3r}$.

\emph{Collection.} The right side of \eqref{8.15:eq-multiplier-mass-c} is at most
$|\langle D',\gamma\rangle|+|\langle\mu,\gamma\rangle|+|\mathfrak{d}|
+|\langle\mathfrak{n},d\xi\rangle|$. Bound it by steps (iii) and (iv), and bound the left side
from below by \eqref{8.15:eq-multiplier-mass-2}. Subtract
$\frac14\sum_p\zeta|D'-D'_K|^2$ and use $\varepsilon\le 1$. This gives
\begin{equation*}
\sum_p\zeta\Big[\frac{\tau\Lambda}{16}+\frac14|D'-D'_K|^2\Big]
\le 32\varepsilon\|D'\|_{3r}^2+24\varepsilon\,\tau\sum_{p\text{ in }Q_{2r}}\Lambda
+C\varpi^2\tau^2r^4 .
\end{equation*}
The summand on the left is non-negative, $\zeta\ge 0$, and $\zeta=1$ on $Q_r$. So the left
side is at least the left side of \eqref{8.15:eq-multiplier-mass-b}. All summands of
$\mathcal{Y}(3r;F_0)$ are non-negative and $Q_{2r}\subset Q_{3r}$, so
$\tau\sum_{p\text{ in }Q_{2r}}\Lambda\le\mathcal{Y}(3r;F_0)$. Insert
\eqref{8.15:eq-multiplier-mass-3} and use $\varepsilon\le 1$. The right side is then at most
\begin{equation*}
152\,\varepsilon\,\mathcal{Y}(3r;F_0)+C\varpi^2\tau^2r^4 .
\end{equation*}
This is \eqref{8.15:eq-multiplier-mass-b} with an absolute constant $C_J$.
\end{proof}

We note how the three ingredients of the proof enter.

When the excess is small in the sense of \eqref{8.15:eq-multiplier-mass-a}, the average of
translates $\hat{F}_K$ does not merely respect the caps. By step (i) it keeps a margin of at
least $\tau/8$ below the cap at every plaquette of $Q_{2r}$, whatever the field does there. This
margin gives the multiplier term its good sign with the rate $\tau\Lambda/16$. Step (i) plays the
part of the absorption step of the model theorem on the exit value, here with a margin of at
least $\tau/8$. That is a quantitative gain which is not available globally in the exit-value
model.

The proof never compares $\hat{F}$ with a competitor. The $2$-form
$(1-\zeta)\hat{F}+\zeta\hat{F}_K+\gamma$ need not respect the caps where $\zeta$ is small. The
identity
\eqref{8.15:eq-multiplier-mass-c} holds for an arbitrary test field, and feasibility of the test
field is never used.

The shift $\kappa$ is the small fraction $\varepsilon$ of the side of the cube, which makes the
cutoff cheap. In step (iii) the commutator $\gamma$ carries a factor
(shift)$\times$(field)$/$(side), not (diameter)$\times$(cap)$/$(side).

\begin{lemma}[Relative 1-forms bounded by their exterior derivative]\label{8.15:lem-relative-one-forms}
Let $r\ge1$ be an integer and $c\in\mathbb{Z}^4$. On $\mathbb{Z}^4$, with unit vectors $e_1,\dots,e_4$ and forward
differences $\partial_\mu f(z)=f(z+e_\mu)-f(z)$, consider $\mathsf{V}$-valued forms, with
$(d\xi)_{\mu\nu}=\partial_\mu\xi_\nu-\partial_\nu\xi_\mu$. For a 1-form $\xi$ put $\|\xi\|^2=\sum_{z,\mu}|\xi_\mu(z)|^2$;
for 2-forms put $\langle X,Y\rangle=\sum_p\langle X(p),Y(p)\rangle$, the sum over plaquettes $p=(z,\mu\nu)$, $\mu<\nu$,
and $\|X\|^2=\langle X,X\rangle$. Let $Q_r=Q_r(c)$ be the cube about $c$, and let $\mathfrak{X}_r$ be the space of
1-forms supported on the bonds having at least one endpoint in $Q_{r-1}$. Then for every $\xi\in\mathfrak{X}_r$ there is
a $\mathsf{V}$-valued function $\lambda$ supported in $Q_{r-1}$ such that
\begin{equation*}
\xi':=\xi-d\lambda\in\mathfrak{X}_r\qquad\text{and}\qquad \|\xi'\|\le r\,\|d\xi\| .
\end{equation*}
\end{lemma}

The lemma is used below only when the defect function $\varpi$ of Definition~\ref{8.15:def-capped-data} is not
identically zero.

For $\xi\in\mathfrak{X}_r$ the 2-form $d\xi$ is supported on plaquettes in $Q_r$, so $\|d\xi\|=\|d\xi\|_r$. On finitely
supported 1-forms, $d^*$ is the adjoint of $d$:
\begin{equation}\label{8.15:eq-relative-one-forms-1}
\langle\xi,d\lambda\rangle=\sum_z\langle(d^*\xi)(z),\lambda(z)\rangle,\qquad
(d^*\xi)(z)=\sum_{\mu}\bigl(\xi_\mu(z-e_\mu)-\xi_\mu(z)\bigr).
\end{equation}
On finitely supported 2-forms, $d^*$ is likewise the adjoint of $d$ acting on 1-forms.

\begin{proof}
\emph{Choice of $\lambda$.} The forms $d\lambda$, with $\lambda$ supported in $Q_{r-1}$, make up a finite-dimensional
linear subspace of the 1-forms. We choose $\lambda$ so that $d\lambda$ is the orthogonal projection of $\xi$ onto this
subspace; thus $\lambda$ minimises $\|\xi-d\lambda\|$.

Since $(d\lambda)_\mu(z)=\lambda(z+e_\mu)-\lambda(z)$, the form $d\lambda$ is supported on bonds with an endpoint in
$Q_{r-1}$. Hence $d\lambda\in\mathfrak{X}_r$, and so $\xi'\in\mathfrak{X}_r$.

By minimality, $\langle\xi',d\eta\rangle=0$ for every $\eta$ supported in $Q_{r-1}$. Taking $\eta$ equal to $v$ at $z$
and zero elsewhere, with $z\in Q_{r-1}$ and $v\in\mathsf{V}$, in \eqref{8.15:eq-relative-one-forms-1} gives
$(d^*\xi')(z)=0$ for all
$z\in Q_{r-1}$. Moreover $(dd\lambda)_{\mu\nu}=\partial_\mu\partial_\nu\lambda-\partial_\nu\partial_\mu\lambda=0$, so
$d\xi'=d\xi$.

\emph{The Weitzenb\"ock identity.} Extend $\xi'$ by zero to all bonds of $\mathbb{Z}^4$. We use the lattice
Weitzenb\"ock identity $dd^*+d^*d=-\Delta$, which holds componentwise on 1-forms on $\mathbb{Z}^4$. Here
\begin{equation*}
(\Delta f)(z)=\sum_i\bigl(f(z+e_i)+f(z-e_i)-2f(z)\bigr).
\end{equation*}
To verify the identity, write $\partial^-_\mu f(z)=f(z)-f(z-e_\mu)$, so that $\partial^-_\mu\partial_\mu f(z)$ equals
$f(z+e_\mu)+f(z-e_\mu)-2f(z)$, and summation by parts reads
$\sum_z\langle a(z),\partial_\mu b(z)\rangle=-\sum_z\langle\partial^-_\mu a(z),b(z)\rangle$ for finitely supported
$a,b$. By \eqref{8.15:eq-relative-one-forms-1}, $d^*\xi=-\sum_\mu\partial^-_\mu\xi_\mu$ on 1-forms. On 2-forms, put
$X_{\nu\mu}=-X_{\mu\nu}$ and $X_{\mu\mu}=0$; then
$\langle X,d\xi\rangle=\sum_{\mu,\nu}\sum_z\langle X_{\mu\nu}(z),\partial_\mu\xi_\nu(z)\rangle$, and summation by parts
gives $(d^*X)_\nu=-\sum_\mu\partial^-_\mu X_{\mu\nu}$. Since $\partial^-_\mu\partial_\nu=\partial_\nu\partial^-_\mu$,
\begin{align*}
(d^*d\xi)_\nu+(dd^*\xi)_\nu
&=-\sum_\mu\partial^-_\mu\bigl(\partial_\mu\xi_\nu-\partial_\nu\xi_\mu\bigr)
-\partial_\nu\sum_\mu\partial^-_\mu\xi_\mu\\
&=-\sum_\mu\partial^-_\mu\partial_\mu\xi_\nu=-\Delta\xi_\nu .
\end{align*}
Summation by parts also gives $-\langle f,\Delta f\rangle=\sum_i\|\partial_i f\|^2$ for finitely supported $f$. Pairing
the identity with $\xi'$ gives
\begin{equation}\label{8.15:eq-relative-one-forms-2}
\|d\xi'\|^2+\|d^*\xi'\|^2=\sum_{\mu,i}\|\partial_i\xi'_\mu\|^2 .
\end{equation}

\emph{Boundary term.} If $z\notin Q_r$, every bond with endpoint $z$ has its other endpoint at $|\cdot|_\infty$-distance
at least $r$ from $c$. Such a bond does not touch $Q_{r-1}$, so $(d^*\xi')(z)=0$. Together with the first step, this
shows that $d^*\xi'$ is supported on $Q_r\setminus Q_{r-1}$.

Let $z\in Q_r\setminus Q_{r-1}$. A neighbour $z\pm e_j$ of $z$ lies in $Q_{r-1}$ only if $j$ is the unique coordinate
with $|z_j-c_j|=r$ and the step points inward. Hence $z$ is joined to $Q_{r-1}$ by at most one bond, which we call
$b_z$. The bonds at $z$ that carry non-zero values of $\xi'$ are among those meeting $Q_{r-1}$. If there is no such bond,
$(d^*\xi')(z)=0$. Otherwise, by \eqref{8.15:eq-relative-one-forms-1}, $|(d^*\xi')(z)|=|\xi'(b_z)|$.

Let $\mu$ be the direction of $b_z$. There are two cases.
\begin{itemize}
\item If $b_z=(z-e_\mu,\mu)$, then $z_\mu-c_\mu=r$. The next collinear bond $(z,\mu)$ has neither endpoint in
$Q_{r-1}$, so $\xi'_\mu(z)=0$ and $\partial_\mu\xi'_\mu(z-e_\mu)=-\xi'(b_z)$.
\item If $b_z=(z,\mu)$, then $z_\mu-c_\mu=-r$. The next collinear bond $(z-e_\mu,\mu)$ has neither endpoint in
$Q_{r-1}$, so $\xi'_\mu(z-e_\mu)=0$ and $\partial_\mu\xi'_\mu(z-e_\mu)=\xi'(b_z)$.
\end{itemize}
In both cases $|\partial_\mu\xi'_\mu(z-e_\mu)|=|(d^*\xi')(z)|$. The pair $(\mu,z-e_\mu)$ determines $z$, so distinct
sites give distinct terms, and therefore
\begin{equation*}
\sum_\mu\|\partial_\mu\xi'_\mu\|^2\ge\|d^*\xi'\|^2 .
\end{equation*}
Inserting this into \eqref{8.15:eq-relative-one-forms-2} yields
\begin{equation}\label{8.15:eq-relative-one-forms-3}
\|d\xi'\|^2\ge\sum_\mu\sum_{i\ne\mu}\|\partial_i\xi'_\mu\|^2 .
\end{equation}

\emph{One-dimensional Poincar\'e inequality.} Let $i\ne\mu$. If $\xi'_\mu(z)\ne0$, then one endpoint of $(z,\mu)$ lies in
$Q_{r-1}$. Both endpoints have $i$-th coordinate $z_i$, so $|z_i-c_i|\le r-1$.

Consequently, on each line parallel to $e_i$, the function $t\mapsto\xi'_\mu(z+te_i)$ is supported on $2r-1$ consecutive
integers. For a function $f$ on $\mathbb{Z}$ supported on $n$ consecutive points, $\sum_t|f(t+1)-f(t)|^2$ is the
quadratic form of the $n\times n$ tridiagonal matrix with $2$ on the diagonal and $-1$ beside it. This matrix has
eigenvalues $4\sin^2\bigl(j\pi/(2(n+1))\bigr)$, $1\le j\le n$; we apply this componentwise in an orthonormal basis of
$\mathsf{V}$. With $n=2r-1$, and using $\sin x\ge 2x/\pi$ on $[0,\pi/2]$, we obtain
\begin{equation*}
\sum_t|\partial f(t)|^2\ge4\sin^2\Bigl(\frac{\pi}{4r}\Bigr)\sum_t|f(t)|^2\ge r^{-2}\sum_t|f(t)|^2 .
\end{equation*}
Summing over all lines parallel to $e_i$ gives $\|\partial_i\xi'_\mu\|^2\ge r^{-2}\|\xi'_\mu\|^2$.

\emph{Conclusion.} For each $\mu$, keep in \eqref{8.15:eq-relative-one-forms-3} one direction $i\ne\mu$ and drop the
other, non-negative, terms. This gives $\|d\xi'\|^2\ge r^{-2}\|\xi'\|^2$. Since $d\xi'=d\xi$, we conclude
$\|\xi'\|\le r\|d\xi\|$.
\end{proof}

\begin{lemma}[The averaged field is close to a harmonic form]\label{8.15:lem-averaged-field}
Assume the hypotheses of Lemma~\ref{8.15:lem-multiplier-mass}. Let $D'=dA-F_0$ and let $D'_K$ be
its average over the translates by $\mathsf{K}$, as there. Then there is a $2$-form $H$ on the
plaquettes in $Q_r=Q_r(c)$ with the following properties.
\begin{itemize}
\item $H$ is closed on the $3$-cells in $Q_{r-1}$.
\item $\langle H,d\xi\rangle=0$ for all $\xi\in\mathfrak{X}_r$.
\item Each component $H_{\mu\nu}$ is $\Delta$-harmonic at the base sites in $Q_{r-2}$. This uses
the criterion that a closed $2$-form is $\Delta$-harmonic, orientation by orientation, at every
plaquette all four of whose edges $b$ satisfy $(d^{*}H)(b)=0$.
\end{itemize}
Moreover, with an absolute constant $C$,
\begin{equation}\label{8.15:eq-averaged-field-a}
\|D'_K-H\|_r\le 2(2\kappa+1)^{-2}\sum_{p\in Q_{2r}}\Lambda(p)+C\varpi\tau r^{2}.
\end{equation}
\end{lemma}

\begin{proof}
Pairings $\langle\cdot,\cdot\rangle$ of $2$-forms are sums over plaquettes. On $1$-forms,
$\|\cdot\|$ is the $\ell^{2}$ norm over bonds. For a function $X$ on plaquettes, $X_K$ is the
average $X_K(p)=\sum_{h}K(h)X(p+h)$ of Lemma~\ref{8.15:lem-translation-potentials}. In particular
$\Lambda_K$, $\mu_K$ and $\mathfrak{n}_K$ are the averages of $\Lambda$, of $\mu$ and of the defect
$\mathfrak{n}:=\hat{F}-dA$. We write $\mathfrak{d}(\xi):=\langle\hat{F}+\mu,d\xi\rangle$.

If $\xi\in\mathfrak{X}_r$, then $d\xi$ vanishes except on plaquettes having an edge with an
endpoint in $Q_{r-1}$, and every such plaquette lies in $Q_r$. Hence
\[
\mathcal{E}_r:=\{d\xi:\xi\in\mathfrak{X}_r\}
\]
is a finite-dimensional space of $2$-forms on the plaquettes in $Q_r$.

\emph{Construction of $H$.} Let $w$ be the orthogonal projection of the restriction of $D'_K$ to
$Q_r$ onto $\mathcal{E}_r$, for the inner product of $\|\cdot\|_r$, and put $H:=D'_K-w$. Since
$w$ is an orthogonal projection, $\langle H,d\xi\rangle=0$ for every $\xi\in\mathfrak{X}_r$.

The form $D'_K$ is closed, being an average of translates of $dA-F_0$. The form $w$ is exact.
Hence $H$ is closed on the $3$-cells in $Q_{r-1}$.

Let $b$ be a bond with an endpoint in $Q_{r-1}$. Applying the orthogonality relation to a form
$\xi\in\mathfrak{X}_r$ supported on $b$ alone, and using that all plaquettes containing $b$ lie in
$Q_r$, gives $(d^{*}H)(b)=0$. Every edge of a plaquette with base site in $Q_{r-2}$ has both
endpoints in $Q_{r-1}$. The harmonicity criterion recalled in the statement therefore applies at
these plaquettes. Finally $D'_K-H=w$, so it remains to bound $\|w\|_r$.

\emph{A potential for $w$.} Since $w\in\mathcal{E}_r$, we may write $w=d\xi$ with
$\xi\in\mathfrak{X}_r$. Lemma~\ref{8.15:lem-relative-one-forms} yields $\lambda$ supported in
$Q_{r-1}$ such that $\xi_w:=\xi-d\lambda\in\mathfrak{X}_r$ and
$\|\xi_w\|\le r\|d\xi\|=r\|w\|_r$. Since $d\,d\lambda=0$, we also have $d\xi_w=w$.

\emph{The identity for $\|w\|_r^{2}$.} Because $w$ is the projection of $D'_K$, we have
$\|w\|_r^{2}=\langle D'_K,d\xi_w\rangle$. Put $\xi^{h}:=\xi_w(\cdot-h)$ for $h\in\mathsf{K}$. The
substitution $q=p+h$ gives
\[
\langle D'_K,d\xi_w\rangle=\sum_{h}K(h)\langle D',d\xi^{h}\rangle .
\]
Since $|h|_\infty\le\kappa$, each $\xi^{h}$ is supported on bonds with an endpoint in
$Q_{r+\kappa-1}(c)$, so $\xi^{h}\in\mathfrak{X}_{r+\kappa}$ about $c$. As $\kappa\le r-2$, this
space is contained in $\mathfrak{X}_{3r}$. For each $h$,
\begin{equation}\label{8.15:eq-averaged-field-1}
\langle D',d\xi^{h}\rangle=\langle\hat{F}+\mu,d\xi^{h}\rangle-\langle\mu,d\xi^{h}\rangle
-\langle\mathfrak{n},d\xi^{h}\rangle .
\end{equation}
Summing \eqref{8.15:eq-averaged-field-1} with the weights $K(h)$ and undoing the same
substitution in the last two terms gives
\begin{equation}\label{8.15:eq-averaged-field-2}
\|w\|_r^{2}=-\langle\mu_K,w\rangle-\langle\mathfrak{n}_K,w\rangle
+\sum_{h}K(h)\,\mathfrak{d}(\xi^{h}).
\end{equation}

\emph{The multiplier term.} By \eqref{8.15:eq-capped-data-a}, $\mu=\Lambda\hat{F}/\tau$ and
$|\hat{F}|\le\tau$, so $|\mu|\le\Lambda$ and, after averaging, $|\mu_K|\le\Lambda_K$. Let
$p\in Q_r$. The plaquettes $p+h$, $h\in\mathsf{K}$, are distinct and lie in
$Q_{r+\kappa}\subset Q_{2r}$, and $K(h)=(2\kappa+1)^{-4}$. Hence
\[
\Lambda_K(p)\le(2\kappa+1)^{-4}\sum_{Q_{2r}}\Lambda .
\]
By the same inclusion,
\[
\sum_{p\in Q_r}\Lambda_K(p)=\sum_hK(h)\sum_{p\in Q_r}\Lambda(p+h)\le\sum_{Q_{2r}}\Lambda .
\]
Therefore
\begin{equation}\label{8.15:eq-averaged-field-3}
\|\Lambda_K\|_r^{2}\le\sup_{Q_r}\Lambda_K\cdot\sum_{Q_r}\Lambda_K
\le(2\kappa+1)^{-4}\Big(\sum_{Q_{2r}}\Lambda\Big)^{2}.
\end{equation}
The averaging thus converts the $\ell^{1}$ mass of the multipliers into a pointwise bound. Only
this mass enters, whatever the concentration of $\Lambda$.

\emph{The defect term.} By \eqref{8.15:eq-capped-data-b}, $|\mathfrak{n}(p)|\le\varpi\tau$ for
$p$ in $Q_{3r}$, so $|\mathfrak{n}_K|\le\varpi\tau$ on $Q_r$. Since $Q_r$ contains at most
$6(2r+1)^{4}$ plaquettes, $\|\mathfrak{n}_K\|_r\le C\varpi\tau r^{2}$ with an absolute $C$.

\emph{The equation term.} Each $\xi^{h}$ lies in $\mathfrak{X}_{3r}$, so
\eqref{8.15:eq-capped-data-b} applies to it. Substituting $q=p-h$ and using
$|\xi^{h}|_{\partial p}=|\xi_w|_{\partial(p-h)}$, we get
\[
\Big|\sum_hK(h)\,\mathfrak{d}(\xi^{h})\Big|
\le\frac{\varpi}{r}\sum_{q}(|\hat{F}|+\Lambda)_K(q)\,|\xi_w|_{\partial q}.
\]
Only plaquettes $q\in Q_r$ contribute, since $\xi_w\in\mathfrak{X}_r$. There
$(|\hat{F}|+\Lambda)_K\le\tau+\Lambda_K$, by \eqref{8.15:eq-capped-data-a}. Each bond lies in six
plaquettes, so $\sum_q|\xi_w|_{\partial q}^{2}\le 6\|\xi_w\|^{2}$. By the Cauchy--Schwarz
inequality, the count of plaquettes, and $\|\xi_w\|\le r\|w\|_r$,
\begin{equation}\label{8.15:eq-averaged-field-4}
\begin{split}
\Big|\sum_hK(h)\,\mathfrak{d}(\xi^{h})\Big|
&\le\frac{\varpi}{r}\big(C\tau r^{2}+\|\Lambda_K\|_r\big)\cdot 4\|\xi_w\|\\
&\le 4\varpi\big(C\tau r^{2}+\|\Lambda_K\|_r\big)\|w\|_r .
\end{split}
\end{equation}

\emph{Conclusion.} If $w=0$ there is nothing to prove. Otherwise insert $|\mu_K|\le\Lambda_K$, the
bound on $\|\mathfrak{n}_K\|_r$ and \eqref{8.15:eq-averaged-field-4} into
\eqref{8.15:eq-averaged-field-2}. Then divide by $\|w\|_r$ and use $4\varpi\le1$, which holds
since $\varpi\le1/32$. This gives
\[
\|w\|_r\le 2\|\Lambda_K\|_r+C\varpi\tau r^{2},
\]
with a new absolute $C$. Together with \eqref{8.15:eq-averaged-field-3} and $D'_K-H=w$, this is
\eqref{8.15:eq-averaged-field-a}.
\end{proof}

\begin{lemma}[Interior decay for lattice-harmonic functions]\label{8.15:lem-harmonic-decay}
Let $c\in\mathbb{Z}^4$, let $s\ge 64$ be an integer, and write $Q_r=Q_r(c)$. Let
$u:Q_s\to\mathsf{V}$ be $\Delta$-harmonic at the sites of $Q_{s-1}$. There is an absolute constant
$C_H$ such that
\begin{equation*}
|u(c)|^2\le C_H\,s^{-4}\sum_{Q_s}|u|^2 ,
\end{equation*}
and such that for $1\le t\le s/16$
\begin{equation}\label{8.15:eq-harmonic-decay-a}
\sum_{Q_t}|u-u(c)|^2\le C_H\Bigl(\frac{t}{s}\Bigr)^{6}\sum_{Q_s}|u|^2 .
\end{equation}
\end{lemma}

Here $\Delta$ is the lattice Laplacian, $\Delta u(z)=\sum_{z'\sim z}(u(z')-u(z))$, so that
$\Delta u(z)$ is defined for $z\in Q_{s-1}$. The inner product of $\mathsf{V}$ is
$\langle\cdot,\cdot\rangle$. We write $e_1,\dots,e_4$ for the unit vectors of $\mathbb{Z}^4$. For a
bond $b$ with endpoints $z,z'$ we put $\partial\varphi(b)=\varphi(z')-\varphi(z)$, and
$\partial_i u(z)=u(z+e_i)-u(z)$. In this proof the mean-value inequality for nonnegative
subharmonic functions is used in the following form: if $v\ge0$ and $\Delta v\ge0$ on
$Q_{2r}(x_0)$, with $r\ge1$ an integer, then
$v(x_0)\le C_{\mathrm{MV}}r^{-4}\sum_{Q_{2r+1}(x_0)}v$.

\begin{proof}
\emph{Subharmonicity and the bound at the centre.} The identity
$|a|^2-|b|^2=|a-b|^2+2\langle b,a-b\rangle$, applied with $a=u(z')$ and $b=u(z)$ and summed over
$z'\sim z$, gives
\begin{equation*}
\Delta|u|^2(z)=\sum_{z'\sim z}|u(z')-u(z)|^2+2\langle u(z),\Delta u(z)\rangle .
\end{equation*}
Hence $\Delta|u|^2\ge0$ on $Q_{s-1}$. Put $r=\lfloor (s-1)/2\rfloor$. Then $2r\le s-1$ and
$2r+1\le s$, and $r\ge s/2-1\ge s/4$. The mean-value inequality applied to $v=|u|^2$ at $x_0=c$
therefore gives
\begin{equation*}
|u(c)|^2\le C_{\mathrm{MV}}r^{-4}\sum_{Q_s}|u|^2\le 256\,C_{\mathrm{MV}}\,s^{-4}\sum_{Q_s}|u|^2 .
\end{equation*}

\emph{Caccioppoli inequality.} We take a cutoff $\eta:\mathbb{Z}^4\to[0,1]$ which equals $1$ on
$Q_{\lfloor s/2\rfloor}$, vanishes off $Q_{s-2}$, and satisfies $|\partial\eta|\le 4/s$ on every
bond. An example is
$\eta(z)=\min\{1,(s-2-|z-c|_\infty)_+/(s-2-\lfloor s/2\rfloor)\}$, since
$s-2-\lfloor s/2\rfloor\ge s/2-2\ge s/4$.

The function $\eta^2u$ vanishes off $Q_{s-2}$, and $\Delta u=0$ on $Q_{s-2}\subset Q_{s-1}$.
Each bond contributes $\langle\varphi(z),u(z')-u(z)\rangle+\langle\varphi(z'),u(z)-u(z')\rangle$
to $\sum_z\langle\varphi,\Delta u\rangle$. Applying this with $\varphi=\eta^2u$ gives
\begin{equation*}
0=\sum_z\eta^2\langle u,\Delta u\rangle=-\sum_{b}\langle\partial(\eta^2u),\partial u\rangle ,
\end{equation*}
the sum running over bonds with both endpoints in $Q_s$. On a bond, primes denote values at the
far endpoint. Then
\begin{equation*}
\partial(\eta^2u)=\tfrac12(\eta'^2+\eta^2)\,\partial u+\tfrac12(u'+u)(\eta'+\eta)\,\partial\eta ,
\end{equation*}
as expanding the right side shows. Substituting and using the Cauchy--Schwarz inequality gives
\begin{equation*}
\tfrac12\sum_b(\eta'^2+\eta^2)|\partial u|^2
\le\sum_b\tfrac12|u'+u|\,|\partial\eta|\,(\eta'+\eta)\,|\partial u| .
\end{equation*}
We now bound each term on the right. We use $xy\le\frac14x^2+y^2$ with
$x=(\eta'+\eta)|\partial u|$ and $y=|u'+u||\partial\eta|$, and then
$(\eta'+\eta)^2\le2(\eta'^2+\eta^2)$. This gives
\begin{equation*}
\sum_b\tfrac12|u'+u|\,|\partial\eta|\,(\eta'+\eta)\,|\partial u|
\le\tfrac14\sum_b(\eta'^2+\eta^2)|\partial u|^2+\tfrac12\sum_b|u'+u|^2|\partial\eta|^2 .
\end{equation*}
Hence
$\frac14\sum_b(\eta'^2+\eta^2)|\partial u|^2\le\frac12\sum_b|u'+u|^2|\partial\eta|^2$.

Consider a bond with an endpoint in $Q_{\lfloor s/2\rfloor-1}$. Both its endpoints lie in
$Q_{\lfloor s/2\rfloor}$, so $\eta'=\eta=1$ there. On the right we use $|\partial\eta|\le4/s$ and
$|u'+u|^2\le2(|u'|^2+|u|^2)$. Every site is an endpoint of at most $8$ bonds. Together these give
\begin{equation}\label{8.15:eq-harmonic-decay-1}
\sum_{b\cap Q_{\lfloor s/2\rfloor-1}\neq\emptyset}|\partial u(b)|^2
\le 256\,s^{-2}\sum_{Q_s}|u|^2 .
\end{equation}

\emph{Pointwise bound on differences.} Fix $i$. For $z\in Q_{s-2}$ and $z'\sim z$ we have
$z',z'+e_i\in Q_s$. Hence $\partial_iu$ is defined on the neighbours of $z$, and
$\Delta\partial_iu(z)=\Delta u(z+e_i)-\Delta u(z)=0$, because $z,z+e_i\in Q_{s-1}$. Thus
$\partial_iu$ is harmonic on $Q_{s-2}$. By the identity of the first step, $|\partial_iu|^2$ is
nonnegative and subharmonic there.

Let $z\in Q_{\lfloor s/16\rfloor}$ and $r=\lfloor s/16\rfloor$, so that $r\ge s/32$. Then
$|w-c|_\infty\le 3s/16+1\le\lfloor s/2\rfloor-2$ for $w\in Q_{2r+1}(z)$. Consequently
$Q_{2r}(z)\subset Q_{s-2}$, and every bond $(w,w+e_i)$ with $w\in Q_{2r+1}(z)$ meets
$Q_{\lfloor s/2\rfloor-1}$. The mean-value inequality applied to $|\partial_iu|^2$ at $z$ on the
cube $Q_{2r+1}(z)$, of half-side about $s/8$, together with \eqref{8.15:eq-harmonic-decay-1},
gives
\begin{equation}\label{8.15:eq-harmonic-decay-2}
|\partial_iu(z)|^2\le C_{\mathrm{MV}}\,32^4s^{-4}\cdot256\,s^{-2}\sum_{Q_s}|u|^2 ,
\qquad z\in Q_{\lfloor s/16\rfloor} .
\end{equation}

\emph{Summation over $Q_t$.} Let $1\le t\le s/16$ and $z\in Q_t$. Join $c$ to $z$ by a
coordinatewise monotone path of $|z-c|_1\le4t$ bonds. The path stays in $Q_t$, so each of its
bonds has the form $(w,w+e_i)$ with $w\in Q_t\subset Q_{\lfloor s/16\rfloor}$. Write
$C'=(2^{28}C_{\mathrm{MV}})^{1/2}$. Then, by \eqref{8.15:eq-harmonic-decay-2},
\begin{equation*}
|u(z)-u(c)|\le|z-c|_1\max|\partial u|\le 4t\cdot C's^{-3}\Bigl(\sum_{Q_s}|u|^2\Bigr)^{1/2}.
\end{equation*}
Since $Q_t$ has $(2\lfloor t\rfloor+1)^4\le81t^4$ sites, summing over $Q_t$ gives
\begin{equation*}
\sum_{Q_t}|u-u(c)|^2\le81t^4\cdot16t^2C'^2s^{-6}\sum_{Q_s}|u|^2
=1296\cdot2^{28}C_{\mathrm{MV}}\Bigl(\frac ts\Bigr)^6\sum_{Q_s}|u|^2 .
\end{equation*}
Both assertions therefore hold with $C_H=1296\cdot2^{28}C_{\mathrm{MV}}$. This constant exceeds
$256\,C_{\mathrm{MV}}$, so it also covers the bound at the centre, and
\eqref{8.15:eq-harmonic-decay-a} follows.
\end{proof}

\begin{lemma}[One-step decay of the hybrid excess]\label{8.15:lem-excess-decay}
Assume the hypotheses of Lemma~\ref{8.15:lem-multiplier-mass}, with its notation: the cubes
$Q_r=Q_r(c)$ about the site $c$, the radius $\kappa$, $\varepsilon:=\kappa/r$,
$\varpi:=\varpi(r)$, the constant $2$-form $F_0$ and $D':=dA-F_0$. Assume in addition that
$r\ge 66$ and $\theta\in[1/r,1/32]$, and write $y:=y(r;F_0)$ and
$Q_{\theta r}:=Q_{\lfloor\theta r\rfloor}(c)$. Let $H$ be the $2$-form constructed in the proof of
Lemma~\ref{8.15:lem-averaged-field}, and let $F_0':=F_0+H(c)$, where $H(c)$ is the constant
$2$-form formed by the six values of $H$ at the plaquettes based at $c$. Then, with an absolute
constant $C_S$,
\begin{equation}\label{8.15:eq-excess-decay-a}
\begin{gathered}
\mathcal{Y}(\theta r;F_0')\le C_S\bigl[\theta^6+\varepsilon+\varepsilon^{-4}y\bigr]\,
\mathcal{Y}(3r;F_0)+C_S\,\varpi^2\tau^2r^4,\\
|F_0'-F_0|\le C_S\bigl(y^{1/2}+\varpi\bigr)\tau .
\end{gathered}
\end{equation}
\end{lemma}

\begin{proof}
Throughout, $C$ denotes an absolute constant, not the same at each occurrence, and
$\mathcal{Y}:=\mathcal{Y}(3r;F_0)$, so that $y=\mathcal{Y}/(\tau^2r^4)$ by
\eqref{8.15:eq-capped-data-c}. For a $2$-form $X$ we write $\|X\|_{\theta r}$ for
$\|X\|_{\lfloor\theta r\rfloor}$. Since $\theta\ge 1/r$, we have $\lfloor\theta r\rfloor\ge1$; since
$\theta\le 1/32$, we have $Q_{\theta r}\subset Q_r$.

\emph{Decomposition.} Let $\mathfrak{n}:=\hat{F}-dA$ be the defect, and let $D'_K$ be the average,
with respect to the uniform probability $K$ on $\mathsf{K}$, of the translates of $D'$ by the
vectors $h\in\mathsf{K}$. Since $D'+\mathfrak{n}=\hat{F}-F_0$, on the plaquettes in
$Q_{\theta r}$ we have
\begin{equation*}
\hat{F}-F_0'=(D'-D'_K)+(D'_K-H)+(H-H(c))+\mathfrak{n}.
\end{equation*}
By the Cauchy--Schwarz inequality, $|a_1+a_2+a_3+a_4|^2\le 4\sum_i|a_i|^2$; hence
$\tfrac12|\hat{F}-F_0'|^2$ is at most twice the sum of the four squares. Summing over the plaquettes
in $Q_{\theta r}$ and adding the multiplier term,
\begin{equation}\label{8.15:eq-excess-decay-1}
\begin{aligned}
\mathcal{Y}(\theta r;F_0')\le{}&2\bigl[\|D'-D'_K\|_r^2+\|D'_K-H\|_r^2
+\|H-H(c)\|_{\theta r}^2+\|\mathfrak{n}\|_{\theta r}^2\bigr]\\
&+\tau\sum_{p\in Q_{\theta r}}\Lambda(p),
\end{aligned}
\end{equation}
where the first two norms were enlarged from $Q_{\theta r}$ to $Q_r$. We bound the five terms.

\emph{First term.} By \eqref{8.15:eq-multiplier-mass-b},
$\tfrac14\|D'-D'_K\|_r^2\le C_J[\varepsilon\mathcal{Y}+\varpi^2\tau^2r^4]$, so
$\|D'-D'_K\|_r^2\le 4C_J[\varepsilon\mathcal{Y}+\varpi^2\tau^2r^4]$.

\emph{Second term.} Squaring \eqref{8.15:eq-averaged-field-a} with
$(a+b)^2\le2a^2+2b^2$ gives
\begin{equation*}
\|D'_K-H\|_r^2\le 8(2\kappa+1)^{-4}\Bigl(\sum_{Q_{2r}}\Lambda\Bigr)^2+C\varpi^2\tau^2r^4 .
\end{equation*}
Both summands of $\mathcal{Y}$ are non-negative and $Q_{2r}\subset Q_{3r}$, so
$\tau\sum_{Q_{2r}}\Lambda\le\mathcal{Y}$; and $2\kappa+1\ge\kappa=\varepsilon r$. Hence
\begin{equation}\label{8.15:eq-excess-decay-2}
\|D'_K-H\|_r^2\le 8(\varepsilon r)^{-4}\tau^{-2}\mathcal{Y}^2+C\varpi^2\tau^2r^4
=8\varepsilon^{-4}y\,\mathcal{Y}+C\varpi^2\tau^2r^4 .
\end{equation}

\emph{A bound for $\|D'\|_{2r}$.} Each $h\in\mathsf{K}$ has $|h|_\infty\le\kappa\le r-2$, so the
translate by $h$ of a plaquette in $Q_r$ lies in $Q_{2r}$; by the triangle inequality for the
average, $\|D'_K\|_r\le\|D'\|_{2r}$. Since $D'=(\hat{F}-F_0)-\mathfrak{n}$ and, by the pointwise
defect bound of Definition~\ref{8.15:def-capped-data}, $|\mathfrak{n}(p)|\le\varpi\tau$ for $p$ in
$Q_{3r}$, while $Q_{3r}$ contains at most $Cr^4$ plaquettes,
\begin{equation}\label{8.15:eq-excess-decay-3}
\|D'_K\|_r^2\le\|D'\|_{2r}^2\le 2\|\hat{F}-F_0\|_{2r}^2+2\|\mathfrak{n}\|_{2r}^2
\le 4\mathcal{Y}+C\varpi^2\tau^2r^4 .
\end{equation}
In the same way $\|\mathfrak{n}\|_{\theta r}^2\le C\varpi^2\tau^2r^4$, which bounds the fourth term
of \eqref{8.15:eq-excess-decay-1}.

\emph{A bound for $\|H\|_r$.} In the construction of Lemma~\ref{8.15:lem-averaged-field},
$H=D'_K-w$, where $w$ is the orthogonal projection of the restriction of $D'_K$ to $Q_r$ onto the
space of exact forms $\{d\xi:\xi\in\mathfrak{X}_r\}$; thus $\|w\|=\|D'_K-H\|_r$, and $H$, being the
residual of this orthogonal projection, is orthogonal to $w$ in the $\ell^2$ norm over the
plaquettes in $Q_r$. By the triangle inequality $\|H\|_r\le\|D'_K\|_r+\|w\|$, so by
\eqref{8.15:eq-excess-decay-2} and \eqref{8.15:eq-excess-decay-3},
\begin{equation}\label{8.15:eq-excess-decay-4}
\|H\|_r^2\le C\bigl(1+\varepsilon^{-4}y\bigr)\mathcal{Y}+C\varpi^2\tau^2r^4 .
\end{equation}
The orthogonality of $H$ and $w$ gives in addition $\|D'_K\|_r^2=\|H\|_r^2+\|w\|^2$, hence, by
\eqref{8.15:eq-excess-decay-3},
\begin{equation}\label{8.15:eq-excess-decay-5}
\|H\|_r^2\le\|D'_K\|_r^2\le 4\mathcal{Y}+C\varpi^2\tau^2r^4 .
\end{equation}
We use \eqref{8.15:eq-excess-decay-4} for the third term of \eqref{8.15:eq-excess-decay-1} and
\eqref{8.15:eq-excess-decay-5} for the value $H(c)$.

\emph{Third term and $H(c)$.} By Lemma~\ref{8.15:lem-averaged-field} each component $H_{\mu\nu}$,
read as a $\mathsf{V}$-valued function of the base site, is $\Delta$-harmonic at the base sites in
$Q_{r-2}$. Put $s:=r-2$ and $t:=\lfloor\theta r\rfloor$. Then $s\ge64$ since $r\ge66$; $H_{\mu\nu}$
is defined on $Q_s$, because the plaquettes based at sites of $Q_{r-2}$ lie in $Q_r$; and it is
$\Delta$-harmonic at the sites of $Q_{s-1}\subset Q_{r-2}$. Moreover
\begin{equation*}
1\le t\le\theta r\le \frac{r}{32}\le\frac{s}{16},\qquad
\frac{t}{s}\le\frac{\theta r}{r-2}\le\frac{33}{32}\,\theta,\qquad
s^{-4}\le\Bigl(\frac{33}{32}\Bigr)^4r^{-4}:
\end{equation*}
the first chain uses $\theta\in[1/r,1/32]$ and, in its last step, $r\ge4$; the second and third use
$r/(r-2)\le\tfrac{33}{32}$, which holds because $r\ge66$. Lemma~\ref{8.15:lem-harmonic-decay}
applies component by component, and summation of \eqref{8.15:eq-harmonic-decay-a} and of the
bound for the value at $c$ in Lemma~\ref{8.15:lem-harmonic-decay} over the six components gives
\begin{equation*}
\|H-H(c)\|_{\theta r}^2\le C\theta^6\|H\|_r^2,\qquad |H(c)|^2\le Cr^{-4}\|H\|_r^2 .
\end{equation*}
With \eqref{8.15:eq-excess-decay-4} and $\theta\le1$, the third term of
\eqref{8.15:eq-excess-decay-1} is at most
$C\theta^6(1+\varepsilon^{-4}y)\mathcal{Y}+C\varpi^2\tau^2r^4$.

\emph{Multiplier term.} Since $Q_{\theta r}\subset Q_r$ and $\Lambda\ge0$,
\eqref{8.15:eq-multiplier-mass-b} gives
$\tau\sum_{Q_{\theta r}}\Lambda\le 16C_J[\varepsilon\mathcal{Y}+\varpi^2\tau^2r^4]$.

\emph{Conclusion.} Inserting the five bounds into \eqref{8.15:eq-excess-decay-1}, the coefficient
of $\mathcal{Y}$ is at most
\begin{equation*}
24C_J\,\varepsilon+16\varepsilon^{-4}y+C\theta^6+C\theta^6\varepsilon^{-4}y ,
\end{equation*}
and $\theta^6\varepsilon^{-4}y\le\varepsilon^{-4}y$ in the bracket; all remaining terms are at most
$C\varpi^2\tau^2r^4$. This is the first inequality of \eqref{8.15:eq-excess-decay-a}. For the
second, $F_0'-F_0=H(c)$, and by \eqref{8.15:eq-excess-decay-5}
\begin{equation*}
|F_0'-F_0|^2\le Cr^{-4}\bigl(4\mathcal{Y}+C\varpi^2\tau^2r^4\bigr)\le C\,y\,\tau^2+C\varpi^2\tau^2 ,
\end{equation*}
whence $|F_0'-F_0|\le C(y^{1/2}+\varpi)\tau$. Taking $C_S$ the largest of the constants obtained
proves \eqref{8.15:eq-excess-decay-a}.
\end{proof}

\begin{theorem}[Model: $\varepsilon$-regularity of capped lattice fields]
\label{8.15:thm-epsilon-regularity}
The model is the lattice setting of Definition~\ref{8.15:def-capped-data}. It consists of the
lattice $\mathbb{Z}^4$, a colour space $\mathsf{V}$ with a Euclidean structure, a centre
$c\in\mathbb{Z}^4$, an integer $R\ge1$ and a cap height $\tau>0$. The data on $Q_{3R}(c)$ are a
$1$-form $A$, a $2$-form $\hat{F}$, a function $\Lambda\ge0$ on plaquettes and a non-decreasing
defect function $\varpi:[1,R]\to[0,1]$. They are subject to the caps and complementary slackness
\eqref{8.15:eq-capped-data-a} and to the defect bounds \eqref{8.15:eq-capped-data-b} for
$1\le r\le R$. For a constant $2$-form $F_0$, $\mathcal{Y}(r;F_0)$ and $y(r;F_0)$ are the hybrid
excess and the normalised excess of \eqref{8.15:eq-capped-data-c}.

There are constants $\theta_E\in(0,1/32]$, $\varepsilon_E\in(0,1/4]$, $\rho_E\in\mathbb{N}$,
$y_E,c_E\in(0,1)$ and $C_E<\infty$ with the following property. They depend only on the
dimension $4$, through the constant $C_{\mathrm{MV}}$ of the mean-value inequality for
lattice-subharmonic functions and through counting. They do not depend on the dimension of
$\mathsf{V}$, on $\tau$, $R$, $c$, or on the data. The property is this. Suppose $R\ge\rho_E$,
$|F_0|\le\tau/2$, and that, for some number $\varpi_R\in[0,c_E]$,
\begin{equation}\label{8.15:eq-epsilon-regularity-a}
\begin{gathered}
\text{(h1)}\quad y(R;F_0)\le y_E,\\
\text{(h2)}\quad \varpi(r)\le\varpi_R\cdot\frac{r}{R}\quad\text{for }\rho_E\le r\le R .
\end{gathered}
\end{equation}
Then every plaquette $p$ based at $c$ obeys
\begin{equation}\label{8.15:eq-epsilon-regularity-1}
|\hat{F}(p)-F_0|\le C_E\bigl(y(R;F_0)^{1/2}+\varpi_R\bigr)\tau\le\frac{\tau}{4};
\end{equation}
in particular $|\hat{F}(p)|<\tau$ and $\Lambda(p)=0$.
\end{theorem}

\begin{proof}
The proof is an iteration of the one-step decay of the hybrid excess
(Lemma~\ref{8.15:lem-excess-decay}). Let $C_S$ be the constant of that lemma. Throughout,
$Q_{\theta r}(c)$ means $Q_{\lfloor\theta r\rfloor}(c)$.

\emph{Choice of constants, in this order.} Let $\theta_E$ be the largest $\theta\in(0,1/32]$ with
$6^4C_S\theta^2\le1/8$. Let $\varepsilon_E$ be the largest $\varepsilon\in(0,1/4]$ with
$6^4C_S\theta_E^{-4}\varepsilon\le1/8$. Both maxima exist: each set is an interval of the form
$(0,a]$ with $a>0$. Let $\rho_E$ be the least integer $\ge\max(66,\,6/\theta_E,\,4/\varepsilon_E)$.
Put
\begin{equation*}
\rho_E^{+}:=\frac{3(\rho_E+1)}{\theta_E},\qquad C_{\sharp}:=6^4C_S\theta_E^{-4}.
\end{equation*}
With these choices,
\begin{equation}\label{8.15:eq-epsilon-regularity-2}
C_{\sharp}\theta_E^6=6^4C_S\theta_E^2\le\tfrac18,\qquad C_{\sharp}\varepsilon_E\le\tfrac18 .
\end{equation}

\emph{Radii.} The letters $r_k$, $\kappa_k$, $y_k$, $\varpi_k$ introduced below are local to this
proof. Put $r_0:=R$ and $r_{k+1}:=\lfloor\theta_Er_k/3\rfloor$. Since $3r_{k+1}$ is an integer
not exceeding $\theta_Er_k$, we have $3r_{k+1}\le\lfloor\theta_Er_k\rfloor$. Since $\theta_E\le1/32$,
we have $r_{k+1}\le\theta_Er_k/3\le r_k/96$. Suppose $r_k\ge\rho_E$. Then $\theta_Er_k\ge6$, so
$r_{k+1}\ge\theta_Er_k/3-1\ge\theta_Er_k/6$.

Put $\kappa_k:=\lfloor\varepsilon_Er_k\rfloor$ and $\varepsilon_k:=\kappa_k/r_k$. Suppose again
$r_k\ge\rho_E$. Then $\varepsilon_Er_k\ge4$, hence, using also $\varepsilon_E\le1/4$,
\begin{gather*}
\kappa_k\ge\varepsilon_Er_k-1\ge\varepsilon_Er_k/2\ge2,\\
\kappa_k\le\varepsilon_Er_k\le r_k/4\le r_k-2 .
\end{gather*}
Thus $\varepsilon_k\in[\varepsilon_E/2,\varepsilon_E]$,
$\kappa_k\ge2$ and $\kappa_k\le r_k-2$.

Let $k^{\ast}$ be the last $k$ with $r_k\ge\rho_E$. It exists, because $r_0=R\ge\rho_E$ and
$r_k\le96^{-k}R$. Since $r_{k^{\ast}+1}\le\rho_E-1$,
\begin{equation*}
\frac{\theta_Er_{k^{\ast}}}{3}<r_{k^{\ast}+1}+1\le\rho_E,
\end{equation*}
so $r_{k^{\ast}}<3\rho_E/\theta_E<\rho_E^{+}$.

\emph{Reference values.} Put $F_0^{(0)}:=F_0$. Let $F_0^{(k+1)}$ be the constant $2$-form $F_0'$
produced by Lemma~\ref{8.15:lem-excess-decay} from the reference $F_0^{(k)}$ at radius $r_k$,
mollification radius $\kappa_k$ and $\theta=\theta_E$. This is $F_0^{(k)}+H(c)$, with $H$ the
$2$-form of Lemma~\ref{8.15:lem-averaged-field}. It is defined whenever
Lemma~\ref{8.15:lem-excess-decay} applies at step $k$.

Put $y_k:=y(r_k;F_0^{(k)})$ and $\varpi_k:=\varpi(r_k)$; in particular $y_0=y(R;F_0)$. For
$k\le k^{\ast}$ we have
\begin{equation*}
\rho_E\le r_k\le96^{-k}R\le R .
\end{equation*}
So (h2) of
\eqref{8.15:eq-epsilon-regularity-a} gives $\varpi_k\le\varpi_Rr_k/R\le\varpi_R\,96^{-k}$.

\emph{The recursion.} Let $k\le k^{\ast}$, and suppose Lemma~\ref{8.15:lem-excess-decay} applies at
step $k$. Then
\begin{equation}\label{8.15:eq-epsilon-regularity-3}
\begin{aligned}
y_{k+1}&=\frac{\mathcal{Y}(3r_{k+1};F_0^{(k+1)})}{\tau^2r_{k+1}^4}
\le\Bigl(\frac{6}{\theta_E}\Bigr)^4\frac{\mathcal{Y}(\theta_Er_k;F_0^{(k+1)})}{\tau^2r_k^4}\\
&\le C_{\sharp}\bigl[\theta_E^6+\varepsilon_E+16\varepsilon_E^{-4}y_k\bigr]y_k+C_{\sharp}\varpi_k^2 .
\end{aligned}
\end{equation}
The equality is the definition of $y_{k+1}$.

The first inequality has two sources. The summands of $\mathcal{Y}$ are non-negative, and
$Q_{3r_{k+1}}(c)\subset Q_{\lfloor\theta_Er_k\rfloor}(c)$; so $\mathcal{Y}(\cdot\,;F_0^{(k+1)})$
increases from radius $3r_{k+1}$ to radius $\theta_Er_k$. Further, $r_{k+1}\ge\theta_Er_k/6$ gives
$r_{k+1}^{-4}\le(6/\theta_E)^4r_k^{-4}$.

The second inequality is \eqref{8.15:eq-excess-decay-a} with $r=r_k$, $\varepsilon=\varepsilon_k$,
$\theta=\theta_E$ and $y=y_k$. There $\mathcal{Y}(3r_k;F_0^{(k)})=y_k\tau^2r_k^4$ and
$\varpi=\varpi_k$. Moreover $\varepsilon_k\le\varepsilon_E$ and
$\varepsilon_k^{-4}\le(\varepsilon_E/2)^{-4}=16\varepsilon_E^{-4}$. Finally
$(6/\theta_E)^4C_S=C_{\sharp}$.

By \eqref{8.15:eq-epsilon-regularity-2}, the first two members of the bracket, multiplied by
$C_{\sharp}y_k$, are each at most $y_k/8$. The third is at most $y_k/8$ if
$16C_{\sharp}\varepsilon_E^{-4}y_k\le1/8$. Under that condition we obtain the recursion in the
first line below. The second line is the claim $(I_k)$:
\begin{equation}\label{8.15:eq-epsilon-regularity-b}
\begin{gathered}
y_{k+1}\le\tfrac38\,y_k+C_{\sharp}\varpi_k^2,\\
(I_k)\qquad y_k\le2^{-k}Y',\qquad Y':=y_0+8C_{\sharp}\varpi_R^2 .
\end{gathered}
\end{equation}

\emph{The claim $(I_k)$.} $(I_0)$ is $y_0\le Y'$, which is trivial.

Assume $(I_k)$, and assume the recursion holds at step $k$. Since $\varpi_k^2\le\varpi_R^2\,96^{-2k}$,
\begin{equation*}
y_{k+1}\le\tfrac38\cdot2^{-k}Y'+C_{\sharp}\varpi_R^2\,96^{-2k}\le2^{-k-1}Y'.
\end{equation*}
The last inequality holds because $Y'\ge8C_{\sharp}\varpi_R^2$, so that
\begin{equation*}
\tfrac18\cdot2^{-k}Y'\ge2^{-k}C_{\sharp}\varpi_R^2\ge C_{\sharp}\varpi_R^2\,96^{-2k},
\end{equation*}
and because $\tfrac38+\tfrac18=\tfrac12$. This is $(I_{k+1})$.

\emph{Consequence.} If $(I_k)$ holds for $0\le k\le k^{\ast}$, put
$S_\infty:=\sum_{k=0}^{k^{\ast}}(y_k^{1/2}+\varpi_k)$. Using
$\sum_k2^{-k/2}=(1-2^{-1/2})^{-1}\le4$ and $\sum_k96^{-k}=96/95\le2$,
\begin{equation}\label{8.15:eq-epsilon-regularity-4}
S_\infty\le4Y'^{1/2}+2\varpi_R .
\end{equation}

\emph{Smallness.} Choose $y_E,c_E\in(0,1)$ so small that, with $Y:=y_E+8C_{\sharp}c_E^2$, the
conditions below hold. Here $C$ is the absolute constant in the bound
$\|D'\|_{3r}^2\le(4y+C\varpi^2)\tau^2r^4$ established in the proof of
Lemma~\ref{8.15:lem-multiplier-mass}.
\begin{itemize}
\item[(s1)] $16C_{\sharp}\varepsilon_E^{-4}Y\le1/8$;
\item[(s2)] $\varepsilon_E^{-4}(4Y+Cc_E^2)\le1/256$;
\item[(s3)] $c_E\le1/32$;
\item[(s4)] $C_S(4Y^{1/2}+2c_E)\le1/8$;
\item[(s5)] $2Y(\rho_E^{+})^4\le1/64$.
\end{itemize}
Such a choice is possible because $\theta_E,\varepsilon_E,\rho_E,C_{\sharp},C_S$ and $C$ are
already fixed, and every left-hand side tends to $0$ as $y_E,c_E\to0$. By (h1) and (h2) of
\eqref{8.15:eq-epsilon-regularity-a},
\begin{equation*}
Y'=y_0+8C_{\sharp}\varpi_R^2\le y_E+8C_{\sharp}c_E^2=Y .
\end{equation*}

\emph{Propagation.} We show by induction on $k\in\{0,\dots,k^{\ast}\}$ two things:
Lemmas~\ref{8.15:lem-multiplier-mass}, \ref{8.15:lem-averaged-field}
and~\ref{8.15:lem-excess-decay} apply at every step $j<k$, and $(I_j)$ holds for $j\le k$. At
$k=0$ this is $(I_0)$. Let
$k\le k^{\ast}$ satisfy the induction hypothesis. Then $y_j\le Y'\le Y$ for $j\le k$. We check the
hypotheses of Lemma~\ref{8.15:lem-multiplier-mass} at radius $r=r_k$ with $\kappa=\kappa_k$,
reference $F_0^{(k)}$ and $D'=dA-F_0^{(k)}$.
\begin{itemize}
\item The bound on the radius holds:
\begin{equation*}
4\le66\le\rho_E\le r_k\le R .
\end{equation*}
\item The bound $1\le\kappa_k\le r_k-2$ holds by the choice of $\kappa_k$.
\item The bound on the defect holds by (s3):
\begin{equation*}
\varpi(r_k)=\varpi_k\le\varpi_R\le c_E\le1/32 .
\end{equation*}
\item For the reference bound, Lemma~\ref{8.15:lem-excess-decay} gives
$|F_0^{(j+1)}-F_0^{(j)}|\le C_S(y_j^{1/2}+\varpi_j)\tau$ for $j<k$. The computation that proves
\eqref{8.15:eq-epsilon-regularity-4}, applied to the indices $j<k$ and using only $(I_j)$ for
$j<k$, gives
\begin{equation*}
\sum_{j<k}\bigl(y_j^{1/2}+\varpi_j\bigr)\le4Y'^{1/2}+2\varpi_R .
\end{equation*}
Summing the bounds on $|F_0^{(j+1)}-F_0^{(j)}|$ and using $Y'\le Y$, $\varpi_R\le c_E$ and (s4),
we get
\begin{equation*}
\begin{aligned}
|F_0^{(k)}|&\le\frac{\tau}{2}+C_S\sum_{j<k}\bigl(y_j^{1/2}+\varpi_j\bigr)\tau
\le\frac{\tau}{2}+C_S\bigl(4Y'^{1/2}+2\varpi_R\bigr)\tau\\
&\le\frac{\tau}{2}+C_S(4Y^{1/2}+2c_E)\tau\le\frac58\tau\le\frac34\tau .
\end{aligned}
\end{equation*}
\item Hypothesis \eqref{8.15:eq-multiplier-mass-a} follows from
$2\kappa_k+1\ge2\varepsilon_Er_k-1\ge\varepsilon_Er_k$ and from the bound of $\|D'\|_{3r_k}^2$
quoted above, together with (s2):
\begin{equation*}
(2\kappa_k+1)^{-4}\|D'\|_{3r_k}^2\le\varepsilon_E^{-4}(4y_k+C\varpi_k^2)\tau^2
\le\varepsilon_E^{-4}(4Y+Cc_E^2)\tau^2\le\Bigl(\frac{\tau}{16}\Bigr)^2 .
\end{equation*}
\end{itemize}
Hence Lemma~\ref{8.15:lem-multiplier-mass} and Lemma~\ref{8.15:lem-averaged-field} apply at
step $k$. Since also $r_k\ge66$ and $\theta_E\in[1/r_k,1/32]$ (as $r_k\ge\rho_E\ge6/\theta_E$),
Lemma~\ref{8.15:lem-excess-decay} applies at step $k$. If $k<k^{\ast}$, condition (s1) gives
$16C_{\sharp}\varepsilon_E^{-4}y_k\le1/8$. So the recursion in
\eqref{8.15:eq-epsilon-regularity-b} holds at step $k$, and the claim gives $(I_{k+1})$. This
completes the induction. Thus the three lemmas apply at every step $k\le k^{\ast}$, $(I_k)$ holds
for all $k\le k^{\ast}$, and \eqref{8.15:eq-epsilon-regularity-4} holds.

\emph{The exact step at unit scale.} A plaquette $p$ based at $c$ lies in $Q_{3r_{k^{\ast}}}(c)$.
The summands of $\mathcal{Y}$ are non-negative. Using $(I_{k^{\ast}})$ and $r_{k^{\ast}}<\rho_E^{+}$,
\begin{equation*}
\tfrac12|\hat{F}(p)-F_0^{(k^{\ast})}|^2\le\mathcal{Y}(3r_{k^{\ast}};F_0^{(k^{\ast})})
=y_{k^{\ast}}\tau^2r_{k^{\ast}}^4\le Y'(\rho_E^{+})^4\tau^2 .
\end{equation*}
Also $|F_0^{(k^{\ast})}-F_0|\le C_SS_\infty\tau$, by summing the bounds
$|F_0^{(j+1)}-F_0^{(j)}|\le C_S(y_j^{1/2}+\varpi_j)\tau$ for $j<k^{\ast}$. Put
\begin{equation*}
C_E:=\bigl(2^{1/2}(\rho_E^{+})^2+4C_S\bigr)\bigl(1+(8C_{\sharp})^{1/2}\bigr)+2C_S .
\end{equation*}
This depends only on $C_S,\theta_E,\varepsilon_E$. We use \eqref{8.15:eq-epsilon-regularity-4}
and $Y'^{1/2}\le y_0^{1/2}+(8C_{\sharp})^{1/2}\varpi_R$. Together with the two bounds above they
give
\begin{equation}\label{8.15:eq-epsilon-regularity-5}
\begin{aligned}
|\hat{F}(p)-F_0|&\le\bigl[(2Y')^{1/2}(\rho_E^{+})^2+C_SS_\infty\bigr]\tau\\
&\le\bigl(2^{1/2}(\rho_E^{+})^2+4C_S\bigr)Y'^{1/2}\tau+2C_S\varpi_R\tau
\le C_E\bigl(y_0^{1/2}+\varpi_R\bigr)\tau .
\end{aligned}
\end{equation}
By (s5), $(2Y')^{1/2}(\rho_E^{+})^2\le(2Y(\rho_E^{+})^4)^{1/2}\le1/8$. By
\eqref{8.15:eq-epsilon-regularity-4} and (s4), $C_SS_\infty\le C_S(4Y^{1/2}+2c_E)\le1/8$.
So the bracket in \eqref{8.15:eq-epsilon-regularity-5} is at most $\frac18+\frac18$, and
$|\hat{F}(p)-F_0|\le\tau/4$. Since $|F_0|\le\tau/2$, it follows that
$|\hat{F}(p)|\le\frac34\tau<\tau$. Complementary slackness in \eqref{8.15:eq-capped-data-a}
then gives $\Lambda(p)=0$.
\end{proof}

\begin{corollary}[Model: inactivity of a whole cube under small excess]
\label{8.15:cor-cube-inactivity}
The model is the one of the model theorem, Theorem~\ref{8.15:thm-epsilon-regularity}: capped lattice
fields on cubes of $\mathbb{Z}^4$. The data are the given $2$-form $\hat{F}$ with cap height $\tau$,
the multiplier magnitude $\Lambda$, the defect function $\varpi$, and the hybrid excess
$\mathcal{Y}(r;F_0)$ with its normalisation $y(r;F_0)$, as in \eqref{8.15:eq-capped-data-c}. We write
$\mathcal{Y}_{c}(r;F_0)$ and $y_{c}(r;F_0)$ for the excess about the centre $c$. Let the data, the
constants and the conditions $R\ge\rho_E$, $|F_0|\le\tau/2$ be as in
Theorem~\ref{8.15:thm-epsilon-regularity}. Suppose that condition (h2) of
\eqref{8.15:eq-epsilon-regularity-a} holds about every centre $c'\in Q_R(c)$, and that
\begin{equation*}
\mathcal{Y}_{c}(4R;F_0)\le y_E\,\tau^2R^4 .
\end{equation*}
Then no plaquette in $Q_R(c)$ is active, and $|\hat{F}-F_0|\le\tau/4$ there.
\end{corollary}

\begin{proof}
Fix a centre $c'\in Q_R(c)$, so that $|c'-c|_\infty\le R$. For $z\in Q_{3R}(c')$ we have
\begin{equation*}
|z-c|_\infty\le|z-c'|_\infty+|c'-c|_\infty\le 3R+R=4R ,
\end{equation*}
so $Q_{3R}(c')\subset Q_{4R}(c)$.

The hybrid excess about $c'$ at radius $R$ is formed from the data on $Q_{3R}(c')$, and these data
lie inside $Q_{4R}(c)$. Hence $\mathcal{Y}_{c'}(R;F_0)\le\mathcal{Y}_{c}(4R;F_0)$, and therefore
$\mathcal{Y}_{c'}(R;F_0)\le y_E\,\tau^2R^4$. By the normalisation in \eqref{8.15:eq-capped-data-c},
this is $y_{c'}(R;F_0)\le y_E$. Thus condition (h1) of \eqref{8.15:eq-epsilon-regularity-a} holds
about $c'$. Condition (h2) holds about $c'$ by hypothesis, and the conditions $R\ge\rho_E$ and
$|F_0|\le\tau/2$ do not depend on the centre.

We now apply the model theorem, Theorem~\ref{8.15:thm-epsilon-regularity}, with centre $c'$. It
gives, for every plaquette $p$ based at $c'$,
\begin{equation*}
|\hat{F}(p)-F_0|\le\tau/4,\qquad |\hat{F}(p)|<\tau,\qquad \Lambda(p)=0 .
\end{equation*}
So $p$ is not active. Every plaquette in $Q_R(c)$ is based at some centre $c'\in Q_R(c)$, and $c'$
was arbitrary. Hence no plaquette in $Q_R(c)$ is active, and $|\hat{F}-F_0|\le\tau/4$ on $Q_R(c)$.
\end{proof}

\begin{remark}[Remarks on $\varepsilon$-regularity: what is used and variants]
\label{8.15:rem-regularity-remarks}
The notation is that of Lemma~\ref{8.15:lem-multiplier-mass}, Lemma~\ref{8.15:lem-excess-decay}
and the model theorem, Theorem~\ref{8.15:thm-epsilon-regularity}.
\begin{enumerate}
\item[(1)] \emph{What is used.} The arguments use the first-order (Karush--Kuhn--Tucker)
system of the capped data, \eqref{8.15:eq-capped-data-a} and \eqref{8.15:eq-capped-data-b},
and nothing else. They use no minimality and no convexity of the problem. They use no
uniqueness of multipliers: any $\Lambda$ satisfying the second condition of
\eqref{8.15:eq-capped-data-a}, the second condition of \eqref{8.15:eq-capped-data-b} and the
first condition of \eqref{8.15:eq-epsilon-regularity-a} will do. They use no order, and no sign
of $\mu$ beyond the second condition of \eqref{8.15:eq-capped-data-a}. They use no column or
row structure, since the four directions are on the same footing. They use no bound on the
kernel of the operator $\Pi$. They use no boundary regularity: all
comparisons are interior, and the Dirichlet projection of the harmonic approximation lemma is
used in $\ell^2$ only.
\item[(2)] \emph{Colour norms.} Let $\|\cdot\|_{\mathrm{c}}$ be any norm on $\mathsf{V}$ with
$c_-|X|\le\|X\|_{\mathrm{c}}\le c_+|X|$ for $X\in\mathsf{V}$, and let the cap read
$\|\hat{F}(p)\|_{\mathrm{c}}\le\tau$. Then $\mu=\Lambda\Theta$ with
$\Theta\in\partial\|\cdot\|_{\mathrm{c}}(\hat{F})$, $\langle\Theta,\hat{F}\rangle=\tau$ and
$\langle\Theta,X\rangle\le\|X\|_{\mathrm{c}}$ for $X\in\mathsf{V}$. The step of the proof of
Lemma~\ref{8.15:lem-multiplier-mass} that bounds $\langle\mu,\hat{F}-\hat{F}_K\rangle$ from
below then reads
\begin{equation}\label{8.15:eq-regularity-remarks-1}
\langle\mu,\hat{F}-\hat{F}_K\rangle\ \ge\ \Lambda\bigl(\tau-\|\hat{F}_K\|_{\mathrm{c}}\bigr),
\end{equation}
and the proofs go through word for word, with constants depending on $c_+/c_-$.
\item[(3)] \emph{Orientation- or position-dependent caps.} Let the cap at $p$ be $C(p)$, with
$C(p)\in[\tau,\varkappa\tau]$ on $Q_{3R}(c)$. The step of the proof of
Lemma~\ref{8.15:lem-multiplier-mass} that bounds $\langle\mu,\hat{F}-\hat{F}_K\rangle$ from
below needs only $|\hat{F}_K|\le\tfrac78\tau\le\tfrac78 C(p)$. The step of that proof that
bounds $D'$ pointwise has
$\sup|D'|\le(\varkappa+1)\tau$. The constants depend on $\varkappa$, with $\varepsilon_E$
proportional to $1/\varkappa$. Translation invariance of the caps is not needed once the
excess is small.
\item[(4)] \emph{Uncapped orientations.} On orientations that carry no cap, the first
condition of \eqref{8.15:eq-capped-data-a} may be replaced by the weak tangential bound,
with any constant $\varkappa$:
\begin{equation}\label{8.15:eq-regularity-remarks-a}
|\hat{F}(p)|\le\varkappa\tau\qquad\text{for the plaquettes $p$ of these orientations in }
Q_{3R}(c).
\end{equation}
In the young band of the setting to which the model theorem is applied, and for the windows
of Ba{\l}aban's construction, every orientation carries a cap. The tangential caps are what
make the commutator term of the multipliers cheap.
\item[(5)] \emph{Why no logarithm.} Two quantities are carried from scale to scale. The first
is the excess, which decays. The second is the reference constant, whose total drift is at
most $C_S\tau\sum_k\bigl(y_k^{1/2}+\varpi_k\bigr)<\infty$. An iteration of the normalised
energy, which does not decay, through
the release identity, applied on suitably chosen radii, pays $(\log r)^{1/2}$ for the boundary
layer. Here no field obtained from the release identity is evaluated near a boundary.
\item[(6)] \emph{The needle witness.} Let $F^{\sharp}$ be the needle witness configuration.
It satisfies the first
condition of \eqref{8.15:eq-capped-data-a}, feasibility and every energy inequality. It
violates the first-order system at the bonds of the needle, and the second condition of
\eqref{8.15:eq-capped-data-b}, a hypothesis of the model theorem,
Theorem~\ref{8.15:thm-epsilon-regularity}, fails for it.
\item[(7)] \emph{Needles.} Consider a needle standing on a weak kink of strength $\nu$. Up to
a constant factor, its excess is at least $\nu^2$ at every scale between its width and its
height. By Lemma~\ref{8.15:lem-excess-decay}, the excess of a field satisfying the
first-order system decays. The two are incompatible as soon as the normalised excess at the
top scale satisfies $y(R;F_0)\le y_E$, the first condition of
\eqref{8.15:eq-epsilon-regularity-a}; the hybrid excess from which $y(R;F_0)$ is formed
contains the needle's multiplier mass. No scale-invariant input
detects the needle. What detects it is the iteration of a scale-invariant inequality with a
gain.
\item[(8)] \emph{The exact step.} A step whose constant is to be uniform in the number of
scales must incur no loss. Here that step is
the unit-scale embedding $\ell^2\subset\ell^\infty$ together with complementary slackness, and
the losses of the steps before it form a convergent series.
\item[(9)] \emph{Defects that do not decay.} Suppose the second condition of
\eqref{8.15:eq-epsilon-regularity-a} is replaced by $\varpi(r)\le\varpi_*$ for all
$r\in[1,R]$, with
no decay, where $\varpi_*\le c_E$. Then, under the remaining hypotheses of the model theorem,
Theorem~\ref{8.15:thm-epsilon-regularity}, and the additional condition
\begin{equation}\label{8.15:eq-regularity-remarks-b}
C_S\bigl(2C_{\sharp}^{1/2}+1\bigr)\varpi_*\bigl(1+\log(R/\rho_E)\bigr)\le\tfrac1{16},
\end{equation}
every plaquette $p$ based at $c$ obeys
\begin{gather*}
|\hat{F}(p)-F_0|\le\bar{C}_E\Bigl(y(R;F_0)^{1/2}+\bigl(1+\log(R/\rho_E)\bigr)\varpi_*\Bigr)\tau,\\
|\hat{F}(p)-F_0|\le\tfrac5{16}\tau,
\end{gather*}
with a constant $\bar{C}_E$ depending only on the dimension $4$, fixed in the proof below; in
particular $|\hat{F}(p)|<\tau$ and $\Lambda(p)=0$.
Here the logarithm multiplies the defect, not a margin. This variant is used only with a
defect of the size of Ba{\l}aban's relative error $\vartheta_1$ of local minimisers, which is
exponentially small.
\end{enumerate}
\end{remark}

\begin{proof}
\emph{Item (2).} Where $\Lambda(p)=0$, both sides of \eqref{8.15:eq-regularity-remarks-1}
vanish. Elsewhere $\mu=\Lambda\Theta$ and $\Lambda\ge0$. The relations
$\langle\Theta,\hat{F}\rangle=\tau$ and
$\langle\Theta,\hat{F}_K\rangle\le\|\hat{F}_K\|_{\mathrm{c}}$ then give
\begin{equation*}
\langle\mu,\hat{F}-\hat{F}_K\rangle=\Lambda\tau-\Lambda\langle\Theta,\hat{F}_K\rangle
\ge\Lambda\bigl(\tau-\|\hat{F}_K\|_{\mathrm{c}}\bigr).
\end{equation*}

\emph{Item (9).} We use the notation of the proof of the model theorem. The radii are
$r_0=R$ and $r_{k+1}=\lfloor\theta_E r_k/3\rfloor$. The index $k^{\ast}$ is the last $k$ with
$r_k\ge\rho_E$. We write $y_k=y(r_k;F_0^{(k)})$ and $\varpi_k=\varpi(r_k)$, and we put
\begin{equation*}
S_\infty=\sum_{k=0}^{k^{\ast}}\bigl(y_k^{1/2}+\varpi_k\bigr).
\end{equation*}

Since $\theta_E\le1/32$, we have $r_{k+1}\le\theta_E r_k/3\le r_k/96$. We write
$Y=y_E+8C_{\sharp}c_E^2$; in the proof of the model theorem the smallness conditions on $y_E$
and $c_E$ are imposed in terms of $Y$ and $c_E$.

\emph{The number of steps.} From $r_{k+1}\le r_k/96$ we get $r_k\le96^{-k}R$. Since
$r_{k^{\ast}}\ge\rho_E$, it follows that $96^{k^{\ast}}\le R/\rho_E$, that is,
\begin{equation*}
k^{\ast}+1\le1+\frac{\log(R/\rho_E)}{\log96}\le1+\log(R/\rho_E).
\end{equation*}
The last inequality uses $\log96>1$ and $R\ge\rho_E$. Consequently, under
\eqref{8.15:eq-regularity-remarks-b},
\begin{equation}\label{8.15:eq-regularity-remarks-3}
C_S(k^{\ast}+1)\bigl(2C_{\sharp}^{1/2}+1\bigr)\varpi_*\le\tfrac1{16}.
\end{equation}

\emph{The induction.} The recursion \eqref{8.15:eq-epsilon-regularity-b} was derived whenever
Lemma~\ref{8.15:lem-excess-decay} applies at step $k$ and
$16C_{\sharp}\varepsilon_E^{-4}y_k\le\tfrac18$, and that derivation does not use the
decay of $\varpi$. With $\varpi_k\le\varpi_*$ it gives, in that case,
\begin{equation*}
y_{k+1}\le\tfrac38y_k+C_{\sharp}\varpi_*^2 .
\end{equation*}
We show by induction on $k$, for $0\le k\le k^{\ast}$, that for $0\le j\le k$
\begin{align*}
y_j&\le(\tfrac38)^jy_0+C_{\sharp}\varpi_*^2\sum_{i=0}^{j-1}(\tfrac38)^i,\\
|F_0^{(j)}-F_0|&\le C_S\tau\sum_{i=0}^{j-1}\bigl(y_i^{1/2}+\varpi_i\bigr).
\end{align*}
For $k=0$ both bounds hold trivially. Assume them for some $k<k^{\ast}$.

For $j\le k$, since $(\tfrac38)^j\le2^{-j}$ and $\sum_{i\ge0}(\tfrac38)^i=\tfrac85\le2$, the
first bound gives
\begin{equation}\label{8.15:eq-regularity-remarks-2}
y_j\le2^{-j}y_0+2C_{\sharp}\varpi_*^2 .
\end{equation}
As $y_0\le y_E$ by the first condition of \eqref{8.15:eq-epsilon-regularity-a} and
$\varpi_*\le c_E$, this yields $y_j\le y_E+2C_{\sharp}c_E^2\le Y$. From
\eqref{8.15:eq-regularity-remarks-2} and $(a+b)^{1/2}\le a^{1/2}+b^{1/2}$ we also get
\begin{equation*}
y_j^{1/2}\le2^{-j/2}y_0^{1/2}+(2C_{\sharp})^{1/2}\varpi_*
\le2^{-j/2}y_0^{1/2}+2C_{\sharp}^{1/2}\varpi_* .
\end{equation*}
Since $\sum_{j\ge0}2^{-j/2}=2+\sqrt2\le4$, $\varpi_j\le\varpi_*$ and $k+1\le k^{\ast}+1$,
summing over $0\le j\le k$ gives
\begin{equation*}
\sum_{j=0}^{k}\bigl(y_j^{1/2}+\varpi_j\bigr)
\le4y_0^{1/2}+(k^{\ast}+1)\bigl(2C_{\sharp}^{1/2}+1\bigr)\varpi_* .
\end{equation*}
The smallness condition $C_S(4Y^{1/2}+2c_E)\le\tfrac18$ of the proof of the model theorem gives
$4C_Sy_0^{1/2}\le\tfrac18$, because $y_0\le Y$. Together with
\eqref{8.15:eq-regularity-remarks-3}, the second bound at $j=k$ and $|F_0|\le\tau/2$ therefore
give
\begin{equation*}
|F_0^{(k)}-F_0|\le\bigl(\tfrac18+\tfrac1{16}\bigr)\tau=\tfrac3{16}\tau,\qquad
|F_0^{(k)}|\le\tfrac{11}{16}\tau\le\tfrac34\tau .
\end{equation*}

Thus at step $k$ we have $y_k\le Y$, $\varpi_k\le c_E$ and $|F_0^{(k)}|\le\tfrac34\tau$. In
the proof of the model theorem the hypotheses of Lemma~\ref{8.15:lem-multiplier-mass} and
Lemma~\ref{8.15:lem-excess-decay} at $(r_k,\kappa_k,\theta_E)$, including the bound
$(2\kappa_k+1)^{-4}\|D'\|_{3r_k}^2\le(\tau/16)^2$ for $D'=dA-F_0^{(k)}$ and the bound
$\varpi_k\le1/32$, were obtained from exactly these
three bounds through the smallness conditions, the conditions on $r_k$, $\kappa_k$ and
$\theta_E$ not involving $\varpi$; hence both lemmas apply at step $k$. Moreover
$16C_{\sharp}\varepsilon_E^{-4}y_k\le16C_{\sharp}\varepsilon_E^{-4}Y\le\tfrac18$ by the
smallness condition of that proof. The recursion above therefore holds at step $k$, and
multiplying the first bound at $j=k$ by $\tfrac38$ and adding $C_{\sharp}\varpi_*^2$ gives it at
$j=k+1$. Lemma~\ref{8.15:lem-excess-decay} also gives
$|F_0^{(k+1)}-F_0^{(k)}|\le C_S(y_k^{1/2}+\varpi_k)\tau$, and adding this to the second bound
at $j=k$ gives it at $j=k+1$. This completes the induction. Hence both bounds hold for
$0\le j\le k^{\ast}$, and the derivation of \eqref{8.15:eq-regularity-remarks-2}, of
$y_j\le Y$ and of the summed estimate above applies word for word with $k=k^{\ast}$.

\emph{The drift.} By the induction, the summed estimate above holds with $k=k^{\ast}$, that is,
\begin{equation*}
S_\infty\le4y_0^{1/2}+(k^{\ast}+1)\bigl(2C_{\sharp}^{1/2}+1\bigr)\varpi_* ,
\end{equation*}
so that, as above, $C_SS_\infty\le\tfrac18+\tfrac1{16}=\tfrac3{16}$ and
$|F_0^{(k^{\ast})}-F_0|\le C_SS_\infty\tau\le\tfrac3{16}\tau$. Only the drift differs from the
proof of the model theorem.

\emph{The exact step.} A plaquette $p$ based at $c$ lies in $Q_{3r_{k^{\ast}}}(c)$. As in the
proof of the model theorem,
\begin{equation*}
\tfrac12|\hat{F}(p)-F_0^{(k^{\ast})}|^2\le\mathcal{Y}(3r_{k^{\ast}};F_0^{(k^{\ast})})
=y_{k^{\ast}}\tau^2r_{k^{\ast}}^4 .
\end{equation*}
First, using $y_{k^{\ast}}\le Y$ and $r_{k^{\ast}}<\rho_E^{+}$, the right-hand side is at most
$Y(\rho_E^{+})^4\tau^2$, and the smallness condition $2Y(\rho_E^{+})^4\le\tfrac1{64}$ gives
$|\hat{F}(p)-F_0^{(k^{\ast})}|\le\tfrac18\tau$. Hence
\begin{equation*}
|\hat{F}(p)-F_0|\le\bigl(\tfrac18+\tfrac3{16}\bigr)\tau=\tfrac5{16}\tau,\qquad
|\hat{F}(p)|\le\tfrac12\tau+\tfrac5{16}\tau=\tfrac{13}{16}\tau<\tau,
\end{equation*}
and the second condition of \eqref{8.15:eq-capped-data-a} gives $\Lambda(p)=0$.

Second, using \eqref{8.15:eq-regularity-remarks-2} at $j=k^{\ast}$, $2^{-k^{\ast}}\le1$,
$r_{k^{\ast}}<\rho_E^{+}$ and $(a+b)^{1/2}\le a^{1/2}+b^{1/2}$, we get
\begin{equation*}
|\hat{F}(p)-F_0^{(k^{\ast})}|\le(2y_{k^{\ast}})^{1/2}r_{k^{\ast}}^2\tau
\le\bigl((2y_0)^{1/2}+2C_{\sharp}^{1/2}\varpi_*\bigr)(\rho_E^{+})^2\tau .
\end{equation*}
By the drift estimate and $k^{\ast}+1\le1+\log(R/\rho_E)$,
\begin{equation*}
|F_0^{(k^{\ast})}-F_0|\le C_SS_\infty\tau
\le C_S\Bigl(4y_0^{1/2}+\bigl(1+\log(R/\rho_E)\bigr)\bigl(2C_{\sharp}^{1/2}+1\bigr)\varpi_*\Bigr)
\tau .
\end{equation*}
We put
\begin{equation*}
\bar{C}_E=\max\Bigl\{\sqrt2(\rho_E^{+})^2+4C_S,\ 2C_{\sharp}^{1/2}(\rho_E^{+})^2
+C_S\bigl(2C_{\sharp}^{1/2}+1\bigr)\Bigr\};
\end{equation*}
it depends only on the dimension, because $C_S$ is absolute and $C_{\sharp}$, $\rho_E^{+}$ are
fixed from $C_S$, $\theta_E$ and $\rho_E$ in the proof of the model theorem. Adding the two
last estimates, using $1\le1+\log(R/\rho_E)$ and $y_0=y(R;F_0)$, we obtain
\begin{equation*}
|\hat{F}(p)-F_0|\le\bar{C}_E\Bigl(y(R;F_0)^{1/2}+\bigl(1+\log(R/\rho_E)\bigr)\varpi_*\Bigr)\tau .
\end{equation*}
\end{proof}

\begin{definition}[Notation for the capped slab problems]\label{8.16:not-slab-notation}
Throughout what follows we work in the slab setting fixed earlier in this chapter. That is, we
are given the slab $\mathfrak{S}^4=(\mathbb{Z}/N)^3\times\{0,\dots,H\}$, where, in the slab
problems only, $N$ denotes the period of the slab in its first three coordinates (not the rank of
the gauge group) and $H$ its height (not the Hamiltonian); its two
faces $A^0$ and $A^H$; the face curvature bound $\ell$ of that setting; and the transversal
face mismatch $\ell_{\times}$ of the two faces $A^0$, $A^H$.

\emph{Capped problems.} For a cap function $C$ on the plaquettes of the slab, $P^{C}$ denotes
the minimisation problem of the slab setting with caps given by $C$. We write $\hat{F}^{C}$ for
the curvature of its minimisers, which is the same for all of them, the minimisers of a capped
slab problem having one common curvature in the slab setting. For a cap height $\tau$,
$P^{\tau}$ denotes the problem $P^{C}$ with
$C\equiv\tau$ on every plaquette of the slab.

\emph{Combined face constants.} We put
\begin{equation}\label{8.16:eq-slab-notation-1}
\ell_c:=\max\Bigl(\ell,\ \frac{\ell_{\times}}{H}\Bigr),\qquad
\mathfrak{E}_c:=6N^3H\Bigl(\ell^2+\frac{\ell_{\times}^2}{H^2}\Bigr).
\end{equation}

\emph{Base row.} The \emph{base row} of a plaquette is the fourth coordinate of its base
site, the terms base site, normal and tangential plaquette, level and row being those of the
slab setting. Thus a normal plaquette of level $y$ has base row $y$, and so does a tangential
plaquette of row $y$.
\end{definition}

\begin{theorem}[Model: exact inactivity away from the slab faces]\label{8.16:thm-slab-inactivity}
The model is the capped slab problem of Definition~\ref{8.16:not-slab-notation}. On the slab
$(\mathbb{Z}/N)^3\times\{0,\dots,H\}$ one minimises the quadratic energy
$E_4(F)=\tfrac12\|F\|^2=\tfrac12\sum_pF(p)^2$ of the curvature $F=dA$ of a real-valued
$1$-form $A$. The minimisation runs over those $A$ whose tangential bonds on the two faces are
held at the face data $A^0$, $A^H$, and whose curvature obeys the caps $|F(p)|\le C(p)$, with
$C(p)\in(0,\infty]$ given for each plaquette. Every problem $P^{C}$ of the capped slab setting on
which Definition~\ref{8.16:not-slab-notation} is built has the following two properties. They
belong to that setting and are used here as stated with it; they are not hypotheses of the
present theorem.
\begin{itemize}
\item[] \emph{Uniqueness of the curvature.} All minimisers of $P^{C}$ have one and the same
curvature $\hat{F}^{C}$.
\item[] \emph{Multipliers.} There is a function $\Lambda\ge0$ on plaquettes, with $\Lambda(p)=0$
unless $|\hat{F}^{C}(p)|=C(p)<\infty$, with the following property. Put
$s_p:=\operatorname{sign}\hat{F}^{C}(p)$ and $\mu(p):=\Lambda(p)s_p$. Then
$\langle\hat{F}^{C}+\mu,d\xi\rangle=0$ for every $1$-form $\xi$ vanishing on the held bonds.
\end{itemize}
Any such $\Lambda$ is called a family of \emph{multipliers} of $P^{C}$. The \emph{row} of a bond
is the fourth coordinate of its base site.

Let $\tau>0$. Let $\rho_E$, $y_E$, $C_E$ be the constants of the model theorem,
Theorem~\ref{8.15:thm-epsilon-regularity}. Define $R_V$ by the second entry of the following
display, and assume the three inequalities it contains:
\begin{equation}\label{8.16:eq-slab-inactivity-a}
\begin{gathered}
\ell_c\le\tfrac{\tau}{2},\qquad
R_V:=\Bigl\lceil\max\Bigl(\rho_E,\Bigl(\frac{\mathfrak{E}_c}{y_E\tau^2}\Bigr)^{1/4}\Bigr)\Bigr\rceil,
\\
6R_V+3\le N,\qquad 6R_V+1\le H.
\end{gathered}
\end{equation}
Then the following hold.
\begin{itemize}
\item[(a)] \emph{Energy and multiplier mass at the natural scale.} We have
$\|\hat{F}^{\tau}\|^2\le\mathfrak{E}_c$. Moreover, for any multipliers $\Lambda$ of $P^{\tau}$,
$\tau\sum_p\Lambda(p)\le\mathfrak{E}_c/2$.
\item[(b)] \emph{Exact inactivity.} Let $p$ be a plaquette whose base row lies in
$[3R_V,H-1-3R_V]$. Then
\[
|\hat{F}^{\tau}(p)|\le C_E\bigl(\mathfrak{E}_c/R_V^4\bigr)^{1/2}\le\tau/4 ,
\]
and $\Lambda(p)=0$ for any multipliers $\Lambda$ of $P^{\tau}$.
\item[(c)] \emph{One field for all problems.} Let $C$ be a cap function with $C\ge\tau$
everywhere and $C=\tau$ on the plaquettes of base row $<3R_V$ or $>H-1-3R_V$. Then
$\hat{F}^{C}=\hat{F}^{\tau}$.

In particular, fix $W_y\ge3R_V$. The following problems all have the same minimising curvature
$\hat{F}^{\tau}$:
\begin{itemize}
\item[--] the problem with cap $\tau$ within rows $[0,W_y]\cup[H-W_y,H]$ and no cap elsewhere,
which we call the young-only problem;
\item[--] every truncation of the young-only problem at a height $a\in[3R_V,W_y]$, by which we
mean the problem with cap $\tau$ within rows $[0,a]\cup[H-a,H]$ and no cap elsewhere;
\item[--] every problem obtained from these by adding caps $\ge\tau$ on deeper plaquettes, for
instance nested rings with caps $\tau\mathsf{T},\tau\mathsf{T}^2,\dots$, where $\mathsf{T}\ge1$
is a fixed ratio, so that these caps are $\ge\tau$, with the innermost region, inside all the
rings, left without cap.
\end{itemize}
This common curvature satisfies the lattice Maxwell equation $d^*\hat{F}^{\tau}=0$ exactly at
every bond of row in $[3R_V+1,H-3R_V-1]$, where $d^*$ denotes the adjoint of $d$ for the
$\ell^2$ pairings of $1$-forms and $2$-forms. Every plaquette containing such a bond has
$|\hat{F}^{\tau}|\le\tau/4$.
\end{itemize}
\end{theorem}

\begin{proof}
Fix a minimiser $\hat{A}$ of $P^{\tau}$, so that $d\hat{A}=\hat{F}^{\tau}$. Fix a family $\Lambda$
of multipliers of $P^{\tau}$, and put $s_p=\operatorname{sign}\hat{F}^{\tau}(p)$ and
$\mu=\Lambda s$. Since every cap of $P^{\tau}$ equals $\tau$, we have $\Lambda(p)\neq0$ only where
$|\hat{F}^{\tau}(p)|=\tau$. At such plaquettes $s_p\hat{F}^{\tau}(p)=\tau$.

(a) We use the linear interpolation between the faces. It is given by
\[
A^{\mathrm{comp}}_j:=(1-y/H)A^0_j+(y/H)A^H_j\quad(j=1,2,3),\qquad A^{\mathrm{comp}}_4:=\lambda/H,
\]
where $\lambda$ is a function on $(\mathbb{Z}/N)^3$ at which the infimum defining $\ell_{\times}$
is attained; the infimum is the $\|\cdot\|_\infty$-distance from $A^H-A^0$ to the
finite-dimensional space of the $d\lambda$, hence it is attained. By the bound for the linear
face interpolation that accompanies the slab setting, its curvature $F^{\mathrm{comp}}$ obeys
\[
|F^{\mathrm{comp}}_{ij}|\le\ell,\qquad |F^{\mathrm{comp}}_{4j}|\le\ell_{\times}/H,\qquad
\sum_p|F^{\mathrm{comp}}(p)|^2\le 6N^3H\bigl(\ell^2+\ell_{\times}^2/H^2\bigr)=\mathfrak{E}_c .
\]
At rows $0$ and $H$ the tangential components of $A^{\mathrm{comp}}$ equal $A^0_j$ and $A^H_j$.
Hence $A^{\mathrm{comp}}$ takes the held values. Both $\ell$ and $\ell_{\times}/H$ are at most
$\ell_c$, and $\ell_c\le\tau/2$ by \eqref{8.16:eq-slab-inactivity-a}. Therefore
$|F^{\mathrm{comp}}(p)|\le\tau/2\le\tau$ for every plaquette $p$, and $A^{\mathrm{comp}}$ is
feasible for $P^{\tau}$. As $\hat{A}$ minimises $E_4=\frac12\|\cdot\|^2$ over the feasible set,
\[
\|\hat{F}^{\tau}\|^2\le\|F^{\mathrm{comp}}\|^2\le\mathfrak{E}_c .
\]

Now apply the multiplier property of $P^{\tau}$ to $\xi:=\hat{A}-A^{\mathrm{comp}}$. This $\xi$
vanishes on the held bonds, because both potentials take the held values there, and
$d\xi=\hat{F}^{\tau}-F^{\mathrm{comp}}$. The identity
$\langle\hat{F}^{\tau}+\mu,d\xi\rangle=0$ gives
\[
\langle\mu,\hat{F}^{\tau}-F^{\mathrm{comp}}\rangle
=-\langle\hat{F}^{\tau},\hat{F}^{\tau}-F^{\mathrm{comp}}\rangle
=-\|\hat{F}^{\tau}\|^2+\langle\hat{F}^{\tau},F^{\mathrm{comp}}\rangle
\le\tfrac14\|F^{\mathrm{comp}}\|^2 .
\]
The last step uses the Cauchy--Schwarz inequality together with $-u^2+uv\le v^2/4$ for real
$u,v$, applied with $u=\|\hat{F}^{\tau}\|$ and $v=\|F^{\mathrm{comp}}\|$.

On the other side, $\Lambda(p)\ne0$ forces $s_p\hat{F}^{\tau}(p)=\tau$. Also
$s_pF^{\mathrm{comp}}(p)\le|F^{\mathrm{comp}}(p)|\le\tau/2$. Hence
\[
\langle\mu,\hat{F}^{\tau}-F^{\mathrm{comp}}\rangle
=\sum_p\Lambda(p)\bigl[\tau-s_pF^{\mathrm{comp}}(p)\bigr]\ge\frac{\tau}{2}\sum_p\Lambda(p).
\]
Combining the two displays gives
\[
\frac{\tau}{2}\sum_p\Lambda(p)\le\frac14\|F^{\mathrm{comp}}\|^2\le\frac14\mathfrak{E}_c ,
\]
that is, $\tau\sum_p\Lambda(p)\le\mathfrak{E}_c/2$.

(b) Let $p$ have base site $c=(x,y)$, with $x\in(\mathbb{Z}/N)^3$ and $3R_V\le y\le H-1-3R_V$.
The projection of the cube $Q_{3R_V+1}(c)$ onto the first three coordinates has side
$6R_V+3\le N$, by \eqref{8.16:eq-slab-inactivity-a}. Hence the cube injects into
$(\mathbb{Z}/N)^3\times\mathbb{Z}$, and we read it, and the data on it, in $\mathbb{Z}^4$.

The sites of $Q_{3R_V}(c)$ have fourth coordinates in $[y-3R_V,y+3R_V]\subset[0,H-1]$. So the
plaquettes in $Q_{3R_V}(c)$ have base rows in $[0,H-1]$. They exist in the slab and carry the cap
$\tau$. For every plaquette $p'$ in $Q_{3R_V}(c)$, feasibility gives
$|\hat{F}^{\tau}(p')|\le\tau$, which is condition (e0) of
\eqref{8.15:eq-capped-data-a}. The multiplier property gives $\Lambda(p')=0$ unless
$|\hat{F}^{\tau}(p')|=\tau$, which is condition (e1) there. On the support of $\Lambda$ we have
$\mu=\Lambda s=\Lambda\hat{F}^{\tau}/\tau$, in agreement with
Definition~\ref{8.15:def-capped-data}.

A bond with an endpoint in $Q_{3R_V-1}(c)$ has that endpoint at a row in
$[y-3R_V+1,y+3R_V-1]\subset[1,H-2]$. Such a bond is therefore normal, or tangential of a row in
$[1,H-2]$; in either case it is free, not held. Let $1\le r\le R_V$. Every
$\xi\in\mathfrak{X}_{3r}$ is supported on bonds with an endpoint in
$Q_{3r-1}(c)\subset Q_{3R_V-1}(c)$, so it vanishes on the held bonds. The multiplier property then
gives $\langle\hat{F}^{\tau}+\mu,d\xi\rangle=0$. This is condition (D2)$_r$ of
\eqref{8.15:eq-capped-data-b} with $\varpi\equiv0$. Since $\hat{F}^{\tau}=d\hat{A}$,
condition (D1)$_r$ of \eqref{8.15:eq-capped-data-b} holds with $\varpi\equiv0$ and $A=\hat{A}$.
Consequently (h2) of \eqref{8.15:eq-epsilon-regularity-a} holds with $\varpi_R=0\le c_E$.

For (h1) we use the definition \eqref{8.15:eq-capped-data-c} of the normalised excess with
$F_0=0$, then part (a), then $R_V^4\ge\mathfrak{E}_c/(y_E\tau^2)$, which holds by the choice of
$R_V$ in \eqref{8.16:eq-slab-inactivity-a}:
\[
y(R_V;0)\le\frac{\tfrac12\|\hat{F}^{\tau}\|^2+\tau\sum_p\Lambda(p)}{\tau^2R_V^4}
\le\frac{\mathfrak{E}_c}{\tau^2R_V^4}\le y_E .
\]
The remaining hypotheses also hold: $R_V\ge\rho_E$ by the choice of $R_V$, and $|F_0|=0\le\tau/2$.

We may therefore apply the model theorem, Theorem~\ref{8.15:thm-epsilon-regularity}, with colour
space $\mathbb{R}$, $R:=R_V$ and $F_0:=0$. It gives, for the plaquette $p$ based at $c$,
\[
|\hat{F}^{\tau}(p)|\le C_E\,y(R_V;0)^{1/2}\,\tau\le\tau/4 ,
\]
together with $|\hat{F}^{\tau}(p)|<\tau$ and $\Lambda(p)=0$. By the bound on $y(R_V;0)$ above,
\[
C_E\,y(R_V;0)^{1/2}\tau\le C_E\Bigl(\frac{\mathfrak{E}_c}{\tau^2R_V^4}\Bigr)^{1/2}\tau
=C_E\Bigl(\frac{\mathfrak{E}_c}{R_V^4}\Bigr)^{1/2}.
\]
The family $\Lambda$ of multipliers was arbitrary, and so was
$p$ with base row in $[3R_V,H-1-3R_V]$. This proves (b).

(c) The potential $\hat{A}$ takes the held values. It also satisfies
$|\hat{F}^{\tau}|\le\tau\le C$, so it is feasible for $P^{C}$. By (b), $\Lambda$ vanishes on
every plaquette of base row in $[3R_V,H-1-3R_V]$. Hence the support of $\Lambda$ consists of
plaquettes of base row $<3R_V$ or $>H-1-3R_V$, where $C=\tau$.

Let $A$ be feasible for $P^{C}$ and $F=dA$. Then $\xi:=A-\hat{A}$ vanishes on the held bonds, and
$F=\hat{F}^{\tau}+d\xi$. Expanding the quadratic energy and using the multiplier property of
$P^{\tau}$, which gives $\langle\hat{F}^{\tau},d\xi\rangle=-\langle\mu,d\xi\rangle$, we get
\begin{align*}
E_4(F)-E_4(\hat{F}^{\tau})
&=\langle\hat{F}^{\tau},d\xi\rangle+\tfrac12\|d\xi\|^2
=-\langle\mu,F-\hat{F}^{\tau}\rangle+\tfrac12\|d\xi\|^2\\
&=\sum_p\Lambda(p)\bigl[\tau-s_pF(p)\bigr]+\tfrac12\|d\xi\|^2\ \ge\ \tfrac12\|d\xi\|^2 .
\end{align*}
The third expression uses $s_p\hat{F}^{\tau}(p)=\tau$ wherever $\Lambda(p)\neq0$. The final
inequality holds because $s_pF(p)\le|F(p)|\le C(p)=\tau$ on the support of $\Lambda$. Thus
$\hat{A}$ minimises $P^{C}$. By the uniqueness of the curvature, $\hat{F}^{\tau}$ is the minimising
curvature of $P^{C}$, that is, $\hat{F}^{C}=\hat{F}^{\tau}$.

We now check that the listed problems satisfy the hypotheses of (c). Each of them has cap function
$\ge\tau$ everywhere, where an absent cap counts as $C=\infty$. Put $a:=W_y$ for the untruncated
instance. Consider a plaquette of base row
$<3R_V$: its rows lie in $[0,3R_V]\subset[0,a]\subset[0,W_y]$. Consider a plaquette of base row
$>H-1-3R_V$: its rows lie in $[H-3R_V,H]\subset[H-a,H]\subset[H-W_y,H]$. In both cases the
plaquette carries the cap $\tau$, since $W_y\ge a\ge3R_V$. Adding caps $\ge\tau$ on deeper
plaquettes does not change these caps. So (c) applies to each of these problems.

Finally, let $b$ be a bond that is not held, and suppose $\Lambda$ vanishes on all plaquettes
containing $b$. Take for $\xi$ the $1$-form equal to one on $b$ and zero elsewhere. By the
definition of the adjoint $d^*$, for a $2$-form $G$ the value $(d^*G)(b)$ is the sum of $\pm G(p)$
over the plaquettes $p$ containing $b$, the sign being the orientation of $b$ in $\partial p$.
Since $\mu=\Lambda s$ vanishes on these plaquettes,
\[
\langle\hat{F}^{\tau}+\mu,d\xi\rangle=\bigl(d^*(\hat{F}^{\tau}+\mu)\bigr)(b)
=(d^*\hat{F}^{\tau})(b).
\]
The multiplier property then reads $(d^*\hat{F}^{\tau})(b)=0$.

Now let $b$ be a bond of row $y'\in[3R_V+1,H-3R_V-1]$. Such a bond is not held. A plaquette
containing it has base row $y'$ or $y'-1$, hence base row in $[3R_V,H-1-3R_V]$. By (b), $\Lambda$
vanishes on all such plaquettes, and $|\hat{F}^{\tau}|\le\tau/4$ on them. Therefore
$d^*\hat{F}^{\tau}=0$ at $b$.
\end{proof}

\begin{corollary}[Model: exit, no cascade and stopping in the slab]\label{8.16:cor-slab-exit}
The model is that of Theorem~\ref{8.16:thm-slab-inactivity}. It consists of the convex problems
$P^{C}$ on the slab $\mathfrak{S}^4=(\mathbb{Z}/N)^3\times\{0,\dots,H\}$: the energy of the
curvature of a potential is minimised, with the held face bonds fixed, under the caps
$|F(p)|\le C(p)$. The minimising curvature of $P^{C}$ is $\hat{F}^{C}$, and $P^{\tau}$ is the
problem with cap $\tau$ on every plaquette. Let $\tau>0$ satisfy (V0), let $R_V$ be the integer
of \eqref{8.16:eq-slab-inactivity-a}, and let (V1) and (V2) hold.

Let $W_y$ be a non-negative integer with $W_y<H-W_y$. The young-only instance of
Theorem~\ref{8.16:thm-slab-inactivity}\textup{(c)} is the problem with cap $\tau$ within rows
$[0,W_y]\cup[H-W_y,H]$ and no cap elsewhere; thus $W_y$ is the fine depth of the zone capped at
$\tau$, and the two capped zones are disjoint.
Assume
\begin{equation}\label{8.16:eq-slab-exit-1}
3R_V\le W_y .
\end{equation}

The $\tau$-capped zone is the \emph{outer ring}; the \emph{deeper rings} are the nested zones beyond
it. They carry the caps
$\tau\mathsf{T},\tau\mathsf{T}^2,\dots$, where $\mathsf{T}\ge1$ is the ratio of successive caps,
so that every deeper ring carries a cap $\ge\tau$.
The plaquettes of a ring, each tested against its cap, are its \emph{windows}. The \emph{seed} is
the distinguished plaquette, which is either left uncapped or carries the cap
$\tau\mathsf{T}^{j_\ast}$, where $j_\ast\ge0$ is the stopping index of the ladder, that is, the
index of its stopping statement; here the \emph{ladder} is the chain of implications of the exit
construction that leads from the exit bound through its successive \emph{rungs} to the stopping
statement, and from the stopping statement to the half-skin bound at the canonical comparison
point. To \emph{truncate at height $a$} is
to keep the cap $\tau$ on the rows within height $a$ of the faces and to drop every other cap; the
dropped caps are the \emph{dropped tests}. Let $\eta$ be the fine height of one row. Let $n_f$ be
the row count of the ladder, in the sense fixed with the ladder: the least $n$ for which its $n$-th
rung holds; in this model the
$n$-th rung reads: the minimiser truncated at height $n\eta$ passes every dropped test.

In the items below, \emph{needles}, their \emph{crossing} a region and their \emph{charge} are
those fixed with the exit bound, the charge being the multiplier charge $\nabla^{\ast}\mu$; a needle
is \emph{active} if one of its plaquettes carries a non-zero multiplier or has curvature at its
cap, inactivity of a plaquette meaning, as in Theorem~\ref{8.16:thm-slab-inactivity}\textup{(b)},
that it carries no multiplier and has curvature below its cap. Then the following hold in this
model: the slab statement, the exit bound, the rungs of the ladder at heights at least $3R_V$
(no cascade), the stopping statement, the
half-skin bound at the canonical comparison point, and the charge statement, each read in the
terms of the model; moreover the front, defined in item \textup{(vii)}, has the size given there.
\begin{itemize}
\item[(i)] \emph{The slab statement.} No active needle crosses the region of base rows
$[3R_V,H-1-3R_V]$. This holds for that whole region, not only for one thin slab.
\item[(ii)] \emph{The exit bound.} At the first free normal level of the young-only instance,
which is the level $W_y$, the minimising curvature is at most $\tau/4$ in absolute value; here
normal bonds are those
transverse to the faces of $\mathfrak{S}^4$, and a level is free in the sense fixed with the exit
bound. This level is the distinguished layer of two statements of the exit construction: the layer
at which the first of them attains its worst bound, and the sheltered level of the second, each in
the sense fixed with those statements.
\item[(iii)] \emph{The rungs of the ladder at heights at least $3R_V$ (no cascade).} For every
height $a\ge3R_V$, the
minimiser truncated at $a$ passes every dropped test. In rows of fine height $\eta$ this gives
\begin{equation}\label{8.16:eq-slab-exit-2}
n_f\le n_V:=\lceil 3R_V/\eta\rceil .
\end{equation}
\item[(iv)] \emph{The stopping statement.} No window of a deeper ring (these carry caps
$\ge\tau$) carries a multiplier. No ring is crossed.
\item[(v)] \emph{The half-skin bound at the canonical comparison point.} Let $p$ be a seed
plaquette, capped or not, with base row in
$[3R_V,H-1-3R_V]$. Then the minimising curvature satisfies
\begin{equation*}
|\hat{F}(p)|\le C_E(\mathfrak{E}_c/R_V^4)^{1/2}\le\tau/4 .
\end{equation*}
Measured against its own cap $\tau\mathsf{T}^{j_\ast}$, this is at most
$\tfrac14\mathsf{T}^{-j_\ast}$.
\item[(vi)] \emph{The charge statement.} There is no active needle, charged or not, above base
row $3R_V$, the words ``above base row $3R_V$'' being read with respect to both faces of the
slab, that is, as ``in the region of base rows $[3R_V,H-1-3R_V]$''.
\item[(vii)] \emph{Size of the front and the conditions.} The integer $R_V$ satisfies
\begin{equation}\label{8.16:eq-slab-exit-3}
R_V\le\max\Bigl(\rho_E,\ (6/y_E)^{1/4}(N^3H)^{1/4}(2^{1/2}\ell_c/\tau)^{1/2}\Bigr)+1 ,
\end{equation}
and hence
\begin{equation}\label{8.16:eq-slab-exit-4}
3R_V\le 3(12/y_E)^{1/4}(N/H)^{3/4}(\ell_c/\tau)^{1/2}\,H+3\rho_E+3 .
\end{equation}
The \emph{front} is the layer of depth $3R_V$ at each face, outside which part (b) of the model
theorem, Theorem~\ref{8.16:thm-slab-inactivity}, gives exact inactivity. By
\eqref{8.16:eq-slab-exit-4} its depth is at most a fixed fraction of the slab, proportional to
$(\ell_c/\tau)^{1/2}$, plus $3\rho_E+3$ lattice units. No
logarithm of $N$, of $H$, or of anything else enters.

\emph{Interpretation of the conditions.} Identify the model's parameters as follows. The cap
exceeds the face curvature bound by one ratio, $\tau=\ell\mathsf{T}$. The sides $N$ and $H$ are
comparable. The quantity $W_y$ is the fine depth of the zone capped at $\tau$, that is, of the
outer ring. Here $N$, $H$ and $W_y$ are fixed multiples of one fine unit; the outer ring consists
of $W_y/\eta$ rows of fine height $\eta$, and the ratios $N/H$ and $H/W_y$ are those fixed in the
renormalisation setting so identified. Then
\begin{equation*}
n_V\le3(12/y_E)^{1/4}\Bigl(\frac{N}{H}\Bigr)^{3/4}\Bigl(\frac{\ell_c}{\ell}\Bigr)^{1/2}
\frac{H}{\eta}\,\mathsf{T}^{-1/2}+\frac{3\rho_E+3}{\eta}+1 ,
\end{equation*}
a number of rows bounded independently of the cutoff and of the coupling: since $\eta$, like $N$,
$H$ and $W_y$, is a fixed multiple of the fine unit, the ratio $H/\eta$ does not depend on that
unit, and the term $(3\rho_E+3)/\eta$ is at most its value when the fine unit is one lattice unit.
Put
\begin{equation*}
C_{\mathrm{VS}}:=36(12/y_E)^{1/2},
\end{equation*}
and consider the conditions
\begin{equation}\label{8.16:eq-slab-exit-a}
\mathsf{T}\ge C_{\mathrm{VS}}\,\frac{\ell_c}{\ell}\,\Bigl(\frac{N}{H}\Bigr)^{3/2}
\Bigl(\frac{H}{W_y}\Bigr)^{2},
\end{equation}
\begin{equation}\label{8.16:eq-slab-exit-5}
6(\rho_E+2)\le W_y,\qquad 6(\rho_E+2)+3\le N .
\end{equation}
Condition \eqref{8.16:eq-slab-exit-a} is one-sided: it is a lower bound on $\mathsf{T}$, and in
this identification, where $\mathsf{T}$ carries the factor $L^{2}$, a floor on $L$. It contains no
logarithm. The common fine unit cancels from the ratios
$N/H$ and $H/W_y$. The conditions \eqref{8.16:eq-slab-exit-5} are lower bounds on $W_y$ and $N$:
the outer ring and the periodic directions of $\mathfrak{S}^4$ each contain at least a fixed number
of lattice units. If \eqref{8.16:eq-slab-exit-a} and \eqref{8.16:eq-slab-exit-5} hold, then
\eqref{8.16:eq-slab-exit-1} holds and $n_V\le W_y/\eta$, the number of rows of the outer ring.
\end{itemize}
\end{corollary}

\begin{proof}
Items (i)--(vi) are read off the model theorem, Theorem~\ref{8.16:thm-slab-inactivity}; item (vii)
is a computation from the definition of $R_V$. The hypotheses (V0)--(V2) of that theorem are those
assumed here.

Part (b) of that theorem gives the following. Every plaquette $p$ with base row in
$[3R_V,H-1-3R_V]$ satisfies
\begin{equation*}
|\hat{F}^{\tau}(p)|\le C_E(\mathfrak{E}_c/R_V^4)^{1/2}\le\tau/4
\qquad\text{and}\qquad \Lambda(p)=0 .
\end{equation*}

Part (c) gives the following. Let $C$ be a cap function with $C\ge\tau$ everywhere and $C=\tau$ on
the plaquettes of base row $<3R_V$ or $>H-1-3R_V$. Then $\hat{F}^{C}=\hat{F}^{\tau}$.

(i) By (b), no plaquette of base rows $[3R_V,H-1-3R_V]$ carries a non-zero multiplier. Since
$\tau/4<\tau$, no such plaquette reaches its cap either. So no active needle meets the region, and
in particular none crosses it.

(ii) By (c), with $W_y\ge3R_V$ from \eqref{8.16:eq-slab-exit-1}, the young-only instance has
minimising curvature $\hat{F}^{\tau}$. The level $W_y$ has base row at least $3R_V$. Since the
capped zones are disjoint, $3R_V\le W_y<H-W_y$ gives $3R_V+W_y<H$, that is,
$W_y\le H-1-3R_V$; so the level lies in the region of (b). By (b) the curvature there is at most
$\tau/4$ in absolute value.

(iii) Let $a\ge3R_V$. The problem truncated at height $a$ has cap $\tau$ on every plaquette of base
row $<3R_V$ or $>H-1-3R_V$, and elsewhere a cap $\tau$ or none. Its cap function therefore meets
the condition of (c), and its minimising curvature is $\hat{F}^{\tau}$. Since $\hat{F}^{\tau}$ is
feasible for $P^{\tau}$, we have $|\hat{F}^{\tau}|\le\tau$ everywhere. Every dropped test is a cap
$\ge\tau$, so it is passed. Now take $n=n_V$. Then $n_V\eta\ge3R_V$ by the definition of $n_V$, so
the truncation at height $n_V\eta$ passes every dropped test, that is, the $n_V$-th rung holds.
By the definition of $n_f$ this gives \eqref{8.16:eq-slab-exit-2}.

(iv) Adding the nested rings with caps $\tau\mathsf{T},\tau\mathsf{T}^2,\dots\ge\tau$ (recall
$\mathsf{T}\ge1$) on deeper
plaquettes, with the seed capped or left uncapped, yields a problem to which (c) applies. Its
minimising curvature is $\hat{F}^{\tau}$. The windows of the deeper rings lie at base rows in
$[3R_V,H-1-3R_V]$. On those windows $|\hat{F}^{\tau}|$ is at most $\tau/4$ by (b), which is
strictly below their caps, so by complementary slackness every choice of multipliers of that
problem vanishes there. So no window of a deeper ring
carries a multiplier or reaches its cap, and no deeper ring is crossed. The outer ring is not
crossed either: it has depth $W_y\ge3R_V$, so a crossing of it passes through its rows of base row
in $[3R_V,W_y]$, which lie in the region of (b) by the computation in (ii); there no plaquette
carries a multiplier or reaches its cap.

(v) Whether the seed is capped by $\tau\mathsf{T}^{j_\ast}\ge\tau$ or uncapped, (c) gives
$\hat{F}=\hat{F}^{\tau}$. Part (b) then gives the displayed bound at the seed. Dividing $\tau/4$ by
the cap $\tau\mathsf{T}^{j_\ast}$ gives $\tfrac14\mathsf{T}^{-j_\ast}$.

(vi) By (b) and (c), every plaquette of base row in $[3R_V,H-1-3R_V]$ has zero multiplier and
curvature at most $\tau/4<\tau$, whatever caps $\ge\tau$ are placed there. Hence no needle there is
active, whether or not its multiplier charge vanishes.

(vii) By definition, $R_V$ is the least integer
$\ge\max(\rho_E,(\mathfrak{E}_c/(y_E\tau^2))^{1/4})$. Hence
\begin{equation*}
R_V\le\max\bigl(\rho_E,(\mathfrak{E}_c/(y_E\tau^2))^{1/4}\bigr)+1 .
\end{equation*}
The face-interpolation energy is $\mathfrak{E}_c=6N^3H(\ell^2+\ell_\times^2/H^2)$, and
$\ell_c=\max(\ell,\ell_\times/H)$. It follows that
\begin{equation*}
\ell^2+\ell_\times^2/H^2\le2\ell_c^2 ,
\qquad
\frac{\mathfrak{E}_c}{y_E\tau^2}\le\frac{6}{y_E}\,N^3H\,\frac{2\ell_c^2}{\tau^2}.
\end{equation*}
Taking fourth roots gives \eqref{8.16:eq-slab-exit-3}.
We have $(6/y_E)^{1/4}2^{1/4}=(12/y_E)^{1/4}$ and $(N^3H)^{1/4}=(N/H)^{3/4}H$. Bounding the maximum
by the sum therefore gives \eqref{8.16:eq-slab-exit-4}.

By the definition of $n_V$ in \eqref{8.16:eq-slab-exit-2}, $n_V\le3R_V/\eta+1$. Inserting
\eqref{8.16:eq-slab-exit-4} and $(\ell_c/\tau)^{1/2}=(\ell_c/\ell)^{1/2}\mathsf{T}^{-1/2}$, which
holds since $\tau=\ell\mathsf{T}$, gives the displayed bound on $n_V$.

We check that the conditions give \eqref{8.16:eq-slab-exit-1}. Since $\tau=\ell\mathsf{T}$,
condition \eqref{8.16:eq-slab-exit-a} is equivalent to
\begin{equation*}
\ell_c/\tau\le W_y^2\big/\bigl(36(12/y_E)^{1/2}(N/H)^{3/2}H^2\bigr),
\end{equation*}
that is, to
\begin{equation*}
3(12/y_E)^{1/4}(N/H)^{3/4}(\ell_c/\tau)^{1/2}H\le W_y/2 .
\end{equation*}
The first condition in \eqref{8.16:eq-slab-exit-5} gives $3(\rho_E+1)\le W_y/2$. Inserting both
bounds into \eqref{8.16:eq-slab-exit-4} yields
\begin{equation*}
3R_V\le W_y/2+W_y/2=W_y .
\end{equation*}
Under \eqref{8.16:eq-slab-exit-a} and \eqref{8.16:eq-slab-exit-5} we therefore have
$3R_V\le W_y$, so $n_V\le W_y/\eta$, the number of rows of the outer ring.
\end{proof}

\begin{theorem}[The exponentially thin active sliver]\label{8.16:thm-active-sliver}
In the slab setting of the model theorem, Theorem~\ref{8.16:thm-slab-inactivity} (face data with
curvature bound $\ell$, a cap $\tau>0$, the problem $P^{\tau}$ in which every plaquette carries the cap
$\tau$, its minimising curvature $\hat{F}^{\tau}$ and its multipliers $\Lambda$), let $b$ be an integer
with $\sigma_S(b)\le\tau/2$. Assume the statements (E0)--(E2) below as hypotheses; in them $\sigma_S(b)$,
$\lambda_S$ and $\mathfrak{e}(b)$ are given quantities, together with a given reference curvature
$\tilde{F}$, which obeys $|\tilde{F}(p)|\le\ell\lambda_S(y)$ whenever the base row of $p$ lies in
$[y,H-1-y]$ with $y\ge b$. The young-only instance is the problem of part (c) of that model theorem
with cap $\tau$ on the plaquettes of the rows in $[0,W_y]\cup[H-W_y,H]$ and no cap elsewhere.
\begin{itemize}
\item[(E0)] There is a potential $A^S$ with the prescribed values on the held bonds, feasible for the
young-only instance, whose curvature $F^S$ obeys $|F^S|\le\sigma_S(b)$ on the plaquettes of base row
outside $[b,H-b]$ (the two boundary layers), $F^S=\tilde{F}$ on the plaquettes of base rows in
$[b,H-b]$, where $|\tilde{F}|\le\ell\lambda_S(b)\le\sigma_S(b)$, and $\|F^S-\tilde{F}\|^2\le\mathfrak{e}(b)$.
\end{itemize}
In (E1) and (E2), $P$ denotes the young-only instance, or any problem on the slab with the same
prescribed values on the held bonds, in which every capped plaquette carries the cap $\tau$, and for
which $A^S$ is feasible with $|F^S|\le\sigma_S(b)$ on every capped plaquette of $P$; $\hat{F}_P$ and
$\Lambda_P$ denote its minimising curvature and its multipliers.
\begin{itemize}
\item[(E1)] For every such problem $P$, $\|\hat{F}_P-\tilde{F}\|^2\le\mathfrak{e}(b)$.
\item[(E2)] For every such problem $P$,
$\sum_p\Lambda_P(p)\le\mathfrak{e}(b)/\bigl(2(\tau-\sigma_S(b))\bigr)$.
\end{itemize}
Let $y^{\flat}$ be the least $y\ge b$ with
$6\cdot 7^4\,\ell^2\lambda_S(y)^2\le y_E\tau^2/2$. Let $R_\beta$ be the least integer with
\begin{equation*}
R_\beta\ \ge\ \max\Bigl(\rho_E,\ \bigl(4\mathfrak{e}(b)/(y_E\tau^2)\bigr)^{1/4}\Bigr),
\end{equation*}
and assume $6R_\beta+3\le N$ and $2(y^{\flat}+3R_\beta)+1\le H$. Then parts (b) and (c) of the model
theorem, Theorem~\ref{8.16:thm-slab-inactivity}, hold with $3R_V$ replaced by $y^{\flat}+3R_\beta$ and
$R_V$ by $R_\beta$, in the following form: every
plaquette $p$ with base row in $[y^{\flat}+3R_\beta,\,H-1-y^{\flat}-3R_\beta]$ has
\begin{equation*}
|\hat{F}^{\tau}(p)|\le C_E\,y(R_\beta;0)^{1/2}\tau\le\tau/4
\quad\text{and}\quad \Lambda(p)=0,
\end{equation*}
where $y(R_\beta;0)$ is the normalised excess \eqref{8.15:eq-capped-data-c} of $P^{\tau}$ about the
base point of $p$;
and for every cap function $C$ with $C\ge\tau$ everywhere and $C=\tau$ on the plaquettes
of base row $<y^{\flat}+3R_\beta$ or $>H-1-y^{\flat}-3R_\beta$ one has $\hat{F}^{C}=\hat{F}^{\tau}$,
together with the consequences listed in part (c) for $W_y\ge y^{\flat}+3R_\beta$.
\end{theorem}

\begin{proof}
By (E0), $|F^S|\le\sigma_S(b)$ on the plaquettes of base row outside $[b,H-b]$, and on those of base
row in $[b,H-b]$ one has $F^S=\tilde{F}$ with
$|\tilde{F}|\le\ell\lambda_S(b)\le\sigma_S(b)$. Hence $|F^S(p)|\le\sigma_S(b)\le\tau/2\le\tau$ on every
plaquette. Since $A^S$ takes the prescribed values on the held bonds and all plaquettes carry the cap
$\tau$ in $P^{\tau}$, the potential $A^S$ is feasible for $P^{\tau}$, with $|F^S|\le\sigma_S(b)$ on every
capped plaquette. Thus $P^{\tau}$ is a problem $P$ as in (E1) and (E2), with $\hat{F}_P=\hat{F}^{\tau}$
and $\Lambda_P=\Lambda$, and (E1), (E2)
give $\|\hat{F}^{\tau}-\tilde{F}\|^2\le\mathfrak{e}(b)$ and
$\sum_p\Lambda(p)\le\mathfrak{e}(b)/\bigl(2(\tau-\sigma_S(b))\bigr)$; as $\sigma_S(b)\le\tau/2$ gives
$2(\tau-\sigma_S(b))\ge\tau$, also $\tau\sum_p\Lambda(p)\le\mathfrak{e}(b)$.

Let $c$ be a point of row $y\in[y^{\flat}+3R_\beta,\,H-1-y^{\flat}-3R_\beta]$. This interval is non-empty
because $2(y^{\flat}+3R_\beta)+1\le H$. The plaquettes based in $Q_{3R_\beta}(c)$ have base rows in
$[y^{\flat},H-1-y^{\flat}]$, so the bound on the reference curvature gives
$|\tilde{F}(p)|\le\ell\lambda_S(y^{\flat})$ on them. There are
at most $6(6R_\beta+1)^4$ of them. The hybrid excess of \eqref{8.15:eq-capped-data-c} with $F_0=0$ is at
most half the curvature energy of $\hat{F}^{\tau}$ in the cube plus $\tau\sum_p\Lambda(p)$; this is the
bound used in the proof of part (b) of the model theorem, Theorem~\ref{8.16:thm-slab-inactivity}. Using
$\tfrac12|\hat{F}^{\tau}|^2\le|\tilde{F}|^2+|\hat{F}^{\tau}-\tilde{F}|^2$ on each plaquette, we obtain
\begin{equation}\label{8.16:eq-active-sliver-1}
\begin{aligned}
\mathcal{Y}(3R_\beta;0)
&\le \sum_{p\in Q_{3R_\beta}(c)}|\tilde{F}(p)|^2+\|\hat{F}^{\tau}-\tilde{F}\|^2+\tau\sum_p\Lambda(p)\\
&\le 6(6R_\beta+1)^4\ell^2\lambda_S(y^{\flat})^2+2\mathfrak{e}(b)
\ \le\ y_E\tau^2R_\beta^4 .
\end{aligned}
\end{equation}
The last inequality is justified as follows. First, $R_\beta\ge\rho_E\ge1$, since $\rho_E$ is a natural
number and the natural numbers are taken positive; so $6R_\beta+1\le7R_\beta$.
With the definition of $y^{\flat}$ this yields
\begin{equation*}
6(6R_\beta+1)^4\ell^2\lambda_S(y^{\flat})^2\le 6\cdot7^4\ell^2\lambda_S(y^{\flat})^2R_\beta^4
\le \tfrac12 y_E\tau^2R_\beta^4 .
\end{equation*}
Second, $R_\beta^4\ge4\mathfrak{e}(b)/(y_E\tau^2)$ gives $2\mathfrak{e}(b)\le\tfrac12y_E\tau^2R_\beta^4$.

The bound \eqref{8.16:eq-active-sliver-1} is the bound $y(R_\beta;0)\le y_E$ on the normalised excess of
\eqref{8.15:eq-capped-data-c}. This is the estimate from which the proof of part (b) of the model
theorem, Theorem~\ref{8.16:thm-slab-inactivity}, proceeds, there at radius $R_V$. We carry out the
remaining steps at
radius $R_\beta$. Since $6R_\beta+3\le N$, the cube $Q_{3R_\beta+1}(c)$ injects into
$(\mathbb{Z}/N)^3\times\mathbb{Z}$ and is read in $\mathbb{Z}^4$. The plaquettes in $Q_{3R_\beta}(c)$ have
base rows in $[0,H-1]$, so they exist and carry the cap $\tau$. A bond with an endpoint in
$Q_{3R_\beta-1}(c)$ is normal, or tangential of a row in $[1,H-2]$, hence free. The model theorem,
Theorem~\ref{8.15:thm-epsilon-regularity}, therefore applies with $F_0:=0$ (so $|F_0|\le\tau/2$), $R:=R_\beta\ge\rho_E$ and
$\varpi\equiv0$. It gives $|\hat{F}^{\tau}(p)|\le C_E\,y(R_\beta;0)^{1/2}\tau\le\tau/4$ at every
plaquette $p$ based at $c$, and, as
in part (b) of the model theorem, Theorem~\ref{8.16:thm-slab-inactivity}, $\Lambda(p)=0$. This is
part (b) with $3R_V$ replaced by $y^{\flat}+3R_\beta$.

The proof of part (c) of the model theorem, Theorem~\ref{8.16:thm-slab-inactivity}, uses part (b) through
the facts that the multipliers of $P^{\tau}$ sit on plaquettes where $C=\tau$, and that $\Lambda=0$ and
$|\hat{F}^{\tau}|\le\tau/4$ on the plaquettes of the middle rows. By what has just been shown, these
facts hold with $3R_V$ replaced by $y^{\flat}+3R_\beta$, for every cap function $C$ as in the statement.
That proof therefore gives $\hat{F}^{C}=\hat{F}^{\tau}$,
and with it the consequences of part (c) for $W_y\ge y^{\flat}+3R_\beta$.
\end{proof}

\begin{remark}[Sizes of the sliver under the correspondence with the renormalisation step; the floor]
Consider the correspondence between the slab model and the renormalisation step fixed in this chapter.
In it, $\mathsf{T}$ is the ratio of successive caps, and $\mathsf{w}$, $\mathsf{n}$ and $D_W$ are
parameters of that correspondence. Let $c_{\flat}>0$ and $C_{\flat},C'_{\flat}<\infty$ be constants of
that correspondence, to
which no value is assigned, and let $\lesssim$ denote inequality up to such a constant. One has
$b,\,y^{\flat}\lesssim H e^{-c_{\flat}\mathsf{T}}$ and
$R_\beta\lesssim(N^3b)^{1/4}(\lambda_S/\mathsf{T})^{1/2}$, with $\lambda_S$ written without its
argument. Hence
$y^{\flat}+3R_\beta\le\mathsf{w}\,\mathsf{n}/L^2$ holds under the floor
\begin{equation}\label{8.16:eq-active-sliver-a}
c_{\flat}\,\mathsf{T}\ \ge\ C_{\flat}+C'_{\flat}\log\bigl(D_WL^2/\mathsf{w}\bigr).
\end{equation}
This floor is a lower bound only. It involves no cutoff, and no value is assigned to its constants. It is
of the same kind as the floor of the model theorem on inactivity at full band width. Under
\eqref{8.16:eq-active-sliver-a}, Theorem~\ref{8.16:thm-active-sliver} gives, with $n_f=1$, the
inactivity statement of that model theorem, and the exit bound at rate parameter $3/4$, in the setting
of the model theorem, Theorem~\ref{8.16:thm-slab-inactivity}. This is the statement of part (c) of the
model theorem on inactivity at full band width for arbitrary vectorial face data.
\end{remark}

\begin{remark}[Remarks on the slab theorem: clamps, sharpness, regime]\label{8.16:rem-slab-remarks}
Throughout, $\tau$, $\ell$, $\ell_{\times}$, $\ell_c$ and $R_V$ are as in the model theorem on exact inactivity away from the slab faces, Theorem~\ref{8.16:thm-slab-inactivity}, under its conditions \eqref{8.16:eq-slab-inactivity-a}.
\begin{itemize}
\item[(i)] \emph{Two clamps.} Let the slab have height $H_M$. Suppose that the rows $[0,W]$ carry a clamp of cap $\tau_s$, and that the rows $[H_M-W,H_M]$ carry its mirror image. Restrict the problem to the rows $[W,H_M-W]$, following the bracket of rows of the setting from which the slab problem is drawn. The restricted problem is a slab problem of height $H_M-2W$ in the sense of Theorem~\ref{8.16:thm-slab-inactivity}, with its face data read off the minimiser: its held faces are rows $W$ and $H_M-W$ of the minimiser itself. For this restricted problem one may take
\begin{equation}\label{8.16:eq-slab-remarks-1}
\ell:=\tau_s,\qquad \ell_{\times}^{\mathrm{sub}}\le \ell_{\times}+2W\tau_s ,
\end{equation}
where $\ell_{\times}^{\mathrm{sub}}$ denotes the transversal face mismatch of the restricted problem and $\ell_{\times}$ that of the original one. Nothing of the interior of the clamp enters: whatever the minimiser does there, whether or not its blocks are saturated in the sense of the saturation statement of that setting, it enters only through $\ell=\tau_s$. This reduction has no counterpart in Ba{\l}aban's construction, where the region playing the role of the clamp depends on the field at a greater depth than the depth at which its cap is imposed.
\item[(ii)] \emph{Sharpness is not claimed.} The true depth of the active zone is the thin sliver of Theorem~\ref{8.16:thm-active-sliver}. The model theorem, Theorem~\ref{8.16:thm-slab-inactivity}, gives up this depth in exchange for a proof that uses no kernel estimate at all. The exchange costs nothing: the outer ring consists of $N_r$ rows, each of fine height $w_n$, and any row $n<N_r$ among them serves.
\item[(iii)] \emph{Below the floor.} For $\tau$ of the order of $\ell_c$ or below, nothing is claimed; condition (V0) of \eqref{8.16:eq-slab-inactivity-a} requires $\ell_c\le\tau/2$. This is the regime of clause (iv) in the setting of the model theorem on inactivity at full band width.
\item[(iv)] \emph{Consistency with the model of inactivity at full band width.} The data on which the model theorem on inactivity at full band width is stated form an instance of the present setting. Clauses (c) and (d) of that model theorem are therefore recovered, up to constants, from Theorem~\ref{8.16:thm-slab-inactivity}, a model statement.
\end{itemize}
\end{remark}

\begin{proof}
We verify \eqref{8.16:eq-slab-remarks-1}.

\emph{The face curvature bound.} The clamp caps the tangential plaquettes of its last row $W$. The minimiser is feasible, so every tangential plaquette $p$ of row $W$ satisfies $|\hat{F}(p)|\le\tau_s$. The same holds for row $H_M-W$ by the mirror clamp. These rows are the held faces of the restricted problem, so their curvature is bounded by $\tau_s$, and one may put $\ell:=\tau_s$.

\emph{The mismatch bound.} Work in the axial gauge, in which every normal bond carries the value $0$. Fix a tangential bond position $b$, and write $\hat{A}_b(y)$ for the value of the potential on the copy of $b$ in row $y$. The normal plaquette spanned by the copies of $b$ in rows $y$ and $y+1$ has two normal sides, and both vanish. Its curvature is therefore $\pm\bigl(\hat{A}_b(y+1)-\hat{A}_b(y)\bigr)$.

For $0\le y\le W-1$ this plaquette lies in the clamp, whose cap is $\tau_s$. Telescoping gives
\begin{equation*}
\bigl|\hat{A}_b(W)-\hat{A}_b(0)\bigr|
\le\sum_{y=0}^{W-1}\bigl|\hat{A}_b(y+1)-\hat{A}_b(y)\bigr|
\le W\tau_s .
\end{equation*}
In the same way, using the mirror clamp on the rows $[H_M-W,H_M]$,
\begin{equation*}
\bigl|\hat{A}_b(H_M-W)-\hat{A}_b(H_M)\bigr|\le W\tau_s .
\end{equation*}
In the axial gauge the transversal face mismatch is read as the largest, over tangential bond positions $b$, of the difference of the potential on the two faces: for the original problem, of $|\hat{A}_b(H_M)-\hat{A}_b(0)|$, and for the restricted problem, of $|\hat{A}_b(H_M-W)-\hat{A}_b(W)|$. By the two displays above and the triangle inequality, for every $b$,
\begin{equation*}
\bigl|\hat{A}_b(H_M-W)-\hat{A}_b(W)\bigr|\le \ell_{\times}+2W\tau_s .
\end{equation*}
Hence
\begin{equation*}
\ell_{\times}^{\mathrm{sub}}\le\ell_{\times}+2W\tau_s .
\end{equation*}
\end{proof}

\begin{corollary}[Stopping of the front under a weak tangential bound]\label{8.16:cor-front-stopping}
We work in the upward exit-value model. It is the model of
Lemma~\ref{8.15:lem-band-optimality}, with the family of problems $(P_a)$, $a_0\le a\le\infty$,
where $P_\infty$ caps all levels $y\ge a_0$. In this model the tangential orientations carry no
cap. The cap height is the same number $\tau$ at every capped level. For the minimiser of
$P_\infty$ we write:
\begin{itemize}
\item $E_y$ for its capped component on the row $y$;
\item $B_y$ for its component on the uncapped orientations;
\item $\mu_y$ for the multiplier of the model lemma, Lemma~\ref{8.15:lem-band-optimality}.
\end{itemize}
In this corollary $N$ denotes the period of the cross-section of the model's slab, not the rank
of the gauge group. Let $\varkappa>0$ be a constant. Assume the
following two statements.
\begin{itemize}
\item[(a)] The weak tangential bound in the form
\begin{equation*}
|B_y|\le\varkappa\tau\qquad\text{on every row } y\ge a_0 .
\end{equation*}
This is the weak tangential bound \eqref{8.15:eq-regularity-remarks-a} on the uncapped
orientations, imposed on every row $y\ge a_0$, hence on every cube $Q_{3R}(c)$ with $c$ on a
row $y\ge a_0+3R$ (see the proof).
\item[(b)] The part of clause (f) of the exit-value model theorem used below: if the
multipliers $\mu_y$ of the minimiser of $P_\infty$ vanish on all rows $y\ge a'$, then the front
of the family $(P_a)$ has stopped at $a'$.
\end{itemize}
Put
\begin{equation}\label{8.16:eq-front-stopping-1}
\mathfrak{E}_{\uparrow}:=\sum_{y\ge a_0}\Big[\frac12\|E_y\|^2+\frac12\|B_y\|^2
+\tau\|\mu_y\|_1\Big].
\end{equation}
Here $\|\cdot\|$ and $\|\cdot\|_1$ are the $\ell^2$ and $\ell^1$ norms over the row $y$.
Let $\rho_E(\varkappa)$ and $y_E(\varkappa)$ be the constants of the model theorem,
Theorem~\ref{8.15:thm-epsilon-regularity}, in the variant of
Remark~\ref{8.15:rem-regularity-remarks} for uncapped orientations; in that variant the
constants depend on $\varkappa$. Suppose that
\begin{equation}\label{8.16:eq-front-stopping-2}
R\ge\rho_E(\varkappa),\qquad 6R+3\le N,\qquad
y_E(\varkappa)\,\tau^2R^4\ge\mathfrak{E}_{\uparrow}.
\end{equation}
Then $\mu_y=0$ for every $y\ge a_0+3R$, and the front of the family $(P_a)$ has stopped at
$a_0+3R$.
\end{corollary}

\begin{proof}
The multiplier $\mu_y$ is defined by the model statement
Lemma~\ref{8.15:lem-band-optimality}, the optimality system inside the capped band. That
lemma defines $\mu_y=\mathcal{D}_y-E_y$, with $\mathcal{D}_y$ its driving field.
It also states that $\mu_y$ lies in the normal cone of the box at $E_y$ on the capped levels.
Thus the minimiser of $P_\infty$ satisfies the optimality system on which the model theorem
of $\varepsilon$-regularity, Theorem~\ref{8.15:thm-epsilon-regularity}, operates.

In the upward exit-value model the tangential orientations carry no cap. Hence the cap
condition of Theorem~\ref{8.15:thm-epsilon-regularity} is not available on them.
Remark~\ref{8.15:rem-regularity-remarks} describes two separate variants of that theorem:
\begin{itemize}
\item with caps $C(p)\in[\tau,\varkappa\tau]$ on the cube $Q_{3R}$, the theorem holds with
constants depending on $\varkappa$;
\item on orientations that carry no cap, the cap condition may be replaced by the hypothesis
\eqref{8.15:eq-regularity-remarks-a}.
\end{itemize}
Here the capped orientations carry the cap $\tau$ at every level, and the uncapped orientations
carry the hypothesis \eqref{8.15:eq-regularity-remarks-a}. In the first variant the constant
$\varkappa$ enters the constants of the theorem through the bound
\begin{equation*}
\sup|D'|\le(\varkappa+1)\tau
\end{equation*}
on $Q_{3R}$, where $D'=dA-F_0$ and $A$ is a potential of the curvature on the cube; the bound
by $\tau$ on the capped orientations together with
the bound \eqref{8.15:eq-regularity-remarks-a} by $\varkappa\tau$ on the uncapped ones plays the
same role here, so the constants depend on $\varkappa$. The first variant presupposes
$\varkappa\ge1$; since (a) with $\varkappa<1$ implies (a) with $\varkappa=1$, we write
$\rho_E(\varkappa)$ and $y_E(\varkappa)$ for the constants of the theorem at
$\max\{\varkappa,1\}$. These are the numbers that appear in \eqref{8.16:eq-front-stopping-2}.

Now let $c$ be a point on a row $y\ge a_0+3R$. The cube $Q_{3R}(c)$ is taken in the
$|\cdot|_\infty$ distance. Its rows $y'$ therefore satisfy
\begin{equation*}
y'\ge y-3R\ge a_0 .
\end{equation*}
So hypothesis (a) gives $|B_{y'}|\le\varkappa\tau$ on all of $Q_{3R}(c)$. This is the bound
\eqref{8.15:eq-regularity-remarks-a} on the uncapped orientations, with the constant
$\varkappa$.

We apply the model theorem, Theorem~\ref{8.15:thm-epsilon-regularity}, in the variant just
described, at the point $c$, with radius $R$ and reference value $F_0=0$. Its hypotheses are
the lower bound on the radius, the bound $|F_0|\le\tau/2$, the smallness (h1) of the normalised
excess $y(R;F_0)$, and the condition (h2) on the defect function $\varpi$. They enter as
follows.
\begin{itemize}
\item The lower bound $R\ge\rho_E(\varkappa)$ is the first condition of
\eqref{8.16:eq-front-stopping-2}.
\item The bound $|F_0|\le\tau/2$ holds because $F_0=0$.
\item Hypothesis (h1). The normalised excess $y(R;0)$ is the hybrid excess
$\mathcal{Y}(R;0)$ about $c$ divided by $\tau^2R^4$; the hybrid excess of
Theorem~\ref{8.15:thm-epsilon-regularity} with $F_0=0$ is the sum of the same three
non-negative densities as in \eqref{8.16:eq-front-stopping-1}, over a cube about $c$. Every
term of
\eqref{8.16:eq-front-stopping-1} is non-negative, and the cubes about $c$ on which the theorem
reads its quantities lie in $Q_{3R}(c)$, hence in rows $y'\ge a_0$, over which
\eqref{8.16:eq-front-stopping-1} sums. Therefore $\mathcal{Y}(R;0)\le\mathfrak{E}_{\uparrow}$,
and by the third condition of \eqref{8.16:eq-front-stopping-2}
\begin{equation*}
y(R;0)\le\frac{\mathfrak{E}_{\uparrow}}{\tau^2R^4}\le y_E(\varkappa),
\end{equation*}
which is (h1).
\item The condition $6R+3\le N$ ensures that the cubes used by the theorem about $c$ do not
wrap around the periodic cross-section of side $N$.
\item Hypothesis (h2) is a property of the defect function $\varpi$ of the minimiser of
$P_\infty$. In the upward exit-value model the curvature of the minimiser is the exterior
derivative of its lateral potential; along the rows this is the relation (Far) of the model
lemma, Lemma~\ref{8.15:lem-band-optimality}, namely $\beta_{y+1}-\beta_y=PE_y$. Hence the
defect $\mathfrak{n}=\hat{F}-dA$ vanishes, the defect function $\varpi$ is identically zero,
and (h2) holds with $\varpi_R=0\le c_E$; \eqref{8.16:eq-front-stopping-2} places no further
condition on it.
\end{itemize}
The conclusion of the theorem, with $F_0=0$, is $|\hat{F}(p)|\le\tau/4$ at every plaquette $p$
based at $c$. In particular, on the capped orientations the component of $E_y$ at $c$ has
modulus at most $\tau/4<\tau$, so it lies in the interior of the box. The normal cone of the
box at an interior point is $\{0\}$, and by the complementary-slackness relation of the model
lemma, Lemma~\ref{8.15:lem-band-optimality} (the multiplier lies in that normal cone), the
multiplier vanishes at $c$. Since $c$ was an arbitrary point on a row $y\ge a_0+3R$, this
yields $\mu_y=0$ for $y\ge a_0+3R$.

It remains to apply hypothesis (b). By clause (f) of the exit-value model theorem, the
vanishing of the multipliers $\mu_y$ on all rows $y\ge a_0+3R$ means that the front of the
family $(P_a)$ has stopped at $a_0+3R$.
\end{proof}

In the upward exit-value model the bound of hypothesis (a) of
Corollary~\ref{8.16:cor-front-stopping} is not proved. There the field $B$ solves
\begin{equation*}
\partial_y^2B=\operatorname{curl}\gamma\operatorname{curl}^{\ast}B
\end{equation*}
with a coefficient field $\gamma$, measurable with values in $[0,1]$ (a local letter, unrelated
to the constants $\gamma$ of the glossary and to the commutator form $\gamma_{\mu\nu}$). This is
a system, to which no maximum principle
applies.

In the young band of the setting in which every orientation carries a cap, and in
Ba{\l}aban's construction, where every orientation is capped as well, the bound is part of the
constraints.

The exit-value model theorem gives three things that Theorem~\ref{8.15:thm-epsilon-regularity}
does not: the exact concavity, in the capping level $a$, of the value $m(a)$ of the problem
$P_a$ established in the exit-value model theorem; the two-sided bound of its clause (c); and
the property that a front once stopped stays stopped, without any smallness condition.

\begin{definition}[Plaquette-capped $\mathrm{SU}(2)$ fields and the multiplier rule]
\label{8.16:def-capped-gauge-fields}
Let $G=\mathrm{SU}(2)$. Let $\mathfrak{g}$ be its Lie algebra with an $\operatorname{Ad}$-invariant inner product
$\langle\cdot,\cdot\rangle$ and norm $|\cdot|$. Fix $c\in\mathbb{Z}^4$ and an integer $R\ge1$. We use the cube
$Q_r(c)$, the plaquettes $p=(z,\mu\nu)$ with base site $z$, the ordered bonds
\[
\partial p=(b_1,b_2,b_3,b_4)=\bigl((z,\mu),\,(z+e_\mu,\nu),\,(z+e_\nu,\mu),\,(z,\nu)\bigr),
\]
and the spaces $\mathfrak{X}_r$, all as in Definition~\ref{8.15:def-capped-data}, with $\mathfrak{g}$ in place of the
colour space $\mathsf{V}$. For a bond $b=(x,\mu)$ write $b_-=x$ and $b_+=x+e_\mu$.
The definitions below make sense for any compact subgroup $G$ of a unitary group; only the remarks on class
functions and on caps use $G=\mathrm{SU}(2)$.

\emph{Fields and plaquette logarithms.} A field is a map $U$ from the bonds of the plaquettes in $Q_{3R}(c)$ to $G$.
For $p=(z,\mu\nu)$ in $Q_{3R}(c)$ the holonomy based at $z$ is
\[
U_p:=U_{b_1}U_{b_2}U_{b_3}^{-1}U_{b_4}^{-1},
\]
and $\Phi(p):=\log U_p\in\mathfrak{g}$ is its principal logarithm, defined when $U_p\ne-1$, as is the case under
the caps below.

\emph{Action.} The action is $\sum_p\mathfrak{a}(|\Phi(p)|)$, the sum running over the plaquettes in $Q_{3R}(c)$.
Here $\mathfrak{a}$ is a function on $[0,\infty)$ with
\[
\mathfrak{a}'(t)=t\,a(t),\qquad |a(t)-1|\le C_a t^2 ,
\]
with a constant $C_a\ge0$.
We put $a_p:=a(|\Phi(p)|)$. Wilson's action falls under this description, since on $\mathrm{SU}(2)$ every class
function of $U_p$ is a function of $|\Phi(p)|$.

\emph{Caps.} Fix $\tau$ with $0<\tau\le\tau_G$, where $\tau_G>0$ is the constant of the defect lemma for
$\mathrm{SU}(2)$ in this chapter. The field is required to satisfy
\begin{equation}\label{8.16:eq-capped-gauge-fields-1}
|\Phi(p)|\le\tau\qquad\text{for every plaquette $p$ in $Q_{3R}(c)$, of every orientation.}
\end{equation}
On $\mathrm{SU}(2)$ this is the same as a cap on the operator norm $\|U_p-1\|$, because $\|U_p-1\|$ is a monotone
function of $|\Phi(p)|$.

\emph{Covariant variation.} Let $\xi$ be a $\mathfrak{g}$-valued $1$-form, with $\xi_b$ attached to the initial
site $b_-$ of $b$. Consider the variation $U_b(t)=e^{t\xi_b}U_b$. Then
\begin{equation}\label{8.16:eq-capped-gauge-fields-2}
\frac{d}{dt}\,U_p(t)U_p^{-1}\Big|_{t=0}=(D_U\xi)(p),
\end{equation}
where
\[
(D_U\xi)(p):=\xi_{b_1}+\operatorname{Ad}(U_{b_1})\xi_{b_2}
-\operatorname{Ad}(U_{b_1}U_{b_2}U_{b_3}^{-1})\xi_{b_3}-\operatorname{Ad}(U_p)\xi_{b_4}.
\]
Let $Y\in\mathfrak{g}$, let $\Phi\ne0$ be a logarithm as above, and let $f(u)=\tilde f(|\log u|)$ for $u$ in a
neighbourhood of $e^{\Phi}$, with $\tilde f$ differentiable at $|\Phi|$. Then
\begin{equation}\label{8.16:eq-capped-gauge-fields-3}
\frac{d}{dt}\,f(e^{tY}e^{\Phi})\Big|_{t=0}=\tilde f'(|\Phi|)\,\frac{\langle\Phi,Y\rangle}{|\Phi|}.
\end{equation}

\emph{The multiplier rule.} This is a hypothesis on $U$. It requires that there be numbers $\Lambda(p)\ge0$, with
$\Lambda(p)=0$ unless $|\Phi(p)|=\tau$, such that
\begin{equation}\label{8.16:eq-capped-gauge-fields-a}
\sum_{p\text{ in }Q_{3R}(c)}\Bigl\langle a_p\Phi(p)+\Lambda(p)\,\frac{\Phi(p)}{\tau},\,(D_U\xi)(p)\Bigr\rangle=0
\qquad\text{for every }\xi\in\mathfrak{X}_{3R}.
\end{equation}
It holds at every local minimiser of the action under the caps \eqref{8.16:eq-capped-gauge-fields-1} at which a
constraint qualification holds. An example of such a qualification is the existence of a variation that decreases
all active $|\Phi(p)|$.

\emph{Gauge covariance.} A gauge transformation $\lambda$ from sites to $G$ acts by
\[
U_b\mapsto\lambda(b_-)U_b\lambda(b_+)^{-1},\qquad \xi_b\mapsto\operatorname{Ad}(\lambda(b_-))\xi_b .
\]
Under this action $\Phi(p)\mapsto\operatorname{Ad}(\lambda(z))\Phi(p)$. The quantities $|\Phi(p)|$ and $\Lambda(p)$
are invariant: the numbers $\Lambda(p)$ serve in \eqref{8.16:eq-capped-gauge-fields-a} for $U$ if and only if they
serve for the transformed field.
\end{definition}

\begin{proof}
\emph{Fields and plaquettes.} Every bond carrying a value of some $\xi\in\mathfrak{X}_{3R}$ has both endpoints in
$Q_{3R}(c)$. Its initial site is therefore the base site of a plaquette in $Q_{3R}(c)$ that contains it, so $U_b$ is
defined on it.

\emph{Caps and class functions.} Since $\mathfrak{g}$ is simple, an $\operatorname{Ad}$-invariant inner product on it
is a positive multiple of $-\operatorname{tr}XY$. Now $U_p=e^{\Phi}$ is conjugate to the diagonal matrix with
entries $e^{\pm i\theta}$, $0\le\theta\le\pi$, with $\theta<\pi$ since $\Phi(p)$ is defined. Hence $|\Phi(p)|$ is a
fixed positive multiple of $\theta$. Moreover
\[
\|U_p-1\|=|e^{i\theta}-1|=2\sin(\theta/2),
\]
which increases in $\theta$. A conjugacy class of $\mathrm{SU}(2)$ is determined by $\theta$. So every class
function, in particular $1-\operatorname{Re}\operatorname{tr}U_p$, is a function of $|\Phi(p)|$.

\emph{Proof of \eqref{8.16:eq-capped-gauge-fields-2}.} Write $U_i=U_{b_i}$ and $\xi_i=\xi_{b_i}$. Then
\[
U_p(t)=e^{t\xi_1}U_1e^{t\xi_2}U_2U_3^{-1}e^{-t\xi_3}U_4^{-1}e^{-t\xi_4}.
\]
Differentiating at $t=0$ gives
\[
\xi_1U_p+U_1\xi_2U_2U_3^{-1}U_4^{-1}-U_1U_2U_3^{-1}\xi_3U_4^{-1}-U_p\xi_4 .
\]
Multiply on the right by $U_p^{-1}=U_4U_3U_2^{-1}U_1^{-1}$. The four terms become
\[
\xi_1,\qquad U_1\xi_2U_1^{-1},\qquad
-(U_1U_2U_3^{-1})\,\xi_3\,(U_1U_2U_3^{-1})^{-1},\qquad -U_p\xi_4U_p^{-1} .
\]
This is $(D_U\xi)(p)$.

\emph{Proof of \eqref{8.16:eq-capped-gauge-fields-3}.} The logarithm is analytic near $e^{\Phi}$. By the formula for
its differential there is a power series $\psi$ with $\psi(0)=1$ such that
\[
\log(e^{tY}e^{\Phi})=\Phi+t\,\psi(\operatorname{ad}_\Phi)Y+O(t^2).
\]
Differentiate $\langle\operatorname{Ad}(e^{s\Phi})X,\operatorname{Ad}(e^{s\Phi})Z\rangle=\langle X,Z\rangle$ at
$s=0$. This shows that $\operatorname{ad}_\Phi$ is skew. Together with $\operatorname{ad}_\Phi\Phi=0$ it gives, for
$n\ge1$,
\[
\langle\Phi,\operatorname{ad}_\Phi^nY\rangle=(-1)^n\langle\operatorname{ad}_\Phi^n\Phi,Y\rangle=0,
\]
and hence $\langle\Phi,\psi(\operatorname{ad}_\Phi)Y\rangle=\langle\Phi,Y\rangle$. Therefore
\[
\frac{d}{dt}\bigl|\log(e^{tY}e^{\Phi})\bigr|\Big|_{t=0}=\frac{\langle\Phi,Y\rangle}{|\Phi|},
\]
and the chain rule gives \eqref{8.16:eq-capped-gauge-fields-3}. The expansion is unchanged if $tY$ is replaced by
$tY+O(t^2)$.

\emph{The multiplier rule.} By \eqref{8.16:eq-capped-gauge-fields-2} and \eqref{8.16:eq-capped-gauge-fields-3},
with $\mathfrak{a}'(t)/t=a(t)$, the derivative along $\xi$ of $\mathfrak{a}(|\Phi(p)|)$ at a plaquette with
$\Phi(p)\ne0$ is $a_p\langle\Phi(p),(D_U\xi)(p)\rangle$, and at an active plaquette the derivative of $|\Phi(p)|$
is $\langle\Phi(p),(D_U\xi)(p)\rangle/\tau$. Where $\Phi(p)=0$ the summand $a_p\langle\Phi(p),(D_U\xi)(p)\rangle$
vanishes, and so does the derivative of $\mathfrak{a}(|\Phi(p)|)$, because $\mathfrak{a}'(0)=0$ and
$|\log U_p(t)|=O(t)$. Thus \eqref{8.16:eq-capped-gauge-fields-a} is the Karush--Kuhn--Tucker condition for the
action under the caps \eqref{8.16:eq-capped-gauge-fields-1}.

\emph{Gauge covariance.} Write $U^\lambda$ and $\xi^\lambda$ for the transformed field and $1$-form. Telescoping in
the product defining $U_p$ gives $U_p\mapsto\lambda(z)U_p\lambda(z)^{-1}$. The principal logarithm commutes with
conjugation, so $\Phi(p)\mapsto\operatorname{Ad}(\lambda(z))\Phi(p)$. Further,
\[
\lambda(b_-)e^{t\xi_b}U_b\lambda(b_+)^{-1}=e^{t\operatorname{Ad}(\lambda(b_-))\xi_b}\,U_b^\lambda .
\]
Thus the transformed variation of $U$ is the variation of $U^\lambda$ along $\xi^\lambda$. Together with
\eqref{8.16:eq-capped-gauge-fields-2} this gives
\[
(D_{U^\lambda}\xi^\lambda)(p)=\operatorname{Ad}(\lambda(z))(D_U\xi)(p).
\]
By $\operatorname{Ad}$-invariance, $|\Phi(p)|$ is unchanged and each summand of
\eqref{8.16:eq-capped-gauge-fields-a} is unchanged. Since $\xi\mapsto\xi^\lambda$ is a bijection of
$\mathfrak{X}_{3R}$, the same $\Lambda$ serves for $U$ and for $U^\lambda$.
\end{proof}

\begin{lemma}[$\mathrm{SU}(2)$ fields in axial gauge satisfy the defect bounds]\label{8.16:lem-axial-defects}
Let $U$, $U_p$, $\Phi(p)$, $a(t)$, $a_p$, $C_a$, $\tau$ and $D_U\xi$ be as in
Definition~\ref{8.16:def-capped-gauge-fields}, on the cube $Q_{3R}(c)$, with $G=\mathrm{SU}(2)$ and caps
$|\Phi(p)|\le\tau$ on every plaquette, and assume the multiplier rule
\eqref{8.16:eq-capped-gauge-fields-a}: there are $\Lambda(p)\ge0$, with $\Lambda(p)=0$ unless
$|\Phi(p)|=\tau$, such that
\begin{equation*}
\sum_p\bigl\langle a_p\Phi(p)+\Lambda(p)\Phi(p)/\tau,(D_U\xi)(p)\bigr\rangle=0
\end{equation*}
for every
$\xi\in\mathfrak{X}_{3R}$. There are constants $C_G$, $\tau_G$, $c_G$, depending only on $G$
and the normalisations, such that the following holds. If $\tau\le\tau_G$ and $R\tau\le c_G$, then in
the complete axial gauge about $c$, with $A_b:=\log U_b$, $\hat{F}:=\Phi$ and
$\mathsf{V}:=\mathfrak{g}$, the data $(A,\hat{F},\Lambda)$ satisfy \eqref{8.15:eq-capped-data-a}
and \eqref{8.15:eq-capped-data-b} of Definition~\ref{8.15:def-capped-data} for all $1\le r\le R$,
with
\begin{equation}\label{8.16:eq-axial-defects-1}
\varpi(r):=C_G\,r^2\tau .
\end{equation}
\end{lemma}

\begin{proof}
By the gauge covariance stated in Definition~\ref{8.16:def-capped-gauge-fields}, $|\Phi(p)|$ and
$\Lambda(p)$ are unchanged by the passage to the complete axial gauge about $c$, and the multiplier
rule holds there as well. Hence $|\hat{F}(p)|=|\Phi(p)|\le\tau$ on every plaquette, which is the
cap condition in \eqref{8.15:eq-capped-data-a}, and $\Lambda(p)=0$ unless $|\hat{F}(p)|=\tau$,
which is the complementary slackness condition there. We write
$\mu(p):=\Lambda(p)\hat{F}(p)/\tau$. The cap gives $|\mu(p)|\le\Lambda(p)$.

\emph{The potential.} Throughout, $\|\cdot\|$ denotes the operator norm, of $2\times2$ matrices
acting on $\mathbb{C}^2$ and, below, of linear maps of $\mathfrak{g}$.
Take the nonabelian Stokes bound in the complete axial gauge about $c$. There
the bonds of the axial tree carry $U_b=1$, and every bond $b$ satisfies
$\|U_b-1\|\le|b_--c|_1\sup_p\|U_p-1\|$, where $b_-$ is the initial site of $b$. The cap
$|\Phi(p)|\le\tau$ bounds $\|U_p-1\|$ by a
constant times $\tau$. The initial site of a bond of a plaquette in $Q_{3r}(c)$ lies at
$\ell^1$-distance from $c$ at most a constant times $r$. Since $r\le R$ and $R\tau\le c_G$ with
$c_G$ small, $\|U_b-1\|$ is small for these bonds, so the principal logarithm $A_b=\log U_b$ is
defined and $|A_b|$ is at most a constant times $\|U_b-1\|$. It follows that
\begin{equation}\label{8.16:eq-axial-defects-2}
m_A(r):=\sup\{|A_b|: b \text{ a bond of a plaquette in } Q_{3r}(c)\}\le C_1 r\tau ,
\end{equation}
with $C_1$ depending only on $G$ and the normalisations. The bound is proportional to $r$, so one
gauge serves all the concentric cubes $Q_{3r}(c)$, $1\le r\le R$.

\emph{The first defect bound.} By the Baker--Campbell--Hausdorff formula there are $C_B$ and
$m_0$ such that, whenever $\max_i|X_i|\le m_0$,
\begin{equation*}
\bigl|\log\bigl(e^{X_1}e^{X_2}e^{-X_3}e^{-X_4}\bigr)-(X_1+X_2-X_3-X_4)\bigr|
\le C_B\max_i|X_i|^2 .
\end{equation*}
Let $p$ be a plaquette in $Q_{3r}(c)$ with $\partial p=(b_1,b_2,b_3,b_4)$, and apply this bound
with $X_i=A_{b_i}$. Then $U_p=e^{X_1}e^{X_2}e^{-X_3}e^{-X_4}$ and
$(dA)(p)=X_1+X_2-X_3-X_4$. By \eqref{8.16:eq-axial-defects-2},
$\max_i|X_i|\le C_1R\tau\le C_1c_G$, and this is at most $m_0$ once $c_G\le m_0/C_1$. Therefore
\begin{equation*}
|\Phi(p)-(dA)(p)|\le C_BC_1^2r^2\tau^2 ,
\end{equation*}
which is the first defect bound of \eqref{8.15:eq-capped-data-b} as soon as $C_G\ge C_BC_1^2$.

\emph{The second defect bound.} Let
$\xi\in\mathfrak{X}_{3r}$; then $\xi\in\mathfrak{X}_{3R}$, since $r\le R$. Since
$\mu(p)=\Lambda(p)\Phi(p)/\tau$, subtracting the left side of the multiplier rule, which vanishes,
from $\sum_p\langle\Phi(p)+\mu(p),(D_U\xi)(p)\rangle$ gives
\begin{equation*}
\sum_p\bigl\langle\Phi(p)+\mu(p),(D_U\xi)(p)\bigr\rangle
=\sum_p\bigl\langle(1-a_p)\Phi(p),(D_U\xi)(p)\bigr\rangle .
\end{equation*}
The pairing $\langle\hat{F}+\mu,d\xi\rangle$ in the second defect bound of
\eqref{8.15:eq-capped-data-b} is the sum $\sum_p\langle\hat{F}(p)+\mu(p),(d\xi)(p)\rangle$.
Since $\hat{F}=\Phi$,
\begin{align}\label{8.16:eq-axial-defects-3}
\sum_p\bigl\langle\hat{F}(p)+\mu(p),(d\xi)(p)\bigr\rangle
={}&\sum_p\bigl\langle\hat{F}(p)+\mu(p),(d\xi-D_U\xi)(p)\bigr\rangle\\
&+\sum_p\bigl\langle(1-a_p)\Phi(p),(D_U\xi)(p)\bigr\rangle .\notag
\end{align}
Only plaquettes carrying a bond on which $\xi\neq0$ contribute. Such a bond has an endpoint in
$Q_{3r-1}(c)$, so the plaquette lies in $Q_{3r}(c)$.

For the first term, $(d\xi)(p)=\xi_{b_1}+\xi_{b_2}-\xi_{b_3}-\xi_{b_4}$, and the formula for
$D_U\xi$ in Definition~\ref{8.16:def-capped-gauge-fields} gives
\begin{align*}
(d\xi-D_U\xi)(p)={}&\bigl(1-\operatorname{Ad}(U_{b_1})\bigr)\xi_{b_2}
-\bigl(1-\operatorname{Ad}(U_{b_1}U_{b_2}U_{b_3}^{-1})\bigr)\xi_{b_3}\\
&-\bigl(1-\operatorname{Ad}(U_p)\bigr)\xi_{b_4} .
\end{align*}
Each $\operatorname{Ad}(g)$ is an
isometry, and
\begin{equation*}
\operatorname{Ad}(gh)-1=\operatorname{Ad}(g)\bigl(\operatorname{Ad}(h)-1\bigr)
+\bigl(\operatorname{Ad}(g)-1\bigr).
\end{equation*}
Moreover $\|\operatorname{Ad}(e^X)-1\|$ is bounded by a constant times $|X|$ for $|X|$ bounded.
Apply these facts with $|A_{b_i}|\le m_A(r)$ and $|\Phi(p)|\le\tau$. This yields
\begin{align*}
|(d\xi-D_U\xi)(p)|
&\le\bigl[\|\operatorname{Ad}(U_{b_1})-1\|+\|\operatorname{Ad}(U_{b_1}U_{b_2}U_{b_3}^{-1})-1\|
+\|\operatorname{Ad}(U_p)-1\|\bigr]\,|\xi|_{\partial p}\\
&\le C_2\bigl(m_A(r)+\tau\bigr)|\xi|_{\partial p} .
\end{align*}
Together with $|\mu(p)|\le\Lambda(p)$, the first term of \eqref{8.16:eq-axial-defects-3} is at most
\begin{equation*}
C_2(C_1r\tau+\tau)\sum_p\bigl(|\hat{F}(p)|+\Lambda(p)\bigr)|\xi|_{\partial p} .
\end{equation*}

For the second term we use three facts: $|a(t)-1|\le C_at^2$ with $t=|\Phi(p)|\le\tau$; the bound
$|(D_U\xi)(p)|\le4|\xi|_{\partial p}$, since $\operatorname{Ad}$ is isometric; and $\Phi=\hat{F}$.
They give
\begin{equation*}
|(1-a_p)\Phi(p)|\,|(D_U\xi)(p)|\le4C_a\tau^2|\Phi(p)|\,|\xi|_{\partial p}
=4C_a\tau^2|\hat{F}(p)|\,|\xi|_{\partial p}.
\end{equation*}

Now add the two estimates, using $r\ge1$ and $\tau\le\tau_G$, where $\tau_G$ is fixed so that the
preceding bounds on the logarithm and on $\operatorname{Ad}$ hold whenever $\tau\le\tau_G$, and put
$C_3:=C_2(C_1+1)+4C_a\tau_G$. We obtain, with the sum over $p$ in $Q_{3r}(c)$,
\begin{align*}
\Bigl|\sum_p\bigl\langle\hat{F}(p)+\mu(p),(d\xi)(p)\bigr\rangle\Bigr|
&\le C_3r\tau\sum_p\bigl(|\hat{F}(p)|+\Lambda(p)\bigr)|\xi|_{\partial p}\\
&=\frac{C_3r^2\tau}{r}\sum_p\bigl(|\hat{F}(p)|+\Lambda(p)\bigr)|\xi|_{\partial p} .
\end{align*}
This is the second defect bound of \eqref{8.15:eq-capped-data-b} with $\varpi$ as in
\eqref{8.16:eq-axial-defects-1}, provided $C_G\ge C_3$. Taking
$C_G:=\max\{C_BC_1^2,C_3\}$ completes the proof.
\end{proof}

\begin{theorem}[$\varepsilon$-regularity for plaquette-capped $\mathrm{SU}(2)$ fields]
\label{8.16:thm-gauge-regularity}
Let $U$ be a plaquette-capped field on the bonds of the plaquettes in $Q_{3R}(c)$, in the setting
of Definition~\ref{8.16:def-capped-gauge-fields}, with cap height $\tau$. Assume that $U$ satisfies the
multiplier rule \eqref{8.16:eq-capped-gauge-fields-a} with multipliers $\Lambda(p)\ge 0$. Let
$\rho_E,y_E,c_E$ be the constants of the model theorem, Theorem~\ref{8.15:thm-epsilon-regularity}, and
let $C_G,\tau_G,c_G$ be the constants of Lemma~\ref{8.16:lem-axial-defects}, with $C_G$ taken so that
$C_Gc_G\ge c_E$. Let $R\ge\rho_E$,
$\tau\le\tau_G$,
\begin{equation}\label{8.16:eq-gauge-regularity-a}
C_G R^2\tau\le c_E ,
\end{equation}
and, with the sum over the plaquettes $p$ in $Q_{3R}(c)$,
\begin{equation}\label{8.16:eq-gauge-regularity-1}
\sum_{p\in Q_{3R}(c)}\Bigl[\tfrac12|\Phi(p)|^2+\tau\Lambda(p)\Bigr]\le y_E\,\tau^2R^4 .
\end{equation}
Then the six plaquettes $p$ based at $c$ obey $|\Phi(p)|\le\tau/4$ and carry no multiplier,
$\Lambda(p)=0$. If $U$ is given on the bonds of the plaquettes in $Q_{4R}(c)$ and these hypotheses
hold about every $c'\in Q_R(c)$ (note that $Q_{3R}(c')\subset Q_{4R}(c)$), then $U$ is exact lattice
Yang--Mills at the bonds of $Q_{R-1}(c)$, that is, the action $\sum_p\mathfrak{a}(|\Phi(p)|)$ is
stationary under every variation $U_b\mapsto e^{t\xi_b}U_b$ of the bonds $b$ with an endpoint in
$Q_{R-1}(c)$.
\end{theorem}

\begin{proof}
By Definition~\ref{8.16:def-capped-gauge-fields}, a gauge transformation $g$ acts by
$\xi_b\mapsto\operatorname{Ad}(g(b_-))\xi_b$ and $\Phi(p)\mapsto\operatorname{Ad}(g(z))\Phi(p)$ for
$p$ based at $z$, and the multiplier rule is gauge covariant. Since the inner product on
$\mathfrak{g}$ is $\operatorname{Ad}$-invariant, $|\Phi(p)|$ and $\Lambda(p)$ are invariant. Hence the
hypotheses and the conclusion are gauge invariant, and it suffices to argue in the complete axial
gauge about $c$.

In that gauge we apply Lemma~\ref{8.16:lem-axial-defects}. The normalisation $C_Gc_G\ge c_E$ made in
the statement is permitted: the defect conditions that lemma asserts are upper bounds in which
$\varpi(r)=C_Gr^2\tau$ is the coefficient, so enlarging $C_G$ only weakens its conclusion, and the
lemma holds with the enlarged constant. Its hypotheses are $\tau\le\tau_G$ and
$R\tau\le c_G$. The first holds by assumption; for the second, \eqref{8.16:eq-gauge-regularity-a},
$R\ge\rho_E\ge1$ (the defect function is defined for radii $r\ge1$, and $\rho_E\in\mathbb{N}$ is the
lower end of the range of radii in the model theorem) and $C_Gc_G\ge c_E$ give
\begin{equation*}
R\tau=\frac{R^2\tau}{R}\le\frac{c_E}{C_G R}\le\frac{c_E}{C_G}\le c_G.
\end{equation*}
With $A_b:=\log U_b$ and $\hat{F}:=\Phi$, the lemma gives that the data
$(A,\hat{F},\Lambda)$ satisfy the defect conditions imposed on the capped lattice data in the model
setting of Theorem~\ref{8.15:thm-epsilon-regularity}, for all $1\le r\le R$, with
$\varpi(r)=C_G r^2\tau$.

We verify the hypotheses of the model theorem, Theorem~\ref{8.15:thm-epsilon-regularity}, with
$F_0=0$. First, $|F_0|=0\le\tau/2$. Second, with $\hat{F}=\Phi$ and $F_0=0$, the left side of
\eqref{8.16:eq-gauge-regularity-1} is the hybrid excess $\mathcal{Y}(R;0)$ of
\eqref{8.15:eq-capped-data-c}. Thus \eqref{8.16:eq-gauge-regularity-1} says $y(R;0)\le y_E$, which is
(h1) of \eqref{8.15:eq-epsilon-regularity-a}. Third, put $\varpi_R:=C_G R^2\tau$. For
$\rho_E\le r\le R$ we have $r/R\le 1$, so
\begin{equation*}
\varpi(r)=\varpi_R\,(r/R)^2\le\varpi_R\,(r/R),
\end{equation*}
and $\varpi_R\le c_E$ by \eqref{8.16:eq-gauge-regularity-a}. This is (h2) of
\eqref{8.15:eq-epsilon-regularity-a}. Finally, $R\ge\rho_E$ holds by assumption.

Theorem~\ref{8.15:thm-epsilon-regularity}, which is a model statement, therefore gives, for every
plaquette $p$ based at $c$, $|\hat{F}(p)-F_0|\le\tau/4$, and in particular $\Lambda(p)=0$. Since
$\hat{F}=\Phi$ and $F_0=0$, this is $|\Phi(p)|\le\tau/4$ and $\Lambda(p)=0$ for the six plaquettes
$p=(c,\mu\nu)$, $\mu<\nu$. By the gauge invariance above, the same holds in every gauge.

For the last assertion, we argue as in Corollary~\ref{8.15:cor-cube-inactivity}, a model statement.
Apply the first assertion with $c$ replaced by each $c'\in Q_R(c)$. It gives $\Lambda(p)=0$ for every
plaquette $p$ based at a site of $Q_R(c)$.

Let $\xi\in\mathfrak{X}_{3R}$ vanish except on bonds with an endpoint in $Q_{R-1}(c)$. If
$(D_U\xi)(p)\neq0$, then some bond of $\partial p$ carries a nonzero value of $\xi$, so a corner of
$p$ lies in $Q_{R-1}(c)$. The base of $p$ lies at $|\cdot|_\infty$-distance at most one from each
corner of $p$, so $p$ is based at a site of $Q_R(c)$, and hence $\Lambda(p)=0$.
The multiplier rule \eqref{8.16:eq-capped-gauge-fields-a} therefore reduces to
\begin{equation*}
\sum_p\bigl\langle a_p\Phi(p),(D_U\xi)(p)\bigr\rangle=0 .
\end{equation*}
By the derivative formula of Definition~\ref{8.16:def-capped-gauge-fields} and
$\mathfrak{a}'(t)=t\,a(t)$, the left side equals
\begin{equation*}
\frac{d}{dt}\Big|_{t=0}\sum_p\mathfrak{a}\bigl(|\Phi(p)|\bigr)
\qquad\text{under } U_b\mapsto e^{t\xi_b}U_b .
\end{equation*}
So the action is stationary under every variation of the bonds of $Q_{R-1}(c)$; that is, $U$ is
exact lattice Yang--Mills there.
\end{proof}

\begin{remark}[General compact groups and invariant cap norms]\label{8.16:rem-general-group}
Let $G$ be a compact subgroup of a unitary group, let $\mathfrak{g}$ carry an $\operatorname{Ad}$-invariant inner
product with norm $|\cdot|$, and let $\|\cdot\|_c$ be an $\operatorname{Ad}$-invariant norm on
$\mathfrak{g}$ (for instance an operator norm), with $c_-|X|\le\|X\|_c\le c_+|X|$. Consider the setting of
Theorem~\ref{8.16:thm-gauge-regularity} with $\mathrm{SU}(2)$ replaced by $G$ and with the cap
$|\Phi(p)|\le\tau$ replaced by $\|\Phi(p)\|_c\le\tau$ on every plaquette. Let $\psi$ be the function,
holomorphic near $0$ with $\psi(0)=1$, for which
\begin{equation*}
\log(e^{tY}e^{X})=X+t\,\psi(\operatorname{ad}_X)Y+O(t^2)\qquad(X,Y\in\mathfrak{g},\ |X|\ \text{small}).
\end{equation*}

The multiplier of the cap at $p$ is then $\Lambda(p)\Theta_p$, with
$\Theta_p\in\partial\|\cdot\|_c(\Phi(p))$ a subgradient of $\|\cdot\|_c$ at $\Phi(p)$ (local notation),
up to the factor
\begin{equation*}
\psi(\operatorname{ad}_{\Phi(p)})^{\ast}=1+O(\tau).
\end{equation*}
This factor contributes one further defect, of relative size $\tau\le\varpi(1)$. The
colour-norm variant of the $\varepsilon$-regularity argument in
Remark~\ref{8.15:rem-regularity-remarks} then applies. At no point is a fixed colour component
singled out.
\end{remark}

\begin{proof}
We first identify the multiplier. In the multiplier rule assumed in
Theorem~\ref{8.16:thm-gauge-regularity}, the term $\Lambda(p)\Phi(p)/\tau$ is replaced, for the
convex cap $\|\cdot\|_c$, by $\psi(\operatorname{ad}_{\Phi(p)})^{\ast}\Lambda(p)\Theta_p$. Indeed,
under the variation $U_b\mapsto e^{t\xi_b}U_b$ of that setting, $\Phi(p)$ changes at first order by
$\psi(\operatorname{ad}_{\Phi(p)})(D_U\xi)(p)$, and
\begin{equation*}
\Lambda(p)\langle\Theta_p,\psi(\operatorname{ad}_{\Phi(p)})(D_U\xi)(p)\rangle
=\langle\psi(\operatorname{ad}_{\Phi(p)})^{\ast}\Lambda(p)\Theta_p,(D_U\xi)(p)\rangle .
\end{equation*}

Next we estimate the factor. Since $\psi(0)=1$, we have
$\psi(\operatorname{ad}_{\Phi(p)})^{\ast}-1=O(|\Phi(p)|)$. Moreover $|\Phi(p)|\le\tau/c_-$, so the
factor is $1+O(\tau)$.

We now bound the extra term. Taking $Y=\Theta_p$ in the subgradient inequality
$\langle\Theta_p,Y\rangle\le\|Y\|_c$ gives $|\Theta_p|^2\le\|\Theta_p\|_c\le c_+|\Theta_p|$, hence
$|\Theta_p|\le c_+$. Since $\operatorname{Ad}$ is orthogonal, $|(D_U\xi)(p)|\le4|\xi|_{\partial p}$.
Hence the remainder
$\Lambda(p)\bigl(\psi(\operatorname{ad}_{\Phi(p)})^{\ast}-1\bigr)\Theta_p$, paired with
$(D_U\xi)(p)$, is at most
\begin{equation*}
C\tau\,\Lambda(p)\,|\xi|_{\partial p},
\end{equation*}
with $C$ depending only on $G$ and $c_\pm$. This is one more defect of the form admitted in the
defect hypotheses checked for the fields of Theorem~\ref{8.16:thm-gauge-regularity} in its proof,
where at level $r$ the admitted defect carries the factor $\varpi(r)/r$ with
$\varpi(r)=C_Gr^2\tau$, so that $\varpi(r)/r\ge\varpi(1)$ for $r\ge1$; the new term is of relative
size $\tau\le\varpi(1)$.

The remaining multiplier $\Lambda(p)\Theta_p$ has the properties required by the colour-norm part
of Remark~\ref{8.15:rem-regularity-remarks}. At an active plaquette
$\langle\Theta_p,\Phi(p)\rangle=\tau$, and $\langle\Theta_p,X\rangle\le\|X\|_c$ for every $X$.
The argument there therefore goes through, with constants depending on $c_+/c_-$.

Finally, since $\|\cdot\|_c$ is $\operatorname{Ad}$-invariant, caps, subgradients and the factor
$\psi(\operatorname{ad}_{\Phi(p)})$ transform covariantly under gauge transformations, while
$\|\Phi(p)\|_c$ and $\Lambda(p)$ are invariant; so no colour component is distinguished.
\end{proof}

\begin{remark}[Where the non-abelian corrections enter, and their size]
\label{8.16:rem-nonabelian-corrections}
Consider the model theorem, Theorem~\ref{8.15:thm-epsilon-regularity}, applied to
$\mathrm{SU}(2)$ fields in the complete axial gauge about $c$. The application goes through
Lemma~\ref{8.16:lem-axial-defects}.
The non-abelian structure then enters the proof of the model theorem,
Theorem~\ref{8.15:thm-epsilon-regularity}, at exactly the following four places.
\begin{enumerate}
\item[(n1)] The curvature is not the lattice curl of the potential. This is the first defect
bound in \eqref{8.15:eq-capped-data-b}, and it is a pointwise error of size
$\varpi(r)\tau$ in the field.
\item[(n2)] The first variation is covariant. This is the second defect bound in
\eqref{8.15:eq-capped-data-b}. The error it admits is of size $\varpi(r)/r$ times the sum
over plaquettes of the magnitude $|\hat{F}|+\Lambda$ of field plus multiplier, taken against
the largest bond value $|\xi|_{\partial p}$ of the test $1$-form $\xi$.
\item[(n3)] The action is not quadratic. The reduced derivative $a$ of the action function
satisfies $a-1=O(\tau^2)$, and this error is contained in the second defect bound.
\item[(n4)] The average of translates is taken componentwise, in the one fixed gauge. The
average is admissible as soon as the cap is a norm ball of $\mathfrak{g}$, because this is
the same ball in every gauge. No parallel transport enters, and no error arises.
\end{enumerate}

\emph{Size.} Each of (n1)--(n3) enters the three lemmas that carry one step of the
iteration in the proof of the model theorem, Theorem~\ref{8.15:thm-epsilon-regularity},
and does so additively, as a term $C\varpi(r)^2\tau^2r^4$. In the quantities of that
iteration this reads
\begin{equation}\label{8.16:eq-nonabelian-corrections-1}
y_{k+1}\le\tfrac38\,y_k+C_{\sharp}\varpi_k^2,\qquad \varpi_k\le\varpi_R\cdot 96^{-2k}.
\end{equation}
This is a convergent perturbation of a contracting recursion, and the sign of the
corrections is immaterial.

\emph{The one ceiling.} The one ceiling, in the sense of
Definition~\ref{2.3:not-stages-first}, that the
non-abelian corrections impose is \eqref{8.16:eq-gauge-regularity-a}. To express it in
Ba{\l}aban's units, work on
the fine lattice and write the cap height and the radius as
$\tau=\theta_y\varepsilon_{k_o}\eta^2$ and $R=\rho_0$ fine lattice units; here $k_o$ is the
level of Ba{\l}aban's step at which the cube is read, $\varepsilon_{k_o}$ is the small-field
parameter of that step (not the mollification ratio of the proof of the model theorem,
Theorem~\ref{8.15:thm-epsilon-regularity}), $\eta$ is the spacing of the fine lattice in the
units of level
$k_o$, so that $\rho_0\eta$ is the radius of the cube in those units, and $\theta_y$ is the
factor, relative to the small-field bound $\varepsilon_{k_o}\eta^2$ on fine plaquettes, at
which the cap is set. In these units the quantity that \eqref{8.16:eq-gauge-regularity-a}
bounds is
\begin{equation}\label{8.16:eq-nonabelian-corrections-2}
R^2\tau=(\rho_0\eta)^2\theta_y\varepsilon_{k_o}\le D_W^2L^2\theta'\varepsilon_{k_o},
\end{equation}
where $D_W$ is the constant measuring the size of the hole, that is, of the large-field zone
about which the cube is read, in the units of level $k_o$, and
$\theta'$ is the corresponding cap factor entering the ceiling, both fixed with the hole.
Hence \eqref{8.16:eq-gauge-regularity-a} holds as soon as $\varepsilon_{k_o}$ times the
factor $C_GD_W^2L^2\theta'$ is at most $c_E$.
The right side of \eqref{8.16:eq-nonabelian-corrections-2} is $\varepsilon_{k_o}$ times a
factor, here $D_W^2L^2\theta'$, that depends
polynomially on the parameters of the hole and on $L$. The
condition is therefore of the same kind as two further ceilings on
the large-field zones, each of which
likewise bounds $\varepsilon_{k_o}$ times a factor depending polynomially on the parameters
of the hole, and which are met by
choosing the coupling $g$ after $L$ and $M$; it is one-sided.

\emph{Where the corrections would not be perturbative.} On cubes with $R^2\tau$ of order one
or larger, no gauge is small; this is the content of the lower bound for connections on
cubes. On such cubes
Theorem~\ref{8.16:thm-gauge-regularity} is void. For an isolated bounded zone this range is
excluded by ceilings
already imposed in this book. Outside that case the range belongs to the unbounded-union
depth statement, not to Theorem~\ref{8.16:thm-gauge-regularity}.

\emph{Comparison with the exit-value model.} The corresponding remark on the exit-value
model reaches the same point from the other side. There the non-abelian correction is
of the order of the exit-value model's small parameter times the main term, with sign
immaterial, and the abelian obstruction
remained. Here nothing abelian remains.
\end{remark}

\begin{proof}
\emph{Size.} The recursion in \eqref{8.15:eq-epsilon-regularity-b} was derived in the proof
of the model theorem, Theorem~\ref{8.15:thm-epsilon-regularity}, from data satisfying the
defect bounds of
\eqref{8.15:eq-capped-data-b}. The corrections (n1)--(n3) are exactly the terms carried by
those bounds, and (n4) contributes none. Each of (n1)--(n3) contributes
$C\varpi(r)^2\tau^2r^4$ at radius $r$. The excess $y_k$ is normalised by $\tau^2r_k^4$, so
at $r=r_k$ these contributions become the single term $C_{\sharp}\varpi_k^2$. This is the
first inequality of \eqref{8.16:eq-nonabelian-corrections-1}.

Lemma~\ref{8.16:lem-axial-defects} gives the defect function $\varpi(r)=C_Gr^2\tau$.
Hence
\begin{equation*}
\varpi(r)=\varpi_R\,(r/R)^2,\qquad \varpi_R=C_GR^2\tau.
\end{equation*}
The radii of the iteration start at $R$ and satisfy $r_{k+1}\le r_k/96$, so
$r_k\le 96^{-k}R$. It follows that
\begin{equation*}
\varpi_k=\varpi(r_k)\le\varpi_R\cdot96^{-2k},
\end{equation*}
which is the second inequality of \eqref{8.16:eq-nonabelian-corrections-1}.

The inequality $y_{k+1}\le\frac38 y_k+C_{\sharp}\varpi_k^2$, together with this bound on
$\varpi_k$, is the inequality from which the claim in \eqref{8.15:eq-epsilon-regularity-b}
is propagated by induction. That claim gives
\begin{equation*}
y_k\le 2^{-k}\bigl(y_0+8C_{\sharp}\varpi_R^2\bigr).
\end{equation*}
Moreover $\sum_k\varpi_k\le\varpi_R\sum_k96^{-2k}<\infty$. The perturbation of the
contracting recursion $y_{k+1}\le\frac38y_k$ is therefore convergent. The corrections enter
only through the non-negative term $C_{\sharp}\varpi_k^2$, which bounds their absolute
value, so their sign plays no role.

\emph{The ceiling.} Hypothesis \eqref{8.16:eq-gauge-regularity-a} reads
$C_GR^2\tau\le c_E$, that is, $\varpi_R\le c_E$. Since
$\varpi(r)=\varpi_R(r/R)^2\le\varpi_R(r/R)$ for $r\le R$, this supplies the defect
hypothesis of the model theorem, Theorem~\ref{8.15:thm-epsilon-regularity}. The hypotheses
under which Lemma~\ref{8.16:lem-axial-defects} supplies the defect bounds are
$\tau\le\tau_G$ and $R\tau\le c_G$. The first is a bound on the cap height alone and is also
a hypothesis of Theorem~\ref{8.16:thm-gauge-regularity}. For the second, note that the model
theorem, Theorem~\ref{8.15:thm-epsilon-regularity}, remains true with any smaller value of
$c_E$, since $c_E$ enters its proof only through conditions requiring it to be small, and
that $C_G$, $c_G$ depend only on the group $\mathrm{SU}(2)$ and the normalisations; we may
therefore assume $c_E\le C_Gc_G$. Then, since $R\ge\rho_E\ge1$,
\begin{equation*}
R\tau\le R^2\tau\le c_E/C_G\le c_G .
\end{equation*}
Hence no condition on the non-abelian terms other than
\eqref{8.16:eq-gauge-regularity-a} enters.
Substituting $\tau=\theta_y\varepsilon_{k_o}\eta^2$ and $R=\rho_0$ gives
\begin{equation*}
R^2\tau=\rho_0^2\theta_y\varepsilon_{k_o}\eta^2=(\rho_0\eta)^2\theta_y\varepsilon_{k_o},
\end{equation*}
and by the inequality $(\rho_0\eta)^2\theta_y\le D_W^2L^2\theta'$, taken here by statement
from the setting in which $\rho_0$, $\eta$ and $\theta_y$ are fixed, this is at most
$D_W^2L^2\theta'\varepsilon_{k_o}$. This is
\eqref{8.16:eq-nonabelian-corrections-2}. Multiplying it by $C_G$, we see that
\eqref{8.16:eq-gauge-regularity-a} holds as soon as
$C_GD_W^2L^2\theta'\varepsilon_{k_o}\le c_E$.

\emph{The non-perturbative range.} Hypothesis \eqref{8.16:eq-gauge-regularity-a} forces
$R^2\tau\le c_E/C_G$. On cubes where $R^2\tau$ is of order one or larger this fails, so
Theorem~\ref{8.16:thm-gauge-regularity} makes no assertion there.
\end{proof}

\begin{remark}[The slab theorem for $\mathrm{SU}(2)$ faces. Proof outstanding]\label{8.16:rem-gauge-slab}
The statement meant is that of Model Theorem~\ref{8.16:thm-slab-inactivity}, conclusions
(a)--(c), for $\mathrm{SU}(2)$ faces, that is, for the plaquette-capped $\mathrm{SU}(2)$ fields
subject to the multiplier rule \eqref{8.16:eq-capped-gauge-fields-a} in place of the
$\mathsf{V}$-valued data of Model Theorem~\ref{8.16:thm-slab-inactivity}.

Proof outstanding.

\emph{Route.} Beyond Theorem~\ref{8.16:thm-gauge-regularity}, which is proved as an implication
from the multiplier rule \eqref{8.16:eq-capped-gauge-fields-a}, a proof has to supply three
inputs. Here $N$, $H$, $\tau$ and $\mathfrak{E}_c$ are as in Model
Theorem~\ref{8.16:thm-slab-inactivity}, and we put
\begin{equation*}
\varpi_{\mathrm{glob}} := C_G (N+H)^2 \tau .
\end{equation*}
\begin{itemize}
\item[(1)] \emph{The comparison configuration.} The comparison configuration that enters the
proof of Model Theorem~\ref{8.16:thm-slab-inactivity} is constructed in a global axial gauge.
Its action must be at most $(1 + C\varpi_{\mathrm{glob}})\cdot\tfrac12\mathfrak{E}_c$, for a
constant $C$.
\item[(2)] \emph{The bound on the multiplier mass} (conclusion~(a) of Model
Theorem~\ref{8.16:thm-slab-inactivity}). The multiplier rule
\eqref{8.16:eq-capped-gauge-fields-a} is tested with $\xi$ equal to the logarithm of the
comparison configuration relative to the minimising configuration $\hat{U}$ of the capped
problem. The resulting defect is $\varpi_{\mathrm{glob}}$.
\item[(3)] \emph{Multiplier rule implies minimiser.} For conclusion~(c) of Model
Theorem~\ref{8.16:thm-slab-inactivity}, one must show that a configuration satisfying the
multiplier rule is a minimiser. The argument is to use the statement that the action is convex
along chart segments.
\end{itemize}
All three are to be established under the global form of condition
\eqref{8.16:eq-gauge-regularity-a}.
\end{remark}

\begin{definition}[The abelian block-slaved model]\label{8.16:def-block-slaved-model}
The \emph{abelian block-slaved model} is an abelian, linearised model of the structure of the
constrained minimisers of Ba{\l}aban's renormalisation step. Its fields are forms on a fine
lattice with values in the one-dimensional colour space $\mathsf{V}=\mathbb{R}$, the abelian case
of the data of Definition~\ref{8.15:def-capped-data}. A linear block-averaging map assigns block
data to each field, and caps are imposed on the curvature of the fine field that the block data
determine. In detail, it consists of the following objects.

\emph{Fine and block forms.} The fine lattice is $\mathbb{Z}^4$. Fine $1$-forms $\xi$ take real
values $\xi_\mu(z)$ on the bonds $(z,\mu)$, and fine $2$-forms take real values on the plaquettes
$p=(z,\mu\nu)$, $\mu<\nu$. For a function $f$ on $\mathbb{Z}^4$ we put
$\partial_\mu f(z):=f(z+e_\mu)-f(z)$, and the differential is
$(d\xi)_{\mu\nu}:=\partial_\mu\xi_\nu-\partial_\nu\xi_\mu$. Fine $0$-forms are real functions
$\lambda$ on the sites of $\mathbb{Z}^4$, and their differential is the fine $1$-form
$(d\lambda)_\mu:=\partial_\mu\lambda$. The pairing
$\langle X,Y\rangle:=\sum_p X(p)Y(p)$ is used whenever only finitely many terms are non-zero.
The letters $\mu,\nu$ as indices of forms, of $\partial$ and in the condition $\mu<\nu$ denote
lattice directions; $\mu$ as a $2$-form, written $\mu(p)$ or inside a pairing, denotes the
multiplier of the Karush--Kuhn--Tucker point below, and $\nu_a$ the multiplier of a window.

The block side is $\mathsf{b}\in\mathbb{N}$. Block $0$-forms and block $1$-forms are defined in
the same way on the lattice $\mathsf{b}\mathbb{Z}^4$, with differential $d^{\mathrm{bl}}$ given by
the formulas above with $e_\mu$ replaced by $\mathsf{b}e_\mu$. The map $Q$,
from fine $1$-forms to block $1$-forms, is linear, of finite range and gauge covariant: there is a
linear map $Q^{(0)}$ from fine $0$-forms to block $0$-forms such that
\begin{equation}\label{8.16:eq-block-slaved-model-1}
Q\,d\lambda=d^{\mathrm{bl}}Q^{(0)}\lambda\qquad\text{for every fine $0$-form }\lambda .
\end{equation}

\emph{Slaved fields and the slaved response.} A fine $1$-form $A$ is \emph{slaved} if the first
variation $\langle dA,d\xi\rangle$ of the formal Maxwell energy $\tfrac12\langle dA,dA\rangle$
vanishes along every finitely supported direction $\xi\in\ker Q$, that is, along its fibre
$A+\ker Q$:
\begin{equation}\label{8.16:eq-block-slaved-model-2}
\langle dA,d\xi\rangle=0\qquad\text{for every finitely supported }\xi\in\ker Q ;
\end{equation}
the sum is finite because $d\xi$ is finitely supported. For a block $1$-form $x$, called
\emph{block data}, $\mathcal{H}x$
denotes a slaved field with $Q\mathcal{H}x=x$. It is part of the model that for all block data $x$
such a field exists and that its curvature $d\mathcal{H}x$ does not depend on the choice of the
field; only this curvature is used below. For a
fine $1$-form $\xi$ we put
\begin{equation}\label{8.16:eq-block-slaved-model-3}
\mathcal{R}\xi:=d\mathcal{H}Q\xi ,
\end{equation}
the curvature of the slaved field whose block data are those of $\xi$. The following hold:
\begin{itemize}
\item[(a)] $\mathcal{R}\,d\lambda=0$ for every fine $0$-form $\lambda$;
\item[(b)] $\mathcal{R}\xi=d\xi$ for every slaved $\xi$;
\item[(c)] $d\xi-\mathcal{R}\xi=d\eta$ for some $\eta\in\ker Q$, namely
$\eta=\xi-\mathcal{H}Q\xi$;
\item[(d)] free Maxwell fields are slaved, and in particular so is every potential of constant
curvature. Here a free Maxwell field is a fine $1$-form $A$ with $\langle dA,d\xi\rangle=0$ for
every finitely supported fine $1$-form $\xi$. A constant curvature is a $2$-form $F_0$ given by
one number $F_{0,\mu\nu}$ per orientation.
\end{itemize}

\emph{Variables and constraints.} The variables are the block data $x$. The constraints are
\begin{equation}\label{8.16:eq-block-slaved-model-4}
|(d\mathcal{H}x)(p)|\le\tau
\end{equation}
for the fine plaquettes $p$ of a region, where $\tau>0$ is the cap height.

\emph{Karush--Kuhn--Tucker point.} A \emph{Karush--Kuhn--Tucker point} consists of
block data $\mathring{x}$ and a function
$\Lambda\ge0$ on fine plaquettes. Put $\hat{A}=\mathcal{H}\mathring{x}$, $\hat{F}:=d\hat{A}$ and
$\mu:=\Lambda\hat{F}/\tau$. The requirements are that $\hat{F}$ and $\Lambda$ satisfy the cap
condition and the complementary slackness condition \eqref{8.15:eq-capped-data-a} of
Definition~\ref{8.15:def-capped-data}, and that
\begin{equation}\label{8.16:eq-block-slaved-model-5}
\langle\hat{F}+\mu,\,d\mathcal{H}y\rangle=0\qquad\text{for every finitely supported block
$1$-form }y .
\end{equation}

\emph{Windows.} The constraints on the fine plaquettes of a cube $\tilde{a}$ may also be
aggregated into one constraint, called a \emph{window} $a$ with \emph{test cube} $\tilde{a}$,
which reads
\begin{equation*}
\sup_{p\in\tilde{a}}|(d\mathcal{H}x)(p)|\le\tau .
\end{equation*}
At a Karush--Kuhn--Tucker point the multiplier of this constraint in the sense of Clarke's
generalised gradient (its Clarke multiplier) is taken to be $\nu_a:=\tau\sum_{p\in\tilde{a}}\Lambda(p)$.
\end{definition}

\begin{proof}[Proof of \textup{(a)}--\textup{(d)}]
The proof of (a), (b) and (c) rests on one consequence of the uniqueness of the curvature of a
slaved field with given block data. Namely, if $A$ is any slaved field, then $d\mathcal{H}QA=dA$,
because $A$ and $\mathcal{H}QA$ are both slaved fields with block data $QA$.

(a) For every $z$ and $\mu<\nu$,
\begin{align*}
(d\,d\lambda)_{\mu\nu}(z)
&=\partial_\mu(d\lambda)_\nu(z)-\partial_\nu(d\lambda)_\mu(z)\\
&=\partial_\mu\partial_\nu\lambda(z)-\partial_\nu\partial_\mu\lambda(z)=0 ,
\end{align*}
since the difference operators $\partial_\mu$ and $\partial_\nu$ commute. Hence
$\langle d\,d\lambda,d\xi\rangle=0$ for every finitely supported $\xi$, and $d\lambda$ is slaved.
By \eqref{8.16:eq-block-slaved-model-1} its block data are
$Q\,d\lambda=d^{\mathrm{bl}}Q^{(0)}\lambda$.
The consequence of uniqueness stated above, applied to $A=d\lambda$, gives
\begin{equation*}
\mathcal{R}\,d\lambda=d\mathcal{H}Q\,d\lambda=d\,d\lambda=0 .
\end{equation*}

(b) Let $\xi$ be slaved. The same consequence, applied to $A=\xi$, gives
$\mathcal{R}\xi=d\mathcal{H}Q\xi=d\xi$.

(c) Put $\eta:=\xi-\mathcal{H}Q\xi$. Then $d\xi-\mathcal{R}\xi=d\eta$. Since $Q$ is linear and
$Q\mathcal{H}Q\xi=Q\xi$ by the defining property of $\mathcal{H}$, we get $Q\eta=Q\xi-Q\xi=0$,
that is, $\eta\in\ker Q$.

(d) Let $A$ be a free Maxwell field. Then $\langle dA,d\xi\rangle=0$ for every finitely supported
fine $1$-form $\xi$, in particular for every finitely supported $\xi\in\ker Q$. This is
\eqref{8.16:eq-block-slaved-model-2}, so $A$ is slaved.

Now let $dA=F_0$ be constant and let $\xi$ be finitely supported. Then
\begin{equation*}
\langle F_0,d\xi\rangle=\sum_{\mu<\nu}F_{0,\mu\nu}\sum_{z}
\bigl(\partial_\mu\xi_\nu(z)-\partial_\nu\xi_\mu(z)\bigr)=0 ,
\end{equation*}
because $\sum_z\partial_\mu f(z)=0$ for every finitely supported $f$: the sum telescopes. Hence a
potential of constant curvature is a free Maxwell field, and so it is slaved.
\end{proof}

\begin{definition}[Hypotheses on the slaved curvature response]\label{8.16:def-response-hypotheses}
We work in the abelian block-slaved model of Definition~\ref{8.16:def-block-slaved-model}, with
block side $\mathsf{b}$, block-averaging map $Q$ and slaved curvature response
$\mathcal{R}\xi=d\mathcal{H}Q\xi$. For $a\in\mathsf{b}\mathbb{Z}^4$ write $\mathsf{s}_a$ for
translation by $a$, acting on fine $1$-forms by $(\mathsf{s}_a\eta)(b)=\eta(b-a)$ and on fine
$2$-forms by $(\mathsf{s}_a F)(p)=F(p-a)$.

The following hypotheses are imposed on the region considered. The constants $C_{\mathcal{R}}$
and $\delta_0$ in them are independent of $\mathsf{b}$. In (T2) and (T3), $p$ is a fine plaquette
of the region considered, $q$ a fine plaquette, $b$ a fine bond, and $|b-p|$, $|q-p|$ are
distances on the fine lattice $\mathbb{Z}^4$.
\begin{itemize}
\item[(T1)] Covariance: $\mathcal{R}$ commutes with the translations of $\mathsf{b}\mathbb{Z}^4$.
\item[(T2)] Quasi-local response: for every fine $1$-form $\eta$ and every fine plaquette $p$ of
the region considered, the corresponding inequality of \eqref{8.16:eq-response-hypotheses-a} holds.
\item[(T3)] Regularity below the block scale: for every slaved fine field $A_1$ and every fine
plaquette $p$ of the region considered, the corresponding inequality of
\eqref{8.16:eq-response-hypotheses-a} holds.
\end{itemize}
\begin{equation}\label{8.16:eq-response-hypotheses-a}
\begin{gathered}
\mathcal{R}\,\mathsf{s}_a=\mathsf{s}_a\,\mathcal{R}\quad (a\in\mathsf{b}\mathbb{Z}^4),
\qquad\qquad
|(\mathcal{R}\eta)(p)|\le C_{\mathcal{R}}\,\mathsf{b}^{-5}\sum_{b}
e^{-\delta_0|b-p|/\mathsf{b}}\,|\eta(b)|,\\
|dA_1(p)|^2\le C_{\mathcal{R}}\,\mathsf{b}^{-4}\sum_{q}
e^{-\delta_0|q-p|/\mathsf{b}}\,|dA_1(q)|^2 .
\end{gathered}
\end{equation}
In \eqref{8.16:eq-response-hypotheses-a} the first relation is (T1), the second is (T2), and the
third is (T3).
\end{definition}

\begin{lemma}[Model: the load identity for block-slaved fields]\label{8.16:lem-load-identity}
The model is the abelian block-slaved model of Definition~\ref{8.16:def-block-slaved-model},
with slaving map $\mathcal{H}$ and slaved curvature response $\mathcal{R}\xi=d\mathcal{H}Q\xi$.

Let $(\hat{A},\hat{F},\Lambda)$ be a Karush--Kuhn--Tucker point of the constrained problem of
that definition, with $\hat{A}=\mathcal{H}x$ for its block data $x$, $\hat{F}=d\hat{A}$,
multiplier magnitude $\Lambda\ge 0$ and multiplier $2$-form $\mu=\Lambda\hat{F}/\tau$. Assume
the quasi-local response bound \textup{(T2)} of Definition~\ref{8.16:def-response-hypotheses},
displayed in \eqref{8.16:eq-response-hypotheses-a}. Then for every finitely supported fine
$1$-form $\xi$
\begin{equation}\label{8.16:eq-load-identity-a}
\langle \hat{F}, d\xi\rangle + \langle \mu, \mathcal{R}\xi\rangle = 0 .
\end{equation}
\end{lemma}

\begin{proof}
Fix a finitely supported fine $1$-form $\xi$.

The block $1$-form $Q\xi$ is finitely supported, because $\xi$ is and $Q$ has finite
range. The slaved field $\mathcal{H}Q\xi$ satisfies $Q\mathcal{H}Q\xi=Q\xi$. Hence
$Q(\xi-\mathcal{H}Q\xi)=0$, that is, $\xi-\mathcal{H}Q\xi\in\ker Q$. By the definition
$\mathcal{R}\xi=d\mathcal{H}Q\xi$ we have the decomposition
\begin{equation}\label{8.16:eq-load-identity-1}
d\xi=\mathcal{R}\xi+d(\xi-\mathcal{H}Q\xi),\qquad \xi-\mathcal{H}Q\xi\in\ker Q .
\end{equation}

The field $\hat{A}=\mathcal{H}x$ is slaved, so $\langle d\hat{A},d\eta\rangle=0$ for every
finitely supported $\eta\in\ker Q$. We use this for $\eta=\xi-\mathcal{H}Q\xi$, which lies in
$\ker Q$ by the first step but is in general not finitely supported. The pairing of $\hat{F}$
with $d\eta$, and the identity $\langle\hat{F},d\eta\rangle=0$ for this $\eta$, are justified
by the decay of $\mathcal{H}$ under \textup{(T2)}: by \eqref{8.16:eq-load-identity-1},
$d\eta=d\xi-\mathcal{R}\xi$, where $d\xi$ is finitely supported and \textup{(T2)} bounds
$|(\mathcal{R}\xi)(p)|$ by a finite sum of terms carrying the factor
$e^{-\delta_0|b-p|/\mathsf{b}}$. Pairing \eqref{8.16:eq-load-identity-1} with
$\hat{F}=d\hat{A}$ therefore gives
\begin{equation*}
\langle \hat{F}, d\xi\rangle=\langle \hat{F}, d\mathcal{H}(Q\xi)\rangle .
\end{equation*}

Next we use the Karush--Kuhn--Tucker condition of
Definition~\ref{8.16:def-block-slaved-model}: $\langle\hat{F}+\mu,d\mathcal{H}y\rangle=0$
for every finitely supported block $1$-form $y$. We apply it with $y=Q\xi$, which is
finitely supported by the first step. This gives
\begin{equation*}
\langle \hat{F}, d\mathcal{H}(Q\xi)\rangle=-\langle \mu, d\mathcal{H}(Q\xi)\rangle
=-\langle \mu,\mathcal{R}\xi\rangle ,
\end{equation*}
where the last equality is the definition of $\mathcal{R}$. Combining the two displays
yields \eqref{8.16:eq-load-identity-a}.
\end{proof}

\begin{remark}
Define the block $1$-form $\mu_{U}$ by $\langle\mu_{U},y\rangle:=-\langle\mu,d\mathcal{H}y\rangle$
for finitely supported block $1$-forms $y$. Then \eqref{8.16:eq-load-identity-a} reads
$\langle\hat{F},d\xi\rangle=\langle\mu_{U},Q\xi\rangle$. In this form the load
$\langle\hat{F},d\xi\rangle$ is carried by the block variables; it is the linearised, abelian
form of the expression $J=(DQ)^{*}\mu_{U}$ of the load in the exit-value model,
where $J$ denotes the load and $D$ the covariant variation of that model.

In \eqref{8.16:eq-load-identity-a} the fine load is smeared over the stencils of
$\mathcal{R}$, and these stencils have tails. Consequently, a fine load obtained from $\mu$
through these stencils may be non-zero at a plaquette $p$ where the constraint
$|\hat{F}(p)|\le\tau$ is not active. This does not affect the identity. Complementary
slackness enters \eqref{8.16:eq-load-identity-a} only through the multiplier $\mu$ itself,
and the smearing is carried entirely by the test field $\mathcal{R}\xi$.
\end{remark}

\begin{theorem}[$\varepsilon$-regularity for block-slaved fields]\label{8.16:thm-block-regularity}
Consider the abelian block-slaved model of Definition~\ref{8.16:def-block-slaved-model}, with block
side $\mathsf{b}$, block averaging $Q$, slaving map $\mathcal{H}$ and slaved curvature response
$\mathcal{R}$. Assume the hypotheses \textup{(T1)}, \textup{(T2)} and \textup{(T3)} of
Definition~\ref{8.16:def-response-hypotheses}, displayed in \eqref{8.16:eq-response-hypotheses-a},
on the region considered, with constants $C_{\mathcal{R}}$ and $\delta_0$ independent of $\mathsf{b}$.
Then there are constants $\rho_E'$ and $y_E'$, depending only on $C_{\mathcal{R}}$, $\delta_0$ and the
dimension, such that the following holds. Suppose that $R\ge\mathsf{b}\rho_E'$, that the
Karush--Kuhn--Tucker system of
Definition~\ref{8.16:def-block-slaved-model} holds on $Q_{3R}(c)$, and that the hybrid excess
\eqref{8.15:eq-capped-data-c} with reference $0$ obeys
\[
\mathcal{Y}(3R;0)\le y_E'\tau^2R^4 .
\]
Then every plaquette $p$ in the block at $c$ satisfies $|\hat{F}(p)|\le\tau/4$ and $\Lambda(p)=0$.
\end{theorem}

\begin{proof}
We follow the proof of the model theorem, Theorem~\ref{8.15:thm-epsilon-regularity}, with reference
$F_0=0$ at the top radius. That proof iterates over radii $r_k$ with references $F_0^{(k)}$. It is
built from Lemma~\ref{8.15:lem-multiplier-mass}, Lemma~\ref{8.15:lem-averaged-field} and the one-step
estimate that yields the recursion \eqref{8.15:eq-epsilon-regularity-b}, and it closes with an exact
step at the last radius. Four changes are made; every other step is used as it stands there.

In all four changes the mollification is taken at the block scale. Thus, with
$\mathsf{K}:=[-\kappa,\kappa]^4\cap\mathbb{Z}^4$ and $K$ the uniform probability on $\mathsf{K}$, the
average of translates of a form $X$ is now $X_K(p):=\sum_{h\in\mathsf{K}}K(h)X(p+\mathsf{b}h)$, so that
the translates range over $\mathsf{b}\cdot\mathsf{K}$. The ratio $\varepsilon:=\kappa\mathsf{b}/r$ takes
the place of $\kappa/r$. Below, $C$ denotes constants depending only on $C_{\mathcal{R}}$, $\delta_0$ and
the dimension.

\emph{First change: the test of Lemma~\ref{8.15:lem-multiplier-mass}.} Let $\zeta$ be the cutoff of the
proof of that lemma. It equals $1$ on $Q_r$, vanishes off $Q_{2r-1}$, and satisfies
$|\zeta(z)-\zeta(z')|\le|z-z'|_1/r$. Put $\beta_K:=\sum_{h\in\mathsf{K}}K(h)\beta_{\mathsf{b}h}$, where
$\beta_{\mathsf{b}h}$ is the translation potential of Lemma~\ref{8.15:lem-translation-potentials} for
the shift $\mathsf{b}h$ (that lemma applied with $\kappa\mathsf{b}$ in place of $\kappa$); its
staircases have length $|\mathsf{b}h|_1\le4\kappa\mathsf{b}$. We insert the test field
$\xi:=\zeta\beta_K$ in the model load identity \eqref{8.16:eq-load-identity-a} of the model lemma,
Lemma~\ref{8.16:lem-load-identity}, in place of the first-order system.

The field term of that identity, $\langle\hat{F},d\xi\rangle$, is the plain fine pairing. It is
therefore treated exactly as in the proof of Lemma~\ref{8.15:lem-multiplier-mass}.

For the multiplier term, put $A_K:=\sum_{h\in\mathsf{K}}K(h)\hat{A}(\cdot+\mathsf{b}h)$. By
\textup{(T1)} and linearity, $A_K-\hat{A}$ is slaved. By the properties
\eqref{8.15:eq-translation-potentials-a} of the translation
potentials, $A_K-\hat{A}=\beta_K+d\lambda_K$ for a function $\lambda_K$. By
Definition~\ref{8.16:def-block-slaved-model}, $\mathcal{R}\,d\lambda=0$ for every function $\lambda$,
and $\mathcal{R}\eta=d\eta$ for every slaved $\eta$. Hence
\[
\mathcal{R}\beta_K=d(A_K-\hat{A})=\hat{F}_K-\hat{F}.
\]
Fix a plaquette $p$, and let $\zeta(p)$ be the value of $\zeta$ at its base site. Write
$\xi=\zeta(p)\beta_K+(\zeta-\zeta(p))\beta_K$. Linearity of $\mathcal{R}$ gives
\begin{equation}\label{8.16:eq-block-regularity-1}
(\mathcal{R}\xi)(p)=\zeta(p)(\hat{F}_K-\hat{F})(p)+\bigl(\mathcal{R}[(\zeta-\zeta(p))\beta_K]\bigr)(p).
\end{equation}
We bound the last term using three facts. The first is \textup{(T2)}. The second is
$|\zeta(b)-\zeta(p)|\le|b-p|_1/r$. The third is the bound \textup{(t3)} of
\eqref{8.15:eq-translation-potentials-a}: the staircases now have length at most $4\kappa\mathsf{b}$,
and the curvature difference they sum is bounded by $2\tau$, as in the proof of
Lemma~\ref{8.15:lem-multiplier-mass}. Together these give
\begin{align*}
\bigl|\mathcal{R}[(\zeta-\zeta(p))\beta_K](p)\bigr|
&\le C_{\mathcal{R}}\mathsf{b}^{-5}\sum_be^{-\delta_0|b-p|/\mathsf{b}}
  \frac{|b-p|_1}{r}\cdot4\kappa\mathsf{b}\cdot2\tau\\
&\le C(C_{\mathcal{R}},\delta_0)\,\varepsilon\,\tau .
\end{align*}
The last inequality holds because $\sum_be^{-\delta_0|b-p|/\mathsf{b}}|b-p|_1$, a sum over the bonds
$b$ of the four-dimensional fine lattice, is at most a constant depending on $\delta_0$ times
$\mathsf{b}^5$.

Pair \eqref{8.16:eq-block-regularity-1} with $\mu(p)$ and use $|\mu|=\Lambda$. The first term is the
one that the proof of Lemma~\ref{8.15:lem-multiplier-mass} bounds from below through the slackness
relation $\langle\mu(p),\hat{F}(p)\rangle=\Lambda(p)\tau$ of the Karush--Kuhn--Tucker
system of
Definition~\ref{8.16:def-block-slaved-model}. The
second term is one more commutator, bounded by $C(C_{\mathcal{R}},\delta_0)\varepsilon\tau\Lambda(p)$.
This is the same size as the bound $16\varepsilon\tau$ of the multiplier commutator in that proof, and
it is paid for in the same way. Hence \eqref{8.15:eq-multiplier-mass-b} holds, with $C_J$ now
depending on $C_{\mathcal{R}}$ and $\delta_0$.

\emph{Second change: condition \eqref{8.15:eq-multiplier-mass-a}.} Put $D:=\hat{F}-F_0$. The reference
$F_0$ is a constant curvature, and potentials of constant curvature are slaved
(Definition~\ref{8.16:def-block-slaved-model}). By linearity, $D$ is the curvature of a slaved field.
Since $K$ is uniform on $(2\kappa+1)^4$ points, the Cauchy--Schwarz inequality gives
\[
|D_K(p)|^2\le(2\kappa+1)^{-4}\sum_h|D(p+\mathsf{b}h)|^2 .
\]
We apply \textup{(T3)} to each term and sum over $h$. We first split the resulting sum over $q$ into
the plaquettes in $Q_{3r}$ and the rest. For $q$ in $Q_{3r}$, the sum
$\sum_{h\in\mathsf{K}}e^{-\delta_0|q-p-\mathsf{b}h|/\mathsf{b}}$ is bounded by a constant depending on
$\delta_0$. This gives
\[
|D_K(p)|^2\le C(\mathsf{b}(2\kappa+1))^{-4}\|D\|^2_{3r}+T_{3r},
\]
where the tail $T_{3r}$, the contribution of the plaquettes $q$ outside $Q_{3r}$, satisfies
$T_{3r}\le Ce^{-\delta_0r/\mathsf{b}}\tau^2$.

Condition \eqref{8.15:eq-multiplier-mass-a} serves to bound $|D_K(p)|\le\tau/16$, the bound on which
the lower estimate of the multiplier term rests. It is therefore read here with its left side replaced
by the right side of the last display. Since $\mathsf{b}(2\kappa+1)\ge2\kappa\mathsf{b}=2\varepsilon r$,
this has the same form, a power $(\varepsilon r)^{-4}$ times $\|D\|^2_{3r}$, as in the iteration.

\emph{Third change: Lemma~\ref{8.15:lem-averaged-field}.} As in the proof of that lemma, let $w$ be the
orthogonal projection of $D_K$ restricted to $Q_r$ onto $\{d\xi:\xi\in\mathfrak{X}_r\}$. Since $w$ lies
in this space, choose $\xi_w\in\mathfrak{X}_r$ with $d\xi_w=w$, and put $H:=D_K-w$. Then
$\|w\|^2=\langle D_K,d\xi_w\rangle$.

For $h\in\mathsf{K}$ put $\xi^h:=\xi_w(\cdot-\mathsf{b}h)$, so that
$\langle D_K,d\xi_w\rangle=\sum_hK(h)\langle D,d\xi^h\rangle$. Since $\langle F_0,d\xi^h\rangle=0$, the
model load identity \eqref{8.16:eq-load-identity-a} gives
\[
\langle D,d\xi^h\rangle=\langle\hat{F},d\xi^h\rangle=-\langle\mu,\mathcal{R}\xi^h\rangle .
\]
By \textup{(T1)}, $\mathcal{R}\xi^h=(\mathcal{R}\xi_w)(\cdot-\mathsf{b}h)$. Summing against $K(h)$ yields
\[
\|w\|^2=\langle D_K,d\xi_w\rangle=-\langle\mu_K,\mathcal{R}\xi_w\rangle .
\]
The form $\mathcal{R}\xi_w$ is a slaved curvature. Moreover,
$d\xi_w-\mathcal{R}\xi_w=d(\xi_w-\mathcal{H}Q\xi_w)$ with $\xi_w-\mathcal{H}Q\xi_w\in\ker Q$, and
fibre-criticality makes this orthogonal to every slaved curvature. Thus $\mathcal{R}\xi_w$ is the
orthogonal projection of $d\xi_w=w$ onto the slaved curvatures, and $\|\mathcal{R}\xi_w\|\le\|w\|$.

By $|\mu_K|\le\Lambda_K$ and the Cauchy--Schwarz inequality with weight $\Lambda_K$ we obtain the first
line of the next display. For the second line we apply \textup{(T3)} to the slaved field
$\mathcal{H}Q\xi_w$, whose curvature is $\mathcal{R}\xi_w$, multiply by $\Lambda_K(p)$, sum over $p$,
exchange the sums, and use $\|\mathcal{R}\xi_w\|\le\|w\|$:
\begin{align*}
|\langle\mu_K,\mathcal{R}\xi_w\rangle|
&\le\Bigl(\sum_p\Lambda_K(p)\Bigr)^{1/2}\Bigl(\sum_p\Lambda_K(p)|\mathcal{R}\xi_w(p)|^2\Bigr)^{1/2}\\
&\le\Bigl(\sum_p\Lambda_K(p)\Bigr)^{1/2}\Bigl[C_{\mathcal{R}}\mathsf{b}^{-4}\sup_q\sum_p\Lambda_K(p)
  e^{-\delta_0|p-q|/\mathsf{b}}\Bigr]^{1/2}\|w\|\\
&\le C(\mathsf{b}(2\kappa+1))^{-2}\Bigl(\sum_p\Lambda(p)\Bigr)\|w\| .
\end{align*}
The last line uses two facts. First, $\sum_p\Lambda_K(p)=\sum_p\Lambda(p)$. Second, for every $q$,
\[
\sum_p\Lambda_K(p)e^{-\delta_0|p-q|/\mathsf{b}}
=(2\kappa+1)^{-4}\sum_{p'}\Lambda(p')\sum_he^{-\delta_0|p'-\mathsf{b}h-q|/\mathsf{b}},
\]
and the inner sum over the block translates is bounded by a constant depending on $\delta_0$.

Dividing by $\|w\|$ and using $\mathsf{b}(2\kappa+1)\ge\varepsilon r$, we get
\[
\|D_K-H\|_r=\|w\|\le C(\varepsilon r)^{-2}\sum_p\Lambda(p).
\]
This is the same superlinear bound as the multiplier term of \eqref{8.15:eq-averaged-field-a}: the
point multipliers of the fine lattice act at the block scale.

\emph{Fourth change: the exact step.} The iteration stops at a radius $r_{k^{\ast}}<\mathsf{b}\rho_E^{+}$.
Let $p$ be a plaquette in the block at $c$. The form $\hat{F}-F_0^{(k^{\ast})}$ is a slaved curvature,
and $\tfrac12|\hat{F}-F_0^{(k^{\ast})}|^2$ summed over $Q_{3r_{k^{\ast}}}$ is at most
$\mathcal{Y}(3r_{k^{\ast}};F_0^{(k^{\ast})})$. Applying \textup{(T3)} at $p$ and splitting the sum as in
the second change gives
\begin{align*}
|\hat{F}(p)-F_0^{(k^{\ast})}|^2
&\le C\mathsf{b}^{-4}\cdot2\,\mathcal{Y}(3r_{k^{\ast}};F_0^{(k^{\ast})})+T_{3r_{k^{\ast}}}\\
&\le C\,y_{k^{\ast}}(\rho_E^{+})^4\tau^2+Ce^{-\delta_0\rho_E'}\tau^2 .
\end{align*}
Here $T_{3r_{k^{\ast}}}$ is the tail from the plaquettes outside $Q_{3r_{k^{\ast}}}$, which is at most
$Ce^{-\delta_0\rho_E'}\tau^2$. We have also used
$\mathcal{Y}(3r_{k^{\ast}};F_0^{(k^{\ast})})=y_{k^{\ast}}\tau^2r_{k^{\ast}}^4$ and
$r_{k^{\ast}}<\mathsf{b}\rho_E^{+}$.

This bound replaces the pointwise bound at a single plaquette in the exact step of the model theorem,
Theorem~\ref{8.15:thm-epsilon-regularity}. The embedding of $\ell^2$ into $\ell^\infty$ is now taken at
the scale of one block, through \textup{(T3)}.
\end{proof}

\begin{remark}[The block scale of the exact step; local minimisers]
The exact step in the proof of Theorem~\ref{8.16:thm-block-regularity} is taken at the scale of one
block. Since $\rho_E'$ and $y_E'$ depend only on $C_{\mathcal{R}}$, $\delta_0$ and the dimension, and
$C_{\mathcal{R}}$ and $\delta_0$ are independent of $\mathsf{b}$, the conclusion is uniform in
$R/\mathsf{b}$ and in $\mathsf{b}$.

Suppose the windows $a$ of Definition~\ref{8.16:def-block-slaved-model} are read with the local
minimisers $\mathcal{H}_a$ attached to them in place of the slaving map
$\mathcal{H}$. Then the two differ by the relative error $\vartheta_1$ of local minimisers described
in \cite[(1.19)]{B-CMP119} and the remark following it. This is an additive defect of the kind measured
by the defect function $\varpi$ in the model theorem, Theorem~\ref{8.15:thm-epsilon-regularity}, and it
does not decay with the scale. The variant for defects that do not decay, in
Remark~\ref{8.15:rem-regularity-remarks}, with condition \eqref{8.15:eq-regularity-remarks-b},
therefore applies, with the one-sided ceiling
\begin{equation}\label{8.16:eq-block-regularity-a}
C\vartheta_1\log(R/\mathsf{b})\le\frac1{16}.
\end{equation}
This ceiling is harmless. Indeed, $\vartheta_1\le CB_3e^{-\delta LM_2R_{k_o}}$, with $B_3$ Ba{\l}aban's
regularity constant and $\delta>0$ Ba{\l}aban's decay constant of \cite[(1.19)]{B-CMP119}, and with
$R_{k_o}$ the length entering the exponent of that bound at the level $k_o$ at which the windows are
read; thus $\vartheta_1$ is exponentially small in $R_{k_o}$. On the other hand
$\log(R/\mathsf{b})\le C\log(D_WL)$, with $D_W$ a constant of the two-point witness construction.
\end{remark}

\begin{definition}[Notation for the top step of a merged region]\label{8.17:not-top-step-notation}
The following notation is used for the top step of a merged region.

\emph{Group and metric.} Let $G=\mathrm{SU}(N)$, let $\mathfrak{g}$ be its Lie algebra and
$d(\mathfrak{g}):=\dim\mathfrak{g}$. On $\mathfrak{g}$ put
\begin{equation*}
|Y|:=\bigl(-\tfrac12\operatorname{tr}Y^{2}\bigr)^{1/2}.
\end{equation*}
Let $d(\cdot,\cdot)$ be the bi-invariant Riemannian distance on $G$, with the metric normalised by
$|\cdot|$ at the identity, and $\iota_G$ its injectivity radius, so that $d(\exp Y,1)=|Y|$ for
$|Y|\le\iota_G$. About any $a\in G$ we use the
exponential chart $Y\mapsto a\exp(Y)$, $|Y|<\iota_G$.

\emph{Levels and scales.} The union $\mathcal{N}$ of large-field regions under consideration is
live at the level $K''$. The integer $c_0\ge 2$ is the number of observation levels, and
\begin{equation*}
k_o:=K':=K''+c_0
\end{equation*}
is the top observation level. We put $\varepsilon:=\varepsilon_{k_o}$, where $\varepsilon_k$ is the
small-field threshold at level $k$, and
$\varepsilon_{\min}:=\min(\varepsilon_{k_o-1},\varepsilon_{k_o})$; further $g:=g_{k_o}$, and
$p_0:=p_0(g_{k_o})$, the value at $g_{k_o}$ of Ba{\l}aban's
degree function $p_0(\cdot)$, which grows polynomially, of degree $\mathfrak{p}_0$, in
$\log g^{-2}$; $p_0$ written without argument always means this value.
The following constants are functions of $L$, $N$, the dimension $4$ and Ba{\l}aban's constants:
the window response constants $K_w$ and $C_{\delta_0}$, whose product enters the Lipschitz bound for
windows, and $K_w'$, which enters the second-derivative bound for windows; the radius $\alpha_P$ of
the bondwise neighbourhoods within which the regularity and Lipschitz bounds for $\mathbf{A}$,
$\Psi$ and the windows used in this chapter hold; the
constant $C'$ of the bound $C'\varepsilon$ for the right-invariant derivatives of the background
action $\mathbf{A}$ introduced below; the constants $C$ and $\alpha_0$ of the Lipschitz bound
$C/\alpha_0$ for the localised small terms $\Psi$ introduced below; the Haar-density constant $c_G$
of the exponential chart; and the income constants $\gamma_0''$ and $\gamma'$. We adopt, as a
convention of normalisation of the small-field threshold,
\begin{equation*}
\varepsilon g/g^{2}=p_0 ,
\end{equation*}
that is, $\varepsilon=g\,p_0$; with it the bound $C'\varepsilon g/g^{2}$ for the change of
$\mathbf{A}/g^{2}$ under a bondwise displacement of total size $g$ equals $C'p_0$.

\emph{Regions.} $\mathfrak{z}$ is a merged top-level raw region: a connected component of the union
of the sub-zones $\mathfrak{z}_{\square}$ of the $c_0$ top levels around the recent seed cubes
$\square$ of members of $\mathcal{N}$. We write $n_s$ for the number of seed cubes whose sub-zone
lies in $\mathfrak{z}$, and $H\subset\mathfrak{z}$ for the level-$k_o$ hole of $\mathfrak{z}$; the
depth and the inradius $h(\mathfrak{z})$ of a merged region are those of the relative-region lemma.
The region $\mathfrak{z}$ is \emph{$h$-deep} if $H$ contains a lattice ball of radius $h$.
The number $h_g$ is the largest depth at which the chart ceiling and the per-seed ceilings of the
flat form of the volume theorem for merged regions hold with that theorem's parameter $\sigma_P$
set equal to $\sigma_P^{\mathrm{rel}}(h)$; it is a boundary of
ranges of validity of statements of this chapter and enters none of the definitions below.

\emph{Variables and windows of the top step.} The moved variables are the tower bond variables
$x=(x_b)_{b\in\mathfrak{B}}\in X:=G^{\mathfrak{B}}$ of $\mathfrak{z}$; $n_b:=|\mathfrak{B}|$, and
$\mathfrak{d}=\mathfrak{d}(\mathfrak{z}):=d(\mathfrak{g})\,n_b$ is the number of real coordinates
of $X$. Let
$\mathfrak{d}_*$ be the supremum over $\square$ of the number of real coordinates of one sub-zone,
and $n_*:=\mathfrak{d}_*/d(\mathfrak{g})$. By the counting bounds for merged regions,
\begin{equation*}
\mathfrak{d}(\mathfrak{z})\le n_s\,\mathfrak{d}_*,\qquad
\Omega_{\Psi}(\mathfrak{z})\le n_s\,\Omega_{\Psi*},
\end{equation*}
with $\Psi$ as below. $dH$ is normalised Haar measure on $X$. $W$ is the finite family of window
functionals $f_a(e,\cdot)\colon X\to[0,\infty)$, where $e$ is the external datum; a window
\emph{passes} at $x$ if $f_a(e,x)<1$. The family $W$ consists of the non-seed windows of all levels
that depend on the variables of $H$, the straddling family $\mathfrak{S}$ (among the windows of
$W$, those of $\mathfrak{S}$ are the only ones that depend on the top field of $e$), and the seed
windows $f_{\square}$. For $\theta\in[1,2]$ and $\Lambda'\ge1$ put
\begin{align*}
S^{\theta}_{\mathrm{all}}(e)&:=\{x\in X:\ f_a(e,x)<\theta\ \text{for every non-seed window } a\in W\},\\
S(e)&:=S^{1}_{\mathrm{all}}(e),\\
\mathcal{E}&:=\{x\in X:\ f_{\square}(e,x)\ge 1\ \text{for every sandwiched seed }\square\},
\end{align*}
and let $S^{1,\Lambda'}(e)$ be the set of $x\in X$ with
\begin{equation*}
\begin{aligned}
&f_a(e,x)<1 \ \text{for every non-seed window } a\notin\mathfrak{S},\qquad
f_a(e,x)<\Lambda'\ \text{for } a\in\mathfrak{S},\\
&f_{\square}(e,x)<1\ \text{for every seed }\square .
\end{aligned}
\end{equation*}
$\mathbf{A}(e,x)$ is Ba{\l}aban's background action of the minimal configuration for the nested data
\cite[Theorem~1]{B-CMP102b}; it is defined and continuous on the crude-regular class, which
contains $\operatorname{cl}S^{2}_{\mathrm{all}}(e)$ by the lemma on the crude-regular class of this
chapter and, in the form in which that lemma is used in this chapter, its bondwise
$\alpha_P$-neighbourhood. $\Psi$ collects the localised small terms;
$\Omega_{\Psi}(\mathfrak{z})$ denotes the bound for the oscillation $\sup-\inf$ of $\Psi(e,\cdot)$
over $\operatorname{cl}S^{2}_{\mathrm{all}}(e)$, and $\Omega_{\Psi*}$ the corresponding bound for a
single sub-zone. We put
\begin{align*}
\operatorname{Num}_{\lambda}(e)&:=\int_{S(e)\cap\mathcal{E}}
\exp\bigl(-\mathbf{A}/g^{2}+\Psi\bigr)\,dH,\\
\operatorname{Den}^{\Lambda'}_{\lambda}(e)&:=\int_{S^{1,\Lambda'}(e)}
\exp\bigl(-\mathbf{A}/g^{2}+\Psi\bigr)\,dH.
\end{align*}

\emph{The top step.} The top-step inequality for $\mathfrak{z}$ is
\begin{equation}\label{8.17:eq-top-step-notation-a}
\operatorname{Num}_{\lambda}(e)\le\Bigl[\prod_{\square\subset\mathfrak{z}}\epsilon_{\square}\Bigr]
\operatorname{Den}^{\Lambda'}_{\lambda}(e'),
\end{equation}
with
\begin{equation*}
\epsilon_{\square}\le\exp\bigl(-\tfrac12\gamma_{\mathrm{inc}}\,p_0(g_{j(\square)})^{2}\bigr),\qquad
e'\in\{e,\Xi e\},\qquad \Lambda'\le\Lambda_L,
\end{equation*}
and with no constant depending on $n_s$, $\operatorname{diam}\mathfrak{z}$ or $h(\mathfrak{z})$;
here $j(\square)$ is the level at which the seed cube $\square$ was created, $\Xi e$ is the external
datum transformed by the move $\Xi$ attached to the top step, and $\Lambda_L$ is Ba{\l}aban's bound
for the relaxed cap $\Lambda'$ of the straddling windows.
This is how the top step is stated in the skin version and in the flat version of the
volume theorem for merged regions. The seed income is a number $\Gamma_{\square}\ge0$ for each
sandwiched seed $\square$, with
\begin{equation}\label{8.17:eq-top-step-notation-1}
\Gamma_{\square}/g^{2}\ge\gamma_{\mathrm{inc}}\,p_0(g_{j(\square)})^{2},
\end{equation}
where $\gamma_{\mathrm{inc}}:=\gamma_0''$ when the upper bound for $\operatorname{Num}_{\lambda}(e)$
is the one used in the skin version of the volume theorem, and $\gamma_{\mathrm{inc}}:=\gamma'$ when
the upper bound is the one supplied by the majorant lemma.

\emph{Window response constants.} We put
\begin{equation*}
\operatorname{Lip}:=\frac{K_wC_{\delta_0}}{\varepsilon_{\min}},\qquad
\rho_1:=\frac{2K_wC_{\delta_0}\,g}{\varepsilon_{\min}}=2\operatorname{Lip}\cdot g;
\end{equation*}
$\rho_1$ is the margin that lets configurations within bondwise distance $g$ of a minimiser $x_0$
of $\mathbf{A}(e,\cdot)$ over $\operatorname{cl}S^{\theta}_{\mathrm{all}}(e)$, for the
$\theta\in[1,2]$ at hand,
pass, by the Lipschitz bound for windows alone.
\end{definition}

\begin{definition}[The abelian window model and its plaquette instance]
\label{8.17:def-abelian-window-model}
The \emph{abelian window model} $\mathbb{M}$ consists of the data (m1)--(m3) below; (m4) defines its
\emph{plaquette instance} $\mathbb{M}_{\square}$.
\begin{enumerate}
\item[(m1)] A finite set $\mathfrak{B}$, whose elements are called bonds, with $n:=|\mathfrak{B}|$.
We put $V:=\mathbb{R}^{\mathfrak{B}}$, with the norm $|w|_\infty:=\max_{b}|w_b|$ and Lebesgue
measure $dy$.
\item[(m2)] Windows: a finite index set $I$ and, for each $i\in I$, an affine function
\begin{equation}\label{8.17:eq-abelian-window-model-1}
\ell_i(y)=\ell_i(0)+\sum_{b\in\mathfrak{B}} l_{i,b}\,y_b
\end{equation}
whose linear part $l_i$ is not zero, together with a cap $c_i>0$. We put
\[
\varphi_i:=\|l_i\|_1:=\sum_{b}|l_{i,b}|,\qquad
\det(i):=\{b:\ l_{i,b}\neq0\},\qquad f_i:=|\ell_i|/c_i .
\]
The number $\varphi_i$ is the Lipschitz constant of $\ell_i$ for $|\cdot|_\infty$. The set
$S:=\{y:\ f_i(y)<1\ \text{for all } i\}$ is assumed non-empty. For $y$ with $f_i(y)\le1$ the
\emph{margin} of the window $i$ at $y$ is
$\operatorname{margin}_i(y):=c_i-|\ell_i(y)|\ge0$. The windows are \emph{wide at scale $g$} if
$c_i\ge3g\varphi_i$ for all $i$.
\item[(m3)] Action: a convex, continuously differentiable function $A\colon V\to\mathbb{R}$ with
\begin{equation}\label{8.17:eq-abelian-window-model-2}
|\partial_bA(y)-\partial_bA(y')|\le\Lambda_1\,|y-y'|_\infty
\qquad\text{for all } y,y'\in V \text{ and all } b .
\end{equation}
A parameter $x>0$, and $g:=x^{-1/2}$. For a point $x_{\mathrm{ref}}\in V$ we put
\[
G_b:=|\partial_bA(x_{\mathrm{ref}})|,\qquad P_b:=x\,G_b\,g=G_b/g,\qquad P:=\max_b P_b .
\]
If the small-field threshold is normalised as $\varepsilon=p_0g$, $p_0$ being the value of the
degree function $p_0(\cdot)$ at the coupling $g$, then a bound $G_b\le C'\varepsilon$ for all $b$,
with some constant $C'$, gives $P\le C'p_0$ with the same constant.
\item[(m4)] The plaquette instance $\mathbb{M}_{\square}$: a finite set of oriented bonds of the
cubic lattice $\mathbb{Z}^d$, $d\ge2$, split into free bonds, which form $\mathfrak{B}$, and held
bonds, which carry fixed real values. Let $\mathcal{P}$ be the set of plaquettes with at least one
free bond. We assume that for every $p\in\mathcal{P}$ all four bonds of $\partial p$ belong to the
given set of bonds, free or held; then
\begin{equation}\label{8.17:eq-abelian-window-model-3}
F_p(y):=\sum_{b\in\partial p}\sigma(p,b)\cdot(\text{value on } b)
\end{equation}
is defined, where $\sigma(p,b)\in\{1,-1\}$ records whether the orientation of $b$ agrees with the
boundary orientation of $p$, the value on a free bond $b$ is $y_b$ and the value on a held bond is
its fixed value. Then $F_p$ is an affine
function of $y$ whose linear part $l_p$ is supported on the free bonds of $p$, with coefficients
$\pm1$ and $\varphi_p\in\{1,2,3,4\}$. The windows are $|F_p|\le\kappa_p$, with caps $\kappa_p>0$;
in the notation of (m2), $I=\mathcal{P}$, $\ell_p=F_p$ and $c_p=\kappa_p$. The action is
\[
A(y):=\tfrac12\sum_{p\in\mathcal{P}}F_p(y)^2 .
\]
We assume that the linear map $y\mapsto(l_p(y))_{p\in\mathcal{P}}$ is injective (for instance
when $\mathfrak{B}$ is the complement of a spanning forest of a region with held boundary).
\end{enumerate}
In $\mathbb{M}_{\square}$ one has
\begin{equation}\label{8.17:eq-abelian-window-model-4}
\partial_bA(y)=\sum_{p\ni b}\sigma(p,b)\,F_p(y),
\end{equation}
so that $G_b\le2(d-1)\max_{p\ni b}\kappa_p$ whenever $x_{\mathrm{ref}}\in\operatorname{cl}S$, and
\eqref{8.17:eq-abelian-window-model-2} holds with $\Lambda_1=8(d-1)$; moreover
$\operatorname{cl}S$ is compact and $A$ attains its minimum on $\operatorname{cl}S$.
\end{definition}

\begin{proof}
\emph{The Lipschitz constant of $\ell_i$.} By \eqref{8.17:eq-abelian-window-model-1},
\[
|\ell_i(y)-\ell_i(y')|=\Bigl|\sum_b l_{i,b}(y_b-y'_b)\Bigr|\le\varphi_i\,|y-y'|_\infty ,
\]
with equality when $y-y'$ has entries $\operatorname{sign}(l_{i,b})$. The inequality
$\operatorname{margin}_i(y)\ge0$ is the hypothesis $f_i(y)\le1$ multiplied by $c_i$.

\emph{$\mathbb{M}_{\square}$ is an instance of $\mathbb{M}$.} For $p\in\mathcal{P}$ the coefficient
of $y_b$ in \eqref{8.17:eq-abelian-window-model-3} is $\sigma(p,b)=\pm1$ if $b$ is a free bond of
$\partial p$ and $0$ otherwise, since the four bonds of $\partial p$ are distinct. Hence $l_p$ is
supported on the free bonds of $p$, $\varphi_p$ is their number, which lies in
$\{1,2,3,4\}$ because $p$ has at least one free bond, and in particular $l_p\neq0$. The action is a
sum of squares of affine functions, hence convex and continuously differentiable. By the chain
rule, $\partial_bA(y)=\sum_{p\in\mathcal{P}}F_p(y)\,\partial_bF_p(y)$ with
$\partial_bF_p=\sigma(p,b)$ if $b\in\partial p$ and $0$ otherwise; every plaquette through a free
bond $b$ lies in $\mathcal{P}$, which gives \eqref{8.17:eq-abelian-window-model-4}.

\emph{The bound on $G_b$.} A bond of $\mathbb{Z}^d$ in direction $\mu$ lies in exactly two
plaquettes of each of the $d-1$ coordinate planes spanned by $\mu$ and $\nu\neq\mu$, hence in
$2(d-1)$ plaquettes. On $S$ one has $|F_p|<\kappa_p$, so $|F_p|\le\kappa_p$ on
$\operatorname{cl}S$ by
continuity. For $x_{\mathrm{ref}}\in\operatorname{cl}S$, \eqref{8.17:eq-abelian-window-model-4}
gives
\[
G_b\le\sum_{p\ni b}|F_p(x_{\mathrm{ref}})|\le2(d-1)\max_{p\ni b}\kappa_p .
\]

\emph{The value of $\Lambda_1$.} By \eqref{8.17:eq-abelian-window-model-4} and
$F_p(y)-F_p(y')=l_p(y-y')$,
\[
|\partial_bA(y)-\partial_bA(y')|\le\sum_{p\ni b}|l_p(y-y')|
\le\sum_{p\ni b}\varphi_p\,|y-y'|_\infty\le2(d-1)\cdot4\,|y-y'|_\infty ,
\]
since each of the at most $2(d-1)$ plaquettes through $b$ reads at most four free bonds with unit
coefficients. This is \eqref{8.17:eq-abelian-window-model-2} with $\Lambda_1=8(d-1)$.

\emph{Compactness and minimisers.} Write $T(y):=(l_p(y))_{p\in\mathcal{P}}$. The set
$\operatorname{cl}S$ is closed and contained in
$\{y:\ |F_p(y)|\le\kappa_p\ \text{for all } p\}$, on
which $|l_p(y)|=|F_p(y)-F_p(0)|\le\kappa_p+|F_p(0)|$. Since $T$ is linear and injective on the
finite-dimensional space $V$, there is $C_T<\infty$ with $|y|_\infty\le C_T\max_p|l_p(y)|$, so
\[
|y|_\infty\le C_T\max_{p}\bigl(\kappa_p+|F_p(0)|\bigr)\qquad\text{for } y\in\operatorname{cl}S .
\]
Thus $\operatorname{cl}S$ is closed and bounded in $V$, hence compact; it is non-empty because
$S$ is. The continuous function $A$ therefore attains its minimum on $\operatorname{cl}S$.
\end{proof}

\begin{lemma}[Closure of the open window set]\label{8.17:lem-window-set-closure}
In the abelian window model of Definition~\ref{8.17:def-abelian-window-model}, let
$S=\{y\in V: f_i(y)<1 \text{ for all } i\in I\}$. If $S\neq\emptyset$, then
\begin{equation*}
\operatorname{cl} S=\{y\in V: f_i(y)\le 1 \text{ for all } i\in I\},
\end{equation*}
and $\operatorname{cl} S\setminus S$ is a null set for the Lebesgue measure $dy$ on $V$.
\end{lemma}

\begin{proof}
Fix $i\in I$. The function $\ell_i$ is affine, and the absolute value is convex. Hence
$|\ell_i|$ is convex and continuous, and so is $f_i=|\ell_i|/c_i$, because $c_i>0$.

The set $S_{\le}=\{y: f_i(y)\le 1 \text{ for all } i\in I\}$ is the intersection of the
finitely many sublevel sets $\{f_i\le 1\}$. Each of these is closed by continuity of $f_i$,
so $S_{\le}$ is closed. Since $S\subset S_{\le}$, we obtain $\operatorname{cl} S\subset S_{\le}$.

Conversely, let $y\in S_{\le}$ and choose $y_0\in S$, which is possible because $S\neq\emptyset$.
For every $t\in(0,1]$ and every $i\in I$, convexity of $f_i$ gives
\begin{equation*}
f_i\bigl((1-t)y+ty_0\bigr)\le(1-t)f_i(y)+t f_i(y_0)<(1-t)+t=1 .
\end{equation*}
The strict inequality holds because $(1-t)f_i(y)\le 1-t$, while $tf_i(y_0)<t$ since $t>0$
and $f_i(y_0)<1$. Hence $(1-t)y+ty_0\in S$ for every $t\in(0,1]$. Moreover,
\begin{equation*}
\bigl|\bigl((1-t)y+ty_0\bigr)-y\bigr|_\infty=t\,|y_0-y|_\infty ,
\end{equation*}
which tends to $0$ as $t\to 0$. Thus $y$ is a limit of points of $S$,
so $y\in\operatorname{cl} S$. This proves $S_{\le}\subset\operatorname{cl} S$, and hence
$\operatorname{cl} S=S_{\le}$.

For the second assertion, note first that $S$ is open. It is a finite intersection of the
sets $\{f_i<1\}$, each of which is open because $f_i$ is continuous. Next, $S$ is convex.
Indeed, each $\{f_i<1\}$ is convex: if $f_i(u)<1$, $f_i(v)<1$ and $t\in[0,1]$, then
$f_i((1-t)u+tv)\le(1-t)f_i(u)+tf_i(v)<1$. Since $S$ is open, $\operatorname{cl}
S\setminus S$ is the boundary of $S$. So $S$ is a convex set with nonempty interior
(namely $S$ itself), and the boundary of a convex set with nonempty interior is
Lebesgue-null. Hence $\operatorname{cl} S\setminus S$ is Lebesgue-null.
\end{proof}

\begin{lemma}[First-order conditions at a constrained minimiser]\label{8.17:lem-first-order-minimiser}
Let the bond set $\mathfrak{B}$, the space $V=\mathbb{R}^{\mathfrak{B}}$, the windows $(\ell_i,c_i)_{i\in I}$
with the quantities $l_i$, $\varphi_i$, $\det(i)$, $f_i$, the set $S$ and the action $A$ be as in
Definition~\ref{8.17:def-abelian-window-model}. Write
$\langle v,w\rangle:=\sum_{b\in\mathfrak{B}}v_bw_b$ on $V$, regard $l_i=(l_{i,b})_{b\in\mathfrak{B}}$ as
a vector of $V$, and put $l_i(w):=\langle l_i,w\rangle$, so that $\ell_i(y+w)=\ell_i(y)+l_i(w)$. Let
$I'\subset I$, let
\begin{equation*}
S':=\{y\in V: f_i(y)<1 \text{ for all } i\in I'\},
\end{equation*}
so that $S\subset S'$, and let $x_0$ minimise $A$ over $\operatorname{cl} S'$. Put
$\mathcal{A}:=\{i\in I': f_i(x_0)=1\}$, the \emph{tight set}, and
$s_i:=\operatorname{sgn}\ell_i(x_0)\in\{\pm1\}$ for $i\in\mathcal{A}$. Then:
\begin{itemize}
\item[(a)] $\langle\nabla A(x_0),y-x_0\rangle\ge0$ for every $y\in\operatorname{cl} S'$;
\item[(b)] there are numbers $\mu_i\ge0$, $i\in\mathcal{A}$, such that
\begin{equation}\label{8.17:eq-first-order-minimiser-1}
\nabla A(x_0)=-\sum_{i\in\mathcal{A}}\mu_i s_i l_i ;
\end{equation}
\item[(c)] $\partial_bA(x_0)=0$ for every $b\notin\bigcup_{i\in\mathcal{A}}\det(i)$.
\end{itemize}
\end{lemma}

\begin{proof}
For $i\in\mathcal{A}$ we have $|\ell_i(x_0)|=c_i>0$, so $s_i$ is defined and
$s_i\ell_i(x_0)=c_i$. The set $S'$ is the intersection of the sublevel sets
$\{f_i<1\}$, $i\in I'$, of the convex functions $f_i$; hence $S'$ and its closure are convex.
The windows indexed by $I'$ form a window family in the sense of
Definition~\ref{8.17:def-abelian-window-model} whose open window set $S'$ contains $S\neq\emptyset$,
so Lemma~\ref{8.17:lem-window-set-closure}, applied to this family, gives
\begin{equation*}
\operatorname{cl} S'=\{y\in V: f_i(y)\le1 \text{ for all } i\in I'\}.
\end{equation*}

(a) Let $y\in\operatorname{cl} S'$. For $t\in[0,1]$ the point $x_0+t(y-x_0)$ lies in
$\operatorname{cl} S'$ by convexity, so $A(x_0+t(y-x_0))\ge A(x_0)$ by the minimality of $x_0$. The
function $t\mapsto A(x_0+t(y-x_0))$ is continuously differentiable with derivative
$\langle\nabla A(x_0),y-x_0\rangle$ at $t=0$; this derivative is the limit as $t\to0^+$ of the
difference quotients $t^{-1}[A(x_0+t(y-x_0))-A(x_0)]\ge0$, and is therefore non-negative.

(b) Let $T:=\{w\in V: s_il_i(w)\le0 \text{ for all } i\in\mathcal{A}\}$, a convex cone. For
$y\in\operatorname{cl} S'$ and $i\in\mathcal{A}$,
\begin{equation*}
s_i\ell_i(y)\le|\ell_i(y)|\le c_i=s_i\ell_i(x_0),
\end{equation*}
so $s_il_i(y-x_0)=s_i\ell_i(y)-s_i\ell_i(x_0)\le0$; thus $\operatorname{cl} S'-x_0\subset T$.
Conversely, with the convention that a minimum over the empty set is $+\infty$, put
\begin{equation*}
\eta:=\min\Bigl\{\min_{i\in I'\setminus\mathcal{A}}\frac{c_i-|\ell_i(x_0)|}{\varphi_i},\
\min_{i\in I'}\frac{2c_i}{\varphi_i}\Bigr\}.
\end{equation*}
Here $\eta>0$: each $\varphi_i>0$ because $l_i\neq0$, and for $i\in I'\setminus\mathcal{A}$ we have
$f_i(x_0)\le1$ and $f_i(x_0)\neq1$, hence $c_i-|\ell_i(x_0)|>0$. Let $w\in T$ with $|w|_\infty\le\eta$.
Since $\varphi_i$ is the Lipschitz constant of $\ell_i$ for $|\cdot|_\infty$, $|l_i(w)|\le\varphi_i\eta$.
For $i\in I'\setminus\mathcal{A}$ this gives
\begin{equation*}
|\ell_i(x_0+w)|\le|\ell_i(x_0)|+\eta\varphi_i\le c_i .
\end{equation*}
For $i\in\mathcal{A}$, $s_i\ell_i(x_0+w)=c_i+s_il_i(w)\le c_i$ because $w\in T$, and
\begin{equation*}
s_i\ell_i(x_0+w)\ge c_i-\eta\varphi_i\ge c_i-2c_i=-c_i ,
\end{equation*}
so again $|\ell_i(x_0+w)|\le c_i$. Hence $f_i(x_0+w)\le1$ for all $i\in I'$, that is,
$x_0+w\in\operatorname{cl} S'$, and (a) with $y=x_0+w$ gives $\langle\nabla A(x_0),w\rangle\ge0$. For an
arbitrary $w\in T\setminus\{0\}$ choose $t>0$ with $t|w|_\infty\le\eta$; then $tw\in T$, so
$t\langle\nabla A(x_0),w\rangle\ge0$. Therefore $\langle\nabla A(x_0),w\rangle\ge0$ for every $w\in T$.

We now use Farkas' lemma, the classical duality statement of finite-dimensional linear
programming: if $q,m_i\in V$ and $\langle q,w\rangle\ge0$ for every $w$ with $\langle m_i,w\rangle\le0$
for all $i$, then $q=-\sum_i\lambda_im_i$ with $\lambda_i\ge0$. Applied with $q:=\nabla A(x_0)$ and
$m_i:=s_il_i$, $i\in\mathcal{A}$, whose hypothesis is the property of $T$ just established, it gives
$\lambda_i\ge0$ with $\nabla A(x_0)=-\sum_{i\in\mathcal{A}}\lambda_is_il_i$; with $\mu_i:=\lambda_i$
this is \eqref{8.17:eq-first-order-minimiser-1}.

(c) Taking the $b$-th coordinate of \eqref{8.17:eq-first-order-minimiser-1},
$\partial_bA(x_0)=-\sum_{i\in\mathcal{A}}\mu_is_il_{i,b}$. If
$b\notin\bigcup_{i\in\mathcal{A}}\det(i)$, then $l_{i,b}=0$ for every $i\in\mathcal{A}$, and hence
$\partial_bA(x_0)=0$.
\end{proof}

\begin{theorem}[Model: lower bound by a short move]\label{8.17:thm-model-short-move}
We work in the abelian window model of Definition~\ref{8.17:def-abelian-window-model}: a finite
set $\mathfrak{B}$ of bonds with $n:=|\mathfrak{B}|$, the space $V=\mathbb{R}^{\mathfrak{B}}$ with the
norm $|\cdot|_\infty$ and Lebesgue measure $dy$, affine windows $\ell_i$ $(i\in I)$ with linear
parts $l_i$, caps $c_i>0$, norms $\varphi_i=\|l_i\|_1$ and supports $\det(i)$, the open window set
$S$ with closure $\operatorname{cl} S$, a convex continuously differentiable action $A$ whose
partial derivatives are $\Lambda_1$-Lipschitz for $|\cdot|_\infty$, and a parameter $x>0$ with
$g=x^{-1/2}$.

Let $x_{\mathrm{ref}}\in V$, $\delta\in(0,g]$, and $u\in V$ with $|u|_\infty\le 1$ such that
$z^*:=x_{\mathrm{ref}}+\delta u\in\operatorname{cl} S$. Let $G_b:=|\partial_bA(x_{\mathrm{ref}})|$,
$P_b:=G_b/g$ and $P:=\max_bP_b$, as in Definition~\ref{8.17:def-abelian-window-model}. For
$i\in I$ and $b\in\mathfrak{B}$ define
\begin{equation*}
\tau_i:=\min\Bigl\{\frac{g}{\delta},\ \frac{c_i-|\ell_i(z^*)|}{\delta\varphi_i}\Bigr\}
\in\Bigl[0,\frac{g}{\delta}\Bigr],
\qquad
r_b:=\delta\cdot\min\{\tau_i : b\in\det(i)\},
\end{equation*}
the minimum over the empty set being $g/\delta$, so that $0\le r_b\le g$; here
$c_i-|\ell_i(z^*)|\ge0$ is the margin of the window $i$ at $z^*$. Put
$\varrho_b:=\delta|u_b|+r_b\le 2g$ and $R:=z^*+\prod_{b}[-r_b,r_b]$. Then $R\subset\operatorname{cl} S$ and
\begin{equation}\label{8.17:eq-model-short-move-a}
\int_S e^{-xA(y)}\,dy\ \ge\ \int_R e^{-xA(y)}\,dy
\ \ge\ e^{-xA(x_{\mathrm{ref}})}\prod_{b\in\mathfrak{B}}(2r_b)\,
\exp\Bigl(-\sum_{b\in\mathfrak{B}}\bigl[xG_b\varrho_b+4\Lambda_1\bigr]\Bigr).
\end{equation}
If moreover $x_{\mathrm{ref}}=x_0$ is a minimiser of $A$ over $\operatorname{cl} S'$ for some sub-family
$I'\subset I$, as in Lemma~\ref{8.17:lem-first-order-minimiser}, then $r_b\le\delta$ and
$\varrho_b\le2\delta$ for every $b$ with $G_b>0$, whence:
\begin{enumerate}
\item[(i)] for $\delta=g$,
\begin{equation}\label{8.17:eq-model-short-move-1}
\int_S e^{-xA}\,dy\ \ge\ e^{-xA(x_0)}\,g^n\prod_b\frac{2r_b}{g}\,
\exp\Bigl(-\sum_b\bigl[2P_b+4\Lambda_1\bigr]\Bigr),
\end{equation}
that is, relative to the reference $e^{-xA(x_0)}g^n$, a loss of $\log(g/(2r_b))+4\Lambda_1$ per
coordinate, plus $2P_b\le2P$ on each coordinate read by a tight window of $I'$, and nothing
further on the others;
\item[(ii)] for $\delta=g/\max(1,P)$,
\begin{equation}\label{8.17:eq-model-short-move-b}
\int_S e^{-xA}\,dy\ \ge\ e^{-xA(x_0)}\,g^n\prod_b\frac{2r_b}{g}\,
\exp\Bigl(-\sum_b\bigl[2\cdot\mathbf{1}\{G_b>0\}+4\Lambda_1\bigr]\Bigr),
\end{equation}
where now
\begin{equation*}
\frac{g}{2r_b}=\frac{\max(1,P)}{2\cdot\min\{\tau_i : b\in\det(i)\}};
\end{equation*}
with $\tau_b:=\min\{\tau_i : b\in\det(i)\}$, this is a loss of
\begin{equation*}
\log\max(1,P)+\log\frac{1}{2\tau_b}+2+4\Lambda_1
\end{equation*}
per coordinate read by some window with $\tau_i<g/\delta$, and of $4\Lambda_1-\log2$ on every
other coordinate.
\end{enumerate}
\end{theorem}

\begin{proof}
\emph{Preliminaries.} Since $z^*\in\operatorname{cl} S$, Lemma~\ref{8.17:lem-window-set-closure} gives
$f_i(z^*)\le1$, that is $|\ell_i(z^*)|\le c_i$, for every $i\in I$; hence the margin
$c_i-|\ell_i(z^*)|$ is non-negative and $\tau_i\in[0,g/\delta]$. Consequently
$0\le r_b\le\delta\cdot g/\delta=g$, with $r_b=g$ when no window reads $b$, and
$\varrho_b\le\delta|u_b|+g\le\delta+g\le2g$ because $|u|_\infty\le1$ and $\delta\le g$.

\emph{The box lies in the closure.} Let $y=z^*+w$ with $|w_b|\le r_b$ for every $b$. For
$i\in I$ the function $\ell_i$ is affine with linear part $l_i$, so
$\ell_i(y)=\ell_i(z^*)+\sum_{b\in\det(i)}l_{i,b}w_b$. For $b\in\det(i)$ the index $i$ enters the
minimum defining $r_b$, so $r_b\le\delta\tau_i$. Therefore
\begin{align*}
|\ell_i(y)|
&\le|\ell_i(z^*)|+\sum_{b\in\det(i)}|l_{i,b}|\,r_b
\le|\ell_i(z^*)|+\delta\tau_i\varphi_i\\
&\le|\ell_i(z^*)|+\bigl(c_i-|\ell_i(z^*)|\bigr)=c_i,
\end{align*}
where the last inequality is $\tau_i\le(c_i-|\ell_i(z^*)|)/(\delta\varphi_i)$, which holds by
the definition of $\tau_i$ as a minimum. Thus $f_i(y)\le1$ for all $i$, and $y\in\operatorname{cl} S$ by
Lemma~\ref{8.17:lem-window-set-closure}; so $R\subset\operatorname{cl} S$. By the same lemma
$\operatorname{cl} S\setminus S$ is Lebesgue-null, and the integrand is positive, so
\begin{equation*}
\int_S e^{-xA}\,dy=\int_{\operatorname{cl} S}e^{-xA}\,dy\ \ge\ \int_R e^{-xA}\,dy,
\end{equation*}
which is the first inequality of \eqref{8.17:eq-model-short-move-a}.

\emph{The action on the box.} Let $y\in R$. Then $y-x_{\mathrm{ref}}=\delta u+w$ with
$|w_b|\le r_b$, so $|y_b-x_{\mathrm{ref},b}|\le\delta|u_b|+r_b=\varrho_b$ for every $b$, and
$|y-x_{\mathrm{ref}}|_\infty\le\delta+g\le2g$. Since $A$ is continuously differentiable, the
fundamental theorem of calculus along the segment from $x_{\mathrm{ref}}$ to $y$ gives
\begin{align*}
A(y)-A(x_{\mathrm{ref}})
&=\int_0^1\sum_b\partial_bA\bigl(x_{\mathrm{ref}}+t(y-x_{\mathrm{ref}})\bigr)
\,(y_b-x_{\mathrm{ref},b})\,dt\\
&\le\sum_b\sup_{t\in[0,1]}\bigl|\partial_bA\bigl(x_{\mathrm{ref}}+t(y-x_{\mathrm{ref}})\bigr)\bigr|
\,\varrho_b\\
&\le\sum_b\bigl(G_b+\Lambda_1|y-x_{\mathrm{ref}}|_\infty\bigr)\varrho_b
\ \le\ \sum_b\bigl(G_b\varrho_b+2\Lambda_1g\cdot2g\bigr).
\end{align*}
In the third line we used the Lipschitz bound on $\partial_bA$ of
Definition~\ref{8.17:def-abelian-window-model}, which gives
$|\partial_bA(x_{\mathrm{ref}}+t(y-x_{\mathrm{ref}}))|\le G_b+\Lambda_1t|y-x_{\mathrm{ref}}|_\infty$
for $t\in[0,1]$; in the last step $|y-x_{\mathrm{ref}}|_\infty\le2g$ and $\varrho_b\le2g$.
Multiplying by $x=g^{-2}$,
\begin{equation}\label{8.17:eq-model-short-move-2}
x\bigl(A(y)-A(x_{\mathrm{ref}})\bigr)\ \le\ \sum_b\bigl[xG_b\varrho_b+4\Lambda_1\bigr]
\qquad(y\in R).
\end{equation}

\emph{The volume.} The box $R$ has volume $\operatorname{vol}(R)=\prod_b(2r_b)$. Hence
$\int_Re^{-xA}\,dy\ge\operatorname{vol}(R)\cdot\inf_Re^{-xA}$, and by
\eqref{8.17:eq-model-short-move-2} the infimum is at least
$e^{-xA(x_{\mathrm{ref}})}\exp(-\sum_b[xG_b\varrho_b+4\Lambda_1])$. This gives the second
inequality of \eqref{8.17:eq-model-short-move-a}.

\emph{The minimiser clause.} Let $x_{\mathrm{ref}}=x_0$ minimise $A$ over $\operatorname{cl} S'$, and let
$\mathcal{A}$ be the tight set of Lemma~\ref{8.17:lem-first-order-minimiser}. Suppose $G_b>0$,
that is $\partial_bA(x_0)\ne0$. By Lemma~\ref{8.17:lem-first-order-minimiser}(c),
$b\in\det(i)$ for some $i\in\mathcal{A}$, i.e.\ for some window $i$ tight in $I'\subset I$. For
such $i$ we have $f_i(x_0)=1$, so the margin at $x_0$ vanishes: $|\ell_i(x_0)|=c_i$. Since
$\varphi_i$ is the Lipschitz constant of $\ell_i$ for $|\cdot|_\infty$ and
$|z^*-x_0|_\infty=\delta|u|_\infty\le\delta$,
\begin{equation*}
|\ell_i(z^*)|\ \ge\ |\ell_i(x_0)|-\delta\varphi_i=c_i-\delta\varphi_i,
\end{equation*}
so $c_i-|\ell_i(z^*)|\le\delta\varphi_i$, hence $\tau_i\le1$. As $i\in I$ and $b\in\det(i)$,
$r_b\le\delta\tau_i\le\delta$, and $\varrho_b\le\delta|u_b|+\delta\le2\delta$. Then, using
$P_b=xG_bg$,
\begin{equation*}
xG_b\varrho_b\ \le\ 2xG_b\delta=\frac{2P_b\delta}{g},
\end{equation*}
which equals $2P_b$ at $\delta=g$ and is at most $2P_b/\max(1,P)\le2$ at
$\delta=g/\max(1,P)$. If $G_b=0$ the term $xG_b\varrho_b$ vanishes, and so do $P_b$ and
$\mathbf{1}\{G_b>0\}$. Hence $\sum_bxG_b\varrho_b\le\sum_b2P_b$ in case (i) and
$\sum_bxG_b\varrho_b\le\sum_b2\cdot\mathbf{1}\{G_b>0\}$ in case (ii); both values of $\delta$
lie in $(0,g]$. Inserting these bounds in \eqref{8.17:eq-model-short-move-a}, with
$x_{\mathrm{ref}}=x_0$ and $\prod_b(2r_b)=g^n\prod_b(2r_b/g)$, gives
\eqref{8.17:eq-model-short-move-1} and \eqref{8.17:eq-model-short-move-b}.

In case (i), a coordinate not read by any tight window of $I'$ has $G_b=0$ by
Lemma~\ref{8.17:lem-first-order-minimiser}(c), so $P_b=0$ there. In case (ii),
$r_b=\delta\tau_b=g\tau_b/\max(1,P)$, which is the displayed expression for $g/(2r_b)$; when no
window reads $b$ the minimum is $g/\delta$, $r_b=g$ and $g/(2r_b)=1/2$.
\end{proof}

\begin{proposition}[Duality criterion for the opening constant]\label{8.17:prop-opening-constant-duality}
Work in the abelian window model of Definition~\ref{8.17:def-abelian-window-model}, and for
$m\in\mathbb{R}^{\mathfrak{B}}$ and $u\in V$ write $\langle m,u\rangle:=\sum_{b}m_b u_b$ and
$\|m\|_1:=\sum_b|m_b|$. Identify each linear part $l_i$ with its coefficient vector
$(l_{i,b})_{b\in\mathfrak{B}}$, so that $l_i(u)=\langle l_i,u\rangle$ and $\varphi_i=\|l_i\|_1$.
Let $\mathcal{A}'\subset I$ be a set of indices equipped with signs $s_i\in\{\pm1\}$, and let
$\tau>0$. The following are equivalent:
\begin{equation}\label{8.17:eq-opening-constant-duality-b}
\begin{aligned}
&\text{(a)}\quad \text{there is } u\in V \text{ with } |u|_\infty\le 1
  \text{ and } -s_i l_i(u)\ge\tau\varphi_i \text{ for every } i\in\mathcal{A}';\\
&\text{(b)}\quad \text{for every } \lambda\in\mathbb{R}^{\mathcal{A}'}_{\ge0}:\quad
  \Bigl\|\sum_{i\in\mathcal{A}'}\lambda_i s_i l_i\Bigr\|_1
  \ge \tau\sum_{i\in\mathcal{A}'}\lambda_i\varphi_i .
\end{aligned}
\end{equation}
Consequently the \emph{opening constant}
\begin{equation}\label{8.17:eq-opening-constant-duality-a}
\mathfrak{i}(\mathcal{A}',s):=\sup\{\tau\ge0:\ \text{(a) holds}\}
=\inf\Bigl\{\frac{\|\sum_i\lambda_i s_i l_i\|_1}{\sum_i\lambda_i\varphi_i}:\
\lambda\ge0,\ \lambda\ne0\Bigr\}
\end{equation}
lies in $[0,1]$, and when $\mathfrak{i}(\mathcal{A}',s)>0$ the supremum in
\eqref{8.17:eq-opening-constant-duality-a} is attained.
\end{proposition}

\begin{proof}
We use throughout the elementary duality: for $m\in\mathbb{R}^{\mathfrak{B}}$,
$\sup_{|u|_\infty\le1}\langle m,u\rangle=\|m\|_1$, the supremum being attained at $u_b=1$ where
$m_b\ge0$ and $u_b=-1$ where $m_b<0$; the bound $\langle m,u\rangle\le\|m\|_1$ for $|u|_\infty\le1$
is the termwise estimate $m_bu_b\le|m_b|$.

(a)$\Rightarrow$(b). Let $u$ be as in (a) and let $\lambda\in\mathbb{R}^{\mathcal{A}'}_{\ge0}$.
Since $\|\cdot\|_1$ is even and $|u|_\infty\le1$, and since each $\lambda_i\ge0$,
\begin{equation*}
\Bigl\|\sum_i\lambda_i s_i l_i\Bigr\|_1
\ge\Bigl\langle-\sum_i\lambda_i s_i l_i,\,u\Bigr\rangle
=\sum_i\lambda_i\bigl(-s_i l_i(u)\bigr)
\ge\tau\sum_i\lambda_i\varphi_i .
\end{equation*}

(b)$\Rightarrow$(a). Suppose (a) fails. Put
\begin{equation*}
\mathcal{P}_\tau:=\bigl\{\bigl(-s_i l_i(u)-\tau\varphi_i\bigr)_{i\in\mathcal{A}'}:\ |u|_\infty\le1\bigr\}
\subset\mathbb{R}^{\mathcal{A}'} .
\end{equation*}
It is the image of the cube $\{|u|_\infty\le1\}$ under an affine map, hence compact and convex.
The failure of (a) says exactly that no point of $\mathcal{P}_\tau$ has all coordinates
nonnegative, i.e.\ $\mathcal{P}_\tau$ is disjoint from the closed convex cone
$\mathbb{R}^{\mathcal{A}'}_{\ge0}$. By the separating hyperplane theorem for a compact convex set
and a disjoint closed convex set there are $\lambda\in\mathbb{R}^{\mathcal{A}'}$, $\lambda\ne0$, and
$\eta\in\mathbb{R}$ with
\begin{equation*}
\langle\lambda,k\rangle<\eta\le\langle\lambda,c\rangle
\qquad\text{for all } k\in\mathcal{P}_\tau,\ c\in\mathbb{R}^{\mathcal{A}'}_{\ge0},
\end{equation*}
where $\langle\cdot,\cdot\rangle$ is now the standard pairing on $\mathbb{R}^{\mathcal{A}'}$.
Taking $c=0$ gives $\eta\le0$. If $\lambda_j<0$ for some $j$, then $c=te_j$ with $t\to\infty$ makes
$\langle\lambda,c\rangle=t\lambda_j$ unbounded below, contradicting $\eta\le\langle\lambda,c\rangle$;
hence $\lambda\ge0$ componentwise. Thus for every $u$ with $|u|_\infty\le1$,
\begin{equation*}
\sum_i\lambda_i\bigl(-s_i l_i(u)\bigr)<\tau\sum_i\lambda_i\varphi_i+\eta
\le\tau\sum_i\lambda_i\varphi_i .
\end{equation*}
The supremum over $u$ of the left side is $\|\sum_i\lambda_i s_i l_i\|_1$ and is attained by
the duality recalled above, so the strict inequality persists:
$\|\sum_i\lambda_i s_i l_i\|_1<\tau\sum_i\lambda_i\varphi_i$. This contradicts (b) for this
$\lambda\in\mathbb{R}^{\mathcal{A}'}_{\ge0}$.

Identification of $\mathfrak{i}(\mathcal{A}',s)$. Denote by $\mathfrak{i}_*$ the infimum on the right
of \eqref{8.17:eq-opening-constant-duality-a}. For $\lambda\ge0$, $\lambda\ne0$, we have
$\sum_i\lambda_i\varphi_i>0$, because $\varphi_i=\|l_i\|_1>0$ by $l_i\ne0$; for $\lambda=0$ the
inequality in (b) is trivial. Hence, for $\tau>0$, (b) holds if and only if
$\tau\le\mathfrak{i}_*$, and by the equivalence just proved the same is true of (a). At $\tau=0$,
(a) holds with $u=0$. So the set of $\tau\ge0$ for which (a) holds is $[0,\mathfrak{i}_*]$, and its
supremum is $\mathfrak{i}_*$. Clearly $\mathfrak{i}_*\ge0$. By the triangle inequality,
$\|\sum_i\lambda_i s_i l_i\|_1\le\sum_i\lambda_i\|l_i\|_1=\sum_i\lambda_i\varphi_i$, so every
quotient in \eqref{8.17:eq-opening-constant-duality-a} is at most $1$, and
$\mathfrak{i}(\mathcal{A}',s)\le1$.

Attainment. The set
\begin{equation*}
\mathcal{Q}:=\bigl\{(u,\tau)\in V\times[0,1]:\ |u|_\infty\le1,\
-s_i l_i(u)\ge\tau\varphi_i \text{ for all } i\in\mathcal{A}'\bigr\}
\end{equation*}
is closed, being cut out by finitely many non-strict inequalities between continuous functions,
and bounded, hence compact; it contains $(0,0)$. Since every $\tau\ge0$ for which (a) holds satisfies
$\tau\le\mathfrak{i}(\mathcal{A}',s)\le1$, the maximum of the continuous function $(u,\tau)\mapsto\tau$
on $\mathcal{Q}$ equals $\mathfrak{i}(\mathcal{A}',s)$ and is attained at some $(u^*,\mathfrak{i}(\mathcal{A}',s))$.
When $\mathfrak{i}(\mathcal{A}',s)>0$, this $u^*$ satisfies (a) with $\tau=\mathfrak{i}(\mathcal{A}',s)$.
\end{proof}

\begin{corollary}[Model: short-move radii from the opening constant]\label{8.17:cor-model-radii}
We work in the abelian affine-window model $\mathbb{M}$ of
Theorem~\ref{8.17:thm-model-short-move}: the affine windows $\ell_i$, $i\in I$, on $V$, with
linear parts $l_i$, caps $c_i$, $\varphi_i$ the $\ell^1$-norm of $l_i$,
$\operatorname{margin}_i=c_i-|\ell_i|$, the function $A$, and the notation of that theorem.
Assume that the windows are wide at scale $g$. Let $I'\subset I$, and let $x_0$ minimise $A$
over $\operatorname{cl} S'$, with $x_0\in\operatorname{cl} S$. For $\delta\in(0,g]$ put
\begin{equation*}
\mathcal{A}_{\delta}:=\{\,i\in I:\ \operatorname{margin}_i(x_0)<2\delta\varphi_i\,\},
\end{equation*}
the windows of the whole family that are tight at scale $\delta$. For
$i\in\mathcal{A}_{\delta}$ put $s_i:=\operatorname{sgn}\ell_i(x_0)$. This sign is nonzero, since
\begin{equation*}
|\ell_i(x_0)|>c_i-2\delta\varphi_i\ge g\varphi_i>0 .
\end{equation*}
Put $\mathfrak{i}_{\delta}:=\mathfrak{i}(\mathcal{A}_{\delta},s)$, the opening constant
\eqref{8.17:eq-opening-constant-duality-a}. If $\mathfrak{i}_{\delta}>0$, there is $u\in V$
with $|u|_\infty\le1$ such that, in Theorem~\ref{8.17:thm-model-short-move},
$z^*:=x_0+\delta u\in\operatorname{cl} S$ and
\begin{equation*}
\tau_i\ge\min(1,\mathfrak{i}_{\delta})=\mathfrak{i}_{\delta}\qquad\text{for every } i\in I ;
\end{equation*}
hence $r_b\ge\delta\,\mathfrak{i}_{\delta}$ for every $b\in\mathfrak{B}$.
\end{corollary}

\begin{proof}
Since $\mathfrak{i}_{\delta}>0$, Proposition~\ref{8.17:prop-opening-constant-duality} states
that the supremum defining $\mathfrak{i}(\mathcal{A}_{\delta},s)$ is attained. Hence
condition~(a) in \eqref{8.17:eq-opening-constant-duality-b} holds with
$\mathcal{A}'=\mathcal{A}_{\delta}$ and
$\tau=\mathfrak{i}_{\delta}$. We choose $u\in V$ accordingly: $|u|_\infty\le1$, and
$-s_il_i(u)\ge\mathfrak{i}_{\delta}\varphi_i$ for every $i\in\mathcal{A}_{\delta}$. Put
$z^*:=x_0+\delta u$. Then $\ell_i(z^*)=\ell_i(x_0)+\delta\, l_i(u)$. Because $\varphi_i$ is the
$\ell^1$-norm of the coefficients of $l_i$, we have
$|l_i(u)|\le\varphi_i|u|_\infty\le\varphi_i$.

Let $i\in\mathcal{A}_{\delta}$. By the choice of $s_i$ we have
$s_i\ell_i(x_0)=|\ell_i(x_0)|$, and $|\ell_i(x_0)|\le c_i$ because $x_0\in\operatorname{cl} S$.
Therefore
\begin{equation*}
s_i\ell_i(z^*)=s_i\ell_i(x_0)-\delta\bigl(-s_il_i(u)\bigr)
\le c_i-\delta\,\mathfrak{i}_{\delta}\varphi_i .
\end{equation*}
In the other direction, $i\in\mathcal{A}_{\delta}$ gives $|\ell_i(x_0)|>c_i-2\delta\varphi_i$,
and wideness gives $c_i-3\delta\varphi_i\ge0$. Hence
\begin{equation*}
s_i\ell_i(z^*)\ge s_i\ell_i(x_0)-\delta\varphi_i>c_i-3\delta\varphi_i\ge0 .
\end{equation*}
So $|\ell_i(z^*)|=s_i\ell_i(z^*)$, and the first bound gives
$\operatorname{margin}_i(z^*)\ge\delta\,\mathfrak{i}_{\delta}\varphi_i$. That is,
$\operatorname{margin}_i(z^*)/(\delta\varphi_i)\ge\mathfrak{i}_{\delta}$. Since also
$g/\delta\ge1$, the definition of $\tau_i$ in Theorem~\ref{8.17:thm-model-short-move} gives
$\tau_i\ge\min(1,\mathfrak{i}_{\delta})$. This equals $\mathfrak{i}_{\delta}$, because
Proposition~\ref{8.17:prop-opening-constant-duality} gives $\mathfrak{i}_{\delta}\in[0,1]$.

Let $i\notin\mathcal{A}_{\delta}$. Then $\operatorname{margin}_i(x_0)\ge2\delta\varphi_i$. Since
$\bigl||\ell_i(z^*)|-|\ell_i(x_0)|\bigr|\le\delta|l_i(u)|\le\delta\varphi_i$, we get
\begin{equation*}
\operatorname{margin}_i(z^*)\ge\operatorname{margin}_i(x_0)-\delta\varphi_i\ge\delta\varphi_i .
\end{equation*}
Together with $g/\delta\ge1$, this gives $\tau_i\ge1\ge\mathfrak{i}_{\delta}$.

In particular, every margin at $z^*$ is nonnegative, so $z^*\in\operatorname{cl} S$. The
hypotheses of Theorem~\ref{8.17:thm-model-short-move} therefore hold with $x_{\mathrm{ref}}=x_0$
and this $u$, and its rates $\tau_i$ are the ones bounded above. Finally, fix $b\in\mathfrak{B}$.
By the definition of $r_b$ in Theorem~\ref{8.17:thm-model-short-move}, in which the minimum
over the empty set is $g/\delta$, and since $\mathfrak{i}_{\delta}\le1$,
\begin{equation*}
r_b=
\begin{cases}
\delta\min\{\tau_i:b\in\det(i)\}\ge\delta\,\mathfrak{i}_{\delta},
& \text{if } b\in\det(i)\text{ for some } i\in I,\\[2pt]
\delta\cdot(g/\delta)=g\ge\delta\ge\delta\,\mathfrak{i}_{\delta},
& \text{otherwise.}
\end{cases}
\end{equation*}
\end{proof}

\begin{corollary}[Model: the opening constant as an isoperimetric constant]
\label{8.17:cor-model-isoperimetric}
The model is the plaquette instance $\mathbb{M}_{\square}$ of the abelian window model
(Definition~\ref{8.17:def-abelian-window-model}). In it the free bonds $\mathfrak{B}$ of a
finite set of oriented bonds of the cubic lattice carry real variables $y$. Each plaquette
$p\in P$ carries the affine window $F_p$, whose linear part $l_p$ has coefficients
$\varepsilon(p,b)\in\{\pm1\}$ on the free bonds of $p$. We orient all plaquettes lying in one
coordinate plane by one fixed orientation of that plane.

In $\mathbb{M}_{\square}$ let $\mathcal{A}'\subset P$ be a set of plaquettes equipped with
signs $s_p\in\{\pm1\}$. Put $\varepsilon(p,b):=0$ when $b$ is not a bond of $p$. Let
$\partial$ be the boundary map on real chains, $\partial p=\sum_{b}\varepsilon(p,b)\,b$, the
sum running over the four bonds $b$ of $p$. For
$\lambda\in\mathbb{R}^{\mathcal{A}'}_{\ge0}$ let
$c_{\lambda}:=\sum_{p\in\mathcal{A}'}\lambda_p s_p\,p$, a real $2$-chain, and put
\begin{equation*}
\|\partial c_{\lambda}\|_{1,\mathrm{fr}}:=\sum_{b\in\mathfrak{B}}\bigl|(\partial c_{\lambda})_b\bigr|,
\end{equation*}
the sum running over the free bonds. Then
\begin{equation}\label{8.17:eq-model-isoperimetric-1}
\Bigl\|\sum_{p\in\mathcal{A}'}\lambda_p s_p l_p\Bigr\|_1=\|\partial c_{\lambda}\|_{1,\mathrm{fr}},
\end{equation}
so that
\begin{equation}\label{8.17:eq-model-isoperimetric-a}
\mathfrak{i}(\mathcal{A}',s)
=\inf_{\lambda\ge0,\ \lambda\ne0}
\frac{\|\partial c_{\lambda}\|_{1,\mathrm{fr}}}{\sum_{p\in\mathcal{A}'}\lambda_p\varphi_p}.
\end{equation}
\begin{itemize}
\item[(i)] Suppose the plaquettes of $\mathcal{A}'$ have pairwise disjoint sets of free bonds.
Then $\mathfrak{i}(\mathcal{A}',s)=1$: isolated tight windows cost no opening rate.
\item[(ii)] Let $S_r$ be a flat square of side $r$, all of whose bonds are free, and suppose
$\mathcal{A}'$ contains all plaquettes of $S_r$, with one constant sign. Then
$\mathfrak{i}(\mathcal{A}',s)\le 1/r$.
\end{itemize}
\end{corollary}

\begin{proof}
We first prove \eqref{8.17:eq-model-isoperimetric-1}. By
Definition~\ref{8.17:def-abelian-window-model}, $F_p(y)$ is the sum over the bonds $b$ of $p$
of $\varepsilon(p,b)$ times the value on $b$. The held bonds contribute only a constant, so the
linear part of $F_p$ is
\begin{equation*}
l_p(y)=\sum_{b}\varepsilon(p,b)\,y_b ,
\end{equation*}
the sum running over the free bonds $b$ of $p$.
Hence the linear part of $\sum_{p}\lambda_p s_p F_p$ is
\begin{equation*}
y\mapsto\sum_{p\in\mathcal{A}'}\lambda_p s_p\,l_p(y)
=\sum_{b\in\mathfrak{B}}\Bigl(\sum_{p\in\mathcal{A}'}\lambda_p s_p\,\varepsilon(p,b)\Bigr)y_b
=\sum_{b\in\mathfrak{B}}(\partial c_{\lambda})_b\,y_b .
\end{equation*}
The last equality holds because, by linearity of $\partial$ and the formula for $\partial p$,
$(\partial c_{\lambda})_b=\sum_{p}\lambda_p s_p\,\varepsilon(p,b)$. This functional is therefore
the pairing $\langle\partial c_{\lambda},y\rangle$ restricted to the free bonds. Its
$\ell^1$-norm $\sum_{b\in\mathfrak{B}}|(\partial c_{\lambda})_b|$, which is its norm as a
functional on $(V,|\cdot|_\infty)$, equals $\|\partial c_{\lambda}\|_{1,\mathrm{fr}}$. This is
\eqref{8.17:eq-model-isoperimetric-1}.

In $\mathbb{M}_{\square}$ the index set of windows is $P$, so
Proposition~\ref{8.17:prop-opening-constant-duality} applies to $\mathcal{A}'\subset P$ with the
signs $s_p$. By that proposition, $\mathfrak{i}(\mathcal{A}',s)$, defined in
\eqref{8.17:eq-opening-constant-duality-a} through condition (a) of
\eqref{8.17:eq-opening-constant-duality-b}, is the infimum over $\lambda\ge0$, $\lambda\ne0$ of
\begin{equation*}
\Bigl\|\sum_{p}\lambda_p s_p l_p\Bigr\|_1\Big/\Bigl(\sum_{p}\lambda_p\varphi_p\Bigr).
\end{equation*}
Inserting
\eqref{8.17:eq-model-isoperimetric-1} into this infimum gives
\eqref{8.17:eq-model-isoperimetric-a}.

The coefficients of $l_p$ are $\pm1$ on the free bonds of $p$ and vanish elsewhere. Hence
$\det(p)$ is the set of free bonds of $p$, and $\varphi_p=\|l_p\|_1$ is their number. Each
$p\in P$ has at least one free bond, so $\varphi_p\ge1$, and the denominator in
\eqref{8.17:eq-model-isoperimetric-a} is positive for $\lambda\ge0$, $\lambda\ne0$.

(i) Assume the sets $\det(p)$, $p\in\mathcal{A}'$, are pairwise disjoint. Then no two terms of
$\sum_p\lambda_p s_p l_p$ have a common bond, and since $\lambda_p\ge0$ and $|s_p|=1$,
\begin{equation*}
\Bigl\|\sum_{p\in\mathcal{A}'}\lambda_p s_p l_p\Bigr\|_1
=\sum_{p\in\mathcal{A}'}\lambda_p\|l_p\|_1=\sum_{p\in\mathcal{A}'}\lambda_p\varphi_p .
\end{equation*}
Every quotient in the infimum therefore equals $1$, and $\mathfrak{i}(\mathcal{A}',s)=1$.

(ii) Let $s_0\in\{\pm1\}$ be the common value of $s_p$ for the plaquettes $p$ of $S_r$. Set
$\lambda_p:=1$ for the
plaquettes $p$ of $S_r$ and $\lambda_p:=0$ for the other plaquettes of $\mathcal{A}'$. Then
$\lambda\ge0$, $\lambda\ne0$, and $c_{\lambda}=s_0\sum_{p\subset S_r}p$.

An interior bond of $S_r$ lies in exactly two plaquettes of $S_r$, which lie on opposite sides
of the bond in the plane of $S_r$; the plaquettes of $S_r$ carrying the orientation of that
plane fixed above, these two induce opposite orientations on the bond, so its
coefficient in
$\partial\bigl(\sum_{p\subset S_r}p\bigr)$ cancels.
Each of the $4r$ boundary bonds of $S_r$ lies in exactly one plaquette of $S_r$, so its
coefficient is $\pm1$. All these bonds are free, hence
$\|\partial c_{\lambda}\|_{1,\mathrm{fr}}=4r$.

The square $S_r$ has $r^2$ plaquettes, each with four free bonds, so $\varphi_p=4$ for each of
them and $\sum_{p\subset S_r}\varphi_p=4r^2$. The quotient in
\eqref{8.17:eq-model-isoperimetric-a} at this $\lambda$ is therefore $4r/(4r^2)=1/r$, and the
infimum satisfies $\mathfrak{i}(\mathcal{A}',s)\le1/r$.
\end{proof}

Item (ii) is the abelian, infinitesimal form of the statement on opening a flat square of side
$r$ whose windows are tight with a common sign in the non-abelian setting. There, opening such a
square by the relative amount $\tau_{\mathrm{op}}$ forces the displacement on each bond to be
at least $\tau_{\mathrm{op}}\cdot(\text{flux density})\cdot r/4$, the flux density being taken
in the sense of that statement. Here, opening at rate $\tau$ per
unit displacement forces $\tau\le1/r$.

\begin{proposition}[Model: the per-coordinate loss cannot be removed]
\label{8.17:prop-model-loss-sharp}
We work in the abelian window model of Definition~\ref{8.17:def-abelian-window-model}, in the
following instance. Let $n\ge 2$, $\sigma>0$, $\bar\sigma\ge\sigma$, $\psi\in(\sigma,\bar\sigma]$,
$x>0$ and $g:=x^{-1/2}$. Take $\mathfrak{B}:=\{1,\dots,n\}$ and the $n+1$ windows
\begin{equation*}
\ell_k(y):=y_k-y_{k-1}\quad(k=1,\dots,n;\ y_0:=0),\qquad
\ell_{n+1}(y):=n\sigma+\psi-y_n ,
\end{equation*}
with cap $c_k:=\sigma$ for $k\le n$ and cap $c_{n+1}:=\bar\sigma$. Each window reads at most two
coordinates, so $\varphi_i\le 2$. Let the action be
$A(y):=\tfrac12\sum_{k=1}^{n+1}\ell_k(y)^2$, so that $\Lambda_1=4$. Then
$x_0:=(k\sigma)_{k\le n}$ is the unique minimiser of $A$ over $\operatorname{cl} S$; the windows
$\ell_1,\dots,\ell_n$ are tight at $x_0$ with sign $+1$; with $x_{\mathrm{ref}}:=x_0$ one has
$G_b=0$ for $b<n$ and $G_n=\psi-\sigma$; the opening constant $\mathfrak{i}$ of the tight set with
these signs satisfies $\mathfrak{i}\le 1/(2n-1)$; and
\begin{equation}\label{8.17:eq-model-loss-sharp-a}
\int_S e^{-xA(y)}\,dy\;\le\;\sqrt{\pi}\;e^{-xA(x_0)}\,g^n
\Bigl(\frac{2e}{n-1}\Bigr)^{(n-1)/2}.
\end{equation}
Consequently no inequality of the form
\begin{equation}\label{8.17:eq-model-loss-sharp-3}
\int_S e^{-xA(y)}\,dy\;\ge\;e^{-xA(x_0)}\,g^n\,e^{-\mathfrak{K}n},
\end{equation}
with $\mathfrak{K}$ depending only on
$\bigl(\Lambda_1,\,P,\,\max_i\varphi_i,\,\min_i c_i/(g\varphi_i)\bigr)$, holds in the model.
Indeed, choose $\sigma:=\eta g$, $\bar\sigma:=2\eta g$ and $\psi:=3\eta g/2$ with $\eta\ge 3$
fixed; then $P=\eta/2$, and these four quantities stay bounded, while the
loss per coordinate, that is, $-n^{-1}$ times the logarithm of the ratio of
$\int_S e^{-xA(y)}\,dy$ to $e^{-xA(x_0)}g^n$, is at least
\begin{equation*}
\frac{n-1}{2n}\log\frac{n-1}{2e}-\frac{\log\sqrt{\pi}}{n},
\end{equation*}
which is unbounded in $n$.
\end{proposition}

\begin{proof}
\emph{Change of variables.} For $y\in V$ put $v_k:=\sigma-\ell_k(y)$, $k=1,\dots,n$. Summing
$\ell_j(y)=y_j-y_{j-1}$ over $j\le k$ gives $y_k=\sum_{j\le k}\ell_j(y)$, hence
$y_k=k\sigma-\sum_{j\le k}v_j$. The affine map $y\mapsto v$ is therefore bijective; its linear
part is lower triangular with every diagonal entry equal to $-1$, so its Jacobian determinant has
absolute value one. Moreover
$\ell_{n+1}(y)=n\sigma+\psi-y_n=\psi+\sum_{k\le n}v_k$. On $S$ we have $|\ell_k(y)|<\sigma$ for
$k\le n$, that is, $v_k\in(0,2\sigma)$; by continuity $v_k\in[0,2\sigma]$ on $\operatorname{cl} S$.

\emph{Expansion of the action.} Writing $\Sigma v:=\sum_{k\le n}v_k$,
\begin{align}
A&=\tfrac12\sum_{k\le n}(\sigma-v_k)^2+\tfrac12(\psi+\Sigma v)^2\notag\\
&=A(x_0)+(\psi-\sigma)\,\Sigma v+\tfrac12\sum_{k\le n}v_k^2+\tfrac12(\Sigma v)^2,
\label{8.17:eq-model-loss-sharp-1}
\end{align}
where $A(x_0)=\tfrac12 n\sigma^2+\tfrac12\psi^2$ and $x_0=(k\sigma)_{k\le n}$ is the point $v=0$.
Since $\psi>\sigma$, each of the three
correction terms in \eqref{8.17:eq-model-loss-sharp-1} is non-negative for $v\ge 0$, and the term
$\frac12\sum_k v_k^2$ vanishes only at $v=0$. Hence $A(y)>A(x_0)$ for every
$y\in\operatorname{cl} S$ with $y\ne x_0$. It remains to see that $x_0\in\operatorname{cl} S$. The
model requires $S\ne\emptyset$; pick $y\in S$ with coordinates $v\in(0,2\sigma)^n$. For
$\nu\in(0,1]$ the point with coordinates $\nu v$ satisfies $\nu v_k\in(0,2\sigma)$, and
$\psi+\nu\Sigma v$ is positive and at most $\psi+\Sigma v<\bar\sigma$, so it lies in $S$; letting
$\nu\to0$ gives
$x_0\in\operatorname{cl} S$. Thus $x_0$ is the unique minimiser of $A$ over $\operatorname{cl} S$.
The same computation shows $\psi<\psi+\Sigma v<\bar\sigma$, so $\psi<\bar\sigma$. At $x_0$ we have
$\ell_k(x_0)=\sigma=c_k$ for $k\le n$, so $f_k(x_0)=1$ with $\operatorname{sgn}\ell_k(x_0)=+1$,
while $\ell_{n+1}(x_0)=\psi<\bar\sigma$ gives $f_{n+1}(x_0)<1$. The tight set of
Lemma~\ref{8.17:lem-first-order-minimiser}, taken with $I'=I$, is therefore $\{1,\dots,n\}$ with
all signs $+1$.

\emph{Gradient and $\Lambda_1$.} $A$ is a sum of squares of affine functions, hence convex and
continuously differentiable. The coordinate $y_k$ occurs with coefficient $+1$ in $\ell_k$ and with
coefficient $-1$ in $\ell_{k+1}$ (for $k=n$ this is $\ell_{n+1}=n\sigma+\psi-y_n$), and in no other
window. Hence $\partial_{y_k}A=\ell_k-\ell_{k+1}$. At $x_0$ this equals $\sigma-\sigma=0$ for
$k<n$ and $\sigma-\psi$ for $k=n$, so $G_k=0$ for $k<n$ and $G_n=\psi-\sigma$. The function
$\partial_{y_k}A(y)-\partial_{y_k}A(y')$ is the linear form
$2(y_k-y'_k)-(y_{k-1}-y'_{k-1})-(y_{k+1}-y'_{k+1})$ for $k<n$ (with $y_0=y'_0=0$) and
$2(y_n-y'_n)-(y_{n-1}-y'_{n-1})$ for $k=n$. Its coefficients have absolute values summing to at
most $4$, so $|\partial_bA(y)-\partial_bA(y')|\le 4|y-y'|_\infty$ and $\Lambda_1=4$.

\emph{The mass.} Passing to the coordinates $v$, which have unit Jacobian, and discarding the
non-negative terms $(\psi-\sigma)\Sigma v$ and $\frac12\sum_kv_k^2$ in
\eqref{8.17:eq-model-loss-sharp-1}, we obtain, using $x=g^{-2}$,
\begin{align}
\int_S e^{-xA(y)}\,dy
&\le e^{-xA(x_0)}\int_{[0,2\sigma]^n}e^{-x(\Sigma v)^2/2}\,dv\notag\\
&\le e^{-xA(x_0)}\int_{[0,\infty)^n}e^{-(\Sigma v)^2/(2g^2)}\,dv .
\label{8.17:eq-model-loss-sharp-2}
\end{align}
The volume of $\{v\in[0,\infty)^n:\Sigma v\le t\}$ is $t^n/n!$. Hence the image of Lebesgue measure
on $[0,\infty)^n$ under $v\mapsto\Sigma v$ has density $t^{n-1}/(n-1)!$ on $(0,\infty)$, and the
last integral in \eqref{8.17:eq-model-loss-sharp-2} equals
\begin{equation*}
\int_0^\infty e^{-t^2/(2g^2)}\,\frac{t^{n-1}}{(n-1)!}\,dt .
\end{equation*}
Write $t^{n-1}e^{-t^2/(2g^2)}=t^{n-1}e^{-t^2/(4g^2)}\cdot e^{-t^2/(4g^2)}$. The derivative of
$\log\bigl(t^{n-1}e^{-t^2/(4g^2)}\bigr)$ is $(n-1)/t-t/(2g^2)$, so the first factor attains its
supremum over $t>0$ at $t^2=2(n-1)g^2$. There it equals
\begin{equation*}
\bigl(2(n-1)g^2\bigr)^{(n-1)/2}e^{-(n-1)/2}
=\Bigl(\frac{2(n-1)g^2}{e}\Bigr)^{(n-1)/2}.
\end{equation*}
Moreover $\int_0^\infty e^{-t^2/(4g^2)}\,dt=g\sqrt{\pi}$. Finally, since
$e^{n-1}\ge(n-1)^{n-1}/(n-1)!$ (a single term of the exponential series), we have
$(n-1)!\ge((n-1)/e)^{n-1}$. Combining,
\begin{align*}
\int_0^\infty e^{-t^2/(2g^2)}\,\frac{t^{n-1}}{(n-1)!}\,dt
&\le g^n\sqrt{\pi}\,\Bigl(\frac{2(n-1)}{e}\Bigr)^{(n-1)/2}
\Bigl(\frac{e}{n-1}\Bigr)^{n-1}\\
&=g^n\sqrt{\pi}\,\Bigl(\frac{2e}{n-1}\Bigr)^{(n-1)/2},
\end{align*}
which together with \eqref{8.17:eq-model-loss-sharp-2} is
\eqref{8.17:eq-model-loss-sharp-a}.

\emph{The opening constant.} The tight set is $\{1,\dots,n\}$ with signs $s_k=+1$. By
Proposition~\ref{8.17:prop-opening-constant-duality}, the opening constant is the infimum in
\eqref{8.17:eq-opening-constant-duality-a} over $\lambda\ge0$, $\lambda\ne0$. We take
$\lambda_k=1$ for all $k\le n$. The linear parts telescope:
$\sum_{k\le n}l_k$ is the linear form $y\mapsto y_n$, whose $\ell^1$-norm is $1$. On the other hand
$\varphi_1=1$ and $\varphi_k=2$ for $2\le k\le n$, so $\sum_{k\le n}\varphi_k=1+2(n-1)=2n-1$.
Hence $\mathfrak{i}\le1/(2n-1)$.

\emph{The consequence.} With $\sigma=\eta g$, $\bar\sigma=2\eta g$ and $\psi=3\eta g/2$ we have
$\sigma\le\bar\sigma$ and $\sigma<\psi<\bar\sigma$, so every statement above applies. We have
$\Lambda_1=4$, and $\max_i\varphi_i=2$ because $n\ge2$. By the gradient computation $P=G_n/g=\eta/2$.
Also
\begin{equation*}
\min_i\frac{c_i}{g\varphi_i}=\min\Bigl\{\eta,\ \frac{\eta}{2},\ 2\eta\Bigr\}=\frac{\eta}{2}.
\end{equation*}
These four values do not depend on $n$, so any admissible $\mathfrak{K}$ would be one fixed number
for the whole family. Taking logarithms in \eqref{8.17:eq-model-loss-sharp-a} and dividing by $-n$
shows that the loss per coordinate is at least
\begin{equation*}
\frac{n-1}{2n}\log\frac{n-1}{2e}-\frac{\log\sqrt{\pi}}{n},
\end{equation*}
which tends to $+\infty$ as $n\to\infty$. An inequality of the form
\eqref{8.17:eq-model-loss-sharp-3} would bound this loss
by the fixed number $\mathfrak{K}$ for every $n\ge2$, which is impossible.
\end{proof}

\begin{remark}
The thinness that produces \eqref{8.17:eq-model-loss-sharp-a} belongs to the window set itself near
$x_0$: the feasible directions at $x_0$ form a simplex, and this simplex drains through the single
window $\ell_{n+1}$, which is the only window not tight at $x_0$, since
$\ell_{n+1}(x_0)=\psi<\bar\sigma$. Consider a comparison of two
integrals over the same window set in which the measure of that set is kept. Both sides of such a
comparison carry the same thinness, and the comparison does not see it. A lower bound normalised by
$e^{-xA(x_0)}g^n$ alone has no such compensation, and so it must pay the loss.
\end{remark}

\begin{corollary}[Model: the top step in the plaquette instance]\label{8.17:cor-model-top-step}
Work in the plaquette instance $\mathbb{M}_{\square}$ of the abelian window model of
Definition~\ref{8.17:def-abelian-window-model}, so that $f_p=|F_p|/\kappa_p$ and
$A=\frac12\sum_pF_p^2$, the sum running over the plaquettes with at least one free bond.
Let $I=I_{\mathrm{ns}}\sqcup J$, where $J$ is a set of pairwise
distinct seed plaquettes $\square$ whose caps satisfy $\kappa_{\square}\ge 4g\varphi_{\square}$, and put
\begin{align*}
S_{\mathrm{ns}}&:=\{y\in V:\ f_i(y)<1\ \text{for all } i\in I_{\mathrm{ns}}\},\\
S&:=S_{\mathrm{ns}}\cap\{y\in V:\ f_{\square}(y)<1\ \text{for all } \square\in J\},\\
\mathcal{E}&:=\{y\in V:\ f_{\square}(y)\ge 1\ \text{for all } \square\in J\}.
\end{align*}
Assume that the windows are wide at scale $g$ in the sense of
Definition~\ref{8.17:def-abelian-window-model}. Let $x_0$ minimise $A$ over
$\operatorname{cl} S_{\mathrm{ns}}$, and assume
\begin{equation}\label{8.17:eq-model-top-step-1}
f_{\square}(x_0)\le\tfrac12\qquad\text{for all }\square\in J .
\end{equation}
Put $\Gamma_{\square}:=\kappa_{\square}^2/8$. Then
\begin{equation}\label{8.17:eq-model-top-step-2}
A(y)\ge A(x_0)+\sum_{\square\in J}\Gamma_{\square}
\qquad\text{for } y\in\operatorname{cl} S_{\mathrm{ns}}\cap\mathcal{E}.
\end{equation}
Moreover, let $\delta\in(0,g]$ and $u\in V$ with $|u|_\infty\le1$ and $x_0+\delta u\in\operatorname{cl} S$,
and let $r_b$, $\varrho_b$ be the radii and displacements of Theorem~\ref{8.17:thm-model-short-move},
formed for the full family $I$ with $x_{\mathrm{ref}}=x_0$. Then, with $\operatorname{vol}$ the
Lebesgue measure on $V$ and $\operatorname{Den}:=\int_S e^{-xA(y)}\,dy$,
\begin{equation}\label{8.17:eq-model-top-step-3}
\begin{split}
\operatorname{Num}:=\int_{S_{\mathrm{ns}}\cap\mathcal{E}}e^{-xA(y)}\,dy
\le{}&\operatorname{vol}(\operatorname{cl} S_{\mathrm{ns}})\prod_{b\in\mathfrak{B}}(2r_b)^{-1}\\
&\cdot\exp\Bigl(\sum_{b\in\mathfrak{B}}\bigl[xG_b\varrho_b+4\Lambda_1\bigr]
-x\sum_{\square\in J}\Gamma_{\square}\Bigr)\cdot\operatorname{Den}.
\end{split}
\end{equation}
The seeds do not reduce the radii: the opening rates of
Theorem~\ref{8.17:thm-model-short-move} satisfy $\tau_{\square}\ge1$ for every $\square\in J$.
\end{corollary}

\begin{proof}
\emph{The bound \eqref{8.17:eq-model-top-step-2}.} For each such plaquette $p$ the function $F_p$ is
affine with linear part $l_p$, so $F_p(y)=F_p(x_0)+l_p(y-x_0)$ and
\[
\tfrac12F_p(y)^2-\tfrac12F_p(x_0)^2=F_p(x_0)\,l_p(y-x_0)+\tfrac12\,l_p(y-x_0)^2 .
\]
Differentiating $A=\frac12\sum_pF_p^2$ gives $\nabla A(x_0)=\sum_pF_p(x_0)\,l_p$. Summing over $p$,
$A$ is quadratic with
\[
A(y)-A(x_0)=\langle\nabla A(x_0),y-x_0\rangle+\tfrac12\sum_{p}\bigl(F_p(y)-F_p(x_0)\bigr)^2 .
\]
The first term is nonnegative for $y\in\operatorname{cl} S_{\mathrm{ns}}$ by part (a) of
Lemma~\ref{8.17:lem-first-order-minimiser}, applied with $I'=I_{\mathrm{ns}}$, $S'=S_{\mathrm{ns}}$,
since $x_0$ minimises $A$ over $\operatorname{cl} S_{\mathrm{ns}}$. In the second term we keep only
the plaquettes of $J$. For $y\in\mathcal{E}$ we have $|F_{\square}(y)|\ge\kappa_{\square}$, and
\eqref{8.17:eq-model-top-step-1} says $|F_{\square}(x_0)|\le\kappa_{\square}/2$. Hence
$|F_{\square}(y)-F_{\square}(x_0)|\ge\kappa_{\square}/2$ and
\[
\tfrac12\sum_{p}\bigl(F_p(y)-F_p(x_0)\bigr)^2
\ge\tfrac12\sum_{\square\in J}\bigl(F_{\square}(y)-F_{\square}(x_0)\bigr)^2
\ge\tfrac12\sum_{\square\in J}\bigl(\kappa_{\square}-\tfrac12\kappa_{\square}\bigr)^2
=\sum_{\square\in J}\Gamma_{\square}.
\]
This proves \eqref{8.17:eq-model-top-step-2}.

\emph{The bound on $\operatorname{Num}$.} Since $S_{\mathrm{ns}}\cap\mathcal{E}\subset
\operatorname{cl} S_{\mathrm{ns}}\cap\mathcal{E}$ and $x>0$, \eqref{8.17:eq-model-top-step-2} gives
\[
\operatorname{Num}
\le\operatorname{vol}(\operatorname{cl} S_{\mathrm{ns}}\cap\mathcal{E})\,
e^{-x(A(x_0)+\sum_{\square}\Gamma_{\square})}
\le\operatorname{vol}(\operatorname{cl} S_{\mathrm{ns}})\,
e^{-x(A(x_0)+\sum_{\square}\Gamma_{\square})}.
\]

\emph{The bound on $\operatorname{Den}$.} The point $z^*:=x_0+\delta u$ lies in $\operatorname{cl} S$
with $\delta\in(0,g]$ and $|u|_\infty\le1$. Therefore Theorem~\ref{8.17:thm-model-short-move}
applies to the full family $I$ with $x_{\mathrm{ref}}=x_0$, which minimises $A$ over
$\operatorname{cl} S'$ for the sub-family $I'=I_{\mathrm{ns}}$. Its main display
\eqref{8.17:eq-model-short-move-a} gives
\[
\operatorname{Den}\ge e^{-xA(x_0)}\prod_{b\in\mathfrak{B}}(2r_b)\,
\exp\Bigl(-\sum_{b\in\mathfrak{B}}\bigl[xG_b\varrho_b+4\Lambda_1\bigr]\Bigr).
\]
Dividing this bound by $\prod_b(2r_b)\exp(-\sum_b[xG_b\varrho_b+4\Lambda_1])$ bounds
$e^{-xA(x_0)}$ from above. Inserting that bound into the bound on $\operatorname{Num}$ gives
\eqref{8.17:eq-model-top-step-3}.

\emph{The seed rates.} For $i\in I$ and $y$ with $f_i(y)\le1$ the number $c_i-|\ell_i(y)|\ge0$ is
the distance of $|\ell_i(y)|$ from the cap $c_i$, as fixed in
Definition~\ref{8.17:def-abelian-window-model}; Theorem~\ref{8.17:thm-model-short-move} defines the
opening rate $\tau_i$ as the minimum of $g/\delta$ and $(c_i-|\ell_i(z^*)|)/(\delta\varphi_i)$.
In $\mathbb{M}_{\square}$ one has $\ell_{\square}=F_{\square}$ and
$c_{\square}=\kappa_{\square}$. Let $\square\in J$. By \eqref{8.17:eq-model-top-step-1} and the
hypothesis on the caps, using $\delta\le g$,
\[
\kappa_{\square}-|F_{\square}(x_0)|
\ge\tfrac12\kappa_{\square}\ge2g\varphi_{\square}\ge2\delta\varphi_{\square}.
\]
The point $z^*$ lies in $\operatorname{cl} S$, so $f_{\square}(z^*)\le1$ and
$\kappa_{\square}-|F_{\square}(z^*)|\ge0$. Since $\varphi_{\square}$ is the Lipschitz
constant of $F_{\square}$ for $|\cdot|_\infty$, we have
$|F_{\square}(z^*)|\le|F_{\square}(x_0)|+\delta\varphi_{\square}|u|_\infty
\le|F_{\square}(x_0)|+\delta\varphi_{\square}$. Therefore
\[
\kappa_{\square}-|F_{\square}(z^*)|
\ge\kappa_{\square}-|F_{\square}(x_0)|-\delta\varphi_{\square}
\ge\delta\varphi_{\square}.
\]
By the definition of $\tau_{\square}$ in
Theorem~\ref{8.17:thm-model-short-move}, and since $g/\delta\ge1$,
\[
\tau_{\square}=\min\Bigl\{\frac{g}{\delta},\,
\frac{\kappa_{\square}-|F_{\square}(z^*)|}{\delta\varphi_{\square}}\Bigr\}
\ge\min\Bigl\{\frac{g}{\delta},1\Bigr\}=1 .
\qedhere
\]
\end{proof}

On the factor $\operatorname{vol}(\operatorname{cl} S_{\mathrm{ns}})$ in
\eqref{8.17:eq-model-top-step-3}: in the group setting the Haar measure is a probability measure and
the corresponding factor is at most one. In the model $\mathbb{M}_{\square}$ it is the one place
where the a priori range of the coordinates enters. The ratio
$\operatorname{vol}(\operatorname{cl} S_{\mathrm{ns}})/\prod_b(2r_b)$ is the volume deficit of the
bound, measured with the radii $r_b$ in place of balls of radius $g$.

\begin{remark}[Model: the loss per coordinate]\label{8.17:rem-model-loss-per-coordinate}
We work in the abelian affine-window model $\mathbb{M}$, in the setting and notation of
Theorem~\ref{8.17:thm-model-short-move}. The reference point is $x_{\mathrm{ref}}=x_0$, a
minimiser of $A$, and $n$ denotes the number of coordinates $b\in\mathfrak{B}$. A window $i$
is \emph{tight at scale $\delta$} if $\tau_i<g/\delta$. For a coordinate $b$ we put
$\tau_b:=\min\{\tau_i : b\in\det(i)\}$, the least rate among the windows reading $b$.

Suppose a lower bound has the form
\begin{equation}\label{8.17:eq-model-loss-per-coordinate-1}
\int_S e^{-xA(y)}\,dy\;\ge\;e^{-xA(x_0)}\,g^n\,\exp\Bigl(-\sum_{b\in\mathfrak{B}}\eta_b\Bigr).
\end{equation}
We say that it \emph{loses} $\eta_b$ on the coordinate $b$ relative to the reference
$e^{-xA(x_0)}g^n$.

The bound of Theorem~\ref{8.17:thm-model-short-move}\,(ii), in which $\delta=g/\max(1,P)$,
loses at most, on a coordinate $b$ read by some window tight at scale $\delta$, the amount
\begin{equation}\label{8.17:eq-model-loss-per-coordinate-2}
\log\max(1,P)+\log\frac{1}{2\tau_b}+2+4\Lambda_1 .
\end{equation}
The constant $2$ in \eqref{8.17:eq-model-loss-per-coordinate-2} stands for the term
$2\cdot\mathbf{1}\{G_b>0\}$ of that bound replaced by its largest value; with this replacement
the loss is exactly \eqref{8.17:eq-model-loss-per-coordinate-2}.
On every other coordinate it loses at most $4\Lambda_1$.

Thus the bound loses the following.
\begin{itemize}
\item On each coordinate read by a tight window, $\log\max(1,P)$ (and not $\max(1,P)$) is
charged. In the normalisation of the scaled gradient sizes used for the crude bound of this
chapter one has $P\le C'p_0$, where $p_0$ is the value of the degree function $p_0(\cdot)$ at
the coupling $g_{k_o}$ of the top observation level $k_o$ and $C'$ is the constant of that
normalisation; there this charge is at most $\log\max(1,C'p_0)$, not $C'p_0$. The charge falls
on each such coordinate, not only once per tight window.
\item On each coordinate read by no tight window, at most the constant $4\Lambda_1$ is charged.
\item A term $\log(1/\tau_b)$ is charged wherever tight windows interact through shared
coordinates, and $\tau_b=1$ exactly when they do not
(Corollary~\ref{8.17:cor-model-isoperimetric}\,(i)).
\end{itemize}
Relative to the crude majorant $\operatorname{vol}(\operatorname{cl}S)\,e^{-xA(x_0)}$, the
further factor $\operatorname{vol}(\operatorname{cl}S)/g^n$ is incurred on every coordinate,
tight or not. This factor is the full $g$-ball deficit: when $\operatorname{cl}S$ is the
product $\prod_b I_b$ of intervals $I_b$, the a-priori ranges of the coordinates, its
logarithm is the sum over $b$ of the logarithm of the coordinate's a-priori range $|I_b|$
over $g$.
\end{remark}

\begin{proof}
Take $\delta=g/\max(1,P)$ in Theorem~\ref{8.17:thm-model-short-move}. Its clause (ii),
display~\eqref{8.17:eq-model-short-move-b}, gives
\begin{equation*}
\int_S e^{-xA}\;\ge\;e^{-xA(x_0)}\,g^n\prod_{b}\frac{2r_b}{g}\,
\exp\Bigl(-\sum_b\bigl[2\cdot\mathbf{1}\{G_b>0\}+4\Lambda_1\bigr]\Bigr),
\end{equation*}
with $g/(2r_b)=\max(1,P)/(2\tau_b)$.

Write $\prod_b(2r_b/g)=\exp\bigl(-\sum_b\log(g/(2r_b))\bigr)$. The bound then has the form
\eqref{8.17:eq-model-loss-per-coordinate-1}, with $\eta_b$ equal to
$\log(g/(2r_b))+2\cdot\mathbf{1}\{G_b>0\}+4\Lambda_1$.

Let $b$ be read by a window tight at scale $\delta$. Then
\begin{equation*}
\log\frac{g}{2r_b}=\log\max(1,P)+\log\frac{1}{2\tau_b}.
\end{equation*}
The indicator is at most $1$. Replacing it by $1$ weakens the bound: $\eta_b$ is at most the
amount \eqref{8.17:eq-model-loss-per-coordinate-2}, with equality when $G_b>0$.

Let $b$ be read by no tight window. Every window reading $b$ then has $\tau_i=g/\delta$, and
the minimum over the empty set is also $g/\delta$. Hence $r_b=\delta\cdot g/\delta=g$ and
$\log(g/(2r_b))=-\log 2$. Clause (ii) of Theorem~\ref{8.17:thm-model-short-move} states that
the loss on every such coordinate is $4\Lambda_1-\log 2$; compared with the expression for
$\eta_b$ above, this is the assertion of that clause that the term $2\cdot\mathbf{1}\{G_b>0\}$
vanishes on a coordinate read by no tight window. The loss $4\Lambda_1-\log 2$ is at most
$4\Lambda_1$.

The itemised consequences are read off from these values. The first term of
\eqref{8.17:eq-model-loss-per-coordinate-2} is a logarithm and is present on each coordinate
read by a tight window. In the normalisation $P\le C'p_0$ recalled in the first item it is at
most $\log\max(1,C'p_0)$, since the logarithm is increasing. On coordinates read by no tight
window the loss is bounded by the constant $4\Lambda_1$.

The rate enters only through $\log(1/(2\tau_b))$. In the plaquette instance
$\mathbb{M}_{\square}$, Corollary~\ref{8.17:cor-model-isoperimetric}\,(i) states that if the
plaquettes of a signed tight set $\mathcal{A}'$ have pairwise disjoint sets of free bonds, its
opening constant $\mathfrak{i}(\mathcal{A}',s)$ equals $1$: isolated tight windows cost no rate.

Finally, the reference $e^{-xA(x_0)}g^n$ and the crude majorant
$\operatorname{vol}(\operatorname{cl}S)\,e^{-xA(x_0)}$ differ by the factor
$\operatorname{vol}(\operatorname{cl}S)/g^n$. This factor does not depend on which windows are
tight, so it is incurred on every coordinate. When $\operatorname{cl}S=\prod_bI_b$ one has
$\operatorname{vol}(\operatorname{cl}S)=\prod_b|I_b|$, and therefore
\begin{equation*}
\log\frac{\operatorname{vol}(\operatorname{cl}S)}{g^n}=\sum_{b\in\mathfrak{B}}\log\frac{|I_b|}{g},
\end{equation*}
one term per coordinate.
\end{proof}

\begin{definition}[The opening-rate condition at a comparison point]\label{8.17:def-opening-rate-condition}
Fix the merged top-level region $\mathfrak{z}$, $\lambda$ and the configuration $e$ as in
Notation~\ref{8.17:not-top-step-notation}. Fix also a number $\theta\in[1,2]$ and a reference point
$x_{\mathrm{ref}}\in\operatorname{cl}S^{\theta}_{\mathrm{all}}(e)$. In every application below,
$x_{\mathrm{ref}}$ is the comparison point $x_0$ of the upper bound for the top step
(Notation~\ref{8.17:not-top-step-notation}). Let
$\delta\in(0,g]$, where $g=g_{k_o}$.

The \emph{opening-rate condition} with step $\delta$ and rates $\tau=(\tau_a)$, for $\mathfrak{z}$,
$e$, $x_{\mathrm{ref}}$ and a configuration $z^*\in X$, is the requirement that there be numbers
$\tau_a\in(0,1]$, $a\in W\setminus\mathfrak{S}$, such that
\begin{equation}\label{8.17:eq-opening-rate-condition-a}
\begin{gathered}
\sup_{b\in\mathfrak{B}} d\bigl(z^*_b,x_{\mathrm{ref},b}\bigr)\le\delta,\\
f_a(e,z^*)\le 1-\tau_a\cdot\operatorname{Lip}\cdot\delta\qquad\text{for every }a\in W\setminus\mathfrak{S}.
\end{gathered}
\end{equation}
The seed windows $f_{\square}$ are among the windows $a\in W\setminus\mathfrak{S}$ in
\eqref{8.17:eq-opening-rate-condition-a}. The numbers $\tau_a$ are the \emph{opening rates}.

Given such data, we assign to each bond $b\in\mathfrak{B}$ a rate and a radius:
\begin{equation}\label{8.17:eq-opening-rate-condition-1}
\tau_b:=\min\{\tau_a:\ a\in W\setminus\mathfrak{S},\ b\in\det(a)\},\qquad
r_b:=\tfrac12\,\delta\,\tau_b\le\tfrac12 .
\end{equation}
Here $\tau_b:=1$ if no window $a\in W\setminus\mathfrak{S}$ has $b\in\det(a)$.
\end{definition}

Let $x_{\mathrm{ref}}=x_0$, so that the seed bound at the comparison point,
$f_{\square}(e,x_{\mathrm{ref}})\le\frac12$ for every seed $\square$, is available. We consider the
one-sided condition
\begin{equation}\label{8.17:eq-opening-rate-condition-b}
4K_wC_{\delta_0}\,g_{k_o}\le\varepsilon_{\min}
\end{equation}
on the coupling at the top observation level, equivalently $4\operatorname{Lip}\cdot g_{k_o}\le1$ with
$\operatorname{Lip}=K_wC_{\delta_0}/\varepsilon_{\min}$ the window response constant of
Notation~\ref{8.17:not-top-step-notation}.

\emph{Claim.} The rates $\tau_b$ lie in $(0,1]$ and $0<r_b\le\frac12\delta$. Under
\eqref{8.17:eq-opening-rate-condition-b}, for every $z^*\in X$ satisfying the first line of
\eqref{8.17:eq-opening-rate-condition-a} and every seed $\square$, one has
$f_{\square}(e,z^*)\le1-\operatorname{Lip}\cdot\delta$; that is, every seed window satisfies the
inequality of \eqref{8.17:eq-opening-rate-condition-a} with rate $\tau_{\square}=1$.

\begin{proof}[Proof of the claim]
\emph{The rates and radii.} The family $W$ is finite and every $\tau_a$ lies in $(0,1]$. Hence the
minimum in \eqref{8.17:eq-opening-rate-condition-1} exists and lies in $(0,1]$; the default value $1$
also lies there. With $\delta\le g$ this gives
\begin{equation*}
0<r_b\le\tfrac12\delta\le\tfrac12 g ,
\end{equation*}
and hence $r_b\le\frac12$, the bound recorded in \eqref{8.17:eq-opening-rate-condition-1},
whenever $g_{k_o}\le1$.

\emph{The seed windows.} Let $\square$ be a seed and let $z^*\in X$ satisfy
$\sup_{b\in\mathfrak{B}}d(z^*_b,x_{\mathrm{ref},b})\le\delta$. Since $x_{\mathrm{ref}}=x_0$, the seed
bound at the comparison point gives $f_{\square}(e,x_{\mathrm{ref}})\le\frac12$.

We apply the linear window-response bound to the window $\square$ at $z=x_{\mathrm{ref}}=x_0$ and
$x'=z^*$. By the first line of \eqref{8.17:eq-opening-rate-condition-a}, the bondwise displacement
satisfies
\begin{equation*}
\eta:=\sup_{b\in\mathfrak{B}}d\bigl(z^*_b,x_{\mathrm{ref},b}\bigr)\le\delta\le g .
\end{equation*}
This is the range in which the bound is applied at the comparison point in
Notation~\ref{8.17:not-top-step-notation}: the room $\rho_1=2\operatorname{Lip}\cdot g$ fixed there is
the room that this bound gives at every configuration within $g$ of $x_0$. Hence
\begin{equation*}
f_{\square}(e,z^*)\le\tfrac12+\operatorname{Lip}\cdot\delta .
\end{equation*}

Recall from Notation~\ref{8.17:not-top-step-notation} that
$\operatorname{Lip}=K_wC_{\delta_0}/\varepsilon_{\min}$, and recall $\delta\le g_{k_o}$. Then
\eqref{8.17:eq-opening-rate-condition-b} gives
\begin{equation*}
\operatorname{Lip}\cdot\delta\le\frac{K_wC_{\delta_0}\,g_{k_o}}{\varepsilon_{\min}}\le\frac14 .
\end{equation*}
Thus $2\operatorname{Lip}\cdot\delta\le\frac12$, and therefore
\begin{equation*}
f_{\square}(e,z^*)\le\tfrac12+\operatorname{Lip}\cdot\delta\le1-\operatorname{Lip}\cdot\delta .
\end{equation*}
This is the inequality of \eqref{8.17:eq-opening-rate-condition-a} for $a=\square$ with
$\tau_{\square}=1$.
\end{proof}

Take $\delta=g$ and $z^*=x_{\mathrm{ref}}=x_0$, and recall $\rho_1=2\operatorname{Lip}\cdot g$
(Notation~\ref{8.17:not-top-step-notation}). The inequality
\begin{equation*}
f_a(e,x_0)\le1-2\operatorname{Lip}\cdot g=1-\rho_1
\end{equation*}
is the formal case $\tau_a=2$ of the inequality of \eqref{8.17:eq-opening-rate-condition-a}; it is
the hypothesis of the single-ball device. When it is imposed at every window
$a\in W\setminus\mathfrak{S}$, it implies the case $\tau_a=1$ and so yields
\eqref{8.17:eq-opening-rate-condition-a} with all rates equal to $1$.

The opening-rate condition is weaker. It does not ask for room at $x_{\mathrm{ref}}$: it allows
$x_{\mathrm{ref}}$ to be tight at a window $a$, that is, it does not require
$f_a(e,x_{\mathrm{ref}})\le1-\rho_1$. It asks only that $a$ be opened at rate $\tau_a$ by the short
move from $x_{\mathrm{ref}}$ to $z^*$, that is, that
$f_a(e,z^*)\le1-\tau_a\cdot\operatorname{Lip}\cdot\delta$. A small rate $\tau_a$ has a price, which
is displayed in the bound derived from the condition in what follows.

The straddling windows $a\in\mathfrak{S}$ are not subject to the condition. In the relaxed integral
$\operatorname{Den}^{\Lambda'}_{\lambda}(e)$ they are relaxed to the cap $\Lambda'$. The changes of
their level-$k_o$ labels are then summed by the bound on the straddling windows used with the relaxed
integral. Alternatively, one may include $\mathfrak{S}$ in \eqref{8.17:eq-opening-rate-condition-a}
and take $\Lambda'=1$. The comparison is then pointwise, and no label changes.

\begin{lemma}[The crude short-move device]\label{8.17:lem-crude-short-move}
Let $\mathfrak{z}$ be a merged top-level raw region in the sense of
Notation~\ref{8.17:not-top-step-notation}, with any number $n_s$ of seed cubes, any diameter, any
lateral extent, any inradius and any shape. Let $e$ be an admissible crude-regular exterior,
$\lambda$ an admissible content, $\theta\in[1,2]$ and $x_{\mathrm{ref}}\in\operatorname{cl}
S^{\theta}_{\mathrm{all}}(e)$. Let $g=g_{k_o}\le\min(1,2\alpha_P/3)$ and $\delta\in(0,g]$. Here
$\alpha_P$, $C'$, $C$, $\alpha_0$ and $\operatorname{Lip}$ are the constants fixed in
Notation~\ref{8.17:not-top-step-notation}. Assume
the following statements (i)--(vi).
\begin{itemize}
\item[(i)] \emph{The window-response lemma, linear form.} Let $a\in W$ be a window, and let $z$
be a configuration in the crude-regular class with $f_{a'}(e,z)\le 2$ for every non-seed window
$a'\in W$. Let
$x'\in X$ satisfy $\eta:=\sup_{b\in\mathfrak{B}}d(x'_b,z_b)\le\alpha_P$. Then
\[
f_a(e,x')\le f_a(e,z)+\operatorname{Lip}\cdot\eta .
\]
\item[(ii)] \emph{Lipschitz bound for the background action.} Here $\mathbf{A}(e,\cdot)$ is the
background action of the minimal configuration for the nested data of
\cite[Theorem~1]{B-CMP102b}. Let $z\in\operatorname{cl}
S^{2}_{\mathrm{all}}(e)$, and let $x'\in X$ satisfy $\sup_{b}d(x'_b,z_b)\le\alpha_P$; the sum
below is finite, since $\mathfrak{B}$ is finite. Then
$x'$ is crude-regular and
\[
|\mathbf{A}(e,x')-\mathbf{A}(e,z)|\le C'\varepsilon\sum_{b\in\mathfrak{B}}d(x'_b,z_b).
\]
\item[(iii)] \emph{Lipschitz and oscillation bound for the localised small terms.} For $z$ and
$x'$ as in (ii),
\[
|\Psi(e,x')-\Psi(e,z)|\le (C/\alpha_0)\sum_{b\in\mathfrak{B}}d(x'_b,z_b).
\]
Moreover $\sup\Psi-\inf\Psi$ over $\operatorname{cl}S^{2}_{\mathrm{all}}(e)$ is at most
$\Omega_{\Psi}(\mathfrak{z})$.
\item[(iv)] \emph{Haar measure in the exponential chart.} In the exponential chart
about any point of $G$, normalised Haar measure has a density $j(Y)$ with
respect to a fixed translation-invariant measure $dY$ on $\mathfrak{g}$. This density satisfies
$j(0)=1$, $j\le 1$, and $j(Y)\ge e^{-c_G|Y|^2}$ for $|Y|\le 1$. Moreover
$\int_{|Y|\le r}dY\ge (r/\iota_G)^{d(\mathfrak{g})}$ for $0<r\le\iota_G$.
\item[(v)] \emph{The comparison point on the closed window set.} For $\theta'\in[1,2]$ the
function $\mathbf{A}(e,\cdot)$
is continuous on the compact set $\operatorname{cl}S^{\theta'}_{\mathrm{all}}(e)$ and attains its
infimum there. Every point of $\operatorname{cl}S^{\theta'}_{\mathrm{all}}(e)$ is crude-regular.
\item[(vi)] \emph{Opening rates.} Let $z^*\in X$ be a configuration with
$\sup_{b\in\mathfrak{B}}d(z^*_b,x_{\mathrm{ref},b})\le\delta$, and let
$\tau_a\in(0,1]$, $a\in W\setminus\mathfrak{S}$, be numbers such that the opening-rate inequality
\eqref{8.17:eq-opening-rate-condition-a} holds, namely
\[
f_a(e,z^*)\le 1-\tau_a\operatorname{Lip}\delta\qquad\text{for every }a\in W\setminus\mathfrak{S}
\]
(seed windows included); that is, the opening-rate condition of
Definition~\ref{8.17:def-opening-rate-condition} with parameters $\delta$ and
$\tau=(\tau_a)_{a\in W\setminus\mathfrak{S}}$ holds at $(\mathfrak{z},e;x_{\mathrm{ref}},z^*)$.
The rates $\tau_b$ and the radii $r_b=\tfrac12\delta\tau_b\le\tfrac12$,
$b\in\mathfrak{B}$, are those of Definition~\ref{8.17:def-opening-rate-condition}.
\end{itemize}
Hypothesis (i) is the window-response lemma in its linear form, with its condition on the base
imposed on the non-seed windows only; it is used here in that form.
Put $\Lambda':=\theta+\tfrac32\operatorname{Lip}\delta$ and
\begin{equation}\label{8.17:eq-crude-short-move-a}
\begin{split}
E_{\delta}(\mathfrak{z},e):={}&\sum_{b\in\mathfrak{B}}d(\mathfrak{g})
\log\frac{2\iota_G}{\delta\tau_b}\\
&+n_b\Bigl[\frac{c_G}{4}+\frac32\,\frac{C'\varepsilon\delta}{g^2}
+\frac{3C}{2\alpha_0}\,\delta\Bigr]+\Omega_{\Psi}(\mathfrak{z}).
\end{split}
\end{equation}
Then
\begin{equation}\label{8.17:eq-crude-short-move-b}
\operatorname{Den}^{\Lambda'}_{\lambda}(e)\ge \exp\bigl(-E_{\delta}(\mathfrak{z},e)\bigr)\cdot
\exp\Bigl(-\frac{\mathbf{A}(e,x_{\mathrm{ref}})}{g^2}
+\sup_{\operatorname{cl}S^{\theta}_{\mathrm{all}}(e)}\Psi\Bigr).
\end{equation}
With the choice $\delta:=g/\max(1,C'p_0(g_{k_o}))$ one has
\begin{equation}\label{8.17:eq-crude-short-move-c}
\begin{split}
E_{\delta}\le{}&\sum_{b\in\mathfrak{B}}d(\mathfrak{g})\Bigl[\log\frac{2\iota_G}{g_{k_o}}
+\log\max\bigl(1,C'p_0(g_{k_o})\bigr)+\log\frac{1}{\tau_b}\Bigr]\\
&+n_b\Bigl[\frac{c_G}{4}+\frac32+\frac{3C}{2\alpha_0}\Bigr]+\Omega_{\Psi}(\mathfrak{z}).
\end{split}
\end{equation}
\end{lemma}

\begin{proof}
Throughout, $f_a(x)$ abbreviates $f_a(e,x)$. We use two facts about the set families. First,
$\theta\le 2$ gives $S^{\theta}_{\mathrm{all}}(e)\subset S^{2}_{\mathrm{all}}(e)$, hence
$\operatorname{cl}S^{\theta}_{\mathrm{all}}(e)\subset\operatorname{cl}S^{2}_{\mathrm{all}}(e)$.
Second, the window functionals are continuous on $X$. Hence every point of
$\operatorname{cl}S^{\theta}_{\mathrm{all}}(e)$ satisfies $f_a\le\theta$ for each non-seed
window $a$.

Step~1: the point $z^*$ is crude-regular and all its windows are at most $2$. By (v), applied
with $\theta'=\theta$, the point $x_{\mathrm{ref}}$ is crude-regular. By the preceding paragraph,
$f_a(x_{\mathrm{ref}})\le\theta\le 2$ for every non-seed window $a$; hence $x_{\mathrm{ref}}$
satisfies the conditions on the base $z$ in (i).

Consider first $a\in\mathfrak{S}$. We apply (i) with base $z=x_{\mathrm{ref}}$ and $x'=z^*$. By
(vi),
\[
\eta=\sup_b d(z^*_b,x_{\mathrm{ref},b})\le\delta\le g\le\tfrac23\alpha_P\le\alpha_P .
\]
Hence $f_a(z^*)\le\theta+\operatorname{Lip}\delta$ for $a\in\mathfrak{S}$, and this is all that
is needed for $\mathfrak{S}$. For $a\in W\setminus\mathfrak{S}$, hypothesis (vi) itself gives
$f_a(z^*)\le 1$.

No bound on the seed windows at $x_{\mathrm{ref}}$ is needed: the condition on the base in (i)
concerns only the non-seed windows, and (vi) bounds the seed windows at $z^*$ directly. Since
$x_{\mathrm{ref}}\in\operatorname{cl}S^{2}_{\mathrm{all}}(e)$ and
$\sup_b d(z^*_b,x_{\mathrm{ref},b})\le\delta\le\alpha_P$, the regularity clause of (ii) shows
that $z^*$ is crude-regular.

Step~2: the product set. For $b\in\mathfrak{B}$ put
\[
B_b:=\{z^*_b\exp(Y):Y\in\mathfrak{g},\ |Y|\le r_b\},\qquad
\mathfrak{K}:=\prod_{b\in\mathfrak{B}}B_b\subset X .
\]
Let $x'\in\mathfrak{K}$, so that $x'_b=z^*_b\exp(Y_b)$ with $|Y_b|\le r_b$. The distance $d$ is
bi-invariant, so $d(x'_b,z^*_b)=d(\exp(Y_b),1)$. The curve $t\mapsto\exp(tY_b)$, $t\in[0,1]$,
has length $|Y_b|$, hence
\[
d(x'_b,z^*_b)\le|Y_b|\le r_b\le\tfrac12\delta .
\]
By the triangle inequality and (vi),
\[
d(x'_b,x_{\mathrm{ref},b})\le\tfrac12\delta+\delta=\tfrac32\delta\le\tfrac32 g\le\alpha_P ,
\]
where the last inequality is $g\le 2\alpha_P/3$.

Step~3: $\mathfrak{K}\subset S^{1,\Lambda'}(e)$. Let $x'\in\mathfrak{K}$.

Consider first $a\in W\setminus\mathfrak{S}$. For every bond $b$ read by $a$, that is
$b\in\det(a)$, we have $\tau_b\le\tau_a$ by the definition of $\tau_b$. Hence
$r_b\le\tfrac12\delta\tau_a$. Define $x''\in X$ by $x''_b:=x'_b$ for $b\in\det(a)$ and
$x''_b:=z^*_b$ otherwise. Then $f_a(x')=f_a(x'')$, because $f_a$ reads only the bonds in
$\det(a)$, and
\[
\sup_b d(x''_b,z^*_b)\le\tfrac12\delta\tau_a\le\alpha_P .
\]
By Step~1 the point $z^*$ satisfies the conditions on the base $z$ in (i). Applying (i) with
base $z^*$ to $x''$,
and then (vi), gives
\[
f_a(x')\le f_a(z^*)+\operatorname{Lip}\cdot\tfrac12\delta\tau_a
\le 1-\tfrac12\tau_a\operatorname{Lip}\delta<1 .
\]
This covers the seed windows and the non-seed windows outside $\mathfrak{S}$.

Consider next $a\in\mathfrak{S}$. Step~2 gives
\[
\sup_b d(x'_b,x_{\mathrm{ref},b})\le\tfrac32\delta\le\alpha_P .
\]
By (i) with base $x_{\mathrm{ref}}$, which satisfies the conditions on the base $z$ in (i) by
Step~1,
\[
f_a(x')\le f_a(x_{\mathrm{ref}})+\tfrac32\operatorname{Lip}\delta
\le\theta+\tfrac32\operatorname{Lip}\delta=\Lambda' .
\]
The strictness required in $S^{1,\Lambda'}(e)$ is immaterial: either the level set
$\{f_a=\Lambda'\}$ of the continuous function $f_a$ is accounted for by reading $\Lambda'$ as
$\Lambda'+0$, or $\Lambda'$ is enlarged infinitesimally; we keep the letter $\Lambda'$.

Step~4: Haar measure of $\mathfrak{K}$. We have
\[
r_b\le\tfrac12\delta\le\tfrac12 g\le\tfrac12,
\]
so the lower bound on the density in (iv) is used on $|Y|\le 1$. Using the chart about $z^*_b$ and
$r_b=\tfrac12\delta\tau_b\le\tfrac12\delta$, we obtain that the normalised Haar measure of $B_b$
in $G$ equals $\int_{|Y|\le r_b}j(Y)\,dY$, and
\[
\int_{|Y|\le r_b}j(Y)\,dY\ge e^{-c_Gr_b^2}\Bigl(\frac{r_b}{\iota_G}\Bigr)^{d(\mathfrak{g})}
\ge e^{-c_G\delta^2/4}\Bigl(\frac{\delta\tau_b}{2\iota_G}\Bigr)^{d(\mathfrak{g})}.
\]
The measure $dH$ is normalised Haar measure on $X=G^{\mathfrak{B}}$, that is, the product of
normalised Haar measures on the factors. Hence $\int_{\mathfrak{K}}dH$ is the product over
$b\in\mathfrak{B}$ of the normalised Haar measures of the sets $B_b$.

Step~5: weights on $\mathfrak{K}$. Let $x'\in\mathfrak{K}$. We have
$x_{\mathrm{ref}}\in\operatorname{cl}S^{2}_{\mathrm{all}}(e)$, and Step~2 gives, for every $b$,
\[
d(x'_b,x_{\mathrm{ref},b})\le\tfrac32\delta\le\alpha_P .
\]
Hypothesis (ii) with base
$x_{\mathrm{ref}}$ therefore gives
\[
\mathbf{A}(e,x')\le\mathbf{A}(e,x_{\mathrm{ref}})+C'\varepsilon\sum_{b}d(x'_b,x_{\mathrm{ref},b})
\le\mathbf{A}(e,x_{\mathrm{ref}})+\tfrac32 C'\varepsilon\delta\,n_b .
\]
Hypothesis (iii) with the same base gives
\[
\Psi(e,x')\ge\Psi(e,x_{\mathrm{ref}})-\frac{3C}{2\alpha_0}\delta n_b .
\]
The oscillation clause of
(iii) and the inclusion $x_{\mathrm{ref}}\in\operatorname{cl}S^{\theta}_{\mathrm{all}}(e)\subset
\operatorname{cl}S^{2}_{\mathrm{all}}(e)$ give
\[
\Psi(e,x_{\mathrm{ref}})\ge\inf_{\operatorname{cl}S^{2}_{\mathrm{all}}(e)}\Psi
\ge\sup_{\operatorname{cl}S^{2}_{\mathrm{all}}(e)}\Psi-\Omega_{\Psi}(\mathfrak{z})
\ge\sup_{\operatorname{cl}S^{\theta}_{\mathrm{all}}(e)}\Psi-\Omega_{\Psi}(\mathfrak{z}).
\]
Hence
\[
\Psi(e,x')\ge\sup_{\operatorname{cl}S^{\theta}_{\mathrm{all}}(e)}\Psi-\Omega_{\Psi}(\mathfrak{z})
-\frac{3C}{2\alpha_0}\delta n_b .
\]

Step~6: integration. By definition,
$\operatorname{Den}^{\Lambda'}_{\lambda}(e)$ is the integral of $e^{-\mathbf{A}/g^2+\Psi}$ over
$S^{1,\Lambda'}(e)$ against $dH$. The integrand is positive and
$\mathfrak{K}\subset S^{1,\Lambda'}(e)$ by Step~3. Step~5 then gives
\[
\begin{split}
\operatorname{Den}^{\Lambda'}_{\lambda}(e)\ge{}&\int_{\mathfrak{K}}e^{-\mathbf{A}/g^2+\Psi}\,dH\\
\ge{}&\Bigl(\int_{\mathfrak{K}}dH\Bigr)\exp\Bigl(-\frac{\mathbf{A}(e,x_{\mathrm{ref}})}{g^2}
-\frac32\,\frac{C'\varepsilon\delta n_b}{g^2}
+\sup_{\operatorname{cl}S^{\theta}_{\mathrm{all}}(e)}\Psi-\Omega_{\Psi}(\mathfrak{z})
-\frac{3C}{2\alpha_0}\delta n_b\Bigr).
\end{split}
\]
Now insert Step~4. Since $\delta^2\le g^2\le 1$ and the product has $n_b$ factors,
\[
\int_{\mathfrak{K}}dH\ge\prod_{b\in\mathfrak{B}}e^{-c_G\delta^2/4}
\Bigl(\frac{\delta\tau_b}{2\iota_G}\Bigr)^{d(\mathfrak{g})}
\ge\exp\Bigl(-\sum_{b\in\mathfrak{B}}d(\mathfrak{g})\log\frac{2\iota_G}{\delta\tau_b}
-n_b\frac{c_G}{4}\Bigr).
\]
The two displays together give \eqref{8.17:eq-crude-short-move-b}, with $E_{\delta}$ as in
\eqref{8.17:eq-crude-short-move-a}.

Finally, let $\delta=g/\max(1,C'p_0)$ with $p_0=p_0(g_{k_o})$. Since $\max(1,C'p_0)\ge 1$, this
$\delta$ lies in $(0,g]$. By the normalisation $\varepsilon g/g^2=p_0$ of the small-field
threshold $\varepsilon=\varepsilon_{k_o}$ fixed in Notation~\ref{8.17:not-top-step-notation},
\[
\frac32\,\frac{C'\varepsilon\delta}{g^2}=\frac32\,\frac{C'\varepsilon g/g^2}{\max(1,C'p_0)}
=\frac32\,\frac{C'p_0}{\max(1,C'p_0)}\le\frac32 .
\]
Since $\delta\le g\le 1$, we also have $\tfrac{3C}{2\alpha_0}\delta\le\tfrac{3C}{2\alpha_0}$.
Moreover,
\[
\log\frac{2\iota_G}{\delta\tau_b}=\log\frac{2\iota_G}{g}+\log\max(1,C'p_0)+\log\frac1{\tau_b}.
\]
Inserting these three relations into \eqref{8.17:eq-crude-short-move-a} gives
\eqref{8.17:eq-crude-short-move-c}.
\end{proof}

\begin{corollary}[The top step from the crude short move]\label{8.17:cor-top-step-crude}
Let $\mathfrak{z}$, $e$, $\lambda$, $\theta\in[1,2]$, $g=g_{k_o}$ and $\delta\in(0,g]$ be as in the
crude short-move device (Lemma~\ref{8.17:lem-crude-short-move}), and let $\Lambda'$ be as defined
there. Assume:
\begin{itemize}
\item[(a)] the hypotheses of Lemma~\ref{8.17:lem-crude-short-move} hold with
$x_{\mathrm{ref}}:=x_0$, where $x_0$ is a minimiser of $\mathbf{A}(e,\cdot)$ over
$\operatorname{cl}S^{\theta}_{\mathrm{all}}(e)$, the comparison point of the upper bound in~(b);
\item[(b)] the upper bound on the sourced fibre integral at the comparison point $x_0$:
\begin{equation*}
\operatorname{Num}_{\lambda}(e)\le \Bigl(\int_{S(e)\cap\mathcal{E}}dH\Bigr)\,
\exp\Bigl(-\frac{\mathbf{A}(e,x_0)+\sum_{\square}\Gamma_{\square}}{g^2}
+\sup_{\operatorname{cl}S^{\theta}_{\mathrm{all}}(e)}\Psi\Bigr),
\end{equation*}
where $dH$ is the normalised Haar measure on $X$, $\Gamma_{\square}$ is the income of the
seed $\square$, and the sum runs over the seeds $\square\subset\mathfrak{z}$;
\item[(c)] for some configuration $z^*\in X$ and some rates $\tau=(\tau_a)$, the opening-rate
condition $\mathcal{A}_{\delta,\tau}(\mathfrak{z},e;x_0,z^*)$ of
Definition~\ref{8.17:def-opening-rate-condition} holds, that is,
\eqref{8.17:eq-opening-rate-condition-a} with $x_{\mathrm{ref}}=x_0$; with these rates let
$E_{\delta}(\mathfrak{z},e)$ be as in \eqref{8.17:eq-crude-short-move-a}.
\end{itemize}
Then, pointwise in $e$,
\begin{equation}\label{8.17:eq-top-step-crude-a}
\operatorname{Num}_{\lambda}(e)\le
\exp\Bigl(E_{\delta}(\mathfrak{z},e)-\frac{\sum_{\square}\Gamma_{\square}}{g^2}\Bigr)\,
\operatorname{Den}^{\Lambda'}_{\lambda}(e).
\end{equation}
For a class $[\lambda_0]^{s}$ of admissible contents $\lambda$, with non-negative weights
$F_{\lambda}$ in $e$ against which the fibre integrals of the contents are integrated, let
$[\lambda_0]^{\mathrm{str}}$ denote the class of contents $\lambda'$ obtained from it by
re-indexing the level-$k_o$ labels at the straddling windows, with weights $F_{\lambda'}$.
Assume moreover:
\begin{itemize}
\item[(d)] the inequality
\begin{equation}\label{8.17:eq-top-step-crude-b}
\theta-1+\frac{3}{2}\,\frac{K_wC_{\delta_0}g_{k_o}}{\varepsilon_{\min}}\le\Lambda_L-1,
\end{equation}
where $K_wC_{\delta_0}/\varepsilon_{\min}=\operatorname{Lip}$ is the window response constant;
\item[(e)] the re-indexing at the straddling windows: whenever the relaxed cap satisfies
$\Lambda'\le\Lambda_L$, re-indexing the level-$k_o$ labels at the straddling windows, with no
change of variables in $e$, yields
\begin{equation*}
\sum_{\lambda\in[\lambda_0]^{s}}\int F_{\lambda}\operatorname{Den}^{\Lambda'}_{\lambda}
\le
\sum_{\lambda'\in[\lambda_0]^{\mathrm{str}}}\int F_{\lambda'}\operatorname{Den}_{\lambda'}.
\end{equation*}
\end{itemize}
Then, class by class,
\begin{equation}\label{8.17:eq-top-step-crude-1}
\sum_{\lambda\in[\lambda_0]^{s}}\int F_{\lambda}\operatorname{Num}_{\lambda}
\le\exp\Bigl(\sup_{e}E_{\delta}-\frac{\sum_{\square}\Gamma_{\square}}{g^2}\Bigr)
\sum_{\lambda'\in[\lambda_0]^{\mathrm{str}}}\int F_{\lambda'}\operatorname{Den}_{\lambda'}.
\end{equation}
In~(c) the inequalities at the seed windows may be replaced by the hypotheses
$f_{\square}(e,x_0)\le\tfrac12$ at every seed $\square$ and
$4K_wC_{\delta_0}g_{k_o}\le\varepsilon_{\min}$.
\end{corollary}

\begin{proof}
If the inequalities of~(c) at the seed windows are replaced by the last two hypotheses of the
statement, they are recovered with $\tau_{\square}=1$. Indeed, the Lipschitz bound of the window
functionals gives $f_{\square}(e,z)\le f_{\square}(e,x)+\operatorname{Lip}\,\sup_{b}d(z_b,x_b)$,
and $\sup_{b\in\mathfrak{B}}d(z^*_b,x_{0,b})\le\delta$ by~(c); hence
\begin{equation*}
f_{\square}(e,z^*)\le f_{\square}(e,x_0)+\operatorname{Lip}\,\delta
\le\tfrac12+\operatorname{Lip}\,\delta\le1-\operatorname{Lip}\,\delta,
\end{equation*}
the last inequality because $\delta\le g_{k_o}$ gives
\begin{equation*}
4\operatorname{Lip}\,\delta\le\frac{4K_wC_{\delta_0}g_{k_o}}{\varepsilon_{\min}}\le1 .
\end{equation*}
So~(c) holds at every seed with $\tau_{\square}=1$, and we may argue with~(c) as stated.

Since $dH$ is a probability measure on $X$, $\int_{S(e)\cap\mathcal{E}}dH\le 1$, and
hypothesis~(b) gives
\begin{equation}\label{8.17:eq-top-step-crude-2}
\operatorname{Num}_{\lambda}(e)\le
\exp\Bigl(-\frac{\mathbf{A}(e,x_0)+\sum_{\square}\Gamma_{\square}}{g^2}
+\sup_{\operatorname{cl}S^{\theta}_{\mathrm{all}}(e)}\Psi\Bigr).
\end{equation}
By~(a) and~(c) the hypotheses of Lemma~\ref{8.17:lem-crude-short-move} hold with
$x_{\mathrm{ref}}=x_0$, so its conclusion \eqref{8.17:eq-crude-short-move-b}, read at
$x_{\mathrm{ref}}=x_0$, gives
\begin{equation}\label{8.17:eq-top-step-crude-3}
\operatorname{Den}^{\Lambda'}_{\lambda}(e)\ge
\exp\bigl(-E_{\delta}(\mathfrak{z},e)\bigr)\,
\exp\Bigl(-\frac{\mathbf{A}(e,x_0)}{g^2}
+\sup_{\operatorname{cl}S^{\theta}_{\mathrm{all}}(e)}\Psi\Bigr)>0 .
\end{equation}
The factor $\exp(-\mathbf{A}(e,x_0)/g^2+\sup\Psi)$ is the same in
\eqref{8.17:eq-top-step-crude-2} and \eqref{8.17:eq-top-step-crude-3}; dividing the first by
the second yields
\begin{equation*}
\frac{\operatorname{Num}_{\lambda}(e)}{\operatorname{Den}^{\Lambda'}_{\lambda}(e)}
\le\exp\Bigl(E_{\delta}(\mathfrak{z},e)-\frac{\sum_{\square}\Gamma_{\square}}{g^2}\Bigr),
\end{equation*}
which is \eqref{8.17:eq-top-step-crude-a}.

For the integrated form, $E_{\delta}(\mathfrak{z},e)\le\sup_eE_{\delta}$, so
\eqref{8.17:eq-top-step-crude-a} gives, for every $\lambda\in[\lambda_0]^{s}$ and every $e$,
\begin{equation*}
\operatorname{Num}_{\lambda}(e)\le
\exp\Bigl(\sup_{e}E_{\delta}-\frac{\sum_{\square}\Gamma_{\square}}{g^2}\Bigr)\,
\operatorname{Den}^{\Lambda'}_{\lambda}(e);
\end{equation*}
integrating against $F_{\lambda}$ and summing over $\lambda\in[\lambda_0]^{s}$ bounds the left
side of \eqref{8.17:eq-top-step-crude-1} by the same exponential times
$\sum_{\lambda\in[\lambda_0]^{s}}\int F_{\lambda}\operatorname{Den}^{\Lambda'}_{\lambda}$.
The relaxed cap obeys $\Lambda'\le\Lambda_L$: indeed
$\Lambda'=\theta+\tfrac32\operatorname{Lip}\delta$ with
$\operatorname{Lip}=K_wC_{\delta_0}/\varepsilon_{\min}$ and $\delta\le g_{k_o}$, hence
\begin{equation*}
\Lambda'-1\le\theta-1+\frac{3}{2}\,\frac{K_wC_{\delta_0}g_{k_o}}{\varepsilon_{\min}}
\le\Lambda_L-1
\end{equation*}
by \eqref{8.17:eq-top-step-crude-b}. Since $\Lambda'\le\Lambda_L$, hypothesis~(e) applies: it
bounds $\sum_{\lambda\in[\lambda_0]^{s}}\int F_{\lambda}\operatorname{Den}^{\Lambda'}_{\lambda}$
by the right side of \eqref{8.17:eq-top-step-crude-1}, with no change of variables in $e$.
\end{proof}

The majorant of $\operatorname{Num}_{\lambda}(e)$ in hypothesis~(b) cannot be replaced by one
that keeps the factor $\exp(-A_{Z^{+}}(U^{\mathrm{free}}(e))/g^2)$, the action of the free minimal
configuration on the enlarged region, as in the large-field bounds of Ba{\l}aban's construction:
the minorant \eqref{8.17:eq-top-step-crude-3} is matched at $x_0$, and the quotient would then
contain the gap $\mathbf{A}(e,x_0)-A_{Z^{+}}(U^{\mathrm{free}}(e))$, on which the minorant gives no
information. For regions at depth, hypothesis~(b) is available in the form that holds without a
coordinate chart on $X$.

\begin{proposition}[The short move from first-order descent directions]
\label{8.17:prop-short-move-construction}
Let $\mathfrak{z}$ be a merged top-level raw region and $e$ an admissible crude-regular exterior,
as in Lemma~\ref{8.17:lem-crude-short-move}. Let
$x_{\mathrm{ref}}\in\operatorname{cl}S^{\theta}_{\mathrm{all}}(e)$, let
$\delta\in(0,\min(g,\alpha_P/4)]$ with $g=g_{k_o}$ and $\alpha_P$ the radius in the relative
variable within which the bounds of hypotheses (i)--(iii) below are asserted, and let
$Y\in\mathfrak{g}^{\mathfrak{B}}$ with $\sup_b|Y_b|\le1$. Assume the following.
\begin{itemize}
\item[(i)] \emph{The second-order window bound.} This is the second-order form in which the
analyticity lemma of this chapter, which rests on \cite[Proposition~9]{B-CMP102b} (analyticity
in the relative variable about an arbitrary regular background), is used here: there is a
constant $K_2$ (the second-derivative constant, held fixed below) such that for every
crude-regular $z$ with $f_{a'}(e,z)\le2$ for all $a'\in W$, every
$Y'\in\mathfrak{g}^{\mathfrak{B}}$ with $\sup_b|Y'_b|\le1$ and every $a\in W$, the function
$t\mapsto f_a(e,(z_b\exp(tY'_b))_b)$ is twice differentiable on $|t|\le\alpha_P/4$, with
second derivative bounded by $K_2/\varepsilon_{\min}$.
\item[(ii)] \emph{The window Lipschitz bound, in its linear form.} For every $a\in W$, every
crude-regular $z$ with
$f_{a'}(e,z)\le2$ for all $a'\in W$, and every $x'\in X$ with
$\eta:=\sup_b d(x'_b,z_b)\le\alpha_P$, one has
$f_a(e,x')\le f_a(e,z)+\operatorname{Lip}\cdot\eta$.
\item[(iii)] \emph{The regularity clause of the action Lipschitz bound.} Every $x'$ for which
some $z\in\operatorname{cl}S^{2}_{\mathrm{all}}(e)$ satisfies
$\sup_b d(x'_b,z_b)\le\alpha_P$ and $\sum_b d(x'_b,z_b)<\infty$ is crude-regular.
\item[(iv)] $\theta=1$, so that $f_a(e,x_{\mathrm{ref}})\le1$ for every non-seed
$a\in W\setminus\mathfrak{S}$; the seed half-window condition
$f_{\square}(e,x_{\mathrm{ref}})\le\frac12$ holds for every seed $\square$, and the one-sided
condition \eqref{8.17:eq-opening-rate-condition-b} on $g_{k_o}$ holds.
\end{itemize}
For $a\in W\setminus\mathfrak{S}$ put
$D_a:=\frac{d}{dt}\big|_{t=0}f_a\bigl(e,(x_{\mathrm{ref},b}\exp(tY_b))_b\bigr)$, call $a$
\emph{tight at scale $\delta$} if $f_a(e,x_{\mathrm{ref}})>1-2\operatorname{Lip}\delta$, and
suppose that for every tight $a$ there is a number $\tau_a$ with
$D_a\le-\tau_a\operatorname{Lip}$ and
\begin{equation}\label{8.17:eq-short-move-construction-a}
\tau_a\ \ge\ \tau_{\mathrm{lin}}(\delta):=\frac{K_2\,\delta}{K_wC_{\delta_0}} ,
\end{equation}
where $\operatorname{Lip}=K_wC_{\delta_0}/\varepsilon_{\min}$ is the window response constant,
so that $\tau_{\mathrm{lin}}(\delta)=K_2\delta/(\varepsilon_{\min}\operatorname{Lip})$.
Then $z^*:=(x_{\mathrm{ref},b}\exp(\delta Y_b))_b$ satisfies the opening-rate condition
$\mathcal{A}_{\delta,\tau/2}(\mathfrak{z},e;x_{\mathrm{ref}},z^*)$ of
Definition~\ref{8.17:def-opening-rate-condition}, with rate $\tau_a/2$ at every tight $a$
and rate $1$ at every other $a\in W\setminus\mathfrak{S}$.
\end{proposition}

\begin{proof}
For every $b$, $d(z^*_b,x_{\mathrm{ref},b})=\delta|Y_b|\le\delta$, and $\delta\le\alpha_P/4$
by the choice of $\delta$. Since
$\theta=1$, $x_{\mathrm{ref}}\in\operatorname{cl}S^{1}_{\mathrm{all}}(e)\subset
\operatorname{cl}S^{2}_{\mathrm{all}}(e)$. By hypothesis (iii),
applied with $x'=z=x_{\mathrm{ref}}$, the point $x_{\mathrm{ref}}$ is crude-regular, and, by
the definition of the relaxed window set $\operatorname{cl}S^{1}_{\mathrm{all}}(e)$, every
window of $x_{\mathrm{ref}}$, straddling ones included, is at most $1$, in particular at
most $2$; so hypotheses (i) and (ii) apply at $z=x_{\mathrm{ref}}$.

Let $a$ be tight. Taylor's formula on $[0,\delta]$ with the bound of hypothesis (i) and with
$f_a(e,x_{\mathrm{ref}})\le1$ from hypothesis (iv) gives
\begin{align*}
f_a(e,z^*)&\le f_a(e,x_{\mathrm{ref}})+\delta D_a+\tfrac12\delta^2K_2/\varepsilon_{\min}\\
&\le 1-\tau_a\operatorname{Lip}\delta+\tfrac12K_2\delta^2/\varepsilon_{\min}
\ \le\ 1-\tfrac12\tau_a\operatorname{Lip}\delta .
\end{align*}
The last inequality holds because, with
$\operatorname{Lip}=K_wC_{\delta_0}/\varepsilon_{\min}$, the inequality
$\tfrac12K_2\delta^2/\varepsilon_{\min}\le\tfrac12\tau_aK_wC_{\delta_0}\delta/\varepsilon_{\min}$
is equivalent to $\tau_a\ge K_2\delta/(K_wC_{\delta_0})$, which is
\eqref{8.17:eq-short-move-construction-a}.

Let $a$ be a non-tight non-seed window, so $f_a(e,x_{\mathrm{ref}})\le1-2\operatorname{Lip}
\delta$. By hypothesis (ii) with $\eta\le\delta\le\alpha_P$,
$f_a(e,z^*)\le1-2\operatorname{Lip}\delta+\operatorname{Lip}\delta=1-\operatorname{Lip}\delta$.

Let $\square$ be a seed. By hypotheses (ii) and (iv),
$f_{\square}(e,z^*)\le\frac12+\operatorname{Lip}\delta\le1-\operatorname{Lip}\delta$, the last
inequality because \eqref{8.17:eq-opening-rate-condition-b} and $\delta\le g_{k_o}$ give
\begin{equation*}
\operatorname{Lip}\delta\ \le\ \operatorname{Lip}g_{k_o}
\ =\ \frac{K_wC_{\delta_0}g_{k_o}}{\varepsilon_{\min}}\ \le\ \frac14 .
\end{equation*}
So the rate at every seed is $1$, and
\eqref{8.17:eq-opening-rate-condition-a} holds with the stated rates.
\end{proof}

\noindent\emph{Size of the floor.} At $\delta=g/\max(1,C'p_0)$, where $p_0=p_0(g)$ is the
value of the degree function $p_0(\cdot)$ at $g=g_{k_o}$ and $C'$ is the constant of
Lemma~\ref{8.17:lem-crude-short-move}, the floor
$\tau_{\mathrm{lin}}$ is of order $g/p_0$, up to the constants of the construction: the
short move of Proposition~\ref{8.17:prop-short-move-construction} serves tight sets whose
linearised opening constant exceeds a quantity of order $g/p_0$ --- by the arithmetic of
part~(ii) of \eqref{8.17:eq-model-isoperimetric-a} in the abelian affine-window model
$\mathbb{M}$, these are the coherently
tight flat pieces of side up to order $p_0/g$ --- with no reference to depth. Whether
$\tau_{\mathrm{lin}}$ or the income threshold $\tau_g(\square)$ of
\eqref{8.17:eq-per-seed-budget-c} binds first is decided by the sign of
$2\mathfrak{p}_0-\mathfrak{r}_{\mathfrak{d}}-1$, which the one-sided conditions on the
auxiliary parameters do not determine (the two degrees $\mathfrak{p}_0$ and
$\mathfrak{r}_{\mathfrak{d}}$ are compared where the per-seed budget is set up); both are
boundaries of a range of rates, and neither
is imposed as a condition.

\begin{remark}[The short move needs no chart, interpolation or path]
\label{8.17:rem-no-chart-no-path}
The crude short-move device, Lemma~\ref{8.17:lem-crude-short-move}, and its bound
\eqref{8.17:eq-crude-short-move-b} do not write any configuration in a gauge chart, do not
interpolate between distant configurations, and do not use any path inside the relaxed window
set $S^{\theta}_{\mathrm{all}}(e)$. Consequently none of the three obstructions described below
bears on it.
\end{remark}

\begin{proof}
Throughout, $g=g_{k_o}$, $\delta\in(0,g]$, and $x_{\mathrm{ref}}$ is as
in Lemma~\ref{8.17:lem-crude-short-move}. We write $x'=(x'_b)_{b\in\mathfrak{B}}\in X$ for a
configuration at which the proof of that lemma evaluates the action and the window functionals
in comparison with $x_{\mathrm{ref}}$.

\emph{No chart.} Consider a region containing a flat square of side $r$ with
\begin{equation*}
r\ \ge\ r_*(\mathsf{f})=(\pi/\mathsf{f})^{1/2}.
\end{equation*}
The first obstruction is a lower bound for flat charts: on such a region, no gauge makes the
bond logarithms of a test field with parameter $\mathsf{f}>0$ smaller than
$(\mathsf{f}/4)\,r_*$. This is a statement about writing a configuration in a small flat chart,
at in-plane scale $r_*$ of the order of $\varepsilon^{-1/2}$. The device writes no configuration
in a chart. The only bondwise information it uses is that, for every $b\in\mathfrak{B}$, the
gauge-covariant distance on $G$ between $x'_b$ and $x_{\mathrm{ref},b}$ is at most $\tfrac32\,g$,
and this requires no smallness of the coordinates of $x_{\mathrm{ref}}$ itself.

\emph{No interpolation.} The second obstruction is an interpolation identity. With
$A_s$ the linear interpolation $(1-s)A_0+sA_1$ and $\xi=A_1-A_0$, it reads
\begin{equation*}
F(A_s)=(1-s)F(A_0)+sF(A_1)-s(1-s)\,\xi\wedge\xi .
\end{equation*}
It gives the cost of interpolating between a configuration $x_0$ and a distant admissible
configuration $x$. There $|\xi|$ is of the order of the depth of the configurations compared
times the cap $c_i$ of the affine window. At depth $h_g$ this cost exhausts the whole
allowance available to the comparison in which the identity is used.

The device interpolates between no two configurations. Its only comparisons are two Lipschitz
bounds, each taken at bondwise displacement at most $\tfrac32 g$:
\begin{itemize}
\item the Lipschitz bound of the action $\mathbf{A}(e,\cdot)$;
\item the first-order response bound of the window functionals $f_a$.
\end{itemize}
A quadratic defect at that scale enters only in
Proposition~\ref{8.17:prop-short-move-construction}, and there it equals
$\tfrac12 K_2\delta^2/\varepsilon_{\min}$. Measured against the room scale
$\operatorname{Lip}\,\delta$, its relative size is
\begin{equation*}
\frac{K_2\,\delta^2/(2\varepsilon_{\min})}{\operatorname{Lip}\,\delta}
=\frac{K_2\,\delta}{2\,\varepsilon_{\min}\operatorname{Lip}}
\ \le\ \frac{K_2\,g}{2\,\varepsilon_{\min}\operatorname{Lip}},
\end{equation*}
since $\delta\le g$. By the definition of the window response constant,
$\operatorname{Lip}=K_wC_{\delta_0}/\varepsilon_{\min}$, so the right side equals
$K_2g/(2K_wC_{\delta_0})$, which is of order $g$.
The bound does not depend on the depth of the hole.

\emph{No path.} The statement on sectors of the relaxed window set asserts that, on a deep ball:
\begin{itemize}
\item that set has at least two connected components;
\item every path joining two of them crosses a large plaquette.
\end{itemize}
The device uses a single set, the product
\begin{equation*}
\mathfrak{K}=\prod_{b\in\mathfrak{B}}B_b
\end{equation*}
of the proof of Lemma~\ref{8.17:lem-crude-short-move}. This set is connected and contains $z^*$.
It joins no two points of $S(e)$. Hence the device does not see the sectors.
\end{proof}

\begin{lemma}[Rates forced by a coherently tight flat square]\label{8.17:lem-flat-square-rates}
Let $r\ge1$ be an integer and let $S_r$ be a flat square of side $r$, all of whose bonds lie in
$\mathfrak{B}$, that is, an $r\times r$ array
of plaquettes in one coordinate plane. Its boundary $\partial S_r$ consists of $4r$ bonds. For a
configuration $V\in X=G^{\mathfrak{B}}$ let $\operatorname{hol}V$ be the ordered product of the
$V_b^{\pm1}$ along $\partial S_r$, starting and ending at a fixed corner of $S_r$, called the base point,
in a fixed orientation; $d_G(\operatorname{hol}V,1)$ does not depend on these choices, since a change of
base point conjugates $\operatorname{hol}V$ and a change of orientation inverts it. Let $V(\partial p)$ be
the corresponding product around a
plaquette $p$. Here $d_G$ is the bi-invariant distance on $G$.
\begin{itemize}
\item[(a)] Let $V,V'\in X$ and let $\tau_{\mathrm{op}}$ be a number. Suppose that
$\Phi:=d_G(\operatorname{hol}V,1)\le\pi$ and that
\[
d_G(V'(\partial p),1)\le(1-\tau_{\mathrm{op}})\Phi/r^2\qquad\text{for every plaquette }p\text{ of }S_r.
\]
Then
\begin{equation}\label{8.17:eq-flat-square-rates-1}
\max_{b\in\partial S_r}d_G(V_b,V'_b)\ \ge\ \frac{\tau_{\mathrm{op}}\,\Phi}{4r}.
\end{equation}
\item[(b)] Let $z^*$, $\delta>0$ and the rates $\tau_{a}$, $a\in W\setminus\mathfrak{S}$, satisfy the
opening-rate condition of
Definition~\ref{8.17:def-opening-rate-condition} at $(\mathfrak{z},e;x_{\mathrm{ref}},z^*)$, with the
region $\mathfrak{z}$, the datum $e$ and the
reference configuration $x_{\mathrm{ref}}$ fixed there. Let $S_r\subset H$. Assume the plaquette-model
reading of the windows: each plaquette $p$ of $S_r$ carries a window $a_p\in W\setminus\mathfrak{S}$ with
cap $f>0$, and $d_G(V(\partial p),1)\le f\cdot f_{a_p}(e,V)$ for every $V\in X$. Assume further that these
windows are coherently tight at $x_{\mathrm{ref}}$, in the sense that
$d_G(\operatorname{hol}x_{\mathrm{ref}},1)=r^2f\le\pi$. Let $a$ be one of the windows $a_p$,
$p\subset S_r$, of least rate, so that $\tau_a=\min_{p\subset S_r}\tau_{a_p}$. Then
\begin{equation}\label{8.17:eq-flat-square-rates-a}
\tau_a\ \le\ \frac{4\,\varepsilon_{\min}}{K_wC_{\delta_0}\,f\,r}.
\end{equation}
\end{itemize}
\end{lemma}

\begin{proof}
\emph{Part (a).} We use two consequences of bi-invariance. First, for $a_i,b_i\in G$,
\[
d_G(a_1a_2,b_1b_2)\le d_G(a_1a_2,a_1b_2)+d_G(a_1b_2,b_1b_2)=d_G(a_2,b_2)+d_G(a_1,b_1).
\]
By induction the distance between two products of $n$ factors is at most the sum of the distances of
corresponding factors. Second, $d_G(a^{-1},b^{-1})=d_G(b,a)$, and $d_G(gcg^{-1},1)=d_G(c,1)$.

The triangle inequality gives
\[
\Phi=d_G(\operatorname{hol}V,1)\le d_G(\operatorname{hol}V,\operatorname{hol}V')
+d_G(\operatorname{hol}V',1).
\]
The products $\operatorname{hol}V$ and $\operatorname{hol}V'$ each have $4r$ factors, $V_b^{\pm1}$ and
${V'_b}^{\pm1}$, with the same signs. By the two facts above,
\[
d_G(\operatorname{hol}V,\operatorname{hol}V')\le\sum_{b\in\partial S_r}d_G(V_b,V'_b)
\le4r\max_{b\in\partial S_r}d_G(V_b,V'_b).
\]
By the non-abelian Stokes theorem on the lattice, $\operatorname{hol}V'$ is an ordered product of the
$r^2$ plaquette holonomies $V'(\partial p)$, $p\subset S_r$, each conjugated by the holonomy of a lattice
path from the base point to a corner of $p$. Bi-invariance and the product estimate therefore give
\[
d_G(\operatorname{hol}V',1)\le\sum_{p\subset S_r}d_G(V'(\partial p),1)
\le r^2\cdot\frac{(1-\tau_{\mathrm{op}})\Phi}{r^2}.
\]
Combining the three estimates,
\[
\Phi\le4r\max_{b\in\partial S_r}d_G(V_b,V'_b)+(1-\tau_{\mathrm{op}})\Phi .
\]
This is $\tau_{\mathrm{op}}\Phi\le4r\max_bd_G(V_b,V'_b)$, which is
\eqref{8.17:eq-flat-square-rates-1}.

\emph{Part (b).} Apply part (a) with $(V,V')=(x_{\mathrm{ref}},z^*)$, $\Phi=r^2f$ and
$\tau_{\mathrm{op}}:=\tau_a\operatorname{Lip}\delta$. This is the cap-relative room demanded by
that condition.

We check the plaquette hypothesis of part (a). For each plaquette $p$ of $S_r$, the plaquette-model
reading and the opening-rate condition give
\[
d_G(z^*(\partial p),1)\le f\cdot f_{a_p}(e,z^*)\le f\cdot(1-\tau_{a_p}\operatorname{Lip}\cdot\delta)
\le(1-\tau_a\operatorname{Lip}\cdot\delta)f.
\]
The last inequality uses $\tau_a\le\tau_{a_p}$. The right-hand side equals $(1-\tau_{\mathrm{op}})\Phi/r^2$.

By Definition~\ref{8.17:def-opening-rate-condition} every bond displacement is at most $\delta$.
Hence \eqref{8.17:eq-flat-square-rates-1} gives
\[
\delta\ \ge\ \max_{b\in\partial S_r}d_G(x_{\mathrm{ref},b},z^*_b)\ \ge\
\frac{\tau_a\operatorname{Lip}\cdot\delta\cdot r^2f}{4r}
=\frac{\tau_a\operatorname{Lip}\cdot\delta\cdot f\cdot r}{4}.
\]
Since $\delta>0$, it cancels, leaving $\tau_a\le4/(\operatorname{Lip}\cdot f\cdot r)$. Inserting
$\operatorname{Lip}=K_wC_{\delta_0}/\varepsilon_{\min}$ gives \eqref{8.17:eq-flat-square-rates-a}.
\end{proof}

The bound \eqref{8.17:eq-flat-square-rates-a} does not rule out the crude short-move device at any side
$r$. Instead it fixes the cost of that device. A coherently tight flat piece of side $r$ in the tight set
of $x_{\mathrm{ref}}$ forces rates of order at most $1/r$ on the bonds its windows read. Each such bond
therefore contributes a loss $d(\mathfrak{g})\log r$ to the exponent $E_{\delta}(\mathfrak{z},e)$.
This is the non-abelian counterpart of clause (ii) of the model corollary,
Corollary~\ref{8.17:cor-model-isoperimetric}. It is also the geometric meaning of the admissible range
of opening rates.

\begin{lemma}[Summation of bond charges over the sub-zones]\label{8.17:lem-subzone-summation}
Let $\mathfrak{z}$ be a merged top-level raw region, with its sub-zones $\mathfrak{z}_{\square}$, its
number $n_s$ of seed cubes, its bond set $\mathfrak{B}$, and the numbers $\mathfrak{d}_*$,
$n_*=\mathfrak{d}_*/d(\mathfrak{g})$, $\Omega_{\Psi}(\mathfrak{z})$ and $\Omega_{\Psi*}$, all as in
Notation~\ref{8.17:not-top-step-notation}. Then every bond of $\mathfrak{z}$ lies in some sub-zone
$\mathfrak{z}_{\square}$; fix an assignment of each bond to one such sub-zone. Each sub-zone has at most
$n_*$ bonds, $\mathfrak{z}$ has at most $n_s n_*$ bonds, and
$\Omega_{\Psi}(\mathfrak{z})\le n_s\Omega_{\Psi*}$. Consequently, for any per-bond quantities
$q_b\ge 0$,
\begin{equation}\label{8.17:eq-subzone-summation-1}
\sum_{b\in\mathfrak{B}} q_b
\;\le\; \sum_{\square}\ \sum_{b\ \text{assigned to}\ \square} q_b
\;\le\; \sum_{\square} n_*\cdot\max_{b\in\mathfrak{z}_{\square}} q_b ,
\end{equation}
where $\square$ runs over the $n_s$ seed cubes whose sub-zone lies in $\mathfrak{z}$.
\end{lemma}

\begin{proof}
By Notation~\ref{8.17:not-top-step-notation}, $\mathfrak{z}$ is a connected component of the union of
the sub-zones $\mathfrak{z}_{\square}$, so each of its bonds lies in at least one sub-zone. This makes an
assignment of each bond to a single sub-zone possible.

Each bond variable takes values in $G$, a manifold of dimension $d(\mathfrak{g})$, so it carries
$d(\mathfrak{g})$ real coordinates. The coordinate count of one sub-zone is therefore $d(\mathfrak{g})$ times
its number of bonds. That count is at most $\mathfrak{d}_*$, since $\mathfrak{d}_*$ is the supremum over
$\square$ of these counts. Dividing by $d(\mathfrak{g})$ shows that a sub-zone has at most
$\mathfrak{d}_*/d(\mathfrak{g})=n_*$ bonds.

The count bound recorded with Notation~\ref{8.17:not-top-step-notation} gives
$\mathfrak{d}(\mathfrak{z})\le n_s\mathfrak{d}_*$ and $\Omega_{\Psi}(\mathfrak{z})\le n_s\Omega_{\Psi*}$.
Since $\mathfrak{d}=d(\mathfrak{g})\,n_b$, dividing the first inequality by $d(\mathfrak{g})$ gives
$n_b\le n_s n_*$. Thus $\mathfrak{z}$ has at most $n_s n_*$ bonds, and the second inequality is the stated
bound on $\Omega_{\Psi}(\mathfrak{z})$.

For \eqref{8.17:eq-subzone-summation-1}, each $b\in\mathfrak{B}$ is assigned to exactly one $\square$.
The sets of bonds assigned to the various $\square$ therefore partition $\mathfrak{B}$, so the left-hand
sum equals the double sum, and in particular is bounded by it. For fixed $\square$, the bonds assigned
to $\square$ lie in $\mathfrak{z}_{\square}$ and number at most $n_*$. Since each $q_b\ge0$ is at most
$\max_{b\in\mathfrak{z}_{\square}}q_b$, the inner sum is at most $n_*\max_{b\in\mathfrak{z}_{\square}}q_b$.
Summing over $\square$ gives the second inequality.
\end{proof}

Depth enters these counts through $n_s$, not through the number of bonds per seed cube. By the covering
count, an $h$-deep $\mathfrak{z}$ has
\begin{equation*}
n_s\;\ge\;\Bigl(\frac{2h+1}{2r_z+1}\Bigr)^{d},
\end{equation*}
where $r_z$ is the radius entering the covering count.

\begin{proposition}[The per-seed budget: coupling ceiling and rate range]
\label{8.17:prop-per-seed-budget}
Assume the hypotheses of the top step from the crude short move,
Corollary~\ref{8.17:cor-top-step-crude}, with $g=g_{k_o}$ and with the choice
\[
\delta:=\frac{g}{\max\bigl(1,C'p_0(g_{k_o})\bigr)} .
\]
Assign every bond of $\mathfrak{z}$ to one sub-zone $\mathfrak{z}_{\square}$ as in
Lemma~\ref{8.17:lem-subzone-summation}. For each seed $\square$ of $\mathfrak{z}$ put
$\tau_*(\square):=\min\{\tau_b : b \text{ assigned to } \square\}$ and
\begin{equation}\label{8.17:eq-per-seed-budget-a}
\begin{split}
E_{\square}:={}&\mathfrak{d}_*\Bigl[\log\frac{2\iota_G}{g_{k_o}}
 +\log\max\bigl(1,C'p_0(g_{k_o})\bigr)\Bigr]
 +\mathfrak{d}_*\log\frac{1}{\tau_*(\square)}\\
&+n_*\Bigl[\frac{c_G}{4}+\frac32+\frac{3C}{2\alpha_0}\Bigr]+\Omega_{\Psi*}.
\end{split}
\end{equation}
Let $\Gamma_{\square}\ge0$ be the income of the sandwiched seed $\square$, which satisfies
$\Gamma_{\square}/g^2\ge\gamma_{\mathrm{inc}}\,p_0(g_{j(\square)})^2$. Then:
\begin{itemize}
\item[(a)] $E_{\delta}(\mathfrak{z},e)\le\sum_{\square}E_{\square}$;
\item[(b)] pointwise in $e$, with $\epsilon_{\square}:=\exp(-\Gamma_{\square}/g^2+E_{\square})$,
\[
\operatorname{Num}_{\lambda}(e)\le\Bigl[\prod_{\square}\epsilon_{\square}\Bigr]
\operatorname{Den}^{\Lambda'}_{\lambda}(e);
\]
\item[(c)] for a sandwiched seed $\square$, the bound
$\epsilon_{\square}\le\exp\bigl(-\tfrac12\gamma_{\mathrm{inc}}p_0(g_{j(\square)})^2\bigr)$, which is the
form required in \eqref{8.17:eq-top-step-notation-a}, holds as soon as
$E_{\square}\le\frac12\gamma_{\mathrm{inc}}p_0(g_{j(\square)})^2$. In particular it holds for
every sandwiched seed of $\mathfrak{z}$ if the coupling ceiling
\begin{equation}\label{8.17:eq-per-seed-budget-b}
\begin{split}
&\mathfrak{d}_*\Bigl[\log\frac{2\iota_G}{g_{k_o}}+\log\max\bigl(1,C'p_0(g_{k_o})\bigr)\Bigr]
+\frac{\mathfrak{d}_*}{d(\mathfrak{g})}\Bigl[\frac{c_G}{4}+\frac32+\frac{3C}{2\alpha_0}\Bigr]
+\Omega_{\Psi*}\\
&\qquad\le\frac14\gamma_{\mathrm{inc}}\,p_0(g_j)^2
\end{split}
\end{equation}
holds for every seed level $j$ with $K''<j\le K'$, and the rate range for the pair
$(\mathfrak{z},e)$,
\begin{equation}\label{8.17:eq-per-seed-budget-c}
\begin{gathered}
\mathfrak{d}_*\log\frac{1}{\tau_*(\square)}\le\frac14\gamma_{\mathrm{inc}}\,p_0(g_{j(\square)})^2,
\quad\text{i.e.}\\
\tau_*(\square)\ge\tau_g(\square):=
\exp\Bigl(-\frac{\gamma_{\mathrm{inc}}\,p_0(g_{j(\square)})^2}{4\mathfrak{d}_*}\Bigr),
\end{gathered}
\end{equation}
holds for every sandwiched seed $\square$ of $\mathfrak{z}$.
\end{itemize}
\end{proposition}

\begin{proof}
\emph{Part (a).} Since $\max(1,C'p_0(g_{k_o}))\ge1$, we have $0<\delta\le g$, so the crude
short-move device, Lemma~\ref{8.17:lem-crude-short-move}, applies with this $\delta$, and its
bound \eqref{8.17:eq-crude-short-move-c} gives
\[
E_{\delta}(\mathfrak{z},e)\le\sum_{b\in\mathfrak{B}}q_b+\Omega_{\Psi}(\mathfrak{z}),
\]
where, writing $n_b[\cdots]$ as the sum of $[\cdots]$ over the $n_b$ bonds of $\mathfrak{B}$,
\[
\begin{split}
q_b:={}&d(\mathfrak{g})\Bigl[\log\frac{2\iota_G}{g_{k_o}}
+\log\max\bigl(1,C'p_0(g_{k_o})\bigr)+\log\frac1{\tau_b}\Bigr]\\
&+\frac{c_G}{4}+\frac32+\frac{3C}{2\alpha_0}.
\end{split}
\]
Each bond is assigned to exactly one sub-zone, so $\sum_{b\in\mathfrak{B}}q_b$ is the sum over
$\square$ of the sums over the bonds assigned to $\square$. If $b$ is assigned to $\square$, then
$\tau_b\ge\tau_*(\square)$ by definition of $\tau_*(\square)$, hence
$\log(1/\tau_b)\le\log(1/\tau_*(\square))$, and $q_b\le q_{\square}$, where $q_{\square}$ is $q_b$
with $\tau_b$ replaced by $\tau_*(\square)$. By Lemma~\ref{8.17:lem-subzone-summation} at most
$n_*=\mathfrak{d}_*/d(\mathfrak{g})$ bonds are assigned to $\square$, so their contribution is at
most $n_*q_{\square}$; and $n_*d(\mathfrak{g})=\mathfrak{d}_*$. The same lemma gives
$\Omega_{\Psi}(\mathfrak{z})\le n_s\Omega_{\Psi*}$, that is, one term $\Omega_{\Psi*}$ for each of the
$n_s$ seeds. Adding, the contribution of each seed is exactly $E_{\square}$ of
\eqref{8.17:eq-per-seed-budget-a}, which proves (a).

\emph{Part (b).} Corollary~\ref{8.17:cor-top-step-crude} gives, pointwise in $e$,
\[
\operatorname{Num}_{\lambda}(e)\le\exp\Bigl(E_{\delta}(\mathfrak{z},e)
-\sum_{\square}\Gamma_{\square}/g^2\Bigr)\operatorname{Den}^{\Lambda'}_{\lambda}(e).
\]
By Lemma~\ref{8.17:lem-crude-short-move}, $\operatorname{Den}^{\Lambda'}_{\lambda}(e)$ is bounded
below by a product of exponentials, hence is positive. Since the exponential is increasing,
part (a) gives
\[
\exp\Bigl(E_{\delta}(\mathfrak{z},e)-\sum_{\square}\frac{\Gamma_{\square}}{g^2}\Bigr)
\le\exp\Bigl(\sum_{\square}\Bigl(E_{\square}-\frac{\Gamma_{\square}}{g^2}\Bigr)\Bigr)
=\prod_{\square}\epsilon_{\square},
\]
and (b) follows.

\emph{Part (c).} By the income bound,
$\epsilon_{\square}\le\exp\bigl(-\gamma_{\mathrm{inc}}p_0(g_{j(\square)})^2+E_{\square}\bigr)$.
If $E_{\square}\le\frac12\gamma_{\mathrm{inc}}p_0(g_{j(\square)})^2$, the right side is at most
$\exp(-\frac12\gamma_{\mathrm{inc}}p_0(g_{j(\square)})^2)$. Since $n_*=\mathfrak{d}_*/d(\mathfrak{g})$,
$E_{\square}$ is the left side of \eqref{8.17:eq-per-seed-budget-b} with $j=j(\square)$ plus the
term $\mathfrak{d}_*\log(1/\tau_*(\square))$. The former is at most
$\frac14\gamma_{\mathrm{inc}}p_0(g_{j(\square)})^2$ by the coupling ceiling applied at the
seed level $j=j(\square)$, and the latter is at most the same quantity by the first line of
\eqref{8.17:eq-per-seed-budget-c}; so
$E_{\square}\le\frac12\gamma_{\mathrm{inc}}p_0(g_{j(\square)})^2$. Finally, as
$\mathfrak{d}_*>0$, dividing the first line of \eqref{8.17:eq-per-seed-budget-c} by
$-\mathfrak{d}_*$ and exponentiating shows that it is equivalent to
$\tau_*(\square)\ge\tau_g(\square)$.
\end{proof}

\begin{remark}[The coupling ceiling is implied by the skin ceiling]\label{8.17:rem-coupling-ceiling}
We use the notation for the top step of a merged region fixed in
Notation~\ref{8.17:not-top-step-notation}; in particular $g=g_{k_o}$ and $p_0=p_0(g_{k_o})$.
Write $\mathfrak{r}_{\mathfrak{d}}$ for the degree of $\mathfrak{d}_*$ as a polynomial in
$\log g^{-2}$, with coefficients depending on $L$, $M$, $c_0$. These coefficients enter
through the size of a sub-zone, which is at most $C(c_0,L)M^d\check{R}^d$ times a
polynomial in $\log g^{-2}$. Let $r_{\mathrm{pr}}$ be the degree of the radius $\check{R}$ in
$\log g^{-2}$. The count of real coordinates per sub-zone made in the unbounded-union depth
statement gives
\begin{equation*}
\mathfrak{r}_{\mathfrak{d}}=d\,r_{\mathrm{pr}} ;
\end{equation*}
we use this relation as stated in that statement.
Let $\mathfrak{p}_0$ be the degree in $\log g^{-2}$ of the degree function $p_0(\cdot)$.
For every seed level $j$ with $K''<j\le k_o$, the \emph{skin ceiling} of the top step is an
inequality whose left member contains the terms $\mathfrak{d}_*\log(4\sigma_P/g_{k_o})$,
$2C'\mathfrak{d}_*p_0(g_{k_o})$, $\Omega_{\Psi*}$ and the constant $6$, and whose right member
is $\tfrac{7}{16}\gamma_0''\,p_0(g_j)^2$. Here $\sigma_P$ is the chart radius, and
$\gamma_0''$ is a constant depending only on $d$, $L$, $N$ and Ba{\l}aban's constants.
The \emph{degree relation} of the skin ceiling is stated in the unbounded-union depth
statement in the form
$\deg\mathfrak{d}<\mathfrak{p}_0$; the degree it compares with $\mathfrak{p}_0$ is
$d\,r_{\mathrm{pr}}=\mathfrak{r}_{\mathfrak{d}}$, and we write it as
$\mathfrak{r}_{\mathfrak{d}}<\mathfrak{p}_0$. Then the following hold.
\begin{itemize}
\item[(a)] The terms of the left member of the coupling ceiling
\eqref{8.17:eq-per-seed-budget-b} have the following degrees in $\log g^{-2}$:
\begin{itemize}
\item the term $\mathfrak{d}_*\log(1/g_{k_o})$ has degree $\mathfrak{r}_{\mathfrak{d}}+1$;
\item the term $\mathfrak{d}_*\log(C'p_0)$ has degree $\mathfrak{r}_{\mathfrak{d}}$, up to a
factor $\log\log g^{-2}$;
\item the term $(\mathfrak{d}_*/d(\mathfrak{g}))\cdot[c_G/4+3/2+3C/(2\alpha_0)]$ has degree
$\mathfrak{r}_{\mathfrak{d}}$.
\end{itemize}
The degree of $\Omega_{\Psi*}$ is not determined here. The right member has degree
$2\mathfrak{p}_0$.
\item[(b)] Suppose the skin ceiling holds. Then, up to the constants inside the logarithms,
\eqref{8.17:eq-per-seed-budget-b} holds for every seed level $j$ with $K''<j\le k_o$ and for
$g$ small after $L$ and $M$.
\item[(c)] The inequality \eqref{8.17:eq-per-seed-budget-b} bounds none of $L$, $K$, $m$,
$M$, $n_s$, $|\mathcal{N}|$, $h(\mathfrak{z})$ or $\operatorname{diam}\mathfrak{z}$ from below.
It bounds no coupling from below either. It is a one-sided ceiling.
\item[(d)] The coupling ceiling of the top step in absolute form is an inequality whose left
member has degree $(d+1)r_{\mathrm{pr}}$ and whose right member has degree
$2\mathfrak{p}_0$. The degree relation compares degree $d\,r_{\mathrm{pr}}$ with
$\mathfrak{p}_0$. The coupling ceiling \eqref{8.17:eq-per-seed-budget-b} is weaker than both,
in the sense that it follows from the skin ceiling.
\item[(e)] Suppose that, instead of the window Lipschitz bound used in
Proposition~\ref{8.17:prop-per-seed-budget}, only its coarse form were available. The coarse
form states that a displacement of at most $g$ raises a window functional by at most
$\rho_1/2$. Then, under the degree relation $\mathfrak{r}_{\mathfrak{d}}<\mathfrak{p}_0$, the
conclusion of (b) is unchanged.
\end{itemize}
\end{remark}

\begin{proof}
(a) Since $\log(1/g_{k_o})=\tfrac12\log g^{-2}$, we have
\begin{equation*}
\mathfrak{d}_*\log(1/g_{k_o})=\tfrac12\,\mathfrak{d}_*\log g^{-2},
\end{equation*}
and this has degree $\mathfrak{r}_{\mathfrak{d}}+1$.

Since $p_0$ is a polynomial of degree $\mathfrak{p}_0$ in $\log g^{-2}$, its logarithm
$\log(C'p_0)$ is of order $\log\log g^{-2}$. Hence $\mathfrak{d}_*\log(C'p_0)$ has degree
$\mathfrak{r}_{\mathfrak{d}}$ up to a factor $\log\log g^{-2}$. The same holds with
$\log\max(1,C'p_0)$ in place of $\log(C'p_0)$.

The third term is $\mathfrak{d}_*$ times the number
$[c_G/4+3/2+3C/(2\alpha_0)]/d(\mathfrak{g})$, which does not depend on $g$; so it has degree
$\mathfrak{r}_{\mathfrak{d}}$.

The quantity $\Omega_{\Psi*}$ occurs unchanged in the left member of the skin ceiling.
Its degree is not determined here; its standing in \eqref{8.17:eq-per-seed-budget-b} is the
one it has in the skin ceiling.

The right member is $\tfrac14\gamma_{\mathrm{inc}}p_0(g_j)^2$. It is the square of a
polynomial of degree $\mathfrak{p}_0$ in $\log g_j^{-2}$, so it has degree
$2\mathfrak{p}_0$.

(b) The skin ceiling holds as a ceiling in the regime $\Phi\downarrow0$. It holds there
because either a polynomial in $\log g^{-2}$ times a positive power of $g$ is bounded by a
constant, or $\mathfrak{r}_{\mathfrak{d}}<\mathfrak{p}_0$. In either case $g$ is chosen after
$L$ and $M$, and the choice is uniform in the level by the lower bound on the flow of the
running couplings. The left member of the skin ceiling contains
$\mathfrak{d}_*\log(4\sigma_P/g_{k_o})$, which
by (a) has degree $\mathfrak{r}_{\mathfrak{d}}+1$. Its right member has degree
$2\mathfrak{p}_0$. Hence the skin ceiling carries the degree inequality
\begin{equation}\label{8.17:eq-coupling-ceiling-2}
\mathfrak{r}_{\mathfrak{d}}+1<2\mathfrak{p}_0 .
\end{equation}
The left member of \eqref{8.17:eq-per-seed-budget-b} is obtained from that of the skin
ceiling by two replacements:
\begin{itemize}
\item $\log(4\sigma_P/g)$ is replaced by $\log(2\iota_G/g)$, which involves no chart radius;
\item $2C'\mathfrak{d}_*p_0(g_{k_o})$, of degree $\mathfrak{r}_{\mathfrak{d}}+\mathfrak{p}_0$,
is replaced by the smaller term $\mathfrak{d}_*\log\max(1,C'p_0)$.
\end{itemize}
The second replacement is smaller because $\log\max(1,y)\le y\le 2y$ for $y>0$; applied with
$y=C'p_0$ and multiplied by $\mathfrak{d}_*$, this gives
\begin{equation*}
\mathfrak{d}_*\log\max(1,C'p_0)\le 2C'\mathfrak{d}_*p_0 .
\end{equation*}
By (a) and \eqref{8.17:eq-coupling-ceiling-2}, every term of
\eqref{8.17:eq-per-seed-budget-b} whose degree is determined has degree below that of the
right member. The remaining term is $\Omega_{\Psi*}$, whose degree is not determined here;
it occurs unchanged in the left member of the skin ceiling, and, as in (a), its standing in
\eqref{8.17:eq-per-seed-budget-b} is taken to be its standing in the skin ceiling.
So, up to the constants inside the logarithms, \eqref{8.17:eq-per-seed-budget-b} holds at
every seed level, for $g$ small after $L$ and $M$, whenever the skin ceiling does.

(c) The members of \eqref{8.17:eq-per-seed-budget-b} involve only $\mathfrak{d}_*$,
$g_{k_o}$, $g_j$, $p_0$, $\Omega_{\Psi*}$ and fixed constants. The count $\mathfrak{d}_*$ is a
supremum over single sub-zones, so $n_s$, $|\mathcal{N}|$, $h(\mathfrak{z})$ and
$\operatorname{diam}\mathfrak{z}$ do not appear; nor do the number of steps $K$ or the torus
exponent $m$, which occur in neither member. The parameters $L$ and $M$ enter only through
constants that are fixed before $g$, namely the coefficients of $\mathfrak{d}_*$ and the
constants $c_G$, $C$, $C'$, $\alpha_0$ (the constant $\iota_G$ depends on $N$ only), and $g$
is chosen after them. Thus the inequality is a condition on the coupling alone, from above.

(d) By the first display of the remark, $\mathfrak{r}_{\mathfrak{d}}=d\,r_{\mathrm{pr}}$. So
the degree relation is the comparison of degree $d\,r_{\mathrm{pr}}$ with $\mathfrak{p}_0$.
As stated in the unbounded-union depth statement, the coupling ceiling in absolute form compares
degree $(d+1)r_{\mathrm{pr}}$ with $2\mathfrak{p}_0$. By (b),
\eqref{8.17:eq-per-seed-budget-b} follows from the skin ceiling.

(e) Under the coarse form one takes the displacement parameter $\delta$ of
Proposition~\ref{8.17:prop-per-seed-budget} equal to $g$ instead of $g/\max(1,C'p_0)$. Then
the term
\begin{equation*}
\tfrac32\,C'\varepsilon\delta g^{-2}n_b=\tfrac32\,C'p_0\,n_b
\end{equation*}
returns to the left member. It has degree $\mathfrak{r}_{\mathfrak{d}}+\mathfrak{p}_0$. Under
the degree relation $\mathfrak{r}_{\mathfrak{d}}<\mathfrak{p}_0$, this degree is below
$2\mathfrak{p}_0$. The returned term is of the same kind as the term
$2C'\mathfrak{d}_*p_0(g_{k_o})$ of the skin ceiling. Hence the argument of (b) applies
unchanged.
\end{proof}

\begin{remark}[The rate range caps lateral coherence, not the coupling]
\label{8.17:rem-rate-range-coherence}
Let $(\mathfrak{z},e)$ be a configuration pair and $\square$ a sandwiched seed of $\mathfrak{z}$. Let
$\tau_*(\square)$, $\tau_g(\square)$ and the rate range $(\mathrm{R}\text{-}\tau)$ be as in
Proposition~\ref{8.17:prop-per-seed-budget} and \eqref{8.17:eq-per-seed-budget-c}; here $j(\square)$
is the index attached to the seed $\square$ there, so that $g_{j(\square)}$ is the running coupling
at that index. Let
$x_0(e)$ be the minimising configuration at $e$.
\begin{itemize}
\item[(a)] The number $\tau_*(\square)$ is a datum of the pair $(\mathfrak{z},e)$: it is a reciprocal
lateral coherence scale of the signed tight set of $x_0(e)$ on the bonds of the sub-zone
$\mathfrak{z}_{\square}$. It has no expression in $g$. Read as a condition, $(\mathrm{R}\text{-}\tau)$ is
\begin{equation}\label{8.17:eq-rate-range-coherence-1}
\frac{1}{\tau_*(\square)}\le
\exp\Bigl(\frac{\gamma_{\mathrm{inc}}\,p_0(g_{j(\square)})^2}{4\mathfrak{d}_*}\Bigr).
\end{equation}
This is an upper bound, at fixed $g$, on a geometric quantity of the configuration pair. It caps
the lateral coherence of the signed tight set of $x_0(e)$. It is a condition on $(\mathfrak{z},e)$;
it is not a condition
on $g$, $L$, $K$ or $m$.
\item[(b)] The right member of \eqref{8.17:eq-rate-range-coherence-1} tends to infinity as $g\to0$,
as the depth boundary $h_g$ and the second range boundary $H^{*}$ of the crude bound do. Like them,
it is the boundary of the range on which the corollary of the crude short-move device quoted in
Proposition~\ref{8.17:prop-per-seed-budget} delivers $(\mathrm{TS}_{\mathfrak{z}})$.
\item[(c)] Put
\begin{equation}\label{8.17:eq-rate-range-coherence-a}
r_g:=\frac{4\varepsilon_{\min}}{K_wC_{\delta_0}\cdot f\cdot\tau_g(\square)} .
\end{equation}
Here $K_w$, $C_{\delta_0}$ and the plaquette cap $f$ are as in
Lemma~\ref{8.17:lem-flat-square-rates}.
Every coherently tight flat piece of side $r>r_g$ that lies in the tight set of $x_0(e)$ and reads
bonds of $\square$ lies outside $(\mathrm{R}\text{-}\tau)$.
\item[(d)] The loss term $\mathfrak{d}_*\log(1/\tau_*(\square))$ of $E_{\square}$, which
$(\mathrm{R}\text{-}\tau)$ bounds, cannot be removed by a better estimate
within the organisation of the crude bound.
\end{itemize}
Hence the inequality ``loss $\le$ income'' is a genuine condition for the $\tau$-free part of
$E_{\square}$. For the $\tau$-part it is either false or capped. The crude short-move device
therefore does not serve the whole deep family of regions $\mathfrak{z}$.
\end{remark}

\begin{proof}
(a) By definition, $\tau_*(\square)=\min\{\tau_b: b \text{ assigned to } \square\}$ is a minimum of
opening rates. These rates are read off the signed tight set of $x_0(e)$ on the bonds of
$\mathfrak{z}_{\square}$.

Two results show that this minimum is a reciprocal lateral coherence scale of that set. The first
is Lemma~\ref{8.17:lem-flat-square-rates}. The second is the model corollary,
Corollary~\ref{8.17:cor-model-isoperimetric}. That corollary expresses the opening constant as the
isoperimetric ratio \eqref{8.17:eq-model-isoperimetric-a}. By (i) of that corollary the opening
constant equals $1$ when the tight windows have pairwise disjoint sets of free bonds; (ii) of that
corollary bounds it from above when the tight set contains, with one sign, all plaquettes of a flat
square $S_r$ whose bonds are all free.

None of these objects involves $g$. The rate range \eqref{8.17:eq-per-seed-budget-c} reads
\begin{equation*}
\mathfrak{d}_*\log\bigl(1/\tau_*(\square)\bigr)\le
\tfrac14\gamma_{\mathrm{inc}}\,p_0(g_{j(\square)})^2 .
\end{equation*}
Since $\mathfrak{d}_*>0$, we may divide by $\mathfrak{d}_*$ and exponentiate. This gives
\eqref{8.17:eq-rate-range-coherence-1}, which is equivalent to
$\tau_*(\square)\ge\tau_g(\square)$.

(b) The function $p_0(\cdot)$ has degree $\mathfrak{p}_0$ in $\log g^{-2}$, and $\mathfrak{d}_*$ has
degree $\mathfrak{r}_{\mathfrak{d}}$. So the exponent in \eqref{8.17:eq-rate-range-coherence-1} has
degree $2\mathfrak{p}_0-\mathfrak{r}_{\mathfrak{d}}>0$, and it tends to infinity as $g\to0$.

(c) Let $S_r$ be a flat square of side $r$ whose plaquette windows (cap $f$) are all tight at
$x_0(e)$ and read bonds of $\square$. Here $K_wC_{\delta_0}/\varepsilon_{\min}=\operatorname{Lip}$ is
the window response constant.

Lemma~\ref{8.17:lem-flat-square-rates} applies non-abelian Stokes around $S_r$. It gives the rate
cap \eqref{8.17:eq-flat-square-rates-a},
\begin{equation*}
\tau_a\le\frac{4\varepsilon_{\min}}{K_wC_{\delta_0}\cdot f\cdot r},
\end{equation*}
in which $\delta$ cancels. Lemma~\ref{8.17:lem-flat-square-rates} records that the piece thereby
forces rates of order $1/r$ on the bonds it reads; the cap enters $\tau_*(\square)$ with the
constant of \eqref{8.17:eq-flat-square-rates-a}, that is, for every bond $b$ of $\square$ read by
$S_r$ the rate $\tau_b$ is at most the right-hand side above. Taking the minimum over these bonds,
\begin{equation*}
\tau_*(\square)\le\frac{4\varepsilon_{\min}}{K_wC_{\delta_0}\cdot f\cdot r}.
\end{equation*}
If $r>r_g$, then by \eqref{8.17:eq-rate-range-coherence-a} the right-hand side is strictly smaller
than $\tau_g(\square)$. Hence $\tau_*(\square)<\tau_g(\square)$, and $(\mathrm{R}\text{-}\tau)$ fails.

(d) We use the model proposition, Proposition~\ref{8.17:prop-model-loss-sharp}. In the affine-window
model it gives \eqref{8.17:eq-model-loss-sharp-a}. With the choice $\kappa=Wg$, $K=2Wg$,
$\varphi=3Wg/2$, $W\ge3$ fixed, made there, the windows are wide and the four quantities
$\Lambda_1$, $P$, $\max_i\varphi_i$ and $\min_i c_i/(g\varphi_i)$ stay bounded, while the loss per
coordinate is at least
\begin{equation*}
\tfrac12\log\tfrac{n-1}{2e}-\tfrac{\log\sqrt\pi}{n},
\end{equation*}
which is unbounded in $n$.

So no loss constant depending only on those four quantities exists. This is the loss carried by the
summand $\mathfrak{d}_*\log(1/\tau_*(\square))$ of $E_{\square}$, which $(\mathrm{R}\text{-}\tau)$
bounds.

Finally, Proposition~\ref{8.17:prop-per-seed-budget} needs
$E_{\square}\le\frac12\gamma_{\mathrm{inc}}p_0(g_{j(\square)})^2$. This requirement splits into a
coupling ceiling for the $\tau$-free part and the bound $(\mathrm{R}\text{-}\tau)$ for the
$\tau$-part. By (c), the second bound is false on every configuration pair carrying a coherently
tight flat piece of side $r>r_g$. By (d), it cannot be improved within the crude organisation.
\end{proof}

\begin{remark}[The full ball deficit on every bond]\label{8.17:rem-ball-deficit}
Let $\mathfrak{z}$ be a merged top-level raw region, $g=g_{k_o}$, and apply the crude short-move
device (Lemma~\ref{8.17:lem-crude-short-move}) with $\delta=g/\max(1,C'p_0)$.

The device bounds the normalised Haar measure of $S(e)\cap\mathcal{E}$ in
$\operatorname{Num}_{\lambda}(e)$ by $1$. In the exponent of its lower bound for
$\operatorname{Den}^{\Lambda'}_{\lambda}(e)$ it charges, on every bond $b\in\mathfrak{B}$, the sum
of the following four terms:
\begin{itemize}
\item[(a)] the $g$-ball deficit $d(\mathfrak{g})\log(2\iota_G/g)+c_G/4$;
\item[(b)] the short-move charge $d(\mathfrak{g})\log\max(1,C'p_0)$;
\item[(c)] the action and $\Psi$ charge $3/2+3C/(2\alpha_0)$;
\item[(d)] the rate charge $d(\mathfrak{g})\log(1/\tau_b)$.
\end{itemize}
The global term $\Omega_{\Psi}(\mathfrak{z})$ of \eqref{8.17:eq-crude-short-move-c} is charged
once. The four charges (a)--(d) are levied on every bond $b\in\mathfrak{B}$, whether or not $b$ is
read by a tight window and whether $b$ is deep or shallow; only the rate $\tau_b$ in (d) reflects
the windows reading $b$, and (a)--(c) do not.

Consequently the device incurs on deep coordinates read by no tight window (the inactive
coordinates) the full deficit of
order $\log g^{-2}$, and not merely $\log(C'p_0)$. This charge is payable. The sharper loss of
the model (Remark~\ref{8.17:rem-model-loss-per-coordinate}) does not transfer to the device.
The short-move charge that does transfer is $d(\mathfrak{g})\log(C'p_0)$ on every bond, and it
is of lower degree in $\log g^{-2}$ than the charge $2C'p_0$ of the move $\Xi$.
\end{remark}

\begin{proof}
\emph{The per-bond charges.} At $\delta=g/\max(1,C'p_0)$ the exponent $E_{\delta}(\mathfrak{z},e)$
of Lemma~\ref{8.17:lem-crude-short-move} is bounded by \eqref{8.17:eq-crude-short-move-c}. In the
sum over $b\in\mathfrak{B}$ of that bound, the bracket multiplying $d(\mathfrak{g})$ is the
expansion of
$\log(2\iota_G/(\delta\tau_b))$ at this value of $\delta$:
\begin{equation}\label{8.17:eq-ball-deficit-1}
\log\frac{2\iota_G}{\delta\tau_b}
=\log\frac{2\iota_G}{g}+\log\max(1,C'p_0)+\log\frac{1}{\tau_b}.
\end{equation}
The terms of \eqref{8.17:eq-ball-deficit-1} give (a) without $c_G/4$, then (b) and (d). The
term $n_b\,[c_G/4+3/2+3C/(2\alpha_0)]$ of \eqref{8.17:eq-crude-short-move-c} contributes
$c_G/4+3/2+3C/(2\alpha_0)$ per bond; this completes (a) and gives (c). The last term of
\eqref{8.17:eq-crude-short-move-c} is $\Omega_{\Psi}(\mathfrak{z})$.

The terms (a)--(c) do not depend on whether $b$ is read by a tight window or on the depth of $b$,
and (d) has the same form on every bond, only the value of $\tau_b$ reflecting the windows that
read $b$. This proves the first assertion. Moreover,
\begin{equation*}
d(\mathfrak{g})\log\frac{2\iota_G}{g}
=\tfrac12\,d(\mathfrak{g})\log g^{-2}+d(\mathfrak{g})\log(2\iota_G).
\end{equation*}
Hence (a) is $\tfrac12\log g^{-2}$ plus a constant per real coordinate, which is the
per-coordinate deficit of the $g$-ball at the current level, of the form $\tfrac12\log g^{-2}$
plus a constant. In particular it is charged on
inactive deep coordinates, and
it is not reduced there to $\log(C'p_0)$.

\emph{Payability.} Sum (a) over the $n_*$ bonds assigned to one sub-zone $\square$. Each bond
carries $d(\mathfrak{g})$ real coordinates, so $n_*d(\mathfrak{g})=\mathfrak{d}_*$ and
\begin{equation*}
\sum_{b\ \text{assigned to}\ \square}\Bigl[d(\mathfrak{g})\log\frac{2\iota_G}{g}+\frac{c_G}{4}\Bigr]
=\tfrac12\,\mathfrak{d}_*\log g^{-2}+\mathfrak{d}_*\log(2\iota_G)+n_*\frac{c_G}{4}.
\end{equation*}
Since $\mathfrak{d}_*$ has degree $\mathfrak{r}_{\mathfrak{d}}$ in $\log g^{-2}$, the sum has
degree $\mathfrak{r}_{\mathfrak{d}}+1$. This is the degree of the charge
$\mathfrak{d}_*\log(4\sigma_P/g)$, where $\sigma_P$ denotes the constant entering the contraction
volume of the matched devices of this chapter; those devices already pay that charge as their
contraction volume.

The statement that fixes the per-coordinate deficit of the $g$-ball also records that, in a
ratio in which the seeds of a region occur on one side only, the deficit is not cancelled. That
is the situation here: the recent seeds of $\mathfrak{z}$ are present in
$\operatorname{Num}_{\lambda}(e)$ only, and their income $\Gamma_{\square}$ pays for the deficit.

\emph{Non-transfer of the model bound.}
Remark~\ref{8.17:rem-model-loss-per-coordinate} gives, in the model, a loss of
$\log(C'p_0)+2$ per coordinate read by a tight window and $O(1)$ per other coordinate. That
bound uses which coordinates carry the gradient of the action at the reference point. The bound
on the action on which Lemma~\ref{8.17:lem-crude-short-move} rests is a Lipschitz bound, and it
does not distinguish those
coordinates from the others. The short move therefore costs (b) on every bond.

At $\delta=g/\max(1,C'p_0)$ this cost is $d(\mathfrak{g})\log(C'p_0)$ per bond whenever
$C'p_0\ge1$. With $p_0=A_0(\log g^{-2})^{\mathfrak{p}_0}$ it equals
\begin{equation*}
d(\mathfrak{g})\bigl[\log(C'A_0)+\mathfrak{p}_0\log\log g^{-2}\bigr].
\end{equation*}
This is of lower degree in $\log g^{-2}$ than the charge $2C'p_0$ of the move $\Xi$, which has
degree $\mathfrak{p}_0$.
\end{proof}

\begin{corollary}[The top step for regions with rates in range]\label{8.17:cor-top-step-rate-range}
Write $g=g_{k_o}$ and put
\begin{equation*}
\delta=\frac{g_{k_o}}{\max(1,C'p_0(g_{k_o}))},
\end{equation*}
where $p_0(\cdot)$ is the degree function of the coupling (Notation~\ref{0.2:not-glossary}).
Here $C'$ and $\alpha_P$ are the constants of Lemma~\ref{8.17:lem-crude-short-move}, and
$K_wC_{\delta_0}/\varepsilon_{\min}$ is the window response constant $\operatorname{Lip}$ of that lemma.
Let $\theta\in[1,2]$. Assume the following.
\begin{itemize}
\item[(a)] The statements from which the crude short-move device
(Lemma~\ref{8.17:lem-crude-short-move}) and the top step from the crude short move
(Corollary~\ref{8.17:cor-top-step-crude}) are derived, namely: the Lipschitz bound for the
window functionals $f_a$, with response constant $\operatorname{Lip}$; the Lipschitz bound for the
gradient of the background action $\mathbf{A}(e,\cdot)$, with constant $\Lambda_1$; the bound on
the oscillation of the localised small terms $\Psi$ by $\Omega_{\Psi}(\mathfrak{z})$; the bound on
the density of the Haar measure in the exponential chart, with constant $c_G$; the statement on
the two lowest tower levels; the upper comparison bound at its comparison point $x_0$; the
strict-class relaxation statement; and the second relaxation statement.
\item[(b)] The one-sided conditions \eqref{8.17:eq-per-seed-budget-b},
\eqref{8.17:eq-opening-rate-condition-b} and \eqref{8.17:eq-top-step-crude-b}, and
$g_{k_o}\le\min(1,2\alpha_P/3)$.
\end{itemize}
Let $\mathfrak{z}$ be any merged top-level raw region (any $n_s$, lateral extent, shape and
inradius; in particular every region that is $h_g$-deep, $h_g$ being the depth boundary). Let $e$
be any admissible crude-regular exterior such that, as required in the hypotheses of
Corollary~\ref{8.17:cor-top-step-crude}, the comparison point $x_0$ of that corollary lies in
$\operatorname{cl}S^{\theta}_{\mathrm{all}}(e)$, and such that the opening-rate condition of
Definition~\ref{8.17:def-opening-rate-condition} holds at $(\mathfrak{z},e;x_0,z^*)$ at this
$\delta$, with some short-moved point $z^*$ and opening rates $\tau$ whose least rates, in the
notation of \eqref{8.17:eq-per-seed-budget-c}, satisfy $\tau_*(\square)\ge\tau_g(\square)$ for
every sandwiched seed $\square$ of $\mathfrak{z}$. Let $\lambda$ be any admissible content.

Then the top step \eqref{8.17:eq-top-step-notation-a} holds with $e'=e$, with every factor of its
product satisfying
\begin{equation}\label{8.17:eq-top-step-rate-range-1}
\epsilon_{\square}\le\exp\bigl(-\tfrac12\gamma_{\mathrm{inc}}\,p_0(g_{j(\square)})^2\bigr),
\end{equation}
where $j(\square)$ is the level of the seed $\square$, and with
\begin{equation*}
\Lambda'=\theta+\tfrac32K_wC_{\delta_0}\delta/\varepsilon_{\min}\le\Lambda_L .
\end{equation*}
\end{corollary}

\begin{proof}
Since $\max(1,C'p_0(g_{k_o}))\ge1$, we have $\delta\in(0,g]$. Together with
$g\le\min(1,2\alpha_P/3)$, $\theta\in[1,2]$, $x_0\in\operatorname{cl}S^{\theta}_{\mathrm{all}}(e)$,
the admissibility of $e$ and of $\lambda$, the opening-rate condition at
$(\mathfrak{z},e;x_0,z^*)$ assumed in the statement,
hypothesis (a) and condition \eqref{8.17:eq-opening-rate-condition-b}, this shows that the
hypotheses of Lemma~\ref{8.17:lem-crude-short-move} hold with $x_{\mathrm{ref}}:=x_0$.
Corollary~\ref{8.17:cor-top-step-crude} therefore gives, pointwise in $e$,
\begin{equation*}
\operatorname{Num}_{\lambda}(e)\le\exp\Bigl(E_{\delta}(\mathfrak{z},e)
-\sum_{\square}\Gamma_{\square}/g^2\Bigr)\operatorname{Den}^{\Lambda'}_{\lambda}(e).
\end{equation*}
At this $\delta$, Proposition~\ref{8.17:prop-per-seed-budget} gives
$E_{\delta}(\mathfrak{z},e)\le\sum_{\square}E_{\square}$, with $E_{\square}$ as in
\eqref{8.17:eq-per-seed-budget-a}. Hence
\begin{equation*}
\operatorname{Num}_{\lambda}(e)\le\Bigl[\prod_{\square}\epsilon_{\square}\Bigr]
\operatorname{Den}^{\Lambda'}_{\lambda}(e),
\qquad \epsilon_{\square}:=\exp\bigl(-\Gamma_{\square}/g^2+E_{\square}\bigr).
\end{equation*}
The seed incomes $\Gamma_{\square}$ and the exponents $E_{\square}$ are indexed by the sandwiched
seeds $\square$ of $\mathfrak{z}$, as in the notation of the top step
\eqref{8.17:eq-top-step-notation-a}; fix one of them.

We bound $E_{\square}$ in two parts. By the definitions of $n_*$ and $\mathfrak{d}_*$ as the numbers
of bonds and of real coordinates of a sub-zone, and since every bond carries $d(\mathfrak{g})$ real
coordinates ($X=G^{\mathfrak{B}}$), we have
$n_*=\mathfrak{d}_*/d(\mathfrak{g})$. The seed level $j(\square)$ lies in $(K'',k_o]$, where $K''$
is the level, below $k_o$, at which the reach-isolated lineage-closed union $\mathcal{N}$ is live.
Condition
\eqref{8.17:eq-per-seed-budget-b} at $j=j(\square)$ therefore bounds the sum of all terms of
$E_{\square}$ other than $\mathfrak{d}_*\log(1/\tau_*(\square))$ by
$\tfrac14\gamma_{\mathrm{inc}}p_0(g_{j(\square)})^2$. Since $\tau_*(\square)\ge\tau_g(\square)$,
\begin{equation*}
\mathfrak{d}_*\log\frac{1}{\tau_*(\square)}\le\mathfrak{d}_*\log\frac{1}{\tau_g(\square)}
=\tfrac14\gamma_{\mathrm{inc}}p_0(g_{j(\square)})^2 .
\end{equation*}
Thus $E_{\square}\le\tfrac12\gamma_{\mathrm{inc}}p_0(g_{j(\square)})^2$. By the seed-income
inequality fixed with the notation of the top step \eqref{8.17:eq-top-step-notation-a},
$\Gamma_{\square}/g^2\ge\gamma_{\mathrm{inc}}p_0(g_{j(\square)})^2$, and so
\eqref{8.17:eq-top-step-rate-range-1} follows. The right-hand side depends on $\square$ only
through $j(\square)$; it depends on none of $n_s$, $\operatorname{diam}\mathfrak{z}$,
$h(\mathfrak{z})$.

Finally, $\Lambda'=\theta+\tfrac32\operatorname{Lip}\delta$ with
$\operatorname{Lip}=K_wC_{\delta_0}/\varepsilon_{\min}$. Since $\delta\le g_{k_o}$, condition
\eqref{8.17:eq-top-step-crude-b} gives
\begin{equation*}
\Lambda'\le\theta+\tfrac32K_wC_{\delta_0}g_{k_o}/\varepsilon_{\min}\le\Lambda_L .
\end{equation*}
This is the top step \eqref{8.17:eq-top-step-notation-a} with $e'=e$.
\end{proof}

\begin{proposition}[The top step for deep regions with coherent tight sets; proof outstanding]
\label{8.17:prop-deep-coherent-top-step}
Let $K$ and $m$ be given, and let $K''\le K'-c_0$ be levels, where $c_0$ is the number of top
levels over which the sub-zones $\mathfrak{z}_{\square}$ forming a merged top-level raw region
are taken; put $k_o:=K'$. Let $\mathcal{N}$ be a
reach-isolated, lineage-closed union live at level $K''$, and let $\mathfrak{z}$ be an
$h_g$-deep merged top-level raw region of $\mathcal{N}$. Put
\begin{equation*}
\delta:=\frac{g_{k_o}}{\max(1,C'p_0(g_{k_o}))} ,
\end{equation*}
the radius at which Corollary~\ref{8.17:cor-top-step-rate-range} applies the opening-rate
condition, $C'$ being the constant fixed there.
Consider every admissible content $\lambda$ and every admissible crude-regular $e$, with
comparison point $x_0(e)$, with the
following property: no $z^*\in X$ with $\sup_{b\in\mathfrak{B}}d(z^*_b,x_0(e)_b)\le\delta$
satisfies the opening-rate condition of Definition~\ref{8.17:def-opening-rate-condition}
at the comparison point $x_0(e)$, in the form \eqref{8.17:eq-opening-rate-condition-a},
with rates such that $\tau_*(\square)\ge\tau_g(\square)$ at every sandwiched seed
$\square$ in the sense of Corollary~\ref{8.17:cor-top-step-rate-range},
where $\tau_g(\square)$ is the threshold of \eqref{8.17:eq-per-seed-budget-c}.
In the plaquette model $\mathbb{M}_{\square}$ this is equivalent to the property that the
tight set of $x_0(e)$ at depth contains signed coherent pieces of isoperimetric scale
larger than $r_g$ of \eqref{8.17:eq-rate-range-coherence-a}. Then the top step
\eqref{8.17:eq-top-step-notation-a} holds for $\mathfrak{z}$, $\lambda$ and $e$ with
\begin{gather*}
\epsilon_{\square}\le\exp\Bigl(-\tfrac12\gamma_{\mathrm{inc}}\,p_0(g_{j(\square)})^2\Bigr)
\quad\text{for every seed }\square,\\
\Lambda'\le\Lambda_L,
\end{gather*}
where $j(\square)$ is as in Corollary~\ref{8.17:cor-top-step-rate-range}, and with no
constant depending on $n_s$, on $\operatorname{diam}\mathfrak{z}$ or on $h(\mathfrak{z})$.
\end{proposition}

Proof outstanding.

\textit{Route.} For the inequality of the statement, its form averaged over $e$ in an event
$B'$ whose definition does not read the lineage of $\mathcal{N}$ suffices for the use made of
this proposition in the old-age tail statement. Since the pairs $(\lambda,e)$ of
the statement lie outside the rate
range \eqref{8.17:eq-per-seed-budget-c}, Corollary~\ref{8.17:cor-top-step-rate-range} does
not apply to them. Two routes are plausible. The first is an annealed
form of the estimate. The
second is a comparison device at depth which uses no chart and keeps the normalised Haar
measure $\int_X\mathbf{1}\{x\in S(e)\}\,dH(x)$ of the strict window set on both sides of the
inequality. A complementary statement, which would supply the rate range at depth, is the
following: for every $h_g$-deep $\mathfrak{z}$ and every admissible crude-regular $e$, the
tight set of $x_0(e)$ at depth greater than $h_1-2\Delta$, where $h_1$ is the half-width
threshold for tight young columns in the description of kinked exteriors and $\Delta$ is the
depth offset attached to it, admits opening rates $\tau_a$ with
$\tau_*(\square)\ge\tau_g(\square)$. In the abelian model $\mathbb{M}$, in the terms of the
isoperimetric proposition of that model, it reads: the signed tight set of the constrained
minimiser has opening
constant at least $\tau_g$ on the scale of each sub-zone. Whether the family of
$(\mathcal{N},\mathfrak{z},\lambda,e)$ in the hypothesis is empty or nonempty is not asserted.

\begin{definition}[Partition functions, complex sources and reflection cells]
\label{8.22:not-partition-sources-cells}
Throughout this section $N=2$, the hypotheses \textup{(H1)}--\textup{(H7)} of Theorem~\ref{3.1:thm-main}
are in force, $\varrho>0$ is the constant of \textup{(H7)}, $g\le\gamma_J$, $0<g_0\le g_-(g)$,
$K=K(g_0,g)$ and $x_0=g_0^{-2}$. The bare lattice on $\mathbb{T}_m$ (Definition~\ref{1.1:def-lattice})
has spacing $\varepsilon=L^{-K}$ and sites $\varepsilon(\mathbb{Z}+\tfrac12)^4$.

\emph{Partition functions.} For a complex function $v$ on $\mathbb{T}_m$ we put
\begin{equation}\label{8.22:eq-partition-sources-cells-1}
Z_m(v)=\int\exp\bigl[-x_0A(U)+\mathcal{O}(v)\bigr]\,dU ,
\end{equation}
with $A(U)$ the Wilson action of Definition~\ref{1.1:def-wilson-action}, $\mathcal{O}(v)$ the smeared
curvature field of Definition~\ref{1.3:def-curvature-field}, and $dU$ the normalised Haar measure on the
bond variables of $\mathbb{T}_m$; no normalising constant is included. We write
$\langle\cdot\rangle_m$ for the expectation in the probability measure
$Z_m(0)^{-1}\exp[-x_0A(U)]\,dU$.

\emph{Complex sources.} $\mathcal{W}$ is the class of complex Lipschitz functions $v$ with
$\|v\|_{\sharp}<\varrho$. For $m\ge m_{\mathbb{T}}(g_-(g))$ the correlation theorem supplies on
$\mathbb{T}_m$ the following.
\begin{itemize}
\item[(a)] Numbers $e_j(X;v)$, for $j\le K$ and $X$ in the family $\mathcal{D}_j$ of localisation
domains at level $j$, with $e_j(X;0)=0$ for all $j$ and $X$, such that, with
\begin{equation}\label{8.22:eq-partition-sources-cells-2}
\mathcal{E}_m(v)=\sum_{j\le K}\ \sum_{X\in\mathcal{D}_j}e_j(X;v),
\end{equation}
so that $\mathcal{E}_m(0)=0$, one has, for $v\in\mathcal{W}$,
\begin{equation}\label{8.22:eq-partition-sources-cells-3}
\frac{Z_m(v)}{Z_m(0)}=e^{\mathcal{E}_m(v)}\,\frac{\widetilde{Z}_m(v)}{\widetilde{Z}_m(0)} ;
\end{equation}
for finitely many real Lipschitz functions $f_i$ and $v=\sum_i z_if_i$, the numbers $e_j(X;v)$ and
$\widetilde{Z}_m(v)$ depend holomorphically on $z=(z_1,z_2,\dots)$ in the set where
$\sum_i|z_i|\,\|f_i\|_{\sharp}<\varrho$, on which $v\in\mathcal{W}$.
\item[(b)] For $v\in\mathcal{W}$, with the constant $\delta$ of the correlation theorem,
\begin{gather*}
|e_j(X;v)|\le C_{\mathrm{vac}}\,e^{-(1-\delta)\kappa d_j(X)},\qquad
|\widetilde{Z}_m(v)|\le e^{E_+^{\sharp}(2L^m)^4},\\
\widetilde{Z}_m(0)\ge e^{-C_{\mathrm{low}}(2L^m)^4}.
\end{gather*}
\item[(c)] $e_j(X;v)$ depends on $v$ only through the values of $v$ at the barycentres of the
plaquettes of the enlarged domain $X^{+}$ attached to $X$ by the correlation theorem.
\item[(d), (e)] The geometric and the counting statements of the correlation theorem.
\end{itemize}
Here $\widetilde{Z}_m$ is the normalised partition function of the correlation theorem,
\begin{equation}\label{8.22:eq-partition-sources-cells-4}
\widetilde{Z}_m(v)=e^{-E_m-\mathcal{E}_m(v)}Z_m(v),\qquad
E_m=\sum_{j<K}E^{(j)}_{\mathrm{vac}}(\mathbb{T}_m,\mathbb{T}_m),
\end{equation}
where $E_m$ is Ba{\l}aban's normalisation constant of the bare density
$\rho_0=\exp[-x_0A(U)-E_m]$ on $\mathbb{T}_m$ (the density $\rho_0$ of
Definition~\ref{1.5:def-densities}) and $E^{(j)}_{\mathrm{vac}}(\mathbb{T}_m,\mathbb{T}_m)$
is its part from level $j$. Since $E_m$ does not depend on $v$ and $\mathcal{E}_m(0)=0$, dividing
\eqref{8.22:eq-partition-sources-cells-4} at $v$ by \eqref{8.22:eq-partition-sources-cells-4} at $v=0$
(where $\widetilde{Z}_m(0)>0$ by (b)) gives \eqref{8.22:eq-partition-sources-cells-3}.
The constants $E_+^{\sharp}$ and $C_{\mathrm{low}}$ do not depend on $K$, $g_0$ or $m$, and we put
\begin{equation}\label{8.22:eq-partition-sources-cells-5}
\mathsf{B}=E_+^{\sharp}+C_{\mathrm{low}} ;
\end{equation}
then $\mathsf{B}\ge0$ and $\mathsf{B}$ does not depend on $K$, $g_0$ or $m$. Moreover, by the
logarithmic degree bound for the stability constants $E_+^{\sharp}$, $C_{\mathrm{low}}$, one has
$\mathsf{B}\le C_{\mathsf{B}}(1+\log g^{-2})$ with the constant $C_{\mathsf{B}}$ of that bound.
The factor $\mathfrak{r}\in\{1,2\}$ records how (a)--(b) are applied, in the volume transfer
below, to complex weights obtained by extending sources mirror-periodically over the cells defined
below: $\mathfrak{r}=1$ when such a weight is used itself as an element of $\mathcal{W}$, and
$\mathfrak{r}=2$ when (a)--(b) are applied instead to an associated system of four real functions at
the pair of parameters $(z,\bar z)$. For real Lipschitz $F_1,F_2$ we write $w_z=z_1F_1+z_2F_2$, and the
complex parameters $z\in\mathbb{C}^2$ of such weights are accordingly restricted to the polydisc
\begin{equation*}
\mathbb{D}_{\mathfrak{r}}=\Bigl\{z\in\mathbb{C}^2:\ \sum_{i=1}^{2}|z_i|\,\|F_i\|_{\sharp}
<\varrho/\mathfrak{r}\Bigr\},
\end{equation*}
so that $\mathfrak{r}=2$ halves the admissible radius. The two readings are those fixed in the
per-ball volume-free bound.

\emph{Cells.} $R>0$ is a radius (below $R=1$). The cell side $\ell$ is the one used by the per-ball
volume-free bound for $n=2$ sources of radius $R$: the least power of $L$ not smaller than $16(R+1)$.
The value $\ell(2,R,g)\ge L^{m_{\mathbb{T}}(g_-(g))}$ of clause \textup{(iii-b)} of
Theorem~\ref{3.1:thm-main} may be used instead; only the size of the depth floor $s_T(g)$ of the
volume transfer changes. $C_{\mathrm{face}}(R)$ is the face constant of the per-ball volume-free bound.
We put
\begin{equation}\label{8.22:eq-partition-sources-cells-6}
m_{\star}(g)=\max\bigl\{m_{\mathbb{T}}(g_-(g)),\ \log_L\ell\bigr\}.
\end{equation}
A function on $\mathbb{R}^4$ supported in a ball of radius $<L^m/2$ is transported to $\mathbb{T}_m$
through the covering map $\mathbb{R}^4\to\mathbb{T}_m$, which is injective on such a ball. The cells
are $c_k=\ell k+[0,\ell]^4$, $k\in\mathbb{Z}^4$ (here $k$ is a cell index, not a level). For
$\mu\in\{1,2,3,4\}$ and $a\in\mathbb{Z}$,
$\theta_{\mu,a}$ is the reflection $x_\mu\mapsto2a\ell-x_\mu$, the other coordinates fixed;
$\mathfrak{G}_{\ell}$ is the group they generate, acting on $\mathbb{R}^4$ and on every torus whose
sides are even multiples of $\ell$. For $\theta\in\mathfrak{G}_{\ell}$ the parity
$\epsilon(\theta)=\pm1$ is the determinant of the linear part of $\theta$, so that
$\epsilon(\theta)=(-1)^r$ when $\theta$ is a product of $r$ generators.

At $R=1$, with $\ell$ the least power of $L$ not smaller than $32$, one has
$m_{\star}(g)=m_{\mathbb{T}}(g_-(g))$: $\mathbb{T}_{m_{\star}(g)}$ is the smallest
torus of the family admitted by \textup{(H5)}.
\end{definition}

\begin{proof}
\emph{$\mathsf{B}\ge0$.} Fix any $m\ge m_{\mathbb{T}}(g_-(g))$. Since $\|0\|_{\sharp}=0<\varrho$, the
source $v=0$ lies in $\mathcal{W}$, and property (b) at $v=0$ gives, with $V=(2L^m)^4>0$,
\begin{equation*}
e^{-C_{\mathrm{low}}V}\le\widetilde{Z}_m(0)=|\widetilde{Z}_m(0)|\le e^{E_+^{\sharp}V}.
\end{equation*}
Taking logarithms and dividing by $V$ gives $-C_{\mathrm{low}}\le E_+^{\sharp}$, that is
$\mathsf{B}\ge0$. Independence of $K$, $g_0$ and $m$ holds because it holds for both summands.

\emph{$m_{\star}(g)=m_{\mathbb{T}}(g_-(g))$ at $R=1$.} Here $\ell$ is the least power of $L$ not
smaller than $32$. Since $1<32$, one has $\ell\ge L$, and $\ell/L$ is a power of $L$, smaller than $32$
by minimality; hence $\ell<32L$. As $L$ is an odd integer larger than one
(Definition~\ref{1.1:def-lattice}), $L\ge3$ and $L^{10}\ge3^{10}>100$; with $M\ge1$, a power of $L$,
\begin{equation*}
\ell<32L\le100L\le L^{11}\le L^{11}M .
\end{equation*}
By the definition $m_{\mathbb{T}}(g')=11+\log_L M+\log_L R(g')$, where $R(\cdot)$ is the volume function
of Definition~\ref{1.2:def-flow-constants} (not the ball radius), whose values satisfy $R(g')\ge1$, the
right side is at most $L^{m_{\mathbb{T}}(g_-(g))}$. Hence $\log_L\ell\le m_{\mathbb{T}}(g_-(g))$, and
the maximum in \eqref{8.22:eq-partition-sources-cells-6} equals $m_{\mathbb{T}}(g_-(g))$. The tori of
the family are the $\mathbb{T}_m$ with $m\ge m_{\mathbb{T}}(g_-(g))$, so $\mathbb{T}_{m_{\star}(g)}$ is
the smallest.
\end{proof}

\begin{lemma}[Trace sub-additivity of mirror-periodic partition functions]
\label{8.22:lem-trace-subadditivity}
Write $\varepsilon:=L^{-K}$ for the bare spacing, as in
Notation~\ref{8.22:not-partition-sources-cells}, let $x_0=g_0^{-2}$, and let $\ell>0$ be a length with
$\ell/\varepsilon\in\mathbb{N}$ (the lemma is applied with $\ell$ the cell side of
Notation~\ref{8.22:not-partition-sources-cells}, for which this holds). For
$\mathsf{n}=(\mathsf{n}_0,\dots,\mathsf{n}_3)$ with every $\mathsf{n}_\mu$ an
even positive integer, let
\begin{equation*}
\mathbb{T}[\mathsf{n}]:=\prod_{\mu=0}^{3}\mathbb{R}/(\mathsf{n}_\mu\ell)\mathbb{Z},
\end{equation*}
with bare sites $\varepsilon(\mathbb{Z}+\tfrac12)^4$, and for a bounded
$W\colon\mathbb{T}[\mathsf{n}]\to\mathbb{C}$ (equivalently, a bounded $W$ on $\mathbb{R}^4$ with period
$\mathsf{n}_\mu\ell$ in the direction $\mu$) put
\begin{equation*}
Z_{\mathsf{n}}(W):=\int\exp\bigl[-x_0A(U)+\mathcal{O}(W)\bigr]\,dU
\quad\text{on }\mathbb{T}[\mathsf{n}],
\end{equation*}
where $A(U)$ and $\mathcal{O}(W)=\sum_pW(x(p))[1-\operatorname{Re}\operatorname{tr}U(\partial p)]$ are given by the
formulas of Definitions~\ref{1.1:def-wilson-action} and~\ref{1.3:def-curvature-field}, read on
$\mathbb{T}[\mathsf{n}]$ with the complex weight $W$, and $dU$ is the normalised Haar measure. For
$a\in\mathbb{Z}$ let $\theta_{\mu,a}$ be the reflection $x_\mu\mapsto 2a\ell-x_\mu$ (the other coordinates
fixed). A bounded $W\colon\mathbb{R}^4\to\mathbb{C}$ is called \emph{mirror-periodic} if
\begin{itemize}
\item[(M1)] $W\circ\theta_{\mu,a}=\overline{W}$ for all $\mu$ and all $a\in\mathbb{Z}$;
\item[(M2)] $W(x(p))=0$ for every plaquette $p$ containing a bond that crosses a plane
$\{x_\mu\in\ell\mathbb{Z}\}$.
\end{itemize}
By (M1), $W$ is $2\ell$-periodic in every direction and descends to every $\mathbb{T}[\mathsf{n}]$.

Let $W$ be mirror-periodic. For every $\mathsf{n}$ as above, every direction $\mu$ and every integer
$k\ge1$, let $\mathsf{n}^{(\mu,k)}$ be $\mathsf{n}$ with $\mathsf{n}_\mu$ replaced by $k\mathsf{n}_\mu$. Then
\begin{equation}\label{8.22:eq-trace-subadditivity-1}
0\le Z_{\mathsf{n}^{(\mu,k)}}(W)\le Z_{\mathsf{n}}(W)^{k}.
\end{equation}
Hence
\begin{equation}\label{8.22:eq-trace-subadditivity-2}
0\le Z_{k\mathsf{n}}(W)\le Z_{\mathsf{n}}(W)^{k^4}.
\end{equation}
\end{lemma}

\begin{proof}
The argument rests on the positivity of the Wilson transfer operator \cite{OSe78}, written out in
the form used below.

\emph{Periodicity.} The composition $\theta_{\mu,a+1}\circ\theta_{\mu,a}$ sends $x_\mu$ to
$2(a+1)\ell-(2a\ell-x_\mu)=x_\mu+2\ell$ and fixes the other coordinates; that is, it is the
translation by $2\ell$ in the direction $\mu$. Applying (M1) twice, for every $x\in\mathbb{R}^4$
\begin{equation*}
W\bigl(\theta_{\mu,a+1}(\theta_{\mu,a}x)\bigr)
=\overline{W(\theta_{\mu,a}x)}=\overline{\overline{W(x)}}=W(x),
\end{equation*}
so $W$ is $2\ell$-periodic in the direction $\mu$. Since every side $\mathsf{n}_\mu\ell$ of
$\mathbb{T}[\mathsf{n}]$ is a multiple of $2\ell$ ($\mathsf{n}_\mu$ being even), $W$ descends to
$\mathbb{T}[\mathsf{n}]$; the same holds for $\mathbb{T}[\mathsf{n}^{(\mu,k)}]$, since the entries of
$\mathsf{n}^{(\mu,k)}$ are again even.

Since (M1), (M2) and the definition of $Z_{\mathsf{n}}$ are symmetric in the four directions, it
suffices to prove \eqref{8.22:eq-trace-subadditivity-1} for $\mu=0$, which we call time. We write
$x^0$ for the time coordinate; it is not to be confused with the inverse bare coupling $x_0=g_0^{-2}$.
For a plaquette $p$ put $a_p(U):=1-\operatorname{Re}\operatorname{tr}U(\partial p)$, so that
\begin{equation*}
\exp\bigl[-x_0A(U)+\mathcal{O}(W)\bigr]=\prod_p\exp\bigl[-(x_0-W(x(p)))\,a_p(U)\bigr].
\end{equation*}

\emph{Layers and transfer operators.} A \emph{layer} is the set of sites with time coordinate
$x^0=\varepsilon(i-\tfrac12)$, $i\in\mathbb{Z}/(\mathsf{n}_0\ell/\varepsilon)\mathbb{Z}$. Let $\Sigma$ be the
cross-section $3$-torus $\prod_{\nu=1}^3\mathbb{R}/(\mathsf{n}_\nu\ell)\mathbb{Z}$ with sites
$\varepsilon(\mathbb{Z}+\tfrac12)^3$, let $\mathsf{E}$ be its set of bonds, each $b\in\mathsf{E}$ oriented
from an endpoint $b_-$ to an endpoint $b_+$, and let
$\mathcal{H}_{\Sigma}:=L^2(\mathrm{SU}(N)^{\mathsf{E}},\text{Haar})$. Translating in $x^0$ identifies the
spatial bonds of every layer with $\mathsf{E}$, so the configuration on a layer is an element
$U\in\mathrm{SU}(N)^{\mathsf{E}}$. Every other bond is time-like: it joins layer $i$ to layer $i+1$ over a
site $x\in\Sigma$, and we write $V_x$ for its variable. Every plaquette is of exactly one of two kinds:
either it is spatial and lies in one layer $i$, or it is time-like and is spanned by a bond
$b\in\mathsf{E}$ and the time direction between layers $i$ and $i+1$; we write $p(b,i)$ for it. Up to
cyclic order and orientation, which do not change $\operatorname{Re}\operatorname{tr}$, the parallel transporter around
$\partial p(b,i)$ is $U_bV_{b_+}U_b'^{-1}V_{b_-}^{-1}$, where $U$ and $U'$ are the configurations on
layers $i$ and $i+1$.

For a layer $i$ let $\mathsf{M}_i=\mathsf{M}_i[W]$ be multiplication on $\mathcal{H}_{\Sigma}$ by
\begin{equation*}
\exp\Bigl[-\sum_{p\text{ spatial in layer }i}\bigl(x_0-W(x(p))\bigr)a_p(U)\Bigr],
\end{equation*}
and for consecutive layers $i,i+1$ let $\mathsf{K}_i=\mathsf{K}_i[W]$ be the integral operator on
$\mathcal{H}_{\Sigma}$ with the continuous kernel
\begin{equation}\label{8.22:eq-trace-subadditivity-3}
\begin{split}
\mathsf{K}_i[W](U,U')
:={}&\int\prod_{x\in\Sigma}dV_x\prod_{b\in\mathsf{E}}
\exp\Bigl\{-\bigl(x_0-W(x(p(b,i)))\bigr)\\
&\qquad\times\bigl[1-\operatorname{Re}\operatorname{tr}(U_bV_{b_+}U_b'^{-1}V_{b_-}^{-1})\bigr]\Bigr\}.
\end{split}
\end{equation}
Since every plaquette is of one of the two kinds, Fubini's theorem applied to the integral over the
layer configurations and the time-like variables gives
\begin{equation}\label{8.22:eq-trace-subadditivity-4}
Z_{\mathsf{n}}(W)=\operatorname{Tr}\bigl(\mathsf{M}_1\mathsf{K}_1\mathsf{M}_2\mathsf{K}_2\cdots\bigr),
\end{equation}
the ordered product running once around the time circle. Here there are
$\mathsf{n}_0\ell/\varepsilon\ge2$ layers, so the cyclic product contains at least two of the operators
$\mathsf{K}_i$, each Hilbert--Schmidt since its kernel is continuous on a compact space, while the
$\mathsf{M}_i$ are bounded. Grouping the cyclic product into two factors, each containing one
$\mathsf{K}_i$, exhibits it as a product of two Hilbert--Schmidt operators $\mathsf{F}_1\mathsf{F}_2$;
it is therefore trace class, and its trace is
$\int\!\!\int\mathsf{F}_1(U,U')\mathsf{F}_2(U',U)\,dU\,dU'$, which after unfolding the kernels is
the integral defining $Z_{\mathsf{n}}(W)$.

\emph{Step 1: adjoints.} Since $x_0$ and $a_p(U)$ are real, the complex conjugate of the function
defining $\mathsf{M}_i[W]$ is the function defining $\mathsf{M}_i[\overline{W}]$; hence
$\mathsf{M}_i[W]^*=\mathsf{M}_i[\overline{W}]$. The adjoint of $\mathsf{K}_i[W]$ has the kernel
$\overline{\mathsf{K}_i[W](U',U)}$. In $\mathsf{K}_i[W](U',U)$ the integrand contains
$\operatorname{Re}\operatorname{tr}(U'_bV_{b_+}U_b^{-1}V_{b_-}^{-1})$. Substitute $V\mapsto V^{-1}$, which
leaves the Haar measure invariant; the trace becomes
$\operatorname{Re}\operatorname{tr}(U'_bV_{b_+}^{-1}U_b^{-1}V_{b_-})$. For $Y\in\mathrm{SU}(N)$,
$\operatorname{tr}Y^{-1}=\operatorname{tr}Y^*=\overline{\operatorname{tr}Y}$, so
$\operatorname{Re}\operatorname{tr}(Y^{-1})=\operatorname{Re}\operatorname{tr}Y$. With
$Y=U'_bV_{b_+}^{-1}U_b^{-1}V_{b_-}$ and cyclicity of the trace,
\begin{equation*}
\operatorname{Re}\operatorname{tr}\bigl(U'_bV_{b_+}^{-1}U_b^{-1}V_{b_-}\bigr)
=\operatorname{Re}\operatorname{tr}\bigl(V_{b_-}^{-1}U_bV_{b_+}U_b'^{-1}\bigr)
=\operatorname{Re}\operatorname{tr}\bigl(U_bV_{b_+}U_b'^{-1}V_{b_-}^{-1}\bigr).
\end{equation*}
Thus the kernel \eqref{8.22:eq-trace-subadditivity-3} is symmetric in $(U,U')$, for any values of
the coefficients, and taking the complex conjugate replaces $W$ by $\overline{W}$ in the
coefficients. Hence $\mathsf{K}_i[W]^*=\mathsf{K}_i[\overline{W}]$.

\emph{Step 2: the face operator.} Since $\ell/\varepsilon\in\mathbb{N}$, a plane $\{x^0=a\ell\}$
contains no site; it lies between two consecutive layers, and every time-like plaquette between
these two layers contains a time-like bond crossing it. By (M2), $W$ vanishes at the barycentres of
these plaquettes, so the operator between these two layers is the free one, $\mathsf{K}^0$, with
kernel
\begin{equation*}
\mathsf{K}^0(U,U')=\int\prod_{x\in\Sigma}dV_x\prod_{b\in\mathsf{E}}
\varphi\bigl(U_bV_{b_+}U_b'^{-1}V_{b_-}^{-1}\bigr),
\qquad \varphi(Y):=e^{-x_0(1-\operatorname{Re}\operatorname{tr}Y)};
\end{equation*}
it is the same operator at every face plane. Let $\mathsf{C}$ be convolution on
$\mathrm{SU}(N)^{\mathsf{E}}$ by $\Phi(U):=\prod_{b\in\mathsf{E}}\varphi(U_b)$, that is
$(\mathsf{C}\psi)(U)=\int dU'\prod_b\varphi(U_bU_b'^{-1})\psi(U')$; for
$V\in\mathrm{SU}(N)^{\Sigma}$ let $(R_V\psi)(U):=\psi\bigl((V_{b_-}^{-1}U_bV_{b_+})_b\bigr)$; and let
$\mathsf{P}:=\int dV\,R_V$, the orthogonal projection onto the gauge-invariant functions. Then
\begin{equation*}
(\mathsf{C}\mathsf{P}\psi)(U)=\int dV\int dU'\prod_b\varphi(U_bU_b'^{-1})\,
\psi\bigl((V_{b_-}^{-1}U'_bV_{b_+})_b\bigr),
\end{equation*}
and the substitution $U'_b=V_{b_-}U''_bV_{b_+}^{-1}$, which preserves Haar measure, turns
$\varphi(U_bU_b'^{-1})$ into $\varphi(U_bV_{b_+}U_b''^{-1}V_{b_-}^{-1})$ and the argument of $\psi$ into
$U''$; so $\mathsf{K}^0=\mathsf{C}\mathsf{P}$. Since $\varphi$ is a class function, the substitution
$U'_b=V_{b_-}^{-1}U''_bV_{b_+}$, which preserves Haar measure, gives in $(R_V\mathsf{C}\psi)(U)$
\begin{equation*}
\varphi\bigl(V_{b_-}^{-1}U_bV_{b_+}U_b'^{-1}\bigr)
=\varphi\bigl(V_{b_-}^{-1}U_bU_b''^{-1}V_{b_-}\bigr)
=\varphi\bigl(U_bU_b''^{-1}\bigr),
\end{equation*}
while the argument $U'$ of $\psi$ becomes $(V_{b_-}^{-1}U''_bV_{b_+})_b$, so that
$\psi(U')=(R_V\psi)(U'')$; whence $R_V\mathsf{C}=\mathsf{C}R_V$ for every $V$, and so
$\mathsf{C}\mathsf{P}=\mathsf{P}\mathsf{C}$. As
$\mathsf{P}^2=\mathsf{P}$, we obtain $\mathsf{K}^0=\mathsf{C}\mathsf{P}=\mathsf{P}\mathsf{C}\mathsf{P}$.

We show that $\mathsf{C}\ge0$. Let $\chi$ be the character of the defining representation of
$\mathrm{SU}(N)$ on $\mathbb{C}^N$. Since $\operatorname{tr}$ is normalised to take the value $1$ at the
identity, $\operatorname{Re}\operatorname{tr}Y=(\chi(Y)+\overline{\chi(Y)})/(2N)$, and
$\chi_\pi:=\chi+\overline{\chi}$ is the character of the representation $\pi$ given by the direct sum of
the defining representation and its complex conjugate. Hence
\begin{equation*}
\varphi(Y)=e^{-x_0}\sum_{n\ge0}\frac{1}{n!}\Bigl(\frac{x_0}{2N}\Bigr)^{n}\chi_\pi(Y)^n ,
\end{equation*}
the series converging uniformly on $\mathrm{SU}(N)$, and $\chi_\pi^n$ is the character of the tensor
power $\pi^{\otimes n}$, which is a finite direct sum of irreducible representations. Collecting
terms, $\varphi=\sum_\sigma c_\sigma\chi_\sigma$, the sum over the irreducible representations $\sigma$
of $\mathrm{SU}(N)$ with characters $\chi_\sigma$ and coefficients $c_\sigma\ge0$; evaluating at the
identity gives $\sum_\sigma c_\sigma\dim\sigma=\varphi(\mathbf{1})<\infty$. On $L^2(\mathrm{SU}(N))$
convolution by $\chi_\sigma$ equals $(\dim\sigma)^{-1}\Pi_\sigma$, where $\Pi_\sigma$ is the orthogonal
projection onto the $\sigma$-isotypic subspace (Schur orthogonality); hence convolution by $\varphi$ on
$L^2(\mathrm{SU}(N))$ is the norm-convergent sum $\sum_\sigma c_\sigma(\dim\sigma)^{-1}\Pi_\sigma$, a
non-negative operator. Under the identification
$\mathcal{H}_\Sigma=\bigotimes_{b\in\mathsf{E}}L^2(\mathrm{SU}(N))$, convolution by the product function
$\Phi$ is the tensor product over $b\in\mathsf{E}$ of the convolutions by $\varphi$, and a tensor product
of non-negative operators is non-negative; so $\mathsf{C}\ge0$. Hence
$\langle\psi,\mathsf{K}^0\psi\rangle=\langle\mathsf{P}\psi,\mathsf{C}\mathsf{P}\psi\rangle\ge0$, i.e.
$\mathsf{K}^0\ge0$. Its kernel is continuous, so by Mercer's theorem $\mathsf{K}^0$ is trace class, and
$\mathsf{S}:=(\mathsf{K}^0)^{1/2}$ is a self-adjoint Hilbert--Schmidt operator.

\emph{Step 3: mirror structure.} For $a\in\{1,\dots,\mathsf{n}_0\}$, read modulo $\mathsf{n}_0$, the
\emph{slab} $a$ consists of the $\ell/\varepsilon$ layers strictly between the planes $x^0=a\ell$ and
$x^0=(a+1)\ell$. Let $\mathsf{B}_a:=\mathsf{M}\mathsf{K}\mathsf{M}\cdots\mathsf{K}\mathsf{M}$ be the ordered product of
the multiplications $\mathsf{M}_i$ of the layers of slab $a$ and of the operators $\mathsf{K}_i$ joining
two layers of slab $a$; it is bounded. Between slab $a-1$ and slab $a$ stands the face operator
$\mathsf{K}^0=\mathsf{S}\mathsf{S}$. Put $\mathsf{A}_a:=\mathsf{S}\mathsf{B}_a\mathsf{S}$, a trace class operator. In
\eqref{8.22:eq-trace-subadditivity-4}, begin the cyclic product at a face operator, write each
$\mathsf{K}^0$ as $\mathsf{S}\mathsf{S}$, and move the first factor $\mathsf{S}$ to the end of the product
(admissible since the remaining product contains $\mathsf{K}^0$ and is therefore trace class, while
$\mathsf{S}$ is bounded); this gives
\begin{equation}\label{8.22:eq-trace-subadditivity-5}
Z_{\mathsf{n}}(W)=\operatorname{Tr}\Bigl(\prod_{a=1}^{\mathsf{n}_0}\mathsf{A}_a\Bigr).
\end{equation}
The reflection $\theta_{0,a+1}\colon x^0\mapsto2(a+1)\ell-x^0$ fixes the spatial coordinates and maps
the layer $x^0=\varepsilon(i-\tfrac12)$ to the layer $x^0=\varepsilon(j-\tfrac12)$ with
$j=2(a+1)\ell/\varepsilon-i+1$; thus it maps the layers of slab $a$ onto those of slab $a+1$ in
reversed order, each spatial plaquette of a layer onto the plaquette over the same spatial position
in the image layer, and $p(b,i)$ onto the time-like plaquette over the same $b$ between the image
layers. By (M1) the value of $W$ at the barycentre of each image plaquette is the complex conjugate
of its value at the barycentre of the original plaquette. Hence, if $i_1,\dots,i_{\ell/\varepsilon}$
are the layers of slab $a$ and
\begin{equation*}
\mathsf{B}_a=\mathsf{M}_{i_1}[W]\mathsf{K}_{i_1}[W]\mathsf{M}_{i_2}[W]\cdots
\mathsf{K}_{i_{\ell/\varepsilon-1}}[W]\mathsf{M}_{i_{\ell/\varepsilon}}[W],
\end{equation*}
then, using Step~1 (the symmetry of the kernels in $(U,U')$ and the adjoint formulas
$\mathsf{M}_i[W]^*=\mathsf{M}_i[\overline{W}]$, $\mathsf{K}_i[W]^*=\mathsf{K}_i[\overline{W}]$),
\begin{align*}
\mathsf{B}_{a+1}
&=\mathsf{M}_{i_{\ell/\varepsilon}}[\overline{W}]\mathsf{K}_{i_{\ell/\varepsilon-1}}[\overline{W}]\cdots
\mathsf{K}_{i_1}[\overline{W}]\mathsf{M}_{i_1}[\overline{W}]\\
&=\mathsf{M}_{i_{\ell/\varepsilon}}[W]^*\mathsf{K}_{i_{\ell/\varepsilon-1}}[W]^*\cdots
\mathsf{K}_{i_1}[W]^*\mathsf{M}_{i_1}[W]^*\\
&=\mathsf{B}_a^*.
\end{align*}
As $\mathsf{S}$ is self-adjoint, $\mathsf{A}_{a+1}=\mathsf{S}\mathsf{B}_a^*\mathsf{S}=\mathsf{A}_a^*$, and so
$\mathsf{A}_{a+2}=\mathsf{A}_a$. Since $\mathsf{n}_0$ is even,
\begin{equation*}
\prod_{a=1}^{\mathsf{n}_0}\mathsf{A}_a
=\mathsf{A}_1\mathsf{A}_1^*\mathsf{A}_1\mathsf{A}_1^*\cdots\mathsf{A}_1\mathsf{A}_1^*
=(\mathsf{A}_1\mathsf{A}_1^*)^{\mathsf{n}_0/2}=:\mathcal{T},
\end{equation*}
and $\mathcal{T}\ge0$ is trace class. The operator $\mathsf{A}_1$ depends only on $x_0$, on the
cross-section $\Sigma$ and on the values of $W$ at the plaquettes of slab $1$ and of its two
bounding faces; these are the same on $\mathbb{T}[\mathsf{n}^{(0,k)}]$, which has the same cross-section
and the time extent $k\mathsf{n}_0\ell$, again an even multiple of $\ell$. The same computation on
$\mathbb{T}[\mathsf{n}^{(0,k)}]$ therefore gives the product $(\mathsf{A}_1\mathsf{A}_1^*)^{k\mathsf{n}_0/2}$, and by
\eqref{8.22:eq-trace-subadditivity-5}
\begin{equation}\label{8.22:eq-trace-subadditivity-6}
Z_{\mathsf{n}^{(0,k)}}(W)=\operatorname{Tr}(\mathcal{T}^k),\qquad
Z_{\mathsf{n}}(W)=\operatorname{Tr}\mathcal{T}.
\end{equation}

\emph{Step 4: the eigenvalue inequality.} Let $\lambda_1,\lambda_2,\dots\ge0$ be the eigenvalues of
the non-negative trace class operator $\mathcal{T}$, repeated according to multiplicity. By the
spectral theorem $\operatorname{Tr}\mathcal{T}^k=\sum_i\lambda_i^k$, and expanding $(\sum_i\lambda_i)^k$
produces the terms $\lambda_i^k$ together with further non-negative terms, so
\begin{equation*}
0\le\sum_i\lambda_i^k\le\Bigl(\sum_i\lambda_i\Bigr)^k .
\end{equation*}
By \eqref{8.22:eq-trace-subadditivity-6} this is \eqref{8.22:eq-trace-subadditivity-1} for $\mu=0$;
in particular $Z_{\mathsf{n}}(W)=\operatorname{Tr}\mathcal{T}\ge0$ (the case $k=1$).

\emph{Step 5: iteration over the directions.} The hypotheses are symmetric in the directions, and
\eqref{8.22:eq-trace-subadditivity-1} has been proved for every quadruple of even positive
integers, including rectangular ones. Apply it successively in the directions $0,1,2,3$ to the
quadruples $(\mathsf{n}_0,k\mathsf{n}_1,k\mathsf{n}_2,k\mathsf{n}_3)$,
$(\mathsf{n}_0,\mathsf{n}_1,k\mathsf{n}_2,k\mathsf{n}_3)$, $(\mathsf{n}_0,\mathsf{n}_1,\mathsf{n}_2,k\mathsf{n}_3)$ and $\mathsf{n}$,
all of whose entries are even. Since all the numbers involved are non-negative and $t\mapsto t^k$
is non-decreasing on $[0,\infty)$,
\begin{align*}
0\le Z_{k\mathsf{n}}(W)
&\le Z_{(\mathsf{n}_0,k\mathsf{n}_1,k\mathsf{n}_2,k\mathsf{n}_3)}(W)^{k}
\le Z_{(\mathsf{n}_0,\mathsf{n}_1,k\mathsf{n}_2,k\mathsf{n}_3)}(W)^{k^2}\\
&\le Z_{(\mathsf{n}_0,\mathsf{n}_1,\mathsf{n}_2,k\mathsf{n}_3)}(W)^{k^3}
\le Z_{\mathsf{n}}(W)^{k^4},
\end{align*}
which is \eqref{8.22:eq-trace-subadditivity-2}.
\end{proof}

\begin{remark}[Scope of the trace inequality]
(i) The proof of Lemma~\ref{8.22:lem-trace-subadditivity} uses no property of Ba{\l}aban's
construction; of the gauge group $\mathrm{SU}(N)$ it uses only that it is compact and that the
single-plaquette weight $\varphi$ is positive definite, and it holds verbatim for any compact gauge
group with a positive-definite plaquette weight.

(ii) The weight $W$ may be complex and of any size; only (M1) and (M2) are used.

(iii) At $W=0$ the integrand is positive, so $Z_{\mathsf{n}}(0)>0$. Write
$|\mathbb{T}[\mathsf{n}]|=\prod_\mu\mathsf{n}_\mu\ell$ for the volume of $\mathbb{T}[\mathsf{n}]$. Since
$|\mathbb{T}[\mathsf{n}^{(\mu,k)}]|=k\,|\mathbb{T}[\mathsf{n}]|$ and $|\mathbb{T}[k\mathsf{n}]|=k^4|\mathbb{T}[\mathsf{n}]|$,
inequalities \eqref{8.22:eq-trace-subadditivity-1} and \eqref{8.22:eq-trace-subadditivity-2} at
$W=0$ say that the quantity $|\mathbb{T}[\mathsf{n}]|^{-1}\log Z_{\mathsf{n}}(0)$, the logarithm per unit
volume of the partition function without normalisation constant, does not increase when
$\mathsf{n}$ is replaced by $\mathsf{n}^{(\mu,k)}$ or by $k\mathsf{n}$.
\end{remark}

\begin{lemma}[Torus-independence of the vacuum-energy density]\label{8.22:lem-vacuum-energy-density}
Let $N=2$, and let the standing assumptions and the notation of
Definition~\ref{8.22:not-partition-sources-cells} be in force. In this lemma
$E_m$ has the meaning fixed in hypothesis~(a). Throughout this lemma and its proof
$m'$ denotes a second torus exponent; the glossary's exponent of the cube side is given no letter
here, and $L^{m_{\mathbb{T}}(g')}$ is written directly as $L^{11}M\,R(g')$.
Fix integers $m,m'\ge m_{\mathbb{T}}(g_-(g))$; in hypotheses (a)--(c) below $m''$ stands for
either of them, and $j$ ranges over $0\le j<K$.
Assume the following.
\begin{itemize}
\item[(a)] $E_{m''}$ denotes Ba{\l}aban's normalisation constant of the bare density on
$\mathbb{T}_{m''}$, as in the normalisation of $\widetilde{Z}_{m''}$ in
Definition~\ref{8.22:not-partition-sources-cells}, less the Gaussian
normalisers of the $K$ steps, the Gaussian normaliser of step $j$ being the logarithm of the
Gaussian integral of step $j$ at the trivial background; the Gaussian normalisers of the two tori
are compared by the determinant lemma. This number is
$E_{m''}=\sum_{j<K}E^{(j)}_{\mathrm{vac}}(\mathbb{T}_{m''},\mathbb{T}_{m''})$, and for every $j<K$
the level-$j$ part $E^{(j)}_{\mathrm{vac}}(\mathbb{T}_{m''},\mathbb{T}_{m''})$ of $E_{m''}$ splits
into three parts,
\begin{equation*}
E^{(j)}_{\mathrm{vac}}(\mathbb{T}_{m''},\mathbb{T}_{m''})=E^{(j)}_{\mathrm{vac},1}(\mathbb{T}_{m''})
+E^{(j)}_{\mathrm{vac},2}(\mathbb{T}_{m''})+E^{(j)}_{\mathrm{vac},3}(\mathbb{T}_{m''}),
\end{equation*}
where $E^{(j)}_{\mathrm{vac},1}(\mathbb{T}_{m''})$ is a sum of normalisation numbers of step $j$ of
the renormalisation transformation \cite[(1.7)]{B-CMP109}, built
from $\log g_j$, $\log z_j$ and constants of $\mathrm{SU}(2)$, each multiplied by a lattice count of
$\mathbb{T}_{m''}$ at level $j$ or $j+1$; here $z_j$ is the second running constant of step $j$ that
enters these normalisation numbers (it is not to be confused with the source coefficients $z_i$
or with a site $z$);
\begin{equation*}
E^{(j)}_{\mathrm{vac},2}(\mathbb{T}_{m''})=\sum_{z}\sum_{X\ni z}\mathbf{E}^{(j+1)}(X,1,z),
\qquad
E^{(j)}_{\mathrm{vac},3}(\mathbb{T}_{m''})=\sum_{X}\mathbf{R}^{(j+1)}(X,1),
\end{equation*}
with $z$ running over the sites of level $j+1$ on $\mathbb{T}_{m''}$, with $X$ running over the
localisation domains $\mathcal{D}_{j+1}$ on $\mathbb{T}_{m''}$, and with the effective-action terms
$\mathbf{E}^{(j+1)}$, $\mathbf{R}^{(j+1)}$ of Ba{\l}aban's format
\cite[(2.27)(i)--(iv)]{B-CMP119} evaluated at the trivial configuration $U=1$.
\item[(b)] Ba{\l}aban's format at zero source, in the form stated in
\cite[(2.27)(i)--(iv), (2.29), (2.31)--(2.33)]{B-CMP119}: by
\cite[(2.29)]{B-CMP119}, at the trivial configuration $U=1$ the unrestricted sum
\begin{equation*}
\mathbf{E}^{(j+1)}(1,z;\mathbb{T}_{m''})
:=\sum_{X\in\mathcal{D}_{j+1},\,X\ni z}\mathbf{E}^{(j+1)}(X,1,z)
\end{equation*}
does not depend on $z$; the $\mathbf{E}$-terms obey $|\mathbf{E}^{(j+1)}(X,1,z)|\le E_0e^{-\kappa d_{j+1}(X)}$;
the $\mathbf{R}$-terms are covariant, \cite[(2.32)--(2.33)]{B-CMP119}, under the translations
preserving the partition of the level-$(j+1)$ lattice into cubes of side $M$, and obey the bound
\cite[(2.31)]{B-CMP119}, which is used in the form
\begin{equation*}
|\mathbf{R}^{(j+1)}(X,1)|\le e^{-\kappa d_{j+1}(X)} .
\end{equation*}
The constant $K_*$ is that of \cite[(1.26)]{B-CMP116}, read at the rate $\kappa/2$:
$\sum_{X\ni z}e^{-\kappa d_{j+1}(X)/2}\le K_*$; \cite[(1.26)]{B-CMP116} permits this rate because
$\kappa$ is chosen in the order of choice of $\mathfrak{p}$
(Definition~\ref{2.3:not-stages-middle}) at least twice the threshold rate of that estimate in
dimension four. The
constants $E_0$ of \cite[(2.27)(i)--(iv)]{B-CMP119} and
$K_*$ of \cite[(1.26)]{B-CMP116} depend on $L$ and $\mathfrak{p}$ only; $\kappa$ is one of the
parameters $\mathfrak{p}$.
\item[(c)] Clause~(0) of Theorem~\ref{3.1:thm-main}, and the torus-independence of non-wrapping
terms in the format of clause~(i-a) of Theorem~\ref{3.1:thm-main} as supplied by parts (i) and
(iii)(b) of clause~(0)'s torus-independence lemma together
with the lift lemma, in the following form. The torus-independence lemma is proved as an
implication from lemmas proved earlier and Ba{\l}aban's definitions
\cite[p.~394; Corollary~3.8, p.~410]{B-CMP99b}; the convergence and the
analytic extension of the expansions on the joined localisation domains enter it as stated in
\cite[p.~192, (1.70)]{B-CMP122a} and \cite[p.~377, (1.68)--(1.69)]{B-CMP122b}, the hypothesis
of those statements on the towers of localisation domains, namely their admissibility for the
theorems of \cite{B-CMP99b}, being supplied by part~(b) of the locality lemma for the pieces of
the renormalisation operation.
It is assumed, as part of clause~(0) and of its
torus-independence lemma, that the running constants $g_j$ and $z_j$ of the steps $j<K$ do not
depend on the torus. Let
$\mathrm{pr}_{m''}\colon\mathbb{R}^4\to\mathbb{T}_{m''}$ be the
covering map. A set
$X\in\mathcal{D}_{j+1}$ on $\mathbb{T}_{m''}$ is called \emph{non-wrapping} (for the pair $m,m'$) if
$X=\mathrm{pr}_{m''}(\widehat{X})$ for a set $\widehat{X}$ contained in a closed cube of $\mathbb{R}^4$
of side at most one quarter of the side $2L^{\min(m,m')}$ of the smaller torus; $\mathrm{pr}_{m''}$ is
injective on such a cube.
Fix a site $z$ of level $j+1$ on $\mathbb{T}_m$ and a lift $\hat{z}\in\mathbb{R}^4$ of it; the site
$\mathrm{pr}_{m'}(\hat{z})$ of $\mathbb{T}_{m'}$ is also written $z$. A non-wrapping $X\ni z$, on
either torus, has exactly one lift
$\widehat{X}\ni\hat{z}$ of this kind: two such lifts lie in the closed cube of side
$L^{\min(m,m')}$ centred at $\hat{z}$, on which $\mathrm{pr}_{m''}$ is injective for either torus,
and they have the same image $X$, so they coincide. It is assumed --- this is the content of parts
(i) and (iii)(b) of the torus-independence lemma together with the lift lemma --- that
$\widehat{X}\mapsto\mathrm{pr}_m(\widehat{X})$ and
$\widehat{X}\mapsto\mathrm{pr}_{m'}(\widehat{X})$ set up a bijection between the non-wrapping
$X\in\mathcal{D}_{j+1}$ containing $z$ on $\mathbb{T}_m$ and those on $\mathbb{T}_{m'}$, under
which corresponding sets carry the same numbers $\mathbf{E}^{(j+1)}(X,1,z)$ and
$\mathbf{R}^{(j+1)}(X,1)$; and that every $X\in\mathcal{D}_{j+1}$ on $\mathbb{T}_{m''}$ with $X\ni z$
and
\begin{equation*}
d_{j+1}(X)<\frac{2L^{\min(m,m')}L^{K-j-1}}{4M}-2
\end{equation*}
is non-wrapping.
\end{itemize}
Then the vacuum-energy densities of
$\mathbb{T}_m$ and $\mathbb{T}_{m'}$ satisfy
\begin{equation}\label{8.22:eq-vacuum-energy-density-1}
\begin{split}
\delta_{\mathrm{vac}}
&:=\Bigl|\frac{E_m}{(2L^m)^4}-\frac{E_{m'}}{(2L^{m'})^4}\Bigr|\\
&\le\bar{\delta}_{\mathrm{vac}}
:=2(E_0+1)K_*\sum_{u\ge0}L^{4u}\exp\Bigl[-\frac{\kappa}{2}\Bigl(\frac{1}{2}L^{11}L^u-2\Bigr)\Bigr]
<\infty ,
\end{split}
\end{equation}
and $\bar{\delta}_{\mathrm{vac}}$ depends only on $L$ and $\mathfrak{p}$; it does not depend on $K$,
$g_0$, $g$, $m$ or $m'$.
\end{lemma}

\begin{proof}
For a level $j$ and a torus $\mathbb{T}_{m''}$, $m''\in\{m,m'\}$, let $\nu_{j+1}(\mathbb{T}_{m''})$ be
the number of sites of level $j+1$ on $\mathbb{T}_{m''}$.
Fix $j$ with $0\le j<K$. Level $j+1$ has spacing $L^{j+1-K}$ and $\mathbb{T}_{m''}$ has side
$2L^{m''}$, so
\begin{equation}\label{8.22:eq-vacuum-energy-density-2}
\nu_{j+1}(\mathbb{T}_{m''})=(2L^{m''})^4L^{4(K-j-1)} ;
\end{equation}
that is, per unit of volume there are $L^{4(K-j-1)}$ sites of level $j+1$, on every torus. By the
formula for the volume function $m_{\mathbb{T}}$ in Definition~\ref{1.2:def-flow-constants}, for
every coupling $g'$ one has
\begin{equation*}
L^{m_{\mathbb{T}}(g')}=L^{11}M\,R(g')\ge L^{11}M ,
\end{equation*}
since $R\ge1$ (Definition~\ref{1.2:def-flow-constants}); with $g'=g_-(g)$ and
$\min(m,m')\ge m_{\mathbb{T}}(g_-(g))$,
\begin{equation}\label{8.22:eq-vacuum-energy-density-3}
L^{\min(m,m')}\ge L^{m_{\mathbb{T}}(g_-(g))}\ge L^{11}M .
\end{equation}
By
hypothesis~(a), the difference of densities in \eqref{8.22:eq-vacuum-energy-density-1} is the sum
over $j<K$ of the differences of the densities of $E^{(j)}_{\mathrm{vac},1}$,
$E^{(j)}_{\mathrm{vac},2}$ and $E^{(j)}_{\mathrm{vac},3}$. We treat the three parts in turn.

\emph{The first part.} On a torus partitioned into blocks every lattice count at level $j$ or $j+1$
is an integer independent of $m''$ times $\nu_{j+1}(\mathbb{T}_{m''})$ (for instance, the number of
level-$j$ sites is $L^4\nu_{j+1}(\mathbb{T}_{m''})$). The numbers multiplying these counts, the
normalisation numbers of step $j$ \cite[(1.7)]{B-CMP109}, are built
from $\log g_j$, $\log z_j$ and constants of $\mathrm{SU}(2)$, and $g_j$, $z_j$ do not depend on the
torus by hypothesis~(c). Hence the quotient
$E^{(j)}_{\mathrm{vac},1}(\mathbb{T}_{m''})/\nu_{j+1}(\mathbb{T}_{m''})$
is the same for $m''=m$ and $m''=m'$: this part is exactly extensive. By
\eqref{8.22:eq-vacuum-energy-density-2} its density is $L^{4(K-j-1)}$ times this quotient on every
torus, and its contribution to $\delta_{\mathrm{vac}}$ is zero.

This is the one point at which the argument uses that no number of the first part is a
global, non-extensive normalisation number, such as a zero-mode factor or a residual-gauge factor per
torus. The averaging constraints, the gauge fixing and the normalisation integrals of Ba{\l}aban's
step are per block, so no such number occurs among the three parts of hypothesis~(a); the one
torus-wide number per level, the Gaussian normaliser of the step, is excluded by hypothesis~(a).
Were hypothesis~(a) to admit even a bounded
number of torus-wide numbers per level, they would contribute to $\delta_{\mathrm{vac}}$ a term of
order $K(2L^{\min(m,m')})^{-4}$, which is not uniform in $K$ at fixed $m'$.

\emph{The second part.} By \cite[(2.29)]{B-CMP119} at $U=1$ (hypothesis~(b)), the sum
$\mathbf{E}^{(j+1)}(1,z;\mathbb{T}_{m''})$ of hypothesis~(b) does not depend on $z$, so
\begin{equation*}
E^{(j)}_{\mathrm{vac},2}(\mathbb{T}_{m''})
=\nu_{j+1}(\mathbb{T}_{m''})\,\mathbf{E}^{(j+1)}(1,z;\mathbb{T}_{m''})
\end{equation*}
and, by \eqref{8.22:eq-vacuum-energy-density-2}, the density of $E^{(j)}_{\mathrm{vac},2}$ on
$\mathbb{T}_{m''}$ is $L^{4(K-j-1)}\mathbf{E}^{(j+1)}(1,z;\mathbb{T}_{m''})$. Put
\begin{equation}\label{8.22:eq-vacuum-energy-density-4}
D_j:=\frac{2L^{\min(m,m')}L^{K-j-1}}{4M}-2 ,
\end{equation}
one quarter of the side $2L^{\min(m,m')}$ of the smaller torus measured in units of $ML^{j+1-K}$,
less two.
Take the site $z$ on the two tori through a lift $\hat{z}$, as in hypothesis~(c). By the last
assertion of hypothesis~(c), every $X\ni z$ with $d_{j+1}(X)<D_j$ is non-wrapping on both tori; by
hypothesis~(c) the non-wrapping $X\ni z$ on $\mathbb{T}_m$ and on $\mathbb{T}_{m'}$ correspond
bijectively and corresponding terms are the same numbers, so they cancel in
$\mathbf{E}^{(j+1)}(1,z;\mathbb{T}_m)-\mathbf{E}^{(j+1)}(1,z;\mathbb{T}_{m'})$. What remains is
the sum over the wrapping $X\in\mathcal{D}_{j+1}$ on $\mathbb{T}_m$ with $X\ni z$, minus the
corresponding sum on $\mathbb{T}_{m'}$, and every such $X$ has $d_{j+1}(X)\ge D_j$. Each term of
either sum is bounded by the bound on
the $\mathbf{E}$-terms in hypothesis~(b). For $d_{j+1}(X)\ge D_j$ we have
$e^{-\kappa d_{j+1}(X)}\le e^{-\kappa D_j/2}e^{-\kappa d_{j+1}(X)/2}$, and the sum over $X\ni z$ of
the last factor is at most $K_*$ by \cite[(1.26)]{B-CMP116} read at the rate $\kappa/2$, which
applies at this rate by hypothesis~(b). Hence each
of the two remaining sums, one on each torus, is at most $E_0K_*e^{-\kappa D_j/2}$ in modulus, and
\begin{equation}\label{8.22:eq-vacuum-energy-density-5}
|\mathbf{E}^{(j+1)}(1,z;\mathbb{T}_m)-\mathbf{E}^{(j+1)}(1,z;\mathbb{T}_{m'})|
\le 2E_0K_*e^{-\kappa D_j/2} .
\end{equation}

\emph{The third part.} For $X\in\mathcal{D}_{j+1}$ let $|X|$ be the number of sites of level
$j+1$ in $X$, so that $|X|\ge1$. Since each $X$ is counted once at each of its $|X|$ sites,
\begin{equation*}
E^{(j)}_{\mathrm{vac},3}(\mathbb{T}_{m''})=\sum_z r_{m''}(z),
\qquad
r_{m''}(z):=\sum_{X\in\mathcal{D}_{j+1},\,X\ni z}|X|^{-1}\mathbf{R}^{(j+1)}(X,1),
\end{equation*}
with $z$ running over the sites of level $j+1$ on $\mathbb{T}_{m''}$. The partition of the
level-$(j+1)$ lattice into cubes of side $M$ divides both tori: $M$ is a
power of $L$, and $M\le L^{\min(m,m')}$ by \eqref{8.22:eq-vacuum-energy-density-3}, so $M$ divides
the side
$2L^{m''}L^{K-j-1}$, $m''\in\{m,m'\}$, in level-$(j+1)$ units. A translation by a vector $\eta$
preserving this partition maps $\mathcal{D}_{j+1}$ to itself, $X\mapsto X+\eta$, and preserves $|X|$,
and by the covariance
\cite[(2.32)--(2.33)]{B-CMP119} of the $\mathbf{R}$-terms at $U=1$ (hypothesis~(b)) it satisfies
$\mathbf{R}^{(j+1)}(X+\eta,1)=\mathbf{R}^{(j+1)}(X,1)$; hence $r_{m''}(z+\eta)=r_{m''}(z)$. So
$r_{m''}$ takes the same values at corresponding sites of all cubes of side $M$, and the density of
$E^{(j)}_{\mathrm{vac},3}$ on $\mathbb{T}_{m''}$ is $L^{4(K-j-1)}M^{-4}$ times
$\sum_{z\in\square}r_{m''}(z)$
for any one cube $\square$ of side $M$, with its $M^4$ sites of level $j+1$; we take $\square$ to be
the image on both tori of one cube of $\mathbb{R}^4$, so that its sites on the two tori correspond
through their lifts in that cube, as in hypothesis~(c). The non-wrapping terms again coincide on the
two tori by hypothesis~(c), and so
do the numbers $|X|$ of corresponding sets, $\mathrm{pr}_m$ and $\mathrm{pr}_{m'}$ being injective on
the lifts. The
bound \cite[(2.31)]{B-CMP119}, in the form displayed in hypothesis~(b), takes the place of the bound
on the $\mathbf{E}$-terms; since $|X|^{-1}\le1$, the sum over $X\ni z$ of
$|X|^{-1}e^{-\kappa d_{j+1}(X)/2}$ is still at most $K_*$, and the argument that gave
\eqref{8.22:eq-vacuum-energy-density-5}, with $E_0$ replaced by $1$, gives for every site $z$ of
$\square$
\begin{equation*}
|r_m(z)-r_{m'}(z)|\le 2K_*e^{-\kappa D_j/2} ,
\end{equation*}
and averaging over the $M^4$ sites of $\square$ keeps this bound. Together with
\eqref{8.22:eq-vacuum-energy-density-5}, the
contributions of the second and third parts of level $j$ to the difference of densities are at
most $2(E_0+1)K_*e^{-(\kappa/2)D_j}$ per site of level $j+1$.

\emph{Summation over levels.} Since there are $L^{4(K-j-1)}$ sites of level $j+1$ per unit of
volume, and the first part contributes nothing, the substitution $u=K-1-j$ gives
\begin{equation}\label{8.22:eq-vacuum-energy-density-6}
\delta_{\mathrm{vac}}
\le\sum_{u=0}^{K-1}L^{4u}\,2(E_0+1)K_*\exp\Bigl[-\frac{\kappa}{2}D_{K-1-u}\Bigr],
\end{equation}
with $D_{K-1-u}$ given by \eqref{8.22:eq-vacuum-energy-density-4} at $j=K-1-u$, where
$L^{K-j-1}=L^u$.
By \eqref{8.22:eq-vacuum-energy-density-3}, $D_{K-1-u}\ge \tfrac{1}{2}L^{11}L^u-2$. The exponential is
decreasing in $D_{K-1-u}$, so each term of
\eqref{8.22:eq-vacuum-energy-density-6} is at most the corresponding term of the series defining
$\bar{\delta}_{\mathrm{vac}}$, and extending the finite sum to all $u\ge0$ gives
$\delta_{\mathrm{vac}}\le\bar{\delta}_{\mathrm{vac}}$. Since $L^u$ grows geometrically in $u$, the
terms $L^{4u}\exp[-(\kappa/2)(\tfrac{1}{2}L^{11}L^u-2)]$ decay faster than every exponential in $u$,
and the
series converges. It contains only $L$, $\kappa$, $E_0$ and $K_*$, which depend only on $L$ and
$\mathfrak{p}$ by hypothesis~(b); in particular $\bar{\delta}_{\mathrm{vac}}$ does not depend on
$K$, $g_0$, $g$, $m$ or $m'$.
\end{proof}

\begin{remark}[Scope of the torus-independence of the vacuum-energy density]
Lemma~\ref{8.22:lem-vacuum-energy-density} concerns the zero-source constant $E_m$ only. It says
nothing about the sourced vacuum energies $\mathcal{E}_m(v)$: no covariance is asserted for
non-constant sources. It is a statement about the number $E_m$; it says nothing about the bare
measure on $\mathbb{T}_m$ or the effective densities $\rho_k$, about the contribution of any single
large-field history $h$, or about the parts of $E_m$ attributable to configurations with a
prescribed large-field history.
Hypothesis~(a) is a hypothesis on how $E_m$ is itemised: every number entering
$E^{(j)}_{\mathrm{vac}}$ is either a normalisation number of the first part or one of the localised
terms of the second and third parts. The determinant lemma, which compares the Gaussian normalisers
of the steps excluded by hypothesis~(a), contributes a further constant, depending on $L$ and
$\mathfrak{p}$ only (in particular not on $K$, $g_0$, $g$, $m$ or $m'$), to the zero-source
comparison of the two tori.
A bound on $|E_m|$ by a constant times the number of bare sites would not suffice: divided by the
volume it leaves a quantity of order $L^{4K}$. What the proof uses is the exact extensivity of
the first part and the cancellation of the non-wrapping terms of the second and third parts.
\end{remark}

\begin{theorem}[The volume-transfer bound for connected functions]\label{8.22:thm-volume-transfer}
Work in the setting of Notation~\ref{8.22:not-partition-sources-cells}. In particular $\varrho>0$
is fixed and $\mathfrak{r}\in\{1,2\}$ is fixed there:
$\mathfrak{r}=1$ when \textup{(T1)}--\textup{(T2)} are applied to the complex periodised weight
itself, and $\mathfrak{r}=2$ when they are applied to the real system of four functions formed
from it; $R>0$ is a radius, $\ell$ is the cell side of the per-ball volume-free bound for two
sources of radius $R$, $C_{\mathrm{face}}(R)$ is the face constant of that bound,
$\mathsf{B}=E_+^{\sharp}+C_{\mathrm{low}}$, and
\begin{equation*}
m_{\star}(g)=\max\{m_{\mathbb{T}}(g_-(g)),\,\log_L\ell\}.
\end{equation*}
Assume the following statements:
\begin{itemize}
\item[(a)] reflection positivity of the lattice measure;
\item[(b)] the chessboard lemma, the placement lemma, the periodised-weight lemma, the face
lemma and the Schwarz--Cauchy lemma that accompany the per-ball volume-free bound;
\item[(c)] parts (i) and (ii) of the per-ball volume-free bound, together with its identity
expressing the real part of its counterterm $\mathcal{E}''$ as the cell average of the real part
of the bulk part of the sourced vacuum-energy functional of $\mathbb{T}_{m'}$ at the periodised
weight, plus the real part of the face part of the counterterm;
\item[(d)] properties \textup{(T1)}--\textup{(T5)} of the correlation theorem on
$\mathbb{T}_{m'}$, and its lower stability bound at zero source on $\mathbb{T}_m$ and on
$\mathbb{T}_{m'}$;
\item[(e)] Lemma~\ref{8.22:lem-trace-subadditivity};
\item[(f)] Lemma~\ref{8.22:lem-vacuum-energy-density}.
\end{itemize}
Let $g\le\gamma_J$, $g_0\in\,]0,g_-(g)]$, $K=K(g_0,g)$ and $m\ge m'\ge m_{\star}(g)$; here $m'$
is a second torus exponent, not the exponent $m'$ of $M=L^{m'}$. Let $F_1$
and $F_2$ be real Lipschitz functions on $\mathbb{R}^4$ supported in one ball of radius $R$,
transported to $\mathbb{T}_m$ and to $\mathbb{T}_{m'}$.
No condition is imposed on the signs of $F_1,F_2$, on the number of components of their
supports, or on the overlap or distance of the supports. Then
\begin{equation}\label{8.22:eq-volume-transfer-b}
\begin{gathered}
\bigl|S^T_{2;K,m}(F_1,F_2)-S^T_{2;K,m'}(F_1,F_2)\bigr|
\le C_{\mathrm{VT}}\,\|F_1\|_{\sharp}\|F_2\|_{\sharp},\\
C_{\mathrm{VT}}:=2\log 2\cdot(2Q_{\star}+2)^2\Bigl(\frac{2\mathfrak{r}}{\varrho}\Bigr)^2,
\qquad
Q_{\star}:=\exp\bigl[(\mathsf{B}+\bar{\delta}_{\mathrm{vac}})\ell^4+4C_{\mathrm{face}}(R)\bigr].
\end{gathered}
\end{equation}
The constant $C_{\mathrm{VT}}$ depends only on $(L,\mathfrak{p},g,\varrho;R)$. It does not depend
on $K$, $g_0$, $m$ or $m'$, on the position of the ball, or on the supports of $F_1,F_2$.
\end{theorem}

\begin{proof}
\emph{Placement.} The ball, of diameter $2R$, does not wrap on either torus: $\ell$, the least
power of $L$ with $\ell\ge16(R+1)$, exceeds $2R$, and $L^{m}\ge L^{m'}\ge\ell$ since
$m'\ge\log_L\ell$. We apply the placement lemma with both balls equal to the given ball. It
provides a translation by a vector of $L^{-K}\mathbb{Z}^4$ after which the ball lies in the
interior of one cell $c$. Here a cell is a cube $\ell\,\mathbf{j}+[0,\ell]^4$ with
$\mathbf{j}\in\mathbb{Z}^4$.
After the translation the ball is at distance at least $R+2$ from every face plane
$\{x_\mu\in\ell\mathbb{Z}\}$.

This translation maps the bare sites $L^{-K}(\mathbb{Z}+\tfrac12)^4$ onto themselves. By
Lemma~\ref{1.1:lem-invariance} it therefore leaves both bare measures invariant, on
$\mathbb{T}_m$ and on $\mathbb{T}_{m'}$. Hence it changes neither $S^T_{2;K,m}(F_1,F_2)$ nor
$S^T_{2;K,m'}(F_1,F_2)$. It also leaves $\|F_1\|_{\sharp}$ and $\|F_2\|_{\sharp}$ unchanged,
because the sup norm and the Lipschitz constant are translation invariant. From now on the ball
has this position. Only one cell, $c$, carries sources, and it carries both of them.

Put
\begin{equation*}
\mathbb{D}_{\mathfrak{r}}:=\Bigl\{z\in\mathbb{C}^2:\ |z_1|\|F_1\|_{\sharp}+|z_2|\|F_2\|_{\sharp}
<\frac{\varrho}{\mathfrak{r}}\Bigr\},
\qquad w_z:=z_1F_1+z_2F_2 .
\end{equation*}
For a complex function $\xi$ write $\xi^{(+1)}:=\xi$ and $\xi^{(-1)}:=\bar{\xi}$. For
$\eta\in\mathfrak{G}_{\ell}$ let $\epsilon(\eta)=\pm1$ be its parity. On $\mathbb{R}^4$ put
\begin{equation*}
W_z:=\sum_{\eta\in\mathfrak{G}_{\ell}}\bigl(w_z\circ\eta^{-1}\bigr)^{(\epsilon(\eta))}.
\end{equation*}
The summands have pairwise disjoint supports, namely the images of the ball under
$\mathfrak{G}_{\ell}$. So at each point at most one term of the sum is non-zero.

We show that $W_z$ is mirror-periodic in the sense of
Lemma~\ref{8.22:lem-trace-subadditivity}. Condition (M1) holds by construction. Let $\theta$ be
one of the reflections $\theta_{\mu,a}$; it is an involution of parity $-1$. In the sum for
$W_z\circ\theta$ substitute $\eta=\theta\eta'$. Then $\eta^{-1}\circ\theta=\eta'^{-1}$ and
$\epsilon(\eta)=-\epsilon(\eta')$, so
\begin{equation*}
W_z\circ\theta=\sum_{\eta'\in\mathfrak{G}_{\ell}}\bigl(w_z\circ\eta'^{-1}\bigr)^{(-\epsilon(\eta'))}
=\overline{W_z}.
\end{equation*}
Condition (M2) holds because of the margin $R+2$. The reflections $\theta_{\mu,a}$ are isometries
and map the family of planes $\{x_\nu\in\ell\mathbb{Z}\}$ onto itself. Hence every image of the
ball lies at distance at least $R+2$ from all these planes. A plaquette containing a bond that
crosses such a plane has its barycentre within distance $2L^{-K}$ of the plane, and
$R+2>2L^{-K}$. So $W_z$ vanishes at the barycentre of every such plaquette.

By (M1), $W_z$ is $2\ell$-periodic in every direction. For $m''\ge\log_L\ell$ the cell count
$2L^{m''}/\ell$ of $\mathbb{T}_{m''}$ is an even integer, since $\ell$ is a power of $L$.
Therefore $W_z$ descends to $\mathbb{T}_{m''}$. There it is the periodisation of $w_z$ from the
cell $c$ considered in the periodised-weight lemma. That lemma gives
\begin{equation*}
\|W_z\|_{\sharp}\le|z_1|\|F_1\|_{\sharp}+|z_2|\|F_2\|_{\sharp}.
\end{equation*}
In particular, for $z\in\mathbb{D}_{\mathfrak{r}}$ we have
$\|W_z\|_{\sharp}<\varrho/\mathfrak{r}\le\varrho$, so $W_z\in\mathcal{W}$.

\emph{Chessboard bound on $\mathbb{T}_m$.} This step does not use the correlation theorem. By
definition of $\langle\cdot\rangle_m$,
\begin{equation*}
\frac{Z_m(w_z)}{Z_m(0)}=\langle\Phi_z\rangle_m,\qquad \Phi_z:=e^{\mathcal{O}(w_z)} .
\end{equation*}
Since $w_z$ is supported in the ball inside $c$, $\Phi_z$ belongs to the cell algebra
$\mathfrak{A}_c$. The chessboard lemma rests on reflection positivity of the lattice measure. It
gives $|\langle\Phi_z\rangle_m|\le\|\Phi_z\|_{\mathrm{ch}}$. The periodised-weight lemma
identifies the chessboard seminorm with a partition function of the mirror-periodic source.
Together these give
\begin{equation}\label{8.22:eq-volume-transfer-1}
\frac{|Z_m(w_z)|}{Z_m(0)}\le\|\Phi_z\|_{\mathrm{ch}}
=\Bigl[\frac{Z_m(W_z)}{Z_m(0)}\Bigr]^{1/|\mathcal{C}_m|},
\qquad |\mathcal{C}_m|=\Bigl(\frac{2L^m}{\ell}\Bigr)^4 .
\end{equation}
Here, for $m''\ge\log_L\ell$, $\mathcal{C}_{m''}$ denotes the set of cells of
$\mathbb{T}_{m''}$.

\emph{From $\mathbb{T}_m$ to $\mathbb{T}_{m'}$.} Put $k:=L^{m-m'}$, an integer $\ge1$, and
$\mathsf{n}:=(2L^{m'}/\ell)(1,1,1,1)$. This quadruple consists of even cell counts. It gives
$\mathbb{T}[\mathsf{n}]=\mathbb{T}_{m'}$ and $\mathbb{T}[k\mathsf{n}]=\mathbb{T}_m$, because
$k\cdot2L^{m'}=2L^m$. We apply Lemma~\ref{8.22:lem-trace-subadditivity} to the mirror-periodic
$W_z$ four times, once in each direction, passing through the intermediate rectangular tori:
put $\mathsf{n}^{(0)}:=\mathsf{n}$ and, for $j=1,\dots,4$, let $\mathsf{n}^{(j)}$ be obtained
from $\mathsf{n}^{(j-1)}$ by multiplying its entry in the $j$-th direction by $k$, so that
$\mathsf{n}^{(4)}=k\mathsf{n}$. The lemma gives, for $j=1,\dots,4$,
\begin{equation*}
Z_{\mathsf{n}^{(j)}}(W_z)\le Z_{\mathsf{n}^{(j-1)}}(W_z)^{k}.
\end{equation*}
Since $L$ is odd, $k=L^{m-m'}$ is odd, so $t\mapsto t^{k}$ is increasing on $\mathbb{R}$. The
four inequalities between real numbers therefore compose, and give
\begin{equation*}
Z_m(W_z)\le Z_{m'}(W_z)^{k^4}.
\end{equation*}
The numbers $Z_m(W_z)$ and $Z_{m'}(W_z)$ are real and non-negative: the identity in
\eqref{8.22:eq-volume-transfer-1}, read on $\mathbb{T}_{m''}$ for $m''\in\{m,m'\}$, expresses
$Z_{m''}(W_z)/Z_{m''}(0)$ as the $|\mathcal{C}_{m''}|$-th power of a chessboard seminorm, and
$Z_{m''}(0)>0$ because its integrand is positive.
Since $|\mathcal{C}_m|=k^4|\mathcal{C}_{m'}|$, taking $|\mathcal{C}_m|$-th roots yields
$Z_m(W_z)^{1/|\mathcal{C}_m|}\le Z_{m'}(W_z)^{1/|\mathcal{C}_{m'}|}$. Inserted into
\eqref{8.22:eq-volume-transfer-1}, this gives
\begin{equation}\label{8.22:eq-volume-transfer-2}
\frac{|Z_m(w_z)|}{Z_m(0)}\le Z_{m'}(W_z)^{1/|\mathcal{C}_{m'}|}\,Z_m(0)^{-1/|\mathcal{C}_m|}.
\end{equation}

\emph{Upper bound on $\mathbb{T}_{m'}$, lower bound at zero source on $\mathbb{T}_m$.} Both
$m'$ and $m$ are at least $m_{\mathbb{T}}(g_-(g))$, so (T1)--(T5) of the correlation theorem
are available on $\mathbb{T}_{m'}$. Fix $z\in\mathbb{D}_{\mathfrak{r}}$; then
$W_z\in\mathcal{W}$ by the placement step. Recall that $Z_{m'}(W_z)\ge0$. The normalisation
$\widetilde{Z}_{m'}=e^{-E_{m'}-\mathcal{E}_{m'}}Z_{m'}$ of (T1) and the upper stability bound
$|\widetilde{Z}_{m'}(W_z)|\le e^{E_+^{\sharp}(2L^{m'})^4}$, one of the properties
\textup{(T1)}--\textup{(T5)} stated in Notation~\ref{8.22:not-partition-sources-cells}, give
\begin{equation*}
\begin{split}
Z_{m'}(W_z)=|Z_{m'}(W_z)|
&=\exp\bigl[E_{m'}+\operatorname{Re}\mathcal{E}_{m'}(W_z)\bigr]\,|\widetilde{Z}_{m'}(W_z)|\\
&\le\exp\bigl[E_{m'}+\operatorname{Re}\mathcal{E}_{m'}(W_z)+E_+^{\sharp}(2L^{m'})^4\bigr].
\end{split}
\end{equation*}
For $\mathfrak{r}=2$ the upper stability bound is applied, as fixed in
Notation~\ref{8.22:not-partition-sources-cells}, to the real system of four functions formed
from $W_z$ at $(z,\bar z)$, on the polydisc of radius $\varrho/2$.
At zero source on $\mathbb{T}_m$ the normalisation reads $Z_m(0)=e^{E_m}\widetilde{Z}_m(0)$. The
lower stability bound at zero source $\widetilde{Z}_m(0)\ge e^{-C_{\mathrm{low}}(2L^m)^4}$,
stated in Notation~\ref{8.22:not-partition-sources-cells}, then gives
\begin{equation*}
Z_m(0)\ge\exp\bigl[E_m-C_{\mathrm{low}}(2L^m)^4\bigr].
\end{equation*}
Insert both bounds into \eqref{8.22:eq-volume-transfer-2} and use
$(2L^{m''})^4/|\mathcal{C}_{m''}|=\ell^4$ for $m''\in\{m,m'\}$. The exponent on the right becomes
\begin{equation*}
\begin{split}
&|\mathcal{C}_{m'}|^{-1}\bigl[E_{m'}+\operatorname{Re}\mathcal{E}_{m'}(W_z)
+E_+^{\sharp}(2L^{m'})^4\bigr]
-|\mathcal{C}_m|^{-1}\bigl[E_m-C_{\mathrm{low}}(2L^m)^4\bigr]\\
&\quad=|\mathcal{C}_{m'}|^{-1}\operatorname{Re}\mathcal{E}_{m'}(W_z)
+(E_+^{\sharp}+C_{\mathrm{low}})\ell^4
+\ell^4\Bigl(\frac{E_{m'}}{(2L^{m'})^4}-\frac{E_m}{(2L^m)^4}\Bigr).
\end{split}
\end{equation*}
Let $\delta_{\mathrm{vac}}=|E_m/(2L^m)^4-E_{m'}/(2L^{m'})^4|$ be the vacuum-density difference
of Lemma~\ref{8.22:lem-vacuum-energy-density}. The
last bracket is at most $\delta_{\mathrm{vac}}$, and $E_+^{\sharp}+C_{\mathrm{low}}=\mathsf{B}$.
Therefore, for $z\in\mathbb{D}_{\mathfrak{r}}$,
\begin{equation}\label{8.22:eq-volume-transfer-a}
\frac{|Z_m(w_z)|}{Z_m(0)}\le\exp\Bigl\{|\mathcal{C}_{m'}|^{-1}\operatorname{Re}
\mathcal{E}_{m'}(W_z)+(\mathsf{B}+\delta_{\mathrm{vac}})\ell^4\Bigr\}.
\end{equation}

\emph{One counterterm for both tori.} Let $\mathcal{E}''$ be the counterterm of the per-ball
volume-free bound, formed on $\mathbb{T}_{m'}$ for the system $(F_1,F_2)$:
\begin{equation*}
\mathcal{E}''(z):=|\mathcal{C}_{m'}|^{-1}\sum_{\eta\in\mathfrak{G}_{\ell}(\mathbb{T}_{m'})}
\Bigl[\mathcal{E}^{\mathrm{bulk}}_{\eta c}\bigl((w_z\circ\eta^{-1})^{(\epsilon(\eta))}\bigr)
\Bigr]^{(\epsilon(\eta))}+\mathcal{E}^{\partial}(w_z).
\end{equation*}
The notation is as follows:
\begin{itemize}
\item $\mathfrak{G}_{\ell}(\mathbb{T}_{m'})$ is the group generated by the reflections
$\theta_{\mu,a}$ acting on $\mathbb{T}_{m'}$;
\item $\mathcal{E}^{\mathrm{bulk}}_{c'}$ is the bulk part attached to the cell $c'$, and
$\mathcal{E}^{\partial}$ is the face part, of that theorem;
\item all the terms $e_j$ involved are those of $\mathbb{T}_{m'}$.
\end{itemize}
By part (i) of the per-ball volume-free bound, $\mathcal{E}''$ is holomorphic on
$\mathbb{D}_{\mathfrak{r}}$ and $\mathcal{E}''(0)=0$.

By the identity assumed in (c), the following holds on $\mathbb{D}_{\mathfrak{r}}$:
\begin{equation*}
\operatorname{Re}\mathcal{E}''(z)=|\mathcal{C}_{m'}|^{-1}\operatorname{Re}
\mathcal{E}^{\mathrm{bulk}}_{m'}(W_z)+\operatorname{Re}\mathcal{E}^{\partial}(w_z).
\end{equation*}
Here $\mathcal{E}_{m'}=\mathcal{E}^{\mathrm{bulk}}_{m'}+\mathcal{E}^{\partial}_{m'}$ is the
splitting of the sourced vacuum-energy functional of $\mathbb{T}_{m'}$ into bulk and face parts
used in that theorem. The symbol $\mathcal{E}^{\partial}(w_z)$, without torus index, denotes the
face part of the counterterm $\mathcal{E}''$, formed from the single source $w_z$; it is a
different quantity from the face part $\mathcal{E}^{\partial}_{m'}(W_z)$ of
$\mathcal{E}_{m'}(W_z)$. The face lemma, applied with both balls equal to the one ball, gives
\begin{equation*}
|\mathcal{C}_{m'}|^{-1}|\mathcal{E}^{\partial}_{m'}(W_z)|\le C_{\mathrm{face}}(R),
\qquad |\mathcal{E}^{\partial}(w_z)|\le C_{\mathrm{face}}(R).
\end{equation*}
Subtracting the identity from
$|\mathcal{C}_{m'}|^{-1}\operatorname{Re}\mathcal{E}_{m'}(W_z)$, the bulk parts cancel. Using the
two face bounds on the remaining terms, we obtain
\begin{equation}\label{8.22:eq-volume-transfer-3}
|\mathcal{C}_{m'}|^{-1}\operatorname{Re}\mathcal{E}_{m'}(W_z)-\operatorname{Re}\mathcal{E}''(z)
\le2C_{\mathrm{face}}(R)\le4C_{\mathrm{face}}(R).
\end{equation}
Define
\begin{equation}\label{8.22:eq-volume-transfer-4}
\mathsf{h}_m(z):=e^{-\mathcal{E}''(z)}\,\frac{Z_m(w_z)}{Z_m(0)},
\qquad
\mathsf{h}_{m'}(z):=e^{-\mathcal{E}''(z)}\,\frac{Z_{m'}(w_z)}{Z_{m'}(0)} .
\end{equation}

\emph{Holomorphy and bounds for $\mathsf{h}_m$ and $\mathsf{h}_{m'}$.} At finite $K$ the bare
lattice is finite. Moreover $\mathcal{O}(w_z)=z_1\mathcal{O}(F_1)+z_2\mathcal{O}(F_2)$ is linear
in $z$ and bounded uniformly in $U$. Hence $z\mapsto Z_{m''}(w_z)$ is entire for
$m''\in\{m,m'\}$. Consequently $\mathsf{h}_m$ and $\mathsf{h}_{m'}$ are holomorphic on
$\mathbb{D}_{\mathfrak{r}}$, and both equal $1$ at $z=0$.

For $z\in\mathbb{D}_{\mathfrak{r}}$ we combine three inputs:
\begin{itemize}
\item \eqref{8.22:eq-volume-transfer-a};
\item \eqref{8.22:eq-volume-transfer-3};
\item $\delta_{\mathrm{vac}}\le\bar{\delta}_{\mathrm{vac}}$, which is
Lemma~\ref{8.22:lem-vacuum-energy-density} and applies since
$m,m'\ge m_{\mathbb{T}}(g_-(g))$.
\end{itemize}
They give
\begin{equation*}
\begin{split}
|\mathsf{h}_m(z)|&=e^{-\operatorname{Re}\mathcal{E}''(z)}\,\frac{|Z_m(w_z)|}{Z_m(0)}
\le\exp\Bigl\{4C_{\mathrm{face}}(R)+(\mathsf{B}+\delta_{\mathrm{vac}})\ell^4\Bigr\}\\
&\le\exp\bigl[(\mathsf{B}+\bar{\delta}_{\mathrm{vac}})\ell^4+4C_{\mathrm{face}}(R)\bigr]
=Q_{\star}.
\end{split}
\end{equation*}
For $\mathsf{h}_{m'}$, part (ii) of the per-ball volume-free bound gives
\begin{equation*}
|\mathsf{h}_{m'}(z)|\le\exp\bigl[\mathsf{B}\ell^4+4C_{\mathrm{face}}(R)\bigr]\le Q_{\star}.
\end{equation*}
The same bound also follows from the argument above with $m=m'$. In that case $k=1$, the
transfer between tori is the identity, and $\delta_{\mathrm{vac}}=0$.

\emph{The Schwarz--Cauchy estimate; the counterterm cancels.} Let $m''\in\{m,m'\}$. Since
$Z_{m''}(0)>0$ and $\mathsf{h}_{m''}(0)=1$, both $\log Z_{m''}(w_z)$ and the principal
logarithm $\operatorname{Log}\mathsf{h}_{m''}(z)$ are defined for $z$ near $0$. There, by
\eqref{8.22:eq-volume-transfer-4},
\begin{equation*}
\log Z_{m''}(w_z)=\log Z_{m''}(0)+\mathcal{E}''(z)+\operatorname{Log}\mathsf{h}_{m''}(z).
\end{equation*}
The connected two-point function is the mixed derivative
$\partial_{z_1}\partial_{z_2}$ of $\log Z_{m''}(w_z)$ at $z=0$. Hence
\begin{equation*}
S^T_{2;K,m''}(F_1,F_2)=\partial_1\partial_2\mathcal{E}''(0)
+\partial_1\partial_2\operatorname{Log}\mathsf{h}_{m''}(0),
\end{equation*}
and the first term is the same number for $m''=m$ and $m''=m'$.

Each $\mathsf{h}_{m''}$ is holomorphic on $\mathbb{D}_{\mathfrak{r}}$, equals $1$ at $0$, and is
bounded in modulus by $Q_{\star}$ there. The Schwarz--Cauchy lemma therefore gives
\begin{equation*}
|\partial_1\partial_2\operatorname{Log}\mathsf{h}_{m''}(0)|
\le\log2\cdot(2Q_{\star}+2)^2\Bigl(\frac{2\mathfrak{r}}{\varrho}\Bigr)^2
\|F_1\|_{\sharp}\|F_2\|_{\sharp}.
\end{equation*}
We subtract the two representations. The common term $\partial_1\partial_2\mathcal{E}''(0)$
cancels, and the triangle inequality applied to the two remaining terms gives
\eqref{8.22:eq-volume-transfer-b}.

\emph{Dependence of the constant.} The formula for $C_{\mathrm{VT}}$ involves only $\mathsf{B}$,
$\bar{\delta}_{\mathrm{vac}}$, $\ell$, $C_{\mathrm{face}}(R)$, $\mathfrak{r}$ and $\varrho$:
\begin{itemize}
\item $\mathsf{B}=E_+^{\sharp}+C_{\mathrm{low}}$ is free of $K$, $g_0$ and of the torus, hence of
$m$ and of $m'$ (Notation~\ref{8.22:not-partition-sources-cells});
\item $\bar{\delta}_{\mathrm{vac}}$ depends only on $(L,\mathfrak{p})$
(Lemma~\ref{8.22:lem-vacuum-energy-density});
\item $\ell$, the least power of $L$ with $\ell\ge16(R+1)$, depends only on $(L,R)$;
\item $C_{\mathrm{face}}(R)$, the face constant of the per-ball volume-free bound, is a series
over all scales majorising the finite sums that occur at each $K$, and is free of $K$, $g_0$, of
the torus (hence of $m$ and $m'$), of $\ell$ and of positions.
\end{itemize}
The position of the ball was fixed by a translation that changes neither side. The supports of
$F_1,F_2$ enter only through $\|F_1\|_{\sharp}$, $\|F_2\|_{\sharp}$ and $R$. The only quantity
that could depend on their distance, $\partial_1\partial_2\mathcal{E}''(0)$, cancels exactly.

Of the torus $\mathbb{T}_m$, the proof uses only reflection positivity of the bare Wilson measure
and the lower stability bound at zero source: $Z_m(W_z)$ enters only through the inequality
$Z_m(W_z)\le Z_{m'}(W_z)^{k^4}$ of Lemma~\ref{8.22:lem-trace-subadditivity}.
\end{proof}

\begin{lemma}[Averaging over signed translates]\label{8.22:lem-signed-translates}
Let $\varepsilon=L^{-K}$ be the bare lattice spacing of
Notation~\ref{8.22:not-partition-sources-cells}, and let $R>0$ and $C\ge 0$. Let $D$ be a real
bilinear form on pairs of real Lipschitz functions with compact support on $\mathbb{R}^4$ such that
\begin{itemize}
\item[(a)] $D(F_1(\cdot-\tau),F_2(\cdot-\tau))=D(F_1,F_2)$ for every $\tau\in\varepsilon\mathbb{Z}^4$;
\item[(b)] $|D(F_1,F_2)|\le C\,\|F_1\|_{\sharp}\|F_2\|_{\sharp}$ whenever $F_1$ and $F_2$ are both
supported in one ball of radius $R$.
\end{itemize}
Let $f_1,f_2$ be real Lipschitz functions supported in one ball $B(y,r)$, where $r$ is a radius
with $\varepsilon\le r\le R/8$. Then
\begin{equation}\label{8.22:eq-signed-translates-1}
|D(f_1,f_2)|\le C\,c_{\Gamma}^{-1}\Bigl(\frac{r}{R}\Bigr)^4\|f_1\|_{\sharp}\|f_2\|_{\sharp},
\end{equation}
where $c_{\Gamma}>0$ is a pure number; one may take $c_{\Gamma}=32^{-4}$.
\end{lemma}

Here $B(y,r)$ is the closed Euclidean ball of centre $y$ and radius $r$, $|\cdot|$ the Euclidean norm
and $|\cdot|_\infty$ the maximum norm on $\mathbb{R}^4$.
Recall $\|f\|_{\sharp}=\max\{\|f\|_\infty,\,40M\operatorname{Lip}f\}$.

\begin{proof}
\emph{The lattice of translates.} Put $r':=\varepsilon\lceil r/\varepsilon\rceil$. Since
$r\ge\varepsilon$, we have $\lceil r/\varepsilon\rceil\ge 1$, so $r'$ is a positive integer
multiple of $\varepsilon$, and $r\le r'<r+\varepsilon\le 2r$. Let
\begin{equation*}
\Gamma:=\{\tau\in 4r'\mathbb{Z}^4:\ |\tau|_\infty\le R/4\}.
\end{equation*}
As $4r'\in\varepsilon\mathbb{Z}$, we have $\Gamma\subset\varepsilon\mathbb{Z}^4$. In each coordinate the
admissible values are $4r'j$ with $j\in\mathbb{Z}$, $|j|\le R/(16r')$; there are
$2\lfloor a\rfloor+1$ of them, with $a:=R/(16r')\ge R/(32r)$. Since $2\lfloor a\rfloor+1\ge a$ for
every $a\ge0$ (for $a<1$ the left side is $1$, and for $a\ge1$ it exceeds $2a-1\ge a$),
\begin{equation*}
|\Gamma|=(2\lfloor a\rfloor+1)^4\ge\Bigl(\frac{R}{32r}\Bigr)^4=c_{\Gamma}\Bigl(\frac{R}{r}\Bigr)^4,
\qquad c_{\Gamma}:=32^{-4}.
\end{equation*}

\emph{Geometry of the copies.} Write $B_\tau:=B(y+\tau,r)$ for $\tau\in\Gamma$. For
$\tau\ne\tau'$ in $\Gamma$ we have $|\tau-\tau'|\ge|\tau-\tau'|_\infty\ge 4r'\ge 4r$, so the closed
balls $B_\tau$, $B_{\tau'}$ are at distance at least $2r>0$; in particular they are pairwise disjoint.
Moreover $|\tau|\le 2|\tau|_\infty\le R/2$ in $\mathbb{R}^4$, hence
$B_\tau\subset B(y,R/2+r)\subset B(y,R)$ because $r\le R/8$.

\emph{The signed arrays.} For $\sigma=(\sigma_\tau)_{\tau\in\Gamma}\in\{\pm1\}^{\Gamma}$ and
$i=1,2$ put
\begin{equation*}
F_i^{\sigma}:=\sum_{\tau\in\Gamma}\sigma_\tau\,f_i(\cdot-\tau).
\end{equation*}
This is a finite sum of real Lipschitz functions, supported in $\bigcup_\tau B_\tau\subset B(y,R)$.
The function $f_i(\cdot-\tau)$ vanishes outside $B_\tau$, and, being continuous, also on the sphere
$\partial B_\tau$. Hence on $B_\tau$ we have $F_i^{\sigma}=\sigma_\tau f_i(\cdot-\tau)$, and
$F_i^{\sigma}=0$ outside $\bigcup_\tau B_\tau$. It follows that $\|F_i^{\sigma}\|_\infty=\|f_i\|_\infty$.

We show $\operatorname{Lip}F_i^{\sigma}=\operatorname{Lip}f_i$. First, for $x\in B_\tau$, let $p$ be
a point of $\partial B_\tau$ nearest to $x$; since $f_i(p-\tau)=0$,
\begin{equation*}
|F_i^{\sigma}(x)|\le\operatorname{Lip}f_i\cdot\operatorname{dist}(x,\partial B_\tau).
\end{equation*}
Now let $x,x'\in\mathbb{R}^4$. If both lie in one $B_\tau$, then
\begin{equation*}
|F_i^{\sigma}(x)-F_i^{\sigma}(x')|=|f_i(x-\tau)-f_i(x'-\tau)|\le\operatorname{Lip}f_i\,|x-x'|.
\end{equation*}
If both lie outside every ball, the difference is $0$. If $x\in B_\tau$ and $x'$ lies in no ball, the
segment from $x$ to $x'$ meets $\partial B_\tau$, so $\operatorname{dist}(x,\partial B_\tau)\le|x-x'|$,
and
\begin{equation*}
|F_i^{\sigma}(x)-F_i^{\sigma}(x')|=|F_i^{\sigma}(x)|\le\operatorname{Lip}f_i\,|x-x'|.
\end{equation*}
Finally let
$x\in B_\tau$, $x'\in B_{\tau'}$ with $\tau\ne\tau'$. By convexity the segment $[x,x']$ meets $B_\tau$
in an initial subsegment $[x,p]$ with $p\in\partial B_\tau$, and meets $B_{\tau'}$ in a final
subsegment $[q,x']$ with $q\in\partial B_{\tau'}$; these are disjoint because the balls are, so
$|x-x'|\ge|x-p|+|q-x'|$. Therefore
\begin{align*}
|F_i^{\sigma}(x)-F_i^{\sigma}(x')|
&\le|F_i^{\sigma}(x)|+|F_i^{\sigma}(x')|\\
&\le\operatorname{Lip}f_i\,\bigl[\operatorname{dist}(x,\partial B_\tau)
+\operatorname{dist}(x',\partial B_{\tau'})\bigr]\le\operatorname{Lip}f_i\,|x-x'|.
\end{align*}
Thus $\operatorname{Lip}F_i^{\sigma}\le\operatorname{Lip}f_i$. Conversely, $\operatorname{Lip}f_i$ is
the supremum of the difference quotients of $f_i$ over pairs of points of $B(y,r)$: if $x\in
B(y,r)$ and $x'\notin B(y,r)$, then $|f_i(x)-f_i(x')|=|f_i(x)-f_i(p)|$ for a point $p$ of the
segment on the sphere, with $|x-p|\le|x-x'|$. Translating such pairs by $\tau$ into $B_\tau$, where
$F_i^{\sigma}=\pm f_i(\cdot-\tau)$, gives $\operatorname{Lip}F_i^{\sigma}\ge\operatorname{Lip}f_i$.
Together,
\begin{equation}\label{8.22:eq-signed-translates-2}
\|F_i^{\sigma}\|_{\sharp}=\|f_i\|_{\sharp}\qquad(i=1,2,\ \sigma\in\{\pm1\}^{\Gamma}).
\end{equation}

\emph{Averaging over signs.} By bilinearity,
\begin{equation*}
D(F_1^{\sigma},F_2^{\sigma})=\sum_{\tau,\tau'\in\Gamma}\sigma_\tau\sigma_{\tau'}\,
D\bigl(f_1(\cdot-\tau),f_2(\cdot-\tau')\bigr).
\end{equation*}
Let $\mathsf{E}_{\sigma}$ denote the average $2^{-|\Gamma|}\sum_{\sigma\in\{\pm1\}^{\Gamma}}$, a
finite average over independent fair signs. Then $\mathsf{E}_{\sigma}\sigma_\tau\sigma_{\tau'}=1$ if
$\tau=\tau'$ and $0$ otherwise. Hence, using (a) for each $\tau\in\Gamma\subset\varepsilon\mathbb{Z}^4$,
\begin{equation*}
\mathsf{E}_{\sigma}D(F_1^{\sigma},F_2^{\sigma})
=\sum_{\tau\in\Gamma}D\bigl(f_1(\cdot-\tau),f_2(\cdot-\tau)\bigr)=|\Gamma|\,D(f_1,f_2).
\end{equation*}
For every $\sigma$, both $F_1^{\sigma}$ and $F_2^{\sigma}$ are supported in the ball $B(y,R)$ of
radius $R$, so (b) and \eqref{8.22:eq-signed-translates-2} give
$|D(F_1^{\sigma},F_2^{\sigma})|\le C\|f_1\|_{\sharp}\|f_2\|_{\sharp}$; the average obeys the same bound.
Consequently
\begin{equation*}
|D(f_1,f_2)|\le\frac{C}{|\Gamma|}\|f_1\|_{\sharp}\|f_2\|_{\sharp}
\le C\,c_{\Gamma}^{-1}\Bigl(\frac{r}{R}\Bigr)^4\|f_1\|_{\sharp}\|f_2\|_{\sharp},
\end{equation*}
which is \eqref{8.22:eq-signed-translates-1}.
\end{proof}

\begin{remark}[Several arguments; count of powers]
Let $n\ge1$ and let $D$ be a real $n$-linear form on real Lipschitz functions with compact support,
invariant under simultaneous translation of all arguments by $\tau\in\varepsilon\mathbb{Z}^4$, with
$|D(F_1,\dots,F_n)|\le C\prod_i\|F_i\|_{\sharp}$ whenever all $F_i$ are supported in one ball of radius
$R$. For $f_1,\dots,f_n$ supported in one ball $B(y,r)$, $\varepsilon\le r\le R/8$, take
$n-1$ independent sign families $\sigma^{(1)},\dots,\sigma^{(n-1)}\in\{\pm1\}^{\Gamma}$, put
$F_i^{\sigma}:=\sum_{\tau}\sigma^{(i)}_\tau f_i(\cdot-\tau)$ for $i<n$ and
$F_n^{\sigma}:=\sum_\tau\bigl(\prod_{i<n}\sigma^{(i)}_\tau\bigr)f_n(\cdot-\tau)$. Each coefficient is
a sign, so \eqref{8.22:eq-signed-translates-2} holds for every $F_i^{\sigma}$; expanding, the average of
$\prod_{i<n}\sigma^{(i)}_{\tau_i}\sigma^{(i)}_{\tau_n}$ vanishes unless $\tau_1=\dots=\tau_n$, so
$\mathsf{E}_{\sigma}D(F_1^{\sigma},\dots,F_n^{\sigma})=|\Gamma|\,D(f_1,\dots,f_n)$, and the same bound
$|D(f_1,\dots,f_n)|\le C\,c_{\Gamma}^{-1}(r/R)^4\prod_i\|f_i\|_{\sharp}$ follows.

In four dimensions, let $\vartheta\in\,]0,1]$, $0<\mathsf{d}\le1/40$ with $\varepsilon\le\mathsf{d}$,
and let $f_1,f_2$ be real Lipschitz functions supported in one ball of radius $5\mathsf{d}$ with
$\operatorname{Lip}f_i\le\|f_i\|_\infty/(\vartheta\mathsf{d})$, the Lipschitz condition of the class
$\mathfrak{T}(\mathsf{d};\vartheta)$. Take $R=1$ and $r=5\mathsf{d}$, so that $\varepsilon\le r\le R/8$.
Since $40M/(\vartheta\mathsf{d})\ge1$, each factor obeys
$\|f_i\|_{\sharp}\le 40M\|f_i\|_\infty/(\vartheta\mathsf{d})$, which contributes $\mathsf{d}^{-2}$ for
the two factors together, while $c_{\Gamma}^{-1}(r/R)^4=160^4\mathsf{d}^4$. The bound of the lemma is
therefore
\begin{equation*}
|D(f_1,f_2)|\le C\cdot160^4\Bigl(\frac{40M}{\vartheta}\Bigr)^2\mathsf{d}^{2}\,
\|f_1\|_\infty\|f_2\|_\infty .
\end{equation*}
For $n\ge2$ sources of this kind the net power is $\mathsf{d}^{4-n}$, a gain for $n=2,3$ only.
\end{remark}

\begin{corollary}[Short-distance volume independence of the two-point function]
\label{8.22:cor-short-distance-volume}
Assume the hypotheses of Theorem~\ref{8.22:thm-volume-transfer} with $R:=1$:
\[
g\le\gamma_J,\qquad g_0\in\,]0,g_-(g)],\qquad K=K(g_0,g),\qquad m\ge m'\ge m_{\star}(g).
\]
Let $\varepsilon:=L^{-K}$ be the spacing of the bare lattice of Definition~\ref{1.1:def-lattice}.
Let $\vartheta\in\,]0,1]$. Let $\mathsf{d}\le 1/40$ with $\varepsilon\le\mathsf{d}$. Let $f_1,f_2$ be
real Lipschitz functions on
$\mathbb{R}^4$, of arbitrary sign, supported in one ball of radius $5\mathsf{d}$, and such that
\[
\operatorname{Lip} f_i\le\frac{\|f_i\|_\infty}{\vartheta\mathsf{d}},\qquad i=1,2.
\]
Then
\begin{equation}\label{8.22:eq-short-distance-volume-a}
\begin{split}
\bigl|S^T_{2;K,m}(f_1,f_2)-S^T_{2;K,m'}(f_1,f_2)\bigr|
&\le C_T(g)\,\mathsf{d}^2\,\|f_1\|_\infty\|f_2\|_\infty,\\
C_T(g)&:=5^4\,c_{\Gamma}^{-1}\Bigl(\frac{40M}{\vartheta}\Bigr)^2C_{\mathrm{VT}},
\end{split}
\end{equation}
uniformly in $K$, $g_0$, the position of the ball, and all $m\ge m'\ge m_{\star}(g)$. Here
$C_{\mathrm{VT}}$ is the constant of \eqref{8.22:eq-volume-transfer-b} at $R=1$, and $c_{\Gamma}>0$
is the pure number of Lemma~\ref{8.22:lem-signed-translates}. For three sources: let
$f_1,f_2,f_3$ be real Lipschitz functions on $\mathbb{R}^4$, of arbitrary sign, all supported in
one ball of radius $5\mathsf{d}$, with $\operatorname{Lip}f_i\le\|f_i\|_\infty/(\vartheta\mathsf{d})$
for $i=1,2,3$. Then
\begin{equation}\label{8.22:eq-short-distance-volume-1}
\bigl|S^T_{3;K,m}(f_1,f_2,f_3)-S^T_{3;K,m'}(f_1,f_2,f_3)\bigr|
\le C_{T,3}(g)\,\mathsf{d}\prod_{i=1}^{3}\|f_i\|_\infty ,
\end{equation}
where $C_{T,3}(g):=5^4\,c_{\Gamma}^{-1}(40M/\vartheta)^3C_{\mathrm{VT},3}$, and $C_{\mathrm{VT},3}$
is the constant of the three-source case of Theorem~\ref{8.22:thm-volume-transfer} at $R=1$.
\end{corollary}

\begin{proof}
For real Lipschitz functions $F_1,F_2$ on $\mathbb{R}^4$ with compact support put
\[
D(F_1,F_2):=S^T_{2;K,m}(F_1,F_2)-S^T_{2;K,m'}(F_1,F_2).
\]
Here a compactly supported function $F$ on $\mathbb{R}^4$ is read on $\mathbb{T}_{m''}$,
$m''\in\{m,m'\}$, through its periodisation $\sum_{k\in 2L^{m''}\mathbb{Z}^4}F(\cdot-k)$. This is
linear in $F$ and commutes with translations; the lattice $2L^{m''}\mathbb{Z}^4$ of periods is
contained in $\varepsilon\mathbb{Z}^4$, since $2L^{m''}=2L^{m''+K}\varepsilon$. If $\operatorname{supp}F$
lies in a ball of radius less than $L^{m''}$, the translates of that ball by non-zero periods are
disjoint from it, and the periodisation coincides with the transport of $F$ to $\mathbb{T}_{m''}$
used in Theorem~\ref{8.22:thm-volume-transfer}. Since $L^{m''}\ge L^{m'}>1$, this is the case for
every function supported in a ball of radius $1$, in particular for $f_1,f_2$ and for every
function to which the bound of Theorem~\ref{8.22:thm-volume-transfer} is applied below.
We verify the hypotheses of Lemma~\ref{8.22:lem-signed-translates} (averaging over signed
translates) for $D$, with $R=1$ and $C=C_{\mathrm{VT}}$.

\emph{Bilinearity.} The observable $\mathcal{O}(F)$ depends linearly on $F$, and on each torus the
connected two-point function $S^T_{2;K,m''}(F_1,F_2)$ is the second cumulant of the pair
$\mathcal{O}(F_1),\mathcal{O}(F_2)$ under the bare measure, that is, their covariance. Hence it
is real and bilinear in $(F_1,F_2)$ for $m''=m$ and for $m''=m'$, and so is the difference $D$.

\emph{Hypothesis \textup{(a)}: invariance under $\varepsilon\mathbb{Z}^4$.} Let
$\tau\in\varepsilon\mathbb{Z}^4$. The translation by $\tau$ maps the bare lattice of
$\mathbb{T}_m$ onto itself, and a plaquette $p$ onto a plaquette $p+\tau$ with barycentre
$x(p)+\tau$. Therefore $\mathcal{O}(F(\cdot-\tau))$ is the composition of $\mathcal{O}(F)$ with
the translation by $\tau$ of the bond variables. By Lemma~\ref{1.1:lem-invariance} the bare
Wilson measure on $\mathbb{T}_m$ is invariant under this lattice isometry, so all moments of
$\mathcal{O}(F_1(\cdot-\tau)),\mathcal{O}(F_2(\cdot-\tau))$ equal those of
$\mathcal{O}(F_1),\mathcal{O}(F_2)$. In particular
$S^T_{2;K,m}(F_1(\cdot-\tau),F_2(\cdot-\tau))=S^T_{2;K,m}(F_1,F_2)$. The same holds on
$\mathbb{T}_{m'}$, and hence (a) holds for $D$.

\emph{Hypothesis \textup{(b)}.} If $F_1,F_2$ are real Lipschitz and supported in one ball of
radius $R=1$, then \eqref{8.22:eq-volume-transfer-b} gives
$|D(F_1,F_2)|\le C_{\mathrm{VT}}\|F_1\|_{\sharp}\|F_2\|_{\sharp}$. This is the bound
\eqref{8.22:eq-volume-transfer-b} of Theorem~\ref{8.22:thm-volume-transfer} at $R=1$; no
condition on the distance between the supports enters.

\emph{Application.} Take $r_0:=5\mathsf{d}$. Since $\varepsilon\le\mathsf{d}$, we have
$\varepsilon\le r_0$; and $r_0=5\mathsf{d}\le 1/8$, that is, $r_0\le R/8$.
Lemma~\ref{8.22:lem-signed-translates} therefore yields
\[
|D(f_1,f_2)|\le C_{\mathrm{VT}}\,c_{\Gamma}^{-1}(5\mathsf{d})^4\,
\|f_1\|_{\sharp}\|f_2\|_{\sharp}.
\]
By definition, $\|f_i\|_{\sharp}=\max\{\|f_i\|_\infty,40M\operatorname{Lip}f_i\}$. We have
$40M\operatorname{Lip}f_i\le 40M\|f_i\|_\infty/(\vartheta\mathsf{d})$ by hypothesis. Since
$M\ge1$, $\vartheta\le1$ and $\mathsf{d}\le1/40$, also $40M/(\vartheta\mathsf{d})\ge1$. Hence
$\|f_i\|_{\sharp}\le 40M\|f_i\|_\infty/(\vartheta\mathsf{d})$, and
\[
\begin{aligned}
|D(f_1,f_2)|&\le C_{\mathrm{VT}}\,c_{\Gamma}^{-1}\,5^4\mathsf{d}^4
\Bigl(\frac{40M}{\vartheta\mathsf{d}}\Bigr)^2\|f_1\|_\infty\|f_2\|_\infty\\
&=C_T(g)\,\mathsf{d}^2\|f_1\|_\infty\|f_2\|_\infty ,
\end{aligned}
\]
which is \eqref{8.22:eq-short-distance-volume-a}.

\emph{Three sources.} For real Lipschitz $F_1,F_2,F_3$ with compact support put
\[
D_3(F_1,F_2,F_3):=S^T_{3;K,m}(F_1,F_2,F_3)-S^T_{3;K,m'}(F_1,F_2,F_3).
\]
On each torus $S^T_{3;K,m''}$ is the third cumulant of $\mathcal{O}(F_1),\mathcal{O}(F_2),
\mathcal{O}(F_3)$; since $\mathcal{O}(F)$ is real and linear in $F$ and cumulants are
multilinear, $D_3$ is real and trilinear. The argument given above for hypothesis (a) shows that
$D_3$ is invariant under simultaneous translation of its three arguments by
$\tau\in\varepsilon\mathbb{Z}^4$. The three-source case of Theorem~\ref{8.22:thm-volume-transfer}
at $R=1$ gives $|D_3(F_1,F_2,F_3)|\le C_{\mathrm{VT},3}\prod_{i}\|F_i\|_{\sharp}$ whenever the
$F_i$ are real Lipschitz and supported in one ball of radius $1$.
Lemma~\ref{8.22:lem-signed-translates} is
stated for bilinear forms; for the trilinear form $D_3$ we carry out its averaging with two
independent families of signs.

Let $B(y,5\mathsf{d})$ be the ball containing the supports of $f_1,f_2,f_3$, and put again
$r_0:=5\mathsf{d}$, so that $\varepsilon\le r_0\le R/8$ with $R=1$. Let
$\Gamma\subset\varepsilon\mathbb{Z}^4$ be the finite set of translates constructed in the proof of
Lemma~\ref{8.22:lem-signed-translates} for the ball $B(y,r_0)$ and $R=1$. That construction uses
only $\varepsilon\le r_0\le R/8$ and has three properties: the balls $B(y+\tau,r_0)$,
$\tau\in\Gamma$, lie in $B(y,1)$; for every real Lipschitz $f$ supported in $B(y,r_0)$ and every
family of signs $(\eta_\tau)_{\tau\in\Gamma}\in\{\pm1\}^{\Gamma}$, the function
$\sum_{\tau\in\Gamma}\eta_\tau f(\cdot-\tau)$ has $\|\cdot\|_{\sharp}$-norm equal to
$\|f\|_{\sharp}$; and the number $|\Gamma|$ of translates satisfies
$|\Gamma|\ge c_{\Gamma}(R/r_0)^4$. For $\sigma,\rho\in\{\pm1\}^{\Gamma}$ put
\[
F_1^{\sigma}:=\sum_{\tau\in\Gamma}\sigma_\tau f_1(\cdot-\tau),\qquad
F_2^{\rho}:=\sum_{\tau\in\Gamma}\rho_\tau f_2(\cdot-\tau),\qquad
F_3^{\sigma\rho}:=\sum_{\tau\in\Gamma}\sigma_\tau\rho_\tau f_3(\cdot-\tau).
\]
Since $(\sigma_\tau\rho_\tau)_{\tau}$ is again a family of signs, these three functions are real
Lipschitz, supported in $B(y,1)$, and have $\|\cdot\|_{\sharp}$-norms $\|f_1\|_{\sharp}$,
$\|f_2\|_{\sharp}$, $\|f_3\|_{\sharp}$. Hence, for every $\sigma,\rho$,
\[
|D_3(F_1^{\sigma},F_2^{\rho},F_3^{\sigma\rho})|
\le C_{\mathrm{VT},3}\prod_{i=1}^{3}\|f_i\|_{\sharp}.
\]
By trilinearity,
\[
D_3(F_1^{\sigma},F_2^{\rho},F_3^{\sigma\rho})
=\sum_{\tau,\tau',\tau''\in\Gamma}\sigma_\tau\sigma_{\tau''}\rho_{\tau'}\rho_{\tau''}\,
D_3\bigl(f_1(\cdot-\tau),f_2(\cdot-\tau'),f_3(\cdot-\tau'')\bigr).
\]
For $\tau\ne\tau''$ the sum of $\sigma_\tau\sigma_{\tau''}$ over $\sigma\in\{\pm1\}^{\Gamma}$
factorises and contains the factor $\sum_{\sigma_\tau=\pm1}\sigma_\tau=0$, while for $\tau=\tau''$
the product equals $1$; hence $2^{-|\Gamma|}\sum_{\sigma}\sigma_\tau\sigma_{\tau''}$ equals $1$ if
$\tau=\tau''$ and $0$ otherwise, and likewise for $\rho$. Averaging the last display over all
pairs $(\sigma,\rho)$ therefore leaves only the terms with $\tau=\tau'=\tau''$, and by the
invariance of $D_3$ under $\varepsilon\mathbb{Z}^4$ each of these equals $D_3(f_1,f_2,f_3)$. The
average is thus $|\Gamma|\,D_3(f_1,f_2,f_3)$, and it is bounded in modulus by the same bound as
each of its terms. Hence
\begin{equation}\label{8.22:eq-short-distance-volume-2}
|D_3(f_1,f_2,f_3)|\le |\Gamma|^{-1}C_{\mathrm{VT},3}\prod_{i=1}^{3}\|f_i\|_{\sharp}
\le C_{\mathrm{VT},3}\,c_{\Gamma}^{-1}(5\mathsf{d})^4\prod_{i=1}^{3}\|f_i\|_{\sharp}.
\end{equation}
With the bound $\|f_i\|_{\sharp}\le 40M\|f_i\|_\infty/(\vartheta\mathsf{d})$
obtained above for each of the three functions, we get
\[
\begin{aligned}
|D_3(f_1,f_2,f_3)|&\le C_{\mathrm{VT},3}\,c_{\Gamma}^{-1}\,5^4\mathsf{d}^4
\Bigl(\frac{40M}{\vartheta\mathsf{d}}\Bigr)^3\prod_{i=1}^{3}\|f_i\|_\infty\\
&=C_{T,3}(g)\,\mathsf{d}\prod_{i=1}^{3}\|f_i\|_\infty ,
\end{aligned}
\]
which is \eqref{8.22:eq-short-distance-volume-1}.

\emph{Uniformity.} The constant $C_{\mathrm{VT}}$ at $R=1$ depends only on
$(L,\mathfrak{p},g,\varrho)$. In particular it does not depend on $K$, $g_0$, $m$, $m'$, the
position or the supports. The number $c_{\Gamma}$ is a pure number, and $M$, a fixed power of
$L$, is one of Ba{\l}aban's auxiliary parameters $\mathfrak{p}$
(Notation~\ref{2.1:not-prefix}), chosen after $L$ and independent of the two torus exponents
$m\ge m'$ compared here. Hence $C_T(g)$ depends only on $L$, $\mathfrak{p}$, $g$, $\varrho$
and $\vartheta$, and the bound holds uniformly as stated. The same holds for
$C_{\mathrm{VT},3}$, which in Theorem~\ref{8.22:thm-volume-transfer} at $R=1$ depends only on
$(L,\mathfrak{p},g,\varrho)$, hence for
$C_{T,3}(g)$, and \eqref{8.22:eq-short-distance-volume-1} is uniform in the same sense.
\end{proof}

\begin{definition}[Consistency with the operator product expansion]\label{8.22:def-expansion-consistency}
Let $\gamma_W^{\mathbb{T}}$ be the coupling ceiling and $m_V(g)$ the horizon exponent of the two-point
witness theorem on a single torus. Let $m_{\star}(g)=\max\{m_{\mathbb{T}}(g_-(g)),\log_L\ell\}$ be the
exponent of the reference torus in the volume-transfer theorem, where $\ell$, the least power of $L$
with $\ell\ge 32$, is the cell side of the chessboard for two sources in one ball of radius one. Let
$C_T(g)$ be the constant of the short-distance volume independence of
the two-point function (Corollary~\ref{8.22:cor-short-distance-volume}). Let $b_0$, $c_{\mathcal{K}}$,
$B^{\sharp}$ and the witness constant $C_W(s)$ be as in Notation~\ref{0.2:not-glossary}.
\begin{itemize}
\item[(a)] The \emph{coupling ceiling of the volume transfer} is
\begin{equation}\label{8.22:eq-expansion-consistency-a}
\gamma_{\star}^{V}:=\sup\bigl\{\eta\le\gamma_W^{\mathbb{T}}:\ m_{\star}(g)\le m_V(g)
\ \text{for all } 0<g\le\eta\bigr\},
\end{equation}
and we put $\gamma_W^{\mathrm{VT}}:=\gamma_{\star}^{V}$. Then $\gamma_{\star}^{V}>0$. The associated
condition on the coupling is $g\le\gamma_W^{\mathrm{VT}}$.
\item[(b)] For each $g$ in the range of Corollary~\ref{8.22:cor-short-distance-volume} (so that
$C_T(g)$ is defined), the \emph{depth floor of the volume transfer}
$s_T(g)=s_T(g;\vartheta,\varrho)$ is the least integer $s_0\ge0$ such that for every integer
$s\ge s_0$
\begin{equation}\label{8.22:eq-expansion-consistency-b}
\begin{gathered}
L\,C_W(s)\,L^{-s}\le\frac{1}{40},\\
C_T(g)\,L^{2}\,C_W(s)^{2}\,L^{-2s}\le\frac{1}{9}\,b_0\,c_{\mathcal{K}}
\bigl(g^{-2}+B^{\sharp}(s+1)\bigr)^{-2}.
\end{gathered}
\end{equation}
This integer exists; it is free of $K$, $g_0$, $m$, $\mathsf{d}$ and of the position of the test
functions, and it is polylogarithmic in $g^{-2}$: it is bounded by a constant times a power of
$1+\log_L g^{-2}$, the constant and the exponent depending on $L$, $\mathfrak{p}$, $\vartheta$ and
$\varrho$ only.
\end{itemize}
\end{definition}

\begin{proof}
(a) By its definition in the two-point witness theorem on a single torus, the horizon exponent
satisfies $m_V(g)\ge c\,g^{-2}-C_0$ with constants $c>0$ and $C_0\ge0$ independent of $g$; in
particular it increases to $\infty$ as $g\downarrow0$. For sources in one ball of radius one, as in
Corollary~\ref{8.22:cor-short-distance-volume}, the cell side $\ell$ fixed above satisfies
\begin{align*}
\ell&\le 32L\le L^{11}M\le L^{m_{\mathbb{T}}(g_-(g))},\\
m_{\star}(g)&=\max\bigl\{m_{\mathbb{T}}(g_-(g)),\log_L\ell\bigr\}=m_{\mathbb{T}}(g_-(g))\\
&=11+\log_L M+\log_L R(g_-(g)),
\end{align*}
the middle inequality because $L\ge3$ and $M\ge1$, the last because $R(\cdot)\ge1$.
Let $r_{\mathrm{pr}}>0$ be the fixed exponent in the growth bound
$\log_L R(g')\le 1+r_{\mathrm{pr}}\log_L\log g'^{-2}$ of the function $R$ in the torus exponent
$m_{\mathbb{T}}$. Since $g_-(g)^{-2}=g^{-2}+B^{\sharp}$,
\begin{equation*}
m_{\star}(g)\le 12+\log_L M+r_{\mathrm{pr}}\log_L\log\bigl(g^{-2}+B^{\sharp}\bigr),
\end{equation*}
which grows like $\log_L$ of a power of $\log g^{-2}$. Hence there is $\eta>0$ with
$m_{\star}(g)\le m_V(g)$ for all $0<g\le\eta$. Replacing $\eta$ by $\min\{\eta,\gamma_W^{\mathbb{T}}\}$,
we obtain a positive element of the set in \eqref{8.22:eq-expansion-consistency-a}, so
$\gamma_{\star}^{V}>0$.

The condition defining that set is required for all $0<g\le\eta$. Hence, together with any $\eta$,
the set contains every smaller positive number, and the condition $g\le\gamma_W^{\mathrm{VT}}$ is
one-sided, a ceiling on $g$. In the order of choice of the constants it is fixed after
$\gamma_W^{\mathbb{T}}$, as
the last of the ceilings on the coupling; nothing chosen earlier depends on it.

(b) The second inequality of \eqref{8.22:eq-expansion-consistency-b} is equivalent, after
multiplication by $9\bigl(g^{-2}+B^{\sharp}(s+1)\bigr)^{2}$, to
\begin{equation}\label{8.22:eq-expansion-consistency-1}
9\,C_T(g)\,L^{2}\,C_W(s)^{2}\,L^{-2s}\bigl(g^{-2}+B^{\sharp}(s+1)\bigr)^{2}\le b_0\,c_{\mathcal{K}} .
\end{equation}
The constant $C_W(s)$ is polylogarithmic in $\bar{x}_s=g^{-2}+B^{\sharp}(s+1)$, which is linear in
$s$. The left sides of the first inequality of
\eqref{8.22:eq-expansion-consistency-b} and of \eqref{8.22:eq-expansion-consistency-1} carry the factors
$L^{-s}$ and $L^{-2s}$. Against them stand therefore a factor polylogarithmic in $s$ and, in
\eqref{8.22:eq-expansion-consistency-1}, the factor $\bar{x}_s^{\,2}$, quadratic in $s$. Hence both
inequalities hold for every sufficiently
large $s$, and a least $s_0$ with the required property exists. Since $C_W(0)\ge1$, the first
inequality fails at $s=0$, so $s_T(g)\ge1$.

By construction both inequalities hold at every depth of the half-line $\{s\ge s_T(g)\}$; no
monotonicity in $s$ of their left sides is needed or claimed. Of the quantities $K$, $g_0$, $m$,
$\mathsf{d}$, the position of the test functions and the coupling $g$, the two conditions depend only
on $g$: through $g^{-2}$ on the right of
\eqref{8.22:eq-expansion-consistency-b} and inside $\bar{x}_s$ in $C_W(s)$, and through $C_T(g)$,
which depends on $g$ through the stability constant $\mathsf{B}$ (and through the cell side whenever
that side depends on $g$; see below).
The constant of Corollary~\ref{8.22:cor-short-distance-volume} obeys
\begin{equation*}
\log C_T(g)\le 2\bigl(\mathsf{B}+|\bar{\delta}_{\mathrm{vac}}|\bigr)\ell^{4}+C' ,
\end{equation*}
where $\mathsf{B}$ and $\bar{\delta}_{\mathrm{vac}}$ are the stability constant and the
vacuum-density bound used there, $\ell$ is the cell side of the chessboard entering $C_T(g)$, and $C'$
depends on $L$, $\mathfrak{p}$, $\vartheta$ and $\varrho$
only. Indeed, $C_T(g)$ is $(2Q_{\star}+2)^{2}$ times a factor depending on $M$, $\vartheta$, $\varrho$
and the factor $\mathfrak{r}\in\{1,2\}$ dividing the radius of the polydisc of holomorphy only, where
\begin{equation*}
Q_{\star}=\exp\bigl[(\mathsf{B}+\bar{\delta}_{\mathrm{vac}})\ell^{4}+4C_{\mathrm{face}}(1)\bigr]
\end{equation*}
is the polydisc bound and $C_{\mathrm{face}}(1)$, the face constant at radius one, does not depend on
$g$. Since $2Q_{\star}+2\le4\max\{Q_{\star},1\}$ and $\mathsf{B}\ge0$,
\begin{equation*}
\log(2Q_{\star}+2)^{2}\le\log16+2\bigl(\mathsf{B}+|\bar{\delta}_{\mathrm{vac}}|\bigr)\ell^{4}
+8\,|C_{\mathrm{face}}(1)|,
\end{equation*}
and the $g$-free terms are absorbed in $C'$.
By the degree bound on the stability constant, $\mathsf{B}\le C_{\mathsf{B}}(1+\log g^{-2})$ with
the constant $C_{\mathsf{B}}$ of that bound, which does not depend on $g$, and
$\bar{\delta}_{\mathrm{vac}}$ depends on $L$ and $\mathfrak{p}$ only; hence
$\log_L C_T(g)\le C''\bigl(1+\ell^{4}\log_L g^{-2}\bigr)$ with $C''$ of the same dependence as $C'$.
Taking logarithms to base $L$ in \eqref{8.22:eq-expansion-consistency-1}, that inequality holds as soon
as
\begin{equation*}
\begin{split}
2s\ge{}&\log_L\frac{9\,C_T(g)\,L^{2}}{b_0\,c_{\mathcal{K}}}+2\log_L C_W(s)\\
&+2\log_L\bigl(g^{-2}+B^{\sharp}(s+1)\bigr);
\end{split}
\end{equation*}
the last two terms depend on $s$ only logarithmically, and the same holds for the first inequality of
\eqref{8.22:eq-expansion-consistency-b} after taking $\log_L$. With the bound on $\log_L C_T(g)$ this
gives
\begin{equation*}
s_T(g)\le C\bigl(1+\ell^{4}\log_L g^{-2}\bigr),
\end{equation*}
with $C$ depending on $L$, $\mathfrak{p}$, $\vartheta$ and $\varrho$ only.
If the cell side entering $C_T(g)$ is the $g$-free side $\ell$ fixed above and the reference units of
Theorem~\ref{3.1:thm-main} are taken to coincide with the units of the last lattice of the correlation
theorem, this bound is linear in $\log_L g^{-2}$. In general the conversion between these two units
multiplies lengths by $L^{b}$, with an integer $b=b(g)\ge0$ independent of $K$ and $m$ and with
$L^{b}$ at most a constant times $L^{11}M(\log g^{-2})^{r_{\mathrm{pr}}}$; the cell side entering
$C_T(g)$ is then at most $32L^{b+1}$, the bound for $\log_L C_T(g)$ holds with $\ell^{4}$ replaced by
$(32L^{b+1})^{4}$, and $s_T(g)$ is polylogarithmic in $g^{-2}$ of degree at most
$4r_{\mathrm{pr}}+1$. The same degree results if the chessboard is run with the cell side
$\ell(2,1,g)$ of Theorem~\ref{3.1:thm-main}. In either case $s_T(g)$ is of the form stated in (b).

The quantities entering \eqref{8.22:eq-expansion-consistency-b} are $L$, $C_W(s)$, $b_0$,
$c_{\mathcal{K}}$, $B^{\sharp}$, $g$ and $C_T(g)$. By Corollary~\ref{8.22:cor-short-distance-volume}
the last is independent of $K$, $g_0$, $\mathsf{d}$, the position and of both torus exponents of
Corollary~\ref{8.22:cor-short-distance-volume}; $C_W(s)$ and
$c_{\mathcal{K}}$ depend on $L$,
$\mathfrak{p}$, $\vartheta$ and $g$ only, $B^{\sharp}$ depends on $L$ and $\mathfrak{p}$, and
$b_0$ is a fixed number. So $s_T(g)$ is free of $K$, $g_0$, $m$,
$\mathsf{d}$ and position.

Since $s(\mathsf{d})$ is non-increasing in $\mathsf{d}$ (by its definition in
Notation~\ref{0.2:not-glossary}, the set over which the minimum is taken grows with $\mathsf{d}$), and
the admissible depths form the half-line $\{s\ge s_T(g)\}$, the condition $s(\mathsf{d})\ge s_T(g)$
holds exactly for $\mathsf{d}$ in an interval with left end $0$; it is a ceiling on the separation.
Both the condition $g\le\gamma_W^{\mathrm{VT}}$ and the condition $s\ge s_T(g)$ are one-sided:
an upper bound on $g$ and a lower bound on $s$. Neither bounds any of the parameters
$L,M,M_1,\kappa,\dots$ from above, so neither is a cap in the sense of
Definition~\ref{2.3:def-cap}.
\end{proof}

\begin{corollary}[The volume-transfer coupling ceiling and depth floor]
\label{8.22:cor-transfer-depth-floor}
Work in the setting of clause~(iv) of Theorem~\ref{3.1:thm-main}: gauge group $\mathrm{SU}(N)$ with $N=2$,
inside the regime of clause~(iii), with $\vartheta$ and the constants $k_W$, $c_{\mathcal{K}}$, $C_{\mathcal{K}}$ of
that clause. Let $\gamma_W^{\mathrm{VT}}$ be the coupling ceiling
of~\eqref{8.22:eq-expansion-consistency-a} and $s_T(g)$ the depth floor
of~\eqref{8.22:eq-expansion-consistency-b} (Definition~\ref{8.22:def-expansion-consistency}). Assume the
following two statements.
\begin{itemize}
\item[(a)] The two-point witness theorem on a single torus below the horizon, that is, for couplings $g$
with $m_{\mathbb{T}}(g_-(g))\le m_V(g)$, where $m_V(g)$ is the horizon exponent of Definition~\ref{8.18:def-volume-horizon},
in the following form. For every $g\le\gamma_W^{\mathbb{T}}$ such that $m':=m_{\mathbb{T}}(g_-(g))$ satisfies
$m'\le m_V(g)$ (here and below $m'$ denotes a torus exponent, as in
Corollary~\ref{8.22:cor-short-distance-volume}), the volume depth $s_{\mathrm{vol}}(g,m')$ of that theorem is
zero, and on the torus $\mathbb{T}_{m'}$, for every $g_0\in\,]0,g_-(g)]$ with $K=K(g_0,g)$, every $\mathsf{d}$
with $s_W(g)\le s(\mathsf{d})=:s$ and $K-s\ge k_W$, and every $(f_1,f_2)\in\mathfrak{T}(\mathsf{d};\vartheta)$ on
$\mathbb{T}_{m'}$,
\begin{equation}\label{8.22:eq-transfer-depth-floor-1}
S^T_{2;K,m'}(f_1,f_2)
=\zeta_{K-s}'^{\,2}\,b_0\,g_{K-s}^4\,\mathcal{K}_s^{(m')}(f_1,f_2)\,(1+R'),\qquad |R'|\le\tfrac14,
\end{equation}
where $\mathcal{K}_s^{(m')}$ is the kernel $\mathcal{K}_s$ read on the bare lattice of $\mathbb{T}_{m'}$,
$\zeta'_{K-s}\in[\frac23,\frac43]$, and
\begin{equation}\label{8.22:eq-transfer-depth-floor-8}
c_{\mathcal{K}}\|f_1\|_\infty\|f_2\|_\infty\le\mathcal{K}_s^{(m')}(f_1,f_2)\le
C_{\mathcal{K}}\|f_1\|_\infty\|f_2\|_\infty .
\end{equation}
\item[(b)] Short-distance volume independence of the two-point function,
Corollary~\ref{8.22:cor-short-distance-volume}.
\end{itemize}
Then for every $g\le\gamma_W^{\mathrm{VT}}$, every $g_0\in\,]0,g_-(g)]$ with $K=K(g_0,g)$, every
$m\ge m_{\mathbb{T}}(g_-(g))$, every $\mathsf{d}$ with $\max\{s_W(g),s_T(g)\}\le s(\mathsf{d})=:s$ and
$K-s\ge k_W$, and every $(f_1,f_2)\in\mathfrak{T}(\mathsf{d};\vartheta)$ on $\mathbb{T}_m$,
\begin{equation}\label{8.22:eq-transfer-depth-floor-a}
S^T_{2;K,m}(f_1,f_2)=\zeta_{K-s}'^{\,2}\,b_0\,g_{K-s}^4\,\mathcal{K}_s^{\star}(f_1,f_2)\,(1+R_K),
\qquad |R_K|\le\tfrac12,
\end{equation}
where $\mathcal{K}_s^{\star}$ is $\mathcal{K}_s$ read on the reference torus $\mathbb{T}_{m_{\star}(g)}$,
$\zeta'_{K-s}\in[\frac23,\frac43]$ as in~(a), and
$\mathcal{K}_s^{\star}$ obeys the kernel bounds~\eqref{8.22:eq-transfer-depth-floor-8} with
$\mathcal{K}_s^{\star}$ in place of $\mathcal{K}_s^{(m')}$. Hence, uniformly in $K$, $g_0$, $m$ and the position,
\begin{equation}\label{8.22:eq-transfer-depth-floor-2}
\tfrac29\,b_0\,g_{K-s}^4\,\mathcal{K}_s^{\star}(f_1,f_2)\le S^T_{2;K,m}(f_1,f_2)
\le\tfrac83\,b_0\,g_{K-s}^4\,\mathcal{K}_s^{\star}(f_1,f_2),
\end{equation}
and the coupling form of the two-point law, the last display of clause~(iv) of
Theorem~\ref{3.1:thm-main}, holds:
\begin{equation}\label{8.22:eq-transfer-depth-floor-3}
\begin{split}
\tfrac29\,b_0\,c_{\mathcal{K}}\,\|f_1\|_\infty\|f_2\|_\infty\,\bigl(g^{-2}+B^{\sharp}(s+1)\bigr)^{-2}
&\le S^T_{2;K,m}(f_1,f_2)\\
&\le\tfrac83\,b_0\,C_{\mathcal{K}}\,\|f_1\|_\infty\|f_2\|_\infty\,
\bigl(g^{-2}+\lambda^{\sharp}s-c_5\bigr)^{-2}.
\end{split}
\end{equation}
\end{corollary}

\begin{proof}
Fix $g$, $g_0$, $K$, $m$, $\mathsf{d}$, $s$ and $(f_1,f_2)$ as in the conclusion.

\emph{The reference torus.} Put $m':=m_{\star}(g)=\max\{m_{\mathbb{T}}(g_-(g)),\log_L\ell\}$. With the ball
radius $R=1$ of Corollary~\ref{8.22:cor-short-distance-volume}, the cell side $\ell$ of the volume-transfer
theorem is the least power of $L$ with $\ell\ge16(R+1)=32$; hence $\ell/L<32$. Since $L^{10}\ge100$,
$M\ge1$, and $m_{\mathbb{T}}(g')=11+\log_LM+\log_LR(g')$ with $R(g')\ge1$,
\begin{equation*}
\ell\le 32L\le L^{11}\le L^{11}M\le L^{m_{\mathbb{T}}(g_-(g))},
\end{equation*}
so that $m_{\star}(g)=m_{\mathbb{T}}(g_-(g))$ (see also~\eqref{8.22:eq-expansion-consistency-a}).
Thus $m'=m_{\mathbb{T}}(g_-(g))\le m$ by the hypothesis on $m$, and $L^{m'}\ge L^{11}$. Since
$g\le\gamma_W^{\mathrm{VT}}=\gamma_{\star}^{V}$, the ceiling~\eqref{8.22:eq-expansion-consistency-a} gives
$m_{\star}(g)\le m_V(g)$, that is $m'\le m_V(g)$; and $\gamma_{\star}^V\le\gamma_W^{\mathbb{T}}$ by its
definition as a supremum over $\gamma\le\gamma_W^{\mathbb{T}}$, so $g\le\gamma_W^{\mathbb{T}}$. Moreover
$s\ge\max\{s_W(g),s_T(g)\}\ge s_W(g)$ and $K-s\ge k_W$.

\emph{The depth and the scale.} Recall that $s(\mathsf{d})$ is the least $s\ge0$ with
$C_W(s)L^{-s}\le\mathsf{d}$, and that $C_W\ge1$ is non-decreasing. The first condition
in~\eqref{8.22:eq-expansion-consistency-b} fails at $s=0$, since $L\,C_W(0)\ge L>\frac1{40}$; as it must hold
for every $s\ge s_T(g)$, we have $s_T(g)\ge1$, hence $s\ge1$. Minimality of $s(\mathsf{d})$ then gives
$C_W(s-1)L^{-(s-1)}>\mathsf{d}$, and since $C_W$ is non-decreasing,
\begin{equation}\label{8.22:eq-transfer-depth-floor-6}
\mathsf{d}<C_W(s-1)\,L^{-(s-1)}\le L\,C_W(s)\,L^{-s}.
\end{equation}
Since $s\ge s_T(g)$, the first condition in~\eqref{8.22:eq-expansion-consistency-b} gives
$L\,C_W(s)L^{-s}\le\frac1{40}$, hence $\mathsf{d}<\frac1{40}$. Since $k_W\ge1$ and $C_W(s)\ge1$, the bare
spacing satisfies
\begin{equation*}
L^{-K}\le L^{-s-1}<C_W(s)\,L^{-s}\le\mathsf{d},
\end{equation*}
the last step by the definition of $s(\mathsf{d})$.

\emph{The pair in one ball and on the small torus.} A pair of $\mathfrak{T}(\mathsf{d};\vartheta)$ on
$\mathbb{T}_m$ has its supports in two balls of radius $\mathsf{d}$ at distance at most $6\mathsf{d}$. Their
centres are then at distance at most $8\mathsf{d}$. Every point of either ball lies within
$4\mathsf{d}+\mathsf{d}=5\mathsf{d}$ of the midpoint $y$ of the centres, so both supports lie in the closed ball
$\bar B(y,5\mathsf{d})$. Since $5\mathsf{d}<1<L^{m'}/2$, the pair restricted to $\bar B(y,5\mathsf{d})$ lifts to
$\mathbb{R}^4$; two lifts differ by a translation in $2L^m\mathbb{Z}^4$, which is contained in
$2L^{m'}\mathbb{Z}^4$ because $m\ge m'$, so the transported pair on $\mathbb{T}_{m'}$, and the transport back to
$\mathbb{T}_m$, which is the given pair, do not depend on the lift. On $\mathbb{T}_{m'}$ the transported pair
again belongs to $\mathfrak{T}(\mathsf{d};\vartheta)$, every condition of that class being local to the ball
and the torus distance on it being the Euclidean distance of the lift. We use the same letters for the
transported pair.

\emph{The small torus.} By the two preceding steps, $g\le\gamma_W^{\mathbb{T}}$, $m'\le m_V(g)$,
$s\ge s_W(g)$, $K-s\ge k_W$, and $(f_1,f_2)\in\mathfrak{T}(\mathsf{d};\vartheta)$ on $\mathbb{T}_{m'}$. Thus
hypothesis~(a) applies on $\mathbb{T}_{m'}$, with volume depth $s_{\mathrm{vol}}(g,m')=0$,
and~\eqref{8.22:eq-transfer-depth-floor-1} reads
\begin{equation}\label{8.22:eq-transfer-depth-floor-4}
S^T_{2;K,m'}(f_1,f_2)=\mathfrak{M}^{\star}(1+R'),\qquad |R'|\le\tfrac14,\qquad
\mathfrak{M}^{\star}:=\zeta_{K-s}'^{\,2}\,b_0\,g_{K-s}^4\,\mathcal{K}_s^{\star}(f_1,f_2),
\end{equation}
because $\mathcal{K}_s^{(m')}=\mathcal{K}_s^{\star}$ for $m'=m_{\star}(g)$. The same hypothesis gives the
kernel bounds~\eqref{8.22:eq-transfer-depth-floor-8} for $\mathcal{K}_s^{(m')}=\mathcal{K}_s^{\star}$, as
asserted in the statement.

\emph{Lower bound on the main term.} Hypothesis~(a) gives $\zeta'_{K-s}\ge\frac23$, and clause~(0) of
Theorem~\ref{3.1:thm-main} gives the bound $g_{K-s}^{4}\ge\bigl(g^{-2}+B^{\sharp}(s+1)\bigr)^{-2}$ on the
running coupling. With the lower kernel bound this yields
\begin{equation}\label{8.22:eq-transfer-depth-floor-5}
\mathfrak{M}^{\star}\ge\tfrac49\,b_0\,c_{\mathcal{K}}\,\bigl(g^{-2}+B^{\sharp}(s+1)\bigr)^{-2}
\,\|f_1\|_\infty\|f_2\|_\infty .
\end{equation}

\emph{The transfer to $\mathbb{T}_m$.} The functions $f_1,f_2$ are real Lipschitz test functions of the class
$\mathfrak{T}(\mathsf{d};\vartheta)$, and $\operatorname{Lip}f_i\le\|f_i\|_\infty/(\vartheta\mathsf{d})$, as
Corollary~\ref{8.22:cor-short-distance-volume} requires; their supports lie in one ball of radius
$5\mathsf{d}$ on both tori, and $L^{-K}<\mathsf{d}<\frac1{40}$ by the preceding steps. Finally
$m\ge m'=m_{\star}(g)$, and $g\le\gamma_W^{\mathbb{T}}\le\gamma_J$, $g_0\in\,]0,g_-(g)]$, $K=K(g_0,g)$, so the
coupling hypotheses of Corollary~\ref{8.22:cor-short-distance-volume} hold as well. Thus
Corollary~\ref{8.22:cor-short-distance-volume} applies to $(f_1,f_2)$ and the pair of exponents $m\ge m'$.
Combined with~\eqref{8.22:eq-transfer-depth-floor-6}, with the second condition
in~\eqref{8.22:eq-expansion-consistency-b} (valid because $s\ge s_T(g)$), and
with~\eqref{8.22:eq-transfer-depth-floor-5}, it gives
\begin{equation}\label{8.22:eq-transfer-depth-floor-7}
\begin{split}
\bigl|S^T_{2;K,m}(f_1,f_2)-S^T_{2;K,m'}(f_1,f_2)\bigr|
&\le C_T(g)\,\mathsf{d}^2\,\|f_1\|_\infty\|f_2\|_\infty\\
&\le C_T(g)\,L^2C_W(s)^2L^{-2s}\,\|f_1\|_\infty\|f_2\|_\infty\\
&\le\tfrac19\,b_0\,c_{\mathcal{K}}\,\bigl(g^{-2}+B^{\sharp}(s+1)\bigr)^{-2}\|f_1\|_\infty\|f_2\|_\infty\\
&\le\tfrac14\,\mathfrak{M}^{\star},
\end{split}
\end{equation}
the last step because $\frac19=\frac14\cdot\frac49$.

\emph{The relative error.} Since $f_1,f_2\not\equiv0$ for pairs of $\mathfrak{T}(\mathsf{d};\vartheta)$,
$\mathfrak{M}^{\star}>0$ by~\eqref{8.22:eq-transfer-depth-floor-5}. Put
\begin{equation*}
R_K:=R'+\frac{S^T_{2;K,m}(f_1,f_2)-S^T_{2;K,m'}(f_1,f_2)}{\mathfrak{M}^{\star}} .
\end{equation*}
Then $S^T_{2;K,m}(f_1,f_2)=\mathfrak{M}^{\star}(1+R_K)$
by~\eqref{8.22:eq-transfer-depth-floor-4}. By~\eqref{8.22:eq-transfer-depth-floor-4}
and~\eqref{8.22:eq-transfer-depth-floor-7}, $|R_K|\le\frac14+\frac14=\frac12$. This
is~\eqref{8.22:eq-transfer-depth-floor-a}. None of the constants used (those of hypothesis~(a), of
Corollary~\ref{8.22:cor-short-distance-volume}, and of~\eqref{8.22:eq-expansion-consistency-b}) depends on
$K$, $g_0$, $m$ or the position, and both hypotheses hold uniformly in these. Hence the bound
$|R_K|\le\frac12$ is uniform in them.

\emph{The consequences.} Hypothesis~(a) gives
$\zeta'_{K-s}\in[\frac23,\frac43]$. Hence $\zeta_{K-s}'^{\,2}\in[\frac49,\frac{16}9]$, while
$1+R_K\in[\frac12,\frac32]$ and $\mathcal{K}_s^{\star}(f_1,f_2)\ge0$. Since $\frac49\cdot\frac12=\frac29$
and $\frac{16}9\cdot\frac32=\frac83$, \eqref{8.22:eq-transfer-depth-floor-a}
gives~\eqref{8.22:eq-transfer-depth-floor-2}. Inserting in~\eqref{8.22:eq-transfer-depth-floor-2} the kernel
bounds and the two-sided bound of clause~(0) on the running coupling,
\begin{equation*}
\bigl(g^{-2}+B^{\sharp}(s+1)\bigr)^{-2}\le g_{K-s}^4\le\bigl(g^{-2}+\lambda^{\sharp}s-c_5\bigr)^{-2},
\end{equation*}
gives~\eqref{8.22:eq-transfer-depth-floor-3}. All these bounds are uniform in $K$, $g_0$, $m$ and the
position.
\end{proof}

\begin{remark}[The two-point law uniformly in the volume]\label{8.22:rem-two-point-uniform}
Write $\mathcal{K}_s^{(m)}$ for the free-lattice kernel $\mathcal{K}_s$ of clause~(iv) of
Theorem~\ref{3.1:thm-main} formed on the bare lattice of $\mathbb{T}_m$ itself. The same kernel
formed on the bare lattice of the reference torus $\mathbb{T}_{m_{\star}(g)}$ is the kernel
$\mathcal{K}_s^{\star}$ of Corollary~\ref{8.22:cor-transfer-depth-floor}, which gives the two-point
law, uniformly in $K$, $g_0$, $m\ge m_{\mathbb{T}}(g_-(g))$ and the position, with the kernel
$\mathcal{K}_s^{\star}$.
The two kernels obey the same two-sided bounds: under the hypotheses of
Corollary~\ref{8.22:cor-transfer-depth-floor}, with $s=s(\mathsf{d})$, for every
$(f_1,f_2)\in\mathfrak{T}(\mathsf{d};\vartheta)$,
\begin{equation*}
c_{\mathcal{K}}\,\|f_1\|_\infty\|f_2\|_\infty\;\le\;\mathcal{K}_s^{(m)}(f_1,f_2)\;\le\;
C_{\mathcal{K}}\,\|f_1\|_\infty\|f_2\|_\infty ,
\end{equation*}
and the same holds for $\mathcal{K}_s^{\star}$ with the same constants. They also have the same
continuum counterpart, the one given by the continuum kernel, namely (here $N=2$)
\begin{equation*}
3(N^2-1)\,\pi^{-4}\iint f_1(x)\,f_2(y)\,|x-y|^{-8}\,dx\,dy .
\end{equation*}
Corollary~\ref{8.22:cor-transfer-depth-floor} gives
\begin{equation*}
S^T_{2;K,m}(f_1,f_2)=\zeta'^{2}_{K-s}\,b_0\,g^4_{K-s}\,\mathcal{K}_s^{\star}(f_1,f_2)\,(1+R_K),
\qquad |R_K|\le\tfrac12 .
\end{equation*}
If this display is taken with
either kernel, the sandwich of $S^T_{2;K,m}(f_1,f_2)$ between $\tfrac29\,b_0\,g^4_{K-s}$ and
$\tfrac83\,b_0\,g^4_{K-s}$ times the kernel holds with that kernel; the coupling sandwich
\begin{align*}
\tfrac29\,b_0\,c_{\mathcal{K}}\,\|f_1\|_\infty\|f_2\|_\infty\,\bigl(g^{-2}+B^{\sharp}(s+1)\bigr)^{-2}
&\le S^T_{2;K,m}(f_1,f_2)\\
&\le \tfrac83\,b_0\,C_{\mathcal{K}}\,\|f_1\|_\infty\|f_2\|_\infty\,
\bigl(g^{-2}+\lambda^{\sharp}s-c_5\bigr)^{-2},
\end{align*}
and the properties of the limits listed with clause~(iv) in Theorem~\ref{3.1:thm-main}, depend on
the kernel only through these bounds and this continuum counterpart, and hold unchanged with either
kernel.

Suppose that, for every $m\ge m_{\mathbb{T}}(g_-(g))$, every $\mathsf{d}$ with $s(\mathsf{d})$
beyond one further depth floor not depending on $g$, and every
$(f_1,f_2)\in\mathfrak{T}(\mathsf{d};\vartheta)$, with $s=s(\mathsf{d})$,
\begin{equation}\label{8.22:eq-two-point-uniform-1}
\Bigl|\frac{\mathcal{K}_s^{(m)}(f_1,f_2)}{\mathcal{K}_s^{\star}(f_1,f_2)}-1\Bigr|\;\le\;\frac{1}{16}.
\end{equation}
Then, once that floor is added to the depth condition of
Corollary~\ref{8.22:cor-transfer-depth-floor} and the constant $9$ in the condition
\begin{equation*}
9\,C_T(g)\,L^2\,C_W(s)^2\,L^{-2s}\,\bigl(g^{-2}+B^{\sharp}(s+1)\bigr)^2\;\le\;b_0\,c_{\mathcal{K}},
\end{equation*}
which together with $L\,C_W(s)\,L^{-s}\le\tfrac1{40}$ defines the depth floor $s_T(g)$, is replaced
by $18$, the display of Corollary~\ref{8.22:cor-transfer-depth-floor} holds with
$\mathcal{K}_s^{(m)}$ in place of $\mathcal{K}_s^{\star}$ and relative error at most $\tfrac12$;
the volume-transfer error is then allowed the fraction $\tfrac18$ of the main term instead of
$\tfrac14$.

The bound \eqref{8.22:eq-two-point-uniform-1} concerns the free lattice field only. A proof would
compare both $\mathcal{K}_s^{(m)}$ and $\mathcal{K}_s^{\star}$ with the truncated kernel used in the
proof of the two-point witness theorem, within the tolerance of the comparison between the kernel
and its truncation made in that proof. It would use the two-sided kernel bound and the exact
periodisation of the step covariances of the truncated kernel, which are block-translation invariant
and have range, counted in cells, equal to $L$ raised to the truncation depth, small compared with
the side of either torus. We do not prove
\eqref{8.22:eq-two-point-uniform-1}, and nothing in this book depends on it:
clause~(iv) is used throughout with the kernel $\mathcal{K}_s^{\star}$.
\end{remark}

\begin{proof}
We prove the assertions of the first two paragraphs of the remark, those of the second under the
assumption stated there.
Let $m\ge m_{\mathbb{T}}(g_-(g))$. Both $\mathbb{T}_m$ and the reference torus
$\mathbb{T}_{m_{\star}(g)}$ belong to the family of tori admitted in
Corollary~\ref{8.22:cor-transfer-depth-floor}.

The two-sided kernel bound stated with the two-point witness theorem holds for $\mathcal{K}_s$
formed on the bare lattice of any torus of this family, with constants $c_{\mathcal{K}}$ and
$C_{\mathcal{K}}$ that do not depend on $K$, $g_0$, the torus exponent or the position; applied once
on $\mathbb{T}_m$ and once on $\mathbb{T}_{m_{\star}(g)}$, it gives the displayed bounds for
$\mathcal{K}_s^{(m)}$ and for $\mathcal{K}_s^{\star}$ with the same pair of constants. By the
continuum kernel, $\mathcal{K}_s$ formed on any torus of the family has one and the same continuum
counterpart, the one displayed in the remark.

The sandwich with the factors $\tfrac29$ and $\tfrac83$ is derived from the display of
Corollary~\ref{8.22:cor-transfer-depth-floor}, the window $\zeta'_{K-s}\in[\tfrac23,\tfrac43]$ and
the non-negativity of the kernel, and it carries the kernel of that display. The coupling sandwich is
derived from the display, the window, the two-sided kernel bound and clause~(0); the properties of
every limit listed with clause~(iv) in Theorem~\ref{3.1:thm-main} are derived from the display,
clause~(0), clause~(ii), reflection positivity of the lattice measure and the continuum kernel. Of
the kernel, these last derivations use only the two-sided kernel bound and the continuum counterpart.
These two inputs are the same for $\mathcal{K}_s^{(m)}$ and $\mathcal{K}_s^{\star}$. Hence the
coupling sandwich and every property of the limits listed with clause~(iv) hold unchanged whichever
of $\mathcal{K}_s^{(m)}$, $\mathcal{K}_s^{\star}$ the display carries.

Assume now the hypotheses of Corollary~\ref{8.22:cor-transfer-depth-floor} with the floor $s_T(g)$
replaced by the floor obtained from the same two conditions with $18$ in place of $9$ (a floor at
least as large, since the condition with $18$ implies the one with $9$), and with the further floor
added to the depth condition, and suppose that
\eqref{8.22:eq-two-point-uniform-1} holds for the pair $(f_1,f_2)$. Put
\begin{equation*}
\mathfrak{M}^{\star}:=\zeta'^{2}_{K-s}\,b_0\,g^4_{K-s}\,\mathcal{K}_s^{\star}(f_1,f_2);
\end{equation*}
it is positive, since $\zeta'_{K-s}\ge\tfrac23$, $b_0>0$ and
$\mathcal{K}_s^{\star}(f_1,f_2)\ge c_{\mathcal{K}}\|f_1\|_\infty\|f_2\|_\infty>0$. In the proof of
Corollary~\ref{8.22:cor-transfer-depth-floor} the two-point law on the reference torus gives
$S^T_{2;K,m_{\star}(g)}(f_1,f_2)=\mathfrak{M}^{\star}(1+R')$ with $|R'|\le\tfrac14$, and the two
conditions defining the floor, together with the volume-transfer bound, give
\begin{equation*}
\bigl|S^T_{2;K,m}(f_1,f_2)-S^T_{2;K,m_{\star}(g)}(f_1,f_2)\bigr|\le\tfrac14\,\mathfrak{M}^{\star}
\end{equation*}
when the constant is $9$. With $18$ in place of $9$ in the condition displayed in the remark the
admitted difference is halved, and the same argument gives the bound $\tfrac18\mathfrak{M}^{\star}$.
Hence $S^T_{2;K,m}(f_1,f_2)=\mathfrak{M}^{\star}(1+R_K)$ with
$|R_K|\le\tfrac14+\tfrac18=\tfrac38$. The bound \eqref{8.22:eq-two-point-uniform-1} gives
$\mathcal{K}_s^{(m)}(f_1,f_2)/\mathcal{K}_s^{\star}(f_1,f_2)\in[\tfrac{15}{16},\tfrac{17}{16}]$,
therefore
\begin{align*}
\frac{S^T_{2;K,m}(f_1,f_2)}{\zeta'^{2}_{K-s}\,b_0\,g^4_{K-s}\,\mathcal{K}_s^{(m)}(f_1,f_2)}
&=\frac{\mathcal{K}_s^{\star}(f_1,f_2)}{\mathcal{K}_s^{(m)}(f_1,f_2)}\,(1+R_K)\\
&\in\Bigl[\frac{16}{17}\cdot\frac58,\;\frac{16}{15}\cdot\frac{11}{8}\Bigr]
=\Bigl[\frac{10}{17},\;\frac{22}{15}\Bigr],
\end{align*}
and the relative error against $\mathcal{K}_s^{(m)}$ is at most $\tfrac{7}{15}<\tfrac12$.
\end{proof}

\begin{corollary}[The kernel read on the reference torus]\label{8.22:cor-reference-torus-kernel}
Let $g$, $\vartheta$ and the constants $b_0$, $c_{\mathcal{K}}$, $C_{\mathcal{K}}$, $B^{\sharp}$,
$\lambda^{\sharp}$, $c_5$, $k_W$ be as in Corollary~\ref{8.22:cor-transfer-depth-floor}; in
particular $N=2$ and $g\le\gamma_W^{\mathrm{VT}}$.
Here $s_T(g)$ is the depth floor
of \eqref{8.22:eq-expansion-consistency-b}, and $s_W(g)$ is the depth floor of the two-point
witness theorem. Let $S^T_{n;\infty}$ be any limit in the sense of clause
(ii-$\mathbb{R}^4$) of Theorem~\ref{3.1:thm-main} at the reference coupling $g$, taken along a
sequence $(g_0^{(j)},K_j,m_j)$ with $g_0^{(j)}\in\,]0,g_-(g)]$, $K_j=K(g_0^{(j)},g)$,
$K_j\to\infty$ and $m_j\to\infty$. Let $\mathsf{d}$ satisfy
\begin{equation*}
\max\{s_W(g),s_T(g)\}\le s(\mathsf{d})=:s,
\end{equation*}
that is, let $\mathsf{d}$ lie below the separation ceilings corresponding to these two depth floors.
Then for every pair $(f_1,f_2)$ of $C^1$ functions on $\mathbb{R}^4$ belonging to
$\mathfrak{T}(\mathsf{d};\vartheta)$,
\begin{equation}\label{8.22:eq-reference-torus-kernel-1}
\begin{split}
\tfrac{2}{9}\,b_0\,c_{\mathcal{K}}\,\|f_1\|_\infty\|f_2\|_\infty
\bigl(g^{-2}+B^{\sharp}(s+1)\bigr)^{-2}
&\le S^T_{2;\infty}(f_1,f_2)\\
&\le \tfrac{8}{3}\,b_0\,C_{\mathcal{K}}\,\|f_1\|_\infty\|f_2\|_\infty
\bigl(g^{-2}+\lambda^{\sharp}s-c_5\bigr)^{-2}.
\end{split}
\end{equation}
\end{corollary}

\begin{proof}
Fix $\mathsf{d}$, $s=s(\mathsf{d})$ and a pair $(f_1,f_2)$ of $C^1$ functions on $\mathbb{R}^4$
belonging to $\mathfrak{T}(\mathsf{d};\vartheta)$
as in the statement. Let $I=[a,b]$ be the closed interval whose endpoints $a$ and $b$ are the
left-hand and the right-hand members of \eqref{8.22:eq-reference-torus-kernel-1}. We show three
facts:
\begin{itemize}
\item the number $S^T_{2;K_j,m_j}(f_1,f_2)$ lies in $I$ for all but finitely many $j$;
\item the set $I$ does not depend on $j$;
\item the set $I$ is closed.
\end{itemize}

\emph{Membership for the members of the sequence.} Consider an index $j$ with
$m_j\ge m_{\mathbb{T}}(g_-(g))$ and $K_j-s\ge k_W$, and such that the pair $(f_1,f_2)$,
transported to $\mathbb{T}_{m_j}$ as in clause (ii-$\mathbb{R}^4$), belongs to
$\mathfrak{T}(\mathsf{d};\vartheta)$ on $\mathbb{T}_{m_j}$. We check the hypotheses of
Corollary~\ref{8.22:cor-transfer-depth-floor} for
the triple $(g_0^{(j)},K_j,m_j)$ and the pair $(f_1,f_2)$.
\begin{itemize}
\item By hypothesis, $g_0^{(j)}\in\,]0,g_-(g)]$ and $K_j=K(g_0^{(j)},g)$.
\item By the choice of $j$, $m_j\ge m_{\mathbb{T}}(g_-(g))$.
\item The depth floor $\max\{s_W(g),s_T(g)\}\le s(\mathsf{d})=s$ holds by hypothesis.
\item The inequality $K_j-s\ge k_W$ is the ultraviolet floor on $K$ required there.
\item By the choice of $j$, the transported pair $(f_1,f_2)$ belongs to
$\mathfrak{T}(\mathsf{d};\vartheta)$ on $\mathbb{T}_{m_j}$.
\end{itemize}
Corollary~\ref{8.22:cor-transfer-depth-floor} therefore gives the representation
\eqref{8.22:eq-transfer-depth-floor-a} of $S^T_{2;K_j,m_j}(f_1,f_2)$, with the kernel
$\mathcal{K}_s^{\star}$ read on the reference torus $\mathbb{T}_{m_{\star}(g)}$ of
Definition~\ref{8.22:def-expansion-consistency} and with
$|R_{K_j}|\le\frac12$. That corollary also gives the kernel bounds
\begin{equation*}
c_{\mathcal{K}}\|f_1\|_\infty\|f_2\|_\infty\le\mathcal{K}_s^{\star}(f_1,f_2)
\le C_{\mathcal{K}}\|f_1\|_\infty\|f_2\|_\infty,
\end{equation*}
and, as its consequence, the sandwich
\begin{equation*}
\tfrac{2}{9}\,b_0\,g^4_{K_j-s}\,\mathcal{K}_s^{\star}(f_1,f_2)
\le S^T_{2;K_j,m_j}(f_1,f_2)
\le\tfrac{8}{3}\,b_0\,g^4_{K_j-s}\,\mathcal{K}_s^{\star}(f_1,f_2).
\end{equation*}

Corollary~\ref{8.22:cor-transfer-depth-floor} also gives the bound obtained by inserting into this
sandwich the kernel bounds and the bounds of clause (0) of Theorem~\ref{3.1:thm-main} on the
running coupling $g_{K_j-s}$: these are exactly the two inequalities of
\eqref{8.22:eq-reference-torus-kernel-1}, with $S^T_{2;K_j,m_j}(f_1,f_2)$ in place of
$S^T_{2;\infty}(f_1,f_2)$. That is, $S^T_{2;K_j,m_j}(f_1,f_2)\in I$.

Since $K_j\to\infty$ and $m_j\to\infty$ along the sequence, the conditions
$m_j\ge m_{\mathbb{T}}(g_-(g))$
and $K_j-s\ge k_W$ hold for all $j$ from some index on. Once the side $2L^{m_j}$ of
$\mathbb{T}_{m_j}$ is large compared with the supports of $f_1$ and $f_2$, the transport to
$\mathbb{T}_{m_j}$ leaves the supports, their distance, the sup norms, the integrals and the
Lipschitz constants unchanged, so the transported pair belongs to
$\mathfrak{T}(\mathsf{d};\vartheta)$ on $\mathbb{T}_{m_j}$ for all $j$ from some index on. So the
membership just proved holds for all but finitely many $j$.

\emph{The interval is fixed and closed.} The endpoints of $I$ are built from $b_0$,
$c_{\mathcal{K}}$, $C_{\mathcal{K}}$, $g$, $s$, $B^{\sharp}$, $\lambda^{\sharp}$, $c_5$ and the
sup norms of $f_1$ and $f_2$. None of these involves $g_0^{(j)}$, $K_j$ or $m_j$, so $I$ does
not depend on $j$. The interval $I$ is closed by its definition.

\emph{Passage to the limit.} By definition, $S^T_{2;\infty}(f_1,f_2)$ is the limit of the
numbers $S^T_{2;K_j,m_j}(f_1,f_2)$. All but finitely many of these numbers lie in the fixed
closed set $I$. Hence the limit lies in $I$, which is \eqref{8.22:eq-reference-torus-kernel-1}.

No limits are exchanged in this argument. The two limits $K\to\infty$ and $m\to\infty$ are
taken jointly, as in clause (ii-$\mathbb{R}^4$).
\end{proof}

\medskip
\noindent\emph{Consequences in the infinite-volume limit.}
Consider any limit of clause (ii-$\mathbb{R}^4$) of Theorem~\ref{3.1:thm-main}, any $\mathsf{d}$
with $\max\{s_W(g),s_T(g)\}\le s(\mathsf{d})=:s$, and any pair $(f_1,f_2)$
of $C^1$ functions on $\mathbb{R}^4$ belonging to $\mathfrak{T}(\mathsf{d};\vartheta)$, all as in
Corollary~\ref{8.22:cor-reference-torus-kernel}, with neither function identically zero.
The left-hand member of \eqref{8.22:eq-reference-torus-kernel-1} is then strictly positive,
because $b_0>0$ (at $N=2$, $b_0=3/2$) and $c_{\mathcal{K}}>0$. Hence $S^T_{2;\infty}$ does not
vanish identically on such pairs.

Pairs in $\mathfrak{T}(\mathsf{d};\vartheta)$ have disjoint supports, since the distance between
the two supports is at least $\mathsf{d}>0$. By the dichotomy for the curvature Hilbert space, in
a limit of clause (ii-$\mathbb{R}^4$) the identical vanishing of $S^T_{2;\infty}$ on pairs with
disjoint supports is equivalent to $\dim\mathcal{H}^{\mathrm{curv}}=1$. By the preceding paragraph
this vanishing fails; hence $\dim\mathcal{H}^{\mathrm{curv}}\neq 1$, and since
$\mathcal{H}^{\mathrm{curv}}$ contains the vacuum vector $\Omega$, in every such limit at a
reference coupling $g\le\gamma_W^{\mathrm{VT}}$, under the hypotheses of
Corollary~\ref{8.22:cor-reference-torus-kernel},
$\dim\mathcal{H}^{\mathrm{curv}}\ge 2$.

Now fix a pair $(f_1,f_2)$ of $C^1$ functions on $\mathbb{R}^4$ belonging to
$\mathfrak{T}(1;\vartheta)$, and for $0<\mathsf{d}<1$ put $f_i^{\mathsf{d}}(x):=f_i(x/\mathsf{d})$,
$i=1,2$. The dilated pair $(f_1^{\mathsf{d}},f_2^{\mathsf{d}})$ belongs to
$\mathfrak{T}(\mathsf{d};\vartheta)$: each support is a ball of radius $1$ dilated by
$\mathsf{d}$, the distance between the supports is multiplied by $\mathsf{d}$,
$\|f_i^{\mathsf{d}}\|_\infty=\|f_i\|_\infty$, $\int f_i^{\mathsf{d}}=\mathsf{d}^4\int f_i$ and
$\operatorname{Lip}f_i^{\mathsf{d}}=\mathsf{d}^{-1}\operatorname{Lip}f_i$. For $\mathsf{d}$ with
$\max\{s_W(g),s_T(g)\}\le s(\mathsf{d})$ put
\begin{equation*}
S(\mathsf{d}):=\frac{S^T_{2;\infty}(f_1^{\mathsf{d}},f_2^{\mathsf{d}})}
{\|f_1\|_\infty\|f_2\|_\infty}.
\end{equation*}
We read both bounds of \eqref{8.22:eq-reference-torus-kernel-1} at $s=s(\mathsf{d})$ on the
dilated pairs, in every limit. Since $C_W\ge1$, every $s\ge0$ with $C_W(s)L^{-s}\le\mathsf{d}$
satisfies $L^{-s}\le\mathsf{d}$, so $s(\mathsf{d})\ge\log_L(1/\mathsf{d})$; in particular
$s(\mathsf{d})\ge1$, and the minimality of $s(\mathsf{d})$ gives
$C_W(s(\mathsf{d})-1)L^{-s(\mathsf{d})+1}>\mathsf{d}$. That is,
\begin{equation*}
\log_L(1/\mathsf{d})\le s(\mathsf{d})
<1+\log_L(1/\mathsf{d})+\log_L C_W\bigl(s(\mathsf{d})-1\bigr).
\end{equation*}
Since $s(\mathsf{d})\ge\log_L(1/\mathsf{d})$, the condition
$\max\{s_W(g),s_T(g)\}\le s(\mathsf{d})$ holds for all sufficiently small $\mathsf{d}$, so that
$S(\mathsf{d})$ is defined for such $\mathsf{d}$.
Since $C_W$ is non-decreasing and polylogarithmic in $\bar{x}_s=g^{-2}+B^{\sharp}(s+1)$, the last
term grows at most logarithmically in $s(\mathsf{d})$, so $s(\mathsf{d})$ is at most a constant
multiple of $\log(1/\mathsf{d})$ for all sufficiently small $\mathsf{d}$. Inserting this bound into
the left member of \eqref{8.22:eq-reference-torus-kernel-1}, and $s(\mathsf{d})\ge\log_L(1/\mathsf{d})$
into its right member, gives, for all sufficiently small $\mathsf{d}$,
\begin{equation*}
S(\mathsf{d})\ge c\,\bigl(\log(1/\mathsf{d})\bigr)^{-2},
\qquad
S(\mathsf{d})\le C\,\bigl(\log(1/\mathsf{d})\bigr)^{-2}.
\end{equation*}
Thus $S(\mathsf{d})>0$, and $S(\mathsf{d})\to0$ as $\mathsf{d}\downarrow0$ more slowly than every
power of $\mathsf{d}$: $S(\mathsf{d})\,\mathsf{d}^{-\varepsilon}\to\infty$ for every
$\varepsilon>0$.
Here $c$ and $C$ depend on $g$, $L$, $b_0$, $c_{\mathcal{K}}$, $C_{\mathcal{K}}$, $B^{\sharp}$,
$\lambda^{\sharp}$, $c_5$ and, for $c$, on the function $C_W$, but not on $\mathsf{d}$; the upper
bound uses only $s(\mathsf{d})\ge\log_L(1/\mathsf{d})$, so $C$ does not involve $C_W$.

\begin{remark}[The two-point law in infinite volume]\label{8.22:rem-infinite-volume-law}
The hypotheses are the following:
\begin{itemize}
\item[(a)] the short-distance volume independence of the two-point function, Corollary~\ref{8.22:cor-short-distance-volume}, together with its hypotheses on $g$ and its exponent $m_{\star}(g)$ of the reference torus;
\item[(b)] the infinite-volume limits of clause (ii-$\mathbb{R}^4$) of Theorem~\ref{3.1:thm-main};
\item[(c)] the limits of clause (ii-$\mathbb{T}$-b) of Theorem~\ref{3.1:thm-main} on the reference torus $\mathbb{T}_{m_{\star}(g)}$ of (a);
\item[(d)] for the last assertion only, clause (v-c)(2) of Theorem~\ref{3.1:thm-main}.
\end{itemize}
Let $\sigma\in\operatorname{Lim}_{\infty}(g)$ be a limit of clause (ii-$\mathbb{R}^4$) taken along a sequence $(g_0^{(j)},K_j,m_j)$, and let $\sigma_{\star}$ be a limit of clause (ii-$\mathbb{T}$-b) on the torus $\mathbb{T}_{m_{\star}(g)}$ taken along a subsequence of $(g_0^{(j)},K_j)$. Write $S^T_2[\sigma]$ and $S^T_2[\sigma_{\star}]$ for their connected two-point functions; on pairs of real functions of class $C^1$ with compact, disjoint supports these are the limits of the numbers $S^T_{2;K_j,m_j}$, respectively $S^T_{2;K_j,m_{\star}(g)}$, along the sequences defining $\sigma$ and $\sigma_{\star}$.

Let $\vartheta\in\,]0,1]$ be the parameter of Corollary~\ref{8.22:cor-short-distance-volume} and let $0<\mathsf{d}\le 1/40$. Let $f_1,f_2$ be real functions of class $C^1$ on $\mathbb{R}^4$, of either sign, with compact, disjoint supports, both supported in one ball of radius $5\mathsf{d}$, with $\operatorname{Lip} f_i\le\|f_i\|_\infty/(\vartheta\mathsf{d})$, transported to $\mathbb{T}_{m_{\star}(g)}$ and to each $\mathbb{T}_{m_j}$ with $m_j\ge m_{\star}(g)$; since
\begin{equation*}
10\mathsf{d}\le 1/4<L^{m_{\star}(g)}\le L^{m_j},
\end{equation*}
this transport is unambiguous. Then
\begin{equation}\label{8.22:eq-infinite-volume-law-1}
\bigl|S^T_2[\sigma](f_1,f_2)-S^T_2[\sigma_{\star}](f_1,f_2)\bigr|
\le C_T(g)\,\mathsf{d}^2\,\|f_1\|_\infty\|f_2\|_\infty .
\end{equation}
The constant $C_T(g)$ is that of Corollary~\ref{8.22:cor-short-distance-volume}.

Under the proviso of clause (v-c)(2) of Theorem~\ref{3.1:thm-main}, $\sigma_{\star}$ is one of the limits $S^{(\hat w)}$ of that clause.

The remark compares infinite-volume limits with limits on one fixed torus at short distances only. It does not assert that limits of clause (ii-$\mathbb{R}^4$) taken along different sequences coincide. Nor does it assert that the Schwinger functions converge as the volume grows.
\end{remark}

\begin{proof}
Restrict the sequence $(g_0^{(j)},K_j,m_j)$ to the subsequence along which $\sigma_{\star}$ is taken. Along this restriction the numbers $S^T_{2;K_j,m_j}$ still converge to $S^T_2[\sigma]$. By definition of $\sigma_{\star}$, the numbers $S^T_{2;K_j,m_{\star}(g)}$ converge to $S^T_2[\sigma_{\star}]$.

Along a limit of clause (ii-$\mathbb{R}^4$) the exponents $K_j$ and $m_j$ tend to infinity. Hence, for all but finitely many indices $j$ of the subsequence, the lattice spacing satisfies $L^{-K_j}\le\mathsf{d}$ and $m_j\ge m_{\star}(g)$.

Fix such an index $j$. The pair $(f_1,f_2)$ meets the hypotheses of Corollary~\ref{8.22:cor-short-distance-volume}: the bound $\mathsf{d}\le 1/40$, the spacing bound $L^{-K_j}\le\mathsf{d}$, support in one ball of radius $5\mathsf{d}$, and the Lipschitz bound. Apply that corollary with $K=K_j$, $g_0=g_0^{(j)}$, $m=m_j$ and $m'=m_{\star}(g)$; this is admissible since $m_j\ge m_{\star}(g)$, and since every member $(g_0^{(j)},K_j)$ of a sequence of clause (ii-$\mathbb{R}^4$) has $g_0^{(j)}\in\,]0,g_-(g)]$ and $K_j=K(g_0^{(j)},g)$, as hypothesis (a) requires. Its bound \eqref{8.22:eq-short-distance-volume-a} then holds for each such $j$.

The right-hand side does not depend on $j$, by the uniformity of Corollary~\ref{8.22:cor-short-distance-volume} in $K$, in $g_0$ and in all $m\ge m'\ge m_{\star}(g)$. The left-hand side converges, along the subsequence, to the left-hand side of \eqref{8.22:eq-infinite-volume-law-1}, and a non-strict inequality with a fixed right-hand side is preserved under limits; this gives \eqref{8.22:eq-infinite-volume-law-1}.

The final assertion is clause (v-c)(2) of Theorem~\ref{3.1:thm-main}, read for the limit $\sigma_{\star}$ of clause (ii-$\mathbb{T}$-b) on $\mathbb{T}_{m_{\star}(g)}$.
\end{proof}

\begin{remark}[Large-field regions beyond the weak-coupling horizon]\label{8.18:prop-uniform-rarity}
No clause of Theorem~\ref{3.1:thm-main} uses a bound, uniform in the volume, on the normalised
weight of large-field regions that are created at some renormalisation step and are still live,
that is never healed, at the last level of the run of renormalisation steps prolonged on the same
bare lattice, on the tori of the family whose exponent $m$ exceeds the weak-coupling horizon
$m_V(g)$ fixed in the proposition of this chapter on the closure of the large-field age series in
finite volume. Such a bound is not proved here and is not among the statements on which
Theorem~\ref{3.1:thm-main} rests.
\end{remark}

\begin{definition}[Tori, members, runs and the history measure]\label{8.18:not-members-runs}
\emph{The family.} As in Definition~\ref{1.1:def-lattice}, $\mathbb{T}_m$ is the torus of side
$2L^m$ carrying the bare lattice of spacing $L^{-K}$, with
$\mathsf{N}:=2L^{m+K}$ bare sites per side. The numbers $x_k=g_k^{-2}$ are the running inverse
couplings of Definition~\ref{1.2:def-couplings}, subject to their flow bounds, for which we use the
flow constants of Definition~\ref{1.2:def-flow-constants}. We recall
\begin{equation*}
K=K(g_0,g):=\min\{k:\ x_k\le g^{-2}+B^{\sharp}\},\qquad
g_-(g):=\bigl(g^{-2}+B^{\sharp}\bigr)^{-1/2},
\end{equation*}
and
\begin{equation*}
m_{\mathbb{T}}(g'):=11+m'+\log_L R(g'),
\end{equation*}
where $M=L^{m'}$ and $R(y)$ is the least power of $L$ that is $\ge(\log y)^{r}$. Hypothesis (H5)
of Notation~\ref{2.2:not-hypotheses}, $m\ge m_{\mathbb{T}}(g_-(g))$, is a lower bound on the torus
exponent only; hypothesis (H6) reads $0<g_0\le g_-(g)$. The degree function is
$p_0(g)=A_0(\log g^{-2})^{p_0}$, with $p_0\ge 5r$ and $r\ge2$. We put
\begin{equation}\label{8.18:eq-members-runs-1}
\mathfrak{P}(y):=A_0(\log y)^{p_0},
\qquad\text{so that}\qquad
p_0(g_k)=\mathfrak{P}(x_k).
\end{equation}
Along the flow, for all levels $j\le k$,
\begin{equation}\label{8.18:eq-members-runs-a}
\mathfrak{P}(x_j)\ \ge\ \mathfrak{P}(x_k)-1 ,
\end{equation}
and, with $\underline{x}_u:=g^{-2}+\lambda^{\sharp}u-c_5$,
\begin{equation}\label{8.18:eq-members-runs-c}
\underline{x}_u\ \le\ x_{K-u}.
\end{equation}

\emph{Members and runs.} A \emph{bare datum} is the pair $(x_0=g_0^{-2},\ \mathsf{N}=2L^{m+K})$. A
\emph{member} over it is a pair $(g',N_T')$ such that, with $X':=g'^{-2}$ and
\begin{equation*}
K(x_0,g'):=\text{the largest } k \text{ with } x_j>X' \text{ for all } j\le k ,
\end{equation*}
one has $\mathsf{N}=N_T'L^{K(x_0,g')}$. The \emph{admissible runs} of the member stop at level
$K(x_0,g')-m''$ for an integer $m''\ge m_{\mathbb{T}}(g')$, on a lattice with $N_T'L^{m''}$ sites
per side. The function $K(x_0,\cdot)$ attached to members is a different function from the
first-passage level $K(g_0,\cdot)$ displayed above. We write $g_V$ for a coupling with
\begin{equation}\label{8.18:eq-members-runs-2}
K_V:=K(x_0,g_V)=K+m_{\mathbb{T}}(g_V),
\end{equation}
and put $a_V:=m-m_{\mathbb{T}}(g_V)$ and $N_T:=2L^{a_V}$. Then $a_V\ge0$, $(g_V,N_T)$ is a member
over the bare datum, and the run of $K$ steps ending on $\mathbb{T}_m$ used in
Theorem~\ref{3.1:thm-main} is its longest run. For a reference coupling $\gamma_*$ we put
$K_*:=K(x_0,\gamma_*)$ and $X_*:=\gamma_*^{-2}$; then $x_{K_*}\le X_*+b_+$, and along the
reference flow, for all levels $i\le j\le K_*$,
\begin{equation}\label{8.18:eq-members-runs-b}
\lambda(j-i)-2c_2\ \le\ x_i-x_j\ \le\ 3\lambda(j-i)+2c_2 ,
\end{equation}
where $\lambda$ denotes the rate constant of the reference flow.
We put $c_{\mathrm{gap}}:=\lceil 2c_2/\lambda\rceil+1$. A level $k^{\diamond}<K_*$ is
\emph{attainable} if and only if $x_{k^{\diamond}+1}<\min_{j\le k^{\diamond}}x_j$.

\emph{The history measure.} Ba{\l}aban's operations $\mathbf{R}$ and $\mathbf{T}$ are applied step
by step as a chain, starting from the bare configuration $U$ distributed as
$\rho_0\,dU/\int\rho_0$, and every term of every partition of unity is kept. A \emph{history}
$\omega$ is one chain of such choices, not resummed; $\Omega$ is the finite set of histories, and
$\tau_{\omega}\ge0$ is the total integral of the term of $\omega$ at the last level of the run. We
put
\begin{equation}\label{8.18:eq-members-runs-3}
\tilde{\nu}(E):=\Bigl(\sum_{\omega\in\Omega}\tau_{\omega}\Bigr)^{-1}
\sum_{\omega\in E}\tau_{\omega}\quad(E\subset\Omega),
\end{equation}
and, for every level $k$, with $\rho_k$ and $V_k$ as in
Definitions~\ref{1.5:def-densities} and~\ref{1.5:def-block-average},
\begin{equation*}
\sum_{\omega\in\Omega}\tau_{\omega}=\int\rho_0=\int\rho_k\,dV_k .
\end{equation*}
For a history $\omega$ and a
level $k$, $\mathcal{O}_k(\omega)$ is its large-field set at level $k$; a connected component of
$\mathcal{O}_k(\omega)$ is \emph{live} if it still carries its large-field factor at level $k$. For
a live component $Z$, $j(Z)$ is its oldest birth index, the index of the first large-field region
contained in $Z$ as in \cite[(1.79)]{B-CMP122b}; components that are merged or refreshed inherit
it, as in \cite[(1.85)]{B-CMP122b}; and $\operatorname{age}_k(Z):=k-j(Z)$. For a level-$k$ cell
$\square$ and an integer $a\ge0$,
\begin{align*}
\mathfrak{lv}_k(\square)&:=\{\omega:\ \square\subset\text{a live component of }
\mathcal{O}_k(\omega)\},\\
\mathfrak{lv}_k^{=a}(\square)&:=\{\omega:\ \square\subset\text{a live component } Z
\text{ of } \mathcal{O}_k(\omega)\text{ with } \operatorname{age}_k(Z)=a\},
\end{align*}
and
\begin{equation*}
q^{\mathrm{lv}}_k:=C_{\mathrm{lv}}\,e^{-c_{\mathrm{lv}}\mathfrak{P}(x_k)},\qquad
q^{=a}_k:=\max_{\square}\tilde{\nu}\bigl(\mathfrak{lv}_k^{=a}(\square)\bigr).
\end{equation*}

\emph{Witness letters.} The witness level is $k_0:=K-s$. With the running scale
$\check{R}_k:=R(\min_{i\le k}x_i)$, the age threshold is
$\mathfrak{l}:=L\,\check{R}_{k_0}+c_L$, where $c_L$ is the additive constant of the age threshold.
The letters $\mathfrak{T}(\mathsf{d};\vartheta)$, $b_0$,
$\mathcal{K}_s$, $\zeta'$, $s(\mathsf{d})$, $s_W(g)$, $r(s)$, $k_W$, $\gamma_W$, $C_W$,
$c_{\mathcal{K}}$, $C_{\mathcal{K}}$ have exactly the meaning fixed in clause~(iv) of
Theorem~\ref{3.1:thm-main}.
\end{definition}

\begin{proof}
The identity in \eqref{8.18:eq-members-runs-1} is a definitional unwinding:
\begin{equation*}
p_0(g_k)=A_0(\log g_k^{-2})^{p_0}=A_0(\log x_k)^{p_0}=\mathfrak{P}(x_k).
\end{equation*}

Inequality \eqref{8.18:eq-members-runs-a} is the drift bound for the degree function established
in the proof of the two-point witness theorem; it is used here as stated there. Inequality
\eqref{8.18:eq-members-runs-c} is the corollary of clause~(0) of
Theorem~\ref{3.1:thm-main}.

The run of $K$ steps on $\mathbb{T}_m$ used in Theorem~\ref{3.1:thm-main} is the longest run of the
member $(g_V,N_T)$. First, $a_V\ge0$ by (H5), since $m_{\mathbb{T}}$ is non-increasing in the
coupling; hence $N_T=2L^{a_V}$ is an integer. Second, $(g_V,N_T)$ is a member over the bare datum,
because by \eqref{8.18:eq-members-runs-2}
\begin{equation*}
N_T\,L^{K(x_0,g_V)}=2L^{a_V}L^{K+m_{\mathbb{T}}(g_V)}
=2L^{m-m_{\mathbb{T}}(g_V)+K+m_{\mathbb{T}}(g_V)}=2L^{m+K}=\mathsf{N}.
\end{equation*}
Third, every admissible run of this member stops at $K_V-m''$ with $m''\ge m_{\mathbb{T}}(g_V)$,
hence at a level $\le K_V-m_{\mathbb{T}}(g_V)$. The choice $m''=m_{\mathbb{T}}(g_V)$ attains this
level, $K_V-m_{\mathbb{T}}(g_V)=K$, on a lattice with
\begin{equation*}
N_T\,L^{m_{\mathbb{T}}(g_V)}=2L^{a_V+m_{\mathbb{T}}(g_V)}=2L^{m}
\end{equation*}
sites per side, which is the run of $K$ steps ending on $\mathbb{T}_m$.

The two-sided bound \eqref{8.18:eq-members-runs-b} along the reference flow, which does not depend
on the torus, and the bound $x_{K_*}\le X_*+b_+$ are taken by statement from the torus-independent
two-sided bound on the reference flow; they hold for the member at $\gamma_*$ up to the level $K_*$.

The weights satisfy $\tau_{\omega}\ge0$ by the positivity of the history weights, and, for every
$k$,
\begin{equation*}
\sum_{\omega}\tau_{\omega}=\int\rho_0=\int\rho_k\,dV_k
\end{equation*}
by the conservation of the total integral along the run. Hence the normalising sum in
\eqref{8.18:eq-members-runs-3} equals $\int\rho_0$, and $\tilde{\nu}$ is a non-negative set
function on $\Omega$. Since $\int\rho_0>0$, the bare configuration being distributed as
$\rho_0\,dU/\int\rho_0$, $\tilde{\nu}$ is a probability measure on the finite set $\Omega$.
\end{proof}

\begin{definition}[Volume horizon and volume-reading depth floor]\label{8.18:def-volume-horizon}
We use the tori, members and runs of Notation~\ref{8.18:not-members-runs}. In particular $B^{\sharp}$,
$g_-(g)=(g^{-2}+B^{\sharp})^{-1/2}$ and the torus exponent function $m_{\mathbb{T}}(\cdot)$ are as
there. We also use the reference-member constants $X_*$, $b_+$, $c_{\mathrm{gap}}$ and the family
constants $c_2$ and $\lambda>0$ fixed there. Further:
\begin{itemize}
\item $\theta>0$ is the reserve fraction, that is, the fraction entering the flat reserve carried by
the large-field regions live at the terminal level in the proof of the two-point witness theorem,
which fixes $\theta$;
\item $c_a$ is the free constant of the age estimates of that proof, with $q=e^{-c_a}$ the per-step
rate whose powers are the survival factors there; it satisfies $c_a\ge 9\log L+\log 2$;
\item $r(s)$ is the extraction depth of clause~(iv) of Theorem~\ref{3.1:thm-main};
\item $\underline{x}_u$ are the lower comparison couplings;
\item $\mathfrak{P}(y)=A_0(\log y)^{p_0}$, so that the degree function satisfies
$p_0(g_k)=\mathfrak{P}(x_k)$.
\end{itemize}
The depth $s$ ranges over the integers $s\ge 0$.

\emph{Volume horizon.} For $g>0$ put
\begin{equation}\label{8.18:eq-volume-horizon-a}
m_V(g):=\Bigl\lfloor \frac{g^{-2}-X_*-b_+-2c_2}{3\lambda}\Bigr\rfloor-1 .
\end{equation}

\emph{Exponent of the global quotient.} For a member with terminal inverse coupling $g'^{-2}$,
let $E_+^{(\theta)}(g')$ and $C_{\mathrm{low}}(g')$ be the upper and the lower part of the exponent
of the global quotient. Both depend on the member only through its terminal inverse coupling. Put
\begin{equation}\label{8.18:eq-volume-horizon-b}
\begin{gathered}
\mathsf{B}(g'):=E_+^{(\theta)}(g')+C_{\mathrm{low}}(g'),\\
\text{which satisfies}\quad \mathsf{B}(g')\le C_{\mathsf{B}}\bigl(1+\log g'^{-2}\bigr),
\end{gathered}
\end{equation}
where the constant $C_{\mathsf{B}}$ depends only on $(L,\mathfrak{p},\gamma,N)$.

\emph{Volume number.} For $g>0$ and a torus exponent $m$ put
\begin{equation}\label{8.18:eq-volume-horizon-c}
\bar{V}(g,m):=16\,C_{\mathsf{B}}\,\bigl(1+\log(g^{-2}+B^{\sharp})\bigr)\,
L^{4\left(m-m_V(g)+c_{\mathrm{gap}}+m_{\mathbb{T}}(g_-(g))\right)} .
\end{equation}

\emph{Depth floor.} The volume-reading depth floor $s_{\mathrm{vol}}(g,m)$ is defined by the
following display, whose first line defines the number $p_C(g)$ for every $g>0$:
\begin{equation}\label{8.18:eq-volume-horizon-d}
\begin{aligned}
p_C(g)&:=\max\bigl\{0,\ m_V(g)-c_{\mathrm{gap}}-m_{\mathbb{T}}(g_-(g))\bigr\},\\
s_{\mathrm{vol}}(g,m)&:=0\qquad\text{if } m\le m_V(g),\\
s_{\mathrm{vol}}(g,m)&:=\min\bigl\{s_0\ge 0:\ c_a\bigl(s-r(s)+p_C(g)\bigr)\ge \bar{V}(g,m)+\theta\\
&\qquad\qquad\ \text{for every } s\ge s_0\bigr\}\qquad\text{if } m>m_V(g).
\end{aligned}
\end{equation}
The left member of the inequality in the last line, $c_a\bigl(s-r(s)+p_C(g)\bigr)$, is affine in
$s-r(s)$ with slope $c_a\ge 9\log L+\log 2>0$.

\emph{Alternative depth floor.} We also use the floor $s_{\mathrm{vol}}^{(\mathfrak{P})}(g,m)$
defined by
\begin{equation}\label{8.18:eq-volume-horizon-e}
\begin{aligned}
s_{\mathrm{vol}}^{(\mathfrak{P})}(g,m)&:=0\qquad\text{if } m\le m_V(g),\\
s_{\mathrm{vol}}^{(\mathfrak{P})}(g,m)&:=\min\bigl\{s_0\ge 0:\
\tfrac12\theta\bigl(\mathfrak{P}(\underline{x}_{s-r(s)})-2\bigr)\ge \bar{V}(g,m)\\
&\qquad\qquad\ \text{for every } s\ge s_0\bigr\}\qquad\text{if } m>m_V(g).
\end{aligned}
\end{equation}
The two-point witness theorem asserts the existence of a depth threshold depending on $m$ as well
as on $g$; either of $s_{\mathrm{vol}}$ and $s_{\mathrm{vol}}^{(\mathfrak{P})}$ serves as that
threshold.

\emph{Properties.} The minima in \eqref{8.18:eq-volume-horizon-d} and
\eqref{8.18:eq-volume-horizon-e} exist. For each $(g,m)$ the set of admissible $s_0$ in either of
them is a half-line. The quantities \eqref{8.18:eq-volume-horizon-a}--\eqref{8.18:eq-volume-horizon-e}
depend neither on $K$, nor on $g_0$, nor on a position in the torus. Finally, $m_V(g)$ increases to
$\infty$ like $g^{-2}/(3\lambda)$ as $g\downarrow 0$.
\end{definition}

\begin{proof}
\emph{The volume horizon.} Write $y(g)=(g^{-2}-X_*-b_+-2c_2)/(3\lambda)$, with $\lambda>0$. From
$y-1<\lfloor y\rfloor\le y$ we get
\begin{equation*}
y(g)-2<m_V(g)\le y(g)-1 .
\end{equation*}
As $g$ decreases, $y(g)$ increases, and so does $\lfloor y(g)\rfloor$. Hence $m_V(g)$ is
non-decreasing as $g\downarrow 0$. Moreover $m_V(g)\big/\bigl(g^{-2}/(3\lambda)\bigr)\to 1$, because
$y(g)\big/\bigl(g^{-2}/(3\lambda)\bigr)\to 1$ and the two bounds above differ from $y(g)$ by at most
$2$. The definition depends only on $g$ and the constants $X_*,b_+,c_2,\lambda$; in particular it
does not depend on $K$.

\emph{The bound in \eqref{8.18:eq-volume-horizon-b}.} This inequality is the degree bound for the
exponent of the global quotient, used here as stated. Its two parts are as follows.
\begin{itemize}
\item The upper part $E_+^{(\theta)}(g')$ is of the order $1+\log g'^{-2}$ of clause~(i-b) of
Theorem~\ref{3.1:thm-main}.
\item The lower part $C_{\mathrm{low}}(g')$ is bounded by means of the lower bound of
clause~(i-b) of Theorem~\ref{3.1:thm-main} together with the normalisation of the Haar measure.
\end{itemize}
The constant $C_{\mathsf{B}}$ depends only on $(L,\mathfrak{p},\gamma,N)$. Both parts depend on the
member only through its terminal inverse coupling. Hence $\mathsf{B}$, as a function of $g'$, and
$C_{\mathsf{B}}$ depend neither on $K$, nor on $g_0$, nor on a position.

\emph{Independence of $K$, $g_0$ and position.} The right side of \eqref{8.18:eq-volume-horizon-c}
is built from $C_{\mathsf{B}}$, $B^{\sharp}$, $L$, $c_{\mathrm{gap}}$, $m$, $m_V(g)$ and
$m_{\mathbb{T}}(g_-(g))$. The function $g_-(g)$ depends only on $g$ and $B^{\sharp}$. Hence
$\bar{V}(g,m)$ depends on neither $K$ nor $g_0$ nor a position.

The same holds for $p_C(g)$. None of $c_a$, $\theta$, $r(s)$, $\underline{x}_u$ and $\mathfrak{P}$
depends on $K$, on $g_0$ or on a position. It then holds for the inequalities in
\eqref{8.18:eq-volume-horizon-d} and \eqref{8.18:eq-volume-horizon-e}, whose remaining entries are
$c_a$, $\theta$, $r(s)$, $\underline{x}_{s-r(s)}$ and $\mathfrak{P}$. It therefore holds for
$s_{\mathrm{vol}}$ and $s_{\mathrm{vol}}^{(\mathfrak{P})}$.

\emph{Existence of the minima and the half-line property.} Fix $(g,m)$ with $m>m_V(g)$; otherwise
both floors are $0$ and there is nothing to prove. For an integer $s\ge 0$ let $(I_s)$ denote the
inequality in \eqref{8.18:eq-volume-horizon-d} at $s$, or the one in
\eqref{8.18:eq-volume-horizon-e} at $s$. Let $\mathcal{S}$ be the set of integers $s_0\ge 0$ such
that $(I_s)$ holds for every $s\ge s_0$.

First, $\mathcal{S}$ is upward closed. If $s_0\in\mathcal{S}$ and $s_0'\ge s_0$, then every
$s\ge s_0'$ satisfies $s\ge s_0$, so $(I_s)$ holds and $s_0'\in\mathcal{S}$. Once $\mathcal{S}$ is
non-empty, it has a least element, being a non-empty set of non-negative integers, and by upward
closure it is then the half-line $\{s_0\ge \min\mathcal{S}\}$ of integers. This uses no
monotonicity of $s\mapsto s-r(s)$.

It remains to show that $\mathcal{S}$ is non-empty. Let $\mathcal{E}$ be the set of integers
$s\ge 0$ at which $(I_s)$ fails. If $\mathcal{E}$ is bounded, then either it is empty and
$0\in\mathcal{S}$, or it is non-empty and $1+\max\mathcal{E}\in\mathcal{S}$. In the first case
$\min\mathcal{S}=0$. In the second case $\min\mathcal{S}=1+\max\mathcal{E}$, because no
$s_0\le\max\mathcal{E}$ lies in $\mathcal{S}$: for such $s_0$ the inequality
$(I_{\max\mathcal{E}})$, with $\max\mathcal{E}\ge s_0$, fails. It therefore suffices to show that
$\mathcal{E}$ is bounded.

Since $r(s)$ is doubly logarithmic in $\bar{x}_s$ (clause~(iv) of
Theorem~\ref{3.1:thm-main}), we have $s-r(s)\to\infty$ as $s\to\infty$.

For \eqref{8.18:eq-volume-horizon-d}: $c_a\ge 9\log L+\log 2>0$, and $p_C(g)$ and
$\bar{V}(g,m)+\theta$ do not depend on $s$. Hence $c_a\bigl(s-r(s)+p_C(g)\bigr)\to\infty$, and the
left member exceeds $\bar{V}(g,m)+\theta$ for all sufficiently large $s$. So $\mathcal{E}$ is
bounded.

For \eqref{8.18:eq-volume-horizon-e}: $\underline{x}_u$ increases to $\infty$ as $u\to\infty$, and
$\mathfrak{P}(y)=A_0(\log y)^{p_0}$ increases to $\infty$ as $y\to\infty$. Therefore
$\mathfrak{P}(\underline{x}_{s-r(s)})\to\infty$ as $s\to\infty$. Since $\theta>0$, it follows that
$\tfrac12\theta\bigl(\mathfrak{P}(\underline{x}_{s-r(s)})-2\bigr)\to\infty$. The left member thus
exceeds the $s$-independent number $\bar{V}(g,m)$ for all sufficiently large $s$, and $\mathcal{E}$
is bounded.
\end{proof}

\begin{lemma}[Prolonging the run on the same lattice]\label{8.18:lem-run-prolongation}
Fix a family member $(g_0,g,m)$ and use the notation of Notation~\ref{8.18:not-members-runs}:
$x_0=g_0^{-2}$ is the bare datum and $x_k=g_k^{-2}$ are the running inverse couplings;
$K=K(g_0,g)$ is the first-passage level, so that $x_K\le g^{-2}+B^{\sharp}$;
$g_-(g)=(g^{-2}+B^{\sharp})^{-1/2}$. In the letters of Notation~\ref{8.18:not-members-runs}
the member $(g_0,g,m)$ is the member $(g_V,N_T)$, $g_V$ being the coupling which that
parametrisation attaches to it, with $N_T=2L^{a_V}$, $a_V\ge 0$, and
$K_V=K(x_0,g_V)$ its first-passage level in that parametrisation, its longest run stopping at
$K=K_V-m_{\mathbb{T}}(g_V)$ on a last lattice
of side $2L^{m}=N_T\,L^{m_{\mathbb{T}}(g_V)}$, so that $m=a_V+m_{\mathbb{T}}(g_V)$.
Attainable levels and the letters $K_*$ and $c_{\mathrm{gap}}$ are as in
Notation~\ref{8.18:not-members-runs}; an attainable level $k$ satisfies
$x_{k+1}<\min_{j\le k}x_j$. The measure
$\tilde{\nu}$ is the normalised measure on histories of the longest run of $(g_0,g,m)$.
An event is \emph{of levels $\le k$} if it is determined by the outcomes of the history
at levels $\le k$. Assume:
\begin{enumerate}
\item[(A)] $K_V<K_*$, and $K\ge k_0+r(s)$, where $k_0$ is the level at which the two-point
witness theorem reads its age events and $r(s)$ is its extraction depth; the second
inequality is the consequence of the ultraviolet and depth conditions of that theorem.
\item[(B)] (Re-targeting lemma, part (a).) For a member $(g',N_T')$ with
$K':=K(x_0,g')<K_*$, the level $K'$ is attainable, and for every
$k_1\in[K',K_*-c_{\mathrm{gap}}-1]$ there is an attainable level in
$[k_1,k_1+c_{\mathrm{gap}}]$.
\item[(C)] (Re-targeting lemma, part (b).) If moreover $k^{\diamond}\in[K',K_*[$ is
attainable and $L^{k^{\diamond}-K'}$ divides $N_T'$, then
$g^{\diamond}:=x_{k^{\diamond}+1}^{-1/2}$ and $N_T^{\diamond}:=N_T'L^{-(k^{\diamond}-K')}$
define a member $(g^{\diamond},N_T^{\diamond})$ over the same bare datum with
$K(x_0,g^{\diamond})=k^{\diamond}$, $g'\le g^{\diamond}<\gamma_*$,
$m_{\mathbb{T}}(g^{\diamond})\le m_{\mathbb{T}}(g')$, and its longest run stops at a level
$k^{\diamond}-m_{\mathbb{T}}(g^{\diamond})$ with
\begin{equation*}
k^{\diamond}-m_{\mathbb{T}}(g^{\diamond})\ \ge\ \bigl(K'-m_{\mathbb{T}}(g')\bigr)
+\bigl(k^{\diamond}-K'\bigr).
\end{equation*}
When $(g',N_T')=(g_V,N_T)$, this member is, in the parametrisation by $(g_0,g,m)$, the member
$(g_0,g^{\diamond},m-(k^{\diamond}-K_V))$ on the same bare lattice.
\item[(D)] (Runs across targets.) No object of step $k$ of a run reads the target
coupling, the cutoff or the torus exponent; hence two runs over the same bare datum
coincide through the last step of the shorter one, and for an event of levels $\le$ that
step the normalised measures of the two runs give it the same value.
\item[(E)] (Runs within one member.) Within one member, the normalised measure of an event
of levels $\le$ the stopping level does not depend on the run that represents it.
\item[(F)] (Fixed-torus lemma, part (1).) For every attainable $k^{\diamond}$ with
$p^{\diamond}:=k^{\diamond}-K_V\le a_V$, one has
$m-p^{\diamond}\ge m_{\mathbb{T}}(g_-(g^{\diamond}))$, $g^{\diamond}$ as in \textup{(C)}.
\end{enumerate}
Put $\hat{\kappa}:=\min\{K+m,\,K_*-1\}$, so that
$\hat{\kappa}=\min\{K_V+a_V,\,K_*-1\}$. Then:
\begin{enumerate}
\item[(i)] There are attainable levels in $[K_V,\hat{\kappa}]$; the largest, $k^{\dagger}$,
satisfies $\hat{\kappa}-k^{\dagger}\le c_{\mathrm{gap}}$.
\item[(ii)] With $g^{\dagger}:=x_{k^{\dagger}+1}^{-1/2}$ and
$N_T^{\dagger}:=2L^{K+m-k^{\dagger}}$, $(g^{\dagger},N_T^{\dagger})$ is a member over the
same bare datum, with $K(x_0,g^{\dagger})=k^{\dagger}$, whose longest run, denoted
$(\dagger)$, stops at a level $K^{\dagger\prime}$ with
\begin{equation*}
K^{\dagger\prime}:=k^{\dagger}-m_{\mathbb{T}}(g^{\dagger})\ \ge\ K\ \ge\ k_0+r(s),
\end{equation*}
on a last lattice $\Lambda^{\dagger}$ of side $N^{\dagger}$, hence with $(N^{\dagger})^4$
sites, where
\begin{equation}\label{8.18:eq-run-prolongation-a}
N^{\dagger}:=N_T^{\dagger}\,L^{m_{\mathbb{T}}(g^{\dagger})}
=2\,L^{K+m-k^{\dagger}}\,L^{m_{\mathbb{T}}(g^{\dagger})} .
\end{equation}
\item[(iii)] With $p:=k^{\dagger}-K_V$: in the parametrisation by $(g_0,g,m)$ the
prolonged member is $(g_0,g^{\dagger},m-p)$ on the same bare lattice of
$2L^{m+K}=2L^{(m-p)+(K+p)}$ sites per side, that is, with torus exponent $m-p$ and
cutoff $K+p$; the lattice of level $K+p$ has $2L^{m-p}$ sites
per side, and $m-p\ge m_{\mathbb{T}}(g_-(g^{\dagger}))$. The run $(\dagger)$ stops at the
level and on the side
\begin{equation*}
\begin{gathered}
K^{\dagger\prime}=K+p+m_{\mathbb{T}}(g_V)-m_{\mathbb{T}}(g^{\dagger})\ \ge\ K+p,\\
N^{\dagger}=2L^{m-p}\,L^{m_{\mathbb{T}}(g^{\dagger})-m_{\mathbb{T}}(g_V)}\ \le\ 2L^{m-p},
\end{gathered}
\end{equation*}
with equality in both places exactly when $m_{\mathbb{T}}(g^{\dagger})=m_{\mathbb{T}}(g_V)$;
thus $(\dagger)$ passes level $K+p$ and stops
$m_{\mathbb{T}}(g_V)-m_{\mathbb{T}}(g^{\dagger})\ge0$ levels after it.
Relative to $(g_0,g,m)$ the bare datum and the bare lattice are unchanged; only the
stopping level moves.
\item[(iv)] With $X^{\dagger}:=x_{k^{\dagger}+1}$,
\begin{equation}\label{8.18:eq-run-prolongation-b}
X^{\dagger}<\min_{j\le k^{\dagger}}x_j\le x_K\le g^{-2}+B^{\sharp},
\end{equation}
so $g^{\dagger}>g_-(g)$, $m_{\mathbb{T}}(g^{\dagger})\le m_{\mathbb{T}}(g_-(g))$,
$\mathsf{B}(g^{\dagger})\le C_{\mathsf{B}}\bigl(1+\log(g^{-2}+B^{\sharp})\bigr)$, and
$x_j>X^{\dagger}$ for every $j\le k^{\dagger}$; in particular
$\mathfrak{P}(x_j)>\mathfrak{P}(X^{\dagger})$ for every $j\le K$.
\item[(v)] With $\tilde{\nu}^{\dagger}$ the normalised measure of the run $(\dagger)$, for
every event $E$ of levels $\le K$,
\begin{equation}\label{8.18:eq-run-prolongation-c}
\tilde{\nu}(E)=\tilde{\nu}^{\dagger}(E);
\end{equation}
in particular for every event entering the two-point witness theorem and the Gaussian
extraction, these being of levels $\le k_0+r(s)\le K$.
\end{enumerate}
\end{lemma}

\begin{proof}
The two expressions for $\hat{\kappa}$ agree because
\begin{equation*}
K+m=\bigl(K_V-m_{\mathbb{T}}(g_V)\bigr)+\bigl(a_V+m_{\mathbb{T}}(g_V)\bigr)=K_V+a_V .
\end{equation*}

(i) By (A), $K_V\le K_*-1$, and $a_V\ge0$, so $\hat{\kappa}\ge K_V$. By (B), which applies
to $(g_V,N_T)$ since $K_V<K_*$, the level $K_V$ is attainable; hence
attainable levels in $[K_V,\hat{\kappa}]$ exist and $k^{\dagger}$ is well defined. Put
$k_1:=\max\{K_V,\hat{\kappa}-c_{\mathrm{gap}}\}$. If $\hat{\kappa}-c_{\mathrm{gap}}<K_V$,
then $k^{\dagger}\ge K_V>\hat{\kappa}-c_{\mathrm{gap}}$, that is
$\hat{\kappa}-k^{\dagger}<c_{\mathrm{gap}}$. Otherwise $k_1=\hat{\kappa}-c_{\mathrm{gap}}$,
and $k_1\in[K_V,K_*-c_{\mathrm{gap}}-1]$ because $\hat{\kappa}\le K_*-1$; by (B) there is an
attainable level in
\begin{equation*}
[k_1,k_1+c_{\mathrm{gap}}]=[\hat{\kappa}-c_{\mathrm{gap}},\hat{\kappa}]
\subset[K_V,\hat{\kappa}],
\end{equation*}
so $k^{\dagger}\ge\hat{\kappa}-c_{\mathrm{gap}}$.

(ii) Since $k^{\dagger}-K_V\le\hat{\kappa}-K_V\le a_V$, the power $L^{k^{\dagger}-K_V}$
divides $N_T=2L^{a_V}$, the quotient being $2L^{a_V-(k^{\dagger}-K_V)}$, again twice a
power of the odd integer $L$. As
$k^{\dagger}\le\hat{\kappa}\le K_*-1$, the level $k^{\dagger}\in[K_V,K_*[$ is attainable,
and (C) with $(g',N_T')=(g_V,N_T)$, $k^{\diamond}=k^{\dagger}$ gives the member
$(g^{\dagger},N_T^{\dagger})$ over the same bare datum, with $K(x_0,g^{\dagger})=k^{\dagger}$
and
\begin{equation*}
N_T^{\dagger}=2L^{a_V-k^{\dagger}+K_V}=2L^{K+m-k^{\dagger}} .
\end{equation*}
By (C) its longest run stops at
\begin{equation*}
\begin{split}
K^{\dagger\prime}=k^{\dagger}-m_{\mathbb{T}}(g^{\dagger})
&\ge\bigl(K_V-m_{\mathbb{T}}(g_V)\bigr)+\bigl(k^{\dagger}-K_V\bigr)\\
&\ge K_V-m_{\mathbb{T}}(g_V)=K,
\end{split}
\end{equation*}
and $K\ge k_0+r(s)$ by (A). The last lattice of that run has side
$N_T^{\dagger}L^{m_{\mathbb{T}}(g^{\dagger})}$, which is
\eqref{8.18:eq-run-prolongation-a}.

(iii) Put $p:=k^{\dagger}-K_V$; then $0\le p\le a_V$ by (ii), and $k^{\dagger}$ is
attainable. By (C), applied as in (ii), the member $(g^{\dagger},N_T^{\dagger})$ is, in the
parametrisation by $(g_0,g,m)$, the member $(g_0,g^{\dagger},m-p)$ on the same bare lattice,
whose side is $2L^{m+K}=2L^{(m-p)+(K+p)}$. In that parametrisation the bare lattice of a
member has side twice $L$ raised to the sum of the member's torus exponent and its cutoff
(Notation~\ref{8.18:not-members-runs}); as the torus exponent of the prolonged member is
$m-p$, its cutoff in this parametrisation is $K+p$. Each step of a run divides the side of
the lattice by the blocking factor $L$, so the lattice of level $k$ has $2L^{m+K-k}$ sites
per side; at level $K+p$ this is $2L^{m-p}$. For the run $(\dagger)$ we use
$K+m=K_V+a_V$ and $m=a_V+m_{\mathbb{T}}(g_V)$:
\begin{equation*}
K+m-k^{\dagger}=a_V-p=m-p-m_{\mathbb{T}}(g_V),\qquad
K^{\dagger\prime}=K_V+p-m_{\mathbb{T}}(g^{\dagger})
=K+p+m_{\mathbb{T}}(g_V)-m_{\mathbb{T}}(g^{\dagger}),
\end{equation*}
and the first identity inserted in \eqref{8.18:eq-run-prolongation-a} gives
$N^{\dagger}=2L^{m-p}L^{m_{\mathbb{T}}(g^{\dagger})-m_{\mathbb{T}}(g_V)}$. By (C),
$m_{\mathbb{T}}(g^{\dagger})\le m_{\mathbb{T}}(g_V)$, whence $K^{\dagger\prime}\ge K+p$ and
$N^{\dagger}\le 2L^{m-p}$, with equality in both exactly when
$m_{\mathbb{T}}(g^{\dagger})=m_{\mathbb{T}}(g_V)$; in all cases
$N^{\dagger}L^{K^{\dagger\prime}}=2L^{m+K}$, the side of the common bare lattice, so the
run $(\dagger)$ passes level $K+p$ and stops
$K^{\dagger\prime}-(K+p)=m_{\mathbb{T}}(g_V)-m_{\mathbb{T}}(g^{\dagger})$ levels after it.
Since $p\le a_V$ and $k^{\dagger}$ is
attainable, (F) gives $m-p\ge m_{\mathbb{T}}(g_-(g^{\dagger}))$. Finally, the bare datum
and the bare lattice are those of $(g_0,g,m)$, and only the stopping level has changed,
from $K$ to $K^{\dagger\prime}$.

(iv) Since $k^{\dagger}$ is attainable, $X^{\dagger}=x_{k^{\dagger}+1}<x_j$ for every
$j\le k^{\dagger}$. Since $m_{\mathbb{T}}\ge0$, (ii) gives $K\le K^{\dagger\prime}\le
k^{\dagger}$, so $\min_{j\le k^{\dagger}}x_j\le x_K$, and $x_K\le g^{-2}+B^{\sharp}$ by
first passage; this is \eqref{8.18:eq-run-prolongation-b}. Hence
\begin{equation*}
g^{\dagger}=(X^{\dagger})^{-1/2}>(g^{-2}+B^{\sharp})^{-1/2}=g_-(g).
\end{equation*}
The function
$m_{\mathbb{T}}(g')=11+m'+\log_L R(g')$ depends on $g'$ only through the least power of $L$
not below a power of $\log g'^{-2}$, which does not increase when $g'$ increases; so
$m_{\mathbb{T}}(g^{\dagger})\le m_{\mathbb{T}}(g_-(g))$. Since
$(g^{\dagger})^{-2}=X^{\dagger}<g^{-2}+B^{\sharp}$, the bound of $\mathsf{B}(g')$ by
$C_{\mathsf{B}}(1+\log g'^{-2})$ of Notation~\ref{8.18:not-members-runs} gives
$\mathsf{B}(g^{\dagger})\le C_{\mathsf{B}}(1+\log(g^{-2}+B^{\sharp}))$. For $j\le K\le
k^{\dagger}$ we have $x_j>X^{\dagger}$, and $\mathfrak{P}(y)=A_0(\log y)^{p_0}$ is
increasing, so $\mathfrak{P}(x_j)>\mathfrak{P}(X^{\dagger})$. No monotonicity of
$k\mapsto x_k$ has been used.

(v) The longest run of $(g_0,g,m)$ and the run $(\dagger)$ are runs over the same bare
datum $x_0$, stopping at $K$ and at $K^{\dagger\prime}\ge K$ respectively. By (D) they
coincide through step $K$, and every event $E$ of levels $\le K$ receives the same value
under their normalised measures; by (E) these values are $\tilde{\nu}(E)$ and
$\tilde{\nu}^{\dagger}(E)$, whichever run represents each member. This is
\eqref{8.18:eq-run-prolongation-c}. The events entering the two-point witness theorem and
the Gaussian extraction are of levels $\le k_0+r(s)\le K$, so
\eqref{8.18:eq-run-prolongation-c} applies to each of them.
\end{proof}

\begin{lemma}[Tori below and beyond the volume horizon]\label{8.18:lem-horizon-cases}
Fix a family member $(g_0,g,m)$ as in Notation~\ref{8.18:not-members-runs}, with its first-passage level $K$. Let $X_*$, $b_+$, $c_2$, $\lambda$ be the family constants entering \eqref{8.18:eq-volume-horizon-a}, and let $K_*$, $\gamma_*$, $c_{\mathrm{gap}}$, $K_V$, $g_V$ be the letters of the reference member fixed in Notation~\ref{8.18:not-members-runs}. Put $\hat{\kappa}:=\min\{K+m,K_*-1\}$. Whenever the prolonged member of Lemma~\ref{8.18:lem-run-prolongation} is defined, let $k^{\dagger}$, $g^{\dagger}$, $N_T^{\dagger}$, $N^{\dagger}$ and $X^{\dagger}$ be as there. Let $m_V(g)$, $\mathsf{B}(g')$ and $\bar{V}(g,m)$ be defined by \eqref{8.18:eq-volume-horizon-a}, \eqref{8.18:eq-volume-horizon-b} and \eqref{8.18:eq-volume-horizon-c}. Assume:
\begin{itemize}
\item[(A)] part~(a) of the lemma on attainable levels and run lengths of the reference member;
\item[(C)] part~(c) of the same lemma;
\item[(R)] the two-sided pinning \eqref{8.18:eq-members-runs-b}.
\end{itemize}
Then the following hold.
\begin{itemize}
\item[(i)] \emph{Tori used up before the coupling reaches $\gamma_*$.} Suppose $K+m\le K_*-1$. Then $\hat{\kappa}=K+m\ge K_V$, so the prolonged member is defined, and $N_T^{\dagger}\le 2L^{c_{\mathrm{gap}}}$. Moreover
\begin{equation}\label{8.18:eq-horizon-cases-a}
\begin{aligned}
(N^{\dagger})^4
&\le 16\,L^{4c_{\mathrm{gap}}}\,L^{4m_{\mathbb{T}}(g^{\dagger})}
 =16\,L^{4c_{\mathrm{gap}}+44}M^4R(g^{\dagger})^4\\
&\le 16\,L^{4c_{\mathrm{gap}}+48}M^4(\log X^{\dagger})^{4r}.
\end{aligned}
\end{equation}
The right-hand side is a polylogarithm in $X^{\dagger}$ of degree $4r$ whose coefficient $16L^{4c_{\mathrm{gap}}+48}M^4$ does not involve $m$; in particular no power $L^{m}$ occurs in it.
\item[(ii)] \emph{A sufficient condition not depending on $K$.} One has
\begin{equation*}
K_*-K\;\ge\;\frac{g^{-2}-X_*-b_+-2c_2}{3\lambda}\;\ge\;m_V(g)+1,
\end{equation*}
and consequently
\begin{equation}\label{8.18:eq-horizon-cases-b}
m\le m_V(g)\quad\Longrightarrow\quad K+m\le K_*-1 .
\end{equation}
\item[(iii)] \emph{Tori beyond the horizon.} Suppose $K+m\ge K_*$, and suppose $g$ satisfies the one-sided ceiling
\begin{equation}\label{8.18:eq-horizon-cases-d}
g^{-2}-X_*-b_+-2c_2\;\ge\;3\lambda\bigl(m_{\mathbb{T}}(g_-(g))+2c_{\mathrm{gap}}+2\bigr).
\end{equation}
Then $K_*-K\ge m_{\mathbb{T}}(g_V)+2c_{\mathrm{gap}}+2$ and $\hat{\kappa}=K_*-1\ge K_V$, so the prolonged member is defined, with $k^{\dagger}\ge K_*-1-c_{\mathrm{gap}}$ and
\begin{equation*}
N_T^{\dagger}=2L^{K+m-k^{\dagger}}\le 2L^{m-(K_*-K)+1+c_{\mathrm{gap}}}\le 2L^{m-m_V(g)+c_{\mathrm{gap}}}.
\end{equation*}
Moreover
\begin{equation}\label{8.18:eq-horizon-cases-c}
\mathsf{B}(g^{\dagger})\,(N^{\dagger})^4\;\le\;\bar{V}(g,m).
\end{equation}
The right-hand side depends only on $(g,m)$ and on the family constants. It does not depend on $K$, on $g_0$ or on position, and it grows like $L^{4(m-m_V(g))}$.
\end{itemize}
\end{lemma}

\begin{proof}
\emph{Part (ii).} The member and the reference member fixed in Notation~\ref{8.18:not-members-runs} satisfy $x_K>g^{-2}$ and $x_{K_*}\le X_*+b_+$. Applying (C) with these two bounds gives
\begin{equation*}
K_*-K\ge \frac{g^{-2}-X_*-b_+-2c_2}{3\lambda}.
\end{equation*}
By \eqref{8.18:eq-volume-horizon-a}, $m_V(g)$ is the integer part of the right-hand side minus one. The integer part of a number does not exceed the number, so the right-hand side is at least $m_V(g)+1$. If $m\le m_V(g)$, then
\begin{equation*}
K+m\le K+m_V(g)\le K_*-1,
\end{equation*}
which is \eqref{8.18:eq-horizon-cases-b}.

\emph{Part (i).} Since $K+m\le K_*-1$, the minimum defining $\hat{\kappa}$ is $K+m$. By hypothesis (H5) of Notation~\ref{2.2:not-hypotheses}, $m\ge m_{\mathbb{T}}(g_-(g))$; since $K_V=K+m_{\mathbb{T}}(g_V)$ and $m_{\mathbb{T}}(g_V)\le m_{\mathbb{T}}(g_-(g))$, it follows that $\hat{\kappa}=K+m\ge K_V$. So the prolonged member of Lemma~\ref{8.18:lem-run-prolongation} is defined in this case without any ceiling on $g$.

That lemma, whose level bound rests on (A), gives $\hat{\kappa}-k^{\dagger}\le c_{\mathrm{gap}}$ and $N_T^{\dagger}=2L^{K+m-k^{\dagger}}$. Hence
\begin{equation*}
N_T^{\dagger}=2L^{\hat{\kappa}-k^{\dagger}}\le 2L^{c_{\mathrm{gap}}}.
\end{equation*}
Since the side of the last lattice satisfies $N^{\dagger}\le N_T^{\dagger}L^{m_{\mathbb{T}}(g^{\dagger})}$ (Lemma~\ref{8.18:lem-run-prolongation}), this bound yields $(N^{\dagger})^4\le 16L^{4c_{\mathrm{gap}}}L^{4m_{\mathbb{T}}(g^{\dagger})}$, the first inequality of \eqref{8.18:eq-horizon-cases-a}.

For the equality, recall from Notation~\ref{8.18:not-members-runs} that
\begin{equation*}
m_{\mathbb{T}}(g^{\dagger})=11+m'+\log_L R(g^{\dagger}), \qquad M=L^{m'}.
\end{equation*}
Therefore $L^{4m_{\mathbb{T}}(g^{\dagger})}=L^{44}M^4R(g^{\dagger})^4$.

For the last inequality, $R(g^{\dagger})$ is the least power of $L$ that is at least $(\log (g^{\dagger})^{-2})^{r}$. Hence $R(g^{\dagger})/L<(\log (g^{\dagger})^{-2})^{r}$, that is, $R(g^{\dagger})^4\le L^4(\log (g^{\dagger})^{-2})^{4r}$. With $(g^{\dagger})^{-2}$ bounded by $X^{\dagger}$ of Lemma~\ref{8.18:lem-run-prolongation}, this gives the last inequality of \eqref{8.18:eq-horizon-cases-a}. No letter in it involves $m$.

\emph{Part (iii).} Combine the first inequality of part (ii) with \eqref{8.18:eq-horizon-cases-d}. This gives
\begin{equation*}
K_*-K\ge m_{\mathbb{T}}(g_-(g))+2c_{\mathrm{gap}}+2,
\end{equation*}
whence, since $m_{\mathbb{T}}(g_V)\le m_{\mathbb{T}}(g_-(g))$, $K_*-K\ge m_{\mathbb{T}}(g_V)+2c_{\mathrm{gap}}+2$. Since $K+m\ge K_*$, the minimum defining $\hat{\kappa}$ is $K_*-1$.

Since $K_V=K+m_{\mathbb{T}}(g_V)$, the requirement $\hat{\kappa}\ge K_V$ of the prolongation reads $K_*-1\ge K+m_{\mathbb{T}}(g_V)$. It holds by the last bound. Moreover the level
\begin{equation*}
k_1:=\max\{K_V,\hat{\kappa}-c_{\mathrm{gap}}\}
\end{equation*}
satisfies $k_1\le K_*-c_{\mathrm{gap}}-1$, so (A) applies at $k_1$. Thus Lemma~\ref{8.18:lem-run-prolongation} provides the prolonged member, with
\begin{equation*}
k^{\dagger}\ge \hat{\kappa}-c_{\mathrm{gap}}=K_*-1-c_{\mathrm{gap}}.
\end{equation*}
Using $K_*-K\ge m_V(g)+1$ from part (ii), we obtain
\begin{equation*}
\begin{aligned}
N_T^{\dagger}=2L^{K+m-k^{\dagger}}
&\le 2L^{m-(K_*-K)+1+c_{\mathrm{gap}}}\\
&\le 2L^{m-m_V(g)+c_{\mathrm{gap}}}.
\end{aligned}
\end{equation*}
Combining this with \eqref{8.18:eq-run-prolongation-b} gives \eqref{8.18:eq-horizon-cases-c}, with $\bar{V}(g,m)$ as in \eqref{8.18:eq-volume-horizon-c}.

By its definition, $\bar{V}(g,m)$ depends on $g$, $m$ and the family constants only. It does not involve $K$, $g_0$ or position, and it contains the factor $L^{4(m-m_V(g))}$.
\end{proof}

\begin{remark}
The floor of (H5) is
\begin{equation*}
m_{\mathbb{T}}(g_-(g))=11+m'+\log_L R(g_-(g))=O(\log_L\log g^{-2}).
\end{equation*}
Hence the range $m_{\mathbb{T}}(g_-(g))\le m\le m_V(g)$ is non-empty under a one-sided ceiling on $g$, in which a power of $g^{-2}$ is compared with an iterated logarithm. Non-emptiness of this range is used nowhere. If the range is empty, a statement asserted uniformly over it holds vacuously.

The ceiling \eqref{8.18:eq-horizon-cases-d} is likewise a one-sided ceiling on $g$ of this kind, chosen after $\gamma_*$. It is not needed in part~(i) of Lemma~\ref{8.18:lem-horizon-cases}.

In physical units, write $\ell_g$ for the unit length at coupling $g$. Let $\ell_*:=L^{K_*-K}\ell_g$ be the length at which the flow started at the bare scale first reaches $\gamma_*$. If $m\le m_V(g)$, the antecedent of \eqref{8.18:eq-horizon-cases-b}, then
\begin{equation*}
2L^m\ell_g\le 2L^{m_V(g)}\ell_g\le L^{K_*-K}\ell_g=\ell_*,
\end{equation*}
where the last inequality uses $K_*-K\ge m_V(g)+1$ from part (ii) of Lemma~\ref{8.18:lem-horizon-cases} and $L\ge 2$. Under the change-of-units identity along $g_0\downarrow 0$ at fixed $g$, $\ell_g$ is the fixed unit. By hypothesis (R) of Lemma~\ref{8.18:lem-horizon-cases}, that is by \eqref{8.18:eq-members-runs-b}, $K_*-K$ is bounded above and below in terms of $(g,\gamma_*)$ alone; this is the only place where hypothesis (R) enters. So part~(i) of Lemma~\ref{8.18:lem-horizon-cases} concerns tori smaller than a fixed physical length.
\end{remark}

\begin{remark}[The stability quotient read at one lattice]\label{8.18:rem-stability-quotient}
This remark explains what the last lattice of the prolonged member is needed for. The conversion
of un-normalised weights into probabilities used here is the quotient of the two global stability
bounds of \cite[Theorem~1, (0.1)]{B-CMP122b}, both read on one lattice $\Lambda$ with $\#\Lambda$
sites.

The first is the \emph{upper global stability bound}. The sum
$\Sigma=\sum_h\mathfrak{M}^{(\theta)}(h)$ of the stripped majorants over the histories $h$ on
$\Lambda$ satisfies
\begin{equation}\label{8.18:eq-stability-quotient-1}
\Sigma\le e^{E_+^{(\theta)}\,\#\Lambda}.
\end{equation}
This is the bound on sums of stripped majorants, which has the shape of the upper bound of clause
(i-b) of Theorem~\ref{3.1:thm-main}; its exponent $E_+^{(\theta)}$ does not depend on $K$ and is
$O(1+\log g'^{-2})$. Here and below $\mathsf{B}=\mathsf{B}(g')$,
$E_+^{(\theta)}=E_+^{(\theta)}(g')$ and $C_{\mathrm{low}}=C_{\mathrm{low}}(g')$ are evaluated at
the terminal coupling $g'$ of the member whose lattice is read, as in
Definition~\ref{8.18:def-volume-horizon}; this is the shape in which the upper bound enters the
bound \eqref{8.18:eq-volume-horizon-b} on $\mathsf{B}$.

The second is the \emph{lower global stability bound}. The normalising integral $Z=\int\rho_k$ of
the effective density at the level $k$ whose lattice is $\Lambda$ satisfies
\begin{equation}\label{8.18:eq-stability-quotient-2}
Z\ge e^{-C_{\mathrm{low}}\,\#\Lambda};
\end{equation}
at the last level $K^{\dagger\prime}$ this is the bound on $Z=\int\rho_{K^{\dagger\prime}}$.
This is the lower bound of clause (i-b) of Theorem~\ref{3.1:thm-main} combined with a lower bound
on the Haar measure of the ball of small-field configurations, as recorded in the lower part of the
bound \eqref{8.18:eq-volume-horizon-b}.

Both bounds must be read at the same lattice for their quotient to be
$e^{\mathsf{B}\,\#\Lambda}$, where $\mathsf{B}=E_+^{(\theta)}+C_{\mathrm{low}}$ as in
\eqref{8.18:eq-volume-horizon-b}. The coarsest lattice that the prolonged run of
Lemma~\ref{8.18:lem-run-prolongation} reaches over the given bare coupling and bare lattice is its
last lattice $\Lambda^{\dagger}$, at level $K^{\dagger\prime}$, with $(N^{\dagger})^4$ sites.

At an earlier level $k$, write $N_k$ for the number of sites per side of the level-$k$ lattice;
the quotient read there is $e^{\mathsf{B}N_k^4}$, where
\begin{equation}\label{8.18:eq-stability-quotient-3}
N_k^4=(N^{\dagger})^4L^{4(K^{\dagger\prime}-k)} .
\end{equation}
In the exponent this exceeds the quotient at $\Lambda^{\dagger}$ by the factor
$L^{4(K^{\dagger\prime}-k)}$. Prolonging the run, as in
Lemma~\ref{8.18:lem-run-prolongation}, changes the terminal volume, and this is its purpose.
\end{remark}

\begin{proof}
\emph{The quotient.} The bound \eqref{8.18:eq-stability-quotient-2} gives $Z>0$. Dividing
\eqref{8.18:eq-stability-quotient-1} by \eqref{8.18:eq-stability-quotient-2}, both read on the
same lattice $\Lambda$, we obtain
\begin{equation*}
\frac{\Sigma}{Z}\le\frac{e^{E_+^{(\theta)}\#\Lambda}}{e^{-C_{\mathrm{low}}\#\Lambda}}
=e^{(E_+^{(\theta)}+C_{\mathrm{low}})\#\Lambda}=e^{\mathsf{B}\,\#\Lambda}.
\end{equation*}
The last equality is the definition of $\mathsf{B}$ in \eqref{8.18:eq-volume-horizon-b}.

If the two bounds were read on lattices with different numbers of sites $n_1\ne n_2$, the division
would give $e^{E_+^{(\theta)}n_1+C_{\mathrm{low}}n_2}$, which is not the quotient
$e^{\mathsf{B}\,\#\Lambda}$ of a single lattice.

\emph{The count of sites.} By Definition~\ref{1.5:def-block-average}, the block lattice at level
$k+1$ has spacing $L$ times that at level $k$ on the same torus. Hence the side lengths satisfy
$N_k=L\,N_{k+1}$. At the last level we have $N_{K^{\dagger\prime}}=N^{\dagger}$. Descending from
$K^{\dagger\prime}$ to $k$, induction gives $N_k=N^{\dagger}L^{K^{\dagger\prime}-k}$, and
raising to the fourth power gives \eqref{8.18:eq-stability-quotient-3}.

The lattice at level $k$ therefore has $N_k^4$ sites. Taking it as $\Lambda$ in the first step,
the lower bound read there gives only $Z\ge e^{-C_{\mathrm{low}}N_k^4}$, and
the quotient read there is $e^{\mathsf{B}N_k^4}$. Its exponent equals
$\mathsf{B}(N^{\dagger})^4L^{4(K^{\dagger\prime}-k)}$, which is $L^{4(K^{\dagger\prime}-k)}$
times the exponent $\mathsf{B}(N^{\dagger})^4$ of the quotient at $\Lambda^{\dagger}$.
\end{proof}

\begin{lemma}[Un-normalised age-and-reserve bound at the last lattice]\label{8.18:lem-age-reserve-bound}
Let $c_a\ge\log 2$ be a free number and put $q:=e^{-c_a}$, so that $0<q\le\tfrac12$. Let
$\mathsf{c}$ be the enlarged entropy constant of geometric forgetting of old regions, and let $\theta$
satisfy
\begin{equation}\label{8.18:eq-age-reserve-bound-1}
0<\theta<\tfrac12\,(1-3\beta_0),
\end{equation}
with $\beta_0$ the reserve fraction of the flat reserve bounds. Assume that $\gamma\le\gamma_L(c_a)$,
and that $\gamma$ is also at most the ceiling under which the flat reserve bounds hold.

Consider the longest run, over the bare datum, of any member in the sense of
Notation~\ref{8.18:not-members-runs}. Let $\hat{K}$ be its last level, let $\hat{N}^4$ be the number
of sites of its last lattice, let $\Omega$ be its set of histories at level $\hat{K}$ and $\tau_{\omega}$
the weight of $\omega\in\Omega$. Recall that $\mathfrak{P}(y)=A_0(\log y)^{p_0}$, so that
$p_0(g_k)=\mathfrak{P}(x_k)$ along the run. For a live component $X$ at level $k$, let $j(X)$ be the
level at which the oldest seed of $X$ was created; its age at level $k$ is $k-j(X)$. For a
terminal-live component $X$ of $\omega$, that is, a component still live at level $\hat{K}$, write
$a_X:=\hat{K}-j(X)$. For a history $\omega$ put
\begin{equation}\label{8.18:eq-age-reserve-bound-2}
w_{\theta}(\omega):=\exp\Bigl(-\theta\sum_{Z\in\mathcal{Z}^{\mathrm{live}}(\omega)}p_0(g_{j(Z)})\Bigr),
\end{equation}
where $\mathcal{Z}^{\mathrm{live}}(\omega)$ is the set of seeds of the terminal-live components of
$\omega$ and, for a seed $Z$, $j(Z)$ is the level at which $Z$ was created. Assume the following.
\begin{itemize}
\item[(a)] \emph{Mechanism of geometric forgetting of old regions.} Ba{\l}aban's induction
\cite[(1.80)--(1.88)]{B-CMP122b}, with the unnamed constants multiplying
$M^{d}R_j^{d+1}d'_j(Z_j)$, and similar terms, taken at the value $\mathsf{c}$, holds.
The constant enters it at four places: after \cite[(1.82)]{B-CMP122b}; \cite[(1.83)]{B-CMP122b};
the preparatory case; and the merging case \cite[(1.85)--(1.88)]{B-CMP122b}. At each of these
places the requirement is a lower bound on the degree function $p_0(\cdot)$, met by taking $p_0$ large
and $\gamma\le\gamma_L(c_a)$. With $\kappa^{[\mathsf{c}]}_k(X)\ge0$ the exponent obeying
\cite[(1.80)]{B-CMP122b} at the constant $\mathsf{c}$, every live component $X$ of age $k-j(X)$ at
level $k$, in any history, carries in the resulting majorant a factor at most
\begin{equation*}
\exp\bigl(-\kappa^{[\mathsf{c}]}_k(X)-2p_0(g_{j(X)})\bigr)\,q^{\,k-j(X)+1}.
\end{equation*}
The combinatorics of \cite{B-CMP122b} use of this factor only the bound
$\exp(-\kappa^{[\mathsf{c}]}_k(X)-2p_0(g_{j(X)}))$.
\item[(b)] \emph{The flat reserve bounds on terminal-live seeds.} Run the same induction with the
coefficient $2-\theta$ in place of $2$ on the seed factor. The change enters at the same four places,
and it is met by taking $p_0$ large and $\gamma$ below the ceiling of the flat reserve bounds. The
induction then yields non-negative $\theta$-stripped majorants $\mathfrak{M}^{(\theta)}(\omega)$,
$\omega\in\Omega$ (that is, majorants of the weight with the reserve factor $w_{\theta}$ taken out),
with the following two properties.
\begin{itemize}
\item[(b1)] For every $\omega\in\Omega$ one has
$\tau_{\omega}\le w_{\theta}(\omega)\,\mathfrak{M}^{(\theta)}(\omega)$. This is the bound, term by term,
of the integral over the live components, and it holds uniformly in the exterior field.
\item[(b2)] One has
\begin{equation*}
\sum_{\omega\in\Omega}\mathfrak{M}^{(\theta)}(\omega)\le e^{E_+^{(\theta)}\hat{N}^4}.
\end{equation*}
The number $E_+^{(\theta)}$ is independent of $K$, of the side of the torus and of positions, and it
depends on the member only through its terminal inverse coupling. The combinatorial proof of this
inequality uses, of the factor carried by a terminal-live component $X$, only that it is at most
$\exp(-\kappa_{\hat{K}}(X)-(2-\theta)p_0(g_{j(X)}))$, where $\kappa_{\hat{K}}(X)\ge0$ is the exponent
obeying \cite[(1.80)]{B-CMP122b} at the admissible value of the constants at which the induction is
run.
\end{itemize}
\item[(c)] The inequality \eqref{8.18:eq-members-runs-a}: $\mathfrak{P}(x_j)\ge\mathfrak{P}(x_k)-1$
for $j\le k$ along the run.
\end{itemize}
Then for every integer $a\ge0$, every $b\le\hat{K}-a$, and every event
\begin{equation*}
E\subset\bigl\{\omega\in\Omega:\ \omega\text{ has at level }\hat{K}\text{ a terminal-live component }X
\text{ with }j(X)\le b\bigr\},
\end{equation*}
one has
\begin{align}
\sum_{\omega\in E}\tau_{\omega}
&\le e^{\theta}\,q^{\hat{K}-b}\,\exp\bigl(-\theta\,\mathfrak{P}(x_b)\bigr)\,
e^{E_+^{(\theta)}\hat{N}^4}\label{8.18:eq-age-reserve-bound-a}\\
&\le e^{\theta}\,q^{a}\,\exp\bigl(-\theta\,\mathfrak{P}(x_b)\bigr)\,
e^{E_+^{(\theta)}\hat{N}^4}.\label{8.18:eq-age-reserve-bound-b}
\end{align}
\end{lemma}

\begin{proof}
\emph{Step (i): one induction with both modifications.} We run the induction
\cite[(1.80)--(1.88)]{B-CMP122b} once, with the entropy constant $\mathsf{c}$ of hypothesis~(a) and
the seed coefficient $2-\theta$ of hypothesis~(b) at the same time.

Both modifications are of one kind. An absolute constant of the induction is taken at another
admissible value, and the change is paid for, as in \cite{B-CMP122b}, by taking $p_0$ large and
$\gamma$ small. Both modifications also enter at the same four places listed in hypothesis~(a). At
each place, carrying them out together requires the lower bound on the degree function $p_0(\cdot)$
demanded by each of them. That bound is secured by $p_0$ large together with $\gamma$ below both
ceilings, namely $\gamma\le\gamma_L(c_a)$ and $\gamma$ at most the ceiling of the flat reserve bounds.
This is assumed in the statement. Taking the two modifications together therefore changes nothing
except the ceiling on $\gamma$.

The resulting $\theta$-stripped majorant $\mathfrak{M}^{(\theta)}(\omega)$ carries, for each
terminal-live component $X$ of $\omega$, a factor at most
\begin{equation}\label{8.18:eq-age-reserve-bound-3}
\exp\bigl(-\kappa^{[\mathsf{c}]}_{\hat{K}}(X)-(2-\theta)p_0(g_{j(X)})\bigr)\,q^{\,a_X+1}.
\end{equation}
This follows from hypothesis~(a) applied at level $k=\hat{K}$, where the age of $X$ is
$a_X=\hat{K}-j(X)$, with the seed coefficient $2-\theta$. In the combined run the exponent of
\cite[(1.80)]{B-CMP122b} in hypothesis~(b2) is $\kappa^{[\mathsf{c}]}_{\hat{K}}(X)$. By the last
sentence of hypothesis~(b2), the combinatorics of the proof of the sum bound therefore use of this
factor only the bound
$\exp(-\kappa^{[\mathsf{c}]}_{\hat{K}}(X)-(2-\theta)p_0(g_{j(X)}))$.

\emph{Step (ii): a real weight $\mu\ge1$ on old components.} Put $\mu:=q^{-a}$. Since $0<q<1$ and
$a\ge0$, we have $\mu\ge1$. For $\omega\in\Omega$ let
\begin{equation*}
n(\omega):=\#\bigl\{X:\ X\text{ a terminal-live component of }\omega\text{ with }a_X\ge a\bigr\},
\qquad
\mathfrak{M}^{(\theta,\mu)}(\omega):=\mathfrak{M}^{(\theta)}(\omega)\,\mu^{n(\omega)}.
\end{equation*}
In $\mathfrak{M}^{(\theta,\mu)}(\omega)$ we attach one factor $\mu$ to each terminal-live component
$X$ with $a_X\ge a$. For such $X$, multiplying \eqref{8.18:eq-age-reserve-bound-3} by $\mu$ gives
\begin{align*}
\mu\,\exp\bigl(-\kappa^{[\mathsf{c}]}_{\hat{K}}(X)-(2-\theta)p_0(g_{j(X)})\bigr)\,q^{\,a_X+1}
&=\exp\bigl(-\kappa^{[\mathsf{c}]}_{\hat{K}}(X)-(2-\theta)p_0(g_{j(X)})\bigr)\,q^{\,a_X-a+1}\\
&\le\exp\bigl(-\kappa^{[\mathsf{c}]}_{\hat{K}}(X)-(2-\theta)p_0(g_{j(X)})\bigr)\,q,
\end{align*}
because $a_X-a\ge0$ and $0<q<1$. Since $q<1$, this is at most
$\exp(-\kappa^{[\mathsf{c}]}_{\hat{K}}(X)-(2-\theta)p_0(g_{j(X)}))$. Components with $a_X<a$, and
all other factors of the majorant, are unchanged.

Thus every per-component factor of the non-negative majorants $\mathfrak{M}^{(\theta,\mu)}$ stays
within the bound that, by Step~(i), is all the proof of the sum bound in hypothesis~(b2) uses. That
proof therefore applies without change to the majorants $\mathfrak{M}^{(\theta,\mu)}$, and gives
\begin{equation}\label{8.18:eq-age-reserve-bound-4}
\sum_{\omega\in\Omega}\mathfrak{M}^{(\theta,\mu)}(\omega)\le e^{E_+^{(\theta)}\hat{N}^4},
\end{equation}
with the same number $E_+^{(\theta)}$. This is the observation made in the corollary, at the last
lattice, of geometric forgetting of old regions: there each old factor is multiplied by a number
$\mu$ with $|\mu|q^{a+1}\le1$, and the global upper bound of the correlation theorem persists. Here
the same observation is applied to non-negative majorants and real $\mu\ge1$.

\emph{Step (iii): using non-negativity.} Let $\omega$ satisfy $n(\omega)\ge1$. Since
$\mathfrak{M}^{(\theta,\mu)}(\omega)\ge0$ and $\mu\ge1$, we have $\mu^{-n(\omega)}\le\mu^{-1}$, and so
\begin{equation*}
\mathfrak{M}^{(\theta)}(\omega)=\mu^{-n(\omega)}\,\mathfrak{M}^{(\theta,\mu)}(\omega)
\le\mu^{-1}\,\mathfrak{M}^{(\theta,\mu)}(\omega)=q^{a}\,\mathfrak{M}^{(\theta,\mu)}(\omega).
\end{equation*}
We sum over these $\omega$. We then enlarge the range of summation to all of $\Omega$, which is
allowed because $\mathfrak{M}^{(\theta,\mu)}\ge0$. Finally we apply
\eqref{8.18:eq-age-reserve-bound-4}. This gives
\begin{equation}\label{8.18:eq-age-reserve-bound-5}
\sum_{\omega:\,n(\omega)\ge1}\mathfrak{M}^{(\theta)}(\omega)
\le q^{a}\sum_{\omega\in\Omega}\mathfrak{M}^{(\theta,\mu)}(\omega)
\le q^{a}\,e^{E_+^{(\theta)}\hat{N}^4}.
\end{equation}

\emph{Step (iv): the reserve.} Let $\omega\in E$. Then $\omega$ has a terminal-live component $X$
with $j(X)\le b\le\hat{K}-a$. Hence $a_X=\hat{K}-j(X)\ge\hat{K}-b\ge a$, and therefore
$n(\omega)\ge1$.

The oldest seed $Z_0$ of $X$ was created at level $j(X)\le b$, and $Z_0$ belongs to
$\mathcal{Z}^{\mathrm{live}}(\omega)$. We keep in the sum of \eqref{8.18:eq-age-reserve-bound-2} only
the term of $Z_0$; the other terms are non-negative and $\theta>0$. This gives
\begin{equation*}
w_{\theta}(\omega)\le\exp\bigl(-\theta p_0(g_{j(X)})\bigr)=\exp\bigl(-\theta\,\mathfrak{P}(x_{j(X)})\bigr).
\end{equation*}

Next we apply hypothesis~(c), that is \eqref{8.18:eq-members-runs-a}, with $j=j(X)\le b$ in the
role of the smaller index and $b$ in the role of the larger. It gives
$\mathfrak{P}(x_{j(X)})\ge\mathfrak{P}(x_b)-1$, and since $\theta>0$,
\begin{equation}\label{8.18:eq-age-reserve-bound-6}
\sup_{\omega\in E}w_{\theta}(\omega)\le e^{\theta}\exp\bigl(-\theta\,\mathfrak{P}(x_b)\bigr).
\end{equation}

\emph{Step (v): assembly.} By hypothesis~(b1), applied to each $\omega\in E$, we have
$\tau_{\omega}\le w_{\theta}(\omega)\mathfrak{M}^{(\theta)}(\omega)$. We bound $w_{\theta}$ by its
supremum over $E$. We then use that $E\subset\{n\ge1\}$ by Step~(iv) and that
$\mathfrak{M}^{(\theta)}\ge0$. Finally we apply \eqref{8.18:eq-age-reserve-bound-6} and
\eqref{8.18:eq-age-reserve-bound-5}. This gives
\begin{align*}
\sum_{\omega\in E}\tau_{\omega}
&\le\sup_{\omega\in E}w_{\theta}(\omega)\sum_{\omega:\,n(\omega)\ge1}\mathfrak{M}^{(\theta)}(\omega)\\
&\le e^{\theta}\,e^{-\theta\mathfrak{P}(x_b)}\,q^{a}\,e^{E_+^{(\theta)}\hat{N}^4}.
\end{align*}
This is the right-hand side of \eqref{8.18:eq-age-reserve-bound-b}.

Steps~(ii) and~(iii) used of the integer $a$ only that $a\ge0$. Step~(iv) used of it only that
$a_X\ge a$ holds on $E$. Both properties hold for $\hat{K}-b$ in place of $a$. Indeed,
$\hat{K}-b\ge a\ge0$ because $b\le\hat{K}-a$. Moreover, on $E$ the component $X$ satisfies
$a_X=\hat{K}-j(X)\ge\hat{K}-b$. Running Steps~(ii)--(iii) with $a$ replaced by $\hat{K}-b$, and then
Steps~(iv)--(v) unchanged, gives \eqref{8.18:eq-age-reserve-bound-a}. Finally,
$\hat{K}-b\ge a$ and $0<q<1$ give $q^{\hat{K}-b}\le q^{a}$. So the right-hand side of
\eqref{8.18:eq-age-reserve-bound-a} is at most that of \eqref{8.18:eq-age-reserve-bound-b}, which
completes the chain.
\end{proof}

\begin{remark}[Remarks on the age-and-reserve bound]\label{8.18:rem-age-bound-remarks}
The following remarks concern the un-normalised age-and-reserve bound at the last lattice,
Lemma~\ref{8.18:lem-age-reserve-bound}. Throughout, $E$ is an event as in that lemma, $\mu$ is
its dial, $\theta$ and $\beta_0$ are its reserve fraction and the parameter entering the range
of $\theta$, and $q=e^{-c_a}$. Further, $a\ge0$ and $b\le\hat{K}-a$ are the age and birth-level
parameters of that lemma, $\hat{K}$ and $\hat{N}$ are its last level and the side of its last
lattice, and $x_b=g_b^{-2}$.
\begin{itemize}
\item[(a)] The bound of Lemma~\ref{8.18:lem-age-reserve-bound} carries no factor counting
positions. The second of the two reserve inequalities from which that lemma is derived
already sums over all the terms in question. Consequently the bound holds whether or not the
event $E$ fixes a cell, with no factor counting positions in either case.
\item[(b)] The proof of Lemma~\ref{8.18:lem-age-reserve-bound} uses no Schwarz lemma, and its
bound carries no factor $2$. The majorants entering it are non-negative and the dial $\mu$ is
real. The age-out corollary allows a complex dial ranging over a disc $\mathbb{D}$ in the
complex plane, as the setting of the one-lattice Lipschitz step requires.
Lemma~\ref{8.18:lem-age-reserve-bound} corresponds to the value $\delta=0$ of the parameter
$\delta$ of that setting; there a real dial suffices, and the complex dial is more than is
needed.
\item[(c)] In the bound of Lemma~\ref{8.18:lem-age-reserve-bound} nothing marks a seed that is
healed, and nothing marks the age of a healed large-field region. This is as the statement on
the rarity of young large-field blocks requires: a mark on the age or on the reserve of a
healed region would set the vacuum parts of the exponentiated $\mathbf{R}$-terms against the
counterterm, which carries no mark, and would leave at each level a contribution whose sum
over the levels is not bounded uniformly in $K$. The
parts of the sum in which an old region is healed at a later step do not belong to $E$. They
are estimated separately, by the healed old-region bound.
\item[(d)] The proposition on old regions stated with a single dial uses, for a component read
at level $k$ whose age there is $k-n$, the dial
\begin{equation*}
\bar{\mu}=q^{-(k-n)}\exp\bigl(2p_0^{\sharp}(n)-2(1+\beta_0)^{-1}p_0(g_k)\bigr),
\end{equation*}
where $p_0^{\sharp}(\cdot)$ is the exponent function of that proposition and the level $n$ is
not to be confused with the count $n(h)$. This dial carries both the age and the reserve
surplus in a single quantity. With $\bar{\mu}$ in
place of $\mu$, the bound of Lemma~\ref{8.18:lem-age-reserve-bound} follows from three
statements: geometric forgetting of old regions, the age-out corollary, and the surplus
provided by the proposition on old regions stated with a single dial, and the reserve
fraction $\theta$ is then replaced by $2\beta_0/(1+\beta_0)$.
This is the same lemma in a different formulation. In this book we keep the flat reserve
fraction $\theta$ of Lemma~\ref{8.18:lem-age-reserve-bound}.
\item[(e)] A second form of the bound is obtained from the age lemma. That lemma is
deterministic. For a live component $C$ at level $k$ it gives
\begin{equation*}
\operatorname{age}_k(C)\le\sum_{Z}\bigl(\check{R}_{j(Z)}+T_0\bigr)+(\text{size term}),
\end{equation*}
the sum running over the seeds $Z$ of $C$ and $T_0$ being the constant of the age lemma.
Thus a large age forces either large-field events or size, at a density linear in the age.
Combine the age lemma with the reserve taken over all seeds, in place of the age dial. The
result is
\begin{equation*}
\sum_{\omega\in E}\tau_{\omega}
\le\exp\bigl(-\max\{\theta\,\mathfrak{P}(x_b),\,\varrho_*\,a\}\bigr)\,e^{E_+\hat{N}^4},
\end{equation*}
with a rate $\varrho_*$ that satisfies
\begin{equation*}
\varrho_*\asymp\frac{\theta\,p_0(\cdot)}{L\check{R}_{\cdot}},\qquad \varrho_*\gg 9\log L ,
\end{equation*}
where $p_0(\cdot)$ is the degree function and $\check{R}_{\cdot}$ the running scale, the
argument of the one and the level of the other not being fixed in this remark.
The second relation holds under a ceiling on the reference coupling $\gamma$ of degree
$p_0-r>0$, in the exponents $p_0$ and $r$ of the glossary, imposed on
the range $y\ge X_*$ of the inverse coupling $y$. The estimates of this chapter that take
Lemma~\ref{8.18:lem-age-reserve-bound} as an input accept either form of the bound without
change; only the form of Lemma~\ref{8.18:lem-age-reserve-bound} is used below.
\end{itemize}
\end{remark}

\begin{lemma}[The fate of an old live region]\label{8.18:lem-region-fate}
Let $(g_0,g,m)$ be a member as in Notation~\ref{8.18:not-members-runs}, and let $(\dagger)$ be its
prolonged run as in Lemma~\ref{8.18:lem-run-prolongation}, with last level
$K^{\dagger\prime}\ge K$ and normalised history measure $\tilde{\nu}^{\dagger}$; $\tilde{\nu}$ is
the normalised history measure of the original run. Assume:
\begin{itemize}
\item[(a)] the exhaustive dichotomy for live components: at every later step of a run, a live
component either is carried into a live component of the next level, possibly merged with others,
with the merged component keeping the oldest index among those merged as in
\cite[(1.85)]{B-CMP122b}, or is removed by the operation $\mathbf{R}$ \cite{B-CMP119}, and no third
alternative occurs;
\item[(b)] the identity \eqref{8.18:eq-run-prolongation-c} of Lemma~\ref{8.18:lem-run-prolongation};
\item[(c)] the comparison of seed age and territory age: the birth index of the oldest seed of a
live component is at most the oldest layer of its territory, since a cell leaves the small-field
region only at or after the birth of the oldest seed of its component, and re-creation of a
region only makes its territory younger.
\end{itemize}
Fix a level $k_0\le K$, a level-$k_0$ cell $\square$, and an age $a>\mathfrak{l}$, where
$\mathfrak{l}$ is the age threshold. Put $E_a(\square):=\mathfrak{lv}^{=a}_{k_0}(\square)$, the event that
$\square$ lies in a live component $Z$ of age $a$ at level $k_0$, so that $j(Z)=k_0-a$. The
\emph{descendant} of $Z$ at a later level is the live component into which $Z$ is carried under
hypothesis (a), and the \emph{lineage} of $Z$ is the sequence of its descendants. Define
$\mathrm{HEAL}^{\dagger}$ as the event that, in the run $(\dagger)$, the lineage of $Z$ is removed by
$\mathbf{R}$ at some level $k_h$ with $k_0<k_h\le K^{\dagger\prime}$, and $\mathrm{END}^{\dagger}$ as
the event that it is live at $K^{\dagger\prime}$: there is a live component $X$ at level
$K^{\dagger\prime}$ that contains the descendant of $Z$ and has $j(X)\le j(Z)=k_0-a$, where $j(X)$
is the birth index of the oldest seed of $X$ and $j(Z)=k_0-a$ is the oldest layer of the territory
of $Z$. Then on
$E_a(\square)$ exactly one of $\mathrm{HEAL}^{\dagger}$, $\mathrm{END}^{\dagger}$ occurs; on
$E_a(\square)\cap\mathrm{END}^{\dagger}$,
\begin{equation}\label{8.18:eq-region-fate-1}
\operatorname{age}_{K^{\dagger\prime}}(X)\ge a+(K^{\dagger\prime}-k_0)\ge a ;
\end{equation}
and, writing
$q^{\bullet,=a}_{k_0}:=\max_{\square}\tilde{\nu}^{\dagger}(E_a(\square)\cap\bullet)$ for
$\bullet\in\{\mathrm{HEAL}^{\dagger},\mathrm{END}^{\dagger}\}$ (denoted
$q^{\mathrm{heal}\dagger,=a}_{k_0}$, $q^{\mathrm{END}\dagger,=a}_{k_0}$),
\begin{equation}\label{8.18:eq-region-fate-a}
q^{=a}_{k_0}\le q^{\mathrm{heal}\dagger,=a}_{k_0}+q^{\mathrm{END}\dagger,=a}_{k_0},
\end{equation}
where $q^{=a}_{k_0}$ denotes the one-cell live weight of age $a$ of the original member, that is,
$\max_{\square}\tilde{\nu}(E_a(\square))$.
\end{lemma}

\begin{proof}
Work on $E_a(\square)$ in the run $(\dagger)$, and let $Z\ni\square$ be the live component of age
$a$ at level $k_0$, so $j(Z)=k_0-a$. By hypothesis (a), applied at each of the levels
$k_0,k_0+1,\dots$, the lineage of $Z$ persists as a live component, possibly merged and then
carrying the oldest index, until $\mathbf{R}$ removes it. Either this removal happens at some level
$k_h$ with $k_0<k_h\le K^{\dagger\prime}$, which is $\mathrm{HEAL}^{\dagger}$, or it does not happen
up to the last level $K^{\dagger\prime}$. In the latter case the descendant of $Z$ lies in a live
component $X$ at level $K^{\dagger\prime}$. Since merging keeps the oldest index (hypothesis (a)),
the oldest seed of $X$ is born no later than the oldest seed of $Z$, and by hypothesis (c) the
birth index of the oldest seed of $Z$ is at most the oldest layer $k_0-a$ of the territory of $Z$;
hence $j(X)\le j(Z)=k_0-a$, and this is $\mathrm{END}^{\dagger}$. The two cases are exclusive, since a
lineage removed at $k_h\le K^{\dagger\prime}$ has no descendant live at $K^{\dagger\prime}$.

On $\mathrm{END}^{\dagger}$, as $\operatorname{age}_{K^{\dagger\prime}}(X)=K^{\dagger\prime}-j(X)$
and $K^{\dagger\prime}\ge K\ge k_0$,
\begin{equation*}
\operatorname{age}_{K^{\dagger\prime}}(X)\ge K^{\dagger\prime}-(k_0-a)
=a+(K^{\dagger\prime}-k_0)\ge a,
\end{equation*}
which is \eqref{8.18:eq-region-fate-1}.

Since $E_a(\square)\subset\mathrm{HEAL}^{\dagger}\cup\mathrm{END}^{\dagger}$, subadditivity gives
for every cell $\square$
\begin{equation*}
\tilde{\nu}^{\dagger}(E_a(\square))
\le\tilde{\nu}^{\dagger}(E_a(\square)\cap\mathrm{HEAL}^{\dagger})
+\tilde{\nu}^{\dagger}(E_a(\square)\cap\mathrm{END}^{\dagger})
\le q^{\mathrm{heal}\dagger,=a}_{k_0}+q^{\mathrm{END}\dagger,=a}_{k_0}.
\end{equation*}
The event $E_a(\square)$ is an event of level $k_0\le K$: which cells lie in which live component,
with which oldest index, at level $k_0$ is determined by the outcomes of the run at levels $\le k_0$.
Hence \eqref{8.18:eq-run-prolongation-c} of Lemma~\ref{8.18:lem-run-prolongation}, which is proved
as an implication from the statements named there, gives
$\tilde{\nu}^{\dagger}(E_a(\square))=\tilde{\nu}(E_a(\square))$. Taking the maximum over $\square$
yields \eqref{8.18:eq-region-fate-a}.
\end{proof}

\begin{lemma}[Conversion to probabilities at the last lattice]\label{8.18:lem-probability-conversion}
Let $(g^{\dagger},N_T^{\dagger})$ be the prolonged member of Lemma~\ref{8.18:lem-run-prolongation}, with
last level $K^{\dagger\prime}$ and last lattice $\Lambda^{\dagger}$ of $(N^{\dagger})^4$ sites. Let
$\tilde{\nu}^{\dagger}$, $\tau_{\omega}$ and $Z=\int\rho_{K^{\dagger\prime}}$ be the normalised measure
on histories, the weights of histories and the normalising integral of the prolonged run, as in
Lemma~\ref{8.18:lem-region-fate}. Let $k_0\le K^{\dagger\prime}$ be a level, $a\ge0$, $\square$ a
cell of the level-$k_0$ lattice, and let $E_a(\square)$ and $\mathrm{END}^{\dagger}$ be the events of
Lemma~\ref{8.18:lem-region-fate}. Assume:
\begin{enumerate}
\item[(A)] the un-normalised age-and-reserve bound at the last lattice,
Lemma~\ref{8.18:lem-age-reserve-bound}, in its second form \eqref{8.18:eq-age-reserve-bound-b};
\item[(B)] the lower global stability bound at $\Lambda^{\dagger}$,
\begin{equation}\label{8.18:eq-probability-conversion-1}
Z=\int\rho_{K^{\dagger\prime}}\ \ge\ \exp\bigl(-C_{\mathrm{low}}(g^{\dagger})\,(N^{\dagger})^4\bigr),
\end{equation}
obtained from the lower stability bound of clause (i-b) of Theorem~\ref{3.1:thm-main}
(\cite[Theorem~1, (0.1)]{B-CMP122b}), with the constant $C_{\mathrm{low}}$ of
Definition~\ref{8.18:def-volume-horizon}, display \eqref{8.18:eq-volume-horizon-b}, which depends
on the member only through $x_{K^{\dagger\prime}}$.
\end{enumerate}
Then
\begin{equation}\label{8.18:eq-probability-conversion-a}
\tilde{\nu}^{\dagger}\bigl(E_a(\square)\cap\mathrm{END}^{\dagger}\bigr)
\ \le\ e^{\theta}\,q^{a}\,\exp\bigl(-\theta\,\mathfrak{P}(x_{k_0-a})\bigr)\,
\exp\bigl(\mathsf{B}(g^{\dagger})\,(N^{\dagger})^4\bigr),
\end{equation}
where
\begin{equation}\label{8.18:eq-probability-conversion-2}
\mathsf{B}(g^{\dagger})=E_+^{(\theta)}(g^{\dagger})+C_{\mathrm{low}}(g^{\dagger})
\ \le\ C_{\mathsf{B}}\bigl(1+\log X^{\dagger}\bigr),
\end{equation}
with $X^{\dagger}=(g^{\dagger})^{-2}$, the terminal inverse coupling of the prolonged member of
Lemma~\ref{8.18:lem-run-prolongation}.
\end{lemma}

\begin{proof}
We apply Lemma~\ref{8.18:lem-age-reserve-bound} to the longest run of the member
$(g^{\dagger},N_T^{\dagger})$ over the bare datum, with last level $\hat{K}=K^{\dagger\prime}$, side
$\hat{N}=N^{\dagger}$, the given $a\ge0$, level parameter $b=k_0-a$, and the event
$E=E_a(\square)\cap\mathrm{END}^{\dagger}$. The hypotheses on the level parameter and the event hold.
First, $k_0\le K^{\dagger\prime}$ gives $k_0-a\le K^{\dagger\prime}-a=\hat{K}-a$. Second, by the
definition of $\mathrm{END}^{\dagger}$ in Lemma~\ref{8.18:lem-region-fate}, every history
$\omega\in\mathrm{END}^{\dagger}$ has at level $K^{\dagger\prime}$ a terminal-live component $X$
with $j(X)\le k_0-a$; hence $E$ lies in the set of histories having at $\hat{K}$ a terminal-live
component $X$ with $j(X)\le k_0-a$, as Lemma~\ref{8.18:lem-age-reserve-bound} requires.
The second form \eqref{8.18:eq-age-reserve-bound-b} then gives
\begin{align}
\sum_{\omega\in E_a(\square)\cap\mathrm{END}^{\dagger}}\tau_{\omega}
&\le e^{\theta}\,q^{\hat{K}-b}\,\exp\bigl(-\theta\,\mathfrak{P}(x_{k_0-a})\bigr)\,
\exp\bigl(E_+^{(\theta)}(g^{\dagger})\,(N^{\dagger})^4\bigr)\notag\\
&\le e^{\theta}\,q^{a}\,\exp\bigl(-\theta\,\mathfrak{P}(x_{k_0-a})\bigr)\,
\exp\bigl(E_+^{(\theta)}(g^{\dagger})\,(N^{\dagger})^4\bigr).
\label{8.18:eq-probability-conversion-3}
\end{align}
In the second inequality we used $\hat{K}-b=K^{\dagger\prime}-k_0+a\ge a$, which holds because
$k_0\le K^{\dagger\prime}$, together with $q=e^{-c_a}<1$; hence $q^{\hat{K}-b}\le q^{a}$.

For the denominator we use hypothesis (B). The lower bound
\eqref{8.18:eq-probability-conversion-1} rests on three ingredients. One restricts the integral
$\int\rho_{K^{\dagger\prime}}$ to the set of configurations $U_{K^{\dagger\prime}}$ on
$\Lambda^{\dagger}$ each of whose bond variables differs from the identity by less than
$g_{K^{\dagger\prime}}$. On that set the lower bound of clause (i-b) of Theorem~\ref{3.1:thm-main}
applies. The Haar mass of the set is bounded below, per site, by a positive constant times a fixed
power of $g_{K^{\dagger\prime}}$.
Finally, on that set the small-field characteristic function $\chi_{K^{\dagger\prime}}$ equals $1$,
and the Wilson action $A(U_{K^{\dagger\prime}})$ is bounded by a constant times
$g_{K^{\dagger\prime}}^{2}\,(N^{\dagger})^4$. These are the three facts entering the constant
$C_{\mathrm{low}}$
of Definition~\ref{8.18:def-volume-horizon}, and the constant $C_{\mathrm{low}}$ so obtained
depends on the member only through $x_{K^{\dagger\prime}}$.

Dividing \eqref{8.18:eq-probability-conversion-3} by $Z$ and using
\eqref{8.18:eq-probability-conversion-1} gives
\begin{align*}
\tilde{\nu}^{\dagger}\bigl(E_a(\square)\cap\mathrm{END}^{\dagger}\bigr)
&=Z^{-1}\sum_{\omega\in E_a(\square)\cap\mathrm{END}^{\dagger}}\tau_{\omega}\\
&\le e^{\theta}\,q^{a}\,\exp\bigl(-\theta\,\mathfrak{P}(x_{k_0-a})\bigr)\,
\exp\bigl((E_+^{(\theta)}(g^{\dagger})+C_{\mathrm{low}}(g^{\dagger}))\,(N^{\dagger})^4\bigr).
\end{align*}
This is \eqref{8.18:eq-probability-conversion-a} with
$\mathsf{B}(g^{\dagger})=E_+^{(\theta)}(g^{\dagger})+C_{\mathrm{low}}(g^{\dagger})$. The bound
$\mathsf{B}(g^{\dagger})\le C_{\mathsf{B}}(1+\log X^{\dagger})$ in
\eqref{8.18:eq-probability-conversion-2} is the degree bound
\eqref{8.18:eq-volume-horizon-b} of Definition~\ref{8.18:def-volume-horizon} at
$g'=g^{\dagger}$.

Inequality \eqref{8.18:eq-probability-conversion-a} is the one step in the estimates of this
section at which a probability is obtained from un-normalised weights; the volume enters only there,
through the factor $(N^{\dagger})^4$ in its exponent.
\end{proof}

\begin{remark}[The degree comparison on small tori]\label{8.18:rem-degree-comparison}
Consider the prolonged member $(g^{\dagger},N_T^{\dagger})$ of Lemma~\ref{8.18:lem-run-prolongation}
(prolonging the run on the same lattice), with last lattice of side $N^{\dagger}$, and put
$X^{\dagger}=(g^{\dagger})^{-2}$. Let the cell $\square$, the level $k_0$ and the age $a$ be as in
Lemma~\ref{8.18:lem-probability-conversion} (conversion to probabilities at the last lattice),
so that the event $E_a(\square)$ concerns a live component whose oldest large-field region was
created at step $k_0-a$.
The pre-payment of that one creation event is $\exp(-\theta\,\mathfrak{P}(x_{k_0-a}))$. The
exponent of the global quotient in the same lemma is $\mathsf{B}(g^{\dagger})\,(N^{\dagger})^4$.
The threshold coupling $\gamma_*$ of the reference member (the member whose letters $K_*$ and
$X_*$ enter Definition~\ref{8.18:def-volume-horizon} and Lemma~\ref{8.18:lem-horizon-cases}) is
subject to the one-sided ceiling
\begin{equation}\label{8.18:eq-degree-comparison-a}
C_{\mathsf{B}}\,(1+\log y)\,\Bigl(2L^{c_{\mathrm{gap}}}\,L^{m_{\mathbb{T}}(y^{-1/2})}\Bigr)^{4}
\le \tfrac12\,\theta\,\bigl(\mathfrak{P}(y)-2\bigr)
\qquad\text{for every } y\ge X_* .
\end{equation}
This is the general volume ceiling on $\gamma_*$, taken with the factor $2$ multiplying
$L^{c_{\mathrm{gap}}}$ in its left member.
It is a condition on $\gamma_*$ (through the threshold $X_*$, which is fixed from $\gamma_*$ in
the definition of the reference member), chosen after $L$, $M$,
$\mathfrak{p}$, $C_{\mathsf{B}}$, $\theta$ and $c_{\mathrm{gap}}$. It is displayed here and is
adopted into no hypothesis.
On Case~V of Lemma~\ref{8.18:lem-horizon-cases} (tori below and beyond the volume horizon) we have
\begin{equation}\label{8.18:eq-degree-comparison-1}
\mathsf{B}(g^{\dagger})\,(N^{\dagger})^4\le\tfrac12\,\theta\,\bigl(\mathfrak{P}(x_{k_0-a})-2\bigr).
\end{equation}
In words, half of the creation pre-payment absorbs the quotient.

Beyond the horizon the situation is different. There $N_T^{\dagger}=2L^{K+m-k^{\dagger}}$ is free in
$m$, and $(N^{\dagger})^4\asymp L^{4(m-m_V(g))}$ up to factors depending only on $L$, the flow
constants and the letters of the reference member. No
function of $X^{\dagger}$ alone bounds this quantity. The quotient is then a genuine volume
factor. Only depth can compensate it, that is, a larger value of $\mathfrak{P}(x_{k_0})$, and this
is the route taken for such tori later in this chapter, through the volume-reading depth floor of
Definition~\ref{8.18:def-volume-horizon}. The alternative would be an input that
is not of the global kind, namely one of the inputs of the old-age tail statement; no such input
is used here.

Consider the inputs consisting of
\begin{itemize}
\item positivity;
\item the upper and the lower bound on the normalising integral, whose quotient is
$\exp\bigl(\mathsf{B}(g^{\dagger})\,(N^{\dagger})^4\bigr)$;
\item a reserve or age factor for each live component;
\item symmetry.
\end{itemize}
That the dichotomy is sharp for this input list is shown by model statements of this book.
Those statements concern the inputs only: they
say that the window in which a volume-free bound follows from these inputs is Case~V and nothing
wider. They do not concern the strength of the two-point witness proved in this chapter.
\end{remark}

\begin{proof}
\emph{The pre-payment.} By \eqref{8.18:eq-run-prolongation-b}, the levels of
Lemma~\ref{8.18:lem-probability-conversion} satisfy $k_0-a\le K\le k^{\dagger}$. By
Lemma~\ref{8.18:lem-run-prolongation}, $X^{\dagger}=x_{k^{\dagger}+1}$. The running inverse
couplings $x_k$ (Definition~\ref{1.2:def-couplings}) decrease strictly along the run, and
$\mathfrak{P}$ is increasing; hence
\begin{align}
x_{k_0-a}&\ge x_{k^{\dagger}}>x_{k^{\dagger}+1}=X^{\dagger},\notag\\
\mathfrak{P}(x_{k_0-a})&>\mathfrak{P}(X^{\dagger}).\label{8.18:eq-degree-comparison-2}
\end{align}
Since $\mathfrak{P}(y)=A_0(\log y)^{p_0}$, the creation pre-payment has degree $p_0$ in
$\log X^{\dagger}$.

\emph{The quotient.} By Definition~\ref{8.18:def-volume-horizon} (volume horizon and
volume-reading depth floor), $\mathsf{B}(g')\le C_{\mathsf{B}}(1+\log g'^{-2})$. Taking
$g'=g^{\dagger}$ gives
\begin{equation*}
\mathsf{B}(g^{\dagger})\,(N^{\dagger})^4\le C_{\mathsf{B}}\,(1+\log X^{\dagger})\,(N^{\dagger})^4 .
\end{equation*}
On Case~V, \eqref{8.18:eq-horizon-cases-a} gives
\begin{equation}\label{8.18:eq-degree-comparison-3}
(N^{\dagger})^4\le 16\,L^{4c_{\mathrm{gap}}}\,L^{4m_{\mathbb{T}}(g^{\dagger})}
\le 16\,L^{4c_{\mathrm{gap}}+48}\,M^4\,(\log X^{\dagger})^{4r}.
\end{equation}
This is a polylogarithm of degree $4r$ in $\log X^{\dagger}$. The exponent of the quotient
therefore has degree $4r+1$ in $\log X^{\dagger}$.

\emph{Comparison of degrees.} By the definition of the torus exponent function in
Notation~\ref{0.2:not-glossary}, $m_{\mathbb{T}}(g')=11+m'+\log_L R(g')$ with $M=L^{m'}$, so
$L^{4m_{\mathbb{T}}(g')}=L^{44}M^4R(g')^4$ at every coupling $g'$; the second inequality in
\eqref{8.18:eq-degree-comparison-3} is this identity together with the bound
$R(g')\le L(\log g'^{-2})^{r}$, which holds at every coupling $g'$ and through which the exponent
$r$ of $R(\cdot)$ enters; at $g'=g^{\dagger}$ it reads $R(g^{\dagger})\le L(\log X^{\dagger})^{r}$.
Applied with $y^{-1/2}$ in place of $g'$, the same bound bounds
the left member of \eqref{8.18:eq-degree-comparison-a} above by a polylogarithm of degree $4r+1$ in
$\log y$, while the right member has degree $p_0$ in $\log y$. The exponent $r$ of $R(\cdot)$
satisfies $r\ge 2$, and the glossary (Notation~\ref{0.2:not-glossary}) records $p_0\ge 5r$; hence
\begin{equation*}
p_0\ge 5r\ge 4r+2>4r+1 .
\end{equation*}
So the right member of \eqref{8.18:eq-degree-comparison-a} grows strictly faster in $\log y$ than
the left member, and \eqref{8.18:eq-degree-comparison-a} holds for all $y$ beyond a threshold
depending only on $L$, $M$, $\mathfrak{p}$, $C_{\mathsf{B}}$, $\theta$ and $c_{\mathrm{gap}}$;
imposing it for all $y\ge X_*$ is the one-sided ceiling on $\gamma_*$ described above.

\emph{The absorption.} Since $g^{\dagger}=(X^{\dagger})^{-1/2}$, the first bound in
\eqref{8.18:eq-degree-comparison-3} can be written as
\begin{equation*}
(N^{\dagger})^4\le\Bigl(2L^{c_{\mathrm{gap}}}\,L^{m_{\mathbb{T}}((X^{\dagger})^{-1/2})}\Bigr)^{4}.
\end{equation*}
Hence $\mathsf{B}(g^{\dagger})(N^{\dagger})^4$ is at most the left member of
\eqref{8.18:eq-degree-comparison-a} at $y=X^{\dagger}$. On Case~V,
Lemma~\ref{8.18:lem-horizon-cases} and Lemma~\ref{8.18:lem-run-prolongation} give
\begin{equation*}
k^{\dagger}\le\hat{\kappa}\qquad\text{and}\qquad\hat{\kappa}=K+m\le K_*-1 .
\end{equation*}
Hence $k^{\dagger}+1\le K_*$, and since the $x_k$ decrease along the run,
\begin{equation*}
X^{\dagger}=x_{k^{\dagger}+1}\ge x_{K_*}\ge X_*,
\end{equation*}
the last inequality by the choice of the level $K_*$ of the reference member. So
\eqref{8.18:eq-degree-comparison-a} applies at $y=X^{\dagger}$. By
\eqref{8.18:eq-degree-comparison-a} at $y=X^{\dagger}$, and then by
\eqref{8.18:eq-degree-comparison-2},
\begin{equation*}
\mathsf{B}(g^{\dagger})\,(N^{\dagger})^4
\le\tfrac12\,\theta\,\bigl(\mathfrak{P}(X^{\dagger})-2\bigr)
<\tfrac12\,\theta\,\bigl(\mathfrak{P}(x_{k_0-a})-2\bigr).
\end{equation*}
This is \eqref{8.18:eq-degree-comparison-1}.
So on Case~V half of the pre-payment of the one creation event absorbs the global quotient.
\end{proof}

\begin{proposition}[Never-healed age slices below the horizon]\label{8.18:prop-never-healed-slices}
Assume the following.
\begin{itemize}
\item[(a)] The statements from which the conclusion is derived:
\begin{itemize}
\item the un-normalised age-and-reserve bound at the last lattice
(Lemma~\ref{8.18:lem-age-reserve-bound});
\item the lower bound for the normalising integral together with the degree bound for
$\mathsf{B}$ (Definition~\ref{8.18:def-volume-horizon});
\item the re-targeting statement for members of the family;
\item rarity of young large-field blocks;
\item the torus statement on which Lemma~\ref{8.18:lem-horizon-cases} rests;
\item the display \eqref{8.18:eq-members-runs-a}.
\end{itemize}
\item[(b)] The degree comparison on small tori, \eqref{8.18:eq-degree-comparison-a}.
\item[(c)] The ceiling $\gamma_*\le\gamma_L(c_a)$ on the reference coupling.
\item[(d)] The reserve window for the reserve fraction $\theta$.
\item[(e)] The floor on the free number $c_a$:
\begin{equation}\label{8.18:eq-never-healed-slices-a}
c_a\ \ge\ 9\log L+\log 2 .
\end{equation}
\end{itemize}
Here $c_a$ is chosen after $L$. Condition \eqref{8.18:eq-never-healed-slices-a} is a lower bound only. Since $q=e^{-c_a}$, it gives $q\le(2L^9)^{-1}$. Enlarging $c_a$ moves only the ceiling $\gamma_L(c_a)$ in~(c).

Suppose we are in Case~V of Lemma~\ref{8.18:lem-horizon-cases}, that is, $K+m\le K_*-1$. By \eqref{8.18:eq-horizon-cases-b} this holds in particular for every $m\le m_V(g)$. Then for every $a\ge0$ and every cell $\square$ of level $k_0$,
\begin{equation}\label{8.18:eq-never-healed-slices-b}
\tilde{\nu}\bigl(\mathfrak{lv}^{=a}_{k_0}(\square)\cap\mathrm{END}^{\dagger}\bigr)
\ \le\ e^{3\theta/2}\,(2L^9)^{-a}\,
\exp\bigl(-\tfrac12\theta\,\mathfrak{P}(x_{k_0})\bigr).
\end{equation}
The bound holds uniformly in $K$, in $g_0$, in the position of $\square$, and in $m$ throughout Case~V. No condition on $s$ is imposed beyond the one in clause~(iv) of Theorem~0.
\end{proposition}

\begin{proof}
Fix $m$ in Case~V, $a\ge0$ and a cell $\square$ of level $k_0$. Let $(g^{\dagger},N_T^{\dagger})$ be the prolonged member of Lemma~\ref{8.18:lem-run-prolongation}. Its last level is $K^{\dagger\prime}$, its last lattice is $\Lambda^{\dagger}$ with $N^{\dagger}$ sites per side, and $X^{\dagger}=(g^{\dagger})^{-2}$. Let $E_a(\square)$ and $\mathrm{END}^{\dagger}$ be the events of Lemma~\ref{8.18:lem-region-fate}.

\emph{Step 1: the conversion.} The conversion to probabilities at the last lattice, \eqref{8.18:eq-probability-conversion-a}, is Lemma~\ref{8.18:lem-probability-conversion}. It is obtained from Lemma~\ref{8.18:lem-age-reserve-bound} and the lower bound for the normalising integral at $\Lambda^{\dagger}$. It gives
\begin{equation}\label{8.18:eq-never-healed-slices-1}
\tilde{\nu}^{\dagger}\bigl(E_a(\square)\cap\mathrm{END}^{\dagger}\bigr)
\le e^{\theta}\,q^{a}\,\exp\bigl(-\theta\,\mathfrak{P}(x_{k_0-a})\bigr)\,
\exp\bigl(\mathsf{B}(g^{\dagger})\,(N^{\dagger})^4\bigr).
\end{equation}
We split the exponent of the middle factor into two equal halves:
\begin{equation*}
\theta\,\mathfrak{P}(x_{k_0-a})
=\tfrac12\theta\,\mathfrak{P}(x_{k_0-a})+\tfrac12\theta\,\mathfrak{P}(x_{k_0-a}).
\end{equation*}

\emph{Step 2: the first half pays for the volume.} Displays \eqref{8.18:eq-run-prolongation-b} and \eqref{8.18:eq-horizon-cases-a} give three facts:
\begin{itemize}
\item $\tfrac12\theta\,\mathfrak{P}(x_{k_0-a})\ge\tfrac12\theta\,\mathfrak{P}(X^{\dagger})$;
\item $X^{\dagger}\ge X_*$;
\item on Case~V the side satisfies $N^{\dagger}\le 2L^{c_{\mathrm{gap}}}\,L^{m_{\mathbb{T}}(g^{\dagger})}$.
\end{itemize}
The degree bound for $\mathsf{B}$ gives
$\mathsf{B}(g^{\dagger})\le C_{\mathsf{B}}(1+\log X^{\dagger})$.
Now apply \eqref{8.18:eq-degree-comparison-a} at $y=X^{\dagger}$. This is admissible since $X^{\dagger}\ge X_*$, and $y^{-1/2}=g^{\dagger}$. We obtain
\begin{align*}
\mathsf{B}(g^{\dagger})\,(N^{\dagger})^4
&\le C_{\mathsf{B}}\bigl(1+\log X^{\dagger}\bigr)
\bigl(2L^{c_{\mathrm{gap}}}\,L^{m_{\mathbb{T}}(g^{\dagger})}\bigr)^4\\
&\le \tfrac12\theta\bigl(\mathfrak{P}(X^{\dagger})-2\bigr)
=\tfrac12\theta\,\mathfrak{P}(X^{\dagger})-\theta .
\end{align*}
Hence
\begin{equation}\label{8.18:eq-never-healed-slices-2}
\tfrac12\theta\,\mathfrak{P}(x_{k_0-a})\ \ge\ \mathsf{B}(g^{\dagger})\,(N^{\dagger})^4+\theta .
\end{equation}

\emph{Step 3: the second half.} By \eqref{8.18:eq-members-runs-a}, applied between the levels $k_0-a$ and $k_0$,
\begin{equation}\label{8.18:eq-never-healed-slices-3}
\tfrac12\theta\,\mathfrak{P}(x_{k_0-a})\ \ge\ \tfrac12\theta\bigl(\mathfrak{P}(x_{k_0})-1\bigr).
\end{equation}

\emph{Step 4: the age factor.} Since $q=e^{-c_a}$, the floor \eqref{8.18:eq-never-healed-slices-a} gives
\begin{equation*}
q^{a}=e^{-a c_a}\le(2L^9)^{-a}.
\end{equation*}

\emph{Step 5: collecting the factors.} Insert \eqref{8.18:eq-never-healed-slices-2}, \eqref{8.18:eq-never-healed-slices-3} and Step~4 into \eqref{8.18:eq-never-healed-slices-1}. The factor $\exp(\mathsf{B}(g^{\dagger})(N^{\dagger})^4)$ cancels against the first half, and we get
\begin{align*}
\tilde{\nu}^{\dagger}\bigl(E_a(\square)\cap\mathrm{END}^{\dagger}\bigr)
&\le e^{\theta}\,(2L^9)^{-a}\,e^{-\theta}\,e^{\theta/2}\,
\exp\bigl(-\tfrac12\theta\,\mathfrak{P}(x_{k_0})\bigr)\\
&\le e^{3\theta/2}\,(2L^9)^{-a}\,\exp\bigl(-\tfrac12\theta\,\mathfrak{P}(x_{k_0})\bigr).
\end{align*}
The three exponential constants have the following origins:
\begin{itemize}
\item $e^{\theta}$ is the factor of \eqref{8.18:eq-probability-conversion-a}, which comes from \eqref{8.18:eq-members-runs-a} inside Lemma~\ref{8.18:lem-age-reserve-bound};
\item $e^{-\theta}$ is the term $+\theta$ kept in \eqref{8.18:eq-never-healed-slices-2}, which comes from the $-2$ in \eqref{8.18:eq-degree-comparison-a};
\item $e^{\theta/2}$ comes from \eqref{8.18:eq-never-healed-slices-3}.
\end{itemize}
Their product $e^{\theta/2}$ is at most $e^{3\theta/2}$ because $\theta\ge0$.

\emph{Step 6: return to the original member.} The event $E_a(\square)=\mathfrak{lv}^{=a}_{k_0}(\square)$ is an event of level $k_0\le K$. Display \eqref{8.18:eq-run-prolongation-c} therefore identifies the left side of the last inequality with the left side of \eqref{8.18:eq-never-healed-slices-b}, read for the original member. This proves \eqref{8.18:eq-never-healed-slices-b}.

The right side of \eqref{8.18:eq-never-healed-slices-b} depends only on $\theta$, $L$, $a$ and $x_{k_0}$. This gives the asserted uniformity in $K$, $g_0$, the position of $\square$, and $m$ within Case~V.
\end{proof}

The factor $(2L^9)^{-a}$ in \eqref{8.18:eq-never-healed-slices-b} decays faster than $L^{-9a}$ in the age $a$. The rate per level is $c_a$, and $c_a$ is free upward. Replacing the floor \eqref{8.18:eq-never-healed-slices-a} by $c_a\ge C\log L+\log 2$ gives the factor $(2L^{C})^{-a}$ for any constant $C$. The only price is the ceiling $\gamma_L(c_a)$.

\begin{corollary}[The never-healed age series below the horizon]\label{8.18:cor-never-healed-series}
Assume the hypotheses of Proposition~\ref{8.18:prop-never-healed-slices}, including the floor
\eqref{8.18:eq-never-healed-slices-a} on the free number $c_a$, and place ourselves on Case~V, that
is, $K+m\le K_*-1$. Then the age series of the one-cell weights $q^{\mathrm{END}\dagger,=a}_{k_0}$
beyond the age threshold $\mathfrak{l}$ satisfies
\begin{equation}\label{8.18:eq-never-healed-series-a}
\sum_{a>\mathfrak{l}} L^{8(a+1)}\, q^{\mathrm{END}\dagger,=a}_{k_0}
\le e^{3\theta/2}\, L^{8(\mathfrak{l}+1)}\, e^{-\frac12\theta\,\mathfrak{P}(x_{k_0})}.
\end{equation}
Consequently, put $c_{\mathrm{lv}}:=\theta/2$, let $C_{\mathrm{lv}}\ge 1$ be the constant with which
the corollary to the rarity of young large-field blocks holds on Case~V at this exponent, and put
\begin{equation*}
q^{\mathrm{lv}}_{k_0}:=C_{\mathrm{lv}}\,e^{-c_{\mathrm{lv}}\mathfrak{P}(x_{k_0})},
\qquad
A^{\mathrm{END}}_{\mathrm{lv}}:=\frac{e^{3\theta/2}}{C_{\mathrm{lv}}}\le e^{3\theta/2}.
\end{equation*}
Then
\begin{equation}\label{8.18:eq-never-healed-series-1}
\sum_{a>\mathfrak{l}} L^{8(a+1)}\, q^{\mathrm{END}\dagger,=a}_{k_0}
\le A^{\mathrm{END}}_{\mathrm{lv}}\, L^{8(\mathfrak{l}+1)}\, q^{\mathrm{lv}}_{k_0},
\end{equation}
with $A^{\mathrm{END}}_{\mathrm{lv}}$ a constant; that is, the $\mathrm{END}^{\dagger}$ part of the
age hypothesis of the additive large-field theorem holds at level $k_0$, with the same one-cell
weight $q^{\mathrm{lv}}_{k_0}$ as the one that the corollary to the rarity of young large-field
blocks delivers on Case~V.
\end{corollary}

\begin{proof}
By Proposition~\ref{8.18:prop-never-healed-slices}, which is proved as an implication from the
statements named there, the slice bound \eqref{8.18:eq-never-healed-slices-b} holds on Case~V
for every $a\ge 0$ and every level-$k_0$ cell $\square$; read for the one-cell weights it gives,
for every $a\ge0$,
\begin{equation*}
q^{\mathrm{END}\dagger,=a}_{k_0}
\le e^{3\theta/2}\,(2L^{9})^{-a}\,e^{-\frac12\theta\,\mathfrak{P}(x_{k_0})}.
\end{equation*}
The factor $(2L^9)^{-a}$ is where the floor \eqref{8.18:eq-never-healed-slices-a} enters: it makes
$q=e^{-c_a}\le(2L^9)^{-1}$. We multiply by $L^{8(a+1)}=L^{8}L^{8a}$ and sum over the integers
$a>\mathfrak{l}$, that is, over $a\ge\mathfrak{l}+1$. Since
$L^{8a}(2L^9)^{-a}=(2L)^{-a}$, we obtain
\begin{align*}
\sum_{a>\mathfrak{l}} L^{8(a+1)}\, q^{\mathrm{END}\dagger,=a}_{k_0}
&\le e^{3\theta/2}e^{-\frac12\theta\mathfrak{P}(x_{k_0})}\,
L^{8}\sum_{a\ge\mathfrak{l}+1}L^{8a}(2L^{9})^{-a}\\
&= e^{3\theta/2}e^{-\frac12\theta\mathfrak{P}(x_{k_0})}\,
L^{8}\sum_{a\ge\mathfrak{l}+1}(2L)^{-a}\\
&\le e^{3\theta/2}e^{-\frac12\theta\mathfrak{P}(x_{k_0})}\,L^{8}\,(2L)^{-\mathfrak{l}}\\
&\le e^{3\theta/2}\,L^{8(\mathfrak{l}+1)}\,e^{-\frac12\theta\mathfrak{P}(x_{k_0})}.
\end{align*}
The third line is the geometric sum:
\begin{equation*}
\sum_{a\ge\mathfrak{l}+1}(2L)^{-a}=\frac{(2L)^{-\mathfrak{l}}}{2L-1}\le(2L)^{-\mathfrak{l}},
\end{equation*}
because $2L-1\ge1$. The fourth line is the inequality $L^{8}(2L)^{-\mathfrak{l}}\le
L^{8(\mathfrak{l}+1)}$, which is $(2L^{9})^{-\mathfrak{l}}\le1$ and holds since $\mathfrak{l}\ge0$
and $L\ge1$. This proves \eqref{8.18:eq-never-healed-series-a}.

For \eqref{8.18:eq-never-healed-series-1}, the definitions of $q^{\mathrm{lv}}_{k_0}$ and
$A^{\mathrm{END}}_{\mathrm{lv}}$ at $c_{\mathrm{lv}}=\theta/2$ give
\begin{equation*}
A^{\mathrm{END}}_{\mathrm{lv}}\,q^{\mathrm{lv}}_{k_0}
=\frac{e^{3\theta/2}}{C_{\mathrm{lv}}}\,C_{\mathrm{lv}}\,e^{-\frac12\theta\mathfrak{P}(x_{k_0})}
=e^{3\theta/2}\,e^{-\frac12\theta\mathfrak{P}(x_{k_0})},
\end{equation*}
so the right side of \eqref{8.18:eq-never-healed-series-1} equals that of
\eqref{8.18:eq-never-healed-series-a}. Since $C_{\mathrm{lv}}\ge1$, we have
$A^{\mathrm{END}}_{\mathrm{lv}}\le e^{3\theta/2}$, a constant. The age hypothesis of the additive
large-field theorem asks for a bound of the form of \eqref{8.18:eq-never-healed-series-1} with a
constant, or a polynomial in $\log x_{k_0}$, in front of $L^{8(\mathfrak{l}+1)}q^{\mathrm{lv}}_{k_0}$;
the $\mathrm{END}^{\dagger}$ part of it is therefore \eqref{8.18:eq-never-healed-series-1}.
Finally, the corollary to the rarity of young large-field blocks delivers on Case~V the one-cell
weight $C_{\mathrm{lv}}e^{-c_{\mathrm{lv}}\mathfrak{P}(x_{k_0})}$ at the same exponent
$c_{\mathrm{lv}}=\theta/2$, and the statements that use these weights accept any fixed exponent in
$(0,\theta]$, which contains $\theta/2$. Hence the additive large-field theorem reads one and the
same $q^{\mathrm{lv}}_{k_0}$ for both of its antecedents.
\end{proof}

\begin{lemma}[Healed-later age slices of the prolonged run]\label{8.18:lem-healed-slices}
Let $(\dagger)$ be the prolonged member of Lemma~\ref{8.18:lem-run-prolongation}, with coupling
$g^{\dagger}$, side $N_T^{\dagger}$ and last level $K^{\dagger\prime}$ of its longest run. For a level
$k_0\le K^{\dagger\prime}$, a level-$k_0$ cell $\square$ and an age $a\ge0$, let $E_a(\square)$,
$\mathrm{HEAL}^{\dagger}$ and $q^{\mathrm{heal}\dagger,=a}_{k_0}$ be as in
Lemma~\ref{8.18:lem-region-fate}, so that
\begin{equation*}
q^{\mathrm{heal}\dagger,=a}_{k_0}=\max_{\square}\tilde{\nu}^{\dagger}\bigl(E_a(\square)\cap
\mathrm{HEAL}^{\dagger}\bigr).
\end{equation*}
Assume the following two hypotheses.
\begin{itemize}
\item[(a)] The healed old-region bound holds in the following form. For every member, in the
sense of~\eqref{8.18:eq-members-runs-a}, whose reference coupling lies below the ceilings of that
bound (ceilings chosen after $L$), at every level $k$ not exceeding the last level of the member's
longest run, with no depth floor, uniformly in the position of the cell and uniformly in $K$, $m$
and $N_T$,
\begin{equation}\label{8.18:eq-healed-slices-1}
\sum_{a>\mathfrak{l}}L^{8(a+1)}\,q^{\mathrm{heal},=a}_k
\le A^{\mathrm{heal}}_{\mathrm{lv}}\,L^{8(\mathfrak{l}+1)}\,q^{\mathrm{lv}}_k ,
\end{equation}
the age-$a$ slice carrying the per-age decay factor
$\exp\bigl(-\max\{\theta\,p_0(g_k),\,\varrho_*(k)\,a\}\bigr)$; here $q^{\mathrm{heal},=a}_k$ is the
maximum over level-$k$ cells of the weight of the age-$a$ event intersected with the event that the
live component is removed by $\mathbf{R}$ at some later level of the run.
Alternatively, the same bound is supplied by the healed old-region bound in the form derived from
geometric forgetting of old regions, under the floor $q\,L^{8}\le\tfrac12$ on $q=e^{-c_a}$.
\item[(b)] Throughout, the ceiling $\gamma_*$ on the reference coupling is taken to be
\begin{equation}\label{8.18:eq-healed-slices-a}
\gamma_*^{\dagger}:=\min\bigl\{\gamma_J,\ (\gamma_V)^{\flat},\ \gamma_L(c_a),\
\gamma_{\mathrm{heal}}\bigr\},
\end{equation}
where $(\gamma_V)^{\flat}$ denotes the ceilings on the reference coupling under which
Proposition~\ref{8.18:prop-never-healed-slices} is stated, and $\gamma_{\mathrm{heal}}$ is the
smallest of the ceilings on the reference coupling under which the healed old-region bound, the
large-field block bound, healed part, and one further proposition whose ceilings are imposed
together with these hold; in particular $\gamma_*=\gamma_*^{\dagger}\le\gamma_{\mathrm{heal}}$.
\end{itemize}
Then for every level $k_0\le K^{\dagger\prime}$,
\begin{equation}\label{8.18:eq-healed-slices-2}
\sum_{a>\mathfrak{l}}L^{8(a+1)}\,q^{\mathrm{heal}\dagger,=a}_{k_0}
\le A^{\mathrm{heal}}_{\mathrm{lv}}\,L^{8(\mathfrak{l}+1)}\,q^{\mathrm{lv}}_{k_0},
\qquad \mathfrak{l}=\mathfrak{l}(x_{k_0}),
\end{equation}
with the per-age decay factor $\exp\bigl(-\max\{\theta\,p_0(g_{k_0}),\,\varrho_*(k_0)\,a\}\bigr)$ on
the age-$a$ slice. The left-hand side of~\eqref{8.18:eq-healed-slices-2} accounts for all live
components of age $a>\mathfrak{l}$ at level $k_0$ that are healed at some level in
$\left]k_0,K^{\dagger\prime}\right]$, in particular for those healed at a level in
$\left]K,K^{\dagger\prime}\right]$, beyond the terminal level $K$ of the original run. The volume
does not enter~\eqref{8.18:eq-healed-slices-2}.
\end{lemma}

\begin{proof}
The proof applies hypothesis~(a) to the member $(\dagger)$ at the level $k_0$; we check that it
applies and identify its terms.

\emph{$(\dagger)$ is a member.} By Lemma~\ref{8.18:lem-run-prolongation}, $(\dagger)$ is a member
$(g^{\dagger},N_T^{\dagger})$ in the sense of~\eqref{8.18:eq-members-runs-a}, over the same bare
lattice as the original member; its admissibility is supplied, in the proof of that lemma, by the
admissibility lemma for re-targeted members. Its longest run ends at the level
$K^{\dagger\prime}$.

\emph{The ceilings hold along $(\dagger)$.} By the construction of $(\dagger)$ in
Lemma~\ref{8.18:lem-run-prolongation}, together with hypothesis~(b)
and~\eqref{8.18:eq-healed-slices-a}, the running couplings of the prolonged run satisfy
\begin{equation*}
g_k<g^{\dagger}<\gamma_*=\gamma_*^{\dagger}\le\gamma_{\mathrm{heal}}
\qquad\text{for all } k\le k^{\dagger}.
\end{equation*}
Hence $g^{\dagger}<\gamma_{\mathrm{heal}}$, that is, the reference coupling
of $(\dagger)$ lies below the ceilings of the healed old-region bound and of the large-field block
bound, healed part. This is the reason why the ceiling $\gamma_{\mathrm{heal}}$ is included
in~\eqref{8.18:eq-healed-slices-a}.

\emph{The level is admissible and no uniformity is lost.} The level $k_0$ satisfies
$k_0\le K^{\dagger\prime}$, the last level of the longest run of $(\dagger)$. Hypothesis~(a) is
stated at every such level with no depth floor, so no lower bound on $K^{\dagger\prime}-k_0$ is
needed. It is uniform in $K$, $m$ and $N_T$, so the replacement of the side $N_T$ of the original
member by the side $N_T^{\dagger}$ of $(\dagger)$ and of the terminal level $K$ by
$K^{\dagger\prime}$ does not affect the constant $A^{\mathrm{heal}}_{\mathrm{lv}}$; it is uniform in
the position of the cell, so the maximum over $\square$ in the definition of
$q^{\mathrm{heal}\dagger,=a}_{k_0}$ is covered.

\emph{Identification of the terms.} For the member $(\dagger)$ the history measure is
$\tilde{\nu}^{\dagger}$, and the healed event of hypothesis~(a) at level $k_0$ is the event that the
live component is removed by $\mathbf{R}$ at some level $k_h\in\left]k_0,K^{\dagger\prime}\right]$;
this is the event $\mathrm{HEAL}^{\dagger}$ of Lemma~\ref{8.18:lem-region-fate}. Therefore the
quantity $q^{\mathrm{heal},=a}_{k_0}$ of hypothesis~(a), read for $(\dagger)$, is
$q^{\mathrm{heal}\dagger,=a}_{k_0}$. The numbers $\mathfrak{l}=\mathfrak{l}(x_{k_0})$ and
$q^{\mathrm{lv}}_{k_0}$ on the right of~\eqref{8.18:eq-healed-slices-1} are read from the running
inverse coupling $x_{k_0}$, which, by Lemma~\ref{8.18:lem-run-prolongation}, is the same in the
original run and in the prolonged run.
Substituting these identifications in~\eqref{8.18:eq-healed-slices-1} at $k=k_0$
gives~\eqref{8.18:eq-healed-slices-2}, together with the per-age decay factor.

\emph{Lineages healed beyond $K$.} Since $\mathrm{HEAL}^{\dagger}$ contains every lineage removed
by $\mathbf{R}$ at a level in $\left]k_0,K^{\dagger\prime}\right]$, it contains those removed at a
level in $\left]K,K^{\dagger\prime}\right]$. Read in the original run, which stops at $K$, these are
live at the terminal level; in the prolonged run they lie on the healed side and are bounded
by~\eqref{8.18:eq-healed-slices-2}, with the same constant $A^{\mathrm{heal}}_{\mathrm{lv}}$.

\emph{The volume.} Hypothesis~(a) is uniform in $m$ and $N_T$, so no quantity depending on the
volume enters~\eqref{8.18:eq-healed-slices-2}. The underlying reason is that the healed
old-region bound rests on Ba{\l}aban's normalised quotient \cite[(1.75)]{B-CMP122b} and on
positivity at the healing step, that is, on a comparison, local to the healed region, between the
configurations in which the large-field factor is present and those in which it is absent.

\emph{The alternative supplier.} If hypothesis~(a) is supplied by the healed old-region bound in the
form derived from geometric forgetting of old regions, its floor $qL^{8}\le\tfrac12$ must hold. It
does under the floor~\eqref{8.18:eq-never-healed-slices-a}, $c_a\ge 9\log L+\log 2$: this gives
$q=e^{-c_a}\le(2L^{9})^{-1}$, hence
\begin{equation*}
q\,L^{8}\le\frac{L^{8}}{2L^{9}}=\frac{1}{2L}\le\frac12 .
\end{equation*}
The remaining steps are unchanged.

\emph{The ceilings not included.} In~\eqref{8.18:eq-healed-slices-a} the ceilings of rarity of
young large-field blocks do not appear: they are needed only where one particular case of
Corollary~\ref{8.8:cor-young-block-rarity} is invoked, and the arguments in which $\gamma_*^{\dagger}$
is used do not invoke that case.
\end{proof}

\begin{proposition}[The age antecedent below the horizon]\label{8.18:prop-age-antecedent}
Assume Case~V, with the level $k_0$, the age threshold $\mathfrak{l}$, the one-cell
age-$a$ weights $q^{=a}_{k_0}$ of the original member and the one-cell live weight
$q^{\mathrm{lv}}_{k_0}$ as in Lemma~\ref{8.18:lem-region-fate} and
Corollary~\ref{8.18:cor-never-healed-series}. Let $A^{\mathrm{heal}}_{\mathrm{lv}}$ be the constant of
Lemma~\ref{8.18:lem-healed-slices}, or, where that lemma is used in its polylogarithmic variant,
the polynomial in $\log x_{k_0}$ given there. Then
\begin{equation}\label{8.18:eq-age-antecedent-1}
\begin{split}
\sum_{a>\mathfrak{l}}L^{8(a+1)}\,q^{=a}_{k_0}
&\le \bigl(A^{\mathrm{heal}}_{\mathrm{lv}}+e^{3\theta/2}\bigr)\,L^{8(\mathfrak{l}+1)}\,
q^{\mathrm{lv}}_{k_0}\\
&=: A_{\mathrm{lv}}\,L^{8(\mathfrak{l}+1)}\,q^{\mathrm{lv}}_{k_0}.
\end{split}
\end{equation}
Here $A_{\mathrm{lv}}$ is a constant; where Lemma~\ref{8.18:lem-healed-slices} is used in its
polylogarithmic variant, $A_{\mathrm{lv}}$ is the corresponding polynomial in $\log x_{k_0}$.
The bound \eqref{8.18:eq-age-antecedent-1} holds uniformly in $K$, $g_0$, the position of the cell
and $m$ throughout Case~V, and no depth floor is imposed. It is the age hypothesis of the additive
theorem, in exactly the form in which that theorem takes it.
\end{proposition}

\begin{proof}
Let $a>\mathfrak{l}$. Lemma~\ref{8.18:lem-region-fate} splits the event $E_a(\square)$ of the
prolonged run according to the fate of the region: it is either removed by $\mathbf{R}$ at some level
$k_h\in\,]k_0,K^{\dagger\prime}]$ (the event $\mathrm{HEAL}^{\dagger}$), or it is still live at the
last level $K^{\dagger\prime}$ (the event $\mathrm{END}^{\dagger}$). By
\eqref{8.18:eq-region-fate-a}, $q^{=a}_{k_0}\le q^{\mathrm{heal}\dagger,=a}_{k_0}
+q^{\mathrm{END}\dagger,=a}_{k_0}$,
and the same lemma states that the left member is the age-$a$ weight $q^{=a}_{k_0}$ of the original
member, because $E_a(\square)$ is an event of level $k_0\le K$. We multiply by $L^{8(a+1)}\ge 0$
and sum over $a>\mathfrak{l}$.

For the never-healed share, \eqref{8.18:eq-never-healed-series-a} of
Corollary~\ref{8.18:cor-never-healed-series} bounds the sum over $a>\mathfrak{l}$ of
$L^{8(a+1)}\,q^{\mathrm{END}\dagger,=a}_{k_0}$ by
$A^{\mathrm{END}}_{\mathrm{lv}}\,L^{8(\mathfrak{l}+1)}\,q^{\mathrm{lv}}_{k_0}$, with
$A^{\mathrm{END}}_{\mathrm{lv}}\le e^{3\theta/2}$.
For the healed share we apply Lemma~\ref{8.18:lem-healed-slices} to the prolonged run at the level
$k=k_0$. This is admissible, since that lemma holds at every level up to the last level
$K^{\dagger\prime}$ of the prolonged run and $k_0\le K\le K^{\dagger\prime}$. Then
\eqref{8.18:eq-healed-slices-a} at $k=k_0$ bounds the sum over $a>\mathfrak{l}$ of
$L^{8(a+1)}\,q^{\mathrm{heal}\dagger,=a}_{k_0}$ by
$A^{\mathrm{heal}}_{\mathrm{lv}}\,L^{8(\mathfrak{l}+1)}\,q^{\mathrm{lv}}_{k_0}$.
Adding the two bounds gives \eqref{8.18:eq-age-antecedent-1}.

As for uniformity, each of the three inputs is uniform in $K$, $g_0$, the position and $m$ on
Case~V, and none of them imposes a depth floor. Hence the same holds for their sum. The constant
$A_{\mathrm{lv}}$ is the sum of $e^{3\theta/2}$ and $A^{\mathrm{heal}}_{\mathrm{lv}}$, so it is
a constant, or a polynomial in $\log x_{k_0}$ when $A^{\mathrm{heal}}_{\mathrm{lv}}$ is.
\end{proof}

\begin{proposition}[Clause (iv) on tori below the horizon]\label{8.18:prop-witness-below-horizon}
Assume the following statements.
\begin{itemize}
\item[(1)] The additive bound on the live member $R^{\mathrm{BAD}\prime}$ of the two-point witness:
under its standing assumptions, its two antecedents at level $k_0$ (the one-cell live bound and
the age antecedent) and its own smallness conditions, the live member satisfies
$|R^{\mathrm{BAD}\prime}|\le\mathfrak{M}/32$; hence, once the remaining members of
$S^T_{2;K}-\mathfrak{M}$ are bounded in modulus, the bounds
\eqref{8.18:eq-witness-below-horizon-2} below and the display of clause~(iv) of
Theorem~\ref{3.1:thm-main} hold.
\item[(2)] Corollary~\ref{8.8:cor-young-block-rarity}, on the rarity of young large-field blocks, in
Case~V and under the ceiling
\eqref{8.18:eq-degree-comparison-a}: at every level $j\le K$ of the run and for every level-$j$
cell $\square$,
\begin{equation}\label{8.18:eq-witness-below-horizon-1}
\tilde{\nu}\bigl(\mathfrak{lv}_j(\square)\bigr)\le q^{\mathrm{lv}}_j .
\end{equation}
\item[(3)] The correlation theorem, applied with the torus exponent of its statement set equal
to $s$ (that exponent is not the index $m$ of the torus $\mathbb{T}_m$ considered here).
\item[(4)] The Gaussian extraction, parts (a), (b) and (c).
\item[(5)] The two-sided kernel bounds for the main term $\mathfrak{M}$.
\item[(6)] Clause (iii-c) of Theorem~\ref{3.1:thm-main}.
\item[(7)] The lemma bounding $N_{k_0}$.
\end{itemize}
Put
\begin{equation}\label{8.18:eq-witness-below-horizon-a}
\gamma_W^{\mathbb{T}}:=\min\{\gamma_W,\ \gamma_*^{\dagger}\},
\end{equation}
where $\gamma_*^{\dagger}$ is the ceiling \eqref{8.18:eq-healed-slices-a} and the second entry of
the minimum is taken with the margin of \eqref{8.18:eq-horizon-cases-d}.
Then for every $g\le\gamma_W^{\mathbb{T}}$, every $g_0$ with $0<g_0\le g_-(g)$, every $K$ and $m$
for which the torus $\mathbb{T}_m$ is in Case~V of Lemma~\ref{8.18:lem-horizon-cases}, that is,
for which \eqref{8.18:eq-horizon-cases-a} holds, every $\mathsf{d}$ satisfying the depth condition
and the ultraviolet condition of clause~(iv), with $s:=s(\mathsf{d})$, and every
$(f_1,f_2)\in\mathfrak{T}(\mathsf{d};\vartheta)$,
\begin{equation}\label{8.18:eq-witness-below-horizon-2}
\bigl|S^T_{2;K}-\mathfrak{M}\bigr|\le\tfrac{6}{32}\,\mathfrak{M},
\qquad |R_K|\le\tfrac14 ,
\end{equation}
and the display of clause~(iv) of Theorem~\ref{3.1:thm-main} holds with every one of its
constants. These
conclusions are uniform in $K$, in $g_0$, in the position of the supports of $f_1,f_2$ on
$\mathbb{T}_m$, and in $m$, throughout Case~V.
\end{proposition}

\begin{proof}
We verify the hypotheses of the additive bound (1) and then collect the remaining members of
the witness.

\emph{The one-cell live bound.} By hypothesis~(2), the bound
\eqref{8.18:eq-witness-below-horizon-1} holds in Case~V, under the ceiling
\eqref{8.18:eq-degree-comparison-a}, at every level $j\le K$ of the run. It
is used twice: by the additive bound at level $j=k_0$, and by the Gaussian extraction at the
levels $j\le k_0+r(s)$. The inequality \eqref{8.18:eq-witness-below-horizon-1} does not involve
the age of any large-field region. We recall, for the uniformity statement below, where
hypothesis~(2) comes from. The corollary in hypothesis~(2) is obtained from the base lemma of the
rarity argument in its Case~V form, reached through the degree comparison on small tori and valid
under its ceiling \eqref{8.18:eq-degree-comparison-a}; its inputs are the following:
\begin{itemize}
\item the one-cell proposition, proved as an implication from the inputs of the large-field
block bound, healed part, together with clause~(i-d) of Theorem~\ref{3.1:thm-main};
\item clause~(i-d) itself;
\item the lower bound in clause~(i-b);
\item a further lemma of the rarity argument;
\item rarity of young large-field blocks;
\item the lemma on the running couplings.
\end{itemize}
None of these inputs depends on the volume in Case~V. The bound
\eqref{8.18:eq-witness-below-horizon-1} is also what the age-slice device of
Proposition~\ref{8.18:prop-age-antecedent} gives when no age threshold is imposed (all ages
summed), combined with the same one-cell proposition.

\emph{The age antecedent.} By Proposition~\ref{8.18:prop-age-antecedent}, in Case~V there is
a constant $A_{\mathrm{lv}}$ such that
\begin{equation*}
\sum_{a>\mathfrak{l}}L^{8(a+1)}q^{=a}_{k_0}\le A_{\mathrm{lv}}\,L^{8(\mathfrak{l}+1)}\,
q^{\mathrm{lv}}_{k_0}.
\end{equation*}
This bound involves the same single one-cell weight $q^{\mathrm{lv}}_{k_0}$ as
\eqref{8.18:eq-witness-below-horizon-1} at $j=k_0$, taken with $c_{\mathrm{lv}}=\theta/2$.
Both antecedents of the additive bound therefore hold in Case~V. They hold with one weight
$q^{\mathrm{lv}}_{k_0}$ and a constant $A_{\mathrm{lv}}$.

\emph{The ceilings of the additive bound.} The additive bound has three smallness conditions.
The first is its source ceiling. In it, the source-smallness function of the additive bound is
evaluated at $c_{\mathrm{lv}}=\theta/2$, and the factor $(1+A_{\mathrm{lv}})$ enters its total
source exponent. With this factor inserted, the total source exponent remains of the same kind
as before, a polynomial in $\log x_{k_0}$ of degree of order $8L\check{R}_{k_0}$, which the
ceiling compares with the exponent $p_0$ of the degree function. The second is a further ceiling
on the reference
coupling $g$, and the third is its condition on the constant $C_0$. All three are already among
the conditions of clause~(iv). Hence hypothesis~(1) applies and gives
$|R^{\mathrm{BAD}\prime}|\le\mathfrak{M}/32$.

\emph{The other members.} The decomposition of $S^T_{2;K}-\mathfrak{M}$ has five members besides
$R^{\mathrm{BAD}\prime}$: $R^{\mathrm{prof}}$, $R^{\mathrm{loc}}$, $R^{\mathrm{GX}}$,
$R^{\mathrm{irr}}$ and the shift member, which vanishes. They, together with the auxiliary
quantities entering the additive bound, are controlled as follows.
\begin{itemize}
\item $R^{\mathrm{prof}}$ and $R^{\mathrm{loc}}$ are bounded by the correlation theorem,
hypothesis~(3), applied with its torus exponent equal to $s$.
\item $R^{\mathrm{GX}}$ and $R^{\mathrm{irr}}$ are bounded by parts (a) and (c) of the Gaussian
extraction, together with its part (b), hypothesis~(4). These are applied with
\eqref{8.18:eq-witness-below-horizon-1} at the levels $j\le k_0+r(s)$.
\item The two-sided bounds on $\mathfrak{M}$ are hypothesis~(5).
\item The running shift $\zeta'$ is controlled by clause~(iii-c), hypothesis~(6).
\item $N_{k_0}$ is bounded by hypothesis~(7).
\end{itemize}
None of these statements depends on the age of a large-field region or on the volume. With
these members bounded in modulus and $|R^{\mathrm{BAD}\prime}|\le\mathfrak{M}/32$,
hypothesis~(1) gives \eqref{8.18:eq-witness-below-horizon-2} and the display of clause~(iv).

\emph{Uniformity.} No constant used above depends on $m$. The only quantity depending on $m$
that enters the argument, through Proposition~\ref{8.18:prop-age-antecedent}, is the side
$N_T^{\dagger}$ of the last lattice of the prolonged run. By
Lemma~\ref{8.18:lem-horizon-cases}, in Case~V it satisfies
\begin{equation*}
N_T^{\dagger}=2L^{\hat{\kappa}-k^{\dagger}}\le 2L^{c_{\mathrm{gap}}}.
\end{equation*}
Therefore \eqref{8.18:eq-witness-below-horizon-2} and the display of clause~(iv) hold for all
$g$, $g_0$, $\mathsf{d}$ and $(f_1,f_2)$ as stated, on every torus of Case~V. The bounds are
uniform in $K$, $g_0$, the position of the supports, and $m$.
\end{proof}

Throughout, $s_{\mathrm{vol}}(g,m)$ denotes the polylogarithmic depth floor $s_{\mathrm{vol}}^{(\mathfrak{P})}(g,m)$ of \eqref{8.18:eq-volume-horizon-e}. Here $\mathfrak{P}(y)=A_0(\log y)^{p_0}$, so that the degree function satisfies $p_0(g_k)=\mathfrak{P}(x_k)$.

We work in the case $K+m\ge K_*$. In this case the conversion \eqref{8.18:eq-probability-conversion-a} still holds, and by \eqref{8.18:eq-horizon-cases-c} its exponent $\mathsf{B}(g^{\dagger})(N^{\dagger})^4$ is at most $\bar{V}(g,m)$. At the terminal coupling $X^{\dagger}\approx X_*$ the decaying factor $\exp(-\theta\mathfrak{P}(\cdot))$ would be a constant, and it could not compensate $\bar{V}(g,m)$. In the conversion, however, the factor $\exp(-\theta\mathfrak{P}(x_{k_0-a}))$ is evaluated at the creation level $k_0-a$. Here $k_0$ is the level at which the conversion is read and $a$ is the age of the large-field component, so that $k_0-a$ is the level at which the component was created. By \eqref{8.18:eq-members-runs-a} we have $x_{k_0-a}\ge x_{k_0}$. With $k_0=K-s$, \eqref{8.18:eq-members-runs-c} gives $\mathfrak{P}(x_{k_0})\ge\mathfrak{P}(\underline{x}_s)$, and the right-hand side grows without bound with the depth $s$. The following lemma makes this precise.

\begin{lemma}[The depth floor pays the volume factor]\label{8.18:lem-depth-floor}
Assume
\[
s\ge s_{\mathrm{vol}}^{(\mathfrak{P})}(g,m).
\]
Let $s'$ be an integer in $[s-r(s),s]$ and put $j=K-s'$. Then the first three inequalities of the chain
\begin{equation}\label{8.18:eq-depth-floor-a}
\begin{split}
\tfrac12\theta\bigl(\mathfrak{P}(x_j)-2\bigr)
&\ge\tfrac12\theta\bigl(\mathfrak{P}(\underline{x}_{s'})-2\bigr)
\ge\tfrac12\theta\bigl(\mathfrak{P}(\underline{x}_{s-r(s)})-2\bigr)\\
&\ge\bar{V}(g,m)\ge\mathsf{B}(g^{\dagger})\,(N^{\dagger})^4
\end{split}
\end{equation}
hold. Assume moreover that $K+m\ge K_*$. Let $g^{\dagger}$ and $N^{\dagger}$ be the coupling and the last-lattice side of the prolonged member in the conversion of Lemma~\ref{8.18:lem-probability-conversion}, as bounded in Lemma~\ref{8.18:lem-horizon-cases}. Then the fourth inequality holds as well, and so the whole chain holds.
\end{lemma}

\begin{proof}
We first record two facts about $\mathfrak{P}$. The function $y\mapsto\mathfrak{P}(y)=A_0(\log y)^{p_0}$ is defined on $[1,\infty)$. It is evaluated below only at its arguments $x_j$, $\underline{x}_{s'}$ and $\underline{x}_{s-r(s)}$. These are inverse squared couplings at which the degree function $p_0(\cdot)$ is read. Since $A_0>0$ and $p_0>0$, the function $\mathfrak{P}$ is nondecreasing on $[1,\infty)$. Since $\theta>0$, the map
\[
y\mapsto\tfrac12\theta\bigl(\mathfrak{P}(y)-2\bigr)
\]
is nondecreasing there as well. Fix an integer $s'\in[s-r(s),s]$ and put $j=K-s'$.

\emph{First inequality.} Apply \eqref{8.18:eq-members-runs-c} with $u=s'$. This gives
\[
\underline{x}_{s'}\le x_{K-s'}=x_j .
\]
The monotonicity just noted then yields the first inequality in \eqref{8.18:eq-depth-floor-a}.

\emph{Second inequality.} The comparison couplings $u\mapsto\underline{x}_u$ of the glossary (Notation~\ref{0.2:not-glossary}) are nondecreasing in the depth $u$. Since $s'\ge s-r(s)$, it follows that $\underline{x}_{s'}\ge\underline{x}_{s-r(s)}$. The same monotonicity then yields the second inequality.

\emph{Third inequality.} By the definition of the depth floor $s_{\mathrm{vol}}^{(\mathfrak{P})}(g,m)$ in \eqref{8.18:eq-volume-horizon-e}, the hypothesis $s\ge s_{\mathrm{vol}}^{(\mathfrak{P})}(g,m)$ gives
\[
\tfrac12\theta\bigl(\mathfrak{P}(\underline{x}_{s-r(s)})-2\bigr)\ge\bar{V}(g,m),
\]
where $\bar{V}(g,m)$ is the volume number of \eqref{8.18:eq-volume-horizon-c}.

\emph{Fourth inequality.} Assume moreover that $K+m\ge K_*$. Then the torus lies beyond the volume horizon. Lemma~\ref{8.18:lem-horizon-cases}, in the form \eqref{8.18:eq-horizon-cases-c}, gives
\[
\mathsf{B}(g^{\dagger})(N^{\dagger})^4\le\bar{V}(g,m).
\]

Chaining the first three inequalities gives the first three inequalities of \eqref{8.18:eq-depth-floor-a} for every $s'\in[s-r(s),s]$. Under the additional assumption $K+m\ge K_*$, the fourth inequality extends this chain, which gives \eqref{8.18:eq-depth-floor-a} in full.
\end{proof}

\begin{proposition}[Terminal-live slices under the polylogarithmic floor]\label{8.18:prop-polylog-floor}
Assume the following.
\begin{itemize}
\item[(a)] Lemma~\ref{8.18:lem-age-reserve-bound} (un-normalised age-and-reserve bound at the last
lattice), the conversion \eqref{8.18:eq-probability-conversion-a} of
Lemma~\ref{8.18:lem-probability-conversion}, rarity of young large-field blocks, and the
remaining statements and one-sided conditions from which
Proposition~\ref{8.18:prop-never-healed-slices} (never-healed age slices below the horizon) is
proved, with one exception: the degree comparison \eqref{8.18:eq-degree-comparison-a} is not
assumed with $X^{\dagger}$ in place of its argument $y$.
\item[(b)] Lemma~\ref{8.18:lem-depth-floor} (the depth floor pays the volume factor), in the form
\eqref{8.18:eq-depth-floor-a}.
\item[(c)] The condition \eqref{8.18:eq-horizon-cases-d} of Lemma~\ref{8.18:lem-horizon-cases}
(tori below and beyond the volume horizon).
\end{itemize}
Let $s_{\mathrm{vol}}^{(\mathfrak{P})}(g,m)$ be the volume-reading depth floor
\eqref{8.18:eq-volume-horizon-e}, and let the level $k_0$, the reserve fraction $\theta$, the
normalised measure $\tilde{\nu}$ on histories, the events $\mathfrak{lv}^{=a}_{k_0}(\square)$,
$\mathfrak{lv}_j(\mathfrak{c})$ and $\mathrm{END}^{\dagger}$, and the last level
$K^{\dagger\prime}$ of the prolonged member be as in
Proposition~\ref{8.18:prop-never-healed-slices} and
Lemma~\ref{8.18:lem-probability-conversion}. For a level $j$ and a level-$j$ cell
$\mathfrak{c}$, let $\mathrm{END}^{\dagger}_j$ be the event that the live component containing
$\mathfrak{c}$ at level $j$, followed through the later levels of the prolonged member, is still
live, that is, still carries its large-field factor, at level $K^{\dagger\prime}$; this is the
analogue at level $j$ of the event $\mathrm{END}^{\dagger}$.

Then for every $m\ge m_{\mathbb{T}}(g_-(g))$, every $s$ with
$s\ge s_{\mathrm{vol}}^{(\mathfrak{P})}(g,m)$, every $a\ge 0$, every level-$k_0$ cell $\square$,
every level $j$ with $K-s\le j\le K-s+r(s)$ and every level-$j$ cell $\mathfrak{c}$,
\begin{equation}\label{8.18:eq-polylog-floor-a}
\begin{gathered}
\tilde{\nu}\bigl(\mathfrak{lv}^{=a}_{k_0}(\square)\cap\mathrm{END}^{\dagger}\bigr)
\le e^{3\theta/2}\,(2L^{9})^{-a}\,e^{-\frac12\theta\,\mathfrak{P}(x_{k_0})},\\
\tilde{\nu}\bigl(\mathfrak{lv}_{j}(\mathfrak{c})\cap\mathrm{END}^{\dagger}_j\bigr)
\le e^{3\theta/2}\,e^{-\frac12\theta\,\mathfrak{P}(x_{j})}.
\end{gathered}
\end{equation}
Both bounds hold uniformly in $K$, in $g_0$ and in the position of $\square$ and
$\mathfrak{c}$. The volume enters only through the condition
$s\ge s_{\mathrm{vol}}^{(\mathfrak{P})}(g,m)$.
\end{proposition}

\begin{proof}
Fix $m\ge m_{\mathbb{T}}(g_-(g))$ and $s\ge s_{\mathrm{vol}}^{(\mathfrak{P})}(g,m)$.

\emph{The prolonged member.} The argument runs on the prolonged member, with last level
$K^{\dagger\prime}$, last lattice $\Lambda^{\dagger}$ of side $N^{\dagger}$, and coupling
$g^{\dagger}$. Hypothesis (c), the condition \eqref{8.18:eq-horizon-cases-d} of
Lemma~\ref{8.18:lem-horizon-cases}, is used for this prolongation.

\emph{The first bound.} Fix $a\ge 0$ and a level-$k_0$ cell $\square$. We repeat the proof of
Proposition~\ref{8.18:prop-never-healed-slices} with the choice $b:=k_0-a$. Its inputs are those
listed in hypothesis (a). In particular, Lemma~\ref{8.18:lem-age-reserve-bound} is applied to
the prolonged member with $\hat{K}=K^{\dagger\prime}$, $\hat{N}=N^{\dagger}$, this $b$ and the
event $E_a(\square)\cap\mathrm{END}^{\dagger}$, where $E_a(\square)$ is the age-$a$ event of
Lemma~\ref{8.18:lem-probability-conversion}, that is, the event
$\mathfrak{lv}^{=a}_{k_0}(\square)$ that $\square$ lies in a live component of age $a$ at level
$k_0$. The division by the normalising integral then
gives the conversion \eqref{8.18:eq-probability-conversion-a}, which carries the factor
$\exp\bigl(\mathsf{B}(g^{\dagger})N^{\dagger4}\bigr)$.

That proof uses the degree comparison \eqref{8.18:eq-degree-comparison-a} with $X^{\dagger}$ in
place of its argument $y$
at one point: to pay this factor out of the first half of the reserve. Here that use is
replaced by hypothesis (b). Since $s\ge s_{\mathrm{vol}}^{(\mathfrak{P})}(g,m)$,
Lemma~\ref{8.18:lem-depth-floor} gives, for every $s'\in[s-r(s),s]$ and every level $i=K-s'$,
\[
\tfrac12\theta\bigl(\mathfrak{P}(x_i)-2\bigr)\ge \bar{V}(g,m)\ge \mathsf{B}(g^{\dagger})N^{\dagger4}.
\]
So the first half of the reserve pays $\exp\bigl(\mathsf{B}(g^{\dagger})N^{\dagger4}\bigr)$
without any use of \eqref{8.18:eq-degree-comparison-a} at $y=X^{\dagger}$. All other steps of
the proof of Proposition~\ref{8.18:prop-never-healed-slices} are taken over unchanged. Their
conclusion is the first line of \eqref{8.18:eq-polylog-floor-a}.

\emph{The second bound.} Fix a level $j$ with $K-s\le j\le K-s+r(s)$ and a level-$j$ cell
$\mathfrak{c}$. We repeat the same argument with $b:=j$ and $a:=0$, applied to the event
$\mathfrak{lv}_j(\mathfrak{c})\cap\mathrm{END}^{\dagger}_j$. Put $s':=K-j$. Then
$s-r(s)\le s'\le s$, so $j=K-s'$ is one of the levels covered by
\eqref{8.18:eq-depth-floor-a}. The volume factor is therefore paid by
Lemma~\ref{8.18:lem-depth-floor} exactly as before. With $a=0$ the factor $(2L^9)^{-a}$ equals
one, and the argument yields the second line of \eqref{8.18:eq-polylog-floor-a}.

\emph{Tori below the horizon.} Suppose the torus lies in the first case of
Lemma~\ref{8.18:lem-horizon-cases} for the actual $K$, that is, $K+m\le K_*-1$. Then
Proposition~\ref{8.18:prop-never-healed-slices} applies as stated. With the substitutions
above, its proof gives \eqref{8.18:eq-polylog-floor-a} without appeal to the floor
$s\ge s_{\mathrm{vol}}^{(\mathfrak{P})}(g,m)$. That condition is then not needed, and assuming
it does not affect the argument.
\end{proof}

\begin{proposition}[Terminal-live slices under the linear floor]\label{8.18:prop-linear-floor}
Assume the inputs of Proposition~\ref{8.18:prop-never-healed-slices} and the further conditions listed in (f),
namely:
\begin{itemize}
\item[(a)] the un-normalised age-and-reserve bound at the last lattice, Lemma~\ref{8.18:lem-age-reserve-bound},
in its first form \eqref{8.18:eq-age-reserve-bound-a};
\item[(b)] the lower bound on the normalising integral used in Lemma~\ref{8.18:lem-probability-conversion},
and the degree bound on $\mathsf{B}$ of Definition~\ref{8.18:def-volume-horizon};
\item[(c)] part~(b) of the re-targeting lemma for runs, by which an attainable level $k^{\dagger}$ yields the
prolonged member $(\dagger)$ over the same lattice, with last level $K^{\dagger\prime}\ge K+(k^{\dagger}-K_V)$;
together with Lemmas~\ref{8.18:lem-run-prolongation} and~\ref{8.18:lem-horizon-cases};
\item[(d)] rarity of young large-field blocks, and the transfer of events and of the normalised history
measure from the prolonged run $(\dagger)$ to the original run used in the proof of
Proposition~\ref{8.18:prop-never-healed-slices};
\item[(e)] the display \eqref{8.18:eq-members-runs-a};
\item[(f)] the ceiling $\gamma_*\le\gamma_L(c_a)$ on the reference coupling, the ceiling $(\gamma_V)^{\flat}$ of
Proposition~\ref{8.18:prop-never-healed-slices}, the window on $\theta$, the floor
\eqref{8.18:eq-never-healed-slices-a} on $c_a$, the ceiling \eqref{8.18:eq-horizon-cases-d}, and the other
ceilings on the reference coupling under which Proposition~\ref{8.18:prop-polylog-floor} holds.
\end{itemize}
Let $m\ge m_{\mathbb{T}}(g_-(g))$, and let $s\ge s_{\mathrm{vol}}(g,m)$, with $s_{\mathrm{vol}}(g,m)$ in its linear
form \eqref{8.18:eq-volume-horizon-d}, so that $s$ satisfies the linear floor
\begin{equation}\label{8.18:eq-linear-floor-1}
c_a\bigl(s-r(s)+p_C(g)\bigr)\;\ge\;\bar{V}(g,m)+\theta .
\end{equation}
Put $k_0=K-s$. Then for every $a\ge0$ and every level-$k_0$ cell $\square$
\begin{equation}\label{8.18:eq-linear-floor-2}
\tilde{\nu}\bigl(\mathfrak{lv}^{=a}_{k_0}(\square)\cap\mathrm{END}^{\dagger}\bigr)
\le e^{3\theta/2}\,(2L^9)^{-a}\,e^{-\frac12\theta\mathfrak{P}(x_{k_0})},
\end{equation}
and for every level $j\in[K-s,\,K-s+r(s)]$ and every level-$j$ cell $\mathfrak{c}$
\begin{equation}\label{8.18:eq-linear-floor-3}
\tilde{\nu}\bigl(\mathfrak{lv}_j(\mathfrak{c})\cap\mathrm{END}^{\dagger}_j\bigr)
\le e^{3\theta/2}\,e^{-\frac12\theta\mathfrak{P}(x_j)},
\end{equation}
uniformly in $K$, $g_0$ and the position of the cell. These are the displays
\eqref{8.18:eq-polylog-floor-a} of Proposition~\ref{8.18:prop-polylog-floor}.
\end{proposition}

\begin{proof}
The inputs, the ceilings, and the place of every constant in the order of choice are those of
Proposition~\ref{8.18:prop-polylog-floor}. Below the volume horizon, $K+m\le K_*-1$, the bounds are obtained as in
Proposition~\ref{8.18:prop-never-healed-slices}, whose hypotheses, including the ceiling $(\gamma_V)^{\flat}$, are
among those assumed in (a)--(f). The argument differs
only in the case of Lemma~\ref{8.18:lem-horizon-cases} beyond the volume horizon,
\eqref{8.18:eq-horizon-cases-c}, which we now treat.

\emph{Step 1: length of the prolongation.} Let $k^{\dagger}$, $K_V$, $g_V$ and the prolonged member $(\dagger)$ be
as in Lemma~\ref{8.18:lem-run-prolongation}, with $\hat{\kappa}=\min\{K+m,K_*-1\}$. In the present case
$\hat{\kappa}=K_*-1$, and Lemma~\ref{8.18:lem-run-prolongation} gives $\hat{\kappa}-k^{\dagger}\le c_{\mathrm{gap}}$,
that is $k^{\dagger}\ge K_*-1-c_{\mathrm{gap}}$. Since $K_V=K+m_{\mathbb{T}}(g_V)$
(Lemma~\ref{8.18:lem-run-prolongation}), this gives
\begin{equation}\label{8.18:eq-linear-floor-4}
k^{\dagger}-K_V\;\ge\;(K_*-1-c_{\mathrm{gap}})-\bigl(K+m_{\mathbb{T}}(g_V)\bigr)
\;\ge\;m_V(g)-c_{\mathrm{gap}}-m_{\mathbb{T}}(g_-(g)),
\end{equation}
the second inequality using $K_*-K\ge m_V(g)+1$ from Lemma~\ref{8.18:lem-horizon-cases} and the comparison of the
torus exponents $m_{\mathbb{T}}(g_V)\le m_{\mathbb{T}}(g_-(g))$; here and below the comparisons
$m_{\mathbb{T}}(g^{\dagger})\le m_{\mathbb{T}}(g_V)\le m_{\mathbb{T}}(g_-(g))$ hold because
$m_{\mathbb{T}}(g')=11+m'+\log_L R(g')$ is non-increasing in $g'$ and $g_-(g)\le g_V\le g^{\dagger}$.
By the definition of $p_C(g)$ in
\eqref{8.18:eq-volume-horizon-d}, \eqref{8.18:eq-linear-floor-4} gives $k^{\dagger}-K_V\ge p_C(g)$. Part~(b) of the
re-targeting lemma for runs then gives, for the last level $K^{\dagger\prime}$ of the prolonged run,
\begin{equation*}
K^{\dagger\prime}\;\ge\;K+(k^{\dagger}-K_V)\;\ge\;K+p_C(g).
\end{equation*}
Since $k_0+r(s)=K-(s-r(s))$, every level $j\le k_0+r(s)$ satisfies
\begin{equation}\label{8.18:eq-linear-floor-5}
K^{\dagger\prime}-j\;\ge\;\bigl(s-r(s)\bigr)+p_C(g).
\end{equation}
For the side of the last lattice, using $m_{\mathbb{T}}(g^{\dagger})\le m_{\mathbb{T}}(g_V)$ and
$K_V=K+m_{\mathbb{T}}(g_V)$,
\begin{equation*}
\begin{aligned}
N^{\dagger}=2L^{K+m-k^{\dagger}+m_{\mathbb{T}}(g^{\dagger})}
&\le 2L^{K+m-k^{\dagger}+m_{\mathbb{T}}(g_V)}=2L^{m-(k^{\dagger}-K_V)}\\
&\le 2L^{m-p_C(g)}.
\end{aligned}
\end{equation*}
Hence $(N^{\dagger})^4\le 16L^{4(m-p_C(g))}$. By the definition of $p_C(g)$ in
\eqref{8.18:eq-volume-horizon-d}, $p_C(g)\ge m_V(g)-c_{\mathrm{gap}}-m_{\mathbb{T}}(g_-(g))$, so that
\begin{equation*}
4\bigl(m-p_C(g)\bigr)\;\le\;4\bigl(m-m_V(g)+c_{\mathrm{gap}}+m_{\mathbb{T}}(g_-(g))\bigr).
\end{equation*}
Moreover $\log X^{\dagger}\le\log(g^{-2}+B^{\sharp})$, the terminal inverse coupling $X^{\dagger}$ of
$(\dagger)$, as in Lemma~\ref{8.18:lem-horizon-cases}, being at most $g_-(g)^{-2}=g^{-2}+B^{\sharp}$. Together with
the degree bound $\mathsf{B}(g^{\dagger})\le C_{\mathsf{B}}(1+\log X^{\dagger})$ of
Definition~\ref{8.18:def-volume-horizon}, the definition \eqref{8.18:eq-volume-horizon-c} of $\bar{V}(g,m)$ gives
$\bar{V}(g,m)\ge\mathsf{B}(g^{\dagger})(N^{\dagger})^4$.

\emph{Step 2: the age-and-reserve bound, keeping the factor $q^{K^{\dagger\prime}-b}$.} Apply the first form
\eqref{8.18:eq-age-reserve-bound-a} of Lemma~\ref{8.18:lem-age-reserve-bound} to the member $(\dagger)$, with
$\hat{K}=K^{\dagger\prime}$, $\hat{N}=N^{\dagger}$, and with $b:=k_0-a$ and
$E:=E_a(\square)\cap\mathrm{END}^{\dagger}$ as in Lemma~\ref{8.18:lem-probability-conversion} for the first
display, respectively with $b:=j$ and $E$ the event $\mathfrak{lv}_j(\mathfrak{c})\cap\mathrm{END}^{\dagger}_j$
of the second display, $j\in[K-s,K-s+r(s)]$. For the corresponding event $E$ it gives
\begin{equation*}
\sum_{\omega\in E}\tau_{\omega}\;\le\;e^{\theta}\,q^{K^{\dagger\prime}-b}\,e^{-\theta\mathfrak{P}(x_b)}\,
e^{E_+^{(\theta)}(N^{\dagger})^4},
\end{equation*}
where $q^{K^{\dagger\prime}-b}=q^a\,q^{K^{\dagger\prime}-k_0}$ in the first case and
$q^{K^{\dagger\prime}-b}=q^{K^{\dagger\prime}-j}$ in the second. Since $q=e^{-c_a}\le1$ and $k_0\le j$,
\eqref{8.18:eq-linear-floor-5} gives
\begin{equation*}
q^{K^{\dagger\prime}-k_0}\;\le\;q^{K^{\dagger\prime}-j}\;\le\;e^{-c_a(s-r(s)+p_C(g))}.
\end{equation*}

\emph{Step 3: the division.} Divide by the normalising integral as in
Lemma~\ref{8.18:lem-probability-conversion}, but keep the factor $q^{K^{\dagger\prime}-b}$ instead of bounding
$q^{K^{\dagger\prime}-k_0}$ by $1$; the division by $Z\ge e^{-C_{\mathrm{low}}(g^{\dagger})(N^{\dagger})^4}$ turns
the factor $e^{E_+^{(\theta)}(N^{\dagger})^4}$ of Step~2 into the quotient
$e^{\mathsf{B}(g^{\dagger})(N^{\dagger})^4}$, $\mathsf{B}=E_+^{(\theta)}+C_{\mathrm{low}}$, and the passage from the
measure $\tilde{\nu}^{\dagger}$ of the prolonged run to $\tilde{\nu}$ is that of the proof of
Proposition~\ref{8.18:prop-never-healed-slices}. By \eqref{8.18:eq-linear-floor-1} and Step~1,
\begin{equation*}
c_a\bigl(s-r(s)+p_C(g)\bigr)\;\ge\;\bar{V}(g,m)+\theta\;\ge\;\mathsf{B}(g^{\dagger})(N^{\dagger})^4+\theta,
\end{equation*}
so the factor $q^{K^{\dagger\prime}-k_0}$ (respectively $q^{K^{\dagger\prime}-j}$) alone compensates the quotient
$e^{\mathsf{B}(g^{\dagger})(N^{\dagger})^4}$:
$e^{-c_a(s-r(s)+p_C(g))}e^{\mathsf{B}(g^{\dagger})(N^{\dagger})^4}\le e^{-\theta}$. Hence
\begin{equation*}
\begin{aligned}
\tilde{\nu}\bigl(\mathfrak{lv}^{=a}_{k_0}(\square)\cap\mathrm{END}^{\dagger}\bigr)
&\le e^{\theta}\,q^a\,e^{-\theta\mathfrak{P}(x_{k_0-a})}\,e^{-\theta}
\le e^{\theta}\,(2L^9)^{-a}\,e^{-\theta\mathfrak{P}(x_{k_0})}\\
&\le e^{3\theta/2}\,(2L^9)^{-a}\,e^{-\frac12\theta\mathfrak{P}(x_{k_0})},
\end{aligned}
\end{equation*}
where the second inequality uses $e^{-\theta}\le1$, $q\le(2L^9)^{-1}$, which is the floor
\eqref{8.18:eq-never-healed-slices-a}, and \eqref{8.18:eq-members-runs-a}; the third uses $\theta\ge0$ and
$\mathfrak{P}(x_{k_0})\ge0$. This is \eqref{8.18:eq-linear-floor-2}. For the second display, at level $j$, the
same steps give
\begin{equation*}
\tilde{\nu}\bigl(\mathfrak{lv}_j(\mathfrak{c})\cap\mathrm{END}^{\dagger}_j\bigr)
\le e^{\theta}\,e^{-\theta\mathfrak{P}(x_j)}\,e^{-\theta}
\le e^{3\theta/2}\,e^{-\frac12\theta\mathfrak{P}(x_j)},
\end{equation*}
which is \eqref{8.18:eq-linear-floor-3}. Here no part of the reserve $e^{-\theta\mathfrak{P}(x_{k_0-a})}$ is spent
on the quotient; the last inequality only weakens it to the form of \eqref{8.18:eq-polylog-floor-a}.
Since the left member of \eqref{8.18:eq-linear-floor-1} is linear in $s$, the linear form of
$s_{\mathrm{vol}}(g,m)$ is of size $\bar{V}(g,m)/c_a$, against $\exp\bigl((\bar{V}/A_0)^{1/p_0}\bigr)$ for the floor
of Proposition~\ref{8.18:prop-polylog-floor}.
\end{proof}

\begin{proposition}[Clause (iv) beyond the volume horizon]\label{8.18:prop-witness-beyond-horizon}
Let the constants, the couplings $g_0$, $g$ and the test functions be as in clause~(iv) of
Theorem~\ref{3.1:thm-main}, with the reference coupling under the ceilings of that clause.
Write $s=s(\mathsf{d})$ for the witness depth and $k_0=K-s$ for
the witness level. Let $m_V(g)$ be the volume horizon of
Definition~\ref{8.18:def-volume-horizon}, and let $s_{\mathrm{vol}}(g,m)$ be the volume-reading
depth floor \eqref{8.18:eq-volume-horizon-d}. Assume:
\begin{itemize}
\item the healed old-region bound, applied to the prolonged run, so that
Lemma~\ref{8.18:lem-healed-slices} yields the healed-later age slices
\eqref{8.18:eq-healed-slices-a};
\item the one-point healed bound of the large-field block bound, healed part, at every level of
every run of the family;
\item the live-member theorem, which bounds the live member $R^{\mathrm{BAD}\prime}$ of the
witness, and with it clause~(iv)'s display, from the one-cell live bound at level $k_0$ and the
age bound at level $k_0$;
\item the correlation theorem, applied with its torus exponent equal to the witness depth $s$;
\item the Gaussian extraction, its clauses (a), (b) and~(c);
\item the kernel statement for the witness kernel $\mathcal{K}_s$;
\item clause~(iii-c) of Theorem~\ref{3.1:thm-main};
\item the lemma on rarity of young large-field blocks, which enters the bounds of the five
members other than $R^{\mathrm{BAD}\prime}$ in the proof of
Proposition~\ref{8.18:prop-witness-below-horizon}.
\end{itemize}
Then for every torus $\mathbb{T}_m$ with $m\ge m_{\mathbb{T}}(g_-(g))$ and every separation
$\mathsf{d}$ with
\begin{equation}\label{8.18:eq-witness-beyond-horizon-1}
\max\{s_W(g),\,s_{\mathrm{vol}}(g,m)\}\le s(\mathsf{d})
\end{equation}
and satisfying the ultraviolet condition on $s(\mathsf{d})$ imposed in clause~(iv) of
Theorem~\ref{3.1:thm-main}, the display of
clause~(iv) holds on $\mathbb{T}_m$, that is, the two-point witness on every torus
\eqref{8.18:eq-torus-witness-a} holds, uniformly in $K$, in $g_0$ and in the position of the
supports of the test functions. Condition \eqref{8.18:eq-witness-beyond-horizon-1} is the
condition $\mathsf{d}\le\mathsf{d}_W^{\mathbb{T}}(g,m)$ on the separation, where
\begin{equation}\label{8.18:eq-witness-beyond-horizon-2}
\mathsf{d}_W^{\mathbb{T}}(g,m):=\text{the largest }\mathsf{d}\text{ with }
s(\mathsf{d})\ge\max\{s_W(g),\,s_{\mathrm{vol}}(g,m)\};
\end{equation}
this threshold involves the volume, through $s_{\mathrm{vol}}(g,m)$, when $m>m_V(g)$, and only
then.
\end{proposition}

\begin{proof}
Clause~(iv)'s display is obtained, as in the proof of
Proposition~\ref{8.18:prop-witness-below-horizon}, from two bounds on the normalised measure
$\tilde{\nu}$ on histories: the one-cell live bound
\begin{equation}\label{8.18:eq-witness-beyond-horizon-3}
\tilde{\nu}\bigl(\mathfrak{lv}_j(\square)\bigr)\le q^{\mathrm{lv}}_j
\quad\text{for every level } j\le k_0+r(s) \text{ and every level-}j\text{ cell }\square,
\end{equation}
and the age bound at the witness level
\begin{equation}\label{8.18:eq-witness-beyond-horizon-4}
\sum_{a>\mathfrak{l}}L^{8(a+1)}q^{=a}_{k_0}
\le A_{\mathrm{lv}}\,L^{8(\mathfrak{l}+1)}\,q^{\mathrm{lv}}_{k_0}.
\end{equation}
We establish both on every torus with $m\ge m_{\mathbb{T}}(g_-(g))$ under
\eqref{8.18:eq-witness-beyond-horizon-1}.

\emph{The age bound.} Lemma~\ref{8.18:lem-healed-slices}, which is the healed old-region bound
applied to the prolonged run, bounds the healed-later age slices
\eqref{8.18:eq-healed-slices-a}. That bound holds at every level up to the last level of the
run, in particular at $k_0$; it carries no depth floor, and the volume does not enter it. The
terminal-live age slices are bounded by \eqref{8.18:eq-polylog-floor-a} under the depth floor
$s(\mathsf{d})\ge s_{\mathrm{vol}}(g,m)$, which is part of
\eqref{8.18:eq-witness-beyond-horizon-1}. Every live component of age $a$ at level $k_0$ is
either healed at a later level of the prolonged run or still live at its last level. Write
$q^{\mathrm{heal}\dagger,=a}_{k_0}$ and $q^{\mathrm{END}\dagger,=a}_{k_0}$ for the weights of
the two parts into which this dichotomy splits the age-$a$ event $E_a(\square)$ at level $k_0$;
accordingly
\begin{equation*}
q^{=a}_{k_0}\le q^{\mathrm{heal}\dagger,=a}_{k_0}+q^{\mathrm{END}\dagger,=a}_{k_0},
\end{equation*}
the first term being the healed-later slice and the second the terminal-live slice. Adding the
two bounds gives
\eqref{8.18:eq-witness-beyond-horizon-4}, with $A_{\mathrm{lv}}$ the sum of their constants.

\emph{The one-cell bound.} Fix a level $j\le k_0+r(s)$. Every live component containing
$\square$ at level $j$ is either healed at a later level of the prolonged run or live at its
last level, so $\tilde{\nu}(\mathfrak{lv}_j(\square))$ is at most the sum of the weights of a
healed part and a terminal-live part. The healed part of the left side
of \eqref{8.18:eq-witness-beyond-horizon-3} is bounded by the one-point healed bound of the
large-field block bound, healed part. That bound holds at every level of every run of the family
and does not depend on the volume. The terminal-live part is bounded by the second display of the
statement on terminal-live slices under the polylogarithmic floor, whose first display is
\eqref{8.18:eq-polylog-floor-a}; that second display bounds, at levels $j\le k_0+r(s)$, the
weight of the event that $\square$ lies in a component still live at the last level of the
prolonged run, again under the depth floor of \eqref{8.18:eq-witness-beyond-horizon-1}. Adding
the two gives \eqref{8.18:eq-witness-beyond-horizon-3}. Beyond the horizon the one-cell bound
\eqref{8.18:eq-witness-beyond-horizon-3} therefore needs no input other than these two bounds.

Alternatively, \eqref{8.18:eq-witness-beyond-horizon-3} follows from the corollary on rarity of
young large-field blocks in its case of tori beyond the volume horizon, under that corollary's
further hypotheses and without the floor $s_{\mathrm{vol}}(g,m)$; either route suffices.

\emph{Assembly.} With \eqref{8.18:eq-witness-beyond-horizon-3} at $j=k_0$ and
\eqref{8.18:eq-witness-beyond-horizon-4}, the live-member theorem gives
$|R^{\mathrm{BAD}\prime}|\le\mathfrak{M}/32$ under the ceilings of clause~(iv); its standing
hypotheses hold on $\mathbb{T}_m$ exactly as verified in the proof of
Proposition~\ref{8.18:prop-witness-below-horizon}, the volume not entering them. The five
remaining members are bounded exactly as in the proof of
Proposition~\ref{8.18:prop-witness-below-horizon}. There the main term $\mathfrak{M}$ and the
members $R^{\mathrm{prof}}$ and $R^{\mathrm{loc}}$ involve no event on histories; they are
controlled by the kernel statement
for $\mathcal{K}_s$ and by the correlation theorem applied with its torus exponent equal to $s$.
The members $R^{\mathrm{GX}}$ and $R^{\mathrm{irr}}$ are
controlled by the Gaussian extraction, clauses (a), (b) and~(c), and depend on the measure only
through the one-cell weight \eqref{8.18:eq-witness-beyond-horizon-3} at levels
$j\le k_0+r(s)$, which has just been established; and the shift member vanishes identically, as
in that proof. Clause~(iii-c) of Theorem~\ref{3.1:thm-main}
and the lemma on rarity of young large-field blocks enter there as they do in that proof. The
six modulus bounds combine to
\begin{equation*}
|S^T_{2;K}-\mathfrak{M}|\le\tfrac{6}{32}\,\mathfrak{M},\qquad |R_K|\le\tfrac14,
\end{equation*}
and this is clause~(iv)'s display on $\mathbb{T}_m$. Each input is uniform in $K$, in $g_0$
and in the position of the supports. The volume enters only through the depth floor
$s_{\mathrm{vol}}(g,m)$.

\emph{The threshold.} By \eqref{8.18:eq-witness-beyond-horizon-2} and the definition
\eqref{8.18:eq-volume-horizon-d} of the volume-reading depth floor, the volume enters the
threshold $\mathsf{d}_W^{\mathbb{T}}(g,m)$, through $s_{\mathrm{vol}}(g,m)$, when $m>m_V(g)$,
and only then.
\end{proof}

\begin{theorem}[The two-point witness on every torus]\label{8.18:thm-torus-witness}
Let $N$, $L$, $\mathfrak{p}$ be as in clause (iv) of Theorem~0 (Theorem~\ref{3.1:thm-main}). The test
class $\mathfrak{T}(\mathsf{d};\vartheta)$, the numbers $b_0$, $s(\mathsf{d})$, $s_W(g)$, $r(s)$,
$k_W$, $C_W$, $c_{\mathcal{K}}$, $C_{\mathcal{K}}$, the kernel $\mathcal{K}_s$ and the running
shift $\zeta'$ are those of clause (iv) of Theorem~\ref{3.1:thm-main}, unchanged. The volume horizon
$m_V(g)$ and the volume-reading depth floor $s_{\mathrm{vol}}(g,m)$ are those of
Definition~\ref{8.18:def-volume-horizon}, and $\gamma_W^{\mathbb{T}}\le\gamma_W$ is the threshold
fixed in \eqref{8.18:eq-witness-below-horizon-a}. None of $\gamma_W^{\mathbb{T}}$, $C_W$, $k_W$,
$c_{\mathcal{K}}$, $C_{\mathcal{K}}$ depends on $K$, on $g_0$, on the position of the supports of
the test functions, or on $m$.

Assume the hypotheses of Lemma~\ref{8.18:lem-horizon-cases},
Proposition~\ref{8.18:prop-witness-below-horizon} and
Proposition~\ref{8.18:prop-witness-beyond-horizon}, as listed there, together with clause (0) of
Theorem~0, reflection positivity of the lattice measure and the continuum kernel. For the
consequences in part~(d) below, assume in addition clause (ii-$\mathbb{T}$) of Theorem~0.

Consider the following one-sided conditions on $g$, $m$, $\mathsf{d}$ and $K$:
\begin{align}
&g\le\gamma_W^{\mathbb{T}},\label{8.18:eq-torus-witness-1}\\
&\max\{s_W(g),\,s_{\mathrm{vol}}(g,m)\}\le s(\mathsf{d}),\label{8.18:eq-torus-witness-a}\\
&K-s(\mathsf{d})\ge k_W,\label{8.18:eq-torus-witness-2}\\
&m\ge m_{\mathbb{T}}(g_-(g)).\label{8.18:eq-torus-witness-3}
\end{align}
Since $s(\mathsf{d})$ is non-increasing in $\mathsf{d}$, condition \eqref{8.18:eq-torus-witness-a}
is equivalent to
\begin{equation*}
\mathsf{d}\le\mathsf{d}_W^{\mathbb{T}}(g,m):=
\sup\bigl\{\mathsf{d}:\ \max\{s_W(g),\,s_{\mathrm{vol}}(g,m)\}\le s(\mathsf{d})\bigr\}.
\end{equation*}
Since $s_{\mathrm{vol}}(g,m)=0$ for $m\le m_V(g)$ by \eqref{8.18:eq-volume-horizon-d}, on those tori
condition \eqref{8.18:eq-torus-witness-a} coincides with the depth
condition $s_W(g)\le s(\mathsf{d})$ of clause (iv). Condition \eqref{8.18:eq-torus-witness-3} is the
torus clause of hypothesis (H5) (Notation~\ref{2.2:not-hypotheses}), verbatim; nothing more is
required of $m$.

Let $g$ satisfy \eqref{8.18:eq-torus-witness-1}, let $g_0\in\,]0,g_-(g)]$ and $K=K(g_0,g)$, let $m$
satisfy \eqref{8.18:eq-torus-witness-3}, let $\mathsf{d}$ satisfy \eqref{8.18:eq-torus-witness-a}
and \eqref{8.18:eq-torus-witness-2}, put $s:=s(\mathsf{d})$, and let
$(f_1,f_2)\in\mathfrak{T}(\mathsf{d};\vartheta)$. Then the following hold.
\begin{itemize}
\item[(a)] On the torus $\mathbb{T}_m$,
\begin{equation}\label{8.18:eq-torus-witness-4}
S^T_{2;K}(f_1,f_2)=(\zeta'_{K-s})^2\,b_0\,g_{K-s}^4\,\mathcal{K}_s(f_1,f_2)\,
\bigl(1+R_K(f_1,f_2)\bigr),
\end{equation}
where
\begin{equation}\label{8.18:eq-torus-witness-5}
|R_K(f_1,f_2)|\le\tfrac12\ \bigl(\text{indeed }\le\tfrac14\bigr),\qquad
\zeta'_{K-s}\in\bigl[\tfrac23,\tfrac43\bigr],
\end{equation}
and
\begin{equation}\label{8.18:eq-torus-witness-6}
c_{\mathcal{K}}\|f_1\|_\infty\|f_2\|_\infty\le\mathcal{K}_s(f_1,f_2)
\le C_{\mathcal{K}}\|f_1\|_\infty\|f_2\|_\infty .
\end{equation}
\item[(b)] Consequently
\begin{equation}\label{8.18:eq-torus-witness-7}
\tfrac29\,b_0\,g_{K-s}^4\,\mathcal{K}_s(f_1,f_2)\le S^T_{2;K}(f_1,f_2)
\le\tfrac83\,b_0\,g_{K-s}^4\,\mathcal{K}_s(f_1,f_2),
\end{equation}
and, by clause (0) of Theorem~0,
\begin{equation}\label{8.18:eq-torus-witness-8}
\begin{split}
\tfrac29\,b_0\,c_{\mathcal{K}}\,\|f_1\|_\infty\|f_2\|_\infty
\bigl(g^{-2}+B^{\sharp}(s+1)\bigr)^{-2}
&\le S^T_{2;K}(f_1,f_2)\\
&\le\tfrac83\,b_0\,C_{\mathcal{K}}\,\|f_1\|_\infty\|f_2\|_\infty
\bigl(g^{-2}+\lambda^{\sharp}s-c_5\bigr)^{-2}.
\end{split}
\end{equation}
\item[(c)] These statements hold uniformly in $K$, $g_0$ and the position of the supports, on every
torus $\mathbb{T}_m$ with $m\ge m_{\mathbb{T}}(g_-(g))$. They hold uniformly also in $m$ on
$\{m_{\mathbb{T}}(g_-(g))\le m\le m_V(g)\}$. For $m>m_V(g)$ the volume enters through
$s_{\mathrm{vol}}(g,m)$ and through nothing else.
\item[(d)] For every torus limit of clause (ii-$\mathbb{T}$-a), (ii-$\mathbb{T}$-b) of Theorem~0, on
every torus of the family:
\begin{itemize}
\item the smeared curvature has strictly positive variance;
\item the limit is not scale invariant under any assignment of a dimension to the curvature;
\item it is reproduced by no free Maxwell field at any charge;
\item the coupling read off the two-point function at separation $\mathsf{d}$ is positive and tends
to $0$ logarithmically as $\mathsf{d}\downarrow0$, at the two-sided rate of
\eqref{8.18:eq-torus-witness-8}.
\end{itemize}
Each of these is a statement at fixed $m$, either at one small $\mathsf{d}$ or as
$\mathsf{d}\downarrow0$.
\end{itemize}

Not asserted: anything about the limits on $\mathbb{R}^4$ of clause (ii-$\mathbb{R}^4$) of
Theorem~0. In particular, it is not asserted that the curvature Hilbert space
$\mathcal{H}^{\mathrm{curv}}$ of those limits has dimension at least two; the bounds above at fixed
$\mathsf{d}$ uniformly in $m$ beyond $m_V(g)$ are not asserted. It
is further not asserted that $\inf_m\mathsf{d}_W^{\mathbb{T}}(g,m)>0$, nor that the floor
$s_{\mathrm{vol}}$ is necessary. Nor is anything asserted that clause (iv) of
Theorem~\ref{3.1:thm-main} lists as not asserted, in particular:
\begin{itemize}
\item that the limit fails to be Gaussian;
\item $O(4)$ invariance;
\item a mass gap;
\item anything for block observables, Wilson loops, or the object denoted $\mathfrak{A}$ in
clause (iv) of Theorem~\ref{3.1:thm-main};
\item anything for one source, for overlapping or signed sources, for separations larger than
$6\mathsf{d}$, or as $\vartheta\to0$;
\item that $\zeta'_{K-s}\to1$, beyond the bound on it in \eqref{8.18:eq-torus-witness-5};
\item the case $N\ge3$.
\end{itemize}
\end{theorem}

\begin{proof}
Fix $g$, $g_0$, $K=K(g_0,g)$, $m$, $\mathsf{d}$, $s=s(\mathsf{d})$ and $(f_1,f_2)$ as in the
statement, so that \eqref{8.18:eq-torus-witness-1}, \eqref{8.18:eq-torus-witness-a},
\eqref{8.18:eq-torus-witness-2} and \eqref{8.18:eq-torus-witness-3} hold.

The inequality $m\le m_V(g)$ is not a hypothesis of the theorem. It describes the range of torus
exponents on which the floor $s_{\mathrm{vol}}(g,m)$ of
Definition~\ref{8.18:def-volume-horizon} vanishes. Every $m\ge m_{\mathbb{T}}(g_-(g))$ falls into
exactly one of the two cases $m\le m_V(g)$ and $m>m_V(g)$, which we treat in turn.

\emph{Part (a) for $m\le m_V(g)$.}
By \eqref{8.18:eq-horizon-cases-b} of Lemma~\ref{8.18:lem-horizon-cases}, the run on
$\mathbb{T}_m$ then ends before the coupling reaches $\gamma_*$, that is, $K+m\le K_*-1$. On
such tori $s_{\mathrm{vol}}(g,m)=0$ by \eqref{8.18:eq-volume-horizon-d}, so
\eqref{8.18:eq-torus-witness-a} reduces to $s_W(g)\le s(\mathsf{d})$.
Proposition~\ref{8.18:prop-witness-below-horizon} applies on such tori under
\eqref{8.18:eq-torus-witness-1}, with $\gamma_W^{\mathbb{T}}$ as in
\eqref{8.18:eq-witness-below-horizon-a}, under $s_W(g)\le s(\mathsf{d})$,
\eqref{8.18:eq-torus-witness-2}, \eqref{8.18:eq-torus-witness-3}, $K+m\le K_*-1$ and
$(f_1,f_2)\in\mathfrak{T}(\mathsf{d};\vartheta)$. It gives
the expansion \eqref{8.18:eq-torus-witness-4} with $|R_K(f_1,f_2)|\le\tfrac14$,
$\zeta'_{K-s}\in[\tfrac23,\tfrac43]$, and the kernel bounds \eqref{8.18:eq-torus-witness-6}.

\emph{Part (a) for $m>m_V(g)$.}
Proposition~\ref{8.18:prop-witness-beyond-horizon} gives the same conclusions, under the depth
condition \eqref{8.18:eq-torus-witness-a}, which now contains $s_{\mathrm{vol}}(g,m)\le s(\mathsf{d})$,
and under \eqref{8.18:eq-torus-witness-1}, \eqref{8.18:eq-torus-witness-2}. In both cases
$|R_K|\le\tfrac14\le\tfrac12$, which is \eqref{8.18:eq-torus-witness-5}. This proves part~(a).

\emph{Part (b).}
By \eqref{8.18:eq-torus-witness-5},
\begin{equation*}
\tfrac49\le(\zeta'_{K-s})^2\le\tfrac{16}9,\qquad
\tfrac12\le1+R_K(f_1,f_2)\le\tfrac32 .
\end{equation*}
Both factors are positive, so
\begin{equation*}
\tfrac29=\tfrac49\cdot\tfrac12\le(\zeta'_{K-s})^2\bigl(1+R_K(f_1,f_2)\bigr)
\le\tfrac{16}9\cdot\tfrac32=\tfrac83 .
\end{equation*}
The factor $b_0=(N^2-1)/2$ is positive and $g_{K-s}^4$ is positive. By the lower bound in
\eqref{8.18:eq-torus-witness-6}, with the constant $c_{\mathcal{K}}$ of clause (iv) positive,
$\mathcal{K}_s(f_1,f_2)\ge0$. Multiplying the last display by
$b_0\,g_{K-s}^4\,\mathcal{K}_s(f_1,f_2)\ge0$ and using \eqref{8.18:eq-torus-witness-4} gives
\eqref{8.18:eq-torus-witness-7}.

Clause (0) of Theorem~0, read at level $K-s$, gives
\begin{equation*}
\bigl(g^{-2}+B^{\sharp}(s+1)\bigr)^{-2}\le g_{K-s}^4\le\bigl(g^{-2}+\lambda^{\sharp}s-c_5\bigr)^{-2}.
\end{equation*}
Insert these bounds for $g_{K-s}^4$ in \eqref{8.18:eq-torus-witness-7}. Then insert the two sides of
\eqref{8.18:eq-torus-witness-6} for $\mathcal{K}_s(f_1,f_2)$, lower with lower and upper with upper;
all factors are non-negative. This gives \eqref{8.18:eq-torus-witness-8} and proves part~(b).

\emph{Part (c).}
The constants $\gamma_W^{\mathbb{T}}$, $C_W$, $k_W$, $c_{\mathcal{K}}$, $C_{\mathcal{K}}$ depend on
none of $K$, $g_0$, the position of the supports, or $m$. In both cases above, the bounds
\eqref{8.18:eq-torus-witness-5}--\eqref{8.18:eq-torus-witness-8} involve only these constants
together with $b_0$, $B^{\sharp}$, $\lambda^{\sharp}$, $c_5$, $g$ and $s$. They therefore hold
uniformly in $K$, $g_0$ and position on every torus with $m\ge m_{\mathbb{T}}(g_-(g))$.

The torus exponent enters the hypotheses only through \eqref{8.18:eq-torus-witness-3} and through
$s_{\mathrm{vol}}(g,m)$ in \eqref{8.18:eq-torus-witness-a}. On
$\{m_{\mathbb{T}}(g_-(g))\le m\le m_V(g)\}$ the floor $s_{\mathrm{vol}}(g,m)$ vanishes, so nothing in
the conditions or in the conclusions depends on $m$ there. For $m>m_V(g)$ the only trace of the
volume is $s_{\mathrm{vol}}(g,m)$ in \eqref{8.18:eq-torus-witness-a}. This proves part~(c).

\emph{Part (d).}
Fix a torus $\mathbb{T}_m$ of the family, $m\ge m_{\mathbb{T}}(g_-(g))$, and $g$ satisfying
\eqref{8.18:eq-torus-witness-1}. Each of the four consequences is a statement at this fixed $m$,
either at one small separation $\mathsf{d}$ or as $\mathsf{d}\downarrow0$.

At fixed $m$, condition \eqref{8.18:eq-torus-witness-a} restricts $\mathsf{d}$ only through the
ceiling $\mathsf{d}\le\mathsf{d}_W^{\mathbb{T}}(g,m)$. So the $m$-dependence of this ceiling does not
restrict these statements. For such $\mathsf{d}$, and $K$ satisfying
\eqref{8.18:eq-torus-witness-2}, the bound \eqref{8.18:eq-torus-witness-8}, whose two sides do not
depend on $K$, $g_0$ or the position of the supports, holds by parts~(b) and~(c). It
therefore passes to every torus limit of clause (ii-$\mathbb{T}$-a), (ii-$\mathbb{T}$-b) of Theorem~0.

The four consequences are obtained from this limiting bound, clause (0) of Theorem~0,
clause (ii-$\mathbb{T}$) of Theorem~0, reflection positivity of the lattice measure and the
continuum kernel by the argument which derives the corresponding consequences listed in clause (iv)
of Theorem~\ref{3.1:thm-main}, applied at the fixed $m$ to the torus limits of clause
(ii-$\mathbb{T}$-a), (ii-$\mathbb{T}$-b) in place of the limits on $\mathbb{R}^4$.
\end{proof}

\begin{remark}[The witness on every torus and the volume-uniform witness bound]\label{8.18:rem-witness-relation}
We compare the two-point witness on every torus, Theorem~\ref{8.18:thm-torus-witness}, with the
following statement, which we call the \emph{volume-uniform witness bound}: for $g\le\gamma_W$, the
conclusion of Theorem~\ref{8.18:thm-torus-witness}, with the same kernel $\mathcal{K}_s$, holds on
every torus of the family when \eqref{8.18:eq-torus-witness-a} is replaced by the depth condition
\begin{equation}\label{8.18:eq-witness-relation-1}
s_W(g)\le s(\mathsf{d}) ,
\end{equation}
the ultraviolet condition and the torus clause of (H5) being unchanged. Let $m_V(g)$ be the volume
horizon of \eqref{8.18:eq-volume-horizon-a} and $s_{\mathrm{vol}}(g,m)$ the volume-reading depth
floor of \eqref{8.18:eq-volume-horizon-d}. We say that the floor can be removed beyond the volume
horizon if, for every $m>m_V(g)$, Theorem~\ref{8.18:thm-torus-witness} remains true when
\eqref{8.18:eq-torus-witness-a} is replaced by \eqref{8.18:eq-witness-relation-1}. Then the
volume-uniform witness bound implies Theorem~\ref{8.18:thm-torus-witness}. Moreover, for couplings
$g\le\gamma_W^{\mathbb{T}}$, the volume-uniform witness bound is equivalent to
Theorem~\ref{8.18:thm-torus-witness} together with the removal of the floor beyond the volume
horizon.

The removal of the floor beyond the volume horizon would follow from the old-age tail statement,
which is not proved, in its form on the range $K+m\ge K_*$ with constants uniform in $m$; here
$K_*$ is the level of the reference member. This implication is neither used nor asserted here.

Theorem~\ref{8.18:thm-torus-witness} by itself therefore gives no uniformity in $m$ beyond $m_V(g)$
under \eqref{8.18:eq-witness-relation-1}. Clause~(iv) of Theorem~\ref{3.1:thm-main} does not rest on
the removal of the floor. It is obtained from the two-point witness on a single torus below the
volume horizon by the volume transfer of Corollary~\ref{8.22:cor-transfer-depth-floor}, with the depth
floors stated in that clause and with the kernel read on the reference torus, and its consequences
for the limits on $\mathbb{R}^4$ are drawn after Corollary~\ref{8.22:cor-reference-torus-kernel}.
\end{remark}

\begin{proof}
Both statements conclude the same two-point bound on $\mathbb{T}_m$, for the same test class
$\mathfrak{T}(\mathsf{d};\vartheta)$, the same kernel $\mathcal{K}_s$ and the same depths; they
differ only in their hypotheses.

First, suppose the hypotheses of Theorem~\ref{8.18:thm-torus-witness} hold. The coupling
threshold satisfies $\gamma_W^{\mathbb{T}}\le\gamma_W$, so $g\le\gamma_W^{\mathbb{T}}$ implies
$g\le\gamma_W$. Condition \eqref{8.18:eq-torus-witness-a} bounds the maximum of $s_W(g)$ and
$s_{\mathrm{vol}}(g,m)$ by $s(\mathsf{d})$, and so it implies \eqref{8.18:eq-witness-relation-1},
which is the depth condition of the volume-uniform witness bound. The ultraviolet condition
$K-s(\mathsf{d})\ge k_W$ and the torus clause of (H5) are common to both statements. The
volume-uniform witness bound therefore applies and gives the conclusion of
Theorem~\ref{8.18:thm-torus-witness}. For $m>m_V(g)$ this argument uses of the depth hypothesis only
\eqref{8.18:eq-witness-relation-1}, so the volume-uniform witness bound also gives the removal of
the floor beyond the volume horizon.

Second, let $g\le\gamma_W^{\mathbb{T}}$ and assume Theorem~\ref{8.18:thm-torus-witness} and the
removal of the floor beyond the volume horizon. For $m\le m_V(g)$ the floor
$s_{\mathrm{vol}}(g,m)$ vanishes by \eqref{8.18:eq-volume-horizon-d}; on these tori, condition
\eqref{8.18:eq-torus-witness-a} is therefore \eqref{8.18:eq-witness-relation-1}, and
Theorem~\ref{8.18:thm-torus-witness} gives the bound. For $m>m_V(g)$ the removal of the floor
gives the bound under \eqref{8.18:eq-witness-relation-1}. Hence the witness bound holds on every
torus under \eqref{8.18:eq-witness-relation-1}, uniformly in $m$, which is the volume-uniform
witness bound for these couplings.
\end{proof}

\begin{lemma}[The source lemma (live marked slots and the annealed shadow of the smeared source); Proof outstanding]\label{8.21:lem-source}
Let $k_0=K-s$ be the level at the witness depth $s$. Let $f_1,f_2$ be the two test functions of the witness, with source radius $\rho$, shape constant $\vartheta$ and supports $S_i=\operatorname{supp}f_i$. Put $w=z_1f_1+z_2f_2$.

For a history $h$ of Ba{\l}aban's expansion at level $k_0$ and a block field $V$, let $\tau^w_h$ be the weight of $h$ when the source $w$ is inserted, so that $\tau^0_h=\tau_h$. The effect of the smeared source $\mathcal{O}(f_i)$ on the block law at the witness depth is
\begin{equation}\label{8.21:eq-source-1}
\Phi_i^h=\partial_{z_i}\log\tau^w_h\big|_{z=0}.
\end{equation}
This is the conditional mean of $\mathcal{O}(f_i)$ given $(h,V)$.

As in the rest of this chapter, a large-field component $Z$ is live at level $k_0$ if it still carries its large-field factor there, and its age at level $k_0$ is $k_0-j(Z)$, where $j(Z)$ is the earliest step at which a point of $Z$ entered the large-field region (for a component formed by merging, the oldest creation step among its parts). Let $S_i^+$ be the set of cells of level $k_0$ within a fixed multiple of $M$ cells of $S_i$. A live component $Z$ of $h$ is \emph{near} $S_i$ if the cells of level $k_0$ on which its operator and the live terms attached to it read field or source lie at distance less than $\rho/8$, in cells, from $S_i^+$. Then, for every history $h$, on the support of $\tau_h$, the following hold.
\begin{enumerate}
\item[(a)] There is a decomposition
\begin{equation}\label{8.21:eq-source-2}
\Phi_i^h=c^0_i+G_i(V)+\mathfrak{L}_i^h(V).
\end{equation}
Here $c^0_i$ depends neither on $h$ nor on $V$. The term $G_i$ is history-free, that is, a function of $V$ alone, and it satisfies $|G_i|\le C_G\,\mathfrak{n}(f_i)\,\varepsilon^2$, where $\varepsilon$ is the regularity parameter of the spaces of regular configurations in Ba{\l}aban's small-field region.
\item[(b)] The large-field part splits as $\mathfrak{L}_i^h=\mathfrak{L}_i^{\mathrm{near},h}+\mathfrak{F}_i^h$. The second term satisfies
\begin{equation}\label{8.21:eq-source-3}
\sup_V|\mathfrak{F}_i^h|\le\mathfrak{f}_i\qquad\text{for every history }h.
\end{equation}
Moreover $\mathfrak{L}_i^{\mathrm{near},h}\equiv0$ unless $h$ has a live component near $S_i$.
\item[(c)] Suppose every live component of $h$ near $S_i$ has age at most $\mathfrak{l}$. Then
\begin{equation}\label{8.21:eq-source-4}
\sup_V|\mathfrak{L}_i^{\mathrm{near},h}|\le\mathfrak{s}^\sharp_i .
\end{equation}
\item[(d)] Let $\mathfrak{O}_i$ be the complementary event: $h$ has a live component near $S_i$ of age greater than $\mathfrak{l}$. Let $\mathfrak{a}_i(h)$ be the largest age at level $k_0$ of a live component near $S_i$. For $h\in\mathfrak{O}_i$,
\begin{equation}\label{8.21:eq-source-5}
|\mathfrak{L}_i^h(V)|\le C_{\mathrm{live}}\,L^{4(\mathfrak{a}_i(h)+1)}\,
\mathfrak{p}_{\mathrm{src}}\,\rho^4\,\mathfrak{n}(f_i)/\vartheta .
\end{equation}
\end{enumerate}
The bounds in \textup{(b)} and \textup{(c)} are
\begin{align*}
\mathfrak{s}^\sharp_i&=C_{\mathrm{live}}\,\bar\Theta\,\mathfrak{p}_{\mathrm{src}}\,\rho^4\,
\mathfrak{n}(f_i)/\vartheta,\\
\mathfrak{f}_i&=C_{\mathrm{src}}C_{\mathrm{far}}\,\|f_i\|_\infty\,\rho^8\,e^{-\varkappa\rho/16},
\end{align*}
where $\bar\Theta=L^{4(\mathfrak{l}(\bar x_s)+1)}$, with $\mathfrak{l}(\cdot)$ the age threshold as a function of the inverse coupling evaluated at $\bar x_s$, so that $\bar\Theta\ge L^{4(\mathfrak{l}+1)}$; $\varkappa$ is the decay rate of the two-point witness theorem; and $\mathfrak{p}_{\mathrm{src}}=\mathfrak{p}_{\mathrm{src}}(\bar x_s)$ is an explicit power of $\log\bar x_s$. The constants $C_G$, $C_{\mathrm{live}}$, $C_{\mathrm{src}}$, $C_{\mathrm{far}}$ depend only on $L$, $M$, $\vartheta$ and Ba{\l}aban's auxiliary parameters and constants.
\end{lemma}

Proof outstanding.

\noindent\emph{Route.}
A proof is to read the decomposition \eqref{8.21:eq-source-2} off Ba{\l}aban's representation of the density at level $k_0$ with the bounded local source $w$ inserted. In that representation the weight $\tau^w_h$ should be the product of the characteristic functions of the history with an integral, against a positive measure independent of $z$ on the variables of all live components jointly, in the form stated in \cite[(2.19)--(2.22)]{B-CMP119}, of an exponential into whose exponent all dependence on $w$ enters, analytic in $z$ on a polydisc around $0$ whose radii are inversely proportional to $\|f_i\|_\infty$. The live large-field operators then act as positive functionals: normalising the integral gives a probability measure, and by positivity alone the mean of each $z_i$-derivative of the exponent is bounded by its supremum on the support of that measure. On that support the configuration inside a live component is to be regular at the scale of the step that created it \cite[(2.17)--(2.22)]{B-CMP119}, which should bound the energy per cell of level $k_0$ by a constant times $L^{4(\mathfrak{a}+1)}$ for a component of age $\mathfrak{a}$; together with the irrelevance bounds \cite[(0.28)--(0.30)]{B-CMP109} and Cauchy estimates on the polydisc this is to give \eqref{8.21:eq-source-5}. The young near components are to be paid by the source radius and the age weight: the age weight of every near young component is at most $L^{4(\mathfrak{l}+1)}\le\bar\Theta$, and the cells near $S_i$ number at most a constant times $\rho^4$, whence the shape of $\mathfrak{s}^\sharp_i$ in \eqref{8.21:eq-source-4}. The far components and the deep labels are to be summed into $\mathfrak{F}_i^h$: their contribution is counted by anchor cells in $S_i^+$ and bounded by the exponential decay of the localised terms in the distance to $S_i$, and the resulting bound $\mathfrak{f}_i$ in \eqref{8.21:eq-source-3} is to hold whatever the number of far components. The history-free part $G_i$ and the constant $c^0_i$ come from the exterior small-field terms.

The lemma makes no statement about the probability of the event $\mathfrak{O}_i$. It also gives no bound on $\Phi_i^h-c^0_i$ that is uniform over all pairs $(h,V)$.

\begin{remark}\label{8.21:lem-continuum-kernel}
The statement and proof of this lemma are not written out here.
\end{remark}

\begin{lemma}[The flat shell bound (input of the alternative bookkeeping only). Proof outstanding]
\label{8.21:lem-flat-shell}
Let $k_0$, the depth $\mathfrak{l}$, the histories $h$ with their labels, the measure
$\tilde\nu$ on histories, the factors $\tilde q_j$, the parameters $\check R$ and
$\check R_j$, the constants $\kappa_c$ and $\varrho$, and the regions $\mathcal{N}_i$ and
$\mathcal{N}_i^{++}$, $i=1,2$, be those of the two-point witness theorem, in the variant of
its bookkeeping in which healed labels are kept in the history index.
There are constants $C_{\mathrm{sh}}$ and $c_{\mathrm{sh}}$ such that the bounds below hold
with the following definitions.
Let $\mathcal{C}$ be the $\sigma$-algebra generated by $V_{k_0}$ and the labels of $h$ of
level at least $k_0-\mathfrak{l}$.
Let $\mathcal{D}_i$ be the set of healed labels of $h$ of level less than
$k_0-\mathfrak{l}$ lying under $\mathcal{N}_i$; they are called \emph{deep}.
For $Y\in\mathcal{D}_i$ of level $j=k_0-\nu$ put
\begin{equation*}
\mu^{\sharp}_Y=C_{\mathrm{sh}}\,(M\check R_j)^{c_{\mathrm{sh}}}\,x_j\,
(\log x_j)^{c_{\mathrm{sh}}}\,L^{-4\nu},
\qquad
\chi_Y=\exp\Bigl(-\frac{\kappa_c\,\varrho\,L^{\nu}}{400\,M\check R_j}\Bigr).
\end{equation*}
Let $B^{\times}$ be the event that some label of $h$ of level at least $k_0-\mathfrak{l}$
meets both $\mathcal{N}_1^{++}$ and $\mathcal{N}_2^{++}$.
For a random variable $X$ put
\begin{equation*}
\mathbb{E}^{\natural}|X|=\sup_{A}\mathbb{E}_{\tilde\nu}\bigl[\,|X|\bigm| A\,\bigr],
\end{equation*}
where $A$ ranges over the $\sigma$-algebra generated by the labels meeting
$\mathcal{N}_1^{++}$.
The statement concerns bounded local functionals of the configurations of deep labels under
$\mathcal{N}_1^{++}$ and $\mathcal{N}_2^{++}$; its atoms are the following:
\begin{itemize}
\item[(a)] for $Y\in\mathcal{D}_1$ of level $j$, and for every bounded $X$ that is
measurable with respect to the $\sigma$-algebra generated by the whole field $V_{k_0}$
(not only its restriction to $\mathcal{N}_2^{++}$) and the labels meeting
$\mathcal{N}_2^{++}$,
\begin{equation}\label{8.21:eq-flat-shell-1}
\bigl|\operatorname{Cov}_{\tilde\nu}(\mathbf{1}_{\{Y\in h\}},X)\bigr|
\le \tilde q_j\,\bigl(\mu^{\sharp}_Y+\chi_Y\bigr)\,\mathbb{E}^{\natural}|X| ;
\end{equation}
\item[(b)] for $Y\in\mathcal{D}_1$ of level $j$ and $Y'\in\mathcal{D}_2$ of level $j'$,
\begin{equation}\label{8.21:eq-flat-shell-2}
\bigl|\operatorname{Cov}_{\tilde\nu}(\mathbf{1}_{\{Y\in h\}},\mathbf{1}_{\{Y'\in h\}})\bigr|
\le \tilde q_j\,\tilde q_{j'}\,
\bigl(\mu^{\sharp}_Y\mu^{\sharp}_{Y'}+\min(\chi_Y,\chi_{Y'})\bigr),
\end{equation}
the covariance being taken under $\tilde\nu$ on the whole space of histories, the event
$B^{\times}$ not excluded.
\end{itemize}
\end{lemma}

The argument of this chapter does not use Lemma~\ref{8.21:lem-flat-shell}, since it sums
healed labels out of the history index, so that $\mathcal{D}_1=\mathcal{D}_2=\emptyset$
there; the lemma is the input of the alternative bookkeeping that keeps healed labels in
the index, where \eqref{8.21:eq-flat-shell-2} is used once, and it is recorded for
completeness.

Proof outstanding.

\emph{Route.} A proof would have to establish two properties of $\tilde\nu$.
First, conditionally on $\mathcal{C}$, the configurations $h\cap\mathcal{D}_1$ and
$h\cap\mathcal{D}_2$ are independent.
This would follow from the factorisation over connected components in the decomposition of
the correlation theorem, together with the fact that every deep label has diameter less
than $\varrho/8$.
Second, for every deep label $Y$ there is a $\sigma$-algebra
$\mathcal{C}^{\mathrm{loc}}_Y\subset\mathcal{C}$, generated by the data of $\mathcal{C}$
near $Y$, such that
\begin{equation*}
\bigl|\tilde\nu(Y\in h\mid\mathcal{C})-\tilde\nu(Y\in h\mid\mathcal{C}^{\mathrm{loc}}_Y)\bigr|
\le \tilde\nu(Y\in h\mid\mathcal{C}^{\mathrm{loc}}_Y)\,\mu^{\sharp}_Y .
\end{equation*}
Here the gauge-invariant terms attached to $Y$ are to be shown to be of second order in the
coarse datum, with the localisation radii of \cite[(2.28)]{B-CMP119}, and the chains of
polymers joining $Y$ to distant data are to be shown to carry the factor $\chi_Y$.
In (a) the term $\chi_Y$ arises only through the labels on which $X$ depends, not through
$V_{k_0}$.
The second property would also give the bound
\begin{equation*}
\tilde\nu(B^{\times})\le \varrho^{8}\,\tilde q_{k_0}\,
e^{-\kappa_c\varrho/(400M\check R)} .
\end{equation*}
The inputs are the large-field block bound, healed part, the kernel bound of this chapter,
and Ba{\l}aban's polymer-gas representation of the effective density
\cite[\S1]{B-CMP122b}, read under the annealed measure $\tilde\nu$.
Every sum over labels has to close with the weights $\tilde q_j$ alone.
Bounds \eqref{8.21:eq-flat-shell-1} and \eqref{8.21:eq-flat-shell-2} would then follow by
combining the two properties.

\begin{corollary}[Consequences of the two-point display on every torus: two-sided bounds,
positivity, and the four properties of every finite-volume limit]\label{8.21:cor-consequences-torus}
Assume the following.
\begin{itemize}
\item[(A1)] The two-point witness theorem on a fixed torus, in the following form: for every
$g\le\gamma_W$, every $g_0\in(0,g_-(g)]$ with $K=K(g_0,g)$, every torus $\mathbb{T}_m$ of the
family of Theorem~\ref{3.1:thm-main}, every
$\mathsf{d}$ satisfying the conditions of that statement (the depth floor and the ultraviolet
condition relative to the spacing $L^{-K}$), with $s=s(\mathsf{d})$, and every
$(f_1,f_2)\in\mathfrak{T}(\mathsf{d};\vartheta)$,
\begin{equation}\label{8.21:eq-consequences-torus-1}
S^T_{2;K}(f_1,f_2)=\zeta'^{\,2}_{K-s}\,b_0\,g^4_{K-s}\,\mathcal{K}_{K,s}(f_1,f_2)\,
\bigl(1+R_K(f_1,f_2)\bigr),
\end{equation}
with $|R_K(f_1,f_2)|\le\frac12$ (indeed $|R_K(f_1,f_2)|\le\frac14$),
$\zeta'_{K-s}\in[\frac23,\frac43]$, and, with constants $0<c_{\mathcal{K}}\le C_{\mathcal{K}}$,
the kernel bounds
\begin{equation*}
c_{\mathcal{K}}\|f_1\|_\infty\|f_2\|_\infty\le\mathcal{K}_{K,s}(f_1,f_2)
\le C_{\mathcal{K}}\|f_1\|_\infty\|f_2\|_\infty .
\end{equation*}
\item[(A2)] Clause (0) of Theorem~\ref{3.1:thm-main}, in the form: for $s$ as in (A1),
\begin{equation}\label{8.21:eq-consequences-torus-2}
g^{-2}+\lambda_\sharp s-c_5\le g^{-2}_{K-s}\le g^{-2}+B_\sharp(s+1),
\end{equation}
the left member of \eqref{8.21:eq-consequences-torus-2} being positive.
\item[(A3)] Clause (ii-$\mathbb{T}$) of Theorem~\ref{3.1:thm-main}: on every torus of the
family the limits $S^T_{n;m}$, $n\ge2$, of (ii-$\mathbb{T}$-b) exist along sequences
$K_i\to\infty$ of first-passage cutoffs, and, in the regime of clause (v), the connected
two-point function of every member $S^{(\hat w)}$ of the tuned family (ii-$\mathbb{T}$-a) is a
limit of two-point functions $S^T_{2;K_i}$, $K_i\to\infty$, along data to which (A1) and (A2)
apply, at fixed $g$, $m$ and $\mathsf{d}$.
\item[(A4)] Reflection positivity of the lattice measure (Chapter~\ref{ch:6}).
\item[(A5)] The continuum kernel (Lemma~\ref{8.21:lem-continuum-kernel}, whose proof is
outstanding), in the form: along the sequences of (A3), every limit point of the kernel factors
$\mathcal{K}_{K_i,s}(f_1,f_2)$ of \eqref{8.21:eq-consequences-torus-1} lies in the interval
\begin{equation*}
\bigl[\,c_{\mathcal{K}}\vartheta^2\|f_1\|_\infty\|f_2\|_\infty,\;
100\,C_{\mathcal{K}}\|f_1\|_\infty\|f_2\|_\infty\,\bigr].
\end{equation*}
\end{itemize}
Then the following hold.
\begin{itemize}
\item[(a)] For all data as in (A1), uniformly in $K$, $g_0$, $m$ and the position of the
supports,
\begin{equation}\label{8.21:eq-consequences-torus-3}
\tfrac29\,b_0\,g^4_{K-s}\,\mathcal{K}_{K,s}(f_1,f_2)\le S^T_{2;K}(f_1,f_2)
\le\tfrac83\,b_0\,g^4_{K-s}\,\mathcal{K}_{K,s}(f_1,f_2),
\end{equation}
$S^T_{2;K}(f_1,f_2)>0$, and
\begin{equation}\label{8.21:eq-consequences-torus-4}
\begin{split}
\tfrac29\,b_0\,c_{\mathcal{K}}\|f_1\|_\infty\|f_2\|_\infty\bigl(g^{-2}+B_\sharp(s+1)\bigr)^{-2}
&\le S^T_{2;K}(f_1,f_2)\\
&\le\tfrac83\,b_0\,C_{\mathcal{K}}\|f_1\|_\infty\|f_2\|_\infty
\bigl(g^{-2}+\lambda_\sharp s-c_5\bigr)^{-2}.
\end{split}
\end{equation}
\item[(b)] For every $\mathsf{d}$ and every $(f_1,f_2)\in\mathfrak{T}(\mathsf{d};\vartheta)$ as in
(A1), with $s=s(\mathsf{d})$, and for $S$ a limit $S^T_{2;m}$ of (ii-$\mathbb{T}$-b), or the
connected two-point function of a member $S^{(\hat w)}$ of (ii-$\mathbb{T}$-a),
\begin{equation}\label{8.21:eq-consequences-torus-5}
\begin{split}
\tfrac29\,b_0\,c_{\mathcal{K}}\vartheta^2\|f_1\|_\infty\|f_2\|_\infty
\bigl(g^{-2}+B_\sharp(s+1)\bigr)^{-2}&\le S(f_1,f_2)\\
&\le\tfrac{800}{3}\,b_0\,C_{\mathcal{K}}\|f_1\|_\infty\|f_2\|_\infty
\bigl(g^{-2}+\lambda_\sharp s-c_5\bigr)^{-2};
\end{split}
\end{equation}
in particular $S(f_1,f_2)>0$, and every limit point $g_{\infty,s}$ of $g_{K_i-s}$ along the
approximating sequence satisfies
\begin{equation*}
g^{-2}+\lambda_\sharp s-c_5\le g^{-2}_{\infty,s}\le g^{-2}+B_\sharp(s+1).
\end{equation*}
\item[(c)] On every torus of the family, for every such $S$: (c1) the smeared curvature has
strictly positive variance in the reflection-positive form of the torus, so that
the curvature Hilbert space $\mathcal{H}^{\mathrm{curv}}_{\mathbb{T}_m}$ of the torus
(Notation~\ref{0.2:not-glossary}) satisfies $\dim\mathcal{H}^{\mathrm{curv}}_{\mathbb{T}_m}\ge2$;
(c2) $S$ is not
scale invariant under any
assignment of a scaling dimension to the curvature; (c3) $S$ is reproduced by no free Maxwell
field on $\mathbb{T}_m$ at any charge; (c4) the coupling read off $S$ at separation
$\mathsf{d}$ is strictly positive and tends to $0$ logarithmically as $\mathsf{d}\downarrow0$,
at the two-sided rate of clause (0) (asymptotic freedom in a gauge-invariant correlation).
\end{itemize}
Nothing is asserted here for the limits on $\mathbb{R}^4$; for these, including
$\dim\mathcal{H}^{\mathrm{curv}}\ge2$ on $\mathbb{R}^4$, see clause~(iv) of
Theorem~\ref{3.1:thm-main} and Corollary~\ref{8.22:cor-reference-torus-kernel}. Also not asserted: that the limit fails to be
a Gaussian field, $O(4)$ invariance, a mass gap; anything on block observables, on Wilson loops,
or on the object denoted $\mathfrak{A}$ in the list of non-assertions of
Theorem~\ref{3.1:thm-main};
anything for a single source, for overlapping or signed sources, for separations larger than
$6\mathsf{d}$, as $\vartheta\to0$, or at depths outside the range of (A1); that
$\zeta'_{K-s}\to1$ beyond the bound on $|\zeta'_{K-s}-1|$ stated in the two-point witness
theorem; anything for $N\ge3$; that the $(g,m)$-dependent ceiling on
$\mathsf{d}$ is bounded below by a positive number uniformly in $m$; that a volume-dependent
term in the depth floor of (A1) is needed (by clause (iv) of Theorem~\ref{3.1:thm-main} it can
be taken to be zero); and the sharp one-loop asymptotics. The
coefficient $b_0$ enters \eqref{8.21:eq-consequences-torus-3} only within the fixed factor
$(8/3)/(2/9)=12$.
\end{corollary}

\begin{proof}
(a) In \eqref{8.21:eq-consequences-torus-1} one has $\zeta'^{\,2}_{K-s}\in[\frac49,\frac{16}9]$,
$1+R_K\in[\frac12,\frac32]$, and $b_0g^4_{K-s}\mathcal{K}_{K,s}\ge0$ by the lower kernel bound.
Since $\frac49\cdot\frac12=\frac29$ and $\frac{16}9\cdot\frac32=\frac83$, this gives
\eqref{8.21:eq-consequences-torus-3}; with $|R_K|\le\frac14$ the same computation would give
$\frac13$ and $\frac{20}9$. The constants in (A1) do not depend on $K$, $g_0$, $m$ or position.
The lower kernel bound and $c_{\mathcal{K}}>0$, the members of $\mathfrak{T}(\mathsf{d};\vartheta)$
not vanishing identically, give $S^T_{2;K}>0$. Raising the three members of
\eqref{8.21:eq-consequences-torus-2} to the power $-2$ (all are positive) bounds $g^4_{K-s}$
between $(g^{-2}+B_\sharp(s+1))^{-2}$ and $(g^{-2}+\lambda_\sharp s-c_5)^{-2}$; inserting this and
the kernel bounds into \eqref{8.21:eq-consequences-torus-3} gives
\eqref{8.21:eq-consequences-torus-4}.

(b) Along a sequence $K_i\to\infty$ of (A3), \eqref{8.21:eq-consequences-torus-3} and
\eqref{8.21:eq-consequences-torus-2} hold, with $s$ fixed, for all $i$ large enough that
$\mathsf{d}$ meets the ultraviolet condition of (A1) at cutoff $K_i$ (for (ii-$\mathbb{T}$-a)
this is part of (A3)). The numbers
$g^{-2}_{K_i-s}$ lie in the fixed interval of \eqref{8.21:eq-consequences-torus-2}, whose left
end is positive. By \eqref{8.21:eq-consequences-torus-3} the kernel factors
$\mathcal{K}_{K_i,s}(f_1,f_2)$ lie between
\begin{equation*}
\tfrac38\,b_0^{-1}g^{-4}_{K_i-s}S^T_{2;K_i}(f_1,f_2)\quad\text{and}\quad
\tfrac92\,b_0^{-1}g^{-4}_{K_i-s}S^T_{2;K_i}(f_1,f_2),
\end{equation*}
so they form a bounded sequence, since $S^T_{2;K_i}(f_1,f_2)$ converges by (A3) and
$g^{-2}_{K_i-s}$ lies in a fixed interval.
Pass to a subsequence along which $S^T_{2;K_i}(f_1,f_2)\to S(f_1,f_2)$ (by (A3)), $g_{K_i-s}$
converges to a limit point $g_{\infty,s}$, and the kernel factors converge to a limit point.
Every limit point $g_{\infty,s}$ obeys \eqref{8.21:eq-consequences-torus-2} with
$g_{\infty,s}$ in place of $g_{K-s}$, the interval being closed. By the continuum kernel
(Lemma~\ref{8.21:lem-continuum-kernel}), whose proof is outstanding, the limit point of the
kernel factors lies between $c_{\mathcal{K}}\vartheta^2\|f_1\|_\infty\|f_2\|_\infty$ and
$100C_{\mathcal{K}}\|f_1\|_\infty\|f_2\|_\infty$. Letting $i\to\infty$ in
\eqref{8.21:eq-consequences-torus-3} bounds $S(f_1,f_2)$ between $\frac29b_0g^4_{\infty,s}$ and
$\frac83b_0g^4_{\infty,s}$ times this limit point; raising the bounds on $g^{-2}_{\infty,s}$ to
the power $-2$ and inserting the bounds on the limit point gives
\eqref{8.21:eq-consequences-torus-5}, whose left member is positive, so $S(f_1,f_2)>0$.

(c) Each property is a statement as $\mathsf{d}\downarrow0$, or at one small $\mathsf{d}$, at
fixed $m$. Since $s(\mathsf{d})\uparrow\infty$ as $\mathsf{d}\downarrow0$, logarithmically in
$1/\mathsf{d}$ and with no upper cap, $s(\mathsf{d})$ exceeds every finite depth floor, so the
$(g,m)$-dependent ceiling on $\mathsf{d}$ costs nothing. Write
$\mathcal{Q}=S(f_1,f_2)/(\|f_1\|_\infty\|f_2\|_\infty)$. By
\eqref{8.21:eq-consequences-torus-5}, $\mathcal{Q}>0$ and $\mathcal{Q}\to0$ as
$\mathsf{d}\downarrow0$. The lower bound is of order $s(\mathsf{d})^{-2}$, so
$\mathsf{d}^{-\varepsilon}\mathcal{Q}\to\infty$ for every $\varepsilon>0$.

(c1) Let $\iota$ be the reflection of $\mathbb{T}_m$ in a hyperplane, and
$(f,\iota f)\in\mathfrak{T}(\mathsf{d};\vartheta)$ with $f$ real-valued and supported on one
side; such pairs exist, taking $f$ supported in a ball of radius $\mathsf{d}$ at a distance
from the hyperplane for which the supports of $f$ and $\iota f$ are separated as
$\mathfrak{T}(\mathsf{d};\vartheta)$ requires. Since
$\operatorname{Re}\operatorname{tr}$ is unchanged by reversing the orientation of $\partial p$,
$\iota\mathcal{O}(f)=\mathcal{O}(\iota f)$. By (A4) the form
$\langle A,B\rangle=\langle\iota(\bar A)B\rangle$, where $\langle\,\cdot\,\rangle$ is
expectation in the lattice measure at cutoff $K$ and $\bar A$ is the complex conjugate of $A$,
is positive semidefinite at every $K$. Put
$A=\mathcal{O}(f)-\langle\mathcal{O}(f)\rangle\,1$ and let $\Omega$ be the class of $1$. Then
$\langle\Omega,\Omega\rangle=\langle1\rangle=1$ and
$\langle\Omega,A\rangle=\langle A\rangle=0$. Since $\mathcal{O}(f)$ is real and
$\langle\mathcal{O}(\iota f)\rangle=\langle\mathcal{O}(f)\rangle$ by the invariance of the
lattice measure under $\iota$ (Lemma~\ref{1.1:lem-invariance}),
\begin{equation*}
\langle A,A\rangle=\langle\mathcal{O}(\iota f)\mathcal{O}(f)\rangle
-\langle\mathcal{O}(\iota f)\rangle\langle\mathcal{O}(f)\rangle=S^T_{2;K}(\iota f,f).
\end{equation*}
The reflection-positive form of the torus limit, from which
$\mathcal{H}^{\mathrm{curv}}_{\mathbb{T}_m}$ and $\Omega$ are obtained
(Notation~\ref{0.2:not-glossary}), takes on $\Omega$ and $A$ the limits along $K_i$ of these
three values, which exist by (A3): $\langle\Omega,\Omega\rangle=1$, $\langle\Omega,A\rangle=0$,
and $\langle A,A\rangle=S(\iota f,f)>0$ by (b). Hence the classes of $\Omega$ and $A$ are
non-zero and orthogonal, hence linearly independent, so
$\dim\mathcal{H}^{\mathrm{curv}}_{\mathbb{T}_m}\ge2$.

(c2) Let $(f_1^\sigma,f_2^\sigma)$, $\sigma\in(0,1]$, be dilates of a pair in
$\mathfrak{T}(\mathsf{d};\vartheta)$ lying in $\mathfrak{T}(\sigma\mathsf{d};\vartheta)$; sup
norms are unchanged. Write $\mathcal{Q}(\sigma)$ for the quotient $\mathcal{Q}$ evaluated at the
pair $(f_1^\sigma,f_2^\sigma)$. Scale invariance with some scaling dimension of the curvature
means $\mathcal{Q}(\sigma)=\sigma^a\mathcal{Q}(1)$ for a real $a$, with $\mathcal{Q}(1)>0$.
If $a>0$ this decays like a power, contrary to the lower bound. If $a<0$ it diverges,
contrary to $\mathcal{Q}\to0$. If $a=0$ it is a positive constant, again contrary to
$\mathcal{Q}\to0$.

(c3) For the free Maxwell field on $\mathbb{T}_m$ at any charge, $\mathcal{Q}(\sigma)$ along the
dilates tends to a strictly positive constant as $\sigma\downarrow0$; this is incompatible with
$\mathcal{Q}\to0$.

(c4) Put $g_{\mathrm{eff}}=(\mathcal{Q}/b_0)^{1/4}>0$. Then \eqref{8.21:eq-consequences-torus-5}
gives
\begin{equation*}
\Bigl(\tfrac{3}{800\,C_{\mathcal{K}}}\Bigr)^{1/2}\bigl(g^{-2}+\lambda_\sharp s-c_5\bigr)
\le g_{\mathrm{eff}}^{-2}
\le\Bigl(\tfrac{9}{2\,c_{\mathcal{K}}\vartheta^2}\Bigr)^{1/2}\bigl(g^{-2}+B_\sharp(s+1)\bigr).
\end{equation*}
Thus $g_{\mathrm{eff}}\to0$ as $\mathsf{d}\downarrow0$, and $g_{\mathrm{eff}}^{-2}$ grows
linearly in $s(\mathsf{d})$ with slopes governed by $\lambda_\sharp$ and $B_\sharp$, the two
slopes of clause (0); this is asymptotic freedom in a gauge-invariant correlation. Hence
$\mathcal{Q}$ is
positive and tends to zero more slowly than every power of $\mathsf{d}$, so $\mathcal{Q}$
satisfies the
two conditions that Definition~\ref{3.1:def-non-trivial} places on this quantity (not
identically zero; decay slower than every power of $\mathsf{d}$), its finite-cutoff approximants
having the form \eqref{8.21:eq-consequences-torus-1}; this is a statement on the torus
$\mathbb{T}_m$ only.
\end{proof}

\paragraph{Inputs used by statement.}
The proof of clause~(iv) of Theorem~\ref{3.1:thm-main} is the only proof of a clause given in
this chapter. For that clause the list below names two kinds of input. Inputs of the first kind
are the published statements of Ba{\l}aban that the proof of clause~(iv) uses by statement, with
their exact references. Inputs of the second kind are the statements of this book used in that
proof whose own proofs are incomplete or outstanding. The clauses are listed in the order of
Part~I.

\begin{description}
\item[(0)] Clause~(0) is not among the clauses whose proof this chapter gives; there is nothing
to list.
\item[(i)] As for clause~(0).
\item[(iii)] As for clause~(0).
\item[(ii)] As for clause~(0).
\item[(iv)]
\emph{Published statements cited by statement}
\begin{itemize}
\item The free propagators of the renormalisation transformations and their decomposition step by
step, \cite[(1.3), (1.5), (1.16)--(1.19), (1.23)--(1.24)]{B-CMP95}.
\item Decay of the fluctuation propagators, \cite[Propositions~2.3, 2.6 and~2.7]{B-CMP96}.
\item The averaging operations and their properties,
\cite[(19)--(20)]{B-CMP98}.
\item Decay bounds for the propagators in a background field,
\cite[(3.156); \S\S C--E]{B-CMP99b}.
\item The minimiser of the variational problem and its regularity,
\cite[Theorem~1 with (6)--(10)]{B-CMP102b}.
\item Analyticity of the minimiser in the data,
\cite[Proposition~9 with (190), pp.~308--309; (172); (182)]{B-CMP102b}.
\item The small-field renormalisation step and its fluctuation covariance,
\cite[\S2, (2.1), (2.3)--(2.13), (2.16), pp.~267--268]{B-CMP109}.
\item The expansion of the effective action in the small-field approximation,
\cite[Theorem~1; \S0; (0.12), (0.14)--(0.26), (0.28)--(0.30)]{B-CMP109}.
\item The localised terms of the effective action and their forms on constant flux,
\cite[\S\S1--3; (1.3), (1.7), (1.11)--(1.16), (1.18), (1.19); (3.1), (3.2), (3.9), (3.32);
\S4, (4.4), (4.8), (4.9), (4.14), (4.15), (4.17), (4.18)]{B-CMP109}.
\item The coupling constant renormalisation,
\cite[Theorem~3; (1.20)--(1.22); (3.61); (5.1)--(5.7), (5.16)]{B-CMP109}.
\item The cluster expansion of the effective action and its bounds,
\cite[Lemma~2 with (1.26), (1.36), (1.41)--(1.43); Lemma~3;
(2.3); (2.5)--(2.6); App.~C]{B-CMP116}.
\item The representation of the effective densities and their bounds,
\cite[Theorem~2; Corollary~3; (0.1); \S2;
(2.1)--(2.3), (2.8), (2.10)--(2.37); pp.~247, 254--256]{B-CMP119}.
\item The regularity conditions and the small-field characteristic functions,
\cite[\S1; (1.1), (1.5), (1.8), (1.9), (1.11), (1.12), (1.18), (1.19)]{B-CMP119}.
\item The renormalisation step on the small-field regions,
\cite[\S3; (3.1)--(3.25)]{B-CMP119}.
\item The coupling increment and the running couplings,
\cite[(3.40), (3.50), (3.56)--(3.61), (3.65); (6.1)]{B-CMP119}.
\item The basic step of the large-field operation and the history measure,
\cite[\S0; (0.3)--(0.5); (1.1), (1.3)--(1.4); (1.98); (1.100)--(1.102)]{B-CMP122a}.
\item The large-field regions and their creation,
\cite[\S1; (1.12)--(1.20), (1.22)--(1.24), (1.27), (1.28), (1.55),
(1.75), (1.76), (1.79), (1.82), (1.88); pp.~175, 177]{B-CMP122a}.
\item Localisation and bounds for the large-field operation,
\cite[\S1; (1.2)--(1.10); (1.33)--(1.35); (1.98)--(1.102); pp.~355, 378, 380, 383,
391]{B-CMP122b}.
\item The bounds of the large-field terms over the large-field regions,
\cite[(1.70)--(1.73); (1.79)--(1.89)]{B-CMP122b}.
\item Admissibility of the prepared large-field data, printed without proof in
\cite[pp.~178--179]{B-CMP122a}; cited as printed.
\item The identification of the flat quadratic form of the small-field step with that of the free
field, printed without proof in \cite[\S1]{B-CMP95}; cited as printed.
\item Positivity and bound of the normalised island operation, printed without proof in
\cite[\S1]{B-CMP122b}; cited as printed.
\item Convergence and analytic extension, with their bounds, of the expansions on the joined
localisation domains, \cite[p.~192, (1.70)]{B-CMP122a} and
\cite[p.~377, (1.68)--(1.69)]{B-CMP122b}; cited as printed.
\end{itemize}
\emph{Statements of this book whose proofs are incomplete or outstanding}
\begin{itemize}
\item The small-piece statement --- stated; proof incomplete.
\item Identity of representations for the shifted frozen run --- proof outstanding.
\item Radius tails of the majorant masses --- stated; proof incomplete.
\item The second-derivative room lemma --- stated; proof incomplete.
\item Covariance of the remainder column, small in the running coupling --- stated; proof
incomplete.
\item Monotonicity of the common prefactor under island removal --- stated; proof incomplete.
\item Later steps do not read earlier healed islands --- stated; proof incomplete.
\item Annealing and conditioning of the removal bound --- stated; proof incomplete.
\item Translation covariance of the label histories --- stated; proof incomplete.
\item Independence of the steps from the target coupling --- stated; proof incomplete.
\item The age lemma --- stated; proof incomplete.
\item Rarity of never-healed old regions above a volume-dependent depth floor --- stated; proof
incomplete.
\item The source lemma --- proof outstanding.
\end{itemize}
\item[(v)] As for clause~(0).
\item[(vi)] As for clause~(0).
\end{description}

No statement on which the two-point witness theorem or the volume-transfer bound for connected
functions rests, whatever its standing --- those listed above and all others proved, or proved as
an implication, in this and earlier chapters --- has clause~(iv), or either of those two
theorems, among its inputs; the statements of this chapter that assemble clause~(iv) from those
two theorems are used in the proof of none of the statements on which the two theorems rest. The
published statements listed above are statements of the cited papers and rest on nothing proved
in this book. No cycle of dependency therefore passes through clause~(iv). The constants entering
the proofs of this chapter
are chosen in the order of Definitions~\ref{2.3:not-stages-first}, \ref{2.3:not-stages-middle}
and~\ref{2.3:not-stages-last}, continued for the constants of the witness construction by the orders
displayed in this chapter. By Lemma~\ref{2.4:lem-no-cycle} and those displays the order has no
cycle, and every condition in it is one-sided.

%% ---- end inlined file: chapters/c08 ----
